\documentclass[11pt]{amsart}
\usepackage[margin=1.1in]{geometry}
\usepackage{amsmath,amssymb,mathrsfs}
\numberwithin{equation}{section}
\newtheorem{theorem}{Theorem}[section]
\newtheorem{proposition}[theorem]{Proposition}
\newtheorem{lemma}[theorem]{Lemma}
\newtheorem{corollary}[theorem]{Corollary}
\theoremstyle{definition}
\newtheorem{definition}[theorem]{Definition}
\newtheorem{example}[theorem]{Example}
\theoremstyle{remark}
\newtheorem{remark}[theorem]{Remark}
\makeatletter
\newcommand{\itemmark}[1]{\protected@edef\@currentlabel{#1}}
\makeatother
\begin{document}
\title{Noncompact lattice QED in three dimensions: limit points of the curvature field and their fourth cumulant}
\author{}
\date{}

\begin{abstract}
Two laws of the curvature $F$ of noncompact lattice QED on three-dimensional tori, with four-component Wilson fermions integrated out, average the constant modes of the gauge field over one period of the fermion determinant; the second also sums the flux sectors. For all charges $e>0$ and lattice masses $>-\frac{1}{2}$ the laws of $F$ are tight; every continuum and infinite-volume limit point is carried by closed two-forms, with Laplace transform below the free Maxwell one, reflection positive at every coordinate plane and invariant under translations and the cubic group. There are volume-independent bare mass shifts such that every test form $f=dd^{*}\beta$ with $b(f)>0$, $b(f)/24$ being the quartic term of the free Dirac $\log\det$ in $d^{*}\beta$, has a constant $C>0$ with $\left|\kappa_4(F(f))-e^4\cdot b(f)\right|\le C\cdot e^{4+\sigma}$, $\sigma>0$, for all sufficiently small charges $e>0$, in every limit point and almost every ergodic component. All these laws are non-Gaussian; no such limit point is a mixture of translation-invariant Gaussian laws whose Laplace transform at every real test $g$ is at most $\exp\left(\frac{1}{2}\|g\|_2^2\right)$.
\end{abstract}
\maketitle

\section{Introduction}\label{sec:introduction}

Let $L$ be an odd integer, which is fixed in Section~\ref{sec:parameters-scales}, and let $\bar{m}\in(0,1]$ be a number, the mass parameter. Put $h:=L^{\left\lfloor\log_L\bar{m}^{-1}\right\rfloor}$ and $\Lambda_{\eta,R}=\eta\mathbb{Z}^3/R\mathbb{Z}^3$, $\eta=L^{-N}<R=h\cdot L^T\ \left(N,T\in\mathbb{N}_0\right)$. In lattice units the gauge field is a real one-cochain $A$ weighted by $\exp\left(-\frac{1}{2}\|dA\|^2\right)\cdot\det\left(D_A+\boldsymbol{\mu}_N\right)$, $D_A$ being the Wilson operator with Wilson parameter one and link phases $\exp(iqA)$, $q=e\cdot\eta^{\frac{1}{2}}$, $e>0$ the charge, on four-component spinors antiperiodic in every direction, and $\boldsymbol{\mu}_N=\bar{m}\cdot\eta+m_{0,N}$ the lattice mass, with $T$-independent bare mass shift $m_{0,N}$. The harmonic-averaged law ($\mathrm{av}$) of the curvature $F^u$ averages the constant part of $A$ over one period of the determinant; the all-flux law ($\mathrm{af}$) also sums the flux sectors. Let $P^{\mathrm{nat}}_{N,T}$, $\mathrm{nat}\in\{\mathrm{af},\mathrm{av}\}$, be the law of the periodic step function $\eta^{-\frac{3}{2}}\cdot F^u$ in $\mathcal{S}'(\mathbb{R}^3;\Lambda^2)$; limit points are weak limits on $H_{-4}$ as $N,T\to\infty$ in any manner. The complete definitions of these objects are given in Section~\ref{sec:setup}.

By Part~I of Theorem~\ref{thm:main-theorem}, for every charge, every sequence $\left(m_{0,N}\right)$ with $\boldsymbol{\mu}_N>-\frac{1}{2}$ and both laws, the family $\left(P^{\mathrm{nat}}_{N,T}\right)$ is tight on $H_{-4}$. Every limit point $\mu$ satisfies $\int\mathrm{d}\mu(F)\,\exp F(f)\le\exp\left(\frac{1}{2}C_0(f,f)\right)$ for real $f\in\mathcal{S}(\mathbb{R}^3;\Lambda^2)$, $C_0$ the free Maxwell covariance, is concentrated on closed two-forms, satisfies the linear growth condition on all tensor tests, is invariant under translations and the cubic group, and is reflection positive at every plane $\{x_j=c\}$. Almost every component $\mu_{\xi}$ of its ergodic decomposition under $\mathbb{R}^3$ satisfies~\eqref{eq:setup-G}, $\int\mathrm{d}\mu_{\xi}(F)\,\exp F(f)\le\exp\left(\frac{1}{2}\|f\|_2^2\right)$, is concentrated on closed two-forms, is reflection positive at every such plane, and clusters under translations $\tau_a$ as $\min_l|a_l|\to\infty$.

Part~II concerns probes, real $\beta\ne0$ in $C_c^{\infty}(\mathbb{R}^3;\Lambda^2)$, with $f=dd^{*}\beta$ and the free coefficient
\[
b(f)=-6\operatorname{Tr}\left[\left(S_{\bar{m}}\,i\gamma\cdot d^{*}\beta\right)^4\right],\qquad S_{\bar{m}}=\left(\bar{m}+\gamma\cdot\partial\right)^{-1}
\]
(Section~\ref{sec:free-coefficient}); probes with $b(f)>0$ are given in Theorem~\ref{thm:free-coefficient-1}(a). There are real numbers $m_{0,N}(e)$, $\left|m_{0,N}(e)\right|<\frac{1}{2}$, functions of $\left(N,e,\bar{m}\right)$ and of the renormalisation group scheme but not of $T$, such that for every probe with $b(f)>0$ and every point $x_0$ as in Section~\ref{sec:probes} there are $e_2>0$ and $C_{\mathrm{kap}}<\infty$, depending only on the parameters of Section~\ref{sec:parameters-scales} and on $\left(\bar{m},\beta,x_0\right)$, such that for $0<e\le e_2$ every limit point $\mu$ of either family with $m_{0,N}=m_{0,N}(e)$ satisfies
\[
\left|\kappa_4^{\mu}(f)-e^4\cdot b(f)\right|\le C_{\mathrm{kap}}\cdot e^{4+\sigma},
\]
$\kappa_4^{\mu}(f)$ being the fourth cumulant of $F(f)$ under $\mu$, and so does almost every ergodic component; $\sigma>0$ depends on the parameters only. Every such $\mu$ is neither Gaussian nor a mixture of translation-invariant Gaussian laws satisfying~\eqref{eq:setup-G}, and almost every ergodic component is not Gaussian.

Ba{\l}aban's block spin renormalisation group for lattice gauge fields~\cite{Balaban} and lattice fermions in an external gauge field~\cite{BalabanOCarrollSchor}, in Dimock's form for three-dimensional QED~\cite{DimockUR,DimockMB,DimockUVS}, is the scheme of Section~\ref{sec:renormalisation-group-scheme}: block averaging, axial gauge for fluctuation fields, Landau gauge for background fields, fermionic polymer activities with H\"older derivatives. For Gaussian domination by reflection and infrared bounds see~\cite{FSS76}; for reflection positivity and chessboard estimates,~\cite{FILS78}; for positivity of the transfer matrix of Wilson fermions,~\cite{Luescher77,OsterwalderSeiler78}; for reconstruction under the linear growth condition,~\cite{OS2}; for ergodic components of reflection positive laws,~\cite{Froehlich76} and~\cite[Section~19.7]{GlimmJaffe}.

Gaussian domination $E^{\mathrm{nat}}_n\exp\langle g,F^u\rangle\le\exp\left(\frac{1}{2}\|g\|^2\right)$ (both laws, all tori $n$, real $g\in C^2(n)$), the lattice symmetries and reflection positivity give Part~I without expansion. Chessboard estimates and subadditivity in each side length of $\log\boldsymbol{\omega}_n$ on products of all mirror copies of one cell observable, with convexity of $\log Z_n=\log\boldsymbol{\omega}_n(1)$, bound all-flux expectations of products of observables in distinct cells on tori of side $k\cdot R^{*}(e)$, $k\in\{1,L,L+1,\dots\}$, by numbers formed for $k\in\{1,L\}$ (Theorem~\ref{thm:transfer-all-flux}). On these two tori the renormalisation group, with an external connection in the determinant, has no large-field region: Gaussian domination and a union bound estimate the complement of the small-field set.

Section~\ref{sec:setup} fixes the objects of the paper and the data of Theorem~\ref{thm:main-theorem}; these data, the parameters among them, are chosen in the twelve stages listed in Section~\ref{sec:parameters-scales}. Section~\ref{sec:main-theorem} states Theorem~\ref{thm:main-theorem} and Corollaries~\ref{cor:ergodic-components} and~\ref{cor:zero-harmonic-law}. Section~\ref{sec:estimates} states Estimates~1--10, each followed by the places where it is stated in full and proved, and deduces Theorem~\ref{thm:main-theorem} from them in outline. Sections~\ref{sec:gaussian-bounds-curvature}--\ref{sec:fourth-cumulants-mixtures} use no expansion, and their results hold for every charge. Section~\ref{sec:gaussian-bounds-curvature} establishes Gaussian domination for the harmonic-averaged law on a finite torus and reflection positivity of the laws with a Proca term, of which the harmonic-averaged law is the limit; reflection positivity of the harmonic-averaged law itself is Proposition~\ref{prop:reflection-positivity}. Section~\ref{sec:continuum-infinite-volume} treats limit points of admissible families of lattice laws, Section~\ref{sec:ergodic-components-laws} the ergodic decomposition of a law that satisfies~\eqref{eq:setup-H}, and Section~\ref{sec:fourth-cumulants-mixtures} fourth cumulants of mixtures. Sections~\ref{sec:partition-functions-all} and~\ref{sec:all-flux-functional} concern the all-flux law: trace formulas, Gaussian domination, convexity of $\log Z_n$ in the side lengths, flux sectors and external fields.

The renormalisation group on the reference tori occupies Section~\ref{sec:renormalisation-group-scheme} and Sections~\ref{sec:wilson-operators-block}--\ref{sec:four-jets-one}. Section~\ref{sec:renormalisation-group-scheme} collects the definitions of the scheme, all statements about it, and the constants together with the conditions imposed on them. Apart from a few short deductions it contains no proofs; each of its statements is followed by the place where it is proved. These proofs are given in Sections~\ref{sec:wilson-operators-block}--\ref{sec:four-jets-one}, with the exception of the proofs of the two theorems of Section~\ref{sec:free-coefficients-bare}, which are given in Section~\ref{sec:free-coefficient-chessboard}. The subjects of Sections~\ref{sec:wilson-operators-block}--\ref{sec:four-jets-one} are Wilson operators with block penalties at zero field (Section~\ref{sec:wilson-operators-block}, applied in Section~\ref{sec:operators-fermion-field}), polymer activities (Section~\ref{sec:polymer-activities-algebra}), the identities of one step (Section~\ref{sec:algebraic-gaussian-identities}), the operators of the photon field (Section~\ref{sec:operators-photon-field}) and of the fermion field (Section~\ref{sec:operators-fermion-field}), one step of the renormalisation group (Section~\ref{sec:renormalisation-group-step}), the range of the running masses (Section~\ref{sec:range-running-masses}), the flow, the bare mass shifts and the partition functions of the reference tori (Section~\ref{sec:renormalisation-flow-tuned}), and four-jets (Section~\ref{sec:four-jets-one}). Section~\ref{sec:free-coefficient-chessboard} contains the free coefficient, the chessboard estimates and the reduction of all tori to the two reference tori. It proves Part~I of Theorem~\ref{thm:main-theorem} in Section~\ref{sec:admissibility-limit-points}, Part~II in Section~\ref{sec:proof-part-ii-2} and the corollaries in Section~\ref{sec:proofs-corollaries-main}; Section~\ref{sec:examples-counterexamples} gives examples showing which conclusions do not follow from Theorem~\ref{thm:main-theorem} and its corollaries.

Each of Sections~\ref{sec:gaussian-bounds-curvature}--\ref{sec:wilson-operators-block} fixes its own notation at its beginning, in particular the letters for the laws and the labelling of the directions, and states there, or at the place named in its first sentence, how its objects are identified with those of Section~\ref{sec:setup}; for Section~\ref{sec:wilson-operators-block} the identification is with the objects of Section~\ref{sec:renormalisation-group-scheme}. Sections~\ref{sec:polymer-activities-algebra}--\ref{sec:four-jets-one} use the notation of Section~\ref{sec:renormalisation-group-scheme}, to which Section~\ref{sec:algebraic-gaussian-identities} adds notation of its own for lattices, cochains and blocks (Section~\ref{sec:preliminaries}). Section~\ref{sec:free-coefficient-chessboard} uses the notation of Section~\ref{sec:setup}; Section~\ref{sec:notation-definitions} recalls it and identifies with it the objects of Sections~\ref{sec:renormalisation-group-scheme}--\ref{sec:all-flux-functional}. For this reason some objects, among them the Wilson operator, the two lattice laws, the hypothesis~\eqref{eq:setup-H}, the ergodic decomposition and the probes, are defined in more than one section, and some facts are proved in more than one section, each time in the form needed there.

The sections are ordered by subject, and the proofs depend on one another as follows. Section~\ref{sec:continuum-infinite-volume} uses Section~\ref{sec:gaussian-bounds-curvature}; Sections~\ref{sec:three-examples-admissible} and~\ref{sec:pinned-averaged-laws} use Theorem~\ref{thm:free-coefficient-1} as well. Sections~\ref{sec:partition-functions-all} and~\ref{sec:all-flux-functional} use Section~\ref{sec:gaussian-bounds-curvature} and Section~\ref{sec:discretised-probes-lattice}. Sections~\ref{sec:admissibility-limit-points}--\ref{sec:partition-functions-cubic} use Sections~\ref{sec:gaussian-bounds-curvature}--\ref{sec:all-flux-functional}. Section~\ref{sec:operators-fermion-field} uses Section~\ref{sec:wilson-operators-block}. Section~\ref{sec:renormalisation-group-step} uses Sections~\ref{sec:polymer-activities-algebra}--\ref{sec:operators-fermion-field}, and within Sections~\ref{sec:polymer-activities-algebra}--\ref{sec:renormalisation-group-step} single statements are used in both directions. Section~\ref{sec:renormalisation-flow-tuned} uses Section~\ref{sec:gaussian-bounds-curvature} and Sections~\ref{sec:polymer-activities-algebra}--\ref{sec:renormalisation-group-step}; conversely, Propositions~\ref{prop:depth-stop-level} and~\ref{prop:consequences} are used in the proof of Proposition~\ref{prop:free-determinant}, which is given in Section~\ref{sec:polymer-activities-algebra}. Section~\ref{sec:four-jets-one} uses Section~\ref{sec:gaussian-bounds-curvature} and Sections~\ref{sec:polymer-activities-algebra}--\ref{sec:renormalisation-flow-tuned}. Section~\ref{sec:proof-theorem} uses Sections~\ref{sec:renormalisation-flow-tuned}, \ref{sec:four-jets-one}, \ref{sec:discretised-probes-lattice} and~\ref{sec:continuum-infinite-volume-2}, and Section~\ref{sec:proof-part-ii} draws on all sections.

\section{Setup}\label{sec:setup}

In this section the objects in terms of which Theorem~\ref{thm:main-theorem} is stated are defined, and the data of that theorem are fixed; a fact that is stated here without proof is followed, in parentheses, by a reference to the place where it is proved. The odd integer $L$ and the mass parameter $\bar{m}$, which occur from Section~\ref{sec:lattice-tori} on, are chosen at stages~\ref{it:order-choice-4} and~\ref{it:order-choice-11} of Section~\ref{sec:parameters-scales}.

\subsection{Lattice tori}\label{sec:lattice-tori}

For the odd integer $L$ and the mass parameter $\bar{m}\in(0,1]$ let
\[
h:=L^{\left\lfloor\log_L\bar{m}^{-1}\right\rfloor};
\]
this power of $L$ is called the \emph{scale}. The exponent $j:=\left\lfloor\log_L\bar{m}^{-1}\right\rfloor$ is an integer with $j\ge0$ and $L^j\le\bar{m}^{-1}<L^{j+1}$, as $0<\bar{m}\le1$; hence
\[
h\ge1,\qquad\frac{1}{L}<\bar{m}\cdot h\le1.
\]
For integers $N,T\ge0$ put
\[
\eta:=L^{-N},\qquad R:=h\cdot L^T,\qquad n:=\frac{R}{\eta}=h\cdot L^{T+N},\qquad\Lambda_{\eta,R}:=\eta\mathbb{Z}^3/R\mathbb{Z}^3;
\]
$\eta$ is the lattice spacing and $R$ the side of the torus $\Lambda_{\eta,R}$. The number $n$ is a power of the odd integer $L$, hence odd, and $n=1$ occurs only for $h=1$ and $N=T=0$. This case is excluded: a pair $(N,T)$ of integers $N,T\ge0$ with $n\ge2$ is called a \emph{cutoff}, and only cutoffs are considered.

\paragraph{Cells and cochains.}
Let $n=(n_1,n_2,n_3)$ be a triple of integers $n_{\mu}\ge2$, and put
\[
\Lambda_n:=\prod_{\mu=1}^3\mathbb{Z}/n_{\mu}\mathbb{Z},\qquad|n|:=n_1\cdot n_2\cdot n_3;
\]
a single integer $n\ge2$ also stands for the triple $(n,n,n)$, and $\Lambda_n$ is then the cubic torus with $n$ sites per side. For $\mu=1,2,3$ let $e_{\mu}$ denote the $\mu$-th unit vector of $\mathbb{Z}^3\subset\mathbb{R}^3$ as well as its image in $\Lambda_n$. For $k\in\{0,1,2,3\}$ a $k$-cell of $\Lambda_n$ is a pair $(x;I)$ of a site $x\in\Lambda_n$ and a set $I=\{i_1,\dots,i_k\}\subset\{1,2,3\}$ with $i_1<\dots<i_k$, so that $\#|I|=k$; it is the cell spanned at $x$ by $e_{i_1},\dots,e_{i_k}$, positively oriented by $e_{i_1}\wedge\dots\wedge e_{i_k}$. The one-cells are the links, the two-cells the plaquettes, and $(x;\mu\nu):=(x;\{\mu,\nu\})$ for $\mu<\nu$. Let $C^k(n)$ be the space of the real $k$-cochains, that is, of the real functions on the $k$-cells, with the counting pairing and the norm
\[
\langle u,v\rangle:=\sum_{x\in\Lambda_n}\ \sum_{I\subset\{1,2,3\}:\ \#|I|=k}u(x;I)\cdot v(x;I),\qquad\|u\|:=\langle u,u\rangle^{\frac{1}{2}};
\]
let $C^k(n;\mathbb{Z})\subset C^k(n)$ be the set of the integer-valued $k$-cochains, and call the elements of $C^k(n)+iC^k(n)$ complex $k$-cochains. The values of $\boldsymbol{\lambda}\in C^0(n)$, $A\in C^1(n)$ and $F\in C^2(n)$ at the cells $(x;\emptyset)$, $(x;\{\mu\})$ and $(x;\mu\nu)$ are written $\boldsymbol{\lambda}(x)$, $A_{\mu}(x)$ and $F_{\mu\nu}(x)$, and the value of a three-cochain at $(x;\{1,2,3\})$ carries the index $123$. The coboundary $d\colon C^k(n)\to C^{k+1}(n)$, $k=0,1,2$, is given by
\begin{align*}
(d\boldsymbol{\lambda})_{\mu}(x)&=\boldsymbol{\lambda}(x+e_{\mu})-\boldsymbol{\lambda}(x),\\
(dA)_{\mu\nu}(x)&=A_{\mu}(x)+A_{\nu}(x+e_{\mu})-A_{\mu}(x+e_{\nu})-A_{\nu}(x)\qquad(\mu<\nu),\\
(dF)_{123}(x)&=\left(F_{23}(x+e_1)-F_{23}(x)\right)-\left(F_{13}(x+e_2)-F_{13}(x)\right)\\
&\qquad+\left(F_{12}(x+e_3)-F_{12}(x)\right),
\end{align*}
and a cochain $u$ with $du=0$ is called closed.

Let $\eta\in(0,1]$ and $R\ge1$ be such that $n:=R/\eta$ is an integer $\ge2$, as at every cutoff, and let $\Lambda_{\eta,R}:=\eta\mathbb{Z}^3/R\mathbb{Z}^3$. The $k$-cells $(x;I)$ of $\Lambda_{\eta,R}$, with $x\in\Lambda_{\eta,R}$, the space $C^k(\Lambda_{\eta,R})$ of its real $k$-cochains with the counting pairing $\langle\cdot,\cdot\rangle$ and the norm $\|\cdot\|$, and its coboundary $d$ are defined as for $\Lambda_n$, with $\eta\cdot e_{\mu}$ in place of $e_{\mu}$: the cell $(x;I)$ is spanned at $x$ by the vectors $\eta\cdot e_i$, $i\in I$, and $d$ is given by the formulas above with $x+\eta\cdot e_{\mu}$ in place of $x+e_{\mu}$. The bijection $x\mapsto\eta^{-1}\cdot x$ of $\Lambda_{\eta,R}$ onto the cubic torus $\Lambda_n$ carries these objects onto those of $\Lambda_n$; through it they are read in lattice units.

An oriented nearest-neighbour bond, or simply a bond, is an ordered pair $(x,x')$ of points of $\mathbb{Z}^3$ with $|x'-x|_2=1$, where $|\cdot|_2$ is the Euclidean norm. Every real or complex one-cochain $B$ defines a function on the bonds, again denoted $B$, by
\[
B(x,x+e_{\mu}):=B_{\mu}(x),\qquad B(x',x):=-B(x,x')\qquad\left(x\in\mathbb{Z}^3,\ \mu=1,2,3\right),
\]
$B_{\mu}(x)$ being the value of $B$ at the link $(\bar{x};\{\mu\})$, where $\bar{x}$ is the image of $x$ in $\Lambda_n$. This function is periodic: its value does not change when both end points of a bond are translated by $n_{\mu}\cdot e_{\mu}$, $\mu=1,2,3$.

The space of the one-cochains is the direct sum
\[
C^1(n)=V_{\mathrm{g}}\oplus V_{\mathrm{h}}\oplus V_{\mathrm{t}},
\]
orthogonal for the counting pairing, where $V_{\mathrm{g}}:=dC^0(n)$ is the space of the exact one-cochains, $V_{\mathrm{h}}$ is the space of the constant one-cochains, that is, of the $A\in C^1(n)$ for which $A_{\mu}(x)$ does not depend on $x$ ($\mu=1,2,3$), and $V_{\mathrm{t}}$ is the orthogonal complement of $V_{\mathrm{g}}\oplus V_{\mathrm{h}}$ in $C^1(n)$. (The spaces $V_{\mathrm{g}}$ and $V_{\mathrm{h}}$ are orthogonal, because $\sum_{x\in\Lambda_n}\left(\boldsymbol{\lambda}(x+e_{\mu})-\boldsymbol{\lambda}(x)\right)=0$ for every $\boldsymbol{\lambda}\in C^0(n)$.) The kernel of $d$ on $C^1(n)$ is $V_{\mathrm{g}}\oplus V_{\mathrm{h}}$ (Lemma~\ref{lem:harmonic-cochains-coercivity}, Remark~\ref{rem:tori-unequal-sides}); hence $d$ is injective on $V_{\mathrm{t}}$.

\subsection{Spinors and the Wilson determinant}\label{sec:spinors-wilson-determinant}

Let $\sigma_1,\sigma_2,\sigma_3$ be the Pauli matrices, and for $\mu=1,2,3$ let
\[
\gamma_{\mu}:=\operatorname{diag}\left(\sigma_{\mu},-\sigma_{\mu}\right)
\]
be the block-diagonal $4\times4$ matrix with the diagonal blocks $\sigma_{\mu}$ and $-\sigma_{\mu}$. As the Pauli matrices are Hermitian and satisfy $\sigma_j\cdot\sigma_k+\sigma_k\cdot\sigma_j=2\delta_{jk}$, the matrices $\gamma_{\mu}$ are Hermitian and
\[
\gamma_{\mu}\cdot\gamma_{\nu}+\gamma_{\nu}\cdot\gamma_{\mu}=2\delta_{\mu\nu}\qquad\left(\mu,\nu\in\{1,2,3\}\right).
\]
For $v\in\mathbb{R}^3$ put $\gamma\cdot v:=\sum_{\mu}v_{\mu}\cdot\gamma_{\mu}$.

A spinor on $\Lambda_n$ is a function $\psi\colon\mathbb{Z}^3\to\mathbb{C}^4$ that is antiperiodic in every direction:
\[
\psi(x+n_{\mu}\cdot e_{\mu})=-\psi(x)\qquad\left(x\in\mathbb{Z}^3,\ \mu=1,2,3\right).
\]
The spinors form a complex vector space of dimension $4|n|$, since a spinor is determined by its values on $\prod_{\mu}\{0,\dots,n_{\mu}-1\}$ and these values can be prescribed arbitrarily.

Let $\boldsymbol{\mu}$ be a real number; it is called the lattice mass. Let $B$ be a complex one-cochain. The function on the bonds of $\mathbb{Z}^3$ that $B$ defines is periodic (Section~\ref{sec:lattice-tori}), and with its values $B(x,x')$ one puts, for a spinor $\psi$ and $x\in\mathbb{Z}^3$,
\[
\left(D[B]\psi\right)(x):=3\psi(x)-\sum_{x'\in\mathbb{Z}^3:\ |x'-x|_2=1}\frac{1}{2}\left(1-\gamma\cdot(x'-x)\right)\cdot\exp\left(iB(x,x')\right)\cdot\psi(x').
\]
As $B$ is periodic and $\psi$ antiperiodic, $D[B]\psi$ is again a spinor; thus $D[B]$ is a linear operator on the space of the spinors, and
\[
\mathrm{Dt}(B):=\det\left(D[B]+\boldsymbol{\mu}\right)
\]
is the determinant of $D[B]+\boldsymbol{\mu}$ on this space. One calls $D[B]$ the Wilson operator; its Wilson parameter equals one. Let $q>0$; this number is called the lattice charge, and every real one-cochain $A$ defines the operator $D_A:=D[q\cdot A]$.

Let $B$ be a real one-cochain. Then
\begin{gather*}
\mathrm{Dt}(B)\ge0,\qquad\mathrm{Dt}(B+d\boldsymbol{\lambda})=\mathrm{Dt}(B),\\
\mathrm{Dt}(B+2\pi\cdot\mathfrak{z})=\mathrm{Dt}(B),\qquad\mathrm{Dt}(-B)=\mathrm{Dt}(B)
\end{gather*}
for every zero-cochain $\boldsymbol{\lambda}\in C^0(n)$ and every integer one-cochain $\mathfrak{z}\in C^1(n;\mathbb{Z})$; and $\mathrm{Dt}$ is invariant under the lattice isometries: if $r$ is an orthogonal $3\times3$ matrix with integer entries such that $r\left(\prod_{\mu=1}^3n_{\mu}\mathbb{Z}\right)=\prod_{\mu=1}^3n_{\mu}\mathbb{Z}$, and if $c\in\mathbb{Z}^3$, then the map $g\colon x\mapsto rx+c$ of $\mathbb{Z}^3$ induces an isometry of $\Lambda_n$, and $\mathrm{Dt}\left(g_{*}B\right)=\mathrm{Dt}(B)$, where $g_{*}B$ is the one-cochain with $\left(g_{*}B\right)(gx,gx')=B(x,x')$ for every bond $(x,x')$ of $\mathbb{Z}^3$ (Theorem~\ref{thm:free-coefficient-1}(d)). This one-cochain is well defined. Indeed, $g$ permutes the bonds of $\mathbb{Z}^3$, since $r$ is orthogonal and $r^{-1}=r^T$ has integer entries; the function $(gx,gx')\mapsto B(x,x')$ on the bonds is odd, and it is periodic because $r$ maps $\prod_{\mu=1}^3n_{\mu}\mathbb{Z}$ onto itself; and an odd periodic function on the bonds is the function of exactly one one-cochain.

\subsection{The lattice laws}\label{sec:lattice-laws}

In this subsection $n=(n_1,n_2,n_3)$ is a triple of integers $n_{\mu}\ge2$, and the notation of Sections~\ref{sec:lattice-tori} and~\ref{sec:spinors-wilson-determinant} is used for the torus $\Lambda_n$, in lattice units. Let $q>0$ be the lattice charge and $\boldsymbol{\mu}>-\frac{1}{2}$ the lattice mass, and put $r_q:=2\pi/q$.

\paragraph{The all-flux functional.}
Let $\mathfrak{F}(n)$ denote the set of the complex valued Borel functions $\Phi$ on $C^2(n)$ that grow at most exponentially: $|\Phi(F)|\le c\cdot\exp\left(c\|F\|\right)$ for all $F\in C^2(n)$, with a constant $c$ that depends on $\Phi$. For $\mu<\nu$ let $\iota$ denote the remaining direction, $\{\iota\}=\{1,2,3\}\setminus\{\mu,\nu\}$; for each $t\in\mathbb{Z}/n_{\iota}\mathbb{Z}$ the set of the plaquettes $(x;\mu\nu)$ with $x_{\iota}=t$ is a coordinate two-torus in the direction $(\mu\nu)$. To a closed $\mathfrak{z}\in C^2(n;\mathbb{Z})$ one attaches its \emph{flux}
\[
w(\mathfrak{z}):=\left(w_{12}(\mathfrak{z}),w_{13}(\mathfrak{z}),w_{23}(\mathfrak{z})\right)\in\mathbb{Z}^3,\qquad w_{\mu\nu}(\mathfrak{z}):=\sum_{x\in\Lambda_n:\ x_{\iota}=t}\mathfrak{z}_{\mu\nu}(x),
\]
the component $w_{\mu\nu}(\mathfrak{z})$ being the sum of $\mathfrak{z}$ over the plaquettes of a coordinate two-torus in the direction $(\mu\nu)$. By~\eqref{eq:partition-functions-2}, this sum is the same for all $t$, that is, for all the parallel two-tori; and $\mathfrak{z}=dm$ for some $m\in C^1(n;\mathbb{Z})$ if and only if $w(\mathfrak{z})=0$.

Denote by $\mathrm{d}A$ the Lebesgue measure on $C^1(n)$, and by $[0,r_q)^{\mathrm{links}}$ the box formed by the $A\in C^1(n)$ whose value $A_{\mu}(x)$ at each link $(x;\{\mu\})$ lies in $[0,r_q)$. For $\Phi\in\mathfrak{F}(n)$ put
\begin{multline*}
\boldsymbol{\omega}_n(\Phi):=\int_{[0,r_q)^{\mathrm{links}}}\mathrm{d}A\sum_{\substack{\mathfrak{z}\in C^2(n;\mathbb{Z})\\ d\mathfrak{z}=0}}\Phi\left(dA+r_q\cdot\mathfrak{z}\right)\\
\cdot\exp\left(-\frac{1}{2}\left\|dA+r_q\cdot\mathfrak{z}\right\|^2\right)\cdot\mathrm{Dt}(qA)\\
=\sum_{w\in\mathbb{Z}^3}\boldsymbol{\omega}_n^{(w)}(\Phi),
\end{multline*}
where, for $w\in\mathbb{Z}^3$, the number $\boldsymbol{\omega}_n^{(w)}(\Phi)$ is given by the same integral, with the sum extended only over the closed $\mathfrak{z}$ of flux $w(\mathfrak{z})=w$. For every $\Phi\in\mathfrak{F}(n)$ the series is absolutely convergent, and
\[
Z_n:=\boldsymbol{\omega}_n(1)>0,\qquad\mathbf{Z}_n:=\boldsymbol{\omega}_n^{(0)}(1)>0
\]
(Sections~\ref{sec:summation-all-fluxes} and~\ref{sec:decomposition-flux}). The argument of $\Phi$ is the \emph{curvature in lattice units}
\[
F^u:=dA+r_q\cdot\mathfrak{z}\in C^2(n),
\]
a closed two-cochain, as $d\circ d=0$ and $d\mathfrak{z}=0$.

\paragraph{The two laws.}
For $\Phi\in\mathfrak{F}(n)$ put
\begin{align}
E^{\mathrm{af}}_n(\Phi)&:=\frac{\boldsymbol{\omega}_n(\Phi)}{Z_n},\tag{af}\label{eq:setup-af}\\
E^{\mathrm{av}}_n(\Phi)&:=\frac{\boldsymbol{\omega}_n^{(0)}(\Phi)}{\mathbf{Z}_n}.\tag{av}\label{eq:setup-av}
\end{align}
Since $\mathrm{Dt}\ge0$ on the real one-cochains (Section~\ref{sec:spinors-wilson-determinant}), $E^{\mathrm{af}}_n(\Phi)$ is the expectation of $\Phi(F^u)$ when the pair $(A,\mathfrak{z})$, with $A\in[0,r_q)^{\mathrm{links}}$ and $\mathfrak{z}\in C^2(n;\mathbb{Z})$ closed, is distributed with the density $Z_n^{-1}\cdot\exp\left(-\frac{1}{2}\|F^u\|^2\right)\cdot\mathrm{Dt}(qA)$ with respect to the product of $\mathrm{d}A$ and the counting measure; and $E^{\mathrm{av}}_n(\Phi)$ is the same expectation with $\mathfrak{z}$ restricted to the flux $0$ and with $\mathbf{Z}_n$ in place of $Z_n$. The two laws of the random closed two-cochain $F^u$ obtained in this way are the \emph{all-flux law} and the \emph{harmonic-averaged law} on $\Lambda_n$; they are again denoted $E^{\mathrm{af}}_n$ and $E^{\mathrm{av}}_n$, and the upper index $\mathrm{nat}$ stands for $\mathrm{af}$ or for $\mathrm{av}$.

Let $\sigma$ stand for the Lebesgue measure on $V_{\mathrm{t}}$ and, for $\theta\in[-\pi,\pi]^3$, let $H_{\theta}\in V_{\mathrm{h}}$ stand for the constant one-cochain whose value at each link $(x;\{\mu\})$ is $\theta_{\mu}/(q\cdot n_{\mu})$. By Lemma~\ref{lem:flux-sectors}(b), the harmonic-averaged law has a second description: $E^{\mathrm{av}}_n$ coincides with the law of the coboundary $F^u=dA$ of a random $A\in V_{\mathrm{t}}$ whose law is the measure
\[
\left(\int_{[-\pi,\pi]^3}\mathrm{d}\theta\,\exp\left(-\frac{1}{2}\|dA\|^2\right)\cdot\det\left(D_{A+H_{\theta}}+\boldsymbol{\mu}\right)\right)\mathrm{d}\sigma(A)
\]
divided by its total mass.

In both laws the curvature $F^u$ carries the Maxwell weight $\exp\left(-\frac{1}{2}\|F^u\|^2\right)$, and the fermions, integrated out, enter through the determinant alone. The gauge field of both laws is noncompact, that is, real valued. In the second description of the harmonic-averaged law it is the variable $A$ of the real vector space $V_{\mathrm{t}}$: the component in $V_{\mathrm{g}}$ of a one-cochain, its exact part, does not occur, owing to gauge invariance, and its component in $V_{\mathrm{h}}$, the constant part, which enters neither $dA$ nor, therefore, the Maxwell weight, is replaced by $H_{\theta}$ and averaged over $\theta\in[-\pi,\pi]^3$, one period in each variable of the function $\theta\mapsto\det\left(D_{A+H_{\theta}}+\boldsymbol{\mu}\right)$. The all-flux law sums in addition over all the fluxes $w\in\mathbb{Z}^3$ of the torus, and its gauge field is a real one-cochain on every contractible region, in the following sense. Let $U$ be a union of closed cells of $\Lambda_n$ that is contractible, and let $\mathfrak{z}\in C^2(n;\mathbb{Z})$ be closed. For some $m\in C^1(n;\mathbb{Z})$ the cochain $\mathfrak{z}-dm$ vanishes on the plaquettes of $U$; the pairs $(A,\mathfrak{z})$ and $\left(A+r_q\cdot m,\mathfrak{z}-dm\right)$ have the same curvature $F^u$ and, as $\mathrm{Dt}(qA+2\pi\cdot m)=\mathrm{Dt}(qA)$, the same determinant. On $U$ the pair $(A,\mathfrak{z})$ is therefore equivalent to $(A',0)$, where $A':=A+r_q\cdot m$ is a real one-cochain that need not lie in $[0,r_q)^{\mathrm{links}}$; there the integrand of the all-flux law reduces to the flux-free weight $\exp\left(-\frac{1}{2}\|dA'\|^2\right)\cdot\mathrm{Dt}(qA')$ of the noncompact gauge field (Section~\ref{sec:summation-all-fluxes}).

\paragraph{The zero-harmonic law.}
Let now $n\ge3$ be an odd integer, and let $\Lambda_n$ be the cubic torus $(n,n,n)$, which has $n$ sites on each side; the torus of every cutoff, with $n=h\cdot L^{T+N}$, is of this kind. Then
\[
0<\int_{V_{\mathrm{t}}}\mathrm{d}\sigma(A)\,\exp\left(-\frac{1}{2}\|dA\|^2\right)\cdot\det\left(D_A+\boldsymbol{\mu}\right)<\infty
\]
(Section~\ref{sec:wilson-operator-lattice}). The \emph{zero-harmonic law} on $\Lambda_n$ is the law of $F^u=dA$ when $A$ is distributed on $V_{\mathrm{t}}$ according to the probability measure proportional to $\exp\left(-\frac{1}{2}\|dA\|^2\right)\cdot\det\left(D_A+\boldsymbol{\mu}\right)\mathrm{d}\sigma(A)$.

\subsection{The curvature field and its limit points}\label{sec:curvature-field-limit}

\paragraph{Cells and cochains in physical units.}
For a cutoff $(N,T)$ let $\eta=L^{-N}$, $R=h\cdot L^T$ and $n=R/\eta$ be as in Section~\ref{sec:lattice-tori}, and let the cochains of $\Lambda_{\eta,R}$ be read in lattice units, as cochains of the cubic torus $\Lambda_n$, through the bijection fixed there; thus $C^k(\Lambda_{\eta,R})=C^k(n)$. Let $\mathbb{T}_R:=\mathbb{R}^3/R\mathbb{Z}^3$, a group that contains $\Lambda_{\eta,R}$. For $x\in\Lambda_{\eta,R}$ and $\mu<\nu$ the plaquette $\boldsymbol{\sigma}=(x;\mu\nu)$ has the \emph{centre} and the \emph{averaging cube}
\[
c(\boldsymbol{\sigma}):=x+\frac{\eta}{2}\left(e_{\mu}+e_{\nu}\right)\in\mathbb{T}_R,\qquad Q(\boldsymbol{\sigma}):=c(\boldsymbol{\sigma})+\left[-\frac{\eta}{2},\frac{\eta}{2}\right)^3\subset\mathbb{T}_R;
\]
for fixed $\mu<\nu$ the cubes $Q(x;\mu\nu)$, $x\in\Lambda_{\eta,R}$, form a partition of $\mathbb{T}_R$. For $k$-cochains $u,v$ on $\Lambda_{\eta,R}$ let
\[
\langle u,v\rangle_{\eta}:=\eta^3\sum_{\boldsymbol{\sigma}}u(\boldsymbol{\sigma})\cdot v(\boldsymbol{\sigma}),\qquad|u|_{\eta}:=\langle u,u\rangle_{\eta}^{\frac{1}{2}},
\]
the sum being taken over all $k$-cells $\boldsymbol{\sigma}$; the pairing is extended to complex cochains by the same formula. Let $d_{\eta}:=\eta^{-1}\cdot d$, and let $d_{\eta}^{*}$ be the adjoint of $d_{\eta}$ with respect to $\langle\cdot,\cdot\rangle_{\eta}$.

\paragraph{Tests and distributions.}
Let $\mathcal{S}=\mathcal{S}(\mathbb{R}^3;\Lambda^2)$ be the space of the two-forms $f=\left(f_{\mu\nu}\right)_{\mu<\nu}$ on $\mathbb{R}^3$ whose coefficients are Schwartz functions, real or complex; $f$ is called real if its coefficients are real valued. Let $\mathcal{S}'=\mathcal{S}'(\mathbb{R}^3;\Lambda^2)$ be the dual of the real space $\mathcal{S}$: its elements are the continuous linear forms $F$ on the real elements of $\mathcal{S}$, each extended complex-linearly to all of $\mathcal{S}$, so that $F(f)$ is real for real $f$. The space $\mathcal{S}'$ carries the cylinder $\sigma$-algebra, the $\sigma$-algebra generated by the maps $F\mapsto F(f)$, $f\in\mathcal{S}$. The norm of $L^2(\mathbb{R}^3;\Lambda^2)$ is
\[
\|f\|_2:=\left(\sum_{\mu<\nu}\int_{\mathbb{R}^3}\mathrm{d}x\,\left|f_{\mu\nu}(x)\right|^2\right)^{\frac{1}{2}}.
\]

\paragraph{The curvature field.}
At a cutoff $(N,T)$ the \emph{physical curvature} is $F:=\eta^{-\frac{3}{2}}\cdot F^u\in C^2(\Lambda_{\eta,R})$, where $F^u$ is the curvature in lattice units on the torus $n=R/\eta$ (Section~\ref{sec:lattice-laws}). For $f\in\mathcal{S}$, $F\in C^2(\Lambda_{\eta,R})$ and a plaquette $\boldsymbol{\sigma}=(x;\mu\nu)$ put
\[
f_{\eta}(\boldsymbol{\sigma}):=\eta^{-3}\int_{Q(\boldsymbol{\sigma})}\mathrm{d}y\,\sum_{k\in\mathbb{Z}^3}f_{\mu\nu}(y+R\cdot k),\qquad X_{\eta,R}(f):=\langle f_{\eta},F\rangle_{\eta},
\]
the periodised coefficient $y\mapsto\sum_{k\in\mathbb{Z}^3}f_{\mu\nu}(y+R\cdot k)$ being a function on $\mathbb{T}_R$. The map $f\mapsto f_{\eta}$ of $\mathcal{S}$ into the space $C^2(\Lambda_{\eta,R})+iC^2(\Lambda_{\eta,R})$ of the complex two-cochains, which takes real tests to real cochains, is the \emph{test map}, and $X_{\eta,R}(f)$ is the \emph{observable}. Let $F^{\sharp}=\left(F^{\sharp}_{\mu\nu}\right)_{\mu<\nu}$ be the two-form on $\mathbb{R}^3$ with the coefficients
\[
F^{\sharp}_{\mu\nu}(y):=F(\boldsymbol{\sigma})\qquad\text{if the image of }y\text{ in }\mathbb{T}_R\text{ lies in }Q(\boldsymbol{\sigma}),\ \boldsymbol{\sigma}=(x;\mu\nu);
\]
for the physical curvature $F$ it is the \emph{curvature field}. Its coefficients are bounded and $R$-periodic; therefore $F^{\sharp}$ is an element of $\mathcal{S}'$, and
\[
F^{\sharp}(f):=\int_{\mathbb{R}^3}\mathrm{d}y\,\sum_{\mu<\nu}f_{\mu\nu}(y)\cdot F^{\sharp}_{\mu\nu}(y)=X_{\eta,R}(f)\qquad(f\in\mathcal{S}).
\]

\paragraph{The spaces $H_s$.}
Let $\mathbb{N}_0$ be the set of the integers $\ge0$, and $\mathcal{K}:=\mathbb{N}_0^3\times\{12,13,23\}$. For $k\in\mathbb{N}_0^3$ let $h_k$ be the Hermite function on $\mathbb{R}^3$ of multi-index $k$, normalised by $\|h_k\|_{L^2(\mathbb{R}^3)}=1$; for $\kappa=(k,c)\in\mathcal{K}$ let $h_{\kappa}$ be the two-form with the component $h_k$ at the index $c$ and the component $0$ at the two other indices. The $h_{\kappa}$, $\kappa\in\mathcal{K}$, form an orthonormal basis of $L^2(\mathbb{R}^3;\Lambda^2)$, and they are eigenvectors of $1+|x|_2^2-\Delta$, applied componentwise:
\[
\left(1+|x|_2^2-\Delta\right)h_{\kappa}=\lambda_{\kappa}\cdot h_{\kappa},\qquad\lambda_{\kappa}:=2|k|_1+4,\qquad|k|_1:=k_1+k_2+k_3.
\]
For $s\in\mathbb{R}$ and $F\in\mathcal{S}'$ let
\[
\|F\|_s^2:=\sum_{\kappa\in\mathcal{K}}\lambda_{\kappa}^{2s}\cdot\left|F(h_{\kappa})\right|^2\in[0,\infty],\qquad H_s:=\left\{F\in\mathcal{S}':\ \|F\|_s<\infty\right\}.
\]
Each $H_s$, with the norm $\|\cdot\|_s$, is a separable Hilbert space (Section~\ref{sec:hermite-sobolev-spaces}); and $H_s\subset H_{s'}$, with $\|F\|_{s'}\le\|F\|_s$, whenever $s'\le s$, because $\lambda_{\kappa}\ge1$.

\paragraph{Charge and lattice mass.}
Let $e>0$ be the charge; at a cutoff $(N,T)$ the lattice charge is $q:=e\cdot\eta^{\frac{1}{2}}$. A sequence $N\mapsto m_{0,N}$, $N\ge0$, of real numbers, the \emph{bare mass shifts}, is called a \emph{mass assignment} if
\[
\boldsymbol{\mu}_N:=\bar{m}\cdot L^{-N}+m_{0,N}>-\frac{1}{2}\qquad\text{for every }N\ge0,
\]
where $\bar{m}$ is the mass parameter. The number $\boldsymbol{\mu}_N$ is the \emph{lattice mass} at each of the cutoffs $(N,T)$, $T\ge0$; neither $\boldsymbol{\mu}_N$ nor the bare mass shift $m_{0,N}$ depends on $T$.

\paragraph{The laws of the curvature field.}
Let $N\mapsto m_{0,N}$ be a mass assignment and $\mathrm{nat}\in\{\mathrm{af},\mathrm{av}\}$. For a cutoff $(N,T)$ let $P^{\mathrm{nat}}_{N,T}$ be the law of the curvature field $F^{\sharp}$, formed with $F=\eta^{-\frac{3}{2}}\cdot F^u$, when $F^u$ has the law $E^{\mathrm{nat}}_n$ of Section~\ref{sec:lattice-laws} with
\[
n=h\cdot L^{T+N},\qquad q=e\cdot\eta^{\frac{1}{2}},\qquad\boldsymbol{\mu}=\boldsymbol{\mu}_N.
\]
For every $F\in C^2(\Lambda_{\eta,R})$ the distribution $F^{\sharp}$ belongs to $H_{-3}$ (Section~\ref{sec:existence-limit-points}), and $F^{\sharp}$ depends linearly on $F$. Every linear map of a space of finite dimension into a normed space is continuous; so $F\mapsto F^{\sharp}$ is a continuous map of $C^2(\Lambda_{\eta,R})$ into $H_{-3}$, and $P^{\mathrm{nat}}_{N,T}$, the image of the law of $F$ under this map, is thus a Borel probability measure on $H_{-3}$. As $H_{-3}\subset H_{-4}$, with $\|\cdot\|_{-4}\le\|\cdot\|_{-3}$ on $H_{-3}$, the inclusion of $H_{-3}$ in $H_{-4}$ is continuous, and $P^{\mathrm{nat}}_{N,T}$ is also a Borel probability measure on $H_{-4}$. If the weight $\exp\left(-\frac{1}{2}\|F^u\|^2\right)\cdot\mathrm{Dt}(qA)$ is multiplied by a constant factor $\exp\left(-\boldsymbol{\varepsilon}_{0,N}\cdot n^3\right)$ with $\boldsymbol{\varepsilon}_{0,N}\in\mathbb{R}$ (a vacuum energy), the factor cancels against the normalisation, and neither $E^{\mathrm{nat}}_n$ nor $P^{\mathrm{nat}}_{N,T}$ changes.

\paragraph{Limit points.}
Let $\mathrm{nat}\in\{\mathrm{af},\mathrm{av}\}$. A \emph{limit point} of the family $\left(P^{\mathrm{nat}}_{N,T}\right)$ is, by definition, a Borel probability measure $\mu$ on $H_{-4}$ to which $P^{\mathrm{nat}}_{N_i,T_i}$ converges weakly on $H_{-4}$ along a sequence of cutoffs $(N_i,T_i)$, $i=1,2,\dots$, such that $N_i\to\infty$ and $T_i\to\infty$; beyond these two limits the sequences $(N_i)$ and $(T_i)$ are subject to no condition, and no relation between them is required (Definition~\ref{def:limit-points}). The set $H_{-4}$ is a member of the cylinder $\sigma$-algebra of $\mathcal{S}'$; moreover, the sets $B\cap H_{-4}$, with $B$ in the cylinder $\sigma$-algebra, are exactly the Borel subsets of $H_{-4}$ (Section~\ref{sec:hermite-sobolev-spaces}). So a limit point $\mu$ is also a probability measure on $\mathcal{S}'$ with its cylinder $\sigma$-algebra: the measure $B\mapsto\mu\left(B\cap H_{-4}\right)$.

\subsection{Laws on $\mathcal{S}'$}\label{sec:laws}

A law on $\mathcal{S}'$ is a probability measure on $\mathcal{S}'$ with its cylinder $\sigma$-algebra (Section~\ref{sec:curvature-field-limit}); the sets that belong to this $\sigma$-algebra are called cylinder sets. For a law $\rho$ let $L^2(\rho)$ be the Hilbert space of the complex functions on $\mathcal{S}'$ that are square integrable for $\rho$, taken modulo $\rho$-null functions; for a $\sigma$-algebra $\mathcal{G}$ contained in the cylinder $\sigma$-algebra, $L^2(\mathcal{S}',\mathcal{G},\rho)$ denotes the closed subspace of $L^2(\rho)$ formed by the classes that contain a $\mathcal{G}$-measurable function. The vectors $e_1,e_2,e_3$ form the standard basis of $\mathbb{R}^3$.

\paragraph{Actions.}
For $b\in\mathbb{R}^3$ and for $r$ in the cubic group $O_{\mathrm{cub}}:=O(3)\cap\mathrm{GL}(3,\mathbb{Z})$ let $\tau_b$ and $U_r$ be the maps of $\mathcal{S}$ into itself given by
\[
(\tau_bf)(y):=f(y-b),\qquad(U_rf)(y):=\left(\wedge^2r\right)f\left(r^{-1}y\right)\qquad\left(f\in\mathcal{S},\ y\in\mathbb{R}^3\right),
\]
where $\wedge^2r$ is the linear map of $\Lambda^2$ induced by $r$, extended complex-linearly to complex tests; thus $U_rf$ is the pull-back of the two-form $f$ by $r^{-1}$. For $\nu\in\{1,2,3\}$ let $r_{\nu}\in O_{\mathrm{cub}}$ be the reflection that changes the sign of the $\nu$-th coordinate, and let $\mathsf{R}_{\nu}:=U_{r_{\nu}}$. It changes the sign of those components of $f$ whose index contains $\nu$:
\[
\left(\mathsf{R}_{\nu}f\right)_{\alpha\beta}(y)=\begin{cases}-f_{\alpha\beta}(r_{\nu}y)&\text{if }\nu\in\{\alpha,\beta\},\\ f_{\alpha\beta}(r_{\nu}y)&\text{otherwise}\end{cases}\qquad(1\le\alpha<\beta\le3).
\]
For $c\in\mathbb{R}$ the map $\mathsf{R}_{\nu,c}:=\tau_{c\cdot e_{\nu}}\mathsf{R}_{\nu}\tau_{-c\cdot e_{\nu}}$ is the reflection of tests in the plane $\{x_{\nu}=c\}$. The group of the eight diagonal matrices with diagonal entries $\pm1$ is denoted $N_3$; it is a subgroup of $O_{\mathrm{cub}}$. The maps $\tau_b$, $U_r$, $\mathsf{R}_{\nu}$ and $\mathsf{R}_{\nu,c}$ are continuous linear bijections of $\mathcal{S}$ onto itself, and they act on $\mathcal{S}'$ by transposition (Section~\ref{sec:test-forms-test}):
\begin{align*}
(\tau_b^{*}F)(f)&:=F(\tau_bf),&(U_r^{*}F)(f)&:=F(U_rf),\\
(\mathsf{R}_{\nu}^{*}F)(f)&:=F(\mathsf{R}_{\nu}f),&(\mathsf{R}_{\nu,c}^{*}F)(f)&:=F(\mathsf{R}_{\nu,c}f)\qquad\left(F\in\mathcal{S}',\ f\in\mathcal{S}\right).
\end{align*}
The maps $\tau_b^{*}$, $U_r^{*}$, $\mathsf{R}_{\nu}^{*}$, $\mathsf{R}_{\nu,c}^{*}$ and $F\mapsto-F$ are bijections of $\mathcal{S}'$ onto itself which, together with their inverses, are measurable for the cylinder $\sigma$-algebra (Section~\ref{sec:euclidean-maps-actions}). A law $\mu$ is invariant under one of them if $\mu$ is equal to its image under this map, that is, if the inverse image of every cylinder set $A$ under the map has the measure $\mu(A)$. The law $\mu$ is said to be invariant under the translations of $\mathbb{R}^3$, under the translations along $e_j$, under $O_{\mathrm{cub}}$ or under $N_3$ if it is invariant under $\tau_b^{*}$ for every $b\in\mathbb{R}^3$, under $\tau_{t\cdot e_j}^{*}$ for every $t\in\mathbb{R}$, under $U_r^{*}$ for every $r\in O_{\mathrm{cub}}$, or under $U_r^{*}$ for every $r\in N_3$, respectively. For $r\in O_{\mathrm{cub}}$ and a law $\rho$, the image of $\rho$ under $U_{r^{-1}}^{*}$, that is, under the map $F\mapsto F\circ U_{r^{-1}}$, is denoted $r_{*}\rho$:
\[
\left(r_{*}\rho\right)(A):=\rho\left(\left\{F\in\mathcal{S}':\ F\circ U_{r^{-1}}\in A\right\}\right)\qquad\text{for every cylinder set }A
\]
(Section~\ref{sec:euclidean-maps-actions}).

\paragraph{Gaussian bound, invariance, reflection positivity.}
A law $\mu$ satisfies the Gaussian bound~\eqref{eq:setup-G} if
\begin{equation}\tag{G}\label{eq:setup-G}
\int\mathrm{d}\mu(F)\,\exp F(f)\le\exp\left(\frac{1}{2}\,\|f\|_2^2\right)\qquad\text{for every real }f\in\mathcal{S}.
\end{equation}
Let $j\in\{1,2,3\}$. The law $\mu$ satisfies~\eqref{eq:setup-RPj} if
\begin{equation}\tag{$RP_j$}\label{eq:setup-RPj}
\sum_{k,l=1}^K\overline{z_k}\cdot z_l\cdot\int\mathrm{d}\mu(F)\,\exp\left(i\left(F(f_l)-F\left(\mathsf{R}_jf_k\right)\right)\right)\ge0
\end{equation}
for every integer $K\ge1$, all $z_1,\dots,z_K\in\mathbb{C}$ and all real $f_1,\dots,f_K\in C_c^{\infty}(\mathbb{R}^3;\Lambda^2)$ with $\operatorname{supp}f_k\subset\{x_j>0\}$ for $k=1,\dots,K$. For $c\in\mathbb{R}$ the law $\mu$ is reflection positive at the plane $\{x_j=c\}$ if the inequality~\eqref{eq:setup-RPj} holds with $\mathsf{R}_{j,c}$ in place of $\mathsf{R}_j$ for every integer $K\ge1$, all $z_1,\dots,z_K\in\mathbb{C}$ and all real $f_1,\dots,f_K\in C_c^{\infty}(\mathbb{R}^3;\Lambda^2)$ with $\operatorname{supp}f_k\subset\{x_j>c\}$ for $k=1,\dots,K$ (Section~\ref{sec:properties-laws}); since $\mathsf{R}_{j,0}=\mathsf{R}_j$, for $c=0$ this is~\eqref{eq:setup-RPj}. The law $\mu$ satisfies the hypothesis~\eqref{eq:setup-H} if
\begin{multline}\tag{H}\label{eq:setup-H}
\text{$\mu$ satisfies~\eqref{eq:setup-G}, is invariant under the translations of $\mathbb{R}^3$}\\
\text{and under $O_{\mathrm{cub}}$, and satisfies~\eqref{eq:setup-RPj} for $j=1,2,3$.}
\end{multline}

\paragraph{Moments and cumulants.}
Let $\mu$ be a law under which $F(f)$ has finite moments of all orders for every real $f\in\mathcal{S}$; by Lemma~\ref{lem:moment-bounds} this is the case if $\mu$ satisfies~\eqref{eq:setup-G}. For an integer $n\ge0$ and $f_1,\dots,f_n\in\mathcal{S}$ put
\[
S_n(f_1,\dots,f_n):=\int\mathrm{d}\mu(F)\,F(f_1)\cdots F(f_n),
\]
so that $S_0=1$. For $n\ge1$ the joint cumulant of the random variables $F(f_1),\dots,F(f_n)$ under $\mu$ is
\[
\kappa_n(f_1,\dots,f_n):=\sum_{\pi}(-1)^{\#|\pi|-1}\cdot\left(\#|\pi|-1\right)!\cdot\prod_{B\in\pi}S_{\#|B|}\left((f_i)_{i\in B}\right),
\]
where $\pi$ ranges over the partitions of $\{1,\dots,n\}$ into non-empty blocks $B$; and $\kappa_n(f):=\kappa_n(f,\dots,f)$. For a random variable $X$ with finite moments of all orders under an expectation $E$, its cumulant $\kappa_m(X;E)$ of order $m$ is the joint cumulant, formed with $E$, of $m$ copies of $X$. When the law with which a moment or a cumulant is formed has to be named, it is written as an upper index, as in $S_2^{\mu}$ and $\kappa_4^{\mu}$. If $\mu$ is invariant under $F\mapsto-F$, its moments of odd order vanish, and
\[
\kappa_4(f)=S_4(f,f,f,f)-3\,S_2(f,f)^2\qquad(f\in\mathcal{S}).
\]
Written without its lower index, $E^{\mathrm{nat}}$, $\mathrm{nat}\in\{\mathrm{af},\mathrm{av}\}$, is the expectation $E^{\mathrm{nat}}_n$ of Section~\ref{sec:lattice-laws} with the parameters of Section~\ref{sec:curvature-field-limit} at the cutoff $(N,T)$ under consideration, that is, with $n=h\cdot L^{T+N}$, $q=e\cdot\eta^{\frac{1}{2}}$ and $\boldsymbol{\mu}=\boldsymbol{\mu}_N$.

\paragraph{Free covariance and closedness.}
For a two-form $g=(g_{12},g_{13},g_{23})$ on $\mathbb{R}^3$ let $\ast g:=\left(g_{23},-g_{13},g_{12}\right)$ be its dual vector field. Let $\langle\cdot,\cdot\rangle_{L^2}$ be the scalar product of $L^2(\mathbb{R}^3;\Lambda^2)$, whose norm is $\|\cdot\|_2$, and let $\Pi_c$ be the orthogonal projection of $L^2(\mathbb{R}^3;\Lambda^2)$ such that, for every $g$ in this space,
\[
\widehat{\ast\Pi_cg}(p)=\left(1-\frac{p\,p^T}{|p|_2^2}\right)\widehat{\ast g}(p)\qquad\text{for almost every }p\in\mathbb{R}^3,
\]
the Fourier transform of a vector field being taken componentwise. The free Maxwell covariance is
\[
C_0(f,g):=\left\langle f,\Pi_cg\right\rangle_{L^2}\qquad\left(f,g\in L^2(\mathbb{R}^3;\Lambda^2)\right)
\]
(Section~\ref{sec:properties-limit-points}). For $g\in\mathcal{S}(\mathbb{R}^3)$ let $d^{*}g\in\mathcal{S}$ be the two-form whose dual vector field is $-\nabla g$, that is, $\ast d^{*}g=-\nabla g$. For $F\in\mathcal{S}'$ let $dF$ be the linear form $g\mapsto F(d^{*}g)$ on $\mathcal{S}(\mathbb{R}^3)$; the distribution $F$ is closed if $dF=0$. The set of the closed $F\in\mathcal{S}'$ is a cylinder set (Section~\ref{sec:closed-distributions}).

\paragraph{Tensor tests.}
For an integer $n\ge1$ let $\mathscr{S}_n:=\mathcal{S}\left(\mathbb{R}^{3n};(\Lambda^2)^{\otimes n}\right)$ be the space of the functions on $\mathbb{R}^{3n}=(\mathbb{R}^3)^n$ with values in the complexification of $(\Lambda^2)^{\otimes n}$ whose components, with respect to the basis formed by the $n$-fold tensor products of the vectors $\mathrm{d}x_1\wedge\mathrm{d}x_2$, $\mathrm{d}x_1\wedge\mathrm{d}x_3$, $\mathrm{d}x_2\wedge\mathrm{d}x_3$, are complex Schwartz functions; and let $\mathscr{S}_0:=\mathbb{C}$. The space $\mathscr{S}_n$ carries the seminorms
\[
q_m(\Psi):=\max_{|\alpha|\le m}\ \sup_{x\in\mathbb{R}^{3n}}\left(1+|x|_2^2\right)^{\frac{m}{2}}\cdot\left|\partial^{\alpha}\Psi(x)\right|\qquad\left(\Psi\in\mathscr{S}_n,\ m\ge0\text{ an integer}\right),
\]
where $\alpha$ ranges over the multi-indices in $\{0,1,2,\dots\}^{3n}$, $|\alpha|$ is the order $\alpha_1+\dots+\alpha_{3n}$ of $\alpha$, $|x|_2$ is the Euclidean norm of $x\in\mathbb{R}^{3n}$, and $\left|\partial^{\alpha}\Psi(x)\right|$ is the largest of the moduli of the components of $\partial^{\alpha}\Psi(x)$. For $f_1,\dots,f_n\in\mathcal{S}$ let $f_1\otimes\dots\otimes f_n\in\mathscr{S}_n$ be the function $(y_1,\dots,y_n)\mapsto f_1(y_1)\otimes\dots\otimes f_n(y_n)$; for integers $k,l\ge1$, $\Psi\in\mathscr{S}_k$ and $\Phi\in\mathscr{S}_l$ let $\Psi\otimes\Phi\in\mathscr{S}_{k+l}$ be the function
\[
(y_1,\dots,y_{k+l})\mapsto\Psi(y_1,\dots,y_k)\otimes\Phi(y_{k+1},\dots,y_{k+l});
\]
and a tensor product one of whose factors is a number $\Psi_0\in\mathscr{S}_0$ is the product of the other factor with $\Psi_0$. For a direction $\nu_0\in\{1,2,3\}$ and $\Psi\in\mathscr{S}_n$ let
\[
\left(\mathsf{R}_{\nu_0}^{\otimes n}\Psi\right)(y_1,\dots,y_n):=\left(\wedge^2r_{\nu_0}\right)^{\otimes n}\Psi\left(r_{\nu_0}y_1,\dots,r_{\nu_0}y_n\right),\qquad\Theta_n\Psi:=\overline{\mathsf{R}_{\nu_0}^{\otimes n}\Psi},
\]
the bar denoting complex conjugation of the components; and let $\Theta_0\Psi_0:=\overline{\Psi_0}$ for $\Psi_0\in\mathscr{S}_0$ (Section~\ref{sec:moments-tensor-tests}). The direction $\nu_0$ is not shown in the notation $\Theta_n$.

\paragraph{Reflections.}
For $\nu\in\{1,2,3\}$ and $c\in\mathbb{R}$ let $\Sigma^{+}_{\nu,c}$ be the $\sigma$-algebra on $\mathcal{S}'$ generated by the functions $F\mapsto F(f)$ with $f\in C_c^{\infty}(\mathbb{R}^3;\Lambda^2)$ real and $\operatorname{supp}f\subset\{x_{\nu}>c\}$. For a function $A\colon\mathcal{S}'\to\mathbb{C}$ let $\Theta_{\nu,c}A\colon\mathcal{S}'\to\mathbb{C}$ be the function
\[
\left(\Theta_{\nu,c}A\right)(F):=\overline{A\left(\mathsf{R}_{\nu,c}^{*}F\right)}\qquad\left(F\in\mathcal{S}'\right).
\]

\paragraph{Gaussian laws and mixtures.}
A law $\mu$ is Gaussian if, for every real $f\in\mathcal{S}$, the random variable $F\mapsto F(f)$ is normal under $\mu$; degenerate normal variables, of variance zero, are admitted. Every finite family $\left(F(f_1),\dots,F(f_n)\right)$ with real $f_1,\dots,f_n\in\mathcal{S}$ is then jointly normal, each real linear combination $\sum_it_i\cdot F(f_i)=F\left(\sum_it_i\cdot f_i\right)$ being normal; in particular a Gaussian law has finite moments of all orders. A law $\mu$ under which every $F(f)$ is integrable is centred if $\int\mathrm{d}\mu(F)\,F(f)=0$ for every real $f\in\mathcal{S}$. By Wick's rule all fourth cumulants $\kappa_4(f_1,f_2,f_3,f_4)$ of a centred Gaussian law vanish (Section~\ref{sec:cumulants-gaussian-laws}), and a law that satisfies~\eqref{eq:setup-G} is centred (Lemma~\ref{lem:moment-bounds}(a)). Hence $\kappa_4=0$ for every Gaussian law that satisfies~\eqref{eq:setup-G}.

Let a property of laws on $\mathcal{S}'$ be given. A law $\mu$ is a mixture of laws with this property if there are a probability space $(\boldsymbol{\Lambda},m_{\mathrm{x}})$ and laws $\boldsymbol{\gamma}_{\mathfrak{l}}$ on $\mathcal{S}'$, $\mathfrak{l}\in\boldsymbol{\Lambda}$, each of which has the property, such that for every cylinder set $A$ the function $\mathfrak{l}\mapsto\boldsymbol{\gamma}_{\mathfrak{l}}(A)$ is measurable and
\[
\mu(A)=\int\mathrm{d}m_{\mathrm{x}}(\mathfrak{l})\,\boldsymbol{\gamma}_{\mathfrak{l}}(A).
\]

\paragraph{The combination $\mathrm{Ext}$.}
Let $\mu$ be a law under which $F(f)$ has finite moments of all orders for every real $f\in\mathcal{S}$. For $f\in\mathcal{S}$ and $x\in\mathbb{R}^3$ put
\[
\mathrm{Ext}_x(f):=\kappa_4(f)-3\,\kappa_4(f,f,\tau_xf,\tau_xf),
\]
written $\mathrm{Ext}^{\mu}_x(f)$ when the law has to be named. If $\mu$ is invariant under the translations of $\mathbb{R}^3$, then $\kappa_4(\tau_xf)=\kappa_4(f)$, and since $\kappa_4$ is symmetric and $4$-linear,
\[
\mathrm{Ext}_x(f)=2\,\kappa_4(f)-\frac{1}{4}\left[\kappa_4(f+\tau_xf)+\kappa_4(f-\tau_xf)\right]
\]
(the identity~\eqref{eq:fourth-cumulants-1} of Section~\ref{sec:cumulants-gaussian-laws}). For a centred Gaussian law one has $\mathrm{Ext}_x(f)=0$ for all $f\in\mathcal{S}$ and $x\in\mathbb{R}^3$, all its fourth cumulants being zero (Section~\ref{sec:cumulants-gaussian-laws}). Let $\mu$ be a mixture of translation-invariant Gaussian laws that satisfy~\eqref{eq:setup-G}. Then $\mu$ is invariant under the translations of $\mathbb{R}^3$ and satisfies~\eqref{eq:setup-G}, because the identity that defines a mixture extends from cylinder sets to non-negative measurable functions on $\mathcal{S}'$; and
\[
\mathrm{Ext}^{\mu}_x(f)\le0\qquad\text{for every real }f\in\mathcal{S}\text{ and every }x\in\mathbb{R}^3
\]
(Theorem~\ref{thm:non-gaussian-components}, last assertion; Step~5 of Section~\ref{sec:proof-part-ii-2}).

\paragraph{Ergodic decomposition.}
Let $\Gamma$ be one of the groups $\mathbb{R}^3$ and $\mathbb{R}\cdot e_j$, $j\in\{1,2,3\}$, and let $\rho$ be a law that is invariant under $\tau_b^{*}$ for every $b\in\Gamma$. The law $\rho$ is ergodic under $\Gamma$ if the constants are the only elements $\Phi$ of $L^2(\rho)$ with $\Phi\circ\tau_b^{*}=\Phi$ for every $b\in\Gamma$ (Section~\ref{sec:invariant-sets-invariant}); for $\Gamma=\mathbb{R}^3$ and for $\Gamma=\mathbb{R}\cdot e_j$ the law is called ergodic under $\mathbb{R}^3$ and ergodic under the translations along $e_j$, respectively.

Let $\mu$ be a law that is invariant under the translations of $\mathbb{R}^3$, and let $\mathcal{I}$ be the $\sigma$-algebra of the cylinder sets $A$ with $\tau_b^{*}(A)=A$ for every $b\in\mathbb{R}^3$ (Section~\ref{sec:invariant-sets-invariant}). On $L^2(\mu)$ the conditional expectation given $\mathcal{I}$, like the conditional expectation given the $\sigma$-algebra of Lemma~\ref{lem:averages-translations-2}, which consists of the cylinder sets $A$ with $\mu\left((\tau_b^{*})^{-1}(A)\,\triangle\,A\right)=0$ for every $b\in\mathbb{R}^3$, is the orthogonal projection onto the subspace of the $\Phi\in L^2(\mu)$ with $\Phi\circ\tau_b^{*}=\Phi$ for every $b\in\mathbb{R}^3$ (Lemma~\ref{lem:averages-translations}(iv), Lemma~\ref{lem:averages-translations-2}(b)).

Let $\mathcal{M}$ be the set of the laws on $\mathcal{S}'$, with the $\sigma$-algebra generated by the maps $\rho\mapsto\rho(A)$, $A$ a cylinder set. The laws that are invariant and ergodic under $\mathbb{R}^3$ form a measurable subset of $\mathcal{M}$, and there is exactly one probability measure $\nu$ on $\mathcal{M}$ that is carried by this subset and satisfies
\[
\mu(A)=\int\mathrm{d}\nu(\xi)\,\mu_{\xi}(A)\qquad\text{for every cylinder set }A,\quad\text{where }\mu_{\xi}:=\xi .
\]
Moreover $\nu$ is the image of $\mu$ under the map $\omega\mapsto\mu_{\omega}$, where $(\mu_{\omega})_{\omega\in\mathcal{S}'}$ is a regular version of the conditional law of $\mu$ given $\mathcal{I}$: every $\mu_{\omega}$ is a law on $\mathcal{S}'$, and for every cylinder set $A$ the function $\omega\mapsto\mu_{\omega}(A)$ is $\mathcal{I}$-measurable and equal to $E_{\mu}\left[1_A\mid\mathcal{I}\right](\omega)$ for $\mu$-almost every $\omega\in\mathcal{S}'$ (Theorem~\ref{thm:ergodic-decomposition} with $J=\{1,2,3\}$; Definition~\ref{def:ergodicity-ergodic-decompositions}, Proposition~\ref{prop:existence-conditional-expectations}). One writes $\mu=\int\mathrm{d}\nu(\xi)\,\mu_{\xi}$, and the laws $\mu_{\xi}$ are the ergodic components of $\mu$. When $\mu_{\xi}$ has finite moments of all orders, $S^{\xi}_n$, $\kappa_4^{\xi}$, $\mathrm{Ext}^{\xi}_x$ and $C_{\xi}$ denote the moments, the fourth cumulant, the combination $\mathrm{Ext}_x$ and the covariance $C_{\xi}(f,g):=\kappa_2^{\mu_{\xi}}(f,g)$ formed with $\mu_{\xi}$.

In the same way, let $j\in\{1,2,3\}$, and let $\mu$ be invariant under the translations along $e_j$. By Theorem~\ref{thm:ergodic-decomposition} with $J=\{j\}$ there is exactly one probability measure $\nu^j$ on $\mathcal{M}$ that is carried by the measurable set of the laws which are invariant and ergodic under the translations along $e_j$ and that satisfies $\mu(A)=\int\mathrm{d}\nu^j(\rho)\,\rho(A)$ for every cylinder set $A$; this is the decomposition $\mu=\int\mathrm{d}\nu^j(\rho)\,\rho$ of $\mu$ into laws that are invariant and ergodic under the translations along $e_j$. For a law $\rho$, $\kappa_4^{\rho}$ is its fourth cumulant, and $\rho(A):=\int\mathrm{d}\rho\,A$ for every $\rho$-integrable function $A$ on $\mathcal{S}'$.

\paragraph{Reconstruction.}
Let $j\in\{1,2,3\}$, and let $\rho$ be a law that is invariant under the translations along $e_j$ and under $\mathsf{R}_j^{*}$ and that satisfies~\eqref{eq:setup-RPj}. The reconstruction of $\rho$ in the direction $j$ consists of the following objects. The form
\[
(A,B)\mapsto\int\mathrm{d}\rho\,(\Theta_{j,0}A)\cdot B
\]
is Hermitian and positive semidefinite on $L^2(\mathcal{S}',\Sigma^{+}_{j,0},\rho)$ (Proposition~\ref{prop:reconstruction-direction}(i)). Let $\mathscr{H}_j(\rho)$ be the Hilbert space obtained by completing the quotient of $L^2(\mathcal{S}',\Sigma^{+}_{j,0},\rho)$ by the null space of this form, with the scalar product induced by the form, and let $[A]\in\mathscr{H}_j(\rho)$ be the class of $A$. There is exactly one self-adjoint operator $\mathsf{H}_j=\mathsf{H}_j(\rho)\ge0$ on $\mathscr{H}_j(\rho)$ such that
\[
\exp\left(-t\cdot\mathsf{H}_j\right)[A]=\left[A\circ\tau_{t\cdot e_j}^{*}\right]\qquad\left(t\ge0,\ A\in L^2(\mathcal{S}',\Sigma^{+}_{j,0},\rho)\right),
\]
and one puts $\Omega_j:=[1]$. If $\rho$ is invariant under the translations of $\mathbb{R}^3$, then for every $y\in\mathbb{R}^3$ orthogonal to $e_j$ there is exactly one unitary operator $V_j(y)$ on $\mathscr{H}_j(\rho)$ with $V_j(y)[A]=\left[A\circ\tau_y^{*}\right]$ for every $A\in L^2(\mathcal{S}',\Sigma^{+}_{j,0},\rho)$; one has $V_j(y)\,V_j(y')=V_j(y+y')$, so that the $V_j(y)$ form commuting unitary groups, and each $V_j(y)$ commutes with each $\exp\left(-t\cdot\mathsf{H}_j\right)$ (Corollary~\ref{cor:transfer-operators-coordinate}; Proposition~\ref{prop:reconstruction-direction}).

\subsection{The free coefficient}\label{sec:free-coefficient}

The mass parameter $\bar{m}\in(0,1]$ is that of Section~\ref{sec:lattice-tori}, and $\gamma_1,\gamma_2,\gamma_3$ are the matrices of Section~\ref{sec:spinors-wilson-determinant}.

\paragraph{The coefficient on $\mathbb{R}^3$.}
Let $\mathsf{M}>0$, write $\gamma\cdot\partial:=\sum_{\mu=1}^3\gamma_{\mu}\partial_{\mu}$, and let
\[
S_{\mathsf{M}}:=\left(\mathsf{M}+\gamma\cdot\partial\right)^{-1},
\]
an operator on $L^2(\mathbb{R}^3;\mathbb{C}^4)$. In terms of the Fourier transform $\hat{\psi}(p)=\int_{\mathbb{R}^3}\mathrm{d}x\,\exp(-ip\cdot x)\,\psi(x)$ it is the Fourier multiplier $\left(\mathsf{M}+i\gamma\cdot p\right)^{-1}$, that is, $\widehat{S_{\mathsf{M}}\psi}(p)=\left(\mathsf{M}+i\gamma\cdot p\right)^{-1}\hat{\psi}(p)$; and it is the convolution with the matrix kernel
\[
z\mapsto\left(\mathsf{M}-\gamma\cdot\partial\right)\frac{\exp\left(-\mathsf{M}\cdot|z|_2\right)}{4\pi\,|z|_2}\qquad\left(z\in\mathbb{R}^3\setminus\{0\}\right).
\]
Let $a\in C_c^{\infty}(\mathbb{R}^3;\mathbb{R}^3)$ be real, and let the matrix function $i\gamma\cdot a=i\sum_{\mu}a_{\mu}\gamma_{\mu}$ act on $L^2(\mathbb{R}^3;\mathbb{C}^4)$ by multiplication. The operator $S_{\mathsf{M}}\left(i\gamma\cdot a\right)S_{\mathsf{M}}$ is a Hilbert--Schmidt operator (Lemma~\ref{lem:hilbert-schmidt-bound}); so is $S_{\mathsf{M}}\left(i\gamma\cdot a\right)S_{\mathsf{M}}\left(i\gamma\cdot a\right)$, the multiplication by $i\gamma\cdot a$ being bounded. It follows that the square
\[
\left(S_{\mathsf{M}}\,i\gamma\cdot a\right)^4=\left[S_{\mathsf{M}}\left(i\gamma\cdot a\right)S_{\mathsf{M}}\left(i\gamma\cdot a\right)\right]^2
\]
of this operator is of trace class (Section~\ref{sec:free-fourth-order}), and one puts
\[
q_{4,\mathsf{M}}(a):=-\frac{1}{4}\operatorname{Tr}_{L^2(\mathbb{R}^3;\mathbb{C}^4)}\left[\left(S_{\mathsf{M}}\,i\gamma\cdot a\right)^4\right].
\]

A \emph{probe} is a real two-form $\beta\in C_c^{\infty}(\mathbb{R}^3;\Lambda^2)$ with $\beta\ne0$, whose components are extended to all pairs of directions by $\beta_{\nu\mu}:=-\beta_{\mu\nu}$. To a probe $\beta$ belong the one-form $a\in C_c^{\infty}(\mathbb{R}^3;\mathbb{R}^3)$, the real two-form $f\in C_c^{\infty}(\mathbb{R}^3;\Lambda^2)$ and the number $b(f)$ defined by
\[
a:=d^{*}\beta,\quad a_{\nu}=-\sum_{\mu=1}^3\partial_{\mu}\beta_{\mu\nu};\qquad f:=dd^{*}\beta,\quad f_{\mu\nu}=\partial_{\mu}a_{\nu}-\partial_{\nu}a_{\mu};
\]
\[
b(f):=24\,q_{4,\bar{m}}(a)=-6\operatorname{Tr}\left[\left(S_{\bar{m}}\,i\gamma\cdot a\right)^4\right];
\]
the number $b(f)$ is the \emph{free coefficient} of $f$. The two-form $f$ determines $a$: if $\beta_1$ is a probe with $dd^{*}\beta_1=f$, the one-form $a-d^{*}\beta_1$ is smooth, compactly supported, closed and coclosed, hence harmonic, hence zero. Thus $b(f)$ depends on $f$ alone, and it is a real number (Section~\ref{sec:free-fourth-order}, Lemma~\ref{lem:hilbert-schmidt-bound}).

With $K:=S_{\bar{m}}\,i\gamma\cdot a$, so that $\bar{m}+\gamma\cdot\left(\partial+ita\right)=\left(\bar{m}+\gamma\cdot\partial\right)\left(1+tK\right)$ for $t\in\mathbb{R}$, the number $b(f)$ (the one-loop light-by-light coefficient of the free Dirac field of mass $\bar{m}$ in the external potential $a$) is $4!$ times the coefficient $-\frac{1}{4}\operatorname{Tr}K^4$ of $t^4$ in the series
\[
\sum_{n\ge1}\frac{(-1)^{n+1}}{n}\,t^n\operatorname{Tr}K^n
\]
of $\log\det\left(1+tK\right)$; of this series only the term $n=4$ is used, and $K^4=\left(S_{\bar{m}}\,i\gamma\cdot a\right)^4$ is of trace class, as shown above.

\paragraph{The coefficient on a lattice torus.}
Let $\eta\in(0,1]$ and $R\ge1$ be such that $n:=R/\eta$ is an integer $\ge2$, and let $\Lambda_{\eta,R}:=\eta\mathbb{Z}^3/R\mathbb{Z}^3$. The cells of $\Lambda_{\eta,R}$, the spaces $C^k(\Lambda_{\eta,R})$ of its real $k$-cochains, the pairing $\langle\cdot,\cdot\rangle_{\eta}$, the coboundary $d$, the operator $d_{\eta}=\eta^{-1}d$ and its adjoint $d_{\eta}^{*}$ are defined by the formulas of Sections~\ref{sec:lattice-tori} and~\ref{sec:curvature-field-limit}. The map $x\mapsto\eta x$ is a bijection of $\Lambda_n$ onto $\Lambda_{\eta,R}$ which carries cells to cells; through it a cochain on $\Lambda_{\eta,R}$ is read in lattice units, as a cochain on $\Lambda_n$. In Theorem~\ref{thm:main-theorem} these objects occur for $\eta=L^{-N}$ and $R=R^{*}(e)$, and then $\Lambda_{\eta,R}$ is one of the tori of Section~\ref{sec:lattice-tori} (Section~\ref{sec:parameters-scales}).

For a smooth function $g$ on $\mathbb{R}^3$ put
\[
J_{\eta}g(x):=\eta^{-3}\int_{\left[-\frac{\eta}{2},\frac{\eta}{2}\right]^3}\mathrm{d}y\,g(x+y)\qquad\left(x\in\mathbb{R}^3\right).
\]
Let $\beta$ be a probe whose support lies in a translate of the cube $[0,R)^3$, that is, in a period cube of the torus of side $R$, and let $\beta^R:=\sum_{k\in\mathbb{Z}^3}\beta(\cdot+Rk)$ be its periodisation with the period $R$ in every direction. The real two-cochain $u$ and the real one-cochain $a_{\eta}$ on $\Lambda_{\eta,R}$ are
\[
u_{\mu\nu}(x):=\left(J_{\eta}\beta^R_{\mu\nu}\right)\left(x+\frac{\eta}{2}\left(e_{\mu}+e_{\nu}\right)\right)\quad\left(x\in\Lambda_{\eta,R},\ \mu<\nu\right),\qquad a_{\eta}:=d_{\eta}^{*}u\in C^1(\Lambda_{\eta,R}),
\]
where $u_{\mu\nu}(x)$ is the value of $u$ at the plaquette $(x;\mu\nu)$. Let $D[\cdot]$ be the operator of Section~\ref{sec:spinors-wilson-determinant} on the spinors of the cubic torus $\Lambda_n$, which has $n=R/\eta$ sites per side, and read $a_{\eta}$ as a one-cochain on $\Lambda_n$. Put
\[
N_0(t):=\frac{\det\left(D[t\eta a_{\eta}]+\bar{m}\cdot\eta\right)}{\det\left(D[0]+\bar{m}\cdot\eta\right)}\quad\left(t\in\mathbb{C}\right),\qquad b_{\eta,R}(0)[\beta]:=\left.\frac{\mathrm{d}^4}{\mathrm{d}t^4}\log N_0(t)\right|_{t=0}.
\]
As $\bar{m}\cdot\eta>0$, the denominator is positive, and $N_0$ is an entire function with $N_0(0)=1$ (Theorem~\ref{thm:free-coefficient-1}(d1), (d2)); so $\log N_0$, with the principal branch, is analytic in a neighbourhood of $t=0$. The function $N_0$ is the \emph{free ratio}, and $b_{\eta,R}(0)[\beta]$ is the \emph{free coefficient on the lattice torus}, written $b_{\eta,R}(0)$ when the probe is understood. When $\eta\to0$ and $R\to\infty$ in any manner, no relation between $\eta$ and $R$ being imposed,
\[
b_{\eta,R}(0)[\beta]\to b(f)\qquad\left(f=dd^{*}\beta\right)
\]
(Theorem~\ref{thm:free-coefficient-1}(c)). With the charge $e$ as its argument, $N_0(e)$ is the threshold of Section~\ref{sec:parameters-scales}, and the threshold is always written with this argument (Definition~\ref{def:notation}(f)). Every other occurrence of the symbol $N_0$ in Sections~\ref{sec:setup}--\ref{sec:estimates} refers to the free ratio: $N_0(t)$, with the complex variable $t$; its value $N_0(0)$; its derivatives $N_0^{(m)}(0)$ (Section~\ref{sec:estimates}); and $N_0$ without an argument, which is the function itself, as in $\log N_0$.

\subsection{Parameters and scales}\label{sec:parameters-scales}

The quantities on which the statements below depend are fixed one after the other, in twelve \emph{stages}: for $n=1,\dots,12$, item~$n$ of the following list states what is fixed at stage~$n$. What is fixed at a stage depends only on what has been fixed at the earlier stages. The complete lists of the constants fixed at the several stages, and of the conditions imposed on them, are given in Definition~\ref{def:order-choice} and in Section~\ref{sec:conditions}.
\begin{enumerate}
\item Absolute constants, among them $\bar{\varepsilon}:=\frac{1}{200}$ and the number $\kappa_0$.
\item The H\"older order $\alpha\in\left(1-\bar{\varepsilon},1\right)$.
\item A number $\kappa\ge7\kappa_0+7$, the rate of decay of the polymer activities of Section~\ref{sec:spaces-activities}.
\item An odd integer $L$, the scaling factor of the block averaging, which satisfies the conditions that Definition~\ref{def:order-choice} imposes at this stage; these conditions involve only $\kappa$ and absolute constants.
\item Numbers that are functions of $\alpha$, $\kappa$, $L$ and of absolute constants. Among them are
\[
N^{*}:=8L^3,\qquad a_L:=\frac{L^3-1}{(L-1)^3},
\]
and the least integer $m_0\ge0$ with $L^{-m_0}\le\mathfrak{m}_0$, where $\mathfrak{m}_0$ is the right end point of the mass range $\left(0,\mathfrak{m}_0\right]$ of Theorem~\ref{thm:fermion-operators-weakened} (Definition~\ref{def:mass-bound-stop}). That theorem gives $\mathfrak{m}_0=1$; hence $m_0=0$. The letter $m_0$ is kept in the formulas.
\item A power $M=L^m$ of $L$, with an integer $m$, the side of the cubes that compose the polymers of Section~\ref{sec:lattice-geometry}; it satisfies the conditions that Definition~\ref{def:order-choice} states for this stage. These conditions involve only $\alpha$, $\kappa$, $L$ and absolute constants; they include the inequality $M\ge L^{m_0+2}$; and none of them involves the number $\varepsilon$ of stage~\ref{it:order-choice-8} or the charge $e$.
\item Further numbers, functions of the quantities fixed at the preceding stages.
\item Integers $r\ge2$, $p_0>3r+2$ and $p>2p_0$, the exponents of the logarithmic thresholds of Section~\ref{sec:running-quantities-tori}; a number $\varepsilon\in\left(1-\alpha,\bar{\varepsilon}\right]$; and a number $s\in(0,1)$ such that
\[
\sigma:=s\cdot\left(\frac{1}{4}-9\varepsilon\right)-12\varepsilon-\frac{2a_L}{N^{*}}>0.
\]
Numbers $s$ with this property exist (Section~\ref{sec:order-choice}). Put moreover
\[
\mathfrak{e}:=s\cdot\left(\frac{1}{4}-9\varepsilon\right)-12\varepsilon,\qquad\mathfrak{e}_2:=6\varepsilon,
\]
so that $\sigma=\mathfrak{e}-\frac{2a_L}{N^{*}}$.
\item A number $e_0>0$, the ceiling of the running charge of Section~\ref{sec:running-quantities-tori}, which satisfies the conditions laid down for this stage in Definition~\ref{def:order-choice}.
\item The number
\[
e_1\in\left(0,\min\left(e_0,\frac{1}{3}\cdot L^{-\frac{m_0}{2}}\right)\right]
\]
that Section~\ref{sec:constants-conditions-torus} defines by the inequalities~\eqref{eq:renormalisation-flow-e1}--\eqref{eq:renormalisation-flow-e8}.
\item The mass parameter $\bar{m}\in(0,1]$, which determines the scale $h$ of Section~\ref{sec:lattice-tori} and the numbers $h'$ and $e_G$ defined below; then the probe $\beta$ and the point $x_0$ of Section~\ref{sec:probes}; then the auxiliary numbers of the three probes $\beta_0$, $\beta_{+}$, $\beta_{-}$ of Section~\ref{sec:probes} and the further numbers that Definition~\ref{def:order-choice} names at this stage, in the order given there; then the constant $C_{\mathrm{kap}}$ and, last, the number $e_2$ of Theorem~\ref{thm:main-theorem}.
\item The charge $e\in\left(0,e_2\right]$, which determines the numbers $R^{*}(e)$, $T^{*}(e)$ and $N_0(e)$ defined below and the bare mass shifts $m_{0,N}(e)$, $N\ge0$, of Theorem~\ref{thm:main-theorem}; then the cutoff $(N,T)$.
\end{enumerate}
The quantities fixed at stages~\ref{it:order-choice-1}--\ref{it:order-choice-10} are called the \emph{parameters}; in Sections~\ref{sec:renormalisation-group-scheme}--\ref{sec:free-coefficient-chessboard} they are called the constants of these stages (Definition~\ref{def:constants}). Each parameter is fixed at its stage in such a way that the conditions that Definition~\ref{def:order-choice} imposes at that stage are satisfied; consequently the parameters satisfy the conditions of the first kind in Section~\ref{sec:conditions}. No parameter depends on $\bar{m}$, on the charge $e$, on $N$, on $T$ or on a probe, nor on any of the further quantities of the renormalisation group scheme that Definition~\ref{def:constants} lists (the level, the side of a torus, the twist, the running mass, the real pair and the external connection). The quantities fixed at stage~\ref{it:order-choice-11} may depend on the parameters, on $\bar{m}$ and on the probe $\beta$ and the point $x_0$ of Section~\ref{sec:probes}, and on nothing else. Accordingly, a constant or a threshold that is said to depend on certain quantities is a function of these quantities and of nothing else.

\paragraph{Scales.}
Let $0<\bar{m}\le1$, let $h$ be the scale of Section~\ref{sec:lattice-tori}, and let $e>0$. Put
\[
h':=h\cdot L^{-m_0},\qquad e_K:=e\cdot h'^{\frac{1}{2}},\qquad v_K:=-\log e_K;
\]
as $m_0=0$, one has $h'=h$. Let $r_K(e)$ be the least power of $L$ that is not less than $v_K^r$,
\[
r_K(e):=\min\left\{L^i:\ i\ge0\text{ an integer},\ L^i\ge v_K^r\right\},
\]
and put
\[
R_{\mathrm{top}}(e):=h'\cdot M\cdot r_K(e),\qquad R^{*}(e):=\frac{R_{\mathrm{top}}(e)}{L},\qquad T^{*}(e):=\log_L\left(\frac{R^{*}(e)}{h}\right).
\]
The threshold $N_0(e)$ is the least integer $N_1\ge0$ such that
\[
R^{*}(e)\cdot L^N\ge L\cdot M\cdot r_0\qquad\text{for every integer }N\ge N_1,
\]
where the number
\[
r_0:=\min\left\{L^i:\ i\ge0\text{ an integer},\ L^i\ge\left(-\log\left(e\cdot L^{-\frac{N}{2}}\right)\right)^r\right\}
\]
depends on $N$. This threshold is the number $N_{\mathrm{min}}\left(R^{*}(e),e\right)$ of Definition~\ref{def:tori-stop-level}. Finally let
\[
e_G=e_G(\bar{m}):=\sup\left\{e'>0:\ e'\cdot h'^{\frac{1}{2}}\le e_1,\ e'\cdot h'^{\frac{1}{2}}\cdot\log h'\le1\right\}.
\]
One has $0<e_G\le e_1\cdot h'^{-\frac{1}{2}}$, and both conditions in the definition of $e_G$ hold for every $e'\in\left(0,e_G\right]$ (Section~\ref{sec:constants-exponents-scales}).

\paragraph{Reference tori.}
If $e\cdot h'^{\frac{1}{2}}\le e_1$, then $T^{*}(e)$ is an integer with $T^{*}(e)\ge1$, and $R^{*}(e)=h\cdot L^{T^{*}(e)}$ (part~\ref{it:free-coefficient-3-1} of Theorem~\ref{thm:free-coefficient-3}). Then $R^{*}(e)$ and $L\cdot R^{*}(e)=R_{\mathrm{top}}(e)$ are sides of tori of Section~\ref{sec:lattice-tori}, and for $\eta=L^{-N}$ the two tori $\Lambda_{\eta,R^{*}(e)}$ and $\Lambda_{\eta,R_{\mathrm{top}}(e)}$ are called the \emph{reference tori}.

Let $e>0$ be such that $v_K\ge1$; this holds whenever $e\cdot h'^{\frac{1}{2}}\le e_1$ (Proposition~\ref{prop:consequences}). Then
\[
v_K^r\le r_K(e)<L\cdot v_K^r.
\]
Indeed, the first inequality holds by the definition of $r_K(e)$. If $r_K(e)=1$, then $r_K(e)<L\le L\cdot v_K^r$. If $r_K(e)\ge L$, then $L^{-1}\cdot r_K(e)$ is a power of $L$, so that $L^{-1}\cdot r_K(e)<v_K^r$ by the minimality of $r_K(e)$. Consequently
\[
h'\cdot M\cdot v_K^r\le R_{\mathrm{top}}(e)=L\cdot R^{*}(e)<L\cdot h'\cdot M\cdot v_K^r,
\]
where $h'=h$, $v_K=-\log\left(e\cdot h^{\frac{1}{2}}\right)$, and $\frac{1}{L\cdot\bar{m}}<h\le\frac{1}{\bar{m}}$ because $\frac{1}{L}<\bar{m}\cdot h\le1$ (Section~\ref{sec:lattice-tori}).

\subsection{Probes}\label{sec:probes}

Besides the parameters and the mass parameter, the data of Part~II of Theorem~\ref{thm:main-theorem} are a probe $\beta$ and a point $x_0\in\mathbb{R}^3$, subject to the following requirements.

Let $\bar{m}\in(0,1]$, and let $\beta$ be a probe (Section~\ref{sec:free-coefficient}) such that the two-form $f=dd^{*}\beta$ has a positive free coefficient:
\[
b(f)>0.
\]
Let $r_{\beta}$ be the radius of a ball that contains $\operatorname{supp}\beta$, let $a=d^{*}\beta$ be the one-form of Section~\ref{sec:free-coefficient}, and put
\[
\|a\|_{\infty}:=\sup_{x\in\mathbb{R}^3}|a(x)|_2.
\]
For every $\bar{m}\in(0,1]$ there are probes with $b(f)>0$ (Theorem~\ref{thm:free-coefficient-1}(a)).

Let $x_0\in\mathbb{R}^3$ be a point such that $d_0:=|x_0|_2-2r_{\beta}>0$ and
\begin{equation}\label{eq:setup-1}
\frac{6}{\bar{m}^2}\cdot\|a\|_{\infty}^4\cdot\left(\frac{4\pi\cdot r_{\beta}^3}{3}\right)^2\cdot\left(\bar{m}^2+\left(\bar{m}+\frac{1}{d_0}\right)^2\right)\cdot\frac{\exp\left(-2\bar{m}\cdot d_0\right)}{4\pi^2\cdot d_0^2}\le\frac{b(f)}{12}.
\end{equation}
Besides $d_0$, the inequality~\eqref{eq:setup-1} involves only $\bar{m}$, $r_{\beta}$, $\|a\|_{\infty}$ and $b(f)$. Its left side tends to $0$ as $d_0\to\infty$, and its right side is positive; hence~\eqref{eq:setup-1} holds for all sufficiently large $|x_0|_2$. It is the condition~\eqref{eq:fourth-cumulants-11} of Theorem~\ref{thm:non-gaussian-components-2}, with $m=\bar{m}$ and $r=r_{\beta}$.

Let $\tau_{x_0}$ be the translation of Section~\ref{sec:laws}, and put
\[
\beta_0:=\beta,\qquad\beta_{\pm}:=\beta\pm\tau_{x_0}\beta,\qquad g_0:=f,\qquad g_{\pm}:=f\pm\tau_{x_0}f.
\]
The supports of $\beta$ and of $\tau_{x_0}\beta$ lie in two balls of radius $r_{\beta}$ whose centres are at the distance $|x_0|_2>2r_{\beta}$; they are therefore disjoint, so that $\beta_{+}\ne0$ and $\beta_{-}\ne0$, and $\beta_0$, $\beta_{+}$, $\beta_{-}$ are probes. As $dd^{*}$ commutes with $\tau_{x_0}$, the three tests $g_0$, $g_{+}$, $g_{-}$ satisfy $g_0=dd^{*}\beta_0$ and $g_{\pm}=dd^{*}\beta_{\pm}$.

\section{Main theorem}\label{sec:main-theorem}

Throughout this section the mass parameter $\bar{m}\in(0,1]$ and the parameters, that is, the quantities of stages~\ref{it:order-choice-1}--\ref{it:order-choice-10} of Section~\ref{sec:parameters-scales}, are held fixed, and the notation is that of Section~\ref{sec:setup}.

\begin{theorem}[Main theorem]\label{thm:main-theorem}
\emph{Part I.} Let $e>0$, let $N\mapsto m_{0,N}$ be a mass assignment, and let $\mathrm{nat}\in\{\mathrm{af},\mathrm{av}\}$; for every cutoff $(N,T)$ let $P^{\mathrm{nat}}_{N,T}$ be the law of Section~\ref{sec:curvature-field-limit} formed with these data.
\begin{enumerate}
\item[(a)] The laws $P^{\mathrm{nat}}_{N,T}$ satisfy
\[
\sup_{(N,T)}\int\mathrm{d}P^{\mathrm{nat}}_{N,T}(F)\,\|F\|_{-3}^2\le61,
\]
where $(N,T)$ ranges over all cutoffs. The family $\left(P^{\mathrm{nat}}_{N,T}\right)$ is tight as a family of laws on $H_{-4}$, and also as a family of laws on $\mathcal{S}'$, for the strong as well as for the weak-$*$ topology of $\mathcal{S}'$. Every sequence $(N_i,T_i)_{i\ge1}$ of cutoffs has a subsequence along which $P^{\mathrm{nat}}_{N_i,T_i}$ converges weakly on $H_{-4}$.
\item[(b)] Each limit point $\mu$ of $\left(P^{\mathrm{nat}}_{N,T}\right)$ has the properties (P1)--(P8) below.
\begin{enumerate}
\item[(P1)]\itemmark{P1}\label{it:main-theorem-P1} For real $f\in\mathcal{S}$,
\[
\int\mathrm{d}\mu(F)\,\exp F(f)\le\exp\left(\frac{1}{2}C_0(f,f)\right)\le\exp\left(\frac{1}{2}\|f\|_2^2\right).
\]
The map $f\mapsto F(f)$ has an extension to a bounded linear map $X\colon L^2(\mathbb{R}^3;\Lambda^2)\to L^2(\mu)$, and for real $f\in L^2(\mathbb{R}^3;\Lambda^2)$ the two inequalities remain true with $X(f)$ in place of $F(f)$. For every real $f\in L^2(\mathbb{R}^3;\Lambda^2)$ and every $z\in\mathbb{C}$ one has
\[
\left|\int\mathrm{d}\mu\,\exp\left(z\cdot X(f)\right)\right|\le\exp\left(\frac{1}{2}(\operatorname{Re}z)^2\cdot\|f\|_2^2\right),
\]
and the integral on the left differs from $0$ whenever $|z|\cdot\|f\|_2<0.9005$. Let $f_1,\dots,f_n\in L^2(\mathbb{R}^3;\Lambda^2)$. On the set of the $z=(z_1,\dots,z_n)\in\mathbb{C}^n$ with $\sum_{i=1}^n|z_i|\cdot\|f_i\|_2\le\frac{2}{5}$ the principal branch of the logarithm defines an analytic function
\[
W(z):=\log\int\mathrm{d}\mu\,\exp\left(\sum_{i=1}^nz_i\cdot X(f_i)\right),
\]
and $|W(z)|\le\log2$ on this set.
\item[(P2)]\itemmark{P2}\label{it:main-theorem-P2} Let $n\ge1$ and let $f_1,\dots,f_n\in L^2(\mathbb{R}^3;\Lambda^2)$ be real. The moment $S_n(f_1,\dots,f_n)$ and the joint cumulant $\kappa_n(f_1,\dots,f_n)$ of $X(f_1),\dots,X(f_n)$ under $\mu$ satisfy
\begin{align*}
\left|S_n(f_1,\dots,f_n)\right|&\le2\left(2^n\cdot n!\right)^{\frac{1}{2}}\cdot\prod_{i=1}^n\|f_i\|_2,\\
\left|\kappa_n(f_1,\dots,f_n)\right|&\le(\log2)\cdot\left(\frac{5\exp(1)}{2}\right)^n\cdot n!\cdot\prod_{i=1}^n\|f_i\|_2.
\end{align*}
Moreover $S_n$ has exactly one extension to a continuous linear functional on $\mathscr{S}_n$, again written $S_n$, whose value at $f_1\otimes\dots\otimes f_n$ is $S_n(f_1,\dots,f_n)$ whenever $f_1,\dots,f_n\in\mathcal{S}$. This functional obeys the linear growth condition of Osterwalder and Schrader on the whole of $\mathscr{S}_n$, in the form
\[
\left|S_n(\Psi)\right|\le2\cdot66^n\cdot(n!)^{\frac{1}{2}}\cdot q_{6n}(\Psi)\qquad\left(\Psi\in\mathscr{S}_n\right).
\]
\item[(P3)]\itemmark{P3}\label{it:main-theorem-P3} One has $S_0=1$; each $S_n$ is symmetric in its arguments and real at real tests. The law $\mu$ is invariant under the map $F\mapsto-F$ of $\mathcal{S}'$, and therefore $S_n=0$ for odd $n$.
\item[(P4)]\itemmark{P4}\label{it:main-theorem-P4} Every probability measure on $\mathcal{S}'$, with its cylinder $\sigma$-algebra, whose moments on the tests in $C_c^{\infty}(\mathbb{R}^3;\Lambda^2)$ are the $S_n$ is equal to $\mu$. One has $\mu(H_{-t})=1$ for every $t>\frac{3}{2}$, and $\mu$ is carried by the closed distributions:
\[
\mu\left(\left\{F\in\mathcal{S}':dF=0\right\}\right)=1.
\]
\item[(P5)]\itemmark{P5}\label{it:main-theorem-P5} The law $\mu$ is invariant under every translation of $\mathbb{R}^3$, and the operators $A\mapsto A\circ\tau_b^{*}$, $b\in\mathbb{R}^3$, form a strongly continuous unitary group on $L^2(\mu)$. The law $\mu$ is invariant under $O_{\mathrm{cub}}$ as well, the cubic group acting on two-forms through $r\mapsto U_r$.
\item[(P6)]\itemmark{P6}\label{it:main-theorem-P6} Let $\nu_0\in\{1,2,3\}$ and $c\in\mathbb{R}$. For every $A\in L^2\left(\mathcal{S}',\Sigma^{+}_{\nu_0,c},\mu\right)$,
\[
\int\mathrm{d}\mu\,\left(\Theta_{\nu_0,c}A\right)\cdot A\ge0 .
\]
Let $\Theta_m$ be formed with the direction $\nu_0$, and let $(\Psi_n)_{n\ge0}$ be a sequence with $\Psi_n\in\mathscr{S}_n$, only finitely many of whose terms differ from $0$, such that every $\Psi_n$ is supported in the set of the $(y_1,\dots,y_n)\in(\mathbb{R}^3)^n$ with $(y_i)_{\nu_0}\ge0$ for $i=1,\dots,n$. Then
\[
\sum_{m,n\ge0}S_{m+n}\left(\Theta_m\Psi_m\otimes\Psi_n\right)\ge0 .
\]
\item[(P7)]\itemmark{P7}\label{it:main-theorem-P7} For each $\nu_0\in\{1,2,3\}$ the reconstruction of Section~\ref{sec:laws} in the direction $\nu_0$ applies to $\mu$. It furnishes the Hilbert space $\mathscr{H}_{\nu_0}(\mu)$; a self-adjoint operator $\mathsf{H}_{\nu_0}\ge0$ on it, with the semigroup $\left(\exp\left(-t\cdot\mathsf{H}_{\nu_0}\right)\right)_{t\ge0}$; the unit vector $\Omega_{\nu_0}$, for which $\mathsf{H}_{\nu_0}\Omega_{\nu_0}=0$; and the unitary groups of spatial translations $V_{\nu_0}(y)$, $y\perp e_{\nu_0}$, which commute with one another and with the semigroup.
\item[(P8)]\itemmark{P8}\label{it:main-theorem-P8} Let $(N_i,T_i)_{i\ge1}$ be a sequence of cutoffs with $N_i\to\infty$ and $T_i\to\infty$ along which $P^{\mathrm{nat}}_{N_i,T_i}\to\mu$ weakly on $H_{-4}$. Put $\eta_i:=L^{-N_i}$, $R_i:=h\cdot L^{T_i}$ and $X_i:=X_{\eta_i,R_i}$, and let $E_i$ be the expectation $E^{\mathrm{nat}}$ at the cutoff $(N_i,T_i)$. Since $X_i(f)=F^{\sharp}(f)$ and the law of $F^{\sharp}$ under $E_i$ is $P^{\mathrm{nat}}_{N_i,T_i}$, one has
\[
E_i\,\Phi\left(X_i(f)\right)=\int\mathrm{d}P^{\mathrm{nat}}_{N_i,T_i}(F)\,\Phi\left(F(f)\right)
\]
for every $f\in\mathcal{S}$ and every Borel map $\Phi\colon\mathbb{C}\to[0,\infty]$; for $f\in\mathcal{S}$ the identity holds as well for every Borel map $\Phi\colon\mathbb{C}\to\mathbb{C}$ for which one of its two sides converges absolutely. Let $f_1,\dots,f_k\in\mathcal{S}$. Then
\[
E_i\exp\left(\sum_{j=1}^kz_j\cdot X_i(f_j)\right)\to\int\mathrm{d}\mu(F)\,\exp\left(\sum_{j=1}^kz_j\cdot F(f_j)\right)
\]
as $i\to\infty$, uniformly in $z=(z_1,\dots,z_k)$ on every compact subset of $\mathbb{C}^k$; and all moments and all joint cumulants of $X_i(f_1),\dots,X_i(f_k)$ under $E_i$ converge to those of $F(f_1),\dots,F(f_k)$ under $\mu$.
\end{enumerate}
\item[(c)] Each limit point $\mu$ satisfies~\eqref{eq:setup-H}; its ergodic decomposition under $\mathbb{R}^3$ (Section~\ref{sec:laws}) is written $\mu=\int\mathrm{d}\nu(\xi)\,\mu_{\xi}$. For $\nu$-almost every $\xi$ the ergodic component $\mu_{\xi}$ is invariant and ergodic under the translations of $\mathbb{R}^3$ and invariant under $N_3$; it satisfies~\eqref{eq:setup-G}, with the same right side $\exp\left(\frac{1}{2}\|f\|_2^2\right)$ as $\mu$; one has $\mu_{\xi}\left(\left\{F\in\mathcal{S}':dF=0\right\}\right)=1$; the law $\mu_{\xi}$ is reflection positive at each plane $\{x_j=c\}$, $j\in\{1,2,3\}$, $c\in\mathbb{R}$; and it has the properties 1--4 below. These assertions hold together outside one $\nu$-null set: there is a set $\Xi_{\mu}$ with $\nu\left(\Xi_{\mu}\right)=1$ such that all of them hold for every $\xi\in\Xi_{\mu}$. The properties 5 and 6 below hold with the quantifiers over $\xi$ that are stated in them.
\begin{enumerate}
\item[1.] One has $S^{\xi}_1=0$, and, for $n\ge1$ and real $f_1,\dots,f_n\in\mathcal{S}$,
\[
\left|S^{\xi}_n(f_1,\dots,f_n)\right|\le n\cdot2^{\frac{n}{2}}\cdot\Gamma\left(\frac{n}{2}\right)\cdot\prod_{p=1}^n\|f_p\|_2,
\]
where $\Gamma$ is Euler's Gamma function.
\item[2.] The law $\mu_{\xi}$ is invariant under the translations of $\mathbb{R}^3$ and under $N_3$, and it is reflection positive at each plane $\{x_j=c\}$, $j\in\{1,2,3\}$, $c\in\mathbb{R}$.
\item[3.] Let $j\in\{1,2,3\}$. By property 2, the reconstruction of Section~\ref{sec:laws} in the direction $j$ applies to $\mu_{\xi}$; let $\mathscr{H}_j(\mu_{\xi})$, $\mathsf{H}_j$, $\Omega_j$ and $V_j(y)$, $y\perp e_j$, be its objects. The vectors of $\mathscr{H}_j(\mu_{\xi})$ that are invariant under $\exp\left(-t\cdot\mathsf{H}_j\right)$ for every $t\ge0$ and under $V_j(y)$ for every $y\perp e_j$ are exactly the multiples of $\Omega_j$, and $\dim\ker\mathsf{H}_j\in\{1,\infty\}$.
\item[4.] Let $n,n'\ge1$ and $f_1,\dots,f_n,g_1,\dots,g_{n'}\in\mathcal{S}$. As $a\in\mathbb{R}^3$ tends to infinity in such a way that $\min_l|a_l|\to\infty$,
\[
S^{\xi}_{n+n'}\left(f_1,\dots,f_n,\tau_ag_1,\dots,\tau_ag_{n'}\right)\to S^{\xi}_n(f_1,\dots,f_n)\cdot S^{\xi}_{n'}(g_1,\dots,g_{n'}),
\]
and the same product is the limit, as $t\to\infty$, of the means over cubes
\[
\frac{1}{t^3}\int_{[0,t]^3}\mathrm{d}a\,S^{\xi}_{n+n'}\left(f_1,\dots,f_n,\tau_ag_1,\dots,\tau_ag_{n'}\right).
\]
For $f,g\in\mathcal{S}$ one has $S^{\xi}_2(f,\tau_ag)\to0$ as $|a|\to\infty$ in any manner. The moments also converge along the coordinate axes and in the coordinate planes: if $K\subset\{1,2,3\}$ has one or two elements and $\left(a^{(k)}\right)_{k\ge1}$ is a sequence in $\mathbb{R}^3$ such that $\left|a^{(k)}_l\right|\to\infty$ for $l\in K$ and the sequence $\left(a^{(k)}_l\right)_{k\ge1}$ converges for each $l\notin K$, then $S^{\xi}_{n+n'}\left(f_1,\dots,f_n,\tau_{a^{(k)}}g_1,\dots,\tau_{a^{(k)}}g_{n'}\right)$ converges as $k\to\infty$, to the limit given by Theorem~\ref{thm:invariant-vectors-clustering}(i).
\item[5.] For each cylinder set $A$ with $\mu(A)=1$ one has $\mu_{\xi}(A)=1$ for $\nu$-almost every $\xi$, the $\nu$-null set of the exceptional $\xi$ depending on $A$.
\item[6.] Let $w\in O_{\mathrm{cub}}$. For every $\xi\in\Xi_{\mu}$ the law $w_{*}\mu_{\xi}$ (Section~\ref{sec:laws}) has the properties 1--4, and it depends on $w$ only through the class of $w$ modulo $N_3$. For each cylinder set $A$ with $\mu(A)=1$ one has $\left(w_{*}\mu_{\xi}\right)(A)=1$ for $\nu$-almost every $\xi$, the $\nu$-null set of the exceptional $\xi$ depending on $A$. The map $\xi\mapsto w_{*}\mu_{\xi}$ leaves the measure $\nu$ invariant.
\end{enumerate}
\item[(d)] Let $\mu^{(N)}$, for each $N$ of an infinite set of integers, be a Borel probability measure on $H_{-4}$ which is the weak limit on $H_{-4}$ of $P^{\mathrm{nat}}_{N,T_j(N)}$ as $j\to\infty$, for some sequence $\left(T_j(N)\right)_{j\ge1}$ with $T_j(N)\to\infty$; and let $\mu$ be the weak limit on $H_{-4}$ of $\mu^{(N_i)}$ along a sequence $N_i\to\infty$ of elements of this set. Then $\mu$ is a limit point of $\left(P^{\mathrm{nat}}_{N,T}\right)$ in the sense of Section~\ref{sec:curvature-field-limit}. The same is true with the roles of $N$ and $T$ interchanged.
\end{enumerate}

\emph{Part II.} There is a family of real numbers $m_{0,N}(e)$, one for each integer $N\ge0$ and each $e>0$ with $e\cdot h'^{\frac{1}{2}}\le e_1$, with the following properties. One has $\left|m_{0,N}(e)\right|<\frac{1}{2}$ for all $N$ and $e$. The numbers $m_{0,N}(e)$ do not depend on the volume: $m_{0,N}(e)$ is a function of $N$, $e$, $\bar{m}$ and of the renormalisation group scheme, that is, of the parameters, of the basis $\{e_i^0\}$ fixed for $L$ in Section~\ref{sec:averaging-fluctuations-gauge} and of the systems of primitives (Definition~\ref{def:terminology}(10)). (The family is that of Theorem~\ref{thm:free-coefficient-3}: for $N\ge N_0(e)$ the number $m_{0,N}(e)$ is the first entry of the pair $P_N(e)$ of Definition~\ref{def:tuned-pair}, and $m_{0,N}(e)=0$ for $0\le N<N_0(e)$. Since $\left|m_{0,N}(e)\right|<\frac{1}{2}$ and $\bar{m}\cdot\eta>0$, the sequence $N\mapsto m_{0,N}(e)$ is a mass assignment for each such $e$.) To every pair $(\beta,x_0)$ of data as in Section~\ref{sec:probes}, that is, to every probe $\beta$ with $b(f)>0$ and every point $x_0\in\mathbb{R}^3$ with $d_0>0$ that satisfies~\eqref{eq:setup-1}, there correspond numbers $e_2\in\left(0,e_G\right]$, so that $e_2\le e_1\cdot h'^{-\frac{1}{2}}$, and $C_{\mathrm{kap}}<\infty$, which depend only on the parameters and on $\bar{m}$, $\beta$ and $x_0$, such that the assertions (i)--(vii) below hold for every charge $e\in(0,e_2]$ and every $\mathrm{nat}\in\{\mathrm{af},\mathrm{av}\}$; in them $P^{\mathrm{nat}}_{N,T}$ are the laws of Section~\ref{sec:curvature-field-limit} formed with the charge $e$ and the mass assignment $N\mapsto m_{0,N}(e)$, that is, with $q=e\cdot\eta^{\frac{1}{2}}$ and the lattice mass $\boldsymbol{\mu}_N=\bar{m}\cdot\eta+m_{0,N}(e)$ on the torus $\Lambda_{\eta,R}$, $\eta=L^{-N}$, $R=h\cdot L^T$. In (iii)--(vii) $\mu$ is an arbitrary limit point of $\left(P^{\mathrm{nat}}_{N,T}\right)$, and in (v)--(vii) $\mu=\int\mathrm{d}\nu(\xi)\,\mu_{\xi}$ is its ergodic decomposition under $\mathbb{R}^3$.
\begin{enumerate}
\item[(i)] If $(N_i,T_i)_{i\ge1}$ is a sequence of cutoffs with $N_i\to\infty$ and $T_i\to\infty$, then some subsequence of $\left(P^{\mathrm{nat}}_{N_i,T_i}\right)_{i\ge1}$ is weakly convergent on $H_{-4}$. Each limit point $\mu$ of $\left(P^{\mathrm{nat}}_{N,T}\right)$ has the properties (P1)--(P8), and the assertions of (c) of Part~I hold for $\mu$.
\item[(ii)] Let $N\ge N_0(e)$ and $T\ge T^{*}(e)$ be integers, $\eta=L^{-N}$ and $R=h\cdot L^T$, and let $g=dd^{*}\beta'$ with $\beta'\in\{\beta_0,\beta_{+},\beta_{-}\}$, so that $g\in\{g_0,g_{+},g_{-}\}$. Then
\begin{gather*}
\left|\kappa_4\left(X_{\eta,R}(g);E^{\mathrm{nat}}\right)-e^4\cdot b_{\eta,R^{*}(e)}(0)[\beta']\right|\le\frac{1}{2}C_{\mathrm{kap}}\cdot e^{4+\sigma}+C_{\mathrm{kap}}\cdot\eta^2,\\
\left|\kappa_4\left(X_{\eta,R}(g);E^{\mathrm{nat}}\right)-e^4\cdot b(g)\right|\le\frac{1}{2}C_{\mathrm{kap}}\cdot e^{4+\sigma}+C_{\mathrm{kap}}\cdot\left(\eta^2+e^4\cdot\eta^{\frac{1}{2}}\right).
\end{gather*}
\item[(iii)] One has
\[
\left|\kappa_4^{\mu}(g)-e^4\cdot b(g)\right|\le C_{\mathrm{kap}}\cdot e^{4+\sigma}\qquad\text{for }g\in\{g_0,g_{+},g_{-}\},
\]
and
\[
\left|S_2^{\mu}(f,f)-\|f\|_2^2\right|\le C_{\mathrm{kap}}\cdot e^{2-\mathfrak{e}_2}.
\]
\item[(iv)] One has $\kappa_4^{\mu}(f)\ge\frac{1}{2}e^4\cdot b(f)>0$ and $\mathrm{Ext}^{\mu}_{x_0}(f)\ge\frac{1}{2}e^4\cdot b(f)$. Consequently $\mu$ is neither a Gaussian law nor a mixture of Gaussian laws that are invariant under the translations of $\mathbb{R}^3$ and satisfy~\eqref{eq:setup-G}.
\item[(v)] One has
\[
\nu\left(\left\{\xi:\mathrm{Ext}^{\xi}_{x_0}(f)\ge\frac{e^4\cdot b(f)}{4}\right\}\right)\ge\frac{e^4\cdot b(f)}{4\cdot17.09\cdot\|f\|_2^4}.
\]
Corollary~\ref{cor:properties-almost-every}, applied to $\mu$ and to its ergodic decomposition under $\mathbb{R}^3$, furnishes a set $\Xi_0$ of full $\nu$-measure such that, for $\xi\in\Xi_0$, the law $\mu_{\xi}$ is invariant under the translations of $\mathbb{R}^3$, ergodic under $\mathbb{R}^3$, and satisfies~\eqref{eq:setup-G}. If $\xi\in\Xi_0$ and $\mathrm{Ext}^{\xi}_{x_0}(f)\ge\frac{1}{4}e^4\cdot b(f)$, then $\mu_{\xi}$ is not Gaussian, and
\[
\max\left\{\left|\kappa_4^{\xi}(g)\right|:g\in\{g_0,g_{+},g_{-}\}\right\}\ge\frac{e^4\cdot b(f)}{10}.
\]
\item[(vi)] For $\nu$-almost every $\xi$,
\[
\left|\kappa_4^{\xi}(f)-e^4\cdot b(f)\right|\le C_{\mathrm{kap}}\cdot e^{4+\sigma},\qquad\left|C_{\xi}(f,f)-\|f\|_2^2\right|\le C_{\mathrm{kap}}\cdot e^{2-\mathfrak{e}_2}.
\]
In particular, for $\nu$-almost every $\xi$ one has $\kappa_4^{\xi}(f)\ge\frac{1}{2}e^4\cdot b(f)>0$, and the law $\mu_{\xi}$ is not Gaussian.
\item[(vii)] For $\nu$-almost every $\xi$ one has $\mathrm{Ext}^{\xi}_{x_0}(f)\ge\frac{1}{2}e^4\cdot b(f)$ and, for each of the three tests $g\in\{g_0,g_{+},g_{-}\}$,
\[
\left|\kappa_4^{\xi}(g)-e^4\cdot b(g)\right|\le C_{\mathrm{kap}}\cdot e^{4+\sigma},\qquad\left|C_{\xi}(g,g)-\|g\|_2^2\right|\le C_{\mathrm{kap}}\cdot e^{2-\mathfrak{e}_2}.
\]
\end{enumerate}
\end{theorem}

\begin{corollary}[Ergodic components]\label{cor:ergodic-components}
Let $\beta$ and $x_0$ be data as in Section~\ref{sec:probes}, with $f=dd^{*}\beta$, and let $C_{\mathrm{kap}}$ and $e_2$ be the numbers of Part~II of Theorem~\ref{thm:main-theorem} that are fixed for these data in Section~\ref{sec:parameters-scales}. Let $e\in(0,e_2]$ and $\mathrm{nat}\in\{\mathrm{af},\mathrm{av}\}$. Let $\mu$ be a limit point (Section~\ref{sec:curvature-field-limit}) of the family $\left(P^{\mathrm{nat}}_{N,T}\right)$ formed with $q=e\cdot\eta^{\frac{1}{2}}$ and the lattice masses $\boldsymbol{\mu}_N=\bar{m}\cdot\eta+m_{0,N}(e)$, the numbers $m_{0,N}(e)$ being those of Part~II of Theorem~\ref{thm:main-theorem}, and let $\mu=\int\mathrm{d}\nu(\xi)\,\mu_{\xi}$ be its ergodic decomposition under $\mathbb{R}^3$ (Section~\ref{sec:laws}).
\begin{enumerate}
\item[(a)] For $\nu$-almost every $\xi$ the ergodic component $\mu_{\xi}$ is a law on $\mathcal{S}'(\mathbb{R}^3;\Lambda^2)$ with the following properties. It is concentrated on the closed two-forms, that is, $\mu_{\xi}\left(\left\{F\in\mathcal{S}':dF=0\right\}\right)=1$; it is invariant and ergodic under the translations of $\mathbb{R}^3$, and invariant under $N_3$; and it is reflection positive at each plane $\{x_j=c\}$, $j\in\{1,2,3\}$, $c\in\mathbb{R}$. It satisfies~\eqref{eq:setup-G}, and consequently it has the properties (P1) and (P2) of Part~I of Theorem~\ref{thm:main-theorem}, with $\mu_{\xi}$ in place of $\mu$ and with the scalar product $\langle\cdot,\cdot\rangle_{L^2}$ in place of the free Maxwell covariance $C_0$ in (P1). It has the reconstruction (P7) of Part~I, with $\mu_{\xi}$ in place of $\mu$. It clusters in the sense of property~4 of Part~I(c), both in the mean and when $\min_l|a_l|\to\infty$: for all integers $n,n'\ge1$ and all $f_1,\dots,f_n,g_1,\dots,g_{n'}\in\mathcal{S}$ one has, as $a=(a_1,a_2,a_3)\in\mathbb{R}^3$ tends to infinity in such a way that $\min_l|a_l|\to\infty$,
\begin{multline*}
S^{\xi}_{n+n'}\left(f_1,\dots,f_n,\tau_ag_1,\dots,\tau_ag_{n'}\right)\\
\to S^{\xi}_n\left(f_1,\dots,f_n\right)\cdot S^{\xi}_{n'}\left(g_1,\dots,g_{n'}\right),
\end{multline*}
and the mean of the left-hand side over $a\in[0,t]^3$ tends to the right-hand side as $t\to\infty$. Finally, $\mu_{\xi}$ is not Gaussian, and $\kappa_4^{\xi}(f)\ge\frac{1}{2}e^4\cdot b(f)$.
\item[(b)] Let $j\in\{1,2,3\}$, and let $\mu=\int\mathrm{d}\nu^j(\rho)\,\rho$ be the ergodic decomposition of $\mu$ under the translations along $e_j$ (Section~\ref{sec:laws}). For $\nu^j$-almost every $\rho$ the following holds. The law $\rho$ is concentrated on the closed two-forms, that is, $\rho\left(\left\{F\in\mathcal{S}':dF=0\right\}\right)=1$; it is invariant under $\mathsf{R}_j^{*}$ and under the translations along $e_j$; it is reflection positive at each plane $\{x_j=c\}$, $c\in\mathbb{R}$; and it satisfies~\eqref{eq:setup-G}. By these properties the reconstruction of Section~\ref{sec:laws} in the direction $j$ applies to $\rho$, and its operator $\mathsf{H}_j(\rho)$ satisfies $\ker\mathsf{H}_j(\rho)=\mathbb{C}\,\Omega_j$. The law $\rho$ clusters along $e_j$: for all $A,B\in L^2(\rho)$,
\[
\rho\left(\overline{A}\cdot\left(B\circ\tau_{s\cdot e_j}^{*}\right)\right)\to\overline{\rho(A)}\cdot\rho(B)\qquad\left(s\in\mathbb{R},\ |s|\to\infty\right),
\]
where $\rho(A)=\int\mathrm{d}\rho\,A$. Moreover $\left|\kappa_4^{\rho}(f)-e^4\cdot b(f)\right|\le C_{\mathrm{kap}}\cdot e^{4+\sigma}$, and $\rho$ is not Gaussian.
\end{enumerate}
\end{corollary}

\begin{corollary}[The zero-harmonic law]\label{cor:zero-harmonic-law}
Let $\beta$ and $x_0$ be data as in Section~\ref{sec:probes}, let $C_{\mathrm{kap}}$ and $e_2$ be the numbers of Part~II of Theorem~\ref{thm:main-theorem} that are fixed for these data in Section~\ref{sec:parameters-scales}, let $e\in(0,e_2]$, and let $P^{\mathrm{av}}_{N,T}$ be the laws of Section~\ref{sec:curvature-field-limit} formed with $q=e\cdot\eta^{\frac{1}{2}}$ and the lattice masses $\boldsymbol{\mu}_N=\bar{m}\cdot\eta+m_{0,N}(e)$, the numbers $m_{0,N}(e)$ being those of Part~II of Theorem~\ref{thm:main-theorem}. For every cutoff $(N,T)$ denote by $P^0_{N,T}$ the law of the curvature field $F^{\sharp}$, where $F=\eta^{-\frac{3}{2}}\cdot F^u$ and $F^u$ has the zero-harmonic law of Section~\ref{sec:lattice-laws} on the cubic torus $\Lambda_n$, $n=h\cdot L^{T+N}$, with the lattice mass $\boldsymbol{\mu}_N$ and the lattice charge $q=e\cdot\eta^{\frac{1}{2}}$. For every integer $N$ with $\boldsymbol{\mu}_N>0$, the Borel probability measures on $H_{-4}$ that are weak limits on $H_{-4}$ of sequences $\left(P^0_{N,T_k}\right)_{k\ge1}$ with $T_k\to\infty$ are exactly the weak limits on $H_{-4}$ of sequences $\left(P^{\mathrm{av}}_{N,T_k}\right)_{k\ge1}$ with $T_k\to\infty$. For each $N$ of an infinite set of integers $N$ with $\boldsymbol{\mu}_N>0$ let $\mu^{(N)}$ be such a weak limit, and let $\mu$ be the weak limit on $H_{-4}$ of a sequence $\left(\mu^{(N_i)}\right)_{i\ge1}$ whose indices $N_i$ belong to this set and tend to $\infty$. Then $\mu$ is one of the limit points of the family $\left(P^{\mathrm{av}}_{N,T}\right)$ (Section~\ref{sec:curvature-field-limit}), and it has the properties (i) and (iii)--(vii) of Part~II of Theorem~\ref{thm:main-theorem}. (If $0\le N<N_0(e)$, then $m_{0,N}(e)=0$ by Part~II of Theorem~\ref{thm:main-theorem}, so that $\boldsymbol{\mu}_N=\bar{m}\cdot L^{-N}>0$; thus the first assertion applies to each of these finitely many $N$. For $N\ge N_0(e)$ and $\eta=L^{-N}$, each bound on $m_{0,N}(e)$ in terms of $e$ and $\eta$ that is proved is either a bound $\left|m_{0,N}(e)\right|\le c\cdot\left(e\cdot\eta^{\frac{1}{2}}\right)^{\frac{3}{4}-8\varepsilon}$ with a number $c>0$ independent of $N$, for instance $c=\frac{1}{8}$ in Proposition~\ref{prop:pair-vanishing-terminal} and $c=\frac{1}{12}$ in its proof, both applied on the torus $R_{\mathrm{top}}(e)$, or a consequence of such a bound; for each such $c$ the right-hand side exceeds $\bar{m}\cdot\eta$ for all sufficiently large $N$ (Remark~\ref{rem:free-coefficient-1}(3)), so that these bounds do not determine the sign of $\boldsymbol{\mu}_N$ for large $N$. In particular it is open whether $\boldsymbol{\mu}_N>0$ for infinitely many $N$, that is, whether the mass assignment of Part~II admits an infinite set of integers $N$ as in the second assertion.)
\end{corollary}

\section{The estimates}\label{sec:estimates}

Theorem~\ref{thm:main-theorem} and Corollaries~\ref{cor:ergodic-components} and~\ref{cor:zero-harmonic-law} are deduced from Estimates~1--10 below. These are results of Sections~\ref{sec:gaussian-bounds-curvature}--\ref{sec:all-flux-functional} and~\ref{sec:free-coefficient-chessboard}, restated here in the notation of Section~\ref{sec:setup} and, in the case of Estimates~3, 6, 7 and~8, with the cells and mirror copies of Section~\ref{sec:reference-tori-cells}, the probes and their placement of Section~\ref{sec:probes-discretisation} and the reference numbers of Section~\ref{sec:numbers-attached-reference}; the parenthesis at the end of each estimate names the statements that contain it and the sections in which these are proved. Nothing is proved in this section except the deduction in its last paragraph.

\paragraph{Estimate~1 (Gaussian domination).}
For every torus $n=(n_1,n_2,n_3)$ with integers $n_{\mu}\ge2$, every $q>0$, every $\boldsymbol{\mu}>-\frac{1}{2}$, every $\mathrm{nat}\in\{\mathrm{af},\mathrm{av}\}$ and every real two-cochain $g\in C^2(n)$ one has
\[
E^{\mathrm{nat}}_n\exp\langle g,F^u\rangle\le\exp\left(\frac{1}{2}\|g\|^2\right).
\]
(For the all-flux law, Theorem~\ref{thm:gaussian-domination} and Corollary~\ref{cor:partition-functions-1}(a), where $\boldsymbol{\mu}>-1$, and Section~\ref{sec:recalled-results}; for the harmonic-averaged law, inequality~\eqref{eq:gaussian-bounds-5} of Theorem~\ref{thm:limit-vanishing-proca}, which holds on cubic tori for $\boldsymbol{\mu}>-1$ with $\|Pg\|$ in place of $\|g\|$, $P$ being the orthogonal projection on the range of $d$, Remark~\ref{rem:tori-unequal-sides} for tori with unequal sides, Lemma~\ref{lem:flux-sectors}(b) for the identification of the law and, in physical units at a cutoff, Proposition~\ref{prop:gaussian-domination}; Theorem~\ref{thm:gaussian-domination-positive}.)

\paragraph{Estimate~2 (lattice symmetries, reflection positivity, chessboard estimate).}
For $q>0$ and $\boldsymbol{\mu}>-\frac{1}{2}$ the laws $E^{\mathrm{af}}_n$ and $E^{\mathrm{av}}_n$ of a cubic torus $n$ are invariant under the lattice translations, under $O_{\mathrm{cub}}$ and under $F^u\mapsto-F^u$; and at every cutoff $(N,T)$, on the cubic torus $n=R/\eta$ identified with $\Lambda_{\eta,R}$, with the physical curvature $F=\eta^{-\frac{3}{2}}F^u$ and with $\tilde{m}:=\frac{n-1}{2}$, one has, for $\mathrm{nat}\in\{\mathrm{af},\mathrm{av}\}$, $\nu_0\in\{1,2,3\}$, every finite family $(g_a)$ of real two-cochains on $\Lambda_{\eta,R}$ supported on $\mathcal{P}^{+\circ}_{\nu_0}$ and every family $(c_a)$ of complex numbers with the same indices,
\[
\sum_{a,b}\overline{c_a}\cdot c_b\cdot E^{\mathrm{nat}}_n\exp\left(i\left\langle g_b-\hat{\mathsf{R}}_{\nu_0}g_a,F\right\rangle_{\eta}\right)\ge0,
\]
where, as in Sections~\ref{sec:tori-cells-cochains} and~\ref{sec:admissible-families-lattice}, the coordinates of a point of $\Lambda_{\eta,R}$ are those of its representative in $\eta\cdot\{-\tilde{m},\dots,\tilde{m}\}^3$, the set $\mathcal{P}^{+\circ}_{\nu_0}$ consists of the plaquettes all of whose vertices have their $\nu_0$-th coordinate in $\{0,\eta,\dots,\tilde{m}\cdot\eta\}$ and which lie in neither of the planes $\{y_{\nu_0}=0\}$ and $\{y_{\nu_0}=\tilde{m}\cdot\eta\}$, and $\hat{\mathsf{R}}_{\nu_0}$ is the reflection $y_{\nu_0}\mapsto-y_{\nu_0}$ acting on two-cochains, with the sign $-$ on the plaquettes that contain the direction $\nu_0$. For the same $q$ and $\boldsymbol{\mu}$, if a torus $n$ carries the cell structure of a width $\boldsymbol{\lambda}=m_{\mathrm{w}}+\frac{1}{2}$, $m_{\mathrm{w}}\ge2$ an integer, of Definition~\ref{def:decomposition-torus-cells}, so that $n_{\mu}=\boldsymbol{\lambda}\cdot N_{\mu}$ with an even positive integer $N_{\mu}$ for every direction $\mu$ and the cells are bounded alternately by site planes and link planes, and if $N_n$ is the number of its cells, then for $\boldsymbol{\omega}\in\left\{\boldsymbol{\omega}_n,\boldsymbol{\omega}_n^{(0)}\right\}$, every $O$ in the class $\mathfrak{A}^{+}$ of that definition and every family $(O_c)$ of elements of $\mathfrak{A}^{+}$ indexed by the cells,
\[
\boldsymbol{\omega}\left(O^{\sharp}\right)\ge0,\qquad\left|\boldsymbol{\omega}\left(\prod_cO_c^{[c]}\right)\right|\le\prod_c\boldsymbol{\omega}\left(O_c^{\sharp}\right)^{\frac{1}{N_n}},
\]
$O^{[c]}$ being the mirror copy of $O$ in the cell $c$ and $O^{\sharp}$ the product of all mirror copies of $O$.
(Theorem~\ref{thm:free-coefficient-1}(d), in particular (d5), for the symmetries; for reflection positivity, Proposition~\ref{prop:gaussian-bounds-2}, which is the reflection positivity of the Proca parent of Section~\ref{sec:proca-parent-harmonic}, a measure on the one-cochains formed with a Proca mass $\varepsilon>0$, Proposition~\ref{prop:reflection-positivity}, which derives from it the reflection positivity of the harmonic-averaged law in the limit $\varepsilon\downarrow0$ by the identity~\eqref{eq:gaussian-bounds-4}, and Corollary~\ref{cor:averaged-laws-admissible}, for the harmonic-averaged law, and Section~\ref{sec:admissibility-limit-points}, with Corollary~\ref{cor:reflection-positivity-fixed} and Lemma~\ref{lem:dependence-boundary-slices}, for the all-flux law, the inequality being condition~\eqref{eq:continuum-infinite-Ad3} of Definition~\ref{def:admissible-families}; Proposition~\ref{prop:chessboard-estimate} for the chessboard estimate.)

\paragraph{Estimate~3 (all volumes from two tori).}
Let $N\ge0$ and $T^{*}\ge1$ be integers, $\eta=L^{-N}$, $R^{*}=h\cdot L^{T^{*}}$ and $q=e\cdot\eta^{\frac{1}{2}}$, where the charge satisfies $0<e\le\frac{1}{3}$ and $6N^{*}\cdot e^2\cdot R^{*}\cdot\log\left(e^{-1}\right)\le\pi^2$, which is~\eqref{eq:free-coefficient-6}; let the lattice mass satisfy $\boldsymbol{\mu}_N>-\frac{1}{2}$, let $\beta$ be a probe that fits $(N,T^{*})$, that is $2\rho_{\beta}+3\eta\le l_c=\frac{R^{*}}{2L}$, and let $X=\langle d\hat{\mathfrak{a}},F^u\rangle$, its copies $X_c=\langle d\hat{\mathfrak{a}}_c,F^u\rangle$, one for each cell $c$, the number $\boldsymbol{\sigma}\ge0$ with $\boldsymbol{\sigma}^2=\|d\hat{\mathfrak{a}}\|^2$ and the polynomials $H_2(x)=x^2-\boldsymbol{\sigma}^2$ and $H_4(x)=x^4-6\boldsymbol{\sigma}^2\cdot x^2+3\boldsymbol{\sigma}^4$ be formed with the placement of Section~\ref{sec:probes-discretisation}; for an integer $k\ge1$ the torus $k$ is the cubic torus of side $k\cdot R^{*}$ of Section~\ref{sec:reference-tori-cells}, cut into $(2kL)^3$ cells, with the laws $E^{\mathrm{nat}}_k$ and the partition function $Z_k$ formed with $q$ and $\boldsymbol{\mu}_N$; for real numbers $c_2,c_4$ and for $m\in\{2,4\}$ put
\begin{gather*}
\dot{M}_4:=600\,(N^{*})^2\cdot\boldsymbol{\sigma}^4\cdot e^{-4}+|c_4|,\qquad\dot{M}_2:=12N^{*}\cdot\boldsymbol{\sigma}^2\cdot e^{-2}+|c_2|,\\
Y:=\log^{+}\left((N^{*})^2\cdot\left(600\,\boldsymbol{\sigma}^4+12\,\boldsymbol{\sigma}^2+|c_4|+|c_2|\right)\right),\qquad\check{X}:=1+\frac{Y+1}{6\log\left(e^{-1}\right)},\\
\boldsymbol{\Gamma}^{\mathrm{af}}:=\exp(\delta),\qquad\boldsymbol{\Theta}^{\mathrm{af}}_m:=\Theta^{(0)}_m+2e^6\cdot\dot{M}_m,\\
\boldsymbol{\Gamma}^{\mathrm{av}}:=\exp\left(\delta+\frac{3\log\left(4+16\,\check{X}\right)+6}{N^{*}}\right),\qquad\boldsymbol{\Theta}^{\mathrm{av}}_m:=\max\left(\Theta^{(0)}_m,e^2\right),
\end{gather*}
where $\log^{+}\colon[0,\infty)\to[0,\infty)$ is given by $\log^{+}0:=0$ and $\log^{+}x:=\max(0,\log x)$ for $x>0$, and where
\[
\delta=\frac{L^3\log Z_1-\log Z_L}{(L-1)^3\cdot N^{*}},\qquad\Theta^{(0)}_m=\left|\left(\prod_c\left(\partial_{t_c}^m-c_m\right)\mathring{\Phi}\right)(0)\right|^{\frac{1}{N^{*}}},
\]
the product being taken over the $N^{*}$ cells $c$ of the torus $1$, are the numbers of Section~\ref{sec:numbers-attached-reference}; then~\eqref{eq:free-coefficient-10} holds: for $\mathrm{nat}\in\{\mathrm{af},\mathrm{av}\}$, every integer $k\in\{1\}\cup\{L,L+1,\dots\}$, $m\in\{2,4\}$ and every set $S$ of cells of the torus $k$,
\[
\left|E^{\mathrm{nat}}_k\prod_{c\in S}\left(e^{-m}\cdot H_m(X_c)-c_m\right)\right|\le\left(\boldsymbol{\Gamma}^{\mathrm{nat}}\cdot\boldsymbol{\Theta}^{\mathrm{nat}}_m\right)^{\#|S|},
\]
the right side being formed on the two reference tori $k=1$ and $k=L$ alone. Moreover $\delta$ satisfies~\eqref{eq:free-coefficient-9},
\[
0\le\delta\le\frac{a_L}{N^{*}}\cdot\log\frac{2\pi}{e^2\cdot R^{*}}+\frac{K_L+D^{\mathrm{mat}}+1}{(L-1)^3\cdot N^{*}},
\]
with the numbers $D^{\mathrm{mat}}$ and $K_L$ of Section~\ref{sec:numbers-attached-reference}, and, for each such $\mathrm{nat}$ and $k$, the cumulants $\kappa_2=\kappa_2\left(X;E^{\mathrm{nat}}_k\right)$ and $\kappa_4=\kappa_4\left(X;E^{\mathrm{nat}}_k\right)$ satisfy~\eqref{eq:free-coefficient-7} and~\eqref{eq:free-coefficient-8},
\begin{gather*}
\left|e^{-2}\cdot\left(\kappa_2-\boldsymbol{\sigma}^2\right)-c_2\right|\le\boldsymbol{\Gamma}^{\mathrm{nat}}\cdot\boldsymbol{\Theta}^{\mathrm{nat}}_2,\\
\left|e^{-4}\cdot\kappa_4-\left(c_4-3c_2^2\right)\right|\le\boldsymbol{\Gamma}^{\mathrm{nat}}\cdot\boldsymbol{\Theta}^{\mathrm{nat}}_4+3\boldsymbol{\Gamma}^{\mathrm{nat}}\cdot\boldsymbol{\Theta}^{\mathrm{nat}}_2\cdot\left(2|c_2|+\boldsymbol{\Gamma}^{\mathrm{nat}}\cdot\boldsymbol{\Theta}^{\mathrm{nat}}_2\right).
\end{gather*}
(Theorem~\ref{thm:free-coefficient-2}; Step~2 of its proof in Section~\ref{sec:proof-cumulant-bounds} contains the identity $EX^4-3\left(EX^2\right)^2=EH_4(X)-3\left(EH_2(X)\right)^2$, valid for every real random variable $X$ with finite fourth moment and every value of $\boldsymbol{\sigma}^2$.)

\paragraph{Estimate~4 (tightness and limit points).}
For every cutoff $(N,T)$ let $\mathrm{Pr}_{N,T}$ be a Borel probability measure on the space $\ker d$ of the closed two-cochains of $\Lambda_{\eta,R}$ such that \eqref{eq:continuum-infinite-Ad1} $\mathrm{Pr}_{N,T}$ is invariant under the lattice translations and under $O_{\mathrm{cub}}$, \eqref{eq:continuum-infinite-Ad2} $\int\mathrm{d}\mathrm{Pr}_{N,T}(F)\,\exp\langle g,F\rangle_{\eta}\le\exp\left(\frac{1}{2}|g|_{\eta}^2\right)$ for every real two-cochain $g$, and \eqref{eq:continuum-infinite-Ad3} the inequality of Estimate~2 holds with the integral with respect to $\mathrm{Pr}_{N,T}$ in place of $E^{\mathrm{nat}}_n$, no relation between the measures at different cutoffs being required; then $\sup_{(N,T)}\int\mathrm{d}\mathrm{Pr}_{N,T}(F)\,\|F^{\sharp}\|_{-3}^2\le61$, the laws of $F^{\sharp}$ under the measures $\mathrm{Pr}_{N,T}$ are tight on $H_{-4}$, and at every cutoff, for an integer $n\ge1$ and real $f_1,\dots,f_n\in\mathcal{S}$, with the expectation $E$ and the joint cumulant $\kappa_n$ with respect to $\mathrm{Pr}_{N,T}$,
\begin{gather*}
\left|E\prod_{i=1}^nX_{\eta,R}(f_i)\right|\le2\left(2^n\cdot n!\right)^{\frac{1}{2}}\cdot\prod_{i=1}^n|(f_i)_{\eta}|_{\eta},\\
\left|\kappa_n\left(X_{\eta,R}(f_1),\dots,X_{\eta,R}(f_n)\right)\right|\le(\log2)\cdot\left(\frac{5\exp(1)}{2}\right)^n\cdot n!\cdot\prod_{i=1}^n|(f_i)_{\eta}|_{\eta},
\end{gather*}
where, with $Q_R:=\left[-\frac{R}{2},\frac{R}{2}\right)^3$ and with $|y|_2^2f$ the function $y\mapsto|y|_2^2f(y)$,
\[
|f_{\eta}|_{\eta}\le p_R(f):=\|f\|_2+26R^{-2}\cdot\left\|\,|y|_2^2f\right\|_{L^2(\mathbb{R}^3\setminus Q_R)},
\]
$p_R(f)\to\|f\|_2$ as $R\to\infty$, and $|f_{\eta}|_{\eta}\le\|f\|_2$ if $\operatorname{supp}f\subset y_0+\left(-\frac{R}{2},\frac{R}{2}\right)^3$ for some $y_0\in\mathbb{R}^3$. Every Borel probability measure $\mu$ on $H_{-4}$ to which the laws of $F^{\sharp}$ under $\mathrm{Pr}_{N_i,T_i}$ converge weakly on $H_{-4}$ along a sequence of cutoffs with $N_i\to\infty$ and $T_i\to\infty$, no relation between $N_i$ and $T_i$ being imposed, has the properties (P1), (P2) and (P4)--(P8) of Part~I of Theorem~\ref{thm:main-theorem}, (P8) with the expectation with respect to $\mathrm{Pr}_{N_i,T_i}$ in place of $E_i$; its moments satisfy $S_0=1$, and every $S_n$ is symmetric and real on real tests; $\mu$ is invariant under $F\mapsto-F$ if every $\mathrm{Pr}_{N_i,T_i}$ is; and for the laws $P^{\mathrm{nat}}_{N,T}$ the iterated weak limits on $H_{-4}$ described in Part~I(d) of Theorem~\ref{thm:main-theorem}, first as $T\to\infty$ and then as $N\to\infty$, or in the other order, are limit points.
(Proposition~\ref{prop:admissible-families-curvature}; Theorems~\ref{thm:estimates-fixed-cutoff}, \ref{thm:second-moments-tightness} and Lemma~\ref{lem:norm-bounds-lattice}; Theorems~\ref{thm:continuum-infinite-1}, \ref{thm:moments-limit-point} and Corollary~\ref{cor:transfer-operators-coordinate}; Remark~\ref{rem:repeated-limits} for iterated limits; Proposition~\ref{prop:sub-gaussian-random} on linear forms with a Gaussian bound, Lemma~\ref{lem:lattice-symmetries-cell} on the test map and the symmetries, Lemma~\ref{lem:hermite-coefficients-tensor} on tensor tests.)

\paragraph{Estimate~5 (ergodic components).}
Let $\mu$ satisfy~\eqref{eq:setup-H}; then on the set of the laws on $\mathcal{S}'$, endowed with the $\sigma$-algebra generated by the maps $\rho\mapsto\rho(A)$, $A$ a cylinder set, there is exactly one probability measure $\nu$ that is carried by the laws invariant and ergodic under $\mathbb{R}^3$ and satisfies $\mu=\int\mathrm{d}\nu(\xi)\,\mu_{\xi}$, where $\mu_{\xi}:=\xi$; for $\nu$-almost every $\xi$ the law $\mu_{\xi}$ satisfies~\eqref{eq:setup-G} with the same right side, is invariant under $\mathbb{R}^3$ and under $N_3$, is reflection positive at every plane $\{x_j=c\}$, $j\in\{1,2,3\}$, $c\in\mathbb{R}$, and has the properties 1--4 of Part~I(c) of Theorem~\ref{thm:main-theorem}, and for every $w\in O_{\mathrm{cub}}$ the law $w_{*}\mu_{\xi}$ has the properties 1--4 and depends on $w$ only through the class of $w$ modulo $N_3$, all these assertions holding outside one $\nu$-null set; for every cylinder set $A$ with $\mu(A)=1$ one has $\mu_{\xi}(A)=1$ for $\nu$-almost every $\xi$, and $\left(w_{*}\mu_{\xi}\right)(A)=1$ for every $w\in O_{\mathrm{cub}}$ and $\nu$-almost every $\xi$, the $\nu$-null set of the exceptional $\xi$ depending on $A$; for every $w\in O_{\mathrm{cub}}$ the map $\xi\mapsto w_{*}\mu_{\xi}$ leaves $\nu$ invariant; and, for $j\in\{1,2,3\}$ and the decomposition $\mu=\int\mathrm{d}\nu^j(\rho)\,\rho$ of $\mu$ into laws invariant and ergodic under the translations along $e_j$, $\nu^j$-almost every $\rho$ is invariant under these translations and under $\mathsf{R}_j^{*}$, is reflection positive at every plane $\{x_j=c\}$, $c\in\mathbb{R}$, satisfies~\eqref{eq:setup-G}, has $\ker\mathsf{H}_j(\rho)=\mathbb{C}\,\Omega_j$ and clusters along $e_j$ in the sense of Corollary~\ref{cor:ergodic-components}(b), and, for every cylinder set $A$ with $\mu(A)=1$, $\nu^j$-almost every $\rho$ satisfies $\rho(A)=1$, the $\nu^j$-null set of the exceptional $\rho$ depending on $A$. Let $(\boldsymbol{\Omega},\boldsymbol{\Sigma},\mu)$ be a probability space, $d\ge1$ an integer, $(T_y)_{y\in\mathbb{Z}^d}$ an action of $\mathbb{Z}^d$ on $\boldsymbol{\Omega}$ by measurable maps that preserve $\mu$, and $U(y)G:=G\circ T_{-y}$; let $G$ be a real function with $G\in L^J(\mu)$ for every integer $J\ge1$, and let $\mathsf{u}\ge0$ be such that
\[
\left|\int\mathrm{d}\mu\,\prod_{i=1}^JU(y_i)G\right|\le\mathsf{u}^J
\]
for every integer $J\ge1$ and all pairwise distinct $y_1,\dots,y_J\in\mathbb{Z}^d$; then the orthogonal projection $P_0G$ of $G$ on the subspace of the vectors of $L^2(\mu)$ that are invariant under all $U(y)$ satisfies $|P_0G|\le\mathsf{u}$ almost everywhere; and, for a law $\mu$ on $\mathcal{S}'$ that is invariant under the translations of $\mathbb{R}^3$, with its ergodic decomposition $\mu=\int\mathrm{d}\nu(\xi)\,\mu_{\xi}$, the conditional expectation $E\left[\,\cdot\mid\mathcal{I}\right]$ under $\mu$ is the orthogonal projection of $L^2(\mu)$ on the vectors that are invariant under $\mathbb{R}^3$, and for $G\in L^2(\mu)$ and $\nu$-almost every $\xi$ the function $G$ lies in $L^2(\mu_{\xi})$ and $E\left[G\mid\mathcal{I}\right]$ coincides $\mu_{\xi}$-almost surely with the constant $E_{\xi}G:=\int\mathrm{d}\mu_{\xi}\,G$. Let further $j\in\{1,2,3\}$, let $\mu$ be a law on $\mathcal{S}'$ that is invariant under the translations along $e_j$, with its decomposition $\mu=\int\mathrm{d}\nu^j(\rho)\,\rho$ into laws invariant and ergodic under the translations along $e_j$, and let $\mathcal{I}^j$ be the $\sigma$-algebra of the cylinder sets $A$ with $\tau_b^{*}(A)=A$ for every $b\in\mathbb{R}\cdot e_j$. Then the conditional expectation $E\left[\,\cdot\mid\mathcal{I}^j\right]$ under $\mu$ is the orthogonal projection of $L^2(\mu)$ on the vectors that are invariant under the translations along $e_j$. There is a regular version $(\mu^j_{\omega})_{\omega\in\mathcal{S}'}$ of the conditional law of $\mu$ given $\mathcal{I}^j$ such that $\nu^j$ is the image of $\mu$ under the map $\omega\mapsto\mu^j_{\omega}$; here every $\mu^j_{\omega}$ is a law on $\mathcal{S}'$, and for every cylinder set $A$ the function $\omega\mapsto\mu^j_{\omega}(A)$ is $\mathcal{I}^j$-measurable and equal to $E\left[1_A\mid\mathcal{I}^j\right](\omega)$ for $\mu$-almost every $\omega$. For every $\mu$-integrable function $G$ on $\mathcal{S}'$ and $\mu$-almost every $\omega$ the function $G$ is $\mu^j_{\omega}$-integrable and $\mu^j_{\omega}(G)=E\left[G\mid\mathcal{I}^j\right](\omega)$. Consequently, if $G$ is a real $\mu$-integrable function on $\mathcal{S}'$ and $\mathsf{u}\ge0$ is a number with $\left|E\left[G\mid\mathcal{I}^j\right]\right|\le\mathsf{u}$ $\mu$-almost surely, then, for $\nu^j$-almost every $\rho$, the function $G$ is $\rho$-integrable and $\left|\rho(G)\right|\le\mathsf{u}$.
(Corollary~\ref{cor:ergodic-components-2}, Theorem~\ref{thm:ergodic-decomposition}, Remark~\ref{rem:components-all-translations}, Theorem~\ref{thm:cubic-symmetry-components}(i) and Corollary~\ref{cor:properties-almost-every} for the components under $\mathbb{R}^3$, Theorem~\ref{thm:decomposition-along-coordinate} and Theorem~\ref{thm:ergodic-decomposition}(a), (d) with $J=\{j\}$ for one axis; Lemma~\ref{lem:uniform-bound-invariant}, Lemma~\ref{lem:averages-translations-2}(b), (d), Proposition~\ref{prop:existence-conditional-expectations}(b), Lemma~\ref{lem:averages-translations}(iv), also with $J=\{j\}$, and for one axis identity~\eqref{eq:ergodic-components-1} with $J=\{j\}$ and the first paragraph of the proof of Lemma~\ref{lem:measurability-defining-properties}; Theorem~\ref{thm:invariant-vectors-clustering}, Proposition~\ref{prop:reconstruction-direction}, Lemma~\ref{lem:consequences-gaussian-bound}, Lemmas~\ref{lem:moment-bounds} and~\ref{lem:exponentials-determination-law}(b).)

\paragraph{Estimate~6 (free coefficient).}
For every $a\in C_c^{\infty}(\mathbb{R}^3;\mathbb{R}^3)$ there is a number $C_a<\infty$, depending on $a$ only, such that
\[
\left|q_{4,\mathsf{M}}(a)-\frac{1}{240\pi\cdot\mathsf{M}^5}\int_{\mathbb{R}^3}\mathrm{d}x\,|\operatorname{curl}a(x)|_2^4\right|\le C_a\cdot\mathsf{M}^{-7}\qquad\text{for every }\mathsf{M}\ge1;
\]
for the function $\rho$ with $\rho(x):=\exp\left(-\frac{1}{1-|x|_2^2}\right)$ if $|x|_2<1$ and $\rho(x):=0$ otherwise, the field $a:=\left(\partial_2\rho,-\partial_1\rho,0\right)$ with its constant $C_a$, the number $I_{\beta}:=\int_{\mathbb{R}^3}\mathrm{d}x\,|\operatorname{curl}a(x)|_2^4$ and the probes $\beta^{(l)}$, $l\ge1$, given by $\beta^{(l)}_{12}(x):=l^2\rho(x/l)$ and $\beta^{(l)}_{13}:=\beta^{(l)}_{23}:=0$, with $f_l:=dd^{*}\beta^{(l)}$, one has $I_{\beta}>0$ and
\[
\left|b(f_l)-\frac{l^3\cdot I_{\beta}}{10\pi\cdot\bar{m}^5}\right|\le24\,C_a\cdot l\cdot\bar{m}^{-7}\qquad\text{whenever }\bar{m}\cdot l\ge1,
\]
so that $b(f_l)>0$ as soon as $\bar{m}\cdot l\ge1$ and $\left(\bar{m}\cdot l\right)^2>240\pi\cdot C_a/I_{\beta}$. Let $\beta$ be a probe, with the number $D_{\beta}$ and the cube $Q_{\beta}$ of Section~\ref{sec:probes-discretisation}; there are $C_1,C_2<\infty$, depending only on $\bar{m}$ and $\beta$, such that for all $\eta\in(0,1]$ and $R>0$ with $R/\eta$ an integer and $Q_{\beta}$ in the interior of a period cube of the torus of side $R$ the following holds: the cochain $a_{\eta}$ of Section~\ref{sec:free-coefficient} is orthogonal to the exact and to the constant one-cochains; the two-form $f^{(\eta)}$ of Section~\ref{sec:probes-discretisation} is smooth and supported in $\operatorname{supp}\beta+[-\eta,\eta]^3\subset Q_{\beta}$, with $\left(f^{(\eta)}\right)_{\eta}=d_{\eta}a_{\eta}$ and $\left\|f^{(\eta)}-f\right\|_{H^s}\le\eta^2\cdot\|\beta\|_{H^{s+4}}$ for every real $s\ge0$, the norms being those of Section~\ref{sec:lattices-tests-function}; if $R\ge D_{\beta}+1$, then
\[
\left|b_{\eta,R}(0)[\beta]-b_R(f)\right|\le C_1\cdot\eta^{\frac{1}{2}},\qquad\left|b_R(f)-b(f)\right|\le C_2\cdot\exp\left(-\bar{m}\cdot\left(R-D_{\beta}\right)\right),
\]
where $b_R(f):=-6\operatorname{Tr}\left[\left(S^R\,i\gamma\cdot a\right)^4\right]$, formed with the periodisation of $a=d^{*}\beta$ and with $S^R=\left(\bar{m}+\gamma\cdot\partial\right)^{-1}$ on the spinor fields $\psi$ with $\psi(x+R\cdot e_{\mu})=-\psi(x)$, $\mu=1,2,3$, the operator in brackets being of trace class, is the free coefficient on the torus of side $R$ (Section~\ref{sec:free-fourth-order}); and if $\eta=L^{-N}$ for an integer $N\ge0$ and $\beta$ fits $(N,T^{*})$ for an integer $T^{*}\ge1$, then, with the placement $v$ of Section~\ref{sec:probes-discretisation} and the objects $\hat{\mathfrak{a}}$, $X$, $\boldsymbol{\sigma}$ formed with it on $\Lambda_{\eta,R}$ as in that section, one has $X=X_{\eta,R}\left(\tau_vf^{(\eta)}\right)$ and
\[
\boldsymbol{\sigma}\le\left\|f^{(\eta)}\right\|_2\le\|f\|_2+\|\beta\|_{H^4},\qquad\left|\boldsymbol{\sigma}-\|f\|_2\right|\le\eta^2\cdot\|\beta\|_{H^4}+21\,\eta\cdot|\beta|_{C^3}\cdot D_{\beta}^{\frac{3}{2}},
\]
and if in addition $R=k\cdot h\cdot L^{T^{*}}$ with an integer $k\ge1$, so that $\Lambda_{\eta,R}$ is the torus $k$ of Section~\ref{sec:reference-tori-cells}, then the support of $d\hat{\mathfrak{a}}$ is contained in $P^o(\mathfrak{c}_0)$, and the cochains $d\hat{\mathfrak{a}}_c$ of Section~\ref{sec:probes-discretisation}, $c$ running over the cells of the torus $k$, have pairwise disjoint supports and satisfy $\left\|d\hat{\mathfrak{a}}_c\right\|=\boldsymbol{\sigma}$.
(Theorem~\ref{thm:free-coefficient-1}(a), proved in Section~\ref{sec:coefficient-fourth-order}; Theorem~\ref{thm:free-coefficient-1}(b1)--(b3), (b5), proved in Section~\ref{sec:discretised-probe}; Theorem~\ref{thm:free-coefficient-1}(c1)--(c3), proved in Section~\ref{sec:continuum-infinite-volume-2}.)

\paragraph{Estimate~7 (bare mass shifts independent of the torus).}
Let $0<\bar{m}\le1$ and $e>0$ with $e\cdot h'^{\frac{1}{2}}\le e_1$, let $m_{0,N}(e)$ be the first entry of the pair $P_N(e)=\left(m_{0,N},\boldsymbol{\varepsilon}_{0,N}\right)$ of Definition~\ref{def:tuned-pair} if $N\ge N_0(e)$ and $m_{0,N}(e):=0$ if $0\le N<N_0(e)$, and put
\[
C_{\mathrm{mat}}:=\left(L^3+1\right)\cdot\left(\varepsilon_5\left(L^{-m_0-1}\right)+1\right),\qquad\varepsilon_5(x):=312\,\log\left(\left(1-\exp\left(-\frac{x}{4}\right)\right)^{-1}\right)\quad(x>0);
\]
then $e\le\frac{1}{3}$, the number $T^{*}(e)$ is an integer $\ge1$ with $R^{*}(e)=h\cdot L^{T^{*}(e)}$, the length $l_c=R^{*}(e)/(2L)$ of Section~\ref{sec:reference-tori-cells} tends to infinity as $e\to0$, and the numbers $m_{0,N}(e)$, $N\ge0$, are real with $\left|m_{0,N}(e)\right|<\frac{1}{8}$, so that $N\mapsto m_{0,N}(e)$ is a mass assignment; these numbers are functions of $N$, $e$, $\bar{m}$ and of the scheme of Definition~\ref{def:terminology}(10), and they do not depend on a torus. For every integer $N\ge N_0(e)$, with $T^{*}=T^{*}(e)$ and the lattice mass $\boldsymbol{\mu}_N=\bar{m}\cdot\eta+m_{0,N}(e)$, one has $D^{\mathrm{mat}}\le C_{\mathrm{mat}}$ for the number $D^{\mathrm{mat}}=L^3\log\rho_{n_{*}}-\log\rho_{L\cdot n_{*}}$ of Section~\ref{sec:numbers-attached-reference}, in which $\rho_n$, the partition function of that section on the cubic torus $n$ averaged over the twists, is taken at the numbers $n_{*}$ and $L\cdot n_{*}$ of sites per side of the reference tori; and if $e\le e_G$, then
\[
6N^{*}\cdot e^2\cdot R^{*}(e)\cdot\log\left(e^{-1}\right)\le\pi^2,
\]
which is the inequality~\eqref{eq:free-coefficient-6} for $R^{*}=R^{*}(e)$. (Part~\ref{it:free-coefficient-3-1} of Theorem~\ref{thm:free-coefficient-3}, proved in Section~\ref{sec:proof-theorem}.)

\paragraph{Estimate~8 (four-jets).}
One has $0<\mathfrak{e}\le\frac{1}{4}$, $0\le\mathfrak{e}_2<1$ and $\sigma>0$; and for every probe $\beta$ there are numbers $e_1(\beta)\in\left(0,e_G\right]$ and $C^G_{\beta}<\infty$, depending only on the parameters, on $\bar{m}$ and on $\beta$, such that for $0<e\le e_1(\beta)$ and every integer $N\ge N_0(e)$, with $T^{*}=T^{*}(e)$ and the lattice mass $\boldsymbol{\mu}_N=\bar{m}\cdot\eta+m_{0,N}(e)$, the probe $\beta$ fits $\left(N,T^{*}(e)\right)$, the derivatives
\[
c_m=c_m(N,e,\beta):=N_0^{(m)}(0)\qquad(m=2,4)
\]
at $t=0$ of the free ratio $N_0(t)$ of Section~\ref{sec:free-coefficient} on $\Lambda_{\eta,R^{*}(e)}$, formed with $\tau_v\beta$, where $v$ is the placement of Section~\ref{sec:probes-discretisation}, are real, and the numbers $\Theta^{(0)}_2$, $\Theta^{(0)}_4$ of Section~\ref{sec:numbers-attached-reference}, formed with these $c_2$, $c_4$, satisfy
\begin{gather*}
\Theta^{(0)}_4\le C^G_{\beta}\cdot e^{\mathfrak{e}},\qquad\left(1+|c_2|\right)\cdot\Theta^{(0)}_2\le C^G_{\beta}\cdot e^{\mathfrak{e}},\\
|c_2|+|c_4|^{\frac{1}{2}}\le C^G_{\beta}\cdot e^{-\mathfrak{e}_2}.
\end{gather*}
Moreover $c_4-3c_2^2=b_{\eta,R^{*}(e)}(0)[\tau_v\beta]$, and
\begin{gather*}
\left|c_4-3c_2^2-b\left(dd^{*}\beta\right)\right|\le C^G_{\beta}\cdot\left(e+\eta^{\frac{1}{2}}\right),\\
\limsup_{N\to\infty}\left|c_4(N,e,\beta)-3c_2(N,e,\beta)^2-b\left(dd^{*}\beta\right)\right|\le C^G_{\beta}\cdot e^{\sigma}.
\end{gather*}
(Part~\ref{it:free-coefficient-3-2} of Theorem~\ref{thm:free-coefficient-3}, proved in Section~\ref{sec:proof-theorem}.)

\paragraph{Estimate~9 (mixtures and three tests).}
Let $\mu$ be a translation-invariant law on $\mathcal{S}'$ that satisfies~\eqref{eq:setup-G}, let $(\Xi,\nu,(\mu_{\xi}))$ be an ergodic decomposition of $\mu$ in the sense of Definition~\ref{def:ergodicity-ergodic-decompositions}, and let $\Xi_0$ be the set of full $\nu$-measure of Corollary~\ref{cor:properties-almost-every}; then, for every real $f\in\mathcal{S}$ and every $x\in\mathbb{R}^3$, one has $\int_{\Xi_0}\mathrm{d}\nu(\xi)\,\mathrm{Ext}^{\xi}_x(f)\ge\mathrm{Ext}^{\mu}_x(f)$ and $\mathrm{Ext}^{\xi}_x(f)\le\left(\frac{3}{2}\exp(2)+6\right)\cdot\|f\|_2^4$ for $\xi\in\Xi_0$, and if $\delta_{\mathrm{Ext}}:=\mathrm{Ext}^{\mu}_x(f)>0$, then
\[
\nu\left(\left\{\xi\in\Xi_0:\ \mathrm{Ext}^{\xi}_x(f)\ge\frac{\delta_{\mathrm{Ext}}}{2}\right\}\right)\ge\frac{\delta_{\mathrm{Ext}}}{2\,\|f\|_2^4}\cdot\left(\frac{3}{2}\exp(2)+6\right)^{-1},
\]
every $\xi$ in this set satisfies $\max\left\{\left|\kappa_4^{\xi}(f)\right|,\left|\kappa_4^{\xi}(f+\tau_xf)\right|,\left|\kappa_4^{\xi}(f-\tau_xf)\right|\right\}\ge\delta_{\mathrm{Ext}}/5$ and the law $\mu_{\xi}$ is not Gaussian, and $\mu$ is not a mixture of translation-invariant Gaussian laws that satisfy~\eqref{eq:setup-G}. Let further $(\beta,x_0)$ be data of Section~\ref{sec:probes}, with $f$ and the tests $g_0=f$, $g_{+}$, $g_{-}$ of that section, let $e>0$ and $\sigma'>0$, and let
\[
\left|\kappa_4^{\mu}(g)-e^4\cdot b(g)\right|\le C_g\cdot e^{4+\sigma'}\qquad\left(g\in\{g_0,g_{+},g_{-}\}\right)
\]
with numbers $C_{g_0},C_{g_{+}},C_{g_{-}}<\infty$ such that $\left(2C_{g_0}+\frac{1}{4}\left(C_{g_{+}}+C_{g_{-}}\right)\right)\cdot e^{\sigma'}\le\frac{1}{4}b(f)$; then $\mathrm{Ext}^{\mu}_{x_0}(f)\ge\frac{1}{2}e^4\cdot b(f)$, $\int_{\Xi_0}\mathrm{d}\nu(\xi)\,\mathrm{Ext}^{\xi}_{x_0}(f)\ge\frac{1}{2}e^4\cdot b(f)$,
\[
\nu\left(\left\{\xi\in\Xi_0:\ \mathrm{Ext}^{\xi}_{x_0}(f)\ge\frac{1}{4}e^4\cdot b(f)\right\}\right)\ge\frac{e^4\cdot b(f)}{4\,\|f\|_2^4}\cdot\left(\frac{3}{2}\exp(2)+6\right)^{-1},
\]
where $\frac{3}{2}\exp(2)+6<17.09$, and for every $\xi$ in this set the law $\mu_{\xi}$ is not Gaussian, with $\max\left\{\left|\kappa_4^{\xi}(g)\right|:\ g\in\{g_0,g_{+},g_{-}\}\right\}\ge\frac{1}{10}e^4\cdot b(f)$. (Theorem~\ref{thm:non-gaussian-components}; Theorem~\ref{thm:non-gaussian-components-2}, whose hypothesis is~\eqref{eq:fourth-cumulants-4g} and whose condition on $x_0$ is~\eqref{eq:setup-1}; that a Gaussian law with~\eqref{eq:setup-G} is centred and has vanishing fourth cumulants: Lemma~\ref{lem:moment-bounds}(a), Section~\ref{sec:cumulants-gaussian-laws}.)

\paragraph{Estimate~10 (zero-harmonic law).}
Let $(N,T)$ be a cutoff, $n=R/\eta$, and let $q>0$ and $\boldsymbol{\mu}>0$ be such that $\tau_{*}:=n\cdot\operatorname{arcsinh}\left(\frac{\boldsymbol{\mu}}{12}\right)\ge1$; then, with $\boldsymbol{\delta}_n:=110\,n^3\cdot\exp\left(-\tau_{*}\right)$, the harmonic-averaged law $E^{\mathrm{av}}_n$ has, with respect to the zero-harmonic law of Section~\ref{sec:lattice-laws} with the same $q$ and $\boldsymbol{\mu}$, a density with values in $\left[\exp\left(-4\boldsymbol{\delta}_n\right),\exp\left(4\boldsymbol{\delta}_n\right)\right]$, and for every Borel function $\Phi\colon C^2(n)\to\mathbb{C}$ with $|\Phi|\le1$ the expectations of $\Phi$ under these two laws differ by at most $\exp\left(4\boldsymbol{\delta}_n\right)-1$. (Theorem~\ref{thm:insensitivity-twist}; the two laws: Section~\ref{sec:wilson-operator-lattice}.)

\paragraph{Deduction of the Main theorem.}
Part~I of Theorem~\ref{thm:main-theorem} is established in Section~\ref{sec:admissibility-limit-points}. By Estimates~1 and~2, the law of the closed two-cochain $F$ under either lattice law satisfies \eqref{eq:continuum-infinite-Ad1}--\eqref{eq:continuum-infinite-Ad3} of Estimate~4: for the harmonic-averaged law through Corollary~\ref{cor:averaged-laws-admissible}, for the all-flux law with \eqref{eq:continuum-infinite-Ad2} from Estimate~1 and \eqref{eq:continuum-infinite-Ad1}, \eqref{eq:continuum-infinite-Ad3} from Estimate~2. Hence Estimate~4 yields (a), (P1), (P2), (P4)--(P8) and (P3) apart from the invariance under $F\mapsto-F$; this invariance follows from Estimate~2 through (P8) and Lemma~\ref{lem:exponentials-determination-law}(b) of Estimate~5. Since (P1), (P5) and (P6) give the hypothesis~\eqref{eq:setup-H} of Section~\ref{sec:laws}, Estimate~5 yields (c), with $\mu_{\xi}\left(\{F\in\mathcal{S}':dF=0\}\right)=1$ by (P4) and property~5; Remark~\ref{rem:repeated-limits}, part of Estimate~4, proves (d). Estimates~7 and~8 form Theorem~\ref{thm:free-coefficient-3}, proved in Section~\ref{sec:proof-theorem} with Estimates~2 and~6 from the results of Sections~\ref{sec:renormalisation-flow-tuned} and~\ref{sec:four-jets-one}. Part~II is proved in Section~\ref{sec:proof-part-ii-2}; its statements on $m_{0,N}(e)$ are contained in Estimate~7. In Steps~0--2, for $0<e\le e_2$ and $N\ge N_0(e)$, Estimates~7 and~8 make Estimate~3 applicable to the probes $\beta_0,\beta_{+},\beta_{-}$, with $T^{*}=T^{*}(e)$ and the numbers $c_2,c_4$ of Estimate~8, and \eqref{eq:free-coefficient-9} of Estimate~3 together with Estimates~6, 7 and~8 gives
\[
\boldsymbol{\Gamma}^{\mathrm{nat}}\le C\cdot e^{-\frac{2a_L}{N^{*}}},\qquad\boldsymbol{\Theta}^{\mathrm{nat}}_4\le C\cdot e^{\mathfrak{e}},\qquad\left(1+|c_2|\right)\cdot\boldsymbol{\Theta}^{\mathrm{nat}}_2\le C\cdot e^{\mathfrak{e}},
\]
with $C$ depending only on the parameters, $\bar{m}$, $\beta$, $x_0$; so $e^4\cdot\boldsymbol{\Gamma}^{\mathrm{nat}}\cdot\boldsymbol{\Theta}^{\mathrm{nat}}_4\le C^2\cdot e^{4+\sigma}$. Step~3 obtains (ii) from \eqref{eq:free-coefficient-7} and \eqref{eq:free-coefficient-8} of Estimate~3 on the tori of side $R=h\cdot L^T=k\cdot R^{*}(e)$ of that estimate, with $k=L^{T-T^{*}(e)}$ and $T\ge T^{*}(e)$; Lemmas~\ref{lem:norm-bounds-lattice} and~\ref{lem:lattice-symmetries-cell} of Estimate~4, together with Estimates~1, 2 and~6, replace the variable $X$ of Estimate~3 by $X_{\eta,R}(g)$, and Estimate~8 relates $c_4-3c_2^2$ to the free coefficients in (ii). Step~4 takes (i) from Part~I, and (iii) by (P8) from (ii) and, for $S_2^{\mu}$, from \eqref{eq:free-coefficient-20} and \eqref{eq:free-coefficient-21}, since $\boldsymbol{\sigma}\to\|f\|_2$ as $\eta\to0$ by Estimate~6. In Step~5, Estimate~9 turns (iii) into (iv) and (v). Step~6 applies \eqref{eq:free-coefficient-10} of Estimate~3 to the translates of one cell, replaces $f^{(\eta)}$ by $f$ by means of Estimates~1, 2, 4 and~6, and obtains \eqref{eq:free-coefficient-23} by (P8); Estimate~5 and the identity of Step~2 of Section~\ref{sec:proof-cumulant-bounds} stated in Estimate~3, together with Estimate~8, turn \eqref{eq:free-coefficient-23} into (vi); (vii) follows in the same way, with Estimate~9 applied to the components. The number $e_2$ is chosen after $C_{\mathrm{kap}}$, so that \eqref{eq:free-coefficient-1}--\eqref{eq:free-coefficient-4} of Section~\ref{sec:constants-exponents-scales} hold. The corollaries are proved in Section~\ref{sec:proofs-corollaries-main}: Corollary~\ref{cor:ergodic-components}(a) from Estimates~4 and~5, Part~I(c) and (vi); Corollary~\ref{cor:ergodic-components}(b) from Estimate~5, with $\mu\left(\{F\in\mathcal{S}':dF=0\}\right)=1$ by property~\eqref{it:main-theorem-P4} of Part~I, and from \eqref{eq:free-coefficient-23} and Estimate~8, as in Step~6 of Section~\ref{sec:proof-part-ii-2}; Corollary~\ref{cor:zero-harmonic-law} from Estimate~10, Part~I(d) and Part~II.

\section{The renormalisation group scheme and its constants}\label{sec:renormalisation-group-scheme}

This section contains, for Sections~\ref{sec:polymer-activities-algebra}--\ref{sec:free-coefficient-chessboard}, the definitions of the renormalisation group scheme (Sections~\ref{sec:constants-notation-terminology}--\ref{sec:domains-norms-step}); the statements on the scheme, namely the properties of the operators and of the operations on activities (Section~\ref{sec:properties-scheme}) and the theorems on one step, on the operators, on the flow, on the jets and on the free coefficients (Section~\ref{sec:one-step-operators}); the constants, which are fixed in the first eleven of the twelve stages of the order of choice, and the conditions imposed on them (Section~\ref{sec:constants}); and the classical results that are quoted (Section~\ref{sec:classical-results}). Apart from a few short deductions, the section contains no proofs: each lemma, proposition and theorem is followed by a sentence that names the place of its proof. Where the proof requires a fuller statement or another notation, the section that contains it states a result of its own, and the statement of this section is said, at its place, to be a part or a consequence of that result; otherwise the statement is proved in the form given here, the later section either stating it again with its proof or heading its argument as the proof of the statement of this section. The proofs are distributed as follows.
\begin{itemize}
\item Section~\ref{sec:polymer-activities-algebra}: the statements of Sections~\ref{sec:norms-polymers-operations} and~\ref{sec:hypotheses-propositions-activities}.
\item Section~\ref{sec:algebraic-gaussian-identities}: the statements of Sections~\ref{sec:identities-exact-operators} and~\ref{sec:hypotheses-identities-exact}.
\item Section~\ref{sec:operators-photon-field}: the statements of Sections~\ref{sec:properties-exact-photon} and~\ref{sec:photon-operators}; in Section~\ref{sec:weakened-photon-operators}, the properties of Section~\ref{sec:properties-systems-primitives} for the two photon families.
\item Section~\ref{sec:operators-fermion-field}: the statements of Sections~\ref{sec:fermion-operators-bounds} and~\ref{sec:fermion-operators}; in Sections~\ref{sec:weakened-walk-expansions}--\ref{sec:level-zero}, the properties of Section~\ref{sec:properties-systems-primitives} for the three fermion families.
\item Section~\ref{sec:renormalisation-group-step}: the theorem of Section~\ref{sec:one-step-large}; in Section~\ref{sec:elementary-properties-field}, the lemmas of Sections~\ref{sec:domains-bond-functions} and~\ref{sec:pairs-sites-holder} other than Lemma~\ref{lem:external-connections}.
\item Section~\ref{sec:range-running-masses}: the statements of Sections~\ref{sec:range-masses} and~\ref{sec:mass-range}.
\item Section~\ref{sec:renormalisation-flow-tuned}: the statements of Section~\ref{sec:trajectories-pair-bounds}; in Section~\ref{sec:projection-primitives}, the properties of Section~\ref{sec:properties-systems-primitives} for the systems of primitives used in the flow, property~\ref{it:renormalisation-group-18} and Proposition~\ref{prop:rows-component}.
\item Section~\ref{sec:four-jets-one}: the statements of Section~\ref{sec:probes-jets-reference}; in Section~\ref{sec:stop-level}, the argument that yields Lemma~\ref{lem:external-connections}.
\item Section~\ref{sec:free-coefficient-chessboard}: the statements of Section~\ref{sec:free-coefficients-bare}.
\end{itemize}
Section~\ref{sec:lattice-measures-laws} defines once more, with the directions labelled $0$, $1$, $2$, the Wilson operator of Section~\ref{sec:spinors-wilson-determinant} and the all-flux and harmonic-averaged laws of Section~\ref{sec:lattice-laws}. In Sections~\ref{sec:polymer-activities-algebra}--\ref{sec:free-coefficient-chessboard}, stage~$n$ is the $n$-th item of the list of Definition~\ref{def:order-choice}, a property cited by its number is the item with that number in the list of Section~\ref{sec:properties-systems-primitives}, and the hypotheses cited by their numbers are those of Definitions~\ref{def:hypothesis-one-step}, \ref{def:hypothesis-flow} and~\ref{def:hypothesis-reference-torus}.

\subsection{Constants, notation and terminology}\label{sec:constants-notation-terminology}

\begin{definition}[Constants]\label{def:constants}
Among the quantities fixed in the stages of Definition~\ref{def:order-choice}, the number $\bar{m}$, the probe $\beta$ and the point $x_0$, fixed in stage~\ref{it:order-choice-11}, and the charge $e$, the side $R$ and the integer $N$, fixed in stage~\ref{it:order-choice-12}, are not constants; the other numbers fixed in stages~\ref{it:order-choice-1}--\ref{it:order-choice-11} are called \emph{constants}, a constant fixed in stage~$n$ being a \emph{constant of stage~$n$}. A constant is fixed once and for all: either it is chosen at its stage under conditions that involve only quantities fixed before it in the order of choice of Section~\ref{sec:order-choice}, or it is defined from such quantities. It depends on no quantity fixed after it; in particular, neither a constant nor one of $\bar{m}$, $\beta$ and $x_0$ depends on $e$, $R$ or $N$. A number that is defined at its place by a formula in constants of stages~\ref{it:order-choice-1} to~$n$ is a constant of stage~$n$, also where Definition~\ref{def:order-choice} does not name it. For instance, the constants $C_V$ and $a$ of~\eqref{eq:operators-fermion-4} and $C_{\mathrm{Gam}}$ of~\eqref{eq:operators-fermion-8} are functions of $(b,\alpha,L)$, the numbers $C_K$ and $C_W$ of Section~\ref{sec:constants-3} are functions of constants of stages~\ref{it:order-choice-1}--\ref{it:order-choice-7}, and the numbers $z_2$ and $M_L$ of Section~\ref{sec:constants-3} are constants of stage~\ref{it:order-choice-11}. Standing alone, the phrase \emph{the constants} denotes the constants of stages~\ref{it:order-choice-1}--\ref{it:order-choice-10} where a number is said to depend on the constants, to be a function of them, to be determined by them or to be formed from them, and in the expressions \emph{condition on the constants} and \emph{requirement on the constants}; in this sense, a constant of stage~\ref{it:order-choice-11} depends only on the constants and on $\bar{m}$, $\beta$ and $x_0$. At other places the context fixes its meaning; for instance, a function orthogonal to the constants is one orthogonal to the constant functions. The word \emph{constant} is also used for a number written with an index or with arguments, such as the constant $C_a$ of Theorem~\ref{thm:free-coefficients-lattice}(a) or the numbers $e_1(\beta)$ and $C^G_{\beta}$ of Theorem~\ref{thm:bare-mass-shifts}, and for a number that is said to depend on certain objects, such as a constant $c$ that depends on a function $\Phi$. Such a number is a function of the objects so named and of the constants. It depends on $\bar{m}$, $\beta$, $x_0$ or a constant of stage~\ref{it:order-choice-11} only where its place names this dependence, as Theorem~\ref{thm:bare-mass-shifts} names $\beta$ and $\bar{m}$ for $e_1(\beta)$ and $C^G_{\beta}$, and on $e$, $R$, $N$ or a quantity of the following list only where its place names that quantity. No constant fixed in one of the stages~\ref{it:order-choice-1}--\ref{it:order-choice-10} depends on any of the following quantities:
\begin{itemize}
\item the integer $N$, where $\eta=L^{-N}$;
\item the level $k$;
\item the side $l$ of the torus;
\item the twist $\theta$ (Section~\ref{sec:lattice-measures-laws}), and with it the twist $c^{(k)}_{\theta}$ of level $k$ (Definition~\ref{def:twist-external-connection});
\item the running mass $\mathfrak{m}\in\left(0,\mathfrak{m}_0\right]$;
\item the real pair $\left(m_{\mathrm{in}},\boldsymbol{\varepsilon}_{\mathrm{in}}\right)$ of Definition~\ref{def:notation};
\item the external connection $a$ (Definition~\ref{def:twist-external-connection}), and the connection $\mathscr{C}\in\mathfrak{X}_k$ in the second slot (Section~\ref{sec:slots-connection-covering}), which along the flow is $a^{[k]}$ at the level $k$ (Definition~\ref{def:twist-external-connection}) and which in the statements on a single step is an arbitrary element of $\mathfrak{X}_k$ (Definition~\ref{def:domains});
\item the number $\bar{m}\in(0,1]$.
\end{itemize}
A constant fixed in stage~\ref{it:order-choice-11} may be a function of $\bar{m}$, but of none of the other quantities in this list.
\end{definition}

\begin{definition}[Sup norm; distance on a torus]\label{def:sup-norm-distance}
For $x=\left(x_0,x_1,x_2\right)\in\mathbb R^3$, the symbol $|x|$, written without an index, is the sup norm $\max_{\mu}\left|x_{\mu}\right|$; thus $|x-x'|$ is the sup distance of two points $x,x'\in\mathbb R^3$. Let $l>0$. On the torus $\mathbb R^3/l\mathbb Z^3$, and on each of its subsets, in particular on the lattice tori, $d(\cdot,\cdot)$ denotes the sup distance
\[
d(x,x')=\min_{z\in\mathbb Z^3}\left|\tilde x-\tilde x'+l\cdot z\right|,
\]
where $\tilde x,\tilde x'\in\mathbb R^3$ are representatives of $x$ and $x'$; the minimum is attained, and it does not depend on the choice of the representatives. For two points $x,x'$ of the torus, or of one of its subsets, $|x-x'|$ denotes the sup norm of a shortest representative of the class $x-x'\in\mathbb R^3/l\mathbb Z^3$:
\[
|x-x'|:=\min\left\{|w|:\ w\in\mathbb R^3,\ w+l\mathbb Z^3=x-x'\right\}=d(x,x'),
\]
the equality holding because the representatives of $x-x'$ are the vectors $\tilde x-\tilde x'+l\cdot z$ with $z\in\mathbb Z^3$.
\end{definition}

\begin{definition}[Notation]\label{def:notation}
\begin{enumerate}
\item[(a)] \emph{Charge, exponential, scales, twist, distances.} The charge is denoted $e$. The exponential function, also in phase factors $\exp(\pm i(\cdot))$, is written $\exp$, and Euler's number is $\exp(1)$. The scales of Section~\ref{sec:running-quantities-tori} are $h$ and $h'$. The \emph{twist} is the parameter $\theta\in\mathbb R^3$ introduced in Section~\ref{sec:lattice-measures-laws}, and the twist of level $k$ is the constant real connection $c^{(k)}_{\theta}$ of Definition~\ref{def:twist-external-connection}, whose value on every positively oriented bond of direction $\mu$ is $\frac{\theta_{\mu}}{e_k\cdot l_k}$; at the level $0$ it is the one-cochain $c^{(0)}_{\theta}=H_{\theta}$ of Section~\ref{sec:lattice-measures-laws}. Written without an index, $|\cdot|$ is the sup norm, and $d(\cdot,\cdot)$ is the sup distance on a torus, both as in Definition~\ref{def:sup-norm-distance}.

\item[(b)] \emph{Mass parameters.} Standing alone, $m_0=m_0(L)$ is the integer of Definition~\ref{def:mass-bound-stop}, that is, the least integer $j\ge0$ with $L^{-j}\le\mathfrak{m}_0$; it is the integer for which the stop scale $h'$ equals $h\cdot L^{-m_0}$. The symbol $m_0$ has another meaning at four places, at each of which the family to which it belongs is introduced. First, in Section~\ref{sec:flow-fixed-torus}, in the proofs of Propositions~\ref{prop:existence-trajectory} and~\ref{prop:pair-vanishing-terminal}, $m_0$ is the term of index $k=0$ of a sequence $(m_k)_{0\le k\le k^{*}}$ of mass parameters, $k^{*}$ being the stop level of Definition~\ref{def:tori-stop-level}: for the trajectory of a real pair it is $m_{\mathrm{in}}$, in agreement with Section~\ref{sec:flow}, where the symbol $m_k$ stands for $m_{\mathrm{in}}$ when $k=0$; for the fixed point $(m_k,E_k)_{0\le k\le k^{*}}$ of the map $\mathrm{Fx}$ in the proof of Proposition~\ref{prop:pair-vanishing-terminal} it is the number by which $m_{\mathrm{in}}$ is defined there. Second, in the proof of Lemma~\ref{lem:dependence-twist-uniform}, $m_0$ is the component of index $0$ of an integer vector $m=(m_0,m_1,m_2)\in\mathfrak{M}$. Third, in the proof of Lemma~\ref{lem:integral-paths}, $m_0=0$ is the first term of the sequence $(m_t)$ of integer cochains defined there by $m_{t+1}=m_t+k_t$. Fourth, in the proof of assertion~(iii) of Proposition~\ref{prop:trace-formula-bloch}, $m_0$ is the term of index $j=0$ of the family $(m_j)$ of integer functions on the plaquettes of the slice that form the second components of the variables $x_j=(A_j,m_j)$ introduced there, and $m_0(p)$ is its value at a plaquette $p$. Real pairs in general are denoted $\left(m_{\mathrm{in}},\boldsymbol{\varepsilon}_{\mathrm{in}}\right)$, and Definition~\ref{def:tuned-pair} gives the tuned pair $P_N(e)=\left(m_{0,N},\boldsymbol{\varepsilon}_{0,N}\right)$. The right end point of the mass range $\left(0,\mathfrak{m}_0\right]$ is always called $\mathfrak{m}_0$. The running mass of level $k$ is the number $\bar{m}\cdot a_k$, where $a_k=L^{k-N}$ (Section~\ref{sec:running-quantities-tori}); it is denoted $\mathfrak{m}_k$ for $k\ge1$, whereas at level $0$ it is written $\bar{m}\cdot\eta$.

\item[(c)] \emph{Exponents and radii.} The H\"older order of stage~\ref{it:order-choice-2} is denoted $\alpha$. The parameter polydisk $\mathrm{Pd}_k$ of Definition~\ref{def:parameter-cubes-polydisks} has the radius $M^{\alpha_0}$, where $\alpha_0$ is the exponent chosen in stage~\ref{it:order-choice-5}. The extraction coefficients of Section~\ref{sec:local-parts-adjustment} are written $\boldsymbol{\alpha}_0(E,X)$ and $\boldsymbol{\alpha}_2(E,X)$. In the definition of $r_k$ (Section~\ref{sec:running-quantities-tori}) the exponent is the integer $r$. The jet radius of Section~\ref{sec:probes-jets-reference} is $\mathbf{r}$, and $h_G:=1+g$, with the number $g$ of Section~\ref{sec:constants-3}.

\item[(d)] \emph{Directions.} In Sections~\ref{sec:renormalisation-group-scheme} and~\ref{sec:polymer-activities-algebra}--\ref{sec:four-jets-one} the directions carry the labels $\mu\in\{0,1,2\}$, except in the definition of the free coefficients at the beginning of Section~\ref{sec:free-coefficients-bare}, which repeats the definition of Section~\ref{sec:free-fourth-order} with the Dirac matrices $\gamma_{\mu}$ of Section~\ref{sec:letters-norms-dirac} and therefore labels the directions $1,2,3$; there $\gamma_{\mu}$ is the matrix $\gamma_{\mu-1}$ of Section~\ref{sec:dirac-matrices-lattice}. Section~\ref{sec:free-coefficient-chessboard} labels them $1,2,3$, except in Sections~\ref{sec:factorisation-time-slices} and~\ref{sec:trace-formula-chessboard}, where the labels are $0,1,2$. Each of Sections~\ref{sec:gaussian-bounds-curvature}--\ref{sec:wilson-operators-block} fixes its own labelling of the directions.

\item[(e)] \emph{The indices $0$ and $*$.} Since, as a number, $e_0$ is reserved for the ceiling fixed in stage~\ref{it:order-choice-9}, the running charge $e\cdot\eta^{\frac{1}{2}}$ of level $0$ is called $q$, and the level-$0$ quantities $v_0$, $r_0$, $p_{0,0}$, $w^{\mathrm{m}}_0$, $w^{\mathrm{E}}_0$ of Section~\ref{sec:running-quantities-tori} are formed with $q$. A formula stated for an arbitrary level $k$ is read at $k=0$ with $q$ in place of $e_k$, never with the ceiling. In the same way, $p_0$ is kept for the exponent of $v_k$ in $p_{0,k}=v_k^{p_0}$; at level $0$ the number $p_k=v_k^p$ therefore appears as $v_0^p$. The operator $H_0$ of level $0$ is the identity, whereas in $H_0=0$ the symbol denotes the one-cochain $H_{\theta}$ of Section~\ref{sec:lattice-measures-laws} at $\theta=0$; the number that occurs in condition~\eqref{eq:renormalisation-group-A7} is called $\mathbf{H}_0$. The stop level is the level $k^{*}=k^{*}(R)$ of Definition~\ref{def:tori-stop-level} (see also Definition~\ref{def:terminology}(6)). Attached to a running quantity of Section~\ref{sec:running-quantities-tori} or to another object of the scheme that carries the level $k$ as an index, a subscript $*$ marks, as in Definition~\ref{def:tori-stop-level}, the value at $k=k^{*}(R)$, the symbol $n_{*}$ excepted; for instance, $e_{*}=e_{k^{*}}$, $l_{*}=l_{k^{*}}$, $r_{*}=r_{k^{*}}$, $v_{*}=v_{k^{*}}$, $b_{*}=b_{k^{*}}$ and $E_{*}=E_{k^{*}}$. The symbol $n_{*}$ always stands for $R^{*}/\eta$, so that the site count of the unit lattice at the stop level appears as $l_{*}^3$. Attached to a map, such as a lattice isometry $g$, the isometry $\mathsf{g}_c$ of Section~\ref{sec:cells-reference-tori}, the translation by a vector $v$, or a matrix $r$ or $w$, a subscript $*$ denotes the push-forward by this map: on cochains, $g_{*}:=\left(g^{-1}\right)^{*}$, where $g^{*}$ is the pull-back (Section~\ref{sec:cells-reference-tori}), and, for a matrix $r$, $r_{*}\beta$ is the push-forward of $\beta$ by $r$. Thus $v_{*}$ and $r_{*}$ denote, at some places, the numbers $v_{k^{*}}$ and $r_{k^{*}}$ and, at others, operators of push-forward. Further symbols with a subscript $*$ are declared at their places and are not values at the stop level; among them are $x_{*}$, $\mathscr{A}_{*}$, $S_{*}(\alpha)$, $C_{*}$, $\tau_{*}$, $f_{*}$ and $q_{*}$. The number $j_{*}$ of Proposition~\ref{prop:depth-stop-level}(b) is a depth, not a level; by Proposition~\ref{prop:depth-stop-level}(c), $j_{*}=K-k^{*}(R)$ for $N\ge N_{\mathrm{min}}(R,e)$, with the level $K$ of Section~\ref{sec:running-quantities-tori}, so that its subscript $*$ also refers to the stop level. Besides the sense $E_{k^{*}}$, the symbol $E_{*}$ has the sense $E_1$ declared in Section~\ref{sec:transfer-larger-cubic}, where also $\bar{E}_{*}:=\bar{E}_1$; and in Example~\ref{ex:prescribed-fourth-cumulant}, $b_{*}$ is a given positive number, not $b_{k^{*}}$.

\item[(f)] \emph{Further fixed symbols.} The domain of Section~\ref{sec:domains-bond-functions} is $R_k$, whereas the orthogonal projection that occurs in the Landau condition of Section~\ref{sec:averaging-fluctuations-gauge} is $\mathscr{R}_k$. The class of sides (Section~\ref{sec:lattice-geometry}) is $\mathrm{Cl}$, and the cells of Section~\ref{sec:cells-reference-tori} form the set $\mathrm{Cel}$. With the charge $e$ as its argument, $N_0(e)=N_{\mathrm{min}}\left(R^{*}(e),e\right)$ is the threshold of Theorem~\ref{thm:bare-mass-shifts}, and the threshold is never written without this argument. The free ratio of Section~\ref{sec:probes-jets-reference} is a function of a complex variable; it is written $N_0(t)$ or $N_0(\cdot)$, with other values of the variable, as in $N_0(0)$, $N_0(-t)$ and $N_0(t_c)$, and with its derivatives $N_0^{(m)}(0)$, and, where the function itself is meant, it is written $N_0$, without an argument, as in $\log N_0$, $\frac{1}{N_0}$, $N_0-1$ or the four-jet of $N_0$. In Section~\ref{sec:trace-formula-chessboard}, and only there, $N_0$ without an argument is the number of slabs, that is, the number $N_{\mu}$ of Definition~\ref{def:decomposition-torus-cells} for $\mu=0$; neither the free ratio nor the threshold occurs in that section. Two further meanings are declared where they occur: in Sections~\ref{sec:sums-aliases} and~\ref{sec:bounds-complex-strip}, the numbers $N_0$, $N_1$, $N_2$ are the values at $p$ of the functions $N_{\mu}(p)=U_0^2\cdot u_{\mu}^2+D_0\cdot X_{\mu}(p)$ of the momentum $p$; and in Step~1 of the proof of Proposition~\ref{prop:existence-conditional-expectations}(a), $N_0$ is a null set. The axial functional of Section~\ref{sec:averaging-fluctuations-gauge} is denoted $\tau$; with a point $x$ as index, $\tau_x$ is a translation; and $\tau_s$ denotes the product $c_1\cdot l_{*}$, with $c_1$ as in Section~\ref{sec:constants-3}. In Sections~\ref{sec:renormalisation-group-scheme} and~\ref{sec:polymer-activities-algebra}--\ref{sec:four-jets-one}, a number written $C_0$, without arguments, is the constant of stage~\ref{it:order-choice-7}, with the exception of three places, at each of which $C_0$ is a datum of the statement and is declared there. In clause~(iv) of item~\eqref{it:weighted-coefficient-norms-N7} of Proposition~\ref{prop:weighted-coefficient-norms}, and in clause~(iv) of the paragraph \emph{Norms on $\mathrm{Gr}[\Psi]$} of Theorem~\ref{thm:polymer-activities-1} together with part~(iv) of its proof, $C_0>0$ is an arbitrary number, which bounds the row sums $\left|H^{(1)}\right|_{1,\infty}$ and $\left|H^{(2)}\right|_{1,\infty}$ of the operators $H^{(1)}$ and $H^{(2)}$ of that clause. In Theorem~\ref{thm:unit-lattice-bounds} and its proof, $C_0$ is any number with $C_0\ge\max\left(C_0^{\mathrm{unit}},C_F\right)$, and it enters the weight $a_2=C_0^{-1}\cdot e_k^{-\frac{1}{4}+\varepsilon}$ of that theorem; here $C_F$ and $\gamma_F$ are the constants of Propositions~\ref{prop:decay-fine-propagator}--\ref{prop:bounds}, and $C_0^{\mathrm{unit}}=4b_{+}+16\,C_F\cdot c(\gamma_F)$ is the number of Section~\ref{sec:absolute-constants}. The paragraph that follows the proof of Theorem~\ref{thm:unit-lattice-bounds} takes this number $C_0$ equal to the constant of stage~\ref{it:order-choice-7} and deduces Proposition~\ref{prop:bounds-unit-lattice}, in which $C_0$ is the constant of stage~\ref{it:order-choice-7}. The choice is admissible: there $\left(C_F,\gamma_F\right)=\left(C^F,\gamma^F\right)$, and the constant of stage~\ref{it:order-choice-7} is the largest of three numbers, two of which are $C^F=C_F$ and $4b_{+}+16\,C^F\cdot c\left(\gamma^F\right)=C_0^{\mathrm{unit}}$. The bilinear form $C_0(\cdot,\cdot)$ of Section~\ref{sec:lattices-tests-function}, with its two arguments, is the covariance of the free Maxwell field.

\item[(g)] \emph{Letters with several meanings.} For each of the following letters the meaning is determined by the context: by the section in which the letter occurs, by the presence of an argument, by the kind of object, number or vector, that it denotes, or by a declaration made at the place where the letter is used, that is, by a defining equation, such as $\psi^{\#}:=H^{\#}_k(\mathscr{T})\Psi$, or by the membership of the letter in a family declared there, such as that of $c_1$ in the family of the numbers $c_m:=\phi_0^{(m)}(0)$.
\begin{itemize}
\item In Section~\ref{sec:renormalisation-group-step}, $C_W$ stands for the number $4C_w^2\cdot c_{\gamma}$ (Section~\ref{sec:notation}); in Section~\ref{sec:probes-jets-reference} and in Section~\ref{sec:four-jets-one} it stands for the constant defined in Section~\ref{sec:constants-3}.
\item In Section~\ref{sec:renormalisation-group-step}, $\kappa_1,\dots,\kappa_4$ are the rates of Section~\ref{sec:notation}. Followed by an argument, $\kappa_n(\cdot)$ is a cumulant: for a random variable $X$ and a law, or an expectation, $E$, the symbol $\kappa_m(X;E)$ denotes the cumulant of order $m$ of $X$ under $E$.
\item The letter $\vartheta$ denotes a remainder in Theorem~\ref{thm:estimates-one-torus} and throughout Section~\ref{sec:renormalisation-flow-tuned}, but the rate defined in Section~\ref{sec:constants-3} in Section~\ref{sec:probes-jets-reference} and in Section~\ref{sec:four-jets-one}.
\item As numbers written without an argument, $e_1$ and $e_2$ always denote the constants of stages~\ref{it:order-choice-10} and~\ref{it:order-choice-11}, and never running charges; with a probe $\beta$ as its argument, $e_1(\beta)$ is the number of Theorem~\ref{thm:bare-mass-shifts}, which depends on $\beta$. As a vector, $e_{\mu}$ is the unit vector of the direction labelled $\mu$, in the labelling of the directions used at that place; this applies to $e_0$ as well.
\item The standard Gaussian measure on $V_k$ (Section~\ref{sec:averaging-fluctuations-gauge}) is denoted $\mu_I$; with the argument $W$, $\mathrm{d}\mu_I(W)$ is the Grassmann Gaussian functional of Proposition~\ref{prop:weighted-coefficient-norms}.
\item The generators $\psi$, $\chi$ of Section~\ref{sec:grassmann-algebras-norms} form the pair $\psi^{\#}=(\psi,\chi)$, and the symbol $\psi^{\#}$ denotes these generators, except where $\psi^{\#}:=H^{\#}_k(\mathscr{T})\Psi$ is declared at the place where the symbol occurs. This is the substitution of Section~\ref{sec:fermion-integrals-blocking}, by which the four fields $\psi=H_k(\mathscr{T})\Psi$, $\chi=\delta_{\alpha,\mathscr{T}}\psi$, $\bar{\psi}=\bar{H}_k(\mathscr{T})\bar{\Psi}$ and $\bar{\chi}=\delta_{\alpha,-\mathscr{T}}\bar{\psi}$ are inserted for the generators; there $\psi^{\#}$ is this value. The substitution is declared in this way in Section~\ref{sec:fermion-integrals-blocking}, in Section~\ref{sec:bounds-unit-lattice} (Theorem~\ref{thm:unit-lattice-bounds}(c) and its proof), in Section~\ref{sec:formula-one-renormalisation}, in Sections~\ref{sec:representation-stop-level} and~\ref{sec:integral-stop-level}, and in Section~\ref{sec:four-jets-one} from clause~\eqref{it:trajectory-fermionic-operators-SL6} of Proposition~\ref{prop:trajectory-fermionic-operators} on. Throughout Section~\ref{sec:renormalisation-group-step}, in particular in Section~\ref{sec:proof-clause-theorem}, and in Section~\ref{sec:comparison-steps-two}, the symbol $\psi^{\#}$ denotes the generators; the values that occur in Proposition~\ref{prop:step-external-connection}, in Theorem~\ref{thm:one-step-formula} and in Section~\ref{sec:proof-clause-theorem} carry a further mark, namely $\psi'^{\#}:=H^{\#}_{k+1}(\mathscr{T}')\Psi'$, $\psi^{0,\#}:=\left(\psi'^{\#}\right)_L$ and, in the proof of Theorem~\ref{thm:one-step-formula}, $\psi^{n,\#}$.
\item In conditions~\eqref{eq:renormalisation-group-19} and~\eqref{eq:renormalisation-group-21} and in Section~\ref{sec:operators-fermion-field}, $C_V$ is the constant introduced in Lemma~\ref{lem:perturbations-commutators-holder}; in Section~\ref{sec:four-jets-one} it is the constant introduced in Section~\ref{sec:constants-3}.
\item Within Section~\ref{sec:operators-fermion-field}, $C_{\mathrm{Gam}}$ is the constant~\eqref{eq:operators-fermion-8}; in stage~\ref{it:order-choice-11} and in Section~\ref{sec:free-coefficient-chessboard} it is the constant introduced in Step~1 of Section~\ref{sec:proof-part-ii-2}.
\item In stage~\ref{it:order-choice-7}, $C_B$ is the constant~\eqref{eq:operators-photon-19}; in Section~\ref{sec:proof-part-ii} it is the number $\max_{\beta'}C^G_{\beta'}$.
\item In Section~\ref{sec:constants-order}, in Proposition~\ref{prop:fluctuation-space}, in Lemma~\ref{lem:dependence-twist-uniform} and in Section~\ref{sec:constants-3}, the letter $c_1$ denotes a constant defined at that place. Further, $c_1,\dots,c_4$ denote the numbers $c_m:=\phi_0^{(m)}(0)$, $m=1,2,3,4$, the derivatives at $0$ of the four-jet $\phi_0$ of the free ratio, for which $c_1=c_3=0$; they occur in Theorems~\ref{thm:free-ratio}, \ref{thm:products-centred-derivatives} and~\ref{thm:bare-mass-shifts}, in Lemma~\ref{lem:free-wilson-ratio}(e) and in Proposition~\ref{prop:products-centred-derivatives}, and in the proofs of the last two. In Sections~\ref{sec:renormalisation-group-scheme} and~\ref{sec:four-jets-one} both numbers are written without an argument, and they are told apart as follows. Where $c_1$ stands together with the family $(c_m)$, which is then declared at that place, namely in the relation $c_1=c_3=0$ and in the list $c_1,\dots,c_4$, it is the member of this family. Standing alone, as in $e_{*}\cdot|A_1|_{\infty}\le c_1$, in $c_1\le\mathfrak{m}_{-}/48$ and in $\tau_s=c_1\cdot l_{*}$, it is the positive constant $\min\left\{\frac{1}{8},\frac{\mathfrak{m}_{-}}{48},\frac{1}{80\,\mathbf{S}}\right\}$ of Section~\ref{sec:constants-3}; this applies in item~(f); in Section~\ref{sec:probes-jets-reference}, in the paragraph on the other constants and the number $\mathfrak{m}_{-}$, where $c_1$ stands alone in the list of the eight numbers $\boldsymbol{\theta}_{\mathrm{CT}}$, $\vartheta$, $g$, $h_G$, $g_1$, $\boldsymbol{\alpha}$, $c_1$, $C_W$ of Section~\ref{sec:constants-3}, and in $\tau_s=c_1\cdot l_{*}$; and in Section~\ref{sec:four-jets-one}.
\item The external connection of Definition~\ref{def:twist-external-connection} is denoted $a$; in Section~\ref{sec:operators-fermion-field}, and in conditions~\eqref{eq:renormalisation-group-19} and~\eqref{eq:renormalisation-group-21}, the same letter stands for the constant defined in~\eqref{eq:operators-fermion-4}.
\item The fermion operators of Section~\ref{sec:exact-fermion-operators} carry a connection as their argument; where this is declared at the place, the argument is suppressed, and the operator is then written $D_k$, or $D_{*}$ at the level $k^{*}$. With a connection $A$ as its index, $D_A$ is the Wilson operator of Section~\ref{sec:exact-fermion-operators} on $T_k$, which at the level $0$ and for the side $n$ is the Wilson operator of Section~\ref{sec:lattice-measures-laws}. Without an argument, $D_0$ is $D_A$ at $A=0$, on the lattice on which $D_A$ is taken at the place where $D_0$ occurs; for instance, in $\det\left(D_0+\bar{m}\cdot\eta\right)$ this is the lattice of level $0$ and side $n$, in Section~\ref{sec:operators-fermion-field} it is the lattice of level $k$, for an arbitrary level $k\ge0$, and in the proof of Lemma~\ref{lem:fourier-decomposition} it is the lattice $\boldsymbol{\eta}\mathbb{Z}^3$ of that proof. With a connection $A$ as its argument, $D_0(A)=D_A+\mathfrak{m}$, where $\mathfrak{m}=\bar{m}\cdot\eta$, is the operator of level $0$ of Section~\ref{sec:exact-fermion-operators}, so that $D_0(0)=D_0+\bar{m}\cdot\eta$. Other senses of $D_0$ are declared at their places; among them, in Sections~\ref{sec:sums-aliases} and~\ref{sec:bounds-complex-strip} the symbol $D_0$ denotes the number $\hat{D}(p)$, which is its sense in item~(f), and in Section~\ref{sec:determinant-identities} it denotes the operator $R_z$ of that section, by the definition $D_0:=R_z$ in the case $k=0$ of Lemma~\ref{lem:block-determinants}(iii). Where $D_k$ stands as a set, as in $X\in D_k$ or in a sum over the polymers $X\in D_k$, it is the polymer set of Section~\ref{sec:lattice-geometry}; in particular, $D_0$ is also the polymer set of level $0$. Without a connection as argument, $T_k$ is the lattice and $B_k(\cdot)$ a block, both as in Section~\ref{sec:lattice-geometry}, while $H_k$ and $Q_k$ are the minimiser and the averaging operator of Section~\ref{sec:averaging-fluctuations-gauge}.
\item As letters, $L$ and $R$ denote numbers; in compositions such as $R\,\mathscr{E}$, $L\,E$ and $L\,R\,\mathscr{E}$ they denote the operations of Section~\ref{sec:local-parts-adjustment}.
\end{itemize}

\item[(h)] \emph{Words.} Several words carry technical meanings, fixed at the places indicated; where more than one meaning is listed, the context decides between them.
\begin{itemize}
\item Besides its ordinary meaning, the word \emph{pair} designates, in a technical sense, the real pairs $\left(m_{\mathrm{in}},\boldsymbol{\varepsilon}_{\mathrm{in}}\right)$; the data $(m,E)$ and $(m,\mathscr{E})$ of a step (Section~\ref{sec:spaces-activities} and Definition~\ref{def:renormalisation-step}); and the ordered pairs of sites, the elements of $\mathrm{Pz}_k$ (Definition~\ref{def:pairs-sites-holder}).
\item Admissibility is defined for pairs of sites relative to a polymer (Definition~\ref{def:pairs-sites-holder}), for unions $\tilde{X}$ of $M$-cubes (Definition~\ref{def:admissibility-flatness-agreement}), for families of laws (Definition~\ref{def:admissible-families}), and for the number $s$ chosen in stage~\ref{it:order-choice-8}; the \emph{admitted} placements are those of Definition~\ref{def:probes}.
\item The word \emph{cell} refers either to a cell of a lattice, that is, to a site, a bond, a plaquette or a cube as in Definition~\ref{def:barycentres}, or to a cell $\mathfrak{c}_c$ of the cell structure of Section~\ref{sec:cells-reference-tori}.
\item The \emph{probes} are the real forms $\beta\in C_c^{\infty}\left(\mathbb R^3;\Lambda^2\right)$ with $\beta\ne0$ (Definition~\ref{def:probes} and Section~\ref{sec:probes-discretisation}); the lattice probe attached to $\beta$ is the cochain $\hat{\mathfrak{a}}$ of Definition~\ref{def:probes}.
\item The word \emph{flat} has two different senses. A polymer is flat in the sense of Definition~\ref{def:admissibility-flatness-agreement}(b), see also Theorem~\ref{thm:projection-agreement-flat}(ii); a connection $A$ is flat if its curvature $dA$ vanishes, and the flat connection of Section~\ref{sec:averaging-blocks-flat} is the constant real connection $c$ defined there.
\item Besides the orthogonal projections, the covering maps, such as $\mathrm{pr}_i$ of Section~\ref{sec:two-tori} and $\mathrm{pj}$ of Definition~\ref{def:periodisation-projection}, are called \emph{projections}, and so is the projection of a system of primitives defined in Definition~\ref{def:periodisation-projection}.
\item The \emph{walks} are those of Definition~\ref{def:walks}.
\item The term \emph{slice data} has one meaning in Section~\ref{sec:slices-slice-operators} and another in Section~\ref{sec:slice-factorisation-determinant}; in the second meaning the slice data are those of Lemma~\ref{lem:slice-factorisation-determinant-2}.
\item The kernel $|O|(y,y')$ of Section~\ref{sec:fermionic-operators-stop} is called a \emph{majorant}.
\end{itemize}
\end{enumerate}
\end{definition}

\begin{definition}[Barycentres]\label{def:barycentres}
Whenever a function of the position is to be evaluated on a cell, that is, on a site, a bond, a plaquette or a cube, its value is taken at the barycentre of that cell. The covariance under reflections requires this convention. Moreover, the component $A_{\nu}(x)=A\left(x,x+\eta_ke_{\nu}\right)$ of a bond function $A$ on the fine lattice $T_k$ (Sections~\ref{sec:running-quantities-tori}, \ref{sec:lattice-geometry} and~\ref{sec:averaging-fluctuations-gauge}) is regarded as a function of the unoriented bond $\left\{x,x+\eta_ke_{\nu}\right\}$, and the difference quotient $\left(\partial_{\nu'}A_{\nu}\right)(x)$ as a function of the cell spanned by the two parallel bonds $\left(x,x+\eta_ke_{\nu}\right)$ and $\left(x+\eta_ke_{\nu'},x+\eta_ke_{\nu'}+\eta_ke_{\nu}\right)$; this cell is a plaquette if $\nu'\ne\nu$ and a segment of two bonds if $\nu'=\nu$, and its barycentre is $x+\frac{\eta_k}{2}\cdot\left(e_{\nu}+e_{\nu'}\right)$. The distance of two such cells with the same pair $(\nu',\nu)$ is the sup distance of their barycentres. On the unit lattice the same conventions hold with $1$ in place of $\eta_k$.
\end{definition}

\begin{definition}[Terminology]\label{def:terminology}
The following terms have, in all sections, the meanings fixed below; where a clause gives a term more than one meaning, it states in which sections, or for which objects, each of these meanings applies.
\begin{enumerate}
\item[(1)] \emph{Dependence on the values at a set of places.} Let $F$ be defined by a displayed formula, where $F$ is either a number, a kernel, an element of a Grassmann algebra or a row of an operator; then $F$ is a function of the objects appearing in this formula. Fix one of these objects and call it $u$. It may be a field, a connection, an operator, a family of operators, a family of parameters or a datum, and it takes values at places, which are rows, indices, bonds, sites, parameter cubes or polymers. Let $W$ be a set of places. The five phrases ``$F$ reads $u$ on $W$'', ``$u$ is read on $W$'', ``$u$ enters $F$ through its values on $W$ only'', ``$F$ depends on $u$ through its values on $W$ only'' and ``$u$ occurs in $F$ at the places of $W$ only'' are synonymous. Each of them states that, all other objects being held fixed, $F$ takes equal values at any two admitted values of $u$ which coincide on $W$. The admitted values of $u$ make up the domain that the statement in question specifies; when $u$ consists of primitives, they are the systems of primitives having the properties that this statement assumes.

\item[(2)] \emph{Operators and objects given by one formula on $\mathbb{L}_k$.} Let $v\in(L\cdot M)\mathbb{Z}^3$. As $(L\cdot M)\mathbb{Z}^3$ is contained in $\mathbb{Z}^3$, in $L\mathbb{Z}^3$ and in $M\mathbb{Z}^3$, the translation by $v$ maps each of the following grids onto itself: the fine lattice $\mathbb{L}_k$ (Section~\ref{sec:lattice-geometry}), the unit lattice $\mathbb{Z}^3$, the set of the $L$-blocks together with the basis elements $\Upsilon$ of Section~\ref{sec:averaging-fluctuations-gauge}, and the set of the $M$-cubes. In this clause an \emph{index set} is a set $I$ of objects of $\mathbb{L}_k$, called indices, that every translation by a vector of $(L\cdot M)\mathbb{Z}^3$ maps onto itself, each index $x\in I$ being given a finite set $\operatorname{supp}x\subset\mathbb{L}_k$, its support, such that $\operatorname{supp}(x+v)=\operatorname{supp}x+v$ for all $x\in I$ and all $v\in(L\cdot M)\mathbb{Z}^3$. Consequently the following sets are index sets: the sets of the cells (sites, bonds, plaquettes and cubes) of $\mathbb{L}_k$, of $\mathbb{Z}^3$ and of $L\mathbb{Z}^3$, the support of a cell being the set of its corners, where a bond function, which is odd under the reversal of bonds, is regarded as a function on the positively oriented bonds; the set of the basis elements $\Upsilon$; the set of the pairs $(y,x)$ with $y\in L\mathbb{Z}^3$ and $x\in B(y)$, at which the rows of $\tau$ stand (Section~\ref{sec:averaging-fluctuations-gauge}), with the support $\{y,x\}$; and the set of the pairs $(x,x')\in\mathbb{L}_k\times\mathbb{L}_k$ with $0<|x-x'|\le M$, with the support $\{x,x'\}$, which lie above the pairs of $\mathrm{Pz}_k$ (Sections~\ref{sec:grassmann-algebras-norms} and~\ref{sec:antiperiodicity}). For sites, bonds and basis elements these supports are those of Section~\ref{sec:lattice-geometry}, where these three types of objects are called indices. The index sets of this clause are thus more general than the indices of Section~\ref{sec:lattice-geometry}; this does not affect the definition of the indices given there and recalled in Sections~\ref{sec:notation} and~\ref{sec:two-tori}. An index may carry, in addition, a spinor index and the index $\boldsymbol{\omega}$ of Section~\ref{sec:grassmann-algebras-norms}, on which neither the translations nor the supports act. Let $l$ be a multiple of $L\cdot M$. As in Definition~\ref{def:periodisation-projection}, a function on the indices of the torus $T_k=\mathbb{L}_k/l\mathbb{Z}^3$ (Section~\ref{sec:lattice-geometry}) is regarded as an $l\mathbb{Z}^3$-periodic function on the corresponding index set of $\mathbb{L}_k$, and a spinor function on $T_k$ as an antiperiodic function $f$, that is, $f(x+l\cdot e_{\nu})=-f(x)$ for $\nu=0,1,2$, the translate of a pair being formed by translating both of its points; in condition~(iii) below the word spinors stands for all the antiperiodic functions on index sets. The translation $T_v$ acts on the functions on every index set by
\[
\left(T_vf\right)(x):=f(x-v).
\]
A linear map $O$ from the functions on an index set $I$ to the functions on an index set $I'$ is written with finite sums,
\[
(Of)(x)=\sum_{x'\in I}O(x,x')\cdot f(x')\qquad\left(x\in I'\right);
\]
the \emph{row} of $O$ at $x$ is the function $x'\mapsto O(x,x')$, its support is the set of the $x'$ with $O(x,x')\ne0$, and weights such as $\eta_k^3$ are counted into the entries $O(x,x')$, as in Section~\ref{sec:two-tori}. A bilinear form enters through its kernel, which is treated as an operator: the form $\langle\bar{\psi},\psi\rangle_{\square}$ has the kernel $\eta_k^3\cdot\mathbf{1}_{\{x=x'\}}\cdot\mathbf{1}_{\square}(x)$ times the unit matrix, the $M$-cube $\square$ taking the place of the set $Y$ below; and a constant $4\times4$ matrix is the operator that acts at each site. Let a class of sides be given, each of them a multiple of $L\cdot M$, let $I$ and $I'$ be index sets, and for every side $l$ of the class let $\mathrm{Op}$ be a linear map from the $l\mathbb{Z}^3$-periodic functions on $I$ to those on $I'$, or from the antiperiodic functions on $I$ to those on $I'$. The map $\mathrm{Op}$ may depend on a bond function $B$ on $T_k$ and on a set $Y$ that is a union of $M$-cubes of $T_k$; it is then written $\mathrm{Op}(B,Y)$. One says that $\mathrm{Op}$ is \emph{given by one formula on} $\mathbb{L}_k$ if to every bond function $\tilde{B}$ on $\mathbb{L}_k$ and to every union $\tilde{Y}$ of $M$-cubes of $\mathbb{L}_k$ there corresponds a linear operator $\mathrm{Op}_{\mathbb{L}}\left(\tilde{B},\tilde{Y}\right)$ from the space of all functions on $I$ to the space of all functions on $I'$, in such a way that the following three conditions hold, together with the requirement of locality stated after them.
\begin{enumerate}
\item[(i)] Each row of $\mathrm{Op}_{\mathbb{L}}\left(\tilde{B},\tilde{Y}\right)$ has finite support.
\item[(ii)] For every $v\in(L\cdot M)\mathbb{Z}^3$,
\[
T_v\circ\mathrm{Op}_{\mathbb{L}}\left(\tilde{B},\tilde{Y}\right)=\mathrm{Op}_{\mathbb{L}}\left(T_v\tilde{B},\tilde{Y}+v\right)\circ T_v.
\]
\item[(iii)] Let $l$ be a side of the class, and let $\tilde{B}$ and $\tilde{Y}$ be the $l\mathbb{Z}^3$-periodic extensions of $B$ and of $Y$ to $\mathbb{L}_k$. Then $\mathrm{Op}(B,Y)$ coincides with $\mathrm{Op}_{\mathbb{L}}\left(\tilde{B},\tilde{Y}\right)$ acting on the periodic functions, and on the antiperiodic functions in the case of spinors.
\end{enumerate}
The requirement of locality is the following: there is a number $\rho\ge0$ such that, for every bond function $\tilde{B}$ on $\mathbb{L}_k$, every union $\tilde{Y}$ of $M$-cubes of $\mathbb{L}_k$ and every row index $x\in I'$, first, $\mathrm{Op}_{\mathbb{L}}\left(\tilde{B},\tilde{Y}\right)(x,x')=0$ for all $x'\in I$ with
\[
\min\left\{|y-y'|:\ y\in\operatorname{supp}x,\ y'\in\operatorname{supp}x'\right\}>\rho,
\]
and, second, the row at $x$ reads $\tilde{B}$ on the bonds both of whose end points have sup distance at most $\rho$ from $\operatorname{supp}x$ and reads $\tilde{Y}$ on the sites at sup distance at most $\rho$ from $\operatorname{supp}x$, in the sense of clause~(1); that is, the rows at $x$ of $\mathrm{Op}_{\mathbb{L}}\left(\tilde{B},\tilde{Y}\right)$ and of $\mathrm{Op}_{\mathbb{L}}\left(\tilde{B}',\tilde{Y}'\right)$ coincide whenever $\tilde{B}=\tilde{B}'$ on these bonds and $\tilde{Y}$, $\tilde{Y}'$ contain the same ones of these sites. Every such number $\rho$ is called a \emph{range} of the formula.

The restriction in (iii) is meaningful. Indeed, $l\mathbb{Z}^3\subset(L\cdot M)\mathbb{Z}^3$, since $l$ is a multiple of $L\cdot M$, and $T_v\tilde{B}=\tilde{B}$, $\tilde{Y}+v=\tilde{Y}$ for $v\in l\mathbb{Z}^3$; by (ii), $\mathrm{Op}_{\mathbb{L}}\left(\tilde{B},\tilde{Y}\right)$ commutes with $T_v$ for these $v$. By (i), $\mathrm{Op}_{\mathbb{L}}\left(\tilde{B},\tilde{Y}\right)$ is defined on all functions; consequently it maps periodic functions to periodic functions and antiperiodic functions to antiperiodic functions. In the situation of Definition~\ref{def:periodisation-projection}, with $l_1$ and $l_2$ in the class, one has
\[
\mathrm{Op}\left(\mathrm{pj}^{*}B,\mathrm{pj}^{-1}Y\right)\circ\mathrm{pj}^{*}=\mathrm{pj}^{*}\circ\mathrm{Op}(B,Y),
\]
where the operator on the left is the one for the side $l_2$ and the operator on the right the one for the side $l_1$. In fact the $l_2\mathbb{Z}^3$-periodic extensions of $\mathrm{pj}^{*}B$ and of $\mathrm{pj}^{-1}Y$ are the $l_1\mathbb{Z}^3$-periodic extensions of $B$ and of $Y$, so that both sides are restrictions of one and the same operator $\mathrm{Op}_{\mathbb{L}}\left(\tilde{B},\tilde{Y}\right)$.

Let $l$ be a side of the class, let $\tilde{B}$ and $\tilde{Y}$ be the $l\mathbb{Z}^3$-periodic extensions of $B$ and of $Y$, let $\boldsymbol{\chi}(n):=1$ for $n\in\mathbb{Z}^3$ if $\mathrm{Op}$ acts on periodic functions and $\boldsymbol{\chi}(n):=(-1)^{n_0+n_1+n_2}$ if it acts on antiperiodic ones, and let $\mathscr{F}$ be a set of representatives of $I$ modulo $l\mathbb{Z}^3$. Then, for $x\in I'$ and every function $f$ on which $\mathrm{Op}(B,Y)$ acts,
\[
\left(\mathrm{Op}(B,Y)f\right)(x)=\sum_{x'\in\mathscr{F}}K_l(x,x')\cdot f(x'),\qquad K_l(x,x'):=\sum_{n\in\mathbb{Z}^3}\boldsymbol{\chi}(n)\cdot\mathrm{Op}_{\mathbb{L}}\left(\tilde{B},\tilde{Y}\right)\left(x,x'+l\cdot n\right),
\]
where only finitely many terms of the sum over $n$ differ from zero, by~(i); and $K_l$ is the kernel of $\mathrm{Op}(B,Y)$ in the sense of Section~\ref{sec:two-tori}. Indeed, by~(iii) the left side is the finite sum of the products $\mathrm{Op}_{\mathbb{L}}\left(\tilde{B},\tilde{Y}\right)(x,x'')\cdot f(x'')$ over all $x''\in I$; each $x''\in I$ equals $x'+l\cdot n$ for exactly one pair $(x',n)\in\mathscr{F}\times\mathbb{Z}^3$, and $f(x'+l\cdot n)=\boldsymbol{\chi}(n)\cdot f(x')$. Now let $l_1$, $l_2$ be any two sides of the class, let $(B_i,Y_i)$ be data on the torus of side $l_i$ and $(\tilde{B}_i,\tilde{Y}_i)$ their periodic extensions, $i=1,2$, let $\rho$ be a range of the formula, and let $x\in I'$ be such that $\tilde{B}_1=\tilde{B}_2$ on the bonds both of whose end points have sup distance at most $\rho$ from $\operatorname{supp}x$ and that $\tilde{Y}_1$, $\tilde{Y}_2$ contain the same sites at sup distance at most $\rho$ from $\operatorname{supp}x$. Then the rows at $x$ of $\mathrm{Op}_{\mathbb{L}}\left(\tilde{B}_1,\tilde{Y}_1\right)$ and of $\mathrm{Op}_{\mathbb{L}}\left(\tilde{B}_2,\tilde{Y}_2\right)$ coincide, by the second part of the requirement of locality. Let, moreover, $x'\in I$, and let $\operatorname{supp}x\cup\operatorname{supp}x'$ be contained in a set $U$ of sites of $\mathbb{L}_k$ with $|z-z'|+\rho<l$ for all $z,z'\in U$, where $l$ is a side of the class. Then, for data $(B,Y)$ on the torus of side $l$ with the periodic extensions $\tilde{B}$, $\tilde{Y}$, one has $K_l(x,x')=\mathrm{Op}_{\mathbb{L}}\left(\tilde{B},\tilde{Y}\right)(x,x')$: for $n\ne0$, $y\in\operatorname{supp}x$ and $y'\in\operatorname{supp}x'$ one has
\[
\left|y-\left(y'+l\cdot n\right)\right|\ge l-|y-y'|>\rho,
\]
and $\operatorname{supp}(x'+l\cdot n)=\operatorname{supp}x'+l\cdot n$, so that the term with this $n$ vanishes by the first part of the requirement of locality. Consequently, if $x$ and $x'$ satisfy both hypotheses, the second one for $l=l_1$ and for $l=l_2$, then the kernels of $\mathrm{Op}(B_1,Y_1)$ and of $\mathrm{Op}(B_2,Y_2)$ take the same value at $(x,x')$.

The phrase is used also for families of objects attached to polymers, in the following sense. Let the level $k$, a class of sides, each of them a multiple of $L\cdot M$, a number $\rho\ge0$ and a type of data be fixed; a statement that uses the phrase names its number $\rho$ and its data. Let $l$ range over the class, and let $X$ range over the polymers of $T_k=\mathbb{L}_k/l\mathbb{Z}^3$ that have a lift $\tilde{X}$ (clause~(9)) with
\[
|z-z'|+2\rho+M<l\qquad\text{for all }z,z'\in\tilde{X}.
\]
Let $X^{\rho}$ be the set of the sites of $T_k$ at sup distance at most $\rho$ from $X$. The data $u$ are finitely many functions of one or several indices contained in $X^{\rho}$, for instance a bond function on the bonds contained in $X^{\rho}$, or kernels all of whose arguments are contained in $X^{\rho}$; they may depend on further parameters. Consider a family of objects $O_l(X;u)$, each of which is a number, a kernel with its arguments contained in $X^{\rho}$, an element of the Grassmann algebra over $\mathbb{C}$ whose generators are the generators of Section~\ref{sec:grassmann-algebras-norms} with indices contained in $X^{\rho}$, or a homomorphism between such algebras. Let $\tilde{X}^{\rho}$ be the set of the sites of $\mathbb{L}_k$ at sup distance at most $\rho$ from $\tilde{X}$. The sup distance of two points of $\tilde{X}^{\rho}$ exceeds that of two suitable points of $\tilde{X}$ by at most $2\rho$, and is therefore less than $l-M$; hence $\tilde{X}^{\rho}$ is a lift of its image, by clause~(9), and this image is $X^{\rho}$: a site of $\tilde{X}^{\rho}$ is mapped to a site at sup distance at most $\rho$ from $X$, and if a site $x$ of $T_k$ satisfies $d(x,y)\le\rho$ for some $y\in X$, then $\tilde{y}+w\in\tilde{X}^{\rho}$ is mapped onto $x$, where $\tilde{y}$ is the preimage of $y$ in $\tilde{X}$ and $w$ is a shortest representative of $x-y$. By clause~(9), applied with the margin $M$, the covering map is injective on $\tilde{X}^{\rho}$ and preserves, together with its inverse, nearest neighbours, bonds and cells, the pairs of $\mathrm{Pz}_k$ together with their six paths, and the $M$-cubes together with their common faces. The functions on the indices of $T_k$ being read as periodic or antiperiodic functions on the index sets of $\mathbb{L}_k$, the \emph{transported datum} $\tilde{u}$ is the restriction of $u$ to the indices contained in $\tilde{X}^{\rho}$; and the generators of $T_k$ whose indices are contained in $\tilde{X}^{\rho}$ are independent and are labelled by these indices, as in Definition~\ref{def:admissibility-flatness-agreement}(c). The family is \emph{given by one formula on} $\mathbb{L}_k$ in terms of the data $u$, with the range $\rho$, if there is a map $\Phi$, one and the same for all sides of the class, defined on the pairs $(\tilde{X},\tilde{u})$ formed by a non-empty finite union $\tilde{X}$ of $M$-cubes of $\mathbb{L}_k$ that is connected through faces and by data $\tilde{u}$ of the given type on the indices contained in $\tilde{X}^{\rho}$, such that
\[
\Phi\left(\tilde{X}+v,T_v\tilde{u}\right)=T_v\,\Phi\left(\tilde{X},\tilde{u}\right)\qquad\text{for all }v\in(L\cdot M)\mathbb{Z}^3,
\]
and such that $O_l(X;u)=\Phi(\tilde{X},\tilde{u})$ for every side $l$ of the class, every polymer $X$ and every datum $u$ as above, and every lift $\tilde{X}$ of $X$; by clause~(9), such a lift is the set of the sites of a union of $M$-cubes of $\mathbb{L}_k$ that is connected through faces. Here $T_v$ is the identity on numbers; on functions of indices and on kernels it acts by $\left(T_vf\right)(x_1,\dots,x_p)=f(x_1-v,\dots,x_p-v)$; on a Grassmann algebra it is the isomorphism that sends the generator labelled $x$ to the generator labelled $x+v$; and on a homomorphism $h$ it acts by $h\mapsto T_v\circ h\circ T_v^{-1}$. In the equation $O_l(X;u)=\Phi(\tilde{X},\tilde{u})$ kernels are compared at the arguments contained in $\tilde{X}^{\rho}$, and elements and homomorphisms through the labelling of the generators by the indices contained in $\tilde{X}^{\rho}$. For the homomorphism $S_L$ which replaces the generators $\psi(x)$, $\chi(x,x')$ of the level $k$ by $\psi_L(x)$, $\chi_L(x,x')$ (Section~\ref{sec:local-parts-adjustment}), and which leads from the level $k$ to the level $k+1$, an index $x$ is sent to $x/L$, and the translation by $v$ at the level $k$ is matched with the translation by the vector $v/L\in M\mathbb{Z}^3$ at the level $k+1$. By clause~(9), two lifts of $X$ differ by a vector $v\in l\mathbb{Z}^3$, polymers being connected through nearest neighbours, and $v\in(L\cdot M)\mathbb{Z}^3$, as $l$ is a multiple of $L\cdot M$. The restriction of a periodic datum to the indices contained in $\tilde{X}^{\rho}+v$ is the image under $T_v$ of its restriction to the indices contained in $\tilde{X}^{\rho}$; for a function that is antiperiodic in each of $p$ arguments the factor $\sigma^p$ appears, where $\sigma:=(-1)^{\frac{v_0+v_1+v_2}{l}}$; and the generator labelled $x+v$ is $\sigma$ times the generator labelled $x$. Hence, by the covariance of $\Phi$, if $O_l(X;u)=\Phi(\tilde{X},\tilde{u})$ holds for one lift $\tilde{X}$ of $X$, then it holds for every lift of $X$, provided that the antiperiodic data and kernels have an even number of arguments, that the elements are even and that the homomorphisms map generators to odd elements (Section~\ref{sec:antiperiodicity}). Now let $l_1$, $l_2$ be sides of the class, let $\mathrm{pr}_i$ be the projection of $\mathbb{L}_k$ onto the torus of side $l_i$, and let $\tilde{X}$ be a non-empty finite union of $M$-cubes of $\mathbb{L}_k$, connected through faces, such that $|z-z'|+2\rho+M$ is less than $l_1$ and less than $l_2$ for any two sites $z$, $z'$ of $\tilde{X}$; then $\mathrm{pr}_i\tilde{X}$ is a polymer, of which $\tilde{X}$ is a lift, by clause~(9). Let $u_1$, $u_2$ be data on the two tori whose transported data coincide, $\tilde{u}_1=\tilde{u}_2$. Then $O_{l_1}(\mathrm{pr}_1\tilde{X};u_1)$ and $O_{l_2}(\mathrm{pr}_2\tilde{X};u_2)$ are both given by $\Phi(\tilde{X},\tilde{u}_1)$: the numbers are equal, the kernels are equal at the arguments contained in $\tilde{X}^{\rho}$, and the elements correspond to each other, and the homomorphisms are intertwined, under the isomorphism that sends each generator of the first torus labelled by an index contained in $\tilde{X}^{\rho}$ to the generator of the second torus with the same label. If, in the situation of Section~\ref{sec:admissible-sets-local}, the set $\tilde{X}$ is admissible (Definition~\ref{def:admissibility-flatness-agreement}(a)) and $\rho\le6d_U\cdot M$, then this margin holds on both tori, since any two sites $z$, $z'$ of $\tilde{X}$ satisfy $|z-z'|<\operatorname{diam}(\tilde{X})\le l_1-14d_U\cdot M$, whence
\[
|z-z'|+2\rho+M<l_1-14d_U\cdot M+12d_U\cdot M+M<l_1\le l_2;
\]
moreover, $\tilde{X}^{\rho}$ lies in $N(\tilde{X})$, because a site at sup distance at most $\rho$ from $\tilde{X}$ lies in an $M$-cube whose centre has sup distance less than $(6d_U+1)\cdot M$, hence at most $6d_U\cdot M$, from the centre of an $M$-cube of $\mathrm{box}(\tilde{X})$; and the isomorphism just described is the restriction of the isomorphism $\mathfrak{i}$ of Definition~\ref{def:admissibility-flatness-agreement}(c).

\item[(3)] \emph{Exact and weakened operators; kinds.} The \emph{exact} operators are the operators of Sections~\ref{sec:averaging-fluctuations-gauge} and~\ref{sec:exact-fermion-operators} and of Definition~\ref{def:pairs-sites-holder}, together with all operators built from these without recourse to primitives, for example the operators $T^{(i)}$ and $\bar{T}^{(i)}$ in~\eqref{eq:renormalisation-group-58}. Given a system of primitives (Section~\ref{sec:primitive-operators}), each member of one of its families~\eqref{it:primitives-p1}--\eqref{it:primitives-p5} is called a \emph{primitive}, or a \emph{weakened operator}. In Section~\ref{sec:weakened-walk-expansions} the same names apply also to the operators $S_y(s,A)$ of Definition~\ref{def:weakened-sums-walks}, out of which the families~\eqref{it:primitives-p3} and~\eqref{it:primitives-p4} are built. The name \emph{weakened fluctuation field} is given to the field $\mathscr{Z}(s')$ of Definition~\ref{def:renormalisation-step}. The \emph{polymer propagators} are the members $S_{k,y}(X,A)$ of the family~\eqref{it:primitives-p5}; in Section~\ref{sec:operators-fermion-field} they are written $S_y(X,A)$. A family with parameters has the \emph{kind} of its parameter cubes (Definition~\ref{def:parameter-cubes-polydisks}): the family~\eqref{it:primitives-p1} is of kind $\pi'_k$, the family~\eqref{it:primitives-p2} is of kind $\pi''_k$, and the family~\eqref{it:primitives-p3} and the family $\Gamma_k$ are of kind $\pi_k$. The indices of a family of kind $\pi'_k$, $\pi''_k$, $\pi_k$ are bonds, basis elements and sites, respectively.

\item[(4)] \emph{Block fields, background, fluctuation field, harmonic part.} A real bond function on the unit lattice $T^0_k$ is called a \emph{block field} of level $k$. The variables $A$, $A'$ of Section~\ref{sec:fermion-integrals-blocking} are block fields, and so are, along the flow, the fields $A_k$ of Definition~\ref{def:coordinates-slice-small}. For a block field $A$ of level $k$ the field $H_kA$ is called the \emph{background}, and for $\zeta\in V_k$ the field $\mathrm{Sq}_k\zeta$ is called the \emph{fluctuation field} (Definition~\ref{def:fluctuation-field-controlled}). For a bond function on a torus, the \emph{harmonic part} is its component in $V_{\mathrm{h}}$ (Section~\ref{sec:lattice-measures-laws}), that is, the orthogonal projection of the bond function onto the space of constant bond functions; this projection is identified with the vector whose entries are the means of the three components of the bond function.

\item[(5)] \emph{Legs.} The factors $\psi(\xi)$ and $\chi(\boldsymbol{\zeta})$ of a product of generators of Section~\ref{sec:grassmann-algebras-norms} are its \emph{legs}; for a kernel, the legs are its arguments $\xi$ and $\boldsymbol{\zeta}$. A leg is \emph{barred} if the corresponding generator is $\bar{\psi}$ or $\bar{\chi}$, and a leg $\chi(\boldsymbol{\zeta})$ is called a \emph{H\"older leg}. In Sections~\ref{sec:dependence-stop-level} and~\ref{sec:fermion-integral-stop} the kernels $\mathfrak{H}_A(\boldsymbol{\upsilon};u)$, through which a leg is expressed in terms of $\Psi$, are called legs as well.

\item[(6)] \emph{Stop level and stop scale.} Let $R\in\mathrm{Tor}(e)$ (Definition~\ref{def:tori-stop-level}) and $N\ge N_{\mathrm{min}}(R,e)$, and assume, as in Theorem~\ref{thm:scales-stop-levels}(a), that $v_K\ge1$ and $\left(1+\frac{\mathrm{lg}}{v_K}\right)^r\le L$. Then the level $k^{*}(R)$ of Definitions~\ref{def:tori-stop-level} and~\ref{def:gauge-fixing-stop} exists by Theorem~\ref{thm:scales-stop-levels}(a), and $1\le k^{*}(R)\le K$ by Proposition~\ref{prop:depth-stop-level}(c) (Section~\ref{sec:running-quantities-tori}); this level is the \emph{stop level} on $R$. An object is called \emph{terminal} when it belongs to the level $k^{*}(R)$. The \emph{stop scale} is the number $h'=h\cdot L^{-m_0}$ of Definition~\ref{def:mass-bound-stop}; it equals $a_K$, where $K=N+\log_Lh'$ is the level of Section~\ref{sec:running-quantities-tori}, and it depends neither on $R$ nor on $e$. The number $a_{*}=a_{k^{*}(R)}$ is the spacing, in physical units, of the unit lattice of level $k^{*}(R)$; one has $a_{*}\le h'$, because $a_k=L^{k-N}$ increases with $k$ and $k^{*}\le K$ (Section~\ref{sec:running-quantities-tori}; Proposition~\ref{prop:depth-stop-level}(c)).

\item[(7)] \emph{Small-field set and large fields.} The set $\mathfrak{G}$ of Definition~\ref{def:coordinates-slice-small} is the \emph{small-field set}, and a \emph{large field} is an element of $F_N\setminus\mathfrak{G}$; in Lemma~\ref{lem:gaussian-bound-large} the large fields are the elements of the larger set $F_N\setminus\mathfrak{G}^{-}$. The qualification \emph{without large fields} (Sections~\ref{sec:renormalisation-group-step} and~\ref{sec:one-step-large}) means that the domains of Section~\ref{sec:domains-bond-functions} serve as the domains of definition of the step map and that the integral over the fluctuation variable is taken over $\mathrm{Box}_k$ (Section~\ref{sec:averaging-fluctuations-gauge}).

\item[(8)] \emph{Marks and marked vertices.} In Section~\ref{sec:four-jets-one} the word \emph{mark} designates one of the variables $t_c$, $c\in\mathrm{Cel}$, introduced in Section~\ref{sec:probes-jets-reference}. Among the factors of a tree expansion, those that are derivatives with respect to a mark are called \emph{marked vertices}; the two distinguished factors in Lemma~\ref{lem:tree-expansion}(c) are of this type.

\item[(9)] \emph{Sets that do not wrap; lifts.} Let $W$ be a set of sites of the torus $T_k=\mathbb{L}_k/l\mathbb{Z}^3$ (Section~\ref{sec:lattice-geometry}), for instance a ball, a cube or a polymer, and let $\eta_k=L^{-k}$ be the spacing of $\mathbb{L}_k$ (Section~\ref{sec:running-quantities-tori}). The sup-diameter of a subset of $\mathbb{L}_k$ is the largest sup distance of two of its points. A \emph{lift} of $W$ is a subset $\tilde{W}$ of $\mathbb{L}_k$ that the covering map $\mathbb{L}_k\to T_k$ maps onto $W$ and whose sup-diameter is less than $l-\eta_k$; one says that $W$ \emph{does not wrap} if it has a lift. Let $\tilde{W}$ be a lift of $W$. The covering map is injective on $\tilde{W}$, since two distinct points of $\mathbb{L}_k$ with the same image differ by a non-zero vector of $l\mathbb{Z}^3$ and are therefore at sup-distance at least $l$; for $x\in W$ let $\tilde{x}$ denote its preimage in $\tilde{W}$.

Let $s\ge0$ be such that the sup-diameter of $\tilde{W}$ is less than $l-s$, let $x,x'\in W$, and let $w\in\mathbb{L}_k$ be a representative of the class $x-x'$ with $|w|\le s$. Then $w=\tilde{x}-\tilde{x}'$, because the vector $w-\left(\tilde{x}-\tilde{x}'\right)$ belongs to $l\mathbb{Z}^3$ and has a sup norm less than $s+(l-s)=l$. In particular $\left|\tilde{x}-\tilde{x}'\right|=d(x,x')$ whenever $d(x,x')\le s$, since Definition~\ref{def:sup-norm-distance} provides a representative of sup norm $d(x,x')$. As the sup-diameter of $\tilde{W}$ is less than $l-\eta_k$, the case $s=\eta_k$ applies; hence $x'=x\pm\eta_ke_{\mu}$ in $T_k$ if and only if $\tilde{x}'=\tilde{x}\pm\eta_ke_{\mu}$. Consequently the bonds, plaquettes and cubes of $T_k$ all of whose corners lie in $W$ are exactly the images of the bonds, plaquettes and cubes of $\mathbb{L}_k$ all of whose corners lie in $\tilde{W}$; the coboundary $d$ and all difference quotients of functions on $W$ can be computed on $\tilde{W}$; the map $x\mapsto\tilde{x}$ is an $\mathbb{R}^3$-valued function on $W$, the coordinate on the lift; and for every $c\in\mathbb{R}^3$ the function $x\mapsto c\cdot\tilde{x}$ satisfies $d\left(c\cdot\tilde{x}\right)=c$ on every bond of $T_k$ with both end points in $W$, that is, $\partial_{\mu}\left(c\cdot\tilde{x}\right)=c_{\mu}$ there. If the sup-diameter of $\tilde{W}$ is less than $l-M$, the case $s=M$ applies, and therefore the pairs of $\mathrm{Pz}_k$ with both sites in $W$ are the images of the pairs $\left(\tilde{x},\tilde{x}'\right)$ of points of $\tilde{W}$ with $0<\left|\tilde{x}-\tilde{x}'\right|\le M$; each such pair $\left(\tilde{x},\tilde{x}'\right)$ is a shortest lift of $(x,x')$ in the sense of Definition~\ref{def:pairs-sites-holder}, so that the six paths $\Gamma^{\pi}(x,x')$ are the images of the six rectilinear paths of $\mathbb{L}_k$ from $\tilde{x}$ to $\tilde{x}'$.

The whole torus wraps: a subset of $\mathbb{L}_k$ that is mapped onto $T_k$ takes in each coordinate $l/\eta_k$ different values of $\eta_k\mathbb{Z}$, so that its sup-diameter is at least $l-\eta_k$. A subset $W'$ of a set $W$ that does not wrap does not wrap either, the preimage of $W'$ in a lift of $W$ being a lift of $W'$. If $W$ is connected through nearest neighbours, as balls and cubes are, any two lifts $\tilde{W}$ and $\hat{W}$ of $W$ differ by a vector of $l\mathbb{Z}^3$. Indeed, for $x\in W$ let $\tilde{x}$ and $\hat{x}$ be its preimages in $\tilde{W}$ and in $\hat{W}$; then $\hat{x}-\tilde{x}\in l\mathbb{Z}^3$, and for nearest neighbours $x$, $y$ in $W$ the preceding paragraph gives $\tilde{y}-\tilde{x}=\hat{y}-\hat{x}$; so $\hat{x}-\tilde{x}$ does not depend on $x\in W$.

Let now $l$ be a multiple of $M$ with $l>2M$, as every element of $\mathrm{Cl}$ is. Then the preimage in $\tilde{W}$ of an $M$-cube contained in $W$ is the set of the sites of an $M$-cube of $\mathbb{L}_k$, and two $M$-cubes contained in $W$ have a common face if and only if their preimages have one. Indeed, the centre $c$ of an $M$-cube $\square\subset W$ is a site of $\square$, and $\tilde{c}\in M\mathbb{Z}^3$, as $l\mathbb{Z}^3\subset M\mathbb{Z}^3$; every site $c+z$ of $\square$, where $|z|<\frac{M}{2}$, is joined to $c$ by nearest-neighbour steps inside $\square$, so that its preimage is $\tilde{c}+z$; hence the preimage of $\square$ is the set of the sites of the $M$-cube of $\mathbb{L}_k$ centred at $\tilde{c}$. Two $M$-cubes in $W$ with a common face contain nearest neighbours $x$, $x'$ on the two sides of this face; then $\tilde{x}$, $\tilde{x}'$ are nearest neighbours in two different $M$-cubes of $\mathbb{L}_k$, which therefore have a common face; conversely, as $l>2M$, the images of two $M$-cubes of $\mathbb{L}_k$ with a common face have a common face. In particular, a lift of a union of $M$-cubes is the set of the sites of a union of $M$-cubes of $\mathbb{L}_k$. A polymer is connected through nearest neighbours, since the sites of an $M$-cube form a box and two $M$-cubes with the centres $c$ and $c+M\cdot e_{\mu}$ contain the nearest neighbours $c+\frac{M-\eta_k}{2}\cdot e_{\mu}$ and $c+\frac{M+\eta_k}{2}\cdot e_{\mu}$; so the lift of a polymer is unique up to a translation by a vector of $l\mathbb{Z}^3$. For a polymer $X$ with $2M\cdot\left(d_M(X)+2\right)<l$ it is the set of the sites of the union $\tilde{X}$ of Proposition~\ref{prop:polymers}(t10), any two points of which have sup distance at most $M\cdot\left(d_M(X)+2\right)<\frac{l}{2}$ by Proposition~\ref{prop:polymers}(t7); and for a flat polymer the lift $\tilde{X}$ of Definition~\ref{def:admissibility-flatness-agreement}(b) is a lift in the sense of this clause on both tori of Section~\ref{sec:admissible-sets-local}, since its sites have sup-diameter at most $S-\eta_k$, where $S\le l_1-14d_U\cdot M$ is the largest side of $\mathrm{box}(\tilde{X})$.

Let $s<l$ be an integer multiple of $\eta_k$. The cube of side $s$ centred at a site $c$ of $T_k$, that is, the image of the set $\left\{\tilde{x}\in\mathbb{L}_k:\left|\tilde{x}-\tilde{c}\right|<\frac{s}{2}\right\}$, where $\tilde{c}$ is a representative of $c$, does not wrap: these sets are its lifts, and their sup-diameter is at most $s-\eta_k$. A ball $K(x_0,r)$ of Section~\ref{sec:balls-torus-lifts} does not wrap if $2r+\eta_k<l$, and its lifts are the sets $\tilde{x}_0+\left\{z\in\mathbb{L}_k:|z|\le r\right\}$, $\tilde{x}_0$ a representative of $x_0$. A kernel that is antiperiodic in each argument is evaluated with all its arguments in one lift of a set $W$; if $W$ is connected through nearest neighbours and the number of arguments is even, the value does not depend on the lift, by the uniqueness of lifts up to translations and Section~\ref{sec:antiperiodicity}. For balls see Section~\ref{sec:balls-torus-lifts}, for polymers Proposition~\ref{prop:polymers}(t10), and for the pairs of $\mathrm{Pz}_k$ Definition~\ref{def:pairs-sites-holder}.

\item[(10)] \emph{The scheme.} The \emph{scheme} has three parts: the parameters, that is, the constants fixed in the stages~\ref{it:order-choice-1}--\ref{it:order-choice-10} of the order of choice of Section~\ref{sec:order-choice} (Definition~\ref{def:order-choice}; Section~\ref{sec:parameters-scales}); the orthonormal basis $\{e_i^0\}$ of the kernel of $\tau$ in the space of the real functions on the interior bonds of $B(0)$, which is chosen in Section~\ref{sec:averaging-fluctuations-gauge}, where Proposition~\ref{prop:basis-map-box} provides it, once for each $L$, and which is the same at every level $k$, for every side $l$, every $R$ and every $N$; and the systems of primitives that Definition~\ref{def:periodisation-projection} prescribes on the tori $T_k(R_{\mathrm{top}})$ with $k<k_{\mathrm{top}}$ (Lemma~\ref{lem:assembled-systems-primitives}), where $R_{\mathrm{top}}$, $k_{\mathrm{top}}$ and the set $\mathrm{Tor}(e)$ are those of Definition~\ref{def:tori-stop-level}. The basis is a choice: it is neither a constant in the sense of Definition~\ref{def:constants}, which is a number, nor a primitive, and the step map depends on it through the basis $\{\Upsilon\}$ of $V_k$, since $\mathrm{Box}_k$ and $\hat{\mu}$ are defined through the coordinates $\zeta(\Upsilon)$ (Section~\ref{sec:averaging-fluctuations-gauge}). A quantity \emph{depends on the scheme} if it depends on these constants, this basis and these systems. The tori of $\mathrm{Tor}(e)$ are called the tori of the scheme.
\end{enumerate}
\end{definition}

\subsection{The objects of the scheme}\label{sec:objects-scheme}

\subsubsection{Lattice measures and laws}\label{sec:lattice-measures-laws}

\paragraph{Cochains.}
Let $L$ be the odd integer of stage~\ref{it:order-choice-4}, let $0<\bar{m}\le1$, let $N\ge0$ be an integer and $\eta:=L^{-N}$, and let $R>0$ be such that $n:=R/\eta$ is an odd integer with $n\ge3$. Put $\Lambda_n:=(\mathbb{Z}/n\mathbb{Z})^3$. For every $j$ let $C^j$ be the real vector space of the $j$-cochains on $\Lambda_n$, normed by the counting $\ell^2$ norm $|\cdot|$, and let $d$ be the coboundary. The Hodge decomposition on the discrete torus (Section~\ref{sec:results-literature}) gives the orthogonal decomposition
\[
C^1=V_{\mathrm{g}}\oplus V_{\mathrm{h}}\oplus V_{\mathrm{t}},
\]
in which $V_{\mathrm{g}}=dC^0$ is the space of the exact one-cochains, $V_{\mathrm{h}}$ the space of the constant one-cochains and $V_{\mathrm{t}}$ the space of the coclosed one-cochains with mean zero; moreover, the kernel of $d$ on $C^1$ is $V_{\mathrm{g}}\oplus V_{\mathrm{h}}$. Let $F_N:=dC^1\subset C^2$, and let $\gamma_N$ be the standard Gaussian measure of the Euclidean space $F_N$, whose inner product is that of $C^2$. Since $d$ vanishes on $V_{\mathrm{g}}\oplus V_{\mathrm{h}}$ and is injective on $V_{\mathrm{t}}$, each $F\in F_N$ has exactly one preimage $A(F)\in V_{\mathrm{t}}$, that is, $dA(F)=F$.

\paragraph{The Wilson operator.}
Let $e>0$ be the charge, and put $q:=e\cdot\eta^{\frac{1}{2}}$. For $\mu\in\{0,1,2\}$ let $e_{\mu}$ be the unit vector of $\mathbb{Z}^3$ in the direction $\mu$. Let $A$ be a complex function on the oriented bonds of $\Lambda_n$ with $A(x',x)=-A(x,x')$, read as an $n\mathbb{Z}^3$-periodic function on the oriented bonds of $\mathbb{Z}^3$. The functions $f\colon\mathbb{Z}^3\to\mathbb{C}^4$ with $f(x+n\cdot e_{\mu})=-f(x)$ for all $x\in\mathbb{Z}^3$ and all $\mu$ form a finite-dimensional complex vector space, and all determinants below are taken on this space. On it the Wilson operator $D_A$, with Wilson parameter one, four components, the link phases $\exp(iq\cdot A)$ and antiperiodic boundary conditions in all directions, is
\begin{multline*}
(D_Af)(x):=-\sum_{\mu=0}^2\left(\frac{1-\gamma_{\mu}}{2}\cdot\exp\left(iq\cdot A(x,x+e_{\mu})\right)\cdot f(x+e_{\mu})\right.\\
\left.+\frac{1+\gamma_{\mu}}{2}\cdot\exp\left(iq\cdot A(x,x-e_{\mu})\right)\cdot f(x-e_{\mu})-f(x)\right);
\end{multline*}
the right side is again antiperiodic, because $A$ is periodic. Here $\gamma_0$, $\gamma_1$, $\gamma_2$ are the matrices of Section~\ref{sec:dirac-matrices-lattice}. The levels are the integers $k\ge0$. If $L$ divides $n$, then $D_A$ is the operator of Section~\ref{sec:exact-fermion-operators} at the level $0$ on the side $l=n$; this is the case for the sides $R=h'\cdot L^T$ of Section~\ref{sec:running-quantities-tori}, for which $n=L^{K+T}$ and, as $n\ge3$, $K+T\ge1$. The operator $D_A$ coincides with the operator $D[qA]$ of~\eqref{eq:free-coefficient-5}.

\paragraph{Twisted densities and the averaged law.}
Let $(m_{\mathrm{in}},\boldsymbol{\varepsilon}_{\mathrm{in}})\in\mathbb{R}^2$ be a real pair, let $\theta\in\mathbb{R}^3$, and let $a\in C^1\otimes\mathbb{C}$. Let $H_{\theta}\in C^1$ be the constant one-cochain with
\[
H_{\theta}(x,x+e_{\mu}):=\frac{\theta_{\mu}}{q\cdot n}\qquad(\mu=0,1,2),
\]
and define, for $F\in F_N$,
\[
\rho^{\theta,a}_N(F):=\exp\left(-\boldsymbol{\varepsilon}_{\mathrm{in}}\cdot n^3\right)\cdot\frac{\det\left(D_{A(F)+H_{\theta}+a}+\bar{m}\cdot\eta+m_{\mathrm{in}}\right)}{\det\left(D_0+\bar{m}\cdot\eta\right)};
\]
the denominator is positive by Lemma~\ref{lem:product-formula-volume}, applied with $\theta=0$ and the mass $\bar{m}\cdot\eta\in(0,1]$. Let $\nu^{\theta,a}:=\rho^{\theta,a}_N\cdot\gamma_N$ be the complex measure on $F_N$ whose density with respect to $\gamma_N$ is $\rho^{\theta,a}_N$, and put
\[
\nu^{\theta}:=\nu^{\theta,0},\qquad\bar{\nu}^a:=\int_{[-\pi,\pi]^3}\frac{\mathrm{d}\theta}{(2\pi)^3}\,\nu^{\theta,a},\qquad\bar{\nu}:=\bar{\nu}^0,\qquad\bar{\mu}:=\frac{\bar{\nu}}{\bar{\nu}(1)}\quad\left(\bar{m}\cdot\eta+m_{\mathrm{in}}>-1\right),
\]
where $\bar{\nu}(1)$ is the total mass of $\bar{\nu}$. Let $\bar{m}\cdot\eta+m_{\mathrm{in}}>-1$. Then $\nu^{\theta}$ is a non-negative measure for every $\theta\in\mathbb{R}^3$: the denominator of $\rho^{\theta,0}_N$ is positive, and its numerator is non-negative by Theorem~\ref{thm:gaussian-domination-positive}(b) and the identification of the Wilson operator in Section~\ref{sec:characters-twisted-spinors}. By this identification, together with the identification of $H_{\theta}$ and the computation in Section~\ref{sec:proca-parent-harmonic}, there is a number $c>0$ with $\bar{\nu}(G)=c\cdot\overline{Z}[G]$ for every bounded Borel function $G$ on $C^2$, where $\bar{\nu}(G)$ is the integral of the restriction of $G$ to $F_N$ and $\overline{Z}$ is the functional of~\eqref{eq:gaussian-bounds-3} formed with the lattice mass $\bar{m}\cdot\eta+m_{\mathrm{in}}$. Hence $0<\bar{\nu}(1)<\infty$ by Theorem~\ref{thm:limit-vanishing-proca}, and $\bar{\mu}$ is a probability measure on $F_N$ whose expectation is the functional $\overline{\mathsf{E}}$ of~\eqref{eq:gaussian-bounds-3}. In the setting of Lemma~\ref{lem:positivity-large-fields}, in which the pair is controlled on $R$ (Definition~\ref{def:fluctuation-field-controlled}), the condition $\bar{m}\cdot\eta+m_{\mathrm{in}}>-1$ holds (proof of Lemma~\ref{lem:positivity-large-fields}). The law $\bar{\mu}$ is the harmonic-averaged law~\eqref{eq:setup-av}.

\paragraph{The all-flux law.}
Let $\boldsymbol{\mu}:=\bar{m}\cdot\eta+m_{\mathrm{in}}$ be the lattice mass. For $\boldsymbol{\mu}>-1$ the all-flux law~\eqref{eq:setup-af} (Definition~\ref{def:partition-functions-all}; Section~\ref{sec:wilson-determinant-lattice}) is the functional $E_n$ formed with this mass in the following way. Let $G$ be a Borel function on $C^2$ for which there are numbers $c$, $c'$ with $|G(F)|\le c\cdot\exp\left(c'\cdot|F|\right)$ for all $F\in C^2$, and put
\begin{align}
Z_n[G]&:=\sum_{\substack{\mathfrak{z}\in C^2(\Lambda_n;\mathbb{Z})\\ d\mathfrak{z}=0}}\ \int_{\left[0,\frac{2\pi}{q}\right)^{\mathrm{bonds}}}\mathrm{d}A\,\exp\left(-\frac{1}{2}|F|^2\right)\cdot\det\left(D_A+\boldsymbol{\mu}\right)\cdot G(F),\notag\\
F&:=dA+\frac{2\pi}{q}\cdot\mathfrak{z},\qquad E_n[G]:=\frac{Z_n[G]}{Z_n[1]}\quad\left(\boldsymbol{\mu}>-1\right).\notag
\end{align}
Here $C^2(\Lambda_n;\mathbb{Z})$ is the set of the integer-valued two-cochains, and the integral is taken, with respect to the Lebesgue measure $\mathrm{d}A$ of $C^1$, over the set $\left[0,\frac{2\pi}{q}\right)^{\mathrm{bonds}}$ of the one-cochains $A$ whose value $A(x,x+e_{\mu})$ on every bond lies in $\left[0,\frac{2\pi}{q}\right)$. For $\boldsymbol{\mu}>-1$ the denominator of $E_n$ is positive: $Z_n[1]>0$ (Section~\ref{sec:summation-all-fluxes}). Let $Z_{n,0}[G]$ be the sum of those terms of the series defining $Z_n[G]$ in which $\mathfrak{z}$ has flux zero (Section~\ref{sec:tori-cochains}). For $\boldsymbol{\mu}>-1$ one has $Z_{n,0}[1]>0$ (Section~\ref{sec:decomposition-flux}), and this partial sum is divided by its own value at $G=1$:
\[
E_{n,0}[G]:=\frac{Z_{n,0}[G]}{Z_{n,0}[1]}\qquad\left(\boldsymbol{\mu}>-1\right).
\]
The functional $E_{n,0}$ is the harmonic-averaged law~\eqref{eq:setup-av} (Section~\ref{sec:decomposition-flux} and Lemma~\ref{lem:flux-sectors}(b)).

\paragraph{Partition functions.}
For a real number $\mathbf{m}$ put
\begin{align*}
\rho_n(q,\mathbf{m})&:=\int_{[-\pi,\pi]^3}\frac{\mathrm{d}\theta}{(2\pi)^3}\int\mathrm{d}\gamma_N(F)\,\det\left(D_{A(F)+H_{\theta}}+\mathbf{m}\right),\\
\rho^{\mathrm{free}}_n&:=\int_{[-\pi,\pi]^3}\frac{\mathrm{d}\theta}{(2\pi)^3}\,\det\left(D_{H_{\theta}}+\bar{m}\cdot\eta\right).
\end{align*}
Inserting the definition of $\rho^{\theta,0}_N$ into that of $\bar{\nu}(1)$ one obtains
\[
\rho_n\left(q,\bar{m}\cdot\eta+m_{\mathrm{in}}\right)=\exp\left(\boldsymbol{\varepsilon}_{\mathrm{in}}\cdot n^3\right)\cdot\det\left(D_0+\bar{m}\cdot\eta\right)\cdot\bar{\nu}(1).
\]

\paragraph{Coordinates of level $k$.}
For a function $a$ on bonds and a level $k$ put
\[
a^{[k]}(b):=L^{\frac{k}{2}}\cdot a\left(L^k\cdot b\right)
\]
for every bond $b$ such that $a$ is defined on the bond $L^k\cdot b$; thus $a^{[k]}$ is $a$ written in the coordinates of level $k$. Let $a\in C^1\otimes\mathbb{C}$, and let $k$ be a level such that $L^k$ divides $n$. Then the function $a^{[k]}$ is defined on the bonds of the fine lattice $T_k$ of level $k$ and side $l_k=n\cdot L^{-k}$ (Sections~\ref{sec:running-quantities-tori} and~\ref{sec:lattice-geometry}), because $b\mapsto L^k\cdot b$ maps these bonds onto the bonds of $\Lambda_n$ (one has $L^k\cdot\eta_k=1$ and $L^k\cdot l_k=n$). This side equals $R/a_k$, with $a_k=L^{k-N}$ as in Section~\ref{sec:running-quantities-tori}; for the sides $R=h'\cdot L^T$ of that section one has $n=L^{K+T}$, and the levels $k$ for which $L^k$ divides $n$ are the levels $k\le K+T$.

\subsubsection{Running quantities and tori}\label{sec:running-quantities-tori}

Let $M=L^m$ be the integer of stage~\ref{it:order-choice-6}, let $r$, $p_0$, $p$ and $\varepsilon$ be the constants of stage~\ref{it:order-choice-8}, let $b$ be the constant of stage~\ref{it:order-choice-1}, and let $m_0$ be the integer of Definition~\ref{def:mass-bound-stop}. Put
\[
h:=L^{\left\lfloor\log_L\frac{1}{\bar{m}}\right\rfloor},\qquad h':=h\cdot L^{-m_0},\qquad\mathrm{lg}:=\frac{1}{2}\log L,
\]
and, for integers $k$,
\begin{align*}
&a_k:=L^{k-N},\qquad e_k:=e\cdot a_k^{\frac{1}{2}},\qquad v_k:=-\log e_k,\\
&r_k:=\min\left\{L^i:\ i\ge0\text{ an integer},\ L^i\ge v_k^r\right\},\qquad p_k:=v_k^p,\qquad p_{0,k}:=v_k^{p_0},\\
&w^{\mathrm{m}}_k:=e_k^{\frac{3}{4}-8\varepsilon},\qquad w^{\mathrm{E}}_k:=e_k^{\frac{1}{4}-7\varepsilon};
\end{align*}
at $k=0$ one has $e_k=q$, and the numbers of level $0$ are formed with $q$ (Section~\ref{sec:lattice-measures-laws}; Definition~\ref{def:notation}(e)). The running mass of level $k$ is $\bar{m}\cdot a_k$; for $k\ge1$ it is denoted
\[
\mathfrak{m}_k:=\bar{m}\cdot a_k,
\]
whereas at level $0$ it is written $\bar{m}\cdot\eta$, the symbol $\mathfrak{m}_0$ being reserved for the upper end of the mass range (Definitions~\ref{def:notation}(b) and~\ref{def:mass-bound-stop}). Let $K$ be the level with $a_K=h'$, that is, $K:=N+\log_Lh'$. The sides of the tori are the numbers $R=h'\cdot L^T$ with an integer $T\ge0$; for such $R$ put
\[
l_k:=\frac{R}{a_k},\qquad n_k:=l_k^3,\qquad\mathfrak{n}_k:=\log_L\left(\frac{l_k}{M\cdot r_k}\right),\qquad\eta_k:=L^{-k}.
\]
For integers $k\ge1$ let
\[
b_k:=\frac{b\cdot\left(1-L^{-1}\right)}{1-L^{-k}};
\]
equivalently, $b_1=b$ and $b_{k+1}=\frac{b\cdot b_k}{b_k+b\cdot L^{-1}}$. With $b_{-}:=b\cdot\left(1-L^{-1}\right)$ and $b_{+}:=b$ one has $b_{-}\le b_k\le b_{+}$.

\begin{definition}[Tori and stop level]\label{def:tori-stop-level}
Put
\[
R_{\mathrm{top}}(e):=h'\cdot M\cdot r_K(e),\qquad\mathrm{Tor}(e):=\left\{R=h'\cdot L^T:\ T\ge0\text{ an integer},\ R\le R_{\mathrm{top}}(e)\right\};
\]
here $r_K(e)$ is the number $r_K$, which depends on $e$ and, since $a_K=h'$, not on $N$. The reference tori are $R^{*}:=R_{\mathrm{top}}/L$ and $L\cdot R^{*}=R_{\mathrm{top}}$. Let $R\in\mathrm{Tor}(e)$. Since $K=N+\log_Lh'$ is the index of a level, the integer $N$ is always taken with $N\ge0$ and $N+\log_Lh'\ge0$, and the integers $N_1$ below are subject to the same two conditions. The number $N_{\mathrm{min}}(R,e)$ is the least such $N_1$ for which
\[
R\cdot L^N\ge L\cdot M\cdot r_0\qquad\text{for all }N\ge N_1;
\]
the number $r_0$ on the right is formed with the running charge $q=e\cdot L^{-\frac{N}{2}}$ of level $0$ (Definition~\ref{def:notation}(e)), and therefore it changes with $N$. As $l_0=R\cdot L^N$, the displayed inequality is equivalent to $\mathfrak{n}_0\ge1$; thus $\mathfrak{n}_0\ge1$ for every $N\ge N_{\mathrm{min}}(R,e)$. The stop level $k^{*}=k^{*}(R)$ is the least level $k$ with $\mathfrak{n}_k=0$. For the running quantities and the objects of the scheme that are indexed by the level $k$, a subscript $*$ marks the value at $k=k^{*}(R)$, the symbol $n_{*}$ excepted (Definition~\ref{def:notation}(e)). Other uses of a subscript $*$, among them the push-forward $g_{*}$ by a lattice isometry $g$ (Section~\ref{sec:cells-reference-tori}), are declared at their places. Finally, $k_{\mathrm{top}}:=k^{*}(R_{\mathrm{top}})$.
\end{definition}

Let $R\in\mathrm{Tor}(e)$ and $N\ge N_{\mathrm{min}}(R,e)$, and assume the hypothesis of Theorem~\ref{thm:scales-stop-levels}(a), namely $v_K\ge1$ and $\left(1+\frac{\mathrm{lg}}{v_K}\right)^r\le L$. Then the level $k^{*}(R)$ exists by Theorem~\ref{thm:scales-stop-levels}(a), and $1\le k^{*}\le K$ by Proposition~\ref{prop:depth-stop-level}(c). (The integers $\mathfrak{n}_k$ satisfy $\mathfrak{n}_0\ge1$ and $\mathfrak{n}_K=\log_L\left(R/R_{\mathrm{top}}(e)\right)\le0$, and by Theorem~\ref{thm:scales-stop-levels}(a) they change by $-1$ or by $0$ from $k$ to $k+1$ for $0\le k<K$; hence the least $k$ with $\mathfrak{n}_k=0$ lies in $\{1,\dots,K\}$.) As $\mathfrak{n}_k=0$ at $k=k^{*}$, one has $l_{*}=M\cdot r_{*}$. Moreover,
\begin{align*}
&L^{-m_0-1}<\mathfrak{m}_K\le L^{-m_0}\le\mathfrak{m}_0,\qquad\bar{m}\cdot a_k\le\mathfrak{m}_{*}\le\mathfrak{m}_K\quad(k\le k^{*}),\\
&\mathfrak{m}_{*}\cdot l_{*}=\bar{m}\cdot R,\qquad e_{*}\le e_K\le e_{*}\cdot l_{*}^{\frac{1}{2}}.
\end{align*}
All these relations are contained in Proposition~\ref{prop:depth-stop-level}(e); those for the masses are contained in Proposition~\ref{prop:bounds-masses} as well.

\subsubsection{Lattice geometry}\label{sec:lattice-geometry}

\paragraph{Lattices and blocks.}
Let $k$ be a level and $l$ a positive integer. Put
\[
\mathbb{L}_k:=\eta_k\mathbb{Z}^3,\qquad T_k:=\mathbb{L}_k/l\mathbb{Z}^3,\qquad T^0_k:=\mathbb{Z}^3/l\mathbb{Z}^3.
\]
The torus $T_k$ is the fine lattice of level $k$ and side $l$, and $T^0_k$ is the unit lattice; their points are the fine sites and the unit sites. For a unit site $y$, the block $B_k(y)=\Delta_y$ is the set of the fine sites $x$ with $|x-y|<\frac{1}{2}$. The class of sides is
\[
\mathrm{Cl}:=\left\{l:\ \frac{l}{L\cdot M}\text{ is an odd integer }\ge13\right\}.
\]
If $R\in\mathrm{Tor}(e)$ and $k\le k^{*}(R)$, the side $l_k=l_k(R)$ belongs to $\mathrm{Cl}$; this follows from the condition~\eqref{eq:renormalisation-group-16} of Section~\ref{sec:conditions} by Proposition~\ref{prop:depth-stop-level}(f). In the remainder of Section~\ref{sec:lattice-geometry} the side $l$ belongs to $\mathrm{Cl}$. For $y\in L\mathbb{Z}^3$, $B(y)$ is the $L$-block of unit sites centred at $y$, that is, the set of the unit sites $x$ with $|x-y|<\frac{L}{2}$.

\paragraph{Cubes and polymers.}
With $M=L^m$, the $M$-cubes are the cubes of side $M$ centred at the points of $M\mathbb{Z}^3$, and the $LM$-cubes are the cubes of side $L\cdot M$ centred at the points of $LM\mathbb{Z}^3$. The closed $M$-cube with a given centre is the closed cube of side $M$ with this centre in the continuum torus $\mathbb{R}^3/l\mathbb{Z}^3$. A net is a connected finite union of closed straight segments in $\mathbb{R}^3/l\mathbb{Z}^3$, segments reduced to a point being admitted; its length is the sum of the Euclidean lengths of its segments. Let $X$ be a union of $M$-cubes.
\begin{itemize}
\item For an integer $i\ge0$, $X^{\sim i}$ is the enlargement of $X$ by $i$ layers of $M$-cubes, that is, the union of the $M$-cubes whose centres have sup distance at most $i\cdot M$ from the centre of some $M$-cube of $X$.
\item $|X|_M$ is the number of the $M$-cubes of $X$, and $\operatorname{Vol}(X):=M^3\cdot|X|_M$.
\item $X$ is connected through faces if any two of its $M$-cubes are joined by a finite sequence of $M$-cubes of $X$ in which any two consecutive cubes have a common face. A polymer is a non-empty union of $M$-cubes that is connected through faces, and $D_k$ is the set of the polymers.
\item $M\cdot d_M(X)$ is the length of a shortest tree that meets all $M$-cubes of $X$; precisely,
\[
M\cdot d_M(X):=\inf\left\{\text{length of }\mathcal{N}:\ \mathcal{N}\text{ a net that meets every closed }M\text{-cube of }X\right\}.
\]
The set $X$ is small if $d_M(X)\le L$, and large otherwise.
\item $\bar{X}^L$ is the union of the $LM$-cubes that contain an $M$-cube of $X$.
\end{itemize}

\begin{definition}[Enlarged cubes]\label{def:enlarged-cubes}
For an $M$-cube $\square$ put $\square^{\dagger}:=\square^{\sim L}$; it is the cube of side $(2L+1)\cdot M$ with the centre of $\square$. The cube $\square^{\sim5}$ has the side $11M$.
\end{definition}

\begin{definition}[Parameter cubes and polydisks]\label{def:parameter-cubes-polydisks}
Let $\pi_k$ be the set of the $M$-cubes, $\pi'_k$ the set of the cubes of side $2M$ with the same centres, and $\pi''_k$ the set of the cubes of side $2L\cdot M$ centred at the points of $LM\mathbb{Z}^3$. The cubes of $\pi'_k$ and of $\pi''_k$ are closed: the cube with the centre $c$ is $\{x:|x-c|\le M\}$, resp.\ $\{x:|x-c|\le L\cdot M\}$. The elements of $\pi_k$, $\pi'_k$ and $\pi''_k$ are the parameter cubes of the three kinds, and a lattice site is a site of $T_k$. Two cubes of $\pi_k$ are linked if they have a common face; two cubes of $\pi'_k$, or two cubes of $\pi''_k$, are linked if they have a common lattice site. A set of parameter cubes of one kind is link-connected if any two of its cubes are joined by a finite sequence of its cubes in which any two consecutive cubes are linked; the link-components of such a set are its maximal link-connected subsets. Let $W$ be a set of lattice sites; a parameter cube is said to be contained in $W$ when each of its lattice sites belongs to $W$. A parameter cube and an $M$-cube are said to meet when some lattice site belongs to both. For an index set $I$ that is one of $\pi_k$, $\pi'_k$, $\pi''_k$, or a union of some of them, put
\[
\mathrm{Pd}_k:=\left\{s=(s_c)_{c\in I}:\ s_c\in\mathbb{C},\ |s_c|\le M^{\alpha_0}\text{ for all }c\in I\right\},
\]
$\alpha_0$ being the exponent of stage~\ref{it:order-choice-5}; the same symbol serves for all these index sets, the index set being that of the family under consideration. Finally,
\[
N_c:=(2L+1)^3,\qquad c_l:=2L,\qquad d_U:=2L+1.
\]
\end{definition}

Section~\ref{sec:cubes-parameter-cubes} derives the geometric properties of these cubes, the facts~\eqref{it:renormalisation-group-g0}--\eqref{it:renormalisation-group-g4}; in particular, the following two hold. First, let $c'$ be a cube of $\pi'_k$, resp.\ of $\pi''_k$, and let $c$ be its centre; an $M$-cube meets $c'$ if and only if its own centre lies at sup distance at most $M$, resp.\ at most $L\cdot M$, from $c$. Consequently every parameter cube, of any of the three kinds, has a lattice site in common with at most $N_c$ of the $M$-cubes (a cube of $\pi_k$ is itself an $M$-cube and meets no other $M$-cube). Second, if a non-empty set of parameter cubes of a single kind is link-connected, then the union of those $M$-cubes that meet at least one of its cubes is a polymer.

\paragraph{Indices.}
The sites, the bonds and the basis elements $\Upsilon$ of Section~\ref{sec:averaging-fluctuations-gauge} are called indices; the two arguments $x$, $x'$ of the kernel $\mathrm{Op}(x,x')$ of an operator are indices. The support $\operatorname{supp}x$ of an index $x$ is the set of its end points: the site itself if $x$ is a site, its two end points if $x$ is a bond, and the end points of the bonds in the support of $\Upsilon$ if $x=\Upsilon$. An index $x$ meets a set $W$ of sites if $\operatorname{supp}x$ meets $W$, and $x$ is contained in $W$, or lies in $W$, if $\operatorname{supp}x\subset W$. For two indices $x$, $x'$ the number
\[
d(x,x'):=\min\left\{d(y,y'):\ y\in\operatorname{supp}x,\ y'\in\operatorname{supp}x'\right\}
\]
is the sup distance on the torus (Definition~\ref{def:sup-norm-distance}) of their supports. For an index $x$ let $\square(x)$ denote the union of all $M$-cubes met by $x$. For a site $z$ let $K_z:=\{x:|x-z|\le2L\}$, the set of the sites $x$ with sup distance at most $2L$ from $z$.

\subsubsection{Dirac matrices and lattice symmetries}\label{sec:dirac-matrices-lattice}

Let $\sigma_1$, $\sigma_2$, $\sigma_3$ be the Pauli matrices and let $1$ also denote the unit $2\times2$ matrix. Define the $4\times4$ matrices
\[
\gamma_{\mu}:=\sigma_3\otimes\sigma_{\mu+1}\quad(\mu=0,1,2),\qquad\gamma_3:=\sigma_1\otimes1,\qquad\mathsf{C}:=i\gamma_1,\qquad\varkappa:=\gamma_3\cdot\mathsf{C}.
\]
Let $\mathrm{Sp}$ be the group generated by $\gamma_0$, $\gamma_1$, $\gamma_2$, $\gamma_3$. An element $S\in\mathrm{Sp}$ is a product of generators, each of which commutes or anticommutes with $\gamma_{\nu}$; hence, for $\nu\in\{0,1,2\}$,
\[
S^{-1}\gamma_{\nu}S=r_{\nu}(S)\cdot\gamma_{\nu}\qquad\text{with a sign }r_{\nu}(S)\in\{1,-1\}.
\]
The diagonal $3\times3$ matrices with diagonal entries $\pm1$ form the group of signs, and $r$ is the map of $\mathrm{Sp}$ into this group given by
\[
r(S):=\operatorname{diag}\left(r_0(S),r_1(S),r_2(S)\right).
\]
The elements of $\mathrm{Sp}$ that commute with $\gamma_3$ form the subgroup
\[
\mathrm{Sp}_3:=\left\{S\in\mathrm{Sp}:\ S\gamma_3=\gamma_3S\right\}.
\]
For a matrix $N$ let $N^T$ be its transpose and $\overline{N}$ the matrix of the complex conjugates of its entries. The following relations are proved in Section~\ref{sec:gamma-matrices-symmetry}. The four matrices $\gamma_{\mu}$ are Hermitian, the matrices $\mathsf{C}$ and $\varkappa$ are real, and
\begin{align*}
&\gamma_{\mu}\gamma_{\nu}+\gamma_{\nu}\gamma_{\mu}=2\delta_{\mu\nu}\qquad(0\le\mu,\nu\le3),\\
&\mathsf{C}^T=\mathsf{C}^{-1}=-\mathsf{C},\qquad-\gamma_{\mu}^T=\mathsf{C}^{-1}\cdot\gamma_{\mu}\cdot\mathsf{C}\qquad(0\le\mu\le3),\\
&\varkappa^T=\varkappa=\varkappa^{-1},\qquad\overline{\gamma_{\mu}}=\varkappa^{-1}\cdot\gamma_{\mu}\cdot\varkappa\qquad(0\le\mu\le2),\\
&\overline{\gamma_3}=-\varkappa^{-1}\cdot\gamma_3\cdot\varkappa,\qquad\varkappa\cdot\mathsf{C}=-\mathsf{C}\cdot\varkappa.
\end{align*}
Moreover, $r$ maps $\mathrm{Sp}_3$ onto the group of signs.

Let $l$, a multiple of $M$, be the side of the torus $T_k$. The group $\mathbb{G}_k$ consists of the pairs $(S,a)$ with $S\in\mathrm{Sp}$ and $a\in M\mathbb{Z}^3/l\mathbb{Z}^3$, multiplied according to
\[
(S,a)(S',a'):=\left(SS',\,a+r(S)a'\right);
\]
an element $g=(S,a)$ acts on the points $x$ of $\mathbb{R}^3/l\mathbb{Z}^3$ by
\[
gx:=r(S)x+a.
\]
If $l$ is a multiple of $L\cdot M$, as every element of $\mathrm{Cl}$ is, the pairs with $a\in LM\mathbb{Z}^3/l\mathbb{Z}^3$ form the step group $\mathbb{G}_k^L$. The subgroups $\mathbb{G}_{k,3}$ of $\mathbb{G}_k$ and $\mathbb{G}^L_{k,3}$ of $\mathbb{G}_k^L$ consist of the elements $(S,a)$ with $S\in\mathrm{Sp}_3$.

The element $g=(S,a)$ acts on functions, on families of parameters and on sets by the following rules:
\begin{itemize}
\item for an integer $i$ with $-1\le i\le k$ and a bond function $A$ on $L^{-i}\mathbb{Z}^3/l\mathbb{Z}^3$: $(gA)(x,x'):=A\left(g^{-1}x,g^{-1}x'\right)$;
\item on scalars by $(g\boldsymbol{\lambda})(x):=\boldsymbol{\lambda}\left(g^{-1}x\right)$;
\item for a function $F$ of the oriented plaquettes $p$ that takes opposite values on the two orientations of each plaquette: $(gF)(p):=F\left(g^{-1}p\right)$;
\item on spinor functions by $(gf)(x):=S\cdot f\left(g^{-1}x\right)$;
\item on the generators of Section~\ref{sec:grassmann-algebras-norms} by
\[
(g\psi)(x):=S\cdot\psi\left(g^{-1}x\right),\qquad(g\bar{\psi})(x):=\left(S^{-1}\right)^T\cdot\bar{\psi}\left(g^{-1}x\right),
\]
\[
(g\chi)(x,x'):=S\cdot\chi\left(g^{-1}x,g^{-1}x'\right),\qquad(g\bar{\chi})(x,x'):=\left(S^{-1}\right)^T\cdot\bar{\chi}\left(g^{-1}x,g^{-1}x'\right),
\]
and by the same formulas on the generators $\Psi$ and $W$;
\item on the families $s=(s_c)$ indexed by parameter cubes $c$ (Definition~\ref{def:parameter-cubes-polydisks}) by $(gs)_c:=s_{g^{-1}c}$, where $g$ is any element of $\mathbb{G}_k$ if every index is a cube of $\pi_k$ or of $\pi'_k$, and $g$ is an element of the step group $\mathbb{G}_k^L$ if cubes of $\pi''_k$ are among the indices, so that in both cases $g^{-1}c$ is a cube of the same kind as $c$ (Proposition~\ref{prop:symmetries-extraction}(a));
\item on sets by $X\mapsto gX$.
\end{itemize}
When $g$ acts on antiperiodic functions, its translation $a$ is taken as a vector of $M\mathbb{Z}^3$. A translation by a vector of $l\mathbb{Z}^3$, and likewise the element $-1$ of $\mathrm{Sp}$, multiplies an antiperiodic function by a sign. Since the group acts on operators by conjugation, these elements act trivially on operators; they also act trivially on the even elements of the Grassmann algebras. Elementary properties of the actions are stated in Proposition~\ref{prop:symmetries-extraction}(a)--(d).

\begin{definition}[Required lattice symmetry]\label{def:required-lattice-symmetry}
The group of lattice symmetries required of the polymer activities (Definition~\ref{def:renormalisation-group-1}) is $\mathbb{G}_k$. The group required of the primitives (Section~\ref{sec:primitive-operators}) depends on the family: for the families~\eqref{it:primitives-p1}--\eqref{it:primitives-p4}, that is, for $H_k(s')$, $C_k^{\frac{1}{2}}(s')$, $H_k(s'',A)$, $\Gamma_k(s'',A)$ and $M_k(X,A)$, it is the step group $\mathbb{G}_k^L$, and for the family~\eqref{it:primitives-p5} of the operators $S_{k,y}(X,A)$ it is the subgroup $\mathbb{G}^L_{k,3}$ of the elements $(S,a)$ with $S\in\mathrm{Sp}_3$; these are the requirements of property~\ref{it:renormalisation-group-14}. An element $S\in\mathrm{Sp}$ that does not belong to $\mathrm{Sp}_3$ anticommutes with $\gamma_3$ (each generator of $\mathrm{Sp}$ commutes or anticommutes with $\gamma_3$), so that conjugation by $S$ replaces $i\gamma_3y$ by $-i\gamma_3y$, that is, $y$ by $-y$: for $g=(S,a)$ with $S\notin\mathrm{Sp}_3$ the exact operators satisfy $S_{k,-y}(gA)\cdot g=g\cdot S_{k,y}(A)$ (Proposition~\ref{prop:symmetries-fermion-operators}(b)), whereas the family~\eqref{it:primitives-p5} is indexed by $y\ge0$. Since $r$ maps $\mathrm{Sp}_3$ onto the group of signs, the elements of $\mathbb{G}^L_{k,3}$ induce the same maps $x\mapsto r(S)x+a$ of the torus as the elements of $\mathbb{G}_k^L$; the geometric lattice symmetries required are therefore the same for all five families. Invariance under permutations of the coordinates, or under translations by unit vectors, is not part of the requirement.
\end{definition}

\subsubsection{Averaging and fluctuations of the gauge field}\label{sec:averaging-fluctuations-gauge}

\paragraph{Difference operators and pairings.}
A bond function $\mathscr{A}$ on $T_k$ changes sign when a bond is reversed; its components are $\mathscr{A}_{\mu}(x):=\mathscr{A}(x,x+\eta_ke_{\mu})$. For functions $f$ on the sites of $T_k$, in particular for scalars $\lambda$ and for the components of bond functions $\mathscr{A}$, set
\begin{align*}
(\partial_{\mu}f)(x)&:=\eta_k^{-1}\cdot\left(f(x+\eta_ke_{\mu})-f(x)\right),\qquad(d\lambda)_{\mu}:=\partial_{\mu}\lambda,\\
(d\mathscr{A})_{\mu\nu}&:=\partial_{\mu}\mathscr{A}_{\nu}-\partial_{\nu}\mathscr{A}_{\mu}\qquad(\mu<\nu).
\end{align*}
For two functions $f$, $f'$ of the same kind (scalars, bond functions or functions of the plaquettes) let
\[
\langle f,f'\rangle:=\sum_x\eta_k^3\cdot f(x)\cdot f'(x),\qquad|f|_2^2:=\langle f,f\rangle,
\]
where the sum is taken over the sites $x$ and over the components; $\delta:=d^T$ is the transpose of $d$ with respect to these pairings. If $\eta_k$ is replaced by $1$ in all definitions of this paragraph, one obtains the corresponding objects on the unit lattice $T^0_k$: the components, the operators $\partial_{\mu}$, $d$ and $\delta$, the pairings and the norms. On either lattice, the plain sum of the values of a bond function $A$ on the bonds of a path $\Gamma$, taken without a factor $\eta_k$, is written
\[
A(\Gamma):=\sum_{b\in\Gamma}A(b).
\]
The straight path from $x$ to $x'$ is denoted by $\Gamma_{x,x'}$. Each of the six orders $\pi$ of the three directions determines a rectilinear path $\Gamma^{\pi}(y,x)$ from $y$ to $x$: it changes the coordinates one at a time, in the order $\pi$, each to its final value.

\paragraph{Block averaging and axial sums.}
Let $y\in T^0_k$. For a scalar $\lambda$, a bond function $\mathscr{A}$ and a function $F$ of the plaquettes of $T_k$ define
\begin{align*}
(Q_k\lambda)(y)&:=\sum_{x\in B_k(y)}\eta_k^3\cdot\lambda(x),\\
(Q_k\mathscr{A})_{\mu}(y)&:=\sum_{x\in B_k(y)}\eta_k^3\sum_{j=0}^{L^k-1}\eta_k\cdot\mathscr{A}_{\mu}(x+j\eta_ke_{\mu}),\\
\left(Q_k^{(2)}F\right)_{\mu\nu}(y)&:=\sum_{x\in B_k(y)}\eta_k^3\sum_{j,j'=0}^{L^k-1}\eta_k^2\cdot F_{\mu\nu}(x+j\eta_ke_{\mu}+j'\eta_ke_{\nu}).
\end{align*}
In the rest of this subsection $l$ is a multiple of $L$, so that the $L$-blocks $B(y)$, $y\in L\mathbb{Z}^3/l\mathbb{Z}^3$, partition $T^0_k$. On the unit lattice the $L$-blocks, the bonds, the faces and the paths of the rest of this subsection are those of $\mathbb{Z}^3$, counted modulo $l\mathbb{Z}^3$: a bond function on $T^0_k$ is read as an $l\mathbb{Z}^3$-periodic bond function on $\mathbb{Z}^3$, and for $y\in L\mathbb{Z}^3$ the block $B(y)$ is the set of the points $x\in\mathbb{Z}^3$ with $|x-y|<\frac{L}{2}$. In particular, for $x\in\mathbb{Z}^3$ and a direction $\mu$ the straight path $\Gamma_{x,x+Le_{\mu}}$ is the path of the $L$ bonds $\left(x+ie_{\mu},x+(i+1)e_{\mu}\right)$, $0\le i<L$. If $l\ge3L$, the blocks, bonds, faces and straight paths so defined are, modulo $l\mathbb{Z}^3$, those formed on the torus $T^0_k$ itself. For a bond function $A$ on $T^0_k$ and $y\in L\mathbb{Z}^3/l\mathbb{Z}^3$ the average of one step is
\[
(QA)(y,y+Le_{\mu}):=L^{-4}\sum_{x\in B(y)}A\left(\Gamma_{x,x+Le_{\mu}}\right),
\]
and for a bond function $A$ and two sites $y$, $x$ the axial functional is the mean of $A$ over the six rectilinear paths from $y$ to $x$:
\[
(\tau A)(y,x):=\frac{1}{3!}\sum_{\pi}A\left(\Gamma^{\pi}(y,x)\right).
\]
Let $Z$ be a bond function on the unit lattice $T^0_k$. The condition $\tau Z=0$ is the system of equations
\[
(\tau Z)(y,x)=0\qquad\left(y\in L\mathbb{Z}^3/l\mathbb{Z}^3,\ x\in B(y)\right),
\]
one for each centre $y$ of an $L$-block and each site $x$ of that block. In the space of the scalars on $T_k$, the kernel $\ker Q_k$ of $Q_k$ is mapped by $\delta d$ onto a subspace $\delta d(\ker Q_k)$; the projection onto this subspace that is orthogonal with respect to the pairing $\langle\cdot,\cdot\rangle$ is denoted by $\mathscr{R}_k$.

\paragraph{Interior bonds, face bonds and the space $V_k$.}
A bond whose two ends lie in one $L$-block $B(y)$, $y\in L\mathbb{Z}^3$, is called interior. Every other positively oriented bond, of direction $\mu$ say, leads from exactly one block $B(y)$ into the block $B(y+Le_{\mu})$; it is called a face bond of the face $\langle y,y+Le_{\mu}\rangle$, the face being determined by the block $B(y)$ that the bond leaves. The faces are indexed by the pairs $(y,\mu)$ of a point $y\in L\mathbb{Z}^3/l\mathbb{Z}^3$ and a direction $\mu\in\{0,1,2\}$. Each face has $L^2$ face bonds, among them its central bond, which lies on the straight line joining $y$ and $y+Le_{\mu}$. The real bond functions $\mathsf{Z}$ on $T^0_k$ whose value at the central bond of every face is zero and which satisfy $\tau\mathsf{Z}=0$ form a real vector space $V_k$, the fluctuation space; it is equipped with the counting scalar product. As the paths $\Gamma^{\pi}(y,x)$ with $x\in B(y)$ consist of interior bonds of $B(y)$, the condition $\tau\mathsf{Z}=0$ constrains only the values of $\mathsf{Z}$ on the interior bonds, block by block. An element $\mathsf{Z}$ of $V_k$ is therefore given by the pair $\mathsf{Z}=(\mathsf{Z}_1,\mathsf{Z}_2)$ of its restrictions: $\mathsf{Z}_1$, the restriction to the interior bonds, satisfies $\tau\mathsf{Z}_1=0$, and $\mathsf{Z}_2$, the restriction to the face bonds that are not central, is an arbitrary real function.

Let $\{e_i^0\}$ be an orthonormal basis of the kernel of $\tau$ in the space of the real functions on the interior bonds of $B(0)$, each element of which is an eigenvector of every one of the reflections $x_{\mu}\mapsto-x_{\mu}$, $\mu=0,1,2$; such a basis exists by Proposition~\ref{prop:basis-map-box} and by Proposition~\ref{prop:fluctuation-space}(i). For a given $L$ this choice is made once: the same vectors $e_i^0$ are taken at every level $k$ and for every side $l$. For $y\in L\mathbb{Z}^3/l\mathbb{Z}^3$ let $e_i^y$ be the translate of $e_i^0$ by $y$, and for a positively oriented bond $b$ let $\delta_b$ be the bond function that equals $1$ on $b$ and $0$ on every other positively oriented bond. The basis of $V_k$ is
\begin{equation}\label{eq:renormalisation-group-1}
\{\Upsilon\}:=\{e_i^y\}\cup\left\{\delta_b:\ b\text{ a positively oriented face bond that is not central}\right\};
\end{equation}
it is orthonormal and has $d_V\cdot n_k$ elements (Proposition~\ref{prop:basis-map-box}; for every multiple $l$ of $L$, Proposition~\ref{prop:fluctuation-space}(ii) and Lemma~\ref{lem:face-averages-map}(c)), where $d_V:=2-2L^{-3}$ and $n_k=l^3$ is the number of the sites of $T^0_k$. For $\mathsf{Z}\in V_k$ write $\mathsf{Z}(\Upsilon):=\langle\mathsf{Z},\Upsilon\rangle$.

For $\mathsf{Z}\in V_k$ let $C\mathsf{Z}$ be the bond function on $T^0_k$ that agrees with $\mathsf{Z}$ on every bond that is not central, and whose value at the central bond $b_c$ of the face $c=\langle y,y+Le_{\mu}\rangle$ is
\[
(C\mathsf{Z})(b_c):=-\sum_{\substack{b\text{ face bond of }c\\ b\ne b_c}}\mathsf{Z}(b)-\sum_{\substack{b\text{ interior bond of direction }\mu\\ \text{of }B(y)\text{ or of }B(y+Le_{\mu})}}\frac{n_c(b)}{L}\cdot\mathsf{Z}(b).
\]
The weight $n_c(b)$ counts the sites $x\in B(y)$ whose straight path $\Gamma_{x,x+Le_{\mu}}$, made of $L$ bonds of direction $\mu$, traverses $b$; these are the weights of~\eqref{eq:operators-photon-16}. By Proposition~\ref{prop:fluctuation-space}(iv), $C$ is a linear bijection of $V_k$ onto the space of the real bond functions $Z$ on $T^0_k$ with $QZ=0$ and $\tau Z=0$.

\paragraph{Background fields, the fluctuation covariance and the one-step minimiser.}
For a real bond function $A$ on $T^0_k$ the background field $H_kA$ is the minimiser of $|d\mathscr{A}|_2^2$ among the real bond functions $\mathscr{A}$ on $T_k$ with $Q_k\mathscr{A}=A$ and $\mathscr{R}_k\delta\mathscr{A}=0$; this minimiser exists and is unique by Proposition~\ref{prop:background-field} and, for every side $l$, by Proposition~\ref{prop:minimiser}. Further, $H_0:=1$, and
\begin{align*}
\Delta_k&:=(dH_k)^TdH_k,\qquad\tilde{\Delta}_k:=C^T\Delta_kC,\qquad C_k:=\tilde{\Delta}_k^{-1},\\
C_k^{\frac{1}{2}}&:=\pi^{-1}\int_0^{\infty}\mathrm{d}x\,x^{-\frac{1}{2}}\left(\tilde{\Delta}_k+x\right)^{-1},\qquad\mathrm{Sq}_k:=C\,C_k^{\frac{1}{2}}.
\end{align*}
The operator $\tilde{\Delta}_k$ is positive definite on $V_k$ (Proposition~\ref{prop:spectral-bounds-decay}; also Proposition~\ref{prop:fluctuation-operator}(a)); hence $C_k$ exists, and $C_k^{\frac{1}{2}}$ is its positive square root.

For a real bond function $A'$ on $T^0_{k+1}$, the unit lattice of side $l/L$, define $A'_L$ on the bonds $b$ of $L\mathbb{Z}^3/l\mathbb{Z}^3$ by $A'_L(b):=L^{-\frac{1}{2}}A'(b/L)$, and let $A^{\min}_k(A')$ be the minimiser of $\langle A,\Delta_kA\rangle$ among the real bond functions $A$ on $T^0_k$ with $QA=A'_L$ and $\tau A=0$; it exists and is unique by Proposition~\ref{prop:minimisers-prescribed-block} and, for every $k\ge0$ and every multiple $l$ of $L$, by Theorem~\ref{thm:minimisation-one-step}(a). Let $\boldsymbol{\omega}_k(A')$ be the real scalar on $T_k$ of Proposition~\ref{prop:minimisers-prescribed-block}, fixed by the normalisation $Q_{k+1}\boldsymbol{\omega}_k(A')=0$; for every $k\ge0$ and every multiple $l$ of $L$ the scalar with this normalisation exists and is unique by Theorem~\ref{thm:minimisation-one-step}(b) (for the role of its additive constant see Section~\ref{sec:fermion-integrals-blocking}). In this normalisation $Q_{k+1}$ stands for the composition $QQ_k$, which maps the scalars on $T_k$ to the scalars on $L\mathbb{Z}^3/l\mathbb{Z}^3$; on a scalar $\lambda$ on $T^0_k$ the operator $Q$ is the mean over the $L$-blocks, $(Q\lambda)(y):=L^{-3}\sum_{x\in B(y)}\lambda(x)$ for $y\in L\mathbb{Z}^3/l\mathbb{Z}^3$ (Section~\ref{sec:averaging-operators}).

Let $\mu_I$ be the standard Gaussian measure on $V_k$, and
\begin{align*}
\mathrm{Box}_k&:=\left\{\zeta\in V_k:\ |\zeta(\Upsilon)|\le p_{0,k}\text{ for every }\Upsilon\right\},\qquad\mathrm{d}\hat{\mu}:=\frac{\mathbf{1}_{\mathrm{Box}_k}}{\mu_I(\mathrm{Box}_k)}\,\mathrm{d}\mu_I,\\
d\varepsilon^0_k&:=-n_k^{-1}\log\mu_I(\mathrm{Box}_k).
\end{align*}

Block averaging is performed one step at a time, with $Q$ on bond fields and with $Q(A)$ on spinors (Section~\ref{sec:exact-fermion-operators}). The gauge is fixed differently for the two kinds of fields: the fluctuation fields $Z$ satisfy $QZ=0$ and, in every $L$-block, the axial condition $\tau Z=0$, symmetrised by the mean over the six orders $\pi$; the background fields $\mathscr{A}=H_kA$ satisfy the Landau condition $\mathscr{R}_k\delta\mathscr{A}=0$.

\begin{definition}[Fluctuation field; controlled pairs]\label{def:fluctuation-field-controlled}
The root $C_k^{\frac{1}{2}}$ is called the exact root. In the step of level $k$ the fluctuation field is $\mathrm{Sq}_k\zeta$, $\zeta\in V_k$, formed with the exact root. A real pair $(m_{\mathrm{in}},\boldsymbol{\varepsilon}_{\mathrm{in}})$ is called \emph{controlled on} $R$ when two requirements are met: the trajectory of the pair on $R$, in the sense of Section~\ref{sec:flow}, is defined at each level $k$ with $0\le k\le k^{*}(R)$, and at every one of these levels the mass $m_k$ of the trajectory obeys $|m_k|\le w^{\mathrm{m}}_k$.
\end{definition}

\begin{definition}[Coordinates on the slice; small fields]\label{def:coordinates-slice-small}
On $\Lambda_n$ let $Q_j$ be the $j$-fold composition of $Q$ (Section~\ref{sec:averaging-operators}), and let the slice be
\[
\mathrm{Sl}:=\left\{A_0\in C^1(\Lambda_n):\ \tau Q_jA_0=0\text{ for }0\le j<k^{*},\text{ and }Q_{k^{*}}A_0\text{ is coclosed with mean zero}\right\}.
\]
Here $Q_jA_0$ is a bond function on the lattice $L^j\mathbb{Z}^3/n\mathbb{Z}^3$, and $\tau$ and the $L$-blocks are formed on this lattice, the blocks being centred at the points of $L^{j+1}\mathbb{Z}^3/n\mathbb{Z}^3$; for the fields $A_j$ defined below, $\tau Q_jA_0=0$ is equivalent to $\tau A_j=0$ (Section~\ref{sec:slice-space-curvatures}). By Proposition~\ref{prop:orthogonal-coordinates}, the restriction of $d$ to $\mathrm{Sl}$ is a linear bijection onto $F_N$. Every $F\in F_N$ has therefore exactly one potential in $\mathrm{Sl}$, that is, exactly one $A_0\in\mathrm{Sl}$ with $dA_0=F$; it is denoted by $A_0=A_0(F)$. The difference $A_0(F)-A(F)$ is exact: it is closed, and it has mean zero, because $A(F)\in V_{\mathrm{t}}$ has mean zero and because, by Proposition~\ref{prop:averaging-operators}(c) with $a=1$, $j=k^{*}$ and $l=n$ (admissible, as $n=L^{k^{*}}\cdot l_{*}$ with the integer $l_{*}=M\cdot r_{*}$), the field $A_0(F)$ has the mean of $Q_{k^{*}}A_0(F)$, which is zero by the definition of $\mathrm{Sl}$.

For $0\le j\le k^{*}$ let $A_j:=\left(Q_jA_0\right)^{[j]}$, the field $Q_jA_0$ in the scaling of level $j$ (Section~\ref{sec:lattice-measures-laws}); it is a bond function on $T^0_j$, and $(A_{j+1})_L=QA_j$ for $j<k^{*}$ (see~\eqref{eq:algebraic-gaussian-15}). Put $A_{*}:=A_{k^{*}}$. For $j<k^{*}$ the element $\zeta_j\in V_j$ is defined by
\[
A_j=A^{\min}_j(A_{j+1})+\mathrm{Sq}_j\zeta_j;
\]
it exists and is unique by Theorem~\ref{thm:hierarchical-coordinates}(a). Let $\mathrm{d}\mathsf{m}_{*}$ be the law of $A_{*}$ under $\gamma_N$, and, for a bond function $A$ on $T^0_{k^{*}}$, let $\chi_{*}(A):=\mathbf{1}\left\{|dA|_{\infty}\le p_{*}\right\}$. The small-field set is $\mathfrak{G}$; the set $\mathfrak{G}^{-}$ is given by the same conditions with each bound lowered by $1$:
\begin{align*}
\mathfrak{G}&:=\left\{F\in F_N:\ |\zeta_j(\Upsilon)|\le p_{0,j}\text{ for all }\Upsilon\text{ and all }j<k^{*};\ |dA_{*}|_{\infty}\le p_{*}\right\},\\
\mathfrak{G}^{-}&:=\left\{F\in F_N:\ |\zeta_j(\Upsilon)|\le p_{0,j}-1\text{ for all }\Upsilon\text{ and all }j<k^{*};\ |dA_{*}|_{\infty}\le p_{*}-1\right\}.
\end{align*}
Each of the linear functionals $F\mapsto\zeta_j(\Upsilon)$ on $F_N$ has norm one (Proposition~\ref{prop:orthogonal-coordinates}).

Finally, for a real bond function $A$ on $T^0_{k^{*}}$ and a family $\zeta=(\zeta_j)_{j<k^{*}}$ with $\zeta_j\in V_j$, the fields $A^{\downarrow}_j$ are defined by downward recursion,
\[
A^{\downarrow}_{k^{*}}:=A,\qquad A^{\downarrow}_j:=A^{\min}_j(A^{\downarrow}_{j+1})+\mathrm{Sq}_j\zeta_j\quad(j=k^{*}-1,\dots,0),
\]
and $A^{\downarrow}_0(A,\zeta)$ denotes the last of them.
\end{definition}

\subsubsection{Exact fermion operators}\label{sec:exact-fermion-operators}

Let $k\ge0$ be a level, let the side $l$ be a positive multiple of $L$ (as is every element of $\mathrm{Cl}$), and let $T_k=\mathbb{L}_k/l\mathbb{Z}^3$ and $T^0_k=\mathbb{Z}^3/l\mathbb{Z}^3$ be the fine and the unit lattice of Section~\ref{sec:lattice-geometry}; since $L$ is odd and divides $l$, the $L$-blocks $B(y)$, $y\in L\mathbb{Z}^3/l\mathbb{Z}^3$, partition $T^0_k$. Let $\mathfrak{m}$ be the mass of level $k$, that is, $\mathfrak{m}=\mathfrak{m}_k$ for $k\ge1$ and $\mathfrak{m}=\bar{m}\cdot\eta$ at the level $0$ (Section~\ref{sec:running-quantities-tori}), and let $A$ be a connection on $T_k$, that is, a complex function on the oriented bonds of $T_k$ with $A(x',x)=-A(x,x')$. A spinor on $T_k$ (on $T^0_k$) is a function $f$ on $\mathbb{L}_k$ (on $\mathbb{Z}^3$) with values in $\mathbb{C}^4$ such that $f(x+l\,e_{\nu})=-f(x)$ for $\nu=0,1,2$. With the matrices $\gamma_{\mu}$ of Section~\ref{sec:dirac-matrices-lattice}, the Wilson operator $D_A$ acts on the spinors on $T_k$ by
\begin{multline*}
\left(D_Af\right)(x):=-\eta_k^{-1}\sum_{\mu=0}^2\left\{\frac{1-\gamma_{\mu}}{2}\cdot\exp\left(i\,e_k\eta_k\cdot A\left(x,x+\eta_ke_{\mu}\right)\right)\cdot f\left(x+\eta_ke_{\mu}\right)\right.\\
\left.{}+\frac{1+\gamma_{\mu}}{2}\cdot\exp\left(i\,e_k\eta_k\cdot A\left(x,x-\eta_ke_{\mu}\right)\right)\cdot f\left(x-\eta_ke_{\mu}\right)-f(x)\right\}.
\end{multline*}
Let $Q(A)$, $Q^T(-A)$ and $P(A)=Q^T(-A)\cdot Q(A)$ be the one-step operators of Section~\ref{sec:fermion-integrals-blocking}; thus $P(A)$ acts on the spinors on $T^0_k$.

Now let $k\ge1$. For $0\le j\le k$ the \emph{blocks of side} $L^j\eta_k$ are the cubes of this side centred at the points of $L^j\eta_k\mathbb{Z}^3$, that is, the sets
\[
\left\{x\in T_k:\ |x-z|<\frac{L^j\eta_k}{2}\right\}\qquad\left(z\in L^j\eta_k\mathbb{Z}^3/l\mathbb{Z}^3\right);
\]
they partition $T_k$, as $L$ is odd and $l\mathbb{Z}^3\subset L^j\eta_k\mathbb{Z}^3$. For a site $x$ of $T_k$ let $y$ be the unit site with $x\in B_k(y)$, and let $y_j$ be the centre of the block of side $L^j\eta_k$ that contains $x$. The blocks of side $\eta_k$ are single sites, and those of side $L^k\eta_k=1$ are the blocks $B_k(y)$; hence $y_0=x$ and $y_k=y$. With the axial functional $\tau$ of Section~\ref{sec:averaging-fluctuations-gauge} put
\[
\left(\tau_kA\right)(y,x):=\sum_{j=0}^{k-1}\left(\tau A\right)\left(y_{j+1},y_j\right).
\]
The averaging operator $Q_k(A)$ maps the spinors on $T_k$ to the spinors on $T^0_k$, and $Q_k^T(-A)$ maps the spinors on $T^0_k$ to the spinors on $T_k$:
\begin{align*}
\left(Q_k(A)f\right)(y)&:=\eta_k^3\sum_{x\in B_k(y)}\exp\left(i\,e_k\eta_k\cdot\left(\tau_kA\right)(y,x)\right)\cdot f(x),\\
\left(Q_k^T(-A)v\right)(x)&:=\exp\left(-i\,e_k\eta_k\cdot\left(\tau_kA\right)(y,x)\right)\cdot v(y)\qquad\left(x\in B_k(y)\right).
\end{align*}
Put
\[
P_k(A):=Q_k^T(-A)\cdot Q_k(A),\qquad P_{k+1}(A):=Q_k^T(-A)\cdot P(A)\cdot Q_k(A).
\]

Let $b$ and $b_k$ be the numbers of Section~\ref{sec:running-quantities-tori}, and let $y$ be real. The $4\times4$ matrices $b_k+i\gamma_3y$ and $b_k+b\,L^{-1}+i\gamma_3y$ are invertible, because $\gamma_3$ is Hermitian with $\gamma_3^2=1$ and $b_k>0$; put
\begin{align*}
\boldsymbol{\alpha}_y&:=b_k\cdot i\gamma_3y\cdot\left(b_k+i\gamma_3y\right)^{-1},\\
\boldsymbol{\beta}_y&:=b_k^2\cdot b\,L^{-1}\cdot\left(b_k+b\,L^{-1}+i\gamma_3y\right)^{-1}\cdot\left(b_k+i\gamma_3y\right)^{-1}.
\end{align*}
With the matrix $\mathsf{C}$ of Section~\ref{sec:dirac-matrices-lattice}, the following operators are defined whenever the inverses that occur exist:
\begin{align*}
S_k(A)&:=\left(D_A+\mathfrak{m}+b_k\cdot P_k(A)\right)^{-1},\\
H_k(A)&:=b_k\cdot S_k(A)\cdot Q_k^T(-A),\qquad\bar{H}_k(A):=\mathsf{C}^{-1}\cdot H_k(-A)\cdot\mathsf{C},\\
D_k(A)&:=b_k-b_k^2\cdot Q_k(A)\cdot S_k(A)\cdot Q_k^T(-A),\qquad M_k(A):=b_k-D_k(A),\\
\dot{M}_k(A)&:=b_k^2\cdot Q_k(A)\cdot S_k(A)^2\cdot Q_k^T(-A),\\
\Gamma_k(A)&:=\left(D_k(A)+b\,L^{-1}\cdot P(A)\right)^{-1},\qquad\delta Z_k(A):=\det\left(D_k(A)+b\,L^{-1}\cdot P(A)\right),\\
\Gamma_{k,y}(A)&:=\left(D_k(A)+b\,L^{-1}\cdot P(A)+i\gamma_3y\right)^{-1},\\
B_{k,y}(A)&:=\left(b_k+i\gamma_3y\right)^{-1}\cdot\left(1-P(A)\right)+\left(b_k+b\,L^{-1}+i\gamma_3y\right)^{-1}\cdot P(A),\qquad B_k(A):=B_{k,0}(A),\\
S_{k,y}(A)&:=\left(D_A+\mathfrak{m}+\boldsymbol{\alpha}_y\cdot P_k(A)+\boldsymbol{\beta}_y\cdot P_{k+1}(A)\right)^{-1},\qquad S^0_{k+1}(A):=S_{k,0}(A),\\
T_k(A)&:=b_k^{-1}\cdot Q_k(A)\cdot\left(D_A+\mathfrak{m}+b_k\cdot P_k(A)\right),\qquad\bar{T}_k(A):=\mathsf{C}^{-1}\cdot T_k(-A)\cdot\mathsf{C}.
\end{align*}
Here a number stands for the corresponding multiple of the identity, and a $4\times4$ matrix acts at every site. The operators $P_k(A)$, $P_{k+1}(A)$, $S_k(A)$ and $S_{k,y}(A)$ act on the spinors on $T_k$; $H_k(A)$ and $\bar{H}_k(A)$ map the spinors on $T^0_k$ to the spinors on $T_k$, and $T_k(A)$, $\bar{T}_k(A)$ map the spinors on $T_k$ to the spinors on $T^0_k$; $D_k(A)$, $M_k(A)$, $\dot{M}_k(A)$, $\Gamma_k(A)$, $\Gamma_{k,y}(A)$ and $B_{k,y}(A)$ act on the spinors on $T^0_k$, and $\delta Z_k(A)$ is a number.

Every $L$-block of the unit lattice has $L^3$ sites, so that $Q(A)\cdot Q^T(-A)=1$ and $P(A)$ is idempotent; its kernel is scalar, so that $P(A)$ commutes with $\gamma_3$. Consequently
\[
B_{k,y}(A)=\left(b_k+b\,L^{-1}\cdot P(A)+i\gamma_3y\right)^{-1}.
\]
Moreover $\boldsymbol{\alpha}_0=0$, and $\boldsymbol{\beta}_0=b_{k+1}L^{-1}$ by the recursion $b_{k+1}=b\cdot b_k\cdot\left(b_k+b\,L^{-1}\right)^{-1}$ of Section~\ref{sec:running-quantities-tori}.

At the level $0$ one has $\eta_0=1$, and the mass is $\mathfrak{m}=\bar{m}\cdot\eta$. The lattices $T_0$ and $T^0_0$ coincide, and $\psi=\Psi$ for the fields of Section~\ref{sec:grassmann-algebras-norms}. The operators of level $0$ are, $y$ being real,
\begin{gather*}
H_0(A):=1,\qquad T_0(A):=1,\qquad D_0(A):=D_A+\mathfrak{m},\\
\bar{H}_0(A):=\mathsf{C}^{-1}\cdot H_0(-A)\cdot\mathsf{C}=1,\qquad\bar{T}_0(A):=\mathsf{C}^{-1}\cdot T_0(-A)\cdot\mathsf{C}=1,\\
\Gamma_0(A):=\left(D_A+\mathfrak{m}+b\,L^{-1}\cdot P(A)\right)^{-1},\qquad S_{0,y}(A):=\left(D_A+\mathfrak{m}+b\,L^{-1}\cdot P(A)+i\gamma_3y\right)^{-1},\\
S^0_1(A):=S_{0,0}(A)=\Gamma_0(A),\\
\delta Z_0(A):=\det\left(D_A+\mathfrak{m}+b\,L^{-1}\cdot P(A)\right).
\end{gather*}
Let now the operators above be formed with data $(k,l,\mathfrak{m},e_k)$ of the scale-free setting of Section~\ref{sec:properties-scheme}: an integer $k\ge0$, a side $l\in\mathrm{Cl}$, a mass $\mathfrak{m}\in(0,\mathfrak{m}_0]$ (Definition~\ref{def:mass-bound-stop}) and a charge $e_k\in(0,e_0]$ (stage~\ref{it:order-choice-9}), the mass and the charge taking the places of those of level $k$; and let $A\in\tilde{G}_k$ (Definition~\ref{def:domains}). Then every inverse that occurs in the definitions above exists. Indeed, the hypotheses of Theorem~\ref{thm:expansion-walks} and of Proposition~\ref{prop:operators-fermion-1} are satisfied, by stages~\ref{it:order-choice-4}, \ref{it:order-choice-6} and~\ref{it:order-choice-9}, condition~\eqref{eq:renormalisation-group-19} and $\mathfrak{m}_0\le1$. By Theorem~\ref{thm:expansion-walks}(a) the operator $F_y(A)$ of Section~\ref{sec:operators-2} is bijective for every real $y$ and, if $k\ge1$, also for $y=\infty$. For $k\ge1$ this operator is $D_A+\mathfrak{m}+\boldsymbol{\alpha}_y\cdot P_k(A)+\boldsymbol{\beta}_y\cdot P_{k+1}(A)$ if $y$ is real and $D_A+\mathfrak{m}+b_k\cdot P_k(A)$ if $y=\infty$; for $k=0$ it is $D_A+\mathfrak{m}+b\,L^{-1}\cdot P(A)+i\gamma_3y$. Hence, for $k\ge1$, the inverse $S_k(A)$ exists, as in Proposition~\ref{prop:decay-fine-propagator}, and $S_{k,y}(A)$ exists for every real $y$, for $y\ge0$ also by Proposition~\ref{prop:decay-resolvents}; and at the level $0$ the inverses $S_{0,y}(A)$, $y$ real, and $\Gamma_0(A)=S_{0,0}(A)$ exist. For $k\ge1$ and every real $y$ the inverse $\Gamma_{k,y}(A)$ exists by Proposition~\ref{prop:operators-fermion-1}(b); for $y=0$ this is the existence of $\Gamma_k(A)$, which amounts to $\delta Z_k(A)\ne0$ (Propositions~\ref{prop:kernel-bounds-2} and~\ref{prop:operators-fermion-1}(d)). Finally, $\tilde{G}_k$ is invariant under $A\mapsto-A$ by Lemma~\ref{lem:invariances}, so that the inverse $S_k(-A)$, which enters $\bar{H}_k(A)$, exists as well.

\subsubsection{Slots of the connection and covering tori}\label{sec:slots-connection-covering}

Several functions of the scheme depend on two connections, which occupy two distinct places among their arguments, called the \emph{first slot} and the \emph{second slot}. The first slot is occupied by the connection that is integrated, the second by an external connection (Definition~\ref{def:twist-external-connection}), which is a parameter. Thus a polymer activity $\mathscr{E}(X;\mathscr{A},\mathscr{C},\psi,\chi)$ of Definition~\ref{def:renormalisation-group-1} has $\mathscr{A}$ in the first slot and $\mathscr{C}$ in the second, and in the renormalisation step the fermionic operators are evaluated at the sum of the connections in the two slots, according to~\eqref{eq:renormalisation-group-4}.

\begin{definition}[Twist and external connection]\label{def:twist-external-connection}
Let $\theta=\left(\theta_0,\theta_1,\theta_2\right)\in\mathbb{R}^3$. The \emph{twist} of level $k$ is the constant real connection $c^{(k)}_{\theta}$ that takes the value $\frac{\theta_{\mu}}{e_k\cdot l_k}$ on every positively oriented bond of direction $\mu$, $\mu=0,1,2$. Since $l_0=n$, the charge of level $0$ is $q$, and $e_{k+1}\cdot l_{k+1}=L^{-\frac{1}{2}}\cdot e_k\cdot l_k$ (Section~\ref{sec:running-quantities-tori}), one has
\[
c^{(0)}_{\theta}=H_{\theta},\qquad L^{-\frac{1}{2}}\cdot c^{(k+1)}_{\theta}=c^{(k)}_{\theta},
\]
where $H_{\theta}$ is the one-cochain of Section~\ref{sec:lattice-measures-laws}. The twist is the harmonic part of the block field and is carried by the first slot (Section~\ref{sec:renormalisation-step}): with $c_{\theta}:=c^{(k^{*})}_{\theta}$, the terminal block field is $A_{*}+c_{\theta}$ (Proposition~\ref{prop:closed-shifts-terminal}, Theorem~\ref{thm:representation-twist-external}). An \emph{external connection} is a complex one-cochain $a\in C^1\otimes\mathbb{C}$ on $\Lambda_n$ (Section~\ref{sec:lattice-measures-laws}); it is carried by the second slot, at the level $k$ in the form $\mathscr{C}=a^{[k]}$ of Section~\ref{sec:lattice-measures-laws}, and it is a parameter, not a variable of integration. Theorem~\ref{thm:representation-twist-external}, the representation with twist and slot, applies to every $a$ with $a^{[j]}\in\mathfrak{X}_j$ (Definition~\ref{def:domains}) for all $j\le k^{*}$.
\end{definition}

\begin{definition}[Periodisation and projection]\label{def:periodisation-projection}
Let $k$ be a level, let $l_1,l_2\in\mathrm{Cl}$ with $l_2=J\cdot l_1$ for an odd integer $J$, let $T_i:=\mathbb{L}_k/l_i\mathbb{Z}^3$ for $i=1,2$, and let $\mathrm{pj}\colon T_2\to T_1$ be the covering map. An object of level $k$ that depends on the side, such as $V_k$, $\mathrm{Pd}_k$, $\tilde{G}_k$ or $D_k$, carries the argument $(T_i)$ when it is formed with the side $l_i$. A function on the cells of $T_i$ (Definition~\ref{def:barycentres}), an element of $V_k(T_i)$ or a parameter family on $T_i$ is regarded as an $l_i\mathbb{Z}^3$-periodic function on the corresponding objects of $\mathbb{L}_k$, and a spinor on $T_i$ as a function $f$ on $\mathbb{L}_k$ with $f(x+l_ie_{\nu})=-f(x)$, $\nu=0,1,2$. In this reading each object on $T_1$ is an object on $T_2$, as $J$ is odd (for a spinor, $f(x+l_2e_{\nu})=(-1)^Jf(x)=-f(x)$); the resulting inclusion is written $\mathrm{pj}^{*}$. On a parameter family $s$ on $T_1$ it reads $\left(\mathrm{pj}^{*}s\right)_c=s_{\mathrm{pj}(c)}$.

Let a system of primitives on $T_2$ (Section~\ref{sec:primitive-operators}) with property~\ref{it:renormalisation-group-14} be given, and let $\mathrm{Op}^{(2)}$ denote any one of its families. The \emph{projection} of the system to $T_1$ (Definition~\ref{def:projection}) consists of the families $\mathrm{Op}^{\mathrm{pj}}$ defined, for $s\in\mathrm{Pd}_k(T_1)$, $A\in\tilde{G}_k(T_1)$ (Definitions~\ref{def:parameter-cubes-polydisks} and~\ref{def:domains}), $X\in D_k(T_1)$, $y\ge0$ and every function $f$ on $T_1$ on which the family acts, by
\begin{align*}
\mathrm{pj}^{*}\left(\mathrm{Op}^{\mathrm{pj}}(s,A)\cdot f\right)&:=\mathrm{Op}^{(2)}\left(\mathrm{pj}^{*}s,\mathrm{pj}^{*}A\right)\cdot\mathrm{pj}^{*}f&&\text{for the families \eqref{it:primitives-p1}--\eqref{it:primitives-p3} and }\Gamma_k,\\
\mathrm{pj}^{*}\left(\mathrm{Op}^{\mathrm{pj}}(X,A)\cdot f\right)&:=\sum_{X'\in D_k(T_2):\ \mathrm{pj}(X')=X}\mathrm{Op}^{(2)}\left(X',\mathrm{pj}^{*}A\right)\cdot\mathrm{pj}^{*}f&&\text{for }M_k\text{ and }S_{k,y};
\end{align*}
for the families~\eqref{it:primitives-p1} and~\eqref{it:primitives-p2} the argument $A$ is absent. The right sides are defined: $\mathrm{pj}^{*}s\in\mathrm{Pd}_k(T_2)$, the entries of $\mathrm{pj}^{*}s$ being entries of $s$, and $\mathrm{pj}^{*}A\in\tilde{G}_k(T_2)$ by Lemma~\ref{lem:periodic-extension}. For $w\in\mathbb{Z}^3$ let $g_w:=(1,l_1\cdot w)$; as $l_1\in\mathrm{Cl}$ is a multiple of $L\cdot M$, this is an element of $\mathbb{G}^L_{k,3}(T_2)\subset\mathbb{G}_k^L(T_2)$ (Section~\ref{sec:dirac-matrices-lattice}), and it acts on the functions on $T_2$ as the translation by $l_1\cdot w$. It fixes $\mathrm{pj}^{*}s$ and $\mathrm{pj}^{*}A$ and permutes the polymers $X'\in D_k(T_2)$ with $\mathrm{pj}(X')=X$; by property~\ref{it:renormalisation-group-14}, the operators $\mathrm{Op}^{(2)}\left(\mathrm{pj}^{*}s,\mathrm{pj}^{*}A\right)$ and $\sum_{X'}\mathrm{Op}^{(2)}\left(X',\mathrm{pj}^{*}A\right)$ of the display therefore commute with $g_w$. A spinor $h$ on $T_2$ lies in the image of $\mathrm{pj}^{*}$ if and only if $g_wh=(-1)^{w_0+w_1+w_2}\cdot h$ for all $w\in\mathbb{Z}^3$, and a bond function or an element $h$ of $V_k(T_2)$ if and only if $g_wh=h$ for all $w\in\mathbb{Z}^3$; moreover each family maps spinors to spinors, bond functions to bond functions and $V_k$ into $V_k$. Hence these operators map functions in the image of $\mathrm{pj}^{*}$ to functions in the image of $\mathrm{pj}^{*}$, and since the inclusion $\mathrm{pj}^{*}$ is injective, the display determines $\mathrm{Op}^{\mathrm{pj}}(s,A)\cdot f$ and $\mathrm{Op}^{\mathrm{pj}}(X,A)\cdot f$, which are linear in $f$: the projection is well defined.

By Theorems~\ref{thm:fermion-operators-weakened} and~\ref{thm:photon-operators-weakened} there are systems of primitives on the tori whose sides lie in $\mathrm{Cl}$. The primitives of the scheme are the following; $T_k(R)$ denotes the lattice $T_k$ of side $l_k(R)$, and $\mathrm{Tor}(e)$, $R_{\mathrm{top}}$, $k_{\mathrm{top}}$ are those of Definition~\ref{def:tori-stop-level}. On $T_k(R_{\mathrm{top}})$, $k<k_{\mathrm{top}}$, the scheme uses the systems built in Theorem~\ref{thm:properties-weakened-photon} and Proposition~\ref{prop:family-level-zero} (the families~\eqref{eq:operators-photon-21} and~\eqref{eq:operators-photon-22}; at the level $0$, the family $H_0(s')$ of property~\ref{it:renormalisation-group-17}) and in Theorems~\ref{thm:properties-weakened-families} and~\ref{thm:properties-level-zero} (the families of Definitions~\ref{def:weakened-families-level} and~\ref{def:families-level-zero}), put together in Lemma~\ref{lem:assembled-systems-primitives}, with the rate $\kappa'=\kappa_w+\kappa_0$ per unit of $d_M$; here $\kappa'$ is the rate of a parameter set and $\kappa_w$ the number defined in Section~\ref{sec:primitive-operators}, and $\kappa_0$ is the constant of stage~\ref{it:order-choice-1}. For $R\in\mathrm{Tor}(e)$ with $R<R_{\mathrm{top}}$ and for the levels $k<k^{*}(R)$, the primitives on $T_k(R)$ are, by definition, the projections of those on $T_k(R_{\mathrm{top}})$, with $J=\frac{R_{\mathrm{top}}}{R}$, which is a power of $L$ and hence odd; the primitives on $T_k(R_{\mathrm{top}})$ are available at all these levels, because $k^{*}(R)\le k_{\mathrm{top}}$ (Proposition~\ref{prop:depth-stop-level}(d)). By Propositions~\ref{prop:periodisation} and~\ref{prop:coverings-tori}, $\mathrm{pj}^{*}$ intertwines the exact operators on $T_1$ with those on $T_2$.
\end{definition}

\subsubsection{Initial and terminal levels}\label{sec:initial-terminal-levels}

\begin{definition}[The step of level zero]\label{def:step-level-zero}
In the first step, at the level $k=0$, one takes as the families~\eqref{it:primitives-p1} and~\eqref{it:primitives-p3} the diagonal families of property~\ref{it:renormalisation-group-17}, as the localised function $G^{(1)}$ the display~\eqref{eq:renormalisation-group-60} of Section~\ref{sec:localisation-exponent}, and for the determinant the function $\acute{E}_0$ of Definition~\ref{def:contour-functional-functions}. The datum of the first step is a pair $(m,0)$, that is, a pair $(m,\mathscr{E})$ in the sense of Definition~\ref{def:renormalisation-step} with the activity $\mathscr{E}=0$; at the level $0$ every element of the set $\mathrm{Ball}^x_0(4)$ of Section~\ref{sec:spaces-activities} is of this form.
\end{definition}

\begin{definition}[Gauge fixing at the stop level]\label{def:gauge-fixing-stop}
On a torus $R$ the flow of Section~\ref{sec:flow} ends at the level $k^{*}(R)$. Let $A$ be a real block field on $T^0_{k^{*}}$ with mean zero. Proposition~\ref{prop:gauge-fixing-controlled}, applied with $l=l_{*}=M\cdot r_{*}\ge2$, gives a function $\lambda$ on the sites with $\left|A-d\lambda\right|_{\infty}\le2l_{*}^2\cdot\left|dA\right|_{\infty}$. Since $A$ is real, one has $A-d\left(\operatorname{Re}\lambda\right)=\operatorname{Re}\left(A-d\lambda\right)$, and $|\operatorname{Re}z|\le|z|$ for every complex number $z$; so the replacement of $\lambda$ by its real part does not increase $\left|A-d\lambda\right|_{\infty}$. Hence $\lambda_A:=\operatorname{Re}\lambda$ is a real function on the sites, and the real bond function $A^g:=A-d\lambda_A$ satisfies
\[
\left|A^g\right|_{\infty}\le2l_{*}^2\cdot\left|dA\right|_{\infty}.
\]
The \emph{terminal integrand} is the function $A\mapsto Y^x_{*}\left(A+c_{\theta};\mathscr{C}_{*}\right)$ under the integral of Theorem~\ref{thm:representation-twist-external}(c); it is computed at the representative $A^g$ in place of $A$. The gauge invariance of Theorem~\ref{thm:representation-twist-external}(b), applied to the field $A+c_{\theta}$ and the function $-\lambda_A$, permits this: where $\chi_{*}(A)=1$, that is, $|dA|_{\infty}\le p_{*}$ (Definition~\ref{def:coordinates-slice-small}), both $A+c_{\theta}$ and $A^g+c_{\theta}$ lie in $\mathrm{Dom}_{k^{*}}$ by Proposition~\ref{prop:fields-domains}(i), because $d\left(A^g+c_{\theta}\right)=d\left(A+c_{\theta}\right)=dA$, and hence
\[
Y^x_{*}\left(A^g+c_{\theta};\mathscr{C}_{*}\right)=Y^x_{*}\left(A+c_{\theta};\mathscr{C}_{*}\right).
\]
\end{definition}

\subsubsection{Cells of the reference tori}\label{sec:cells-reference-tori}

Let $R^{*}$ be the side of the reference torus (Definition~\ref{def:tori-stop-level}), and put
\[
n_{*}:=\frac{R^{*}}{\eta},\qquad\mathrm{per}:=\frac{n_{*}}{L}.
\]
As $R^{*}=R_{\mathrm{top}}/L=h'\cdot M\cdot r_K\cdot L^{-1}$, $\eta=L^{-N}$ and $h'=a_K=L^{K-N}$ (Section~\ref{sec:running-quantities-tori}), one has
\[
\mathrm{per}=h'\cdot L^N\cdot M\cdot r_K\cdot L^{-2}=L^{K+m-2}\cdot r_K.
\]
Here $K\ge0$ (Definition~\ref{def:tori-stop-level}), $m\ge2$ because $M=L^m\ge L^2$ (stage~\ref{it:order-choice-6}), and $r_K$ is a power of $L$ with a non-negative integer exponent (Section~\ref{sec:running-quantities-tori}); so $\mathrm{per}$ is a power of $L$ with a non-negative exponent, and it is odd, $L$ being odd (stage~\ref{it:order-choice-4}). Let $m_{\mathrm{p}}\ge0$ be the integer with $\mathrm{per}=2m_{\mathrm{p}}+1$. For an integer $k\ge1$, the \emph{torus $k$} is the torus with $k\cdot n_{*}$ sites on each side, that is, of side $k\cdot R^{*}$; in lattice units its sites are the points of $\left(\mathbb{Z}/kn_{*}\mathbb{Z}\right)^3$. The torus $1$ is the reference torus of side $R^{*}$, and the torus $L$, of side $L\cdot R^{*}=R_{\mathrm{top}}$, is the second reference torus.

The site planes $\left\{x_{\mu}\in\mathrm{per}\mathbb{Z}\right\}$ and the link planes $\left\{x_{\mu}\in\mathrm{per}\mathbb{Z}+m_{\mathrm{p}}+\frac{1}{2}\right\}$, $\mu=0,1,2$, cut the torus $k$ into $(2kL)^3$ cells. In each direction they cut $\mathbb{Z}/kn_{*}\mathbb{Z}$ into the intervals
\[
I_{2i}:=\left\{i\cdot\mathrm{per},\dots,i\cdot\mathrm{per}+m_{\mathrm{p}}\right\},\qquad I_{2i+1}:=\left\{i\cdot\mathrm{per}+m_{\mathrm{p}}+1,\dots,(i+1)\cdot\mathrm{per}\right\}\qquad\left(i\in\mathbb{Z}/kL\mathbb{Z}\right),
\]
and the cells are
\[
\mathfrak{c}_c:=\prod_{\mu=0}^2I_{c_{\mu}},\qquad c=\left(c_0,c_1,c_2\right)\in\left(\mathbb{Z}/2kL\mathbb{Z}\right)^3.
\]
The set of the cells of the torus $1$ is $\mathrm{Cel}$; it has $N^{*}:=8L^3$ elements. Since $R^{*}=L\cdot\mathrm{per}\cdot\eta$, one has
\[
l_c:=\frac{R^{*}}{2L}=\left(m_{\mathrm{p}}+\frac{1}{2}\right)\eta.
\]
The base cell is $\mathfrak{c}_0=\left\{0,\dots,m_{\mathrm{p}}\right\}^3$. For $c\in\left(\mathbb{Z}/2kL\mathbb{Z}\right)^3$ let $\mathsf{g}_c$ be the element of the group generated by the reflections in the site planes and the link planes that carries $\mathfrak{c}_0$ onto $\mathfrak{c}_c$; on the coordinate $x_{\mu}$ it acts by
\[
x_{\mu}\mapsto i\cdot\mathrm{per}+x_{\mu}\quad\text{if }c_{\mu}=2i,\qquad x_{\mu}\mapsto(i+1)\cdot\mathrm{per}-x_{\mu}\quad\text{if }c_{\mu}=2i+1.
\]
The pull-back of cochains by a lattice isometry $g$ is denoted by $g^{*}$, and $g_{*}:=\left(g^{-1}\right)^{*}$.

Let $P^o(\mathfrak{c}_c)$ be the set of the plaquettes whose vertices all lie in $\mathfrak{c}_c$ and which lie in none of the six planes $\left\{x_{\mu}=t\right\}$, $t$ an end point of $I_{c_{\mu}}$, $\mu=0,1,2$. The \emph{cell algebra} $\mathfrak{A}(\mathfrak{c}_c)$ consists of the Borel functions $\Phi$ of $\left(F_p\right)_{p\in P^o(\mathfrak{c}_c)}$, $F$ a real two-cochain on the torus $k$, for which there is a number $C_{\Phi}$ with
\[
\left|\Phi(F)\right|\le C_{\Phi}\cdot\exp\left(C_{\Phi}\cdot|F|\right)
\]
for all $F$; put $\mathfrak{A}:=\mathfrak{A}(\mathfrak{c}_0)$. The \emph{mirror copy} of $O\in\mathfrak{A}$ in the cell $\mathfrak{c}_c$ (Section~\ref{sec:reference-tori-cells}) is
\[
O^{[c]}(F):=\begin{cases}O\left(\mathsf{g}_c^{*}F\right)&\text{if }\sum_{\mu}c_{\mu}\text{ is even},\\ \overline{O\left(\mathsf{g}_c^{*}F\right)}&\text{if }\sum_{\mu}c_{\mu}\text{ is odd};\end{cases}
\]
the parity of $\sum_{\mu}c_{\mu}$ is well defined, since $2kL$ is even.

\subsection{Domains, norms and the step map}\label{sec:domains-norms-step}

\subsubsection{Domains of bond functions}\label{sec:domains-bond-functions}

Let $A$ be a bond function on $T_k$, with the components $A_{\mu}(x):=A(x,x+\eta_k\cdot e_{\mu})$ of Section~\ref{sec:averaging-fluctuations-gauge}. For $\mu,\nu\in\{0,1,2\}$ its difference quotients are
\[
(\partial_{\nu}A_{\mu})(x):=\eta_k^{-1}\cdot\left(A_{\mu}(x+\eta_k\cdot e_{\nu})-A_{\mu}(x)\right),
\]
and, $\alpha$ being the H\"older order of stage~\ref{it:order-choice-2}, one puts
\[
|\delta_{\alpha}\partial A|_{\infty}:=\sup_{\mu,\nu}\ \sup_{0<|x-x'|\le 1}|x-x'|^{-\alpha}\cdot\left|(\partial_{\nu}A_{\mu})(x')-(\partial_{\nu}A_{\mu})(x)\right|,
\]
the inner supremum being taken over the sites $x$, $x'$ of $T_k$ with $0<|x-x'|\le 1$. Let $|A|_{\infty}$ and $|\partial A|_{\infty}$ denote the suprema of $|A_{\mu}(x)|$ and of $|(\partial_{\nu}A_{\mu})(x)|$ over all directions and all sites; the numbers $|A|_{\infty}$, $|\partial A|_{\infty}$ and $|\delta_{\alpha}\partial A|_{\infty}$ are called the \emph{three norms} of $A$. For a set $W$ of sites, the three norms of $A$ on $W$ are formed in the same way, with only those bonds and those differences that lie inside $W$. For a function $F$ of the plaquettes of $T_k$, or of the unit lattice $T^0_k$, given by its components $F_{\mu\nu}(x)$, $\mu<\nu$, as is $dA$ in Section~\ref{sec:averaging-fluctuations-gauge}, $|F|_{\infty}$ denotes the supremum of the moduli of its values, that is, of $|F_{\mu\nu}(x)|$ over $\mu<\nu$ and over all sites $x$.

\begin{definition}[Domains]\label{def:domains}
Let $0<c\le 1$. All inequalities in this definition are strict. The representations $A_0=A'+d\lambda$ that occur in (b), (c) and (e) are local: the real functions $\lambda$ and $A'$ in such a representation are defined on the enlarged cube in question only. For an $M$-cube $\square$ the cubes $\square^{\dagger}=\square^{\sim L}$ and $\square^{\sim 5}$ are those of Definition~\ref{def:enlarged-cubes}.
\begin{enumerate}
\item[(a)] $R_k$ consists of the complex bond functions $A$ on $T_k$ with
\[
|A|_{\infty}<e_k^{-\frac{3}{4}+\varepsilon},\qquad|\partial A|_{\infty}<e_k^{-\frac{3}{4}+2\varepsilon},\qquad|\delta_{\alpha}\partial A|_{\infty}<e_k^{-\frac{3}{4}+3\varepsilon}.
\]
The three numbers on the right are the \emph{thresholds} of $R_k$. For a real number $t$, the set $t\cdot R_k$ consists of all complex bond functions on $T_k$ whose three norms are smaller than $t$ times the corresponding thresholds. If $t>0$, these are the functions $t\cdot A$ with $A\in R_k$; if $t\le 0$, the set $t\cdot R_k$ is empty.
\item[(b)] $c\cdot\tilde{R}_k$ is the set of the bond functions on $T_k$ of the form $A_0+A_1$, where $A_1\in c\cdot R_k$ and $A_0$ is a real bond function with the following property: for every $M$-cube $\square$, some real function $\lambda$ on the sites of $\square^{\dagger}$ and some real bond function $A'$ on $\square^{\dagger}$ satisfy $A_0=A'+d\lambda$ on every bond of $\square^{\dagger}$, while the norms $|A'|_{\infty}$, $|\partial A'|_{\infty}$ and $|\delta_{\alpha}\partial A'|_{\infty}$, formed on $\square^{\dagger}$, are smaller than $c$ times the corresponding thresholds of $R_k$. The set $1\cdot\tilde{R}_k$ is denoted by $\tilde{R}_k$.
\item[(c)] $\tilde{R}^{-}_k$ is defined as $\tilde{R}_k$, the bounds on $A'$ being those for $c=1$, but with the condition $A_1\in\left(1-e_k^{\varepsilon}\right)\cdot R_k$ in place of $A_1\in R_k$.
\item[(d)] $G_k$ is the set of the complex bond functions $A$ on $T_k$ with $|A|_{\infty}<e_k^{-1+\varepsilon}$.
\item[(e)] $\tilde{G}_k$ is the set of the bond functions on $T_k$ of the form $A_0+A_1$, where $A_1\in G_k$ and $A_0$ is a real bond function with the following property: for every $M$-cube $\square$, some real function $\lambda$ on the sites of $\square^{\sim 5}$ and some real bond function $A'$ on $\square^{\sim 5}$ satisfy $A_0=A'+d\lambda$ on every bond of $\square^{\sim 5}$, and $|A'|_{\infty}<e_k^{-1+\varepsilon}$ on $\square^{\sim 5}$.
\item[(f)] $\mathfrak{X}_k$ is the set of the complex bond functions $\mathscr{C}$ on $T_k$ with
\[
|\mathscr{C}|_{\infty}<e_k^{-1+2\varepsilon},\qquad|\partial\mathscr{C}|_{\infty}<e_k^{-1+3\varepsilon}.
\]
For a real number $t$, the set $t\cdot\mathfrak{X}_k$ consists of all complex bond functions $\mathscr{C}$ on $T_k$ with $|\mathscr{C}|_{\infty}<t\cdot e_k^{-1+2\varepsilon}$ and $|\partial\mathscr{C}|_{\infty}<t\cdot e_k^{-1+3\varepsilon}$.
\item[(g)] $\mathrm{Dom}_k$ is the set of the real bond functions $A$ on $T^0_k$, with arbitrary harmonic part, such that $H_kA\in\frac{1}{4}\tilde{R}_k$, where $H_k$ is the operator of Section~\ref{sec:averaging-fluctuations-gauge}.
\end{enumerate}
\end{definition}

The elements of these sets are bond functions: a splitting $A=A_0+A_1$ with the stated properties, and the local representations of $A_0$, are only required to exist, and they are not part of the element. The three norms are continuous on the finite-dimensional complex vector space formed by all complex bond functions on $T_k$; since the inequalities in (a), (d) and (f) are strict, the sets $R_k$, $G_k$, $\mathfrak{X}_k$ and $c\cdot R_k$ are open in this vector space. The sets $c\cdot\tilde{R}_k$ and $\tilde{G}_k$ are open in it as well (Lemma~\ref{lem:openness}): if $A=A_0+A_1$ is split as in (b), respectively (e), and $D$ is a complex bond function for which $A_1+D$ belongs to $c\cdot R_k$, respectively to $G_k$, which happens for all $D$ near $0$, then $A_0+(A_1+D)$ is a splitting of $A+D$ of the same kind. Analyticity on one of these open sets means holomorphy on it.

The fermion operators of Section~\ref{sec:exact-fermion-operators}, and those primitives (Section~\ref{sec:primitive-operators}) that have a connection among their arguments, are considered at the connections $A\in\tilde{G}_k$; the set $\tilde{G}_k$ is their domain in the variable $A$. The conditions that define $\tilde{G}_k$ bound sup norms only and impose no bound on the curvature; in particular, the sharp cut-off of an element of $G_k$ outside a polymer, which replaces its values on the bonds outside the polymer by $0$, is again an element of $G_k$.

The domains have the following elementary properties.

\begin{lemma}[Inclusions and bounds]\label{lem:inclusions-bounds}
Let $k\ge 0$, $l\in\mathrm{Cl}$ and $0<e_k\le e_0$.
\begin{enumerate}
\item[(a)] The sets $\tilde{R}_k$ and $\tilde{R}_k+\mathfrak{X}_k$ are contained in $\tilde{G}_k$.
\item[(b)] The set $\frac{1}{4}R_k$ is contained in $\mathfrak{X}_k$.
\item[(c)] Let $B$ be a complex bond function on $T_k$ with $|B|_{\infty}\le 2e_k^{-\frac{3}{4}+\varepsilon}$, let $s\in\mathbb{C}$ with $|s|\le\frac{1}{4}e_k^{-\frac{1}{4}}$, and let $\mathscr{C}\in\mathfrak{X}_k$. Then $s\cdot B+\mathscr{C}\in G_k$.
\item[(d)] Every $\mathscr{T}\in\tilde{R}_k+\mathfrak{X}_k$ satisfies
\[
e_k\cdot|\operatorname{Im}\mathscr{T}|_{\infty}\le e_k^{\frac{1}{4}+\varepsilon}+e_k^{2\varepsilon},\qquad e_k\cdot|d\mathscr{T}|_{\infty}\le 4e_k^{\frac{1}{4}+2\varepsilon}+2e_k^{3\varepsilon}.
\]
\end{enumerate}
\end{lemma}

This is proved as Lemma~\ref{lem:inclusions-domains}.

\begin{lemma}[Rescaling]\label{lem:rescaling}
Let $k\ge 0$, let both $l$ and $l/L$ belong to $\mathrm{Cl}$, and let $0<e_k\le e_0$. The domains of level $k+1$ are those of the torus $T_{k+1}$ of side $l/L$, formed with the number $e_{k+1}:=L^{\frac{1}{2}}\cdot e_k$. For a bond function $\mathscr{A}$ on $T_{k+1}$ let $\mathscr{A}_L$ be the bond function on $T_k$ given by $\mathscr{A}_L(b):=L^{-\frac{1}{2}}\cdot\mathscr{A}(b/L)$ (Section~\ref{sec:local-parts-adjustment}).
\begin{enumerate}
\item[(a)] If $\mathscr{A}\in\tilde{R}_{k+1}$, then $\mathscr{A}_L\in L^{-\frac{3}{4}}\cdot\tilde{R}_k$; moreover $L^{-\frac{3}{4}}\cdot\tilde{R}_k\subset\tilde{R}^{-}_k$.
\item[(b)] If $\mathscr{C}\in\mathfrak{X}_{k+1}$, then $\mathscr{C}_L\in L^{-\frac{3}{4}}\cdot\mathfrak{X}_k$.
\end{enumerate}
\end{lemma}

This is a special case of Lemma~\ref{lem:scaling-domains}.

\begin{lemma}[Invariances]\label{lem:invariances}
Let $k\ge 0$, $l\in\mathrm{Cl}$ and $0<e_k\le e_0$. Each of the sets $R_k$, $c\cdot\tilde{R}_k$ with $0<c\le 1$, $\tilde{R}^{-}_k$, $G_k$, $\tilde{G}_k$ and $\mathfrak{X}_k$ of Definition~\ref{def:domains} is invariant under the action of $\mathbb{G}_k$ (Section~\ref{sec:dirac-matrices-lattice}), under $A\mapsto-A$ and under $A\mapsto\overline{A}$; under hypothesis~\ref{def:hypothesis-one-step}, the set $\mathrm{Dom}_k$ is invariant under the same three transformations. In addition, each of the sets $c\cdot\tilde{R}_k$, $\tilde{R}^{-}_k$ and $\tilde{G}_k$ contains, with every element $A$, all functions $A-d\lambda$, where $\lambda$ is a real-valued function of the sites of $T_k$ (gauge transformations), and all sums $A+\Omega$, where $\Omega$ is a real constant bond function (since $l\in\mathrm{Cl}$, and $L\ge5$ by stage~\ref{it:order-choice-4}, one has $11M\le(2L+1)\cdot M<13L\cdot M\le l$; hence the cubes $\square^{\dagger}$ and $\square^{\sim 5}\subset\square^{\dagger}$, of the sides $(2L+1)\cdot M$ and $11M$, do not wrap in the sense of Definition~\ref{def:terminology}(9), each of them, of side $s<l$, having as a lift a cube of $\mathbb{L}_k$ of sup-diameter $s-\eta_k<l-\eta_k$; two points of such a lift whose images are nearest neighbours in $T_k$ differ by $\pm\eta_k\cdot e_{\mu}+v$ with $v\in l\mathbb{Z}^3$, and $v\ne0$ would make their sup-distance at least $l-\eta_k$; so every bond of $T_k$ with both end points in one of these cubes is the image of a bond of the lifted cube, and on these bonds $\Omega=d\lambda_{\Omega}$ for the real function $\lambda_{\Omega}(x):=\sum_{\mu}\Omega_{\mu}\cdot\tilde{x}_{\mu}$, where $\Omega_{\mu}$ is the constant value of the component $\Omega_{\mu}(\cdot)$ and $\tilde{x}\in\mathbb{L}_k$ is the point of the lift above the site $x$ of the cube).
\end{lemma}

This is Lemma~\ref{lem:symmetries-domains}, where the proof is given.

\begin{lemma}[Periodic extension]\label{lem:periodic-extension}
Let $k$, $l_1$, $l_2$, $T_1$, $T_2$ and $\mathrm{pj}$ be as in Definition~\ref{def:periodisation-projection}, let $e_k>0$, and let the sets of Definition~\ref{def:domains} on $T_1$ and on $T_2$ be formed with this level $k$ and this number $e_k$. Then the map $\mathrm{pj}^{*}$ carries each of the sets $R_k$, $c\cdot\tilde{R}_k$ with $0<c\le 1$, $\tilde{R}^{-}_k$, $G_k$, $\tilde{G}_k$ and $\mathfrak{X}_k$ of $T_1$ into the set of the same name of $T_2$.
\end{lemma}

The same statement reappears as Lemma~\ref{lem:periodisation}, with its proof.

\begin{lemma}[External connections]\label{lem:external-connections}
Assume hypothesis~\ref{def:hypothesis-flow}. For an external connection $a\in C^1\otimes\mathbb{C}$ (Definition~\ref{def:twist-external-connection}) put
\[
|a|_{\mathrm{ext}}:=\left|a^{[k^{*}]}\right|_{\infty}+\left|\partial a^{[k^{*}]}\right|_{\infty}+\left|\delta_{\alpha}\partial a^{[k^{*}]}\right|_{\infty},
\]
the total of the three norms of the function $a^{[k^{*}]}$ of Section~\ref{sec:lattice-measures-laws}. If $|a|_{\mathrm{ext}}<e_{*}^{-1+3\varepsilon}$, then $a^{[k]}\in\mathfrak{X}_k$ for every $k\le k^{*}$.
\end{lemma}

\begin{proof}
Let $0\le k\le k^{*}$. The sequence $a_k=L^{k-N}$ is increasing in $k$, and therefore $e_k=e\cdot a_k^{\frac{1}{2}}\le e_{*}$. Further,
\[
e_{*}\le e_K\le e_1\le e_0<1,
\]
where the four inequalities hold, in this order, by Section~\ref{sec:running-quantities-tori}, by hypothesis~\ref{def:hypothesis-flow}, by stage~\ref{it:order-choice-10} and by condition~\eqref{eq:renormalisation-group-A8}. With the inequalities $e_k\le e_{*}\le e_0<1$ so established, the computation in the proof of Lemma~\ref{lem:scaling-external-connection}, which derives the norms of $a^{[k]}$ from those of $a^{[k^{*}]}$ by scaling and compares them with the thresholds of $\mathfrak{X}_k$, applies on the torus $R$ of hypothesis~\ref{def:hypothesis-flow}; it gives $a^{[k]}\in\mathfrak{X}_k$.
\end{proof}

\subsubsection{Pairs of sites and H\"older derivatives}\label{sec:pairs-sites-holder}

\begin{definition}[Pairs of sites; H\"older derivative]\label{def:pairs-sites-holder}
Let $l\in\mathrm{Cl}$; then $l\ge 13L\cdot M>2M$. Put $l_H:=M$, and let $\mathrm{Pz}_k$ be the set of the pairs $(x,x')$ of sites of $T_k$, taken in this order, with $0<|x-x'|\le l_H$. Let $(x,x')\in\mathrm{Pz}_k$. Its shortest lift, a pair of points of $\mathbb{L}_k$ that project to $x$ and $x'$ and have the sup-distance $|x-x'|$ from each other, is unique up to a common translation by a vector of $l\mathbb{Z}^3$, because $2M<l$. For each of the six orders $\pi$ of the coordinates, $\Gamma^{\pi}(x,x')$ denotes the rectilinear path from $x$ to $x'$ that brings the coordinates to their final values in the order $\pi$ (Section~\ref{sec:averaging-fluctuations-gauge}), formed with this lift. For a bond function $B$ on $T_k$ the number $(\tau B)(x,x')$ is defined by
\[
\eta_k\cdot(\tau B)(x,x'):=\frac{1}{3!}\sum_{\pi}\ \sum_{b\in\Gamma^{\pi}(x,x')}\eta_k\cdot B(b),
\]
and the \emph{H\"older derivative} of a spinor function $f$ along $B$ by
\[
(\delta_{\alpha,B}f)(x,x'):=|x-x'|^{-\alpha}\cdot\left(\exp\left\{i\cdot e_k\cdot\eta_k\cdot(\tau B)(x,x')\right\}\cdot f(x')-f(x)\right).
\]
In this display $f$ is regarded as an antiperiodic function on $\mathbb{L}_k$ (Section~\ref{sec:grassmann-algebras-norms}), and $f(x)$, $f(x')$ are its values at the two points of the shortest lift of $(x,x')$ with which the paths $\Gamma^{\pi}(x,x')$ are formed. Thus $(\delta_{\alpha,B}f)(x,x')$ is a function of the shortest lift. A common translation of the two points of the lift by $l\cdot e_{\nu}$, $\nu\in\{0,1,2\}$, translates the paths with them and leaves $(\tau B)(x,x')$ unchanged, $B$ being an $l\mathbb{Z}^3$-periodic function on the bonds of $\mathbb{L}_k$; it leaves $|x-x'|$ unchanged as well, and it reverses the signs of both values of $f$. Hence $(\delta_{\alpha,B}f)(x,x')$ changes its sign under this translation, as do the generators $\chi$ (Section~\ref{sec:grassmann-algebras-norms}). As two shortest lifts of $(x,x')$ differ by a common translation by a vector of $l\mathbb{Z}^3$, the modulus $\left|(\delta_{\alpha,B}f)(x,x')\right|$ is a function on $\mathrm{Pz}_k$. On barred fields (Section~\ref{sec:grassmann-algebras-norms}) $\delta_{\alpha,-B}$ takes the place of $\delta_{\alpha,B}$.

A pair $(x,x')\in\mathrm{Pz}_k$ is \emph{admissible} for a polymer $X$ if it belongs to one of the following two kinds:
\begin{enumerate}
\item[(I)]\itemmark{I}\label{it:pairs-sites-holder-I} $x$ and $x'$ belong to $X$, and each of the six rectilinear paths $\Gamma^{\pi}(x,x')$ is contained in $X$;
\item[(II)]\itemmark{II}\label{it:pairs-sites-holder-II} $x'$ is a nearest neighbour of $x$ in $T_k$, and $x\in X$ or $x'\in X$.
\end{enumerate}
For an arbitrary union $Z$ of $M$-cubes, the pairs admissible for $Z$ are defined by the same two conditions, with $Z$ in place of $X$.
\end{definition}

\begin{lemma}[Phase factors]\label{lem:phase-factors}
Let $k\ge 0$, $l\in\mathrm{Cl}$ and $0<e_k\le e_0$; by condition~\eqref{eq:renormalisation-group-A8}, $3M\cdot e_k^{\varepsilon}\le 3M\cdot e_0^{\varepsilon}\le 1$. Let $B\in\tilde{G}_k$ and $(x,x')\in\mathrm{Pz}_k$. Then
\[
\left|\exp\left\{\pm i\cdot e_k\cdot\eta_k\cdot(\tau B)(x,x')\right\}\right|\le\exp\left(3M\cdot e_k\cdot|\operatorname{Im}B|_{\infty}\right)\le\exp(1).
\]
If moreover $|x-x'|\ge 1$, then every spinor function $f$ satisfies $\left|(\delta_{\alpha,B}f)(x,x')\right|\le 4|f|_{\infty}$. Here the norm $|\cdot|$ on $\mathbb{C}^4$ and the number $|f|_{\infty}$, the supremum of $|f|$ over the sites (over the pairs of $\mathrm{Pz}_k$ for a function on $\mathrm{Pz}_k$), are those of the beginning of Section~\ref{sec:fermion-operators-bounds}; the last inequality holds for each of the two norms on $\mathbb{C}^4$ admitted there, as it follows from the display above, the triangle inequality, $|x-x'|^{-\alpha}\le1$ and $\exp(1)+1\le4$.
\end{lemma}

This is proved as Lemma~\ref{lem:phases-differences-along}.

\subsubsection{Grassmann algebras and norms of kernels}\label{sec:grassmann-algebras-norms}

Let $\boldsymbol{\beta}$ denote a spinor index and let $\boldsymbol{\omega}\in\{0,1\}$. The generators are the independent symbols
\begin{gather*}
\psi(\xi),\qquad\xi=(x,\boldsymbol{\beta},\boldsymbol{\omega}),\quad x\in T_k;\\
\chi(\boldsymbol{\zeta}),\qquad\boldsymbol{\zeta}=(x,x',\boldsymbol{\beta},\boldsymbol{\omega}),\quad(x,x')\in\mathrm{Pz}_k;\\
\Psi(u)\quad\text{and}\quad W(u),\qquad u=(y,\boldsymbol{\beta},\boldsymbol{\omega}),\quad y\in T^0_k.
\end{gather*}
A generator is \emph{unbarred} if $\boldsymbol{\omega}=0$ and \emph{barred} if $\boldsymbol{\omega}=1$, and $\psi^{\#}:=(\psi,\chi)$. The symbols $\mathrm{Gr}[\psi,\chi]$, $\mathrm{Gr}[\Psi]$, and the like, denote the Grassmann algebras over $\mathbb{C}$ whose generators are those of the fields named in the brackets. The indices $\xi$ and $\boldsymbol{\zeta}$ carry the weights $\eta_k^3$ and $\eta_k^6$,
\[
\int\mathrm{d}\xi\,g(\xi):=\sum_{\xi}\eta_k^3\cdot g(\xi),\qquad\int\mathrm{d}\boldsymbol{\zeta}\,g(\boldsymbol{\zeta}):=\sum_{\boldsymbol{\zeta}}\eta_k^6\cdot g(\boldsymbol{\zeta}),
\]
and the indices $u$ carry the counting measure.

\paragraph{Antiperiodicity.}
The sites in the arguments of the spinor functions on $T_k$ and of the generators $\psi$, $\chi$ are points of $\mathbb{L}_k$, and those in the arguments of the spinor functions on $T^0_k$ and of the generators $\Psi$, $W$ are points of $\mathbb{Z}^3$. A translation by $l\cdot e_{\nu}$, $\nu=0,1,2$, reverses the sign:
\[
f(x+l\cdot e_{\nu})=-f(x),\qquad\chi(x+l\cdot e_{\nu},x'+l\cdot e_{\nu})=-\chi(x,x'),
\]
for every spinor function $f$; the generators $\psi$, $\Psi$ and $W$ satisfy the relation written for $f$, and the indices $\boldsymbol{\beta}$, $\boldsymbol{\omega}$ are left unchanged. Each pair of $\mathrm{Pz}_k$ enters the arguments of $\chi$ through a shortest lift (Definition~\ref{def:pairs-sites-holder}). Fix, for the translations by the vectors of $l\mathbb{Z}^3$, a fundamental domain in $\mathbb{L}_k$ and one in $\mathbb{Z}^3$. The generators $\psi$, $\Psi$ and $W$ whose site lies in the corresponding fundamental domain, together with the generators $\chi$ whose first site lies in the fundamental domain in $\mathbb{L}_k$, are independent; with their sites read on the torus, they are the generators $\psi(\xi)$, $\chi(\boldsymbol{\zeta})$, $\Psi(u)$ and $W(u)$ of the list, with $x\in T_k$, $(x,x')\in\mathrm{Pz}_k$ and $y\in T^0_k$.

\paragraph{Kernels.}
An element of one of these Grassmann algebras is a finite sum of products of generators with complex coefficients. Its \emph{kernels} are the coefficient functions of such an expansion when it is written as a sum, or as an integral with respect to $\mathrm{d}\xi$ and $\mathrm{d}\boldsymbol{\zeta}$, over the arguments of the generators in the products. By the convention of Section~\ref{sec:antiperiodicity}, a kernel is regarded as a function of arguments in $\mathbb{L}_k$, an argument of $\chi$ being a pair of points of $\mathbb{L}_k$ that forms a shortest lift, respectively of arguments $u$ with sites in $\mathbb{Z}^3$ for the generators $\Psi$ and $W$; and it is antiperiodic in each argument, that is, it changes its sign when the site, or the pair of sites, of one argument is translated by $l\cdot e_{\nu}$. Hence the product of a kernel with the generators of its arguments is invariant under these translations; it is a function on the torus.

\paragraph{Norms.}
Let $E=(E_{nm})_{n,m\ge 0}$ be a family of kernels $E_{nm}(\xi_1,\dots,\xi_n;\boldsymbol{\zeta}_1,\dots,\boldsymbol{\zeta}_m)$, each of them antisymmetric in $\xi_1,\dots,\xi_n$ and, separately, in $\boldsymbol{\zeta}_1,\dots,\boldsymbol{\zeta}_m$. For $h=(h_1,h_2)$ with $h_1,h_2>0$ put
\[
\|E_{nm}\|:=\int\mathrm{d}\xi\,\mathrm{d}\boldsymbol{\zeta}\,|E_{nm}|,\qquad\|E\|_h:=\sum_{n,m}\frac{h_1^n\cdot h_2^m}{n!\cdot m!}\cdot\|E_{nm}\|,
\]
the integral being taken over all $n+m$ arguments. At level $k$ the weights are
\[
h_k:=\left(e_k^{-\frac{1}{4}},\,e_k^{-\frac{1}{4}+\varepsilon}\right).
\]
Let $a>0$ be a real number, and let $F\in\mathrm{Gr}[\Psi]$ be written with antisymmetric kernels $F_n$ as
\[
F=\sum_n\frac{1}{n!}\sum_{u_1,\dots,u_n}F_n(u_1,\dots,u_n)\cdot\Psi(u_1)\cdots\Psi(u_n).
\]
The norm $|F|_a$ is defined by
\[
|F|_a:=\sum_n\frac{a^n}{n!}\sum_{u_1,\dots,u_n}\left|F_n(u_1,\dots,u_n)\right|.
\]
Let the independent indices $u$ be totally ordered. As $F_n$ is antisymmetric, its value at arguments $u_1<\dots<u_n$ is the coefficient of $F$ at the ordered monomial $\Psi(u_1)\cdots\Psi(u_n)$, and
\[
|F|_a=\sum_na^n\sum_{u_1<\dots<u_n}\left|F_n(u_1,\dots,u_n)\right|.
\]
Thus $|F|_a$ is the sum of the absolute values of the coefficients of $F$ with respect to the basis formed by the ordered monomials, each multiplied by $a^n$, where $n$ is the degree of the monomial; in particular, $|F|_a$ is a non-decreasing function of $a$.

\paragraph{Choice of the fundamental domain.}
Replacing the fundamental domain by another one changes the signs of some of the independent generators. Since the product of a kernel with the generators of its arguments is a function on the torus, and since the modulus of a kernel does not change when one of its arguments is translated by a vector of $l\mathbb{Z}^3$, neither the moduli of the kernels nor the norms $\|E_{nm}\|$, $\|E\|_h$ and $|F|_a$ depend on the choice of the fundamental domain. An even element is invariant under $\psi\to-\psi$, that is, under the replacement of every generator by its negative (Section~\ref{sec:antiperiodicity}).

\subsubsection{Spaces of activities}\label{sec:spaces-activities}

Fix a level $k\ge0$ and a side $l$ in the class $\mathrm{Cl}$ of Section~\ref{sec:lattice-geometry}, and let $T_k$ be the fine lattice of level $k$ and side $l$. The generators of $\mathrm{Gr}[\psi,\chi]$ are $\psi(\xi)$, $\xi=(x,\boldsymbol{\beta},\boldsymbol{\omega})$, and $\chi(\boldsymbol{\zeta})$, $\boldsymbol{\zeta}=(x,x',\boldsymbol{\beta},\boldsymbol{\omega})$, as in Section~\ref{sec:grassmann-algebras-norms}. For a site $x$ and a pair $(x,x')\in\mathrm{Pz}_k$, the columns, indexed by the spinor index $\boldsymbol{\beta}$, of the unbarred generators ($\boldsymbol{\omega}=0$) are written $\psi(x)$ and $\chi(x,x')$, and those of the barred generators ($\boldsymbol{\omega}=1$) are written $\bar{\psi}(x)$ and $\bar{\chi}(x,x')$. In the same way, for a site $y$ of the unit lattice $T^0_k$, the column of the unbarred generators $\Psi(y,\boldsymbol{\beta},0)$ of Section~\ref{sec:grassmann-algebras-norms} is written $\Psi(y)$, and that of the barred generators $\Psi(y,\boldsymbol{\beta},1)$ is written $\bar{\Psi}(y)$.

\begin{definition}\label{def:renormalisation-group-1}
Consider families $\mathscr{E}=\left\{\mathscr{E}(X;\mathscr{A},\mathscr{C},\psi,\chi):\ X\in D_k\right\}$, indexed by the polymers of Section~\ref{sec:lattice-geometry}, in which for every $X$ the map $(\mathscr{A},\mathscr{C})\mapsto\mathscr{E}(X;\mathscr{A},\mathscr{C},\psi,\chi)$ is defined on $\tilde{R}_k\times\mathfrak{X}_k$ (Definition~\ref{def:domains}) and has even values in $\mathrm{Gr}[\psi,\chi]$ of the form
\begin{multline*}
\mathscr{E}(X;\mathscr{A},\mathscr{C},\psi,\chi)=\sum_{n,m\ge0}\frac{1}{n!\cdot m!}\int\mathrm{d}\xi\,\mathrm{d}\boldsymbol{\zeta}\;\mathscr{E}_{nm}\left(X;\mathscr{A},\mathscr{C};\xi_1,\dots,\xi_n;\boldsymbol{\zeta}_1,\dots,\boldsymbol{\zeta}_m\right)\\
\cdot\psi(\xi_1)\cdots\psi(\xi_n)\cdot\chi(\boldsymbol{\zeta}_1)\cdots\chi(\boldsymbol{\zeta}_m);
\end{multline*}
here the integration extends over all the arguments $\xi_1,\dots,\xi_n,\boldsymbol{\zeta}_1,\dots,\boldsymbol{\zeta}_m$, with the weights of Section~\ref{sec:grassmann-algebras-norms}, the kernels $\mathscr{E}_{nm}$ are antisymmetric in the $\xi_i$ and, separately, in the $\boldsymbol{\zeta}_j$, and in each monomial of such a value the barred legs are exactly as many as the unbarred ones. If entries $\psi'$, $\chi'$ other than $\psi$, $\chi$ stand in the last two places, $\psi'$ and $\chi'$ are to be substituted for the generators. The space $K^{x,C}_k$ consists of the families of this kind with the following properties.
\begin{enumerate}
\item[(k1)]\itemmark{k1}\label{it:renormalisation-group-1-k1} (Analyticity.) For fixed $X$, $n$, $m$ and fixed arguments, the kernel $\mathscr{E}_{nm}(X;\mathscr{A},\mathscr{C};\cdot)$ is an analytic function of $(\mathscr{A},\mathscr{C})$ on $\tilde{R}_k\times\mathfrak{X}_k$.
\item[(k2)]\itemmark{k2}\label{it:renormalisation-group-1-k2} (Localisation in $X$.) The kernel $\mathscr{E}_{nm}\left(X;\mathscr{A},\mathscr{C};\xi_1,\dots,\xi_n;\boldsymbol{\zeta}_1,\dots,\boldsymbol{\zeta}_m\right)$ is zero unless the site of every $\xi_i$ lies in $X$ and the pair of sites of every $\boldsymbol{\zeta}_j$ is admissible for $X$ (Definition~\ref{def:pairs-sites-holder}). Moreover, it does not change when $\mathscr{A}$ and $\mathscr{C}$ are altered on bonds without an end point in $X$: of the two connections, only the values $\mathscr{A}(b)$, $\mathscr{C}(b)$ at the bonds $b$ meeting $X$ enter.
\item[(k3)]\itemmark{k3}\label{it:renormalisation-group-1-k3} (Lattice symmetry.) For every $g\in\mathbb{G}_k$, acting as in Section~\ref{sec:dirac-matrices-lattice}, every $X\in D_k$ and every $(\mathscr{A},\mathscr{C})\in\tilde{R}_k\times\mathfrak{X}_k$,
\[
\mathscr{E}(gX;g\mathscr{A},g\mathscr{C},g\psi,g\chi)=\mathscr{E}(X;\mathscr{A},\mathscr{C},\psi,\chi),
\]
and the kernels of the two sides obey the identity that corresponds to this one.
\item[(k4)]\itemmark{k4}\label{it:renormalisation-group-1-k4} (Gauge invariance near $X$.) Let $(\mathscr{A},\mathscr{C})$ and $(\mathscr{A}^1,\mathscr{C}^1)$ be elements of $\tilde{R}_k\times\mathfrak{X}_k$, and let $\boldsymbol{\lambda}$, $\boldsymbol{\lambda}'$ be real functions on the sites of $T_k$ with
\[
\mathscr{A}^1(b)=\mathscr{A}(b)-(d\boldsymbol{\lambda})(b),\qquad\mathscr{C}^1(b)=\mathscr{C}(b)-(d\boldsymbol{\lambda}')(b)
\]
for every bond $b$ at distance at most $1$ from $X$; here a bond counts as being at distance at most $1$ from $X$ when each of its two end points lies within sup-distance $1$ of $X$. Then, with $\boldsymbol{\Lambda}:=\boldsymbol{\lambda}+\boldsymbol{\lambda}'$,
\[
\mathscr{E}\left(X;\mathscr{A}^1,\mathscr{C}^1,\psi^{\boldsymbol{\Lambda}},\chi^{\boldsymbol{\Lambda}}\right)=\mathscr{E}(X;\mathscr{A},\mathscr{C},\psi,\chi),
\]
where the entries $\psi^{\boldsymbol{\Lambda}}$, $\chi^{\boldsymbol{\Lambda}}$ denote the substitution
\begin{align*}
\psi(x)&\mapsto\exp\left(i\cdot e_k\cdot\boldsymbol{\Lambda}(x)\right)\cdot\psi(x),&\bar{\psi}(x)&\mapsto\exp\left(-i\cdot e_k\cdot\boldsymbol{\Lambda}(x)\right)\cdot\bar{\psi}(x),\\
\chi(x,x')&\mapsto\exp\left(i\cdot e_k\cdot\boldsymbol{\Lambda}(x)\right)\cdot\chi(x,x'),&\bar{\chi}(x,x')&\mapsto\exp\left(-i\cdot e_k\cdot\boldsymbol{\Lambda}(x)\right)\cdot\bar{\chi}(x,x').
\end{align*}
Thus barred and unbarred legs acquire opposite phases, and a leg at a pair $(x,x')$ acquires the phase that belongs to its first site $x$.
\item[(k5)]\itemmark{k5}\label{it:renormalisation-group-1-k5} (Invariance under charge conjugation.) For every $X\in D_k$ and every $(\mathscr{A},\mathscr{C})\in\tilde{R}_k\times\mathfrak{X}_k$,
\[
\mathscr{E}(X;-\mathscr{A},-\mathscr{C},\mathsf{C}\psi,\mathsf{C}\chi)=\mathscr{E}(X;\mathscr{A},\mathscr{C},\psi,\chi),
\]
where the entries $\mathsf{C}\psi$, $\mathsf{C}\chi$ mean that the generators are replaced according to
\[
\psi\mapsto\mathsf{C}\cdot\bar{\psi},\qquad\bar{\psi}\mapsto-\mathsf{C}\cdot\psi,\qquad\chi\mapsto\mathsf{C}\cdot\bar{\chi},\qquad\bar{\chi}\mapsto-\mathsf{C}\cdot\chi.
\]
\item[(k6)]\itemmark{k6}\label{it:renormalisation-group-1-k6} (Decay.) With the constant $\kappa$ of stage~\ref{it:order-choice-3} and the norm $\|\cdot\|_{h_k}$ of Section~\ref{sec:grassmann-algebras-norms},
\[
\|\mathscr{E}\|^x_k:=\sup_{X\in D_k}\exp\left(\kappa\cdot d_M(X)\right)\cdot\sup_{(\mathscr{A},\mathscr{C})\in\tilde{R}_k\times\mathfrak{X}_k}\left\|\mathscr{E}(X;\mathscr{A},\mathscr{C})\right\|_{h_k}<\infty.
\]
Since $\|\cdot\|_{h_k}$ is a sum over the bidegrees $(n,m)$, the supremum in $(\mathscr{A},\mathscr{C})$ is taken of the whole sum, not term by term.
\end{enumerate}
Equipped with $\|\cdot\|^x_k$, the space $K^{x,C}_k$ is a complex Banach space: the set $D_k$ is finite; for each $X$, a uniform limit of analytic maps from the open set $\tilde{R}_k\times\mathfrak{X}_k$ to the finite-dimensional algebra $\mathrm{Gr}[\psi,\chi]$ is analytic; and the requirements other than~\eqref{it:renormalisation-group-1-k1} and~\eqref{it:renormalisation-group-1-k6} are linear conditions that persist under pointwise limits. If the dependence on the side is to be shown, the space is written $K^{x,C}_k(l)$.

Those families in $K^{x,C}_k$ that have, in addition, the property below form a real linear subspace, denoted $K^x_k$.
\begin{enumerate}
\item[(k7)]\itemmark{k7}\label{it:renormalisation-group-1-k7} (Real structure.) Let $\mathbf{K}$ denote the antilinear automorphism of the ring $\mathrm{Gr}[\psi,\chi]$ under which each coefficient goes over into its complex conjugate and the generators $\psi$, $\bar{\psi}$, $\chi$, $\bar{\chi}$ go over into $\varkappa\cdot\psi$, $\varkappa\cdot\bar{\psi}$, $\varkappa\cdot\chi$, $\varkappa\cdot\bar{\chi}$, with the matrix $\varkappa$ of Section~\ref{sec:dirac-matrices-lattice}. For every $X\in D_k$ and every $(\mathscr{A},\mathscr{C})\in\tilde{R}_k\times\mathfrak{X}_k$,
\[
\mathbf{K}\left[\mathscr{E}\left(X;-\overline{\mathscr{A}},-\overline{\mathscr{C}}\right)\right]=\mathscr{E}(X;\mathscr{A},\mathscr{C}).
\]
In terms of kernels: transforming the kernels at $\left(-\overline{\mathscr{A}},-\overline{\mathscr{C}}\right)$ with $\varkappa$ and conjugating them gives back the kernels at $(\mathscr{A},\mathscr{C})$.
\end{enumerate}
Suppressing the second slot, that is, putting $\mathscr{C}=0$ everywhere, the same requirements define the complex Banach space $K^C_k$ of families $E=\left\{E(X;\mathscr{A},\psi,\chi):\ X\in D_k\right\}$, $\mathscr{A}\in\tilde{R}_k$, with the norm $\|\cdot\|_k$, and its real linear subspace $K_k$ of the families with~\eqref{it:renormalisation-group-1-k7}. The elements of these four spaces are called activities.
\end{definition}

For $\mathscr{E}\in K^{x,C}_k$ let $\mathscr{E}|_0:=\mathscr{E}(\cdot\,;\cdot\,,0,\cdot\,)$ be the family with the value $0$ in the second slot, that is, $\mathscr{E}|_0(X;\mathscr{A},\psi,\chi):=\mathscr{E}(X;\mathscr{A},0,\psi,\chi)$. As $0\in\mathfrak{X}_k$, and as each requirement on $\mathscr{E}|_0$ is the requirement on $\mathscr{E}$ at $\mathscr{C}=0$ (in~\eqref{it:renormalisation-group-1-k4} at $\mathscr{C}=\mathscr{C}^1=0$), the family $\mathscr{E}|_0$ lies in $K^C_k$, and in $K_k$ if $\mathscr{E}\in K^x_k$. Moreover
\[
\left\|\mathscr{E}|_0\right\|_k\le\|\mathscr{E}\|^x_k,
\]
because the supremum on the left is taken over a subset of the pairs that occur in~\eqref{it:renormalisation-group-1-k6}.

Activities at $(\mathscr{A},\mathscr{C})\in\tilde{R}_k\times\mathfrak{X}_k$ are evaluated at fields that are built from the sum $\mathscr{T}:=\mathscr{A}+\mathscr{C}$, called the total connection: the generators are replaced by
\[
\psi=H_k(\mathscr{T})\Psi,\qquad\chi=\delta_{\alpha,\mathscr{T}}\psi,\qquad\bar{\psi}=\bar{H}_k(\mathscr{T})\bar{\Psi},\qquad\bar{\chi}=\delta_{\alpha,-\mathscr{T}}\bar{\psi},
\]
where $H_k(\cdot)$ and $\bar{H}_k(\cdot)$ are the operators of Section~\ref{sec:exact-fermion-operators}, and $\delta_{\alpha,\mathscr{T}}$, $\delta_{\alpha,-\mathscr{T}}$ are the H\"older derivatives of Definition~\ref{def:pairs-sites-holder}. For $k\ge1$ the operators $H_k(\mathscr{T})$ and $\bar{H}_k(\mathscr{T})$ are formed with the inverses $S_k(\mathscr{T})$ and $S_k(-\mathscr{T})$, respectively. Let the data $(k,l,\mathfrak{m},e_k)$ with which the operators of Section~\ref{sec:exact-fermion-operators} are formed belong to the scale-free setting of Section~\ref{sec:properties-scheme}, that is, let in addition the mass $\mathfrak{m}$ lie in $(0,\mathfrak{m}_0]$ (Definition~\ref{def:mass-bound-stop}) and the charge $e_k$ in $(0,e_0]$ (stage~\ref{it:order-choice-9}). Then $\mathscr{T}\in\tilde{R}_k+\mathfrak{X}_k\subset\tilde{G}_k$ by Lemma~\ref{lem:inclusions-bounds}(a), and both operators exist at $\mathscr{T}$ by the last paragraph of Section~\ref{sec:exact-fermion-operators}. The data of Section~\ref{sec:fermion-integrals-blocking}, where the activities are evaluated in this way, belong to this setting.

Let $c$ be a real number. For $k\ge1$, with the numbers $w^{\mathrm{m}}_k$ and $w^{\mathrm{E}}_k$ of Section~\ref{sec:running-quantities-tori}, let
\[
\mathrm{Ball}_k(c):=\left\{(m,E)\in\mathbb{C}\times K^C_k:\ |m|\le c\cdot w^{\mathrm{m}}_k,\ \|E\|_k\le c\cdot w^{\mathrm{E}}_k\right\},
\]
and let $\mathrm{Ball}^x_k(c)$ be the set of the pairs $(m,\mathscr{E})\in\mathbb{C}\times K^{x,C}_k$ that satisfy the same two inequalities, with $\|\mathscr{E}\|^x_k$ in place of $\|E\|_k$. For $k=0$ both sets are, by definition, equal to
\[
\left\{(m,0):\ m\in\mathbb{C},\ |m|\le c\cdot w^{\mathrm{m}}_0\right\}.
\]

\subsubsection{Local parts, adjustment and change of scale}\label{sec:local-parts-adjustment}

In this subsection $E=\left\{E(X;\mathscr{A},\mathscr{C},\psi,\chi):\ X\in D_k\right\}$ is a family of level $k$ with kernels $E_{nm}$, written as in Definition~\ref{def:renormalisation-group-1}, and $E(X)$ stands for $E(X;\mathscr{A},\mathscr{C},\psi,\chi)$. For a family without the second slot all definitions below are made in the same way, with $\mathscr{C}$ omitted. Sites carry the weight $\eta_k^3$, as the indices $\xi$ do in Section~\ref{sec:grassmann-algebras-norms}: for a set $W$ of sites of $T_k$ and a function $f$ on $W$ with values in a finite-dimensional complex vector space,
\[
\int_W\mathrm{d}x\,f(x):=\sum_{x\in W}\eta_k^3\cdot f(x),\qquad\int\mathrm{d}x\,f(x):=\int_{T_k}\mathrm{d}x\,f(x),
\]
and an integral over several sites is the iterated integral.

\paragraph{Lifts of small polymers.}
As the side $l$ of $T_k$ belongs to $\mathrm{Cl}$ (Section~\ref{sec:spaces-activities}), one has $l\ge13L\cdot M$. Let $X$ be small, that is, $d_M(X)\le L$. By Proposition~\ref{prop:polymers}(t7), two points of $X$ have sup-distance at most $M\cdot\left(d_M(X)+2\right)\le M\cdot(L+2)$, so that $X$ is contained in the ball $\left\{x:\ d(x,x_0)\le M\cdot(L+2)\right\}$ about any of its sites $x_0$. This ball does not wrap (Definition~\ref{def:terminology}(9)), since
\[
2M\cdot(L+2)+M=2L\cdot M+5M<13L\cdot M\le l.
\]
In the next two paragraphs kernels are integrated over their arguments without being multiplied by the corresponding generators, at two places: in the definition of the matrix $\boldsymbol{\alpha}_2(E,X)$ and in the adjustment~\eqref{eq:renormalisation-group-2}. At both places $X$ is small, and all the arguments, together with all the points $p_i$ that the adjustment attaches to chains of $M$-cubes of $X$, are chosen in a single lift of such a ball. The resulting lift of $X$ is the same for every choice of the site $x_0$, up to a translation by a vector of $l\mathbb{Z}^3$; and if the lift is changed for all arguments at once, every kernel whose number of arguments is even keeps its value (Section~\ref{sec:antiperiodicity}).

\paragraph{Removal of the local parts.}
For a small polymer $X$ let $E_{00}(X;0,0)$ be the kernel of $E(X)$ without arguments, taken at $\mathscr{A}=\mathscr{C}=0$, and let $E_{20}(X;0,0;x,y)$ be the $4\times4$ matrix whose entry in the row $\boldsymbol{\beta}$ and the column $\boldsymbol{\beta}'$ is the kernel $E_{20}\left(X;0,0;(x,\boldsymbol{\beta},1),(y,\boldsymbol{\beta}',0)\right)$, which has a barred argument at the site $x$ and an unbarred one at the site $y$. Assume that, for every small $X$, the matrix $E_{20}(X;0,0;x,y)$ vanishes unless both $x$ and $y$ lie in $X$. This is the first clause of the property~\eqref{it:renormalisation-group-1-k2} for the kernel $E_{20}$ at $\mathscr{A}=\mathscr{C}=0$; hence every element of $K^{x,C}_k$ and every element of $K^C_k$ satisfies the assumption. Let $X$ be small. The extraction numbers of $E$ at $X$ are the number $\boldsymbol{\alpha}_0(E,X)$ and the $4\times4$ matrix $\boldsymbol{\alpha}_2(E,X)$ defined by
\[
\boldsymbol{\alpha}_0(E,X):=\operatorname{Vol}(X)^{-1}\cdot E_{00}(X;0,0),\qquad\boldsymbol{\alpha}_2(E,X):=\operatorname{Vol}(X)^{-1}\int\mathrm{d}x\,\mathrm{d}y\;E_{20}(X;0,0;x,y),
\]
where $\operatorname{Vol}(X)=M^3\cdot|X|_M$ (Section~\ref{sec:lattice-geometry}); as an $M$-cube contains $\left(M\cdot L^k\right)^3$ sites, $\operatorname{Vol}(X)=\int_X\mathrm{d}x\,1$. By the assumption on $E_{20}$, only sites $x$, $y$ of $X$ contribute to the integral that defines $\boldsymbol{\alpha}_2(E,X)$, and for these the arguments are taken in the lift of $X$ described above. For an $M$-cube $\square$ put
\[
\varepsilon(E):=-\sum_{X\ \text{small},\ X\supset\square}\boldsymbol{\alpha}_0(E,X),\qquad m(E)\cdot\mathbf{1}:=-\sum_{X\ \text{small},\ X\supset\square}\boldsymbol{\alpha}_2(E,X),
\]
where $\mathbf{1}$ is the unit $4\times4$ matrix. Under hypothesis~\ref{def:hypothesis-one-step} without primitives (Section~\ref{sec:hypotheses-propositions-activities}), Proposition~\ref{prop:symmetries-extraction} holds by Theorem~\ref{thm:renormalisation-group-1}(a); for every $E$ in $K^{x,C}_k$ or in $K^C_k$, Proposition~\ref{prop:symmetries-extraction}(f) then shows that neither of the two sums depends on the choice of $\square$ and that the second one is a multiple of $\mathbf{1}$. Under this hypothesis, and for such $E$, the two sums define the complex numbers $\varepsilon(E)$ and $m(E)$. The family $R\,E$ is obtained from $E$ by subtracting these local parts on the small polymers:
\[
(R\,E)(X):=E(X)-\boldsymbol{\alpha}_0(E,X)\cdot\operatorname{Vol}(X)-\int_X\mathrm{d}x\;\bar{\psi}(x)^T\cdot\boldsymbol{\alpha}_2(E,X)\cdot\psi(x)\qquad\text{for small $X$,}
\]
and $(R\,E)(X):=E(X)$ for large $X$. The two subtracted terms do not depend on the connections.

\paragraph{Adjustment of the bilinear part.}
Fix $(\mathscr{A},\mathscr{C})\in\tilde{R}_k\times\mathfrak{X}_k$, let $\mathscr{T}=\mathscr{A}+\mathscr{C}$ be the total connection, and let $X$ be a small polymer. With $\tau$ as in Definition~\ref{def:pairs-sites-holder} put
\[
u(p,p'):=\exp\left\{i\cdot e_k\cdot\eta_k\cdot(\tau\mathscr{T})(p,p')\right\}\quad\text{for }(p,p')\in\mathrm{Pz}_k,\qquad u(p,p):=1.
\]
For sites $y$, $y_0$ of $X$ consider the chains $(\square_0,\dots,\square_n)$ of $M$-cubes of $X$ such that $\square_0$ contains $y$, $\square_n$ contains $y_0$ and any two consecutive cubes share a face, and whose length $n$ is minimal among all such chains. With each of them associate
\[
p_0:=y,\qquad p_{i+1}:=\text{the centre of }\square_i\quad(0\le i\le n),\qquad p_{n+2}:=y_0,\qquad U_i:=\prod_{0\le j<i}u(p_j,p_{j+1}),
\]
an empty product being $1$, and write $\operatorname{Mn}$ for the arithmetic mean over these chains. Let $\boldsymbol{\Theta}=\boldsymbol{\Theta}_{X,\mathscr{T}}$ be the homomorphism into $\mathrm{Gr}[\psi,\chi]$ of the subalgebra spanned by the monomials in those generators $\psi(\xi)$ whose sites lie in $X$, which acts on these generators, componentwise in the spinor index, by
\begin{align*}
\boldsymbol{\Theta}\,\psi(y)&:=\operatorname{Vol}(X)^{-1}\int_X\mathrm{d}y_0\;\operatorname{Mn}\left\{U_{n+2}\cdot\psi(y_0)-\sum_{i=0}^{n+1}U_i\cdot\left|p_i-p_{i+1}\right|^{\alpha}\cdot\chi(p_i,p_{i+1})\right\},\\
\boldsymbol{\Theta}\,\bar{\psi}(x)&:=\operatorname{Vol}(X)^{-1}\int_X\mathrm{d}x_0\;\operatorname{Mn}\left\{U_{n+2}^{-1}\cdot\bar{\psi}(x_0)-\sum_{i=0}^{n+1}U_i^{-1}\cdot\left|p_i-p_{i+1}\right|^{\alpha}\cdot\bar{\chi}(p_i,p_{i+1})\right\}.
\end{align*}
In the second line the chains lead from the cube of $x$ to the cube of $x_0$, with $p_0=x$ and $p_{n+2}=x_0$, and $U_i^{-1}=\prod_{0\le j<i}u(p_j,p_{j+1})^{-1}$; so the second line results from the first by the replacements $u\to u^{-1}$, $\psi\to\bar{\psi}$, $\chi\to\bar{\chi}$. Summands with $p_i=p_{i+1}$ are dropped from both sums.

Denote by $E^{[2,0]}(X)$ the sum of the terms of $E(X)$ that contain two factors $\psi$ and no factor $\chi$, that is, the component of $E(X)$ of bidegree $(2,0)$. If $E$ has the property~\eqref{it:renormalisation-group-1-k2}, then $E^{[2,0]}(X)$ belongs to the domain of $\boldsymbol{\Theta}_{X,\mathscr{T}}$ for every small $X$. The adjustment $E\mapsto E^{\mathrm{nat}}$ is given, at every $(\mathscr{A},\mathscr{C})$, by
\begin{equation}\label{eq:renormalisation-group-2}
E^{\mathrm{nat}}(X):=\begin{cases}E(X)-E^{[2,0]}(X)+\boldsymbol{\Theta}_{X,\mathscr{T}}\left(E^{[2,0]}(X)\right)&\text{for small $X$,}\\E(X)&\text{for large $X$;}\end{cases}
\end{equation}
it coincides with the adjustment of Definition~\ref{def:transport-mean}.

\paragraph{Change of scale.}
By the definition of $\mathrm{Cl}$ in Section~\ref{sec:lattice-geometry}, the side $l$ of Section~\ref{sec:spaces-activities} is $l=n\cdot L\cdot M$ with an odd integer $n\ge13$. Hence $l/L=n\cdot M$ is a positive integer and a multiple of $M$, and $l/L\ge13M>2M$. The side $l/L$ need not belong to $\mathrm{Cl}$: for $l=13L\cdot M$ one has $l/L=13M$ and $\frac{13M}{L\cdot M}=\frac{13}{L}<13$. The following objects of level $k+1$ are defined for every side that is a positive integer, a multiple of $M$ and larger than $2M$, and they are formed with the side $l/L$: the fine lattice $T_{k+1}=\mathbb{L}_{k+1}/(l/L)\mathbb{Z}^3$ of level $k+1$ and side $l/L$ (Section~\ref{sec:lattice-geometry}, where the fine lattice is defined for every positive integer side); the group $\mathbb{G}_{k+1}$ of Section~\ref{sec:dirac-matrices-lattice}, which requires a side that is a multiple of $M$; the set $\mathrm{Pz}_{k+1}$ of the ordered pairs $(y,y')$ of sites of $T_{k+1}$ with $0<|y-y'|\le M$, the shortest lifts of these pairs, which are unique up to a common translation by a vector of $(l/L)\mathbb{Z}^3$ because $2M<l/L$, and the paths $\Gamma^{\pi}$, the numbers $(\tau B)(y,y')$ for the bond functions $B$ on $T_{k+1}$ and the H\"older derivatives, formed with these lifts as in Definition~\ref{def:pairs-sites-holder}, whose construction uses of the side of the torus only that it exceeds $2M$; and the generators $\psi$, $\chi$ of level $k+1$, antiperiodic under the translations by the vectors $(l/L)\cdot e_{\nu}$, $\nu=0,1,2$, together with the Grassmann algebra $\mathrm{Gr}[\psi,\chi]$ of level $k+1$ (Section~\ref{sec:grassmann-algebras-norms}). The map $x\mapsto x/L$ carries $\mathbb{L}_k$ onto $\mathbb{L}_{k+1}$, since $\eta_{k+1}=\eta_k/L$, and $l\mathbb{Z}^3$ onto $(l/L)\mathbb{Z}^3$; hence it induces a bijection of the sites of $T_k$ onto those of $T_{k+1}$, which carries bonds to bonds and divides the sup distance by $L$. In particular, every $(x,x')\in\mathrm{Pz}_k$ satisfies $0<|x/L-x'/L|\le M/L\le M$, that is, $(x/L,x'/L)\in\mathrm{Pz}_{k+1}$. For a bond function $A$ on $T_{k+1}$ let $A_L$ be the bond function on $T_k$ given by
\[
A_L(b):=L^{-\frac{1}{2}}\cdot A(b/L)\qquad\text{for the bonds $b$ of $T_k$,}
\]
and for the generators $\psi$, $\chi$ of level $k+1$, componentwise and for barred and unbarred generators alike, put
\[
\psi_L(x):=L^{-1}\cdot\psi(x/L),\qquad\chi_L(x,x'):=L^{-1-\alpha}\cdot\chi(x/L,x'/L)
\]
for the sites $x$ and the pairs $(x,x')\in\mathrm{Pz}_k$ of $T_k$. As $(x+l\cdot e_{\nu})/L=x/L+(l/L)\cdot e_{\nu}$, the elements $\psi_L(x)$ and $\chi_L(x,x')$ change sign when $x$, respectively $x$ and $x'$ together, are translated by $l\cdot e_{\nu}$, as the generators of level $k$ do. Thus $A_L$, $\psi_L$ and $\chi_L$ are defined for every $l\in\mathrm{Cl}$. In the remainder of this paragraph the side $l/L$ belongs to $\mathrm{Cl}$ as well, and $D_{k+1}$ denotes the set of the polymers of $T_{k+1}$ (Section~\ref{sec:lattice-geometry}, where cubes and polymers are defined for sides in $\mathrm{Cl}$). The reblocked family $L\,E$ is the family of level $k+1$ given by
\[
(L\,E)(Y;\mathscr{A},\mathscr{C},\psi,\chi):=\sum_{X\in D_k:\ \bar{X}^L=LY}E\left(X;\mathscr{A}_L,\mathscr{C}_L,\psi_L,\chi_L\right)\qquad\left(Y\in D_{k+1}\right),
\]
where $LY:=\{Ly:\ y\in Y\}$ and $\bar{X}^L$ is the union of the $LM$-cubes that contain an $M$-cube of $X$ (Section~\ref{sec:lattice-geometry}). If, moreover, $0<e_k\le e_0$ and the domains of level $k+1$ are formed on $T_{k+1}$ with the number $e_{k+1}:=L^{\frac{1}{2}}\cdot e_k$, then the hypotheses of Lemma~\ref{lem:rescaling} are fulfilled, and that lemma, together with the inclusions $\tilde{R}^{-}_k\subset\tilde{R}_k$ and $L^{-\frac{3}{4}}\cdot\mathfrak{X}_k\subset\mathfrak{X}_k$ (Definition~\ref{def:domains}), gives $\left(\mathscr{A}_L,\mathscr{C}_L\right)\in\tilde{R}_k\times\mathfrak{X}_k$ for every $(\mathscr{A},\mathscr{C})\in\tilde{R}_{k+1}\times\mathfrak{X}_{k+1}$. Finally, extraction, adjustment and reblocking are composed to
\[
L\,R\,E:=L\left((R\,E)^{\mathrm{nat}}\right).
\]

\subsubsection{Fermion integrals and blocking}\label{sec:fermion-integrals-blocking}

Throughout this subsection the data $(k,l,\mathfrak{m},e_k)$ are those of the scale-free setting of Section~\ref{sec:properties-scheme}: an integer $k\ge0$, a side $l\in\mathrm{Cl}$, a mass $\mathfrak{m}\in(0,\mathfrak{m}_0]$ and a charge $e_k\in(0,e_0]$, with no relation among them. The lattices $T_k$, $T^0_k$, the operators of Sections~\ref{sec:averaging-fluctuations-gauge} and~\ref{sec:exact-fermion-operators} and the domains of Section~\ref{sec:domains-bond-functions} are formed with these data, the mass and the charge taking the places of those of level $k$, as in the last paragraph of Section~\ref{sec:exact-fermion-operators}; and $\mathrm{Gr}[\Psi]$ is the Grassmann algebra generated by the fields $\Psi$, $\bar{\Psi}$ on $T^0_k$ (Section~\ref{sec:grassmann-algebras-norms}). Wherever objects of level $k+1$ occur, namely the generators $\Psi'$, $\bar{\Psi}'$, the fields on the block centres, the numbers $n_{k+1}$ and $c_k$, the fields $A'$ and $\mathscr{C}'$, the sets $\mathrm{Dom}_{k+1}$ and $\mathfrak{X}_{k+1}$ and the block transformation $\mathbf{T}^x_k$, it is required in addition that $l/L\in\mathrm{Cl}$. The objects of level $k+1$ are formed with the side $l/L$, the mass $L\cdot\mathfrak{m}$ and the charge $e_{k+1}=L^{\frac{1}{2}}\cdot e_k$; in particular $\Psi'$, $\bar{\Psi}'$ denote the generators of level $k+1$, which are fields on $T^0_{k+1}$. The data of the levels $k$ and $k+1$ of hypothesis~\ref{def:hypothesis-one-step}, with or without primitives, belong to this setting, since there $e_k\le e_{k+1}\le e_0$, $L\cdot\mathfrak{m}\le\mathfrak{m}_0$, and $l$ and $l/L$ lie in $\mathrm{Cl}$. So do the data $(k^{*},l_{*},\mathfrak{m}_{*},e_{*})$ of the stop level (Theorem~\ref{thm:scales-stop-levels}(c)), which need not satisfy hypothesis~\ref{def:hypothesis-one-step} and with which the function $y^x_{k^{*}}$ of Theorem~\ref{thm:representation-twist-external} is formed; the paragraph on the action and the density of level $k$ below contains no object of level $k+1$.

\paragraph{Pairings and Berezin integrals.}
For fields $\bar{\psi}$, $\psi$ on $T_k$ put
\[
\left\langle\bar{\psi},\psi\right\rangle:=\sum_{x\in T_k}\eta_k^3\sum_{\boldsymbol{\beta}}\bar{\psi}_{\boldsymbol{\beta}}(x)\cdot\psi_{\boldsymbol{\beta}}(x),
\]
the inner sum being taken over the spinor index $\boldsymbol{\beta}$; for an $M$-cube $\square$, $\left\langle\bar{\psi},\psi\right\rangle_{\square}$ denotes the same sum with $x$ restricted to the sites of $\square$. For an operator $O$ on the spinors on $T^0_k$, with kernel $O(y,y')$, put
\[
\left\langle\bar{\Psi},O\Psi\right\rangle:=\sum_{y,y'\in T^0_k}\bar{\Psi}(y)\cdot O(y,y')\cdot\Psi(y').
\]
The Berezin integral $\int\mathrm{D}\Psi$ over the generators $\Psi$, $\bar{\Psi}$ is normalised by
\[
\int\mathrm{D}\Psi\,\exp\left(-\left\langle\bar{\Psi},O\Psi\right\rangle\right)=\det O
\]
for every such operator $O$; the Berezin integral $\int\mathrm{D}\Psi'$ over the generators $\Psi'$, $\bar{\Psi}'$ is normalised in the same way at level $k+1$.

\paragraph{Fields on the block centres.}
For $x\in L\mathbb{Z}^3/l\mathbb{Z}^3$ put
\[
\Psi'_L(x):=L^{-1}\cdot\Psi'\left(\frac{x}{L}\right),\qquad\bar{\Psi}'_L(x):=L^{-1}\cdot\bar{\Psi}'\left(\frac{x}{L}\right),
\]
and for fields $f$, $g$ on $L\mathbb{Z}^3/l\mathbb{Z}^3$ put
\[
\langle f,g\rangle:=L^3\sum_{x\in L\mathbb{Z}^3/l\mathbb{Z}^3}\ \sum_{\boldsymbol{\beta}}f_{\boldsymbol{\beta}}(x)\cdot g_{\boldsymbol{\beta}}(x);
\]
this is the pairing in the exponent of the block transformation below. Let $c_k:=b^{-4n_{k+1}}$. Then
\[
b\cdot L^{-1}\cdot\left\langle\bar{\Psi}'_L,\Psi'_L\right\rangle=b\cdot\sum_{y'\in T^0_{k+1}}\bar{\Psi}'(y')\cdot\Psi'(y'),
\]
and for all fields $f$, $\bar{f}$ on $L\mathbb{Z}^3/l\mathbb{Z}^3$ with values in the odd part of $\mathrm{Gr}[\Psi]$
\[
c_k\int\mathrm{D}\Psi'\,\exp\left\{-b\cdot L^{-1}\cdot\left\langle\bar{\Psi}'_L-\bar{f},\Psi'_L-f\right\rangle\right\}=1;
\]
indeed, by the translation invariance of the Berezin integral the left side equals $c_k\int\mathrm{D}\Psi'\,\exp\left(-b\cdot\sum_{y'\in T^0_{k+1}}\bar{\Psi}'(y')\cdot\Psi'(y')\right)=c_k\cdot b^{4n_{k+1}}$, the spinors on $T^0_{k+1}$ forming a space of dimension $4n_{k+1}$.

\paragraph{Averaging of spinors over $L$-blocks.}
Let $B$ be a connection on the bonds of the fine lattice $T_k$. For $y\in L\mathbb{Z}^3/l\mathbb{Z}^3$ and a unit site $x$ of the $L$-block $B(y)$ of Section~\ref{sec:lattice-geometry} put
\[
(\tau B)(y,x):=\frac{1}{3!}\sum_{\pi}\ \sum_{b\subset\Gamma^{\pi}(y,x)}B(b),
\]
where, for each of the six orders $\pi$, $\Gamma^{\pi}(y,x)$ is the rectilinear path of the unit lattice from $y$ to $x$ of Section~\ref{sec:averaging-fluctuations-gauge}, and the inner sum is taken over the fine bonds $b$ of which this path consists. Define
\begin{gather*}
\left(Q(B)\Psi\right)(y):=L^{-3}\sum_{x\in B(y)}\exp\left\{i\cdot e_k\cdot\eta_k\cdot(\tau B)(y,x)\right\}\cdot\Psi(x)\qquad\left(y\in L\mathbb{Z}^3/l\mathbb{Z}^3\right),\\
\left(Q^T(-B)v\right)(x):=\exp\left\{-i\cdot e_k\cdot\eta_k\cdot(\tau B)(y,x)\right\}\cdot v(y)\qquad\left(x\in B(y)\right),\qquad P(B):=Q^T(-B)\cdot Q(B).
\end{gather*}

\paragraph{Action and density of level $k$.}
For $B\in\tilde{G}_k$ every inverse that occurs in $H_k(B)$, in $\bar{H}_k(B)=\mathsf{C}^{-1}\cdot H_k(-B)\cdot\mathsf{C}$ and in $D_k(B)$ exists, by the last paragraph of Section~\ref{sec:exact-fermion-operators}, where $-B\in\tilde{G}_k$ by Lemma~\ref{lem:invariances}; let $H^{\#}_k(B):=\left(H_k(B),\delta_{\alpha,B}H_k(B)\right)$ on the unbarred fields and $H^{\#}_k(B):=\left(\bar{H}_k(B),\delta_{\alpha,-B}\bar{H}_k(B)\right)$ on the barred fields. Let $(\mathscr{A},\mathscr{C})\in\tilde{R}_k\times\mathfrak{X}_k$ and $\mathscr{T}:=\mathscr{A}+\mathscr{C}$; then $\mathscr{T}\in\tilde{R}_k+\mathfrak{X}_k\subset\tilde{G}_k$ by Lemma~\ref{lem:inclusions-bounds}(a), whose hypotheses $k\ge0$, $l\in\mathrm{Cl}$ and $0<e_k\le e_0$ are part of the setting of this subsection. Let $\psi^{\#}:=H^{\#}_k(\mathscr{T})\Psi$ stand for the four fields
\[
\psi=H_k(\mathscr{T})\Psi,\qquad\chi=\delta_{\alpha,\mathscr{T}}\psi,\qquad\bar{\psi}=\bar{H}_k(\mathscr{T})\bar{\Psi},\qquad\bar{\chi}=\delta_{\alpha,-\mathscr{T}}\bar{\psi}.
\]
For $m\in\mathbb{C}$ and $\mathscr{E}\in K^{x,C}_k$ (Definition~\ref{def:renormalisation-group-1}) define the elements
\begin{align*}
I^x_k[m,\mathscr{E}](\mathscr{A},\mathscr{C};\Psi)&:=-m\cdot\left\langle\bar{\psi},\psi\right\rangle+\sum_{X\in D_k}\mathscr{E}\left(X;\mathscr{A},\mathscr{C},\psi^{\#}\right),\\
\mathrm{Act}^x_k[m,\mathscr{E}](\mathscr{A},\mathscr{C};\Psi)&:=-\left\langle\bar{\Psi},D_k(\mathscr{T})\Psi\right\rangle+I^x_k[m,\mathscr{E}](\mathscr{A},\mathscr{C};\Psi)
\end{align*}
of $\mathrm{Gr}[\Psi]$; here $\mathscr{E}(X;\mathscr{A},\mathscr{C},\psi^{\#})$ is the element obtained from $\mathscr{E}(X;\mathscr{A},\mathscr{C},\psi,\chi)$ when the fields $\psi^{\#}$ are inserted for the generators. Thus $\psi^{\#}$, $I^x_k$ and $\mathrm{Act}^x_k$ are defined on $\tilde{R}_k\times\mathfrak{X}_k$. Let $\mathfrak{P}_k=\left\{\mathfrak{P}_k(X,\cdot):X\in D_k\right\}$ be a given family of complex functions $\mathfrak{P}_k(X,\cdot)$ on $\mathfrak{X}_k$, otherwise arbitrary. The density $\mathscr{F}^x_k$ and the function $y^x_k$ below are formed with this family, which enters them only through the factor $\exp\left\{\sum_{X\in D_k}\mathfrak{P}_k(X,\mathscr{C})\right\}$, a complex number that depends on the second slot alone. Along a flow the family is the family $\mathfrak{P}_k$ of the trajectory with slot of Section~\ref{sec:flow}, and this family is meant wherever no other family is named; it is defined only in the setting of that section, on a torus $R$ with systems of primitives at the levels $j<k$, except at level $0$, where it is $\mathfrak{P}_0=0$ and no flow is required. For $A\in\mathrm{Dom}_k$ one has $H_kA\in\frac{1}{4}\tilde{R}_k\subset\tilde{R}_k$ (Definition~\ref{def:domains}(g) and (b): bounds with the factor $\frac{1}{4}$ imply those with the factor $1$), so that $(H_kA,\mathscr{C})\in\tilde{R}_k\times\mathfrak{X}_k$ for every $\mathscr{C}\in\mathfrak{X}_k$, and the following objects are defined on $\mathrm{Dom}_k\times\mathfrak{X}_k$. For a number $\boldsymbol{\varepsilon}$, for $A\in\mathrm{Dom}_k$ and for $\mathscr{C}\in\mathfrak{X}_k$ put
\begin{multline*}
\mathscr{F}^x_k[\boldsymbol{\varepsilon},m,\mathscr{E}](A,\mathscr{C};\Psi):=\left(\det D_k(0)\right)^{-1}\\
\cdot\exp\left\{\mathrm{Act}^x_k[m,\mathscr{E}](H_kA,\mathscr{C};\Psi)-\boldsymbol{\varepsilon}\cdot n_k+\sum_{X\in D_k}\mathfrak{P}_k(X,\mathscr{C})\right\}
\end{multline*}
and
\[
y^x_k[\boldsymbol{\varepsilon},m,\mathscr{E}](A;\mathscr{C}):=\int\mathrm{D}\Psi\,\mathscr{F}^x_k[\boldsymbol{\varepsilon},m,\mathscr{E}](A,\mathscr{C};\Psi).
\]
Here $\det D_k(0)\ne0$. For $k\ge1$ this follows from Proposition~\ref{prop:block-covariance-real}, applied at the real connection $0$; that proposition holds in the scale-free setting with $k\ge1$, by Theorem~\ref{thm:renormalisation-group-1}(c). For $k=0$ the operator of level $0$ is $D_0(A)=D_A+\mathfrak{m}$ (Section~\ref{sec:exact-fermion-operators}). In this subsection $D_0$ without argument denotes the Wilson operator $D_A$ of Section~\ref{sec:exact-fermion-operators} at level $0$, taken at $A=0$ (for the side $n$, the operator $D_A$ of Section~\ref{sec:lattice-measures-laws} at $A=0$); it is neither the operator $D_0(A)$ of level $0$ nor the set of the polymers of level $0$, which is what $D_k$ without argument denotes in the sums above (Definition~\ref{def:notation}(g)). Thus $D_0(0)=D_0+\mathfrak{m}$. Since $\operatorname{Re}D_0\ge0$ by Lemma~\ref{lem:identities-block-operators}(e), applied at the real connection $0$, one has $\operatorname{Re}D_0(0)\ge\mathfrak{m}>0$; hence $D_0(0)$ is injective, and therefore invertible. For the data of level $0$ of the flow the mass is $\mathfrak{m}=\bar{m}\cdot\eta$, as in the formula for $y^x_0$ below.

\paragraph{Blocking.}
Let $A'$ be a real bond function on $T^0_{k+1}$ and $\mathscr{C}'$ a complex bond function on $T_{k+1}$. Put
\[
\mathscr{C}:=\mathscr{C}'_L,\qquad\mathscr{B}_k:=H_kA^{\min}_k(A'),\qquad\boldsymbol{\omega}_k:=\boldsymbol{\omega}_k(A'),
\]
where $\mathscr{C}'_L$ is formed with the scaling of Section~\ref{sec:local-parts-adjustment}, and $A^{\min}_k(A')$ and the real function $\boldsymbol{\omega}_k(A')$ are those of Section~\ref{sec:averaging-fluctuations-gauge}, the latter with the normalisation $Q_{k+1}\boldsymbol{\omega}_k=0$ of Theorem~\ref{thm:minimisation-one-step}(b). Let $A\mapsto F(A,\mathscr{C};\Psi)$ be a map with values in $\mathrm{Gr}[\Psi]$ that is defined at the points $A^{\min}_k(A')+\mathrm{Sq}_k\zeta$, $\zeta\in\mathrm{Box}_k$, and for which the map $\zeta\mapsto F\left(A^{\min}_k(A')+\mathrm{Sq}_k\zeta,\mathscr{C};\Psi\right)$ from $\mathrm{Box}_k$ to $\mathrm{Gr}[\Psi]$ is integrable with respect to $\mu_I$ (the space $\mathrm{Gr}[\Psi]$ is a complex vector space of finite dimension, and integrability refers to the coefficients of the monomials). For such a map define
\begin{multline*}
\left(\mathbf{T}^x_kF\right)(A',\mathscr{C}';\Psi'):=c_k\int_{\mathrm{Box}_k}\mathrm{d}\mu_I(\zeta)\int\mathrm{D}\Psi\,\exp\left\{-b\cdot L^{-1}\cdot\left\langle\exp\left(-i\cdot e_k\cdot\boldsymbol{\omega}_k\right)\cdot\bar{\Psi}'_L\right.\right.\\
\left.\left.-Q\left(-\mathscr{B}_k-\mathscr{C}\right)\bar{\Psi},\ \exp\left(i\cdot e_k\cdot\boldsymbol{\omega}_k\right)\cdot\Psi'_L-Q\left(\mathscr{B}_k+\mathscr{C}\right)\Psi\right\rangle\right\}\\
\cdot F\left(A^{\min}_k(A')+\mathrm{Sq}_k\zeta,\mathscr{C};\Psi\right),
\end{multline*}
where $\mathrm{Box}_k$, the standard Gaussian measure $\mu_I$ on $V_k$ and $\mathrm{Sq}_k$ are those of Section~\ref{sec:averaging-fluctuations-gauge}, the function $\boldsymbol{\omega}_k$ is evaluated at the sites of $L\mathbb{Z}^3$, and $\langle\cdot,\cdot\rangle$ is the pairing on $L\mathbb{Z}^3/l\mathbb{Z}^3$ defined above. If, for every $A$, each monomial of $F(A,\mathscr{C};\Psi)$ contains as many barred as unbarred generators, then $\mathbf{T}^x_kF$ does not change when a constant is added to $\boldsymbol{\omega}_k$.

Let the data satisfy hypothesis~\ref{def:hypothesis-one-step} without primitives (Section~\ref{sec:hypotheses-propositions-activities}), let the hypotheses of Proposition~\ref{prop:fields-domains}(ii) be satisfied, let $A'\in\mathrm{Dom}_{k+1}$ and $\mathscr{C}'\in\mathfrak{X}_{k+1}$, and let $A\mapsto F(A,\mathscr{C}'_L;\Psi)$ be a bounded Borel map on $\mathrm{Dom}_k$ whose values lie in the even part of $\mathrm{Gr}[\Psi]$. By Proposition~\ref{prop:fields-domains}(ii) the points $A^{\min}_k(A')+\mathrm{Sq}_k\zeta$, $\zeta\in\mathrm{Box}_k$, lie in $\mathrm{Dom}_k$, so that $\mathbf{T}^x_kF$ is defined, and Lemma~\ref{lem:three-berezin-integrals}(b) gives
\begin{equation}\label{eq:renormalisation-group-3}
\int\mathrm{D}\Psi'\,\left(\mathbf{T}^x_kF\right)(A',\mathscr{C}';\Psi')=\int_{\mathrm{Box}_k}\mathrm{d}\mu_I(\zeta)\int\mathrm{D}\Psi\,F\left(A^{\min}_k(A')+\mathrm{Sq}_k\zeta,\mathscr{C}'_L;\Psi\right).
\end{equation}

\paragraph{Vanishing second slot; level zero.}
When the second slot vanishes, that is, at $\mathscr{C}=0$, and at $\mathscr{C}'=0$ for the block transformation, the objects $\mathscr{F}^x_k$, $y^x_k$ and $\mathbf{T}^x_k$ are written $\mathscr{F}_k$, $y_k$ and $\mathbf{T}_k$; in $\mathscr{F}_k$ and $y_k$ the family $\mathfrak{P}_k$ is one with $\mathfrak{P}_k(X,0)=0$ for every $X\in D_k$. The factor $\exp\left\{\sum_{X\in D_k}\mathfrak{P}_k(X,0)\right\}$ is then equal to $1$, so that $\mathscr{F}_k$ and $y_k$ do not depend on the family, and they contain $\mathscr{E}$ only through $\mathscr{E}|_0$. The families $\mathfrak{P}_k$ of Section~\ref{sec:flow} have this property. At level $0$, where the side is $n$, the mass is $\bar{m}\cdot\eta$ and the family is $\mathfrak{P}_0=0$, and where the data of level $0$ satisfy hypothesis~\ref{def:hypothesis-one-step} without primitives, one has for every real pair $(m_{\mathrm{in}},\boldsymbol{\varepsilon}_{\mathrm{in}})$, every $A_0\in\mathrm{Dom}_0$ and every $a\in\mathfrak{X}_0$, by Lemma~\ref{lem:three-berezin-integrals}(c),
\[
y^x_0\left[\boldsymbol{\varepsilon}_{\mathrm{in}},m_{\mathrm{in}},0\right](A_0;a)=\exp\left(-\boldsymbol{\varepsilon}_{\mathrm{in}}\cdot n^3\right)\cdot\frac{\det\left(D_{A_0+a}+\bar{m}\cdot\eta+m_{\mathrm{in}}\right)}{\det\left(D_0+\bar{m}\cdot\eta\right)}.
\]

\subsubsection{The renormalisation step}\label{sec:renormalisation-step}

\begin{definition}[Renormalisation step]\label{def:renormalisation-step}
Let hypothesis~\ref{def:hypothesis-one-step} hold. The data $(k,l,\mathfrak{m},e_k)$ are those of this hypothesis, and the system of primitives of level $k$ is the one that it provides, a system on $T_k$ in the sense of Section~\ref{sec:primitive-operators} with the properties~\ref{it:renormalisation-group-1}--\ref{it:renormalisation-group-17} of Section~\ref{sec:properties-systems-primitives} for the parameters $\mathrm{par}$; the objects of level $k+1$ are formed with the side $l/L$, the mass $L\cdot\mathfrak{m}$ and the charge $e_{k+1}=L^{\frac{1}{2}}\cdot e_k$. The data of the step are the pairs $(m,\mathscr{E})\in\mathrm{Ball}^x_k(4)$; for such a pair let $m':=m+m(\mathscr{E})$ and $\mathscr{E}^r:=R\,\mathscr{E}$. These are defined: for every $\mathscr{E}\in K^{x,C}_k$ the number $m(\mathscr{E})$ of Section~\ref{sec:local-parts-adjustment} is defined and $R\,\mathscr{E}\in K^{x,C}_k$, by Proposition~\ref{prop:symmetries-extraction}(f), which holds under hypothesis~\ref{def:hypothesis-one-step} without primitives by Theorem~\ref{thm:renormalisation-group-1}(a). Formed with the given system, the step map of level $k$ is defined by the seven displays \eqref{eq:renormalisation-group-56}, \eqref{eq:renormalisation-group-57}, \eqref{eq:renormalisation-group-59}, \eqref{eq:renormalisation-group-58}, \eqref{eq:renormalisation-group-60}, \eqref{eq:renormalisation-group-61} and~\eqref{eq:renormalisation-group-62} of Section~\ref{sec:definition-step-map}, together with the determinant, which is formed according to~\eqref{eq:renormalisation-group-7} in Definition~\ref{def:contour-functional-functions} and to Definition~\ref{def:determinant-functions}. The contour integrals of these displays are taken over the circles $|t|=\rho_1$, $|t|=\rho_2$ and $|t|=\rho_4$, where
\[
\rho_1:=\frac{1}{2}\cdot e_k^{-\frac{3}{4}+\frac{9\varepsilon}{2}},\qquad\rho_2:=e_k^{-\frac{3}{4}+7\varepsilon},\qquad\rho_4:=e_k^{-\frac{1}{4}+3\varepsilon},
\]
and the weakened fluctuation field is
\[
\mathscr{Z}(s'):=H_k(s')\cdot C\cdot C_k^{\frac{1}{2}}(s')\cdot\mathsf{Z}\qquad\left(\mathsf{Z}\in V_k\right),
\]
built from the families~\eqref{it:primitives-p1} and~\eqref{it:primitives-p2} and from the map $C$ of Section~\ref{sec:averaging-fluctuations-gauge}.

The displays contain two connections: $A\in\tilde{R}^{-}_k$, which occupies the first slot, and $\mathscr{C}\in\mathfrak{X}_k$, which occupies the second; their sum is $\mathscr{T}:=A+\mathscr{C}$. For the step from level $k$ to level $k+1$, $(\mathscr{A}',\mathscr{C}')$ being the pair of connections of level $k+1$, one takes $A=\mathscr{A}'_L$ and $\mathscr{C}=\mathscr{C}'_L$. Let $O$ be any of the exact or weakened fermionic operators $M_k(X,\cdot)$, $H_k(s'',\cdot)$, $\bar{H}_k$, $\Gamma_k(s'',\cdot)$, $T^{(i)}$, $\bar{T}^{(i)}$ and $\delta_{\alpha,\cdot}$, or one of the link phases $u(x,x')$ in~\eqref{eq:renormalisation-group-58}, and let $B$ be the connection that $O$ receives through the first slot; for the activity $\mathscr{E}^r$ the connection received through the first slot, again written $B$, is either $A$ or $A+t\cdot\mathscr{Z}(s')$. The operator $O$, the activity and the determinant enter the displays as
\begin{equation}\label{eq:renormalisation-group-4}
O\left(B+\mathscr{C}\right),\qquad\mathscr{E}^r\left(X;B,\mathscr{C},\cdot\right)\qquad\text{and}\qquad\delta E^{\mathrm{det}}_k\left(X;A,\mathscr{C}\right)\qquad\left(X\in D_k\right),
\end{equation}
respectively, with $\delta E^{\mathrm{det}}_k$ as in Definition~\ref{def:determinant-functions}; the cluster expansion, the operation $\mathrm{Loc}$, the contours, the measures and the photon operators are independent of $\mathscr{C}$.

From the displays one obtains $E^{(i),x}$ and $E^{\#,x}_k(m,\mathscr{E})$, and one sets $E^{*,x}_k(m,\mathscr{E}):=L\,E^{\#,x}_k(m,\mathscr{E})$, $L$ being the operation of Section~\ref{sec:local-parts-adjustment}. In the case without slot of Section~\ref{sec:activities-slot}, in which the second connection is $\mathscr{C}=0$, the index $x$ is dropped, and these objects are written $E^{(i)}$, $E^{\#}_k$ and $E^{*}_k(m,E)$.

The integrals over $y$ in~\eqref{eq:renormalisation-group-7} converge absolutely (Definition~\ref{def:contour-functional-functions}), since the properties~\ref{it:renormalisation-group-1}, \ref{it:renormalisation-group-2}, \ref{it:renormalisation-group-7} and~\ref{it:renormalisation-group-12}--\ref{it:renormalisation-group-14} for $\mathrm{par}$ that Proposition~\ref{prop:local-expansion-determinant} requires are among the properties~\ref{it:renormalisation-group-1}--\ref{it:renormalisation-group-17}. For $(m,\mathscr{E})\in\mathrm{Ball}^x_k(4)$, $A\in\tilde{R}^{-}_k$ and $\mathscr{C}\in\mathfrak{X}_k$ the series in~\eqref{eq:renormalisation-group-61} (Section~\ref{sec:fluctuation-integral-new}) converge, by Theorem~\ref{thm:step-small-fields}(a). Thus $\mathrm{Ball}^x_k(4)$ is the set on which the step map is defined.
\end{definition}

\begin{definition}[The contour functional $\oint_{\rho}$; the functions $G^{(1)}$, $G^{(3)}$, $G^{(5)}$ and $\acute{E}_k$]\label{def:contour-functional-functions}
Let $\rho>0$ and $\rho\ne1$. For a continuous function $f$ on the circle $\{t\in\mathbb{C}:|t|=\rho\}$ with values in a complex vector space of finite dimension let
\[
\oint_{\rho}f:=\frac{1}{2\pi i}\oint_{|t|=\rho}\mathrm{d}t\,\frac{f(t)}{t(t-1)}.
\]
In the displays of Section~\ref{sec:definition-step-map} the signs of the functions $G^{(1)}$, $G^{(3)}$, $G^{(5)}$ and of the determinant are fixed as follows. For $k\ge1$ the function $G^{(1)}$ is
\begin{equation}\label{eq:renormalisation-group-5}
G^{(1)}=+\oint_{\rho_1}\left\langle\bar{\Psi},M_k\left(X,\mathscr{T}+t\cdot\mathscr{Z}(s')\right)\Psi\right\rangle.
\end{equation}
For $k=0$ the function $G^{(1)}$ is the function~\eqref{eq:renormalisation-group-60} of the single $M$-cubes (Definition~\ref{def:step-level-zero}), in which every sign is written out: its sum over the $M$-cubes, taken at $\chi=\delta_{\alpha,\mathscr{T}}\psi$ and $\bar{\chi}=\delta_{\alpha,-\mathscr{T}}\bar{\psi}$, is $-\left\langle\bar{\psi},\left(D_{\mathscr{T}+\mathscr{Z}(s')}-D_{\mathscr{T}}\right)\psi\right\rangle$, written as half the sum of two re-indexed forms of this pairing, so that its overall sign is $+$. The functions $G^{(3)}$ and $G^{(5)}$ are built from the function
\begin{equation}\label{eq:renormalisation-group-6}
\square\longmapsto-m'\cdot\left\langle\bar{\psi},\psi\right\rangle_{\square}
\end{equation}
on the $M$-cubes $\square$. The determinant is built from the functions $\acute{E}_k(X,A)$, $X\in D_k$, $A\in\tilde{G}_k$, given by
\begin{align}
\acute{E}_k(X,A)&:=-i\cdot b_k^2\int_0^{\infty}\mathrm{d}y\,\operatorname{Tr}\left[\gamma_3\cdot B_{k,y}(A)\cdot Q_k(A)\cdot S_{k,y}(X,A)\right.\notag\\
&\qquad\qquad\left.\cdot\,Q_k^T(-A)\cdot B_{k,y}(A)\right]\qquad(k\ge1),\notag\\
\acute{E}_0(X,A)&:=\int_1^{\infty}\frac{\mathrm{d}y}{y}\,\operatorname{Tr}\left[\left(D_A+\mathfrak{m}+b\cdot L^{-1}\cdot P(A)\right)\cdot S_{0,y}(X,A)\right]\notag\\
&\qquad\qquad-i\int_0^1\mathrm{d}y\,\operatorname{Tr}\left[\gamma_3\cdot S_{0,y}(X,A)\right].\label{eq:renormalisation-group-7}
\end{align}
For a system of primitives with the properties named in Proposition~\ref{prop:local-expansion-determinant}, these integrals over $y$ converge absolutely; this is established in the proof of Proposition~\ref{prop:local-expansion-determinant}(d1) (Proposition~\ref{prop:polymer-expansion-determinant}).
\end{definition}

\begin{remark}
The sign in~\eqref{eq:renormalisation-group-7} agrees with formula (183) of~\cite{DimockMB} and with formulas (374) and (576) of~\cite{DimockUR}; the latter two place the matrix $\gamma_3$ as a factor to the left of the trace, not under it. Formula (186) of~\cite{DimockMB} has the sign opposite to that of~\eqref{eq:renormalisation-group-7}, and thus also to that of formula (183) of the same paper. The sign $+$ in~\eqref{eq:renormalisation-group-5} is the one that formula (159) of~\cite{DimockUVS}, together with its formulas (169) and (170), gives; formulas (171), (186) and (188) of~\cite{DimockUVS} have the opposite sign.
\end{remark}

\begin{definition}[Determinant functions]\label{def:determinant-functions}
The family~\eqref{it:primitives-p5} is gauge covariant (property~\ref{it:renormalisation-group-13}); no mean over gauge transformations is taken in the following functions. For $X\in D_k$ and $A\in\tilde{G}_k$ let
\begin{align*}
E^d_k(X,A)&:=\frac{1}{4}\left(\acute{E}_k(X,A)+\acute{E}_k(X,-A)+\overline{\acute{E}_k\left(X,\overline{A}\right)}+\overline{\acute{E}_k\left(X,-\overline{A}\right)}\right),\\
E^{\mathrm{det}}_k(X,A)&:=E^d_k(X,A)-E^d_k(X,0),
\end{align*}
and, for connections $\mathscr{A}$, $\mathscr{C}$ such that $\mathscr{A}+\mathscr{C}$ and $\mathscr{C}$ lie in $\tilde{G}_k$,
\[
\delta E^{\mathrm{det}}_k(X;\mathscr{A},\mathscr{C}):=E^d_k(X,\mathscr{A}+\mathscr{C})-E^d_k(X,\mathscr{C}).
\]
\end{definition}

\subsubsection{The flow}\label{sec:flow}

Let $e>0$, let $R\in\mathrm{Tor}(e)$ (Definition~\ref{def:tori-stop-level}) and $N\ge N_{\mathrm{min}}(R,e)$, and let the hypothesis of Theorem~\ref{thm:scales-stop-levels}(a) hold, namely $v_K\ge1$ and $\left(1+\frac{\mathrm{lg}}{v_K}\right)^r\le L$; then the stop level $k^{*}=k^{*}(R)$ exists, and $1\le k^{*}\le K$ (Section~\ref{sec:running-quantities-tori}). At every level $k<k^{*}$ let a system of primitives of level $k$ on $T_k(R)$ be given such that the data $(k,l_k,\bar{m}\cdot a_k,e_k)$, together with this system, satisfy hypothesis~\ref{def:hypothesis-one-step}; at $k=0$ the mass is $\bar{m}\cdot\eta$ and the charge is $q$. The objects of level $k+1$ of that hypothesis are the objects of level $k+1$ on $R$, because $l_{k+1}=l_k/L$, $\bar{m}\cdot a_{k+1}=L\cdot\bar{m}\cdot a_k$ and $e_{k+1}=L^{\frac{1}{2}}\cdot e_k$. In particular the data of level $k^{*}$, which are the data of level $k+1$ of hypothesis~\ref{def:hypothesis-one-step} at $k=k^{*}-1$, belong to the scale-free setting of Section~\ref{sec:properties-scheme}: $l_{*}\in\mathrm{Cl}$, $\mathfrak{m}_{*}\le\mathfrak{m}_0$ and $e_{*}\le e_0$. Hence the objects of Section~\ref{sec:fermion-integrals-blocking} that contain no object of level $k+1$ are defined at every level $k\le k^{*}$, and the others at every level $k<k^{*}$. Under hypothesis~\ref{def:hypothesis-flow} all this holds for the primitives of Definition~\ref{def:periodisation-projection}: Theorem~\ref{thm:scales-stop-levels}(b) gives $v_K\ge8$ and $\left(1+\frac{3\mathrm{lg}}{v_K}\right)^r<L$, hence the hypothesis of Theorem~\ref{thm:scales-stop-levels}(a), and Theorem~\ref{thm:scales-stop-levels}(c) gives hypothesis~\ref{def:hypothesis-one-step} at the levels $k<k^{*}(R)$.

Let $(m_{\mathrm{in}},\boldsymbol{\varepsilon}_{\mathrm{in}})$ be a real pair. Put $\boldsymbol{\varepsilon}^t_0:=\boldsymbol{\varepsilon}_{\mathrm{in}}$, $E_0:=0$, $\mathscr{E}_0:=0$ and $\mathfrak{P}_0:=0$, let the mass parameter of level $0$ be $m_{\mathrm{in}}$, so that the symbol $m_k$ stands for $m_{\mathrm{in}}$ when $k=0$, and define recursively
\begin{align*}
\boldsymbol{\varepsilon}^t_{k+1}&:=L^3\left(\boldsymbol{\varepsilon}^t_k+\varepsilon(E_k)+d\varepsilon^0_k\right),\\
m_{k+1}&:=L\left(m_k+m(E_k)\right),\\
E_{k+1}&:=L\,R\,E_k+L\,E^{\mathrm{det}}_k+E^{*}_k(m_k,E_k),\\
\mathscr{E}_{k+1}&:=L\,R\,\mathscr{E}_k+L\,\delta E^{\mathrm{det}}_k+E^{*,x}_k(m_k,\mathscr{E}_k),\\
\mathfrak{P}_{k+1}(Y,\mathscr{C})&:=\sum_{X\in D_k:\,\bar{X}^L=LY}\left(\mathfrak{P}_k(X,\mathscr{C}_L)+E^{\mathrm{det}}_k(X,\mathscr{C}_L)\right)\qquad\left(Y\in D_{k+1},\ \mathscr{C}\in\mathfrak{X}_{k+1}\right).
\end{align*}
Here $\varepsilon(\cdot)$, $m(\cdot)$ and the operations $R$, $L$ and $L\,R$ are those of Section~\ref{sec:local-parts-adjustment}, $d\varepsilon^0_k$ is the number of Section~\ref{sec:averaging-fluctuations-gauge}, $E^{\mathrm{det}}_k$ and $\delta E^{\mathrm{det}}_k$ are the functions of Definition~\ref{def:determinant-functions}, and $E^{*}_k$, $E^{*,x}_k$ are the maps of Definition~\ref{def:renormalisation-step}.

For $\mathscr{C}\in\mathfrak{X}_{k+1}$ the bond function $\mathscr{C}_L$ of Section~\ref{sec:local-parts-adjustment} lies in $\mathfrak{X}_k$. Indeed $|\mathscr{C}_L|_{\infty}=L^{-\frac{1}{2}}\cdot|\mathscr{C}|_{\infty}$ and $|\partial\mathscr{C}_L|_{\infty}=L^{-\frac{3}{2}}\cdot|\partial\mathscr{C}|_{\infty}$, and with $e_{k+1}=L^{\frac{1}{2}}\cdot e_k$ the inequalities that define $\mathfrak{X}_{k+1}$ give
\[
|\mathscr{C}_L|_{\infty}<L^{-1+\varepsilon}\cdot e_k^{-1+2\varepsilon},\qquad|\partial\mathscr{C}_L|_{\infty}<L^{-2+\frac{3\varepsilon}{2}}\cdot e_k^{-1+3\varepsilon},
\]
where both powers of $L$ are smaller than $1$. Moreover $\mathfrak{X}_k\subset\tilde{G}_k$ for $k<k^{*}$, by $0\in\tilde{R}_k$ and Lemma~\ref{lem:inclusions-bounds}(a), which applies because $l_k\in\mathrm{Cl}$ and $0<e_k\le e_0$ by hypothesis~\ref{def:hypothesis-one-step}. By induction on $k$ it follows that, for every $X\in D_k$, $\mathfrak{P}_k(X,\cdot)$ is a complex function on $\mathfrak{X}_k$, and that the arguments $\mathscr{C}_L$ in the recursion for $\mathfrak{P}_{k+1}$ lie in the domains of $\mathfrak{P}_k(X,\cdot)$ and of $E^{\mathrm{det}}_k(X,\cdot)$. The recursion for $\mathfrak{P}_{k+1}$ does not contain $(m_k,\mathscr{E}_k)$, and of the system of level $k$ it contains only the family~\eqref{it:primitives-p5}, through $E^{\mathrm{det}}_k$. Hence, the data of the levels $j<k^{*}$ being as above, it defines $\mathfrak{P}_k$, $k\le k^{*}$, equally for arbitrary families~\eqref{it:primitives-p5} of the levels $j<k$ for which the functions $E^{\mathrm{det}}_j(X,\cdot)$ of Definition~\ref{def:determinant-functions} are defined on $\tilde{G}_j$; in this generality $\mathfrak{P}_k$ is formed in Proposition~\ref{prop:free-determinant} (Theorem~\ref{thm:renormalisation-group-1}(d)). Finally, $\mathfrak{P}_k(X,0)=0$ for every $k\le k^{*}$ and every $X\in D_k$. This holds for $k=0$, as $\mathfrak{P}_0=0$; and if it holds at $k<k^{*}$, then it holds at $k+1$, because the bond function $\mathscr{C}_L$ vanishes for $\mathscr{C}=0$ and $E^{\mathrm{det}}_k(X,0)=E^d_k(X,0)-E^d_k(X,0)=0$ by Definition~\ref{def:determinant-functions}.

One calls $(\boldsymbol{\varepsilon}^t_k,m_k,E_k)_k$ the \emph{trajectory} and $(\mathscr{E}_k,\mathfrak{P}_k)_k$ the \emph{trajectory with slot}. The terms of index $k+1$ are given by the recursion whenever $(m_k,\mathscr{E}_k)$ belongs to $\mathrm{Ball}^x_k(4)$, the set on which the step map of level $k$ is defined. The trajectory is said to \emph{exist up to level}~$k_1$ when $(m_k,\mathscr{E}_k)\in\mathrm{Ball}^x_k(4)$ holds for every $k<k_1$.

Let $k_1\le k^{*}$, and let the trajectory exist up to level $k_1$. Then for every $k\le k_1$ the terms of index $k$ are defined, $\mathscr{E}_k\in K^{x,C}_k$, $E_k\in K^C_k$, and
\[
E_k=\mathscr{E}_k|_0.
\]
Consequently, for $k<k_1$, the pair $(m_k,E_k)$ lies in $\mathrm{Ball}_k(4)$, the domain of $E^{*}_k$, because $\left\|\mathscr{E}_k|_0\right\|_k\le\|\mathscr{E}_k\|^x_k$ (Section~\ref{sec:spaces-activities}; for $k=0$ both balls are the set of the pairs $(m,0)$ with $|m|\le4\cdot w^{\mathrm{m}}_0$), and $\varepsilon(E_k)=\varepsilon(\mathscr{E}_k)$, $m(E_k)=m(\mathscr{E}_k)$, by Proposition~\ref{prop:symmetries-extraction}(f). Thus the five lines of the recursion are meaningful as soon as $(m_k,\mathscr{E}_k)\in\mathrm{Ball}^x_k(4)$, the condition $(m_{k+1},\mathscr{E}_{k+1})\in\mathrm{Ball}^x_{k+1}(4)$ has a meaning, and the first two lines give the numbers $\boldsymbol{\varepsilon}^{+}$ and $m^{+}$ that Theorem~\ref{thm:step-small-fields}(c) forms with $\mathscr{E}_k$.

The assertion is proved by induction on $k$. For $k=0$ one has $E_0=\mathscr{E}_0=0$. Let $k<k_1$, let the assertion hold at $k$, and let $(m_k,\mathscr{E}_k)\in\mathrm{Ball}^x_k(4)$; hypothesis~\ref{def:hypothesis-one-step} holds at the level $k$ by the setting of this subsection. First, $L\,R\,\mathscr{E}_k\in K^{x,C}_{k+1}$ and $\left(L\,R\,\mathscr{E}_k\right)|_0=L\,R\left(\mathscr{E}_k|_0\right)$, by Proposition~\ref{prop:contraction}, which holds here by Theorem~\ref{thm:renormalisation-group-1}(a). Second, $L\,\delta E^{\mathrm{det}}_k\in K^x_{k+1}\subset K^{x,C}_{k+1}$ by Proposition~\ref{prop:local-expansion-determinant}(d5), whose conditions hold under hypothesis~\ref{def:hypothesis-one-step}: the inequality $l_k\ge M\cdot r_k$, the relations $l_k/L\in\mathrm{Cl}$ and $e_{k+1}\le e_0$ and the properties~\ref{it:renormalisation-group-1}, \ref{it:renormalisation-group-2}, \ref{it:renormalisation-group-7} and~\ref{it:renormalisation-group-12}--\ref{it:renormalisation-group-14} for $\mathrm{par}$ are part of that hypothesis, and the two further inequalities of Proposition~\ref{prop:local-expansion-determinant}(d4) follow from $e_k\le e_0$, as stated there; and $\left(L\,\delta E^{\mathrm{det}}_k\right)|_0=L\,E^{\mathrm{det}}_k$, since the second slot $0$ scales to $0$ and
\[
\delta E^{\mathrm{det}}_k\left(X;\mathscr{A}_L,0\right)=E^d_k\left(X,\mathscr{A}_L\right)-E^d_k(X,0)=E^{\mathrm{det}}_k\left(X,\mathscr{A}_L\right)\qquad\left(\mathscr{A}\in\tilde{R}_{k+1}\right)
\]
by Definition~\ref{def:determinant-functions}. Third, $E^{*,x}_k(m_k,\mathscr{E}_k)\in K^{x,C}_{k+1}$ by Theorem~\ref{thm:step-small-fields}(a), and $E^{*,x}_k(m_k,\mathscr{E}_k)|_0=E^{*}_k\left(m_k,\mathscr{E}_k|_0\right)$ by Theorem~\ref{thm:step-small-fields}(b), which holds on all of $\mathrm{Ball}^x_k(4)$. Fourth, the map $\mathscr{E}\mapsto\mathscr{E}|_0$ is linear. Adding the first three statements one obtains $\mathscr{E}_{k+1}\in K^{x,C}_{k+1}$ and, by the assertion at $k$,
\[
\mathscr{E}_{k+1}|_0=L\,R\,E_k+L\,E^{\mathrm{det}}_k+E^{*}_k(m_k,E_k)=E_{k+1},
\]
which lies in $K^C_{k+1}$ (Section~\ref{sec:spaces-activities}).

\subsection{Properties of the scheme}\label{sec:properties-scheme}

In Sections~\ref{sec:primitive-operators}--\ref{sec:fermion-operators-bounds} the setting is \emph{scale-free}. Let $k\ge0$ be an integer, $l\in\mathrm{Cl}$, $\mathfrak{m}\in(0,\mathfrak{m}_0]$ a real mass and $e_k\in(0,e_0]$; the lattices $T_k$, $T^0_k$, the operators of Sections~\ref{sec:averaging-fluctuations-gauge} and~\ref{sec:exact-fermion-operators} and the domains of Section~\ref{sec:domains-bond-functions} are those built with $(k,l,\mathfrak{m},e_k)$. No relation among $k$, $l$, $\mathfrak{m}$ and $e_k$ is required, and every constant is independent of them. By Lemma~\ref{lem:invariances}, applied to $0\in\tilde{G}_k$, every real constant connection belongs to $\tilde{G}_k$; the twist is therefore covered by this setting.

\subsubsection{Primitive operators}\label{sec:primitive-operators}

\begin{definition}[Primitives]
Let $k\ge1$. A \emph{system of primitives of level $k$ on $T_k$} is given by five families of linear maps, whose indices are taken from $\mathrm{Pd}_k$ (the parameters $s$), from $\tilde{G}_k$ (the connections $A$), from $D_k$ (the polymers $X$) and from $[0,\infty)$ (the variable $y$), as follows.
\begin{enumerate}
\item[(p1)]\itemmark{p1}\label{it:primitives-p1} The operators $H_k(s')$, where $s'=(s'_c)_{c\in\pi'_k}$, from the bond functions on $T^0_k$ to the bond functions on $T_k$.
\item[(p2)]\itemmark{p2}\label{it:primitives-p2} The operators $C_k^{\frac{1}{2}}(s')$ on $V_k$, where $s'=(s'_c)_{c\in\pi''_k}$.
\item[(p3)]\itemmark{p3}\label{it:primitives-p3} The operators $H_k(s'',A)$, where $s''=(s''_c)_{c\in\pi_k}$, from the spinors on $T^0_k$ to the spinors on $T_k$; together with them one has the operators $\bar{H}_k(s'',A):=\mathsf{C}^{-1}\cdot H_k(s'',-A)\cdot\mathsf{C}$.
\item[(p4)]\itemmark{p4}\label{it:primitives-p4} The operators $\Gamma_k(s'',A)$, which depend on the same variables $s''$ as the operators of~\eqref{it:primitives-p3}, and the operators $M_k(X,A)$, all of them on the spinors on $T^0_k$.
\item[(p5)]\itemmark{p5}\label{it:primitives-p5} The operators $S_{k,y}(X,A)$ on the spinors on $T_k$, which are Borel measurable functions of $y$.
\end{enumerate}
At level $0$ a system of primitives is formed by a family~\eqref{it:primitives-p2} together with operators $\Gamma_0(s'',A)$ and operators $S_{0,y}(X,A)$, $y\ge0$, on the spinors on $T_0$, which form the family~\eqref{it:primitives-p5} of level $0$ and are Borel measurable functions of $y$; the families~\eqref{it:primitives-p1} and~\eqref{it:primitives-p3} of level $0$ are those of property~\ref{it:renormalisation-group-17}.
\end{definition}

Families~\eqref{it:primitives-p1} and~\eqref{it:primitives-p2} are provided by Theorem~\ref{thm:photon-operators-weakened}, and families~\eqref{it:primitives-p3}, \eqref{it:primitives-p4} and~\eqref{it:primitives-p5} by Theorem~\ref{thm:fermion-operators-weakened}.

A \emph{parameter set} is a quadruple $\mathrm{par}'=(\alpha_0,\gamma,\kappa';C)$: the number $\alpha_0$ fixes the radius $M^{\alpha_0}$ of the parameter polydisk $\mathrm{Pd}_k$, $\gamma$ is a rate per unit of length, $\kappa'$ a rate per unit of $d_M$, and $C$ a constant. Two parameter sets occur:
\[
\mathrm{par}^0:=\left(\alpha_0,\gamma_w,\kappa_w+\kappa_0;C_w^1\right),\qquad\mathrm{par}:=\left(\alpha_0,\gamma_w,\kappa_w;C_w\right),
\]
with
\[
C_w:=64K_0\cdot C_w^1,\qquad\kappa_w:=\kappa+\kappa_0+3.
\]
Here $K_0=\frac{\exp(\kappa_0)}{36}$ is the absolute constant of stage~\ref{it:order-choice-1}, that is, the constant of Proposition~\ref{prop:polymers}(t2). The set $\mathrm{par}^0$ is the one of Theorems~\ref{thm:fermion-operators-weakened} and~\ref{thm:photon-operators-weakened}, with the numbers $\alpha_0$, $\gamma_w$ and $C_w^1$ of Section~\ref{sec:photon-operators}; the set $\mathrm{par}$ is the one of the systems after projection (property~\ref{it:renormalisation-group-18}) and the one in the hypothesis of Theorem~\ref{thm:step-small-fields}.

\subsubsection{Properties of systems of primitives}\label{sec:properties-systems-primitives}

Let $\mathrm{par}'=(\alpha_0,\gamma,\kappa';C)$ be parameters as in Section~\ref{sec:primitive-operators}: $M^{\alpha_0}$ is the radius of the parameter polydisk $\mathrm{Pd}_k$ of Definition~\ref{def:parameter-cubes-polydisks}, $\gamma$ is a rate of decay per unit of length, $\kappa'$ a rate of decay per unit of $d_M$, and $C$ is a constant.

The kernels of the operators are indexed by sites, bonds and basis elements $\Upsilon$, as in Section~\ref{sec:lattice-geometry}; for an index $x$ and a set $W$ of sites, ``$x$ meets $W$'' and ``$x$ is contained in $W$'' have the meaning given there, and $d(x,x')$ is the sup-distance of the supports of $x$ and $x'$. The kernel of an operator on spinors has $4\times4$ blocks as entries. In the assertions below, $\lvert\cdot\rvert$ is the modulus of a number; for vectors of $\mathbb{C}^4$ it is one of the two norms admitted at the beginning of Section~\ref{sec:fermion-operators-bounds}, the Euclidean norm $\lvert\cdot\rvert_2$ and the largest modulus of a component $\lvert\cdot\rvert_c$,
\[
\lvert v\rvert_2:=\left(\sum_{\boldsymbol{\beta}}\lvert v_{\boldsymbol{\beta}}\rvert^2\right)^{\frac{1}{2}},\qquad\lvert v\rvert_c:=\max_{\boldsymbol{\beta}}\lvert v_{\boldsymbol{\beta}}\rvert\qquad\left(v\in\mathbb{C}^4\right),
\]
and for $4\times4$ blocks it is the operator norm belonging to the norm taken on $\mathbb{C}^4$; for $\lvert\cdot\rvert_c$ this is the largest row sum $\max_{\boldsymbol{\beta}}\sum_{\boldsymbol{\beta}'}\lvert m_{\boldsymbol{\beta}\boldsymbol{\beta}'}\rvert$ of a block $m=\left(m_{\boldsymbol{\beta}\boldsymbol{\beta}'}\right)$. For a function $f$ on a set $Z$ of sites, or of pairs of sites, $\sup_Z\lvert f\rvert$ denotes the supremum of $\lvert f(z)\rvert$ over $z\in Z$. Norms of vectors of $\mathbb{C}^4$ and of blocks occur in the properties~\ref{it:renormalisation-group-5}, \ref{it:renormalisation-group-6} and~\ref{it:renormalisation-group-7}, and in property~\ref{it:renormalisation-group-17}: in the bound for $S_{0,y}(X,A)(x,x')$ displayed there, and in the properties~\ref{it:renormalisation-group-5} and~\ref{it:renormalisation-group-6} as property~\ref{it:renormalisation-group-17} asserts them for $H_0(s'',A)$ and $\Gamma_0$. Each of these assertions is understood for both norms: it is the conjunction of the two assertions that result when every such norm, and every supremum formed with it, is read with $\lvert\cdot\rvert_2$ throughout, respectively with $\lvert\cdot\rvert_c$ throughout, with one and the same constant for the two norms, namely the constant $C$ of $\mathrm{par}'$, respectively, for $H_0(s'',A)$ in property~\ref{it:renormalisation-group-17}, the constant $4\exp(2\gamma)\cdot M^{27\alpha_0}$ named there. Consequently, by the definition of the properties for $\mathrm{par}'$ given below, a system, a family or a set of families has property~\ref{it:renormalisation-group-5}, \ref{it:renormalisation-group-6} or~\ref{it:renormalisation-group-7}, or the part of property~\ref{it:renormalisation-group-17} just described, for $\mathrm{par}'$ if and only if the bounds in question hold for both norms with these constants and the numbers $\gamma$, $\kappa'$ of $\mathrm{par}'$; in particular, each of these properties provides the bounds for either norm. Since $\lvert v\rvert_c\le\lvert v\rvert_2\le2\lvert v\rvert_c$ for $v\in\mathbb{C}^4$ (Section~\ref{sec:norms-pairings-kernels}), such a bound with a constant $C$ for one of the two norms implies the same bound with the constant $2C$ for the other.

The kind of a family with parameters is the kind of its parameter cubes (Definition~\ref{def:terminology}(3)): $\pi'_k$ for the family~\eqref{it:primitives-p1}, $\pi''_k$ for the family~\eqref{it:primitives-p2}, and $\pi_k$ for the family~\eqref{it:primitives-p3} and for $\Gamma_k$. For a set $\mathscr{Y}$ of parameter cubes of one kind, with union $Y$, a link-component of $Y$ is the union of the cubes of a link-component of $\mathscr{Y}$ in the sense of Definition~\ref{def:parameter-cubes-polydisks}.

The families~\eqref{it:primitives-p1} and~\eqref{it:primitives-p2} do not depend on a connection; for them the argument $A$ is omitted. In an assertion on several families, $\mathrm{Op}(s,A)$ denotes the member at $(s,A)$ of one of the families~\eqref{it:primitives-p1}--\eqref{it:primitives-p3} and $\Gamma_k$, that is, $H_k(s)$, $C_k^{\frac{1}{2}}(s)$, $H_k(s,A)$ or $\Gamma_k(s,A)$, the parameter family $s$ being indexed by the parameter cubes of the kind of that family.

For a linear map $O$ from the bond functions on $T^0_k$ to the bond functions on $T_k$, the kernel is defined by
\[
(OA)(b)=\sum_{b'}O(b,b')\cdot A(b'),
\]
where $b$ is a positively oriented bond of $T_k$ and the sum is over the positively oriented bonds $b'$ of $T^0_k$. Let $\mathfrak{l}\in\{0,1,2\}$, and let $b$, $\tilde b$ be positively oriented bonds of $T_k$ of the same direction whose initial points $x$, $\tilde x$ satisfy $0<\lvert x-\tilde x\rvert\le1$. The differences of $O$ in its first argument are
\begin{align*}
\left(\partial_{\mathfrak{l}}O\right)(b,b')&:=\eta_k^{-1}\cdot\left(O\left(b+\eta_k\cdot e_{\mathfrak{l}},b'\right)-O(b,b')\right),\\
\left(\delta_{\alpha}\partial_{\mathfrak{l}}O\right)\left((b,\tilde b);b'\right)&:=\lvert x-\tilde x\rvert^{-\alpha}\cdot\left(\left(\partial_{\mathfrak{l}}O\right)(\tilde b,b')-\left(\partial_{\mathfrak{l}}O\right)(b,b')\right),
\end{align*}
$b+\eta_k\cdot e_{\mathfrak{l}}$ being the translate of $b$ by $\eta_k\cdot e_{\mathfrak{l}}$.

Let a system of primitives of level $k$ on $T_k$ (Section~\ref{sec:primitive-operators}) be given. It has the property stated in one of the items~\ref{it:renormalisation-group-1}--\ref{it:renormalisation-group-17} below for $\mathrm{par}'$ if the assertion of that item holds for all $s\in\mathrm{Pd}_k$, $A\in\tilde{G}_k$, $X\in D_k$ and $y\ge0$, with the numbers $\gamma$, $\kappa'$ and $C$ of $\mathrm{par}'$. Likewise, a family among~\eqref{it:primitives-p1}--\eqref{it:primitives-p5}, or a set of such families, has the property stated in an item for $\mathrm{par}'$ if the part of the assertion of that item that concerns these families holds, under the same quantifiers and with the same numbers. Item~\ref{it:renormalisation-group-17} concerns the level $k=0$ only and imposes no condition at the levels $k\ge1$; thus every system of level $k\ge1$ has property~\ref{it:renormalisation-group-17}. By Lemma~\ref{lem:invariances}, the set $\tilde{G}_k$ contains, together with $A$, the connections $A-d\lambda$ for real functions $\lambda$ on the sites, $gA$ for $g\in\mathbb{G}_k$, $-A$ and $\overline{A}$; so all operators in the relations below are defined.
\begin{enumerate}
\item\label{it:renormalisation-group-1} \emph{Normalisation.} Let $1$ also denote the parameter family all of whose entries are equal to $1$. Then
\[
H_k(1)=H_k,\qquad C_k^{\frac{1}{2}}(1)=C_k^{\frac{1}{2}},\qquad H_k(1,A)=H_k(A),\qquad \Gamma_k(1,A)=\Gamma_k(A),
\]
\[
\sum_{X\in D_k}M_k(X,A)=M_k(A),\qquad \sum_{X\in D_k}S_{k,y}(X,A)=S_{k,y}(A),
\]
where the right-hand sides are the exact operators of Sections~\ref{sec:averaging-fluctuations-gauge} and~\ref{sec:exact-fermion-operators}.

\item\label{it:renormalisation-group-2} \emph{Analyticity.} For each of the families~\eqref{it:primitives-p1}--\eqref{it:primitives-p3} and $\Gamma_k$, the map $(s,A)\mapsto\mathrm{Op}(s,A)$, which for the families~\eqref{it:primitives-p1} and~\eqref{it:primitives-p2} does not depend on $A$, is continuous on $\mathrm{Pd}_k\times\tilde{G}_k$ and analytic in $(s,A)$ on the interior of this set; and for every fixed $s\in\mathrm{Pd}_k$, the points of the boundary of $\mathrm{Pd}_k$ included, the map $A\mapsto\mathrm{Op}(s,A)$ is analytic on $\tilde{G}_k$. For all $X\in D_k$ and $y\ge0$ the maps $A\mapsto M_k(X,A)$ and $A\mapsto S_{k,y}(X,A)$ are analytic on $\tilde{G}_k$.

\item\label{it:renormalisation-group-3} \emph{Kernel bounds for the family~\eqref{it:primitives-p1}.} Let $b$ be a positively oriented bond of $T_k$ with initial point $x$, let $b'$ be a positively oriented bond of $T^0_k$, let $\mathfrak{l}\in\{0,1,2\}$, and let $\tilde b$ be a positively oriented bond of $T_k$ of the direction of $b$ whose initial point $\tilde x$ satisfies $0<\lvert x-\tilde x\rvert\le1$. Then each of the three numbers
\[
\left\lvert H_k(s)(b,b')\right\rvert,\qquad \left\lvert\left(\partial_{\mathfrak{l}}H_k(s)\right)(b,b')\right\rvert,\qquad \left\lvert\left(\delta_{\alpha}\partial_{\mathfrak{l}}H_k(s)\right)\left((b,\tilde b);b'\right)\right\rvert
\]
is at most $C\cdot\exp\left(-\gamma\cdot d(b,b')\right)$.

\item\label{it:renormalisation-group-4} \emph{Kernel bound for the family~\eqref{it:primitives-p2}.} For all elements $\Upsilon$, $\Upsilon'$ of the basis~\eqref{eq:renormalisation-group-1} of $V_k$,
\[
\left\lvert C_k^{\frac{1}{2}}(s)(\Upsilon,\Upsilon')\right\rvert\le C\cdot\exp\left(-\gamma\cdot d(\Upsilon,\Upsilon')\right).
\]

\item\label{it:renormalisation-group-5} \emph{Bounds for the family~\eqref{it:primitives-p3}.} For every spinor $f$ on $T^0_k$,
\[
\sup_{T_k}\left\lvert H_k(s,A)\,f\right\rvert\le C\cdot\sup_{T^0_k}\lvert f\rvert,\qquad \sup_P\left\lvert\delta_{\alpha,A}H_k(s,A)\,f\right\rvert\le C\cdot\sup_{T^0_k}\lvert f\rvert,
\]
where $P$ is the set of the pairs $(x,x')\in\mathrm{Pz}_k$ with $d(x,x')<1$, and $\delta_{\alpha,A}$ is the operator of Definition~\ref{def:pairs-sites-holder}.

\item\label{it:renormalisation-group-6} \emph{Kernel bounds for the family~\eqref{it:primitives-p4}.} For all sites $x$, $x'$ of $T^0_k$,
\[
\left\lvert\Gamma_k(s,A)(x,x')\right\rvert\le C\cdot\exp\left(-\gamma\cdot d(x,x')\right),\qquad \left\lvert M_k(X,A)(x,x')\right\rvert\le C\cdot\exp\left(-\gamma\cdot d(x,x')-\kappa'\cdot d_M(X)\right).
\]

\item\label{it:renormalisation-group-7} \emph{Bound for the family~\eqref{it:primitives-p5}.} For all $y\ge0$, all sites $x$, $x'$ of $T^0_k$ and every spinor $f$ on $T_k$,
\[
\sup_{T_k}\left\lvert\mathbf{1}_{\Delta_x}\cdot S_{k,y}(X,A)\cdot\mathbf{1}_{\Delta_{x'}}\,f\right\rvert\le C\cdot\exp\left(-\gamma\cdot d(x,x')-\kappa'\cdot d_M(X)\right)\cdot\sup_{T_k}\lvert f\rvert,
\]
where $\Delta_x=B_k(x)$ is the block of Section~\ref{sec:lattice-geometry}. For $k=0$ the assertion of this item is that the kernel bound displayed in property~\ref{it:renormalisation-group-17}, which carries the factor $(1+y)^{-1}$, holds for all $y\ge0$ and all sites $x$, $x'$ of $T^0_0$. That bound implies the inequality above: for $k=0$ one has $\Delta_x=\{x\}$, so that the left-hand side of the inequality is at most $\left\lvert S_{0,y}(X,A)(x,x')\right\rvert\cdot\sup_{T_0}\lvert f\rvert$, and $(1+y)^{-1}\le1$.

\item\label{it:renormalisation-group-8} \emph{Locality of rows.} Let $\mathrm{Op}$ be one of the families~\eqref{it:primitives-p1}--\eqref{it:primitives-p3} and $\Gamma_k$. Let $\mathscr{Y}$ be a set of parameter cubes of the kind of $\mathrm{Op}$, let $Y$ be the union of the cubes of $\mathscr{Y}$, let $s\in\mathrm{Pd}_k$ satisfy $s_c=0$ for every parameter cube $c$ of this kind that does not belong to $\mathscr{Y}$, and let $Y_1$ be a link-component of $Y$. Let $x$ be a row index of $\mathrm{Op}$ that is contained in $Y_1$. Then $\mathrm{Op}(s,A)(x,x')=0$ unless $x'$ is contained in $Y_1$, and
\[
\mathrm{Op}(s,A)(x,\cdot\,)=\mathrm{Op}(\tilde s,\tilde A)(x,\cdot\,)
\]
for every $\tilde s\in\mathrm{Pd}_k$ that vanishes at the cubes not belonging to $\mathscr{Y}$ and coincides with $s$ at the cubes of $Y_1$, and for every $\tilde A\in\tilde{G}_k$ that coincides with $A$ on the bonds that meet $Y_1$.

\item\label{it:renormalisation-group-9} \emph{Locality of columns.} Let $\mathrm{Op}$, $\mathscr{Y}$, $Y$, $s$ and $Y_1$ be as in property~\ref{it:renormalisation-group-8}, and let $x'$ be a column index of $\mathrm{Op}$ that is contained in $Y_1$. Then $\mathrm{Op}(s,A)(x,x')=0$ unless $x$ is contained in $Y_1$, and the column $\mathrm{Op}(s,A)(\cdot\,,x')$ depends on $s$ only through the values $s_c$ at the cubes $c$ of $Y_1$ and on $A$ only through its values on the bonds that meet $Y_1$; that is, $\mathrm{Op}(s,A)(\cdot\,,x')=\mathrm{Op}(\tilde s,\tilde A)(\cdot\,,x')$ for all $\tilde s$ and $\tilde A$ as in property~\ref{it:renormalisation-group-8}.

\item\label{it:renormalisation-group-10} \emph{Rows outside the link-components.} Let $\mathrm{Op}$ be one of the families~\eqref{it:primitives-p1}--\eqref{it:primitives-p3}, and let $\mathscr{Y}$, $Y$ and $s$ be as in property~\ref{it:renormalisation-group-8}. If a row index $x$ of $\mathrm{Op}$ is contained in no link-component of $Y$, then $\mathrm{Op}(s,A)(x,\cdot\,)=0$.

\item\label{it:renormalisation-group-11} \emph{Rows at fine sites.} Let $\mathscr{Y}$, a set of cubes of $\pi_k$, and $s$ be as in property~\ref{it:renormalisation-group-8} for the family~\eqref{it:primitives-p3}, and let $z$ be a site of $T_k$, with $K_z$ as in Section~\ref{sec:lattice-geometry}. If $\mathscr{Y}$ does not contain all cubes of $\pi_k$ that meet $K_z$, then $H_k(s,A)(z,\cdot\,)=0$.

\item\label{it:renormalisation-group-12} \emph{Locality of the polymer families.} Let $\mathrm{Op}(X,A)$ stand for $M_k(X,A)$ or for $S_{k,y}(X,A)$. Then $\mathrm{Op}(X,A)(x,x')=0$ unless both indices $x$ and $x'$ lie in $X$, and $\mathrm{Op}(X,A)=\mathrm{Op}(X,\tilde A)$ for every $\tilde A\in\tilde{G}_k$ that coincides with $A$ on the bonds that meet $X$.

\item\label{it:renormalisation-group-13} \emph{Gauge covariance.} Let $\lambda$ be a real function on the sites of $T_k$, and let $\lambda^u$ be its restriction to $T^0_k$. With the exponentials acting as multiplication operators,
\begin{align*}
H_k(s,A-d\lambda)&=\exp\left(i e_k\cdot\lambda\right)\cdot H_k(s,A)\cdot\exp\left(-i e_k\cdot\lambda^u\right),\\
\Gamma_k(s,A-d\lambda)&=\exp\left(i e_k\cdot\lambda^u\right)\cdot\Gamma_k(s,A)\cdot\exp\left(-i e_k\cdot\lambda^u\right),\\
M_k(X,A-d\lambda)&=\exp\left(i e_k\cdot\lambda^u\right)\cdot M_k(X,A)\cdot\exp\left(-i e_k\cdot\lambda^u\right),\\
S_{k,y}(X,A-d\lambda)&=\exp\left(i e_k\cdot\lambda\right)\cdot S_{k,y}(X,A)\cdot\exp\left(-i e_k\cdot\lambda\right).
\end{align*}

\item\label{it:renormalisation-group-14} \emph{Covariance under the step group.} For every $g\in\mathbb{G}_k^L$, acting as in Section~\ref{sec:dirac-matrices-lattice}, and for $\mathrm{Op}$ any of the families~\eqref{it:primitives-p1}--\eqref{it:primitives-p3} and $\Gamma_k$,
\[
\mathrm{Op}(gs,gA)\cdot g=g\cdot\mathrm{Op}(s,A),\qquad M_k(gX,gA)\cdot g=g\cdot M_k(X,A);
\]
moreover $S_{k,y}(gX,gA)\cdot g=g\cdot S_{k,y}(X,A)$ for every $g\in\mathbb{G}^L_{k,3}$.

\item\label{it:renormalisation-group-15} \emph{Charge conjugation.} Let the transpose $O^T$ of an operator $O$ on the spinors on $T^0_k$ be taken with respect to the bilinear pairing of Section~\ref{sec:norms-pairings-kernels},
\[
\langle f,g\rangle:=\sum_yf(y)\cdot g(y),
\]
the sum extending over the sites $y$ of $T^0_k$ and over the four components (the product of two spinors is $l\mathbb{Z}^3$-periodic). Then
\[
\Gamma_k(s,A)^T=\mathsf{C}^{-1}\cdot\Gamma_k(s,-A)\cdot\mathsf{C},\qquad M_k(X,A)^T=\mathsf{C}^{-1}\cdot M_k(X,-A)\cdot\mathsf{C}.
\]

\item\label{it:renormalisation-group-16} \emph{Complex conjugation.} For a linear operator $O$ between spaces of spinors let $\overline{O}$ be the operator $f\mapsto\overline{O\,\overline{f}}$, whose kernel is the complex conjugate of the kernel of $O$; moreover $\overline{s}:=\left(\overline{s_c}\right)_c$, and $\overline{A}$ is the complex conjugate of the connection $A$. For the families~\eqref{it:primitives-p3} and~\eqref{it:primitives-p4},
\begin{align*}
\overline{H_k\left(\overline{s},\overline{A}\right)}&=\varkappa^{-1}\cdot H_k(s,-A)\cdot\varkappa,\qquad \overline{\Gamma_k\left(\overline{s},\overline{A}\right)}=\varkappa^{-1}\cdot\Gamma_k(s,-A)\cdot\varkappa,\\
\overline{M_k\left(X,\overline{A}\right)}&=\varkappa^{-1}\cdot M_k(X,-A)\cdot\varkappa;
\end{align*}
for real $s$ the operators $H_k(s)$ and $C_k^{\frac{1}{2}}(s)$ of the families~\eqref{it:primitives-p1} and~\eqref{it:primitives-p2} are real.

\item\label{it:renormalisation-group-17} \emph{Level zero.} The diagonal families of level $0$, which are the families~\eqref{it:primitives-p1} and~\eqref{it:primitives-p3} at $k=0$ (Section~\ref{sec:primitive-operators}), are defined by
\[
H_0(s')(b,b'):=\delta_{bb'}\cdot\prod_cs'_c,\qquad H_0(s'',A)(x,x'):=\delta_{xx'}\cdot\prod_cs''_c,
\]
the first product extending over the cubes $c\in\pi'_0$ that contain $b$, the second over the cubes $c\in\pi_0$ that meet $K_x$. These two families have the properties~\ref{it:renormalisation-group-1}, \ref{it:renormalisation-group-2}, \ref{it:renormalisation-group-8}--\ref{it:renormalisation-group-11}, \ref{it:renormalisation-group-13}, \ref{it:renormalisation-group-14} and~\ref{it:renormalisation-group-16}, and the properties~\ref{it:renormalisation-group-3} and~\ref{it:renormalisation-group-5} with the constant $4\exp(2\gamma)\cdot M^{27\alpha_0}$ in place of $C$. The families~\eqref{it:primitives-p2}, $\Gamma_0$ and $S_{0,y}$ have the properties~\ref{it:renormalisation-group-1}, \ref{it:renormalisation-group-2}, \ref{it:renormalisation-group-4}, \ref{it:renormalisation-group-6}--\ref{it:renormalisation-group-9} and~\ref{it:renormalisation-group-12}--\ref{it:renormalisation-group-16}, and the family~\eqref{it:primitives-p2} has property~\ref{it:renormalisation-group-10}; here property~\ref{it:renormalisation-group-7} is, by its definition for $k=0$, the kernel bound
\[
\left\lvert S_{0,y}(X,A)(x,x')\right\rvert\le C\cdot(1+y)^{-1}\cdot\exp\left(-\gamma\cdot d(x,x')-\kappa'\cdot d_M(X)\right).
\]
A system of level $0$ has no family $M_0$, and the assertions on $M_k$ concern the levels $k\ge1$ only.

\item\label{it:renormalisation-group-18} \emph{Projection.} Let $l_1,l_2\in\mathrm{Cl}$ with $l_2=J\cdot l_1$ for an odd integer $J$, let $M\ge L^2$ and $\gamma_w\cdot M\ge1$, and let $T_1$, $T_2$ and the covering map $\mathrm{pj}\colon T_2\to T_1$ be as in Definition~\ref{def:periodisation-projection}. If a system of primitives of level $k$ on $T_2$ has the properties~\ref{it:renormalisation-group-1}--\ref{it:renormalisation-group-17} for $\mathrm{par}^0$ (Section~\ref{sec:primitive-operators}), then its projection to $T_1$ (Definition~\ref{def:periodisation-projection}) is well defined, its kernels are given, for indices $x$, $x'$ of $\mathbb{L}_k$, by
\begin{equation}\label{eq:renormalisation-group-8}
\mathrm{Op}^{\mathrm{pj}}(\,\cdot\,)(x,x')=\sum_{w\in(\mathbb{Z}/J\mathbb{Z})^3}\boldsymbol{\chi}(w)\cdot\widetilde{\mathrm{Op}}(\,\cdot\,)\left(x,x'+l_1\cdot w\right),
\end{equation}
and it has the properties~\ref{it:renormalisation-group-1}--\ref{it:renormalisation-group-17} for $\mathrm{par}$; at level $0$ the projections of the two families of property~\ref{it:renormalisation-group-17} on $T_2$ are the two families of property~\ref{it:renormalisation-group-17} on $T_1$. Here the notation is that of Definition~\ref{def:periodisation-projection}: $\mathrm{Op}^{(2)}$ is a family of the system on $T_2$ and $\mathrm{Op}^{\mathrm{pj}}$ its projection; $(\,\cdot\,)$ stands for $(s,A)$, resp. for $(X,A)$, with $s\in\mathrm{Pd}_k(T_1)$, $A\in\tilde{G}_k(T_1)$ and $X\in D_k(T_1)$;
\[
\widetilde{\mathrm{Op}}(s,A):=\mathrm{Op}^{(2)}\left(\mathrm{pj}^{*}s,\mathrm{pj}^{*}A\right),\qquad \widetilde{\mathrm{Op}}(X,A):=\sum_{X'\in D_k(T_2)\colon\mathrm{pj}(X')=X}\mathrm{Op}^{(2)}\left(X',\mathrm{pj}^{*}A\right);
\]
and $\boldsymbol{\chi}\colon\mathbb{Z}^3\to\{1,-1\}$, the character of the functions on which the family acts, is given by $\boldsymbol{\chi}(w):=1$ for the families that act on bond functions or on elements of $V_k$, and by $\boldsymbol{\chi}(w):=(-1)^{w_0+w_1+w_2}$ for the families that act on spinors. The kernels in~\eqref{eq:renormalisation-group-8} are those of Section~\ref{sec:two-tori}: for $i=1,2$, the kernel of an operator $O$ on $T_i$ from one space of functions to another, both of character $\boldsymbol{\chi}$, is the function $O(x,x')$ of two indices $x$, $x'$ of $\mathbb{L}_k$ such that
\[
O\left(x+l_i\cdot e_{\nu},x'\right)=O\left(x,x'+l_i\cdot e_{\nu}\right)=\boldsymbol{\chi}(e_{\nu})\cdot O(x,x')\quad(\nu=0,1,2),\qquad(Of)(x)=\sum_{x'}O(x,x')\cdot f(x'),
\]
the sum extending over a set of representatives $x'$ of the column indices of $\mathbb{L}_k$ modulo $l_i\mathbb{Z}^3$; the factors $\eta_k^3$ that weight sums over the fine lattice are part of the kernel. The sum in~\eqref{eq:renormalisation-group-8} extends over a set of representatives $w\in\mathbb{Z}^3$ of the classes of $(\mathbb{Z}/J\mathbb{Z})^3$, and each of its terms is independent of the representative of its class: if $w$ is replaced by $w+J\cdot e_{\nu}$, then $\widetilde{\mathrm{Op}}(\,\cdot\,)\left(x,x'+l_1\cdot w\right)$ acquires the factor $\boldsymbol{\chi}(e_{\nu})$, since $J\cdot l_1=l_2$, and $\boldsymbol{\chi}(w)$ acquires the factor $\boldsymbol{\chi}(J\cdot e_{\nu})=\boldsymbol{\chi}(e_{\nu})$, since $J$ is odd; moreover $\boldsymbol{\chi}(e_{\nu})^2=1$.
\end{enumerate}

For the families~\eqref{it:primitives-p1} and~\eqref{it:primitives-p2} of Theorem~\ref{thm:properties-weakened-photon}, the properties~\ref{it:renormalisation-group-1}--\ref{it:renormalisation-group-4}, \ref{it:renormalisation-group-8}--\ref{it:renormalisation-group-10}, \ref{it:renormalisation-group-14} and~\ref{it:renormalisation-group-16} are contained in that theorem, and property~\ref{it:renormalisation-group-17} for $H_0(s')$ is contained in Proposition~\ref{prop:family-level-zero}. For the families~\eqref{it:primitives-p3}--\eqref{it:primitives-p5} of Theorem~\ref{thm:properties-weakened-families}, the properties~\ref{it:renormalisation-group-1}, \ref{it:renormalisation-group-2} and~\ref{it:renormalisation-group-5}--\ref{it:renormalisation-group-16} are contained in that theorem, and property~\ref{it:renormalisation-group-17} for $H_0(s'',A)$, $\Gamma_0$ and $S_{0,y}$ is contained in Theorem~\ref{thm:properties-level-zero}. These two theorems give the bounds of the properties~\ref{it:renormalisation-group-5}--\ref{it:renormalisation-group-7}, and the kernel bounds of property~\ref{it:renormalisation-group-17} for $\Gamma_0$ and $S_{0,y}$, with the constant $C_1$ of Theorem~\ref{thm:properties-weakened-families} for the Euclidean norm of Section~\ref{sec:norms-pairings-kernels}; for both norms, that is, as properties in the sense fixed at the beginning of this subsection, they give them with the constant $2C_1$, and hence with the constant $C^F\ge2C_1$ of Theorem~\ref{thm:fermion-operators}; this is Theorem~\ref{thm:fermion-operators}(ii), which is Theorem~\ref{thm:fermion-operators-weakened}(ii). For the family $H_0(s'',A)$, whose kernel is diagonal with multiples of the unit matrix as entries, the bound of property~\ref{it:renormalisation-group-5} given by Theorem~\ref{thm:properties-level-zero}(b) holds for both norms with one and the same constant; and the properties~\ref{it:renormalisation-group-1}, \ref{it:renormalisation-group-2} and~\ref{it:renormalisation-group-8}--\ref{it:renormalisation-group-16} do not involve a norm on $\mathbb{C}^4$. That the systems of primitives of Definition~\ref{def:periodisation-projection} on the tori $T_k(R_{\mathrm{top}})$ have the properties~\ref{it:renormalisation-group-1}--\ref{it:renormalisation-group-17} for $\mathrm{par}^0$, and hence for $\mathrm{par}$, is a special case of Lemma~\ref{lem:assembled-systems-primitives}. Property~\ref{it:renormalisation-group-18} is a consequence of Theorem~\ref{thm:properties-projection}.

Let $W$ be a union of $M$-cubes. A parameter family $s$ is said to \emph{vanish off} $W$ if $s_c=0$ at each parameter cube $c$ not contained in $W$. The family $s\cdot\mathbf{1}_W$ is defined by
\[
\left(s\cdot\mathbf{1}_W\right)_c:=s_c\ \text{ if $c$ is contained in $W$},\qquad \left(s\cdot\mathbf{1}_W\right)_c:=0\ \text{ otherwise}.
\]
A \emph{component} of $W$ is a maximal union of $M$-cubes contained in $W$ that is connected through faces. For an index $x$, $\square(x)$ is the union of the $M$-cubes that $x$ meets, and $\square(x)^{\sim c_l}$ is its enlargement by $c_l=2L$ layers of $M$-cubes (Section~\ref{sec:lattice-geometry} and Definition~\ref{def:parameter-cubes-polydisks}).

\begin{proposition}[Rows over a component]\label{prop:rows-component}
Let $l\in\mathrm{Cl}$ and $M>4L$, and let a system of primitives of level $k$ on $T_k$ have the properties~\ref{it:renormalisation-group-8} and~\ref{it:renormalisation-group-12}. Let $\mathrm{Op}$ be one of the families~\eqref{it:primitives-p1}--\eqref{it:primitives-p3} and $\Gamma_k$ of this system, $W$ a union of $M$-cubes, $s\in\mathrm{Pd}_k$ a parameter family of the kind of $\mathrm{Op}$ that vanishes off $W$, $W_1$ a component of $W$, and $x$ a row index of $\mathrm{Op}$ such that $\square(x)^{\sim c_l}$ is contained in $W_1$. Let $A,A_1\in\tilde{G}_k$ coincide on the bonds that meet $W_1$. Then $\mathrm{Op}(s,A)(x,x')=0$ unless $x'$ is contained in $W_1$, and
\[
\mathrm{Op}(s,A)(x,\cdot\,)=\mathrm{Op}\left(s\cdot\mathbf{1}_{W_1},A_1\right)(x,\cdot\,).
\]
\end{proposition}

The assertion is contained in Lemma~\ref{lem:locality-rows}.

\subsubsection{Properties of the exact photon operators}\label{sec:properties-exact-photon}

\begin{proposition}[The background field $H_kA$]\label{prop:background-field}
Let $A$ be a real bond function on $T^0_k$. The variational problem of Section~\ref{sec:averaging-fluctuations-gauge}, that is, the minimisation of $|d\mathscr{A}|_2^2$ over the real bond functions $\mathscr{A}$ on $T_k$ subject to $Q_k\mathscr{A}=A$ and $\mathscr{R}_k\delta\mathscr{A}=0$, has exactly one solution $H_kA$. The map $A\mapsto H_kA$ is linear and real, and $Q_kH_k=1$. Every constant bond function $c$ satisfies $H_kc=c$, and $dA=0$ implies $dH_kA=0$.
\end{proposition}

This is a special case of Proposition~\ref{prop:minimiser}.

\begin{proposition}[Kernel bounds for $H_k$]\label{prop:kernel-bounds}
The constants $C_P\ge1$ and $\gamma_P\in(0,1]$ defined in~\eqref{eq:operators-photon-18} depend only on $L$ and $\alpha$. For all $k\ge0$ and all $l\in\mathrm{Cl}$ property~\ref{it:renormalisation-group-3} holds at $s=1$, that is, for $H_k$ in place of $H_k(s)$, with $C=C_P$ and $\gamma=\gamma_P$. With the kernels of property~\ref{it:renormalisation-group-3} this reads
\[
\left|H_k(b,b')\right|,\quad\left|(\partial_{\mathfrak{l}}H_k)(b,b')\right|,\quad\left|(\delta_{\alpha}\partial_{\mathfrak{l}}H_k)\left((b,\tilde{b});b'\right)\right|\ \le\ C_P\cdot\exp\left(-\gamma_P\cdot d(b,b')\right)
\]
for $\mathfrak{l}=0,1,2$, for all positively oriented fine bonds $b$, $\tilde{b}$ of $T_k$ of the same direction whose initial points $x$, $\tilde{x}$ satisfy $0<|x-\tilde{x}|\le1$, and for all positively oriented unit bonds $b'$.
\end{proposition}

This is a consequence of Theorem~\ref{thm:decay-kernels}.

\begin{proposition}[Spectral bounds and decay of $\tilde{\Delta}_k$]\label{prop:spectral-bounds-decay}
Let $C_P$ and $\gamma_P$ be the constants of Proposition~\ref{prop:kernel-bounds}. There is a constant $c_P$ such that
\[
c_P^{-1}\cdot|\mathsf{Z}|_2^2\ \le\ \left\langle\mathsf{Z},\tilde{\Delta}_k\mathsf{Z}\right\rangle\ \le\ c_P\cdot|\mathsf{Z}|_2^2\qquad\text{for every }\mathsf{Z}\in V_k,
\]
and the matrix elements of $\tilde{\Delta}_k$ in the basis $\{\Upsilon\}$ satisfy
\[
\left|\tilde{\Delta}_k(\Upsilon,\Upsilon')\right|\le C_P\cdot\exp\left(-\gamma_P\cdot d(\Upsilon,\Upsilon')\right).
\]
\end{proposition}

\begin{proof}
Proposition~\ref{prop:bounds-fluctuation-form} gives the two-sided bound with the number $c_P$ defined in~\eqref{eq:operators-photon-17}, together with
\[
\left|\tilde{\Delta}_k(\Upsilon,\Upsilon')\right|\le C_{\mathrm{Dt}}\cdot\exp\left(-\gamma_H\cdot d(\Upsilon,\Upsilon')\right),
\]
where $C_{\mathrm{Dt}}$ is defined in~\eqref{eq:operators-photon-17} and $\gamma_H$ in~\eqref{eq:operators-photon-8}. The definition~\eqref{eq:operators-photon-18} of $C_P$ and $\gamma_P$ gives $C_{\mathrm{Dt}}\le C_P$ and $\gamma_P\le\gamma_H$; as $d(\Upsilon,\Upsilon')\ge0$, the bound on the matrix elements follows.
\end{proof}

\begin{proposition}[Decay of the exact root]\label{prop:decay-exact-root}
With the constants $C_P$ and $\gamma_P$ of Proposition~\ref{prop:kernel-bounds}, the matrix elements of the exact root $C_k^{\frac{1}{2}}$ in the basis $\{\Upsilon\}$ satisfy
\[
\left|C_k^{\frac{1}{2}}(\Upsilon,\Upsilon')\right|\le C_P\cdot\exp\left(-\gamma_P\cdot d(\Upsilon,\Upsilon')\right).
\]
\end{proposition}

\begin{proof}
Proposition~\ref{prop:decay-resolvents-root} gives
\[
\left|C_k^{\frac{1}{2}}(\Upsilon,\Upsilon')\right|\le(2c_P)^{\frac{1}{2}}\cdot\exp\left(-\gamma_P\cdot d(\Upsilon,\Upsilon')\right),
\]
with $c_P$ defined in~\eqref{eq:operators-photon-17} and $\gamma_P$ in~\eqref{eq:operators-photon-18}, and $(2c_P)^{\frac{1}{2}}\le C_P$ by the definition~\eqref{eq:operators-photon-18} of $C_P$.
\end{proof}

\begin{proposition}[The basis, the map $C$ and the box]\label{prop:basis-map-box}
For $\mu=0,1,2$ let $\mathfrak{r}_{\mu}$ be the reflection $x_{\mu}\mapsto-x_{\mu}$. It maps the $L$-block $B(0)$ onto itself, and it acts on the real functions $A$ on the interior bonds of $B(0)$ by $(\mathfrak{r}_{\mu}A)(x,x'):=A(\mathfrak{r}_{\mu}x,\mathfrak{r}_{\mu}x')$. The three operators $\mathfrak{r}_0$, $\mathfrak{r}_1$, $\mathfrak{r}_2$ so defined are commuting orthogonal involutions, and each of them maps onto itself the space $\ker\tau$ of those functions $A$ with $(\tau A)(0,x)=0$ for all $x\in B(0)$. Consequently $\ker\tau$ has an orthonormal basis that consists of joint eigenvectors of $\mathfrak{r}_0$, $\mathfrak{r}_1$, $\mathfrak{r}_2$. One such basis, the basis $\{e_i^0\}$ of Section~\ref{sec:averaging-fluctuations-gauge}, is fixed once for each $L$.

The family $\{\Upsilon\}$ of~\eqref{eq:renormalisation-group-1} is an orthonormal basis of $V_k$ with $d_V\cdot n_k$ elements. Each $g\in\mathbb{G}_k$ is an isometry of $V_k$ onto $V_k$, and it maps the finite set $\{\pm\Upsilon\}$ bijectively onto itself. The set $\mathrm{Box}_k$ and the measures $\mu_I$ and $\hat{\mu}$ are invariant under $\mathbb{G}_k$ and under $\mathsf{Z}\mapsto-\mathsf{Z}$.

For every $\Upsilon$ the set $\operatorname{supp}\Upsilon$ is contained in one $L$-block or consists of the two end points of one bond. Let $b$ be a bond of $T^0_k$. For $\mathsf{Z}\in V_k$ the number $(C\mathsf{Z})(b)$ depends only on the coordinates $\mathsf{Z}(\Upsilon)$ of the basis elements $\Upsilon$ that are contained in one of the (at most two) $L$-blocks containing an end point of $b$ and, if there are two such blocks, on the coordinates $\mathsf{Z}(\delta_{b'})$ with $b'$ a positively oriented non-central face bond of their common face. There is a constant $c_C=c_C(L)$, depending only on $L$, such that
\[
|C\mathsf{Z}|_{\infty}\le c_C\cdot\sup_{\Upsilon}\left|\mathsf{Z}(\Upsilon)\right|\qquad\text{for every }\mathsf{Z}\in V_k.
\]

Finally,
\[
d\varepsilon^0_k=-d_V\cdot\log\left((2\pi)^{-\frac{1}{2}}\int_{-p_{0,k}}^{p_{0,k}}\mathrm{d}t\,\exp\left(-\frac{t^2}{2}\right)\right),
\]
the argument of the logarithm being the probability that a standard normal random variable has modulus at most $p_{0,k}$. If $p_{0,k}\ge1$, which is the case for every $e_k\in(0,e_0]$ by the conditions~\eqref{eq:renormalisation-group-A9}, then
\[
0\le d\varepsilon^0_k\le12\exp\left(-\frac{p_{0,k}^2}{2}\right).
\]
\end{proposition}

The assertion is contained in Proposition~\ref{prop:fluctuation-space}.

\begin{proposition}[Symmetries and reality]\label{prop:symmetries-reality}
For every $g\in\mathbb{G}_k$
\[
gH_k=H_kg,\qquad gC=Cg,\qquad g\tilde{\Delta}_k=\tilde{\Delta}_kg,\qquad gC_k^{\frac{1}{2}}=C_k^{\frac{1}{2}}g,\qquad g\mathrm{Sq}_k=\mathrm{Sq}_kg.
\]
Let $g=(S,a)\in\mathbb{G}_k^L$, and let $g'=(S,a/L)$ be the element of the group $\mathbb{G}_{k+1}$ of side $l/L$ that Proposition~\ref{prop:symmetries-extraction}(d) assigns to $g$. Then
\[
A^{\min}_k(g'A')=g\,A^{\min}_k(A')
\]
for every real bond function $A'$ on the unit lattice $T^0_{k+1}$ of side $l/L$. The operators $H_k$, $C$, $\tilde{\Delta}_k$, $C_k^{\frac{1}{2}}$, $\mathrm{Sq}_k$ and the map $A^{\min}_k$ are real. The set $\mathrm{Box}_k$ and the measures $\mu_I$ and $\hat{\mu}$ are invariant under $\mathbb{G}_k$ and under $\mathsf{Z}\mapsto-\mathsf{Z}$.
\end{proposition}

These assertions are contained in Proposition~\ref{prop:lattice-symmetries}.

\begin{proposition}[Local regularity of $H_kA$]\label{prop:local-regularity}
Put $s_d:=(2L+1)M$, the side of the cubes $\square^{\dagger}$ of Definition~\ref{def:enlarged-cubes}; then $s_d<l$, because $l\in\mathrm{Cl}$. Let $A$ be a real bond function on $T^0_k$, with arbitrary harmonic part, and let $\square$ be an $M$-cube. There is a real function $\nu$ on the sites of $\square^{\dagger}$ such that each of the three norms of Section~\ref{sec:domains-bond-functions} of $H_kA-d\nu$ on $\square^{\dagger}$ is at most $C_B\cdot|dA|_{\infty}$; here $C_B=C_B(C_P,\gamma_P,s_d)$ is the constant defined in~\eqref{eq:operators-photon-19}.
\end{proposition}

This is proved as Proposition~\ref{prop:regularity-enlarged-cubes}.

\begin{proposition}[Fields in the domains]\label{prop:fields-domains}
Assume that
\[
8C_B\cdot p_k\le e_k^{-\frac{3}{4}+3\varepsilon},\qquad 8e_k^{\frac{3}{4}-\frac{7\varepsilon}{2}}\le1,
\]
with the constant $C_B$ of Proposition~\ref{prop:local-regularity}, and that the conditions~\eqref{eq:renormalisation-group-A1} and~\eqref{eq:renormalisation-group-A2} of Section~\ref{sec:conditions} hold. Then the following assertions hold.
\begin{enumerate}
\item[(i)] The set of the real bond functions $A$ on $T^0_k$ with $|dA|_{\infty}\le p_k$ is contained in the domain $\mathrm{Dom}_k$ of Definition~\ref{def:domains}; it contains $0$ and the constant bond functions.
\item[(ii)] Let moreover $l/L\in\mathrm{Cl}$, and let $\mathrm{Dom}_{k+1}$ be formed with the side $l/L$ and with the charge $e_{k+1}=L^{\frac{1}{2}}\cdot e_k$. Then $A^{\min}_k(A')+\mathrm{Sq}_k\zeta\in\mathrm{Dom}_k$ for all $A'\in\mathrm{Dom}_{k+1}$ and all $\zeta\in\mathrm{Box}_k$.
\item[(iii)] For every $\mathsf{Z}\in\mathrm{Box}_k$ each of the three norms of $\mathscr{Z}:=H_kC\,C_k^{\frac{1}{2}}\mathsf{Z}$ is at most $\frac{1}{2}e_k^{-\frac{\varepsilon}{2}}$.
\end{enumerate}
\end{proposition}

All three assertions are special cases of Proposition~\ref{prop:backgrounds-domains}.

\begin{proposition}[Periodisation]\label{prop:periodisation}
Let $l_1,l_2\in\mathrm{Cl}$ with $l_2=J\cdot l_1$ for an odd integer $J$, as in property~\ref{it:renormalisation-group-18}, and let $\mathrm{pj}^{*}$ be the inclusion of Definition~\ref{def:periodisation-projection}. For each of the operators
\[
O=d,\ \delta,\ Q,\ Q_k,\ C,\ H_k,\ \tilde{\Delta}_k,\ C_k^{\frac{1}{2}},\ A^{\min}_k
\]
let $O^{(i)}$ denote the operator $O$ formed on the torus of side $l_i$, $i=1,2$. Then
\[
O^{(2)}\mathrm{pj}^{*}=\mathrm{pj}^{*}O^{(1)};
\]
for $O=A^{\min}_k$ the inclusion $\mathrm{pj}^{*}$ on the left side is that of the unit lattices of level $k+1$, of sides $l_1/L$ and $l_2/L$.
\end{proposition}

The assertion is contained in Proposition~\ref{prop:periodisation-photon-operators}.

\begin{proposition}[Gauge fixing controlled by the curvature]\label{prop:gauge-fixing-controlled}
Let $l\ge2$ be an integer, and let $A$ be a real one-cochain on $(\mathbb{Z}/l\mathbb{Z})^3$ with mean zero. There is a function $\lambda$ on the sites of $(\mathbb{Z}/l\mathbb{Z})^3$ such that
\[
|A-d\lambda|_{\infty}\le2l^2\cdot|dA|_{\infty}.
\]
\end{proposition}

The same statement reappears as Corollary~\ref{cor:bounded-gauge-torus}, with its proof.

\begin{proposition}[Axial gauge]\label{prop:axial-gauge}
Let $l$ be a multiple of $L$, and let $A$ be a real bond function on $T^0_k$ with $\tau A=0$. Then every bond $b$ satisfies
\[
|A(b)|\le\frac{3}{2}L^2\cdot|dA|_{\infty}+L\cdot|QA|_{\infty},
\]
where $|dA|_{\infty}$ denotes here the supremum of $|dA|$ over those plaquettes that lie in the union of the $L$-blocks met by $b$, of which there are one or two, and $|QA|_{\infty}$ the supremum of $|QA|$ over the face of $b$ (for an interior bond $b$ the second term is absent). In particular,
\[
|Z|_{\infty}\le\frac{3}{2}L^2\cdot|dZ|_{\infty}
\]
for every real bond function $Z$ on $T^0_k$ with $QZ=0$ and $\tau Z=0$.
\end{proposition}

Both bounds are contained in Corollary~\ref{cor:local-bounds-axial}.

\subsubsection{Fermion operators: bounds, symmetries and the determinant}\label{sec:fermion-operators-bounds}

The setting is that of Section~\ref{sec:properties-scheme}. In the statements of this subsection $C_F\ge1$, $\gamma_F\in(0,1]$ and $\mathfrak{m}_0\in(0,1]$ are constants, of which $\gamma_F$ and $\mathfrak{m}_0$ are functions of $(b,\alpha,L)$ alone; the mass $\mathfrak{m}$ of the setting ranges over $(0,\mathfrak{m}_0]$. The letter $A$ always stands for an element of $\tilde{G}_k$.

The letters $y$ and $y'$ carry two meanings. As subscripts of $\Delta$ and of $\mathbf{1}$, as arguments of the distance $d$ and as arguments of a kernel they are unit sites, that is, sites of $T^0_k$. As indices of $S$, $B$ and $\Gamma$, as in $S_{k,y}$, $B_{k,y}$, $\Gamma_{k,y}$ (Section~\ref{sec:exact-fermion-operators}), and in the factor $\left(1+|y|\right)^{-1}$, they stand for the resolvent parameter: $y$ is a real number, and $y'$ is one of the two numbers $y$, $-y$.

If $F$ is a function on $\mathrm{Pz}_k$ and $y$ is a unit site, then $\mathbf{1}_{\Delta_y}F$ denotes the function $(x,x')\mapsto\mathbf{1}_{\Delta_y}(x)\cdot F(x,x')$; a pair is thus attached to the cube $\Delta_y$ that contains its first site.

For the norm $|\cdot|$ on $\mathbb{C}^4$ either the Euclidean norm or the maximum of the moduli of the four components may be taken; a complex $4\times4$ matrix then carries the operator norm belonging to this choice. For a spinor $f$ on $T_k$ or on $T^0_k$, and for a function $f$ on $\mathrm{Pz}_k$ with values in $\mathbb{C}^4$, $|f|_{\infty}$ is the supremum of $|f|$ over the sites, respectively over the pairs. The constant $C^F$ of Theorem~\ref{thm:fermion-operators-weakened} and the constant $C_T$ of Proposition~\ref{prop:left-inverse} are valid for both choices of the norm (Section~\ref{sec:norms-pairings-kernels}).

\begin{proposition}[Decay of the fine propagator]\label{prop:decay-fine-propagator}
Let $k\ge1$. The operator $D_A+\mathfrak{m}+b_k\cdot P_k(A)$ on the spinors on $T_k$ has an inverse $S_k(A)$, which depends analytically on $A\in\tilde{G}_k$. For all unit sites $y$, $y'$ and every spinor $f$ on $T_k$ one has
\begin{align*}
\left|\mathbf{1}_{\Delta_y}\cdot S_k(A)\cdot\mathbf{1}_{\Delta_{y'}}\,f\right|_{\infty}&\le C_F\cdot\exp\left(-\gamma_F\cdot d(y,y')\right)\cdot|f|_{\infty},\\
\left|\mathbf{1}_{\Delta_y}\cdot\delta_{\alpha,A}S_k(A)\cdot\mathbf{1}_{\Delta_{y'}}\,f\right|_{\infty}&\le C_F\cdot\exp\left(-\gamma_F\cdot d(y,y')\right)\cdot|f|_{\infty},
\end{align*}
where in the second inequality the supremum on the left is restricted to the pairs $(x,x')\in\mathrm{Pz}_k$ with $|x-x'|<1$.
\end{proposition}

This is a special case of Theorem~\ref{thm:expansion-walks}.

\begin{proposition}[Decay of the resolvents $S_{k,y}$]\label{prop:decay-resolvents}
Let $k\ge1$ and $y\ge0$. The operator $D_A+\mathfrak{m}+\boldsymbol{\alpha}_y\cdot P_k(A)+\boldsymbol{\beta}_y\cdot P_{k+1}(A)$ has an inverse $S_{k,y}(A)$, which depends analytically on $A\in\tilde{G}_k$, and the two inequalities of Proposition~\ref{prop:decay-fine-propagator} remain true when $S_k(A)$ is replaced by $S_{k,y}(A)$, with one and the same pair of constants $C_F$, $\gamma_F$ for all $y\ge0$. Moreover $S^0_{k+1}(A)=S_{k,0}(A)$.
\end{proposition}

This, too, is a special case of Theorem~\ref{thm:expansion-walks}.

\begin{proposition}[Kernel bounds for $M_k$, $\Gamma_k$ and $B_{k,y}$]\label{prop:kernel-bounds-2}
Let $k\ge1$. For all unit sites $x$, $x'$,
\[
\left|M_k(A)(x,x')\right|\le C_F\cdot\exp\left(-\gamma_F\cdot d(x,x')\right),\qquad\left|\Gamma_k(A)(x,x')\right|\le C_F\cdot\exp\left(-\gamma_F\cdot d(x,x')\right);
\]
for every real $y$ the kernel $B_{k,y}(A)(x,x')$ is zero unless some $L$-block contains both $x$ and $x'$, and
\[
\left|B_{k,y}(A)(x,x')\right|\le C_F\cdot\left(1+|y|\right)^{-1}.
\]
Further, $\delta Z_k(A)\ne0$ and, with $B_k(A)=B_{k,0}(A)$,
\[
\Gamma_k(A)=B_k(A)+b_k^2\cdot B_k(A)\cdot Q_k(A)\cdot S^0_{k+1}(A)\cdot Q_k^T(-A)\cdot B_k(A).
\]
At the level $0$ one has, for all sites $x$, $x'$ of $T_0$ and every real $y$,
\[
\left|\Gamma_0(A)(x,x')\right|\le C_F\cdot\exp\left(-\gamma_F\cdot d(x,x')\right),\qquad\left|S_{0,y}(A)(x,x')\right|\le C_F\cdot\left(1+|y|\right)^{-1}\cdot\exp\left(-\gamma_F\cdot d(x,x')\right).
\]
\end{proposition}

For $k\ge1$ the assertions follow from the results of Section~\ref{sec:operators-unit-lattice}; the bounds at the level $0$ are a special case of Theorem~\ref{thm:expansion-walks}.

\begin{proposition}[Bounds for $H_k(A)$ and $\bar{H}_k(A)$]\label{prop:bounds}
For every spinor $u$ on $T^0_k$
\[
\left|H_k^{\#}(A)\,u\right|_{\infty}\le C_F\cdot|u|_{\infty};
\]
here $H_k^{\#}(A)\,u$ is the pair of the spinor $H_k(A)\,u$ on $T_k$ and the function $\delta_{\alpha,A}H_k(A)\,u$ on $\mathrm{Pz}_k$ (Section~\ref{sec:fermion-integrals-blocking}), and the supremum on the left is taken over all sites of $T_k$ for the first entry and over all pairs of $\mathrm{Pz}_k$, those with $|x-x'|\ge1$ included, for the second. For a unit site $y'$ and $v\in\mathbb{C}^4$ let $\mathbf{1}_{y'}v$ be the spinor on $T^0_k$ that takes the value $v$ at $y'$ and vanishes at every other unit site. For all unit sites $y$, $y'$ and all $v\in\mathbb{C}^4$
\begin{align*}
\left|\mathbf{1}_{\Delta_y}\cdot H_k(A)\,\mathbf{1}_{y'}v\right|_{\infty}&\le C_F\cdot\exp\left(-\gamma_F\cdot d(y,y')\right)\cdot|v|,\\
\left|\mathbf{1}_{\Delta_y}\cdot\delta_{\alpha,A}H_k(A)\,\mathbf{1}_{y'}v\right|_{\infty}&\le C_F\cdot\exp\left(-\gamma_F\cdot d(y,y')\right)\cdot|v|,
\end{align*}
where in the second inequality only the pairs with $|x-x'|<1$ enter the supremum. Let $(x,x')\in\mathrm{Pz}_k$ with $1\le|x-x'|\le M$, and let $z$, $z'$ be the unit sites with $x\in\Delta_z$ and $x'\in\Delta_{z'}$. Then for every unit site $y'$ and every $v\in\mathbb{C}^4$
\[
\left|\left(\delta_{\alpha,A}H_k(A)\,\mathbf{1}_{y'}v\right)(x,x')\right|\le C_F\cdot\exp\left(-\gamma_F\cdot\min\left\{d(z,y'),d(z',y')\right\}\right)\cdot|v|.
\]
Each of these bounds stays valid when $H_k(A)$ is replaced by $\bar{H}_k(A)$ and $\delta_{\alpha,A}$ by $\delta_{\alpha,-A}$.
\end{proposition}

The assertion is contained in Proposition~\ref{prop:operators}.

\begin{proposition}[The left inverse $T_k(A)$]\label{prop:left-inverse}
One has $T_k(A)\cdot H_k(A)=1$ and $\bar{T}_k(A)\cdot\bar{H}_k(A)=1$.

Now let $k\ge1$, $A\in\tilde{R}_k+\mathfrak{X}_k$ and $\mathbf{Y}\in D_k$, and let $T^{(1)}(A)$ and $T^{(2)}(A)$ be the operators of the display~\eqref{eq:renormalisation-group-58} in Section~\ref{sec:localisation-exponent}, formed with the polymer $\mathbf{Y}$. There is a constant $C_T$, depending on $L$ and $b_{-}$ only, such that, for the norm $|\cdot|_{1,\infty}$ of Proposition~\ref{prop:weighted-coefficient-norms},
\[
\left|T^{(1)}(A)\right|_{1,\infty}\le C_T,\qquad\left|T^{(2)}(A)\right|_{1,\infty}\le C_T.
\]
The row of $T^{(1)}(A)$ and the row of $T^{(2)}(A)$ at a unit site $y$ read $A$ on the bonds that have an end point in $B_k(y)$, in the sense of Definition~\ref{def:terminology}(1). If $\psi$ is a spinor on $T_k$ and $\chi:=\delta_{\alpha,A}\psi$, then $T^{(1)}(A)\,\psi+T^{(2)}(A)\,\chi$ and $T_k(A)\,\psi$ take the same value at every unit site of $\mathbf{Y}$.
\end{proposition}

This is a consequence of Proposition~\ref{prop:identities-bounds}.

\begin{proposition}[Symmetries of the fermion operators]\label{prop:symmetries-fermion-operators}
Let $k\ge1$ and let $y$ be real.
\begin{enumerate}
\item[(a)] \emph{Gauge covariance.} Let $\boldsymbol{\omega}$ be a real connection on $T_k$ which is closed, $d\boldsymbol{\omega}=0$, and let $u$ be a function of modulus one on the sites of $T_k$ such that
\[
u(x')=u(x)\cdot\exp\left\{i\,e_k\cdot\eta_k\cdot\boldsymbol{\omega}(x,x')\right\}
\]
for every bond $(x,x')$ of $T_k$. Denote by $u^0$ and $u^1$ the restrictions of $u$ to $T^0_k$ and to $L\mathbb{Z}^3/l\mathbb{Z}^3$, and by the same letters $u$, $u^0$, $u^1$ the operators of multiplication by these functions. The hypotheses on $(\boldsymbol{\omega},u)$ are satisfied, for instance, by $\boldsymbol{\omega}=d\lambda$ and $u=\exp\left(i\,e_k\cdot\lambda\right)$, where $\lambda$ is a real function on the sites of $T_k$, and by a real constant connection $\boldsymbol{\omega}$, that is, $\boldsymbol{\omega}(x,x+\eta_ke_{\mu})=\boldsymbol{\omega}_{\mu}$ for all sites $x$ and $\mu=0,1,2$, with $e_k\cdot l\cdot\boldsymbol{\omega}_{\mu}\in2\pi\mathbb{Z}$ for $\mu=0,1,2$, together with
\[
u(x)=\exp\left(i\,e_k\sum_{\mu=0}^2\boldsymbol{\omega}_{\mu}\cdot x_{\mu}\right)\qquad\left(x\in\mathbb{L}_k\right),
\]
which is a function on the sites of $T_k$ because $e_k\cdot l\cdot\boldsymbol{\omega}_{\mu}\in2\pi\mathbb{Z}$. Under these hypotheses $A-\boldsymbol{\omega}\in\tilde{G}_k$, and
\begin{align*}
D_{A-\boldsymbol{\omega}}&=u\cdot D_A\cdot u^{-1},\\
Q_k(A-\boldsymbol{\omega})&=u^0\cdot Q_k(A)\cdot u^{-1},\qquad Q_k^T\left(-(A-\boldsymbol{\omega})\right)=u\cdot Q_k^T(-A)\cdot\left(u^0\right)^{-1},\\
Q(A-\boldsymbol{\omega})&=u^1\cdot Q(A)\cdot\left(u^0\right)^{-1},\qquad Q^T\left(-(A-\boldsymbol{\omega})\right)=u^0\cdot Q^T(-A)\cdot\left(u^1\right)^{-1},\\
O(A-\boldsymbol{\omega})&=u\cdot O(A)\cdot u^{-1}\qquad\text{for }O=P_k,\ P_{k+1},\ S_k,\ S_{k,y},\\
H_k(A-\boldsymbol{\omega})&=u\cdot H_k(A)\cdot\left(u^0\right)^{-1},\qquad T_k(A-\boldsymbol{\omega})=u^0\cdot T_k(A)\cdot u^{-1},\\
O(A-\boldsymbol{\omega})&=u^0\cdot O(A)\cdot\left(u^0\right)^{-1}\qquad\text{for }O=P,\ D_k,\ M_k,\ \dot{M}_k,\ B_{k,y},\ \Gamma_{k,y},\\
\delta Z_k(A-\boldsymbol{\omega})&=\delta Z_k(A).
\end{align*}
Moreover, for every spinor $f$ on $T_k$ and every $(x,x')\in\mathrm{Pz}_k$,
\[
\left(\delta_{\alpha,A-\boldsymbol{\omega}}(u\cdot f)\right)(x,x')=u(x)\cdot\left(\delta_{\alpha,A}f\right)(x,x').
\]
\item[(b)] \emph{Lattice symmetries.} Let $g=(S,a)\in\mathbb{G}_k$, acting as in Section~\ref{sec:dirac-matrices-lattice}. Then $O(gA)\cdot g=g\cdot O(A)$ for each of the operators
\begin{gather*}
O(A)=D_A,\ Q_k(A),\ Q_k^T(-A),\ Q(A),\ Q^T(-A),\ P_k(A),\ P(A),\ P_{k+1}(A),\ S_k(A),\\
S^0_{k+1}(A),\ H_k(A),\ T_k(A),\ D_k(A),\ M_k(A),\ \dot{M}_k(A),\ B_k(A),\ \Gamma_k(A).
\end{gather*}
Let $y':=y$ if $S\in\mathrm{Sp}_3$ and $y':=-y$ if $S\notin\mathrm{Sp}_3$. Then $O_{y'}(gA)\cdot g=g\cdot O_y(A)$ for $O_y=S_{k,y}$, $B_{k,y}$, $\Gamma_{k,y}$, and $\delta Z_k(gA)=\delta Z_k(A)$.
\item[(c)] \emph{Transposition.} Transposes being formed with respect to the bilinear pairings of Section~\ref{sec:norms-pairings-kernels}, one has $O(A)^T=\mathsf{C}^{-1}\cdot O(-A)\cdot\mathsf{C}$ for $O(A)=D_A$, $P_k(A)$, $P_{k+1}(A)$, $S_k(A)$, $S^0_{k+1}(A)$, $P(A)$, $D_k(A)$, $M_k(A)$, $\dot{M}_k(A)$, $B_k(A)$, $\Gamma_k(A)$, and $O_y(A)^T=\mathsf{C}^{-1}\cdot O_{-y}(-A)\cdot\mathsf{C}$ for $O_y=S_{k,y}$, $B_{k,y}$, $\Gamma_{k,y}$. Further $Q_k(A)^T=Q_k^T(A)$, $Q(A)^T=Q^T(A)$, and
\[
\bar{H}_k(A)=b_k\cdot S_k(A)^T\cdot Q_k^T(A),\qquad\bar{T}_k(A)=b_k^{-1}\cdot Q_k(-A)\cdot\left(D_A+\mathfrak{m}+b_k\cdot P_k(A)\right)^T.
\]
\item[(d)] \emph{Complex conjugation.} For an operator $O$ on spinors let $\overline{O}$ be the operator $f\mapsto\overline{O\,\overline{f}}$. Then $\overline{O(\overline{A})}=\varkappa^{-1}\cdot O(-A)\cdot\varkappa$ for each operator $O(A)$ of the list in (b), and $\overline{O_y(\overline{A})}=\varkappa^{-1}\cdot O_y(-A)\cdot\varkappa$ for $O_y=S_{k,y}$, $B_{k,y}$, $\Gamma_{k,y}$, with the same $y$ on the two sides.
\item[(e)] \emph{The determinant.} $\delta Z_k(-A)=\delta Z_k(A)$ and $\overline{\delta Z_k(\overline{A})}=\delta Z_k(A)$.
\end{enumerate}
At the level $0$ the assertions (a)--(e) hold for the operators $D_A$, $Q(A)$, $P(A)$, $S_{0,y}(A)$, $\Gamma_0(A)$ and for the number $\delta Z_0(A)$.
\end{proposition}

For the operators $D_A$, $Q_k(A)$, $Q_k^T(-A)$, $Q(A)$, $Q^T(-A)$, $P_k(A)$, $P(A)$, $P_{k+1}(A)$ and for $\delta_{\alpha,A}$ the assertions are contained in Lemma~\ref{lem:symmetries-elementary-operators}; the remaining assertions, those at the level $0$ included, are contained in Proposition~\ref{prop:symmetries-fermion-operators-2}.

\begin{proposition}[Local expansion of the determinant $\delta Z_k$]\label{prop:local-expansion-determinant}
Let a system of primitives of level $k$ on $T_k$ (Section~\ref{sec:primitive-operators}) have the properties~\ref{it:renormalisation-group-1}, \ref{it:renormalisation-group-2}, \ref{it:renormalisation-group-7} and~\ref{it:renormalisation-group-12}--\ref{it:renormalisation-group-14} for the parameter set $\mathrm{par}$, where for $k=0$ property~\ref{it:renormalisation-group-7} is taken in the form stated in property~\ref{it:renormalisation-group-17}, that is, as the bound
\[
\left|S_{0,y}(X,A)(x,x')\right|\le C_w\cdot\left(1+y\right)^{-1}\cdot\exp\left(-\gamma_w\cdot d(x,x')-\kappa_w\cdot d_M(X)\right)
\]
for all $y\ge0$, $X\in D_0$, $A\in\tilde{G}_0$ and all sites $x$, $x'$ of $T_0$; and let $\acute{E}_k$, $E^d_k$, $E^{\mathrm{det}}_k$ and $\delta E^{\mathrm{det}}_k$ be formed with its polymer propagators $S_{k,y}(X,A)$ as in Definitions~\ref{def:contour-functional-functions} and~\ref{def:determinant-functions}. There is a constant $C_d$, a function of $(b,c_X,C_w)$ given by~\eqref{eq:operators-fermion-10} with the number $C_w$ of $\mathrm{par}$, such that the following assertions hold.
\begin{enumerate}
\item[(d0)]\itemmark{d0}\label{it:local-expansion-determinant-d0} Put $c^d_0:=0$ and
\[
c^d_k:=4\left[\left(1-L^{-3}\right)\cdot\log b_k+L^{-3}\cdot\log\left(b_k+b\,L^{-1}\right)\right]\qquad(k\ge1).
\]
Then, $n_k=l^3$ being the number of the unit sites,
\[
\delta Z_k(A)=\exp\left\{c^d_k\cdot n_k+\sum_{X\in D_k}\acute{E}_k(X,A)\right\},\qquad\frac{\delta Z_k(A)}{\delta Z_k(0)}=\exp\left\{\sum_{X\in D_k}E^{\mathrm{det}}_k(X,A)\right\}.
\]
\item[(d1)]\itemmark{d1}\label{it:local-expansion-determinant-d1} For every polymer $X$ the maps $\acute{E}_k(X,\cdot)$ and $E^d_k(X,\cdot)$ are analytic on $\tilde{G}_k$, and on $\tilde{G}_k$ their moduli do not exceed $C_d\cdot M^3\cdot\exp\left(-(\kappa+1)\cdot d_M(X)\right)$.
\item[(d2)]\itemmark{d2}\label{it:local-expansion-determinant-d2} For every polymer $X$ the numbers $\acute{E}_k(X,A)$ and $E^d_k(X,A)$ depend on $A$ through its values on the bonds that meet $X$ only (Definition~\ref{def:terminology}(1)).
\item[(d3)]\itemmark{d3}\label{it:local-expansion-determinant-d3} For every polymer $X$, every real function $\lambda$ on the sites of $T_k$ and every $g\in\mathbb{G}_k^L$, each of the four numbers
\[
E^d_k(X,A-d\lambda),\qquad E^d_k(X,-A),\qquad\overline{E^d_k(X,\overline{A})},\qquad E^d_k(gX,gA)
\]
is equal to $E^d_k(X,A)$.
\item[(d4)]\itemmark{d4}\label{it:local-expansion-determinant-d4} Suppose that
\[
l\ge M\cdot r_k,\qquad v_k\ge9,\qquad 8M\cdot\left(\frac{v_k}{4}+4\right)\cdot e_k^{\varepsilon}\le1;
\]
by condition~\eqref{eq:renormalisation-group-A10}, the second and the third of these inequalities are satisfied whenever $e_k\le e_0$. Then for every polymer $X$ and every pair $(\mathscr{A},\mathscr{C})\in\tilde{R}_k\times\mathfrak{X}_k$
\[
\left|\delta E^{\mathrm{det}}_k(X;\mathscr{A},\mathscr{C})\right|\le8\,C_d\cdot M^3\cdot e_k^{\frac{1}{4}}\cdot\exp\left(-\kappa\cdot d_M(X)\right);
\]
for $\mathscr{C}=0$ the left side is $\left|E^{\mathrm{det}}_k(X,\mathscr{A})\right|$.
\item[(d5)]\itemmark{d5}\label{it:local-expansion-determinant-d5} Suppose, besides the conditions of (d4), that $l/L\in\mathrm{Cl}$ and $e_{k+1}=L^{\frac{1}{2}}\cdot e_k\le e_0$. Then $L\,E^{\mathrm{det}}_k\in K_{k+1}$ and $L\,\delta E^{\mathrm{det}}_k\in K^x_{k+1}$, where $L$ is the operation of Section~\ref{sec:local-parts-adjustment} and $K_{k+1}$, $K^x_{k+1}$ are the spaces of Definition~\ref{def:renormalisation-group-1} of the level $k+1$, formed with the side $l/L$ and the charge $e_{k+1}$.
\item[(d6)]\itemmark{d6}\label{it:local-expansion-determinant-d6} Under the conditions of (d5),
\[
\left\|L\,E^{\mathrm{det}}_k\right\|_{k+1}\le e_{k+1}^{\frac{1}{4}-2\varepsilon},\qquad\left\|L\,\delta E^{\mathrm{det}}_k\right\|^x_{k+1}\le e_{k+1}^{\frac{1}{4}-2\varepsilon}.
\]
\end{enumerate}
\end{proposition}

This is a special case of Proposition~\ref{prop:polymer-expansion-determinant}. The hypotheses of Theorem~\ref{thm:expansion-walks}, which that proposition assumes, are satisfied in the setting of Section~\ref{sec:properties-scheme}, by stages~\ref{it:order-choice-4}, \ref{it:order-choice-6} and~\ref{it:order-choice-9}, condition~\eqref{eq:renormalisation-group-19} and $\mathfrak{m}_0\le1$ (Section~\ref{sec:exact-fermion-operators}). Condition~\eqref{eq:renormalisation-group-23}, which clause~\eqref{it:polymer-expansion-determinant-e} of that proposition assumes, is part of stage~\ref{it:order-choice-9}. The inequality $L\cdot\left(\kappa-\kappa_0-1\right)\ge\kappa$, which clause~\eqref{it:polymer-expansion-determinant-g} assumes, follows from stage~\ref{it:order-choice-4}: as $\kappa_0\ge0$, one has $L\cdot\left(\kappa-\kappa_0-1\right)\ge L\cdot\left(\kappa-6\kappa_0-4\right)\ge\kappa+1$. Clause~\eqref{it:polymer-expansion-determinant-g} assumes, besides, condition~\eqref{eq:renormalisation-group-24}, in which $C_d$ is the number~\eqref{eq:operators-fermion-10} formed with the constant $C_w$ of the parameter set at hand. For $\mathrm{par}$ this number is the constant $C_d$ of the proposition, and with this constant condition~\eqref{eq:renormalisation-group-24} is part of stage~\ref{it:order-choice-9} and of condition~\eqref{eq:renormalisation-group-A10}; so it holds in the setting of Section~\ref{sec:properties-scheme}.

\begin{proposition}[Coverings of tori]\label{prop:coverings-tori}
Let $l_2=J\cdot l_1$ be as in property~\ref{it:renormalisation-group-18}, let $T_1$, $T_2$ and the inclusion $\mathrm{pj}^{*}$ be as in Definition~\ref{def:periodisation-projection}, and for an operator $O$ of Section~\ref{sec:exact-fermion-operators} let $O^{(i)}$ denote the operator $O$ formed on $T_i$, $i=1,2$. For every $A\in\tilde{G}_k(T_1)$ and every such $O$,
\[
O^{(2)}\left(\mathrm{pj}^{*}A\right)\cdot\mathrm{pj}^{*}=\mathrm{pj}^{*}\cdot O^{(1)}(A),
\]
and the same identity holds with the H\"older derivative $\delta_{\alpha,A}$ in place of $O(A)$.
\end{proposition}

These identities are contained in Proposition~\ref{prop:operators-covering-torus}.

\subsubsection{Identities for the exact operators}\label{sec:identities-exact-operators}

In this section all operators are exact (Definition~\ref{def:terminology}(3)). The block fields (Definition~\ref{def:terminology}(4)) are real, and their harmonic parts are arbitrary. The level $k\ge0$ is arbitrary; the level $k=0$ of the first step is included. Objects of the level $k+1$ are formed with the side $l/L$, the mass $L\cdot\mathfrak{m}$ and the charge $e_{k+1}=L^{\frac{1}{2}}\cdot e_k$, where $l$, $\mathfrak{m}$ and $e_k$ denote the side, the mass and the charge of the level $k$. Each of the following propositions holds under the hypotheses that Theorem~\ref{thm:renormalisation-group-2} assigns to it.

\begin{proposition}[Block averages]\label{prop:block-averages}
For every bond function $A$ on $T^0_k$ the bond function $QA$ on $L\mathbb{Z}^3/l\mathbb{Z}^3$ has the same mean as $A$. One has $Q_kc=c$ for every constant scalar function $c$ on $T_k$ and for every constant bond function $c$ on $T_k$. For every scalar $\lambda$ and every bond function $\mathscr{A}$ on $T_k$
\[
Q_k\,d\lambda=d\,Q_k\lambda,\qquad d\,Q_k\mathscr{A}=Q^{(2)}_k\,d\mathscr{A} .
\]
Let $|\cdot|_{\ell^2}$ be the counting $\ell^2$ norm of the functions on the plaquettes of $T^0_k$. For every function $F=(F_{\mu\nu})_{\mu<\nu}$ on the plaquettes of $T_k$, whose norm $|F|_2$ is formed on $T_k$ in the scaling of the level $k$ (Section~\ref{sec:averaging-fluctuations-gauge}),
\[
\left|Q^{(2)}_kF\right|_{\ell^2}\le C_Q\cdot|F|_2,\qquad\text{where }C_Q:=1 .
\]
\end{proposition}

This is a special case of Proposition~\ref{prop:averaging-operators}.

\begin{proposition}[Fluctuation fields]\label{prop:fluctuation-fields}
By Proposition~\ref{prop:spectral-bounds-decay}, and also by Proposition~\ref{prop:fluctuation-operator}(a), the operator $\tilde{\Delta}_k$ is positive definite on $V_k$; hence $C_k$, $C_k^{\frac{1}{2}}$ and $\mathrm{Sq}_k=C\,C_k^{\frac{1}{2}}$ are defined. One has
\[
Q\,\mathrm{Sq}_k=0,\qquad\tau\,\mathrm{Sq}_k=0,
\]
and for every $\zeta\in V_k$ the mean of the bond function $\mathrm{Sq}_k\zeta$ on $T^0_k$ vanishes.
\end{proposition}

The assertion is contained in Proposition~\ref{prop:fluctuation-operator}.

\begin{proposition}[Minimisers with prescribed block averages]\label{prop:minimisers-prescribed-block}
Let $A'$ be a real bond function on $T^0_{k+1}$, with arbitrary harmonic part, and let $(\cdot)_L$ be the scaling of bond functions of Section~\ref{sec:local-parts-adjustment}.
\begin{enumerate}
\item[(i)] The minimiser $A^{\min}_k(A')$ of Section~\ref{sec:averaging-fluctuations-gauge}, that is, the minimiser of $\langle A,\Delta_kA\rangle$ over the real bond functions $A$ on $T^0_k$ with $QA=A'_L$ and $\tau A=0$, exists and is unique, and it depends linearly on $A'$.
\item[(ii)] For every bond function $Z$ on $T^0_k$ with $QZ=0$ and $\tau Z=0$ one has $\left\langle A^{\min}_k(A'),\Delta_kZ\right\rangle=0$.
\item[(iii)] There is a real function $\boldsymbol{\omega}_k(A')$ on the sites of $T_k$, depending linearly on $A'$, such that
\[
H_k\,A^{\min}_k(A')=\left(H_{k+1}A'\right)_L-d\,\boldsymbol{\omega}_k(A').
\]
\item[(iv)] $\left|dH_k\,A^{\min}_k(A')\right|_2=\left|dH_{k+1}A'\right|_2$.
\item[(v)] If $A'$ is closed and has the mean $c'\in\mathbb{R}^3$, then $A^{\min}_k(A')$ is closed and has the mean $L^{-\frac{1}{2}}\cdot c'$.
\end{enumerate}
\end{proposition}

All five assertions are contained in Theorem~\ref{thm:minimisation-one-step}.

\begin{proposition}[Orthogonal coordinates on $F_N$]\label{prop:orthogonal-coordinates}
Let $F_N$ and $\gamma_N$ be as in Section~\ref{sec:lattice-measures-laws}; the scalar product $\langle\cdot,\cdot\rangle$ and the norm $|\cdot|$ of $F_N$ are those of the real two-cochains on $\Lambda_n$ with the counting $\ell^2$ norm. On each of the spaces $V_j$, $0\le j<k^{*}$, the norm $|\cdot|$ is the one that belongs to the counting scalar product of $V_j$ (Section~\ref{sec:averaging-fluctuations-gauge}), that is, the norm $|\cdot|_2$ of bond functions on the unit lattice $T^0_j$; since the basis~\eqref{eq:renormalisation-group-1} of $V_j$ is orthonormal, $|z|^2=\sum_{\Upsilon}z(\Upsilon)^2$ for every $z\in V_j$. Let $V_{\mathrm{t}}(T^0_{k^{*}})$ be the space of the coclosed real bond functions on $T^0_{k^{*}}$ with mean zero.
\begin{enumerate}
\item[(i)] The coboundary $d$ maps $\mathrm{Sl}$ isomorphically onto $F_N$.
\item[(ii)] For $F\in F_N$ let $A_0(F)$ be the element of $\mathrm{Sl}$ with $dA_0(F)=F$, and let $A_{*}$ and $\zeta_j\in V_j$, $0\le j<k^{*}$, be formed from $A_0=A_0(F)$ as in Definition~\ref{def:coordinates-slice-small}. The map
\[
F\mapsto\left(A_{*};\zeta_0,\dots,\zeta_{k^{*}-1}\right)
\]
is a linear bijection of $F_N$ onto $V_{\mathrm{t}}(T^0_{k^{*}})\times\prod_{j<k^{*}}V_j$. Its inverse is $(A_{*};\zeta)\mapsto d\,A^{\downarrow}_0(A_{*},\zeta)$, and for every $F\in F_N$
\[
|F|^2=\left|dH_{k^{*}}A_{*}\right|_2^2+\sum_{j<k^{*}}|\zeta_j|^2.
\]
\item[(iii)] Consequently there is an orthonormal family $(\phi_{j,\Upsilon})$ in $F_N$, indexed by the integers $0\le j<k^{*}$ and the elements $\Upsilon$ of the basis~\eqref{eq:renormalisation-group-1} of $V_j$, such that $\zeta_j(\Upsilon)=\langle\phi_{j,\Upsilon},F\rangle$ for every $F\in F_N$. This family is orthogonal to the linear functionals of $A_{*}$: if $f$ is a linear form on $V_{\mathrm{t}}(T^0_{k^{*}})$ and $\phi\in F_N$ satisfies $f(A_{*})=\langle\phi,F\rangle$ for every $F\in F_N$, then $\phi$ is orthogonal to every $\phi_{j,\Upsilon}$.
\item[(iv)] For every plaquette $p$ of $T^0_{k^{*}}$ there is $\phi_p\in F_N$ with $|\phi_p|\le C_Q$, the constant of Proposition~\ref{prop:block-averages}, such that $(dA_{*})(p)=\langle\phi_p,F\rangle$ for every $F\in F_N$.
\item[(v)] Under $\gamma_N$ the random variables $\zeta_j(\Upsilon)$, where $0\le j<k^{*}$ and $\Upsilon$ ranges over the basis of $V_j$, are independent standard normal variables, and they are independent of $A_{*}$.
\end{enumerate}
\end{proposition}

This is the special case $k=k^{*}$ of Theorem~\ref{thm:hierarchical-coordinates-curvatures}.

\begin{proposition}[Closed shifts of the terminal block field]\label{prop:closed-shifts-terminal}
Let $A^c$ be a closed real bond function on $T^0_{k^{*}}$ with the mean $c\in\mathbb{R}^3$, and let $L^{-\frac{k^{*}}{2}}\cdot c$ denote the constant bond function on $\Lambda_n$ with the components $L^{-\frac{k^{*}}{2}}\cdot c_{\mu}$. For every real bond function $A$ on $T^0_{k^{*}}$ and every family $\zeta=(\zeta_j)_{j<k^{*}}$ with $\zeta_j\in V_j$ there is a function $\lambda$ on the sites of $\Lambda_n$ such that
\[
A^{\downarrow}_0(A+A^c,\zeta)=A^{\downarrow}_0(A,\zeta)+L^{-\frac{k^{*}}{2}}\cdot c+d\lambda .
\]
In particular, for $\theta\in\mathbb{R}^3$, the twist $c_{\theta}$ of Definition~\ref{def:twist-external-connection} and the constant bond function $H_{\theta}$ of Section~\ref{sec:lattice-measures-laws}, there is a function $\lambda$ on the sites of $\Lambda_n$ such that
\[
A^{\downarrow}_0(A+c_{\theta},\zeta)=A^{\downarrow}_0(A,\zeta)+H_{\theta}+d\lambda .
\]
\end{proposition}

The proposition is the special case $k=k^{*}$ of Corollary~\ref{cor:adding-closed-field}.

\begin{proposition}[Comparison of the fermion operators of levels $k$ and $k+1$]\label{prop:comparison-fermion-operators}
Let $B'$ be a complex bond function on $T_{k+1}$ such that $B:=B'_L$ lies in $\tilde{G}_k$; the operators of the level $k+1$ (Section~\ref{sec:exact-fermion-operators}) are taken at $B'$. If $k=0$, the operators $\bar{H}_0$, $\bar{T}_0$ and $S^0_1$ are given by the formulas that define $\bar{H}_k$, $\bar{T}_k$ and $S^0_{k+1}$ in Section~\ref{sec:exact-fermion-operators}, taken at $k=0$: for a connection $A$ on $T_0$
\[
\bar{H}_0(A):=\mathsf{C}^{-1}\cdot H_0(-A)\cdot\mathsf{C}=1,\qquad\bar{T}_0(A):=\mathsf{C}^{-1}\cdot T_0(-A)\cdot\mathsf{C}=1,\qquad S^0_1(A):=S_{0,0}(A)=\Gamma_0(A),
\]
the last operator being defined whenever this inverse exists. A field on a lattice is an antiperiodic function on its sites whose four components lie in the odd part of a Grassmann algebra over $\mathbb{C}$. For a spinor function or a field $f$ on $T_{k+1}$, respectively on $T^0_{k+1}$, let $f_L$ be the function of the same kind on $T_k$, respectively on $L\mathbb{Z}^3/l\mathbb{Z}^3$, given by $f_L(x):=L^{-1}\cdot f(x/L)$. The pairings are those of Section~\ref{sec:fermion-integrals-blocking}.
\begin{enumerate}
\item[(a)] \emph{Intertwining relations.} For every spinor function $f$ on $T_{k+1}$ and every spinor function $v$ on $T^0_{k+1}$
\begin{align*}
S^0_{k+1}(B)f_L&=L\cdot\left(S_{k+1}(B')f\right)_L,\\
bL^{-1}\cdot H_k(B)\cdot\Gamma_k(B)\cdot Q^T(-B)v_L&=\left(H_{k+1}(B')v\right)_L,\\
\left(bL^{-1}-\left(bL^{-1}\right)^2\cdot Q(B)\cdot\Gamma_k(B)\cdot Q^T(-B)\right)v_L&=L^{-1}\cdot\left(D_{k+1}(B')v\right)_L.
\end{align*}
\item[(b)] \emph{Completion of the square.} Let $\Psi'$, $\bar{\Psi}'$ be fields on $T^0_{k+1}$, let $W'$, $\bar{W}'$ be fields on $T^0_k$, and put
\begin{gather*}
\psi^0:=\left(H_{k+1}(B')\Psi'\right)_L,\qquad\Psi^c:=T_k(B)\psi^0,\\
\bar{\psi}^0:=\left(\bar{H}_{k+1}(B')\bar{\Psi}'\right)_L,\qquad\bar{\Psi}^c:=\bar{T}_k(B)\bar{\psi}^0.
\end{gather*}
Then $\psi^0=H_k(B)\Psi^c$ and $\bar{\psi}^0=\bar{H}_k(B)\bar{\Psi}^c$, and for $\Psi:=\Psi^c+W'$ and $\bar{\Psi}:=\bar{\Psi}^c+\bar{W}'$
\begin{multline*}
\left\langle\bar{\Psi},D_k(B)\Psi\right\rangle+bL^{-1}\cdot\left\langle\bar{\Psi}'_L-Q(-B)\bar{\Psi},\,\Psi'_L-Q(B)\Psi\right\rangle\\
=\left\langle\bar{\Psi}',D_{k+1}(B')\Psi'\right\rangle+\left\langle\bar{W}',\left(D_k(B)+bL^{-1}\cdot P(B)\right)W'\right\rangle,
\end{multline*}
where the second pairing on the left side is the one on $L\mathbb{Z}^3/l\mathbb{Z}^3$, the first pairing on the right side is the one of the level $k+1$, formed on $T^0_{k+1}$, and the remaining two are formed on $T^0_k$.
\item[(c)] \emph{Determinants.} With $c'_k:=b^{4n_{k+1}}$,
\[
\det D_{k+1}(B')=c'_k\cdot\frac{\det D_k(B)}{\delta Z_k(B)} .
\]
\item[(d)] Under hypothesis~\ref{def:hypothesis-one-step} without primitives, which Theorem~\ref{thm:renormalisation-group-2}(b) assigns to this proposition, the data $(k,l,\mathfrak{m},e_k)$ belong to the scale-free setting of Section~\ref{sec:properties-scheme}, since $l\in\mathrm{Cl}$, $\mathfrak{m}\le\mathfrak{m}_0/L\le\mathfrak{m}_0$ and $e_k\le e_{k+1}\le e_0$. The inverses that occur in (a)--(c) exist. At the level $k$, the connection $-B$ lies in $\tilde{G}_k$ together with $B$ (Lemma~\ref{lem:invariances}); by the last paragraph of Section~\ref{sec:exact-fermion-operators}, which rests on Theorem~\ref{thm:expansion-walks}(a) and Proposition~\ref{prop:operators-fermion-1}(b),~(d), the inverses $S_k(\pm B)$ (if $k\ge1$), $S^0_{k+1}(\pm B)=S_{k,0}(\pm B)$ and $\Gamma_k(\pm B)$ exist, and $\delta Z_k(B)=\det\left(\Gamma_k(B)^{-1}\right)\ne0$. For $k\ge1$ this is contained also in Propositions~\ref{prop:decay-fine-propagator}--\ref{prop:kernel-bounds-2}, applied at $B$ and at $-B$; for $k=0$ one has $S^0_1(B)=S_{0,0}(B)=\Gamma_0(B)=\left(D_B+\mathfrak{m}+b\,L^{-1}\cdot P(B)\right)^{-1}$, and likewise at $-B$. These are the inverses that occur in $H_k(B)$, $D_k(B)$, $\Gamma_k(B)$, $S^0_{k+1}(B)$ and $1/\delta Z_k(B)$, and, through $S_k(-B)$, in $\bar{H}_k(B)$. At the level $k+1$ the inverses are $S_{k+1}(B')$, which enters $H_{k+1}(B')$ and $D_{k+1}(B')$, and $S_{k+1}(-B')$, which enters $\bar{H}_{k+1}(B')$. For every spinor function $f$ on $T_{k+1}$ one has, by Step~1 of the proof of Theorem~\ref{thm:fermion-operators-two},
\[
\left(D_B+\mathfrak{m}+b_{k+1}L^{-1}\cdot P_{k+1}(B)\right)f_L=L^{-1}\cdot\left(\left(D_{B'}+L\cdot\mathfrak{m}+b_{k+1}\cdot P_{k+1}(B')\right)f\right)_L,
\]
where $P_1(B)$ stands for $P(B)$ if $k=0$. By the definition $S^0_{k+1}(B)=S_{k,0}(B)$, the operator on the left side is $S^0_{k+1}(B)^{-1}$: for $k\ge1$ because $\boldsymbol{\alpha}_0=0$ and $\boldsymbol{\beta}_0=b_{k+1}L^{-1}$ (Section~\ref{sec:exact-fermion-operators}), for $k=0$ because $b_1=b$ (Section~\ref{sec:running-quantities-tori}). This operator is invertible by what precedes, and $f\mapsto f_L$ is a bijection of the spinor functions on $T_{k+1}$ onto the spinor functions on $T_k$; hence $D_{B'}+L\cdot\mathfrak{m}+b_{k+1}\cdot P_{k+1}(B')$ is invertible, that is, $S_{k+1}(B')$ exists (Theorem~\ref{thm:fermion-operators-two}(a)). The same argument at $-B=(-B')_L$ gives the existence of $S_{k+1}(-B')$ (Step~4 of the proof of Theorem~\ref{thm:fermion-operators-two}). Finally, (a)--(c) are identities between analytic functions of $B$ on $\tilde{G}_k$.
\end{enumerate}
\end{proposition}

This is a consequence of Theorem~\ref{thm:fermion-operators-two}.

\begin{proposition}[The step with an external connection]\label{prop:step-external-connection}
Let $k\ge0$, let $\boldsymbol{\varepsilon}$ and $m$ be real numbers with $|m|\le w^{\mathrm{m}}_k$, let $\mathscr{E}\in K^x_k$ with $\|\mathscr{E}\|^x_k\le w^{\mathrm{E}}_k$, where $\mathscr{E}=0$ if $k=0$, and let $A'\in\mathrm{Dom}_{k+1}$ and $\mathscr{C}'\in\mathfrak{X}_{k+1}$. With the numbers $\varepsilon(\cdot)$ and $m(\cdot)$, the operation $R$ and the scalings $(\cdot)_L$ of Section~\ref{sec:local-parts-adjustment}, and with the operators $H^{\#}_{k+1}(\cdot)$ of Section~\ref{sec:fermion-integrals-blocking}, define
\begin{gather*}
m':=m+m(\mathscr{E}),\qquad\mathscr{E}^r:=R\,\mathscr{E},\qquad m^{+}:=L\cdot m',\qquad\boldsymbol{\varepsilon}^{+}:=L^3\cdot\left(\boldsymbol{\varepsilon}+\varepsilon(\mathscr{E})+d\varepsilon^0_k\right),\\
\mathscr{A}':=H_{k+1}A',\qquad\mathscr{T}':=\mathscr{A}'+\mathscr{C}',\qquad\psi'^{\#}:=H^{\#}_{k+1}(\mathscr{T}')\Psi',\\
A:=\mathscr{A}'_L,\qquad\mathscr{C}:=\mathscr{C}'_L,\qquad\mathscr{T}:=A+\mathscr{C},\qquad\psi^{0,\#}:=\left(\psi'^{\#}\right)_L.
\end{gather*}
Let $\psi'$ and $\bar{\psi}'$ be the unbarred and the barred spinor field in $\psi'^{\#}$, and $\psi^0$, $\bar{\psi}^0$ those in $\psi^{0,\#}$. Let $W$, $\bar{W}$ be the generators of Section~\ref{sec:grassmann-algebras-norms}, and for $\mathsf{Z}\in\mathrm{Box}_k$ let
\begin{gather*}
\Psi^c:=T_k(\mathscr{T})\psi^0,\qquad\bar{\Psi}^c:=\bar{T}_k(\mathscr{T})\bar{\psi}^0,\qquad\mathscr{Z}:=H_k\,C\,C_k^{\frac{1}{2}}\mathsf{Z},\\
\Psi^f:=\Psi^c+\Gamma_k(\mathscr{T})W,\qquad\bar{\Psi}^f:=\bar{\Psi}^c+\bar{W},
\end{gather*}
\begin{multline*}
V^{\mathrm{ex},x}_k:=\mathrm{Act}^x_k[m',\mathscr{E}^r](A+\mathscr{Z},\mathscr{C};\Psi^f)-\mathrm{Act}^x_k[m',\mathscr{E}^r](A,\mathscr{C};\Psi^f)\\
+I^x_k[m',\mathscr{E}^r](A,\mathscr{C};\Psi^f)-I^x_k[m',\mathscr{E}^r](A,\mathscr{C};\Psi^c),
\end{multline*}
where the last argument $\Psi^f$, resp.\ $\Psi^c$, of the functionals $\mathrm{Act}^x_k$ and $I^x_k$ of Section~\ref{sec:fermion-integrals-blocking} means that the pair $(\Psi,\bar{\Psi})$ is replaced by $(\Psi^f,\bar{\Psi}^f)$, resp.\ by $(\Psi^c,\bar{\Psi}^c)$, and
\[
\Xi^x_k:=\int\mathrm{d}\mu_I(W)\,\mathrm{d}\hat{\mu}(\mathsf{Z})\,\exp\left(V^{\mathrm{ex},x}_k\right),
\]
where $\mathrm{d}\mu_I(W)$ is the Grassmann Gaussian functional of Proposition~\ref{prop:weighted-coefficient-norms} and $\hat{\mu}$ is the probability measure on $V_k$ of Section~\ref{sec:averaging-fluctuations-gauge}. Let $\mathfrak{P}_k=\left\{\mathfrak{P}_k(X,\cdot\,):\ X\in D_k\right\}$ be an arbitrary family of complex functions on $\mathfrak{X}_k$, and let $\mathscr{F}^x_k$ be the density of Section~\ref{sec:fermion-integrals-blocking} formed with this family; in particular, $\mathfrak{P}_k$ may be the family of level $k$ of a trajectory with slot (Section~\ref{sec:flow}). Then, with the pairings of Section~\ref{sec:fermion-integrals-blocking} formed at the level $k+1$,
\begin{multline}\label{eq:renormalisation-group-9}
\left(\mathbf{T}^x_k\mathscr{F}^x_k[\boldsymbol{\varepsilon},m,\mathscr{E}]\right)(A',\mathscr{C}';\Psi')=\left(\det D_{k+1}(0)\right)^{-1}\\
\cdot\exp\left\{-\left\langle\bar{\Psi}',D_{k+1}(\mathscr{T}')\Psi'\right\rangle-m^{+}\cdot\left\langle\bar{\psi}',\psi'\right\rangle-\boldsymbol{\varepsilon}^{+}\cdot n_{k+1}\right.\\
\left.+\sum_{X\in D_k}\mathscr{E}^r\left(X;A,\mathscr{C},\psi^{0,\#}\right)+\sum_{X\in D_k}\mathfrak{P}_k(X,\mathscr{C})\right\}\cdot\frac{\delta Z_k(\mathscr{T})}{\delta Z_k(0)}\cdot\Xi^x_k.
\end{multline}
The identity~\eqref{eq:renormalisation-group-9} holds as well for complex numbers $\boldsymbol{\varepsilon}$ and $m$ and for $\mathscr{E}\in K^{x,C}_k$ with $\mathscr{E}=0$ if $k=0$, without the bounds on $m$ and on $\mathscr{E}$, for all $A'\in\mathrm{Dom}_{k+1}$ and $\mathscr{C}'\in\mathfrak{X}_{k+1}$ and for every family $\mathfrak{P}_k$ as above.
\end{proposition}

This is proved as Theorem~\ref{thm:one-step-formula}.

\subsubsection{Norms, polymers and operations on activities}\label{sec:norms-polymers-operations}

The statements of this subsection hold under the hypotheses specified in Theorem~\ref{thm:renormalisation-group-1}. The primitives of Section~\ref{sec:primitive-operators} occur in none of them except Proposition~\ref{prop:free-determinant}, in which they enter through $E^{\mathrm{det}}_j$ (Definition~\ref{def:determinant-functions}) and $\mathfrak{P}_k$ (Section~\ref{sec:flow}), by Proposition~\ref{prop:local-expansion-determinant}(d0).

\begin{proposition}[Weighted coefficient norms]\label{prop:weighted-coefficient-norms}
Let $S_1,\dots,S_p$ be finite, totally ordered index sets with weights $\nu_i\colon S_i\to(0,\infty)$, and for a function $f$ on $S_i$ write $\int\mathrm{d}s\,f(s):=\sum_{s\in S_i}\nu_i(s)\cdot f(s)$. Let $\mathrm{Gr}[\boldsymbol{\phi}_1,\dots,\boldsymbol{\phi}_p]$ be the Grassmann algebra over $\mathbb{C}$ with the independent generators $\boldsymbol{\phi}_i(s)$, $s\in S_i$, $1\le i\le p$. For $I\subseteq S_i$ let $\boldsymbol{\phi}_i(I)$ be the product of the generators $\boldsymbol{\phi}_i(s)$, $s\in I$, taken in increasing order, the empty product being $1$. Every $F\in\mathrm{Gr}[\boldsymbol{\phi}_1,\dots,\boldsymbol{\phi}_p]$ has exactly one representation
\[
F=\sum_{I_1\subseteq S_1,\dots,I_p\subseteq S_p}c_F(I_1,\dots,I_p)\cdot\boldsymbol{\phi}_1(I_1)\cdots\boldsymbol{\phi}_p(I_p)
\]
with complex coefficients $c_F(I_1,\dots,I_p)$, the sum being taken over all $p$-tuples of subsets. For $g=(g_1,\dots,g_p)$ with $g_i>0$ put
\[
\|F\|_g:=\sum_{I_1,\dots,I_p}\left|c_F(I_1,\dots,I_p)\right|\cdot g_1^{\#|I_1|}\cdots g_p^{\#|I_p|}.
\]
Let $\{\xi\}$ be the set of all indices $\xi=(x,\boldsymbol{\beta},\boldsymbol{\omega})$ with $x\in T_k$, let $\{\boldsymbol{\zeta}\}$ be the set of all indices $\boldsymbol{\zeta}=(x,x',\boldsymbol{\beta},\boldsymbol{\omega})$ with $(x,x')\in\mathrm{Pz}_k$, and let $\{u\}$ be the set of all indices $u=(y,\boldsymbol{\beta},\boldsymbol{\omega})$ with $y\in T^0_k$; here $\boldsymbol{\beta}$ runs through the spinor indices and $\boldsymbol{\omega}$ through $\{0,1\}$, as in Section~\ref{sec:grassmann-algebras-norms}, so that these are the index sets of the generators $\psi(\xi)$, of the generators $\chi(\boldsymbol{\zeta})$, and of the generators $\Psi(u)$ and $W(u)$. Each of the three sets carries a total order, which may be chosen arbitrarily: a change of the order multiplies every coefficient $c_F(I_1,\dots,I_p)$ by $1$ or by $-1$, and therefore leaves $\|F\|_g$ unchanged. For $p=2$, with $S_1=\{\xi\}$, $\nu_1=\eta_k^3$, $S_2=\{\boldsymbol{\zeta}\}$, $\nu_2=\eta_k^6$ and $g=h$, the norm $\|\cdot\|_g$ is the norm $\|E\|_h$ of an element $E\in\mathrm{Gr}[\psi,\chi]$ defined in Section~\ref{sec:grassmann-algebras-norms}; for $p=1$, with $S_1=\{u\}$, the counting weight and $g=a$, it is the norm $|F|_a$ of an element $F\in\mathrm{Gr}[\Psi]$ defined in Section~\ref{sec:grassmann-algebras-norms}.

For an operator $O$ from the functions on an index set $S$ with weights $\nu$ to the functions on an index set $T$, acting by $(Of)(t)=\sum_{s\in S}O(t,s)\cdot f(s)\cdot\nu(s)$, put
\[
|O|_{1,\infty}:=\sup_{t\in T}\sum_{s\in S}\left|O(t,s)\right|\cdot\nu(s).
\]

Let $W$ and $\bar{W}$ denote the unbarred and the barred generators $W(u)$ of Section~\ref{sec:grassmann-algebras-norms}, with the components $W_{\boldsymbol{\beta}}(y):=W(y,\boldsymbol{\beta},0)$ and $\bar{W}_{\boldsymbol{\beta}}(y):=W(y,\boldsymbol{\beta},1)$; let $\left\langle\bar{W},OW\right\rangle$ be the pairing of Section~\ref{sec:fermion-integrals-blocking}, for an operator $O$ on the spinors on $T^0_k$, that is, a complex matrix with rows and columns indexed by the pairs $(y,\boldsymbol{\beta})$; and let $\int\mathrm{D}W$ be the Berezin integral over these generators, normalised as $\int\mathrm{D}\Psi$ in Section~\ref{sec:fermion-integrals-blocking}. The algebra $\mathrm{Gr}[\boldsymbol{\phi}_1,\dots,\boldsymbol{\phi}_p,W]$ is formed with $W$ as a further family of generators, indexed by $\{u\}$ with the counting weight. For $F\in\mathrm{Gr}[\boldsymbol{\phi}_1,\dots,\boldsymbol{\phi}_p,W]$ and an invertible complex matrix $\operatorname{Cov}$ with rows and columns indexed by the pairs $(y,\boldsymbol{\beta})$, the elements $\int\mathrm{d}\mu_I(W)\,F$ and $\int\mathrm{d}\mu_{\operatorname{Cov}}(W)\,F$ of $\mathrm{Gr}[\boldsymbol{\phi}_1,\dots,\boldsymbol{\phi}_p]$ are defined by
\begin{align*}
\int\mathrm{d}\mu_I(W)\,F&:=\int\mathrm{D}W\,\exp\left(-\left\langle\bar{W},W\right\rangle\right)\cdot F,\\
\int\mathrm{d}\mu_{\operatorname{Cov}}(W)\,F&:=\det(\operatorname{Cov})\cdot\int\mathrm{D}W\,\exp\left(-\left\langle\bar{W},\operatorname{Cov}^{-1}W\right\rangle\right)\cdot F.
\end{align*}
The same formulas, with $\Psi$, $\bar{\Psi}$ and $\int\mathrm{D}\Psi$ in place of $W$, $\bar{W}$ and $\int\mathrm{D}W$, define $\int\mathrm{d}\mu_I(\Psi)$ and $\int\mathrm{d}\mu_{\operatorname{Cov}}(\Psi)$. Then the following assertions hold.
\begin{enumerate}
\item[(N0)]\itemmark{N0}\label{it:weighted-coefficient-norms-N0} \emph{Kernels.} For every $n=(n_1,\dots,n_p)$ with integers $n_i\ge0$ let $K_n$ be an arbitrary complex function on $S_1^{n_1}\times\dots\times S_p^{n_p}$. For an element $s=\left(s^1_1,\dots,s^1_{n_1};\dots;s^p_1,\dots,s^p_{n_p}\right)$ of this set put
\[
\boldsymbol{\phi}(s):=\boldsymbol{\phi}_1\left(s^1_1\right)\cdots\boldsymbol{\phi}_1\left(s^1_{n_1}\right)\cdots\boldsymbol{\phi}_p\left(s^p_1\right)\cdots\boldsymbol{\phi}_p\left(s^p_{n_p}\right),
\]
and let $\int\mathrm{d}s$ integrate every argument $s^i_j$ over $S_i$ with the weight $\nu_i$. If
\[
F=\sum_n\left(\prod_{i=1}^p\frac{1}{n_i!}\right)\cdot\int\mathrm{d}s\,K_n(s)\cdot\boldsymbol{\phi}(s),
\]
then
\[
\|F\|_g\le\sum_n\left(\prod_{i=1}^p\frac{g_i^{n_i}}{n_i!}\right)\cdot\int\mathrm{d}s\,\left|K_n(s)\right|,
\]
with equality if every $K_n$ is antisymmetric in each of its $p$ groups of arguments. The antisymmetric kernels of such a representation are determined by $F$.
\item[(N1)]\itemmark{N1}\label{it:weighted-coefficient-norms-N1} \emph{Products.} For $F,G\in\mathrm{Gr}[\boldsymbol{\phi}_1,\dots,\boldsymbol{\phi}_p]$,
\[
\|F\cdot G\|_g\le\|F\|_g\cdot\|G\|_g.
\]
The even elements, normed by $\|\cdot\|_g$, constitute a commutative Banach algebra with a unit element, and every even $F$ satisfies $\|\exp F-1\|_g\le\exp\left(\|F\|_g\right)-1$.
\item[(N2)]\itemmark{N2}\label{it:weighted-coefficient-norms-N2} \emph{Linear substitutions.} Let $\boldsymbol{\theta}_1,\dots,\boldsymbol{\theta}_q$ be families of independent generators, indexed by finite, totally ordered sets $S'_1,\dots,S'_q$ with weights $\nu'_j\colon S'_j\to(0,\infty)$, and let $g'=(g'_1,\dots,g'_q)$ with $g'_j>0$. For $1\le i\le p$ and $1\le j\le q$ let $O_{ij}$ be an operator from the functions on $S'_j$, with the weights $\nu'_j$, to the functions on $S_i$, and let $C_{ij}$ be numbers with $|O_{ij}|_{1,\infty}\le C_{ij}$. Let $F\mapsto F'$ be the homomorphism of $\mathrm{Gr}[\boldsymbol{\phi}_1,\dots,\boldsymbol{\phi}_p]$ to $\mathrm{Gr}[\boldsymbol{\theta}_1,\dots,\boldsymbol{\theta}_q]$ with
\[
\boldsymbol{\phi}_i(t)\mapsto\sum_{j=1}^q\left(O_{ij}\boldsymbol{\theta}_j\right)(t)=\sum_{j=1}^q\sum_{s\in S'_j}O_{ij}(t,s)\cdot\boldsymbol{\theta}_j(s)\cdot\nu'_j(s)\qquad\left(1\le i\le p,\ t\in S_i\right).
\]
If $\sum_{j=1}^qC_{ij}\cdot g'_j\le g_i$ for $1\le i\le p$, then
\[
\|F'\|_{g'}\le\|F\|_g\qquad\left(F\in\mathrm{Gr}[\boldsymbol{\phi}_1,\dots,\boldsymbol{\phi}_p]\right).
\]
\item[(N3)]\itemmark{N3}\label{it:weighted-coefficient-norms-N3} \emph{Integration.} For $F\in\mathrm{Gr}[\boldsymbol{\phi}_1,\dots,\boldsymbol{\phi}_p,W]$,
\[
\left\|\int\mathrm{d}\mu_I(W)\,F\right\|_{(g_1,\dots,g_p)}\le\|F\|_{(g_1,\dots,g_p,1)}.
\]
\item[(N4)]\itemmark{N4}\label{it:weighted-coefficient-norms-N4} \emph{Factorisation.} An element $F\in\mathrm{Gr}[\boldsymbol{\phi}_1,\dots,\boldsymbol{\phi}_p,W]$ is said to contain a generator $W(y,\cdot,\cdot)$ at the unit site $y$ if, for some $\boldsymbol{\beta}$ and $\boldsymbol{\omega}$, the generator $W(y,\boldsymbol{\beta},\boldsymbol{\omega})$ is a factor of a basis monomial that has a non-zero coefficient in $F$. Let $F_1,F_2\in\mathrm{Gr}[\boldsymbol{\phi}_1,\dots,\boldsymbol{\phi}_p,W]$, and assume that the set of the unit sites at which $F_1$ contains a generator $W(y,\cdot,\cdot)$ is disjoint from the set of the unit sites at which $F_2$ contains one. Then
\[
\int\mathrm{d}\mu_I(W)\,\left(F_1\cdot F_2\right)=\left(\int\mathrm{d}\mu_I(W)\,F_1\right)\cdot\left(\int\mathrm{d}\mu_I(W)\,F_2\right).
\]
\item[(N4')]\itemmark{N4'}\label{it:weighted-coefficient-norms-N4p} \emph{Invariance.} Let $U$ be an invertible complex matrix with rows and columns indexed by the pairs $(y,\boldsymbol{\beta})$, and let $s_U$ be the homomorphism of $\mathrm{Gr}[\boldsymbol{\phi}_1,\dots,\boldsymbol{\phi}_p,W]$ to itself with $W\mapsto UW$ and $\bar{W}\mapsto\left(U^{-1}\right)^T\bar{W}$ that fixes the generators $\boldsymbol{\phi}_i(s)$. Then
\[
\int\mathrm{d}\mu_I(W)\,s_U(F)=\int\mathrm{d}\mu_I(W)\,F\qquad\left(F\in\mathrm{Gr}[\boldsymbol{\phi}_1,\dots,\boldsymbol{\phi}_p,W]\right).
\]
\item[(N5)]\itemmark{N5}\label{it:weighted-coefficient-norms-N5} \emph{Contour integrals.} Let $\rho>1$, and let $f$ be an analytic function on a neighbourhood of the disc $\{t\in\mathbb{C}:|t|\le\rho\}$ with values in a normed space of finite dimension. Let $\oint_{\rho}f$ be as in Definition~\ref{def:contour-functional-functions}, and put
\[
\oint^{(2)}_{\rho}f:=\frac{1}{2\pi i}\oint_{|t|=\rho}\mathrm{d}t\,\frac{f(t)}{t^2\cdot(t-1)}.
\]
Then
\begin{align*}
\oint_{\rho}f&=f(1)-f(0),&\left\|\oint_{\rho}f\right\|&\le\frac{1}{\rho-1}\cdot\sup_{|t|=\rho}\|f(t)\|,\\
\oint^{(2)}_{\rho}f&=f(1)-f(0)-f'(0),&\left\|\oint^{(2)}_{\rho}f\right\|&\le\frac{1}{\rho\cdot(\rho-1)}\cdot\sup_{|t|=\rho}\|f(t)\|.
\end{align*}
Let, moreover, $\boldsymbol{\theta}_1,\dots,\boldsymbol{\theta}_p$ be families of independent generators, $\boldsymbol{\theta}_i$ indexed by $S_i$, and let $F(\boldsymbol{\phi}+\boldsymbol{\theta})$ be the image of $F\in\mathrm{Gr}[\boldsymbol{\phi}_1,\dots,\boldsymbol{\phi}_p]$ under the homomorphism into $\mathrm{Gr}[\boldsymbol{\phi}_1,\dots,\boldsymbol{\phi}_p,\boldsymbol{\theta}_1,\dots,\boldsymbol{\theta}_p]$ with $\boldsymbol{\phi}_i(s)\mapsto\boldsymbol{\phi}_i(s)+\boldsymbol{\theta}_i(s)$, and $F(\boldsymbol{\phi}):=F$. If the weights $g'=(g'_1,\dots,g'_p)$ of the families $\boldsymbol{\phi}_i$ and $g''=(g''_1,\dots,g''_p)$ of the families $\boldsymbol{\theta}_i$ satisfy $g'_i+\rho\cdot g''_i\le g_i$ for all $i$, then
\[
\left\|F(\boldsymbol{\phi}+\boldsymbol{\theta})-F(\boldsymbol{\phi})\right\|_{(g',g'')}\le\frac{\|F\|_g}{\rho-1}.
\]
\item[(N6)]\itemmark{N6}\label{it:weighted-coefficient-norms-N6} \emph{Monotonicity in the weights.} If $0<g'_i\le g_i$ for all $i$, then $\|F\|_{g'}\le\|F\|_g$. If, in addition, $r_1,\dots,r_p$ are integers with $r_i\ge0$, and every basis monomial $\boldsymbol{\phi}_1(I_1)\cdots\boldsymbol{\phi}_p(I_p)$ with $c_F(I_1,\dots,I_p)\ne0$ has the degree $\#|I_i|\ge r_i$ in $\boldsymbol{\phi}_i$ for every $i$, then
\[
\|F\|_{g'}\le\prod_{i=1}^p\left(\frac{g'_i}{g_i}\right)^{r_i}\cdot\|F\|_g.
\]
\item[(N7)]\itemmark{N7}\label{it:weighted-coefficient-norms-N7} \emph{Norms on $\mathrm{Gr}[\Psi]$.} Let $a>0$ be a number, and let $|\cdot|_a$ be the norm of Section~\ref{sec:grassmann-algebras-norms}.
\begin{enumerate}
\item[(i)] $|F\cdot G|_a\le|F|_a\cdot|G|_a$ for all $F,G\in\mathrm{Gr}[\Psi]$.
\item[(ii)] $\left|\int\mathrm{d}\mu_I(\Psi)\,F\right|\le|F|_1$ for every $F\in\mathrm{Gr}[\Psi]$.
\item[(iii)] For every invertible complex matrix $\operatorname{Cov}$ with rows and columns indexed by the pairs $(y,\boldsymbol{\beta})$ and every $f=f(\bar{\Psi},\Psi)\in\mathrm{Gr}[\Psi]$,
\[
\int\mathrm{d}\mu_{\operatorname{Cov}}(\Psi)\,f(\bar{\Psi},\Psi)=\int\mathrm{d}\mu_I(\Psi)\,f(\bar{\Psi},\operatorname{Cov}\Psi),
\]
where $f(\bar{\Psi},\operatorname{Cov}\Psi)$ is the image of $f$ under the homomorphism with $\Psi\mapsto\operatorname{Cov}\Psi$ and $\bar{\Psi}\mapsto\bar{\Psi}$. Every complex matrix $O$ with rows and columns indexed by the pairs $(y,\boldsymbol{\beta})$ satisfies
\[
\int\mathrm{D}\Psi\,\exp\left(-\left\langle\bar{\Psi},O\Psi\right\rangle\right)=\det O.
\]
\item[(iv)] Let $h_1,h_2>0$, let $C_0>0$ be an arbitrary number, and let $H^{(1)}$ and $H^{(2)}$ be operators from the functions on $\{u\}$, with the counting weight, to the functions on $\{\xi\}$, resp.\ on $\{\boldsymbol{\zeta}\}$, such that $\left|H^{(i)}\right|_{1,\infty}\le C_0$ for $i=1,2$. For $E\in\mathrm{Gr}[\psi,\chi]$ let $E\left(H^{(1)}\Psi,H^{(2)}\Psi\right)\in\mathrm{Gr}[\Psi]$ be the image of $E$ under the homomorphism with $\psi(\xi)\mapsto\left(H^{(1)}\Psi\right)(\xi)$ and $\chi(\boldsymbol{\zeta})\mapsto\left(H^{(2)}\Psi\right)(\boldsymbol{\zeta})$. If $C_0\cdot a\le\min(h_1,h_2)$, then
\[
\left|E\left(H^{(1)}\Psi,H^{(2)}\Psi\right)\right|_a\le\|E\|_{(h_1,h_2)}\qquad\left(E\in\mathrm{Gr}[\psi,\chi]\right).
\]
\end{enumerate}
\end{enumerate}
Let $H$ be an operator from the functions on $\{u\}$ to the functions on $\{\xi\}$, let $B$ be a bond function on $T_k$, and let $\delta_{\alpha,B}H$ be the composition of $H$ with the H\"older derivative $\delta_{\alpha,B}$ of Definition~\ref{def:pairs-sites-holder}, an operator from the functions on $\{u\}$ to the functions on $\{\boldsymbol{\zeta}\}$. In~\eqref{it:weighted-coefficient-norms-N2} with the substitution $\chi\mapsto\delta_{\alpha,B}H\Psi$, that is, with $\boldsymbol{\phi}_i=\chi$, $\boldsymbol{\theta}_j=\Psi$ and $O_{ij}=\delta_{\alpha,B}H$, and in~\eqref{it:weighted-coefficient-norms-N7}(iv), the supremum in the definition of $\left|O_{ij}\right|_{1,\infty}$, resp.\ of $\left|H^{(2)}\right|_{1,\infty}$, runs over all indices $\boldsymbol{\zeta}$, that is, over all pairs $(x,x')$ of $\mathrm{Pz}_k$, as in Proposition~\ref{prop:bounds}.
\end{proposition}

This is a consequence of Theorem~\ref{thm:polymer-activities-1}.

\begin{proposition}[Polymers]\label{prop:polymers}
Let $c_X$, $\kappa_0$, $K_0$, $c_{\mathrm{rb}}$ be the absolute constants of stage~\ref{it:order-choice-1}, and let $k\ge0$ and $l\in\mathrm{Cl}$. The $M$-cubes, the polymers, the set $D_k$, the numbers $|X|_M$ and $d_M(X)$, the nets, the small polymers and the sets $\bar{X}^L$ are those of Section~\ref{sec:lattice-geometry}, formed on the torus $T_k$ of side $l$. The $M$-cubes of $\mathbb{R}^3$ are the cubes of side $M$ centred at the points of $M\mathbb{Z}^3$; for a finite union $\tilde{X}$ of $M$-cubes of $\mathbb{R}^3$ that is connected through faces, $d_M(\tilde{X})$ is defined as in Section~\ref{sec:lattice-geometry}, with nets in $\mathbb{R}^3$ in place of nets in $\mathbb{R}^3/l\mathbb{Z}^3$. Two polymers are called \emph{incompatible} when their union is connected through faces; otherwise they are called \emph{compatible}. For polymers $X_1,\dots,X_n$ put $\rho^T(X_1):=1$ and, if $n\ge2$,
\[
\rho^T(X_1,\dots,X_n):=\sum_G(-1)^{\#|G|},
\]
where $G$ ranges over the connected graphs with the vertex set $\{1,\dots,n\}$ each of whose edges $\{i,j\}$ joins two incompatible polymers $X_i$, $X_j$, and $\#|G|$ is the number of the edges of $G$. Then the following assertions hold.
\begin{enumerate}
\item[(t0)]\itemmark{t0}\label{it:polymers-t0} $d_M(\square)=0$ for every $M$-cube $\square$, and $d_M(X)\le|X|_M-1$ for every $X\in D_k$.
\item[(t1)]\itemmark{t1}\label{it:polymers-t1} $|X|_M\le c_X\cdot\left(d_M(X)+1\right)$ for every $X\in D_k$.
\item[(t2)]\itemmark{t2}\label{it:polymers-t2} For every $M$-cube $\square$,
\[
\sum_{X\in D_k:\,X\supseteq\square}\exp\left(-\kappa_0\cdot d_M(X)\right)\le K_0;
\]
consequently, for every union $Y$ of $M$-cubes,
\[
\sum_{X\in D_k:\,X\subseteq Y}\exp\left(-\kappa_0\cdot d_M(X)\right)\le K_0\cdot|Y|_M.
\]
\item[(t3)]\itemmark{t3}\label{it:polymers-t3} If $X$ and $Y'$ are polymers with $X\subseteq Y'$, then $d_M(Y')\le d_M(X)+|Y'\setminus X|_M$.
\item[(t3')]\itemmark{t3'}\label{it:polymers-t3p} If $X$ and $Y'$ are as in~\eqref{it:polymers-t3}, then $d_M(X)\le d_M(Y')$.
\item[(t4)]\itemmark{t4}\label{it:polymers-t4} Let $X_1,\dots,X_n$ be polymers, and assume that the graph with the vertex set $\{1,\dots,n\}$ in which two vertices $i\ne j$ are joined whenever $X_i\cup X_j$ is connected through faces is connected. Then $X_1\cup\dots\cup X_n$ is a polymer, and
\[
d_M\left(X_1\cup\dots\cup X_n\right)\le\sum_{i=1}^nd_M(X_i)+3(n-1).
\]
\item[(t5)]\itemmark{t5}\label{it:polymers-t5} Assume, in addition, $l/L\in\mathrm{Cl}$. Let $D_{k+1}$ be the set of the polymers of level $k+1$, formed on the torus of side $l/L$, and for $Y\in D_{k+1}$ let $LY$ be the image of $Y$ under the map $x\mapsto Lx$ of this torus onto that of level $k$. Each $X\in D_k$ determines a unique $Y\in D_{k+1}$ with $\bar{X}^L=LY$, and for this $Y$ one has $L\cdot d_M(Y)\le d_M(X)$.
\item[(t6)]\itemmark{t6}\label{it:polymers-t6} Every real number $d\ge0$ satisfies $d+1\le\exp(d)$ and $d+1\le2\exp\left(\frac{d}{2}\right)$.
\item[(t7)]\itemmark{t7}\label{it:polymers-t7} Let $X\in D_k$. Any two sites $x$, $x'$ of $X$ satisfy
\[
d(x,x')\le M\cdot\left(d_M(X)+2\right),
\]
$d$ being the sup distance on the torus (Definition~\ref{def:sup-norm-distance}). The same inequality holds for any two points $x$, $x'$ of the closed $M$-cubes of $X$, with the sup distance of $\mathbb{R}^3/l\mathbb{Z}^3$. Moreover, if $\tilde{X}$ is a finite union of $M$-cubes of $\mathbb{R}^3$ that is connected through faces, then any two points $x$, $x'$ of its closed cubes satisfy $|x-x'|\le M\cdot\left(d_M(\tilde{X})+2\right)$, where $d_M(\tilde{X})$ is formed with nets in $\mathbb{R}^3$.
\item[(t8)]\itemmark{t8}\label{it:polymers-t8} In the situation of~\eqref{it:polymers-t5}, every $Y\in D_{k+1}$ and every $\kappa'\ge\kappa_0$ satisfy
\[
\sum_{X\in D_k:\,\bar{X}^L=LY}\exp\left(-\kappa'\cdot d_M(X)\right)\le c_{\mathrm{rb}}\cdot L^3\cdot\exp\left(-L\cdot\left(\kappa'-\kappa_0-1\right)\cdot d_M(Y)\right).
\]
\item[(t9)]\itemmark{t9}\label{it:polymers-t9} For every polymer $Y'$ there is a set of at most $7\cdot|Y'|_M$ $M$-cubes such that every polymer that is incompatible with $Y'$ contains one of these cubes.
\item[(t10)]\itemmark{t10}\label{it:polymers-t10} \emph{Lifts.} Let $X\in D_k$ satisfy $2M\cdot\left(d_M(X)+2\right)<l$. Then there is a finite union $\tilde{X}$ of $M$-cubes of $\mathbb{R}^3$, connected through faces, such that the projection of $\mathbb{R}^3$ onto $\mathbb{R}^3/l\mathbb{Z}^3$ is injective on $\tilde{X}$ and maps $\tilde{X}$ onto $X$; any two such unions differ by a translation by an element of $l\mathbb{Z}^3$, and $d_M(\tilde{X})=d_M(X)$, where $d_M(\tilde{X})$ is formed with nets in $\mathbb{R}^3$. Conversely, let $\tilde{X}$ be a non-empty finite union of $M$-cubes of $\mathbb{R}^3$ that is connected through faces and satisfies $2M\cdot\left(d_M(\tilde{X})+2\right)<l$; then the projection is injective on $\tilde{X}$, it maps $\tilde{X}$ onto a polymer, and the number $d_M$ of this polymer equals $d_M(\tilde{X})$. In particular, let $\square$ be an $M$-cube, and fix a lift of $\square$, that is, an $M$-cube of $\mathbb{R}^3$ whose image under the projection is $\square$. Through the projection, the small polymers that contain $\square$ correspond bijectively to the finite unions $\tilde{X}$ of $M$-cubes of $\mathbb{R}^3$ that are connected through faces, contain the fixed lift and satisfy $d_M(\tilde{X})\le L$; this set of unions does not depend on $l$.
\item[(CE1)]\itemmark{CE1}\label{it:polymers-CE1} In the ring of formal power series in commuting variables $z_X$, $X\in D_k$,
\[
\sum_{\mathcal{X}}\ \prod_{X\in\mathcal{X}}z_X=\exp\left(\sum_{n\ge1}\frac{1}{n!}\sum_{(X_1,\dots,X_n)\in D_k^n}\rho^T(X_1,\dots,X_n)\cdot\prod_{i=1}^nz_{X_i}\right),
\]
where $\mathcal{X}$ ranges over the finite sets of distinct, pairwise compatible polymers, the empty set included.
\item[(CE2)]\itemmark{CE2}\label{it:polymers-CE2} For all polymers $X_1,\dots,X_n$, the number $\left|\rho^T(X_1,\dots,X_n)\right|$ does not exceed the number of the trees with the vertex set $\{1,\dots,n\}$ in which every edge $\{i,j\}$ joins two incompatible polymers $X_i$, $X_j$.
\end{enumerate}
\end{proposition}

The assertions are contained in Theorem~\ref{thm:polymer-activities-2}.

\begin{proposition}[Symmetries and extraction]\label{prop:symmetries-extraction}
Assume the hypotheses that Theorem~\ref{thm:renormalisation-group-1} states for this proposition, and let $g=(S,a)\in\mathbb{G}_k$ act as in Section~\ref{sec:dirac-matrices-lattice}.
\begin{enumerate}
\item[(a)] The map $x\mapsto gx$ is an isometry for the sup norm, for the $\ell^1$ norm and for the Euclidean norm. It maps each of the lattices $L^{-i}\mathbb{Z}^3$, $-1\le i\le k$, onto itself, every block onto a block and the centre of a block onto the centre of its image, and the rectilinear path $\Gamma^{\pi}(y,x)$ onto the path $\Gamma^{\pi}(gy,gx)$ with the same order $\pi$. It maps $M$-cubes onto $M$-cubes and polymers onto polymers, and for every polymer $X$ and every integer $i\ge0$
\[
|gX|_M=|X|_M,\qquad d_M(gX)=d_M(X),\qquad(gX)^{\sim i}=g\left(X^{\sim i}\right).
\]
The cubes of $\pi_k$, resp.\ of $\pi'_k$, are mapped onto cubes of the same kind, and linked cubes onto linked cubes; $K_{gz}=g\,K_z$ for every site $z$. The map $(x,x')\mapsto(gx,gx')$ carries $\mathrm{Pz}_k$ onto $\mathrm{Pz}_k$, and a pair that is admissible for $X$ of kind~(I), resp.\ of kind~(II), onto a pair that is admissible for $gX$ of the same kind (Definition~\ref{def:pairs-sites-holder}). If $g\in\mathbb{G}_k^L$, then $x\mapsto gx$ maps $LM$-cubes onto $LM$-cubes and the cubes of $\pi''_k$ onto cubes of $\pi''_k$, and $g\left(\bar{X}^L\right)$ is the union of the $LM$-cubes that contain an $M$-cube of $gX$.
\item[(b)] The map $A\mapsto gA$ preserves the sup norm, the $\ell^1$ norm and the $\ell^2$ norm, the pairings of Section~\ref{sec:averaging-fluctuations-gauge} and the three norms of Section~\ref{sec:domains-bond-functions}, and
\[
(\tau\,gA)(gy,gx)=(\tau A)(y,x),\qquad g(d\boldsymbol{\lambda})=d\left(\boldsymbol{\lambda}\circ g^{-1}\right)
\]
for every bond function $A$ and every scalar $\boldsymbol{\lambda}$, where $\tau$ is the axial functional of Section~\ref{sec:averaging-fluctuations-gauge}. The elements of $\mathrm{Sp}$ and the matrices $\mathsf{C}$ and $\varkappa$ are monomial: each row and each column of such a matrix contains exactly one non-zero entry, and this entry has modulus one. Multiplication by these matrices preserves $|f|_{\infty}$ for every spinor function $f$. The four transformations of the generators that occur in the properties~\eqref{it:renormalisation-group-1-k3}, \eqref{it:renormalisation-group-1-k4}, \eqref{it:renormalisation-group-1-k5} and~\eqref{it:renormalisation-group-1-k7} of Definition~\ref{def:renormalisation-group-1} preserve the norms of Section~\ref{sec:grassmann-algebras-norms} and the bidegree.
\item[(c)] For every $c\in(0,1]$ each of the sets $R_k$, $c\cdot\tilde{R}_k$, $\tilde{R}^{-}_k$, $G_k$, $\tilde{G}_k$ and $\mathfrak{X}_k$ of Definition~\ref{def:domains} is invariant under the three maps $A\mapsto gA$, $A\mapsto-A$ and $A\mapsto\overline{A}$. The sets $c\cdot\tilde{R}_k$, $\tilde{R}^{-}_k$ and $\tilde{G}_k$ are, moreover, invariant under $A\mapsto A-d\boldsymbol{\lambda}$ for every real function $\boldsymbol{\lambda}$ on the sites, and under the addition of real constant bond functions.
\item[(d)] Let $g\in\mathbb{G}_k^L$, and put $g':=(S,a/L)$. The map $g\mapsto g'$ is an isomorphism of $\mathbb{G}_k^L$ onto the group $\mathbb{G}_{k+1}$ of the torus of side $l/L$. With the scalings of Section~\ref{sec:local-parts-adjustment},
\[
(g'A)_L=g\left(A_L\right),\qquad(g'\psi)_L=g\left(\psi_L\right),\qquad(g'\chi)_L=g\left(\chi_L\right),\qquad L\left(g'Y\right)=g\left(LY\right)
\]
for every bond function $A$ of the level $k+1$, for the generators $\psi$ and $\chi$ of the level $k+1$, and for every subset $Y$ of the torus of the level $k+1$, where $LY:=\{Ly:\ y\in Y\}$.
\item[(e)] \emph{Kernels of low degree.} Let $X\in D_k$ be small, and let $\mathscr{E}=\left\{\mathscr{E}(X';\mathscr{A},\mathscr{C}):\ X'\in D_k\right\}$ be a family of elements of $\mathrm{Gr}[\psi,\chi]$, with kernels written as in Definition~\ref{def:renormalisation-group-1}, such that for every $X'\in D_k$ the $4\times4$ matrix $\mathscr{E}_{20}(X';0,0;x,y)$ of Section~\ref{sec:local-parts-adjustment} vanishes unless both $x$ and $y$ lie in $X'$; this is the first clause of the property~\eqref{it:renormalisation-group-1-k2}. Let $\boldsymbol{\alpha}_0(\mathscr{E},X)$ and $\boldsymbol{\alpha}_2(\mathscr{E},X)$ be the extraction numbers of Section~\ref{sec:local-parts-adjustment}. If $\mathscr{E}$ has the property~\eqref{it:renormalisation-group-1-k3}, then
\[
\boldsymbol{\alpha}_0(\mathscr{E},gX)=\boldsymbol{\alpha}_0(\mathscr{E},X),\qquad\boldsymbol{\alpha}_2(\mathscr{E},gX)=S\cdot\boldsymbol{\alpha}_2(\mathscr{E},X)\cdot S^{-1}.
\]
If $\mathscr{E}$ has the property~\eqref{it:renormalisation-group-1-k5}, then
\[
\boldsymbol{\alpha}_2(\mathscr{E},X)=\mathsf{C}^{-1}\cdot\boldsymbol{\alpha}_2(\mathscr{E},X)^T\cdot\mathsf{C}.
\]
If $\mathscr{E}$ has the property~\eqref{it:renormalisation-group-1-k7}, then $\boldsymbol{\alpha}_0(\mathscr{E},X)$ is real and
\[
\boldsymbol{\alpha}_2(\mathscr{E},X)=\varkappa^{-1}\cdot\overline{\boldsymbol{\alpha}_2(\mathscr{E},X)}\cdot\varkappa.
\]
\item[(f)] \emph{Extraction.} Let $\mathscr{E}\in K^{x,C}_k$, and let $c_{\mathrm{fl}}\ge1$ be the constant of stage~\ref{it:order-choice-1}. The two sums over the small polymers that contain an $M$-cube $\square$, by which $\varepsilon(\mathscr{E})$ and $m(\mathscr{E})\cdot\mathbf{1}$ are defined in Section~\ref{sec:local-parts-adjustment}, do not depend on $\square$, and the second of them is a multiple of the unit matrix $\mathbf{1}$; thus $\varepsilon(\mathscr{E})$ and $m(\mathscr{E})$ are well defined complex numbers. The maps $\mathscr{E}\mapsto\varepsilon(\mathscr{E})$ and $\mathscr{E}\mapsto m(\mathscr{E})$ are complex linear, their values are real if $\mathscr{E}\in K^x_k$, and
\begin{gather*}
\varepsilon(\mathscr{E})=\varepsilon(\mathscr{E}|_0),\qquad m(\mathscr{E})=m(\mathscr{E}|_0),\\
|\varepsilon(\mathscr{E})|\le c_{\mathrm{fl}}\cdot\|\mathscr{E}\|^x_k,\qquad|m(\mathscr{E})|\le c_{\mathrm{fl}}\cdot e_k^{\frac{1}{2}}\cdot\|\mathscr{E}\|^x_k.
\end{gather*}
For every $(\mathscr{A},\mathscr{C})\in\tilde{R}_k\times\mathfrak{X}_k$ one has, in $\mathrm{Gr}[\psi,\chi]$,
\[
\sum_{X\in D_k}\mathscr{E}(X;\mathscr{A},\mathscr{C})=-\varepsilon(\mathscr{E})\cdot n_k-m(\mathscr{E})\cdot\left\langle\bar{\psi},\psi\right\rangle+\sum_{X\in D_k}(R\,\mathscr{E})(X;\mathscr{A},\mathscr{C}),
\]
where $n_k=l^3$ is the number of the sites of the unit lattice $T^0_k$, and $\left\langle\bar{\psi},\psi\right\rangle$ is the pairing of Section~\ref{sec:fermion-integrals-blocking}. The map $R$ of Section~\ref{sec:local-parts-adjustment} maps $K^{x,C}_k$ into itself and $K^x_k$ into itself, every $\mathscr{E}\in K^{x,C}_k$ satisfies $\|R\,\mathscr{E}\|^x_k\le c_{\mathrm{fl}}\cdot\|\mathscr{E}\|^x_k$, and $(R\,\mathscr{E})|_0=R\,(\mathscr{E}|_0)$.
\item[(g)] \emph{Reblocking.} Let $E'=\left\{E'(X;\mathscr{A},\mathscr{C}):\ X\in D_k,\ (\mathscr{A},\mathscr{C})\in\tilde{R}_k\times\mathfrak{X}_k\right\}$ be a family of elements of $\mathrm{Gr}[\psi,\chi]$ of the level $k$, and let $L\,E'$ be the reblocked family of the level $k+1$ (Section~\ref{sec:local-parts-adjustment}). If $E'$ has the property~\eqref{it:renormalisation-group-1-k3} for the elements of $\mathbb{G}_k^L$, then $L\,E'$ has the property~\eqref{it:renormalisation-group-1-k3} at the level $k+1$, for the group $\mathbb{G}_{k+1}$. If $E'$ has the property~\eqref{it:renormalisation-group-1-k4}, resp.~\eqref{it:renormalisation-group-1-k5}, \eqref{it:renormalisation-group-1-k7}, then $L\,E'$ has the same property at the level $k+1$.
\end{enumerate}
The same assertions hold without the second slot, that is, with $K^C_k$, $K_k$ and $\|\cdot\|_k$ in place of $K^{x,C}_k$, $K^x_k$ and $\|\cdot\|^x_k$.
\end{proposition}

This follows from the results of Section~\ref{sec:lattice-symmetries-local}.

\begin{proposition}[Adjustment]\label{prop:adjustment}
Assume the hypotheses that Theorem~\ref{thm:renormalisation-group-1} states for this proposition. Let $E\in K^{x,C}_k$, let $X\in D_k$ be small, let $(\mathscr{A},\mathscr{C})\in\tilde{R}_k\times\mathfrak{X}_k$, put $\mathscr{T}:=\mathscr{A}+\mathscr{C}$, and let $E^{\mathrm{nat}}$ be the adjustment~\eqref{eq:renormalisation-group-2} of $E$.
\begin{enumerate}
\item[(a)] \emph{Identity.} Under the substitution $\chi=\delta_{\alpha,\mathscr{T}}\psi$, $\bar{\chi}=\delta_{\alpha,-\mathscr{T}}\bar{\psi}$ the elements $E^{\mathrm{nat}}(X;\mathscr{A},\mathscr{C},\psi,\chi)$ and $E(X;\mathscr{A},\mathscr{C},\psi,\chi)$ become equal.
\item[(b)] \emph{Low degrees.} If $\int\mathrm{d}x\,\mathrm{d}y\,E_{20}(X;0,0;x,y)=0$, then $E^{\mathrm{nat}}_{20}(X;0,0;\cdot)=0$. For every $(n,m)\notin\{(2,0),(1,1),(0,2)\}$ the kernels satisfy $E^{\mathrm{nat}}_{nm}=E_{nm}$.
\item[(c)] \emph{Locality.} The family $E^{\mathrm{nat}}$ has the property~\eqref{it:renormalisation-group-1-k2}. Every leg $\chi(\boldsymbol{\zeta})$ of $E^{\mathrm{nat}}(X;\mathscr{A},\mathscr{C})$ that is not a leg of $E(X;\mathscr{A},\mathscr{C})$ is of kind~(I) for $X$, that is, the pair of sites of $\boldsymbol{\zeta}$ is admissible for $X$ of kind~(I) (Definition~\ref{def:pairs-sites-holder}).
\item[(d)] \emph{Bound.} One has
\[
\left\|E^{\mathrm{nat}}(X;\mathscr{A},\mathscr{C})\right\|_{h_k}\le10\cdot\left\|E(X;\mathscr{A},\mathscr{C})\right\|_{h_k}.
\]
\item[(e)] \emph{Symmetries.} The family $E^{\mathrm{nat}}$ belongs to $K^{x,C}_k$, and to $K^x_k$ if $E$ does. The map $E\mapsto E^{\mathrm{nat}}$ is linear, and $E^{\mathrm{nat}}|_0=(E|_0)^{\mathrm{nat}}$.
\end{enumerate}
Assertions~(a)--(d), and those of~(e) other than the identity $E^{\mathrm{nat}}|_0=(E|_0)^{\mathrm{nat}}$, hold also without the second slot, that is, for $E\in K^C_k$, a small $X\in D_k$ and $\mathscr{A}\in\tilde{R}_k$, with $\mathscr{T}:=\mathscr{A}$, with the adjustment~\eqref{eq:renormalisation-group-2} formed with the total connection $\mathscr{A}$ (Section~\ref{sec:local-parts-adjustment}), and with $E(X;\mathscr{A})$, $E^{\mathrm{nat}}(X;\mathscr{A})$, $E_{20}(X;0;x,y)$, $E^{\mathrm{nat}}_{20}(X;0;\cdot)$, $K^C_k$ and $K_k$ in place of $E(X;\mathscr{A},\mathscr{C})$, $E^{\mathrm{nat}}(X;\mathscr{A},\mathscr{C})$, $E_{20}(X;0,0;x,y)$, $E^{\mathrm{nat}}_{20}(X;0,0;\cdot)$, $K^{x,C}_k$ and $K^x_k$.
\end{proposition}

The assertion is contained in Theorem~\ref{thm:properties-adjustment}.

\begin{proposition}[Connections near a small polymer]\label{prop:connections-near-small}
Assume the hypotheses that Theorem~\ref{thm:renormalisation-group-1} states for this proposition. Let $X\in D_k$ be small, let $x_0$ be a site of $X$, let $(\mathscr{A},\mathscr{C})\in\tilde{R}_{k+1}\times\mathfrak{X}_{k+1}$, and put $\rho:=2LM$. Then the assertions (i)--(iii) hold for a suitable choice of a real function $\mathring{m}$ on the sites of $T_k$, of two complex bond functions $\hat{b}$ and $\mathscr{C}^b$ on $T_k$, and of two constant complex bond functions $c_1$ and $c_2$. The choice can be made so that, in addition, $\mathring{m}$, $\hat{b}$ and $c_1$ depend only on $\mathscr{A}$ and $x_0$, $\mathscr{C}^b$ and $c_2$ depend only on $\mathscr{C}$ and $x_0$, and $\mathscr{C}^b=0$ and $c_2=0$ if $\mathscr{C}=0$.
\begin{enumerate}
\item[(i)] On every bond whose two end points lie in the set $\left\{x:\ d(x,x_0)\le\rho\right\}$
\[
\mathscr{A}_L-d\mathring{m}=\hat{b}+c_1,\qquad\mathscr{C}_L=\mathscr{C}^b+c_2;
\]
every bond that has an end point at distance at most $1$ from $X$ is such a bond.
\item[(ii)] For every complex number $t$ with $|t|\le L^{\frac{7}{4}}$ the pair $\left(t\cdot\hat{b},t\cdot\mathscr{C}^b\right)$ belongs to $R_k\times\mathfrak{X}_k$, and
\[
|c_1|\le L^{-\frac{7}{8}+\frac{\varepsilon}{2}}\cdot e_k^{-\frac{3}{4}+\varepsilon},\qquad|c_2|\le L^{-1+\varepsilon}\cdot e_k^{-1+2\varepsilon}.
\]
\item[(iii)] Let $E$ be a family of the level $k$, written with kernels $E_{nm}$ as in Definition~\ref{def:renormalisation-group-1}, that has the properties~\eqref{it:renormalisation-group-1-k1}, \eqref{it:renormalisation-group-1-k2} and~\eqref{it:renormalisation-group-1-k4}. For all $n,m\ge0$ and all arguments $\xi=(\xi_1,\dots,\xi_n)$ and $\boldsymbol{\zeta}=(\boldsymbol{\zeta}_1,\dots,\boldsymbol{\zeta}_m)$
\[
\left|E_{nm}\left(X;\mathscr{A}_L,\mathscr{C}_L;\xi,\boldsymbol{\zeta}\right)\right|\le\exp\left(2(n+m)\right)\cdot\left|E_{nm}\left(X;\hat{b},\mathscr{C}^b;\xi,\boldsymbol{\zeta}\right)\right|,
\]
and $E_{00}\left(X;\mathscr{A}_L,\mathscr{C}_L\right)=E_{00}\left(X;\hat{b},\mathscr{C}^b\right)$. More precisely, the modulus on the left is the modulus on the right times one factor of size at most $\exp(1)$ for each argument and each slot. The comparison holds also without the second slot. Let $\mathscr{C}=0$, so that $\mathscr{C}^b=0$ and $c_2=0$ for the choice described before~(i), and let $E$ be a family of the level $k$ without the second slot, defined on $\tilde{R}_k$ and written with kernels $E_{nm}$ as in Definition~\ref{def:renormalisation-group-1} with the second slot suppressed, that has the properties~\eqref{it:renormalisation-group-1-k1}, \eqref{it:renormalisation-group-1-k2} and~\eqref{it:renormalisation-group-1-k4} in the form without the second slot in which they define $K^C_k$ (Definition~\ref{def:renormalisation-group-1}), the last with $\boldsymbol{\lambda}'=0$. Then for all $n,m\ge0$ and all arguments $\xi$ and $\boldsymbol{\zeta}$
\[
\left|E_{nm}\left(X;\mathscr{A}_L;\xi,\boldsymbol{\zeta}\right)\right|\le\exp\left(2(n+m)\right)\cdot\left|E_{nm}\left(X;\hat{b};\xi,\boldsymbol{\zeta}\right)\right|,
\]
and $E_{00}\left(X;\mathscr{A}_L\right)=E_{00}\left(X;\hat{b}\right)$; more precisely, the modulus on the left is the modulus on the right times one factor of size at most $\exp(1)$ for each argument.
\end{enumerate}
\end{proposition}

These assertions are contained in Theorem~\ref{thm:polymer-activities-3}.

\begin{proposition}[Contraction]\label{prop:contraction}
Assume the hypotheses that Theorem~\ref{thm:renormalisation-group-1} states for this proposition, let $L\,R\,\mathscr{E}=L\left((R\,\mathscr{E})^{\mathrm{nat}}\right)$ be as in Section~\ref{sec:local-parts-adjustment}, and let $c_{\mathrm{fl}}$ be the constant of stage~\ref{it:order-choice-1}. If $1-\varepsilon<\alpha<1$, then $L\,R$ maps $K^{x,C}_k$ into $K^{x,C}_{k+1}$ and $K^x_k$ into $K^x_{k+1}$, and
\[
\left\|L\,R\,\mathscr{E}\right\|^x_{k+1}\le c_{\mathrm{fl}}\cdot L^{-\frac{1}{4}+2\varepsilon}\cdot\|\mathscr{E}\|^x_k\qquad\left(\mathscr{E}\in K^{x,C}_k\right).
\]
Moreover $(L\,R\,\mathscr{E})|_0=L\,R\,(\mathscr{E}|_0)$ for every $\mathscr{E}\in K^{x,C}_k$. The same assertions hold without the second slot, that is, with $K^C_k$, $K_k$, $\|\cdot\|_k$, $K^C_{k+1}$, $K_{k+1}$, $\|\cdot\|_{k+1}$ in place of $K^{x,C}_k$, $K^x_k$, $\|\cdot\|^x_k$, $K^{x,C}_{k+1}$, $K^x_{k+1}$, $\|\cdot\|^x_{k+1}$.
\end{proposition}

The proposition is contained in Theorem~\ref{thm:contraction}.

\begin{remark}
The factor $c_{\mathrm{fl}}\cdot L^{-\frac{1}{4}+2\varepsilon}$ in the bound of Proposition~\ref{prop:contraction} may exceed $1$. Measured in the units of the windows $w^{\mathrm{E}}_k=e_k^{\frac{1}{4}-7\varepsilon}$ of Section~\ref{sec:running-quantities-tori}, however, the map $L\,R$ is a contraction. Indeed $e_{k+1}=L^{\frac{1}{2}}\cdot e_k$ gives $w^{\mathrm{E}}_{k+1}=L^{\frac{1}{8}-\frac{7\varepsilon}{2}}\cdot w^{\mathrm{E}}_k$, so that
\[
c_{\mathrm{fl}}\cdot L^{-\frac{1}{4}+2\varepsilon}\cdot\frac{w^{\mathrm{E}}_k}{w^{\mathrm{E}}_{k+1}}=c_{\mathrm{fl}}\cdot L^{-\frac{3}{8}+\frac{11\varepsilon}{2}}\le\frac{1}{16}
\]
by stage~\ref{it:order-choice-4}, since $\frac{11\varepsilon}{2}\le6\bar{\varepsilon}$; thus Proposition~\ref{prop:contraction} gives
\[
\frac{\left\|L\,R\,\mathscr{E}\right\|^x_{k+1}}{w^{\mathrm{E}}_{k+1}}\le\frac{1}{16}\cdot\frac{\|\mathscr{E}\|^x_k}{w^{\mathrm{E}}_k}.
\]
\end{remark}

\begin{proposition}[Bounds on the unit lattice]\label{prop:bounds-unit-lattice}
Assume the hypotheses that Theorem~\ref{thm:renormalisation-group-1} states for this proposition. Let $C_0$ be the constant of stage~\ref{it:order-choice-7}, let $h_k=(h_{k,1},h_{k,2})$ be the weights of Section~\ref{sec:grassmann-algebras-norms}, and put
\[
a_1:=e_k^{-\frac{1}{4}}=h_{k,1},\qquad a_2:=C_0^{-1}\cdot e_k^{-\frac{1}{4}+\varepsilon}=C_0^{-1}\cdot h_{k,2};
\]
then $a_2\le a_1$. The norms $|\cdot|_{a_1}$ and $|\cdot|_{a_2}$ on $\mathrm{Gr}[\Psi]$ are those of Section~\ref{sec:grassmann-algebras-norms}. Let $A\in\tilde{G}_k$.
\begin{enumerate}
\item[(a)] The element $\left\langle\bar{\Psi},D_k(A)\Psi\right\rangle$ of $\mathrm{Gr}[\Psi]$ is an analytic function of $A$ on $\tilde{G}_k$, and
\[
\left|\left\langle\bar{\Psi},D_k(A)\Psi\right\rangle\right|_{a_1}\le C_0\cdot e_k^{-\frac{1}{2}}\cdot l^3.
\]
\item[(b)] At $\psi=H_k(A)\Psi$ and $\bar{\psi}=\bar{H}_k(A)\bar{\Psi}$,
\[
\left|\left\langle\bar{\psi},\psi\right\rangle\right|_{a_2}\le4e_k^{-\frac{1}{2}+2\varepsilon}\cdot l^3\le4e_k^{-\frac{1}{2}}\cdot l^3.
\]
\item[(c)] For every $\mathscr{E}\in K^{x,C}_k$, every $X\in D_k$ and every $(\mathscr{A},\mathscr{C})\in\tilde{R}_k\times\mathfrak{X}_k$, with $\mathscr{T}:=\mathscr{A}+\mathscr{C}$ and the fields $H_k^{\#}(\mathscr{T})\Psi$ of Section~\ref{sec:fermion-integrals-blocking} inserted for the generators,
\[
\left|\mathscr{E}\left(X;\mathscr{A},\mathscr{C},H_k^{\#}(\mathscr{T})\Psi\right)\right|_{a_2}\le\left\|\mathscr{E}(X;\mathscr{A},\mathscr{C})\right\|_{h_k}.
\]
\end{enumerate}
\end{proposition}

This is a special case of Theorem~\ref{thm:unit-lattice-bounds}.

\begin{proposition}[Block covariance at real connections]\label{prop:block-covariance-real}
Assume the hypotheses that Theorem~\ref{thm:renormalisation-group-1} states for this proposition, and let $\mathfrak{m}$ be the mass of Section~\ref{sec:exact-fermion-operators}. The spinor functions on $T_k$ and those on $T^0_k$ carry the scalar products of Section~\ref{sec:block-covariance-real}, with the weight $\eta_k^3$ at each site of $T_k$ and with the counting measure on $T^0_k$:
\[
(f,g)=\eta_k^3\sum_{x\in T_k}\sum_{\boldsymbol{\beta}}\overline{f_{\boldsymbol{\beta}}(x)}\cdot g_{\boldsymbol{\beta}}(x),\qquad(v,w)=\sum_{y\in T^0_k}\sum_{\boldsymbol{\beta}}\overline{v_{\boldsymbol{\beta}}(y)}\cdot w_{\boldsymbol{\beta}}(y).
\]
Adjoints and operator norms refer to these scalar products; for an operator $O$ and a real number $c$, $\operatorname{Re}O\ge c$ means that $\operatorname{Re}(f,O\,f)\ge c\cdot(f,f)$ for every $f$; and a $4\times4$ block is normed as an operator on $\mathbb{C}^4$ with the Euclidean norm.
Let $A$ be a real bond function on $T_k$. Then $\operatorname{Re}\left(D_A+\mathfrak{m}\right)\ge\mathfrak{m}$, and $D_k(A)$ is invertible, with
\[
D_k(A)^{-1}=b_k^{-1}+Q_k(A)\cdot\left(D_A+\mathfrak{m}\right)^{-1}\cdot Q_k(A)^{*},\qquad Q_k(A)^{*}=Q_k^T(-A),
\]
where $Q_k(A)^{*}$ is the adjoint of $Q_k(A)$. The operator norm of $D_k(A)^{-1}$ is at most $b_k^{-1}+\mathfrak{m}^{-1}$, and for all sites $y$ and $y'$ of $T^0_k$ the $4\times4$ blocks of $D_k(A)^{-1}$ satisfy
\[
\left|D_k(A)^{-1}(y,y')\right|\le b_k^{-1}\cdot\mathbf{1}[y=y']+\frac{2\exp(1)}{\mathfrak{m}}\cdot\exp\left(-\frac{\mathfrak{m}}{12\exp(1)}\cdot d(y,y')\right).
\]
\end{proposition}

This is a special case of Theorem~\ref{thm:block-covariance}; the condition $\mathfrak{m}\le12\exp(1)$ of that theorem holds, since $\mathfrak{m}\le\mathfrak{m}_0\le1$.

\begin{proposition}[The free determinant]\label{prop:free-determinant}
Assume the hypotheses that Theorem~\ref{thm:renormalisation-group-1} states for this proposition; the functions $\mathfrak{P}_k$ of Section~\ref{sec:flow} and $E^{\mathrm{det}}_j$ of Definition~\ref{def:determinant-functions} are formed with the families~\eqref{it:primitives-p5} of the levels $j<k$ that occur in these hypotheses. Let $a\in C^1\otimes\mathbb{C}$ be an external connection (Definition~\ref{def:twist-external-connection}) such that its expressions $a^{[j]}$ in the coordinates of level $j$ (Section~\ref{sec:lattice-measures-laws}) satisfy $a^{[j]}\in\mathfrak{X}_j$ for $0\le j\le k$, and let $D_a$ and $D_0$ be the operators $D_A$ of Section~\ref{sec:lattice-measures-laws} at $A=a$ and at $A=0$. Then
\[
\frac{\det\left(D_a+\bar{m}\cdot\eta\right)}{\det\left(D_0+\bar{m}\cdot\eta\right)}=\exp\left\{\sum_{X\in D_k}\mathfrak{P}_k\left(X,a^{[k]}\right)\right\}\cdot\frac{\det D_k\left(a^{[k]}\right)}{\det D_k(0)}
\]
and
\[
\sum_{X\in D_k}\mathfrak{P}_k\left(X,a^{[k]}\right)=\sum_{j<k}\ \sum_{X\in D_j}E^{\mathrm{det}}_j\left(X,a^{[j]}\right).
\]
\end{proposition}

This is proved in Section~\ref{sec:determinant-free-operator}.

\begin{proposition}[Locality]\label{prop:locality}
Assume the hypotheses that Theorem~\ref{thm:renormalisation-group-1} states for this proposition. Let $Z$ be a union of $M$-cubes of $T_k$, and let a Grassmann algebra over $\mathbb{C}$ of the kind of Proposition~\ref{prop:weighted-coefficient-norms} be fixed, with the following independent generators: the $\psi(\xi)$ and the $\chi(\boldsymbol{\zeta})$ of Section~\ref{sec:grassmann-algebras-norms}, and finitely many further families of independent generators, possibly none, each of them indexed by a finite set; for instance, the further generators may be the $\Psi(u)$ and the $W(u)$ of Section~\ref{sec:grassmann-algebras-norms}, or those of one of these two families. This algebra is a complex vector space of finite dimension, since it has finitely many generators, and the limits in~(a) are pointwise limits in this space. If the generators $W(u)$ are among the further generators, then the integral $\int\mathrm{d}\mu_I(W)$ of Proposition~\ref{prop:weighted-coefficient-norms}, which occurs in~(a), maps this algebra into the algebra of the same kind formed without them. A map $F$ from a set of pairs $(\mathscr{A},\mathscr{C})$ of connections to this algebra is called \emph{local in} $Z$ if, of the generators $\psi(\xi)$ and $\chi(\boldsymbol{\zeta})$, its values contain only the $\psi(\xi)$ with $\xi=(x,\boldsymbol{\beta},\boldsymbol{\omega})$ and $x\in Z$ and the $\chi(\boldsymbol{\zeta})$ with $\boldsymbol{\zeta}=(x,x',\boldsymbol{\beta},\boldsymbol{\omega})$ and $(x,x')$ admissible for $Z$ (Definition~\ref{def:pairs-sites-holder}), the further generators being subject to no restriction, and if $F(\mathscr{A},\mathscr{C})=F(\mathscr{A}',\mathscr{C}')$ whenever $(\mathscr{A},\mathscr{C})$ and $(\mathscr{A}',\mathscr{C}')$ both belong to the set on which $F$ is defined and $\mathscr{A}(b)=\mathscr{A}'(b)$, $\mathscr{C}(b)=\mathscr{C}'(b)$ for every bond $b$ of $T_k$ that meets $Z$, that is, that has an end point in $Z$; this is the notion of Definition~\ref{def:local-maps}. In the terms of Definition~\ref{def:terminology}(1), with the pair $(\mathscr{A},\mathscr{C})$ as the object of that clause, the second condition states that $F$ reads this pair on the bonds that meet $Z$. For a family $\mathscr{E}$ as in Definition~\ref{def:renormalisation-group-1}, the property~\eqref{it:renormalisation-group-1-k2} states that, for every $X\in D_k$, the map $(\mathscr{A},\mathscr{C})\mapsto\mathscr{E}(X;\mathscr{A},\mathscr{C})$ on $\tilde{R}_k\times\mathfrak{X}_k$ is local in $X$.
\begin{enumerate}
\item[(a)] The maps local in $Z$, defined on a fixed set of pairs of connections, form an algebra under the pointwise operations. This algebra is closed under limits, under the exponential, under norm-convergent series, under derivatives and integrals with respect to auxiliary parameters, and under the integral $\int\mathrm{d}\mu_I(W)$ of Proposition~\ref{prop:weighted-coefficient-norms}.
\item[(b)] Let $Z_1,\dots,Z_n$ be unions of $M$-cubes of $T_k$. If $F_i$ is local in $Z_i$ for $i=1,\dots,n$, then the product $F_1\cdots F_n$ is local in $Z_1\cup\dots\cup Z_n$.
\item[(c)] If a family $\mathscr{E}$ has the property~\eqref{it:renormalisation-group-1-k2}, then so have $R\,\mathscr{E}$ and $\mathscr{E}^{\mathrm{nat}}$, and $L\,\mathscr{E}$ has the property~\eqref{it:renormalisation-group-1-k2} at the level $k+1$.
\end{enumerate}
\end{proposition}

Assertions (a)--(c) are contained in Theorem~\ref{thm:closure-properties-local}.

\subsubsection{The range of the masses}\label{sec:range-masses}

\begin{proposition}[Bounds on the masses]\label{prop:bounds-masses}
Let $e>0$ satisfy
\[
v_K\ge1\qquad\text{and}\qquad\left(1+\frac{\mathrm{lg}}{v_K}\right)^r\le L,
\]
which is the hypothesis of Theorem~\ref{thm:scales-stop-levels}(a). Let $R\in\mathrm{Tor}(e)$ and $N\ge N_{\mathrm{min}}(R,e)$, and let $k^{*}=k^{*}(R)$ be the stop level of Definition~\ref{def:tori-stop-level}, which exists by Theorem~\ref{thm:scales-stop-levels}(a). For $0\le k\le k^{*}$ let $\bar{m}\cdot a_k$ be the mass of level $k$ (Section~\ref{sec:running-quantities-tori}); it is denoted by $\mathfrak{m}_k$ for $k\ge1$ and by $\bar{m}\cdot\eta$ for $k=0$, the symbol $\mathfrak{m}_0$ being reserved for the upper end of the mass range (Definition~\ref{def:notation}). Then, for $0\le k\le k^{*}$,
\begin{gather*}
0<\bar{m}\cdot a_k\le\mathfrak{m}_{*}\le\mathfrak{m}_K\le L^{-m_0}\le\mathfrak{m}_0,\qquad\mathfrak{m}_K>L^{-m_0-1},\\
\mathfrak{m}_{*}\cdot l_{*}=\bar{m}\cdot R\ge\bar{m}\cdot h'>L^{-m_0-1}.
\end{gather*}
If moreover $\left(1+\frac{3\mathrm{lg}}{v_K}\right)^r<L$, then for $R=R^{*}$
\[
L^{-m_0-3}<\mathfrak{m}_{*}\le L^{-m_0-1}.
\]
These assertions hold without any relation of the form $\bar{m}\cdot R\le1$.
\end{proposition}

These bounds are contained in Proposition~\ref{prop:bounds-running-masses}.

The number $\mathfrak{m}_0$ is the constant of Propositions~\ref{prop:decay-fine-propagator}--\ref{prop:bounds}: it satisfies $0<\mathfrak{m}_0\le1$ and depends only on $b$, $\alpha$ and $L$, and Theorem~\ref{thm:fermion-operators-weakened} gives $\mathfrak{m}_0=1$. The integer $m_0$ is the least integer $m\ge0$ with $L^{-m}\le\mathfrak{m}_0$ (Definition~\ref{def:mass-bound-stop}). Hence $m_0=0$ and $h'=h\cdot L^{-m_0}=h$; nevertheless the formulas are written with the symbols $\mathfrak{m}_0$, $m_0$ and $h'$.

\subsection{One step, the operators, the flow and the jets}\label{sec:one-step-operators}

\paragraph{Standing hypotheses.}
The data of a single step and the data of the flow are fixed as follows.

\begin{definition}[Hypothesis for one step]\label{def:hypothesis-one-step}
The constants of stages~\ref{it:order-choice-1}--\ref{it:order-choice-9} are chosen as prescribed in Section~\ref{sec:order-choice}, and they satisfy conditions~\eqref{eq:renormalisation-group-A0}--\eqref{eq:renormalisation-group-A10}, \eqref{eq:renormalisation-group-19}, \eqref{eq:renormalisation-group-20}--\eqref{eq:renormalisation-group-24} and~\eqref{eq:renormalisation-group-25} of Section~\ref{sec:conditions}. Moreover the following data are given.
\begin{itemize}
\item An integer $k\ge 0$.
\item A number $e_k>0$ such that
\[
e_{k+1}:=L^{\frac{1}{2}}\cdot e_k\le e_0.
\]
The numbers $v_k$, $r_k$, $p_k$, $p_{0,k}$, $w^{\mathrm{m}}_k$ and $w^{\mathrm{E}}_k$ are formed with this $e_k$ by the formulas of Section~\ref{sec:running-quantities-tori}.
\item A side $l$ such that $l\in\mathrm{Cl}$, $l/L\in\mathrm{Cl}$ and $l\ge M\cdot r_k$.
\item A mass $\mathfrak{m}\in\left(0,\mathfrak{m}_0/L\right]$.
\item A system of primitives of level $k$ on $T_k$ in the sense of Section~\ref{sec:primitive-operators}, which has the properties~\ref{it:renormalisation-group-1}--\ref{it:renormalisation-group-17} of Section~\ref{sec:properties-systems-primitives} for the parameters $\mathrm{par}$.
\end{itemize}
The objects of level $k+1$ are formed with the side $l/L$, the mass $L\cdot\mathfrak{m}$ and the number $e_{k+1}$.
\end{definition}

\begin{definition}[Hypothesis for the flow]\label{def:hypothesis-flow}
The constants of stages~\ref{it:order-choice-1}--\ref{it:order-choice-10} are chosen as prescribed in Section~\ref{sec:order-choice}; they satisfy the conditions named in Definition~\ref{def:hypothesis-one-step} and, in addition, condition~\eqref{eq:renormalisation-group-26}. Moreover $0<\bar{m}\le 1$; the charge $e>0$ obeys
\[
e_K=e\cdot h'^{\frac{1}{2}}\le e_1;
\]
$R\in\mathrm{Tor}(e)$ (Definition~\ref{def:tori-stop-level}); $N\ge N_{\mathrm{min}}(R,e)$; and the systems of primitives are those of Definition~\ref{def:periodisation-projection}.
\end{definition}

Under hypothesis~\ref{def:hypothesis-flow}, conditions~\eqref{eq:renormalisation-group-16}, \eqref{eq:renormalisation-group-fl} and~\eqref{eq:renormalisation-group-T1}--\eqref{eq:renormalisation-group-T6} of Section~\ref{sec:conditions} are fulfilled, by Theorem~\ref{thm:scales-stop-levels}.

\subsubsection{One step without large fields}\label{sec:one-step-large}

\begin{theorem}[A step on small fields with an external connection]\label{thm:step-small-fields}
Assume hypothesis~\ref{def:hypothesis-one-step}, and let the step map of level $k$ be formed as in Section~\ref{sec:renormalisation-step} with the system of primitives that this hypothesis provides.
\begin{enumerate}
\item[(a)]\itemmark{a}\label{it:step-small-fields-a} \emph{Convergence, bound and analyticity.} Let $(m,\mathscr{E})\in\mathrm{Ball}^x_k(4)$; for $k=0$ the ball contains only pairs with $\mathscr{E}=0$. Let $A\in\tilde{R}^{-}_k$ and $\mathscr{C}\in\mathfrak{X}_k$. Then the series in~\eqref{eq:renormalisation-group-61} of Section~\ref{sec:fluctuation-integral-new} converge in norm; their convergence is absolute with respect to the norm $\|\cdot\|_{h_k/4}$ of Section~\ref{sec:grassmann-algebras-norms} (see Section~\ref{sec:activity-before-reblocking}). The family $E^{*,x}_k(m,\mathscr{E})$ is an element of $K^{x,C}_{k+1}$, and
\[
\left\|E^{*,x}_k(m,\mathscr{E})\right\|^x_{k+1}\le e_k^{\frac{1}{4}-6\varepsilon}.
\]
If $m$ is real and $\mathscr{E}\in K^x_k$, then $E^{*,x}_k(m,\mathscr{E})\in K^x_{k+1}$. The map $(m,\mathscr{E})\mapsto E^{*,x}_k(m,\mathscr{E})$ is analytic along complex lines, in the following sense. Let $(m_a,\mathscr{E}_a)$ and $(m_b,\mathscr{E}_b)$ be elements of $\mathbb{C}\times K^{x,C}_k$. The set of those $t\in\mathbb{C}$ for which $\left(m_a+t\cdot m_b,\,\mathscr{E}_a+t\cdot\mathscr{E}_b\right)$ lies in $\mathrm{Ball}^x_k(4)$ is closed; on this set the function
\[
t\longmapsto E^{*,x}_k\left(m_a+t\cdot m_b,\,\mathscr{E}_a+t\cdot\mathscr{E}_b\right)
\]
is continuous, and in its interior the function is analytic. Finally, the number $d\varepsilon^0_k$ of Section~\ref{sec:averaging-fluctuations-gauge} obeys $0\le d\varepsilon^0_k\le e_k^7$.

\item[(b)]\itemmark{b}\label{it:step-small-fields-b} \emph{The step at vanishing second slot.} Let $(m,\mathscr{E})\in\mathrm{Ball}^x_k(4)$. Then $(m,\mathscr{E}|_0)\in\mathrm{Ball}_k(4)$, as $\left\|\mathscr{E}|_0\right\|_k\le\|\mathscr{E}\|^x_k$, and
\[
E^{*,x}_k(m,\mathscr{E})|_0=E^{*}_k\left(m,\mathscr{E}|_0\right).
\]
The map $E^{*}_k$ carries $\mathrm{Ball}_k(4)$ into $K^C_{k+1}$, and it carries the pairs $(m,E)\in\mathrm{Ball}_k(4)$ with real $m$ and with $E\in K_k$ into $K_{k+1}$. For every $(m,E)\in\mathrm{Ball}_k(4)$,
\[
\left\|E^{*}_k(m,E)\right\|_{k+1}\le e_k^{\frac{1}{4}-6\varepsilon}.
\]
For $(m_a,E_a)$ and $(m_b,E_b)$ in $\mathbb{C}\times K^C_k$, the function $t\mapsto E^{*}_k\left(m_a+t\cdot m_b,\,E_a+t\cdot E_b\right)$ is continuous on the closed set of those $t\in\mathbb{C}$ for which $\left(m_a+t\cdot m_b,\,E_a+t\cdot E_b\right)\in\mathrm{Ball}_k(4)$, and analytic in the interior of that set.

\item[(c)]\itemmark{c}\label{it:step-small-fields-c} \emph{The identity of one step.} Let $\boldsymbol{\varepsilon}$ and $m$ be real numbers with $|m|\le w^{\mathrm{m}}_k$, and let $\mathscr{E}\in K^x_k$ with $\|\mathscr{E}\|^x_k\le w^{\mathrm{E}}_k$, where $\mathscr{E}=0$ if $k=0$; thus $(m,\mathscr{E})\in\mathrm{Ball}^x_k(1)$. Let $\mathfrak{P}_k=\left\{\mathfrak{P}_k(X,\cdot\,):\ X\in D_k\right\}$ be an arbitrary family of functions $\mathfrak{P}_k(X,\cdot\,)$ from $\mathfrak{X}_k$ to $\mathbb{C}$. Define
\begin{gather*}
\boldsymbol{\varepsilon}^{+}:=L^3\cdot\left(\boldsymbol{\varepsilon}+\varepsilon(\mathscr{E})+d\varepsilon^0_k\right),\qquad m^{+}:=L\cdot\left(m+m(\mathscr{E})\right),\\
\mathscr{E}^{+}:=L R\,\mathscr{E}+L\,\delta E^{\mathrm{det}}_k+E^{*,x}_k(m,\mathscr{E}),
\end{gather*}
and let $\mathfrak{P}_{k+1}=\left\{\mathfrak{P}_{k+1}(Y,\cdot\,):\ Y\in D_{k+1}\right\}$ be the family of functions on $\mathfrak{X}_{k+1}$ given by
\[
\mathfrak{P}_{k+1}(Y,\mathscr{C}'):=\sum_{X\in D_k:\ \bar{X}^L=LY}\left(\mathfrak{P}_k\left(X,\mathscr{C}'_L\right)+E^{\mathrm{det}}_k\left(X,\mathscr{C}'_L\right)\right)\qquad\left(\mathscr{C}'\in\mathfrak{X}_{k+1}\right).
\]
Let the density $\mathscr{F}^x_k$ of Section~\ref{sec:fermion-integrals-blocking} be formed with the family $\mathfrak{P}_k$, and the density $\mathscr{F}^x_{k+1}$ of level $k+1$ with the family $\mathfrak{P}_{k+1}$. Then for every $A'\in\mathrm{Dom}_{k+1}$, with an arbitrary harmonic part, and for every $\mathscr{C}'\in\mathfrak{X}_{k+1}$ one has
\begin{align}
\left(\mathbf{T}^x_k\,\mathscr{F}^x_k[\boldsymbol{\varepsilon},m,\mathscr{E}]\right)(A',\mathscr{C}';\,\cdot\,)&=\mathscr{F}^x_{k+1}[\boldsymbol{\varepsilon}^{+},m^{+},\mathscr{E}^{+}](A',\mathscr{C}';\,\cdot\,),\label{eq:renormalisation-group-10}\\
y^x_{k+1}[\boldsymbol{\varepsilon}^{+},m^{+},\mathscr{E}^{+}](A';\mathscr{C}')&=\int_{\mathrm{Box}_k}\mathrm{d}\mu_I(\zeta)\,y^x_k[\boldsymbol{\varepsilon},m,\mathscr{E}]\left(A^{\min}_k(A')+\mathrm{Sq}_k\zeta;\,\mathscr{C}'_L\right).\label{eq:renormalisation-group-11}
\end{align}
Suppose in addition that $\mathfrak{P}_k(X,0)=0$ for every $X\in D_k$. Then $\mathfrak{P}_{k+1}(Y,0)=0$ for every $Y\in D_{k+1}$, because $E^{\mathrm{det}}_k(X,0)=0$, and at $\mathscr{C}'=0$ the identities~\eqref{eq:renormalisation-group-10} and~\eqref{eq:renormalisation-group-11} are the identities for $\mathscr{F}_k$, $y_k$ and $\mathbf{T}_k$ with the data $(\boldsymbol{\varepsilon},m,\mathscr{E}|_0)$ and their image $(\boldsymbol{\varepsilon}^{+},m^{+},\mathscr{E}^{+}|_0)$.

\item[(d)]\itemmark{d}\label{it:step-small-fields-d} \emph{The data of level $k+1$.} Under the hypotheses of~(c), and with its notation, $\mathscr{E}^{+}$ is an element of $K^x_{k+1}$ with
\[
\mathscr{E}^{+}|_0=L R\left(\mathscr{E}|_0\right)+L\,E^{\mathrm{det}}_k+E^{*}_k\left(m,\mathscr{E}|_0\right),
\]
and the following inequalities hold:
\begin{gather*}
\left\|\mathscr{E}^{+}\right\|^x_{k+1}\le c_{\mathrm{fl}}\cdot L^{-\frac{1}{4}+2\varepsilon}\cdot\|\mathscr{E}\|^x_k+e_{k+1}^{\frac{1}{4}-2\varepsilon}+e_k^{\frac{1}{4}-6\varepsilon},\\
\left|m^{+}-L\cdot m\right|\le L\cdot c_{\mathrm{fl}}\cdot e_k^{\frac{1}{2}}\cdot\|\mathscr{E}\|^x_k,\qquad\left|\boldsymbol{\varepsilon}^{+}-L^3\cdot\boldsymbol{\varepsilon}\right|\le L^3\cdot\left(c_{\mathrm{fl}}\cdot\|\mathscr{E}\|^x_k+e_k^7\right).
\end{gather*}

\item[(e)]\itemmark{e}\label{it:step-small-fields-e} \emph{Lipschitz bound in the data.} For $(m_1,\mathscr{E}_1)$ and $(m_2,\mathscr{E}_2)$ in $\mathrm{Ball}^x_k(1)$,
\[
\left\|E^{*,x}_k(m_1,\mathscr{E}_1)-E^{*,x}_k(m_2,\mathscr{E}_2)\right\|^x_{k+1}\le e_k^{\frac{1}{4}-6\varepsilon}\cdot\left(\frac{|m_1-m_2|}{w^{\mathrm{m}}_k}+\frac{\|\mathscr{E}_1-\mathscr{E}_2\|^x_k}{w^{\mathrm{E}}_k}\right).
\]
For $(m_1,E_1)$ and $(m_2,E_2)$ in $\mathrm{Ball}_k(1)$ the analogous bound holds without the second slot:
\[
\left\|E^{*}_k(m_1,E_1)-E^{*}_k(m_2,E_2)\right\|_{k+1}\le e_k^{\frac{1}{4}-6\varepsilon}\cdot\left(\frac{|m_1-m_2|}{w^{\mathrm{m}}_k}+\frac{\|E_1-E_2\|_k}{w^{\mathrm{E}}_k}\right).
\]

\item[(f)]\itemmark{f}\label{it:step-small-fields-f} \emph{Symmetries.} Let $(m,\mathscr{E})\in\mathrm{Ball}^x_k(4)$. The family $E^{*,x}_k(m,\mathscr{E})$ has property~\eqref{it:renormalisation-group-1-k3} at level $k+1$, that is, with respect to the group $\mathbb{G}_{k+1}$ of Section~\ref{sec:dirac-matrices-lattice}; it has property~\eqref{it:renormalisation-group-1-k4} with respect to both connections, and property~\eqref{it:renormalisation-group-1-k5}; and in each of its monomials the barred legs are as many as the unbarred legs. If moreover $m$ is real and $\mathscr{E}\in K^x_k$, the family has property~\eqref{it:renormalisation-group-1-k7}.

\item[(g)]\itemmark{g}\label{it:step-small-fields-g} \emph{Locality and analyticity.} Let $(m,\mathscr{E})\in\mathrm{Ball}^x_k(4)$. The family $E^{*,x}_k(m,\mathscr{E})$ has properties~\eqref{it:renormalisation-group-1-k1} and~\eqref{it:renormalisation-group-1-k2} at level $k+1$. Thus, for every $Y\in D_{k+1}$, the kernels of $E^{*,x}_k(m,\mathscr{E})(Y;\cdot\,)$ are analytic on $\tilde{R}_{k+1}\times\mathfrak{X}_{k+1}$, they vanish unless each pair among their arguments is of kind~\eqref{it:pairs-sites-holder-I} or~\eqref{it:pairs-sites-holder-II} for $Y$ (Definition~\ref{def:pairs-sites-holder}), and they read the two connections on the bonds that meet $Y$. Furthermore, let $\mathbf{Y}\in D_k$. In the terminology of Definition~\ref{def:terminology}, the following assertions hold (Proposition~\ref{prop:local-constituents-step}).
\begin{itemize}
\item The objects $E^{(i),x}(\mathbf{Y})$, $E^{(\le 5),x}(\mathbf{Y})$, $\mathsf{K}(\mathbf{Y})$, $\mathsf{H}(\mathbf{Y})$ and $E^{\#,x}_k(\mathbf{Y})$ of Section~\ref{sec:definition-step-map} are sums over polymers and over parameter cubes that are contained in $\mathbf{Y}$.
\item The exact operators in the displays that define these objects are given by one formula on $\mathbb{L}_k$.
\item The primitives enter these objects only through their values at parameter families that vanish outside the parameter cubes contained in $\mathbf{Y}$, and only through their rows at indices contained in $\mathbf{Y}^{\sim 1}$.
\item These objects read the connections on the bonds that meet $\mathbf{Y}$.
\item Every integration in these displays is either an integration over auxiliary parameters or an integration against a product measure all of whose factors are copies of one measure.
\end{itemize}

\item[(h)]\itemmark{h}\label{it:step-small-fields-h} \emph{Uniformity.} Every condition on the constants on which (a)--(g) rest is a condition of Section~\ref{sec:order-choice} or of Section~\ref{sec:conditions}. No constant and none of these conditions depends on $k$, on $l$, on $\mathfrak{m}$, on the harmonic part of $A'$, on $\mathscr{C}'$, on the system of primitives other than through $\mathrm{par}$, or on the data $(m,\mathscr{E})$.
\end{enumerate}
\end{theorem}

This is established in Section~\ref{sec:renormalisation-group-step}.

\subsubsection{Fermion operators}\label{sec:fermion-operators}

\begin{theorem}[Fermion operators and their weakened families]\label{thm:fermion-operators-weakened}
Let $b>0$, and let $\alpha\in(1-\bar{\varepsilon},1)$, $\kappa$ and $L$ be as in Section~\ref{sec:order-choice}. There exist numbers $\mathfrak{m}_0\in(0,1]$, $\alpha^F\in\left(0,\frac{1}{4}\right]$ and $\gamma^F\in(0,1]$, functions of $(b,\alpha,L)$ alone, and numbers $M^F$, $C^F$ and $e^F_0$ that depend on further arguments in the following way. The number $M^F(b,\alpha,L,\alpha_0,\kappa')$ is defined for $\alpha_0\in\left(0,\alpha^F\right]$ and $\kappa'>0$. The number $C^F(b,\alpha,L,M,\alpha_0,\kappa')$ is defined for such $\alpha_0$ and $\kappa'$ and for every power $M=L^m$, $m\in\mathbb{N}$, with $M\ge M^F(b,\alpha,L,\alpha_0,\kappa')$. The number $e^F_0(b,\alpha,L,M,\alpha_0,\kappa',\varepsilon)>0$ is defined for such $\alpha_0$, $\kappa'$ and $M$ and for every $\varepsilon\in(0,\bar{\varepsilon}]$. These numbers have the following property. Let $\alpha_0$, $\kappa'$, $M$ and $\varepsilon$ be as just described, and write $C^F$ and $e^F_0$ for $C^F(b,\alpha,L,M,\alpha_0,\kappa')$ and $e^F_0(b,\alpha,L,M,\alpha_0,\kappa',\varepsilon)$. Let $k\ge0$ be an integer, let $l\in\mathrm{Cl}$, $\mathfrak{m}\in(0,\mathfrak{m}_0]$ and $e_k\in\left(0,e^F_0\right]$, and let the lattices $T_k$ and $T^0_k$, the operators of Sections~\ref{sec:averaging-fluctuations-gauge} and~\ref{sec:exact-fermion-operators} and the domains of Section~\ref{sec:domains-bond-functions} be formed with $(k,l,\mathfrak{m},e_k)$. Then the following assertions hold.
\begin{enumerate}
\item[(i)] Propositions~\ref{prop:decay-fine-propagator}--\ref{prop:symmetries-fermion-operators} and~\ref{prop:coverings-tori} hold with $(C_F,\gamma_F)=(C^F,\gamma^F)$.
\item[(ii)] There is a system of families~\eqref{it:primitives-p3}, \eqref{it:primitives-p4} and~\eqref{it:primitives-p5} of level $k$ on $T_k$ with the properties~\ref{it:renormalisation-group-1}, \ref{it:renormalisation-group-2} and~\ref{it:renormalisation-group-5}--\ref{it:renormalisation-group-17} for the parameter set $(\alpha_0,\gamma^F,\kappa';C^F)$. For $k=0$ this system consists of families $\Gamma_0(s'',A)$ and $S_{0,y}(X,A)$, the family $H_0(s'',A)$ of property~\ref{it:renormalisation-group-17} taking the place of the family~\eqref{it:primitives-p3}.
\item[(iii)] Assume $\kappa'\ge\kappa_w$, where $\kappa_w=\kappa+\kappa_0+3$ is the number of Section~\ref{sec:primitive-operators}. Then the family~\eqref{it:primitives-p5} of the system of~(ii) has the properties~\ref{it:renormalisation-group-1}, \ref{it:renormalisation-group-2}, \ref{it:renormalisation-group-7}, \ref{it:renormalisation-group-12}, \ref{it:renormalisation-group-13} and~\ref{it:renormalisation-group-14} for the parameter set $(\alpha_0,\gamma^F,\kappa_w;C^F)$, property~\ref{it:renormalisation-group-7} being taken for $k=0$ in the form stated in property~\ref{it:renormalisation-group-17}, because among these properties only property~\ref{it:renormalisation-group-7} contains numbers of the parameter set, and its bound with the rate $\kappa'$ implies the bound with the rate $\kappa_w\le\kappa'$, as $d_M(X)\ge0$. Let $\acute{E}_k$, $E^d_k$, $E^{\mathrm{det}}_k$ and $\delta E^{\mathrm{det}}_k$ be formed with this family. For these functions the assertions~\eqref{it:local-expansion-determinant-d0}--\eqref{it:local-expansion-determinant-d6} of Proposition~\ref{prop:local-expansion-determinant} hold when $\mathrm{par}$ is replaced by $(\alpha_0,\gamma^F,\kappa_w;C^F)$, the constant $C_d$ by the number~\eqref{eq:operators-fermion-10} formed with $C_w=C^F$, and $e_0$ by $e^F_0$; thus $e_k$ ranges over $\left(0,e^F_0\right]$, and the hypothesis of~\eqref{it:local-expansion-determinant-d5} and~\eqref{it:local-expansion-determinant-d6} on the charges is $e_{k+1}=L^{\frac{1}{2}}\cdot e_k\le e^F_0$. Here the inequalities $l\ge M\cdot r_k$, $v_k\ge9$ and $8M\cdot\left(\frac{v_k}{4}+4\right)\cdot e_k^{\varepsilon}\le1$ are hypotheses of~\eqref{it:local-expansion-determinant-d4}, \eqref{it:local-expansion-determinant-d5} and~\eqref{it:local-expansion-determinant-d6}, the condition $l/L\in\mathrm{Cl}$ is a further hypothesis of~\eqref{it:local-expansion-determinant-d5} and~\eqref{it:local-expansion-determinant-d6}, and the sentence of~\eqref{it:local-expansion-determinant-d4} that derives two of these inequalities from condition~\eqref{eq:renormalisation-group-A10} does not belong to the assertion. All objects are formed with the numbers $M$, $\alpha_0$ and $\varepsilon$ of the theorem: $v_k$ and $r_k$ are formed with $e_k$ by the formulas of Section~\ref{sec:running-quantities-tori}, $r$ being the integer of stage~\ref{it:order-choice-8}, and $\kappa$, $\kappa_0$, $c_X$ and $c_{\mathrm{rb}}$ are the numbers of stages~\ref{it:order-choice-1} and~\ref{it:order-choice-3}. These assertions hold under the hypotheses of the theorem, without any further condition on the constants.
\end{enumerate}
Proposition~\ref{prop:local-expansion-determinant} holds for every system of primitives that has the properties required in its statement, each clause of that proposition under its own hypotheses alone; it is stated for the parameter set $\mathrm{par}$ of Section~\ref{sec:primitive-operators}, with $C_d$ formed with the constant $C_w$ of $\mathrm{par}$, and for the ceiling $e_0$ of stage~\ref{it:order-choice-9}. The conditions on the constants on which its clauses rest are absent from its statement, because they are part of the choice of the constants: the hypotheses of Theorem~\ref{thm:expansion-walks} hold by stages~\ref{it:order-choice-4}, \ref{it:order-choice-6} and~\ref{it:order-choice-9}; condition~\eqref{eq:renormalisation-group-23}, on which~\eqref{it:local-expansion-determinant-d4}--\eqref{it:local-expansion-determinant-d6} rest, holds by stage~\ref{it:order-choice-9}; and~\eqref{it:local-expansion-determinant-d6} rests, in addition, on the inequality $L(\kappa-\kappa_0-1)\ge\kappa$, which follows from the inequality $L(\kappa-6\kappa_0-4)-1\ge\kappa$ of stage~\ref{it:order-choice-4}, and on condition~\eqref{eq:renormalisation-group-24}, in which $C_d$ is the constant of the proposition, formed with the constant $C_w$ of $\mathrm{par}$; in this form condition~\eqref{eq:renormalisation-group-24} is part of stage~\ref{it:order-choice-9} and of condition~\eqref{eq:renormalisation-group-A10}. The general form of Proposition~\ref{prop:local-expansion-determinant} is Proposition~\ref{prop:polymer-expansion-determinant}: there the rate $\gamma_w>0$ and the constant $C_w>0$ of the parameter set $(\alpha_0,\gamma_w,\kappa_w;C_w)$ are arbitrary, $C_d$ is computed by~\eqref{eq:operators-fermion-10} from the constant $C_w$ at hand, and the conditions just named are hypotheses, those of Theorem~\ref{thm:expansion-walks} in its head, condition~\eqref{eq:renormalisation-group-23} in clause~\eqref{it:polymer-expansion-determinant-e}, and the inequality $L(\kappa-\kappa_0-1)\ge\kappa$ and condition~\eqref{eq:renormalisation-group-24} in clause~\eqref{it:polymer-expansion-determinant-g}. For the system of~(ii) with $\kappa'\ge\kappa_w$ one has $C_w=C^F$. Conditions~\eqref{eq:renormalisation-group-20}--\eqref{eq:renormalisation-group-23}, and condition~\eqref{eq:renormalisation-group-24} with $C_w=C^F$, hold with $e^F_0$ in place of $e_0$, by the definition of $e^F_0$ in Section~\ref{sec:conditions}; the hypotheses of Theorem~\ref{thm:expansion-walks} on $L$ and $M$ follow from $L\ge5$, $L$ odd, and $M\ge M^F(b,\alpha,L,\alpha_0,\kappa')$; and the inequality $L(\kappa-\kappa_0-1)\ge\kappa$ holds because $L\ge5$ and $\kappa\ge7\kappa_0+7$ (stage~\ref{it:order-choice-3}):
\[
L(\kappa-\kappa_0-1)-\kappa=(L-1)\cdot\kappa-L\cdot(\kappa_0+1)\ge(6L-7)\cdot(\kappa_0+1)>0.
\]
For this reason assertion~(iii) holds without further conditions on the constants. Moreover, for the constants of Section~\ref{sec:order-choice} and $\kappa'=\kappa_w+\kappa_0$ one has $e_0\le e^F_0$ by condition~\eqref{eq:renormalisation-group-A10}, and the system of~(ii) has the properties required in Proposition~\ref{prop:local-expansion-determinant} also for $\mathrm{par}$, because $\gamma_w\le\gamma^F$, $\kappa_w\le\kappa_w+\kappa_0$ and $C^F\le C_w$ (Lemma~\ref{lem:assembled-systems-primitives}); so Proposition~\ref{prop:local-expansion-determinant} applies to this system as stated, for every $e_k\in(0,e_0]$.

The numbers of the theorem are $\mathfrak{m}_0=1$ and $\alpha^F=\frac{1}{54}$; moreover $\gamma^F\le\frac{1}{6L}$, and $M^F(b,\alpha,L,\alpha_0,\kappa')$ is the number defined in condition~\eqref{eq:renormalisation-group-19} of Section~\ref{sec:conditions}. All assertions of the theorem hold for every $\alpha\in(0,1)$ and every odd integer $L\ge L^F:=5$.
\end{theorem}

This is a consequence of Theorem~\ref{thm:fermion-operators}.

\subsubsection{Photon operators}\label{sec:photon-operators}

\begin{theorem}[Photon operators and their weakened families]\label{thm:photon-operators-weakened}
Let $\alpha$ and $L$ be as in Section~\ref{sec:order-choice}. There exist numbers $\alpha^P\in\left(0,\frac{1}{4}\right]$, $\gamma_P\in(0,1]$, $\gamma^P\in(0,1]$, $C_P$, $c_P$ and $c_C$, functions of $(\alpha,L)$ alone, and numbers $M^P$ and $C^P$ that depend on further arguments in the following way. The number $M^P(\alpha,L,\alpha_0)$ is defined for $\alpha_0\in\left(0,\alpha^P\right]$, and the number $C^P(\alpha,L,M,\alpha_0)$ is defined for such $\alpha_0$ and for every power $M=L^m$, $m\in\mathbb{N}$, with $M\ge M^P(\alpha,L,\alpha_0)$. These numbers have the following property. Let $\alpha_0$ and $M$ be as just described, write $C^P$ for $C^P(\alpha,L,M,\alpha_0)$, let $k\ge0$ be an integer, and let $l\in\mathrm{Cl}$. Then the following assertions hold.
\begin{enumerate}
\item[(i)] Propositions~\ref{prop:background-field}--\ref{prop:local-regularity}, \ref{prop:periodisation} and~\ref{prop:gauge-fixing-controlled} hold, the constants $C_P$, $\gamma_P$, $c_P$ and $c_C$ that occur in them being the numbers above.
\item[(ii)] There are families~\eqref{it:primitives-p1} and~\eqref{it:primitives-p2} of level $k$, the family~\eqref{it:primitives-p1} being for $k=0$ the family $H_0(s')$ of property~\ref{it:renormalisation-group-17}, with the properties~\ref{it:renormalisation-group-1}--\ref{it:renormalisation-group-4}, \ref{it:renormalisation-group-8}--\ref{it:renormalisation-group-10}, \ref{it:renormalisation-group-14} and~\ref{it:renormalisation-group-16} for the parameter set $(\alpha_0,\gamma^P,\cdot\,;C^P)$; the third entry, a rate per unit of $d_M$, is left open, since it occurs in none of these properties. The assertions of property~\ref{it:renormalisation-group-17} on $H_0(s')$ hold.
\item[(iii)] Proposition~\ref{prop:fields-domains} holds under the conditions stated in it.
\end{enumerate}
Moreover, Proposition~\ref{prop:axial-gauge} holds. The numbers of the theorem are $\alpha^P=\frac{1}{4}$, $\gamma^P=\frac{\gamma_P}{2}$ and $c_C=3L^{\frac{9}{2}}$, and $M^P(\alpha,L,\alpha_0)$ is the number defined in~\eqref{eq:operators-photon-20}. The threshold of the theorem for $L$ is $L^P:=5$: its assertions hold for every odd integer $L\ge L^P$, $\alpha$ being as above.
\end{theorem}

This is proved as Theorem~\ref{thm:photon-operators}.

\begin{definition}[Parameters of the primitives]
Let $b$, $\alpha$, $\kappa$ and $L$ be the constants of stages~\ref{it:order-choice-1}--\ref{it:order-choice-4} of Section~\ref{sec:order-choice}; in particular $b=1$. For these values of $b$, $\alpha$ and $L$, let $\alpha^F$, $\gamma^F$ and $C^F$ be the numbers of Theorem~\ref{thm:fermion-operators-weakened}, and let $\alpha^P$, $\gamma^P$ and $C^P$ be those of Theorem~\ref{thm:photon-operators-weakened}. Put
\[
\alpha_0:=\min\left(\alpha^F,\alpha^P\right),\qquad\gamma_w:=\min\left(\gamma^F,\gamma^P\right).
\]
Let $M$ be the number of stage~\ref{it:order-choice-6}; it satisfies $M\ge M^F(b,\alpha,L,\alpha_0,\kappa_w+\kappa_0)$ and $M\ge M^P(\alpha,L,\alpha_0)$. Put
\begin{gather*}
C_w^0:=\max\left(C^F(b,\alpha,L,M,\alpha_0,\kappa_w+\kappa_0),\ C^P(\alpha,L,M,\alpha_0)\right),\\
C_w^1:=\max\left(C_w^0,\ 4\exp(2\gamma_w)\cdot M^{27\alpha_0}\right);
\end{gather*}
thus the number $C^F$ in $C_w^0$ is taken at $\kappa'=\kappa_w+\kappa_0$.
\end{definition}

\subsubsection{The mass range}\label{sec:mass-range}

\begin{theorem}[Range of the running masses]\label{thm:range-running-masses}
Proposition~\ref{prop:bounds-masses} holds under the hypotheses of Theorem~\ref{thm:scales-stop-levels}, that is, whenever $v_K\ge1$ and $\left(1+\frac{\mathrm{lg}}{v_K}\right)^r\le L$.
\end{theorem}

The theorem is a consequence of Proposition~\ref{prop:bounds-running-masses}.

\subsubsection{Hypotheses of the propositions on activities}\label{sec:hypotheses-propositions-activities}

Hypothesis~\ref{def:hypothesis-one-step} \emph{without primitives} consists of the requirements of hypothesis~\ref{def:hypothesis-one-step} other than the datum of a system of primitives of level $k$ on $T_k$ with the properties~\ref{it:renormalisation-group-1}--\ref{it:renormalisation-group-17} for $\mathrm{par}$. Hypothesis~\ref{def:hypothesis-flow} \emph{without primitives} consists of the requirements of hypothesis~\ref{def:hypothesis-flow} other than the requirement that the systems of primitives be those of Definition~\ref{def:periodisation-projection}.

\begin{theorem}\label{thm:renormalisation-group-1}
\begin{enumerate}
\item[(a)] Propositions~\ref{prop:weighted-coefficient-norms}--\ref{prop:contraction} and~\ref{prop:locality} hold under hypothesis~\ref{def:hypothesis-one-step} without primitives.
\item[(b)] Proposition~\ref{prop:weighted-coefficient-norms} is a statement on abstract Grassmann algebras and holds without any hypothesis on the data $(k,l,\mathfrak{m},e_k)$. Proposition~\ref{prop:polymers}, with the exception of its clauses~\eqref{it:polymers-t5} and~\eqref{it:polymers-t8}, and the assertions (a)--(d) of Proposition~\ref{prop:symmetries-extraction} hold for every integer $k\ge0$ under the sole condition $l\in\mathrm{Cl}$ on the side; in particular, they hold in the scale-free setting of Section~\ref{sec:properties-scheme}. In assertion~(d) the torus of level $k+1$ has the side $l/L$, which for $l\in\mathrm{Cl}$ is $M$ times an odd integer $\ge13$, hence a multiple of $M$ larger than $2M$, and which need not belong to $\mathrm{Cl}$; the objects of level $k+1$ that occur in~(d), namely the group $\mathbb{G}_{k+1}$, the bond functions, the generators $\psi$, $\chi$ and the scalings $A_L$, $\psi_L$, $\chi_L$, are those defined under the sole condition $l\in\mathrm{Cl}$ in the paragraph on the change of scale in Section~\ref{sec:local-parts-adjustment}.
\item[(c)] Propositions~\ref{prop:bounds-unit-lattice} and~\ref{prop:block-covariance-real} hold in the scale-free setting of Section~\ref{sec:properties-scheme} with $k\ge1$, that is, for every integer $k\ge1$, every $l\in\mathrm{Cl}$, every $\mathfrak{m}\in(0,\mathfrak{m}_0]$ and every $e_k\in(0,e_0]$.
\item[(d)] Proposition~\ref{prop:free-determinant} holds under hypothesis~\ref{def:hypothesis-flow} without primitives for every $k\le k^{*}(R)$, the only datum on primitives being, for each level $j<k$, a family~\eqref{it:primitives-p5} of level $j$ on $T_j(R)$ for which clause~\eqref{it:local-expansion-determinant-d0} of Proposition~\ref{prop:local-expansion-determinant} holds; the functions $E^{\mathrm{det}}_j$ and $\mathfrak{P}_k$ of Proposition~\ref{prop:free-determinant} are formed with these families, of which the families~\eqref{it:primitives-p5} of the systems of Definition~\ref{def:periodisation-projection} are a special case. Equivalently, Proposition~\ref{prop:free-determinant}, with $E^{\mathrm{det}}_j$ and $\mathfrak{P}_k$ formed with such families, holds under hypothesis~\ref{def:hypothesis-flow} for every $k\le k^{*}(R)$ and every choice of these families: data that satisfy hypothesis~\ref{def:hypothesis-flow} without primitives satisfy hypothesis~\ref{def:hypothesis-flow} once the systems of primitives are taken to be those of Definition~\ref{def:periodisation-projection}, and these systems do not occur in the assertions of Proposition~\ref{prop:free-determinant}.
\end{enumerate}
\end{theorem}

This is established in Section~\ref{sec:settings-basic-notation}, which lays down the settings in which these propositions are proved.

\begin{remark}
Assume hypothesis~\ref{def:hypothesis-flow}. At the level $k^{*}(R)$ the side $l_{*}$ need not satisfy $l_{*}/L\in\mathrm{Cl}$, and the mass $\mathfrak{m}_{*}$ need not satisfy $\mathfrak{m}_{*}\le\mathfrak{m}_0/L$; both conditions are part of hypothesis~\ref{def:hypothesis-one-step}. The data $(k^{*}(R),l_{*},\mathfrak{m}_{*},e_{*})$ of this level belong to the scale-free setting of Section~\ref{sec:properties-scheme}, by Theorem~\ref{thm:scales-stop-levels}.
\end{remark}

\subsubsection{Hypotheses of the identities for the exact operators}\label{sec:hypotheses-identities-exact}

\begin{theorem}\label{thm:renormalisation-group-2}
\begin{enumerate}
\item[(a)] Proposition~\ref{prop:block-averages} holds for every integer $k\ge0$ and every side $l\in L\mathbb{N}$.
\item[(b)] Propositions~\ref{prop:fluctuation-fields}, \ref{prop:minimisers-prescribed-block}, \ref{prop:comparison-fermion-operators} and~\ref{prop:step-external-connection} hold under hypothesis~\ref{def:hypothesis-one-step} without primitives (Section~\ref{sec:hypotheses-propositions-activities}).
\item[(c)] Propositions~\ref{prop:orthogonal-coordinates} and~\ref{prop:closed-shifts-terminal} hold under hypothesis~\ref{def:hypothesis-flow} without primitives.
\end{enumerate}
\end{theorem}

Assertion~(a) is a special case of Proposition~\ref{prop:averaging-operators}. Assertion~(b) is established for Propositions~\ref{prop:fluctuation-fields} and~\ref{prop:minimisers-prescribed-block} in Section~\ref{sec:one-averaging-step}, for Proposition~\ref{prop:comparison-fermion-operators} in Section~\ref{sec:one-level-next}, and for Proposition~\ref{prop:step-external-connection} in Section~\ref{sec:identity-one-step}, the family $\mathfrak{P}_k$ of that proposition being introduced in Section~\ref{sec:densities-block-transformation}. For~(c), the data with which Propositions~\ref{prop:orthogonal-coordinates} and~\ref{prop:closed-shifts-terminal} are proved are fixed at the beginning of Section~\ref{sec:coordinates-axial-fields}.

\subsubsection{Trajectories, the pair $P_N(e)$ and bounds on $\log\rho_n$}\label{sec:trajectories-pair-bounds}

The conditions~\eqref{eq:renormalisation-group-16}, \eqref{eq:renormalisation-group-fl}, \eqref{eq:renormalisation-group-T1}--\eqref{eq:renormalisation-group-T6} and~\eqref{eq:renormalisation-group-17}, as well as the numbers $\varepsilon_1,\dots,\varepsilon_5$, that occur in Theorems~\ref{thm:scales-stop-levels}--\ref{thm:partition-functions-tori} are given in Section~\ref{sec:conditions}.

\begin{theorem}[Scales, stop levels and the number $e_1$]\label{thm:scales-stop-levels}
\begin{enumerate}
\item[(a)] Let $v_K\ge1$ and $\left(1+\frac{\mathrm{lg}}{v_K}\right)^r\le L$. Then
\[
r_{k+1}\in\left\{r_k,\frac{r_k}{L}\right\}\qquad\text{and}\qquad\mathfrak{n}_{k+1}-\mathfrak{n}_k\in\{-1,0\}\qquad\text{for }0\le k<K.
\]
Consequently, for every $R\in\mathrm{Tor}(e)$ and every $N\ge N_{\mathrm{min}}(R,e)$ the stop level $k^{*}(R)$ of Definition~\ref{def:tori-stop-level} exists; below, $k^{*}(R)$ and the quantities with the subscript $*$ are taken for such $N$. The numbers $K-k^{*}(R)$, $a_{*}$, $e_{*}$, $l_{*}$ and $\mathfrak{m}_{*}$ are the same for all $N\ge N_{\mathrm{min}}(R,e)$. The function $R\mapsto N_{\mathrm{min}}(R,e)$ is non-increasing on $\mathrm{Tor}(e)$, and $k^{*}(R)\le k^{*}(R')$ for all $R\le R'$ in $\mathrm{Tor}(e)$. If moreover
\[
\left(1+\frac{3\mathrm{lg}}{v_K}\right)^r<L,
\]
then $k^{*}(R^{*})\in\{K-1,K-2\}$; on the torus $R^{*}$ one has
\[
a_{*}\in\left\{\frac{h'}{L},\frac{h'}{L^2}\right\}\qquad\text{and}\qquad r_{*}\in\left\{r_K,L\cdot r_K\right\};
\]
and $k_{\mathrm{top}}\in\{K,K-1\}$.
\item[(b)] There is a number $e_1>0$, namely the number of stage~\ref{it:order-choice-10}, which is defined by condition~\eqref{eq:renormalisation-group-26}, with the following properties. Let $e>0$ satisfy $e_K=e\cdot h'^{\frac{1}{2}}\le e_1$. Then $e\le\frac{1}{3}$ and $v_K\ge8$; every $R\in\mathrm{Tor}(e)$ satisfies, for all $N\ge N_{\mathrm{min}}(R,e)$, the conditions~\eqref{eq:renormalisation-group-16}, \eqref{eq:renormalisation-group-fl} and~\eqref{eq:renormalisation-group-T1}--\eqref{eq:renormalisation-group-T6}; and
\[
\left(1+\frac{3\mathrm{lg}}{v_K}\right)^r<L.
\]
Moreover, condition~\eqref{eq:renormalisation-group-17} holds for every $e>0$ that satisfies $e_K\le e_1$ and condition~\eqref{eq:renormalisation-group-18}, that is,
\[
e_K\cdot\log h'\le1.
\]
\item[(c)] Let hypothesis~\ref{def:hypothesis-flow} hold. Then, at every level $k<k^{*}(R)$, the data $(k,l_k,\mathfrak{m}_k,e_k)$, together with the primitives of Definition~\ref{def:periodisation-projection} on $T_k(R)$, satisfy hypothesis~\ref{def:hypothesis-one-step}; at $k=0$ the running mass is $\bar{m}\cdot\eta$ (Definition~\ref{def:notation}(b)) and the running charge is $q$ (Definition~\ref{def:notation}(e)), and the properties of the primitives that hypothesis~\ref{def:hypothesis-one-step} requires are provided by Theorems~\ref{thm:fermion-operators-weakened} and~\ref{thm:photon-operators-weakened} and by property~\ref{it:renormalisation-group-18}. At the level $k=k^{*}(R)$ the data $(k^{*},l_{*},\mathfrak{m}_{*},e_{*})$ belong to the scale-free setting of Section~\ref{sec:properties-scheme}: one has $l_{*}\in\mathrm{Cl}$, $\mathfrak{m}_{*}\in(0,\mathfrak{m}_0]$ and $e_{*}\le e_0$.
\end{enumerate}
\end{theorem}

Part~(a) is contained in Proposition~\ref{prop:depth-stop-level}, part~(b) in Proposition~\ref{prop:consequences}, and part~(c) in Proposition~\ref{prop:data-along-flow}.

\begin{theorem}[Projection; agreement on flat polymers]\label{thm:projection-agreement-flat}
Property~\ref{it:renormalisation-group-18} and Proposition~\ref{prop:rows-component} hold. Moreover the following holds.

Let hypothesis~\ref{def:hypothesis-flow} hold, let $R<R_{\mathrm{top}}$, and put $R':=R_{\mathrm{top}}$. For a level $k\le k^{*}(R)$ let
\[
T_1:=T_k(R),\qquad T_2:=T_k(R')
\]
be the fine lattices of level $k$ with the sides $l_k(R)$ and $l_k(R')=J\cdot l_k(R)$, where $J=R'/R$ is a power of $L$; both sides lie in $\mathrm{Cl}$, by Section~\ref{sec:lattice-geometry}, since $k\le k^{*}(R)\le k_{\mathrm{top}}$ (Theorem~\ref{thm:scales-stop-levels}(a)). Let $\mathrm{pr}_i\colon\mathbb{L}_k\to T_i$, $i=1,2$, be the projections, and let $\mathrm{pj}\colon T_2\to T_1$ and $\mathrm{pj}^{*}$ be the covering map and the inclusion of Definition~\ref{def:periodisation-projection}.
\begin{enumerate}
\item[(i)] Let $\tilde{X}$ be a non-empty finite union of $M$-cubes of $\mathbb{L}_k$. Among the sets $I_1\times I_2\times I_3$, with intervals $I_1$, $I_2$, $I_3$, that are unions of $M$-cubes, let $\mathrm{box}(\tilde{X})$ be the least one containing $\tilde{X}$, and put
\[
N(\tilde{X}):=\mathrm{box}(\tilde{X})^{\sim6d_U}.
\]
The set $\tilde{X}$ is \emph{admissible} if the largest side of $\mathrm{box}(\tilde{X})$ is at most $l_k(R)-14\,d_U\cdot M$.
\item[(ii)] A polymer of $T_i$ is \emph{flat} if it equals $\mathrm{pr}_i(\tilde{X})$ for an admissible set $\tilde{X}$ that is connected through faces.
\item[(iii)] For $i=1,2$ let $F^{(i)}=\left\{F^{(i)}(X;\mathscr{A},\mathscr{C})\right\}$ be a family of even elements of the Grassmann algebra $\mathrm{Gr}[\psi,\chi]$ on $T_i$, where $X$ ranges over the polymers of $T_i$ and $(\mathscr{A},\mathscr{C})$ over $\tilde{R}_k\times\mathfrak{X}_k$, or over $\tilde{R}^{-}_k\times\mathfrak{X}_k$, formed on $T_i$; let both families have the locality property~\eqref{it:renormalisation-group-1-k2} of Definition~\ref{def:renormalisation-group-1}. The generators on $T_1$ and on $T_2$ are indexed by points of $\mathbb{L}_k$ as in Section~\ref{sec:grassmann-algebras-norms}; the index of a generator $\psi(\xi)$, $\xi=(x,\boldsymbol{\beta},\boldsymbol{\omega})$, lies in a set $W$ of sites of $\mathbb{L}_k$ if $x\in W$, and that of a generator $\chi(\boldsymbol{\zeta})$, $\boldsymbol{\zeta}=(x,x',\boldsymbol{\beta},\boldsymbol{\omega})$, if $x\in W$ and $x'\in W$. The families \emph{agree on flat polymers}, written $F^{(1)}\sim F^{(2)}$, if the following holds for every admissible $\tilde{X}$ that is connected through faces. Whenever $(\mathscr{A}_1,\mathscr{C}_1)$ on $T_1$ and $(\mathscr{A}_2,\mathscr{C}_2)$ on $T_2$ lie in the domains of the families and satisfy
\[
\mathscr{A}_1(\mathrm{pr}_1b)=\mathscr{A}_2(\mathrm{pr}_2b)\qquad\text{and}\qquad\mathscr{C}_1(\mathrm{pr}_1b)=\mathscr{C}_2(\mathrm{pr}_2b)
\]
for every bond $b$ of $\mathbb{L}_k$ contained in $N(\tilde{X})$, then $F^{(1)}(\mathrm{pr}_1\tilde{X};\mathscr{A}_1,\mathscr{C}_1)$ is carried to $F^{(2)}(\mathrm{pr}_2\tilde{X};\mathscr{A}_2,\mathscr{C}_2)$ by the isomorphism of algebras that maps every generator on $T_1$ with index in $N(\tilde{X})$ to the generator on $T_2$ with the same index. For families without the second connection the same definition is used with $\mathscr{C}_1=\mathscr{C}_2=0$.
\item[(iv)] For $i=1,2$ let $\mathfrak{P}^{(i)}=\left\{\mathfrak{P}^{(i)}(X,\cdot):\ X\in D_k(T_i)\right\}$ be a family of complex functions on $\mathfrak{X}_k(T_i)$. These families agree, written $\mathfrak{P}^{(1)}\sim\mathfrak{P}^{(2)}$, if
\[
\mathfrak{P}^{(2)}\left(\mathrm{pr}_2\tilde{X},\mathrm{pj}^{*}\mathscr{C}\right)=\mathfrak{P}^{(1)}\left(\mathrm{pr}_1\tilde{X},\mathscr{C}\right)
\]
for every admissible $\tilde{X}$ that is connected through faces and every $\mathscr{C}\in\mathfrak{X}_k(T_1)$; here $\mathrm{pj}^{*}\mathscr{C}\in\mathfrak{X}_k(T_2)$ by Lemma~\ref{lem:periodic-extension}.
\end{enumerate}
Lemma~\ref{lem:kernels-near-admissible} holds; at the levels $k<k^{*}(R)$ it compares the kernels of the primitives of Definition~\ref{def:periodisation-projection} on $T_1$ with those on $T_2$ near the images $\mathrm{pr}_1(\tilde{Y})$ and $\mathrm{pr}_2(\tilde{Y})$ of an admissible set $\tilde{Y}$.

Let $k<k^{*}(R)$. The objects without a prime are formed on the torus $R$, with the primitives of Definition~\ref{def:periodisation-projection} on $R$, and the objects with a prime are formed in the same way on the torus $R'$. Let $m$ be a real number with $|m|\le w^{\mathrm{m}}_k$, and let $\mathscr{E}\in K^x_k$ on $T_k(R)$ and $\mathscr{E}'\in K^x_k$ on $T_k(R')$ satisfy
\[
\|\mathscr{E}\|^x_k\le w^{\mathrm{E}}_k,\qquad\|\mathscr{E}'\|^x_k\le w^{\mathrm{E}}_k,\qquad\mathscr{E}\sim\mathscr{E}',
\]
where $\mathscr{E}=\mathscr{E}'=0$ if $k=0$. Then the following hold.
\begin{enumerate}
\item[(a)] $\varepsilon(\mathscr{E})=\varepsilon(\mathscr{E}')$ and $m(\mathscr{E})=m(\mathscr{E}')$, and the number $d\varepsilon^0_k$ is the same on both tori.
\item[(b)] Let
\[
\mathscr{E}^{+}:=LR\mathscr{E}+L\delta E^{\mathrm{det}}_k+E^{*,x}_k(m,\mathscr{E}),
\]
formed with the operations of Section~\ref{sec:local-parts-adjustment}, the family $\delta E^{\mathrm{det}}_k$ of Definition~\ref{def:determinant-functions} and the map $E^{*,x}_k$ of Definition~\ref{def:renormalisation-step}, and let $\mathscr{E}'^{+}$ be formed from $\mathscr{E}'$ in the same way. Then $\mathscr{E}^{+}\sim\mathscr{E}'^{+}$ at the level $k+1$, that is, on the lattices $T_{k+1}(R)$ and $T_{k+1}(R')$.
\item[(c)] Let $\mathfrak{P}_k$ and $\mathfrak{P}'_k$ be families of complex functions on $\mathfrak{X}_k$, indexed by the polymers of $T_k(R)$ and of $T_k(R')$ respectively, and let $\mathfrak{P}_{k+1}$ be formed from $\mathfrak{P}_k$ as in clause~\eqref{it:step-small-fields-c} of Theorem~\ref{thm:step-small-fields}, that is,
\[
\mathfrak{P}_{k+1}(Y,\mathscr{C}):=\sum_{X\in D_k:\ \bar{X}^L=LY}\left(\mathfrak{P}_k(X,\mathscr{C}_L)+E^{\mathrm{det}}_k(X,\mathscr{C}_L)\right)\qquad\left(Y\in D_{k+1},\ \mathscr{C}\in\mathfrak{X}_{k+1}\right),
\]
and let $\mathfrak{P}'_{k+1}$ be formed from $\mathfrak{P}'_k$ in the same way. If $\mathfrak{P}_k\sim\mathfrak{P}'_k$, then $\mathfrak{P}_{k+1}\sim\mathfrak{P}'_{k+1}$.
\end{enumerate}
\end{theorem}

The notions (i)--(iv) reappear in Definition~\ref{def:admissibility-flatness-agreement}, and assertions (a)--(c) are proved as Proposition~\ref{prop:agreement-after-one}.

\begin{theorem}[Trajectories on a fixed torus]\label{thm:trajectories-fixed-torus}
Let hypothesis~\ref{def:hypothesis-flow} hold. The trajectories below are those of Section~\ref{sec:flow} on the torus $R$, formed with the primitives of Definition~\ref{def:periodisation-projection}, and $k^{*}=k^{*}(R)$.
\begin{enumerate}
\item[(a)] Let $0\le k\le k^{*}$, and let $(m_{\mathrm{in}},\boldsymbol{\varepsilon}_{\mathrm{in}})$ be a real pair with the following property, for $i=0,1,\dots,k-1$ in this order: the trajectory of the pair has a mass $m_i$ at the level $i$, and $|m_i|\le w^{\mathrm{m}}_i$. Then the data $E_k\in K_k$ and $\mathscr{E}_k\in K^x_k$ of the level $k$ exist, $\mathscr{E}_k|_0=E_k$, and
\[
\|E_k\|_k\le\|\mathscr{E}_k\|^x_k\le4\,e_k^{\varepsilon}\cdot w^{\mathrm{E}}_k.
\]
\item[(b)] Among the real pairs that are controlled on $R$ (Definition~\ref{def:fluctuation-field-controlled}), exactly one has a trajectory with $m_{k^{*}}=0$ and $\boldsymbol{\varepsilon}^t_{k^{*}}=0$. At every level $k$ with $0\le k\le k^{*}$ the trajectory of this pair obeys
\[
|m_k|\le\frac{w^{\mathrm{m}}_k}{8},\qquad\boldsymbol{\varepsilon}^t_k=-\sum_{i=k}^{k^{*}-1}L^{-3(i-k)}\cdot\left(\varepsilon(E_i)+d\varepsilon^0_i\right),\qquad\left|\boldsymbol{\varepsilon}^t_k\right|\le5\,c_{\mathrm{fl}}\cdot e_k^{\varepsilon}\cdot w^{\mathrm{E}}_k,
\]
where $c_{\mathrm{fl}}$ is the constant of stage~\ref{it:order-choice-1}.
\end{enumerate}
\end{theorem}

Part~(a) is contained in Proposition~\ref{prop:existence-trajectory}, and part~(b) is proved as Proposition~\ref{prop:pair-vanishing-terminal}.

\begin{definition}[The tuned pair]\label{def:tuned-pair}
Let $e>0$ satisfy $e_K=e\cdot h'^{\frac{1}{2}}\le e_1$, and let $N\ge N_{\mathrm{min}}(R_{\mathrm{top}},e)$; then hypothesis~\ref{def:hypothesis-flow} holds with $R=R_{\mathrm{top}}$. The pair
\[
P_N(e)=\left(m_{0,N},\boldsymbol{\varepsilon}_{0,N}\right)
\]
is the real pair of Theorem~\ref{thm:trajectories-fixed-torus}(b) on the torus $R_{\mathrm{top}}$.
\end{definition}

\begin{theorem}[One pair for all tori]\label{thm:one-pair-all}
Let $e>0$ satisfy $e_K\le e_1$, let $R\in\mathrm{Tor}(e)$, and let $N\ge N_{\mathrm{min}}(R,e)$; then $N\ge N_{\mathrm{min}}(R_{\mathrm{top}},e)$ by Theorem~\ref{thm:scales-stop-levels}, so that $P_N(e)$ is defined. The pair $P_N(e)$ is controlled on $R$. Denote by $\left(\boldsymbol{\varepsilon}^t_k(R),m_k(R),E_k(R)\right)$ and $\left(\mathscr{E}_k(R),\mathfrak{P}_k(R)\right)$ its trajectory and its trajectory with slot on $R$ (Section~\ref{sec:flow}), and by the same symbols with $R_{\mathrm{top}}$ in place of $R$ those on $R_{\mathrm{top}}$. Then, for $0\le k\le k^{*}(R)$,
\begin{gather*}
m_k(R)=m_k(R_{\mathrm{top}}),\qquad\boldsymbol{\varepsilon}^t_k(R)=\boldsymbol{\varepsilon}^t_k(R_{\mathrm{top}}),\\
E_k(R)\sim E_k(R_{\mathrm{top}}),\qquad\mathscr{E}_k(R)\sim\mathscr{E}_k(R_{\mathrm{top}}),\qquad\mathfrak{P}_k(R)\sim\mathfrak{P}_k(R_{\mathrm{top}}),
\end{gather*}
where $\sim$ is the agreement on flat polymers of Theorem~\ref{thm:projection-agreement-flat} at the level $k$ if $R<R_{\mathrm{top}}$, and equality if $R=R_{\mathrm{top}}$. Consequently
\[
|m_k(R)|\le\frac{w^{\mathrm{m}}_k}{8}\quad\left(0\le k\le k^{*}(R)\right),\qquad\left|\boldsymbol{\varepsilon}^t_{k^{*}(R)}(R)\right|\le5\,c_{\mathrm{fl}}\cdot e_{*}^{\varepsilon}\cdot w^{\mathrm{E}}_{*}.
\]
The pair $P_N(e)$ is real; it depends on $N$, $e$, $\bar{m}$ and on the scheme (Definition~\ref{def:terminology}(10)), but neither on $R$ nor on the twist $\theta$.
\end{theorem}

The assertions are contained in Theorem~\ref{thm:agreement-trajectories}.

\begin{theorem}[Representation with twist and external connection]\label{thm:representation-twist-external}
Let hypothesis~\ref{def:hypothesis-flow} hold, and let $(m_{\mathrm{in}},\boldsymbol{\varepsilon}_{\mathrm{in}})$ be a real pair that is controlled on $R$, with the trajectory $(\boldsymbol{\varepsilon}^t_k,m_k,E_k)$ and the trajectory with slot $(\mathscr{E}_k,\mathfrak{P}_k)$ on $R$. For $A\in\mathrm{Dom}_{k^{*}}$ and $\mathscr{C}\in\mathfrak{X}_{k^{*}}$ put
\[
Y^x_{*}(A;\mathscr{C}):=y^x_{k^{*}}\left[\boldsymbol{\varepsilon}^t_{*},m_{*},\mathscr{E}_{*}\right](A;\mathscr{C}),
\]
the function $y^x_{k^{*}}$ of Section~\ref{sec:fermion-integrals-blocking} being formed with the family $\mathfrak{P}_{k^{*}}$. Let $a\in C^1\otimes\mathbb{C}$ satisfy $a^{[j]}\in\mathfrak{X}_j$ for all $j\le k^{*}$, and put $\mathscr{C}_{*}:=a^{[k^{*}]}$. Let $A_{*}$, $\zeta=(\zeta_j)_{j<k^{*}}$ with $\zeta_j\in V_j$, $A^{\downarrow}_0$, $\mathrm{d}\mathsf{m}_{*}$, $\chi_{*}$ and $\mathfrak{G}$ be as in Definition~\ref{def:coordinates-slice-small}, and let $c_{\theta}:=c^{(k^{*})}_{\theta}$ be the twist of Definition~\ref{def:twist-external-connection}.
\begin{enumerate}
\item[(a)] For every $A\in\mathrm{Dom}_{k^{*}}$,
\[
Y^x_{*}(A;\mathscr{C}_{*})=\int\prod_{j<k^{*}}\mathrm{d}\mu_I(\zeta_j)\;\prod_{j<k^{*}}\mathbf{1}_{\mathrm{Box}_j}(\zeta_j)\cdot y^x_0\left[\boldsymbol{\varepsilon}_{\mathrm{in}},m_{\mathrm{in}},0\right]\left(A^{\downarrow}_0(A,\zeta);a\right),
\]
the integral being taken over $\zeta\in\prod_{j<k^{*}}V_j$, with the standard Gaussian measure $\mu_I$ on each $V_j$. The integrand is read as $0$ where one of the factors $\mathbf{1}_{\mathrm{Box}_j}(\zeta_j)$ vanishes. If $\zeta_j\in\mathrm{Box}_j$ for every $j<k^{*}$, then $A^{\downarrow}_0(A,\zeta)\in\mathrm{Dom}_0$, by Proposition~\ref{prop:fields-domains}(ii) applied at the levels $j=k^{*}-1,\dots,0$ in turn (hypothesis~\ref{def:hypothesis-one-step} holds at these levels by Theorem~\ref{thm:scales-stop-levels}(c)).
\item[(b)] Let $\lambda$ be a real function on the sites of $T^0_{k^{*}}$. If $A$ and $A+d\lambda$ lie in $\mathrm{Dom}_{k^{*}}$, then
\[
Y^x_{*}(A+d\lambda;\mathscr{C}_{*})=Y^x_{*}(A;\mathscr{C}_{*}).
\]
\item[(c)] For every $\theta\in\mathbb{R}^3$ and every bounded Borel function $\Phi$ on the real bond functions on $T^0_{k^{*}}$,
\[
\nu^{\theta,a}\left(\mathbf{1}_{\mathfrak{G}}\cdot\Phi\right)=\int\mathrm{d}\mathsf{m}_{*}(A)\,\Phi(A)\cdot\chi_{*}(A)\cdot Y^x_{*}\left(A+c_{\theta};\mathscr{C}_{*}\right).
\]
Here, on the left, $\Phi$ stands for the function $F\mapsto\Phi\left(A_{*}(F)\right)$ on $F_N$, and $\nu^{\theta,a}(f)$ is the integral of $f$ with respect to $\nu^{\theta,a}$; on the right, the integrand is read as $0$ where $\chi_{*}(A)=0$, and where $\chi_{*}(A)=1$ the field $A+c_{\theta}$ lies in $\mathrm{Dom}_{k^{*}}$ by Proposition~\ref{prop:fields-domains}(i), because $d(A+c_{\theta})=dA$.
\end{enumerate}
\end{theorem}

This is proved as Theorem~\ref{thm:representation-twist-external-2}.

\begin{theorem}[Estimates on one torus]\label{thm:estimates-one-torus}
Let hypothesis~\ref{def:hypothesis-flow} hold, let $(m_{\mathrm{in}},\boldsymbol{\varepsilon}_{\mathrm{in}})$ be a real pair that is controlled on $R$, and let $Y^x_{*}$ and $c_{\theta}$ be formed with this pair as in Theorem~\ref{thm:representation-twist-external}. Put
\[
Y_{*}(A):=Y^x_{*}(A;0)\quad\left(A\in\mathrm{Dom}_{k^{*}}\right),\qquad D_{*}:=D_{k^{*}},\qquad R_{\theta}:=\frac{\det D_{*}(c_{\theta})}{\det D_{*}(0)}\quad\left(\theta\in\mathbb{R}^3\right),
\]
where $D_{k^{*}}$ is the operator $D_k$ of Section~\ref{sec:exact-fermion-operators} at the level $k^{*}$ (one has $\det D_{*}(0)\ne0$ by Proposition~\ref{prop:block-covariance-real}); with $n=R/\eta$ write $\rho_n$ for $\rho_n\left(q,\bar{m}\cdot\eta+m_{\mathrm{in}}\right)$. For $\mathbf{m}>0$ and $p\in\mathbb{R}^3$ put
\[
w_{\mathbf{m}}(p):=\left(\mathbf{m}+\sum_{\nu=0}^2\left(1-\cos p_{\nu}\right)\right)^2+\sum_{\nu=0}^2\sin^2p_{\nu},\qquad p^0(\mathbf{m}):=\frac{2}{(2\pi)^3}\int_{[-\pi,\pi]^3}\mathrm{d}^3p\,\log w_{\mathbf{m}}(p).
\]
Let $\gamma_1$ and $C_{\Sigma}$ be the constants of stage~\ref{it:order-choice-7}.
\begin{enumerate}
\item[(a)] \emph{The terminal integrand.} Let $\theta\in\mathbb{R}^3$, and let $A$ be a real bond function on $T^0_{k^{*}}$ with mean zero and $|dA|_{\infty}\le p_{*}$; then $A+c_{\theta}$ lies in $\mathrm{Dom}_{k^{*}}$ by Proposition~\ref{prop:fields-domains}(i), because $d(A+c_{\theta})=dA$. One has
\[
Y_{*}(A+c_{\theta})=\exp\left(-\boldsymbol{\varepsilon}^t_{*}\cdot l_{*}^3\right)\cdot R_{\theta}\cdot\left(1+\vartheta\right)
\]
with a number $\vartheta$ that satisfies $|\vartheta|\le\varepsilon_1$.
\item[(b)] \emph{Large fields.} For every $\theta\in\mathbb{R}^3$ the measure $\nu^{\theta}$ is non-negative; the number $\bar{\nu}(1)$ is positive and finite; and
\[
\bar{\mu}\left(F_N\setminus\mathfrak{G}\right)\le\varepsilon_3,\qquad\mathrm{d}\mathsf{m}_{*}\left(\left\{A:\ \chi_{*}(A)=0\right\}\right)\le\varepsilon_4.
\]
\item[(c)] \emph{The partition function.} For every $\theta\in\mathbb{R}^3$ one has
\begin{multline}\label{eq:renormalisation-group-12}
\exp\left(-\boldsymbol{\varepsilon}^t_{*}\cdot l_{*}^3\right)\cdot R_{\theta}\cdot\left(1-\varepsilon_1\right)\cdot\left(1-\varepsilon_4\right)\le\nu^{\theta}(\mathfrak{G})\\
\le\exp\left(-\boldsymbol{\varepsilon}^t_{*}\cdot l_{*}^3\right)\cdot R_{\theta}\cdot\left(1+\varepsilon_1\right).
\end{multline}
Moreover,
\begin{equation}\label{eq:renormalisation-group-13}
\bar{\nu}(\mathfrak{G})\le\bar{\nu}(1)\le\frac{\bar{\nu}(\mathfrak{G})}{1-\varepsilon_3},
\end{equation}
\begin{multline}\label{eq:renormalisation-group-14}
\left|\log\rho_n-\log\rho^{\mathrm{free}}_n-\boldsymbol{\varepsilon}_{\mathrm{in}}\cdot n^3+\boldsymbol{\varepsilon}^t_{*}\cdot l_{*}^3\right|\\
\le2\,\varepsilon_1+\varepsilon_2+2\,\varepsilon_3+2\,\varepsilon_4\le2\,e_{*}^{\frac{1}{4}-9\varepsilon},
\end{multline}
\begin{multline}\label{eq:renormalisation-group-15}
\left|\log\rho_n-\left(p^0(\bar{m}\cdot\eta)+\boldsymbol{\varepsilon}_{\mathrm{in}}\right)\cdot n^3+\boldsymbol{\varepsilon}^t_{*}\cdot l_{*}^3\right|\\
\le\varepsilon_5(\bar{m}\cdot R)+2\,e_{*}^{\frac{1}{4}-9\varepsilon}.
\end{multline}
\item[(d)] \emph{Determinants with a constant connection.} Let $n\ge1$ be any integer, let $\theta\in\mathbb{R}^3$, $0<\mathbf{m}\le1$ and $q>0$, and let $D_{H_{\theta}}$ be the operator of Section~\ref{sec:lattice-measures-laws} on $\Lambda_n$, formed with this $q$; if $n\le2$, the one-cochain $H_{\theta}$ is regarded, in the manner of Section~\ref{sec:characters-twisted-spinors}, as a function on the links of $\mathbb{Z}^3$ with period $n$ in each direction. Then
\[
\det\left(D_{H_{\theta}}+\mathbf{m}\right)>0\qquad\text{and}\qquad\left|\log\det\left(D_{H_{\theta}}+\mathbf{m}\right)-p^0(\mathbf{m})\cdot n^3\right|\le\varepsilon_5(\mathbf{m}\cdot n).
\]
For $n=R/\eta$ one has
\[
\left|\log\rho^{\mathrm{free}}_n-p^0(\bar{m}\cdot\eta)\cdot n^3\right|\le\varepsilon_5(\bar{m}\cdot R).
\]
For every $\theta\in\mathbb{R}^3$ one has $R_{\theta}>0$ and
\[
\left|\log R_{\theta}\right|\le2\,\varepsilon_5\left(\mathfrak{m}_{*}\cdot l_{*}\right)+C_{\Sigma}\cdot l_{*}^{-1}\cdot\exp\left(-\gamma_1\cdot l_{*}\right).
\]
\end{enumerate}
\end{theorem}

Part~(a) is proved as Proposition~\ref{prop:estimate-stop-level}, and part~(b) is contained in Lemma~\ref{lem:positivity-large-fields}; part~(c) reappears as Theorem~\ref{thm:twist-average-one}, with its proof, and part~(d) follows from the results of Section~\ref{sec:determinants-constant-connections}.

\begin{theorem}[Partition functions on the tori of $\mathrm{Tor}(e)$]\label{thm:partition-functions-tori}
Let $e>0$ satisfy $e_K\le e_1$, let $R\in\mathrm{Tor}(e)$ and $N\ge N_{\mathrm{min}}(R,e)$, put $n:=R/\eta$, let $P_N(e)=(m_{0,N},\boldsymbol{\varepsilon}_{0,N})$ be the pair of Definition~\ref{def:tuned-pair}, and put
\[
\rho_n:=\rho_n\left(e\cdot\eta^{\frac{1}{2}},\bar{m}\cdot\eta+m_{0,N}\right),\qquad p_F:=p^0(\bar{m}\cdot\eta)+\boldsymbol{\varepsilon}_{0,N},
\]
with the function $p^0$ of Theorem~\ref{thm:estimates-one-torus} and the number $\rho^{\mathrm{free}}_n$ of Section~\ref{sec:lattice-measures-laws}. Then
\begin{align*}
&\left|\log\rho_n-\log\rho^{\mathrm{free}}_n-\boldsymbol{\varepsilon}_{0,N}\cdot n^3\right|\le3\,e_{*}^{\frac{1}{4}-9\varepsilon},\\
&\left|\log\rho_n-p_F\cdot n^3\right|\le\varepsilon_5(\bar{m}\cdot R)+3\,e_{*}^{\frac{1}{4}-9\varepsilon}\le\varepsilon_5\left(L^{-m_0-1}\right)+3\,e_K^{\frac{1}{4}-9\varepsilon}.
\end{align*}
Let moreover $N\ge N_{\mathrm{min}}(R^{*},e)$, put $n_{*}:=R^{*}/\eta$, and form $\rho_{n_{*}}$ and $\rho_{L\cdot n_{*}}$ with the same two arguments as $\rho_n$. Then
\begin{multline*}
D^{\mathrm{mat}}:=L^3\cdot\log\rho_{n_{*}}-\log\rho_{L\cdot n_{*}}\\
\le\left(L^3+1\right)\cdot\left(\varepsilon_5\left(\bar{m}\cdot R^{*}\right)+3\,e_K^{\frac{1}{4}-9\varepsilon}\right),
\end{multline*}
and $-D^{\mathrm{mat}}$ satisfies the same bound.
\end{theorem}

These inequalities are contained in Theorem~\ref{thm:partition-function-pair}.

\subsubsection{Probes and jets on the reference torus}\label{sec:probes-jets-reference}

\begin{definition}[Hypothesis on the reference torus]\label{def:hypothesis-reference-torus}
Hypothesis~\ref{def:hypothesis-flow} is satisfied with the reference torus $R=R^{*}$ of Definition~\ref{def:tori-stop-level}, and the pair is $\left(m_{\mathrm{in}},\boldsymbol{\varepsilon}_{\mathrm{in}}\right)=P_N(e)$, the pair of Definition~\ref{def:tuned-pair}.
\end{definition}

The set $\mathrm{Cel}$ of the $N^{*}=8L^3$ cells of the reference torus (the torus $1$), the base cell $\mathfrak{c}_0$, the maps $\mathsf{g}_c$ and the length $l_c=R^{*}/(2L)$ are those of Section~\ref{sec:cells-reference-tori}.

\begin{definition}[Probes]\label{def:probes}
Assume hypothesis~\ref{def:hypothesis-reference-torus}. A \emph{probe} is a real two-form $\beta\in C_c^{\infty}\left(\mathbb{R}^3;\Lambda^2\right)$ with $\beta\ne0$. Let $\beta$ be a probe, and let $\rho_{\beta}$ be the least number $\rho$ for which some point $y\in\mathbb{R}^3$ satisfies
\[
\operatorname{supp}\beta+[-2,2]^3\subset y+[-\rho,\rho]^3,
\]
let $y_{\beta}$ be a point $y$ for which this inclusion holds with $\rho=\rho_{\beta}$, and put $D_{\beta}:=2\rho_{\beta}$.

Points of the reference torus are written in physical coordinates, in which the base cell $\mathfrak{c}_0$ is the cube $[0,l_c]^3$; these are $\eta$ times the lattice coordinates of Section~\ref{sec:cells-reference-tori}, so that the sites of $\Lambda_{n_{*}}$ are the points of $\eta\mathbb{Z}^3$, taken modulo $R^{*}\mathbb{Z}^3$. A vector $v\in\eta\mathbb{Z}^3$ is an \emph{admitted placement} if
\[
v+y_{\beta}+[-\rho_{\beta},\rho_{\beta}]^3\subset\left[\frac{l_c}{8},\frac{7l_c}{8}\right]^3,
\]
that is, if~\eqref{eq:four-jets-1} holds. The \emph{centred placement} of Section~\ref{sec:probes-discretisation} is a vector $v\in\eta\mathbb{Z}^3$ with
\[
\left|v+y_{\beta}-\frac{l_c}{2}\cdot(1,1,1)\right|\le\frac{\eta}{2},
\]
$|\cdot|$ denoting the sup norm. It is admitted as soon as $\rho_{\beta}+\frac{\eta}{2}\le\frac{3}{8}\cdot l_c$, because each coordinate of a point of $v+y_{\beta}+[-\rho_{\beta},\rho_{\beta}]^3$ then differs from $\frac{l_c}{2}$ by at most $\rho_{\beta}+\frac{\eta}{2}$.

Let $v\in\eta\mathbb{Z}^3$. For a smooth function $g$ on $\mathbb{R}^3$, a point $x\in\mathbb{R}^3$ and a direction $j\in\{0,1,2\}$ let
\[
\delta_{\eta,j}g(x):=\frac{g\left(x+\frac{\eta}{2}e_j\right)-g\left(x-\frac{\eta}{2}e_j\right)}{\eta},\qquad J_{\eta}g(x):=\eta^{-3}\int_{\left[-\frac{\eta}{2},\frac{\eta}{2}\right]^3}\mathrm{d}y\,g(x+y).
\]
Let $d$ be the coboundary on the cochains of $\Lambda_{n_{*}}$, put $d_{\eta}:=\eta^{-1}d$, and let $d_{\eta}^{*}$ be the adjoint of $d_{\eta}$ with respect to the pairings $\eta^3\sum$ on the cochains of each degree. Let $\tau_v\beta$ be the translate $x\mapsto\beta(x-v)$, periodised with the period $R^{*}$ in each coordinate. For a site $x$ of $\Lambda_{n_{*}}$ and directions $\mu,\nu$ put
\[
u_{\mu\nu}(x):=\left(J_{\eta}(\tau_v\beta)_{\mu\nu}\right)\left(x+\frac{\eta}{2}\left(e_{\mu}+e_{\nu}\right)\right);
\]
this defines a two-cochain $u$ on $\Lambda_{n_{*}}$. Let
\[
a_{\eta}:=d_{\eta}^{*}u,\qquad\hat{\mathfrak{a}}:=\eta^{\frac{1}{2}}\cdot a_{\eta}\in C^1(\Lambda_{n_{*}}),\qquad\boldsymbol{\sigma}^2:=|d\hat{\mathfrak{a}}|_2^2,\qquad f:=dd^{*}\beta,
\]
where, in the definition of $f$, the operators $d$ and $d^{*}$ are the exterior derivative and its formal adjoint on forms on $\mathbb{R}^3$. Finally, for $c\in\mathrm{Cel}$ and $F\in F_N$ let
\[
\hat{\mathfrak{a}}_c:=(\mathsf{g}_c)_{*}\hat{\mathfrak{a}},\qquad X_c:=\left\langle d\hat{\mathfrak{a}}_c,F\right\rangle,
\]
with the push-forward of Section~\ref{sec:cells-reference-tori} and the counting scalar product of two-cochains, and let $X:=X_0$ be the variable of the base cell $\mathfrak{c}_0$.
\end{definition}

\paragraph{Two constants of the probe.}
Let
\[
C_{\beta}:=\sup|\hat{\mathfrak{a}}|_{\mathrm{ext}},\qquad\mathbf{D}:=\sup\boldsymbol{\sigma}^2,
\]
where $|\cdot|_{\mathrm{ext}}$ is the sum of norms of Lemma~\ref{lem:external-connections}, and where each supremum is formed over $N$, over the two values $a_{*}\in\left\{h'/L,\,h'/L^2\right\}$, over every admitted placement and over the forms $r_{*}\beta$ put in place of $\beta$, $r$ ranging over the diagonal sign matrices; here $r_{*}\beta$ is the push-forward of $\beta$ under the linear map $x\mapsto rx$. Both suprema are finite, by the bounds of Lemma~\ref{lem:properties-probes}(c). For an admitted placement, $|\hat{\mathfrak{a}}|_{\mathrm{ext}}$ is a function of $(\beta,N,a_{*})$ and $\boldsymbol{\sigma}^2$ is a function of $(\beta,N)$ alone; moreover $N\ge N_{\mathrm{min}}(R^{*},e)$ holds if and only if $N\ge0$ and $N+\log_La_{*}\ge1$, and the suprema are taken over the pairs $(N,a_{*})$ with these two properties. Consequently $C_{\beta}$ and $\mathbf{D}$ depend neither on $e$, nor on $R^{*}$, nor on the placement (Section~\ref{sec:probes-mirror-copies}).

\paragraph{Jet radius, variables and supports.}
Put
\[
t_x:=e_{*}^{3\varepsilon}\cdot\left(1+h'^{\frac{1}{2}}\cdot C_{\beta}\right)^{-1},\qquad\mathbf{r}:=\frac{t_x}{4},\qquad\hat{c}_c:=e^{-1}\cdot\hat{\mathfrak{a}}_c^{[k^{*}]}\quad\left(c\in\mathrm{Cel}\right);
\]
the number $\mathbf{r}$ is the jet radius, and $\hat{c}_c$ is a real bond function on $T_{k^{*}}$. For $t=(t_c)_{c\in\mathrm{Cel}}\in\mathbb{C}^{\mathrm{Cel}}$ and $\rho>0$ let
\[
a(t):=\sum_{c\in\mathrm{Cel}}\frac{t_c}{e}\cdot\hat{\mathfrak{a}}_c,\qquad\mathscr{C}(t):=\sum_{c\in\mathrm{Cel}}t_c\cdot\hat{c}_c,\qquad\mathbb{D}(\rho):=\left\{t\in\mathbb{C}^{\mathrm{Cel}}:\ |t_c|\le\rho\ \text{for every}\ c\in\mathrm{Cel}\right\};
\]
thus $\mathscr{C}(t)=a(t)^{[k^{*}]}$. For $c\in\mathrm{Cel}$ let $P_c$ be the set of the unit sites $y\in T^0_{k^{*}}$ for which the block $\Delta_y$ of Section~\ref{sec:lattice-geometry} contains an end point of a bond $b$ with $\hat{c}_c(b)\ne0$, and let
\[
P_c^{+}:=\left\{y\in T^0_{k^{*}}:\ d(y,P_c)\le M+1\right\},
\]
$d$ denoting the sup distance on the torus. Put
\[
\mathsf{p}:=\left(\frac{D_{\beta}}{a_{*}}+2M+4\right)^3,\qquad d_{\mathrm{sep}}:=\min_{c\ne c'}d\left(P_c,P_{c'}\right),
\]
with the convention $d_{\mathrm{sep}}:=\infty$ if $\hat{\mathfrak{a}}=0$. Condition~\eqref{eq:four-jets-sm2} requires that $d_{\mathrm{sep}}\ge6M+8$ and that the placement be admitted.

\paragraph{Separation and remainder constants.}
Put $p_{0*}:=p_{0,k^{*}}=v_{*}^{p_0}$, the value at the stop level $k^{*}=k^{*}(R^{*})$ of the number $p_{0,k}=v_k^{p_0}$ of Section~\ref{sec:running-quantities-tori}. With the constants $C_W$ and $\vartheta$ of Section~\ref{sec:constants-3} and the constant $C_Q$ of stage~\ref{it:order-choice-5}, let
\begin{align*}
\varpi_{\mathrm{sep}}&:=2^{N^{*}}\cdot C_W\cdot\mathsf{p}\cdot\exp\left(-\vartheta\cdot\left(d_{\mathrm{sep}}-4M-6\right)\right),\\
\varpi_0&:=18\,l_{*}^3\cdot\exp\left(-\frac{p_{0*}}{2}\right)+6\,l_{*}^3\cdot\exp\left(-\frac{(p_{*}-1)^2}{2C_0^2}\right),\\
z_2&:=\min\left\{\left(N^{*}\right)^{-\frac{1}{2}},\ \frac{1}{C_{\beta}},\ \frac{1}{\max(2,C_Q)\cdot\sqrt{N^{*}\cdot\mathbf{D}}}\right\},\qquad t_2:=e\cdot z_2,\\
\bar{\kappa}&:=6\,e_{*}^{\frac{1}{4}-9\varepsilon},\qquad\varpi_L:=\varpi_0^{\frac{1}{2}}\cdot\exp\left(\frac{\mathbf{D}}{2}\right),\qquad M_L:=2+\exp\left(\frac{\mathbf{D}}{2}\right),\\
\boldsymbol{\lambda}_L&:=\left(\varpi_L^{\frac{1}{3}}\cdot M_L^{\frac{2}{3}}\right)^{\frac{1}{N^{*}}}\cdot5\cdot\left(\frac{2N^{*}\cdot\mathbf{r}}{t_2}\right)^4.
\end{align*}
In the minimum that defines $z_2$, an entry whose denominator vanishes is read as $\infty$.

\paragraph{Generating functions.}
For $t\in\mathbb{C}^{\mathrm{Cel}}$ let
\[
\mathbf{Z}(t):=\bar{\nu}^{a(t)}(F_N),\qquad\mathring{\Phi}(t):=\frac{\mathbf{Z}(t)}{\mathbf{Z}(0)},
\]
where $\bar{\nu}^a$ is the measure of Section~\ref{sec:lattice-measures-laws}, formed on $\Lambda_{n_{*}}$ with the pair $P_N(e)$. The function $\mathbf{Z}$ is entire (Section~\ref{sec:averaged-law-partition}), and $\mathbf{Z}(0)>0$ by Lemma~\ref{lem:signs-symmetries-wilson}(a); hence $\mathring{\Phi}$ is an entire function on $\mathbb{C}^{\mathrm{Cel}}$. For $t\in\mathbb{C}$ let
\[
N_0(t):=\frac{\det\left(D_{(t/e)\cdot\hat{\mathfrak{a}}}+\bar{m}\cdot\eta\right)}{\det\left(D_0+\bar{m}\cdot\eta\right)},\qquad b_{\eta,R^{*}}(0):=\left.\frac{\mathrm{d}^4}{\mathrm{d}t^4}\log N_0(t)\right|_{t=0};
\]
the function $N_0$ is entire and $N_0(0)=1$ (Section~\ref{sec:averaged-law-partition}), so that $\log N_0$ is analytic in a neighbourhood of $t=0$.

\paragraph{Jets.}
Let
\[
\mathcal{J}:=\mathbb{C}\left[t_c:c\in\mathrm{Cel}\right]\big/\left(t_c^5:c\in\mathrm{Cel}\right).
\]
Every $P\in\mathcal{J}$ is the class of exactly one polynomial $\sum_{\mathbf{a}}P_{\mathbf{a}}\,t^{\mathbf{a}}$, in which $\mathbf{a}=(\mathbf{a}_c)_{c\in\mathrm{Cel}}$ ranges over $\{0,1,2,3,4\}^{\mathrm{Cel}}$ and $t^{\mathbf{a}}:=\prod_{c\in\mathrm{Cel}}t_c^{\mathbf{a}_c}$. For $\rho>0$ put
\[
\|P\|_{\rho}:=\sum_{\mathbf{a}}|P_{\mathbf{a}}|\cdot\rho^{|\mathbf{a}|},\qquad|\mathbf{a}|:=\sum_{c\in\mathrm{Cel}}\mathbf{a}_c;
\]
for a polynomial in one variable of degree at most four the number $\|\cdot\|_{\rho}$ is formed in the same way. A function that is analytic in a neighbourhood of $t=0$ in $\mathbb{C}^{\mathrm{Cel}}$ is identified with its image in $\mathcal{J}$, namely the class of the polynomial formed by those terms of its Taylor series at $t=0$ in which every variable has degree at most four.

\paragraph{The other constants, the smallness conditions and the number $\mathfrak{m}_{-}$.}
The other constants, among them $B_0$, are likewise those of Section~\ref{sec:constants-3}, and the smallness conditions are the conditions~\eqref{eq:four-jets-sm1}--\eqref{eq:four-jets-sm8} of Section~\ref{sec:smallness-conditions}. The running mass $\mathfrak{m}_{*}$ of the stop level does not occur in the definitions of Section~\ref{sec:constants-3}; wherever a mass enters them, it is the number
\[
\mathfrak{m}_{-}:=L^{-m_0-3},
\]
which satisfies $\mathfrak{m}_{-}<\mathfrak{m}_{*}$ (Lemma~\ref{lem:scales-stop-level}(a)). The number $\mathfrak{m}_{-}$ is a function of $L$ and of the integer $m_0$ of Definition~\ref{def:mass-bound-stop} alone; hence it does not depend on $e$. The running mass $\mathfrak{m}_{*}=\bar{m}\cdot a_{*}$, by contrast, may change with $e$, since $a_{*}$ is one of the two numbers $h'/L$ and $h'/L^2$ (Lemma~\ref{lem:scales-stop-level}(a)). The eight numbers $\boldsymbol{\theta}_{\mathrm{CT}}$, $\vartheta$, $g$, $h_G$, $g_1$, $\boldsymbol{\alpha}$, $c_1$ and $C_W$ of Section~\ref{sec:constants-3} are formed from $\mathfrak{m}_{-}$ and the constants of stages~\ref{it:order-choice-1}--\ref{it:order-choice-7} alone, and therefore they do not depend on $e$. Five further quantities of Section~\ref{sec:constants-3} are formed with some of these eight numbers, namely $\tau_s=c_1\cdot l_{*}$, $\bar{r}$ (which contains $h_G$), $B_1$ (which contains $g_1$), $B_0=C_{\mathrm{fr}}\cdot e_{*}^{4\varepsilon}+B_{\mathrm{wr}}+B_1$ and $\varpi_{\mathrm{sep}}$ (which contains $C_W$ and $\vartheta$); these five quantities do depend on $e$ (Section~\ref{sec:constants-3}). The charge $e$ enters them through $l_{*}$, $e_{*}$, $r_{*}$, the number $\mathsf{p}$ (that is, through $a_{*}$) and $d_{\mathrm{sep}}$, and the number $\mathfrak{m}_{-}$ adds no dependence on $e$ to them.

\begin{theorem}[The stop level of the reference torus]\label{thm:stop-level-reference}
Assume hypothesis~\ref{def:hypothesis-reference-torus}. Then the assertions~\eqref{it:trajectory-fermionic-operators-SL1}--\eqref{it:trajectory-fermionic-operators-SL6} of Proposition~\ref{prop:trajectory-fermionic-operators} hold. The forest formula of Lemma~\ref{lem:forest-formula-brydges}, for polynomials in variables indexed by the two-element subsets of $\{1,\dots,n\}$, $n\ge1$, holds, and so does the weighted tree expansion of Lemmas~\ref{lem:wick-rule-gram} and~\ref{lem:tree-expansion}. Moreover, for every admitted placement and every $t\in\mathbb{D}(t_x)$,
\[
a(t)^{[j]}\in\mathfrak{X}_j\qquad\text{for all }j\le k^{*}.
\]
\end{theorem}

The first and the last assertion are contained in Proposition~\ref{prop:trajectory-fermionic-operators}; the forest formula and the weighted tree expansion are established in Section~\ref{sec:six-auxiliary-lemmas}.

\begin{theorem}[Representation in the ring of jets]\label{thm:representation-ring-jets}
Assume hypothesis~\ref{def:hypothesis-reference-torus} and conditions~\eqref{eq:four-jets-sm1}--\eqref{eq:four-jets-sm7}. For every $N\ge N_{\mathrm{min}}(R^{*},e)$ there are a measure $\boldsymbol{\varphi}$ on $[-\pi,\pi]^3\times V_{\mathrm{t}}(T^0_{k^{*}})$, concentrated on the set of the pairs $(\theta,A_{*})$ with $\chi_{*}(A_{*})=1$, elements $k_S(\theta,A_{*};\cdot)$ of $\mathcal{J}$, indexed by the non-empty subsets $S$ of $\mathrm{Cel}$ and defined for $\theta\in\mathbb{R}^3$ and for $A_{*}\in V_{\mathrm{t}}(T^0_{k^{*}})$ with $\chi_{*}(A_{*})=1$ (Definition~\ref{def:coordinates-slice-small}, Proposition~\ref{prop:orthogonal-coordinates}), and an element $\mathrm{Lr}$ of $\mathcal{J}$, all of them constructed in Sections~\ref{sec:fields-outside-small} and~\ref{sec:ratio-partition-functions}, such that in $\mathcal{J}$
\[
\mathring{\Phi}(t)=\prod_{c\in\mathrm{Cel}}\phi_0(t_c)\cdot\int\boldsymbol{\varphi}(\mathrm{d}\theta\,\mathrm{d}A_{*})\,\prod_{\emptyset\ne S\subset\mathrm{Cel}}\left(1+k_S(\theta,A_{*};t)\right)+\mathrm{Lr}(t),
\]
where the integrand, a function with values in $\mathcal{J}$ that is defined $\boldsymbol{\varphi}$-almost everywhere, is integrated coefficientwise, and such that the following hold.
\begin{enumerate}
\item[(a)] $\phi_0$ is the four-jet of $N_0$ at $t=0$, that is, the Taylor polynomial of degree four of $N_0$ at $0$; it is even, and $\phi_0(0)=1$.
\item[(b)] $\boldsymbol{\varphi}$ is the measure of Section~\ref{sec:fields-outside-small}; its density with respect to $(2\pi)^{-3}\,\mathrm{d}\theta\,\mathrm{d}\mathsf{m}_{*}(A_{*})$ contains the factor $\chi_{*}(A_{*})$, so that $\boldsymbol{\varphi}$ is concentrated on the set where $\chi_{*}(A_{*})=1$. It is a positive measure of total mass $\bar{\mu}(\mathfrak{G})$, and $1-\varpi_0\le\bar{\mu}(\mathfrak{G})\le1$.
\item[(c)] For every non-empty $S\subset\mathrm{Cel}$ the coefficients of $k_S(\theta,A_{*};\cdot)$ are real and depend continuously on $(\theta,A_{*})$, and every monomial that occurs in $k_S(\theta,A_{*};\cdot)$ contains each of the variables $t_c$, $c\in S$, and no other variable.
\item[(d)] For every $c\in\mathrm{Cel}$,
\[
\sup_{\theta,A_{*}}\left\|k_{\{c\}}(\theta,A_{*};\cdot)\right\|_{\mathbf{r}}\le\bar{\kappa}.
\]
\item[(e)] For every $S\subset\mathrm{Cel}$ with $\#|S|\ge2$,
\[
\sup_{\theta,A_{*}}\left\|k_S(\theta,A_{*};\cdot)\right\|_{\mathbf{r}}\le2\,\varpi_{\mathrm{sep}}\le\bar{\kappa}^{N^{*}}.
\]
\item[(f)] The coefficients of $\mathrm{Lr}$ are real, and $\|\mathrm{Lr}\|_{\mathbf{r}}\le\boldsymbol{\lambda}_L^{N^{*}}$.
\end{enumerate}
In (d) and (e) the suprema are taken over the pairs $(\theta,A_{*})$ at which $k_S$ is defined. The bounds $\varpi_0$, $\bar{\kappa}$, $\bar{\kappa}^{N^{*}}$ and $\boldsymbol{\lambda}_L^{N^{*}}$ and the number $\mathbf{r}$ take the same values for all $N\ge N_{\mathrm{min}}(R^{*},e)$, all $\theta$ and all admitted placements, and they do not change when $\beta$ is replaced by $r_{*}\beta$, $r$ a diagonal sign matrix. For all these choices $\varpi_{\mathrm{sep}}$ is given by one and the same function of $d_{\mathrm{sep}}$, and
\[
\varpi_{\mathrm{sep}}\le2^{N^{*}}\cdot C_W\cdot\mathsf{p}\cdot\exp\left(-\vartheta\cdot\left(\frac{M\cdot r_{*}}{8L}-4M-7\right)\right).
\]
\end{theorem}

The theorem is contained in Theorem~\ref{thm:factorised-representation}.

\begin{theorem}[The free ratio]\label{thm:free-ratio}
Assume hypothesis~\ref{def:hypothesis-reference-torus} and conditions~\eqref{eq:four-jets-sm2} and~\eqref{eq:four-jets-sm4}. There is a function $\mathrm{lgN}$ on the disc $\left\{t\in\mathbb{C}:\ |t|\le t_x\right\}$ which is analytic and even there and real at real points, such that for $|t|\le t_x$
\[
N_0(t)=\exp\left(\mathrm{lgN}(t)\right),\qquad\left|\mathrm{lgN}(t)\right|\le B_0\le\frac{1}{4}.
\]
Let $\phi_0$ be the four-jet of $N_0$ at $t=0$, as in Theorem~\ref{thm:representation-ring-jets}(a), and put $c_m:=\phi_0^{(m)}(0)$ for $m\in\{1,2,3,4\}$. Then $c_1=c_3=0$, and
\[
|c_2|\le\frac{\mathbf{r}^{-2}}{16},\qquad|c_4|\le\frac{\mathbf{r}^{-4}}{16},\qquad\|\phi_0\|_{\mathbf{r}}\le2,\qquad c_4-3c_2^2=b_{\eta,R^{*}}(0).
\]
\end{theorem}

All these assertions are contained in Lemma~\ref{lem:free-wilson-ratio}.

\begin{theorem}[Products of centred derivatives]\label{thm:products-centred-derivatives}
Assume hypothesis~\ref{def:hypothesis-reference-torus} and conditions~\eqref{eq:four-jets-sm1}--\eqref{eq:four-jets-sm7}, as in Theorem~\ref{thm:representation-ring-jets}. Let $s$ be the number of stage~\ref{it:order-choice-8}, and let $c_1,\dots,c_4$ be the numbers of Theorem~\ref{thm:free-ratio}. For $m\in\{1,2,3,4\}$ put
\[
\Theta^{(0)}_m:=\left|\left(\prod_{c\in\mathrm{Cel}}\left(\partial_{t_c}^m-c_m\right)\mathring{\Phi}\right)(0)\right|^{\frac{1}{N^{*}}}.
\]
Then, for $m\in\{1,2,3,4\}$,
\[
\Theta^{(0)}_m\le m!\cdot\mathbf{r}^{-m}\cdot\left(2\exp(1)\cdot\bar{\kappa}^s+\boldsymbol{\lambda}_L\right),
\]
and $\Theta^{(0)}_m\le8\,m!\cdot\mathbf{r}^{-m}\cdot\bar{\kappa}^s$ whenever $\boldsymbol{\lambda}_L\le\bar{\kappa}^s$, an inequality that condition~\eqref{eq:four-jets-sm8} requires. If $\boldsymbol{\lambda}_L\le\bar{\kappa}^s$, then, with
\[
\mathfrak{e}:=s\cdot\left(\frac{1}{4}-9\varepsilon\right)-12\varepsilon,\qquad\bar{C}:=1+h'^{\frac{1}{2}}\cdot C_{\beta},
\]
one has
\begin{gather*}
\Theta^{(0)}_4\le3\cdot10^5\cdot\bar{C}^4\cdot e_{*}^{\mathfrak{e}},\qquad\left(1+|c_2|\right)\cdot\Theta^{(0)}_2\le3\cdot10^4\cdot\bar{C}^4\cdot e_{*}^{\mathfrak{e}},\\
|c_2|+|c_4|^{\frac{1}{2}}\le8\,\bar{C}^2\cdot e_{*}^{-6\varepsilon},\qquad e\cdot\left(\frac{h'}{L^2}\right)^{\frac{1}{2}}\le e_{*}\le e\cdot h'^{\frac{1}{2}}.
\end{gather*}
\end{theorem}

This is obtained by combining the results of Section~\ref{sec:numbers}.

\begin{theorem}[The smallness conditions for small charge]\label{thm:smallness-conditions-small}
Let the constants of stages~\ref{it:order-choice-1}--\ref{it:order-choice-10} be chosen as in hypothesis~\ref{def:hypothesis-flow}, let $0<\bar{m}\le1$, and let $\beta$ be a probe in the sense of Definition~\ref{def:probes}. There is a number
\[
e_{\sharp}\in\left(0,e_1\cdot h'^{-\frac{1}{2}}\right],
\]
depending only on the constants of stages~\ref{it:order-choice-1}--\ref{it:order-choice-10}, on $\bar{m}$ and on $\beta$, and on $\beta$ only through $C_{\beta}$, $\mathbf{D}$ and $D_{\beta}$, with the following property. Let $0<e\le e_{\sharp}$; then $e\cdot h'^{\frac{1}{2}}\le e_1$, so that hypothesis~\ref{def:hypothesis-reference-torus} holds for every $N\ge N_{\mathrm{min}}(R^{*},e)$, where $R^{*}=R^{*}(e)$. For every such $N$ the centred placement is admitted, one has
\[
2\rho_{\beta}+3\eta\le l_c,
\]
and the conditions~\eqref{eq:four-jets-sm1}--\eqref{eq:four-jets-sm8} hold for every admitted placement.
\end{theorem}

This is proved as Lemma~\ref{lem:threshold-charge}.

\begin{theorem}[The fourth cumulant of one probe]\label{thm:fourth-cumulant-one}
Assume hypothesis~\ref{def:hypothesis-reference-torus} and conditions~\eqref{eq:four-jets-sm1}--\eqref{eq:four-jets-sm7}, as in Theorem~\ref{thm:representation-ring-jets}, and let $N\ge N_{\mathrm{min}}(R^{*},e)$. Let $\kappa_4(X;\bar{\mu})$ be the fourth cumulant of the variable $X$ of Definition~\ref{def:probes} under the law $\bar{\mu}$ of Section~\ref{sec:lattice-measures-laws}, formed on $\Lambda_{n_{*}}$ with the pair $P_N(e)$. If $\boldsymbol{\lambda}_L^{N^{*}}\le\frac{1}{16}$, then
\[
\left|\kappa_4(X;\bar{\mu})-e^4\cdot b_{\eta,R^{*}}(0)\right|\le96\,e^4\cdot\mathbf{r}^{-4}\cdot\left(\bar{\kappa}+2\,\boldsymbol{\lambda}_L^{N^{*}}\right).
\]
\end{theorem}

The same statement reappears as Corollary~\ref{cor:fourth-cumulant-one}, with its proof.

\subsubsection{Free coefficients, bare mass shifts, four-jets and lattice probes}\label{sec:free-coefficients-bare}

Section~\ref{sec:wilson-determinant-lattice} contains the definitions of three notions: the lattice laws $P^{\mathrm{nat}}_{N,T}$, $\mathrm{nat}\in\{\mathrm{af},\mathrm{av}\}$; mass assignments; and limit points of the families $\left(P^{\mathrm{nat}}_{N,T}\right)$. The free coefficients are those of Section~\ref{sec:free-fourth-order}; they are defined as follows. Write $\gamma\cdot\partial:=\sum_{\mu}\gamma_{\mu}\partial_{\mu}$ and $\gamma\cdot a:=\sum_{\mu}a_{\mu}\gamma_{\mu}$, the Dirac matrices $\gamma_{\mu}$ being those of Section~\ref{sec:letters-norms-dirac}. For $\mathsf{M}>0$ let
\[
S_{\mathsf{M}}:=\left(\mathsf{M}+\gamma\cdot\partial\right)^{-1}\qquad\text{on }L^2\left(\mathbb{R}^3;\mathbb{C}^4\right),
\]
and for a real $a\in C_c^{\infty}\left(\mathbb{R}^3;\mathbb{R}^3\right)$ put
\[
q_{4,\mathsf{M}}(a):=-\frac{1}{4}\operatorname{Tr}\left[\left(S_{\mathsf{M}}\,i\gamma\cdot a\right)^4\right].
\]
For a real two-form $\beta\in C_c^{\infty}\left(\mathbb{R}^3;\Lambda^2\right)$, with $a:=d^{*}\beta$ and $f:=da$, let
\[
b(f):=24\,q_{4,\bar{m}}(a)=-6\operatorname{Tr}\left[\left(S_{\bar{m}}\,i\gamma\cdot a\right)^4\right].
\]
Let $R>0$. A function $\psi\in L^2_{\mathrm{loc}}\left(\mathbb{R}^3;\mathbb{C}^4\right)$ is called $R$-antiperiodic if $\psi\left(\cdot+R\,e_{\mu}\right)=-\psi$ for every direction $\mu$. On the space of the $R$-antiperiodic functions, normed by the $L^2$ norm over one period cube, let $S^R:=\left(\bar{m}+\gamma\cdot\partial\right)^{-1}$. Provided that the support of $a$ lies in the interior of a period cube, the letter $a$ is used as well for the $R$-periodic function $\sum_{z\in\mathbb{Z}^3}a\left(\cdot+R\,z\right)$, and one puts
\[
b_R(f):=-6\operatorname{Tr}\left[\left(S^R\,i\gamma\cdot a\right)^4\right].
\]
Finally, let $e>0$, let $\eta\in(0,1]$ be such that $n:=R/\eta$ is an integer with $n\ge2$, and put $q:=e\cdot\eta^{\frac{1}{2}}$. For a complex one-cochain $A$ on the torus $\Lambda_n$ of Section~\ref{sec:wilson-determinant-lattice} let $D_A:=D[qA]$, where $D[\cdot]$ is the operator of~\eqref{eq:free-coefficient-5}; for $\eta=L^{-N}$ and odd $n$ this is the Wilson operator $D_A$ of Section~\ref{sec:lattice-measures-laws}. Let $a_{\eta}$ be a one-cochain on the torus $\Lambda_{\eta,R}=\eta\mathbb{Z}^3/R\mathbb{Z}^3$ of Section~\ref{sec:lattices-tests-function}, regarded in lattice units as a one-cochain on $\Lambda_n$, let $\hat{\mathfrak{a}}:=\eta^{\frac{1}{2}}\cdot a_{\eta}$, and put, for $t\in\mathbb{C}$,
\[
N_0(t):=\frac{\det\left(D_{e^{-1}\cdot t\cdot\hat{\mathfrak{a}}}+\bar{m}\cdot\eta\right)}{\det\left(D_0+\bar{m}\cdot\eta\right)},\qquad b_{\eta,R}(0):=\left.\frac{\mathrm{d}^4}{\mathrm{d}t^4}\log N_0(t)\right|_{t=0}.
\]
As $q\cdot e^{-1}\cdot t\cdot\hat{\mathfrak{a}}=t\cdot\eta\cdot a_{\eta}$, the function $N_0$ and the number $b_{\eta,R}(0)$ are those of Section~\ref{sec:free-fourth-order}. When $\beta$ is a probe (Section~\ref{sec:probes-discretisation}), $a_{\eta}$ is the one-cochain of Section~\ref{sec:probes-discretisation} built on $\Lambda_{\eta,R}$ from $\beta$, or from a translate $\tau_v\beta=\beta(\cdot-v)$ with $v\in\eta\mathbb{Z}^3$. The cochain built from $\tau_v\beta$ is the translate by the lattice vector $v/\eta$ of the one built from $\beta$, since every step of the construction of Section~\ref{sec:probes-discretisation} commutes with the translations by vectors of $\eta\mathbb{Z}^3$; by Theorem~\ref{thm:free-coefficients-lattice}(d), $b_{\eta,R}(0)$ is therefore the same for all $v$. For $R=R^{*}$ and $\eta=L^{-N}$ the cochain $\hat{\mathfrak{a}}$ built from $\tau_v\beta$ is the cochain $\hat{\mathfrak{a}}$ of Definition~\ref{def:probes} formed with the same vector $v$.

\begin{theorem}[The bare mass shifts and the four-jets on the reference torus]\label{thm:bare-mass-shifts}
Let $0<\bar{m}\le1$. For $e>0$ with $e\cdot h'^{\frac{1}{2}}\le e_1$ define
\begin{align*}
T^{*}(e)&:=\log_L\left(\frac{R^{*}(e)}{h}\right)=m-m_0-1+\log_Lr_K(e),\qquad N_0(e):=N_{\mathrm{min}}\left(R^{*}(e),e\right),\\
m_{0,N}(e)&:=\begin{cases}\text{the first entry of }P_N(e)&\text{if }N\ge N_0(e),\\ 0&\text{if }0\le N<N_0(e),\end{cases}
\end{align*}
and put
\begin{gather*}
C_{\mathrm{mat}}:=\left(L^3+1\right)\cdot\left(\varepsilon_5\left(L^{-m_0-1}\right)+1\right),\qquad c_m:=\phi_0^{(m)}(0)\quad(m=1,2,3,4),\\
\mathfrak{e}:=s\cdot\left(\frac{1}{4}-9\varepsilon\right)-12\varepsilon,\qquad\mathfrak{e}_2:=6\varepsilon,\\
e_G=e_G(\bar{m}):=\sup\left\{e>0:\ e\cdot h'^{\frac{1}{2}}\le e_1\ \text{and}\ e\cdot h'^{\frac{1}{2}}\cdot\log h'\le1\right\}.
\end{gather*}
Here $e_1$ is the constant of stage~\ref{it:order-choice-10}; $h$, $h'$ and $r_K(e)$ are the numbers of Section~\ref{sec:running-quantities-tori}; $R^{*}(e)=R_{\mathrm{top}}(e)/L$ is the side of the reference torus and $N_{\mathrm{min}}$ the threshold of Definition~\ref{def:tori-stop-level}; $m$ is the exponent in $M=L^m$ (stage~\ref{it:order-choice-6}) and $m_0$ the integer of Definition~\ref{def:mass-bound-stop}; $P_N(e)$ is the pair of Definition~\ref{def:tuned-pair}; $\varepsilon_5$ is the function of Section~\ref{sec:conditions}; $\phi_0$ is the four-jet at $t=0$ of the free ratio $N_0(t)$ (Theorem~\ref{thm:representation-ring-jets}(a)); and $s$, $\varepsilon$ and $\sigma$ are the numbers of stage~\ref{it:order-choice-8}. Then the following assertions hold.
\begin{enumerate}
\item\label{it:bare-mass-shifts-1} Let $e>0$ with $e\cdot h'^{\frac{1}{2}}\le e_1$. Then $e\le\frac{1}{3}$. The number $T^{*}(e)$ is an integer with $T^{*}(e)\ge1$, and $R^{*}(e)=h\cdot L^{T^{*}(e)}$. For every integer $N\ge0$ the number $m_{0,N}(e)$ is real, with $\left|m_{0,N}(e)\right|<\frac{1}{8}$. These numbers are functions of $N$, $e$, $\bar{m}$ and of the scheme in the sense of Definition~\ref{def:terminology}(10), and they are independent of the torus. The length $l_c=R^{*}(e)/(2L)$ of Section~\ref{sec:cells-reference-tori} tends to infinity as $e\to0$. For every integer $N\ge N_0(e)$ the number $D^{\mathrm{mat}}$ of Theorem~\ref{thm:partition-functions-tori}, formed on the reference torus $R^{*}=R^{*}(e)$ with the lattice mass $\boldsymbol{\mu}_N=\bar{m}\cdot\eta+m_{0,N}(e)$, satisfies
\[
D^{\mathrm{mat}}\le C_{\mathrm{mat}}.
\]
If moreover $e\le e_G$, then condition~\eqref{eq:renormalisation-group-17} holds for $R^{*}=R^{*}(e)$.
\item\label{it:bare-mass-shifts-2} One has $0<\mathfrak{e}\le\frac{1}{4}$, $0\le\mathfrak{e}_2<1$ and $\sigma>0$. For every probe $\beta$, that is, every real $\beta\in C_c^{\infty}\left(\mathbb{R}^3;\Lambda^2\right)$ with $\beta\ne0$ (Definition~\ref{def:probes}), there are numbers $e_1(\beta)\in\left(0,e_G\right]$ and $C^G_{\beta}<\infty$, whose values are determined by $\beta$, $\bar{m}$ and the constants, with the following property. Let $0<e\le e_1(\beta)$ (then $e\cdot h'^{\frac{1}{2}}\le e_1$, by the definition of $e_G$), let $N\ge N_0(e)$ be an integer, and let the objects of Section~\ref{sec:probes-jets-reference} be formed on the reference torus $R^{*}=R^{*}(e)$ with the lattice mass $\boldsymbol{\mu}_N=\bar{m}\cdot\eta+m_{0,N}(e)$. Then
\[
2\rho_{\beta}+3\eta\le l_c,
\]
and, for the centred placement of Definition~\ref{def:probes}, with which $\phi_0$, the numbers $c_2$, $c_4$ and the numbers $\Theta^{(0)}_2$, $\Theta^{(0)}_4$ of Theorem~\ref{thm:products-centred-derivatives} are formed, the following hold.
\begin{enumerate}
\item[(a)] One has
\[
\Theta^{(0)}_4\le C^G_{\beta}\cdot e^{\mathfrak{e}},\qquad\left(1+|c_2|\right)\cdot\Theta^{(0)}_2\le C^G_{\beta}\cdot e^{\mathfrak{e}}.
\]
\item[(b)] One has
\[
|c_2|+|c_4|^{\frac{1}{2}}\le C^G_{\beta}\cdot e^{-\mathfrak{e}_2}.
\]
\item[(c)] One has $c_4-3c_2^2=b_{\eta,R^{*}(e)}(0)$, and
\[
\left|c_4-3c_2^2-b\left(d\,d^{*}\beta\right)\right|\le C^G_{\beta}\cdot\left(e+\eta^{\frac{1}{2}}\right).
\]
\end{enumerate}
\end{enumerate}
\end{theorem}

Both assertions are contained in Theorem~\ref{thm:free-coefficient-3}.

\begin{theorem}[Free coefficients, lattice probes and symmetries of the determinant]\label{thm:free-coefficients-lattice}
\begin{enumerate}
\item[(a)] \emph{Large mass.} Each real $a\in C_c^{\infty}\left(\mathbb{R}^3;\mathbb{R}^3\right)$ admits a constant $C_a<\infty$ for which
\[
\left|q_{4,\mathsf{M}}(a)-\frac{1}{240\pi\cdot\mathsf{M}^5}\int_{\mathbb{R}^3}\mathrm{d}x\,\left|\operatorname{curl}a(x)\right|_2^4\right|\le C_a\cdot\mathsf{M}^{-7}\qquad\text{whenever }\mathsf{M}\ge1,
\]
$|\cdot|_2$ denoting the Euclidean norm on $\mathbb{R}^3$. Let now $a$ be the field of Theorem~\ref{thm:free-coefficient-1}(a), let $C_a$ be the number that the preceding assertion attaches to this field, let $\beta^{(l)}$, $l\ge1$, be the probes of Theorem~\ref{thm:free-coefficient-1}(a), and put
\[
f_l:=dd^{*}\beta^{(l)},\qquad I_{\beta}:=\int_{\mathbb{R}^3}\mathrm{d}x\,\left|\operatorname{curl}a(x)\right|_2^4.
\]
Then $I_{\beta}>0$, and $b(f_l)>0$ for every $l\ge1$ such that
\[
\bar{m}\cdot l\ge1\qquad\text{and}\qquad\left(\bar{m}\cdot l\right)^2>\frac{240\pi\cdot C_a}{I_{\beta}}.
\]
\item[(b)] \emph{Lattice probes.} Let $0<\eta\le1$, let $R>0$ be such that $R/\eta$ is an integer, and let $\beta$ be a probe (Definition~\ref{def:probes}) whose cube $Q_{\beta}=y_{\beta}+\left[-\rho_{\beta},\rho_{\beta}\right]^3$ (Section~\ref{sec:probes-discretisation}) lies in the interior of a period cube of the torus $\Lambda_{\eta,R}$ of Section~\ref{sec:lattices-tests-function}; points of $\Lambda_{\eta,R}$ are represented in this period cube. Let $f:=dd^{*}\beta$, and let $f^{(\eta)}$ and $a_{\eta}\in C^1(\Lambda_{\eta,R})$ be formed from $\beta$ as in Section~\ref{sec:probes-discretisation}, with $d_{\eta}=\eta^{-1}d$. Then $a_{\eta}$ is orthogonal to $\ker d_{\eta}$ with respect to the pairing $\langle\cdot,\cdot\rangle_{\eta}$ of Section~\ref{sec:lattices-tests-function}, and each bond at which $a_{\eta}$ does not vanish has its two end points in the set
\[
\operatorname{supp}\beta+\left[-\frac{3\eta}{2},\frac{3\eta}{2}\right]^3.
\]
Let $\left(f^{(\eta)}\right)_{\eta}$ be the two-cochain to which the test map~\eqref{eq:continuum-infinite-2} carries $f^{(\eta)}$, and let $\|\cdot\|_{H^s}$ be the Sobolev norms of Section~\ref{sec:lattices-tests-function}. One has
\[
\left(f^{(\eta)}\right)_{\eta}=d_{\eta}a_{\eta},\qquad\left\|f^{(\eta)}-f\right\|_{H^s}\le\eta^2\cdot\|\beta\|_{H^{s+4}}\quad\text{for every real }s\ge0.
\]
Moreover, the assertions (b2), (b4) and (b5) of Theorem~\ref{thm:free-coefficient-1} hold. Finally, let $\hat{\mathfrak{a}}$ and $\boldsymbol{\sigma}^2$ be the cochain and the number that Definition~\ref{def:probes} attaches to the probe $\beta$, and let $|\cdot|_{\mathrm{ext}}$ be the norm of Lemma~\ref{lem:external-connections}. Then the numbers $C_{\beta}$ and $\mathbf{D}$ of Section~\ref{sec:probes-jets-reference} are finite:
\[
C_{\beta}=\sup|\hat{\mathfrak{a}}|_{\mathrm{ext}}<\infty,\qquad\mathbf{D}=\sup\boldsymbol{\sigma}^2<\infty,
\]
where, as in Section~\ref{sec:probes-jets-reference}, each supremum is taken over the pairs $(N,a_{*})$ of an integer $N\ge0$ and a number $a_{*}\in\left\{h'/L,\,h'/L^2\right\}$ with $N+\log_La_{*}\ge1$, over the admitted placements, and over the forms $r_{*}\beta$ put in place of $\beta$, $r$ ranging over the diagonal sign matrices.
\item[(c)] \emph{Lattice, torus and $\mathbb{R}^3$.} Let $\beta$ be a probe, let $f:=dd^{*}\beta$, and let $R\ge D_{\beta}+1$, where $D_{\beta}=2\rho_{\beta}$ is the number of Definition~\ref{def:probes} (so that the cube $Q_{\beta}$, whose side $D_{\beta}$ is smaller than $R$, lies in the interior of a period cube of the torus of side $R$). There are numbers $C_1,C_2<\infty$, depending only on $\bar{m}$ and $\beta$, such that
\[
\left|b_{\eta,R}(0)-b_R(f)\right|\le C_1\cdot\eta^{\frac{1}{2}}
\]
for every $\eta\in(0,1]$ such that $R/\eta$ is an integer, the number $b_{\eta,R}(0)$ being formed as in Section~\ref{sec:free-fourth-order} with the cochain $a_{\eta}$ of~(b), and such that
\[
\left|b_R(f)-b(f)\right|\le C_2\cdot\exp\left(-\bar{m}\cdot\left(R-D_{\beta}\right)\right).
\]
Whenever the numbers $R\cdot L^N$ are integers from some $N$ on, which holds for every $R=h\cdot L^T$, the first of these bounds may be applied with $\eta=L^{-N}$ for all large $N$, and it yields
\[
b_R(f)=\lim_{N\to\infty}b_{L^{-N},R}(0).
\]
\item[(d)] \emph{Symmetries of the Wilson determinant.} Let $n=(n_{\mu})_{\mu}$ be a triple of integers $n_{\mu}\ge2$, indexed by the three directions $\mu$, let $\Lambda_n$ be the torus of Section~\ref{sec:wilson-determinant-lattice}, let $\boldsymbol{\mu}$ be a real number, and for a complex one-cochain $B$ on $\Lambda_n$ put
\[
\mathrm{Dt}_n(B):=\det\left(D[B]+\boldsymbol{\mu}\right),
\]
$D[B]$ being the operator of~\eqref{eq:free-coefficient-5}. Let $r$ be an element of the cubic group $O_{\mathrm{cub}}=O(3)\cap\mathrm{GL}(3,\mathbb{Z})$, the group of the orthogonal $3\times3$ matrices with integer entries (Section~\ref{sec:lattices-tests-function}), let $c\in\mathbb{Z}^3$, and let $n'=(n'_{\mu})_{\mu}$ be a triple of integers $n'_{\mu}\ge2$ with $r\left(\prod_{\mu}n_{\mu}\mathbb{Z}\right)=\prod_{\mu}n'_{\mu}\mathbb{Z}$; the map $g\colon x\mapsto rx+c$ induces a lattice isometry of $\Lambda_n$ onto $\Lambda_{n'}$, which carries a one-cochain $B$ on $\Lambda_n$ to the one-cochain $g_{*}B$ on $\Lambda_{n'}$ with $\left(g_{*}B\right)(gx,gx'):=B(x,x')$. Then, for every complex one-cochain $B$, every complex zero-cochain $\boldsymbol{\lambda}$ and every integer one-cochain $\mathfrak{z}$ on $\Lambda_n$,
\begin{gather*}
\mathrm{Dt}_n\left(B+d\boldsymbol{\lambda}\right)=\mathrm{Dt}_n(B),\qquad\mathrm{Dt}_n\left(B+2\pi\cdot\mathfrak{z}\right)=\mathrm{Dt}_n(B),\\
\mathrm{Dt}_n(-B)=\mathrm{Dt}_n(B),\qquad\mathrm{Dt}_{n'}\left(g_{*}B\right)=\mathrm{Dt}_n(B),
\end{gather*}
and $\mathrm{Dt}_n(B)\ge0$ for every real $B$. Consequently, if $\boldsymbol{\mu}>-1$, then on a cubic torus the lattice translations, the elements of the cubic group $O_{\mathrm{cub}}$ and the sign change $F\mapsto-F$ leave both the all-flux law~\eqref{eq:setup-af} and the harmonic-averaged law~\eqref{eq:setup-av} invariant; for such $\boldsymbol{\mu}$ the numbers $Z_n$ and $\mathbf{Z}_n$ by which these two laws are normalised are positive (Section~\ref{sec:wilson-determinant-lattice}). Moreover, $g$ carries the sectors of fixed flux of~\eqref{eq:setup-af} on $\Lambda_n$ to those on $\Lambda_{n'}$ (Theorem~\ref{thm:free-coefficient-1}(d5)), in the following sense. Since $r$ is an orthogonal matrix with integer entries, there are a permutation $\mu\mapsto\mu'$ of the directions and signs $\epsilon_{\mu}\in\{1,-1\}$ such that $re_{\mu}=\epsilon_{\mu}\cdot e_{\mu'}$ for every direction $\mu$. Each $w=(w_{\mu\nu})_{\mu<\nu}\in\mathbb{Z}^3$ is extended to all ordered pairs of distinct directions by $w_{\nu\mu}:=-w_{\mu\nu}$, and $g\cdot w\in\mathbb{Z}^3$ is defined by requiring of its extension that
\[
(g\cdot w)_{\mu'\nu'}=\epsilon_{\mu}\cdot\epsilon_{\nu}\cdot w_{\mu\nu}\qquad\text{for all distinct directions }\mu,\nu.
\]
The map $w\mapsto g\cdot w$ is a bijection of $\mathbb{Z}^3$ onto itself, and
\[
\boldsymbol{\omega}_{n'}^{(g\cdot w)}\left(\Phi\circ g^{*}\right)=\boldsymbol{\omega}_n^{(w)}(\Phi)\qquad\text{for every }w\in\mathbb{Z}^3\text{ and every }\Phi\in\mathfrak{F}(n).
\]
Here $\mathfrak{F}(n)$ is the class of functions on $C^2(n)$ introduced in Section~\ref{sec:wilson-determinant-lattice}; $\boldsymbol{\omega}_n^{(w)}$ is the functional of flux $w$ on $\Lambda_n$ defined there, formed with any $q>0$ and with the number $\boldsymbol{\mu}$ above, and $\boldsymbol{\omega}_{n'}^{(g\cdot w)}$ is formed in the same way on $\Lambda_{n'}$, with the same $q$ and $\boldsymbol{\mu}$; and $g^{*}$ is the pull-back of two-cochains from $\Lambda_{n'}$ to $\Lambda_n$. Finally, take $n=R/\eta$, so that $\Lambda_n$ is the cubic torus on which $b_{\eta,R}(0)$ is formed. Then $n'=n$ (as the entries of $n'$ are positive and $re_{\mu}=\epsilon_{\mu}\cdot e_{\mu'}$, the condition imposed on $n'$ amounts to $n'_{\mu'}=n_{\mu}$ for every direction $\mu$), and the number $b_{\eta,R}(0)$ is the same whether it is computed from $a_{\eta}$ or from $g_{*}a_{\eta}$.
\end{enumerate}
\end{theorem}

Part~(a), the assertions of~(b) other than the finiteness of $C_{\beta}$ and $\mathbf{D}$, and parts~(c) and~(d) are contained in Theorem~\ref{thm:free-coefficient-1}. The finiteness of $C_{\beta}$ and $\mathbf{D}$ is a consequence of Lemma~\ref{lem:properties-probes}(c).

\subsection{The constants}\label{sec:constants}

\subsubsection{The order of choice}\label{sec:order-choice}

\begin{definition}[Order of choice]\label{def:order-choice}
The constants, as well as the number $\bar{m}$, the probe $\beta$ and the point $x_0$ of stage~\ref{it:order-choice-11} and the charge $e$, the side $R$ and the integer $N$ of stage~\ref{it:order-choice-12}, which are not constants (Definition~\ref{def:constants}), are chosen in the twelve stages of Section~\ref{sec:parameters-scales}, in the order of the following list: for $n=1,\dots,12$, the $n$-th item is stage~$n$, and it gives the quantities chosen at that stage together with the conditions imposed on them. What is chosen at a stage depends only on what has been chosen at the earlier stages.
\begin{enumerate}
\item\label{it:order-choice-1} Absolute constants: the number $\bar{\varepsilon}:=\frac{1}{200}$; the constants $c_X$, $\kappa_0$, $K_0$ and $c_{\mathrm{rb}}$ of Proposition~\ref{prop:polymers}; the constant $c_{\mathrm{fl}}\ge1$ of Propositions~\ref{prop:symmetries-extraction} and~\ref{prop:contraction}; and the number $b:=1$. The values of $c_X$, $\kappa_0$, $K_0$, $c_{\mathrm{rb}}$ and $c_{\mathrm{fl}}$, fixed in Section~\ref{sec:absolute-constants}, are
\[
c_X=54,\qquad\kappa_0=54\log72,\qquad K_0=\frac{\exp(\kappa_0)}{36},\qquad c_{\mathrm{rb}}=c_X\cdot K_0,\qquad c_{\mathrm{fl}}=61\,c_X\cdot K_0.
\]
\item\label{it:order-choice-2} The H\"older order $\alpha\in\left(1-\bar{\varepsilon},1\right)$.
\item\label{it:order-choice-3} A number $\kappa$ with $\kappa\ge7\kappa_0+7$.
\item\label{it:order-choice-4} An odd integer $L$ such that the inequalities
\begin{gather*}
L^{\frac{1}{8}-\frac{3\bar{\varepsilon}}{2}}\ge2400,\qquad c_{\mathrm{fl}}\cdot L^{-\frac{3}{8}+6\bar{\varepsilon}}\le\frac{1}{16},\\
L\cdot\left(\kappa-6\kappa_0-4\right)-1\ge\kappa,\qquad c_{\mathrm{rb}}\cdot L^3\cdot\exp\left(-\left(L\cdot\left(\kappa-\kappa_0-1\right)-\kappa\right)\right)\le\frac{1}{2}
\end{gather*}
hold for every $\varepsilon\in\left(0,\bar{\varepsilon}\right]$, and such that $L\ge L^F$ and $L\ge L^P$, where $L^F=L^P=5$ are the thresholds of Theorems~\ref{thm:fermion-operators-weakened} and~\ref{thm:photon-operators-weakened}. For every $\varepsilon\in\left(0,\bar{\varepsilon}\right]$ the first of the four displayed inequalities implies
\[
L\ge4^8,\qquad L^{-\frac{5}{8}}\le\frac{1}{4},\qquad L^{-3+\frac{1}{8}}\le\frac{1}{2},\qquad13.6\le L^{\frac{1}{4}-\frac{3\varepsilon}{2}},\qquad L^{-\frac{3}{4}}\le\frac{1}{2}.
\]
\item\label{it:order-choice-5} Functions of $(b,\alpha,\kappa,L)$: the numbers $b_{-}$ and $b_{+}$ of Section~\ref{sec:running-quantities-tori}; the numbers
\begin{gather*}
N^{*}=8L^3,\qquad a_L:=\frac{L^3-1}{(L-1)^3},\\
K_L:=\left(1+L+L^2\right)\cdot\left(\log L+L\cdot G_L\right),\quad\text{where}\quad G_L:=\left(1+\frac{L}{3.2}\right)\cdot\frac{2L^2}{2.56}
\end{gather*}
(Sections~\ref{sec:cells-reference-tori} and~\ref{sec:numbers-attached-reference}); the integers $N_c$, $c_l$ and $d_U$ of Definition~\ref{def:parameter-cubes-polydisks}; the constants $c_C$, $C_Q$ and $C_T$ of Propositions~\ref{prop:basis-map-box}, \ref{prop:block-averages} and~\ref{prop:left-inverse}; the number $\mathfrak{m}_0$, the integer $m_0$ and the scale $h'$ of Definition~\ref{def:mass-bound-stop} (of $h'$ only the defining formula belongs to this stage, since the scale $h$, and with it $h'$, is fixed together with $\bar{m}$ in stage~\ref{it:order-choice-11}); the numbers $\alpha^F$ of Theorem~\ref{thm:fermion-operators-weakened}, $\alpha^P$ of Theorem~\ref{thm:photon-operators-weakened} and $\alpha_0$ of Section~\ref{sec:photon-operators}; the numbers $\gamma^F$ of Theorem~\ref{thm:fermion-operators-weakened}, $\gamma_P$ of Proposition~\ref{prop:kernel-bounds}, $\gamma^P$ of Theorem~\ref{thm:photon-operators-weakened} and $\gamma_w$ of Section~\ref{sec:photon-operators}; the constants $C_P$ and $c_P$ of Propositions~\ref{prop:kernel-bounds} and~\ref{prop:spectral-bounds-decay}; and
\[
\kappa_w:=\kappa+\kappa_0+3,\qquad\hat{\kappa}:=N_c\cdot\kappa+\log\left(6c_X\right).
\]
\item\label{it:order-choice-6} A power $M=L^m$ of $L$, with an integer $m$, such that
\begin{gather*}
M\ge L^2,\qquad M>4L,\qquad\gamma_w\cdot M\ge16,\qquad M^{\alpha_0}\ge1+\exp(\hat{\kappa}),\\
M\ge M^F\left(b,\alpha,L,\alpha_0,\kappa_w+\kappa_0\right),\qquad M\ge M^P\left(\alpha,L,\alpha_0\right),\qquad M\ge L^{m_0+2},
\end{gather*}
where $M^F$ is the number defined in condition~\eqref{eq:renormalisation-group-19} of Section~\ref{sec:conditions} and $M^P$ is the number~\eqref{eq:operators-photon-20}. None of these conditions on $M$ involves $C_w$, $\varepsilon$ or $e$.
\item\label{it:order-choice-7} Functions of the constants chosen so far: the constants $C^F$ of Theorem~\ref{thm:fermion-operators-weakened} and $C^P$ of Theorem~\ref{thm:photon-operators-weakened}, the constants $C_w^0$ and $C_w^1$ of Section~\ref{sec:photon-operators}, and the constant $C_w$ of Section~\ref{sec:primitive-operators}, which depend on $\kappa_w$; the parameter sets $\mathrm{par}^0$ and $\mathrm{par}$ of Section~\ref{sec:primitive-operators}; the number
\[
c_{\gamma}:=\sup_x\sum_{x'}\exp\left(-\gamma_w\cdot d(x,x')\right),
\]
where the supremum is taken over all levels $k\ge0$, all sides $l\in\mathrm{Cl}$ and all indices $x$, and where $x'$ ranges over the unit sites with a spinor index, over the unit bonds, or over the basis $\{\Upsilon\}$; the number $c_{\gamma}$ depends on $(\gamma_w,L)$ only (Section~\ref{sec:notation}); the constants $C_B$ of Proposition~\ref{prop:local-regularity} and $C_d$ of Proposition~\ref{prop:local-expansion-determinant}; the number
\[
c_{\kappa}:=\sup_{k,\,l,\,\square}\ \sum_{X\in D_k,\ X\supseteq\square}\exp\left(-\kappa\cdot d_M(X)\right),
\]
the supremum being taken over all levels $k\ge0$, all sides $l\in\mathrm{Cl}$ and all $M$-cubes $\square$; it satisfies $c_{\kappa}\le K_0$, by Proposition~\ref{prop:polymers}(t2) and $\kappa\ge\kappa_0$; the constant $C_0$, the largest of the constants that occur in Proposition~\ref{prop:kernel-bounds} (summed over the second index), Proposition~\ref{prop:block-averages}, Proposition~\ref{prop:bounds} and Proposition~\ref{prop:bounds-unit-lattice}, that is,
\[
C_0:=\max\left\{C_P\cdot c_{\mathrm{bd}}\left(\gamma_P\right),\ C^F,\ 4b_{+}+16\,C^F\cdot c\left(\gamma^F\right)\right\},
\]
where $c(g):=\sum_{z\in\mathbb{Z}^3}\exp\left(-g\cdot|z|\right)$ and $c_{\mathrm{bd}}(g):=3\exp(3g)\cdot c(g)$ for $g>0$ (Section~\ref{sec:constants-conditions-torus}); the constants $\gamma_1$ and $C_{\Sigma}$ of Lemma~\ref{lem:dependence-twist-uniform}, for the interval $\left[b_{-},b_{+}\right]$; and
\[
c_S:=10^5\,c_X^5\cdot K_0^5\cdot\exp(6\kappa).
\]
\item\label{it:order-choice-8} Integers $r\ge2$, $p_0>3r+2$ and $p>2p_0$; a number $\varepsilon\in\left(1-\alpha,\bar{\varepsilon}\right]$; and an \emph{admissible} number $s\in(0,1)$, that is, a number $s\in(0,1)$ with
\[
\sigma:=s\cdot\left(\frac{1}{4}-9\varepsilon\right)-12\varepsilon-\frac{2a_L}{N^{*}}>0.
\]
\item\label{it:order-choice-9} A number $e_0>0$ that satisfies every condition of Section~\ref{sec:conditions} in which $e_0$ occurs, namely the conditions~\eqref{eq:renormalisation-group-A1}--\eqref{eq:renormalisation-group-A10}, the conditions~\eqref{eq:renormalisation-group-20}--\eqref{eq:renormalisation-group-24} with the constant $C_w$ of $\mathrm{par}$, and condition~\eqref{eq:renormalisation-group-25}.
\item\label{it:order-choice-10} A number
\[
e_1\in\left(0,\min\left(e_0,\frac{1}{3}\,L^{-\frac{m_0}{2}}\right)\right]
\]
with the following property (Theorem~\ref{thm:scales-stop-levels}): if $e>0$ and $e_K=e\cdot h'^{\frac{1}{2}}\le e_1$, then every $R\in\mathrm{Tor}(e)$ satisfies, for all $N\ge N_{\mathrm{min}}(R,e)$, the conditions~\eqref{eq:renormalisation-group-16}, \eqref{eq:renormalisation-group-fl} and~\eqref{eq:renormalisation-group-T1}--\eqref{eq:renormalisation-group-T6}, and $\left(1+3\mathrm{lg}/v_K\right)^r<L$. In terms of $e_K$ and $h'$, condition~\eqref{eq:renormalisation-group-17} reads
\[
6N^{*}\cdot\frac{M\cdot r_K}{L}\cdot e_K^2\cdot\left(v_K+\frac{1}{2}\log h'\right)\le\pi^2,
\]
since $R^{*}=h'\cdot M\cdot r_K/L$ (Definition~\ref{def:tori-stop-level}) and $\log\frac{1}{e}=v_K+\frac{1}{2}\log h'$. This inequality holds if $e_K\le e_1$ and $e_K\cdot\log h'\le1$, the second requirement being condition~\eqref{eq:renormalisation-group-18}; in particular, it holds if $e_K\le e_1$ and $e\cdot h'\le1$ (if $h'\ge1$, then $\log h'\le h'^{\frac{1}{2}}$, so that $e_K\cdot\log h'\le e_K\cdot h'^{\frac{1}{2}}=e\cdot h'$; if $h'<1$, then $\log h'<0$). Condition~\eqref{eq:renormalisation-group-18} belongs to stage~\ref{it:order-choice-11}: in Theorem~\ref{thm:bare-mass-shifts} it is part of the definition of $e_G$, and in Theorem~\ref{thm:main-theorem} it holds because $e_2\le e_G$. If $e_K\le e_1$, then $e\le e_1\cdot L^{\frac{m_0}{2}}\le\frac{1}{3}$, because $h\ge1$. Ranges of the charge are always expressed by the inequality $e\cdot h'^{\frac{1}{2}}\le e_1$; the number $e_1$ does not depend on $\bar{m}$. It is defined in Section~\ref{sec:constants-conditions-torus} by the inequalities~\eqref{eq:renormalisation-flow-e1}--\eqref{eq:renormalisation-flow-e8}.
\item\label{it:order-choice-11} The number $\bar{m}\in(0,1]$, and with it the scales $h$, $h'$ and the number $e_G=e_G(\bar{m})$ of Theorem~\ref{thm:bare-mass-shifts}; the probe $\beta$ (Definition~\ref{def:probes}) and the point $x_0\in\mathbb{R}^3$ of Theorem~\ref{thm:main-theorem}; then the numbers $C_{\beta}$ and $\mathbf{D}$ of Section~\ref{sec:probes-jets-reference}, the numbers $D_{\beta}$ and $\rho_{\beta}$ of Definition~\ref{def:probes}, and the constants $C_1$, $C_2$ of Theorem~\ref{thm:free-coefficients-lattice}(c); then, in this order, $e_{\sharp}$ (Theorem~\ref{thm:smallness-conditions-small}) and $C^G_{\beta}$ (Theorem~\ref{thm:bare-mass-shifts}); $C_{\mathrm{Gam}}$ and $C_Y$ (Step~1 in Section~\ref{sec:proof-part-ii-2}); $C_{\mathrm{kap}}$; and $e_2$ (Theorem~\ref{thm:main-theorem}).
\item\label{it:order-choice-12} The charge $e$; then the side $R$ of the torus; then $N$.
\end{enumerate}
\end{definition}

\begin{definition}[Mass bound and stop scale]\label{def:mass-bound-stop}
Let $\mathfrak{m}_0\in(0,1]$ be the number of Theorem~\ref{thm:fermion-operators-weakened}; its value is $\mathfrak{m}_0=1$. Let $m_0$ be the least integer $j\ge0$ with $L^{-j}\le\mathfrak{m}_0$, and put
\[
h':=h\cdot L^{-m_0},
\]
where $h$ is the scale of Section~\ref{sec:running-quantities-tori}.
\end{definition}

The H\"older order $\alpha$ is chosen before stage~\ref{it:order-choice-5} because the constants of Propositions~\ref{prop:decay-fine-propagator} and~\ref{prop:kernel-bounds} depend on it, and $\varepsilon$ is restricted to the interval $\left(1-\alpha,\bar{\varepsilon}\right]$ because Proposition~\ref{prop:contraction} requires $1-\varepsilon<\alpha$. Admissible numbers $s$ exist. Indeed, for $s=\frac{1}{2}$ and $0<\varepsilon\le\bar{\varepsilon}=\frac{1}{200}$,
\[
s\cdot\left(\frac{1}{4}-9\varepsilon\right)-12\varepsilon=\frac{1}{8}-\frac{33\varepsilon}{2}\ge\frac{17}{400}=0.0425,
\]
whereas
\[
\frac{2a_L}{N^{*}}=\frac{L^3-1}{4L^3\cdot(L-1)^3}<\frac{1}{4\,(L-1)^3}\le\frac{1}{256},
\]
as $L\ge5$; hence $\sigma>0$ for $s=\frac{1}{2}$.

Odd integers $L$ as required in stage~\ref{it:order-choice-4} exist. Indeed,
\[
\frac{1}{8}-\frac{3\bar{\varepsilon}}{2}=0.1175>0,\qquad\frac{3}{8}-6\bar{\varepsilon}=0.345>0,\qquad\kappa-6\kappa_0-4\ge\kappa_0+3>0,\qquad\kappa-\kappa_0-1>0,
\]
the last two by $\kappa\ge7\kappa_0+7$. Hence, as $L\to\infty$, the left side of the first inequality displayed in stage~\ref{it:order-choice-4} and the number $L\cdot\left(\kappa-6\kappa_0-4\right)-1-\kappa$ tend to $\infty$, while the left side of the second inequality and that of the fourth, which equals $c_{\mathrm{rb}}\cdot\exp(\kappa)\cdot L^3\cdot\exp\left(-L\cdot\left(\kappa-\kappa_0-1\right)\right)$, tend to $0$. So the four inequalities, together with $L\ge L^F$ and $L\ge L^P$, hold for every sufficiently large $L$. Powers $M=L^m$ as required in stage~\ref{it:order-choice-6} exist, since $L^m\to\infty$ as $m\to\infty$ (because $L\ge5$) and each condition of that stage is a lower bound on $M$ by one of the numbers
\begin{gather*}
L^2,\qquad4L,\qquad\frac{16}{\gamma_w},\qquad\left(1+\exp(\hat{\kappa})\right)^{\frac{1}{\alpha_0}},\\
M^F\left(b,\alpha,L,\alpha_0,\kappa_w+\kappa_0\right),\qquad M^P\left(\alpha,L,\alpha_0\right),\qquad L^{m_0+2},
\end{gather*}
which are finite and depend on the constants of stages~\ref{it:order-choice-1}--\ref{it:order-choice-5} only. Numbers $e_0$ as required in stage~\ref{it:order-choice-9} exist. Each inequality of the conditions of that stage is either imposed for every $e_k\in(0,e_0]$ and holds for all $e_k$ below a positive number, or it is an inequality for $e_0$ that holds for all $e_0$ below a positive number; so every $e_0>0$ below the least of these finitely many positive numbers satisfies all of them. Indeed, after multiplication by a power of $e_k$ or of $e_0$ where necessary, each of these inequalities either bounds from above by a positive constant a sum of positive multiples of positive powers of $e_k$ or of $e_0$, or a polynomial in $v_k$ times a positive power of $e_k$, or it is one of the lower bounds
\[
\frac{1}{2}\,v_k^{2p_0}\ge7v_k+\log12,\qquad v_k\ge9,\qquad\left(-\log e_0\right)^2\ge10L,
\]
the first of which is the second inequality of~\eqref{eq:renormalisation-group-A1} after taking logarithms; the first two hold for all large $v_k$, the third for all small $e_0$. The exponents of $e_k$ and of $e_0$ that occur are positive since $\varepsilon\le\bar{\varepsilon}=\frac{1}{200}$, and a polynomial in $v_k$ times a positive power $e_k^a=\exp\left(-a\cdot v_k\right)$ tends to $0$ as $e_k\to0$.

\subsubsection{The conditions}\label{sec:conditions}

The conditions of this subsection are of three kinds.
\begin{enumerate}
\item[(i)] Conditions on the constants of stages~\ref{it:order-choice-1}--\ref{it:order-choice-10}: the conditions~\eqref{eq:renormalisation-group-A0}--\eqref{eq:renormalisation-group-A10}, condition~\eqref{eq:renormalisation-group-19} at $\kappa'=\kappa_w+\kappa_0$, the conditions~\eqref{eq:renormalisation-group-20}--\eqref{eq:renormalisation-group-24}, \eqref{eq:renormalisation-group-25}, and condition~\eqref{eq:renormalisation-group-26}, which defines $e_1$.
\item[(ii)] Conditions on $(R,e,N)$: the conditions~\eqref{eq:renormalisation-group-16}, \eqref{eq:renormalisation-group-fl} and~\eqref{eq:renormalisation-group-T1}--\eqref{eq:renormalisation-group-T6}. They follow from $e_K\le e_1$ (Theorem~\ref{thm:scales-stop-levels}).
\item[(iii)] Conditions of stage~\ref{it:order-choice-11}, which involve $\bar{m}$ or the probes: the conditions~\eqref{eq:renormalisation-group-17}, \eqref{eq:renormalisation-group-18}, \eqref{eq:four-jets-sm1}--\eqref{eq:four-jets-sm8}, \eqref{eq:renormalisation-group-27}--\eqref{eq:renormalisation-group-29} and~\eqref{eq:free-coefficient-1}--\eqref{eq:free-coefficient-4}. A statement that assumes one of them names it.
\end{enumerate}
The phrase \emph{the constants satisfy the conditions of Section~\ref{sec:conditions}} refers to the conditions of the first kind.

For $R\in\mathrm{Tor}(e)$ and $N\ge N_{\mathrm{min}}(R,e)$, the quantities with the index $*$ are taken at the level $k^{*}=k^{*}(R)$ of Definition~\ref{def:tori-stop-level}. The first condition on $(R,e,N)$ concerns the number $r_{*}=r_{k^{*}}$:
\begin{equation}\label{eq:renormalisation-group-16}
r_{*}>31L+16=3L+14\,d_U+2.
\end{equation}
By Proposition~\ref{prop:depth-stop-level}(f), condition~\eqref{eq:renormalisation-group-16} gives $l_k(R)\in\mathrm{Cl}$ for $k\le k^{*}$; by Lemma~\ref{lem:agreement-criterion-small}(b), it gives that the small polymers are flat. The inequality $l_k\ge M\cdot r_k$ for $k\le k^{*}$ holds by the definition of $k^{*}$ (Theorem~\ref{thm:scales-stop-levels}).

The requirements
\begin{multline}\tag{A0}\label{eq:renormalisation-group-A0}
L\cdot\left(\kappa-6\kappa_0-4\right)-1\ge\kappa,\qquad\kappa_w=\kappa+\kappa_0+3,\\
M\ge L^2,\qquad M>4L,\qquad\gamma_w\cdot M\ge1,\qquad M^{\alpha_0}\ge1+\exp(\hat{\kappa})
\end{multline}
are contained in stages~\ref{it:order-choice-4}--\ref{it:order-choice-6}.

The conditions~\eqref{eq:renormalisation-group-A1}--\eqref{eq:renormalisation-group-A7} are imposed for every $e_k\in(0,e_0]$, with $v_k=-\log e_k$ and $p_{0,k}=v_k^{p_0}$:
\begin{align}
&p_{0,k}\le e_k^{-\frac{\varepsilon}{4}},\qquad12\exp\left(-\frac{p_{0,k}^2}{2}\right)\le e_k^7,\qquad e_k^{\frac{5\varepsilon}{2}}\le\frac{1}{2},\tag{A1}\label{eq:renormalisation-group-A1}\\
&\qquad c_{\mathrm{fl}}\cdot e_k^{\varepsilon}\le1,\qquad3M\cdot e_k^{\varepsilon}\le1;\notag\\
&2\,C_w^2\cdot c_{\gamma}^2\cdot c_C\cdot e_k^{\frac{\varepsilon}{4}}\le1;\tag{A2}\label{eq:renormalisation-group-A2}\\
&16\,c_X\cdot c_{\gamma}\cdot C_w\cdot C_T^2\cdot M^3\cdot e_k^{\frac{\varepsilon}{2}}\le1,\qquad96\,M^3\cdot e_k^{\frac{1}{4}}\le1;\tag{A3}\label{eq:renormalisation-group-A3}\\
&32\,C_w\cdot C_T\cdot e_k^{\frac{3\varepsilon}{2}}\le1;\tag{A4}\label{eq:renormalisation-group-A4}\\
&88\,M^3\cdot e_k^{\frac{1}{4}-6\varepsilon}\le1;\tag{A5}\label{eq:renormalisation-group-A5}\\
&16\,C_w^2\cdot c_{\gamma}\cdot e_k^{2\varepsilon}\le3;\tag{A6}\label{eq:renormalisation-group-A6}\\
&c_S\cdot e_k^{\frac{1}{4}-5\varepsilon}\le1,\qquad c_S\cdot L^3\cdot e_k^{\varepsilon}\le1,\qquad\mathbf{H}_0:=36\,c_X^2\cdot K_0^2\cdot e_k^{\frac{1}{4}-5\varepsilon}\le1.\tag{A7}\label{eq:renormalisation-group-A7}
\end{align}
Furthermore, $e_0$ is subject to the conditions~\eqref{eq:renormalisation-group-A8}, which concern the second slot, to the conditions~\eqref{eq:renormalisation-group-A9}, which are those for Propositions~\ref{prop:adjustment}--\ref{prop:contraction}, and to condition~\eqref{eq:renormalisation-group-A10}:
\begin{align}
&e_0^{\varepsilon}\le\frac{1}{4},\qquad3M\cdot\left(e_0^{\frac{1}{4}+\varepsilon}+e_0^{2\varepsilon}\right)\le1,\qquad3M\cdot e_0^{\varepsilon}\le1,\tag{A8}\label{eq:renormalisation-group-A8}\\
&\qquad e_0^{-\frac{3}{4}+\varepsilon}+e_0^{-1+2\varepsilon}\le e_0^{-1+\varepsilon};\notag\\
&160\,L^{1+\frac{3\varepsilon}{2}}\cdot M\cdot e_0^{\varepsilon}\le1,\qquad90\,c_X\cdot L\cdot M\cdot e_0^{\varepsilon}\le1,\qquad9M\cdot L^{\varepsilon}\cdot e_0^{2\varepsilon}\le1,\tag{A9}\label{eq:renormalisation-group-A9}\\
&\qquad\left(-\log e_0\right)^2\ge10L;\notag\\
&e_0\le e^F_0,\quad\text{and, for every }e_k\in(0,e_0],\text{ with }v_k=-\log e_k\text{ and }p_k=v_k^p,\tag{A10}\label{eq:renormalisation-group-A10}\\
&\qquad v_k\ge9,\qquad8M\cdot\left(\frac{v_k}{4}+4\right)\cdot e_k^{\varepsilon}\le1,\notag\\
&\qquad8\,C_B\cdot p_k\le e_k^{-\frac{3}{4}+3\varepsilon},\qquad8\,e_k^{\frac{3}{4}-\frac{7\varepsilon}{2}}\le1;\notag\\
&\qquad8\,c_{\mathrm{rb}}\cdot L^3\cdot C_d\cdot M^3\cdot e_0^{2\varepsilon}\le1.\notag
\end{align}
Here $e^F_0$ is the number $e^F_0(b,\alpha,L,M,\alpha_0,\kappa',\varepsilon)$ of Theorem~\ref{thm:fermion-operators-weakened} at $\kappa'=\kappa_w+\kappa_0$, formed with the constants $b$, $\alpha$, $L$, $\alpha_0$, $M$ and $\varepsilon$ of stages~\ref{it:order-choice-1}--\ref{it:order-choice-8}; it is described after condition~\eqref{eq:renormalisation-group-24} below. The exponent $p$ in $p_k=v_k^p$ is that of stage~\ref{it:order-choice-8}. The number $C^F=C^F(b,\alpha,L,M,\alpha_0,\kappa')$ of Theorem~\ref{thm:fermion-operators-weakened} depends on $\kappa'$, and so does $e^F_0$, whose description after condition~\eqref{eq:renormalisation-group-24} involves $C^F$. In stage~\ref{it:order-choice-7}, in the formula for $C_0$ and in the description of the number $e^F_0$ of~\eqref{eq:renormalisation-group-A10}, the number $C^F$ is taken at $\kappa'=\kappa_w+\kappa_0$ and formed with the constants $b$, $\alpha$, $L$, $M$ and $\alpha_0$ of stages~\ref{it:order-choice-1}--\ref{it:order-choice-6}; $\kappa_w+\kappa_0$ is also the value of $\kappa'$ at which $M^F$ enters stage~\ref{it:order-choice-6} and $C^F$ enters the constant $C_w^0$ of Section~\ref{sec:photon-operators}. The first two of the four inequalities that~\eqref{eq:renormalisation-group-A10} imposes for every $e_k\in(0,e_0]$, namely $v_k\ge9$ and $8M\cdot\left(\frac{v_k}{4}+4\right)\cdot e_k^{\varepsilon}\le1$, are the second and the third inequality of Proposition~\ref{prop:local-expansion-determinant}(d4). The last two, $8\,C_B\cdot p_k\le e_k^{-\frac{3}{4}+3\varepsilon}$ and $8\,e_k^{\frac{3}{4}-\frac{7\varepsilon}{2}}\le1$, are the two displayed inequalities of Proposition~\ref{prop:fields-domains}, with the constant $C_B$ of Proposition~\ref{prop:local-regularity} (stage~\ref{it:order-choice-7}); the other conditions of Proposition~\ref{prop:fields-domains} on the charge are~\eqref{eq:renormalisation-group-A1} and~\eqref{eq:renormalisation-group-A2}. In the last inequality of~\eqref{eq:renormalisation-group-A10}, $C_d$ is the constant of stage~\ref{it:order-choice-7}, that is, the number~\eqref{eq:operators-fermion-10} formed with the constant $C_w$ of $\mathrm{par}$, so that this inequality is condition~\eqref{eq:renormalisation-group-24} with the constant $C_w$ of $\mathrm{par}$. It is the condition on $e_0$ on which Proposition~\ref{prop:local-expansion-determinant}(d6) rests; in Proposition~\ref{prop:polymer-expansion-determinant}, of which Proposition~\ref{prop:local-expansion-determinant} is a special case, condition~\eqref{eq:renormalisation-group-24} is a hypothesis of clause~(g). Clause~(d4) of Proposition~\ref{prop:local-expansion-determinant} assumes moreover $l\ge M\cdot r_k$, its clauses~(d5) and~(d6) assume in addition $l/L\in\mathrm{Cl}$ and $e_{k+1}=L^{\frac{1}{2}}\cdot e_k\le e_0$, and Proposition~\ref{prop:fields-domains}(ii) assumes $l/L\in\mathrm{Cl}$. These are requirements on the data $(k,l,e_k)$, and they are part of Definition~\ref{def:hypothesis-one-step}, whereas~\eqref{eq:renormalisation-group-A10} is a condition on the constants alone. Condition~\eqref{eq:renormalisation-group-A10} holds for every sufficiently small $e_0>0$, and it remains true when $e_0$ is decreased. Indeed, for $t\in(0,1)$ put $v=-\log t$. For $e_k=t$ the four inequalities that~\eqref{eq:renormalisation-group-A10} imposes for every $e_k\in(0,e_0]$ state that $v\ge9$ and that each of the three numbers
\begin{gather*}
8M\cdot\left(\frac{v}{4}+4\right)\cdot\exp\left(-\varepsilon\cdot v\right),\qquad8\,C_B\cdot v^p\cdot\exp\left(-\left(\frac{3}{4}-3\varepsilon\right)\cdot v\right),\\
8\exp\left(-\left(\frac{3}{4}-\frac{7\varepsilon}{2}\right)\cdot v\right)
\end{gather*}
is at most $1$. These numbers tend to $0$ as $v\to\infty$, because $\varepsilon>0$ and, as $\varepsilon\le\frac{1}{200}$, both $\frac{3}{4}-3\varepsilon$ and $\frac{3}{4}-\frac{7\varepsilon}{2}$ are positive. Hence there is a number $v^{\circ}\ge9$ such that all three are at most $1$ for every $v\ge v^{\circ}$, and the four inequalities hold for every $e_k\in(0,e_0]$ as soon as $e_0\le\exp\left(-v^{\circ}\right)$. The two remaining inequalities of~\eqref{eq:renormalisation-group-A10} state that $e_0\le e^F_0$ and $e_0\le\left(8c_{\mathrm{rb}}\cdot L^3\cdot C_d\cdot M^3\right)^{-\frac{1}{2\varepsilon}}$, and both bounds are positive numbers fixed before stage~\ref{it:order-choice-9}. The other conditions of stage~\ref{it:order-choice-9} have the same two properties: \eqref{eq:renormalisation-group-A1}--\eqref{eq:renormalisation-group-A7} are imposed for every $e_k\in(0,e_0]$ and hold for all sufficiently small $e_k$, and each inequality of~\eqref{eq:renormalisation-group-A8}, \eqref{eq:renormalisation-group-A9} and~\eqref{eq:renormalisation-group-20}--\eqref{eq:renormalisation-group-25} holds for all sufficiently small $e_0$ and, if it holds for some $e_0\in(0,1)$, it holds for every smaller positive number (for the fourth inequality of~\eqref{eq:renormalisation-group-A8} this is seen after multiplication by $e_0^{1-\varepsilon}$). Therefore all conditions of stage~\ref{it:order-choice-9}, condition~\eqref{eq:renormalisation-group-A10} among them, hold together for every sufficiently small $e_0>0$.

For $R$ and $N$ as above, the remaining conditions on $(R,e,N)$ are
\begin{align}
&\sum_{k=0}^{k^{*}}\mathfrak{s}_k\le\frac{1}{16}\quad\text{with}\quad\mathfrak{s}_k:=\left(c_{\mathrm{fl}}+2\right)\cdot e_k^{\varepsilon},\qquad v_{*}\ge8;\tag{fl}\label{eq:renormalisation-group-fl}\\
&v_{*}\ge8,\qquad p_0\cdot v_{*}^{p_0-1}\ge28,\qquad w^{\mathrm{m}}_0<1;\tag{T1}\label{eq:renormalisation-group-T1}\\
&8\,C_0^2\cdot l_{*}^6\cdot\left(b_{-}^{-1}+\mathfrak{m}_{*}^{-1}\right)\le e_{*}^{-\frac{\varepsilon}{2}},\tag{T2}\label{eq:renormalisation-group-T2}\\
&\qquad4\left(2C_0+c_{\kappa}+5c_{\mathrm{fl}}+6\right)\cdot l_{*}^3\le e_{*}^{-\varepsilon},\qquad e_{*}^{\frac{1}{4}-9\varepsilon}\le\frac{1}{8};\notag\\
&4\,C_0\cdot l_{*}^2\cdot p_{*}<e_{*}^{-\varepsilon};\tag{T3}\label{eq:renormalisation-group-T3}\\
&12\,l_{*}^3\cdot\left(\exp\left(-\frac{p_{0*}}{2}\right)+\exp\left(-\frac{p_{*}^2}{2C_0^2}\right)\right)\le e_{*};\tag{T4}\label{eq:renormalisation-group-T4}\\
&C_{\Sigma}\cdot\exp\left(-\gamma_1\cdot l_{*}\right)\le e_{*};\tag{T5}\label{eq:renormalisation-group-T5}\\
&L^2\cdot e^2\cdot R\le1.\tag{T6}\label{eq:renormalisation-group-T6}
\end{align}
In~\eqref{eq:renormalisation-group-T1} the number $w^{\mathrm{m}}_0$ is formed with the charge $q$ of level $0$ (Definition~\ref{def:notation}(e)).

The error terms that occur in Section~\ref{sec:trajectories-pair-bounds} are the numbers
\begin{gather*}
\varepsilon_1:=\frac{1}{2}\,e_{*}^{\frac{1}{4}-9\varepsilon},\qquad\varepsilon_2:=C_{\Sigma}\cdot l_{*}^{-1}\cdot\exp\left(-\gamma_1\cdot l_{*}\right),\\
\varepsilon_3:=6\,l_{*}^3\cdot\exp\left(-\frac{p_{0*}}{2}\right)+6\,l_{*}^3\cdot\exp\left(-\frac{p_{*}^2}{2C_0^2}\right),\\
\varepsilon_4:=6\,l_{*}^3\cdot\exp\left(-\frac{p_{*}^2}{2C_0^2}\right),
\end{gather*}
and the values of the function $\varepsilon_5\colon(0,\infty)\to(0,\infty)$,
\[
\varepsilon_5(x):=312\cdot\log\left(\left(1-\exp\left(-\frac{x}{4}\right)\right)^{-1}\right).
\]

The following two conditions belong to stage~\ref{it:order-choice-11}:
\begin{equation}\label{eq:renormalisation-group-17}
6N^{*}\cdot e^2\cdot R^{*}\cdot\log\frac{1}{e}\le\pi^2,
\end{equation}
\begin{equation}\label{eq:renormalisation-group-18}
e_K\cdot\log h'\le1.
\end{equation}
Together with $e_K\le e_1$, condition~\eqref{eq:renormalisation-group-18} implies~\eqref{eq:renormalisation-group-17} (Theorem~\ref{thm:scales-stop-levels} and stage~\ref{it:order-choice-10}).

The conditions~\eqref{eq:four-jets-sm1}--\eqref{eq:four-jets-sm8} are the smallness conditions of Section~\ref{sec:smallness-conditions}; they hold for $e\le e_{\sharp}$ (Theorem~\ref{thm:smallness-conditions-small}). Hypothesis~\ref{def:hypothesis-reference-torus} of Section~\ref{sec:probes-jets-reference} is assumed throughout Section~\ref{sec:four-jets-one}.

The remaining conditions are stated also in Sections~\ref{sec:operators-photon-field}, \ref{sec:operators-fermion-field}, \ref{sec:renormalisation-flow-tuned}, \ref{sec:four-jets-one} and~\ref{sec:free-coefficient-chessboard}. Each of them except the first requires $e_0$, $e_1$, $e_{\sharp}$ or $e_2$ to be small in terms of constants chosen earlier, and none of them adds a condition on $\kappa$, $L$ or $M$ to stages~\ref{it:order-choice-3}, \ref{it:order-choice-4} and~\ref{it:order-choice-6}. The first one is, for $\kappa'>0$,
\begin{multline}\label{eq:renormalisation-group-19}
M=L^m\ge M^F\left(b,\alpha,L,\alpha_0,\kappa'\right)\\
:=\max\left\{4L+1,\ \left(108\pi\cdot C_V\cdot a\right)^2,\ 108\pi\cdot C_V\cdot a\cdot\exp\left(27\kappa'\right)\right\},
\end{multline}
where $C_V$ and $a$ are the constants of~\eqref{eq:operators-fermion-4}, functions of $(b,\alpha,L)$. For an arbitrary $\kappa'>0$ the display defines the number $M^F(b,\alpha,L,\alpha_0,\kappa')$ and states the hypothesis on $M$ of Theorem~\ref{thm:fermion-operators-weakened}. As a condition on the constants, \eqref{eq:renormalisation-group-19} is imposed at $\kappa'=\kappa_w+\kappa_0$, and in this role it is the inequality $M\ge M^F\left(b,\alpha,L,\alpha_0,\kappa_w+\kappa_0\right)$ of stage~\ref{it:order-choice-6}. Next come the conditions
\begin{gather}
6L\cdot e_0^{\varepsilon}\le1,\label{eq:renormalisation-group-20}\\
4\,C_V\cdot a\cdot e_0^{\varepsilon}\le1,\label{eq:renormalisation-group-21}\\
3M\cdot e_0^{\varepsilon}\le1,\label{eq:renormalisation-group-22}\\
e_0^{\varepsilon}\le\frac{1}{4},\qquad e_0^{\frac{1}{4}}+e_0^{\varepsilon}\le1,\qquad4e_0^{\frac{1}{4}+2\varepsilon}+2e_0^{3\varepsilon}\le1,\label{eq:renormalisation-group-23}\\
8\,c_{\mathrm{rb}}\cdot L^3\cdot C_d\cdot M^3\cdot e_0^{2\varepsilon}\le1,\label{eq:renormalisation-group-24}
\end{gather}
where, in~\eqref{eq:renormalisation-group-24}, $C_d$ is the number~\eqref{eq:operators-fermion-10} formed with the constant $C_w$ of the parameter set for which Proposition~\ref{prop:polymer-expansion-determinant} is applied. In stage~\ref{it:order-choice-9}, in Definition~\ref{def:hypothesis-one-step} and in~\eqref{eq:renormalisation-group-A10} this is the constant $C_w$ of $\mathrm{par}$, so that $C_d$ is the constant of Proposition~\ref{prop:local-expansion-determinant} (stage~\ref{it:order-choice-7}); in the definition of $e^F_0$ it is $C^F$. These five conditions belong to stage~\ref{it:order-choice-9}. The number $e^F_0$ of Theorem~\ref{thm:fermion-operators-weakened} is the largest $t\in(0,1]$ such that~\eqref{eq:renormalisation-group-20}--\eqref{eq:renormalisation-group-23} hold with $t$ in place of $e_0$ and~\eqref{eq:renormalisation-group-24} holds with $t$ in place of $e_0$ and with $C_w=C^F$. The condition
\begin{equation}\label{eq:renormalisation-group-25}
e_0\le\exp(-1)
\end{equation}
belongs to stage~\ref{it:order-choice-9} and is contained in~\eqref{eq:renormalisation-group-A9}: the first inequality of~\eqref{eq:renormalisation-group-A9} gives $e_0^{\varepsilon}\le\frac{1}{160}$, hence $e_0\le160^{-\frac{1}{\varepsilon}}<\exp(-1)$. The condition
\begin{equation}\label{eq:renormalisation-group-26}
e_1\ \text{is the number defined by the inequalities~\eqref{eq:renormalisation-flow-e1}--\eqref{eq:renormalisation-flow-e8} of Section~\ref{sec:constants-conditions-torus}}
\end{equation}
belongs to stage~\ref{it:order-choice-10}. The conditions
\begin{gather}
\left(4+\frac{M^3}{2}\right)\cdot e_{*}^{\varepsilon}\le1,\label{eq:renormalisation-group-27}\\
4\,C_0\cdot l_{*}^3\cdot g_1\le e_{*}^{-\frac{1}{4}+\varepsilon},\label{eq:renormalisation-group-28}\\
48\,e_{*}^{2\varepsilon}\le\mathfrak{m}_{-}\label{eq:renormalisation-group-29}
\end{gather}
are parts of the conditions~\eqref{eq:four-jets-sm1}, \eqref{eq:four-jets-sm3} and~\eqref{eq:four-jets-sm4}; here $g_1$ is the number of Section~\ref{sec:constants-3}, and $\mathfrak{m}_{-}=L^{-m_0-3}$ is the mass of Section~\ref{sec:probes-jets-reference}. They belong to stage~\ref{it:order-choice-11}. Finally, the number $e_2$ is subject to the four conditions~\eqref{eq:free-coefficient-1}--\eqref{eq:free-coefficient-4}, which belong to stage~\ref{it:order-choice-11}:
\begin{gather*}
e_2\le e_G(\bar{m}),\\
e_2\le e_{\sharp}(\beta')\quad\text{for the three probes}\quad\beta'=\beta,\ \beta+\tau_{x_0}\beta,\ \beta-\tau_{x_0}\beta,\\
\frac{C_Y+1}{6\log\left(e_2^{-1}\right)}+\frac{\mathfrak{e}_2}{3}\le1,\\
\frac{5}{2}\,C_{\mathrm{kap}}\cdot e_2^{\sigma}\le\frac{b(f)}{4};
\end{gather*}
here $C_Y$ is the constant of Step~1 in Section~\ref{sec:proof-part-ii-2}, $\sigma$ is the number defined in stage~\ref{it:order-choice-8}, $\mathfrak{e}_2=6\varepsilon$, and $b(f)$ is the free coefficient (Section~\ref{sec:free-coefficients-bare}) of the two-form $f=dd^{*}\beta$.

\subsection{Classical results}\label{sec:classical-results}

\subsubsection{Results from the literature}\label{sec:results-literature}

The principal classical results that are quoted in their standard form are the following, with references.

\paragraph{Complex analysis.}
Cauchy's integral theorem; Cauchy's integral formula and the Cauchy estimates, in one and in several complex variables; the identity theorem; Schwarz's lemma; Morera's theorem; Weierstrass' theorem on uniform limits of analytic functions; the theorem of Vitali and Porter. See W.~Rudin, \emph{Real and Complex Analysis}, third edition, Chapters~10 and~12; L.~H\"ormander, \emph{An Introduction to Complex Analysis in Several Variables}, third edition, Chapter~II; J.~L.~Schiff, \emph{Normal Families}, Chapter~2.

\paragraph{Real analysis and Fourier analysis.}
Fubini's theorem and its analogue for Berezin integrals; the dominated convergence theorem and Fatou's lemma; the inequalities of Jensen, Cauchy--Schwarz, H\"older, Minkowski and Young; Taylor's formula; Plancherel's theorem; Fourier series of smooth periodic functions; the inversion formula for the finite Fourier transform. See Rudin, loc.\ cit., Chapters~1, 3, 8 and~9; E.~M.~Stein and G.~Weiss, \emph{Introduction to Fourier Analysis on Euclidean Spaces}, Chapters~I and~VII.

\paragraph{Matrices and operators.}
The Neumann series; Schur's test; the contraction mapping theorem of Banach; Cramer's rule; Schur complements, the Woodbury identity, Sylvester's determinant identity and the matrix determinant lemma; the identity
\[
\det(1+Y)=\exp\operatorname{Tr}\log(1+Y)
\]
for every complex square matrix $Y$ of operator norm less than $1$, where $\log(1+Y):=\sum_{j=1}^{\infty}(-1)^{j+1}j^{-1}\cdot Y^j$; the spectral theorem for Hermitian and for real symmetric matrices, and the formula
\[
x^{-\frac{1}{2}}=\frac{1}{\pi}\int_0^{\infty}\mathrm{d}t\,t^{-\frac{1}{2}}\cdot(x+t)^{-1}\qquad(x>0);
\]
Gram's determinant inequality; the Schur product theorem; Hilbert--Schmidt operators and trace class operators, the cyclicity of the trace, and H\"older's inequality, also for the Schatten classes; the Kato--Seiler--Simon inequality; the spectral theorem for self-adjoint operators (see Corollary~\ref{cor:transfer-operators-coordinate} and Proposition~\ref{prop:reconstruction-direction}); Stone's theorem, together with the spectral theorem for finitely many commuting self-adjoint operators, in the following form (see the proof of Proposition~\ref{prop:trace-formula-bloch}): for every strongly continuous unitary representation $y\mapsto\tau_y$ of $\mathbb{R}^l$, $l\ge1$, on a separable Hilbert space $\mathcal{H}$ with the scalar product $\langle\cdot,\cdot\rangle$ there are commuting self-adjoint operators $E_1,\dots,E_l$ on $\mathcal{H}$ and a projection-valued measure $P$ on the Borel sets of $\mathbb{R}^l$, their joint spectral measure, such that $\tau_y=\exp\left(-i\sum_{b=1}^ly_bE_b\right)$, that is,
\[
\langle f,\tau_yg\rangle=\int\mathrm{d}\langle f,P_{\lambda}g\rangle\,\exp\left(-i\sum_{b=1}^ly_b\lambda_b\right)\qquad\left(f,g\in\mathcal{H},\ y\in\mathbb{R}^l\right),
\]
and $\exp\left(-\frac{1}{2}\sum_{b=1}^lE_b^2\right)$ denotes the operator $\int\mathrm{d}P_{\lambda}\,\exp\left(-\frac{1}{2}|\lambda|^2\right)$; the description of the elements of the Schwartz space $\mathcal{S}$ and of the space $\mathcal{S}'$ of tempered distributions by their Hermite coefficients. See R.~A.~Horn and C.~R.~Johnson, \emph{Matrix Analysis}, second edition; M.~Reed and B.~Simon, \emph{Methods of Modern Mathematical Physics~I}, Chapters~II and~V--VIII, in particular Theorems~V.13, V.14, VI.22, VI.23, VII.3, VIII.4 and~VIII.5; B.~Simon, \emph{Trace Ideals and Their Applications}, second edition, Theorems~2.8, 3.1 and~4.1 and Corollary~3.8.

\paragraph{Combinatorics and topology.}
M\"obius inversion over the lattice of the subsets of a finite set; Cayley's formula for the number of labelled trees with given vertex degrees; Euler's theorem on Euler circuits; Kirchhoff's matrix-tree theorem; the Hodge decomposition on the discrete torus; for the three-dimensional torus $\mathbb{T}^3$, the isomorphism $H^2(\mathbb{T}^3;\mathbb{Z})\cong\mathbb{Z}^3$ given by the three fluxes, that is, by the values of a cohomology class on the fundamental classes of the three coordinate two-tori. See R.~P.~Stanley, \emph{Enumerative Combinatorics}, Volume~1, Chapters~2 and~3; J.~W.~Moon, \emph{Counting Labelled Trees}; B.~Eckmann, \emph{Harmonische Funktionen und Randwertaufgaben in einem Komplex}, Comment.\ Math.\ Helv.\ 17 (1945); N.~Biggs, \emph{Algebraic Graph Theory}, second edition, Chapter~6; A.~Hatcher, \emph{Algebraic Topology}, Theorem~3.2 (the universal coefficient theorem) and Section~3.B (the K\"unneth formula).

\paragraph{Grassmann integration.}
Berezin integration; the Gaussian integral $\int\mathrm{d}\mu_{\operatorname{Cov}}(W)\,F$ of Proposition~\ref{prop:weighted-coefficient-norms}, for an invertible covariance matrix $\operatorname{Cov}$; Wick's theorem for Gaussian Grassmann integrals. See F.~A.~Berezin, \emph{The Method of Second Quantization}; J.~Feldman, H.~Kn\"orrer and E.~Trubowitz, \emph{Fermionic Functional Integrals and the Renormalization Group}, Chapter~I. Section~\ref{sec:spinor-spaces-berezin} contains the rules \eqref{eq:algebraic-gaussian-Ber1}--\eqref{eq:algebraic-gaussian-Ber5} of Berezin integration with their proofs: the invariance under shifts of the generators, the Gaussian integral with sources, the invariance under the scaling of $\Psi$ by a periodic function $u$ without zeros and of $\bar{\Psi}$ by $u^{-1}$, the change of covariance, and the interchange of integrals. Section~\ref{sec:norms-gaussian-integrals} contains the definitions of the maps $\int\mathrm{D}W$, $\mathrm{d}\mu_I(W)$ and $\mathrm{d}\mu_{\operatorname{Cov}}(W)$, the formula~\eqref{eq:polymer-activities-m} for the moments of the basis monomials with its proof, and, as parts of Theorem~\ref{thm:polymer-activities-1} with their proofs, the bound, the factorisation, the invariance and the change of covariance, that is, items~\eqref{it:weighted-coefficient-norms-N3}, \eqref{it:weighted-coefficient-norms-N4}, \eqref{it:weighted-coefficient-norms-N4p} and~\eqref{it:weighted-coefficient-norms-N7}(iii) of Proposition~\ref{prop:weighted-coefficient-norms}. Wick's rule in the form of a determinant, for $\mathrm{d}\mu_C$ with an invertible matrix $C$, is Lemma~\ref{lem:wick-rule-gram}(a) and~(b). Wick's theorem in the form of a Pfaffian, for $\mathrm{d}\mu_I$ and products of arbitrary linear combinations of the barred and the unbarred generators, is stated, and derived from~\eqref{eq:polymer-activities-m}, in the paragraph \emph{Classical facts} of Section~\ref{sec:recalled-auxiliary-results}; it is used in the proof of Lemma~\ref{lem:substitutions-integral}(a).

\paragraph{Probability.}
Tail bounds for Gaussian variables, and union bounds; the finite-dimensional Cameron--Martin formula; Dynkin's lemma; weak convergence of probability measures and Prokhorov's theorem (P.~Billingsley, \emph{Convergence of Probability Measures}, second edition, Theorems~2.7, 3.5, 5.1 and~6.8); the existence of regular conditional distributions (O.~Kallenberg, \emph{Foundations of Modern Probability}, second edition, Theorem~6.3; R.~M.~Dudley, \emph{Real Analysis and Probability}, Theorem~10.2.2); Minlos' theorem~\cite[Theorem~3.4.2]{GlimmJaffe}, in the following form (see the proofs of Lemma~\ref{lem:gaussian-mixtures}, Lemma~\ref{lem:gaussian-laws-bounded}(a) and Lemma~\ref{lem:white-noise-telegraph}): let $\mathcal{S}\left(\mathbb{R}^3;\Lambda^2\right)$ be the real vector space of Schwartz maps of $\mathbb{R}^3$ into the space $\Lambda^2$ of real alternating two-tensors on $\mathbb{R}^3$, with its usual topology, and let $\mathcal{S}'\left(\mathbb{R}^3;\Lambda^2\right)$ be the space of continuous linear maps $\omega\colon\mathcal{S}\left(\mathbb{R}^3;\Lambda^2\right)\to\mathbb{R}$, endowed with its cylinder $\sigma$-algebra, the smallest $\sigma$-algebra for which the maps $\omega\mapsto\omega(f)$, $f\in\mathcal{S}\left(\mathbb{R}^3;\Lambda^2\right)$, are measurable; then every continuous function $\varphi\colon\mathcal{S}\left(\mathbb{R}^3;\Lambda^2\right)\to\mathbb{C}$ with $\varphi(0)=1$ which is positive definite, that is,
\[
\sum_{p,q=1}^n\overline{c_p}\cdot c_q\cdot\varphi\left(f_q-f_p\right)\ge0\qquad\left(n\ge1,\ c\in\mathbb{C}^n,\ f_1,\dots,f_n\in\mathcal{S}\left(\mathbb{R}^3;\Lambda^2\right)\right),
\]
is the characteristic functional of a probability measure $\rho$ on $\mathcal{S}'\left(\mathbb{R}^3;\Lambda^2\right)$ with this $\sigma$-algebra, that is,
\[
\int\mathrm{d}\rho(\omega)\,\exp\left(i\,\omega(f)\right)=\varphi(f)\qquad\left(f\in\mathcal{S}\left(\mathbb{R}^3;\Lambda^2\right)\right).
\]

\paragraph{Classical results proved in full.}
Complete proofs are given for the following classical results: the argument of Combes and Thomas, for instance in Lemma~\ref{lem:combes-thomas-estimate}, in the proof of Theorem~\ref{thm:block-covariance}, in Theorem~\ref{thm:exponentially-weighted-bound} and in the proof of Lemma~\ref{lem:zero-field-level}; the Mayer expansion and the tree bound for its coefficients, in clauses~\eqref{it:polymers-CE1} and~\eqref{it:polymers-CE2} of Theorem~\ref{thm:polymer-activities-2} and of Proposition~\ref{prop:polymers}; the forest formula, in Theorem~\ref{thm:stop-level-reference} and Lemma~\ref{lem:forest-formula-brydges}; the ergodic decomposition, with the uniqueness of the decomposing measure, in Theorem~\ref{thm:ergodic-decomposition}(a), and the existence of an ergodic decomposition, in Proposition~\ref{prop:existence-conditional-expectations}(a); the mean ergodic theorem, that is, the strong convergence on $L^2$ of the averages over cubes of the unitary operators of a measure-preserving action to the orthogonal projection onto the invariant vectors, in Lemma~\ref{lem:averages-translations}(iii), in Lemma~\ref{lem:averages-translations-2}(c) and~(d), and in Step~(1) of the proof of Lemma~\ref{lem:uniform-bound-invariant}; the two-constants lemma, in Lemma~\ref{lem:two-constants-estimate}; Mehler's kernel, in the proof of Lemma~\ref{lem:free-coefficient-1}.

\paragraph{Origin of the methods.}
In the following list each method is followed, in parentheses, by the places at which it is carried out with complete proofs, and then by the papers from which it is taken. Gaussian domination by means of reflections (Section~\ref{sec:gaussian-bounds-curvature}): J.~Fr\"ohlich, B.~Simon and T.~Spencer, \emph{Infrared bounds, phase transitions and continuous symmetry breaking}, Commun.\ Math.\ Phys.\ 50 (1976) 79--95. Chessboard estimates (Section~\ref{sec:trace-formula-chessboard}, by means of the trace formula of Proposition~\ref{prop:trace-formula-bloch} and of H\"older's inequality for Schatten classes, Lemma~\ref{lem:schatten-norms-traces}(a)): J.~Fr\"ohlich, R.~Israel, E.~H.~Lieb and B.~Simon, \emph{Phase transitions and reflection positivity~I}, Commun.\ Math.\ Phys.\ 62 (1978) 1--34. Positivity of the transfer matrix of Wilson fermions (Sections~\ref{sec:determinant-wilson-operator} and~\ref{sec:second-quantisation-trace}, Lemma~\ref{lem:slice-factorisation-determinant-2}): M.~L\"uscher~\cite{Luescher77}; K.~Osterwalder and E.~Seiler, \emph{Gauge field theories on a lattice}, Ann.\ Phys.\ 110 (1978) 440--471. The tree bound for the coefficients of the Mayer expansion (clause~\eqref{it:polymers-CE2} of Theorem~\ref{thm:polymer-activities-2} and of Proposition~\ref{prop:polymers}): O.~Penrose, \emph{Convergence of fugacity expansions for classical systems}, in: \emph{Statistical Mechanics, Foundations and Applications} (T.~A.~Bak, ed.), Benjamin, 1967. The forest formula (Lemma~\ref{lem:forest-formula-brydges}): D.~C.~Brydges and T.~Kennedy, \emph{Mayer expansions and the Hamilton--Jacobi equation}, J.\ Stat.\ Phys.\ 48 (1987) 19--49; A.~Abdesselam and V.~Rivasseau, \emph{Trees, forests and jungles: a botanical garden for cluster expansions}, in: \emph{Constructive Physics}, Lecture Notes in Physics 446, Springer, 1995. Exponential decay by conjugation with exponential weights (for instance Lemma~\ref{lem:combes-thomas-estimate}, the proof of Theorem~\ref{thm:block-covariance}, Theorem~\ref{thm:exponentially-weighted-bound} and the proof of Lemma~\ref{lem:zero-field-level}): J.~M.~Combes and L.~Thomas, \emph{Asymptotic behaviour of eigenfunctions for multiparticle Schr\"odinger operators}, Commun.\ Math.\ Phys.\ 34 (1973) 251--270. Local inverses and walks (Sections~\ref{sec:local-inverses-mirror}, \ref{sec:inverses-torus} and~\ref{sec:weakened-walk-expansions}): T.~Ba{\l}aban~\cite{Balaban}, for the representation of the inverse of a lattice operator with a block penalty by a parametrix built from local inverses and by a sum over walks, carried out for the covariant Laplacian with a block penalty in \emph{Regularity and decay of lattice Green's functions}, Commun.\ Math.\ Phys.\ 89 (1983) 571--597. For the lattice Dirac operator with a block penalty in an external gauge field, that is, for the operator $F_y(A)=D_A+\mathfrak{m}+\operatorname{Pen}_y(A)$ of~\eqref{eq:operators-fermion-2}, whose inverse is expanded in walks in these sections, the local inverses and the expansion in walks go back to T.~Ba{\l}aban, M.~O'Carroll and R.~Schor, \emph{Properties of block renormalization group operators for Euclidean fermions in an external gauge field}, J.\ Math.\ Phys.\ 32 (1991) 3199--3208, the second paper of~\cite{BalabanOCarrollSchor}. The following parts have the form given to them by J.~Dimock: the properties required of a local inverse, Proposition~\ref{prop:properties-local-inverse}(b), (c), after~\cite[(128), (131), (132)]{DimockUR}; the parametrix, the commutator operators, the Neumann series and the product over a walk, Lemma~\ref{lem:parametrix-identities}, Definition~\ref{def:walks} and Theorem~\ref{thm:expansion-walks}, after~\cite[(135)--(140)]{DimockUR} and~\cite[(100)--(105)]{DimockMB}; the weakening parameters, Definition~\ref{def:weakened-sums-walks}, after~\cite[(153)--(158)]{DimockUR}. The construction of the local inverses by means of mirror tori and folded fields, Definitions~\ref{def:mirror-torus-fold} and~\ref{def:folded-fields-local}, is not taken from these papers; it takes the place of their construction, which uses soft boundary conditions and a decomposition into shells of several scales, that is, an expansion in walks on several scales.

For real $a\in C_c^{\infty}\left(\mathbb{R}^3;\mathbb{R}^3\right)$, the leading term of the coefficient $q_{4,\mathsf{M}}(a)$ of Section~\ref{sec:free-coefficients-bare} as $\mathsf{M}\to\infty$, given in Theorem~\ref{thm:free-coefficients-lattice}(a), is of the type found by W.~Heisenberg and H.~Euler, \emph{Folgerungen aus der Diracschen Theorie des Positrons}, Z.\ Phys.\ 98 (1936) 714--732.

\section{Gaussian bounds for the curvature and reflection positivity}\label{sec:gaussian-bounds-curvature}

On the cubic lattice torus with $n\ge2$ sites per side, $n$ even or odd, and for every lattice mass $\boldsymbol{\mu}>-1$, Gaussian domination holds for the partition functions with sources (Theorem~\ref{thm:gaussian-domination-positive}), for the curvature under the law with the Proca term $\frac{\varepsilon}{2}\|A\|^2$, $\varepsilon>0$ (Theorem~\ref{thm:laplace-transform-curvature}), and, in the limit $\varepsilon\downarrow0$, for the curvature under the harmonic-averaged law (Theorem~\ref{thm:limit-vanishing-proca}), with consequences for linear functionals of the curvature (Corollary~\ref{cor:linear-functionals-curvature}); the law with the Proca term is moreover reflection positive (Proposition~\ref{prop:gaussian-bounds-2}), and these results give Estimate~1 for the harmonic-averaged law and a part of Estimate~2. The notation, in which the Clifford triple and the character are arbitrary, is fixed in the following subsection, and the Wilson operator of Section~\ref{sec:spinors-wilson-determinant} and the harmonic-averaged law of Section~\ref{sec:lattice-laws}, in the form in which Section~\ref{sec:lattice-measures-laws} writes them, are identified with special cases of the objects defined there, the operator in Section~\ref{sec:characters-twisted-spinors} and the law in Section~\ref{sec:proca-parent-harmonic}.

\subsection{Notation and standing assumptions}

Throughout, $n\ge2$ is an integer, even or odd, $\varepsilon>0$ and $q\in\mathbb R$ are real numbers, and $\boldsymbol{\mu}>-1$ is the total lattice mass. The letters $\mu,\nu,\mu'$ denote directions, that is, elements of $\{1,2,3\}$. Each of the letters $\lambda$, $\omega$, $F$, $\Phi$, $g$ has several meanings, each of which is declared at the place where it occurs: $\lambda$ is a zero-cochain, and from Section~\ref{sec:slices-slice-operators} on it also denotes the layer data; $\omega$ is a point of $C^2$, and from Section~\ref{sec:slices-slice-operators} on it also denotes the slice data; $F$ is a test function, the curvature, or an element of the class $\mathcal A_+$ of Proposition~\ref{prop:gaussian-bounds-2}; $\Phi$ is a spinor (Section~\ref{sec:characters-twisted-spinors}), and $\Phi_\pm$ are the operator-valued functions defined by integrals in the proofs of Propositions~\ref{prop:gaussian-bounds-1} and~\ref{prop:gaussian-bounds-2}; $g$ is a two-cochain or an endomorphism.

\subsubsection{Cochains on the periodic lattice}\label{sec:cochains-periodic-lattice}

Let $\Lambda=(\mathbb Z/n\mathbb Z)^3$, and let $\hat{1},\hat{2},\hat{3}$ denote the unit vectors of $\mathbb Z^3$ and their images in $\Lambda$. Put $C^0=\mathbb R^\Lambda$. Let $C^1$ be the space of the real functions $A=\left(A_\mu(x)\right)$ on the $3n^3$ links $(x,\mu)$, $x\in\Lambda$, $\mu\in\{1,2,3\}$, and let $C^2$ be the space of the real functions $\mathfrak{h}=\left(\mathfrak{h}_{\mu\nu}(x)\right)_{\mu<\nu}$ on the $3n^3$ plaquettes $(x,\mu\nu)$, $x\in\Lambda$, $1\le\mu<\nu\le3$. Each of the spaces $C^0$, $C^1$, $C^2$ carries the counting inner product $\langle\cdot,\cdot\rangle$ and the corresponding norm $\|\cdot\|$, and $\mathrm{d}A$ denotes Lebesgue measure on $C^1$. The coboundaries $d_0\colon C^0\to C^1$ and $d_1\colon C^1\to C^2$ are given by
\[
(d_0\lambda)_\mu(x)=\lambda(x+\hat\mu)-\lambda(x),\qquad(d_1A)_{\mu\nu}(x)=A_\mu(x)+A_\nu(x+\hat\mu)-A_\mu(x+\hat\nu)-A_\nu(x).
\]
Let $P\colon C^2\to C^2$ be the orthogonal projection onto $\operatorname{ran} d_1$, and put
\[
V_{\mathrm{g}}=\operatorname{ran} d_0,\qquad V_{\mathrm{h}}=\ker d_1\cap{V_{\mathrm{g}}}^{\perp},\qquad V_{\mathrm{t}}=\left(V_{\mathrm{g}}\oplus V_{\mathrm{h}}\right)^{\perp}=\left(\ker d_1\right)^{\perp};
\]
the last equality holds because $d_1\circ d_0=0$, so that $\ker d_1=V_{\mathrm{g}}\oplus V_{\mathrm{h}}$. These are the spaces $V_{\mathrm{g}}$, $V_{\mathrm{h}}$, $V_{\mathrm{t}}$ of Section~\ref{sec:lattice-measures-laws}; by Lemma~\ref{lem:harmonic-cochains-coercivity}, $V_{\mathrm{h}}$ consists of the constant cochains.

\subsubsection{Clifford triples}\label{sec:clifford-triples}

A \emph{Clifford triple} is a triple $(\gamma_1,\gamma_2,\gamma_3)$ of Hermitian $s\times s$ matrices such that
\[
\gamma_\mu\cdot\gamma_\nu+\gamma_\nu\cdot\gamma_\mu=2\delta_{\mu\nu}\cdot1\qquad\left(\mu,\nu\in\{1,2,3\}\right).
\]
For every Clifford triple the integer $s$ is even (Section~\ref{sec:slices-slice-operators}), and $\rho$ denotes the integer with $s=2\rho$. The \emph{four-component representation} is the Clifford triple with $s=4$ and
\[
\gamma_\mu=\operatorname{diag}\left(\sigma_\mu,-\sigma_\mu\right)\qquad\left(\mu\in\{1,2,3\}\right),
\]
with the Pauli matrices $\sigma_1,\sigma_2,\sigma_3$. Section~\ref{sec:dirac-matrices-lattice} labels the directions by $0,1,2$; for $\mu=1,2,3$ the matrix $\gamma_\mu$ of the four-component representation is the matrix $\gamma_{\mu-1}$ of that section.

\subsubsection{Characters, twisted spinors and the Wilson operator}\label{sec:characters-twisted-spinors}

A \emph{character} is a triple $\chi=(\chi_1,\chi_2,\chi_3)$ of complex numbers with $|\chi_\mu|=1$ for $\mu=1,2,3$; the antiperiodic character is $\chi^{\mathrm{ap}}=(-1,-1,-1)$. For a character $\chi$ let
\[
\mathcal S_\chi=\left\{\Psi\colon\mathbb Z^3\to\mathbb C^s\ :\ \Psi(x+n\hat\mu)=\chi_\mu\cdot\Psi(x)\text{ for all }x\in\mathbb Z^3\text{ and }\mu\in\{1,2,3\}\right\}.
\]
Its elements are called \emph{spinors twisted by}~$\chi$, or simply \emph{spinors}, and $\mathcal S_\chi$ carries the inner product
\[
\langle\Phi,\Psi\rangle=\sum_{x\in\{0,\dots,n-1\}^3}\Phi(x)^*\cdot\Psi(x)\qquad\left(\Phi,\Psi\in\mathcal S_\chi\right).
\]
An element of $\mathcal S_\chi$ is determined by its values on $\{0,\dots,n-1\}^3$, and these values can be prescribed arbitrarily; hence $\dim\mathcal S_\chi=s\cdot n^3$. Since $|\chi_\mu|=1$, the function $x\mapsto\Phi(x)^*\cdot\Psi(x)$ is periodic with period $n$ in each direction, so that the sum may be taken over any translate of $\{0,\dots,n-1\}^3$.

A one-cochain $A\in C^1$ is read as a function on the links of $\mathbb Z^3$ which is periodic with period $n$ in each direction. The \emph{covariant shifts} are the operators
\[
\left(T_\mu\Psi\right)(x)=\exp\left(iq\cdot A_\mu(x)\right)\cdot\Psi(x+\hat\mu),\qquad\left(T_\mu^*\Psi\right)(x)=\exp\left(-iq\cdot A_\mu(x-\hat\mu)\right)\cdot\Psi(x-\hat\mu)
\]
on $\mathcal S_\chi$; they map $\mathcal S_\chi$ into itself because $A$ is periodic. They are unitary and adjoint to each other: $T_\mu\cdot T_\mu^*=T_\mu^*\cdot T_\mu=1$, and a shift of the summation variable by $\hat\mu$ in $\langle\Phi,T_\mu\Psi\rangle$, together with the periodicity just noted, gives $\langle\Phi,T_\mu\Psi\rangle=\langle T_\mu^*\Phi,\Psi\rangle$. The \emph{Wilson operator}, with Wilson parameter one, is the operator on $\mathcal S_\chi$
\begin{multline}\label{eq:gaussian-bounds-1}
D_A=\sum_{\mu=1}^3\left[\frac12\gamma_\mu\cdot\left(T_\mu-T_\mu^*\right)+\frac12\left(2-T_\mu-T_\mu^*\right)\right]\\
=3-\sum_{\mu=1}^3\left[P^-_\mu\cdot T_\mu+P^+_\mu\cdot T_\mu^*\right],\qquad P^\pm_\mu=\frac12\left(1\pm\gamma_\mu\right),
\end{multline}
in which $\gamma_\mu$ acts on the values of a spinor. Besides on $A$, the operators $T_\mu$ and $D_A$ depend on $q$ and on $\chi$, and $D_A$ depends on the Clifford triple; this is not shown in the notation. The Wilson operator $D_A$ of Section~\ref{sec:lattice-measures-laws}, which is the level-zero case of the operator of Section~\ref{sec:exact-fermion-operators}, is the operator \eqref{eq:gaussian-bounds-1} in the case where the Clifford triple is the four-component representation, $\chi=\chi^{\mathrm{ap}}$ and $q=e\cdot\sqrt{\eta}$, the charge $e$ and the lattice spacing $\eta$ being those of Section~\ref{sec:lattice-measures-laws}. Finally, put
\[
\Delta^\chi(A)=\det\nolimits_{\mathcal S_\chi}\left(D_A+\boldsymbol{\mu}\right)\qquad\left(A\in C^1\right).
\]

\subsubsection{Partition functions with sources}\label{sec:partition-functions-sources}

For a character $\chi$, a two-cochain $\mathfrak{h}\in C^2$ and a one-cochain $b\in C^1$ put
\begin{multline}\label{eq:gaussian-bounds-2}
Z^\chi_\varepsilon(\mathfrak{h},b)=\int_{C^1}\mathrm{d}A\,\exp\left(-\frac12\left\|d_1A+\mathfrak{h}\right\|^2-\frac{\varepsilon}{2}\left\|A+b\right\|^2\right)\cdot\Delta^\chi(A),\\
Z^\chi_\varepsilon(\mathfrak{h})=Z^\chi_\varepsilon(\mathfrak{h},0);
\end{multline}
the integral converges absolutely, since $\varepsilon>0$ and $\Delta^\chi$ is continuous and bounded (Lemma~\ref{lem:gaussian-bounds-1}(i)). A superscript $\mathrm{ap}$ in place of $\chi$ refers to the character $\chi^{\mathrm{ap}}$; for instance $Z^{\mathrm{ap}}_\varepsilon=Z^{\chi^{\mathrm{ap}}}_\varepsilon$ and $\Delta^{\mathrm{ap}}=\Delta^{\chi^{\mathrm{ap}}}$. With $d_1^*\colon C^2\to C^1$ the adjoint of $d_1$, the operator
\[
\mathcal C_\varepsilon=\left(d_1^*d_1+\varepsilon\right)^{-1}\colon C^1\to C^1
\]
is called the \emph{free Proca covariance}; the inverse exists because $d_1^*d_1+\varepsilon\ge\varepsilon>0$.

\subsubsection{The Proca parent and the harmonic-averaged law}\label{sec:proca-parent-harmonic}

Let $\chi$ be a character such that $\Delta^\chi(A)\ge0$ for all $A\in C^1$ and $Z^\chi_\varepsilon(0)>0$. The probability measure
\[
Z^\chi_\varepsilon(0)^{-1}\cdot\exp\left(-\frac12\|d_1A\|^2-\frac{\varepsilon}{2}\|A\|^2\right)\cdot\Delta^\chi(A)\,\mathrm{d}A
\]
on $C^1$ is called the \emph{Proca parent} with character $\chi$, and $\mathsf E^\chi_\varepsilon$ denotes its expectation; $\mathsf E^{\mathrm{ap}}_\varepsilon=\mathsf E^{\chi^{\mathrm{ap}}}_\varepsilon$.

Now let $q\ne0$. For $\theta\in\mathbb R^3$ let $c^\theta\in C^1$ be the constant cochain
\[
c^\theta_\mu(x)=\frac{\theta_\mu}{q\cdot n}\qquad\left(x\in\Lambda,\ \mu\in\{1,2,3\}\right);
\]
it is the cochain $H_\theta$ of Section~\ref{sec:lattice-measures-laws}. Let $\sigma_{V_{\mathrm{t}}}$ be Lebesgue measure on the Euclidean space $V_{\mathrm{t}}$, whose inner product is that of $C^1$. In the four-component representation put
\begin{multline}\label{eq:gaussian-bounds-3}
\overline{\mathsf E}[F]=\frac{\overline Z[F]}{\overline Z[1]},\\
\overline Z[F]=\int_{[-\pi,\pi]^3}\frac{\mathrm{d}\theta}{(2\pi)^3}\int_{V_{\mathrm{t}}}\mathrm{d}\sigma_{V_{\mathrm{t}}}(A)\,F(d_1A)\cdot\exp\left(-\frac12\|d_1A\|^2\right)\cdot\Delta^{\mathrm{ap}}\left(A+c^\theta\right),
\end{multline}
for every Borel function $F\colon C^2\to\mathbb C$ for which the double integral defining $\overline Z[F]$ converges absolutely. By Theorem~\ref{thm:limit-vanishing-proca}, $0<\overline Z[1]<\infty$, and $\overline Z[F]$ converges absolutely whenever $|F(\omega)|\le c_1\cdot\exp\left(c_2\cdot\|\omega\|\right)$ for all $\omega\in C^2$, with constants $c_1,c_2$. Since moreover $\Delta^{\mathrm{ap}}\ge0$ by Theorem~\ref{thm:gaussian-domination-positive}(b), $\overline{\mathsf E}$ is the expectation of a probability measure on $C^2$; this measure is called the \emph{harmonic-averaged law}.

\paragraph{Identification with the averaged law of Section~\ref{sec:lattice-measures-laws}.}
Let $L$, $\bar{m}$, $N$, $\eta=L^{-N}$ and the odd integer $n=R/\eta$ be as in Section~\ref{sec:lattice-measures-laws}, let $e>0$ and $q=e\cdot\sqrt{\eta}$, and let $\boldsymbol{\mu}=\bar{m}\cdot\eta+m_{0,N}$, where $m_{0,N}$ is a real number with $m_{0,N}>-1-\bar{m}\cdot\eta$. Then $\overline{\mathsf E}$ is the expectation of the harmonic-averaged law $\bar{\mu}$ of Section~\ref{sec:lattice-measures-laws}, formed with the bare mass shift $m_{\mathrm{in}}=m_{0,N}$ and an arbitrary vacuum energy $\boldsymbol{\varepsilon}_{\mathrm{in}}$.

Indeed, by Lemma~\ref{lem:harmonic-cochains-coercivity} the map $d_1$ is a linear bijection of $V_{\mathrm{t}}$ onto $\operatorname{ran} d_1$. It carries $\sigma_{V_{\mathrm{t}}}$ to a constant multiple of Lebesgue measure on the Euclidean space $\operatorname{ran} d_1$, and therefore the measure $\exp\left(-\frac12\|d_1A\|^2\right)\mathrm{d}\sigma_{V_{\mathrm{t}}}(A)$ to a constant multiple of the standard Gaussian measure of $\operatorname{ran} d_1$. By its definition in Section~\ref{sec:lattice-measures-laws}, $\bar{\mu}$ is the normalisation of the average, with respect to $(2\pi)^{-3}\,\mathrm{d}\theta$ on $[-\pi,\pi]^3$, of the measures on $\operatorname{ran} d_1$ whose densities with respect to this Gaussian measure are
\[
\omega\mapsto\exp\left(-\boldsymbol{\varepsilon}_{\mathrm{in}}\cdot n^3\right)\cdot\frac{\det\left(D_{A(\omega)+H_\theta}+\bar{m}\cdot\eta+m_{\mathrm{in}}\right)}{\det\left(D_0+\bar{m}\cdot\eta\right)},
\]
where $A(\omega)$ is the element of $V_{\mathrm{t}}$ with $d_1A(\omega)=\omega$. For $m_{\mathrm{in}}=m_{0,N}$ one has $\bar{m}\cdot\eta+m_{\mathrm{in}}=\boldsymbol{\mu}$, and the determinant in the numerator equals $\Delta^{\mathrm{ap}}\left(A(\omega)+c^\theta\right)$, by the identifications of the Wilson operator and of $H_\theta$ in Sections~\ref{sec:characters-twisted-spinors} and~\ref{sec:proca-parent-harmonic}. The remaining factor, namely the vacuum factor $\exp\left(-\boldsymbol{\varepsilon}_{\mathrm{in}}\cdot n^3\right)$ divided by the determinant of the free operator $D_0$ at the mass $\bar{m}\cdot\eta$, depends neither on $\omega$ nor on $\theta$, and it is positive: its denominator is the number $\Delta^{\mathrm{ap}}(0)$ formed with the lattice mass $\bar{m}\cdot\eta>-1$ in place of $\boldsymbol{\mu}$, which is positive by Corollary~\ref{cor:positivity}. Consequently the average of these measures is a positive constant multiple of the functional $F\mapsto\overline Z[F]$ of \eqref{eq:gaussian-bounds-3}, the constant cancels in the normalisation, and the expectation of $F$ under $\bar{\mu}$ is $\overline Z[F]/\overline Z[1]=\overline{\mathsf E}[F]$.

\subsection{Gaussian domination at positive and vanishing Proca mass, and linear functionals of the curvature}\label{sec:gaussian-domination-positive}

\begin{theorem}[Gaussian domination at positive Proca mass]\label{thm:gaussian-domination-positive}
Let $n\ge2$, $\varepsilon>0$, $q\in\mathbb R$ and $\boldsymbol{\mu}>-1$, and let $(\gamma_1,\gamma_2,\gamma_3)$ be an arbitrary Clifford triple.
\begin{enumerate}
\item[(a)] Let $\chi$ be a character for which there are directions $\mu\ne\nu$ with $\chi_\mu=\chi_\nu=-1$. Then $Z^\chi_\varepsilon(0)\ge0$, and every $\mathfrak{h}\in C^2$ satisfies
\[
\left|Z^\chi_\varepsilon(\mathfrak{h})\right|\le Z^\chi_\varepsilon(0).
\]
\item[(b)] In the four-component representation, $\Delta^\chi(A)\ge0$ for every character $\chi$ and every $A\in C^1$; $Z^\chi_\varepsilon(\mathfrak{h},b)\ge0$ for every character $\chi$, every $\mathfrak{h}\in C^2$ and every $b\in C^1$; and $Z^\chi_\varepsilon(0)>0$ whenever $q\ne0$ or $\chi_\mu=-1$ for some direction $\mu$.
\end{enumerate}
\end{theorem}

\begin{theorem}[Laplace transform of the curvature]\label{thm:laplace-transform-curvature}
Let the Clifford triple be the four-component representation, and let $n$, $\varepsilon$, $q$, $\boldsymbol{\mu}$ be as in Theorem~\ref{thm:gaussian-domination-positive}.
Let $\chi$ be a character for which there are directions $\mu\ne\nu$ with $\chi_\mu=\chi_\nu=-1$. By Theorem~\ref{thm:gaussian-domination-positive}(b) the Proca parent with character $\chi$ is defined, and its expectation satisfies, for every $g\in C^2$,
\[
\mathsf E^\chi_\varepsilon\left[\exp\left\langle g,d_1A\right\rangle\right]\le\exp\left(\frac12\|Pg\|^2\right)\le\exp\left(\frac12\|g\|^2\right).
\]
\end{theorem}

\begin{theorem}[The limit of vanishing Proca mass]\label{thm:limit-vanishing-proca}
Let the Clifford triple be the four-component representation, $n\ge2$, $q\ne0$ and $\boldsymbol{\mu}>-1$. Then $\overline Z[1]$ is finite and positive. Let $F\colon C^2\to\mathbb C$ be a Borel function, and assume that there are constants $c_1,c_2$ with
\[
|F(\omega)|\le c_1\cdot\exp\left(c_2\cdot\|\omega\|\right)\qquad\left(\omega\in C^2\right).
\]
Then the double integral that defines $\overline Z[F]$ in \eqref{eq:gaussian-bounds-3} converges absolutely. As $q\ne0$, Theorem~\ref{thm:gaussian-domination-positive}(b) shows that the Proca parent with character $\chi$ is defined for every character $\chi$ and every $\varepsilon>0$, and one has
\begin{equation}\label{eq:gaussian-bounds-4}
\lim_{\varepsilon\downarrow0}\mathsf E^\chi_\varepsilon\left[F(d_1A)\right]=\overline{\mathsf E}[F]\qquad\text{for every character }\chi.
\end{equation}
For every $g\in C^2$
\begin{equation}\label{eq:gaussian-bounds-5}
\overline{\mathsf E}\left[\exp\langle g,\cdot\rangle\right]\le\exp\left(\frac12\|Pg\|^2\right)\le\exp\left(\frac12\|g\|^2\right).
\end{equation}
The free field, which corresponds to $q=0$, has the expectation
\begin{align*}
\mathsf E^{\mathrm{free}}[F]&=\left(\int_{V_{\mathrm{t}}}\mathrm{d}\sigma_{V_{\mathrm{t}}}(A)\,\exp\left(-\frac12\|d_1A\|^2\right)\right)^{-1}\\
&\qquad\cdot\int_{V_{\mathrm{t}}}\mathrm{d}\sigma_{V_{\mathrm{t}}}(A)\,F(d_1A)\cdot\exp\left(-\frac12\|d_1A\|^2\right)
\end{align*}
for every $F$ as above, and $\mathsf E^{\mathrm{free}}\left[\exp\langle g,\cdot\rangle\right]=\exp\left(\frac12\|Pg\|^2\right)$ for every $g\in C^2$. Thus the first inequality in \eqref{eq:gaussian-bounds-5} becomes an equality when $\overline{\mathsf E}$ is replaced by $\mathsf E^{\mathrm{free}}$.
\end{theorem}

\begin{corollary}[Linear functionals of the curvature]\label{cor:linear-functionals-curvature}
Let $\mathsf E$ be either the expectation $\overline{\mathsf E}$ of the harmonic-averaged law, under the hypotheses of Theorem~\ref{thm:limit-vanishing-proca}, or the expectation $\mathsf E^\chi_\varepsilon$ of the Proca parent, under the hypotheses of Theorem~\ref{thm:laplace-transform-curvature}, and let $\mathsf P$ be the corresponding probability measure. Let $F$ denote the curvature: under $\mathsf E^\chi_\varepsilon$ it is the random variable $F=d_1A$ with values in $C^2$; under $\overline{\mathsf E}$, a law on $C^2$ which is concentrated on the range of $d_1$, it is the identity map of $C^2$. For $g\in C^2$ put
\[
X_g=\langle g,F\rangle,\qquad\sigma_g=\|Pg\|.
\]
Then the following hold for every $g\in C^2$.
\begin{enumerate}
\item[(i)] $\mathsf EX_g=0$ and $\mathsf EX_g^2\le\sigma_g^2$. (This is the infrared bound: by the formula for $\mathsf E^{\mathrm{free}}\exp\langle g,\cdot\rangle$ in Theorem~\ref{thm:limit-vanishing-proca}, $\sigma_g^2$ is the second moment of $\langle g,\cdot\rangle$ under the free law $\mathsf E^{\mathrm{free}}$, so that the second moment of $F$ is dominated, as a quadratic form in $g$, by that of the free law.)
\item[(ii)] For every $z\in\mathbb C$,
\[
\left|\mathsf E\exp\left(zX_g\right)\right|\le\exp\left(\frac12\left(\operatorname{Re}z\right)^2\sigma_g^2\right).
\]
\item[(iii)] If $\sigma_g>0$, then for every real $y\ge0$ and every integer $k\ge1$
\[
\mathsf P\left(X_g\ge y\right)\le\exp\left(-\frac{y^2}{2\sigma_g^2}\right),\qquad
\mathsf E\left|X_g\right|^{2k}\le2\cdot k!\cdot\left(2\sigma_g^2\right)^k.
\]
\end{enumerate}
\end{corollary}

\begin{remark}[Tori with unequal sides]\label{rem:tori-unequal-sides}
Let $n_1,n_2,n_3\ge2$ be integers of arbitrary parities, and let $\Lambda$ now denote the torus $\prod_{\mu=1}^3\mathbb Z/n_\mu\mathbb Z$, so that $\#|\Lambda|=n_1n_2n_3$. No argument of Sections~\ref{sec:cochains-periodic-lattice}--\ref{sec:proca-parent-harmonic} and~\ref{sec:gaussian-domination-positive}--\ref{sec:vanishing-proca-mass} makes use of the fact that the three sides of the cubic torus have the same length; in Sections~\ref{sec:cochains-periodic-lattice}--\ref{sec:proca-parent-harmonic} and~\ref{sec:determinant-wilson-operator}--\ref{sec:vanishing-proca-mass}, and in the statements of Theorems~\ref{thm:gaussian-domination-positive}--\ref{thm:limit-vanishing-proca} and of Corollary~\ref{cor:linear-functionals-curvature}, the common length $n$ of the three sides plays no part except in the condition $n\ge2$ and in the sections, lemmas and proofs named in (1)--(3) below. Consequently, when the substitutions (1)--(3) are carried out in the statements and in the proofs, the arguments given for the cubic torus establish, for the torus $\Lambda$, every statement of Sections~\ref{sec:determinant-wilson-operator}--\ref{sec:vanishing-proca-mass} as well as Theorems~\ref{thm:gaussian-domination-positive}--\ref{thm:limit-vanishing-proca} and Corollary~\ref{cor:linear-functionals-curvature}.
\begin{enumerate}
\item[(1)] In Sections~\ref{sec:cochains-periodic-lattice}, \ref{sec:characters-twisted-spinors} and~\ref{sec:proca-parent-harmonic}, in Lemmas~\ref{lem:gaussian-bounds-1} and~\ref{lem:gaussian-decay-continuity} together with their proofs, and in Section~\ref{sec:vanishing-proca-mass}: $n^3$ is replaced by $\#|\Lambda|$; the condition that defines $\mathcal S_\chi$ becomes $\Psi(x+n_\mu\hat\mu)=\chi_\mu\Psi(x)$, and the fundamental domain $\{0,\dots,n-1\}^3$ becomes $\prod_{\mu=1}^3\{0,\dots,n_\mu-1\}$; one puts $c^\theta_\mu=\theta_\mu/(qn_\mu)$; and in the proof of Lemma~\ref{lem:gaussian-bounds-1}(iv) the operator $\mathsf M$ is
\[
(\mathsf M\Psi)(x)=\exp\left(i\sum_{\mu=1}^3\frac{\vartheta_\mu x_\mu}{n_\mu}\right)\cdot\Psi(x).
\]
\item[(2)] In Sections~\ref{sec:slices-slice-operators}--\ref{sec:factorisation-slices} and~\ref{sec:second-quantisation-trace}--\ref{sec:inequalities-mirrors}, and in the proof of Theorem~\ref{thm:gaussian-domination-positive} in Section~\ref{sec:proof-gaussian-domination}, with $\nu$ the time direction: $n_\nu$, which is the number of slices, stands in place of $n$; $\Sigma=\prod_{i\ne\nu}\mathbb Z/n_i\mathbb Z$, $v=\#|\Lambda|/n_\nu$, $l=2v$ and $r=\rho v$; the condition that defines $\mathcal V$ becomes $\phi(\mathbf x+n_i\hat\imath)=\chi_i\phi(\mathbf x)$; and in Sections~\ref{sec:slices-slice-operators} and~\ref{sec:second-quantisation-trace} a slice has $\prod_{i\ne\nu}\{0,\dots,n_i-1\}$ as its fundamental domain.
\item[(3)] In Section~\ref{sec:vanishing-proca-mass}: the argument $k$ of a discrete Fourier transform has the components $k_\mu\in\frac{2\pi}{n_\mu}\left(\mathbb Z/n_\mu\mathbb Z\right)$, $\mu=1,2,3$; one has $\|c\|^2=\#|\Lambda|\cdot|c|^2$ and $\theta_\mu=qn_\mu c_\mu$;
\[
\gamma_\varepsilon(\mathrm dc)=\left(\frac{\varepsilon\cdot\#|\Lambda|}{2\pi}\right)^{\frac32}\exp\left(-\frac{\varepsilon}{2}\,\#|\Lambda|\cdot|c|^2\right)\mathrm d^3c;
\]
for every $k\in\mathbb Z^3$,
\[
\int_{\mathbb R^3}\gamma_\varepsilon(\mathrm dc)\,\exp\left(i\sum_{\mu=1}^3k_\mu qn_\mu c_\mu\right)=\exp\left(-\frac{q^2}{2\varepsilon\cdot\#|\Lambda|}\sum_{\mu=1}^3n_\mu^2k_\mu^2\right);
\]
and
\[
\delta_\varepsilon=\sum_{k\in\mathbb Z^3\setminus\{0\}}\exp\left(-\frac{q^2}{2\varepsilon\cdot\#|\Lambda|}\sum_{\mu=1}^3n_\mu^2k_\mu^2\right).
\]
\end{enumerate}
\end{remark}

\subsection{The determinant of the Wilson operator}\label{sec:determinant-wilson-operator}

\subsubsection{Bounds, invariances and positivity}\label{sec:bounds-invariances-positivity}

\begin{lemma}\label{lem:gaussian-bounds-1}
Let $A\in C^1$, and let $\chi$ be a character.
\begin{enumerate}
\item[(i)] One has $\left|\Delta^\chi(A)\right|\le\left(9+|\boldsymbol{\mu}|\right)^{sn^3}$, and the function $A\mapsto\Delta^\chi(A)$ on $C^1$ is continuous.
\item[(ii)] For every $\lambda\in C^0$, $\Delta^\chi\left(A+d_0\lambda\right)=\Delta^\chi(A)$.
\item[(iii)] If $q\ne0$, then $\Delta^\chi(A+\beta)=\Delta^\chi(A)$ for every $\beta\in C^1$ with values in $\frac{2\pi}{q}\mathbb Z$.
\item[(iv)] Let $q\ne0$. For $\vartheta\in\mathbb R^3$ let $\chi\exp(i\vartheta)$ denote the character with the components $\chi_\mu\cdot\exp(i\vartheta_\mu)$, $\mu=1,2,3$. Then
\[
\Delta^{\chi\exp(i\vartheta)}(A)=\Delta^\chi\left(A+c^\vartheta\right)\qquad\left(\vartheta\in\mathbb R^3\right).
\]
Moreover, the function $\theta\mapsto\Delta^{\mathrm{ap}}\left(A+c^\theta\right)$ on $\mathbb R^3$ is a trigonometric polynomial: there are complex numbers $\hat f_k$, $k\in\mathbb Z^3$, of which only finitely many differ from $0$, such that
\[
\Delta^{\mathrm{ap}}\left(A+c^\theta\right)=\sum_{k\in\mathbb Z^3}\hat f_k\cdot\exp(ik\cdot\theta)\qquad\left(\theta\in\mathbb R^3\right).
\]
\item[(v)] In the four-component representation,
\[
\Delta^\chi(A)=\left|\det\left(D^{(2)}_A+\boldsymbol{\mu}\right)\right|^2\ge0,
\]
where $D^{(2)}_A$ is the operator \eqref{eq:gaussian-bounds-1} with each $\gamma_\mu$ replaced by the Pauli matrix $\sigma_\mu$, acting on the space $\mathcal S^{(2)}_\chi$ of the functions $\Psi:\mathbb Z^3\to\mathbb C^2$ with $\Psi(x+n\hat\mu)=\chi_\mu\Psi(x)$ for $\mu=1,2,3$, and the determinant is taken on $\mathcal S^{(2)}_\chi$.
\end{enumerate}
\end{lemma}

\begin{proof}
For $A\in C^1$ and $\mu\in\{1,2,3\}$ let $T^A_\mu$ denote the covariant shift $T_\mu$ of Section~\ref{sec:characters-twisted-spinors} formed with $A$, so that $D_A$ is built from the operators $T^A_\mu$ and $\left(T^A_\mu\right)^*$ by \eqref{eq:gaussian-bounds-1}. Norms of operators on $\mathcal S_\chi$ refer to the inner product of Section~\ref{sec:characters-twisted-spinors}.

(i) The matrices $\gamma_\mu$ are Hermitian with $\gamma_\mu^2=1$, hence unitary, and the shifts $T_\mu$ are unitary. Therefore
\[
\left\|\tfrac12\gamma_\mu\left(T_\mu-T_\mu^*\right)\right\|\le1,\qquad\left\|\tfrac12\left(2-T_\mu-T_\mu^*\right)\right\|\le2,
\]
and summation over the three directions gives $\left\|D_A+\boldsymbol{\mu}\right\|\le9+|\boldsymbol{\mu}|$. For a linear map $M$ of a finite-dimensional Hilbert space $\mathcal H$ into itself one has $|\det M|\le\|M\|^{\dim\mathcal H}$, because $|\det M|$ is the product of the singular values of $M$, each of which is at most $\|M\|$. Since $\dim\mathcal S_\chi=sn^3$, the bound follows. An element of $\mathcal S_\chi$ is determined by its values on $\{0,\dots,n-1\}^3$; let $\mathcal S_\chi$ carry the basis $\left(\Psi_{y,a}\right)$ indexed by $(y,a)\in\{0,\dots,n-1\}^3\times\{1,\dots,s\}$, where $\Psi_{y,a}(y)$ is the $a$-th standard basis vector of $\mathbb C^s$ and $\Psi_{y,a}(y')=0$ for the other points $y'$ of $\{0,\dots,n-1\}^3$. In this basis the matrix entries of $D_A$ are continuous functions of $A$, and $\Delta^\chi(A)$ is a polynomial in these entries. Hence $A\mapsto\Delta^\chi(A)$ is continuous.

(ii) Let $\lambda\in C^0$, read as a periodic function on $\mathbb Z^3$, and let $(\mathsf g\Psi)(x)=\exp\left(iq\lambda(x)\right)\cdot\Psi(x)$. As $\lambda$ is periodic, $\mathsf g$ maps $\mathcal S_\chi$ into itself, and it is unitary. For $\Psi\in\mathcal S_\chi$ and $x\in\mathbb Z^3$,
\[
\left(\mathsf g^{-1}T^A_\mu\mathsf g\Psi\right)(x)=\exp\left(iq\left(A_\mu(x)+\lambda(x+\hat\mu)-\lambda(x)\right)\right)\cdot\Psi(x+\hat\mu)=\left(T^{A+d_0\lambda}_\mu\Psi\right)(x),
\]
that is, $\mathsf g^{-1}T^A_\mu\mathsf g=T^{A+d_0\lambda}_\mu$; taking adjoints, and using that $\mathsf g$ is unitary, $\mathsf g^{-1}\left(T^A_\mu\right)^*\mathsf g=\left(T^{A+d_0\lambda}_\mu\right)^*$. Since $\mathsf g$ commutes with the matrices $\gamma_\mu$, it follows that $\mathsf g^{-1}D_A\mathsf g=D_{A+d_0\lambda}$, and therefore $\Delta^\chi(A)=\Delta^\chi\left(A+d_0\lambda\right)$.

(iii) Let $q\ne0$, and let $\beta\in C^1$ take values in $\frac{2\pi}{q}\mathbb Z$. For every link $(x,\mu)$ of $\mathbb Z^3$ the number $q\beta_\mu(x)$ lies in $2\pi\mathbb Z$, so that $\exp\left(\pm iq\left(A_\mu(x)+\beta_\mu(x)\right)\right)=\exp\left(\pm iqA_\mu(x)\right)$. Each of the operators $T^A_\mu$, $\left(T^A_\mu\right)^*$ shifts the argument of a spinor by $\pm\hat\mu$ and multiplies the shifted spinor by factors of this form; therefore $T^{A+\beta}_\mu=T^A_\mu$ and $\left(T^{A+\beta}_\mu\right)^*=\left(T^A_\mu\right)^*$ for $\mu=1,2,3$. Hence $D_{A+\beta}=D_A$, and $\Delta^\chi(A+\beta)=\Delta^\chi(A)$.

(iv) Let $(\mathsf M\Psi)(x)=\exp\left(\frac{i\,\vartheta\cdot x}{n}\right)\cdot\Psi(x)$. For $\Psi\in\mathcal S_\chi$ one has $(\mathsf M\Psi)(x+n\hat\mu)=\chi_\mu\exp(i\vartheta_\mu)\cdot(\mathsf M\Psi)(x)$, so $\mathsf M$ maps $\mathcal S_\chi$ unitarily onto $\mathcal S_{\chi\exp(i\vartheta)}$. For $\Psi\in\mathcal S_\chi$, with $T^A_\mu$ acting on $\mathcal S_{\chi\exp(i\vartheta)}$,
\[
\left(\mathsf M^{-1}T^A_\mu\mathsf M\Psi\right)(x)=\exp\left(iq\left(A_\mu(x)+\frac{\vartheta_\mu}{qn}\right)\right)\cdot\Psi(x+\hat\mu)=\left(T^{A+c^\vartheta}_\mu\Psi\right)(x).
\]
Taking adjoints as in (ii), $\mathsf M^{-1}D_A\mathsf M=D_{A+c^\vartheta}$, where $D_A$ acts on $\mathcal S_{\chi\exp(i\vartheta)}$ and $D_{A+c^\vartheta}$ acts on $\mathcal S_\chi$; this gives the first assertion. Now let $\Psi\in\mathcal S_\chi$ and $x\in\{0,\dots,n-1\}^3$. By \eqref{eq:gaussian-bounds-1} and the definition of the covariant shifts,
\begin{align*}
(D_A\Psi)(x)&=3\Psi(x)-\sum_{\mu=1}^3\exp\left(iqA_\mu(x)\right)\cdot P^-_\mu\cdot\Psi(x+\hat\mu)\\
&\qquad-\sum_{\mu=1}^3\exp\left(-iqA_\mu(x-\hat\mu)\right)\cdot P^+_\mu\cdot\Psi(x-\hat\mu),
\end{align*}
and no coefficient in this formula involves $\chi$. If $x_\mu\le n-2$, the point $x+\hat\mu$ lies in $\{0,\dots,n-1\}^3$; if $x_\mu=n-1$, the definition of $\mathcal S_\chi$ gives $\Psi(x+\hat\mu)=\chi_\mu\cdot\Psi(x+\hat\mu-n\hat\mu)$, and $x+\hat\mu-n\hat\mu$ lies in $\{0,\dots,n-1\}^3$. Likewise, $x-\hat\mu$ lies in $\{0,\dots,n-1\}^3$ if $x_\mu\ge1$, while for $x_\mu=0$ one has $\Psi(x-\hat\mu)=\overline{\chi_\mu}\cdot\Psi(x-\hat\mu+n\hat\mu)$ with $x-\hat\mu+n\hat\mu\in\{0,\dots,n-1\}^3$; here $\overline{\chi_\mu}=\chi_\mu^{-1}$, as $|\chi_\mu|=1$. The coordinates of an element of $\mathcal S_\chi$ in the basis $\left(\Psi_{y,a}\right)$ of (i) are the components of its values at the points of $\{0,\dots,n-1\}^3$. It follows that every entry of the matrix of $D_A+\boldsymbol{\mu}$ in this basis is a Laurent polynomial in $\chi_1,\chi_2,\chi_3$ with coefficients independent of $\chi$, and hence so is its determinant $\Delta^\chi(A)$. By the first assertion, $\Delta^{\mathrm{ap}}\left(A+c^\theta\right)=\Delta^{\chi^{\mathrm{ap}}\exp(i\theta)}(A)$ is obtained from this Laurent polynomial by the substitution $\chi_\mu=-\exp(i\theta_\mu)$, and is therefore a trigonometric polynomial in $\theta$.

(v) In the four-component representation $\gamma_\mu=\operatorname{diag}\left(\sigma_\mu,-\sigma_\mu\right)$; accordingly $\mathcal S_\chi=\mathcal S^{(2)}_\chi\oplus\mathcal S^{(2)}_\chi$, and
\[
D_A+\boldsymbol{\mu}=\operatorname{diag}\left(D^{(2)}_A+\boldsymbol{\mu},\ \widetilde D+\boldsymbol{\mu}\right),\qquad
\widetilde D=\sum_{\mu=1}^3\left[-\tfrac12\sigma_\mu\left(T_\mu-T_\mu^*\right)+\tfrac12\left(2-T_\mu-T_\mu^*\right)\right].
\]
Since $T_\mu^*$ is the adjoint of $T_\mu$, the operator $T_\mu-T_\mu^*$ is skew-adjoint and $2-T_\mu-T_\mu^*$ is self-adjoint. Moreover, $T_\mu-T_\mu^*$ commutes with $\sigma_\mu$, because $\sigma_\mu$ acts on the values of a spinor, whereas $T_\mu$ and $T_\mu^*$ shift its argument and multiply by scalar factors. Together with $\sigma_\mu^*=\sigma_\mu$ this gives $\left(\sigma_\mu\left(T_\mu-T_\mu^*\right)\right)^*=-\left(T_\mu-T_\mu^*\right)\sigma_\mu=-\sigma_\mu\left(T_\mu-T_\mu^*\right)$, and therefore $\widetilde D=\left(D^{(2)}_A\right)^*$. As $\boldsymbol{\mu}$ is real, $\det\left(\widetilde D+\boldsymbol{\mu}\right)=\overline{\det\left(D^{(2)}_A+\boldsymbol{\mu}\right)}$, and
\[
\Delta^\chi(A)=\det\left(D^{(2)}_A+\boldsymbol{\mu}\right)\cdot\det\left(\widetilde D+\boldsymbol{\mu}\right)=\left|\det\left(D^{(2)}_A+\boldsymbol{\mu}\right)\right|^2.
\]
\end{proof}

\subsubsection{Slices and slice operators}\label{sec:slices-slice-operators}

Let $\nu\in\{1,2,3\}$ be a direction, called the \emph{time direction}, let $j<j'$ be the two other directions, called \emph{spatial}, and let $\chi$ be a character. A site is written $x=(t,\mathbf x)$ with $t=x_\nu\in\mathbb Z/n\mathbb Z$ and $\mathbf x=(x_j,x_{j'})\in\Sigma:=(\mathbb Z/n\mathbb Z)^2$. For $i\in\{j,j'\}$ the unit vector of the direction $i$ is written $\hat\imath$, in $\Lambda$ as well as in $\Sigma$ and in $\mathbb Z^2$. For $t\in\mathbb Z/n\mathbb Z$ the \emph{slice} $t$ consists of the $v:=n^2$ sites $x$ with $x_\nu=t$, of the $l:=2n^2$ links $(x,i)$ with $x_\nu=t$ and $i\in\{j,j'\}$, and of the $v$ plaquettes $(x,jj')$ with $x_\nu=t$. The space $\mathbb R^v$ is identified with the space of real functions on $\Sigma$, and $\mathbb R^l$ with the space of real functions on $\Sigma\times\{j,j'\}$; both carry the Euclidean inner product, and $|\cdot|$ denotes the corresponding norm.

\paragraph{Fields on the slices.}
For $A\in C^1$ and $t\in\mathbb Z/n\mathbb Z$ let $A_t\in\mathbb R^l$ and $a_t\in\mathbb R^v$ be defined by
\[
A_t(\mathbf x,i):=A_i(t,\mathbf x)\quad\left(i\in\{j,j'\}\right),\qquad a_t(\mathbf x):=A_\nu(t,\mathbf x).
\]
In the same way a cochain $b\in C^1$ defines $b^s_t\in\mathbb R^l$ and $b^0_t\in\mathbb R^v$, namely $b^s_t(\mathbf x,i):=b_i(t,\mathbf x)$ and $b^0_t(\mathbf x):=b_\nu(t,\mathbf x)$. The spatial differentials $d_s\colon\mathbb R^v\to\mathbb R^l$ and $d_s\colon\mathbb R^l\to\mathbb R^v$ are defined by
\begin{align*}
(d_sa)(\mathbf x,i)&:=a(\mathbf x+\hat\imath)-a(\mathbf x)&&\left(a\in\mathbb R^v\right),\\
(d_sA)(\mathbf x)&:=A(\mathbf x,j)+A(\mathbf x+\hat\jmath,j')-A(\mathbf x+\widehat{j'},j)-A(\mathbf x,j')&&\left(A\in\mathbb R^l\right),
\end{align*}
and $d_s^*\colon\mathbb R^l\to\mathbb R^v$ denotes the adjoint of the first of them. For $\mathfrak{h}\in C^2$ and $t\in\mathbb Z/n\mathbb Z$, the \emph{magnetic part} $\mathfrak{h}^m_t\in\mathbb R^v$ and the \emph{electric part} $\mathfrak{h}^e_t\in\mathbb R^l$ of $\mathfrak{h}$ are
\[
\mathfrak{h}^m_t(\mathbf x):=\mathfrak{h}_{jj'}(t,\mathbf x),\qquad
\mathfrak{h}^e_t(\mathbf x,i):=\begin{cases}\mathfrak{h}_{\nu i}(t,\mathbf x)&\text{if }\nu<i,\\ -\mathfrak{h}_{i\nu}(t,\mathbf x)&\text{if }i<\nu.\end{cases}
\]
The definition of $d_1$ in Section~\ref{sec:cochains-periodic-lattice} gives $(d_1A)_{jj'}(t,\mathbf x)=(d_sA_t)(\mathbf x)$, and, for $x=(t,\mathbf x)$ and $i\in\{j,j'\}$ with $\nu<i$,
\[
(d_1A)_{\nu i}(t,\mathbf x)=\left[A_i(x+\hat\nu)-A_i(x)\right]-\left[A_\nu(x+\hat\imath)-A_\nu(x)\right]=\left(A_{t+1}-A_t-d_sa_t\right)(\mathbf x,i);
\]
for $i\in\{j,j'\}$ with $i<\nu$ the same computation gives $(d_1A)_{i\nu}(t,\mathbf x)=-\left(A_{t+1}-A_t-d_sa_t\right)(\mathbf x,i)$, and this sign is matched by the sign in the definition of $\mathfrak{h}^e_t$. Consequently, for $A,b\in C^1$ and $\mathfrak{h}\in C^2$, the indices $t$ being taken in $\mathbb Z/n\mathbb Z$,
\begin{gather}
\|d_1A+\mathfrak{h}\|^2=\sum_{t\in\mathbb Z/n\mathbb Z}\left(\left|A_{t+1}-A_t-d_sa_t+\mathfrak{h}^e_t\right|^2+\left|d_sA_t+\mathfrak{h}^m_t\right|^2\right),\notag\\
\|A+b\|^2=\sum_{t\in\mathbb Z/n\mathbb Z}\left(\left|A_t+b^s_t\right|^2+\left|a_t+b^0_t\right|^2\right).\label{eq:gaussian-bounds-6}
\end{gather}

\paragraph{Slice data and layer data.}
For $\mathfrak{h}\in C^2$, $b\in C^1$ and $t\in\mathbb Z/n\mathbb Z$, the pair $\omega_t:=\left(\mathfrak{h}^m_t,b^s_t\right)\in\mathbb R^v\times\mathbb R^l$ is called the \emph{slice data}, and the pair $\lambda_t:=\left(\mathfrak{h}^e_t,b^0_t\right)\in\mathbb R^l\times\mathbb R^v$ the \emph{layer data}, at time $t$. Here the \emph{layer} $t$ is the set of the time-like links and plaquettes between the slices $t$ and $t+1$, that is, of the links $(x,\nu)$ and of the plaquettes spanned at $x$ by $\nu$ and a direction $i\in\{j,j'\}$, with $x_\nu=t$. The map $(\mathfrak{h},b)\mapsto(\omega_t,\lambda_t)_{t\in\mathbb Z/n\mathbb Z}$ is a bijection of $C^2\times C^1$ onto $\left(\mathbb R^v\times\mathbb R^l\times\mathbb R^l\times\mathbb R^v\right)^{\mathbb Z/n\mathbb Z}$; it permutes the coordinates and changes the signs of some of them.

\paragraph{A spinor basis adapted to the time direction.}
The matrix $\gamma_\nu$ is Hermitian with $\gamma_\nu^2=1$, so $\mathbb C^s$ is the orthogonal sum of its eigenspaces for the eigenvalues $1$ and $-1$. The matrix $\gamma_j$ is invertible, because $\gamma_j^2=1$, and it anticommutes with $\gamma_\nu$; it therefore maps each of these two eigenspaces injectively into the other. Hence the two eigenspaces have the same dimension; consequently $s$ is even, and both have the dimension $\rho=s/2$. Let $w$ be a unitary $s\times s$ matrix with $w\gamma_\nu w^*=\operatorname{diag}(1_\rho,-1_\rho)$. For $i\in\{j,j'\}$ the Hermitian matrix $w\gamma_iw^*$ anticommutes with $\operatorname{diag}(1_\rho,-1_\rho)$, so that its diagonal blocks vanish; being Hermitian, it has the form
\[
w\gamma_iw^*=\begin{pmatrix}0&\sigma_i\\ \sigma_i^*&0\end{pmatrix},\qquad\|\sigma_i\|\le1,
\]
with $\rho\times\rho$ matrices $\sigma_i$; the bound holds because $w\gamma_iw^*$, a Hermitian matrix with square $1$, is unitary. For $i\in\{j,j'\}$ the notation $\sigma_i$ refers to the block in this display, and not to the Pauli matrix of Section~\ref{sec:clifford-triples} with the same index. The matrices $w\gamma_\mu w^*$, $\mu=1,2,3$, again form a Clifford triple. The unitary operator $1\otimes w$ on $\mathcal S_\chi$, given by $\left((1\otimes w)\Psi\right)(x):=w\Psi(x)$, commutes with every $T_\mu$, and $(1\otimes w)D_A(1\otimes w)^*$ is the operator \eqref{eq:gaussian-bounds-1} formed with the matrices $w\gamma_\mu w^*$. A determinant being invariant under conjugation, the replacement of $(\gamma_\mu)$ by $(w\gamma_\mu w^*)$ does not change $\Delta^\chi$. From now on $\gamma_\nu$ and $\gamma_i$, $i\in\{j,j'\}$, are assumed to have the above form.

\paragraph{Operators on a slice.}
Let
\[
\mathcal V:=\left\{\phi\colon\mathbb Z^2\to\mathbb C^\rho:\ \phi(\mathbf x+n\hat\imath)=\chi_i\,\phi(\mathbf x)\text{ for all }\mathbf x\in\mathbb Z^2\text{ and }i\in\{j,j'\}\right\},
\]
the coordinates of $\mathbb Z^2$ being indexed by $j$ and $j'$. An element of $\mathcal V$ is determined by its values on $\{0,\dots,n-1\}^2$, and these values may be prescribed arbitrarily; hence
\[
r:=\dim\mathcal V=\rho n^2.
\]
The formula
\[
\langle\phi,\phi'\rangle:=\sum_{\mathbf x\in\{0,\dots,n-1\}^2}\phi(\mathbf x)^*\phi'(\mathbf x)\qquad\left(\phi,\phi'\in\mathcal V\right)
\]
defines an inner product on $\mathcal V$; on each of the spaces $\mathcal K:=\mathcal V\oplus\mathcal V$ and $\mathcal V^n$, the inner product of two elements is the sum of the inner products of their components. Since $|\chi_i|=1$ for $i\in\{j,j'\}$, the function $\mathbf x\mapsto\phi(\mathbf x)^*\phi'(\mathbf x)$ on $\mathbb Z^2$ is periodic, in the sense that it factors through the quotient map $\mathbb Z^2\to\Sigma$; consequently the sum that defines $\langle\phi,\phi'\rangle$ may be taken over any system of representatives of $\Sigma$ in $\mathbb Z^2$. Through this quotient map, functions on $\Sigma$ and on $\Sigma\times\{j,j'\}$ are read as periodic functions on $\mathbb Z^2$ and on $\mathbb Z^2\times\{j,j'\}$. For a slice field $A\in\mathbb R^l$, for $a\in\mathbb R^v$ and for $i\in\{j,j'\}$ let
\[
\left(\tau_i(A)\phi\right)(\mathbf x):=\exp\left(iqA(\mathbf x,i)\right)\cdot\phi(\mathbf x+\hat\imath),\qquad\left(u(a)\phi\right)(\mathbf x):=\exp\left(iqa(\mathbf x)\right)\cdot\phi(\mathbf x);
\]
these operators map $\mathcal V$ into itself, $A$ and $a$ being periodic, and they are unitary on $\mathcal V$ by the periodicity just noted. Put
\begin{equation}\label{eq:gaussian-bounds-7}
B(A):=(3+\boldsymbol{\mu})-\frac12\sum_{i\ne\nu}\left(\tau_i+\tau_i^*\right),\qquad C(A):=\frac12\sum_{i\ne\nu}\sigma_i\left(\tau_i-\tau_i^*\right),
\end{equation}
where $\tau_i=\tau_i(A)$ and the matrices $\sigma_i$ act on the values of the elements of $\mathcal V$, and, with $B=B(A)$ and $C=C(A)$ (the operator $B$ is invertible by \eqref{eq:gaussian-bounds-9} below),
\begin{align}
S(A)&:=\begin{pmatrix}B^{-1}&-B^{-1}C\\ -C^*B^{-1}&B+C^*B^{-1}C\end{pmatrix}\notag\\
&=\begin{pmatrix}1&0\\ -C^*&1\end{pmatrix}\begin{pmatrix}B^{-1}&0\\ 0&B\end{pmatrix}\begin{pmatrix}1&-C\\ 0&1\end{pmatrix}\qquad\text{on }\mathcal K=\mathcal V\oplus\mathcal V.\label{eq:gaussian-bounds-8}
\end{align}
Each $\tau_i$ is unitary, so $-1\le\frac12\left(\tau_i+\tau_i^*\right)\le1$; it follows that
\begin{multline}\label{eq:gaussian-bounds-9}
0<1+\boldsymbol{\mu}\le B(A)\le5+\boldsymbol{\mu},\qquad\|C(A)\|\le2,\\
S(A)=S(A)^*>0,\qquad\|S(A)\|\le s_*:=9\max\left\{5+\boldsymbol{\mu},(1+\boldsymbol{\mu})^{-1}\right\}.
\end{multline}
Indeed, $B(A)$ is self-adjoint, its two bounds follow by adding the bounds for the two directions $i\in\{j,j'\}$, and $1+\boldsymbol{\mu}>0$ because $\boldsymbol{\mu}>-1$. The bound for $C(A)$ follows from $\|\sigma_i\|\le1$ and $\|\tau_i-\tau_i^*\|\le2$. The second form in \eqref{eq:gaussian-bounds-8} is $\mathsf L^*\cdot\operatorname{diag}\left(B^{-1},B\right)\cdot\mathsf L$ with
\[
\mathsf L:=\begin{pmatrix}1&-C\\ 0&1\end{pmatrix},\qquad\|\mathsf L\|\le1+\|C\|\le3,
\]
and $\mathsf L$ is invertible. As $\operatorname{diag}\left(B^{-1},B\right)$ is self-adjoint and positive definite, so is $S(A)$, and $\|S(A)\|\le\|\mathsf L\|^2\cdot\max\left\{\|B^{-1}\|,\|B\|\right\}\le s_*$.

\subsubsection{Determinants of cyclic block operators}\label{sec:determinants-cyclic-block}

\begin{lemma}[Cyclic block determinants]\label{lem:cyclic-block-determinants}
Let $\mathcal K$ be a finite-dimensional vector space, let $n\ge1$, let $g_0,\dots,g_{n-1}\in\operatorname{End}\mathcal K$, and let $\mathbb W$ be the endomorphism of $\mathcal K^n$ given by
\[
(\mathbb Ww)_t:=g_{t-1}w_{t-1}\quad(1\le t\le n-1),\qquad(\mathbb Ww)_0:=g_{n-1}w_{n-1},
\]
for $w=(w_0,\dots,w_{n-1})\in\mathcal K^n$. Then
\[
\det(1-\mathbb W)=\det\left(1-g_{n-1}g_{n-2}\cdots g_0\right).
\]
\end{lemma}

\begin{proof}
In a basis of $\mathcal K$ the endomorphisms of $\mathcal K$ and of $\mathcal K^n$ are square matrices. For a square matrix $W$ and for $z\in\mathbb C$ of small modulus one has the expansion
\[
\log\det(1-zW)=-\sum_{k\ge1}\frac{z^k}{k}\cdot\operatorname{tr}W^k,
\]
the logarithm being the branch that vanishes at $z=0$. The operator $\mathbb W^k$ maps the summand of index $t$ of $\mathcal K^n$ into the summand whose index is congruent to $t+k$ modulo $n$; hence $\operatorname{tr}\mathbb W^k=0$ unless $n$ divides $k$. The operator $\mathbb W^n$ is block diagonal, with the diagonal blocks
\[
G_t:=g_{t-1}\cdots g_0\,g_{n-1}\cdots g_t\qquad(0\le t\le n-1),
\]
a product without factors being $1$; in particular $G_0=g_{n-1}\cdots g_0$. Each $G_t$ is a cyclic rotation of $G_0$: one has $G_t=XY$ and $G_0=YX$ with $X:=g_{t-1}\cdots g_0$ and $Y:=g_{n-1}\cdots g_t$. By the cyclicity of the trace, $\operatorname{tr}G_t^k=\operatorname{tr}G_0^k$, and therefore $\operatorname{tr}\mathbb W^{nk}=n\cdot\operatorname{tr}G_0^k$. It follows that, for $|z|$ small,
\[
\log\det(1-z\mathbb W)=-\sum_{k\ge1}\frac{z^k}{k}\cdot\operatorname{tr}\mathbb W^k=-\sum_{k\ge1}\frac{z^{nk}}{k}\cdot\operatorname{tr}G_0^k=\log\det\left(1-z^nG_0\right).
\]
Both $\det(1-z\mathbb W)$ and $\det\left(1-z^nG_0\right)$ are polynomials in $z$, and they agree for all $z$ of small modulus. Hence they agree for all $z\in\mathbb C$, in particular for $z=1$.
\end{proof}

\subsubsection{Factorisation over the slices}\label{sec:factorisation-slices}

\begin{lemma}[Slice factorisation of the determinant]\label{lem:slice-factorisation-determinant}
Let $A\in C^1$, let $\chi$ be a character and $\nu$ a direction, and let the notation of Section~\ref{sec:slices-slice-operators} be formed with the time direction $\nu$ and the character $\chi$. Then
\begin{align}
\Delta^\chi(A)&=(-\chi_\nu)^r\cdot\prod_{t=0}^{n-1}\left[\det B(A_t)\cdot\det u(a_t)\right]\cdot\det\nolimits_{\mathcal K}\left(1-\bar\chi_\nu\,g_{n-1}g_{n-2}\cdots g_0\right),\label{eq:gaussian-bounds-10}\\
g_t&:=\left(u(a_t)^*\oplus u(a_t)^*\right)\cdot S(A_t)\qquad(0\le t\le n-1),\notag
\end{align}
the determinants of $B(A_t)$ and of $u(a_t)$ being taken on $\mathcal V$.
\end{lemma}

\begin{proof}
For $\Psi\in\mathcal S_\chi$, $t\in\mathbb Z$ and $\mathbf x\in\mathbb Z^2$ let $\psi^+_t(\mathbf x)$ and $\psi^-_t(\mathbf x)$ be the upper and the lower $\rho$ components of $\Psi(t,\mathbf x)$, in the basis fixed in Section~\ref{sec:slices-slice-operators}. Then $\psi^\pm_t\in\mathcal V$ and $\psi^\pm_{t+n}=\chi_\nu\psi^\pm_t$, and the map $\Psi\mapsto(\psi^+,\psi^-)$, where $\psi^\pm:=\left(\psi^\pm_t\right)_{t=0}^{n-1}\in\mathcal V^n$, identifies $\mathcal S_\chi$ with $\mathcal V^n\oplus\mathcal V^n$; it is unitary for the inner products of Sections~\ref{sec:characters-twisted-spinors} and~\ref{sec:slices-slice-operators}. On $\mathcal V^n$ let
\[
\mathbb B:=\operatorname{diag}\left(B(A_0),\dots,B(A_{n-1})\right),\qquad\mathbb C:=\operatorname{diag}\left(C(A_0),\dots,C(A_{n-1})\right),
\]
and let $\mathbb T$ be given by
\[
(\mathbb T\psi)_t:=u(a_t)\psi_{t+1}\qquad(0\le t\le n-1),\qquad\text{where }\psi_n:=\chi_\nu\psi_0 .
\]
The operator $\mathbb T$ is unitary, being a cyclic shift of the components, one of them multiplied by the number $\chi_\nu$ of modulus one, followed by the unitary operators $u(a_t)$; its adjoint is $(\mathbb T^*\psi)_t=u(a_{t-1})^*\psi_{t-1}$, where $\psi_{-1}:=\bar\chi_\nu\psi_{n-1}$ and $a_{-1}=a_{n-1}$.

Under the identification, $T_\nu$ acts on $\psi^+$ and on $\psi^-$ as $\mathbb T$, and $T_i$, $i\in\{j,j'\}$, acts on both as $\operatorname{diag}\left(\tau_i(A_0),\dots,\tau_i(A_{n-1})\right)$. The two expressions in \eqref{eq:gaussian-bounds-1} agree summand by summand; so the summand of the direction $\nu$ is $1-P^-_\nu T_\nu-P^+_\nu T_\nu^*$, where, in the basis of Section~\ref{sec:slices-slice-operators}, $P^+_\nu=\operatorname{diag}(1_\rho,0)$ and $P^-_\nu=\operatorname{diag}(0,1_\rho)$, and it acts as $1-\mathbb T^*$ on $\psi^+$ and as $1-\mathbb T$ on $\psi^-$. The summands of the two directions $i\ne\nu$ contribute $2-\frac12\sum_{i\ne\nu}\left(\tau_i+\tau_i^*\right)$ on the diagonal and, in the slice $t$, the off-diagonal part
\[
\frac12\sum_{i\ne\nu}\gamma_i\left(\tau_i-\tau_i^*\right)=\begin{pmatrix}0&C\\ C'&0\end{pmatrix},\qquad C':=\frac12\sum_{i\ne\nu}\sigma_i^*\left(\tau_i-\tau_i^*\right)=-C^*,
\]
with $\tau_i=\tau_i(A_t)$ and $C=C(A_t)$; the last equality holds because $\sigma_i^*$ commutes with $\tau_i$ and with $\tau_i^*$. By \eqref{eq:gaussian-bounds-7}, the diagonal parts together with $\boldsymbol{\mu}$ add up to $\mathbb B-\mathbb T^*$ and $\mathbb B-\mathbb T$. Therefore, on $\mathcal V^n\oplus\mathcal V^n$,
\[
D_A+\boldsymbol{\mu}=\mathbb D:=\begin{pmatrix}\mathbb B-\mathbb T^*&\mathbb C\\ -\mathbb C^*&\mathbb B-\mathbb T\end{pmatrix}.
\]

Let $\hat{\mathbb T}:=\mathbb T\oplus\mathbb T$ and $\mathbb S:=\operatorname{diag}\left(S(A_0),\dots,S(A_{n-1})\right)$, both acting on $\mathcal K^n$, which is identified with $\mathcal V^n\oplus\mathcal V^n$ by $\left(\left(k^+_t,k^-_t\right)\right)_t\mapsto\left(\left(k^+_t\right)_t,\left(k^-_t\right)_t\right)$, and let
\[
\mathbb N:=\begin{pmatrix}\mathbb B&0\\ -\mathbb C^*&-1\end{pmatrix}.
\]
By the first form in \eqref{eq:gaussian-bounds-8}, the blocks of $\mathbb S$ on $\mathcal V^n\oplus\mathcal V^n$ are $\mathbb B^{-1}$, $-\mathbb B^{-1}\mathbb C$, $-\mathbb C^*\mathbb B^{-1}$ and $\mathbb B+\mathbb C^*\mathbb B^{-1}\mathbb C$. Using also $\mathbb T^*\mathbb T=1$, one finds
\begin{gather*}
\mathbb D\begin{pmatrix}\mathbb T&0\\ 0&1\end{pmatrix}=\begin{pmatrix}\mathbb B\mathbb T-1&\mathbb C\\ -\mathbb C^*\mathbb T&\mathbb B-\mathbb T\end{pmatrix},\\
\mathbb N\hat{\mathbb T}=\begin{pmatrix}\mathbb B\mathbb T&0\\ -\mathbb C^*\mathbb T&-\mathbb T\end{pmatrix},\qquad
\mathbb N\mathbb S=\begin{pmatrix}1&-\mathbb C\\ 0&-\mathbb B\end{pmatrix},
\end{gather*}
so that, $\hat{\mathbb T}$ being unitary,
\[
\mathbb D(\mathbb T\oplus1)=\mathbb N\left(\hat{\mathbb T}-\mathbb S\right)=\mathbb N\hat{\mathbb T}\left(1-\hat{\mathbb T}^*\mathbb S\right).
\]
The operator $\mathbb N$ is block triangular, so $\det\mathbb N=(-1)^{nr}\cdot\det\mathbb B$; moreover $\det\hat{\mathbb T}=(\det\mathbb T)^2$ and $\det(\mathbb T\oplus1)=\det\mathbb T\ne0$. Taking determinants, which are multiplicative, therefore gives
\[
\det\mathbb D=(-1)^{nr}\cdot\det\mathbb B\cdot\det\mathbb T\cdot\det\left(1-\hat{\mathbb T}^*\mathbb S\right).
\]

Let $e_0,\dots,e_{n-1}$ be the standard basis vectors of the space of complex column vectors of length $n$, and write $\psi=\sum_te_t\otimes\psi_t$ for $\psi\in\mathcal V^n$. Then
\[
\mathbb T=\operatorname{diag}\left(u(a_0),\dots,u(a_{n-1})\right)\cdot\left(\varsigma\otimes1_{\mathcal V}\right),
\]
where $\varsigma$ is the $n\times n$ matrix with $\varsigma e_{t+1}=e_t$ for $0\le t\le n-2$ and $\varsigma e_0=\chi_\nu e_{n-1}$. It is the matrix of a cyclic permutation of $n$ elements with one entry multiplied by $\chi_\nu$, so $\det\varsigma=(-1)^{n-1}\chi_\nu$, and, as $\dim\mathcal V=r$,
\[
\det\mathbb T=\left((-1)^{n-1}\chi_\nu\right)^r\cdot\prod_{t=0}^{n-1}\det u(a_t).
\]
Finally, for $w=(w_0,\dots,w_{n-1})\in\mathcal K^n$,
\[
\left(\hat{\mathbb T}^*\mathbb Sw\right)_t=g_{t-1}w_{t-1}\quad(1\le t\le n-1),\qquad\left(\hat{\mathbb T}^*\mathbb Sw\right)_0=\bar\chi_\nu\,g_{n-1}w_{n-1}.
\]
Lemma~\ref{lem:cyclic-block-determinants}, applied to the endomorphisms $g_0,\dots,g_{n-2},\bar\chi_\nu g_{n-1}$ of $\mathcal K$, gives
\[
\det\left(1-\hat{\mathbb T}^*\mathbb S\right)=\det\left(1-\bar\chi_\nu\,g_{n-1}\cdots g_0\right).
\]
Since $\Delta^\chi(A)=\det\mathbb D$ and $\det\mathbb B=\prod_t\det B(A_t)$, and since the signs combine to $(-1)^{nr}\cdot(-1)^{(n-1)r}\cdot\chi_\nu^r=(-\chi_\nu)^r$, the identity \eqref{eq:gaussian-bounds-10} follows.
\end{proof}

\begin{corollary}[Positivity]\label{cor:positivity}
\begin{enumerate}
\item[(i)] Let $\chi$ be a character and $\nu$ a direction with $\chi_\nu=-1$, and let the notation of Section~\ref{sec:slices-slice-operators} be formed with the time direction $\nu$ and the character $\chi$. If $A\in C^1$ and $A_*\in\mathbb R^l$ satisfy $a_t=0$ and $A_t=A_*$ for all $t\in\mathbb Z/n\mathbb Z$, then
\[
\Delta^\chi(A)=\left(\det B(A_*)\right)^n\cdot\det\left(1+S(A_*)^n\right)>0.
\]
\item[(ii)] Let the Clifford triple be the four-component representation, and let $\chi$ be a character. Either of the two conditions
\[
\chi_\mu=-1\text{ for some direction }\mu,\qquad q\ne0,
\]
implies $Z^\chi_\varepsilon(0)>0$.
\end{enumerate}
\end{corollary}

\begin{proof}
(i) Under the hypotheses of (i), $u(a_t)=u(0)=1$ and $g_t=S(A_*)$ for every $t$, while $(-\chi_\nu)^r=1$ and $\bar\chi_\nu=-1$; so the formula is a special case of \eqref{eq:gaussian-bounds-10}. By \eqref{eq:gaussian-bounds-9}, $B(A_*)>0$ and $S(A_*)>0$; hence $\det B(A_*)>0$, $S(A_*)^n>0$ and $\det\left(1+S(A_*)^n\right)>0$.

(ii) Let the Clifford triple be the four-component representation, and let $\chi$ be an arbitrary character. Then $\Delta^\chi\ge0$ by Lemma~\ref{lem:gaussian-bounds-1}(v), and $\Delta^\chi$ is continuous by Lemma~\ref{lem:gaussian-bounds-1}(i). Assume first that $\chi_\mu=-1$ for a direction $\mu$. Assertion (i), applied to the character $\chi$, to $\mu$ in place of $\nu$ and to $A=0$, for which $a_t=0$ and $A_t=0$ for all $t$, gives $\Delta^\chi(0)>0$. A continuous nonnegative function which is positive at one point is positive on a neighbourhood of that point; hence the integrand of \eqref{eq:gaussian-bounds-2} at $\mathfrak{h}=0$ and $b=0$, namely
\[
\exp\left(-\frac12\|d_1A\|^2-\frac{\varepsilon}{2}\|A\|^2\right)\cdot\Delta^\chi(A),
\]
is nonnegative on $C^1$ and positive on a nonempty open set, and $Z^\chi_\varepsilon(0)>0$. Assume next that $q\ne0$, and choose $\vartheta\in\mathbb R^3$ with $\chi=\chi^{\mathrm{ap}}\exp(i\vartheta)$, that is, $\chi_\mu=-\exp(i\vartheta_\mu)$ for $\mu=1,2,3$; this is possible because $|\chi_\mu|=1$. Lemma~\ref{lem:gaussian-bounds-1}(iv), applied to the character $\chi^{\mathrm{ap}}$, gives $\Delta^\chi(A)=\Delta^{\mathrm{ap}}\left(A+c^\vartheta\right)$ for all $A\in C^1$, with $c^\vartheta$ as in Section~\ref{sec:proca-parent-harmonic}. Since every component of $\chi^{\mathrm{ap}}$ equals $-1$, assertion (i), applied to the character $\chi^{\mathrm{ap}}$, to any direction as the time direction and to $A=0$, gives $\Delta^{\mathrm{ap}}(0)>0$; thus $\Delta^\chi$ is positive at the point $-c^\vartheta$, and the argument just given yields $Z^\chi_\varepsilon(0)>0$.
\end{proof}

\subsection{Second quantisation and a trace formula}\label{sec:second-quantisation-trace}

Let $\mathcal{K}$ be a Hilbert space of finite dimension $\varkappa$, and let $(e_i)_{1\le i\le\varkappa}$ be an orthonormal basis of $\mathcal{K}$. For $1\le k\le\varkappa$ let $\wedge^k\mathcal{K}$ be the $k$-th exterior power of $\mathcal{K}$, let $\wedge^0\mathcal{K}$ be the field of scalars, and put
\[
\mathcal{F}=\bigoplus_{k=0}^{\varkappa}\wedge^k\mathcal{K}.
\]
For $1\le k\le\varkappa$ the space $\wedge^k\mathcal{K}$ carries the inner product determined by
\[
\left\langle u_1\wedge\dots\wedge u_k,\,v_1\wedge\dots\wedge v_k\right\rangle=\det\left(\langle u_a,v_b\rangle\right)_{a,b=1}^k\qquad(u_1,\dots,u_k,v_1,\dots,v_k\in\mathcal{K}),
\]
the space $\wedge^0\mathcal{K}$ carries the inner product $\langle z,z'\rangle=\bar{z}z'$, and $\mathcal{F}$ carries the inner product for which the summands $\wedge^k\mathcal{K}$, $0\le k\le\varkappa$, are mutually orthogonal and which restricts to these inner products on them. For a subset $J=\{j_1<\dots<j_k\}$ of $\{1,\dots,\varkappa\}$ put $e_J=e_{j_1}\wedge\dots\wedge e_{j_k}\in\wedge^k\mathcal{K}$, and put $e_{\emptyset}=1\in\wedge^0\mathcal{K}$. Whatever the orthonormal basis $(e_i)$ of $\mathcal{K}$, the family $(e_J)_{J\subset\{1,\dots,\varkappa\}}$ is an orthonormal basis of $\mathcal{F}$. Indeed, if $\#|I|=\#|J|$, the matrix $\left(\langle e_i,e_j\rangle\right)_{i\in I,\,j\in J}$ is the unit matrix when $I=J$ and has a zero row when $I\ne J$; and the number of the sets $J$ with $\#|J|=k$ is $\binom{\varkappa}{k}=\dim\wedge^k\mathcal{K}$.

For a linear map $g\colon\mathcal{K}\to\mathcal{K}$ let $\Gamma(g)\colon\mathcal{F}\to\mathcal{F}$ be the linear map which is the identity on $\wedge^0\mathcal{K}$ and satisfies
\[
\Gamma(g)\left(u_1\wedge\dots\wedge u_k\right)=gu_1\wedge\dots\wedge gu_k\qquad(1\le k\le\varkappa,\ u_1,\dots,u_k\in\mathcal{K});
\]
such a map exists and is unique, since the right side is multilinear and alternating in $(u_1,\dots,u_k)$ and the vectors $u_1\wedge\dots\wedge u_k$ span $\wedge^k\mathcal{K}$. For $I,J\subset\{1,\dots,\varkappa\}$ let $g_{IJ}$ be the matrix $\left(\langle e_i,ge_j\rangle\right)_{i\in I,\,j\in J}$, with rows and columns in increasing order; the determinant of the empty matrix is $1$. As $\Gamma(g)e_J=ge_{j_1}\wedge\dots\wedge ge_{j_k}$, the definition of the inner product shows that the matrix elements of $\Gamma(g)$ are the minors of $g$:
\[
\left\langle e_I,\Gamma(g)e_J\right\rangle=\det(g_{IJ})\quad\text{if }\#|I|=\#|J|,\qquad\left\langle e_I,\Gamma(g)e_J\right\rangle=0\quad\text{if }\#|I|\ne\#|J|.
\]
The degree operator $\mathsf{N}$ is the operator on $\mathcal{F}$ which equals $k$ on $\wedge^k\mathcal{K}$ for $0\le k\le\varkappa$. For a scalar $z$, $z^{\mathsf{N}}$ is the operator on $\mathcal{F}$ which equals $z^k$ on $\wedge^k\mathcal{K}$, where $z^0=1$. The identity maps of $\mathcal{K}$ and of $\mathcal{F}$ are both denoted by $1$.

\begin{lemma}[Properties of second quantisation]\label{lem:properties-second-quantisation}
Let $g,g'\colon\mathcal{K}\to\mathcal{K}$ be linear maps; $\|\cdot\|$ denotes operator norms.
\begin{enumerate}
\item[(a)] $\Gamma(gg')=\Gamma(g)\cdot\Gamma(g')$, and $\Gamma(1)=1$.
\item[(b)] $\Gamma(g)^{*}=\Gamma(g^{*})$.
\item[(c)] If $g\ge0$, then $\Gamma(g)\ge0$.
\item[(d)] $\operatorname{Tr}_{\mathcal{F}}\Gamma(g)=\det\nolimits_{\mathcal{K}}(1+g)$.
\item[(e)] $\|\Gamma(g)\|\le\max(1,\|g\|)^{\varkappa}$.
\item[(f)] If $z_1,\dots,z_{\varkappa}$ are scalars and $ge_i=z_ie_i$ for $1\le i\le\varkappa$, then $\Gamma(g)e_J=\left(\prod_{i\in J}z_i\right)e_J$ for every $J\subset\{1,\dots,\varkappa\}$. In particular, $\Gamma(z\cdot1)=z^{\mathsf{N}}$ for every scalar $z$.
\item[(g)] $\mathsf{N}$ commutes with $\Gamma(g)$.
\end{enumerate}
\end{lemma}

\begin{proof}
\emph{(a), (f), (g).} For $1\le k\le\varkappa$ and $u_1,\dots,u_k\in\mathcal{K}$,
\[
\Gamma(g)\cdot\Gamma(g')\left(u_1\wedge\dots\wedge u_k\right)=\Gamma(g)\left(g'u_1\wedge\dots\wedge g'u_k\right)=gg'u_1\wedge\dots\wedge gg'u_k=\Gamma(gg')\left(u_1\wedge\dots\wedge u_k\right).
\]
Both $\Gamma(g)\cdot\Gamma(g')$ and $\Gamma(gg')$ are the identity on $\wedge^0\mathcal{K}$, and the vectors $u_1\wedge\dots\wedge u_k$ span $\wedge^k\mathcal{K}$; this proves the first identity of (a). The identity $\Gamma(1)=1$ is immediate from the definition. Under the hypothesis of (f), multilinearity of the exterior product gives, for $J=\{j_1<\dots<j_k\}$ with $k\ge1$,
\[
\Gamma(g)e_J=z_{j_1}e_{j_1}\wedge\dots\wedge z_{j_k}e_{j_k}=\left(\prod_{i\in J}z_i\right)e_J,
\]
while for $J=\emptyset$ both sides equal $e_{\emptyset}$, the empty product being $1$. For $g=z\cdot1$ the factor is $z^{\#|J|}$, which is the value of $z^{\mathsf{N}}$ on $e_J\in\wedge^{\#|J|}\mathcal{K}$; as $(e_J)$ is a basis of $\mathcal{F}$, it follows that $\Gamma(z\cdot1)=z^{\mathsf{N}}$. Finally, $\Gamma(g)$ maps each summand $\wedge^k\mathcal{K}$ into itself, and on each summand $\mathsf{N}$ is a multiple of the identity; this is (g).

\emph{(b).} Let $I,J\subset\{1,\dots,\varkappa\}$ with $\#|I|=\#|J|$. Since $\langle e_i,g^{*}e_j\rangle=\overline{\langle e_j,ge_i\rangle}$, the matrix $(g^{*})_{IJ}$ is the conjugate transpose of $g_{JI}$, and therefore
\[
\left\langle e_I,\Gamma(g)^{*}e_J\right\rangle=\overline{\left\langle e_J,\Gamma(g)e_I\right\rangle}=\overline{\det(g_{JI})}=\det\left((g^{*})_{IJ}\right)=\left\langle e_I,\Gamma(g^{*})e_J\right\rangle.
\]
If $\#|I|\ne\#|J|$, the first and the last of these terms vanish. As $(e_J)$ is an orthonormal basis of $\mathcal{F}$, (b) follows.

\emph{(c).} Let $g\ge0$, and let $g^{\frac{1}{2}}\ge0$ be its square root, which is self-adjoint. By (a) and (b),
\[
\Gamma(g)=\Gamma\left(g^{\frac{1}{2}}\right)\cdot\Gamma\left(g^{\frac{1}{2}}\right)=\Gamma\left(g^{\frac{1}{2}}\right)^{*}\cdot\Gamma\left(g^{\frac{1}{2}}\right)\ge0.
\]

\emph{(d).} By the formula for the matrix elements of $\Gamma(g)$,
\[
\operatorname{Tr}_{\mathcal{F}}\Gamma(g)=\sum_J\left\langle e_J,\Gamma(g)e_J\right\rangle=\sum_J\det(g_{JJ}),
\]
both sums extending over all subsets $J$ of $\{1,\dots,\varkappa\}$. On the other hand, let $M=\left(\langle e_i,ge_j\rangle\right)_{i,j=1}^{\varkappa}$ be the matrix of $g$. The matrix of $1+g$ is $1+M$, whose $j$-th column is the sum of the $j$-th column of $M$ and the $j$-th column of the unit matrix. As the determinant is linear in each column, $\det\nolimits_{\mathcal{K}}(1+g)=\sum_J\det M^{(J)}$, where $M^{(J)}$ is the matrix whose $j$-th column is that of $M$ for $j\in J$ and that of the unit matrix for $j\notin J$. Expanding $\det M^{(J)}$ successively along the columns of index $j\notin J$, each of which is a column of the unit matrix, deletes the row and the column of index $j$ with the sign $+1$; hence $\det M^{(J)}=\det(g_{JJ})$. Thus $\det\nolimits_{\mathcal{K}}(1+g)$ is the sum of all principal minors $\det(g_{JJ})$ of $g$, which is $\operatorname{Tr}_{\mathcal{F}}\Gamma(g)$.

\emph{(e).} By (b) and (a), $\Gamma(g)^{*}\cdot\Gamma(g)=\Gamma(g^{*})\cdot\Gamma(g)=\Gamma(g^{*}g)$. By the spectral theorem there are an orthonormal basis $(e'_i)_{1\le i\le\varkappa}$ of $\mathcal{K}$ and numbers $s_i\ge0$ with $g^{*}ge'_i=s_i^2e'_i$; here $s_i^2\le\|g^{*}g\|=\|g\|^2$. The orthonormal basis $(e_i)$ having been arbitrary, (f) applies to $(e'_i)$: the operator $\Gamma(g^{*}g)$ is diagonal in the orthonormal basis $(e'_J)$ of $\mathcal{F}$, with the eigenvalues
\[
\prod_{i\in J}s_i^2\le\max(1,\|g\|)^{2\#|J|}\le\max(1,\|g\|)^{2\varkappa}.
\]
Therefore $\|\Gamma(g)\|^2=\|\Gamma(g)^{*}\cdot\Gamma(g)\|=\|\Gamma(g^{*}g)\|\le\max(1,\|g\|)^{2\varkappa}$.
\end{proof}

From now on $\mathcal{K}=\mathcal{V}\oplus\mathcal{V}$ is the space of Section~\ref{sec:slices-slice-operators}, formed with the time direction $\nu$ and the character $\chi$, with the inner product defined there; thus $\varkappa=2r$. An element of $\mathcal{K}$ is read as a function on $\mathbb{Z}^2$ with values in $\mathbb{C}^{\rho}\oplus\mathbb{C}^{\rho}$, and it is determined by its values on $\{0,\dots,n-1\}^2$. In what follows $\mathbf{x}$ ranges over $\{0,\dots,n-1\}^2$, the set of representatives of the $v$ sites of $\Sigma$. For $\mathbf{x}\in\{0,\dots,n-1\}^2$ and $\alpha\in\{1,\dots,2\rho\}$ let $e_{\mathbf{x},\alpha}\in\mathcal{K}$ be the element whose value at $\mathbf{x}$ is the $\alpha$-th unit vector of $\mathbb{C}^{\rho}\oplus\mathbb{C}^{\rho}$ and whose value at every other point of $\{0,\dots,n-1\}^2$ is $0$. These $2\rho v=2r=\varkappa$ vectors form an orthonormal basis of $\mathcal{K}$. The index set $\{0,\dots,n-1\}^2\times\{1,\dots,2\rho\}$ is ordered in an arbitrary fixed way, and $(e_J)$ denotes the orthonormal basis of $\mathcal{F}$ formed with the basis $(e_{\mathbf{x},\alpha})$, $J$ ranging over the subsets of this index set. Let $Q_{\mathbf{x}}$, $\langle a,Q\rangle$ for $a\in\mathbb{R}^v$, and $Q_{\mathrm{tot}}$ be the operators on $\mathcal{F}$ given by
\begin{align}
Q_{\mathbf{x}}e_J&=\left(\rho-\#|\{\alpha:(\mathbf{x},\alpha)\in J\}|\right)e_J,\notag\\
\langle a,Q\rangle&=\sum_{\mathbf{x}}a(\mathbf{x})\cdot Q_{\mathbf{x}},\qquad Q_{\mathrm{tot}}=\sum_{\mathbf{x}}Q_{\mathbf{x}}=r-\mathsf{N}.\label{eq:gaussian-bounds-11}
\end{align}
The last identity holds because $\sum_{\mathbf{x}}\#|\{\alpha:(\mathbf{x},\alpha)\in J\}|=\#|J|$ and $\rho v=r$. The operators in \eqref{eq:gaussian-bounds-11} are diagonal in the orthonormal basis $(e_J)$, with real eigenvalues; hence they are self-adjoint and commute with each other. A linear map $T\colon\mathcal{F}\to\mathcal{F}$ is said to be \emph{diagonal} if every vector $e_J$ is an eigenvector of $T$, that is, if $Te_J\in\mathbb{C}e_J$ for every $J$. Any two diagonal operators commute.

\begin{proposition}[Trace formula for the determinant]\label{prop:trace-formula-determinant}
Let $\theta\in\mathbb{R}$ be such that $\chi_{\nu}=-\exp(i\theta)$. Then, for every $A\in C^1$,
\begin{equation}\label{eq:gaussian-bounds-12}
\Delta^{\chi}(A)=\operatorname{Tr}_{\mathcal{F}}\left[\exp(i\theta Q_{\mathrm{tot}})\cdot\prod_{t=n-1}^0\exp(iq\langle a_t,Q\rangle)\cdot\det B(A_t)\cdot\Gamma(S(A_t))\right],
\end{equation}
where $A_t\in\mathbb{R}^l$ and $a_t\in\mathbb{R}^v$ are formed from $A$ as in Section~\ref{sec:slices-slice-operators}, the operators $B(A_t)$ and $S(A_t)$ are given by \eqref{eq:gaussian-bounds-7} and \eqref{eq:gaussian-bounds-8}, and the factors of the product stand from left to right in the order $t=n-1,n-2,\dots,0$. The operator $\exp(i\theta Q_{\mathrm{tot}})$ commutes with each factor of the product.
\end{proposition}

\begin{proof}
Let $a\in\mathbb{R}^v$. The operator $u(a)$ of Section~\ref{sec:slices-slice-operators} multiplies the value at $\mathbf{x}$ of an element of $\mathcal{V}$ by $\exp(iqa(\mathbf{x}))$. Hence $e_{\mathbf{x},\alpha}$ is an eigenvector of $u(a)\oplus u(a)$ with the eigenvalue $\exp(iqa(\mathbf{x}))$, and of $u(a)^{*}\oplus u(a)^{*}$ with the eigenvalue $\exp(-iqa(\mathbf{x}))$. The space $\mathcal{V}$ has an orthonormal basis of eigenvectors of $u(a)$ in which each $\mathbf{x}$ carries $\rho$ vectors, all with the eigenvalue $\exp(iqa(\mathbf{x}))$; therefore
\[
\det\nolimits_{\mathcal{V}}u(a)=\prod_{\mathbf{x}}\exp\left(iq\rho\,a(\mathbf{x})\right).
\]
By Lemma~\ref{lem:properties-second-quantisation}(f), applied to $u(a)^{*}\oplus u(a)^{*}$ and the basis $(e_{\mathbf{x},\alpha})$, and by \eqref{eq:gaussian-bounds-11}, one has for every $J$
\[
\det\nolimits_{\mathcal{V}}u(a)\cdot\Gamma\left(u(a)^{*}\oplus u(a)^{*}\right)e_J=\exp\left(iq\sum_{\mathbf{x}}a(\mathbf{x})\left(\rho-\#|\{\alpha:(\mathbf{x},\alpha)\in J\}|\right)\right)e_J=\exp(iq\langle a,Q\rangle)\,e_J,
\]
that is,
\[
\det\nolimits_{\mathcal{V}}u(a)\cdot\Gamma\left(u(a)^{*}\oplus u(a)^{*}\right)=\exp(iq\langle a,Q\rangle).
\]

The hypothesis on $\theta$ gives $(-\chi_{\nu})^r=\exp(i\theta r)$ and $\bar{\chi}_{\nu}=-\exp(-i\theta)$. Let $g_t=\left(u(a_t)^{*}\oplus u(a_t)^{*}\right)S(A_t)$, $0\le t\le n-1$, be the operators of \eqref{eq:gaussian-bounds-10}. By Lemma~\ref{lem:properties-second-quantisation}(d), then (a), then (f),
\begin{align*}
\det\nolimits_{\mathcal{K}}\left(1-\bar{\chi}_{\nu}\,g_{n-1}g_{n-2}\cdots g_0\right)&=\det\nolimits_{\mathcal{K}}\left(1+\exp(-i\theta)\cdot g_{n-1}g_{n-2}\cdots g_0\right)\\
&=\operatorname{Tr}_{\mathcal{F}}\Gamma\left(\exp(-i\theta)\cdot g_{n-1}g_{n-2}\cdots g_0\right)\\
&=\operatorname{Tr}_{\mathcal{F}}\left[\Gamma\left(\exp(-i\theta)\cdot1\right)\cdot\Gamma(g_{n-1})\cdot\Gamma(g_{n-2})\cdots\Gamma(g_0)\right]\\
&=\operatorname{Tr}_{\mathcal{F}}\left[\exp(-i\theta\mathsf{N})\cdot\Gamma(g_{n-1})\cdot\Gamma(g_{n-2})\cdots\Gamma(g_0)\right];
\end{align*}
in the last step $\Gamma(\exp(-i\theta)\cdot1)=\left(\exp(-i\theta)\right)^{\mathsf{N}}$ by Lemma~\ref{lem:properties-second-quantisation}(f), and this operator is $\exp(-i\theta\mathsf{N})$, both being equal to $\exp(-i\theta k)$ on $\wedge^k\mathcal{K}$. Moreover, $\Gamma(g_t)=\Gamma\left(u(a_t)^{*}\oplus u(a_t)^{*}\right)\cdot\Gamma(S(A_t))$ by Lemma~\ref{lem:properties-second-quantisation}(a).

These identities are now substituted into the formula \eqref{eq:gaussian-bounds-10} of Lemma~\ref{lem:slice-factorisation-determinant}. The numbers $\det B(A_t)$ and $\det\nolimits_{\mathcal{V}}u(a_t)$ may be moved into the product, next to $\Gamma(g_t)$; by the identity of the first paragraph, with $a=a_t$,
\[
\det B(A_t)\cdot\det\nolimits_{\mathcal{V}}u(a_t)\cdot\Gamma(g_t)=\exp(iq\langle a_t,Q\rangle)\cdot\det B(A_t)\cdot\Gamma(S(A_t)).
\]
It follows that
\[
\Delta^{\chi}(A)=\exp(i\theta r)\cdot\operatorname{Tr}_{\mathcal{F}}\left[\exp(-i\theta\mathsf{N})\cdot\prod_{t=n-1}^0\exp(iq\langle a_t,Q\rangle)\cdot\det B(A_t)\cdot\Gamma(S(A_t))\right],
\]
the product being ordered as in \eqref{eq:gaussian-bounds-12}. Since $\exp(i\theta r)\cdot\exp(-i\theta\mathsf{N})=\exp\left(i\theta(r-\mathsf{N})\right)=\exp(i\theta Q_{\mathrm{tot}})$ by \eqref{eq:gaussian-bounds-11}, this is \eqref{eq:gaussian-bounds-12}.

Finally, $Q_{\mathrm{tot}}=r-\mathsf{N}$ commutes with every $\Gamma(S(A_t))$ by Lemma~\ref{lem:properties-second-quantisation}(g); being diagonal, it commutes with every diagonal operator, in particular with every $\exp(iq\langle a_t,Q\rangle)$; and the factors $\det B(A_t)$ are numbers. Hence $\exp(i\theta Q_{\mathrm{tot}})$ commutes with each factor of the product.
\end{proof}

\subsection{Layer kernels and the chain formula}\label{sec:layer-kernels-chain}

\subsubsection{The kernel of a layer}\label{sec:kernel-layer}

Let the time direction $\nu$ be as in Section~\ref{sec:slices-slice-operators}. On $\mathbb{R}^l$ and on $\mathbb{R}^v$, $|\cdot|$ and $\langle\cdot,\cdot\rangle$ denote the Euclidean norm and inner product; $\|\cdot\|$ denotes the operator norm of operators on $\mathcal{F}$. \emph{Layer data} are pairs $\lambda=(\mathfrak{h}^e,b^0)\in\mathbb{R}^l\times\mathbb{R}^v$, and for such $\lambda$ one puts $-\lambda=(-\mathfrak{h}^e,-b^0)$; the layer data $\lambda_t=(\mathfrak{h}^e_t,b^0_t)$ of Section~\ref{sec:slices-slice-operators} are of this kind. For layer data $\lambda=(\mathfrak{h}^e,b^0)$ and for $y,p\in\mathbb{R}^l$ define operators $k_{\lambda}(y)$, $\varpi(p)$ and $\mathsf{u}_{\lambda}(p)$ on $\mathcal{F}$ by
\begin{equation}\label{eq:gaussian-bounds-13}
k_{\lambda}(y)=\int_{\mathbb{R}^v}\mathrm{d}a\,\exp\left(-\frac{1}{2}|y-d_sa+\mathfrak{h}^e|^2-\frac{\varepsilon}{2}|a+b^0|^2\right)\cdot\exp(iq\langle a,Q\rangle)
\end{equation}
and
\begin{gather}
\varpi(p)=(2\pi)^{-\frac{l}{2}}\cdot\left(\frac{2\pi}{\varepsilon}\right)^{\frac{v}{2}}\cdot\exp\left(-\frac{1}{2}|p|^2-\frac{1}{2\varepsilon}\sum_{\mathbf{x}}\left(qQ_{\mathbf{x}}-(d_s^{*}p)(\mathbf{x})\right)^2\right),\notag\\
\mathsf{u}_{\lambda}(p)=\exp\left(i\langle p,\mathfrak{h}^e+d_sb^0\rangle\right)\cdot\exp\left(-iq\langle b^0,Q\rangle\right),\label{eq:gaussian-bounds-14}
\end{gather}
where $Q_{\mathbf{x}}$, $\langle a,Q\rangle$ and the range of $\mathbf{x}$ are those of \eqref{eq:gaussian-bounds-11}, and $d_s$, $d_s^{*}$ are the spatial differentials of Section~\ref{sec:slices-slice-operators}. For every $a\in\mathbb{R}^v$ the operator $\langle a,Q\rangle$ is diagonal and self-adjoint, so that $\exp(iq\langle a,Q\rangle)$ is diagonal and unitary. The integrand in \eqref{eq:gaussian-bounds-13} is therefore a diagonal operator of norm at most $\exp\left(-\frac{\varepsilon}{2}|a+b^0|^2\right)$, which is an integrable function of $a\in\mathbb{R}^v$; the integral converges absolutely, and $k_{\lambda}(y)$ is diagonal.

\begin{lemma}[Fourier form of the layer kernel]\label{lem:fourier-form-layer}
Let $\lambda=(\mathfrak{h}^e,b^0)\in\mathbb{R}^l\times\mathbb{R}^v$ be layer data.
\begin{enumerate}
\item[(a)] For every $p\in\mathbb{R}^l$ the operator $\varpi(p)$ is diagonal and positive definite, and
\[
\|\varpi(p)\|\le c_{\varpi}\cdot\exp\left(-\frac{1}{2}|p|^2\right),\qquad c_{\varpi}=(2\pi)^{-\frac{l}{2}}\cdot\left(\frac{2\pi}{\varepsilon}\right)^{\frac{v}{2}}.
\]
The operator $\mathsf{u}_{\lambda}(p)$ is diagonal and unitary, and $\mathsf{u}_0(p)=1$ for the zero layer data $0=(0,0)$.
\item[(b)] For every $y\in\mathbb{R}^l$,
\begin{equation}\label{eq:gaussian-bounds-15}
k_{\lambda}(y)=\int_{\mathbb{R}^l}\mathrm{d}p\,\exp(i\langle p,y\rangle)\cdot\mathsf{u}_{\lambda}(p)\cdot\varpi(p).
\end{equation}
\item[(c)] $\|k_{\lambda}(y)\|\le\left(\frac{2\pi}{\varepsilon}\right)^{\frac{v}{2}}$ for every $y\in\mathbb{R}^l$.
\item[(d)] $k_{\lambda}(y)^{*}=k_{-\lambda}(-y)$ for every $y\in\mathbb{R}^l$.
\item[(e)] The map $y\mapsto k_{\lambda}(y)$ from $\mathbb{R}^l$ to the operators on $\mathcal{F}$ is continuous.
\end{enumerate}
\end{lemma}

\begin{proof}
\emph{(a).} The operator in the exponent of $\varpi(p)$ in \eqref{eq:gaussian-bounds-14} is $-\frac{1}{2}|p|^2$ minus $\frac{1}{2\varepsilon}$ times a sum of squares of diagonal self-adjoint operators. It is therefore diagonal, self-adjoint and bounded above by $-\frac{1}{2}|p|^2$, and its exponential is a diagonal operator whose eigenvalues lie in the interval $\left(0,\exp\left(-\frac{1}{2}|p|^2\right)\right]$. Multiplication by the positive number $c_{\varpi}$ gives the assertions on $\varpi(p)$. In $\mathsf{u}_{\lambda}(p)$ the first factor is a complex number of modulus one, and the second factor is diagonal and unitary, because $\langle b^0,Q\rangle$ is diagonal and self-adjoint. For $\lambda=(0,0)$ both factors are equal to $1$.

\emph{(b).} The substitution $a=a'-b^0$ in \eqref{eq:gaussian-bounds-13} replaces $y-d_sa+\mathfrak{h}^e$ by $(y+\mathfrak{h}^e+d_sb^0)-d_sa'$, $a+b^0$ by $a'$, and $\exp(iq\langle a,Q\rangle)$ by $\exp(-iq\langle b^0,Q\rangle)\cdot\exp(iq\langle a',Q\rangle)$. Hence
\[
k_{\lambda}(y)=\exp(-iq\langle b^0,Q\rangle)\cdot k_{(0,0)}\left(y+\mathfrak{h}^e+d_sb^0\right).
\]
Let $e_J$ be one of the basis vectors of $\mathcal{F}$ of Section~\ref{sec:second-quantisation-trace}. By \eqref{eq:gaussian-bounds-11}, each $Q_{\mathbf{x}}$ acts on $e_J$ as multiplication by an integer $\mathsf{q}_{\mathbf{x}}$; let $\mathsf{q}\in\mathbb{R}^v$ be the vector with the components $\mathsf{q}_{\mathbf{x}}$. Then $\exp(iq\langle a',Q\rangle)e_J=\exp(iq\langle a',\mathsf{q}\rangle)e_J$ for $a'\in\mathbb{R}^v$. For $z\in\mathbb{R}^l$ the Gaussian integral
\[
\exp\left(-\frac{1}{2}|z|^2\right)=(2\pi)^{-\frac{l}{2}}\int_{\mathbb{R}^l}\mathrm{d}p\,\exp\left(-\frac{1}{2}|p|^2+i\langle p,z\rangle\right)
\]
holds; insert it, with $z=y-d_sa'$, into \eqref{eq:gaussian-bounds-13} for the layer data $(0,0)$. Applied to $e_J$, this yields a double integral over $(p,a')\in\mathbb{R}^l\times\mathbb{R}^v$ which converges absolutely, the norm of its integrand being $(2\pi)^{-\frac{l}{2}}\cdot\exp\left(-\frac{1}{2}|p|^2-\frac{\varepsilon}{2}|a'|^2\right)$. Fubini's theorem and the identity $\langle p,d_sa'\rangle=\langle d_s^{*}p,a'\rangle$ give
\begin{align*}
k_{(0,0)}(y)e_J&=(2\pi)^{-\frac{l}{2}}\int_{\mathbb{R}^l}\mathrm{d}p\,\exp\left(-\frac{1}{2}|p|^2+i\langle p,y\rangle\right)\\
&\qquad\cdot\int_{\mathbb{R}^v}\mathrm{d}a'\,\exp\left(-\frac{\varepsilon}{2}|a'|^2+i\langle a',q\mathsf{q}-d_s^{*}p\rangle\right)e_J.
\end{align*}
The inner integral is the value at $w=q\mathsf{q}-d_s^{*}p$ of the Gaussian integral
\[
\int_{\mathbb{R}^v}\mathrm{d}a'\,\exp\left(-\frac{\varepsilon}{2}|a'|^2+i\langle a',w\rangle\right)=\left(\frac{2\pi}{\varepsilon}\right)^{\frac{v}{2}}\cdot\exp\left(-\frac{1}{2\varepsilon}|w|^2\right)\qquad\left(w\in\mathbb{R}^v\right).
\]
The operator $\varpi(p)$ acts on $e_J$ as multiplication by the number obtained from \eqref{eq:gaussian-bounds-14} when each $Q_{\mathbf{x}}$ is replaced by $\mathsf{q}_{\mathbf{x}}$, and $|q\mathsf{q}-d_s^{*}p|^2=\sum_{\mathbf{x}}\left(q\mathsf{q}_{\mathbf{x}}-(d_s^{*}p)(\mathbf{x})\right)^2$. It follows that
\[
k_{(0,0)}(y)e_J=\int_{\mathbb{R}^l}\mathrm{d}p\,\exp(i\langle p,y\rangle)\cdot\varpi(p)e_J.
\]
The integral $\int_{\mathbb{R}^l}\mathrm{d}p\,\exp(i\langle p,y\rangle)\cdot\varpi(p)$ converges absolutely by (a), and the vectors $e_J$ form a basis of $\mathcal{F}$; hence $k_{(0,0)}(y)$ equals this integral, for every $y\in\mathbb{R}^l$. Inserting it into the expression for $k_{\lambda}(y)$ found above and recalling the definition of $\mathsf{u}_{\lambda}(p)$ in \eqref{eq:gaussian-bounds-14}, one obtains
\begin{align*}
k_{\lambda}(y)&=\int_{\mathbb{R}^l}\mathrm{d}p\,\exp\left(i\langle p,y+\mathfrak{h}^e+d_sb^0\rangle\right)\cdot\exp(-iq\langle b^0,Q\rangle)\cdot\varpi(p)\\
&=\int_{\mathbb{R}^l}\mathrm{d}p\,\exp(i\langle p,y\rangle)\cdot\mathsf{u}_{\lambda}(p)\cdot\varpi(p),
\end{align*}
which is \eqref{eq:gaussian-bounds-15}.

\emph{(c).} By the bound on the norm of the integrand of \eqref{eq:gaussian-bounds-13} noted after \eqref{eq:gaussian-bounds-14},
\[
\|k_{\lambda}(y)\|\le\int_{\mathbb{R}^v}\mathrm{d}a\,\exp\left(-\frac{\varepsilon}{2}|a+b^0|^2\right)=\left(\frac{2\pi}{\varepsilon}\right)^{\frac{v}{2}}.
\]

\emph{(d).} The adjoint of $\exp(iq\langle a,Q\rangle)$ is $\exp(-iq\langle a,Q\rangle)$, and the scalar factor of the integrand in \eqref{eq:gaussian-bounds-13} is real. Taking adjoints in \eqref{eq:gaussian-bounds-13} and substituting $-a$ for $a$, one obtains
\begin{align*}
k_{\lambda}(y)^{*}&=\int_{\mathbb{R}^v}\mathrm{d}a\,\exp\left(-\frac{1}{2}|y+d_sa+\mathfrak{h}^e|^2-\frac{\varepsilon}{2}|-a+b^0|^2\right)\cdot\exp(iq\langle a,Q\rangle)\\
&=\int_{\mathbb{R}^v}\mathrm{d}a\,\exp\left(-\frac{1}{2}|-y-d_sa-\mathfrak{h}^e|^2-\frac{\varepsilon}{2}|a-b^0|^2\right)\cdot\exp(iq\langle a,Q\rangle),
\end{align*}
and by \eqref{eq:gaussian-bounds-13} the last integral is $k_{-\lambda}(-y)$.

\emph{(e).} For each $a\in\mathbb{R}^v$ the integrand of \eqref{eq:gaussian-bounds-13} depends continuously on $y$, and its norm is at most $\exp\left(-\frac{\varepsilon}{2}|a+b^0|^2\right)$, an integrable function of $a$ that does not depend on $y$. The dominated convergence theorem shows that $y\mapsto k_{\lambda}(y)$ is continuous.
\end{proof}

\subsubsection{The weight of a slice}\label{sec:weight-slice}

\emph{Slice data} are pairs $\omega=(\mathfrak{h}^m,b^s)\in\mathbb{R}^v\times\mathbb{R}^l$; the slice data $\omega_t=(\mathfrak{h}^m_t,b^s_t)$ of Section~\ref{sec:slices-slice-operators} are of this kind. For slice data $\omega=(\mathfrak{h}^m,b^s)$ and $A\in\mathbb{R}^l$ let $\mathcal{W}_{\omega}(A)$ be the operator on $\mathcal{F}$ given by
\begin{equation}\label{eq:gaussian-bounds-16}
\mathcal{W}_{\omega}(A)=\exp\left(-\frac{1}{2}|d_sA+\mathfrak{h}^m|^2-\frac{\varepsilon}{2}|A+b^s|^2\right)\cdot\det B(A)\cdot\Gamma(S(A)),
\end{equation}
where $B(A)$ and $S(A)$ are the operators \eqref{eq:gaussian-bounds-7} and \eqref{eq:gaussian-bounds-8} on $\mathcal{V}$ and on $\mathcal{K}$, and $\Gamma$ is the map of Section~\ref{sec:second-quantisation-trace}.

By \eqref{eq:gaussian-bounds-9}, $B(A)$ is a self-adjoint operator on the $r$-dimensional space $\mathcal{V}$ with spectrum in $[1+\boldsymbol{\mu},5+\boldsymbol{\mu}]$, and $1+\boldsymbol{\mu}>0$; hence $0<\det B(A)\le(5+\boldsymbol{\mu})^r$. Also by \eqref{eq:gaussian-bounds-9}, $S(A)>0$ and $\|S(A)\|\le s_{*}$, where $s_{*}\ge9\cdot(5+\boldsymbol{\mu})>1$. Lemma~\ref{lem:properties-second-quantisation}(c), (e), applied on $\mathcal{K}$, whose dimension is $\varkappa=2r$, yields $\Gamma(S(A))\ge0$ and $\|\Gamma(S(A))\|\le s_{*}^{2r}$. Consequently $\mathcal{W}_{\omega}(A)\ge0$, and
\begin{equation}\label{eq:gaussian-bounds-17}
\|\mathcal{W}_{\omega}(A)\|\le c_W\cdot\exp\left(-\frac{\varepsilon}{2}|A+b^s|^2\right),\qquad c_W=(5+\boldsymbol{\mu})^r\cdot s_{*}^{2r}.
\end{equation}
Moreover $\mathcal{W}_{\omega}(A)$ depends continuously on $A$ (the matrix entries of the operators $\tau_i(A)$ are continuous in $A$, $B(A)^{-1}$ is continuous because $B(A)\ge1+\boldsymbol{\mu}>0$, and the matrix elements of $\Gamma(S(A))$ are minors of $S(A)$). The spatial characters $\chi_j,\chi_{j'}$ enter $\mathcal{W}_{\omega}(A)$ through the space $\mathcal{V}$; the character $\chi_{\nu}$ does not enter it.

\subsubsection{The chain formula}\label{sec:chain-formula}

\begin{proposition}[Chain formula for the partition function]\label{prop:chain-formula-partition}
Let $\chi$ be a character and $\theta\in\mathbb{R}$ with $\chi_{\nu}=-\exp(i\theta)$, let $(\mathfrak{h},b)\in C^2\times C^1$, and let $(\omega_t,\lambda_t)_{t\in\mathbb{Z}/n}$ be the slice data and the layer data of $(\mathfrak{h},b)$ for the time direction $\nu$ (Section~\ref{sec:slices-slice-operators}). For $t_0\in\mathbb{Z}/n$ and operators $X_t$ on $\mathcal{F}$, $t\in\mathbb{Z}/n$, put
\[
\prod_t^{\leftarrow}X_t=X_{t_0+n-1}\cdot X_{t_0+n-2}\cdots X_{t_0+1}\cdot X_{t_0},
\]
the indices being read modulo $n$. Then
\begin{equation}\label{eq:gaussian-bounds-18}
Z^{\chi}_{\varepsilon}(\mathfrak{h},b)=\int_{(\mathbb{R}^l)^n}\prod_t\mathrm{d}A_t\,\operatorname{Tr}_{\mathcal{F}}\left[\exp(i\theta Q_{\mathrm{tot}})\cdot\prod_t^{\leftarrow}\left(k_{\lambda_t}(A_{t+1}-A_t)\cdot\mathcal{W}_{\omega_t}(A_t)\right)\right].
\end{equation}
The absolute value of the integrand in \eqref{eq:gaussian-bounds-18} does not exceed
\[
\left(\frac{2\pi}{\varepsilon}\right)^{\frac{nv}{2}}\cdot c_W^n\cdot\dim\mathcal{F}\cdot\prod_t\exp\left(-\frac{\varepsilon}{2}|A_t+b^s_t|^2\right),
\]
and the value of the right side of \eqref{eq:gaussian-bounds-18} is independent of $t_0$.
\end{proposition}

\begin{proof}
For $A\in C^1$ let $A_t\in\mathbb{R}^l$ and $a_t\in\mathbb{R}^v$, $t\in\mathbb{Z}/n$, be formed from $A$ as in Section~\ref{sec:slices-slice-operators}. The map $A\mapsto\left((A_t)_t,(a_t)_t\right)$ merely permutes coordinates; hence Lebesgue measure on $C^1$ is the product of the Lebesgue measures $\mathrm{d}A_t$ on $\mathbb{R}^l$ and $\mathrm{d}a_t$ on $\mathbb{R}^v$, $t\in\mathbb{Z}/n$. By \eqref{eq:gaussian-bounds-6}, the exponential factor of the integrand of \eqref{eq:gaussian-bounds-2} is the product over $t\in\mathbb{Z}/n$ of the numbers
\begin{align*}
&\exp\left(-\frac{1}{2}|A_{t+1}-A_t-d_sa_t+\mathfrak{h}^e_t|^2-\frac{\varepsilon}{2}|a_t+b^0_t|^2\right)\\
\text{and}\quad&\exp\left(-\frac{1}{2}|d_sA_t+\mathfrak{h}^m_t|^2-\frac{\varepsilon}{2}|A_t+b^s_t|^2\right),
\end{align*}
and $\Delta^{\chi}(A)$ is given by \eqref{eq:gaussian-bounds-12}, by Proposition~\ref{prop:trace-formula-determinant}. Attach, for each $t$, the first of these numbers to the factor $\exp(iq\langle a_t,Q\rangle)$ of \eqref{eq:gaussian-bounds-12} and the second to the factor $\det B(A_t)\cdot\Gamma(S(A_t))$. By \eqref{eq:gaussian-bounds-16} the integrand of \eqref{eq:gaussian-bounds-2} becomes
\begin{multline*}
\operatorname{Tr}_{\mathcal{F}}\left[\exp(i\theta Q_{\mathrm{tot}})\cdot\prod_{t=n-1}^0\left(\exp\left(-\frac{1}{2}|A_{t+1}-A_t-d_sa_t+\mathfrak{h}^e_t|^2-\frac{\varepsilon}{2}|a_t+b^0_t|^2\right)\right.\right.\\
\left.\vphantom{\prod_{t=n-1}^0}\left.\vphantom{\left(\frac{1}{2}\right)}\cdot\exp(iq\langle a_t,Q\rangle)\cdot\mathcal{W}_{\omega_t}(A_t)\right)\right],
\end{multline*}
the factors standing from left to right in the order $t=n-1,\dots,1,0$, and $A_n=A_0$. The operators $\exp(i\theta Q_{\mathrm{tot}})$ and $\exp(iq\langle a_t,Q\rangle)$ are unitary. By \eqref{eq:gaussian-bounds-17}, the product under the trace therefore has norm at most
\[
c_W^n\cdot\prod_t\exp\left(-\frac{\varepsilon}{2}\left(|a_t+b^0_t|^2+|A_t+b^s_t|^2\right)\right),
\]
and the integrand of \eqref{eq:gaussian-bounds-2} is bounded by $\dim\mathcal{F}$ times this function, which is integrable over $C^1$. By Fubini's theorem the integrations over $a_0,\dots,a_{n-1}$ may be carried out first. The trace of the product is linear in each factor, and $a_t$ occurs only in the $t$-th factor, in the operator
\[
\exp\left(-\frac{1}{2}|A_{t+1}-A_t-d_sa_t+\mathfrak{h}^e_t|^2-\frac{\varepsilon}{2}|a_t+b^0_t|^2\right)\cdot\exp(iq\langle a_t,Q\rangle),
\]
whose integral over $a_t\in\mathbb{R}^v$ is $k_{\lambda_t}(A_{t+1}-A_t)$ by \eqref{eq:gaussian-bounds-13}, with $\lambda_t=(\mathfrak{h}^e_t,b^0_t)$. This proves \eqref{eq:gaussian-bounds-18} for $t_0=0$.

Put $X_t=k_{\lambda_t}(A_{t+1}-A_t)\cdot\mathcal{W}_{\omega_t}(A_t)$. The operator $\exp(i\theta Q_{\mathrm{tot}})$ commutes with every $X_t$: being diagonal, it commutes with the diagonal operator $k_{\lambda_t}(A_{t+1}-A_t)$, and $Q_{\mathrm{tot}}=r-\mathsf{N}$ commutes with $\Gamma(S(A_t))$ by Lemma~\ref{lem:properties-second-quantisation}(g), hence with $\mathcal{W}_{\omega_t}(A_t)$. Let $t_0$ be represented by an integer in $\{0,\dots,n-1\}$, and put $Y=X_{t_0-1}\cdots X_0$ and $Y'=X_{n-1}\cdots X_{t_0}$ (with $Y=1$ if $t_0=0$). The ordered product for $t_0$ is $Y\cdot Y'$, and the one for $0$ is $Y'\cdot Y$. Cyclicity of the trace and the commutation just shown give
\[
\operatorname{Tr}_{\mathcal{F}}\left[\exp(i\theta Q_{\mathrm{tot}})\cdot Y\cdot Y'\right]=\operatorname{Tr}_{\mathcal{F}}\left[Y'\cdot\exp(i\theta Q_{\mathrm{tot}})\cdot Y\right]=\operatorname{Tr}_{\mathcal{F}}\left[\exp(i\theta Q_{\mathrm{tot}})\cdot Y'\cdot Y\right].
\]
So the integrand of \eqref{eq:gaussian-bounds-18}, and with it the right side of \eqref{eq:gaussian-bounds-18}, does not depend on $t_0$.

Finally, $\|\exp(i\theta Q_{\mathrm{tot}})\|=1$, $\|k_{\lambda_t}(A_{t+1}-A_t)\|\le\left(\frac{2\pi}{\varepsilon}\right)^{\frac{v}{2}}$ by Lemma~\ref{lem:fourier-form-layer}(c), and $\|\mathcal{W}_{\omega_t}(A_t)\|\le c_W\cdot\exp\left(-\frac{\varepsilon}{2}|A_t+b^s_t|^2\right)$ by \eqref{eq:gaussian-bounds-17}. Since $\left|\operatorname{Tr}_{\mathcal{F}}X\right|\le\dim\mathcal{F}\cdot\|X\|$ for every operator $X$ on $\mathcal{F}$, the integrand of \eqref{eq:gaussian-bounds-18} obeys the stated bound.
\end{proof}

A \emph{family of data} is a family $\mathsf{d}=(\omega_t,\lambda_t)_{t\in\mathbb{Z}/n}$ of slice data $\omega_t\in\mathbb{R}^v\times\mathbb{R}^l$ and layer data $\lambda_t\in\mathbb{R}^l\times\mathbb{R}^v$. With the spatial characters $\chi_j,\chi_{j'}$ fixed, for $\theta\in\mathbb{R}$ and a family of data $\mathsf{d}$ let $\mathcal{Z}^{\theta}[\mathsf{d}]=\mathcal{Z}^{\theta}[(\omega_t,\lambda_t)_t]$ denote the right side of \eqref{eq:gaussian-bounds-18}. Every family of data is the family of data of exactly one element of $C^2\times C^1$ (Section~\ref{sec:slices-slice-operators}), and for every $\theta\in\mathbb{R}$ the character whose components are $\chi_j,\chi_{j'}$ in the spatial directions and $-\exp(i\theta)$ in the direction $\nu$ satisfies the hypothesis of Proposition~\ref{prop:chain-formula-partition}. Hence that proposition applies, and $\mathcal{Z}^{\theta}[\mathsf{d}]$ does not depend on $t_0$. Consequently, for every $t'\in\mathbb{Z}/n$,
\begin{equation}\label{eq:gaussian-bounds-19}
\mathcal{Z}^{\theta}[(\omega_{t+t'},\lambda_{t+t'})_t]=\mathcal{Z}^{\theta}[(\omega_t,\lambda_t)_t].
\end{equation}
Indeed, once the integration variables $A_t$ are renamed $A_{t+t'}$, the right side of \eqref{eq:gaussian-bounds-18} for the family $(\omega_{t+t'},\lambda_{t+t'})_t$ and the index $t_0$ becomes the right side of \eqref{eq:gaussian-bounds-18} for the family $(\omega_t,\lambda_t)_t$ and the index $t_0+t'$.

\subsection{Inequalities at mirrors}\label{sec:inequalities-mirrors}

\subsubsection{Reflected data}\label{sec:reflected-data}

Let the time direction $\nu$ be fixed. For a layer $\jmath\in\mathbb{Z}/n$, the \emph{mirror} at $\jmath$ is the reflection that maps the slice $t$ to the slice $2\jmath+1-t$ and the layer $t$ to the layer $2\jmath-t$, for $t\in\mathbb{Z}/n$; as the layer $t$ joins the slices $t$ and $t+1$, its image joins the images $2\jmath+1-t$ and $2\jmath-t$ of these slices. Let $m$ be the integer with $n=2m+1$ when $n$ is odd, and with $n=2m$ when $n$ is even. A slice $t$ is fixed by the mirror if and only if $2t=2\jmath+1$ in $\mathbb{Z}/n$, and a layer $t$ if and only if $2t=2\jmath$ in $\mathbb{Z}/n$. Thus the layer $\jmath$, a plane of links, is fixed for every $n$, and the other fixed planes are a single slice, a plane of sites, when $n$ is odd, and a single layer different from $\jmath$ when $n$ is even.

By \eqref{eq:gaussian-bounds-19}, $\mathcal{Z}^{\theta}$ is invariant under cyclic relabellings of the slices and layers; accordingly, the definitions below are stated for the mirror at the layer $m$. When $\jmath\ne m$, they are applied not to the family of data $(\omega_t,\lambda_t)_t$ but to the shifted family $(\omega_{t+\jmath-m},\lambda_{t+\jmath-m})_t$, in which the layer $\jmath$ bears the label $m$. The labels thus attached to the slices and layers, with representatives in the ranges fixed below, form the \emph{normal labelling} for $\jmath$.

\paragraph{Odd $n$.} Let $n=2m+1$ with $m\ge1$. The representatives $-m,\dots,m$ of the classes of $\mathbb{Z}/n$ serve as labels both for the slices and for the layers. For a label $t\le m-1$ the layer $t$ lies between the slices $t$ and $t+1$, while the layer $m$ lies between the slices $m$ and $-m$, because $m+1=-m$ in $\mathbb{Z}/n$. The mirror maps the slice $t$ to the slice $-t$ and the layer $t$ to the layer $-t-1$; the layer $m$ is mapped to itself. The fixed slice is $0$, and the fixed layer is $m$. For a family of data $\mathsf{d}=(\omega_t,\lambda_t)_t$ define the families of data $\mathsf{d}^{\pm}=(\omega^{\pm}_t,\lambda^{\pm}_t)_t$, $t\in\{-m,\dots,m\}$, by
\[
\omega^{+}_t=\omega_{|t|},\qquad\lambda^{+}_t=\begin{cases}\lambda_t,&0\le t\le m-1,\\-\lambda_{-t-1},&-m\le t\le-1,\\0,&t=m,\end{cases}
\]
\[
\omega^{-}_t=\omega_{-|t|},\qquad\lambda^{-}_t=\begin{cases}\lambda_t,&-m\le t\le-1,\\-\lambda_{-t-1},&0\le t\le m-1,\\0,&t=m.\end{cases}
\]

\paragraph{Even $n$.} Let $n=2m$ with $m\ge1$. Here the slices and the layers alike carry as labels the representatives $-m+1,\dots,m$ of the classes of $\mathbb{Z}/n$. Again the layer $t$ lies between the slices $t$ and $t+1$ when $t\le m-1$, and the layer $m$ lies between the slices $m$ and $-m+1$, because $m+1=-m+1$ in $\mathbb{Z}/n$. The mirror maps the slice $t$ to the slice $1-t$ and the layer $t$ to the layer $-t$; the layer $m$ is mapped to itself. The fixed layers are $0$ and $m$. For a family of data $\mathsf{d}=(\omega_t,\lambda_t)_t$ define $\mathsf{d}^{\pm}=(\omega^{\pm}_t,\lambda^{\pm}_t)_t$, $t\in\{-m+1,\dots,m\}$, by
\[
\omega^{+}_t=\begin{cases}\omega_t,&1\le t\le m,\\\omega_{1-t},&-m+1\le t\le0,\end{cases}\qquad\lambda^{+}_t=\begin{cases}\lambda_t,&1\le t\le m-1,\\-\lambda_{-t},&-m+1\le t\le-1,\\0,&t\in\{0,m\},\end{cases}
\]
\[
\omega^{-}_t=\begin{cases}\omega_t,&-m+1\le t\le0,\\\omega_{1-t},&1\le t\le m,\end{cases}\qquad\lambda^{-}_t=\begin{cases}\lambda_t,&-m+1\le t\le-1,\\-\lambda_{-t},&1\le t\le m-1,\\0,&t\in\{0,m\}.\end{cases}
\]

\paragraph{Reflected cochains and counts.} The map $(\mathfrak{h},b)\mapsto(\omega_t,\lambda_t)_t$ of Section~\ref{sec:slices-slice-operators} is a bijection of $C^2\times C^1$ onto the set of families of data. Hence for $(\mathfrak{h},b)\in C^2\times C^1$ there are unique elements $(\mathfrak{h},b)^{+}$ and $(\mathfrak{h},b)^{-}$ of $C^2\times C^1$ whose families of data are $\mathsf{d}^{+}$ and $\mathsf{d}^{-}$, where $\mathsf{d}$ is the family of data of $(\mathfrak{h},b)$ in the normal labelling for $\jmath$. When the direction and the layer are to be indicated, one writes $(\mathfrak{h},b)^{\pm}=(\mathfrak{h},b)^{\pm;\nu,\jmath}$.

For a family $\mathsf{c}$ of real components, namely a slice datum $\omega_t$, a layer datum $\lambda_t$ or a family of data $\mathsf{d}$, let $\mathsf{n}(\mathsf{c})$ be the number of its nonzero components. For $(\mathfrak{h},b)\in C^2\times C^1$ put
\[
\mathsf{n}(\mathfrak{h},b)=\#\left|\{p:\mathfrak{h}_p\ne0\}\right|+\#\left|\{\ell:b_{\ell}\ne0\}\right|,
\]
where $p$ ranges over the plaquettes and $\ell$ over the links of $\Lambda$, and let $\mathsf{n}_{\mathrm{fix}}(\mathfrak{h},b)$ be the sum of $\mathsf{n}(\lambda_t)$ over those layers $t$ that are fixed by the mirror, $\lambda_t$ denoting the layer data of $(\mathfrak{h},b)$; in the normal labelling, $\mathsf{n}_{\mathrm{fix}}(\mathfrak{h},b)=\mathsf{n}(\lambda_m)$ for odd $n$, and $\mathsf{n}_{\mathrm{fix}}(\mathfrak{h},b)=\mathsf{n}(\lambda_0)+\mathsf{n}(\lambda_m)$ for even $n$. The bijection of Section~\ref{sec:slices-slice-operators} permutes coordinates and changes some signs, and a relabelling permutes the members of a family of data. Therefore, if $\mathsf{d}=(\omega_t,\lambda_t)_t$ is the family of data of $(\mathfrak{h},b)$ in the normal labelling for $\jmath$,
\[
\mathsf{n}(\mathfrak{h},b)=\mathsf{n}(\mathsf{d})=\sum_t\left[\mathsf{n}(\omega_t)+\mathsf{n}(\lambda_t)\right],
\]
and $\mathsf{n}\left((\mathfrak{h},b)^{\pm}\right)=\mathsf{n}(\mathsf{d}^{\pm})$. A count of the components in the definitions of $\mathsf{d}^{\pm}$ gives
\begin{equation}\label{eq:gaussian-bounds-20}
\mathsf{n}\left((\mathfrak{h},b)^{+}\right)+\mathsf{n}\left((\mathfrak{h},b)^{-}\right)=2\,\mathsf{n}(\mathfrak{h},b)-2\,\mathsf{n}_{\mathrm{fix}}(\mathfrak{h},b).
\end{equation}
Indeed, since $\mathsf{n}(-\lambda)=\mathsf{n}(\lambda)$, the definitions give for odd $n$
\begin{align*}
\mathsf{n}(\mathsf{d}^{+})&=\mathsf{n}(\omega_0)+2\sum_{t=1}^m\mathsf{n}(\omega_t)+2\sum_{t=0}^{m-1}\mathsf{n}(\lambda_t),\\
\mathsf{n}(\mathsf{d}^{-})&=\mathsf{n}(\omega_0)+2\sum_{t=1}^m\mathsf{n}(\omega_{-t})+2\sum_{t=-m}^{-1}\mathsf{n}(\lambda_t),
\end{align*}
and for even $n$
\begin{align*}
\mathsf{n}(\mathsf{d}^{+})&=2\sum_{t=1}^m\mathsf{n}(\omega_t)+2\sum_{t=1}^{m-1}\mathsf{n}(\lambda_t),\\
\mathsf{n}(\mathsf{d}^{-})&=2\sum_{t=-m+1}^0\mathsf{n}(\omega_t)+2\sum_{t=-m+1}^{-1}\mathsf{n}(\lambda_t).
\end{align*}
In both cases $\mathsf{n}(\mathsf{d}^{+})+\mathsf{n}(\mathsf{d}^{-})$ is twice the sum of $\mathsf{n}(\omega_t)$ over all slices plus twice the sum of $\mathsf{n}(\lambda_t)$ over the layers that are not fixed, that is, $2\,\mathsf{n}(\mathsf{d})-2\,\mathsf{n}_{\mathrm{fix}}(\mathfrak{h},b)$.

Finally, $(0,0)^{\pm}=(0,0)$, since all data of $(0,0)$ vanish. If $b=0$, the components $b^s_t$ and $b^0_t$ of all data of $(\mathfrak{h},0)$ vanish, and so do the corresponding components of the data in $\mathsf{d}^{\pm}$; thus $(\mathfrak{h},0)^{\pm}=(\mathfrak{h}^{\pm},0)$ with two-cochains $\mathfrak{h}^{\pm}\in C^2$, which this identity defines.

\subsubsection{The inequality across a mirror}\label{sec:inequality-across-mirror}

\begin{proposition}\label{prop:gaussian-bounds-1}
Let $\chi$ be a character, let $\nu\in\{1,2,3\}$ be a direction, and let $\chi'$ be the character with $\chi'_{\nu}=-1$ and $\chi'_i=\chi_i$ for $i\ne\nu$. Let $\jmath\in\mathbb{Z}/n$ be a layer for the time direction $\nu$, and for $(\mathfrak{h},b)\in C^2\times C^1$ let $(\mathfrak{h},b)^{\pm}=(\mathfrak{h},b)^{\pm;\nu,\jmath}\in C^2\times C^1$ be the elements defined in Section~\ref{sec:reflected-data}. Then for every $(\mathfrak{h},b)\in C^2\times C^1$
\begin{equation}\label{eq:gaussian-bounds-21}
Z^{\chi'}_{\varepsilon}\left((\mathfrak{h},b)^{\pm}\right)\ge0,\qquad\left|Z^{\chi}_{\varepsilon}(\mathfrak{h},b)\right|^2\le Z^{\chi'}_{\varepsilon}\left((\mathfrak{h},b)^{+}\right)\cdot Z^{\chi'}_{\varepsilon}\left((\mathfrak{h},b)^{-}\right).
\end{equation}
\end{proposition}

\begin{proof}
Traces are taken on $\mathcal{F}$, and $\|\cdot\|$ is the operator norm. Let $\theta\in\mathbb{R}$ be such that $\chi_{\nu}=-\exp(i\theta)$. The characters $\chi$ and $\chi'$ have the same components in the two directions different from $\nu$. Hence the space $\mathcal{V}$ of Section~\ref{sec:slices-slice-operators}, and with it the space $\mathcal{F}$, the operators $Q_{\mathbf{x}}$, $k_{\lambda}(y)$, $\varpi(p)$, $\mathsf{u}_{\lambda}(p)$ and the slice weights $\mathcal{W}_{\omega}(A)$ (Section~\ref{sec:weight-slice}), are the same for $\chi'$ as for $\chi$; moreover $\chi'_{\nu}=-1=-\exp(i\cdot0)$. Let $\mathcal{Z}^{\theta}$ be the functional of Section~\ref{sec:chain-formula}, given by the right side of \eqref{eq:gaussian-bounds-18}. By Proposition~\ref{prop:chain-formula-partition}, applied to $\chi$ and to $\chi'$, the three numbers that occur in \eqref{eq:gaussian-bounds-21} are the value of $\mathcal{Z}^{\theta}$ at the family of data of $(\mathfrak{h},b)$ and the values of $\mathcal{Z}^0$ at the families of data of $(\mathfrak{h},b)^{+}$ and of $(\mathfrak{h},b)^{-}$. By \eqref{eq:gaussian-bounds-19} none of these values changes when the slices and layers are relabelled by $t\mapsto t+\jmath-m$ as in Section~\ref{sec:reflected-data}. It therefore suffices to prove \eqref{eq:gaussian-bounds-21} in the normal labelling of Section~\ref{sec:reflected-data}, in which $\jmath=m$. In this labelling let $\mathsf{d}=(\omega_t,\lambda_t)_t$ be the family of data of $(\mathfrak{h},b)$; then $\mathsf{d}^{+}$ and $\mathsf{d}^{-}$ are the families of data of $(\mathfrak{h},b)^{+}$ and of $(\mathfrak{h},b)^{-}$, and
\[
Z^{\chi}_{\varepsilon}(\mathfrak{h},b)=\mathcal{Z}^{\theta}[\mathsf{d}],\qquad Z^{\chi'}_{\varepsilon}\left((\mathfrak{h},b)^{\pm}\right)=\mathcal{Z}^0[\mathsf{d}^{\pm}].
\]
For variables $A_t\in\mathbb{R}^l$, the indices being read modulo $n$, write
\[
k_t=k_{\lambda_t}\left(A_{t+1}-A_t\right),\qquad\mathcal{W}_t=\mathcal{W}_{\omega_t}(A_t).
\]

Recall that the operators $\varpi(p)$ and $\mathsf{u}_{\lambda}(p)$ are diagonal and depend continuously on $p\in\mathbb{R}^l$, that $\varpi(p)$ is positive definite and that $\mathsf{u}_{\lambda}(p)$ is unitary (Lemma~\ref{lem:fourier-form-layer}(a) and \eqref{eq:gaussian-bounds-14}); that $\exp\left(i\theta Q_{\mathrm{tot}}\right)$ is diagonal and unitary by \eqref{eq:gaussian-bounds-11}; that any two diagonal operators commute, the positive square root $\varpi(p)^{\frac{1}{2}}$ being again diagonal; and that every slice weight $\mathcal{W}_{\omega}(A)$ is positive semidefinite, in particular self-adjoint, and depends continuously on $A$ (Section~\ref{sec:weight-slice}). The map $M\mapsto M^{\frac{1}{2}}$ is continuous on the set of positive semidefinite operators on $\mathcal{F}$.

Let $(\Omega,\varrho)$ be a measure space, and let $X$ and $Y$ be measurable maps from $\Omega$ to the linear operators on $\mathcal{F}$ such that $\int_{\Omega}\mathrm{d}\varrho\,\operatorname{Tr}(XX^{*})<\infty$ and $\int_{\Omega}\mathrm{d}\varrho\,\operatorname{Tr}(Y^{*}Y)<\infty$. Then $\operatorname{Tr}(XY)$ is $\varrho$-integrable, and
\begin{equation}\label{eq:gaussian-bounds-22}
\left|\int_{\Omega}\mathrm{d}\varrho\,\operatorname{Tr}(XY)\right|^2\le\int_{\Omega}\mathrm{d}\varrho\,\operatorname{Tr}(XX^{*})\cdot\int_{\Omega}\mathrm{d}\varrho\,\operatorname{Tr}(Y^{*}Y).
\end{equation}
Indeed, \eqref{eq:gaussian-bounds-22} is the Cauchy--Schwarz inequality for the inner product $(X_1,X_2)\mapsto\int_{\Omega}\mathrm{d}\varrho\,\operatorname{Tr}(X_1^{*}X_2)$ on the measurable maps $Z$ from $\Omega$ to the linear operators on $\mathcal{F}$ with $\int_{\Omega}\mathrm{d}\varrho\,\operatorname{Tr}(Z^{*}Z)<\infty$, applied to $X_1=X^{*}$ and $X_2=Y$.

\emph{Odd $n$.} Let $n=2m+1$ with $m\ge1$. Then $t$ ranges over $\{-m,\dots,m\}$, and $A_{m+1}=A_{-m}$. Proposition~\ref{prop:chain-formula-partition}, with $t_0=-m$ in \eqref{eq:gaussian-bounds-18}, gives
\begin{multline*}
Z^{\chi}_{\varepsilon}(\mathfrak{h},b)=\int_{(\mathbb{R}^l)^n}\prod_{t=-m}^m\mathrm{d}A_t\,\operatorname{Tr}\Bigl[\exp\left(i\theta Q_{\mathrm{tot}}\right)\cdot k_m\cdot\mathcal{W}_m\cdot k_{m-1}\cdot\mathcal{W}_{m-1}\cdots\mathcal{W}_1\cdot k_0\\
\cdot\mathcal{W}_0\cdot k_{-1}\cdot\mathcal{W}_{-1}\cdots k_{-m}\cdot\mathcal{W}_{-m}\Bigr],
\end{multline*}
where $k_m=k_{\lambda_m}(A_{-m}-A_m)$. For $p,A_0\in\mathbb{R}^l$ define operators on $\mathcal{F}$ by
\begin{align*}
\Phi_{+}(p,A_0)&=\int_{(\mathbb{R}^l)^m}\prod_{t=1}^m\mathrm{d}A_t\,\exp\left(-i\langle p,A_m\rangle\right)\cdot\mathcal{W}_m\cdot k_{m-1}\cdot\mathcal{W}_{m-1}\cdots\mathcal{W}_1\cdot k_0,\\
\Phi_{-}(A_0,p)&=\int_{(\mathbb{R}^l)^m}\prod_{t=1}^m\mathrm{d}A_{-t}\,\exp\left(i\langle p,A_{-m}\rangle\right)\cdot k_{-1}\cdot\mathcal{W}_{-1}\cdots k_{-m}\cdot\mathcal{W}_{-m}.
\end{align*}
Each integrand contains $m$ factors $k_t$, each of norm at most $\left(\frac{2\pi}{\varepsilon}\right)^{\frac{v}{2}}$ by Lemma~\ref{lem:fourier-form-layer}(c), and $m$ factors $\mathcal{W}_t$, each of norm at most $c_W\cdot\exp\left(-\frac{\varepsilon}{2}\left|A_t+b^s_t\right|^2\right)$ by \eqref{eq:gaussian-bounds-17}. So the norm of the integrand is bounded by $\left(\left(\frac{2\pi}{\varepsilon}\right)^{\frac{v}{2}}\cdot c_W\right)^m$ times the product of the functions $\exp\left(-\frac{\varepsilon}{2}\left|A_t+b^s_t\right|^2\right)$ over the $m$ variables of integration. This bound does not depend on $(p,A_0)$, and the integral of each of these $m$ functions over $\mathbb{R}^l$ is $\left(\frac{2\pi}{\varepsilon}\right)^{\frac{l}{2}}$. Hence the two integrals converge absolutely, and
\[
\sup_{\mathbb{R}^l\times\mathbb{R}^l}\left\|\Phi_{\pm}\right\|\le\left(\left(\frac{2\pi}{\varepsilon}\right)^{\frac{v}{2}}\cdot c_W\cdot\left(\frac{2\pi}{\varepsilon}\right)^{\frac{l}{2}}\right)^m<\infty.
\]
The variable $A_0$ enters the integrands only through $k_0=k_{\lambda_0}(A_1-A_0)$, respectively $k_{-1}=k_{\lambda_{-1}}(A_0-A_{-1})$, and $p$ only through a phase factor. By Lemma~\ref{lem:fourier-form-layer}(e) the integrands are therefore continuous in $(p,A_0)$, and dominated convergence, with the majorant above, which does not depend on $(p,A_0)$, shows that $\Phi_{+}$ and $\Phi_{-}$ are continuous on $\mathbb{R}^l\times\mathbb{R}^l$.

Now insert \eqref{eq:gaussian-bounds-15} for $k_m$, with $y=A_{-m}-A_m$ and $\exp\left(i\langle p,y\rangle\right)=\exp\left(i\langle p,A_{-m}\rangle\right)\cdot\exp\left(-i\langle p,A_m\rangle\right)$. The integral over $(p,(A_t)_t)\in\mathbb{R}^l\times(\mathbb{R}^l)^n$ so obtained converges absolutely by Lemma~\ref{lem:fourier-form-layer}(a),(c) and \eqref{eq:gaussian-bounds-17}. By Fubini's theorem the integrations over $A_1,\dots,A_m$ and over $A_{-1},\dots,A_{-m}$ may be carried out first; they produce $\Phi_{+}$ and $\Phi_{-}$, and one obtains
\begin{equation}\label{eq:gaussian-bounds-23}
Z^{\chi}_{\varepsilon}(\mathfrak{h},b)=\int_{\mathbb{R}^l\times\mathbb{R}^l}\mathrm{d}p\,\mathrm{d}A_0\,\operatorname{Tr}\left[\mathsf{w}(p)\cdot\varpi(p)\cdot\Phi_{+}(p,A_0)\cdot\mathcal{W}_{\omega_0}(A_0)\cdot\Phi_{-}(A_0,p)\right],
\end{equation}
where $\mathsf{w}(p)=\exp\left(i\theta Q_{\mathrm{tot}}\right)\cdot\mathsf{u}_{\lambda_m}(p)$.

The operator $\mathsf{w}(p)$ is diagonal and unitary; being diagonal, it commutes with $\varpi(p)^{\frac{1}{2}}$. On $\Omega=\mathbb{R}^l\times\mathbb{R}^l$, with the Lebesgue measure $\mathrm{d}p\,\mathrm{d}A_0$, put
\[
X=\varpi(p)^{\frac{1}{2}}\cdot\Phi_{+}(p,A_0)\cdot\mathcal{W}_{\omega_0}(A_0)^{\frac{1}{2}},\qquad Y=\mathcal{W}_{\omega_0}(A_0)^{\frac{1}{2}}\cdot\Phi_{-}(A_0,p)\cdot\varpi(p)^{\frac{1}{2}}\cdot\mathsf{w}(p).
\]
These are continuous functions of $(p,A_0)$, because $\Phi_{\pm}$, $\varpi$, $\mathsf{w}$ and $\mathcal{W}_{\omega_0}$ are continuous and the square root is continuous on positive semidefinite operators; in particular $X$ and $Y$ are measurable. By cyclicity of the trace, and since $\mathsf{w}(p)$ commutes with $\varpi(p)^{\frac{1}{2}}$, $\operatorname{Tr}(XY)$ is the integrand of \eqref{eq:gaussian-bounds-23}. Again by cyclicity, and since $\mathsf{w}(p)\cdot\mathsf{w}(p)^{*}=1$,
\begin{align*}
\operatorname{Tr}(XX^{*})&=\operatorname{Tr}\left[\varpi(p)\cdot\Phi_{+}(p,A_0)\cdot\mathcal{W}_{\omega_0}(A_0)\cdot\Phi_{+}(p,A_0)^{*}\right],\\
\operatorname{Tr}(Y^{*}Y)&=\operatorname{Tr}\left[\varpi(p)\cdot\Phi_{-}(A_0,p)^{*}\cdot\mathcal{W}_{\omega_0}(A_0)\cdot\Phi_{-}(A_0,p)\right].
\end{align*}
Let
\begin{align*}
I_{+}&=\int_{\mathbb{R}^l\times\mathbb{R}^l}\mathrm{d}p\,\mathrm{d}A_0\,\operatorname{Tr}\left[\varpi(p)\cdot\Phi_{+}(p,A_0)\cdot\mathcal{W}_{\omega_0}(A_0)\cdot\Phi_{+}(p,A_0)^{*}\right]\ge0,\\
I_{-}&=\int_{\mathbb{R}^l\times\mathbb{R}^l}\mathrm{d}p\,\mathrm{d}A_0\,\operatorname{Tr}\left[\varpi(p)\cdot\Phi_{-}(A_0,p)^{*}\cdot\mathcal{W}_{\omega_0}(A_0)\cdot\Phi_{-}(A_0,p)\right]\ge0.
\end{align*}
The integrands are $\operatorname{Tr}(XX^{*})\ge0$ and $\operatorname{Tr}(Y^{*}Y)\ge0$, and both integrals are finite by Lemma~\ref{lem:fourier-form-layer}(a), \eqref{eq:gaussian-bounds-17} and the bound on $\Phi_{\pm}$. Therefore \eqref{eq:gaussian-bounds-22} applies and gives
\[
\left|Z^{\chi}_{\varepsilon}(\mathfrak{h},b)\right|^2\le I_{+}\cdot I_{-}.
\]

For a family of data $\mathsf{d}'=(\omega'_t,\lambda'_t)_t$, labelled by $t\in\{-m,\dots,m\}$, let $\Phi_{\pm}[\mathsf{d}']$ be the functions defined as $\Phi_{\pm}$ above, with $\omega'_t,\lambda'_t$ in place of $\omega_t,\lambda_t$; thus $\Phi_{\pm}=\Phi_{\pm}[\mathsf{d}]$. The computation that leads from \eqref{eq:gaussian-bounds-18} to \eqref{eq:gaussian-bounds-23} uses no property of $\mathsf{d}$ and $\theta$. Hence, for every such family $\mathsf{d}'$ and every $\theta'\in\mathbb{R}$, the number $\mathcal{Z}^{\theta'}[\mathsf{d}']$ equals the right side of \eqref{eq:gaussian-bounds-23} with $\Phi_{\pm}[\mathsf{d}']$, $\mathcal{W}_{\omega'_0}$ and $\exp\left(i\theta'Q_{\mathrm{tot}}\right)\cdot\mathsf{u}_{\lambda'_m}(p)$ in place of $\Phi_{\pm}$, $\mathcal{W}_{\omega_0}$ and $\mathsf{w}(p)$.

\emph{The value of $I_{+}$.} In $\Phi_{+}(p,A_0)^{*}$ rename the variables of integration $A_1,\dots,A_m$ as $A_{-1},\dots,A_{-m}$. Taking the adjoint reverses the order of the factors, and the slice weights are self-adjoint; therefore
\begin{multline*}
\Phi_{+}(p,A_0)^{*}=\int_{(\mathbb{R}^l)^m}\prod_{t=1}^m\mathrm{d}A_{-t}\,\exp\left(i\langle p,A_{-m}\rangle\right)\cdot k_{\lambda_0}(A_{-1}-A_0)^{*}\cdot\mathcal{W}_{\omega_1}(A_{-1})\\
\cdot k_{\lambda_1}(A_{-2}-A_{-1})^{*}\cdots k_{\lambda_{m-1}}(A_{-m}-A_{-m+1})^{*}\cdot\mathcal{W}_{\omega_m}(A_{-m}).
\end{multline*}
By the definition of $\mathsf{d}^{+}$ one has $\lambda^{+}_{-t-1}=-\lambda_t$ for $0\le t\le m-1$ and $\omega^{+}_{-t}=\omega_t$ for $1\le t\le m$; together with Lemma~\ref{lem:fourier-form-layer}(d) this gives
\begin{align*}
k_{\lambda_t}(A_{-t-1}-A_{-t})^{*}&=k_{-\lambda_t}(A_{-t}-A_{-t-1})=k_{\lambda^{+}_{-t-1}}\left(A_{(-t-1)+1}-A_{-t-1}\right)&&(0\le t\le m-1),\\
\mathcal{W}_{\omega_t}(A_{-t})&=\mathcal{W}_{\omega^{+}_{-t}}(A_{-t})&&(1\le t\le m).
\end{align*}
Consequently $\Phi_{+}(p,A_0)^{*}=\Phi_{-}[\mathsf{d}^{+}](A_0,p)$. On the other hand $\omega^{+}_t=\omega_t$ for $0\le t\le m$ and $\lambda^{+}_t=\lambda_t$ for $0\le t\le m-1$, so that $\Phi_{+}[\mathsf{d}^{+}]=\Phi_{+}$ and $\mathcal{W}_{\omega^{+}_0}=\mathcal{W}_{\omega_0}$; and $\lambda^{+}_m=0$, so that $\mathsf{u}_{\lambda^{+}_m}(p)=1$ by Lemma~\ref{lem:fourier-form-layer}(a). Thus the right side of \eqref{eq:gaussian-bounds-23}, formed with the family $\mathsf{d}^{+}$ and with $\theta'=0$, is the integral that defines $I_{+}$, and
\[
I_{+}=\mathcal{Z}^0[\mathsf{d}^{+}]=Z^{\chi'}_{\varepsilon}\left((\mathfrak{h},b)^{+}\right).
\]

\emph{The value of $I_{-}$.} In $\Phi_{-}(A_0,p)^{*}$ rename the variables of integration $A_{-1},\dots,A_{-m}$ as $A_1,\dots,A_m$. Since the adjoint reverses the order of the factors and the slice weights are self-adjoint,
\begin{multline*}
\Phi_{-}(A_0,p)^{*}=\int_{(\mathbb{R}^l)^m}\prod_{t=1}^m\mathrm{d}A_t\,\exp\left(-i\langle p,A_m\rangle\right)\cdot\mathcal{W}_{\omega_{-m}}(A_m)\cdot k_{\lambda_{-m}}(A_{m-1}-A_m)^{*}\\
\cdots\mathcal{W}_{\omega_{-1}}(A_1)\cdot k_{\lambda_{-1}}(A_0-A_1)^{*}.
\end{multline*}
By the definition of $\mathsf{d}^{-}$ one has $\lambda^{-}_t=-\lambda_{-t-1}$ for $0\le t\le m-1$ and $\omega^{-}_t=\omega_{-t}$ for $1\le t\le m$; together with Lemma~\ref{lem:fourier-form-layer}(d) this gives
\begin{align*}
k_{\lambda_{-t-1}}(A_t-A_{t+1})^{*}&=k_{-\lambda_{-t-1}}(A_{t+1}-A_t)=k_{\lambda^{-}_t}(A_{t+1}-A_t)&&(0\le t\le m-1),\\
\mathcal{W}_{\omega_{-t}}(A_t)&=\mathcal{W}_{\omega^{-}_t}(A_t)&&(1\le t\le m).
\end{align*}
Consequently $\Phi_{-}(A_0,p)^{*}=\Phi_{+}[\mathsf{d}^{-}](p,A_0)$. Moreover $\omega^{-}_t=\omega_t$ for $-m\le t\le0$ and $\lambda^{-}_t=\lambda_t$ for $-m\le t\le-1$, so that $\Phi_{-}[\mathsf{d}^{-}]=\Phi_{-}$ and $\mathcal{W}_{\omega^{-}_0}=\mathcal{W}_{\omega_0}$; and $\lambda^{-}_m=0$, so that $\mathsf{u}_{\lambda^{-}_m}(p)=1$ by Lemma~\ref{lem:fourier-form-layer}(a). Thus the right side of \eqref{eq:gaussian-bounds-23}, formed with the family $\mathsf{d}^{-}$ and with $\theta'=0$, is the integral that defines $I_{-}$, and
\[
I_{-}=\mathcal{Z}^0[\mathsf{d}^{-}]=Z^{\chi'}_{\varepsilon}\left((\mathfrak{h},b)^{-}\right).
\]
Since $I_{+}\ge0$, $I_{-}\ge0$ and $\left|Z^{\chi}_{\varepsilon}(\mathfrak{h},b)\right|^2\le I_{+}\cdot I_{-}$, this proves \eqref{eq:gaussian-bounds-21} for odd $n$.

\emph{Even $n$.} Let $n=2m$ with $m\ge1$. Then $t$ ranges over $\{-m+1,\dots,m\}$, and $A_{m+1}=A_{-m+1}$. Proposition~\ref{prop:chain-formula-partition}, with $t_0=-m+1$ in \eqref{eq:gaussian-bounds-18}, gives
\begin{multline*}
Z^{\chi}_{\varepsilon}(\mathfrak{h},b)=\int_{(\mathbb{R}^l)^n}\prod_{t=-m+1}^m\mathrm{d}A_t\,\operatorname{Tr}\Bigl[\exp\left(i\theta Q_{\mathrm{tot}}\right)\cdot k_m\cdot\mathcal{W}_m\cdot k_{m-1}\cdots k_1\cdot\mathcal{W}_1\\
\cdot k_0\cdot\mathcal{W}_0\cdot k_{-1}\cdot\mathcal{W}_{-1}\cdots k_{-m+1}\cdot\mathcal{W}_{-m+1}\Bigr],
\end{multline*}
where $k_m=k_{\lambda_m}(A_{-m+1}-A_m)$ and $k_0=k_{\lambda_0}(A_1-A_0)$. For $p,p'\in\mathbb{R}^l$ define operators on $\mathcal{F}$ by
\begin{align*}
\Phi_{+}(p,p')&=\int_{(\mathbb{R}^l)^m}\prod_{t=1}^m\mathrm{d}A_t\,\exp\left(-i\langle p,A_m\rangle\right)\cdot\mathcal{W}_m\cdot k_{m-1}\cdots k_1\cdot\mathcal{W}_1\cdot\exp\left(i\langle p',A_1\rangle\right),\\
\Phi_{-}(p',p)&=\int_{(\mathbb{R}^l)^m}\prod_{t=-m+1}^0\mathrm{d}A_t\,\exp\left(-i\langle p',A_0\rangle\right)\cdot\mathcal{W}_0\cdot k_{-1}\cdots k_{-m+1}\cdot\mathcal{W}_{-m+1}\cdot\exp\left(i\langle p,A_{-m+1}\rangle\right).
\end{align*}
As in the odd case, Lemma~\ref{lem:fourier-form-layer}(c) and \eqref{eq:gaussian-bounds-17} bound the norms of the integrands by an integrable function of the variables of integration which does not depend on $(p,p')$. Hence the integrals converge absolutely and $\Phi_{+}$, $\Phi_{-}$ are bounded on $\mathbb{R}^l\times\mathbb{R}^l$. The variables $p,p'$ enter the integrands only through the phase factors displayed, which are continuous and of modulus one; so $\Phi_{+}$ and $\Phi_{-}$ are continuous by dominated convergence.

Insert \eqref{eq:gaussian-bounds-15} for $k_m$, with the variable $p$, and for $k_0$, with the variable $p'$, and write $\exp\left(i\langle p,A_{-m+1}-A_m\rangle\right)=\exp\left(i\langle p,A_{-m+1}\rangle\right)\cdot\exp\left(-i\langle p,A_m\rangle\right)$ and $\exp\left(i\langle p',A_1-A_0\rangle\right)=\exp\left(i\langle p',A_1\rangle\right)\cdot\exp\left(-i\langle p',A_0\rangle\right)$. The integral over $(p,p',(A_t)_t)$ so obtained converges absolutely by Lemma~\ref{lem:fourier-form-layer}(a),(c) and \eqref{eq:gaussian-bounds-17}, and Fubini's theorem gives
\begin{equation}\label{eq:gaussian-bounds-24}
Z^{\chi}_{\varepsilon}(\mathfrak{h},b)=\int_{\mathbb{R}^l\times\mathbb{R}^l}\mathrm{d}p\,\mathrm{d}p'\,\operatorname{Tr}\left[\mathsf{w}(p)\cdot\varpi(p)\cdot\Phi_{+}(p,p')\cdot\mathsf{u}_{\lambda_0}(p')\cdot\varpi(p')\cdot\Phi_{-}(p',p)\right],
\end{equation}
with $\mathsf{w}(p)=\exp\left(i\theta Q_{\mathrm{tot}}\right)\cdot\mathsf{u}_{\lambda_m}(p)$ as before. On $\Omega=\mathbb{R}^l\times\mathbb{R}^l$, with the Lebesgue measure $\mathrm{d}p\,\mathrm{d}p'$, put
\[
X=\varpi(p)^{\frac{1}{2}}\cdot\Phi_{+}(p,p')\cdot\varpi(p')^{\frac{1}{2}},\qquad Y=\mathsf{u}_{\lambda_0}(p')\cdot\varpi(p')^{\frac{1}{2}}\cdot\Phi_{-}(p',p)\cdot\varpi(p)^{\frac{1}{2}}\cdot\mathsf{w}(p).
\]
Both are continuous, by the continuity of $\Phi_{\pm}$, $\varpi$, $\mathsf{u}_{\lambda_0}$, $\mathsf{w}$ and of the square root, hence measurable. As $\mathsf{u}_{\lambda_0}(p')$ commutes with $\varpi(p')^{\frac{1}{2}}$ and $\mathsf{w}(p)$ with $\varpi(p)^{\frac{1}{2}}$, cyclicity of the trace shows that $\operatorname{Tr}(XY)$ is the integrand of \eqref{eq:gaussian-bounds-24}; by cyclicity again, and as $\mathsf{u}_{\lambda_0}(p')$ and $\mathsf{w}(p)$ are unitary, $\operatorname{Tr}(XX^{*})$ and $\operatorname{Tr}(Y^{*}Y)$ are the integrands of
\begin{align*}
I_{+}&=\int_{\mathbb{R}^l\times\mathbb{R}^l}\mathrm{d}p\,\mathrm{d}p'\,\operatorname{Tr}\left[\varpi(p)\cdot\Phi_{+}(p,p')\cdot\varpi(p')\cdot\Phi_{+}(p,p')^{*}\right],\\
I_{-}&=\int_{\mathbb{R}^l\times\mathbb{R}^l}\mathrm{d}p\,\mathrm{d}p'\,\operatorname{Tr}\left[\varpi(p)\cdot\Phi_{-}(p',p)^{*}\cdot\varpi(p')\cdot\Phi_{-}(p',p)\right].
\end{align*}
Thus $I_{+}\ge0$ and $I_{-}\ge0$, and both are finite by Lemma~\ref{lem:fourier-form-layer}(a) and the boundedness of $\Phi_{\pm}$. Now \eqref{eq:gaussian-bounds-22} gives $\left|Z^{\chi}_{\varepsilon}(\mathfrak{h},b)\right|^2\le I_{+}\cdot I_{-}$.

For a family of data $\mathsf{d}'=(\omega'_t,\lambda'_t)_t$, labelled by $t\in\{-m+1,\dots,m\}$, let $\Phi_{\pm}[\mathsf{d}']$ be the functions defined as $\Phi_{\pm}$ above, with $\omega'_t,\lambda'_t$ in place of $\omega_t,\lambda_t$. As the computation that leads from \eqref{eq:gaussian-bounds-18} to \eqref{eq:gaussian-bounds-24} uses no property of $\mathsf{d}$ and $\theta$, for every such family $\mathsf{d}'$ and every $\theta'\in\mathbb{R}$ the number $\mathcal{Z}^{\theta'}[\mathsf{d}']$ equals the right side of \eqref{eq:gaussian-bounds-24} with $\Phi_{\pm}[\mathsf{d}']$, $\exp\left(i\theta'Q_{\mathrm{tot}}\right)\cdot\mathsf{u}_{\lambda'_m}(p)$ and $\mathsf{u}_{\lambda'_0}(p')$ in place of $\Phi_{\pm}$, $\mathsf{w}(p)$ and $\mathsf{u}_{\lambda_0}(p')$.

\emph{The value of $I_{+}$.} In $\Phi_{+}(p,p')^{*}$ rename the variable of integration $A_t$ as $A_{1-t}$, for $1\le t\le m$. Since the adjoint reverses the order of the factors and the slice weights are self-adjoint,
\begin{multline*}
\Phi_{+}(p,p')^{*}=\int_{(\mathbb{R}^l)^m}\prod_{t=1}^m\mathrm{d}A_{1-t}\,\exp\left(-i\langle p',A_0\rangle\right)\cdot\mathcal{W}_{\omega_1}(A_0)\cdot k_{\lambda_1}(A_{-1}-A_0)^{*}\cdot\mathcal{W}_{\omega_2}(A_{-1})\\
\cdots k_{\lambda_{m-1}}(A_{-m+1}-A_{-m+2})^{*}\cdot\mathcal{W}_{\omega_m}(A_{-m+1})\cdot\exp\left(i\langle p,A_{-m+1}\rangle\right).
\end{multline*}
By the definition of $\mathsf{d}^{+}$ one has $\lambda^{+}_{-t}=-\lambda_t$ for $1\le t\le m-1$ and $\omega^{+}_{1-t}=\omega_t$ for $1\le t\le m$; together with Lemma~\ref{lem:fourier-form-layer}(d) this gives
\begin{align*}
k_{\lambda_t}(A_{-t}-A_{1-t})^{*}&=k_{-\lambda_t}(A_{1-t}-A_{-t})=k_{\lambda^{+}_{-t}}(A_{-t+1}-A_{-t})&&(1\le t\le m-1),\\
\mathcal{W}_{\omega_t}(A_{1-t})&=\mathcal{W}_{\omega^{+}_{1-t}}(A_{1-t})&&(1\le t\le m).
\end{align*}
Consequently $\Phi_{+}(p,p')^{*}=\Phi_{-}[\mathsf{d}^{+}](p',p)$. Since $\omega^{+}_t=\omega_t$ for $1\le t\le m$ and $\lambda^{+}_t=\lambda_t$ for $1\le t\le m-1$, also $\Phi_{+}[\mathsf{d}^{+}]=\Phi_{+}$; and $\lambda^{+}_0=\lambda^{+}_m=0$, so that $\mathsf{u}_{\lambda^{+}_0}=\mathsf{u}_{\lambda^{+}_m}=1$ by Lemma~\ref{lem:fourier-form-layer}(a). Thus the right side of \eqref{eq:gaussian-bounds-24}, formed with the family $\mathsf{d}^{+}$ and with $\theta'=0$, is the integral that defines $I_{+}$, and
\[
I_{+}=\mathcal{Z}^0[\mathsf{d}^{+}]=Z^{\chi'}_{\varepsilon}\left((\mathfrak{h},b)^{+}\right).
\]

\emph{The value of $I_{-}$.} In $\Phi_{-}(p',p)^{*}$ rename the variable of integration $A_u$ as $A_{1-u}$, for $-m+1\le u\le0$. Since the adjoint reverses the order of the factors and the slice weights are self-adjoint,
\begin{multline*}
\Phi_{-}(p',p)^{*}=\int_{(\mathbb{R}^l)^m}\prod_{t=1}^m\mathrm{d}A_t\,\exp\left(-i\langle p,A_m\rangle\right)\cdot\mathcal{W}_{\omega_{-m+1}}(A_m)\cdot k_{\lambda_{-m+1}}(A_{m-1}-A_m)^{*}\\
\cdots\mathcal{W}_{\omega_{-1}}(A_2)\cdot k_{\lambda_{-1}}(A_1-A_2)^{*}\cdot\mathcal{W}_{\omega_0}(A_1)\cdot\exp\left(i\langle p',A_1\rangle\right).
\end{multline*}
By the definition of $\mathsf{d}^{-}$ one has $\lambda^{-}_t=-\lambda_{-t}$ for $1\le t\le m-1$ and $\omega^{-}_t=\omega_{1-t}$ for $1\le t\le m$; together with Lemma~\ref{lem:fourier-form-layer}(d) this gives
\begin{align*}
k_{\lambda_{-t}}(A_t-A_{t+1})^{*}&=k_{-\lambda_{-t}}(A_{t+1}-A_t)=k_{\lambda^{-}_t}(A_{t+1}-A_t)&&(1\le t\le m-1),\\
\mathcal{W}_{\omega_{1-t}}(A_t)&=\mathcal{W}_{\omega^{-}_t}(A_t)&&(1\le t\le m).
\end{align*}
Consequently $\Phi_{-}(p',p)^{*}=\Phi_{+}[\mathsf{d}^{-}](p,p')$. Since $\omega^{-}_t=\omega_t$ for $-m+1\le t\le0$ and $\lambda^{-}_t=\lambda_t$ for $-m+1\le t\le-1$, also $\Phi_{-}[\mathsf{d}^{-}]=\Phi_{-}$; and $\lambda^{-}_0=\lambda^{-}_m=0$, so that $\mathsf{u}_{\lambda^{-}_0}=\mathsf{u}_{\lambda^{-}_m}=1$ by Lemma~\ref{lem:fourier-form-layer}(a). Thus the right side of \eqref{eq:gaussian-bounds-24}, formed with the family $\mathsf{d}^{-}$ and with $\theta'=0$, is the integral that defines $I_{-}$, and
\[
I_{-}=\mathcal{Z}^0[\mathsf{d}^{-}]=Z^{\chi'}_{\varepsilon}\left((\mathfrak{h},b)^{-}\right).
\]
Since $I_{\pm}\ge0$ and $\left|Z^{\chi}_{\varepsilon}(\mathfrak{h},b)\right|^2\le I_{+}\cdot I_{-}$, this proves \eqref{eq:gaussian-bounds-21} for even $n$.
\end{proof}

\subsubsection{Reflection positivity of the Proca parent}\label{sec:reflection-positivity-proca}

In Proposition~\ref{prop:gaussian-bounds-2} and in its proof the letter $\vartheta$ stands for a reflection of $C^1$. In Lemma~\ref{lem:gaussian-bounds-1}(iv) and in Section~\ref{sec:vanishing-proca-mass} it stands for a vector of $\mathbb R^3$ by which a character $\chi$ is multiplied: $\chi\exp(i\vartheta)$ is the character with the components $\chi_\mu\cdot\exp(i\vartheta_\mu)$, $\mu=1,2,3$.

\begin{proposition}\label{prop:gaussian-bounds-2}
Let $\nu$ be the time direction, let $\chi$ be a character with $\chi_\nu=-1$, and number the slices and the layers by the normal labelling of Section~\ref{sec:reflected-data}. For $A\in C^1$ let $(A_t)_t$ and $(a_t)_t$ be its components in the slices and in the layers (Section~\ref{sec:slices-slice-operators}); indices are read modulo $n$.

If $n=2m+1$, let $\mathcal A_+$ be the set of bounded Borel functions $F\colon C^1\to\mathbb C$ that depend on $A$ only through $(A_t)_{0\le t\le m}$ and $(a_t)_{0\le t\le m-1}$, and let $\vartheta\colon C^1\to C^1$ be the linear map for which $\vartheta A$ has the components
\[
(\vartheta A)_t=A_{-t}\quad\text{in the slices},\qquad(\vartheta a)_t=-a_{-t-1}\quad\text{in the layers}.
\]
If $n=2m$, let $\mathcal A_+$ be the set of bounded Borel functions $F\colon C^1\to\mathbb C$ that depend on $A$ only through $(A_t)_{1\le t\le m}$ and $(a_t)_{1\le t\le m-1}$, and let $\vartheta\colon C^1\to C^1$ be the linear map for which $\vartheta A$ has the components
\[
(\vartheta A)_t=A_{1-t}\quad\text{in the slices},\qquad(\vartheta a)_t=-a_{-t}\quad\text{in the layers}.
\]
Then every $F\in\mathcal A_+$ satisfies
\[
\int_{C^1}\mathrm{d}A\,\overline{F(\vartheta A)}\cdot F(A)\cdot\exp\left(-\frac12\|d_1A\|^2-\frac{\varepsilon}{2}\|A\|^2\right)\cdot\Delta^\chi(A)\ge0.
\]
\end{proposition}

\begin{proof}
Let $I(F)$ be the integral in the assertion. For $A',A\in\mathbb R^l$ and $a\in\mathbb R^v$ let $\tilde k(A',A;a)$ be the operator on $\mathcal F$ given by
\[
\tilde k(A',A;a)=\exp\left(-\frac12\left|A'-A-d_sa\right|^2-\frac{\varepsilon}{2}|a|^2\right)\cdot\exp(iq\langle a,Q\rangle),
\]
with $\langle a,Q\rangle$ as in \eqref{eq:gaussian-bounds-11}. As $\langle a,Q\rangle$ is self-adjoint, $\exp(iq\langle a,Q\rangle)$ is unitary, so that $\left\|\tilde k(A',A;a)\right\|\le\exp\left(-\frac{\varepsilon}{2}|a|^2\right)$; and \eqref{eq:gaussian-bounds-13} with $\lambda=(0,0)$ says that
\[
\int_{\mathbb R^v}\mathrm{d}a\,\tilde k(A',A;a)=k_{(0,0)}(A'-A).
\]
For $A\in C^1$ and $t\in\mathbb Z/n$ write $\tilde k_t=\tilde k(A_{t+1},A_t;a_t)$ and $\mathcal W_t=\mathcal W_{(0,0)}(A_t)$, the slice weight \eqref{eq:gaussian-bounds-16} with $\omega=(0,0)$. By Section~\ref{sec:weight-slice}, $\mathcal W_t\ge0$; in particular $\mathcal W_t^*=\mathcal W_t$. By \eqref{eq:gaussian-bounds-17}, $\|\mathcal W_t\|\le c_W\cdot\exp\left(-\frac{\varepsilon}{2}|A_t|^2\right)$.

Since $\chi_\nu=-1$, Proposition~\ref{prop:trace-formula-determinant} applies with $\theta=0$. Insert \eqref{eq:gaussian-bounds-6} with $\mathfrak{h}=0$, $b=0$ and \eqref{eq:gaussian-bounds-12} with $\theta=0$ into the integrand of $I(F)$, and attach the Gaussian factors of the layer $t$ to $\exp(iq\langle a_t,Q\rangle)$ and those of the slice $t$ to $\det B(A_t)\cdot\Gamma\left(S(A_t)\right)$, as in the proof of Proposition~\ref{prop:chain-formula-partition}; by \eqref{eq:gaussian-bounds-16} this gives, for every $A\in C^1$ and every $t_0\in\mathbb Z/n$,
\begin{multline*}
\overline{F(\vartheta A)}\cdot F(A)\cdot\exp\left(-\frac12\|d_1A\|^2-\frac{\varepsilon}{2}\|A\|^2\right)\cdot\Delta^\chi(A)\\
=\operatorname{Tr}_{\mathcal F}\left[\overline{F(\vartheta A)}\cdot F(A)\cdot\tilde k_{t_0+n-1}\cdot\mathcal W_{t_0+n-1}\cdots\tilde k_{t_0}\cdot\mathcal W_{t_0}\right],
\end{multline*}
the choice of $t_0$ being free because the trace is cyclic. The product under the trace has norm at most
\[
M(A)=\sup|F|^2\cdot c_W^n\cdot\prod_{t\in\mathbb Z/n}\exp\left(-\frac{\varepsilon}{2}\left(|a_t|^2+|A_t|^2\right)\right),
\]
and consequently the absolute value of the integrand of $I(F)$ is at most $\dim\mathcal F\cdot M(A)$, which is integrable over $C^1$.

The integral $I(F)$ is now computed along the derivation of \eqref{eq:gaussian-bounds-23} if $n$ is odd, and of \eqref{eq:gaussian-bounds-24} if $n$ is even, in the proof of Proposition~\ref{prop:gaussian-bounds-1}, at $\mathfrak{h}=0$, $b=0$, $\theta=0$, with two changes. First, for every layer $t$ that the mirror does not fix, the variable $a_t$ is kept: the operator $\tilde k_t$ takes the place of $k_{(0,0)}(A_{t+1}-A_t)$, and the integration over $a_t$ is added. Second, the factor $F(A)$ is put under the integral that defines $\Phi_+$, and the factor $\overline{F(\vartheta A)}$ under the integral that defines $\Phi_-$; if $n$ is odd, both of them may involve the variable $A_0$ of the fixed slice. Neither $F$ nor $F\circ\vartheta$ involves the variable $a_t$ of a layer $t$ that the mirror fixes, namely $a_m$ if $n$ is odd, and $a_0$ and $a_m$ if $n$ is even. By Fubini's theorem these variables are integrated first, as in the proof of Proposition~\ref{prop:chain-formula-partition}; this replaces $\tilde k_t$ by $k_{(0,0)}(A_{t+1}-A_t)$ for each fixed layer $t$. Then \eqref{eq:gaussian-bounds-15} with $\lambda=(0,0)$, in which $\mathsf u_\lambda=1$ by Lemma~\ref{lem:fourier-form-layer}(a), is inserted:
\[
k_{(0,0)}(y)=\int_{\mathbb R^l}\mathrm{d}p\,\exp(i\langle p,y\rangle)\cdot\varpi(p)\qquad\left(y\in\mathbb R^l\right),
\]
with the integration variable $p$ for the layer $m$ and $p'$ for the layer $0$. By Lemma~\ref{lem:fourier-form-layer}(a), $\|\varpi(p)\|\le c_\varpi\cdot\exp\left(-\frac12|p|^2\right)$; hence, after this insertion, the norm of the product under the trace is bounded by $M(A)$ with the factor $\exp\left(-\frac{\varepsilon}{2}|a_m|^2\right)$ replaced by $c_\varpi\cdot\exp\left(-\frac12|p|^2\right)$ and, for even $n$, the factor $\exp\left(-\frac{\varepsilon}{2}|a_0|^2\right)$ replaced by $c_\varpi\cdot\exp\left(-\frac12|p'|^2\right)$. Therefore all the integrals below converge absolutely, and Fubini's theorem applies to them.

\emph{The case $n=2m+1$.} Here $F(A)$ depends only on $A_0,\dots,A_m$ and $a_0,\dots,a_{m-1}$, and $F(\vartheta A)$ only on $A_0,A_{-1},\dots,A_{-m}$ and $a_{-1},\dots,a_{-m}$. For $p,A_0\in\mathbb R^l$ define operators $\Phi_+(p,A_0)$ and $\Phi_-(A_0,p)$ on $\mathcal F$ by
\begin{align*}
\Phi_+(p,A_0)&=\int\prod_{t=1}^m\mathrm{d}A_t\prod_{t=0}^{m-1}\mathrm{d}a_t\,\exp(-i\langle p,A_m\rangle)\cdot F(A)\\
&\qquad\cdot\mathcal W_m\cdot\tilde k_{m-1}\cdot\mathcal W_{m-1}\cdots\mathcal W_1\cdot\tilde k_0,\\
\Phi_-(A_0,p)&=\int\prod_{t=1}^m\mathrm{d}A_{-t}\prod_{t=-m}^{-1}\mathrm{d}a_t\,\exp(i\langle p,A_{-m}\rangle)\cdot\overline{F(\vartheta A)}\\
&\qquad\cdot\tilde k_{-1}\cdot\mathcal W_{-1}\cdots\tilde k_{-m}\cdot\mathcal W_{-m}.
\end{align*}
Take $t_0=-m$ in the trace formula for the integrand. Then $\tilde k_m=\tilde k(A_{-m},A_m;a_m)$, which the integration over $a_m$ and the insertion of \eqref{eq:gaussian-bounds-15} turn into the integral over $p$ of $\exp(i\langle p,A_{-m}-A_m\rangle)\cdot\varpi(p)$, and Fubini's theorem yields
\[
I(F)=\int_{\mathbb R^l\times\mathbb R^l}\mathrm{d}p\,\mathrm{d}A_0\,\operatorname{Tr}\left[\varpi(p)\cdot\Phi_+(p,A_0)\cdot\mathcal W_{(0,0)}(A_0)\cdot\Phi_-(A_0,p)\right].
\]
Taking adjoints under the integral sign and using $\mathcal W_t^*=\mathcal W_t$,
\begin{multline*}
\Phi_+(p,A_0)^*=\int\prod_{t=1}^m\mathrm{d}A_t\prod_{t=0}^{m-1}\mathrm{d}a_t\,\exp(i\langle p,A_m\rangle)\cdot\overline{F(A)}\\
\cdot\tilde k_0^*\cdot\mathcal W_1\cdot\tilde k_1^*\cdots\tilde k_{m-1}^*\cdot\mathcal W_m.
\end{multline*}
In this integral denote the variable $A_t$ by $A_{-t}$ for $1\le t\le m$, and the variable $a_t$ by $-a_{-t-1}$ for $0\le t\le m-1$; this change of variables preserves Lebesgue measure. By the definition of $\vartheta$ it turns $\overline{F(A)}$ into $\overline{F(\vartheta A)}$; it turns $\exp(i\langle p,A_m\rangle)$ into $\exp(i\langle p,A_{-m}\rangle)$ and $\mathcal W_t$ into $\mathcal W_{-t}$. Moreover
\[
\tilde k(A',A;a)^*=\tilde k(A,A';-a),
\]
because $|A'-A-d_sa|=|A-A'-d_s(-a)|$ and the adjoint of $\exp(iq\langle a,Q\rangle)$ is $\exp(iq\langle -a,Q\rangle)$. For $0\le t\le m-1$ this identity gives $\tilde k_t^*=\tilde k(A_t,A_{t+1};-a_t)$, and the change of variables transforms this operator into $\tilde k(A_{-t},A_{-t-1};a_{-t-1})$, that is, into $\tilde k_{-t-1}$. It follows that $\Phi_+(p,A_0)^*=\Phi_-(A_0,p)$, and hence
\[
I(F)=\int_{\mathbb R^l\times\mathbb R^l}\mathrm{d}p\,\mathrm{d}A_0\,\operatorname{Tr}\left[\varpi(p)\cdot\Phi_+(p,A_0)\cdot\mathcal W_{(0,0)}(A_0)\cdot\Phi_+(p,A_0)^*\right]\ge0;
\]
indeed, $\varpi(p)$ and $\mathcal W_{(0,0)}(A_0)$ are positive semidefinite (Lemma~\ref{lem:fourier-form-layer}(a), Section~\ref{sec:weight-slice}), and by cyclicity the trace equals $\operatorname{Tr}(XX^*)\ge0$ with $X=\varpi(p)^{\frac12}\cdot\Phi_+(p,A_0)\cdot\mathcal W_{(0,0)}(A_0)^{\frac12}$.

\emph{The case $n=2m$.} Here $F(A)$ depends only on $A_1,\dots,A_m$ and $a_1,\dots,a_{m-1}$, and $F(\vartheta A)$ only on $A_0,A_{-1},\dots,A_{-m+1}$ and $a_{-1},\dots,a_{-m+1}$. For $p,p'\in\mathbb R^l$ define operators $\Phi_+(p,p')$ and $\Phi_-(p',p)$ on $\mathcal F$ by
\begin{align*}
\Phi_+(p,p')&=\int\prod_{t=1}^m\mathrm{d}A_t\prod_{t=1}^{m-1}\mathrm{d}a_t\,\exp(-i\langle p,A_m\rangle)\cdot F(A)\\
&\qquad\cdot\mathcal W_m\cdot\tilde k_{m-1}\cdots\tilde k_1\cdot\mathcal W_1\cdot\exp(i\langle p',A_1\rangle),\\
\Phi_-(p',p)&=\int\prod_{t=-m+1}^0\mathrm{d}A_t\prod_{t=-m+1}^{-1}\mathrm{d}a_t\,\exp(-i\langle p',A_0\rangle)\cdot\overline{F(\vartheta A)}\\
&\qquad\cdot\mathcal W_0\cdot\tilde k_{-1}\cdots\tilde k_{-m+1}\cdot\mathcal W_{-m+1}\cdot\exp(i\langle p,A_{-m+1}\rangle).
\end{align*}
Take $t_0=-m+1$ in the trace formula for the integrand. Then $\tilde k_m=\tilde k(A_{-m+1},A_m;a_m)$ and $\tilde k_0=\tilde k(A_1,A_0;a_0)$; the integrations over $a_m$ and $a_0$ and the insertion of \eqref{eq:gaussian-bounds-15} turn them into the integrals over $p$ of $\exp(i\langle p,A_{-m+1}-A_m\rangle)\cdot\varpi(p)$ and over $p'$ of $\exp(i\langle p',A_1-A_0\rangle)\cdot\varpi(p')$, and Fubini's theorem yields
\[
I(F)=\int_{\mathbb R^l\times\mathbb R^l}\mathrm{d}p\,\mathrm{d}p'\,\operatorname{Tr}\left[\varpi(p)\cdot\Phi_+(p,p')\cdot\varpi(p')\cdot\Phi_-(p',p)\right].
\]
In the integral that represents $\Phi_+(p,p')^*$, obtained by taking adjoints under the integral sign, denote the variable $A_t$ by $A_{1-t}$ for $1\le t\le m$, and the variable $a_t$ by $-a_{-t}$ for $1\le t\le m-1$. This change of variables preserves Lebesgue measure and turns $\overline{F(A)}$ into $\overline{F(\vartheta A)}$, $\mathcal W_t=\mathcal W_t^*$ into $\mathcal W_{1-t}$, the factor $\tilde k_t^*$ into $\tilde k(A_{1-t},A_{-t};a_{-t})=\tilde k_{-t}$, and $\exp(i\langle p,A_m\rangle)\cdot\exp(-i\langle p',A_1\rangle)$ into $\exp(i\langle p,A_{-m+1}\rangle)\cdot\exp(-i\langle p',A_0\rangle)$. It follows that $\Phi_+(p,p')^*=\Phi_-(p',p)$, and hence
\[
I(F)=\int_{\mathbb R^l\times\mathbb R^l}\mathrm{d}p\,\mathrm{d}p'\,\operatorname{Tr}\left[\varpi(p)\cdot\Phi_+(p,p')\cdot\varpi(p')\cdot\Phi_+(p,p')^*\right]\ge0,
\]
since $\varpi(p)$ and $\varpi(p')$ are positive definite (Lemma~\ref{lem:fourier-form-layer}(a)) and by cyclicity the trace equals $\operatorname{Tr}(YY^*)\ge0$ with $Y=\varpi(p)^{\frac12}\cdot\Phi_+(p,p')\cdot\varpi(p')^{\frac12}$.
\end{proof}

\subsection{Proof of Gaussian domination}\label{sec:proof-gaussian-domination}

\begin{lemma}[Gaussian decay and continuity]\label{lem:gaussian-decay-continuity}
Let $\chi$ be a character.
\begin{enumerate}
\item[(i)] For all $\mathfrak{h}\in C^2$ and $b\in C^1$ one has
\[
\left|Z^{\chi}_{\varepsilon}(\mathfrak{h},b)\right|\le c_0\cdot\exp\left(-\frac{\varepsilon}{2(16+\varepsilon)}\left\|\mathfrak{h}-d_1b\right\|^2\right),
\]
where
\[
c_0=\left(9+|\boldsymbol{\mu}|\right)^{sn^3}\cdot\int_{C^1}\mathrm{d}A\,\exp\left(-\frac{1}{2}\left\|d_1A\right\|^2-\frac{\varepsilon}{2}\|A\|^2\right).
\]
\item[(ii)] The function $C^2\times C^1\to\mathbb{C}$, $(\mathfrak{h},b)\mapsto Z^{\chi}_{\varepsilon}(\mathfrak{h},b)$, is continuous.
\item[(iii)] The function $C^2\to[0,\infty)$, $\mathfrak{h}\mapsto\left|Z^{\chi}_{\varepsilon}(\mathfrak{h})\right|$, attains its supremum.
\end{enumerate}
\end{lemma}

\begin{proof}
(i) Fix $\mathfrak{h}\in C^2$ and $b\in C^1$. By Lemma~\ref{lem:gaussian-bounds-1}(i), $\left|\Delta^{\chi}(A)\right|\le\left(9+|\boldsymbol{\mu}|\right)^{sn^3}$ for every $A\in C^1$, and therefore \eqref{eq:gaussian-bounds-2} yields
\[
\left|Z^{\chi}_{\varepsilon}(\mathfrak{h},b)\right|\le\left(9+|\boldsymbol{\mu}|\right)^{sn^3}\cdot\int_{C^1}\mathrm{d}A\,\exp\left(-\frac{1}{2}\left\|d_1A+\mathfrak{h}\right\|^2-\frac{\varepsilon}{2}\|A+b\|^2\right).
\]
Let $k=\mathfrak{h}-d_1b\in C^2$. The translation $A=A'-b$ leaves Lebesgue measure on $C^1$ invariant, and it turns $d_1A+\mathfrak{h}$ into $d_1A'+k$ and $A+b$ into $A'$. Let $\mathcal{C}_{\varepsilon}=\left(d_1^{*}d_1+\varepsilon\right)^{-1}$ be the covariance on $C^1$ introduced in Section~\ref{sec:partition-functions-sources}; it is self-adjoint and positive definite. Completing the square one finds, for every $A'\in C^1$,
\begin{align*}
&\frac{1}{2}\left\|d_1A'+k\right\|^2+\frac{\varepsilon}{2}\|A'\|^2\\
&\qquad=\frac{1}{2}\left\langle A'+\mathcal{C}_{\varepsilon}d_1^{*}k,\,\mathcal{C}_{\varepsilon}^{-1}\left(A'+\mathcal{C}_{\varepsilon}d_1^{*}k\right)\right\rangle+\frac{1}{2}\left\langle k,\left(1-d_1\mathcal{C}_{\varepsilon}d_1^{*}\right)k\right\rangle ,
\end{align*}
since each side equals $\frac{1}{2}\left\langle A',\mathcal{C}_{\varepsilon}^{-1}A'\right\rangle+\left\langle A',d_1^{*}k\right\rangle+\frac{1}{2}\|k\|^2$. The translation $A''=A'+\mathcal{C}_{\varepsilon}d_1^{*}k$ again preserves Lebesgue measure, and $\left\langle A,\mathcal{C}_{\varepsilon}^{-1}A\right\rangle=\left\|d_1A\right\|^2+\varepsilon\|A\|^2$ for $A\in C^1$. It follows that
\[
\left|Z^{\chi}_{\varepsilon}(\mathfrak{h},b)\right|\le c_0\cdot\exp\left(-\frac{1}{2}\left\langle k,\left(1-d_1\mathcal{C}_{\varepsilon}d_1^{*}\right)k\right\rangle\right).
\]
The constant $c_0$ is finite, the integrand in its definition being at most $\exp\left(-\frac{\varepsilon}{2}\|A\|^2\right)$. From $d_1\left(d_1^{*}d_1+\varepsilon\right)=\left(d_1d_1^{*}+\varepsilon\right)d_1$ one obtains $d_1\left(d_1^{*}d_1+\varepsilon\right)^{-1}=\left(d_1d_1^{*}+\varepsilon\right)^{-1}d_1$; as $d_1d_1^{*}$ commutes with $\left(d_1d_1^{*}+\varepsilon\right)^{-1}$, this gives, for operators on $C^2$,
\[
1-d_1\left(d_1^{*}d_1+\varepsilon\right)^{-1}d_1^{*}=1-d_1d_1^{*}\left(d_1d_1^{*}+\varepsilon\right)^{-1}=\varepsilon\left(d_1d_1^{*}+\varepsilon\right)^{-1}\ge\frac{\varepsilon}{\left\|d_1\right\|^2+\varepsilon},
\]
where the inequality holds because the spectrum of the self-adjoint operator $d_1d_1^{*}$ lies in the interval $\left[0,\left\|d_1\right\|^2\right]$. It remains to bound the operator norm of $d_1\colon C^1\to C^2$. For a plaquette $p$ let $\partial p$ be the set of its four links, and write $A_{\ell}$ for the value of $A\in C^1$ on the link $\ell$. The number $\left(d_1A\right)_p$ is a sum of the four terms $\pm A_{\ell}$, $\ell\in\partial p$, so the Cauchy--Schwarz inequality gives $\left(d_1A\right)_p^2\le4\sum_{\ell\in\partial p}A_{\ell}^2$. Each link belongs to exactly four plaquettes; summation over $p$ therefore yields
\[
\left\|d_1A\right\|^2\le4\sum_p\sum_{\ell\in\partial p}A_{\ell}^2=16\,\|A\|^2\qquad\left(A\in C^1\right).
\]
Thus $\left\|d_1\right\|^2\le16$, whence $\left\langle k,\left(1-d_1\mathcal{C}_{\varepsilon}d_1^{*}\right)k\right\rangle\ge\frac{\varepsilon}{16+\varepsilon}\,\|k\|^2$, and (i) is proved.

(ii) Let $\left(\mathfrak{h}_j,b_j\right)_{j\ge1}$ be a sequence in $C^2\times C^1$ converging to $(\mathfrak{h},b)$. Being convergent, it is bounded; put $\beta=\sup_j\left\|b_j\right\|$. For $A\in C^1$ one has $\left\|A+b_j\right\|\ge\left(\|A\|-\beta\right)_{+}$, where $(\cdot)_{+}$ denotes the positive part. Together with Lemma~\ref{lem:gaussian-bounds-1}(i), this shows that the integrand of \eqref{eq:gaussian-bounds-2} at $\left(\mathfrak{h}_j,b_j\right)$ is bounded in absolute value by
\[
\left(9+|\boldsymbol{\mu}|\right)^{sn^3}\cdot\exp\left(-\frac{\varepsilon}{2}\left(\|A\|-\beta\right)_{+}^2\right),
\]
a function of $A$ alone which is integrable over $C^1$ (for $\|A\|\ge2\beta$ it is at most $\left(9+|\boldsymbol{\mu}|\right)^{sn^3}\cdot\exp\left(-\frac{\varepsilon}{8}\|A\|^2\right)$). For each fixed $A$ the integrand at $\left(\mathfrak{h}_j,b_j\right)$ converges to the integrand at $(\mathfrak{h},b)$, since $\Delta^{\chi}(A)$ does not depend on the pair. By the dominated convergence theorem, $Z^{\chi}_{\varepsilon}\left(\mathfrak{h}_j,b_j\right)\to Z^{\chi}_{\varepsilon}(\mathfrak{h},b)$.

(iii) Let $M=\sup_{\mathfrak{h}\in C^2}\left|Z^{\chi}_{\varepsilon}(\mathfrak{h})\right|$; by (i), $M\le c_0<\infty$. If $M=0$, the supremum is attained at every point of $C^2$. Let $M>0$. By (i) with $b=0$,
\[
\left|Z^{\chi}_{\varepsilon}(\mathfrak{h})\right|\le c_0\cdot\exp\left(-\frac{\varepsilon}{2(16+\varepsilon)}\|\mathfrak{h}\|^2\right)\qquad\left(\mathfrak{h}\in C^2\right).
\]
Choose $R_0>0$ so large that $c_0\cdot\exp\left(-\frac{\varepsilon R_0^2}{2(16+\varepsilon)}\right)<M$. Then $\left|Z^{\chi}_{\varepsilon}(\mathfrak{h})\right|$ is bounded by this number, which is smaller than $M$, for every $\mathfrak{h}$ outside the closed ball of radius $R_0$ about $0$ in $C^2$; consequently the supremum of $\mathfrak{h}\mapsto\left|Z^{\chi}_{\varepsilon}(\mathfrak{h})\right|$ over this ball equals $M$. The ball is compact, $C^2$ being finite-dimensional, and the function is continuous by (ii). A continuous function attains its maximum on a compact set; hence $\left|Z^{\chi}_{\varepsilon}(\mathfrak{h}_0)\right|=M$ for some $\mathfrak{h}_0$ in the ball.
\end{proof}

\begin{proof}[Proof of Theorem~\ref{thm:gaussian-domination-positive}(a)]
Let $\mathcal{D}=\left\{\mu\in\{1,2,3\}:\chi_{\mu}=-1\right\}$; by hypothesis $\#|\mathcal{D}|\ge2$. If $\nu\in\mathcal{D}$, then the character $\chi'$ that Proposition~\ref{prop:gaussian-bounds-1} associates with $\chi$ and $\nu$ is $\chi$ itself. Apply \eqref{eq:gaussian-bounds-21} with a direction of $\mathcal{D}$, an arbitrary layer and $(\mathfrak{h},b)=(0,0)$: since $(0,0)^{\pm}=(0,0)$ (Section~\ref{sec:reflected-data}), its first assertion reads $Z^{\chi}_{\varepsilon}(0)\ge0$.

Let $M=\sup_{\mathfrak{h}\in C^2}\left|Z^{\chi}_{\varepsilon}(\mathfrak{h})\right|$, a finite number by Lemma~\ref{lem:gaussian-decay-continuity}(i). If $M=0$, then $Z^{\chi}_{\varepsilon}$ vanishes identically on $C^2$, and the inequality to be proved holds with both sides equal to $0$. Let now $M>0$. The set
\[
\mathcal{M}=\left\{\mathfrak{h}\in C^2:\left|Z^{\chi}_{\varepsilon}(\mathfrak{h})\right|=M\right\}
\]
is not empty by Lemma~\ref{lem:gaussian-decay-continuity}(iii). Since $\mathsf{n}(\cdot,0)$ takes values in the nonnegative integers, there is $\mathfrak{h}\in\mathcal{M}$ such that $\mathsf{n}(\mathfrak{h},0)\le\mathsf{n}(\mathfrak{h}',0)$ for every $\mathfrak{h}'\in\mathcal{M}$; fix such an $\mathfrak{h}$.

Suppose that $\mathsf{n}(\mathfrak{h},0)\ge1$, and let $x\in\Lambda$ and directions $\mu<\mu'$ be such that $\mathfrak{h}_{\mu\mu'}(x)\ne0$. The sets $\{\mu,\mu'\}$ and $\mathcal{D}$ are subsets of $\{1,2,3\}$ with at least two elements each, so they have an element in common; call it $\nu$, and let $i$ be the element of $\{\mu,\mu'\}$ different from $\nu$. Take $\nu$ as the time direction in Section~\ref{sec:slices-slice-operators} and write $x=(t,\mathbf{x})$ with $t=x_{\nu}$. The plaquette $(x,\mu\mu')$ contains the direction $\nu$, hence is time-like, and by the definition of the electric part of the layer data in Section~\ref{sec:slices-slice-operators},
\[
\mathfrak{h}_{\mu\mu'}(x)=\pm\,\mathfrak{h}^e_t(\mathbf{x},i).
\]
Thus the plaquette lies in the layer $\jmath=x_{\nu}$, and $\mathfrak{h}_{\mu\mu'}(x)$ is a nonzero component of the layer datum of this layer. The mirror attached to $\jmath$ fixes the layer $\jmath$ (Section~\ref{sec:reflected-data}); for this mirror, therefore, $\mathsf{n}_{\mathrm{fix}}(\mathfrak{h},0)\ge1$. Moreover $(\mathfrak{h},0)^{\pm}=(\mathfrak{h}^{\pm},0)$ (Section~\ref{sec:reflected-data}), and $\chi'=\chi$ because $\nu\in\mathcal{D}$. Applied to $(\mathfrak{h},0)$, \eqref{eq:gaussian-bounds-21} gives $Z^{\chi}_{\varepsilon}(\mathfrak{h}^{+})\ge0$, $Z^{\chi}_{\varepsilon}(\mathfrak{h}^{-})\ge0$ and
\[
M^2=\left|Z^{\chi}_{\varepsilon}(\mathfrak{h})\right|^2\le Z^{\chi}_{\varepsilon}(\mathfrak{h}^{+})\cdot Z^{\chi}_{\varepsilon}(\mathfrak{h}^{-})\le M\cdot M,
\]
the last inequality by the definition of $M$. As $0\le Z^{\chi}_{\varepsilon}(\mathfrak{h}^{\pm})\le M$ and $M>0$, it follows that $Z^{\chi}_{\varepsilon}(\mathfrak{h}^{+})=Z^{\chi}_{\varepsilon}(\mathfrak{h}^{-})=M$; in particular $\mathfrak{h}^{+}\in\mathcal{M}$ and $\mathfrak{h}^{-}\in\mathcal{M}$. On the other hand, \eqref{eq:gaussian-bounds-20} gives
\[
\mathsf{n}(\mathfrak{h}^{+},0)+\mathsf{n}(\mathfrak{h}^{-},0)=2\,\mathsf{n}(\mathfrak{h},0)-2\,\mathsf{n}_{\mathrm{fix}}(\mathfrak{h},0)\le2\,\mathsf{n}(\mathfrak{h},0)-2,
\]
so that $\mathsf{n}(\mathfrak{h}^{+},0)<\mathsf{n}(\mathfrak{h},0)$ or $\mathsf{n}(\mathfrak{h}^{-},0)<\mathsf{n}(\mathfrak{h},0)$: one of $\mathfrak{h}^{+}$, $\mathfrak{h}^{-}$ is an element of $\mathcal{M}$ with fewer nonzero components than $\mathfrak{h}$. This contradicts the choice of $\mathfrak{h}$.

Therefore $\mathsf{n}(\mathfrak{h},0)=0$, that is, $\mathfrak{h}=0$. Consequently $0\in\mathcal{M}$ and $M=\left|Z^{\chi}_{\varepsilon}(0)\right|=Z^{\chi}_{\varepsilon}(0)$, which is the asserted inequality $\left|Z^{\chi}_{\varepsilon}(\mathfrak{h}')\right|\le Z^{\chi}_{\varepsilon}(0)$ for all $\mathfrak{h}'\in C^2$.
\end{proof}

\begin{proof}[Proof of Theorem~\ref{thm:gaussian-domination-positive}(b)]
In the four-component representation, Lemma~\ref{lem:gaussian-bounds-1}(v) gives $\Delta^{\chi}\ge0$. The integrand of \eqref{eq:gaussian-bounds-2} is then nonnegative, and $Z^{\chi}_{\varepsilon}(\mathfrak{h},b)\ge0$ for all $\mathfrak{h}\in C^2$ and $b\in C^1$. If $q\ne0$, or if $\chi_{\mu}=-1$ for some direction $\mu$, then $Z^{\chi}_{\varepsilon}(0)>0$ by Corollary~\ref{cor:positivity}.
\end{proof}

\subsection{Exponential moments of the curvature}\label{sec:exponential-moments-curvature}

\begin{proof}[Proof of Theorem~\ref{thm:laplace-transform-curvature}]
Let $\chi$ be a character with $\chi_\mu=-1$ for at least two directions $\mu$, and let $\mathfrak{h}\in C^2$. For every $A\in C^1$ one has
\[
\left\|d_1A+\mathfrak{h}\right\|^2=\left\|d_1A\right\|^2+2\langle\mathfrak{h},d_1A\rangle+\|\mathfrak{h}\|^2 ,
\]
and inserting this identity into the definition \eqref{eq:gaussian-bounds-2} of $Z^\chi_\varepsilon(\mathfrak{h})=Z^\chi_\varepsilon(\mathfrak{h},0)$ gives
\begin{multline*}
Z^\chi_\varepsilon(\mathfrak{h})=\exp\left(-\frac12\|\mathfrak{h}\|^2\right)\\
\cdot\int_{C^1}\mathrm{d}A\,\exp\left(-\langle\mathfrak{h},d_1A\rangle\right)\cdot\exp\left(-\frac12\|d_1A\|^2-\frac{\varepsilon}{2}\|A\|^2\right)\cdot\Delta^\chi(A).
\end{multline*}
Since the Clifford triple is the four-component representation and $\chi_\mu=-1$ for some $\mu$, Corollary~\ref{cor:positivity} gives $Z^\chi_\varepsilon(0)>0$. Multiplying the last display by $\exp\left(\frac12\|\mathfrak{h}\|^2\right)\cdot Z^\chi_\varepsilon(0)^{-1}$ and comparing with the definition of $\mathsf E^\chi_\varepsilon$ in Section~\ref{sec:proca-parent-harmonic}, one obtains
\[
\mathsf E^\chi_\varepsilon\left[\exp\left(-\langle\mathfrak{h},d_1A\rangle\right)\right]=\exp\left(\frac12\|\mathfrak{h}\|^2\right)\cdot\frac{Z^\chi_\varepsilon(\mathfrak{h})}{Z^\chi_\varepsilon(0)}\le\exp\left(\frac12\|\mathfrak{h}\|^2\right),
\]
where the inequality holds because $\left|Z^\chi_\varepsilon(\mathfrak{h})\right|\le Z^\chi_\varepsilon(0)$ by Theorem~\ref{thm:gaussian-domination-positive}(a).

Now let $g\in C^2$, and put $\mathfrak{h}=-Pg$. Since $d_1A$ lies in $\operatorname{ran} d_1$ and $P$ is the orthogonal projection of $C^2$ onto $\operatorname{ran} d_1$, one has $\langle g,d_1A\rangle=\langle Pg,d_1A\rangle$ for every $A\in C^1$. Together with $\|Pg\|\le\|g\|$ this gives
\[
\mathsf E^\chi_\varepsilon\left[\exp\langle g,d_1A\rangle\right]=\mathsf E^\chi_\varepsilon\left[\exp\langle Pg,d_1A\rangle\right]\le\exp\left(\frac12\|Pg\|^2\right)\le\exp\left(\frac12\|g\|^2\right).
\]
\end{proof}

\begin{proof}[Proof of Corollary~\ref{cor:linear-functionals-curvature}]
Let $\mathsf E$, $\mathsf P$, $F$, $X_g$ and $\sigma_g$ be as in Corollary~\ref{cor:linear-functionals-curvature}, and let $g\in C^2$. For $s\in\mathbb R$ one has $X_{sg}=sX_g$ and $\sigma_{sg}=|s|\cdot\sigma_g$. Hence Theorem~\ref{thm:laplace-transform-curvature} if $\mathsf E=\mathsf E^\chi_\varepsilon$, and \eqref{eq:gaussian-bounds-5} if $\mathsf E=\overline{\mathsf E}$, each applied to $sg$ in place of $g$, give
\[
\phi(s):=\mathsf E\left[\exp(sX_g)\right]\le\exp\left(\frac12s^2\sigma_g^2\right)\qquad(s\in\mathbb R).
\]
As $\exp\left(|sX_g|\right)\le\exp(sX_g)+\exp(-sX_g)$, it follows that
\[
\mathsf E\left[\exp\left(|s|\cdot|X_g|\right)\right]\le\phi(s)+\phi(-s)<\infty\qquad(s\in\mathbb R).
\]
For $z\in\mathbb C$ one has $\left|\exp(zX_g)\right|=\exp\left((\operatorname{Re}z)\cdot X_g\right)\le\exp\left(|z|\cdot|X_g|\right)$, so that $\phi(z):=\mathsf E\left[\exp(zX_g)\right]$ is defined for every $z\in\mathbb C$; this extends $\phi$ to $\mathbb C$. Let $r>0$, let $k\ge0$ be an integer, and let $z,w\in\mathbb C$ with $|z|\le r$ and $0<|w|\le1$. Since $|X_g|^k\le k!\cdot\exp\left(|X_g|\right)$ and $\left|\exp(wX_g)-1\right|\le|w|\cdot|X_g|\cdot\exp\left(|X_g|\right)\le|w|\cdot\exp\left(2|X_g|\right)$, one has
\[
\left|\frac{X_g^k\cdot\exp\left((z+w)X_g\right)-X_g^k\cdot\exp(zX_g)}{w}\right|\le k!\cdot\exp\left((r+3)\cdot|X_g|\right),
\]
and the right side is integrable by the preceding bound. By dominated convergence and induction on $k$, the function $\phi$ is complex differentiable at every point of $\mathbb C$, hence entire, and it may be differentiated under $\mathsf E$:
\[
\phi^{(k)}(z)=\mathsf E\left[X_g^k\cdot\exp(zX_g)\right]\qquad\left(z\in\mathbb C,\ k\ge0\right).
\]

(i) For $s\in\mathbb R$ put $\psi(s)=\exp\left(\frac12s^2\sigma_g^2\right)-\phi(s)$. Then $\psi$ is infinitely differentiable on $\mathbb R$, one has $\psi\ge0$ by the bound on $\phi$, and $\psi(0)=1-\phi(0)=0$. Thus $\psi$ attains its minimum at $s=0$, so that $\psi'(0)=0$ and $\psi''(0)\ge0$. Since
\[
\psi'(0)=-\phi'(0)=-\mathsf EX_g,\qquad\psi''(0)=\sigma_g^2-\phi''(0)=\sigma_g^2-\mathsf EX_g^2,
\]
it follows that $\mathsf EX_g=0$ and $\mathsf EX_g^2\le\sigma_g^2$.

(ii) Let $z\in\mathbb C$. As $\mathsf P$ is a positive measure and $\left|\exp(zX_g)\right|=\exp\left((\operatorname{Re}z)\cdot X_g\right)$,
\[
\left|\mathsf E\left[\exp(zX_g)\right]\right|\le\mathsf E\left[\exp\left((\operatorname{Re}z)\cdot X_g\right)\right]=\phi(\operatorname{Re}z)\le\exp\left(\frac12\left(\operatorname{Re}z\right)^2\sigma_g^2\right).
\]

(iii) Let $\sigma_g>0$ and $y\ge0$, and put $s=y/\sigma_g^2$. As $s\ge0$, the event $\left\{X_g\ge y\right\}$ is contained in $\left\{\exp(sX_g)\ge\exp(sy)\right\}$, and Markov's inequality together with the bound on $\phi$ gives
\[
\mathsf P\left(X_g\ge y\right)\le\exp(-sy)\cdot\phi(s)\le\exp\left(-sy+\frac12s^2\sigma_g^2\right)=\exp\left(-\frac{y^2}{2\sigma_g^2}\right).
\]
Since $-X_g=X_{-g}$ and $\sigma_{-g}=\sigma_g$, the same bound holds for $\mathsf P\left(-X_g\ge y\right)$, and therefore
\[
\mathsf P\left(\left|X_g\right|\ge y\right)\le2\exp\left(-\frac{y^2}{2\sigma_g^2}\right)\qquad(y\ge0).
\]
For every real random variable $X$ and every real number $a>0$ one has
\[
\mathsf E|X|^a=\int_0^\infty\mathrm{d}y\,a\cdot y^{a-1}\cdot\mathsf P\left(|X|\ge y\right),
\]
by Tonelli's theorem applied to $|X|^a=\int_0^\infty\mathrm{d}y\,a\cdot y^{a-1}\cdot\mathbf 1_{\{|X|\ge y\}}$. Let $k\ge1$ be an integer. The substitution $u=y^2/(2\sigma_g^2)$ gives
\[
\int_0^\infty\mathrm{d}y\,y^{2k-1}\cdot\exp\left(-\frac{y^2}{2\sigma_g^2}\right)=\frac12\left(2\sigma_g^2\right)^k\int_0^\infty\mathrm{d}u\,u^{k-1}\cdot\exp(-u)=\frac12\left(2\sigma_g^2\right)^k\cdot(k-1)!\,.
\]
By the formula above with $a=2k$, by the bound for $\mathsf P\left(\left|X_g\right|\ge y\right)$, and by the last display,
\begin{align*}
\mathsf E\left|X_g\right|^{2k}&=\int_0^\infty\mathrm{d}y\,2k\cdot y^{2k-1}\cdot\mathsf P\left(\left|X_g\right|\ge y\right)\le4k\int_0^\infty\mathrm{d}y\,y^{2k-1}\cdot\exp\left(-\frac{y^2}{2\sigma_g^2}\right)\\
&=4k\cdot\frac12\left(2\sigma_g^2\right)^k\cdot(k-1)!=2\cdot k!\cdot\left(2\sigma_g^2\right)^k.
\end{align*}
\end{proof}

\subsection{Vanishing Proca mass}\label{sec:vanishing-proca-mass}

\begin{lemma}[Harmonic cochains; coercivity of $d_1$]\label{lem:harmonic-cochains-coercivity}
The space $V_{\mathrm{h}}$ of Section~\ref{sec:cochains-periodic-lattice} consists of the constant cochains, that is, of the cochains $c\in C^1$ such that, for each direction $\mu$, the value $c_\mu(x)$ does not depend on $x\in\Lambda$; such a cochain is identified with the vector $c=(c_1,c_2,c_3)\in\mathbb R^3$ of its values, so that $c_\mu(x)=c_\mu$. Let $\lambda_{V_{\mathrm{t}}}$ be the smallest nonzero eigenvalue of the operator $d_1^*d_1$ on $C^1$. Then $\lambda_{V_{\mathrm{t}}}>0$, and
\[
\|d_1A\|^2\ge\lambda_{V_{\mathrm{t}}}\cdot\|A\|^2\qquad\left(A\in V_{\mathrm{t}}\right).
\]
\end{lemma}

\begin{proof}
Since $V_{\mathrm{g}}$ is the range of $d_0$, its orthogonal complement in $C^1$ is the kernel of the adjoint $d_0^*\colon C^1\to C^0$, and therefore $V_{\mathrm{h}}=\ker d_1\cap\ker d_0^*$. Summation by parts gives
\[
\left(d_0^*A\right)(x)=\sum_{\mu=1}^3\left(A_\mu(x-\hat\mu)-A_\mu(x)\right)\qquad\left(A\in C^1,\ x\in\Lambda\right).
\]
For $k\in\tfrac{2\pi}{n}\left(\mathbb Z/n\mathbb Z\right)^3$ and $x\in\Lambda$ the number $\exp(ik\cdot x)$ is well defined. The discrete Fourier transform of a function $f\colon\Lambda\to\mathbb C$ is the function
\[
\hat f(k)=\sum_{x\in\Lambda}\exp(-ik\cdot x)\cdot f(x)\qquad\left(k\in\tfrac{2\pi}{n}\left(\mathbb Z/n\mathbb Z\right)^3\right).
\]
Summation of geometric series shows that $\sum_{x\in\Lambda}\exp(-ik\cdot x)$ equals $n^3$ for $k=0$ and $0$ for $k\ne0$, and that $\sum_k\exp(ik\cdot x)$, the sum being taken over $\tfrac{2\pi}{n}\left(\mathbb Z/n\mathbb Z\right)^3$, equals $n^3$ for $x=0$ and $0$ for $x\in\Lambda\setminus\{0\}$. Hence the transform of the constant function $1$ equals $n^3$ at $k=0$ and vanishes at every $k\ne0$, and the map $f\mapsto\hat f$ is linear and bijective, with inverse
\[
f(x)=n^{-3}\sum_{k\in\frac{2\pi}{n}\left(\mathbb Z/n\mathbb Z\right)^3}\exp(ik\cdot x)\cdot\hat f(k)\qquad\left(x\in\Lambda\right).
\]
Consequently $f$ is constant if and only if $\hat f(k)=0$ for every $k\ne0$. The transform of $x\mapsto f(x\pm\hat\mu)$ is $k\mapsto\exp(\pm ik_\mu)\cdot\hat f(k)$. For $A\in C^1$ let $\hat A_\mu$ denote the transform of the function $A_\mu$ on $\Lambda$, and put
\[
\zeta_\mu(k)=\exp(ik_\mu)-1\qquad\left(\mu\in\{1,2,3\}\right).
\]
By the definition of $d_1$ in Section~\ref{sec:cochains-periodic-lattice} and the formula for $d_0^*$, the transforms of $(d_1A)_{\mu\nu}$, for $\mu<\nu$, and of $d_0^*A$ are the functions
\[
k\mapsto\zeta_\mu(k)\cdot\hat A_\nu(k)-\zeta_\nu(k)\cdot\hat A_\mu(k),\qquad
k\mapsto\sum_{\mu=1}^3\overline{\zeta_\mu(k)}\cdot\hat A_\mu(k).
\]
Let $A\in V_{\mathrm{h}}$, fix $k$, and write $\zeta_\mu$ for $\zeta_\mu(k)$. Both transforms vanish at $k$, so that $\zeta_\mu\cdot\hat A_\nu(k)=\zeta_\nu\cdot\hat A_\mu(k)$ for all $\mu,\nu\in\{1,2,3\}$: for $\mu<\nu$ by the first formula, for $\mu=\nu$ trivially, and for $\mu>\nu$ by exchanging the roles of $\mu$ and $\nu$. Multiplication by $\overline{\zeta_\mu}$ and summation over $\mu$ give
\[
|\zeta|^2\cdot\hat A_\nu(k)=\zeta_\nu\cdot\sum_{\mu=1}^3\overline{\zeta_\mu}\cdot\hat A_\mu(k)=0,\qquad|\zeta|^2=\sum_{\mu=1}^3|\zeta_\mu|^2 .
\]
If $k\ne0$, then $k_\mu\ne0$ for some $\mu$, so that $\zeta_\mu\ne0$ and $|\zeta|^2>0$; hence $\hat A_\nu(k)=0$ for $\nu=1,2,3$. At $k=0$ all $\zeta_\mu$ vanish, and the two conditions impose nothing. It follows that a cochain $A\in C^1$ lies in $V_{\mathrm{h}}$ if and only if $\hat A_\nu(k)=0$ for all $k\ne0$ and all $\nu$, that is, if and only if $A$ is constant.

If $d_1^*d_1A=0$, then $\|d_1A\|^2=\langle A,d_1^*d_1A\rangle=0$; hence $\ker\left(d_1^*d_1\right)=\ker d_1$, and by Section~\ref{sec:cochains-periodic-lattice}
\[
V_{\mathrm{t}}=\left(\ker d_1\right)^\perp=\left(\ker d_1^*d_1\right)^\perp .
\]
The operator $d_1^*d_1$ is symmetric and positive semidefinite, and it is not zero ($d_1$ does not vanish on the cochain which equals $1$ on one link and $0$ on all others). Its nonzero eigenvalues are therefore positive, and $\lambda_{V_{\mathrm{t}}}>0$. By the spectral theorem, $V_{\mathrm{t}}$ is the orthogonal sum of the eigenspaces of $d_1^*d_1$ that belong to its nonzero eigenvalues, each of which is at least $\lambda_{V_{\mathrm{t}}}$; expanding $A\in V_{\mathrm{t}}$ in an orthonormal basis of eigenvectors one obtains $\|d_1A\|^2=\langle A,d_1^*d_1A\rangle\ge\lambda_{V_{\mathrm{t}}}\cdot\|A\|^2$.
\end{proof}

\begin{proof}[Proof of \eqref{eq:gaussian-bounds-4}]
Let the hypotheses of Theorem~\ref{thm:limit-vanishing-proca} hold, let $\chi$ be a character, and let $F\colon C^2\to\mathbb C$ be a Borel function with $|F(\omega)|\le c_1\cdot\exp\left(c_2\cdot\|\omega\|\right)$ for all $\omega\in C^2$, with constants $c_1,c_2$ as in Theorem~\ref{thm:limit-vanishing-proca}. Since $|\chi_\mu|=1$ for each $\mu$, there is $\vartheta\in\mathbb R^3$ with $\chi=\chi^{\mathrm{ap}}\exp(i\vartheta)$, in the notation of Lemma~\ref{lem:gaussian-bounds-1}(iv). Put
\[
\Delta_*=\left(9+|\boldsymbol{\mu}|\right)^{4n^3}.
\]
In the four-component representation $s=4$, and Lemma~\ref{lem:gaussian-bounds-1}(i),(v) gives $0\le\Delta^\chi\le\Delta_*$ and $0\le\Delta^{\mathrm{ap}}\le\Delta_*$ on $C^1$. For $\varepsilon>0$ put
\[
N_\varepsilon(F)=\int_{C^1}\mathrm dA\,F(d_1A)\cdot\exp\left(-\frac12\|d_1A\|^2-\frac{\varepsilon}{2}\|A\|^2\right)\cdot\Delta^\chi(A).
\]
As $c_2t-\frac12t^2\le\frac12c_2^2$ for real $t$, the integrand is bounded in absolute value by $c_1\cdot\Delta_*\cdot\exp\left(\frac12c_2^2-\frac{\varepsilon}{2}\|A\|^2\right)$, an integrable function of $A$; so the integral converges absolutely.

The decomposition $C^1=V_{\mathrm{g}}\oplus V_{\mathrm{h}}\oplus V_{\mathrm{t}}$ is orthogonal, and every $A\in C^1$ is written uniquely as $A=A_{V_{\mathrm{g}}}+c+A_{V_{\mathrm{t}}}$ with $A_{V_{\mathrm{g}}}\in V_{\mathrm{g}}$, $c\in V_{\mathrm{h}}$ and $A_{V_{\mathrm{t}}}\in V_{\mathrm{t}}$. Accordingly, Lebesgue measure on $C^1$ is the product of the Lebesgue measures $\sigma_{V_{\mathrm{g}}}$, $\sigma_{V_{\mathrm{h}}}$, $\sigma_{V_{\mathrm{t}}}$ of the Euclidean spaces $V_{\mathrm{g}}$, $V_{\mathrm{h}}$, $V_{\mathrm{t}}$, and
\[
\|A\|^2=\|A_{V_{\mathrm{g}}}\|^2+\|c\|^2+\|A_{V_{\mathrm{t}}}\|^2 .
\]
As $V_{\mathrm{g}}\oplus V_{\mathrm{h}}=\ker d_1$, one has $d_1A=d_1A_{V_{\mathrm{t}}}$. As $A_{V_{\mathrm{g}}}$ lies in the range of $d_0$, Lemma~\ref{lem:gaussian-bounds-1}(ii) and then Lemma~\ref{lem:gaussian-bounds-1}(iv) give
\[
\Delta^\chi(A)=\Delta^\chi\left(c+A_{V_{\mathrm{t}}}\right)=\Delta^{\mathrm{ap}}\left(A_{V_{\mathrm{t}}}+c+c^\vartheta\right).
\]
Lemma~\ref{lem:harmonic-cochains-coercivity} shows that the value of $c$ on a link $(x,\mu)$ depends only on the direction $\mu$; as in that lemma, $c$ also denotes the vector of $\mathbb R^3$ whose entries are its values in the three directions, and $|c|$ the Euclidean norm of this vector. Grouping the $3n^3$ links $(x,\mu)$ according to $x\in\Lambda$ gives $\|c\|^2=n^3\cdot|c|^2$, and the definition of $c^\theta$ in Section~\ref{sec:proca-parent-harmonic} gives $c^{qnc}=c$. The map $\theta\mapsto c^\theta$ is linear, so that $c+c^\vartheta=c^{qnc+\vartheta}$. For $A\in C^1$ let $f_A\colon\mathbb R^3\to[0,\infty)$ be the function
\[
f_A(\theta)=\Delta^{\mathrm{ap}}\left(A+c^\theta\right)\qquad\left(\theta\in\mathbb R^3\right);
\]
thus $\Delta^\chi(A)=f_{A_{V_{\mathrm{t}}}}(qnc+\vartheta)$. Let $\gamma_\varepsilon$ be the probability measure on $\mathbb R^3$ given by
\[
\gamma_\varepsilon(\mathrm dc)=\left(\frac{\varepsilon n^3}{2\pi}\right)^{\frac32}\exp\left(-\frac{\varepsilon}{2}\,n^3\cdot|c|^2\right)\mathrm d^3c,
\]
where $\mathrm d^3c$ is Lebesgue measure on $\mathbb R^3$. In the coordinates $c\in\mathbb R^3$ one has $\mathrm d\sigma_{V_{\mathrm{h}}}(c)=\left(n^3\right)^{\frac32}\mathrm d^3c$, because $\|c\|^2=n^3\cdot|c|^2$; hence
\[
\exp\left(-\frac{\varepsilon}{2}\|c\|^2\right)\mathrm d\sigma_{V_{\mathrm{h}}}(c)=\left(\frac{2\pi}{\varepsilon}\right)^{\frac32}\gamma_\varepsilon(\mathrm dc).
\]
Moreover, by a Gaussian integral,
\[
\int_{V_{\mathrm{g}}}\mathrm d\sigma_{V_{\mathrm{g}}}(A_{V_{\mathrm{g}}})\,\exp\left(-\frac{\varepsilon}{2}\|A_{V_{\mathrm{g}}}\|^2\right)=\left(\frac{2\pi}{\varepsilon}\right)^{\frac{\dim V_{\mathrm{g}}}{2}}.
\]
Fubini's theorem, which applies because $N_\varepsilon(F)$ converges absolutely, therefore gives
\begin{multline}\label{eq:gaussian-bounds-25}
N_\varepsilon(F)=\left(\frac{2\pi}{\varepsilon}\right)^{\frac{\dim V_{\mathrm{g}}+3}{2}}\int_{V_{\mathrm{t}}}\mathrm d\sigma_{V_{\mathrm{t}}}(A_{V_{\mathrm{t}}})\,F(d_1A_{V_{\mathrm{t}}})\cdot\exp\left(-\frac12\|d_1A_{V_{\mathrm{t}}}\|^2-\frac{\varepsilon}{2}\|A_{V_{\mathrm{t}}}\|^2\right)\\
\cdot\left[\int_{\mathbb R^3}\gamma_\varepsilon(\mathrm dc)\,f_{A_{V_{\mathrm{t}}}}(qnc+\vartheta)\right].
\end{multline}

By Lemma~\ref{lem:gaussian-bounds-1}(iv), for each $A\in C^1$ there are complex numbers $\hat f_{A,k}$, $k\in\mathbb Z^3$, only finitely many of which differ from $0$, such that
\[
f_A(\theta)=\sum_{k\in\mathbb Z^3}\hat f_{A,k}\cdot\exp(ik\cdot\theta)\qquad\left(\theta\in\mathbb R^3\right).
\]
By the orthogonality of the functions $\theta\mapsto\exp(ik\cdot\theta)$ on $[-\pi,\pi]^3$,
\[
\hat f_{A,k}=\int_{[-\pi,\pi]^3}\frac{\mathrm d\theta}{(2\pi)^3}\,\exp(-ik\cdot\theta)\cdot f_A(\theta)\qquad\left(k\in\mathbb Z^3\right).
\]
Hence each Fourier coefficient is bounded by the supremum of the function, $\left|\hat f_{A,k}\right|\le\sup|f_A|\le\Delta_*$, and
\[
\hat f_{A,0}=\int_{[-\pi,\pi]^3}\frac{\mathrm d\theta}{(2\pi)^3}\,\Delta^{\mathrm{ap}}\left(A+c^\theta\right).
\]
The measure $\gamma_\varepsilon$ is the centred Gaussian measure on $\mathbb R^3$ whose covariance is $\left(\varepsilon n^3\right)^{-1}$ times the identity; hence, for $k\in\mathbb Z^3$,
\[
\int_{\mathbb R^3}\gamma_\varepsilon(\mathrm dc)\,\exp\left(ik\cdot(qnc)\right)=\exp\left(-\frac{q^2n^2|k|^2}{2\varepsilon n^3}\right)=\exp\left(-\frac{q^2|k|^2}{2\varepsilon n}\right).
\]
Integration of the finite sum term by term yields
\[
\int_{\mathbb R^3}\gamma_\varepsilon(\mathrm dc)\,f_A(qnc+\vartheta)=\sum_{k\in\mathbb Z^3}\hat f_{A,k}\cdot\exp(ik\cdot\vartheta)\cdot\exp\left(-\frac{q^2|k|^2}{2\varepsilon n}\right),
\]
and the term with $k=0$ is $\hat f_{A,0}$. Consequently, for all $A\in V_{\mathrm{t}}$ and all $\vartheta\in\mathbb R^3$,
\[
\left|\int_{\mathbb R^3}\gamma_\varepsilon(\mathrm dc)\,f_A(qnc+\vartheta)-\hat f_{A,0}\right|\le\Delta_*\cdot\delta_\varepsilon,\qquad
\delta_\varepsilon=\sum_{k\in\mathbb Z^3\setminus\{0\}}\exp\left(-\frac{q^2|k|^2}{2\varepsilon n}\right).
\]
Since $q\ne0$, one has $\delta_\varepsilon\to0$ as $\varepsilon\downarrow0$: each term decreases to $0$, and for $\varepsilon\le1$ the terms are bounded by those of the convergent series $\delta_1$. Thus, as $\varepsilon\downarrow0$, the bracket in \eqref{eq:gaussian-bounds-25} converges to $\hat f_{A_{V_{\mathrm{t}}},0}$, uniformly in $A_{V_{\mathrm{t}}}\in V_{\mathrm{t}}$ and $\vartheta\in\mathbb R^3$.

Since $\exp\left(-\frac{\varepsilon}{2}\|A_{V_{\mathrm{t}}}\|^2\right)\le1$, since $\gamma_\varepsilon$ is a probability measure and since $0\le f_A\le\Delta_*$, the integrand of the integral over $V_{\mathrm{t}}$ in \eqref{eq:gaussian-bounds-25} is bounded in absolute value by the function
\[
A\mapsto c_1\cdot\Delta_*\cdot\exp\left(c_2\cdot\|d_1A\|-\frac12\|d_1A\|^2\right)\qquad\left(A\in V_{\mathrm{t}}\right),
\]
which does not depend on $\varepsilon$. This function is integrable on $V_{\mathrm{t}}$: by the inequality $c_2t-\frac12t^2\le c_2^2-\frac14t^2$ for real $t$ and by Lemma~\ref{lem:harmonic-cochains-coercivity}, it is at most $c_1\cdot\Delta_*\cdot\exp\left(c_2^2-\frac14\lambda_{V_{\mathrm{t}}}\cdot\|A\|^2\right)$. For all $\theta\in[-\pi,\pi]^3$ and $A\in V_{\mathrm{t}}$, the growth bound on $F$ at $\omega=d_1A$ and the inequalities $0\le\Delta^{\mathrm{ap}}\le\Delta_*$ give
\[
\left|F(d_1A)\cdot\exp\left(-\frac12\|d_1A\|^2\right)\cdot\Delta^{\mathrm{ap}}\left(A+c^\theta\right)\right|\le c_1\cdot\Delta_*\cdot\exp\left(c_2\cdot\|d_1A\|-\frac12\|d_1A\|^2\right).
\]
On the left stands the modulus of the integrand of $\overline Z[F]$ in \eqref{eq:gaussian-bounds-3}, and on the right the integrable majorant obtained above for \eqref{eq:gaussian-bounds-25}, in which $\theta$ does not occur. As $(2\pi)^{-3}\,\mathrm d\theta$ has total mass one on $[-\pi,\pi]^3$, integrating this inequality with respect to $(2\pi)^{-3}\,\mathrm d\theta$ and $\sigma_{V_{\mathrm{t}}}$ bounds the double integral of the left side by the integral of the right side over $V_{\mathrm{t}}$, which is finite. Thus $\overline Z[F]$ is an absolutely convergent double integral. Since the constant function $1$ obeys the growth bound of Theorem~\ref{thm:limit-vanishing-proca}, the case $F=1$ gives $\overline Z[1]<\infty$. By the dominated convergence theorem applied to \eqref{eq:gaussian-bounds-25}, and by Fubini's theorem, by which the integration over $\theta$ in \eqref{eq:gaussian-bounds-3} may be carried out first,
\[
\lim_{\varepsilon\downarrow0}\left(\frac{2\pi}{\varepsilon}\right)^{-\frac{\dim V_{\mathrm{g}}+3}{2}}\cdot N_\varepsilon(F)=\int_{V_{\mathrm{t}}}\mathrm d\sigma_{V_{\mathrm{t}}}(A)\,F(d_1A)\cdot\exp\left(-\frac12\|d_1A\|^2\right)\cdot\hat f_{A,0}=\overline Z[F].
\]

For $F=1$ this limit is positive. Indeed, the integrand $(\theta,A)\mapsto\exp\left(-\frac12\|d_1A\|^2\right)\cdot\Delta^{\mathrm{ap}}\left(A+c^\theta\right)$ of $\overline Z[1]$ is continuous and nonnegative on $[-\pi,\pi]^3\times V_{\mathrm{t}}$ by Lemma~\ref{lem:gaussian-bounds-1}(i),(v), and at $(\theta,A)=(0,0)$ it equals $\Delta^{\mathrm{ap}}(0)$, which is positive by Corollary~\ref{cor:positivity}. Hence it is positive on an open neighbourhood of $(0,0)$ in $[-\pi,\pi]^3\times V_{\mathrm{t}}$, a set of positive measure, and $\overline Z[1]>0$.

Finally, $N_\varepsilon(1)=Z^\chi_\varepsilon(0)>0$ by Corollary~\ref{cor:positivity}, since $q\ne0$, and $\mathsf E^\chi_\varepsilon\left[F(d_1A)\right]=N_\varepsilon(F)/N_\varepsilon(1)$ by the definition of the Proca parent in Section~\ref{sec:proca-parent-harmonic}. Dividing the limit for $F$ by the limit for $1$, one obtains
\[
\lim_{\varepsilon\downarrow0}\mathsf E^\chi_\varepsilon\left[F(d_1A)\right]=\frac{\overline Z[F]}{\overline Z[1]}=\overline{\mathsf E}[F].
\]
\end{proof}

\begin{proof}[Proof of \eqref{eq:gaussian-bounds-5}]
Let $g\in C^2$. The function $F(\omega)=\exp\langle g,\omega\rangle$ on $C^2$ is continuous, and $|F(\omega)|\le\exp\left(\|g\|\cdot\|\omega\|\right)$ by the Cauchy--Schwarz inequality; so \eqref{eq:gaussian-bounds-4} applies to $F$ and to the character $\chi^{\mathrm{ap}}$. All three components of $\chi^{\mathrm{ap}}$ equal $-1$, and Theorem~\ref{thm:laplace-transform-curvature} gives, for every $\varepsilon>0$,
\[
\mathsf E^{\mathrm{ap}}_\varepsilon\left[\exp\langle g,d_1A\rangle\right]\le\exp\left(\frac12\|Pg\|^2\right)\le\exp\left(\frac12\|g\|^2\right).
\]
By \eqref{eq:gaussian-bounds-4}, the left side converges to $\overline{\mathsf E}\left[\exp\langle g,\cdot\rangle\right]$ as $\varepsilon\downarrow0$, and \eqref{eq:gaussian-bounds-5} follows.
\end{proof}

\begin{proof}[Proof of the formula for $\mathsf E^{\mathrm{free}}$ in Theorem~\ref{thm:limit-vanishing-proca}]
By the proof of Lemma~\ref{lem:harmonic-cochains-coercivity}, $V_{\mathrm{t}}$ is the orthogonal complement of the kernel of the symmetric operator $d_1^*d_1$, hence equal to its range. So $d_1^*d_1$ maps $V_{\mathrm{t}}$ onto itself, and its restriction $d_1^*d_1|_{V_{\mathrm{t}}}$ is an invertible operator on $V_{\mathrm{t}}$. Let $g\in C^2$. The cochain $d_1^*g$ is orthogonal to $\ker d_1$, so that it lies in $V_{\mathrm{t}}$, and $\langle g,d_1A\rangle=\langle d_1^*g,A\rangle$ for $A\in C^1$. Put
\[
u=\left(d_1^*d_1|_{V_{\mathrm{t}}}\right)^{-1}d_1^*g\in V_{\mathrm{t}} .
\]
For $A\in V_{\mathrm{t}}$, expansion of $\left\|d_1(A-u)\right\|^2=\left\langle A-u,d_1^*d_1(A-u)\right\rangle$ together with $d_1^*d_1u=d_1^*g$ gives
\[
\langle g,d_1A\rangle-\frac12\|d_1A\|^2=\frac12\langle d_1^*g,u\rangle-\frac12\left\|d_1(A-u)\right\|^2 .
\]
The translation $A\mapsto A+u$ of $V_{\mathrm{t}}$ preserves $\sigma_{V_{\mathrm{t}}}$. It follows that
\begin{align*}
\int_{V_{\mathrm{t}}}\mathrm d\sigma_{V_{\mathrm{t}}}(A)\,\exp\left(\langle g,d_1A\rangle-\frac12\|d_1A\|^2\right)
&=\exp\left(\frac12\left\langle d_1^*g,\left(d_1^*d_1|_{V_{\mathrm{t}}}\right)^{-1}d_1^*g\right\rangle\right)\\
&\qquad\cdot\int_{V_{\mathrm{t}}}\mathrm d\sigma_{V_{\mathrm{t}}}(A)\,\exp\left(-\frac12\|d_1A\|^2\right),
\end{align*}
where both Gaussian integrals are finite by Lemma~\ref{lem:harmonic-cochains-coercivity}. Moreover $d_1\left(d_1^*d_1|_{V_{\mathrm{t}}}\right)^{-1}d_1^*=P$. Indeed, both sides vanish on $\ker d_1^*$, which is the orthogonal complement of $\operatorname{ran} d_1$ in $C^2$. Every $\omega\in\operatorname{ran} d_1$ is of the form $\omega=d_1A$ with $A\in V_{\mathrm{t}}$, because $d_1$ vanishes on $V_{\mathrm{g}}\oplus V_{\mathrm{h}}$; then $d_1^*\omega=d_1^*d_1A$, and the left side maps $\omega$ to $d_1A=\omega=P\omega$. Consequently
\[
\left\langle d_1^*g,\left(d_1^*d_1|_{V_{\mathrm{t}}}\right)^{-1}d_1^*g\right\rangle=\langle g,Pg\rangle=\|Pg\|^2,
\]
and division by $\int_{V_{\mathrm{t}}}\mathrm d\sigma_{V_{\mathrm{t}}}(A)\,\exp\left(-\frac12\|d_1A\|^2\right)$ gives $\mathsf E^{\mathrm{free}}\left[\exp\langle g,\cdot\rangle\right]=\exp\left(\frac12\|Pg\|^2\right)$.
\end{proof}

The inequality \eqref{eq:gaussian-bounds-5} can be written as an inequality between partition functions with an external two-cochain, of the type \eqref{eq:gaussian-bounds-2}. For $\mathfrak{h}\in C^2$ put
\[
\overline Z(\mathfrak{h})=\exp\left(-\frac12\|\mathfrak{h}\|^2\right)\cdot\overline Z\left[\exp\left(-\langle\mathfrak{h},\cdot\rangle\right)\right].
\]
Expanding $\|d_1A+\mathfrak{h}\|^2=\|d_1A\|^2+2\langle\mathfrak{h},d_1A\rangle+\|\mathfrak{h}\|^2$ in \eqref{eq:gaussian-bounds-3}, one finds
\[
\overline Z(\mathfrak{h})=\int_{[-\pi,\pi]^3}\frac{\mathrm d\theta}{(2\pi)^3}\int_{V_{\mathrm{t}}}\mathrm d\sigma_{V_{\mathrm{t}}}(A)\,\exp\left(-\frac12\left\|d_1A+\mathfrak{h}\right\|^2\right)\cdot\Delta^{\mathrm{ap}}\left(A+c^\theta\right).
\]
Since $\overline Z(0)=\overline Z[1]$ and $\|\mathfrak{h}\|^2=\|P\mathfrak{h}\|^2+\|(1-P)\mathfrak{h}\|^2$, the inequality \eqref{eq:gaussian-bounds-5} for $g=-\mathfrak{h}$ reads
\[
\overline Z(\mathfrak{h})\le\exp\left(-\frac12\left\|(1-P)\mathfrak{h}\right\|^2\right)\cdot\overline Z(0)\le\overline Z(0)\qquad\left(\mathfrak{h}\in C^2\right).
\]

\subsection{Twisted laws and the zero-harmonic law}\label{sec:twisted-laws-zero}

Throughout this subsection the Clifford triple is the four-component representation, and $n\ge2$, $q\ne0$, $\boldsymbol{\mu}>-1$.

\begin{definition}[Twisted partition functions and laws]
Let $\theta\in\mathbb R^3$, and let $F\colon C^2\to\mathbb C$ be a Borel function such that either $F$ takes its values in $[0,\infty)$, or there are constants $c_1,c_2$ with
\[
|F(\omega)|\le c_1\cdot\exp\left(c_2\cdot\|\omega\|\right)\qquad\left(\omega\in C^2\right).
\]
Put
\[
Z_\theta[F]=\int_{V_{\mathrm{t}}}\mathrm{d}\sigma_{V_{\mathrm{t}}}(A)\,F(d_1A)\cdot\exp\left(-\frac12\|d_1A\|^2\right)\cdot\Delta^{\mathrm{ap}}\left(A+c^\theta\right),
\]
an element of $[0,\infty]$ if $F\ge0$, and, provided that $Z_\theta[1]>0$,
\[
\mathsf E_\theta[F]=\frac{Z_\theta[F]}{Z_\theta[1]}.
\]
\end{definition}

By Lemma~\ref{lem:gaussian-bounds-1}(i),(v), and since $s=4$,
\[
0\le\Delta^{\mathrm{ap}}(A)\le\Delta_*:=\left(9+|\boldsymbol{\mu}|\right)^{4n^3}\qquad\left(A\in C^1\right).
\]
Let $F$ obey the exponential bound of the definition, with constants $c_1,c_2$. By this bound and the last display, the modulus of the integrand of $Z_\theta[F]$ at $A\in V_{\mathrm{t}}$ is at most the left side of
\[
c_1\cdot\Delta_*\cdot\exp\left(c_2\cdot\|d_1A\|-\frac12\|d_1A\|^2\right)\le c_1\cdot\Delta_*\cdot\exp\left(c_2^2-\frac14\lambda_{V_{\mathrm{t}}}\cdot\|A\|^2\right)\qquad\left(A\in V_{\mathrm{t}}\right),
\]
and this inequality holds by $c_2t-\frac12t^2\le c_2^2-\frac14t^2$ for real $t$ and by Lemma~\ref{lem:harmonic-cochains-coercivity}, where $\lambda_{V_{\mathrm{t}}}>0$ is the number defined there. The right side does not depend on $\theta$, and it is integrable on $V_{\mathrm{t}}$; hence $Z_\theta[F]$ converges absolutely. In particular $0\le Z_\theta[1]<\infty$. If $Z_\theta[1]>0$, then $\mathsf E_\theta[F]$ is a complex number when $F$ obeys the exponential bound and an element of $[0,\infty]$ when $F\ge0$, and $\mathsf E_\theta$ is the expectation of a probability measure on $C^2$, namely of the image under $d_1$ of the probability measure
\[
Z_\theta[1]^{-1}\cdot\exp\left(-\frac12\|d_1A\|^2\right)\cdot\Delta^{\mathrm{ap}}\left(A+c^\theta\right)\,\mathrm{d}\sigma_{V_{\mathrm{t}}}(A)
\]
on $V_{\mathrm{t}}$. This probability measure on $C^2$ is called the \emph{law with twist}~$\theta$.

For a Borel function $F\colon C^2\to[0,\infty)$ the double integral in \eqref{eq:gaussian-bounds-3} has a nonnegative integrand; it is denoted by $\overline Z[F]\in[0,\infty]$ also when it diverges. For every $F$ as in the definition, the inner integral in \eqref{eq:gaussian-bounds-3} is $Z_\theta[F]$, and $\theta\mapsto Z_\theta[F]$ is a Borel function on $[-\pi,\pi]^3$ (by Tonelli's theorem if $F\ge0$, and otherwise by Fubini's theorem, the double integral converging absolutely by Theorem~\ref{thm:limit-vanishing-proca}). Therefore
\[
\overline Z[F]=\int_{[-\pi,\pi]^3}\frac{\mathrm{d}\theta}{(2\pi)^3}\,Z_\theta[F].
\]

One has $Z_0[1]>0$. Indeed, $c^0=0$, so that the integrand of $Z_0[1]$ is the function $A\mapsto\exp\left(-\frac12\|d_1A\|^2\right)\cdot\Delta^{\mathrm{ap}}(A)$ on $V_{\mathrm{t}}$. This function is continuous and nonnegative by Lemma~\ref{lem:gaussian-bounds-1}(i),(v), and its value at $A=0$ is $\Delta^{\mathrm{ap}}(0)$, which is positive by Corollary~\ref{cor:positivity}. Hence it is positive on a neighbourhood of $0$ in $V_{\mathrm{t}}$, a set of positive measure, and $Z_0[1]>0$. Thus the law with twist $0$ is defined.

\begin{remark}[The zero-harmonic law]
By Lemma~\ref{lem:harmonic-cochains-coercivity} the space $V_{\mathrm{h}}$ consists of the constant cochains; in particular $c^\theta\in V_{\mathrm{h}}$ for every $\theta\in\mathbb R^3$. As the decomposition $C^1=V_{\mathrm{g}}\oplus V_{\mathrm{h}}\oplus V_{\mathrm{t}}$ is orthogonal, for $A\in V_{\mathrm{t}}$ the cochain $A+c^\theta$ has the component $0$ in $V_{\mathrm{g}}$, the component $c^\theta$ in $V_{\mathrm{h}}$ and the component $A$ in $V_{\mathrm{t}}$. In \eqref{eq:gaussian-bounds-3} the determinant $\Delta^{\mathrm{ap}}\left(A+c^\theta\right)$ is averaged over $\theta\in[-\pi,\pi]^3$ with respect to $(2\pi)^{-3}\,\mathrm{d}\theta$, whereas in $Z_\theta$ the constant component $c^\theta$ is held fixed. For $\theta=0$ this component vanishes, and
\[
Z_0[F]=\int_{V_{\mathrm{t}}}\mathrm{d}\sigma_{V_{\mathrm{t}}}(A)\,F(d_1A)\cdot\exp\left(-\frac12\|d_1A\|^2\right)\cdot\Delta^{\mathrm{ap}}(A).
\]
Consequently $Z_0$, and with it the law with twist $0$, is obtained from the inner integral in \eqref{eq:gaussian-bounds-3} by replacing the component $c^\theta$ of the one-cochain $A+c^\theta$ in $V_{\mathrm{h}}$ by $0$. This law is called the \emph{zero-harmonic law}, and its expectation is $\mathsf E_0$.
\end{remark}

\begin{lemma}[Two forms of a Gaussian bound for the zero-harmonic law]
The following two assertions are equivalent.
\begin{enumerate}
\item[(i)] For every $g\in C^2$,
\[
\mathsf E_0\left[\exp\langle g,\cdot\rangle\right]\le\exp\left(\frac12\|Pg\|^2\right).
\]
\item[(ii)] For every $\mathfrak{h}\in C^2$,
\begin{multline*}
\int_{V_{\mathrm{t}}}\mathrm{d}\sigma_{V_{\mathrm{t}}}(A)\,\exp\left(-\frac12\left\|d_1A+\mathfrak{h}\right\|^2\right)\cdot\Delta^{\mathrm{ap}}(A)\\
\le\int_{V_{\mathrm{t}}}\mathrm{d}\sigma_{V_{\mathrm{t}}}(A)\,\exp\left(-\frac12\|d_1A\|^2\right)\cdot\Delta^{\mathrm{ap}}(A).
\end{multline*}
\end{enumerate}
\end{lemma}

\begin{proof}
For $\mathfrak{h}\in C^2$ let $Z_0(\mathfrak{h})$ denote the left side of the inequality in (ii). Inserting
\[
\left\|d_1A+\mathfrak{h}\right\|^2=\|d_1A\|^2+2\langle\mathfrak{h},d_1A\rangle+\|\mathfrak{h}\|^2
\]
and comparing with the definition of $Z_0$, in which $c^0=0$, one finds
\[
Z_0(\mathfrak{h})=\exp\left(-\frac12\|\mathfrak{h}\|^2\right)\cdot Z_0\left[\exp\left(-\langle\mathfrak{h},\cdot\rangle\right)\right].
\]
By the Cauchy--Schwarz inequality the function $\omega\mapsto\exp\left(-\langle\mathfrak{h},\omega\rangle\right)$ obeys the exponential bound with $c_1=1$ and $c_2=\|\mathfrak{h}\|$; therefore $Z_0(\mathfrak{h})$ is a finite nonnegative number. The right side of the inequality in (ii) is $Z_0(0)=Z_0[1]$, which is finite and positive.

First, let $g\in C^2$. For every $A\in C^1$ the cochain $d_1A$ lies in $\operatorname{ran} d_1$, and $P$ is the orthogonal projection of $C^2$ onto $\operatorname{ran} d_1$; hence $\langle g,d_1A\rangle=\langle Pg,d_1A\rangle$, and completing the square gives
\[
\langle g,d_1A\rangle-\frac12\|d_1A\|^2=\frac12\|Pg\|^2-\frac12\left\|d_1A-Pg\right\|^2 .
\]
Taking the exponential of both sides, multiplying by $\Delta^{\mathrm{ap}}(A)$ and integrating over $V_{\mathrm{t}}$ with respect to $\sigma_{V_{\mathrm{t}}}$, one obtains $Z_0\left[\exp\langle g,\cdot\rangle\right]=\exp\left(\frac12\|Pg\|^2\right)\cdot Z_0(-Pg)$, and division by $Z_0[1]=Z_0(0)$ gives
\[
\mathsf E_0\left[\exp\langle g,\cdot\rangle\right]=\exp\left(\frac12\|Pg\|^2\right)\cdot\frac{Z_0(-Pg)}{Z_0(0)}.
\]
Consequently the inequality in (i) holds for $g$ if and only if $Z_0(-Pg)\le Z_0(0)$.

Secondly, let $\mathfrak{h}\in C^2$ and $A\in V_{\mathrm{t}}$. The cochain $d_1A+P\mathfrak{h}$ lies in $\operatorname{ran} d_1$, whereas $(1-P)\mathfrak{h}$ is orthogonal to $\operatorname{ran} d_1$. Hence
\[
\left\|d_1A+\mathfrak{h}\right\|^2=\left\|d_1A+P\mathfrak{h}\right\|^2+\left\|(1-P)\mathfrak{h}\right\|^2,
\]
and inserting this identity into the definition of $Z_0(\mathfrak{h})$ gives
\[
Z_0(\mathfrak{h})=\exp\left(-\frac12\left\|(1-P)\mathfrak{h}\right\|^2\right)\cdot Z_0(P\mathfrak{h})\le Z_0(P\mathfrak{h}).
\]

Now assume (i), and let $\mathfrak{h}\in C^2$. For $g=-\mathfrak{h}$ one has $-Pg=P\mathfrak{h}$, and the first step gives $Z_0(P\mathfrak{h})\le Z_0(0)$; together with the second step, $Z_0(\mathfrak{h})\le Z_0(P\mathfrak{h})\le Z_0(0)$, which is the inequality of (ii) for $\mathfrak{h}$. Conversely, assume (ii), and let $g\in C^2$. The inequality of (ii) for $\mathfrak{h}=-Pg$ states that $Z_0(-Pg)\le Z_0(0)$, and the first step yields the inequality of (i) for $g$.
\end{proof}

\section{Continuum and infinite-volume limit points}\label{sec:continuum-infinite-volume}

The members of an admissible family of lattice laws (Definition~\ref{def:admissible-families}) obey the bounds at a fixed cutoff of Theorem~\ref{thm:estimates-fixed-cutoff}, the family is tight (Theorem~\ref{thm:second-moments-tightness}), and all its limit points have the properties stated in Theorem~\ref{thm:continuum-infinite-1}, Theorem~\ref{thm:moments-limit-point} and Corollary~\ref{cor:transfer-operators-coordinate}; the averaged laws form an admissible family (Corollary~\ref{cor:averaged-laws-admissible}) and are compared with the pinned laws in Theorem~\ref{thm:insensitivity-twist}; these results contain Estimate~10, a part of Estimate~2, and Estimate~4 apart from its assertion that the iterated weak limits are limit points, which is proved in Remark~\ref{rem:repeated-limits}. Section~\ref{sec:lattice-laws-admissible} sets out again the tori, the cochains, the test map and the spaces $H_s$ of Section~\ref{sec:setup}; for $q>0$ and $\boldsymbol{\mu}>-\frac{1}{2}$ the averaged law and the pinned law defined there are the harmonic-averaged law and the zero-harmonic law of Section~\ref{sec:lattice-laws}, written for the curvature in physical units.

\subsection{Lattice laws and admissible families}\label{sec:lattice-laws-admissible}

\subsubsection{Tori, cells and cochains}\label{sec:tori-cells-cochains}

Let $L\ge3$ be an odd integer and let $0<\bar{m}\le1$. Put
\[
h=L^{\left\lfloor\log_L{\bar{m}}^{-1}\right\rfloor};
\]
as $\bar{m}\le1$, the exponent is a non-negative integer, and $h\ge1$. For integers $N,T\ge0$ let
\[
\eta=L^{-N},\qquad R=h\cdot L^T,\qquad n=\frac{R}{\eta},\qquad m=\frac{n-1}{2}.
\]
Then $R\ge1$, and $n=L^{\left\lfloor\log_L{\bar{m}}^{-1}\right\rfloor+T+N}$ is a power of the odd integer $L$; hence $n$ is odd and $m$ is an integer. The lattice is
\[
\Lambda=\Lambda_{\eta,R}=\eta\mathbb Z^3/R\mathbb Z^3 ;
\]
every point of $\Lambda$ has exactly one representative in $\eta\cdot\{-m,\dots,m\}^3$, and the coordinates of a point of $\Lambda$ are those of its representative. Further, $\mathbb T_R=\mathbb R^3/R\mathbb Z^3$ and $Q_R=\left[-\frac{R}{2},\frac{R}{2}\right)^3$.

A pair $(N,T)$ is called a \emph{cutoff} if $n\ge2$. Since $n$ is a power of $L$, the case $n=1$ occurs only for $h=1$ and $N=T=0$; for a cutoff one has $n\ge L\ge3$ and $m\ge1$. From now on only cutoffs are considered.

Let $e_1,e_2,e_3$ be the standard basis of $\mathbb R^3$. For $x\in\Lambda$ and $I=\{i_1,\dots,i_k\}\subset\{1,2,3\}$ with $i_1<\dots<i_k$, so that $\#|I|=k$, the \emph{$k$-cell} $\sigma=(x;I)$ is the subset
\[
x+\sum_{j\in I}[0,\eta]\cdot e_j
\]
of $\mathbb T_R$, oriented by $e_{i_1}\wedge\dots\wedge e_{i_k}$. Its \emph{centre} and its \emph{averaging cube} are
\[
c(\sigma)=x+\frac{\eta}{2}\sum_{j\in I}e_j,\qquad Q(\sigma)=c(\sigma)+\left[-\frac{\eta}{2},\frac{\eta}{2}\right)^3 .
\]
For each fixed $I$ the cubes $Q(x;I)$, $x\in\Lambda$, form a partition of $\mathbb T_R$. One-cells are also called links, and two-cells plaquettes; the plaquette $(x;\{\mu,\nu\})$ with $\mu<\nu$ is written $(x;\mu\nu)$.

For $k=0,1,2,3$ let $C^k=C^k(\Lambda)$ be the real vector space of all functions from the set of $k$-cells to $\mathbb R$; its elements are the $k$-cochains. It carries the pairing and the norm
\[
\langle u,v\rangle_\eta=\eta^3\sum_\sigma u(\sigma)\cdot v(\sigma),\qquad|u|_\eta=\langle u,u\rangle_\eta^{\frac12},
\]
the sum being taken over all $k$-cells $\sigma$; for two-cochains it is the sum over $x\in\Lambda$ and over the pairs $\mu<\nu$ of the terms at the plaquettes $(x;\mu\nu)$. Without the index $\eta$,
\[
\langle u,v\rangle=\sum_\sigma u(\sigma)\cdot v(\sigma),\qquad\|u\|=\langle u,u\rangle^{\frac12}
\]
denote the counting pairing and the counting norm. The elements of the complexification $C^k+iC^k$ are called complex $k$-cochains. Both pairings are extended to complex cochains by the same formulas, without complex conjugation, and so become complex-bilinear. The norm $|\cdot|_\eta$ is extended to complex cochains by
\[
|u+iv|_\eta^2=|u|_\eta^2+|v|_\eta^2\qquad(u,v\in C^k),
\]
so that $|u+iv|_\eta^2=\langle u-iv,u+iv\rangle_\eta$. For $k=0,1,2$ let $d\colon C^k\to C^{k+1}$ be the coboundary of the lattice with unit spacing; on $C^0$ and $C^1$ it is given by the formulas for $d_0$ and $d_1$ in Section~\ref{sec:cochains-periodic-lattice}. Moreover $d_\eta=\eta^{-1}\cdot d$, and $d_\eta^*$ is the adjoint of $d_\eta$ with respect to $\langle\cdot,\cdot\rangle_\eta$.

The space of one-cochains is the orthogonal direct sum
\[
C^1=V_{\mathrm{g}}\oplus V_{\mathrm{h}}\oplus V_{\mathrm{t}} ,
\]
in which $V_{\mathrm{g}}=\operatorname{ran} d_0$ is the range of $d_0=d|_{C^0}$; $V_{\mathrm{h}}$ is the span of the three constant directional cochains, the one of direction $\nu$ being equal to $1$ on every link $(x;\{\nu\})$, $x\in\Lambda$, and to $0$ on all other links; and $V_{\mathrm{t}}=\left(V_{\mathrm{g}}\oplus V_{\mathrm{h}}\right)^{\perp}$. Since $d\circ d_0=0$ and Lemma~\ref{lem:harmonic-cochains-coercivity} identifies $\ker\left(d|_{C^1}\right)\cap{V_{\mathrm{g}}}^{\perp}$ with $V_{\mathrm{h}}$, one has
\[
V_{\mathrm{g}}\oplus V_{\mathrm{h}}=\ker\left(d|_{C^1}\right);
\]
in particular $d$ is injective on $V_{\mathrm{t}}$.

\subsubsection{The Wilson operator and the lattice laws}\label{sec:wilson-operator-lattice}

The gauge field is a one-cochain $A\in C^1$, in lattice units; its curvature in physical units is the two-cochain
\[
F=F_\eta=\eta^{-\frac32}\cdot dA\in C^2,\qquad\text{so that}\qquad\frac12\|dA\|^2=\frac12|F|_\eta^2 .
\]

Let $\sigma_1,\sigma_2,\sigma_3$ be the Pauli matrices. Write $\mathbb C^4=\mathbb C^2\otimes\mathbb C^2$, the first factor labelling two blocks of size two, and let $\tau_1,\tau_3$ be the Pauli matrices $\sigma_1,\sigma_3$ acting on the first factor. Put
\[
\gamma_\nu=\operatorname{diag}(\sigma_\nu,-\sigma_\nu)=\tau_3\otimes\sigma_\nu\quad(\nu=1,2,3),\qquad\gamma_4=\tau_1\otimes1 .
\]

Let $q\ge0$ and $\boldsymbol{\mu}\in\mathbb R$. In the description of the Wilson operator the points of $\Lambda$ are written in unit coordinates: $x\in\{-m,\dots,m\}^3$ stands for the point $\eta\cdot x$ of $\Lambda$. An \emph{oriented bond} is a pair $b=(x\to x\pm e_\nu)$ with $\nu\in\{1,2,3\}$; its initial point is $\operatorname{start}(b)=x$, and its final point $\operatorname{end}(b)$ is $x\pm e_\nu$ with the $\nu$-coordinate reduced modulo $n$ to $\{-m,\dots,m\}$. The reversed bond is denoted $-b$. A one-cochain $A$ is read as a function on oriented bonds: $A(x\to x+e_\nu)$ is the value of $A$ on the link $(\eta\cdot x;\{\nu\})$, and $A(-b)=-A(b)$. For $b=(x\to x\pm e_\nu)$ put $\gamma(b)=\pm\gamma_\nu$. The bond $b$ is said to cross the seam if one of the $\nu$-coordinates of $\operatorname{start}(b)$ and $\operatorname{end}(b)$ equals $m$ and the other equals $-m$; this condition does not depend on the order of the two endpoints, so that $-b$ crosses the seam if and only if $b$ does. Put
\[
\varsigma(b)=\begin{cases}-1&\text{if $b$ crosses the seam},\\+1&\text{otherwise};\end{cases}
\]
then $\varsigma(-b)=\varsigma(b)$. The Wilson operator $D_A$ with Wilson parameter one, link phases $\exp(iq\cdot A)$ and antiperiodic shifts is the operator on $\ell^2(\Lambda;\mathbb C^4)$ given by
\begin{equation}\label{eq:continuum-infinite-1}
(D_A\psi)(x)=3\psi(x)-\sum_{b\colon\operatorname{start}(b)=x}\frac12\left(1-\gamma(b)\right)\cdot\varsigma(b)\cdot\exp\left(iq\cdot A(b)\right)\cdot\psi\left(\operatorname{end}(b)\right),
\end{equation}
the sum being over the six oriented bonds with initial point $x$. Restriction to the box $\{-m,\dots,m\}^3$, which contains exactly one point of each coset of $n\mathbb Z^3$ in $\mathbb Z^3$, maps the space $\mathcal S_{\chi^{\mathrm{ap}}}$ of Section~\ref{sec:characters-twisted-spinors}, taken with $s=4$, bijectively onto $\ell^2(\Lambda;\mathbb C^4)$. This bijection transforms the operator \eqref{eq:gaussian-bounds-1}, formed with the same $q$ and $A$, with the matrices $\gamma_1,\gamma_2,\gamma_3$ above and with the character $\chi^{\mathrm{ap}}$, into $D_A$: the summands $P^-_\nu\cdot T_\nu$ and $P^+_\nu\cdot T_\nu^*$ of \eqref{eq:gaussian-bounds-1} correspond to the bonds $(x\to x+e_\nu)$ and $(x\to x-e_\nu)$, and a shift that leaves the box, that is, a shift along a bond $b$ that crosses the seam, acquires from the antiperiodicity the factor $-1=\varsigma(b)$.

For $\theta\in\mathbb R^3$ let $D^\theta_A$ be the operator obtained from \eqref{eq:continuum-infinite-1} when the term of each bond $b=(x\to x\pm e_\nu)$ is multiplied by $\exp\left(\pm\frac{i\theta_\nu}{n}\right)$. If $q>0$, then $D^\theta_A=D_{A+H_\theta}$, where $H_\theta\in V_{\mathrm{h}}$ is the cochain with the value $\frac{\theta_\nu}{q\cdot n}$ on every link $(x;\{\nu\})$.

One has
\[
\gamma_4\cdot D^\theta_A\cdot\gamma_4=\left(D^\theta_A\right)^*,
\]
because $\gamma_4$ anticommutes with $\gamma_1,\gamma_2,\gamma_3$, while the adjoint of the term of a bond $b$ is the term of $-b$ with $\frac12\left(1-\gamma(-b)\right)$ replaced by $\frac12\left(1+\gamma(-b)\right)$. The matrices $\gamma_1,\gamma_2,\gamma_3$ are block diagonal, so that $D^\theta_A=\operatorname{diag}(D_1,D_2)$ with respect to the two blocks, and $\gamma_4$ interchanges the two blocks. Hence $D_2=D_1^*$, and for real $A$, $\theta$ and $\boldsymbol{\mu}$
\[
\det\left(D^\theta_A+\boldsymbol{\mu}\right)=\left|\det\left(D_1+\boldsymbol{\mu}\right)\right|^2\ge0 ;
\]
this is Lemma~\ref{lem:gaussian-bounds-1}(v).

For $\theta\in\mathbb R^3$ define $w_\theta\colon V_{\mathrm{t}}\to[0,\infty)$ and the number $Z_\theta$ by
\[
w_\theta(A)=\exp\left(-\frac12\|dA\|^2\right)\cdot\det\left(D^\theta_A+\boldsymbol{\mu}\right),\qquad Z_\theta=\int_{V_{\mathrm{t}}}\mathrm{d}A\,w_\theta(A),
\]
and define $\bar w\colon V_{\mathrm{t}}\to[0,\infty)$ and the number $\bar Z$ by
\[
\bar w(A)=\int_{[-\pi,\pi]^3}\frac{\mathrm{d}\theta}{(2\pi)^3}\,w_\theta(A),\qquad\bar Z=\int_{V_{\mathrm{t}}}\mathrm{d}A\,\bar w(A),
\]
where $\mathrm{d}A$ is Lebesgue measure on $V_{\mathrm{t}}$.

One has $0<\bar Z<\infty$. Indeed, for each of the six choices of a direction $\nu$ and a sign, the terms of $D^\theta_A$ that belong to the bonds $b=(x\to x\pm e_\nu)$, $x\in\Lambda$, add up to an operator of norm at most one (the projection $\frac12\left(1\mp\gamma_\nu\right)$ composed with a shift multiplied by phases of modulus one), so that $\left\|D^\theta_A+\boldsymbol{\mu}\right\|\le9+|\boldsymbol{\mu}|$ and
\[
\det\left(D^\theta_A+\boldsymbol{\mu}\right)\le\left(9+|\boldsymbol{\mu}|\right)^{4n^3};
\]
and $d$ is injective on the finite-dimensional space $V_{\mathrm{t}}$, so that $\exp\left(-\frac12\|dA\|^2\right)$ is integrable over $V_{\mathrm{t}}$. This gives $\bar Z<\infty$. The function $(\theta,A)\mapsto w_\theta(A)$ is continuous and non-negative, and at $A=0$ its value is
\[
\prod_p\left[\left(\boldsymbol{\mu}+\sum_\nu\left(1-\cos p_\nu\right)\right)^2+\sum_\nu\sin^2p_\nu\right]^2,\qquad p_\nu=\frac{2\pi\left(j_\nu+\frac12\right)+\theta_\nu}{n},
\]
the product being over the $n^3$ points $p=(p_1,p_2,p_3)$ with $(j_1,j_2,j_3)\in\{0,\dots,n-1\}^3$ (a discrete Fourier transform makes $D^\theta_0+\boldsymbol{\mu}$ unitarily equivalent to the direct sum over $p$ of the $4\times4$ matrices $\boldsymbol{\mu}+\sum_\nu\left(1-\cos p_\nu\right)+i\sum_\nu\sin p_\nu\cdot\gamma_\nu$, and the determinant of each of them is the corresponding factor). This value is positive when no $\theta_\nu$ lies in $\pi\mathbb Z$. Hence $w_\theta(A)>0$ on a non-empty open subset of $[-\pi,\pi]^3\times V_{\mathrm{t}}$, and $\bar Z>0$. Moreover $Z_0>0$ if $\boldsymbol{\mu}>-6$: for $A=0$ and $\theta=0$, $n$ being odd, the only $p$ with $\sin p_\nu=0$ for $\nu=1,2,3$ is $p=(\pi,\pi,\pi)$, and there the bracket equals $(\boldsymbol{\mu}+6)^2>0$; so $w_0(0)>0$, and $Z_0>0$ by continuity.

For $\theta\in\mathbb R^3$ with $Z_\theta>0$ let $\mathsf P^\theta=\mathsf P^\theta_{\eta,R;q,\boldsymbol{\mu}}$ be the law on $C^2$ of $F=\eta^{-\frac32}\cdot dA$ when $A$ is distributed on $V_{\mathrm{t}}$ with the density $Z_\theta^{-1}\cdot w_\theta$ with respect to $\mathrm{d}A$, and let $E^\theta$ be the expectation under $\mathsf P^\theta$. The law $\mathsf P^0$ is the \emph{pinned law}, or zero-harmonic law; its expectation is $E^0$.

The \emph{averaged law} $\bar{\mathsf P}=\bar{\mathsf P}_{\eta,R;q,\boldsymbol{\mu}}$ is the law on $C^2$ of $F$ when $A$ is distributed on $V_{\mathrm{t}}$ with the density $\bar Z^{-1}\cdot\bar w$; $\bar E$ denotes the expectation under $\bar{\mathsf P}$. By Fubini's theorem
\[
\bar{\mathsf P}=\int_{[-\pi,\pi]^3}\mathrm{d}\omega(\theta)\,\mathsf P^\theta,\qquad\mathrm{d}\omega(\theta)=\frac{Z_\theta}{\int_{[-\pi,\pi]^3}\mathrm{d}\vartheta\,Z_\vartheta}\,\mathrm{d}\theta ;
\]
here $\omega$ is a probability measure on $[-\pi,\pi]^3$, and the set of the $\theta$ with $Z_\theta=0$ is an $\omega$-null set. For $q>0$ the expectation $\bar E$ is the expectation $\overline{\mathsf E}$ of \eqref{eq:gaussian-bounds-3}, applied to functions of $F=\eta^{-\frac32}\cdot dA$; the cochain $c^\theta$ in \eqref{eq:gaussian-bounds-3} is $H_\theta$.

Multiplying all the weights $w_\theta$ by one positive constant, such as a vacuum factor, changes none of these laws.

A \emph{parameter assignment} is a map which assigns to every cutoff $(N,T)$ a pair $(q,\boldsymbol{\mu})\in[0,\infty)\times(-1,\infty)$; no relation between its values at different cutoffs is required. In Theorem~\ref{thm:main-theorem} the parameter assignment is
\[
q=e\cdot\sqrt{\eta},\qquad\boldsymbol{\mu}=\bar{m}\cdot\eta+m_{0,N},
\]
with the charge $e$ and the numbers $m_{0,N}$ of that theorem. The parameter assignment is arbitrary in each subsequent assertion on the averaged law, and no constant in such an assertion depends on it.

\subsubsection{Test forms and the test map}\label{sec:test-forms-test}

Let $\mathcal S(\mathbb R^3;\Lambda^2)$ be the space of two-forms $f=\sum_{\mu<\nu}f_{\mu\nu}\,\mathrm{d}y_\mu\wedge\mathrm{d}y_\nu$ on $\mathbb R^3$ whose coefficients $f_{\mu\nu}$ are Schwartz functions, real or complex. The elements of $\mathcal S(\mathbb R^3;\Lambda^2)$ are called tests, and a test is real if its coefficients are real-valued. For a test $f$ the periodisation
\[
f^R(y)=\sum_{k\in\mathbb Z^3}f(y+R\cdot k)\qquad(y\in\mathbb R^3)
\]
is $R$-periodic, hence a two-form on $\mathbb T_R$, and one puts
\begin{align}
f_\eta(\sigma)&=\eta^{-3}\int_{Q(\sigma)}\mathrm{d}y\,f^R_{\mu\nu}(y)\qquad\left(\sigma=(x;\mu\nu)\right),\notag\\
X_{\eta,R}(f)&=\langle f_\eta,F\rangle_\eta\qquad\left(F\in C^2\right).\label{eq:continuum-infinite-2}
\end{align}
Thus $f_\eta(\sigma)$ is the mean value of the component $f^R_{\mu\nu}$ of the periodised test over the averaging cube $Q(\sigma)$, the cube of side $\eta$ centred at the centre of the plaquette $\sigma$. Hence $f_\eta\in C^2+iC^2$ is a complex two-cochain, and $f_\eta\in C^2$ if $f$ is real; the map $f\mapsto f_\eta$ from $\mathcal S(\mathbb R^3;\Lambda^2)$ to $C^2+iC^2$ is complex-linear, and it is called the test map. For each test $f$, \eqref{eq:continuum-infinite-2} defines the function
\[
X_{\eta,R}(f)\colon C^2\to\mathbb C,\qquad F\mapsto\langle f_\eta,F\rangle_\eta ,
\]
formed with the complex-bilinear pairing of Section~\ref{sec:tori-cells-cochains}; for given $F\in C^2$ its value at $F$ is again written $X_{\eta,R}(f)$. This function is linear in $F$, depends complex-linearly on $f$, and takes real values if $f$ is real. If $\operatorname{supp}f\subset\left(-\frac{R}{2},\frac{R}{2}\right)^3$, then $f^R=f$ on $Q_R$.

For $F\in C^2$ let $F^\sharp=(F^\sharp_{\mu\nu})_{\mu<\nu}$ be the two-form on $\mathbb R^3$ with the coefficients
\[
F^\sharp_{\mu\nu}(y)=F(\sigma)\qquad\text{for }y\in Q(\sigma)+R\mathbb Z^3,\ \sigma=(x;\mu\nu),
\]
which are well defined because the cubes $Q(x;\mu\nu)$, $x\in\Lambda$, partition $\mathbb T_R$. They are bounded $R$-periodic functions, so that $F^\sharp$ is an element of the space $\mathcal S'(\mathbb R^3;\Lambda^2)$ of tempered distributions with values in two-forms, and
\begin{equation}\label{eq:continuum-infinite-3}
F^\sharp(f):=\int_{\mathbb R^3}\mathrm{d}y\,\sum_{\mu<\nu}f_{\mu\nu}(y)\cdot F^\sharp_{\mu\nu}(y)=X_{\eta,R}(f)\qquad\left(f\in\mathcal S(\mathbb R^3;\Lambda^2)\right);
\end{equation}
indeed, by periodicity the integral equals $\sum_{\mu<\nu}\int_{Q_R}\mathrm{d}y\,f^R_{\mu\nu}(y)\cdot F^\sharp_{\mu\nu}(y)$, and modulo $R\mathbb Z^3$ the cube $Q_R$ is the disjoint union of the cubes $Q(x;\mu\nu)$, $x\in\Lambda$, on each of which $F^\sharp_{\mu\nu}$ is constant; hence the integral equals
\[
\sum_\sigma F(\sigma)\cdot\int_{Q(\sigma)}\mathrm{d}y\,f^R_{\mu\nu}(y)=\eta^3\sum_\sigma f_\eta(\sigma)\cdot F(\sigma)=\langle f_\eta,F\rangle_\eta ,
\]
the sums being taken over all plaquettes $\sigma=(x;\mu\nu)$.

The translations and the cubic group act on tests as follows. For $b\in\mathbb R^3$ and $r\in O_{\mathrm{cub}}=O(3)\cap\mathrm{GL}(3,\mathbb Z)$ let
\[
(\tau_bf)(y)=f(y-b),\qquad(U_rf)(y)=\left(\wedge^2r\right)f(r^{-1}y)\qquad(y\in\mathbb R^3),
\]
where the value $f(y)$ is read as the element $\sum_{\alpha<\beta}f_{\alpha\beta}(y)\cdot e_\alpha\wedge e_\beta$ of $\wedge^2\mathbb C^3$, that is, $\mathrm{d}y_\alpha\wedge\mathrm{d}y_\beta$ is identified with $e_\alpha\wedge e_\beta$, and where $\wedge^2r$ is the second exterior power of $r$, the linear map of $\wedge^2\mathbb C^3$ with $\left(\wedge^2r\right)(e_\alpha\wedge e_\beta)=re_\alpha\wedge re_\beta$. Equivalently, $U_rf$ is the pull-back $\left(r^{-1}\right)^*f$ of the two-form $f$ by $r^{-1}$. Indeed, this pull-back evaluates the coefficients at $r^{-1}y$ and replaces $\mathrm{d}y_\alpha$ by the differential of the $\alpha$-th coordinate of $r^{-1}y$, whose coefficient vector is the $\alpha$-th row of $r^{-1}$, that is, the $\alpha$-th column $re_\alpha$ of the orthogonal matrix $r$. For $\nu\in\{1,2,3\}$ let $r_\nu\in O_{\mathrm{cub}}$ be the reflection $y_\nu\mapsto-y_\nu$, and let $\mathsf R_\nu=U_{r_\nu}$. As $r_\nu^{-1}=r_\nu$, this is the pull-back $r_\nu^*$ of two-forms by $r_\nu$; it changes the sign of exactly the components whose indices contain $\nu$: for $\alpha<\beta$,
\[
(\mathsf R_\nu f)_{\alpha\beta}(y)=\begin{cases}-f_{\alpha\beta}(r_\nu y)&\text{if }\nu\in\{\alpha,\beta\},\\ f_{\alpha\beta}(r_\nu y)&\text{otherwise.}\end{cases}
\]
These maps act on $\mathcal S'(\mathbb R^3;\Lambda^2)$ by transposition: for $F\in\mathcal S'(\mathbb R^3;\Lambda^2)$ and $f\in\mathcal S(\mathbb R^3;\Lambda^2)$,
\[
(\tau_b^*F)(f)=F(\tau_bf),\qquad(U_r^*F)(f)=F(U_rf),\qquad(\mathsf R_\nu^*F)(f)=F(\mathsf R_\nu f).
\]

Let $\eta>0$ and $R\ge1$ be real numbers such that $n=R/\eta$ is an integer. On the torus $\eta\mathbb Z^3/R\mathbb Z^3$ the following are defined by the same formulas: the cells with their centres and averaging cubes; the real and the complex cochains with the pairings $\langle\cdot,\cdot\rangle_\eta$ and $\langle\cdot,\cdot\rangle$ and with the norms of Section~\ref{sec:tori-cells-cochains}; the coboundary $d$ and the operators $d_\eta$ and $d_\eta^*$; the test map and the functions $X_{\eta,R}(f)$ of \eqref{eq:continuum-infinite-2}; and the function $F^\sharp$. For each fixed $I$ the cubes $Q(x;I)$ partition $\mathbb T_R$, \eqref{eq:continuum-infinite-3} holds, and Lemmas~\ref{lem:norm-bounds-lattice} and~\ref{lem:lattice-symmetries-cell} hold with the same proofs; here $\eta\mathbb Z^3$ and $R\mathbb Z^3$ are invariant under $O_{\mathrm{cub}}$ for every such pair $\eta$, $R$. If moreover $n\ge2$, the subspaces $V_{\mathrm{g}}$, $V_{\mathrm{h}}$, $V_{\mathrm{t}}$ of $C^1$ are defined as in Section~\ref{sec:tori-cells-cochains}, and the argument given there, with Lemma~\ref{lem:harmonic-cochains-coercivity}, which holds for every integer $n\ge2$, even or odd, shows that $C^1$ is the orthogonal direct sum $V_{\mathrm{g}}\oplus V_{\mathrm{h}}\oplus V_{\mathrm{t}}$ and that $V_{\mathrm{g}}\oplus V_{\mathrm{h}}=\ker\left(d|_{C^1}\right)$.

\subsubsection{Admissible families of lattice laws}\label{sec:admissible-families-lattice}

Let $\nu\in\{1,2,3\}$. Let $\mathcal P^+_\nu$ be the set of the plaquettes of $\Lambda$ whose four vertices all have their $\nu$-th coordinate in $\{0,\eta,\dots,m\eta\}$, and let $\mathcal P^{+\circ}_\nu\subset\mathcal P^+_\nu$ consist of those elements of $\mathcal P^+_\nu$ which lie in neither of the two planes $\{y_\nu=0\}$ and $\{y_\nu=m\eta\}$. The reflection $r_\nu$ maps every cell $\sigma=(x;I)$ onto a cell $r_\nu\sigma$ with the same set of directions $I$. Define $\hat{\mathsf R}_\nu\colon C^2\to C^2$ by
\[
\left(\hat{\mathsf R}_\nu g\right)(\sigma)=
\begin{cases}
-g(r_\nu\sigma)&\text{if }\nu\in I,\\
g(r_\nu\sigma)&\text{if }\nu\notin I,
\end{cases}
\qquad\sigma=(x;I)\text{ a plaquette}.
\]
Since $r_\nu$ is an involution and preserves the set of directions of a cell, $\hat{\mathsf R}_\nu$ is an involution; it permutes the values of $g$ up to signs, so it is orthogonal. For $v\in\eta\mathbb Z^3$ and $r\in O_{\mathrm{cub}}$ let $v_*$ and $r_*$ be the actions on cochains of the translation by $v$ and of $r$, defined in Lemma~\ref{lem:lattice-symmetries-cell}; they are orthogonal and commute with $d$, so that they map the kernel of $d$ in $C^2$ onto itself.

\begin{definition}[Admissible families]\label{def:admissible-families}
Let $(\mathrm{Pr}_{N,T})$ be a family that assigns to each cutoff $(N,T)$ a Borel probability measure $\mathrm{Pr}_{N,T}$ on the space of the closed two-cochains of $\Lambda_{\eta,R}$,
\[
\ker d=\left\{F\in C^2(\Lambda_{\eta,R}):dF=0\right\}.
\]
Such a family is \emph{admissible} if each of its members satisfies the three conditions below. First, $\mathrm{Pr}_{N,T}$ is invariant under the lattice translations and under the cubic group:
\begin{equation}\tag{Ad1}\label{eq:continuum-infinite-Ad1}
\mathrm{Pr}_{N,T}\circ v_*^{-1}=\mathrm{Pr}_{N,T}\quad\text{and}\quad\mathrm{Pr}_{N,T}\circ r_*^{-1}=\mathrm{Pr}_{N,T}\qquad\text{for all }v\in\eta\mathbb Z^3,\ r\in O_{\mathrm{cub}}.
\end{equation}
Second, for every real $g\in C^2(\Lambda)$,
\begin{equation}\tag{Ad2}\label{eq:continuum-infinite-Ad2}
\int\mathrm{d}\mathrm{Pr}_{N,T}(F)\,\exp\left(\langle g,F\rangle_\eta\right)\le\exp\left(\frac12|g|_\eta^2\right).
\end{equation}
Third, for every $\nu\in\{1,2,3\}$, every integer $s\ge1$, all real $g_1,\dots,g_s\in C^2(\Lambda)$ supported on $\mathcal P^{+\circ}_\nu$ and every $c=(c_1,\dots,c_s)\in\mathbb C^s$,
\begin{equation}\tag{Ad3}\label{eq:continuum-infinite-Ad3}
\sum_{a,b=1}^s\bar c_ac_b\int\mathrm{d}\mathrm{Pr}_{N,T}(F)\,\exp\left(i\left\langle g_b-\hat{\mathsf R}_\nu g_a,F\right\rangle_\eta\right)\ge0.
\end{equation}
\end{definition}

The names \eqref{eq:continuum-infinite-Ad1}, \eqref{eq:continuum-infinite-Ad2} and \eqref{eq:continuum-infinite-Ad3} are also given to the same properties of a single Borel probability measure on $C^2(\Lambda)$. By \eqref{eq:continuum-infinite-2}, $X_{\eta,R}(f)=\langle f_\eta,F\rangle_\eta$. Hence, for a Borel probability measure $\mathrm{Pr}_{N,T}$ on $\ker d$ or on $C^2(\Lambda)$, property \eqref{eq:continuum-infinite-Ad2} is equivalent to
\[
\int\mathrm{d}\mathrm{Pr}_{N,T}(F)\,\exp\left(X_{\eta,R}(f)\right)\le\exp\left(\frac12|f_\eta|_\eta^2\right)\qquad\text{for every real }f\in\mathcal S(\mathbb R^3;\Lambda^2).
\]
Indeed, \eqref{eq:continuum-infinite-Ad2} with $g=f_\eta$ is this inequality; conversely, the inequality implies \eqref{eq:continuum-infinite-Ad2}, because $f\mapsto f_\eta$ maps the real elements of $C^\infty_c(\mathbb R^3;\Lambda^2)$ onto $C^2(\Lambda)$. To see the latter, let $\sigma=(x;\mu\nu)$ be a plaquette, let $\tilde c\in\mathbb R^3$ be a representative of its centre $c(\sigma)\in\mathbb T_R$, and put $\tilde Q=\tilde c+\left[-\frac{\eta}{2},\frac{\eta}{2}\right)^3\subset\mathbb R^3$. The side $\eta=\frac{R}{n}$ of $\tilde Q$ is smaller than $R$; hence the translates $\tilde Q+R\cdot k$, $k\in\mathbb Z^3$, are pairwise disjoint, and the quotient map $\mathbb R^3\to\mathbb T_R$ maps $\tilde Q$ bijectively onto $Q(\sigma)$ and carries Lebesgue measure on $\tilde Q$ to Lebesgue measure on $Q(\sigma)$. Let $\phi\ge0$ be a smooth function on $\mathbb R^3$, not identically zero, whose support lies in the interior of $\tilde Q$, and let $f\in C^\infty_c(\mathbb R^3;\Lambda^2)$ be the real two-form with $f_{\mu\nu}=\phi$ and with vanishing other components. The component $f^R_{\mu\nu}(y)=\sum_{k\in\mathbb Z^3}\phi(y+R\cdot k)$ of the periodisation coincides with $\phi$ on $\tilde Q$ and vanishes outside $\tilde Q+R\mathbb Z^3$; as a function on $\mathbb T_R$ it therefore vanishes outside $Q(\sigma)$, and the other components of $f^R$ vanish. Since the cubes $Q(x';\mu\nu)$, $x'\in\Lambda$, are pairwise disjoint in $\mathbb T_R$, \eqref{eq:continuum-infinite-2} shows that $f_\eta$ vanishes at every plaquette other than $\sigma$, while
\[
f_\eta(\sigma)=\eta^{-3}\int_{\tilde Q}\mathrm{d}y\,\phi(y)=\eta^{-3}\int_{\mathbb R^3}\mathrm{d}y\,\phi(y)>0.
\]
Linear combinations of such two-forms give every element of $C^2(\Lambda)$.

\begin{proposition}[Gaussian domination]\label{prop:gaussian-domination}
Let $(N,T)$ be a cutoff, $q\ge0$ and $\boldsymbol{\mu}>-1$. Then for every real $g\in C^2(\Lambda)$
\[
\bar E\exp\left(\langle g,F\rangle_\eta\right)\le\exp\left(\frac12|g|_\eta^2\right).
\]
\end{proposition}

\begin{proof}
As $F=\eta^{-\frac32}dA$, one has $\langle g,F\rangle_\eta=\langle\eta^{\frac32}g,dA\rangle$ and $|g|_\eta=\|\eta^{\frac32}g\|$, the pairing and the norm on the right being the counting ones.

Let $q>0$. By Section~\ref{sec:wilson-operator-lattice}, the image of $\bar{\mathsf P}$ under the map $F\mapsto\eta^{\frac32}F$ of $C^2$ is the harmonic-averaged law of Section~\ref{sec:proca-parent-harmonic}, formed in the four-component representation with the same $q$ and the lattice mass $\boldsymbol{\mu}$; its expectation $\overline{\mathsf E}$ is given by \eqref{eq:gaussian-bounds-3}. Together with the first paragraph this gives
\[
\bar E\exp\left(\langle g,F\rangle_\eta\right)=\overline{\mathsf E}\left[\exp\left\langle\eta^{\frac32}g,\cdot\right\rangle\right].
\]
The hypotheses of Theorem~\ref{thm:limit-vanishing-proca} are fulfilled, as $n\ge3$, $q\ne0$ and $\boldsymbol{\mu}>-1$, and its inequality \eqref{eq:gaussian-bounds-5}, applied to the cochain $\eta^{\frac32}g$, bounds the right side by $\exp\left(\frac12\left\|\eta^{\frac32}g\right\|^2\right)=\exp\left(\frac12|g|_\eta^2\right)$.

Let $q=0$. Then $D^\theta_A=D^\theta_0$ for every $A$, and therefore
\[
\bar w(A)=c_0\cdot\exp\left(-\frac12\|dA\|^2\right)=c_0\cdot\exp\left(-\frac12|F|_\eta^2\right),\qquad c_0=\int_{[-\pi,\pi]^3}\frac{\mathrm{d}\theta}{(2\pi)^3}\,\det\left(D^\theta_0+\boldsymbol{\mu}\right).
\]
The constant $c_0$ is positive, because $\bar Z=c_0\int_{V_{\mathrm{t}}}\mathrm{d}A\,\exp\left(-\frac12\|dA\|^2\right)$ is positive (Section~\ref{sec:wilson-operator-lattice}). The map $A\mapsto\eta^{-\frac32}dA$ is a linear bijection of $V_{\mathrm{t}}$ onto $\operatorname{ran}\,d$, since $d$ vanishes on $V_{\mathrm{g}}\oplus V_{\mathrm{h}}$ and is injective on $V_{\mathrm{t}}$ (Section~\ref{sec:tori-cells-cochains}). Consequently the law $\bar{\mathsf P}$ of $F$ has a density proportional to $\exp\left(-\frac12|F|_\eta^2\right)$ with respect to the Lebesgue measure of $\operatorname{ran}\,d$; that is, under $\bar{\mathsf P}$ the random variable $F$ is a standard Gaussian vector of the Euclidean space $\left(\operatorname{ran}\,d,\langle\cdot,\cdot\rangle_\eta\right)$. Let $\Pi\colon C^2\to C^2$ be the orthogonal projection onto $\operatorname{ran}\,d$. Since $\langle g,F\rangle_\eta=\langle\Pi g,F\rangle_\eta$ for $F\in\operatorname{ran}\,d$, and $|\Pi g|_\eta\le|g|_\eta$,
\[
\bar E\exp\left(\langle g,F\rangle_\eta\right)=\exp\left(\frac12|\Pi g|_\eta^2\right)\le\exp\left(\frac12|g|_\eta^2\right).
\]
\end{proof}

\begin{proposition}[Reflection positivity]\label{prop:reflection-positivity}
Let $(N,T)$ be a cutoff, $q\ge0$, $\boldsymbol{\mu}>-1$ and $\nu\in\{1,2,3\}$. Then for every integer $s\ge1$, all real $g_1,\dots,g_s\in C^2(\Lambda)$ supported on $\mathcal P^+_\nu$ and every $c=(c_1,\dots,c_s)\in\mathbb C^s$,
\[
\sum_{a,b=1}^s\bar c_ac_b\,\bar E\exp\left(i\left\langle g_b-\hat{\mathsf R}_\nu g_a,F\right\rangle_\eta\right)\ge0.
\]
\end{proposition}

\begin{proof}
Regard $\nu$ as the time direction. Since $n=2m+1$ is odd, the labels $t\in\{-m,\dots,m\}$ of the slices $\{y_\nu=t\eta\}$ form the normal labelling of the odd case in Section~\ref{sec:reflected-data}, the mirror, which maps the slice $t$ to the slice $-t$, being attached to the layer $m$. With the notation of Sections~\ref{sec:slices-slice-operators} and~\ref{sec:reflected-data}, a one-cochain is written $A=(A_t,a_t)_{t\in\{-m,\dots,m\}}$, where $A_t$ is the restriction of $A$ to the links of the slice $t$, and $a_t$ is its restriction to the $\nu$-links that join the slice $t$ to the slice $t+1$ (for $t=m$, the slice $m$ to the slice $-m$). Let $d_s$ be the coboundary within a slice, and put
\[
\mathcal E_t=A_{t+1}-A_t-d_sa_t,
\]
indices being taken modulo $n$. As in Section~\ref{sec:slices-slice-operators}, the links and the plaquettes of each slice are labelled by their coordinates other than the $\nu$-th and by their directions; through these labels $\mathcal E_t$ is a function on the links of the slice $t$, and $d_sA_t$ a function on its plaquettes. By that section, $dA$ equals $d_sA_t$ on the plaquettes lying in the slice $t$. A plaquette that lies in no slice has the form $(x;I)$ with $\nu\in I$ and $x$ in some slice $t$, and it connects the slice $t$ with the slice $t+1$; if $\alpha$ denotes the element of $I$ other than $\nu$, Section~\ref{sec:slices-slice-operators} gives
\[
dA(x;I)=\begin{cases}\mathcal E_t(x;\{\alpha\})&\text{if }\nu<\alpha,\\-\mathcal E_t(x;\{\alpha\})&\text{if }\alpha<\nu.\end{cases}
\]
So the sign depends on the plaquette only through its set of directions $I$. The elements of $\mathcal P^+_\nu$ are the plaquettes lying in one of the slices $0,1,\dots,m$, together with the plaquettes connecting the slice $t$ with the slice $t+1$ for some $t\in\{0,\dots,m-1\}$. Consequently the function $P\colon C^1\to\mathbb C$,
\[
P(A)=\sum_{b=1}^sc_b\exp\left(i\langle g_b,F\rangle_\eta\right),\qquad F=\eta^{-\frac32}dA,
\]
is bounded and continuous, and it depends on $A$ only through $(A_t)_{0\le t\le m}$ and $(a_t)_{0\le t\le m-1}$; in other words, $P$ belongs to the class $\mathcal A_+$ of Proposition~\ref{prop:gaussian-bounds-2}.

Let $\vartheta\colon C^1\to C^1$ be the reflection given by $(\vartheta A)_t=A_{-t}$ and $(\vartheta a)_t=-a_{-t-1}$. Then
\[
d_s(\vartheta A)_t=d_sA_{-t},\qquad\mathcal E_t(\vartheta A)=-\mathcal E_{-t-1}(A).
\]
The reflection $r_\nu$ changes no coordinate other than the $\nu$-th, and it keeps the set of directions of every cell. It maps the plaquettes of the slice $t$ onto those of the slice $-t$, and a plaquette $(x;I)$ with $\nu\in I$ connecting the slices $t$ and $t+1$ onto the plaquette with the same set of directions $I$ connecting the slices $-t-1$ and $-t$; by the description of $dA$ above, these two plaquettes carry the same sign. Hence the first identity says that $d(\vartheta A)(\sigma)=dA(r_\nu\sigma)$ for every plaquette $\sigma=(x;I)$ with $\nu\notin I$, and the second says that $d(\vartheta A)(\sigma)=-dA(r_\nu\sigma)$ for every plaquette $\sigma=(x;I)$ with $\nu\in I$. By the definition of $\hat{\mathsf R}_\nu$, this means
\[
d(\vartheta A)=\hat{\mathsf R}_\nu\,dA\qquad(A\in C^1).
\]
Being an orthogonal involution, $\hat{\mathsf R}_\nu$ is symmetric. As the $g_a$ are real and $\eta^{-\frac32}d(\vartheta A)=\hat{\mathsf R}_\nu F$, it follows that
\[
\overline{P(\vartheta A)}=\sum_{a=1}^s\bar c_a\exp\left(-i\left\langle g_a,\hat{\mathsf R}_\nu F\right\rangle_\eta\right)=\sum_{a=1}^s\bar c_a\exp\left(-i\left\langle\hat{\mathsf R}_\nu g_a,F\right\rangle_\eta\right),
\]
and hence
\[
\overline{P(\vartheta A)}\cdot P(A)=\sum_{a,b=1}^s\bar c_ac_b\exp\left(i\left\langle g_b-\hat{\mathsf R}_\nu g_a,F\right\rangle_\eta\right)=\Phi(dA),
\]
where $\Phi\colon C^2\to\mathbb C$ is the bounded continuous function
\[
\Phi(\omega)=\sum_{a,b=1}^s\bar c_ac_b\exp\left(i\eta^{-\frac32}\left\langle g_b-\hat{\mathsf R}_\nu g_a,\omega\right\rangle_\eta\right).
\]
The left side of the assertion is $\bar E\left[\Phi(dA)\right]$.

For $\varepsilon>0$ let $\mathsf E^{\mathrm{ap}}_\varepsilon$ be the expectation of the Proca parent with the antiperiodic character $\chi^{\mathrm{ap}}$ (Section~\ref{sec:proca-parent-harmonic}), that is, of the probability measure
\[
Z^{\mathrm{ap}}_\varepsilon(0)^{-1}\exp\left(-\frac12\|dA\|^2-\frac{\varepsilon}{2}\|A\|^2\right)\cdot\det\left(D_A+\boldsymbol{\mu}\right)\mathrm{d}A
\]
on $C^1$. Here $\det(D_A+\boldsymbol{\mu})$ is the determinant $\Delta^{\chi^{\mathrm{ap}}}(A)$ of Section~\ref{sec:characters-twisted-spinors}, because \eqref{eq:continuum-infinite-1} is the operator \eqref{eq:gaussian-bounds-1} of the four-component representation with the character $\chi^{\mathrm{ap}}$ (Section~\ref{sec:wilson-operator-lattice}). All components of $\chi^{\mathrm{ap}}$ equal $-1$, so Theorem~\ref{thm:gaussian-domination-positive}(b) gives $Z^{\mathrm{ap}}_\varepsilon(0)>0$, whatever the real number $q$. As $\chi^{\mathrm{ap}}_\nu=-1$ and $P\in\mathcal A_+$, Proposition~\ref{prop:gaussian-bounds-2} states that the integral of $\overline{P(\vartheta A)}\cdot P(A)=\Phi(dA)$ with respect to the measure $\exp\left(-\frac12\|dA\|^2-\frac{\varepsilon}{2}\|A\|^2\right)\cdot\det\left(D_A+\boldsymbol{\mu}\right)\mathrm{d}A$ on $C^1$ is non-negative. That integral equals $Z^{\mathrm{ap}}_\varepsilon(0)\cdot\mathsf E^{\mathrm{ap}}_\varepsilon\left[\Phi(dA)\right]$; consequently
\[
\mathsf E^{\mathrm{ap}}_\varepsilon\left[\Phi(dA)\right]\ge0\qquad\left(\varepsilon>0,\ q\in\mathbb R\right).
\]

Let $q>0$. As in the proof of Proposition~\ref{prop:gaussian-domination}, the image of $\bar{\mathsf P}$ under $F\mapsto\eta^{\frac32}F$ is the harmonic-averaged law, whose expectation is $\overline{\mathsf E}$; since $\Phi(dA)=\Phi\left(\eta^{\frac32}F\right)$, it follows that $\bar E\left[\Phi(dA)\right]=\overline{\mathsf E}[\Phi]$. The hypotheses of Theorem~\ref{thm:limit-vanishing-proca} are fulfilled as in that proof, and $\Phi$ is bounded and Borel. The identity \eqref{eq:gaussian-bounds-4} of that theorem, for the character $\chi^{\mathrm{ap}}$, therefore exhibits $\overline{\mathsf E}[\Phi]$ as the limit of the non-negative numbers $\mathsf E^{\mathrm{ap}}_\varepsilon\left[\Phi(dA)\right]$:
\[
\bar E\left[\Phi(dA)\right]=\overline{\mathsf E}[\Phi]=\lim_{\varepsilon\downarrow0}\mathsf E^{\mathrm{ap}}_\varepsilon\left[\Phi(dA)\right]\ge0.
\]

Let $q=0$. Then the determinant is the constant $\det(D_0+\boldsymbol{\mu})$, which is positive (Section~\ref{sec:wilson-operator-lattice}). Write $A=A'+A''$ with $A'\in V_{\mathrm{g}}\oplus V_{\mathrm{h}}$ and $A''\in V_{\mathrm{t}}$ (Section~\ref{sec:tori-cells-cochains}). Then $\|A\|^2=\|A'\|^2+\|A''\|^2$ and $dA=dA''$, so $\Phi(dA)$ depends on $A''$ only, and under $\mathsf E^{\mathrm{ap}}_\varepsilon$ the law of $A''$ has, with respect to the Lebesgue measure of $V_{\mathrm{t}}$, a density proportional to the function
\[
u_\varepsilon(A)=\exp\left(-\frac12\|dA\|^2-\frac{\varepsilon}{2}\|A\|^2\right)\qquad(A\in V_{\mathrm{t}}).
\]
As $\varepsilon\downarrow0$, $u_\varepsilon$ increases to $u_0(A)=\exp\left(-\frac12\|dA\|^2\right)$, which is integrable on $V_{\mathrm{t}}$, since $d$ is injective on $V_{\mathrm{t}}$, and which is proportional to $\bar w$ on $V_{\mathrm{t}}$ (proof of Proposition~\ref{prop:gaussian-domination}). For a bounded Borel function $\Psi\colon C^2\to\mathbb C$ and $\varepsilon\ge0$ put
\[
J_\varepsilon(\Psi)=\int_{V_{\mathrm{t}}}\mathrm{d}A\,\Psi(dA)\cdot u_\varepsilon(A).
\]
By dominated convergence, with the dominating function $\sup|\Psi|\cdot u_0$, one has $J_\varepsilon(\Psi)\to J_0(\Psi)$ as $\varepsilon\downarrow0$. Moreover $\mathsf E^{\mathrm{ap}}_\varepsilon\left[\Phi(dA)\right]=J_\varepsilon(\Phi)/J_\varepsilon(1)$ for $\varepsilon>0$, and $\bar E\left[\Phi(dA)\right]=J_0(\Phi)/J_0(1)$ with $J_0(1)>0$. Hence the non-negative numbers $\mathsf E^{\mathrm{ap}}_\varepsilon\left[\Phi(dA)\right]$ converge to $\bar E\left[\Phi(dA)\right]$, the left side of the assertion, which is therefore non-negative.
\end{proof}

\begin{corollary}[The averaged laws are admissible]\label{cor:averaged-laws-admissible}
Let a parameter assignment be given, and let $(q,\boldsymbol{\mu})$ denote its value at the cutoff $(N,T)$. Then the averaged laws $\bar{\mathsf P}_{\eta,R;q,\boldsymbol{\mu}}$, one for each cutoff $(N,T)$, form an admissible family.
\end{corollary}

\begin{proof}
At every cutoff the value $(q,\boldsymbol{\mu})$ of a parameter assignment satisfies $q\ge0$ and $\boldsymbol{\mu}>-1$ (Section~\ref{sec:wilson-operator-lattice}). Under $\bar{\mathsf P}_{\eta,R;q,\boldsymbol{\mu}}$ one has $F=\eta^{-\frac32}dA\in\operatorname{ran}\,d$, and $\operatorname{ran}\,d\subset\ker d$ because $d\circ d=0$; so $\bar{\mathsf P}_{\eta,R;q,\boldsymbol{\mu}}$ is a Borel probability measure on $\ker d$. Lemma~\ref{lem:invariances-averaged-law}(b) yields the invariance \eqref{eq:continuum-infinite-Ad1}, and Proposition~\ref{prop:gaussian-domination} yields the bound \eqref{eq:continuum-infinite-Ad2}. A real cochain supported on $\mathcal P^{+\circ}_\nu$ is supported on the larger set $\mathcal P^+_\nu$; therefore \eqref{eq:continuum-infinite-Ad3} is a special case of Proposition~\ref{prop:reflection-positivity}.
\end{proof}

\subsubsection{Hermite--Sobolev spaces}\label{sec:hermite-sobolev-spaces}

Let $h_k$, $k\in\mathbb N_0^3$, be the Hermite functions on $\mathbb R^3$, normalised in $L^2(\mathbb R^3)$, and let $\mathcal K=\mathbb N_0^3\times\{12,13,23\}$. For $\kappa=(k,c)\in\mathcal K$ put $h_\kappa=h_k\cdot e_c$, where $e_c$ is the constant two-form whose component $c$ equals $1$ and whose other two components vanish. The family $(h_\kappa)_{\kappa\in\mathcal K}$ is an orthonormal basis of $L^2(\mathbb R^3;\Lambda^2)$, and the operator $\mathsf A=1+|x|^2-\Delta$, acting on each component, satisfies
\[
\mathsf A h_\kappa=\lambda_\kappa h_\kappa,\qquad\lambda_\kappa=2|k|_1+4\qquad\left(\kappa=(k,c)\in\mathcal K\right).
\]
Write $\mathcal S=\mathcal S(\mathbb R^3;\Lambda^2)$ and $\mathcal S'=\mathcal S'(\mathbb R^3;\Lambda^2)$. For $s\in\mathbb R$ and $T\in\mathcal S'$ put
\[
\|T\|_s^2=\sum_{\kappa\in\mathcal K}\lambda_\kappa^{2s}\cdot|T(h_\kappa)|^2\in[0,\infty],\qquad H_s=\left\{T\in\mathcal S':\|T\|_s<\infty\right\}.
\]
By the $N$-representation theorems for $\mathcal S$ and for $\mathcal S'$, one has $\mathcal S=\bigcap_{s\in\mathbb R}H_s$, the topology of $\mathcal S$ is the one defined by the norms $\|\cdot\|_s$, $s\in\mathbb R$, one has $\mathcal S'=\bigcup_{s\in\mathbb R}H_{-s}$, and the Hermite expansion of every element of $\mathcal S$ converges in $\mathcal S$. Consequently $T\mapsto\left(T(h_\kappa)\right)_{\kappa\in\mathcal K}$ is an isometry of $H_s$, normed by $\|\cdot\|_s$, onto the weighted $\ell^2$ space of the families $(t_\kappa)_{\kappa\in\mathcal K}$ with $\sum_\kappa\lambda_\kappa^{2s}\cdot|t_\kappa|^2<\infty$; such a family is the image of the distribution $f\mapsto\sum_{(k,c)\in\mathcal K}t_{(k,c)}\int\mathrm{d}x\,h_k(x)\cdot f_c(x)$, which is continuous on $\mathcal S$ by the Cauchy--Schwarz inequality. So $H_s$ is a separable Hilbert space. The same statements hold for the Hermite functions on $\mathbb R^{3\ell}$, for every integer $\ell\ge1$. For $T\in H_{-s}$ and $f\in\mathcal S$ the Hermite expansion of $f$ and the Cauchy--Schwarz inequality give
\[
|T(f)|\le\|T\|_{-s}\cdot\|f\|_s.
\]

Let $\mathcal C$ be the cylinder $\sigma$-algebra of $\mathcal S'$, that is, the $\sigma$-algebra generated by the maps $F\mapsto F(f)$, $f\in\mathcal S$. For every $s\in\mathbb R$ one has $H_{-s}\in\mathcal C$, because $F\mapsto\|F\|_{-s}^2$ is a countable sum of non-negative $\mathcal C$-measurable functions; and the trace $\left\{B\cap H_{-s}:B\in\mathcal C\right\}$ of $\mathcal C$ on $H_{-s}$ is the Borel $\sigma$-algebra of $H_{-s}$. Indeed, by the last inequality each map $F\mapsto F(f)$, $f\in\mathcal S$, is continuous on $H_{-s}$. Conversely, the Borel sets of the separable Hilbert space $H_{-s}$ are generated by its continuous linear functionals; by the isometry above, these are the maps $F\mapsto F(f)=\sum_\kappa F(h_\kappa)\cdot f(h_\kappa)$ with $f\in H_s$, and each of them is the pointwise limit on $H_{-s}$ of the maps $F\mapsto F(f_j)$, where the $f_j\in\mathcal S$ are finite partial sums of the Hermite expansion of $f$.

For $a\in\mathbb R$ put
\begin{multline}\label{eq:continuum-infinite-4}
\mathsf Z_a=\sum_{\kappa\in\mathcal K}\lambda_\kappa^{-a}=3\sum_{j\ge0}\frac{(j+1)(j+2)}{2}\cdot(2j+4)^{-a}<\infty\iff a>3;\\
\mathsf Z_4=\frac{3}{32}\sum_{j\ge0}\frac{j+1}{(j+2)^3}\le\frac{3}{32}\left(\frac{\pi^2}{6}-1\right)<\frac{1}{16}.
\end{multline}
Here the factor $3$ is the number of components, $\frac{(j+1)(j+2)}{2}$ is the number of the multi-indices $k\in\mathbb N_0^3$ with $|k|_1=j$, and the terms of the series are of order $j^{2-a}$; the first inequality for $\mathsf Z_4$ follows from $(j+1)\cdot(j+2)^{-3}\le(j+2)^{-2}$ and $\sum_{j\ge0}(j+2)^{-2}=\frac{\pi^2}{6}-1$, and the last one from $\pi^2<10$.

\subsection{Averages of tests and the averaged law}\label{sec:averages-tests-averaged}

\begin{lemma}[Norm bounds for lattice averages]\label{lem:norm-bounds-lattice}
For $R\ge1$ let $Q_R=[-R/2,R/2)^3$ and define $p_R\colon\mathcal S(\mathbb R^3;\Lambda^2)\to[0,\infty)$ by
\[
p_R(f)=\|f\|_{L^2}+26R^{-2}\cdot\|\,|y|^2f\,\|_{L^2(\mathbb R^3\setminus Q_R)},
\]
where $|y|^2f$ denotes the function $y\mapsto|y|^2f(y)$ on $\mathbb R^3$. Let $f\in\mathcal S(\mathbb R^3;\Lambda^2)$, let $\eta$ and $R\ge1$ be the spacing and the period of Section~\ref{sec:tori-cells-cochains}, let $f^R(y)=\sum_{k\in\mathbb Z^3}f(y+R\cdot k)$, and let $f_\eta$ be given by \eqref{eq:continuum-infinite-2}. Then
\[
|f_\eta|_\eta\le\|f^R\|_{L^2(\mathbb T_R)}\le p_R(f)\le p_1(f),\qquad\lim_{R\to\infty}p_R(f)=\|f\|_{L^2}.
\]
If $\operatorname{supp}f\subset y_0+(-R/2,R/2)^3$ for some $y_0\in\mathbb R^3$, then $|f_\eta|_\eta\le\|f\|_{L^2}$. For the Hermite functions $h_\kappa$ and the eigenvalues $\lambda_\kappa$ of Section~\ref{sec:hermite-sobolev-spaces} one has $p_1(h_\kappa)\le31\lambda_\kappa$ for every $\kappa\in\mathcal K$.
\end{lemma}

\begin{proof}
For $z\in\mathbb R^3$ write $|z|_\infty=\max_{1\le i\le3}|z_i|$.

Let $\sigma=(x;\mu\nu)$ be a plaquette. By \eqref{eq:continuum-infinite-2}, $f_\eta(\sigma)$ is the mean value of $f^R_{\mu\nu}$ over the cube $Q(\sigma)$, whose volume is $\eta^3$; the Cauchy--Schwarz inequality therefore gives
\[
|f_\eta(\sigma)|^2\le\eta^{-3}\int_{Q(\sigma)}\mathrm{d}y\,|f^R_{\mu\nu}(y)|^2.
\]
For fixed $\mu<\nu$ the cubes $Q(x;\mu\nu)$, $x\in\Lambda$, tile $\mathbb T_R$ (Section~\ref{sec:tori-cells-cochains}). Multiplication by $\eta^3$ and summation over $x\in\Lambda$ and over the pairs $\mu<\nu$ yield
\[
|f_\eta|_\eta^2=\eta^3\sum_{\mu<\nu}\sum_{x\in\Lambda}|f_\eta(x;\mu\nu)|^2\le\sum_{\mu<\nu}\int_{\mathbb T_R}\mathrm{d}y\,|f^R_{\mu\nu}(y)|^2=\|f^R\|_{L^2(\mathbb T_R)}^2,
\]
which is the first inequality.

The cube $Q_R$ is a fundamental domain of $R\mathbb Z^3$, so that $\|f^R\|_{L^2(\mathbb T_R)}=\|f^R\|_{L^2(Q_R)}$, and the cubes $Q_R+R\cdot k$, $k\in\mathbb Z^3$, tile $\mathbb R^3$. Suppose that $\operatorname{supp}f\subset y_0+(-R/2,R/2)^3$ for some $y_0\in\mathbb R^3$. For $k\ne k'$ in $\mathbb Z^3$ the open cubes $y_0+R\cdot k+(-R/2,R/2)^3$ and $y_0+R\cdot k'+(-R/2,R/2)^3$ are disjoint, because $|R\cdot(k-k')|_\infty\ge R$. Hence the sets $\operatorname{supp}f+R\cdot k$, $k\in\mathbb Z^3$, are pairwise disjoint, at each $y\in\mathbb R^3$ at most one term of the series $\sum_kf(y+R\cdot k)$ differs from zero, and
\[
|f^R(y)|^2=\sum_{k\in\mathbb Z^3}|f(y+R\cdot k)|^2\qquad(y\in\mathbb R^3).
\]
Integration over $Q_R$ therefore gives
\[
\|f^R\|_{L^2(\mathbb T_R)}^2=\sum_{k\in\mathbb Z^3}\int_{Q_R+R\cdot k}\mathrm{d}y\,|f(y)|^2=\|f\|_{L^2}^2,
\]
and the first inequality yields $|f_\eta|_\eta\le\|f\|_{L^2}$.

For the second inequality put $a_k=\|f\|_{L^2(Q_R+R\cdot k)}$ for $k\in\mathbb Z^3$; then $a_0\le\|f\|_{L^2}$. On $Q_R$ the function $f^R$ is the sum of the functions $f(\cdot+R\cdot k)$, $k\in\mathbb Z^3$, and the norm of $f(\cdot+R\cdot k)$ in $L^2(Q_R)$ is $a_k$. Minkowski's inequality gives
\[
\|f^R\|_{L^2(\mathbb T_R)}=\|f^R\|_{L^2(Q_R)}\le\sum_{k\in\mathbb Z^3}a_k.
\]
Let $k\ne0$ and $y\in Q_R+R\cdot k$. Then $|y-R\cdot k|_\infty\le R/2$ and $|k|_\infty\ge1$, so that
\[
|y|\ge|y|_\infty\ge R\cdot|k|_\infty-\frac{R}{2}\ge\frac{R\cdot|k|_\infty}{2}.
\]
Thus $1\le4R^{-2}\cdot|k|_\infty^{-2}\cdot|y|^2$ on $Q_R+R\cdot k$, and
\[
a_k\le4R^{-2}\cdot|k|_\infty^{-2}\cdot\|\,|y|^2f\,\|_{L^2(Q_R+R\cdot k)}.
\]
The cubes $Q_R+R\cdot k$, $k\ne0$, tile $\mathbb R^3\setminus Q_R$. The Cauchy--Schwarz inequality for the sum over $k\ne0$ therefore gives
\[
\sum_{k\ne0}a_k\le4R^{-2}\cdot\left(\sum_{k\ne0}|k|_\infty^{-4}\right)^{\frac12}\cdot\|\,|y|^2f\,\|_{L^2(\mathbb R^3\setminus Q_R)}.
\]
For every integer $j\ge1$,
\[
\#\left|\{k\in\mathbb Z^3:|k|_\infty=j\}\right|=(2j+1)^3-(2j-1)^3=24j^2+2,
\]
and $\sum_{j\ge1}j^{-2}=\frac{\pi^2}{6}$, $\sum_{j\ge1}j^{-4}=\frac{\pi^4}{90}$. Consequently
\[
\sum_{k\ne0}|k|_\infty^{-4}=\sum_{j\ge1}(24j^2+2)\cdot j^{-4}=4\pi^2+\frac{\pi^4}{45}<41.65,
\]
where the last inequality holds because $\pi^2<9.87$. Moreover $4\cdot41.65^{\frac12}<26$, since $6.5^2=42.25$. Together with $a_0\le\|f\|_{L^2}$ this proves $\|f^R\|_{L^2(\mathbb T_R)}\le p_R(f)$.

For $R\ge1$ one has $R^{-2}\le1$ and $Q_1\subset Q_R$, hence $\mathbb R^3\setminus Q_R\subset\mathbb R^3\setminus Q_1$ and $p_R(f)\le p_1(f)$. Every $y\in\mathbb R^3$ lies in $Q_R$ as soon as $R>2|y|_\infty$, and the function $y\mapsto|y|^4\cdot|f(y)|^2$ is integrable, $f$ being a Schwartz function. By dominated convergence,
\[
\lim_{R\to\infty}\int_{\mathbb R^3\setminus Q_R}\mathrm{d}y\,|y|^4\cdot|f(y)|^2=0,
\]
and therefore $p_R(f)\to\|f\|_{L^2}$ as $R\to\infty$.

It remains to estimate $p_1(h_\kappa)$. Let $\mathsf M$ be the operator of multiplication by $1+|y|^2$, and let $P_j=-i\partial_j$ for $j=1,2,3$; these operators act componentwise on $\mathcal S(\mathbb R^3;\Lambda^2)$, map this space into itself, and are symmetric for the scalar product $\langle\cdot,\cdot\rangle_{L^2}$ of $L^2(\mathbb R^3;\Lambda^2)$. Since $\sum_{j=1}^3P_j^2=-\Delta$, the operator $\mathsf A$ of Section~\ref{sec:hermite-sobolev-spaces} is $\mathsf A=\mathsf M-\Delta=\mathsf M+\sum_{j=1}^3P_j^2$. The commutator $[P_j,\mathsf M]$ is multiplication by $-2i\,y_j$, so $[P_j,[P_j,\mathsf M]]=-2$ and
\[
\mathsf M P_j^2+P_j^2\mathsf M=2P_j\mathsf M P_j+[P_j,[P_j,\mathsf M]]=2P_j\mathsf M P_j-2\qquad(j=1,2,3).
\]
Expanding $\|(\mathsf M-\Delta)u\|_{L^2}^2$ and summing these identities over $j$, one obtains for $u\in\mathcal S(\mathbb R^3;\Lambda^2)$
\[
\|\mathsf A u\|_{L^2}^2=\|\mathsf M u\|_{L^2}^2+\|\Delta u\|_{L^2}^2+2\sum_{j=1}^3\langle P_ju,\mathsf M P_ju\rangle_{L^2}-6\|u\|_{L^2}^2\ge\|\mathsf M u\|_{L^2}^2-6\|u\|_{L^2}^2,
\]
where the inequality holds because $\langle P_ju,\mathsf M P_ju\rangle_{L^2}\ge0$, as $\mathsf M\ge1$. Moreover $\|\mathsf A u\|_{L^2}\ge4\|u\|_{L^2}$, since $\mathsf A$ is symmetric and the $h_\kappa$ form an orthonormal basis of $L^2(\mathbb R^3;\Lambda^2)$ with $\mathsf A h_\kappa=\lambda_\kappa h_\kappa$ and $\lambda_\kappa\ge4$ (Section~\ref{sec:hermite-sobolev-spaces}). Hence
\[
\|\mathsf M u\|_{L^2}^2\le\|\mathsf A u\|_{L^2}^2+6\|u\|_{L^2}^2\le\frac{11}{8}\|\mathsf A u\|_{L^2}^2.
\]
Take $u=h_\kappa$, so that $\|h_\kappa\|_{L^2}=1$ and $\|\mathsf A h_\kappa\|_{L^2}=\lambda_\kappa$. Since $|y|^2\le1+|y|^2$, since $26\cdot\left(\frac{11}{8}\right)^{\frac12}\le30.5$ (indeed $26^2\cdot11=7436\le7442=8\cdot30.5^2$), and since $1\le\frac{\lambda_\kappa}{4}$,
\[
p_1(h_\kappa)\le1+26\|\mathsf M h_\kappa\|_{L^2}\le1+26\cdot\left(\frac{11}{8}\right)^{\frac12}\lambda_\kappa\le\left(\frac14+30.5\right)\lambda_\kappa\le31\lambda_\kappa.
\]
\end{proof}

\begin{lemma}[Lattice symmetries and cell averages]\label{lem:lattice-symmetries-cell}
Let $v\in\eta\mathbb Z^3$ and $r\in O_{\mathrm{cub}}$. For $0\le k\le3$ and a $k$-cell $\sigma=(x;I)$ of $\Lambda$ (Section~\ref{sec:tori-cells-cochains}) with $I=\{i_1<\dots<i_k\}$, let
\[
o(\sigma)=e_{i_1}\wedge\dots\wedge e_{i_k}\in\wedge^k\mathbb R^3
\]
be its orientation, with $o(\sigma)=1\in\wedge^0\mathbb R^3=\mathbb R$ for $k=0$, and let $\sigma+v$ and $r\sigma$ denote the images of $\sigma$ under $y\mapsto y+v$ and $y\mapsto ry$. These two maps of $\mathbb T_R$ permute the $k$-cells of $\Lambda$; they map centres to centres, $c(\sigma+v)=c(\sigma)+v$ and $c(r\sigma)=r\,c(\sigma)$; and they map averaging cubes to averaging cubes up to sets of measure zero: $Q(\sigma+v)=Q(\sigma)+v$, and $r\,Q(\sigma)$ differs from $Q(r\sigma)$ by a set of measure zero. Moreover, for every $k$-cell $\sigma$ there is a sign $\varepsilon(r,\sigma)\in\{1,-1\}$ with
\[
\left(\wedge^kr\right)(o(\sigma))=\varepsilon(r,\sigma)\cdot o(r\sigma).
\]
Let $v_*,r_*\colon C^k\to C^k$ be the linear maps given by
\[
(v_*u)(\sigma+v)=u(\sigma),\qquad(r_*u)(r\sigma)=\varepsilon(r,\sigma)\cdot u(\sigma)\qquad(\sigma\text{ a $k$-cell}),
\]
and by the same formulas on the complexification of $C^k$. Then $v_*$ and $r_*$ are orthogonal for $\langle\cdot,\cdot\rangle_\eta$; they commute with the coboundary, that is, $d\,v_*=v_*\,d$ and $d\,r_*=r_*\,d$ on $C^k$ for $0\le k\le2$; and for every $f\in\mathcal S(\mathbb R^3;\Lambda^2)$ and $\nu\in\{1,2,3\}$, with the actions on tests of Section~\ref{sec:test-forms-test} and the reflection $\hat{\mathsf R}_\nu$ of Section~\ref{sec:admissible-families-lattice}, extended to the complexification of $C^2$ by the same formula,
\[
(\tau_vf)_\eta=v_*f_\eta,\qquad(U_rf)_\eta=r_*f_\eta,\qquad(\mathsf R_\nu f)_\eta=\hat{\mathsf R}_\nu f_\eta.
\]
\end{lemma}

\begin{proof}
\emph{Lattice and periodisation.} The matrix $r$ and its inverse, the transpose of $r$, have integer entries. Hence $r$ maps $\eta\mathbb Z^3$ and $R\mathbb Z^3$ onto themselves, and $y\mapsto ry$ and $y\mapsto y+v$ are bijections of $\mathbb T_R$ which preserve Lebesgue measure and map $\Lambda$ onto $\Lambda$. The columns of $r$ are orthonormal vectors with integer entries. Hence there are a permutation $\pi$ of $\{1,2,3\}$ and signs $\varepsilon_1,\varepsilon_2,\varepsilon_3\in\{1,-1\}$ such that
\[
re_j=\varepsilon_j\cdot e_{\pi(j)}\qquad(j=1,2,3),
\]
that is, $(ry)_{\pi(j)}=\varepsilon_j\cdot y_j$ for $y\in\mathbb R^3$ and $j=1,2,3$. In particular, if $n$ is odd, then $r$ maps the set $\eta\cdot\{-m,\dots,m\}^3$ of the centred representatives of $\Lambda$ onto itself. Periodisation commutes with both actions on tests: $(\tau_vf)^R=\tau_v\left(f^R\right)$, and, since $\ell\mapsto r^{-1}\ell$ is a bijection of $\mathbb Z^3$,
\[
(U_rf)^R(y)=\left(\wedge^2r\right)\sum_{\ell\in\mathbb Z^3}f\left(r^{-1}y+R\cdot r^{-1}\ell\right)=\left(\wedge^2r\right)f^R\left(r^{-1}y\right)=\left(U_rf^R\right)(y)\qquad\left(y\in\mathbb R^3\right).
\]

\emph{Cells, centres, cubes and signs.} Let $\sigma=(x;I)$ be a $k$-cell. Since $[0,\eta]\cdot\left(-e_{\pi(j)}\right)=-\eta\cdot e_{\pi(j)}+[0,\eta]\cdot e_{\pi(j)}$ for $j\in I$,
\[
r\sigma=rx+\sum_{j\in I}[0,\eta]\cdot\varepsilon_j\cdot e_{\pi(j)}=x'+\sum_{j\in I}[0,\eta]\cdot e_{\pi(j)},\qquad x'=rx-\eta\cdot\sum_{j\in I,\ \varepsilon_j=-1}e_{\pi(j)}\in\Lambda .
\]
So $r\sigma$ is the $k$-cell $(x';\pi(I))$ on the directions $\pi(I)$, and
\[
c(r\sigma)=x'+\frac{\eta}{2}\sum_{j\in I}e_{\pi(j)}=rx+\frac{\eta}{2}\sum_{j\in I}\varepsilon_j\cdot e_{\pi(j)}=r\,c(\sigma).
\]
The same holds for $r^{-1}\in O_{\mathrm{cub}}$, so $\sigma\mapsto r\sigma$ is a permutation of the $k$-cells with inverse $\sigma\mapsto r^{-1}\sigma$. The closed cube $[-\eta/2,\eta/2]^3$ is invariant under $r$; therefore $r\,Q(\sigma)=r\,c(\sigma)+r[-\eta/2,\eta/2)^3$ and $Q(r\sigma)=c(r\sigma)+[-\eta/2,\eta/2)^3$ differ by a set of measure zero. Translations act in the same way without signs: $\sigma+v=(x+v;I)$, so $c(\sigma+v)=c(\sigma)+v$ and $Q(\sigma+v)=Q(\sigma)+v$, and $\sigma\mapsto\sigma+v$ is a permutation of the $k$-cells with inverse $\sigma\mapsto\sigma-v$. Finally, for $I=\{i_1<\dots<i_k\}$,
\[
\left(\wedge^kr\right)(o(\sigma))=re_{i_1}\wedge\dots\wedge re_{i_k}=\varepsilon_{i_1}\cdots\varepsilon_{i_k}\cdot e_{\pi(i_1)}\wedge\dots\wedge e_{\pi(i_k)},
\]
and $o(r\sigma)$ is the wedge product of the same vectors $e_{\pi(i_1)},\dots,e_{\pi(i_k)}$ taken in increasing order of their indices. Hence the right side equals $\varepsilon(r,\sigma)\cdot o(r\sigma)$ for a sign $\varepsilon(r,\sigma)\in\{1,-1\}$. In particular, for a plaquette $\sigma=(x;ij)$ with $i<j$,
\[
\left(\wedge^2r\right)(e_i\wedge e_j)=\varepsilon_i\cdot\varepsilon_j\cdot e_{\pi(i)}\wedge e_{\pi(j)}=\varepsilon(r,\sigma)\cdot o(r\sigma).
\]

\emph{Orthogonality.} As $\sigma$ ranges over the $k$-cells, so do $\sigma+v$ and $r\sigma$. Hence the sums that define $\langle v_*u,v_*w\rangle_\eta$ and $\langle r_*u,r_*w\rangle_\eta$ may be taken over the cells $\sigma+v$ and $r\sigma$, respectively, and the definitions of $v_*$ and $r_*$ give, with sums over the $k$-cells $\sigma$,
\[
\langle r_*u,r_*w\rangle_\eta=\eta^3\sum_\sigma\varepsilon(r,\sigma)^2\cdot u(\sigma)\cdot w(\sigma)=\langle u,w\rangle_\eta,\qquad\langle v_*u,v_*w\rangle_\eta=\eta^3\sum_\sigma u(\sigma)\cdot w(\sigma)=\langle u,w\rangle_\eta .
\]

\emph{Commutation with the coboundary.} Let $0\le k\le2$ and let $\rho=(x;J)$ be a $(k+1)$-cell, $J=\{j_1<\dots<j_{k+1}\}$. For $1\le p\le k+1$ let $\rho^+_p$ and $\rho^-_p$ be the $k$-cells
\[
\rho^+_p=\left(x+\eta\cdot e_{j_p};J\setminus\{j_p\}\right),\qquad\rho^-_p=\left(x;J\setminus\{j_p\}\right).
\]
For $u\in C^k$ the coboundary is
\[
(du)(\rho)=\sum_{p=1}^{k+1}(-1)^{p-1}\left[u(\rho^+_p)-u(\rho^-_p)\right];
\]
for $k=0,1$ these are the formulas for $d_0$, $d_1$ in Section~\ref{sec:cochains-periodic-lattice}, and for $k=2$ this is $\eta$ times the formula for $d_\eta=\eta^{-1}d$ on $C^2$ in the proof of Lemma~\ref{lem:discrete-codifferentials-three}.

If $n=1$, then $\Lambda$ consists of a single point, so that $\rho^+_p=\rho^-_p$ for every $p$, and $d=0$ on $C^0$, $C^1$ and $C^2$; in this case both sides of $d\,v_*=v_*\,d$ and of $d\,r_*=r_*\,d$ vanish. Let now $n\ge2$. Then $0<\eta<R$, so that $\eta\cdot e_j\notin R\mathbb Z^3$ for $j=1,2,3$. It follows that $\rho^+_p\ne\rho^-_p$, because $x+\eta\cdot e_{j_p}\ne x$ in $\Lambda$; and two of the cells $\rho^\pm_p$ with different indices $p$ have different sets of directions. Thus the $2k+2$ cells $\rho^\pm_p$, $1\le p\le k+1$, are pairwise distinct; they are the \emph{faces} of $\rho$. It also follows that $\frac{\eta}{2}e_j\ne-\frac{\eta}{2}e_j$ in $\mathbb T_R$ for $j=1,2,3$, the difference $\eta\cdot e_j$ of the two points not lying in $R\mathbb Z^3$. A $k$-cell $\sigma$ is a face of $\rho$ exactly when, for some $j\in J$, its directions are $J\setminus\{j\}$ and $c(\sigma)-c(\rho)=\pm\frac{\eta}{2}e_j$ in $\mathbb T_R$; here $j$ and the sign are determined by $\sigma$, the face $\rho^+_p$ corresponding to $j=j_p$ and the sign $+$, the face $\rho^-_p$ to $j=j_p$ and the sign $-$. Hence $(du)(\rho)=\sum_\sigma[\rho:\sigma]\cdot u(\sigma)$, the sum over the faces $\sigma$ of $\rho$, with incidence numbers $[\rho:\rho^\pm_p]=\pm(-1)^{p-1}$. Let $\mathsf n(\rho,\rho^\pm_p)=\pm e_{j_p}$ be the outward unit normal of $\rho$ along the face $\rho^\pm_p$, which is well defined because the faces are pairwise distinct; then $c(\sigma)-c(\rho)=\frac{\eta}{2}\,\mathsf n(\rho,\sigma)$ in $\mathbb T_R$ for every face $\sigma$ of $\rho$. Moving $e_{j_p}$ past the $p-1$ factors $e_{j_1},\dots,e_{j_{p-1}}$ gives $e_{j_p}\wedge o(\rho^\pm_p)=(-1)^{p-1}\cdot o(\rho)$, and therefore
\[
\mathsf n(\rho,\sigma)\wedge o(\sigma)=[\rho:\sigma]\cdot o(\rho)\qquad\text{for every face $\sigma$ of $\rho$}.
\]
This identity determines the incidence numbers. Let $r\in O_{\mathrm{cub}}$ and let $\sigma$ be a face of $\rho$, with directions $J\setminus\{j\}$ and $c(\sigma)-c(\rho)=\pm\frac{\eta}{2}e_j$. Then $r\sigma$ has the directions $\pi(J)\setminus\{\pi(j)\}$, the cell $r\rho$ has the directions $\pi(J)$, and
\[
c(r\sigma)-c(r\rho)=r\left(c(\sigma)-c(\rho)\right)=\pm\frac{\eta}{2}\,\varepsilon_j\cdot e_{\pi(j)}.
\]
So $r\sigma$ is a face of $r\rho$, and its outward unit normal is $\mathsf n(r\rho,r\sigma)=\pm\varepsilon_j\cdot e_{\pi(j)}=r\,\mathsf n(\rho,\sigma)$; as the same holds for $r^{-1}$, the map $\sigma\mapsto r\sigma$ is a bijection from the faces of $\rho$ onto the faces of $r\rho$, and $r$ maps outward normals to outward normals. Apply $\wedge^{k+1}r$ to the identity for $[\rho:\sigma]$. Since $\left(\wedge^{k+1}r\right)(\xi\wedge\omega)=r\xi\wedge\left(\wedge^kr\right)(\omega)$ for $\xi\in\mathbb R^3$ and $\omega\in\wedge^k\mathbb R^3$, the definition of the signs gives
\[
\varepsilon(r,\sigma)\cdot\mathsf n(r\rho,r\sigma)\wedge o(r\sigma)=\varepsilon(r,\rho)\cdot[\rho:\sigma]\cdot o(r\rho).
\]
By the same identity for the face $r\sigma$ of $r\rho$, the left side equals $\varepsilon(r,\sigma)\cdot[r\rho:r\sigma]\cdot o(r\rho)$. Hence, whatever the sign of $\det r$,
\[
[r\rho:r\sigma]=\varepsilon(r,\rho)\cdot\varepsilon(r,\sigma)\cdot[\rho:\sigma].
\]
For $u\in C^k$ it follows, with sums over the faces $\sigma$ of $\rho$, that
\begin{align*}
(d\,r_*u)(r\rho)&=\sum_\sigma[r\rho:r\sigma]\cdot(r_*u)(r\sigma)=\sum_\sigma\varepsilon(r,\rho)\cdot\varepsilon(r,\sigma)^2\cdot[\rho:\sigma]\cdot u(\sigma)\\
&=\varepsilon(r,\rho)\cdot(du)(\rho)=(r_*\,du)(r\rho).
\end{align*}
Every $(k+1)$-cell has the form $r\rho$, so $d\,r_*=r_*\,d$. The translation by $v$ involves no signs: $\rho^\pm_p+v=(\rho+v)^\pm_p$, so that $\sigma\mapsto\sigma+v$ carries the faces of $\rho$ onto the faces of $\rho+v$ with the same outward normals and orientations. Hence $[\rho+v:\sigma+v]=[\rho:\sigma]$, and
\[
(d\,v_*u)(\rho+v)=\sum_\sigma[\rho:\sigma]\cdot u(\sigma)=(du)(\rho)=(v_*\,du)(\rho+v).
\]
Every $(k+1)$-cell has the form $\rho+v$, so $d\,v_*=v_*\,d$.

\emph{Cell averages.} Identify the two-form $\mathrm{d}y_\alpha\wedge\mathrm{d}y_\beta$ with the bivector $e_\alpha\wedge e_\beta\in\wedge^2\mathbb R^3$ for $1\le\alpha<\beta\le3$, that is, the covector $\mathrm{d}y_\alpha$ with the vector $e_\alpha$. Under this identification a linear automorphism of $\mathbb R^3$ carries covectors forward by its inverse transpose, which for the orthogonal matrix $r$ is $r$ itself. Hence the map induced by $r$ on two-forms in the definition of $U_r$ (Section~\ref{sec:test-forms-test}) is the second exterior power $\wedge^2r$ of $r$, extended complex-linearly to complex two-forms. Let $\sigma=(x;ij)$ be a plaquette, $i<j$, let $o(r\sigma)=e_a\wedge e_b$ with $a<b$, and write $f=\sum_{i'<j'}f_{i'j'}\cdot e_{i'}\wedge e_{j'}$. In
\[
(U_rf)(y)=\sum_{i'<j'}f_{i'j'}\left(r^{-1}y\right)\cdot\left(\wedge^2r\right)(e_{i'}\wedge e_{j'})
\]
the bivector $\left(\wedge^2r\right)(e_{i'}\wedge e_{j'})=\pm e_{\pi(i')}\wedge e_{\pi(j')}$ is a multiple of $e_a\wedge e_b$ only for $\{i',j'\}=\{i,j\}$, because $\pi$ is injective. By the identity for plaquettes proved above, the component of $U_rf$ on the cell $r\sigma$ is therefore
\[
(U_rf)_{ab}(y)=\varepsilon(r,\sigma)\cdot f_{ij}\left(r^{-1}y\right)\qquad\left(y\in\mathbb R^3\right),
\]
and the same holds with $f^R$ in place of $f$. The definition \eqref{eq:continuum-infinite-2}, the commutation of periodisation with $U_r$, and the substitution $y=ry'$, which preserves Lebesgue measure on $\mathbb T_R$ and maps $Q(\sigma)$ onto $Q(r\sigma)$ up to a set of measure zero, give
\begin{align*}
(U_rf)_\eta(r\sigma)&=\eta^{-3}\int_{Q(r\sigma)}\mathrm{d}y\,\left(U_rf^R\right)_{ab}(y)=\varepsilon(r,\sigma)\cdot\eta^{-3}\int_{Q(r\sigma)}\mathrm{d}y\,f^R_{ij}\left(r^{-1}y\right)\\
&=\varepsilon(r,\sigma)\cdot\eta^{-3}\int_{Q(\sigma)}\mathrm{d}y'\,f^R_{ij}(y')=\varepsilon(r,\sigma)\cdot f_\eta(\sigma)=(r_*f_\eta)(r\sigma).
\end{align*}
As $\sigma$ ranges over the plaquettes, so does $r\sigma$; hence $(U_rf)_\eta=r_*f_\eta$. The same computation without signs, with $(\tau_vf)^R=\tau_v\left(f^R\right)$ and $Q(\sigma+v)=Q(\sigma)+v$, gives
\[
(\tau_vf)_\eta(\sigma+v)=\eta^{-3}\int_{Q(\sigma)+v}\mathrm{d}y\,f^R_{ij}(y-v)=\eta^{-3}\int_{Q(\sigma)}\mathrm{d}y\,f^R_{ij}(y)=f_\eta(\sigma)=(v_*f_\eta)(\sigma+v),
\]
that is, $(\tau_vf)_\eta=v_*f_\eta$.

\emph{Reflections.} For $r=r_\nu$ the permutation $\pi$ is the identity, $\varepsilon_\nu=-1$, and $\varepsilon_j=1$ for $j\ne\nu$. So $r_\nu\sigma$ has the same directions as $\sigma$, and $\varepsilon(r_\nu,\sigma)=-1$ exactly when $\sigma$ contains the direction $\nu$, that is, exactly when $r_\nu$ reverses the orientation of $\sigma$. Since $r_\nu$ is an involution, for $g$ in $C^2$ or in its complexification
\[
\left((r_\nu)_*g\right)(\sigma)=\varepsilon(r_\nu,r_\nu\sigma)\cdot g(r_\nu\sigma),
\]
where $\varepsilon(r_\nu,r_\nu\sigma)=-1$ exactly when $\sigma$ contains the direction $\nu$. By the definition of $\hat{\mathsf R}_\nu$ in Section~\ref{sec:admissible-families-lattice}, and by the same formula on the complexification of $C^2$, the right side is $(\hat{\mathsf R}_\nu g)(\sigma)$. Hence $\hat{\mathsf R}_\nu=(r_\nu)_*$ on $C^2$ and on its complexification, and, as $\mathsf R_\nu=U_{r_\nu}$, the identity for $U_r$ gives $(\mathsf R_\nu f)_\eta=(r_\nu)_*f_\eta=\hat{\mathsf R}_\nu f_\eta$.
\end{proof}

\begin{lemma}[Invariances of the averaged law]\label{lem:invariances-averaged-law}
Let $(N,T)$ be a cutoff, $q\ge0$ and $\boldsymbol{\mu}\in\mathbb R$. Recall from Section~\ref{sec:characters-twisted-spinors} that a character is a triple $\chi=(\chi_1,\chi_2,\chi_3)\in\mathbb C^3$ with $|\chi_1|=|\chi_2|=|\chi_3|=1$, and that $\chi^{\mathrm{ap}}=(-1,-1,-1)$. For a character $\chi$ and $\alpha\in\mathbb R^3$ let $\chi\exp(i\alpha)$ denote the character $\left(\chi_1\exp(i\alpha_1),\chi_2\exp(i\alpha_2),\chi_3\exp(i\alpha_3)\right)$. For a character $\chi$, $A\in C^1$ and $\theta\in\mathbb R^3$, let $D^\theta_A[\chi]$ be the operator on $\ell^2(\Lambda;\mathbb C^4)$ obtained from the operator $D^\theta_A$ of Section~\ref{sec:wilson-operator-lattice} as follows: for $\nu\in\{1,2,3\}$ and every bond $b$ in the direction $\nu$ that crosses the seam, the factor $\varsigma(b)=-1$ of \eqref{eq:continuum-infinite-1} is replaced by $\chi_\nu$ if $b=(x\to x+e_\nu)$ and by $\bar\chi_\nu$ if $b=(x\to x-e_\nu)$. Since $\chi_\nu=\bar\chi_\nu=-1$ for $\chi=\chi^{\mathrm{ap}}$, one has $D^\theta_A[\chi^{\mathrm{ap}}]=D^\theta_A$.

As in Section~\ref{sec:wilson-operator-lattice}, where $\chi=\chi^{\mathrm{ap}}$, restriction to the box $\{-m,\dots,m\}^3$ maps the space $\mathcal S_\chi$ of Section~\ref{sec:characters-twisted-spinors}, formed with $s=4$, bijectively onto $\ell^2(\Lambda;\mathbb C^4)$, and it transforms the operator \eqref{eq:gaussian-bounds-1} with the character $\chi$, formed with the same $q$ and $A$ and with the matrices $\gamma_1,\gamma_2,\gamma_3$ of Section~\ref{sec:wilson-operator-lattice}, into $D^0_A[\chi]$. Compared with the case $\chi=\chi^{\mathrm{ap}}$, the only change concerns the factors acquired by the covariant shifts $T_\nu$ and $T^*_\nu$ along the bonds $(x\to x+e_\nu)$ and $(x\to x-e_\nu)$, respectively, that cross the seam: these factors are $\chi_\nu$ and $\chi_\nu^{-1}=\bar\chi_\nu$, because such a shift evaluates a spinor at the final point of the bond plus $n\cdot e_\nu$, respectively minus $n\cdot e_\nu$, and every $\Psi\in\mathcal S_\chi$ satisfies $\Psi(y\pm n\cdot e_\nu)=\chi_\nu^{\pm1}\cdot\Psi(y)$ for $y\in\mathbb Z^3$. Moreover, $D^\theta_A[\chi]$ arises from $D^0_A[\chi]$ when the term of each bond $b=(x\to x\pm e_\nu)$ is multiplied by $\exp\left(\pm\frac{i\theta_\nu}{n}\right)$, as $D^\theta_A$ arises from $D_A$.
\begin{enumerate}
\item[(a)] For every character $\chi$, the law $\bar{\mathsf P}$ is unchanged if $D^\theta_A$ is replaced by $D^\theta_A[\chi]$ in the definitions of Section~\ref{sec:wilson-operator-lattice}. More precisely, for all $A\in C^1$ and all $\theta,\alpha\in\mathbb R^3$,
\[
\det\left(D^\theta_A\left[\chi^{\mathrm{ap}}\exp(i\alpha)\right]+\boldsymbol{\mu}\right)=\det\left(D^{\theta+\alpha}_A+\boldsymbol{\mu}\right),
\]
where the operator on the left carries the seam factors $-\exp(i\alpha_\nu)$ on the bonds $(x\to x+e_\nu)$ and $-\exp(-i\alpha_\nu)$ on the bonds $(x\to x-e_\nu)$ that cross the seam. In particular, the function $\theta\mapsto\det\left(D^\theta_A+\boldsymbol{\mu}\right)$ is $2\pi$-periodic in each of the variables $\theta_1,\theta_2,\theta_3$.
\item[(b)] $\bar{\mathsf P}$ is invariant under $v_*$ for every $v\in\eta\mathbb Z^3$ and under $r_*$ for every $r\in O_{\mathrm{cub}}$.
\item[(c)] $\bar{\mathsf P}$ is invariant under the map $F\mapsto-F$ of $C^2(\Lambda)$.
\item[(d)] Let $\mathrm{Pr}$ be a probability measure on $C^2(\Lambda)$ with property~\eqref{eq:continuum-infinite-Ad1} of Definition~\ref{def:admissible-families}, that is, invariant under $v_*$ for every $v\in\eta\mathbb Z^3$ and under $r_*$ for every $r\in O_{\mathrm{cub}}$. For $f\in\mathcal S(\mathbb R^3;\Lambda^2)$ let $X(f)$ denote the function $F\mapsto X_{\eta,R}(f)=\langle f_\eta,F\rangle_\eta$ on $C^2(\Lambda)$, see \eqref{eq:continuum-infinite-2}. Let $s\ge1$ be an integer, let $f_1,\dots,f_s\in\mathcal S(\mathbb R^3;\Lambda^2)$, $v\in\eta\mathbb Z^3$ and $r\in O_{\mathrm{cub}}$. Let $\mathbb K=\mathbb C$, or $\mathbb K=\mathbb R$ if $f_1,\dots,f_s$ are real, and let $\Phi\colon\mathbb K^s\to[0,\infty)$ be a Borel function or $\Phi\colon\mathbb K^s\to\mathbb C$ a bounded Borel function. Then
\begin{align*}
\int\mathrm{d}\mathrm{Pr}\;\Phi\left(X(\tau_vf_1),\dots,X(\tau_vf_s)\right)&=\int\mathrm{d}\mathrm{Pr}\;\Phi\left(X(f_1),\dots,X(f_s)\right)\\
&=\int\mathrm{d}\mathrm{Pr}\;\Phi\left(X(U_rf_1),\dots,X(U_rf_s)\right).
\end{align*}
\end{enumerate}
\end{lemma}

\begin{proof}
As in \eqref{eq:continuum-infinite-1}, sites are written in unit coordinates: $x\in\{-m,\dots,m\}^3$ stands for the site $\eta x$ of $\Lambda$, a bond is $b=(x\to x\pm e_\nu)$, and its endpoint is reduced modulo $n$ to $\{-m,\dots,m\}^3$. A bond in the direction $\nu$ crosses the seam if the $\nu$-coordinates of its two endpoints are $m$ and $-m$. In $D^\theta_A[\chi]$ the bond $b=(x\to x\pm e_\nu)$ enters with the matrix $-\tfrac12\left(1-\gamma(b)\right)$, multiplied by its seam factor (equal to $1$ if $b$ does not cross the seam), by $\exp\left(iqA(b)\right)$ and by its twist factor $\exp\left(\pm\frac{i\theta_\nu}{n}\right)$.

\emph{(a)} For $\alpha\in\mathbb R^3$ let $g_\alpha$ be the unitary operator on $\ell^2(\Lambda;\mathbb C^4)$ of multiplication by the function
\[
g_\alpha(x)=\exp\left(\frac in\sum_{\nu=1}^3\alpha_\nu x_\nu\right),\qquad x\in\{-m,\dots,m\}^3.
\]
Let $x\in\{-m,\dots,m\}^3$ and $\nu\in\{1,2,3\}$. If $x_\nu<m$, then $g_\alpha(x+e_\nu)/g_\alpha(x)=\exp\left(\frac{i\alpha_\nu}{n}\right)$. If $x_\nu=m$, the bond $(x\to x+e_\nu)$ crosses the seam, the site $x+e_\nu$ has the $\nu$-coordinate $-m$, and, because $2m=n-1$,
\[
\frac{g_\alpha(x+e_\nu)}{g_\alpha(x)}=\exp\left(-\frac{2im\alpha_\nu}{n}\right)=\exp\left(\frac{i\alpha_\nu}{n}\right)\cdot\exp(-i\alpha_\nu).
\]
In $g_\alpha^{-1}D^\theta_A[\chi]g_\alpha$ the term of each bond $b=(x\to y)$ is that of $D^\theta_A[\chi]$ multiplied by $g_\alpha(y)/g_\alpha(x)$. By the two ratios just computed, and by their inverses, which belong to the bonds $(x\to x-e_\nu)$, the twist factor $\exp\left(\pm\frac{i\theta_\nu}{n}\right)$ of every bond is thereby multiplied by $\exp\left(\pm\frac{i\alpha_\nu}{n}\right)$, and on the bonds that cross the seam the seam factors $\chi_\nu$ and $\bar\chi_\nu$ are in addition multiplied by $\exp(-i\alpha_\nu)$ and by $\exp(i\alpha_\nu)$, respectively. Hence
\[
g_\alpha^{-1}D^\theta_A[\chi]g_\alpha=D^{\theta+\alpha}_A\left[\chi\exp(-i\alpha)\right]
\]
for every character $\chi$, every $A\in C^1$ and all $\theta,\alpha\in\mathbb R^3$. For $\chi=\chi^{\mathrm{ap}}\exp(i\alpha)$ the right side is $D^{\theta+\alpha}_A[\chi^{\mathrm{ap}}]=D^{\theta+\alpha}_A$, and $g_\alpha$ commutes with the multiple $\boldsymbol{\mu}$ of the identity; this gives the identity of determinants in (a). For $\alpha=2\pi e_\nu$ one has $\chi^{\mathrm{ap}}\exp(i\alpha)=\chi^{\mathrm{ap}}$, so that the seam factors do not change, and the identity reads $\det\left(D^\theta_A+\boldsymbol{\mu}\right)=\det\left(D^{\theta+2\pi e_\nu}_A+\boldsymbol{\mu}\right)$; this is the periodicity.

Every character can be written as $\chi=\chi^{\mathrm{ap}}\exp(i\alpha)$ with $\alpha\in\mathbb R^3$: choose $\alpha_\nu$ with $\exp(i\alpha_\nu)=-\chi_\nu$. Let $\bar w[\chi]$ be the function $\bar w$ of Section~\ref{sec:wilson-operator-lattice} formed with $D^\theta_A[\chi]$ in place of $D^\theta_A$. For $A\in V_{\mathrm{t}}$ the identity of determinants and the substitution $\theta'=\theta+\alpha$ give
\begin{align*}
\bar w[\chi](A)&=\frac{1}{(2\pi)^3}\int_{[-\pi,\pi]^3}\mathrm{d}\theta\,\exp\left(-\tfrac12\|dA\|^2\right)\cdot\det\left(D^\theta_A[\chi]+\boldsymbol{\mu}\right)\\
&=\frac{1}{(2\pi)^3}\int_{\alpha+[-\pi,\pi]^3}\mathrm{d}\theta'\,\exp\left(-\tfrac12\|dA\|^2\right)\cdot\det\left(D^{\theta'}_A+\boldsymbol{\mu}\right)=\bar w(A),
\end{align*}
where the last equality holds because the integrand is $2\pi$-periodic in each of $\theta'_1,\theta'_2,\theta'_3$, so that its integral over the translate $\alpha+[-\pi,\pi]^3$ of $[-\pi,\pi]^3$ equals its integral over $[-\pi,\pi]^3$. Thus $\bar w$ does not depend on the character, and neither do $\bar Z=\int_{V_{\mathrm{t}}}\mathrm{d}A\,\bar w(A)$ and $\bar{\mathsf P}$.

\emph{(b)} Put $\chi^{\mathrm{per}}=(1,1,1)$. By (a) and its proof, $\bar w$ and $\bar{\mathsf P}$ do not change if $D^\theta_A$ is replaced by $D^\theta_A[\chi^{\mathrm{per}}]$, the operator \eqref{eq:continuum-infinite-1} with the twist $\theta$ and with all seam factors equal to $1$. In this part of the proof $\bar w$ is formed with $D^\theta_A[\chi^{\mathrm{per}}]$.

Let $r\in O_{\mathrm{cub}}$. There are a permutation $\pi$ of $\{1,2,3\}$ and signs $\varepsilon_1,\varepsilon_2,\varepsilon_3\in\{1,-1\}$ with $re_\nu=\varepsilon_\nu e_{\pi(\nu)}$ for $\nu=1,2,3$; thus the entries of $r$ are $r_{\lambda\nu}=\varepsilon_\nu$ if $\lambda=\pi(\nu)$ and $r_{\lambda\nu}=0$ otherwise. Write $\mathbb C^4=\mathbb C^2\otimes\mathbb C^2$ as in Section~\ref{sec:wilson-operator-lattice}, so that $\gamma_\nu=\tau_3\otimes\sigma_\nu$ for $\nu=1,2,3$ and $\gamma_4=\tau_1\otimes1$. Suppose first that $\det r=1$. The map which assigns to $u\in\mathrm{SU}(2)$ the matrix, in the basis $\sigma_1,\sigma_2,\sigma_3$, of the linear map $X\mapsto uXu^{-1}$ of the real span of $\sigma_1,\sigma_2,\sigma_3$ onto itself is the two-fold covering $\mathrm{SU}(2)\to\mathrm{SO}(3)$. It is onto; so there is $u_r\in\mathrm{SU}(2)$ with
\[
u_r\sigma_\nu u_r^{-1}=\sum_{\lambda=1}^3r_{\lambda\nu}\sigma_\lambda,\qquad\nu=1,2,3,
\]
and one puts $V_r=1\otimes u_r$. If $\det r=-1$, then $r'=-r\in\mathrm{SO}(3)\cap O_{\mathrm{cub}}$, and one puts $V_r=\gamma_4V_{r'}$. Since $\gamma_4$ is unitary, $\gamma_4^2=1$, and $\gamma_4$ anticommutes with $\gamma_1,\gamma_2,\gamma_3$, in both cases $V_r$ is a unitary $4\times4$ matrix with
\[
V_r\gamma_\nu V_r^{-1}=\sum_{\lambda=1}^3r_{\lambda\nu}\gamma_\lambda,\qquad\nu=1,2,3.
\]
The map $x\mapsto rx$ is a bijection of $\{-m,\dots,m\}^3$ onto itself. For a bond $b=(x\to x\pm e_\nu)$ let $rb=\left(rx\to rx\pm\varepsilon_\nu e_{\pi(\nu)}\right)$ be its image. Then $\gamma(b)=\pm\gamma_\nu$ and $\gamma(rb)=\pm\varepsilon_\nu\gamma_{\pi(\nu)}=\pm\sum_\lambda r_{\lambda\nu}\gamma_\lambda$, so that the last display says
\[
V_r\gamma(b)V_r^{-1}=\gamma(rb).
\]
Let $\mathsf V$ be the unitary operator on $\ell^2(\Lambda;\mathbb C^4)$ defined by $(\mathsf V\psi)(rx)=V_r\psi(x)$. The map $b\mapsto rb$ is a bijection from the bonds that start at $x$ onto the bonds that start at $rx$. For every oriented bond $b$ one has $(r_*A)(rb)=A(b)$. Indeed, let $b=(x\to x+e_\nu)$ and let $\sigma$ be its $1$-cell, oriented by $e_\nu$. By Lemma~\ref{lem:lattice-symmetries-cell}, $re_\nu=\varepsilon(r,\sigma)\cdot o(r\sigma)$, so the bond $rb=(rx\to rx+re_\nu)$ is oriented like the cell $r\sigma$ if $\varepsilon(r,\sigma)=1$ and oppositely if $\varepsilon(r,\sigma)=-1$; hence
\[
(r_*A)(rb)=\varepsilon(r,\sigma)\cdot(r_*A)(r\sigma)=\varepsilon(r,\sigma)^2\cdot A(\sigma)=A(b),
\]
and both sides change sign when $b$ is reversed. Moreover the twist factor $\exp\left(\pm\frac{i\theta_\nu}{n}\right)$ of $b=(x\to x\pm e_\nu)$ is the twist factor of $rb$ for the twist $\theta'\in\mathbb R^3$ defined by
\[
\theta'_{\pi(\nu)}=\varepsilon_\nu\theta_\nu,\qquad\nu=1,2,3.
\]
Consequently, for $\psi\in\ell^2(\Lambda;\mathbb C^4)$, the term of $b$ in $\left(D^\theta_A[\chi^{\mathrm{per}}]\mathsf V^{-1}\psi\right)(x)$, multiplied on the left by $V_r$, is the term of $rb$ in $\left(D^{\theta'}_{r_*A}[\chi^{\mathrm{per}}]\psi\right)(rx)$, and \eqref{eq:continuum-infinite-1} gives
\[
\mathsf V\left(D^\theta_A[\chi^{\mathrm{per}}]+\boldsymbol{\mu}\right)\mathsf V^{-1}=D^{\theta'}_{r_*A}[\chi^{\mathrm{per}}]+\boldsymbol{\mu}.
\]
The map $\theta\mapsto\theta'$ permutes the coordinates up to signs; it maps $[-\pi,\pi]^3$ onto itself and preserves $\mathrm{d}\theta$. By Lemma~\ref{lem:lattice-symmetries-cell}, $r_*$ is orthogonal on $C^1$ and commutes with $d$. Hence $r_*$ maps $V_{\mathrm{g}}=d_0(C^0)$ onto itself. It maps $V_{\mathrm{h}}$ onto itself, because it carries the constant cochain of the direction $\nu$ (equal to $1$ on the $1$-cells of the direction $\nu$ and to $0$ on the others) to $\varepsilon_\nu$ times the constant cochain of the direction $\pi(\nu)$. Being orthogonal, $r_*$ therefore maps $V_{\mathrm{t}}=(V_{\mathrm{g}}\oplus V_{\mathrm{h}})^\perp$ onto itself and preserves the Lebesgue measure of $V_{\mathrm{t}}$, and $\|dr_*A\|=\|r_*dA\|=\|dA\|$. For $A\in V_{\mathrm{t}}$, the substitution $\theta\mapsto\theta'$ and the conjugation identity give
\[
\bar w(r_*A)=\frac{1}{(2\pi)^3}\int_{[-\pi,\pi]^3}\mathrm{d}\theta\,\exp\left(-\tfrac12\|dA\|^2\right)\cdot\det\left(D^{\theta'}_{r_*A}[\chi^{\mathrm{per}}]+\boldsymbol{\mu}\right)=\bar w(A),
\]
and $F(r_*A)=\eta^{-\frac32}\,dr_*A=r_*F(A)$. So $r_*$ maps the probability measure $\bar Z^{-1}\bar w(A)\,\mathrm{d}A$ on $V_{\mathrm{t}}$ onto itself, and, as $F(r_*A)=r_*F(A)$, the law of $r_*F$ under $\bar{\mathsf P}$ is $\bar{\mathsf P}$.

For $v\in\eta\mathbb Z^3$ the same argument applies with $V_r=1$, that is, with $\mathsf V$ the translation $(\mathsf V\psi)(x+\eta^{-1}v)=\psi(x)$, with $v_*$ in place of $r_*$, and with $\theta'=\theta$.

\emph{(c)} Here $D^\theta_A$ is the operator \eqref{eq:continuum-infinite-1} with the twist $\theta$ and the seam factors $\varsigma$. In $D^\theta_A$ the matrix $-\tfrac12\left(1-\gamma(b)\right)$ of an oriented bond $b=(x\to x\pm e_\nu)$ is multiplied by three scalars, and their product is
\[
u(b)=\varsigma(b)\cdot\exp\left(iqA(b)\right)\cdot\exp\left(\pm\frac{i\theta_\nu}{n}\right).
\]
The reversed bond $-b=(x\pm e_\nu\to x)$ has $\varsigma(-b)=\varsigma(b)\in\{1,-1\}$, $A(-b)=-A(b)$ and the twist factor $\exp\left(\mp\frac{i\theta_\nu}{n}\right)$; hence $u(-b)=u(b)^{-1}$. For an operator $M$ on $\ell^2(\Lambda;\mathbb C^4)$ and sites $x,y\in\{-m,\dots,m\}^3$ let $M_{xy}$ be the $4\times4$ matrix such that $(M\psi)(x)=\sum_yM_{xy}\cdot\psi(y)$ for every $\psi$; the transposed operator $M^T$, whose matrix is the transpose of that of $M$, has the blocks $\left(M^T\right)_{yx}=\left(M_{xy}\right)^T$. The final point of a bond $(x\to x\pm e_\nu)$ differs from $x$ only in the $\nu$-coordinate, by $\pm1$ modulo $n$, and the residues of $1$, $-1$ and $0$ modulo $n$ are distinct because $n\ge3$. Hence the six oriented bonds with initial point $x$ end at six distinct sites, none of which is $x$: for given sites $x$ and $y$ there is at most one oriented bond from $x$ to $y$, and none if $y=x$. By \eqref{eq:continuum-infinite-1} it follows that $\left(D^\theta_A\right)_{xx}=3\cdot1$ for every site $x$, that $\left(D^\theta_A\right)_{xy}=-\tfrac12\left(1-\gamma(b)\right)\cdot u(b)$ whenever $b=(x\to y)$ is an oriented bond, and that $\left(D^\theta_A\right)_{xy}=0$ for all other pairs of sites. Transposing, one obtains that the diagonal blocks of $\left(D^\theta_A\right)^T$ are again $3\cdot1$, that for every oriented bond $b=(x\to y)$
\[
\left(\left(D^\theta_A\right)^T\right)_{yx}=-\tfrac12\left(1-\gamma(b)\right)^T\cdot u(b),
\]
the indices $y$ and $x$ being the initial and the final point of the reversed bond $-b=(y\to x)$, and that all other blocks of $\left(D^\theta_A\right)^T$ vanish.

Let $\mathsf C$ be the operator on $\ell^2(\Lambda;\mathbb C^4)$ which acts at every site by the matrix $\operatorname{diag}(\sigma_2,\sigma_2)$. The Pauli matrices satisfy
\[
\sigma_2\sigma_\nu^T\sigma_2=-\sigma_\nu,\qquad\nu=1,2,3,
\]
since $\sigma_1$ and $\sigma_3$ are symmetric and anticommute with $\sigma_2$, while $\sigma_2$ is antisymmetric; moreover $\sigma_2^2=1$. As $\gamma_\nu=\operatorname{diag}(\sigma_\nu,-\sigma_\nu)$, it follows that $\mathsf C\gamma_\nu^T\mathsf C^{-1}=-\gamma_\nu$, hence
\[
\mathsf C\left(1-\gamma(b)\right)^T\mathsf C^{-1}=1+\gamma(b)=1-\gamma(-b).
\]
Now let $b=(x\to x\pm e_\nu)$ be an oriented bond with final point $y$. The operator $D^{-\theta}_{-A}$ is formed with $-A$ and $-\theta$ in place of $A$ and $\theta$, and the reversed bond $-b=(y\to x)$ is the bond $(y\to y\mp e_\nu)$ of the direction $\nu$, whose sign is opposite to that of $b$; so in $D^{-\theta}_{-A}$ the matrix $-\tfrac12\left(1-\gamma(-b)\right)$ of $-b$ is multiplied by
\[
\varsigma(-b)\cdot\exp\left(iq\cdot(-A)(-b)\right)\cdot\exp\left(\mp\frac{i\cdot(-\theta_\nu)}{n}\right)=\varsigma(b)\cdot\exp\left(-iqA(-b)\right)\cdot\exp\left(\pm\frac{i\theta_\nu}{n}\right)=u(b).
\]
On the other hand, by the identity for $\mathsf C\left(1-\gamma(b)\right)^T\mathsf C^{-1}$ above, the block of $\mathsf C\left(D^\theta_A\right)^T\mathsf C^{-1}$ with the indices $y$, $x$ is
\[
\operatorname{diag}(\sigma_2,\sigma_2)\cdot\left(\left(D^\theta_A\right)^T\right)_{yx}\cdot\operatorname{diag}(\sigma_2,\sigma_2)^{-1}=-\tfrac12\left(1-\gamma(-b)\right)\cdot u(b).
\]
So the two operators have the same block with the indices of $-b$, for every oriented bond $b$. By the description of the blocks above, applied to $D^\theta_A$ and to $D^{-\theta}_{-A}$, both operators have the diagonal blocks $3\cdot1$, since $\operatorname{diag}(\sigma_2,\sigma_2)$ commutes with $3\cdot1$, and both have zero blocks at all remaining pairs of sites, because every oriented bond is the reversal of an oriented bond. Therefore
\[
\mathsf C\left(D^\theta_A\right)^T\mathsf C^{-1}=D^{-\theta}_{-A},
\]
so that $\det\left(D^{-\theta}_{-A}+\boldsymbol{\mu}\right)=\det\left(D^\theta_A+\boldsymbol{\mu}\right)$ and, since $\|d(-A)\|=\|dA\|$, $w_{-\theta}(-A)=w_\theta(A)$ for $A\in V_{\mathrm{t}}$. The map $\theta\mapsto-\theta$ maps $[-\pi,\pi]^3$ onto itself and preserves $\mathrm{d}\theta$; the map $A\mapsto-A$ maps $V_{\mathrm{t}}$ onto itself and preserves its Lebesgue measure and $\|dA\|$. Hence $\bar w(-A)=\bar w(A)$, and $F(-A)=-F(A)$. As at the end of the proof of (b), the law of $-F$ under $\bar{\mathsf P}$ is $\bar{\mathsf P}$.

\emph{(d)} If $f_1,\dots,f_s$ are real, then so are the tests $\tau_vf_i$ and $U_rf_i$, and $X(g)$ is real-valued for every real test $g$, because $g_\eta$ is then real and $C^2(\Lambda)$ is a real vector space; so in either case every argument of $\Phi$ in the assertion lies in $\mathbb K^s$. By Lemma~\ref{lem:lattice-symmetries-cell}, $(U_rf)_\eta=r_*f_\eta$ and $(\tau_vf)_\eta=v_*f_\eta$ for every $f\in\mathcal S(\mathbb R^3;\Lambda^2)$, and $r_*$, $v_*$ are orthogonal. Hence, for every $F\in C^2(\Lambda)$,
\[
X(U_rf)(r_*F)=\langle r_*f_\eta,r_*F\rangle_\eta=X(f)(F),\qquad X(\tau_vf)(v_*F)=\langle v_*f_\eta,v_*F\rangle_\eta=X(f)(F).
\]
Since $\mathrm{Pr}$ is invariant under $r_*$,
\begin{align*}
\int\mathrm{d}\mathrm{Pr}(F)\;\Phi\left(X(U_rf_1)(F),\dots,X(U_rf_s)(F)\right)&=\int\mathrm{d}\mathrm{Pr}(F)\;\Phi\left(X(U_rf_1)(r_*F),\dots,X(U_rf_s)(r_*F)\right)\\
&=\int\mathrm{d}\mathrm{Pr}(F)\;\Phi\left(X(f_1)(F),\dots,X(f_s)(F)\right),
\end{align*}
and the same computation with $\tau_v$ and $v_*$ in place of $U_r$ and $r_*$, using the invariance of $\mathrm{Pr}$ under $v_*$, gives the remaining equality.
\end{proof}

\begin{lemma}[Discrete codifferentials of three-forms]\label{lem:discrete-codifferentials-three}
Let $(N,T)$ be a cutoff, with $\eta$, $R$ and $\Lambda=\Lambda_{\eta,R}$ as in Section~\ref{sec:tori-cells-cochains}, and let $g\in\mathcal S(\mathbb R^3)$, regarded as the three-form $g\,\mathrm{d}y_1\wedge\mathrm{d}y_2\wedge\mathrm{d}y_3$. For $j\in\{1,2,3\}$ let $\delta_{\eta,j}g\in\mathcal S(\mathbb R^3)$ be the centred difference quotient
\[
(\delta_{\eta,j}g)(y)=\eta^{-1}\left(g\left(y+\frac{\eta}{2}e_j\right)-g\left(y-\frac{\eta}{2}e_j\right)\right)\qquad(y\in\mathbb R^3),
\]
and let $u^{(\eta)}$ and $d^*g$ be the elements of $\mathcal S(\mathbb R^3;\Lambda^2)$ whose components, listed in the order $23,13,12$, are
\begin{align*}
\left(u^{(\eta)}_{23},u^{(\eta)}_{13},u^{(\eta)}_{12}\right)&=\left(-\delta_{\eta,1}g,\ \delta_{\eta,2}g,\ -\delta_{\eta,3}g\right),\\
\left((d^*g)_{23},(d^*g)_{13},(d^*g)_{12}\right)&=\left(-\partial_1g,\ \partial_2g,\ -\partial_3g\right).
\end{align*}
Then $X_{\eta,R}(u^{(\eta)})$ vanishes on the closed two-cochains: $X_{\eta,R}(u^{(\eta)})(F)=\left\langle(u^{(\eta)})_\eta,F\right\rangle_\eta=0$ for every $F\in\ker d\subset C^2(\Lambda)$. Moreover
\[
\left\|u^{(\eta)}-d^*g\right\|_{L^2}\le\frac{\eta^2}{24}\sum_{j=1}^3\left\|\partial_j^3g\right\|_{L^2}.
\]
\end{lemma}

\begin{proof}
Let $g^R(y)=\sum_{k\in\mathbb Z^3}g(y+R\cdot k)$ be the periodisation of $g$, formed as for two-forms in Section~\ref{sec:test-forms-test}. Write $(x;123)$ for the three-cell $(x;\{1,2,3\})$, and let $g_\eta$ be the function on the three-cells given by
\[
g_\eta(x;123)=\eta^{-3}\int_{Q(x;123)}\mathrm{d}y\,g^R(y)\qquad(x\in\Lambda).
\]
For $F\in C^2(\Lambda)$ and $x\in\Lambda$ write $F_{\mu\nu}(x)=F(x;\mu\nu)$. On three-cells the coboundary of $F$ is
\begin{align*}
(d_\eta F)(x;123)&=\eta^{-1}\cdot(dF)(x;123)\\
&=\eta^{-1}\left(F_{23}(x+\eta e_1)-F_{23}(x)-F_{13}(x+\eta e_2)+F_{13}(x)+F_{12}(x+\eta e_3)-F_{12}(x)\right).
\end{align*}
From now on $\{\mu,\nu,j\}=\{1,2,3\}$ and $\mu<\nu$. In this notation $u^{(\eta)}_{\mu\nu}=(-1)^j\delta_{\eta,j}g$ and $(d^*g)_{\mu\nu}=(-1)^j\partial_jg$, and the part of $(d_\eta F)(x;123)$ that contains $F_{\mu\nu}$ is $(-1)^{j+1}\eta^{-1}\left(F_{\mu\nu}(x+\eta e_j)-F_{\mu\nu}(x)\right)$. For functions $a$, $b$ on $\Lambda$ one has
\[
\sum_{x\in\Lambda}a(x)\cdot\left(b(x+\eta e_j)-b(x)\right)=\sum_{x\in\Lambda}\left(a(x-\eta e_j)-a(x)\right)\cdot b(x),
\]
since $x\mapsto x-\eta e_j$ is a bijection of $\Lambda$. Applying this summation by parts to the three parts of $\langle g_\eta,d_\eta F\rangle_\eta$, one obtains the values of $d_\eta^*g_\eta$, where $d_\eta^*$ is the adjoint of $d_\eta$ with respect to $\langle\cdot,\cdot\rangle_\eta$:
\[
(d_\eta^*g_\eta)(x;\mu\nu)=(-1)^j\eta^{-1}\left(g_\eta(x;123)-g_\eta(x-\eta e_j;123)\right)\qquad(x\in\Lambda).
\]
The centres of Section~\ref{sec:tori-cells-cochains} satisfy $c(x;123)=c(x;\mu\nu)+\frac{\eta}{2}e_j$. Hence $Q(x;123)=Q(x;\mu\nu)+\frac{\eta}{2}e_j$, and $Q(x-\eta e_j;123)$ coincides with $Q(x;\mu\nu)-\frac{\eta}{2}e_j$ up to a translation by a vector of $R\mathbb Z^3$, which does not change integrals of the $R$-periodic function $g^R$. The integral of a function over a translated cube is the integral of the translated function over the cube, and periodisation commutes with translations, so that $\delta_{\eta,j}(g^R)=(\delta_{\eta,j}g)^R$. Therefore
\begin{align*}
\eta^{-1}\left(g_\eta(x;123)-g_\eta(x-\eta e_j;123)\right)&=\eta^{-3}\int_{Q(x;\mu\nu)}\mathrm{d}y\,\eta^{-1}\left(g^R\left(y+\frac{\eta}{2}e_j\right)-g^R\left(y-\frac{\eta}{2}e_j\right)\right)\\
&=\eta^{-3}\int_{Q(x;\mu\nu)}\mathrm{d}y\,(\delta_{\eta,j}g)^R(y).
\end{align*}
After multiplication by $(-1)^j$ the left side becomes $(d_\eta^*g_\eta)(x;\mu\nu)$, and by \eqref{eq:continuum-infinite-2} the right side becomes $(u^{(\eta)})_\eta(x;\mu\nu)$. Thus $(u^{(\eta)})_\eta=d_\eta^*g_\eta$, and for $F\in\ker d$
\[
X_{\eta,R}(u^{(\eta)})(F)=\left\langle(u^{(\eta)})_\eta,F\right\rangle_\eta=\left\langle d_\eta^*g_\eta,F\right\rangle_\eta=\left\langle g_\eta,d_\eta F\right\rangle_\eta=0,
\]
because $d_\eta F=\eta^{-1}\cdot dF=0$.

For the estimate let $\hat h(p)=\int_{\mathbb R^3}\mathrm{d}y\,\exp(-ip\cdot y)\,h(y)$ be the Fourier transform of $h\in\mathcal S(\mathbb R^3)$. Then
\[
\widehat{\delta_{\eta,j}h}(p)=\frac{2i}{\eta}\sin\left(\frac{\eta p_j}{2}\right)\cdot\hat h(p),\qquad\widehat{\partial_jh}(p)=ip_j\,\hat h(p),\qquad\widehat{\partial_j^3h}(p)=-ip_j^3\,\hat h(p).
\]
So $\delta_{\eta,j}-\partial_j$ acts as multiplication by $i\left(\frac{2}{\eta}\sin\left(\frac{\eta p_j}{2}\right)-p_j\right)$, a function of modulus at most $\frac{\eta^2|p_j|^3}{24}$, because $|\sin s-s|\le\frac{|s|^3}{6}$ for real $s$. By Plancherel's theorem, $\left\|(\delta_{\eta,j}-\partial_j)g\right\|_{L^2}\le\frac{\eta^2}{24}\left\|\partial_j^3g\right\|_{L^2}$. The components of $u^{(\eta)}-d^*g$ are the functions $(-1)^j(\delta_{\eta,j}-\partial_j)g$, and therefore
\begin{align*}
\left\|u^{(\eta)}-d^*g\right\|_{L^2}&=\left(\sum_{j=1}^3\left\|(\delta_{\eta,j}-\partial_j)g\right\|_{L^2}^2\right)^{\frac12}\\
&\le\sum_{j=1}^3\left\|(\delta_{\eta,j}-\partial_j)g\right\|_{L^2}\le\frac{\eta^2}{24}\sum_{j=1}^3\left\|\partial_j^3g\right\|_{L^2}.
\end{align*}
\end{proof}

\begin{lemma}[Tests supported in a half-space]\label{lem:tests-supported-half}
Let $(N,T)$ be a cutoff, with $\eta$ and $R$ as in Section~\ref{sec:tori-cells-cochains}. Let $\nu\in\{1,2,3\}$, let $\rho$ and $\delta>0$ be real numbers, and let $f\in C_c^\infty(\mathbb R^3;\Lambda^2)$ satisfy
\[
\operatorname{supp}f\subset\left\{y\in\mathbb R^3:y_\nu\ge\delta\right\}\cap[-\rho,\rho]^3.
\]
If $\eta\le\delta$ and $R\ge2\rho+3\eta$, then $f_\eta$ is supported on the set $\mathcal P^{+\circ}_\nu$ of Section~\ref{sec:admissible-families-lattice}: $f_\eta(\sigma)=0$ for every plaquette $\sigma\notin\mathcal P^{+\circ}_\nu$.
\end{lemma}

\begin{proof}
Let $n=R/\eta$ and $m=\frac{n-1}{2}$ be as in Section~\ref{sec:tori-cells-cochains}. Then $m\eta=\frac{R-\eta}{2}$, and the hypothesis on $R$ gives $m\eta\ge\rho+\eta$; in particular $\rho+\eta<\frac{R}{2}$.

Let $\sigma=(x;\alpha\beta)$, $\alpha<\beta$, be a plaquette with $f_\eta(\sigma)\ne0$; here $x$ is identified with its representative in $\eta\cdot\{-m,\dots,m\}^3$. By \eqref{eq:continuum-infinite-2} the cube $Q(\sigma)$ meets $\operatorname{supp}f+R\cdot k$ for some $k\in\mathbb Z^3$. Since $-m\eta\le c(\sigma)_\nu\le m\eta+\frac{\eta}{2}=\frac{R}{2}$, the $\nu$-coordinates of the points of $Q(\sigma)$ lie in $\left[-\frac{R}{2},\frac{R}{2}+\frac{\eta}{2}\right)$. The $\nu$-coordinates of the points of $\operatorname{supp}f+R\cdot k$ lie in $[\delta,\rho]+Rk_\nu$; for $k_\nu\ne0$ this set is contained in $(-\infty,\rho-R]\cup[\delta+R,\infty)$, which is disjoint from $\left[-\frac{R}{2},\frac{R}{2}+\frac{\eta}{2}\right)$ because $\rho<\frac{R}{2}$, $\delta>0$ and $\eta\le R$. Hence $k_\nu=0$, and some point $y\in Q(\sigma)$ satisfies $\delta\le y_\nu\le\rho$. Together with $c(\sigma)_\nu-\frac{\eta}{2}\le y_\nu<c(\sigma)_\nu+\frac{\eta}{2}$ this gives
\[
\delta-\frac{\eta}{2}<c(\sigma)_\nu\le\rho+\frac{\eta}{2}.
\]

Suppose first that $\nu\notin\{\alpha,\beta\}$. Then $c(\sigma)_\nu=x_\nu\in\eta\mathbb Z$; from $x_\nu>\delta-\frac{\eta}{2}\ge\frac{\eta}{2}$ it follows that $x_\nu\ge\eta$, and $x_\nu\le\rho+\frac{\eta}{2}<m\eta$. Every vertex of $\sigma$ has the $\nu$-coordinate $x_\nu\in\{\eta,\dots,(m-1)\eta\}$, so $\sigma\in\mathcal P^+_\nu$; moreover $\sigma$ lies in the plane $\{y_\nu=x_\nu\}$, strictly between the planes $\{y_\nu=0\}$ and $\{y_\nu=m\eta\}$, and is contained in neither of them. Thus $\sigma\in\mathcal P^{+\circ}_\nu$.

Suppose next that $\nu\in\{\alpha,\beta\}$. Then $c(\sigma)_\nu=x_\nu+\frac{\eta}{2}$, hence $x_\nu>\delta-\eta\ge0$, so that $x_\nu\ge\eta$, and $x_\nu+\eta\le\rho+\eta\le m\eta$. The vertices of $\sigma$ have the $\nu$-coordinates $x_\nu$ and $x_\nu+\eta$, both in $\{\eta,\dots,m\eta\}$, so $\sigma\in\mathcal P^+_\nu$; and $\sigma$ has an edge in the direction $\nu$, so it is contained in no plane $\{y_\nu=a\}$, $a\in\mathbb R$. Thus $\sigma\in\mathcal P^{+\circ}_\nu$.
\end{proof}

\subsection{Exponential moments and cumulants at a fixed cutoff}\label{sec:exponential-moments-cumulants}

\begin{proposition}[Sub-Gaussian random linear forms]\label{prop:sub-gaussian-random}
Let $(\Omega,P)$ be a probability space with expectation $E$, let $\mathcal T$ be a real vector space with a Hilbert seminorm $|\cdot|$, that is, the seminorm of a positive semidefinite symmetric bilinear form on $\mathcal T$, and let $X$ be a linear map from $\mathcal T$ to the real random variables on $(\Omega,P)$ such that
\begin{equation}\label{eq:continuum-infinite-5}
E\exp(X(g))\le\exp\left(\frac{|g|^2}{2}\right)\qquad(g\in\mathcal T).
\end{equation}
Let $\mathcal T_{\mathbb C}=\mathcal T+i\mathcal T$ be the complexification of $\mathcal T$, and for $g,h\in\mathcal T$ put
\[
X(g+ih)=X(g)+iX(h),\qquad|g+ih|=\left(|g|^2+|h|^2\right)^{\frac12},\qquad\operatorname{Re}(g+ih)=g,\qquad\operatorname{Im}(g+ih)=h.
\]
For integers $n\ge0$ put
\[
\sigma_n=2\cdot2^{\frac{n}{2}}\cdot\Gamma\left(\frac{n}{2}+1\right).
\]
Then $\sigma_n\le2\left(2^n\cdot n!\right)^{\frac12}$ for every integer $n\ge0$, and the following assertions hold.
\begin{enumerate}
\item[(1)] For $g\in\mathcal T$ and real $t$,
\[
E\exp(t|X(g)|)\le2\exp\left(\frac{t^2|g|^2}{2}\right).
\]
For an integer $s\ge1$ and $f_1,\dots,f_s\in\mathcal T_{\mathbb C}$ the function $z\mapsto E\exp\left(\sum_{i=1}^sz_iX(f_i)\right)$ is entire on $\mathbb C^s$. For $f\in\mathcal T_{\mathbb C}$,
\[
|E\exp(X(f))|\le\exp\left(\frac{|\operatorname{Re}f|^2}{2}\right);
\]
in particular $|E\exp(zX(g))|\le\exp\left(\frac{(\operatorname{Re}z)^2|g|^2}{2}\right)$ for $g\in\mathcal T$ and $z\in\mathbb C$.
\item[(2)] For $g\in\mathcal T$ one has $EX(g)=0$ and $EX(g)^2\le|g|^2$, and for real $t$
\[
1\le E\exp(tX(g))\le\exp\left(\frac{t^2|g|^2}{2}\right).
\]
\item[(3)] Let $g\in\mathcal T$. For $t>0$,
\[
P(|X(g)|\ge t)\le2\exp\left(-\frac{t^2}{2|g|^2}\right),
\]
where the right side is read as $0$ if $|g|=0$. For $p>0$,
\[
E|X(g)|^p\le2\cdot2^{\frac{p}{2}}\cdot\Gamma\left(\frac{p}{2}+1\right)\cdot|g|^p.
\]
For real $a$ with $0\le2a|g|^2<1$,
\[
E\exp\left(aX(g)^2\right)\le\left(1-2a|g|^2\right)^{-\frac12}.
\]
\item[(4)] For an integer $n\ge1$, $g_1,\dots,g_n\in\mathcal T$ and $f_1,\dots,f_n\in\mathcal T_{\mathbb C}$,
\[
E\prod_{i=1}^n|X(g_i)|\le\sigma_n\prod_{i=1}^n|g_i|,\qquad E\prod_{i=1}^n|X(f_i)|\le2^{\frac{n}{2}}\sigma_n\prod_{i=1}^n|f_i|.
\]
\item[(5)] For $g\in\mathcal T$ and $z\in\mathbb C$,
\[
|E\exp(zX(g))-1|\le2\left(\exp\left(\frac{|z|^2|g|^2}{2}\right)-1\right).
\]
Consequently $|E\exp(zX(g))-1|\le\frac12$ if $|z|\cdot|g|\le\frac23$, and $E\exp(zX(g))\ne0$ if $|z|\cdot|g|<\left(2\log\frac32\right)^{\frac12}=0.9005\ldots$
\item[(6)] For $f\in\mathcal T_{\mathbb C}$,
\[
|E\exp(X(f))-1|\le\exp\left(2|f|^2\right)-1,
\]
and the right side is at most $\frac12$ if $|f|\le\frac25$.
\item[(7)] Let $n\ge1$ be an integer and $f_1,\dots,f_n\in\mathcal T_{\mathbb C}$. The function
\[
W(z)=\log E\exp\left(\sum_{i=1}^nz_iX(f_i)\right),
\]
formed with the principal branch of the logarithm, is analytic and satisfies $|W(z)|\le\log2$ on the set of the $z\in\mathbb C^n$ with $\sum_{i=1}^n|z_i|\cdot|f_i|\le\frac25$. The number $\kappa_n(f_1,\dots,f_n)=\partial_{z_1}\cdots\partial_{z_n}W(0)$, the joint cumulant of $X(f_1),\dots,X(f_n)$, satisfies
\[
|\kappa_n(f_1,\dots,f_n)|\le(\log2)\cdot\left(\frac{5n}{2}\right)^n\cdot\prod_{i=1}^n|f_i|\le(\log2)\cdot\left(\frac{5\exp(1)}{2}\right)^n\cdot n!\cdot\prod_{i=1}^n|f_i|.
\]
For $g\in\mathcal T$,
\[
|\kappa_n(g,\dots,g)|\le(\log2)\cdot n!\cdot\left(\frac32\right)^n|g|^n,\qquad\kappa_1(g)=0,\qquad0\le\kappa_2(g,g)\le|g|^2.
\]
\end{enumerate}
\end{proposition}

\begin{proof}
(1) For real $t$ and $x$ one has $\exp(t|x|)\le\exp(tx)+\exp(-tx)$. Taking $x=X(g)$ and applying \eqref{eq:continuum-infinite-5} to $tg$ and to $-tg$, whose seminorm is $|t|\cdot|g|$, one obtains the first inequality. Let $s\ge1$, $f_1,\dots,f_s\in\mathcal T_{\mathbb C}$ and $r>0$. If $|z_i|\le r$ for $i=1,\dots,s$, then
\[
\left|\exp\left(\sum_{i=1}^sz_iX(f_i)\right)\right|\le\prod_{i=1}^s\exp\left(r|X(\operatorname{Re}f_i)|\right)\cdot\exp\left(r|X(\operatorname{Im}f_i)|\right).
\]
The right side does not depend on $z$, and it is integrable: by H\"older's inequality with $2s$ exponents equal to $2s$, its expectation is at most the product of the $2s$ numbers $\left(E\exp(2sr|X(h)|)\right)^{\frac{1}{2s}}$ with $h=\operatorname{Re}f_i$ or $h=\operatorname{Im}f_i$, $1\le i\le s$, and these numbers are finite by the first inequality. Put $M(z)=E\exp\left(\sum_{i=1}^sz_iX(f_i)\right)$. By dominated convergence $M$ is continuous on $\mathbb C^s$. Fix all variables except $z_i$ and let $\Delta$ be a closed triangle in the $z_i$-plane. By Fubini's theorem and Cauchy's theorem for the entire function $z_i\mapsto\exp\left(\sum_{k=1}^sz_kX(f_k)\right)$,
\[
\oint_{\partial\Delta}\mathrm{d}z_i\,M(z)=E\oint_{\partial\Delta}\mathrm{d}z_i\,\exp\left(\sum_{k=1}^sz_kX(f_k)\right)=0,
\]
so $M$ is analytic in $z_i$ by Morera's theorem. Thus $M$ is continuous on $\mathbb C^s$ and analytic in each of $z_1,\dots,z_s$ when the remaining variables are held fixed, and Osgood's lemma shows that $M$ is entire. Finally, $|\exp(X(f))|=\exp(X(\operatorname{Re}f))$ for $f\in\mathcal T_{\mathbb C}$, so that $|E\exp(X(f))|\le E\exp(X(\operatorname{Re}f))\le\exp\left(\frac{|\operatorname{Re}f|^2}{2}\right)$ by \eqref{eq:continuum-infinite-5}; for $f=zg$ with $g\in\mathcal T$ and $z\in\mathbb C$ one has $\operatorname{Re}f=(\operatorname{Re}z)g$.

(2) Let $g\in\mathcal T$, $Y=X(g)$, and $G(t)=E\exp(tY)$ for $t\in\mathbb C$. By part (1), $E\exp(|t|\cdot|Y|)<\infty$, so the exponential series may be integrated term by term:
\[
G(t)=\sum_{k\ge0}\frac{t^k}{k!}EY^k=1+t\,EY+\frac{t^2}{2}EY^2+O(t^3)\qquad(t\to0).
\]
For real $t$, \eqref{eq:continuum-infinite-5} for $tg$ gives
\[
G(t)\le\exp\left(\frac{t^2|g|^2}{2}\right)=1+\frac{t^2}{2}|g|^2+O(t^4)\qquad(t\to0).
\]
Dividing $G(t)-1$ by $t$ and letting $t\downarrow0$, respectively $t\uparrow0$, one obtains $EY\le0$ and $EY\ge0$; hence $EY=0$. Dividing then $G(t)-1$ by $t^2$ and letting $t\to0$ gives $EY^2\le|g|^2$. By Jensen's inequality $G(t)\ge\exp(t\,EY)=1$ for real $t$, and the upper bound for $G(t)$ is \eqref{eq:continuum-infinite-5} for $tg$.

(3) Let $g\in\mathcal T$ and $Y=X(g)$. If $|g|=0$, then $EY^2=0$ by part (2), so $Y=0$ almost surely, and the three assertions hold: both sides of the first two vanish, and both sides of the third are equal to $1$. Let $|g|>0$ and $t>0$. For $\lambda\ge0$, Markov's inequality for $\exp(\lambda Y)$ and \eqref{eq:continuum-infinite-5} for $\lambda g$ give
\[
P(Y\ge t)\le\exp(-\lambda t)\cdot E\exp(\lambda Y)\le\exp\left(-\lambda t+\frac{\lambda^2|g|^2}{2}\right),
\]
and the choice $\lambda=\frac{t}{|g|^2}$ yields $P(Y\ge t)\le\exp\left(-\frac{t^2}{2|g|^2}\right)$. The same bound holds for $-Y=X(-g)$, and $P(|Y|\ge t)\le P(Y\ge t)+P(-Y\ge t)$ gives the first assertion. For $p>0$, Tonelli's theorem, the first assertion and the substitution $t=|g|\cdot(2v)^{\frac12}$ give
\[
E|Y|^p=\int_0^\infty\mathrm{d}t\,p\,t^{p-1}P(|Y|>t)\le2p\int_0^\infty\mathrm{d}t\,t^{p-1}\exp\left(-\frac{t^2}{2|g|^2}\right)=p\cdot2^{\frac{p}{2}}\cdot\Gamma\left(\frac{p}{2}\right)\cdot|g|^p,
\]
and $p\cdot\Gamma\left(\frac{p}{2}\right)=2\,\Gamma\left(\frac{p}{2}+1\right)$; this is the second assertion. In the third assertion $a\ge0$, because $|g|>0$. Let $\xi$ be a standard normal random variable independent of $Y$, realised on the product of $(\Omega,P)$ with a second probability space, and let $E_\xi$ and $E_Y$ denote the expectations with respect to $\xi$ and to $Y$. Since $E_\xi\exp(c\xi)=\exp\left(\frac{c^2}{2}\right)$ for real $c$, Tonelli's theorem and \eqref{eq:continuum-infinite-5} for the elements $(2a)^{\frac12}\xi g$ of $\mathcal T$ give
\[
E\exp\left(aY^2\right)=E_\xi E_Y\exp\left((2a)^{\frac12}\xi Y\right)\le E_\xi\exp\left(a\xi^2|g|^2\right)=\left(1-2a|g|^2\right)^{-\frac12},
\]
where the last equality holds because $2a|g|^2<1$.

(4) For every integer $n\ge0$ the point $\frac{n}{2}+1$ is the midpoint of $1$ and $n+1$, and $\log\Gamma$ is convex on $(0,\infty)$; hence $\Gamma\left(\frac{n}{2}+1\right)^2\le\Gamma(1)\cdot\Gamma(n+1)=n!$, which is the inequality $\sigma_n\le2\left(2^n\cdot n!\right)^{\frac12}$. Now let $n\ge1$. By part (3) with $p=n$, $E|X(g)|^n\le\sigma_n|g|^n$ for $g\in\mathcal T$, and H\"older's inequality with $n$ exponents equal to $n$ gives
\[
E\prod_{i=1}^n|X(g_i)|\le\prod_{i=1}^n\left(E|X(g_i)|^n\right)^{\frac{1}{n}}\le\sigma_n\prod_{i=1}^n|g_i|.
\]
For $f=g+ih$ with $g,h\in\mathcal T$, Minkowski's inequality in $L^n(P)$ gives
\[
\|X(f)\|_{L^n}\le\|X(g)\|_{L^n}+\|X(h)\|_{L^n}\le\sigma_n^{\frac{1}{n}}\left(|g|+|h|\right)\le\sigma_n^{\frac{1}{n}}\sqrt2\,|f|,
\]
and H\"older's inequality, applied in the same way, gives $E\prod_{i=1}^n|X(f_i)|\le\prod_{i=1}^n\|X(f_i)\|_{L^n}\le2^{\frac{n}{2}}\sigma_n\prod_{i=1}^n|f_i|$.

(5) Let $\varphi(t)=\exp(t)-1-t=\sum_{j\ge2}\frac{t^j}{j!}$ for real $t$. For $t\ge0$,
\[
\varphi(t)=(\cosh t-1)+(\sinh t-t)\le2(\cosh t-1),
\]
since $\sinh t-t$ and $\cosh t-1$ vanish at $t=0$ and $(\sinh t-t)'=\cosh t-1\le\sinh t=(\cosh t-1)'$, the middle inequality being $\cosh t-\sinh t=\exp(-t)\le1$. Comparison of the power series gives $|\exp(w)-1-w|\le\varphi(|w|)$ for $w\in\mathbb C$. Let $g\in\mathcal T$, $Y=X(g)$ and $z\in\mathbb C$. By part (2), $E(zY)=0$, hence $E\exp(zY)-1=E\left(\exp(zY)-1-zY\right)$; together with \eqref{eq:continuum-infinite-5} for $|z|g$ and for $-|z|g$ this gives
\begin{align*}
|E\exp(zY)-1|&\le E\varphi(|z|\cdot|Y|)\le2E\left(\cosh(|z|Y)-1\right)\\
&=E\left(\exp(|z|Y)+\exp(-|z|Y)-2\right)\le2\left(\exp\left(\frac{|z|^2|g|^2}{2}\right)-1\right).
\end{align*}
If $|z|\cdot|g|\le\frac23$, the right side is at most $2\left(\exp\left(\frac29\right)-1\right)=0.49769\ldots<\frac12$. If $|z|\cdot|g|<\left(2\log\frac32\right)^{\frac12}$, then $\exp\left(\frac{|z|^2|g|^2}{2}\right)<\frac32$, so that $|E\exp(zY)-1|<1$ and $E\exp(zY)\ne0$.

(6) Let $f=g+ih$ with $g,h\in\mathcal T$, and $w=X(f)=X(g)+iX(h)$. By part (2), $Ew=0$, and the inequality $|\exp(w)-1-w|\le\varphi(|w|)$ from the proof of part (5) gives
\[
|E\exp(X(f))-1|=\left|E\left(\exp(w)-1-w\right)\right|\le E\varphi(|w|).
\]
Since $|w|\le\frac12\left(2|X(g)|\right)+\frac12\left(2|X(h)|\right)$ and $\varphi$ is increasing and convex on $[0,\infty)$, the bound $\varphi(t)\le2(\cosh t-1)$, $t\ge0$, gives
\[
\varphi(|w|)\le\frac12\varphi\left(2|X(g)|\right)+\frac12\varphi\left(2|X(h)|\right)\le\left(\cosh(2X(g))-1\right)+\left(\cosh(2X(h))-1\right).
\]
By \eqref{eq:continuum-infinite-5} for $\pm2g$ and for $\pm2h$, $E\cosh(2X(g))\le\exp\left(2|g|^2\right)$ and $E\cosh(2X(h))\le\exp\left(2|h|^2\right)$. Hence
\[
|E\exp(X(f))-1|\le\left(\exp\left(2|g|^2\right)-1\right)+\left(\exp\left(2|h|^2\right)-1\right)\le\exp\left(2|f|^2\right)-1,
\]
the last inequality because $(\exp(\alpha)-1)\cdot(\exp(\beta)-1)\ge0$ for $\alpha,\beta\ge0$ and $|f|^2=|g|^2+|h|^2$. If $|f|\le\frac25$, then $2|f|^2\le0.32$, and $\exp(0.32)-1=0.3771\ldots<\frac12$.

(7) For $z\in\mathbb C^n$ put
\[
M(z)=E\exp\left(\sum_{i=1}^nz_iX(f_i)\right)=E\exp\left(X\left(\sum_{i=1}^nz_if_i\right)\right);
\]
$M$ is entire by part (1). The seminorm $|\cdot|$ on $\mathcal T_{\mathbb C}$ is the seminorm of the Hermitian extension of the symmetric form that defines $|\cdot|$ on $\mathcal T$; so it is a Hilbert seminorm on $\mathcal T_{\mathbb C}$, and $\left|\sum_{i=1}^nz_if_i\right|\le\sum_{i=1}^n|z_i|\cdot|f_i|$. By part (6), $|M(z)-1|\le\frac12$ whenever $\sum_i|z_i|\cdot|f_i|\le\frac25$. These points lie in the open set $\{z\in\mathbb C^n:|M(z)-1|<1\}$, on which $W=\log M$ is analytic; since $|\log(1+u)|\le\sum_{k\ge1}\frac{|u|^k}{k}=-\log(1-|u|)$ for $|u|<1$, they satisfy $|W(z)|\le-\log\left(1-\frac12\right)=\log2$.

If $|f_i|=0$ for some $i$, then $X(f_i)=0$ almost surely by part (2), $W$ does not depend on $z_i$, and $\kappa_n(f_1,\dots,f_n)=0$. Otherwise put $r_i=\frac{2}{5n|f_i|}$. Every point of the closed polydisk $\{z\in\mathbb C^n:|z_i|\le r_i\text{ for }1\le i\le n\}$ satisfies $\sum_i|z_i|\cdot|f_i|\le\frac25$, and Cauchy's integral formula on this polydisk gives
\[
|\kappa_n(f_1,\dots,f_n)|\le(\log2)\cdot\prod_{i=1}^nr_i^{-1}=(\log2)\cdot\left(\frac{5n}{2}\right)^n\cdot\prod_{i=1}^n|f_i|.
\]
The second bound for $\kappa_n(f_1,\dots,f_n)$ follows from $n^n\le\exp(n)\cdot n!$, which holds because $\frac{n^n}{n!}$ is a term of the power series of $\exp(n)$.

Let $g\in\mathcal T$. The function $\zeta\mapsto E\exp(\zeta X(g))$ is entire by part (1), so the set $U_g$ of the $\zeta\in\mathbb C$ with $\left|E\exp(\zeta X(g))-1\right|<1$ is open; it contains $0$, where the function takes the value $1$. For $\zeta\in U_g$ the number $E\exp(\zeta X(g))$ lies in the disc $\{w\in\mathbb C:|w-1|<1\}$, which does not contain $0$ and on which the principal branch of the logarithm is analytic; hence
\[
W_g(\zeta)=\log E\exp(\zeta X(g))\qquad(\zeta\in U_g)
\]
defines an analytic function $W_g\colon U_g\to\mathbb C$. By part (5), $|E\exp(\zeta X(g))-1|\le\frac12$ whenever $|\zeta|\cdot|g|\le\frac23$; these $\zeta$ therefore belong to $U_g$, and the bound for $|\log(1+u)|$ used above gives $|W_g(\zeta)|\le\log2$ for them. For $f_1=\dots=f_n=g$ one has $M(z)=E\exp\left((z_1+\dots+z_n)X(g)\right)$ for $z\in\mathbb C^n$; so every $z$ in the open set $\{z\in\mathbb C^n:|M(z)-1|<1\}$, which contains $0$, satisfies $z_1+\dots+z_n\in U_g$ and $W(z)=W_g(z_1+\dots+z_n)$. Hence $\kappa_n(g,\dots,g)=W_g^{(n)}(0)$. If $|g|=0$, then $X(g)=0$ almost surely by part (2), and this derivative vanishes; otherwise Cauchy's formula on the circle $|\zeta|=\frac{2}{3|g|}$ gives $\left|W_g^{(n)}(0)\right|\le(\log2)\cdot n!\cdot\left(\frac{3|g|}{2}\right)^n$. Finally, by part (2), $\kappa_1(g)=W_g'(0)=EX(g)=0$ and
\[
\kappa_2(g,g)=W_g''(0)=EX(g)^2-\left(EX(g)\right)^2=EX(g)^2\in\left[0,|g|^2\right].
\]
\end{proof}

\begin{theorem}[Estimates at a fixed cutoff]\label{thm:estimates-fixed-cutoff}
Let $(N,T)$ be a cutoff and $\Lambda=\Lambda_{\eta,R}$. Let $\mathrm{Pr}$ be a Borel probability measure on $C^2(\Lambda)$ with property~\eqref{eq:continuum-infinite-Ad2}, and let $E$ be the expectation with respect to $\mathrm{Pr}$. Then every assertion of Proposition~\ref{prop:sub-gaussian-random} holds for the probability space $\left(C^2(\Lambda),\mathrm{Pr}\right)$, the space $\mathcal T=C^2(\Lambda)$ of real two-cochains with the norm $|g|=|g|_{\eta}$, and the linear map $X$ from $C^2(\Lambda)$ to real random variables given by $X(g)(F)=\langle g,F\rangle_{\eta}$ for $F\in C^2(\Lambda)$. Since $f\mapsto f_{\eta}$ is linear and $X_{\eta,R}(f)=X(f_{\eta})$, these assertions hold in particular for the space of real tests $f\in\mathcal S(\mathbb R^3;\Lambda^2)$, the Hilbert seminorm $|f|:=|f_{\eta}|_{\eta}$ and the map $f\mapsto X_{\eta,R}(f)$; this seminorm satisfies
\[
|f_{\eta}|_{\eta}\le p_R(f),\qquad |f_{\eta}|_{\eta}\le\|f\|_{L^2}\ \text{ if }\ \operatorname{supp}f\subset y_0+\left(-\frac{R}{2},\frac{R}{2}\right)^3\ \text{ for some }y_0\in\mathbb R^3.
\]
Explicitly, for real tests $f,f_1,\dots,f_n$ with $n\ge1$, and for $z\in\mathbb C$,
\begin{gather*}
\left|E\exp\left(z\cdot X_{\eta,R}(f)\right)\right|\le\exp\left(\frac{1}{2}(\operatorname{Re}z)^2\cdot|f_{\eta}|_{\eta}^2\right),\\
\left|E\exp\left(z\cdot X_{\eta,R}(f)\right)-1\right|\le\frac{1}{2}\quad\text{if }\ |z|\cdot|f_{\eta}|_{\eta}\le\frac{2}{3},\\
\left|E\prod_{i=1}^nX_{\eta,R}(f_i)\right|\le2\sqrt{2^n\,n!}\cdot\prod_{i=1}^n|(f_i)_{\eta}|_{\eta},\\
\left|\kappa_n(f_1,\dots,f_n)\right|\le(\log2)\cdot\left(\frac{5}{2}\right)^n\cdot\exp(n)\cdot n!\cdot\prod_{i=1}^n|(f_i)_{\eta}|_{\eta},
\end{gather*}
where $\kappa_n(f_1,\dots,f_n)$ is the cumulant of Proposition~\ref{prop:sub-gaussian-random}(7) formed with the map $f\mapsto X_{\eta,R}(f)$.

If moreover $\mathrm{Pr}(\operatorname{ran}\,d)=1$, as is the case for the averaged law $\bar{\mathsf P}$, under which $F=\eta^{-\frac{3}{2}}\cdot dA$, and if $\Pi$ denotes the orthogonal projection of $C^2(\Lambda)$ onto $\operatorname{ran}\,d$, then for every real $g\in C^2(\Lambda)$
\[
E\exp\langle g,F\rangle_{\eta}\le\exp\left(\frac{1}{2}|\Pi g|_{\eta}^2\right)=\int\mathrm{d}\mathsf P^{\mathrm{free}}(F)\,\exp\langle g,F\rangle_{\eta},
\]
where $\mathsf P^{\mathrm{free}}$ denotes the averaged law at charge $q=0$, with an arbitrary lattice mass $\boldsymbol{\mu}\in\mathbb R$.
\end{theorem}

\begin{proof}
For the space $C^2(\Lambda)$ with the norm $|\cdot|_{\eta}$ and the map $X$, hypothesis~\eqref{eq:continuum-infinite-5} of Proposition~\ref{prop:sub-gaussian-random} is property~\eqref{eq:continuum-infinite-Ad2}. By~\eqref{eq:continuum-infinite-2} one has $X_{\eta,R}(f)=\langle f_{\eta},F\rangle_{\eta}=X(f_{\eta})$, and $f\mapsto f_{\eta}$ is linear; hence $f\mapsto|f_{\eta}|_{\eta}$ is a Hilbert seminorm on the real tests, and~\eqref{eq:continuum-infinite-5} for the map $f\mapsto X_{\eta,R}(f)$ is property~\eqref{eq:continuum-infinite-Ad2} for the cochain $g=f_{\eta}$. The two bounds for $|f_{\eta}|_{\eta}$ are those of Lemma~\ref{lem:norm-bounds-lattice}. The first two displayed estimates are Proposition~\ref{prop:sub-gaussian-random}(1) and~(5). The third follows from Proposition~\ref{prop:sub-gaussian-random}(4), from the bound $\sigma_n\le2\sqrt{2^n\,n!}$ stated there, and from
\[
\left|E\prod_{i=1}^nX_{\eta,R}(f_i)\right|\le E\prod_{i=1}^n\left|X_{\eta,R}(f_i)\right|.
\]
The fourth is Proposition~\ref{prop:sub-gaussian-random}(7).

Suppose now that $\mathrm{Pr}(\operatorname{ran}\,d)=1$. Every $F\in\operatorname{ran}\,d$ satisfies $\Pi F=F$, and $\Pi$ is symmetric for $\langle\cdot,\cdot\rangle_{\eta}$; hence $\langle g,F\rangle_{\eta}=\langle\Pi g,F\rangle_{\eta}$ for $\mathrm{Pr}$-almost every $F$, and property~\eqref{eq:continuum-infinite-Ad2} for the real cochain $\Pi g$ yields the inequality. At $q=0$ the weight $\bar w$ is a positive constant times $\exp\left(-\frac{1}{2}\|dA\|^2\right)=\exp\left(-\frac{1}{2}|F|_{\eta}^2\right)$, and $A\mapsto F=\eta^{-\frac{3}{2}}\cdot dA$ is a linear bijection of $V_{\mathrm{t}}$ onto $\operatorname{ran}\,d$ (Section~\ref{sec:wilson-operator-lattice} and the proof of Proposition~\ref{prop:gaussian-domination}). Consequently, for every value of $\boldsymbol{\mu}$, the random vector $F$ is under $\mathsf P^{\mathrm{free}}$ the standard Gaussian vector of the Euclidean space $\left(\operatorname{ran}\,d,|\cdot|_{\eta}\right)$; its Laplace transform at $g$ is therefore $\exp\left(\frac{1}{2}|\Pi g|_{\eta}^2\right)$.
\end{proof}

\subsection{Existence of limit points}\label{sec:existence-limit-points}

Let $(\mathrm{Pr}_{N,T})$ be an admissible family (Definition~\ref{def:admissible-families}), so that $\mathrm{Pr}_{N,T}$ is, for every cutoff $(N,T)$, a Borel probability measure on $\ker d\subset C^2(\Lambda_{\eta,R})$, where $\eta=L^{-N}$ and $R=h\cdot L^T$. For a cutoff $(N,T)$ let $\mathsf P_{N,T}$ be the law under $\mathrm{Pr}_{N,T}$ of the distribution $F^{\sharp}$ of Section~\ref{sec:test-forms-test}. The map $F\mapsto F^{\sharp}$ is linear, and it takes its values in $H_{-3}$, because
\[
\|F^{\sharp}\|_{-3}\le8|F|_{\eta}\qquad\left(F\in C^2(\Lambda_{\eta,R})\right),
\]
a bound established at the end of the proof of Theorem~\ref{thm:second-moments-tightness}(1). Being a linear map on the finite-dimensional space $C^2(\Lambda_{\eta,R})$, it is continuous; thus $\mathsf P_{N,T}$ is a Borel probability measure on $H_{-3}\subset H_{-4}$.

\begin{theorem}[Second moments and tightness]\label{thm:second-moments-tightness}
Let $(\mathrm{Pr}_{N,T})$ be an admissible family, and for every cutoff $(N,T)$ let $\mathsf P_{N,T}$ be the law of $F^{\sharp}$ under $\mathrm{Pr}_{N,T}$. Then the following hold.
\begin{enumerate}
\item[(1)] One has
\[
\sup_{N,T}\int\mathrm{d}\mathsf P_{N,T}(F)\,\|F\|_{-3}^2\le961\,\mathsf Z_4\le61,
\]
the supremum being taken over all cutoffs $(N,T)$.
\item[(2)] For $\varrho>0$ let $K_{\varrho}=\left\{S\in\mathcal S'(\mathbb R^3;\Lambda^2):\|S\|_{-3}\le\varrho\right\}$. Then $K_{\varrho}$ is a compact subset of $H_{-4}$, and $\mathsf P_{N,T}(H_{-4}\setminus K_{\varrho})\le61\varrho^{-2}$ for every cutoff $(N,T)$. Consequently the family $\{\mathsf P_{N,T}\}$ is tight on $H_{-4}$, and also on $\mathcal S'(\mathbb R^3;\Lambda^2)$ endowed with its strong topology or with its weak-$*$ topology.
\item[(3)] From every sequence $(N_k,T_k)_{k\in\mathbb N}$ of cutoffs one can extract a subsequence along which $\mathsf P_{N_k,T_k}$ converges weakly on $H_{-4}$, and therefore on $\mathcal S'(\mathbb R^3;\Lambda^2)$, to a Borel probability measure $\mu$ on $H_{-4}$ with $\mu(H_{-3})=1$.
\end{enumerate}
\end{theorem}

\begin{proof}
\begin{enumerate}
\item[(1)] Let $E$ be the expectation with respect to $\mathrm{Pr}_{N,T}$, regarded as a measure on $C^2(\Lambda_{\eta,R})$. The family being admissible, this measure has property~\eqref{eq:continuum-infinite-Ad2} (Definition~\ref{def:admissible-families}), so that Theorem~\ref{thm:estimates-fixed-cutoff} applies to it. The Hermite functions $h_{\kappa}$ are real, and $F^{\sharp}(h_{\kappa})=X_{\eta,R}(h_{\kappa})$ by~\eqref{eq:continuum-infinite-3}. By Theorem~\ref{thm:estimates-fixed-cutoff}, in the form of Proposition~\ref{prop:sub-gaussian-random}(2), and by Lemma~\ref{lem:norm-bounds-lattice},
\[
E\left[X_{\eta,R}(h_{\kappa})^2\right]\le|(h_{\kappa})_{\eta}|_{\eta}^2\le p_R(h_{\kappa})^2\le p_1(h_{\kappa})^2\le961\,\lambda_{\kappa}^2\qquad(\kappa\in\mathcal K).
\]
With the definition of $\|\cdot\|_{-3}$ in Section~\ref{sec:hermite-sobolev-spaces}, and since all terms of the series are nonnegative, so that sum and integral may be interchanged, it follows that
\begin{align}
\int\mathrm{d}\mathsf P_{N,T}(F)\,\|F\|_{-3}^2
&=\sum_{\kappa\in\mathcal K}\lambda_{\kappa}^{-6}\cdot E\left[X_{\eta,R}(h_{\kappa})^2\right]\notag\\
&\le\sum_{\kappa\in\mathcal K}\lambda_{\kappa}^{-6}\cdot p_1(h_{\kappa})^2\le961\sum_{\kappa\in\mathcal K}\lambda_{\kappa}^{-4}.\label{eq:continuum-infinite-6}
\end{align}
By~\eqref{eq:continuum-infinite-4} the last sum equals $\mathsf Z_4<\frac{1}{16}$, and $\frac{961}{16}<61$; this is~(1). For a fixed $F\in C^2(\Lambda_{\eta,R})$ the same computation, with the Cauchy--Schwarz bound
\[
|F^{\sharp}(h_{\kappa})|=\left|\langle(h_{\kappa})_{\eta},F\rangle_{\eta}\right|\le|(h_{\kappa})_{\eta}|_{\eta}\cdot|F|_{\eta}
\]
in place of the expectation, gives $\|F^{\sharp}\|_{-3}^2\le961\,\mathsf Z_4\cdot|F|_{\eta}^2\le61\,|F|_{\eta}^2$, and hence, as $61<64$, $\|F^{\sharp}\|_{-3}\le8|F|_{\eta}$.
\item[(2)] Let $(S_j)_{j\in\mathbb N}$ be a sequence in $K_{\varrho}$. For all $j$ and $\kappa$ one has $|S_j(h_{\kappa})|\le\varrho\cdot\lambda_{\kappa}^3$, and $\mathcal K$ is countable; a diagonal argument therefore gives a subsequence, again denoted by $(S_j)$, and numbers $t_{\kappa}$ such that $S_j(h_{\kappa})\to t_{\kappa}$ for every $\kappa\in\mathcal K$. By Fatou's lemma $\sum_{\kappa}\lambda_{\kappa}^{-6}\cdot|t_{\kappa}|^2\le\varrho^2$. Put
\[
S(f)=\sum_{\kappa\in\mathcal K}t_{\kappa}\cdot\langle h_{\kappa},f\rangle_{L^2}\qquad\left(f\in\mathcal S(\mathbb R^3;\Lambda^2)\right);
\]
the series converges absolutely, with $|S(f)|\le\varrho\cdot\|f\|_3$ by the Cauchy--Schwarz inequality, and $S(h_{\kappa})=t_{\kappa}$, so that $S\in K_{\varrho}$. Let $M>0$. The indices $\kappa$ with $\lambda_{\kappa}\le M$ are finite in number, and the terms with $\lambda_{\kappa}>M$ contribute to $\|S_j-S\|_{-4}^2$ at most
\[
M^{-2}\cdot\sum_{\kappa\in\mathcal K}\lambda_{\kappa}^{-6}\cdot|S_j(h_{\kappa})-t_{\kappa}|^2\le M^{-2}\cdot\left(\|S_j\|_{-3}+\|S\|_{-3}\right)^2\le4\varrho^2\cdot M^{-2}.
\]
Hence
\[
\|S_j-S\|_{-4}^2\le\sum_{\lambda_{\kappa}\le M}\lambda_{\kappa}^{-8}\cdot|S_j(h_{\kappa})-t_{\kappa}|^2+4\varrho^2\cdot M^{-2},
\]
where the finite sum tends to $0$ as $j\to\infty$. So $\limsup_j\|S_j-S\|_{-4}^2\le4\varrho^2\cdot M^{-2}$ for every $M>0$, that is, $S_j\to S$ in $H_{-4}$. Thus $K_{\varrho}$ is sequentially compact, hence compact, in the metric space $H_{-4}$. Chebyshev's inequality and~(1) give
\[
\mathsf P_{N,T}(H_{-4}\setminus K_{\varrho})=\mathsf P_{N,T}\left(\left\{F:\|F\|_{-3}>\varrho\right\}\right)\le\varrho^{-2}\cdot\int\mathrm{d}\mathsf P_{N,T}(F)\,\|F\|_{-3}^2\le61\varrho^{-2}.
\]
For a bounded set $B\subset\mathcal S(\mathbb R^3;\Lambda^2)$ the number $\sup_{f\in B}\|f\|_4$ is finite, and by Section~\ref{sec:hermite-sobolev-spaces}
\[
\sup_{f\in B}|F(f)|\le\|F\|_{-4}\cdot\sup_{f\in B}\|f\|_4\qquad(F\in H_{-4}).
\]
So the inclusion $H_{-4}\to\mathcal S'(\mathbb R^3;\Lambda^2)$ is continuous for the strong topology, and a fortiori for the weak-$*$ topology. Therefore $K_{\varrho}$ is compact in $\mathcal S'(\mathbb R^3;\Lambda^2)$ for both topologies, which gives the tightness on $\mathcal S'(\mathbb R^3;\Lambda^2)$.
\item[(3)] By~(2) the family $\{\mathsf P_{N,T}\}$ is tight on the separable Hilbert space $H_{-4}$. By Prokhorov's theorem (a tight family of Borel probability measures on a metric space is relatively compact for weak convergence), every sequence $\left(\mathsf P_{N_k,T_k}\right)_k$ has a subsequence that converges weakly on $H_{-4}$ to a Borel probability measure $\mu$ on $H_{-4}$. A bounded function on $\mathcal S'(\mathbb R^3;\Lambda^2)$ that is continuous for the strong topology, in particular one that is continuous for the weak-$*$ topology, restricts to a bounded continuous function on $H_{-4}$, the inclusion being continuous; hence the convergence also holds on $\mathcal S'(\mathbb R^3;\Lambda^2)$. Finally, $K_{\varrho}$ is closed in $H_{-4}$, being compact, and the portmanteau theorem gives, along the subsequence,
\[
\mu(K_{\varrho})\ge\limsup_{k\to\infty}\mathsf P_{N_k,T_k}(K_{\varrho})\ge1-61\varrho^{-2}\qquad(\varrho>0).
\]
Since $K_{\varrho}\subset H_{-3}$ for every $\varrho$, it follows that $\mu(H_{-3})=1$.
\end{enumerate}
\end{proof}

\begin{definition}[Limit points]\label{def:limit-points}
Let $(\mathrm{Pr}_{N,T})$ be an admissible family and let $\mu$ be a Borel probability measure on $H_{-4}$. One calls $\mu$ a \emph{limit point} of $(\mathrm{Pr}_{N,T})$ if there exists a sequence $(N_k,T_k)_{k\in\mathbb N}$ of cutoffs, with $N_k\to\infty$ and $T_k\to\infty$ as $k\to\infty$, along which $\mathsf P_{N_k,T_k}$ converges weakly on $H_{-4}$ to $\mu$. The sequences $(N_k)$ and $(T_k)$ may tend to infinity at arbitrary rates, each irrespective of the other. Along such a sequence put $\eta_k=L^{-N_k}$, $R_k=h\cdot L^{T_k}$ and $X_k=X_{\eta_k,R_k}$, and let $E_k$ denote the expectation with respect to $\mathrm{Pr}_{N_k,T_k}$. By~\eqref{eq:continuum-infinite-3}, $X_k(f)=F^{\sharp}(f)$ for every test $f$; since $\mathsf P_{N_k,T_k}$ is the law of $F^{\sharp}$ under $\mathrm{Pr}_{N_k,T_k}$, it follows that
\[
E_k\,\Phi\left(X_k(f_1),\dots,X_k(f_s)\right)=\int\mathrm{d}\mathsf P_{N_k,T_k}(F)\,\Phi\left(F(f_1),\dots,F(f_s)\right)
\]
for real tests $f_1,\dots,f_s$ and every Borel function $\Phi$ on $\mathbb R^s$ with values in $[0,\infty]$, or with complex values and such that one of the two sides converges absolutely. By Section~\ref{sec:hermite-sobolev-spaces}, $H_{-4}$ belongs to the cylinder $\sigma$-algebra $\mathcal C$ of $\mathcal S'(\mathbb R^3;\Lambda^2)$, and the trace of $\mathcal C$ on $H_{-4}$ is the Borel $\sigma$-algebra of $H_{-4}$; accordingly a limit point $\mu$ is also regarded as the probability measure $B\mapsto\mu(B\cap H_{-4})$ on $\left(\mathcal S'(\mathbb R^3;\Lambda^2),\mathcal C\right)$. By Corollary~\ref{cor:averaged-laws-admissible}, each parameter assignment yields an admissible family $\left(\bar{\mathsf P}_{\eta,R;q,\boldsymbol{\mu}}\right)_{N,T}$ of averaged laws, $(q,\boldsymbol{\mu})$ being the value of the assignment at the cutoff $(N,T)$. A Borel probability measure on $H_{-4}$ is called a \emph{limit point of the averaged laws} if there is a parameter assignment for whose family of averaged laws it is a limit point.
\end{definition}

Let $(\mathrm{Pr}_{N,T})$ be an admissible family. By Theorem~\ref{thm:second-moments-tightness}(3), every sequence $(N_k,T_k)_{k\in\mathbb N}$ of cutoffs with $N_k\to\infty$ and $T_k\to\infty$, whatever the relation between the two sequences, has a subsequence along which $\mathsf P_{N_k,T_k}$ converges weakly on $H_{-4}$; the limit is a limit point of the family.

\subsection{Properties of the limit points}\label{sec:properties-limit-points}

For a two-form $f=(f_{12},f_{13},f_{23})$ on $\mathbb R^3$ its dual vector field is
\[
\ast f=\left(f_{23},\,-f_{13},\,f_{12}\right).
\]
Let $u\mapsto\hat u$ denote the Fourier transform on $L^2(\mathbb R^3;\mathbb C^3)$, taken componentwise. The orthogonal projection $\Pi_c$ of $L^2(\mathbb R^3;\Lambda^2)$ is defined by the multiplier
\[
\widehat{(\ast\Pi_cf)}_j(p)=\sum_{l=1}^3\left(\delta_{jl}-\frac{p_j\cdot p_l}{|p|^2}\right)\cdot\widehat{(\ast f)}_l(p)\qquad\left(j\in\{1,2,3\},\ p\in\mathbb R^3\setminus\{0\}\right)
\]
(for each $p\ne0$ the matrix is the orthogonal projection of $\mathbb C^3$ onto the orthogonal complement of $p$, and $f\mapsto\ast f$ preserves the $L^2$ norm). The free Maxwell covariance is
\[
C_0(f,g)=\left\langle f,\Pi_cg\right\rangle_{L^2}\qquad\left(f,g\in L^2(\mathbb R^3;\Lambda^2)\right).
\]
For $g\in\mathcal S(\mathbb R^3)$ let $d^{*}g\in\mathcal S(\mathbb R^3;\Lambda^2)$ be the two-form of Lemma~\ref{lem:discrete-codifferentials-three}; thus $\ast d^{*}g=-\nabla g$. For $F\in\mathcal S'$ let $dF$ be the linear form on $\mathcal S(\mathbb R^3)$ given by
\[
(dF)(g)=F\left(d^{*}g\right)\qquad\left(g\in\mathcal S(\mathbb R^3)\right).
\]

From here on, with the same convention in Sections~\ref{sec:test-forms-distributions} and~\ref{sec:test-forms-distributions-2}, an element $F\in\mathcal S'$ is a continuous linear form on the real vector space of the real tests $f\in\mathcal S=\mathcal S(\mathbb R^3;\Lambda^2)$, so that $F(f)\in\mathbb R$ for every real test $f$. On complex tests $F$ acts through the real and the imaginary part,
\[
F(f)=F(\operatorname{Re}f)+iF(\operatorname{Im}f)\qquad(f\in\mathcal S),
\]
which makes $f\mapsto F(f)$ complex-linear. Let $\mathcal D_0$ be a linear subspace of the real vector space of the real tests; when $\mathcal D_0$ is all of this space, one writes $\mathcal D_0=\mathcal S$. Let
\[
\Sigma(\mathcal D_0)=\sigma\left(F\mapsto F(f)\,:\,f\in\mathcal D_0\right)\subset\mathcal C
\]
be the $\sigma$-algebra on $\mathcal S'$ generated by the functions $F\mapsto F(f)$ with $f\in\mathcal D_0$.

Let $\mu$ be a probability measure on $(\mathcal S',\mathcal C)$, and put
\[
\Phi_\mu(f)=\int\mathrm{d}\mu(F)\,\exp\left(iF(f)\right)\qquad(f\in\mathcal S\ \text{real}).
\]
With the actions $\tau_b$, $U_r$ and $\mathsf R_\nu=U_{r_\nu}$ on tests of Section~\ref{sec:test-forms-test}, consider the following conditions on $\mu$:
\begin{equation}\tag{L1}\label{eq:continuum-infinite-L1}
\int\mathrm{d}\mu(F)\,\exp\left(F(f)\right)\le\exp\left(\frac{1}{2}\|f\|_{L^2}^2\right)\qquad\text{for every real }f\in\mathcal S;
\end{equation}
\begin{equation}\tag{L2}\label{eq:continuum-infinite-L2}
\mu\left(\left\{F\in\mathcal S':F\left(d^{*}g\right)=0\right\}\right)=1\qquad\text{for every real }g\in C_c^{\infty}(\mathbb R^3);
\end{equation}
\begin{equation}\tag{L3}\label{eq:continuum-infinite-L3}
\Phi_\mu(\tau_bf)=\Phi_\mu(f)=\Phi_\mu(U_rf)\qquad\text{for every real }f\in\mathcal S,\ b\in\mathbb R^3,\ r\in O_{\mathrm{cub}};
\end{equation}
and, for every $\nu\in\{1,2,3\}$, every integer $s\ge1$, every $c\in\mathbb C^s$ and all real $f_1,\dots,f_s\in C_c^{\infty}(\mathbb R^3;\Lambda^2)$ with $\operatorname{supp}f_i\subset\{y\in\mathbb R^3:y_\nu>0\}$ for $1\le i\le s$,
\begin{equation}\tag{L4}\label{eq:continuum-infinite-L4}
\sum_{i,j=1}^s\overline{c_i}\cdot c_j\cdot\Phi_\mu\left(f_j-\mathsf R_\nu f_i\right)\ge0.
\end{equation}

\begin{lemma}[Totality of the exponentials]\label{lem:totality-exponentials}
Let $\mu$ be a probability measure on $(\mathcal S',\mathcal C)$, and let $\mathcal D_0$ be a linear subspace of the real space $\mathcal S$. Then the linear span of the functions $F\mapsto\exp\left(iF(f)\right)$, $f\in\mathcal D_0$, is dense in $L^2\left(\mathcal S',\Sigma(\mathcal D_0),\mu\right)$.
\end{lemma}

\begin{proof}
Let $\Psi\in L^2\left(\mathcal S',\Sigma(\mathcal D_0),\mu\right)$ be orthogonal to the function $F\mapsto\exp\left(iF(f)\right)$ for every $f\in\mathcal D_0$. It suffices to show that $\Psi=0$ $\mu$-almost everywhere.

Fix $f_1,\dots,f_s\in\mathcal D_0$. Let $\mathcal G\subset\Sigma(\mathcal D_0)$ be the $\sigma$-algebra generated by $F(f_1),\dots,F(f_s)$, and let $\rho$ be the law under $\mu$ of the vector $\left(F(f_1),\dots,F(f_s)\right)$, a probability measure on $\mathbb R^s$. Every $\mathcal G$-measurable function is a Borel function of this vector; hence
\[
E_\mu\left(\Psi\mid\mathcal G\right)=\psi\left(F(f_1),\dots,F(f_s)\right)
\]
with a Borel function $\psi\in L^1(\rho)$. For $t\in\mathbb R^s$ the test $\sum_it_i\cdot f_i$ lies in $\mathcal D_0$, and $F\left(\sum_it_i\cdot f_i\right)=\sum_it_i\cdot F(f_i)$ is real. Orthogonality to $\exp\left(iF\left(\sum_it_i\cdot f_i\right)\right)$ therefore gives
\[
0=\int\mathrm{d}\mu(F)\,\Psi(F)\cdot\exp\left(-i\sum_{i=1}^st_i\cdot F(f_i)\right)=\int\mathrm{d}\rho(x)\,\psi(x)\cdot\exp\left(-i\langle t,x\rangle\right)\qquad\left(t\in\mathbb R^s\right).
\]
Thus the finite complex measure $\psi\cdot\rho$ on $\mathbb R^s$ has vanishing Fourier transform, and by the uniqueness theorem for Fourier transforms of finite complex measures on $\mathbb R^s$ one has $\psi=0$ $\rho$-almost everywhere. It follows that $\int_B\mathrm{d}\mu\,\Psi=0$ for every $B\in\mathcal G$.

The union of the $\sigma$-algebras $\mathcal G$ so obtained, over all finite families $f_1,\dots,f_s$ in $\mathcal D_0$, is a $\pi$-system (a set of the $\sigma$-algebra of one family and a set of the $\sigma$-algebra of another both belong to the $\sigma$-algebra of the joined family), and it generates $\Sigma(\mathcal D_0)$. The sets $B\in\Sigma(\mathcal D_0)$ with $\int_B\mathrm{d}\mu\,\Psi=0$ form a $\lambda$-system: they include $\mathcal S'$, they are closed under proper differences, and they are closed under increasing unions by dominated convergence, $\Psi$ being integrable. By Dynkin's $\pi$-$\lambda$ theorem, $\int_B\mathrm{d}\mu\,\Psi=0$ for every $B\in\Sigma(\mathcal D_0)$. Since $\Psi$ is $\Sigma(\mathcal D_0)$-measurable, $\Psi=0$ $\mu$-almost everywhere.
\end{proof}

\begin{theorem}\label{thm:continuum-infinite-1}
Let $(\mathrm{Pr}_{N,T})$ be an admissible family and $\mu$ one of its limit points; let $(N_k,T_k)$, $\eta_k$, $R_k$, $E_k$ and $X_k$ be as in Definition~\ref{def:limit-points}, so that $\mathsf P_{N_k,T_k}\to\mu$ weakly on $H_{-4}$.
\begin{enumerate}
\item[(1)] \emph{Exponential bound; extension to $L^2$.} For every real $f\in\mathcal S$,
\[
\int\mathrm{d}\mu(F)\,\exp\left(F(f)\right)\le\exp\left(\frac{1}{2}\|f\|_{L^2}^2\right).
\]
There is exactly one bounded linear map $X\colon L^2(\mathbb R^3;\Lambda^2)\to L^2(\mu)$ with $X(f)=F(f)$ for $f\in\mathcal S$, and
\[
\int\mathrm{d}\mu\,\exp\left(X(f)\right)\le\exp\left(\frac{1}{2}\|f\|_{L^2}^2\right)\qquad\text{for every real }f\in L^2(\mathbb R^3;\Lambda^2).
\]
By the last inequality, Proposition~\ref{prop:sub-gaussian-random} applies to the probability space $(\mathcal S',\mathcal C,\mu)$, to the real space $\mathcal T=L^2(\mathbb R^3;\Lambda^2)$ with the norm $|f|=\|f\|_{L^2}$, and to the map $X$.
\item[(2)] \emph{Convergence of moments and cumulants.} Let $f_1,\dots,f_s\in\mathcal S$ be complex tests and, for $z\in\mathbb C^s$,
\[
G_k(z)=E_k\exp\left(\sum_{i=1}^sz_i\cdot X_k(f_i)\right),\qquad G(z)=\int\mathrm{d}\mu(F)\,\exp\left(\sum_{i=1}^sz_i\cdot F(f_i)\right).
\]
Then $G_k\to G$ locally uniformly on $\mathbb C^s$, and all moments converge:
\[
E_k\prod_{i=1}^sX_k(f_i)^{\alpha_i}\longrightarrow\int\mathrm{d}\mu(F)\,\prod_{i=1}^sF(f_i)^{\alpha_i}\qquad\left(\alpha\in\mathbb N_0^s\right).
\]
The principal logarithms $W_k=\log G_k$ converge to $W=\log G$ locally uniformly on $\left\{z\in\mathbb C^s:\sum_i|z_i|\cdot\|f_i\|_{L^2}<\frac{2}{5}\right\}$, and all cumulants converge: every derivative of $W_k$ at $z=0$ converges to the corresponding derivative of $W$.
\item[(3)] \emph{Support in $H_{-s}$.} Let $s>\frac{3}{2}$. Then
\[
\int\mathrm{d}\mu(F)\,\|F\|_{-s}^2\le\mathsf Z_{2s},\qquad\int\mathrm{d}\mu(F)\,\exp\left(a\cdot\|F\|_{-s}^2\right)\le\left(1-2a\cdot\mathsf Z_{2s}\right)^{-\frac{1}{2}}\quad\text{for }0\le2a\cdot\mathsf Z_{2s}<1.
\]
In particular $\mu(H_{-s})=1$.
\item[(4)] \emph{The constraint $dF=0$; bound by the covariance $C_0$.} One has $\mu\left(\{F\in\mathcal S':dF=0\}\right)=1$, and $X(f)=X(\Pi_cf)$ $\mu$-almost surely for every $f\in L^2(\mathbb R^3;\Lambda^2)$. For every real $f\in L^2(\mathbb R^3;\Lambda^2)$,
\[
\int\mathrm{d}\mu\,\exp\left(X(f)\right)\le\exp\left(\frac{1}{2}C_0(f,f)\right);\quad\text{in particular}\quad0\le\int\mathrm{d}\mu\,X(f)^2\le C_0(f,f)\le\|f\|_{L^2}^2.
\]
\item[(5)] \emph{Translation invariance.} One has $\mu\circ(\tau_b^{*})^{-1}=\mu$ for every $b\in\mathbb R^3$, and $(\mathcal U_b\Psi)(F)=\Psi(\tau_b^{*}F)$ defines a strongly continuous unitary representation $b\mapsto\mathcal U_b$ of $\mathbb R^3$ on $L^2(\mu)$.
\item[(6)] \emph{Invariance under the cubic group.} One has $\mu\circ(U_r^{*})^{-1}=\mu$ for every $r\in O_{\mathrm{cub}}$.
\item[(7)] \emph{Reflection positivity.} For $\nu\in\{1,2,3\}$ and $c\in\mathbb R$ let $\mathsf R_{\nu,c}=\tau_{c\cdot e_\nu}\mathsf R_\nu\tau_{-c\cdot e_\nu}$; let $\mathcal D^{+}_{\nu,c}$ be the space of real $f\in C_c^{\infty}(\mathbb R^3;\Lambda^2)$ with $\operatorname{supp}f\subset\{y_\nu>c\}$, and $\Sigma^{+}_{\nu,c}=\Sigma(\mathcal D^{+}_{\nu,c})$; for a function $A$ on $\mathcal S'$ let $(\Theta_{\nu,c}A)(F)=\overline{A(\mathsf R^{*}_{\nu,c}F)}$, where $(\mathsf R^{*}_{\nu,c}F)(f)=F(\mathsf R_{\nu,c}f)$. Then
\[
\int\mathrm{d}\mu\,(\Theta_{\nu,c}A)\cdot A\ge0\qquad\text{for every }A\in L^2\left(\mathcal S',\Sigma^{+}_{\nu,c},\mu\right).
\]
For every $f\in L^2(\mathbb R^3;\Lambda^2)$ that vanishes almost everywhere on $\{y_\nu<c\}$, the function $X(f)$ coincides $\mu$-almost surely with a $\Sigma^{+}_{\nu,c}$-measurable function. Hence the inequality holds for all polynomials and all exponentials in finitely many such $X(f)$, in particular in finitely many $F(f)$ with $f\in\mathcal S$ and $\operatorname{supp}f\subset\{y_\nu\ge c\}$.
\item[(8)] \emph{The moments determine $\mu$.} For $n\ge1$ and real $f_1,\dots,f_n\in C_c^{\infty}(\mathbb R^3;\Lambda^2)$ let
\[
S_n(f_1,\dots,f_n)=\int\mathrm{d}\mu(F)\,\prod_{i=1}^nF(f_i).
\]
A probability measure $\mu'$ on $(\mathcal S',\mathcal C)$ such that $\int\mathrm{d}\mu'(F)\,\prod_{i=1}^nF(f_i)=S_n(f_1,\dots,f_n)$ for all such $n$ and $f_1,\dots,f_n$ is equal to $\mu$.
\item[(9)] \emph{Symmetry under $F\mapsto-F$.} Suppose that every law $\mathrm{Pr}_{N_k,T_k}$ is invariant under the map $F\mapsto-F$. Then so is $\mu$, and $S_n=0$ for every odd $n$. This hypothesis is satisfied when $\mu$ is a limit point of the averaged laws (Definition~\ref{def:limit-points}), by Lemma~\ref{lem:invariances-averaged-law}(c).
\end{enumerate}
\end{theorem}

\begin{proof}[Proof of Theorem~\ref{thm:continuum-infinite-1}, parts (1)--(3)]
Since $T_k\to\infty$, the periods $R_k=h\cdot L^{T_k}$ tend to infinity. For $f\in\mathcal S$ the map $F\mapsto F(f)$ is a continuous linear form on $H_{-4}$, because $|F(f)|\le\|F\|_{-4}\cdot\|f\|_4$ (Section~\ref{sec:hermite-sobolev-spaces}). By the continuous mapping theorem, for every continuous map $\phi$ from $H_{-4}$ into a metric space the image measures $\mathsf P_{N_k,T_k}\circ\phi^{-1}$ tend weakly to $\mu\circ\phi^{-1}$. By Definition~\ref{def:limit-points}, which rests on~\eqref{eq:continuum-infinite-3}, the law of $\left(X_k(g_1),\dots,X_k(g_m)\right)$ under $\mathrm{Pr}_{N_k,T_k}$ is the law of $\left(F(g_1),\dots,F(g_m)\right)$ under $\mathsf P_{N_k,T_k}$ for real tests $g_1,\dots,g_m$; applied to the real and imaginary parts of complex tests, this identifies $E_k\Phi\left(X_k(f_1),\dots,X_k(f_s)\right)$ with $\int\mathrm{d}\mathsf P_{N_k,T_k}(F)\,\Phi\left(F(f_1),\dots,F(f_s)\right)$ also for complex tests $f_1,\dots,f_s$, whenever $\Phi$ is a Borel function on $\mathbb C^s$ that is non-negative or for which one side converges absolutely. Each $\mathrm{Pr}_{N_k,T_k}$ has property~\eqref{eq:continuum-infinite-Ad2} (Definition~\ref{def:admissible-families}), so Theorem~\ref{thm:estimates-fixed-cutoff} applies to it. For a test $f$ and $k\in\mathbb N$ let $f_{\eta_k}$ be the cochain $f_\eta$ of~\eqref{eq:continuum-infinite-2} at the cutoff $(N_k,T_k)$, that is, with $\eta=\eta_k$ and $R=R_k$; for an indexed test $f_i$ this cochain is written $f_{i,\eta_k}=(f_i)_{\eta_k}$.

\emph{Part (1).} Let $f\in\mathcal S$ be real and $M>0$. The function $F\mapsto\min\left(\exp(F(f)),M\right)$ is bounded and continuous on $H_{-4}$. By Theorem~\ref{thm:estimates-fixed-cutoff} and Lemma~\ref{lem:norm-bounds-lattice},
\[
E_k\exp\left(X_k(f)\right)\le\exp\left(\frac{1}{2}|f_{\eta_k}|_{\eta_k}^2\right)\le\exp\left(\frac{1}{2}p_{R_k}(f)^2\right),
\]
and $p_{R_k}(f)\to\|f\|_{L^2}$ as $R_k\to\infty$ (Lemma~\ref{lem:norm-bounds-lattice}). Therefore
\[
\int\mathrm{d}\mu(F)\,\min\left(\exp(F(f)),M\right)=\lim_{k\to\infty}E_k\min\left(\exp(X_k(f)),M\right)\le\liminf_{k\to\infty}\exp\left(\frac{1}{2}p_{R_k}(f)^2\right)=\exp\left(\frac{1}{2}\|f\|_{L^2}^2\right).
\]
Letting $M\uparrow\infty$, monotone convergence gives the first inequality of~(1). In other words, $\mu$ satisfies~\eqref{eq:continuum-infinite-L1}.

For the probability space $(\mathcal S',\mathcal C,\mu)$, the real space $\mathcal T=\mathcal S$ with the Hilbert seminorm $\|\cdot\|_{L^2}$ and the linear map $f\mapsto F(f)$, the hypothesis~\eqref{eq:continuum-infinite-5} of Proposition~\ref{prop:sub-gaussian-random} is~\eqref{eq:continuum-infinite-L1}. Part~(2) of that proposition gives
\[
\int\mathrm{d}\mu(F)\,F(f)^2\le\|f\|_{L^2}^2\qquad(f\in\mathcal S\ \text{real}).
\]
Thus $f\mapsto F(f)$ is a bounded linear map from $\mathcal S$, endowed with the norm of $L^2(\mathbb R^3;\Lambda^2)$, into $L^2(\mu)$. Since $\mathcal S$ is dense in $L^2(\mathbb R^3;\Lambda^2)$, this map has a unique bounded linear extension $X\colon L^2(\mathbb R^3;\Lambda^2)\to L^2(\mu)$. Let $f\in L^2(\mathbb R^3;\Lambda^2)$ be real, and let $f_j\in\mathcal S$ be real tests with $f_j\to f$ in $L^2$. Then $F(f_j)=X(f_j)\to X(f)$ in $L^2(\mu)$, so a subsequence $X(f_{j_l})$ converges to $X(f)$ $\mu$-almost surely, and Fatou's lemma together with~\eqref{eq:continuum-infinite-L1} gives
\[
\int\mathrm{d}\mu\,\exp\left(X(f)\right)\le\liminf_{l\to\infty}\int\mathrm{d}\mu(F)\,\exp\left(F(f_{j_l})\right)\le\lim_{l\to\infty}\exp\left(\frac{1}{2}\|f_{j_l}\|_{L^2}^2\right)=\exp\left(\frac{1}{2}\|f\|_{L^2}^2\right).
\]
Hence~\eqref{eq:continuum-infinite-5} holds for $X$ on the real space $L^2(\mathbb R^3;\Lambda^2)$ with $|f|=\|f\|_{L^2}$, and Proposition~\ref{prop:sub-gaussian-random} applies under $\mu$.

\emph{Part (2).} Let $f_1,\dots,f_s\in\mathcal S$ be complex, let $z\in\mathbb C^s$, and let $g_z=\operatorname{Re}\sum_iz_i\cdot f_i$, a real test. Since $X_k$ is complex-linear and real on real tests,
\[
\left|\exp\left(\sum_{i=1}^sz_i\cdot X_k(f_i)\right)\right|^2=\exp\left(2X_k(g_z)\right).
\]
By Theorem~\ref{thm:estimates-fixed-cutoff} and Lemma~\ref{lem:norm-bounds-lattice}, with $|(g_z)_{\eta_k}|_{\eta_k}\le p_{R_k}(g_z)\le p_1(g_z)$,
\[
\sup_kE_k\left|\exp\left(\sum_{i=1}^sz_i\cdot X_k(f_i)\right)\right|^2=\sup_kE_k\exp\left(2X_k(g_z)\right)\le\exp\left(2p_1(g_z)^2\right),
\]
and the right side is bounded for $z$ in a compact subset of $\mathbb C^s$, $p_1$ being a seminorm. The real and imaginary parts of $F\mapsto\exp\left(\sum_iz_i\cdot F(f_i)\right)$ are continuous on $H_{-4}$, so their laws under $\mathsf P_{N_k,T_k}$ converge weakly to their laws under $\mu$; by the bound just obtained, these real random variables are uniformly integrable. For a uniformly integrable sequence of real random variables that converges in law, the expectations converge to the expectation of the limit law; hence $G_k(z)\to G(z)$ for every $z\in\mathbb C^s$.

By Proposition~\ref{prop:sub-gaussian-random}(1), applied under $\mathrm{Pr}_{N_k,T_k}$ by means of Theorem~\ref{thm:estimates-fixed-cutoff}, the functions $G_k$ are entire, and $|G_k(z)|\le\exp\left(p_1(g_z)^2\right)$ by the Cauchy--Schwarz inequality; so the $G_k$ are bounded on each compact set uniformly in $k$. By the Vitali--Porter theorem the convergence $G_k\to G$ is locally uniform on $\mathbb C^s$, and all derivatives converge locally uniformly as well. (A direct argument: Cauchy's estimates make the $G_k$ equi-Lipschitz on every compact set, which turns pointwise into locally uniform convergence; Cauchy's integral formula then yields the convergence of the derivatives.) The derivatives of $G_k$ and of $G$ at $z=0$ are the moments $E_k\prod_iX_k(f_i)^{\alpha_i}$ and $\int\mathrm{d}\mu(F)\,\prod_iF(f_i)^{\alpha_i}$ (the exponential moments of Proposition~\ref{prop:sub-gaussian-random}(1) allow differentiation under the integral); so all moments converge.

Let $K$ be a compact subset of $\left\{z\in\mathbb C^s:\sum_i|z_i|\cdot\|f_i\|_{L^2}<\frac{2}{5}\right\}$. The proof of Lemma~\ref{lem:norm-bounds-lattice} applies verbatim to complex tests; hence
\[
\sum_{i=1}^s|z_i|\cdot|f_{i,\eta_k}|_{\eta_k}\le\sum_{i=1}^s|z_i|\cdot p_{R_k}(f_i),
\]
and, since $p_{R_k}(f_i)\to\|f_i\|_{L^2}$ for each $i$ and $K$ is compact, the right side is at most $\frac{2}{5}$ on $K$ for all large $k$. For such $k$ the triangle inequality gives $\left|\left(\sum_iz_i\cdot f_i\right)_{\eta_k}\right|_{\eta_k}\le\frac{2}{5}$ on $K$, and Proposition~\ref{prop:sub-gaussian-random}(6) under $\mathrm{Pr}_{N_k,T_k}$ (Theorem~\ref{thm:estimates-fixed-cutoff}), applied to the test $\sum_iz_i\cdot f_i$, gives $|G_k-1|\le\frac{1}{2}$ on $K$; hence also $|G-1|\le\frac{1}{2}$ on $K$. The principal branch of the logarithm is uniformly continuous on $\left\{w\in\mathbb C:|w-1|\le\frac{1}{2}\right\}$, so $W_k=\log G_k$ converges to $\log G=W$ uniformly on $K$.

Now let $r>0$ satisfy $r\cdot\sum_i\|f_i\|_{L^2}<\frac{2}{5}$, and let
\[
D_r=\left\{z\in\mathbb C^s:|z_i|\le r\ \text{for}\ 1\le i\le s\right\}
\]
be the closed polydisc of radius $r$ centred at $0$. It is a compact subset of the open set $\left\{z\in\mathbb C^s:\sum_i|z_i|\cdot\|f_i\|_{L^2}<\frac{2}{5}\right\}$, so what has been shown for $K$ applies to $K=D_r$: for all large $k$ one has $|G_k-1|\le\frac{1}{2}$ on $D_r$, and $W_k\to W$ uniformly on $D_r$. For such $k$ the function $W_k=\log G_k$ is analytic on a neighbourhood of $D_r$, because $G_k$ is entire and the principal branch of the logarithm is analytic on the half-plane $\left\{w\in\mathbb C:\operatorname{Re}w>0\right\}$, which contains the disc $\left\{w\in\mathbb C:|w-1|\le\frac{1}{2}\right\}$. For the same reason $W=\log G$ is analytic on a neighbourhood of $D_r$: the function $G$ is entire, as a locally uniform limit of entire functions, and $|G-1|\le\frac{1}{2}$ on $D_r$. Cauchy's integral formula on $D_r$ expresses each derivative at $0$ of $W_k$, and of $W$, as the integral of that function over the torus $\left\{z\in\mathbb C^s:|z_1|=\dots=|z_s|=r\right\}$ against a bounded kernel that depends only on the order of the derivative. By the uniform convergence on $D_r$, the derivatives of $W_k$ at $0$, which are the cumulants, therefore converge to the corresponding derivatives of $W$.

\emph{Part (3).} Let $s>\frac{3}{2}$; then $\mathsf Z_{2s}<\infty$ by~\eqref{eq:continuum-infinite-4}. The Hermite functions $h_\kappa$ of Section~\ref{sec:hermite-sobolev-spaces} are real with $\|h_\kappa\|_{L^2}=1$, so part~(1) and Proposition~\ref{prop:sub-gaussian-random}(2) give $\int\mathrm{d}\mu(F)\,F(h_\kappa)^2\le1$ for every $\kappa\in\mathcal K$. By monotone convergence,
\[
\int\mathrm{d}\mu(F)\,\|F\|_{-s}^2=\sum_{\kappa\in\mathcal K}\lambda_\kappa^{-2s}\cdot\int\mathrm{d}\mu(F)\,F(h_\kappa)^2\le\sum_{\kappa\in\mathcal K}\lambda_\kappa^{-2s}=\mathsf Z_{2s}.
\]
Let $0\le2a\cdot\mathsf Z_{2s}<1$, and let $\pi_\kappa=\frac{\lambda_\kappa^{-2s}}{\mathsf Z_{2s}}$ for $\kappa\in\mathcal K$; these are probability weights on $\mathcal K$. Jensen's inequality for the convex function $\exp$ gives
\[
\exp\left(a\cdot\|F\|_{-s}^2\right)=\exp\left(\sum_{\kappa\in\mathcal K}\pi_\kappa\cdot a\cdot\mathsf Z_{2s}\cdot F(h_\kappa)^2\right)\le\sum_{\kappa\in\mathcal K}\pi_\kappa\cdot\exp\left(a\cdot\mathsf Z_{2s}\cdot F(h_\kappa)^2\right).
\]
Integrate term by term, the terms being non-negative, and apply Proposition~\ref{prop:sub-gaussian-random}(3) under $\mu$ to the test $h_\kappa$, for which $2a\cdot\mathsf Z_{2s}\cdot\|h_\kappa\|_{L^2}^2<1$:
\[
\int\mathrm{d}\mu(F)\,\exp\left(a\cdot\|F\|_{-s}^2\right)\le\sum_{\kappa\in\mathcal K}\pi_\kappa\cdot\left(1-2a\cdot\mathsf Z_{2s}\right)^{-\frac{1}{2}}=\left(1-2a\cdot\mathsf Z_{2s}\right)^{-\frac{1}{2}}.
\]
Finally, the finiteness of $\int\mathrm{d}\mu(F)\,\|F\|_{-s}^2$ implies $\|F\|_{-s}<\infty$ for $\mu$-almost every $F$, that is, $\mu(H_{-s})=1$.
\end{proof}

\begin{proof}[Proof of Theorem~\ref{thm:continuum-infinite-1}, parts (4)--(9)]
Let $(N_k,T_k)$, $\eta_k$, $R_k$, $E_k$ and $X_k$ be as in the statement of the theorem. Then $\eta_k=L^{-N_k}\to0$ and $R_k=h\cdot L^{T_k}\to\infty$. For real $f\in\mathcal S$ the linear form $F\mapsto F(f)$ is continuous on $H_{-4}$, since $|F(f)|\le\|F\|_{-4}\cdot\|f\|_4$ (Section~\ref{sec:hermite-sobolev-spaces}). Let $\chi\colon\mathbb R\to\mathbb C$ be bounded and continuous. Then $F\mapsto\chi\left(F(f)\right)$ is bounded and continuous on $H_{-4}$, and the weak convergence $\mathsf P_{N_k,T_k}\to\mu$, together with the identity for $E_k$ in Definition~\ref{def:limit-points}, gives the limit relation
\[
\int\mathrm{d}\mu(F)\,\chi\left(F(f)\right)=\lim_{k\to\infty}\int\mathrm{d}\mathsf P_{N_k,T_k}(F)\,\chi\left(F(f)\right)=\lim_{k\to\infty}E_k\,\chi\left(X_k(f)\right)\qquad(f\in\mathcal S\ \text{real}).
\]
In particular
\[
\Phi_\mu(f)=\lim_{k\to\infty}E_k\exp\left(iX_k(f)\right)\qquad(f\in\mathcal S\ \text{real}).
\]

\emph{Part (4).} Let $g\in C_c^\infty(\mathbb R^3)$ be real, and let $d^*g$ and $u^{(\eta)}$ be the two-forms of Lemma~\ref{lem:discrete-codifferentials-three}. Both $d^*g$ and $u^{(\eta_k)}$ vanish at every point whose distance from $\operatorname{supp}g$ exceeds $\frac{\eta_k}{2}\le\frac12$; since $R_k\to\infty$, both are supported in $(-R_k/2,R_k/2)^3$ for all large $k$. The measure $\mathrm{Pr}_{N_k,T_k}$ is concentrated on $\ker d$ (Definition~\ref{def:admissible-families}), and $X_k\left(u^{(\eta_k)}\right)$ vanishes on $\ker d$ by Lemma~\ref{lem:discrete-codifferentials-three}; hence $X_k(d^*g)=X_k\left(d^*g-u^{(\eta_k)}\right)$ almost surely under $\mathrm{Pr}_{N_k,T_k}$. Put
\[
C_g=\left(\frac{1}{24}\sum_{l=1}^3\left\|\partial_l^3g\right\|_{L^2}\right)^2 .
\]
Then, for all large $k$,
\[
E_k\left[X_k(d^*g)^2\right]=E_k\left[X_k\left(d^*g-u^{(\eta_k)}\right)^2\right]\le\left|\left(d^*g-u^{(\eta_k)}\right)_{\eta_k}\right|_{\eta_k}^2\le\left\|d^*g-u^{(\eta_k)}\right\|_{L^2}^2\le C_g\cdot\eta_k^4 .
\]
Here the first inequality is Proposition~\ref{prop:sub-gaussian-random}(2) under $\mathrm{Pr}_{N_k,T_k}$, which holds by Theorem~\ref{thm:estimates-fixed-cutoff} because $\mathrm{Pr}_{N_k,T_k}$ has property~\eqref{eq:continuum-infinite-Ad2}; the second is the bound of Lemma~\ref{lem:norm-bounds-lattice} for tests supported in $(-R_k/2,R_k/2)^3$; the third is the estimate of Lemma~\ref{lem:discrete-codifferentials-three}. Let $M>0$. The limit relation, applied to the bounded continuous function $\chi(x)=\min\left(x^2,M\right)$, gives
\[
\int\mathrm{d}\mu(F)\,\min\left(F(d^*g)^2,M\right)=\lim_{k\to\infty}E_k\min\left(X_k(d^*g)^2,M\right)\le\limsup_{k\to\infty}C_g\cdot\eta_k^4=0 .
\]
As $M>0$ is arbitrary, $F(d^*g)=0$ for $\mu$-almost every $F$. Thus $\mu$ satisfies~\eqref{eq:continuum-infinite-L2}.

The real space $\mathcal S(\mathbb R^3)$ is a separable metric space in which the real functions of $C_c^\infty(\mathbb R^3)$ are dense. Hence there is a sequence $(\zeta_m)_{m\ge1}$ of real functions in $C_c^\infty(\mathbb R^3)$ which is dense in $\mathcal S(\mathbb R^3)$. By~\eqref{eq:continuum-infinite-L2} the event
\[
\bigcap_{m\ge1}\left\{F\in\mathcal S':F\left(d^*\zeta_m\right)=0\right\}
\]
has $\mu$-probability one. Let $F$ belong to this event, let $g\in\mathcal S(\mathbb R^3)$ be real, and choose indices $m_i$ with $\zeta_{m_i}\to g$ in $\mathcal S(\mathbb R^3)$. Since $d^*\colon\mathcal S(\mathbb R^3)\to\mathcal S(\mathbb R^3;\Lambda^2)$ and $F\colon\mathcal S(\mathbb R^3;\Lambda^2)\to\mathbb R$ are continuous, $F(d^*g)=\lim_{i\to\infty}F\left(d^*\zeta_{m_i}\right)=0$; by linearity $F(d^*g)=0$ also for complex $g\in\mathcal S(\mathbb R^3)$. Hence the event above is the set $\{F\in\mathcal S':dF=0\}$, and $\mu\left(\{F\in\mathcal S':dF=0\}\right)=1$.

Let $f\in L^2(\mathbb R^3;\Lambda^2)$ be real, and let $b=\ast f$ be its dual vector field. Normalise the Fourier transform as
\[
\hat w(p)=(2\pi)^{-\frac32}\int\mathrm{d}y\,\exp\left(-i\,p\cdot y\right)\cdot w(y),
\]
so that it is unitary on $L^2$ by Plancherel's theorem and $\widehat{\partial_lw}(p)=i\,p_l\cdot\hat w(p)$; let $\check w$ denote the inverse transform. For $p\in\mathbb R^3\setminus\{0\}$ put
\[
\varphi(p)=-\frac{i}{|p|^2}\sum_{l=1}^3p_l\cdot\hat b_l(p).
\]
The operator $1-\Pi_c$ acts on $\hat b$ as multiplication by the matrix $\left(\frac{p_jp_l}{|p|^2}\right)_{j,l}$; hence
\[
\widehat{\ast(1-\Pi_c)f}(p)=\left(\sum_{l=1}^3p_l\cdot\hat b_l(p)\right)\frac{p}{|p|^2}=i\,\varphi(p)\,p\qquad(p\ne0).
\]
By the Cauchy--Schwarz inequality in $\mathbb C^3$ one has $|p|^2\cdot|\varphi(p)|^2\le|\hat b(p)|^2$, so that
\[
\int\mathrm{d}p\,|p|^2\cdot|\varphi(p)|^2\le\left\|\hat b\right\|_{L^2}^2 ;
\]
moreover $\varphi(-p)=\overline{\varphi(p)}$, because $b$ is real. The space $C_c^\infty(\mathbb R^3\setminus\{0\})$ of complex functions is dense in $L^2\left(\mathbb R^3,|p|^2\,\mathrm{d}p\right)$; choose $\varphi_j$ in it with $\varphi_j\to\varphi$ in $L^2\left(\mathbb R^3,|p|^2\,\mathrm{d}p\right)$. The map $\psi\mapsto\overline{\psi(-\,\cdot\,)}$ is an isometry of this space which maps $C_c^\infty(\mathbb R^3\setminus\{0\})$ onto itself and fixes $\varphi$; after replacing $\varphi_j$ by $\frac12\left(\varphi_j+\overline{\varphi_j(-\,\cdot\,)}\right)$, which is at least as close to $\varphi$, one has in addition $\varphi_j(-p)=\overline{\varphi_j(p)}$. Then $g_j=\check\varphi_j$ is a real function in $\mathcal S(\mathbb R^3)$, $\widehat{\nabla g_j}(p)=i\,\varphi_j(p)\,p$, and Plancherel's theorem gives
\[
\left\|\nabla g_j-\ast(1-\Pi_c)f\right\|_{L^2}^2=\int\mathrm{d}p\,|p|^2\cdot\left|\varphi_j(p)-\varphi(p)\right|^2\to0\qquad(j\to\infty).
\]
Since $\ast d^*g_j=-\nabla g_j$ and $f\mapsto\ast f$ preserves the $L^2$ norm, $-d^*g_j\to(1-\Pi_c)f$ in $L^2(\mathbb R^3;\Lambda^2)$. By the preceding paragraph $X(d^*g_j)=F(d^*g_j)=0$ for $\mu$-almost every $F$, and $X\colon L^2(\mathbb R^3;\Lambda^2)\to L^2(\mu)$ is bounded by Theorem~\ref{thm:continuum-infinite-1}(1). Hence
\[
X\left((1-\Pi_c)f\right)=\lim_{j\to\infty}X\left(-d^*g_j\right)=0\qquad\text{in }L^2(\mu),
\]
that is, $X(f)=X(\Pi_cf)$ $\mu$-almost surely; by linearity the same holds for complex $f$.

The matrix $\left(\delta_{jl}-\frac{p_jp_l}{|p|^2}\right)_{j,l}$ that defines $\Pi_c$ is real and even in $p$, so $\Pi_cf$ is real for real $f$; and since $\Pi_c$ is an orthogonal projection, $\|\Pi_cf\|_{L^2}^2=\langle f,\Pi_cf\rangle_{L^2}=C_0(f,f)$. Theorem~\ref{thm:continuum-infinite-1}(1), applied to $\Pi_cf$, gives
\[
\int\mathrm{d}\mu\,\exp\left(X(f)\right)=\int\mathrm{d}\mu\,\exp\left(X(\Pi_cf)\right)\le\exp\left(\frac12\|\Pi_cf\|_{L^2}^2\right)=\exp\left(\frac12C_0(f,f)\right).
\]
Thus~\eqref{eq:continuum-infinite-5} holds under $\mu$ for the real space $L^2(\mathbb R^3;\Lambda^2)$, the map $X$ and the Hilbert seminorm $f\mapsto C_0(f,f)^{\frac12}=\|\Pi_cf\|_{L^2}$, and Proposition~\ref{prop:sub-gaussian-random}(2) gives $0\le\int\mathrm{d}\mu\,X(f)^2\le C_0(f,f)$. Finally $C_0(f,f)=\|\Pi_cf\|_{L^2}^2\le\|f\|_{L^2}^2$.

\emph{Part (5).} Let $f\in\mathcal S$ be real and $b\in\mathbb R^3$, and for every $k$ choose $b_k\in\eta_k\mathbb Z^3$ with $|b_k-b|_\infty\le\frac{\eta_k}{2}$. Since $\mathrm{Pr}_{N_k,T_k}$ has property~\eqref{eq:continuum-infinite-Ad1}, Lemma~\ref{lem:invariances-averaged-law}(d), applied with $v=b_k$ to the bounded function $x\mapsto\exp(ix)$, gives
\[
E_k\exp\left(iX_k(\tau_{b_k}f)\right)=E_k\exp\left(iX_k(f)\right).
\]
For real $x$ and $y$ one has $|\exp(ix)-\exp(iy)|\le|x-y|$. Hence the Cauchy--Schwarz inequality, Proposition~\ref{prop:sub-gaussian-random}(2) under $\mathrm{Pr}_{N_k,T_k}$ (Theorem~\ref{thm:estimates-fixed-cutoff}) and Lemma~\ref{lem:norm-bounds-lattice} give
\begin{align*}
\left|E_k\exp\left(iX_k(\tau_bf)\right)-E_k\exp\left(iX_k(\tau_{b_k}f)\right)\right|&\le\left(E_k\left[X_k\left(\tau_bf-\tau_{b_k}f\right)^2\right]\right)^{\frac12}\\
&\le\left|\left(\tau_bf-\tau_{b_k}f\right)_{\eta_k}\right|_{\eta_k}\le p_1\left(\tau_bf-\tau_{b_k}f\right).
\end{align*}
Since $\eta_k\to0$, one has $b_k\to b$. The map $b'\mapsto\tau_{b'}f$ is continuous from $\mathbb R^3$ into $\mathcal S$, and $p_1$ is a continuous seminorm on $\mathcal S$. Indeed, for $w\in\mathcal S$ the proof of Lemma~\ref{lem:norm-bounds-lattice} gives $\left\|\left(1+|y|^2\right)\cdot w\right\|_{L^2}\le\left(\frac{11}{8}\right)^{\frac12}\cdot\|\mathsf Aw\|_{L^2}$ and $\|w\|_{L^2}\le\frac14\|\mathsf Aw\|_{L^2}$, and $\|\mathsf Aw\|_{L^2}=\|w\|_1$ by Parseval's identity, as $\langle h_\kappa,\mathsf Aw\rangle_{L^2}=\lambda_\kappa\cdot\langle h_\kappa,w\rangle_{L^2}$ (Section~\ref{sec:hermite-sobolev-spaces}); therefore
\[
p_1(w)\le\|w\|_{L^2}+26\cdot\left\|\left(1+|y|^2\right)\cdot w\right\|_{L^2}\le27\cdot\left(\frac{11}{8}\right)^{\frac12}\cdot\|w\|_1 .
\]
Consequently the right side of the preceding estimate tends to $0$ as $k\to\infty$, and the limit relation for $\Phi_\mu$ yields
\[
\Phi_\mu(\tau_bf)=\lim_{k\to\infty}E_k\exp\left(iX_k(\tau_bf)\right)=\lim_{k\to\infty}E_k\exp\left(iX_k(\tau_{b_k}f)\right)=\lim_{k\to\infty}E_k\exp\left(iX_k(f)\right)=\Phi_\mu(f).
\]
This is the first of the two equalities in~\eqref{eq:continuum-infinite-L3}.

Let $f_1,\dots,f_s\in\mathcal S$ be real and $t\in\mathbb R^s$. The equality just proved, for $f=\sum_{i=1}^st_i\cdot f_i$, states that the random vectors $\left(F(\tau_bf_1),\dots,F(\tau_bf_s)\right)$ and $\left(F(f_1),\dots,F(f_s)\right)$ have the same characteristic function under $\mu$; by the uniqueness theorem for characteristic functions on $\mathbb R^s$ they have the same law. As $(\tau_b^*F)(f_i)=F(\tau_bf_i)$, the measures $\mu\circ(\tau_b^*)^{-1}$ and $\mu$ coincide on the cylinder sets $\left\{F:\left(F(f_1),\dots,F(f_s)\right)\in B\right\}$, $B\subset\mathbb R^s$ a Borel set. These sets form a $\pi$-system, and the $\sigma$-algebra generated by them is $\mathcal C$. By Dynkin's $\pi$-$\lambda$ theorem, $\mu\circ(\tau_b^*)^{-1}=\mu$.

Consequently each $\mathcal U_b$ is an isometry of $L^2(\mu)$. For $b,b'\in\mathbb R^3$ and $f\in\mathcal S$ one has $(\tau_{b'}^*\tau_b^*F)(f)=F(\tau_b\tau_{b'}f)=F(\tau_{b+b'}f)$; hence $\mathcal U_b\mathcal U_{b'}=\mathcal U_{b+b'}$, and together with $\mathcal U_0=1$ it follows that $\mathcal U_b$ is unitary with inverse $\mathcal U_{-b}$. For real $f\in\mathcal S$, the inequality $|\exp(ix)-\exp(iy)|\le|x-y|$ and Proposition~\ref{prop:sub-gaussian-random}(2) under $\mu$ (Theorem~\ref{thm:continuum-infinite-1}(1)) give
\[
\left\|\mathcal U_b\exp\left(iF(f)\right)-\exp\left(iF(f)\right)\right\|_{L^2(\mu)}^2\le\int\mathrm{d}\mu(F)\,F\left(\tau_bf-f\right)^2\le\left\|\tau_bf-f\right\|_{L^2}^2\to0\qquad(b\to0).
\]
By Lemma~\ref{lem:totality-exponentials}, applied to the real space $\mathcal D_0=\mathcal S$, for which $\Sigma(\mathcal S)=\mathcal C$, the linear span of the functions $F\mapsto\exp\left(iF(f)\right)$, $f\in\mathcal S$ real, is dense in $L^2(\mu)$. Since $\|\mathcal U_b\|=1$ for every $b$, the set of the $\Psi\in L^2(\mu)$ with $\left\|\mathcal U_b\Psi-\Psi\right\|_{L^2(\mu)}\to0$ as $b\to0$ is a closed linear subspace of $L^2(\mu)$; it contains these functions and is therefore all of $L^2(\mu)$. Finally $\left\|\mathcal U_b\Psi-\mathcal U_{b_0}\Psi\right\|_{L^2(\mu)}=\left\|\mathcal U_{b-b_0}\Psi-\Psi\right\|_{L^2(\mu)}$ for $b,b_0\in\mathbb R^3$, so $b\mapsto\mathcal U_b\Psi$ is continuous for every $\Psi\in L^2(\mu)$.

\emph{Part (6).} Let $r\in O_{\mathrm{cub}}$ and let $f\in\mathcal S$ be real. Since $\mathrm{Pr}_{N_k,T_k}$ has property~\eqref{eq:continuum-infinite-Ad1}, Lemma~\ref{lem:invariances-averaged-law}(d) gives
\[
E_k\exp\left(iX_k(U_rf)\right)=E_k\exp\left(iX_k(f)\right)
\]
for every $k$, with no error term. By the limit relation for $\Phi_\mu$, the left side tends to $\Phi_\mu(U_rf)$ and the right side to $\Phi_\mu(f)$; hence $\Phi_\mu(U_rf)=\Phi_\mu(f)$. This is the second of the two equalities in~\eqref{eq:continuum-infinite-L3}. Let $f_1,\dots,f_s\in\mathcal S$ be real and $t\in\mathbb R^s$. This equality for $f=\sum_{i=1}^st_i\cdot f_i$ shows that $\left(F(U_rf_1),\dots,F(U_rf_s)\right)$ and $\left(F(f_1),\dots,F(f_s)\right)$ have the same characteristic function, hence the same law, under $\mu$. Since $(U_r^*F)(f_i)=F(U_rf_i)$, the measures $\mu\circ(U_r^*)^{-1}$ and $\mu$ coincide on the cylinder sets, and Dynkin's $\pi$-$\lambda$ theorem gives $\mu\circ(U_r^*)^{-1}=\mu$.

\emph{Part (7).} Fix $\nu\in\{1,2,3\}$, and let first $c=0$. Write $\mathsf R=\mathsf R_\nu$, $\Theta=\Theta_{\nu,0}$, $\mathcal D^+=\mathcal D^+_{\nu,0}$ and $\Sigma^+=\Sigma^+_{\nu,0}$. Let $f_1,\dots,f_s\in\mathcal D^+$ and $c_1,\dots,c_s\in\mathbb C$. The supports of $f_1,\dots,f_s$ are compact subsets of the open half-space $\{y_\nu>0\}$; hence there are $\delta>0$ and $\rho>0$ with
\[
\operatorname{supp}f_i\subset\{y\in\mathbb R^3:y_\nu\ge\delta\}\cap[-\rho,\rho]^3\qquad(i=1,\dots,s).
\]
For all large $k$ one has $\eta_k\le\delta$ and $R_k\ge2\rho+3\eta_k$. For such $k$ and for $i=1,\dots,s$ let $f_{i,\eta_k}=(f_i)_{\eta_k}$ be the real cochain that~\eqref{eq:continuum-infinite-2} assigns to $f_i$ at the cutoff $(N_k,T_k)$, that is, with $\eta=\eta_k$ and $R=R_k$. Lemma~\ref{lem:tests-supported-half} shows that the cochains $f_{1,\eta_k},\dots,f_{s,\eta_k}$ are supported on $\mathcal P^{+\circ}_\nu$, and Lemma~\ref{lem:lattice-symmetries-cell} gives $(\mathsf Rf_i)_{\eta_k}=\hat{\mathsf R}_\nu f_{i,\eta_k}$, so that
\[
X_k\left(f_j-\mathsf Rf_i\right)=\left\langle f_{j,\eta_k}-\hat{\mathsf R}_\nu f_{i,\eta_k},F\right\rangle_{\eta_k}\qquad(i,j=1,\dots,s).
\]
Property~\eqref{eq:continuum-infinite-Ad3} of $\mathrm{Pr}_{N_k,T_k}$, applied to the cochains $f_{1,\eta_k},\dots,f_{s,\eta_k}$ and to the coefficients $c_1,\dots,c_s$, therefore gives
\[
\sum_{i,j=1}^s\overline{c_i}\cdot c_j\cdot E_k\exp\left(iX_k\left(f_j-\mathsf Rf_i\right)\right)\ge0 .
\]
By the limit relation for $\Phi_\mu$, each expectation tends to $\Phi_\mu\left(f_j-\mathsf Rf_i\right)$ as $k\to\infty$, whence
\[
\sum_{i,j=1}^s\overline{c_i}\cdot c_j\cdot\Phi_\mu\left(f_j-\mathsf Rf_i\right)\ge0 .
\]
This is~\eqref{eq:continuum-infinite-L4}. For $A=\sum_{j=1}^sc_j\cdot\exp\left(iF(f_j)\right)$ one has
\[
(\Theta A)(F)=\overline{A(\mathsf R^*F)}=\sum_{i=1}^s\overline{c_i}\cdot\exp\left(-iF(\mathsf Rf_i)\right),
\]
so the last inequality reads $\int\mathrm{d}\mu\,(\Theta A)\cdot A\ge0$.

By part~(6) for $r=r_\nu$, the map $\mathsf R^*=U_{r_\nu}^*$ preserves $\mu$; moreover $\mathsf R^*\mathsf R^*F=F$, because $\mathsf R\mathsf Rf=f$ for every test $f$. Hence $\Theta$ is an antiunitary involution of $L^2(\mu)$, and the sesquilinear form
\[
\mathsf b(A,B)=\int\mathrm{d}\mu\,(\Theta A)\cdot B\qquad\left(A,B\in L^2(\mu)\right)
\]
satisfies, by the Cauchy--Schwarz inequality and $\|\Theta A\|_{L^2(\mu)}=\|A\|_{L^2(\mu)}$,
\[
\left|\mathsf b(A,B)\right|\le\|A\|_{L^2(\mu)}\cdot\|B\|_{L^2(\mu)}\qquad\left(A,B\in L^2(\mu)\right).
\]
In particular $A\mapsto\mathsf b(A,A)$ is continuous on $L^2(\mu)$. By Lemma~\ref{lem:totality-exponentials}, applied to the real space $\mathcal D_0=\mathcal D^+$, the functions $A$ above are dense in $L^2\left(\mathcal S',\Sigma^+,\mu\right)$. Hence $\int\mathrm{d}\mu\,(\Theta A)\cdot A\ge0$ for every $A\in L^2\left(\mathcal S',\Sigma^+,\mu\right)$.

Let now $c\in\mathbb R$ be arbitrary, and let $A\in L^2\left(\mathcal S',\Sigma^+_{\nu,c},\mu\right)$ be $\Sigma^+_{\nu,c}$-measurable. One has $\mathcal D^+_{\nu,c}=\left\{\tau_{c\cdot e_\nu}f_0:f_0\in\mathcal D^+_{\nu,0}\right\}$ and $F(\tau_{c\cdot e_\nu}f_0)=(\tau^*_{c\cdot e_\nu}F)(f_0)$. Hence $\Sigma^+_{\nu,c}$ consists of the sets $\left(\tau^*_{c\cdot e_\nu}\right)^{-1}(G)$ with $G\in\Sigma^+_{\nu,0}$. Since $\tau^*_{c\cdot e_\nu}\tau^*_{-c\cdot e_\nu}$ and $\tau^*_{-c\cdot e_\nu}\tau^*_{c\cdot e_\nu}$ are the identity of $\mathcal S'$, the function $B=A\circ\tau^*_{-c\cdot e_\nu}$ is $\Sigma^+_{\nu,0}$-measurable and $A=B\circ\tau^*_{c\cdot e_\nu}$. By part~(5) the measure $\mu$ is invariant under $\tau^*_{c\cdot e_\nu}$, so $B\in L^2\left(\mathcal S',\Sigma^+_{\nu,0},\mu\right)$. For $f_0\in\mathcal S$,
\[
\left(\tau^*_{c\cdot e_\nu}\mathsf R^*_{\nu,c}F\right)(f_0)=F\left(\mathsf R_{\nu,c}\tau_{c\cdot e_\nu}f_0\right)=F\left(\tau_{c\cdot e_\nu}\mathsf R_\nu f_0\right)=\left(\mathsf R^*_\nu\tau^*_{c\cdot e_\nu}F\right)(f_0),
\]
that is, $\tau^*_{c\cdot e_\nu}\mathsf R^*_{\nu,c}=\mathsf R^*_\nu\tau^*_{c\cdot e_\nu}$. Therefore $A(\mathsf R^*_{\nu,c}F)=B(\mathsf R^*_\nu\tau^*_{c\cdot e_\nu}F)$ and $A(F)=B(\tau^*_{c\cdot e_\nu}F)$, and the invariance of $\mu$ under $\tau^*_{c\cdot e_\nu}$, together with the case $c=0$, gives
\[
\int\mathrm{d}\mu\,(\Theta_{\nu,c}A)\cdot A=\int\mathrm{d}\mu(F)\,\overline{B\left(\mathsf R^*_\nu\tau^*_{c\cdot e_\nu}F\right)}\cdot B\left(\tau^*_{c\cdot e_\nu}F\right)=\int\mathrm{d}\mu(F)\,\overline{B\left(\mathsf R^*_\nu F\right)}\cdot B(F)=\int\mathrm{d}\mu\,(\Theta_{\nu,0}B)\cdot B\ge0 .
\]

Let $c\in\mathbb R$, and let $f\in L^2(\mathbb R^3;\Lambda^2)$ be real and vanish almost everywhere on $\{y_\nu<c\}$. As the plane $\{y_\nu=c\}$ is a null set, $f$ is the limit in $L^2$ of real functions in $C_c^\infty(\mathbb R^3;\Lambda^2)$ whose supports lie in the open half-space $\{y_\nu>c\}$, that is, of elements $f_j\in\mathcal D^+_{\nu,c}$. By Theorem~\ref{thm:continuum-infinite-1}(1), $F(f_j)=X(f_j)\to X(f)$ in $L^2(\mu)$, and a subsequence converges $\mu$-almost surely. So $X(f)$ is the $\mu$-almost sure limit of $\Sigma^+_{\nu,c}$-measurable functions and coincides $\mu$-almost surely with a $\Sigma^+_{\nu,c}$-measurable function; by linearity the same holds for complex $f$. By Proposition~\ref{prop:sub-gaussian-random}(1) and~(4), which hold under $\mu$ by Theorem~\ref{thm:continuum-infinite-1}(1), polynomials and exponentials in finitely many such $X(f)$ belong to $L^2(\mu)$, hence to $L^2\left(\mathcal S',\Sigma^+_{\nu,c},\mu\right)$, and the inequality $\int\mathrm{d}\mu\,(\Theta_{\nu,c}A)\cdot A\ge0$ applies to them. For $f\in\mathcal S$ with $\operatorname{supp}f\subset\{y_\nu\ge c\}$ one has $X(f)=F(f)$, and $f$ vanishes on $\{y_\nu<c\}$.

\emph{Part (8).} Let $\mu'$ be a probability measure on $(\mathcal S',\mathcal C)$ such that, for every $n\ge1$ and all real $f_1,\dots,f_n\in C_c^\infty(\mathbb R^3;\Lambda^2)$, the function $F\mapsto\prod_{i=1}^nF(f_i)$ is $\mu'$-integrable and
\[
\int\mathrm{d}\mu'(F)\,\prod_{i=1}^nF(f_i)=S_n(f_1,\dots,f_n).
\]
Let $f\in C_c^\infty(\mathbb R^3;\Lambda^2)$ be real and $t\in\mathbb R$. The series $\cosh x=\sum_{j\ge0}\frac{x^{2j}}{(2j)!}$ has non-negative terms for real $x$, and its term $j=0$ is the constant $1$, whose integral against the probability measure $\mu'$ is $1$. Monotone convergence and the hypothesis on $\mu'$, which covers the terms $j\ge1$, give the equality below. For $j\ge1$, Proposition~\ref{prop:sub-gaussian-random}(4) under $\mu$ (Theorem~\ref{thm:continuum-infinite-1}(1)), applied with $n=2j$ and $g_1=\dots=g_{2j}=f$, gives $S_{2j}(f,\dots,f)\le\sigma_{2j}\cdot\|f\|_{L^2}^{2j}$, and for $j=0$ one has $1\le2=\sigma_0$; this gives the inequality:
\[
\int\mathrm{d}\mu'(F)\,\cosh\left(t\cdot F(f)\right)=1+\sum_{j\ge1}\frac{t^{2j}}{(2j)!}\cdot S_{2j}(f,\dots,f)\le\sum_{j\ge0}\frac{t^{2j}}{(2j)!}\cdot\sigma_{2j}\cdot\|f\|_{L^2}^{2j}<\infty,
\]
where the last series converges for every $t$ because
\[
\frac{\sigma_{2j}}{(2j)!}=\frac{2\cdot2^j\cdot j!}{(2j)!}.
\]
The same computation with $\mu$ in place of $\mu'$ gives the same bound for $\int\mathrm{d}\mu(F)\,\cosh\left(t\cdot F(f)\right)$.

Let $f_1,\dots,f_s\in C_c^\infty(\mathbb R^3;\Lambda^2)$ be real. For $t\in\mathbb R^s$, H\"older's inequality and $\exp|x|\le2\cosh x$ give
\[
\int\mathrm{d}\mu'(F)\,\exp\left(\sum_{i=1}^s|t_i|\cdot|F(f_i)|\right)\le\prod_{i=1}^s\left(\int\mathrm{d}\mu'(F)\,2\cosh\left(s\cdot t_i\cdot F(f_i)\right)\right)^{\frac1s}<\infty,
\]
and likewise for $\mu$. For $z\in\mathbb C^s$ and $\alpha\in\mathbb N_0^s$ write $z^\alpha=\prod_iz_i^{\alpha_i}$, $\alpha!=\prod_i\alpha_i!$ and $|\alpha|=\sum_i\alpha_i$. Since
\[
\sum_{\alpha\in\mathbb N_0^s}\frac{|z^\alpha|}{\alpha!}\prod_{i=1}^s|F(f_i)|^{\alpha_i}=\exp\left(\sum_{i=1}^s|z_i|\cdot|F(f_i)|\right)
\]
is $\mu'$-integrable, the exponential series may be integrated term by term:
\[
\int\mathrm{d}\mu'(F)\,\exp\left(i\sum_{i=1}^sz_i\cdot F(f_i)\right)=\sum_{\alpha\in\mathbb N_0^s}\frac{i^{|\alpha|}\cdot z^\alpha}{\alpha!}\int\mathrm{d}\mu'(F)\,\prod_{i=1}^sF(f_i)^{\alpha_i}\qquad\left(z\in\mathbb C^s\right),
\]
and the same identity holds for $\mu$. For $\alpha\ne0$ the integral on the right is the moment $S_{|\alpha|}$ with the argument $f_i$ repeated $\alpha_i$ times, for $\mu'$ and for $\mu$ alike, and for $\alpha=0$ it equals $1$. Hence the functions $z\mapsto\int\mathrm{d}\mu'(F)\,\exp\left(i\sum_iz_i\cdot F(f_i)\right)$ and $z\mapsto\int\mathrm{d}\mu(F)\,\exp\left(i\sum_iz_i\cdot F(f_i)\right)$ are given by the same power series, convergent on all of $\mathbb C^s$, and coincide. For $z=t\in\mathbb R^s$ they are the characteristic functions of $\left(F(f_1),\dots,F(f_s)\right)$ under $\mu'$ and under $\mu$; by the uniqueness theorem for characteristic functions on $\mathbb R^s$ this vector has the same law under $\mu$ and under $\mu'$. Thus $\mu$ and $\mu'$ coincide on the cylinder sets formed with real tests in $C_c^\infty(\mathbb R^3;\Lambda^2)$; these sets form a $\pi$-system, the $\sigma$-algebra generated by them is $\Sigma\left(C_c^\infty\right)$, and Dynkin's $\pi$-$\lambda$ theorem gives $\mu=\mu'$ on $\Sigma\left(C_c^\infty\right)$. For $F\in\mathcal S'$ and real $f\in\mathcal S$ choose real $f_j\in C_c^\infty(\mathbb R^3;\Lambda^2)$ with $f_j\to f$ in $\mathcal S$; then $F(f)=\lim_{j\to\infty}F(f_j)$. Hence every map $F\mapsto F(f)$, $f\in\mathcal S$, is $\Sigma\left(C_c^\infty\right)$-measurable, so that $\Sigma\left(C_c^\infty\right)=\mathcal C$ and $\mu'=\mu$.

\emph{Part (9).} Suppose that every $\mathrm{Pr}_{N_k,T_k}$ is invariant under $F\mapsto-F$. For real $f\in\mathcal S$ the function $X_k(f)$ is linear on $C^2(\Lambda_{\eta_k,R_k})$, so $X_k(f)(-F)=-X_k(f)(F)$, and therefore
\[
E_k\exp\left(iX_k(f)\right)=E_k\exp\left(-iX_k(f)\right).
\]
The limit relation for $\Phi_\mu$, applied to $f$ and to $-f$, gives $\Phi_\mu(-f)=\Phi_\mu(f)$. Let $f_1,\dots,f_s\in\mathcal S$ be real and $t\in\mathbb R^s$. This equality for $f=\sum_{i=1}^st_i\cdot f_i$ shows that $\left(-F(f_1),\dots,-F(f_s)\right)$ and $\left(F(f_1),\dots,F(f_s)\right)$ have the same characteristic function, hence the same law, under $\mu$. So the image of $\mu$ under $F\mapsto-F$ coincides with $\mu$ on the cylinder sets, and by Dynkin's $\pi$-$\lambda$ theorem on $\mathcal C$: the measure $\mu$ is invariant under $F\mapsto-F$. Consequently, for $n\ge1$ and real $f_1,\dots,f_n\in C_c^\infty(\mathbb R^3;\Lambda^2)$,
\[
S_n(f_1,\dots,f_n)=\int\mathrm{d}\mu(F)\,\prod_{i=1}^n\left(-F(f_i)\right)=(-1)^n\cdot S_n(f_1,\dots,f_n),
\]
and $S_n=0$ for odd $n$. If $\mu$ is a limit point of the averaged laws, then every $\mathrm{Pr}_{N_k,T_k}$ is an averaged law, and these laws are invariant under $F\mapsto-F$ by Lemma~\ref{lem:invariances-averaged-law}(c).
\end{proof}

\subsection{Moments on tensor tests}\label{sec:moments-tensor-tests}

Let $n\ge1$ be an integer. A point of $\mathbb R^{3n}=\left(\mathbb R^3\right)^n$ is written $x=(x_1,\dots,x_n)$ or $y=(y_1,\dots,y_n)$ with $x_i,y_i\in\mathbb R^3$; $|x|$ denotes the Euclidean norm of $x\in\mathbb R^{3n}$, and, for $\nu\in\{1,2,3\}$, $(y_i)_\nu$ denotes the $\nu$-th coordinate of $y_i$. The vectors $e_{c_1}\otimes\dots\otimes e_{c_n}$, $c_1,\dots,c_n\in\{12,13,23\}$, form a basis of $\left(\Lambda^2\right)^{\otimes n}$, and the $3^n$ coefficients of a vector with respect to this basis are its components; for $\mathbf c=(c_1,\dots,c_n)\in\{12,13,23\}^n$ the component at $e_{c_1}\otimes\dots\otimes e_{c_n}$ of a vector $v$ of $\left(\Lambda^2\right)^{\otimes n}$, or of its complexification, is written $v_{\mathbf c}$, and for a function $\Phi$ on $\mathbb R^{3n}$ with values in this complexification the function $x\mapsto\Phi(x)_{\mathbf c}$ is written $\Phi_{\mathbf c}$. Let
\[
\mathcal S_n=\mathcal S\left(\mathbb R^{3n};\left(\Lambda^2\right)^{\otimes n}\right)
\]
be the space of the functions on $\mathbb R^{3n}$ with values in the complexification of $\left(\Lambda^2\right)^{\otimes n}$ whose components are complex Schwartz functions. For $\Psi\in\mathcal S_n$ and $x\in\mathbb R^{3n}$ let $\left|\Psi(x)\right|$ be the largest of the moduli of the components of $\Psi(x)$; the same convention applies to the derivatives $\partial^\alpha\Psi$, $\alpha\in\mathbb N_0^{3n}$, whose order is $|\alpha|=\alpha_1+\dots+\alpha_{3n}$. For $m\in\mathbb N_0$ define $q_m\colon\mathcal S_n\to[0,\infty)$ by
\[
q_m(\Psi)=\max_{|\alpha|\le m}\ \sup_{x\in\mathbb R^{3n}}\left(1+|x|^2\right)^{\frac m2}\cdot\left|\partial^\alpha\Psi(x)\right|;
\]
the seminorms $q_m$, $m\in\mathbb N_0$, define the topology of $\mathcal S_n$.

Let $L^2\left(\mathbb R^{3n};\left(\Lambda^2\right)^{\otimes n}\right)$ be the complex Hilbert space of the measurable functions $\Phi$ on $\mathbb R^{3n}$ with values in the complexification of $\left(\Lambda^2\right)^{\otimes n}$, taken modulo equality almost everywhere, for which
\[
\|\Phi\|_{L^2}^2=\int_{\mathbb R^{3n}}\mathrm{d}x\,\sum_{\mathbf c\in\{12,13,23\}^n}\left|\Phi_{\mathbf c}(x)\right|^2
\]
is finite. For $\Phi,\Psi$ in this space the scalar product
\[
\left\langle\Phi,\Psi\right\rangle_{L^2}=\int_{\mathbb R^{3n}}\mathrm{d}x\,\sum_{\mathbf c\in\{12,13,23\}^n}\overline{\Phi_{\mathbf c}(x)}\cdot\Psi_{\mathbf c}(x)
\]
is antilinear in $\Phi$ and linear in $\Psi$, and $\mathcal S_n$ is a subspace of $L^2\left(\mathbb R^{3n};\left(\Lambda^2\right)^{\otimes n}\right)$.

For $f_1,\dots,f_n\in\mathcal S(\mathbb R^3;\Lambda^2)$, real or complex, let $f_1\otimes\dots\otimes f_n$ be the function $x\mapsto f_1(x_1)\otimes\dots\otimes f_n(x_n)$ on $\mathbb R^{3n}$; its component of index $\mathbf c=(c_1,\dots,c_n)$ is the function $x\mapsto\prod_{i=1}^n\left(f_i\right)_{c_i}(x_i)$. For integers $m,m'\ge1$, $\Psi\in\mathcal S_m$ and $\Phi\in\mathcal S_{m'}$ let $\Psi\otimes\Phi$ be the function on $\mathbb R^{3(m+m')}$ given by
\[
\left(\Psi\otimes\Phi\right)(y_1,\dots,y_{m+m'})=\Psi(y_1,\dots,y_m)\otimes\Phi(y_{m+1},\dots,y_{m+m'});
\]
its component of index $(\mathbf c,\mathbf c')\in\{12,13,23\}^m\times\{12,13,23\}^{m'}$ is the function $y\mapsto\Psi_{\mathbf c}(y_1,\dots,y_m)\cdot\Phi_{\mathbf c'}(y_{m+1},\dots,y_{m+m'})$. A product of Schwartz functions of disjoint groups of variables is a Schwartz function of all of them; hence $f_1\otimes\dots\otimes f_n\in\mathcal S_n$ and $\Psi\otimes\Phi\in\mathcal S_{m+m'}$.

Let $\mathcal K$, the Hermite functions $h_\kappa$ ($\kappa\in\mathcal K$), the operator $\mathsf A$ and its eigenvalues $\lambda_\kappa$ be as in Section~\ref{sec:hermite-sobolev-spaces}. For $\boldsymbol\kappa=(\kappa_1,\dots,\kappa_n)\in\mathcal K^n$ and $\Psi\in\mathcal S_n$ put
\[
h_{\boldsymbol\kappa}(x)=h_{\kappa_1}(x_1)\otimes\dots\otimes h_{\kappa_n}(x_n)\quad\left(x\in\mathbb R^{3n}\right),\qquad c_{\boldsymbol\kappa}(\Psi)=\left\langle h_{\boldsymbol\kappa},\Psi\right\rangle_{L^2},
\]
so that $h_{\boldsymbol\kappa}=h_{\kappa_1}\otimes\dots\otimes h_{\kappa_n}$. The functions $h_{\boldsymbol\kappa}$ are real; hence
\[
c_{\boldsymbol\kappa}(\Psi)=\int_{\mathbb R^{3n}}\mathrm{d}x\,\sum_{\mathbf c\in\{12,13,23\}^n}\left(h_{\boldsymbol\kappa}\right)_{\mathbf c}(x)\cdot\Psi_{\mathbf c}(x),
\]
and $c_{\boldsymbol\kappa}$ is a complex linear form on $\mathcal S_n$. The family $\left(h_{\boldsymbol\kappa}\right)_{\boldsymbol\kappa\in\mathcal K^n}$ is an orthonormal basis of $L^2\left(\mathbb R^{3n};\left(\Lambda^2\right)^{\otimes n}\right)$, since $\left(h_\kappa\right)_{\kappa\in\mathcal K}$ is an orthonormal basis of $L^2(\mathbb R^3;\Lambda^2)$.

The operators of Section~\ref{sec:test-forms-test} act on $\mathcal S_n$ in each variable: for $b\in\mathbb R^3$, $r\in O_{\mathrm{cub}}$ and $\Psi\in\mathcal S_n$ let
\[
\left(\tau_b^{\otimes n}\Psi\right)(y)=\Psi(y_1-b,\dots,y_n-b),\qquad\left(U_r^{\otimes n}\Psi\right)(y)=\left(\wedge^2r\right)^{\otimes n}\Psi\left(r^{-1}y_1,\dots,r^{-1}y_n\right),
\]
and let $\mathsf R_\nu^{\otimes n}=U_{r_\nu}^{\otimes n}$ for $\nu\in\{1,2,3\}$. For $t\in\mathbb N_0$ let $\left(\mathsf A^t\right)^{\otimes n}$ be the operator on $\mathcal S_n$ that applies $\mathsf A^t$, componentwise, in each of the variables $x_1,\dots,x_n$; the $n$ operators so obtained commute, and $\mathsf A^{\otimes n}$ stands for $\left(\mathsf A^1\right)^{\otimes n}$. These operators map $\mathcal S_n$ continuously into itself, and on products they act factor by factor; for instance $\tau_b^{\otimes n}(f_1\otimes\dots\otimes f_n)=\tau_bf_1\otimes\dots\otimes\tau_bf_n$ for $f_1,\dots,f_n\in\mathcal S(\mathbb R^3;\Lambda^2)$.

For a direction $\nu\in\{1,2,3\}$ let $\Theta_n\colon\mathcal S_n\to\mathcal S_n$ be the map
\[
\Theta_n\Psi=\overline{\mathsf R_\nu^{\otimes n}\Psi}.
\]
Finally put $\mathcal S_0=\mathbb C$ and $\Theta_0\Psi_0=\overline{\Psi_0}$ for $\Psi_0\in\mathcal S_0$, and let the tensor product of $\Psi_0\in\mathcal S_0$ with $\Psi\in\mathcal S_n$, $n\ge0$, in either order, be the product: $\Psi_0\otimes\Psi=\Psi\otimes\Psi_0=\Psi_0\cdot\Psi$.

\begin{lemma}[Hermite coefficients of tensor tests]\label{lem:hermite-coefficients-tensor}
For every integer $n\ge1$ and every $\Psi\in\mathcal S_n$,
\[
\sum_{\boldsymbol\kappa\in\mathcal K^n}\left|c_{\boldsymbol\kappa}(\Psi)\right|\le\mathsf Z_4^{\frac n2}\cdot\left\|\left(\mathsf A^2\right)^{\otimes n}\Psi\right\|_{L^2}\le4^{-n}\cdot\left(34\pi\sqrt3\right)^n\cdot q_{6n}(\Psi).
\]
\end{lemma}

\begin{proof}
The operator $\mathsf A$ is symmetric on $\mathcal S(\mathbb R^3;\Lambda^2)$ for the scalar product of $L^2(\mathbb R^3;\Lambda^2)$ (integration by parts), and $\mathsf A h_\kappa=\lambda_\kappa\cdot h_\kappa$ with real $\lambda_\kappa$. Hence, for $t\in\mathbb N_0$, $\boldsymbol\kappa\in\mathcal K^n$ and $\Psi\in\mathcal S_n$,
\[
c_{\boldsymbol\kappa}\left(\left(\mathsf A^t\right)^{\otimes n}\Psi\right)=\left\langle\left(\mathsf A^t\right)^{\otimes n}h_{\boldsymbol\kappa},\Psi\right\rangle_{L^2}=\prod_{i=1}^n\lambda_{\kappa_i}^t\cdot c_{\boldsymbol\kappa}(\Psi).
\]
The Cauchy--Schwarz inequality, applied with the weights $\prod_i\lambda_{\kappa_i}^{-2}$ and $\prod_i\lambda_{\kappa_i}^2$, gives
\[
\sum_{\boldsymbol\kappa\in\mathcal K^n}\left|c_{\boldsymbol\kappa}(\Psi)\right|\le\left(\sum_{\boldsymbol\kappa\in\mathcal K^n}\prod_{i=1}^n\lambda_{\kappa_i}^{-4}\right)^{\frac12}\cdot\left(\sum_{\boldsymbol\kappa\in\mathcal K^n}\prod_{i=1}^n\lambda_{\kappa_i}^4\cdot\left|c_{\boldsymbol\kappa}(\Psi)\right|^2\right)^{\frac12}.
\]
By \eqref{eq:continuum-infinite-4} the first factor equals $\left(\sum_{\kappa\in\mathcal K}\lambda_\kappa^{-4}\right)^{\frac n2}=\mathsf Z_4^{\frac n2}$. By the identity above with $t=2$ and by Parseval's identity for the orthonormal basis $\left(h_{\boldsymbol\kappa}\right)_{\boldsymbol\kappa\in\mathcal K^n}$, the second factor equals
\[
\left(\sum_{\boldsymbol\kappa\in\mathcal K^n}\left|c_{\boldsymbol\kappa}\left(\left(\mathsf A^2\right)^{\otimes n}\Psi\right)\right|^2\right)^{\frac12}=\left\|\left(\mathsf A^2\right)^{\otimes n}\Psi\right\|_{L^2}.
\]
This is the first inequality.

Let $\mathsf M$ be the function $z\mapsto1+|z|^2$ of the variable $z=(z_1,z_2,z_3)\in\mathbb R^3$; the same letter denotes the operator of multiplication by this function, acting componentwise, so that $\mathsf A=\mathsf M-\Delta$. Since $\nabla\mathsf M(z)=2z$ and $\Delta\mathsf M=6$, the Leibniz rule gives $\Delta(\mathsf Mu)=\mathsf M\Delta u+4z\cdot\nabla u+6u$ for $u\in\mathcal S(\mathbb R^3;\Lambda^2)$, and therefore
\[
\mathsf A^2u=\mathsf M^2u-\mathsf M\Delta u-\Delta(\mathsf Mu)+\Delta^2u=\mathsf M^2u-2\mathsf M\Delta u-4z\cdot\nabla u-6u+\Delta^2u.
\]
In coordinates, with $\partial_j$ the derivative with respect to $z_j$,
\[
\mathsf A^2u=\mathsf M^2u-2\sum_{j=1}^3\mathsf M\cdot\partial_j^2u-4\sum_{j=1}^3z_j\cdot\partial_ju-6u+\sum_{j,l=1}^3\partial_j^2\partial_l^2u.
\]
Thus $\mathsf A^2=\sum_{p=1}^{17}a_p\cdot\partial^{\beta_p}$ with multi-indices $\beta_p\in\mathbb N_0^3$ of order $|\beta_p|\le4$ and with functions $a_p$ on $\mathbb R^3$ such that $|a_p|\le C_p\cdot\mathsf M^2$. As $1\le\mathsf M$ and $|z_j|\le|z|\le\mathsf M(z)$, one may take $C_p=1$ for the term $\mathsf M^2$, $C_p=2$ for each of the three terms $-2\mathsf M\cdot\partial_j^2$, $C_p=4$ for each of the three terms $-4z_j\cdot\partial_j$, $C_p=6$ for the term $-6$, and $C_p=1$ for each of the nine terms $\partial_j^2\partial_l^2$; then
\[
\sum_{p=1}^{17}C_p=1+6+12+6+9=34.
\]
The operator $\left(\mathsf A^2\right)^{\otimes n}$ is the product of the operators $\mathsf A^2$ acting in the variables $x_1,\dots,x_n$. A function of $x_i$ commutes with the derivatives in the other variables; hence
\[
\left(\mathsf A^2\right)^{\otimes n}=\sum_{p_1,\dots,p_n=1}^{17}\prod_{i=1}^na_{p_i}(x_i)\cdot\partial^{(\beta_{p_1},\dots,\beta_{p_n})},
\]
where the multi-index $(\beta_{p_1},\dots,\beta_{p_n})\in\mathbb N_0^{3n}$ has order at most $4n$. Since $\sum_{p_1,\dots,p_n}\prod_iC_{p_i}=34^n$, every component of $\left(\mathsf A^2\right)^{\otimes n}\Psi$ satisfies
\[
\left|\left(\mathsf A^2\right)^{\otimes n}\Psi(x)\right|\le34^n\cdot\prod_{i=1}^n\left(1+|x_i|^2\right)^2\cdot\max_{|\alpha|\le4n}\left|\partial^\alpha\Psi(x)\right|\qquad\left(x\in\mathbb R^{3n}\right).
\]
On the other hand $1+|x_i|^2\le1+|x|^2$ for each $i$, so that $\prod_{i=1}^n\left(1+|x_i|^2\right)\le\left(1+|x|^2\right)^n$, and the definition of $q_{6n}$ gives
\[
\max_{|\alpha|\le6n}\left|\partial^\alpha\Psi(x)\right|\le q_{6n}(\Psi)\cdot\left(1+|x|^2\right)^{-3n}\le q_{6n}(\Psi)\cdot\prod_{i=1}^n\left(1+|x_i|^2\right)^{-3}.
\]
As $4n\le6n$, every component of $\left(\mathsf A^2\right)^{\otimes n}\Psi$ is bounded at $x$ by $34^n\cdot q_{6n}(\Psi)\cdot\prod_{i=1}^n\left(1+|x_i|^2\right)^{-1}$. In polar coordinates
\[
\int_{\mathbb R^3}\mathrm{d}z\,\left(1+|z|^2\right)^{-2}=4\pi\int_0^\infty\mathrm{d}\rho\,\rho^2\cdot\left(1+\rho^2\right)^{-2}=\pi^2,
\]
and there are $3^n$ components; therefore
\[
\left\|\left(\mathsf A^2\right)^{\otimes n}\Psi\right\|_{L^2}^2\le3^n\cdot34^{2n}\cdot\pi^{2n}\cdot q_{6n}(\Psi)^2.
\]
Finally $\mathsf Z_4<\frac1{16}$ by \eqref{eq:continuum-infinite-4}, so that $\mathsf Z_4^{\frac n2}<4^{-n}$, and the second inequality follows.
\end{proof}

\begin{theorem}[Moments of a limit point]\label{thm:moments-limit-point}
Let $\mu$ be a limit point of an admissible family (Definitions~\ref{def:admissible-families} and~\ref{def:limit-points}). Let $X\colon L^2(\mathbb R^3;\Lambda^2)\to L^2(\mu)$ be the bounded linear map of Theorem~\ref{thm:continuum-infinite-1}(1), extended to the complex space $L^2(\mathbb R^3;\Lambda^2)_{\mathbb C}=L^2(\mathbb R^3;\Lambda^2)+i\,L^2(\mathbb R^3;\Lambda^2)$, normed by $\|g+ih\|_{L^2}^2=\|g\|_{L^2}^2+\|h\|_{L^2}^2$, through $X(g+ih)=X(g)+iX(h)$. Let $\sigma_n=2\cdot2^{\frac n2}\cdot\Gamma\left(\frac n2+1\right)$ be the numbers of Proposition~\ref{prop:sub-gaussian-random}. For an integer $n\ge0$ and $f_1,\dots,f_n\in L^2(\mathbb R^3;\Lambda^2)_{\mathbb C}$ the Schwinger functions of $\mu$ are
\[
S_n(f_1,\dots,f_n)=\int\mathrm{d}\mu\,\prod_{i=1}^nX(f_i),
\]
the empty product being $1$.
\begin{enumerate}
\item[(1)] $S_0=1$ and $S_1=0$; each $S_n$ is multilinear and symmetric, and it is real when $f_1,\dots,f_n$ are real. Moreover
\[
\left|S_n(f_1,\dots,f_n)\right|\le2\cdot\left(2^n\cdot n!\right)^{\frac12}\cdot\prod_{i=1}^n\|f_i\|_{L^2}\qquad\text{for real }f_1,\dots,f_n,
\]
\[
\left|S_n(f_1,\dots,f_n)\right|\le2\cdot\left(4^n\cdot n!\right)^{\frac12}\cdot\prod_{i=1}^n\|f_i\|_{L^2}\qquad\text{for complex }f_1,\dots,f_n.
\]
\item[(2)] Let $n\ge1$ and $\Psi\in\mathcal S_n$. For every $F\in\mathcal S'$ the series
\[
Q_\Psi(F)=\sum_{\boldsymbol\kappa\in\mathcal K^n}c_{\boldsymbol\kappa}(\Psi)\cdot\prod_{i=1}^nF(h_{\kappa_i})
\]
converges absolutely, and $Q_\Psi(F)=F^{\otimes n}(\Psi)$; here $F^{\otimes n}$ is the $n$-th tensor power of $F$, that is, the continuous linear functional on $\mathcal S_n$ with $F^{\otimes n}(f_1\otimes\dots\otimes f_n)=\prod_{i=1}^nF(f_i)$ for $f_1,\dots,f_n\in\mathcal S(\mathbb R^3;\Lambda^2)$. The function $Q_\Psi\colon\mathcal S'\to\mathbb C$ satisfies
\[
\|Q_\Psi\|_{L^2(\mu)}\le\sigma_{2n}^{\frac12}\cdot\sum_{\boldsymbol\kappa\in\mathcal K^n}\left|c_{\boldsymbol\kappa}(\Psi)\right|.
\]
For $\Psi_0\in\mathcal S_0$ let $Q_{\Psi_0}$ be the constant function $\Psi_0$ on $\mathcal S'$, and for $n\ge0$ and $\Psi\in\mathcal S_n$ put $S_n(\Psi)=\int\mathrm{d}\mu\,Q_\Psi$. For $n\ge1$ the map $\Psi\mapsto S_n(\Psi)$ is the unique continuous linear functional on $\mathcal S_n$ such that $S_n(f_1\otimes\dots\otimes f_n)=S_n(f_1,\dots,f_n)$ for all $f_1,\dots,f_n\in\mathcal S(\mathbb R^3;\Lambda^2)$, and
\[
\left|S_n(\Psi)\right|\le\sigma_n\cdot\mathsf Z_4^{\frac n2}\cdot\left\|\left(\mathsf A^2\right)^{\otimes n}\Psi\right\|_{L^2}\le2\cdot66^n\cdot(n!)^{\frac12}\cdot q_{6n}(\Psi)\qquad\left(\Psi\in\mathcal S_n\right).
\]
Condition \textup{(E0$'$)} of \cite{OS2}, the linear growth condition, involves an integer $s$ and a sequence of constants, which \cite{OS2} writes as $\sigma_n$. The bound of $\left|S_n(\Psi)\right|$ by the last member of the display is this condition for $s=6$, the constants of \cite{OS2} being here the numbers $2\cdot66^n\cdot(n!)^{\frac12}$, not the numbers $\sigma_n$ of Proposition~\ref{prop:sub-gaussian-random}. It is valid for every $\Psi\in\mathcal S_n$, without any restriction on $\Psi$ at the points where two of the variables $y_1,\dots,y_n$ coincide.
\item[(3)] Let $\left(N_k,T_k\right)_{k\in\mathbb N}$ be a sequence of cutoffs as in Definition~\ref{def:limit-points}, with $N_k\to\infty$, $T_k\to\infty$ and $\mathsf P_{N_k,T_k}\to\mu$ weakly on $H_{-4}$, let $E_k$ be the expectation with respect to $\mathrm{Pr}_{N_k,T_k}$, and for $F\in C^2(\Lambda_{\eta_k,R_k})$ let $F^\sharp\in\mathcal S'$ be as in Section~\ref{sec:test-forms-test}. Let $n\ge1$ and $\Psi\in\mathcal S_n$. For every $k$ the function $F\mapsto\left(F^\sharp\right)^{\otimes n}(\Psi)$ is integrable with respect to $\mathrm{Pr}_{N_k,T_k}$, and
\[
\left|E_k\left[\left(F^\sharp\right)^{\otimes n}(\Psi)\right]\right|\le\sigma_n\cdot31^n\cdot\mathsf Z_4^{\frac n2}\cdot\left\|\left(\mathsf A^3\right)^{\otimes n}\Psi\right\|_{L^2};
\]
moreover
\[
\lim_{k\to\infty}E_k\left[\left(F^\sharp\right)^{\otimes n}(\Psi)\right]=S_n(\Psi).
\]
\item[(4)] For $n\ge1$, $\Psi\in\mathcal S_n$, $b\in\mathbb R^3$ and $r\in O_{\mathrm{cub}}$,
\[
S_n\left(\tau_b^{\otimes n}\Psi\right)=S_n(\Psi)=S_n\left(U_r^{\otimes n}\Psi\right),
\]
and $S_n\left(d^*g\otimes\Psi'\right)=0$ for all $g\in\mathcal S(\mathbb R^3)$ and $\Psi'\in\mathcal S_{n-1}$, where $d^*g$ is the two-form of Lemma~\ref{lem:discrete-codifferentials-three}.
\item[(5)] Let $\nu\in\{1,2,3\}$, and let the maps $\Theta_m$, $m\ge0$, be formed with this direction $\nu$. Let $\ell\ge0$ be an integer, let $\Psi_0\in\mathcal S_0$, and for $1\le n\le\ell$ let $\Psi_n\in\mathcal S_n$ be such that
\[
\operatorname{supp}\Psi_n\subset\left\{y\in\mathbb R^{3n}:(y_i)_\nu\ge0\text{ for }i=1,\dots,n\right\}.
\]
Then
\[
\sum_{m=0}^\ell\sum_{n=0}^\ell S_{m+n}\left(\Theta_m\Psi_m\otimes\Psi_n\right)\ge0.
\]
\end{enumerate}
\end{theorem}

\begin{proof}
By Theorem~\ref{thm:continuum-infinite-1}(1), every assertion of Proposition~\ref{prop:sub-gaussian-random} holds for the probability space $(\mathcal S',\mathcal C,\mu)$, the space $\mathcal T=L^2(\mathbb R^3;\Lambda^2)$ normed by $\|\cdot\|_{L^2}$, and the map $X$; moreover, for $f\in\mathcal S(\mathbb R^3;\Lambda^2)$ one has $X(f)=F(f)$ for $\mu$-almost every $F$ (for complex $f$ by linearity).
\begin{enumerate}
\item[(1)] Since $\mu$ is a probability measure, $S_0=1$. Let $n\ge1$. By Proposition~\ref{prop:sub-gaussian-random}(4) the product $\prod_iX(f_i)$ is $\mu$-integrable, and
\[
\int\mathrm{d}\mu\,\prod_{i=1}^n\left|X(f_i)\right|\le\sigma_n\cdot\prod_{i=1}^n\|f_i\|_{L^2}\qquad\text{for real }f_1,\dots,f_n,
\]
\[
\int\mathrm{d}\mu\,\prod_{i=1}^n\left|X(f_i)\right|\le2^{\frac n2}\cdot\sigma_n\cdot\prod_{i=1}^n\|f_i\|_{L^2}\qquad\text{for complex }f_1,\dots,f_n.
\]
As $\sigma_n\le2\cdot\left(2^n\cdot n!\right)^{\frac12}$ (Proposition~\ref{prop:sub-gaussian-random}) and $2^{\frac n2}\cdot2\cdot\left(2^n\cdot n!\right)^{\frac12}=2\cdot\left(4^n\cdot n!\right)^{\frac12}$, the two bounds of (1) follow; for $n=0$ they read $1\le2$. Multilinearity follows from the linearity of $X$, symmetry from the commutativity of the product, and reality from the fact that $X(f)$ is real-valued for real $f$. By Proposition~\ref{prop:sub-gaussian-random}(2), $\int\mathrm{d}\mu\,X(f)=0$ for real $f$, hence by linearity for complex $f$; so $S_1=0$.
\item[(2)] Let $n\ge1$, $\Psi\in\mathcal S_n$ and $F\in\mathcal S'$. By Section~\ref{sec:hermite-sobolev-spaces} there is $s\ge0$ with $F\in H_{-s}$, and the definition of $\|\cdot\|_{-s}$ gives $\left|F(h_\kappa)\right|\le\|F\|_{-s}\cdot\lambda_\kappa^s$ for $\kappa\in\mathcal K$. Let $t\ge s$ be an integer. By the identity for $c_{\boldsymbol\kappa}\left(\left(\mathsf A^t\right)^{\otimes n}\Psi\right)$ in the proof of Lemma~\ref{lem:hermite-coefficients-tensor}, by $\lambda_\kappa\ge1$, and by Lemma~\ref{lem:hermite-coefficients-tensor} applied to $\left(\mathsf A^t\right)^{\otimes n}\Psi\in\mathcal S_n$,
\[
\sum_{\boldsymbol\kappa\in\mathcal K^n}\left|c_{\boldsymbol\kappa}(\Psi)\right|\cdot\prod_{i=1}^n\lambda_{\kappa_i}^s\le\sum_{\boldsymbol\kappa\in\mathcal K^n}\left|c_{\boldsymbol\kappa}\left(\left(\mathsf A^t\right)^{\otimes n}\Psi\right)\right|\le\mathsf Z_4^{\frac n2}\cdot\left\|\left(\mathsf A^{t+2}\right)^{\otimes n}\Psi\right\|_{L^2}<\infty.
\]
Hence
\[
\sum_{\boldsymbol\kappa\in\mathcal K^n}\left|c_{\boldsymbol\kappa}(\Psi)\right|\cdot\prod_{i=1}^n\left|F(h_{\kappa_i})\right|\le\|F\|_{-s}^n\cdot\sum_{\boldsymbol\kappa\in\mathcal K^n}\left|c_{\boldsymbol\kappa}(\Psi)\right|\cdot\prod_{i=1}^n\lambda_{\kappa_i}^s<\infty,
\]
and the series $Q_\Psi(F)$ converges absolutely. Let $J_1\subset J_2\subset\dots$ be finite subsets of $\mathcal K^n$ whose union is $\mathcal K^n$. By Section~\ref{sec:hermite-sobolev-spaces}, applied to the Hermite functions on $\mathbb R^{3n}$, the partial sums $\sum_{\boldsymbol\kappa\in J_j}c_{\boldsymbol\kappa}(\Psi)\cdot h_{\boldsymbol\kappa}$ converge to $\Psi$ in $\mathcal S_n$ as $j\to\infty$. Since $F^{\otimes n}$ is continuous on $\mathcal S_n$ and $F^{\otimes n}(h_{\boldsymbol\kappa})=\prod_iF(h_{\kappa_i})$, it follows that $F^{\otimes n}(\Psi)=Q_\Psi(F)$.

The function $Q_\Psi$ is $\mathcal C$-measurable, as a pointwise limit of polynomials in the functions $F\mapsto F(h_\kappa)$. Let $\boldsymbol\kappa,\boldsymbol\kappa'\in\mathcal K^n$. The $2n$ functions $h_{\kappa_1},\dots,h_{\kappa_n},h_{\kappa'_1},\dots,h_{\kappa'_n}$ are real elements of $L^2(\mathbb R^3;\Lambda^2)$ of norm $1$, so that Proposition~\ref{prop:sub-gaussian-random}(4) gives
\[
\int\mathrm{d}\mu(F)\,\prod_{i=1}^n\left|F(h_{\kappa_i})\right|\cdot\prod_{i=1}^n\left|F(h_{\kappa'_i})\right|\le\sigma_{2n}.
\]
By Tonelli's theorem,
\begin{align*}
\int\mathrm{d}\mu\,\left|Q_\Psi\right|^2&\le\sum_{\boldsymbol\kappa,\boldsymbol\kappa'\in\mathcal K^n}\left|c_{\boldsymbol\kappa}(\Psi)\right|\cdot\left|c_{\boldsymbol\kappa'}(\Psi)\right|\cdot\int\mathrm{d}\mu(F)\,\prod_{i=1}^n\left|F(h_{\kappa_i})\right|\cdot\prod_{i=1}^n\left|F(h_{\kappa'_i})\right|\\
&\le\sigma_{2n}\cdot\left(\sum_{\boldsymbol\kappa\in\mathcal K^n}\left|c_{\boldsymbol\kappa}(\Psi)\right|\right)^2,
\end{align*}
which is the bound for $\|Q_\Psi\|_{L^2(\mu)}$; its right side is finite by Lemma~\ref{lem:hermite-coefficients-tensor}. In the same way, Proposition~\ref{prop:sub-gaussian-random}(4) with the $n$ functions $h_{\kappa_1},\dots,h_{\kappa_n}$ gives
\[
\int\mathrm{d}\mu\,\left|Q_\Psi\right|\le\sigma_n\cdot\sum_{\boldsymbol\kappa\in\mathcal K^n}\left|c_{\boldsymbol\kappa}(\Psi)\right|.
\]
The map $\Psi\mapsto S_n(\Psi)$ is linear, since $\Psi\mapsto F^{\otimes n}(\Psi)$ is linear for every $F$. By the last bound, by Lemma~\ref{lem:hermite-coefficients-tensor} and by $\sigma_n\le2\cdot\left(2^n\cdot n!\right)^{\frac12}$,
\begin{align*}
\left|S_n(\Psi)\right|&\le\sigma_n\cdot\sum_{\boldsymbol\kappa\in\mathcal K^n}\left|c_{\boldsymbol\kappa}(\Psi)\right|\le\sigma_n\cdot\mathsf Z_4^{\frac n2}\cdot\left\|\left(\mathsf A^2\right)^{\otimes n}\Psi\right\|_{L^2}\\
&\le2\cdot\left(2^n\cdot n!\right)^{\frac12}\cdot4^{-n}\cdot\left(34\pi\sqrt3\right)^n\cdot q_{6n}(\Psi)=2\cdot\left(\frac{\sqrt2\cdot34\pi\sqrt3}{4}\right)^n\cdot(n!)^{\frac12}\cdot q_{6n}(\Psi)\\
&\le2\cdot66^n\cdot(n!)^{\frac12}\cdot q_{6n}(\Psi),
\end{align*}
the last step because $\frac{\sqrt2\cdot34\pi\sqrt3}{4}=8.5\cdot\pi\sqrt6$, $\pi\sqrt6<7.7$ and $8.5\cdot7.7<66$. As $q_{6n}$ is a continuous seminorm on $\mathcal S_n$, the functional $S_n$ is continuous. For $f_1,\dots,f_n\in\mathcal S(\mathbb R^3;\Lambda^2)$ one has $Q_{f_1\otimes\dots\otimes f_n}(F)=F^{\otimes n}(f_1\otimes\dots\otimes f_n)=\prod_iF(f_i)$, and $F(f_i)=X(f_i)$ for $\mu$-almost every $F$; integration gives $S_n(f_1\otimes\dots\otimes f_n)=S_n(f_1,\dots,f_n)$. Two continuous linear functionals on $\mathcal S_n$ with this property coincide on every $h_{\boldsymbol\kappa}$, hence on the finite Hermite sums, which are dense in $\mathcal S_n$ (Section~\ref{sec:hermite-sobolev-spaces}); this proves the uniqueness.
\item[(3)] Let $n\ge1$ and $\Psi\in\mathcal S_n$. The family is admissible, so that $\mathrm{Pr}_{N_k,T_k}$ has property~\eqref{eq:continuum-infinite-Ad2}, and by Theorem~\ref{thm:estimates-fixed-cutoff} the assertions of Proposition~\ref{prop:sub-gaussian-random} hold under $\mathrm{Pr}_{N_k,T_k}$ for the real tests $f$, the seminorm $|f_{\eta_k}|_{\eta_k}$ and the map $f\mapsto X_k(f)$. By Lemma~\ref{lem:norm-bounds-lattice}, $\left|(h_\kappa)_{\eta_k}\right|_{\eta_k}\le p_{R_k}(h_\kappa)\le p_1(h_\kappa)\le31\lambda_\kappa$ for $\kappa\in\mathcal K$. Hence Proposition~\ref{prop:sub-gaussian-random}(4) yields
\[
E_k\prod_{i=1}^n\left|X_k(h_{\kappa_i})\right|\le\sigma_n\cdot\prod_{i=1}^n31\lambda_{\kappa_i}\qquad\left(\boldsymbol\kappa\in\mathcal K^n\right).
\]
By the identity for $c_{\boldsymbol\kappa}\left(\mathsf A^{\otimes n}\Psi\right)$ in the proof of Lemma~\ref{lem:hermite-coefficients-tensor} and by Lemma~\ref{lem:hermite-coefficients-tensor} applied to $\mathsf A^{\otimes n}\Psi$,
\[
\sum_{\boldsymbol\kappa\in\mathcal K^n}\left|c_{\boldsymbol\kappa}(\Psi)\right|\cdot\prod_{i=1}^n\lambda_{\kappa_i}=\sum_{\boldsymbol\kappa\in\mathcal K^n}\left|c_{\boldsymbol\kappa}\left(\mathsf A^{\otimes n}\Psi\right)\right|\le\mathsf Z_4^{\frac n2}\cdot\left\|\left(\mathsf A^3\right)^{\otimes n}\Psi\right\|_{L^2}.
\]
For every $F\in C^2(\Lambda_{\eta_k,R_k})$ one has $F^\sharp\in\mathcal S'$, so that part (2) and \eqref{eq:continuum-infinite-3} give
\[
\left(F^\sharp\right)^{\otimes n}(\Psi)=Q_\Psi\left(F^\sharp\right)=\sum_{\boldsymbol\kappa\in\mathcal K^n}c_{\boldsymbol\kappa}(\Psi)\cdot\prod_{i=1}^nX_k(h_{\kappa_i})(F).
\]
The two estimates give
\[
E_k\sum_{\boldsymbol\kappa\in\mathcal K^n}\left|c_{\boldsymbol\kappa}(\Psi)\right|\cdot\prod_{i=1}^n\left|X_k(h_{\kappa_i})\right|\le\sigma_n\cdot31^n\cdot\mathsf Z_4^{\frac n2}\cdot\left\|\left(\mathsf A^3\right)^{\otimes n}\Psi\right\|_{L^2}<\infty.
\]
Therefore $F\mapsto\left(F^\sharp\right)^{\otimes n}(\Psi)$ is integrable, Fubini's theorem allows the expectation to be taken term by term,
\[
E_k\left[\left(F^\sharp\right)^{\otimes n}(\Psi)\right]=\sum_{\boldsymbol\kappa\in\mathcal K^n}c_{\boldsymbol\kappa}(\Psi)\cdot E_k\prod_{i=1}^nX_k(h_{\kappa_i}),
\]
and the bound of (3) follows. For each $\boldsymbol\kappa$, Theorem~\ref{thm:continuum-infinite-1}(2) gives $E_k\prod_iX_k(h_{\kappa_i})\to\int\mathrm{d}\mu(F)\,\prod_iF(h_{\kappa_i})$ as $k\to\infty$. The terms of the last series are bounded in modulus by $\left|c_{\boldsymbol\kappa}(\Psi)\right|\cdot\sigma_n\cdot\prod_i31\lambda_{\kappa_i}$, which does not depend on $k$ and is summable over $\boldsymbol\kappa$. By dominated convergence for series,
\[
\lim_{k\to\infty}E_k\left[\left(F^\sharp\right)^{\otimes n}(\Psi)\right]=\sum_{\boldsymbol\kappa\in\mathcal K^n}c_{\boldsymbol\kappa}(\Psi)\cdot\int\mathrm{d}\mu(F)\,\prod_{i=1}^nF(h_{\kappa_i})=\int\mathrm{d}\mu\,Q_\Psi=S_n(\Psi);
\]
the middle equality is Fubini's theorem, which applies because $\int\mathrm{d}\mu(F)\,\sum_{\boldsymbol\kappa}\left|c_{\boldsymbol\kappa}(\Psi)\right|\cdot\prod_i\left|F(h_{\kappa_i})\right|\le\sigma_n\cdot\sum_{\boldsymbol\kappa}\left|c_{\boldsymbol\kappa}(\Psi)\right|<\infty$ by the proof of (2).
\item[(4)] Let $n\ge1$, $\Psi\in\mathcal S_n$, $b\in\mathbb R^3$ and $F\in\mathcal S'$. The maps $\Psi\mapsto F^{\otimes n}\left(\tau_b^{\otimes n}\Psi\right)$ and $\Psi\mapsto\left(\tau_b^*F\right)^{\otimes n}(\Psi)$ are continuous linear functionals on $\mathcal S_n$ which agree on the products $f_1\otimes\dots\otimes f_n$, both being equal to $\prod_iF(\tau_bf_i)$ there; the products span a dense subspace of $\mathcal S_n$, since they include the functions $h_{\boldsymbol\kappa}$, whose finite linear combinations are dense in $\mathcal S_n$ (Section~\ref{sec:hermite-sobolev-spaces}). By (2) it follows that $Q_{\tau_b^{\otimes n}\Psi}(F)=Q_\Psi\left(\tau_b^*F\right)$. As $\mu\circ\left(\tau_b^*\right)^{-1}=\mu$ by Theorem~\ref{thm:continuum-infinite-1}(5),
\[
S_n\left(\tau_b^{\otimes n}\Psi\right)=\int\mathrm{d}\mu(F)\,Q_\Psi\left(\tau_b^*F\right)=\int\mathrm{d}\mu\,Q_\Psi=S_n(\Psi).
\]
The same argument, with $U_r$, $U_r^*$ in the place of $\tau_b$, $\tau_b^*$ and with Theorem~\ref{thm:continuum-infinite-1}(6), gives $S_n\left(U_r^{\otimes n}\Psi\right)=S_n(\Psi)$ for $r\in O_{\mathrm{cub}}$.

For integers $m,m'\ge0$, $\Psi\in\mathcal S_m$ and $\Phi\in\mathcal S_{m'}$ one has the product rule
\[
Q_\Psi\cdot Q_\Phi=Q_{\Psi\otimes\Phi}\qquad\text{on }\mathcal S'.
\]
For $m=0$ or $m'=0$ this is the convention for $\mathcal S_0$. For $m,m'\ge1$ the functions $h_{\boldsymbol\kappa}\otimes h_{\boldsymbol\kappa'}$, $(\boldsymbol\kappa,\boldsymbol\kappa')\in\mathcal K^m\times\mathcal K^{m'}=\mathcal K^{m+m'}$, are the functions $h_{(\boldsymbol\kappa,\boldsymbol\kappa')}$ of $\mathcal S_{m+m'}$, and by Fubini's theorem the coefficient of $\Psi\otimes\Phi$ at $(\boldsymbol\kappa,\boldsymbol\kappa')$ is $c_{\boldsymbol\kappa}(\Psi)\cdot c_{\boldsymbol\kappa'}(\Phi)$; so the product of the two absolutely convergent series of (2) is the series of $Q_{\Psi\otimes\Phi}$. In particular, for $g\in\mathcal S(\mathbb R^3)$ and $\Psi'\in\mathcal S_{n-1}$,
\[
Q_{d^*g\otimes\Psi'}(F)=Q_{d^*g}(F)\cdot Q_{\Psi'}(F)=F\left(d^*g\right)\cdot Q_{\Psi'}(F)\qquad(F\in\mathcal S'),
\]
because $F^{\otimes1}=F$. By Theorem~\ref{thm:continuum-infinite-1}(4), $F(d^*g)=0$ for $\mu$-almost every $F$ (for complex $g$ by linearity). Hence $Q_{d^*g\otimes\Psi'}=0$ $\mu$-almost everywhere, and $S_n\left(d^*g\otimes\Psi'\right)=0$.
\item[(5)] Throughout, $\Sigma^+_{\nu,0}$, $L^2\left(\mathcal S',\Sigma^+_{\nu,0},\mu\right)$ and $\Theta_{\nu,0}$ are as in Theorem~\ref{thm:continuum-infinite-1}(7); note that $\mathsf R_{\nu,0}=\mathsf R_\nu$.

\emph{Step 1: if $n\ge1$ and $\Psi\in\mathcal S_n$ satisfies $\operatorname{supp}\Psi\subset\left\{y\in\mathbb R^{3n}:(y_i)_\nu\ge0\text{ for }i=1,\dots,n\right\}$, then $Q_\Psi$ coincides $\mu$-almost everywhere with a $\Sigma^+_{\nu,0}$-measurable function.} For $\varepsilon>0$ let $\Psi_\varepsilon=\tau_{\varepsilon\cdot e_\nu}^{\otimes n}\Psi$. Then
\[
\operatorname{supp}\Psi_\varepsilon\subset\left\{y\in\mathbb R^{3n}:(y_i)_\nu\ge\varepsilon\text{ for }i=1,\dots,n\right\},
\]
and $\Psi_\varepsilon\to\Psi$ in $\mathcal S_n$ as $\varepsilon\downarrow0$. Let $\varkappa_\varepsilon\in C^\infty(\mathbb R)$ be a real function with $\varkappa_\varepsilon=0$ on $\left(-\infty,\frac{\varepsilon}{2}\right]$ and $\varkappa_\varepsilon=1$ on $[\varepsilon,\infty)$; put $k_\varepsilon(y)=\varkappa_\varepsilon(y_\nu)$ for $y\in\mathbb R^3$ and $k_\varepsilon^{\otimes n}(y)=\prod_{i=1}^nk_\varepsilon(y_i)$ for $y\in\mathbb R^{3n}$. Then $k_\varepsilon^{\otimes n}=1$ on the support of $\Psi_\varepsilon$, so that $k_\varepsilon^{\otimes n}\cdot\Psi_\varepsilon=\Psi_\varepsilon$. Let $J_1\subset J_2\subset\dots$ be finite subsets of $\mathcal K^n$ with union $\mathcal K^n$, and put $\Psi_{\varepsilon,j}=\sum_{\boldsymbol\kappa\in J_j}c_{\boldsymbol\kappa}(\Psi_\varepsilon)\cdot h_{\boldsymbol\kappa}$; then $\Psi_{\varepsilon,j}\to\Psi_\varepsilon$ in $\mathcal S_n$ as $j\to\infty$ (Section~\ref{sec:hermite-sobolev-spaces}). The function $k_\varepsilon^{\otimes n}$ is smooth and bounded together with all its derivatives, so that multiplication by it is continuous on $\mathcal S_n$, and $k_\varepsilon^{\otimes n}\cdot\Psi_{\varepsilon,j}\to k_\varepsilon^{\otimes n}\cdot\Psi_\varepsilon=\Psi_\varepsilon$ in $\mathcal S_n$ as $j\to\infty$. Since $k_\varepsilon^{\otimes n}\cdot h_{\boldsymbol\kappa}=\left(k_\varepsilon\cdot h_{\kappa_1}\right)\otimes\dots\otimes\left(k_\varepsilon\cdot h_{\kappa_n}\right)$, part (2) gives
\[
Q_{k_\varepsilon^{\otimes n}\cdot\Psi_{\varepsilon,j}}(F)=\sum_{\boldsymbol\kappa\in J_j}c_{\boldsymbol\kappa}(\Psi_\varepsilon)\cdot\prod_{i=1}^nF\left(k_\varepsilon\cdot h_{\kappa_i}\right)\qquad(F\in\mathcal S'),
\]
a polynomial in the functions $F\mapsto F\left(k_\varepsilon\cdot h_\kappa\right)$, $\kappa\in\mathcal K$. Each $k_\varepsilon\cdot h_\kappa$ is a real Schwartz test supported in $\left\{y\in\mathbb R^3:y_\nu\ge\frac{\varepsilon}{2}\right\}$; by Theorem~\ref{thm:continuum-infinite-1}(7) the function $F\mapsto F\left(k_\varepsilon\cdot h_\kappa\right)$ coincides $\mu$-almost everywhere with a $\Sigma^+_{\nu,0}$-measurable function, and therefore so does $Q_{k_\varepsilon^{\otimes n}\cdot\Psi_{\varepsilon,j}}$. By (2) and Lemma~\ref{lem:hermite-coefficients-tensor},
\[
\|Q_\Phi\|_{L^2(\mu)}\le\sigma_{2n}^{\frac12}\cdot4^{-n}\cdot\left(34\pi\sqrt3\right)^n\cdot q_{6n}(\Phi)\qquad\left(\Phi\in\mathcal S_n\right),
\]
so that $\Phi\mapsto Q_\Phi$ is a continuous linear map from $\mathcal S_n$ to $L^2(\mu)$. Consequently $Q_{\Psi_\varepsilon}=\lim_{j\to\infty}Q_{k_\varepsilon^{\otimes n}\cdot\Psi_{\varepsilon,j}}$ and $Q_\Psi=\lim_{\varepsilon\downarrow0}Q_{\Psi_\varepsilon}$ in $L^2(\mu)$. The elements of $L^2(\mu)$ that have a $\Sigma^+_{\nu,0}$-measurable representative form a closed subspace of $L^2(\mu)$, namely the image of the complete space $L^2\left(\mathcal S',\Sigma^+_{\nu,0},\mu\right)$ under an isometry; hence $Q_\Psi$ has such a representative.

\emph{Step 2.} Let $m\ge1$ and $\Psi\in\mathcal S_m$. The argument given in (4) for $\tau_b$, applied to $U_{r_\nu}=\mathsf R_\nu$, gives $Q_\Psi\left(\mathsf R_\nu^*F\right)=Q_{\mathsf R_\nu^{\otimes m}\Psi}(F)$ for $F\in\mathcal S'$. Since the functions $h_{\boldsymbol\kappa}$ are real and $F(f)$ is real for every real test $f$, one has $c_{\boldsymbol\kappa}\left(\overline\Phi\right)=\overline{c_{\boldsymbol\kappa}(\Phi)}$ and $Q_{\overline\Phi}(F)=\overline{Q_\Phi(F)}$ for $\Phi\in\mathcal S_m$. Hence
\[
\left(\Theta_{\nu,0}Q_\Psi\right)(F)=\overline{Q_\Psi\left(\mathsf R_\nu^*F\right)}=\overline{Q_{\mathsf R_\nu^{\otimes m}\Psi}(F)}=Q_{\Theta_m\Psi}(F)\qquad(F\in\mathcal S').
\]
For $m=0$ and $\Psi_0\in\mathcal S_0$ the identity $\Theta_{\nu,0}Q_{\Psi_0}=Q_{\Theta_0\Psi_0}$ holds as well, both sides being the constant $\overline{\Psi_0}$. Put $P=\sum_{n=0}^\ell Q_{\Psi_n}$. By (2), by Step 1, and since $Q_{\Psi_0}$ is constant, $P$ is an element of $L^2(\mu)$ with a $\Sigma^+_{\nu,0}$-measurable representative, that is, an element of $L^2\left(\mathcal S',\Sigma^+_{\nu,0},\mu\right)$. Theorem~\ref{thm:continuum-infinite-1}(7) with $c=0$, the additivity and antilinearity of $\Theta_{\nu,0}$, and the product rule of (4) give
\begin{align*}
0\le\int\mathrm{d}\mu\,\left(\Theta_{\nu,0}P\right)\cdot P&=\sum_{m=0}^\ell\sum_{n=0}^\ell\int\mathrm{d}\mu\,Q_{\Theta_m\Psi_m}\cdot Q_{\Psi_n}=\sum_{m=0}^\ell\sum_{n=0}^\ell\int\mathrm{d}\mu\,Q_{\Theta_m\Psi_m\otimes\Psi_n}\\
&=\sum_{m=0}^\ell\sum_{n=0}^\ell S_{m+n}\left(\Theta_m\Psi_m\otimes\Psi_n\right).
\end{align*}
\end{enumerate}
\end{proof}

\begin{corollary}[Transfer operators in a coordinate direction]\label{cor:transfer-operators-coordinate}
Let $\mu$ be a limit point of an admissible family (Definition~\ref{def:limit-points}), let $\nu\in\{1,2,3\}$, and let $L^2(\mu)=L^2(\mathcal S',\mathcal C,\mu)$, a complex Hilbert space with norm $\|\cdot\|_{L^2(\mu)}$. Let $\Sigma^+_{\nu,0}$ and $\Theta_{\nu,0}$ be as in Theorem~\ref{thm:continuum-infinite-1}(7) with $c=0$, so that $(\Theta_{\nu,0}A)(F)=\overline{A\left(\mathsf R_\nu^*F\right)}$, and let $(\mathcal U_aA)(F)=A\left(\tau_a^*F\right)$, $a\in\mathbb R^3$, be the unitary representation of Theorem~\ref{thm:continuum-infinite-1}(5). Let $\mathcal E_+\subset L^2(\mu)$ be the set of classes which have a $\Sigma^+_{\nu,0}$-measurable representative; it is a closed subspace, since it is the isometric image of the complete space $L^2(\mathcal S',\Sigma^+_{\nu,0},\mu)$. For $A,B\in L^2(\mu)$ put
\[
\mathsf b(A,B)=\int\mathrm{d}\mu\,\left(\Theta_{\nu,0}A\right)\cdot B,
\]
and let $\mathcal N=\{A\in\mathcal E_+:\mathsf b(A,A)=0\}$. By Theorem~\ref{thm:continuum-infinite-1}(7), $\mathsf b(A,A)\ge0$ for $A\in\mathcal E_+$; consequently $\mathcal N$ is a linear subspace of $\mathcal E_+$, and, if $[A]$ denotes the class of $A\in\mathcal E_+$ in $\mathcal E_+/\mathcal N$, the formula $\left\langle[A],[B]\right\rangle_{\mathcal H_\nu}=\mathsf b(A,B)$ defines an inner product on $\mathcal E_+/\mathcal N$. Let $\mathcal H_\nu$ be the completion of $\mathcal E_+/\mathcal N$ with respect to this inner product, $\|\cdot\|_{\mathcal H_\nu}$ the norm of $\mathcal H_\nu$, and $\Omega=[1]$.
\begin{enumerate}
\item[(1)] For $t\ge0$ the assignment $[A]\mapsto\left[\mathcal U_{t\cdot e_\nu}A\right]$, $A\in\mathcal E_+$, is well defined and extends to a bounded operator $P_t$ on $\mathcal H_\nu$. The family $(P_t)_{t\ge0}$ is a strongly continuous semigroup of self-adjoint contractions, and there is a self-adjoint operator $\mathsf H_\nu\ge0$ on $\mathcal H_\nu$ with $P_t=\exp\left(-t\cdot\mathsf H_\nu\right)$ for every $t\ge0$. Moreover $\|\Omega\|_{\mathcal H_\nu}=1$ and $\mathsf H_\nu\Omega=0$.
\item[(2)] Let $\lambda\in\{1,2,3\}$, $\lambda\ne\nu$. For $s\in\mathbb R$ the assignment $[A]\mapsto\left[\mathcal U_{s\cdot e_\lambda}A\right]$, $A\in\mathcal E_+$, is well defined and extends to a unitary operator $\hat{\mathcal U}_{s\cdot e_\lambda}$ on $\mathcal H_\nu$. The family $\left(\hat{\mathcal U}_{s\cdot e_\lambda}\right)_{s\in\mathbb R}$ is a strongly continuous unitary group; $\hat{\mathcal U}_{s\cdot e_\lambda}P_t=P_t\hat{\mathcal U}_{s\cdot e_\lambda}$ for $s\in\mathbb R$ and $t\ge0$; and $\hat{\mathcal U}_{s\cdot e_\lambda}\Omega=\Omega$.
\item[(3)] For $k\ge1$ let $\mathbb R^{3k}_{\nu,+}$ be the set of the points $(y_1,\dots,y_k)\in\left(\mathbb R^3\right)^k$ with $(y_i)_\nu\ge0$ for $1\le i\le k$. Let $m,n\ge0$ be integers, let $\Psi\in\mathcal S_m$ and $\Phi\in\mathcal S_n$ (Section~\ref{sec:moments-tensor-tests}) satisfy $\mathrm{supp}\,\Psi\subset\mathbb R^{3m}_{\nu,+}$ if $m\ge1$ and $\mathrm{supp}\,\Phi\subset\mathbb R^{3n}_{\nu,+}$ if $n\ge1$, and let $\Theta_m\colon\mathcal S_m\to\mathcal S_m$ be the map of that section formed with the direction $\nu$. Then the functions $Q_\Psi$, $Q_\Phi$ of Theorem~\ref{thm:moments-limit-point}(2) belong to $\mathcal E_+$, and for every $t\ge0$
\[
S_{m+n}\left(\Theta_m\Psi\otimes\tau_{t\cdot e_\nu}^{\otimes n}\Phi\right)=\left\langle[Q_\Psi],\exp\left(-t\cdot\mathsf H_\nu\right)[Q_\Phi]\right\rangle_{\mathcal H_\nu},
\]
where $\tau_{t\cdot e_\nu}^{\otimes0}$ is the identity of $\mathcal S_0=\mathbb C$. As $t\to\infty$, this number converges to $\left\langle[Q_\Psi],E_0[Q_\Phi]\right\rangle_{\mathcal H_\nu}$, where $E_0$ is the orthogonal projection of $\mathcal H_\nu$ onto $\ker\mathsf H_\nu$.
\end{enumerate}
\end{corollary}

\begin{proof}
Write $\Theta=\Theta_{\nu,0}$, and let $\mathcal D^+_{\nu,0}$ be as in Theorem~\ref{thm:continuum-infinite-1}(7); thus $\Sigma^+_{\nu,0}$ is the $\sigma$-algebra generated by the functions $F\mapsto F(f)$, $f\in\mathcal D^+_{\nu,0}$.

\emph{The form $\mathsf b$.} One has $\mathsf R_\nu=U_{r_\nu}$ with $r_\nu\in O_{\mathrm{cub}}$, so $\mathsf R_\nu^*$ preserves $\mu$ by Theorem~\ref{thm:continuum-infinite-1}(6). Hence $\Theta$ is a well-defined antilinear isometry of $L^2(\mu)$, the integral defining $\mathsf b(A,B)$ converges absolutely, and the Cauchy--Schwarz inequality in $L^2(\mu)$ gives
\[
\left|\mathsf b(A,B)\right|\le\|\Theta A\|_{L^2(\mu)}\cdot\|B\|_{L^2(\mu)}=\|A\|_{L^2(\mu)}\cdot\|B\|_{L^2(\mu)}\qquad\left(A,B\in L^2(\mu)\right).
\]
The form $\mathsf b$ is antilinear in its first and linear in its second argument. Its diagonal values $\mathsf b(A,A)$, $A\in\mathcal E_+$, are real and non-negative by Theorem~\ref{thm:continuum-infinite-1}(7). By polarisation $\mathsf b(B,A)=\overline{\mathsf b(A,B)}$ for $A,B\in\mathcal E_+$, and the Cauchy--Schwarz inequality $\left|\mathsf b(A,B)\right|^2\le\mathsf b(A,A)\cdot\mathsf b(B,B)$ holds for $A,B\in\mathcal E_+$. Therefore $\mathcal N$ consists of the $A\in\mathcal E_+$ with $\mathsf b(A,B)=0$ for all $B\in\mathcal E_+$, so it is a linear subspace; $\mathsf b$ induces an inner product on $\mathcal E_+/\mathcal N$; and
\[
\left\|[A]\right\|_{\mathcal H_\nu}^2=\mathsf b(A,A)\le\|A\|_{L^2(\mu)}^2\qquad\left(A\in\mathcal E_+\right).
\]
As $\Theta1=1$, one has $\|\Omega\|_{\mathcal H_\nu}^2=\mathsf b(1,1)=\mu(\mathcal S')=1$.

\emph{The reflection and the translations.} Let $a\in\mathbb R^3$. Since $r_\nu^{-1}=r_\nu$, for every test $f$ and every $y\in\mathbb R^3$
\[
(\mathsf R_\nu\tau_af)(y)=\left(\wedge^2r_\nu\right)f(r_\nu y-a)=\left(\tau_{r_\nu a}\mathsf R_\nu f\right)(y),
\]
that is, $\mathsf R_\nu\tau_a=\tau_{r_\nu a}\mathsf R_\nu$. In particular $\mathsf R_\nu\tau_{t\cdot e_\nu}=\tau_{-t\cdot e_\nu}\mathsf R_\nu$ for $t\in\mathbb R$, and $\mathsf R_\nu\tau_{s\cdot e_\lambda}=\tau_{s\cdot e_\lambda}\mathsf R_\nu$ for $s\in\mathbb R$ and $\lambda\ne\nu$. For $F\in\mathcal S'$ it follows that
\[
\left(\tau_a^*\mathsf R_\nu^*F\right)(f)=F(\mathsf R_\nu\tau_af)=F\left(\tau_{r_\nu a}\mathsf R_\nu f\right)=\left(\mathsf R_\nu^*\tau_{r_\nu a}^*F\right)(f),
\]
and therefore
\[
\left(\Theta\,\mathcal U_aA\right)(F)=\overline{A\left(\tau_a^*\mathsf R_\nu^*F\right)}=\overline{A\left(\mathsf R_\nu^*\tau_{r_\nu a}^*F\right)}=\left(\mathcal U_{r_\nu a}\Theta A\right)(F).
\]
Thus $\Theta\,\mathcal U_a=\mathcal U_{r_\nu a}\Theta$; in particular $\Theta\,\mathcal U_{t\cdot e_\nu}=\mathcal U_{-t\cdot e_\nu}\Theta$ for $t\in\mathbb R$.

\emph{Translations that preserve $\mathcal E_+$.} Let $a\in\mathbb R^3$ with $a_\nu\ge0$. For $f\in\mathcal D^+_{\nu,0}$ the test $\tau_af$ is real and smooth, and its support $\mathrm{supp}\,f+a$ is a compact subset of $\{y_\nu>a_\nu\}\subset\{y_\nu>0\}$; hence $\tau_a$ maps $\mathcal D^+_{\nu,0}$ into itself. So, for each $f\in\mathcal D^+_{\nu,0}$, the function $F\mapsto\left(\tau_a^*F\right)(f)=F(\tau_af)$ is $\Sigma^+_{\nu,0}$-measurable, and so is $\mathcal U_a\exp\left(iF(f)\right)=\exp\left(iF(\tau_af)\right)$. Since $\Sigma^+_{\nu,0}$ is generated by the functions $F\mapsto F(f)$, $f\in\mathcal D^+_{\nu,0}$, the map $\tau_a^*$ is measurable from $(\mathcal S',\Sigma^+_{\nu,0})$ to $(\mathcal S',\Sigma^+_{\nu,0})$, and $A\circ\tau_a^*$ is $\Sigma^+_{\nu,0}$-measurable whenever $A$ is. As $A\circ\tau_a^*$ is a representative of $\mathcal U_aA$, one obtains
\[
\mathcal U_a\mathcal E_+\subset\mathcal E_+\qquad\left(a\in\mathbb R^3,\ a_\nu\ge0\right);
\]
in particular $\mathcal U_{t\cdot e_\nu}\mathcal E_+\subset\mathcal E_+$ for $t\ge0$.

\emph{Part (1).} Let $A,B\in L^2(\mu)$ and $t\in\mathbb R$. The product $\left(\Theta A\right)\cdot\left(\mathcal U_{t\cdot e_\nu}B\right)$ is integrable, $\mathcal U_{-t\cdot e_\nu}$ is multiplicative, and $\mu$ is invariant under $\tau^*_{-t\cdot e_\nu}$ (Theorem~\ref{thm:continuum-infinite-1}(5)). Together with $\Theta\,\mathcal U_{t\cdot e_\nu}=\mathcal U_{-t\cdot e_\nu}\Theta$ this gives
\begin{align*}
\mathsf b\left(A,\mathcal U_{t\cdot e_\nu}B\right)&=\int\mathrm{d}\mu\,\mathcal U_{-t\cdot e_\nu}\left[\left(\Theta A\right)\cdot\left(\mathcal U_{t\cdot e_\nu}B\right)\right]=\int\mathrm{d}\mu\,\left(\mathcal U_{-t\cdot e_\nu}\Theta A\right)\cdot B\\
&=\int\mathrm{d}\mu\,\left(\Theta\,\mathcal U_{t\cdot e_\nu}A\right)\cdot B=\mathsf b\left(\mathcal U_{t\cdot e_\nu}A,B\right).
\end{align*}
Now let $A\in\mathcal E_+$ and $t\ge0$. Then $\mathcal U_{t\cdot e_\nu}A$ and $\mathcal U_{2t\cdot e_\nu}A$ belong to $\mathcal E_+$. The identity just proved, with $B=\mathcal U_{t\cdot e_\nu}A$, and the Cauchy--Schwarz inequality for $\mathsf b$ on $\mathcal E_+$ give
\[
\mathsf b\left(\mathcal U_{t\cdot e_\nu}A,\mathcal U_{t\cdot e_\nu}A\right)=\mathsf b\left(A,\mathcal U_{2t\cdot e_\nu}A\right)\le\mathsf b(A,A)^{\frac12}\cdot\mathsf b\left(\mathcal U_{2t\cdot e_\nu}A,\mathcal U_{2t\cdot e_\nu}A\right)^{\frac12}.
\]
It follows by induction on $j$ that, for every $j\ge1$,
\begin{align*}
\mathsf b\left(\mathcal U_{t\cdot e_\nu}A,\mathcal U_{t\cdot e_\nu}A\right)&\le\mathsf b(A,A)^{1-2^{-j}}\cdot\mathsf b\left(\mathcal U_{2^j\cdot t\cdot e_\nu}A,\mathcal U_{2^j\cdot t\cdot e_\nu}A\right)^{2^{-j}}\\
&\le\mathsf b(A,A)^{1-2^{-j}}\cdot\|A\|_{L^2(\mu)}^{2^{1-j}}.
\end{align*}
Indeed, for $j=1$ the first inequality is the preceding one. If it holds for $j$, the preceding inequality with $2^j\cdot t$ in place of $t$ bounds $\mathsf b\left(\mathcal U_{2^j\cdot t\cdot e_\nu}A,\mathcal U_{2^j\cdot t\cdot e_\nu}A\right)$ by $\mathsf b(A,A)^{\frac12}\cdot\mathsf b\left(\mathcal U_{2^{j+1}\cdot t\cdot e_\nu}A,\mathcal U_{2^{j+1}\cdot t\cdot e_\nu}A\right)^{\frac12}$, and inserting this bound gives the first inequality for $j+1$. The second inequality follows from the bound on $\left|\mathsf b\right|$ proved above and from $\left\|\mathcal U_{2^j\cdot t\cdot e_\nu}A\right\|_{L^2(\mu)}=\|A\|_{L^2(\mu)}$. For $A=0$ there is nothing to prove. If $A\ne0$, then, as $j\to\infty$, $\|A\|_{L^2(\mu)}^{2^{1-j}}\to1$ and $\mathsf b(A,A)^{1-2^{-j}}\to\mathsf b(A,A)$, the latter also when $\mathsf b(A,A)=0$. Hence
\[
\mathsf b\left(\mathcal U_{t\cdot e_\nu}A,\mathcal U_{t\cdot e_\nu}A\right)\le\mathsf b(A,A)\qquad\left(A\in\mathcal E_+,\ t\ge0\right).
\]
In particular $\mathcal U_{t\cdot e_\nu}\mathcal N\subset\mathcal N$ for $t\ge0$. Therefore $[A]\mapsto\left[\mathcal U_{t\cdot e_\nu}A\right]$ is a well-defined linear map of $\mathcal E_+/\mathcal N$ into itself which does not increase the norm, and it extends uniquely to a contraction $P_t$ of $\mathcal H_\nu$. For $A,B\in\mathcal E_+$
\[
\left\langle[A],P_t[B]\right\rangle_{\mathcal H_\nu}=\mathsf b\left(A,\mathcal U_{t\cdot e_\nu}B\right)=\mathsf b\left(\mathcal U_{t\cdot e_\nu}A,B\right)=\left\langle P_t[A],[B]\right\rangle_{\mathcal H_\nu};
\]
by continuity $P_t$ is symmetric on $\mathcal H_\nu$, hence self-adjoint, being bounded. $P_0$ is the identity, and $P_sP_t=P_{s+t}$ for $s,t\ge0$, because $\mathcal U_{s\cdot e_\nu}\mathcal U_{t\cdot e_\nu}=\mathcal U_{(s+t)\cdot e_\nu}$ and both sides are continuous. For $A\in\mathcal E_+$
\[
\left\|P_t[A]-[A]\right\|_{\mathcal H_\nu}^2=\mathsf b\left(\mathcal U_{t\cdot e_\nu}A-A,\mathcal U_{t\cdot e_\nu}A-A\right)\le\left\|\mathcal U_{t\cdot e_\nu}A-A\right\|_{L^2(\mu)}^2,
\]
which tends to $0$ as $t\downarrow0$ by the strong continuity in Theorem~\ref{thm:continuum-infinite-1}(5). Since $\|P_t\|\le1$ and $\mathcal E_+/\mathcal N$ is dense in $\mathcal H_\nu$, it follows that $P_t\psi\to\psi$ as $t\downarrow0$ for every $\psi\in\mathcal H_\nu$. The identities
\begin{gather*}
P_{t+h}\psi-P_t\psi=P_t\left(P_h\psi-\psi\right)\qquad(h\ge0),\\
P_{t-h}\psi-P_t\psi=P_{t-h}\left(\psi-P_h\psi\right)\qquad(0\le h\le t)
\end{gather*}
then show that $t\mapsto P_t\psi$ is continuous on $[0,\infty)$. So $(P_t)_{t\ge0}$ is a strongly continuous semigroup of self-adjoint contractions.

For $s\ge0$ one has $P_{2s}=P_s^2=P_s^*P_s\ge0$; hence $P_t\ge0$ for all $t\ge0$. If $s>0$ and $P_{2s}\psi=0$, then $\left\|P_s\psi\right\|_{\mathcal H_\nu}^2=\left\langle\psi,P_{2s}\psi\right\rangle_{\mathcal H_\nu}=0$. Consequently, if $P_t\psi=0$ for some $t>0$, then $P_{2^{-j}\cdot t}\psi=0$ for all $j\ge0$ by induction, and $\psi=\lim_{j\to\infty}P_{2^{-j}\cdot t}\psi=0$. In particular $P_1$ is injective.

The Hilbert space $\mathcal H_\nu$ is separable. To see this, let $h_\kappa$, $\kappa\in\mathcal K$, be the Hermite functions of Section~\ref{sec:hermite-sobolev-spaces}, which are real tests, and let $J_1\subset J_2\subset\dots$ be finite subsets of the countable set $\mathcal K$ with union $\mathcal K$. For $f\in\mathcal S$ the partial sums $\sum_{\kappa\in J_l}\left\langle h_\kappa,f\right\rangle_{L^2}\cdot h_\kappa$ of the Hermite expansion of $f$ converge to $f$ in $\mathcal S$ as $l\to\infty$ (Section~\ref{sec:hermite-sobolev-spaces}), and every $F\in\mathcal S'$ is continuous on $\mathcal S$; therefore
\[
F(f)=\lim_{l\to\infty}\sum_{\kappa\in J_l}\left\langle h_\kappa,f\right\rangle_{L^2}\cdot F(h_\kappa)\qquad\left(F\in\mathcal S',\ f\in\mathcal S\right).
\]
Hence the functions $F\mapsto F(f)$, $f\in\mathcal S$, which generate $\mathcal C$, are measurable with respect to the $\sigma$-algebra generated by the countably many real functions $F\mapsto F(h_\kappa)$, $\kappa\in\mathcal K$. As the half-lines $(-\infty,q)$, $q\in\mathbb Q$, generate the Borel $\sigma$-algebra of $\mathbb R$, and as the $h_\kappa$ are themselves tests, it follows that $\mathcal C$ is the $\sigma$-algebra generated by the countable family $\mathcal G_0$ of the sets $\left\{F\in\mathcal S':F(h_\kappa)<q\right\}$, $\kappa\in\mathcal K$, $q\in\mathbb Q$. Let $\mathcal A_0$ be the smallest algebra of subsets of $\mathcal S'$ that contains $\mathcal G_0$; it is countable, since its elements are the finite unions of finite intersections of sets of $\mathcal G_0$ and of their complements. Let $D$ be the countable set of the finite linear combinations, with coefficients in $\mathbb Q+i\mathbb Q$, of the indicator functions $\mathbf 1_G$, $G\in\mathcal A_0$, and let $\bar D$ be its closure in $L^2(\mu)$, a closed complex linear subspace, since $\mathbb Q+i\mathbb Q$ is dense in $\mathbb C$. The sets $G'\in\mathcal C$ with $\mathbf 1_{G'}\in\bar D$ include the sets of $\mathcal A_0$, and they form a monotone class: if sets $G'_i$ of this kind increase or decrease to $G'$, then $\mathbf 1_{G'_i}\to\mathbf 1_{G'}$ in $L^2(\mu)$ by dominated convergence, $\mu$ being finite. By the monotone class theorem they are all the sets of $\mathcal C$. So $\bar D$ contains the simple functions, which are dense in $L^2(\mu)$; thus $\bar D=L^2(\mu)$, and $L^2(\mu)$ is separable. Its subset $\mathcal E_+$ is separable as well. Finally, the linear map $A\mapsto[A]$ of $\mathcal E_+$ onto $\mathcal E_+/\mathcal N$ does not increase the norm, and $\mathcal E_+/\mathcal N$ is dense in $\mathcal H_\nu$; so it maps a countable dense subset of $\mathcal E_+$ onto a countable dense subset of $\mathcal H_\nu$.

The spectral theorem for bounded self-adjoint operators, in its multiplication operator form, applies to $P_1$ on the separable Hilbert space $\mathcal H_\nu$: there are a finite measure space $(Y,\mathcal Y,\rho)$, a unitary operator $W\colon\mathcal H_\nu\to L^2(Y,\rho)$ and a bounded measurable function $p\colon Y\to\mathbb R$ such that $WP_1W^{-1}$ is the operator of multiplication by $p$ on $L^2(Y,\rho)$. For $\varphi\in L^2(Y,\rho)$ and $\psi=W^{-1}\varphi$ one has $\left\langle\psi,P_1\psi\right\rangle_{\mathcal H_\nu}=\int\mathrm{d}\rho\,p\cdot|\varphi|^2$ and $\|\psi\|_{\mathcal H_\nu}^2=\int\mathrm{d}\rho\,|\varphi|^2$, while $0\le\left\langle\psi,P_1\psi\right\rangle_{\mathcal H_\nu}\le\|\psi\|_{\mathcal H_\nu}^2$ because $P_1\ge0$ and $\|P_1\|\le1$. With $\varphi$ the indicator function of $\{p<0\}$, respectively of $\{p>1\}$, which lies in $L^2(Y,\rho)$ because $\rho$ is finite, these inequalities give $\int_{\{p<0\}}\mathrm{d}\rho\,p\ge0$, respectively $\int_{\{p>1\}}\mathrm{d}\rho\,(p-1)\le0$; as the integrand is negative, respectively positive, on the domain of integration, $\rho(\{p<0\})=0$ and $\rho(\{p>1\})=0$. Moreover $P_1W^{-1}\mathbf 1_{\{p=0\}}=W^{-1}\left(p\cdot\mathbf 1_{\{p=0\}}\right)=0$; as $P_1$ is injective, $\mathbf 1_{\{p=0\}}=0$ in $L^2(Y,\rho)$, that is, $\rho(\{p=0\})=0$. Hence $0<p\le1$ $\rho$-almost everywhere, and after a change of $p$ on a $\rho$-null set, which does not alter the multiplication operator, $0<p\le1$ everywhere on $Y$. For a measurable function $g\colon Y\to\mathbb R$ let $\mathsf M_g$ be the operator $\psi\mapsto W^{-1}\left(g\cdot W\psi\right)$, defined on the $\psi\in\mathcal H_\nu$ with $g\cdot W\psi\in L^2(Y,\rho)$; it is self-adjoint, non-negative if $g\ge0$, and bounded with norm at most $\sup|g|$ if $g$ is bounded. Define
\[
\mathsf H_\nu=-\log P_1=\mathsf M_{-\log p},\qquad\exp\left(-t\cdot\mathsf H_\nu\right)=\mathsf M_{p^t}\quad(t\ge0).
\]
Then $\mathsf H_\nu$ is self-adjoint and $\mathsf H_\nu\ge0$. For $j\ge0$ one has $P_{2^{-j}}=\mathsf M_{p^{2^{-j}}}$: for $j=0$ this is the choice of $p$, and if it holds for $j$, then $P_{2^{-j-1}}$ and $\mathsf M_{p^{2^{-j-1}}}$ are non-negative bounded operators whose squares both equal $P_{2^{-j}}$, so they coincide by the uniqueness of non-negative square roots. Hence
\[
P_{k\cdot2^{-j}}=\left(P_{2^{-j}}\right)^k=\mathsf M_{p^{k\cdot2^{-j}}}=\exp\left(-k\cdot2^{-j}\cdot\mathsf H_\nu\right)\qquad(k,j\ge0\text{ integers}).
\]
For each $\psi\in\mathcal H_\nu$ the maps $t\mapsto P_t\psi$ and $t\mapsto\mathsf M_{p^t}\psi$ are continuous on $[0,\infty)$, the second by dominated convergence, since $t\mapsto p^t$ is continuous and $0<p^t\le1$. As they agree at the numbers $k\cdot2^{-j}$, which are dense in $[0,\infty)$, one obtains $P_t=\exp\left(-t\cdot\mathsf H_\nu\right)$ for every $t\ge0$.

Finally $\mathcal U_{t\cdot e_\nu}1=1$, so $P_t\Omega=\Omega$ for $t\ge0$; in particular $P_1\Omega=\Omega$, that is, $p\cdot W\Omega=W\Omega$. Thus $p=1$ almost everywhere on the set where $W\Omega\ne0$, and $(\log p)\cdot W\Omega=0$: the vector $\Omega$ lies in the domain of $\mathsf H_\nu$ and $\mathsf H_\nu\Omega=0$. Together with $\|\Omega\|_{\mathcal H_\nu}=1$ this proves (1).

\emph{Part (2).} Let $\lambda\ne\nu$, $s\in\mathbb R$ and $a=s\cdot e_\lambda$. Then $r_\nu a=a$ and $a_\nu=0$. By what was shown above, $\Theta\,\mathcal U_a=\mathcal U_a\Theta$ and $\mathcal U_{\pm a}\mathcal E_+\subset\mathcal E_+$; since $\mathcal U_{-a}$ is the inverse of $\mathcal U_a$, the operator $\mathcal U_a$ maps $\mathcal E_+$ onto itself. Moreover $\mathcal U_a\mathcal U_{t\cdot e_\nu}=\mathcal U_{a+t\cdot e_\nu}=\mathcal U_{t\cdot e_\nu}\mathcal U_a$ for $t\in\mathbb R$. By the invariance of $\mu$ under $\tau_a^*$ (Theorem~\ref{thm:continuum-infinite-1}(5)), for $A,B\in L^2(\mu)$
\[
\mathsf b\left(\mathcal U_aA,\mathcal U_aB\right)=\int\mathrm{d}\mu\,\left(\Theta\,\mathcal U_aA\right)\cdot\left(\mathcal U_aB\right)=\int\mathrm{d}\mu\,\mathcal U_a\left[\left(\Theta A\right)\cdot B\right]=\mathsf b(A,B),
\]
so $\mathcal U_a$ preserves $\mathsf b$. Consequently $\mathcal U_a\mathcal N=\mathcal N$, and $[A]\mapsto\left[\mathcal U_aA\right]$ is a well-defined linear bijection of $\mathcal E_+/\mathcal N$ onto itself which preserves the inner product; it extends to a unitary operator $\hat{\mathcal U}_{s\cdot e_\lambda}$ of $\mathcal H_\nu$, which for $s=0$ is the identity. The identities $\hat{\mathcal U}_{s\cdot e_\lambda}\hat{\mathcal U}_{s'\cdot e_\lambda}=\hat{\mathcal U}_{(s+s')\cdot e_\lambda}$ for $s'\in\mathbb R$ and $\hat{\mathcal U}_{s\cdot e_\lambda}P_t=P_t\hat{\mathcal U}_{s\cdot e_\lambda}$ for $t\ge0$ hold on $\mathcal E_+/\mathcal N$ by the corresponding identities for $\mathcal U$, hence on $\mathcal H_\nu$ by continuity; and $\hat{\mathcal U}_{s\cdot e_\lambda}\Omega=\Omega$ because $\mathcal U_a1=1$. For $A\in\mathcal E_+$
\[
\left\|\hat{\mathcal U}_{s\cdot e_\lambda}[A]-[A]\right\|_{\mathcal H_\nu}^2\le\left\|\mathcal U_{s\cdot e_\lambda}A-A\right\|_{L^2(\mu)}^2,
\]
which tends to $0$ as $s\to0$ (Theorem~\ref{thm:continuum-infinite-1}(5)). As the operators $\hat{\mathcal U}_{s\cdot e_\lambda}$ are unitary and $\mathcal E_+/\mathcal N$ is dense in $\mathcal H_\nu$, $\hat{\mathcal U}_{s\cdot e_\lambda}\psi\to\psi$ as $s\to0$ for every $\psi\in\mathcal H_\nu$, and by the group property $s\mapsto\hat{\mathcal U}_{s\cdot e_\lambda}\psi$ is continuous on $\mathbb R$.

\emph{Part (3).} By the proof of Theorem~\ref{thm:moments-limit-point}(5), for $k\ge1$ and $\Xi\in\mathcal S_k$ with $\mathrm{supp}\,\Xi\subset\mathbb R^{3k}_{\nu,+}$ the function $Q_\Xi$ has a $\Sigma^+_{\nu,0}$-measurable representative. By the same proof, in which the maps $\Theta_m$ are formed with the direction $\nu$, and by the product rule in the proof of Theorem~\ref{thm:moments-limit-point}(4),
\[
\Theta Q_{\Psi'}=Q_{\Theta_m\Psi'},\qquad Q_{\Psi'}\cdot Q_{\Phi'}=Q_{\Psi'\otimes\Phi'}\qquad\left(\Psi'\in\mathcal S_m,\ \Phi'\in\mathcal S_n\right).
\]
For $k=0$, and for $m=0$ or $n=0$, these statements hold by the conventions of Section~\ref{sec:moments-tensor-tests}, namely $\mathcal S_0=\mathbb C$, $Q_{\Psi_0}=\Psi_0$, $\Theta_0\Psi_0=\overline{\Psi_0}$ and $\Psi_0\otimes\Xi=\Xi\otimes\Psi_0=\Psi_0\cdot\Xi$ for $\Psi_0\in\mathcal S_0$ and $\Xi\in\mathcal S_k$, $k\ge0$, together with the linearity of $\Xi\mapsto Q_\Xi$; in particular $Q_\Xi$ is constant for $k=0$. Hence $Q_\Psi,Q_\Phi\in\mathcal E_+$. Next, $\mathcal U_{t\cdot e_\nu}Q_\Phi=Q_{\tau^{\otimes n}_{t\cdot e_\nu}\Phi}$. For $n=0$ both sides equal $\Phi$. For $n\ge1$ and $F\in\mathcal S'$, Theorem~\ref{thm:moments-limit-point}(2), applied to $\tau_{t\cdot e_\nu}^*F$ and to $F$, gives
\[
Q_\Phi\left(\tau_{t\cdot e_\nu}^*F\right)=\left(\tau_{t\cdot e_\nu}^*F\right)^{\otimes n}(\Phi),\qquad Q_{\tau^{\otimes n}_{t\cdot e_\nu}\Phi}(F)=F^{\otimes n}\left(\tau^{\otimes n}_{t\cdot e_\nu}\Phi\right);
\]
the two right sides are continuous linear functionals of $\Phi\in\mathcal S_n$ which agree on the products $f_1\otimes\dots\otimes f_n$, where both equal $\prod_iF\left(\tau_{t\cdot e_\nu}f_i\right)$, and the finite Hermite sums are dense in $\mathcal S_n$ (Section~\ref{sec:hermite-sobolev-spaces}); so they coincide. Since $S_k(\Xi)=\int\mathrm{d}\mu\,Q_\Xi$ for $\Xi\in\mathcal S_k$ (Theorem~\ref{thm:moments-limit-point}(2)), and $\mathcal U_{t\cdot e_\nu}Q_\Phi\in\mathcal E_+$ for $t\ge0$, part (1) gives
\begin{align*}
S_{m+n}\left(\Theta_m\Psi\otimes\tau^{\otimes n}_{t\cdot e_\nu}\Phi\right)&=\int\mathrm{d}\mu\,Q_{\Theta_m\Psi\otimes\tau^{\otimes n}_{t\cdot e_\nu}\Phi}=\int\mathrm{d}\mu\,Q_{\Theta_m\Psi}\cdot Q_{\tau^{\otimes n}_{t\cdot e_\nu}\Phi}\\
&=\int\mathrm{d}\mu\,\left(\Theta Q_\Psi\right)\cdot\left(\mathcal U_{t\cdot e_\nu}Q_\Phi\right)=\mathsf b\left(Q_\Psi,\mathcal U_{t\cdot e_\nu}Q_\Phi\right)\\
&=\left\langle[Q_\Psi],P_t[Q_\Phi]\right\rangle_{\mathcal H_\nu}=\left\langle[Q_\Psi],\exp\left(-t\cdot\mathsf H_\nu\right)[Q_\Phi]\right\rangle_{\mathcal H_\nu}.
\end{align*}
In the notation of the proof of (1), $\exp\left(-t\cdot\mathsf H_\nu\right)=\mathsf M_{p^t}$, and as $t\to\infty$ the functions $p^t$ converge pointwise to the indicator function $\mathbf 1_{\{p=1\}}$, with $0<p^t\le1$. By dominated convergence $\exp\left(-t\cdot\mathsf H_\nu\right)\psi\to\mathsf M_{\mathbf 1_{\{p=1\}}}\psi$ for every $\psi\in\mathcal H_\nu$. The operator $\mathsf M_{\mathbf 1_{\{p=1\}}}$ is the orthogonal projection onto the subspace of the $\psi$ with $W\psi=0$ almost everywhere on $\{p<1\}$, and this subspace is $\ker\mathsf M_{-\log p}=\ker\mathsf H_\nu$. So $\mathsf M_{\mathbf 1_{\{p=1\}}}=E_0$, and the right side of the last display converges to $\left\langle[Q_\Psi],E_0[Q_\Phi]\right\rangle_{\mathcal H_\nu}$.
\end{proof}

\begin{proposition}[Laws with the four properties of limit points]\label{prop:laws-four-properties}
Assume that a probability measure $\mu$ on $(\mathcal S',\mathcal C)$ fulfils the conditions \eqref{eq:continuum-infinite-L1}, \eqref{eq:continuum-infinite-L2}, \eqref{eq:continuum-infinite-L3} and \eqref{eq:continuum-infinite-L4} introduced in Section~\ref{sec:properties-limit-points}. In the statements listed below, let $\mu$ take the place of the limit point, and let every object which they form from the limit point be formed from $\mu$; among these objects are the map $X$, the moments $S_n$ of Theorem~\ref{thm:continuum-infinite-1}(8) and the Schwinger functions $S_n$ of Theorem~\ref{thm:moments-limit-point}, the representation $b\mapsto\mathcal U_b$ and the maps $\Theta_{\nu,c}$ on $L^2(\mu)$, and the space $\mathcal H_\nu$ with the operators $P_t$ and $\mathsf H_\nu$. Then each of these statements is true:
\begin{enumerate}
\item[(a)] Theorem~\ref{thm:continuum-infinite-1}, parts (1), (3), (4), (5), (6), (7) and (8);
\item[(b)] Theorem~\ref{thm:moments-limit-point}, parts (1), (2), (4) and (5);
\item[(c)] every assertion of Corollary~\ref{cor:transfer-operators-coordinate}.
\end{enumerate}
Assume in addition that
\[
\Phi_\mu(-f)=\Phi_\mu(f)\qquad\text{for every real }f\in\mathcal S.
\]
Then the image of $\mu$ under the map $F\mapsto-F$ of $\mathcal S'$ is $\mu$, and the moments of Theorem~\ref{thm:continuum-infinite-1}(8), formed with $\mu$, vanish in every odd order:
\[
S_n(f_1,\dots,f_n)=0\qquad\text{for every odd }n\ge1\text{ and all real }f_1,\dots,f_n\in C_c^\infty(\mathbb R^3;\Lambda^2).
\]
\end{proposition}

\begin{proof}
In Theorem~\ref{thm:continuum-infinite-1}, Theorem~\ref{thm:moments-limit-point} and Corollary~\ref{cor:transfer-operators-coordinate} the measure is a limit point of an admissible family, and the proof of Theorem~\ref{thm:continuum-infinite-1} employs cutoffs $(N_k,T_k)$ for which $\mathsf P_{N_k,T_k}$ converges weakly on $H_{-4}$ to that limit point (Definition~\ref{def:limit-points}). Apart from the proof of part (2) and from the proof of the last sentence of part (9), these cutoffs, the quantities attached to them in Definition~\ref{def:limit-points} and the space $H_{-4}$ enter the proof of Theorem~\ref{thm:continuum-infinite-1} only in the derivations of the following properties of the limit point, which take place in parts (1), (4), (5), (6), (7) and (9), respectively: \eqref{eq:continuum-infinite-L1}; \eqref{eq:continuum-infinite-L2}; the first equality in \eqref{eq:continuum-infinite-L3}; the second equality in \eqref{eq:continuum-infinite-L3}; \eqref{eq:continuum-infinite-L4}; and the equality of the values of its characteristic functional at $f$ and at $-f$ for every real $f\in\mathcal S$. The other steps on which the proposition rests are the parts of the proofs of Theorem~\ref{thm:continuum-infinite-1}(1), (4), (5), (6) and (7) that follow these derivations; the deduction, in the proof of Theorem~\ref{thm:continuum-infinite-1}(9), of the invariance under $F\mapsto-F$ and of the vanishing of the odd moments; the proofs of Theorem~\ref{thm:continuum-infinite-1}(3) and (8); the proofs of Theorem~\ref{thm:moments-limit-point}(1), (2), (4) and (5); and the proof of Corollary~\ref{cor:transfer-operators-coordinate}. These steps are arguments about a probability measure on the measurable space $(\mathcal S',\mathcal C)$, and they use this measure only through the properties named in the paragraphs below. Applied to the measure $\mu$ of the proposition, they prove it, as follows.

\emph{Theorem~\ref{thm:continuum-infinite-1}(1).} The first inequality of this assertion is \eqref{eq:continuum-infinite-L1}. From it the proof of Theorem~\ref{thm:continuum-infinite-1}(1) derives the remaining assertions. For the probability space $(\mathcal S',\mathcal C,\mu)$, the real space $\mathcal S$ with the seminorm $\|\cdot\|_{L^2}$ and the map $f\mapsto F(f)$, the hypothesis \eqref{eq:continuum-infinite-5} of Proposition~\ref{prop:sub-gaussian-random} is \eqref{eq:continuum-infinite-L1}, and part (2) of that proposition gives
\[
\int\mathrm{d}\mu(F)\,F(f)^2\le\|f\|_{L^2}^2\qquad(f\in\mathcal S\ \text{real}).
\]
Since $\mathcal S$ is dense in $L^2(\mathbb R^3;\Lambda^2)$, it follows that $f\mapsto F(f)$ has exactly one bounded linear extension $X\colon L^2(\mathbb R^3;\Lambda^2)\to L^2(\mu)$. For real $f\in L^2(\mathbb R^3;\Lambda^2)$ and real $f_j\in\mathcal S$ with $f_j\to f$ in $L^2$, Fatou's lemma, applied along a subsequence on which $F(f_j)\to X(f)$ $\mu$-almost surely, together with \eqref{eq:continuum-infinite-L1}, gives
\[
\int\mathrm{d}\mu\,\exp\left(X(f)\right)\le\exp\left(\frac12\|f\|_{L^2}^2\right).
\]
Thus Theorem~\ref{thm:continuum-infinite-1}(1) holds for $\mu$. The last inequality is the hypothesis \eqref{eq:continuum-infinite-5} of Proposition~\ref{prop:sub-gaussian-random} for the probability space $(\mathcal S',\mathcal C,\mu)$, the real space $\mathcal T=L^2(\mathbb R^3;\Lambda^2)$ with $|f|=\|f\|_{L^2}$ and the map $X$; so every assertion of that proposition may be used under $\mu$ in this setting.

\emph{Theorem~\ref{thm:continuum-infinite-1}(3).} Its proof uses part (1), parts (2) and (3) of Proposition~\ref{prop:sub-gaussian-random} under $\mu$, the sums~\eqref{eq:continuum-infinite-4}, monotone convergence and Jensen's inequality.

\emph{Theorem~\ref{thm:continuum-infinite-1}(4).} The proof of this part first obtains \eqref{eq:continuum-infinite-L2}. Its remainder uses \eqref{eq:continuum-infinite-L2}, the continuity of $d^*\colon\mathcal S(\mathbb R^3)\to\mathcal S(\mathbb R^3;\Lambda^2)$ and of every $F\in\mathcal S'$ on $\mathcal S$, part (1), and Proposition~\ref{prop:sub-gaussian-random}(2) under $\mu$. Indeed, with a sequence of real functions in $C_c^\infty(\mathbb R^3)$ which is dense in $\mathcal S(\mathbb R^3)$, \eqref{eq:continuum-infinite-L2} and these continuities give $\mu\left(\left\{F\in\mathcal S':dF=0\right\}\right)=1$; this, Plancherel's theorem and the boundedness of $X$ give $X(f)=X(\Pi_cf)$ $\mu$-almost surely for $f\in L^2(\mathbb R^3;\Lambda^2)$; and the exponential bound of part (1) for the test $\Pi_cf$, together with Proposition~\ref{prop:sub-gaussian-random}(2), gives the bounds by $C_0(f,f)$.

\emph{Theorem~\ref{thm:continuum-infinite-1}(5) and (6).} The proofs of these parts first obtain the first and the second equality in \eqref{eq:continuum-infinite-L3}, respectively. Their remainders use \eqref{eq:continuum-infinite-L3}, part (1) and Lemma~\ref{lem:totality-exponentials}. For real $f_1,\dots,f_s\in\mathcal S$ and $t\in\mathbb R^s$, \eqref{eq:continuum-infinite-L3} for the test $\sum_it_i\cdot f_i$ states that each of the vectors $\left(F(\tau_bf_1),\dots,F(\tau_bf_s)\right)$ and $\left(F(U_rf_1),\dots,F(U_rf_s)\right)$ has under $\mu$ the characteristic function of $\left(F(f_1),\dots,F(f_s)\right)$; the uniqueness theorem for characteristic functions on $\mathbb R^s$ and Dynkin's $\pi$-$\lambda$ theorem then give $\mu\circ\left(\tau_b^*\right)^{-1}=\mu$ and $\mu\circ\left(U_r^*\right)^{-1}=\mu$. The unitarity of the operators $\mathcal U_b$ follows from the first invariance, and their strong continuity from Proposition~\ref{prop:sub-gaussian-random}(2) under $\mu$ and from Lemma~\ref{lem:totality-exponentials}.

\emph{Theorem~\ref{thm:continuum-infinite-1}(7).} Let $\nu\in\{1,2,3\}$, and let $\mathcal D^+_{\nu,c}$, $\Sigma^+_{\nu,c}$ and $\Theta_{\nu,c}$, $c\in\mathbb R$, be as in Theorem~\ref{thm:continuum-infinite-1}(7). For $c=0$ the proof of this part first obtains \eqref{eq:continuum-infinite-L4} for real tests $f_1,\dots,f_s\in\mathcal D^+_{\nu,0}$ and complex coefficients. Its remainder uses \eqref{eq:continuum-infinite-L4}, parts (1), (5) and (6), Lemma~\ref{lem:totality-exponentials}, and Proposition~\ref{prop:sub-gaussian-random}(1), (4) under $\mu$. For real $f_1,\dots,f_s\in\mathcal D^+_{\nu,0}$, $z\in\mathbb C^s$ and $A=\sum_{j=1}^sz_j\cdot\exp\left(iF(f_j)\right)$ one has
\[
\int\mathrm{d}\mu\,\left(\Theta_{\nu,0}A\right)\cdot A=\sum_{i,j=1}^s\overline{z_i}\cdot z_j\cdot\Phi_\mu\left(f_j-\mathsf R_\nu f_i\right),
\]
and the right side is non-negative by \eqref{eq:continuum-infinite-L4}. Part (6), taken with $r=r_\nu$, shows that the image of $\mu$ under $\mathsf R_\nu^*=U_{r_\nu}^*$ is $\mu$. Since $\mathsf R_\nu^*$ is its own inverse, $\Theta_{\nu,0}$ is then an antiunitary operator on $L^2(\mu)$ with $\Theta_{\nu,0}^2=1$. For $A,B\in L^2(\mu)$ put $\mathsf b(A,B)=\int\mathrm{d}\mu\,\left(\Theta_{\nu,0}A\right)\cdot B$; the Cauchy--Schwarz inequality gives $\left|\mathsf b(A,B)\right|\le\|A\|_{L^2(\mu)}\cdot\|B\|_{L^2(\mu)}$, and so $A\mapsto\mathsf b(A,A)$ is continuous on $L^2(\mu)$. Lemma~\ref{lem:totality-exponentials}, applied to $\mathcal D_0=\mathcal D^+_{\nu,0}$, extends the inequality $\mathsf b(A,A)\ge0$ to all $A\in L^2\left(\mathcal S',\Sigma^+_{\nu,0},\mu\right)$. Part (5) yields the case of an arbitrary $c\in\mathbb R$, and part (1) with Proposition~\ref{prop:sub-gaussian-random}(1), (4) yields the statements on the functions $X(f)$ and on the polynomials and exponentials in finitely many of them.

\emph{Theorem~\ref{thm:continuum-infinite-1}(8).} Its proof uses Proposition~\ref{prop:sub-gaussian-random}(4) under $\mu$, H\"older's inequality, the uniqueness theorem for characteristic functions on $\mathbb R^s$, Dynkin's $\pi$-$\lambda$ theorem and the density of $C_c^\infty(\mathbb R^3;\Lambda^2)$ in $\mathcal S$.

\emph{Theorem~\ref{thm:moments-limit-point}(1), (2), (4), (5).} By Theorem~\ref{thm:continuum-infinite-1}(1) for $\mu$, the map $X$ exists; so the Schwinger functions $S_n$ of $\mu$ on $L^2(\mathbb R^3;\Lambda^2)_{\mathbb C}$ are defined as in Theorem~\ref{thm:moments-limit-point}. The functions $Q_\Psi$ of Theorem~\ref{thm:moments-limit-point}(2) are given at each $F\in\mathcal S'$ by their series, without reference to a measure, and the functionals $\Psi\mapsto S_n(\Psi)=\int\mathrm{d}\mu\,Q_\Psi$ are formed with $\mu$. The proofs of these four assertions use Proposition~\ref{prop:sub-gaussian-random} under $\mu$, the parts (1) and (4)--(7) of Theorem~\ref{thm:continuum-infinite-1}, Lemma~\ref{lem:hermite-coefficients-tensor} and the Hermite expansions of Section~\ref{sec:hermite-sobolev-spaces}; by what precedes, all of them are available for $\mu$.

\emph{Corollary~\ref{cor:transfer-operators-coordinate}.} Its proof uses the parts (5), (6) and (7) of Theorem~\ref{thm:continuum-infinite-1}; the part (2) of Theorem~\ref{thm:moments-limit-point}; three facts established in the proofs of Theorem~\ref{thm:moments-limit-point}(4), (5), namely the identities
\[
Q_\Psi\cdot Q_\Phi=Q_{\Psi\otimes\Phi},\qquad\Theta_{\nu,0}Q_\Psi=Q_{\Theta_m\Psi}\qquad\left(\Psi\in\mathcal S_m,\ \Phi\in\mathcal S_n\right),
\]
where the maps $\Theta_m$ of Section~\ref{sec:moments-tensor-tests} are formed with the direction $\nu$, and the existence of a $\Sigma^+_{\nu,0}$-measurable representative of $Q_\Psi$ for every $\Psi\in\mathcal S_m$, $m\ge1$, with $\operatorname{supp}\Psi\subset\left\{y\in\mathbb R^{3m}:(y_i)_\nu\ge0\text{ for }i=1,\dots,m\right\}$; the Hermite expansions of Section~\ref{sec:hermite-sobolev-spaces}; and the spectral theorem for bounded self-adjoint operators. By what precedes, all of these are available for $\mu$. This completes the proof of (a), (b) and (c).

\emph{The case $\Phi_\mu(-f)=\Phi_\mu(f)$.} Assume this identity for every real $f\in\mathcal S$. As in the proof of Theorem~\ref{thm:continuum-infinite-1}(9), the invariance of $\mu$ under $F\mapsto-F$ is deduced from this identity alone, and the vanishing of the odd moments from the invariance together with the integrability of the products $\prod_iF(f_i)$, which comes from \eqref{eq:continuum-infinite-L1} through Theorem~\ref{thm:continuum-infinite-1}(1) and Proposition~\ref{prop:sub-gaussian-random}(4). Let $f_1,\dots,f_s\in\mathcal S$ be real and $t\in\mathbb R^s$. The hypothesis, applied to the test $\sum_it_i\cdot f_i$, gives
\[
\int\mathrm{d}\mu(F)\,\exp\left(-i\sum_{i=1}^st_i\cdot F(f_i)\right)=\int\mathrm{d}\mu(F)\,\exp\left(i\sum_{i=1}^st_i\cdot F(f_i)\right).
\]
Hence the random vectors $\left(-F(f_1),\dots,-F(f_s)\right)$ and $\left(F(f_1),\dots,F(f_s)\right)$ have the same characteristic function under $\mu$, and, by the uniqueness theorem for characteristic functions on $\mathbb R^s$, the same law. So the image of $\mu$ under $F\mapsto-F$ agrees with $\mu$ on the sets $\left\{F\in\mathcal S':\left(F(f_1),\dots,F(f_s)\right)\in B\right\}$, $B\subset\mathbb R^s$ a Borel set. These sets form a $\pi$-system, the $\sigma$-algebra generated by them is $\mathcal C$, and Dynkin's $\pi$-$\lambda$ theorem shows that $\mu$ is invariant under $F\mapsto-F$. Now let $n\ge1$ and let $f_1,\dots,f_n\in C_c^\infty(\mathbb R^3;\Lambda^2)$ be real. The function $F\mapsto\prod_{i=1}^nF(f_i)$ is $\mu$-integrable by Theorem~\ref{thm:continuum-infinite-1}(1) and Proposition~\ref{prop:sub-gaussian-random}(4), and the invariance just proved gives
\[
S_n(f_1,\dots,f_n)=\int\mathrm{d}\mu(F)\,\prod_{i=1}^n\left(-F(f_i)\right)=(-1)^n\cdot S_n(f_1,\dots,f_n).
\]
For odd $n$ it follows that $S_n(f_1,\dots,f_n)=0$.
\end{proof}

\begin{remark}[Equivalent forms of the conditions]
Let $\mu$ be a probability measure on $(\mathcal S',\mathcal C)$.

The measure $\mu$ satisfies \eqref{eq:continuum-infinite-L3} if and only if $\mu\circ\left(\tau_b^*\right)^{-1}=\mu$ for every $b\in\mathbb R^3$ and $\mu\circ\left(U_r^*\right)^{-1}=\mu$ for every $r\in O_{\mathrm{cub}}$. Indeed, if \eqref{eq:continuum-infinite-L3} holds, the arguments of the proofs of Theorem~\ref{thm:continuum-infinite-1}(5), (6) which follow the derivation of \eqref{eq:continuum-infinite-L3} give these two invariances, as recalled in the proof of Proposition~\ref{prop:laws-four-properties}. Conversely, if $\mu\circ\left(\tau_b^*\right)^{-1}=\mu$, then for every real $f\in\mathcal S$
\[
\Phi_\mu(\tau_bf)=\int\mathrm{d}\mu(F)\,\exp\left(i\left(\tau_b^*F\right)(f)\right)=\int\mathrm{d}\mu(F)\,\exp\left(iF(f)\right)=\Phi_\mu(f),
\]
and in the same way $\mu\circ\left(U_r^*\right)^{-1}=\mu$ implies $\Phi_\mu(U_rf)=\Phi_\mu(f)$.

The measure $\mu$ satisfies \eqref{eq:continuum-infinite-L2} if and only if $\mu\left(\left\{F\in\mathcal S':dF=0\right\}\right)=1$. Indeed, if \eqref{eq:continuum-infinite-L2} holds, the argument of the proof of Theorem~\ref{thm:continuum-infinite-1}(4) which follows the derivation of \eqref{eq:continuum-infinite-L2} gives $\mu\left(\left\{F\in\mathcal S':dF=0\right\}\right)=1$. Conversely, since $(dF)(g)=F\left(d^*g\right)$, the set $\left\{F\in\mathcal S':dF=0\right\}$ is contained in $\left\{F\in\mathcal S':F\left(d^*g\right)=0\right\}$ for every real $g\in C_c^\infty(\mathbb R^3)$.

Consequently, $\mu$ satisfies \eqref{eq:continuum-infinite-L1}, \eqref{eq:continuum-infinite-L3} and \eqref{eq:continuum-infinite-L4} if and only if it satisfies the hypothesis \eqref{eq:setup-H} of Section~\ref{sec:hypothesis-law}; there the measurable space $(\mathcal S',\mathcal C)$ is denoted $(Q,\Sigma)$, and $X(f)$ is the function $F\mapsto F(f)$ (Section~\ref{sec:test-forms-distributions}). Indeed, \eqref{eq:continuum-infinite-L1} is the bound \eqref{eq:setup-G}. In Section~\ref{sec:euclidean-maps-actions} a Euclidean map $g$ acts on $\mathcal S'$ by $(gF)(f)=F\left(g^{-1}\cdot f\right)$; for the translation $x\mapsto x+a$, $a\in\mathbb R^3$, and for the map $x\mapsto wx$, $w\in O_{\mathrm{cub}}$, the action $f\mapsto g\cdot f$ on tests is $\tau_a$ and $U_w$, respectively, so that the action on $\mathcal S'$ is $\tau_{-a}^*$ and $U_{w^{-1}}^*$. As $a\mapsto-a$ and $w\mapsto w^{-1}$ are bijections of $\mathbb R^3$ and of $O_{\mathrm{cub}}$, the first paragraph shows that \eqref{eq:continuum-infinite-L3} is equivalent to \eqref{eq:ergodic-components-TJ} for $J=\{1,2,3\}$ together with $w_*\mu=\mu$ for every $w\in O_{\mathrm{cub}}$. Finally, for $\nu\in\{1,2,3\}$ the reflection $\theta_\nu$ of Section~\ref{sec:euclidean-maps-actions} acts on tests as $\mathsf R_\nu$, and the space $\mathcal D(\{x_\nu>0\})$ of Section~\ref{sec:test-forms-distributions} consists of the real $f\in C_c^\infty(\mathbb R^3;\Lambda^2)$ with $\operatorname{supp}f\subset\{y\in\mathbb R^3:y_\nu>0\}$; since $\int\mathrm{d}\mu(F)\,\exp\left(i\left(F(f_q)-F(\mathsf R_\nu f_p)\right)\right)=\Phi_\mu\left(f_q-\mathsf R_\nu f_p\right)$ for such tests $f_p$, $f_q$, the inequality \eqref{eq:continuum-infinite-L4} for the direction $\nu$ is \eqref{eq:setup-RPj} for $j=\nu$.
\end{remark}

\subsection{Three examples of admissible families}\label{sec:three-examples-admissible}

\begin{lemma}[Poincar\'e inequality on a cube]
Let $\ell>0$ and $z\in\mathbb R^3$, let $Q=z+[0,\ell]^3$, and let $u\colon\mathbb R^3\to\mathbb R$ be continuously differentiable, with gradient $\nabla u=(\partial_1u,\partial_2u,\partial_3u)$. Let
\[
\langle u\rangle_Q=\ell^{-3}\int_Q\mathrm{d}y\,u(y)
\]
be the mean value of $u$ over $Q$. Then
\[
\int_Q\mathrm{d}y\,\left(u(y)-\langle u\rangle_Q\right)^2\le\frac{\ell^2}{\pi^2}\int_Q\mathrm{d}y\,\left|\nabla u(y)\right|^2 .
\]
\end{lemma}

\begin{proof}
After a translation one may assume $z=0$. For $v,w\in L^2(Q)$ write $(v,w)_Q=\int_Q\mathrm{d}y\,v(y)\cdot w(y)$ and $\|w\|_Q=(w,w)_Q^{\frac12}$. For $k\in\mathbb N_0$ and $t\in[0,\ell]$ put
\[
c_k(t)=\cos\frac{\pi kt}{\ell},\qquad s_k(t)=\sin\frac{\pi kt}{\ell}.
\]
The product formulas $2c_k\cdot c_{k'}=c_{|k-k'|}+c_{k+k'}$ and $2s_k\cdot s_{k'}=c_{|k-k'|}-c_{k+k'}$, together with $\int_0^\ell\mathrm{d}t\,c_j(t)=0$ for every integer $j\ge1$, give
\[
\int_0^\ell\mathrm{d}t\,c_k(t)\cdot c_{k'}(t)=
\begin{cases}
0&\text{if }k\ne k',\\
\ell&\text{if }k=k'=0,\\
\frac{\ell}{2}&\text{if }k=k'\ge1,
\end{cases}
\qquad
\int_0^\ell\mathrm{d}t\,s_k(t)\cdot s_{k'}(t)=
\begin{cases}
0&\text{if }k\ne k',\\
\frac{\ell}{2}&\text{if }k=k',
\end{cases}
\]
where $k,k'\in\mathbb N_0$ in the first formula and $k,k'\ge1$ in the second.

The functions $c_k$, $k\in\mathbb N_0$, form an orthogonal basis of $L^2(0,\ell)$. Indeed, let $v\in L^2(0,\ell)$ be orthogonal to all of them, and let $\tilde v$ be the even function on $(-\ell,\ell)$ that coincides with $v$ on $(0,\ell)$. Then $\tilde v$ is orthogonal in $L^2(-\ell,\ell)$ to the functions $t\mapsto\cos\frac{\pi kt}{\ell}$, $k\ge0$, and, being even, to the odd functions $t\mapsto\sin\frac{\pi kt}{\ell}$, $k\ge1$. These functions make up the trigonometric system of period $2\ell$, which is complete in $L^2(-\ell,\ell)$; hence $\tilde v=0$, and $v=0$.

For $k=(k_1,k_2,k_3)\in\mathbb N_0^3$ and $y\in Q$ let $\phi_k(y)=c_{k_1}(y_1)\cdot c_{k_2}(y_2)\cdot c_{k_3}(y_3)$. The products of the elements of an orthogonal basis of $L^2(0,\ell)$ form an orthogonal basis of $L^2\left((0,\ell)^3\right)$. So $(\phi_k)_{k\in\mathbb N_0^3}$ is an orthogonal basis of $L^2(Q)$, and Parseval's identity reads
\[
\|v\|_Q^2=\sum_{k\in\mathbb N_0^3}\frac{(v,\phi_k)_Q^2}{\|\phi_k\|_Q^2}\qquad\left(v\in L^2(Q)\right).
\]
Apply it to $v=u-\langle u\rangle_Q$. Since $\phi_0=1$, one has $(v,\phi_0)_Q=0$. If $k\ne0$, some $k_j$ is positive, so that $\int_Q\mathrm{d}y\,\phi_k(y)=0$ and $(v,\phi_k)_Q=(u,\phi_k)_Q$. Therefore
\[
\int_Q\mathrm{d}y\,\left(u(y)-\langle u\rangle_Q\right)^2=\sum_{k\in\mathbb N_0^3\setminus\{0\}}\frac{(u,\phi_k)_Q^2}{\|\phi_k\|_Q^2}.
\]

Now fix $j\in\{1,2,3\}$. For $k\in\mathbb N_0^3$ with $k_j\ge1$ let $\psi_{k,j}$ be the function obtained from $\phi_k$ when the factor $c_{k_j}(y_j)$ is replaced by $s_{k_j}(y_j)$. By the integrals computed above, the functions $\psi_{k,j}$, $k\in\mathbb N_0^3$ with $k_j\ge1$, are pairwise orthogonal in $L^2(Q)$, and $\|\psi_{k,j}\|_Q=\|\phi_k\|_Q$. Since $s_{k_j}(0)=s_{k_j}(\ell)=0$ and $s_{k_j}'=\frac{\pi k_j}{\ell}\cdot c_{k_j}$, an integration by parts in the variable $y_j$, followed by an integration over the two other variables (Fubini's theorem), gives
\[
(\partial_ju,\psi_{k,j})_Q=-\frac{\pi k_j}{\ell}\cdot(u,\phi_k)_Q .
\]
Bessel's inequality for the orthogonal family $(\psi_{k,j})_{k_j\ge1}$ therefore yields
\[
\int_Q\mathrm{d}y\,\left(\partial_ju(y)\right)^2\ge\sum_{k\in\mathbb N_0^3,\ k_j\ge1}\frac{(\partial_ju,\psi_{k,j})_Q^2}{\|\psi_{k,j}\|_Q^2}=\frac{\pi^2}{\ell^2}\sum_{k\in\mathbb N_0^3,\ k_j\ge1}k_j^2\cdot\frac{(u,\phi_k)_Q^2}{\|\phi_k\|_Q^2}.
\]
Add these three inequalities, $j=1,2,3$. For $k\ne0$ the term $\frac{(u,\phi_k)_Q^2}{\|\phi_k\|_Q^2}$ then carries the factor $\frac{\pi^2}{\ell^2}\left(k_1^2+k_2^2+k_3^2\right)$, which is at least $\frac{\pi^2}{\ell^2}$, as $k_1,k_2,k_3$ are integers that do not all vanish. Together with the expansion of the left side of the assertion obtained above, this gives
\[
\int_Q\mathrm{d}y\,\left|\nabla u(y)\right|^2\ge\frac{\pi^2}{\ell^2}\sum_{k\in\mathbb N_0^3\setminus\{0\}}\frac{(u,\phi_k)_Q^2}{\|\phi_k\|_Q^2}=\frac{\pi^2}{\ell^2}\int_Q\mathrm{d}y\,\left(u(y)-\langle u\rangle_Q\right)^2 .
\]
\end{proof}

\begin{proposition}[Scaled free laws and their mixtures]\label{prop:scaled-free-laws}
For a cutoff $(N,T)$ let $\mathsf P^{\mathrm{free}}_{\eta,R}$ be the averaged law at charge $q=0$, as in Theorem~\ref{thm:estimates-fixed-cutoff}, and for $c\ge0$ let $\mathsf P^{(c)}_{\eta,R}$ be the law of $\sqrt{c}\cdot F$ under $\mathsf P^{\mathrm{free}}_{\eta,R}$, that is, the image of $\mathsf P^{\mathrm{free}}_{\eta,R}$ under the map $F\mapsto\sqrt{c}\cdot F$ of $C^2(\Lambda_{\eta,R})$. Let $0\le a\le b\le1$, and let $(\mathrm{Pr}_{N,T})$ be one of the following three families, indexed by the cutoffs $(N,T)$:
\begin{equation}\tag{E1}\label{eq:continuum-infinite-E1}
\mathrm{Pr}_{N,T}=\mathsf P^{(1)}_{\eta,R};
\end{equation}
\begin{equation}\tag{E2}\label{eq:continuum-infinite-E2}
\mathrm{Pr}_{N,T}=\frac12\mathsf P^{(a)}_{\eta,R}+\frac12\mathsf P^{(b)}_{\eta,R};
\end{equation}
\begin{equation}\tag{E3}\label{eq:continuum-infinite-E3}
\mathrm{Pr}_{N,T}=\mathsf P^{(a)}_{\eta,R}\quad\text{if $N$ is even},\qquad\mathrm{Pr}_{N,T}=\mathsf P^{(b)}_{\eta,R}\quad\text{if $N$ is odd}.
\end{equation}
Then every $\mathrm{Pr}_{N,T}$ is a Borel probability measure on $\operatorname{ran}\,d\subset\ker d\subset C^2(\Lambda_{\eta,R})$, and the family $(\mathrm{Pr}_{N,T})$ is admissible (Definition~\ref{def:admissible-families}); consequently Theorems~\ref{thm:estimates-fixed-cutoff}--\ref{thm:moments-limit-point} and Corollary~\ref{cor:transfer-operators-coordinate} apply to it. Moreover, the following hold.
\begin{enumerate}
\item[(i)] Every limit point $\mu$ of the family~\eqref{eq:continuum-infinite-E1} is a centred Gaussian measure: for all real $f_1,\dots,f_s\in\mathcal S(\mathbb R^3;\Lambda^2)$ the law of $\left(F(f_1),\dots,F(f_s)\right)$ under $\mu$ is a centred Gaussian law on $\mathbb R^s$. For every cutoff, $\mathsf P^{(0)}_{\eta,R}=\delta_0$ is the point mass at the zero cochain; accordingly the family $\left(\mathsf P^{(0)}_{\eta,R}\right)$, which is the family~\eqref{eq:continuum-infinite-E2} for $a=b=0$, has exactly one limit point, namely the point mass at the zero distribution $F\equiv0$.
\item[(ii)] Let $a<b$. Let $\beta\in C_c^\infty(\mathbb R^3;\Lambda^2)$ be real, and let $f=dd^*\beta$; here, as in Section~\ref{sec:probes-discretisation}, the components of $\beta$ are extended by $\beta_{\nu\mu}=-\beta_{\mu\nu}$, $d^*\beta$ is the one-form with the components $(d^*\beta)_\nu=-\sum_{\mu=1}^3\partial_\mu\beta_{\mu\nu}$, and $(dd^*\beta)_{\mu\nu}=\partial_\mu(d^*\beta)_\nu-\partial_\nu(d^*\beta)_\mu$. Suppose that $f\ne0$. Then every limit point of the family~\eqref{eq:continuum-infinite-E2} satisfies, with its Schwinger functions $S_n$ (Theorem~\ref{thm:moments-limit-point}),
\[
S_4(f,f,\tau_xf,\tau_xf)-S_2(f,f)^2\ge\frac{(b-a)^2}{4}\cdot\|f\|_{L^2}^4>0\qquad\text{for every }x\in\mathbb R^3.
\]
In particular, $S_4(f,f,\tau_xf,\tau_xf)$ does not converge to $S_2(f,f)\cdot S_2(\tau_xf,\tau_xf)=S_2(f,f)^2$ as $|x|\to\infty$.
\item[(iii)] Let $a<b$. Then the family~\eqref{eq:continuum-infinite-E3} has two different limit points.
\end{enumerate}
\end{proposition}

\begin{proof}
Let $(N,T)$ be a cutoff and $c\ge0$, and let $\mathsf Q^{(c)}_{N,T}$ denote the law under $\mathsf P^{(c)}_{\eta,R}$ of the distribution $F^\sharp$ of Section~\ref{sec:test-forms-test}; by Section~\ref{sec:existence-limit-points}, it is a Borel probability measure on $H_{-3}\subset H_{-4}$. As $F\mapsto F^\sharp$ is linear, $\left(\sqrt c\cdot F\right)^\sharp=\sqrt c\cdot F^\sharp$; hence $\mathsf Q^{(c)}_{N,T}$ is the image of $\mathsf Q^{(1)}_{N,T}$ under the map $S\mapsto\sqrt c\cdot S$ of $H_{-4}$. The law $\mathsf P_{N,T}$ of $F^\sharp$ under $\mathrm{Pr}_{N,T}$ (Section~\ref{sec:existence-limit-points}) is $\mathsf Q^{(1)}_{N,T}$ for the family~\eqref{eq:continuum-infinite-E1}, it is $\frac12\mathsf Q^{(a)}_{N,T}+\frac12\mathsf Q^{(b)}_{N,T}$ for the family~\eqref{eq:continuum-infinite-E2}, and for the family~\eqref{eq:continuum-infinite-E3} it is $\mathsf Q^{(a)}_{N,T}$ if $N$ is even and $\mathsf Q^{(b)}_{N,T}$ if $N$ is odd.

\emph{Admissibility.} By the proof of Theorem~\ref{thm:estimates-fixed-cutoff}, $\mathsf P^{\mathrm{free}}_{\eta,R}$ coincides with the averaged law $\bar{\mathsf P}_{\eta,R;0,\boldsymbol{\mu}}$ for every value of $\boldsymbol{\mu}$, and under it $F$ is the standard Gaussian vector of the Euclidean space $\left(\operatorname{ran}\,d,|\cdot|_\eta\right)$. Thus $\left(\mathsf P^{\mathrm{free}}_{\eta,R}\right)_{N,T}$ is the family of the averaged laws of the parameter assignment $(N,T)\mapsto(0,0)$, and it is admissible by Corollary~\ref{cor:averaged-laws-admissible}. Let $0\le c\le1$. The map $F\mapsto\sqrt c\cdot F$ sends $\operatorname{ran}\,d$ into itself, so that $\mathsf P^{(c)}_{\eta,R}$ is a Borel probability measure on $\operatorname{ran}\,d$. The maps $v_*$, $v\in\eta\mathbb Z^3$, and $r_*$, $r\in O_{\mathrm{cub}}$, are linear and therefore commute with $F\mapsto\sqrt c\cdot F$; since $\mathsf P^{\mathrm{free}}_{\eta,R}$ is invariant under them, so is its image $\mathsf P^{(c)}_{\eta,R}$, and this is property~\eqref{eq:continuum-infinite-Ad1}. For real $g\in C^2(\Lambda_{\eta,R})$, property~\eqref{eq:continuum-infinite-Ad2} of $\mathsf P^{\mathrm{free}}_{\eta,R}$, applied to the cochain $\sqrt c\cdot g$, yields
\[
\int\mathrm{d}\mathsf P^{(c)}_{\eta,R}(F)\,\exp\langle g,F\rangle_\eta=\int\mathrm{d}\mathsf P^{\mathrm{free}}_{\eta,R}(F)\,\exp\left\langle\sqrt c\cdot g,F\right\rangle_\eta\le\exp\left(\frac{c}{2}|g|_\eta^2\right)\le\exp\left(\frac12|g|_\eta^2\right),
\]
which is property~\eqref{eq:continuum-infinite-Ad2} of $\mathsf P^{(c)}_{\eta,R}$. Let $\nu\in\{1,2,3\}$, let $g_1,\dots,g_s\in C^2(\Lambda_{\eta,R})$ be real and supported on $\mathcal P^{+\circ}_\nu$, and let $w\in\mathbb C^s$. The cochains $\sqrt c\cdot g_1,\dots,\sqrt c\cdot g_s$ are real and supported on $\mathcal P^{+\circ}_\nu$, and $\hat{\mathsf R}_\nu$ is linear; property~\eqref{eq:continuum-infinite-Ad3} of $\mathsf P^{\mathrm{free}}_{\eta,R}$ for these cochains therefore gives
\begin{align*}
&\sum_{p,q=1}^s\overline{w_p}\cdot w_q\int\mathrm{d}\mathsf P^{(c)}_{\eta,R}(F)\,\exp\left(i\left\langle g_q-\hat{\mathsf R}_\nu g_p,F\right\rangle_\eta\right)\\
&\qquad=\sum_{p,q=1}^s\overline{w_p}\cdot w_q\int\mathrm{d}\mathsf P^{\mathrm{free}}_{\eta,R}(F)\,\exp\left(i\left\langle\sqrt c\cdot g_q-\hat{\mathsf R}_\nu\left(\sqrt c\cdot g_p\right),F\right\rangle_\eta\right)\ge0,
\end{align*}
which is property~\eqref{eq:continuum-infinite-Ad3} of $\mathsf P^{(c)}_{\eta,R}$. A convex combination of two Borel probability measures on $\operatorname{ran}\,d$ with the properties~\eqref{eq:continuum-infinite-Ad1}, \eqref{eq:continuum-infinite-Ad2} and~\eqref{eq:continuum-infinite-Ad3} has these properties again: invariance under a map is preserved by convex combinations, and the left sides of~\eqref{eq:continuum-infinite-Ad2} and~\eqref{eq:continuum-infinite-Ad3} depend linearly on the measure. In particular $\frac12\mathsf P^{(a)}_{\eta,R}+\frac12\mathsf P^{(b)}_{\eta,R}$ has the three properties. As the conditions of Definition~\ref{def:admissible-families} concern each cutoff separately, the families~\eqref{eq:continuum-infinite-E1}, \eqref{eq:continuum-infinite-E2} and~\eqref{eq:continuum-infinite-E3} are admissible.

\emph{Gaussian limit points.} Let $\mu$ be a limit point of the family~\eqref{eq:continuum-infinite-E1}, and let $(N_k,T_k)$, $\eta_k$, $R_k$, $E_k$ and $X_k$ be as in Definition~\ref{def:limit-points}; here $E_k$ is the expectation with respect to $\mathsf P^{\mathrm{free}}_{\eta_k,R_k}$. For a real test $g$ one has $X_k(g)=\left\langle g_{\eta_k},F\right\rangle_{\eta_k}$ by~\eqref{eq:continuum-infinite-2}, and $F$ is the standard Gaussian vector of $\left(\operatorname{ran}\,d,|\cdot|_{\eta_k}\right)$; hence $X_k(g)$ is a centred Gaussian random variable, and
\[
E_k\exp\left(iX_k(g)\right)=\exp\left(-\frac12E_k\left[X_k(g)^2\right]\right).
\]
Let $f_1,\dots,f_s\in\mathcal S(\mathbb R^3;\Lambda^2)$ be real, let $t\in\mathbb R^s$, and apply this identity to $g=\sum_{j=1}^st_j\cdot f_j$. By Theorem~\ref{thm:continuum-infinite-1}(2), as $k\to\infty$ the left side tends to $\int\mathrm{d}\mu(F)\,\exp\left(iF(g)\right)$, and $E_k\left[X_k(g)^2\right]$ tends to $\int\mathrm{d}\mu(F)\,F(g)^2$. Consequently
\[
\int\mathrm{d}\mu(F)\,\exp\left(i\sum_{j=1}^st_j\cdot F(f_j)\right)=\exp\left(-\frac12\sum_{j,l=1}^st_j\cdot t_l\cdot\int\mathrm{d}\mu(F)\,F(f_j)\cdot F(f_l)\right)\qquad\left(t\in\mathbb R^s\right).
\]
The matrix $\left(\int\mathrm{d}\mu(F)\,F(f_j)\cdot F(f_l)\right)_{j,l}$ is symmetric and positive semidefinite, and the right side is the characteristic function of the centred Gaussian law on $\mathbb R^s$ with this covariance matrix. By the uniqueness theorem for characteristic functions, this Gaussian law is the law of $\left(F(f_1),\dots,F(f_s)\right)$ under $\mu$.

For every cutoff, $\mathsf P^{(0)}_{\eta,R}$ is the image of a probability measure under the constant map $F\mapsto0$, that is, the point mass $\delta_0$, and $\mathsf Q^{(0)}_{N,T}$ is the point mass at $0\in H_{-4}$. Every sequence of cutoffs, in particular every sequence with $N_k\to\infty$ and $T_k\to\infty$, therefore gives this point mass as the weak limit, and it is the only limit point of the family $\left(\mathsf P^{(0)}_{\eta,R}\right)$; as a measure on $(\mathcal S',\mathcal C)$ it is the point mass at the zero distribution. This proves~(i).

\emph{The variance of a limit point of~\eqref{eq:continuum-infinite-E1}.} Let $\mu_1$ be a limit point of the family~\eqref{eq:continuum-infinite-E1}, let $(N_k,T_k)$, $\eta_k$, $R_k$, $E_k$ and $X_k$ be as in Definition~\ref{def:limit-points} for this family and $\mu_1$, and let
\[
C(g,h)=\int\mathrm{d}\mu_1(F)\,F(g)\cdot F(h)\qquad\left(g,h\in\mathcal S(\mathbb R^3;\Lambda^2)\text{ real}\right)
\]
be the covariance of $\mu_1$. Let $\beta\in C_c^\infty(\mathbb R^3;\Lambda^2)$ be real and different from zero, and let $f=dd^*\beta$, as in~(ii). Then
\[
C(f,f)=\|f\|_{L^2}^2 .
\]
For the proof, let $a'$ and $f^{(\eta)}$, for $0<\eta\le1$, be formed from $\beta$ as in Section~\ref{sec:probes-discretisation}:
\[
a'_\nu=-\sum_{\mu=1}^3\delta_{\eta,\mu}\beta_{\mu\nu},\qquad f^{(\eta)}_{\mu\nu}=\delta_{\eta,\mu}a'_\nu-\delta_{\eta,\nu}a'_\mu ,
\]
where $\delta_{\eta,j}$ is the centred difference quotient of Lemma~\ref{lem:discrete-codifferentials-three}. Let $D_\beta$ and $Q_\beta=y_\beta+\left[-\frac{D_\beta}{2},\frac{D_\beta}{2}\right]^3$ be the number and the cube attached to $\beta$ in Section~\ref{sec:probes-discretisation}; thus $\operatorname{supp}\beta+[-2,2]^3\subset Q_\beta$. Since $R_k\to\infty$, there is $k_0$ such that $R_k>D_\beta$ for all $k\ge k_0$. For such $k$ the cube $Q_\beta$ lies in the interior $y_\beta+\left(-\frac{R_k}{2},\frac{R_k}{2}\right)^3$ of a period cube of the torus of side $R_k$; moreover $0<\eta_k\le1$, and $R_k/\eta_k$ is an integer. Hence Theorem~\ref{thm:free-coefficient-1}(b1)--(b3) applies to $\beta$ with $\eta=\eta_k$ and $R=R_k$, and gives, for $k\ge k_0$,
\begin{gather*}
\left(f^{(\eta_k)}\right)_{\eta_k}=d_{\eta_k}a_{\eta_k}=\eta_k^{-1}\cdot d\,a_{\eta_k}\in\operatorname{ran}\,d,\\
\operatorname{supp}f^{(\eta_k)}\subset\operatorname{supp}\beta+[-\eta_k,\eta_k]^3\subset Q_\beta,\\
\left\|f^{(\eta_k)}-f\right\|_{H^1}\le\eta_k^2\cdot\|\beta\|_{H^5},
\end{gather*}
where $a_{\eta_k}$ is the one-cochain of Section~\ref{sec:probes-discretisation} formed on $\Lambda_{\eta_k,R_k}$, the two-form $f^{(\eta_k)}$ is smooth, and $\|\cdot\|_{H^s}$ are the norms of Section~\ref{sec:lattices-tests-function}. Also $\operatorname{supp}f\subset\operatorname{supp}\beta\subset Q_\beta$. Consequently, for $k\ge k_0$, the real tests $f$, $f^{(\eta_k)}$ and $f-f^{(\eta_k)}$ are supported in $y_\beta+\left(-\frac{R_k}{2},\frac{R_k}{2}\right)^3$.

Under $\mathsf P^{\mathrm{free}}_{\eta_k,R_k}$ the vector $F$ is the standard Gaussian vector of $\left(\operatorname{ran}\,d,|\cdot|_{\eta_k}\right)$, so that $\langle h,F\rangle_{\eta_k}$ has the variance $|h|_{\eta_k}^2$ for every $h\in\operatorname{ran}\,d$; this is the last assertion of Theorem~\ref{thm:estimates-fixed-cutoff}, which holds with equality for the law $\mathsf P^{\mathrm{free}}_{\eta_k,R_k}$ itself. For $h=\left(f^{(\eta_k)}\right)_{\eta_k}$, by~\eqref{eq:continuum-infinite-2},
\[
E_k\left[X_k\left(f^{(\eta_k)}\right)^2\right]=\left|\left(f^{(\eta_k)}\right)_{\eta_k}\right|_{\eta_k}^2\qquad(k\ge k_0).
\]
By Theorem~\ref{thm:estimates-fixed-cutoff}, in the form of Proposition~\ref{prop:sub-gaussian-random}(2), and by Lemma~\ref{lem:norm-bounds-lattice},
\[
E_k\left[X_k\left(f-f^{(\eta_k)}\right)^2\right]\le\left|\left(f-f^{(\eta_k)}\right)_{\eta_k}\right|_{\eta_k}^2\le\left\|f-f^{(\eta_k)}\right\|_{L^2}^2\qquad(k\ge k_0).
\]

For a two-form $h$ on $\mathbb R^3$ or on $\mathbb T_R$ with square integrable components and first derivatives put
\[
\|h\|^2=\sum_{\mu<\nu}\int\mathrm{d}y\,h_{\mu\nu}(y)^2,\qquad\|\nabla h\|^2=\sum_{\mu<\nu}\sum_{j=1}^3\int\mathrm{d}y\,\left(\partial_jh_{\mu\nu}(y)\right)^2,
\]
the integrals being taken over $\mathbb R^3$ or over $\mathbb T_R$; the two cases are indicated by the subscripts $L^2$ and $L^2(\mathbb T_R)$. Let $(N,T)$ be a cutoff, and let $g\in C_c^\infty(\mathbb R^3;\Lambda^2)$ be real, with $\operatorname{supp}g\subset y_0+\left(-\frac{R}{2},\frac{R}{2}\right)^3$ for some $y_0\in\mathbb R^3$. Let $\bar g$ be the two-form on $\mathbb T_R$ whose component $\bar g_{\mu\nu}$, $\mu<\nu$, takes the value $g_\eta(x;\mu\nu)$ on the cube $Q(x;\mu\nu)$, for every $x\in\Lambda$; since for fixed $\mu<\nu$ these cubes tile $\mathbb T_R$ (Section~\ref{sec:tori-cells-cochains}), $\bar g$ is defined almost everywhere. Then
\[
\|\bar g\|_{L^2(\mathbb T_R)}=|g_\eta|_\eta,\qquad\left\|g^R\right\|_{L^2(\mathbb T_R)}=\|g\|_{L^2},\qquad\left\|g^R-\bar g\right\|_{L^2(\mathbb T_R)}\le\frac{\eta}{\pi}\cdot\|\nabla g\|_{L^2}\le\frac{\sqrt3\cdot\eta}{\pi}\cdot\|\nabla g\|_{L^2}.
\]
Indeed, each cube $Q(x;\mu\nu)$ has the volume $\eta^3$, so that $\|\bar g\|_{L^2(\mathbb T_R)}^2=\eta^3\sum_{\mu<\nu}\sum_{x\in\Lambda}g_\eta(x;\mu\nu)^2=|g_\eta|_\eta^2$. As in the proof of Lemma~\ref{lem:norm-bounds-lattice}, the sets $\operatorname{supp}g+R\cdot k$, $k\in\mathbb Z^3$, are pairwise disjoint; hence at every point at most one term of the series $g^R=\sum_{k\in\mathbb Z^3}g(\cdot+R\cdot k)$ is different from zero, and the same holds for the series $\partial_jg^R=\sum_{k\in\mathbb Z^3}(\partial_jg)(\cdot+R\cdot k)$, obtained by termwise differentiation of the locally finite sum. Integration over the fundamental domain $Q_R$ gives $\left\|g^R\right\|_{L^2(\mathbb T_R)}=\|g\|_{L^2}$ and $\left\|\nabla g^R\right\|_{L^2(\mathbb T_R)}=\|\nabla g\|_{L^2}$. Let $\sigma=(x;\mu\nu)$ be a plaquette. By~\eqref{eq:continuum-infinite-2}, the value of $\bar g_{\mu\nu}$ on $Q(\sigma)$ is the mean value of $g^R_{\mu\nu}$ over $Q(\sigma)$. Let $\tilde Q=\tilde y+\left[-\frac{\eta}{2},\frac{\eta}{2}\right)^3\subset\mathbb R^3$ be a cube that the quotient map $\mathbb R^3\to\mathbb T_R$ carries bijectively onto $Q(\sigma)$; as $\eta<R$, this map preserves the Lebesgue measure of $\tilde Q$. The Poincar\'e inequality on a cube, proved above, applied to the smooth function $g^R_{\mu\nu}$ on $\mathbb R^3$ and to the closure of $\tilde Q$, which differs from $\tilde Q$ by a null set, gives
\[
\int_{Q(\sigma)}\mathrm{d}y\,\left(g^R_{\mu\nu}(y)-\bar g_{\mu\nu}(y)\right)^2\le\frac{\eta^2}{\pi^2}\int_{Q(\sigma)}\mathrm{d}y\,\left|\nabla g^R_{\mu\nu}(y)\right|^2 .
\]
Summation over $x\in\Lambda$ and over $\mu<\nu$ yields $\left\|g^R-\bar g\right\|_{L^2(\mathbb T_R)}\le\frac{\eta}{\pi}\cdot\left\|\nabla g^R\right\|_{L^2(\mathbb T_R)}=\frac{\eta}{\pi}\cdot\|\nabla g\|_{L^2}$, and $\frac{\eta}{\pi}\le\frac{\sqrt3\cdot\eta}{\pi}$. By the triangle inequality in $L^2(\mathbb T_R)$ it follows that
\[
\left||g_\eta|_\eta-\|g\|_{L^2}\right|=\left|\|\bar g\|_{L^2(\mathbb T_R)}-\left\|g^R\right\|_{L^2(\mathbb T_R)}\right|\le\left\|\bar g-g^R\right\|_{L^2(\mathbb T_R)}\le\frac{\sqrt3\cdot\eta}{\pi}\cdot\|\nabla g\|_{L^2}.
\]

For $k\ge k_0$ apply this to $g=f^{(\eta_k)}$ at the cutoff $(N_k,T_k)$. By Plancherel's theorem, with the Fourier transform of Section~\ref{sec:lattices-tests-function}, $\|h\|_{L^2}^2+\|\nabla h\|_{L^2}^2=\|h\|_{H^1}^2$ for every $h\in C_c^\infty(\mathbb R^3;\Lambda^2)$; together with the bound on $\left\|f^{(\eta_k)}-f\right\|_{H^1}$ this gives
\[
\left\|f^{(\eta_k)}-f\right\|_{L^2}\le\eta_k^2\cdot\|\beta\|_{H^5},\qquad\left\|\nabla f^{(\eta_k)}\right\|_{L^2}\le\|\nabla f\|_{L^2}+\eta_k^2\cdot\|\beta\|_{H^5}.
\]
Hence, for $k\ge k_0$,
\[
\left|\left|\left(f^{(\eta_k)}\right)_{\eta_k}\right|_{\eta_k}-\|f\|_{L^2}\right|\le\frac{\sqrt3\cdot\eta_k}{\pi}\cdot\left(\|\nabla f\|_{L^2}+\eta_k^2\cdot\|\beta\|_{H^5}\right)+\eta_k^2\cdot\|\beta\|_{H^5}.
\]
The numbers $\|\nabla f\|_{L^2}$ and $\|\beta\|_{H^5}$ are finite, $\beta$ being smooth with compact support, and $\eta_k\to0$ as $N_k\to\infty$; therefore $\left|\left(f^{(\eta_k)}\right)_{\eta_k}\right|_{\eta_k}\to\|f\|_{L^2}$. Minkowski's inequality in $L^2\left(\mathsf P^{\mathrm{free}}_{\eta_k,R_k}\right)$ and the two relations for second moments obtained above give, for $k\ge k_0$,
\[
\left(E_k\left[X_k(f)^2\right]\right)^{\frac12}\ge\left(E_k\left[X_k\left(f^{(\eta_k)}\right)^2\right]\right)^{\frac12}-\left(E_k\left[X_k\left(f-f^{(\eta_k)}\right)^2\right]\right)^{\frac12}\ge\left|\left(f^{(\eta_k)}\right)_{\eta_k}\right|_{\eta_k}-\left\|f-f^{(\eta_k)}\right\|_{L^2},
\]
and the right side tends to $\|f\|_{L^2}$. By Theorem~\ref{thm:continuum-infinite-1}(2), $E_k\left[X_k(f)^2\right]$ tends to $\int\mathrm{d}\mu_1(F)\,F(f)^2=C(f,f)$; hence $C(f,f)\ge\|f\|_{L^2}^2$. On the other hand, Theorem~\ref{thm:continuum-infinite-1}(4) gives $C(f,f)=\int\mathrm{d}\mu_1\,X(f)^2\le C_0(f,f)\le\|f\|_{L^2}^2$. So $C(f,f)=\|f\|_{L^2}^2$.

\emph{Limit points of~\eqref{eq:continuum-infinite-E2}.} Let $\mu_1$ be a limit point of the family~\eqref{eq:continuum-infinite-E1} along a sequence of cutoffs $(N_k,T_k)$, so that $\mathsf Q^{(1)}_{N_k,T_k}\to\mu_1$ weakly on $H_{-4}$. For $c\ge0$ let $\mu_c$ be the image of $\mu_1$ under the continuous map $S\mapsto\sqrt c\cdot S$ of $H_{-4}$, that is, the law of $\sqrt c\cdot F$ under $\mu_1$. By the continuous mapping theorem, $\mathsf Q^{(c)}_{N_k,T_k}\to\mu_c$ weakly on $H_{-4}$. Hence $\frac12\mathsf Q^{(a)}_{N_k,T_k}+\frac12\mathsf Q^{(b)}_{N_k,T_k}\to\frac12\mu_a+\frac12\mu_b$ along the same sequence, and $\frac12\mu_a+\frac12\mu_b$ is a limit point of~\eqref{eq:continuum-infinite-E2}. Conversely, let $\mu$ be a limit point of~\eqref{eq:continuum-infinite-E2}, along a sequence $(N_k,T_k)$. By Theorem~\ref{thm:second-moments-tightness}(3), applied to the family~\eqref{eq:continuum-infinite-E1}, there is a subsequence $(N_{k_l},T_{k_l})_{l\in\mathbb N}$ along which $\mathsf Q^{(1)}_{N_{k_l},T_{k_l}}$ converges weakly on $H_{-4}$ to a measure $\mu_1$; since $N_{k_l}\to\infty$ and $T_{k_l}\to\infty$, $\mu_1$ is a limit point of~\eqref{eq:continuum-infinite-E1}. By what precedes, the laws $\frac12\mathsf Q^{(a)}_{N_{k_l},T_{k_l}}+\frac12\mathsf Q^{(b)}_{N_{k_l},T_{k_l}}$ converge weakly to $\frac12\mu_a+\frac12\mu_b$; they also converge to $\mu$, and weak limits on the metric space $H_{-4}$ are unique. Thus $\mu=\frac12\mu_a+\frac12\mu_b$, with a limit point $\mu_1$ of~\eqref{eq:continuum-infinite-E1}.

For real tests $f_1,\dots,f_n$ and $c\ge0$,
\[
\int\mathrm{d}\mu_c(F)\,\prod_{i=1}^nF(f_i)=c^{\frac n2}\cdot\int\mathrm{d}\mu_1(F)\,\prod_{i=1}^nF(f_i),
\]
the integrals converging absolutely by Theorem~\ref{thm:continuum-infinite-1}(1) and Proposition~\ref{prop:sub-gaussian-random}(4). Therefore the Schwinger functions of $\mu=\frac12\mu_a+\frac12\mu_b$ are
\[
S_n(f_1,\dots,f_n)=\frac{a^{\frac n2}+b^{\frac n2}}{2}\cdot\int\mathrm{d}\mu_1(F)\,\prod_{i=1}^nF(f_i)\qquad\left(f_1,\dots,f_n\in\mathcal S(\mathbb R^3;\Lambda^2)\text{ real}\right).
\]
Let $a<b$, let $f=dd^*\beta\ne0$ be as in~(ii), so that $\beta\ne0$, and let $x\in\mathbb R^3$. By~(i) the vector $\left(F(f),F(\tau_xf)\right)$ is centred Gaussian under $\mu_1$, and Wick's rule gives
\[
\int\mathrm{d}\mu_1(F)\,F(f)^2\cdot F(\tau_xf)^2=C(f,f)\cdot C(\tau_xf,\tau_xf)+2C(f,\tau_xf)^2 .
\]
By Theorem~\ref{thm:continuum-infinite-1}(5), $\mu_1$ is invariant under $\tau_x^*$, so that $C(\tau_xf,\tau_xf)=C(f,f)$. It follows that
\[
S_4(f,f,\tau_xf,\tau_xf)=\frac{a^2+b^2}{2}\cdot\left(C(f,f)^2+2C(f,\tau_xf)^2\right),\qquad S_2(f,f)=\frac{a+b}{2}\cdot C(f,f),
\]
and, by the identity $C(f,f)=\|f\|_{L^2}^2$,
\[
S_4(f,f,\tau_xf,\tau_xf)-S_2(f,f)^2\ge\left(\frac{a^2+b^2}{2}-\frac{(a+b)^2}{4}\right)\cdot C(f,f)^2=\frac{(b-a)^2}{4}\cdot\|f\|_{L^2}^4 .
\]
As $f$ is continuous and not identically zero, $\|f\|_{L^2}>0$, and the right side is positive because $a<b$. Finally, $S_2(\tau_xf,\tau_xf)=S_2(f,f)$ by Theorem~\ref{thm:continuum-infinite-1}(5). So the difference $S_4(f,f,\tau_xf,\tau_xf)-S_2(f,f)\cdot S_2(\tau_xf,\tau_xf)$ is at least the positive number $\frac{(b-a)^2}{4}\cdot\|f\|_{L^2}^4$ for every $x\in\mathbb R^3$, and it does not tend to zero as $|x|\to\infty$. This proves~(ii).

\emph{Two limit points of~\eqref{eq:continuum-infinite-E3}.} Let $a<b$. Let $\varphi\in C_c^\infty(\mathbb R^3)$ be real and not identically zero, and let $\beta$ be the real two-form with $\beta_{12}=\varphi$ and $\beta_{13}=\beta_{23}=0$. Then $(d^*\beta)_1=\partial_2\varphi$, $(d^*\beta)_2=-\partial_1\varphi$ and $(d^*\beta)_3=0$, and $f=dd^*\beta$ has the component $f_{12}=-\left(\partial_1^2+\partial_2^2\right)\varphi$, whose Fourier transform is $p\mapsto\left(p_1^2+p_2^2\right)\cdot\hat\varphi(p)$. The function $\hat\varphi$ is continuous and not identically zero; it is therefore different from zero on a non-empty open set, which is not contained in the line $\{p_1=p_2=0\}$. Hence $f\ne0$. The sequences $(N,T)=(2k,k)$ and $(N,T)=(2k+1,k)$, $k\in\mathbb N$, consist of cutoffs, and both of their entries tend to infinity. By Theorem~\ref{thm:second-moments-tightness}(3), applied to the family~\eqref{eq:continuum-infinite-E1}, they have subsequences along which $\mathsf Q^{(1)}_{N,T}$ converges weakly on $H_{-4}$, to limit points $\mu_1$ and $\mu'_1$ of~\eqref{eq:continuum-infinite-E1}, respectively. Along the first subsequence $N$ is even, the law of $F^\sharp$ under the member of~\eqref{eq:continuum-infinite-E3} is $\mathsf Q^{(a)}_{N,T}$, and, as shown above, it converges to the law $\mu_a$ of $\sqrt a\cdot F$ under $\mu_1$. Along the second subsequence $N$ is odd, and the law of $F^\sharp$ converges in the same way to the law $\mu'_b$ of $\sqrt b\cdot F$ under $\mu'_1$. Thus $\mu_a$ and $\mu'_b$ are limit points of~\eqref{eq:continuum-infinite-E3}. The identity $C(f,f)=\|f\|_{L^2}^2$, applied to $\mu_1$ and to $\mu'_1$, gives
\[
\int\mathrm{d}\mu_a(F)\,F(f)^2=a\cdot\|f\|_{L^2}^2,\qquad\int\mathrm{d}\mu'_b(F)\,F(f)^2=b\cdot\|f\|_{L^2}^2 .
\]
These two numbers differ, since $a<b$ and $\|f\|_{L^2}>0$. Hence $\mu_a\ne\mu'_b$, which proves~(iii).
\end{proof}

\begin{corollary}[Properties not implied by admissibility]
The defining properties~\eqref{eq:continuum-infinite-Ad1}, \eqref{eq:continuum-infinite-Ad2} and~\eqref{eq:continuum-infinite-Ad3} of an admissible family do not imply that the family has only one limit point, nor that its limit points are not Gaussian, nor that they differ from the point mass at the zero distribution, nor that $S_4(f,f,\tau_xf,\tau_xf)-S_2(f,f)^2$ tends to zero as $|x|\to\infty$ for every real $f\in C_c^\infty(\mathbb R^3;\Lambda^2)$. More precisely, there exist admissible families with each of the following properties.
\begin{enumerate}
\item[(1)] The family has two different limit points.
\item[(2)] The family has limit points, and all of them are centred Gaussian measures.
\item[(3)] The family has exactly one limit point, the point mass at the zero distribution $F\equiv0$.
\item[(4)] The family has a limit point $\mu$ such that, for some real $f\in C_c^\infty(\mathbb R^3;\Lambda^2)$, the Schwinger functions $S_n$ of $\mu$ satisfy
\[
\inf_{x\in\mathbb R^3}\left(S_4(f,f,\tau_xf,\tau_xf)-S_2(f,f)^2\right)>0 .
\]
\end{enumerate}
\end{corollary}

\begin{proof}
Let $0\le a<b\le1$, for instance $a=0$ and $b=1$. By Proposition~\ref{prop:scaled-free-laws} the families~\eqref{eq:continuum-infinite-E1}, \eqref{eq:continuum-infinite-E2}, \eqref{eq:continuum-infinite-E3} and $\left(\mathsf P^{(0)}_{\eta,R}\right)$ are admissible. Property~(1) holds for the family~\eqref{eq:continuum-infinite-E3}, by Proposition~\ref{prop:scaled-free-laws}(iii). The family~\eqref{eq:continuum-infinite-E1} has limit points, by Theorem~\ref{thm:second-moments-tightness}(3) applied to the sequence of cutoffs $(k,k)$, $k\in\mathbb N$, and all of them are centred Gaussian measures by Proposition~\ref{prop:scaled-free-laws}(i); this is~(2). Property~(3) holds for the family $\left(\mathsf P^{(0)}_{\eta,R}\right)$, by Proposition~\ref{prop:scaled-free-laws}(i). For~(4), let $\mu$ be a limit point of the family~\eqref{eq:continuum-infinite-E2}, which exists by Theorem~\ref{thm:second-moments-tightness}(3), and let $f=dd^*\beta\ne0$ be the two-form constructed in the proof of Proposition~\ref{prop:scaled-free-laws}(iii); by Proposition~\ref{prop:scaled-free-laws}(ii) the infimum is at least $\frac{(b-a)^2}{4}\cdot\|f\|_{L^2}^4>0$.
\end{proof}

\subsection{Pinned and averaged laws}\label{sec:pinned-averaged-laws}

\subsubsection{The pinned law on a fixed lattice}\label{sec:pinned-law-fixed}

In Section~\ref{sec:pinned-averaged-laws} a parameter assignment is fixed once and for all, and the notation of Section~\ref{sec:existence-limit-points} is used for the admissible family of the averaged laws of this assignment (Corollary~\ref{cor:averaged-laws-admissible}). Thus, for a cutoff $(N,T)$, $\mathsf P_{N,T}$ is the law of $F^\sharp$ under $\bar{\mathsf P}_{\eta,R;q,\boldsymbol{\mu}}$, where $(q,\boldsymbol{\mu})$ is the value of the parameter assignment at $(N,T)$; and for a sequence $(N_k,T_k)_{k\in\mathbb N}$ of cutoffs, $E_k$ (Definition~\ref{def:limit-points}) and $\bar E_k$ both denote the expectation with respect to the averaged law at the cutoff $(N_k,T_k)$.

Let $(N,T)$ be a cutoff, and let $q\ge0$ and $\boldsymbol{\mu}>-1$. By Section~\ref{sec:wilson-operator-lattice} the number $Z_0$ is finite, since the determinant is bounded and $d$ is injective on $V_{\mathrm{t}}$, and it is positive, since $\boldsymbol{\mu}>-6$; hence the pinned law $\mathsf P^0=\mathsf P^0_{\eta,R;q,\boldsymbol{\mu}}$ is defined. As $F=\eta^{-\frac32}\cdot dA$ lies in $\operatorname{ran}\,d$ for every $A\in V_{\mathrm{t}}$, the law $\mathsf P^0$ is a Borel probability measure on $\operatorname{ran}\,d$. It is invariant under the lattice translations and under the cubic group:
\[
\mathsf P^0\circ v_*^{-1}=\mathsf P^0\quad\text{and}\quad\mathsf P^0\circ r_*^{-1}=\mathsf P^0\qquad\text{for all }v\in\eta\mathbb Z^3\text{ and }r\in O_{\mathrm{cub}} ,
\]
where $v_*$ and $r_*$ are the maps of Lemma~\ref{lem:lattice-symmetries-cell}.

To prove the invariance, let $g_*$ be either $v_*$ for some $v\in\eta\mathbb Z^3$ or $r_*$ for some $r\in O_{\mathrm{cub}}$, acting on $C^0$, $C^1$ and $C^2$, and let $g$ be the isometry of $\Lambda$ induced by $y\mapsto y+v$, respectively by $y\mapsto ry$. In unit coordinates $g$ is the map $x\mapsto x+\eta^{-1}\cdot v$, respectively the map $x\mapsto rx$; in both cases it has the form $x\mapsto r'x+c$ with $r'\in O_{\mathrm{cub}}$, $c\in\mathbb Z^3$ and $r'(n\mathbb Z^3)=n\mathbb Z^3$, as in Theorem~\ref{thm:free-coefficient-1}(d4). One-cochains are regarded as odd functions of the oriented bonds, as in Section~\ref{sec:wilson-operator-lattice}, and $g(b)$ denotes the image of an oriented bond $b$ under $g$. Then
\[
(g_*A)\left(g(b)\right)=A(b)\qquad\text{for every $A\in C^1$ and every oriented bond $b$.}
\]
For $v_*$ this is the definition. For $r_*$ it suffices, both sides being odd under the reversal of $b$, to take for $b$ the oriented bond along a link $\sigma$ in the direction of its orientation $o(\sigma)$. Then $g(b)$ is the oriented bond along the link $r\sigma$ in the direction $r\left(o(\sigma)\right)=\varepsilon(r,\sigma)\cdot o(r\sigma)$, with the sign $\varepsilon(r,\sigma)$ of Lemma~\ref{lem:lattice-symmetries-cell}, and the definition of $r_*$ gives
\[
(r_*A)\left(g(b)\right)=\varepsilon(r,\sigma)\cdot(r_*A)(r\sigma)=\varepsilon(r,\sigma)^2\cdot A(\sigma)=A(b).
\]
Theorem~\ref{thm:free-coefficient-1}(d4), applied with $n'=n$ to the one-cochain $q\cdot A$ and to the isometry $g$, therefore gives
\[
\mathrm{Dt}\left(g_*(q\cdot A)\right)=\mathrm{Dt}(q\cdot A)\qquad\left(A\in C^1\right),
\]
where $\mathrm{Dt}$ is the number of \eqref{eq:free-coefficient-5}, formed with the lattice mass $\boldsymbol{\mu}$ on the torus $\Lambda_n$ of Section~\ref{sec:wilson-determinant-lattice} with $n_1=n_2=n_3=n$; in unit coordinates this torus is $\Lambda$. Restricting spinors to the representatives $\{-m,\dots,m\}^3$ identifies the space of the spinors of Section~\ref{sec:wilson-determinant-lattice} with $\ell^2(\Lambda;\mathbb C^4)$, and under this identification $D[q\cdot A]$ is the operator $D_A=D^0_A$ of \eqref{eq:continuum-infinite-1}: both operators attach to an oriented bond $b$ the matrix $\frac12\left(1-\gamma(b)\right)$ and the phase $\exp\left(iq\cdot A(b)\right)$, and the antiperiodicity of the spinors yields the factor $\varsigma(b)$ on the bonds that cross the seam. For $q>0$ this is the identification $D_A=D[q\cdot A]$ of Section~\ref{sec:wilson-determinant-lattice}; for $q=0$ both operators are $D[0]$. Hence $\det\left(D_A+\boldsymbol{\mu}\right)=\mathrm{Dt}(q\cdot A)$ for every $A\in C^1$, and, as $g_*(q\cdot A)=q\cdot g_*A$,
\[
\det\left(D_{g_*A}+\boldsymbol{\mu}\right)=\det\left(D_A+\boldsymbol{\mu}\right)\qquad\left(A\in C^1\right).
\]

By Lemma~\ref{lem:lattice-symmetries-cell}, $g_*$ is orthogonal on $C^1$ and on $C^2$, for $\langle\cdot,\cdot\rangle_\eta$ and hence for the counting pairing, and it commutes with $d$. Hence $g_*$ maps $V_{\mathrm{g}}=d_0(C^0)$ onto itself, since $g_*\circ d_0=d_0\circ g_*$ and $g_*$ is a bijection of $C^0$. It maps $V_{\mathrm{h}}$ onto itself, since it carries each constant directional cochain to plus or minus a constant directional cochain: $v_*$ fixes each of them, and if $re_\nu=\pm e_\lambda$, then $r_*$ carries the constant cochain of the direction $\nu$ to the constant cochain of the direction $\lambda$ multiplied by the same sign. Being orthogonal on $C^1$, the map $g_*$ therefore maps $V_{\mathrm{t}}=\left(V_{\mathrm{g}}\oplus V_{\mathrm{h}}\right)^\perp$ onto itself and preserves the Lebesgue measure $\mathrm{d}A$ of $V_{\mathrm{t}}$. Moreover $\|d(g_*A)\|=\|g_*(dA)\|=\|dA\|$ for $A\in C^1$. Together with the identity of the determinants this gives, for the weight $w_0$ of Section~\ref{sec:wilson-operator-lattice},
\[
w_0(g_*A)=w_0(A)\qquad\left(A\in V_{\mathrm{t}}\right).
\]
Let $\Phi$ be a bounded Borel function on $C^2$. As $\eta^{-\frac32}\cdot d(g_*A)=g_*\left(\eta^{-\frac32}\cdot dA\right)$, the substitution $A\mapsto g_*A$ of $V_{\mathrm{t}}$, which preserves $\mathrm{d}A$ and $w_0$, yields
\begin{align*}
\int\mathrm{d}\mathsf P^0(F)\,\Phi\left(g_*F\right)&=\frac{1}{Z_0}\int_{V_{\mathrm{t}}}\mathrm{d}A\,w_0(A)\cdot\Phi\left(\eta^{-\frac32}\cdot d(g_*A)\right)\\
&=\frac{1}{Z_0}\int_{V_{\mathrm{t}}}\mathrm{d}A\,w_0(A)\cdot\Phi\left(\eta^{-\frac32}\cdot dA\right)=\int\mathrm{d}\mathsf P^0(F)\,\Phi(F).
\end{align*}
Thus $\mathsf P^0\circ g_*^{-1}=\mathsf P^0$. In particular, $\mathsf P^0$ is invariant under $\hat{\mathsf R}_\nu$ for $\nu=1,2,3$, because $\hat{\mathsf R}_\nu=(r_\nu)_*$ on $C^2$ (Lemma~\ref{lem:lattice-symmetries-cell} and its proof).

By Section~\ref{sec:wilson-operator-lattice},
\[
\bar{\mathsf P}=\int_{[-\pi,\pi]^3}\mathrm{d}\omega(\theta)\,\mathsf P^\theta ,
\]
and $\mathsf P^0$ is the fibre over $\theta=0$ of this decomposition of the averaged law. The probability measure $\omega$ has a density with respect to the Lebesgue measure of $[-\pi,\pi]^3$; hence $\omega(\{0\})=0$, and the right side does not change if $\mathsf P^0$ is replaced by any other Borel probability measure on $\operatorname{ran}\,d$. Consequently, a property of $\bar{\mathsf P}$ gives information on $\mathsf P^0$ only by way of a control of how $\mathsf P^\theta$ depends on $\theta$.

\begin{example}[A pinned law without reflection positivity]\label{ex:pinned-law-reflection}
Let $n_1,n_2,n_3$ be odd positive integers, and put $m_\nu=\frac{n_\nu-1}{2}$ for $\nu=1,2,3$. On the rectangular torus $\prod_{\nu=1}^3\mathbb Z/n_\nu\mathbb Z$ the constructions of Sections~\ref{sec:tori-cells-cochains}, \ref{sec:wilson-operator-lattice} and~\ref{sec:admissible-families-lattice} are carried out by the same formulas, with three modifications. The $\nu$-th coordinate of a site is taken in $\{-m_\nu,\dots,m_\nu\}$, and $m_\nu$ is used in place of $m$ for the direction $\nu$. A bond in the direction $\nu$ crosses the seam if one of the $\nu$-th coordinates of its two endpoints equals $m_\nu$ and the other equals $-m_\nu$. For $\theta\in\mathbb R^3$, the factor $\exp\left(\pm\frac{i\theta_\nu}{n}\right)$ by which $D^\theta_A$ multiplies the term of a bond in the direction $\nu$ becomes $\exp\left(\pm\frac{i\theta_\nu}{n_\nu}\right)$, with the same sign. In this example $\Lambda$ denotes the torus
\[
\Lambda=\mathbb Z/1\mathbb Z\times\mathbb Z/1\mathbb Z\times\mathbb Z/3\mathbb Z
\]
with the lattice spacing $\eta=1$; thus $(n_1,n_2,n_3)=(1,1,3)$. The reflection is taken in the direction $\nu=3$, for which $m=m_3=1$. Hence $F=dA$, the pairing $\langle\cdot,\cdot\rangle_\eta$ is the counting pairing, and $\{0,\eta,\dots,m\eta\}=\{0,1\}$ in the definition of $\mathcal P^+_3$. The three sites of $\Lambda$ are identified with their third coordinates $t\in\{-1,0,1\}$; sums and differences of these labels are taken modulo $3$ and represented in $\{-1,0,1\}$, and $(t\to t\pm e_\nu)$ denotes the oriented bonds that start at the site $t$. For $\alpha\in\{1,2\}$ one has $n_\alpha=1$ and $m_\alpha=0$; so both bonds $(t\to t\pm e_\alpha)$ lead from $t$ back to $t$, and both cross the seam. The bond $(t\to t+e_3)$ crosses the seam if and only if $t=1$. For $A\in C^1$ and $t\in\{-1,0,1\}$ put
\[
A_{t,\alpha}=A(t\to t+e_\alpha)\quad(\alpha=1,2),\qquad a_t=A(t\to t+e_3).
\]

One has
\[
V_{\mathrm{t}}=\left\{A\in C^1:\ a_t=0\text{ for }t\in\{-1,0,1\},\ \ \sum_{t=-1}^1A_{t,\alpha}=0\text{ for }\alpha=1,2\right\}.
\]
Indeed, $V_{\mathrm{g}}$ consists of the cochains $d_0\beta$, $\beta\in C^0$, which vanish on the links in the directions $1$ and $2$ and take the values $a_t=\beta(t+1)-\beta(t)$; hence $V_{\mathrm{g}}\oplus V_{\mathrm{h}}$ is the space of the cochains whose values $A_{t,1}$ and $A_{t,2}$ do not depend on $t$, and the space displayed is its orthogonal complement. The formula for $d_1$ in Section~\ref{sec:cochains-periodic-lattice}, together with $t+e_1=t+e_2=t$, gives for every $A\in C^1$
\[
dA(t;12)=0,\qquad dA(t;\alpha3)=-\left(A_{t+1,\alpha}-A_{t,\alpha}\right)\qquad(\alpha=1,2),
\]
the value $a_t$ entering $dA(t;\alpha3)$ twice, with opposite signs. For $A\in V_{\mathrm{t}}$ put $x=(x_1,x_2)$ and $y=(y_1,y_2)$, where
\[
x_\alpha=A_{1,\alpha}-A_{0,\alpha},\qquad y_\alpha=A_{-1,\alpha}-A_{0,\alpha}\qquad(\alpha=1,2).
\]
Since $A_{1,\alpha}+A_{0,\alpha}+A_{-1,\alpha}=0$ on $V_{\mathrm{t}}$, the map $A\mapsto(x,y)$ is a linear bijection of $V_{\mathrm{t}}$ onto $\mathbb R^2\times\mathbb R^2$; its inverse $(x,y)\mapsto A(x,y)$ is given by
\[
A_{0,\alpha}=-\frac{x_\alpha+y_\alpha}{3},\qquad A_{1,\alpha}=x_\alpha+A_{0,\alpha},\qquad A_{-1,\alpha}=y_\alpha+A_{0,\alpha}\qquad(\alpha=1,2).
\]
So $(x,y)$ are linear coordinates on $V_{\mathrm{t}}$. For $A=A(x,y)$ the coboundary $dA$ takes the values $-x_\alpha$, $y_\alpha$ and $x_\alpha-y_\alpha$ on the plaquettes $(0;\alpha3)$, $(-1;\alpha3)$ and $(1;\alpha3)$, and it vanishes on the plaquettes $(t;12)$; consequently
\[
\|dA\|^2=|x|^2+|y|^2+|x-y|^2\qquad\left(A=A(x,y)\in V_{\mathrm{t}}\right).
\]

For $A\in V_{\mathrm{t}}$ the operator $D_A$ of \eqref{eq:continuum-infinite-1} acts on $\ell^2(\Lambda;\mathbb C^4)$, the space of dimension $12$ of the functions $\psi\colon\{-1,0,1\}\to\mathbb C^4$, by
\begin{multline*}
(D_A\psi)(t)=\left[3+\sum_{\alpha=1}^2\left(\cos\left(q\cdot A_{t,\alpha}\right)-i\gamma_\alpha\cdot\sin\left(q\cdot A_{t,\alpha}\right)\right)\right]\psi(t)\\
-\frac12\left(1-\gamma_3\right)\cdot\varsigma_t\cdot\psi(t+1)-\frac12\left(1+\gamma_3\right)\cdot\varsigma_{t-1}\cdot\psi(t-1),
\end{multline*}
where $\varsigma_1=-1$ and $\varsigma_{-1}=\varsigma_0=1$, the labels being taken modulo $3$. Indeed, the bonds $(t\to t+e_\alpha)$ and $(t\to t-e_\alpha)$ carry the values $A_{t,\alpha}$ and $-A_{t,\alpha}$ of $A$, the matrices $\gamma_\alpha$ and $-\gamma_\alpha$, and the sign $-1$; so their two terms in \eqref{eq:continuum-infinite-1} add up to
\begin{multline*}
\frac12\left(1-\gamma_\alpha\right)\cdot\exp\left(iq\cdot A_{t,\alpha}\right)\cdot\psi(t)+\frac12\left(1+\gamma_\alpha\right)\cdot\exp\left(-iq\cdot A_{t,\alpha}\right)\cdot\psi(t)\\
=\left(\cos\left(q\cdot A_{t,\alpha}\right)-i\gamma_\alpha\cdot\sin\left(q\cdot A_{t,\alpha}\right)\right)\psi(t).
\end{multline*}
In the direction $3$ the cochain $A$ vanishes; the bond $(t\to t+e_3)$ carries the matrix $\gamma_3$ and the sign $\varsigma_t$, and the bond $(t\to t-e_3)$ carries the matrix $-\gamma_3$ and the sign $\varsigma_{t-1}$ of the bond between $t-1$ and $t$.

Let $q\ge0$ and $\boldsymbol{\mu}>-1$, and let $w_0$, $Z_0$, $\mathsf P^0$, $E^0$, $\bar w$, $\bar Z$ and $\bar{\mathsf P}$ be the objects of Section~\ref{sec:wilson-operator-lattice} for the torus $\Lambda$, formed with $q$ and $\boldsymbol{\mu}$. As in Section~\ref{sec:wilson-operator-lattice}, the weights are continuous and non-negative, and $0<Z_0<\infty$ and $0<\bar Z<\infty$: the integrals are finite because the determinants are bounded and the quadratic form $\|dA\|^2$ is positive definite on $V_{\mathrm{t}}$ by the formula above, and the argument given there for $Z_0>0$ applies because $n_1,n_2,n_3$ are odd and $\boldsymbol{\mu}>-6$. Let $J>0$ be the absolute value of the determinant of the map $(x,y)\mapsto A(x,y)$ with respect to an orthonormal basis of $V_{\mathrm{t}}$, so that $\int_{V_{\mathrm{t}}}\mathrm{d}A\,f(A)=J\cdot\int_{\mathbb R^2\times\mathbb R^2}\mathrm{d}x\,\mathrm{d}y\,f\left(A(x,y)\right)$ for every integrable function $f$ on $V_{\mathrm{t}}$. If $A$ is distributed on $V_{\mathrm{t}}$ with the density $Z_0^{-1}\cdot w_0$, as in the definition of $\mathsf P^0$, the pair $(x,y)$ has the continuous density
\[
\rho^0(x,y)=\frac{J}{Z_0}\cdot\exp\left(-\frac12\left(|x|^2+|y|^2+|x-y|^2\right)\right)\cdot\det\left(D_{A(x,y)}+\boldsymbol{\mu}\right)
\]
with respect to the Lebesgue measure of $\mathbb R^2\times\mathbb R^2$. As $x_\alpha=-F(0;\alpha3)$ and $y_\alpha=F(-1;\alpha3)$, the pair $(x,y)$ is a function of $F$, and $E^0\Phi(x,y)=\int\mathrm{d}x\,\mathrm{d}y\,\Phi(x,y)\cdot\rho^0(x,y)$ for every bounded Borel function $\Phi$ on $\mathbb R^2\times\mathbb R^2$.

The density $\rho^0$ is symmetric: $\rho^0(y,x)=\rho^0(x,y)$ for all $x,y\in\mathbb R^2$. To see this, let $\mathsf V$ be the unitary operator on $\ell^2(\Lambda;\mathbb C^4)$ with $(\mathsf V\psi)(-t)=\gamma_3\gamma_4\psi(t)$. The unitary matrix $\gamma_3\gamma_4$ commutes with $\gamma_1$ and $\gamma_2$ and anticommutes with $\gamma_3$, because $\gamma_4$ anticommutes with $\gamma_1,\gamma_2,\gamma_3$, the matrix $\gamma_3$ anticommutes with $\gamma_1$ and $\gamma_2$, and $\gamma_3^2=1$. Hence conjugation by $\gamma_3\gamma_4$ fixes $\gamma_1$ and $\gamma_2$ and interchanges $\frac12\left(1-\gamma_3\right)$ and $\frac12\left(1+\gamma_3\right)$. Moreover $\varsigma_{-t-1}=\varsigma_t$ and $\varsigma_{-t}=\varsigma_{t-1}$ for $t\in\{-1,0,1\}$. For $A\in V_{\mathrm{t}}$ the formula for $D_A$ therefore gives
\begin{multline*}
\left(\mathsf VD_A\mathsf V^{-1}\psi\right)(t)=\left[3+\sum_{\alpha=1}^2\left(\cos\left(q\cdot A_{-t,\alpha}\right)-i\gamma_\alpha\cdot\sin\left(q\cdot A_{-t,\alpha}\right)\right)\right]\psi(t)\\
-\frac12\left(1+\gamma_3\right)\cdot\varsigma_{-t}\cdot\psi(t-1)-\frac12\left(1-\gamma_3\right)\cdot\varsigma_{-t-1}\cdot\psi(t+1)=\left(D_{A'}\psi\right)(t),
\end{multline*}
where $A'\in V_{\mathrm{t}}$ is the cochain with $A'_{t,\alpha}=A_{-t,\alpha}$ for $t\in\{-1,0,1\}$ and $\alpha=1,2$. By the definition of the coordinates, the cochain $A'$ attached to $A=A(x,y)$ is $A(y,x)$. Consequently
\[
\det\left(D_{A(y,x)}+\boldsymbol{\mu}\right)=\det\left(\mathsf V\left(D_{A(x,y)}+\boldsymbol{\mu}\right)\mathsf V^{-1}\right)=\det\left(D_{A(x,y)}+\boldsymbol{\mu}\right),
\]
and the other factors of $\rho^0(x,y)$ are symmetric in $x$ and $y$.

The set $\mathcal P^+_3$ consists of the plaquettes $(0;12)$, $(1;12)$, $(0;13)$ and $(0;23)$, because the vertices of $(t;12)$ have the third coordinate $t$ and those of $(t;\alpha3)$ have the third coordinates $t$ and $t+1$. Hence $(0;12)$ lies in the plane of the points with third coordinate $0$, and $(1;12)$ in the plane of the points with third coordinate $m\eta=1$; these are the two planes excluded in the definition of $\mathcal P^{+\circ}_3$, and therefore $\mathcal P^{+\circ}_3=\left\{(0;13),(0;23)\right\}$. On $(0;\alpha3)$ one has $F=-x_\alpha$, and $F$ vanishes on $(0;12)$ and on $(1;12)$. The reflection $r_3$ maps $(0;\alpha3)$ onto $(-1;\alpha3)$ and $(t;12)$ onto $(-t;12)$; therefore $\left(\hat{\mathsf R}_3F\right)(0;\alpha3)=-F(-1;\alpha3)=-y_\alpha$, and $\hat{\mathsf R}_3F$ vanishes on $(0;12)$ and on $(1;12)$. In this sense $\hat{\mathsf R}_3$ exchanges $x$ and $y$: if $g\in C^2(\Lambda)$ is real and supported on $\mathcal P^+_3$, and $k=\left(-g(0;13),-g(0;23)\right)\in\mathbb R^2$, then
\[
\langle g,F\rangle_\eta=k\cdot x,\qquad\left\langle\hat{\mathsf R}_3g,F\right\rangle_\eta=\left\langle g,\hat{\mathsf R}_3F\right\rangle_\eta=k\cdot y,
\]
since $\hat{\mathsf R}_3$ is an orthogonal involution, hence symmetric. Every $k\in\mathbb R^2$ arises in this way from a cochain $g$ supported on $\mathcal P^{+\circ}_3$.

The inequality of Proposition~\ref{prop:reflection-positivity} for the pinned law, in the direction $\nu=3$, is the assertion that
\[
\sum_{a,b=1}^s\bar c_ac_b\,E^0\exp\left(i\left\langle g_b-\hat{\mathsf R}_3g_a,F\right\rangle_\eta\right)\ge0
\]
for every integer $s\ge1$, all real $g_1,\dots,g_s\in C^2(\Lambda)$ supported on $\mathcal P^+_3$ and every $c=(c_1,\dots,c_s)\in\mathbb C^s$. For a finite sum
\[
P(z)=\sum_{b=1}^sc_b\cdot\exp\left(ik_b\cdot z\right)\qquad\left(z\in\mathbb R^2\right),
\]
with $s\ge1$, $c\in\mathbb C^s$ and $k_1,\dots,k_s\in\mathbb R^2$, put
\[
\mathsf B(P)=\int_{\mathbb R^2\times\mathbb R^2}\mathrm{d}x\,\mathrm{d}y\,\overline{P(y)}\cdot P(x)\cdot\rho^0(x,y).
\]
If $k_b$ is attached to $g_b$ as above, for $b=1,\dots,s$, the left side of the inequality equals
\[
\sum_{a,b=1}^s\bar c_ac_b\,E^0\exp\left(ik_b\cdot x-ik_a\cdot y\right)=E^0\left[\overline{P(y)}\cdot P(x)\right]=\mathsf B(P).
\]
Hence the inequality of Proposition~\ref{prop:reflection-positivity} for the pinned law, in the direction $3$, holds if and only if $\mathsf B(P)\ge0$ for all finite sums $P$ of this form.

Suppose that it holds. Let $\lambda$ be the law of $x$ under $\mathsf P^0$, a probability measure on $\mathbb R^2$; by the symmetry of $\rho^0$, the law of $y$ is $\lambda$ as well. For $P,Q\in L^2(\mathbb R^2,\lambda)$ the Cauchy--Schwarz inequality gives
\begin{align*}
\int_{\mathbb R^2\times\mathbb R^2}\mathrm{d}x\,\mathrm{d}y\,|P(y)|\cdot|Q(x)|\cdot\rho^0(x,y)&\le\left(E^0|P(y)|^2\right)^{\frac12}\cdot\left(E^0|Q(x)|^2\right)^{\frac12}\\
&=\|P\|_{L^2(\lambda)}\cdot\|Q\|_{L^2(\lambda)}.
\end{align*}
So the integral that defines $\mathsf B(P)$ converges absolutely for every $P\in L^2(\mathbb R^2,\lambda)$, and the identity $\overline{P(y)}P(x)-\overline{Q(y)}Q(x)=\overline{(P-Q)(y)}\cdot P(x)+\overline{Q(y)}\cdot(P-Q)(x)$ gives
\[
\left|\mathsf B(P)-\mathsf B(Q)\right|\le\|P-Q\|_{L^2(\lambda)}\cdot\left(\|P\|_{L^2(\lambda)}+\|Q\|_{L^2(\lambda)}\right)\qquad\left(P,Q\in L^2(\mathbb R^2,\lambda)\right);
\]
thus $\mathsf B$ is continuous on $L^2(\mathbb R^2,\lambda)$. The finite sums $P$ are dense in $L^2(\mathbb R^2,\lambda)$: if $u\in L^2(\mathbb R^2,\lambda)$ is orthogonal to all of them, the finite complex measure $u\cdot\lambda$ has a vanishing Fourier transform, and $u=0$ $\lambda$-almost everywhere by the uniqueness theorem for Fourier transforms of finite complex measures on $\mathbb R^2$, as in the proof of Lemma~\ref{lem:totality-exponentials}. Consequently $\mathsf B(P)\ge0$ for every $P\in L^2(\mathbb R^2,\lambda)$, in particular for every bounded Borel function $P\colon\mathbb R^2\to\mathbb C$.

Let $\xi,\zeta\in\mathbb R^2$ and $c=(c_1,c_2)\in\mathbb C^2$, put $p_1=\xi$ and $p_2=\zeta$, and for $\varepsilon>0$ let $\varphi_\varepsilon\colon\mathbb R^2\to[0,\infty)$ be the function equal to $\left(\pi\varepsilon^2\right)^{-1}$ on the open disc of radius $\varepsilon$ about $0$ and equal to $0$ outside it. The bounded Borel function $P_\varepsilon(z)=c_1\cdot\varphi_\varepsilon(z-\xi)+c_2\cdot\varphi_\varepsilon(z-\zeta)$ satisfies
\[
0\le\mathsf B(P_\varepsilon)=\sum_{a,b=1}^2\bar c_ac_b\int_{\mathbb R^2\times\mathbb R^2}\mathrm{d}x\,\mathrm{d}y\,\varphi_\varepsilon(y-p_a)\cdot\varphi_\varepsilon(x-p_b)\cdot\rho^0(x,y).
\]
For each pair $(a,b)$ the measure $\varphi_\varepsilon(y-p_a)\cdot\varphi_\varepsilon(x-p_b)\,\mathrm{d}x\,\mathrm{d}y$ is a probability measure carried by the set of the points $(x,y)$ with $|x-p_b|<\varepsilon$ and $|y-p_a|<\varepsilon$; as $\rho^0$ is continuous, the corresponding integral tends to $\rho^0(p_b,p_a)$ as $\varepsilon\to0$. Hence
\[
\sum_{a,b=1}^2\bar c_ac_b\,\rho^0(p_a,p_b)=\sum_{a,b=1}^2\bar c_ac_b\,\rho^0(p_b,p_a)\ge0\qquad\left(c\in\mathbb C^2\right),
\]
the equality by the symmetry of $\rho^0$. The real symmetric matrix $\left(\rho^0(p_a,p_b)\right)_{a,b=1,2}$ is therefore positive semi-definite, and its determinant $\rho^0(\xi,\xi)\cdot\rho^0(\zeta,\zeta)-\rho^0(\xi,\zeta)^2$ is non-negative. As $\xi$ and $\zeta$ are arbitrary, the inequality of Proposition~\ref{prop:reflection-positivity} for the pinned law, in the direction $3$, implies
\[
\rho^0(x,y)^2\le\rho^0(x,x)\cdot\rho^0(y,y)\qquad\text{for all }x,y\in\mathbb R^2.
\]

Consider the following two choices of the parameters and of the points; in both of them $x=\frac{2\pi}{q}\cdot(1,1)$ and $y=\frac{4\pi}{q}\cdot(1,1)$:
\begin{enumerate}
\item[(i)] $q=6$, $\boldsymbol{\mu}=\frac12$, $x=\left(\frac\pi3,\frac\pi3\right)$ and $y=\left(\frac{2\pi}3,\frac{2\pi}3\right)$;
\item[(ii)] $q=4$, $\boldsymbol{\mu}=-0.45$, $x=\left(\frac\pi2,\frac\pi2\right)$ and $y=(\pi,\pi)$.
\end{enumerate}
For all $x,y\in\mathbb R^2$ with $\rho^0(x,x)\cdot\rho^0(y,y)\ne0$ the formula for $\rho^0$ gives
\[
\frac{\rho^0(x,y)^2}{\rho^0(x,x)\cdot\rho^0(y,y)}=\exp\left(-|x-y|^2\right)\cdot\frac{\det\left(D_{A(x,y)}+\boldsymbol{\mu}\right)^2}{\det\left(D_{A(x,x)}+\boldsymbol{\mu}\right)\cdot\det\left(D_{A(y,y)}+\boldsymbol{\mu}\right)},
\]
because the factors $\frac{J}{Z_0}$ cancel and the exponents add up to $-\left(|x|^2+|y|^2+|x-y|^2\right)+|x|^2+|y|^2=-|x-y|^2$.

For $A\in V_{\mathrm{t}}$ and $\theta\in\mathbb R^3$ the operator $D^\theta_A$ is obtained from the formula for $D_A$ above when $q\cdot A_{t,\alpha}$ is replaced by $q\cdot A_{t,\alpha}+\theta_\alpha$, for $t\in\{-1,0,1\}$ and $\alpha=1,2$, and the terms with $\psi(t+1)$ and $\psi(t-1)$ are multiplied by $\exp\left(\frac{i\theta_3}{3}\right)$ and by $\exp\left(-\frac{i\theta_3}{3}\right)$, respectively. Indeed, by the third modification, with $n_1=n_2=1$ and $n_3=3$, the terms of the bonds $(t\to t\pm e_\alpha)$ acquire the factors $\exp\left(\pm i\theta_\alpha\right)$, which turn $\exp\left(\pm iq\cdot A_{t,\alpha}\right)$ into $\exp\left(\pm i\left(q\cdot A_{t,\alpha}+\theta_\alpha\right)\right)$, and the terms of the bonds $(t\to t\pm e_3)$ acquire the factors $\exp\left(\pm\frac{i\theta_3}{3}\right)$. Thus $D^\theta_A$ involves $A$, $\theta_1$ and $\theta_2$ only through the numbers $\cos\left(q\cdot A_{t,\alpha}+\theta_\alpha\right)$ and $\sin\left(q\cdot A_{t,\alpha}+\theta_\alpha\right)$.

Suppose that $q>0$. At the points $x=\frac{2\pi}{q}\cdot(1,1)$ and $y=\frac{4\pi}{q}\cdot(1,1)$ the formula for $A(x,y)$ shows that the triple $\left(q\cdot A_{0,\alpha},q\cdot A_{1,\alpha},q\cdot A_{-1,\alpha}\right)$, which is the same for $\alpha=1$ and for $\alpha=2$, equals
\[
\left(-2\pi,0,2\pi\right),\qquad\left(-\frac{4\pi}{3},\frac{2\pi}{3},\frac{2\pi}{3}\right),\qquad\left(-\frac{8\pi}{3},\frac{4\pi}{3},\frac{4\pi}{3}\right)
\]
for $A=A(x,y)$, $A=A(x,x)$ and $A=A(y,y)$, respectively. Modulo $2\pi$, all three entries of the first triple are congruent to $0$, those of the second to $\frac{2\pi}{3}$, and those of the third to $-\frac{2\pi}{3}$. Hence, with $\theta^\ast=\left(\frac{2\pi}{3},\frac{2\pi}{3},0\right)$,
\[
D^\theta_{A(x,y)}=D^\theta_0,\qquad D^\theta_{A(x,x)}=D^{\theta+\theta^\ast}_0,\qquad D^\theta_{A(y,y)}=D^{\theta-\theta^\ast}_0\qquad\left(\theta\in\mathbb R^3\right).
\]

Let $\theta\in\mathbb R^3$, and put $\omega_j=\frac{\pi(2j+1)+\theta_3}{3}$ for $j\in\{0,1,2\}$. For $j\in\{0,1,2\}$ and $\chi\in\mathbb C^4$ let $\psi_{j,\chi}\in\ell^2(\Lambda;\mathbb C^4)$ be the function $t\mapsto\exp\left(\frac{i\pi(2j+1)t}{3}\right)\cdot\chi$. For $\psi=\psi_{j,\chi}$ and every $t\in\{-1,0,1\}$ one has
\[
\varsigma_t\cdot\psi(t+1)=\exp\left(\frac{i\pi(2j+1)}{3}\right)\cdot\psi(t),\qquad\varsigma_{t-1}\cdot\psi(t-1)=\exp\left(-\frac{i\pi(2j+1)}{3}\right)\cdot\psi(t):
\]
for $t\in\{-1,0\}$ in the first identity and for $t\in\{0,1\}$ in the second, the sign $\varsigma_t$, respectively $\varsigma_{t-1}$, equals $1$ and the identity is the definition of $\psi$, while in the remaining two cases, $t=1$ in the first identity and $t=-1$ in the second, this sign is $\varsigma_1=-1$, the site $t\pm1$ is $-t$, and the identity follows from $\exp\left(i\pi(2j+1)\right)=-1$. The three functions $t\mapsto\exp\left(\frac{i\pi(2j+1)t}{3}\right)$ on $\{-1,0,1\}$, $j=0,1,2$, are pairwise orthogonal, since $\sum_{t=-1}^1\exp\left(\frac{2\pi i(j-k)t}{3}\right)=0$ for distinct $j,k\in\{0,1,2\}$; hence $\ell^2(\Lambda;\mathbb C^4)$ is the orthogonal direct sum of the three subspaces $\left\{\psi_{j,\chi}:\chi\in\mathbb C^4\right\}$, each of dimension $4$. By the description of $D^\theta_0$ above and the two identities, $D^\theta_0+\boldsymbol{\mu}$ maps $\psi_{j,\chi}$ to $\psi_{j,M_j\chi}$, where
\[
M_j=\boldsymbol{\mu}+3+\sum_{\alpha=1}^2\left(\cos\theta_\alpha-i\gamma_\alpha\cdot\sin\theta_\alpha\right)-\cos\omega_j+i\gamma_3\cdot\sin\omega_j,
\]
because $\frac12\left(1-\gamma_3\right)\cdot\exp\left(i\omega_j\right)+\frac12\left(1+\gamma_3\right)\cdot\exp\left(-i\omega_j\right)=\cos\omega_j-i\gamma_3\cdot\sin\omega_j$. Consequently $\det\left(D^\theta_0+\boldsymbol{\mu}\right)=\det M_0\cdot\det M_1\cdot\det M_2$. One has $M_j=v_0+i\sum_{\nu=1}^3v_\nu\gamma_\nu$ with the real numbers $v_0=\boldsymbol{\mu}+3+\cos\theta_1+\cos\theta_2-\cos\omega_j$, $v_1=-\sin\theta_1$, $v_2=-\sin\theta_2$ and $v_3=\sin\omega_j$. The matrices $\gamma_1,\gamma_2,\gamma_3$ are Hermitian, have trace $0$ and square $1$, and anticommute in pairs; therefore the Hermitian matrix $\sum_\nu v_\nu\gamma_\nu$ has trace $0$ and square $|v|^2=v_1^2+v_2^2+v_3^2$, its eigenvalues, counted with multiplicity, are $|v|,|v|,-|v|,-|v|$, and $\det M_j=\left(v_0^2+|v|^2\right)^2$. It follows that
\[
\det\left(D^\theta_0+\boldsymbol{\mu}\right)=\prod_{j=0}^2\left[\left(\boldsymbol{\mu}+3+\cos\theta_1+\cos\theta_2-\cos\omega_j\right)^2+\sin^2\theta_1+\sin^2\theta_2+\sin^2\omega_j\right]^2\qquad\left(\theta\in\mathbb R^3\right);
\]
with $n_\nu$ in place of $n$, this is the formula for $w_\theta(0)=\det\left(D^\theta_0+\boldsymbol{\mu}\right)$ in Section~\ref{sec:wilson-operator-lattice}, the three points being $\left(\pi+\theta_1,\pi+\theta_2,\omega_j\right)$, $j=0,1,2$, since $1-\cos\left(\pi+\theta_\alpha\right)=1+\cos\theta_\alpha$ and $\sin^2\left(\pi+\theta_\alpha\right)=\sin^2\theta_\alpha$.

For $\theta=0$ one has $\cos\theta_\alpha=1$ and $\sin\theta_\alpha=0$, and $\left(\omega_0,\omega_1,\omega_2\right)=\left(\frac\pi3,\pi,\frac{5\pi}3\right)$, so that $\cos\omega_j$ takes the values $\frac12,-1,\frac12$ and $\sin^2\omega_j$ the values $\frac34,0,\frac34$. For $\theta=\theta^\ast$ and for $\theta=-\theta^\ast$ the numbers $\omega_j$ are the same, while $\cos\theta_\alpha=-\frac12$ and $\sin^2\theta_\alpha=\frac34$ for $\alpha=1,2$. Hence
\[
\det\left(D_0+\boldsymbol{\mu}\right)=\left(\boldsymbol{\mu}+6\right)^4\cdot\left(\left(\boldsymbol{\mu}+\frac92\right)^2+\frac34\right)^4,
\]
\[
\det\left(D^{\theta^\ast}_0+\boldsymbol{\mu}\right)=\det\left(D^{-\theta^\ast}_0+\boldsymbol{\mu}\right)=\left(\left(\boldsymbol{\mu}+3\right)^2+\frac32\right)^2\cdot\left(\left(\boldsymbol{\mu}+\frac32\right)^2+\frac94\right)^4>0.
\]
At the points $x=\frac{2\pi}{q}\cdot(1,1)$ and $y=\frac{4\pi}{q}\cdot(1,1)$ one has $|x-y|^2=\frac{8\pi^2}{q^2}$; by the identities $D_{A(x,y)}=D_0$, $D_{A(x,x)}=D^{\theta^\ast}_0$ and $D_{A(y,y)}=D^{-\theta^\ast}_0$, the case $\theta=0$ of those above, the denominator of the quotient is positive and
\[
\frac{\rho^0(x,y)^2}{\rho^0(x,x)\cdot\rho^0(y,y)}=\exp\left(-\frac{8\pi^2}{q^2}\right)\cdot K(\boldsymbol{\mu})^2,\qquad K(\boldsymbol{\mu})=\frac{\det\left(D_0+\boldsymbol{\mu}\right)}{\det\left(D^{\theta^\ast}_0+\boldsymbol{\mu}\right)}.
\]
For $\boldsymbol{\mu}>-1$ one has $\left(\boldsymbol{\mu}+6\right)^2>\left(\boldsymbol{\mu}+3\right)^2+\frac32$ and $\left(\boldsymbol{\mu}+\frac92\right)^2+\frac34>\left(\boldsymbol{\mu}+\frac32\right)^2+\frac94$, the differences of the two sides being $6\boldsymbol{\mu}+\frac{51}{2}$ and $6\boldsymbol{\mu}+\frac{33}{2}$; comparing the two closed forms factor by factor, one obtains $K(\boldsymbol{\mu})>1$. Consequently the inequality $\rho^0(x,y)^2\le\rho^0(x,x)\cdot\rho^0(y,y)$ fails at the points $x=\frac{2\pi}{q}\cdot(1,1)$ and $y=\frac{4\pi}{q}\cdot(1,1)$ for every $\boldsymbol{\mu}>-1$ and every $q>\frac{2\pi}{\sqrt{\log K(\boldsymbol{\mu})}}$. In the cases (i) and (ii) one has
\[
K\left(\frac12\right)=\frac{13^4\cdot103^4}{55^2\cdot25^4}=2720.42\ldots,\qquad K(-0.45)=\frac{111^4\cdot6861^4}{3201^2\cdot1341^4}=10152.08\ldots,
\]
and therefore
\[
\frac{\rho^0(x,y)^2}{\rho^0(x,x)\cdot\rho^0(y,y)}=
\begin{cases}
\exp\left(-\frac{2\pi^2}{9}\right)\cdot K\left(\frac12\right)^2=8.2558\ldots\cdot10^5&\text{in the case (i)},\\
\exp\left(-\frac{\pi^2}{2}\right)\cdot K(-0.45)^2=7.4122\ldots\cdot10^5&\text{in the case (ii)}.
\end{cases}
\]
Here $K\left(\frac12\right)^2>7\cdot10^6$ and $K(-0.45)^2>10^8$, while $\pi^2<10$ and $\exp(1)<3$ give $\exp\left(\frac{2\pi^2}{9}\right)<\exp(3)<27$ and $\exp\left(\frac{\pi^2}{2}\right)<\exp(5)<243$.
Both values are larger than $1$, whereas the inequality of Proposition~\ref{prop:reflection-positivity} for the pinned law, in the direction $3$, would bound them by $1$. For the density $\bar\rho(x,y)=\frac{J}{\bar Z}\cdot\bar w\left(A(x,y)\right)$ of $(x,y)$ under the averaged law $\bar{\mathsf P}$ of the same torus, the corresponding quotient
\[
\frac{\bar\rho(x,y)^2}{\bar\rho(x,x)\cdot\bar\rho(y,y)}
\]
equals $\exp\left(-\frac{8\pi^2}{q^2}\right)<1$ at the points $x=\frac{2\pi}{q}\cdot(1,1)$ and $y=\frac{4\pi}{q}\cdot(1,1)$, for every $q>0$ and every $\boldsymbol{\mu}>-1$; its value is $\exp\left(-\frac{2\pi^2}{9}\right)=0.1115\ldots$ in the case (i) and $\exp\left(-\frac{\pi^2}{2}\right)=7.191\ldots\cdot10^{-3}$ in the case (ii). Indeed, by the definitions of $w_\theta$ and $\bar w$ in Section~\ref{sec:wilson-operator-lattice} and the formula for $\|dA\|^2$, for all $x,y\in\mathbb R^2$
\begin{gather*}
\bar\rho(x,y)=\frac{J}{\bar Z}\cdot\exp\left(-\frac12\left(|x|^2+|y|^2+|x-y|^2\right)\right)\cdot I\left(A(x,y)\right),\\
I(A)=\int_{[-\pi,\pi]^3}\frac{\mathrm{d}\theta}{(2\pi)^3}\,\det\left(D^\theta_A+\boldsymbol{\mu}\right)\qquad\left(A\in V_{\mathrm{t}}\right);
\end{gather*}
hence, as for $\rho^0$,
\[
\frac{\bar\rho(x,y)^2}{\bar\rho(x,x)\cdot\bar\rho(y,y)}=\exp\left(-|x-y|^2\right)\cdot\frac{I\left(A(x,y)\right)^2}{I\left(A(x,x)\right)\cdot I\left(A(y,y)\right)}
\]
whenever $I\left(A(x,x)\right)\cdot I\left(A(y,y)\right)\ne0$. At the points $x=\frac{2\pi}{q}\cdot(1,1)$ and $y=\frac{4\pi}{q}\cdot(1,1)$, the identities $D^\theta_{A(x,y)}=D^\theta_0$, $D^\theta_{A(x,x)}=D^{\theta+\theta^\ast}_0$ and $D^\theta_{A(y,y)}=D^{\theta-\theta^\ast}_0$ show that $I\left(A(x,y)\right)$, $I\left(A(x,x)\right)$ and $I\left(A(y,y)\right)$ are the means over $[-\pi,\pi]^3$ of the functions $\theta\mapsto\det\left(D^\theta_0+\boldsymbol{\mu}\right)$, $\theta\mapsto\det\left(D^{\theta+\theta^\ast}_0+\boldsymbol{\mu}\right)$ and $\theta\mapsto\det\left(D^{\theta-\theta^\ast}_0+\boldsymbol{\mu}\right)$. By the product formula for $\det\left(D^\theta_0+\boldsymbol{\mu}\right)$ above, the function $\theta\mapsto\det\left(D^\theta_0+\boldsymbol{\mu}\right)$ depends on $\theta_1$ and $\theta_2$ only through $\cos\theta_1$, $\sin\theta_1$, $\cos\theta_2$ and $\sin\theta_2$, so that it has period $2\pi$ in each of these two variables; as the third component of $\theta^\ast$ vanishes, the three means coincide, and their common value is $I(0)$. Moreover $I(0)>0$: by the same product formula, the function $\theta\mapsto\det\left(D^\theta_0+\boldsymbol{\mu}\right)$ is continuous and non-negative, and at every $\theta$ with $\sin\theta_1\ne0$ each of its three factors is at least $\sin^4\theta_1>0$. Therefore the quotient equals $\exp\left(-|x-y|^2\right)=\exp\left(-\frac{8\pi^2}{q^2}\right)$.
\end{example}

\subsubsection{A criterion for the convergence of the pinned laws}\label{sec:criterion-convergence-pinned}

Let $(N_k,T_k)_{k\in\mathbb N}$ be a sequence of cutoffs, and let $\eta_k$, $R_k$ and $X_k=X_{\eta_k,R_k}$ be as in Definition~\ref{def:limit-points}. For $k\in\mathbb N$ let $(q_k,\boldsymbol{\mu}_k)$ be the value of the parameter assignment at the cutoff $(N_k,T_k)$. For $\theta\in\mathbb R^3$ let $Z_{\theta,k}$ denote the number $Z_\theta$ of Section~\ref{sec:wilson-operator-lattice}, formed on the lattice $\Lambda_{\eta_k,R_k}$ with $q=q_k$ and $\boldsymbol{\mu}=\boldsymbol{\mu}_k$, and, if $Z_{\theta,k}>0$, put $\mathsf P^\theta_k=\mathsf P^\theta_{\eta_k,R_k;q_k,\boldsymbol{\mu}_k}$. As $\boldsymbol{\mu}_k>-1$, the number $Z_{0,k}$ is positive (Section~\ref{sec:pinned-law-fixed}); the expectation under the pinned law $\mathsf P^0_k$ is denoted by $E^0_k$. Let $\omega_k$ be the probability measure $\omega$ of Section~\ref{sec:wilson-operator-lattice} on $[-\pi,\pi]^3$, formed with the same data. The $\theta$ with $Z_{\theta,k}=0$ form an $\omega_k$-null set, and by Section~\ref{sec:wilson-operator-lattice}
\[
\bar{\mathsf P}_{\eta_k,R_k;q_k,\boldsymbol{\mu}_k}=\int_{[-\pi,\pi]^3}\mathrm{d}\omega_k(\theta)\,\mathsf P^\theta_k ;
\]
as in Section~\ref{sec:pinned-law-fixed}, $\bar E_k$ is the expectation under this averaged law.

Let $s\ge1$ be an integer and let $f_1,\dots,f_s\in C_c^\infty(\mathbb R^3;\Lambda^2)$ be real tests. For $k\in\mathbb N$, $z\in\mathbb C^s$ and $\theta\in\mathbb R^3$ with $Z_{\theta,k}>0$ put
\[
G^\theta_k(z)=\int\mathrm{d}\mathsf P^\theta_k(F)\,\exp\left(\sum_{i=1}^sz_i\cdot X_k(f_i)\right),\qquad\bar G_k(z)=\bar E_k\exp\left(\sum_{i=1}^sz_i\cdot X_k(f_i)\right).
\]
For $\alpha\in\mathbb N_0^s$ let $|\alpha|=\alpha_1+\dots+\alpha_s$, $\alpha!=\alpha_1!\cdots\alpha_s!$, $z^\alpha=z_1^{\alpha_1}\cdots z_s^{\alpha_s}$ and $\partial^\alpha=\partial_{z_1}^{\alpha_1}\cdots\partial_{z_s}^{\alpha_s}$. The functions $G^\theta_k$ and $\bar G_k$ are entire on $\mathbb C^s$, their derivatives at $0$ are the moments
\[
\partial^\alpha G^\theta_k(0)=\int\mathrm{d}\mathsf P^\theta_k(F)\,\prod_{i=1}^sX_k(f_i)^{\alpha_i},\qquad\partial^\alpha\bar G_k(0)=\bar E_k\prod_{i=1}^sX_k(f_i)^{\alpha_i}\qquad\left(\alpha\in\mathbb N_0^s\right),
\]
and
\[
\bar G_k(z)=\int_{[-\pi,\pi]^3}\mathrm{d}\omega_k(\theta)\,G^\theta_k(z)\qquad\left(z\in\mathbb C^s\right).
\]
To prove this, fix $k$ and write $\eta=\eta_k$, $R=R_k$, $n=R/\eta$, $q=q_k$ and $\boldsymbol{\mu}=\boldsymbol{\mu}_k$. By Section~\ref{sec:wilson-operator-lattice} one has $\det\left(D^\theta_A+\boldsymbol{\mu}\right)\le\left(9+|\boldsymbol{\mu}|\right)^{4n^3}$ for all $\theta\in\mathbb R^3$, so that the weights of that section satisfy
\[
\max\left(w_\theta(A),\bar w(A)\right)\le\left(9+|\boldsymbol{\mu}|\right)^{4n^3}\cdot\exp\left(-\frac12\|dA\|^2\right)\qquad\left(\theta\in\mathbb R^3,\ A\in V_{\mathrm{t}}\right).
\]
Hence each of the laws $\mathsf P=\mathsf P^\theta_k$, $Z_{\theta,k}>0$, and $\mathsf P=\bar{\mathsf P}_{\eta,R;q,\boldsymbol{\mu}}$ is the image, under the map $A\mapsto F=\eta^{-\frac32}\cdot dA$, of a probability measure on $V_{\mathrm{t}}$ whose density with respect to $\mathrm{d}A$ is at most a constant times $\exp\left(-\frac12\|dA\|^2\right)$. For a real test $f$ and $F=\eta^{-\frac32}\cdot dA$, \eqref{eq:continuum-infinite-2} and the Cauchy--Schwarz inequality for the counting pairing give
\[
\left|X_k(f)\right|=\left|\langle f_\eta,F\rangle_\eta\right|=\left|\left\langle\eta^{\frac32}\cdot f_\eta,dA\right\rangle\right|\le|f_\eta|_\eta\cdot\|dA\| .
\]
Let $r\ge0$ and $c=\sum_{i=1}^s|(f_i)_\eta|_\eta$; then $\sum_i\left|X_k(f_i)\right|\le c\cdot\|dA\|$ for $F=\eta^{-\frac32}\cdot dA$. Since $r\cdot c\cdot t-\frac12t^2\le r^2\cdot c^2-\frac14t^2$ for real $t$, it follows that
\[
\exp\left(r\cdot\sum_{i=1}^s\left|X_k(f_i)\right|\right)\cdot\exp\left(-\frac12\|dA\|^2\right)\le\exp\left(r^2\cdot c^2\right)\cdot\exp\left(-\frac14\|dA\|^2\right)\qquad\left(A\in V_{\mathrm{t}}\right),
\]
and the right side is integrable over $V_{\mathrm{t}}$ with respect to $\mathrm{d}A$, because $d$ is injective on the finite-dimensional space $V_{\mathrm{t}}$ (Section~\ref{sec:tori-cells-cochains}), so that $A\mapsto\|dA\|$ is a norm on $V_{\mathrm{t}}$. Therefore
\[
\int\mathrm{d}\mathsf P(F)\,\exp\left(r\cdot\sum_{i=1}^s\left|X_k(f_i)\right|\right)<\infty\qquad(r\ge0)
\]
for each of these laws $\mathsf P$. Let $z\in\mathbb C^s$. The series $\sum_{\alpha\in\mathbb N_0^s}\frac{z^\alpha}{\alpha!}\cdot\prod_iX_k(f_i)^{\alpha_i}$ has the sum $\exp\left(\sum_iz_i\cdot X_k(f_i)\right)$, and the series of the absolute values of its terms has the sum $\exp\left(\sum_i|z_i|\cdot\left|X_k(f_i)\right|\right)$, which is $\mathsf P$-integrable. By Fubini's theorem
\[
\int\mathrm{d}\mathsf P(F)\,\exp\left(\sum_{i=1}^sz_i\cdot X_k(f_i)\right)=\sum_{\alpha\in\mathbb N_0^s}\frac{z^\alpha}{\alpha!}\cdot\int\mathrm{d}\mathsf P(F)\,\prod_{i=1}^sX_k(f_i)^{\alpha_i},
\]
the power series on the right converging absolutely for every $z\in\mathbb C^s$. Thus $G^\theta_k$ and $\bar G_k$ are entire, with the stated derivatives at $0$. Finally, the decomposition of the averaged law displayed above gives
\[
\bar E_k\Phi=\int_{[-\pi,\pi]^3}\mathrm{d}\omega_k(\theta)\,\int\mathrm{d}\mathsf P^\theta_k(F)\,\Phi(F)
\]
for every Borel function $\Phi\colon C^2\to[0,\infty]$, and hence for every Borel function $\Phi\colon C^2\to\mathbb C$ with $\bar E_k|\Phi|<\infty$; the choice $\Phi=\exp\left(\sum_iz_i\cdot X_k(f_i)\right)$ yields the last identity.

\begin{proposition}[Pinned and averaged generating functions]\label{prop:pinned-averaged-generating}
Let $(N_k,T_k)_{k\in\mathbb N}$ be a sequence of cutoffs with $N_k\to\infty$ and $T_k\to\infty$ such that $\mathsf P_{N_k,T_k}$ converges weakly on $H_{-4}$ to a Borel probability measure $\mu$; thus $\mu$ is a limit point of the averaged laws (Definition~\ref{def:limit-points}). For an integer $s\ge1$ and $r>0$ let
\[
D_r=\left\{z\in\mathbb C^s:|z_i|\le r\text{ for }1\le i\le s\right\};
\]
for real tests $f_1,\dots,f_s\in C_c^\infty(\mathbb R^3;\Lambda^2)$ let $G^\theta_k$ and $\bar G_k$ be the functions defined above, and let
\[
S_s(f_1,\dots,f_s)=\int\mathrm{d}\mu(F)\,\prod_{i=1}^sF(f_i)
\]
be the Schwinger function of $\mu$ (Theorems~\ref{thm:continuum-infinite-1}(8) and~\ref{thm:moments-limit-point}). Consider the following two conditions on the sequence $(N_k,T_k)_{k\in\mathbb N}$. The first is that for every integer $s\ge1$ and all real $f_1,\dots,f_s\in C_c^\infty(\mathbb R^3;\Lambda^2)$ there is $\rho>0$ with
\begin{equation}\tag{TW}\label{eq:continuum-infinite-TW}
\lim_{k\to\infty}\,\sup_{z\in D_\rho}\left|G^0_k(z)-\bar G_k(z)\right|=0 .
\end{equation}
The second is that for every integer $s\ge1$ and all real $f_1,\dots,f_s\in C_c^\infty(\mathbb R^3;\Lambda^2)$ there is $\rho'>0$ with
\begin{equation}\tag{TW$'$}\label{eq:continuum-infinite-TWp}
\sup_{k\in\mathbb N}\,\sup_{z\in D_{\rho'}}\left|G^0_k(z)\right|<\infty\qquad\text{and}\qquad\lim_{k\to\infty}E^0_k\prod_{i=1}^sX_k(f_i)=S_s(f_1,\dots,f_s).
\end{equation}
\begin{enumerate}
\item[(1)] The conditions~\eqref{eq:continuum-infinite-TW} and~\eqref{eq:continuum-infinite-TWp} are equivalent.
\item[(2)] If they hold, then for every integer $s\ge1$ and all real $f_1,\dots,f_s\in C_c^\infty(\mathbb R^3;\Lambda^2)$ the law of the random vector $\left(X_k(f_1),\dots,X_k(f_s)\right)$ under $\mathsf P^0_k$ converges weakly, as $k\to\infty$, to the law of $\left(F(f_1),\dots,F(f_s)\right)$ under $\mu$.
\item[(3)] If they hold, then the cumulants of the pinned laws converge to those of $\mu$: for every integer $s\ge1$ and all real $f_1,\dots,f_s\in C_c^\infty(\mathbb R^3;\Lambda^2)$ the joint cumulant of $X_k(f_1),\dots,X_k(f_s)$ under $\mathsf P^0_k$ converges to the joint cumulant of $F(f_1),\dots,F(f_s)$ under $\mu$, that is,
\[
\lim_{k\to\infty}\partial_{z_1}\cdots\partial_{z_s}\left(\log G^0_k\right)(0)=\partial_{z_1}\cdots\partial_{z_s}\left(\log G\right)(0),
\]
where $G(z)=\int\mathrm{d}\mu(F)\,\exp\left(\sum_iz_i\cdot F(f_i)\right)$ is the function of Theorem~\ref{thm:continuum-infinite-1}(2), and $\log G^0_k$ and $\log G$ are the principal branches, which are analytic near $0$ because $G^0_k(0)=G(0)=1$.
\item[(4)] If they hold, then for every integer $s\ge1$ the multilinear form
\[
(f_1,\dots,f_s)\mapsto\lim_{k\to\infty}E^0_k\prod_{i=1}^sX_k(f_i),
\]
defined for real $f_1,\dots,f_s\in C_c^\infty(\mathbb R^3;\Lambda^2)$, has exactly one extension to a multilinear form on the real space $L^2(\mathbb R^3;\Lambda^2)$ that is continuous for the norm of $L^2$, namely the restriction to this space of the Schwinger function $S_s$ of Theorem~\ref{thm:moments-limit-point}. The measure $\mu$ and its Schwinger functions have all the properties stated in Theorems~\ref{thm:continuum-infinite-1} and~\ref{thm:moments-limit-point} and in Corollary~\ref{cor:transfer-operators-coordinate}.
\item[(5)] Condition~\eqref{eq:continuum-infinite-TW} holds if for every integer $s\ge1$ and all real $f_1,\dots,f_s\in C_c^\infty(\mathbb R^3;\Lambda^2)$ there is $\rho>0$ with
\[
\lim_{k\to\infty}\,\sup_\theta\,\sup_{z\in D_\rho}\left|G^\theta_k(z)-G^0_k(z)\right|=0,
\]
the supremum being taken over the $\theta\in[-\pi,\pi]^3$ with $Z_{\theta,k}>0$.
\end{enumerate}
\end{proposition}

\begin{proof}
By Corollary~\ref{cor:averaged-laws-admissible} the averaged laws of the parameter assignment form an admissible family. The measure $\mu$ is a limit point of this family along the sequence $(N_k,T_k)$, and the expectation $E_k$ of Definition~\ref{def:limit-points} is $\bar E_k$. Hence Theorem~\ref{thm:estimates-fixed-cutoff} applies to the averaged law at every cutoff, and Theorems~\ref{thm:continuum-infinite-1} and~\ref{thm:moments-limit-point} and Corollary~\ref{cor:transfer-operators-coordinate} apply to $\mu$ and to the sequence $(N_k,T_k)$; this is the last assertion of~(4).

For an entire function $Y$ on $\mathbb C^s$, a number $r>0$ and $\alpha\in\mathbb N_0^s$, Cauchy's integral formula on the torus $\left\{z\in\mathbb C^s:|z_1|=\dots=|z_s|=r\right\}$ gives Cauchy's estimate
\[
\left|\partial^\alpha Y(0)\right|\le\alpha!\cdot r^{-|\alpha|}\cdot\sup_{z\in D_r}\left|Y(z)\right| .
\]
In Steps~1 to~3 an integer $s\ge1$ and real tests $f_1,\dots,f_s\in C_c^\infty(\mathbb R^3;\Lambda^2)$ are fixed, and $G^\theta_k$ and $\bar G_k$ are formed with them.

\emph{Step 1: a bound for $\bar G_k$.} Since $T_k\to\infty$, the periods $R_k=h\cdot L^{T_k}$ tend to infinity. The supports of $f_1,\dots,f_s$ being compact, there is $k_0\in\mathbb N$ such that $\operatorname{supp}f_i\subset\left(-\frac{R_k}{2},\frac{R_k}{2}\right)^3$ for $1\le i\le s$ and $k\ge k_0$. Let $k\ge k_0$ and $z\in\mathbb C^s$, and let $g=\sum_{i=1}^s\operatorname{Re}(z_i)\cdot f_i$; this is a real test supported in $\left(-\frac{R_k}{2},\frac{R_k}{2}\right)^3$. As the functions $X_k(f_i)$ are real,
\[
\left|\exp\left(\sum_{i=1}^sz_i\cdot X_k(f_i)\right)\right|=\exp\left(X_k(g)\right).
\]
Theorem~\ref{thm:estimates-fixed-cutoff}, applied to the averaged law at the cutoff $(N_k,T_k)$, and Lemma~\ref{lem:norm-bounds-lattice} therefore give
\begin{align*}
\left|\bar G_k(z)\right|&\le\bar E_k\exp\left(X_k(g)\right)\le\exp\left(\frac12|g_{\eta_k}|_{\eta_k}^2\right)\\
&\le\exp\left(\frac12\|g\|_{L^2}^2\right)\le\exp\left(\frac12\left(\sum_{i=1}^s|z_i|\cdot\|f_i\|_{L^2}\right)^2\right).
\end{align*}
For each of the finitely many $k<k_0$ the entire function $\bar G_k$ is bounded on every compact subset of $\mathbb C^s$. Consequently
\[
\sup_{k\in\mathbb N}\,\sup_{z\in D_r}\left|\bar G_k(z)\right|<\infty\qquad(r>0).
\]

\emph{Step 2: \eqref{eq:continuum-infinite-TW} implies \eqref{eq:continuum-infinite-TWp}.} Let $\rho>0$ be as in~\eqref{eq:continuum-infinite-TW}. There is $k_1\in\mathbb N$ with $\sup_{z\in D_\rho}\left|G^0_k(z)-\bar G_k(z)\right|\le1$ for $k\ge k_1$, and each of the finitely many entire functions $G^0_k$, $k<k_1$, is bounded on $D_\rho$. Together with Step~1 this gives $\sup_k\sup_{z\in D_\rho}\left|G^0_k(z)\right|<\infty$, the first half of~\eqref{eq:continuum-infinite-TWp} with $\rho'=\rho$. Cauchy's estimate on $D_\rho$, applied to $G^0_k-\bar G_k$, gives
\[
\left|\partial^\alpha\left(G^0_k-\bar G_k\right)(0)\right|\le\alpha!\cdot\rho^{-|\alpha|}\cdot\sup_{z\in D_\rho}\left|G^0_k(z)-\bar G_k(z)\right|\qquad\left(\alpha\in\mathbb N_0^s\right),
\]
and by~\eqref{eq:continuum-infinite-TW} the right side tends to $0$ as $k\to\infty$. For $\alpha=(1,\dots,1)$ the formula for the derivatives at $0$ proved before the proposition turns this into
\[
\lim_{k\to\infty}\left(E^0_k\prod_{i=1}^sX_k(f_i)-\bar E_k\prod_{i=1}^sX_k(f_i)\right)=0 .
\]
By Theorem~\ref{thm:continuum-infinite-1}(2), $\bar E_k\prod_iX_k(f_i)$ converges to $\int\mathrm{d}\mu(F)\,\prod_iF(f_i)=S_s(f_1,\dots,f_s)$. This is the second half of~\eqref{eq:continuum-infinite-TWp}.

\emph{Step 3: \eqref{eq:continuum-infinite-TWp} implies \eqref{eq:continuum-infinite-TW}.} Let $\rho'>0$ be as in~\eqref{eq:continuum-infinite-TWp}. By~\eqref{eq:continuum-infinite-TWp} and Step~1,
\[
M=\sup_{k\in\mathbb N}\,\sup_{z\in D_{\rho'}}\left|G^0_k(z)-\bar G_k(z)\right|<\infty .
\]
For $k\in\mathbb N$ and $\alpha\in\mathbb N_0^s$ let $a_{k,\alpha}=\frac{1}{\alpha!}\cdot\partial^\alpha\left(G^0_k-\bar G_k\right)(0)$, so that
\[
G^0_k(z)-\bar G_k(z)=\sum_{\alpha\in\mathbb N_0^s}a_{k,\alpha}\cdot z^\alpha\qquad\left(z\in\mathbb C^s\right),
\]
the series converging absolutely. Each $a_{k,\alpha}$ tends to $0$ as $k\to\infty$. Indeed, $a_{k,0}=G^0_k(0)-\bar G_k(0)=1-1=0$. If $\alpha\ne0$, let $j=|\alpha|$, and let $(g_1,\dots,g_j)$ be a family of real tests in which $f_i$ occurs exactly $\alpha_i$ times, $1\le i\le s$. By the formula for the derivatives at $0$,
\[
\alpha!\cdot a_{k,\alpha}=E^0_k\prod_{l=1}^jX_k(g_l)-\bar E_k\prod_{l=1}^jX_k(g_l);
\]
the first term tends to $S_j(g_1,\dots,g_j)$ by~\eqref{eq:continuum-infinite-TWp}, applied to the family $(g_1,\dots,g_j)$, and the second tends to $\int\mathrm{d}\mu(F)\,\prod_lF(g_l)=S_j(g_1,\dots,g_j)$ by Theorem~\ref{thm:continuum-infinite-1}(2). Cauchy's estimate on $D_{\rho'}$ gives
\[
\left|a_{k,\alpha}\right|\le M\cdot(\rho')^{-|\alpha|}\qquad\left(k\in\mathbb N,\ \alpha\in\mathbb N_0^s\right).
\]
Put $r=\frac{\rho'}{2}$. For $z\in D_r$
\[
\left|G^0_k(z)-\bar G_k(z)\right|\le\sum_{\alpha\in\mathbb N_0^s}\left|a_{k,\alpha}\right|\cdot r^{|\alpha|} .
\]
The terms of the series on the right tend to $0$ as $k\to\infty$, and they are bounded by $M\cdot2^{-|\alpha|}$, independently of $k$, where $\sum_{\alpha\in\mathbb N_0^s}2^{-|\alpha|}=2^s$. By dominated convergence the series tends to $0$. Hence $\sup_{z\in D_r}\left|G^0_k(z)-\bar G_k(z)\right|\to0$, which is~\eqref{eq:continuum-infinite-TW} with $\rho=r$. Steps~2 and~3 prove~(1).

\emph{Step 4: the laws of the vectors.} Suppose that~\eqref{eq:continuum-infinite-TWp} holds, let $s\ge1$ and real $f_1,\dots,f_s\in C_c^\infty(\mathbb R^3;\Lambda^2)$ be given, let $t\in\mathbb R^s$, and let $f=\sum_{i=1}^st_i\cdot f_i$, a real test in $C_c^\infty(\mathbb R^3;\Lambda^2)$. For $j\ge1$, condition~\eqref{eq:continuum-infinite-TWp}, applied to the family that consists of $j$ copies of $f$, gives
\[
\lim_{k\to\infty}E^0_kX_k(f)^j=S_j(f,\dots,f)=\int\mathrm{d}\mu(F)\,F(f)^j .
\]
By Theorem~\ref{thm:continuum-infinite-1}(1), Proposition~\ref{prop:sub-gaussian-random}(1) holds under $\mu$ for the space $L^2(\mathbb R^3;\Lambda^2)$ with the norm of $L^2$; hence
\[
\int\mathrm{d}\mu(F)\,\exp\left(u\cdot\left|F(f)\right|\right)\le2\cdot\exp\left(\frac12u^2\cdot\|f\|_{L^2}^2\right)<\infty\qquad\left(u\in\mathbb R\right).
\]
So the law of $F(f)$ under $\mu$ has moments of all orders, and its moment generating function is finite on all of $\mathbb R$. A probability measure on $\mathbb R$ whose moment generating function is finite on a neighbourhood of $0$ is determined by its moments; and if real random variables with moments of all orders have moments that converge to those of such a measure, then their laws converge weakly to it (method of moments). The random variable $X_k(f)$ has moments of all orders under $\mathsf P^0_k$, as shown before the proposition. Therefore the law of $X_k(f)$ under $\mathsf P^0_k$ converges weakly to the law of $F(f)$ under $\mu$, and, the function $x\mapsto\exp(ix)$ being bounded and continuous on $\mathbb R$,
\begin{align*}
E^0_k\exp\left(i\sum_{l=1}^st_l\cdot X_k(f_l)\right)&=E^0_k\exp\left(iX_k(f)\right)\\
&\longrightarrow\int\mathrm{d}\mu(F)\,\exp\left(iF(f)\right)=\int\mathrm{d}\mu(F)\,\exp\left(i\sum_{l=1}^st_l\cdot F(f_l)\right)\qquad(k\to\infty).
\end{align*}
Thus the characteristic function of the law of $\left(X_k(f_1),\dots,X_k(f_s)\right)$ under $\mathsf P^0_k$ converges at every point $t\in\mathbb R^s$ to the characteristic function of the law of $\left(F(f_1),\dots,F(f_s)\right)$ under $\mu$, and L\'evy's continuity theorem on $\mathbb R^s$ gives~(2).

\emph{Step 5: the cumulants.} Let $Y$ be analytic on a neighbourhood of $0$ in $\mathbb C^s$ with $Y(0)=1$; then the principal branch of $\log Y$ is analytic on a neighbourhood of $0$. For a partition $\pi$ of $\{1,\dots,s\}$ into non-empty blocks let $\#|\pi|$ be the number of its blocks, and for $B\subset\{1,\dots,s\}$ let $\partial_B=\prod_{i\in B}\partial_{z_i}$. For a function $\phi$ of one complex variable that is analytic near $Y(0)$ one has
\[
\partial_{z_1}\cdots\partial_{z_s}\,\phi\left(Y(z)\right)=\sum_\pi\phi^{(\#|\pi|)}\left(Y(z)\right)\cdot\prod_{B\in\pi}\partial_BY(z)
\]
near $0$, the sum being taken over all partitions $\pi$ of $\{1,\dots,s\}$ into non-empty blocks. This follows by induction on the number of derivatives: the derivative $\partial_{z_{j+1}}$ of a term that belongs to a partition of $\{1,\dots,j\}$ falls either on the factor $\phi^{(\#|\pi|)}(Y)$, which adds the block $\{j+1\}$ to $\pi$, or on one of the factors $\partial_BY$, which adds $j+1$ to the block $B$, and every partition of $\{1,\dots,j+1\}$ arises in exactly one of these ways. For $\phi=\log$ one has $\phi^{(m)}(1)=(-1)^{m-1}\cdot(m-1)!$, and therefore
\[
\partial_{z_1}\cdots\partial_{z_s}\left(\log Y\right)(0)=\sum_\pi(-1)^{\#|\pi|-1}\cdot\left(\#|\pi|-1\right)!\cdot\prod_{B\in\pi}\partial_BY(0).
\]
Suppose now that~\eqref{eq:continuum-infinite-TWp} holds, and let $s\ge1$ and real $f_1,\dots,f_s\in C_c^\infty(\mathbb R^3;\Lambda^2)$ be given. By the formula for the derivatives at $0$, $\partial_BG^0_k(0)=E^0_k\prod_{i\in B}X_k(f_i)$ for every non-empty $B\subset\{1,\dots,s\}$. By Theorem~\ref{thm:continuum-infinite-1}(1), Proposition~\ref{prop:sub-gaussian-random}(1) and H\"older's inequality, $\int\mathrm{d}\mu(F)\,\exp\left(r\cdot\sum_i\left|F(f_i)\right|\right)<\infty$ for every $r\ge0$; so the argument given before the proposition, with $\mu$ in place of $\mathsf P$ and $F(f_i)$ in place of $X_k(f_i)$, shows that $G$ is entire and that $\partial_BG(0)=\int\mathrm{d}\mu(F)\,\prod_{i\in B}F(f_i)$. Condition~\eqref{eq:continuum-infinite-TWp}, applied to the family $(f_i)_{i\in B}$, gives $\partial_BG^0_k(0)\to\partial_BG(0)$ for every non-empty $B\subset\{1,\dots,s\}$. The last display, applied to $Y=G^0_k$ and to $Y=G$, yields~(3).

\emph{Step 6: the extension.} Suppose that~\eqref{eq:continuum-infinite-TWp} holds, and let $s\ge1$. By~\eqref{eq:continuum-infinite-TWp} the limit $\lim_kE^0_k\prod_iX_k(f_i)$ exists for all real $f_1,\dots,f_s\in C_c^\infty(\mathbb R^3;\Lambda^2)$ and equals $S_s(f_1,\dots,f_s)$. For $f\in\mathcal S(\mathbb R^3;\Lambda^2)$ one has $X(f)=F(f)$ $\mu$-almost surely, $X$ being the map of Theorem~\ref{thm:continuum-infinite-1}(1); so on these tests the Schwinger function $S_s$ of Theorem~\ref{thm:moments-limit-point} takes the same values. By Theorem~\ref{thm:moments-limit-point}(1) the restriction of $S_s$ to the real space $L^2(\mathbb R^3;\Lambda^2)$ is multilinear and satisfies
\[
\left|S_s(f_1,\dots,f_s)\right|\le2\cdot\left(2^s\cdot s!\right)^{\frac12}\cdot\prod_{i=1}^s\|f_i\|_{L^2};
\]
hence it is continuous on $L^2(\mathbb R^3;\Lambda^2)^s$ for the norm of $L^2$, and it extends the multilinear form of~(4). If $S'$ is a second multilinear extension that is continuous for the norm of $L^2$, then $S'-S_s$ is continuous on $L^2(\mathbb R^3;\Lambda^2)^s$ and vanishes on the $s$-tuples of real tests in $C_c^\infty(\mathbb R^3;\Lambda^2)$. These $s$-tuples form a dense subset of $L^2(\mathbb R^3;\Lambda^2)^s$, because $C_c^\infty(\mathbb R^3;\Lambda^2)$ is dense in $L^2(\mathbb R^3;\Lambda^2)$; hence $S'=S_s$. Together with the first paragraph of the proof, this proves~(4).

\emph{Step 7: the sufficient condition.} Let $s\ge1$, real $f_1,\dots,f_s\in C_c^\infty(\mathbb R^3;\Lambda^2)$ and $\rho>0$ be as in~(5). For $k\in\mathbb N$ the measure $\omega_k$ is a probability measure, and the $\theta$ with $Z_{\theta,k}=0$ form an $\omega_k$-null set. By the integral formula for $\bar G_k$ proved before the proposition, for $z\in D_\rho$
\begin{align*}
\left|\bar G_k(z)-G^0_k(z)\right|&=\left|\int_{[-\pi,\pi]^3}\mathrm{d}\omega_k(\theta)\,\left(G^\theta_k(z)-G^0_k(z)\right)\right|\\
&\le\int_{[-\pi,\pi]^3}\mathrm{d}\omega_k(\theta)\,\left|G^\theta_k(z)-G^0_k(z)\right|\le\sup_\theta\,\sup_{z'\in D_\rho}\left|G^\theta_k(z')-G^0_k(z')\right|,
\end{align*}
the supremum being taken over the $\theta\in[-\pi,\pi]^3$ with $Z_{\theta,k}>0$. By hypothesis the right side tends to $0$ as $k\to\infty$; so~\eqref{eq:continuum-infinite-TW} holds with this $\rho$, which proves~(5).
\end{proof}

\subsubsection{The pinned law at positive lattice mass}\label{sec:pinned-law-positive}

\begin{theorem}[Insensitivity to the twist]\label{thm:insensitivity-twist}
Let $(N,T)$ be a cutoff, let $\Lambda=\Lambda_{\eta,R}$ be the lattice of Section~\ref{sec:tori-cells-cochains}, and let $n=R/\eta$ be the number of its sites per side. Let $q\ge0$ and $\boldsymbol{\mu}>0$, put
\[
\tau_*=n\cdot\operatorname{arcsinh}\left(\frac{\boldsymbol{\mu}}{12}\right),\qquad\delta=110\cdot n^3\cdot\exp(-\tau_*),
\]
and assume that $\tau_*\ge1$. Then for every $A\in C^1$ and every $\theta\in\mathbb R^3$
\begin{equation}\label{eq:continuum-infinite-7}
\left|\log\det\left(D^\theta_A+\boldsymbol{\mu}\right)-\log\det\left(D^0_A+\boldsymbol{\mu}\right)\right|\le2\delta .
\end{equation}
Consequently, for every $\theta\in\mathbb R^3$ the laws $\mathsf P^\theta$ and $\bar{\mathsf P}$ of Section~\ref{sec:wilson-operator-lattice} have densities with respect to the pinned law $\mathsf P^0$, and these satisfy
\[
\exp(-4\delta)\le\frac{\mathrm{d}\mathsf P^\theta}{\mathrm{d}\mathsf P^0}\le\exp(4\delta),\qquad\exp(-4\delta)\le\frac{\mathrm{d}\bar{\mathsf P}}{\mathrm{d}\mathsf P^0}\le\exp(4\delta).
\]
For every Borel function $\Phi\colon C^2\to[0,\infty]$
\[
\exp(-4\delta)\cdot E^0\Phi\le\bar E\Phi\le\exp(4\delta)\cdot E^0\Phi ,
\]
and
\[
\sup_\Phi\left|E^0\Phi-\bar E\Phi\right|\le\exp(4\delta)-1,
\]
the supremum being taken over all Borel functions $\Phi\colon C^2\to\mathbb C$ with $|\Phi|\le1$. In particular, for every real $g\in C^2(\Lambda)$
\[
E^0\exp\left(\langle g,F\rangle_\eta\right)\le\exp(4\delta)\cdot\exp\left(\frac12|g|_\eta^2\right),
\]
and for every $\nu\in\{1,2,3\}$, every integer $s\ge1$, all real $g_1,\dots,g_s\in C^2(\Lambda)$ supported on $\mathcal P^+_\nu$ and every $c=(c_1,\dots,c_s)\in\mathbb C^s$
\[
\sum_{a,b=1}^s\bar c_a\cdot c_b\cdot E^0\exp\left(i\left\langle g_b-\hat{\mathsf R}_\nu g_a,F\right\rangle_\eta\right)\ge-\left(\exp(4\delta)-1\right)\cdot\left(\sum_{a=1}^s|c_a|\right)^2 .
\]
The double sum on the left of the last inequality is a real number, since the image of the pinned law $\mathsf P^0$ under $\hat{\mathsf R}_\nu$ is $\mathsf P^0$ (Section~\ref{sec:pinned-law-fixed}).
\end{theorem}

\begin{proof}
In Steps 1 to 3 the cochain $A\in C^1$ is fixed; $\langle\cdot,\cdot\rangle$ denotes the scalar product of $\ell^2(\Lambda;\mathbb C^4)$, antilinear in the first argument, $\|\cdot\|$ the corresponding norm and operator norm, and $1$ the identity operator. As in \eqref{eq:continuum-infinite-1}, sites are written in unit coordinates $x\in\{-m,\dots,m\}^3$, coordinates being reduced modulo $n$.

\emph{Step 1: invertibility for real twists.} Let $\theta\in\mathbb R^3$. For $\nu\in\{1,2,3\}$ and a site $x$ put
\[
\zeta_\nu(x)=\varsigma(b)\cdot\exp\left(iq\cdot A(b)\right)\cdot\exp\left(\frac{i\theta_\nu}{n}\right),\qquad b=(x\to x+e_\nu),
\]
and let $S_\nu$ be the operator on $\ell^2(\Lambda;\mathbb C^4)$ given by
\[
(S_\nu\psi)(x)=\zeta_\nu(x)\cdot\psi(x+e_\nu).
\]
As $|\zeta_\nu(x)|=1$ and $x\mapsto x+e_\nu$ permutes the sites, $S_\nu$ is unitary; it acts on the site variable only, so it commutes with $\gamma_1,\gamma_2,\gamma_3$. Its adjoint is
\[
(S_\nu^*\psi)(x)=\overline{\zeta_\nu(x-e_\nu)}\cdot\psi(x-e_\nu),
\]
and, for the bond $b'=(x\to x-e_\nu)$,
\[
\overline{\zeta_\nu(x-e_\nu)}=\varsigma(b')\cdot\exp\left(iq\cdot A(b')\right)\cdot\exp\left(-\frac{i\theta_\nu}{n}\right),
\]
because $b'$ is the reversed bond of $(x-e_\nu\to x)$, $A(-b)=-A(b)$ and $\varsigma(-b)=\varsigma(b)$. Thus $S_\nu$ and $S_\nu^*$ are the hopping operators along the bonds $(x\to x+e_\nu)$ and $(x\to x-e_\nu)$ respectively, with link phase, seam sign and twist included. Since $\gamma(b)=\gamma_\nu$ on bonds of the first kind and $\gamma(b)=-\gamma_\nu$ on bonds of the second kind, \eqref{eq:continuum-infinite-1} with the twist $\theta$ (Section~\ref{sec:wilson-operator-lattice}) gives
\[
D^\theta_A=3-\sum_{\nu=1}^3\left(\frac12\left(1-\gamma_\nu\right)S_\nu+\frac12\left(1+\gamma_\nu\right)S_\nu^*\right)=\frac12\sum_{\nu=1}^3\gamma_\nu\left(S_\nu-S_\nu^*\right)+\frac12\sum_{\nu=1}^3\left(1-S_\nu\right)^*\left(1-S_\nu\right),
\]
where $\left(1-S_\nu\right)^*\left(1-S_\nu\right)=2-S_\nu-S_\nu^*$ by unitarity. The matrix $\gamma_\nu$ is Hermitian and commutes with the anti-Hermitian operator $S_\nu-S_\nu^*$; hence $\gamma_\nu\left(S_\nu-S_\nu^*\right)$ is anti-Hermitian, and $\left\langle u,\gamma_\nu\left(S_\nu-S_\nu^*\right)u\right\rangle$ is purely imaginary. Therefore, for every $u\in\ell^2(\Lambda;\mathbb C^4)$,
\[
\operatorname{Re}\left\langle u,\left(D^\theta_A+\boldsymbol{\mu}\right)u\right\rangle=\boldsymbol{\mu}\cdot\|u\|^2+\frac12\sum_{\nu=1}^3\left\|\left(1-S_\nu\right)u\right\|^2\ge\boldsymbol{\mu}\cdot\|u\|^2 .
\]
By the Cauchy--Schwarz inequality, $\left\|\left(D^\theta_A+\boldsymbol{\mu}\right)u\right\|\ge\boldsymbol{\mu}\cdot\|u\|$. As $\ell^2(\Lambda;\mathbb C^4)$ has finite dimension, $D^\theta_A+\boldsymbol{\mu}$ is invertible, and
\[
\left\|\left(D^\theta_A+\boldsymbol{\mu}\right)^{-1}\right\|\le\frac{1}{\boldsymbol{\mu}}.
\]
In particular $\det\left(D^\theta_A+\boldsymbol{\mu}\right)\ne0$. Since $\det\left(D^\theta_A+\boldsymbol{\mu}\right)\ge0$ by Section~\ref{sec:wilson-operator-lattice}, it follows that $\det\left(D^\theta_A+\boldsymbol{\mu}\right)>0$; so the logarithms in \eqref{eq:continuum-infinite-7} are defined.

\emph{Step 2: twists with an imaginary part.} For $\vartheta=\theta+i\tau\in\mathbb C^3$, $\theta,\tau\in\mathbb R^3$, let $D^\vartheta_A$ be the operator obtained from \eqref{eq:continuum-infinite-1} when the term of each bond $b=(x\to x\pm e_\nu)$ is multiplied by $\exp\left(\pm\frac{i\vartheta_\nu}{n}\right)$; for $\tau=0$ this is the operator $D^\theta_A$ of Section~\ref{sec:wilson-operator-lattice}. The matrix entries of $D^\vartheta_A$ are polynomials in the six functions $\exp\left(\pm\frac{i\vartheta_\nu}{n}\right)$, so that $\vartheta\mapsto\det\left(D^\vartheta_A+\boldsymbol{\mu}\right)$ is an entire function on $\mathbb C^3$. This function is $2\pi$-periodic in each variable: for $\nu\in\{1,2,3\}$ the entire function
\[
\vartheta\mapsto\det\left(D^{\vartheta+2\pi e_\nu}_A+\boldsymbol{\mu}\right)-\det\left(D^\vartheta_A+\boldsymbol{\mu}\right)
\]
vanishes on $\mathbb R^3$ by Lemma~\ref{lem:invariances-averaged-law}(a), and therefore on $\mathbb C^3$ by the identity theorem.

Let $\theta,\tau\in\mathbb R^3$, and let $S_1,S_2,S_3$ be the operators of Step 1 formed with this $\theta$. As
\[
\exp\left(\pm\frac{i(\theta_\nu+i\tau_\nu)}{n}\right)=\exp\left(\pm\frac{i\theta_\nu}{n}\right)\cdot\exp\left(\mp\frac{\tau_\nu}{n}\right),
\]
the perturbation $\mathsf E_\tau=D^{\theta+i\tau}_A-D^\theta_A$ is given by
\[
\mathsf E_\tau=-\sum_{\nu=1}^3\left(\left(\exp\left(-\frac{\tau_\nu}{n}\right)-1\right)\cdot\frac12\left(1-\gamma_\nu\right)S_\nu+\left(\exp\left(\frac{\tau_\nu}{n}\right)-1\right)\cdot\frac12\left(1+\gamma_\nu\right)S_\nu^*\right).
\]
The matrices $\frac12\left(1\mp\gamma_\nu\right)$ are orthogonal projections, $\gamma_\nu$ being Hermitian with $\gamma_\nu^2=1$, and $S_\nu$ is unitary. Since $\left|\exp(-t)-1\right|+\left|\exp(t)-1\right|=2\sinh|t|$ for real $t$, one obtains
\[
\left\|\mathsf E_\tau\right\|\le\sum_{\nu=1}^32\sinh\left(\frac{|\tau_\nu|}{n}\right),
\]
and, by the definition of $\tau_*$,
\[
\left\|\mathsf E_\tau\right\|\le6\sinh\left(\frac{\tau_*}{n}\right)=\frac{\boldsymbol{\mu}}{2}\qquad\text{if }|\tau_\nu|\le\tau_*\text{ for }\nu=1,2,3.
\]
Put $K_{\theta,\tau}=\left(D^\theta_A+\boldsymbol{\mu}\right)^{-1}\mathsf E_\tau$. By Step 1, $\left\|K_{\theta,\tau}\right\|\le\frac12$ whenever $|\tau_\nu|\le\tau_*$ for $\nu=1,2,3$, and then
\[
D^{\theta+i\tau}_A+\boldsymbol{\mu}=\left(D^\theta_A+\boldsymbol{\mu}\right)\left(1+K_{\theta,\tau}\right)
\]
is invertible.

The same estimates give $\left\|K_{\theta,\tau}\right\|\le\frac{6}{\boldsymbol{\mu}}\cdot\sinh\left(\frac{\max_\nu|\tau_\nu|}{n}\right)<1$, and hence $\det\left(D^{\theta+i\tau}_A+\boldsymbol{\mu}\right)\ne0$, for every $\theta+i\tau$ in the open set
\[
\mathcal W=\left\{\theta+i\tau:\ \theta,\tau\in\mathbb R^3,\ |\tau_\nu|<n\cdot\operatorname{arcsinh}\left(\frac{\boldsymbol{\mu}}{6}\right)\text{ for }\nu=1,2,3\right\}.
\]
The set $\mathcal W$ contains the closed tube $\left\{\vartheta\in\mathbb C^3:\ |\operatorname{Im}\vartheta_\nu|\le\tau_*\text{ for }\nu=1,2,3\right\}$, it is convex, and it is invariant under the translations by elements of $\mathbb R^3$. A holomorphic function without zeros on a convex open subset of $\mathbb C^3$ has a holomorphic logarithm. Let therefore $\ell\colon\mathcal W\to\mathbb C$ be holomorphic with
\[
\exp\left(\ell(\vartheta)\right)=\det\left(D^\vartheta_A+\boldsymbol{\mu}\right)\quad(\vartheta\in\mathcal W),\qquad\ell(0)=\log\det\left(D^0_A+\boldsymbol{\mu}\right);
\]
the normalisation is possible because $\det\left(D^0_A+\boldsymbol{\mu}\right)>0$ by Step 1. On $\mathbb R^3$ the function $\theta\mapsto\ell(\theta)-\log\det\left(D^\theta_A+\boldsymbol{\mu}\right)$ is continuous, takes its values in $2\pi i\mathbb Z$ and vanishes at $0$; hence
\[
\ell(\theta)=\log\det\left(D^\theta_A+\boldsymbol{\mu}\right)\in\mathbb R\qquad(\theta\in\mathbb R^3).
\]
For $\nu\in\{1,2,3\}$ the function $\vartheta\mapsto\ell(\vartheta+2\pi e_\nu)-\ell(\vartheta)$ on $\mathcal W$ is continuous and, by the periodicity of the determinant, takes its values in $2\pi i\mathbb Z$. As $\mathcal W$ is connected, this function is constant, and it vanishes on $\mathbb R^3$ by the last display and Lemma~\ref{lem:invariances-averaged-law}(a). So $\ell$ is $2\pi$-periodic in each variable.

For an operator $K$ on $\ell^2(\Lambda;\mathbb C^4)$ with $\|K\|\le\frac12$ put
\[
\log\left(1+K\right)=\sum_{j\ge1}\frac{(-1)^{j+1}}{j}\cdot K^j .
\]
The series converges in norm, uniformly in $K$, and $\exp\left(\log(1+K)\right)=1+K$; together with $\det\left(\exp X\right)=\exp\left(\operatorname{Tr}X\right)$ this gives $\exp\left(\operatorname{Tr}\log(1+K)\right)=\det(1+K)$. Now let $\theta,\tau\in\mathbb R^3$ with $|\tau_\nu|\le\tau_*$ for $\nu=1,2,3$. For $t\in[0,1]$ one has $|t\tau_\nu|\le\tau_*$ as well, so that $\left\|K_{\theta,t\tau}\right\|\le\frac12$ and
\[
\exp\left(\ell(\theta+it\tau)-\ell(\theta)\right)=\frac{\det\left(D^{\theta+it\tau}_A+\boldsymbol{\mu}\right)}{\det\left(D^\theta_A+\boldsymbol{\mu}\right)}=\det\left(1+K_{\theta,t\tau}\right)=\exp\left(\operatorname{Tr}\log\left(1+K_{\theta,t\tau}\right)\right).
\]
Hence the function
\[
t\mapsto\ell(\theta+it\tau)-\ell(\theta)-\operatorname{Tr}\log\left(1+K_{\theta,t\tau}\right),\qquad t\in[0,1],
\]
takes its values in $2\pi i\mathbb Z$. It is continuous, since $K_{\theta,t\tau}$ depends continuously on $t$, and it vanishes at $t=0$, where $K_{\theta,0}=0$. So it vanishes on $[0,1]$. Put $\rho(\theta,\tau)=\operatorname{Tr}\log\left(1+K_{\theta,\tau}\right)$. The space $\ell^2(\Lambda;\mathbb C^4)$ has dimension $4n^3$, so that $\left|\operatorname{Tr}K_{\theta,\tau}^j\right|\le4n^3\cdot\left\|K_{\theta,\tau}\right\|^j\le4n^3\cdot2^{-j}$. The value at $t=1$ therefore gives
\[
\ell(\theta+i\tau)=\ell(\theta)+\rho(\theta,\tau),\qquad|\rho(\theta,\tau)|\le4n^3\cdot\sum_{j\ge1}\frac{2^{-j}}{j}=4n^3\cdot\log2 .
\]

\emph{Step 3: the Fourier coefficients of $\ell$.} For $k\in\mathbb Z^3$ let $|k|_1=|k_1|+|k_2|+|k_3|$ and
\[
c_k=\int_{[-\pi,\pi]^3}\frac{\mathrm{d}\theta}{(2\pi)^3}\,\ell(\theta)\cdot\exp(-ik\cdot\theta).
\]
Let $k\in\mathbb Z^3\setminus\{0\}$, and define $\tau\in\mathbb R^3$ by
\[
\tau_\nu=\begin{cases}-\tau_*&\text{if }k_\nu>0,\\ \tau_*&\text{if }k_\nu<0,\\ 0&\text{if }k_\nu=0,\end{cases}
\]
so that $|\tau_\nu|\le\tau_*$ for $\nu=1,2,3$ and $k\cdot\tau=-\tau_*\cdot|k|_1$. The function $\vartheta\mapsto\ell(\vartheta)\cdot\exp(-ik\cdot\vartheta)$ is holomorphic on $\mathcal W$ and, as $k\in\mathbb Z^3$, $2\pi$-periodic in each variable. If a function of one complex variable is $2\pi$-periodic and holomorphic on a neighbourhood of the closed strip between $\mathbb R$ and $\mathbb R+it$, $t\in\mathbb R$, then its integrals over $[-\pi,\pi]$ and over $[-\pi,\pi]+it$ coincide: this is Cauchy's integral theorem for the rectangle with these two horizontal sides, whose two vertical sides contribute opposite amounts by periodicity. Applying this in the variables $\vartheta_1,\vartheta_2,\vartheta_3$ one after the other, with the shifts $\tau_1,\tau_2,\tau_3$, all points involved lying in the closed tube $\left\{|\operatorname{Im}\vartheta_\nu|\le\tau_*\text{ for }\nu=1,2,3\right\}\subset\mathcal W$, one obtains
\[
c_k=\int_{[-\pi,\pi]^3}\frac{\mathrm{d}\theta}{(2\pi)^3}\,\ell(\theta+i\tau)\cdot\exp\left(-ik\cdot(\theta+i\tau)\right).
\]
By Step 2, $\ell(\theta+i\tau)=\ell(\theta)+\rho(\theta,\tau)$, and $\exp\left(-ik\cdot(\theta+i\tau)\right)=\exp\left(-\tau_*\cdot|k|_1\right)\cdot\exp(-ik\cdot\theta)$. Therefore
\[
c_k=\exp\left(-\tau_*\cdot|k|_1\right)\cdot\left(c_k+\rho_k\right),\qquad\rho_k=\int_{[-\pi,\pi]^3}\frac{\mathrm{d}\theta}{(2\pi)^3}\,\rho(\theta,\tau)\cdot\exp(-ik\cdot\theta),\qquad|\rho_k|\le4n^3\cdot\log2 .
\]
Since $|k|_1\ge1$ and $\tau_*>0$, it follows that
\[
|c_k|\le4n^3\cdot\log2\cdot\frac{\exp\left(-\tau_*\cdot|k|_1\right)}{1-\exp\left(-\tau_*\cdot|k|_1\right)} .
\]
Put $x=\exp(-\tau_*)$; as $\tau_*\ge1$, one has $0<x\le\exp(-1)$. From $\sum_{j\in\mathbb Z}x^{|j|}=\frac{1+x}{1-x}=1+2\sum_{j\ge1}x^j$ one obtains
\[
\sum_{k\in\mathbb Z^3\setminus\{0\}}x^{|k|_1}=\left(\frac{1+x}{1-x}\right)^3-1\le24.9\cdot x ;
\]
indeed, the quotient of the middle term by $x$ is a power series in $x$ with non-negative coefficients, hence increasing on $(0,1)$, and its value at $x=\exp(-1)$ is smaller than $24.83$. Moreover, for $k\ne0$,
\[
\frac{1}{1-\exp\left(-\tau_*\cdot|k|_1\right)}\le\frac{1}{1-x}\le\frac{1}{1-\exp(-1)}\le1.582 .
\]
Consequently
\[
\sum_{k\ne0}|c_k|\le4\log2\cdot1.582\cdot n^3\cdot\sum_{k\ne0}x^{|k|_1}\le4\log2\cdot1.582\cdot24.9\cdot n^3\cdot\exp(-\tau_*)\le\delta,
\]
because $4\log2\cdot1.582\cdot24.9<110$. The restriction of $\ell$ to $\mathbb R^3$ is a smooth function, $2\pi$-periodic in each variable, whose Fourier coefficients are summable; so its Fourier series converges to it absolutely and uniformly, $\ell(\theta)=\sum_{k\in\mathbb Z^3}c_k\cdot\exp(ik\cdot\theta)$ for $\theta\in\mathbb R^3$, and
\[
\sup_{\theta\in\mathbb R^3}\left|\ell(\theta)-c_0\right|\le\sum_{k\ne0}|c_k|\le\delta .
\]
For $\theta\in\mathbb R^3$ one has $\ell(\theta)=\log\det\left(D^\theta_A+\boldsymbol{\mu}\right)$ by Step 2, and $|\ell(\theta)-\ell(0)|\le|\ell(\theta)-c_0|+|c_0-\ell(0)|\le2\delta$. This is \eqref{eq:continuum-infinite-7}.

\emph{Step 4: comparison of the laws.} Let $\theta\in\mathbb R^3$. By Step 1 the weights $w_\theta$ and $w_0$ of Section~\ref{sec:wilson-operator-lattice} are positive on $V_{\mathrm{t}}$, and by \eqref{eq:continuum-infinite-7}
\[
\frac{w_\theta(A)}{w_0(A)}=\frac{\det\left(D^\theta_A+\boldsymbol{\mu}\right)}{\det\left(D^0_A+\boldsymbol{\mu}\right)}\in\left[\exp(-2\delta),\exp(2\delta)\right]\qquad(A\in V_{\mathrm{t}}).
\]
Integration over $V_{\mathrm{t}}$ with respect to $\mathrm{d}A$, and over $[-\pi,\pi]^3$ with respect to $\frac{\mathrm{d}\theta}{(2\pi)^3}$, shows that the numbers
\[
\frac{Z_\theta}{Z_0},\qquad\frac{\bar w(A)}{w_0(A)}\quad(A\in V_{\mathrm{t}}),\qquad\frac{\bar Z}{Z_0}
\]
lie in the same interval $\left[\exp(-2\delta),\exp(2\delta)\right]$; here $0<Z_0<\infty$ by Section~\ref{sec:wilson-operator-lattice}. In particular $Z_\theta>0$, so that $\mathsf P^\theta$ is defined. The map $\iota\colon V_{\mathrm{t}}\to C^2$, $\iota(A)=\eta^{-\frac32}\cdot dA$, is a linear bijection of $V_{\mathrm{t}}$ onto $\operatorname{ran}\,d$ (Section~\ref{sec:tori-cells-cochains}). Let $J_\theta,\bar J\colon C^2\to(0,\infty)$ be the Borel functions which are equal to $1$ outside $\operatorname{ran}\,d$ and satisfy
\[
J_\theta\left(\iota(A)\right)=\frac{Z_0}{Z_\theta}\cdot\frac{w_\theta(A)}{w_0(A)},\qquad\bar J\left(\iota(A)\right)=\frac{Z_0}{\bar Z}\cdot\frac{\bar w(A)}{w_0(A)}\qquad(A\in V_{\mathrm{t}}).
\]
Both take their values in $\left[\exp(-4\delta),\exp(4\delta)\right]$. For every Borel function $\Phi\colon C^2\to[0,\infty]$ the definitions of $\mathsf P^\theta$ and $\bar{\mathsf P}$ give
\[
E^\theta\Phi=\frac{1}{Z_\theta}\int_{V_{\mathrm{t}}}\mathrm{d}A\,w_\theta(A)\cdot\Phi\left(\iota(A)\right)=\frac{1}{Z_0}\int_{V_{\mathrm{t}}}\mathrm{d}A\,w_0(A)\cdot J_\theta\left(\iota(A)\right)\cdot\Phi\left(\iota(A)\right)=E^0\left[J_\theta\cdot\Phi\right]
\]
and
\[
\bar E\Phi=\frac{1}{\bar Z}\int_{V_{\mathrm{t}}}\mathrm{d}A\,\bar w(A)\cdot\Phi\left(\iota(A)\right)=\frac{1}{Z_0}\int_{V_{\mathrm{t}}}\mathrm{d}A\,w_0(A)\cdot\bar J\left(\iota(A)\right)\cdot\Phi\left(\iota(A)\right)=E^0\left[\bar J\cdot\Phi\right].
\]
Hence $J_\theta=\frac{\mathrm{d}\mathsf P^\theta}{\mathrm{d}\mathsf P^0}$ and $\bar J=\frac{\mathrm{d}\bar{\mathsf P}}{\mathrm{d}\mathsf P^0}$; this proves the bounds on the densities and, for $\Phi\ge0$, the inequalities $\exp(-4\delta)\cdot E^0\Phi\le\bar E\Phi\le\exp(4\delta)\cdot E^0\Phi$. If $\Phi\colon C^2\to\mathbb C$ is a Borel function with $|\Phi|\le1$, then
\[
\left|E^0\Phi-\bar E\Phi\right|=\left|E^0\left[\Phi\cdot\left(1-\bar J\right)\right]\right|\le E^0\left|1-\bar J\right|\le\exp(4\delta)-1,
\]
because $|1-y|\le\max\left(\exp(4\delta)-1,\,1-\exp(-4\delta)\right)=\exp(4\delta)-1$ for $y\in\left[\exp(-4\delta),\exp(4\delta)\right]$.

Let $g\in C^2(\Lambda)$ be real. The cutoff $(N,T)$ and the numbers $q\ge0$ and $\boldsymbol{\mu}>0>-1$ satisfy the hypotheses of Proposition~\ref{prop:gaussian-domination}. With $\Phi(F)=\exp\left(\langle g,F\rangle_\eta\right)$, the inequality $\exp(-4\delta)\cdot E^0\Phi\le\bar E\Phi$ and Proposition~\ref{prop:gaussian-domination} give
\[
E^0\exp\left(\langle g,F\rangle_\eta\right)\le\exp(4\delta)\cdot\bar E\exp\left(\langle g,F\rangle_\eta\right)\le\exp(4\delta)\cdot\exp\left(\frac12|g|_\eta^2\right).
\]

Let $\nu$, $s$, $g_1,\dots,g_s$ and $c$ be as in the last assertion, and define $\Psi\colon C^2\to\mathbb C$ by
\[
\Psi(F)=\sum_{a,b=1}^s\bar c_a\cdot c_b\cdot\exp\left(i\left\langle g_b-\hat{\mathsf R}_\nu g_a,F\right\rangle_\eta\right).
\]
Since $|\Psi|\le\left(\sum_{a=1}^s|c_a|\right)^2$, the bound for $|\Phi|\le1$, applied to $\Phi=\left(\sum_{a=1}^s|c_a|\right)^{-2}\cdot\Psi$ if $c\ne0$, gives
\[
\left|E^0\Psi-\bar E\Psi\right|\le\left(\exp(4\delta)-1\right)\cdot\left(\sum_{a=1}^s|c_a|\right)^2 ,
\]
and this holds trivially if $c=0$. By Proposition~\ref{prop:reflection-positivity}, whose hypotheses are satisfied, $\bar E\Psi$ is real and non-negative. The number $E^0\Psi$ is real as well. Indeed, by Section~\ref{sec:pinned-law-fixed} the law $\mathsf P^0$ is invariant under $r_*$ for every $r\in O_{\mathrm{cub}}$, in particular under $(r_\nu)_*$, which equals $\hat{\mathsf R}_\nu$ by the proof of Lemma~\ref{lem:lattice-symmetries-cell}; and $\hat{\mathsf R}_\nu$ is an orthogonal involution of $C^2$ (Section~\ref{sec:admissible-families-lattice}), hence symmetric, so that
\[
\left\langle g_b-\hat{\mathsf R}_\nu g_a,\hat{\mathsf R}_\nu F\right\rangle_\eta=\left\langle\hat{\mathsf R}_\nu g_b-g_a,F\right\rangle_\eta .
\]
Taking complex conjugates, and then replacing $F$ by $\hat{\mathsf R}_\nu F$ under $E^0$, one obtains
\[
\overline{E^0\Psi}=\sum_{a,b=1}^sc_a\cdot\bar c_b\cdot E^0\exp\left(-i\left\langle g_b-\hat{\mathsf R}_\nu g_a,F\right\rangle_\eta\right)=\sum_{a,b=1}^sc_a\cdot\bar c_b\cdot E^0\exp\left(i\left\langle g_a-\hat{\mathsf R}_\nu g_b,F\right\rangle_\eta\right),
\]
and the right side is $E^0\Psi$ with the indices $a$ and $b$ interchanged. Consequently
\[
E^0\Psi\ge\bar E\Psi-\left|E^0\Psi-\bar E\Psi\right|\ge-\left(\exp(4\delta)-1\right)\cdot\left(\sum_{a=1}^s|c_a|\right)^2 ,
\]
which is the last assertion.
\end{proof}

\section{Ergodic components of laws on $\mathcal S'(\mathbb R^3;\Lambda^2)$}\label{sec:ergodic-components-laws}

This section treats laws on $\mathcal S'(\mathbb R^3;\Lambda^2)$, in particular an arbitrary law that satisfies the hypothesis~\eqref{eq:setup-H}, and involves no lattice; for such laws it establishes the ergodic decomposition under the translations of $\mathbb R^3$ and under the translations along one coordinate axis, the description of the vectors invariant under translations, clustering, and the action of the cubic group on the ergodic components, results which together make up a part of Estimate~5; it also contains the examples of Section~\ref{sec:examples}, which are used in Section~\ref{sec:examples-counterexamples} and are built from white noise, telegraph processes and the free Maxwell law. The notation is proper to this section and is fixed in its first subsection, where the cubic group, the Gaussian bound~\eqref{eq:setup-G}, the conditions $(RP_j)$ and the hypothesis~\eqref{eq:setup-H} are stated again; ergodic decompositions, averages over translations, the determination of a law by exponentials, telegraph processes and the free Maxwell law are treated once more, from the beginning and in a different notation, in Section~\ref{sec:fourth-cumulants-mixtures}.

\subsection{Spaces, symmetries and hypotheses}

\subsubsection{Test forms, distributions and laws}\label{sec:test-forms-distributions}

Points of $\mathbb R^3$ are written $x=(x_1,x_2,x_3)$; $e_1,e_2,e_3$ is the standard basis of $\mathbb R^3$, $n\cdot x$ is the Euclidean scalar product of $n,x\in\mathbb R^3$, and $|x|$ is the Euclidean norm of $x$.

Let $\Lambda^2$ be the real vector space of alternating two-tensors on $\mathbb R^3$, endowed with the scalar product for which the tensors $\mathrm{d}x_i\wedge\mathrm{d}x_l$, $1\le i<l\le3$, form an orthonormal basis. The coordinates of $F\in\Lambda^2$ in this basis are denoted $F_{il}$, $i<l$, and one sets $F_{li}=-F_{il}$ for all $i,l\in\{1,2,3\}$. For a map $f\colon\mathbb R^3\to\Lambda^2$ and $x\in\mathbb R^3$ one writes $f_{il}(x)=\left(f(x)\right)_{il}$.

Let $\mathcal S=\mathcal S(\mathbb R^3;\Lambda^2)$ be the real vector space of Schwartz maps $\mathbb R^3\to\Lambda^2$, with its usual topology, and let $\mathcal D=C_c^\infty(\mathbb R^3;\Lambda^2)\subset\mathcal S$ be the subspace of smooth maps with compact support; the elements of $\mathcal S$ are called test forms. For an open set $O\subset\mathbb R^3$ with closure $\bar O$ put
\[
\mathcal D(O)=\left\{f\in\mathcal D:\operatorname{supp}f\subset O\right\},\qquad\mathcal S(\bar O)=\left\{f\in\mathcal S:\operatorname{supp}f\subset\bar O\right\}.
\]
The real Hilbert space $L^2(\mathbb R^3;\Lambda^2)$ of square integrable maps $\mathbb R^3\to\Lambda^2$ contains $\mathcal S$; its scalar product and its norm are
\[
\langle f,g\rangle=\sum_{i<l}\int\mathrm{d}x\,f_{il}(x)\cdot g_{il}(x),\qquad\|f\|=\langle f,f\rangle^{\frac{1}{2}}.
\]
Thus $\|f\|$ is the $L^2$ norm of $f$; it is also written $\|f\|_2$.

Let $Q=\mathcal S'(\mathbb R^3;\Lambda^2)$ be the space of continuous linear maps $\omega\colon\mathcal S\to\mathbb R$. For $f\in\mathcal S$ the function $X(f)\colon Q\to\mathbb R$ is defined by $X(f)(\omega)=\omega(f)$. The cylinder $\sigma$-algebra of $Q$,
\[
\Sigma=\sigma\left(X(f):f\in\mathcal S\right),
\]
is the smallest $\sigma$-algebra on $Q$ for which every $X(f)$, $f\in\mathcal S$, is measurable. For an open set $O\subset\mathbb R^3$ let
\[
\Sigma_O=\sigma\left(X(f):f\in\mathcal D(O)\right)\subset\Sigma.
\]

The set of probability measures on $(Q,\Sigma)$ is denoted $\mathcal M$; its elements are called laws. The set $\mathcal M$ carries the $\sigma$-algebra $\mathcal B_{\mathcal M}$ generated by the maps $\mathcal M\to[0,1]$, $\rho\mapsto\rho(A)$, where $A$ ranges over $\Sigma$. For $\rho\in\mathcal M$ and a measurable function $F$ on $Q$ which is either complex valued and $\rho$-integrable or takes its values in $[0,\infty]$, one writes
\[
\rho(F)=\int\mathrm{d}\rho\,F.
\]
The complex Hilbert space $L^2(\rho)=L^2(Q,\Sigma,\rho;\mathbb C)$ carries the scalar product
\[
\langle F,G\rangle_\rho=\rho\left(\overline{F}\cdot G\right)\qquad\left(F,G\in L^2(\rho)\right),
\]
which is antilinear in the first entry and linear in the second. For a $\sigma$-algebra $\mathcal G\subset\Sigma$, $L^2(\mathcal G,\rho)$ denotes the closed subspace of $L^2(\rho)$ consisting of the classes which contain a $\mathcal G$-measurable function.

\subsubsection{Euclidean maps and their actions}\label{sec:euclidean-maps-actions}

A Euclidean map is a map $g\colon\mathbb R^3\to\mathbb R^3$ of the form $g(x)=Rx+a$ with $R\in O(3)$ and $a\in\mathbb R^3$; the Euclidean maps form a group under composition. For $R\in O(3)$ let $\wedge^2R\colon\Lambda^2\to\Lambda^2$ be the linear map induced by $R$ on alternating two-tensors; in coordinates,
\[
\left((\wedge^2R)F\right)_{ab}=\sum_{i,l=1}^3R_{ai}\cdot R_{bl}\cdot F_{il}\qquad\left(F\in\Lambda^2,\ a,b\in\{1,2,3\}\right),
\]
and $\wedge^2R$ is orthogonal. A Euclidean map $g(x)=Rx+a$ acts on maps $f\colon\mathbb R^3\to\Lambda^2$, on elements $\omega\in Q$ and on functions $F$ on $Q$ by
\[
(g\cdot f)(x)=(\wedge^2R)\,f\left(R^{-1}(x-a)\right),\qquad(g\omega)(f)=\omega\left(g^{-1}\cdot f\right),\qquad(\hat gF)(\omega)=F\left(g^{-1}\omega\right).
\]
Here $g\omega\in Q$, since $f\mapsto g^{-1}\cdot f$ is a continuous linear bijection of $\mathcal S$ onto itself. For Euclidean maps $g,h$ one has $(gh)\cdot f=g\cdot(h\cdot f)$ and $(gh)\omega=g(h\omega)$. The bijection $\omega\mapsto g\omega$ of $Q$ and its inverse $\omega\mapsto g^{-1}\omega$ are measurable for $\Sigma$, because $X(f)(g\omega)=X\left(g^{-1}\cdot f\right)(\omega)$. For $A\subset Q$ and $O\subset\mathbb R^3$ let $gA=\{g\omega:\omega\in A\}$ and $gO=\{g(x):x\in O\}$, and for $\rho\in\mathcal M$ let $g_*\rho\in\mathcal M$ be the image of $\rho$ under $\omega\mapsto g\omega$, that is, $(g_*\rho)(A)=\rho\left(g^{-1}A\right)$ for $A\in\Sigma$.

The following properties hold for all Euclidean maps $g,h$.
\begin{itemize}
\item $\hat gX(f)=X(g\cdot f)$ for $f\in\mathcal S$, and $\widehat{gh}=\hat g\hat h$.
\item $\|g\cdot f\|=\|f\|$ for $f\in L^2(\mathbb R^3;\Lambda^2)$, because $\wedge^2R$ is orthogonal and $x\mapsto R^{-1}(x-a)$ preserves the Lebesgue measure.
\item For every open set $O\subset\mathbb R^3$, $\hat g$ maps the $\Sigma_O$-measurable functions onto the $\Sigma_{gO}$-measurable functions, because $\hat gX(f)=X(g\cdot f)$ and $g\cdot\mathcal D(O)=\mathcal D(gO)$.
\item $\hat g(F\cdot G)=\left(\hat gF\right)\cdot\left(\hat gG\right)$, $\hat g\overline{F}=\overline{\hat gF}$ and $\hat g(\varphi\circ F)=\varphi\circ\hat gF$ for all functions $F,G$ on $Q$ and every function $\varphi$ defined on the range of $F$.
\item If $\rho\in\mathcal M$ and $g_*\rho=\rho$, then $\rho(\hat gF)=\rho(F)$ whenever $\rho(F)$ is defined, and $\hat g$ is a unitary operator of $L^2(\rho)$ with inverse $\widehat{g^{-1}}$.
\end{itemize}

\paragraph{Translations.}
For $a\in\mathbb R^3$ let $\tau_a$ be the translation $\tau_a(x)=x+a$, and let $U(a)=\hat\tau_a$. Then $(\tau_a\cdot f)(x)=f(x-a)$ and $U(a)X(f)=X(\tau_a\cdot f)$.

\paragraph{Coordinate reflections.}
For $j\in\{1,2,3\}$ let $\theta_j\in O(3)$ be the reflection that changes the sign of the coordinate $x_j$ and fixes the two other coordinates, and let $\Theta_j=\hat\theta_j$. For $i,l\in\{1,2,3\}$,
\[
(\theta_j\cdot f)_{il}(x)=\begin{cases}-f_{il}(\theta_jx)&\text{if }j\in\{i,l\},\\ f_{il}(\theta_jx)&\text{if }j\notin\{i,l\}.\end{cases}
\]
For $c\in\mathbb R$ the reflection in the plane $\{x_j=c\}$ is $\theta_{j,c}=\tau_{2c\cdot e_j}\theta_j$, that is, $\theta_{j,c}(x)=\theta_jx+2c\cdot e_j$.

The operator $\Theta_j$ is complex linear. By contrast, the reflection operators $\Theta_{\nu,c}$ of Theorem~\ref{thm:continuum-infinite-1}(7), which also enter property~\eqref{it:main-theorem-P6} of Theorem~\ref{thm:main-theorem}, and the operator denoted $\Theta_j$ in Section~\ref{sec:test-forms-distributions-2} are antilinear: for $\nu=j$ they are the maps $A\mapsto\overline{\hat\theta_{j,c}A}$, with $c=0$ in Section~\ref{sec:test-forms-distributions-2}. Consequently, for $\rho\in\mathcal M$ and measurable functions $A,B$ on $Q$ such that $\Theta_jA$ and $B$ lie in $L^2(\rho)$, the integral $\int\mathrm{d}\rho\,(\Theta_{j,0}A)\cdot B$ formed with the operator $\Theta_{j,0}$ of Theorem~\ref{thm:continuum-infinite-1}(7) is $\rho\left(\overline{\Theta_jA}\cdot B\right)=\langle\Theta_jA,B\rangle_\rho$.

\paragraph{The cubic group.}
Let
\[
O_{\mathrm{cub}}=\left\{w\in O(3):\text{all entries of }w\text{ are integers}\right\}.
\]
Each column of a matrix $w\in O_{\mathrm{cub}}$ is an integer vector of length $1$, hence of the form $\pm e_k$, and different columns are orthogonal. Therefore every $w\in O_{\mathrm{cub}}$ is given by
\[
w\,e_i=\varepsilon_i\cdot e_{\pi(i)}\qquad(i=1,2,3)
\]
with a unique $\varepsilon=(\varepsilon_1,\varepsilon_2,\varepsilon_3)\in\{\pm1\}^3$ and a unique $\pi\in S_3$, the group of permutations of $\{1,2,3\}$, and every pair $(\varepsilon,\pi)$ occurs; so $\#|O_{\mathrm{cub}}|=48$. The Euclidean map $x\mapsto wx$ acts on forms by
\[
(w\cdot f)_{\pi(i)\pi(l)}(x)=\varepsilon_i\cdot\varepsilon_l\cdot f_{il}\left(w^{-1}x\right)\qquad(i,l\in\{1,2,3\}).
\]
The map $w\mapsto\pi$ is a homomorphism of $O_{\mathrm{cub}}$ onto $S_3$. Its kernel $N_3$ consists of the diagonal matrices in $O_{\mathrm{cub}}$; it is a normal subgroup with $\#|N_3|=8$, generated by $\theta_1,\theta_2,\theta_3$. Hence $O_{\mathrm{cub}}/N_3\cong S_3$.

\paragraph{Index sets.}
For $J\subset\{1,2,3\}$ let
\[
\Gamma_J=\left\{a\in\mathbb R^3:a_l=0\text{ for every }l\notin J\right\};
\]
in particular $\Gamma_\emptyset=\{0\}$. In indices, $123$ stands for $\{1,2,3\}$ and $j$ stands for $\{j\}$; thus $\Gamma_{123}=\mathbb R^3$ and $\Gamma_j=\{t\cdot e_j:t\in\mathbb R\}$. Finally, for $j\in\{1,2,3\}$,
\[
\Sigma_j^+=\Sigma_{\{x_j>0\}},\qquad\Sigma_j^-=\Sigma_{\{x_j<0\}}.
\]

\subsubsection{Properties of laws}\label{sec:properties-laws}

Let $\rho\in\mathcal M$, $J\subset\{1,2,3\}$ and $j\in\{1,2,3\}$. The law $\rho$ satisfies the Gaussian bound~\eqref{eq:setup-G} if
\[
\int\mathrm{d}\rho\,\exp X(f)\le\exp\left(\frac{1}{2}\|f\|^2\right)\qquad\text{for every }f\in\mathcal S.
\]
It satisfies $(T_J)$ if
\begin{equation}\tag{$T_J$}\label{eq:ergodic-components-TJ}
(\tau_a)_*\rho=\rho\qquad\text{for every }a\in\Gamma_J,
\end{equation}
and it satisfies $(R_j)$ if
\begin{equation}\tag{$R_j$}\label{eq:ergodic-components-Rj}
(\theta_j)_*\rho=\rho.
\end{equation}
It satisfies $(RP_j)$ if for every integer $n\ge1$, every $c=(c_1,\dots,c_n)\in\mathbb C^n$ and all $f_1,\dots,f_n\in\mathcal D(\{x_j>0\})$
\[
\sum_{p,q=1}^n\overline{c_p}\cdot c_q\cdot\rho\left(\exp\left(i\left(X(f_q)-X(\theta_j\cdot f_p)\right)\right)\right)\ge0.
\]

\paragraph{Reflection positivity at a plane.}
Let $n\in\mathbb R^3$ be a unit vector and $c\in\mathbb R$. The reflection in the plane $\{n\cdot x=c\}$ is the Euclidean map
\[
r_{n,c}(x)=x-2\left(n\cdot x-c\right)\cdot n,
\]
and $r_n=r_{n,0}$. The law $\rho$ is called reflection positive at the plane $\{n\cdot x=c\}$ if
\[
\sum_{p,q=1}^k\overline{z_p}\cdot z_q\cdot\rho\left(\exp\left(i\left(X(f_q)-X(r_{n,c}\cdot f_p)\right)\right)\right)\ge0
\]
for every integer $k\ge1$, all $z_1,\dots,z_k\in\mathbb C$ and all $f_1,\dots,f_k\in\mathcal D(\{n\cdot x>c\})$. For $n=e_j$ one has $r_{e_j,c}=\theta_{j,c}$, and for $n=e_j$, $c=0$ the condition is $(RP_j)$.

The condition concerns the half space $\{n\cdot x>c\}$ on one side of the plane. Since $r_{-n,-c}=r_{n,c}$ and $\{(-n)\cdot x>-c\}=\{n\cdot x<c\}$, reflection positivity at the plane $\{(-n)\cdot x=-c\}$ is the same condition for the half space on the other side. If $(r_{n,c})_*\rho=\rho$, the condition for either half space implies the condition for the other. Indeed, write $r=r_{n,c}$; the map $f\mapsto r\cdot f$ is a bijection of $\mathcal D(\{n\cdot x<c\})$ onto $\mathcal D(\{n\cdot x>c\})$, and $r\cdot(r\cdot f)=f$. For $g_1,\dots,g_k\in\mathcal D(\{n\cdot x<c\})$ and $f_p=r\cdot g_p$ the rules of Section~\ref{sec:euclidean-maps-actions} give
\[
\hat r\exp\left(i\left(X(f_q)-X(r\cdot f_p)\right)\right)=\exp\left(i\left(X(g_q)-X(r\cdot g_p)\right)\right)\qquad(p,q=1,\dots,k),
\]
and $\rho(\hat rF)=\rho(F)$ for bounded measurable $F$, because $r_*\rho=\rho$. Hence the sum formed with $z_1,\dots,z_k$ and $g_1,\dots,g_k$ for the side $\{n\cdot x<c\}$ equals the sum formed with $z_1,\dots,z_k$ and $f_1,\dots,f_k$ for the side $\{n\cdot x>c\}$.

\paragraph{The class $\mathcal C$.}
Let $\mathcal C$ be the set of the laws $\rho\in\mathcal M$ that satisfy $(T_{123})$ and, for $j=1,2,3$, $(R_j)$ and $(RP_j)$.

\subsubsection{Invariant sets, invariant vectors and ergodicity}\label{sec:invariant-sets-invariant}

Let $J\subset\{1,2,3\}$. The $\sigma$-algebra of the $\Gamma_J$-invariant sets is
\[
\mathcal I^J=\left\{A\in\Sigma:\tau_aA=A\text{ for every }a\in\Gamma_J\right\}.
\]
Let $\rho\in\mathcal M$ satisfy $(T_J)$. Then the operators $U(a)$, $a\in\Gamma_J$, are unitary on $L^2(\rho)$ (Section~\ref{sec:euclidean-maps-actions}), so that
\[
\mathcal E^J(\rho)=\left\{F\in L^2(\rho):U(a)F=F\text{ for every }a\in\Gamma_J\right\}
\]
is a closed subspace of $L^2(\rho)$; $P^J_\rho$ denotes the orthogonal projection of $L^2(\rho)$ onto $\mathcal E^J(\rho)$. For $J=\emptyset$ every law satisfies $(T_\emptyset)$, and $\Gamma_\emptyset=\{0\}$ gives $\mathcal I^\emptyset=\Sigma$, $\mathcal E^\emptyset(\rho)=L^2(\rho)$ and $P^\emptyset_\rho=1$, the identity of $L^2(\rho)$.

A law $\rho$ with $(T_J)$ is called $\Gamma_J$-ergodic if $\mathcal E^J(\rho)=\mathbb C\cdot1$, that is, if the constants are the only elements of $L^2(\rho)$ that are invariant under all $U(a)$, $a\in\Gamma_J$. The $\Gamma_J$-ergodic laws with $(T_J)$ form the set $\mathcal M^J_{\mathrm{erg}}\subset\mathcal M$.

A law $\rho$ satisfying $(T_{123})$ is called mixing if
\[
\lim_{|a|\to\infty}\rho\left(\overline{F}\cdot U(a)G\right)=\overline{\rho(F)}\cdot\rho(G)\qquad\text{for all }F,G\in L^2(\rho).
\]
By Lemma~\ref{lem:density-exponentials}(i) it suffices that this hold for all $F,G\in\mathcal A$, where $\mathcal A$ is the complex linear span of the functions $\exp\left(iX(f)\right)$, $f\in\mathcal D_0$, with the countable family $\mathcal D_0\subset\mathcal D$ of Lemma~\ref{lem:countable-family-test}(i). Indeed, $\mathcal A$ is dense in $L^2(\rho)$ (Lemma~\ref{lem:density-exponentials}(i) with $O=\mathbb R^3$, and $\Sigma_{\mathbb R^3}=\Sigma$ by Lemma~\ref{lem:countable-family-test}(ii)), and both sides of the limit relation are continuous in $(F,G)\in L^2(\rho)\times L^2(\rho)$, the left side uniformly in $a$, since $U(a)$ is unitary.

\subsubsection{The hypothesis on the law}\label{sec:hypothesis-law}

Let $\mu\in\mathcal M$. In the notation of this section, the hypothesis~\eqref{eq:setup-H} on $\mu$ is the conjunction of the following conditions:
\begin{multline*}
\text{$\mu$ satisfies \eqref{eq:setup-G} and $(T_{123})$; $w_*\mu=\mu$ for every $w\in O_{\mathrm{cub}}$;}\\
\text{and $\mu$ satisfies $(RP_j)$ for $j=1,2,3$.}
\end{multline*}
The reflections $\theta_1$, $\theta_2$, $\theta_3$ belong to $O_{\mathrm{cub}}$. Hence a law $\mu$ with \eqref{eq:setup-H} satisfies $(R_j)$ for $j=1,2,3$, and it belongs to the class $\mathcal C$ of Section~\ref{sec:properties-laws}.

Let $j\in\{1,2,3\}$, and let $\mu\in\mathcal M$ satisfy \eqref{eq:setup-G}, $(R_j)$ and $(T_j)$, as every law with \eqref{eq:setup-H} does. By Lemma~\ref{lem:reformulations-reflection-positivity}, combined with Lemma~\ref{lem:countable-family-test}(ii) in the case of closed half spaces, the condition $(RP_j)$ on $\mu$ is equivalent to each of the following four properties of $\mu$.
\begin{itemize}
\item One has $\langle\Theta_jG,G\rangle_\mu\ge0$ for every complex polynomial $G$ in finitely many smeared fields $X(f)$ with $f\in\mathcal S$ and $\operatorname{supp}f\subset\{x_j\ge0\}$ (Lemma~\ref{lem:reformulations-reflection-positivity}(4)).
\item One has $\langle\Theta_jG,G\rangle_\mu\ge0$ for every $G\in L^2(\Sigma_j^+,\mu)$.
\item The law $\mu$ is reflection positive, in the sense of Section~\ref{sec:properties-laws}, at every plane $\{x_j=c\}$ with $c\in\mathbb R$.
\item The law $\mu$ is reflection positive at every plane $\{x_j=c\}$, $c\in\mathbb R$, with closed half spaces: for every $c\in\mathbb R$, every integer $k\ge1$, all $z_1,\dots,z_k\in\mathbb C$ and all $f_1,\dots,f_k\in\mathcal S$ with $\operatorname{supp}f_p\subset\{x_j\ge c\}$ for $p=1,\dots,k$,
\[
\sum_{p,q=1}^k\overline{z_p}\cdot z_q\cdot\mu\left(\exp\left(i\left(X(f_q)-X(\theta_{j,c}\cdot f_p)\right)\right)\right)\ge0.
\]
\end{itemize}

Every limit point of the lattice laws $P^{\mathrm{nat}}_{N,T}$, $\mathrm{nat}\in\{\mathrm{af},\mathrm{av}\}$, formed in Part~I of Theorem~\ref{thm:main-theorem} satisfies~\eqref{eq:setup-H}. For both values of $\mathrm{nat}$, the bound~\eqref{eq:setup-G}, the condition $(T_{123})$, the invariance under $O_{\mathrm{cub}}$ and the conditions $(RP_j)$, $j=1,2,3$, are given by the properties~\eqref{it:main-theorem-P1}, \eqref{it:main-theorem-P5} and~\eqref{it:main-theorem-P6} of Part~I(b) of Theorem~\ref{thm:main-theorem}. For a limit point of the averaged lattice laws they are also given by Theorem~\ref{thm:continuum-infinite-1}(1), (5), (6) and~(7), which applies to it because the family of the averaged laws is admissible by Corollary~\ref{cor:averaged-laws-admissible}.

\subsubsection{Closed distributions}\label{sec:closed-distributions}

Let $\ast\colon\Lambda^2\to\mathbb R^3$ be the linear map given by $(\ast F)_k=F_{il}$ for $F\in\Lambda^2$, whenever $(k,i,l)$ is a cyclic permutation of $(1,2,3)$. It is an isometry of the Euclidean space $\Lambda^2$ of Section~\ref{sec:test-forms-distributions} onto $\mathbb R^3$. Applied pointwise, it assigns to every $f\in\mathcal S$ the vector field
\[
\ast f=\left(f_{23},-f_{13},f_{12}\right),
\]
whose components lie in the space $\mathcal S(\mathbb R^3)$ of real Schwartz functions on $\mathbb R^3$. A distribution $\omega\in Q$ is called closed, and one writes $d\omega=0$, if
\[
\omega(f)=0\qquad\text{for every }f\in\mathcal S\text{ such that }\ast f=\nabla\mathsf g\text{ for some }\mathsf g\in\mathcal S(\mathbb R^3).
\]

This is the equation $\operatorname{div}(\ast\omega)=0$ in the sense of distributions. Indeed, let $\mathcal S'(\mathbb R^3)$ be the space of continuous linear maps $\mathcal S(\mathbb R^3)\to\mathbb R$; for $i<l$ let $\omega_{il}\in\mathcal S'(\mathbb R^3)$ map $\varphi\in\mathcal S(\mathbb R^3)$ to $\omega(f)$, where $f\in\mathcal S$ has the component $f_{il}=\varphi$ and its two other components equal to zero, and put $\ast\omega=\left(\omega_{23},-\omega_{13},\omega_{12}\right)$. By linearity,
\[
\omega(f)=\sum_{i<l}\omega_{il}(f_{il})=\sum_{k=1}^3(\ast\omega)_k\left((\ast f)_k\right)\qquad(f\in\mathcal S).
\]
If $\ast f=\nabla\mathsf g$, the right side equals $-\left(\operatorname{div}(\ast\omega)\right)(\mathsf g)$; and every $\mathsf g\in\mathcal S(\mathbb R^3)$ occurs, with $f=\ast^{-1}\nabla\mathsf g\in\mathcal S$. Hence $\omega$ is closed if and only if
\[
\operatorname{div}(\ast\omega)=\partial_1\omega_{23}-\partial_2\omega_{13}+\partial_3\omega_{12}=0\qquad\text{in }\mathcal S'(\mathbb R^3).
\]

The set of the closed distributions belongs to $\Sigma$. To see this, let $(\mathsf g_n)_{n\ge1}$ be a sequence that is dense in the separable space $\mathcal S(\mathbb R^3)$, and let $f_n\in\mathcal S$ be given by $\ast f_n=\nabla\mathsf g_n$. The map $\mathsf g\mapsto\ast^{-1}\nabla\mathsf g$ from $\mathcal S(\mathbb R^3)$ to $\mathcal S$ is continuous, and every $\omega\in Q$ is continuous on $\mathcal S$; so $\omega(f_n)=0$ for all $n$ implies that $\omega$ is closed. Therefore
\[
\{\omega\in Q:d\omega=0\}=\bigcap_{n\ge1}\left\{X(f_n)=0\right\}\in\Sigma.
\]
In particular, let $\rho\in\mathcal M$ satisfy $\rho\left(X(f)^2\right)=0$ for every $f\in\mathcal S$ such that $\ast f=\nabla\mathsf g$ for some $\mathsf g\in\mathcal S(\mathbb R^3)$. Then each of the sets $\{X(f_n)=0\}$ has $\rho$-measure one, and $\rho$ is concentrated on $\{d\omega=0\}$.

\subsection{Ergodic decomposition, invariant vectors and cubic symmetry}\label{sec:ergodic-decomposition-invariant}

\begin{theorem}[Ergodic decomposition]\label{thm:ergodic-decomposition}
Let $J$ be a non-empty subset of $\{1,2,3\}$, and let $m\in\mathcal M$ satisfy $(T_J)$. There exists a probability kernel $(\omega,A)\mapsto m_\omega(A)$ from $(Q,\mathcal I^J)$ to $(Q,\Sigma)$, that is, a family $(m_\omega)_{\omega\in Q}$ of probability measures on $(Q,\Sigma)$ such that $\omega\mapsto m_\omega(A)$ is $\mathcal I^J$-measurable for every $A\in\Sigma$, with the following properties.
\begin{enumerate}
\item[(a)] For every $A\in\Sigma$,
\[
m_\omega(A)=E_m\left[1_A\mid\mathcal I^J\right](\omega)\qquad\text{for $m$-almost every }\omega\in Q,
\]
where $E_m\left[\,\cdot\mid\mathcal I^J\right]$ is the conditional expectation under $m$ given $\mathcal I^J$. The map $\omega\mapsto m_\omega$ from $(Q,\mathcal I^J)$ to $(\mathcal M,\mathcal B_{\mathcal M})$ is measurable, since $\mathcal B_{\mathcal M}$ is generated by the maps $\rho\mapsto\rho(A)$, $A\in\Sigma$. The image $\nu$ of $m$ under this map is a probability measure on $(\mathcal M,\mathcal B_{\mathcal M})$, and
\[
m=\int_{\mathcal M}\nu(\mathrm d\rho)\,\rho,\qquad\mathcal M^J_{\mathrm{erg}}\in\mathcal B_{\mathcal M},\qquad\nu\left(\mathcal M^J_{\mathrm{erg}}\right)=1,
\]
where the first identity means that $m(A)=\int_{\mathcal M}\nu(\mathrm d\rho)\,\rho(A)$ for every $A\in\Sigma$. Moreover, every probability measure $\nu'$ on $(\mathcal M,\mathcal B_{\mathcal M})$ with $m=\int_{\mathcal M}\nu'(\mathrm d\rho)\,\rho$ and $\nu'\left(\mathcal M^J_{\mathrm{erg}}\right)=1$ is equal to $\nu$.
\item[(b)] Let $j\in J$, and let $m$ satisfy $(R_j)$ and $(RP_j)$. Then $\nu$-almost every $\rho$ satisfies $(R_j)$ and is reflection positive at the plane $\{x_j=c\}$ for every $c\in\mathbb R$.
\item[(c)] Let $m$ satisfy \eqref{eq:setup-G}. Then $\nu$-almost every $\rho$ satisfies \eqref{eq:setup-G}, the right side of the bound being unchanged; that is, for $\nu$-almost every $\rho$,
\[
\int\mathrm d\rho\,\exp X(f)\le\exp\left(\frac{1}{2}\|f\|^2\right)\qquad\text{for all }f\in\mathcal S\text{ simultaneously}.
\]
\item[(d)] Let $A\in\Sigma$ satisfy $m(A)=1$. Then $\rho(A)=1$ for $\nu$-almost every $\rho$. This applies, for instance, to the set $A=\{\omega\in Q:d\omega=0\}$ of the closed distributions whenever $m(A)=1$; this set belongs to $\Sigma$ by Section~\ref{sec:closed-distributions}.
\end{enumerate}
\end{theorem}

\begin{remark}[The components under all translations]\label{rem:components-all-translations}
Let $\mu$ satisfy \eqref{eq:setup-H}. Then Theorem~\ref{thm:ergodic-decomposition} applies to $m=\mu$ and $J=\{1,2,3\}$, and $\mu$ satisfies $(R_j)$ and $(RP_j)$ for $j=1,2,3$ (Section~\ref{sec:hypothesis-law}). Write $\Xi=\mathcal M$, denote the elements of $\Xi$ by $\xi$, and put $\mu_\xi=\xi$. In this notation, parts (a), (b) for $j=1,2,3$, and (c) of Theorem~\ref{thm:ergodic-decomposition} state that
\[
\mu=\int_\Xi\nu(\mathrm d\xi)\,\mu_\xi,
\]
and that for $\nu$-almost every $\xi$ the law $\mu_\xi$ is invariant and ergodic under the translations of $\mathbb R^3$, is invariant under $\theta_1$, $\theta_2$ and $\theta_3$, is reflection positive at every plane $\{x_j=c\}$ with $j\in\{1,2,3\}$ and $c\in\mathbb R$, and satisfies \eqref{eq:setup-G}. Let $(\mathrm G)$ also denote the set of the laws $\rho\in\mathcal M$ that satisfy \eqref{eq:setup-G}. Then these properties say that
\[
\mu_\xi\in\mathcal C\cap(\mathrm G)\cap\mathcal M^{123}_{\mathrm{erg}}\qquad\text{for $\nu$-almost every }\xi\in\Xi.
\]
\end{remark}

\begin{theorem}[Invariant vectors and clustering]\label{thm:invariant-vectors-clustering}
Let $\rho\in\mathcal C$. For $j=1,2,3$ let $\left(\mathcal H_j,[\cdot]_j,H_j,V_j,\Omega_j\right)$ be the reconstruction of $\rho$ along $e_j$ given by Proposition~\ref{prop:reconstruction-direction}. Thus $\mathcal H_j$ is the completion of the quotient of $L^2(\Sigma_j^+,\rho)$ by the null space $\left\{G:\langle\Theta_jG,G\rangle_\rho=0\right\}$ of the form $(F,G)\mapsto\langle\Theta_jF,G\rangle_\rho$, and $[\cdot]_j\colon L^2(\Sigma_j^+,\rho)\to\mathcal H_j$ is the canonical map; $H_j$ is the non-negative self-adjoint operator on $\mathcal H_j$ such that
\[
\exp\left(-t\cdot H_j\right)[G]_j=\left[U(t\cdot e_j)G\right]_j\qquad\left(t\ge0,\ G\in L^2(\Sigma_j^+,\rho)\right);
\]
for $y\in\mathbb R^3$ with $y\perp e_j$, $V_j(y)$ is the unitary operator on $\mathcal H_j$ with $V_j(y)[G]_j=[U(y)G]_j$ for $G\in L^2(\Sigma_j^+,\rho)$; and $\Omega_j=[1]_j$.
\begin{enumerate}
\item[(i)] \emph{Limits of translates.} Let $K\subset\{1,2,3\}$, let $(a^{(k)})_{k\ge1}$ be a sequence in $\mathbb R^3$, and let $\bar a_l\in\mathbb R$ for $l\notin K$, such that, as $k\to\infty$, $|a^{(k)}_l|\to\infty$ for every $l\in K$ and $a^{(k)}_l\to\bar a_l$ for every $l\notin K$. Put $\bar a=\sum_{l\notin K}\bar a_l\cdot e_l$ and $P^\emptyset_\rho=1$. Then, for all $F,G\in L^2(\rho)$,
\[
\lim_{k\to\infty}\left\langle F,U(a^{(k)})G\right\rangle_\rho=\left\langle F,P^K_\rho U(\bar a)G\right\rangle_\rho,\qquad P^K_\rho=\prod_{l\in K}P^l_\rho=E_\rho\left[\,\cdot\mid\mathcal I^K\right].
\]
\item[(ii)] \emph{Vectors of zero energy.} Let $j\in\{1,2,3\}$, and let $\rho$ be any law in $\mathcal M$ that satisfies $(T_j)$, $(R_j)$ and $(RP_j)$; for such $\rho$ let $\mathcal H_j$, $[\cdot]_j$, $H_j$ and $\Omega_j$ be formed as above. Then $\mathcal E^j(\rho)=L^2(\mathcal I^j,\rho)$, this space is contained in $L^2(\Sigma_j^+,\rho)$ (Lemma~\ref{lem:reflection-positivity-weights}(i)), and $F\mapsto[F]_j$ maps it isometrically onto $\ker H_j$. Moreover, the following properties are equivalent.
\begin{enumerate}
\item[(1)] $\ker H_j=\mathbb C\cdot\Omega_j$.
\item[(2)] $\rho$ is $\Gamma_j$-ergodic.
\item[(3)] For all $F,G\in L^2(\rho)$,
\[
\lim_{|s|\to\infty}\rho\left(\overline{F}\cdot U(s\cdot e_j)G\right)=\overline{\rho(F)}\cdot\rho(G).
\]
\item[(4)] For all $F,G\in L^2(\rho)$,
\[
\lim_{T\to\infty}\frac{1}{T}\int_0^T\mathrm ds\,\rho\left(\overline{F}\cdot U(s\cdot e_j)G\right)=\overline{\rho(F)}\cdot\rho(G).
\]
\end{enumerate}
\item[(iii)] \emph{Laws that are ergodic under all translations.} Let $\rho\in\mathcal C$ be $\Gamma_{123}$-ergodic. Then the following assertions hold for all $F,G\in L^2(\rho)$ and $j=1,2,3$.
\begin{enumerate}
\item[1.] $\rho\left(\overline{F}\cdot U(a)G\right)\to\overline{\rho(F)}\cdot\rho(G)$ as $\min\left\{|a_1|,|a_2|,|a_3|\right\}\to\infty$. In particular, $\rho\left(\overline{F}\cdot U(s\cdot n)G\right)\to\overline{\rho(F)}\cdot\rho(G)$ as $s\to\infty$, for every $n\in\mathbb R^3$ with $n_1\cdot n_2\cdot n_3\ne0$.
\item[2.] One has
\[
\lim_{T\to\infty}\frac{1}{T^3}\int_{[0,T]^3}\mathrm da\,\rho\left(\overline{F}\cdot U(a)G\right)=\overline{\rho(F)}\cdot\rho(G).
\]
\item[3.] Every $\psi\in\mathcal H_j$ such that $\exp\left(-t\cdot H_j\right)\psi=\psi$ for all $t\ge0$ and $V_j(y)\psi=\psi$ for all $y\perp e_j$ is a multiple of $\Omega_j$.
\item[4.] Both limits $s\to+\infty$ and $s\to-\infty$ of $\rho\left(\overline{F}\cdot U(s\cdot e_j)G\right)$ exist, and
\[
\lim_{s\to\pm\infty}\rho\left(\overline{F}\cdot U(s\cdot e_j)G\right)=\rho\left(\overline{E_\rho\left[F\mid\mathcal I^j\right]}\cdot E_\rho\left[G\mid\mathcal I^j\right]\right).
\]
\item[5.] $\dim\ker H_j\in\{1,\infty\}$.
\end{enumerate}
\item[(iv)] \emph{Moments.} Let $\rho\in\mathcal C$ satisfy \eqref{eq:setup-G}. Then every polynomial in the fields $X(f)$, $f\in\mathcal S$, lies in $L^2(\rho)$, so that (i)--(iii) apply to the Schwinger functions
\[
S^\rho_n(f_1,\dots,f_n)=\rho\left(X(f_1)\cdots X(f_n)\right)\qquad\left(n\ge1,\ f_1,\dots,f_n\in\mathcal S\right),
\]
with products of fields in place of $F$ and $G$. Moreover, $S^\rho_1=0$,
\[
\left|S^\rho_n(f_1,\dots,f_n)\right|\le n\cdot2^{\frac{n}{2}}\cdot\Gamma\left(\frac{n}{2}\right)\cdot\prod_{p=1}^n\|f_p\|\qquad\left(n\ge1,\ f_1,\dots,f_n\in\mathcal S\right),
\]
$P^j_\rho X(f)=0$ for $j=1,2,3$ and every $f\in\mathcal S$, and $\lim_{|a|\to\infty}S^\rho_2(f,\tau_a\cdot g)=0$ for all $f,g\in\mathcal S$.
\end{enumerate}
\end{theorem}

\begin{theorem}[Cubic symmetry of the components]\label{thm:cubic-symmetry-components}
Let $\mu$ satisfy \eqref{eq:setup-H}, and let $(\mu_\omega)_{\omega\in Q}$ and $\nu$ be the probability kernel and the probability measure given by Theorem~\ref{thm:ergodic-decomposition} for $m=\mu$ and $J=\{1,2,3\}$.
\begin{enumerate}
\item[(i)] For $\nu$-almost every $\rho$, the law $\rho$ is invariant under the group $\mathbb R^3\rtimes N_3$ of the Euclidean maps $x\mapsto\varepsilon x+a$ with $\varepsilon\in N_3$ and $a\in\mathbb R^3$.
\item[(ii)] For every $w\in O_{\mathrm{cub}}$ one has $w_*\mu_\omega=\mu_{w\omega}$ for $\mu$-almost every $\omega\in Q$, and the map $\rho\mapsto w_*\rho$ of $\mathcal M$ into itself leaves $\nu$ invariant and maps the set $\mathcal C\cap(\mathrm G)\cap\mathcal M^{123}_{\mathrm{erg}}$ of Remark~\ref{rem:components-all-translations} onto itself. For $\nu$-almost every $\rho$, the law $w_*\rho$ depends on $w\in O_{\mathrm{cub}}$ only through the class of $w$ in $O_{\mathrm{cub}}/N_3\cong S_3$. For $\nu$-almost every $\rho$, the orbit $O_{\mathrm{cub}}\rho=\left\{w_*\rho:w\in O_{\mathrm{cub}}\right\}$ has $1$, $2$, $3$ or $6$ elements.
\end{enumerate}
\end{theorem}

\begin{theorem}[Orbit averages and invariance of the components]\label{thm:orbit-averages-invariance}
Let $\mu\in\mathcal M$ satisfy~\eqref{eq:setup-H}, and let $(\mu_\omega)_{\omega\in Q}$ and $\nu$ be the probability kernel and the probability measure given by Theorem~\ref{thm:ergodic-decomposition} for $m=\mu$ and $J=\{1,2,3\}$, as in Theorem~\ref{thm:cubic-symmetry-components}.
\begin{enumerate}
\item[(i)] For $\omega\in Q$ let $O_{\mathrm{cub}}\mu_\omega=\left\{w_*\mu_\omega:w\in O_{\mathrm{cub}}\right\}$ be the orbit of $\mu_\omega$ under $O_{\mathrm{cub}}$, and let $\bar\mu_\omega\in\mathcal M$ be the average
\[
\bar\mu_\omega=\frac{1}{48}\sum_{w\in O_{\mathrm{cub}}}w_*\mu_\omega.
\]
Then
\[
\mu=\int_Q\mu(\mathrm d\omega)\,\bar\mu_\omega,\qquad\text{that is,}\qquad\mu(A)=\int_Q\mu(\mathrm d\omega)\,\bar\mu_\omega(A)\quad\text{for every }A\in\Sigma;
\]
here the integrand is an $\mathcal I^{123}$-measurable function of $\omega$, since $w^{-1}A\in\Sigma$ for $w\in O_{\mathrm{cub}}$ and
\[
\bar\mu_\omega(A)=\frac{1}{48}\sum_{w\in O_{\mathrm{cub}}}\mu_\omega\left(w^{-1}A\right)\qquad(\omega\in Q).
\]
For $a\in\mathbb R^3$ put $\min_l|a_l|=\min\left\{|a_1|,|a_2|,|a_3|\right\}$. For $\mu$-almost every $\omega\in Q$ the following four properties hold.
\begin{itemize}
\item The law $\bar\mu_\omega$ satisfies~\eqref{eq:setup-H}.
\item The constant functions are the only elements $F$ of $L^2(\bar\mu_\omega)$ which are invariant under the group $\mathbb R^3\rtimes O_{\mathrm{cub}}$ of the Euclidean maps $g\colon x\mapsto wx+a$ with $w\in O_{\mathrm{cub}}$ and $a\in\mathbb R^3$, in the sense that $\hat gF=F$ for all these maps $g$. (Each such $\hat g$ is a unitary operator of $L^2(\bar\mu_\omega)$, by the first property and Section~\ref{sec:euclidean-maps-actions}.)
\item The space $\mathcal E^{123}(\bar\mu_\omega)$ of the elements of $L^2(\bar\mu_\omega)$ that are invariant under all $U(a)$, $a\in\mathbb R^3$, satisfies
\[
\dim\mathcal E^{123}(\bar\mu_\omega)=\#|O_{\mathrm{cub}}\mu_\omega|\le6.
\]
\item For all $F,G\in L^2(\bar\mu_\omega)$,
\[
\lim_{\min_l|a_l|\to\infty}\bar\mu_\omega\left(\overline{F}\cdot U(a)G\right)=\frac{1}{\#|O_{\mathrm{cub}}\mu_\omega|}\cdot\sum_{\rho\in O_{\mathrm{cub}}\mu_\omega}\overline{\rho(F)}\cdot\rho(G).
\]
The numbers $\rho(F)$ and $\rho(G)$ on the right side are defined for every $\rho\in O_{\mathrm{cub}}\mu_\omega$, because $\rho\le48\cdot\bar\mu_\omega$, so that $L^2(\bar\mu_\omega)\subset L^2(\rho)\subset L^1(\rho)$.
\end{itemize}
\item[(ii)] Let $P=P^{123}_\mu$ be the orthogonal projection of $L^2(\mu)$ onto $\mathcal E^{123}(\mu)$ (Section~\ref{sec:invariant-sets-invariant}), and let $w\in O_{\mathrm{cub}}$. Since $w_*\mu=\mu$, the operator $\hat w$ is unitary on $L^2(\mu)$. The following properties are equivalent.
\begin{enumerate}
\item[(1)] $w_*\rho=\rho$ for $\nu$-almost every $\rho\in\mathcal M$.
\item[(2)] $P\hat w=P$.
\item[(3)] For all $F,G$ in the countable set $\mathcal A^0$ of Lemma~\ref{lem:density-exponentials},
\[
\lim_{\min_l|a_l|\to\infty}\mu\left(\overline{F}\cdot U(a)\hat wG\right)=\lim_{\min_l|a_l|\to\infty}\mu\left(\overline{F}\cdot U(a)G\right).
\]
Both limits exist, by Theorem~\ref{thm:invariant-vectors-clustering}(i) with $K=\{1,2,3\}$, applied to the law $\mu$, which belongs to $\mathcal C$ (Section~\ref{sec:hypothesis-law}).
\end{enumerate}
\end{enumerate}
\end{theorem}

\begin{theorem}[Asymmetric components and non-ergodicity along the axes]\label{thm:asymmetric-components-non}
For $\rho\in\mathcal M$ let $\delta_\rho$ be the probability measure on $(\mathcal M,\mathcal B_{\mathcal M})$ given by $\delta_\rho(B)=1$ if $\rho\in B$ and $\delta_\rho(B)=0$ if $\rho\notin B$. For $\rho\in\mathcal C$ and $j\in\{1,2,3\}$ let $H_j(\rho)$ be the non-negative self-adjoint operator $H_j$, and $\Omega_j=[1]_j$ the unit vector, of the reconstruction of $\rho$ along $e_j$ (Theorem~\ref{thm:invariant-vectors-clustering}, Proposition~\ref{prop:reconstruction-direction}); every law with~\eqref{eq:setup-H} belongs to $\mathcal C$ (Section~\ref{sec:hypothesis-law}).

There exist laws $m_A,m_B\in\mathcal M$ such that each of them satisfies~\eqref{eq:setup-H}, each of them is concentrated on the set of the closed distributions, which belongs to $\Sigma$ by Section~\ref{sec:closed-distributions},
\[
m_A\left(\{\omega\in Q:d\omega=0\}\right)=m_B\left(\{\omega\in Q:d\omega=0\}\right)=1,
\]
and the assertions (A) and (B) below hold. Moreover, assertion (C) below holds.
\begin{enumerate}
\item[(A)] Let $\nu_A$ be the probability measure on $(\mathcal M,\mathcal B_{\mathcal M})$ that Theorem~\ref{thm:ergodic-decomposition}(a) assigns to $m=m_A$ and $J=\{1,2,3\}$; it is the decomposition of $m_A$ into $\Gamma_{123}$-ergodic laws. There are six pairwise distinct laws $\rho_1,\dots,\rho_6\in\mathcal M$ such that
\[
\nu_A=\frac{1}{6}\sum_{p=1}^6\delta_{\rho_p}.
\]
In particular $\rho_p\in\mathcal M^{123}_{\mathrm{erg}}$ for $p=1,\dots,6$, since $\nu_A\left(\mathcal M^{123}_{\mathrm{erg}}\right)=1$. Each $\rho_p$ is Gaussian, that is, for every $f\in\mathcal S$ the function $X(f)$ is a normal random variable under $\rho_p$, possibly degenerate; and each $\rho_p$ is mixing in the sense of Section~\ref{sec:invariant-sets-invariant}. For every $w\in O_{\mathrm{cub}}$ the map $\rho\mapsto w_*\rho$ permutes the set $\{\rho_1,\dots,\rho_6\}$, and it fixes each element of this set if $w\in N_3$. The resulting action of $O_{\mathrm{cub}}/N_3\cong S_3$ on $\{\rho_1,\dots,\rho_6\}$ is simply transitive; that is, for all $p,q\in\{1,\dots,6\}$ the set of the $w\in O_{\mathrm{cub}}$ with $w_*\rho_p=\rho_q$ is a single coset of $N_3$. In particular, no $\rho_p$ is invariant under an element of $O_{\mathrm{cub}}\setminus N_3$: one has $w_*\rho_p\ne\rho_p$ for every $p\in\{1,\dots,6\}$ and every $w\in O_{\mathrm{cub}}\setminus N_3$.
\item[(B)] The law $m_B$ is $\Gamma_{123}$-ergodic, that is, $m_B\in\mathcal M^{123}_{\mathrm{erg}}$. Hence its decomposition is trivial: by the uniqueness assertion of Theorem~\ref{thm:ergodic-decomposition}(a), the probability measure that this theorem assigns to $m=m_B$ and $J=\{1,2,3\}$ is $\delta_{m_B}$. Let $j\in\{1,2,3\}$. Then
\[
\dim\ker H_j(m_B)=\infty,
\]
and there exists $f\in\mathcal D$ such that the limit in the following inequality exists and the inequality holds:
\[
\lim_{s\to\infty}m_B\left(\exp\left(iX(f)\right)\cdot\exp\left(iX\left(\tau_{s\cdot e_j}\cdot f\right)\right)\right)>m_B\left(\exp\left(iX(f)\right)\right)^2.
\]
Moreover, there is no probability measure $\nu'$ on $(\mathcal M,\mathcal B_{\mathcal M})$ such that
\[
m_B(A)=\int_{\mathcal M}\nu'(\mathrm d\rho)\,\rho(A)\qquad\text{for every }A\in\Sigma
\]
and such that $\nu'$-almost every $\rho\in\mathcal M$ satisfies $(T_{123})$ and is $\Gamma_j$-ergodic. (Every $\rho\in\mathcal C$ satisfies $(T_j)$, $(R_j)$ and $(RP_j)$, and by Theorem~\ref{thm:invariant-vectors-clustering}(ii) such a law is $\Gamma_j$-ergodic if and only if $\ker H_j(\rho)=\mathbb C\cdot\Omega_j$.)
\item[(C)] Let $m_1$ and $m_2$ be the laws of Sections~\ref{sec:average-six-anisotropic} and~\ref{sec:white-noise-modulated}, which are built from white noise in place of the free Maxwell law of Lemma~\ref{lem:properties-free-maxwell}. Both laws satisfy~\eqref{eq:setup-H}. Assertion~(A) holds with $m_1$ in place of $m_A$, and with the probability measure that Theorem~\ref{thm:ergodic-decomposition}(a) assigns to $m=m_1$ and $J=\{1,2,3\}$ in place of $\nu_A$; assertion~(B) holds with $m_2$ in place of $m_B$. Neither of the two laws is concentrated on the set of the closed distributions:
\[
m_1\left(\{\omega\in Q:d\omega=0\}\right)<1,\qquad m_2\left(\{\omega\in Q:d\omega=0\}\right)<1.
\]
\end{enumerate}
\end{theorem}

\begin{theorem}[Decomposition along a coordinate axis]\label{thm:decomposition-along-coordinate}
Let $m\in\mathcal{C}$ and $j\in\{1,2,3\}$. Since $m$ satisfies $(T_{123})$, it satisfies $(T_j)$, and Theorem~\ref{thm:ergodic-decomposition} applies to the law $m$ and to $J=\{j\}$. Let $(\omega,A)\mapsto m^j_{\omega}(A)$ be the probability kernel from $(Q,\mathcal{I}^j)$ to $(Q,\Sigma)$ and $\nu^j$ the probability measure on $(\mathcal{M},\mathcal{B}_{\mathcal{M}})$ which that theorem provides; thus $\nu^j$ is the image of $m$ under the map $\omega\mapsto m^j_{\omega}$. For a law $\rho\in\mathcal{M}$ which satisfies $(T_j)$ and $(R_j)$ and is reflection positive at the plane $\{x_j=0\}$, that is, satisfies $(RP_j)$, the reconstruction of $\rho$ along $e_j$ in Proposition~\ref{prop:reconstruction-direction} is defined; $H_j(\rho)$ and $\Omega_j$ denote the self-adjoint operator $H_j\ge0$ and the unit vector $\Omega_j=[1]_j$ of this reconstruction.
For $\nu^j$-almost every $\rho\in\mathcal{M}$ the following hold: $\rho$ satisfies $(T_j)$ and $(R_j)$; $\rho$ is reflection positive at every plane $\{x_j=c\}$, $c\in\mathbb{R}$; $\rho$ satisfies~\eqref{eq:setup-G} if $m$ satisfies~\eqref{eq:setup-G}; one has $\ker H_j(\rho)=\mathbb{C}\,\Omega_j$; and, for all $F,G\in L^2(\rho)$,
\[
\rho\left(\bar{F}\cdot U(se_j)G\right)\to\overline{\rho(F)}\cdot\rho(G)\qquad\text{as }s\in\mathbb{R},\ |s|\to\infty.
\]
Moreover, let $g\colon x\mapsto Rx+a$, with $R\in O(3)$ and $a\in\mathbb{R}^3$, be a Euclidean map such that $g_{*}m=m$ and such that conjugation by $g$ maps the group of the translations $\tau_b$, $b\in\Gamma_j$, onto itself:
\[
\left\{g\circ\tau_b\circ g^{-1}:b\in\Gamma_j\right\}=\left\{\tau_b:b\in\Gamma_j\right\}
\]
(as $g\circ\tau_b\circ g^{-1}=\tau_{Rb}$ for $b\in\mathbb{R}^3$, this condition is equivalent to $R\Gamma_j=\Gamma_j$). Then $g_{*}m^j_{\omega}=m^j_{g\omega}$ for $m$-almost every $\omega\in Q$.
\end{theorem}

\begin{theorem}[Vacuum and symmetries of the components]\label{thm:vacuum-symmetries-components}
Let $m\in\mathcal C$ and $j\in\{1,2,3\}$. The law $m$ satisfies $(T_{123})$, hence also $(T_j)$, and Theorem~\ref{thm:ergodic-decomposition} applies to $m$ both for $J=\{j\}$ and for $J=\{1,2,3\}$. As in Theorem~\ref{thm:decomposition-along-coordinate}, let $(\omega,A)\mapsto m^j_{\omega}(A)$ be the probability kernel from $(Q,\mathcal I^j)$ to $(Q,\Sigma)$ and $\nu^j$ the probability measure on $(\mathcal M,\mathcal B_{\mathcal M})$ which Theorem~\ref{thm:ergodic-decomposition} provides for $J=\{j\}$. Let $(\omega,A)\mapsto m_{\omega}(A)$ be the probability kernel from $(Q,\mathcal I^{123})$ to $(Q,\Sigma)$ and $\nu$ the probability measure on $(\mathcal M,\mathcal B_{\mathcal M})$ which Theorem~\ref{thm:ergodic-decomposition} provides for $J=\{1,2,3\}$; thus $\nu$ is the image of $m$ under the map $\omega\mapsto m_{\omega}$. One also writes $m^{123}_{\omega}=m_{\omega}$. For $k\in\{1,2,3\}$ and every law $\rho\in\mathcal M$ which satisfies $(T_k)$, $(R_k)$ and $(RP_k)$, let $H_k(\rho)$ and $\Omega_k$ denote the self-adjoint operator $H_k\ge0$ and the unit vector $\Omega_k=[1]_k$ of the reconstruction of $\rho$ along $e_k$ (Proposition~\ref{prop:reconstruction-direction}); for $k=j$ these are the objects so denoted in Theorem~\ref{thm:decomposition-along-coordinate}. By Theorem~\ref{thm:ergodic-decomposition}(a), $\nu$-almost every $\rho$ satisfies $(T_{123})$. For $k=1,2,3$ the law $m$ satisfies $(R_k)$ and $(RP_k)$, as $m\in\mathcal C$, and Theorem~\ref{thm:ergodic-decomposition}(b), applied to the index $k\in J=\{1,2,3\}$, shows that $\nu$-almost every $\rho$ satisfies $(R_k)$ and is reflection positive at the plane $\{x_k=0\}$, which is the property $(RP_k)$ (Section~\ref{sec:properties-laws}). Hence $\nu$-almost every $\rho$ belongs to $\mathcal C$, and for such $\rho$ the operators $H_k(\rho)$ and the vectors $\Omega_k$, $k=1,2,3$, are defined. Finally, for $J\subset\{1,2,3\}$, $P^J_m$ is the orthogonal projection of $L^2(m)$ onto $\mathcal E^J(m)$ (Section~\ref{sec:invariant-sets-invariant}), and $\mathcal A^0$ is the countable set of functions on $Q$ defined in Lemma~\ref{lem:density-exponentials}.
\begin{enumerate}
\item[(i)] The following four conditions are equivalent.
\begin{enumerate}
\item[(1)] For $\nu$-almost every $\rho$ one has $\ker H_j(\rho)=\mathbb C\cdot\Omega_j$.
\item[(2)] One has $P^j_m=P^{123}_m$.
\item[(3)] For all $F,G\in\mathcal A^0$ and every $y\in\mathbb Q^3$,
\[
\lim_{s\to\infty}m\left(\overline{F}\cdot U(s\cdot e_j+y)G\right)=\lim_{s\to\infty}m\left(\overline{F}\cdot U(s\cdot e_j)G\right),
\]
where $s$ ranges over $\mathbb R$; both limits exist by Theorem~\ref{thm:invariant-vectors-clustering}(i), applied to the law $m\in\mathcal C$ with $K=\{j\}$.
\item[(4)] There is a probability measure $\nu'$ on $(\mathcal M,\mathcal B_{\mathcal M})$ such that
\[
m=\int_{\mathcal M}\nu'(\mathrm d\rho)\,\rho,\qquad\text{that is,}\qquad m(A)=\int_{\mathcal M}\nu'(\mathrm d\rho)\,\rho(A)\quad\text{for every }A\in\Sigma,
\]
and such that $\nu'$-almost every $\rho$ satisfies $(T_{123})$ and is $\Gamma_j$-ergodic.
\end{enumerate}
If (1)--(4) hold, then $\nu^j=\nu$, and every probability measure $\nu'$ with the properties required in~(4) is equal to $\nu$. If (1)--(4) hold for each of the three indices $1,2,3$ in place of $j$, then $\nu$-almost every $\rho$ is mixing in the sense of Section~\ref{sec:invariant-sets-invariant}: for all $F,G\in L^2(\rho)$,
\[
\lim_{|a|\to\infty}\rho\left(\overline{F}\cdot U(a)G\right)=\overline{\rho(F)}\cdot\rho(G),
\]
in whatever way $a\in\mathbb R^3$ tends to infinity.
\item[(ii)] Let $g\in O(3)$, regarded as the Euclidean map $x\mapsto g\,x$, satisfy $g_{*}m=m$. Then $g_{*}m_{\omega}=m_{g\omega}$ for $m$-almost every $\omega\in Q$. Moreover, for $\nu$-almost every $\rho$ the following holds for $l=1,2,3$ and every $c\in\mathbb R$: the law $\rho$ is invariant under $g\theta_lg^{-1}$, and it is reflection positive at the plane
\[
g\{x_l=c\}=\left\{x\in\mathbb R^3:(g\,e_l)\cdot x=c\right\}.
\]
\item[(iii)] Let $g\in O(3)$ satisfy $g_{*}m=m$, and let all three coordinates of the vector $g\,e_j$ be different from zero. Then the conditions (1)--(4) of~(i) hold.
\item[(iv)] Let $g_{*}m=m$ for every $g\in O(3)$. Then $\nu$-almost every $\rho$ has the following properties: $\rho$ is invariant under the group $\mathbb R^3\rtimes O(3)$ of all Euclidean maps; for all $F,G\in L^2(\rho)$ one has $\rho\left(\overline{F}\cdot U(a)G\right)\to\overline{\rho(F)}\cdot\rho(G)$ as $|a|\to\infty$, in whatever way $a\in\mathbb R^3$ tends to infinity; and $\ker H_j(\rho)=\mathbb C\cdot\Omega_j$.
\item[(v)] Let $n\in\mathbb R^3$ be a unit vector such that $(r_n)_{*}m=m$ and such that $m$ is reflection positive at the plane $\{n\cdot x=0\}$ in the sense of Section~\ref{sec:properties-laws}. Then, for $\nu$-almost every $\rho$, one has $(r_n)_{*}\rho=\rho$, and $\rho$ is reflection positive at the plane $\{n\cdot x=c\}$ for every $c\in\mathbb R$. Moreover, let $P^{(n)}_m$ be the orthogonal projection of $L^2(m)$ onto the closed subspace
\[
\left\{F\in L^2(m):U(s\cdot n)F=F\text{ for every }s\in\mathbb R\right\}.
\]
Then $U(s\cdot n)$ converges weakly to $P^{(n)}_m$ as $|s|\to\infty$: for all $F,G\in L^2(m)$,
\[
\lim_{|s|\to\infty}\left\langle F,U(s\cdot n)G\right\rangle_m=\left\langle F,P^{(n)}_mG\right\rangle_m.
\]
\item[(vi)] Assume that each of the six unit vectors
\[
n=\frac{1}{\sqrt{2}}\left(e_i+e_l\right),\qquad n=\frac{1}{\sqrt{2}}\left(e_i-e_l\right)\qquad(1\le i<l\le3)
\]
has the properties required of $n$ in~(v). Then $\nu$-almost every $\rho$ is invariant under the group $\mathbb R^3\rtimes O_{\mathrm{cub}}$ of the Euclidean maps $x\mapsto w\,x+a$ with $w\in O_{\mathrm{cub}}$ and $a\in\mathbb R^3$; one has $P^k_m=P^{123}_m$ for $k=1,2,3$; and $\nu$-almost every $\rho$ satisfies $\ker H_k(\rho)=\mathbb C\cdot\Omega_k$ for $k=1,2,3$ and $\rho\left(\overline{F}\cdot U(a)G\right)\to\overline{\rho(F)}\cdot\rho(G)$ for all $F,G\in L^2(\rho)$ as $|a|\to\infty$, in whatever way $a\in\mathbb R^3$ tends to infinity. If $w_{*}m=m$ for every $w\in O_{\mathrm{cub}}$, which is the case when $m$ satisfies~\eqref{eq:setup-H}, and if $m$ is reflection positive at the plane $\{x_1=x_2\}$, that is, at the plane $\{n_0\cdot x=0\}$ with $n_0=\left(e_1-e_2\right)/\sqrt{2}$, then each of the six vectors above has the properties required of $n$ in~(v).
\end{enumerate}
\end{theorem}

\begin{corollary}[The ergodic components of $\mu$]\label{cor:ergodic-components-2}
Assume that $\mu\in\mathcal M$ satisfies~\eqref{eq:setup-H}, and write $\mu_{\xi}=\xi$ for $\xi\in\mathcal M$. There is exactly one probability measure $\nu$ on $(\mathcal M,\mathcal B_{\mathcal M})$ which is carried by $\mathcal M^{123}_{\mathrm{erg}}$, in the sense that $\nu\left(\mathcal M^{123}_{\mathrm{erg}}\right)=1$, and for which
\[
\mu=\int_{\mathcal M}\mathrm{d}\nu(\xi)\,\mu_{\xi},\qquad\text{that is,}\qquad\mu(A)=\int_{\mathcal M}\mathrm{d}\nu(\xi)\,\mu_{\xi}(A)\quad\text{for every }A\in\Sigma.
\]
The set $\mathcal C\cap(\mathrm G)\cap\mathcal M^{123}_{\mathrm{erg}}$ of Remark~\ref{rem:components-all-translations} belongs to $\mathcal B_{\mathcal M}$ and has $\nu$-measure $1$. For every $\xi$ in this set and every $w\in O_{\mathrm{cub}}$, the laws $\mu_{\xi}$ and $w_{*}\mu_{\xi}$ have the properties 1--4 below, and $w_{*}\mu_{\xi}$ depends on $w$ only through the class $wN_3$ (property~6). The properties 5 and 6 hold with the quantifiers stated in them. In the properties 1--6,
\[
S^{\xi}_n(f_1,\dots,f_n)=S^{\mu_{\xi}}_n(f_1,\dots,f_n)=\mu_{\xi}\left(X(f_1)\cdots X(f_n)\right)\qquad\left(n\ge1,\ f_1,\dots,f_n\in\mathcal S\right)
\]
are the Schwinger functions of $\mu_{\xi}$ in the sense of Theorem~\ref{thm:invariant-vectors-clustering}(iv).
\begin{enumerate}
\item[1.] The law $\mu_{\xi}$ satisfies~\eqref{eq:setup-G}, with the same right side as $\mu$:
\[
\int\mathrm{d}\mu_{\xi}\,\exp X(f)\le\exp\left(\frac{1}{2}\|f\|^2\right)\qquad\text{for every }f\in\mathcal S.
\]
Consequently $S^{\xi}_1=0$, and
\[
\left|S^{\xi}_n(f_1,\dots,f_n)\right|\le n\cdot2^{\frac{n}{2}}\cdot\Gamma\left(\frac{n}{2}\right)\cdot\prod_{p=1}^n\|f_p\|\qquad\left(n\ge1,\ f_1,\dots,f_n\in\mathcal S\right).
\]
\item[2.] The law $\mu_{\xi}$ is invariant under $\mathbb R^3\rtimes N_3$, and it is reflection positive at every plane $\{x_j=c\}$ with $j\in\{1,2,3\}$ and $c\in\mathbb R$.
\item[3.] Let $j\in\{1,2,3\}$. By 2 the law $\mu_{\xi}$ belongs to $\mathcal C$; let $\left(\mathcal H_j(\mu_{\xi}),[\cdot]_j,H_j,V_j,\Omega_j\right)$ be its reconstruction along $e_j$, as in Theorem~\ref{thm:invariant-vectors-clustering} (Proposition~\ref{prop:reconstruction-direction}). The vectors $\psi\in\mathcal H_j(\mu_{\xi})$ with $\exp(-t\cdot H_j)\psi=\psi$ for every $t\ge0$ and $V_j(y)\psi=\psi$ for every $y\perp e_j$ are exactly the multiples of $\Omega_j$. The Hilbert space $\ker H_j$ is isomorphic to $L^2(\mathcal I^j,\mu_{\xi})$, and $\dim\ker H_j\in\{1,\infty\}$.
\item[4.] Let $n,n'\ge1$ and $f_1,\dots,f_n,g_1,\dots,g_{n'}\in\mathcal S$. Then, with $a\in\mathbb R^3$,
\[
S^{\xi}_{n+n'}\left(f_1,\dots,f_n,\tau_a\cdot g_1,\dots,\tau_a\cdot g_{n'}\right)\to S^{\xi}_n(f_1,\dots,f_n)\cdot S^{\xi}_{n'}(g_1,\dots,g_{n'})\qquad\text{as }\min_l|a_l|\to\infty,
\]
and the same limit is attained in the mean over $[0,T]^3$:
\[
\lim_{T\to\infty}\frac{1}{T^3}\int_{[0,T]^3}\mathrm{d}a\,S^{\xi}_{n+n'}\left(f_1,\dots,f_n,\tau_a\cdot g_1,\dots,\tau_a\cdot g_{n'}\right)=S^{\xi}_n(f_1,\dots,f_n)\cdot S^{\xi}_{n'}(g_1,\dots,g_{n'}).
\]
For all $f,g\in\mathcal S$ one has $S^{\xi}_2(f,\tau_a\cdot g)\to0$ as $|a|\to\infty$ in any manner. Along the coordinate axes and in the coordinate planes the limits exist and are given by Theorem~\ref{thm:invariant-vectors-clustering}(i): if $K\subset\{1,2,3\}$ has one or two elements, if the sequence $\left(a^{(k)}\right)_{k\ge1}$ in $\mathbb R^3$ and the vector $\bar{a}$ are as in Theorem~\ref{thm:invariant-vectors-clustering}(i), and if $F=X(f_1)\cdots X(f_n)$ and $G=X(g_1)\cdots X(g_{n'})$, then
\[
\lim_{k\to\infty}S^{\xi}_{n+n'}\left(f_1,\dots,f_n,\tau_{a^{(k)}}\cdot g_1,\dots,\tau_{a^{(k)}}\cdot g_{n'}\right)=\left\langle F,P^K_{\mu_{\xi}}U(\bar{a})G\right\rangle_{\mu_{\xi}}.
\]
\item[5.] For each fixed $A\in\Sigma$ with $\mu(A)=1$ one has $\mu_{\xi}(A)=1$ for $\nu$-almost every $\xi\in\mathcal M$, where the $\nu$-null set of exceptional $\xi$ depends on $A$. For instance, if $\mu$ is concentrated on the set $\{\omega\in Q:d\omega=0\}$ of closed elements of $Q$, which belongs to $\Sigma$ by Section~\ref{sec:closed-distributions}, then $\mu_{\xi}$ is concentrated on this set for $\nu$-almost every $\xi$.
\item[6.] Let $w\in O_{\mathrm{cub}}$. For every $\xi$ in the set $\mathcal C\cap(\mathrm G)\cap\mathcal M^{123}_{\mathrm{erg}}$, the law $w_{*}\mu_{\xi}$ has the properties 1--4, and $w_{*}\mu_{\xi}$ depends on $w$ only through the class of $w$ in $O_{\mathrm{cub}}/N_3\cong S_3$. For each fixed $A\in\Sigma$ with $\mu(A)=1$ one has $\left(w_{*}\mu_{\xi}\right)(A)=1$ for $\nu$-almost every $\xi\in\mathcal M$, where the $\nu$-null set of exceptional $\xi$ depends on $A$; in this sense $w_{*}\mu_{\xi}$ has property~5. The map $\xi\mapsto w_{*}\mu_{\xi}$, that is, the map $\rho\mapsto w_{*}\rho$ of $\mathcal M$ into itself, preserves $\nu$:
\[
\nu\left(\left\{\rho\in\mathcal M:w_{*}\rho\in B\right\}\right)=\nu(B)\qquad\text{for every }B\in\mathcal B_{\mathcal M}.
\]
\end{enumerate}
\end{corollary}

\begin{proof}
By~\eqref{eq:setup-H} the law $\mu$ satisfies $(T_{123})$, so Theorem~\ref{thm:ergodic-decomposition} applies to $m=\mu$ and $J=\{1,2,3\}$. Let $\nu$ be the probability measure on $(\mathcal M,\mathcal B_{\mathcal M})$ provided by Theorem~\ref{thm:ergodic-decomposition}(a). By that statement, $\mathcal M^{123}_{\mathrm{erg}}$ belongs to $\mathcal B_{\mathcal M}$, one has $\nu\left(\mathcal M^{123}_{\mathrm{erg}}\right)=1$ and $\mu=\int_{\mathcal M}\mathrm{d}\nu(\xi)\,\mu_{\xi}$, and every probability measure on $(\mathcal M,\mathcal B_{\mathcal M})$ with these two properties coincides with $\nu$. This proves the existence and the uniqueness of $\nu$.

Since $\mu$ satisfies~\eqref{eq:setup-G} by~\eqref{eq:setup-H}, Theorem~\ref{thm:ergodic-decomposition}(c) yields a $\nu$-null set outside of which $\mu_{\xi}$ satisfies~\eqref{eq:setup-G}, with the right side $\exp\left(\frac{1}{2}\|f\|^2\right)$ and for all $f\in\mathcal S$ simultaneously; this is the inequality in property~1. By Theorem~\ref{thm:cubic-symmetry-components}(i), $\mu_{\xi}$ is invariant under $\mathbb R^3\rtimes N_3$ for $\nu$-almost every $\xi$. Let $j\in\{1,2,3\}$. As $\theta_j\in N_3\subset O_{\mathrm{cub}}$, hypothesis~\eqref{eq:setup-H} gives $(\theta_j)_{*}\mu=\mu$, which is $(R_j)$ for $\mu$, and $(RP_j)$ for $\mu$ is part of~\eqref{eq:setup-H}. Since $j\in J=\{1,2,3\}$, Theorem~\ref{thm:ergodic-decomposition}(b) shows that, for $\nu$-almost every $\xi$, the law $\mu_{\xi}$ satisfies $(R_j)$ and is reflection positive at every plane $\{x_j=c\}$, $c\in\mathbb R$. The union of the three exceptional sets, $j=1,2,3$, is a $\nu$-null set, and property~2 follows. Every law in $\mathcal M^{123}_{\mathrm{erg}}$ satisfies $(T_{123})$, and reflection positivity at the plane $\{x_j=0\}$ is $(RP_j)$ (Section~\ref{sec:properties-laws}). Together with $\nu\left(\mathcal M^{123}_{\mathrm{erg}}\right)=1$ it follows that, for $\nu$-almost every $\xi$,
\[
\mu_{\xi}\in\mathcal C\cap\mathcal M^{123}_{\mathrm{erg}}\qquad\text{and}\qquad\mu_{\xi}\text{ satisfies~\eqref{eq:setup-G}}.
\]

Let now $\rho\in\mathcal C\cap\mathcal M^{123}_{\mathrm{erg}}$ be an arbitrary law that satisfies~\eqref{eq:setup-G}. The law $\rho$ has the properties 1--4, with $\rho$ in place of $\mu_{\xi}$, for the following reasons.

Property~1: the inequality is~\eqref{eq:setup-G} for $\rho$, and Theorem~\ref{thm:invariant-vectors-clustering}(iv) gives $S^{\rho}_1=0$ together with the bound on $\left|S^{\rho}_n(f_1,\dots,f_n)\right|$.

Property~2: the law $\rho$ satisfies $(T_{123})$, $(R_1)$, $(R_2)$ and $(R_3)$, and $N_3$ is the group generated by $\theta_1,\theta_2,\theta_3$ (Section~\ref{sec:euclidean-maps-actions}); hence $\rho$ is invariant under $\mathbb R^3\rtimes N_3$. For $j\in\{1,2,3\}$ the law $\rho$ satisfies $(R_j)$, $(RP_j)$ and $(T_j)$, the last because $\Gamma_j\subset\Gamma_{123}$, so that Lemma~\ref{lem:reformulations-reflection-positivity}(5) makes $\rho$ reflection positive at every plane $\{x_j=c\}$, $c\in\mathbb R$.

Property~3: let $j\in\{1,2,3\}$. As $\rho\in\mathcal C$, the reconstruction $\left(\mathcal H_j(\rho),[\cdot]_j,H_j,V_j,\Omega_j\right)$ of Theorem~\ref{thm:invariant-vectors-clustering} is defined. Since $U(a)1=1$ for every $a\in\mathbb R^3$, the relations of Theorem~\ref{thm:invariant-vectors-clustering} that define the reconstruction give
\[
\exp(-t\cdot H_j)\Omega_j=\left[U(t\cdot e_j)1\right]_j=\Omega_j\quad(t\ge0),\qquad V_j(y)\Omega_j=\left[U(y)1\right]_j=\Omega_j\quad(y\perp e_j),
\]
so the multiples of $\Omega_j$ are invariant vectors. Conversely, $\rho\in\mathcal M^{123}_{\mathrm{erg}}$ is $\Gamma_{123}$-ergodic, that is, $\mathbb R^3$-ergodic, and Theorem~\ref{thm:invariant-vectors-clustering}(iii)3 shows that every vector of $\mathcal H_j(\rho)$ which is invariant under all $\exp(-t\cdot H_j)$, $t\ge0$, and all $V_j(y)$, $y\perp e_j$, is a multiple of $\Omega_j$. By Theorem~\ref{thm:invariant-vectors-clustering}(ii), the map $F\mapsto[F]_j$ is a linear isometry of $\mathcal E^j(\rho)=L^2(\mathcal I^j,\rho)$ onto $\ker H_j$; thus the Hilbert spaces $\ker H_j$ and $L^2(\mathcal I^j,\rho)$ are isomorphic. Finally, $\dim\ker H_j\in\{1,\infty\}$ by Theorem~\ref{thm:invariant-vectors-clustering}(iii)5.

Property~4: by Theorem~\ref{thm:invariant-vectors-clustering}(iv), every polynomial in the functions $X(f)$, $f\in\mathcal S$, lies in $L^2(\rho)$, and the assertions of Theorem~\ref{thm:invariant-vectors-clustering}(i)--(iii) hold for the Schwinger functions of $\rho$. Explicitly, let $n,n'\ge1$ and $f_1,\dots,f_n,g_1,\dots,g_{n'}\in\mathcal S$, and put $F=X(f_1)\cdots X(f_n)$ and $G=X(g_1)\cdots X(g_{n'})$. The functions $X(f)$ are real valued, and $U(a)$ is multiplicative with $U(a)X(g)=X(\tau_a\cdot g)$ (Section~\ref{sec:euclidean-maps-actions}). Hence, for every $a\in\mathbb R^3$,
\[
\left\langle F,U(a)G\right\rangle_{\rho}=\rho\left(\overline{F}\cdot U(a)G\right)=S^{\rho}_{n+n'}\left(f_1,\dots,f_n,\tau_a\cdot g_1,\dots,\tau_a\cdot g_{n'}\right)
\]
and
\[
\overline{\rho(F)}\cdot\rho(G)=S^{\rho}_n(f_1,\dots,f_n)\cdot S^{\rho}_{n'}(g_1,\dots,g_{n'}).
\]
By these identities, the convergence as $\min_l|a_l|\to\infty$ is Theorem~\ref{thm:invariant-vectors-clustering}(iii)1, and the convergence of the means over $[0,T]^3$ is Theorem~\ref{thm:invariant-vectors-clustering}(iii)2. The convergence $S^{\rho}_2(f,\tau_a\cdot g)\to0$ as $|a|\to\infty$ in any manner, for $f,g\in\mathcal S$, is the last assertion of Theorem~\ref{thm:invariant-vectors-clustering}(iv). If $K\subset\{1,2,3\}$, the sequence $\left(a^{(k)}\right)_{k\ge1}$ and the vector $\bar{a}$ are as in Theorem~\ref{thm:invariant-vectors-clustering}(i), then that theorem gives
\[
\lim_{k\to\infty}S^{\rho}_{n+n'}\left(f_1,\dots,f_n,\tau_{a^{(k)}}\cdot g_1,\dots,\tau_{a^{(k)}}\cdot g_{n'}\right)=\lim_{k\to\infty}\left\langle F,U\left(a^{(k)}\right)G\right\rangle_{\rho}=\left\langle F,P^K_{\rho}U(\bar{a})G\right\rangle_{\rho},
\]
in particular when $K$ has one or two elements.

By Lemma~\ref{lem:measurability-defining-properties}, the sets of the laws that satisfy $(T_{123})$, $(R_j)$, $(RP_j)$ and~\eqref{eq:setup-G}, respectively, for $j=1,2,3$, and the set $\mathcal M^{123}_{\mathrm{erg}}$ belong to $\mathcal B_{\mathcal M}$; hence so does their intersection $\mathcal C\cap(\mathrm G)\cap\mathcal M^{123}_{\mathrm{erg}}$, the set of Remark~\ref{rem:components-all-translations}. It has been shown above that, for $\nu$-almost every $\xi$, the law $\mu_{\xi}$ lies in $\mathcal C\cap\mathcal M^{123}_{\mathrm{erg}}$ and satisfies~\eqref{eq:setup-G}. As $\mu_{\xi}=\xi$, the complement of $\mathcal C\cap(\mathrm G)\cap\mathcal M^{123}_{\mathrm{erg}}$ is therefore a $\nu$-null set, and this set has $\nu$-measure $1$. For every $\xi$ in it, $\rho=\mu_{\xi}$ is a law of the kind just considered, so that $\mu_{\xi}$ has the properties 1--4. Hence the properties 1--4 hold for all $\xi$ outside one $\nu$-null set simultaneously.

Property~5 is Theorem~\ref{thm:ergodic-decomposition}(d) for $m=\mu$: if $A\in\Sigma$ and $\mu(A)=1$, then $\mu_{\xi}(A)=1$ for $\nu$-almost every $\xi$. The set $\{\omega\in Q:d\omega=0\}$ belongs to $\Sigma$ by Section~\ref{sec:closed-distributions} and may therefore serve as $A$.

Property~6: let $w\in O_{\mathrm{cub}}$. By Theorem~\ref{thm:cubic-symmetry-components}(ii), the map $\rho\mapsto w_{*}\rho$ maps the set $\mathcal C\cap(\mathrm G)\cap\mathcal M^{123}_{\mathrm{erg}}$ onto itself. Let $\xi$ belong to this set. Then so does $w_{*}\mu_{\xi}$, and $w_{*}\mu_{\xi}$ has the properties 1--4 by what has been shown for an arbitrary law of this set. No exceptional set of $\xi$ arises here beyond the complement of $\mathcal C\cap(\mathrm G)\cap\mathcal M^{123}_{\mathrm{erg}}$: the assertion holds for every $\xi$ of this set and all $48$ elements $w$ of $O_{\mathrm{cub}}$ simultaneously. Moreover, $\mu_{\xi}$ is invariant under $N_3$, by property~2. Let $w'\in O_{\mathrm{cub}}$ lie in the class of $w$ in $O_{\mathrm{cub}}/N_3$. As $N_3$ is a normal subgroup of $O_{\mathrm{cub}}$, the left and the right classes of $w$ modulo $N_3$ coincide, so that $w'=w\varepsilon$ with some $\varepsilon\in N_3$. Since $(w\varepsilon)\omega=w(\varepsilon\omega)$ for $\omega\in Q$ (Section~\ref{sec:euclidean-maps-actions}), one has
\[
w'_{*}\mu_{\xi}=w_{*}\left(\varepsilon_{*}\mu_{\xi}\right)=w_{*}\mu_{\xi}.
\]
Hence, for every $\xi$ of the set, the law $w_{*}\mu_{\xi}$ depends on $w$ only through the class of $w$ in $O_{\mathrm{cub}}/N_3\cong S_3$. By Theorem~\ref{thm:cubic-symmetry-components}(ii), the map $\rho\mapsto w_{*}\rho$ preserves $\nu$; since $\mu_{\xi}=\xi$, the map $\xi\mapsto w_{*}\mu_{\xi}$ is the map $\rho\mapsto w_{*}\rho$, and the assertion of property~6 on $\nu$ follows. Finally, let $A\in\Sigma$ with $\mu(A)=1$. The set $w^{-1}A$ belongs to $\Sigma$, because the map $\omega\mapsto w\omega$ is measurable for $\Sigma$ (Section~\ref{sec:euclidean-maps-actions}), and~\eqref{eq:setup-H} gives
\[
\mu\left(w^{-1}A\right)=\left(w_{*}\mu\right)(A)=\mu(A)=1.
\]
Property~5, applied to $w^{-1}A$, shows that $\left(w_{*}\mu_{\xi}\right)(A)=\mu_{\xi}\left(w^{-1}A\right)=1$ for $\nu$-almost every $\xi$, the exceptional set depending on $A$. Thus $w_{*}\mu_{\xi}$ has property~5.
\end{proof}

\begin{corollary}[Sharpness of Corollary~\ref{cor:ergodic-components-2}]
For a law $\mu\in\mathcal M$ which satisfies~\eqref{eq:setup-H}, let $\nu$ be the probability measure on $(\mathcal M,\mathcal B_{\mathcal M})$ of Corollary~\ref{cor:ergodic-components-2}, and let $\mu_{\xi}=\xi$ for $\xi\in\mathcal M$. For $\nu$-almost every $\xi$ the law $\mu_{\xi}$ belongs to $\mathcal C$ (property~3 of Corollary~\ref{cor:ergodic-components-2}). For such $\xi$ and for $j\in\{1,2,3\}$ let $\left(\mathcal H_j(\mu_{\xi}),[\cdot]_j,H_j(\mu_{\xi}),V_j,\Omega_j\right)$ be the reconstruction of $\mu_{\xi}$ along $e_j$, as in Theorem~\ref{thm:invariant-vectors-clustering} (Proposition~\ref{prop:reconstruction-direction}); thus $H_j(\mu_{\xi})$ is a non-negative self-adjoint operator on $\mathcal H_j(\mu_{\xi})$, and $\Omega_j=[1]_j$. For $j\in\{1,2,3\}$ and $w\in O_{\mathrm{cub}}$ consider the following properties of $\mu$.
\begin{enumerate}
\item[(a$_j$)] One has $\ker H_j(\mu_{\xi})=\mathbb C\,\Omega_j$ for $\nu$-almost every $\xi$.
\item[(b$_j$)] For $\nu$-almost every $\xi$ and for all $F,G\in L^2(\mu_{\xi})$,
\[
\mu_{\xi}\left(\overline{F}\cdot U(s\cdot e_j)G\right)\to\overline{\mu_{\xi}(F)}\cdot\mu_{\xi}(G)\qquad\text{as }s\in\mathbb R,\ |s|\to\infty.
\]
\item[(c$_w$)] One has $w_{*}\mu_{\xi}=\mu_{\xi}$ for $\nu$-almost every $\xi$.
\end{enumerate}
Every law with~\eqref{eq:setup-H} has the property (c$_w$) for each $w\in N_3$, by property~2 of Corollary~\ref{cor:ergodic-components-2}. The laws $m_A$ and $m_B$ of Theorem~\ref{thm:asymmetric-components-non} satisfy~\eqref{eq:setup-H} and are concentrated on the set $\{\omega\in Q:d\omega=0\}$ of the closed distributions (Section~\ref{sec:closed-distributions}); moreover:
\begin{enumerate}
\item[(i)] for $\mu=m_B$ and for every $j\in\{1,2,3\}$, neither (a$_j$) nor (b$_j$) holds;
\item[(ii)] for $\mu=m_A$ and for every $w\in O_{\mathrm{cub}}\setminus N_3$, the property (c$_w$) does not hold.
\end{enumerate}
In particular, for $j\in\{1,2,3\}$ and $w\in O_{\mathrm{cub}}\setminus N_3$, none of the properties (a$_j$), (b$_j$) and (c$_w$) follows from~\eqref{eq:setup-H} together with $\mu\left(\{\omega\in Q:d\omega=0\}\right)=1$.
\end{corollary}

\begin{proof}
A property of $\xi\in\mathcal M$ holds for $\nu$-almost every $\xi$ if it holds for every $\xi$ outside some set $N\in\mathcal B_{\mathcal M}$ with $\nu(N)=0$. Let $\mu$ satisfy~\eqref{eq:setup-H}, let $\nu$ be as in Corollary~\ref{cor:ergodic-components-2}, and let $\rho\in\mathcal M$ and $c>0$ be such that
\[
\nu(N)\ge c\cdot1_N(\rho)\qquad\text{for every }N\in\mathcal B_{\mathcal M}.
\]
Then $\rho$ lies in no set $N\in\mathcal B_{\mathcal M}$ with $\nu(N)=0$. Consequently, a property which holds for $\nu$-almost every $\xi$ holds for $\xi=\rho$, that is, for the law $\mu_{\rho}=\rho$.

Consider first the law $m_B$. By Theorem~\ref{thm:asymmetric-components-non} it satisfies~\eqref{eq:setup-H} and is concentrated on $\{\omega\in Q:d\omega=0\}$, and by Theorem~\ref{thm:asymmetric-components-non}(B) it is $\Gamma_{123}$-ergodic; as it satisfies $(T_{123})$, it belongs to $\mathcal M^{123}_{\mathrm{erg}}$. Let $\mu=m_B$, and let $\delta_{m_B}$ be the Dirac measure at $m_B$ on $(\mathcal M,\mathcal B_{\mathcal M})$, given by $\delta_{m_B}(N)=1_N(m_B)$ for $N\in\mathcal B_{\mathcal M}$. Since $\mathcal M^{123}_{\mathrm{erg}}\in\mathcal B_{\mathcal M}$ by Theorem~\ref{thm:ergodic-decomposition}(a), one has $\delta_{m_B}\left(\mathcal M^{123}_{\mathrm{erg}}\right)=1$, and
\[
\int_{\mathcal M}\mathrm{d}\delta_{m_B}(\xi)\,\mu_{\xi}(A)=\mu_{m_B}(A)=m_B(A)\qquad\text{for every }A\in\Sigma.
\]
By the uniqueness of $\nu$ in Corollary~\ref{cor:ergodic-components-2}, $\nu=\delta_{m_B}$. Hence the first paragraph applies with $\rho=m_B$ and $c=1$: a property among (a$_j$), (b$_j$) can hold only if the corresponding assertion holds for the law $m_B$ itself.

Let $j\in\{1,2,3\}$. By Theorem~\ref{thm:asymmetric-components-non}(B), the space $\ker H_j(m_B)$ has infinite dimension, whereas $\mathbb C\,\Omega_j$ has dimension at most one. Thus $\ker H_j(m_B)\ne\mathbb C\,\Omega_j$, and (a$_j$) does not hold. Theorem~\ref{thm:asymmetric-components-non}(B) also provides $f\in\mathcal D$ for which the limit
\[
L=\lim_{s\to\infty}m_B\left(\exp\left(iX(f)\right)\cdot\exp\left(iX(\tau_{s\cdot e_j}\cdot f)\right)\right)
\]
exists and satisfies $L>m_B\left(\exp\left(iX(f)\right)\right)^2$. Put $F=\exp\left(-iX(f)\right)$ and $G=\exp\left(iX(f)\right)$. These are bounded $\Sigma$-measurable functions on $Q$, hence elements of $L^2(m_B)$, and $\overline{F}=G$ because $X(f)$ is real valued. The operator $U(s\cdot e_j)$ maps $X(f)$ to $X(\tau_{s\cdot e_j}\cdot f)$ and commutes with the composition by the function $t\mapsto\exp(it)$ (Section~\ref{sec:euclidean-maps-actions}); therefore
\[
U(s\cdot e_j)G=\exp\left(iX(\tau_{s\cdot e_j}\cdot f)\right)\qquad(s\in\mathbb R),
\]
and
\[
\lim_{s\to\infty}m_B\left(\overline{F}\cdot U(s\cdot e_j)G\right)=L>m_B\left(\exp\left(iX(f)\right)\right)^2=m_B\left(\overline{F}\right)\cdot m_B(G)=\overline{m_B(F)}\cdot m_B(G).
\]
So $m_B\left(\overline{F}\cdot U(s\cdot e_j)G\right)$ does not converge to $\overline{m_B(F)}\cdot m_B(G)$ as $|s|\to\infty$, and (b$_j$) does not hold. This proves (i).

Consider now the law $m_A$. By Theorem~\ref{thm:asymmetric-components-non} it satisfies~\eqref{eq:setup-H} and is concentrated on $\{\omega\in Q:d\omega=0\}$. By Theorem~\ref{thm:asymmetric-components-non}(A) there are six distinct laws $\rho_1,\dots,\rho_6\in\mathcal M$, none of which is invariant under any element of $O_{\mathrm{cub}}\setminus N_3$, such that the probability measure $\nu'$ given by Theorem~\ref{thm:ergodic-decomposition}(a) for $m=m_A$ and $J=\{1,2,3\}$ is the uniform probability measure on $\{\rho_1,\dots,\rho_6\}$:
\[
\nu'(N)=\frac{1}{6}\sum_{k=1}^61_N(\rho_k)\qquad\left(N\in\mathcal B_{\mathcal M}\right).
\]
Let $\mu=m_A$. By Theorem~\ref{thm:ergodic-decomposition}(a), $\nu'\left(\mathcal M^{123}_{\mathrm{erg}}\right)=1$ and $\mu=\int_{\mathcal M}\nu'(\mathrm{d}\rho)\,\rho$; by the uniqueness of $\nu$ in Corollary~\ref{cor:ergodic-components-2}, $\nu=\nu'$. In particular $\nu(N)\ge\frac{1}{6}\cdot1_N(\rho_1)$ for every $N\in\mathcal B_{\mathcal M}$, and the first paragraph applies with $\rho=\rho_1$ and $c=\frac{1}{6}$. Let $w\in O_{\mathrm{cub}}\setminus N_3$. If (c$_w$) held, it would follow that $w_{*}\rho_1=\rho_1$, which is excluded by the choice of $\rho_1$. Hence (c$_w$) does not hold, and (ii) is proved.

The last assertion follows from (i) and (ii), since $m_A$ and $m_B$ satisfy~\eqref{eq:setup-H} and
\[
m_A\left(\{\omega\in Q:d\omega=0\}\right)=m_B\left(\{\omega\in Q:d\omega=0\}\right)=1.
\]
\end{proof}

\subsection{Auxiliary lemmas}\label{sec:auxiliary-lemmas}

Let $\mathfrak O$ be the family of open subsets of $\mathbb R^3$ that consists of $\mathbb R^3$ itself and of the six open coordinate half spaces:
\[
\mathfrak O=\left\{\mathbb R^3\right\}\cup\left\{\left\{x\in\mathbb R^3:x_j>0\right\},\ \left\{x\in\mathbb R^3:x_j<0\right\}\ :\ j=1,2,3\right\}.
\]

\begin{lemma}[A countable family of test forms]\label{lem:countable-family-test}
\begin{enumerate}
\item[(i)] There is a countable vector space $\mathcal D_0\subset\mathcal D$ over $\mathbb Q$ such that $w\cdot f\in\mathcal D_0$ and $\tau_a\cdot f\in\mathcal D_0$ for all $f\in\mathcal D_0$, $w\in O_{\mathrm{cub}}$ and $a\in\mathbb Q^3$, and such that, for every $O\in\mathfrak O$, the vector space $\mathcal D_0(O)=\mathcal D_0\cap\mathcal D(O)$ over $\mathbb Q$ is dense in $\mathcal S(\bar O)$ for the topology of $\mathcal S$.
\end{enumerate}
Fix such a space $\mathcal D_0$.
\begin{enumerate}
\item[(ii)] For every $O\in\mathfrak O$ one has $\Sigma_O=\sigma\left(X(f):f\in\mathcal D_0(O)\right)$, and $X(f)$ is $\Sigma_O$-measurable for every $f\in\mathcal S(\bar O)$. Moreover, $\Sigma_{\mathbb R^3}=\Sigma$.
\item[(iii)] The measurable space $(Q,\Sigma)$ is standard Borel.
\end{enumerate}
\end{lemma}

\begin{proof}
(i) The space $\mathcal S$ is metrisable and separable, and every subset of a separable metrisable space is separable. Hence, for every $O\in\mathfrak O$, there is a countable set $C_O\subset\mathcal D(O)$ that is dense in $\mathcal D(O)$ for the topology of $\mathcal S$. Let $\mathcal D_0$ be the set of the finite linear combinations with rational coefficients of the forms
\[
\tau_a\cdot(w\cdot f),\qquad f\in\bigcup_{O\in\mathfrak O}C_O,\quad a\in\mathbb Q^3,\quad w\in O_{\mathrm{cub}}.
\]
It is a vector space over $\mathbb Q$ contained in $\mathcal D$, and it is countable, being spanned over $\mathbb Q$ by a countable set. For $w'\in O_{\mathrm{cub}}$ and $b\in\mathbb Q^3$ the relations $w'\tau_aw=\tau_{w'a}w'w$ and $\tau_b\tau_a=\tau_{a+b}$ between Euclidean maps give
\[
w'\cdot\left(\tau_a\cdot(w\cdot f)\right)=\tau_{w'a}\cdot\left((w'w)\cdot f\right),\qquad\tau_b\cdot\left(\tau_a\cdot(w\cdot f)\right)=\tau_{a+b}\cdot(w\cdot f),
\]
and $w'a\in\mathbb Q^3$, the entries of $w'$ being integers. Since the maps $g\mapsto w'\cdot g$ and $g\mapsto\tau_b\cdot g$ are linear, it follows that $\mathcal D_0$ is stable under $O_{\mathrm{cub}}$ and under the translations $\tau_a$, $a\in\mathbb Q^3$. The choice $a=0$, $w=1$ shows that $C_O\subset\mathcal D_0\cap\mathcal D(O)=\mathcal D_0(O)$. Since $C_O$ is dense in $\mathcal D(O)$ and $\mathcal D_0(O)\subset\mathcal D(O)\subset\mathcal S(\bar O)$, it remains to prove that $\mathcal D(O)$ is dense in $\mathcal S(\bar O)$.

Let $\chi\in C_c^\infty(\mathbb R^3)$ be a real function equal to $1$ on a neighbourhood of $0$, and put $\chi_n(x)=\chi\left(\frac{x}{n}\right)$ for $n\ge1$. Let $j\in\{1,2,3\}$, $O=\{x_j>0\}$ and $f\in\mathcal S(\bar O)$, and let
\[
f_n(x)=\chi_n(x)\cdot f\left(x-\frac{1}{n}\cdot e_j\right)\qquad\left(x\in\mathbb R^3,\ n\ge1\right).
\]
The form $f_n$ is smooth, and its support is compact and contained in $\left\{x_j\ge\frac{1}{n}\right\}\subset O$; thus $f_n\in\mathcal D(O)$. Moreover $f_n\to f$ in $\mathcal S$ as $n\to\infty$: the translated forms $x\mapsto f\left(x-\frac{1}{n}\cdot e_j\right)$ converge to $f$ in $\mathcal S$, and the multiplications by $\chi_n$, $n\ge1$, form an equicontinuous family of linear maps of $\mathcal S$ into itself which converges pointwise to the identity. For $O=\{x_j<0\}$ the same argument applies with $-\frac{1}{n}\cdot e_j$ in place of $\frac{1}{n}\cdot e_j$, and for $O=\mathbb R^3$ it applies to the forms $\chi_n\cdot f$.

(ii) Every $\omega\in Q$ is continuous on $\mathcal S$; hence $X(g_k)\to X(g)$ pointwise on $Q$ whenever $g_k\to g$ in $\mathcal S$. Let $O\in\mathfrak O$ and $f\in\mathcal S(\bar O)$. By (i), and because the topology of $\mathcal S$ is metrisable, there are $f_k\in\mathcal D_0(O)$ with $f_k\to f$ in $\mathcal S$ as $k\to\infty$; so $X(f)$ is measurable for $\sigma\left(X(g):g\in\mathcal D_0(O)\right)$. As $\mathcal D_0(O)\subset\mathcal D(O)\subset\mathcal S(\bar O)$, it follows that
\[
\sigma\left(X(g):g\in\mathcal D_0(O)\right)\subset\Sigma_O\subset\sigma\left(X(g):g\in\mathcal S(\bar O)\right)\subset\sigma\left(X(g):g\in\mathcal D_0(O)\right),
\]
so that these $\sigma$-algebras coincide. For $O=\mathbb R^3$ one has $\mathcal S(\bar O)=\mathcal S$, so that $\sigma\left(X(g):g\in\mathcal S(\bar O)\right)=\Sigma$, and therefore $\Sigma_{\mathbb R^3}=\Sigma$.

(iii) Recall that a measurable space is standard Borel if it is isomorphic to a Polish space with its Borel $\sigma$-algebra. Let $\|\cdot\|_{(1)}\le\|\cdot\|_{(2)}\le\cdots$ be norms on $\mathcal S$ which define its topology, and for $n\ge1$ let
\[
Q_n=\left\{\omega\in Q:|\omega(f)|\le n\cdot\|f\|_{(n)}\text{ for every }f\in\mathcal S\right\}.
\]
For $\omega\in Q$, continuity at $0$ gives $C>0$ and $k\ge1$ with $|\omega(f)|\le C\cdot\|f\|_{(k)}$ for all $f\in\mathcal S$, whence $\omega\in Q_n$ for every integer $n\ge\max\{C,k\}$; thus $Q=\bigcup_{n\ge1}Q_n$. Let $\mathbb R^{\mathcal D_0}$ carry the product topology. As $\mathcal D_0$ is countable, $\mathbb R^{\mathcal D_0}$ is a Polish space, and its Borel $\sigma$-algebra is generated by the coordinate maps $z\mapsto z_f$, $f\in\mathcal D_0$. Define
\[
\iota\colon Q\to\mathbb R^{\mathcal D_0},\qquad\iota(\omega)=\left(\omega(f)\right)_{f\in\mathcal D_0}.
\]
By (i) with $O=\mathbb R^3$, the space $\mathcal D_0=\mathcal D_0(\mathbb R^3)$ is dense in $\mathcal S$, and the elements of $Q$ are continuous; so $\iota$ is injective. Since $z_f\circ\iota=X(f)$ for $f\in\mathcal D_0$, part (ii) with $O=\mathbb R^3$ gives
\[
\Sigma=\left\{\iota^{-1}(B):B\subset\mathbb R^{\mathcal D_0}\text{ a Borel set}\right\}.
\]
Next, $\iota(Q_n)$ is the set of the points $z\in\mathbb R^{\mathcal D_0}$ such that
\[
z_{s\cdot f+t\cdot g}=s\cdot z_f+t\cdot z_g\quad\text{and}\quad|z_f|\le n\cdot\|f\|_{(n)}\qquad\text{for all }f,g\in\mathcal D_0\text{ and }s,t\in\mathbb Q.
\]
Every point of $\iota(Q_n)$ satisfies these conditions, by the linearity of the elements of $Q$ and the definition of $Q_n$. Conversely, let $z$ satisfy them. Then $f\mapsto z_f$ is $\mathbb Q$-linear on $\mathcal D_0$, and $|z_f-z_g|=|z_{f-g}|\le n\cdot\|f-g\|_{(n)}$, so that it is uniformly continuous on $\mathcal D_0$ for the norm $\|\cdot\|_{(n)}$. The topology of this norm is coarser than that of $\mathcal S$, so $\mathcal D_0$ is dense in $\mathcal S$ for it as well; hence $f\mapsto z_f$ extends uniquely to a function $\omega\colon\mathcal S\to\mathbb R$ that is continuous for $\|\cdot\|_{(n)}$. By continuity, $\omega$ is $\mathbb Q$-linear and hence real-linear, and $|\omega(f)|\le n\cdot\|f\|_{(n)}$ for every $f\in\mathcal S$. Thus $\omega\in Q_n$ and $\iota(\omega)=z$. Each of the conditions above involves at most three coordinates of $z$ and defines a closed subset of $\mathbb R^{\mathcal D_0}$; therefore $\iota(Q_n)$ is closed, and $\iota(Q)=\bigcup_{n\ge1}\iota(Q_n)$ is an $F_\sigma$ subset, in particular a Borel subset, of $\mathbb R^{\mathcal D_0}$. Together with the injectivity of $\iota$ and the description of $\Sigma$ above, this shows that $\iota$ is an isomorphism of the measurable space $(Q,\Sigma)$ onto $\iota(Q)$ endowed with the trace of the Borel $\sigma$-algebra of $\mathbb R^{\mathcal D_0}$. A Borel subset of a Polish space, endowed with the trace of the Borel $\sigma$-algebra, is a standard Borel space. Hence $(Q,\Sigma)$ is standard Borel.

A measurable space is called a Borel space if it is isomorphic to a Borel subset of $[0,1]$ endowed with the trace of the Borel $\sigma$-algebra of $[0,1]$. Every Polish space with its Borel $\sigma$-algebra is a Borel space. Moreover, a measurable subset of a Borel space, endowed with the trace $\sigma$-algebra, is again a Borel space: if $\kappa$ is an isomorphism of a Borel space $(Y,\mathcal Y)$ onto a Borel set $B\subset[0,1]$ and $A\in\mathcal Y$, then $\kappa(A)$ is a Borel subset of $B$, hence of $[0,1]$, and $\kappa$ restricts to an isomorphism of $A$ onto $\kappa(A)$. Applied to the Polish space $\mathbb R^{\mathcal D_0}$ and to its Borel subset $\iota(Q)$, this shows that $(Q,\Sigma)$ is a Borel space.
\end{proof}

\begin{lemma}[Density of exponentials]\label{lem:density-exponentials}
For $O\in\mathfrak O$ let $\mathcal A(O)$ be the complex linear span of the functions $\exp\left(iX(f)\right)$, $f\in\mathcal D_0(O)$, on $Q$, and let $\mathcal A^0(O)\subset\mathcal A(O)$ be the countable set of the finite linear combinations of these functions with coefficients in $\mathbb Q+i\mathbb Q$. Put $\mathcal A=\mathcal A(\mathbb R^3)$ and $\mathcal A^0=\mathcal A^0(\mathbb R^3)$. Let $\rho\in\mathcal M$.
\begin{enumerate}
\item[(i)] For every $O\in\mathfrak O$ the set $\mathcal A^0(O)$ is dense in $L^2(\Sigma_O,\rho)$.
\item[(ii)] The law $\rho$ is determined by the numbers $\rho\left(\exp\left(iX(f)\right)\right)$, $f\in\mathcal D_0$: if $\rho'\in\mathcal M$ and $\rho'\left(\exp\left(iX(f)\right)\right)=\rho\left(\exp\left(iX(f)\right)\right)$ for every $f\in\mathcal D_0$, then $\rho'=\rho$.
\end{enumerate}
\end{lemma}

\begin{proof}
Let $O\in\mathfrak O$, and let $f_1,f_2,\ldots$ be an enumeration of the countable set $\mathcal D_0(O)$. For $n\ge1$ let
\[
Y_n=\left(X(f_1),\ldots,X(f_n)\right)\colon Q\to\mathbb R^n,\qquad\mathcal F_n=\sigma(Y_n),
\]
and let $\lambda_n$ be the law of $Y_n$ under $\rho$, a probability measure on $\mathbb R^n$. For $t,x\in\mathbb R^n$ write $t\cdot x=\sum_{p=1}^nt_p\cdot x_p$. For $t\in\mathbb Q^n$ the form $g_t=\sum_{p=1}^nt_p\cdot f_p$ lies in $\mathcal D_0(O)$, which is a vector space over $\mathbb Q$ by Lemma~\ref{lem:countable-family-test}(i), and $X(g_t)=t\cdot Y_n$.

(i) The functions $\exp\left(iX(f)\right)$, $f\in\mathcal D_0(O)$, are bounded and $\Sigma_O$-measurable, so $\mathcal A(O)\subset L^2(\Sigma_O,\rho)$. Let $h\in L^2(\Sigma_O,\rho)$ be orthogonal to $\mathcal A(O)$. A $\sigma(Y_n)$-measurable function is a Borel function of $Y_n$; hence $E_\rho[h\mid\mathcal F_n]=\eta_n\circ Y_n$ for a Borel function $\eta_n\colon\mathbb R^n\to\mathbb C$, and $\eta_n\in L^2(\lambda_n)\subset L^1(\lambda_n)$. For $t\in\mathbb Q^n$ the function $\exp\left(iX(g_t)\right)$ is bounded, $\mathcal F_n$-measurable and an element of $\mathcal A(O)$; therefore, by the defining property of the conditional expectation,
\[
\int\lambda_n(\mathrm{d}x)\,\overline{\eta_n(x)}\cdot\exp\left(i\,t\cdot x\right)=\rho\left(\overline{E_\rho[h\mid\mathcal F_n]}\cdot\exp\left(iX(g_t)\right)\right)=\left\langle h,\exp\left(iX(g_t)\right)\right\rangle_\rho=0.
\]
Thus the Fourier transform of the finite complex measure $\overline{\eta_n}\cdot\lambda_n$ on $\mathbb R^n$, which is a continuous function on $\mathbb R^n$, vanishes on the dense set $\mathbb Q^n$, and hence everywhere. By the uniqueness theorem for Fourier transforms of finite complex measures on $\mathbb R^n$, $\eta_n=0$ $\lambda_n$-almost everywhere, that is, $E_\rho[h\mid\mathcal F_n]=0$. The $\sigma$-algebras $\mathcal F_n$ increase with $n$, and $\sigma\left(\bigcup_n\mathcal F_n\right)=\sigma\left(X(f):f\in\mathcal D_0(O)\right)=\Sigma_O$ by Lemma~\ref{lem:countable-family-test}(ii). As $h$ has a $\Sigma_O$-measurable representative, L\'evy's upward theorem gives $E_\rho[h\mid\mathcal F_n]\to E_\rho[h\mid\Sigma_O]=h$ in $L^2(\rho)$. Hence $h=0$, and $\mathcal A(O)$ is dense in $L^2(\Sigma_O,\rho)$. Finally, every element of $\mathcal A(O)$ is a uniform limit on $Q$ of elements of $\mathcal A^0(O)$, since $\left|\exp\left(iX(f)\right)\right|=1$ and the coefficients can be approximated by elements of $\mathbb Q+i\mathbb Q$; so $\mathcal A^0(O)$ is dense in $L^2(\Sigma_O,\rho)$ as well.

(ii) Take $O=\mathbb R^3$ above, so that $\mathcal D_0(O)=\mathcal D_0$, and let $\lambda'_n$ be the law of $Y_n$ under $\rho'$. For $t\in\mathbb Q^n$,
\[
\int\lambda_n(\mathrm{d}x)\,\exp\left(i\,t\cdot x\right)=\rho\left(\exp\left(iX(g_t)\right)\right)=\rho'\left(\exp\left(iX(g_t)\right)\right)=\int\lambda'_n(\mathrm{d}x)\,\exp\left(i\,t\cdot x\right).
\]
The characteristic functions of $\lambda_n$ and $\lambda'_n$ are continuous and coincide on $\mathbb Q^n$, hence on $\mathbb R^n$, and the uniqueness theorem gives $\lambda_n=\lambda'_n$; that is, $\rho$ and $\rho'$ coincide on $\mathcal F_n$, for every $n\ge1$. The family $\bigcup_n\mathcal F_n$ is closed under finite intersections, the $\mathcal F_n$ being increasing, and $\sigma\left(\bigcup_n\mathcal F_n\right)=\sigma\left(X(f):f\in\mathcal D_0\right)=\Sigma$ by Lemma~\ref{lem:countable-family-test}(ii). By Dynkin's $\pi$-$\lambda$ theorem, $\rho'=\rho$.
\end{proof}

\begin{lemma}[Averages over translations]\label{lem:averages-translations}
Let $J\subset\{1,2,3\}$, let $\mathrm{d}a$ denote Lebesgue measure on $\Gamma_J$ in the coordinates $a_l$, $l\in J$ (the unit mass at $0$ if $J=\emptyset$), and put $|b|_1=|b_1|+|b_2|+|b_3|$ for $b\in\mathbb R^3$.
\begin{enumerate}
\item[(i)] The map $\mathbb R^3\times Q\to Q$, $(a,\omega)\mapsto\tau_a\omega$, is measurable, where $\mathbb R^3\times Q$ carries the product of the Borel $\sigma$-algebra of $\mathbb R^3$ and $\Sigma$, and $Q$ carries $\Sigma$. Let $T>0$, and let $F$ be a measurable function on $Q$ which is either bounded and complex valued or takes its values in $[0,\infty]$. Then the function $A^J_TF$ on $Q$ given by
\[
\left(A^J_TF\right)(\omega)=T^{-\#|J|}\cdot\int_{[0,T]^J}\mathrm{d}a\,F(\tau_{-a}\omega),\qquad[0,T]^J=\left\{a\in\Gamma_J:0\le a_l\le T\text{ for }l\in J\right\},
\]
is measurable; and if $F$ is bounded, then for all $b\in\Gamma_J$ and $\omega\in Q$
\[
\left|\left(A^J_TF\right)(\tau_b\omega)-\left(A^J_TF\right)(\omega)\right|\le2|b|_1\cdot T^{-1}\cdot\sup_Q|F|.
\]
\end{enumerate}
Let now $\rho\in\mathcal M$ satisfy $(T_J)$.
\begin{enumerate}
\item[(ii)] The map $a\mapsto U(a)$ is a strongly continuous unitary representation of the group $\Gamma_J$ on $L^2(\rho)$.
\item[(iii)] Let $F\in L^2(\rho)$. For every $T>0$ the function $A^J_TF$ is defined $\rho$-almost everywhere, and $\left\|A^J_TF\right\|_{L^p(\rho)}\le\|F\|_{L^p(\rho)}$ for $1\le p<\infty$; moreover $A^J_TF\to P^J_\rho F$ in $L^2(\rho)$ as $T\to\infty$.
\item[(iv)] $\mathcal E^J(\rho)=L^2(\mathcal I^J,\rho)$ and $P^J_\rho=E_\rho\left[\,\cdot\mid\mathcal I^J\right]$.
\item[(v)] The projections $P^l_\rho$, $l\in J$, commute with each other and with every $U(a)$, $a\in\mathbb R^3$, such that $(\tau_a)_*\rho=\rho$; and $P^J_\rho=\prod_{l\in J}P^l_\rho$.
\end{enumerate}
\end{lemma}

\begin{proof}
(i) The $\sigma$-algebra $\Sigma$ is generated by the functions $X(f)$, $f\in\mathcal S$, and $X(f)(\tau_a\omega)=(\tau_a\omega)(f)=\omega(\tau_{-a}\cdot f)$ (Section~\ref{sec:euclidean-maps-actions}). It therefore suffices to show that, for every $f\in\mathcal S$, the function $\psi\colon\mathbb R^3\times Q\to\mathbb R$, $\psi(a,\omega)=\omega(\tau_{-a}\cdot f)$, is measurable. For fixed $\omega$ it is continuous in $a$, because $a\mapsto\tau_{-a}\cdot f$ is continuous from $\mathbb R^3$ to $\mathcal S$ and $\omega$ is continuous on $\mathcal S$; for fixed $a$ it is the $\Sigma$-measurable function $X(\tau_{-a}\cdot f)$. For $n\ge1$ let
\[
\psi_n(a,\omega)=\sum_{z\in\mathbb Z^3}\mathbf{1}_{z+[0,1)^3}\left(2^n\cdot a\right)\cdot\omega\left(\tau_{-2^{-n}\cdot z}\cdot f\right).
\]
Each term is the product of a Borel function of $a$ and a $\Sigma$-measurable function of $\omega$, and for each $(a,\omega)$ at most one term is different from $0$; hence $\psi_n$ is measurable. If $z\in\mathbb Z^3$ is the point with $2^n\cdot a\in z+[0,1)^3$, then $\left|2^{-n}\cdot z-a\right|_1\le3\cdot2^{-n}$; so the continuity in $a$ gives $\psi_n\to\psi$ pointwise as $n\to\infty$, and $\psi$ is measurable.

Let $F$ be as in (i). The function $(a,\omega)\mapsto F(\tau_{-a}\omega)$ on $\Gamma_J\times Q$ is measurable, as the composition of $F$ with a measurable map, and it is bounded or non-negative; by the theorems of Fubini and Tonelli, $A^J_TF$ is measurable. Now let $F$ be bounded and $b\in\Gamma_J$; write $C=[0,T]^J$, and let $\operatorname{vol}$ denote the measure $\mathrm{d}a$ on $\Gamma_J$. Since $\tau_{-a}\tau_b\omega=\tau_{-(a-b)}\omega$, the substitution $a=a'+b$ gives
\[
\left(A^J_TF\right)(\tau_b\omega)=T^{-\#|J|}\cdot\int_{C-b}\mathrm{d}a'\,F(\tau_{-a'}\omega),
\]
and therefore, with $\triangle$ denoting the symmetric difference,
\[
\left|\left(A^J_TF\right)(\tau_b\omega)-\left(A^J_TF\right)(\omega)\right|\le T^{-\#|J|}\cdot\operatorname{vol}\left(C\,\triangle\,(C-b)\right)\cdot\sup_Q|F|.
\]
A point $a\in C$ lies outside $C-b$ only if $a_l+b_l\notin[0,T]$ for some $l\in J$, and for fixed $l$ the set of these $a\in C$ has volume at most $|b_l|\cdot T^{\#|J|-1}$; hence $\operatorname{vol}\left(C\setminus(C-b)\right)\le|b|_1\cdot T^{\#|J|-1}$. Since $(C-b)\setminus C=\left(C\setminus(C+b)\right)-b$, the same bound, with $-b$ in place of $b$, holds for $\operatorname{vol}\left((C-b)\setminus C\right)$. So $\operatorname{vol}\left(C\,\triangle\,(C-b)\right)\le2|b|_1\cdot T^{\#|J|-1}$ for every $b\in\Gamma_J$, and the bound of (i) follows.

(ii) By $(T_J)$ and Section~\ref{sec:euclidean-maps-actions}, every $U(a)$, $a\in\Gamma_J$, is a unitary operator on $L^2(\rho)$, and $U(a)U(b)=U(a+b)$ because $\tau_a\tau_b=\tau_{a+b}$. Let $f\in\mathcal D_0$ and $F=\exp\left(iX(f)\right)$. By Section~\ref{sec:euclidean-maps-actions}, $U(a)F=\exp\left(iX(\tau_a\cdot f)\right)$; since $|F|=1$ and $X$ is linear,
\[
\left\|U(a)F-F\right\|_{L^2(\rho)}^2=\rho\left(\left|\exp\left(iX(\tau_a\cdot f-f)\right)-1\right|^2\right).
\]
As $a\to0$, $\tau_a\cdot f\to f$ in $\mathcal S$, so that $X(\tau_a\cdot f-f)\to0$ pointwise on $Q$; the integrand is bounded by $4$, and dominated convergence gives $\left\|U(a)F-F\right\|_{L^2(\rho)}\to0$. By linearity, $\left\|U(a)G-G\right\|_{L^2(\rho)}\to0$ as $a\to0$ for every $G\in\mathcal A$. By Lemma~\ref{lem:density-exponentials}(i) with $O=\mathbb R^3$ and by Lemma~\ref{lem:countable-family-test}(ii), $\mathcal A$ is dense in $L^2(\Sigma,\rho)=L^2(\rho)$. Given $G\in L^2(\rho)$ and $\varepsilon>0$, choose $G'\in\mathcal A$ with $\left\|G-G'\right\|_{L^2(\rho)}<\varepsilon$; as the operators $U(a)$ are unitary,
\[
\left\|U(a)G-G\right\|_{L^2(\rho)}\le2\varepsilon+\left\|U(a)G'-G'\right\|_{L^2(\rho)},
\]
which is smaller than $3\varepsilon$ for all $a\in\Gamma_J$ close to $0$. Thus $U(a)\to1$ strongly as $a\to0$, and, by the group law, $a\mapsto U(a)G$ is continuous on $\Gamma_J$ for every $G\in L^2(\rho)$.

(iii) Let $F\in L^2(\rho)$, fix a measurable representative of $F$, and let $1\le p<\infty$. For $a\in\Gamma_J$ the function $\omega\mapsto F(\tau_{-a}\omega)$ represents $U(a)F$, and $(T_J)$ gives $\int\rho(\mathrm{d}\omega)\,|F(\tau_{-a}\omega)|^p=\|F\|_{L^p(\rho)}^p$. By (i) and Tonelli's theorem,
\[
\int\rho(\mathrm{d}\omega)\int_{[0,T]^J}\mathrm{d}a\,|F(\tau_{-a}\omega)|^p=T^{\#|J|}\cdot\|F\|_{L^p(\rho)}^p.
\]
For $p=1$ the right side is finite, $\rho$ being a probability measure. Hence the integral that defines $\left(A^J_TF\right)(\omega)$ converges absolutely for $\rho$-almost every $\omega$; and, applied to the difference of two representatives of $F$, the same identity with $p=1$ shows that $A^J_TF$ depends on the choice of the representative only on a $\rho$-null set. Jensen's inequality for the probability measure $T^{-\#|J|}\cdot\mathrm{d}a$ on $[0,T]^J$ gives
\[
\left|\left(A^J_TF\right)(\omega)\right|^p\le T^{-\#|J|}\cdot\int_{[0,T]^J}\mathrm{d}a\,|F(\tau_{-a}\omega)|^p,
\]
and integration with respect to $\rho$ yields $\left\|A^J_TF\right\|_{L^p(\rho)}\le\|F\|_{L^p(\rho)}$. In particular $A^J_T$ is a linear operator of norm at most $1$ on $L^2(\rho)$. For $G\in L^2(\rho)$ and $a\in\Gamma_J$ the Cauchy--Schwarz inequality and $(T_J)$ give $\int\rho(\mathrm{d}\omega)\,|G(\omega)|\cdot|F(\tau_{-a}\omega)|\le\|G\|_{L^2(\rho)}\cdot\|F\|_{L^2(\rho)}$; so Fubini's theorem applies and yields
\[
\left\langle G,A^J_TF\right\rangle_\rho=T^{-\#|J|}\cdot\int_{[0,T]^J}\mathrm{d}a\,\left\langle G,U(a)F\right\rangle_\rho\qquad\left(F,G\in L^2(\rho)\right).
\]
Let $\mathcal R$ be the closed linear span of the set $\left\{G-U(b)G:G\in L^2(\rho),\ b\in\Gamma_J\right\}$. Since $U(b)^*=U(-b)$, an element $F\in L^2(\rho)$ is orthogonal to $\mathcal R$ if and only if $U(-b)F=F$ for every $b\in\Gamma_J$; thus $\mathcal R^\perp=\mathcal E^J(\rho)$, and $L^2(\rho)=\mathcal E^J(\rho)\oplus\mathcal R$. If $F\in\mathcal E^J(\rho)$, the displayed identity for $\left\langle G,A^J_TF\right\rangle_\rho$ gives $A^J_TF=F$ for every $T>0$. Let $F=G-U(b)G$ with $G\in L^2(\rho)$ and $b\in\Gamma_J$. As $U(a)U(b)=U(a+b)$, the same identity, together with the bound for the volume of $[0,T]^J\,\triangle\,\left([0,T]^J+b\right)$ obtained in the proof of (i) (with $-b$ in place of $b$), gives for every $H\in L^2(\rho)$
\begin{align*}
\left|\left\langle H,A^J_TF\right\rangle_\rho\right|&\le T^{-\#|J|}\cdot\operatorname{vol}\left([0,T]^J\,\triangle\,\left([0,T]^J+b\right)\right)\cdot\|H\|_{L^2(\rho)}\cdot\|G\|_{L^2(\rho)}\\
&\le2|b|_1\cdot T^{-1}\cdot\|H\|_{L^2(\rho)}\cdot\|G\|_{L^2(\rho)}.
\end{align*}
Hence $\left\|A^J_TF\right\|_{L^2(\rho)}\le2|b|_1\cdot T^{-1}\cdot\|G\|_{L^2(\rho)}\to0$ as $T\to\infty$. Since the operators $A^J_T$ have norm at most $1$, it follows that $A^J_TF\to0$ for every $F\in\mathcal R$. For $F\in L^2(\rho)$ one has $F-P^J_\rho F\in\mathcal R$, and therefore
\[
A^J_TF=P^J_\rho F+A^J_T\left(F-P^J_\rho F\right)\to P^J_\rho F\qquad\text{in }L^2(\rho)\text{ as }T\to\infty.
\]

(iv) Let $G$ be an $\mathcal I^J$-measurable function on $Q$, and let $a\in\Gamma_J$ and $\omega\in Q$. The set $\{G=G(\omega)\}$ belongs to $\mathcal I^J$ and contains $\omega$; as $-a\in\Gamma_J$, the map $\tau_{-a}$ carries this set onto itself, so that it contains $\tau_{-a}\omega$, that is, $G(\tau_{-a}\omega)=G(\omega)$. So $U(a)G=G$, and $L^2(\mathcal I^J,\rho)\subset\mathcal E^J(\rho)$.

Conversely, let $F\in\mathcal E^J(\rho)$. By Section~\ref{sec:euclidean-maps-actions} the operators $U(a)$ commute with complex conjugation and with the composition $F\mapsto\varphi\circ F$ by real functions $\varphi$. Hence the real and the imaginary part of $F$ lie in $\mathcal E^J(\rho)$, and for real $F\in\mathcal E^J(\rho)$ so do the functions $F_k=\max\{-k,\min\{F,k\}\}$, $k\ge1$, which converge to $F$ in $L^2(\rho)$ by dominated convergence. Since $L^2(\mathcal I^J,\rho)$ is a closed subspace, it suffices to treat a real bounded $F\in\mathcal E^J(\rho)$; fix a bounded real measurable representative of it, again denoted $F$, and put
\[
\tilde F=\limsup_{n\to\infty}A^J_nF,
\]
where $n$ ranges over the positive integers. By (i), $\tilde F$ is measurable and bounded, and $\tilde F(\tau_b\omega)=\tilde F(\omega)$ for all $b\in\Gamma_J$ and $\omega\in Q$, because
\[
\left|\left(A^J_nF\right)(\tau_b\omega)-\left(A^J_nF\right)(\omega)\right|\le2|b|_1\cdot n^{-1}\cdot\sup_Q|F|\to0\qquad(n\to\infty).
\]
Consequently, for every Borel set $B\subset\mathbb R$, the set $\left\{\tilde F\in B\right\}$ satisfies $\tau_{-b}\left\{\tilde F\in B\right\}=\left\{\tilde F\in B\right\}$ for all $b\in\Gamma_J$, hence lies in $\mathcal I^J$; that is, $\tilde F$ is $\mathcal I^J$-measurable. On the other hand, the set
\[
D=\left\{(a,\omega)\in\Gamma_J\times Q:F(\tau_{-a}\omega)\ne F(\omega)\right\}
\]
is measurable by (i), and for each $a\in\Gamma_J$ the set $\{\omega:(a,\omega)\in D\}$ is a $\rho$-null set, because $U(a)F=F$ in $L^2(\rho)$. By Tonelli's theorem, for $\rho$-almost every $\omega$ the set $\{a\in\Gamma_J:(a,\omega)\in D\}$ is a null set for $\mathrm{d}a$. For such $\omega$ one has $F(\tau_{-a}\omega)=F(\omega)$ for almost every $a\in\Gamma_J$, hence $\left(A^J_nF\right)(\omega)=F(\omega)$ for every $n$, and $\tilde F(\omega)=F(\omega)$. Thus $\tilde F$ is an $\mathcal I^J$-measurable representative of $F$, and $F\in L^2(\mathcal I^J,\rho)$. Finally, for every $\sigma$-algebra $\mathcal G\subset\Sigma$ the orthogonal projection of $L^2(\rho)$ onto $L^2(\mathcal G,\rho)$ is the conditional expectation $E_\rho[\,\cdot\mid\mathcal G]$; with $\mathcal G=\mathcal I^J$ this gives $P^J_\rho=E_\rho\left[\,\cdot\mid\mathcal I^J\right]$.

(v) Let $l\in J$. Since $\Gamma_l\subset\Gamma_J$, the law $\rho$ satisfies $(T_l)$, so that $P^l_\rho$ is defined, and by (iii), with $\{l\}$ in place of $J$, for every $F\in L^2(\rho)$
\[
P^l_\rho F=\lim_{T\to\infty}A^l_TF\quad\text{in }L^2(\rho),\qquad\left(A^l_TF\right)(\omega)=T^{-1}\cdot\int_0^T\mathrm{d}s\,F(\tau_{-s\cdot e_l}\omega).
\]
Let $a\in\mathbb R^3$ with $(\tau_a)_*\rho=\rho$. Then $U(a)$ is unitary on $L^2(\rho)$ (Section~\ref{sec:euclidean-maps-actions}), and $A^l_TU(a)F=U(a)A^l_TF$ because $\tau_{-a}\tau_{-s\cdot e_l}=\tau_{-s\cdot e_l}\tau_{-a}$; letting $T\to\infty$ gives $P^l_\rho U(a)=U(a)P^l_\rho$. Now let $l'\in J$. By what has just been shown, applied to $l'$ and to the translations $\tau_{s\cdot e_l}$, $s\in\mathbb R$, which preserve $\rho$ by $(T_J)$, the projection $P^{l'}_\rho$ commutes with every $U(s\cdot e_l)$. Since $P^{l'}_\rho$ is self-adjoint, the identity for $\left\langle G,A^J_TF\right\rangle_\rho$ in the proof of (iii), with $\{l\}$ in place of $J$, gives
\[
\left\langle G,A^l_TP^{l'}_\rho F\right\rangle_\rho=T^{-1}\cdot\int_0^T\mathrm{d}s\,\left\langle P^{l'}_\rho G,U(s\cdot e_l)F\right\rangle_\rho=\left\langle G,P^{l'}_\rho A^l_TF\right\rangle_\rho\qquad\left(F,G\in L^2(\rho)\right),
\]
so that $A^l_TP^{l'}_\rho=P^{l'}_\rho A^l_T$; letting $T\to\infty$ gives $P^l_\rho P^{l'}_\rho=P^{l'}_\rho P^l_\rho$. A product of commuting orthogonal projections is the orthogonal projection onto the intersection of their ranges, and
\[
\bigcap_{l\in J}\mathcal E^l(\rho)=\mathcal E^J(\rho),
\]
because $U(a)=\prod_{l\in J}U(a_l\cdot e_l)$ for $a\in\Gamma_J$. Hence $P^J_\rho=\prod_{l\in J}P^l_\rho$.
\end{proof}

\begin{lemma}[Clustering criterion and separation of ergodic laws]\label{lem:clustering-criterion-separation}
Let $J\subset\{1,2,3\}$, and let $\mathcal A$ be the complex linear span of the functions $\exp\left(iX(f)\right)$, $f\in\mathcal D_0$, on $Q$, as in Lemma~\ref{lem:density-exponentials}.
\begin{enumerate}
\item[(a)] Let $\rho\in\mathcal M$ satisfy $(T_J)$, and let $\left(a^{(k)}\right)_{k\ge1}$ be a sequence in $\Gamma_J$ such that
\[
\lim_{k\to\infty}\rho\left(\overline{F}\cdot U\left(a^{(k)}\right)G\right)=\overline{\rho(F)}\cdot\rho(G)\qquad\text{for all }F,G\in\mathcal A.
\]
Then $\rho\in\mathcal M^J_{\mathrm{erg}}$.
\item[(b)] Let $n\ge1$ be an integer, and let $\rho_1,\dots,\rho_n\in\mathcal M^J_{\mathrm{erg}}$ be pairwise distinct. Then there exist pairwise disjoint sets $B_1,\dots,B_n\in\mathcal I^J$ such that $\rho_p(B_p)=1$ for $p=1,\dots,n$.
\end{enumerate}
\end{lemma}

\begin{proof}
(a) Let $a\in\Gamma_J$. By $(T_J)$ one has $(\tau_a)_*\rho=\rho$, so that $U(a)$ is a unitary operator of $L^2(\rho)$ (Section~\ref{sec:euclidean-maps-actions}); moreover $\rho\left(\overline{F}\cdot U(a)G\right)=\left\langle F,U(a)G\right\rangle_\rho$ for $F,G\in L^2(\rho)$. For $H\in L^2(\rho)$ one has $\rho(H)=\langle1,H\rangle_\rho$ and $\|1\|_{L^2(\rho)}=1$, $\rho$ being a probability measure; hence $\left|\rho(H)\right|\le\|H\|_{L^2(\rho)}$ by the Cauchy--Schwarz inequality.

The limit relation assumed for $F,G\in\mathcal A$ holds for all $F,G\in L^2(\rho)$. Indeed, the elements of $\mathcal A$ are bounded functions, and $\mathcal A$ contains the set $\mathcal A^0$ of Lemma~\ref{lem:density-exponentials}; so $\mathcal A$ is a dense subspace of $L^2(\rho)$, by Lemma~\ref{lem:density-exponentials}(i) for $O=\mathbb R^3$ and because $\Sigma_{\mathbb R^3}=\Sigma$ (Lemma~\ref{lem:countable-family-test}(ii)). Let $F,G\in L^2(\rho)$ and $\varepsilon>0$. Choose $F',G'\in\mathcal A$ with $\|F-F'\|_{L^2(\rho)}<\varepsilon$ and $\|G-G'\|_{L^2(\rho)}<\varepsilon$, and put
\[
c=\|F\|_{L^2(\rho)}+\|G\|_{L^2(\rho)}+\varepsilon,
\]
so that $\|F'\|_{L^2(\rho)}+\|G\|_{L^2(\rho)}\le c$. For every $a\in\Gamma_J$, the unitarity of $U(a)$ and the Cauchy--Schwarz inequality give
\begin{align*}
\left|\left\langle F,U(a)G\right\rangle_\rho-\left\langle F',U(a)G'\right\rangle_\rho\right|&\le\left|\left\langle F-F',U(a)G\right\rangle_\rho\right|+\left|\left\langle F',U(a)(G-G')\right\rangle_\rho\right|\\
&\le\|F-F'\|_{L^2(\rho)}\cdot\|G\|_{L^2(\rho)}+\|F'\|_{L^2(\rho)}\cdot\|G-G'\|_{L^2(\rho)}\le\varepsilon\cdot c,
\end{align*}
and the bound $\left|\rho(H)\right|\le\|H\|_{L^2(\rho)}$, together with the linearity of $\rho$, gives
\[
\left|\overline{\rho(F)}\cdot\rho(G)-\overline{\rho(F')}\cdot\rho(G')\right|\le\left|\rho(F-F')\right|\cdot\left|\rho(G)\right|+\left|\rho(F')\right|\cdot\left|\rho(G-G')\right|\le\varepsilon\cdot c.
\]
By hypothesis, $\left\langle F',U\left(a^{(k)}\right)G'\right\rangle_\rho=\rho\left(\overline{F'}\cdot U\left(a^{(k)}\right)G'\right)$ tends to $\overline{\rho(F')}\cdot\rho(G')$ as $k\to\infty$. With the two preceding bounds for $a=a^{(k)}$ and the triangle inequality, it follows that
\[
\limsup_{k\to\infty}\left|\left\langle F,U\left(a^{(k)}\right)G\right\rangle_\rho-\overline{\rho(F)}\cdot\rho(G)\right|\le2\varepsilon\cdot c.
\]
Here $\varepsilon>0$ is arbitrary, and $2\varepsilon\cdot c\to0$ as $\varepsilon\to0$. Therefore
\[
\lim_{k\to\infty}\left\langle F,U\left(a^{(k)}\right)G\right\rangle_\rho=\overline{\rho(F)}\cdot\rho(G)\qquad\text{for all }F,G\in L^2(\rho).
\]

Now let $G\in\mathcal E^J(\rho)$. Since $a^{(k)}\in\Gamma_J$, one has $U\left(a^{(k)}\right)G=G$ for every $k\ge1$, and the last display yields, for every $F\in L^2(\rho)$,
\[
\langle F,G\rangle_\rho=\lim_{k\to\infty}\left\langle F,U\left(a^{(k)}\right)G\right\rangle_\rho=\overline{\rho(F)}\cdot\rho(G)=\left\langle F,\rho(G)\cdot1\right\rangle_\rho.
\]
Taking $F=G-\rho(G)\cdot1$ one obtains $\left\|G-\rho(G)\cdot1\right\|_{L^2(\rho)}^2=0$, that is, $G=\rho(G)\cdot1$. Hence $\mathcal E^J(\rho)\subset\mathbb C\cdot1$. Conversely, every constant function on $Q$ is invariant under each operator $U(a)=\hat\tau_a$, by the definition of $\hat\tau_a$ in Section~\ref{sec:euclidean-maps-actions}. Thus $\mathcal E^J(\rho)=\mathbb C\cdot1$, so that $\rho$ is $\Gamma_J$-ergodic; as $\rho$ satisfies $(T_J)$, it belongs to $\mathcal M^J_{\mathrm{erg}}$.

(b) Let $p,q\in\{1,\dots,n\}$ with $p\ne q$. The laws $\rho_p$ and $\rho_q$ are different, so by Lemma~\ref{lem:density-exponentials}(ii) there is $f_{pq}\in\mathcal D_0$ such that
\[
\rho_p\left(\exp\left(iX(f_{pq})\right)\right)\ne\rho_q\left(\exp\left(iX(f_{pq})\right)\right).
\]
The function $\exp\left(iX(f_{pq})\right)$ belongs to the countable set $\mathcal A^0$ of Lemma~\ref{lem:density-exponentials}. Let $\mathcal F\subset\mathcal A^0$ be the finite set of the functions $\exp\left(iX(f_{pq})\right)$, where $(p,q)$ ranges over the ordered pairs of distinct elements of $\{1,\dots,n\}$. Then, for any two distinct $p,q\in\{1,\dots,n\}$, some $F\in\mathcal F$ satisfies $\rho_p(F)\ne\rho_q(F)$.

Every $F\in\mathcal F$ is a bounded $\Sigma$-measurable function on $Q$. For $T>0$ the integral that defines $\left(A^J_TF\right)(\omega)$ in Lemma~\ref{lem:averages-translations}(i) therefore converges at every point $\omega\in Q$; the function $A^J_TF$ so obtained is $\Sigma$-measurable, satisfies $\left|A^J_TF\right|\le\sup_Q|F|$, and, for each $p$, represents the element $A^J_TF$ of $L^2(\rho_p)$ of Lemma~\ref{lem:averages-translations}(iii). Let $p\in\{1,\dots,n\}$. The law $\rho_p$ satisfies $(T_J)$, and $\mathcal E^J(\rho_p)=\mathbb C\cdot1$. As $\langle1,1\rangle_{\rho_p}=1$, the orthogonal projection $P^J_{\rho_p}$ of $L^2(\rho_p)$ onto $\mathbb C\cdot1$ is the operator $H\mapsto\langle1,H\rangle_{\rho_p}\cdot1=\rho_p(H)\cdot1$. Lemma~\ref{lem:averages-translations}(iii), applied to $\rho_p$, now gives
\[
\lim_{T\to\infty}\left\|A^J_TF-\rho_p(F)\cdot1\right\|_{L^2(\rho_p)}=0\qquad\text{for every }F\in\mathcal F.
\]
Since the set $\{1,\dots,n\}\times\mathcal F$ is finite, there are real numbers $T_k\ge k$, $k=1,2,\dots$, such that
\[
\left\|A^J_{T_k}F-\rho_p(F)\cdot1\right\|_{L^2(\rho_p)}\le2^{-k}\qquad\text{for all }k\ge1,\ p\in\{1,\dots,n\}\text{ and }F\in\mathcal F.
\]
Fix $p\in\{1,\dots,n\}$ and $F\in\mathcal F$, and let $h_k=\left|A^J_{T_k}F-\rho_p(F)\right|$, a non-negative $\Sigma$-measurable function on $Q$. By the Cauchy--Schwarz inequality, $\rho_p(h_k)\le\left\|A^J_{T_k}F-\rho_p(F)\cdot1\right\|_{L^2(\rho_p)}\le2^{-k}$, and the monotone convergence theorem gives
\[
\rho_p\left(\sum_{k\ge1}h_k\right)=\sum_{k\ge1}\rho_p(h_k)\le1.
\]
Consequently $\sum_{k\ge1}h_k(\omega)<\infty$, and in particular $h_k(\omega)\to0$ as $k\to\infty$, for $\rho_p$-almost every $\omega\in Q$. Thus, for every $p\in\{1,\dots,n\}$ and every $F\in\mathcal F$,
\[
\lim_{k\to\infty}\left(A^J_{T_k}F\right)(\omega)=\rho_p(F)\qquad\text{for $\rho_p$-almost every }\omega\in Q.
\]
For $p=1,\dots,n$ let
\[
B_p=\left\{\omega\in Q:\lim_{k\to\infty}\left(A^J_{T_k}F\right)(\omega)=\rho_p(F)\text{ for every }F\in\mathcal F\right\}.
\]
These sets have the required properties.

\emph{Measurability.} Each function $A^J_{T_k}F$ is $\Sigma$-measurable, and the set of the points at which a sequence of $\Sigma$-measurable complex functions converges to a given complex number belongs to $\Sigma$. As $\mathcal F$ is finite, $B_p\in\Sigma$.

\emph{Full measure.} The set $B_p$ is the intersection of the finitely many sets $\left\{\omega\in Q:\lim_{k\to\infty}\left(A^J_{T_k}F\right)(\omega)=\rho_p(F)\right\}$, $F\in\mathcal F$, each of which belongs to $\Sigma$ and, by the convergence $\rho_p$-almost everywhere established above, has $\rho_p$-measure $1$. Hence $\rho_p(B_p)=1$.

\emph{Invariance.} Let $b\in\Gamma_J$ and $\omega\in Q$. By Lemma~\ref{lem:averages-translations}(i), with $|b|_1=|b_1|+|b_2|+|b_3|$, one has for $F\in\mathcal F$ and $k\ge1$
\[
\left|\left(A^J_{T_k}F\right)(\tau_b\omega)-\left(A^J_{T_k}F\right)(\omega)\right|\le2|b|_1\cdot T_k^{-1}\cdot\sup_Q|F|,
\]
and the right side tends to $0$ as $k\to\infty$, because $T_k\ge k$. Hence the sequence $\left(\left(A^J_{T_k}F\right)(\tau_b\omega)\right)_{k\ge1}$ converges to $\rho_p(F)$ if and only if the sequence $\left(\left(A^J_{T_k}F\right)(\omega)\right)_{k\ge1}$ does; that is, $\tau_b\omega\in B_p$ if and only if $\omega\in B_p$. Since $\omega\mapsto\tau_b\omega$ is a bijection of $Q$ onto itself, it follows that $\tau_bB_p=B_p$. This holds for every $b\in\Gamma_J$, and $B_p\in\Sigma$; therefore $B_p\in\mathcal I^J$.

\emph{Disjointness.} Let $p\ne q$, and choose $F\in\mathcal F$ with $\rho_p(F)\ne\rho_q(F)$. At a point $\omega\in B_p\cap B_q$ the sequence $\left(\left(A^J_{T_k}F\right)(\omega)\right)_{k\ge1}$ would converge both to $\rho_p(F)$ and to $\rho_q(F)$, which is impossible. Hence $B_p\cap B_q=\emptyset$.
\end{proof}

\begin{lemma}[Consequences of the Gaussian bound]\label{lem:consequences-gaussian-bound}
Let $\rho\in\mathcal M$ satisfy~\eqref{eq:setup-G}.
\begin{enumerate}
\item[(i)] For every $f\in\mathcal S$,
\[
\rho\left(\exp|X(f)|\right)\le2\exp\left(\frac{1}{2}\|f\|^2\right),\qquad\rho\left(X(f)\right)=0,\qquad\rho\left(X(f)^2\right)\le\|f\|^2,
\]
and, for every integer $n\ge1$,
\[
\rho\left(|X(f)|^n\right)\le n\cdot2^{\frac{n}{2}}\cdot\Gamma\left(\frac{n}{2}\right)\cdot\|f\|^n.
\]
\item[(ii)] The map $f\mapsto X(f)$ from $\mathcal S$ to $L^2(\rho)$ has exactly one extension to a linear map $X\colon L^2(\mathbb R^3;\Lambda^2)\to L^2(\rho)$ of norm at most $1$, where $L^2(\mathbb R^3;\Lambda^2)$ is the real Hilbert space of Section~\ref{sec:test-forms-distributions}. For this extension, the inequality in~\eqref{eq:setup-G} and the assertions of (i) hold for every $f\in L^2(\mathbb R^3;\Lambda^2)$. If $O\in\mathfrak O$ and $f\in L^2(\mathbb R^3;\Lambda^2)$ vanishes almost everywhere outside $O$, then $X(f)\in L^2(\Sigma_O,\rho)$.
\item[(iii)] Let $O\in\mathfrak O$. The complex polynomials in the functions $X(f)$, $f\in\mathcal D(O)$, form a dense subspace of $L^2(\Sigma_O,\rho)$. Here a complex polynomial in a family of functions on $Q$ is a linear combination, with complex coefficients, of finitely many finite products of members of the family; the empty product is the constant function $1$.
\end{enumerate}
\end{lemma}

\begin{proof}
(i) Let $f\in\mathcal S$ and $\sigma=\|f\|$. The map $X$ is linear, so that $X(-f)=-X(f)$, and $\exp|x|\le\exp(x)+\exp(-x)$ for $x\in\mathbb R$. The inequality~\eqref{eq:setup-G} for $f$ and for $-f$ therefore gives
\[
\rho\left(\exp|X(f)|\right)\le\rho\left(\exp X(f)\right)+\rho\left(\exp X(-f)\right)\le2\exp\left(\frac{\sigma^2}{2}\right).
\]
By this bound for $tf$ in place of $f$, the function $\exp\left(|t|\cdot|X(f)|\right)$ is $\rho$-integrable for every $t\in\mathbb R$. Hence
\[
\varphi\colon\mathbb R\to(0,\infty),\qquad\varphi(t)=\rho\left(\exp\left(t\cdot X(f)\right)\right),
\]
is well defined, and $\varphi$ is smooth with $\varphi^{(k)}(t)=\rho\left(X(f)^k\cdot\exp\left(t\cdot X(f)\right)\right)$ for $k\ge0$; differentiation under the integral sign is permitted by dominated convergence, because $|x|^k\cdot\exp(s|x|)\le k!\cdot\exp\left((s+1)\cdot|x|\right)$ for $x\in\mathbb R$ and $s\ge0$. In particular $\varphi(0)=1$, $\varphi'(0)=\rho\left(X(f)\right)$ and $\varphi''(0)=\rho\left(X(f)^2\right)$. Taylor's formula at $t=0$ and~\eqref{eq:setup-G} for $tf$ give, as $t\to0$,
\[
1+t\cdot\rho\left(X(f)\right)+\frac{t^2}{2}\cdot\rho\left(X(f)^2\right)+o(t^2)=\varphi(t)\le\exp\left(\frac{t^2\sigma^2}{2}\right)=1+\frac{t^2}{2}\cdot\sigma^2+o(t^2).
\]
Subtracting $1$, dividing by $t>0$ and letting $t\to0$ one obtains $\rho\left(X(f)\right)\le0$; dividing instead by $t<0$, which reverses the inequality, and letting $t\to0$ one obtains $\rho\left(X(f)\right)\ge0$. Thus $\rho\left(X(f)\right)=0$, and then dividing by $t^2$ and letting $t\to0$ gives $\rho\left(X(f)^2\right)\le\sigma^2$.

Let $n\ge1$. If $\sigma=0$, then $\rho\left(X(f)^2\right)=0$, so $X(f)=0$ $\rho$-almost everywhere, and both sides of the asserted bound on $\rho\left(|X(f)|^n\right)$ vanish. Let $\sigma>0$ and $u>0$. For every $t>0$, Chebyshev's inequality for the function $\exp\left(t\cdot X(f)\right)$ and~\eqref{eq:setup-G} for $tf$ give
\[
\rho\left(\{X(f)>u\}\right)\le\exp(-tu)\cdot\rho\left(\exp\left(t\cdot X(f)\right)\right)\le\exp\left(-tu+\frac{t^2\sigma^2}{2}\right).
\]
The infimum of the right side over $t>0$ is attained at $t=u/\sigma^2$ and equals $\exp\left(-\frac{u^2}{2\sigma^2}\right)$. The same argument with $-f$ in place of $f$ bounds $\rho\left(\{-X(f)>u\}\right)$ by the same number, so that $\rho\left(\{|X(f)|>u\}\right)\le2\exp\left(-\frac{u^2}{2\sigma^2}\right)$ for every $u>0$. By Tonelli's theorem,
\begin{align*}
\rho\left(|X(f)|^n\right)&=\int_0^{\infty}\mathrm{d}u\,n\cdot u^{n-1}\cdot\rho\left(\{|X(f)|>u\}\right)\le2n\int_0^{\infty}\mathrm{d}u\,u^{n-1}\cdot\exp\left(-\frac{u^2}{2\sigma^2}\right)\\
&=n\cdot2^{\frac{n}{2}}\cdot\Gamma\left(\frac{n}{2}\right)\cdot\sigma^n,
\end{align*}
where the last integral is evaluated by the substitution $u=\sigma\sqrt{2v}$, $v>0$, and the definition of the Gamma function.

(ii) For $f\in\mathcal S$, (i) gives $\|X(f)\|_{L^2(\rho)}=\rho\left(X(f)^2\right)^{\frac{1}{2}}\le\|f\|$. Thus $f\mapsto X(f)$ is a linear map from $\mathcal S$, normed by $\|\cdot\|$, to $L^2(\rho)$, of norm at most $1$. As $\mathcal S$ is dense in $L^2(\mathbb R^3;\Lambda^2)$, this map has exactly one continuous extension to $L^2(\mathbb R^3;\Lambda^2)$; the extension is linear and of norm at most $1$, and its values are classes of real functions, being limits in $L^2(\rho)$ of real functions. Let $f\in L^2(\mathbb R^3;\Lambda^2)$, and let $f_n\in\mathcal S$ satisfy $\|f_n-f\|\to0$. Then $X(f_n)\to X(f)$ in $L^2(\rho)$, and some subsequence $(f_{n_k})_{k\ge1}$ satisfies $X(f_{n_k})\to X(f)$ $\rho$-almost everywhere. By Fatou's lemma and~\eqref{eq:setup-G} for $f_{n_k}$,
\[
\rho\left(\exp X(f)\right)\le\liminf_{k\to\infty}\rho\left(\exp X(f_{n_k})\right)\le\lim_{k\to\infty}\exp\left(\frac{1}{2}\|f_{n_k}\|^2\right)=\exp\left(\frac{1}{2}\|f\|^2\right).
\]
So the inequality in~\eqref{eq:setup-G} holds for every $f\in L^2(\mathbb R^3;\Lambda^2)$. The proof of (i) rests only on the linearity of $X$ and on this inequality for the multiples $tf$, $t\in\mathbb R$; it therefore applies without change to every $f\in L^2(\mathbb R^3;\Lambda^2)$, and (i) holds for such $f$.

Now let $O\in\mathfrak O$, and let $f\in L^2(\mathbb R^3;\Lambda^2)$ vanish almost everywhere outside $O$. The space $\mathcal D(O)$ is dense in the closed subspace $L^2(O;\Lambda^2)$ of $L^2(\mathbb R^3;\Lambda^2)$ formed by the maps that vanish almost everywhere outside $O$; choose $f_n\in\mathcal D(O)$ with $\|f_n-f\|\to0$. Each $X(f_n)$ is $\Sigma_O$-measurable, by the definition of $\Sigma_O$, and square integrable, by (i); hence $X(f_n)\in L^2(\Sigma_O,\rho)$. Since $X(f_n)\to X(f)$ in $L^2(\rho)$ and $L^2(\Sigma_O,\rho)$ is closed in $L^2(\rho)$, also $X(f)\in L^2(\Sigma_O,\rho)$.

(iii) Let $O\in\mathfrak O$. For $f\in\mathcal D(O)$ the function $X(f)$ is $\Sigma_O$-measurable and, by (i), has finite moments of all orders. Hence, by H\"older's inequality, every complex polynomial in the functions $X(f)$, $f\in\mathcal D(O)$, belongs to $L^2(\Sigma_O,\rho)$. Let $h\in L^2(\Sigma_O,\rho)$ be orthogonal to all these polynomials, and let $f\in\mathcal D_0(O)$. For $N\ge0$ put
\[
s_N=\sum_{k=0}^N\frac{1}{k!}\cdot\left(iX(f)\right)^k.
\]
Each $s_N$ is such a polynomial, $|s_N|\le\exp|X(f)|$, and $s_N\to\exp\left(iX(f)\right)$ pointwise on $Q$ as $N\to\infty$. The function $\exp|X(f)|$ lies in $L^2(\rho)$, because its square $\exp|X(2f)|$ is integrable by (i). By dominated convergence, with the integrable majorant $4\exp|X(2f)|$ of $\left|s_N-\exp\left(iX(f)\right)\right|^2$, the functions $s_N$ converge to $\exp\left(iX(f)\right)$ in $L^2(\rho)$. Therefore
\[
\left\langle h,\exp\left(iX(f)\right)\right\rangle_{\rho}=\lim_{N\to\infty}\langle h,s_N\rangle_{\rho}=0.
\]
So $h$ is orthogonal to $\mathcal A(O)$, and in particular to $\mathcal A^0(O)$, which is dense in $L^2(\Sigma_O,\rho)$ by Lemma~\ref{lem:density-exponentials}(i). Hence $h=0$. The orthogonal complement in $L^2(\Sigma_O,\rho)$ of the space of these polynomials is thus $\{0\}$, which proves (iii).
\end{proof}

\begin{lemma}[Reformulations of reflection positivity]\label{lem:reformulations-reflection-positivity}
Let $j\in\{1,2,3\}$, and let $\rho\in\mathcal M$ satisfy $(R_j)$. The following conditions are equivalent.
\begin{enumerate}
\item[(1)] $\rho$ satisfies $(RP_j)$.
\item[(2)] $\rho\left(\overline{\Theta_jG}\cdot G\right)\ge0$ for every $G\in\mathcal A^0(\{x_j>0\})$.
\item[(3)] $\langle\Theta_jG,G\rangle_{\rho}\ge0$ for every $G\in L^2(\Sigma_j^{+},\rho)$.
\end{enumerate}
If $\rho$ satisfies in addition~\eqref{eq:setup-G}, then (1)--(3) are also equivalent to
\begin{enumerate}
\item[(4)] $\langle\Theta_jG,G\rangle_{\rho}\ge0$ for every complex polynomial $G$ in the functions $X(f)$ with $f\in\mathcal S$ and $\operatorname{supp}f\subset\{x_j\ge0\}$.
\end{enumerate}
If $\rho$ satisfies in addition $(T_j)$, then (1)--(3) are also equivalent to
\begin{enumerate}
\item[(5)] $\rho$ is reflection positive, in the sense of Section~\ref{sec:properties-laws}, at every plane $\{x_j=c\}$ with $c\in\mathbb R$.
\end{enumerate}
\end{lemma}

\begin{proof}
Let $n\ge1$, $c_1,\dots,c_n\in\mathbb C$, $f_1,\dots,f_n\in\mathcal S$, and
\[
G=\sum_{p=1}^nc_p\cdot\exp\left(iX(f_p)\right).
\]
The operator $\Theta_j=\hat\theta_j$ is multiplicative, commutes with complex conjugation and with composition, and maps $X(f)$ to $X(\theta_j\cdot f)$ (Section~\ref{sec:euclidean-maps-actions}). Hence, at every point of $Q$,
\[
\overline{\Theta_jG}\cdot G=\sum_{p,q=1}^n\overline{c_p}\cdot c_q\cdot\exp\left(i\left(X(f_q)-X(\theta_j\cdot f_p)\right)\right),
\]
and $\rho\left(\overline{\Theta_jG}\cdot G\right)=\langle\Theta_jG,G\rangle_{\rho}$ by the definition of the scalar product. Let now $f_1,\dots,f_n\in\mathcal D(\{x_j>0\})$. Then $G$ is bounded and $\Sigma_j^{+}$-measurable (Lemma~\ref{lem:countable-family-test}(ii)), so that $G\in L^2(\Sigma_j^{+},\rho)$, and the sum in the definition of $(RP_j)$ (Section~\ref{sec:properties-laws}) equals $\langle\Theta_jG,G\rangle_{\rho}$. Thus (1) states the inequality of (3) for the functions $G$ of this form; in particular (3) implies (1). The elements of $\mathcal A^0(\{x_j>0\})$ are functions of this form, with $c_p\in\mathbb Q+i\mathbb Q$ and $f_p\in\mathcal D_0(\{x_j>0\})\subset\mathcal D(\{x_j>0\})$; hence (1) implies (2).

Assume (2). By $(R_j)$ and Section~\ref{sec:euclidean-maps-actions}, $\Theta_j$ is a unitary operator of $L^2(\rho)$. Hence $G\mapsto\langle\Theta_jG,G\rangle_{\rho}$ is continuous on $L^2(\rho)$, and the set of the $G\in L^2(\Sigma_j^{+},\rho)$ with $\langle\Theta_jG,G\rangle_{\rho}\ge0$ is closed in $L^2(\Sigma_j^{+},\rho)$. By (2) this set contains $\mathcal A^0(\{x_j>0\})$, which is dense in $L^2(\Sigma_j^{+},\rho)$ by Lemma~\ref{lem:density-exponentials}(i) for $O=\{x_j>0\}$. So the set is all of $L^2(\Sigma_j^{+},\rho)$, which is (3).

Let $\rho$ satisfy in addition~\eqref{eq:setup-G}. Let $f\in\mathcal S$ with $\operatorname{supp}f\subset\{x_j\ge0\}$, the closure of $\{x_j>0\}$. By Lemma~\ref{lem:countable-family-test}(ii) the function $X(f)$ is $\Sigma_j^{+}$-measurable, and by Lemma~\ref{lem:consequences-gaussian-bound}(i) it has finite moments of all orders. By H\"older's inequality every complex polynomial in such functions lies in $L^2(\Sigma_j^{+},\rho)$; so (3) implies (4). Conversely, these polynomials include the complex polynomials in the functions $X(f)$, $f\in\mathcal D(\{x_j>0\})$, which are dense in $L^2(\Sigma_j^{+},\rho)$ by Lemma~\ref{lem:consequences-gaussian-bound}(iii). By the continuity of $G\mapsto\langle\Theta_jG,G\rangle_{\rho}$ on $L^2(\rho)$, (4) implies (3).

Let $\rho$ satisfy in addition $(T_j)$; it remains to show that (1) and (5) are equivalent. Let $c\in\mathbb R$ and $b=c\cdot e_j$. By Section~\ref{sec:euclidean-maps-actions} the reflection in the plane $\{x_j=c\}$ is $\theta_{j,c}=\tau_{2b}\theta_j$. As $\theta_jb=-b$, one has $\tau_b\theta_j\tau_{-b}(x)=\theta_jx+2b$ for $x\in\mathbb R^3$, that is, $\theta_{j,c}=\tau_b\theta_j\tau_{-b}$. Moreover $\tau_b$ maps the half space $\{x_j>0\}$ onto $\{x_j>c\}$. Consequently $f\mapsto\tau_{-b}\cdot f$ is a bijection of $\mathcal D(\{x_j>c\})$ onto $\mathcal D(\{x_j>0\})$, and $g=\tau_{-b}\cdot f$ satisfies
\[
f=\tau_b\cdot g,\qquad\theta_{j,c}\cdot f=\tau_b\cdot(\theta_j\cdot g).
\]
Let $k\ge1$, $z_1,\dots,z_k\in\mathbb C$, $f_1,\dots,f_k\in\mathcal D(\{x_j>c\})$, and $g_p=\tau_{-b}\cdot f_p$ for $p=1,\dots,k$. Since $U(b)X(h)=X(\tau_b\cdot h)$ for $h\in\mathcal S$, and $U(b)=\hat\tau_b$ is multiplicative and commutes with composition (Section~\ref{sec:euclidean-maps-actions}),
\[
\exp\left(i\left(X(f_q)-X(\theta_{j,c}\cdot f_p)\right)\right)=U(b)\exp\left(i\left(X(g_q)-X(\theta_j\cdot g_p)\right)\right)\qquad(p,q=1,\dots,k).
\]
By $(T_j)$ one has $(\tau_b)_{*}\rho=\rho$, hence $\rho\left(U(b)F\right)=\rho(F)$ for every bounded measurable function $F$ on $Q$ (Section~\ref{sec:euclidean-maps-actions}). Therefore
\begin{align*}
&\sum_{p,q=1}^k\overline{z_p}\cdot z_q\cdot\rho\left(\exp\left(i\left(X(f_q)-X(\theta_{j,c}\cdot f_p)\right)\right)\right)\\
&\qquad=\sum_{p,q=1}^k\overline{z_p}\cdot z_q\cdot\rho\left(\exp\left(i\left(X(g_q)-X(\theta_j\cdot g_p)\right)\right)\right).
\end{align*}
When $(f_1,\dots,f_k)$ ranges over $\mathcal D(\{x_j>c\})^k$, the family $(g_1,\dots,g_k)$ ranges over $\mathcal D(\{x_j>0\})^k$. Hence $\rho$ is reflection positive at the plane $\{x_j=c\}$ if and only if it satisfies $(RP_j)$. Since $c\in\mathbb R$ was arbitrary, (1) and (5) are equivalent.
\end{proof}

\begin{lemma}[Measurability of the defining properties]\label{lem:measurability-defining-properties}
Let $\emptyset\ne J\subset\{1,2,3\}$ and $j\in\{1,2,3\}$. The set of the laws $\rho\in\mathcal M$ that satisfy $(T_J)$, the sets of the laws $\rho\in\mathcal M$ that satisfy $(R_j)$, $(RP_j)$ and~\eqref{eq:setup-G}, respectively, and the set $\mathcal M^J_{\mathrm{erg}}$ belong to $\mathcal B_{\mathcal M}$. More precisely, the following holds for every $\rho\in\mathcal M$, and each of the conditions stated consists of countably many relations.
\begin{itemize}
\item $\rho$ satisfies $(T_J)$ if and only if $\rho\left(\exp\left(iX(\tau_a\cdot f)\right)\right)=\rho\left(\exp\left(iX(f)\right)\right)$ for all $a\in\Gamma_J\cap\mathbb Q^3$ and all $f\in\mathcal D_0$.
\item $\rho$ satisfies $(R_j)$ if and only if $\rho\left(\exp\left(iX(\theta_j\cdot f)\right)\right)=\rho\left(\exp\left(iX(f)\right)\right)$ for all $f\in\mathcal D_0$.
\item $\rho$ satisfies $(RP_j)$ if and only if it satisfies condition (2) of Lemma~\ref{lem:reformulations-reflection-positivity}, that is, $\rho\left(\overline{\Theta_jG}\cdot G\right)\ge0$ for every $G\in\mathcal A^0(\{x_j>0\})$; this equivalence holds whether or not $\rho$ satisfies $(R_j)$.
\item $\rho$ satisfies~\eqref{eq:setup-G} if and only if $\rho\left(\exp X(f)\right)\le\exp\left(\frac{1}{2}\|f\|^2\right)$ for all $f\in\mathcal D_0$.
\item If $\rho$ satisfies $(T_J)$, then $\rho\in\mathcal M^J_{\mathrm{erg}}$ if and only if
\[
\lim_{n\to\infty}\rho\left(\overline{F}\cdot A^J_nG\right)=\overline{\rho(F)}\cdot\rho(G)\qquad\text{for all }F,G\in\mathcal A^0,
\]
where $n$ tends to infinity through the positive integers and $A^J_n$ is the average $A^J_T$ of Lemma~\ref{lem:averages-translations}(i) with $T=n$.
\end{itemize}
\end{lemma}

\begin{proof}
Let $\Phi$ be a bounded $\Sigma$-measurable function on $Q$. Then $\rho\mapsto\rho(\Phi)$ is a $\mathcal B_{\mathcal M}$-measurable function on $\mathcal M$. Indeed, for the indicator function of a set in $\Sigma$ this is the definition of $\mathcal B_{\mathcal M}$; the real bounded $\Sigma$-measurable functions with this property form a vector space which contains the constants and, by monotone convergence, is closed under bounded increasing limits; so the monotone class theorem gives the assertion for real $\Phi$, and complex $\Phi$ are treated through their real and imaginary parts. For a $\Sigma$-measurable function $\Phi\colon Q\to[0,\infty]$, the map $\rho\mapsto\rho(\Phi)\in[0,\infty]$ is $\mathcal B_{\mathcal M}$-measurable as the increasing limit of the maps $\rho\mapsto\rho\left(\min\{\Phi,k\}\right)$, $k\in\mathbb N$.

The conditions on the right of the equivalences in the statement are indexed by the countable sets $\Gamma_J\cap\mathbb Q^3$, $\mathcal D_0$, $\mathcal A^0(\{x_j>0\})$ and $\mathcal A^0\times\mathcal A^0$ (Lemmas~\ref{lem:countable-family-test}(i) and~\ref{lem:density-exponentials}). Each single condition is an equality between, or an inequality for, numbers $\rho(\Phi)$ with $\Phi$ bounded or non-negative and $\Sigma$-measurable, or, in the last item, the convergence of a sequence of numbers $\rho(\Phi_n)$ of this kind to the number $\overline{\rho(F)}\cdot\rho(G)$, which is a $\mathcal B_{\mathcal M}$-measurable function of $\rho$ as well; there $\Phi_n=\overline{F}\cdot A^J_nG$ is bounded and $\Sigma$-measurable by Lemma~\ref{lem:averages-translations}(i), applied to the real and the imaginary part of $G$. Since the set of points at which a sequence of measurable functions converges to a given measurable function is measurable, each single condition defines a set in $\mathcal B_{\mathcal M}$, and so does each of the countable families. Once the equivalences are proved, it follows that the sets of the laws with $(T_J)$, $(R_j)$, $(RP_j)$ and~\eqref{eq:setup-G} belong to $\mathcal B_{\mathcal M}$, and so does $\mathcal M^J_{\mathrm{erg}}$, which is the intersection of the set of the laws with $(T_J)$ and the set defined by the condition of the last item.

The property $(T_J)$. For $a\in\mathbb R^3$ and $f\in\mathcal S$ one has $U(a)\exp\left(iX(f)\right)=\exp\left(iX(\tau_a\cdot f)\right)$ and $\left(U(a)F\right)(\omega)=F(\tau_{-a}\omega)$ (Section~\ref{sec:euclidean-maps-actions}); hence
\[
\rho\left(\exp\left(iX(\tau_a\cdot f)\right)\right)=\left((\tau_{-a})_{*}\rho\right)\left(\exp\left(iX(f)\right)\right).
\]
If $\rho$ satisfies $(T_J)$ and $a\in\Gamma_J$, then $-a\in\Gamma_J$, and the right side equals $\rho\left(\exp\left(iX(f)\right)\right)$. Conversely, assume the equalities of the statement. For $f\in\mathcal D_0$ the function $a\mapsto\rho\left(\exp\left(iX(\tau_a\cdot f)\right)\right)$ is continuous on $\mathbb R^3$, by dominated convergence, because $a\mapsto\tau_a\cdot f$ is continuous from $\mathbb R^3$ to $\mathcal S$ and every $\omega\in Q$ is continuous on $\mathcal S$. As $\Gamma_J\cap\mathbb Q^3$ is dense in $\Gamma_J$, the equalities hold for all $a\in\Gamma_J$, that is, $\left((\tau_{-a})_{*}\rho\right)\left(\exp\left(iX(f)\right)\right)=\rho\left(\exp\left(iX(f)\right)\right)$ for all $a\in\Gamma_J$ and $f\in\mathcal D_0$. By Lemma~\ref{lem:density-exponentials}(ii), $(\tau_{-a})_{*}\rho=\rho$ for every $a\in\Gamma_J$, which is $(T_J)$.

The property $(R_j)$. Since $\theta_j^{-1}=\theta_j$ and $\Theta_j\exp\left(iX(f)\right)=\exp\left(iX(\theta_j\cdot f)\right)$, one has in the same way
\[
\rho\left(\exp\left(iX(\theta_j\cdot f)\right)\right)=\left((\theta_j)_{*}\rho\right)\left(\exp\left(iX(f)\right)\right)\qquad(f\in\mathcal S).
\]
Hence $(R_j)$ implies the equalities of the statement; conversely, these equalities for all $f\in\mathcal D_0$ give $(\theta_j)_{*}\rho=\rho$ by Lemma~\ref{lem:density-exponentials}(ii).

The property $(RP_j)$. The identity for $\overline{\Theta_jG}\cdot G$ at the beginning of the proof of Lemma~\ref{lem:reformulations-reflection-positivity} holds at every point of $Q$ and does not involve $(R_j)$. By this identity, condition (2) of Lemma~\ref{lem:reformulations-reflection-positivity} is the inequality in $(RP_j)$ for the families with $f_p\in\mathcal D_0(\{x_j>0\})$ and $c_p\in\mathbb Q+i\mathbb Q$; so $(RP_j)$ implies condition (2). Conversely, assume condition (2), and let $n\ge1$, $c=(c_1,\dots,c_n)\in\mathbb C^n$ and $f_1,\dots,f_n\in\mathcal D(\{x_j>0\})$. By Lemma~\ref{lem:countable-family-test}(i) there are $f^{(k)}_p\in\mathcal D_0(\{x_j>0\})$, $k\ge1$, with $f^{(k)}_p\to f_p$ in $\mathcal S$ as $k\to\infty$, for $p=1,\dots,n$; let moreover $c^{(k)}=(c^{(k)}_1,\dots,c^{(k)}_n)\in(\mathbb Q+i\mathbb Q)^n$ satisfy $c^{(k)}\to c$. Then $\theta_j\cdot f^{(k)}_p\to\theta_j\cdot f_p$ in $\mathcal S$, as $f\mapsto\theta_j\cdot f$ is continuous on $\mathcal S$ (Section~\ref{sec:euclidean-maps-actions}). Since every $\omega\in Q$ is continuous on $\mathcal S$, the functions $\exp\left(i\left(X(f^{(k)}_q)-X(\theta_j\cdot f^{(k)}_p)\right)\right)$, which have modulus $1$, converge pointwise on $Q$ to $\exp\left(i\left(X(f_q)-X(\theta_j\cdot f_p)\right)\right)$, for all $p,q$. By dominated convergence,
\begin{align*}
&\sum_{p,q=1}^n\overline{c^{(k)}_p}\cdot c^{(k)}_q\cdot\rho\left(\exp\left(i\left(X(f^{(k)}_q)-X(\theta_j\cdot f^{(k)}_p)\right)\right)\right)\\
&\qquad\longrightarrow\sum_{p,q=1}^n\overline{c_p}\cdot c_q\cdot\rho\left(\exp\left(i\left(X(f_q)-X(\theta_j\cdot f_p)\right)\right)\right)\qquad(k\to\infty).
\end{align*}
The numbers on the left are non-negative by condition (2), hence so is the limit. This is $(RP_j)$.

The property~\eqref{eq:setup-G}. If $\rho$ satisfies~\eqref{eq:setup-G}, the inequalities of the statement hold in particular for $f\in\mathcal D_0$. Conversely, assume these inequalities, and let $f\in\mathcal S$. By Lemma~\ref{lem:countable-family-test}(i) for $O=\mathbb R^3$ there are $f_n\in\mathcal D_0$ with $f_n\to f$ in $\mathcal S$. Then $X(f_n)\to X(f)$ pointwise on $Q$, as every $\omega\in Q$ is continuous on $\mathcal S$, and $\|f_n\|\to\|f\|$, as convergence in $\mathcal S$ implies convergence in $L^2(\mathbb R^3;\Lambda^2)$. By Fatou's lemma,
\[
\rho\left(\exp X(f)\right)\le\liminf_{n\to\infty}\rho\left(\exp X(f_n)\right)\le\lim_{n\to\infty}\exp\left(\frac{1}{2}\|f_n\|^2\right)=\exp\left(\frac{1}{2}\|f\|^2\right).
\]
Thus $\rho$ satisfies~\eqref{eq:setup-G}.

Ergodicity. Let $\rho$ satisfy $(T_J)$, and let $F,G\in\mathcal A^0$. These functions are bounded, hence in $L^2(\rho)$, and by Lemma~\ref{lem:averages-translations}(iii) one has $A^J_nG\to P^J_{\rho}G$ in $L^2(\rho)$ as $n\to\infty$. Hence
\[
\lim_{n\to\infty}\rho\left(\overline{F}\cdot A^J_nG\right)=\lim_{n\to\infty}\left\langle F,A^J_nG\right\rangle_{\rho}=\left\langle F,P^J_{\rho}G\right\rangle_{\rho}.
\]
Since $\langle1,1\rangle_{\rho}=1$, the orthogonal projection of $L^2(\rho)$ onto $\mathbb C\cdot1$ is the operator $G\mapsto\langle1,G\rangle_{\rho}\cdot1=\rho(G)\cdot1$, and $\left\langle F,\rho(G)\cdot1\right\rangle_{\rho}=\overline{\rho(F)}\cdot\rho(G)$. So the condition of the last item says that $\left\langle F,P^J_{\rho}G\right\rangle_{\rho}=\left\langle F,\rho(G)\cdot1\right\rangle_{\rho}$ for all $F,G\in\mathcal A^0$. The set $\mathcal A^0$ is dense in $L^2(\rho)$, by Lemma~\ref{lem:density-exponentials}(i) for $O=\mathbb R^3$ together with $\Sigma_{\mathbb R^3}=\Sigma$ (Lemma~\ref{lem:countable-family-test}(ii)), and both operators are bounded. Hence the condition holds if and only if $P^J_{\rho}G=\rho(G)\cdot1$ for every $G\in L^2(\rho)$. Since $P^J_{\rho}$ is the orthogonal projection of $L^2(\rho)$ onto $\mathcal E^J(\rho)$, this is the case if and only if $\mathcal E^J(\rho)=\mathbb C\cdot1$, that is, if and only if $\rho$ is $\Gamma_J$-ergodic; as $\rho$ satisfies $(T_J)$, this means $\rho\in\mathcal M^J_{\mathrm{erg}}$.
\end{proof}

\subsection{Functions invariant along a coordinate axis}\label{sec:functions-invariant-along}

\begin{lemma}[Reflection positivity with weights invariant along $e_j$]\label{lem:reflection-positivity-weights}
Let $j\in\{1,2,3\}$, and let $\rho\in\mathcal M$ satisfy $(T_j)$.
\begin{enumerate}
\item[(i)] For every $s\in\mathbb R$,
\[
\mathcal E^j(\rho)\subset L^2\left(\Sigma_{\{x_j>s\}},\rho\right)\cap L^2\left(\Sigma_{\{x_j<s\}},\rho\right).
\]
\end{enumerate}
Assume in addition that $\rho$ satisfies $(R_j)$ and $(RP_j)$.
\begin{enumerate}
\item[(ii)] $\Theta_jF=F$ for every $F\in\mathcal E^j(\rho)$.
\item[(iii)] Let $h\in\mathcal E^j(\rho)$ be bounded. Then $\rho\left(h\cdot\Theta_jG\right)=\rho\left(h\cdot G\right)$ for every $G\in L^1(\rho)$. If moreover $h\ge0$, then
\[
\rho\left(h\cdot\overline{\Theta_jG}\cdot G\right)\ge0\qquad\text{for every }G\in L^2\left(\Sigma_j^+,\rho\right).
\]
\end{enumerate}
\end{lemma}

\begin{proof}
(i) Let $F\in\mathcal E^j(\rho)$ and $s\in\mathbb R$. The space $\mathcal A$ is dense in $L^2(\rho)$, by Lemma~\ref{lem:density-exponentials}(i) with $O=\mathbb R^3$ and because $\Sigma_{\mathbb R^3}=\Sigma$ (Lemma~\ref{lem:countable-family-test}(ii)); choose $G_n\in\mathcal A$ with $\left\|F-G_n\right\|_{L^2(\rho)}\to0$. For each $n$, the definition of $\mathcal A$ in Lemma~\ref{lem:density-exponentials} provides an integer $N\ge1$, numbers $\alpha_1,\dots,\alpha_N\in\mathbb C$ and forms $f_1,\dots,f_N\in\mathcal D_0$, all depending on $n$, such that
\[
G_n=\sum_{q=1}^N\alpha_q\cdot\exp\left(iX(f_q)\right).
\]
As $\mathcal D_0\subset\mathcal D$, the supports of $f_1,\dots,f_N$ are compact; choose $r_n>0$ such that all of them lie in $\{x:|x_j|<r_n\}$. Then every $f_q$ belongs to $\mathcal D(\{x_j>-r_n\})$ and to $\mathcal D(\{x_j<r_n\})$, so that $G_n$ is measurable with respect to $\Sigma_{\{x_j>-r_n\}}$ and with respect to $\Sigma_{\{x_j<r_n\}}$. The translation $\tau_{(r_n+s)\cdot e_j}$ maps $\{x_j>-r_n\}$ onto $\{x_j>s\}$, and $\tau_{(s-r_n)\cdot e_j}$ maps $\{x_j<r_n\}$ onto $\{x_j<s\}$. By Section~\ref{sec:euclidean-maps-actions}, the function $U\left((r_n+s)\cdot e_j\right)G_n$ is $\Sigma_{\{x_j>s\}}$-measurable, and $U\left((s-r_n)\cdot e_j\right)G_n$ is $\Sigma_{\{x_j<s\}}$-measurable. Moreover, by $(T_j)$ and Section~\ref{sec:euclidean-maps-actions}, every $U(t\cdot e_j)$, $t\in\mathbb R$, is unitary on $L^2(\rho)$, and $U(t\cdot e_j)F=F$. Therefore, for $t=r_n+s$ and for $t=s-r_n$,
\[
\left\|F-U(t\cdot e_j)G_n\right\|_{L^2(\rho)}=\left\|U(t\cdot e_j)\left(F-G_n\right)\right\|_{L^2(\rho)}=\left\|F-G_n\right\|_{L^2(\rho)}\to0\qquad(n\to\infty).
\]
The subspaces $L^2\left(\Sigma_{\{x_j>s\}},\rho\right)$ and $L^2\left(\Sigma_{\{x_j<s\}},\rho\right)$ of $L^2(\rho)$ are closed, so $F$ belongs to both.

(ii) For $s\in\mathbb R$ one has $\theta_j\tau_{s\cdot e_j}\theta_j=\tau_{-s\cdot e_j}$, so the rule $\widehat{gh}=\hat g\hat h$ of Section~\ref{sec:euclidean-maps-actions} gives $\Theta_jU(s\cdot e_j)=U(-s\cdot e_j)\Theta_j$. By $(R_j)$ and Section~\ref{sec:euclidean-maps-actions}, $\Theta_j$ is unitary on $L^2(\rho)$; moreover $\Theta_j^2=1$, because $\theta_j^2$ is the identity. For $F\in\mathcal E^j(\rho)$ and $s\in\mathbb R$ one obtains $U(s\cdot e_j)\Theta_jF=\Theta_jU(-s\cdot e_j)F=\Theta_jF$. Hence $\Theta_j$ maps $\mathcal E^j(\rho)$ into itself, and onto itself since $\Theta_j^2=1$. The restriction $\Theta'$ of $\Theta_j$ to $\mathcal E^j(\rho)$ is a unitary operator with $(\Theta')^2=1$; thus $(\Theta')^*=(\Theta')^{-1}=\Theta'$, and $\Theta'$ is self-adjoint. By (i) with $s=0$, $\mathcal E^j(\rho)\subset L^2\left(\Sigma_j^+,\rho\right)$, so Lemma~\ref{lem:reformulations-reflection-positivity}(3) gives $\langle\Theta'F,F\rangle_\rho\ge0$ for every $F\in\mathcal E^j(\rho)$. By the spectral theorem for bounded self-adjoint operators, the spectrum of $\Theta'$ is contained in $\{-1,1\}$, since $(\Theta')^2=1$, and in $[0,\infty)$, since $\Theta'$ is non-negative. Hence the spectrum of $\Theta'$ is $\{1\}$, and $\Theta'=1$.

(iii) Let $h\in\mathcal E^j(\rho)$ be bounded, and let $G\in L^1(\rho)$. By $(R_j)$ and Section~\ref{sec:euclidean-maps-actions}, $\rho\left(\Theta_j\Psi\right)=\rho(\Psi)$ for every $\Psi\in L^1(\rho)$, and $\Theta_j$ is multiplicative; by (ii), $\Theta_jh=h$. Therefore
\[
\rho\left(h\cdot G\right)=\rho\left(\Theta_j\left(h\cdot G\right)\right)=\rho\left(\Theta_jh\cdot\Theta_jG\right)=\rho\left(h\cdot\Theta_jG\right).
\]
Now let in addition $h\ge0$, let $G\in L^2\left(\Sigma_j^+,\rho\right)$, and put $k=\sqrt h$. The function $k$ is bounded, and $k\in\mathcal E^j(\rho)$: since $U(s\cdot e_j)$ commutes with composition by the square root (Section~\ref{sec:euclidean-maps-actions}), $U(s\cdot e_j)k=\sqrt{U(s\cdot e_j)h}=\sqrt h=k$ for every $s\in\mathbb R$. By (i) with $s=0$, $k$ has a $\Sigma_j^+$-measurable representative $k_1$; if $M$ is a bound of $k$, then $\min\left(|k_1|,M\right)$ is a bounded $\Sigma_j^+$-measurable representative of $k$. Consequently $k\cdot G\in L^2\left(\Sigma_j^+,\rho\right)$, and $\Theta_j\left(k\cdot G\right)=\Theta_jk\cdot\Theta_jG=k\cdot\Theta_jG$ by (ii). As $k$ is real and $k^2=h$, Lemma~\ref{lem:reformulations-reflection-positivity}(3) yields
\[
\rho\left(h\cdot\overline{\Theta_jG}\cdot G\right)=\rho\left(\overline{\Theta_j\left(k\cdot G\right)}\cdot k\cdot G\right)=\left\langle\Theta_j\left(k\cdot G\right),k\cdot G\right\rangle_\rho\ge0.
\]
\end{proof}

\subsection{Proof of the decomposition theorem}\label{sec:proof-decomposition-theorem}

\begin{proof}[Proof of Theorem~\ref{thm:ergodic-decomposition}]
Let $J$ be a non-empty subset of $\{1,2,3\}$, and let $m\in\mathcal M$ satisfy $(T_J)$.

\paragraph{The kernel and the measure $\nu$; part (d).}
By Lemma~\ref{lem:countable-family-test}(iii), $(Q,\Sigma)$ is a standard Borel space; in particular it is a Borel space, that is, it is isomorphic as a measurable space to a Borel subset of $[0,1]$. On the probability space $(Q,\Sigma,m)$ let $\xi$ be the identity map of $Q$, regarded as a random element with values in the Borel space $(Q,\Sigma)$, and let $\eta$ be the identity map of $Q$, regarded as a random element with values in $(Q,\mathcal I^J)$; the map $\eta$ is measurable because $\mathcal I^J\subset\Sigma$, and the $\sigma$-algebra generated by $\eta$ is $\mathcal I^J$. The theorem on the existence of regular conditional distributions, applied to $\xi$ and $\eta$, provides a probability kernel $(\omega,A)\mapsto m_\omega(A)$ from $(Q,\mathcal I^J)$ to $(Q,\Sigma)$ such that, for every $A\in\Sigma$,
\[
m_\omega(A)=E_m\left[1_A\mid\mathcal I^J\right](\omega)\qquad\text{for $m$-almost every }\omega\in Q.
\]
By the monotone class theorem, the following holds for every $\Sigma$-measurable function $\Phi\colon Q\to[0,\infty]$ and for every $\Phi\in L^1(m)$:
\begin{equation}\label{eq:ergodic-components-1}
g_\Phi(\omega):=m_\omega(\Phi)=E_m\left[\Phi\mid\mathcal I^J\right](\omega)\quad\text{for $m$-almost every }\omega,\qquad m(\Phi)=\int\mathrm dm(\omega)\,m_\omega(\Phi).
\end{equation}
Here, for $\Phi\ge0$, $g_\Phi$ is an $\mathcal I^J$-measurable function on $Q$ with values in $[0,\infty]$; for bounded complex valued $\Phi$, $g_\Phi$ is, by linearity, a bounded $\mathcal I^J$-measurable function on $Q$; for $\Phi\in L^1(m)$, the number $m_\omega(\Phi)$ is defined for $m$-almost every $\omega$, because $\int\mathrm dm(\omega)\,m_\omega\left(|\Phi|\right)=m\left(|\Phi|\right)<\infty$ by the first case.

The map $\omega\mapsto m_\omega$ from $(Q,\mathcal I^J)$ to $(\mathcal M,\mathcal B_{\mathcal M})$ is measurable, because $\mathcal B_{\mathcal M}$ is generated by the maps $\rho\mapsto\rho(A)$, $A\in\Sigma$, and the functions $\omega\mapsto m_\omega(A)$ are $\mathcal I^J$-measurable. Let $\nu$ be the image of $m$ under this map; it is the probability measure on $(\mathcal M,\mathcal B_{\mathcal M})$ given by
\[
\nu(B)=m\left(\left\{\omega\in Q:m_\omega\in B\right\}\right)\qquad\left(B\in\mathcal B_{\mathcal M}\right).
\]
For $A\in\Sigma$, the second identity of \eqref{eq:ergodic-components-1} with $\Phi=1_A$ and the change of variables for the image measure $\nu$ give $m(A)=\int\mathrm d\nu(\rho)\,\rho(A)$; that is, $m=\int\mathrm d\nu(\rho)\,\rho$.

By Lemma~\ref{lem:measurability-defining-properties}, each of the properties $(T_J)$, \eqref{eq:setup-G}, and $(T_l)$, $(R_l)$, $(RP_l)$ for $l\in\{1,2,3\}$, as well as membership in $\mathcal M^J_{\mathrm{erg}}$, defines a set belonging to $\mathcal B_{\mathcal M}$; in particular $\mathcal M^J_{\mathrm{erg}}\in\mathcal B_{\mathcal M}$. For every set $B$ of this kind, the formula for $\nu(B)$ shows that $\nu$-almost every $\rho$ belongs to $B$ if and only if $m_\omega\in B$ for $m$-almost every $\omega$.

Now let $A\in\Sigma$ satisfy $m(A)=1$. The second identity of \eqref{eq:ergodic-components-1} with $\Phi=1_A$ gives $\int\mathrm dm(\omega)\,m_\omega(A)=1$, while $m_\omega(A)\le1$ for every $\omega$. Hence $m_\omega(A)=1$ for $m$-almost every $\omega$, that is, $\rho(A)=1$ for $\nu$-almost every $\rho$. This is part (d).

\paragraph{Part (a): invariance of the laws $m_\omega$.}
Let $I\in\mathcal I^J$, $a\in\Gamma_J$ and $F\in\mathcal A^0$. By Section~\ref{sec:euclidean-maps-actions}, $U(a)1_I=1_{\tau_aI}=1_I$, and $U(a)$ is multiplicative; by $(T_J)$ and Section~\ref{sec:euclidean-maps-actions}, $m\left(U(a)\Psi\right)=m(\Psi)$ for every $\Psi\in L^1(m)$. Hence
\[
m\left(1_I\cdot U(a)F\right)=m\left(U(a)\left(1_I\cdot F\right)\right)=m\left(1_I\cdot F\right).
\]
Since $I\in\mathcal I^J$ is arbitrary, $E_m\left[U(a)F\mid\mathcal I^J\right]=E_m\left[F\mid\mathcal I^J\right]$ $m$-almost everywhere, and \eqref{eq:ergodic-components-1} gives $m_\omega\left(U(a)F\right)=m_\omega(F)$ for $m$-almost every $\omega$. The pairs $(a,F)$ with $a\in\Gamma_J\cap\mathbb Q^3$ and $F=\exp\left(iX(f)\right)$, $f\in\mathcal D_0$, are countably many; these functions $F$ belong to $\mathcal A^0$, and $U(a)\exp\left(iX(f)\right)=\exp\left(iX(\tau_a\cdot f)\right)$ by Section~\ref{sec:euclidean-maps-actions}. Consequently, for $m$-almost every $\omega$,
\[
m_\omega\left(\exp\left(iX(\tau_a\cdot f)\right)\right)=m_\omega\left(\exp\left(iX(f)\right)\right)\qquad\text{for all }a\in\Gamma_J\cap\mathbb Q^3\text{ and }f\in\mathcal D_0,
\]
and Lemma~\ref{lem:measurability-defining-properties} shows that $m_\omega$ satisfies $(T_J)$.

\paragraph{Part (a): ergodicity of the laws $m_\omega$.}
Write $\mathcal A^0=\{F_1,F_2,\dots\}$. Each $F_p$ is a bounded $\Sigma$-measurable function. For $T>0$ the function $A^J_TF_p$ is defined at every point of $Q$, is $\Sigma$-measurable by Lemma~\ref{lem:averages-translations}(i), and is bounded by $\sup|F_p|$; the function $g_{F_p}$ is bounded and $\mathcal I^J$-measurable. By Lemma~\ref{lem:averages-translations}(iii), (iv) and \eqref{eq:ergodic-components-1}, $A^J_TF_p\to P^J_mF_p=E_m\left[F_p\mid\mathcal I^J\right]=g_{F_p}$ in $L^2(m)$ as $T\to\infty$. Fix an increasing sequence $T_k\to\infty$ such that
\[
\left\|A^J_{T_k}F_p-g_{F_p}\right\|_{L^2(m)}^2\le2^{-k}\qquad\text{for }p\le k.
\]
Then, for every $p$, the series $\sum_k\left\|A^J_{T_k}F_p-g_{F_p}\right\|_{L^2(m)}^2$ converges. By \eqref{eq:ergodic-components-1} and monotone convergence its sum equals $\int\mathrm dm(\omega)\,\sum_km_\omega\left(\left|A^J_{T_k}F_p-g_{F_p}\right|^2\right)$, and by monotone convergence it also equals $\int\mathrm dm(\omega')\,\sum_k\left|A^J_{T_k}F_p(\omega')-g_{F_p}(\omega')\right|^2$. Hence both integrands are finite almost everywhere, so that the terms of both series tend to $0$ almost everywhere; as $p$ ranges over a countable set,
\begin{align}
&\text{for $m$-almost every $\omega$ and all $p$:}\quad m_\omega\left(\left|A^J_{T_k}F_p-g_{F_p}\right|^2\right)\to0\quad(k\to\infty);\notag\\
&\text{for $m$-almost every $\omega'$ and all $p$:}\quad A^J_{T_k}F_p(\omega')\to g_{F_p}(\omega')\quad(k\to\infty).\label{eq:ergodic-components-2}
\end{align}

Fix $p$ and write $g=g_{F_p}$. Since $g$ is bounded and $\mathcal I^J$-measurable, so is $|g|^2$, and each of these two functions coincides $m$-almost everywhere with its own conditional expectation given $\mathcal I^J$. Taking $\Phi=g$ and $\Phi=|g|^2$ in \eqref{eq:ergodic-components-1}, one obtains an $m$-null set outside of which
\[
m_\omega(g)=g(\omega)\qquad\text{and}\qquad m_\omega\left(|g|^2\right)=|g(\omega)|^2.
\]
For every $\omega$ outside this null set, expanding the square and inserting these two identities, one finds
\[
m_\omega\left(\left|g-g(\omega)\right|^2\right)=m_\omega\left(|g|^2\right)-2\operatorname{Re}\left(\overline{g(\omega)}\cdot m_\omega(g)\right)+\left|g(\omega)\right|^2=0.
\]
Together with the first assertion of \eqref{eq:ergodic-components-2} and the triangle inequality in $L^2(m_\omega)$, this shows that for $m$-almost every $\omega$ and all $p$ the functions $A^J_{T_k}F_p$ converge in $L^2(m_\omega)$, as $k\to\infty$, to the constant function $g_{F_p}(\omega)=m_\omega(F_p)$. On the other hand, for $m$-almost every $\omega$ the law $m_\omega$ satisfies $(T_J)$, and then $A^J_{T_k}F_p\to P^J_{m_\omega}F_p$ in $L^2(m_\omega)$ by Lemma~\ref{lem:averages-translations}(iii); here $A^J_{T_k}F_p$ is one and the same function on $Q$, defined at every point, whatever the law. Hence, for $m$-almost every $\omega$, $P^J_{m_\omega}F=m_\omega(F)\cdot1$ for every $F\in\mathcal A^0$. As $\mathcal A^0$ is dense in $L^2(m_\omega)$ (Lemma~\ref{lem:density-exponentials}(i)) and both sides of this identity are continuous in $F\in L^2(m_\omega)$, $P^J_{m_\omega}$ is the orthogonal projection onto $\mathbb C\cdot1$. Thus, for $m$-almost every $\omega$, $\mathcal E^J(m_\omega)=\mathbb C\cdot1$ and $m_\omega\in\mathcal M^J_{\mathrm{erg}}$; that is, $\nu\left(\mathcal M^J_{\mathrm{erg}}\right)=1$.

\paragraph{Part (a): uniqueness of $\nu$.}
Let $\nu'$ be a probability measure on $(\mathcal M,\mathcal B_{\mathcal M})$ with $m=\int\mathrm d\nu'(\rho)\,\rho$ and $\nu'\left(\mathcal M^J_{\mathrm{erg}}\right)=1$. Let $Q_1\in\Sigma$ be the set of the points $\omega'\in Q$ at which the second assertion of \eqref{eq:ergodic-components-2} holds for all $p$. Then $m(Q_1)=1$, so $\int\mathrm d\nu'(\rho)\,\rho(Q_1)=1$, and $\rho(Q_1)=1$ for $\nu'$-almost every $\rho$. Let $\rho\in\mathcal M^J_{\mathrm{erg}}$ satisfy $\rho(Q_1)=1$, and let $p\ge1$. Since $\mathcal E^J(\rho)=\mathbb C\cdot1$, Lemma~\ref{lem:averages-translations}(iii) gives $A^J_{T_k}F_p\to P^J_\rho F_p=\rho(F_p)\cdot1$ in $L^2(\rho)$; at the same time $A^J_{T_k}F_p\to g_{F_p}$ at every point of $Q_1$, hence $\rho$-almost everywhere. A subsequence of a sequence that converges in $L^2(\rho)$ converges $\rho$-almost everywhere to the same limit; therefore $g_{F_p}=\rho(F_p)$ $\rho$-almost everywhere, that is, $m_{\omega'}(F_p)=\rho(F_p)$ for $\rho$-almost every $\omega'$. The functions $\exp\left(iX(f)\right)$, $f\in\mathcal D_0$, are among the $F_p$, and $p$ ranges over a countable set; so Lemma~\ref{lem:density-exponentials}(ii) gives $m_{\omega'}=\rho$ for $\rho$-almost every $\omega'$. Now let $B\in\mathcal B_{\mathcal M}$. The set $\{\omega':m_{\omega'}\in B\}$ belongs to $\mathcal I^J\subset\Sigma$, and what has just been shown gives
\[
\nu(B)=m\left(\left\{\omega':m_{\omega'}\in B\right\}\right)=\int\mathrm d\nu'(\rho)\,\rho\left(\left\{\omega':m_{\omega'}\in B\right\}\right)=\int\mathrm d\nu'(\rho)\,1_B(\rho)=\nu'(B).
\]
Hence $\nu'=\nu$.

\paragraph{Part (b): reflections.}
Let $j\in J$, and let $m$ satisfy $(R_j)$ and $(RP_j)$. As $\Gamma_j\subset\Gamma_J$, the law $m$ satisfies $(T_j)$. Let $h$ be a bounded $\mathcal I^J$-measurable function. Its level sets belong to $\mathcal I^J$, so $U(a)h=h$ for every $a\in\Gamma_J$; thus $h\in\mathcal E^J(m)\subset\mathcal E^j(m)$.

Let $G\in\mathcal A^0$. Lemma~\ref{lem:reflection-positivity-weights}(iii), applied to $\rho=m$, gives $m\left(h\cdot\Theta_jG\right)=m\left(h\cdot G\right)$ for every bounded $\mathcal I^J$-measurable $h$. Hence $E_m\left[\Theta_jG\mid\mathcal I^J\right]=E_m\left[G\mid\mathcal I^J\right]$ $m$-almost everywhere, and \eqref{eq:ergodic-components-1} gives $m_\omega\left(\Theta_jG\right)=m_\omega(G)$ for $m$-almost every $\omega$. For $G=\exp\left(iX(f)\right)$ with $f\in\mathcal D_0$ one has $\Theta_jG=\exp\left(iX(\theta_j\cdot f)\right)$ by Section~\ref{sec:euclidean-maps-actions}. These are countably many conditions, and Lemma~\ref{lem:measurability-defining-properties} shows that $m_\omega$ satisfies $(R_j)$ for $m$-almost every $\omega$.

Now let $G\in\mathcal A^0(\{x_j>0\})$; then $G$ is bounded and $\Sigma_j^+$-measurable. Put $\Phi=\overline{\Theta_jG}\cdot G$ and $\kappa=g_\Phi$. By \eqref{eq:ergodic-components-1}, $\kappa$ is a bounded $\mathcal I^J$-measurable function with $\kappa=E_m\left[\Phi\mid\mathcal I^J\right]$ $m$-almost everywhere. For every bounded $\mathcal I^J$-measurable $h\ge0$, Lemma~\ref{lem:reflection-positivity-weights}(iii) gives
\[
m\left(h\cdot\kappa\right)=m\left(h\cdot\overline{\Theta_jG}\cdot G\right)\ge0.
\]
For $h=1_{\{\operatorname{Im}\kappa>0\}}$ and for $h=1_{\{\operatorname{Im}\kappa<0\}}$ the number $m\left(h\cdot\kappa\right)$ is real, so that
\[
m\left(1_{\{\operatorname{Im}\kappa>0\}}\cdot\operatorname{Im}\kappa\right)=0=m\left(1_{\{\operatorname{Im}\kappa<0\}}\cdot\operatorname{Im}\kappa\right),
\]
and $\operatorname{Im}\kappa=0$ $m$-almost everywhere. For $h=1_{\{\operatorname{Re}\kappa<0\}}$ one obtains $m\left(1_{\{\operatorname{Re}\kappa<0\}}\cdot\operatorname{Re}\kappa\right)\ge0$, so $\operatorname{Re}\kappa\ge0$ $m$-almost everywhere. Thus $\kappa\ge0$ $m$-almost everywhere; by the definition of $\kappa$, $m_\omega\left(\overline{\Theta_jG}\cdot G\right)\ge0$ for $m$-almost every $\omega$.

The set $\mathcal A^0(\{x_j>0\})$ is countable. Hence, for $m$-almost every $\omega$, the law $m_\omega$ satisfies $(R_j)$ and condition (2) of Lemma~\ref{lem:reformulations-reflection-positivity}, and therefore, by that lemma, condition (1), which is $(RP_j)$. Moreover $m_\omega$ satisfies $(T_J)$, hence $(T_j)$, for $m$-almost every $\omega$. By the observation following the definition of $\nu$, $\nu$-almost every $\rho$ satisfies $(T_j)$, $(R_j)$ and $(RP_j)$, and by Lemma~\ref{lem:reformulations-reflection-positivity}(5) every such $\rho$ is reflection positive at every plane $\{x_j=c\}$, $c\in\mathbb R$.

\paragraph{Part (c): the Gaussian bound.}
Let $m$ satisfy \eqref{eq:setup-G}. Fix $f\in\mathcal S$, and put
\[
G=\exp X(f)>0,\qquad\mathsf c=\exp\left(\frac{1}{2}\|f\|^2\right),\qquad\phi=g_G,\qquad d=\#|J|;
\]
by \eqref{eq:ergodic-components-1}, $\phi=E_m\left[G\mid\mathcal I^J\right]$ $m$-almost everywhere. Applying \eqref{eq:setup-G} to $rf$ one finds
\[
\|G\|_{L^r(m)}^r=m\left(\exp X(rf)\right)\le\exp\left(\frac{r^2}{2}\|f\|^2\right)\qquad\text{for every integer }r\ge1.
\]

\emph{Step 1.} Let $n\ge2$ be an integer and $T>0$, and let $\mathrm dx$ be the Lebesgue measure on $(\Gamma_J)^n$. For $x=(x_1,\dots,x_n)\in(\Gamma_J)^n$, Section~\ref{sec:euclidean-maps-actions} gives $\prod_{p=1}^nU(x_p)G=\exp X\left(\sum_{p=1}^n\tau_{x_p}\cdot f\right)$, and
\[
\left\|\sum_{p=1}^n\tau_{x_p}\cdot f\right\|^2=n\cdot\|f\|^2+\sum_{p\ne q}c(x_p-x_q),\qquad c(y)=\left\langle f,\tau_y\cdot f\right\rangle\quad\left(y\in\mathbb R^3\right),
\]
the sum being taken over the ordered pairs $(p,q)$ with $p\ne q$. By the definition of $A^J_T$ in Lemma~\ref{lem:averages-translations}(i), $\left(A^J_TG\right)^n=T^{-dn}\cdot\int_{([0,T]^J)^n}\mathrm dx\,\prod_{p=1}^nU(x_p)G$ at every point of $Q$. By Tonelli's theorem, which applies since the integrand is jointly measurable (Lemma~\ref{lem:averages-translations}(i)), and by \eqref{eq:setup-G},
\begin{align*}
m\left(\left(A^J_TG\right)^n\right)&=T^{-dn}\cdot\int_{([0,T]^J)^n}\mathrm dx\;m\left(\exp X\left(\sum_{p=1}^n\tau_{x_p}\cdot f\right)\right)\\
&\le T^{-dn}\cdot\int_{([0,T]^J)^n}\mathrm dx\;\exp\left(\frac{n}{2}\|f\|^2+\frac{1}{2}\sum_{p\ne q}c(x_p-x_q)\right).
\end{align*}
By the Cauchy--Schwarz inequality and $\|\tau_y\cdot f\|=\|f\|$, one has $|c(y)|\le\|f\|^2$. For $y\in\mathbb R^3$ let $1_{\{|z|\ge\frac{1}{2}|y|\}}f$ be the product of $f$ with the indicator function of $\left\{z\in\mathbb R^3:|z|\ge\frac{1}{2}|y|\right\}$. For all $y,z\in\mathbb R^3$ one has $|z|\ge\frac{1}{2}|y|$ or $|z-y|\ge\frac{1}{2}|y|$. Splitting the integral over $z$ which defines $c(y)$ into the integral over $\left\{|z|\ge\frac{1}{2}|y|\right\}$ and the integral over its complement, on which $|z-y|\ge\frac{1}{2}|y|$, and using dominated convergence, one obtains
\[
|c(y)|\le2\,\|f\|\cdot\left\|1_{\{|z|\ge\frac{1}{2}|y|\}}f\right\|\to0\qquad\left(|y|\to\infty\right).
\]
Let $\varepsilon>0$, and choose $\delta>0$ such that $|c(y)|\le\varepsilon$ whenever $|y|_\infty\ge\delta$, where $|y|_\infty=\max_l|y_l|$; this is possible because $|y|\ge|y|_\infty$. Let $K_\delta$ be the set of the $x\in([0,T]^J)^n$ such that $|x_p-x_q|_\infty<\delta$ for some $p\ne q$. For each of the $\binom{n}{2}$ unordered pairs $\{p,q\}$, the condition $|x_p-x_q|_\infty<\delta$ confines $x_q$ to a cube of side $2\delta$ in $\Gamma_J$ around $x_p$; hence the Lebesgue measure of $K_\delta$ is at most $\binom{n}{2}\cdot(2\delta)^d\cdot T^{d(n-1)}$. Outside $K_\delta$ the exponent in the last integral is at most $\frac{n}{2}\|f\|^2+\frac{n(n-1)}{2}\varepsilon$, and on $K_\delta$ it is at most $\frac{n^2}{2}\|f\|^2$, since $|c|\le\|f\|^2$. Therefore
\begin{equation}\label{eq:ergodic-components-3}
m\left(\left(A^J_TG\right)^n\right)\le\mathsf c^{\,n}\cdot\exp\left(\frac{n(n-1)\varepsilon}{2}\right)+\binom{n}{2}\cdot(2\delta)^d\cdot T^{-d}\cdot\exp\left(\frac{n^2}{2}\|f\|^2\right).
\end{equation}

\emph{Step 2.} As $G\in L^2(m)$, Lemma~\ref{lem:averages-translations}(iii), (iv) gives $A^J_TG\to P^J_mG=E_m\left[G\mid\mathcal I^J\right]=\phi$ in $L^2(m)$ as $T\to\infty$. Moreover $\left\|A^J_TG\right\|_{L^{2n}(m)}\le\|G\|_{L^{2n}(m)}$ by Lemma~\ref{lem:averages-translations}(iii), and $\|\phi\|_{L^{2n}(m)}\le\|G\|_{L^{2n}(m)}$ by Jensen's inequality for conditional expectations. Put $\vartheta=\frac{1}{n-1}\in(0,1]$; then $\frac{1}{n}=\frac{\vartheta}{2}+\frac{1-\vartheta}{2n}$, and H\"older's inequality gives
\[
\left\|A^J_TG-\phi\right\|_{L^n(m)}\le\left\|A^J_TG-\phi\right\|_{L^2(m)}^{\vartheta}\cdot\left(2\|G\|_{L^{2n}(m)}\right)^{1-\vartheta}\to0\qquad(T\to\infty).
\]
Since $A^J_TG\ge0$ and $\phi\ge0$, it follows that $m\left(\left(A^J_TG\right)^n\right)=\left\|A^J_TG\right\|_{L^n(m)}^n\to\|\phi\|_{L^n(m)}^n=m\left(\phi^n\right)$. Let $T\to\infty$ in \eqref{eq:ergodic-components-3}, and then $\varepsilon\to0$: one obtains $m\left(\phi^n\right)\le\mathsf c^{\,n}$, that is, $\|\phi\|_{L^n(m)}\le\mathsf c$, for every integer $n\ge2$. As $m$ is a probability measure, $\|\phi\|_{L^\infty(m)}=\lim_{n\to\infty}\|\phi\|_{L^n(m)}\le\mathsf c$; in other words,
\[
m_\omega\left(\exp X(f)\right)\le\exp\left(\frac{1}{2}\|f\|^2\right)\qquad\text{for $m$-almost every }\omega.
\]

\emph{Step 3.} Steps 1 and 2 apply to each of the countably many $f\in\mathcal D_0$. Hence, for $m$-almost every $\omega$, the inequality of \eqref{eq:setup-G} holds for $m_\omega$ and every $f\in\mathcal D_0$, and Lemma~\ref{lem:measurability-defining-properties} shows that $m_\omega$ satisfies \eqref{eq:setup-G}, for every $f\in\mathcal S$ and with the same right side $\exp\left(\frac{1}{2}\|f\|^2\right)$. By the observation following the definition of $\nu$, $\nu$-almost every $\rho$ satisfies \eqref{eq:setup-G}.
\end{proof}

\subsection{Reconstruction along a coordinate axis and directional mixing}\label{sec:reconstruction-along-coordinate}

\begin{proposition}[Reconstruction in the direction $e_j$]\label{prop:reconstruction-direction}
Let $j\in\{1,2,3\}$ and let $\rho\in\mathcal M$ satisfy $(T_j)$, $(R_j)$ and $(RP_j)$. Put $\mathcal E_{+}=L^2(\Sigma_j^{+},\rho)$, and let $b\colon\mathcal E_{+}\times\mathcal E_{+}\to\mathbb C$ be the form
\[
b(F,G)=\langle\Theta_jF,G\rangle_{\rho}\qquad\left(F,G\in\mathcal E_{+}\right).
\]
\begin{enumerate}
\item[(i)] The form $b$ is sesquilinear, Hermitian and positive semidefinite, and $b(G,G)\le\|G\|^2_{L^2(\rho)}$ for every $G\in\mathcal E_{+}$. The set $\left\{G\in\mathcal E_{+}:b(G,G)=0\right\}$ is a linear subspace of $\mathcal E_{+}$, and $b$ induces a scalar product on the quotient of $\mathcal E_{+}$ by this subspace. Let $\mathcal H_j$ be the completion of this quotient; its scalar product and its norm are written $\langle\cdot,\cdot\rangle$ and $\|\cdot\|$. Let $[\cdot]=[\cdot]_j\colon\mathcal E_{+}\to\mathcal H_j$ be the canonical map, so that $\langle[F],[G]\rangle=b(F,G)$ for $F,G\in\mathcal E_{+}$; it is a linear contraction with dense range. Put $\Omega_j=[1]$. Then $\|\Omega_j\|=1$, and $\langle\Omega_j,[G]\rangle=\rho(G)$ for every $G\in\mathcal E_{+}$.
\item[(ii)] There is exactly one family $\left(T(t)\right)_{t\ge0}$ of bounded operators on $\mathcal H_j$ such that
\[
T(t)[G]=\left[U(t\cdot e_j)G\right]\qquad\left(t\ge0,\ G\in\mathcal E_{+}\right).
\]
Every $T(t)$ is self-adjoint with $\|T(t)\|\le1$; one has $T(0)=1$ and $T(s)T(t)=T(s+t)$ for all $s,t\ge0$; and for every $\psi\in\mathcal H_j$ the map $t\mapsto T(t)\psi$ is continuous on $[0,\infty)$. There is exactly one self-adjoint operator $H_j\ge0$ in $\mathcal H_j$ such that $T(t)=\exp(-t\cdot H_j)$ for every $t\ge0$, namely $H_j=-\log T(1)$; and $H_j\Omega_j=0$.
\item[(iii)] Let $y\in\mathbb R^3$ satisfy $y\perp e_j$ and $(\tau_y)_{*}\rho=\rho$. There is exactly one unitary operator $V_j(y)$ on $\mathcal H_j$ such that $V_j(y)[G]=[U(y)G]$ for every $G\in\mathcal E_{+}$. It commutes with $T(t)$ for every $t\ge0$, and $V_j(y)\Omega_j=\Omega_j$. Put $e_j^{\perp}=\left\{y\in\mathbb R^3:y\perp e_j\right\}$. If $(\tau_y)_{*}\rho=\rho$ for every $y\in e_j^{\perp}$, then $y\mapsto V_j(y)$ is a strongly continuous unitary representation of the additive group $e_j^{\perp}$ on $\mathcal H_j$.
\item[(iv)] Let $P_0$ be the orthogonal projection of $\mathcal H_j$ onto $\ker H_j$. Then $P_0[G]=[P^j_{\rho}G]$ for every $G\in\mathcal E_{+}$, and the map $F\mapsto[F]$ is an isometry of $\mathcal E^j(\rho)$, which is a subspace of $\mathcal E_{+}$ by Lemma~\ref{lem:reflection-positivity-weights}(i), onto $\ker H_j$.
\end{enumerate}
\end{proposition}

\begin{proof}
(i) By $(R_j)$, $(\theta_j)_{*}\rho=\rho$; as moreover $\theta_j\circ\theta_j$ is the identity, the operator $\Theta_j$ is unitary on $L^2(\rho)$ and $\Theta_j^2=1$ (Section~\ref{sec:euclidean-maps-actions}); a unitary involution is self-adjoint. Since $\Theta_j$ is complex linear and $\langle\cdot,\cdot\rangle_{\rho}$ is antilinear in the first entry, $b$ is sesquilinear, and for $F,G\in\mathcal E_{+}$
\[
\overline{b(G,F)}=\langle F,\Theta_jG\rangle_{\rho}=\langle\Theta_jF,G\rangle_{\rho}=b(F,G).
\]
One has $b(G,G)\ge0$ by Lemma~\ref{lem:reformulations-reflection-positivity}(3), and the Cauchy--Schwarz inequality in $L^2(\rho)$ gives
\[
b(G,G)\le\left\|\Theta_jG\right\|_{L^2(\rho)}\cdot\|G\|_{L^2(\rho)}=\|G\|^2_{L^2(\rho)}.
\]
By the Cauchy--Schwarz inequality for the positive semidefinite Hermitian form $b$, $\left|b(F,G)\right|^2\le b(F,F)\cdot b(G,G)$, the equality $b(G,G)=0$ holds if and only if $b(F,G)=0$ for every $F\in\mathcal E_{+}$. Hence these $G$ form a linear subspace, $b(F,G)$ depends only on the classes of $F$ and $G$ modulo this subspace, and $b$ defines a scalar product on the quotient. The map $[\cdot]$ is linear, $\|[G]\|^2=b(G,G)\le\|G\|^2_{L^2(\rho)}$, and its range is dense by the construction of the completion. The constant function $1$ lies in $\mathcal E_{+}$ and $\Theta_j1=1$; therefore $\langle\Omega_j,[G]\rangle=\langle\Theta_j1,G\rangle_{\rho}=\rho(G)$, and the choice $G=1$ gives $\|\Omega_j\|^2=\rho(1)=1$.

(ii) For $t\in\mathbb R$ write $U_t=U(t\cdot e_j)$, and let $t\ge0$. Since $\tau_{t\cdot e_j}\{x_j>0\}=\{x_j>t\}\subset\{x_j>0\}$, the operator $U_t$ maps $\Sigma_j^{+}$-measurable functions to $\Sigma_{\{x_j>t\}}$-measurable, hence $\Sigma_j^{+}$-measurable, functions (Section~\ref{sec:euclidean-maps-actions}); by $(T_j)$ it is unitary on $L^2(\rho)$, with inverse $U_{-t}$. Hence $U_t\mathcal E_{+}\subset\mathcal E_{+}$. For $G\in\mathcal E_{+}$ put $|G|_b=b(G,G)^{\frac{1}{2}}=\|[G]\|$. The identity $\theta_j\circ\tau_{t\cdot e_j}\circ\theta_j=\tau_{-t\cdot e_j}$ gives $\Theta_jU_t=U_{-t}\Theta_j$, and therefore, for $F,G\in\mathcal E_{+}$,
\[
b(F,U_tG)=\langle U_{-t}\Theta_jF,G\rangle_{\rho}=\langle\Theta_jU_tF,G\rangle_{\rho}=b(U_tF,G).
\]
Together with $U_tU_t=U_{2t}$ and the Cauchy--Schwarz inequality for $b$, this yields
\[
|U_tG|_b^2=b(G,U_{2t}G)\le|G|_b\cdot|U_{2t}G|_b.
\]
Applying this inequality with $2^k\cdot t$ in place of $t$, $k=1,\dots,n-1$, one obtains by induction on $n\ge1$
\[
|U_tG|_b\le|G|_b^{1-2^{-n}}\cdot\left|U_{2^nt}G\right|_b^{2^{-n}}\le|G|_b^{1-2^{-n}}\cdot\|G\|_{L^2(\rho)}^{2^{-n}},
\]
where the second inequality uses $|U_sG|_b\le\|U_sG\|_{L^2(\rho)}=\|G\|_{L^2(\rho)}$, by (i) and the unitarity of $U_s$. Letting $n\to\infty$ one obtains $|U_tG|_b\le|G|_b$.

\emph{The semigroup.} By the last inequality, $[G]=[G']$ implies $[U_tG]=[U_tG']$, and $[G]\mapsto[U_tG]$ is a well defined linear contraction on the dense subspace $[\mathcal E_{+}]$ of $\mathcal H_j$. It has exactly one continuous extension $T(t)$ to $\mathcal H_j$, and $\|T(t)\|\le1$; every bounded operator that maps $[G]$ to $[U_tG]$ for all $G\in\mathcal E_{+}$ coincides with $T(t)$. This proves the existence and the uniqueness of the family. By the symmetry of $b$ under $U_t$, $\langle[F],T(t)[G]\rangle=\langle T(t)[F],[G]\rangle$ for $F,G\in\mathcal E_{+}$, so that the bounded operator $T(t)$ is self-adjoint, by density and continuity. On classes, $T(0)[G]=[G]$ and $T(s)T(t)[G]=[U_sU_tG]=[U_{s+t}G]=T(s+t)[G]$; by density, $T(0)=1$ and $T(s)T(t)=T(s+t)$ for $s,t\ge0$. By Lemma~\ref{lem:averages-translations}(ii),
\[
\left\|T(t)[G]-[G]\right\|=|U_tG-G|_b\le\|U_tG-G\|_{L^2(\rho)}\to0\qquad(t\to0).
\]
As $\|T(t)\|\le1$ and $[\mathcal E_{+}]$ is dense, it follows that $T(t)\psi\to\psi$ as $t\to0$ for every $\psi\in\mathcal H_j$; by the semigroup law and $\|T(s)\|\le1$, the map $t\mapsto T(t)\psi$ is continuous on $[0,\infty)$.

\emph{Positivity and injectivity.} For $t\ge0$ one has $T(t)=T(t/2)^2=T(t/2)^{*}\,T(t/2)\ge0$. Let $t>0$ and $\psi\in\mathcal H_j$ with $T(t)\psi=0$. Then $\|T(t/2)\psi\|^2=\langle\psi,T(t)\psi\rangle=0$; by induction $T(2^{-n}\cdot t)\psi=0$ for every $n\in\mathbb N$, and $\psi=\lim_{n\to\infty}T(2^{-n}\cdot t)\psi=0$. Thus every $T(t)$ is injective, and $0\le T(1)\le1$.

\emph{The operator $H_j=-\log T(1)$.} The Hilbert space $\mathcal H_j$ is separable. Indeed, by Lemma~\ref{lem:density-exponentials}(i) the countable set $\mathcal A^0(\{x_j>0\})$ is dense in $L^2(\Sigma_{\{x_j>0\}},\rho)=\mathcal E_{+}$, and its image under the continuous map $[\cdot]$, whose range is dense, is a countable dense subset of $\mathcal H_j$. By the spectral theorem in multiplication operator form, applied to the bounded self-adjoint operator $T(1)$ on the separable Hilbert space $\mathcal H_j$, there are a finite measure space $(Z,\mathcal Z,\kappa)$, a unitary operator $W\colon\mathcal H_j\to L^2(Z,\kappa)$ and a bounded measurable function $h\colon Z\to\mathbb R$ such that
\[
W\,T(1)\,W^{-1}\phi=h\cdot\phi\qquad\left(\phi\in L^2(Z,\kappa)\right).
\]
In what follows a bounded measurable function on $Z$ also denotes the operator of multiplication by it on $L^2(Z,\kappa)$; thus $T(1)=W^{-1}\,h\,W$. Since $0\le T(1)\le1$, one has $0\le h\le1$ $\kappa$-almost everywhere (apply the two inequalities to the indicator functions of $\{h<0\}$ and of $\{h>1\}$, which lie in $L^2(Z,\kappa)$ because $\kappa$ is finite). Since $T(1)$ is injective, $\kappa(\{h=0\})=0$ (multiplication by $h$ annihilates the indicator function of $\{h=0\}$). Hence $\{h\le0\}\cup\{h>1\}$ is a $\kappa$-null set, and the function equal to $1$ on this set and to $h$ elsewhere defines the same operator of multiplication as $h$; replacing $h$ by it, one may assume that $0<h\le1$ at every point of $Z$. Then $-\log h$ is a measurable function from $Z$ to $[0,\infty)$. Multiplication by a real measurable function that is finite almost everywhere is a self-adjoint operator on its natural domain. Consequently the operator $H_j$ in $\mathcal H_j$ with domain and action
\[
\left\{\psi\in\mathcal H_j:(\log h)\cdot W\psi\in L^2(Z,\kappa)\right\},\qquad W\,H_j\psi=-(\log h)\cdot W\psi,
\]
is self-adjoint, and $H_j\ge0$; it is the operator $-\log T(1)$ of the functional calculus. By the functional calculus, $\exp(-t\cdot H_j)=W^{-1}\,h^t\,W$ for $t\ge0$, where $h^t=\exp\left(t\cdot\log h\right)$; in particular $\exp(-H_j)=T(1)$.

\emph{The identity $T(t)=\exp(-t\cdot H_j)$.} Let $n\ge0$. The operators $T(2^{-n})$ and $\exp(-2^{-n}\cdot H_j)=W^{-1}\,h^{2^{-n}}\,W$ are non-negative, and the $2^n$-th power of each of them is $T(1)$, by the semigroup law and because $\left(h^{2^{-n}}\right)^{2^n}=h$. A non-negative bounded operator has exactly one non-negative square root; by induction on $n$, $T(1)$ has exactly one non-negative $2^n$-th root (if $A\ge0$ and $A^{2^n}=T(1)$ with $n\ge1$, then $A^2$ is a non-negative $2^{n-1}$-th root of $T(1)$, and $A$ is the non-negative square root of $A^2$). Hence $T(2^{-n})=\exp(-2^{-n}\cdot H_j)$, and
\[
T(k\cdot2^{-n})=T(2^{-n})^k=\exp\left(-k\cdot2^{-n}\cdot H_j\right)\qquad(k,n\in\mathbb N).
\]
For $\psi\in\mathcal H_j$ the maps $t\mapsto T(t)\psi$ and $t\mapsto\exp(-t\cdot H_j)\psi=W^{-1}\left(h^t\cdot W\psi\right)$ are continuous on $[0,\infty)$, the second by dominated convergence. They agree at the non-negative dyadic rationals, hence everywhere: $T(t)=\exp(-t\cdot H_j)$ for every $t\ge0$.

\emph{Uniqueness of $H_j$.} Let $\chi\colon[0,1]\to[0,1]$ be the continuous function with $\chi(0)=0$ and $\chi(s)=(1-\log s)^{-1}$ for $0<s\le1$. Let $H$ be a self-adjoint operator in $\mathcal H_j$ with $H\ge0$ and $\exp(-H)=T(1)$. By the spectral theorem for self-adjoint operators, bounded or not, in multiplication operator form, there are a finite measure space $(Z',\mathcal Z',\kappa')$, a unitary operator $W'\colon\mathcal H_j\to L^2(Z',\kappa')$ and a measurable function $\eta\colon Z'\to\mathbb R$ such that $W'\,H\,W'^{-1}$ is multiplication by $\eta$ on its natural domain; $\eta\ge0$ $\kappa'$-almost everywhere, because $H\ge0$. By the functional calculus, $W'\,T(1)\,W'^{-1}=W'\exp(-H)\,W'^{-1}$ is multiplication by $\exp(-\eta)$, and $W'\,(1+H)^{-1}\,W'^{-1}$ is multiplication by $(1+\eta)^{-1}=\chi\circ\exp(-\eta)$. For every polynomial $p$, $W'\,p\left(T(1)\right)\,W'^{-1}$ is multiplication by $p\circ\exp(-\eta)$. Approximating $\chi$ uniformly on $[0,1]$ by polynomials (Weierstrass approximation theorem), one obtains
\[
(1+H)^{-1}=\chi\left(T(1)\right),
\]
where $\chi\left(T(1)\right)$ is given by the continuous functional calculus of $T(1)$, whose spectrum lies in $[0,1]$. The operator $H_j$ is of this kind, with $Z'=Z$, $W'=W$ and $\eta=-\log h$. Hence $(1+H)^{-1}=\chi\left(T(1)\right)=(1+H_j)^{-1}$. A self-adjoint operator $H\ge0$ is determined by $(1+H)^{-1}$: its domain is the range of $(1+H)^{-1}$, on which $H=\left((1+H)^{-1}\right)^{-1}-1$. Therefore $H=H_j$. In particular, every self-adjoint operator $H\ge0$ in $\mathcal H_j$ with $T(t)=\exp(-t\cdot H)$ for all $t\ge0$ equals $H_j$.

\emph{The kernel of $H_j$.} For $\psi\in\mathcal H_j$ each of the conditions $H_j\psi=0$ and $T(1)\psi=\psi$ means that $W\psi$ vanishes $\kappa$-almost everywhere on $\{h\ne1\}$; and if they hold, then $T(t)\psi=W^{-1}\left(h^t\cdot W\psi\right)=\psi$ for every $t\ge0$. Hence
\[
\ker H_j=\left\{\psi\in\mathcal H_j:T(t)\psi=\psi\text{ for every }t\ge0\right\},
\]
and the orthogonal projection of $\mathcal H_j$ onto $\ker H_j$ is $W^{-1}\,1_{\{h=1\}}\,W$. Since $U_t1=1$, one has $T(t)\Omega_j=[U_t1]=\Omega_j$ for every $t\ge0$, and therefore $H_j\Omega_j=0$.

(iii) Let $y\perp e_j$ with $(\tau_y)_{*}\rho=\rho$. Then $(\tau_{-y})_{*}\rho=\rho$, and $U(y)$ is unitary on $L^2(\rho)$, with inverse $U(-y)$ (Section~\ref{sec:euclidean-maps-actions}). As $\theta_jy=y$, one has $\theta_j\circ\tau_y\circ\theta_j=\tau_y$, so that $U(y)$ commutes with $\Theta_j$. As $\tau_y\{x_j>0\}=\tau_{-y}\{x_j>0\}=\{x_j>0\}$, the operators $U(y)$ and $U(-y)$ map $\Sigma_j^{+}$-measurable functions to $\Sigma_j^{+}$-measurable functions; hence $U(y)$ maps $\mathcal E_{+}$ onto itself. Therefore, for $F,G\in\mathcal E_{+}$,
\[
b\left(U(y)F,U(y)G\right)=\langle U(y)\Theta_jF,U(y)G\rangle_{\rho}=\langle\Theta_jF,G\rangle_{\rho}=b(F,G).
\]
It follows that $[G]\mapsto[U(y)G]$ is a well defined isometry of $[\mathcal E_{+}]$ onto itself, with inverse $[G]\mapsto[U(-y)G]$. Its continuous extension $V_j(y)$ to $\mathcal H_j$ is a unitary operator, and by density it is the only bounded operator that maps $[G]$ to $[U(y)G]$ for all $G\in\mathcal E_{+}$. From $U(y)U_t=U(y+t\cdot e_j)=U_tU(y)$ one obtains $V_j(y)T(t)[G]=T(t)V_j(y)[G]$ on classes, hence $V_j(y)T(t)=T(t)V_j(y)$ by density; moreover $V_j(y)\Omega_j=[U(y)1]=\Omega_j$.

Let now $(\tau_y)_{*}\rho=\rho$ for every $y\in e_j^{\perp}$. For $y,y'\in e_j^{\perp}$ and $G\in\mathcal E_{+}$,
\[
V_j(y)V_j(y')[G]=\left[U(y)U(y')G\right]=\left[U(y+y')G\right]=V_j(y+y')[G],
\]
and $V_j(0)[G]=[G]$; by density, $V_j(y)V_j(y')=V_j(y+y')$ and $V_j(0)=1$. With $J=\{1,2,3\}\setminus\{j\}$ one has $\Gamma_J=e_j^{\perp}$, so that $\rho$ satisfies $(T_J)$, and Lemma~\ref{lem:averages-translations}(ii) gives
\[
\left\|V_j(y)[G]-[G]\right\|=\left|U(y)G-G\right|_b\le\left\|U(y)G-G\right\|_{L^2(\rho)}\to0\qquad(y\to0).
\]
As $\|V_j(y)\|=1$ and $[\mathcal E_{+}]$ is dense, $V_j(y)\psi\to\psi$ as $y\to0$ for every $\psi\in\mathcal H_j$; by the group law, $y\mapsto V_j(y)\psi$ is continuous on $e_j^{\perp}$.

(iv) By Lemma~\ref{lem:reflection-positivity-weights}(i) with $s=0$, $\mathcal E^j(\rho)\subset L^2(\Sigma_{\{x_j>0\}},\rho)=\mathcal E_{+}$; in particular $P^j_{\rho}G\in\mathcal E_{+}$ for $G\in\mathcal E_{+}$. Let $F,G\in\mathcal E_{+}$. For $t\ge0$,
\[
\langle[F],T(t)[G]\rangle=b(F,U_tG)=\langle\Theta_jF,U_tG\rangle_{\rho}.
\]
Both sides are averaged over $t\in[0,S]$, $S>0$, and $S$ tends to infinity. For $\lambda\ge0$ let
\[
\varphi_S(\lambda)=\frac{1}{S}\int_0^S\mathrm{d}t\,\exp(-t\cdot\lambda);
\]
then $0\le\varphi_S\le1$, and $\varphi_S(\lambda)\to1_{\{0\}}(\lambda)$ as $S\to\infty$, for every $\lambda\ge0$. Since $T(t)=W^{-1}\,h^t\,W$ and $h^t=\exp\left(-t\cdot(-\log h)\right)$, Fubini's theorem gives
\[
\frac{1}{S}\int_0^S\mathrm{d}t\,\langle[F],T(t)[G]\rangle=\left\langle W[F],\varphi_S(-\log h)\cdot W[G]\right\rangle_{L^2(Z,\kappa)},
\]
and by dominated convergence the right side tends to $\left\langle W[F],1_{\{h=1\}}\cdot W[G]\right\rangle_{L^2(Z,\kappa)}=\langle[F],P_0[G]\rangle$, because $P_0=W^{-1}\,1_{\{h=1\}}\,W$ (proof of (ii)). On the other side, Fubini's theorem gives
\[
\frac{1}{S}\int_0^S\mathrm{d}t\,\langle\Theta_jF,U_tG\rangle_{\rho}=\left\langle\Theta_jF,A^j_SG\right\rangle_{\rho},
\]
with the average $A^j_S$ of Lemma~\ref{lem:averages-translations} for $J=\{j\}$, and $A^j_SG\to P^j_{\rho}G$ in $L^2(\rho)$ as $S\to\infty$ by Lemma~\ref{lem:averages-translations}(iii). Comparing the two limits,
\[
\langle[F],P_0[G]\rangle=\langle\Theta_jF,P^j_{\rho}G\rangle_{\rho}=b(F,P^j_{\rho}G)=\left\langle[F],[P^j_{\rho}G]\right\rangle
\]
for every $F\in\mathcal E_{+}$. As $[\mathcal E_{+}]$ is dense in $\mathcal H_j$, $P_0[G]=[P^j_{\rho}G]$.

Let $F\in\mathcal E^j(\rho)$. By Lemma~\ref{lem:reflection-positivity-weights}(ii), $\Theta_jF=F$, so that $\|[F]\|^2=\langle\Theta_jF,F\rangle_{\rho}=\|F\|^2_{L^2(\rho)}$; and $U_tF=F$, so that $T(t)[F]=[U_tF]=[F]$ for every $t\ge0$, that is, $[F]\in\ker H_j$ by the description of $\ker H_j$ in the proof of (ii). Thus the restriction of $[\cdot]$ to $\mathcal E^j(\rho)$ is an isometry into $\ker H_j$. Its image is complete, hence closed, because $\mathcal E^j(\rho)$ is closed in $L^2(\rho)$. The image contains $P_0[\mathcal E_{+}]=\left\{[P^j_{\rho}G]:G\in\mathcal E_{+}\right\}$, which is dense in $\ker H_j=P_0\mathcal H_j$, since $[\mathcal E_{+}]$ is dense in $\mathcal H_j$ and $P_0$ is continuous. Hence the image is $\ker H_j$.
\end{proof}

\begin{proposition}[Uniform mixing along a coordinate axis]\label{prop:uniform-mixing-along}
Let $j\in\{1,2,3\}$, and let $\rho\in\mathcal M$ satisfy $(T_j)$, $(R_j)$ and $(RP_j)$. Let $Y$ be a subset of $e_j^{\perp}=\left\{y\in\mathbb R^3:y\perp e_j\right\}$ such that $0\in Y$, $Y=-Y$, and $(\tau_y)_{*}\rho=\rho$ for every $y\in Y$. Then, for all $F,G\in L^2(\rho)$,
\[
\lim_{|s|\to\infty}\,\sup_{y\in Y}\left|\left\langle F,U(s\cdot e_j+y)\left(1-P^j_{\rho}\right)G\right\rangle_{\rho}\right|=0,
\]
where $s$ ranges over $\mathbb R$. In particular, $U(s\cdot e_j)$ converges weakly to $P^j_{\rho}$ as $|s|\to\infty$: for all $F,G\in L^2(\rho)$,
\[
\lim_{|s|\to\infty}\left\langle F,U(s\cdot e_j)G\right\rangle_{\rho}=\left\langle F,P^j_{\rho}G\right\rangle_{\rho}.
\]
\end{proposition}

\begin{proof}
Let $\mathcal E_{+}=L^2(\Sigma_j^{+},\rho)$, and let $\mathcal H_j$, the canonical map $[\cdot]\colon\mathcal E_{+}\to\mathcal H_j$, the semigroup $T(t)=\exp(-t\cdot H_j)$, $t\ge0$, the unitary operators $V_j(y)$ and the orthogonal projection $P_0$ of $\mathcal H_j$ onto $\ker H_j$ be those of Proposition~\ref{prop:reconstruction-direction}; one has $\langle[F],[G]\rangle=\langle\Theta_jF,G\rangle_{\rho}$ for $F,G\in\mathcal E_{+}$ (Proposition~\ref{prop:reconstruction-direction}(i)). Put $P=P^j_{\rho}$. Let $s\in\mathbb R$ and $y\in Y$. By $(T_j)$ and by the choice of $Y$, $(\tau_{s\cdot e_j+y})_{*}\rho=\rho$; hence $U(s\cdot e_j+y)=U(s\cdot e_j)U(y)$ is a unitary operator of $L^2(\rho)$ which commutes with $P$ (Lemma~\ref{lem:averages-translations}(v)), and the operator $V_j(y)$ is defined (Proposition~\ref{prop:reconstruction-direction}(iii)).

\emph{Case 1: $F=\Theta_jF'$ with $F',G\in\mathcal E_{+}$, and $s\ge0$.} By Lemma~\ref{lem:reflection-positivity-weights}(i), $PG\in\mathcal E^j(\rho)\subset\mathcal E_{+}$. Thus $(1-P)G\in\mathcal E_{+}$, and Proposition~\ref{prop:reconstruction-direction}(iv) gives
\[
\left[(1-P)G\right]=[G]-[PG]=(1-P_0)[G].
\]
The translation $\tau_y$ maps the half space $\{x_j>0\}$ onto itself, and $\tau_{s\cdot e_j}$ maps it into itself, because $s\ge0$; so $U(y)$ and $U(s\cdot e_j)$ map $\mathcal E_{+}$ into itself (Section~\ref{sec:euclidean-maps-actions}). By Proposition~\ref{prop:reconstruction-direction}(ii), (iii),
\[
\left[U(s\cdot e_j+y)(1-P)G\right]=T(s)V_j(y)(1-P_0)[G]=V_j(y)\exp(-s\cdot H_j)(1-P_0)[G],
\]
and therefore, $V_j(y)$ being unitary,
\begin{align*}
\left|\left\langle\Theta_jF',U(s\cdot e_j+y)(1-P)G\right\rangle_{\rho}\right|&=\left|\left\langle[F'],V_j(y)\exp(-s\cdot H_j)(1-P_0)[G]\right\rangle\right|\\
&\le\left\|[F']\right\|\cdot\left\|\exp(-s\cdot H_j)(1-P_0)[G]\right\|.
\end{align*}
The right side does not depend on $y\in Y$, and it tends to $0$ as $s\to\infty$. Indeed, $1-P_0$ is the spectral projection of $H_j\ge0$ for the interval $(0,\infty)$, and the functions $\lambda\mapsto\exp(-s\cdot\lambda)\cdot1_{(0,\infty)}(\lambda)$ on $[0,\infty)$ are bounded by $1$ and tend to $0$ at every point as $s\to\infty$; by the spectral theorem and dominated convergence, $\exp(-s\cdot H_j)(1-P_0)\psi\to0$ for every $\psi\in\mathcal H_j$.

\emph{Case 2: $F,G\in L^2(\rho)$, and $s\to+\infty$.} Let $\varepsilon>0$. The space $\mathcal A$ of Lemma~\ref{lem:density-exponentials} contains $\mathcal A^0$ and is therefore dense in $L^2(\rho)$, by Lemma~\ref{lem:density-exponentials}(i) for $O=\mathbb R^3$ and by $\Sigma_{\mathbb R^3}=\Sigma$ (Lemma~\ref{lem:countable-family-test}(ii)). Choose $F_1,G_1\in\mathcal A$ with $\|F-F_1\|<\varepsilon$ and $\|G-G_1\|<\varepsilon$. Each of $F_1$, $G_1$ is a finite linear combination of functions $\exp\left(iX(f)\right)$ with $f\in\mathcal D_0$, and these test forms have compact support; choose $r>0$ such that all of them have their support in $\{|x_j|<r\}$. For each such $f$ one has $\tau_{r\cdot e_j}\cdot f\in\mathcal D(\{x_j>0\})$ and $\tau_{-r\cdot e_j}\cdot f\in\mathcal D(\{x_j<0\})$, since $(\tau_a\cdot f)(x)=f(x-a)$; as $U(a)X(f)=X(\tau_a\cdot f)$, it follows that
\[
G_2=U(r\cdot e_j)G_1\in\mathcal E_{+},\qquad U(-r\cdot e_j)F_1\in L^2(\Sigma_j^{-},\rho).
\]
By $(R_j)$ and Section~\ref{sec:euclidean-maps-actions}, $\Theta_j$ is a unitary operator of $L^2(\rho)$ with $\Theta_j^2=1$, and it maps $L^2(\Sigma_j^{-},\rho)$ onto $\mathcal E_{+}$, because $\theta_j\{x_j<0\}=\{x_j>0\}$. Hence $F_2=\Theta_jU(-r\cdot e_j)F_1$ belongs to $\mathcal E_{+}$, and $U(-r\cdot e_j)F_1=\Theta_jF_2$. Now $U(-r\cdot e_j)U(s\cdot e_j+y)=U\left((s-2r)\cdot e_j+y\right)U(r\cdot e_j)$, the operator $U(r\cdot e_j)$ commutes with $P$ (Lemma~\ref{lem:averages-translations}(v)), and $U(-r\cdot e_j)$ is unitary; therefore
\[
\left\langle F_1,U(s\cdot e_j+y)(1-P)G_1\right\rangle_{\rho}=\left\langle\Theta_jF_2,U\left((s-2r)\cdot e_j+y\right)(1-P)G_2\right\rangle_{\rho}\qquad(s\in\mathbb R,\ y\in Y).
\]
For $s\ge2r$, Case 1, applied to $F_2$, $G_2$ and $s-2r$ in place of $F'$, $G$ and $s$, bounds the modulus of the right side by a number which does not depend on $y\in Y$ and tends to $0$ as $s\to\infty$. Moreover $\left\|U(s\cdot e_j+y)(1-P)\right\|\le1$ and $\|F_1\|\le\|F\|+\varepsilon$, so that, for all $s\in\mathbb R$ and $y\in Y$,
\begin{multline*}
\left|\left\langle F,U(s\cdot e_j+y)(1-P)G\right\rangle_{\rho}-\left\langle F_1,U(s\cdot e_j+y)(1-P)G_1\right\rangle_{\rho}\right|\\
\le\|F-F_1\|\cdot\|G\|+\|F_1\|\cdot\|G-G_1\|\le\varepsilon\cdot\left(\|F\|+\|G\|+\varepsilon\right).
\end{multline*}
Consequently
\[
\limsup_{s\to+\infty}\,\sup_{y\in Y}\left|\left\langle F,U(s\cdot e_j+y)(1-P)G\right\rangle_{\rho}\right|\le\varepsilon\cdot\left(\|F\|+\|G\|+\varepsilon\right),
\]
and, $\varepsilon>0$ being arbitrary, the left side vanishes.

\emph{Case 3: $F,G\in L^2(\rho)$, and $s\to-\infty$.} The adjoint of $U(s\cdot e_j+y)$ is $U(-s\cdot e_j-y)$, and the self-adjoint operator $1-P$ commutes with both. Hence
\[
\left\langle F,U(s\cdot e_j+y)(1-P)G\right\rangle_{\rho}=\overline{\left\langle G,U\left((-s)\cdot e_j+(-y)\right)(1-P)F\right\rangle_{\rho}}\qquad(s\in\mathbb R,\ y\in Y).
\]
As $Y=-Y$, the supremum over $y\in Y$ of the modulus of the left side equals the supremum over $y'\in Y$ of $\left|\left\langle G,U\left((-s)\cdot e_j+y'\right)(1-P)F\right\rangle_{\rho}\right|$, and this tends to $0$ as $s\to-\infty$ by Case 2, applied to the pair $(G,F)$ in place of $(F,G)$.

This proves the first assertion. The range $\mathcal E^j(\rho)$ of $P$ consists of vectors invariant under all $U(s\cdot e_j)$, so that $U(s\cdot e_j)=P+U(s\cdot e_j)(1-P)$. Since $0\in Y$, the modulus of $\left\langle F,U(s\cdot e_j)(1-P)G\right\rangle_{\rho}$ is at most the supremum in the first assertion, and it tends to $0$ as $|s|\to\infty$; hence $\left\langle F,U(s\cdot e_j)G\right\rangle_{\rho}\to\langle F,PG\rangle_{\rho}$.
\end{proof}

\begin{proof}[Proof of Theorem~\ref{thm:invariant-vectors-clustering}]
\emph{Part (i).} For $k\ge1$ put
\[
b^{(k)}=\sum_{l\in K}a^{(k)}_l\cdot e_l\in\Gamma_K,\qquad c^{(k)}=\sum_{l\notin K}a^{(k)}_l\cdot e_l;
\]
then $a^{(k)}=b^{(k)}+c^{(k)}$, and $c^{(k)}\to\bar a$ as $k\to\infty$. If $K=\emptyset$, then $b^{(k)}=0$, in accordance with $\Gamma_{\emptyset}=\{0\}$ (Section~\ref{sec:euclidean-maps-actions}). Since $\rho$ satisfies $(T_{123})$, the map $a\mapsto U(a)$ is a strongly continuous unitary representation of $\mathbb R^3$ on $L^2(\rho)$ (Lemma~\ref{lem:averages-translations}(ii)). Hence $U\left(c^{(k)}\right)G\to U(\bar a)G$ in $L^2(\rho)$, and, $U\left(b^{(k)}\right)$ being unitary,
\[
\left|\left\langle F,U\left(a^{(k)}\right)G\right\rangle_{\rho}-\left\langle F,U\left(b^{(k)}\right)U(\bar a)G\right\rangle_{\rho}\right|\le\|F\|\cdot\left\|U\left(c^{(k)}\right)G-U(\bar a)G\right\|\to0.
\]
It therefore suffices to show that $U\left(b^{(k)}\right)$ converges weakly to $P^K_{\rho}$, that is, $\left\langle F,U\left(b^{(k)}\right)G\right\rangle_{\rho}\to\left\langle F,P^K_{\rho}G\right\rangle_{\rho}$ for all $F,G\in L^2(\rho)$.

Write $K=\{l_1,\dots,l_p\}$ with $p=\#|K|$, put $P^l=P^l_{\rho}$ for $l=1,2,3$, and let a product of no operators denote the identity. By Lemma~\ref{lem:averages-translations}(v), the projections $P^l$ commute with each other and with every $U(a)$, $a\in\mathbb R^3$, and $P^{l_1}\cdots P^{l_p}=P^K_{\rho}$ (for $K=\emptyset$ both sides are the identity, Section~\ref{sec:invariant-sets-invariant}). The range $\mathcal E^K(\rho)$ of $P^K_{\rho}$ consists of vectors invariant under every $U(b)$, $b\in\Gamma_K$; hence
\[
U(b)P^{l_1}\cdots P^{l_p}=P^K_{\rho}\qquad(b\in\Gamma_K).
\]
The identity operator of $L^2(\rho)$ decomposes as
\[
1=P^{l_1}\cdots P^{l_p}+\sum_{q=1}^pP^{l_1}\cdots P^{l_{q-1}}\left(1-P^{l_q}\right).
\]
With $G_q=P^{l_1}\cdots P^{l_{q-1}}G$, $q=1,\dots,p$, it follows that
\[
\left\langle F,U\left(b^{(k)}\right)G\right\rangle_{\rho}=\left\langle F,P^K_{\rho}G\right\rangle_{\rho}+\sum_{q=1}^p\left\langle F,U\left(b^{(k)}\right)\left(1-P^{l_q}\right)G_q\right\rangle_{\rho}.
\]
Fix $q\in\{1,\dots,p\}$. The vector $b^{(k)}-a^{(k)}_{l_q}\cdot e_{l_q}$ is orthogonal to $e_{l_q}$. The law $\rho$ satisfies $(T_{l_q})$, $(R_{l_q})$ and $(RP_{l_q})$, and the set $Y=e_{l_q}^{\perp}$ satisfies $0\in Y=-Y$ and, by $(T_{123})$, $(\tau_y)_{*}\rho=\rho$ for every $y\in Y$. Proposition~\ref{prop:uniform-mixing-along}, applied with $l_q$ in place of $j$, with this set $Y$ and with $G_q$ in place of $G$, gives
\[
\left|\left\langle F,U\left(b^{(k)}\right)\left(1-P^{l_q}\right)G_q\right\rangle_{\rho}\right|\le\sup_{y\perp e_{l_q}}\left|\left\langle F,U\left(a^{(k)}_{l_q}\cdot e_{l_q}+y\right)\left(1-P^{l_q}\right)G_q\right\rangle_{\rho}\right|\to0\qquad(k\to\infty),
\]
because $\left|a^{(k)}_{l_q}\right|\to\infty$. Hence $\left\langle F,U\left(b^{(k)}\right)G\right\rangle_{\rho}\to\left\langle F,P^K_{\rho}G\right\rangle_{\rho}$. Finally, $P^K_{\rho}=E_{\rho}\left[\,\cdot\mid\mathcal I^K\right]$ by Lemma~\ref{lem:averages-translations}(iv).

\emph{Part (ii).} Let $\rho\in\mathcal M$ satisfy $(T_j)$, $(R_j)$ and $(RP_j)$. By Lemma~\ref{lem:averages-translations}(iv), $\mathcal E^j(\rho)=L^2(\mathcal I^j,\rho)$; by Lemma~\ref{lem:reflection-positivity-weights}(i), this space is contained in $L^2(\Sigma_j^{+},\rho)$; and by Proposition~\ref{prop:reconstruction-direction}(iv), $F\mapsto[F]_j$ maps it isometrically onto $\ker H_j$. This map sends the constant function $1$ to $\Omega_j$. Hence $\ker H_j=\mathbb C\cdot\Omega_j$ if and only if $\mathcal E^j(\rho)=\mathbb C\cdot1$, that is, if and only if $\rho$ is $\Gamma_j$-ergodic; so (1) and (2) are equivalent.

Assume (2). Then $P^j_{\rho}G=\langle1,G\rangle_{\rho}\cdot1=\rho(G)\cdot1$ for $G\in L^2(\rho)$, and the last assertion of Proposition~\ref{prop:uniform-mixing-along}, with $Y=\{0\}$, gives
\[
\rho\left(\overline{F}\cdot U(s\cdot e_j)G\right)=\left\langle F,U(s\cdot e_j)G\right\rangle_{\rho}\to\left\langle F,P^j_{\rho}G\right\rangle_{\rho}=\overline{\rho(F)}\cdot\rho(G)\qquad(|s|\to\infty),
\]
which is (3).

Assume (3), and let $F,G\in L^2(\rho)$. The function $h\colon[0,\infty)\to\mathbb C$, $h(s)=\rho\left(\overline{F}\cdot U(s\cdot e_j)G\right)$, is continuous by Lemma~\ref{lem:averages-translations}(ii) and bounded by $\|F\|\cdot\|G\|$, and $h(s)\to\overline{\rho(F)}\cdot\rho(G)$ as $s\to\infty$. Hence its Ces\`aro means $T^{-1}\int_0^T\mathrm{d}s\,h(s)$ converge to the same limit as $T\to\infty$; this is (4).

Assume (4), and let $F,G\in L^2(\rho)$. By Fubini's theorem,
\[
T^{-1}\int_0^T\mathrm{d}s\,\rho\left(\overline{F}\cdot U(s\cdot e_j)G\right)=\left\langle F,A^j_TG\right\rangle_{\rho},
\]
where $A^j_T$ is the average of Lemma~\ref{lem:averages-translations} for $J=\{j\}$, and $A^j_TG\to P^j_{\rho}G$ in $L^2(\rho)$ as $T\to\infty$, by Lemma~\ref{lem:averages-translations}(iii). So (4) states that $\left\langle F,P^j_{\rho}G\right\rangle_{\rho}=\left\langle F,\rho(G)\cdot1\right\rangle_{\rho}$ for all $F,G\in L^2(\rho)$, that is, $P^j_{\rho}G=\langle1,G\rangle_{\rho}\cdot1$ for every $G$. The range $\mathcal E^j(\rho)$ of $P^j_{\rho}$ is then $\mathbb C\cdot1$, which is (2).

\emph{Part (iii).} Let $\rho\in\mathcal C$ be $\Gamma_{123}$-ergodic. Then $\mathcal E^{123}(\rho)=\mathbb C\cdot1$ and $P^{123}_{\rho}G=\rho(G)\cdot1$ for $G\in L^2(\rho)$. Let $F,G\in L^2(\rho)$ and $j\in\{1,2,3\}$.

\emph{Assertion 1.} Let $\left(a^{(k)}\right)_{k\ge1}$ be a sequence in $\mathbb R^3$ with $\min_l\left|a^{(k)}_l\right|\to\infty$. Part (i), with $K=\{1,2,3\}$ and $\bar a=0$, gives
\[
\rho\left(\overline{F}\cdot U\left(a^{(k)}\right)G\right)=\left\langle F,U\left(a^{(k)}\right)G\right\rangle_{\rho}\to\left\langle F,P^{123}_{\rho}G\right\rangle_{\rho}=\overline{\rho(F)}\cdot\rho(G).
\]
If $n\in\mathbb R^3$ and $n_1\cdot n_2\cdot n_3\ne0$, then $\min_l|s\cdot n_l|=s\cdot\min_l|n_l|\to\infty$ as $s\to\infty$; this gives the assertion along the rays.

\emph{Assertion 2.} By Fubini's theorem,
\[
T^{-3}\int_{[0,T]^3}\mathrm{d}a\,\rho\left(\overline{F}\cdot U(a)G\right)=\left\langle F,A^{123}_TG\right\rangle_{\rho},
\]
and $A^{123}_TG\to P^{123}_{\rho}G=\rho(G)\cdot1$ in $L^2(\rho)$ as $T\to\infty$, by Lemma~\ref{lem:averages-translations}(iii).

\emph{Assertion 3.} Let $\psi\in\mathcal H_j$ satisfy $\exp(-t\cdot H_j)\psi=\psi$ for $t\ge0$ and $V_j(y)\psi=\psi$ for $y\perp e_j$. By the spectral theorem, the identity $\exp(-H_j)\psi=\psi$ implies that the spectral measure of $\psi$ with respect to $H_j$ is carried by $\left\{\lambda\ge0:\exp(-\lambda)=1\right\}=\{0\}$, that is, $\psi\in\ker H_j$. By Proposition~\ref{prop:reconstruction-direction}(iv) there is $F\in\mathcal E^j(\rho)$ with $\psi=[F]_j$. Let $y\perp e_j$. The operator $U(y)$ commutes with every $U(s\cdot e_j)$, so that $U(y)F\in\mathcal E^j(\rho)$, and Proposition~\ref{prop:reconstruction-direction}(iii) gives
\[
\left[U(y)F\right]_j=V_j(y)[F]_j=V_j(y)\psi=\psi=[F]_j.
\]
As $[\cdot]_j$ is isometric, hence injective, on $\mathcal E^j(\rho)$, one obtains $U(y)F=F$. Thus $F$ is invariant under every $U(s\cdot e_j)$, $s\in\mathbb R$, and under every $U(y)$, $y\perp e_j$, hence under every $U(a)$, $a\in\mathbb R^3$. So $F\in\mathcal E^{123}(\rho)=\mathbb C\cdot1$, and $\psi$ is a multiple of $[1]_j=\Omega_j$.

\emph{Assertion 4.} Let $(s_k)_{k\ge1}$ be a real sequence with $s_k\to+\infty$ or with $s_k\to-\infty$. Part (i), with $K=\{j\}$, $a^{(k)}=s_k\cdot e_j$ and $\bar a=0$, gives
\[
\rho\left(\overline{F}\cdot U(s_k\cdot e_j)G\right)\to\left\langle F,P^j_{\rho}G\right\rangle_{\rho}=\left\langle P^j_{\rho}F,P^j_{\rho}G\right\rangle_{\rho}=\rho\left(\overline{E_{\rho}\left[F\mid\mathcal I^j\right]}\cdot E_{\rho}\left[G\mid\mathcal I^j\right]\right),
\]
since $P^j_{\rho}$ is an orthogonal projection and $P^j_{\rho}=E_{\rho}\left[\,\cdot\mid\mathcal I^j\right]$ by part (i).

\emph{Assertion 5.} Let $d=\dim\ker H_j$, and assume $d<\infty$. By part (ii), $d=\dim L^2(\mathcal I^j,\rho)$; moreover $d\ge1$, since $\Omega_j\in\ker H_j$ and $\|\Omega_j\|=1$ (Proposition~\ref{prop:reconstruction-direction}(i), (ii)). If $B_1,\dots,B_k\in\mathcal I^j$ are pairwise disjoint and of positive measure, then $1_{B_1},\dots,1_{B_k}$ are linearly independent elements of $L^2(\mathcal I^j,\rho)$, and therefore $k\le d$. Consequently, among the partitions of $Q$ into finitely many sets of $\mathcal I^j$ of positive measure (the partition $\{Q\}$ is one of them), there is one, $A_1,\dots,A_{d'}$, with the largest number $d'\le d$ of parts. If some $B\in\mathcal I^j$ satisfied $0<\rho(B\cap A_p)<\rho(A_p)$, then the two sets $B\cap A_p$ and $A_p\setminus B$ in place of $A_p$ would give such a partition with more parts. Hence, for every $B\in\mathcal I^j$ and every $p$, either $\rho(B\cap A_p)=0$ or $\rho(A_p\setminus B)=0$; in other words, every element of $\mathcal I^j$ coincides, up to a $\rho$-null set, with a union of some of the sets $A_p$. Applied to the sets $\{\operatorname{Re}u\le c\}$ and $\{\operatorname{Im}u\le c\}$, $c\in\mathbb Q$, this shows that every $\mathcal I^j$-measurable function $u\colon Q\to\mathbb C$ is almost everywhere constant on each $A_p$. Therefore $1_{A_1},\dots,1_{A_{d'}}$ span $L^2(\mathcal I^j,\rho)$, and $d'=d$.

Let $y\in\mathbb R^3$. Translations commute with each other, so $\tau_y$ and $\tau_{-y}$ map $\mathcal I^j$ onto itself; they preserve $\rho$ by $(T_{123})$. Fix $p$. The set $\tau_yA_p\in\mathcal I^j$ has positive measure, so it contains some $A_q$ up to a null set. Then $\tau_{-y}A_q\in\mathcal I^j$ has positive measure and is contained in $A_p$ up to a null set; by the dichotomy above, $\tau_{-y}A_q$ coincides with $A_p$ up to a null set, that is, $\rho\left(\tau_yA_p\,\triangle\,A_q\right)=0$. The index $q$ is unique, the sets $A_q$ being disjoint and of positive measure; write $q=\sigma_y(p)$. The sets $\tau_yA_1,\dots,\tau_yA_d$ are pairwise disjoint, so $\sigma_y$ is injective, hence a permutation of $\{1,\dots,d\}$. Since $U(y)1_A=1_{\tau_yA}$ for $A\in\Sigma$ (Section~\ref{sec:euclidean-maps-actions}),
\[
U(y)1_{A_p}=1_{A_{\sigma_y(p)}}\quad\text{in }L^2(\rho)\qquad(p=1,\dots,d).
\]
For $p\ne q$ one has $\left\|1_{A_p}-1_{A_q}\right\|^2=\rho(A_p)+\rho(A_q)\ge2\cdot\min_r\rho(A_r)>0$. By Lemma~\ref{lem:averages-translations}(ii) there is $\delta>0$ such that $\left\|U(y)1_{A_p}-1_{A_p}\right\|^2<2\cdot\min_r\rho(A_r)$ for $p=1,\dots,d$ and $|y|<\delta$; for such $y$ the permutation $\sigma_y$ fixes every $p$. From $U(y+y')=U(y)U(y')$ one obtains $\sigma_{y+y'}=\sigma_y\circ\sigma_{y'}$; and every $y\in\mathbb R^3$ is the sum of $n$ copies of $n^{-1}\cdot y$, where $\left|n^{-1}\cdot y\right|<\delta$ for large $n$. Hence $\sigma_y$ fixes every $p$, for every $y\in\mathbb R^3$, and $1_{A_p}\in\mathcal E^{123}(\rho)=\mathbb C\cdot1$ for $p=1,\dots,d$. An indicator function which is almost everywhere constant and has positive integral equals $1$ almost everywhere; so $\rho(A_p)=1$ for every $p$, and $d=1$, the sets $A_p$ being disjoint. Therefore $\dim\ker H_j\in\{1,\infty\}$.

\emph{Part (iv).} Let $\rho\in\mathcal C$ satisfy \eqref{eq:setup-G}. By Lemma~\ref{lem:consequences-gaussian-bound}(i), for every $f\in\mathcal S$ and every integer $n\ge1$,
\[
\rho\left(X(f)\right)=0,\qquad\left\|X(f)\right\|_{L^n(\rho)}\le\left(n\cdot2^{\frac{n}{2}}\cdot\Gamma\left(\frac{n}{2}\right)\right)^{\frac{1}{n}}\cdot\|f\|.
\]
In particular $S^{\rho}_1=0$. By H\"older's inequality every finite product of fields $X(f)$, $f\in\mathcal S$, lies in $L^2(\rho)$; so all polynomials in the fields lie in $L^2(\rho)$, and (i)--(iii) apply to them. H\"older's inequality for $n$ factors, each in $L^n(\rho)$, gives for $f_1,\dots,f_n\in\mathcal S$
\[
\left|S^{\rho}_n(f_1,\dots,f_n)\right|\le\prod_{p=1}^n\left\|X(f_p)\right\|_{L^n(\rho)}\le n\cdot2^{\frac{n}{2}}\cdot\Gamma\left(\frac{n}{2}\right)\cdot\prod_{p=1}^n\|f_p\|.
\]

Let $j\in\{1,2,3\}$ and $f\in\mathcal S$. The law $\rho$ satisfies $(T_j)$, so Theorem~\ref{thm:ergodic-decomposition} applies to $m=\rho$ and $J=\{j\}$; let $(\rho_{\omega})_{\omega\in Q}$ and $\nu$ be the kernel and the measure which it provides. As $X(f)\in L^1(\rho)$, Lemma~\ref{lem:averages-translations}(iv) and \eqref{eq:ergodic-components-1} give
\[
P^j_{\rho}X(f)=E_{\rho}\left[X(f)\mid\mathcal I^j\right]=\rho_{\omega}\left(X(f)\right)\qquad\text{for $\rho$-almost every }\omega\in Q.
\]
By Theorem~\ref{thm:ergodic-decomposition}(c), $\nu$-almost every law satisfies \eqref{eq:setup-G}; as $\nu$ is the image of $\rho$ under $\omega\mapsto\rho_{\omega}$ (Theorem~\ref{thm:ergodic-decomposition}(a)), the law $\rho_{\omega}$ satisfies \eqref{eq:setup-G} for $\rho$-almost every $\omega$, and then $\rho_{\omega}\left(X(f)\right)=0$ by Lemma~\ref{lem:consequences-gaussian-bound}(i). Hence $P^j_{\rho}X(f)=0$.

Let $f,g\in\mathcal S$, and let $\left(a^{(k)}\right)_{k\ge1}$ be a sequence in $\mathbb R^3$ with $\left|a^{(k)}\right|\to\infty$. Every subsequence has a further subsequence along which, for each $l\in\{1,2,3\}$, either $\left|a^{(k)}_l\right|\to\infty$ or $a^{(k)}_l$ converges to a number $\bar a_l\in\mathbb R$; the set $K$ of the indices $l$ of the first kind is not empty. As $X(f)$ is real valued and $U(a)X(g)=X(\tau_a\cdot g)$, part (i) gives, along such a subsequence,
\[
S^{\rho}_2\left(f,\tau_{a^{(k)}}\cdot g\right)=\left\langle X(f),U\left(a^{(k)}\right)X(g)\right\rangle_{\rho}\to\left\langle X(f),P^K_{\rho}X\left(\tau_{\bar a}\cdot g\right)\right\rangle_{\rho}=0,
\]
since $P^K_{\rho}=\prod_{l\in K}P^l_{\rho}$ and each $P^l_{\rho}$ annihilates $X\left(\tau_{\bar a}\cdot g\right)$. Thus every subsequence of $\left(S^{\rho}_2\left(f,\tau_{a^{(k)}}\cdot g\right)\right)_{k\ge1}$ has a further subsequence tending to $0$, and the whole sequence tends to $0$.
\end{proof}

\subsection{Covariance of the ergodic decomposition}\label{sec:covariance-ergodic-decomposition}

\begin{lemma}[Euclidean covariance of the kernel]\label{lem:euclidean-covariance-kernel}
Let $J$ be a non-empty subset of $\{1,2,3\}$, let $m\in\mathcal M$ satisfy $(T_J)$, and let $(m_{\omega})_{\omega\in Q}$ and $\nu$ be the probability kernel and the probability measure given by Theorem~\ref{thm:ergodic-decomposition} for $m$ and $J$. Let $g\colon x\mapsto Rx+a$, with $R\in O(3)$ and $a\in\mathbb R^3$, be a Euclidean map such that $g_{*}m=m$ and $R\Gamma_J=\Gamma_J$. Then $g\mathcal I^J=\mathcal I^J$, where $g\mathcal I^J=\left\{gA:A\in\mathcal I^J\right\}$, and $\hat gP^J_m=P^J_m\hat g$ as operators on $L^2(m)$. Moreover, $g_{*}m_{\omega}=m_{g\omega}$ for $m$-almost every $\omega\in Q$, and the map $\rho\mapsto g_{*}\rho$ of $\mathcal M$ into itself leaves $\nu$ invariant.
\end{lemma}

\begin{proof}
For $b\in\mathbb R^3$ one has $g\circ\tau_b\circ g^{-1}=\tau_{Rb}$ as maps of $\mathbb R^3$, hence also as maps of $Q$ (Section~\ref{sec:euclidean-maps-actions}). Let $A\in\mathcal I^J$ and $b\in\Gamma_J$. Then $R^{-1}b\in\Gamma_J$, and
\[
\tau_b(gA)=g\left(\tau_{R^{-1}b}A\right)=gA;
\]
moreover $gA\in\Sigma$, since $gA$ is the preimage of $A$ under the measurable map $\omega\mapsto g^{-1}\omega$. Thus $gA\in\mathcal I^J$. The same argument applies to $g^{-1}\colon x\mapsto R^{-1}x-R^{-1}a$, for which $R^{-1}\Gamma_J=\Gamma_J$. Therefore $g\mathcal I^J=\mathcal I^J$ and $g^{-1}\mathcal I^J=\mathcal I^J$.

As $g_{*}m=m$, the operator $\hat g$ is unitary on $L^2(m)$, and its inverse $\widehat{g^{-1}}$ acts by $\widehat{g^{-1}}F=F\circ g$ (Section~\ref{sec:euclidean-maps-actions}). For a bounded measurable function $\Phi$ on $Q$ let $m_{\bullet}(\Phi)\colon Q\to\mathbb C$ be the function $\omega\mapsto m_{\omega}(\Phi)$. By \eqref{eq:ergodic-components-1}, $m_{\bullet}(\Phi)$ is bounded and $\mathcal I^J$-measurable, and it coincides $m$-almost everywhere with $E_m\left[\Phi\mid\mathcal I^J\right]$. Since $g^{-1}\mathcal I^J=\mathcal I^J$, the function $m_{\bullet}(\Phi)\circ g$ is $\mathcal I^J$-measurable as well. Let $I\in\mathcal I^J$; then $gI\in\mathcal I^J$ and $1_I=1_{gI}\circ g$. Hence
\begin{align*}
m\left(1_I\cdot\left(m_{\bullet}(\Phi)\circ g\right)\right)&=m\left(\left(1_{gI}\cdot m_{\bullet}(\Phi)\right)\circ g\right)=m\left(1_{gI}\cdot m_{\bullet}(\Phi)\right)\\
&=m\left(1_{gI}\cdot\Phi\right)=m\left(\left(1_{gI}\cdot\Phi\right)\circ g\right)=m\left(1_I\cdot(\Phi\circ g)\right):
\end{align*}
the first and the last equality follow from $1_I=1_{gI}\circ g$, the second and the fourth from $g_{*}m=m$, and the third from \eqref{eq:ergodic-components-1}, since $1_{gI}$ is a bounded $\mathcal I^J$-measurable function. As $I\in\mathcal I^J$ is arbitrary, $m_{\bullet}(\Phi)\circ g=E_m\left[\Phi\circ g\mid\mathcal I^J\right]$ $m$-almost everywhere.

By Lemma~\ref{lem:averages-translations}(iv), $P^J_m=E_m\left[\,\cdot\mid\mathcal I^J\right]$. Since composition with $g$ preserves equality $m$-almost everywhere, the last identity says
\[
\widehat{g^{-1}}P^J_m\Phi=\left(P^J_m\Phi\right)\circ g=m_{\bullet}(\Phi)\circ g=P^J_m(\Phi\circ g)=P^J_m\widehat{g^{-1}}\Phi
\]
for every bounded measurable $\Phi$. Such functions are dense in $L^2(m)$, and both sides are bounded operators on $L^2(m)$; hence $\widehat{g^{-1}}P^J_m=P^J_m\widehat{g^{-1}}$, and multiplication by $\hat g$ from the left and from the right gives $\hat gP^J_m=P^J_m\hat g$. Evaluated at $\omega$, the identity $m_{\bullet}(\Phi)\circ g=E_m\left[\Phi\circ g\mid\mathcal I^J\right]$ reads, by \eqref{eq:ergodic-components-1} applied to $\Phi\circ g$,
\[
m_{g\omega}(\Phi)=m_{\omega}(\Phi\circ g)=\left(g_{*}m_{\omega}\right)(\Phi)\qquad\text{for $m$-almost every }\omega\in Q,
\]
with an exceptional null set which depends on $\Phi$.

Apply this to $\Phi=\exp\left(iX(f)\right)$ for the countably many $f\in\mathcal D_0$. Outside an $m$-null set one has $\left(g_{*}m_{\omega}\right)\left(\exp\left(iX(f)\right)\right)=m_{g\omega}\left(\exp\left(iX(f)\right)\right)$ for every $f\in\mathcal D_0$, and there $g_{*}m_{\omega}=m_{g\omega}$ by Lemma~\ref{lem:density-exponentials}(ii).

Finally, $\nu$ is the image of $m$ under the map $\omega\mapsto m_{\omega}$ (Theorem~\ref{thm:ergodic-decomposition}(a)), and the map $\rho\mapsto g_{*}\rho$ is measurable for $\mathcal B_{\mathcal M}$, because $(g_{*}\rho)(A)=\rho\left(g^{-1}A\right)$ for $A\in\Sigma$. The image of $\nu$ under $\rho\mapsto g_{*}\rho$ is therefore the image of $m$ under $\omega\mapsto g_{*}m_{\omega}$. This map coincides $m$-almost everywhere with $\omega\mapsto m_{g\omega}$, and the image of $m$ under $\omega\mapsto m_{g\omega}$ is the image of $g_{*}m=m$ under $\omega\mapsto m_{\omega}$, that is, $\nu$.
\end{proof}

\begin{proof}[Proof of Theorem~\ref{thm:cubic-symmetry-components}]
By \eqref{eq:setup-H}, $\mu$ satisfies $(T_{123})$ and $(RP_j)$ for $j=1,2,3$, and $w_{*}\mu=\mu$ for every $w\in O_{\mathrm{cub}}$; in particular $\mu$ satisfies $(R_j)$ for $j=1,2,3$ (Section~\ref{sec:hypothesis-law}).

\emph{Part (i).} Theorem~\ref{thm:ergodic-decomposition}(a) for $J=\{1,2,3\}$ gives $\nu\left(\mathcal M^{123}_{\mathrm{erg}}\right)=1$, so $\nu$-almost every $\rho$ satisfies $(T_{123})$. Theorem~\ref{thm:ergodic-decomposition}(b), applied with $j=1,2,3$, shows that $\nu$-almost every $\rho$ satisfies $(R_1)$, $(R_2)$ and $(R_3)$. The translations $\tau_a$, $a\in\mathbb R^3$, and the reflections $\theta_1$, $\theta_2$, $\theta_3$ generate the group $\mathbb R^3\rtimes N_3$, because $N_3$ is generated by $\theta_1,\theta_2,\theta_3$ (Section~\ref{sec:euclidean-maps-actions}). A law which is invariant under each element of a set of generators of a group of Euclidean maps is invariant under the whole group, since $(g\circ h)_{*}\rho=g_{*}(h_{*}\rho)$, and $(g^{-1})_{*}\rho=\rho$ whenever $g_{*}\rho=\rho$. This proves (i).

\emph{Part (ii).} Let $w\in O_{\mathrm{cub}}$. The Euclidean map $x\mapsto wx$ satisfies $w_{*}\mu=\mu$ and $w\Gamma_{123}=\mathbb R^3=\Gamma_{123}$. Lemma~\ref{lem:euclidean-covariance-kernel}, applied to $m=\mu$, $J=\{1,2,3\}$ and $g=w$, shows that $w_{*}\mu_{\omega}=\mu_{w\omega}$ for $\mu$-almost every $\omega\in Q$ and that $\rho\mapsto w_{*}\rho$ leaves $\nu$ invariant.

Let $\rho\in\mathcal C\cap(\mathrm G)\cap\mathcal M^{123}_{\mathrm{erg}}$, and write $w\,e_i=\varepsilon_i\cdot e_{\pi(i)}$, $i=1,2,3$, with $\varepsilon\in\{\pm1\}^3$ and $\pi\in S_3$ (Section~\ref{sec:euclidean-maps-actions}). For a function $F$ on $Q$ let $F\circ w$ be the function $\omega\mapsto F(w\omega)$; thus $F\circ w=\widehat{w^{-1}}F$. The map $F\mapsto F\circ w$ is a unitary operator of $L^2(w_{*}\rho)$ onto $L^2(\rho)$, with inverse $F\mapsto F\circ w^{-1}$. For $a\in\mathbb R^3$ and $i\in\{1,2,3\}$ one has $w\circ\tau_a\circ w^{-1}=\tau_{wa}$ and $w\circ\theta_i\circ w^{-1}=\theta_{\pi(i)}$; consequently, for every function $F$ on $Q$,
\[
\left(U(a)F\right)\circ w=U\left(w^{-1}a\right)(F\circ w),\qquad\left(\Theta_{\pi(i)}F\right)\circ w=\Theta_i(F\circ w).
\]
Moreover $(wx)_{\pi(i)}=\varepsilon_i\cdot x_i$ for $x\in\mathbb R^3$, so that $w^{-1}\{x_{\pi(i)}>0\}=\{\varepsilon_i\cdot x_i>0\}$; by Section~\ref{sec:euclidean-maps-actions}, the map $F\mapsto F\circ w$ therefore maps $L^2\left(\Sigma_{\{x_{\pi(i)}>0\}},w_{*}\rho\right)$ onto $L^2\left(\Sigma_{\{\varepsilon_i\cdot x_i>0\}},\rho\right)$.

These relations show that $w_{*}\rho\in\mathcal C\cap(\mathrm G)\cap\mathcal M^{123}_{\mathrm{erg}}$. First, $(\tau_a)_{*}(w_{*}\rho)=w_{*}\left((\tau_{w^{-1}a})_{*}\rho\right)=w_{*}\rho$ for $a\in\mathbb R^3$, and $(\theta_{\pi(i)})_{*}(w_{*}\rho)=w_{*}\left((\theta_i)_{*}\rho\right)=w_{*}\rho$ for $i=1,2,3$; so $w_{*}\rho$ satisfies $(T_{123})$, and $(R_j)$ for $j=1,2,3$. Second, a function $F\in L^2(w_{*}\rho)$ is invariant under all $U(a)$, $a\in\mathbb R^3$, if and only if $F\circ w\in L^2(\rho)$ is; hence $\mathcal E^{123}(w_{*}\rho)=\mathbb C\cdot1$. Third, for $f\in\mathcal S$ one has $X(f)\circ w=X\left(w^{-1}\cdot f\right)$ and $\left\|w^{-1}\cdot f\right\|=\|f\|$ (Section~\ref{sec:euclidean-maps-actions}), whence
\[
\int\mathrm{d}(w_{*}\rho)\,\exp X(f)=\int\mathrm{d}\rho\,\exp X\left(w^{-1}\cdot f\right)\le\exp\left(\frac{1}{2}\left\|w^{-1}\cdot f\right\|^2\right)=\exp\left(\frac{1}{2}\|f\|^2\right),
\]
which is \eqref{eq:setup-G} for $w_{*}\rho$. Fourth, let $i\in\{1,2,3\}$ and $G\in L^2\left(\Sigma_{\pi(i)}^{+},w_{*}\rho\right)$, and put $G_1=G\circ w$, an element of $L^2\left(\Sigma_{\{\varepsilon_i\cdot x_i>0\}},\rho\right)$. Then
\[
\left\langle\Theta_{\pi(i)}G,G\right\rangle_{w_{*}\rho}=\left\langle\Theta_iG_1,G_1\right\rangle_{\rho}.
\]
If $\varepsilon_i=1$, the right side is non-negative by the implication (1)$\Rightarrow$(3) of Lemma~\ref{lem:reformulations-reflection-positivity}, applied to $\rho$. If $\varepsilon_i=-1$, then $G_1\in L^2(\Sigma_i^{-},\rho)$; the function $G_2=\Theta_iG_1$ lies in $L^2(\Sigma_i^{+},\rho)$ and satisfies $G_1=\Theta_iG_2$ (Section~\ref{sec:euclidean-maps-actions}), and, $\Theta_i$ being a unitary involution of $L^2(\rho)$, hence self-adjoint,
\[
\left\langle\Theta_iG_1,G_1\right\rangle_{\rho}=\left\langle G_2,\Theta_iG_2\right\rangle_{\rho}=\overline{\left\langle\Theta_iG_2,G_2\right\rangle_{\rho}}\ge0,
\]
again by the implication (1)$\Rightarrow$(3) of Lemma~\ref{lem:reformulations-reflection-positivity} for $\rho$. So the law $w_{*}\rho$, which satisfies $(R_{\pi(i)})$, satisfies condition (3) of Lemma~\ref{lem:reformulations-reflection-positivity} for the direction $\pi(i)$, and hence $(RP_{\pi(i)})$. As $\pi$ is a bijection of $\{1,2,3\}$, $w_{*}\rho$ satisfies $(RP_j)$ for $j=1,2,3$. Thus $\rho\mapsto w_{*}\rho$ maps $\mathcal C\cap(\mathrm G)\cap\mathcal M^{123}_{\mathrm{erg}}$ into itself; the same holds for $w^{-1}$ in place of $w$, and $\rho=w_{*}\left((w^{-1})_{*}\rho\right)$, so that the map is onto.

Let now $\rho$ be invariant under $\mathbb R^3\rtimes N_3$; by (i), this holds for $\nu$-almost every $\rho$. For $\eta\in N_3$ one has $(w\eta)_{*}\rho=w_{*}(\eta_{*}\rho)=w_{*}\rho$; so $w_{*}\rho$ depends on $w$ only through the coset $wN_3$, that is, through the class of $w$ in $O_{\mathrm{cub}}/N_3\cong S_3$. The stabiliser $\mathrm{Stab}=\left\{w\in O_{\mathrm{cub}}:w_{*}\rho=\rho\right\}$ of $\rho$ is a subgroup of $O_{\mathrm{cub}}$ which contains $N_3$, and $\mathrm{Stab}/N_3$ is a subgroup of $O_{\mathrm{cub}}/N_3\cong S_3$; by Lagrange's theorem, $\#|\mathrm{Stab}/N_3|$ divides $6$, so it is $1$, $2$, $3$ or $6$. The map $w\,\mathrm{Stab}\mapsto w_{*}\rho$ is a bijection of the set of the left cosets of $\mathrm{Stab}$ in $O_{\mathrm{cub}}$ onto the orbit $O_{\mathrm{cub}}\rho=\left\{w_{*}\rho:w\in O_{\mathrm{cub}}\right\}$. With $\#|O_{\mathrm{cub}}|=48$ and $\#|N_3|=8$ it follows that
\[
\#|O_{\mathrm{cub}}\rho|=\frac{\#|O_{\mathrm{cub}}|}{\#|\mathrm{Stab}|}=\frac{48}{8\cdot\#|\mathrm{Stab}/N_3|}=\frac{6}{\#|\mathrm{Stab}/N_3|}\in\{1,2,3,6\}.
\]
\end{proof}

\begin{proof}[Proof of Theorem~\ref{thm:orbit-averages-invariance}]
Let $\mu$, $(\mu_{\omega})_{\omega\in Q}$, $\nu$ and $\bar\mu_{\omega}$, $\omega\in Q$, be as in Theorem~\ref{thm:orbit-averages-invariance}, and let $P=P^{123}_{\mu}$. By Theorem~\ref{thm:ergodic-decomposition}(a) and its proof, the identities~\eqref{eq:ergodic-components-1} hold for $m=\mu$ and $J=\{1,2,3\}$, the map $\omega\mapsto\mu_{\omega}$ is measurable from $(Q,\mathcal I^{123})$ to $(\mathcal M,\mathcal B_{\mathcal M})$, and $\nu$ is the image of $\mu$ under this map.

\emph{Part (i).} Let $Q_0$ be the set of the points $\omega\in Q$ such that $\mu_{\omega}\in\mathcal C\cap(\mathrm G)\cap\mathcal M^{123}_{\mathrm{erg}}$. This set of laws belongs to $\mathcal B_{\mathcal M}$ by Lemma~\ref{lem:measurability-defining-properties}, so that $Q_0\in\mathcal I^{123}$; and $\nu$-almost every law belongs to it, by Theorem~\ref{thm:ergodic-decomposition}(a)--(c) in the form stated in Remark~\ref{rem:components-all-translations}. Hence
\[
\mu(Q_0)=\nu\left(\mathcal C\cap(\mathrm G)\cap\mathcal M^{123}_{\mathrm{erg}}\right)=1.
\]

\emph{The integral formula.} Let $A\in\Sigma$ and $w\in O_{\mathrm{cub}}$. The set $w^{-1}A$ belongs to $\Sigma$ (Section~\ref{sec:euclidean-maps-actions}), and $\left(w_{*}\mu_{\omega}\right)(A)=\mu_{\omega}\left(w^{-1}A\right)$ is an $\mathcal I^{123}$-measurable function of $\omega\in Q$. The second identity of~\eqref{eq:ergodic-components-1}, applied to $\Phi=1_{w^{-1}A}$, and the invariance $w_{*}\mu=\mu$, which is part of~\eqref{eq:setup-H}, give
\[
\int_Q\mu(\mathrm d\omega)\,\left(w_{*}\mu_{\omega}\right)(A)=\int_Q\mu(\mathrm d\omega)\,\mu_{\omega}\left(w^{-1}A\right)=\mu\left(w^{-1}A\right)=\left(w_{*}\mu\right)(A)=\mu(A);
\]
in other words, $\int_Q\mu(\mathrm d\omega)\,w_{*}\mu_{\omega}=w_{*}\mu=\mu$. The mean of these identities over $w\in O_{\mathrm{cub}}$ is
\[
\int_Q\mu(\mathrm d\omega)\,\bar\mu_{\omega}(A)=\mu(A)\qquad(A\in\Sigma),
\]
that is, $\mu=\int_Q\mu(\mathrm d\omega)\,\bar\mu_{\omega}$.

\emph{The orbit of $\mu_{\omega}$.} Fix $\omega\in Q_0$, and write $\mathsf O=O_{\mathrm{cub}}\mu_{\omega}$ for the orbit of $\mu_{\omega}$ and $\bar\mu=\bar\mu_{\omega}$. By Theorem~\ref{thm:cubic-symmetry-components}(ii), every map $\rho\mapsto w_{*}\rho$, $w\in O_{\mathrm{cub}}$, maps $\mathcal C\cap(\mathrm G)\cap\mathcal M^{123}_{\mathrm{erg}}$ onto itself; as $\mu_{\omega}$ belongs to this set, $\mathsf O\subset\mathcal C\cap(\mathrm G)\cap\mathcal M^{123}_{\mathrm{erg}}$. Let
\[
\mathrm{Stab}=\left\{w\in O_{\mathrm{cub}}:w_{*}\mu_{\omega}=\mu_{\omega}\right\}
\]
be the stabiliser of $\mu_{\omega}$, a subgroup of $O_{\mathrm{cub}}$. Since $(vw)_{*}\rho=v_{*}\left(w_{*}\rho\right)$ for $v,w\in O_{\mathrm{cub}}$ and $\rho\in\mathcal M$ (Section~\ref{sec:euclidean-maps-actions}), one has $w_{*}\mu_{\omega}=v_{*}\mu_{\omega}$ if and only if $v^{-1}w\in\mathrm{Stab}$. So, for $\rho=v_{*}\mu_{\omega}\in\mathsf O$, the elements $w\in O_{\mathrm{cub}}$ with $w_{*}\mu_{\omega}=\rho$ form the left coset $v\,\mathrm{Stab}$, which has $\#|\mathrm{Stab}|$ elements; in particular $\#|\mathsf O|\cdot\#|\mathrm{Stab}|=\#|O_{\mathrm{cub}}|=48$. Collecting the terms of the sum that defines $\bar\mu$ according to the value of $w_{*}\mu_{\omega}$, one finds
\[
\bar\mu=\frac{\#|\mathrm{Stab}|}{48}\cdot\sum_{\rho\in\mathsf O}\rho=\frac{1}{\#|\mathsf O|}\cdot\sum_{\rho\in\mathsf O}\rho.
\]
The law $\mu_{\omega}\in\mathcal C$ satisfies $(R_1)$, $(R_2)$ and $(R_3)$, and $N_3$ is generated by $\theta_1,\theta_2,\theta_3$ (Section~\ref{sec:euclidean-maps-actions}); hence $N_3\subset\mathrm{Stab}$, so that $\#|\mathrm{Stab}|\ge\#|N_3|=8$ and $\#|\mathsf O|\le6$.

\emph{The law $\bar\mu$ satisfies~\eqref{eq:setup-H}.} Every $\rho\in\mathsf O$ satisfies~\eqref{eq:setup-G}, $(T_{123})$, and $(RP_j)$ for $j=1,2,3$. Each of these conditions consists of linear equalities or inequalities in the law, one for each $f\in\mathcal S$, for each $a\in\mathbb R^3$, or for each choice of $n$, $c\in\mathbb C^n$ and $f_1,\dots,f_n\in\mathcal D(\{x_j>0\})$; therefore it is preserved under convex combinations of laws. Explicitly, $\bar\mu(\Phi)=\left(\#|\mathsf O|\right)^{-1}\cdot\sum_{\rho\in\mathsf O}\rho(\Phi)$ for every bounded or non-negative measurable function $\Phi$ on $Q$, whence
\[
\int\mathrm d\bar\mu\,\exp X(f)=\frac{1}{\#|\mathsf O|}\cdot\sum_{\rho\in\mathsf O}\int\mathrm d\rho\,\exp X(f)\le\exp\left(\frac{1}{2}\|f\|^2\right)\qquad(f\in\mathcal S);
\]
moreover $(\tau_a)_{*}\bar\mu=\left(\#|\mathsf O|\right)^{-1}\cdot\sum_{\rho\in\mathsf O}(\tau_a)_{*}\rho=\bar\mu$ for $a\in\mathbb R^3$; and, for $j\in\{1,2,3\}$, the sum in $(RP_j)$ formed with $\bar\mu$ is the mean over $\rho\in\mathsf O$ of the sums formed with $\rho$ and the same $n$, $c$, $f_1,\dots,f_n$, each of which is non-negative. Finally, $\bar\mu$ is invariant under $O_{\mathrm{cub}}$ by its definition: for $v\in O_{\mathrm{cub}}$ the map $w\mapsto vw$ is a bijection of $O_{\mathrm{cub}}$ onto itself, and
\[
v_{*}\bar\mu=\frac{1}{48}\sum_{w\in O_{\mathrm{cub}}}(vw)_{*}\mu_{\omega}=\frac{1}{48}\sum_{w\in O_{\mathrm{cub}}}w_{*}\mu_{\omega}=\bar\mu.
\]
Thus $\bar\mu$ satisfies~\eqref{eq:setup-H}. By Section~\ref{sec:euclidean-maps-actions}, the operators $U(a)$, $a\in\mathbb R^3$, and $\hat w$, $w\in O_{\mathrm{cub}}$, are therefore unitary on $L^2(\bar\mu)$. The operators $U(a)$ are unitary on every $L^2(\rho)$, $\rho\in\mathsf O$, as well, and on all these spaces they act by $\left(U(a)F\right)(\omega')=F(\tau_{-a}\omega')$.

\emph{Decomposition of $L^2(\bar\mu)$.} The laws $\rho\in\mathsf O$ are finitely many pairwise distinct elements of $\mathcal M^{123}_{\mathrm{erg}}$. Lemma~\ref{lem:clustering-criterion-separation}(b), applied with $J=\{1,2,3\}$, yields pairwise disjoint sets $B_{\rho}\in\mathcal I^{123}$, $\rho\in\mathsf O$, such that $\rho(B_{\rho})=1$ for every $\rho\in\mathsf O$. For $\rho\in\mathsf O$ and $\rho'\in\mathsf O\setminus\{\rho\}$ the set $B_{\rho}$ is contained in $Q\setminus B_{\rho'}$, so that $\rho'(B_{\rho})=0$. Consequently, for $\rho\in\mathsf O$ and every measurable function $F\colon Q\to\mathbb C$,
\[
\int\mathrm d\bar\mu\,\left|1_{B_{\rho}}\cdot F\right|^2=\frac{1}{\#|\mathsf O|}\cdot\sum_{\rho'\in\mathsf O}\int\mathrm d\rho'\,1_{B_{\rho}}\cdot|F|^2=\frac{1}{\#|\mathsf O|}\cdot\int\mathrm d\rho\,|F|^2.
\]
For $F=1$ this gives $\bar\mu(B_{\rho})=\left(\#|\mathsf O|\right)^{-1}$; the sets $B_{\rho}$ being disjoint, their union has $\bar\mu$-measure $1$, and $\sum_{\rho\in\mathsf O}1_{B_{\rho}}=1$ holds $\bar\mu$-almost everywhere. Moreover $\rho\le\#|\mathsf O|\cdot\bar\mu$ for $\rho\in\mathsf O$; thus every $\bar\mu$-null set is a $\rho$-null set, and every element $F$ of $L^2(\bar\mu)$ determines an element of $L^2(\rho)$, again denoted $F$.

It follows that $L^2(\bar\mu)$ is the orthogonal direct sum of its closed subspaces $1_{B_{\rho}}\cdot L^2(\bar\mu)$, $\rho\in\mathsf O$: every $F\in L^2(\bar\mu)$ is the sum of the functions $1_{B_{\rho}}\cdot F$, $\rho\in\mathsf O$, which have pairwise disjoint supports. Let $\rho\in\mathsf O$. The map that sends $G\in1_{B_{\rho}}\cdot L^2(\bar\mu)$ to its class in $L^2(\rho)$ is a linear bijection of $1_{B_{\rho}}\cdot L^2(\bar\mu)$ onto $L^2(\rho)$; its inverse is $F\mapsto1_{B_{\rho}}\cdot F$, which is well defined by the identity for $\int\mathrm d\bar\mu\,\left|1_{B_{\rho}}\cdot F\right|^2$ above, and that identity also gives
\[
\|G\|_{L^2(\bar\mu)}=\left(\#|\mathsf O|\right)^{-\frac{1}{2}}\cdot\|G\|_{L^2(\rho)}\qquad\left(G\in1_{B_{\rho}}\cdot L^2(\bar\mu)\right).
\]
Thus the summand $1_{B_{\rho}}\cdot L^2(\bar\mu)$ is the space $L^2(\rho)$, with its norm multiplied by $\left(\#|\mathsf O|\right)^{-\frac{1}{2}}$.

\emph{Vectors invariant under the translations.} Let $\rho\in\mathsf O$ and $a\in\mathbb R^3$. As $B_{\rho}\in\mathcal I^{123}$, one has $U(a)1_{B_{\rho}}=1_{\tau_aB_{\rho}}=1_{B_{\rho}}$; since $U(a)$ is multiplicative (Section~\ref{sec:euclidean-maps-actions}), $U(a)\left(1_{B_{\rho}}\cdot F\right)=1_{B_{\rho}}\cdot U(a)F$ for $F\in L^2(\bar\mu)$. Hence $U(a)$ maps each summand $1_{B_{\rho}}\cdot L^2(\bar\mu)$ into itself, and the identification of this summand with $L^2(\rho)$ carries it to the operator $U(a)$ of $L^2(\rho)$: the operators $U(a)$ act diagonally on the decomposition. An element $F\in L^2(\bar\mu)$ is therefore invariant under all $U(a)$, $a\in\mathbb R^3$, if and only if, for every $\rho\in\mathsf O$, the element $F$ of $L^2(\rho)$ lies in $\mathcal E^{123}(\rho)$. Every $\rho\in\mathsf O$ is $\Gamma_{123}$-ergodic, that is, $\mathcal E^{123}(\rho)=\mathbb C\cdot1$; so the condition means that for every $\rho\in\mathsf O$ there is $z_{\rho}\in\mathbb C$ with $1_{B_{\rho}}\cdot F=z_{\rho}\cdot1_{B_{\rho}}$ in $L^2(\bar\mu)$. Consequently
\[
\mathcal E^{123}(\bar\mu)=\operatorname{span}\left\{1_{B_{\rho}}:\rho\in\mathsf O\right\}.
\]
The functions $1_{B_{\rho}}$, $\rho\in\mathsf O$, are pairwise orthogonal elements of $L^2(\bar\mu)$, and none of them is $0$, since $\bar\mu(B_{\rho})>0$. They form a basis of $\mathcal E^{123}(\bar\mu)$, and
\[
\dim\mathcal E^{123}(\bar\mu_{\omega})=\#|\mathsf O|=\#|O_{\mathrm{cub}}\mu_{\omega}|\le6.
\]

\emph{The action of $O_{\mathrm{cub}}$; vectors invariant under $\mathbb R^3\rtimes O_{\mathrm{cub}}$.} Let $w\in O_{\mathrm{cub}}$ and $\rho\in\mathsf O$. By Section~\ref{sec:euclidean-maps-actions}, $\hat w1_{B_{\rho}}=1_{wB_{\rho}}$, where $wB_{\rho}\in\Sigma$, and
\[
\left(w_{*}\rho\right)(wB_{\rho})=\rho\left(w^{-1}wB_{\rho}\right)=\rho(B_{\rho})=1.
\]
Let $\rho''\in\mathsf O$ with $\rho''\ne w_{*}\rho$. The law $(w^{-1})_{*}\rho''$ lies in the orbit $\mathsf O$, and it is different from $\rho$, because $\rho''\ne w_{*}\rho$; therefore
\[
\rho''(wB_{\rho})=\left((w^{-1})_{*}\rho''\right)(B_{\rho})=0.
\]
On the other hand $\left(w_{*}\rho\right)(B_{w_{*}\rho})=1$ and $\rho''(B_{w_{*}\rho})=0$. Hence every law in $\mathsf O$ gives measure $0$ to the symmetric difference $wB_{\rho}\,\triangle\,B_{w_{*}\rho}$, and so does $\bar\mu$; that is,
\[
\hat w1_{B_{\rho}}=1_{B_{w_{*}\rho}}\qquad\text{$\bar\mu$-almost everywhere}.
\]
For fixed $w$ the map $\rho\mapsto w_{*}\rho$ is a permutation of $\mathsf O$, and for any $\rho,\rho'\in\mathsf O$ there is $w\in O_{\mathrm{cub}}$ with $w_{*}\rho=\rho'$, because $\mathsf O$ is an orbit of $O_{\mathrm{cub}}$. So the operators $\hat w$, $w\in O_{\mathrm{cub}}$, permute the basis $\left\{1_{B_{\rho}}:\rho\in\mathsf O\right\}$ of $\mathcal E^{123}(\bar\mu)$, and they permute it transitively.

Now let $F\in L^2(\bar\mu)$ be invariant under $\mathbb R^3\rtimes O_{\mathrm{cub}}$, that is, $\hat gF=F$ for every Euclidean map $g\colon x\mapsto wx+a$ with $w\in O_{\mathrm{cub}}$ and $a\in\mathbb R^3$. Then $U(a)F=F$ for every $a\in\mathbb R^3$, so that $F=\sum_{\rho\in\mathsf O}z_{\rho}\cdot1_{B_{\rho}}$ with numbers $z_{\rho}\in\mathbb C$; and, for every $w\in O_{\mathrm{cub}}$,
\[
\sum_{\rho\in\mathsf O}z_{\rho}\cdot1_{B_{\rho}}=F=\hat wF=\sum_{\rho\in\mathsf O}z_{\rho}\cdot1_{B_{w_{*}\rho}}.
\]
The coefficient of the basis vector $1_{B_{w_{*}\rho}}$ is $z_{w_{*}\rho}$ on the left side and $z_{\rho}$ on the right side; hence $z_{w_{*}\rho}=z_{\rho}$ for all $w\in O_{\mathrm{cub}}$ and $\rho\in\mathsf O$. By transitivity, all the numbers $z_{\rho}$ are equal to one number $z$, and $F=z\cdot\sum_{\rho\in\mathsf O}1_{B_{\rho}}=z$ $\bar\mu$-almost everywhere. The constants are invariant under $\mathbb R^3\rtimes O_{\mathrm{cub}}$; so they are the only elements of $L^2(\bar\mu_{\omega})$ with this invariance.

\emph{The limit.} Let $F,G\in L^2(\bar\mu)$ and $a\in\mathbb R^3$. The function $\overline{F}\cdot U(a)G$ is integrable for $\bar\mu$ and for every $\rho\in\mathsf O$, and
\[
\bar\mu\left(\overline{F}\cdot U(a)G\right)=\frac{1}{\#|\mathsf O|}\cdot\sum_{\rho\in\mathsf O}\rho\left(\overline{F}\cdot U(a)G\right).
\]
Every $\rho\in\mathsf O$ belongs to $\mathcal C$ and is $\Gamma_{123}$-ergodic. Theorem~\ref{thm:invariant-vectors-clustering}(iii)1, applied to $\rho$ and to the elements $F$, $G$ of $L^2(\rho)$, gives $\rho\left(\overline{F}\cdot U(a)G\right)\to\overline{\rho(F)}\cdot\rho(G)$ as $\min_l|a_l|\to\infty$. Summing over $\rho\in\mathsf O$, one obtains
\[
\lim_{\min_l|a_l|\to\infty}\bar\mu_{\omega}\left(\overline{F}\cdot U(a)G\right)=\frac{1}{\#|O_{\mathrm{cub}}\mu_{\omega}|}\cdot\sum_{\rho\in O_{\mathrm{cub}}\mu_{\omega}}\overline{\rho(F)}\cdot\rho(G).
\]
The four properties of (i) are thus established for every $\omega\in Q_0$, and $\mu(Q_0)=1$. This proves (i).

\emph{Part (ii).} Let $w\in O_{\mathrm{cub}}$. By~\eqref{eq:setup-H}, $w_{*}\mu=\mu$; so $\hat w$ is a unitary operator on $L^2(\mu)$ with inverse $\widehat{w^{-1}}$, and $\widehat{w^{-1}}G=G\circ w$ for every function $G$ on $Q$ (Section~\ref{sec:euclidean-maps-actions}). For $\rho\in\mathcal M$ and every bounded measurable function $G$ on $Q$,
\[
\left(w_{*}\rho\right)(G)=\rho(G\circ w)=\rho\left(\widehat{w^{-1}}G\right).
\]
The countable set $\mathcal A^0$ of Lemma~\ref{lem:density-exponentials} contains the functions $\exp\left(iX(f)\right)$, $f\in\mathcal D_0$; by Lemma~\ref{lem:density-exponentials}(ii), $w_{*}\rho=\rho$ holds if and only if $\rho\left(\widehat{w^{-1}}G\right)=\rho(G)$ for every $G\in\mathcal A^0$. For every bounded measurable $\Phi$ the map $\rho\mapsto\rho(\Phi)$ is $\mathcal B_{\mathcal M}$-measurable (for indicator functions of sets in $\Sigma$ by the definition of $\mathcal B_{\mathcal M}$, and in general by the monotone class theorem, as in the proof of Lemma~\ref{lem:measurability-defining-properties}). Hence the set $\mathcal M_w$ of the laws $\rho\in\mathcal M$ with $w_{*}\rho=\rho$ belongs to $\mathcal B_{\mathcal M}$. Since $\nu$ is the image of $\mu$ under $\omega\mapsto\mu_{\omega}$, condition (1) holds if and only if $\mu_{\omega}\in\mathcal M_w$ for $\mu$-almost every $\omega$; as $\mathcal A^0$ is countable, this means that, for every $G\in\mathcal A^0$,
\[
\mu_{\omega}\left(\widehat{w^{-1}}G\right)=\mu_{\omega}(G)\qquad\text{for $\mu$-almost every }\omega\in Q.
\]
By~\eqref{eq:ergodic-components-1} and Lemma~\ref{lem:averages-translations}(iv), every bounded measurable function $\Phi$ on $Q$ satisfies $\mu_{\omega}(\Phi)=E_{\mu}\left[\Phi\mid\mathcal I^{123}\right](\omega)=(P\Phi)(\omega)$ for $\mu$-almost every $\omega$. Therefore (1) holds if and only if
\[
P\widehat{w^{-1}}G=PG\qquad\text{for every }G\in\mathcal A^0.
\]
The set $\mathcal A^0$ is dense in $L^2(\mu)$, by Lemma~\ref{lem:density-exponentials}(i) with $O=\mathbb R^3$ and by $\Sigma_{\mathbb R^3}=\Sigma$ (Lemma~\ref{lem:countable-family-test}(ii)), and $P\widehat{w^{-1}}$ and $P$ are bounded operators on $L^2(\mu)$. So (1) is equivalent to $P\widehat{w^{-1}}=P$. Multiplying this identity from the right by $\hat w$ gives $P=P\hat w$, and multiplying $P=P\hat w$ from the right by $\widehat{w^{-1}}$ gives $P\widehat{w^{-1}}=P$ again. Thus (1) and (2) are equivalent.

The law $\mu$ belongs to $\mathcal C$ (Section~\ref{sec:hypothesis-law}). Let $F,G'\in L^2(\mu)$, and let $\left(a^{(k)}\right)_{k\ge1}$ be a sequence in $\mathbb R^3$ with $\min_l\left|a^{(k)}_l\right|\to\infty$. Theorem~\ref{thm:invariant-vectors-clustering}(i), applied to $\rho=\mu$ and $K=\{1,2,3\}$, for which $\bar a=0$ and $P^K_{\mu}=P$, gives $\left\langle F,U\left(a^{(k)}\right)G'\right\rangle_{\mu}\to\left\langle F,PG'\right\rangle_{\mu}$. As the sequence is arbitrary,
\[
\lim_{\min_l|a_l|\to\infty}\mu\left(\overline{F}\cdot U(a)G'\right)=\lim_{\min_l|a_l|\to\infty}\left\langle F,U(a)G'\right\rangle_{\mu}=\left\langle F,PG'\right\rangle_{\mu}.
\]
Let $F,G\in\mathcal A^0$. With $G'=\hat wG$ and with $G'=G$, this shows that the two limits in (3) exist and are equal to $\left\langle F,P\hat wG\right\rangle_{\mu}$ and to $\left\langle F,PG\right\rangle_{\mu}$, respectively. Hence (3) holds if and only if $\left\langle F,P\hat wG\right\rangle_{\mu}=\left\langle F,PG\right\rangle_{\mu}$ for all $F,G\in\mathcal A^0$. Both sides are continuous in $(F,G)\in L^2(\mu)\times L^2(\mu)$, and $\mathcal A^0$ is dense in $L^2(\mu)$; so (3) is equivalent to $P\hat w=P$, which is (2).
\end{proof}

\subsection{The decomposition along one coordinate axis}\label{sec:decomposition-along-one}

\begin{proof}[Proof of Theorem~\ref{thm:decomposition-along-coordinate}]
Since $m\in\mathcal{C}$, the law $m$ satisfies $(T_{123})$, $(R_j)$ and $(RP_j)$; as $\Gamma_j\subset\Gamma_{123}=\mathbb{R}^3$, it satisfies $(T_j)$ as well. By definition, $m^j_{\omega}$ and $\nu^j$ are the kernel and the measure which Theorem~\ref{thm:ergodic-decomposition} provides for the law $m$ and for $J=\{j\}$. Theorem~\ref{thm:ergodic-decomposition}(a) gives $\nu^j\left(\mathcal{M}^j_{\mathrm{erg}}\right)=1$, that is, $\nu^j$-almost every $\rho$ satisfies $(T_j)$ and is $\Gamma_j$-ergodic. Theorem~\ref{thm:ergodic-decomposition}(b), applied with $j\in J$, shows that $\nu^j$-almost every $\rho$ satisfies $(R_j)$ and is reflection positive at every plane $\{x_j=c\}$, $c\in\mathbb{R}$. If $m$ satisfies~\eqref{eq:setup-G}, then $\nu^j$-almost every $\rho$ satisfies~\eqref{eq:setup-G}, by Theorem~\ref{thm:ergodic-decomposition}(c).

Consequently $\nu^j$-almost every $\rho\in\mathcal{M}$ satisfies $(T_j)$ and $(R_j)$, is $\Gamma_j$-ergodic and is reflection positive at every plane $\{x_j=c\}$, $c\in\mathbb{R}$. Let $\rho$ be such a law. Reflection positivity at the plane $\{x_j=0\}$ is the property $(RP_j)$ (Section~\ref{sec:properties-laws}); thus $\rho$ satisfies $(T_j)$, $(R_j)$ and $(RP_j)$. Theorem~\ref{thm:invariant-vectors-clustering}(ii) applies to every law in $\mathcal{M}$ with these three properties, so it applies to $\rho$. Condition~(2) of Theorem~\ref{thm:invariant-vectors-clustering}(ii) holds, $\rho$ being $\Gamma_j$-ergodic. By the equivalence of~(2) with~(1) and with~(3) in that theorem, $\ker H_j(\rho)=\mathbb{C}\,\Omega_j$, and for all $F,G\in L^2(\rho)$
\[
\rho\left(\bar{F}\cdot U(se_j)G\right)\to\overline{\rho(F)}\cdot\rho(G)\qquad\text{as }s\in\mathbb{R},\ |s|\to\infty.
\]

Let now $g\colon x\mapsto Rx+a$, with $R\in O(3)$ and $a\in\mathbb{R}^3$, be a Euclidean map such that $g_{*}m=m$ and such that conjugation by $g$ maps the group of the translations $\tau_b$, $b\in\Gamma_j$, onto itself. One has $g\circ\tau_b\circ g^{-1}=\tau_{Rb}$ for every $b\in\mathbb{R}^3$, and $b\mapsto\tau_b$ is injective; hence $R\Gamma_j=\Gamma_j$. The law $m$ satisfies $(T_j)$, so Lemma~\ref{lem:euclidean-covariance-kernel} applies to $m$, $g$ and $J=\{j\}$. The kernel of Theorem~\ref{thm:ergodic-decomposition} for this $J$ is $m^j_{\omega}$, and the lemma gives $g_{*}m^j_{\omega}=m^j_{g\omega}$ for $m$-almost every $\omega\in Q$.
\end{proof}

\begin{proof}[Proof of Theorem~\ref{thm:vacuum-symmetries-components}(i)]
As in Theorem~\ref{thm:vacuum-symmetries-components}, let $\nu$ and $(m_{\omega})_{\omega\in Q}$, with $m_{\omega}=m^{123}_{\omega}$, be the probability measure and the probability kernel given by Theorem~\ref{thm:ergodic-decomposition} for the law $m$ and $J=\{1,2,3\}$, and let $\nu^j$ and $(m^j_{\omega})_{\omega\in Q}$ be those given for $J=\{j\}$. The law $m\in\mathcal C$ satisfies $(T_{123})$, and $(R_l)$, $(RP_l)$ for $l=1,2,3$. By Theorem~\ref{thm:ergodic-decomposition}(a), $\nu\left(\mathcal M^{123}_{\mathrm{erg}}\right)=1$; in particular $\nu$-almost every $\rho$ satisfies $(T_{123})$. By Theorem~\ref{thm:ergodic-decomposition}(b), applied with each $l\in\{1,2,3\}$ in the place of $j$, $\nu$-almost every $\rho$ satisfies $(R_l)$ and is reflection positive at the plane $\{x_l=0\}$, that is, satisfies $(RP_l)$ (Section~\ref{sec:properties-laws}), for $l=1,2,3$. Hence
\[
\rho\in\mathcal C\cap\mathcal M^{123}_{\mathrm{erg}}\qquad\text{for $\nu$-almost every }\rho\in\mathcal M.
\]
For every such $\rho$ and every $l\in\{1,2,3\}$ the reconstruction of $\rho$ along $e_l$ (Proposition~\ref{prop:reconstruction-direction}), with the operator $H_l(\rho)$ and the vector $\Omega_l$, is defined, and Theorem~\ref{thm:invariant-vectors-clustering} applies to $\rho$. Since $m$ satisfies $(T_{123})$, the map $a\mapsto U(a)$ is a strongly continuous unitary representation of $\mathbb R^3$ on $L^2(m)$ (Lemma~\ref{lem:averages-translations}(ii)), and the projections $P^1_m$, $P^2_m$, $P^3_m$ commute with each other and with every $U(a)$, $a\in\mathbb R^3$ (Lemma~\ref{lem:averages-translations}(v)). For $F,G\in L^2(m)$ and $a\in\mathbb R^3$ one has $m\left(\overline{F}\cdot U(a)G\right)=\left\langle F,U(a)G\right\rangle_m$.

\emph{The conditions (2) and (3) are equivalent.} Let $F,G\in L^2(m)$ and $y\in\mathbb R^3$, and put $y_{\perp}=y-y_j\cdot e_j$. Let $(s_k)_{k\ge1}$ be a real sequence with $s_k\to\infty$, and put $a^{(k)}=s_k\cdot e_j+y$. Then $\left|a^{(k)}_j\right|=\left|s_k+y_j\right|\to\infty$, while $a^{(k)}_l=y_l$ for $l\ne j$. Theorem~\ref{thm:invariant-vectors-clustering}(i), applied to $\rho=m$ and $K=\{j\}$, with $\bar a_l=y_l$ for $l\ne j$, so that $\bar a=\sum_{l\ne j}y_l\cdot e_l=y_{\perp}$, gives $\left\langle F,U\left(a^{(k)}\right)G\right\rangle_m\to\left\langle F,P^j_mU(y_{\perp})G\right\rangle_m$. The sequence $(s_k)_{k\ge1}$ being arbitrary, the limits in condition~(3) exist, and
\[
\lim_{s\to\infty}m\left(\overline{F}\cdot U(s\cdot e_j+y)G\right)=\left\langle F,P^j_mU(y_{\perp})G\right\rangle_m,\qquad\lim_{s\to\infty}m\left(\overline{F}\cdot U(s\cdot e_j)G\right)=\left\langle F,P^j_mG\right\rangle_m;
\]
the second identity is the case $y=0$ of the first. The range $\mathcal E^j(m)$ of $P^j_m$ consists of vectors which are invariant under every $U(t\cdot e_j)$, $t\in\mathbb R$, and $P^j_m$ commutes with $U(t\cdot e_j)$; hence $P^j_mU(t\cdot e_j)=U(t\cdot e_j)P^j_m=P^j_m$. As $U(y)=U(y_j\cdot e_j)U(y_{\perp})$, it follows that
\[
P^j_mU(y)=P^j_mU(y_{\perp})\qquad\left(y\in\mathbb R^3\right).
\]
Consequently, condition~(3) states that $\left\langle F,P^j_mU(y)G\right\rangle_m=\left\langle F,P^j_mG\right\rangle_m$ for all $F,G\in\mathcal A^0$ and all $y\in\mathbb Q^3$. The set $\mathcal A^0$ is dense in $L^2(m)$, by Lemma~\ref{lem:density-exponentials}(i) with $O=\mathbb R^3$ and because $\Sigma_{\mathbb R^3}=\Sigma$ (Lemma~\ref{lem:countable-family-test}(ii)), and both sides of the last identity are continuous in $(F,G)\in L^2(m)\times L^2(m)$. So condition~(3) is equivalent to $P^j_mU(y)=P^j_m$ for all $y\in\mathbb Q^3$; and, as $y\mapsto U(y)G$ is continuous for every $G\in L^2(m)$ and $\mathbb Q^3$ is dense in $\mathbb R^3$, it is equivalent to
\[
P^j_mU(y)=P^j_m\qquad\text{for every }y\in\mathbb R^3.
\]
The adjoint of $U(y)$ is $U(-y)$. Taking adjoints, one finds that the last property is equivalent to $U(-y)P^j_m=P^j_m$ for every $y\in\mathbb R^3$, that is, to the inclusion $\mathcal E^j(m)\subset\mathcal E^{123}(m)$ of the range of $P^j_m$ in the space of the vectors invariant under all translations. Since $\mathcal E^{123}(m)\subset\mathcal E^j(m)$ in any case, this inclusion means $\mathcal E^j(m)=\mathcal E^{123}(m)$, that is, $P^j_m=P^{123}_m$, which is condition~(2).

\emph{Condition (2) implies condition (1).} Assume~(2). By Lemma~\ref{lem:averages-translations}(iv), $P^j_m=E_m\left[\,\cdot\mid\mathcal I^j\right]$ and $P^{123}_m=E_m\left[\,\cdot\mid\mathcal I^{123}\right]$. Let $f\in\mathcal D_0$ and $\Phi=\exp\left(iX(f)\right)$. Then $E_m\left[\Phi\mid\mathcal I^j\right]=E_m\left[\Phi\mid\mathcal I^{123}\right]$ $m$-almost everywhere, and~\eqref{eq:ergodic-components-1}, for $J=\{j\}$ and for $J=\{1,2,3\}$, gives $m^j_{\omega}(\Phi)=m_{\omega}(\Phi)$ for $m$-almost every $\omega$. As $\mathcal D_0$ is countable, there is an $m$-null set outside of which $m^j_{\omega}\left(\exp\left(iX(f)\right)\right)=m_{\omega}\left(\exp\left(iX(f)\right)\right)$ for every $f\in\mathcal D_0$; outside this set $m^j_{\omega}=m_{\omega}$, by Lemma~\ref{lem:density-exponentials}(ii). By Theorem~\ref{thm:ergodic-decomposition}(a), $\nu^j$ and $\nu$ are the images of $m$ under the maps $\omega\mapsto m^j_{\omega}$ and $\omega\mapsto m_{\omega}$, which coincide $m$-almost everywhere; therefore $\nu^j=\nu$. By Theorem~\ref{thm:decomposition-along-coordinate}, $\ker H_j(\rho)=\mathbb C\cdot\Omega_j$ for $\nu^j$-almost every $\rho$, that is, for $\nu$-almost every $\rho$; this is~(1).

\emph{Condition (1) implies condition (4).} Assume~(1). Then $\nu$-almost every $\rho$ belongs to $\mathcal C\cap\mathcal M^{123}_{\mathrm{erg}}$ and has $\ker H_j(\rho)=\mathbb C\cdot\Omega_j$. Such a law $\rho$ satisfies $(T_{123})$, hence $(T_j)$, and $(R_j)$, $(RP_j)$; by the implication (1)$\Rightarrow$(2) of Theorem~\ref{thm:invariant-vectors-clustering}(ii), it is $\Gamma_j$-ergodic. As $m=\int_{\mathcal M}\mathrm{d}\nu(\rho)\,\rho$ (Theorem~\ref{thm:ergodic-decomposition}(a)), the measure $\nu'=\nu$ has the properties required in~(4).

\emph{Condition (4) implies condition (2).} Let $\nu'$ be a probability measure on $(\mathcal M,\mathcal B_{\mathcal M})$ with the properties required in~(4). Let $\rho\in\mathcal M$ satisfy $(T_{123})$ and be $\Gamma_j$-ergodic. Then $\rho$ satisfies $(T_j)$, as $\Gamma_j\subset\Gamma_{123}$, so that $\rho\in\mathcal M^j_{\mathrm{erg}}$; and $\mathbb C\cdot1\subset\mathcal E^{123}(\rho)\subset\mathcal E^j(\rho)=\mathbb C\cdot1$, so that $\rho\in\mathcal M^{123}_{\mathrm{erg}}$. The sets $\mathcal M^j_{\mathrm{erg}}$ and $\mathcal M^{123}_{\mathrm{erg}}$ belong to $\mathcal B_{\mathcal M}$ (Theorem~\ref{thm:ergodic-decomposition}(a)); hence $\nu'\left(\mathcal M^j_{\mathrm{erg}}\right)=\nu'\left(\mathcal M^{123}_{\mathrm{erg}}\right)=1$. The law $m$ satisfies $(T_j)$ and $(T_{123})$, and $m=\int_{\mathcal M}\mathrm{d}\nu'(\rho)\,\rho$. The uniqueness assertion of Theorem~\ref{thm:ergodic-decomposition}(a), for $J=\{j\}$ and for $J=\{1,2,3\}$, gives
\[
\nu'=\nu^j\qquad\text{and}\qquad\nu'=\nu.
\]
Let $F\in\mathcal A^0$. The function $F$ is bounded and $\Sigma$-measurable, so that $\rho\mapsto\left|\rho(F)\right|^2$ is a bounded $\mathcal B_{\mathcal M}$-measurable function on $\mathcal M$ (proof of Lemma~\ref{lem:measurability-defining-properties}). By Lemma~\ref{lem:averages-translations}(iv) and~\eqref{eq:ergodic-components-1} for $J=\{j\}$, the function $P^j_mF=E_m\left[F\mid\mathcal I^j\right]$ coincides $m$-almost everywhere with $\omega\mapsto m^j_{\omega}(F)$. Since $\nu^j$ is the image of $m$ under $\omega\mapsto m^j_{\omega}$,
\[
\left\|P^j_mF\right\|^2_{L^2(m)}=\int\mathrm{d}m(\omega)\,\left|m^j_{\omega}(F)\right|^2=\int_{\mathcal M}\mathrm{d}\nu^j(\rho)\,\left|\rho(F)\right|^2.
\]
In the same way, with $J=\{1,2,3\}$, one has $\left\|P^{123}_mF\right\|^2_{L^2(m)}=\int_{\mathcal M}\mathrm{d}\nu(\rho)\,\left|\rho(F)\right|^2$. As $\nu^j=\nu$, the two norms coincide. The inclusion $\mathcal E^{123}(m)\subset\mathcal E^j(m)$ gives $P^j_mP^{123}_m=P^{123}_m=P^{123}_mP^j_m$; so $P^j_m-P^{123}_m$ is an orthogonal projection, and $P^j_mF$ is the sum of the two orthogonal vectors $P^{123}_mF$ and $\left(P^j_m-P^{123}_m\right)F$. Hence
\[
\left\|\left(P^j_m-P^{123}_m\right)F\right\|^2_{L^2(m)}=\left\|P^j_mF\right\|^2_{L^2(m)}-\left\|P^{123}_mF\right\|^2_{L^2(m)}=0.
\]
The bounded operator $P^j_m-P^{123}_m$ thus vanishes on the dense set $\mathcal A^0$, hence on $L^2(m)$; this is~(2).

So the conditions (1)--(4) are equivalent, and the preceding paragraph shows that every probability measure $\nu'$ as in~(4) satisfies $\nu'=\nu=\nu^j$.

\emph{Clustering when the conditions hold for $j=1,2,3$.} Assume that (1)--(4) hold with $j$ replaced by $1$, by $2$ and by $3$. By~(1) for these three indices and by the beginning of the proof, there is a $\nu$-null set outside of which $\rho\in\mathcal C\cap\mathcal M^{123}_{\mathrm{erg}}$ and $\ker H_l(\rho)=\mathbb C\cdot\Omega_l$ for $l=1,2,3$. Let $\rho$ be such a law. By the implication (1)$\Rightarrow$(2) of Theorem~\ref{thm:invariant-vectors-clustering}(ii), $\rho$ is $\Gamma_l$-ergodic, that is, $\mathcal E^l(\rho)=\mathbb C\cdot1$, and
\[
P^l_{\rho}G=\langle1,G\rangle_{\rho}\cdot1=\rho(G)\cdot1\qquad\left(G\in L^2(\rho),\ l=1,2,3\right).
\]
Let $F,G\in L^2(\rho)$, and let $\left(a^{(k)}\right)_{k\ge1}$ be a sequence in $\mathbb R^3$ with $\left|a^{(k)}\right|\to\infty$. Every subsequence has a further subsequence along which, for each $l\in\{1,2,3\}$, either $\left|a^{(k)}_l\right|\to\infty$ or $a^{(k)}_l$ converges to a number $\bar a_l\in\mathbb R$; the set $K$ of the indices $l$ of the first kind is not empty, since $\left|a^{(k)}\right|\to\infty$. Along such a subsequence, Theorem~\ref{thm:invariant-vectors-clustering}(i) gives
\[
\left\langle F,U\left(a^{(k)}\right)G\right\rangle_{\rho}\to\left\langle F,P^K_{\rho}U(\bar a)G\right\rangle_{\rho}=\overline{\rho(F)}\cdot\rho\left(U(\bar a)G\right)=\overline{\rho(F)}\cdot\rho(G);
\]
indeed, $P^K_{\rho}=\prod_{l\in K}P^l_{\rho}$ is the orthogonal projection onto $\mathbb C\cdot1$, the set $K$ being non-empty, and $\rho\left(U(\bar a)G\right)=\rho(G)$ by $(T_{123})$. Thus every subsequence of $\left(\rho\left(\overline{F}\cdot U\left(a^{(k)}\right)G\right)\right)_{k\ge1}$ has a further subsequence which tends to $\overline{\rho(F)}\cdot\rho(G)$, and the whole sequence tends to this number. The sequence $\left(a^{(k)}\right)_{k\ge1}$ being arbitrary, $\rho\left(\overline{F}\cdot U(a)G\right)\to\overline{\rho(F)}\cdot\rho(G)$ as $|a|\to\infty$ in any manner.
\end{proof}

\begin{proof}[Proof of Theorem~\ref{thm:vacuum-symmetries-components}(ii)]
Let $g\in O(3)$ satisfy $g_{*}m=m$; as a Euclidean map, $g$ is $x\mapsto gx$. Since $g\Gamma_{123}=\mathbb R^3=\Gamma_{123}$ and $m$ satisfies $(T_{123})$, Lemma~\ref{lem:euclidean-covariance-kernel}, applied to $m$, $J=\{1,2,3\}$ and $g$, gives
\[
g_{*}m_{\omega}=m_{g\omega}\qquad\text{for $m$-almost every }\omega\in Q.
\]

Let $Q'=\left\{\omega\in Q:m_{\omega}\in\mathcal C\right\}$. The set $\mathcal C$ belongs to $\mathcal B_{\mathcal M}$ by Lemma~\ref{lem:measurability-defining-properties}, and $\omega\mapsto m_{\omega}$ is a measurable map from $(Q,\mathcal I^{123})$ to $(\mathcal M,\mathcal B_{\mathcal M})$ by Theorem~\ref{thm:ergodic-decomposition}(a); hence $Q'\in\mathcal I^{123}\subset\Sigma$. By Theorem~\ref{thm:ergodic-decomposition}(a), (b), $\nu$-almost every $\rho$ belongs to $\mathcal C$, as shown at the beginning of the proof of Theorem~\ref{thm:vacuum-symmetries-components}(i); and $\nu$ is the image of $m$ under $\omega\mapsto m_{\omega}$. Therefore $m(Q')=\nu(\mathcal C)=1$. From $g_{*}m=m$ one obtains $(g^{-1})_{*}m=m$, and hence $m(gA)=\left((g^{-1})_{*}m\right)(A)=m(A)$ for $A\in\Sigma$, where $gA\in\Sigma$ (Section~\ref{sec:euclidean-maps-actions}). Let $Z\in\Sigma$ be an $m$-null set such that $g_{*}m_{\omega'}=m_{g\omega'}$ for every $\omega'\in Q\setminus Z$. Then $m\left(gQ'\setminus gZ\right)=m\left(Q'\setminus Z\right)=1$, and for $\omega\in gQ'\setminus gZ$ the point $\omega'=g^{-1}\omega$ lies in $Q'\setminus Z$, so that
\[
m_{\omega}=m_{g\omega'}=g_{*}m_{\omega'}\qquad\text{with }m_{\omega'}\in\mathcal C.
\]
Thus, for $m$-almost every $\omega$, the law $m_{\omega}$ belongs to the set $g_{*}\mathcal C=\left\{g_{*}\rho':\rho'\in\mathcal C\right\}$. This set is the preimage of $\mathcal C$ under the map $\rho\mapsto(g^{-1})_{*}\rho$, which is $\mathcal B_{\mathcal M}$-measurable because $\left((g^{-1})_{*}\rho\right)(A)=\rho(gA)$ for $A\in\Sigma$; so $g_{*}\mathcal C\in\mathcal B_{\mathcal M}$, and
\[
\nu\left(g_{*}\mathcal C\right)=m\left(\left\{\omega\in Q:m_{\omega}\in g_{*}\mathcal C\right\}\right)=1.
\]

It remains to show that every law in $g_{*}\mathcal C$ has the asserted properties. Let first $h\in O(3)$, let $n'\in\mathbb R^3$ be a unit vector, let $c\in\mathbb R$, and put $n=hn'$. As $h$ is orthogonal, $n'\cdot h^{-1}x=n\cdot x$ for $x\in\mathbb R^3$. Hence $h\{n'\cdot x>c\}=\{n\cdot x>c\}$, $h\{n'\cdot x=c\}=\{n\cdot x=c\}$, and
\[
h\circ r_{n',c}\circ h^{-1}=r_{n,c},
\]
since $h\left(h^{-1}x-2\left(n'\cdot h^{-1}x-c\right)\cdot n'\right)=x-2\left(n\cdot x-c\right)\cdot n$ for $x\in\mathbb R^3$. Let $\rho\in\mathcal M$. For every bounded measurable function $\Phi$ on $Q$ one has $(h_{*}\rho)(\Phi)=\rho(\Phi\circ h)$, where $(\Phi\circ h)(\omega)=\Phi(h\omega)$; and $X(f)\circ h=X\left(h^{-1}\cdot f\right)$ for $f\in\mathcal S$, because $(h\omega)(f)=\omega\left(h^{-1}\cdot f\right)$ (Section~\ref{sec:euclidean-maps-actions}). Let $k\ge1$, $z_1,\dots,z_k\in\mathbb C$ and $f_1,\dots,f_k\in\mathcal D(\{n\cdot x>c\})$. Then $h^{-1}\cdot f_p\in\mathcal D(\{n'\cdot x>c\})$ (Section~\ref{sec:euclidean-maps-actions}), and $h^{-1}\cdot\left(r_{n,c}\cdot f_p\right)=r_{n',c}\cdot\left(h^{-1}\cdot f_p\right)$ by the last display. Therefore
\begin{multline*}
\sum_{p,q=1}^k\overline{z_p}\cdot z_q\cdot\left(h_{*}\rho\right)\left(\exp\left(i\left(X(f_q)-X(r_{n,c}\cdot f_p)\right)\right)\right)\\
=\sum_{p,q=1}^k\overline{z_p}\cdot z_q\cdot\rho\left(\exp\left(i\left(X\left(h^{-1}\cdot f_q\right)-X\left(r_{n',c}\cdot\left(h^{-1}\cdot f_p\right)\right)\right)\right)\right).
\end{multline*}
Consequently, if $\rho$ is reflection positive at the plane $\{n'\cdot x=c\}$, then $h_{*}\rho$ is reflection positive at the plane $\{n\cdot x=c\}$, in the sense of Section~\ref{sec:properties-laws}. For $n'=e_l$, $l\in\{1,2,3\}$, one has $r_{e_l,c}=\theta_{l,c}$ and $\{n'\cdot x=c\}=\{x_l=c\}$; then $n=he_l$, the plane $\{n\cdot x=c\}$ is $h\{x_l=c\}$, and $r_{n,c}=h\circ\theta_{l,c}\circ h^{-1}$.

Now let $\rho'\in\mathcal C$ and $l\in\{1,2,3\}$. As $(\theta_l)_{*}\rho'=\rho'$,
\[
\left(g\circ\theta_l\circ g^{-1}\right)_{*}\left(g_{*}\rho'\right)=g_{*}\left((\theta_l)_{*}\rho'\right)=g_{*}\rho',
\]
so $g_{*}\rho'$ is invariant under $g\circ\theta_l\circ g^{-1}$. The law $\rho'$ satisfies $(T_l)$, $(R_l)$ and $(RP_l)$; by Lemma~\ref{lem:reformulations-reflection-positivity}(5) it is reflection positive at every plane $\{x_l=c\}$, $c\in\mathbb R$. The computation above, with $h=g$ and $n'=e_l$, shows that $g_{*}\rho'$ is reflection positive at every plane $g\{x_l=c\}=\{(ge_l)\cdot x=c\}$, $c\in\mathbb R$, the reflection in which is $g\circ\theta_{l,c}\circ g^{-1}$. Since $\nu$-almost every $\rho$ belongs to $g_{*}\mathcal C$, it follows that $\nu$-almost every $\rho$ is invariant under $g\circ\theta_l\circ g^{-1}$ and reflection positive at the planes $g\{x_l=c\}$, for $l=1,2,3$ and $c\in\mathbb R$.
\end{proof}

\begin{proof}[Proof of Theorem~\ref{thm:vacuum-symmetries-components}(iii)]
Let $g\in O(3)$ satisfy $g_{*}m=m$, and let $n=ge_j$ have three non-zero coordinates. Let $(s_k)_{k\ge1}$ be a real sequence with $s_k\to\infty$. Then $\left|s_k\cdot n_l\right|=s_k\cdot|n_l|\to\infty$ for $l=1,2,3$, and Theorem~\ref{thm:invariant-vectors-clustering}(i), applied to $\rho=m$, $K=\{1,2,3\}$ and $a^{(k)}=s_k\cdot n$, gives $\left\langle F,U(s_k\cdot n)G\right\rangle_m\to\left\langle F,P^{123}_mG\right\rangle_m$ for all $F,G\in L^2(m)$. Hence $U(s\cdot n)$ converges weakly to $P^{123}_m$ as $s\to\infty$.

As $g_{*}m=m$, the operator $\hat g$ is unitary on $L^2(m)$, with inverse $\widehat{g^{-1}}$ (Section~\ref{sec:euclidean-maps-actions}), and $g\circ\tau_{s\cdot e_j}\circ g^{-1}=\tau_{s\cdot n}$ gives $\hat gU(s\cdot e_j)\hat g^{-1}=U(s\cdot n)$ for $s\in\mathbb R$. Lemma~\ref{lem:euclidean-covariance-kernel}, applied to $m$, $J=\{1,2,3\}$ and $g$ (here $g\Gamma_{123}=\mathbb R^3=\Gamma_{123}$), gives $\hat gP^{123}_m=P^{123}_m\hat g$. Hence, for $F,G\in L^2(m)$, as $s\to\infty$,
\[
\left\langle F,U(s\cdot e_j)G\right\rangle_m=\left\langle\hat gF,U(s\cdot n)\hat gG\right\rangle_m\to\left\langle\hat gF,P^{123}_m\hat gG\right\rangle_m=\left\langle F,\hat g^{-1}P^{123}_m\hat gG\right\rangle_m=\left\langle F,P^{123}_mG\right\rangle_m.
\]
On the other hand, $m$ satisfies $(T_j)$, $(R_j)$ and $(RP_j)$, and by Proposition~\ref{prop:uniform-mixing-along} the operators $U(s\cdot e_j)$ converge weakly to $P^j_m$ as $|s|\to\infty$. A weak limit being unique, $P^j_m=P^{123}_m$. This is condition~(2) of Theorem~\ref{thm:vacuum-symmetries-components}(i); by Theorem~\ref{thm:vacuum-symmetries-components}(i), the conditions (1)--(4) of that clause hold for this $j$.
\end{proof}

\begin{proof}[Proof of Theorem~\ref{thm:vacuum-symmetries-components}(iv)]
Assume that $g_{*}m=m$ for every $g\in O(3)$. Let $\mathbb S^2=\left\{n\in\mathbb R^3:|n|=1\right\}$, and let $\{n_k:k\ge1\}$ be a countable dense subset of $\mathbb S^2$. For each $k$ choose $g_k\in O(3)$ with $g_ke_1=n_k$. By the computation in the proof of Theorem~\ref{thm:vacuum-symmetries-components}(ii), with $h=g_k$, $n'=e_1$ and $c=0$, the map $g_k\circ\theta_1\circ g_k^{-1}$ is the reflection $r_{n_k}$ in the plane $\{n_k\cdot x=0\}$. By Theorem~\ref{thm:vacuum-symmetries-components}(ii), applied to $g_k$, $\nu$-almost every $\rho$ is invariant under $r_{n_k}$; as there are countably many $k$, $\nu$-almost every $\rho$ is invariant under $r_{n_k}$ for every $k\ge1$. By Theorem~\ref{thm:ergodic-decomposition}(a), (b), $\nu$-almost every $\rho$ belongs to $\mathcal C\cap\mathcal M^{123}_{\mathrm{erg}}$, as shown at the beginning of the proof of Theorem~\ref{thm:vacuum-symmetries-components}(i).

Let $\rho\in\mathcal M$ be invariant under every $r_{n_k}$, $k\ge1$, and let
\[
O(3)_{\rho}=\left\{g\in O(3):g_{*}\rho=\rho\right\}.
\]
Since $(g\circ g')_{*}\rho=g_{*}(g'_{*}\rho)$ for $g,g'\in O(3)$, and $(g^{-1})_{*}\rho=\rho$ whenever $g_{*}\rho=\rho$, the set $O(3)_{\rho}$ is a subgroup of $O(3)$. It is closed. Indeed, for $f\in\mathcal S$ the map $g\mapsto g\cdot f$ from $O(3)$ to $\mathcal S$ is continuous, and every $\omega\in Q$ is continuous on $\mathcal S$; so $g\mapsto\rho\left(\exp\left(iX(g\cdot f)\right)\right)$ is continuous on $O(3)$, by dominated convergence. As $(g_{*}\rho)\left(\exp\left(iX(f)\right)\right)=\rho\left(\exp\left(iX\left(g^{-1}\cdot f\right)\right)\right)$ (proof of Theorem~\ref{thm:vacuum-symmetries-components}(ii)), Lemma~\ref{lem:density-exponentials}(ii) gives
\[
O(3)_{\rho}=\bigcap_{f\in\mathcal D_0}\left\{g\in O(3):\rho\left(\exp\left(iX\left(g^{-1}\cdot f\right)\right)\right)=\rho\left(\exp\left(iX(f)\right)\right)\right\},
\]
an intersection of closed subsets of $O(3)$. The map $n\mapsto r_n$ from $\mathbb S^2$ to $O(3)$, $r_nx=x-2(n\cdot x)\cdot n$, is continuous; so the set of the $n\in\mathbb S^2$ with $r_n\in O(3)_{\rho}$ is closed in $\mathbb S^2$. It contains the dense set $\{n_k:k\ge1\}$, hence equals $\mathbb S^2$: the group $O(3)_{\rho}$ contains every reflection $r_n$, $n\in\mathbb S^2$. By the theorem of Cartan and Dieudonn\'e, every element of $O(3)$ is a product of such reflections. Hence $O(3)_{\rho}=O(3)$. If $\rho$ satisfies in addition $(T_{123})$, then $\rho$ is invariant under every Euclidean map $x\mapsto gx+a$, which is the composition of $g$ and $\tau_a$, with $g\in O(3)$ and $a\in\mathbb R^3$; that is, $\rho$ is invariant under $\mathbb R^3\rtimes O(3)$. Consequently $\nu$-almost every $\rho$ is invariant under $\mathbb R^3\rtimes O(3)$.

Now let $\rho\in\mathcal C\cap\mathcal M^{123}_{\mathrm{erg}}$ be invariant under $O(3)$; $\nu$-almost every $\rho$ is such a law. Let $F,G\in L^2(\rho)$, and let $\left(a^{(k)}\right)_{k\ge1}$ be a sequence in $\mathbb R^3$ with $\left|a^{(k)}\right|\to\infty$. Consider a subsequence. Omitting finitely many terms, one may assume $a^{(k)}\ne0$ and put $u^{(k)}=\left|a^{(k)}\right|^{-1}\cdot a^{(k)}\in\mathbb S^2$; since $\mathbb S^2$ is compact, there is a further subsequence along which $u^{(k)}$ converges to some $n\in\mathbb S^2$. Choose $g\in O(3)$ with $gn=3^{-\frac{1}{2}}\cdot(1,1,1)$. Along this further subsequence $gu^{(k)}\to3^{-\frac{1}{2}}\cdot(1,1,1)$, so that every coordinate of $gu^{(k)}$ is at least $\frac{1}{2}\cdot3^{-\frac{1}{2}}$ for large $k$, and
\[
\min_l\left|\left(ga^{(k)}\right)_l\right|\ge\frac{1}{2}\cdot3^{-\frac{1}{2}}\cdot\left|a^{(k)}\right|\to\infty.
\]
As $g_{*}\rho=\rho$, the operator $\hat g$ is unitary on $L^2(\rho)$, $\rho(\hat g\Phi)=\rho(\Phi)$ for $\Phi\in L^1(\rho)$, and $\hat gU(a)\hat g^{-1}=U(ga)$ for $a\in\mathbb R^3$, because $g\circ\tau_a\circ g^{-1}=\tau_{ga}$. Theorem~\ref{thm:invariant-vectors-clustering}(iii)1, applied to $\hat gF$, $\hat gG$ and the sequence $\left(ga^{(k)}\right)$, gives along the further subsequence
\[
\rho\left(\overline{F}\cdot U\left(a^{(k)}\right)G\right)=\left\langle\hat gF,U\left(ga^{(k)}\right)\hat gG\right\rangle_{\rho}\to\overline{\rho(\hat gF)}\cdot\rho(\hat gG)=\overline{\rho(F)}\cdot\rho(G).
\]
Every subsequence of $\left(\rho\left(\overline{F}\cdot U\left(a^{(k)}\right)G\right)\right)_{k\ge1}$ thus has a further subsequence which tends to $\overline{\rho(F)}\cdot\rho(G)$; hence the whole sequence tends to this number, and $\rho\left(\overline{F}\cdot U(a)G\right)\to\overline{\rho(F)}\cdot\rho(G)$ as $|a|\to\infty$ in any manner. In particular $\rho\left(\overline{F}\cdot U(s\cdot e_j)G\right)\to\overline{\rho(F)}\cdot\rho(G)$ as $|s|\to\infty$, for all $F,G\in L^2(\rho)$. This is condition~(3) of Theorem~\ref{thm:invariant-vectors-clustering}(ii) for the law $\rho$, which satisfies $(T_j)$, $(R_j)$ and $(RP_j)$; that theorem gives $\ker H_j(\rho)=\mathbb C\cdot\Omega_j$.
\end{proof}

\begin{proof}[Proof of Theorem~\ref{thm:vacuum-symmetries-components}(v)]
Let $g\in O(3)$. For $a\in\mathbb R^3$ one has $\tau_a\circ g=g\circ\tau_{g^{-1}a}$; hence, if $\rho\in\mathcal M$ satisfies $(T_{123})$, then $(\tau_a)_{*}(g_{*}\rho)=g_{*}\left((\tau_{g^{-1}a})_{*}\rho\right)=g_{*}\rho$ for every $a\in\mathbb R^3$, and $g_{*}\rho$ satisfies $(T_{123})$. Let $i\in\{1,2,3\}$, and put $m'=(g^{-1})_{*}m$; by what has just been said, applied to $g^{-1}$ and $m$, the law $m'$ satisfies $(T_{123})$. By the computation in the proof of Theorem~\ref{thm:vacuum-symmetries-components}(ii), with $h=g$, $n'=e_i$ and $c=0$, one has $g\circ\theta_i\circ g^{-1}=r_{ge_i}$, that is, $\theta_i\circ g^{-1}=g^{-1}\circ r_{ge_i}$. If $(r_{ge_i})_{*}m=m$, it follows that
\[
(\theta_i)_{*}m'=\left(g^{-1}\right)_{*}\left(r_{ge_i}\right)_{*}m=m',
\]
so that $m'$ satisfies $(R_i)$. If moreover $m$ is reflection positive at the plane $\{(ge_i)\cdot x=0\}$, then the same computation, with $h=g^{-1}$, $n'=ge_i$ and $c=0$, so that $hn'=e_i$, shows that $m'=h_{*}m$ is reflection positive at the plane $\{e_i\cdot x=0\}=\{x_i=0\}$; that is, $m'$ satisfies $(RP_i)$ (Section~\ref{sec:properties-laws}).

Now let $n$ be a unit vector such that $(r_n)_{*}m=m$ and $m$ is reflection positive at the plane $\{n\cdot x=0\}$. Choose $g\in O(3)$ with $ge_1=n$, and put $m'=(g^{-1})_{*}m$. By the preceding paragraph, with $i=1$, the law $m'$ satisfies $(T_{123})$, $(R_1)$ and $(RP_1)$. Let $\nu'$ be the probability measure on $(\mathcal M,\mathcal B_{\mathcal M})$ given by Theorem~\ref{thm:ergodic-decomposition} for the law $m'$ and $J=\{1,2,3\}$.

The map $\rho'\mapsto g_{*}\rho'$ of $\mathcal M$ into itself is a bijection, with inverse $\rho\mapsto(g^{-1})_{*}\rho$, and both maps are $\mathcal B_{\mathcal M}$-measurable, since $(g_{*}\rho')(A)=\rho'(g^{-1}A)$ with $g^{-1}A\in\Sigma$ for $A\in\Sigma$ (Section~\ref{sec:euclidean-maps-actions}). It maps $\mathcal M^{123}_{\mathrm{erg}}$ onto itself. Indeed, let $\rho'\in\mathcal M^{123}_{\mathrm{erg}}$; by the first paragraph of this proof, $g_{*}\rho'$ satisfies $(T_{123})$. As in the proof of Theorem~\ref{thm:cubic-symmetry-components}(ii), the map $F\mapsto F\circ g$ is a unitary operator of $L^2(g_{*}\rho')$ onto $L^2(\rho')$, with inverse $F\mapsto F\circ g^{-1}$, and
\[
\left(U(a)F\right)\circ g=U\left(g^{-1}a\right)(F\circ g)\qquad\left(a\in\mathbb R^3\right),
\]
because $g^{-1}\circ\tau_a\circ g=\tau_{g^{-1}a}$. So $F\in L^2(g_{*}\rho')$ is invariant under all $U(a)$, $a\in\mathbb R^3$, if and only if $F\circ g$ is, and $\mathcal E^{123}(g_{*}\rho')=\mathbb C\cdot1$; thus $g_{*}\rho'\in\mathcal M^{123}_{\mathrm{erg}}$. The same holds with $g^{-1}$ in place of $g$, and the map is onto $\mathcal M^{123}_{\mathrm{erg}}$.

Let $\tilde\nu$ be the image of $\nu'$ under $\rho'\mapsto g_{*}\rho'$. Then $\tilde\nu\left(\mathcal M^{123}_{\mathrm{erg}}\right)=\nu'\left(\mathcal M^{123}_{\mathrm{erg}}\right)=1$ by Theorem~\ref{thm:ergodic-decomposition}(a); and, since $g_{*}m'=m$, for every $A\in\Sigma$
\[
m(A)=m'\left(g^{-1}A\right)=\int_{\mathcal M}\mathrm{d}\nu'(\rho')\,\rho'\left(g^{-1}A\right)=\int_{\mathcal M}\mathrm{d}\nu'(\rho')\,(g_{*}\rho')(A)=\int_{\mathcal M}\mathrm{d}\tilde\nu(\rho)\,\rho(A).
\]
So $m=\int_{\mathcal M}\mathrm{d}\tilde\nu(\rho)\,\rho$, and the uniqueness assertion of Theorem~\ref{thm:ergodic-decomposition}(a), for $m$ and $J=\{1,2,3\}$, gives $\tilde\nu=\nu$: the measure $\nu$ is the image of $\nu'$ under $\rho'\mapsto g_{*}\rho'$.

By Theorem~\ref{thm:ergodic-decomposition}(b), applied to $m'$ with $J=\{1,2,3\}$ and $j=1$, there is a set $B'\in\mathcal B_{\mathcal M}$ with $\nu'(B')=1$ each element $\rho'$ of which satisfies $(R_1)$ and is reflection positive at every plane $\{x_1=c\}$, $c\in\mathbb R$. Let $\rho'\in B'$. As $r_n=g\circ\theta_1\circ g^{-1}$, one has $(r_n)_{*}(g_{*}\rho')=g_{*}\left((\theta_1)_{*}\rho'\right)=g_{*}\rho'$; and the computation in the proof of Theorem~\ref{thm:vacuum-symmetries-components}(ii), with $h=g$ and $n'=e_1$, shows that $g_{*}\rho'$ is reflection positive at every plane $g\{x_1=c\}=\{n\cdot x=c\}$, $c\in\mathbb R$. The set $g_{*}B'=\left\{g_{*}\rho':\rho'\in B'\right\}$ is the preimage of $B'$ under the measurable map $\rho\mapsto(g^{-1})_{*}\rho$; it belongs to $\mathcal B_{\mathcal M}$, and $\nu(g_{*}B')=\nu'(B')=1$. Hence $\nu$-almost every $\rho$ is invariant under $r_n$ and reflection positive at every plane $\{n\cdot x=c\}$, $c\in\mathbb R$.

Finally, for every $\Phi\in L^1(m')$ one has $m(\hat g\Phi)=m\left(\Phi\circ g^{-1}\right)=\left((g^{-1})_{*}m\right)(\Phi)=m'(\Phi)$; hence $\hat g$ is an isometry of $L^2(m')$ into $L^2(m)$, and it is onto, with inverse $\widehat{g^{-1}}$. Since $g\circ\tau_{s\cdot e_1}\circ g^{-1}=\tau_{s\cdot n}$, one has $\hat gU(s\cdot e_1)\hat g^{-1}=U(s\cdot n)$ for $s\in\mathbb R$. The law $m'$ satisfies $(T_1)$, $(R_1)$ and $(RP_1)$, and Proposition~\ref{prop:uniform-mixing-along}, applied to $m'$ and the direction $e_1$, shows that $U(s\cdot e_1)$ converges weakly on $L^2(m')$ to $P^1_{m'}$ as $|s|\to\infty$. For $F,G\in L^2(m)$ it follows that, as $|s|\to\infty$,
\[
\left\langle F,U(s\cdot n)G\right\rangle_m=\left\langle\hat g^{-1}F,U(s\cdot e_1)\hat g^{-1}G\right\rangle_{m'}\to\left\langle\hat g^{-1}F,P^1_{m'}\hat g^{-1}G\right\rangle_{m'}=\left\langle F,\hat gP^1_{m'}\hat g^{-1}G\right\rangle_m.
\]
The operator $\hat gP^1_{m'}\hat g^{-1}$ is the orthogonal projection of $L^2(m)$ onto $\hat g\,\mathcal E^1(m')$; and $F\in\hat g\,\mathcal E^1(m')$ if and only if $U(s\cdot e_1)\hat g^{-1}F=\hat g^{-1}F$ for all $s\in\mathbb R$, that is, if and only if $U(s\cdot n)F=F$ for all $s\in\mathbb R$. Thus $U(s\cdot n)$ converges weakly on $L^2(m)$, as $|s|\to\infty$, to the orthogonal projection onto $\left\{F\in L^2(m):U(s\cdot n)F=F\text{ for every }s\in\mathbb R\right\}$.
\end{proof}

\begin{proof}[Proof of Theorem~\ref{thm:vacuum-symmetries-components}(vi)]
Assume that for each of the six vectors $n=\frac{1}{\sqrt{2}}\left(e_i\pm e_l\right)$, $1\le i<l\le3$, one has $(r_n)_{*}m=m$ and $m$ is reflection positive at the plane $\{n\cdot x=0\}$.

\emph{Invariance of the components.} Let $1\le i<l\le3$ and $n=\frac{1}{\sqrt{2}}\left(e_i-e_l\right)$. Then $n\cdot x=\frac{1}{\sqrt{2}}\left(x_i-x_l\right)$ and $r_n(x)=x-(x_i-x_l)\cdot(e_i-e_l)$ for $x\in\mathbb R^3$: the reflection $r_n$ exchanges the coordinates $x_i$ and $x_l$ and leaves the third coordinate unchanged. By Theorem~\ref{thm:vacuum-symmetries-components}(v), $\nu$-almost every $\rho$ is invariant under these three transpositions of coordinates. By Theorem~\ref{thm:ergodic-decomposition}(a), (b), $\nu$-almost every $\rho$ satisfies $(T_{123})$ and $(R_1)$, $(R_2)$, $(R_3)$. An element $w\in O_{\mathrm{cub}}$ with $w\,e_i=\varepsilon_i\cdot e_{\pi(i)}$ for $i=1,2,3$ is the product of the diagonal matrix in $N_3$ which multiplies $e_{\pi(i)}$ by $\varepsilon_i$ and of the matrix of the permutation $e_i\mapsto e_{\pi(i)}$, which is a product of transpositions of coordinates; and $N_3$ is generated by $\theta_1,\theta_2,\theta_3$ (Section~\ref{sec:euclidean-maps-actions}). Hence every Euclidean map $x\mapsto wx+a$, $w\in O_{\mathrm{cub}}$, $a\in\mathbb R^3$, is a finite product of translations, of the reflections $\theta_1,\theta_2,\theta_3$ and of the three transpositions. A law which is invariant under each factor of such a product is invariant under the product, because $(g\circ g')_{*}\rho=g_{*}(g'_{*}\rho)$. Therefore $\nu$-almost every $\rho$ is invariant under $\mathbb R^3\rtimes O_{\mathrm{cub}}$.

\emph{The projections.} Let $j,l\in\{1,2,3\}$ be distinct, let $l'$ be the third element of $\{1,2,3\}$, write $\{j,l\}=\{i_1,i_2\}$ with $i_1<i_2$, and put
\[
p=\frac{1}{\sqrt{2}}\left(e_{i_1}+e_{i_2}\right),\qquad q=\frac{1}{\sqrt{2}}\left(e_{i_1}-e_{i_2}\right);
\]
these are two of the six vectors. The vectors $p$, $q$, $e_{l'}$ form an orthonormal basis of $\mathbb R^3$; let $g\in O(3)$ be given by $g\,e_j=p$, $g\,e_l=q$ and $g\,e_{l'}=e_{l'}$, and put $m'=(g^{-1})_{*}m$. The hypothesis of Theorem~\ref{thm:vacuum-symmetries-components}(v) holds for $p=ge_j$ and for $q=ge_l$. Moreover $(r_{e_{l'}})_{*}m=(\theta_{l'})_{*}m=m$, and $m$ is reflection positive at the plane $\{e_{l'}\cdot x=0\}=\{x_{l'}=0\}$, since $m$ satisfies $(R_{l'})$ and $(RP_{l'})$. By the first paragraph of the proof of Theorem~\ref{thm:vacuum-symmetries-components}(v), applied with $i=j$, with $i=l$ and with $i=l'$, the law $m'$ satisfies $(T_{123})$, and $(R_i)$, $(RP_i)$ for $i=1,2,3$; that is, $m'\in\mathcal C$.

As $g$ is orthogonal, $g^{-1}e_j=\sum_{i=1}^3\left(ge_i\cdot e_j\right)\cdot e_i=(p\cdot e_j)\cdot e_j+(q\cdot e_j)\cdot e_l=\frac{1}{\sqrt{2}}\left(e_j\pm e_l\right)$, with the sign $+$ if $j=i_1$ and the sign $-$ if $j=i_2$. Let $(s_k)_{k\ge1}$ be a real sequence with $s_k\to\infty$. The points $a^{(k)}=s_k\cdot g^{-1}e_j$ satisfy $\left|a^{(k)}_j\right|=\left|a^{(k)}_l\right|=\frac{s_k}{\sqrt{2}}\to\infty$ and $a^{(k)}_{l'}=0$. Theorem~\ref{thm:invariant-vectors-clustering}(i), applied to $m'\in\mathcal C$ with $K=\{j,l\}$ and $\bar a=0$, gives
\[
\lim_{s\to\infty}\left\langle F',U\left(s\cdot g^{-1}e_j\right)G'\right\rangle_{m'}=\left\langle F',P^{\{j,l\}}_{m'}G'\right\rangle_{m'}\qquad\left(F',G'\in L^2(m')\right).
\]
As in the proof of Theorem~\ref{thm:vacuum-symmetries-components}(v), $\hat g$ is a unitary operator of $L^2(m')$ onto $L^2(m)$; moreover $\hat gU(b)\hat g^{-1}=U(gb)$ for $b\in\mathbb R^3$, because $g\circ\tau_b\circ g^{-1}=\tau_{gb}$. Hence $\hat gU\left(s\cdot g^{-1}e_j\right)\hat g^{-1}=U(s\cdot e_j)$, and for $F,G\in L^2(m)$, as $s\to\infty$,
\[
\left\langle F,U(s\cdot e_j)G\right\rangle_m=\left\langle\hat g^{-1}F,U\left(s\cdot g^{-1}e_j\right)\hat g^{-1}G\right\rangle_{m'}\to\left\langle F,\hat gP^{\{j,l\}}_{m'}\hat g^{-1}G\right\rangle_m.
\]
The operator $\hat gP^{\{j,l\}}_{m'}\hat g^{-1}$ is the orthogonal projection of $L^2(m)$ onto $\hat g\,\mathcal E^{\{j,l\}}(m')$. A vector $F\in L^2(m)$ belongs to this space if and only if $\hat g^{-1}F$ is invariant under all $U(t\cdot e_j)$ and $U(t\cdot e_l)$, $t\in\mathbb R$, that is, if and only if $F$ is invariant under all $U(t\cdot p)$ and $U(t\cdot q)$, $t\in\mathbb R$. The vectors $p$ and $q$ span the plane $\Gamma_{\{j,l\}}$ spanned by $e_j$ and $e_l$, and $U(t\cdot p+t'\cdot q)=U(t\cdot p)U(t'\cdot q)$; so this is invariance under all $U(a)$, $a\in\Gamma_{\{j,l\}}$, and $\hat g\,\mathcal E^{\{j,l\}}(m')=\mathcal E^{\{j,l\}}(m)$. Thus $U(s\cdot e_j)$ converges weakly to $P^{\{j,l\}}_m$ as $s\to\infty$. On the other hand, $m$ satisfies $(T_j)$, $(R_j)$ and $(RP_j)$, and by Proposition~\ref{prop:uniform-mixing-along} the operators $U(s\cdot e_j)$ converge weakly to $P^j_m$ as $|s|\to\infty$. Hence
\[
P^j_m=P^{\{j,l\}}_m\qquad\text{for all distinct }j,l\in\{1,2,3\}.
\]
Let $j\in\{1,2,3\}$, and let $l$ and $l'$ be the two other indices. By Lemma~\ref{lem:averages-translations}(v), the projections $P^1_m$, $P^2_m$, $P^3_m$ commute, and $P^{\{j,l\}}_m=P^j_mP^l_m$, $P^{\{j,l'\}}_m=P^j_mP^{l'}_m$, $P^{123}_m=P^1_mP^2_mP^3_m$. Therefore
\[
P^j_m=P^j_mP^j_m=P^{\{j,l\}}_mP^{\{j,l'\}}_m=P^j_mP^l_mP^j_mP^{l'}_m=P^j_mP^l_mP^{l'}_m=P^{123}_m.
\]
So condition~(2) of Theorem~\ref{thm:vacuum-symmetries-components}(i) holds with $j$ replaced by $1$, by $2$ and by $3$. By Theorem~\ref{thm:vacuum-symmetries-components}(i), condition~(1) of that clause holds for these three indices, that is, $\ker H_j(\rho)=\mathbb C\cdot\Omega_j$ for $j=1,2,3$ and $\nu$-almost every $\rho$; and by the last assertion of Theorem~\ref{thm:vacuum-symmetries-components}(i), for $\nu$-almost every $\rho$ and all $F,G\in L^2(\rho)$, $\rho\left(\overline{F}\cdot U(a)G\right)\to\overline{\rho(F)}\cdot\rho(G)$ as $|a|\to\infty$ in any manner.

\emph{A single plane suffices for a law invariant under the cubic group.} Let $m$ be invariant under $O_{\mathrm{cub}}$ and reflection positive at the plane $\{x_1=x_2\}=\{n_0\cdot x=0\}$, where $n_0=\frac{1}{\sqrt{2}}\left(e_1-e_2\right)$. The reflection $r_{n_0}$ exchanges $x_1$ and $x_2$; it belongs to $O_{\mathrm{cub}}$, so that $(r_{n_0})_{*}m=m$, and by Section~\ref{sec:properties-laws} reflection positivity of $m$ at this plane for the half space $\{n_0\cdot x>0\}=\{x_1>x_2\}$ is equivalent to reflection positivity for the half space $\{x_1<x_2\}$. Let $1\le i<l\le3$, let $\varepsilon,\varepsilon'\in\{1,-1\}$, put $n=\frac{1}{\sqrt{2}}\left(\varepsilon\cdot e_i+\varepsilon'\cdot e_l\right)$, and let $l'$ be the third element of $\{1,2,3\}$. The matrix $w$ with $w\,e_1=\varepsilon\cdot e_i$, $w\,e_2=-\varepsilon'\cdot e_l$ and $w\,e_3=e_{l'}$ belongs to $O_{\mathrm{cub}}$, and
\[
wn_0=\frac{1}{\sqrt{2}}\left(we_1-we_2\right)=n;
\]
so $O_{\mathrm{cub}}$ acts transitively on the twelve vectors $\frac{1}{\sqrt{2}}\left(\pm e_i\pm e_l\right)$, $1\le i<l\le3$. By the computation in the proof of Theorem~\ref{thm:vacuum-symmetries-components}(ii), with $h=w$, $n'=n_0$ and $c=0$, one has $r_n=w\circ r_{n_0}\circ w^{-1}$, an element of $O_{\mathrm{cub}}$; hence $(r_n)_{*}m=m$. By the same computation, the law $w_{*}m$, which equals $m$, is reflection positive at the plane $w\{n_0\cdot x=0\}=\{n\cdot x=0\}$, with the half space $w\{n_0\cdot x>0\}=\{n\cdot x>0\}$. In particular, the hypothesis of Theorem~\ref{thm:vacuum-symmetries-components}(v) holds for each of the six vectors $\frac{1}{\sqrt{2}}\left(e_i\pm e_l\right)$, $1\le i<l\le3$.
\end{proof}

\subsection{Examples}\label{sec:examples}

\subsubsection{Auxiliary results}\label{sec:auxiliary-results}

\begin{lemma}[Gaussian mixtures]\label{lem:gaussian-mixtures}
Let $(\mathcal U,\mathbb P)$ be a probability space, and let $\mathbb E$ denote integration with respect to $\mathbb P$ in the variable $u\in\mathcal U$. For every $u\in\mathcal U$ let $q_u$ be a symmetric bilinear form on the real vector space $\mathcal S$, and write $q_u(f)=q_u(f,f)$ for $f\in\mathcal S$. Assume that $0\le q_u(f)\le\|f\|^2$ for all $u\in\mathcal U$ and $f\in\mathcal S$, and that the function $u\mapsto q_u(f)$ is measurable for every $f\in\mathcal S$. Then there is exactly one law $m\in\mathcal M$ such that
\[
m\left(\exp\left(iX(f)\right)\right)=\mathbb E\exp\left(-\frac{1}{2}\cdot q_u(f)\right)\qquad\left(f\in\mathcal S\right).
\]
This law satisfies
\[
m\left(\exp X(f)\right)=\mathbb E\exp\left(\frac{1}{2}\cdot q_u(f)\right)\qquad\left(f\in\mathcal S\right),
\]
and therefore the Gaussian bound~\eqref{eq:setup-G}.

Let, moreover, $g$ be a Euclidean map, and let $\mathcal U\to\mathcal U$, $u\mapsto u^g$, be a measurable map which preserves $\mathbb P$ and satisfies $q_u(g\cdot f)=q_{u^g}(f)$ for all $u\in\mathcal U$ and $f\in\mathcal S$. Then $g_*m=m$.
\end{lemma}

\begin{proof}
For $f\in\mathcal S$ put
\[
S(f)=\mathbb E\exp\left(-\frac{1}{2}\cdot q_u(f)\right).
\]
The integrand is a measurable function of $u$ with values in $(0,1]$, and $S(0)=1$.

\emph{Positive definiteness.} The function $S$ is positive definite, that is,
\[
\sum_{p,q=1}^n\overline{c_p}\cdot c_q\cdot S(f_q-f_p)\ge0\qquad\left(n\ge1,\ c\in\mathbb C^n,\ f_1,\dots,f_n\in\mathcal S\right).
\]
Indeed, fix $u\in\mathcal U$. The real symmetric matrix $C=\left(q_u(f_p,f_q)\right)_{p,q=1}^n$ is positive semidefinite, because $\sum_{p,q=1}^nt_p\cdot t_q\cdot q_u(f_p,f_q)=q_u\left(\sum_{p=1}^nt_p\cdot f_p\right)\ge0$ for $t\in\mathbb R^n$. Let $N(0,C)$ be the centred Gaussian law on $\mathbb R^n$ with covariance matrix $C$. With $t\cdot y=\sum_{p=1}^nt_p\cdot y_p$, its characteristic function is
\[
\int N(0,C)(\mathrm{d}y)\,\exp\left(i\,t\cdot y\right)=\exp\left(-\frac{1}{2}\cdot q_u\left(\sum_{p=1}^nt_p\cdot f_p\right)\right)\qquad\left(t\in\mathbb R^n\right).
\]
For the vectors $t$ with $t\cdot y=y_q-y_p$ this gives $\exp\left(-\frac{1}{2}\cdot q_u(f_q-f_p)\right)=\int N(0,C)(\mathrm{d}y)\,\exp\left(i\left(y_q-y_p\right)\right)$, and therefore
\[
\sum_{p,q=1}^n\overline{c_p}\cdot c_q\cdot\exp\left(-\frac{1}{2}\cdot q_u(f_q-f_p)\right)=\int N(0,C)(\mathrm{d}y)\,\left|\sum_{q=1}^nc_q\cdot\exp\left(iy_q\right)\right|^2\ge0.
\]
Integration with respect to $\mathbb P$ gives the assertion.

\emph{Continuity.} Let $f,h\in\mathcal S$ and $u\in\mathcal U$. Since $q_u$ is bilinear and symmetric, $q_u(f)-q_u(h)=q_u(f-h,f+h)$. The inequality $\left|\exp(-a)-\exp(-b)\right|\le|a-b|$ for $a,b\ge0$, the Cauchy--Schwarz inequality for the positive semidefinite form $q_u$, and the bound $q_u(k)\le\|k\|^2$ for $k\in\mathcal S$ give
\[
\left|\exp\left(-\frac{1}{2}\cdot q_u(f)\right)-\exp\left(-\frac{1}{2}\cdot q_u(h)\right)\right|\le\frac{1}{2}\cdot\left|q_u(f-h,f+h)\right|\le\frac{1}{2}\cdot\|f-h\|\cdot\|f+h\|.
\]
Hence $\left|S(f)-S(h)\right|\le\frac{1}{2}\cdot\|f-h\|\cdot\|f+h\|$. As the norm $\|\cdot\|$ is continuous on $\mathcal S$, this bound shows that $S$ is continuous on $\mathcal S$.

\emph{Existence and uniqueness.} The space $\mathcal S$ is nuclear. By the theorem of Minlos \cite[Section~3.4]{GlimmJaffe}, a continuous positive definite function on $\mathcal S$ which takes the value $1$ at $0$ is the characteristic functional of exactly one probability measure on $(Q,\Sigma)$. Applied to $S$, this gives a law $m\in\mathcal M$ with $m\left(\exp\left(iX(f)\right)\right)=S(f)$ for every $f\in\mathcal S$. If $m'\in\mathcal M$ has the same property, then $m'\left(\exp\left(iX(f)\right)\right)=m\left(\exp\left(iX(f)\right)\right)$ in particular for every $f\in\mathcal D_0$, and $m'=m$ by Lemma~\ref{lem:density-exponentials}(ii).

\emph{Exponential moments.} Fix $f\in\mathcal S$. For $q\ge0$ let $N(0,q)$ be the centred normal law on $\mathbb R$ with variance $q$, the unit mass at $0$ if $q=0$. Its characteristic function is $t\mapsto\exp\left(-\frac{1}{2}\cdot t^2\cdot q\right)$, and for every Borel set $B\subset\mathbb R$ the function $q\mapsto N(0,q)(B)$ on $[0,\infty)$ is a Borel function. Therefore
\[
\lambda(B)=\mathbb E\left[N\left(0,q_u(f)\right)(B)\right]\qquad\left(B\subset\mathbb R\text{ a Borel set}\right)
\]
defines a probability measure $\lambda$ on $\mathbb R$, the mixture of the laws $N\left(0,q_u(f)\right)$, $u\in\mathcal U$. For every Borel function $\eta$ on $\mathbb R$ which is either bounded and complex valued or takes its values in $[0,\infty]$ one has
\[
\int\lambda(\mathrm{d}x)\,\eta(x)=\mathbb E\int N\left(0,q_u(f)\right)(\mathrm{d}x)\,\eta(x);
\]
for indicator functions this is the definition of $\lambda$, and the general case follows by linearity and monotone convergence. As $X(tf)=t\cdot X(f)$ and $q_u(tf)=t^2\cdot q_u(f)$ for $t\in\mathbb R$,
\[
m\left(\exp\left(i\,t\cdot X(f)\right)\right)=S(tf)=\mathbb E\exp\left(-\frac{1}{2}\cdot t^2\cdot q_u(f)\right)=\int\lambda(\mathrm{d}x)\,\exp\left(i\,t\cdot x\right)\qquad\left(t\in\mathbb R\right).
\]
By the uniqueness theorem for the characteristic functions of probability measures on $\mathbb R$, the law of $X(f)$ under $m$ is $\lambda$. Consequently
\begin{align*}
m\left(\exp X(f)\right)&=\int\lambda(\mathrm{d}x)\,\exp(x)=\mathbb E\int N\left(0,q_u(f)\right)(\mathrm{d}x)\,\exp(x)\\
&=\mathbb E\exp\left(\frac{1}{2}\cdot q_u(f)\right)\le\exp\left(\frac{1}{2}\cdot\|f\|^2\right),
\end{align*}
the inequality because $q_u(f)\le\|f\|^2$ for every $u$. This is~\eqref{eq:setup-G}.

\emph{Invariance.} Let $g$ and $u\mapsto u^g$ be as in the second paragraph of the lemma. By the definitions of Section~\ref{sec:euclidean-maps-actions}, $\left((g^{-1})_*m\right)(F)=m\left(\hat gF\right)$ for every bounded measurable function $F$ on $Q$, and $\hat g\exp\left(iX(f)\right)=\exp\left(iX(g\cdot f)\right)$. Hence, for $f\in\mathcal S$,
\begin{align*}
\left((g^{-1})_*m\right)\left(\exp\left(iX(f)\right)\right)&=m\left(\exp\left(iX(g\cdot f)\right)\right)=\mathbb E\exp\left(-\frac{1}{2}\cdot q_u(g\cdot f)\right)=\mathbb E\exp\left(-\frac{1}{2}\cdot q_{u^g}(f)\right)\\
&=\mathbb E\exp\left(-\frac{1}{2}\cdot q_u(f)\right)=m\left(\exp\left(iX(f)\right)\right);
\end{align*}
the third equality is the hypothesis $q_u(g\cdot f)=q_{u^g}(f)$, and the fourth holds because $u\mapsto u^g$ preserves $\mathbb P$. By the uniqueness shown above, $(g^{-1})_*m=m$. For Euclidean maps $g_1,g_2$ one has $(g_1g_2)_*=(g_1)_*(g_2)_*$, because $(g_1g_2)\omega=g_1(g_2\omega)$ for $\omega\in Q$; hence $g_*m=g_*\left((g^{-1})_*m\right)=m$.
\end{proof}

\begin{definition}[Telegraph processes]
For $t\ge0$ and $z,z'\in\{0,1\}$ let
\[
p_t(z,z')=\frac{1}{2}\cdot\left(1+(-1)^{z+z'}\cdot\exp(-2t)\right),
\]
and for a function $\psi\colon\{0,1\}\to\mathbb C$ let $p_t\psi\colon\{0,1\}\to\mathbb C$ be the function $(p_t\psi)(z)=\sum_{z'\in\{0,1\}}p_t(z,z')\cdot\psi(z')$. A telegraph process is a stationary Markov jump process with state space $\{0,1\}$, with both jump rates equal to $1$ and with right continuous paths. Precisely, it is a family $\varsigma=\left(\varsigma(s)\right)_{s\in\mathbb R}$ of random variables with values in $\{0,1\}$, defined on a probability space with probability $\mathbb P$ and expectation $\mathbb E$, such that every path $s\mapsto\varsigma(s)$ is right continuous and has only finitely many points of discontinuity in every bounded interval, hence is a step function with finitely many jumps in every bounded interval, and such that
\[
\mathbb P\left(\varsigma(s_1)=z_1,\dots,\varsigma(s_n)=z_n\right)=\frac{1}{2}\cdot\prod_{i=1}^{n-1}p_{s_{i+1}-s_i}(z_i,z_{i+1})
\]
for every integer $n\ge1$, all real numbers $s_1<s_2<\dots<s_n$ and all $z_1,\dots,z_n\in\{0,1\}$. For $n=1$ the product is empty; thus $\mathbb P\left(\varsigma(s)=z\right)=\frac{1}{2}$ for all $s\in\mathbb R$ and $z\in\{0,1\}$. For an interval $I\subset\mathbb R$ let $\mathcal F_I=\sigma\left(\varsigma(s):s\in I\right)$, and let $\check\varsigma=\left(\check\varsigma(s)\right)_{s\in\mathbb R}$ be the reversed process, $\check\varsigma(s)=\varsigma(-s)$.
\end{definition}

\paragraph{A realisation.}
Let $\mathcal N$ be a Poisson point process on $\mathbb R$ with intensity one, that is, a random locally finite subset of $\mathbb R$ such that, for pairwise disjoint bounded Borel sets $B_1,\dots,B_k\subset\mathbb R$, the numbers $\#\left|\mathcal N\cap B_1\right|,\dots,\#\left|\mathcal N\cap B_k\right|$ are independent and Poisson distributed, with means equal to the Lebesgue measures of $B_1,\dots,B_k$. Let $\zeta$ be a random variable which is uniform on $\{0,1\}$ and independent of $\mathcal N$. Put $\ell(s)=\#\left|\mathcal N\cap(0,s]\right|$ for $s\ge0$ and $\ell(s)=-\#\left|\mathcal N\cap(s,0]\right|$ for $s<0$. The function $\ell$ is right continuous, it increases by $1$ at each point of $\mathcal N$ and is constant on every interval which contains no point of $\mathcal N$, and
\[
\ell(t)-\ell(s)=\#\left|\mathcal N\cap(s,t]\right|\qquad\left(s<t\right).
\]
For $s\in\mathbb R$ let $\varsigma(s)\in\{0,1\}$ be the residue of $\zeta+\ell(s)$ modulo $2$. Then $\varsigma(0)=\zeta$, every path of $\varsigma$ is right continuous, its points of discontinuity are the points of $\mathcal N$, of which every bounded interval contains finitely many, and
\[
\varsigma(t)-\varsigma(s)\equiv\#\left|\mathcal N\cap(s,t]\right|\pmod{2}\qquad\left(s<t\right).
\]
For every $s\in\mathbb R$ the variable $\varsigma(s)$ is uniform on $\{0,1\}$ and independent of $\mathcal N$. Indeed, $\varsigma(s)=z$ holds if and only if $\zeta$ is the residue of $z-\ell(s)$ modulo $2$; since $\zeta$ is uniform on $\{0,1\}$ and independent of $\mathcal N$, it follows that $\mathbb P\left(\{\varsigma(s)=z\}\cap E\right)=\frac{1}{2}\cdot\mathbb P(E)$ for $z\in\{0,1\}$ and every event $E$ determined by $\mathcal N$.

Now let $s_1<\dots<s_n$ and $z_1,\dots,z_n\in\{0,1\}$, and put $K_i=\#\left|\mathcal N\cap(s_i,s_{i+1}]\right|$ for $1\le i<n$. The random variables $\varsigma(s_1),K_1,\dots,K_{n-1}$ are independent, because $K_1,\dots,K_{n-1}$ are determined by $\mathcal N$ and count its points in pairwise disjoint intervals, while $\varsigma(s_1)$ is independent of $\mathcal N$; and $K_i$ is Poisson distributed with mean $s_{i+1}-s_i$. A Poisson distributed random variable with mean $t$ is even with probability $\exp(-t)\cdot\cosh t=\frac{1}{2}\cdot\left(1+\exp(-2t)\right)$ and odd with probability $\frac{1}{2}\cdot\left(1-\exp(-2t)\right)$; and $\varsigma(s_{i+1})=\varsigma(s_i)$ if and only if $K_i$ is even. Hence
\begin{align*}
\mathbb P\left(\varsigma(s_1)=z_1,\dots,\varsigma(s_n)=z_n\right)&=\mathbb P\left(\varsigma(s_1)=z_1\right)\cdot\prod_{i=1}^{n-1}\mathbb P\left(K_i\equiv z_{i+1}-z_i\pmod{2}\right)\\
&=\frac{1}{2}\cdot\prod_{i=1}^{n-1}p_{s_{i+1}-s_i}(z_i,z_{i+1}).
\end{align*}
Thus $\varsigma$ is a telegraph process; in particular, telegraph processes exist.

\paragraph{The Markov property.}
Let $\varsigma$ be a telegraph process, and let $r\in\mathbb R$. Let $A=\{\varsigma(s_1)=x_1,\dots,\varsigma(s_k)=x_k\}$ with $k\ge0$, real numbers $s_1<\dots<s_k\le r$ and $x_1,\dots,x_k\in\{0,1\}$, and let $B=\{\varsigma(t_1)=y_1,\dots,\varsigma(t_l)=y_l\}$ with $l\ge0$, real numbers $r\le t_1<\dots<t_l$ and $y_1,\dots,y_l\in\{0,1\}$. Let $z\in\{0,1\}$. If $s_k=r$ and $x_k\ne z$, then $A\cap\{\varsigma(r)=z\}$ is empty, and if $t_1=r$ and $y_1\ne z$, then $B\cap\{\varsigma(r)=z\}$ is empty; in these cases both sides of the identity below vanish. In the other cases the defining formula, applied to the times of $A$, the time $r$ and the times of $B$, expresses $\mathbb P\left(A\cap B\cap\{\varsigma(r)=z\}\right)$ as $\frac{1}{2}$ times the product of the factors $p$ that belong to consecutive times up to $r$ and of those that belong to consecutive times from $r$ on, and it expresses $\mathbb P\left(A\cap\{\varsigma(r)=z\}\right)$ and $\mathbb P\left(B\cap\{\varsigma(r)=z\}\right)$ as $\frac{1}{2}$ times the first and $\frac{1}{2}$ times the second of these two products. As $\mathbb P\left(\varsigma(r)=z\right)=\frac{1}{2}$, it follows that
\[
\mathbb P\left(A\cap B\cap\{\varsigma(r)=z\}\right)\cdot\mathbb P\left(\varsigma(r)=z\right)=\mathbb P\left(A\cap\{\varsigma(r)=z\}\right)\cdot\mathbb P\left(B\cap\{\varsigma(r)=z\}\right).
\]
The events of the form of $A$, together with the empty set, form a $\pi$-system which contains the whole space and generates $\mathcal F_{(-\infty,r]}$; the events of the form of $B$, together with the empty set, form a $\pi$-system which contains the whole space and generates $\mathcal F_{[r,\infty)}$. For fixed $B$ both sides of the last identity are finite measures in $A$, and for fixed $A$ they are finite measures in $B$. By Dynkin's $\pi$-$\lambda$ theorem, applied twice, the identity holds for all $A\in\mathcal F_{(-\infty,r]}$ and $B\in\mathcal F_{[r,\infty)}$: the $\sigma$-algebras $\mathcal F_{(-\infty,r]}$ and $\mathcal F_{[r,\infty)}$ are conditionally independent given $\varsigma(r)$. This is the Markov property. Summation of the identity over $z$ gives $\mathbb P(A\cap B)=\mathbb E\left[1_A\cdot\chi_B\left(\varsigma(r)\right)\right]$ with $\chi_B(z)=2\cdot\mathbb P\left(B\cap\{\varsigma(r)=z\}\right)$. Consequently, first for indicator functions and then by linearity and uniform approximation, every bounded complex $\mathcal F_{[r,\infty)}$-measurable random variable $\Psi$ satisfies
\begin{gather*}
\mathbb E\left[\Psi\mid\mathcal F_{(-\infty,r]}\right]=\mathbb E\left[\Psi\mid\varsigma(r)\right]=\chi\left(\varsigma(r)\right),\\
\chi(z)=\mathbb E\left[\Psi\mid\varsigma(r)=z\right]=2\cdot\mathbb E\left[\Psi\cdot1_{\{\varsigma(r)=z\}}\right]\qquad\left(z\in\{0,1\}\right).
\end{gather*}
Let $t\ge0$ and $\psi\colon\{0,1\}\to\mathbb C$. The defining formula for the two times $r<r+t$ if $t>0$, and the values $p_0(z,z)=1$, $p_0(z,1-z)=0$ if $t=0$, give $2\cdot\mathbb E\left[\psi\left(\varsigma(r+t)\right)\cdot1_{\{\varsigma(r)=z\}}\right]=(p_t\psi)(z)$ for $z\in\{0,1\}$; hence
\[
\mathbb E\left[\psi\left(\varsigma(r+t)\right)\mid\mathcal F_{(-\infty,r]}\right]=(p_t\psi)\left(\varsigma(r)\right).
\]

\begin{lemma}[Properties of telegraph processes]
Let $\varsigma$ be a telegraph process.
\begin{enumerate}
\item[(i)] For every $r\in\mathbb R$ the processes $\left(\varsigma(s+r)\right)_{s\in\mathbb R}$ and $\check\varsigma$ have the law of $\varsigma$, that is, for every integer $n\ge1$, all $s_1,\dots,s_n\in\mathbb R$ and all $z_1,\dots,z_n\in\{0,1\}$
\begin{multline}\label{eq:ergodic-components-4}
\mathbb P\left(\varsigma(s_1+r)=z_1,\dots,\varsigma(s_n+r)=z_n\right)=\mathbb P\left(\check\varsigma(s_1)=z_1,\dots,\check\varsigma(s_n)=z_n\right)\\
=\mathbb P\left(\varsigma(s_1)=z_1,\dots,\varsigma(s_n)=z_n\right).
\end{multline}
\item[(ii)] Let $r\in\mathbb R$ and $t\ge0$, and let $\Phi$ and $\Psi$ be bounded complex random variables such that $\Phi$ is $\mathcal F_{(-\infty,r]}$-measurable and $\Psi$ is $\mathcal F_{[r+t,\infty)}$-measurable. Then
\begin{equation}\label{eq:ergodic-components-5}
\left|\mathbb E\left[\Phi\cdot\Psi\right]-\mathbb E[\Phi]\cdot\mathbb E[\Psi]\right|\le\exp(-2t)\cdot\sup|\Phi|\cdot\sup|\Psi|.
\end{equation}
\item[(iii)] Let $I\subset\mathbb R$ be a non-empty bounded interval of length $|I|$, and let $z\in\{0,1\}$. Then
\begin{equation}\label{eq:ergodic-components-6}
\mathbb P\left(\varsigma(s)=z\text{ for every }s\in I\right)=\frac{1}{2}\cdot\exp\left(-|I|\right)>0.
\end{equation}
\item[(iv)] Let $C_c((0,\infty))$ be the set of the real continuous functions on $\mathbb R$ whose support is compact and contained in $(0,\infty)$, and for a Borel function $y\colon\mathbb R\to\{0,1\}$ and $\phi\in C_c((0,\infty))$ put
\[
L_\phi(y)=\int\mathrm{d}s\,\phi(s)\cdot y(s).
\]
Let $G$ be a bounded measurable function of these integrals, that is, $G(y)=g\left(\left(L_\phi(y)\right)_{\phi\in C_c((0,\infty))}\right)$ for every Borel function $y\colon\mathbb R\to\{0,1\}$, where $g\colon\mathbb R^{C_c((0,\infty))}\to\mathbb C$ is bounded and measurable for the product $\sigma$-algebra. The paths of $\varsigma$ and of $\check\varsigma$ are right continuous, respectively left continuous, hence Borel functions, so that $G(\varsigma)$ and $G(\check\varsigma)$ are defined. Then $G(\varsigma)$ is $\mathcal F_{(0,\infty)}$-measurable, and
\begin{equation}\label{eq:ergodic-components-7}
\mathbb E\left[\overline{G(\check\varsigma)}\cdot G(\varsigma)\right]=\mathbb E\left[\left|\mathbb E\left[G(\varsigma)\mid\varsigma(0)\right]\right|^2\right]\ge0.
\end{equation}
\end{enumerate}
\end{lemma}

\begin{proof}
(i) If $s_p=s_q$ and $z_p\ne z_q$ for some $p,q$, the three probabilities in~\eqref{eq:ergodic-components-4} vanish. Otherwise repeated pairs $(s_p,z_p)$ may be omitted and the remaining pairs arranged so that the times increase, and this changes none of the three probabilities. It therefore suffices to treat $s_1<\dots<s_n$. The times $s_1+r<\dots<s_n+r$ have the same differences as $s_1,\dots,s_n$, and the right side of the defining formula depends only on these differences; this gives the first process the finite-dimensional distributions of $\varsigma$. For $\check\varsigma$, the increasing times $-s_n<\dots<-s_1$ carry the values $z_n,\dots,z_1$, and the defining formula together with the symmetry $p_t(z,z')=p_t(z',z)$ gives
\[
\mathbb P\left(\check\varsigma(s_1)=z_1,\dots,\check\varsigma(s_n)=z_n\right)=\frac{1}{2}\cdot\prod_{i=1}^{n-1}p_{s_{i+1}-s_i}(z_{i+1},z_i)=\frac{1}{2}\cdot\prod_{i=1}^{n-1}p_{s_{i+1}-s_i}(z_i,z_{i+1}).
\]
A probability measure on $\{0,1\}^{\mathbb R}$ with the product $\sigma$-algebra is determined by its finite-dimensional marginals, by Dynkin's $\pi$-$\lambda$ theorem. Hence both processes have the law of $\varsigma$.

(ii) For $z\in\{0,1\}$ let $\varphi(z)=\mathbb E\left[\Phi\mid\varsigma(r)=z\right]$ and $\psi(z)=\mathbb E\left[\Psi\mid\varsigma(r+t)=z\right]$; these conditional expectations are defined because $\mathbb P\left(\varsigma(s)=z\right)=\frac{1}{2}$ for every $s$, and $|\varphi(z)|\le\sup|\Phi|$, $|\psi(z)|\le\sup|\Psi|$. By the Markov property at the time $r+t$ one has $\mathbb E\left[\Psi\mid\mathcal F_{(-\infty,r+t]}\right]=\psi\left(\varsigma(r+t)\right)$, and by the Markov property at the time $r$ one has $\mathbb E\left[\psi\left(\varsigma(r+t)\right)\mid\mathcal F_{(-\infty,r]}\right]=(p_t\psi)\left(\varsigma(r)\right)$. As $\Phi$ is measurable for $\mathcal F_{(-\infty,r]}\subset\mathcal F_{(-\infty,r+t]}$, it follows that
\begin{align*}
\mathbb E\left[\Phi\cdot\Psi\right]&=\mathbb E\left[\Phi\cdot\psi\left(\varsigma(r+t)\right)\right]=\mathbb E\left[\Phi\cdot(p_t\psi)\left(\varsigma(r)\right)\right]\\
&=\mathbb E\left[\varphi\left(\varsigma(r)\right)\cdot(p_t\psi)\left(\varsigma(r)\right)\right].
\end{align*}
Since $\varsigma(r+t)$ is uniform on $\{0,1\}$, one has $\mathbb E[\Psi]=\frac{1}{2}\cdot\left(\psi(0)+\psi(1)\right)$, and the definition of $p_t$ gives, for $z\in\{0,1\}$,
\begin{align*}
(p_t\psi)(z)&=\frac{1}{2}\cdot\left(\psi(0)+\psi(1)\right)+\frac{1}{2}\cdot\exp(-2t)\cdot\left(\psi(z)-\psi(1-z)\right)\\
&=\mathbb E[\Psi]+\exp(-2t)\cdot\left(\psi(z)-\mathbb E[\Psi]\right),
\end{align*}
and
\[
\left|\psi(z)-\mathbb E[\Psi]\right|=\frac{1}{2}\cdot\left|\psi(z)-\psi(1-z)\right|\le\sup|\Psi|.
\]
As $\mathbb E\left[\varphi\left(\varsigma(r)\right)\right]=\mathbb E[\Phi]$, one obtains
\[
\mathbb E\left[\Phi\cdot\Psi\right]-\mathbb E[\Phi]\cdot\mathbb E[\Psi]=\exp(-2t)\cdot\mathbb E\left[\varphi\left(\varsigma(r)\right)\cdot\left(\psi\left(\varsigma(r)\right)-\mathbb E[\Psi]\right)\right],
\]
and the absolute value of the right side is at most $\exp(-2t)\cdot\sup|\Phi|\cdot\sup|\Psi|$.

(iii) Let $a\le b$ be the end points of $I$, so that $|I|=b-a$, and let $D$ be the countable set of the points of $I$ that are rational or end points of $I$. A point $s\in I\setminus D$ satisfies $a<s<b$ and is the limit of a decreasing sequence of rational numbers in $(s,b)\subset I$. As the paths are right continuous, the event in~\eqref{eq:ergodic-components-6} is the intersection of the events $\{\varsigma(s)=z\}$, $s\in D$, and its probability is the limit of the probabilities of finite intersections of these events. It is therefore determined by the finite-dimensional distributions, and it is the same for all telegraph processes; so it may be computed for the realisation above. There $\varsigma(s)=z$ for every $s\in I$ if and only if $\varsigma(a)=z$ and $\mathcal N$ has no point in $I\setminus\{a\}$. Indeed, if these two conditions hold and $s\in I$, $s>a$, then $(a,s]\subset I\setminus\{a\}$ and $\varsigma(s)\equiv\varsigma(a)+\#\left|\mathcal N\cap(a,s]\right|=\varsigma(a)$ modulo $2$. Conversely, let $\varsigma(s)=z$ for every $s\in I$. Then $\varsigma(a)=z$, directly if $a\in I$, and by right continuity if $a\notin I$, in which case $(a,b)\subset I$ and $b>a$, as $I$ is not empty. If $\mathcal N$ had a point $t\in I\setminus\{a\}$, then, $\mathcal N$ being locally finite, there would be $s'\in(a,t)$ such that $\mathcal N\cap(s',t]=\{t\}$; then $s'\in I$ and $\varsigma(t)-\varsigma(s')\equiv1$ modulo $2$, which is impossible. The variable $\varsigma(a)$ is uniform on $\{0,1\}$ and independent of $\mathcal N$, and the number of the points of $\mathcal N$ in the set $I\setminus\{a\}$, whose Lebesgue measure is $|I|$, is Poisson distributed with mean $|I|$. Hence the probability in~\eqref{eq:ergodic-components-6} equals $\frac{1}{2}\cdot\exp\left(-|I|\right)$, which is positive.

(iv) Let $y$ be a path of $\varsigma$ or of $\check\varsigma$, and let $\phi\in C_c((0,\infty))$. The function $y$ is right continuous or left continuous, so the set of its points of discontinuity is countable. The function $\phi\cdot y$ is bounded, vanishes outside a compact set and is continuous outside a countable set; hence it is Riemann integrable, and
\[
L_\phi(y)=\lim_{n\to\infty}2^{-n}\cdot\sum_{k\in\mathbb Z}\phi\left(k\cdot2^{-n}\right)\cdot y\left(k\cdot2^{-n}\right).
\]
Only points $k\cdot2^{-n}>0$ contribute to these sums. Therefore $L_\phi(\varsigma)$ is $\mathcal F_{(0,\infty)}$-measurable, and $L_\phi(\check\varsigma)$, which is a limit of combinations of the values $\varsigma\left(-k\cdot2^{-n}\right)$, $k\cdot2^{-n}>0$, is $\mathcal F_{(-\infty,0)}$-measurable. As $g$ is measurable for the product $\sigma$-algebra, $G(\varsigma)$ is $\mathcal F_{(0,\infty)}$-measurable and $G(\check\varsigma)$ is $\mathcal F_{(-\infty,0)}$-measurable.

For $x\in\{0,1\}^{\mathbb R}$ and $\phi\in C_c((0,\infty))$ the Riemann sums $2^{-n}\cdot\sum_{k\in\mathbb Z}\phi\left(k\cdot2^{-n}\right)\cdot x\left(k\cdot2^{-n}\right)$ are bounded in $n$; let $\tilde L_\phi(x)$ be their upper limit as $n\to\infty$. The map
\[
\{0,1\}^{\mathbb R}\to\mathbb C\times\{0,1\},\qquad x\mapsto\left(g\left(\left(\tilde L_\phi(x)\right)_{\phi\in C_c((0,\infty))}\right),\ x(0)\right),
\]
is measurable for the product $\sigma$-algebra of $\{0,1\}^{\mathbb R}$, and by the convergence of the Riemann sums it maps $\varsigma$ to $\left(G(\varsigma),\varsigma(0)\right)$ and $\check\varsigma$ to $\left(G(\check\varsigma),\check\varsigma(0)\right)$. By~\eqref{eq:ergodic-components-4} the processes $\check\varsigma$ and $\varsigma$ have the same law on $\{0,1\}^{\mathbb R}$, so that $\left(G(\check\varsigma),\check\varsigma(0)\right)$ and $\left(G(\varsigma),\varsigma(0)\right)$ have the same law. As $\check\varsigma(0)=\varsigma(0)$, this gives
\[
\mathbb E\left[G(\check\varsigma)\mid\varsigma(0)\right]=\mathbb E\left[G(\varsigma)\mid\varsigma(0)\right].
\]
The random variable $\overline{G(\check\varsigma)}$ is $\mathcal F_{(-\infty,0]}$-measurable and $G(\varsigma)$ is $\mathcal F_{[0,\infty)}$-measurable. By the Markov property with $r=0$, $\mathbb E\left[G(\varsigma)\mid\mathcal F_{(-\infty,0]}\right]=\mathbb E\left[G(\varsigma)\mid\varsigma(0)\right]$, and therefore
\begin{align*}
\mathbb E\left[\overline{G(\check\varsigma)}\cdot G(\varsigma)\right]&=\mathbb E\left[\overline{G(\check\varsigma)}\cdot\mathbb E\left[G(\varsigma)\mid\varsigma(0)\right]\right]\\
&=\mathbb E\left[\overline{\mathbb E\left[G(\check\varsigma)\mid\varsigma(0)\right]}\cdot\mathbb E\left[G(\varsigma)\mid\varsigma(0)\right]\right]\\
&=\mathbb E\left[\left|\mathbb E\left[G(\varsigma)\mid\varsigma(0)\right]\right|^2\right]\ge0;
\end{align*}
the second equality holds because $\mathbb E\left[G(\varsigma)\mid\varsigma(0)\right]$ is a function of $\varsigma(0)$, and the third by the equality of the two conditional expectations shown above.
\end{proof}

\begin{lemma}[Reflection positivity of products]\label{lem:reflection-positivity-products}
For $i=1,2$ let $(Q_i,\mathcal F_i,\rho_i)$ be a probability space, let $\vartheta_i\colon Q_i\to Q_i$ be a measurable involution which preserves $\rho_i$, and let $\mathcal F_{i+}\subset\mathcal F_i$ be a $\sigma$-algebra such that
\[
\int\mathrm{d}\rho_i\,\overline{G\circ\vartheta_i}\cdot G\ge0\qquad\text{for every bounded }\mathcal F_{i+}\text{-measurable function }G\colon Q_i\to\mathbb C.
\]
Let $\vartheta_1\times\vartheta_2$ denote the map $(\omega_1,\omega_2)\mapsto(\vartheta_1\omega_1,\vartheta_2\omega_2)$ of $Q_1\times Q_2$ into itself. Then
\[
\int\mathrm{d}(\rho_1\otimes\rho_2)\,\overline{G\circ(\vartheta_1\times\vartheta_2)}\cdot G\ge0
\]
for every bounded $\mathcal F_{1+}\otimes\mathcal F_{2+}$-measurable function $G\colon Q_1\times Q_2\to\mathbb C$.
\end{lemma}

\begin{proof}
For functions $g$ on $Q_1$ and $h$ on $Q_2$ let $g\otimes h$ be the function $(\omega_1,\omega_2)\mapsto g(\omega_1)\cdot h(\omega_2)$ on $Q_1\times Q_2$. Let first
\[
G=\sum_{a=1}^ng_a\otimes h_a,
\]
where $n\ge1$, the functions $g_a$ are bounded, complex and $\mathcal F_{1+}$-measurable, and the functions $h_a$ are bounded, complex and $\mathcal F_{2+}$-measurable. For $a,b\in\{1,\dots,n\}$ put
\[
M^1_{ab}=\int\mathrm{d}\rho_1\,\overline{g_a\circ\vartheta_1}\cdot g_b,\qquad M^2_{ab}=\int\mathrm{d}\rho_2\,\overline{h_a\circ\vartheta_2}\cdot h_b.
\]
By Fubini's theorem
\[
\int\mathrm{d}(\rho_1\otimes\rho_2)\,\overline{G\circ(\vartheta_1\times\vartheta_2)}\cdot G=\sum_{a,b=1}^nM^1_{ab}\cdot M^2_{ab}.
\]
For $c\in\mathbb C^n$ the hypothesis, applied to the bounded $\mathcal F_{1+}$-measurable function $\sum_{b=1}^nc_b\cdot g_b$, gives $\sum_{a,b=1}^n\overline{c_a}\cdot c_b\cdot M^1_{ab}\ge0$. Thus the matrix $M^1=\left(M^1_{ab}\right)_{a,b=1}^n$ is positive semidefinite, in particular Hermitian, and in the same way so is $M^2$. By the Schur product theorem the matrix $\left(M^1_{ab}\cdot M^2_{ab}\right)_{a,b=1}^n$ is positive semidefinite. The sum of its entries is the value of its quadratic form at the vector $(1,\dots,1)$ and is therefore non-negative. This proves the assertion for these $G$.

Let now $G$ be bounded and $\mathcal F_{1+}\otimes\mathcal F_{2+}$-measurable, and write $\|\cdot\|_{L^2}$ for the norm of $L^2(\rho_1\otimes\rho_2)$ on $Q_1\times Q_2$ with the $\sigma$-algebra $\mathcal F_1\otimes\mathcal F_2$. The image of $\rho_1\otimes\rho_2$ under $\vartheta_1\times\vartheta_2$ coincides with $\rho_1\otimes\rho_2$ on the measurable rectangles, which form a $\pi$-system generating $\mathcal F_1\otimes\mathcal F_2$; so $\vartheta_1\times\vartheta_2$ preserves $\rho_1\otimes\rho_2$, and $G'\mapsto G'\circ(\vartheta_1\times\vartheta_2)$ is an isometry of $L^2(\rho_1\otimes\rho_2)$. For $G'\in L^2(\rho_1\otimes\rho_2)$ put
\[
\beta(G')=\int\mathrm{d}(\rho_1\otimes\rho_2)\,\overline{G'\circ(\vartheta_1\times\vartheta_2)}\cdot G'.
\]
By the Cauchy--Schwarz inequality, $\left|\beta(G')-\beta(G'')\right|\le\left\|G'-G''\right\|_{L^2}\cdot\left(\left\|G'\right\|_{L^2}+\left\|G''\right\|_{L^2}\right)$ for $G',G''\in L^2(\rho_1\otimes\rho_2)$. The finite sums $\sum_ag_a\otimes h_a$ considered above are dense in the closed subspace of $L^2(\rho_1\otimes\rho_2)$ formed by the classes which contain an $\mathcal F_{1+}\otimes\mathcal F_{2+}$-measurable function, and the class of $G$ lies in this subspace. Hence there are such sums $G_k$ with $\left\|G-G_k\right\|_{L^2}\to0$ as $k\to\infty$, and $\beta(G)=\lim_{k\to\infty}\beta(G_k)\ge0$.
\end{proof}

\subsubsection{An average of six anisotropic white noises}\label{sec:average-six-anisotropic}

For $w\in O_{\mathrm{cub}}$ let $\pi_w\in S_3$ be the permutation such that $w\,e_i=\varepsilon_i\cdot e_{\pi_w(i)}$ for $i=1,2,3$, with $\varepsilon\in\{\pm1\}^3$ (Section~\ref{sec:euclidean-maps-actions}). The map $w\mapsto\pi_w$ is a homomorphism of $O_{\mathrm{cub}}$ onto $S_3$, and its kernel is $N_3$.

\begin{example}[A mixture of six white noises]\label{ex:mixture-six-white}
Let $D=(d_{12},d_{13},d_{23})$ be a triple of pairwise distinct numbers in $(0,1]$, and put $d_{li}=d_{il}$ for $1\le i<l\le3$. For $f,g\in\mathcal S$ let
\[
q_D(f,g)=\sum_{i<l}d_{il}\cdot\int\mathrm{d}x\,f_{il}(x)\cdot g_{il}(x),\qquad q_D(f)=q_D(f,f)=\sum_{i<l}d_{il}\cdot\|f_{il}\|^2,
\]
where $\|f_{il}\|^2=\int\mathrm{d}x\,f_{il}(x)^2$. As $\|f\|^2=\sum_{i<l}\|f_{il}\|^2$ and $0<d_{il}\le1$, the quadratic form $q_D$ on $\mathcal S$ satisfies $0\le q_D(f)\le\|f\|^2$. Let $\gamma_D\in\mathcal M$ be the law that Lemma~\ref{lem:gaussian-mixtures} provides for the probability space consisting of one point and for the form $q_D$; thus
\[
\gamma_D\left(\exp\left(iX(f)\right)\right)=\exp\left(-\frac{1}{2}q_D(f)\right)\qquad(f\in\mathcal S).
\]
For $\pi\in S_3$ let $\pi D$ be the triple with the entries
\[
(\pi D)_{ab}=d_{\pi^{-1}(a)\pi^{-1}(b)}\qquad(1\le a<b\le3),
\]
and put $(\pi D)_{ba}=(\pi D)_{ab}$; then $(\pi D)_{\pi(i)\pi(l)}=d_{il}$ for all $i,l\in\{1,2,3\}$ with $i\ne l$. The entries of $\pi D$ are those of $D$ in another order; so $\pi D$ is again a triple of pairwise distinct numbers in $(0,1]$, and $q_{\pi D}$ and $\gamma_{\pi D}$ are defined. Let
\[
m_1=\frac{1}{6}\sum_{\pi\in S_3}\gamma_{\pi D}.
\]
For every triple $D$ as above, every $w\in O_{\mathrm{cub}}$ and every $f\in\mathcal S$,
\[
q_D\left(w^{-1}\cdot f\right)=q_{\pi_wD}(f),\qquad\text{and consequently}\qquad w_*\gamma_D=\gamma_{\pi_wD}.
\]
Moreover, the following assertions hold.
\begin{enumerate}
\item[(1)] For every triple $D$ as above, the law $\gamma_D$ satisfies~\eqref{eq:setup-G} and $(T_{123})$, one has $\varepsilon_*\gamma_D=\gamma_D$ for every $\varepsilon\in N_3$, and $\gamma_D$ satisfies $(RP_j)$ for $j=1,2,3$.
\item[(2)] For every triple $D$ as above and all $f,g\in\mathcal S$,
\[
\lim_{|a|\to\infty}\gamma_D\left(\exp\left(iX(f)\right)\cdot\exp\left(iX(\tau_a\cdot g)\right)\right)=\gamma_D\left(\exp\left(iX(f)\right)\right)\cdot\gamma_D\left(\exp\left(iX(g)\right)\right).
\]
The law $\gamma_D$ is mixing, and $\gamma_D\in\mathcal M^{123}_{\mathrm{erg}}$.
\item[(3)] The six triples $\pi D$, $\pi\in S_3$, are pairwise distinct, and so are the six laws $\gamma_{\pi D}$, $\pi\in S_3$. For $w\in O_{\mathrm{cub}}$ and $\pi\in S_3$ one has $w_*\gamma_{\pi D}=\gamma_{(\pi_w\pi)D}$; hence each $w\in O_{\mathrm{cub}}$ permutes the six laws $\gamma_{\pi D}$, and $w_*m_1=m_1$. The law $m_1$ satisfies~\eqref{eq:setup-H}.
\item[(4)] Let $\nu$ be the probability measure on $(\mathcal M,\mathcal B_{\mathcal M})$ that Theorem~\ref{thm:ergodic-decomposition}(a) assigns to $m=m_1$ and $J=\{1,2,3\}$, and for $\rho\in\mathcal M$ let $\delta_\rho$ be the probability measure on $(\mathcal M,\mathcal B_{\mathcal M})$ concentrated at $\rho$. Then
\[
\nu=\frac{1}{6}\sum_{\pi\in S_3}\delta_{\gamma_{\pi D}}.
\]
For $\pi\in S_3$ and $w\in O_{\mathrm{cub}}$ one has $w_*\gamma_{\pi D}=\gamma_{\pi D}$ if and only if $\pi_w$ is the identity, that is, if and only if $w\in N_3$.
\end{enumerate}
\end{example}

\begin{proof}[Proof of the assertions of Example~\ref{ex:mixture-six-white}]
Let $D$ be a triple of pairwise distinct numbers in $(0,1]$.

\emph{The transformation rule.} Let $w\in O_{\mathrm{cub}}$, let $\varepsilon\in\{\pm1\}^3$ be such that $w\,e_i=\varepsilon_i\cdot e_{\pi_w(i)}$ for $i=1,2,3$, and let $f\in\mathcal S$. The formula for the action of $w$ on forms in Section~\ref{sec:euclidean-maps-actions}, applied to $w^{-1}\cdot f$, gives
\[
f_{\pi_w(i)\pi_w(l)}(x)=\left(w\cdot\left(w^{-1}\cdot f\right)\right)_{\pi_w(i)\pi_w(l)}(x)=\varepsilon_i\cdot\varepsilon_l\cdot\left(w^{-1}\cdot f\right)_{il}\left(w^{-1}x\right)\qquad\left(x\in\mathbb R^3,\ i,l\in\{1,2,3\}\right).
\]
As $x\mapsto w^{-1}x$ preserves the Lebesgue measure, $\left\|\left(w^{-1}\cdot f\right)_{il}\right\|=\left\|f_{\pi_w(i)\pi_w(l)}\right\|$. When $\{i,l\}$ ranges over the three two-element subsets of $\{1,2,3\}$, so does $\{\pi_w(i),\pi_w(l)\}$; and $(\pi_wD)_{ab}$ and $\|f_{ab}\|$ do not change when $a$ and $b$ are interchanged, because $f_{ba}=-f_{ab}$. Therefore
\[
q_D\left(w^{-1}\cdot f\right)=\sum_{i<l}d_{il}\cdot\left\|f_{\pi_w(i)\pi_w(l)}\right\|^2=\sum_{i<l}(\pi_wD)_{\pi_w(i)\pi_w(l)}\cdot\left\|f_{\pi_w(i)\pi_w(l)}\right\|^2=\sum_{a<b}(\pi_wD)_{ab}\cdot\|f_{ab}\|^2=q_{\pi_wD}(f).
\]
By Section~\ref{sec:euclidean-maps-actions}, $X(f)(w\omega)=\omega\left(w^{-1}\cdot f\right)=X\left(w^{-1}\cdot f\right)(\omega)$ for $\omega\in Q$, and $w_*\gamma_D$ is the image of $\gamma_D$ under $\omega\mapsto w\omega$. Hence
\[
\left(w_*\gamma_D\right)\left(\exp\left(iX(f)\right)\right)=\gamma_D\left(\exp\left(iX\left(w^{-1}\cdot f\right)\right)\right)=\exp\left(-\frac{1}{2}q_{\pi_wD}(f)\right)=\gamma_{\pi_wD}\left(\exp\left(iX(f)\right)\right)
\]
for every $f\in\mathcal S$, and Lemma~\ref{lem:density-exponentials}(ii) gives $w_*\gamma_D=\gamma_{\pi_wD}$.

\emph{Assertion (1).} The law $\gamma_D$ satisfies~\eqref{eq:setup-G} by Lemma~\ref{lem:gaussian-mixtures}. Let $g$ be a translation $\tau_a$, $a\in\mathbb R^3$, or an element of $N_3$. Then $\|(g\cdot f)_{il}\|=\|f_{il}\|$ for $f\in\mathcal S$ and $i<l$: for $g=\tau_a$ because $(\tau_a\cdot f)_{il}(x)=f_{il}(x-a)$, and for $g\in N_3$ by the computation above, applied to $w=g^{-1}$, for which $\pi_w$ is the identity. Hence $q_D(g\cdot f)=q_D(f)$ for every $f\in\mathcal S$, and the last assertion of Lemma~\ref{lem:gaussian-mixtures}, applied to the Euclidean map $g$ and to the identity map of the one-point probability space, gives $g_*\gamma_D=\gamma_D$. Thus $\gamma_D$ satisfies $(T_{123})$ and is invariant under every element of $N_3$.

Let $j\in\{1,2,3\}$, $n\ge1$, $c=(c_1,\dots,c_n)\in\mathbb C^n$ and $f_1,\dots,f_n\in\mathcal D(\{x_j>0\})$. For $p,q\in\{1,\dots,n\}$ the form $\theta_j\cdot f_p$ is supported in $\{x_j<0\}$; so the supports of $f_q$ and $\theta_j\cdot f_p$ are disjoint, and $q_D(f_q,\theta_j\cdot f_p)=0$. Since $\theta_j\in N_3$, moreover $q_D(\theta_j\cdot f_p)=q_D(f_p)$. Consequently
\[
q_D\left(f_q-\theta_j\cdot f_p\right)=q_D(f_q)-2\cdot q_D\left(f_q,\theta_j\cdot f_p\right)+q_D\left(\theta_j\cdot f_p\right)=q_D(f_q)+q_D(f_p).
\]
As $X(f_q)-X(\theta_j\cdot f_p)=X\left(f_q-\theta_j\cdot f_p\right)$ on $Q$, it follows that
\begin{align*}
\sum_{p,q=1}^n\overline{c_p}\cdot c_q\cdot\gamma_D\left(\exp\left(i\left(X(f_q)-X(\theta_j\cdot f_p)\right)\right)\right)&=\sum_{p,q=1}^n\overline{c_p}\cdot c_q\cdot\exp\left(-\frac{1}{2}q_D(f_p)\right)\cdot\exp\left(-\frac{1}{2}q_D(f_q)\right)\\
&=\left|\sum_{q=1}^nc_q\cdot\exp\left(-\frac{1}{2}q_D(f_q)\right)\right|^2\ge0.
\end{align*}
This is $(RP_j)$.

\emph{Assertion (2).} Let $f,g\in\mathcal S$ and $a\in\mathbb R^3$. Since $\exp\left(iX(f)\right)\cdot\exp\left(iX(\tau_a\cdot g)\right)=\exp\left(iX(f+\tau_a\cdot g)\right)$, and $q_D(\tau_a\cdot g)=q_D(g)$ by the proof of (1),
\[
\gamma_D\left(\exp\left(iX(f)\right)\cdot\exp\left(iX(\tau_a\cdot g)\right)\right)=\exp\left(-\frac{1}{2}q_D(f)-\frac{1}{2}q_D(g)-q_D\left(f,\tau_a\cdot g\right)\right).
\]
Here
\[
q_D\left(f,\tau_a\cdot g\right)=\sum_{i<l}d_{il}\cdot\int\mathrm{d}x\,f_{il}(x)\cdot g_{il}(x-a)\to0\qquad\text{as }|a|\to\infty,
\]
because $\int\mathrm{d}x\,\phi(x)\cdot\psi(x-a)\to0$ as $|a|\to\infty$ for all $\phi,\psi\in L^2(\mathbb R^3)$: the integral vanishes for large $|a|$ if $\phi$ and $\psi$ vanish outside a bounded set, and the general case follows by approximation in $L^2(\mathbb R^3)$ and the Cauchy--Schwarz inequality. This proves the limit relation in (2).

Let now $F=\sum_pc_p\cdot\exp\left(iX(f_p)\right)$ and $G=\sum_qc'_q\cdot\exp\left(iX(g_q)\right)$ be elements of $\mathcal A$, with finitely many $f_p,g_q\in\mathcal D_0$ and $c_p,c'_q\in\mathbb C$. By Section~\ref{sec:euclidean-maps-actions}, $\overline{F}=\sum_p\overline{c_p}\cdot\exp\left(iX(-f_p)\right)$ and $U(a)G=\sum_qc'_q\cdot\exp\left(iX(\tau_a\cdot g_q)\right)$. Moreover, for $h\in\mathcal S$ the number $\gamma_D\left(\exp\left(iX(h)\right)\right)$ is real and equals $\gamma_D\left(\exp\left(iX(-h)\right)\right)$. The limit relation in (2) therefore gives
\[
\lim_{|a|\to\infty}\gamma_D\left(\overline{F}\cdot U(a)G\right)=\sum_{p,q}\overline{c_p}\cdot c'_q\cdot\overline{\gamma_D\left(\exp\left(iX(f_p)\right)\right)}\cdot\gamma_D\left(\exp\left(iX(g_q)\right)\right)=\overline{\gamma_D(F)}\cdot\gamma_D(G).
\]
By (1), $\gamma_D$ satisfies $(T_{123})$; by the criterion of Section~\ref{sec:invariant-sets-invariant}, $\gamma_D$ is mixing. Lemma~\ref{lem:clustering-criterion-separation}(a), applied with $J=\{1,2,3\}$ and the sequence $a^{(k)}=k\cdot e_1$, $k\ge1$, gives $\gamma_D\in\mathcal M^{123}_{\mathrm{erg}}$.

\emph{Assertion (3).} For a two-element set $P=\{i,l\}\subset\{1,2,3\}$ and $\pi\in S_3$ put $\pi P=\{\pi(i),\pi(l)\}$, $d_P=d_{il}$ and $(\pi D)_P=(\pi D)_{il}$; then $(\pi D)_{\pi P}=d_P$. If $\sigma\in S_3$ satisfies $\sigma P=P$ for all three two-element subsets $P$, then $\sigma$ is the identity, since every $k\in\{1,2,3\}$ is the only common element of the two two-element subsets that contain $k$. Thus $S_3$ acts faithfully on the set of the three two-element subsets of $\{1,2,3\}$, and, having six elements, it acts on this set as the group of all its permutations.

Let $\pi,\pi'\in S_3$ with $\pi D=\pi'D$. For every two-element subset $P$, the number $d_P$ is the entry of the triple $\pi D$ at the place $\pi P$ and the entry of the triple $\pi'D=\pi D$ at the place $\pi'P$; the entries of $\pi D$ being pairwise distinct, $\pi P=\pi'P$. Hence $\pi^{-1}\pi'$ fixes every two-element subset, and $\pi=\pi'$. So the six triples $\pi D$ are pairwise distinct. Let $D'=(d'_{12},d'_{13},d'_{23})$ and $D''=(d''_{12},d''_{13},d''_{23})$ be triples of pairwise distinct numbers in $(0,1]$ with $\gamma_{D'}=\gamma_{D''}$. Then $q_{D'}(f)=q_{D''}(f)$ for every $f\in\mathcal S$. For $a<b$ and an $f\in\mathcal S$ with $f_{ab}\ne0$ whose two other components vanish, this reads $d'_{ab}\cdot\|f_{ab}\|^2=d''_{ab}\cdot\|f_{ab}\|^2$, that is, $d'_{ab}=d''_{ab}$. Thus $D'=D''$, and the six laws $\gamma_{\pi D}$ are pairwise distinct.

Let $w\in O_{\mathrm{cub}}$ and $\pi\in S_3$. For $1\le a<b\le3$,
\[
\left(\pi_w(\pi D)\right)_{ab}=(\pi D)_{\pi_w^{-1}(a)\pi_w^{-1}(b)}=d_{\pi^{-1}\pi_w^{-1}(a)\,\pi^{-1}\pi_w^{-1}(b)}=\left((\pi_w\pi)D\right)_{ab},
\]
so that $\pi_w(\pi D)=(\pi_w\pi)D$, and the transformation rule, applied to the triple $\pi D$, gives $w_*\gamma_{\pi D}=\gamma_{(\pi_w\pi)D}$. Since $\pi\mapsto\pi_w\pi$ is a bijection of $S_3$ onto itself, $w_*$ permutes the six laws $\gamma_{\pi D}$, and
\[
w_*m_1=\frac{1}{6}\sum_{\pi\in S_3}w_*\gamma_{\pi D}=\frac{1}{6}\sum_{\pi\in S_3}\gamma_{(\pi_w\pi)D}=m_1.
\]
By (1), each $\gamma_{\pi D}$ satisfies~\eqref{eq:setup-G}, $(T_{123})$ and $(RP_j)$ for $j=1,2,3$. Each of these conditions is a family of linear equalities or inequalities in the law, and is therefore preserved under convex combinations of laws. Explicitly, $m_1(\Phi)=\frac{1}{6}\sum_{\pi\in S_3}\gamma_{\pi D}(\Phi)$ for every bounded or non-negative measurable function $\Phi$ on $Q$; hence $\int\mathrm{d}m_1\,\exp X(f)\le\exp\left(\frac{1}{2}\|f\|^2\right)$ for $f\in\mathcal S$; $(\tau_a)_*m_1=\frac{1}{6}\sum_{\pi\in S_3}(\tau_a)_*\gamma_{\pi D}=m_1$ for $a\in\mathbb R^3$; and, for $j\in\{1,2,3\}$, the sum in $(RP_j)$ formed with $m_1$ is the mean of the six sums formed with the laws $\gamma_{\pi D}$ and the same $n$, $c$, $f_1,\dots,f_n$, each of which is non-negative. Together with the invariance $w_*m_1=m_1$ for all $w\in O_{\mathrm{cub}}$, this shows that $m_1$ satisfies~\eqref{eq:setup-H}.

\emph{Assertion (4).} By (3), $m_1$ satisfies $(T_{123})$, so that Theorem~\ref{thm:ergodic-decomposition} applies to $m=m_1$ and $J=\{1,2,3\}$. Let $\nu'=\frac{1}{6}\sum_{\pi\in S_3}\delta_{\gamma_{\pi D}}$. This is a probability measure on $(\mathcal M,\mathcal B_{\mathcal M})$, and
\[
\int_{\mathcal M}\nu'(\mathrm{d}\rho)\,\rho(A)=\frac{1}{6}\sum_{\pi\in S_3}\gamma_{\pi D}(A)=m_1(A)\qquad(A\in\Sigma).
\]
Moreover $\mathcal M^{123}_{\mathrm{erg}}\in\mathcal B_{\mathcal M}$ by Theorem~\ref{thm:ergodic-decomposition}(a), and $\gamma_{\pi D}\in\mathcal M^{123}_{\mathrm{erg}}$ for every $\pi\in S_3$, by (2) applied to the triple $\pi D$; hence $\nu'\left(\mathcal M^{123}_{\mathrm{erg}}\right)=1$. By the uniqueness assertion of Theorem~\ref{thm:ergodic-decomposition}(a), $\nu=\nu'$.

Let $\pi\in S_3$ and $w\in O_{\mathrm{cub}}$. By (3), $w_*\gamma_{\pi D}=\gamma_{(\pi_w\pi)D}$, and the laws $\gamma_{\pi'D}$, $\pi'\in S_3$, are pairwise distinct; hence $w_*\gamma_{\pi D}=\gamma_{\pi D}$ if and only if $\pi_w\pi=\pi$, that is, if and only if $\pi_w$ is the identity. The kernel of the homomorphism $w\mapsto\pi_w$ is $N_3$; so $\pi_w$ is the identity if and only if $w\in N_3$.
\end{proof}

\begin{remark}[Fourth cumulant and ergodic components]
By Example~\ref{ex:mixture-six-white}(4), the ergodic components of $m_1$ given by Theorem~\ref{thm:ergodic-decomposition}(a) for $J=\{1,2,3\}$ are the six laws $\gamma_{\pi D}$, $\pi\in S_3$. Each of them is Gaussian: for $f\in\mathcal S$ the characteristic function $t\mapsto\gamma_{\pi D}\left(\exp\left(itX(f)\right)\right)=\exp\left(-\frac{1}{2}t^2\cdot q_{\pi D}(f)\right)$ shows that $X(f)$ is a centred normal random variable with variance $q_{\pi D}(f)$ under $\gamma_{\pi D}$. The law $m_1$ itself is not Gaussian. Indeed, let $f\in\mathcal S$. Under $m_1$ the characteristic function of $X(f)$ is $t\mapsto\frac{1}{6}\sum_{\pi\in S_3}\exp\left(-\frac{1}{2}t^2\cdot q_{\pi D}(f)\right)$; so, by the uniqueness of characteristic functions, the law of $X(f)$ under $m_1$ is the mean of the six centred normal laws with the variances $q_{\pi D}(f)$, and
\[
m_1\left(X(f)^2\right)=\frac{1}{6}\sum_{\pi\in S_3}q_{\pi D}(f),\qquad m_1\left(X(f)^4\right)=\frac{1}{6}\sum_{\pi\in S_3}3\cdot q_{\pi D}(f)^2.
\]
Hence the fourth cumulant of the centred random variable $X(f)$ under $m_1$ is
\[
\kappa_4\left(X(f)\right)=m_1\left(X(f)^4\right)-3\cdot m_1\left(X(f)^2\right)^2=3\cdot\operatorname{Var}_\pi q_{\pi D}(f)\ge0,
\]
where
\[
\operatorname{Var}_\pi q_{\pi D}(f)=\frac{1}{6}\sum_{\pi\in S_3}q_{\pi D}(f)^2-\left(\frac{1}{6}\sum_{\pi\in S_3}q_{\pi D}(f)\right)^2
\]
is the variance of the function $\pi\mapsto q_{\pi D}(f)$ with respect to the uniform probability on $S_3$. If $f_{13}=f_{23}=0$ and $f_{12}\ne0$, then $q_{\pi D}(f)=(\pi D)_{12}\cdot\|f_{12}\|^2=d_{\pi^{-1}(1)\pi^{-1}(2)}\cdot\|f_{12}\|^2$ takes three different values as $\pi$ ranges over $S_3$, because $\{\pi^{-1}(1),\pi^{-1}(2)\}$ ranges over the three two-element subsets of $\{1,2,3\}$ and $\|f_{12}\|>0$; so $\kappa_4\left(X(f)\right)>0$, whereas the fourth cumulant of a centred normal random variable vanishes. Thus $m_1$ is a law which satisfies~\eqref{eq:setup-H}, by Example~\ref{ex:mixture-six-white}(3), whose ergodic components are all Gaussian, and under which $\kappa_4\left(X(f)\right)>0$ for every $f\in\mathcal S$ with $f_{12}\ne0$ and $f_{13}=f_{23}=0$.
\end{remark}

\subsubsection{White noise modulated by telegraph processes}\label{sec:white-noise-modulated}

\begin{example}[White noise in a telegraph medium]\label{ex:white-noise-telegraph}
Let $\varsigma_1,\varsigma_2,\varsigma_3$ be independent telegraph processes, in the sense of Section~\ref{sec:auxiliary-results}, defined on one probability space; expectations with respect to its probability measure are denoted by $\mathbb E$. For $l\in\{1,2,3\}$ and an interval $I\subset\mathbb R$ let $\mathcal F^l_I=\sigma\left(\varsigma_l(s):s\in I\right)$, let $\check\varsigma_l$ be the reversed process, $\check\varsigma_l(s)=\varsigma_l(-s)$, and put
\[
v_l=\frac{1}{2}\left(1+\varsigma_l\right);
\]
every path of $v_l$ takes its values in $\left\{\frac{1}{2},1\right\}$.

Let $C_c(\mathbb R)$ be the space of continuous functions $\mathbb R\to\mathbb R$ with compact support, and, for an open interval $I\subset\mathbb R$, let $C_c(I)$ be the space of those $\phi\in C_c(\mathbb R)$ whose support is a compact subset of $I$. Let $\mathcal V$ be the set of the classes, modulo equality almost everywhere with respect to the Lebesgue measure, of the measurable functions $\mathbb R\to\left[\frac{1}{2},1\right]$, endowed with the $\sigma$-algebra generated by the maps $z\mapsto\int\mathrm{d}s\,\phi(s)\cdot z(s)$, $\phi\in C_c(\mathbb R)$, and let $\mathcal U=\mathcal V^3$, with the product $\sigma$-algebra. The elements of $\mathcal U$ are written $v=(v_1,v_2,v_3)$, and the same letters denote the random element of $\mathcal U$ whose entries are the classes of the paths of the processes $v_1,v_2,v_3$ above. By (i) below, this random element is measurable, and its entries $v_1,v_2,v_3$ are independent random elements of $\mathcal V$ which have one and the same law. The law of $v$ is denoted by $\mathbb P$, and the common law of $v_1$, $v_2$, $v_3$ on $\mathcal V$ by $\mathbb P_1$; by the independence, $\mathbb P$ is the product of three copies of $\mathbb P_1$.

For $v\in\mathcal U$, $a\in\mathbb R^3$, $w\in O_{\mathrm{cub}}$ with $w\,e_i=\varepsilon_i\cdot e_{\pi_w(i)}$, and $z\in\mathcal V$, let $v^a,v^w\in\mathcal U$ and $\check z\in\mathcal V$ be given by
\[
v^a_l(s)=v_l(s+a_l),\qquad v^w_i(s)=v_{\pi_w(i)}(\varepsilon_i\cdot s),\qquad\check z(s)=z(-s)\qquad\left(s\in\mathbb R,\ i,l\in\{1,2,3\}\right);
\]
in particular $v^{\theta_1}=(\check v_1,v_2,v_3)$. For $v\in\mathcal U$ and $f\in\mathcal S$ let
\[
\mathsf a_v(x)=v_1(x_1)\cdot v_2(x_2)\cdot v_3(x_3)\quad\left(x\in\mathbb R^3\right),\qquad q_v(f)=\int\mathrm{d}x\,\mathsf a_v(x)\cdot|f(x)|^2,\qquad\Phi_f(v)=\exp\left(-\frac{1}{2}q_v(f)\right),
\]
where $|f(x)|$ is the norm of $f(x)$ in $\Lambda^2$. The function $\mathsf a_v$ is defined almost everywhere on $\mathbb R^3$ and takes its values in $\left[\frac{1}{8},1\right]$; so $q_v$ is a quadratic form on $\mathcal S$ with $\frac{1}{8}\|f\|^2\le q_v(f)\le\|f\|^2$, and $0<\Phi_f(v)\le1$. For a bounded measurable function $H$ on $\mathcal U$, $\mathbb E_{v_1}H(v)$ denotes the function
\[
(v_2,v_3)\mapsto\int_{\mathcal V}\mathbb P_1(\mathrm{d}z)\,H(z,v_2,v_3)
\]
on $\mathcal V^2$, the expectation with respect to $v_1$ alone; further, $\mathbb E_{v_1,v_2}H(v)$ and $\mathbb E_{v_2,v_3}H(v)$ denote the functions
\[
v_3\mapsto\int_{\mathcal V}\mathbb P_1(\mathrm{d}z_1)\int_{\mathcal V}\mathbb P_1(\mathrm{d}z_2)\,H(z_1,z_2,v_3),\qquad v_1\mapsto\int_{\mathcal V}\mathbb P_1(\mathrm{d}z_2)\int_{\mathcal V}\mathbb P_1(\mathrm{d}z_3)\,H(v_1,z_2,z_3)
\]
on $\mathcal V$. By Fubini's theorem, $\mathbb E\left[H(v)\right]=\mathbb E\left[\mathbb E_{v_1}H(v)\right]=\mathbb E\left[\mathbb E_{v_1,v_2}H(v)\right]=\mathbb E_{v_2,v_3}\mathbb E_{v_1}H(v)$.

The objects so defined have the following properties.
\begin{enumerate}
\item[(i)] Let $\eta\colon\mathbb R\to\{0,1\}$ be a function which is right continuous or left continuous, and let $\phi\in C_c(\mathbb R)$. Then $\eta$ is continuous outside a countable set, it is a Borel function, and
\[
\int\mathrm{d}s\,\phi(s)\cdot\eta(s)=\lim_{n\to\infty}2^{-n}\cdot\sum_{k\in\mathbb Z}\phi\left(k\cdot2^{-n}\right)\cdot\eta\left(k\cdot2^{-n}\right).
\]
By the definition of a telegraph process, every path of $\varsigma_l$ is right continuous, so that every path of $\check\varsigma_l$ is left continuous; thus the display holds for all these paths, $l=1,2,3$. Consequently, if $\phi$ vanishes outside an interval $I$, the random variable $\int\mathrm{d}s\,\phi(s)\cdot\varsigma_l(s)$ is $\mathcal F^l_I$-measurable; the random element $v$ of $\mathcal U$ is measurable; its law $\mathbb P$ is determined by the finite-dimensional distributions of $(\varsigma_1,\varsigma_2,\varsigma_3)$ in the following sense: if $\eta_1,\eta_2,\eta_3$ are processes on some probability space such that every path of each $\eta_l$ is a right continuous or a left continuous function $\mathbb R\to\{0,1\}$, and if $(\eta_1,\eta_2,\eta_3)$ has the same finite-dimensional distributions as $(\varsigma_1,\varsigma_2,\varsigma_3)$, then the classes of the paths of $\frac{1}{2}(1+\eta_l)$, $l=1,2,3$, form a random element of $\mathcal U$ with the law $\mathbb P$; and the entries $v_1,v_2,v_3$ of $v$ are independent random elements of $\mathcal V$ with one and the same law.
\item[(ii)] For every $a\in\mathbb R^3$ and every $w\in O_{\mathrm{cub}}$, the maps $v\mapsto v^a$ and $v\mapsto v^w$ of $\mathcal U$ into itself are measurable and preserve $\mathbb P$.
\item[(iii)] Let $f\in\mathcal S$, and let $I_1,I_2,I_3\subset\mathbb R$ be open intervals such that $f(x)=0$ for every $x\in\mathbb R^3\setminus(I_1\times I_2\times I_3)$. There are functions $h_n$, $n\ge1$, on $\mathbb R^3$, each a finite sum of products $\phi_1(x_1)\cdot\phi_2(x_2)\cdot\phi_3(x_3)$ with $\phi_l\in C_c(I_l)$, such that
\[
\lim_{n\to\infty}\,\sup_{v\in\mathcal U}\left|q_v(f)-\int\mathrm{d}x\,\mathsf a_v(x)\cdot h_n(x)\right|=0.
\]
In particular, $q_v(f)$ is a uniform limit on $\mathcal U$ of polynomials in the integrals $\int\mathrm{d}s\,\phi(s)\cdot v_l(s)$ with $l\in\{1,2,3\}$ and $\phi\in C_c(I_l)$, and $v\mapsto q_v(f)$ is measurable on $\mathcal U$.
\end{enumerate}
By (iii) and the bounds $0\le q_v(f)\le\|f\|^2$, Lemma~\ref{lem:gaussian-mixtures} applies to the probability space $(\mathcal U,\mathbb P)$ and to the quadratic forms $q_v$, $v\in\mathcal U$. Let $m_2\in\mathcal M$ be the law that it provides; thus
\[
m_2\left(\exp\left(iX(f)\right)\right)=\mathbb E\left[\Phi_f(v)\right]\qquad(f\in\mathcal S).
\]
The law $m_2$ has the following properties.
\begin{enumerate}
\item[(1)] For $f\in\mathcal S$, $v\in\mathcal U$, $a\in\mathbb R^3$ and $w\in O_{\mathrm{cub}}$,
\[
q_v(\tau_a\cdot f)=q_{v^a}(f),\qquad q_v(w\cdot f)=q_{v^w}(f).
\]
The law $m_2$ satisfies~\eqref{eq:setup-G} and $(T_{123})$, and $w_*m_2=m_2$ for every $w\in O_{\mathrm{cub}}$.
\item[(2)] The law $m_2$ satisfies $(RP_j)$ for $j=1,2,3$. Consequently $m_2$ satisfies~\eqref{eq:setup-H}.
\item[(3)] Let $r>0$, let $f,g\in\mathcal D$ have their supports in $[-r,r]^3$, and let $a\in\mathbb R^3$ satisfy $\min_l|a_l|>2r$. Then
\[
\left|m_2\left(\exp\left(iX(f)\right)\cdot\exp\left(iX(\tau_a\cdot g)\right)\right)-m_2\left(\exp\left(iX(f)\right)\right)\cdot m_2\left(\exp\left(iX(g)\right)\right)\right|\le\sum_{l=1}^3\exp\left(-2\left(|a_l|-2r\right)\right).
\]
Moreover $m_2\in\mathcal M^{123}_{\mathrm{erg}}$.
\item[(4)] Let $r>0$, let $f\in\mathcal D$, $f\ne0$, have its support in $[-r,r]^3$, and let $\bar\Phi=\mathbb E_{v_1}\Phi_f(v)$, a function of $(v_2,v_3)$. For $s>2r$,
\[
\left|m_2\left(\exp\left(iX(f)\right)\cdot\exp\left(iX(\tau_{s\cdot e_1}\cdot f)\right)\right)-\mathbb E\left[\bar\Phi^2\right]\right|\le\exp\left(-2(s-2r)\right),\qquad m_2\left(\exp\left(iX(f)\right)\right)=\mathbb E\left[\bar\Phi\right],
\]
and
\[
\lim_{s\to\infty}m_2\left(\exp\left(iX(f)\right)\cdot\exp\left(iX(\tau_{s\cdot e_1}\cdot f)\right)\right)=\mathbb E\left[\bar\Phi^2\right]>\left(\mathbb E\left[\bar\Phi\right]\right)^2=m_2\left(\exp\left(iX(f)\right)\right)^2.
\]
\item[(5)] Let $r$ and $f$ be as in (4), and let $\tilde\Phi=\mathbb E_{v_1,v_2}\Phi_f(v)$, a function of $v_3$. Then
\[
\lim_{s\to\infty}m_2\left(\exp\left(iX(f)\right)\cdot\exp\left(iX\left(\tau_{s\cdot(e_1+e_2)}\cdot f\right)\right)\right)=\mathbb E\left[\tilde\Phi^2\right]>m_2\left(\exp\left(iX(f)\right)\right)^2.
\]
\end{enumerate}
\end{example}

\begin{proof}[Proof of the assertions of Example~\ref{ex:white-noise-telegraph}]
\emph{Assertion (i).} Let $\eta$ be right continuous, and let $x\in\mathbb R$ be a point of discontinuity of $\eta$. As $\eta$ is right continuous at $x$ and takes only the values $0$ and $1$, there is $\delta_x>0$ such that $\eta(y)=\eta(x)$ for $x\le y<x+\delta_x$; hence $\eta$ is continuous at every point of the open interval $(x,x+\delta_x)$. If $x'>x$ is another point of discontinuity of $\eta$, it follows that $x'\ge x+\delta_x$. So the intervals $(x,x+\delta_x)$ attached to the points of discontinuity $x$ of $\eta$ are pairwise disjoint; each of them contains a rational number, and therefore the set $N_\eta$ of the points of discontinuity of $\eta$ is countable. If $\eta$ is left continuous, the same reasoning applies to the intervals $(x-\delta_x,x)$. At a point of continuity $x$ with $\eta(x)=1$, the function $\eta$ is equal to $1$ on a neighbourhood of $x$. Hence the set $\{\eta=1\}$ is the union of its interior and of a subset of the countable set $N_\eta$; it is a Borel set, and $\eta$ is a Borel function. The function $\phi\cdot\eta$ is bounded, vanishes outside a bounded interval and is continuous outside the Lebesgue null set $N_\eta$; by Lebesgue's criterion for Riemann integrability it is Riemann integrable, and its Riemann sums for the partitions of $\mathbb R$ into the intervals $\left[k\cdot2^{-n},(k+1)\cdot2^{-n}\right)$, $k\in\mathbb Z$, with the left end points as intermediate points, converge to its integral as $n\to\infty$. These Riemann sums are the sums in the display. If $\phi$ vanishes outside $I$, the $n$-th sum for $\eta=\varsigma_l$ is a finite linear combination of the random variables $\varsigma_l\left(k\cdot2^{-n}\right)$ with $k\cdot2^{-n}\in I$; it is $\mathcal F^l_I$-measurable, and so is its limit $\int\mathrm{d}s\,\phi(s)\cdot\varsigma_l(s)$. As
\[
\int\mathrm{d}s\,\phi(s)\cdot v_l(s)=\frac{1}{2}\int\mathrm{d}s\,\phi(s)+\frac{1}{2}\int\mathrm{d}s\,\phi(s)\cdot\varsigma_l(s)
\]
and the $\sigma$-algebra of $\mathcal U$ is generated by the maps $v\mapsto\int\mathrm{d}s\,\phi(s)\cdot v_l(s)$, $\phi\in C_c(\mathbb R)$, $l\in\{1,2,3\}$, the random element $v$ of $\mathcal U$ is measurable. For the same reason, each entry $v_l$ is a random element of $\mathcal V$ which is measurable with respect to $\mathcal F^l_{\mathbb R}$: composed with $v_l$, the maps that generate the $\sigma$-algebra of $\mathcal V$ become the $\mathcal F^l_{\mathbb R}$-measurable random variables $\int\mathrm{d}s\,\phi(s)\cdot v_l(s)$, $\phi\in C_c(\mathbb R)$. The processes $\varsigma_1,\varsigma_2,\varsigma_3$ being independent, the $\sigma$-algebras $\mathcal F^1_{\mathbb R},\mathcal F^2_{\mathbb R},\mathcal F^3_{\mathbb R}$ are independent, and hence so are $v_1,v_2,v_3$.

Let $\eta_1,\eta_2,\eta_3$ be as in (i), and let $v^\eta$ be formed by the classes of the paths of $\frac{1}{2}(1+\eta_l)$, $l=1,2,3$. Each integral $\int\mathrm{d}s\,\phi(s)\cdot v^\eta_l(s)=\frac{1}{2}\int\mathrm{d}s\,\phi(s)+\frac{1}{2}\int\mathrm{d}s\,\phi(s)\cdot\eta_l(s)$ is the limit of the corresponding Riemann sums, which are random variables; so $v^\eta$ is a random element of $\mathcal U$. Let $k\ge1$, $\phi_1,\dots,\phi_k\in C_c(\mathbb R)$ and $l_1,\dots,l_k\in\{1,2,3\}$, and let $Y$ and $Y^\eta$ be the random vectors in $\mathbb R^k$ with the entries $\int\mathrm{d}s\,\phi_j(s)\cdot v_{l_j}(s)$, resp. $\int\mathrm{d}s\,\phi_j(s)\cdot v^\eta_{l_j}(s)$, $j=1,\dots,k$. By the first assertion, $Y$ is the pointwise limit of the random vectors $Y_n$ with the entries
\[
\frac{1}{2}\int\mathrm{d}s\,\phi_j(s)+\frac{1}{2}\cdot2^{-n}\cdot\sum_{m\in\mathbb Z}\phi_j\left(m\cdot2^{-n}\right)\cdot\varsigma_{l_j}\left(m\cdot2^{-n}\right)\qquad(j=1,\dots,k),
\]
and $Y^\eta$ is the pointwise limit of the random vectors $Y^\eta_n$ with the entries
\[
\frac{1}{2}\int\mathrm{d}s\,\phi_j(s)+\frac{1}{2}\cdot2^{-n}\cdot\sum_{m\in\mathbb Z}\phi_j\left(m\cdot2^{-n}\right)\cdot\eta_{l_j}\left(m\cdot2^{-n}\right)\qquad(j=1,\dots,k).
\]
For each $n$, the vectors $Y_n$ and $Y^\eta_n$ are the same Borel function of finitely many values of $(\varsigma_1,\varsigma_2,\varsigma_3)$, resp. of $(\eta_1,\eta_2,\eta_3)$; so they have the same law. Pointwise convergence implies convergence in law; hence $Y$ and $Y^\eta$ have the same law. The sets
\[
\left\{u\in\mathcal U:\left(\int\mathrm{d}s\,\phi_j(s)\cdot u_{l_j}(s)\right)_{j=1,\dots,k}\in B\right\},
\]
with $k$, $\phi_j$, $l_j$ as above and Borel sets $B\subset\mathbb R^k$, form a family which is closed under finite intersections and generates the $\sigma$-algebra of $\mathcal U$, and the laws of $v$ and of $v^\eta$ coincide on this family. By Dynkin's $\pi$-$\lambda$ theorem, they coincide.

Finally, each of the independent processes $\varsigma_1,\varsigma_2,\varsigma_3$ has the finite-dimensional distributions given by the defining formula of a telegraph process. Therefore the triples $(\varsigma_2,\varsigma_1,\varsigma_3)$ and $(\varsigma_3,\varsigma_2,\varsigma_1)$ have the same finite-dimensional distributions as $(\varsigma_1,\varsigma_2,\varsigma_3)$; their paths are right continuous. By the preceding paragraph, the random elements $(v_2,v_1,v_3)$ and $(v_3,v_2,v_1)$ of $\mathcal U$ have the law $\mathbb P$, which is the law of $(v_1,v_2,v_3)$. The map $\mathcal U\to\mathcal V$, $u\mapsto u_1$, is measurable; taking image laws under it, one finds that $v_2$ and $v_3$ have the same law as $v_1$.

\emph{Assertion (ii).} For $\phi\in C_c(\mathbb R)$, $v\in\mathcal U$ and $i,l\in\{1,2,3\}$, the substitutions $s\mapsto s-a_l$ and $s\mapsto\varepsilon_i\cdot s$ give
\[
\int\mathrm{d}s\,\phi(s)\cdot v^a_l(s)=\int\mathrm{d}s\,\phi(s-a_l)\cdot v_l(s),\qquad\int\mathrm{d}s\,\phi(s)\cdot v^w_i(s)=\int\mathrm{d}s\,\phi(\varepsilon_i\cdot s)\cdot v_{\pi_w(i)}(s),
\]
and the right sides are measurable functions of $v$; hence both maps are measurable. The random element $v^a$ of $\mathcal U$ is formed by the classes of the paths of $\frac{1}{2}(1+\eta_l)$ with $\eta_l=\varsigma_l(\cdot+a_l)$, $l=1,2,3$. By~\eqref{eq:ergodic-components-4}, each $\eta_l$ has the finite-dimensional distributions of a telegraph process, and $\eta_1,\eta_2,\eta_3$ are independent; so $(\eta_1,\eta_2,\eta_3)$ has the same finite-dimensional distributions as $(\varsigma_1,\varsigma_2,\varsigma_3)$. Every path of each $\eta_l$ is right continuous, and (i) shows that $v^a$ has the law $\mathbb P$. The random element $v^w$ is formed by the classes of the paths of $\frac{1}{2}(1+\eta_i)$ with the processes $\eta_i(s)=\varsigma_{\pi_w(i)}(\varepsilon_i\cdot s)$, $i=1,2,3$. Each $\eta_i$ is $\varsigma_{\pi_w(i)}$ or the reversed process $\check\varsigma_{\pi_w(i)}$, and it has the finite-dimensional distributions of a telegraph process by~\eqref{eq:ergodic-components-4}; the processes $\eta_1,\eta_2,\eta_3$ are independent, since $\pi_w$ is a bijection; and every path of $\eta_i$ is right continuous if $\varepsilon_i=1$ and left continuous if $\varepsilon_i=-1$. By (i), $v^w$ has the law $\mathbb P$.

\emph{Assertion (iii).} The function $x\mapsto|f(x)|^2$ is integrable and vanishes outside the open box $B=I_1\times I_2\times I_3$. The finite sums of products $\phi_1(x_1)\cdot\phi_2(x_2)\cdot\phi_3(x_3)$ with $\phi_l\in C_c(I_l)$ are dense in $L^1(B)$. Indeed, the continuous functions with compact support in $B$ are dense in $L^1(B)$. For such a function $g$ choose compact intervals $K_l\subset I_l$ with $\operatorname{supp}g\subset K_1\times K_2\times K_3$, and functions $\chi_l\in C_c(I_l)$ with $0\le\chi_l\le1$ and $\chi_l=1$ on $K_l$, and let $L$ be the product of the supports of $\chi_1,\chi_2,\chi_3$. By the Weierstrass approximation theorem there are polynomials $p_n$ on $\mathbb R^3$ with $\sup_L|p_n-g|\to0$. The functions $\chi_1(x_1)\cdot\chi_2(x_2)\cdot\chi_3(x_3)\cdot p_n(x)$ are finite sums of products of the required kind, vanish outside the compact set $L\subset B$, and differ from $g$ on $L$ by at most $\sup_L|p_n-g|$; so they converge to $g$ in $L^1(B)$. Choose finite sums $h_n$ of products of this kind with $\int\mathrm{d}x\,\left||f(x)|^2-h_n(x)\right|\to0$. Since $0\le\mathsf a_v\le1$,
\[
\left|q_v(f)-\int\mathrm{d}x\,\mathsf a_v(x)\cdot h_n(x)\right|\le\int\mathrm{d}x\,\left||f(x)|^2-h_n(x)\right|\qquad(v\in\mathcal U),
\]
which proves the uniform convergence. If $h_n(x)=\sum_k\phi^k_1(x_1)\cdot\phi^k_2(x_2)\cdot\phi^k_3(x_3)$, Fubini's theorem gives
\[
\int\mathrm{d}x\,\mathsf a_v(x)\cdot h_n(x)=\sum_k\prod_{l=1}^3\int\mathrm{d}s\,\phi^k_l(s)\cdot v_l(s),
\]
a polynomial in integrals of the stated kind; it is a measurable function of $v$, and so is the uniform limit $q_v(f)$.

\emph{Assertion (1).} Let $f\in\mathcal S$ and $v\in\mathcal U$. For $a\in\mathbb R^3$, the substitution $x=y+a$ and the identity $\mathsf a_v(y+a)=\mathsf a_{v^a}(y)$ give
\[
q_v(\tau_a\cdot f)=\int\mathrm{d}x\,\mathsf a_v(x)\cdot|f(x-a)|^2=\int\mathrm{d}y\,\mathsf a_v(y+a)\cdot|f(y)|^2=q_{v^a}(f).
\]
Let $w\in O_{\mathrm{cub}}$ with $w\,e_i=\varepsilon_i\cdot e_{\pi_w(i)}$. Since $\wedge^2w$ is orthogonal, $|(w\cdot f)(x)|=\left|f\left(w^{-1}x\right)\right|$, and the substitution $x=wy$ gives
\[
q_v(w\cdot f)=\int\mathrm{d}y\,\mathsf a_v(wy)\cdot|f(y)|^2.
\]
Now $wy=\sum_iy_i\cdot\varepsilon_i\cdot e_{\pi_w(i)}$, that is, $(wy)_{\pi_w(i)}=\varepsilon_i\cdot y_i$; hence
\[
\mathsf a_v(wy)=\prod_{i=1}^3v_{\pi_w(i)}\left((wy)_{\pi_w(i)}\right)=\prod_{i=1}^3v_{\pi_w(i)}(\varepsilon_i\cdot y_i)=\prod_{i=1}^3v^w_i(y_i)=\mathsf a_{v^w}(y),
\]
and $q_v(w\cdot f)=q_{v^w}(f)$. The law $m_2$ satisfies~\eqref{eq:setup-G} by Lemma~\ref{lem:gaussian-mixtures}. By (ii), the maps $v\mapsto v^a$ and $v\mapsto v^w$ preserve $\mathbb P$; so the last assertion of Lemma~\ref{lem:gaussian-mixtures}, applied to the Euclidean maps $\tau_a$ and $w$ together with these maps of $\mathcal U$, gives $(\tau_a)_*m_2=m_2$ for every $a\in\mathbb R^3$, which is $(T_{123})$, and $w_*m_2=m_2$ for every $w\in O_{\mathrm{cub}}$.

\emph{Assertion (2).} Let $n\ge1$, $c\in\mathbb C^n$ and $f_1,\dots,f_n\in\mathcal D(\{x_1>0\})$. For $p,q\in\{1,\dots,n\}$ the support of $f_q$ lies in $\{x_1>0\}$ and the support of $\theta_1\cdot f_p$ in $\{x_1<0\}$, so that $(f_q)_{il}(x)\cdot(\theta_1\cdot f_p)_{il}(x)=0$ for all $x\in\mathbb R^3$ and $i<l$. Hence, by (1) for $w=\theta_1$ and $v^{\theta_1}=(\check v_1,v_2,v_3)$,
\[
q_v\left(f_q-\theta_1\cdot f_p\right)=q_v(f_q)+q_v(\theta_1\cdot f_p)=q_v(f_q)+q_{(\check v_1,v_2,v_3)}(f_p)\qquad(v\in\mathcal U),
\]
and therefore, as $X(f_q)-X(\theta_1\cdot f_p)=X\left(f_q-\theta_1\cdot f_p\right)$ and $\Phi_{f_p}$ is real,
\[
\sum_{p,q=1}^n\overline{c_p}\cdot c_q\cdot m_2\left(\exp\left(i\left(X(f_q)-X(\theta_1\cdot f_p)\right)\right)\right)=\mathbb E\left[\overline{\sum_{p=1}^nc_p\cdot\Phi_{f_p}(\check v_1,v_2,v_3)}\cdot\sum_{q=1}^nc_q\cdot\Phi_{f_q}(v_1,v_2,v_3)\right].
\]
For fixed $(v_2,v_3)\in\mathcal V^2$ let $G(z)=\sum_{p=1}^nc_p\cdot\Phi_{f_p}(z,v_2,v_3)$ for $z\in\mathcal V$. By the independence of $v_1$ and $(v_2,v_3)$ and Fubini's theorem, the right side equals $\mathbb E_{v_2,v_3}\mathbb E_{v_1}\left[\overline{G(\check v_1)}\cdot G(v_1)\right]$. Each $f_p$ vanishes outside $(0,\infty)\times\mathbb R\times\mathbb R$; by (iii), applied with $I_1=(0,\infty)$ and $I_2=I_3=\mathbb R$, the function $G$ is a bounded measurable function of the integrals $\int\mathrm{d}s\,\phi(s)\cdot z(s)$, $\phi\in C_c((0,\infty))$. As $\int\mathrm{d}s\,\phi(s)\cdot v_1(s)=\frac{1}{2}\int\mathrm{d}s\,\phi(s)+\frac{1}{2}\int\mathrm{d}s\,\phi(s)\cdot\varsigma_1(s)$ and $\int\mathrm{d}s\,\phi(s)\cdot\check v_1(s)=\frac{1}{2}\int\mathrm{d}s\,\phi(s)+\frac{1}{2}\int\mathrm{d}s\,\phi(s)\cdot\check\varsigma_1(s)$, the random variable $G(v_1)$ is a bounded measurable function of the integrals $\int\mathrm{d}s\,\phi(s)\cdot\varsigma_1(s)$, $\phi\in C_c((0,\infty))$, and $G(\check v_1)$ is the same function of the integrals $\int\mathrm{d}s\,\phi(s)\cdot\check\varsigma_1(s)$. Property~\eqref{eq:ergodic-components-7} of the telegraph process $\varsigma_1$ gives $\mathbb E_{v_1}\left[\overline{G(\check v_1)}\cdot G(v_1)\right]\ge0$ for every value of $(v_2,v_3)$. Hence the sum above is non-negative, and $m_2$ satisfies $(RP_1)$.

Let $j\in\{2,3\}$, and let $w\in O_{\mathrm{cub}}$ be the matrix that interchanges $e_1$ and $e_j$ and fixes the third vector of the standard basis. Then $w^{-1}=w$, $w\{x_1>0\}=\{x_j>0\}$ and $w\theta_1w^{-1}=\theta_j$. Let $n\ge1$, $c\in\mathbb C^n$ and $g_1,\dots,g_n\in\mathcal D(\{x_j>0\})$, and put $f_p=w^{-1}\cdot g_p$; then $f_p\in\mathcal D(\{x_1>0\})$, $g_p=w\cdot f_p$ and $\theta_j\cdot g_p=w\cdot(\theta_1\cdot f_p)$. By Section~\ref{sec:euclidean-maps-actions},
\[
\exp\left(i\left(X(g_q)-X(\theta_j\cdot g_p)\right)\right)=\hat w\exp\left(i\left(X(f_q)-X(\theta_1\cdot f_p)\right)\right),
\]
and $m_2(\hat wF)=m_2(F)$ for every bounded measurable function $F$ on $Q$, because $w_*m_2=m_2$ by (1). Hence the sum in $(RP_j)$ formed with $c$ and $g_1,\dots,g_n$ equals the sum in $(RP_1)$ formed with $c$ and $f_1,\dots,f_n$, which is non-negative. Thus $m_2$ satisfies $(RP_j)$ for $j=1,2,3$; together with (1), this shows that $m_2$ satisfies~\eqref{eq:setup-H}.

\emph{Two auxiliary facts.} Let $\varsigma'_1,\varsigma'_2,\varsigma'_3$ be telegraph processes such that the six processes $\varsigma_1,\varsigma_2,\varsigma_3,\varsigma'_1,\varsigma'_2,\varsigma'_3$ are independent (the underlying probability space being enlarged if necessary), and let $v'=(v'_1,v'_2,v'_3)$ be the random element of $\mathcal U$ whose entries are the classes of the paths of $\frac{1}{2}(1+\varsigma'_l)$, $l=1,2,3$. The paths of $\varsigma'_1,\varsigma'_2,\varsigma'_3$ are right continuous, and $(\varsigma'_1,\varsigma'_2,\varsigma'_3)$ has the same finite-dimensional distributions as $(\varsigma_1,\varsigma_2,\varsigma_3)$; so $v'$ has the law $\mathbb P$, by (i), and $v'$ is independent of $v$.

($\alpha$) Let $r>0$, let $f,g\in\mathcal D$ have their supports in $[-r,r]^3$, and let $a\in\mathbb R^3$ satisfy $|a_k|>2r$ for at least one $k\in\{1,2,3\}$. Then
\[
m_2\left(\exp\left(iX(f)\right)\cdot\exp\left(iX(\tau_a\cdot g)\right)\right)=\mathbb E\left[\Phi_f(v)\cdot\Phi_g(v^a)\right].
\]
Indeed, every point $x$ of the support of $f$ satisfies $|x_k|\le r$, and every point $x$ of the support of $\tau_a\cdot g$ satisfies $|x_k-a_k|\le r$; so these supports are disjoint, $f_{il}(x)\cdot(\tau_a\cdot g)_{il}(x)=0$ for all $x$ and $i<l$, and by (1) $q_v(f+\tau_a\cdot g)=q_v(f)+q_v(\tau_a\cdot g)=q_v(f)+q_{v^a}(g)$ for $v\in\mathcal U$. Since $\exp\left(iX(f)\right)\cdot\exp\left(iX(\tau_a\cdot g)\right)=\exp\left(iX(f+\tau_a\cdot g)\right)$, the identity follows from the formula defining $m_2$.

($\beta$) Let $r$, $f$, $g$, $a$ be as in ($\alpha$), and let $l\in\{1,2,3\}$ satisfy $|a_l|>2r$. Let $y$ and $y'$ be random elements of $\mathcal U$ such that $y_l=v_l$, $y'_l=v'_l$ and, for each $k\ne l$, $y_k=y'_k$, where $y_k$ is either $v_k$ or $v'_k$. Then
\[
\left|\mathbb E\left[\Phi_f(v)\cdot\Phi_g\left(y^a\right)\right]-\mathbb E\left[\Phi_f(v)\cdot\Phi_g\left(y'^a\right)\right]\right|\le\exp\left(-2\left(|a_l|-2r\right)\right).
\]
To prove this, note first that $f$ and $g$ vanish outside $(-r,r)^3$, because they are continuous and their supports lie in $[-r,r]^3$. By (iii), applied with $I_1=I_2=I_3=(-r,r)$, the number $\Phi_f(u)$, $u\in\mathcal U$, is a measurable function of the integrals $\int\mathrm{d}s\,\phi(s)\cdot u_k(s)$ with $\phi\in C_c((-r,r))$ and $k\in\{1,2,3\}$, and so is $\Phi_g(u)$. Since $\int\mathrm{d}s\,\phi(s)\cdot u_k(s+a_k)=\int\mathrm{d}s\,\phi(s-a_k)\cdot u_k(s)$, the number $\Phi_g(u^a)$ is a measurable function of the integrals $\int\mathrm{d}s\,\psi(s)\cdot u_k(s)$ with $\psi\in C_c((a_k-r,a_k+r))$ and $k\in\{1,2,3\}$.

The pair $(v_l,v'_l)$ is independent of the family $Z$ of the four random elements $v_k$, $v'_k$, $k\ne l$, of $\mathcal V$, and $v_l$, $v'_l$ are independent, both with the law $\mathbb P_1$. Fix a value of $Z$. For $z\in\mathcal V$ let $H_1(z)$ be the value of $\Phi_f$ at the triple obtained from $v$ by replacing its entry $v_l$ by $z$, and let $H_2(z)$ be the value of $\Phi_g$ at $u^a$, where $u$ is the triple obtained from $y$ by replacing its entry $y_l$ by $z$. Then $\Phi_f(v)\cdot\Phi_g(y^a)=H_1(v_l)\cdot H_2(v_l)$ and $\Phi_f(v)\cdot\Phi_g(y'^a)=H_1(v_l)\cdot H_2(v'_l)$, and $\mathbb E\left[H_1(v_l)\cdot H_2(v'_l)\right]=\mathbb E\left[H_1(v_l)\right]\cdot\mathbb E\left[H_2(v_l)\right]$. By Fubini's theorem it therefore suffices to show that
\[
\left|\mathbb E\left[H_1(v_l)\cdot H_2(v_l)\right]-\mathbb E\left[H_1(v_l)\right]\cdot\mathbb E\left[H_2(v_l)\right]\right|\le\exp\left(-2\left(|a_l|-2r\right)\right)
\]
for every value of $Z$. By the preceding paragraph and (i), the random variable $H_1(v_l)$ is $\mathcal F^l_{(-r,r)}$-measurable and $H_2(v_l)$ is $\mathcal F^l_{(a_l-r,a_l+r)}$-measurable, and both take their values in $(0,1]$. Put $t=|a_l|-2r>0$. If $a_l>2r$, these $\sigma$-algebras are contained in $\mathcal F^l_{(-\infty,r]}$ and in $\mathcal F^l_{[r+t,\infty)}$, respectively; if $a_l<-2r$, they are contained in $\mathcal F^l_{[r'+t,\infty)}$ and in $\mathcal F^l_{(-\infty,r']}$, respectively, where $r'=a_l+r$. In both cases property~\eqref{eq:ergodic-components-5} of the telegraph process $\varsigma_l$, applied at the point $r$, resp. $r'$, gives the inequality, the left side being symmetric in $H_1$ and $H_2$.

\emph{Assertion (3).} For $k=0,1,2,3$ let $v^{(k)}$ be the random element of $\mathcal U$ with the entries $v^{(k)}_l=v'_l$ for $l\le k$ and $v^{(k)}_l=v_l$ for $l>k$; thus $v^{(0)}=v$ and $v^{(3)}=v'$. By ($\alpha$),
\[
m_2\left(\exp\left(iX(f)\right)\cdot\exp\left(iX(\tau_a\cdot g)\right)\right)=\mathbb E\left[\Phi_f(v)\cdot\Phi_g\left(\left(v^{(0)}\right)^a\right)\right].
\]
For $k=1,2,3$, ($\beta$), applied with $l=k$, $y=v^{(k-1)}$ and $y'=v^{(k)}$, gives
\[
\left|\mathbb E\left[\Phi_f(v)\cdot\Phi_g\left(\left(v^{(k-1)}\right)^a\right)\right]-\mathbb E\left[\Phi_f(v)\cdot\Phi_g\left(\left(v^{(k)}\right)^a\right)\right]\right|\le\exp\left(-2\left(|a_k|-2r\right)\right).
\]
Finally, $v'$ is independent of $v$, and $v'^a$ has the law $\mathbb P$ by (ii); hence
\[
\mathbb E\left[\Phi_f(v)\cdot\Phi_g\left(v'^a\right)\right]=\mathbb E\left[\Phi_f(v)\right]\cdot\mathbb E\left[\Phi_g(v)\right]=m_2\left(\exp\left(iX(f)\right)\right)\cdot m_2\left(\exp\left(iX(g)\right)\right).
\]
The triangle inequality gives the bound in (3).

Let $F=\sum_pc_p\cdot\exp\left(iX(f_p)\right)$ and $G=\sum_qc'_q\cdot\exp\left(iX(g_q)\right)$ be elements of $\mathcal A$, with finitely many $f_p,g_q\in\mathcal D_0$ and $c_p,c'_q\in\mathbb C$, and choose $r>0$ such that the supports of all $f_p$ and $g_q$ lie in $[-r,r]^3$. For $h\in\mathcal S$ one has $q_v(-h)=q_v(h)$, so that $m_2\left(\exp\left(iX(-h)\right)\right)=m_2\left(\exp\left(iX(h)\right)\right)$, and this number is real. Let $a^{(k)}=k\cdot(e_1+e_2+e_3)$ for $k\ge1$. By Section~\ref{sec:euclidean-maps-actions}, $\overline{F}=\sum_p\overline{c_p}\cdot\exp\left(iX(-f_p)\right)$ and $U\left(a^{(k)}\right)G=\sum_qc'_q\cdot\exp\left(iX\left(\tau_{a^{(k)}}\cdot g_q\right)\right)$. The bound in (3), applied to the pairs $(-f_p,g_q)$, gives for $k>2r$
\[
\left|m_2\left(\overline{F}\cdot U\left(a^{(k)}\right)G\right)-\overline{m_2(F)}\cdot m_2(G)\right|\le3\cdot\exp\left(-2(k-2r)\right)\cdot\sum_{p,q}|c_p|\cdot|c'_q|,
\]
and the right side tends to $0$ as $k\to\infty$. By (1), $m_2$ satisfies $(T_{123})$; Lemma~\ref{lem:clustering-criterion-separation}(a), applied with $J=\{1,2,3\}$, gives $m_2\in\mathcal M^{123}_{\mathrm{erg}}$.

\emph{Assertion (4).} Let $s>2r$. By ($\alpha$), and by ($\beta$) applied with $g=f$, $a=s\cdot e_1$, $l=1$, $y=v$ and $y'=(v'_1,v_2,v_3)$, for which $y'^a=\left(v'_1(\cdot+s),v_2,v_3\right)$,
\[
\left|m_2\left(\exp\left(iX(f)\right)\cdot\exp\left(iX(\tau_{s\cdot e_1}\cdot f)\right)\right)-\mathbb E\left[\Phi_f(v)\cdot\Phi_f\left(v'_1(\cdot+s),v_2,v_3\right)\right]\right|\le\exp\left(-2(s-2r)\right).
\]
The random element $v'_1(\cdot+s)$ of $\mathcal V$ is independent of $v$, and it has the law $\mathbb P_1$, by (ii) applied to $v'$ and to $a=s\cdot e_1$. So $v_1$ and $v'_1(\cdot+s)$ are independent, independent of $(v_2,v_3)$, and both have the law $\mathbb P_1$; by Fubini's theorem,
\[
\mathbb E\left[\Phi_f(v)\cdot\Phi_f\left(v'_1(\cdot+s),v_2,v_3\right)\right]=\mathbb E\left[\int_{\mathcal V}\mathbb P_1(\mathrm{d}z)\,\Phi_f(z,v_2,v_3)\cdot\int_{\mathcal V}\mathbb P_1(\mathrm{d}z')\,\Phi_f(z',v_2,v_3)\right]=\mathbb E\left[\bar\Phi^2\right].
\]
Moreover $m_2\left(\exp\left(iX(f)\right)\right)=\mathbb E\left[\Phi_f(v)\right]=\mathbb E\left[\bar\Phi\right]$. This proves the first two assertions of (4), and it follows that the limit as $s\to\infty$ exists and equals $\mathbb E\left[\bar\Phi^2\right]$.

For $z\in\mathcal V$ put $J_1(z)=\int\mathrm{d}x\,z(x_1)\cdot|f(x)|^2$; then $J_1(z)\ge\frac{1}{2}\|f\|^2>0$, as $z\ge\frac{1}{2}$ and $f\ne0$. Let $A_1$, resp. $A_0$, be the event that $\varsigma_2(t)=\varsigma_3(t)=1$, resp. $\varsigma_2(t)=\varsigma_3(t)=0$, for all $t\in[-r,r]$. By~\eqref{eq:ergodic-components-6} and the independence of $\varsigma_2$ and $\varsigma_3$, each of these two events has the probability $\frac{1}{4}\exp(-4r)>0$. On $A_1$ one has $v_2=v_3=1$ on $[-r,r]$; as $f$ vanishes outside $[-r,r]^3$, it follows that $\mathsf a_v(x)\cdot|f(x)|^2=v_1(x_1)\cdot|f(x)|^2$ for almost every $x$, so that $q_v(f)=J_1(v_1)$ and
\[
\bar\Phi=\int_{\mathcal V}\mathbb P_1(\mathrm{d}z)\,\exp\left(-\frac{1}{2}J_1(z)\right)=:\beta_1\qquad\text{on }A_1.
\]
On $A_0$ one has $v_2=v_3=\frac{1}{2}$ on $[-r,r]$, hence $q_v(f)=\frac{1}{4}J_1(v_1)$ and
\[
\bar\Phi=\int_{\mathcal V}\mathbb P_1(\mathrm{d}z)\,\exp\left(-\frac{1}{8}J_1(z)\right)=:\beta_0\qquad\text{on }A_0.
\]
Since $J_1(z)>0$, one has $\exp\left(-\frac{1}{8}J_1(z)\right)>\exp\left(-\frac{1}{2}J_1(z)\right)$ for every $z\in\mathcal V$, and therefore $\beta_0>\beta_1$. Thus the random variable $\bar\Phi$ takes the two different values $\beta_1$ and $\beta_0$ on events of positive probability; it is not almost surely constant, and
\[
\mathbb E\left[\bar\Phi^2\right]-\left(\mathbb E\left[\bar\Phi\right]\right)^2=\mathbb E\left[\left(\bar\Phi-\mathbb E\left[\bar\Phi\right]\right)^2\right]>0.
\]
This completes the proof of (4).

\emph{Assertion (5).} Let $s>2r$ and $a=s\cdot(e_1+e_2)$, so that $|a_1|=|a_2|=s>2r$. By ($\alpha$), and by ($\beta$) applied twice with $g=f$, first with $l=1$, $y=v$, $y'=(v'_1,v_2,v_3)$, then with $l=2$, $y=(v'_1,v_2,v_3)$, $y'=(v'_1,v'_2,v_3)$,
\[
\left|m_2\left(\exp\left(iX(f)\right)\cdot\exp\left(iX(\tau_a\cdot f)\right)\right)-\mathbb E\left[\Phi_f(v)\cdot\Phi_f\left(v'_1(\cdot+s),v'_2(\cdot+s),v_3\right)\right]\right|\le2\cdot\exp\left(-2(s-2r)\right).
\]
By (ii), applied to $v'$ and to $a$, the random elements $v'_1(\cdot+s)$ and $v'_2(\cdot+s)$ of $\mathcal V$ are independent and have the law $\mathbb P_1$; they are independent of $v$. So $v_1$, $v_2$, $v'_1(\cdot+s)$, $v'_2(\cdot+s)$ and $v_3$ are independent, and the first four of them have the law $\mathbb P_1$. By Fubini's theorem,
\[
\mathbb E\left[\Phi_f(v)\cdot\Phi_f\left(v'_1(\cdot+s),v'_2(\cdot+s),v_3\right)\right]=\mathbb E\left[\tilde\Phi(v_3)^2\right],
\]
and $m_2\left(\exp\left(iX(f)\right)\right)=\mathbb E\left[\Phi_f(v)\right]=\mathbb E\left[\tilde\Phi\right]$. Hence the limit in (5) exists and equals $\mathbb E\left[\tilde\Phi^2\right]$. For $z_1,z_2\in\mathcal V$ put $J_2(z_1,z_2)=\int\mathrm{d}x\,z_1(x_1)\cdot z_2(x_2)\cdot|f(x)|^2$; then $J_2(z_1,z_2)\ge\frac{1}{4}\|f\|^2>0$. Let $A'_1$, resp. $A'_0$, be the event that $\varsigma_3(t)=1$, resp. $\varsigma_3(t)=0$, for all $t\in[-r,r]$; by~\eqref{eq:ergodic-components-6}, each of them has the probability $\frac{1}{2}\exp(-2r)>0$. On $A'_1$ one has $v_3=1$ on $[-r,r]$, hence $q_v(f)=J_2(v_1,v_2)$ and
\[
\tilde\Phi=\int_{\mathcal V}\mathbb P_1(\mathrm{d}z_1)\int_{\mathcal V}\mathbb P_1(\mathrm{d}z_2)\,\exp\left(-\frac{1}{2}J_2(z_1,z_2)\right)=:\beta'_1\qquad\text{on }A'_1;
\]
on $A'_0$ one has $v_3=\frac{1}{2}$ on $[-r,r]$, hence $q_v(f)=\frac{1}{2}J_2(v_1,v_2)$ and
\[
\tilde\Phi=\int_{\mathcal V}\mathbb P_1(\mathrm{d}z_1)\int_{\mathcal V}\mathbb P_1(\mathrm{d}z_2)\,\exp\left(-\frac{1}{4}J_2(z_1,z_2)\right)=:\beta'_0\qquad\text{on }A'_0.
\]
At every point $(z_1,z_2)\in\mathcal V^2$ the integrand of $\beta'_0$ exceeds that of $\beta'_1$, because $J_2(z_1,z_2)>0$; hence $\beta'_0>\beta'_1$. The random variable $\tilde\Phi$ therefore takes the two different values $\beta'_1$ and $\beta'_0$ on events of positive probability; it is not almost surely constant, and
\[
\mathbb E\left[\tilde\Phi^2\right]-\left(\mathbb E\left[\tilde\Phi\right]\right)^2=\mathbb E\left[\left(\tilde\Phi-\mathbb E\left[\tilde\Phi\right]\right)^2\right]>0,
\]
that is, $\mathbb E\left[\tilde\Phi^2\right]>\left(\mathbb E\left[\tilde\Phi\right]\right)^2=m_2\left(\exp\left(iX(f)\right)\right)^2$.
\end{proof}

\begin{remark}[Clustering when two coordinates diverge]
By Example~\ref{ex:white-noise-telegraph}(2) and~(3), the law $m_2$ satisfies~\eqref{eq:setup-H}, so that it belongs to $\mathcal C$ (Section~\ref{sec:hypothesis-law}), and it is $\Gamma_{123}$-ergodic; thus Theorem~\ref{thm:invariant-vectors-clustering}(iii) applies to $\rho=m_2$. Let $f\in\mathcal D$, $f\ne0$, choose $r>0$ such that the support of $f$ lies in $[-r,r]^3$, and put $F=\exp\left(-iX(f)\right)$ and $G=\exp\left(iX(f)\right)$. Then $\overline{F}=G$, $U(a)G=\exp\left(iX(\tau_a\cdot f)\right)$ for $a\in\mathbb R^3$ (Section~\ref{sec:euclidean-maps-actions}), and $\overline{m_2(F)}\cdot m_2(G)=m_2\left(\exp\left(iX(f)\right)\right)^2$, because $m_2(F)=m_2(G)$ is real. By Example~\ref{ex:white-noise-telegraph}(5),
\[
\lim_{s\to\infty}m_2\left(\overline{F}\cdot U\left(s\cdot(e_1+e_2)\right)G\right)>\overline{m_2(F)}\cdot m_2(G),
\]
although the first two coordinates of $s\cdot(e_1+e_2)$ tend to infinity. Hence, in Theorem~\ref{thm:invariant-vectors-clustering}(iii)1, the condition $\min_l|a_l|\to\infty$ cannot be replaced by the condition that $|a_l|\to\infty$ for two of the three indices $l$.
\end{remark}

\subsubsection{The free Maxwell law and coordinatewise changes of variables}\label{sec:free-maxwell-law}

\paragraph{Two-forms as vector fields.}
Let $\ast\colon\Lambda^2\to\mathbb R^3$ be the isometry of Section~\ref{sec:closed-distributions}; thus $(\ast F)_k=F_{il}$ for $F\in\Lambda^2$ whenever $(k,i,l)$ is a cyclic permutation of $(1,2,3)$. For $R\in O(3)$ the linear map $\ast\circ(\wedge^2R)\circ\ast^{-1}$ of $\mathbb R^3$ is given by the cofactor matrix of $R$, and this matrix equals $(\det R)\cdot R$. Consequently, for every Euclidean map $g(x)=Rx+a$ and every map $f\colon\mathbb R^3\to\Lambda^2$,
\[
\ast(g\cdot f)(x)=(\det R)\cdot R\,(\ast f)\left(g^{-1}x\right)\qquad\left(x\in\mathbb R^3\right).
\]
From now on a map $f\colon\mathbb R^3\to\Lambda^2$ is identified with the vector field $b=\ast f\colon\mathbb R^3\to\mathbb R^3$. In this way $L^2(\mathbb R^3;\Lambda^2)$ is identified with the real Hilbert space $L^2(\mathbb R^3;\mathbb R^3)$, whose scalar product and norm
\[
\langle b',b\rangle=\int\mathrm{d}x\,\sum_{k=1}^3b'_k(x)\cdot b_k(x),\qquad\|b\|=\langle b,b\rangle^{\frac{1}{2}},
\]
correspond to those of Section~\ref{sec:test-forms-distributions}, because $\ast$ is an isometry; the space $\mathcal D$ is identified with $C_c^\infty(\mathbb R^3;\mathbb R^3)$, and, for an open set $O\subset\mathbb R^3$, the space $\mathcal D(O)$ with the space of the smooth vector fields with compact support in $O$. A Euclidean map $g(x)=Rx+a$ acts on vector fields by
\[
(g\cdot b)(x)=(\det R)\cdot R\,b\left(g^{-1}x\right)\qquad\left(x\in\mathbb R^3\right),
\]
and operators on vector fields act on forms through the identification. Let $w\in O_{\mathrm{cub}}$, and let $\varepsilon\in\{\pm1\}^3$ and $\pi\in S_3$ be given by $w\,e_i=\varepsilon_i\cdot e_{\pi(i)}$, $i=1,2,3$, as in Section~\ref{sec:euclidean-maps-actions}. The matrix $w$ maps a vector with the components $b_1,b_2,b_3$ to the vector whose component with the index $\pi(k)$ is $\varepsilon_k\cdot b_k$, for $k=1,2,3$; therefore
\[
(w\cdot b)_{\pi(k)}(x)=(\det w)\cdot\varepsilon_k\cdot b_k\left(w^{-1}x\right)\qquad\left(k=1,2,3,\ x\in\mathbb R^3\right).
\]
With these conventions, a distribution $\omega\in Q$ is closed if and only if $\omega(f)=0$ for every $f\in\mathcal S$ whose vector field $\ast f$ is the gradient $\nabla\mathsf g$ of a function $\mathsf g\in\mathcal S(\mathbb R^3)$ (Section~\ref{sec:closed-distributions}).

\paragraph{The transversal projection.}
The Fourier transform $\hat b$ of $b\in L^2(\mathbb R^3;\mathbb R^3)$ is taken componentwise and normalised by $\hat b(p)=\int\mathrm{d}x\,\exp(-ip\cdot x)\cdot b(x)$ for integrable $b$. For $p\in\mathbb R^3\setminus\{0\}$ the real symmetric matrix $1-|p|^{-2}\cdot p\,p^T$, where $1$ denotes the identity matrix, is the orthogonal projection of $\mathbb R^3$ onto the plane orthogonal to $p$, and it is an even function of $p$. Let $\Pi$ be the operator on $L^2(\mathbb R^3;\mathbb R^3)$ given by
\[
\widehat{\Pi b}(p)=\left(1-|p|^{-2}\cdot p\,p^T\right)\hat b(p)\qquad\left(p\in\mathbb R^3\setminus\{0\}\right).
\]
It is an orthogonal projection of $L^2(\mathbb R^3;\mathbb R^3)$, and it has the following three properties.
\begin{itemize}
\item One has $\Pi(g\cdot b)=g\cdot\Pi b$ for every Euclidean map $g$ and every $b\in L^2(\mathbb R^3;\mathbb R^3)$. Indeed, a translation multiplies Fourier transforms by a scalar function, and for $R\in O(3)$
\[
R^{-1}\left(1-|p|^{-2}\cdot p\,p^T\right)R=1-\left|R^{-1}p\right|^{-2}\cdot\left(R^{-1}p\right)\left(R^{-1}p\right)^T.
\]
\item One has $\Pi\nabla\psi=0$ for every $\psi\in H^1(\mathbb R^3)$, since the Fourier transform of $\nabla\psi$ is the function $p\mapsto i\,\hat\psi(p)\cdot p$.
\item Let $G\colon\mathbb R^3\setminus\{0\}\to(0,\infty)$ be the function $G(y)=\frac{1}{4\pi\cdot|y|}$, and for $h\in C_c^\infty(\mathbb R^3)$ let $(-\Delta)^{-1}h$ be the smooth function
\[
(-\Delta)^{-1}h(y)=\int\mathrm{d}y'\,G(y-y')\cdot h(y')\qquad\left(y\in\mathbb R^3\right).
\]
Then, for all $b,b'\in\mathcal D$,
\[
\langle b',\Pi b\rangle=\langle b',b\rangle-\left\langle\operatorname{div}b',(-\Delta)^{-1}\operatorname{div}b\right\rangle,
\]
where the last scalar product is $\int\mathrm{d}y\,\operatorname{div}b'(y)\cdot(-\Delta)^{-1}\operatorname{div}b(y)$. Indeed, the Fourier transform of $\operatorname{div}b$ is the function $p\mapsto i\,p\cdot\hat b(p)$, and the Fourier transform of $G$ is the function $p\mapsto|p|^{-2}$.
\end{itemize}

\begin{lemma}[Properties of the free Maxwell law]\label{lem:properties-free-maxwell}
For $f,h\in\mathcal S$ put $q(f,h)=\langle\Pi f,\Pi h\rangle$ and $q(f)=q(f,f)=\|\Pi f\|^2$, that is, $q(f)=\|\Pi(\ast f)\|^2$. As $\Pi$ is an orthogonal projection, $q$ is a symmetric bilinear form on $\mathcal S$ with $0\le q(f)\le\|f\|^2$. Let $\gamma\in\mathcal M$ be the law which Lemma~\ref{lem:gaussian-mixtures} provides for a probability space consisting of one point and for the form $q$; thus $\gamma$ is the unique law with
\[
\gamma\left(\exp\left(iX(f)\right)\right)=\exp\left(-\frac{1}{2}\cdot\|\Pi f\|^2\right)\qquad\left(f\in\mathcal S\right).
\]
The law $\gamma$ is called the free Maxwell law.
\begin{enumerate}
\item[(i)] The law $\gamma$ satisfies~\eqref{eq:setup-G}; one has $g_*\gamma=\gamma$ for every Euclidean map $g$; the law $\gamma$ is concentrated on the closed distributions, that is, $\gamma\left(\{\omega\in Q:d\omega=0\}\right)=1$; and $\gamma$ satisfies $(RP_j)$ for $j=1,2,3$.
\item[(ii)] For a measurable vector field $b$ on $\mathbb R^3$ let
\[
\|b\|_{L^1}=\int\mathrm{d}y\,|b(y)|\in[0,\infty],
\]
where $|b(y)|$ is the Euclidean norm of $b(y)\in\mathbb R^3$. Let $\delta>0$, and let $b',b''\in L^2(\mathbb R^3;\mathbb R^3)$ have compact supports such that $|y-y'|\ge\delta$ for all $y\in\operatorname{supp}b'$ and $y'\in\operatorname{supp}b''$. Then
\[
\left|\langle b',\Pi b''\rangle\right|\le\frac{1}{2\pi\cdot\delta^3}\cdot\|b'\|_{L^1}\cdot\|b''\|_{L^1}.
\]
\end{enumerate}
\end{lemma}

\begin{proof}
(ii) Assume first that $b',b''\in\mathcal D$. Their supports are disjoint; hence $\langle b',b''\rangle=0$, and the third property of $\Pi$ gives
\[
\langle b',\Pi b''\rangle=-\int\mathrm{d}y\,\operatorname{div}b'(y)\cdot\int\mathrm{d}y'\,G(y-y')\cdot\operatorname{div}b''(y').
\]
The function $(y,y')\mapsto G(y-y')$ is smooth on the open set of the pairs $(y,y')$ with $|y-y'|>\frac{\delta}{2}$, and this set contains $\operatorname{supp}b'\times\operatorname{supp}b''$. An integration by parts in $y'$, followed by an integration by parts in $y$, therefore yields
\[
\langle b',\Pi b''\rangle=\int\mathrm{d}y\int\mathrm{d}y'\,\sum_{k,l=1}^3b'_k(y)\cdot\left(\partial_k\partial_lG\right)(y-y')\cdot b''_l(y').
\]
For $y\ne0$ the matrix of the second partial derivatives of $G$ at $y$ is
\[
\left(\left(\partial_k\partial_lG\right)(y)\right)_{k,l=1}^3=\frac{1}{4\pi\cdot|y|^5}\cdot\left(3\,y\,y^T-|y|^2\cdot1\right),
\]
where $1$ is the identity matrix. The symmetric matrix $3\,y\,y^T-|y|^2\cdot1$ has the eigenvalue $2|y|^2$ on the line spanned by $y$ and the eigenvalue $-|y|^2$ on the plane orthogonal to $y$, so that its operator norm is $2|y|^2$. Hence the matrix $\left(\left(\partial_k\partial_lG\right)(y)\right)_{k,l}$ has the operator norm $\frac{1}{2\pi\cdot|y|^3}$. As $|y-y'|\ge\delta$ for $y\in\operatorname{supp}b'$ and $y'\in\operatorname{supp}b''$, the integrand above is bounded in absolute value by $\frac{1}{2\pi\cdot\delta^3}\cdot|b'(y)|\cdot|b''(y')|$, and the inequality follows.

In the general case let $\eta\in C_c^\infty(\mathbb R^3)$ be non-negative, with support in the closed unit ball and with $\int\mathrm{d}z\,\eta(z)=1$. For integers $n>\frac{2}{\delta}$ put $\delta_n=\delta-\frac{2}{n}>0$ and
\[
b'_n(y)=\int\mathrm{d}z\,n^3\cdot\eta\left(n\cdot(y-z)\right)\cdot b'(z),\qquad b''_n(y)=\int\mathrm{d}z\,n^3\cdot\eta\left(n\cdot(y-z)\right)\cdot b''(z).
\]
Then $b'_n,b''_n\in\mathcal D$. Every point of $\operatorname{supp}b'_n$ has distance at most $\frac{1}{n}$ from $\operatorname{supp}b'$, and every point of $\operatorname{supp}b''_n$ has distance at most $\frac{1}{n}$ from $\operatorname{supp}b''$; so $|y-y'|\ge\delta_n$ for $y\in\operatorname{supp}b'_n$ and $y'\in\operatorname{supp}b''_n$. Since $\eta\ge0$, one has $|b'_n(y)|\le\int\mathrm{d}z\,n^3\cdot\eta\left(n\cdot(y-z)\right)\cdot|b'(z)|$, whence $\|b'_n\|_{L^1}\le\|b'\|_{L^1}$, and likewise $\|b''_n\|_{L^1}\le\|b''\|_{L^1}$. Moreover $b'_n\to b'$ and $b''_n\to b''$ in $L^2(\mathbb R^3;\mathbb R^3)$ as $n\to\infty$. By the first case,
\[
\left|\langle b'_n,\Pi b''_n\rangle\right|\le\frac{1}{2\pi\cdot\delta_n^3}\cdot\|b'\|_{L^1}\cdot\|b''\|_{L^1},
\]
and as $n\to\infty$ the left side converges to $\left|\langle b',\Pi b''\rangle\right|$, because $\Pi$ is bounded, while $\delta_n\to\delta$.

(i) By Lemma~\ref{lem:gaussian-mixtures}, $\gamma$ satisfies~\eqref{eq:setup-G}. Let $g$ be a Euclidean map. By the first property of $\Pi$, and since $g$ acts isometrically on $L^2(\mathbb R^3;\mathbb R^3)$ (Section~\ref{sec:euclidean-maps-actions}), one has $q(g\cdot f)=\|g\cdot\Pi f\|^2=\|\Pi f\|^2=q(f)$ for every $f\in\mathcal S$. The last assertion of Lemma~\ref{lem:gaussian-mixtures}, applied with the identity map of the one-point probability space, gives $g_*\gamma=\gamma$.

For $f\in\mathcal S$ and $t\in\mathbb R$ one has $\gamma\left(\exp\left(i\,t\cdot X(f)\right)\right)=\gamma\left(\exp\left(iX(tf)\right)\right)=\exp\left(-\frac{1}{2}\cdot t^2\cdot q(f)\right)$. By the uniqueness theorem for characteristic functions, $X(f)$ is under $\gamma$ a centred normal random variable with variance $q(f)$, and $\gamma\left(X(f)^2\right)=q(f)$. Let now $f\in\mathcal S$ satisfy $\ast f=\nabla\mathsf g$ for some $\mathsf g\in\mathcal S(\mathbb R^3)$. Since $\mathsf g\in H^1(\mathbb R^3)$, the second property of $\Pi$ gives $\gamma\left(X(f)^2\right)=\|\Pi\nabla\mathsf g\|^2=0$. By the last assertion of Section~\ref{sec:closed-distributions}, which rests on a sequence dense in $\mathcal S(\mathbb R^3)$ and on the continuity of every $\omega\in Q$ on $\mathcal S$, the law $\gamma$ is concentrated on $\{\omega\in Q:d\omega=0\}$.

It remains to prove $(RP_j)$. Consider first $j=1$. Let $f\in\mathcal D(\{x_1>0\})$, $b=\ast f$ and $h=\operatorname{div}b\in C_c^\infty(\mathbb R^3)$. The support of $b$ is a compact subset of the open half space $\{x_1>0\}$; hence there is $t_0>0$ such that $h(x)=0$ whenever $x_1<t_0$. As $\det\theta_1=-1$, the action of $\theta_1$ on vector fields reads
\[
(\theta_1\cdot b)(x)=\left(b_1(\theta_1x),-b_2(\theta_1x),-b_3(\theta_1x)\right)\qquad\left(x\in\mathbb R^3\right),
\]
whence $\operatorname{div}(\theta_1\cdot b)=-h\circ\theta_1$. The supports of $\theta_1\cdot b$ and of $b$ lie in $\{x_1<0\}$ and in $\{x_1>0\}$, respectively, so that $\langle\theta_1\cdot b,b\rangle=0$, and the third property of $\Pi$ yields
\[
\langle\theta_1\cdot f,\Pi f\rangle=\langle\theta_1\cdot b,\Pi b\rangle=\left\langle h\circ\theta_1,(-\Delta)^{-1}h\right\rangle.
\]
Write the points of $\mathbb R^3$ as $(t,y)$ with $t\in\mathbb R$ and $y\in\mathbb R^2$, so that $G(t,y)=\frac{1}{4\pi}\cdot\left(t^2+|y|^2\right)^{-\frac{1}{2}}$ for $(t,y)\ne0$. For $s\in\mathbb R\setminus\{0\}$ and $y\in\mathbb R^2$ one has
\[
G(s,y)=\int\frac{\mathrm{d}^2k}{(2\pi)^2}\,\exp(ik\cdot y)\cdot\frac{\exp\left(-|k|\cdot|s|\right)}{2|k|},
\]
where $k$ ranges over $\mathbb R^2$ and $k\cdot y$ is the scalar product of $\mathbb R^2$. The integral converges absolutely, because
\[
\int\mathrm{d}^2k\,\frac{\exp\left(-|k|\cdot|s|\right)}{2|k|}=\frac{\pi}{|s|}.
\]
In polar coordinates $k=\kappa\cdot(\cos\alpha,\sin\alpha)$ the integration over $\alpha$ produces $2\pi\cdot J_0(\kappa\cdot|y|)$, where $J_0$ is the Bessel function of order zero, and the formula reduces to Lipschitz's integral $\int_0^\infty\mathrm{d}\kappa\,\exp(-\kappa\cdot|s|)\cdot J_0(\kappa\cdot|y|)=\left(s^2+|y|^2\right)^{-\frac{1}{2}}$.

The function $h\circ\theta_1$ vanishes at $(t,y)$ unless $t\le-t_0$. The substitution $t\mapsto-t$ in the first factor, together with $G(-t,y)=G(t,y)$, therefore gives
\[
\left\langle h\circ\theta_1,(-\Delta)^{-1}h\right\rangle=\int_0^\infty\mathrm{d}t\int_0^\infty\mathrm{d}t'\int\mathrm{d}^2y\int\mathrm{d}^2y'\,h(t,y)\cdot G(t+t',y-y')\cdot h(t',y').
\]
Insert here the formula for $G(s,y-y')$ with $s=t+t'>0$. The resulting integral over $(t,t',y,y',k)$ converges absolutely: integrating the absolute value of its integrand first over $k$, one obtains
\[
\int_0^\infty\mathrm{d}t\int_0^\infty\mathrm{d}t'\,\frac{\|h(t,\cdot)\|_{L^1(\mathbb R^2)}\cdot\|h(t',\cdot)\|_{L^1(\mathbb R^2)}}{4\pi\cdot(t+t')}\le\frac{1}{8\pi\cdot t_0}\cdot\left(\int\mathrm{d}t\,\|h(t,\cdot)\|_{L^1(\mathbb R^2)}\right)^2<\infty,
\]
since $h$ is bounded, has compact support and vanishes for $t<t_0$. By Fubini's theorem the integration over $k$ may be carried out last. For $t\in\mathbb R$ and $k\in\mathbb R^2$ let
\[
\tilde h(t,k)=\int\mathrm{d}^2y\,\exp(-ik\cdot y)\cdot h(t,y);
\]
as $h$ is real, $\int\mathrm{d}^2y\,\exp(ik\cdot y)\cdot h(t,y)=\overline{\tilde h(t,k)}$. Therefore
\[
\langle\theta_1\cdot f,\Pi f\rangle=\left\langle h\circ\theta_1,(-\Delta)^{-1}h\right\rangle=\int\frac{\mathrm{d}^2k}{(2\pi)^2}\,\frac{1}{2|k|}\cdot\left|\int_0^\infty\mathrm{d}t\,\exp\left(-|k|\cdot t\right)\cdot\tilde h(t,k)\right|^2\ge0.
\]

Now let $n\ge1$, $c=(c_1,\dots,c_n)\in\mathbb C^n$ and $f_1,\dots,f_n\in\mathcal D(\{x_1>0\})$, and put $M_{pq}=\langle\theta_1\cdot f_p,\Pi f_q\rangle\in\mathbb R$ for $p,q=1,\dots,n$. The map $f\mapsto\theta_1\cdot f$ is an isometric involution of $L^2(\mathbb R^3;\mathbb R^3)$, hence symmetric, and it commutes with $\Pi$; therefore
\[
M_{qp}=\langle f_q,\Pi(\theta_1\cdot f_p)\rangle=\langle\Pi f_q,\theta_1\cdot f_p\rangle=M_{pq},
\]
and the real matrix $M=\left(M_{pq}\right)_{p,q}$ is symmetric. For $a\in\mathbb R^n$ the form $f=\sum_{p=1}^na_p\cdot f_p$ belongs to $\mathcal D(\{x_1>0\})$, and by bilinearity $\langle\theta_1\cdot f,\Pi f\rangle=\sum_{p,q=1}^na_p\cdot a_q\cdot M_{pq}$. The inequality $\langle\theta_1\cdot f,\Pi f\rangle\ge0$, which has been proved above for every $f\in\mathcal D(\{x_1>0\})$, therefore gives $\sum_{p,q=1}^na_p\cdot a_q\cdot M_{pq}\ge0$; so $M$ is positive semidefinite. By the Schur product theorem the matrices $\left(M_{pq}^m\right)_{p,q}$, $m=0,1,2,\dots$, are positive semidefinite (for $m=0$ this is the matrix all of whose entries equal $1$), and hence so is the matrix
\[
\left(\exp\left(M_{pq}\right)\right)_{p,q}=\sum_{m=0}^{\infty}\frac{1}{m!}\cdot\left(M_{pq}^m\right)_{p,q}.
\]
Being real, symmetric and positive semidefinite, it satisfies $\sum_{p,q=1}^n\overline{z_p}\cdot z_q\cdot\exp\left(M_{pq}\right)\ge0$ for all $z\in\mathbb C^n$. For $p,q\in\{1,\dots,n\}$ one has $X(f_q)-X(\theta_1\cdot f_p)=X(f_q-\theta_1\cdot f_p)$, $\|\Pi(\theta_1\cdot f_p)\|=\|\theta_1\cdot\Pi f_p\|=\|\Pi f_p\|$ and $\langle\Pi f_q,\Pi(\theta_1\cdot f_p)\rangle=\langle\Pi f_q,\theta_1\cdot f_p\rangle=M_{pq}$, so that
\[
\left\|\Pi\left(f_q-\theta_1\cdot f_p\right)\right\|^2=\|\Pi f_q\|^2+\|\Pi f_p\|^2-2M_{pq}.
\]
Hence, with $c'_p=c_p\cdot\exp\left(-\frac{1}{2}\cdot\|\Pi f_p\|^2\right)$,
\[
\sum_{p,q=1}^n\overline{c_p}\cdot c_q\cdot\gamma\left(\exp\left(i\left(X(f_q)-X(\theta_1\cdot f_p)\right)\right)\right)=\sum_{p,q=1}^n\overline{c'_p}\cdot c'_q\cdot\exp\left(M_{pq}\right)\ge0,
\]
which is $(RP_1)$.

Let $j\in\{2,3\}$, and let $w\in O_{\mathrm{cub}}$ be the matrix that interchanges $e_1$ and $e_j$ and fixes the third vector of the standard basis. Then $w\,\theta_1\,w^{-1}=\theta_j$ and $w\{x_1>0\}=\{x_j>0\}$, so that $w\cdot\mathcal D(\{x_1>0\})=\mathcal D(\{x_j>0\})$ (Section~\ref{sec:euclidean-maps-actions}). Let $n\ge1$, $c\in\mathbb C^n$ and $f_1,\dots,f_n\in\mathcal D(\{x_j>0\})$, and put $f'_p=w^{-1}\cdot f_p\in\mathcal D(\{x_1>0\})$ for $p=1,\dots,n$. The operator $\hat w$ is multiplicative, commutes with the composition with the exponential function and satisfies $\hat wX(f)=X(w\cdot f)$ (Section~\ref{sec:euclidean-maps-actions}); since moreover $w\cdot(\theta_1\cdot f'_p)=\theta_j\cdot f_p$,
\[
\hat w\exp\left(i\left(X(f'_q)-X(\theta_1\cdot f'_p)\right)\right)=\exp\left(i\left(X(f_q)-X(\theta_j\cdot f_p)\right)\right)\qquad\left(p,q=1,\dots,n\right).
\]
As $w_*\gamma=\gamma$, one has $\gamma\left(\hat wF\right)=\gamma(F)$ for every bounded measurable function $F$ on $Q$. Consequently the sum in $(RP_j)$ formed with $c$ and $f_1,\dots,f_n$ equals the sum in $(RP_1)$ formed with $c$ and $f'_1,\dots,f'_n$, which is non-negative.
\end{proof}

\paragraph{Coordinatewise changes of variables.}
For an interval $I\subset\mathbb R$ let $C_c(I)$ be the set of the continuous functions $\phi\colon\mathbb R\to\mathbb R$ whose support is a compact subset of $I$. Let $\mathcal U$ be the set of the classes, modulo equality almost everywhere, of the Lebesgue measurable functions $u\colon\mathbb R\to\left[\frac{1}{4},\frac{1}{2}\right]$, endowed with the $\sigma$-algebra generated by the maps
\[
\mathcal U\to\mathbb R,\qquad u\mapsto\int\mathrm{d}s\,\phi(s)\cdot u(s)\qquad\left(\phi\in C_c(\mathbb R)\right).
\]
For an interval $I$, let $\sigma(u|_I)$ denote the $\sigma$-algebra on $\mathcal U$ generated by the maps $u\mapsto\int\mathrm{d}s\,\phi(s)\cdot u(s)$ with $\phi\in C_c(I)$.

Let $\varsigma$ be a telegraph process, with the probability $\mathbb P$, the expectation $\mathbb E$ and the $\sigma$-algebras $\mathcal F_I$ of Section~\ref{sec:auxiliary-results}, and let $u$ be the class of the function $s\mapsto\frac{1}{4}\cdot\left(1+\varsigma(s)\right)$; it is a random element of $\mathcal U$. Let $I$ be an interval and $\phi\in C_c(I)$. Every path of $\varsigma$ is right continuous and has therefore at most countably many points of discontinuity; so $\phi\cdot\varsigma$ is Riemann integrable for every path, and
\[
\int\mathrm{d}s\,\phi(s)\cdot\varsigma(s)=\lim_{n\to\infty}2^{-n}\cdot\sum_{k\in\mathbb Z}\phi\left(k\cdot2^{-n}\right)\cdot\varsigma\left(k\cdot2^{-n}\right).
\]
For each $n$ only the finitely many points $k\cdot2^{-n}$ of $\operatorname{supp}\phi\subset I$ contribute to the sum. Hence $\int\mathrm{d}s\,\phi(s)\cdot u(s)=\frac{1}{4}\int\mathrm{d}s\,\phi(s)+\frac{1}{4}\int\mathrm{d}s\,\phi(s)\cdot\varsigma(s)$ is $\mathcal F_I$-measurable. Next let $\phi_1,\dots,\phi_n\in C_c(\mathbb R)$. The random vector with the components $\int\mathrm{d}s\,\phi_p(s)\cdot u(s)$, $p=1,\dots,n$, takes its values in a bounded set. The value $\varsigma(s)$ is a jointly measurable function of $s\in\mathbb R$ and of the sample point: by right continuity it is the limit, as $j\to\infty$, of the value at $s$ of the process which on each interval $\left[k\cdot2^{-j},(k+1)\cdot2^{-j}\right)$, $k\in\mathbb Z$, is equal to $\varsigma\left((k+1)\cdot2^{-j}\right)$, and these processes are jointly measurable. Hence Fubini's theorem applies, and each mixed moment of the random vector has the form
\[
4^{-m}\cdot\int\mathrm{d}s_1\cdots\int\mathrm{d}s_m\,\phi_{p_1}(s_1)\cdots\phi_{p_m}(s_m)\cdot\mathbb E\left[\left(1+\varsigma(s_1)\right)\cdots\left(1+\varsigma(s_m)\right)\right]
\]
with $m\ge0$ and $p_1,\dots,p_m\in\{1,\dots,n\}$, and the expectations in the integrand are determined by the finite-dimensional distributions of $\varsigma$. The law of a bounded random vector is determined by its mixed moments; therefore the joint law of finitely many of the integrals $\int\mathrm{d}s\,\phi(s)\cdot u(s)$, $\phi\in C_c(\mathbb R)$, is determined by the finite-dimensional distributions of $\varsigma$. The sets
\[
\left\{u'\in\mathcal U:\left(\int\mathrm{d}s\,\phi_p(s)\cdot u'(s)\right)_{p=1,\dots,n}\in B\right\},
\]
with $n\ge1$, $\phi_1,\dots,\phi_n\in C_c(\mathbb R)$ and Borel sets $B\subset\mathbb R^n$, form a $\pi$-system which generates the $\sigma$-algebra of $\mathcal U$; by Dynkin's $\pi$-$\lambda$ theorem the law of $u$ on $\mathcal U$ is thus determined by the finite-dimensional distributions of $\varsigma$ as well. Since the integrals $\int\mathrm{d}s\,\phi(s)\cdot u(s)$ with $\phi\in C_c(I)$ are $\mathcal F_I$-measurable, every bounded $\sigma(u|_I)$-measurable function of $u$ is a bounded $\mathcal F_I$-measurable function of $\varsigma$.

For $v\in\mathcal U$ and $r\in\mathbb R$ let $\check v$ and $v(\cdot+r)$ be the classes of the functions $s\mapsto v(-s)$ and $s\mapsto v(s+r)$, respectively. For $\phi\in C_c(\mathbb R)$ one has
\[
\int\mathrm{d}s\,\phi(s)\cdot\check v(s)=\int\mathrm{d}s\,\phi(-s)\cdot v(s),\qquad\int\mathrm{d}s\,\phi(s)\cdot v(s+r)=\int\mathrm{d}s\,\phi(s-r)\cdot v(s),
\]
and the functions $s\mapsto\phi(-s)$ and $s\mapsto\phi(s-r)$ belong to $C_c(\mathbb R)$; hence the maps $v\mapsto\check v$ and $v\mapsto v(\cdot+r)$ of $\mathcal U$ into itself are measurable. The class $\check u$ is the class of the function $s\mapsto\frac{1}{4}\cdot\left(1+\check\varsigma(s)\right)$. The properties~\eqref{eq:ergodic-components-4}--\eqref{eq:ergodic-components-7} of $\varsigma$ yield the following four statements on $u$.
\begin{itemize}
\item For every $r\in\mathbb R$ the random elements $\check u$ and $u(\cdot+r)$ of $\mathcal U$ have the law of $u$. Indeed, let $n\ge1$ and $\phi_1,\dots,\phi_n\in C_c(\mathbb R)$. By the identities just displayed, the random vector with the components $\int\mathrm{d}s\,\phi_p(s)\cdot\check u(s)$, respectively $\int\mathrm{d}s\,\phi_p(s)\cdot u(s+r)$, $p=1,\dots,n$, is the random vector formed as above with the functions $s\mapsto\phi_p(-s)$, respectively $s\mapsto\phi_p(s-r)$, in place of $\phi_p$. The formula for the mixed moments, applied to these functions and followed by the substitution of $-s_i$, respectively $s_i+r$, for $s_i$, $i=1,\dots,m$, shows that its mixed moments are given by the same formula with $\varsigma(-s_i)$, respectively $\varsigma(s_i+r)$, in place of $\varsigma(s_i)$. By~\eqref{eq:ergodic-components-4} the processes $\check\varsigma$ and $\left(\varsigma(s+r)\right)_{s\in\mathbb R}$ have the finite-dimensional distributions of $\varsigma$, so that these expectations coincide with those of the formula. Hence the random vector in question and the random vector with the components $\int\mathrm{d}s\,\phi_p(s)\cdot u(s)$, $p=1,\dots,n$, both of which take their values in a bounded set, have the same mixed moments and therefore the same law; and Dynkin's $\pi$-$\lambda$ theorem, applied to the $\pi$-system above, gives the assertion.
\item Let $r\in\mathbb R$ and $t\ge0$, and let $\Psi_1$ and $\Psi_2$ be bounded complex functions on $\mathcal U$ which are measurable for $\sigma\left(u|_{(-\infty,r]}\right)$ and for $\sigma\left(u|_{[r+t,\infty)}\right)$, respectively. By the last assertion of the preceding paragraph, $\Psi_1(u)$ is $\mathcal F_{(-\infty,r]}$-measurable and $\Psi_2(u)$ is $\mathcal F_{[r+t,\infty)}$-measurable, and~\eqref{eq:ergodic-components-5} gives
\[
\left|\mathbb E\left[\Psi_1(u)\cdot\Psi_2(u)\right]-\mathbb E\left[\Psi_1(u)\right]\cdot\mathbb E\left[\Psi_2(u)\right]\right|\le\exp(-2t)\cdot\sup|\Psi_1|\cdot\sup|\Psi_2|.
\]
\item Let $I$ be a bounded interval with the end points $a<b$, let $z\in\{0,1\}$, and put $c=\frac{1}{4}\cdot(1+z)$. For every $v\in\mathcal U$ the function $v-c$ is non-negative if $z=0$ and non-positive if $z=1$, so that $v=c$ almost everywhere on $I$ if and only if $\int_a^b\mathrm{d}s\,\left(v(s)-c\right)=0$. The set of these $v$ is therefore measurable: by dominated convergence, the map $v\mapsto\int_a^b\mathrm{d}s\,v(s)$ is the pointwise limit of the maps $v\mapsto\int\mathrm{d}s\,\phi_j(s)\cdot v(s)$, where the functions $\phi_j\in C_c((a,b))$ take their values in $[0,1]$ and converge to $1$ at every point of $(a,b)$. Moreover, $u=c$ almost everywhere on $I$, that is, $\varsigma=z$ almost everywhere on $I$, if and only if $\varsigma(s)=z$ for every $s\in[a,b)$. Indeed, if $\varsigma=z$ almost everywhere on $I$ and $s\in[a,b)$, then $s$ is the limit of a decreasing sequence of points of $(s,b)\subset I$ at which $\varsigma$ takes the value $z$, and $\varsigma(s)=z$ by right continuity; conversely, if $\varsigma=z$ on $[a,b)$, then $\varsigma=z$ on $I$ except possibly at the point $b$. By~\eqref{eq:ergodic-components-6}, applied to the interval $[a,b)$, whose length is $|I|$, this event has the probability $\frac{1}{2}\cdot\exp\left(-|I|\right)>0$.
\item Let $\Psi$ be a bounded complex function on $\mathcal U$ which is measurable for $\sigma\left(u|_{(0,\infty)}\right)$. Then $\mathbb E\left[\overline{\Psi(\check u)}\cdot\Psi(u)\right]\ge0$. Indeed, by the Doob--Dynkin lemma there is a function $\chi\colon\mathbb R^{C_c((0,\infty))}\to\mathbb C$, measurable for the product $\sigma$-algebra, such that
\[
\Psi(v)=\chi\left(\left(\int\mathrm{d}s\,\phi(s)\cdot v(s)\right)_{\phi\in C_c((0,\infty))}\right)\qquad\left(v\in\mathcal U\right),
\]
and $\chi$ may be taken bounded, since its values of modulus larger than $\sup|\Psi|$ may be replaced by $0$. With $L_\phi(y)=\int\mathrm{d}s\,\phi(s)\cdot y(s)$ as in~\eqref{eq:ergodic-components-7}, one has, for $\phi\in C_c((0,\infty))$, $\int\mathrm{d}s\,\phi(s)\cdot u(s)=\frac{1}{4}\int\mathrm{d}s\,\phi(s)+\frac{1}{4}\cdot L_\phi(\varsigma)$ and $\int\mathrm{d}s\,\phi(s)\cdot\check u(s)=\frac{1}{4}\int\mathrm{d}s\,\phi(s)+\frac{1}{4}\cdot L_\phi(\check\varsigma)$. For $x=(x_\phi)_{\phi\in C_c((0,\infty))}$ in $\mathbb R^{C_c((0,\infty))}$ put
\[
\tilde\chi(x)=\chi\left(\left(\frac{1}{4}\int\mathrm{d}s\,\phi(s)+\frac{1}{4}\cdot x_\phi\right)_{\phi\in C_c((0,\infty))}\right);
\]
the function $\tilde\chi$ is bounded and measurable for the product $\sigma$-algebra, as each coordinate of its inner map is an affine function of one coordinate of $x$, and $\Psi(u)$ and $\Psi(\check u)$ are its values at $\left(L_\phi(\varsigma)\right)_\phi$ and at $\left(L_\phi(\check\varsigma)\right)_\phi$. Property~\eqref{eq:ergodic-components-7}, applied to the function $y\mapsto\tilde\chi\left(\left(L_\phi(y)\right)_\phi\right)$, gives the inequality.
\end{itemize}

Let $u=(u_1,u_2,u_3)\in\mathcal U^3$. For $l=1,2,3$ let
\[
\varphi_l(s)=\int_0^s\mathrm{d}t\,u_l(t)\qquad\left(s\in\mathbb R\right).
\]
As $\frac{1}{4}\le u_l\le\frac{1}{2}$, the function $\varphi_l$ is a continuous and strictly increasing bijection of $\mathbb R$ onto $\mathbb R$. Let $\Phi_u\colon\mathbb R^3\to\mathbb R^3$ be the bijection
\[
\Phi_u(x)=\left(\varphi_1(x_1),\varphi_2(x_2),\varphi_3(x_3)\right),
\]
and for $b\in L^2(\mathbb R^3;\mathbb R^3)$ let $\Phi_{u\#}b$ be the vector field on $\mathbb R^3$ with the components
\begin{equation}\label{eq:ergodic-components-8}
\left(\Phi_{u\#}b\right)_k(y)=\frac{b_k(x)}{u_k(x_k)},\qquad x=\Phi_u^{-1}(y)\qquad\left(y\in\mathbb R^3,\ k=1,2,3\right).
\end{equation}
The substitution $y=\Phi_u(x)$ transforms the Lebesgue measure according to $\mathrm{d}y=u_1(x_1)\cdot u_2(x_2)\cdot u_3(x_3)\,\mathrm{d}x$; that is,
\[
\int\mathrm{d}y\,H(y)=\int\mathrm{d}x\,u_1(x_1)\cdot u_2(x_2)\cdot u_3(x_3)\cdot H\left(\Phi_u(x)\right)
\]
for every measurable function $H\colon\mathbb R^3\to[0,\infty]$. By Fubini's theorem this reduces to the substitution rule in one variable for the increasing absolutely continuous functions $\varphi_l$, whose derivatives equal $u_l$ almost everywhere. In particular, $\Phi_{u\#}$ is the transpose of the pullback of two-forms by $\Phi_u$: if $F\in L^2(\mathbb R^3;\Lambda^2)$ and $\Phi_u^*F$ is the form with the components
\[
\left(\Phi_u^*F\right)_{il}(x)=u_i(x_i)\cdot u_l(x_l)\cdot F_{il}\left(\Phi_u(x)\right)\qquad\left(i,l=1,2,3,\ x\in\mathbb R^3\right),
\]
then $\langle\Phi_u^*F,b\rangle=\langle F,\Phi_{u\#}b\rangle$ for every $b\in L^2(\mathbb R^3;\mathbb R^3)$, the scalar products being formed through the identification of forms with vector fields. Finally, for $a\in\mathbb R^3$ and for $w\in O_{\mathrm{cub}}$ with $w\,e_i=\varepsilon_i\cdot e_{\pi(i)}$, $i=1,2,3$, let $u^a,u^w\in\mathcal U^3$ be given by
\[
u^a_l(s)=u_l(s+a_l),\qquad u^w_i(s)=u_{\pi(i)}\left(\varepsilon_i\cdot s\right)\qquad\left(s\in\mathbb R,\ i,l=1,2,3\right).
\]

\begin{lemma}[Coordinatewise changes of variables]\label{lem:coordinatewise-changes-variables}
Let $u\in\mathcal U^3$ and $b\in L^2(\mathbb R^3;\mathbb R^3)$.
\begin{enumerate}
\item[(i)] For $l=1,2,3$ and $s,s'\in\mathbb R$,
\[
\frac{1}{4}\cdot|s-s'|\le\left|\varphi_l(s)-\varphi_l(s')\right|\le\frac{1}{2}\cdot|s-s'|.
\]
The map $\Phi_u$ is a homeomorphism of $\mathbb R^3$ onto itself with $\Phi_u(0)=0$, and it maps each of the half spaces $\{x_l>0\}$ and $\{x_l<0\}$, $l=1,2,3$, onto itself. Moreover, with the $L^1$ norms of Lemma~\ref{lem:properties-free-maxwell}(ii),
\[
\|\Phi_{u\#}b\|\le\|b\|,\qquad\|\Phi_{u\#}b\|_{L^1}\le\frac{1}{2}\cdot\|b\|_{L^1},\qquad\operatorname{supp}\Phi_{u\#}b=\Phi_u(\operatorname{supp}b).
\]
\item[(ii)] For every $a\in\mathbb R^3$ one has $\Phi_{u\#}(\tau_a\cdot b)=\tau_{\Phi_u(a)}\cdot\Phi_{u^a\#}b$.
\item[(iii)] For every $w\in O_{\mathrm{cub}}$ one has $\Phi_{u\#}(w\cdot b)=w\cdot\Phi_{u^w\#}b$.
\item[(iv)] For every $\mathsf g\in\mathcal S(\mathbb R^3)$ one has $\Phi_{u\#}\nabla\mathsf g=\nabla\left(\mathsf g\circ\Phi_u^{-1}\right)$, the gradient on the right being taken in the sense of distributions; consequently $\Pi\Phi_{u\#}\nabla\mathsf g=0$.
\item[(v)] Let $r>0$, and let $b\in\mathcal D$ have its support in $[-r,r]^3$. Then $\Phi_{u\#}b$ has its support in $\left[-\frac{r}{2},\frac{r}{2}\right]^3$, and it depends on $u$ only through the restrictions $u_1|_{[-r,r]}$, $u_2|_{[-r,r]}$, $u_3|_{[-r,r]}$; if moreover $\operatorname{supp}b\subset\{x_l>0\}$ for some $l\in\{1,2,3\}$, then $u_l|_{[-r,r]}$ may be replaced by $u_l|_{(0,r]}$ in this assertion. The map $u\mapsto\Phi_{u\#}b\in L^2(\mathbb R^3;\mathbb R^3)$ is continuous for the convergence of $u_1,u_2,u_3$ in $L^1([-r,r])$: if $u^{(n)}\in\mathcal U^3$ for $n\ge1$ and $u^{(n)}_l\to u_l$ in $L^1([-r,r])$ for $l=1,2,3$, then $\Phi_{u^{(n)}\#}b\to\Phi_{u\#}b$ in $L^2(\mathbb R^3;\mathbb R^3)$.
\end{enumerate}
\end{lemma}

\begin{proof}
(i) Let $l\in\{1,2,3\}$ and $s<s'$. Since $\frac{1}{4}\le u_l\le\frac{1}{2}$,
\[
\frac{1}{4}\cdot(s'-s)\le\varphi_l(s')-\varphi_l(s)=\int_s^{s'}\mathrm{d}t\,u_l(t)\le\frac{1}{2}\cdot(s'-s).
\]
This gives the two inequalities, and it shows that $\varphi_l$ is a strictly increasing bijection of $\mathbb R$ onto $\mathbb R$ with $\varphi_l(0)=0$, positive on $(0,\infty)$ and negative on $(-\infty,0)$, such that $\varphi_l$ and $\varphi_l^{-1}$ are Lipschitz. Hence $\Phi_u$ is a homeomorphism of $\mathbb R^3$ with $\Phi_u(0)=0$ which maps each of the half spaces $\{x_l>0\}$, $\{x_l<0\}$ onto itself; and $\Phi_u$ and $\Phi_u^{-1}$ are Lipschitz, so that both map null sets onto null sets. By~\eqref{eq:ergodic-components-8} and the substitution rule,
\[
\|\Phi_{u\#}b\|^2=\int\mathrm{d}x\,u_1(x_1)\cdot u_2(x_2)\cdot u_3(x_3)\cdot\sum_{k=1}^3\frac{|b_k(x)|^2}{u_k(x_k)^2}=\int\mathrm{d}x\,\sum_{k=1}^3|b_k(x)|^2\cdot\frac{u_i(x_i)\cdot u_l(x_l)}{u_k(x_k)},
\]
where, in the term with the index $k$, $i$ and $l$ are the two indices different from $k$. As $u_i\cdot u_l\le\frac{1}{4}$ and $u_k\ge\frac{1}{4}$, the last fraction is at most $1$, and $\|\Phi_{u\#}b\|\le\|b\|$. As $u_k^{-1}\le4$, the vector $\left(\Phi_{u\#}b\right)\left(\Phi_u(x)\right)$, whose components are $\frac{b_k(x)}{u_k(x_k)}$, has a Euclidean norm at most $4\cdot|b(x)|$; as moreover $u_1\cdot u_2\cdot u_3\le\frac{1}{8}$,
\[
\|\Phi_{u\#}b\|_{L^1}=\int\mathrm{d}x\,u_1(x_1)\cdot u_2(x_2)\cdot u_3(x_3)\cdot\left|\left(\Phi_{u\#}b\right)\left(\Phi_u(x)\right)\right|\le\int\mathrm{d}x\,\frac{1}{8}\cdot4\cdot|b(x)|=\frac{1}{2}\cdot\|b\|_{L^1}.
\]
Finally, let $O\subset\mathbb R^3$ be open. Since the $u_k$ are positive and $\Phi_u$, $\Phi_u^{-1}$ map null sets onto null sets, $\Phi_{u\#}b$ vanishes almost everywhere on $\Phi_u(O)$ if and only if $b$ vanishes almost everywhere on $O$. As $\Phi_u$ is a homeomorphism, the largest open set on which $\Phi_{u\#}b$ vanishes almost everywhere is therefore the image under $\Phi_u$ of the largest open set on which $b$ vanishes almost everywhere; taking complements, one obtains $\operatorname{supp}\Phi_{u\#}b=\Phi_u(\operatorname{supp}b)$.

(ii) Let $a\in\mathbb R^3$. For $l=1,2,3$ and $s\in\mathbb R$,
\[
\varphi_l(s+a_l)=\varphi_l(a_l)+\int_{a_l}^{s+a_l}\mathrm{d}t\,u_l(t)=\varphi_l(a_l)+\int_0^s\mathrm{d}t\,u^a_l(t);
\]
that is, $\Phi_u(x+a)=\Phi_u(a)+\Phi_{u^a}(x)$ for $x\in\mathbb R^3$. Let $y\in\mathbb R^3$ and $x=\Phi_u^{-1}(y)-a$. Then $\Phi_{u^a}(x)=y-\Phi_u(a)$, that is, $x=\Phi_{u^a}^{-1}\left(y-\Phi_u(a)\right)$. As $(\tau_a\cdot b)(x+a)=b(x)$, \eqref{eq:ergodic-components-8} gives, for $k=1,2,3$,
\[
\left(\Phi_{u\#}(\tau_a\cdot b)\right)_k(y)=\frac{b_k(x)}{u_k(x_k+a_k)}=\frac{b_k(x)}{u^a_k(x_k)}=\left(\Phi_{u^a\#}b\right)_k\left(y-\Phi_u(a)\right)=\left(\tau_{\Phi_u(a)}\cdot\Phi_{u^a\#}b\right)_k(y).
\]

(iii) Let $w\,e_i=\varepsilon_i\cdot e_{\pi(i)}$ for $i=1,2,3$. Then $(wx)_{\pi(i)}=\varepsilon_i\cdot x_i$ and $\left(w^{-1}z\right)_i=\varepsilon_i\cdot z_{\pi(i)}$ for $x,z\in\mathbb R^3$. With the substitution $t=\varepsilon_i\cdot t'$ one obtains, for $x\in\mathbb R^3$ and $i=1,2,3$,
\[
\left(w^{-1}\Phi_u(wx)\right)_i=\varepsilon_i\cdot\varphi_{\pi(i)}\left(\varepsilon_i\cdot x_i\right)=\varepsilon_i\cdot\int_0^{\varepsilon_i\cdot x_i}\mathrm{d}t\,u_{\pi(i)}(t)=\int_0^{x_i}\mathrm{d}t'\,u_{\pi(i)}\left(\varepsilon_i\cdot t'\right)=\int_0^{x_i}\mathrm{d}t'\,u^w_i(t');
\]
hence $w^{-1}\circ\Phi_u\circ w=\Phi_{u^w}$. Let $y\in\mathbb R^3$, $x=\Phi_u^{-1}(y)$ and $x'=w^{-1}x$. Then $\Phi_{u^w}(x')=w^{-1}\Phi_u(x)=w^{-1}y$, that is, $x'=\Phi_{u^w}^{-1}\left(w^{-1}y\right)$. By~\eqref{eq:ergodic-components-8} and the formula for $w\cdot b$, for $k=1,2,3$,
\[
\left(\Phi_{u\#}(w\cdot b)\right)_{\pi(k)}(y)=\frac{(w\cdot b)_{\pi(k)}(x)}{u_{\pi(k)}\left(x_{\pi(k)}\right)}=\frac{(\det w)\cdot\varepsilon_k\cdot b_k(x')}{u_{\pi(k)}\left(x_{\pi(k)}\right)}
\]
and
\[
\left(w\cdot\Phi_{u^w\#}b\right)_{\pi(k)}(y)=(\det w)\cdot\varepsilon_k\cdot\left(\Phi_{u^w\#}b\right)_k\left(w^{-1}y\right)=\frac{(\det w)\cdot\varepsilon_k\cdot b_k(x')}{u^w_k(x'_k)}.
\]
Since $x'_k=\varepsilon_k\cdot x_{\pi(k)}$, one has $u^w_k(x'_k)=u_{\pi(k)}\left(\varepsilon_k\cdot x'_k\right)=u_{\pi(k)}\left(x_{\pi(k)}\right)$, so that the two right sides coincide. As $\pi$ is a bijection of $\{1,2,3\}$, this proves (iii).

(iv) Let $\mathsf g\in\mathcal S(\mathbb R^3)$ and $\psi=\mathsf g\circ\Phi_u^{-1}$, and fix representatives of $u_1,u_2,u_3$ with values in $\left[\frac{1}{4},\frac{1}{2}\right]$. By Lebesgue's differentiation theorem there are null sets $E_1,E_2,E_3\subset\mathbb R$ such that $\varphi_l$ is differentiable at every $s\notin E_l$, with $\varphi_l'(s)=u_l(s)\ge\frac{1}{4}$. The sets $\varphi_l(E_l)$ are null sets, $\varphi_l$ being Lipschitz, and at every $t\notin\varphi_l(E_l)$ the inverse function $\varphi_l^{-1}$ is differentiable, with the derivative $\frac{1}{u_l\left(\varphi_l^{-1}(t)\right)}$. Hence, for $k=1,2,3$ and every $y\in\mathbb R^3$ with $y_k\notin\varphi_k(E_k)$, thus for almost every $y\in\mathbb R^3$, the partial derivative of $\psi$ with respect to $y_k$ exists and equals
\[
\frac{\left(\partial_k\mathsf g\right)(x)}{u_k(x_k)},\qquad x=\Phi_u^{-1}(y);
\]
by~\eqref{eq:ergodic-components-8}, the gradient of $\psi$, which is defined almost everywhere, is $\Phi_{u\#}\nabla\mathsf g$. The function $\psi$ is Lipschitz, because $\nabla\mathsf g$ is bounded and $\Phi_u^{-1}$ is Lipschitz by (i); and the gradient of a Lipschitz function, defined almost everywhere, is its gradient in the sense of distributions. This proves the identity in (iv). By (i), $\|\Phi_{u\#}\nabla\mathsf g\|\le\|\nabla\mathsf g\|<\infty$. By (i), moreover, $\left|\varphi_l^{-1}(t)\right|\ge2|t|$ for $t\in\mathbb R$, whence $\left|\Phi_u^{-1}(y)\right|\ge2|y|$ for $y\in\mathbb R^3$; as $\mathsf g$ decreases rapidly, there is a constant $C$ with $|\psi(y)|\le C\cdot(1+2|y|)^{-2}$ for $y\in\mathbb R^3$, and so $\psi\in L^2(\mathbb R^3)$. Thus $\psi\in H^1(\mathbb R^3)$, and the second property of $\Pi$ gives $\Pi\Phi_{u\#}\nabla\mathsf g=\Pi\nabla\psi=0$.

(v) By (i), $\operatorname{supp}\Phi_{u\#}b=\Phi_u(\operatorname{supp}b)\subset\Phi_u\left([-r,r]^3\right)\subset\left[-\frac{r}{2},\frac{r}{2}\right]^3$, because $|\varphi_l(s)|\le\frac{1}{2}\cdot|s|$. Let $K=[-r,r]^3$, and in the case $\operatorname{supp}b\subset\{x_l>0\}$ let instead $K$ be the set of the points $x\in[-r,r]^3$ with $x_l\in(0,r]$. In both cases $\operatorname{supp}b\subset K$, and, as $\varphi_l(s)=\int_0^s\mathrm{d}t\,u_l(t)$, the restriction of $\Phi_u$ to $K$ depends on $u$ only through the restrictions named in (v); hence so do the set $\Phi_u(K)$ and the inverse of the bijection $\Phi_u|_K\colon K\to\Phi_u(K)$. By~\eqref{eq:ergodic-components-8}, $\Phi_{u\#}b$ vanishes outside $\Phi_u(K)$, and at a point $y\in\Phi_u(K)$ its components are $\frac{b_k(x)}{u_k(x_k)}$ with $x=\left(\Phi_u|_K\right)^{-1}(y)\in K$. This proves the assertion on the dependence.

Let now $u^{(n)}$ be as in (v), and let $\varphi^{(n)}_l$ and $\Phi^{(n)}=\Phi_{u^{(n)}}$ be formed with $u^{(n)}$ in place of $u$. For $|s|\le r$ one has $\left|\varphi^{(n)}_l(s)-\varphi_l(s)\right|\le\left\|u^{(n)}_l-u_l\right\|_{L^1([-r,r])}$, so that $\Phi^{(n)}\to\Phi_u$ uniformly on $[-r,r]^3$. Let $v\in C_c(\mathbb R^3;\mathbb R^3)$. By~\eqref{eq:ergodic-components-8} and the substitution rule,
\[
\langle v,\Phi_{u\#}b\rangle=\int\mathrm{d}x\,\sum_{k=1}^3v_k\left(\Phi_u(x)\right)\cdot b_k(x)\cdot u_i(x_i)\cdot u_l(x_l),\qquad\|\Phi_{u\#}b\|^2=\int\mathrm{d}x\,\sum_{k=1}^3|b_k(x)|^2\cdot\frac{u_i(x_i)\cdot u_l(x_l)}{u_k(x_k)},
\]
where, in the terms with the index $k$, $i$ and $l$ are the two indices different from $k$; the same formulas hold with $u^{(n)}$ in place of $u$. The integrands vanish outside $[-r,r]^3$, and $b$ is bounded. On $[-r,r]^3$ the functions $v\circ\Phi^{(n)}$ converge uniformly to $v\circ\Phi_u$, as $v$ is uniformly continuous; and the products $u^{(n)}_i(x_i)\cdot u^{(n)}_l(x_l)$ and the quotients $\frac{u^{(n)}_i(x_i)\cdot u^{(n)}_l(x_l)}{u^{(n)}_k(x_k)}$ converge in $L^1([-r,r]^3)$ to the corresponding expressions formed with $u$, since all the functions $u^{(n)}_m$ and $u_m$, $m=1,2,3$, take their values in $\left[\frac{1}{4},\frac{1}{2}\right]$ and $u^{(n)}_m\to u_m$ in $L^1([-r,r])$. Hence $\langle v,\Phi_{u^{(n)}\#}b\rangle\to\langle v,\Phi_{u\#}b\rangle$ and $\|\Phi_{u^{(n)}\#}b\|\to\|\Phi_{u\#}b\|$. As $\|\Phi_{u^{(n)}\#}b\|\le\|b\|$ for all $n$, by (i), and as $C_c(\mathbb R^3;\mathbb R^3)$ is dense in $L^2(\mathbb R^3;\mathbb R^3)$, the sequence $\Phi_{u^{(n)}\#}b$ converges weakly to $\Phi_{u\#}b$. In a Hilbert space weak convergence together with convergence of the norms implies convergence in norm; so $\Phi_{u^{(n)}\#}b\to\Phi_{u\#}b$ in $L^2(\mathbb R^3;\mathbb R^3)$.
\end{proof}

\begin{lemma}[Mixtures over the changes of variables $\Phi_u$]\label{lem:mixtures-changes-variables}
Let $\mathcal U$ be the measurable space of Section~\ref{sec:free-maxwell-law}, let $\mathcal U^3$ carry the product $\sigma$-algebra, and let $\mathbb P$ be a probability measure on $\mathcal U^3$; let $\mathbb E$ denote integration with respect to $\mathbb P$ in the variable $u=(u_1,u_2,u_3)\in\mathcal U^3$. For $u\in\mathcal U^3$ and $f,g\in L^2(\mathbb R^3;\Lambda^2)$ put
\[
q_u(f,g)=\left\langle\Pi\Phi_{u\#}f,\Pi\Phi_{u\#}g\right\rangle,\qquad q_u(f)=q_u(f,f)=\left\|\Pi\Phi_{u\#}f\right\|^2,
\]
where forms are identified with vector fields through $\ast$, as in Section~\ref{sec:free-maxwell-law}. The probability space $(\mathcal U^3,\mathbb P)$ and the forms $q_u$, $u\in\mathcal U^3$, restricted to $\mathcal S$, satisfy the hypotheses of Lemma~\ref{lem:gaussian-mixtures}. Let $m_{\mathbb P}\in\mathcal M$ be the law which Lemma~\ref{lem:gaussian-mixtures} provides for this probability space and these forms, that is, the law with
\[
m_{\mathbb P}\left(\exp\left(iX(f)\right)\right)=\mathbb E\exp\left(-\frac{1}{2}\cdot q_u(f)\right)\qquad\left(f\in\mathcal S\right).
\]
\begin{enumerate}
\item[(i)] The law $m_{\mathbb P}$ satisfies~\eqref{eq:setup-G} and is concentrated on the set $\{\omega\in Q:d\omega=0\}$ of the closed distributions.
\item[(ii)] If, for every $a\in\mathbb R^3$, the map $u\mapsto u^a$ of $\mathcal U^3$ into itself preserves $\mathbb P$, then $m_{\mathbb P}$ satisfies $(T_{123})$. If $w\in O_{\mathrm{cub}}$ and the map $u\mapsto u^w$ of $\mathcal U^3$ into itself preserves $\mathbb P$, then $w_*m_{\mathbb P}=m_{\mathbb P}$.
\item[(iii)] Let $j\in\{1,2,3\}$. Assume that the coordinates $u_1,u_2,u_3$ are independent under $\mathbb P$, and that the law of $u_j$ under $\mathbb P$ is either the point mass at a constant function belonging to $\mathcal U$ or the law of the class of $\frac{1}{4}\cdot(1+\varsigma)$, where $\varsigma$ is a telegraph process (Section~\ref{sec:auxiliary-results}). Then $m_{\mathbb P}$ satisfies $(RP_j)$.
\end{enumerate}
\end{lemma}

\begin{proof}
Throughout, forms are identified with vector fields as in Section~\ref{sec:free-maxwell-law}, and the translations and the elements of $O_{\mathrm{cub}}$ act on vector fields as described there.

\emph{The hypotheses of Lemma~\ref{lem:gaussian-mixtures}.} Let $u\in\mathcal U^3$. The map $f\mapsto\Pi\Phi_{u\#}f$ is linear, so that $q_u$ is a symmetric bilinear form. Since $\Pi$ is an orthogonal projection, Lemma~\ref{lem:coordinatewise-changes-variables}(i) gives
\[
0\le q_u(f)\le\left\|\Phi_{u\#}f\right\|^2\le\|f\|^2\qquad\left(f\in L^2(\mathbb R^3;\Lambda^2)\right),
\]
and the Cauchy--Schwarz inequality for the positive semidefinite form $q_u$ yields
\[
\left|q_u(f)-q_u(g)\right|=\left|q_u(f-g,f+g)\right|\le\|f-g\|\cdot\|f+g\|\qquad\left(f,g\in L^2(\mathbb R^3;\Lambda^2)\right).
\]
Let $r>0$, and let $f\in\mathcal D$ satisfy $\operatorname{supp}f\subset[-r,r]^3$. By Lemma~\ref{lem:coordinatewise-changes-variables}(v), the field $\Phi_{u\#}f$, and with it the number $q_u(f)$, depends on $u$ only through the restrictions $u_l|_{[-r,r]}$, $l=1,2,3$, and is a continuous function of these restrictions for the norm of $L^1([-r,r])$. For $\phi\in C_c((-r,r))$ let $\ell_\phi$ denote the continuous linear functional $v\mapsto\int\mathrm ds\,\phi(s)\cdot v(s)$ on $L^1([-r,r])$. Approximating the sign of $v$ by such functions $\phi$ with $|\phi|\le1$ (Lusin's theorem and dominated convergence), one finds
\[
\|v\|_{L^1}=\sup\left\{\left|\ell_\phi(v)\right|:\phi\in C_c((-r,r)),\ |\phi|\le1\right\}\qquad\left(v\in L^1([-r,r])\right);
\]
as $\left|\ell_\phi(v)-\ell_{\phi'}(v)\right|\le\sup|\phi-\phi'|\cdot\|v\|_{L^1}$ and $C_c((-r,r))$ is separable for the supremum norm, a suitable countable set of these $\phi$ yields the same supremum. Every open ball of $L^1([-r,r])$ therefore lies in the $\sigma$-algebra generated by the functionals $\ell_\phi$, $\phi\in C_c((-r,r))$; since $L^1([-r,r])$ has a countable dense subset, every open subset of it is a countable union of open balls and lies in this $\sigma$-algebra as well. As the $\ell_\phi$ are continuous, this $\sigma$-algebra coincides with the Borel $\sigma$-algebra of $L^1([-r,r])$, and a map into $L^1([-r,r])$ is Borel measurable as soon as its compositions with all $\ell_\phi$ are measurable. For $\phi\in C_c((-r,r))$ the composition of $u_l\mapsto u_l|_{[-r,r]}$ with $\ell_\phi$ is the function $u_l\mapsto\int\mathrm ds\,\phi(s)\cdot u_l(s)$, one of the functions by which the $\sigma$-algebra of $\mathcal U$ is generated. Hence the maps $u\mapsto u_l|_{[-r,r]}$ are measurable; as $L^1([-r,r])$ is separable, the Borel $\sigma$-algebra of $L^1([-r,r])^3$ is the product of the Borel $\sigma$-algebras of the factors, and therefore $u\mapsto q_u(f)$ is measurable. Now let $f\in\mathcal S$. There are $f_k\in\mathcal D$, $k\ge1$, with $\|f_k-f\|\to0$, each supported in some cube $[-r_k,r_k]^3$. By the bound above, $q_u(f_k)\to q_u(f)$ for every $u\in\mathcal U^3$, so that $u\mapsto q_u(f)$ is measurable as a pointwise limit of measurable functions. Thus the hypotheses of Lemma~\ref{lem:gaussian-mixtures} are satisfied.

(i) The bound~\eqref{eq:setup-G} is part of the assertion of Lemma~\ref{lem:gaussian-mixtures}. Let $f\in\mathcal S$ be such that $\ast f=\nabla\mathsf g$ for some $\mathsf g\in\mathcal S(\mathbb R^3)$. By Lemma~\ref{lem:coordinatewise-changes-variables}(iv) one has $\Pi\Phi_{u\#}\nabla\mathsf g=0$, that is, $q_u(f)=0$ for every $u\in\mathcal U^3$. For $q\ge0$ let $N(0,q)$ be the centred normal law on $\mathbb R$ with variance $q$, the unit mass at $0$ if $q=0$. For $t\in\mathbb R$ the identity defining $m_{\mathbb P}$, applied to $tf$, gives
\[
m_{\mathbb P}\left(\exp\left(i\,t\cdot X(f)\right)\right)=\mathbb E\exp\left(-\frac{1}{2}\cdot t^2\cdot q_u(f)\right)=\mathbb E\int N\left(0,q_u(f)\right)(\mathrm dx)\,\exp\left(i\,t\cdot x\right).
\]
The right side is the characteristic function of the mixture $B\mapsto\mathbb E\left[N\left(0,q_u(f)\right)(B)\right]$ of the laws $N\left(0,q_u(f)\right)$; by the uniqueness theorem for characteristic functions this mixture is the law of $X(f)$ under $m_{\mathbb P}$, and by Tonelli's theorem
\[
m_{\mathbb P}\left(X(f)^2\right)=\mathbb E\int N\left(0,q_u(f)\right)(\mathrm dx)\,x^2=\mathbb E\,q_u(f)=0.
\]
By Section~\ref{sec:closed-distributions}, $m_{\mathbb P}$ is concentrated on $\{\omega\in Q:d\omega=0\}$.

(ii) Let $a\in\mathbb R^3$, and let $w\in O_{\mathrm{cub}}$ with $w\,e_i=\varepsilon_i\cdot e_{\pi(i)}$ for $i=1,2,3$. The maps $u\mapsto u^a$ and $u\mapsto u^w$ of $\mathcal U^3$ into itself are measurable, because for $\phi\in C_c(\mathbb R)$
\[
\int\mathrm ds\,\phi(s)\cdot u_l(s+a_l)=\int\mathrm ds\,\phi(s-a_l)\cdot u_l(s),\qquad\int\mathrm ds\,\phi(s)\cdot u_{\pi(i)}(\varepsilon_i\cdot s)=\int\mathrm ds\,\phi(\varepsilon_i\cdot s)\cdot u_{\pi(i)}(s).
\]
The projection $\Pi$ commutes with the actions of the translations $\tau_c$, $c\in\mathbb R^3$, and of $w$ on vector fields, and these actions preserve the norm (Sections~\ref{sec:euclidean-maps-actions} and~\ref{sec:free-maxwell-law}). Therefore Lemma~\ref{lem:coordinatewise-changes-variables}(ii) and (iii) give, for $u\in\mathcal U^3$ and $f\in L^2(\mathbb R^3;\Lambda^2)$,
\begin{align*}
q_u(\tau_a\cdot f)&=\left\|\Pi\left(\tau_{\Phi_u(a)}\cdot\Phi_{u^a\#}f\right)\right\|^2=\left\|\tau_{\Phi_u(a)}\cdot\Pi\Phi_{u^a\#}f\right\|^2=q_{u^a}(f),\\
q_u(w\cdot f)&=\left\|\Pi\left(w\cdot\Phi_{u^w\#}f\right)\right\|^2=\left\|w\cdot\Pi\Phi_{u^w\#}f\right\|^2=q_{u^w}(f).
\end{align*}
If $u\mapsto u^a$ preserves $\mathbb P$, the second part of Lemma~\ref{lem:gaussian-mixtures}, applied to the Euclidean map $\tau_a$, gives $(\tau_a)_*m_{\mathbb P}=m_{\mathbb P}$; if this holds for every $a\in\mathbb R^3$, then $m_{\mathbb P}$ satisfies $(T_{123})$. In the same way, if $u\mapsto u^w$ preserves $\mathbb P$, then $w_*m_{\mathbb P}=m_{\mathbb P}$.

(iii) Let first $j=1$. By independence, $\mathbb P=\mathbb P_1\otimes\mathbb P_2\otimes\mathbb P_3$, where $\mathbb P_l$ is the law of $u_l$ under $\mathbb P$, a probability measure on $\mathcal U$. Fix $u_2,u_3\in\mathcal U$, and for $u_1\in\mathcal U$ write $u=(u_1,u_2,u_3)$. For $v\in\mathcal U$ let $\check v\in\mathcal U$ be the class of the function $s\mapsto v(-s)$, and put
\[
u^{\theta_1}=\left(\check u_1,u_2,u_3\right);
\]
this is the triple $u^w$ of Section~\ref{sec:free-maxwell-law} for $w=\theta_1$.

\emph{Two reflection positive systems.} Let $\vartheta_1\colon\mathcal U\to\mathcal U$ be the map $v\mapsto\check v$, and let $\mathcal F_{1+}=\sigma\left(u|_{(0,\infty)}\right)$ be the $\sigma$-algebra on $\mathcal U$ defined in Section~\ref{sec:free-maxwell-law} for the interval $(0,\infty)$. The involution $\vartheta_1$ is measurable, since $\int\mathrm ds\,\phi(s)\cdot\check v(s)=\int\mathrm ds\,\phi(-s)\cdot v(s)$ for $\phi\in C_c(\mathbb R)$. It preserves $\mathbb P_1$, and
\[
\int_{\mathcal U}\mathbb P_1(\mathrm dv)\,\overline{G(\check v)}\cdot G(v)\ge0\qquad\text{for every bounded $\mathcal F_{1+}$-measurable function }G\colon\mathcal U\to\mathbb C.
\]
Indeed, if $\mathbb P_1$ is the point mass at a constant function $v_0\in\mathcal U$, then $\check v_0=v_0$, and the integral equals $\left|G(v_0)\right|^2$. If $\mathbb P_1$ is the law of the class $v$ of $\frac{1}{4}\cdot(1+\varsigma)$, then $\check v$ is the class of $\frac{1}{4}\cdot(1+\check\varsigma)$. Let $G\colon\mathcal U\to\mathbb C$ be bounded and $\mathcal F_{1+}$-measurable. Since $\mathcal F_{1+}$ is generated by the maps $v'\mapsto\int\mathrm ds\,\phi(s)\cdot v'(s)$, $\phi\in C_c((0,\infty))$, the function $G$ is a bounded measurable function of these integrals (Doob--Dynkin lemma). For $\phi\in C_c((0,\infty))$ one has
\[
\int\mathrm ds\,\phi(s)\cdot v(s)=\frac{1}{4}\cdot\int\mathrm ds\,\phi(s)+\frac{1}{4}\cdot\int\mathrm ds\,\phi(s)\cdot\varsigma(s),
\]
and the same identity with $\check v$ and $\check\varsigma$ in place of $v$ and $\varsigma$. Hence a single function $G'$ of the kind admitted in~\eqref{eq:ergodic-components-7} satisfies both $G(v)=G'(\varsigma)$ and $G(\check v)=G'(\check\varsigma)$, and the integral $\int_{\mathcal U}\mathbb P_1(\mathrm dv)\,\overline{G(\check v)}\cdot G(v)$ equals $\mathbb E\left[\overline{G'(\check\varsigma)}\cdot G'(\varsigma)\right]$, which is non-negative by~\eqref{eq:ergodic-components-7}. Moreover $\check v$ has the law of $v$, that is, $\vartheta_1$ preserves $\mathbb P_1$, by~\eqref{eq:ergodic-components-4} in the form which this property takes for the classes (Section~\ref{sec:free-maxwell-law}).

Let $\gamma$ be the free Maxwell law of Lemma~\ref{lem:properties-free-maxwell}, let $\vartheta_2\colon Q\to Q$ be the map $\omega\mapsto\theta_1\omega$, and let $\mathcal F_{2+}=\Sigma_1^+$. By Lemma~\ref{lem:properties-free-maxwell}(i), $\gamma$ is invariant under all Euclidean maps, in particular under $\theta_1$, and satisfies $(RP_1)$. Hence the measurable involution $\vartheta_2$ preserves $\gamma$, and for every bounded $\Sigma_1^+$-measurable function $G$ on $Q$ one has $\Theta_1G=G\circ\theta_1$ and, by Lemma~\ref{lem:reformulations-reflection-positivity}(3),
\[
\int\mathrm d\gamma\,\overline{G\circ\theta_1}\cdot G=\left\langle\Theta_1G,G\right\rangle_\gamma\ge0.
\]
Thus the probability spaces $(\mathcal U,\mathbb P_1)$ and $(Q,\Sigma,\gamma)$, with the involutions $\vartheta_1$ and $\vartheta_2$ and the $\sigma$-algebras $\mathcal F_{1+}$ and $\mathcal F_{2+}$, satisfy the hypotheses of Lemma~\ref{lem:reflection-positivity-products}, which gives
\[
\int\left(\mathbb P_1\otimes\gamma\right)\left(\mathrm d(v,\omega)\right)\,\overline{G\left(\check v,\theta_1\omega\right)}\cdot G(v,\omega)\ge0
\]
for every bounded $\mathcal F_{1+}\otimes\mathcal F_{2+}$-measurable function $G$ on $\mathcal U\times Q$. Moreover, the map $\vartheta_1\times\vartheta_2$, $(v,\omega)\mapsto\left(\check v,\theta_1\omega\right)$, preserves $\mathbb P_1\otimes\gamma$, since $\vartheta_1$ preserves $\mathbb P_1$ and $\vartheta_2$ preserves $\gamma$.

\emph{The fields $Y_f$.} Let
\[
L^2_+=\left\{g\in L^2(\mathbb R^3;\Lambda^2):g=0\text{ almost everywhere outside }\{x_1>0\}\right\},
\]
and let $(h_k)_{k\ge1}$ be an orthonormal basis of $L^2_+$ with $h_k\in\mathcal D(\{x_1>0\})$ for every $k$; such a basis is obtained by the Gram--Schmidt procedure from a sequence in $\mathcal D(\{x_1>0\})$ which is dense in $L^2_+$. The law $\gamma$ satisfies~\eqref{eq:setup-G} by Lemma~\ref{lem:properties-free-maxwell}(i). By Lemma~\ref{lem:consequences-gaussian-bound}(i) and (ii), $\gamma\left(X(g)^2\right)\le\|g\|^2$ for $g\in\mathcal S$, and $X$ extends to a linear map $X\colon L^2(\mathbb R^3;\Lambda^2)\to L^2(\gamma)$ of norm at most $1$.

Let $f\in\mathcal D(\{x_1>0\})$, and let $r>0$ be such that $\operatorname{supp}f\subset[-r,r]^3$. By Lemma~\ref{lem:coordinatewise-changes-variables}(i), $\operatorname{supp}\Phi_{u\#}f=\Phi_u(\operatorname{supp}f)$ is a compact subset of $\{x_1>0\}$, so that $\Phi_{u\#}f\in L^2_+$; the same holds with $u^{\theta_1}$ in place of $u$. By Lemma~\ref{lem:coordinatewise-changes-variables}(v), $\Phi_{u\#}f$ depends on $u_1$ only through the restriction $u_1|_{(0,r]}$, and it is a continuous function of this restriction for the norm of $L^1((0,r))$. The reasoning given above for $L^1([-r,r])$, carried out on $L^1((0,r))$ with the functionals $v\mapsto\int\mathrm ds\,\phi(s)\cdot v(s)$, $\phi\in C_c((0,r))$, shows that a map into $L^1((0,r))$ is Borel measurable as soon as its compositions with these functionals are measurable. Composed with the map $u_1\mapsto u_1|_{(0,r)}$, the functional belonging to $\phi\in C_c((0,r))$ becomes $u_1\mapsto\int\mathrm ds\,\phi(s)\cdot u_1(s)$, which is $\mathcal F_{1+}$-measurable because $C_c((0,r))\subset C_c((0,\infty))$; so $u_1\mapsto u_1|_{(0,r)}$ is measurable for $\mathcal F_{1+}$ and the Borel sets of $L^1((0,r))$. Hence the coefficients
\[
\kappa_k(u_1)=\left\langle h_k,\Phi_{u\#}f\right\rangle\qquad\left(k\ge1\right)
\]
are $\mathcal F_{1+}$-measurable functions of $u_1\in\mathcal U$. For $N\ge1$ put
\[
Y^N_f(u_1,\omega)=\sum_{k=1}^N\kappa_k(u_1)\cdot X(h_k)(\omega)\qquad\left((u_1,\omega)\in\mathcal U\times Q\right).
\]
Each $X(h_k)$ is $\Sigma_1^+$-measurable, as $h_k\in\mathcal D(\{x_1>0\})$; so $Y^N_f$ is a real $\mathcal F_{1+}\otimes\mathcal F_{2+}$-measurable function, and $Y^N_f(u_1,\cdot)=X\left(\sum_{k\le N}\kappa_k(u_1)\cdot h_k\right)$. For $1\le M<N$ and $u_1\in\mathcal U$,
\[
\int\gamma(\mathrm d\omega)\,\left|Y^N_f(u_1,\omega)-Y^M_f(u_1,\omega)\right|^2\le\left\|\sum_{k=M+1}^N\kappa_k(u_1)\cdot h_k\right\|^2\le\sum_{k>M}\kappa_k(u_1)^2.
\]
By Bessel's inequality and Lemma~\ref{lem:coordinatewise-changes-variables}(i), the right side is at most $\left\|\Phi_{u\#}f\right\|^2\le\|f\|^2$, and it tends to $0$ as $M\to\infty$. By Fubini's theorem and dominated convergence in $u_1$, $\left(Y^N_f\right)_{N\ge1}$ is a Cauchy sequence in $L^2\left(\mathbb P_1\otimes\gamma\right)$. Let $Y_f$ be a real $\mathcal F_{1+}\otimes\mathcal F_{2+}$-measurable function on $\mathcal U\times Q$ such that $Y^N_f\to Y_f$ in $L^2\left(\mathbb P_1\otimes\gamma\right)$; for instance the limit of a subsequence which converges $\mathbb P_1\otimes\gamma$-almost everywhere, put equal to $0$ where this subsequence does not converge.

For every $u_1\in\mathcal U$ one has $\sum_{k\le N}\kappa_k(u_1)\cdot h_k\to\Phi_{u\#}f$ in $L^2(\mathbb R^3;\Lambda^2)$ as $N\to\infty$, because $\Phi_{u\#}f\in L^2_+$; hence $Y^N_f(u_1,\cdot)\to X\left(\Phi_{u\#}f\right)$ in $L^2(\gamma)$, where $X\left(\Phi_{u\#}f\right)$ is the value of the extension of Lemma~\ref{lem:consequences-gaussian-bound}(ii). On the other hand, Fubini's theorem gives $\int\mathbb P_1(\mathrm du_1)\,\left\|Y^N_f(u_1,\cdot)-Y_f(u_1,\cdot)\right\|^2_{L^2(\gamma)}\to0$, so that, along a subsequence, $Y^N_f(u_1,\cdot)\to Y_f(u_1,\cdot)$ in $L^2(\gamma)$ for $\mathbb P_1$-almost every $u_1$. Consequently
\[
Y_f(u_1,\cdot)=X\left(\Phi_{u\#}f\right)\quad\text{in }L^2(\gamma),\qquad\text{for $\mathbb P_1$-almost every }u_1\in\mathcal U.
\]

\emph{The reflected fields.} Let $u_1\in\mathcal U$, and let $u=(u_1,u_2,u_3)$ and $u^{\theta_1}=(\check u_1,u_2,u_3)$ as above. Section~\ref{sec:euclidean-maps-actions} gives $X(h)(g\omega)=X(g^{-1}\cdot h)(\omega)$ for every Euclidean map $g$, and $\theta_1^{-1}=\theta_1$; thus
\[
X(h)\left(\theta_1\omega\right)=X\left(\theta_1\cdot h\right)(\omega)\qquad\left(h\in\mathcal S,\ \omega\in Q\right).
\]
Replacing $u_1$ by $\check u_1$ in the definition of $\kappa_k$ replaces the triple $u$ by $u^{\theta_1}$, so that $\kappa_k(\check u_1)=\left\langle h_k,\Phi_{u^{\theta_1}\#}f\right\rangle$. Put
\[
g_N=\sum_{k=1}^N\kappa_k(\check u_1)\cdot h_k=\sum_{k=1}^N\left\langle h_k,\Phi_{u^{\theta_1}\#}f\right\rangle\cdot h_k\qquad\left(N\ge1\right).
\]
Applying the identity $X(h)\left(\theta_1\omega\right)=X\left(\theta_1\cdot h\right)(\omega)$ with $h=h_k$, $k\le N$, and using that $h\mapsto X(\theta_1\cdot h)(\omega)$ is linear on $\mathcal S$, one obtains from the definition of $Y^N_f$ that $Y^N_f\left(\check u_1,\theta_1\omega\right)=X\left(\theta_1\cdot g_N\right)(\omega)$ for every $\omega\in Q$ and every $N\ge1$.

Now let $N\to\infty$ on both sides. The functions $g_N$ are the partial sums of the expansion of $\Phi_{u^{\theta_1}\#}f\in L^2_+$ in the orthonormal basis $(h_k)$ of $L^2_+$, so $g_N\to\Phi_{u^{\theta_1}\#}f$ in $L^2(\mathbb R^3;\Lambda^2)$. The reflection $\theta_1$ acts isometrically on $L^2(\mathbb R^3;\Lambda^2)$, and the extended map $X$ has norm at most $1$; consequently $X(\theta_1\cdot g_N)\to X\left(\theta_1\cdot\Phi_{u^{\theta_1}\#}f\right)$ in $L^2(\gamma)$, and this holds for every $u_1\in\mathcal U$. On the other side, composition with the map $\vartheta_1\times\vartheta_2$ is an isometry of $L^2\left(\mathbb P_1\otimes\gamma\right)$, this map being measure preserving. It turns the convergence $Y^N_f\to Y_f$ into the convergence, in $L^2\left(\mathbb P_1\otimes\gamma\right)$, of the functions $(v,\omega)\mapsto Y^N_f\left(\check v,\theta_1\omega\right)$ to the function $(v,\omega)\mapsto Y_f\left(\check v,\theta_1\omega\right)$. Fubini's theorem then provides indices $N_1<N_2<\dots$ such that, for $\mathbb P_1$-almost every $u_1$, $Y^{N_m}_f\left(\check u_1,\theta_1\,\cdot\,\right)\to Y_f\left(\check u_1,\theta_1\,\cdot\,\right)$ in $L^2(\gamma)$ as $m\to\infty$. For every such $u_1$, letting $m\to\infty$ in the identity $Y^{N_m}_f\left(\check u_1,\theta_1\,\cdot\,\right)=X\left(\theta_1\cdot g_{N_m}\right)$ yields $Y_f\left(\check u_1,\theta_1\,\cdot\,\right)=X\left(\theta_1\cdot\Phi_{u^{\theta_1}\#}f\right)$ in $L^2(\gamma)$. Finally, Lemma~\ref{lem:coordinatewise-changes-variables}(iii), applied with $w=\theta_1$, for which $u^w=u^{\theta_1}$, gives $\theta_1\cdot\Phi_{u^{\theta_1}\#}f=\Phi_{u\#}(\theta_1\cdot f)$. Altogether,
\[
Y_f\left(\check u_1,\theta_1\,\cdot\,\right)=X\left(\theta_1\cdot\Phi_{u^{\theta_1}\#}f\right)=X\left(\Phi_{u\#}(\theta_1\cdot f)\right)\quad\text{in }L^2(\gamma),\qquad\text{for $\mathbb P_1$-almost every }u_1\in\mathcal U.
\]

\emph{The characteristic functional of $\gamma$ on $L^2(\mathbb R^3;\Lambda^2)$.} For every $g\in L^2(\mathbb R^3;\Lambda^2)$,
\[
\gamma\left(\exp\left(iX(g)\right)\right)=\exp\left(-\frac{1}{2}\cdot\left\|\Pi g\right\|^2\right).
\]
For $g\in\mathcal S$ this is the definition of $\gamma$ in Lemma~\ref{lem:properties-free-maxwell}. For general $g$ let $g_\ell\in\mathcal S$, $\ell\ge1$, satisfy $\|g_\ell-g\|\to0$. Since $\left|\exp(ia)-\exp(ib)\right|\le|a-b|$ for real $a,b$, the Cauchy--Schwarz inequality and Lemma~\ref{lem:consequences-gaussian-bound}(ii) give
\[
\left|\gamma\left(\exp\left(iX(g_\ell)\right)\right)-\gamma\left(\exp\left(iX(g)\right)\right)\right|\le\gamma\left(\left|X(g_\ell)-X(g)\right|\right)\le\|g_\ell-g\|\to0,
\]
while $\|\Pi g_\ell\|\to\|\Pi g\|$.

\emph{The inequality $(RP_1)$.} Let $n\ge1$, $c=(c_1,\dots,c_n)\in\mathbb C^n$ and $f_1,\dots,f_n\in\mathcal D(\{x_1>0\})$, and let
\[
G=\sum_{p=1}^nc_p\cdot\exp\left(iY_{f_p}\right),
\]
a bounded $\mathcal F_{1+}\otimes\mathcal F_{2+}$-measurable function on $\mathcal U\times Q$. By the two identities for $Y_{f_q}(u_1,\cdot)$ and $Y_{f_p}\left(\check u_1,\theta_1\,\cdot\,\right)$ proved above, and by the linearity of $\Phi_{u\#}$ and of the extended map $X$, one has for $\mathbb P_1$-almost every $u_1$ and $\gamma$-almost every $\omega$
\[
\overline{G\left(\check u_1,\theta_1\omega\right)}\cdot G(u_1,\omega)=\sum_{p,q=1}^n\overline{c_p}\cdot c_q\cdot\exp\left(iX\left(\Phi_{u\#}(f_q-\theta_1\cdot f_p)\right)(\omega)\right).
\]
The inequality derived from Lemma~\ref{lem:reflection-positivity-products}, Fubini's theorem and the characteristic functional of $\gamma$ on $L^2(\mathbb R^3;\Lambda^2)$ now give, with $u=(u_1,u_2,u_3)$,
\begin{align*}
0&\le\int\left(\mathbb P_1\otimes\gamma\right)\left(\mathrm d(u_1,\omega)\right)\,\overline{G\left(\check u_1,\theta_1\omega\right)}\cdot G(u_1,\omega)\\
&=\sum_{p,q=1}^n\overline{c_p}\cdot c_q\cdot\int_{\mathcal U}\mathbb P_1(\mathrm du_1)\,\exp\left(-\frac{1}{2}\cdot q_u\left(f_q-\theta_1\cdot f_p\right)\right).
\end{align*}
This holds for all $u_2,u_3\in\mathcal U$. Integration with respect to $\mathbb P_2\otimes\mathbb P_3$, Fubini's theorem and $\mathbb P=\mathbb P_1\otimes\mathbb P_2\otimes\mathbb P_3$ give
\[
\sum_{p,q=1}^n\overline{c_p}\cdot c_q\cdot\mathbb E\exp\left(-\frac{1}{2}\cdot q_u\left(f_q-\theta_1\cdot f_p\right)\right)\ge0.
\]
Since $f_q-\theta_1\cdot f_p\in\mathcal D$ and $X(f_q-\theta_1\cdot f_p)=X(f_q)-X(\theta_1\cdot f_p)$, the identity defining $m_{\mathbb P}$ turns this inequality into
\[
\sum_{p,q=1}^n\overline{c_p}\cdot c_q\cdot m_{\mathbb P}\left(\exp\left(i\left(X(f_q)-X(\theta_1\cdot f_p)\right)\right)\right)\ge0,
\]
which is $(RP_1)$.

\emph{The directions $j=2,3$.} For $j\in\{2,3\}$ the same argument applies with the coordinate $x_j$ in place of $x_1$: the law $\mathbb P_j$ of $u_j$, the reflection $\theta_j$, the $\sigma$-algebra $\Sigma_j^+$, the half space $\{x_j>0\}$, and the triple $u^{\theta_j}$ obtained from $u$ by replacing $u_j$ with $\check u_j$, take the places of $\mathbb P_1$, $\theta_1$, $\Sigma_1^+$, $\{x_1>0\}$ and $u^{\theta_1}$; Lemma~\ref{lem:coordinatewise-changes-variables}(iii) is applied with $w=\theta_j$, and Lemmas~\ref{lem:properties-free-maxwell}(i) and~\ref{lem:reformulations-reflection-positivity}(3) are applied to the direction $j$.
\end{proof}

\subsubsection{Permuted diagonal rescalings of the free Maxwell law}\label{sec:permuted-diagonal-rescalings}

\begin{example}[A mixture of six rescaled free Maxwell laws]\label{ex:mixture-six-rescaled}
Let $\lambda=(\lambda_1,\lambda_2,\lambda_3)=\left(\frac{1}{4},\frac{1}{3},\frac{1}{2}\right)$. For $\pi\in S_3$ let $u^{(\pi)}\in\mathcal U^3$ be the triple of the constant functions
\[
u^{(\pi)}_l\equiv\lambda_{\pi(l)}\qquad\left(l=1,2,3\right),
\]
let $\delta_{u^{(\pi)}}$ be the point mass at $u^{(\pi)}$, and let $\gamma^{(\pi)}=m_{\delta_{u^{(\pi)}}}$ be the law of Lemma~\ref{lem:mixtures-changes-variables} for $\mathbb P=\delta_{u^{(\pi)}}$; thus
\[
\gamma^{(\pi)}\left(\exp\left(iX(f)\right)\right)=\exp\left(-\frac{1}{2}\cdot q_{u^{(\pi)}}(f)\right)\qquad\left(f\in\mathcal S\right).
\]
For $u=u^{(\pi)}$ one has $\Phi_u=\operatorname{diag}\left(\lambda_{\pi(1)},\lambda_{\pi(2)},\lambda_{\pi(3)}\right)$, a linear bijection of $\mathbb R^3$. Let $T_\pi\colon Q\to Q$ be the map with $(T_\pi\omega)(f)=\omega\left(\Phi_{u\#}f\right)$ for $f\in\mathcal S$; it is the pullback of two-forms by $\Phi_u$, extended to distributions. The image of the free Maxwell law $\gamma$ of Lemma~\ref{lem:properties-free-maxwell} under $T_\pi$ is $\gamma^{(\pi)}$. Put
\[
m_A=\frac{1}{6}\cdot\sum_{\pi\in S_3}\gamma^{(\pi)},
\]
and for $w\in O_{\mathrm{cub}}$ let $\pi_w\in S_3$ be its permutation part, $w\,e_i=\varepsilon_i\cdot e_{\pi_w(i)}$ for $i=1,2,3$, as in Section~\ref{sec:average-six-anisotropic}. Then the following assertions hold.
\begin{enumerate}
\item[(1)] Let $\pi\in S_3$. The law $\gamma^{(\pi)}$ is a centred Gaussian law: for every $f\in\mathcal S$, the function $X(f)$ has under $\gamma^{(\pi)}$ the centred normal law with variance $q_{u^{(\pi)}}(f)$. The law $\gamma^{(\pi)}$ satisfies~\eqref{eq:setup-G}, is concentrated on $\{\omega\in Q:d\omega=0\}$, satisfies $(T_{123})$, is invariant under every element of $N_3$, and satisfies $(RP_j)$ for $j=1,2,3$. It is mixing, and $\gamma^{(\pi)}\in\mathcal M^{123}_{\mathrm{erg}}$.
\item[(2)] For $w\in O_{\mathrm{cub}}$ and $\pi\in S_3$,
\[
\left(u^{(\pi)}\right)^w=u^{(\pi\pi_w)},\qquad(w^{-1})_*\gamma^{(\pi)}=\gamma^{(\pi\pi_w)}.
\]
Hence every $w\in O_{\mathrm{cub}}$ permutes the six laws $\gamma^{(\pi)}$, $\pi\in S_3$. The law $m_A$ satisfies~\eqref{eq:setup-H} and is concentrated on $\{\omega\in Q:d\omega=0\}$.
\item[(3)] The six laws $\gamma^{(\pi)}$, $\pi\in S_3$, are pairwise distinct.
\item[(4)] Let $\nu$ be the probability measure on $(\mathcal M,\mathcal B_{\mathcal M})$ which Theorem~\ref{thm:ergodic-decomposition}(a) assigns to $m=m_A$ and $J=\{1,2,3\}$, and for $\rho\in\mathcal M$ let $\delta_\rho$ be the point mass at $\rho$ on $(\mathcal M,\mathcal B_{\mathcal M})$. Then
\[
\nu=\frac{1}{6}\cdot\sum_{\pi\in S_3}\delta_{\gamma^{(\pi)}}.
\]
For $w\in O_{\mathrm{cub}}$ and $\pi,\pi'\in S_3$ one has $w_*\gamma^{(\pi)}=\gamma^{(\pi')}$ if and only if $\pi_w=\pi'^{-1}\pi$, and the set of these $w$ is a coset of $N_3$. In particular, $w_*\gamma^{(\pi)}=\gamma^{(\pi)}$ if and only if $\pi_w$ is the identity permutation, that is, if and only if $w\in N_3$.
\end{enumerate}
\end{example}

\begin{proof}[Proof of the assertions of Example~\ref{ex:mixture-six-rescaled}]
\emph{The laws $\gamma^{(\pi)}$ as pullbacks.} Let $\pi\in S_3$ and $u=u^{(\pi)}$. Then $\varphi_l(s)=\lambda_{\pi(l)}\cdot s$ for $s\in\mathbb R$ and $l=1,2,3$, so that $\Phi_u=\operatorname{diag}\left(\lambda_{\pi(1)},\lambda_{\pi(2)},\lambda_{\pi(3)}\right)$. By~\eqref{eq:ergodic-components-8}, the components of $\Phi_{u\#}f$ are those of $f\circ\Phi_u^{-1}$, multiplied by constants. So $\Phi_{u\#}$ is a continuous linear map of $\mathcal S$ into itself, $T_\pi\omega$ belongs to $Q$ for $\omega\in Q$, and $T_\pi$ is measurable, because $X(f)\circ T_\pi=X\left(\Phi_{u\#}f\right)$ for $f\in\mathcal S$. Since $\Phi_{u\#}$ is the transpose of the pullback of two-forms by $\Phi_u$ (Section~\ref{sec:free-maxwell-law}), $T_\pi$ is the extension of this pullback to distributions. The image of $\gamma$ under $T_\pi$ has the characteristic functional
\[
f\mapsto\gamma\left(\exp\left(iX\left(\Phi_{u\#}f\right)\right)\right)=\exp\left(-\frac{1}{2}\cdot\left\|\Pi\Phi_{u\#}f\right\|^2\right)=\exp\left(-\frac{1}{2}\cdot q_u(f)\right),
\]
which is that of $\gamma^{(\pi)}$; by Lemma~\ref{lem:density-exponentials}(ii) this image is $\gamma^{(\pi)}$.

\emph{Assertion (1).} Let $\pi\in S_3$ and $u=u^{(\pi)}$. For $f\in\mathcal S$ and $t\in\mathbb R$ the identity defining $\gamma^{(\pi)}$, applied to $tf$, gives $\gamma^{(\pi)}\left(\exp\left(i\,t\cdot X(f)\right)\right)=\exp\left(-\frac{1}{2}\cdot t^2\cdot q_u(f)\right)$; by the uniqueness theorem for characteristic functions, $X(f)$ has under $\gamma^{(\pi)}$ the centred normal law with variance $q_u(f)$. Since $u$ is constant, $u^a=u$ for every $a\in\mathbb R^3$ and $u^\varepsilon=u$ for every $\varepsilon\in N_3$; so $\delta_u$ is preserved by the maps $v\mapsto v^a$ and $v\mapsto v^\varepsilon$ of $\mathcal U^3$, and under $\delta_u$ the three coordinates are deterministic constants, hence independent. By Lemma~\ref{lem:mixtures-changes-variables}(i), (ii) and (iii), $\gamma^{(\pi)}$ satisfies~\eqref{eq:setup-G}, is concentrated on $\{\omega\in Q:d\omega=0\}$, satisfies $(T_{123})$, is invariant under every element of $N_3$, and satisfies $(RP_j)$ for $j=1,2,3$.

Let $f,g\in\mathcal D$, and let $r>0$ be such that $\operatorname{supp}f\cup\operatorname{supp}g\subset[-r,r]^3$. For $a\in\mathbb R^3$, Lemma~\ref{lem:coordinatewise-changes-variables}(ii) and $u^a=u$ give $\Phi_{u\#}(\tau_a\cdot g)=\tau_{\Phi_u(a)}\cdot\Phi_{u\#}g$. As $\Pi$ is a self-adjoint projection which commutes with the translations, and the translations preserve the norm, it follows that
\[
q_u(f+\tau_a\cdot g)=q_u(f)+q_u(g)+2\left\langle\Phi_{u\#}f,\Pi\left(\tau_{\Phi_u(a)}\cdot\Phi_{u\#}g\right)\right\rangle.
\]
By Lemma~\ref{lem:coordinatewise-changes-variables}(i), the supports of $\Phi_{u\#}f$ and $\Phi_{u\#}g$ lie in $\Phi_u\left([-r,r]^3\right)\subset\left[-\frac{r}{2},\frac{r}{2}\right]^3$, their $L^1$ norms are at most $\frac{1}{2}\cdot\|f\|_{L^1}$ and $\frac{1}{2}\cdot\|g\|_{L^1}$, and $\left|\Phi_u(a)\right|\ge\frac{1}{4}\cdot|a|$. Hence, if $|a|>4\sqrt3\cdot r$, the supports of $\Phi_{u\#}f$ and of $\tau_{\Phi_u(a)}\cdot\Phi_{u\#}g$ have distance at least $\delta_a=\frac{1}{4}\cdot|a|-\sqrt3\cdot r>0$, and Lemma~\ref{lem:properties-free-maxwell}(ii) gives
\[
\left|\left\langle\Phi_{u\#}f,\Pi\left(\tau_{\Phi_u(a)}\cdot\Phi_{u\#}g\right)\right\rangle\right|\le\frac{\|f\|_{L^1}\cdot\|g\|_{L^1}}{8\pi\cdot\delta_a^3}\to0\qquad\left(|a|\to\infty\right).
\]
Let $F=\exp\left(iX(f)\right)$ and $G=\exp\left(iX(g)\right)$. Then $\overline F\cdot U(a)G=\exp\left(iX(-f+\tau_a\cdot g)\right)$ by Section~\ref{sec:euclidean-maps-actions}; the numbers $\gamma^{(\pi)}(F)=\exp\left(-\frac{1}{2}\cdot q_u(f)\right)$ and $\gamma^{(\pi)}(G)$ are real; and the two preceding displays, applied with $-f$ in place of $f$, give, as $|a|\to\infty$,
\begin{align*}
\gamma^{(\pi)}\left(\overline F\cdot U(a)G\right)&=\exp\left(-\frac{1}{2}\cdot q_u(-f+\tau_a\cdot g)\right)\\
&\to\exp\left(-\frac{1}{2}\cdot q_u(f)\right)\cdot\exp\left(-\frac{1}{2}\cdot q_u(g)\right)=\overline{\gamma^{(\pi)}(F)}\cdot\gamma^{(\pi)}(G).
\end{align*}
By linearity, $\gamma^{(\pi)}\left(\overline F\cdot U(a)G\right)\to\overline{\gamma^{(\pi)}(F)}\cdot\gamma^{(\pi)}(G)$ as $|a|\to\infty$ for all $F,G$ in the span $\mathcal A$ of Lemma~\ref{lem:density-exponentials}. Thus $\gamma^{(\pi)}$ is mixing (Section~\ref{sec:invariant-sets-invariant}), and Lemma~\ref{lem:clustering-criterion-separation}(a), applied with $J=\{1,2,3\}$ and $a^{(k)}=k\cdot e_1$, gives $\gamma^{(\pi)}\in\mathcal M^{123}_{\mathrm{erg}}$.

\emph{Assertion (2).} Let $w\in O_{\mathrm{cub}}$ and $\pi\in S_3$. By the definition of $u^w$ in Section~\ref{sec:free-maxwell-law}, $\left(u^{(\pi)}\right)^w_i\equiv u^{(\pi)}_{\pi_w(i)}\equiv\lambda_{\pi(\pi_w(i))}$ for $i=1,2,3$, that is, $\left(u^{(\pi)}\right)^w=u^{(\pi\pi_w)}$. For $\rho\in\mathcal M$ and $f\in\mathcal S$ one has $\left((w^{-1})_*\rho\right)\left(\exp\left(iX(f)\right)\right)=\rho\left(\exp\left(iX(w\cdot f)\right)\right)$ (Section~\ref{sec:euclidean-maps-actions}). With the identity $q_u(w\cdot f)=q_{u^w}(f)$ of the proof of Lemma~\ref{lem:mixtures-changes-variables}(ii) it follows that, for $f\in\mathcal S$,
\begin{align*}
\left((w^{-1})_*\gamma^{(\pi)}\right)\left(\exp\left(iX(f)\right)\right)&=\exp\left(-\frac{1}{2}\cdot q_{u^{(\pi)}}(w\cdot f)\right)=\exp\left(-\frac{1}{2}\cdot q_{u^{(\pi\pi_w)}}(f)\right)\\
&=\gamma^{(\pi\pi_w)}\left(\exp\left(iX(f)\right)\right),
\end{align*}
and Lemma~\ref{lem:density-exponentials}(ii) gives $(w^{-1})_*\gamma^{(\pi)}=\gamma^{(\pi\pi_w)}$. As $\pi\mapsto\pi\pi_w$ is a bijection of $S_3$ onto itself, the map $\rho\mapsto(w^{-1})_*\rho$ permutes the six laws $\gamma^{(\pi)}$, $\pi\in S_3$, and $(w^{-1})_*m_A=m_A$. Since $w\in O_{\mathrm{cub}}$ is arbitrary, every $w\in O_{\mathrm{cub}}$ permutes the six laws, and $w_*m_A=m_A$. The bound~\eqref{eq:setup-G}, the property $(T_{123})$, the properties $(RP_j)$, $j=1,2,3$, and the property $\rho\left(\{\omega\in Q:d\omega=0\}\right)=1$ are linear equalities or inequalities in the law $\rho$; they hold for each $\gamma^{(\pi)}$ by (1), hence for the convex combination $m_A$. Therefore $m_A$ satisfies~\eqref{eq:setup-H} and is concentrated on $\{\omega\in Q:d\omega=0\}$.

\emph{Assertion (3).} With the normalisation of the Fourier transform in Section~\ref{sec:free-maxwell-law}, Plancherel's theorem reads $\|b\|^2=(2\pi)^{-3}\cdot\int\mathrm dp\,\left|\hat b(p)\right|^2$ for $b\in L^2(\mathbb R^3;\mathbb R^3)$, and the same normalisation is used for scalar functions. Let $\lambda'=(\lambda'_1,\lambda'_2,\lambda'_3)$ be a permutation of $\lambda$, let $u\in\mathcal U^3$ be the constant triple $u\equiv\lambda'$, and let $f\in\mathcal D$ with $\ast f=b=(b_1,0,0)$. By~\eqref{eq:ergodic-components-8}, the second and the third component of $\Phi_{u\#}b$ vanish, and
\[
\left(\Phi_{u\#}b\right)_1(y)=\frac{1}{\lambda'_1}\cdot b_1\left(\frac{y_1}{\lambda'_1},\frac{y_2}{\lambda'_2},\frac{y_3}{\lambda'_3}\right)\qquad\left(y\in\mathbb R^3\right);
\]
the substitution $y_l=\lambda'_l\cdot x_l$, $l=1,2,3$, gives
\[
\widehat{\left(\Phi_{u\#}b\right)_1}(p)=\lambda'_2\cdot\lambda'_3\cdot\hat b_1\left(\lambda'_1\cdot p_1,\lambda'_2\cdot p_2,\lambda'_3\cdot p_3\right)\qquad\left(p\in\mathbb R^3\right).
\]
The matrix $1-|p|^{-2}\cdot p\,p^T$ is an orthogonal projection whose entry at the place $(1,1)$ is $1-|p|^{-2}\cdot p_1^2$. Hence Plancherel's theorem and the substitution $k_l=\lambda'_l\cdot p_l$, $l=1,2,3$, give
\[
q_u(f)=(2\pi)^{-3}\cdot\int\mathrm dk\,\left|\hat b_1(k)\right|^2\cdot M_{\lambda'}(k),\qquad M_{\lambda'}(k)=\frac{\lambda'_2\cdot\lambda'_3}{\lambda'_1}\cdot\left(1-\frac{\lambda'^{-2}_1\cdot k_1^2}{\sum_{l=1}^3\lambda'^{-2}_l\cdot k_l^2}\right)\quad\left(k\ne0\right).
\]

Let $\lambda'$ and $\lambda''$ be permutations of $\lambda$ for which the corresponding laws coincide, that is, $m_{\delta_{u'}}=m_{\delta_{u''}}$ for the constant triples $u'\equiv\lambda'$ and $u''\equiv\lambda''$. Then $\exp\left(-\frac{1}{2}\cdot q_{u'}(f)\right)=\exp\left(-\frac{1}{2}\cdot q_{u''}(f)\right)$, hence $q_{u'}(f)=q_{u''}(f)$, for every $f\in\mathcal S$. The function $M_{\lambda'}$ depends on $k$ only through $k_1^2,k_2^2,k_3^2$, so that $M_{\lambda'}(-k)=M_{\lambda'}(k)$; the denominator in its definition is positive for $k\ne0$, so that $M_{\lambda'}$ is continuous on $\mathbb R^3\setminus\{0\}$; and the quotient in its definition lies in $[0,1]$, so that $0\le M_{\lambda'}\le\frac{\lambda'_2\cdot\lambda'_3}{\lambda'_1}$. The same holds for $M_{\lambda''}$. Let $k_0\in\mathbb R^3\setminus\{0\}$, let $\chi\in C_c^\infty(\mathbb R^3)$ be real and not identically zero, and for $R>0$ let $f_R\in\mathcal D$ be given by $\ast f_R=(b_{1,R},0,0)$ with
\[
b_{1,R}(x)=\chi\left(\frac{x}{R}\right)\cdot\cos\left(k_0\cdot x\right)\qquad\left(x\in\mathbb R^3\right).
\]
Then
\[
\hat b_{1,R}(k)=\frac{R^3}{2}\cdot\left(\hat\chi\left(R\cdot(k-k_0)\right)+\hat\chi\left(R\cdot(k+k_0)\right)\right)\qquad\left(k\in\mathbb R^3\right).
\]
The function $b_{1,R}$ does not vanish identically, since $\chi$ is continuous and not identically zero and $\cos(k_0\cdot x)\ne0$ outside a null set; hence $\hat b_{1,R}\ne0$, and
\[
\nu_R(\mathrm dk)=\left\|\hat b_{1,R}\right\|_{L^2}^{-2}\cdot\left|\hat b_{1,R}(k)\right|^2\,\mathrm dk
\]
is a probability measure on $\mathbb R^3$. Let $\psi$ be bounded and continuous on $\mathbb R^3$. The substitution $k=\pm k_0+R^{-1}\cdot\xi$ and dominated convergence give
\[
R^3\cdot\int\mathrm dk\,\psi(k)\cdot\left|\hat\chi\left(R\cdot(k\mp k_0)\right)\right|^2\to\psi(\pm k_0)\cdot\int\mathrm d\xi\,\left|\hat\chi(\xi)\right|^2\qquad\left(R\to\infty\right),
\]
and the substitution $k=k_0+R^{-1}\cdot\xi$ gives
\[
R^3\cdot\left|\int\mathrm dk\,\psi(k)\cdot\hat\chi\left(R\cdot(k-k_0)\right)\cdot\overline{\hat\chi\left(R\cdot(k+k_0)\right)}\right|\le\sup|\psi|\cdot\int\mathrm d\xi\,\left|\hat\chi(\xi)\right|\cdot\left|\hat\chi\left(\xi+2R\cdot k_0\right)\right|,
\]
which tends to $0$ by dominated convergence, $\hat\chi$ being a Schwartz function and $k_0\ne0$. Hence $R^{-3}\cdot\int\mathrm dk\,\psi(k)\cdot\left|\hat b_{1,R}(k)\right|^2\to\frac{1}{4}\cdot\left(\psi(k_0)+\psi(-k_0)\right)\cdot\int\mathrm d\xi\,\left|\hat\chi(\xi)\right|^2$; for $\psi=1$ the limit is $\frac{1}{2}\cdot\int\mathrm d\xi\,\left|\hat\chi(\xi)\right|^2$, which is positive as $\chi\ne0$. Dividing the first limit by the second, one finds $\int\nu_R(\mathrm dk)\,\psi(k)\to\frac{1}{2}\cdot\left(\psi(k_0)+\psi(-k_0)\right)$ for every bounded continuous $\psi$; this is the weak convergence of $\nu_R$, as $R\to\infty$, to the probability measure $\frac{1}{2}\cdot\left(\delta_{k_0}+\delta_{-k_0}\right)$. Since $M_{\lambda'}$ is bounded and continuous outside $\{0\}$, a null set for the limit measure, the continuous mapping theorem gives $\int\nu_R(\mathrm dk)\,M_{\lambda'}(k)\to\frac{1}{2}\cdot\left(M_{\lambda'}(k_0)+M_{\lambda'}(-k_0)\right)=M_{\lambda'}(k_0)$. As $\left\|b_{1,R}\right\|^2=(2\pi)^{-3}\cdot\left\|\hat b_{1,R}\right\|_{L^2}^2$, this means
\[
\lim_{R\to\infty}\frac{q_{u'}(f_R)}{\left\|b_{1,R}\right\|^2}=M_{\lambda'}(k_0),
\]
and likewise for $\lambda''$. As $q_{u'}(f_R)=q_{u''}(f_R)$ for every $R$, one obtains $M_{\lambda'}(k_0)=M_{\lambda''}(k_0)$. The point $k_0\in\mathbb R^3\setminus\{0\}$ being arbitrary, $M_{\lambda'}=M_{\lambda''}$ on $\mathbb R^3\setminus\{0\}$.

Every permutation $\lambda'$ of $\lambda$ satisfies $\lambda'_1\cdot\lambda'_2\cdot\lambda'_3=\lambda_1\cdot\lambda_2\cdot\lambda_3$; hence
\[
M_{\lambda'}(e_2)=\frac{\lambda_1\cdot\lambda_2\cdot\lambda_3}{\lambda'^2_1},\qquad M_{\lambda'}(e_1+e_2)=\frac{\lambda_1\cdot\lambda_2\cdot\lambda_3}{\lambda'^2_1}\cdot\frac{\lambda'^{-2}_2}{\lambda'^{-2}_1+\lambda'^{-2}_2}.
\]
The equality $M_{\lambda'}(e_2)=M_{\lambda''}(e_2)$ gives $\lambda'_1=\lambda''_1$. Then the equality at $e_1+e_2$ gives $\lambda'_2=\lambda''_2$, because $t\mapsto t\cdot\left(\lambda'^{-2}_1+t\right)^{-1}$ is strictly increasing on $(0,\infty)$; so also $\lambda'_3=\lambda''_3$, and $\lambda'=\lambda''$. The entries of $\lambda$ are pairwise distinct, so distinct $\pi\in S_3$ give distinct triples $\left(\lambda_{\pi(1)},\lambda_{\pi(2)},\lambda_{\pi(3)}\right)$. Hence the six laws $\gamma^{(\pi)}$ are pairwise distinct.

\emph{Assertion (4).} Let
\[
\nu'=\frac{1}{6}\cdot\sum_{\pi\in S_3}\delta_{\gamma^{(\pi)}}.
\]
By (1), $\gamma^{(\pi)}\in\mathcal M^{123}_{\mathrm{erg}}$ for every $\pi\in S_3$. Hence $\nu'$ is a probability measure on $(\mathcal M,\mathcal B_{\mathcal M})$ with $m_A(A)=\int_{\mathcal M}\nu'(\mathrm d\rho)\,\rho(A)$ for every $A\in\Sigma$ and $\nu'\left(\mathcal M^{123}_{\mathrm{erg}}\right)=1$. By the uniqueness in Theorem~\ref{thm:ergodic-decomposition}(a), $\nu=\nu'$.

Let $w\in O_{\mathrm{cub}}$. As $w\mapsto\pi_w$ is a homomorphism of $O_{\mathrm{cub}}$ onto $S_3$ with kernel $N_3$ (Section~\ref{sec:euclidean-maps-actions}), assertion (2), applied to $w^{-1}$, gives $w_*\gamma^{(\pi)}=\gamma^{(\pi\pi_w^{-1})}$ for $\pi\in S_3$. By (3), for $\pi,\pi'\in S_3$ one has $w_*\gamma^{(\pi)}=\gamma^{(\pi')}$ if and only if $\pi\pi_w^{-1}=\pi'$, that is, if and only if $\pi_w=\pi'^{-1}\pi$; and the set of the $w\in O_{\mathrm{cub}}$ with $\pi_w=\pi'^{-1}\pi$ is a coset of the kernel $N_3$. For $\pi'=\pi$ the condition is that $\pi_w$ be the identity permutation, that is, that $w\in N_3$.

By (1)--(4), the law $m_A$ satisfies~\eqref{eq:setup-H}, is concentrated on $\{\omega\in Q:d\omega=0\}$, and has the properties stated in assertion~(A) of Theorem~\ref{thm:asymmetric-components-non}.
\end{proof}

\subsubsection{The free Maxwell law in a random product medium}\label{sec:free-maxwell-law-2}

\begin{example}[A mixture of free Maxwell laws over telegraph changes of variables]\label{ex:mixture-free-maxwell}
Let $\varsigma_1,\varsigma_2,\varsigma_3$ be independent telegraph processes (Section~\ref{sec:auxiliary-results}), and let
\[
u_l=\frac{1}{4}\cdot\left(1+\varsigma_l\right)\qquad(l=1,2,3).
\]
Each $u_l$ takes its values in $\left\{\frac{1}{4},\frac{1}{2}\right\}$, and its class is an element of the space $\mathcal U$ of Section~\ref{sec:free-maxwell-law}. Let $\mathbb P$ be the law of $u=(u_1,u_2,u_3)$ on $\mathcal U^3$, and let $\mathbb E$ denote the expectation with respect to $\mathbb P$. Let $m_B=m_{\mathbb P}$ be the law of Lemma~\ref{lem:mixtures-changes-variables} for this $\mathbb P$, that is, the law of Lemma~\ref{lem:gaussian-mixtures} for the probability space $(\mathcal U^3,\mathbb P)$ and the quadratic forms $q_u(f)=\left\|\Pi\Phi_{u\#}f\right\|^2$. For $f\in\mathcal S$ and $u\in\mathcal U^3$ let
\[
\Phi_f(u)=\exp\left(-\frac{1}{2}q_u(f)\right)\in(0,1].
\]
By Lemma~\ref{lem:gaussian-mixtures},
\[
m_B\left(\exp\left(iX(f)\right)\right)=\mathbb E\left[\Phi_f(u)\right]\qquad(f\in\mathcal S).
\]
Since $u_1,u_2,u_3$ are independent, $\mathbb P$ is the product of the laws of $u_1$, $u_2$ and $u_3$. For a bounded measurable function $\Psi$ on $\mathcal U^3$, $\mathbb E_{u_1}\Psi(u)$ denotes the integral of $\Psi(u)$ with respect to the law of $u_1$, the components $u_2$ and $u_3$ being fixed; by Fubini's theorem it is a measurable function of $(u_2,u_3)$. Test forms are identified with vector fields as in Section~\ref{sec:free-maxwell-law}, and $\|f\|_{L^1}$ is the $L^1$ norm of Lemma~\ref{lem:properties-free-maxwell}(ii). For $a\in\mathbb R^3$ let $|a|_\infty=\max_l|a_l|$, and let $u\mapsto u^a$ and $u\mapsto u^w$, $w\in O_{\mathrm{cub}}$, be the maps of $\mathcal U^3$ defined in Section~\ref{sec:free-maxwell-law}. The law $m_B$ has the following properties.
\begin{enumerate}
\item[(1)] The law $\mathbb P$ is invariant under the maps $u\mapsto u^a$, $a\in\mathbb R^3$, and $u\mapsto u^w$, $w\in O_{\mathrm{cub}}$. The law $m_B$ satisfies~\eqref{eq:setup-H}, and $m_B\left(\{\omega\in Q:d\omega=0\}\right)=1$.
\item[(2)] Let $r>0$, and let $f,g\in\mathcal D$ be supported in $[-r,r]^3$. For $a\in\mathbb R^3$ with $|a|_\infty>4r$ put
\[
\varepsilon(a)=\frac{\|f\|_{L^1}\cdot\|g\|_{L^1}}{8\pi}\cdot\left(\frac{|a|_\infty}{4}-r\right)^{-3}.
\]
Then every $a\in\mathbb R^3$ with $|a|_\infty>4r$ satisfies
\begin{equation}\label{eq:ergodic-components-9}
\left|m_B\left(\exp\left(iX(f)\right)\cdot\exp\left(iX(\tau_a\cdot g)\right)\right)-\mathbb E\left[\Phi_f(u)\cdot\Phi_g(u^a)\right]\right|\le\exp\left(\varepsilon(a)\right)-1.
\end{equation}
Moreover, $\Phi_f(u)$ is a measurable function of the restrictions $u_l|_{[-r,r]}$, $l=1,2,3$, and, for every $a\in\mathbb R^3$, $\Phi_g(u^a)$ is a measurable function of the restrictions $u_l|_{[a_l-r,a_l+r]}$, $l=1,2,3$.
\item[(3)] Let $r>0$, let $f,g\in\mathcal D$ be supported in $[-r,r]^3$, and let $a\in\mathbb R^3$ satisfy $\min_l|a_l|>2r$. Then
\[
\left|\mathbb E\left[\Phi_f(u)\cdot\Phi_g(u^a)\right]-\mathbb E\left[\Phi_f(u)\right]\cdot\mathbb E\left[\Phi_g(u)\right]\right|\le\sum_{l=1}^3\exp\left(-2\left(|a_l|-2r\right)\right).
\]
Moreover, $m_B$ is $\Gamma_{123}$-ergodic, that is, $m_B\in\mathcal M^{123}_{\mathrm{erg}}$.
\item[(4)] Let $r>0$, let $f\in\mathcal D$ be supported in $[-r,r]^3$, and let $\bar\Phi(u_2,u_3)=\mathbb E_{u_1}\Phi_f(u)$. Then
\[
\lim_{s\to\infty}m_B\left(\exp\left(iX(f)\right)\cdot\exp\left(iX(\tau_{s\cdot e_1}\cdot f)\right)\right)=\mathbb E\left[\bar\Phi^2\right],\qquad m_B\left(\exp\left(iX(f)\right)\right)=\mathbb E\left[\bar\Phi\right].
\]
If, moreover, $\ast f=(b_1,0,0)$, where $b_1\in C_c^\infty\left((-r,r)^3\right)$ is real and not identically zero, then
\[
\operatorname{Var}\bar\Phi=\mathbb E\left[\bar\Phi^2\right]-\left(\mathbb E\left[\bar\Phi\right]\right)^2>0,
\]
and consequently
\begin{equation}\label{eq:ergodic-components-10}
\lim_{s\to\infty}m_B\left(\exp\left(iX(f)\right)\cdot\exp\left(iX(\tau_{s\cdot e_1}\cdot f)\right)\right)>m_B\left(\exp\left(iX(f)\right)\right)^2.
\end{equation}
\end{enumerate}
\end{example}

\begin{proof}[Proof of the assertions of Example~\ref{ex:mixture-free-maxwell}]
\emph{Assertion~(1).} Let $a\in\mathbb R^3$ and $w\in O_{\mathrm{cub}}$, and let $\pi\in S_3$ be the permutation part of $w$, so that $w\,e_i=\pm e_{\pi(i)}$ for $i=1,2,3$ (Section~\ref{sec:euclidean-maps-actions}). By the definitions of Section~\ref{sec:free-maxwell-law}, the components of $u^a$ are the shifted processes $s\mapsto u_l(s+a_l)$, $l=1,2,3$, and the $i$-th component of $u^w$ is $u_{\pi(i)}$ if $w\,e_i=e_{\pi(i)}$, and the reversed process $s\mapsto u_{\pi(i)}(-s)$ if $w\,e_i=-e_{\pi(i)}$. By~\eqref{eq:ergodic-components-4}, which holds for the classes (Section~\ref{sec:free-maxwell-law}), each shifted and each reversed process has the law of the process from which it is formed. The processes $u_1,u_2,u_3$ are independent and have the same law, and the three components of $u^a$, as well as the three components of $u^w$, are functions of three distinct components of $u$. Hence $u^a$ and $u^w$ have the law $\mathbb P$.

The components of $u$ are independent, and each of them is of the form $\frac{1}{4}(1+\varsigma)$ with a telegraph process $\varsigma$. Therefore Lemma~\ref{lem:mixtures-changes-variables} applies to $\mathbb P$. By its part~(i), $m_B$ satisfies~\eqref{eq:setup-G} and $m_B\left(\{\omega\in Q:d\omega=0\}\right)=1$; by its part~(ii) and the invariance of $\mathbb P$ just proved, $m_B$ satisfies $(T_{123})$, and $w_*m_B=m_B$ for every $w\in O_{\mathrm{cub}}$; by its part~(iii), $m_B$ satisfies $(RP_j)$ for $j=1,2,3$. These are the conditions of~\eqref{eq:setup-H}.

\emph{Assertion~(2).} Let $a\in\mathbb R^3$ satisfy $|a|_\infty>4r$, and let $h=f+\tau_a\cdot g$. Since $X$ is linear, $\exp\left(iX(f)\right)\cdot\exp\left(iX(\tau_a\cdot g)\right)=\exp\left(iX(h)\right)$, and Lemma~\ref{lem:gaussian-mixtures} gives
\[
m_B\left(\exp\left(iX(f)\right)\cdot\exp\left(iX(\tau_a\cdot g)\right)\right)=\mathbb E\left[\exp\left(-\frac{1}{2}q_u(h)\right)\right].
\]
The map $\Phi_{u\#}$ is linear, and $\Phi_{u\#}(\tau_a\cdot g)=\tau_{\Phi_u(a)}\cdot\Phi_{u^a\#}g$ by Lemma~\ref{lem:coordinatewise-changes-variables}(ii). The orthogonal projection $\Pi$ satisfies $\Pi^2=\Pi=\Pi^*$ and commutes with the translations, which are isometries of $L^2(\mathbb R^3;\mathbb R^3)$ (Section~\ref{sec:free-maxwell-law}). Hence $\left\|\Pi\,\tau_{\Phi_u(a)}\cdot\Phi_{u^a\#}g\right\|^2=\left\|\Pi\Phi_{u^a\#}g\right\|^2=q_{u^a}(g)$, and the expansion of the square $\left\|\Pi\left(\Phi_{u\#}f+\tau_{\Phi_u(a)}\cdot\Phi_{u^a\#}g\right)\right\|^2$ gives
\[
q_u(h)=q_u(f)+q_{u^a}(g)+2c_a(u),\qquad c_a(u)=\left\langle\Phi_{u\#}f,\Pi\,\tau_{\Phi_u(a)}\cdot\Phi_{u^a\#}g\right\rangle.
\]
By Lemma~\ref{lem:coordinatewise-changes-variables}(v), the fields $\Phi_{u\#}f$ and $\Phi_{u^a\#}g$ are supported in $\left[-\frac{r}{2},\frac{r}{2}\right]^3$; so the two fields that enter $c_a(u)$ are supported in $\left[-\frac{r}{2},\frac{r}{2}\right]^3$ and in $\Phi_u(a)+\left[-\frac{r}{2},\frac{r}{2}\right]^3$, respectively. By Lemma~\ref{lem:coordinatewise-changes-variables}(i), $\Phi_u(0)=0$ and $|\varphi_l(a_l)|\ge\frac{1}{4}|a_l|$ for $l=1,2,3$, whence $|\Phi_u(a)|_\infty\ge\frac{1}{4}|a|_\infty$. For $y\in\left[-\frac{r}{2},\frac{r}{2}\right]^3$ and $y'\in\Phi_u(a)+\left[-\frac{r}{2},\frac{r}{2}\right]^3$ it follows that
\[
|y-y'|\ge|y-y'|_\infty\ge|\Phi_u(a)|_\infty-r\ge\frac{|a|_\infty}{4}-r>0.
\]
Thus the two fields have compact supports at distance at least $\frac{1}{4}|a|_\infty-r$. Lemma~\ref{lem:properties-free-maxwell}(ii), the bound $\|\Phi_{u\#}b\|_{L^1}\le\frac{1}{2}\|b\|_{L^1}$ of Lemma~\ref{lem:coordinatewise-changes-variables}(i), applied to $u$ and to $u^a$, and the invariance of the $L^1$ norm under translations give
\[
\left|c_a(u)\right|\le\frac{1}{2\pi}\cdot\left(\frac{|a|_\infty}{4}-r\right)^{-3}\cdot\left\|\Phi_{u\#}f\right\|_{L^1}\cdot\left\|\Phi_{u^a\#}g\right\|_{L^1}\le\frac{1}{2\pi}\cdot\left(\frac{|a|_\infty}{4}-r\right)^{-3}\cdot\frac{\|f\|_{L^1}}{2}\cdot\frac{\|g\|_{L^1}}{2}=\varepsilon(a).
\]
Since $q_u(f)$, $q_{u^a}(g)$ and $q_u(h)$ are measurable functions of $u$ (proof of Lemma~\ref{lem:mixtures-changes-variables}), so is $c_a(u)$, and
\[
m_B\left(\exp\left(iX(f)\right)\cdot\exp\left(iX(\tau_a\cdot g)\right)\right)-\mathbb E\left[\Phi_f(u)\cdot\Phi_g(u^a)\right]=\mathbb E\left[\Phi_f(u)\cdot\Phi_g(u^a)\cdot\left(\exp\left(-c_a(u)\right)-1\right)\right].
\]
As $0<\Phi_f(u)\cdot\Phi_g(u^a)\le1$, and as $\left|\exp(-c)-1\right|\le\exp(|c|)-1$ for every real number $c$, the bound~\eqref{eq:ergodic-components-9} follows.

For the last assertion of~(2), let $a\in\mathbb R^3$ be arbitrary. By Lemma~\ref{lem:coordinatewise-changes-variables}(v), the field $\Phi_{u\#}f\in L^2$ depends only on the restrictions $u_l|_{[-r,r]}$, $l=1,2,3$, and it depends continuously on them for the norm of $L^1([-r,r])$. The Borel sets of $L^1([-r,r])$ are generated by the functionals $\eta\mapsto\int\mathrm ds\,\phi(s)\cdot\eta(s)$, $\phi\in C_c\left((-r,r)\right)$, as in the proof of Lemma~\ref{lem:mixtures-changes-variables}. Hence $q_u(f)$ and $\Phi_f(u)$ are measurable functions of the restrictions $u_l|_{[-r,r]}$, that is, measurable with respect to the product of the $\sigma$-algebras $\sigma(u_l|_{[-r,r]})$ of Section~\ref{sec:free-maxwell-law}. The same argument, applied to $g$ and to $u^a$ in place of $f$ and $u$, shows that $\Phi_g(u^a)$ is a measurable function of the restrictions of the components $u_l(\cdot+a_l)$ of $u^a$ to $[-r,r]$, that is, of the restrictions $u_l|_{[a_l-r,a_l+r]}$, $l=1,2,3$.

\emph{Assertion~(3).} Let $(u,u')$ be the pair of coordinates on $\mathcal U^3\times\mathcal U^3$, endowed with the probability measure $\mathbb P\otimes\mathbb P$; thus $u'=(u'_1,u'_2,u'_3)$ has the law $\mathbb P$ and is independent of $u$, and the six components of $u$ and $u'$ are independent. In this part, expectations refer to $\mathbb P\otimes\mathbb P$. For $k=0,1,2,3$ let $u^{(k)}$ be the triple obtained from $u$ by replacing its components $u_1,\dots,u_k$ by $u'_1,\dots,u'_k$, and let
\[
Z_k=\mathbb E\left[\Phi_f(u)\cdot\Phi_g\left(\left(u^{(k)}\right)^a\right)\right].
\]
Then $u^{(0)}=u$ and $u^{(3)}=u'$. As $u'$ is independent of $u$, and as $(u')^a$ has the law $\mathbb P$ by~(1),
\[
Z_0=\mathbb E\left[\Phi_f(u)\cdot\Phi_g(u^a)\right],\qquad Z_3=\mathbb E\left[\Phi_f(u)\right]\cdot\mathbb E\left[\Phi_g(u)\right].
\]
Let $l\in\{1,2,3\}$, and fix all components of $u$ and of $u'$ other than $u_l$ and $u'_l$. By the last assertion of~(2), applied to the triples $u$, $u^{(l-1)}$ and $u^{(l)}$, there are functions $\psi_1$ and $\psi_2$ on $\mathcal U$ with values in $[0,1]$ such that
\[
\Phi_f(u)=\psi_1(u_l),\qquad\Phi_g\left(\left(u^{(l-1)}\right)^a\right)=\psi_2(u_l),\qquad\Phi_g\left(\left(u^{(l)}\right)^a\right)=\psi_2(u'_l),
\]
where $\psi_1(u_l)$ is measurable with respect to $\sigma(u_l|_{[-r,r]})$ and $\psi_2(u_l)$ with respect to $\sigma(u_l|_{[a_l-r,a_l+r]})$ (Section~\ref{sec:free-maxwell-law}); the last identity holds because $u^{(l)}$ differs from $u^{(l-1)}$ only in the $l$-th component, which is $u'_l$ instead of $u_l$. By Section~\ref{sec:free-maxwell-law}, $\psi_1(u_l)$ and $\psi_2(u_l)$ are bounded functions of the telegraph process $\varsigma_l$, measurable with respect to $\mathcal F_{[-r,r]}$ and to $\mathcal F_{[a_l-r,a_l+r]}$, respectively. If $a_l>2r$, then $\mathcal F_{[-r,r]}\subset\mathcal F_{(-\infty,r]}$ and $\mathcal F_{[a_l-r,a_l+r]}\subset\mathcal F_{[r+t,\infty)}$ with $t=a_l-2r$. If $a_l<-2r$, then $\mathcal F_{[a_l-r,a_l+r]}\subset\mathcal F_{(-\infty,a_l+r]}$ and $\mathcal F_{[-r,r]}\subset\mathcal F_{[(a_l+r)+t,\infty)}$ with $t=|a_l|-2r$. In both cases~\eqref{eq:ergodic-components-5} gives
\[
\left|\mathbb E\left[\psi_1(u_l)\cdot\psi_2(u_l)\right]-\mathbb E\left[\psi_1(u_l)\right]\cdot\mathbb E\left[\psi_2(u_l)\right]\right|\le\exp\left(-2\left(|a_l|-2r\right)\right),
\]
while $\mathbb E\left[\psi_1(u_l)\cdot\psi_2(u'_l)\right]=\mathbb E\left[\psi_1(u_l)\right]\cdot\mathbb E\left[\psi_2(u_l)\right]$, because $u'_l$ is independent of $u_l$ and has the same law. Integrating over the fixed components (Fubini's theorem), one obtains
\[
\left|Z_{l-1}-Z_l\right|\le\exp\left(-2\left(|a_l|-2r\right)\right)\qquad(l=1,2,3),
\]
and summation over $l$ gives the inequality of~(3).

Let now $r>0$, let $f,g\in\mathcal D$ be supported in $[-r,r]^3$, and let $a\in\mathbb R^3$ satisfy $\min_l|a_l|>4r$. Then $|a|_\infty>4r$ and $\min_l|a_l|>2r$. By~\eqref{eq:ergodic-components-9}, by the inequality just proved, and by $m_B\left(\exp\left(iX(f)\right)\right)=\mathbb E\left[\Phi_f(u)\right]$, $m_B\left(\exp\left(iX(g)\right)\right)=\mathbb E\left[\Phi_g(u)\right]$,
\begin{multline*}
\left|m_B\left(\exp\left(iX(f)\right)\cdot\exp\left(iX(\tau_a\cdot g)\right)\right)-m_B\left(\exp\left(iX(f)\right)\right)\cdot m_B\left(\exp\left(iX(g)\right)\right)\right|\\
\le\exp\left(\varepsilon(a)\right)-1+\sum_{l=1}^3\exp\left(-2\left(|a_l|-2r\right)\right).
\end{multline*}
For $k\ge1$ let $a^{(k)}=k\cdot(e_1+e_2+e_3)$. As $k\to\infty$, $\varepsilon(a^{(k)})\to0$ and the sum on the right tends to $0$. Since every element of $\mathcal D$ is supported in some cube $[-r,r]^3$, it follows that
\[
\lim_{k\to\infty}m_B\left(\exp\left(iX(f)\right)\cdot\exp\left(iX(\tau_{a^{(k)}}\cdot g)\right)\right)=m_B\left(\exp\left(iX(f)\right)\right)\cdot m_B\left(\exp\left(iX(g)\right)\right)\qquad(f,g\in\mathcal D).
\]
Let $F=\sum_pz_p\cdot\exp\left(iX(f_p)\right)$ and $G=\sum_qz'_q\cdot\exp\left(iX(g_q)\right)$ be elements of the space $\mathcal A$ of Lemma~\ref{lem:density-exponentials}, with finitely many coefficients $z_p,z'_q\in\mathbb C$ and test forms $f_p,g_q\in\mathcal D_0\subset\mathcal D$. Then $\overline{F}=\sum_p\overline{z_p}\cdot\exp\left(iX(-f_p)\right)$ and $U(a)G=\sum_qz'_q\cdot\exp\left(iX(\tau_a\cdot g_q)\right)$ (Section~\ref{sec:euclidean-maps-actions}), and $m_B\left(\exp\left(iX(-f_p)\right)\right)=\overline{m_B\left(\exp\left(iX(f_p)\right)\right)}$. The last display therefore gives
\[
\lim_{k\to\infty}m_B\left(\overline{F}\cdot U\left(a^{(k)}\right)G\right)=\sum_{p,q}\overline{z_p}\cdot z'_q\cdot\overline{m_B\left(\exp\left(iX(f_p)\right)\right)}\cdot m_B\left(\exp\left(iX(g_q)\right)\right)=\overline{m_B(F)}\cdot m_B(G).
\]
By~(1), $m_B$ satisfies $(T_{123})$, and $a^{(k)}\in\Gamma_{123}=\mathbb R^3$. Lemma~\ref{lem:clustering-criterion-separation}(a), applied with $J=\{1,2,3\}$, shows that $m_B\in\mathcal M^{123}_{\mathrm{erg}}$.

\emph{Assertion~(4).} Let $s>4r$. Then $|s\cdot e_1|_\infty=s>4r$, and~\eqref{eq:ergodic-components-9}, applied with $g=f$ and $a=s\cdot e_1$, gives
\[
\left|m_B\left(\exp\left(iX(f)\right)\cdot\exp\left(iX(\tau_{s\cdot e_1}\cdot f)\right)\right)-\mathbb E\left[\Phi_f(u)\cdot\Phi_f\left(u^{s\cdot e_1}\right)\right]\right|\le\exp\left(\varepsilon(s\cdot e_1)\right)-1,
\]
where $u^{s\cdot e_1}=\left(u_1(\cdot+s),u_2,u_3\right)$, and where
\[
\varepsilon(s\cdot e_1)=\frac{\|f\|_{L^1}^2}{8\pi}\cdot\left(\frac{s}{4}-r\right)^{-3}
\]
tends to $0$ as $s\to\infty$. Fix $u_2$ and $u_3$. By the last assertion of~(2), $\Phi_f(u)$ and $\Phi_f\left(u^{s\cdot e_1}\right)$ are functions of $u_1$ with values in $[0,1]$, measurable with respect to $\sigma(u_1|_{[-r,r]})$ and to $\sigma(u_1|_{[s-r,s+r]})$, respectively; by Section~\ref{sec:free-maxwell-law}, they are bounded functions of $\varsigma_1$, measurable with respect to $\mathcal F_{(-\infty,r]}$ and to $\mathcal F_{[r+t,\infty)}$, where $t=s-2r$. By~\eqref{eq:ergodic-components-4}, the process $u_1(\cdot+s)$ has the law of $u_1$, so that $\mathbb E_{u_1}\Phi_f\left(u^{s\cdot e_1}\right)=\mathbb E_{u_1}\Phi_f(u)=\bar\Phi(u_2,u_3)$. Now~\eqref{eq:ergodic-components-5}, applied to $\varsigma_1$, gives
\[
\left|\mathbb E_{u_1}\left[\Phi_f(u)\cdot\Phi_f\left(u^{s\cdot e_1}\right)\right]-\bar\Phi(u_2,u_3)^2\right|\le\exp\left(-2(s-2r)\right),
\]
and integration over $(u_2,u_3)$ yields
\[
\left|\mathbb E\left[\Phi_f(u)\cdot\Phi_f\left(u^{s\cdot e_1}\right)\right]-\mathbb E\left[\bar\Phi^2\right]\right|\le\exp\left(-2(s-2r)\right).
\]
Together with the preceding bound this proves the first identity of~(4). The second identity is $m_B\left(\exp\left(iX(f)\right)\right)=\mathbb E\left[\Phi_f(u)\right]=\mathbb E\left[\bar\Phi\right]$, by Fubini's theorem.

Let now $\ast f=(b_1,0,0)$ as in~(4), and write $b=\ast f$. For $c\in\left\{\frac{1}{4},\frac{1}{2}\right\}$ let
\[
B_c=\left\{u\in\mathcal U^3:\int_{-r}^r\mathrm ds\,\left|u_2(s)-c\right|=\int_{-r}^r\mathrm ds\,\left|u_3(s)-c\right|=0\right\}
\]
be the event that $u_2=u_3=c$ almost everywhere on $[-r,r]$. It depends on $(u_2,u_3)$ only, and it has positive probability, by~\eqref{eq:ergodic-components-6} (Section~\ref{sec:free-maxwell-law}) and the independence of $u_2$ and $u_3$.

Let $u\in B_c$. With the functions $\varphi_l(x_l)=\int_0^{x_l}\mathrm ds\,u_l(s)$ of Section~\ref{sec:free-maxwell-law} one has $\varphi_2(x_2)=c\cdot x_2$ and $\varphi_3(x_3)=c\cdot x_3$ for $x_2,x_3\in[-r,r]$. Since $b_2=b_3=0$, the definition~\eqref{eq:ergodic-components-8} shows that the second and the third component of $\Phi_{u\#}b$ vanish and that
\[
\left(\Phi_{u\#}b\right)_1(y)=\frac{b_1(x)}{u_1(x_1)},\qquad x=\Phi_u^{-1}(y)\qquad(y\in\mathbb R^3).
\]
If $|y_2|\le c\cdot r$ and $|y_3|\le c\cdot r$, then $x_2=c^{-1}\cdot y_2$ and $x_3=c^{-1}\cdot y_3$. If $|y_2|>c\cdot r$, then $|x_2|>r$, because $\varphi_2$ is increasing with $\varphi_2(\pm r)=\pm c\cdot r$; then $b_1(x)=0$, and also $|c^{-1}\cdot y_2|>r$. The same holds with $y_3$ in place of $y_2$. As $b_1$ vanishes outside $(-r,r)^3$, it follows that
\[
\left(\Phi_{u\#}b\right)_1(y)=\beta\left(y_1,c^{-1}\cdot y_2,c^{-1}\cdot y_3\right)\quad(y\in\mathbb R^3),\qquad\beta(y_1,x_2,x_3)=\frac{b_1\left(\varphi_1^{-1}(y_1),x_2,x_3\right)}{u_1\left(\varphi_1^{-1}(y_1)\right)}.
\]
The function $\beta$ depends on $u_1$ alone, and not on $u_2$, $u_3$ or $c$. It is bounded, it has compact support, and it does not vanish on the image of the non-empty open set $\{b_1\ne0\}$ under the homeomorphism $(x_1,x_2,x_3)\mapsto\left(\varphi_1(x_1),x_2,x_3\right)$ of $\mathbb R^3$ (Lemma~\ref{lem:coordinatewise-changes-variables}(i)), which is a non-empty open set. Hence $\beta\in L^1(\mathbb R^3)\cap L^2(\mathbb R^3)$, and its Fourier transform
\[
\hat\beta(p)=\int\mathrm dy\,\exp(-ip\cdot y)\cdot\beta(y)\qquad(p\in\mathbb R^3)
\]
is continuous and not identically zero.

With the Fourier transform normalised in this way, Plancherel's theorem reads $\|\eta\|^2=(2\pi)^{-3}\int\mathrm dp\,\left|\hat\eta(p)\right|^2$ for $\eta\in L^2(\mathbb R^3)$, and the projection $\Pi$ acts by the multiplier $1-p\,p^T/|p|^2$ (Section~\ref{sec:free-maxwell-law}). For a field $b'=(b'_1,0,0)\in L^2(\mathbb R^3;\mathbb R^3)$ it follows that
\[
\left\|\Pi b'\right\|^2=\frac{1}{(2\pi)^3}\int\mathrm dp\,\left|\widehat{b'_1}(p)\right|^2\cdot\left(1-\frac{p_1^2}{|p|^2}\right)=\frac{1}{(2\pi)^3}\int\mathrm dp\,\left|\widehat{b'_1}(p)\right|^2\cdot\frac{p_2^2+p_3^2}{|p|^2}.
\]
The first component $y\mapsto\beta\left(y_1,c^{-1}\cdot y_2,c^{-1}\cdot y_3\right)$ of $\Phi_{u\#}b$ has the Fourier transform $p\mapsto c^2\cdot\hat\beta(p_1,c\cdot p_2,c\cdot p_3)$, by the substitution $y_2=c\cdot x_2$, $y_3=c\cdot x_3$. Hence, for $u\in B_c$,
\begin{align*}
q_u(f)&=\frac{1}{(2\pi)^3}\int\mathrm dp\,c^4\cdot\left|\hat\beta(p_1,c\cdot p_2,c\cdot p_3)\right|^2\cdot\frac{p_2^2+p_3^2}{|p|^2}\\
&=\frac{1}{(2\pi)^3}\int\mathrm dp_1\int\mathrm d^2k\,\left|\hat\beta(p_1,k)\right|^2\cdot\frac{|k|^2}{p_1^2+c^{-2}\cdot|k|^2},
\end{align*}
the second expression by the substitution $k=(c\cdot p_2,c\cdot p_3)\in\mathbb R^2$. For $p_1\in\mathbb R$ and $k\in\mathbb R^2\setminus\{0\}$, the factor $|k|^2/\left(p_1^2+c^{-2}\cdot|k|^2\right)$ is strictly larger for $c=\frac{1}{2}$ than for $c=\frac{1}{4}$. The set of the points $(p_1,k)$ with $\hat\beta(p_1,k)\ne0$ is open and non-empty, hence of positive measure, and the points with $k=0$ form a null set. Consequently, for every $u_1$, the value of $q_u(f)$ for $u\in B_{1/2}$, which depends on $u_1$ alone and is denoted $Q_{1/2}(u_1)$, is strictly larger than the value $Q_{1/4}(u_1)$ of $q_u(f)$ for $u\in B_{1/4}$.

It follows that $\bar\Phi$ is constant on each of the events $B_{1/2}$ and $B_{1/4}$, and that its value on $B_{1/2}$ is strictly smaller than its value on $B_{1/4}$:
\[
\mathbb E_{u_1}\exp\left(-\frac{1}{2}Q_{1/2}(u_1)\right)<\mathbb E_{u_1}\exp\left(-\frac{1}{2}Q_{1/4}(u_1)\right),
\]
the inequality being strict because the difference of the two integrands is strictly positive for every $u_1$. As $B_{1/2}$ and $B_{1/4}$ have positive probability, $\bar\Phi$ is not almost surely constant, and $\operatorname{Var}\bar\Phi>0$. With the two identities of~(4) this gives
\[
\lim_{s\to\infty}m_B\left(\exp\left(iX(f)\right)\cdot\exp\left(iX(\tau_{s\cdot e_1}\cdot f)\right)\right)=\mathbb E\left[\bar\Phi^2\right]>\left(\mathbb E\left[\bar\Phi\right]\right)^2=m_B\left(\exp\left(iX(f)\right)\right)^2,
\]
which is~\eqref{eq:ergodic-components-10}.
\end{proof}

\subsubsection{Asymmetric components and degenerate vacua}\label{sec:asymmetric-components-degenerate}

\begin{proof}[Proof of Theorem~\ref{thm:asymmetric-components-non}]
Every law with~\eqref{eq:setup-H} belongs to $\mathcal C$ (Section~\ref{sec:hypothesis-law}); for such a law $\rho$ and for $j\in\{1,2,3\}$, the operator $H_j(\rho)$ and the vector $\Omega_j$ are therefore defined.

\emph{The law $m_A$ and assertion~(A).} Let $\gamma^{(\pi)}$, $\pi\in S_3$, be the six laws of Section~\ref{sec:permuted-diagonal-rescalings}, and for $w\in O_{\mathrm{cub}}$ let $\pi_w\in S_3$ be the permutation part of $w$, so that $w\,e_i=\pm e_{\pi_w(i)}$ for $i=1,2,3$. By Section~\ref{sec:permuted-diagonal-rescalings}(2), $m_A$ satisfies~\eqref{eq:setup-H} and $m_A\left(\{\omega\in Q:d\omega=0\}\right)=1$. By Section~\ref{sec:permuted-diagonal-rescalings}(1), (3) and~(4), the laws $\gamma^{(\pi)}$, $\pi\in S_3$, are pairwise distinct; each of them is a centred Gaussian law, under which $X(f)$ has the centred normal law with variance $q_{u^{(\pi)}}(f)$ for every $f\in\mathcal S$ (proof of Lemma~\ref{lem:gaussian-mixtures}), and each of them is mixing; and
\[
\nu_A=\frac{1}{6}\sum_{\pi\in S_3}\delta_{\gamma^{(\pi)}}.
\]
By Section~\ref{sec:permuted-diagonal-rescalings}(2), $(w^{-1})_*\gamma^{(\pi)}=\gamma^{(\pi\pi_w)}$ for $w\in O_{\mathrm{cub}}$ and $\pi\in S_3$. The map $w\mapsto\pi_w$ is a homomorphism of $O_{\mathrm{cub}}$ onto $S_3$ with kernel $N_3$ (Section~\ref{sec:euclidean-maps-actions}); applied to $w^{-1}$ in place of $w$, the last identity therefore reads
\[
w_*\gamma^{(\pi)}=\gamma^{(\pi\pi_w^{-1})}\qquad\left(w\in O_{\mathrm{cub}},\ \pi\in S_3\right).
\]
Since $\pi\mapsto\pi\pi_w^{-1}$ is a bijection of $S_3$, the map $\rho\mapsto w_*\rho$ permutes the set $\left\{\gamma^{(\pi)}:\pi\in S_3\right\}$; it fixes each element of this set if $w\in N_3$, because $\pi_w$ is then the identity. Let $\pi,\sigma\in S_3$. As the six laws are pairwise distinct, $w_*\gamma^{(\pi)}=\gamma^{(\sigma)}$ holds if and only if $\pi\pi_w^{-1}=\sigma$, that is, if and only if $\pi_w=\sigma^{-1}\pi$; and the set of the $w\in O_{\mathrm{cub}}$ with $\pi_w=\sigma^{-1}\pi$ is a coset of $N_3$. An enumeration $\rho_1,\dots,\rho_6$ of the laws $\gamma^{(\pi)}$ therefore has all the properties required in~(A).

\emph{Assertion~(B).} It is derived from the following claim.

\emph{Claim.} Let $m\in\mathcal M$ satisfy~\eqref{eq:setup-H}, let $m\in\mathcal M^{123}_{\mathrm{erg}}$, and let there be $f\in\mathcal D$ such that the limit
\[
L=\lim_{s\to\infty}m\left(\exp\left(iX(f)\right)\cdot\exp\left(iX(\tau_{s\cdot e_1}\cdot f)\right)\right)
\]
exists and $L>m\left(\exp\left(iX(f)\right)\right)^2$. Then assertion~(B) holds with $m$ in place of $m_B$.

By Section~\ref{sec:free-maxwell-law-2}(1), the law $m_B$ satisfies~\eqref{eq:setup-H} and $m_B\left(\{\omega\in Q:d\omega=0\}\right)=1$; by Section~\ref{sec:free-maxwell-law-2}(3), $m_B\in\mathcal M^{123}_{\mathrm{erg}}$. Let $r>0$, let $b_1\in C_c^\infty\left((-r,r)^3\right)$ be real and not identically zero, and let $f\in\mathcal D$ be given by $\ast f=(b_1,0,0)$; by Section~\ref{sec:free-maxwell-law-2}(4), the limit in~\eqref{eq:ergodic-components-10} exists, and~\eqref{eq:ergodic-components-10} holds. Thus $m_B$ satisfies the hypotheses of the claim, and~(B) follows from the claim.

\emph{Proof of the claim, step~1: the decomposition of $m$ is trivial.} The law $m$ satisfies $(T_{123})$, and $\mathcal M^{123}_{\mathrm{erg}}\in\mathcal B_{\mathcal M}$ by Theorem~\ref{thm:ergodic-decomposition}(a). The probability measure $\delta_m$ on $(\mathcal M,\mathcal B_{\mathcal M})$ satisfies $\delta_m\left(\mathcal M^{123}_{\mathrm{erg}}\right)=1$, because $m\in\mathcal M^{123}_{\mathrm{erg}}$, and $m(A)=\int_{\mathcal M}\delta_m(\mathrm d\rho)\,\rho(A)$ for every $A\in\Sigma$. By the uniqueness assertion of Theorem~\ref{thm:ergodic-decomposition}(a), $\delta_m$ is the probability measure which Theorem~\ref{thm:ergodic-decomposition}(a) assigns to $m$ and $J=\{1,2,3\}$.

\emph{Step~2: a test form for each axis.} For $j\in\{1,2,3\}$ there is $f_j\in\mathcal D$ such that the limit
\[
L_j=\lim_{s\to\infty}m\left(\exp\left(iX(f_j)\right)\cdot\exp\left(iX(\tau_{s\cdot e_j}\cdot f_j)\right)\right)
\]
exists and $L_j>m\left(\exp\left(iX(f_j)\right)\right)^2$. For $j=1$ let $f_1=f$. For $j\in\{2,3\}$ let $w\in O_{\mathrm{cub}}$ be the matrix that exchanges $e_1$ and $e_j$ and fixes the third vector of the standard basis, and let $f_j=w\cdot f\in\mathcal D$. Since $w(x+s\cdot e_1)=wx+s\cdot e_j$ for $x\in\mathbb R^3$, the Euclidean maps $w\circ\tau_{s\cdot e_1}$ and $\tau_{s\cdot e_j}\circ w$ coincide, and Section~\ref{sec:euclidean-maps-actions} gives $\tau_{s\cdot e_j}\cdot f_j=w\cdot\left(\tau_{s\cdot e_1}\cdot f\right)$. The rules of Section~\ref{sec:euclidean-maps-actions} for $\hat w$ then give
\[
\exp\left(iX(f_j)\right)=\hat w\exp\left(iX(f)\right),\qquad\exp\left(iX(f_j)\right)\cdot\exp\left(iX(\tau_{s\cdot e_j}\cdot f_j)\right)=\hat w\left(\exp\left(iX(f)\right)\cdot\exp\left(iX(\tau_{s\cdot e_1}\cdot f)\right)\right),
\]
and $m(\hat wF)=m(F)$ for every bounded measurable function $F$ on $Q$, because $w_*m=m$. Hence $m\left(\exp\left(iX(f_j)\right)\right)=m\left(\exp\left(iX(f)\right)\right)$, and the limit $L_j$ exists and equals $L$; so $L_j>m\left(\exp\left(iX(f_j)\right)\right)^2$.

\emph{Step~3: the kernel of $H_j(m)$.} Let $j\in\{1,2,3\}$. The law $m$ satisfies $(T_j)$, $(R_j)$ and $(RP_j)$, so Theorem~\ref{thm:invariant-vectors-clustering}(ii) applies to it. Let $F=\exp\left(-iX(f_j)\right)$ and $G=\exp\left(iX(f_j)\right)$. Then $\overline{F}=\exp\left(iX(f_j)\right)$, $U(s\cdot e_j)G=\exp\left(iX(\tau_{s\cdot e_j}\cdot f_j)\right)$ and $\overline{m(F)}\cdot m(G)=m\left(\exp\left(iX(f_j)\right)\right)^2$. By step~2, $m\left(\overline{F}\cdot U(s\cdot e_j)G\right)$ tends to $L_j$ as $s\to\infty$, and $L_j\ne\overline{m(F)}\cdot m(G)$. Thus property~(3) of Theorem~\ref{thm:invariant-vectors-clustering}(ii) fails, and, by the equivalence of the properties~(1) and~(3) there, $\ker H_j(m)\ne\mathbb C\cdot\Omega_j$. By Theorem~\ref{thm:invariant-vectors-clustering}(ii), the map $F'\mapsto[F']_j$ maps $\mathcal E^j(m)$ isometrically onto $\ker H_j(m)$. The constant function $1$ belongs to $\mathcal E^j(m)$; hence $\Omega_j=[1]_j$ lies in $\ker H_j(m)$ and has norm $1$. If $\ker H_j(m)$ had dimension $1$, it would be equal to $\mathbb C\cdot\Omega_j$; therefore $\dim\ker H_j(m)\ge2$. Since $m\in\mathcal C$ is $\Gamma_{123}$-ergodic, Theorem~\ref{thm:invariant-vectors-clustering}(iii)5 gives $\dim\ker H_j(m)\in\{1,\infty\}$, and so
\[
\dim\ker H_j(m)=\infty.
\]

\emph{Step~4: no decomposition into $\Gamma_j$-ergodic laws.} Let $j\in\{1,2,3\}$, and apply Theorem~\ref{thm:vacuum-symmetries-components}(i) to $m\in\mathcal C$ and $j$. By step~1, the probability measure $\nu$ of that theorem is $\delta_m$. Property~(1) of Theorem~\ref{thm:vacuum-symmetries-components}(i) fails: it would give $\ker H_j(\rho)=\mathbb C\cdot\Omega_j$ for all $\rho$ in a set of $\delta_m$-measure one, and every such set contains $m$, whereas $\ker H_j(m)\ne\mathbb C\cdot\Omega_j$ by step~3. By the equivalence of the properties~(1) and~(4) in Theorem~\ref{thm:vacuum-symmetries-components}(i), property~(4) fails as well: there is no probability measure $\nu'$ on $(\mathcal M,\mathcal B_{\mathcal M})$ with $m(A)=\int_{\mathcal M}\nu'(\mathrm d\rho)\,\rho(A)$ for every $A\in\Sigma$ such that $\nu'$-almost every $\rho$ satisfies $(T_{123})$ and is $\Gamma_j$-ergodic.

Steps~1 to~4, together with the hypothesis $m\in\mathcal M^{123}_{\mathrm{erg}}$, are the assertions of~(B) for $m$ in place of $m_B$. This proves the claim, and hence~(B).

\emph{The laws $m_1$, $m_2$ and assertion~(C).} Let $D=(d_{12},d_{13},d_{23})$, the laws $\gamma_{D'}$ and the law $m_1=\frac{1}{6}\sum_{\pi\in S_3}\gamma_{\pi D}$ be as in Section~\ref{sec:average-six-anisotropic}, and let $\nu_1$ be the probability measure which Theorem~\ref{thm:ergodic-decomposition}(a) assigns to $m=m_1$ and $J=\{1,2,3\}$. By Section~\ref{sec:average-six-anisotropic}(3), $m_1$ satisfies~\eqref{eq:setup-H}, and the six laws $\gamma_{\pi D}$, $\pi\in S_3$, are pairwise distinct; by Section~\ref{sec:average-six-anisotropic}(4),
\[
\nu_1=\frac{1}{6}\sum_{\pi\in S_3}\delta_{\gamma_{\pi D}}.
\]
Each $\gamma_{\pi D}$ is the law of Lemma~\ref{lem:gaussian-mixtures} for a one-point probability space and the form $q_{\pi D}$; by the proof of that lemma, $X(f)$ has under $\gamma_{\pi D}$ the centred normal law with variance $q_{\pi D}(f)$, for every $f\in\mathcal S$. By Section~\ref{sec:average-six-anisotropic}(2), applied to the triple $\pi D$ in place of $D$, each $\gamma_{\pi D}$ is mixing. By Section~\ref{sec:average-six-anisotropic}, $w_*\gamma_{D'}=\gamma_{\pi_wD'}$ for $w\in O_{\mathrm{cub}}$ and for every triple $D'=(d'_{12},d'_{13},d'_{23})$ of distinct numbers in $(0,1]$. For $\pi,\sigma\in S_3$ and $i<l$, the definition of the triples $\pi D'$ gives
\[
\left(\pi(\sigma D)\right)_{\pi\sigma(i)\,\pi\sigma(l)}=(\sigma D)_{\sigma(i)\,\sigma(l)}=d_{il},
\]
that is, $\pi(\sigma D)=(\pi\sigma)D$. Hence
\[
w_*\gamma_{\sigma D}=\gamma_{(\pi_w\sigma)D}\qquad\left(w\in O_{\mathrm{cub}},\ \sigma\in S_3\right).
\]
As $\sigma\mapsto\pi_w\sigma$ is a bijection of $S_3$, the map $\rho\mapsto w_*\rho$ permutes the set $\left\{\gamma_{\sigma D}:\sigma\in S_3\right\}$, and it fixes each of its elements if $w\in N_3$. For $\pi,\sigma\in S_3$, the six laws being pairwise distinct, $w_*\gamma_{\sigma D}=\gamma_{\pi D}$ holds if and only if $\pi_w\sigma=\pi$, that is, if and only if $\pi_w=\pi\sigma^{-1}$; the set of these $w$ is a coset of $N_3$. Thus~(A) holds with $m_1$ in place of $m_A$ and $\nu_1$ in place of $\nu_A$.

By Section~\ref{sec:white-noise-modulated}(1) and~(2), the law $m_2$ satisfies~\eqref{eq:setup-H}; by Section~\ref{sec:white-noise-modulated}(3), $m_2\in\mathcal M^{123}_{\mathrm{erg}}$; and by Section~\ref{sec:white-noise-modulated}(4) there is $f\in\mathcal D$ such that the limit $\lim_{s\to\infty}m_2\left(\exp\left(iX(f)\right)\cdot\exp\left(iX(\tau_{s\cdot e_1}\cdot f)\right)\right)$ exists and is larger than $m_2\left(\exp\left(iX(f)\right)\right)^2$. By the claim, (B) holds with $m_2$ in place of $m_B$.

Finally, let $\mathsf g\in\mathcal S(\mathbb R^3)$, $\mathsf g\ne0$, and let $f\in\mathcal S$ be given by $\ast f=\nabla\mathsf g$. A Schwartz function with vanishing gradient is constant, hence zero; so $\nabla\mathsf g\ne0$, and $\|f\|=\|\nabla\mathsf g\|>0$, since $\ast$ is an isometry. Every closed $\omega\in Q$ satisfies $\omega(f)=0$ (Section~\ref{sec:closed-distributions}); hence every law $\rho$ with $\rho\left(\{\omega\in Q:d\omega=0\}\right)=1$ satisfies $\rho\left(X(f)^2\right)=0$. By the proof of Lemma~\ref{lem:gaussian-mixtures}, under the law of that lemma the random variable $X(f)$ has the law of the mixture of the centred normal laws with the variances $q_u(f)$, and therefore its second moment is the mean of $q_u(f)$. The law $m_1$ is the average of the six laws $\gamma_{\pi D}$, and $m_2$ is the law of Lemma~\ref{lem:gaussian-mixtures} for the forms $q_v(f)=\int\mathrm dx\,\mathsf a_v(x)\cdot|f(x)|^2$ of Section~\ref{sec:white-noise-modulated}. Therefore
\[
m_1\left(X(f)^2\right)=\frac{1}{6}\sum_{\pi\in S_3}q_{\pi D}(f)>0,\qquad m_2\left(X(f)^2\right)=\mathbb E\,q_v(f)\ge\frac{1}{8}\|f\|^2>0,
\]
where $\mathbb E$ is the expectation with respect to the law of $v$; indeed $q_{\pi D}(f)\ge\min_{i<l}d_{il}\cdot\|f\|^2$ with $\min_{i<l}d_{il}>0$, and $\mathsf a_v\ge\frac{1}{8}$. Consequently
\[
m_1\left(\{\omega\in Q:d\omega=0\}\right)<1,\qquad m_2\left(\{\omega\in Q:d\omega=0\}\right)<1.
\]
This completes the proof of~(C) and of the theorem.
\end{proof}

\begin{remark}[Periodic and mixing modulations]\label{rem:periodic-mixing-modulations}
Let $r>0$, and let $(m,f)$ be one of the following two pairs: $m=m_2$ and $f\in\mathcal D$, $f\ne0$, supported in $[-r,r]^3$, as in Example~\ref{ex:white-noise-telegraph}(4); or $m=m_B$ and $f\in\mathcal D$ given by $\ast f=(b_1,0,0)$, where $b_1\in C_c^\infty\left((-r,r)^3\right)$ is real and not identically zero, as in the second part of Example~\ref{ex:mixture-free-maxwell}(4). The limit of $m\left(\exp\left(iX(f)\right)\cdot\exp\left(iX(\tau_{s\cdot e_1}\cdot f)\right)\right)$ as $s\to\infty$ exceeds $m\left(\exp\left(iX(f)\right)\right)^2$ because of the media $(v_2,v_3)$ of Example~\ref{ex:white-noise-telegraph}, resp.\ $(u_2,u_3)$ of Example~\ref{ex:mixture-free-maxwell}, which do not depend on $x_1$. Indeed, by Example~\ref{ex:white-noise-telegraph}(4), resp.\ Example~\ref{ex:mixture-free-maxwell}(4), this limit is $\mathbb E\left[\bar\Phi^2\right]$, and $m\left(\exp\left(iX(f)\right)\right)^2=\left(\mathbb E\left[\bar\Phi\right]\right)^2$, where $\bar\Phi$, the mean of $\Phi_f$ over $v_1$, resp.\ over $u_1$, is a function of these media alone; the excess is the variance of $\bar\Phi$, which is strictly positive by Example~\ref{ex:white-noise-telegraph}(4), resp.\ by the second part of Example~\ref{ex:mixture-free-maxwell}(4).

Proposition~\ref{prop:uniform-mixing-along} places a restriction on structures of this kind, in the form of the following statement on eigenvectors. Let $j\in\{1,2,3\}$, let $\rho\in\mathcal M$ satisfy $(T_j)$, $(R_j)$ and $(RP_j)$, and let $F\in L^2(\rho)$, $F\ne0$, and $\lambda\in\mathbb R$ satisfy
\[
U(s\cdot e_j)F=\exp(i\lambda\cdot s)\cdot F\qquad\text{for every }s\in\mathbb R.
\]
Then $\lambda=0$. Indeed, $\left\langle F,U(s\cdot e_j)F\right\rangle_\rho=\exp(i\lambda\cdot s)\cdot\|F\|_{L^2(\rho)}^2$ with $\|F\|_{L^2(\rho)}>0$; by Proposition~\ref{prop:uniform-mixing-along} this converges as $s\to\infty$, and $\exp(i\lambda\cdot s)$ converges as $s\to\infty$ only if $\lambda=0$. Consequently, a vector of $L^2(\rho)$ that is an eigenvector of the one-parameter group $\left(U(s\cdot e_j)\right)_{s\in\mathbb R}$, with the eigenvalues $\exp(i\lambda\cdot s)$, $s\in\mathbb R$, is left fixed by every $U(s\cdot e_j)$.

In particular, let $\rho\in\mathcal M$ satisfy $(T_j)$, let $\ell>0$ and $n\in\mathbb Z\setminus\{0\}$, and let $F\in L^2(\rho)$ satisfy $|F|=1$ $\rho$-almost everywhere and
\[
U(s\cdot e_j)F=\exp\left(\frac{2\pi i\cdot n\cdot s}{\ell}\right)\cdot F\qquad\text{for every }s\in\mathbb R.
\]
(This is the situation of a law modulated periodically along $e_j$, with period $\ell$, in which the non-trivial character $\alpha\mapsto\exp\left(\frac{2\pi i\cdot n\cdot\alpha}{\ell}\right)$ of the phase $\alpha\in\mathbb R/\ell\mathbb Z$ of the modulation is a function $F$ of the field, the translation by $s\cdot e_j$ shifting the phase by $s$.) Then $\rho$ does not satisfy $(RP_j)$. Otherwise $\rho$ would satisfy $(R_j)$ by Remark~\ref{rem:reflection-invariance}, and the preceding paragraph, applied with $\lambda=\frac{2\pi n}{\ell}\ne0$, would be contradicted. A mixing modulation, by contrast, is compatible with all of~\eqref{eq:setup-H}: the media of Examples~\ref{ex:white-noise-telegraph} and~\ref{ex:mixture-free-maxwell} are built from telegraph processes, whose correlations decay exponentially by~\eqref{eq:ergodic-components-5}, and the laws $m_2$ and $m_B$ satisfy~\eqref{eq:setup-H} by Example~\ref{ex:white-noise-telegraph}(2) and Example~\ref{ex:mixture-free-maxwell}(1).
\end{remark}

\begin{remark}[Reflection invariance from $(T_j)$ and $(RP_j)$]\label{rem:reflection-invariance}
Let $j\in\{1,2,3\}$, and let $\rho\in\mathcal M$ satisfy the properties $(T_j)$ and $(RP_j)$ of Section~\ref{sec:properties-laws}. Then $\rho$ satisfies $(R_j)$. Consequently, a law $\rho\in\mathcal M$ belongs to the class $\mathcal C$ if and only if it satisfies $(T_{123})$ and $(RP_l)$ for $l=1,2,3$.
\end{remark}

\begin{proof}
Let $\mathcal V$ be a complex vector space, and let $\beta\colon\mathcal V\times\mathcal V\to\mathbb C$ be antilinear in the first and linear in the second argument, with $\beta(F,F)\ge0$ for every $F\in\mathcal V$. Then $\beta$ is Hermitian, that is, $\beta(F,G)=\overline{\beta(G,F)}$ for all $F,G\in\mathcal V$. Indeed, for $F,G\in\mathcal V$ and $c\in\mathbb C$ the number
\[
c\cdot\beta(F,G)+\overline{c}\cdot\beta(G,F)=\beta(F+cG,F+cG)-\beta(F,F)-|c|^2\cdot\beta(G,G)
\]
is real, hence equal to its complex conjugate. For $c=1$, and for $c=i$ after division by $i$, this gives
\[
\beta(F,G)+\beta(G,F)=\overline{\beta(F,G)}+\overline{\beta(G,F)},\qquad\beta(F,G)-\beta(G,F)=\overline{\beta(G,F)}-\overline{\beta(F,G)},
\]
and the sum of the two relations is $2\beta(F,G)=2\overline{\beta(G,F)}$.

Let $\mathcal A_+$ be the complex linear span of the functions $\exp\left(iX(f)\right)$, $f\in\mathcal D(\{x_j>0\})$, on $Q$. Its elements are bounded and $\Sigma$-measurable. The zero form lies in $\mathcal D(\{x_j>0\})$, so the constant function $1=\exp\left(iX(0)\right)$ belongs to $\mathcal A_+$. For $F\in\mathcal A_+$ the function $\Theta_jF$ is again bounded and $\Sigma$-measurable (Section~\ref{sec:euclidean-maps-actions}). Hence
\[
\beta(F,G)=\rho\left(\overline{\Theta_jF}\cdot G\right)\qquad\left(F,G\in\mathcal A_+\right)
\]
defines a map $\beta\colon\mathcal A_+\times\mathcal A_+\to\mathbb C$; it is antilinear in $F$ and linear in $G$, because $\Theta_j$ is complex linear. Let $G\in\mathcal A_+$, and write
\[
G=\sum_{p=1}^nc_p\cdot\exp\left(iX(f_p)\right)
\]
with $n\ge1$, $c_1,\dots,c_n\in\mathbb C$ and $f_1,\dots,f_n\in\mathcal D(\{x_j>0\})$. The operator $\Theta_j=\hat\theta_j$ is multiplicative, commutes with complex conjugation and with composition, and maps $X(f)$ to $X(\theta_j\cdot f)$ for $f\in\mathcal S$ (Section~\ref{sec:euclidean-maps-actions}). Therefore, at every point of $Q$,
\[
\overline{\Theta_jG}\cdot G=\sum_{p,q=1}^n\overline{c_p}\cdot c_q\cdot\exp\left(i\left(X(f_q)-X(\theta_j\cdot f_p)\right)\right),
\]
and integration with respect to $\rho$ gives
\[
\beta(G,G)=\sum_{p,q=1}^n\overline{c_p}\cdot c_q\cdot\rho\left(\exp\left(i\left(X(f_q)-X(\theta_j\cdot f_p)\right)\right)\right)\ge0,
\]
the inequality being $(RP_j)$. Thus $\beta(G,G)\ge0$ for every $G\in\mathcal A_+$, and $\beta$ is Hermitian by the preceding paragraph. Since $\Theta_j1=1$, the relation $\beta(1,G)=\overline{\beta(G,1)}$ reads
\[
\rho(G)=\beta(1,G)=\overline{\beta(G,1)}=\overline{\rho\left(\overline{\Theta_jG}\right)}=\rho\left(\Theta_jG\right)\qquad\left(G\in\mathcal A_+\right).
\]

Now let $G\in\mathcal A$, where $\mathcal A$ is the complex linear span of the functions $\exp\left(iX(f)\right)$, $f\in\mathcal D_0$ (Lemma~\ref{lem:density-exponentials}). Write $G=\sum_{p=1}^nc_p\cdot\exp\left(iX(f_p)\right)$ with $n\ge1$, $c_1,\dots,c_n\in\mathbb C$ and $f_1,\dots,f_n\in\mathcal D_0$. As $\mathcal D_0\subset\mathcal D$ (Lemma~\ref{lem:countable-family-test}(i)), the supports of $f_1,\dots,f_n$ are compact, and there is $r>0$ such that
\[
\operatorname{supp}f_p\subset\left\{x\in\mathbb R^3:|x|<r\right\}\qquad(p=1,\dots,n).
\]
Let $t>r$. For each $p$ the form $\tau_{t\cdot e_j}\cdot f_p$, that is $x\mapsto f_p(x-t\cdot e_j)$, is smooth, and its support is the compact set $\left\{y+t\cdot e_j:y\in\operatorname{supp}f_p\right\}$. Every point $x$ of this set satisfies $|x_j-t|\le|x-t\cdot e_j|<r$, hence $x_j>t-r>0$. So $\tau_{t\cdot e_j}\cdot f_p\in\mathcal D(\{x_j>0\})$ for $p=1,\dots,n$, and, since $U(t\cdot e_j)\exp\left(iX(f_p)\right)=\exp\left(iX(\tau_{t\cdot e_j}\cdot f_p)\right)$ (Section~\ref{sec:euclidean-maps-actions}),
\[
U(t\cdot e_j)G=\sum_{p=1}^nc_p\cdot\exp\left(iX(\tau_{t\cdot e_j}\cdot f_p)\right)\in\mathcal A_+.
\]
As $\theta_je_j=-e_j$, the Euclidean maps $\theta_j$ and $\tau_{t\cdot e_j}$ satisfy $\theta_j\tau_{t\cdot e_j}=\tau_{-t\cdot e_j}\theta_j$, and the rule $\widehat{gh}=\hat g\hat h$ of Section~\ref{sec:euclidean-maps-actions} gives
\[
\Theta_jU(t\cdot e_j)=U(-t\cdot e_j)\Theta_j.
\]
By $(T_j)$, $(\tau_a)_*\rho=\rho$ for $a=t\cdot e_j\in\Gamma_j$ and for $a=-t\cdot e_j\in\Gamma_j$; therefore $\rho\left(U(\pm t\cdot e_j)H\right)=\rho(H)$ for every bounded $\Sigma$-measurable function $H$ on $Q$ (Section~\ref{sec:euclidean-maps-actions}). It follows that
\[
\rho\left(\Theta_jG\right)=\rho\left(U(-t\cdot e_j)\Theta_jG\right)=\rho\left(\Theta_jU(t\cdot e_j)G\right)=\rho\left(U(t\cdot e_j)G\right)=\rho(G).
\]
Here the first and the last equality hold by $(T_j)$, the second by the commutation relation, and the third by the identity $\rho\left(\Theta_jG'\right)=\rho(G')$, valid for every $G'\in\mathcal A_+$, applied to $G'=U(t\cdot e_j)G$. Hence $\rho\left(\Theta_jG\right)=\rho(G)$ for every $G\in\mathcal A$.

Let $f\in\mathcal D_0$. The function $\exp\left(iX(f)\right)$ belongs to $\mathcal A$, and $\Theta_j\exp\left(iX(f)\right)=\exp\left(iX(\theta_j\cdot f)\right)$ (Section~\ref{sec:euclidean-maps-actions}). Consequently
\[
\rho\left(\exp\left(iX(\theta_j\cdot f)\right)\right)=\rho\left(\exp\left(iX(f)\right)\right)\qquad\text{for every }f\in\mathcal D_0,
\]
and by the criterion for $(R_j)$ in Lemma~\ref{lem:measurability-defining-properties} the law $\rho$ satisfies $(R_j)$.

Finally, let $\rho\in\mathcal M$ satisfy $(T_{123})$ and $(RP_l)$ for $l=1,2,3$. For each $l\in\{1,2,3\}$ one has $\Gamma_l\subset\Gamma_{123}$, so that $\rho$ satisfies $(T_l)$; by the first assertion, with $l$ in place of $j$, it satisfies $(R_l)$. Hence $\rho\in\mathcal C$. Conversely, every law in $\mathcal C$ satisfies $(T_{123})$ and $(RP_l)$ for $l=1,2,3$, by the definition of $\mathcal C$ (Section~\ref{sec:properties-laws}).
\end{proof}

\subsection{Bibliographical note}\label{sec:bibliographical-note}

The arguments in the proofs of Lemma~\ref{lem:reflection-positivity-weights} and of Proposition~\ref{prop:reconstruction-direction}(iv) are classical; they go back to Fr\"ohlich~\cite{Froehlich76} and to Glimm and Jaffe~\cite[Section~19.7]{GlimmJaffe}, where the measures are invariant under the Euclidean group. In the form given above, these arguments require of the law only the properties $(T_j)$, $(R_j)$ and $(RP_j)$.

\section{Fourth cumulants of mixtures}\label{sec:fourth-cumulants-mixtures}

For translation-invariant laws on the space $\mathcal S'(\mathbb{R}^3;\Lambda^2)$ of tempered distributions with values in two-forms that satisfy the Gaussian bound~\eqref{eq:fourth-cumulants-H1}, this section establishes bounds on moments, ergodic decompositions in the sense of Definition~\ref{def:ergodicity-ergodic-decompositions} and identities for the fourth cumulants of mixtures, and derives from them the inequalities of Theorem~\ref{thm:non-gaussian-components} for the combination $\mathrm{Ext}$ of Section~\ref{sec:cumulants-gaussian-laws} and those of Theorem~\ref{thm:non-gaussian-components-2} for three separated tests; these results constitute Estimate~9 and a part of Estimate~5. The notation of the section is its own and is fixed in Section~\ref{sec:test-forms-distributions-2}; ergodic decompositions, averages over translations, the determination of a law by its exponentials, telegraph processes and the free field are treated also in Section~\ref{sec:ergodic-components-laws}, and bounds on moments also in Section~\ref{sec:exponential-moments-cumulants}, each section proving in its own notation what it needs; the remark following Definition~\ref{def:ergodicity-ergodic-decompositions} relates the two treatments of ergodic decompositions.

In this section the letter $e$ without index always stands for the charge, whereas $e_1$, $e_2$, $e_3$ are the vectors of the standard basis of $\mathbb{R}^3$; $\exp$ denotes the exponential function, and indices $j\in\{1,2,3\}$ number the coordinate directions. The letter $\mu$ always denotes the law under study, and $E$, $E'$, $E_\xi$ denote the expectations with respect to $\mu$, to a law $\mu'$ and to the laws $\mu_\xi$ of Definition~\ref{def:ergodicity-ergodic-decompositions}, respectively. The letters $L$, $T$, $N$, $M$ are used only with the following local meanings, which bear no relation to the quantities $L$, $T$, $N$, $M$ of Sections~\ref{sec:lattice-measures-laws}, \ref{sec:running-quantities-tori} and~\ref{sec:order-choice}:
\begin{itemize}
\item $L$ is the side of an averaging cube $[0,L]^3$;
\item $T$ denotes the maps $T_g$ of the space $\Omega$ of Section~\ref{sec:test-forms-distributions-2}, the operator $T_a$ of Section~\ref{sec:free-quartic-coefficient}, and the length of an averaging interval $[0,T]$;
\item $N$ denotes null sets, invariant sets, and the two operators $N$, $N'$ in the proof of Lemma~\ref{lem:polarisation-decay-quartic}(b), while $\mathbb{N}$ is the index set of the sequence space $\mathbb{R}^{\mathbb{N}}$;
\item $M$ denotes bounds, and the matrix $M$ in the proof of Lemma~\ref{lem:gaussian-laws-bounded}(d).
\end{itemize}

\subsection{Test forms, distributions and Euclidean motions}\label{sec:test-forms-distributions-2}

Let $\Lambda^2$ be the space of real alternating two-tensors on $\mathbb{R}^3$; a two-form $f$ has the components $f_{ab}$, $1\le a<b\le3$. Let $\mathcal S=\mathcal S(\mathbb{R}^3;\Lambda^2)$ be the space of real Schwartz two-forms $f=(f_{12},f_{13},f_{23})$, and let $\mathcal D=C_c^\infty(\mathbb{R}^3;\Lambda^2)$ be its subspace of smooth two-forms with compact support. For $f,g\in\mathcal S$ and $x\in\mathbb{R}^3$ put
\[
\langle f,g\rangle=\sum_{a<b}\int\mathrm{d}x\,f_{ab}(x)\cdot g_{ab}(x),\qquad |f(x)|^2=\sum_{a<b}f_{ab}(x)^2,
\]
and let $\|f\|_p$ denote the $L^p$ norm of the function $x\mapsto|f(x)|$.

Let $\Omega=\mathcal S'$ be the real dual of $\mathcal S$, that is, the space of continuous linear maps $\omega\colon\mathcal S\to\mathbb{R}$. For $f\in\mathcal S$ the function $X(f)\colon\Omega\to\mathbb{R}$ is given by $X(f)(\omega)=\omega(f)$, and
\[
\mathcal F=\sigma\left(X(f):f\in\mathcal S\right)
\]
is the cylinder $\sigma$-algebra of $\Omega$.

A Euclidean motion is a pair $(b,r)$ with $b\in\mathbb{R}^3$ and $r\in O(3)$; it acts on $\mathbb{R}^3$ by $x\mapsto rx+b$, and on $\mathcal S$ by
\[
(U_{b,r}f)(x)=(\wedge^2r)\,f\left(r^{-1}(x-b)\right),
\]
where $\wedge^2r$ is the linear map of $\Lambda^2$ induced by $r$. For $x\in\mathbb{R}^3$ the translation is $\tau_x=U_{x,1}$, where $1$ is the identity of $O(3)$; thus $(\tau_xf)(y)=f(y-x)$. For $j\in\{1,2,3\}$ let $\theta_j\in O(3)$ be the reflection $x_j\mapsto-x_j$ of the $j$-th coordinate, and put $\theta_jf=U_{0,\theta_j}f$ for $f\in\mathcal S$; in components,
\[
(\theta_jf)_{ab}(x)=s\cdot f_{ab}(\theta_jx),\qquad s=\begin{cases}-1&\text{if }j\in\{a,b\},\\+1&\text{otherwise}.\end{cases}
\]
The cubic group $O_{\mathrm{cub}}\subset O(3)$ is the group of the signed permutations of the coordinates.

For a Euclidean motion $g$ the map $T_g\colon\Omega\to\Omega$ is given by $(T_g\omega)(f)=\omega(U_g^{-1}f)$, and for a function $G$ on $\Omega$ one puts $\mathbb{U}_gG=G\circ T_g^{-1}$. Since $T_g^{-1}\omega=\omega\circ U_g$, one has $\mathbb{U}_gX(f)=X(U_gf)$. For translations one writes $T_x=T_{x,1}$ and $\mathbb{U}_x=\mathbb{U}_{x,1}$. For $j\in\{1,2,3\}$ put
\[
\Theta_jG=\overline{\mathbb{U}_{0,\theta_j}G},\qquad\Sigma_j^+=\sigma\left(X(f):f\in\mathcal D,\ \operatorname{supp}f\subset\{x_j>0\}\right),
\]
and let $A_j^+$ be the complex linear span of the functions $\exp\left(iX(f)\right)$ with $f\in\mathcal D$ and $\operatorname{supp}f\subset\{x_j>0\}$.

Finally, let $\mu'$ be a law on $(\Omega,\mathcal F)$, that is, a probability measure, with expectation $E'$. For $G,H\in L^2(\mu')$ put
\[
\langle G,H\rangle=E'\left[\overline{G}\cdot H\right],\qquad\|G\|=\|G\|_2=\langle G,G\rangle^{\frac{1}{2}}.
\]

\subsection{Hypotheses on the law}\label{sec:hypotheses-law}

Let $\mathbb{R}^3\rtimes O_{\mathrm{cub}}$ denote the group of the Euclidean motions $(b,r)$ with $r\in O_{\mathrm{cub}}$. The law $\mu$ is a probability measure on $(\Omega,\mathcal F)$ with the following three properties:
\begin{equation}\tag{H1}\label{eq:fourth-cumulants-H1}
\int\mathrm{d}\mu\,\exp X(f)\le\exp\left(\frac{1}{2}\|f\|_2^2\right)\quad\text{for every }f\in\mathcal S;
\end{equation}
\begin{equation}\tag{H2}\label{eq:fourth-cumulants-H2}
\mu\circ T_g^{-1}=\mu\quad\text{for every }g\in\mathbb{R}^3\rtimes O_{\mathrm{cub}};
\end{equation}
\begin{equation}\tag{H3}\label{eq:fourth-cumulants-H3}
\int\mathrm{d}\mu\,(\Theta_jG)\cdot G\ge0\quad\text{for }j=1,2,3\text{ and every }G\in A_j^+.
\end{equation}
Together, \eqref{eq:fourth-cumulants-H1}--\eqref{eq:fourth-cumulants-H3} are the hypothesis~\eqref{eq:setup-H} of Section~\ref{sec:hypothesis-law}.

By \eqref{eq:fourth-cumulants-H2}, the inequality \eqref{eq:fourth-cumulants-H3} holds at every plane $\{x_j=c\}$, $c\in\mathbb{R}$, in place of the plane $\{x_j=0\}$. Precisely, let $(2c\cdot e_j,\theta_j)$ be the reflection $x\mapsto\theta_jx+2c\cdot e_j$ in the plane $\{x_j=c\}$; then
\[
\int\mathrm{d}\mu\,\overline{\mathbb{U}_{2ce_j,\theta_j}G}\cdot G\ge0
\]
for every $G$ in the complex linear span of the functions $\exp\left(iX(f)\right)$ with $f\in\mathcal D$ and $\operatorname{supp}f\subset\{x_j>c\}$. Indeed, the translation $\tau_{-ce_j}$ maps the forms supported in $\{x_j>c\}$ onto those supported in $\{x_j>0\}$, so that $G'=\mathbb{U}_{-ce_j}G$ lies in $A_j^+$; since $U_{2ce_j,\theta_j}=\tau_{ce_j}\,\theta_j\,\tau_{-ce_j}$ on $\mathcal S$, one has
\[
\overline{\mathbb{U}_{2ce_j,\theta_j}G}\cdot G=\mathbb{U}_{ce_j}\left[(\Theta_jG')\cdot G'\right];
\]
and by \eqref{eq:fourth-cumulants-H2} the integral of the right side with respect to $\mu$ equals $\int\mathrm{d}\mu\,(\Theta_jG')\cdot G'$, which is non-negative by \eqref{eq:fourth-cumulants-H3}.

The following condition on $\mu$ is an additional hypothesis, assumed explicitly wherever it is used. Let $g\in\mathcal S$, let $e>0$, $\sigma>0$ and $C_g<\infty$ be numbers, and let $b(g)$ be a real number. Condition \eqref{eq:fourth-cumulants-4g} reads
\begin{equation}\tag{$\kappa_4$;\,$g$}\label{eq:fourth-cumulants-4g}
\left|\kappa_4^\mu(g)-e^4\cdot b(g)\right|\le C_g\cdot e^{4+\sigma},
\end{equation}
where $\kappa_4^\mu(g)=E\left[X(g)^4\right]-3\left(E\left[X(g)^2\right]\right)^2$ is the fourth cumulant of $X(g)$ under $\mu$, in the notation of Section~\ref{sec:cumulants-gaussian-laws}. For every limit point $\mu$ of the lattice laws of Theorem~\ref{thm:main-theorem}, condition \eqref{eq:fourth-cumulants-4g} holds at each of the three tests $g$ of Theorem~\ref{thm:non-gaussian-components-2}; this is Theorem~\ref{thm:main-theorem}(iii).

\subsection{Cumulants and Gaussian laws}\label{sec:cumulants-gaussian-laws}

Let $\mu'$ be a law on $(\Omega,\mathcal F)$ with expectation $E'$ such that $E'\left[X(f)^4\right]<\infty$ and $E'X(f)=0$ for every $f\in\mathcal S$; by H\"older's inequality all expectations below are finite. For $f,g,f_1,f_2,f_3,f_4\in\mathcal S$ and $x\in\mathbb{R}^3$ put
\begin{align*}
S_2'(f,g)&=E'\left[X(f)\cdot X(g)\right],\\
\kappa_4'(f_1,f_2,f_3,f_4)&=E'\left[X(f_1)\cdot X(f_2)\cdot X(f_3)\cdot X(f_4)\right]-S_2'(f_1,f_2)\cdot S_2'(f_3,f_4)\\
&\qquad-S_2'(f_1,f_3)\cdot S_2'(f_2,f_4)-S_2'(f_1,f_4)\cdot S_2'(f_2,f_3),\\
\kappa_4'(f)&=\kappa_4'(f,f,f,f)=E'\left[X(f)^4\right]-3\left(E'\left[X(f)^2\right]\right)^2,\\
\mathrm{Ext}'_x(f)&=\kappa_4'(f)-3\,\kappa_4'(f,f,\tau_xf,\tau_xf).
\end{align*}
For a law $\lambda$ on $(\Omega,\mathcal F)$ with the same two properties, $S_2^\lambda$, $\kappa_4^\lambda$ and $\mathrm{Ext}^\lambda_x$ denote the quantities $S_2'$, $\kappa_4'$ and $\mathrm{Ext}'_x$ formed with $\lambda$ in place of $\mu'$. If $(\Xi,\nu,(\mu_\xi))$ is an ergodic decomposition of $\mu$ under $\mathbb{R}^3$ (Definition~\ref{def:ergodicity-ergodic-decompositions}), so that $\mu=\int\mathrm{d}\nu(\xi)\,\mu_\xi$, one writes
\[
C_\xi=S_2^{\mu_\xi},\qquad\kappa_4^\xi=\kappa_4^{\mu_\xi},\qquad\mathrm{Ext}^\xi_x=\mathrm{Ext}^{\mu_\xi}_x
\]
for every $\xi$ such that $\mu_\xi$ has the two properties.

The form $\kappa_4'$ is symmetric and $4$-linear on $\mathcal S$: the function $X(f)$ depends linearly on $f$, and a permutation of $f_1,\dots,f_4$ permutes the three products of values of $S_2'$. In the expansions of $\kappa_4'(f+g)$ and of $\kappa_4'(f-g)$ the terms in which $g$ occurs an odd number of times have opposite signs, which gives the polarisation identity
\[
\kappa_4'(f+g)+\kappa_4'(f-g)=2\,\kappa_4'(f)+2\,\kappa_4'(g)+12\,\kappa_4'(f,f,g,g)\qquad(f,g\in\mathcal S).
\]
Let $\mu'$ moreover be translation invariant, that is, $\mu'\circ T_x^{-1}=\mu'$ for every $x\in\mathbb{R}^3$. Then $X(\tau_xf)=\mathbb{U}_xX(f)$ has under $\mu'$ the same law as $X(f)$, so that $\kappa_4'(\tau_xf)=\kappa_4'(f)$, and the polarisation identity with $g=\tau_xf$ becomes
\[
12\,\kappa_4'(f,f,\tau_xf,\tau_xf)=\kappa_4'(f+\tau_xf)+\kappa_4'(f-\tau_xf)-4\,\kappa_4'(f).
\]
Inserting this into the definition of $\mathrm{Ext}'_x$ one obtains
\begin{equation}\label{eq:fourth-cumulants-1}
\mathrm{Ext}'_x(f)=2\,\kappa_4'(f)-\frac{1}{4}\left[\kappa_4'(f+\tau_xf)+\kappa_4'(f-\tau_xf)\right]\qquad(f\in\mathcal S,\ x\in\mathbb{R}^3).
\end{equation}

A law $\mu'$ on $(\Omega,\mathcal F)$ is called Gaussian if $X(f)$ is a normal random variable, possibly degenerate, for every $f\in\mathcal S$. In that case every finite family $\left(X(f_1),\dots,X(f_n)\right)$ is jointly normal, since every linear combination $\sum_it_i\cdot X(f_i)=X\left(\sum_it_i\cdot f_i\right)$ is normal; in particular all moments are finite. If moreover $E'X(f)=0$ for every $f\in\mathcal S$, as is the case for every law that satisfies \eqref{eq:fourth-cumulants-H1} in place of $\mu$, by Lemma~\ref{lem:moment-bounds}(a), then Wick's formula gives $\kappa_4'(f_1,f_2,f_3,f_4)=0$ for all $f_1,\dots,f_4\in\mathcal S$, and consequently $\mathrm{Ext}'_x(f)=0$ for all $f\in\mathcal S$ and $x\in\mathbb{R}^3$.

Thus $\mathrm{Ext}'_x$ vanishes identically for every centred Gaussian law.

Let $\mu'$ be a translation-invariant law on $(\Omega,\mathcal F)$ with expectation $E'$ such that $E'\left[X(f)^4\right]<\infty$ and $E'X(f)=0$ for every $f\in\mathcal S$, and suppose that $\mu'$ satisfies the clustering condition~\eqref{eq:fourth-cumulants-C1} of Theorem~\ref{thm:clustering-mixing-ergodicity}, with $\mu'$ in place of $\mu$. Then
\[
\lim_{|x|\to\infty}\mathrm{Ext}'_x(f)=\kappa_4'(f)\qquad\text{for every }f\in\mathcal D.
\]
Indeed, let $f\in\mathcal D$. Since $\mu'$ is translation invariant, $S_2'(\tau_xf,\tau_xf)=S_2'(f,f)$, and the definition of $\kappa_4'$ becomes
\[
\kappa_4'(f,f,\tau_xf,\tau_xf)=E'\left[X(f)^2\cdot X(\tau_xf)^2\right]-S_2'(f,f)^2-2\,S_2'(f,\tau_xf)^2.
\]
As $|x|\to\infty$, condition~\eqref{eq:fourth-cumulants-C1} with $m=n=2$, $f_1=f_2=g_1=g_2=f$ and $a=x$ shows that the first term on the right converges to $\left(E'\left[X(f)^2\right]\right)^2=S_2'(f,f)^2$, and the same condition with $m=n=1$, $f_1=g_1=f$ and $a=x$ shows that $S_2'(f,\tau_xf)=E'\left[X(f)\cdot X(\tau_xf)\right]$ converges to $\left(E'X(f)\right)^2=0$. Consequently $\kappa_4'(f,f,\tau_xf,\tau_xf)\to0$, and $\mathrm{Ext}'_x(f)=\kappa_4'(f)-3\,\kappa_4'(f,f,\tau_xf,\tau_xf)\to\kappa_4'(f)$.

Let $\mu'$ be a law on $(\Omega,\mathcal F)$. A \emph{representation of $\mu'$ as a mixture} is a probability space $(\Xi,B,\nu)$ together with a family $(\mu_\xi)_{\xi\in\Xi}$ of laws on $(\Omega,\mathcal F)$ which satisfies~\eqref{eq:fourth-cumulants-ED1} of Definition~\ref{def:ergodicity-ergodic-decompositions} with $\mu'$ in place of $\mu$: for every $A\in\mathcal F$ the function $\xi\mapsto\mu_\xi(A)$ is measurable and $\mu'(A)=\int\mathrm{d}\nu(\xi)\,\mu_\xi(A)$. The law $\mu'$ is called a \emph{mixture of laws with a given property} if it has a representation as a mixture in which every $\mu_\xi$, $\xi\in\Xi$, has this property. In a representation as a mixture let $E_\xi$ be the expectation with respect to $\mu_\xi$. By linearity and monotone convergence, \eqref{eq:fourth-cumulants-ED1} extends from indicator functions to every measurable function $G\colon\Omega\to[0,\infty]$: the function $\xi\mapsto E_\xi[G]$, with values in $[0,\infty]$, is measurable, and
\[
E'[G]=\int\mathrm{d}\nu(\xi)\,E_\xi[G].
\]
For a real function $h\in L^2(\nu)$ put $\operatorname{Var}_\nu(h)=\int\mathrm{d}\nu(\xi)\,h(\xi)^2-\left(\int\mathrm{d}\nu(\xi)\,h(\xi)\right)^2$, which is non-negative by the Cauchy--Schwarz inequality.

Now let $\mu'$ be a law on $(\Omega,\mathcal F)$ with expectation $E'$ such that $E'\left[X(f)^4\right]<\infty$ and $E'X(f)=0$ for every $f\in\mathcal S$, and suppose that $\mu'$ is a mixture of translation-invariant centred Gaussian laws; let $(\Xi,B,\nu)$, $(\mu_\xi)_{\xi\in\Xi}$ be a representation of $\mu'$ as a mixture in which every $\mu_\xi$ is such a law. Then
\[
\mathrm{Ext}'_x(f)=-6\operatorname{Var}_\nu\left(\xi\mapsto S_2^{\mu_\xi}(f,\tau_xf)\right)\le0\qquad(f\in\mathcal S,\ x\in\mathbb{R}^3).
\]
For the proof fix $f\in\mathcal S$ and $x\in\mathbb{R}^3$, put $g=\tau_xf$, and define functions $y$, $z$ on $\Xi$ by $y(\xi)=E_\xi\left[X(f)^2\right]$ and $z(\xi)=E_\xi\left[X(f)\cdot X(g)\right]=S_2^{\mu_\xi}(f,g)$. Each $\mu_\xi$ is a centred Gaussian law; hence all its moments are finite, and Wick's formula gives $\kappa_4^{\mu_\xi}(f_1,f_2,f_3,f_4)=0$ for all $f_1,\dots,f_4\in\mathcal S$. As $\mu_\xi$ is translation invariant, $E_\xi\left[X(g)^2\right]=y(\xi)$. The vanishing of $\kappa_4^{\mu_\xi}$ at $(f,f,f,f)$ and at $(f,f,g,g)$ therefore reads
\[
E_\xi\left[X(f)^4\right]=3\,y(\xi)^2,\qquad E_\xi\left[X(f)^2\cdot X(g)^2\right]=y(\xi)^2+2\,z(\xi)^2\qquad(\xi\in\Xi),
\]
and the Cauchy--Schwarz inequality gives $|z(\xi)|\le y(\xi)$. Since $X(f)\cdot X(g)=\frac{1}{4}\left(X(f+g)^2-X(f-g)^2\right)$, the extended form of~\eqref{eq:fourth-cumulants-ED1} shows that $y$ and
\[
z(\xi)=\frac{1}{4}\left(E_\xi\left[X(f+g)^2\right]-E_\xi\left[X(f-g)^2\right]\right)
\]
are measurable functions of $\xi$; applied to $G=X(f)^4$ it gives $3\int\mathrm{d}\nu(\xi)\,y(\xi)^2=E'\left[X(f)^4\right]<\infty$, so that $y$ and $z$ belong to $L^2(\nu)$. Applied to $G=X(f)^2$, $X(g)^2$, $X(f)^2\cdot X(g)^2$, $X(f+g)^2$ and $X(f-g)^2$, whose $\mu'$-expectations are finite, it gives
\begin{gather*}
S_2'(f,f)=S_2'(g,g)=\int\mathrm{d}\nu(\xi)\,y(\xi),\qquad S_2'(f,g)=\int\mathrm{d}\nu(\xi)\,z(\xi),\\
E'\left[X(f)^2\cdot X(g)^2\right]=\int\mathrm{d}\nu(\xi)\,\left(y(\xi)^2+2\,z(\xi)^2\right).
\end{gather*}
Inserting these values and $E'\left[X(f)^4\right]=3\int\mathrm{d}\nu(\xi)\,y(\xi)^2$ into the definitions of $\kappa_4'(f)$ and $\kappa_4'(f,f,g,g)$ one obtains
\begin{align*}
\kappa_4'(f)&=E'\left[X(f)^4\right]-3\,S_2'(f,f)^2=3\operatorname{Var}_\nu(y),\\
\kappa_4'(f,f,g,g)&=E'\left[X(f)^2\cdot X(g)^2\right]-S_2'(f,f)\cdot S_2'(g,g)-2\,S_2'(f,g)^2=\operatorname{Var}_\nu(y)+2\operatorname{Var}_\nu(z),
\end{align*}
and hence $\mathrm{Ext}'_x(f)=\kappa_4'(f)-3\,\kappa_4'(f,f,g,g)=-6\operatorname{Var}_\nu(z)\le0$.

\subsection{Moments and ergodic decompositions}\label{sec:moments-ergodic-decompositions}

\begin{lemma}[Moment bounds]\label{lem:moment-bounds}
Let $\mu'$ be a probability measure on $(\Omega,\mathcal F)$ with expectation $E'$, and assume that $\mu'$ satisfies \eqref{eq:fourth-cumulants-H1} in place of $\mu$, that is,
\[
E'\exp\left(X(f)\right)\le\exp\left(\frac{\|f\|_2^2}{2}\right)\qquad\text{for every }f\in\mathcal S.
\]
For every integer $m\ge1$ put
\[
c_{2m}=(2m)!\cdot\left(\frac{\exp(1)}{2m}\right)^m.
\]
Let $n\ge1$ and $k\ge1$ be integers, let $t\ge0$, and let $f,f_1,\dots,f_n\in\mathcal S$. Then the following hold.
\begin{enumerate}
\item[(a)] $E'X(f)=0$.
\item[(b)] $E'\left[X(f)^2\right]\le\|f\|_2^2$.
\item[(c)] $E'\left[X(f)^{2k}\right]\le c_{2k}\cdot\|f\|_2^{2k}$; moreover $c_4=\frac{3}{2}\exp(2)<11.09$ and $c_8<537.5$.
\item[(d)] $E'\exp\left(t\cdot|X(f)|\right)\le2\exp\left(\frac{t^2\cdot\|f\|_2^2}{2}\right)$.
\item[(e)] One has
\[
\left|E'\prod_{i=1}^nX(f_i)\right|\le c_{2n}^{\frac{1}{2}}\cdot\prod_{i=1}^n\|f_i\|_2 .
\]
In particular, every joint moment and every joint cumulant, under $\mu'$, of finitely many of the variables $X(g)$, $g\in\mathcal S$, is a continuous function of each of the tests $g$ that enter it, $\mathcal S$ being given the norm $\|\cdot\|_2$.
\end{enumerate}
\end{lemma}

\begin{proof}
Fix $f\in\mathcal S$ and write $X=X(f)$. For each $\omega\in\Omega$ the map $g\mapsto X(g)(\omega)=\omega(g)$ is linear on $\mathcal S$; hence $X(s\cdot f)=s\cdot X$ for $s\in\mathbb R$, and the hypothesis, applied to the tests $s\cdot f$, gives
\[
E'\exp(s\cdot X)\le\exp\left(\frac{s^2\cdot\|f\|_2^2}{2}\right)\qquad(s\in\mathbb R).
\]
Let $t\ge0$. Pointwise on $\Omega$ one has $\exp(t\cdot|X|)\le\exp(t\cdot X)+\exp(-t\cdot X)$ and $\cosh(t\cdot X)=\frac{1}{2}\left(\exp(t\cdot X)+\exp(-t\cdot X)\right)$. The cases $s=t$ and $s=-t$ of the preceding bound therefore yield
\[
E'\exp\left(t\cdot|X|\right)\le2\exp\left(\frac{t^2\cdot\|f\|_2^2}{2}\right),\qquad E'\cosh(t\cdot X)\le\exp\left(\frac{t^2\cdot\|f\|_2^2}{2}\right).
\]
The first of these two inequalities is (d).

(a) Since $|X|\le\exp(|X|)$, the inequality (d) with $t=1$ shows that $X$ is integrable. By Jensen's inequality for the convex function $\exp$ and by the bound for $E'\exp(s\cdot X)$,
\[
\exp\left(s\cdot E'X\right)\le E'\exp(s\cdot X)\le\exp\left(\frac{s^2\cdot\|f\|_2^2}{2}\right)\qquad(s\in\mathbb R),
\]
that is, $s\cdot E'X\le\frac{1}{2}\cdot s^2\cdot\|f\|_2^2$ for every real $s$. Division by $s>0$ and the limit $s\downarrow0$ give $E'X\le0$; division by $s<0$ and the limit $s\uparrow0$ give $E'X\ge0$. Hence $E'X=0$.

(b), (c) For real $u$ the series $\cosh u=\sum_{j\ge0}\frac{u^{2j}}{(2j)!}$ has nonnegative terms, so that
\[
\cosh u\ge1+\frac{u^2}{2},\qquad\cosh u\ge\frac{u^{2k}}{(2k)!}\qquad(u\in\mathbb R).
\]
Let $t>0$. Substituting $u=t\cdot X$ in these two inequalities, taking expectations and using the bound for $E'\cosh(t\cdot X)$, one obtains
\[
1+\frac{t^2}{2}\cdot E'\left[X^2\right]\le\exp\left(\frac{t^2\cdot\|f\|_2^2}{2}\right),\qquad\frac{t^{2k}}{(2k)!}\cdot E'\left[X^{2k}\right]\le\exp\left(\frac{t^2\cdot\|f\|_2^2}{2}\right).
\]
The first of these inequalities states that
\[
E'\left[X^2\right]\le\frac{2}{t^2}\cdot\left(\exp\left(\frac{t^2\cdot\|f\|_2^2}{2}\right)-1\right)\qquad(t>0).
\]
As $t\to0$ the right side tends to $\|f\|_2^2$, because the derivative of $\exp$ at $0$ equals $1$; this proves (b). The second states that
\[
E'\left[X^{2k}\right]\le(2k)!\cdot t^{-2k}\cdot\exp\left(\frac{t^2\cdot\|f\|_2^2}{2}\right)\qquad(t>0).
\]
If $\|f\|_2=0$, the limit $t\to\infty$ in this inequality gives $E'\left[X^{2k}\right]=0$, and (c) holds. If $\|f\|_2>0$, choose $t>0$ with $t^2\cdot\|f\|_2^2=2k$; then $t^{-2k}=(2k)^{-k}\cdot\|f\|_2^{2k}$ and the exponential on the right equals $\exp(k)$, whence
\[
E'\left[X^{2k}\right]\le(2k)!\cdot(2k)^{-k}\cdot\exp(k)\cdot\|f\|_2^{2k}=c_{2k}\cdot\|f\|_2^{2k}.
\]
Finally, $c_4=\frac{4!}{4^2}\exp(2)=\frac{3}{2}\exp(2)$ and $c_8=\frac{8!}{8^4}\exp(4)=\frac{315}{32}\exp(4)$; since $\exp(2)<7.3891$ and $\exp(4)<54.5982$, one obtains $c_4<11.09$ and $c_8<537.5$.

(e) Let $f_1,\dots,f_n\in\mathcal S$. H\"older's inequality with $n$ exponents equal to $n$, the monotonicity in $p$ of the norm of $L^p(\mu')$ (the measure $\mu'$ has total mass $1$), and (c) with $k=n$ give
\begin{align*}
\left|E'\prod_{i=1}^nX(f_i)\right|
&\le\prod_{i=1}^n\|X(f_i)\|_{L^n(\mu')}\le\prod_{i=1}^n\|X(f_i)\|_{L^{2n}(\mu')}\\
&\le\prod_{i=1}^n\left(c_{2n}^{\frac{1}{2n}}\cdot\|f_i\|_2\right)=c_{2n}^{\frac{1}{2}}\cdot\prod_{i=1}^n\|f_i\|_2 .
\end{align*}
Define $\Psi_n\colon\mathcal S^n\to\mathbb R$ by $\Psi_n(g_1,\dots,g_n)=E'\prod_{i=1}^nX(g_i)$. Since $g\mapsto X(g)$ is linear, $\Psi_n$ is $n$-linear, and by the inequality just proved $\left|\Psi_n(g_1,\dots,g_n)\right|\le c_{2n}^{\frac{1}{2}}\cdot\prod_{i=1}^n\|g_i\|_2$. For $g_1,\dots,g_n,h_1,\dots,h_n\in\mathcal S$ the difference $\Psi_n(g_1,\dots,g_n)-\Psi_n(h_1,\dots,h_n)$ is the sum over $j=1,\dots,n$ of the terms $\Psi_n(h_1,\dots,h_{j-1},g_j-h_j,g_{j+1},\dots,g_n)$; therefore
\[
\left|\Psi_n(g_1,\dots,g_n)-\Psi_n(h_1,\dots,h_n)\right|\le c_{2n}^{\frac{1}{2}}\cdot\sum_{j=1}^n\left(\prod_{i<j}\|h_i\|_2\right)\cdot\|g_j-h_j\|_2\cdot\left(\prod_{i>j}\|g_i\|_2\right),
\]
and $\Psi_n$ is continuous on $\mathcal S^n$, each factor $\mathcal S$ carrying the norm $\|\cdot\|_2$. A joint moment of finitely many variables $X(g)$, $g\in\mathcal S$, has the form $\Psi_n(g_1,\dots,g_n)$, where each $g_i$ is one of the tests involved, a test being repeated as often as its variable occurs; hence it is a continuous function of these tests. A joint cumulant of finitely many such variables is a polynomial in their joint moments, and is therefore continuous in the same sense. This proves the last assertion of (e).
\end{proof}

For measurable maps $u,h\colon\mathbb{R}^3\to\Lambda^2$ the pointwise norm $|u(y)|$, the norms $\|u\|_p\in[0,\infty]$ and, when $\|u\|_2$ and $\|h\|_2$ are finite, the scalar product $\langle u,h\rangle$ are defined by the formulas of Section~\ref{sec:test-forms-distributions-2}; for $x\in\mathbb{R}^3$ the translation $\tau_x$ acts on such maps by $(\tau_xu)(y)=u(y-x)$. The space $L^2(\mathbb{R}^3;\Lambda^2)$ of the maps $u$ with $\|u\|_2<\infty$, two maps being identified when they agree almost everywhere, is a real Hilbert space with the scalar product $\langle\cdot,\cdot\rangle$ and the norm $\|\cdot\|_2$. It contains $\mathcal S$, and every $\tau_x$ is a linear bijection of it onto itself that preserves $\|\cdot\|_2$.

\begin{lemma}[Decay of the two-point function]\label{lem:decay-two-point}
Let $\mu'$ be a probability law on $(\Omega,\mathcal F)$ with expectation $E'$ which satisfies \eqref{eq:fourth-cumulants-H1} in place of $\mu$. By Lemma~\ref{lem:moment-bounds}(a),(c), one has $E'X(h)=0$ and $E'\left[X(h)^4\right]<\infty$ for every $h\in\mathcal S$, so that the form $S_2'$ of Section~\ref{sec:cumulants-gaussian-laws} is defined on $\mathcal S\times\mathcal S$. Let $f,g\in\mathcal S$. Then the following hold.
\begin{enumerate}
\item[(a)] There is $u_f\in L^2(\mathbb{R}^3;\Lambda^2)$ with $\|u_f\|_2\le\|f\|_2$ such that
\[
S_2'(f,h)=\langle u_f,h\rangle\qquad\text{for every }h\in\mathcal S.
\]
\item[(b)] $S_2'(f,\tau_xg)\to0$ as $|x|\to\infty$.
\item[(c)] One has
\[
\int_{\mathbb{R}^3}\mathrm{d}x\,S_2'(f,\tau_xg)^2\le\|f\|_2^2\cdot\|g\|_1^2 .
\]
\end{enumerate}
\end{lemma}

\begin{proof}
(a) Let $h\in\mathcal S$. The Cauchy--Schwarz inequality for $\mu'$, followed by Lemma~\ref{lem:moment-bounds}(b) for the tests $f$ and $h$, gives
\[
\left|S_2'(f,h)\right|=\left|E'\left[X(f)\cdot X(h)\right]\right|\le\left(E'\left[X(f)^2\right]\right)^{\frac{1}{2}}\cdot\left(E'\left[X(h)^2\right]\right)^{\frac{1}{2}}\le\|f\|_2\cdot\|h\|_2 .
\]
The map $h\mapsto S_2'(f,h)$ is linear on $\mathcal S$, since $X(h)(\omega)=\omega(h)$ depends linearly on $h$ for every $\omega\in\Omega$. The space $\mathcal S$ contains $\mathcal D$, which is dense in $L^2(\mathbb{R}^3;\Lambda^2)$. By the preceding inequality, the map $h\mapsto S_2'(f,h)$ therefore extends uniquely by continuity to a linear functional $\ell_f\colon L^2(\mathbb{R}^3;\Lambda^2)\to\mathbb{R}$, and
\[
\left|\ell_f(h)\right|\le\|f\|_2\cdot\|h\|_2\qquad\text{for every }h\in L^2(\mathbb{R}^3;\Lambda^2).
\]
By the Riesz representation theorem in the real Hilbert space $L^2(\mathbb{R}^3;\Lambda^2)$ there is $u_f\in L^2(\mathbb{R}^3;\Lambda^2)$ with $\ell_f(h)=\langle u_f,h\rangle$ for every $h\in L^2(\mathbb{R}^3;\Lambda^2)$, and $\|u_f\|_2$ equals the norm of $\ell_f$, which is at most $\|f\|_2$. Since $\ell_f(h)=S_2'(f,h)$ for $h\in\mathcal S$, this proves (a).

(b) Let $u_f$ be as in (a), and let $\varepsilon>0$. Choose $R>0$ and maps $u',g'\colon\mathbb{R}^3\to\Lambda^2$ with $\|u'\|_2<\infty$ and $\|g'\|_2<\infty$ which vanish at every point outside the ball $B_R=\{y\in\mathbb{R}^3:|y|\le R\}$ and satisfy
\[
\|u_f-u'\|_2<\varepsilon,\qquad\|g-g'\|_2<\varepsilon .
\]
Such maps exist: the products of a representative of $u_f$ and of $g$ with the indicator function of $B_R$ have these properties for all sufficiently large $R$, because, by dominated convergence, $\int_{|y|>R}\mathrm{d}y\,|u_f(y)|^2$ and $\int_{|y|>R}\mathrm{d}y\,|g(y)|^2$ tend to $0$ as $R\to\infty$. By the triangle inequality, $\|u'\|_2\le\|u_f\|_2+\varepsilon$.

Let $x\in\mathbb{R}^3$ with $|x|>2R$, and let $y\in\mathbb{R}^3$. If $|y|>R$, then $u'(y)=0$. If $|y|\le R$, then $|y-x|\ge|x|-|y|>R$, and so $(\tau_xg')(y)=g'(y-x)=0$. Hence $u'_{ab}(y)\cdot(\tau_xg')_{ab}(y)=0$ for all $y\in\mathbb{R}^3$ and $a<b$, and $\langle u',\tau_xg'\rangle=0$. The form $\tau_xg$ belongs to $\mathcal S$, and $\tau_x$ is linear; by (a), applied to $h=\tau_xg$,
\[
S_2'(f,\tau_xg)=\langle u_f,\tau_xg\rangle=\langle u_f-u',\tau_xg\rangle+\langle u',\tau_x(g-g')\rangle+\langle u',\tau_xg'\rangle,
\]
and the last term vanishes. As $\tau_x$ preserves $\|\cdot\|_2$, the Cauchy--Schwarz inequality in $L^2(\mathbb{R}^3;\Lambda^2)$ gives
\[
\left|S_2'(f,\tau_xg)\right|\le\|u_f-u'\|_2\cdot\|g\|_2+\|u'\|_2\cdot\|g-g'\|_2\le\varepsilon\cdot\|g\|_2+\|u'\|_2\cdot\varepsilon\le\varepsilon\cdot\left(\|g\|_2+\|u_f\|_2+\varepsilon\right)
\]
for every $x\in\mathbb{R}^3$ with $|x|>2R$. Consequently $\limsup_{|x|\to\infty}\left|S_2'(f,\tau_xg)\right|\le\varepsilon\cdot\left(\|g\|_2+\|u_f\|_2+\varepsilon\right)$ for every $\varepsilon>0$, which proves (b).

(c) Let $u_f$ be as in (a). The function $x\mapsto S_2'(f,\tau_xg)$ is continuous on $\mathbb{R}^3$, so that the integral in (c) is defined: for $x,x'\in\mathbb{R}^3$ one has $\tau_xg-\tau_{x'}g=\tau_{x'}\left(\tau_{x-x'}g-g\right)$, the translation $\tau_{x'}$ preserves $\|\cdot\|_2$, and the first display of the proof of (a), applied to $h=\tau_xg-\tau_{x'}g$, together with the linearity of $S_2'$ in its second argument, gives
\[
\left|S_2'(f,\tau_xg)-S_2'(f,\tau_{x'}g)\right|\le\|f\|_2\cdot\left\|\tau_{x-x'}g-g\right\|_2 ,
\]
which tends to $0$ as $x'\to x$, by the continuity of translations in $L^2(\mathbb{R}^3;\Lambda^2)$.

Let $x\in\mathbb{R}^3$. For every $y\in\mathbb{R}^3$ the Cauchy--Schwarz inequality in $\Lambda^2$, with the norm $|\cdot|$, gives
\[
\left|\sum_{a<b}(u_f)_{ab}(y)\cdot g_{ab}(y-x)\right|\le|u_f(y)|\cdot|g(y-x)| .
\]
Let $\tilde g\colon\mathbb{R}^3\to\Lambda^2$ be the map $\tilde g(z)=g(-z)$, and let $|u_f|$ and $|\tilde g|$ denote the non-negative functions $y\mapsto|u_f(y)|$ and $z\mapsto|\tilde g(z)|$ on $\mathbb{R}^3$; then $\left\||u_f|\right\|_2=\|u_f\|_2$ and $\left\||\tilde g|\right\|_1=\|g\|_1$. For non-negative measurable functions $\varphi$, $\psi$ on $\mathbb{R}^3$, whose $L^p$ norms are denoted by $\|\varphi\|_p$ and $\|\psi\|_p$, let $\varphi*\psi$ be their convolution,
\[
(\varphi*\psi)(x)=\int_{\mathbb{R}^3}\mathrm{d}y\,\varphi(y)\cdot\psi(x-y)\in[0,\infty]\qquad(x\in\mathbb{R}^3).
\]
Since $|g(y-x)|=|\tilde g(x-y)|$, integration of the preceding inequality over $y$ and (a), applied to $h=\tau_xg$, yield
\[
\left|S_2'(f,\tau_xg)\right|=\left|\langle u_f,\tau_xg\rangle\right|\le\int_{\mathbb{R}^3}\mathrm{d}y\,|u_f(y)|\cdot|g(y-x)|=\left(|u_f|*|\tilde g|\right)(x)\qquad(x\in\mathbb{R}^3).
\]
Young's inequality $\|\varphi*\psi\|_2\le\|\varphi\|_2\cdot\|\psi\|_1$ and the bound $\|u_f\|_2\le\|f\|_2$ of (a) now give
\[
\int_{\mathbb{R}^3}\mathrm{d}x\,S_2'(f,\tau_xg)^2\le\left\||u_f|*|\tilde g|\right\|_2^2\le\|u_f\|_2^2\cdot\|g\|_1^2\le\|f\|_2^2\cdot\|g\|_1^2 .
\]
\end{proof}

\begin{lemma}[Exponentials; determination of a law]\label{lem:exponentials-determination-law}
Let $\mu'$ be a probability law on $(\Omega,\mathcal F)$ with expectation $E'$, let $V$ be a linear subspace of $\mathcal S$, and put $\mathcal F_V=\sigma\left(X(f):f\in V\right)$.
\begin{enumerate}
\item[(a)] The complex linear span of the functions $\exp\left(iX(f)\right)$, $f\in V$, is dense in $L^2(\Omega,\mathcal F_V,\mu')$.
\item[(b)] Let $V_0\subset\mathcal S$ be a set, not necessarily a linear subspace, that is sequentially dense in $\mathcal S$: every $f\in\mathcal S$ is the limit in $\mathcal S$ of a sequence of elements of $V_0$. Then $\sigma\left(X(f):f\in V_0\right)=\mathcal F$, and the closure in $L^2(\mu')$ of the complex linear span of the functions $\exp\left(iX(f)\right)$, $f\in V_0$, is $L^2(\Omega,\mathcal F,\mu')$. Moreover, a probability law on $(\Omega,\mathcal F)$ is determined by the values of its characteristic functional on $V_0$: if $\mu''$ is a probability law on $(\Omega,\mathcal F)$ with expectation $E''$ such that
\[
E''\exp\left(iX(f)\right)=E'\exp\left(iX(f)\right)\qquad\text{for every }f\in V_0,
\]
then $\mu''=\mu'$.
\end{enumerate}
\end{lemma}

\begin{proof}
(a) Let $H\in L^2(\Omega,\mathcal F_V,\mu')$ satisfy
\[
E'\left[H\cdot\exp\left(iX(f)\right)\right]=0\qquad\text{for every }f\in V.
\]
With the scalar product $\langle G_1,G_2\rangle=E'\left[\overline{G_1}\cdot G_2\right]$ of $L^2(\mu')$, this condition states that $\overline{H}$ is orthogonal to all functions $\exp\left(iX(f)\right)$, $f\in V$; since $V=-V$, it equally states that $H$ is orthogonal to all of them. Hence (a) is proved once it is shown that $H=0$.

Fix an integer $k\ge1$ and $f_1,\dots,f_k\in V$. Let $Y=\left(X(f_1),\dots,X(f_k)\right)\colon\Omega\to\mathbb{R}^k$, let $\mathcal F_k=\sigma(Y)$, and let $\mu'_Y=\mu'\circ Y^{-1}$ be the law of $Y$ on $\mathbb{R}^k$. Every $\sigma(Y)$-measurable function factors through $Y$ (the factorisation lemma of Doob and Dynkin); hence there is a Borel function $h\colon\mathbb{R}^k\to\mathbb{C}$ with
\[
E'[H\mid\mathcal F_k]=h(Y)\qquad\mu'\text{-almost surely},
\]
and $h$ is $\mu'_Y$-integrable, because $E'[H\mid\mathcal F_k]$ is $\mu'$-integrable. Let $h\,\mu'_Y$ be the finite complex Borel measure on $\mathbb{R}^k$ with density $h$ with respect to $\mu'_Y$. For $t=(t_1,\dots,t_k)\in\mathbb{R}^k$ put $g=\sum_{l=1}^kt_l\cdot f_l$. Then $g\in V$, as $V$ is a linear subspace, and $X(g)=\sum_{l=1}^kt_l\cdot X(f_l)$, as every $\omega\in\Omega$ is linear on $\mathcal S$; in particular $\exp\left(iX(g)\right)$ is bounded and $\mathcal F_k$-measurable. The Fourier transform of $h\,\mu'_Y$ at $t$ is therefore
\[
\int_{\mathbb{R}^k}\mathrm{d}\mu'_Y(y)\,h(y)\cdot\exp\left(i\sum_{l=1}^kt_l\cdot y_l\right)=E'\left[E'[H\mid\mathcal F_k]\cdot\exp\left(iX(g)\right)\right]=E'\left[H\cdot\exp\left(iX(g)\right)\right]=0.
\]
By the uniqueness of the Fourier transform of a finite complex Borel measure on $\mathbb{R}^k$, the measure $h\,\mu'_Y$ vanishes. Thus $E'[H\mid\mathcal F_k]=0$ almost surely; equivalently, $E'\left[H\cdot1_A\right]=0$ for every $A\in\mathcal F_k$.

Let $\mathcal F_{\cup}$ be the union of the $\sigma$-algebras $\sigma\left(X(f_1),\dots,X(f_k)\right)$ over all $k\ge1$ and all $f_1,\dots,f_k\in V$. Any two of these $\sigma$-algebras are contained in a third one, namely the one formed with the family obtained by joining the two families; hence $\mathcal F_{\cup}$ is an algebra of sets, and the $\sigma$-algebra generated by $\mathcal F_{\cup}$ is $\mathcal F_V$. The sets $A\in\mathcal F_V$ with $E'\left[H\cdot1_A\right]=0$ form a Dynkin system (by linearity, and by dominated convergence, $H$ being integrable), and by the preceding paragraph this system contains $\mathcal F_{\cup}$. By Dynkin's $\pi$-$\lambda$ theorem, $E'\left[H\cdot1_A\right]=0$ for every $A\in\mathcal F_V$. Since $H$ is $\mathcal F_V$-measurable, it follows that $H=0$.

(b) Let $f\in\mathcal S$, and let $(f_n)_{n\ge1}$ be a sequence in $V_0$ that converges to $f$ in $\mathcal S$. Each $\omega\in\Omega$ is a continuous linear form on $\mathcal S$; so
\[
X(f_n)(\omega)=\omega(f_n)\to\omega(f)=X(f)(\omega)\qquad\text{for every }\omega\in\Omega .
\]
As a pointwise limit of $\sigma\left(X(g):g\in V_0\right)$-measurable functions, $X(f)$ is measurable for this $\sigma$-algebra. Hence $\mathcal F\subset\sigma\left(X(g):g\in V_0\right)$, and the reverse inclusion holds because $V_0\subset\mathcal S$. Furthermore $\exp\left(iX(f_n)\right)\to\exp\left(iX(f)\right)$ pointwise, and all these functions have modulus $1$; by dominated convergence, $\exp\left(iX(f_n)\right)\to\exp\left(iX(f)\right)$ in $L^2(\mu')$. Consequently the closure in $L^2(\mu')$ of the complex linear span of the functions $\exp\left(iX(g)\right)$, $g\in V_0$, contains $\exp\left(iX(f)\right)$ for every $f\in\mathcal S$. By (a), applied with $V=\mathcal S$, for which $\mathcal F_V=\mathcal F$, this closure is $L^2(\Omega,\mathcal F,\mu')$.

Let $\mu''$ and $E''$ be as in the statement, let $f\in\mathcal S$, and let $(f_n)$ be as above. Dominated convergence under $\mu'$ and under $\mu''$ gives
\[
E''\exp\left(iX(f)\right)=\lim_{n\to\infty}E''\exp\left(iX(f_n)\right)=\lim_{n\to\infty}E'\exp\left(iX(f_n)\right)=E'\exp\left(iX(f)\right).
\]
Now let $k\ge1$ and $f_1,\dots,f_k\in\mathcal S$. For $t\in\mathbb{R}^k$ one has $\sum_{l=1}^kt_l\cdot X(f_l)=X\left(\sum_{l=1}^kt_l\cdot f_l\right)$, so the laws of $\left(X(f_1),\dots,X(f_k)\right)$ under $\mu'$ and under $\mu''$ have the same characteristic function; by the uniqueness of the Fourier transform they coincide. Therefore $\mu'$ and $\mu''$ agree on every cylinder set $\left\{\left(X(f_1),\dots,X(f_k)\right)\in B\right\}$ with $B\subset\mathbb{R}^k$ a Borel set. The cylinder sets form an algebra, and the $\sigma$-algebra generated by it is $\mathcal F$; by Dynkin's $\pi$-$\lambda$ theorem, $\mu''=\mu'$.
\end{proof}

\begin{lemma}[Averages over translations]\label{lem:averages-translations-2}
The map $\mathbb{R}^3\times\Omega\to\Omega$, $(x,\omega)\mapsto T_x\omega$, is jointly measurable, that is, measurable from $\mathbb{R}^3\times\Omega$, equipped with the product of the Borel $\sigma$-algebra of $\mathbb{R}^3$ and $\mathcal F$, to $(\Omega,\mathcal F)$.

Let $\mu'$ be a translation-invariant probability law on $(\Omega,\mathcal F)$, that is, $\mu'\circ T_x^{-1}=\mu'$ for every $x\in\mathbb{R}^3$, and let $E'$ be its expectation. Put
\[
I'=\left\{A\in\mathcal F:\ \mu'\left(T_x^{-1}A\triangle A\right)=0\ \text{for every }x\in\mathbb{R}^3\right\},
\]
\[
H'_{\mathrm{inv}}=\left\{G\in L^2(\mu'):\ \mathbb{U}_xG=G\ \text{for every }x\in\mathbb{R}^3\right\},
\]
and, for $L>0$ and $G\in L^2(\mu')$,
\[
(A_LG)(\omega)=L^{-3}\int_{[0,L]^3}\mathrm{d}x\,G(T_{-x}\omega)\qquad(\omega\in\Omega).
\]
Then the following hold.
\begin{enumerate}
\item[(a)] For $G\in L^2(\mu')$ the function $A_LG$ is defined $\mu'$-almost everywhere and depends only on the class of $G$ in $L^2(\mu')$, and $\|A_LG\|_2\le\|G\|_2$.
\item[(b)] $H'_{\mathrm{inv}}=L^2(\Omega,I',\mu')$; the orthogonal projection of $L^2(\mu')$ onto $H'_{\mathrm{inv}}$ is the conditional expectation $E'[\,\cdot\mid I']$; and $E'[\mathbb{U}_xG\mid I']=E'[G\mid I']$ for all $G\in L^2(\mu')$ and $x\in\mathbb{R}^3$.
\item[(c)] For $G\in L^2(\mu')$, $A_LG\to E'[G\mid I']$ in $L^2(\mu')$ as $L\to\infty$.
\item[(d)] Let $j\in\{1,2,3\}$, let $e_j$ be the $j$-th unit vector of $\mathbb{R}^3$, and put
\[
I'_j=\left\{A\in\mathcal F:\ \mu'\left(T_{te_j}^{-1}A\triangle A\right)=0\ \text{for every }t\in\mathbb{R}\right\},
\]
\[
H'_j=\left\{G\in L^2(\mu'):\ \mathbb{U}_{te_j}G=G\ \text{for every }t\in\mathbb{R}\right\},\qquad P'_j=E'[\,\cdot\mid I'_j].
\]
Then (a)--(c) hold for the one-parameter group $\left(\mathbb{U}_{te_j}\right)_{t\in\mathbb{R}}$ in place of $\left(\mathbb{U}_x\right)_{x\in\mathbb{R}^3}$. That is, for $G\in L^2(\mu')$ and $T>0$ the average $\frac{1}{T}\int_0^T\mathrm{d}t\,\mathbb{U}_{te_j}G$, whose value at $\omega$ is $\frac{1}{T}\int_0^T\mathrm{d}t\,G(T_{-te_j}\omega)$, is defined $\mu'$-almost everywhere, depends only on the class of $G$, and has norm at most $\|G\|_2$ in $L^2(\mu')$; one has $H'_j=L^2(\Omega,I'_j,\mu')$, the orthogonal projection of $L^2(\mu')$ onto $H'_j$ is $P'_j$, and $P'_j\mathbb{U}_{te_j}G=P'_jG$ for $t\in\mathbb{R}$; and
\[
\frac{1}{T}\int_0^T\mathrm{d}t\,\mathbb{U}_{te_j}G\to P'_jG\qquad\text{in }L^2(\mu')\text{ as }T\to\infty .
\]
Moreover $I'\subset I'_j$.
\end{enumerate}
\end{lemma}

\begin{proof}
\emph{Joint measurability.} The $\sigma$-algebra $\mathcal F$ is generated by the functions $X(f)$, $f\in\mathcal S$; so it suffices to prove that $(x,\omega)\mapsto X(f)(T_x\omega)$ is jointly measurable for each $f\in\mathcal S$. Since $\tau_x^{-1}=\tau_{-x}$, one has $X(f)(T_x\omega)=\omega(\tau_{-x}f)$. For fixed $\omega$ this is a continuous function of $x$, because $x\mapsto\tau_{-x}f$ is continuous from $\mathbb{R}^3$ to $\mathcal S$ and $\omega$ is continuous on $\mathcal S$; for fixed $x$ it is the measurable function $X(\tau_{-x}f)$ of $\omega$. It is therefore the pointwise limit, as $n\to\infty$, of the jointly measurable functions
\[
(x,\omega)\mapsto\sum_{q\in2^{-n}\mathbb{Z}^3}1_{q+[0,2^{-n})^3}(x)\cdot\omega(\tau_{-q}f)
\]
(for given $x$ only the term with $x\in q+[0,2^{-n})^3$ is non-zero, and this $q$ tends to $x$ as $n\to\infty$), hence jointly measurable.

\emph{Unitarity.} Let $\mu'$ be translation invariant, and let $x,y\in\mathbb{R}^3$. By the invariance of $\mu'$, the operator $\mathbb{U}_x\colon G\mapsto G\circ T_x^{-1}=G\circ T_{-x}$ maps functions that vanish $\mu'$-almost everywhere to functions that vanish $\mu'$-almost everywhere, and it preserves the norm of $L^2(\mu')$. Moreover $\mathbb{U}_x\mathbb{U}_y=\mathbb{U}_{x+y}$, and $\mathbb{U}_0$ is the identity. Hence $\mathbb{U}_x$ is a unitary operator on $L^2(\mu')$, with inverse and adjoint $\mathbb{U}_{-x}$; it commutes with complex conjugation.

(a) Let $G$ be a measurable representative of an element of $L^2(\mu')$. By the joint measurability, $(x,\omega)\mapsto G(T_{-x}\omega)$ is jointly measurable, and Tonelli's theorem together with the invariance of $\mu'$ gives
\[
\int\mathrm{d}\mu'(\omega)\int_{[0,L]^3}\mathrm{d}x\,\left|G(T_{-x}\omega)\right|^2=\int_{[0,L]^3}\mathrm{d}x\,E'\left[\left|G\circ T_{-x}\right|^2\right]=L^3\cdot\|G\|_2^2<\infty.
\]
Hence, for $\mu'$-almost every $\omega$, the function $x\mapsto G(T_{-x}\omega)$ is square integrable, and so integrable, on $[0,L]^3$; for such $\omega$ the integral that defines $A_LG(\omega)$ exists, and by Fubini's theorem $A_LG$ is measurable. If $\tilde G$ is another representative of the same element, the identity, applied to $G-\tilde G$, shows that for $\mu'$-almost every $\omega$ one has $G(T_{-x}\omega)=\tilde G(T_{-x}\omega)$ for almost every $x\in[0,L]^3$; hence $A_LG=A_L\tilde G$ $\mu'$-almost everywhere. Finally, by Jensen's inequality,
\[
\left|A_LG(\omega)\right|^2\le L^{-3}\int_{[0,L]^3}\mathrm{d}x\,\left|G(T_{-x}\omega)\right|^2,
\]
and integration with respect to $\mu'$, together with the identity above, gives $\|A_LG\|_2^2\le\|G\|_2^2$.

(b) The class $I'$ is a $\sigma$-algebra that contains every $\mu'$-null set of $\mathcal F$: taking preimages under $T_x$ commutes with complements and countable unions, the symmetric difference of two countable unions is contained in the union of the symmetric differences of the corresponding sets, and $\mu'\left(T_x^{-1}N\triangle N\right)\le\mu'\left(T_x^{-1}N\right)+\mu'(N)=2\mu'(N)$ for $N\in\mathcal F$.

Let $A\in I'$ and $x\in\mathbb{R}^3$. Then $\mathbb{U}_x1_A=1_A\circ T_{-x}=1_{T_{-x}^{-1}A}$, and $\mu'\left(T_{-x}^{-1}A\triangle A\right)=0$; so $\mathbb{U}_x1_A=1_A$ in $L^2(\mu')$. By linearity every $I'$-measurable simple function lies in $H'_{\mathrm{inv}}$. These functions are dense in $L^2(\Omega,I',\mu')$, and $H'_{\mathrm{inv}}$ is closed because every $\mathbb{U}_x$ is continuous; therefore $L^2(\Omega,I',\mu')\subset H'_{\mathrm{inv}}$.

Conversely, let $G\in H'_{\mathrm{inv}}$. Its real and imaginary parts belong to $H'_{\mathrm{inv}}$, since every $\mathbb{U}_x$ commutes with complex conjugation; so it suffices to treat a real $G$, of which a real measurable representative, again denoted $G$, is fixed. For $c\in\mathbb{R}$ and $x\in\mathbb{R}^3$,
\[
T_x^{-1}\{G>c\}=\{G\circ T_x>c\},
\]
and $G\circ T_x=\mathbb{U}_{-x}G=G$ $\mu'$-almost everywhere. Hence $T_x^{-1}\{G>c\}$ and $\{G>c\}$ differ by a $\mu'$-null set, that is, $\{G>c\}\in I'$. Thus $G$ is $I'$-measurable, and $H'_{\mathrm{inv}}=L^2(\Omega,I',\mu')$. The orthogonal projection of $L^2(\Omega,\mathcal F,\mu')$ onto its closed subspace $L^2(\Omega,I',\mu')$ is the conditional expectation $E'[\,\cdot\mid I']$. Finally, for $H\in H'_{\mathrm{inv}}$, $G\in L^2(\mu')$ and $x\in\mathbb{R}^3$,
\[
\langle H,\mathbb{U}_xG\rangle=\langle\mathbb{U}_{-x}H,G\rangle=\langle H,G\rangle;
\]
so $\mathbb{U}_xG-G$ is orthogonal to $H'_{\mathrm{inv}}$, and $E'[\mathbb{U}_xG\mid I']=E'[G\mid I']$.

(c) For $G,H\in L^2(\mu')$ and $y\in\mathbb{R}^3$, the unitarity of $\mathbb{U}_y$ gives
\[
\langle G,\mathbb{U}_yH-H\rangle=\langle\mathbb{U}_{-y}G-G,H\rangle .
\]
Hence $G$ is orthogonal to all functions $\mathbb{U}_yH-H$, $y\in\mathbb{R}^3$, $H\in L^2(\mu')$, if and only if $\mathbb{U}_{-y}G=G$ for every $y\in\mathbb{R}^3$, that is, if and only if $G\in H'_{\mathrm{inv}}$. Let $R'$ be the closure in $L^2(\mu')$ of the linear span of these functions; then
\[
L^2(\mu')=H'_{\mathrm{inv}}\oplus R'
\]
is an orthogonal decomposition.

Let $G\in H'_{\mathrm{inv}}$. For each $x\in\mathbb{R}^3$ one has $G\circ T_{-x}=\mathbb{U}_xG=G$ $\mu'$-almost everywhere. The set $\left\{(x,\omega):G(T_{-x}\omega)\ne G(\omega)\right\}$ is jointly measurable, and its sections at fixed $x$ are $\mu'$-null sets; by Tonelli's theorem, for $\mu'$-almost every $\omega$ its section at $\omega$ is a Lebesgue null set, that is, $G(T_{-x}\omega)=G(\omega)$ for almost every $x\in\mathbb{R}^3$. For such $\omega$ one has $A_LG(\omega)=G(\omega)$. Hence $A_LG=G$.

Let now $G=\mathbb{U}_yH-H$ with $y\in\mathbb{R}^3$ and $H\in L^2(\mu')$. Since $\left(\mathbb{U}_yH\right)(T_{-x}\omega)=H(T_{-y}T_{-x}\omega)=H(T_{-(x+y)}\omega)$, the substitution $x\mapsto x-y$ gives, $\mu'$-almost everywhere,
\[
A_LG=L^{-3}\left(\int_{y+[0,L]^3}\mathrm{d}x\,H\circ T_{-x}-\int_{[0,L]^3}\mathrm{d}x\,H\circ T_{-x}\right).
\]
The contributions of $\left(y+[0,L]^3\right)\cap[0,L]^3$ to the two integrals cancel, so that only the symmetric difference of the two cubes remains; moreover $\|H\circ T_{-x}\|_2=\|H\|_2$ for every $x$. By Minkowski's integral inequality,
\[
\|A_LG\|_2\le L^{-3}\cdot\left|\left(y+[0,L]^3\right)\triangle[0,L]^3\right|\cdot\|H\|_2\le\frac{2\left(|y_1|+|y_2|+|y_3|\right)}{L}\cdot\|H\|_2 ,
\]
where $|B|$ denotes the Lebesgue measure of a Borel set $B\subset\mathbb{R}^3$, and likewise of a Borel set $B\subset\mathbb{R}$; the second inequality holds because the symmetric difference has measure $2\left(L^3-\prod_{i=1}^3\max\{L-|y_i|,0\}\right)$, which is at most $2L^2\left(|y_1|+|y_2|+|y_3|\right)$. So $A_LG\to0$ in $L^2(\mu')$ as $L\to\infty$.

By linearity, $A_LG\to0$ for every $G$ in the linear span of the functions $\mathbb{U}_yH-H$; since the linear operators $A_L$ have norm at most $1$ by (a), the same holds for every $G\in R'$. An arbitrary $G\in L^2(\mu')$ is the sum $G=G_{\mathrm{inv}}+G_{\perp}$ of some $G_{\mathrm{inv}}\in H'_{\mathrm{inv}}$ and some $G_{\perp}\in R'$; then $A_LG=G_{\mathrm{inv}}+A_LG_{\perp}\to G_{\mathrm{inv}}$ in $L^2(\mu')$, and $G_{\mathrm{inv}}=E'[G\mid I']$ by (b).

(d) Let $A\in I'$. The condition that defines $I'$, taken at the points $x=te_j$, $t\in\mathbb{R}$, states that $\mu'\left(T_{te_j}^{-1}A\triangle A\right)=0$ for every $t\in\mathbb{R}$; so $A\in I'_j$, and $I'\subset I'_j$. The other assertions of (d) are proved by the arguments of (a)--(c) in one variable, with $\mathbb{R}^3$, $x$, $[0,L]^3$ and $L^{-3}$ replaced by $\mathbb{R}$, $te_j$, $[0,T]$ and $\frac{1}{T}$. In detail: the map $(t,\omega)\mapsto T_{te_j}\omega$ is jointly measurable, as the composition of $(t,\omega)\mapsto(te_j,\omega)$ with $(x,\omega)\mapsto T_x\omega$; Tonelli's theorem and the invariance of $\mu'$ give
\[
\int\mathrm{d}\mu'(\omega)\int_0^T\mathrm{d}t\,\left|G(T_{-te_j}\omega)\right|^2=T\cdot\|G\|_2^2 ,
\]
from which the analogue of (a) follows as in (a), by Jensen's inequality. The proof of (b) uses only the invariance of $\mu'$ under the maps concerned and the unitarity of the corresponding operators; it gives $H'_j=L^2(\Omega,I'_j,\mu')$, the identification of the orthogonal projection onto $H'_j$ with $P'_j$, and $P'_j\mathbb{U}_{te_j}G=P'_jG$. The first step of the proof of (c) gives the orthogonal decomposition $L^2(\mu')=H'_j\oplus R'_j$, where $R'_j$ is the closure of the linear span of the functions $\mathbb{U}_{se_j}H-H$, $s\in\mathbb{R}$, $H\in L^2(\mu')$. For $G\in H'_j$ the average equals $G$, by Tonelli's theorem as in (c). For $G=\mathbb{U}_{se_j}H-H$, Minkowski's integral inequality gives
\[
\left\|\frac{1}{T}\int_0^T\mathrm{d}t\,\mathbb{U}_{te_j}G\right\|_2\le\frac{1}{T}\cdot\left|\left(s+[0,T]\right)\triangle[0,T]\right|\cdot\|H\|_2\le\frac{2|s|}{T}\cdot\|H\|_2\to0\qquad(T\to\infty).
\]
Since the averages are linear operators of norm at most $1$, the convergence to $P'_jG$ for every $G\in L^2(\mu')$ follows as at the end of the proof of (c).
\end{proof}

Whenever a translation-invariant probability law on $(\Omega,\mathcal F)$ is denoted by $\mu$, with expectation $E$, the objects of Lemma~\ref{lem:averages-translations-2} formed with $\mu'=\mu$ are written without primes:
\[
I,\qquad I_j,\qquad H_{\mathrm{inv}},\qquad H_j,\qquad P=E[\,\cdot\mid I],\qquad P_j=E[\,\cdot\mid I_j].
\]
In particular this notation applies to the law $\mu$ of Section~\ref{sec:hypotheses-law}, which is translation invariant by~\eqref{eq:fourth-cumulants-H2}.

\begin{definition}[Ergodicity; ergodic decompositions]\label{def:ergodicity-ergodic-decompositions}
A translation-invariant law $\mu'$ on $(\Omega,\mathcal F)$ is called \emph{ergodic} if the space $H'_{\mathrm{inv}}$ of Lemma~\ref{lem:averages-translations-2}, formed with the law $\mu'$, consists of the constant functions.

Let $\mu$ be a law on $(\Omega,\mathcal F)$. An \emph{ergodic decomposition} of $\mu$ is a probability space $(\Xi,B,\nu)$ together with a family $(\mu_{\xi})_{\xi\in\Xi}$ of laws on $(\Omega,\mathcal F)$ such that
\begin{equation}\tag{ED1}\label{eq:fourth-cumulants-ED1}
\xi\mapsto\mu_{\xi}(A)\ \text{is $B$-measurable, and}\ \mu(A)=\int\mathrm{d}\nu(\xi)\,\mu_{\xi}(A),\quad\text{for every }A\in\mathcal F;
\end{equation}
\begin{equation}\tag{ED2}\label{eq:fourth-cumulants-ED2}
\text{for $\nu$-almost every $\xi\in\Xi$ the law $\mu_{\xi}$ is translation invariant and ergodic.}
\end{equation}
For $\xi\in\Xi$ the expectation with respect to $\mu_{\xi}$ is written $E_{\xi}$.
\end{definition}

\begin{remark}
Let $\mu$ be a translation-invariant law on $(\Omega,\mathcal F)$. Theorem~\ref{thm:ergodic-decomposition}(a), applied to the law $m=\mu$ and to $J=\{1,2,3\}$, furnishes an ergodic decomposition of $\mu$ in the sense of Definition~\ref{def:ergodicity-ergodic-decompositions}: $\Xi$ is the set of all laws on $(\Omega,\mathcal F)$; $B$ is the $\sigma$-algebra on $\Xi$ of Section~\ref{sec:test-forms-distributions}, which is generated by the maps $\rho\mapsto\rho(A)$, $A\in\mathcal F$; $\nu$ is the measure provided by that theorem; and $\mu_{\xi}=\xi$ for $\xi\in\Xi$. Property \eqref{eq:fourth-cumulants-ED1} follows from the definition of $B$ and from the identity $\mu=\int\mathrm{d}\nu(\rho)\,\rho$ of that theorem. Property \eqref{eq:fourth-cumulants-ED2} is the equality $\nu\left(\mathcal{M}^{123}_{\mathrm{erg}}\right)=1$ of that theorem: by Section~\ref{sec:invariant-sets-invariant}, a law $\rho$ lies in $\mathcal{M}^{123}_{\mathrm{erg}}$ exactly when $\rho$ is translation invariant and the constants are the only elements of $L^2(\rho)$ left fixed by every translation of $\mathbb{R}^3$, and this is ergodicity in the sense of Definition~\ref{def:ergodicity-ergodic-decompositions}. Proposition~\ref{prop:existence-conditional-expectations}(b) below shows that the distribution under $\nu$ of a Borel function of finitely many of the numbers $E_{\xi}[G]$, $G\in L^2(\mu)$, does not depend on the choice of the ergodic decomposition of $\mu$.
\end{remark}

\begin{proposition}[Existence; conditional expectations]\label{prop:existence-conditional-expectations}
Let $\mu$ be a translation-invariant law on $(\Omega,\mathcal F)$ with expectation $E$, and let $I$ be the $\sigma$-algebra $I'$ of Lemma~\ref{lem:averages-translations-2} formed with $\mu'=\mu$.
\begin{enumerate}
\item[(a)] The law $\mu$ has an ergodic decomposition.
\item[(b)] Let $(\Xi,B,\nu)$, $(\mu_{\xi})_{\xi\in\Xi}$ be an ergodic decomposition of $\mu$. If $G\in L^2(\mu)$, then, for $\nu$-almost every $\xi\in\Xi$, the function $G$ belongs to $L^2(\mu_{\xi})$ and the conditional expectation $E[G\mid I]$ coincides $\mu_{\xi}$-almost surely with the constant $E_{\xi}[G]$. If $G_1,\dots,G_k\in L^2(\mu)$, it follows that the $\mathbb{C}^k$-valued random vectors $\left(E_{\xi}[G_1],\dots,E_{\xi}[G_k]\right)$ on $(\Xi,B,\nu)$ and $\left(E[G_1\mid I],\dots,E[G_k\mid I]\right)$ on $(\Omega,\mathcal F,\mu)$ have the same distribution. For $G,H\in L^2(\mu)$ one obtains in particular
\[
\int\mathrm{d}\nu(\xi)\,E_{\xi}[G]\cdot E_{\xi}[H]=E\left[E[G\mid I]\cdot E[H\mid I]\right].
\]
\end{enumerate}
\end{proposition}

\begin{proof}[Proof of (b)]
Both statements of \eqref{eq:fourth-cumulants-ED1} carry over, by linearity and monotone convergence, from indicator functions to arbitrary measurable functions $\Phi\colon\Omega\to[0,\infty]$: the map $\xi\mapsto\int\mathrm{d}\mu_{\xi}\,\Phi$ is $B$-measurable, and
\[
\int\mathrm{d}\mu\,\Phi=\int\mathrm{d}\nu(\xi)\int\mathrm{d}\mu_{\xi}\,\Phi .
\]
In particular, a $\mu$-null set is $\mu_{\xi}$-null for $\nu$-almost every $\xi$, and the choice $\Phi=|G|^2$ gives
\[
\int\mathrm{d}\nu(\xi)\,\|G\|^2_{L^2(\mu_{\xi})}=\|G\|^2_{L^2(\mu)}<\infty ,
\]
so that $G\in L^2(\mu_{\xi})$ for $\nu$-almost every $\xi$.
For $t\in\mathbb{R}$ write $t^{+}=\max\{t,0\}$ and $t^{-}=\max\{-t,0\}$. Let $\Phi\colon\Omega\to\mathbb{C}$ be $\mathcal F$-measurable, and let
\[
\Xi_{\Phi}=\left\{\xi\in\Xi:\ \int\mathrm{d}\mu_{\xi}\,|\Phi|<\infty\right\}.
\]
By the extended form of \eqref{eq:fourth-cumulants-ED1}, the map $\xi\mapsto\int\mathrm{d}\mu_{\xi}\,|\Phi|$ is $B$-measurable with values in $[0,\infty]$; the set $\Xi_{\Phi}$, the preimage of $[0,\infty)$ under this map, therefore belongs to $B$. The functions $(\operatorname{Re}\Phi)^{+}$, $(\operatorname{Re}\Phi)^{-}$, $(\operatorname{Im}\Phi)^{+}$, $(\operatorname{Im}\Phi)^{-}$ are $\mathcal F$-measurable, non-negative and bounded by $|\Phi|$. By the extended form of \eqref{eq:fourth-cumulants-ED1} again, their $\mu_{\xi}$-integrals are $B$-measurable functions of $\xi$ with values in $[0,\infty]$; these values are finite for $\xi\in\Xi_{\Phi}$, and then
\[
\int\mathrm{d}\mu_{\xi}\,\Phi=\int\mathrm{d}\mu_{\xi}\,(\operatorname{Re}\Phi)^{+}-\int\mathrm{d}\mu_{\xi}\,(\operatorname{Re}\Phi)^{-}+i\int\mathrm{d}\mu_{\xi}\,(\operatorname{Im}\Phi)^{+}-i\int\mathrm{d}\mu_{\xi}\,(\operatorname{Im}\Phi)^{-}.
\]
Consequently $\xi\mapsto\int\mathrm{d}\mu_{\xi}\,\Phi$ is a $B$-measurable complex function on $\Xi_{\Phi}$. If moreover $\int\mathrm{d}\mu\,|\Phi|<\infty$, then $\int\mathrm{d}\nu(\xi)\int\mathrm{d}\mu_{\xi}\,|\Phi|=\int\mathrm{d}\mu\,|\Phi|<\infty$, so that $\nu\left(\Xi_{\Phi}\right)=1$; the four functions above have finite $\mu$-integrals, and the extended form of \eqref{eq:fourth-cumulants-ED1}, applied to each of them, the last display and the linearity of the integrals give
\[
\int\mathrm{d}\mu\,\Phi=\int\mathrm{d}\nu(\xi)\int\mathrm{d}\mu_{\xi}\,\Phi ,
\]
the integrand on the right being defined on $\Xi_{\Phi}$ and of modulus at most $\int\mathrm{d}\mu_{\xi}\,|\Phi|$. Let $\Phi'\colon\Omega\to\mathbb{C}$ be a second $\mathcal F$-measurable function with $\Phi'=\Phi$ $\mu$-almost everywhere. The $\mu$-null set $\{\Phi'\ne\Phi\}\in\mathcal F$ is $\mu_{\xi}$-null for $\nu$-almost every $\xi$; for each such $\xi$ one has $\xi\in\Xi_{\Phi'}$ exactly when $\xi\in\Xi_{\Phi}$, and then $\int\mathrm{d}\mu_{\xi}\,\Phi'=\int\mathrm{d}\mu_{\xi}\,\Phi$. In particular, if $\Phi$ and $\Phi'$ represent the same element of $L^1(\mu)$, the numbers $\int\mathrm{d}\mu_{\xi}\,\Phi$ and $\int\mathrm{d}\mu_{\xi}\,\Phi'$ are defined and equal for $\nu$-almost every $\xi$. Every representative $\Phi$ of an element of $L^2(\mu)$ represents an element of $L^1(\mu)$, since $\int\mathrm{d}\mu\,|\Phi|\le\|\Phi\|_{L^2(\mu)}$ by the Cauchy--Schwarz inequality and $\mu(\Omega)=1$. For a representative $\Phi$ of an element of $L^1(\mu)$, the number $E_{\xi}[\Phi]$ in (b) is $\int\mathrm{d}\mu_{\xi}\,\Phi$ if $\xi\in\Xi_{\Phi}$, and it is set equal to $0$ if $\xi$ lies in the $\nu$-null set $\Xi\setminus\Xi_{\Phi}$; with this convention $\xi\mapsto E_{\xi}[\Phi]$ is a $B$-measurable complex function on $\Xi$, as $\Xi_{\Phi}\in B$. For $G_1,\dots,G_k\in L^2(\mu)$ the map $\xi\mapsto\left(E_{\xi}[G_1],\dots,E_{\xi}[G_k]\right)$ is therefore $B$-measurable, its components being so, and it is a $\mathbb{C}^k$-valued random vector on $(\Xi,B,\nu)$. Since $\nu\left(\Xi_{G_i}\right)=1$ for each $i$, and two representatives of $G_i$ have equal $\mu_{\xi}$-integrals for $\nu$-almost every $\xi$, replacing the representatives of $G_1,\dots,G_k$ by others, or, for each $i$, the value $0$ on the set $\Xi\setminus\Xi_{G_i}$ by a fixed number $c_i\in\mathbb{C}$, yields again a $B$-measurable map $\Xi\to\mathbb{C}^k$; the set of the $\xi$ at which it differs from this random vector belongs to $B$ and is $\nu$-null, so that the two maps have the same distribution. Let now $\Phi$ be a representative of an element of $L^1(\mu)$ with values in $[0,\infty)$. By the extended form of \eqref{eq:fourth-cumulants-ED1}, the integral $\int\mathrm{d}\mu_{\xi}\,\Phi\in[0,\infty]$ is a $B$-measurable function of $\xi$ on all of $\Xi$; as $|\Phi|=\Phi$, it is finite exactly on $\Xi_{\Phi}$, where it equals $E_{\xi}[\Phi]$, and it takes the value $\infty$ on the $\nu$-null set $\Xi\setminus\Xi_{\Phi}\in B$. Hence, for every Borel set $S\subset[0,\infty]$, the sets
\[
\left\{\xi\in\Xi:\ \int\mathrm{d}\mu_{\xi}\,\Phi\in S\right\}\quad\text{and}\quad\left\{\xi\in\Xi:\ E_{\xi}[\Phi]\in S\right\}
\]
belong to $B$ and differ only by a subset of $\Xi\setminus\Xi_{\Phi}$, so that they have the same $\nu$-measure. Applied to a representative $G$ of an element of $L^2(\mu)$ with values in $[0,\infty)$, this shows that the equality of distributions in (b), for $k=1$, holds also with the $[0,\infty]$-valued function $\xi\mapsto\int\mathrm{d}\mu_{\xi}\,G$ in place of $\xi\mapsto E_{\xi}[G]$:
\[
\nu\left(\left\{\xi\in\Xi:\ \int\mathrm{d}\mu_{\xi}\,G\in S\right\}\right)=\mu\left(\left\{E[G\mid I]\in S\right\}\right)\qquad\text{for every Borel set }S\subset[0,\infty].
\]

Let $\bar{\bar{G}}$ be a representative of $E[G\mid I]$. For $L>0$ the function $A_LG$ of Lemma~\ref{lem:averages-translations-2}, formed with the fixed function $G$, is defined $\mu$-almost everywhere and belongs to $L^2(\mu)$, by Lemma~\ref{lem:averages-translations-2}(a). The identity above with $\Phi=\left|A_LG-\bar{\bar{G}}\right|^2$, together with Lemma~\ref{lem:averages-translations-2}(c), gives
\[
\int\mathrm{d}\nu(\xi)\,\left\|A_LG-\bar{\bar{G}}\right\|^2_{L^2(\mu_{\xi})}=\left\|A_LG-\bar{\bar{G}}\right\|^2_{L^2(\mu)}\longrightarrow0\qquad(L\to\infty).
\]
Choose integers $L_1<L_2<\cdots$ for which the numbers $\left\|A_{L_k}G-\bar{\bar{G}}\right\|^2_{L^2(\mu)}$, $k\ge1$, have a finite sum. By the same identity, the function $\xi\mapsto\sum_k\left\|A_{L_k}G-\bar{\bar{G}}\right\|^2_{L^2(\mu_{\xi})}$ has a finite $\nu$-integral; it is therefore finite for $\nu$-almost every $\xi$, and for every such $\xi$
\[
A_{L_k}G\longrightarrow\bar{\bar{G}}\quad\text{in }L^2(\mu_{\xi})\qquad(k\to\infty).
\]
On the other hand, by \eqref{eq:fourth-cumulants-ED2} and the first paragraph, for $\nu$-almost every $\xi$ the law $\mu_{\xi}$ is translation invariant and ergodic, and $G\in L^2(\mu_{\xi})$. For such $\xi$, Lemma~\ref{lem:averages-translations-2}(a)--(c), applied to $\mu'=\mu_{\xi}$, shows that the same function $A_LG$ is defined $\mu_{\xi}$-almost everywhere and converges in $L^2(\mu_{\xi})$, as $L\to\infty$, to the orthogonal projection of $G$ onto the space $H'_{\mathrm{inv}}$ formed with $\mu_{\xi}$. This space consists of the constants, so the projection is the constant $E_{\xi}[G]$, and
\[
A_LG\longrightarrow E_{\xi}[G]\quad\text{in }L^2(\mu_{\xi})\qquad(L\to\infty).
\]
The two limits along $(L_k)$ coincide. Hence $\bar{\bar{G}}=E_{\xi}[G]$ holds $\mu_{\xi}$-almost surely, for $\nu$-almost every $\xi$, which is the first assertion.

Now let $G_1,\dots,G_k\in L^2(\mu)$, let $\bar{\bar{G}}_1,\dots,\bar{\bar{G}}_k$ be representatives of $E[G_1\mid I],\dots,E[G_k\mid I]$, and let $\varphi\colon\mathbb{C}^k\to\mathbb{R}$ be a bounded Borel function, $\mathbb{C}^k$ carrying the Borel $\sigma$-algebra of $\mathbb{R}^{2k}$. The extended form of \eqref{eq:fourth-cumulants-ED1}, applied to the positive and the negative part of $\varphi\left(\bar{\bar{G}}_1,\dots,\bar{\bar{G}}_k\right)$, and the first assertion give
\begin{align*}
\int\mathrm{d}\mu\,\varphi\left(\bar{\bar{G}}_1,\dots,\bar{\bar{G}}_k\right)&=\int\mathrm{d}\nu(\xi)\int\mathrm{d}\mu_{\xi}(\omega)\,\varphi\left(\bar{\bar{G}}_1(\omega),\dots,\bar{\bar{G}}_k(\omega)\right)\\
&=\int\mathrm{d}\nu(\xi)\,\varphi\left(E_{\xi}[G_1],\dots,E_{\xi}[G_k]\right),
\end{align*}
and this is the asserted equality of laws. Finally, let $G,H\in L^2(\mu)$. By the Cauchy--Schwarz inequality,
\[
\left|E_{\xi}[G]\cdot E_{\xi}[H]\right|\le\|G\|_{L^2(\mu_{\xi})}\cdot\|H\|_{L^2(\mu_{\xi})},
\]
and the right side is $\nu$-integrable, by the Cauchy--Schwarz inequality once more and the identity $\int\mathrm{d}\nu(\xi)\,\|G\|^2_{L^2(\mu_{\xi})}=\|G\|^2_{L^2(\mu)}$ and its analogue for $H$. Thus the function $(z,w)\mapsto z\cdot w$ on $\mathbb{C}^2$ is integrable with respect to the law of $\left(E_{\xi}[G],E_{\xi}[H]\right)$ under $\nu$, which is the law of $\left(E[G\mid I],E[H\mid I]\right)$ under $\mu$; integrating it against these two equal laws gives the last formula.
\end{proof}

\begin{proof}[Proof of (a)]
\emph{Step 0: the measurable space $(\Omega,\mathcal F)$ is standard Borel.} Let $(h_n)_{n\in\mathbb{N}}$ be the Hermite basis of $L^2(\mathbb{R}^3;\Lambda^2)$: its elements are the two-forms of which one component is a product of one-dimensional Hermite functions of $x_1$, $x_2$ and $x_3$, the two other components being zero, and they are enumerated by increasing degree, the degree being the sum of the orders of the three factors. Define $J\colon\Omega\to\mathbb{R}^{\mathbb{N}}$ by $J(\omega)=\left(\omega(h_n)\right)_{n\in\mathbb{N}}$. By the $N$-representation theorem for $\mathcal S$ and its dual, $J$ is a bijection of $\Omega$ onto
\[
s'=\bigcup_{k,l\in\mathbb{N}}\left\{a\in\mathbb{R}^{\mathbb{N}}:\ |a_n|\le l\cdot(1+n)^k\ \text{for every }n\in\mathbb{N}\right\},
\]
and
\[
\omega(f)=\sum_{n\in\mathbb{N}}a_n\cdot\langle h_n,f\rangle\qquad\left(\omega\in\Omega,\ a=J(\omega),\ f\in\mathcal S\right);
\]
since the enumeration is by increasing degree, polynomial bounds in the multi-indices of the Hermite functions and polynomial bounds in $n$ are equivalent, so that the theorem takes this form. Each of the sets whose union is $s'$ is closed in the Polish space $\mathbb{R}^{\mathbb{N}}$; hence $s'$ is a countable union of closed sets, in particular a Borel set. The components of $J$ are the measurable functions $X(h_n)$, so $J$ is measurable. The inverse map $J^{-1}\colon s'\to\Omega$ is measurable as well: for every $f\in\mathcal S$ the function $X(f)\circ J^{-1}$, that is $a\mapsto\sum_na_n\cdot\langle h_n,f\rangle$, is the pointwise limit on $s'$ of finite linear combinations of the coordinates. Consequently $J(A)=\left(J^{-1}\right)^{-1}(A)$ is a Borel subset of $\mathbb{R}^{\mathbb{N}}$ for every $A\in\mathcal F$, and $(\Omega,\mathcal F)$, being isomorphic to the Borel subset $s'$ of a Polish space, is standard Borel.

\emph{Step 1: a regular conditional distribution.} Let $\mathbb{R}^{\mathbb{N}}$ carry its Borel $\sigma$-algebra, and write $\mu(A\mid I)=E[1_A\mid I]$ for $A\in\mathcal F$. The theorem on the existence of regular conditional distributions for random variables with values in a Polish space, applied to the $\mathbb{R}^{\mathbb{N}}$-valued random variable $J$ on $(\Omega,\mathcal F,\mu)$ and to the sub-$\sigma$-algebra $I$ of $\mathcal F$, furnishes a $\mu$-null set $N_1\in\mathcal F$ and a function $Q$, which assigns a number $Q(\omega,Y)\in[0,1]$ to every $\omega\in\Omega$ and every Borel set $Y\subset\mathbb{R}^{\mathbb{N}}$, with the following properties: for every $\omega\in\Omega\setminus N_1$ the set function $Q(\omega,\cdot)$ is a probability measure on $\mathbb{R}^{\mathbb{N}}$; for every Borel set $Y\subset\mathbb{R}^{\mathbb{N}}$ the function $\omega\mapsto Q(\omega,Y)$ is $I$-measurable, and
\[
Q(\omega,Y)=\mu\left(J^{-1}(Y)\,\middle|\,I\right)(\omega)\qquad\text{for $\mu$-almost every }\omega .
\]
By the invariance of $\mu$, every $\mu$-null set $A\in\mathcal F$ belongs to $I$: for $x\in\mathbb{R}^3$ one has $\mu\left(T_x^{-1}A\triangle A\right)\le\mu\left(T_x^{-1}A\right)+\mu(A)=2\mu(A)=0$. In particular $N_1\in I$. Since $J^{-1}(s')=\Omega$, the displayed identity, applied to the Borel set $Y=s'$, gives $Q(\omega,s')=1$ for $\mu$-almost every $\omega$; the set of the $\omega$ with $Q(\omega,s')\ne1$ is therefore $\mu$-null, and it lies in $I$, the function $Q(\cdot,s')$ being $I$-measurable. Let $N_0$ be the union of $N_1$ and of this set. Then $N_0$ is a $\mu$-null set that belongs to $I$, and for every $\omega\in\Omega\setminus N_0$ the set function $Q(\omega,\cdot)$ is a probability measure on $\mathbb{R}^{\mathbb{N}}$ with $Q(\omega,s')=1$. For $\omega\in\Omega$ and $A\in\mathcal F$ put
\[
\mu_{\omega}(A)=\begin{cases}Q\left(\omega,J(A)\right)&\text{if }\omega\notin N_0,\\ \mu(A)&\text{if }\omega\in N_0;\end{cases}
\]
the set $J(A)$ is Borel by Step~0. Since $J$ is a bijection of $\Omega$ onto $s'$ and $Q(\omega,s')=1$ for $\omega\notin N_0$, every $\mu_{\omega}$ is a law on $(\Omega,\mathcal F)$; since $N_0\in I$, the map $(\omega,A)\mapsto\mu_{\omega}(A)$ is a probability kernel from $(\Omega,I)$ to $(\Omega,\mathcal F)$. Moreover $J^{-1}\left(J(A)\right)=A$, so that, for each $A\in\mathcal F$,
\[
\mu_{\omega}(A)=\mu(A\mid I)(\omega)\qquad\text{for $\mu$-almost every }\omega .
\]
By the monotone class theorem and monotone convergence, it follows that for each $G\in L^1(\mu)$
\[
\int\mathrm{d}\mu_{\omega}\,G=E[G\mid I](\omega)\qquad\text{for $\mu$-almost every }\omega .
\]
Let $(\Xi,B,\nu)=\left(\Omega,I,\mu|_I\right)$, with the family $(\mu_{\omega})_{\omega\in\Omega}$. For $A\in\mathcal F$ the map $\omega\mapsto\mu_{\omega}(A)$ is $I$-measurable, and $\int\mathrm{d}\mu(\omega)\,\mu_{\omega}(A)=E\left[\mu(A\mid I)\right]=\mu(A)$. Thus \eqref{eq:fourth-cumulants-ED1} holds.

\emph{Step 2: the laws $\mu_{\omega}$ are translation invariant.} Let $V_0\subset\mathcal S$ be the $\mathbb{Q}$-linear span of the two-forms $\tau_qh_n$, $q\in\mathbb{Q}^3$, $n\in\mathbb{N}$. The set $V_0$ is countable, it is mapped into itself by $\tau_q$ for every $q\in\mathbb{Q}^3$, and it is sequentially dense in $\mathcal S$: by the $N$-representation theorem the finite real linear combinations of the $h_n$ are dense in the metrizable space $\mathcal S$, and real coefficients are limits of rational ones.

Let $f\in V_0$ and $q\in\mathbb{Q}^3$. Since $\mathbb{U}_q$ is the composition with $T_q^{-1}$ and $\mathbb{U}_qX(f)=X(\tau_qf)$, one has $\exp\left(iX(\tau_qf)\right)=\mathbb{U}_q\exp\left(iX(f)\right)$. Step~1 and Lemma~\ref{lem:averages-translations-2}(b) give, for $\mu$-almost every $\omega$,
\[
\int\mathrm{d}\mu_{\omega}\,\exp\left(iX(\tau_qf)\right)=E\left[\mathbb{U}_q\exp\left(iX(f)\right)\,\middle|\,I\right](\omega)=E\left[\exp\left(iX(f)\right)\,\middle|\,I\right](\omega)=\int\mathrm{d}\mu_{\omega}\,\exp\left(iX(f)\right).
\]
As $V_0\times\mathbb{Q}^3$ is countable, there is one $\mu$-null set outside of which these identities hold for all $f\in V_0$ and all $q\in\mathbb{Q}^3$. Let $\omega$ lie outside this set, and define $\Phi_{\omega}\colon\mathcal S\to\mathbb{C}$ by
\[
\Phi_{\omega}(f)=\int\mathrm{d}\mu_{\omega}\,\exp\left(iX(f)\right).
\]
If $f_m\to f$ in $\mathcal S$, then $X(f_m)\to X(f)$ pointwise on $\Omega$, and dominated convergence gives $\Phi_{\omega}(f_m)\to\Phi_{\omega}(f)$; thus $\Phi_{\omega}$ is sequentially continuous. The map $\mathbb{R}^3\times\mathcal S\to\mathcal S$, $(x,f)\mapsto\tau_xf$, is jointly continuous. Given $x\in\mathbb{R}^3$ and $f\in\mathcal S$, choose $q_m\in\mathbb{Q}^3$ and $f_m\in V_0$ with $q_m\to x$ and $f_m\to f$ in $\mathcal S$; then $\tau_{q_m}f_m\to\tau_xf$ in $\mathcal S$, and
\[
\Phi_{\omega}(\tau_xf)=\lim_{m\to\infty}\Phi_{\omega}(\tau_{q_m}f_m)=\lim_{m\to\infty}\Phi_{\omega}(f_m)=\Phi_{\omega}(f).
\]
Because $X(f)\circ T_x=X(\tau_{-x}f)$, the characteristic functional of the law $\mu_{\omega}\circ T_x^{-1}$ is $f\mapsto\Phi_{\omega}(\tau_{-x}f)=\Phi_{\omega}(f)$, the same as that of $\mu_{\omega}$. By Lemma~\ref{lem:exponentials-determination-law}(b), applied with the sequentially dense set $\mathcal S$, one obtains $\mu_{\omega}\circ T_x^{-1}=\mu_{\omega}$ for every $x\in\mathbb{R}^3$.

\emph{Step 3: the laws $\mu_{\omega}$ are ergodic.} Let
\[
G_0=\left\{\sum_{l=1}^pc_l\cdot\exp\left(iX(f_l)\right):\ p\in\mathbb{N},\ c_1,\dots,c_p\in\mathbb{Q}+i\mathbb{Q},\ f_1,\dots,f_p\in V_0\right\}.
\]
This set is countable, and it is dense in $L^2(\mu')$ for every law $\mu'$ on $(\Omega,\mathcal F)$: by Lemma~\ref{lem:exponentials-determination-law}(b) the closure of the complex linear span of the functions $\exp\left(iX(f)\right)$, $f\in V_0$, is $L^2(\mu')$, and, these functions having modulus $1$, complex coefficients may be replaced by nearby elements of $\mathbb{Q}+i\mathbb{Q}$. For $G\in G_0$ let $\bar{\bar{G}}$ be a bounded $I$-measurable representative of $E[G\mid I]$.

Fix an enumeration of $G_0$. By Lemma~\ref{lem:averages-translations-2}(c), integers $L_1<L_2<\cdots$ can be chosen such that $\left\|A_{L_k}G-\bar{\bar{G}}\right\|^2_{L^2(\mu)}\le2^{-k}$ for the first $k$ elements $G$ of the enumeration; then $\sum_k\left\|A_{L_k}G-\bar{\bar{G}}\right\|^2_{L^2(\mu)}<\infty$ for every $G\in G_0$. The family $(\mu_{\omega})_{\omega\in\Omega}$ satisfies \eqref{eq:fourth-cumulants-ED1} with respect to $\left(\Omega,I,\mu|_I\right)$ (Step~1); therefore the argument of the proof of (b) shows that, for $\mu$-almost every $\omega$ and all $G\in G_0$,
\[
A_{L_k}G\longrightarrow\bar{\bar{G}}\quad\text{in }L^2(\mu_{\omega})\qquad(k\to\infty).
\]
Further, Step~1, applied to the bounded $I$-measurable functions $\left|\bar{\bar{G}}-c\right|^2$, gives for $\mu$-almost every $\omega$, all $G\in G_0$ and all $c\in\mathbb{Q}+i\mathbb{Q}$
\[
\int\mathrm{d}\mu_{\omega}\,\left|\bar{\bar{G}}-c\right|^2=E\left[\left|\bar{\bar{G}}-c\right|^2\,\middle|\,I\right](\omega)=\left|\bar{\bar{G}}(\omega)-c\right|^2.
\]
Each side depends continuously on $c\in\mathbb{C}$, the left side by dominated convergence, $\bar{\bar{G}}$ being bounded; hence the identity holds for every $c\in\mathbb{C}$, and the choice $c=\bar{\bar{G}}(\omega)$ shows that $\bar{\bar{G}}=\bar{\bar{G}}(\omega)$ holds $\mu_{\omega}$-almost surely.

Let $\omega$ lie outside the exceptional sets of Steps~2 and~3, and let $G\in G_0$. By Step~2 the law $\mu_{\omega}$ is translation invariant, and Lemma~\ref{lem:averages-translations-2}(b),(c), applied to $\mu'=\mu_{\omega}$, shows that $A_LG$ converges in $L^2(\mu_{\omega})$, as $L\to\infty$, to the orthogonal projection of $G$ onto the space $H'_{\mathrm{inv}}$ of the invariant elements of $L^2(\mu_{\omega})$. Along $(L_k)$ this limit is $\bar{\bar{G}}$, which is $\mu_{\omega}$-almost surely constant. So the projection onto $H'_{\mathrm{inv}}$ maps every element of the dense set $G_0$ to a constant; being continuous, it maps $L^2(\mu_{\omega})$ into the closed subspace of the constants, and, as it is the identity on $H'_{\mathrm{inv}}$, the space $H'_{\mathrm{inv}}$ consists of constants. Thus $\mu_{\omega}$ is ergodic. All the exceptional sets are $\mu$-null sets of $\mathcal F$, which belong to $I$ (Step~1); hence $\mu_{\omega}$ is translation invariant and ergodic for $\nu$-almost every $\omega$, which is \eqref{eq:fourth-cumulants-ED2}.
\end{proof}

\begin{lemma}[A conditional exponential bound]\label{lem:conditional-exponential-bound}
Let $\mu$ be a translation-invariant law on $(\Omega,\mathcal F)$ that satisfies \eqref{eq:fourth-cumulants-H1}, let $E$ be its expectation and $I$ the $\sigma$-algebra $I'$ of Lemma~\ref{lem:averages-translations-2} formed with $\mu'=\mu$, and let $f\in\mathcal S$.
The inequality $E\left[\exp\left(X(f)\right)\mid I\right]\le\exp\left(\frac{1}{2}\|f\|_2^2\right)$ holds $\mu$-almost surely.
\end{lemma}

\begin{proof}
Assume first that $f\in\mathcal D$, and let $r>0$ be such that the support of $f$ lies in a closed ball of radius $r$. The function $G=\exp\left(X(f)\right)$ is non-negative, and it belongs to $L^2(\mu)$, since $G^2=\exp\left(X(2f)\right)$ is integrable by \eqref{eq:fourth-cumulants-H1}. For $x\in\mathbb{R}^3$ one has $G\circ T_{-x}=\exp\left(X(\tau_xf)\right)$, because $(T_{-x}\omega)(f)=\omega(\tau_xf)$. Let $n\ge1$ be an integer and $L>0$. The $n$-th power of $A_LG$ is an integral over $\left([0,L]^3\right)^n$; the linearity of $X$, the theorem of Tonelli and \eqref{eq:fourth-cumulants-H1}, applied to the tests $\sum_{i=1}^n\tau_{x_i}f$, give
\begin{align*}
E\left[(A_LG)^n\right]&=L^{-3n}\int_{([0,L]^3)^n}\mathrm{d}x_1\cdots\mathrm{d}x_n\,E\left[\exp\left(X\left(\sum_{i=1}^n\tau_{x_i}f\right)\right)\right]\\
&\le L^{-3n}\int_{([0,L]^3)^n}\mathrm{d}x_1\cdots\mathrm{d}x_n\,\exp\left(\frac{1}{2}\left\|\sum_{i=1}^n\tau_{x_i}f\right\|_2^2\right).
\end{align*}
If $|x_i-x_l|>2r$ for all $i\ne l$, the supports of the forms $\tau_{x_i}f$ are pairwise disjoint, and $\left\|\sum_i\tau_{x_i}f\right\|_2^2=n\cdot\|f\|_2^2$. The other points of $\left([0,L]^3\right)^n$ lie in the union, over the pairs $i<l$, of the sets $\left\{|x_i-x_l|\le2r\right\}$; this union has measure at most $\binom{n}{2}\cdot\frac{4\pi}{3}(2r)^3\cdot L^{3(n-1)}$, and on it $\left\|\sum_i\tau_{x_i}f\right\|_2^2\le n^2\cdot\|f\|_2^2$ by the triangle inequality. Therefore
\[
E\left[(A_LG)^n\right]\le\exp\left(\frac{n}{2}\|f\|_2^2\right)+\binom{n}{2}\cdot\frac{4\pi}{3}(2r)^3\cdot L^{-3}\cdot\exp\left(\frac{n^2}{2}\|f\|_2^2\right),
\]
and consequently
\[
\limsup_{L\to\infty}E\left[(A_LG)^n\right]\le\exp\left(\frac{n}{2}\|f\|_2^2\right).
\]
By Lemma~\ref{lem:averages-translations-2}(c), $A_LG\to E[G\mid I]$ in $L^2(\mu)$ as $L\to\infty$; hence there are $L_k\to\infty$ such that $A_{L_k}G$ converges $\mu$-almost surely to a representative $\bar{\bar{G}}$ of $E[G\mid I]$, and $\bar{\bar{G}}\ge0$ almost surely, since $A_LG\ge0$. Fatou's lemma gives $E\left[\bar{\bar{G}}^n\right]\le\liminf_{k\to\infty}E\left[(A_{L_k}G)^n\right]\le\exp\left(\frac{n}{2}\|f\|_2^2\right)$, that is,
\[
\left\|\bar{\bar{G}}\right\|_{L^n(\mu)}\le\exp\left(\frac{1}{2}\|f\|_2^2\right)\qquad\text{for every integer }n\ge1 .
\]
Since $\left\|\bar{\bar{G}}\right\|_{L^n(\mu)}\to\left\|\bar{\bar{G}}\right\|_{L^{\infty}(\mu)}$ as $n\to\infty$, it follows that $\bar{\bar{G}}\le\exp\left(\frac{1}{2}\|f\|_2^2\right)$ holds $\mu$-almost surely.

Now let $f\in\mathcal S$ be arbitrary, and choose $f_k\in\mathcal D$ with $f_k\to f$ in $\mathcal S$. Then $X(f_k)\to X(f)$ pointwise on $\Omega$, and $\|f_k\|_2\to\|f\|_2$. The conditional form of Fatou's lemma and the case already treated give, $\mu$-almost surely,
\[
E\left[\exp\left(X(f)\right)\mid I\right]\le\liminf_{k\to\infty}E\left[\exp\left(X(f_k)\right)\mid I\right]\le\liminf_{k\to\infty}\exp\left(\frac{1}{2}\|f_k\|_2^2\right)=\exp\left(\frac{1}{2}\|f\|_2^2\right).
\]
\end{proof}

\begin{lemma}[Consequences of reflection positivity]\label{lem:consequences-reflection-positivity}
Let $\mu$ be a law on $(\Omega,\mathcal F)$ which satisfies~\eqref{eq:fourth-cumulants-H2} and~\eqref{eq:fourth-cumulants-H3}, and let $E$ be its expectation. By~\eqref{eq:fourth-cumulants-H2} the law $\mu$ is translation invariant, and $I$, $I_j$, $H_{\mathrm{inv}}$, $H_j$, $P=E[\,\cdot\mid I]$ and $P_j=E[\,\cdot\mid I_j]$ denote the objects of Lemma~\ref{lem:averages-translations-2} formed with $\mu'=\mu$. For real-valued $Y,Z\in L^2(\mu)$ put
\[
\operatorname{Var}(Y)=E\left[Y^2\right]-\left(E[Y]\right)^2,\qquad\operatorname{Cov}(Y,Z)=E\left[Y\cdot Z\right]-E[Y]\cdot E[Z].
\]
Let $j\in\{1,2,3\}$.
\begin{enumerate}
\item[(a)] Every $G\in L^2(\Omega,\Sigma_j^+,\mu)$ satisfies $E\left[(\Theta_jG)\cdot G\right]\ge0$.
\item[(b)] Let $H$ be a real-valued element of $H_j$, that is, $\mathbb{U}_{te_j}H=H$ for every $t\in\mathbb{R}$. Then $H$ has a $\Sigma_j^+$-measurable representative, and $\Theta_jH=H$. Moreover $H_{\mathrm{inv}}\subset H_j$; consequently each real-valued element $H$ of $H_{\mathrm{inv}}$ satisfies $\Theta_jH=H$ and admits a $\Sigma_j^+$-measurable representative.
\item[(c)] Let $G$ be a real-valued element of $L^2(\Omega,\Sigma_j^+,\mu)$, and define the function $\varphi_G\colon[0,\infty)\to\mathbb{R}$ by
\[
\varphi_G(t)=E\left[(\Theta_jG)\cdot\mathbb{U}_{te_j}G\right].
\]
The function $\varphi_G$ is non-negative and non-increasing, and
\[
\lim_{t\to\infty}\varphi_G(t)=\|P_jG\|_2^2\ge\left\|E[G\mid I]\right\|_2^2 .
\]
Consequently, for every $t\ge0$,
\begin{equation}\label{eq:fourth-cumulants-2}
\operatorname{Var}\left(E[G\mid I]\right)\le\operatorname{Var}\left(P_jG\right)\le\operatorname{Cov}\left(\Theta_jG,\mathbb{U}_{te_j}G\right).
\end{equation}
\end{enumerate}
\end{lemma}

\begin{proof}
\emph{The operators $\Theta_j$ and $\mathbb{U}_{te_j}$.} The reflection $\theta_j$ belongs to $O_{\mathrm{cub}}$; hence~\eqref{eq:fourth-cumulants-H2} gives $\mu\circ T_{0,\theta_j}^{-1}=\mu$, as well as $\mu\circ T_x^{-1}=\mu$ for every $x\in\mathbb{R}^3$. Since $\theta_j^2=1$, the map $U_{0,\theta_j}$ is its own inverse. Thus $(T_{0,\theta_j}\omega)(f)=\omega(\theta_jf)$ for $\omega\in\Omega$ and $f\in\mathcal S$, the map $T_{0,\theta_j}$ is an involution of $\Omega$, and
\[
\Theta_jG=\overline{G\circ T_{0,\theta_j}}
\]
for every function $G$ on $\Omega$. By the invariance of $\mu$ under $T_{0,\theta_j}$, the operator $\Theta_j$ maps functions that vanish $\mu$-almost everywhere to functions that vanish $\mu$-almost everywhere, and it preserves the norm of $L^2(\mu)$. It follows that $\Theta_j$ is an antilinear isometry of $L^2(\mu)$ onto itself, with $\Theta_j\Theta_jG=G$ for $G\in L^2(\mu)$. If $G$ is real-valued, then $\Theta_jG=G\circ T_{0,\theta_j}$ is real-valued; on the real-valued elements of $L^2(\mu)$ the operator $\Theta_j$ is linear, and it is symmetric: for real-valued $G,H\in L^2(\mu)$ the invariance of $\mu$ under the involution $T_{0,\theta_j}$ gives
\[
E\left[(\Theta_jG)\cdot H\right]=E\left[\left((\Theta_jG)\cdot H\right)\circ T_{0,\theta_j}\right]=E\left[G\cdot\Theta_jH\right],
\]
that is, $\langle\Theta_jG,H\rangle=\langle G,\Theta_jH\rangle$.

Let $t\in\mathbb{R}$. As $\theta_je_j=-e_j$, for every $f\in\mathcal S$ both $\theta_j\tau_{te_j}f$ and $\tau_{-te_j}\theta_jf$ are equal to the form $x\mapsto(\wedge^2\theta_j)\,f(\theta_jx-te_j)$. Since $(T_x\omega)(f)=\omega(\tau_{-x}f)$ for $x\in\mathbb{R}^3$, it follows that
\[
(T_{0,\theta_j}T_{te_j}\omega)(f)=\omega(\tau_{-te_j}\theta_jf)=\omega(\theta_j\tau_{te_j}f)=(T_{-te_j}T_{0,\theta_j}\omega)(f)\qquad(\omega\in\Omega,\ f\in\mathcal S),
\]
that is, $T_{0,\theta_j}\circ T_{te_j}=T_{-te_j}\circ T_{0,\theta_j}$. As $\mathbb{U}_{te_j}G=G\circ T_{te_j}^{-1}=G\circ T_{-te_j}$, one obtains, for every function $G$ on $\Omega$, the commutation relation
\[
\Theta_j\mathbb{U}_{te_j}G=\overline{G\circ T_{-te_j}\circ T_{0,\theta_j}}=\overline{G\circ T_{0,\theta_j}\circ T_{te_j}}=\mathbb{U}_{-te_j}\Theta_jG .
\]
Finally, by the translation invariance of $\mu$, every operator $\mathbb{U}_{te_j}$ preserves the norm of $L^2(\mu)$, and $E\left[\mathbb{U}_{te_j}F\right]=E[F]$ for every $\mu$-integrable function $F$.

(a) Let $V_j^+$ be the set of the forms $f\in\mathcal D$ with $\operatorname{supp}f\subset\{x_j>0\}$. It is a linear subspace of $\mathcal S$, because $\operatorname{supp}(a\cdot f+b\cdot g)\subset\operatorname{supp}f\cup\operatorname{supp}g$ for $f,g\in\mathcal D$ and $a,b\in\mathbb{R}$. By definition, $\Sigma_j^+$ is the $\sigma$-algebra generated by the functions $X(f)$, $f\in V_j^+$, and $A_j^+$ is the complex linear span of the functions $\exp\left(iX(f)\right)$, $f\in V_j^+$. Lemma~\ref{lem:exponentials-determination-law}(a), applied with $\mu'=\mu$ and $V=V_j^+$, therefore shows that $A_j^+$ is dense in $L^2(\Omega,\Sigma_j^+,\mu)$.

For $G,G'\in L^2(\mu)$ one has $\Theta_jG-\Theta_jG'=\Theta_j(G-G')$, and hence
\[
E\left[(\Theta_jG)\cdot G\right]-E\left[(\Theta_jG')\cdot G'\right]=E\left[(\Theta_jG)\cdot(G-G')\right]+E\left[\Theta_j(G-G')\cdot G'\right].
\]
The Cauchy--Schwarz inequality and the isometry of $\Theta_j$ give
\[
\left|E\left[(\Theta_jG)\cdot G\right]-E\left[(\Theta_jG')\cdot G'\right]\right|\le\|G-G'\|_2\cdot\left(\|G\|_2+\|G'\|_2\right).
\]
Now let $G\in L^2(\Omega,\Sigma_j^+,\mu)$, and choose $G_n\in A_j^+$, $n\ge1$, with $\|G_n-G\|_2\to0$. By~\eqref{eq:fourth-cumulants-H3}, $E\left[(\Theta_jG_n)\cdot G_n\right]\ge0$ for every $n$, and by the preceding inequality these numbers converge to $E\left[(\Theta_jG)\cdot G\right]$. Hence $E\left[(\Theta_jG)\cdot G\right]\ge0$.

(b) Let $H$ be a real-valued element of $H_j$, and let $\varepsilon>0$. The subspace $\mathcal D$ is sequentially dense in $\mathcal S$: if $\chi\in C_c^\infty(\mathbb{R}^3)$ equals $1$ in a neighbourhood of $0$, then for every $f\in\mathcal S$ the forms $x\mapsto\chi\left(\frac{x}{n}\right)\cdot f(x)$ converge to $f$ in $\mathcal S$ as $n\to\infty$. Lemma~\ref{lem:exponentials-determination-law}(b), applied with $\mu'=\mu$ and $V_0=\mathcal D$, therefore provides an integer $p\ge1$, numbers $c_1,\dots,c_p\in\mathbb{C}$ and forms $f_1,\dots,f_p\in\mathcal D$ such that the function
\[
G_\varepsilon=\sum_{l=1}^pc_l\cdot\exp\left(iX(f_l)\right)
\]
satisfies $\|H-G_\varepsilon\|_2<\varepsilon$. Let $r>0$ be such that the supports of $f_1,\dots,f_p$ lie in the closed ball of radius $r$ about $0$, and let $s>r$. The support of $\tau_{se_j}f_l$ is the translate by $se_j$ of the support of $f_l$; it lies in $\{x_j\ge s-r\}\subset\{x_j>0\}$. Since $\mathbb{U}_{se_j}$ is the composition with $T_{se_j}^{-1}$ and $\mathbb{U}_{se_j}X(f)=X(\tau_{se_j}f)$ for $f\in\mathcal S$, it follows that
\[
\mathbb{U}_{se_j}G_\varepsilon=\sum_{l=1}^pc_l\cdot\exp\left(iX(\tau_{se_j}f_l)\right)\in A_j^+\subset L^2(\Omega,\Sigma_j^+,\mu).
\]
As $\mathbb{U}_{se_j}H=H$ and $\mathbb{U}_{se_j}$ preserves the norm of $L^2(\mu)$,
\[
\left\|H-\mathbb{U}_{se_j}G_\varepsilon\right\|_2=\left\|\mathbb{U}_{se_j}\left(H-G_\varepsilon\right)\right\|_2=\|H-G_\varepsilon\|_2<\varepsilon .
\]
Thus $H$ belongs to the closure of $L^2(\Omega,\Sigma_j^+,\mu)$ in $L^2(\mu)$. This subspace is complete, hence closed; so $H\in L^2(\Omega,\Sigma_j^+,\mu)$, which means that $H$ has a $\Sigma_j^+$-measurable representative.

The function $\Theta_jH$ is real-valued, and for $t\in\mathbb{R}$ the commutation relation, with $-t$ in place of $t$, gives $\mathbb{U}_{te_j}\Theta_jH=\Theta_j\mathbb{U}_{-te_j}H=\Theta_jH$. Hence $\Theta_j$ maps the real-valued elements of $H_j$ to real-valued elements of $H_j$. Put
\[
H_-=\frac{1}{2}\left(H-\Theta_jH\right).
\]
Then $H_-$ is a real-valued element of $H_j$, and $\Theta_jH_-=\frac{1}{2}\left(\Theta_jH-H\right)=-H_-$, because $\Theta_j\Theta_jH=H$. By the first part of the proof of (b), applied to $H_-$ in place of $H$, the function $H_-$ belongs to $L^2(\Omega,\Sigma_j^+,\mu)$, and (a) gives
\[
0\le E\left[(\Theta_jH_-)\cdot H_-\right]=-E\left[H_-^2\right]=-\|H_-\|_2^2 .
\]
Hence $H_-=0$, that is, $\Theta_jH=H$.

Finally, every $H\in H_{\mathrm{inv}}$ satisfies $\mathbb{U}_xH=H$ for every $x\in\mathbb{R}^3$, in particular for $x=te_j$ with $t\in\mathbb{R}$. Therefore $H_{\mathrm{inv}}\subset H_j$, and the assertions of (b) hold for every real-valued $H\in H_{\mathrm{inv}}$.

(c) Let $L^2_{\mathbb{R}}(\Sigma_j^+)$ be the real vector space of the real-valued elements of $L^2(\Omega,\Sigma_j^+,\mu)$, and define the form $\langle\cdot,\cdot\rangle_{\Theta_j}\colon L^2_{\mathbb{R}}(\Sigma_j^+)\times L^2_{\mathbb{R}}(\Sigma_j^+)\to\mathbb{R}$ by
\[
\langle G,H\rangle_{\Theta_j}=E\left[(\Theta_jG)\cdot H\right]\qquad\left(G,H\in L^2_{\mathbb{R}}(\Sigma_j^+)\right).
\]
It is bilinear because $\Theta_j$ is linear on real-valued functions, symmetric by the symmetry of $\Theta_j$, and positive semidefinite by (a). Consequently it obeys the Cauchy--Schwarz inequality
\[
\langle G,H\rangle_{\Theta_j}\le\langle G,G\rangle_{\Theta_j}^{\frac{1}{2}}\cdot\langle H,H\rangle_{\Theta_j}^{\frac{1}{2}}\qquad\left(G,H\in L^2_{\mathbb{R}}(\Sigma_j^+)\right).
\]

For $t\ge0$ the operator $\mathbb{U}_{te_j}$ maps $L^2_{\mathbb{R}}(\Sigma_j^+)$ into itself. Indeed, let $f\in V_j^+$; then $\tau_{te_j}f\in\mathcal D$ has its support in $\{x_j>t\}\subset\{x_j>0\}$, so that $X(f)\circ T_{-te_j}=\mathbb{U}_{te_j}X(f)=X(\tau_{te_j}f)$ is $\Sigma_j^+$-measurable. As $\Sigma_j^+$ is the $\sigma$-algebra generated by these functions $X(f)$, the map $T_{-te_j}$ is measurable from $(\Omega,\Sigma_j^+)$ to $(\Omega,\Sigma_j^+)$, and $\mathbb{U}_{te_j}G=G\circ T_{-te_j}$ is $\Sigma_j^+$-measurable whenever $G$ is; moreover $\mathbb{U}_{te_j}$ preserves real values and the norm of $L^2(\mu)$.

Let $G\in L^2_{\mathbb{R}}(\Sigma_j^+)$ and $s,s'\ge0$. The operator $\mathbb{U}_{-se_j}$, a composition with a map of $\Omega$, is multiplicative, and $\mathbb{U}_{-se_j}\mathbb{U}_{(s+s')e_j}=\mathbb{U}_{s'e_j}$. The commutation relation and the identity $E\left[\mathbb{U}_{-se_j}F\right]=E[F]$ therefore give
\begin{align*}
\left\langle\mathbb{U}_{se_j}G,\mathbb{U}_{s'e_j}G\right\rangle_{\Theta_j}&=E\left[\left(\mathbb{U}_{-se_j}\Theta_jG\right)\cdot\mathbb{U}_{s'e_j}G\right]\\
&=E\left[\mathbb{U}_{-se_j}\left((\Theta_jG)\cdot\mathbb{U}_{(s+s')e_j}G\right)\right]=\varphi_G(s+s').
\end{align*}
For $t\ge0$, the choice $s=s'=\frac{t}{2}$ in this identity gives
\[
\varphi_G(t)=\left\langle\mathbb{U}_{\frac{t}{2}e_j}G,\mathbb{U}_{\frac{t}{2}e_j}G\right\rangle_{\Theta_j}\ge0,
\]
since $\mathbb{U}_{\frac{t}{2}e_j}G\in L^2_{\mathbb{R}}(\Sigma_j^+)$; so $\varphi_G$ is non-negative. The choices $(s,s')=(0,t)$, $(0,0)$ and $(t,t)$, together with the Cauchy--Schwarz inequality for $\langle\cdot,\cdot\rangle_{\Theta_j}$, give
\[
\varphi_G(t)=\left\langle G,\mathbb{U}_{te_j}G\right\rangle_{\Theta_j}\le\langle G,G\rangle_{\Theta_j}^{\frac{1}{2}}\cdot\left\langle\mathbb{U}_{te_j}G,\mathbb{U}_{te_j}G\right\rangle_{\Theta_j}^{\frac{1}{2}}=\varphi_G(0)^{\frac{1}{2}}\cdot\varphi_G(2t)^{\frac{1}{2}}\qquad(t\ge0).
\]
By induction on $n$ one obtains
\[
\varphi_G(t)\le\varphi_G(0)^{1-2^{-n}}\cdot\varphi_G(2^nt)^{2^{-n}}\qquad(t\ge0,\ n\ge1):
\]
for $n=1$ this is the preceding inequality, and the preceding inequality at the point $2^nt$, namely $\varphi_G(2^nt)\le\varphi_G(0)^{\frac{1}{2}}\cdot\varphi_G(2^{n+1}t)^{\frac{1}{2}}$, leads from $n$ to $n+1$, all the numbers involved being non-negative. By the Cauchy--Schwarz inequality in $L^2(\mu)$ and the invariance of the norm under $\Theta_j$ and under $\mathbb{U}_{ue_j}$,
\[
\varphi_G(u)\le\|\Theta_jG\|_2\cdot\left\|\mathbb{U}_{ue_j}G\right\|_2=\|G\|_2^2\qquad(u\ge0).
\]
Therefore
\[
\varphi_G(t)\le\varphi_G(0)^{1-2^{-n}}\cdot\|G\|_2^{2^{1-n}}\qquad(t\ge0,\ n\ge1).
\]
If $G=0$, then $\varphi_G=0$. Otherwise $\|G\|_2^{2^{1-n}}\to1$ as $n\to\infty$, and the right side tends to $\varphi_G(0)$. In both cases
\[
\varphi_G(t)\le\varphi_G(0)\qquad\left(t\ge0,\ G\in L^2_{\mathbb{R}}(\Sigma_j^+)\right).
\]
Now let $0\le s\le t$, and put $G'=\mathbb{U}_{\frac{s}{2}e_j}G$, an element of $L^2_{\mathbb{R}}(\Sigma_j^+)$. The identity for $\langle\cdot,\cdot\rangle_{\Theta_j}$ gives, for $u\ge0$,
\[
\varphi_{G'}(u)=\left\langle G',\mathbb{U}_{ue_j}G'\right\rangle_{\Theta_j}=\left\langle\mathbb{U}_{\frac{s}{2}e_j}G,\mathbb{U}_{\left(u+\frac{s}{2}\right)e_j}G\right\rangle_{\Theta_j}=\varphi_G(u+s).
\]
The inequality $\varphi_{G'}(t-s)\le\varphi_{G'}(0)$ thus reads $\varphi_G(t)\le\varphi_G(s)$, and $\varphi_G$ is non-increasing.

Being non-increasing and non-negative, $\varphi_G$ has a limit $\ell=\lim_{t\to\infty}\varphi_G(t)\ge0$, and $\ell\le\varphi_G(t)\le\varphi_G(0)$ for every $t\ge0$; in particular $\varphi_G$ is a bounded Borel function on $[0,\infty)$. Its Ces\`aro means converge to the same limit:
\[
\lim_{T\to\infty}\frac{1}{T}\int_0^T\mathrm{d}t\,\varphi_G(t)=\ell .
\]
Indeed, given $\eta>0$, choose $t_0>0$ with $\varphi_G(t)\le\ell+\eta$ for $t\ge t_0$; then, for $T>t_0$, the mean lies between $\ell$ and $\ell+\eta+\frac{t_0}{T}\cdot\varphi_G(0)$.

On the other hand, let a measurable representative of $G$ be fixed. By the joint measurability asserted in Lemma~\ref{lem:averages-translations-2}, the function $(t,\omega)\mapsto G(T_{-te_j}\omega)$ on $\mathbb{R}\times\Omega$ is jointly measurable, and
\[
\int_0^T\mathrm{d}t\,E\left[\left|\Theta_jG\right|\cdot\left|\mathbb{U}_{te_j}G\right|\right]\le T\cdot\|G\|_2^2<\infty\qquad(T>0).
\]
Fubini's theorem therefore gives, for $T>0$,
\[
\frac{1}{T}\int_0^T\mathrm{d}t\,\varphi_G(t)=E\left[(\Theta_jG)\cdot\frac{1}{T}\int_0^T\mathrm{d}t\,\mathbb{U}_{te_j}G\right],
\]
where the mean on the right is the function $\omega\mapsto\frac{1}{T}\int_0^T\mathrm{d}t\,G(T_{-te_j}\omega)$ of Lemma~\ref{lem:averages-translations-2}(d). By Lemma~\ref{lem:averages-translations-2}(d), this mean converges to $P_jG$ in $L^2(\mu)$ as $T\to\infty$; by the Cauchy--Schwarz inequality the right side converges to $E\left[(\Theta_jG)\cdot P_jG\right]$. Consequently
\[
\ell=E\left[(\Theta_jG)\cdot P_jG\right]=E\left[G\cdot\Theta_jP_jG\right]=E\left[G\cdot P_jG\right]=\|P_jG\|_2^2 .
\]
Here $P_jG$ is real-valued, as the conditional expectation of a real-valued function. The second equality is the symmetry of $\Theta_j$. By Lemma~\ref{lem:averages-translations-2}(d), $P_j$ is the orthogonal projection of $L^2(\mu)$ onto $H_j$; so $P_jG$ is a real-valued element of $H_j$, and (b) gives $\Theta_jP_jG=P_jG$, which is the third equality. The last one holds because $\langle G,P_jG\rangle=\langle P_jG,P_jG\rangle$ for an orthogonal projection $P_j$.

By Lemma~\ref{lem:averages-translations-2}(b), $P=E[\,\cdot\mid I]$ is the orthogonal projection of $L^2(\mu)$ onto $H_{\mathrm{inv}}$, and $H_{\mathrm{inv}}\subset H_j$ by (b). The function $G-P_jG$ is orthogonal to $H_j$, hence to $H_{\mathrm{inv}}$; so $PG=P(P_jG)$, and, an orthogonal projection having norm at most $1$,
\[
\left\|E[G\mid I]\right\|_2=\left\|P(P_jG)\right\|_2\le\|P_jG\|_2 .
\]
This proves the assertions on $\varphi_G$.

Finally, let $t\ge0$. As $\varphi_G$ is non-increasing with limit $\|P_jG\|_2^2$, one has $\varphi_G(t)\ge\|P_jG\|_2^2$. The conditional expectations satisfy $E\left[E[G\mid I]\right]=E[G]=E\left[P_jG\right]$, and the invariance of $\mu$ under $T_{0,\theta_j}$ and under the translations gives $E\left[\Theta_jG\right]=E[G]=E\left[\mathbb{U}_{te_j}G\right]$. Therefore
\begin{align*}
\operatorname{Var}\left(E[G\mid I]\right)&=\left\|E[G\mid I]\right\|_2^2-\left(E[G]\right)^2\le\|P_jG\|_2^2-\left(E[G]\right)^2=\operatorname{Var}\left(P_jG\right)\\
&\le\varphi_G(t)-E\left[\Theta_jG\right]\cdot E\left[\mathbb{U}_{te_j}G\right]=\operatorname{Cov}\left(\Theta_jG,\mathbb{U}_{te_j}G\right),
\end{align*}
which is~\eqref{eq:fourth-cumulants-2}.
\end{proof}

\begin{corollary}[Properties of almost every component]\label{cor:properties-almost-every}
Let $\mu$ be a translation-invariant law on $(\Omega,\mathcal F)$ that satisfies \eqref{eq:fourth-cumulants-H1}, and let $(\Xi,\nu,(\mu_{\xi}))$ be an ergodic decomposition of $\mu$. There is a set $\Xi_0\subset\Xi$ of full $\nu$-measure such that, for every $\xi\in\Xi_0$, the law $\mu_{\xi}$ is translation invariant and ergodic and satisfies \eqref{eq:fourth-cumulants-H1}, that is, $E_{\xi}\exp\left(X(f)\right)\le\exp\left(\frac{1}{2}\|f\|_2^2\right)$ for every $f\in\mathcal S$. Consequently Lemma~\ref{lem:moment-bounds} holds for every law $\mu_{\xi}$ with $\xi\in\Xi_0$. In particular, for all $\xi\in\Xi_0$ and $f\in\mathcal S$ one has $E_{\xi}X(f)=0$ and, with the constant $c_4$ of Lemma~\ref{lem:moment-bounds}(c),
\[
0\le C_{\xi}(f,f)\le\|f\|_2^2,\qquad E_{\xi}\left[X(f)^4\right]\le c_4\cdot\|f\|_2^4,
\]
as well as
\begin{equation}\label{eq:fourth-cumulants-3}
-2\|f\|_2^4\le\kappa_4^{\xi}(f)\le c_4\cdot\|f\|_2^4.
\end{equation}
\end{corollary}

\begin{proof}
By \eqref{eq:fourth-cumulants-ED2} there is a $\nu$-null set $N_1\subset\Xi$ such that $\mu_{\xi}$ is translation invariant and ergodic for every $\xi\in\Xi\setminus N_1$.

Let $V_0\subset\mathcal S$ be the set introduced in the proof of Proposition~\ref{prop:existence-conditional-expectations}(a), namely the $\mathbb{Q}$-linear span of the functions $\tau_qh_n$ with $q\in\mathbb{Q}^3$, where $(h_n)$ is the Hermite basis of $L^2(\mathbb{R}^3;\Lambda^2)$. The set $V_0$ is countable, and every $f\in\mathcal S$ is the limit in $\mathcal S$ of a sequence of elements of $V_0$.

Let $f\in V_0$, and let $G=\exp\left(X(f)\right)$, a positive function on $\Omega$. Since $G^2=\exp\left(X(2f)\right)$, hypothesis \eqref{eq:fourth-cumulants-H1} for the test $2f$ shows that $G\in L^2(\mu)$. By Lemma~\ref{lem:conditional-exponential-bound}, the inequality $E[G\mid I]\le\exp\left(\frac{1}{2}\|f\|_2^2\right)$ holds $\mu$-almost surely. For every $\xi\in\Xi$, $E_{\xi}G=\int\mathrm{d}\mu_{\xi}\,G$ is the expectation of $G$ with respect to $\mu_{\xi}$ in the sense of Definition~\ref{def:ergodicity-ergodic-decompositions}; it lies in $(0,\infty]$, as $G$ is positive and $\mu_{\xi}$ is a probability measure. Where this integral is finite, it is the number $E_{\xi}[G]$ of Proposition~\ref{prop:existence-conditional-expectations}(b); where it is infinite, one has $E_{\xi}G=\infty$, whereas that number is $0$ by the convention fixed in the proof of Proposition~\ref{prop:existence-conditional-expectations}(b). The function $G$ is $\mathcal F$-measurable with values in $[0,\infty)$, so the extended form of \eqref{eq:fourth-cumulants-ED1} given in Section~\ref{sec:cumulants-gaussian-laws} shows that $\xi\mapsto E_{\xi}G$ is a $B$-measurable map from $\Xi$ to $[0,\infty]$. Consequently the set $\left\{\xi\in\Xi:E_{\xi}G>\exp\left(\frac{1}{2}\|f\|_2^2\right)\right\}$ belongs to $B$, and it contains every $\xi$ with $E_{\xi}G=\infty$. Since $G$ is a representative with values in $[0,\infty)$ of an element of $L^2(\mu)$, the last display of the first paragraph of the proof of Proposition~\ref{prop:existence-conditional-expectations}(b) applies to $G$:
\[
\nu\left(\left\{\xi\in\Xi:\ \int\mathrm{d}\mu_{\xi}\,G\in S\right\}\right)=\mu\left(\left\{E[G\mid I]\in S\right\}\right)\qquad\text{for every Borel set }S\subset[0,\infty].
\]
The choice $S=\left(\exp\left(\frac{1}{2}\|f\|_2^2\right),\infty\right]$ and Lemma~\ref{lem:conditional-exponential-bound} give
\[
\nu\left(\left\{\xi\in\Xi:E_{\xi}G>\exp\left(\frac{1}{2}\|f\|_2^2\right)\right\}\right)=\mu\left(\left\{E[G\mid I]>\exp\left(\frac{1}{2}\|f\|_2^2\right)\right\}\right)=0.
\]
Let $N_2\subset\Xi$ be the union over $f\in V_0$ of the sets on the left. As $V_0$ is countable, $N_2$ is a $\nu$-null set that belongs to $B$; it contains every $\xi$ at which $E_{\xi}\exp\left(X(f)\right)=\infty$ for some $f\in V_0$, and outside $N_2$ the integrals $E_{\xi}\exp\left(X(f)\right)=\int\mathrm{d}\mu_{\xi}\,\exp\left(X(f)\right)$ satisfy
\[
E_{\xi}\exp\left(X(f)\right)\le\exp\left(\frac{1}{2}\|f\|_2^2\right)\qquad\left(\xi\in\Xi\setminus N_2,\ f\in V_0\right).
\]

Now let $\xi\in\Xi\setminus N_2$ and $f\in\mathcal S$, and let $(f_m)_{m\ge1}$ be a sequence in $V_0$ that converges to $f$ in $\mathcal S$. Each $\omega\in\Omega$ is a continuous linear form on $\mathcal S$; hence $X(f_m)(\omega)=\omega(f_m)$ tends to $\omega(f)=X(f)(\omega)$ for every $\omega\in\Omega$. Moreover $\|f_m\|_2\to\|f\|_2$, since convergence in $\mathcal S$ implies convergence in $L^2(\mathbb{R}^3;\Lambda^2)$. Fatou's lemma, applied under $\mu_{\xi}$ to the positive functions $\exp\left(X(f_m)\right)$, which converge pointwise to $\exp\left(X(f)\right)$, gives
\[
E_{\xi}\exp\left(X(f)\right)\le\liminf_{m\to\infty}E_{\xi}\exp\left(X(f_m)\right)\le\liminf_{m\to\infty}\exp\left(\frac{1}{2}\|f_m\|_2^2\right)=\exp\left(\frac{1}{2}\|f\|_2^2\right).
\]
Put $\Xi_0=\Xi\setminus\left(N_1\cup N_2\right)$. This set has full $\nu$-measure, and for every $\xi\in\Xi_0$ the law $\mu_{\xi}$ is translation invariant, ergodic and subject to \eqref{eq:fourth-cumulants-H1}.

Let $\xi\in\Xi_0$ and $f\in\mathcal S$. The only hypothesis that Lemma~\ref{lem:moment-bounds} imposes on the law is \eqref{eq:fourth-cumulants-H1}; so the lemma applies to $\mu_{\xi}$. Part~(a) gives $E_{\xi}X(f)=0$, and part~(c) with $k=2$ gives $E_{\xi}\left[X(f)^4\right]\le c_4\cdot\|f\|_2^4$. Thus $\mu_{\xi}$ has mean zero and finite fourth moments, and the forms $C_{\xi}=S_2^{\mu_{\xi}}$ and $\kappa_4^{\xi}=\kappa_4^{\mu_{\xi}}$ of Section~\ref{sec:cumulants-gaussian-laws} are defined. Part~(b) gives $C_{\xi}(f,f)=E_{\xi}\left[X(f)^2\right]\le\|f\|_2^2$, and $C_{\xi}(f,f)\ge0$, being the mean of a square. By definition of the fourth cumulant,
\[
\kappa_4^{\xi}(f)=E_{\xi}\left[X(f)^4\right]-3\left(E_{\xi}\left[X(f)^2\right]\right)^2=E_{\xi}\left[X(f)^4\right]-3\,C_{\xi}(f,f)^2.
\]
The subtracted term is non-negative, whence $\kappa_4^{\xi}(f)\le E_{\xi}\left[X(f)^4\right]\le c_4\cdot\|f\|_2^4$. By the Cauchy--Schwarz inequality, $E_{\xi}\left[X(f)^4\right]\ge\left(E_{\xi}\left[X(f)^2\right]\right)^2=C_{\xi}(f,f)^2$, so that $\kappa_4^{\xi}(f)\ge-2\,C_{\xi}(f,f)^2$. Finally $-2\,C_{\xi}(f,f)^2\ge-2\|f\|_2^4$, because $0\le C_{\xi}(f,f)\le\|f\|_2^2$. This proves \eqref{eq:fourth-cumulants-3}.
\end{proof}

\subsection{Scale mixtures of the free field}\label{sec:scale-mixtures-free}

\begin{lemma}[Gaussian laws with bounded covariance]\label{lem:gaussian-laws-bounded}
Let $C\colon\mathcal S\times\mathcal S\to\mathbb{R}$ be bilinear, symmetric and positive semidefinite, and assume that $C(f,f)\le\|f\|_2^2$ for every $f\in\mathcal S$.
\begin{enumerate}
\item[(a)] There is exactly one law $\gamma_C$ on $(\Omega,\mathcal F)$ such that
\[
\int\mathrm{d}\gamma_C\,\exp\left(iX(f)\right)=\exp\left(-\frac{1}{2}C(f,f)\right)\qquad\text{for every }f\in\mathcal S .
\]
The law $\gamma_C$ is centred and Gaussian, with covariance $C$: for $f\in\mathcal S$ the variable $X(f)$ is normal with mean $0$ and variance $C(f,f)$, and $\int\mathrm{d}\gamma_C\,X(f)\cdot X(h)=C(f,h)$ for $f,h\in\mathcal S$. Moreover
\[
\int\mathrm{d}\gamma_C\,\exp X(f)=\exp\left(\frac{1}{2}C(f,f)\right)\le\exp\left(\frac{1}{2}\|f\|_2^2\right)\qquad(f\in\mathcal S),
\]
so that $\gamma_C$ satisfies \eqref{eq:fourth-cumulants-H1} in place of $\mu$.
\item[(b)] Let a group of Euclidean motions be given such that $C(U_gf,U_gh)=C(f,h)$ for every $g$ in the group and all $f,h\in\mathcal S$. Then $\gamma_C\circ T_g^{-1}=\gamma_C$ for every $g$ in the group.
\item[(c)] Let $C$ be translation invariant, that is, $C(\tau_xf,\tau_xh)=C(f,h)$ for all $x\in\mathbb{R}^3$ and $f,h\in\mathcal S$. Then $\gamma_C$ is translation invariant and mixing: with the scalar product $\langle\cdot,\cdot\rangle$ of $L^2(\gamma_C)$,
\[
\lim_{|x|\to\infty}\langle G,\mathbb{U}_xH\rangle=\langle G,1\rangle\cdot\langle1,H\rangle\qquad\text{for all }G,H\in L^2(\gamma_C).
\]
In particular $\gamma_C$ is ergodic.
\item[(d)] Let $j\in\{1,2,3\}$. Assume that $C(\theta_jf,\theta_jh)=C(f,h)$ for all $f,h\in\mathcal S$, and that $C(\theta_jf,f)\ge0$ for every $f\in\mathcal D$ with $\operatorname{supp}f\subset\{x_j>0\}$. Then
\[
\int\mathrm{d}\gamma_C\,(\Theta_jG)\cdot G\ge0\qquad\text{for every }G\in A_j^+ .
\]
\end{enumerate}
\end{lemma}

\begin{proof}
(a) Define $\varphi\colon\mathcal S\to\mathbb{R}$ by $\varphi(f)=\exp\left(-\frac{1}{2}C(f,f)\right)$; then $\varphi(0)=1$. The function $\varphi$ is continuous on $\mathcal S$, because the norm $\|\cdot\|_2$ is continuous on $\mathcal S$ and, by the Cauchy--Schwarz inequality for the positive semidefinite form $C$ and by the bound $C(g,g)\le\|g\|_2^2$,
\[
\left|C(f,f)-C(h,h)\right|=\left|C(f-h,f+h)\right|\le C(f-h,f-h)^{\frac{1}{2}}\cdot C(f+h,f+h)^{\frac{1}{2}}\le\|f-h\|_2\cdot\|f+h\|_2\qquad(f,h\in\mathcal S).
\]
The function $\varphi$ is positive definite. Indeed, let $n\ge1$, $f_1,\dots,f_n\in\mathcal S$ and $c_1,\dots,c_n\in\mathbb{C}$; let $g_1,\dots,g_m$ be a basis of the linear span of $f_1,\dots,f_n$, and write $f_k=\sum_{a=1}^mt^{(k)}_a\cdot g_a$ with $t^{(k)}\in\mathbb{R}^m$. The real symmetric matrix $\left(C(g_a,g_{a'})\right)_{a,a'=1}^m$ is positive semidefinite; let $\eta$ be the centred normal law on $\mathbb{R}^m$ with this covariance matrix. Its characteristic function is
\[
\int\mathrm{d}\eta(y)\,\exp(it\cdot y)=\exp\left(-\frac{1}{2}\sum_{a,a'=1}^mt_at_{a'}\cdot C(g_a,g_{a'})\right)=\varphi\left(\sum_{a=1}^mt_a\cdot g_a\right)\qquad(t\in\mathbb{R}^m);
\]
in other words, restricted to the span of $f_1,\dots,f_n$ and written in the coordinates $t$, the function $\varphi$ coincides with the characteristic function of the Gaussian law $\eta$. It follows that
\begin{align*}
\sum_{l,k=1}^n\overline{c_l}\cdot c_k\cdot\varphi(f_k-f_l)&=\int\mathrm{d}\eta(y)\,\sum_{l,k=1}^n\overline{c_l}\cdot c_k\cdot\exp\left(i\left(t^{(k)}-t^{(l)}\right)\cdot y\right)\\
&=\int\mathrm{d}\eta(y)\,\left|\sum_{k=1}^nc_k\cdot\exp\left(it^{(k)}\cdot y\right)\right|^2\ge0 .
\end{align*}
By Minlos' theorem~\cite[Theorem~3.4.2]{GlimmJaffe}, a normalised, positive definite and continuous function on $\mathcal S$ is the characteristic functional of a probability measure on $(\Omega,\mathcal F)$; this provides a law $\gamma_C$ with the required property. If $\gamma'$ is a law on $(\Omega,\mathcal F)$ with the same property, Lemma~\ref{lem:exponentials-determination-law}(b), applied with $V_0=\mathcal S$, gives $\gamma'=\gamma_C$.

Let $f\in\mathcal S$. For $t\in\mathbb{R}$ one has $\int\mathrm{d}\gamma_C\,\exp\left(itX(f)\right)=\varphi(tf)=\exp\left(-\frac{1}{2}t^2\cdot C(f,f)\right)$, which is the characteristic function of the normal law with mean $0$ and variance $C(f,f)$. By the uniqueness of the Fourier transform of a probability measure on $\mathbb{R}$, the variable $X(f)$ has this law under $\gamma_C$; thus $\gamma_C$ is Gaussian in the sense of Section~\ref{sec:cumulants-gaussian-laws}, $\int\mathrm{d}\gamma_C\,X(f)=0$ and $\int\mathrm{d}\gamma_C\,X(f)^2=C(f,f)$. Since $X(f+h)=X(f)+X(h)$ for $f,h\in\mathcal S$,
\[
\int\mathrm{d}\gamma_C\,X(f)\cdot X(h)=\frac{1}{2}\left(C(f+h,f+h)-C(f,f)-C(h,h)\right)=C(f,h).
\]
Finally, the normal law with mean $0$ and variance $v\ge0$ has the moment generating function $t\mapsto\exp\left(\frac{1}{2}v\cdot t^2\right)$; its value at $t=1$, for $v=C(f,f)\le\|f\|_2^2$, gives the last display of (a).

(b) Let $g$ be an element of the group, and let $f\in\mathcal S$. Since $(T_g\omega)(f)=\omega(U_g^{-1}f)$ for $\omega\in\Omega$, one has $X(f)\circ T_g=X(U_g^{-1}f)$; and the hypothesis, applied to the tests $U_g^{-1}f$ and $U_g^{-1}f$, gives $C(U_g^{-1}f,U_g^{-1}f)=C(f,f)$. Hence
\[
\int\mathrm{d}\left(\gamma_C\circ T_g^{-1}\right)\,\exp\left(iX(f)\right)=\int\mathrm{d}\gamma_C\,\exp\left(iX(U_g^{-1}f)\right)=\exp\left(-\frac{1}{2}C(f,f)\right),
\]
and the uniqueness in (a) yields $\gamma_C\circ T_g^{-1}=\gamma_C$.

(c) The translations $(x,1)$, $x\in\mathbb{R}^3$, form a group of Euclidean motions; by (b), $\gamma_C$ is translation invariant. Consequently, for every $x\in\mathbb{R}^3$, the operator $\mathbb{U}_x$ is an isometry of $L^2(\gamma_C)$ onto itself, with inverse $\mathbb{U}_{-x}$. By (a), $\gamma_C$ satisfies \eqref{eq:fourth-cumulants-H1} in place of $\mu$, and its two-point function $S_2^{\gamma_C}$ is $C$; Lemma~\ref{lem:decay-two-point}(b), applied to the law $\gamma_C$, gives
\[
\lim_{|x|\to\infty}C(f,\tau_xh)=0\qquad(f,h\in\mathcal S).
\]
Let $f,h\in\mathcal S$. Since $X(f)+X(\tau_xh)=X(f+\tau_xh)$ and $C(\tau_xh,\tau_xh)=C(h,h)$, part (a) gives, as $|x|\to\infty$,
\begin{align*}
\int\mathrm{d}\gamma_C\,\exp\left(iX(f)\right)\cdot\exp\left(iX(\tau_xh)\right)&=\exp\left(-\frac{1}{2}\left[C(f,f)+C(h,h)+2C(f,\tau_xh)\right]\right)\\
&\to\exp\left(-\frac{1}{2}C(f,f)\right)\cdot\exp\left(-\frac{1}{2}C(h,h)\right)=\int\mathrm{d}\gamma_C\,\exp\left(iX(f)\right)\cdot\int\mathrm{d}\gamma_C\,\exp\left(iX(h)\right).
\end{align*}
As $\mathbb{U}_xX(h)=X(\tau_xh)$, one has $\mathbb{U}_x\exp\left(iX(h)\right)=\exp\left(iX(\tau_xh)\right)$; so, for $G=\exp\left(-iX(f)\right)$ and $H=\exp\left(iX(h)\right)$, the last relation reads $\langle G,\mathbb{U}_xH\rangle\to\langle G,1\rangle\cdot\langle1,H\rangle$. As $f$ ranges over $\mathcal S$, so does $-f$; thus the relation holds for all $G$, $H$ in the set $\left\{\exp\left(iX(f)\right):f\in\mathcal S\right\}$. Both $\langle G,\mathbb{U}_xH\rangle$ and $\langle G,1\rangle\cdot\langle1,H\rangle$ are conjugate linear in $G$ and linear in $H$; hence the relation holds for all $G$, $H$ in the complex linear span of this set, which is dense in $L^2(\gamma_C)$ by Lemma~\ref{lem:exponentials-determination-law}(a), applied with $V=\mathcal S$.

Now let $G,H\in L^2(\gamma_C)$ and $0<\varepsilon\le1$, and choose $G'$, $H'$ in this span with $\|G-G'\|\le\varepsilon$ and $\|H-H'\|\le\varepsilon$. Since $\|\mathbb{U}_xH\|=\|H\|$ and $\|1\|=1$, the Cauchy--Schwarz inequality gives
\[
\left|\langle G,\mathbb{U}_xH\rangle-\langle G',\mathbb{U}_xH'\rangle\right|\le\|G-G'\|\cdot\|H\|+\|G'\|\cdot\|H-H'\|\le\varepsilon\cdot\left(\|G\|+\|H\|+1\right),
\]
and the same bound holds for $\left|\langle G,1\rangle\cdot\langle1,H\rangle-\langle G',1\rangle\cdot\langle1,H'\rangle\right|$. Therefore
\[
\limsup_{|x|\to\infty}\left|\langle G,\mathbb{U}_xH\rangle-\langle G,1\rangle\cdot\langle1,H\rangle\right|\le2\varepsilon\cdot\left(\|G\|+\|H\|+1\right),
\]
and, $\varepsilon$ being arbitrary, $\gamma_C$ is mixing.

Let $G\in L^2(\gamma_C)$ satisfy $\mathbb{U}_xG=G$ for every $x\in\mathbb{R}^3$, and put $c=\int\mathrm{d}\gamma_C\,G=\langle1,G\rangle$. Then $\|G\|^2=\langle G,\mathbb{U}_xG\rangle$ for every $x$, and the mixing property gives $\|G\|^2=\langle G,1\rangle\cdot\langle1,G\rangle=|c|^2$. Hence $\int\mathrm{d}\gamma_C\,|G-c|^2=\|G\|^2-|c|^2=0$, that is, $G=c$ almost surely. So the only elements of $L^2(\gamma_C)$ that are invariant under all translations are the constants, and $\gamma_C$ is ergodic in the sense of Definition~\ref{def:ergodicity-ergodic-decompositions}.

(d) Let $G\in A_j^+$, say $G=\sum_{l=1}^nc_l\cdot\exp\left(iX(f_l)\right)$ with $n\ge1$, $c_l\in\mathbb{C}$, $f_l\in\mathcal D$ and $\operatorname{supp}f_l\subset\{x_j>0\}$. Since $\mathbb{U}_{0,\theta_j}X(f_l)=X(\theta_jf_l)$,
\[
\Theta_jG=\overline{\mathbb{U}_{0,\theta_j}G}=\sum_{l=1}^n\overline{c_l}\cdot\exp\left(-iX(\theta_jf_l)\right),
\]
and part (a) gives
\[
\int\mathrm{d}\gamma_C\,(\Theta_jG)\cdot G=\sum_{l,k=1}^n\overline{c_l}\cdot c_k\int\mathrm{d}\gamma_C\,\exp\left(iX(f_k-\theta_jf_l)\right)=\sum_{l,k=1}^n\overline{c_l}\cdot c_k\cdot\exp\left(-\frac{1}{2}C(f_k-\theta_jf_l,f_k-\theta_jf_l)\right).
\]
For $l,k\in\{1,\dots,n\}$ put
\[
M_{lk}=C(\theta_jf_l,f_k),\qquad d_l=c_l\cdot\exp\left(-\frac{1}{2}C(f_l,f_l)\right).
\]
The hypothesis gives $C(\theta_jf_l,\theta_jf_l)=C(f_l,f_l)$, so that
\[
C(f_k-\theta_jf_l,f_k-\theta_jf_l)=C(f_k,f_k)+C(f_l,f_l)-2M_{lk},
\]
and therefore
\[
\int\mathrm{d}\gamma_C\,(\Theta_jG)\cdot G=\sum_{l,k=1}^n\overline{d_l}\cdot d_k\cdot\exp(M_{lk}).
\]
The matrix $M=(M_{lk})$ is real. It is symmetric: $\theta_j\theta_jf=f$ for $f\in\mathcal S$, so that the hypothesis, applied to the pair $\theta_jf_l$, $f_k$, together with the symmetry of $C$ gives
\[
M_{lk}=C(\theta_j\theta_jf_l,\theta_jf_k)=C(f_l,\theta_jf_k)=C(\theta_jf_k,f_l)=M_{kl}.
\]
It is positive semidefinite: for $a\in\mathbb{R}^n$ the two-form $f=\sum_la_l\cdot f_l$ belongs to $\mathcal D$ and has its support in $\{x_j>0\}$, and $\sum_{l,k}a_la_k\cdot M_{lk}=C(\theta_jf,f)\ge0$. By the Schur product theorem, the entrywise product of two positive semidefinite real symmetric $n\times n$ matrices is positive semidefinite. By induction, the matrix $\left(M_{lk}^m\right)_{l,k=1}^n$ is positive semidefinite for every integer $m\ge1$; for $m=0$ it is the matrix all of whose entries equal $1$, which is positive semidefinite since $\sum_{l,k}a_la_k=\left(\sum_la_l\right)^2$. Hence the real symmetric matrix
\[
\left(\exp(M_{lk})\right)_{l,k=1}^n=\sum_{m=0}^\infty\frac{1}{m!}\left(M_{lk}^m\right)_{l,k=1}^n,
\]
a convergent series of positive semidefinite matrices, is positive semidefinite. Writing $d_l=a_l+ib_l$ with $a,b\in\mathbb{R}^n$ and using the symmetry of this matrix, one obtains
\[
\int\mathrm{d}\gamma_C\,(\Theta_jG)\cdot G=\sum_{l,k=1}^na_la_k\cdot\exp(M_{lk})+\sum_{l,k=1}^nb_lb_k\cdot\exp(M_{lk})\ge0 .
\]
\end{proof}

\begin{lemma}[The free transverse covariance]\label{lem:free-transverse-covariance}
For $f\in\mathcal S$ let $\vec f=(f_{23},-f_{13},f_{12})$ be the vector field associated with $f$; then $\langle f,g\rangle=\int\mathrm{d}x\,\vec f(x)\cdot\vec g(x)$ and $|f(x)|=|\vec f(x)|$ for $f,g\in\mathcal S$ and $x\in\mathbb{R}^3$, and in the same way $L^2(\mathbb{R}^3;\Lambda^2)$ is identified with $L^2(\mathbb{R}^3;\mathbb{R}^3)$. For every Euclidean motion $(b,r)$ the vector field of $U_{b,r}f$ is $x\mapsto\det(r)\cdot r\,\vec f\left(r^{-1}(x-b)\right)$, and $df=0$ if and only if $\operatorname{div}\vec f=0$. Put $j^f=\operatorname{curl}\vec f$. For a Schwartz three-form $h$, whose only component is $h_{123}$, let $d^*h$ be its image under the formal adjoint $d^*$ of the exterior derivative $d$, taken with respect to $\langle\cdot,\cdot\rangle$ and to the scalar product $\int\mathrm{d}x\,h_{123}(x)\cdot h'_{123}(x)$ of three-forms. The Fourier transform of an integrable function $\phi$ on $\mathbb{R}^3$, with values in $\mathbb{C}$ or in $\mathbb{C}^3$, is
\[
\hat\phi(p)=\int\mathrm{d}x\,\exp(-ip\cdot x)\cdot\phi(x)\qquad(p\in\mathbb{R}^3),
\]
and it is extended to square integrable functions by Plancherel's theorem; for $u,v\in\mathbb{C}^3$ put $u\cdot v=\sum_{k=1}^3u_kv_k$, and let $1$ be the unit $3\times3$ matrix. Let $\Pi$ be the orthogonal projection of $L^2(\mathbb{R}^3;\mathbb{R}^3)$ onto the closed subspace of the fields $v$ with $\operatorname{div}v=0$ in the sense of distributions. For $f,g\in\mathcal S$ put
\[
C_0(f,g)=(2\pi)^{-3}\int\mathrm{d}p\,\overline{\widehat{\vec f}(p)}\cdot\left(1-\frac{pp^T}{|p|^2}\right)\widehat{\vec g}(p).
\]
Then $C_0(f,g)=\langle f,\Pi g\rangle$, where $\langle f,\Pi g\rangle=\int\mathrm{d}x\,\vec f(x)\cdot(\Pi\vec g)(x)$, and $C_0$ is a symmetric bilinear form on $\mathcal S$. Moreover the following hold.
\begin{enumerate}
\item[(a)] For $f\in\mathcal S$ one has $0\le C_0(f,f)\le\|f\|_2^2$, with equality on the right if and only if $df=0$, for instance if $f=da$ with $a$ a Schwartz one-form.
\item[(b)] $C_0(U_{b,r}f,U_{b,r}h)=C_0(f,h)$ for every Euclidean motion $(b,r)$ and all $f,h\in\mathcal S$.
\item[(c)] For $f,g\in\mathcal S$,
\[
C_0(f,g)=\sum_{k=1}^3\int\mathrm{d}x\int\mathrm{d}y\,\frac{j^f_k(x)\cdot j^g_k(y)}{4\pi|x-y|}.
\]
\item[(d)] Let $j\in\{1,2,3\}$, and let $f\in\mathcal D$ with $\operatorname{supp}f\subset\{x_j>0\}$. Then $C_0(\theta_jf,f)\ge0$.
\item[(e)] $C_0(d^*h,g)=0$ for every Schwartz three-form $h$ and every $g\in\mathcal S$.
\end{enumerate}
\end{lemma}

\begin{proof}
For $u,v\in\mathbb{R}^3$ the two-form $u\wedge v$, with the components $(u\wedge v)_{kl}=u_kv_l-u_lv_k$, is associated with the vector $u\times v$; these two-forms span $\Lambda^2$, and $(ru)\times(rv)=\det(r)\cdot r(u\times v)$ for $r\in O(3)$. As $(\wedge^2r)(u\wedge v)=(ru)\wedge(rv)$ and $(U_{b,r}f)(x)=(\wedge^2r)f\left(r^{-1}(x-b)\right)$, this gives the formula for the vector field of $U_{b,r}f$. The three-form $df$ has the component $(df)_{123}=\partial_1f_{23}-\partial_2f_{13}+\partial_3f_{12}=\operatorname{div}\vec f$, whence $df=0$ if and only if $\operatorname{div}\vec f=0$. For a Schwartz three-form $h$, integration by parts in $\langle d^*h,g\rangle=\int\mathrm{d}x\,h_{123}(x)\cdot(dg)_{123}(x)$, $g\in\mathcal S$, shows that
\[
(d^*h)_{12}=-\partial_3h_{123},\qquad(d^*h)_{13}=\partial_2h_{123},\qquad(d^*h)_{23}=-\partial_1h_{123};
\]
thus $d^*h\in\mathcal S$, and its vector field is $-\operatorname{grad}h_{123}$. A field $v\in L^2(\mathbb{R}^3;\mathbb{R}^3)$ satisfies $\operatorname{div}v=0$ if and only if $p\cdot\hat v(p)=0$ for almost every $p\in\mathbb{R}^3$, and for $p\ne0$ the real symmetric matrix $1-|p|^{-2}pp^T$ is the orthogonal projection of $\mathbb{C}^3$ onto the orthogonal complement of $p$. Hence $\Pi$ multiplies the Fourier transform by this matrix, and Plancherel's theorem,
\[
\int\mathrm{d}x\,\overline{\phi(x)}\cdot\psi(x)=(2\pi)^{-3}\int\mathrm{d}p\,\overline{\hat\phi(p)}\cdot\hat\psi(p),
\]
gives $C_0(f,g)=\langle f,\Pi g\rangle$. Since $\Pi$ is an orthogonal projection, $\langle f,\Pi g\rangle=\langle\Pi f,\Pi g\rangle=\langle g,\Pi f\rangle$; so $C_0$ is a symmetric bilinear form.

(a) Let $f\in\mathcal S$, put $v=\vec f$, and let $\|\cdot\|$ be the norm of $L^2(\mathbb{R}^3;\mathbb{R}^3)$. As $\Pi$ is an orthogonal projection, $C_0(f,f)=\int\mathrm{d}x\,v(x)\cdot(\Pi v)(x)=\|\Pi v\|^2\ge0$ and
\[
\|f\|_2^2=\|v\|^2=\|\Pi v\|^2+\|v-\Pi v\|^2\ge C_0(f,f),
\]
with equality if and only if $\Pi v=v$, that is, if and only if $\operatorname{div}\vec f=0$, or equivalently $df=0$. For a Schwartz one-form $a$ the two-form $da$ lies in $\mathcal S$, and $d(da)=0$.

(b) Let $(b,r)$ be a Euclidean motion and $f\in\mathcal S$, and let $w$ be the vector field of $U_{b,r}f$. The substitution $x=ry+b$, whose Jacobian has modulus $1$, gives
\[
\hat w(p)=\det(r)\cdot\exp(-ip\cdot b)\cdot r\,\widehat{\vec f}\left(r^{-1}p\right)\qquad(p\in\mathbb{R}^3).
\]
For $p\ne0$ and $q=r^{-1}p=r^Tp$ one has $|q|=|p|$ and
\[
r^T\left(1-\frac{pp^T}{|p|^2}\right)r=1-\frac{qq^T}{|q|^2}.
\]
Since $\det(r)^2=1$ and $\left|\exp(-ip\cdot b)\right|=1$, it follows for $h\in\mathcal S$, by the substitution $p=rq$, that
\begin{align*}
C_0(U_{b,r}f,U_{b,r}h)&=(2\pi)^{-3}\int\mathrm{d}p\,\overline{\widehat{\vec f}\left(r^{-1}p\right)}\cdot r^T\left(1-\frac{pp^T}{|p|^2}\right)r\,\widehat{\vec h}\left(r^{-1}p\right)\\
&=(2\pi)^{-3}\int\mathrm{d}q\,\overline{\widehat{\vec f}(q)}\cdot\left(1-\frac{qq^T}{|q|^2}\right)\widehat{\vec h}(q)=C_0(f,h).
\end{align*}

(c) The Fourier transform turns $\partial_k$ into the multiplication by $ip_k$; hence $\widehat{j^f}(p)=ip\times\widehat{\vec f}(p)$. For $p\in\mathbb{R}^3\setminus\{0\}$ and $v,w\in\mathbb{C}^3$ the identity $|p\times v|^2=|p|^2\cdot|v|^2-|p\cdot v|^2$ has the polarised form $\overline{(p\times v)}\cdot(p\times w)=|p|^2\cdot\left(\overline v\cdot w\right)-(p\cdot\overline v)\cdot(p\cdot w)$, that is,
\[
\frac{\overline{(p\times v)}\cdot(p\times w)}{|p|^2}=\overline v\cdot\left(1-\frac{pp^T}{|p|^2}\right)w .
\]
With $v=\widehat{\vec f}(p)$ and $w=\widehat{\vec g}(p)$ it follows that
\[
C_0(f,g)=(2\pi)^{-3}\int\mathrm{d}p\,\frac{\overline{\widehat{j^f}(p)}\cdot\widehat{j^g}(p)}{|p|^2}.
\]
The function $p\mapsto|p|^{-2}$ is locally integrable on $\mathbb{R}^3$, and it is the Fourier transform, in the sense of tempered distributions, of the locally integrable function $x\mapsto(4\pi|x|)^{-1}$. Hence, for real Schwartz functions $\phi$, $\psi$ on $\mathbb{R}^3$,
\[
(2\pi)^{-3}\int\mathrm{d}p\,\frac{\overline{\hat\phi(p)}\cdot\hat\psi(p)}{|p|^2}=\int\mathrm{d}x\int\mathrm{d}y\,\frac{\phi(x)\cdot\psi(y)}{4\pi|x-y|}.
\]
Applied to the components $\phi=j^f_k$ and $\psi=j^g_k$, which are real Schwartz functions, and summed over $k$, this gives (c).

(e) Let $h$ be a Schwartz three-form and $g\in\mathcal S$. The vector field of $d^*h$ is $w=-\operatorname{grad}h_{123}$. For every field $v\in L^2(\mathbb{R}^3;\mathbb{R}^3)$ with $\operatorname{div}v=0$ in the sense of distributions, $\int\mathrm{d}x\,w(x)\cdot v(x)=(\operatorname{div}v)(h_{123})=0$; so $w$ is orthogonal to the range of $\Pi$, and $\Pi w=0$. As $\Pi$ is self-adjoint,
\[
C_0(d^*h,g)=\int\mathrm{d}x\,w(x)\cdot(\Pi\vec g)(x)=\int\mathrm{d}x\,(\Pi w)(x)\cdot\vec g(x)=0 .
\]

(d) First let $j\in\{2,3\}$, and let $r\in O(3)$ be the matrix that interchanges the coordinates $x_1$ and $x_j$; then $r^{-1}=r$ and $r\theta_1r=\theta_j$. Let $f\in\mathcal D$ with $\operatorname{supp}f\subset\{x_j>0\}$, and put $f'=U_{0,r}f$. Then $f'\in\mathcal D$, $\operatorname{supp}f'=r(\operatorname{supp}f)\subset\{x_1>0\}$, $U_{0,r}f'=f$ and $U_{0,r}\theta_1f'=U_{0,r\theta_1r}f=\theta_jf$; hence (b) gives $C_0(\theta_jf,f)=C_0(\theta_1f',f')$. It therefore suffices to prove (d) for $j=1$.

Let $f\in\mathcal D$ with $\operatorname{supp}f\subset\{x_1>0\}$. As $\theta_1=\operatorname{diag}(-1,1,1)$ has determinant $-1$, the vector field of $\theta_1f$ is $x\mapsto\operatorname{diag}(1,-1,-1)\,\vec f(\theta_1x)$. Write $v=\vec f$. Since $\partial_1\left[v_k(\theta_1x)\right]=-(\partial_1v_k)(\theta_1x)$ and $\partial_l\left[v_k(\theta_1x)\right]=(\partial_lv_k)(\theta_1x)$ for $l=2,3$, the curl of this field has at $x$ the components $-(\partial_2v_3-\partial_3v_2)(\theta_1x)$, $(\partial_3v_1-\partial_1v_3)(\theta_1x)$ and $(\partial_1v_2-\partial_2v_1)(\theta_1x)$; that is,
\[
j^{\theta_1f}(x)=\operatorname{diag}(-1,1,1)\,j^f(\theta_1x)\qquad(x\in\mathbb{R}^3).
\]
Put $J=j^f$. Since $\operatorname{supp}f$ is a compact subset of $\{x_1>0\}$, there is $\varepsilon>0$ with $\operatorname{supp}f\subset\{x_1\ge\varepsilon\}$; the field $J$ is smooth, its support is a compact subset of $\{x_1\ge\varepsilon\}$, and $\operatorname{div}J=\operatorname{div}\operatorname{curl}\vec f=0$. For $x\in\mathbb{R}^3$ write $x=(x_1,x_\perp)$ with $x_\perp=(x_2,x_3)\in\mathbb{R}^2$. By (c), applied to $\theta_1f$ and $f$, by the formula for $j^{\theta_1f}$ and by the substitution $x\mapsto\theta_1x$, which preserves the Lebesgue measure,
\[
C_0(\theta_1f,f)=I(J_2)+I(J_3)-I(J_1),
\]
where, for a bounded measurable real function $u$ on $\mathbb{R}^3$ with compact support in $\{x_1\ge\varepsilon\}$,
\[
I(u)=\int\mathrm{d}x\int\mathrm{d}y\,\frac{u(x)\cdot u(y)}{4\pi|\theta_1x-y|},\qquad|\theta_1x-y|^2=(x_1+y_1)^2+|x_\perp-y_\perp|^2 .
\]
Let $J_0$ be the Bessel function of order zero, $J_0(\lambda)=\frac{1}{2\pi}\int_0^{2\pi}\mathrm{d}\alpha\,\exp(i\lambda\cos\alpha)$ for $\lambda\in\mathbb{R}$. For $t>0$ and $z\in\mathbb{R}^2$ one has
\[
\frac{1}{4\pi\sqrt{t^2+|z|^2}}=\int_{\mathbb{R}^2}\mathrm{d}^2k\,(2\pi)^{-2}\cdot\frac{\exp(-|k|t)\cdot\exp(ik\cdot z)}{2|k|}.
\]
Indeed, in the polar coordinates $k=\kappa\cdot(\cos\alpha,\sin\alpha)$, the angle $\alpha$ being measured from the direction of $z$ if $z\ne0$, one has $\mathrm{d}^2k=\kappa\,\mathrm{d}\kappa\,\mathrm{d}\alpha$ and $k\cdot z=\kappa|z|\cos\alpha$, so that the right side equals $\frac{1}{4\pi}\int_0^\infty\mathrm{d}\kappa\,\exp(-\kappa t)\cdot J_0(\kappa|z|)$; and the Lipschitz integral
\[
\int_0^\infty\mathrm{d}\kappa\,\exp(-\kappa t)\cdot J_0(\kappa|z|)=\frac{1}{\sqrt{t^2+|z|^2}}
\]
gives the identity. Moreover
\[
\int_{\mathbb{R}^2}\mathrm{d}^2k\,(2\pi)^{-2}\cdot\frac{\exp(-|k|t)}{2|k|}=\frac{1}{4\pi t},
\]
so that the integral over $k$ converges absolutely, uniformly for $t\ge2\varepsilon$. If $x$ and $y$ lie in the support of $u$, then $x_1+y_1\ge2\varepsilon$. Hence, with $t=x_1+y_1$ and $z=x_\perp-y_\perp$, Fubini's theorem applies to $I(u)$, and since $u$ is real it gives
\[
I(u)=\int_{\mathbb{R}^2}\mathrm{d}^2k\,(2\pi)^{-2}\cdot\frac{|\tilde u(k)|^2}{2|k|},\qquad\tilde u(k)=\int\mathrm{d}y\,u(y)\cdot\exp(-|k|y_1)\cdot\exp(ik\cdot y_\perp),
\]
where $k=(k_2,k_3)\in\mathbb{R}^2$ and $k\cdot y_\perp=k_2y_2+k_3y_3$. These integrals are finite, as $|\tilde u(k)|\le\exp(-\varepsilon|k|)\cdot\|u\|_1$.

For $l=1,2,3$ let $\tilde J_l$ be the function $\tilde u$ formed with $u=J_l$. Let $k\in\mathbb{R}^2$, $k\ne0$. Multiplying the identity $\operatorname{div}J=0$ by $\exp(-|k|y_1)\cdot\exp(ik\cdot y_\perp)$, integrating over $\mathbb{R}^3$ and integrating by parts, which is legitimate because $J$ has compact support, one finds
\[
0=\int\mathrm{d}y\,(\operatorname{div}J)(y)\cdot\exp(-|k|y_1)\cdot\exp(ik\cdot y_\perp)=|k|\cdot\tilde J_1(k)-i\left(k_2\cdot\tilde J_2(k)+k_3\cdot\tilde J_3(k)\right).
\]
Thus $|k|\cdot\tilde J_1(k)=i\left(k_2\cdot\tilde J_2(k)+k_3\cdot\tilde J_3(k)\right)$, and the Cauchy--Schwarz inequality gives $|k|^2\cdot|\tilde J_1(k)|^2\le|k|^2\cdot\left(|\tilde J_2(k)|^2+|\tilde J_3(k)|^2\right)$, that is, $|\tilde J_1(k)|^2\le|\tilde J_2(k)|^2+|\tilde J_3(k)|^2$. Consequently
\[
C_0(\theta_1f,f)=\int_{\mathbb{R}^2}\mathrm{d}^2k\,(2\pi)^{-2}\cdot\frac{|\tilde J_2(k)|^2+|\tilde J_3(k)|^2-|\tilde J_1(k)|^2}{2|k|}\ge0 .
\]
\end{proof}

\begin{theorem}[Properties of the scale mixtures]\label{thm:properties-scale-mixtures}
For $s\in[0,1]$ the form $s\cdot C_0$ is symmetric and positive semidefinite, and $s\cdot C_0(f,f)\le\|f\|_2^2$ for $f\in\mathcal S$, by Lemma~\ref{lem:free-transverse-covariance}(a); let $\gamma_s=\gamma_{s\cdot C_0}$ be the law given by Lemma~\ref{lem:gaussian-laws-bounded}(a). Let $\rho$ be a Borel probability measure on $[0,1]$. For every $A\in\mathcal F$ the function $s\mapsto\gamma_s(A)$ is Borel measurable on $[0,1]$, and
\[
\mu^\rho(A)=\int\mathrm{d}\rho(s)\,\gamma_s(A)\qquad(A\in\mathcal F)
\]
defines a law $\mu^\rho$ on $(\Omega,\mathcal F)$; it is the law of $s^{\frac{1}{2}}F$, where $s$ and $F$ are independent random elements of $[0,1]$ and of $\Omega$ with the laws $\rho$ and $\gamma_1$. Put
\[
E_\rho[s]=\int\mathrm{d}\rho(s)\,s,\qquad\operatorname{Var}_\rho(s)=\int\mathrm{d}\rho(s)\,s^2-\left(E_\rho[s]\right)^2 .
\]
\begin{enumerate}
\item[(i)] The law $\mu^\rho$ satisfies \eqref{eq:fourth-cumulants-H1}, \eqref{eq:fourth-cumulants-H2} and \eqref{eq:fourth-cumulants-H3} in place of $\mu$. It is invariant under the group $\mathbb{R}^3\rtimes O(3)$ of all Euclidean motions: $\mu^\rho\circ T_g^{-1}=\mu^\rho$ for every Euclidean motion $g$. For every Schwartz three-form $h$ one has $X(d^*h)=0$ $\mu^\rho$-almost surely.
\item[(ii)] The probability space $([0,1],B,\rho)$, with $B$ the Borel $\sigma$-algebra of $[0,1]$, together with the family $(\gamma_s)_{s\in[0,1]}$, forms an ergodic decomposition of $\mu^\rho$, and each of its members $\gamma_s$, $s\in[0,1]$, is a centred Gaussian law.
\item[(iii)] By (i) and Lemma~\ref{lem:moment-bounds}(a), (c), the law $\mu^\rho$ has mean zero and finite moments of every order; let $S_2=S_2^{\mu^\rho}$ and $\kappa_4=\kappa_4^{\mu^\rho}$ be the forms of Section~\ref{sec:cumulants-gaussian-laws}. For $f,g,f_1,f_2,f_3,f_4\in\mathcal S$,
\begin{align*}
S_2(f,g)&=E_\rho[s]\cdot C_0(f,g),\\
\kappa_4(f_1,f_2,f_3,f_4)&=\operatorname{Var}_\rho(s)\cdot\left[C_0(f_1,f_2)\cdot C_0(f_3,f_4)+C_0(f_1,f_3)\cdot C_0(f_2,f_4)+C_0(f_1,f_4)\cdot C_0(f_2,f_3)\right];
\end{align*}
in particular $\kappa_4(f)=3\operatorname{Var}_\rho(s)\cdot C_0(f,f)^2$.
\item[(iv)] Let $f_*\in\mathcal S$ with $C_0(f_*,f_*)>0$, let $b_*>0$, and put
\[
q_*=\frac{1}{C_0(f_*,f_*)}\cdot\left(\frac{b_*}{3}\right)^{\frac{1}{2}} .
\]
For $0<e\le(2q_*)^{-\frac{1}{2}}$ let $\rho_e=\frac{1}{2}\left(\delta_1+\delta_{1-2q_*\cdot e^2}\right)$, where $\delta_t$ denotes the point mass at $t$. Then
\[
\kappa_4^{\mu^{\rho_e}}(f_*)=e^4\cdot b_* .
\]
Consequently the law $\mu^{\rho_e}$ satisfies \eqref{eq:fourth-cumulants-4g} in place of $\mu$, at the test $g=f_*$, with $b(f_*)=b_*$, with $C_{f_*}=0$ and with any $\sigma>0$.
\item[(v)] Let $\mathrm{Ext}_x=\mathrm{Ext}^{\mu^\rho}_x$ be the quantity of Section~\ref{sec:cumulants-gaussian-laws}. For $f\in\mathcal S$ and $x\in\mathbb{R}^3$,
\[
\kappa_4(f,f,\tau_xf,\tau_xf)=\operatorname{Var}_\rho(s)\cdot\left[C_0(f,f)^2+2C_0(f,\tau_xf)^2\right]\ge\frac{1}{3}\kappa_4(f),\qquad\lim_{|x|\to\infty}\kappa_4(f,f,\tau_xf,\tau_xf)=\frac{1}{3}\kappa_4(f),
\]
and
\[
\mathrm{Ext}_x(f)=-6\operatorname{Var}_\rho(s)\cdot C_0(f,\tau_xf)^2\le0 .
\]
\end{enumerate}
\end{theorem}

\begin{proof}
For $t\ge0$ and $\omega\in\Omega$ let $t\omega\in\Omega$ be the functional $f\mapsto t\cdot\omega(f)$, and let $B$ be the Borel $\sigma$-algebra of $[0,1]$. The map
\[
\Phi\colon[0,1]\times\Omega\to\Omega,\qquad\Phi(s,\omega)=s^{\frac{1}{2}}\omega,
\]
is measurable with respect to $B\otimes\mathcal F$ and $\mathcal F$: the $\sigma$-algebra $\mathcal F$ is generated by the functions $X(f)$, $f\in\mathcal S$, and $X(f)\circ\Phi$ is the product of the measurable functions $(s,\omega)\mapsto s^{\frac{1}{2}}$ and $(s,\omega)\mapsto X(f)(\omega)$. For $s\in[0,1]$ the image of $\gamma_1$ under $\omega\mapsto s^{\frac{1}{2}}\omega$ has the characteristic functional
\[
f\mapsto\int\mathrm{d}\gamma_1\,\exp\left(is^{\frac{1}{2}}X(f)\right)=\int\mathrm{d}\gamma_1\,\exp\left(iX\left(s^{\frac{1}{2}}f\right)\right)=\exp\left(-\frac{1}{2}s\cdot C_0(f,f)\right),
\]
and is therefore equal to $\gamma_s$, by the uniqueness in Lemma~\ref{lem:gaussian-laws-bounded}(a). So $\gamma_s(A)=\int\mathrm{d}\gamma_1(\omega)\,1_A\left(\Phi(s,\omega)\right)$ for $A\in\mathcal F$, and Tonelli's theorem, applied to the non-negative measurable function $1_A\circ\Phi$ on $[0,1]\times\Omega$, shows that $s\mapsto\gamma_s(A)$ is Borel measurable. By monotone convergence $\mu^\rho$ is countably additive, and $\mu^\rho(\Omega)=1$. Thus $\mu^\rho$ is a law on $(\Omega,\mathcal F)$, and $([0,1],B,\rho)$ together with $(\gamma_s)_{s\in[0,1]}$ satisfies \eqref{eq:fourth-cumulants-ED1} with $\mu^\rho$ in place of $\mu$. For independent $s$ and $F$ with the laws $\rho$ and $\gamma_1$, the law of $s^{\frac{1}{2}}F$ is the image of $\rho\otimes\gamma_1$ under $\Phi$, and Tonelli's theorem gives
\[
(\rho\otimes\gamma_1)\left(\Phi^{-1}(A)\right)=\int\mathrm{d}\rho(s)\int\mathrm{d}\gamma_1(\omega)\,1_A\left(s^{\frac{1}{2}}\omega\right)=\int\mathrm{d}\rho(s)\,\gamma_s(A)=\mu^\rho(A)\qquad(A\in\mathcal F).
\]
By linearity and monotone convergence, \eqref{eq:fourth-cumulants-ED1} extends from indicator functions to measurable functions $G\colon\Omega\to[0,\infty]$: the function $s\mapsto\int\mathrm{d}\gamma_s\,G$ is measurable, and
\[
\int\mathrm{d}\mu^\rho\,G=\int\mathrm{d}\rho(s)\int\mathrm{d}\gamma_s\,G .
\]
Splitting into real and imaginary parts, and these into positive and negative parts, one obtains the same identity for every measurable $G\colon\Omega\to\mathbb{C}$ with $\int\mathrm{d}\mu^\rho\,|G|<\infty$.

(i) Let $f\in\mathcal S$. For $s\in[0,1]$, Lemma~\ref{lem:free-transverse-covariance}(a) gives $0\le s\cdot C_0(f,f)\le\|f\|_2^2$; so, by Lemma~\ref{lem:gaussian-laws-bounded}(a),
\[
\int\mathrm{d}\mu^\rho\,\exp X(f)=\int\mathrm{d}\rho(s)\,\exp\left(\frac{1}{2}s\cdot C_0(f,f)\right)\le\exp\left(\frac{1}{2}\|f\|_2^2\right),
\]
which is \eqref{eq:fourth-cumulants-H1}. Let $g$ be a Euclidean motion. By Lemma~\ref{lem:free-transverse-covariance}(b), $s\cdot C_0(U_gf,U_gh)=s\cdot C_0(f,h)$ for all $f,h\in\mathcal S$ and $s\in[0,1]$; Lemma~\ref{lem:gaussian-laws-bounded}(b), applied to the group of all Euclidean motions, gives $\gamma_s\circ T_g^{-1}=\gamma_s$, and hence
\[
\mu^\rho\left(T_g^{-1}A\right)=\int\mathrm{d}\rho(s)\,\gamma_s\left(T_g^{-1}A\right)=\int\mathrm{d}\rho(s)\,\gamma_s(A)=\mu^\rho(A)\qquad(A\in\mathcal F).
\]
In particular $\mu^\rho$ satisfies \eqref{eq:fourth-cumulants-H2}. Let $j\in\{1,2,3\}$ and $s\in[0,1]$. The form $s\cdot C_0$ satisfies $s\cdot C_0(\theta_jf,\theta_jh)=s\cdot C_0(f,h)$ for $f,h\in\mathcal S$, by Lemma~\ref{lem:free-transverse-covariance}(b) applied to the motion $(0,\theta_j)$, and $s\cdot C_0(\theta_jf,f)\ge0$ for every $f\in\mathcal D$ with $\operatorname{supp}f\subset\{x_j>0\}$, by Lemma~\ref{lem:free-transverse-covariance}(d). So Lemma~\ref{lem:gaussian-laws-bounded}(d) gives $\int\mathrm{d}\gamma_s\,(\Theta_jG)\cdot G\ge0$ for $G\in A_j^+$. For such $G$ the function $(\Theta_jG)\cdot G$ is bounded and measurable, and the extended form of \eqref{eq:fourth-cumulants-ED1} gives
\[
\int\mathrm{d}\mu^\rho\,(\Theta_jG)\cdot G=\int\mathrm{d}\rho(s)\int\mathrm{d}\gamma_s\,(\Theta_jG)\cdot G\ge0,
\]
which is \eqref{eq:fourth-cumulants-H3}. Finally, let $h$ be a Schwartz three-form. Under $\gamma_s$ the variable $X(d^*h)$ is normal with mean $0$ and variance $s\cdot C_0(d^*h,d^*h)$, which vanishes by Lemma~\ref{lem:free-transverse-covariance}(e). Hence
\[
\int\mathrm{d}\mu^\rho\,X(d^*h)^2=\int\mathrm{d}\rho(s)\,s\cdot C_0(d^*h,d^*h)=0,
\]
and $X(d^*h)=0$ $\mu^\rho$-almost surely.

(ii) Property \eqref{eq:fourth-cumulants-ED1} has been established above. Let $s\in(0,1]$. By Lemma~\ref{lem:free-transverse-covariance}(b) the form $s\cdot C_0$ is translation invariant, and Lemma~\ref{lem:gaussian-laws-bounded}(c) shows that $\gamma_s$ is translation invariant and ergodic. The characteristic functional of $\gamma_0$ is identically $1$; so is that of the point mass at the zero element of $\Omega$, and by the uniqueness in Lemma~\ref{lem:gaussian-laws-bounded}(a) the law $\gamma_0$ is this point mass. It is translation invariant, since every $T_x$ maps the zero element to itself, and every element of $L^2(\gamma_0)$ is a constant, namely its value at the zero element; so $L^2(\gamma_0)$ is one-dimensional and $\gamma_0$ is ergodic. Hence \eqref{eq:fourth-cumulants-ED2} holds, even for every $s\in[0,1]$. By Lemma~\ref{lem:gaussian-laws-bounded}(a), every $\gamma_s$ is a centred Gaussian law.

(iii) By (i), Lemma~\ref{lem:moment-bounds}(c) and H\"older's inequality, every product of finitely many of the variables $X(f)$, $f\in\mathcal S$, is $\mu^\rho$-integrable. Let $f_1,\dots,f_4\in\mathcal S$ and put
\[
Q=C_0(f_1,f_2)\cdot C_0(f_3,f_4)+C_0(f_1,f_3)\cdot C_0(f_2,f_4)+C_0(f_1,f_4)\cdot C_0(f_2,f_3).
\]
For $s\in[0,1]$ the variables $X(f_1),\dots,X(f_4)$ are, under $\gamma_s$, jointly normal and centred (Section~\ref{sec:cumulants-gaussian-laws}), with the covariances $s\cdot C_0(f_a,f_{a'})$, $a,a'\in\{1,2,3,4\}$ (Lemma~\ref{lem:gaussian-laws-bounded}(a)). Wick's formula gives
\[
\int\mathrm{d}\gamma_s\,X(f_1)\cdot X(f_2)\cdot X(f_3)\cdot X(f_4)=s^2\cdot Q .
\]
Integration with respect to $\rho$, by the extended form of \eqref{eq:fourth-cumulants-ED1}, yields
\[
S_2(f,g)=\int\mathrm{d}\rho(s)\,s\cdot C_0(f,g)=E_\rho[s]\cdot C_0(f,g)\qquad(f,g\in\mathcal S)
\]
and
\[
\int\mathrm{d}\mu^\rho\,X(f_1)\cdot X(f_2)\cdot X(f_3)\cdot X(f_4)=\int\mathrm{d}\rho(s)\,s^2\cdot Q .
\]
By the first of these formulas, the three products subtracted in the definition of $\kappa_4(f_1,f_2,f_3,f_4)$ in Section~\ref{sec:cumulants-gaussian-laws} add up to $\left(E_\rho[s]\right)^2\cdot Q$; therefore
\[
\kappa_4(f_1,f_2,f_3,f_4)=\left(\int\mathrm{d}\rho(s)\,s^2-\left(E_\rho[s]\right)^2\right)\cdot Q=\operatorname{Var}_\rho(s)\cdot Q .
\]
For $f_1=f_2=f_3=f_4=f$ one has $Q=3C_0(f,f)^2$.

(iv) The number $q_*$ is positive. Since $0<e\le(2q_*)^{-\frac{1}{2}}$, the number $1-2q_*\cdot e^2$ lies in $[0,1)$, and $\rho_e$ is a probability measure on $[0,1]$. Its mean is $E_{\rho_e}[s]=1-q_*\cdot e^2$, and each of the two points $1$ and $1-2q_*\cdot e^2$ lies at the distance $q_*\cdot e^2$ from it; hence $\operatorname{Var}_{\rho_e}(s)=\left(q_*\cdot e^2\right)^2$. By (iii) and the definition of $q_*$,
\[
\kappa_4^{\mu^{\rho_e}}(f_*)=3\left(q_*\cdot e^2\right)^2\cdot C_0(f_*,f_*)^2=3\cdot\frac{b_*}{3}\cdot e^4=e^4\cdot b_* .
\]
So $\left|\kappa_4^{\mu^{\rho_e}}(f_*)-e^4\cdot b_*\right|=0$, which is \eqref{eq:fourth-cumulants-4g} for the law $\mu^{\rho_e}$ at the test $g=f_*$, with $b(f_*)=b_*$, $C_{f_*}=0$ and any $\sigma>0$.

(v) Let $f\in\mathcal S$ and $x\in\mathbb{R}^3$. By Lemma~\ref{lem:free-transverse-covariance}(b), $C_0(\tau_xf,\tau_xf)=C_0(f,f)$; so (iii), applied to $(f,f,\tau_xf,\tau_xf)$, gives
\[
\kappa_4(f,f,\tau_xf,\tau_xf)=\operatorname{Var}_\rho(s)\cdot\left[C_0(f,f)^2+2C_0(f,\tau_xf)^2\right]\ge\operatorname{Var}_\rho(s)\cdot C_0(f,f)^2=\frac{1}{3}\kappa_4(f).
\]
By Lemma~\ref{lem:gaussian-laws-bounded}(a), the law $\gamma_1$ satisfies \eqref{eq:fourth-cumulants-H1} in place of $\mu$ and has the two-point function $S_2^{\gamma_1}=C_0$; Lemma~\ref{lem:decay-two-point}(b), applied to $\gamma_1$, gives $C_0(f,\tau_xf)\to0$ as $|x|\to\infty$, and hence $\kappa_4(f,f,\tau_xf,\tau_xf)\to\operatorname{Var}_\rho(s)\cdot C_0(f,f)^2=\frac{1}{3}\kappa_4(f)$. Finally, by the definition of $\mathrm{Ext}_x$ in Section~\ref{sec:cumulants-gaussian-laws},
\begin{align*}
\mathrm{Ext}_x(f)&=\kappa_4(f)-3\kappa_4(f,f,\tau_xf,\tau_xf)\\
&=3\operatorname{Var}_\rho(s)\cdot C_0(f,f)^2-3\operatorname{Var}_\rho(s)\cdot\left[C_0(f,f)^2+2C_0(f,\tau_xf)^2\right]=-6\operatorname{Var}_\rho(s)\cdot C_0(f,\tau_xf)^2,
\end{align*}
which is non-positive, because $\operatorname{Var}_\rho(s)\ge0$ by the Cauchy--Schwarz inequality.
\end{proof}

\begin{remark}[Mixtures of free fields]\label{rem:mixtures-free-fields}
The notation is that of Theorem~\ref{thm:properties-scale-mixtures}: $\gamma_s=\gamma_{s\cdot C_0}$ for $s\in[0,1]$, and $\mu^{\rho}=\int\mathrm{d}\rho(s)\,\gamma_s$ for a probability measure $\rho$ on $[0,1]$; further, $E_{\rho}[s]$, $E_{\rho}\left[s^2\right]$ and $\operatorname{Var}_{\rho}(s)$ are the first moment, the second moment and the variance of the variable $s$ under $\rho$. The law $\gamma_1=\gamma_{C_0}$, that is, the centred Gaussian law whose covariance is the form $C_0$ of Lemma~\ref{lem:free-transverse-covariance}, is called the free field. Throughout, $f_*\in\mathcal S$ and $b_*$ are as in Theorem~\ref{thm:properties-scale-mixtures}(iv), so that $C_0(f_*,f_*)>0$ and $b_*>0$, and, as there,
\[
q_*=\frac{1}{C_0(f_*,f_*)}\cdot\left(\frac{b_*}{3}\right)^{\frac{1}{2}},\qquad\rho_e=\frac{1}{2}\left(\delta_1+\delta_{1-2q_*\cdot e^2}\right)\qquad\left(0<e\le(2q_*)^{-\frac{1}{2}}\right).
\]
For these $e$ one has $0\le1-2q_*\cdot e^2<1$ and $\mu^{\rho_e}=\frac{1}{2}\left(\gamma_1+\gamma_{1-2q_*\cdot e^2}\right)$. The two atoms of $\rho_e$ lie at the distance $q_*\cdot e^2$ from the point $1-q_*\cdot e^2$, whence
\[
E_{\rho_e}[s]=1-q_*\cdot e^2,\qquad\operatorname{Var}_{\rho_e}(s)=q_*^2\cdot e^4,\qquad E_{\rho_e}\left[s^2\right]=\operatorname{Var}_{\rho_e}(s)+E_{\rho_e}[s]^2.
\]
\begin{enumerate}
\item[(1)] A test $f\in\mathcal S$ is called closed if $df=0$. For $a\in C_c^\infty(\mathbb{R}^3;\mathbb{R}^3)$ the test $da$ is closed, because $d\circ d=0$. By Lemma~\ref{lem:free-transverse-covariance}(a), every closed test $f$ satisfies
\[
C_0(f,f)=\|f\|_2^2 .
\]
Let $f=dd^*\beta$ be one of the tests of Section~\ref{sec:free-quartic-coefficient}. It equals $da$ with $a=d^*\beta\in C_c^\infty(\mathbb{R}^3;\mathbb{R}^3)$, and it is therefore closed. Assume that $b(f)>0$. Then $f\ne0$. Indeed, if $f=0$, the smooth vector field $a$ has compact support and satisfies $da=f=0$ and $d^*a=d^*d^*\beta=0$; as in Section~\ref{sec:free-quartic-coefficient}, $a$ is then harmonic and vanishes. The operator $T_a$ of that section depends linearly on $a$; hence $T_a=0$, and $b(f)=\mathsf{b}(a)=-6\operatorname{Tr}\left[T_a^4\right]=0$. Consequently $C_0(f,f)=\|f\|_2^2>0$, and the test $f_*=f$ together with the number $b_*=b(f)$ satisfies the hypotheses of Theorem~\ref{thm:properties-scale-mixtures}(iv). For this choice and for $0<e\le(2q_*)^{-\frac{1}{2}}$, the law $\mu^{\rho_e}$ satisfies \eqref{eq:fourth-cumulants-4g} at the test $g=f_*$, with the number $b(f_*)$ of Section~\ref{sec:free-quartic-coefficient}, with $C_{f_*}=0$ and with every $\sigma>0$.

For all $f_*$, $b_*$ as above and every $e\in\left(0,(2q_*)^{-\frac{1}{2}}\right]$, Theorem~\ref{thm:properties-scale-mixtures}(iii) gives
\[
S_2^{\mu^{\rho_e}}(f,g)=E_{\rho_e}[s]\cdot C_0(f,g)=\left(1-q_*\cdot e^2\right)\cdot C_0(f,g)\qquad(f,g\in\mathcal S).
\]
Moreover, $\mu^{\rho_e}$ converges to the free field $\gamma_1$ as $e\to0$, in the following sense. By Lemma~\ref{lem:gaussian-laws-bounded}(a), $\int\mathrm{d}\gamma_s\,\exp\left(iX(f)\right)=\exp\left(-\frac{s}{2}\cdot C_0(f,f)\right)$ for $s\in[0,1]$ and $f\in\mathcal S$. Therefore, for every $f\in\mathcal S$,
\begin{align*}
\int\mathrm{d}\mu^{\rho_e}\,\exp\left(iX(f)\right)&=\frac{1}{2}\exp\left(-\frac{1}{2}C_0(f,f)\right)+\frac{1}{2}\exp\left(-\frac{1-2q_*\cdot e^2}{2}\cdot C_0(f,f)\right)\\
&\longrightarrow\exp\left(-\frac{1}{2}C_0(f,f)\right)=\int\mathrm{d}\gamma_1\,\exp\left(iX(f)\right)\qquad(e\to0).
\end{align*}
Let $f_1,\dots,f_k\in\mathcal S$. Since $\sum_lt_l\cdot X(f_l)=X\left(\sum_lt_l\cdot f_l\right)$ for $t\in\mathbb{R}^k$, the characteristic function of the law of $\left(X(f_1),\dots,X(f_k)\right)$ under $\mu^{\rho_e}$ converges pointwise on $\mathbb{R}^k$ to the characteristic function of its law under $\gamma_1$. By L\'evy's continuity theorem, the first of these laws converges weakly to the second.
\item[(2)] Let $0<e\le(2q_*)^{-\frac{1}{2}}$, put $\mu'=\mu^{\rho_e}$, and let $E'$, $S_2'$ and $\kappa_4'$ be the expectation, the two-point function and the fourth cumulant of $\mu'$, in the notation of Section~\ref{sec:cumulants-gaussian-laws}. By Theorem~\ref{thm:properties-scale-mixtures}(i) the law $\mu'$ satisfies \eqref{eq:fourth-cumulants-H1}; hence it is centred and has finite fourth moments, by Lemma~\ref{lem:moment-bounds}(a),(c).

The law $\mu'$ is not Gaussian. Indeed, the fourth cumulant of a centred Gaussian law vanishes (Section~\ref{sec:cumulants-gaussian-laws}), whereas $\kappa_4'(f_*)=e^4\cdot b_*>0$ by Theorem~\ref{thm:properties-scale-mixtures}(iv).

The law $\mu'$ does not satisfy the clustering condition~\eqref{eq:fourth-cumulants-C1} of Theorem~\ref{thm:clustering-mixing-ergodicity}. To see this, let $f_1,f_2,f_3,f_4\in\mathcal S$. The definition of $\kappa_4'$, the formulas of Theorem~\ref{thm:properties-scale-mixtures}(iii) for $S_2'$ and $\kappa_4'$, and the identity $\operatorname{Var}_{\rho_e}(s)+E_{\rho_e}[s]^2=E_{\rho_e}\left[s^2\right]$ give
\begin{align*}
E'\left[X(f_1)\cdot X(f_2)\cdot X(f_3)\cdot X(f_4)\right]&=\kappa_4'(f_1,f_2,f_3,f_4)+S_2'(f_1,f_2)\cdot S_2'(f_3,f_4)\\
&\qquad+S_2'(f_1,f_3)\cdot S_2'(f_2,f_4)+S_2'(f_1,f_4)\cdot S_2'(f_2,f_3)\\
&=E_{\rho_e}\left[s^2\right]\cdot\left[C_0(f_1,f_2)\cdot C_0(f_3,f_4)+C_0(f_1,f_3)\cdot C_0(f_2,f_4)\right.\\
&\qquad\left.+C_0(f_1,f_4)\cdot C_0(f_2,f_3)\right].
\end{align*}
Let $f\in\mathcal S$ and $x\in\mathbb{R}^3$. The form $C_0$ is invariant under translations (Lemma~\ref{lem:free-transverse-covariance}(b)); the choice $(f_1,f_2,f_3,f_4)=(f,f,\tau_xf,\tau_xf)$ therefore yields
\[
E'\left[X(f)^2\cdot X(\tau_xf)^2\right]=E_{\rho_e}\left[s^2\right]\cdot\left[C_0(f,f)^2+2\,C_0(f,\tau_xf)^2\right].
\]
By Lemma~\ref{lem:gaussian-laws-bounded}(a) and Lemma~\ref{lem:free-transverse-covariance}(a), the free field $\gamma_1$ satisfies \eqref{eq:fourth-cumulants-H1}, and its two-point function is $S_2^{\gamma_1}=C_0$. Lemma~\ref{lem:decay-two-point}(b), applied to $\gamma_1$, gives $C_0(f,\tau_xf)\to0$ as $|x|\to\infty$. It follows that
\[
\lim_{|x|\to\infty}E'\left[X(f)^2\cdot X(\tau_xf)^2\right]=E_{\rho_e}\left[s^2\right]\cdot C_0(f,f)^2,\qquad\left(E'\left[X(f)^2\right]\right)^2=S_2'(f,f)^2=E_{\rho_e}[s]^2\cdot C_0(f,f)^2,
\]
and the first of these two numbers exceeds the second by $\operatorname{Var}_{\rho_e}(s)\cdot C_0(f,f)^2=q_*^2\cdot e^4\cdot C_0(f,f)^2$, which is positive whenever $C_0(f,f)>0$. Now let $f\in\mathcal D$ be a closed test with $f\ne0$, for instance $f=da$ with $a\in C_c^\infty(\mathbb{R}^3;\mathbb{R}^3)$ and $da\ne0$. By (1) above, $C_0(f,f)=\|f\|_2^2>0$. Condition~\eqref{eq:fourth-cumulants-C1} for $\mu'$, with $m=n=2$ and $f_1=f_2=g_1=g_2=f$, would require $E'\left[X(f)^2\cdot X(\tau_xf)^2\right]\to\left(E'\left[X(f)^2\right]\right)^2$ as $|x|\to\infty$; so it fails.

By Theorem~\ref{thm:properties-scale-mixtures}(ii), $([0,1],\rho_e,(\gamma_s))$ is an ergodic decomposition of $\mu'$, and $\rho_e$ is concentrated on the two points $1$ and $1-2q_*\cdot e^2$. For $s\in[0,1]$ the law $\gamma_s$ is the image of the free field $\gamma_1$ under the map $\omega\mapsto s^{\frac{1}{2}}\omega$ of $\Omega$. Indeed, the image law has the characteristic functional
\[
f\longmapsto\int\mathrm{d}\gamma_1\,\exp\left(is^{\frac{1}{2}}X(f)\right)=\int\mathrm{d}\gamma_1\,\exp\left(iX\left(s^{\frac{1}{2}}f\right)\right)=\exp\left(-\frac{s}{2}\cdot C_0(f,f)\right),
\]
which is that of $\gamma_s$, and a law with this characteristic functional is unique by Lemma~\ref{lem:gaussian-laws-bounded}(a). Thus $\rho_e$-almost every ergodic component of $\mu'$ is the image of the free field $\gamma_1$ under $\omega\mapsto s^{\frac{1}{2}}\omega$, with $s=1$ or $s=1-2q_*\cdot e^2$.
\item[(3)] Let $f_*$ and $b_*$ be as above (for instance, by (1), a test $f_*$ of Section~\ref{sec:free-quartic-coefficient} with $b(f_*)>0$, and $b_*=b(f_*)$), and let $0<e\le(2q_*)^{-\frac{1}{2}}$. By Theorem~\ref{thm:properties-scale-mixtures}(i), (ii) and (iv), the law $\mu^{\rho_e}$ satisfies \eqref{eq:fourth-cumulants-H1}, \eqref{eq:fourth-cumulants-H2} and \eqref{eq:fourth-cumulants-H3}, and it satisfies \eqref{eq:fourth-cumulants-4g} at the single test $g=f_*$, with $b(f_*)=b_*$, $C_{f_*}=0$ and every $\sigma>0$; at the same time every component $\gamma_s$, $s\in[0,1]$, of its ergodic decomposition $([0,1],\rho_e,(\gamma_s))$ is a centred Gaussian law. Hence \eqref{eq:fourth-cumulants-H1}--\eqref{eq:fourth-cumulants-H3}, together with \eqref{eq:fourth-cumulants-4g} at one test, do not imply that every ergodic decomposition of the law has a component which is not Gaussian.

Now let $\rho$ be an arbitrary probability measure on $[0,1]$. Theorem~\ref{thm:properties-scale-mixtures}(v) gives
\[
\mathrm{Ext}^{\mu^{\rho}}_x(f)=-6\operatorname{Var}_{\rho}(s)\cdot C_0(f,\tau_xf)^2\le0\qquad(f\in\mathcal S,\ x\in\mathbb{R}^3).
\]
Since $\frac{1}{2}e^4\cdot b(f)>0$ whenever $e>0$ and $b(f)>0$, the inequality of Theorem~\ref{thm:non-gaussian-components-2}(i) fails for $\mu=\mu^{\rho}$, for every $e>0$, every test $f$ of Section~\ref{sec:free-quartic-coefficient} with $b(f)>0$ and every $x_0\in\mathbb{R}^3$. Consequently the following holds. Let $\beta$, $r$, $a=d^*\beta$ and $f=dd^*\beta$ be as in Theorem~\ref{thm:non-gaussian-components-2}, with $b(f)>0$; let $x_0\in\mathbb{R}^3$ satisfy $d_0=|x_0|-2r>0$ and \eqref{eq:fourth-cumulants-11}; and let $e>0$ and $\sigma>0$. Then there are no constants for which $\mu^{\rho}$ satisfies \eqref{eq:fourth-cumulants-4g} at each of the three tests $g=f$, $g=f+\tau_{x_0}f$, $g=f-\tau_{x_0}f$, with the numbers $b(g)$ of Section~\ref{sec:free-quartic-coefficient}, and for which the number $C_*$ of Theorem~\ref{thm:non-gaussian-components-2} satisfies $C_*\cdot e^{\sigma}\le\frac{b(f)}{4}$. Indeed, $\mu^{\rho}$ is translation invariant and satisfies \eqref{eq:fourth-cumulants-H1} by Theorem~\ref{thm:properties-scale-mixtures}(i), and $([0,1],\rho,(\gamma_s))$ is an ergodic decomposition of $\mu^{\rho}$ by Theorem~\ref{thm:properties-scale-mixtures}(ii). For constants of that kind all hypotheses of Theorem~\ref{thm:non-gaussian-components-2} would be satisfied with $\mu=\mu^{\rho}$, and part~(i) of that theorem would give $\mathrm{Ext}^{\mu^{\rho}}_{x_0}(f)\ge\frac{1}{2}e^4\cdot b(f)>0$, in contradiction with the preceding display.
\end{enumerate}
\end{remark}

\subsection{Cumulant identities and non-Gaussian components}\label{sec:cumulant-identities-non}

Throughout this subsection $\mu$ is a translation-invariant law on $(\Omega,\mathcal F)$ which satisfies~\eqref{eq:fourth-cumulants-H1}; the hypotheses~\eqref{eq:fourth-cumulants-H2} and~\eqref{eq:fourth-cumulants-H3} are made only where they are stated. Further, $(\Xi,\nu,(\mu_{\xi}))$ is an arbitrary ergodic decomposition of $\mu$ in the sense of Definition~\ref{def:ergodicity-ergodic-decompositions}, and $\Xi_0\subset\Xi$ is the set of full $\nu$-measure given by Corollary~\ref{cor:properties-almost-every}. Functions of $\xi$ are considered on $\Xi_0$ only, and the letter $\xi$ always denotes an element of $\Xi_0$.

\begin{lemma}[Large values of a variable bounded above]\label{lem:large-values-variable}
Let $P$ be a probability measure on a measurable space, with expectation $E$, and let $Z$ be a measurable function on that space with values in $\mathbb{R}$. Let $M\in\mathbb{R}$ and $\delta>0$, and assume that $Z\le M$ holds $P$-almost surely and that $E[Z]\ge\delta$. Then
\[
P\left(Z\ge\frac{\delta}{2}\right)\ge\frac{\delta}{2M-\delta}\ge\frac{\delta}{2M}.
\]
\end{lemma}

\begin{proof}
Put $p:=P\left(Z\ge\frac{\delta}{2}\right)$, so that $P\left(Z<\frac{\delta}{2}\right)=1-p$. Split $E[Z]$ into the contributions of the event $\left\{Z<\frac{\delta}{2}\right\}$, on which $Z<\frac{\delta}{2}$, and of its complement, on which $Z\le M$ almost surely. This gives
\[
\delta\le E[Z]\le\frac{\delta}{2}\cdot P\left(Z<\frac{\delta}{2}\right)+M\cdot p\le\frac{\delta}{2}+\left(M-\frac{\delta}{2}\right)\cdot p.
\]
Since $M\ge E[Z]\ge\delta$, one has $M-\frac{\delta}{2}\ge\frac{\delta}{2}>0$ and $2M-\delta\ge\delta>0$. The display yields $\frac{\delta}{2}\le\left(M-\frac{\delta}{2}\right)\cdot p$, that is $\delta\le(2M-\delta)\cdot p$, and this is the first inequality of the lemma. The second one follows from $0<2M-\delta\le2M$.
\end{proof}

\begin{theorem}[The fourth cumulant of a mixture]\label{thm:fourth-cumulant-mixture}
For $f_1,f_2,f_3,f_4\in\mathcal S$ and $f\in\mathcal S$ define the real numbers
\[
\Gamma(f_1,f_2;f_3,f_4):=\operatorname{Cov}_{\nu}\left(C_{\xi}(f_1,f_2),C_{\xi}(f_3,f_4)\right),\qquad V(f):=\Gamma(f,f;f,f)=\operatorname{Var}_{\nu}\left(C_{\xi}(f,f)\right),
\]
where $\operatorname{Cov}_{\nu}$ and $\operatorname{Var}_{\nu}$ denote the covariance and the variance of functions of $\xi$ with respect to $\nu$.
\begin{enumerate}
\item[(a)] For all $f_1,f_2,f_3,f_4\in\mathcal S$,
\begin{multline*}
\kappa_4^{\mu}(f_1,f_2,f_3,f_4)=\int\mathrm{d}\nu(\xi)\,\kappa_4^{\xi}(f_1,f_2,f_3,f_4)\\
+\Gamma(f_1,f_2;f_3,f_4)+\Gamma(f_1,f_3;f_2,f_4)+\Gamma(f_1,f_4;f_2,f_3).
\end{multline*}
\item[(b)] For all $f_1,f_2,f_3,f_4\in\mathcal S$ and every $x\in\mathbb{R}^3$,
\[
\Gamma(f_1,f_2;\tau_xf_3,\tau_xf_4)=\Gamma(f_1,f_2;f_3,f_4).
\]
\item[(c)] The following three identities hold for all $f\in\mathcal S$ and $x\in\mathbb{R}^3$:
\begin{align}
\kappa_4^{\mu}(f)&=\int\mathrm{d}\nu(\xi)\,\kappa_4^{\xi}(f)+3V(f),\label{eq:fourth-cumulants-4}\\
\kappa_4^{\mu}(f,f,\tau_xf,\tau_xf)&=\int\mathrm{d}\nu(\xi)\,\kappa_4^{\xi}(f,f,\tau_xf,\tau_xf)\notag\\
&\qquad+V(f)+2\operatorname{Var}_{\nu}\left(C_{\xi}(f,\tau_xf)\right),\label{eq:fourth-cumulants-5}\\
\int\mathrm{d}\nu(\xi)\,\mathrm{Ext}^{\xi}_x(f)&=\mathrm{Ext}^{\mu}_x(f)+6\operatorname{Var}_{\nu}\left(C_{\xi}(f,\tau_xf)\right).\label{eq:fourth-cumulants-6}
\end{align}
\item[(d)] Let $(\Xi,B,\nu)$ be a probability space and $(\mu_{\xi})_{\xi\in\Xi}$ a family of laws on $(\Omega,\mathcal F)$ which satisfies~\eqref{eq:fourth-cumulants-ED1} and in which every $\mu_{\xi}$ is translation invariant and satisfies~\eqref{eq:fourth-cumulants-H1}. Let $C_{\xi}=S_2^{\mu_{\xi}}$, $\kappa_4^{\xi}=\kappa_4^{\mu_{\xi}}$ and $\mathrm{Ext}^{\xi}_x=\mathrm{Ext}^{\mu_{\xi}}_x$, and let $\Gamma$ and $V$ be defined by the formulas above. Then (a), (b), \eqref{eq:fourth-cumulants-4}, \eqref{eq:fourth-cumulants-5} and \eqref{eq:fourth-cumulants-6} hold for this family, with $\Xi$ in place of $\Xi_0$.
\end{enumerate}
\end{theorem}

\begin{proof}
(a) Let $f_1,f_2,f_3,f_4\in\mathcal S$, and put $Y:=X(f_1)\cdot X(f_2)\cdot X(f_3)\cdot X(f_4)$. The law $\mu$ satisfies~\eqref{eq:fourth-cumulants-H1} by hypothesis, and so does $\mu_{\xi}$ for every $\xi\in\Xi_0$, by Corollary~\ref{cor:properties-almost-every}. By Lemma~\ref{lem:moment-bounds}(a) and~(c), applied to each of these laws, the variable $X(g)$ has mean zero and a finite fourth moment for every $g\in\mathcal S$. These are the two properties that Section~\ref{sec:cumulants-gaussian-laws} requires, so that its definitions apply to $\mu$ and to every $\mu_{\xi}$. By Lemma~\ref{lem:moment-bounds}(e) the functions $\xi\mapsto E_{\xi}[Y]$ and $\xi\mapsto C_{\xi}(f_a,f_b)$, $1\le a<b\le4$, are bounded on $\Xi_0$. The identity~\eqref{eq:fourth-cumulants-ED1} extends from indicator functions to nonnegative measurable functions on $\Omega$, and then to every function $G$ with $E_{\mu}\left[|G|\right]<\infty$, where $E_{\mu}$ is the expectation with respect to $\mu$: for such $G$ the function $\xi\mapsto E_{\xi}[G]$ is defined $\nu$-almost everywhere and measurable, and $E_{\mu}[G]=\int\mathrm{d}\nu(\xi)\,E_{\xi}[G]$. The variables $Y$ and $X(f_a)\cdot X(f_b)$ are $\mu$-integrable by Lemma~\ref{lem:moment-bounds}(c) and H\"older's inequality. Hence
\[
E_{\mu}[Y]=\int\mathrm{d}\nu(\xi)\,E_{\xi}[Y],\qquad S_2^{\mu}(f_a,f_b)=\int\mathrm{d}\nu(\xi)\,C_{\xi}(f_a,f_b)\qquad(1\le a<b\le4).
\]
Let $Q$ be the set of the three quadruples $(1,2,3,4)$, $(1,3,2,4)$ and $(1,4,2,3)$; they correspond to the three partitions of $\{1,2,3,4\}$ into two pairs $\{a,b\}$, $\{c,d\}$. The definition of the fourth cumulant, applied to $\mu_{\xi}$ and to $\mu$, reads
\begin{align*}
E_{\xi}[Y]&=\kappa_4^{\xi}(f_1,f_2,f_3,f_4)+\sum_{(a,b,c,d)\in Q}C_{\xi}(f_a,f_b)\cdot C_{\xi}(f_c,f_d),\\
\kappa_4^{\mu}(f_1,f_2,f_3,f_4)&=E_{\mu}[Y]-\sum_{(a,b,c,d)\in Q}S_2^{\mu}(f_a,f_b)\cdot S_2^{\mu}(f_c,f_d).
\end{align*}
Integrating the first identity against $\nu$ and inserting the result, together with the integral formulas for $S_2^{\mu}$, into the second identity, one obtains
\begin{multline*}
\kappa_4^{\mu}(f_1,f_2,f_3,f_4)=\int\mathrm{d}\nu(\xi)\,\left(\kappa_4^{\xi}(f_1,f_2,f_3,f_4)+\sum_{(a,b,c,d)\in Q}C_{\xi}(f_a,f_b)\cdot C_{\xi}(f_c,f_d)\right)\\
-\sum_{(a,b,c,d)\in Q}\left(\int\mathrm{d}\nu(\xi)\,C_{\xi}(f_a,f_b)\right)\cdot\left(\int\mathrm{d}\nu(\xi)\,C_{\xi}(f_c,f_d)\right).
\end{multline*}
All functions of $\xi$ in this formula are bounded on $\Xi_0$, so that the integral of the sum is the sum of the integrals. For $(a,b,c,d)\in Q$ the definition of $\Gamma$ gives
\[
\int\mathrm{d}\nu(\xi)\,C_{\xi}(f_a,f_b)\cdot C_{\xi}(f_c,f_d)-\left(\int\mathrm{d}\nu(\xi)\,C_{\xi}(f_a,f_b)\right)\cdot\left(\int\mathrm{d}\nu(\xi)\,C_{\xi}(f_c,f_d)\right)=\Gamma(f_a,f_b;f_c,f_d),
\]
and the quadruples $(1,2,3,4)$, $(1,3,2,4)$, $(1,4,2,3)$ yield the three terms $\Gamma$ in (a).

(b) Let $x\in\mathbb{R}^3$. For $g\in\mathcal S$ one has $X(\tau_xg)=\mathbb{U}_xX(g)=X(g)\circ T_x^{-1}$ (Section~\ref{sec:test-forms-distributions-2}), and for $\xi\in\Xi_0$ the law $\mu_{\xi}$ is translation invariant. Therefore
\[
C_{\xi}(\tau_xf_3,\tau_xf_4)=E_{\xi}\left[\left(X(f_3)\cdot X(f_4)\right)\circ T_x^{-1}\right]=E_{\xi}\left[X(f_3)\cdot X(f_4)\right]=C_{\xi}(f_3,f_4)\qquad(\xi\in\Xi_0).
\]
Thus the functions $\xi\mapsto C_{\xi}(\tau_xf_3,\tau_xf_4)$ and $\xi\mapsto C_{\xi}(f_3,f_4)$ coincide on $\Xi_0$, and (b) follows from the definition of $\Gamma$.

(c) For $f_1=f_2=f_3=f_4=f$ each of the three terms $\Gamma$ in (a) equals $\Gamma(f,f;f,f)=V(f)$; since $\kappa_4^{\mu}(f)=\kappa_4^{\mu}(f,f,f,f)$ and $\kappa_4^{\xi}(f)=\kappa_4^{\xi}(f,f,f,f)$, this is~\eqref{eq:fourth-cumulants-4}. Next, take $(f_1,f_2,f_3,f_4)=(f,f,\tau_xf,\tau_xf)$ in (a). The first term $\Gamma$ is then $\Gamma(f,f;\tau_xf,\tau_xf)$, which equals $\Gamma(f,f;f,f)=V(f)$ by (b); the second and the third are both equal to $\Gamma(f,\tau_xf;f,\tau_xf)=\operatorname{Var}_{\nu}\left(C_{\xi}(f,\tau_xf)\right)$. This proves~\eqref{eq:fourth-cumulants-5}. Finally, by the definition of $\mathrm{Ext}$ in Section~\ref{sec:cumulants-gaussian-laws},
\[
\mathrm{Ext}^{\mu}_x(f)=\kappa_4^{\mu}(f)-3\kappa_4^{\mu}(f,f,\tau_xf,\tau_xf),\qquad\mathrm{Ext}^{\xi}_x(f)=\kappa_4^{\xi}(f)-3\kappa_4^{\xi}(f,f,\tau_xf,\tau_xf).
\]
Subtracting three times~\eqref{eq:fourth-cumulants-5} from~\eqref{eq:fourth-cumulants-4} therefore yields
\[
\mathrm{Ext}^{\mu}_x(f)=\int\mathrm{d}\nu(\xi)\,\mathrm{Ext}^{\xi}_x(f)+3V(f)-3V(f)-6\operatorname{Var}_{\nu}\left(C_{\xi}(f,\tau_xf)\right),
\]
which is~\eqref{eq:fourth-cumulants-6}.

(d) In the proofs of (a), (b) and (c) the family $(\mu_{\xi})$ entered only through~\eqref{eq:fourth-cumulants-ED1} and through the following properties of the laws $\mu_{\xi}$, $\xi\in\Xi_0$: translation invariance, mean zero, finite fourth moments, and the bounds of Lemma~\ref{lem:moment-bounds}(e). Under the hypotheses of (d) every $\mu_{\xi}$, $\xi\in\Xi$, is translation invariant and satisfies~\eqref{eq:fourth-cumulants-H1}; by Lemma~\ref{lem:moment-bounds}(a), (c) and~(e) it therefore has mean zero and finite fourth moments and satisfies these bounds. The same arguments, with $\Xi$ in place of $\Xi_0$, prove (d).
\end{proof}

\begin{theorem}[Non-Gaussian components of positive measure]\label{thm:non-gaussian-components}
Let $f\in\mathcal S$, and let $c_4=\frac{3}{2}\exp(2)$ be the constant of Lemma~\ref{lem:moment-bounds}(c). The law $\mu$, its ergodic decomposition $(\Xi,\nu,(\mu_\xi))$ and the set $\Xi_0$ are those of Section~\ref{sec:cumulant-identities-non}.
Let $x\in\mathbb{R}^3$, and put
\begin{align*}
\delta_{\mathrm{Ext}}=\mathrm{Ext}^{\mu}_x(f)&=\kappa_4^{\mu}(f)-3\kappa_4^{\mu}(f,f,\tau_x f,\tau_x f)\\
&=2\kappa_4^{\mu}(f)-\frac{1}{4}\cdot\left[\kappa_4^{\mu}(f+\tau_x f)+\kappa_4^{\mu}(f-\tau_x f)\right].
\end{align*}
Then
\[
\int\mathrm{d}\nu(\xi)\,\mathrm{Ext}^{\xi}_x(f)\ge\delta_{\mathrm{Ext}} ,
\]
and $\mathrm{Ext}^{\xi}_x(f)\le(c_4+6)\cdot\|f\|_2^4$ for every $\xi\in\Xi_0$. Suppose that $\delta_{\mathrm{Ext}}>0$, which is to say that $\kappa_4^{\mu}(f+\tau_x f)+\kappa_4^{\mu}(f-\tau_x f)<8\kappa_4^{\mu}(f)$. Then
\[
\nu\left(\left\{\xi\in\Xi_0:\ \mathrm{Ext}^{\xi}_x(f)\ge\frac{\delta_{\mathrm{Ext}}}{2}\right\}\right)\ge\frac{\delta_{\mathrm{Ext}}}{2(c_4+6)\cdot\|f\|_2^4} ,
\]
and each $\xi$ that belongs to the set on the left has two further properties: the law $\mu_\xi$ is not Gaussian, and
\[
\max\left\{\left|\kappa_4^{\xi}(f)\right|,\ \left|\kappa_4^{\xi}(f+\tau_x f)\right|,\ \left|\kappa_4^{\xi}(f-\tau_x f)\right|\right\}\ge\frac{\delta_{\mathrm{Ext}}}{5}.
\]
The law $\mu$ is called a \emph{mixture} of laws with a given property if there exist a probability space $(\Xi,B,\nu)$ and a family $(\mu_\xi)_{\xi\in\Xi}$ of laws on $(\Omega,\mathcal F)$ such that~\eqref{eq:fourth-cumulants-ED1} holds and every $\mu_\xi$, $\xi\in\Xi$, has that property. For the property of being translation invariant and satisfying~\eqref{eq:fourth-cumulants-H1}, a probability space together with a family of this kind is called an \emph{invariant representation} of $\mu$, and one writes $C_\xi=S_2^{\mu_\xi}$, $\kappa_4^{\xi}=\kappa_4^{\mu_\xi}$ and $\mathrm{Ext}^{\xi}_x=\mathrm{Ext}^{\mu_\xi}_x$ in the notation of Section~\ref{sec:cumulants-gaussian-laws}. The preceding assertions remain true when the ergodic decomposition is replaced by an arbitrary invariant representation of $\mu$, and $\Xi_0$ by $\Xi$. Consequently, if $\delta_{\mathrm{Ext}}>0$ for at least one choice of $f\in\mathcal S$ and $x\in\mathbb{R}^3$, then no invariant representation of $\mu$ consists of Gaussian laws only; that is, $\mu$ is not a mixture of Gaussian laws that are translation invariant and satisfy~\eqref{eq:fourth-cumulants-H1}.
\end{theorem}

\begin{proof}
The argument applies to two kinds of families $(\Xi,\nu,(\mu_\xi))$: the ergodic decomposition of Section~\ref{sec:cumulant-identities-non}, together with the set $\Xi_0$ of Corollary~\ref{cor:properties-almost-every}, and an invariant representation of $\mu$, in which case $\Xi_0$ stands for $\Xi$. In both cases, for every $\xi\in\Xi_0$ the law $\mu_\xi$ is translation invariant and satisfies~\eqref{eq:fourth-cumulants-H1}: for the ergodic decomposition by Corollary~\ref{cor:properties-almost-every}, for an invariant representation by definition. Lemma~\ref{lem:moment-bounds}(a), (b), (c), applied to $\mu_\xi$, therefore gives, for every $\xi\in\Xi_0$ and every $g\in\mathcal S$,
\[
E_\xi X(g)=0,\qquad 0\le C_\xi(g,g)\le\|g\|_2^2,\qquad E_\xi\left[X(g)^4\right]\le c_4\cdot\|g\|_2^4.
\]

Of the three expressions for $\delta_{\mathrm{Ext}}$ in the statement, the first two coincide by the definition of $\mathrm{Ext}^{\mu}_x$ in Section~\ref{sec:cumulants-gaussian-laws}, and the first and the third by~\eqref{eq:fourth-cumulants-1} for $\mu$, which applies because $\mu$ is translation invariant.

Formula~\eqref{eq:fourth-cumulants-6} of Theorem~\ref{thm:fourth-cumulant-mixture}(c) holds for the ergodic decomposition and, by Theorem~\ref{thm:fourth-cumulant-mixture}(d), for every invariant representation of $\mu$. As a variance is nonnegative, it gives
\[
\int\mathrm{d}\nu(\xi)\,\mathrm{Ext}^{\xi}_x(f)=\mathrm{Ext}^{\mu}_x(f)+6\operatorname{Var}_{\nu}\left(C_\xi(f,\tau_x f)\right)\ge\mathrm{Ext}^{\mu}_x(f)=\delta_{\mathrm{Ext}} .
\]

Let $\xi\in\Xi_0$, and put $A=X(f)$, $B=X(\tau_x f)$ and $y=C_\xi(f,f)$. By the translation invariance of $\mu_\xi$, also $C_\xi(\tau_x f,\tau_x f)=y$. Because $\mu_\xi$ is centred, the definitions of Section~\ref{sec:cumulants-gaussian-laws} take the form
\[
\kappa_4^{\xi}(f)=E_\xi\left[A^4\right]-3y^2,\qquad\kappa_4^{\xi}(f,f,\tau_x f,\tau_x f)=E_\xi\left[A^2\cdot B^2\right]-y^2-2C_\xi(f,\tau_x f)^2,
\]
and hence
\begin{align*}
\mathrm{Ext}^{\xi}_x(f)&=E_\xi\left[A^4\right]-3y^2-3E_\xi\left[A^2\cdot B^2\right]+3y^2+6C_\xi(f,\tau_x f)^2\\
&\le E_\xi\left[A^4\right]+6y^2\le(c_4+6)\cdot\|f\|_2^4.
\end{align*}
The first inequality holds because $E_\xi\left[A^2\cdot B^2\right]\ge0$ and, by the Cauchy--Schwarz inequality,
\[
C_\xi(f,\tau_x f)^2=\left(E_\xi\left[A\cdot B\right]\right)^2\le E_\xi\left[A^2\right]\cdot E_\xi\left[B^2\right]=y^2;
\]
the second follows from $E_\xi\left[A^4\right]\le c_4\cdot\|f\|_2^4$ and $0\le y\le\|f\|_2^2$.

Assume from now on that $\delta_{\mathrm{Ext}}>0$; by the third expression for $\delta_{\mathrm{Ext}}$, this is equivalent to $\kappa_4^{\mu}(f+\tau_x f)+\kappa_4^{\mu}(f-\tau_x f)<8\kappa_4^{\mu}(f)$. Let $M=(c_4+6)\cdot\|f\|_2^4$, and regard the function $Z(\xi)=\mathrm{Ext}^{\xi}_x(f)$ of $\xi\in\Xi_0$ as a random variable on $\Xi$ with respect to the probability measure $\nu$, for which $\Xi_0$ has full measure. By the two preceding steps, $Z\le M$ on $\Xi_0$ and $\int\mathrm{d}\nu(\xi)\,Z(\xi)\ge\delta_{\mathrm{Ext}}>0$; in particular $M\ge\delta_{\mathrm{Ext}}>0$, so that $\|f\|_2>0$. Lemma~\ref{lem:large-values-variable} with $\delta=\delta_{\mathrm{Ext}}$ yields
\[
\nu\left(\left\{\xi\in\Xi_0:\ \mathrm{Ext}^{\xi}_x(f)\ge\frac{\delta_{\mathrm{Ext}}}{2}\right\}\right)\ge\frac{\delta_{\mathrm{Ext}}}{2M}=\frac{\delta_{\mathrm{Ext}}}{2(c_4+6)\cdot\|f\|_2^4} .
\]

Let $\xi\in\Xi_0$ with $\mathrm{Ext}^{\xi}_x(f)\ge\frac{\delta_{\mathrm{Ext}}}{2}$. The law $\mu_\xi$ is centred; if it were Gaussian, Section~\ref{sec:cumulants-gaussian-laws} would give $\mathrm{Ext}^{\xi}_x(f)=0<\frac{\delta_{\mathrm{Ext}}}{2}$. Hence $\mu_\xi$ is not Gaussian.

Suppose that some invariant representation of $\mu$ consisted of Gaussian laws only. Everything above would apply to it, with $\Xi_0=\Xi$. By the measure bound, the set of the $\xi\in\Xi$ with $\mathrm{Ext}^{\xi}_x(f)\ge\frac{\delta_{\mathrm{Ext}}}{2}$ would have positive $\nu$-measure and would therefore contain some $\xi$; by the preceding paragraph $\mu_\xi$ would not be Gaussian, a contradiction. Hence no invariant representation of $\mu$ consists of Gaussian laws only, which by definition means that $\mu$ is not a mixture of Gaussian laws that are translation invariant and satisfy~\eqref{eq:fourth-cumulants-H1}.

Finally, for a family of either kind let $\xi\in\Xi_0$ with $\mathrm{Ext}^{\xi}_x(f)\ge\frac{\delta_{\mathrm{Ext}}}{2}$, and let $m_\xi$ be the largest of the three numbers $\left|\kappa_4^{\xi}(f)\right|$, $\left|\kappa_4^{\xi}(f+\tau_x f)\right|$, $\left|\kappa_4^{\xi}(f-\tau_x f)\right|$. Since $\mu_\xi$ is translation invariant, \eqref{eq:fourth-cumulants-1} holds for $\mu_\xi$, and the triangle inequality gives
\[
\frac{\delta_{\mathrm{Ext}}}{2}\le\mathrm{Ext}^{\xi}_x(f)\le2\left|\kappa_4^{\xi}(f)\right|+\frac{1}{4}\cdot\left|\kappa_4^{\xi}(f+\tau_x f)\right|+\frac{1}{4}\cdot\left|\kappa_4^{\xi}(f-\tau_x f)\right|\le\frac{5}{2}\cdot m_\xi ,
\]
so that $m_\xi\ge\frac{\delta_{\mathrm{Ext}}}{5}$.
\end{proof}

\begin{lemma}[White noise with telegraph-modulated variance]\label{lem:white-noise-telegraph}
Let $\acute{N}$ be a two-sided Poisson process of rate $\frac{1}{2}$ on $\mathbb{R}$; such a process exists. Thus $\acute{N}(0)=0$, the paths $t\mapsto\acute{N}(t)$ are integer valued and right-continuous and have finitely many jumps in every bounded interval, the increments of $\acute{N}$ are independent and stationary, and the restrictions of $\acute{N}$ to $\{t>0\}$ and to $\{t<0\}$ are independent. Let $z_0$ be a random variable that is independent of $\acute{N}$ and has the uniform law on $\{+1,-1\}$. The \emph{telegraph process} $z$ is defined by
\[
z(t)=z_0\cdot(-1)^{\acute{N}(t)}\qquad(t\in\mathbb{R}).
\]
Let $W$ be the set of the right-continuous functions $w\colon\mathbb{R}\to\{+1,-1\}$ that have finitely many jumps in every bounded interval, endowed with the $\sigma$-algebra $\mathcal W$ generated by the coordinate maps $w\mapsto w(t)$, $t\in\mathbb{R}$; for $u\in\mathbb{R}$ let $\operatorname{sh}_u\colon W\to W$ be the shift, $(\operatorname{sh}_uw)(t)=w(t+u)$. The paths of $z$ belong to $W$, and $z$ is a random variable with values in $(W,\mathcal W)$.

Let $z_1$, $z_2$, $z_3$ be independent random variables with values in $W$, each of which has the law of $z$; from now on $z$ denotes the triple $(z_1,z_2,z_3)$, and $E_z$ the expectation with respect to it. Let $p_0$ and $q_0$ be real numbers with $0<q_0\le\min\{p_0,1-p_0\}$. Put
\[
v_z(x)=p_0+q_0\cdot z_1(x_1)\cdot z_2(x_2)\cdot z_3(x_3)\qquad(x\in\mathbb{R}^3),
\]
a function on $\mathbb{R}^3$ with values in $[p_0-q_0,p_0+q_0]\subset[0,1]$, and
\[
\langle f,v_zg\rangle=\int\mathrm{d}x\,v_z(x)\cdot\sum_{a<b}f_{ab}(x)\cdot g_{ab}(x)\qquad(f,g\in\mathcal S).
\]
Then there exists exactly one law $\mu'$ on $(\Omega,\mathcal F)$ such that
\[
\int\mathrm{d}\mu'\,\exp\left(iX(f)\right)=E_z\exp\left(-\frac{1}{2}\langle f,v_zf\rangle\right)\qquad\text{for every }f\in\mathcal S.
\]
The law $\mu'$ satisfies \eqref{eq:fourth-cumulants-H1}, \eqref{eq:fourth-cumulants-H2} and \eqref{eq:fourth-cumulants-H3} in place of $\mu$, it is ergodic, and for all $f,g\in\mathcal S$
\begin{gather*}
S_2^{\mu'}(f,g)=p_0\cdot\langle f,g\rangle,\\
\kappa_4^{\mu'}(f)=3q_0^2\int\mathrm{d}x\int\mathrm{d}y\,|f(x)|^2\cdot|f(y)|^2\cdot\exp\left(-|x_1-y_1|-|x_2-y_2|-|x_3-y_3|\right).
\end{gather*}
Consequently $\kappa_4^{\mu'}(f)>0$ for every $f\in\mathcal S$ with $f\ne0$, and $\mu'$ is not Gaussian.
\end{lemma}

\begin{proof}
\emph{Step 1: four properties of the telegraph process.} In this step $z$ is the telegraph process itself, and $E$ is the expectation on the probability space on which $\acute{N}$ and $z_0$ are defined. Let $L^2(W)$ be the space $L^2$ of $(W,\mathcal W)$ with respect to the law of $z$, with the norm $\|a\|_2=\left(E\left[|a(z)|^2\right]\right)^{\frac{1}{2}}$.

\emph{Measurability.} For $t\in\mathbb{R}$ and $n\ge1$ let $t_n$ be the least element of $2^{-n}\mathbb{Z}$ that is larger than $t$. Every $w\in W$ is right-continuous, so that $w(t)=\lim_{n\to\infty}w(t_n)$; and $(t,w)\mapsto w(t_n)$ is measurable on $\mathbb{R}\times W$, since $t\mapsto t_n$ is constant on each interval $\left[(k-1)2^{-n},k2^{-n}\right)$, $k\in\mathbb{Z}$. Hence $(t,w)\mapsto w(t)$ is measurable on $\mathbb{R}\times W$, and $(u,w)\mapsto\operatorname{sh}_uw$ is a measurable map of $\mathbb{R}\times W$ into $W$.

\emph{Stationarity.} Let $u\in\mathbb{R}$. For every $t\in\mathbb{R}$
\[
z(u+t)=z(u)\cdot(-1)^{\acute{N}(u+t)-\acute{N}(u)}.
\]
The process $t\mapsto\acute{N}(u+t)-\acute{N}(u)$ vanishes at $t=0$, and its increments are increments of $\acute{N}$; it is therefore a two-sided Poisson process of rate $\frac{1}{2}$ as well. Conditionally on $\acute{N}$, the variable $z(u)=z_0\cdot(-1)^{\acute{N}(u)}$ has the uniform law on $\{+1,-1\}$, because $z_0$ has this law and is independent of $\acute{N}$; hence $z(u)$ is independent of $\acute{N}$. So the pair formed by $z(u)$ and the process $t\mapsto\acute{N}(u+t)-\acute{N}(u)$ has the same law as the pair $(z_0,\acute{N})$, and the identity above shows
\begin{equation}\label{eq:fourth-cumulants-7}
\operatorname{sh}_uz\ \text{has the same law as}\ z\qquad\text{for every }u\in\mathbb{R}.
\end{equation}
In particular every $z(u)$ has the uniform law on $\{+1,-1\}$, so that $E\left[z(u)\right]=0$.

\emph{Reversibility.} For $w\in W$ and $t\in\mathbb{R}$ the left limit $w(t-)=\lim_{s\uparrow t}w(s)$ exists, and $w(t-)=w(t)$ unless $t$ is a jump time of $w$. For $w\in W$ let $\check w$ be the function $t\mapsto w\left((-t)-\right)$ on $\mathbb{R}$. It belongs to $W$, because for every $t\in\mathbb{R}$ the function $w$ is constant on some interval with right end point $-t$; and the map $w\mapsto\check w$ of $W$ into $W$ is measurable, since $\check w(t)=\lim_{n\to\infty}w\left(-t-\frac{1}{n}\right)$ for every $t\in\mathbb{R}$. Put
\[
\check z(t)=\lim_{s\downarrow t}z(-s)=z\left((-t)-\right)\qquad(t\in\mathbb{R}).
\]
Then $\check z$, the image of $z$ under this map, is a random variable with values in $W$, and $\check z(t)=z(-t)$ for all $t\in\mathbb{R}$ except those for which $-t$ is a jump time of $z$; these exceptional $t$ are finitely many in every bounded interval, so that they form a Lebesgue null set. For $t>0$ one has
\[
z(t)=z_0\cdot(-1)^{\acute{N}(t)},\qquad\check z(t)=z_0\cdot(-1)^{\acute{N}(0)-\acute{N}((-t)-)}.
\]
The variable $z_0=z(0)$, the process $\left(\acute{N}(t)\right)_{t>0}$ and the process $\left(\acute{N}(0)-\acute{N}((-t)-)\right)_{t>0}$, which counts the jumps of $\acute{N}$ in $[-t,0]$, are independent, and the last two are Poisson processes of rate $\frac{1}{2}$ on $(0,\infty)$. Therefore
\begin{multline}\label{eq:fourth-cumulants-8}
\text{conditionally on $z(0)$, the processes $\left(z(t)\right)_{t>0}$ and $\left(\check z(t)\right)_{t>0}$}\\
\text{are independent and identically distributed.}
\end{multline}
In particular $\check z$ has the law of $z$. Indeed, for $t>0$ one has $z(-t)=\check z(t-)$ and $\check z(-t)=z(t-)$. Let $\Upsilon$ send a triple $(\sigma,a,b)$, formed by a sign $\sigma$ and two right-continuous functions $a,b\colon(0,\infty)\to\{+1,-1\}$ with finitely many jumps in every bounded interval, to the function on $\mathbb{R}$ that equals $\sigma$ at $0$, $a(t)$ at $t>0$ and $b(|t|-)$ at $t<0$. This map is measurable, and
\[
z=\Upsilon\left(z(0),\left(z(t)\right)_{t>0},\left(\check z(t)\right)_{t>0}\right),\qquad\check z=\Upsilon\left(\check z(0),\left(\check z(t)\right)_{t>0},\left(z(t)\right)_{t>0}\right).
\]
Moreover $\check z(0)=z(0)$ almost surely, since $\acute{N}$ has a jump at $0$ with probability $0$. By~\eqref{eq:fourth-cumulants-8} the two triples to which $\Upsilon$ is applied have the same law; hence so have $z$ and $\check z$.

\emph{Transition probabilities.} Let $h$ be a complex function on $\{+1,-1\}$, put $\bar{h}=\frac{1}{2}\left(h(1)+h(-1)\right)$, and for $t\ge0$ let $R_th$ be the function $\bar{h}+\exp(-t)\cdot\left(h-\bar{h}\right)$ on $\{+1,-1\}$. Let $t\ge0$. The variable $\acute{N}(t)$ is independent of $z_0$ and of the restriction of $\acute{N}$ to $\{s<0\}$, and it has the Poisson law with mean $\frac{t}{2}$; so it is odd with probability
\[
\exp\left(-\frac{t}{2}\right)\cdot\sinh\left(\frac{t}{2}\right)=\frac{1-\exp(-t)}{2}.
\]
Since $z(t)=z_0\cdot(-1)^{\acute{N}(t)}$, it follows that
\[
E\left[h\left(z(t)\right)\,\middle|\,z_0,\ \acute{N}(s),\ s\le0\right]=\frac{1+\exp(-t)}{2}\cdot h(z_0)+\frac{1-\exp(-t)}{2}\cdot h(-z_0)=(R_th)(z_0).
\]
The right side is a function of $z(0)=z_0$, and the $\sigma$-algebra generated by the variables $z(s)=z_0\cdot(-1)^{\acute{N}(s)}$, $s\le0$, contains the one generated by $z(0)$ and is contained in the one generated by $z_0$ and the $\acute{N}(s)$, $s\le0$. Hence, by the tower property of conditional expectations,
\begin{equation}\label{eq:fourth-cumulants-9}
E\left[h\left(z(t)\right)\,\middle|\,z(s),\ s\le0\right]=(R_th)\left(z(0)\right)\qquad(t\ge0).
\end{equation}
For $h(y)=y$ one has $\bar{h}=0$ and $R_th=\exp(-t)\cdot h$. For $s\le t$, \eqref{eq:fourth-cumulants-7} and \eqref{eq:fourth-cumulants-9} therefore give $E\left[z(s)\cdot z(t)\right]=E\left[z(0)\cdot z(t-s)\right]=\exp\left(-(t-s)\right)\cdot E\left[z(0)^2\right]$; thus
\[
E\left[z(s)\cdot z(t)\right]=\exp\left(-|s-t|\right)\qquad(s,t\in\mathbb{R}).
\]

\emph{Mixing in the mean.} For $s\in\mathbb{R}$ let $\mathcal W_{\le s}$ and $\mathcal W_{\ge s}$ be the $\sigma$-algebras on $W$ generated by the maps $w\mapsto w(t)$ with $t\le s$, respectively with $t\ge s$.

First, the Markov property of $z$. Let $s\in\mathbb{R}$, and let $\mathcal G_s$ be the $\sigma$-algebra generated by $z(s)$ and the increments $\acute{N}(t)-\acute{N}(s)$, $t\le s$. As $z(t)=z(s)\cdot(-1)^{\acute{N}(t)-\acute{N}(s)}$ for $t\le s$, the $\sigma$-algebra generated by the $z(t)$, $t\le s$, is contained in $\mathcal G_s$; and $z(s+t)=z(s)\cdot(-1)^{\acute{N}(s+t)-\acute{N}(s)}$ for $t\ge0$. The variable $z(s)$ is independent of $\acute{N}$, as shown in the proof of~\eqref{eq:fourth-cumulants-7}, and the increments of $\acute{N}$ over subintervals of $[s,\infty)$ are independent of its increments over subintervals of $(-\infty,s]$; therefore $z(s)$, the process $\left(\acute{N}(t)-\acute{N}(s)\right)_{t\le s}$ and the process $\left(\acute{N}(s+t)-\acute{N}(s)\right)_{t\ge0}$ are independent, and the last one is independent of $\mathcal G_s$. Let $\Psi$ be a $\mathcal W_{\ge0}$-measurable function on $W$ with $E\left[|\Psi(\operatorname{sh}_sz)|\right]<\infty$. The variable $\Psi(\operatorname{sh}_sz)$ is a function of $z(s)$ and of the process $\left(\acute{N}(s+t)-\acute{N}(s)\right)_{t\ge0}$ alone; so its conditional expectation given $\mathcal G_s$ is a function of $z(s)$. Since $z(s)$ is measurable with respect to the $\sigma$-algebra generated by the $z(t)$, $t\le s$, which is contained in $\mathcal G_s$, it follows that
\[
E\left[\Psi(\operatorname{sh}_sz)\,\middle|\,z(t),\ t\le s\right]=E\left[\Psi(\operatorname{sh}_sz)\,\middle|\,z(s)\right].
\]
This is the Markov property: given the values $z(t)$, $t\le s$, the conditional law of $\left(z(s+t)\right)_{t\ge0}$ depends on $z(s)$ only.

Now let $T_0\ge0$, let $a\in L^2(W)$ be $\mathcal W_{\le T_0}$-measurable, let $b\in L^2(W)$ be $\mathcal W_{\ge-T_0}$-measurable, and let $u>2T_0$. Put $s=u-T_0$, so that $s>T_0$, and $\Psi=b\circ\operatorname{sh}_{T_0}$; this function is $\mathcal W_{\ge0}$-measurable, and $b\circ\operatorname{sh}_u=\Psi\circ\operatorname{sh}_s$. By~\eqref{eq:fourth-cumulants-7} the pair $\left(z(s),\operatorname{sh}_sz\right)$ has the law of $\left(z(0),z\right)$, and the pair $\left(z(-T_0),z\right)$ has the law of $\left(z(0),\operatorname{sh}_{T_0}z\right)$. With
\[
\tilde{b}(y)=E\left[b(z)\,\middle|\,z(-T_0)=y\right]\qquad\left(y\in\{+1,-1\}\right),
\]
the Markov property therefore gives
\[
E\left[b(\operatorname{sh}_uz)\,\middle|\,z(t),\ t\le s\right]=E\left[\Psi(\operatorname{sh}_sz)\,\middle|\,z(s)\right]=\tilde{b}\left(z(s)\right).
\]
The variable $a(z)$ is measurable with respect to the $\sigma$-algebra generated by the $z(t)$, $t\le T_0$, and $T_0<s$. So the preceding identity gives the first of the following two equalities, and~\eqref{eq:fourth-cumulants-9}, applied to the process $\operatorname{sh}_{T_0}z$, which has the law of $z$ by~\eqref{eq:fourth-cumulants-7}, with $s-T_0=u-2T_0$ in place of $t$, gives the second:
\[
E\left[a(z)\cdot b(\operatorname{sh}_uz)\right]=E\left[a(z)\cdot\tilde{b}\left(z(s)\right)\right]=E\left[a(z)\cdot\left(R_{u-2T_0}\tilde{b}\right)\left(z(T_0)\right)\right].
\]
As $u\to\infty$, the functions $R_{u-2T_0}\tilde{b}$ converge uniformly on $\{+1,-1\}$ to the constant $\frac{1}{2}\left(\tilde{b}(1)+\tilde{b}(-1)\right)$, which equals $E\left[b(z)\right]$ because $z(-T_0)$ has the uniform law. Hence $E\left[a(z)\cdot b(\operatorname{sh}_uz)\right]\to E\left[a(z)\right]\cdot E\left[b(z)\right]$ as $u\to\infty$. The function $u\mapsto E\left[a(z)\cdot b(\operatorname{sh}_uz)\right]$ is measurable, by the measurability of $(u,w)\mapsto\operatorname{sh}_uw$ and Fubini's theorem, and it is bounded by $\|a\|_2\cdot\|b\|_2$, by the Cauchy--Schwarz inequality and~\eqref{eq:fourth-cumulants-7}; so its averages over $[0,T]$ converge to the same limit as $T\to\infty$.

If the $\sigma$-algebra of a probability space is generated by the union of an increasing sequence of sub-$\sigma$-algebras, then the square integrable functions that are measurable with respect to some member of the sequence form a dense subspace of $L^2$ of this space: by the martingale convergence theorem, the conditional expectations of a square integrable function with respect to the members of the sequence converge to the function in $L^2$. The $\sigma$-algebras $\mathcal W_{\le n}$, $n\in\mathbb{N}$, increase with $n$, and $\mathcal W$ is generated by their union, because $w\mapsto w(t)$ is $\mathcal W_{\le n}$-measurable as soon as $n\ge t$. Hence the functions in $L^2(W)$ that are $\mathcal W_{\le n}$-measurable for some $n$ form a dense subspace of $L^2(W)$. In the same way the $\sigma$-algebras $\mathcal W_{\ge-n}$, $n\in\mathbb{N}$, increase with $n$, $\mathcal W$ is generated by their union, and the functions in $L^2(W)$ that are $\mathcal W_{\ge-n}$-measurable for some $n$ form a dense subspace of $L^2(W)$. If $a$ is $\mathcal W_{\le n}$-measurable and $b$ is $\mathcal W_{\ge-n'}$-measurable, the preceding paragraph applies with $T_0=\max\{n,n'\}$. For $a,b\in L^2(W)$ and $T>0$ both
\[
\frac{1}{T}\int_0^T\mathrm{d}u\,E\left[a(z)\cdot b(\operatorname{sh}_uz)\right]\qquad\text{and}\qquad E\left[a(z)\right]\cdot E\left[b(z)\right]
\]
are bilinear in $(a,b)$ and bounded in absolute value by $\|a\|_2\cdot\|b\|_2$. If $a',b'\in L^2(W)$ satisfy $\|a-a'\|_2\le\varepsilon$ and $\|b-b'\|_2\le\varepsilon$, each of the two expressions therefore changes by at most $\varepsilon\cdot\left(\|b\|_2+\|a'\|_2\right)$ when $(a,b)$ is replaced by $(a',b')$. Approximating $a$ and $b$ by elements of the two dense subspaces one obtains
\begin{equation}\label{eq:fourth-cumulants-10}
\lim_{T\to\infty}\frac{1}{T}\int_0^T\mathrm{d}u\,E\left[a(z)\cdot b(\operatorname{sh}_uz)\right]=E\left[a(z)\right]\cdot E\left[b(z)\right]\qquad\left(a,b\in L^2(W)\right).
\end{equation}

\emph{Step 2: the law $\mu'$ and its finite-dimensional distributions.} From now on $z=(z_1,z_2,z_3)$ is the triple of the statement. The properties of the telegraph process established in Step~1 are properties of its law on $(W,\mathcal W)$, the process $\check z$ being the image of $z$ under the measurable map $w\mapsto\check w$; so they hold for each of $z_1$, $z_2$, $z_3$, which have this law. By Step~1 the function $(x,z)\mapsto v_z(x)$ is measurable on $\mathbb{R}^3\times W^3$, and $0\le v_z\le1$. By Fubini's theorem $z\mapsto\langle f,v_zg\rangle$ is measurable for $f,g\in\mathcal S$, and
\[
0\le\langle f,v_zf\rangle=\int\mathrm{d}x\,v_z(x)\cdot|f(x)|^2\le\|f\|_2^2,\qquad\left|\langle f,v_zg\rangle\right|\le\|f\|_2\cdot\|g\|_2 .
\]
For $f\in\mathcal S$ put $\Phi'(f)=E_z\exp\left(-\frac{1}{2}\langle f,v_zf\rangle\right)$. For fixed $z$ the form $(f,g)\mapsto\langle f,v_zg\rangle$ on $\mathcal S$ is bilinear, symmetric and positive semidefinite. Hence $f\mapsto\exp\left(-\frac{1}{2}\langle f,v_zf\rangle\right)$ is positive definite, that is, $\sum_{i,k=1}^n\overline{c_i}\cdot c_k\cdot\exp\left(-\frac{1}{2}\langle f_k-f_i,v_z(f_k-f_i)\rangle\right)\ge0$ for $f_1,\dots,f_n\in\mathcal S$ and $c_1,\dots,c_n\in\mathbb{C}$: on the linear span of finitely many elements of $\mathcal S$ it is the characteristic function of a centred normal law. As an average over $z$ of positive definite functionals, $\Phi'$ is positive definite, and $\Phi'(0)=1$. For $f,g\in\mathcal S$ the bound $0\le v_z\le1$, the inequality $\left||f(x)|^2-|g(x)|^2\right|\le|f(x)-g(x)|\cdot\left(|f(x)|+|g(x)|\right)$ and the Cauchy--Schwarz inequality give
\[
\left|\langle f,v_zf\rangle-\langle g,v_zg\rangle\right|\le\|f-g\|_2\cdot\left(\|f\|_2+\|g\|_2\right).
\]
Since $\left|\exp\left(-\frac{\eta}{2}\right)-\exp\left(-\frac{\eta'}{2}\right)\right|\le\frac{1}{2}|\eta-\eta'|$ for $\eta,\eta'\ge0$, and since convergence in $\mathcal S$ implies convergence in $L^2$, the functional $\Phi'$ is continuous on $\mathcal S$. By Minlos' theorem~\cite[Theorem~3.4.2]{GlimmJaffe} there is a law $\mu'$ on $(\Omega,\mathcal F)$ with $\int\mathrm{d}\mu'\,\exp\left(iX(f)\right)=\Phi'(f)$ for every $f\in\mathcal S$, and by Lemma~\ref{lem:exponentials-determination-law}(b), applied with $V_0=\mathcal S$, there is only one such law. Let $E'$ be the expectation with respect to $\mu'$.

Let $n\ge1$ and $f_1,\dots,f_n\in\mathcal S$. For $z\in W^3$ let $K_z$ be the real symmetric positive semidefinite $n\times n$ matrix with the entries $\langle f_k,v_zf_l\rangle$, $1\le k,l\le n$, and let $\lambda_z$ be the centred normal law on $\mathbb{R}^n$ with covariance matrix $K_z$. This law is the image of the standard normal law on $\mathbb{R}^n$ under the linear map $K_z^{\frac{1}{2}}$, the positive semidefinite square root of $K_z$; that is, with the Euclidean norm $|\eta|$ of $\eta\in\mathbb{R}^n$,
\[
\int\mathrm{d}\lambda_z(y)\,\varphi(y)=(2\pi)^{-\frac{n}{2}}\int_{\mathbb{R}^n}\mathrm{d}\eta\,\exp\left(-\frac{1}{2}|\eta|^2\right)\cdot\varphi\left(K_z^{\frac{1}{2}}\eta\right)
\]
for every Borel function $\varphi\colon\mathbb{R}^n\to[0,\infty]$. The entries of $K_z$ are measurable functions of $z$, and the positive semidefinite square root is a continuous function on the closed cone of the positive semidefinite $n\times n$ matrices; so the integrand on the right is a measurable function of $(z,\eta)\in W^3\times\mathbb{R}^n$, and by Tonelli's theorem the left side is a measurable function of $z$. In particular $z\mapsto\lambda_z(Y)$ is measurable for every Borel set $Y\subset\mathbb{R}^n$. For $t\in\mathbb{R}^n$ one has $\sum_lt_l\cdot X(f_l)=X\left(\sum_lt_l\cdot f_l\right)$, and the bilinearity of $(f,g)\mapsto\langle f,v_zg\rangle$ gives
\[
E'\exp\left(i\sum_{l=1}^nt_l\cdot X(f_l)\right)=E_z\exp\left(-\frac{1}{2}\sum_{k,l=1}^nt_k\cdot t_l\cdot\langle f_k,v_zf_l\rangle\right)=E_z\int\mathrm{d}\lambda_z(y)\,\exp\left(i\sum_{l=1}^nt_l\cdot y_l\right).
\]
By the uniqueness of the Fourier transform of a finite measure on $\mathbb{R}^n$, the law of $\left(X(f_1),\dots,X(f_n)\right)$ under $\mu'$ is the mixture $Y\mapsto E_z\left[\lambda_z(Y)\right]$, $Y\subset\mathbb{R}^n$ a Borel set, of the laws $\lambda_z$; that is,
\[
E'\left[\varphi\left(X(f_1),\dots,X(f_n)\right)\right]=E_z\int\mathrm{d}\lambda_z(y)\,\varphi(y)
\]
for every Borel function $\varphi\colon\mathbb{R}^n\to[0,\infty]$.

\emph{Step 3: the exponential bound and the invariance under $\mathbb{R}^3\rtimes O_{\mathrm{cub}}$.} Let $f\in\mathcal S$. Step~2 with $n=1$, $f_1=f$ and $\varphi=\exp$, together with $\langle f,v_zf\rangle\le\|f\|_2^2$, gives
\[
\int\mathrm{d}\mu'\,\exp X(f)=E_z\exp\left(\frac{1}{2}\langle f,v_zf\rangle\right)\le\exp\left(\frac{1}{2}\|f\|_2^2\right),
\]
which is \eqref{eq:fourth-cumulants-H1} for $\mu'$.

For $x\in\mathbb{R}^3$ put $\operatorname{sh}_xz=\left(\operatorname{sh}_{x_1}z_1,\operatorname{sh}_{x_2}z_2,\operatorname{sh}_{x_3}z_3\right)$. By~\eqref{eq:fourth-cumulants-7} and the independence of $z_1$, $z_2$, $z_3$, the triple $\operatorname{sh}_xz$ has the law of $z$. Since $(\tau_xf)(y)=f(y-x)$ and $v_z(y+x)=v_{\operatorname{sh}_xz}(y)$,
\[
\Phi'(\tau_xf)=E_z\exp\left(-\frac{1}{2}\int\mathrm{d}y\,v_z(y+x)\cdot|f(y)|^2\right)=E_z\exp\left(-\frac{1}{2}\langle f,v_{\operatorname{sh}_xz}f\rangle\right)=\Phi'(f).
\]
Let $r\in O_{\mathrm{cub}}$. There are a permutation $\pi$ of $\{1,2,3\}$ and signs $s_1,s_2,s_3\in\{+1,-1\}$, determined by $r$, such that $(ry)_l=s_l\cdot y_{\pi(l)}$ for $y\in\mathbb{R}^3$ and $l=1,2,3$. The map $\wedge^2r$ is orthogonal on $\Lambda^2$, so that $|(U_{0,r}f)(y)|=|f(r^{-1}y)|$, and the substitution $y\mapsto ry$, whose Jacobian has absolute value $1$, gives
\[
\Phi'(U_{0,r}f)=E_z\exp\left(-\frac{1}{2}\int\mathrm{d}y\,v_z(ry)\cdot|f(y)|^2\right),\qquad v_z(ry)=p_0+q_0\cdot\prod_{l=1}^3z_l\left(s_l\cdot y_{\pi(l)}\right).
\]
Let $\check z_l$ be formed from $z_l$ as $\check z$ from $z$ in Step~1, and put $z'_{\pi(l)}=z_l$ if $s_l=1$ and $z'_{\pi(l)}=\check z_l$ if $s_l=-1$. By the independence of $z_1$, $z_2$, $z_3$ and by the consequence of~\eqref{eq:fourth-cumulants-8} that $\check z_l$ has the law of $z_l$, the triple $z'=(z'_1,z'_2,z'_3)$ has the law of $z$; and as $z_l(-t)=\check z_l(t)$ for almost every $t\in\mathbb{R}$, one has $v_z(ry)=v_{z'}(y)$ for almost every $y\in\mathbb{R}^3$. Hence
\[
\Phi'(U_{0,r}f)=E_z\exp\left(-\frac{1}{2}\langle f,v_{z'}f\rangle\right)=\Phi'(f).
\]
Every element $g=(b,r)$ of $\mathbb{R}^3\rtimes O_{\mathrm{cub}}$ satisfies $U_g=\tau_b\,U_{0,r}$; so $\Phi'(U_gf)=\Phi'(f)$ for all such $g$ and all $f\in\mathcal S$. As $X(f)\circ T_g=X(U_g^{-1}f)=X(U_{g^{-1}}f)$, the law $\mu'\circ T_g^{-1}$ has the characteristic functional $f\mapsto\Phi'(U_{g^{-1}}f)=\Phi'(f)$, and Lemma~\ref{lem:exponentials-determination-law}(b), applied with $V_0=\mathcal S$, gives $\mu'\circ T_g^{-1}=\mu'$. This is \eqref{eq:fourth-cumulants-H2} for $\mu'$; in particular $\mu'$ is translation invariant.

\emph{Step 4: reflection positivity.} For a triple $w=(w_1,w_2,w_3)$ of measurable functions $\mathbb{R}\to\{+1,-1\}$ define $v_w$ and $\langle f,v_wg\rangle$ by the formulas of the statement with $w$ in place of $z$. Let first $j=1$. Let $n\ge1$, let $f_1,\dots,f_n\in\mathcal D$ have their supports in $\{x_1>0\}$, let $c_1,\dots,c_n\in\mathbb{C}$, and $G=\sum_{k=1}^nc_k\cdot\exp\left(iX(f_k)\right)$; every element of $A_1^+$ has this form. For $1\le k,l\le n$ the support of $\theta_1f_k$ lies in $\{x_1<0\}$, so that $f_l$ and $\theta_1f_k$ have disjoint supports, and $|(\theta_1f_k)(x)|=|f_k(\theta_1x)|$. Let $\theta z$ be the triple given by $(\theta z)_1(t)=z_1(-t)$, $(\theta z)_2=z_2$, $(\theta z)_3=z_3$, and put
\[
\boldsymbol{\alpha}_k(w)=\exp\left(-\frac{1}{2}\langle f_k,v_wf_k\rangle\right)\qquad(k=1,\dots,n).
\]
Since $v_z(\theta_1y)=v_{\theta z}(y)$, the substitution $x=\theta_1y$ gives $\langle\theta_1f_k,v_z\,\theta_1f_k\rangle=\langle f_k,v_{\theta z}f_k\rangle$. As $X(f_l)-X(\theta_1f_k)=X(f_l-\theta_1f_k)$ and the supports of $f_l$ and $\theta_1f_k$ are disjoint, the definition of $\mu'$ yields
\[
\int\mathrm{d}\mu'\,\exp\left(iX(f_l)-iX(\theta_1f_k)\right)=E_z\exp\left(-\frac{1}{2}\langle f_l,v_zf_l\rangle-\frac{1}{2}\langle\theta_1f_k,v_z\,\theta_1f_k\rangle\right)=E_z\left[\boldsymbol{\alpha}_l(z)\cdot\boldsymbol{\alpha}_k(\theta z)\right].
\]
Each $\boldsymbol{\alpha}_k(z)$ is a real number, and, as $f_k$ vanishes on $\{x_1\le0\}$, it is a function of $z_2$, $z_3$ and the restriction of $z_1$ to $(0,\infty)$. Put $\boldsymbol{\alpha}=\sum_kc_k\cdot\boldsymbol{\alpha}_k$. Since $\mathbb{U}_{0,\theta_1}X(f)=X(\theta_1f)$, one has $\Theta_1G=\sum_k\overline{c_k}\cdot\exp\left(-iX(\theta_1f_k)\right)$, and consequently
\[
\int\mathrm{d}\mu'\,(\Theta_1G)\cdot G=\sum_{k,l=1}^n\overline{c_k}\cdot c_l\cdot E_z\left[\boldsymbol{\alpha}_l(z)\cdot\boldsymbol{\alpha}_k(\theta z)\right]=E_z\left[\overline{\boldsymbol{\alpha}(\theta z)}\cdot\boldsymbol{\alpha}(z)\right].
\]
Let $\check z_1$ be formed from $z_1$ as in Step~1, and $z^*=(\check z_1,z_2,z_3)$. As $z_1(-t)=\check z_1(t)$ for almost every $t$, the functions $v_{\theta z}$ and $v_{z^*}$ agree almost everywhere, and $\boldsymbol{\alpha}(\theta z)=\boldsymbol{\alpha}(z^*)$. So there is a measurable function $\Psi$ with
\[
\boldsymbol{\alpha}(z)=\Psi\left(\left(z_1(t)\right)_{t>0},z_2,z_3\right),\qquad\boldsymbol{\alpha}(\theta z)=\Psi\left(\left(\check z_1(t)\right)_{t>0},z_2,z_3\right).
\]
By~\eqref{eq:fourth-cumulants-8} for $z_1$, and since $(z_2,z_3)$ is independent of $z_1$, the processes $\left(z_1(t)\right)_{t>0}$ and $\left(\check z_1(t)\right)_{t>0}$ are independent and identically distributed conditionally on $\left(z_1(0),z_2,z_3\right)$. Therefore
\[
E_z\left[\overline{\boldsymbol{\alpha}(\theta z)}\cdot\boldsymbol{\alpha}(z)\,\middle|\,z_1(0),z_2,z_3\right]=\left|E_z\left[\boldsymbol{\alpha}(z)\,\middle|\,z_1(0),z_2,z_3\right]\right|^2\ge0,
\]
and taking expectations one obtains $\int\mathrm{d}\mu'\,(\Theta_1G)\cdot G\ge0$. This is \eqref{eq:fourth-cumulants-H3} for $\mu'$ and $j=1$.

Let $j\in\{2,3\}$, and let $r\in O_{\mathrm{cub}}$ be the permutation of the coordinates that exchanges $e_1$ and $e_j$. Then $r\,\theta_j\,r^{-1}=\theta_1$, and $r$ maps $\{x_j>0\}$ onto $\{x_1>0\}$. Let $G\in A_j^+$. Since $\mathbb{U}_{0,r}\exp\left(iX(f)\right)=\exp\left(iX(U_{0,r}f)\right)$ and $\operatorname{supp}U_{0,r}f=r\left(\operatorname{supp}f\right)$, the function $G'=\mathbb{U}_{0,r}G$ lies in $A_1^+$; and $\mathbb{U}_{0,\theta_1}\mathbb{U}_{0,r}=\mathbb{U}_{0,\theta_1r}=\mathbb{U}_{0,r\theta_j}=\mathbb{U}_{0,r}\mathbb{U}_{0,\theta_j}$ gives $\Theta_1G'=\mathbb{U}_{0,r}\Theta_jG$. By \eqref{eq:fourth-cumulants-H2} for $\mu'$ (Step~3) and the case $j=1$,
\[
\int\mathrm{d}\mu'\,(\Theta_jG)\cdot G=\int\mathrm{d}\mu'\,\mathbb{U}_{0,r}\left[(\Theta_jG)\cdot G\right]=\int\mathrm{d}\mu'\,(\Theta_1G')\cdot G'\ge0 .
\]
Thus $\mu'$ satisfies \eqref{eq:fourth-cumulants-H3}.

\emph{Step 5: the invariant subspace of $L^2(\mu')$.} Let $L^2(W^3)$ be the space $L^2$ of $W^3$ with respect to the law of $z$, which is the product of three copies of the law of the telegraph process, and let $\operatorname{sh}_xz$, $x\in\mathbb{R}^3$, be as in Step~3. For $\boldsymbol{\alpha},\boldsymbol{\alpha}'\in L^2(W^3)$ and $L>0$ put
\[
B_L(\boldsymbol{\alpha},\boldsymbol{\alpha}')=L^{-3}\int_{[0,L]^3}\mathrm{d}x\,E_z\left[\boldsymbol{\alpha}(z)\cdot\boldsymbol{\alpha}'(\operatorname{sh}_xz)\right].
\]
The integrand is a measurable function of $x$, by Step~1 and Fubini's theorem, and it is bounded in absolute value by $\|\boldsymbol{\alpha}\|_2\cdot\|\boldsymbol{\alpha}'\|_2$, by the Cauchy--Schwarz inequality and because $\operatorname{sh}_xz$ has the law of $z$ (Step~3). Let first $\boldsymbol{\alpha}(z)=a_1(z_1)\cdot a_2(z_2)\cdot a_3(z_3)$ and $\boldsymbol{\alpha}'(z)=b_1(z_1)\cdot b_2(z_2)\cdot b_3(z_3)$ with $a_l,b_l\in L^2(W)$. The independence of $z_1$, $z_2$, $z_3$, Fubini's theorem and~\eqref{eq:fourth-cumulants-10}, applied to each $z_l$, give
\[
B_L(\boldsymbol{\alpha},\boldsymbol{\alpha}')=\prod_{l=1}^3\frac{1}{L}\int_0^L\mathrm{d}u\,E_z\left[a_l(z_l)\cdot b_l(\operatorname{sh}_uz_l)\right]\longrightarrow\prod_{l=1}^3E_z\left[a_l(z_l)\right]\cdot E_z\left[b_l(z_l)\right]=E_z\left[\boldsymbol{\alpha}(z)\right]\cdot E_z\left[\boldsymbol{\alpha}'(z)\right]
\]
as $L\to\infty$. For a product of finitely many probability spaces, the linear combinations of products of square integrable functions of the single factors are dense in $L^2$ of the product measure. Indeed, the closure of the set of these linear combinations is a closed linear subspace that contains the indicator functions of the measurable rectangles. The measurable sets whose indicator functions belong to this closure form a Dynkin system; as the product $\sigma$-algebra is generated by the $\pi$-system of the measurable rectangles, Dynkin's $\pi$-$\lambda$ theorem shows that this system is the whole product $\sigma$-algebra. So the closure contains all simple functions, which are dense in $L^2$. As the law of $z$ is the product of three copies of the law of the telegraph process, the linear combinations of the functions $a_1(z_1)\cdot a_2(z_2)\cdot a_3(z_3)$ with $a_1,a_2,a_3\in L^2(W)$ are dense in $L^2(W^3)$. Since $B_L$ and $(\boldsymbol{\alpha},\boldsymbol{\alpha}')\mapsto E_z\left[\boldsymbol{\alpha}(z)\right]\cdot E_z\left[\boldsymbol{\alpha}'(z)\right]$ are bilinear and bounded by $\|\boldsymbol{\alpha}\|_2\cdot\|\boldsymbol{\alpha}'\|_2$, uniformly in $L$, the approximation argument that concludes the proof of~\eqref{eq:fourth-cumulants-10} yields
\[
\lim_{L\to\infty}B_L(\boldsymbol{\alpha},\boldsymbol{\alpha}')=E_z\left[\boldsymbol{\alpha}(z)\right]\cdot E_z\left[\boldsymbol{\alpha}'(z)\right]\qquad\left(\boldsymbol{\alpha},\boldsymbol{\alpha}'\in L^2(W^3)\right).
\]

Let $f,g\in\mathcal D$, and put $\boldsymbol{\alpha}_f(z)=\exp\left(-\frac{1}{2}\langle f,v_zf\rangle\right)$ and $\boldsymbol{\alpha}_g(z)=\exp\left(-\frac{1}{2}\langle g,v_zg\rangle\right)$, functions on $W^3$ with values in $(0,1]$. Choose $c(f,g)>0$ such that the supports of $f$ and of $g$ lie in $\left[-\frac{1}{2}c(f,g),\frac{1}{2}c(f,g)\right]^3$, and write $|x|_\infty=\max_l|x_l|$ for $x\in\mathbb{R}^3$. If $|x|_\infty>c(f,g)$, the supports of $f$ and $\tau_xg$ are disjoint; as $\langle\tau_xg,v_z\,\tau_xg\rangle=\int\mathrm{d}y\,v_z(y+x)\cdot|g(y)|^2=\langle g,v_{\operatorname{sh}_xz}g\rangle$, the definition of $\mu'$ then gives
\[
\int\mathrm{d}\mu'\,\exp\left(-iX(f)+iX(\tau_xg)\right)=E_z\exp\left(-\frac{1}{2}\langle f,v_zf\rangle-\frac{1}{2}\langle\tau_xg,v_z\,\tau_xg\rangle\right)=E_z\left[\boldsymbol{\alpha}_f(z)\cdot\boldsymbol{\alpha}_g(\operatorname{sh}_xz)\right].
\]
The points $x\in[0,L]^3$ with $|x|_\infty\le c(f,g)$ form a set of measure at most $c(f,g)^3$, and for every $x\in\mathbb{R}^3$ the two outer members of the last display are bounded by $1$ in absolute value.

Let $I'$, $H'_{\mathrm{inv}}$ and $A_L$ be formed with $\mu'$ as in Lemma~\ref{lem:averages-translations-2}, which applies since $\mu'$ is translation invariant (Step~3), and put $P'=E'[\,\cdot\mid I']$. Let $G=\exp\left(iX(f)\right)$ and $H=\exp\left(iX(g)\right)$. Since $\mathbb{U}_xX(g)=X(\tau_xg)$, one has $\langle G,\mathbb{U}_xH\rangle=E'\left[\exp\left(-iX(f)+iX(\tau_xg)\right)\right]$; moreover $(A_LH)(\omega)=L^{-3}\int_{[0,L]^3}\mathrm{d}x\,(\mathbb{U}_xH)(\omega)$, as $\mathbb{U}_xH=H\circ T_{-x}$. Lemma~\ref{lem:averages-translations-2}(c), Fubini's theorem and the two preceding paragraphs give
\begin{align*}
\langle G,P'H\rangle&=\lim_{L\to\infty}\langle G,A_LH\rangle=\lim_{L\to\infty}L^{-3}\int_{[0,L]^3}\mathrm{d}x\,\langle G,\mathbb{U}_xH\rangle\\
&=\lim_{L\to\infty}B_L(\boldsymbol{\alpha}_f,\boldsymbol{\alpha}_g)=E_z\left[\boldsymbol{\alpha}_f(z)\right]\cdot E_z\left[\boldsymbol{\alpha}_g(z)\right]=\langle G,1\rangle\cdot\langle1,H\rangle .
\end{align*}
In the third equality the two integrands over $[0,L]^3$ differ only on a set of measure at most $c(f,g)^3$, on which both are bounded by $1$ in absolute value, so that the two averages differ by at most $2c(f,g)^3\cdot L^{-3}$; the last equality holds because $\langle G,1\rangle=E'\exp\left(-iX(f)\right)=\Phi'(-f)=E_z\left[\boldsymbol{\alpha}_f(z)\right]$ and $\langle1,H\rangle=E'\exp\left(iX(g)\right)=\Phi'(g)=E_z\left[\boldsymbol{\alpha}_g(z)\right]$.

Both sides of the identity $\langle G,P'H\rangle=\langle G,1\rangle\cdot\langle1,H\rangle$ are antilinear in $G$, linear in $H$, and continuous on $L^2(\mu')\times L^2(\mu')$. Hence the identity holds for all $G$ and $H$ in the complex linear span of the functions $\exp\left(iX(f)\right)$, $f\in\mathcal D$. This span is dense in $L^2(\mu')$ by Lemma~\ref{lem:exponentials-determination-law}(b), applied with $V_0=\mathcal D$, which is sequentially dense in $\mathcal S$; so the identity holds for all $G,H\in L^2(\mu')$, and $P'H=\langle1,H\rangle\cdot1$ for every $H\in L^2(\mu')$. By Lemma~\ref{lem:averages-translations-2}(b), $P'$ is the orthogonal projection of $L^2(\mu')$ onto $H'_{\mathrm{inv}}$. Therefore $H'_{\mathrm{inv}}$ consists of the constant functions, and $\mu'$ is ergodic in the sense of Definition~\ref{def:ergodicity-ergodic-decompositions}.

\emph{Step 6: moments of order two and four.} Let $f,g\in\mathcal S$, and take $n=2$ and $(f_1,f_2)=(f,g)$ in Step~2. The matrix $K_z$ then has the diagonal entries $\langle f,v_zf\rangle\in\left[0,\|f\|_2^2\right]$ and $\langle g,v_zg\rangle\in\left[0,\|g\|_2^2\right]$, and both of its off-diagonal entries equal $\langle f,v_zg\rangle$. Let $\varphi\colon\mathbb{R}^2\to\mathbb{R}$ be a Borel function with $E_z\int\mathrm{d}\lambda_z(y)\,|\varphi(y)|<\infty$. The identity at the end of Step~2, applied to $|\varphi|$ and to the positive and the negative part of $\varphi$, shows that $\varphi\left(X(f),X(g)\right)$ is $\mu'$-integrable and that
\[
E'\left[\varphi\left(X(f),X(g)\right)\right]=E_z\int\mathrm{d}\lambda_z(y)\,\varphi(y).
\]
For every $z\in W^3$ the Cauchy--Schwarz inequality gives
\begin{gather*}
\int\mathrm{d}\lambda_z(y)\,|y_1|\le\langle f,v_zf\rangle^{\frac{1}{2}}\le\|f\|_2,\\
\int\mathrm{d}\lambda_z(y)\,|y_1\cdot y_2|\le\langle f,v_zf\rangle^{\frac{1}{2}}\cdot\langle g,v_zg\rangle^{\frac{1}{2}}\le\|f\|_2\cdot\|g\|_2 ,
\end{gather*}
so that this applies to $\varphi(y)=y_1$ and to $\varphi(y)=y_1\cdot y_2$, whose integrals with respect to $\lambda_z$ are $0$ and $\langle f,v_zg\rangle$. The function $y\mapsto y_1^4$ is non-negative, so that the identity at the end of Step~2 applies to it directly, and Wick's formula gives $\int\mathrm{d}\lambda_z(y)\,y_1^4=3\langle f,v_zf\rangle^2$. Integration over $z$ therefore yields
\[
E'X(f)=0,\qquad E'\left[X(f)\cdot X(g)\right]=E_z\langle f,v_zg\rangle,\qquad E'\left[X(f)^4\right]=3E_z\left[\langle f,v_zf\rangle^2\right]\le3\|f\|_2^4 .
\]
So the quantities of Section~\ref{sec:cumulants-gaussian-laws} are defined for $\mu'$. By~\eqref{eq:fourth-cumulants-7} every $z_l(t)$ has mean $0$; by the independence of $z_1$, $z_2$, $z_3$ it follows that $E_z\left[z_1(x_1)\cdot z_2(x_2)\cdot z_3(x_3)\right]=0$, that is, $E_z\left[v_z(x)\right]=p_0$ for every $x\in\mathbb{R}^3$. Fubini's theorem therefore gives
\[
S_2^{\mu'}(f,g)=E_z\langle f,v_zg\rangle=\int\mathrm{d}x\,E_z\left[v_z(x)\right]\cdot\sum_{a<b}f_{ab}(x)\cdot g_{ab}(x)=p_0\cdot\langle f,g\rangle .
\]
Let $\operatorname{Var}_z$ denote the variance with respect to $E_z$. Then
\[
\kappa_4^{\mu'}(f)=E'\left[X(f)^4\right]-3\left(E'\left[X(f)^2\right]\right)^2=3E_z\left[\langle f,v_zf\rangle^2\right]-3\left(E_z\langle f,v_zf\rangle\right)^2=3\operatorname{Var}_z\left(\langle f,v_zf\rangle\right).
\]
The variable
\[
\langle f,v_zf\rangle-p_0\cdot\|f\|_2^2=q_0\int\mathrm{d}x\,z_1(x_1)\cdot z_2(x_2)\cdot z_3(x_3)\cdot|f(x)|^2
\]
has mean $0$. By Fubini's theorem and the independence of $z_1$, $z_2$, $z_3$,
\[
\operatorname{Var}_z\left(\langle f,v_zf\rangle\right)=q_0^2\int\mathrm{d}x\int\mathrm{d}y\,|f(x)|^2\cdot|f(y)|^2\cdot\prod_{l=1}^3E_z\left[z_l(x_l)\cdot z_l(y_l)\right],
\]
and by the consequence of~\eqref{eq:fourth-cumulants-9} derived in Step~1, $E_z\left[z_l(x_l)\cdot z_l(y_l)\right]=\exp\left(-|x_l-y_l|\right)$ for $l=1,2,3$. This proves the formula for $\kappa_4^{\mu'}(f)$.

Let $f\ne0$, and choose $x_0\in\mathbb{R}^3$ with $f(x_0)\ne0$. The integrand of the double integral is continuous and non-negative, and it is positive at $(x_0,x_0)$; since $q_0>0$, it follows that $\kappa_4^{\mu'}(f)>0$. Finally, $\mu'$ satisfies \eqref{eq:fourth-cumulants-H1} (Step~3); if it were Gaussian, Section~\ref{sec:cumulants-gaussian-laws} would give $\kappa_4^{\mu'}(f)=0$ for every $f\in\mathcal S$. Hence $\mu'$ is not Gaussian.
\end{proof}

\begin{proposition}[Gaussian components and four-point data]\label{prop:gaussian-components-four}
Let $\mu$, $(\Xi,\nu,(\mu_\xi))$ and $\Xi_0$ be as in Section~\ref{sec:cumulant-identities-non}, and let $V$ be the function of Theorem~\ref{thm:fourth-cumulant-mixture}. For $f\in\mathcal S$, an integer $k\ge1$ and $\xi\in\Xi_0$ put
\[
m_k^\xi=E_\xi\left[X(f)^k\right],\qquad V_k(f)=\operatorname{Var}_\nu\left(m_k^\xi\right),\qquad\delta_V=\kappa_4^\mu(f)-3V(f);
\]
thus $m_2^\xi=C_\xi(f,f)$ and $V_2(f)=V(f)$. By~\eqref{eq:fourth-cumulants-4},
\[
\delta_V=\int\mathrm{d}\nu(\xi)\,\kappa_4^\xi(f),
\]
so that $\delta_V\ne0$ if and only if $\int\mathrm{d}\nu(\xi)\,\kappa_4^\xi(f)\ne0$. If $\mu$ satisfies~\eqref{eq:fourth-cumulants-H2} and~\eqref{eq:fourth-cumulants-H3}, then for $j\in\{1,2,3\}$, $f\in\mathcal D$ with $\operatorname{supp}f\subset\{x_j>0\}$ and $t\ge0$ put
\[
D_2(t)=\operatorname{Cov}_\mu\left(X(\theta_jf)^2,X(\tau_{te_j}f)^2\right),\qquad\delta_D(t)=\kappa_4^\mu(f)-3D_2(t),
\]
where $\operatorname{Cov}_\mu(Y,Z)=E[Y\cdot Z]-E[Y]\cdot E[Z]$ for real $Y,Z\in L^2(\mu)$. Finally, for $f\in\mathcal S$ and $x\in\mathbb{R}^3$, let $\delta_{\mathrm{Ext}}=\mathrm{Ext}^\mu_x(f)$ be the number of Theorem~\ref{thm:non-gaussian-components}.
\begin{enumerate}
\item[(i)] If $\mu_\xi$ is Gaussian for $\nu$-almost every $\xi$, then
\[
\kappa_4^\mu(f)=3V(f)\ge0,\qquad\delta_V=0,\qquad\mathrm{Ext}^\mu_x(f)=-6\operatorname{Var}_\nu\left(C_\xi(f,\tau_xf)\right)\le0
\]
for all $f\in\mathcal S$ and $x\in\mathbb{R}^3$; if, in addition, $\mu$ satisfies~\eqref{eq:fourth-cumulants-H2} and~\eqref{eq:fourth-cumulants-H3}, then $\delta_D(t)\le0$ for all $j\in\{1,2,3\}$, all $f\in\mathcal D$ with $\operatorname{supp}f\subset\{x_j>0\}$ and all $t\ge0$. Consequently each of the following three conditions implies that the set of the $\xi\in\Xi_0$ for which $\mu_\xi$ is not Gaussian is contained in no $\nu$-null set: $\delta_{\mathrm{Ext}}>0$ for some $f\in\mathcal S$ and some $x\in\mathbb{R}^3$; $\delta_V\ne0$ for some $f\in\mathcal S$; $\mu$ satisfies~\eqref{eq:fourth-cumulants-H2} and~\eqref{eq:fourth-cumulants-H3}, and $\delta_D(t)>0$ for some $j\in\{1,2,3\}$, some $f\in\mathcal D$ with $\operatorname{supp}f\subset\{x_j>0\}$ and some $t\ge0$. None of these three conditions is satisfied by any of the laws $\mu^\rho$ of Theorem~\ref{thm:properties-scale-mixtures}.
\item[(ii)] Let $\mu'$ be an ergodic law on $(\Omega,\mathcal F)$, with expectation $E'$, which satisfies~\eqref{eq:fourth-cumulants-H1}, \eqref{eq:fourth-cumulants-H2} and~\eqref{eq:fourth-cumulants-H3} in place of $\mu$ and is not Gaussian; let $C'=S_2^{\mu'}$, and let $0\le p<1$. Then $C'$ satisfies the hypotheses of Lemma~\ref{lem:gaussian-laws-bounded}, and the law $\gamma_{C'}$ of that lemma satisfies~\eqref{eq:fourth-cumulants-H1}--\eqref{eq:fourth-cumulants-H3} and is ergodic. The law
\[
\mu=(1-p)\cdot\mu'+p\cdot\gamma_{C'}
\]
satisfies~\eqref{eq:fourth-cumulants-H1}--\eqref{eq:fourth-cumulants-H3}. Let $\Xi=\{1,2\}$, with the $\sigma$-algebra of all its subsets and the probability measure $\nu$ given by $\nu(\{1\})=1-p$ and $\nu(\{2\})=p$, and let $\mu_1=\mu'$ and $\mu_2=\gamma_{C'}$. Then $(\Xi,\nu,(\mu_\xi))$ is an ergodic decomposition of $\mu$, and for this decomposition
\[
S_2^\mu=S_2^{\mu'}=C',\qquad\kappa_4^\mu=(1-p)\cdot\kappa_4^{\mu'},\qquad\Gamma=0,\qquad V=0,\qquad\delta_V=\kappa_4^\mu(f)\quad(f\in\mathcal S),
\]
the second identity being an identity of $4$-linear forms on $\mathcal S$. Consequently, for every $p\in[0,1)$ the law $\mu$ has an ergodic decomposition with the following three properties. All its components have the same two-point function: $C_\xi=C'$ for $\xi\in\Xi$. For all $f_1,f_2,f_3,f_4\in\mathcal S$ and $x\in\mathbb{R}^3$,
\[
\left|\kappa_4^\mu(f_1,f_2,\tau_xf_3,\tau_xf_4)\right|\le\left|\kappa_4^{\mu'}(f_1,f_2,\tau_xf_3,\tau_xf_4)\right|,
\]
so that $\kappa_4^\mu(f_1,f_2,\tau_xf_3,\tau_xf_4)\to0$ as $|x|\to\infty$ for all $f_1,f_2,f_3,f_4\in\mathcal S$ whenever the same holds for $\kappa_4^{\mu'}$. The Gaussian components form a set of $\nu$-measure $p$:
\[
\nu\left(\left\{\xi\in\Xi:\ \mu_\xi\text{ is Gaussian}\right\}\right)=p.
\]
Moreover, for every $f\in\mathcal S$,
\[
V_4(f)=p\cdot(1-p)\cdot\kappa_4^{\mu'}(f)^2,\qquad\nu\left(\left\{\xi\in\Xi:\ \mu_\xi\text{ is not Gaussian}\right\}\right)=1-p,
\]
and, if $\kappa_4^{\mu'}(f)\ne0$,
\[
1-p=\frac{\kappa_4^\mu(f)}{\kappa_4^{\mu'}(f)}.
\]
Thus, when $\mu'$ and a test $f$ with $\kappa_4^{\mu'}(f)\ne0$ are fixed and $p$ varies in $[0,1)$, the $\nu$-measure of the set of the non-Gaussian components is proportional to $\kappa_4^\mu(f)$; the lower bound of Theorem~\ref{thm:non-gaussian-components} for the measure of a set of non-Gaussian components depends on $\delta_{\mathrm{Ext}}$ in the same linear way.
\item[(iii)] Laws $\mu'$ with the properties required in (ii) exist: the law $\mu'$ of Lemma~\ref{lem:white-noise-telegraph} is such a law, and it satisfies $\kappa_4^{\mu'}(f)>0$ for every $f\in\mathcal S$ with $f\ne0$.
\end{enumerate}
\end{proposition}

\begin{proof}
\emph{The definitions.} Let $f\in\mathcal S$ and let $k\ge1$ be an integer. For $\xi\in\Xi_0$ the law $\mu_\xi$ satisfies~\eqref{eq:fourth-cumulants-H1}, by Corollary~\ref{cor:properties-almost-every}; Lemma~\ref{lem:moment-bounds}(e), applied to $\mu_\xi$ with $n=k$ and $f_1=\dots=f_k=f$, gives
\[
\left|m_k^\xi\right|\le c_{2k}^{\frac{1}{2}}\cdot\|f\|_2^k\qquad(\xi\in\Xi_0).
\]
The function $X(f)^k$ is $\mu$-integrable, by Lemma~\ref{lem:moment-bounds}(c) and the Cauchy--Schwarz inequality; as in the proof of Theorem~\ref{thm:fourth-cumulant-mixture}(a), the extension of~\eqref{eq:fourth-cumulants-ED1} to $\mu$-integrable functions shows that $\xi\mapsto m_k^\xi$ is measurable. Hence $m_k^\xi$ is a bounded measurable function of $\xi\in\Xi_0$, and $V_k(f)$ is defined. Since $m_2^\xi=E_\xi\left[X(f)\cdot X(f)\right]=C_\xi(f,f)$, one has $V_2(f)=\operatorname{Var}_\nu\left(C_\xi(f,f)\right)=V(f)$, and~\eqref{eq:fourth-cumulants-4} of Theorem~\ref{thm:fourth-cumulant-mixture}(c) reads $\delta_V=\kappa_4^\mu(f)-3V(f)=\int\mathrm{d}\nu(\xi)\,\kappa_4^\xi(f)$. If $\mu$ satisfies~\eqref{eq:fourth-cumulants-H2} and~\eqref{eq:fourth-cumulants-H3}, the functions $X(\theta_jf)^2$ and $X(\tau_{te_j}f)^2$ belong to $L^2(\mu)$ by Lemma~\ref{lem:moment-bounds}(c) with $k=2$, so that $D_2(t)$ is defined.

(i) Suppose that $\mu_\xi$ is Gaussian for $\nu$-almost every $\xi$, and let $N\subset\Xi$ be a $\nu$-null set such that $\mu_\xi$ is Gaussian for every $\xi\in\Xi_0\setminus N$. For such $\xi$ the law $\mu_\xi$ satisfies~\eqref{eq:fourth-cumulants-H1} and has therefore mean zero (Lemma~\ref{lem:moment-bounds}(a)); by Section~\ref{sec:cumulants-gaussian-laws}, $\kappa_4^\xi(f_1,f_2,f_3,f_4)=0$ for all $f_1,f_2,f_3,f_4\in\mathcal S$, and $\mathrm{Ext}^\xi_x(f)=0$ for all $f\in\mathcal S$ and $x\in\mathbb{R}^3$. As $\Xi_0\setminus N$ has full $\nu$-measure,
\[
\int\mathrm{d}\nu(\xi)\,\kappa_4^\xi(f)=0,\qquad\int\mathrm{d}\nu(\xi)\,\mathrm{Ext}^\xi_x(f)=0\qquad(f\in\mathcal S,\ x\in\mathbb{R}^3).
\]
Now~\eqref{eq:fourth-cumulants-4} gives $\kappa_4^\mu(f)=3V(f)$, which is non-negative because $V(f)$ is a variance, and $\delta_V=\kappa_4^\mu(f)-3V(f)=0$. Formula~\eqref{eq:fourth-cumulants-6} becomes
\[
0=\mathrm{Ext}^\mu_x(f)+6\operatorname{Var}_\nu\left(C_\xi(f,\tau_xf)\right),
\]
which is the asserted expression for $\mathrm{Ext}^\mu_x(f)$; it is non-positive, a variance being non-negative.

Next, let $\mu$ be an arbitrary law as in Section~\ref{sec:cumulant-identities-non} which satisfies~\eqref{eq:fourth-cumulants-H2} and~\eqref{eq:fourth-cumulants-H3}, and let $j\in\{1,2,3\}$, $f\in\mathcal D$ with $\operatorname{supp}f\subset\{x_j>0\}$, and $t\ge0$. The function $G=X(f)^2$ is real, it is $\Sigma_j^+$-measurable by the definition of $\Sigma_j^+$, and it is square integrable under $\mu$ by Lemma~\ref{lem:moment-bounds}(c). Each operator $\mathbb{U}_g$ is the composition with a map of $\Omega$, hence multiplicative, and $\mathbb{U}_gX(h)=X(U_gh)$ for $h\in\mathcal S$ (Section~\ref{sec:test-forms-distributions-2}). Since $X(\theta_jf)$ is real, it follows that
\[
\Theta_jG=\overline{\left(\mathbb{U}_{0,\theta_j}X(f)\right)^2}=X(\theta_jf)^2,\qquad\mathbb{U}_{te_j}G=\left(\mathbb{U}_{te_j}X(f)\right)^2=X(\tau_{te_j}f)^2 .
\]
So Lemma~\ref{lem:consequences-reflection-positivity}(c) applies to $G$, and~\eqref{eq:fourth-cumulants-2} gives
\[
\operatorname{Var}\left(E[G\mid I]\right)\le\operatorname{Cov}\left(\Theta_jG,\mathbb{U}_{te_j}G\right)=D_2(t).
\]
By Proposition~\ref{prop:existence-conditional-expectations}(b), the function $\xi\mapsto E_\xi[G]=C_\xi(f,f)$ has under $\nu$ the same distribution as $E[G\mid I]$ under $\mu$; both are square integrable, so the left side equals $\operatorname{Var}_\nu\left(C_\xi(f,f)\right)=V(f)$. Hence
\[
V(f)\le D_2(t).
\]
If $\mu_\xi$ is Gaussian for $\nu$-almost every $\xi$, the first part of (i) gives $\kappa_4^\mu(f)=3V(f)\le3D_2(t)$, that is, $\delta_D(t)\le0$.

Thus, if $\mu_\xi$ is Gaussian for $\nu$-almost every $\xi$, then $\delta_{\mathrm{Ext}}=\mathrm{Ext}^\mu_x(f)\le0$ for all $f\in\mathcal S$ and $x\in\mathbb{R}^3$, $\delta_V=0$ for all $f\in\mathcal S$, and, under~\eqref{eq:fourth-cumulants-H2} and~\eqref{eq:fourth-cumulants-H3}, $\delta_D(t)\le0$ for all $j$, $f$, $t$ as above. Each of the three conditions listed in (i) contradicts one of these conclusions. Hence, if one of the conditions holds, there is no $\nu$-null set $N$ such that $\mu_\xi$ is Gaussian for every $\xi\in\Xi_0\setminus N$, which is the assertion.

Let $\rho$ be a probability measure on $[0,1]$, and let $\mu^\rho$ be the corresponding law of Theorem~\ref{thm:properties-scale-mixtures}. By Theorem~\ref{thm:properties-scale-mixtures}(i) it satisfies~\eqref{eq:fourth-cumulants-H1}, \eqref{eq:fourth-cumulants-H2} and~\eqref{eq:fourth-cumulants-H3}; in particular it is translation invariant, and all of the above applies to $\mu^\rho$ in place of $\mu$. By Theorem~\ref{thm:properties-scale-mixtures}(ii), the interval $[0,1]$ with the measure $\rho$ and the laws $\gamma_s$, $s\in[0,1]$, of that theorem is an ergodic decomposition of $\mu^\rho$ all of whose components are Gaussian. The numbers $\delta_{\mathrm{Ext}}$ and $\delta_D(t)$ are defined in terms of $\mu^\rho$ alone. So is $\delta_V$: let $f\in\mathcal S$; for every ergodic decomposition of $\mu^\rho$, Proposition~\ref{prop:existence-conditional-expectations}(b), applied to the function $X(f)^2$, which belongs to $L^2(\mu^\rho)$ by Lemma~\ref{lem:moment-bounds}(c), shows that $V(f)$ is the variance of $E\left[X(f)^2\mid I\right]$ under $\mu^\rho$, where $I$ is formed with $\mu^\rho$. The first part of (i), applied to $\mu^\rho$ and to the ergodic decomposition of Theorem~\ref{thm:properties-scale-mixtures}(ii), therefore gives $\delta_{\mathrm{Ext}}\le0$, $\delta_V=0$ and $\delta_D(t)\le0$ for all $f$, $x$, $j$ and $t$ as above, and none of the three conditions holds for $\mu^\rho$.

(ii) \emph{The law $\gamma_{C'}$.} The law $\mu'$ satisfies~\eqref{eq:fourth-cumulants-H1}; by Lemma~\ref{lem:moment-bounds}(a), (b), (c) it has mean zero and finite moments, so that $C'=S_2^{\mu'}$ is defined. The form $C'$ is bilinear, because $X(f)$ depends linearly on $f$, and symmetric; it is positive semidefinite, since $C'(f,f)=E'\left[X(f)^2\right]\ge0$; and $C'(f,f)\le\|f\|_2^2$ for $f\in\mathcal S$ by Lemma~\ref{lem:moment-bounds}(b), applied to $\mu'$. These are the hypotheses of Lemma~\ref{lem:gaussian-laws-bounded}. By part (a) of that lemma, $\gamma_{C'}$ is a centred Gaussian law with covariance $C'$, that is, $S_2^{\gamma_{C'}}=C'$, and it satisfies~\eqref{eq:fourth-cumulants-H1}.

Let $g\in\mathbb{R}^3\rtimes O_{\mathrm{cub}}$ and $f,h\in\mathcal S$. Since $\mathbb{U}_g$ is the composition with $T_g^{-1}$, one has $\mathbb{U}_g\left(X(f)\cdot X(h)\right)=X(U_gf)\cdot X(U_gh)$. By~\eqref{eq:fourth-cumulants-H2} for $\mu'$, applied to the inverse of $g$ in the group $\mathbb{R}^3\rtimes O_{\mathrm{cub}}$, the image of $\mu'$ under $T_g^{-1}$ is $\mu'$; hence
\[
C'(U_gf,U_gh)=E'\left[\left(X(f)\cdot X(h)\right)\circ T_g^{-1}\right]=E'\left[X(f)\cdot X(h)\right]=C'(f,h).
\]
By Lemma~\ref{lem:gaussian-laws-bounded}(b), $\gamma_{C'}$ is invariant under $\mathbb{R}^3\rtimes O_{\mathrm{cub}}$, which is~\eqref{eq:fourth-cumulants-H2}. In particular $C'$ is translation invariant, and Lemma~\ref{lem:gaussian-laws-bounded}(c) shows that $\gamma_{C'}$ is mixing, hence ergodic. Now let $j\in\{1,2,3\}$ and $f\in\mathcal D$ with $\operatorname{supp}f\subset\{x_j>0\}$. The function $X(f)$ is real, $\Sigma_j^+$-measurable and square integrable under $\mu'$, and $\Theta_jX(f)=\overline{\mathbb{U}_{0,\theta_j}X(f)}=X(\theta_jf)$. Lemma~\ref{lem:consequences-reflection-positivity}(a), applied to the law $\mu'$, which satisfies~\eqref{eq:fourth-cumulants-H2} and~\eqref{eq:fourth-cumulants-H3}, gives
\[
C'(\theta_jf,f)=E'\left[\left(\Theta_jX(f)\right)\cdot X(f)\right]\ge0 .
\]
Moreover $C'(\theta_jf,\theta_jh)=C'(f,h)$ for all $f,h\in\mathcal S$, by the invariance of $C'$ under the motion $(0,\theta_j)\in\mathbb{R}^3\rtimes O_{\mathrm{cub}}$. Lemma~\ref{lem:gaussian-laws-bounded}(d) now shows that $\gamma_{C'}$ satisfies~\eqref{eq:fourth-cumulants-H3}.

\emph{The law $\mu$.} Each of the conditions~\eqref{eq:fourth-cumulants-H1}--\eqref{eq:fourth-cumulants-H3} is preserved under convex combinations of laws. Indeed, let $f\in\mathcal S$, $g\in\mathbb{R}^3\rtimes O_{\mathrm{cub}}$, $j\in\{1,2,3\}$ and $G\in A_j^+$; the function $(\Theta_jG)\cdot G$ is bounded and measurable, $G$ being a finite linear combination of functions $\exp\left(iX(h)\right)$, $h\in\mathcal D$. By~\eqref{eq:fourth-cumulants-H1}--\eqref{eq:fourth-cumulants-H3} for $\mu'$ and for $\gamma_{C'}$,
\begin{align*}
\int\mathrm{d}\mu\,\exp X(f)&=(1-p)\cdot\int\mathrm{d}\mu'\,\exp X(f)+p\cdot\int\mathrm{d}\gamma_{C'}\,\exp X(f)\le\exp\left(\frac{1}{2}\|f\|_2^2\right),\\
\mu\circ T_g^{-1}&=(1-p)\cdot\left(\mu'\circ T_g^{-1}\right)+p\cdot\left(\gamma_{C'}\circ T_g^{-1}\right)=\mu,\\
\int\mathrm{d}\mu\,(\Theta_jG)\cdot G&=(1-p)\cdot\int\mathrm{d}\mu'\,(\Theta_jG)\cdot G+p\cdot\int\mathrm{d}\gamma_{C'}\,(\Theta_jG)\cdot G\ge0 .
\end{align*}

\emph{The decomposition.} For $A\in\mathcal F$ the function $\xi\mapsto\mu_\xi(A)$ on $\Xi$ is measurable, every subset of $\Xi$ being measurable, and
\[
\int\mathrm{d}\nu(\xi)\,\mu_\xi(A)=(1-p)\cdot\mu'(A)+p\cdot\gamma_{C'}(A)=\mu(A);
\]
this is~\eqref{eq:fourth-cumulants-ED1}. The laws $\mu_1=\mu'$ and $\mu_2=\gamma_{C'}$ are both translation invariant and ergodic, which is~\eqref{eq:fourth-cumulants-ED2}. So $(\Xi,\nu,(\mu_\xi))$ is an ergodic decomposition of $\mu$ in the sense of Definition~\ref{def:ergodicity-ergodic-decompositions}. As both laws moreover satisfy~\eqref{eq:fourth-cumulants-H1}, the properties stated in Corollary~\ref{cor:properties-almost-every} hold for every $\xi\in\Xi$; one may take $\Xi_0=\Xi$, and Theorem~\ref{thm:fourth-cumulant-mixture}(d) applies to this family, with $\Xi$ in place of $\Xi_0$.

\emph{Two- and four-point functions.} One has $C_\xi=C'$ for $\xi=1$ by definition and for $\xi=2$ by Lemma~\ref{lem:gaussian-laws-bounded}(a). For $f,h\in\mathcal S$ the product $X(f)\cdot X(h)$ is integrable under $\mu'$ and under $\gamma_{C'}$, by Lemma~\ref{lem:moment-bounds}(b) and the Cauchy--Schwarz inequality, so that
\[
S_2^\mu(f,h)=(1-p)\cdot C'(f,h)+p\cdot C'(f,h)=C'(f,h).
\]
For $f_1,f_2\in\mathcal S$ the function $\xi\mapsto C_\xi(f_1,f_2)$ is constant on $\Xi$; hence every covariance $\Gamma(f_1,f_2;f_3,f_4)$ vanishes. In particular $V(f)=0$ and $\delta_V=\kappa_4^\mu(f)-3V(f)=\kappa_4^\mu(f)$ for $f\in\mathcal S$. The law $\gamma_{C'}$ is a centred Gaussian law, so that $\kappa_4^{\gamma_{C'}}=0$ by Section~\ref{sec:cumulants-gaussian-laws}. Theorem~\ref{thm:fourth-cumulant-mixture}(a), in the form given by part (d), now yields
\[
\kappa_4^\mu(f_1,f_2,f_3,f_4)=(1-p)\cdot\kappa_4^{\mu'}(f_1,f_2,f_3,f_4)+p\cdot0\qquad(f_1,f_2,f_3,f_4\in\mathcal S).
\]

\emph{The three properties.} The equality $C_\xi=C'$, $\xi\in\Xi$, is the first one. The identity of $4$-linear forms, evaluated at $(f_1,f_2,\tau_xf_3,\tau_xf_4)$, gives
\[
\left|\kappa_4^\mu(f_1,f_2,\tau_xf_3,\tau_xf_4)\right|=(1-p)\cdot\left|\kappa_4^{\mu'}(f_1,f_2,\tau_xf_3,\tau_xf_4)\right|,
\]
and $0<1-p\le1$; this is the second property. The law $\mu_2=\gamma_{C'}$ is Gaussian, and $\mu_1=\mu'$ is not. Therefore $\left\{\xi\in\Xi:\ \mu_\xi\text{ is Gaussian}\right\}=\{2\}$, a set of $\nu$-measure $p$, which is the third property; its complement $\{1\}$ has $\nu$-measure $1-p$.

\emph{The fourth moments.} Let $f\in\mathcal S$, and let $E_{\gamma_{C'}}$ denote the expectation with respect to $\gamma_{C'}$. By the definition of the fourth cumulant in Section~\ref{sec:cumulants-gaussian-laws}, and since $\kappa_4^{\gamma_{C'}}(f)=0$, the function $\xi\mapsto m_4^\xi$ takes the values
\[
E'\left[X(f)^4\right]=\kappa_4^{\mu'}(f)+3C'(f,f)^2\quad\text{at }\xi=1,\qquad E_{\gamma_{C'}}\left[X(f)^4\right]=3C'(f,f)^2\quad\text{at }\xi=2 .
\]
A real function $h$ on $\Xi$ has the variance
\[
\operatorname{Var}_\nu(h)=(1-p)\cdot h(1)^2+p\cdot h(2)^2-\left((1-p)\cdot h(1)+p\cdot h(2)\right)^2=p\cdot(1-p)\cdot\left(h(1)-h(2)\right)^2 ;
\]
for $h(\xi)=m_4^\xi$ the difference $h(1)-h(2)$ is $\kappa_4^{\mu'}(f)$, and $V_4(f)=p\cdot(1-p)\cdot\kappa_4^{\mu'}(f)^2$ follows. Finally $\kappa_4^\mu(f)=(1-p)\cdot\kappa_4^{\mu'}(f)$. If $\kappa_4^{\mu'}(f)\ne0$, division by $\kappa_4^{\mu'}(f)$ gives $1-p=\frac{\kappa_4^\mu(f)}{\kappa_4^{\mu'}(f)}$, and $1-p$ is the $\nu$-measure of the set of the non-Gaussian components; for fixed $\mu'$ and $f$ it is the multiple $\frac{1}{\kappa_4^{\mu'}(f)}\cdot\kappa_4^\mu(f)$ of $\kappa_4^\mu(f)$. The lower bound of Theorem~\ref{thm:non-gaussian-components} is
\[
\frac{\delta_{\mathrm{Ext}}}{2(c_4+6)\cdot\|f\|_2^4},
\]
a constant multiple of $\delta_{\mathrm{Ext}}$.

(iii) By Lemma~\ref{lem:white-noise-telegraph}, the law $\mu'$ constructed there satisfies~\eqref{eq:fourth-cumulants-H1}, \eqref{eq:fourth-cumulants-H2} and~\eqref{eq:fourth-cumulants-H3}, is ergodic and is not Gaussian; these are the properties required in (ii). The same lemma gives $\kappa_4^{\mu'}(f)>0$ for every $f\in\mathcal S$ with $f\ne0$.
\end{proof}

\subsection{Three separated tests}\label{sec:three-separated-tests}

\subsubsection{The free quartic coefficient}\label{sec:free-quartic-coefficient}

Let $m=\bar{m}>0$. Let $\sigma_1,\sigma_2,\sigma_3$ be the Pauli matrices, and let $\gamma_1,\gamma_2,\gamma_3$ be the Hermitian $4\times4$ matrices given in blocks of size $2$ by
\[
\gamma_k=\operatorname{diag}(\sigma_k,-\sigma_k)\qquad(k=1,2,3).
\]
They satisfy $\gamma_k\gamma_l+\gamma_l\gamma_k=2\delta_{kl}$ for $k,l\in\{1,2,3\}$, and after a renumbering of the coordinate directions they coincide with the matrices of Section~\ref{sec:dirac-matrices-lattice}. For $v=(v_1,v_2,v_3)$ write $\gamma\cdot v=\sum_{k=1}^3v_k\gamma_k$; in the same way $\gamma\cdot\partial=\sum_{k=1}^3\gamma_k\partial_k$, where $\partial_k$ is the partial derivative with respect to the $k$-th coordinate of $\mathbb{R}^3$. On the Hilbert space $L^2(\mathbb{R}^3;\mathbb{C}^4)$ put
\[
S_m=(m+\gamma\cdot\partial)^{-1}.
\]
For a real vector field $a\in C_c^\infty(\mathbb{R}^3;\mathbb{R}^3)$, the notation $i\gamma\cdot a$ also stands for the operator on $L^2(\mathbb{R}^3;\mathbb{C}^4)$ of multiplication by the matrix-valued function $x\mapsto i\gamma\cdot a(x)$, and
\[
T_a=S_m\,(i\gamma\cdot a),\qquad\mathsf{b}(a)=-6\operatorname{Tr}\left[T_a^4\right];
\]
the trace exists by Lemma~\ref{lem:polarisation-decay-quartic}(a).

Let $\beta\in C_c^\infty(\mathbb{R}^3;\Lambda^2)$ be a real two-form; its components $\beta_{kl}$, $1\le k<l\le3$, are extended to all $k,l\in\{1,2,3\}$ by $\beta_{lk}=-\beta_{kl}$. Put
\[
a=d^*\beta,\qquad f=dd^*\beta,\qquad b(f)=\mathsf{b}(a),
\]
where $d$ is the exterior derivative and $d^*$ is its formal adjoint; in components,
\[
a_l=-\sum_{k=1}^3\partial_k\beta_{kl}\quad(l=1,2,3),\qquad f_{kl}=\partial_ka_l-\partial_la_k\quad(1\le k<l\le3).
\]
Thus $a\in C_c^\infty(\mathbb{R}^3;\mathbb{R}^3)$ is real and $f=da$ lies in $\mathcal D$. The number $b(f)$ coincides with the number $b(f)$ of Section~\ref{sec:free-fourth-order}. It depends on $f$ alone, because $f$ determines $a$: if $\beta'\in C_c^\infty(\mathbb{R}^3;\Lambda^2)$ is real and $dd^*\beta'=f$, then $w=d^*\beta-d^*\beta'$ is a smooth vector field with compact support such that $dw=0$ and $d^*w=0$, the latter because $d^*d^*=0$; consequently $\Delta w=-(dd^*+d^*d)w=0$ componentwise, and $w$, being harmonic with compact support, vanishes.

Let $a$ be a vector field on $\mathbb{R}^3$. For $x\in\mathbb{R}^3$ the Euclidean length of the vector $a(x)\in\mathbb{R}^3$ is written $|a(x)|$, and $\|a\|_\infty=\sup_{x\in\mathbb{R}^3}|a(x)|$; the translate $\tau_xa$ is the vector field $y\mapsto a(y-x)$ on $\mathbb{R}^3$; and $|\operatorname{supp}a|$ stands for the three-dimensional Lebesgue measure of the closed set $\operatorname{supp}a$. Likewise $|x_0|$ denotes the Euclidean length of a point $x_0\in\mathbb{R}^3$.

For $1\le p<\infty$ let $\mathfrak S_p$ be the Schatten class of the compact operators on $L^2(\mathbb{R}^3;\mathbb{C}^4)$ whose singular values form a sequence in $\ell^p$. Thus $\mathfrak S_1$ is the trace class, on which $\operatorname{Tr}$ denotes the trace, and $\mathfrak S_2$ is the class of Hilbert--Schmidt operators, with the norm $\|\cdot\|_{\mathrm{HS}}$. The operator norm of a bounded operator, and of a $4\times4$ matrix, is written $\|\cdot\|$, and $\operatorname{tr}$ is the trace of a $4\times4$ matrix.

\begin{lemma}[Polarisation and decay of the quartic trace]\label{lem:polarisation-decay-quartic}
Let $a,a'\in C_c^\infty(\mathbb{R}^3;\mathbb{R}^3)$ be real vector fields.
\begin{enumerate}
\item[(a)] The operator $T_a$ belongs to $\mathfrak S_4$, and $\mathsf{b}(\tau_xa)=\mathsf{b}(a)$ for every $x\in\mathbb{R}^3$. With
\[
\mathsf{b}(a,a,a',a')=-4\operatorname{Tr}\left[T_aT_aT_{a'}T_{a'}\right]-2\operatorname{Tr}\left[T_aT_{a'}T_aT_{a'}\right]
\]
one has
\[
\mathsf{b}(a+a')+\mathsf{b}(a-a')=2\,\mathsf{b}(a)+2\,\mathsf{b}(a')+12\,\mathsf{b}(a,a,a',a').
\]
Moreover, $\mathsf{b}(a)\in\mathbb{R}$ and $\mathsf{b}(a,a,a',a')\in\mathbb{R}$.
\item[(b)] If $\mathfrak{d}=\operatorname{dist}(\operatorname{supp}a,\operatorname{supp}a')>0$, then
\[
\left|\mathsf{b}(a,a,a',a')\right|\le\frac{6}{m^2}\cdot\|a\|_\infty^2\cdot\|a'\|_\infty^2\cdot|\operatorname{supp}a|\cdot|\operatorname{supp}a'|\cdot\left(m^2+\left(m+\frac{1}{\mathfrak{d}}\right)^2\right)\cdot\frac{\exp(-2m\cdot\mathfrak{d})}{4\pi^2\cdot\mathfrak{d}^2}.
\]
\end{enumerate}
\end{lemma}

\begin{proof}
(a) Under the Fourier transform the operator $\partial_k$ becomes the multiplication by $ip_k$, $p\in\mathbb{R}^3$, and the anticommutation relations give $(m+i\gamma\cdot p)(m-i\gamma\cdot p)=m^2+|p|^2$. Hence $S_m$ is the Fourier multiplier with the matrix-valued symbol
\[
p\longmapsto\frac{m-i\gamma\cdot p}{m^2+|p|^2}.
\]
Since $(m-i\gamma\cdot p)^*(m-i\gamma\cdot p)=m^2+|p|^2$, the norm of this matrix equals $(m^2+|p|^2)^{-\frac{1}{2}}$. So each of its entries is bounded in absolute value by $(m^2+|p|^2)^{-\frac{1}{2}}$, a function of $p$ that lies in $L^4(\mathbb{R}^3)$, as $\int\mathrm{d}p\,(m^2+|p|^2)^{-2}=\frac{\pi^2}{m}$. The entries of the matrix-valued function $i\gamma\cdot a$ are bounded and have compact support, so they lie in $L^4(\mathbb{R}^3)$ as well. By the Kato--Seiler--Simon inequality with exponent $4$, the operator $u(x)\,v(-i\nabla)$ on $L^2(\mathbb{R}^3)$, that is, the multiplication by $u$ composed with the Fourier multiplier with symbol $v$, belongs to the Schatten class with exponent $4$ of $L^2(\mathbb{R}^3)$ whenever $u,v\in L^4(\mathbb{R}^3)$. With respect to the four components of $\mathbb{C}^4$, each of the sixteen entries of $T_a$ is a sum of four adjoints of operators of this form. Hence $T_a\in\mathfrak S_4$.

By H\"older's inequality for Schatten classes, a product of four operators of $\mathfrak S_4$ is of trace class, and a product of three lies in $\mathfrak S_{\frac{4}{3}}$; moreover $\operatorname{Tr}AB=\operatorname{Tr}BA$ for $A\in\mathfrak S_4$ and $B\in\mathfrak S_{\frac{4}{3}}$. Hence the traces in the lemma exist, and the trace of a product of four of the operators $T_a$, $T_{a'}$ does not change when the factors are rotated cyclically.

Let $x\in\mathbb{R}^3$. The translation $\psi\mapsto\psi(\cdot-x)$ is a unitary operator on $L^2(\mathbb{R}^3;\mathbb{C}^4)$; it commutes with the Fourier multiplier $S_m$ and conjugates the multiplication by $i\gamma\cdot a$ into the multiplication by $i\gamma\cdot(\tau_xa)$. Hence $T_{\tau_xa}$ is unitarily equivalent to $T_a$, and $\mathsf{b}(\tau_xa)=\mathsf{b}(a)$.

The map $a\mapsto T_a$ is linear, so that $T_{a\pm a'}=T_a\pm T_{a'}$. Expanding $(T_a+T_{a'})^4$ and $(T_a-T_{a'})^4$ gives sums over the sixteen words of length four in the letters $T_a$ and $T_{a'}$, each word carrying the sign $+1$, resp.\ $(-1)^j$, where $j$ is the number of letters $T_{a'}$ in it. In the sum of the two expansions the words with odd $j$ cancel, and every other word appears twice. These words are $T_a^4$, $T_{a'}^4$ and the six words with two letters of each kind. Call the word obtained by moving the first $k$ letters of a word, in their order, to its end the cyclic shift of the word by $k$ places. Among the six words just named, $T_aT_aT_{a'}T_{a'}$, $T_aT_{a'}T_{a'}T_a$, $T_{a'}T_{a'}T_aT_a$ and $T_{a'}T_aT_aT_{a'}$ are the cyclic shifts of $T_aT_aT_{a'}T_{a'}$ by $0$, $1$, $2$ and $3$ places, and the remaining two, $T_aT_{a'}T_aT_{a'}$ and $T_{a'}T_aT_{a'}T_a$, are the shifts of $T_aT_{a'}T_aT_{a'}$ by $0$ and $1$ place. The cyclicity of the trace therefore gives
\begin{multline*}
\operatorname{Tr}\left[(T_a+T_{a'})^4\right]+\operatorname{Tr}\left[(T_a-T_{a'})^4\right]\\
=2\operatorname{Tr}\left[T_a^4\right]+2\operatorname{Tr}\left[T_{a'}^4\right]+8\operatorname{Tr}\left[T_aT_aT_{a'}T_{a'}\right]+4\operatorname{Tr}\left[T_aT_{a'}T_aT_{a'}\right],
\end{multline*}
and multiplication by $-6$ yields the identity of (a).

It remains to prove that $\mathsf{b}(a)$ and $\mathsf{b}(a,a,a',a')$ are real. The matrices $\gamma_k$ are Hermitian, and the operators $\partial_k$ are anti-selfadjoint and commute with them; hence $\gamma\cdot\partial$ is anti-selfadjoint and $S_m^*=(m-\gamma\cdot\partial)^{-1}$. The multiplication by $\gamma\cdot a$ is selfadjoint, since $a$ is real; therefore
\[
T_a^*=-(i\gamma\cdot a)\,S_m^*.
\]
Let $\hat{\Gamma}$ be the $4\times4$ matrix that exchanges the first two coordinates of $\mathbb{C}^4$ with the last two, $\hat{\Gamma}(z_1,z_2,z_3,z_4)=(z_3,z_4,z_1,z_2)$. It is Hermitian and unitary, and $\hat{\Gamma}\gamma_k\hat{\Gamma}=\operatorname{diag}(-\sigma_k,\sigma_k)=-\gamma_k$ for $k=1,2,3$. Consequently $\hat{\Gamma}(m-\gamma\cdot\partial)\hat{\Gamma}=m+\gamma\cdot\partial$, whence $\hat{\Gamma} S_m^*\hat{\Gamma}=S_m$, and $\hat{\Gamma}(i\gamma\cdot a)\hat{\Gamma}=-i\gamma\cdot a$; therefore
\[
\hat{\Gamma} T_a^*\hat{\Gamma}=(i\gamma\cdot a)\,S_m.
\]
Now let $a_1,a_2,a_3,a_4\in C_c^\infty(\mathbb{R}^3;\mathbb{R}^3)$ be real, and put $V_j=i\gamma\cdot a_j$ for $j=1,2,3,4$, so that $T_{a_j}=S_mV_j$ and $V_jS_m=\hat{\Gamma} T_{a_j}^*\hat{\Gamma}\in\mathfrak S_4$. For a trace class operator $X$ one has $\operatorname{Tr}X^*=\overline{\operatorname{Tr}X}$, and the trace is invariant under conjugation by the unitary matrix $\hat{\Gamma}=\hat{\Gamma}^{-1}$. Conjugating by $\hat{\Gamma}$ and inserting $\hat{\Gamma}^2=1$ between the factors, one obtains
\begin{align*}
\overline{\operatorname{Tr}\left[T_{a_1}T_{a_2}T_{a_3}T_{a_4}\right]}
&=\operatorname{Tr}\left[T_{a_4}^*T_{a_3}^*T_{a_2}^*T_{a_1}^*\right]
=\operatorname{Tr}\left[V_4S_mV_3S_mV_2S_mV_1S_m\right]\\
&=\operatorname{Tr}\left[S_mV_4S_mV_3S_mV_2S_mV_1\right]
=\operatorname{Tr}\left[T_{a_4}T_{a_3}T_{a_2}T_{a_1}\right].
\end{align*}
The third equality is the identity $\operatorname{Tr}AB=\operatorname{Tr}BA$, which holds for bounded operators $A$ and $B$ such that $AB$ and $BA$ are of trace class; here $A=V_4S_mV_3S_mV_2S_mV_1$ and $B=S_m$, and both $AB=(V_4S_m)(V_3S_m)(V_2S_m)(V_1S_m)$ and $BA=T_{a_4}T_{a_3}T_{a_2}T_{a_1}$ are products of four operators of $\mathfrak S_4$. Thus the trace of a word of length four in $T_a$ and $T_{a'}$ is the complex conjugate of the trace of the reversed word. Reversal leaves $T_a^4$ unchanged; it turns $T_aT_aT_{a'}T_{a'}$ into $T_{a'}T_{a'}T_aT_a$, the cyclic shift of $T_aT_aT_{a'}T_{a'}$ by two places; and it turns $T_aT_{a'}T_aT_{a'}$ into $T_{a'}T_aT_{a'}T_a$, the cyclic shift of $T_aT_{a'}T_aT_{a'}$ by one place. A cyclic shift does not alter the trace of a word, so the traces of $T_a^4$, $T_aT_aT_{a'}T_{a'}$ and $T_aT_{a'}T_aT_{a'}$ coincide with their own complex conjugates. These three traces are therefore real, and so is each of $\mathsf{b}(a)$ and $\mathsf{b}(a,a,a',a')$, being a combination of them with real coefficients.

(b) Let $\chi$ and $\chi'$ be the indicator functions of $\operatorname{supp}a$ and of $\operatorname{supp}a'$, regarded as multiplication operators on $L^2(\mathbb{R}^3;\mathbb{C}^4)$, and put
\[
N=\chi S_m\chi',\qquad N'=\chi'S_m\chi.
\]
Since $\chi$ and $\chi'$ are scalar functions, the multiplication by $\chi$, and likewise the multiplication by $\chi'$, commutes with the multiplication by each of the matrix-valued functions $i\gamma\cdot a$ and $i\gamma\cdot a'$. As $a$ vanishes outside $\operatorname{supp}a$ and $a'$ vanishes outside $\operatorname{supp}a'$, one obtains
\[
i\gamma\cdot a=\chi\,(i\gamma\cdot a)=(i\gamma\cdot a)\,\chi,\qquad i\gamma\cdot a'=\chi'\,(i\gamma\cdot a')=(i\gamma\cdot a')\,\chi'.
\]
Inserting these identities into $T_aT_{a'}=S_m(i\gamma\cdot a)\,S_m(i\gamma\cdot a')$ and into $T_{a'}T_a=S_m(i\gamma\cdot a')\,S_m(i\gamma\cdot a)$ gives
\[
T_aT_{a'}=S_m(i\gamma\cdot a)\,N\,(i\gamma\cdot a'),\qquad T_{a'}T_a=S_m(i\gamma\cdot a')\,N'\,(i\gamma\cdot a).
\]

By the anticommutation relations $(\gamma\cdot\partial)^2=\sum_{k,l}\gamma_k\gamma_l\partial_k\partial_l=\Delta$, so that $(m+\gamma\cdot\partial)(m-\gamma\cdot\partial)=m^2-\Delta$ and
\[
S_m=(m-\gamma\cdot\partial)(m^2-\Delta)^{-1}.
\]
The operator $(m^2-\Delta)^{-1}$ on $L^2(\mathbb{R}^3)$ is the convolution with the function
\[
G_m(u)=\frac{\exp(-m\cdot|u|)}{4\pi\cdot|u|}\qquad(u\in\mathbb{R}^3,\ u\ne0).
\]
Its partial derivatives $\partial_kG_m(u)=-\left(m+\frac{1}{|u|}\right)\cdot\frac{u_k}{|u|}\cdot G_m(u)$ are locally integrable. Hence $S_m$ is the integral operator with the kernel $S_m(x-y)$, where
\[
S_m(u)=\left[m+\frac{\gamma\cdot u}{|u|}\cdot\left(m+\frac{1}{|u|}\right)\right]\cdot G_m(u)\qquad(u\ne0).
\]
For $u\ne0$ the matrix $S_m(u)$ is Hermitian, and since $(\gamma\cdot u)^2=|u|^2$ and $\operatorname{tr}\gamma_k=0$,
\[
\operatorname{tr}\left[S_m(u)^2\right]=4\left(m^2+\left(m+\frac{1}{|u|}\right)^2\right)\cdot G_m(u)^2,
\]
which is a decreasing function of $|u|$.

The operator $N$ is the integral operator with the kernel $\chi(x)\,S_m(x-y)\,\chi'(y)$. For an integral operator on $L^2(\mathbb{R}^3;\mathbb{C}^4)$ whose kernel $K(x,y)$ takes values in the $4\times4$ matrices, the square of the Hilbert--Schmidt norm equals
\[
\int\mathrm{d}x\int\mathrm{d}y\,\operatorname{tr}\left[K(x,y)^*K(x,y)\right];
\]
for the kernel of $N$ the integrand is $\chi(x)\,\chi'(y)\operatorname{tr}\left[S_m(x-y)^2\right]$, as $S_m(x-y)$ is Hermitian. For $x\in\operatorname{supp}a$ and $y\in\operatorname{supp}a'$ one has $|x-y|\ge\mathfrak{d}$. Since $\operatorname{tr}\left[S_m(u)^2\right]$ decreases in $|u|$, and since $4\cdot\left(4\pi\right)^{-2}=\left(4\pi^2\right)^{-1}$, it follows that
\begin{multline*}
\|N\|_{\mathrm{HS}}^2=\int_{\operatorname{supp}a}\mathrm{d}x\int_{\operatorname{supp}a'}\mathrm{d}y\,\operatorname{tr}\left[S_m(x-y)^2\right]\\
\le|\operatorname{supp}a|\cdot|\operatorname{supp}a'|\cdot\left(m^2+\left(m+\frac{1}{\mathfrak{d}}\right)^2\right)\cdot\frac{\exp(-2m\cdot\mathfrak{d})}{4\pi^2\cdot\mathfrak{d}^2}.
\end{multline*}
The kernel of $N'$ is $\chi'(x)\,S_m(x-y)\,\chi(y)$, and the same bound holds for $\|N'\|_{\mathrm{HS}}^2$.

By the description of $S_m$ as a Fourier multiplier in the proof of (a), $\|S_m\|=\sup_{p\in\mathbb{R}^3}(m^2+|p|^2)^{-\frac{1}{2}}=\frac{1}{m}$. For $v\in\mathbb{R}^3$ the Hermitian matrix $\gamma\cdot v$ satisfies $(\gamma\cdot v)^2=|v|^2$, so that $\|\gamma\cdot v\|=|v|$; hence the multiplication by $i\gamma\cdot a$ has norm at most $\|a\|_\infty$, and likewise for $a'$. The inequality $\|XNY\|_{\mathrm{HS}}\le\|X\|\cdot\|N\|_{\mathrm{HS}}\cdot\|Y\|$, valid for bounded operators $X$, $Y$ and a Hilbert--Schmidt operator $N$, gives
\[
\|T_aT_{a'}\|_{\mathrm{HS}}\le\frac{\|a\|_\infty\cdot\|a'\|_\infty}{m}\cdot\|N\|_{\mathrm{HS}},\qquad\|T_{a'}T_a\|_{\mathrm{HS}}\le\frac{\|a\|_\infty\cdot\|a'\|_\infty}{m}\cdot\|N'\|_{\mathrm{HS}}.
\]
By cyclicity, $\operatorname{Tr}\left[T_aT_aT_{a'}T_{a'}\right]=\operatorname{Tr}\left[(T_aT_{a'})(T_{a'}T_a)\right]$, and $T_aT_{a'}T_aT_{a'}=(T_aT_{a'})^2$. The inequality $|\operatorname{Tr}AB|\le\|A\|_{\mathrm{HS}}\cdot\|B\|_{\mathrm{HS}}$ for Hilbert--Schmidt operators $A$ and $B$ therefore gives
\begin{align*}
\left|\operatorname{Tr}\left[T_aT_aT_{a'}T_{a'}\right]\right|&\le\frac{\|a\|_\infty^2\cdot\|a'\|_\infty^2}{m^2}\cdot\|N\|_{\mathrm{HS}}\cdot\|N'\|_{\mathrm{HS}},\\
\left|\operatorname{Tr}\left[(T_aT_{a'})^2\right]\right|&\le\frac{\|a\|_\infty^2\cdot\|a'\|_\infty^2}{m^2}\cdot\|N\|_{\mathrm{HS}}^2.
\end{align*}
Both $\|N\|_{\mathrm{HS}}\cdot\|N'\|_{\mathrm{HS}}$ and $\|N\|_{\mathrm{HS}}^2$ are bounded by the right side of the bound for $\|N\|_{\mathrm{HS}}^2$ obtained above. Inserting these estimates into the definition of $\mathsf{b}(a,a,a',a')$, whose two traces carry the coefficients $4$ and $2$, and adding the $4+2=6$ resulting terms, one obtains (b).
\end{proof}

\begin{theorem}[Non-Gaussian components from three tests]\label{thm:non-gaussian-components-2}
Let $\mu$ be a translation-invariant law on $(\Omega,\mathcal F)$ that satisfies \eqref{eq:fourth-cumulants-H1}, let $(\Xi,\nu,(\mu_\xi))$ be an ergodic decomposition of $\mu$ (Definition~\ref{def:ergodicity-ergodic-decompositions}), and let $\Xi_0\subset\Xi$ be the set of full $\nu$-measure of Corollary~\ref{cor:properties-almost-every}; as in Section~\ref{sec:cumulant-identities-non}, the components $\mu_\xi$ are considered for $\xi\in\Xi_0$ only. Let $\beta\in C_c^\infty(\mathbb{R}^3;\Lambda^2)$ be real, with $\operatorname{supp}\beta$ contained in a ball of radius $r$, let $a=d^*\beta$ and $f=dd^*\beta$ as in Section~\ref{sec:free-quartic-coefficient}, and assume $b(f)>0$. Let $x_0\in\mathbb{R}^3$ be such that $d_0=|x_0|-2r>0$ and
\begin{equation}\label{eq:fourth-cumulants-11}
\frac{6}{m^2}\cdot\|a\|_\infty^4\cdot\left(\frac{4\pi r^3}{3}\right)^2\cdot\left(m^2+\left(m+\frac{1}{d_0}\right)^2\right)\cdot\frac{\exp(-2m\cdot d_0)}{4\pi^2\cdot d_0^2}\le\frac{b(f)}{12}.
\end{equation}
Let $e>0$ and $\sigma>0$, and assume that \eqref{eq:fourth-cumulants-4g} holds for each of the three tests
\[
g=f,\qquad g=f+\tau_{x_0}f,\qquad g=f-\tau_{x_0}f,
\]
with the constants $C_0$, $C_+$, $C_-$, respectively, in place of $C_g$. These tests are $dd^*\beta$, $dd^*(\beta+\tau_{x_0}\beta)$ and $dd^*(\beta-\tau_{x_0}\beta)$, so that $b(g)$ is defined for them in Section~\ref{sec:free-quartic-coefficient}. The constant $C_0$ is a number and has no relation to the bilinear form $C_0(\cdot,\cdot)$ of Lemma~\ref{lem:free-transverse-covariance}. (Suppose that $e\le1$ and that, for each of the three tests, \eqref{eq:fourth-cumulants-4g} is satisfied with an exponent of its own in place of $\sigma$. Then \eqref{eq:fourth-cumulants-4g} is satisfied for all three tests with the least of the three exponents as the common value of $\sigma$, since $e^{4+\sigma'}\le e^{4+\sigma}$ for $\sigma'\ge\sigma$ when $e\le1$.) Put
\[
C_*=2C_0+\frac{1}{4}\left(C_++C_-\right),
\]
and assume $C_*\cdot e^\sigma\le\frac{b(f)}{4}$. Then the following hold.
\begin{enumerate}
\item[(i)] $\mathrm{Ext}^\mu_{x_0}(f)\ge\frac{1}{2}e^4\cdot b(f)$.
\item[(ii)] $\int_{\Xi_0}\mathrm{d}\nu(\xi)\,\mathrm{Ext}^\xi_{x_0}(f)\ge\frac{1}{2}e^4\cdot b(f)$.
\item[(iii)] With the constant $c_4=\frac{3}{2}\exp(2)$ of Lemma~\ref{lem:moment-bounds}(c),
\[
\nu\left(\left\{\xi\in\Xi_0:\mathrm{Ext}^\xi_{x_0}(f)\ge\frac{1}{4}e^4\cdot b(f)\right\}\right)\ge\frac{e^4\cdot b(f)}{4(c_4+6)\cdot\|f\|_2^4},\qquad c_4+6=\frac{3}{2}\exp(2)+6<17.09.
\]
\item[(iv)] For every $\xi\in\Xi_0$ in the set of (iii) the law $\mu_\xi$ is not Gaussian, and
\[
\max\left\{\left|\kappa_4^\xi(f)\right|,\left|\kappa_4^\xi(f+\tau_{x_0}f)\right|,\left|\kappa_4^\xi(f-\tau_{x_0}f)\right|\right\}\ge\frac{1}{10}e^4\cdot b(f).
\]
\end{enumerate}
\end{theorem}

\begin{proof}
Put $g_\pm=f\pm\tau_{x_0}f$. Translations commute with $d$ and with $d^*$; hence $d^*(\beta\pm\tau_{x_0}\beta)=a\pm\tau_{x_0}a$ and $g_\pm=dd^*(\beta\pm\tau_{x_0}\beta)=d(a\pm\tau_{x_0}a)$, and the definition of Section~\ref{sec:free-quartic-coefficient} gives
\[
b(g_\pm)=\mathsf{b}(a\pm\tau_{x_0}a).
\]

(i) By Lemma~\ref{lem:moment-bounds}(a),(c) the law $\mu$ is centred and has finite fourth moments, and it is translation invariant; so \eqref{eq:fourth-cumulants-1} gives
\[
\mathrm{Ext}^\mu_{x_0}(f)=2\,\kappa_4^\mu(f)-\frac{1}{4}\left[\kappa_4^\mu(g_+)+\kappa_4^\mu(g_-)\right].
\]
Condition \eqref{eq:fourth-cumulants-4g} at the test $f$ gives the lower bound $\kappa_4^\mu(f)\ge e^4\cdot b(f)-C_0\cdot e^{4+\sigma}$, and at the tests $g_\pm$ it gives the upper bounds $\kappa_4^\mu(g_\pm)\le e^4\cdot b(g_\pm)+C_\pm\cdot e^{4+\sigma}$. Therefore
\[
\mathrm{Ext}^\mu_{x_0}(f)\ge e^4\cdot\left\{2\,b(f)-\frac{1}{4}\left[b(g_+)+b(g_-)\right]\right\}-C_*\cdot e^{4+\sigma}.
\]
Lemma~\ref{lem:polarisation-decay-quartic}(a), applied to $a$ and $a'=\tau_{x_0}a$, together with $\mathsf{b}(\tau_{x_0}a)=\mathsf{b}(a)=b(f)$, gives
\[
b(g_+)+b(g_-)=\mathsf{b}(a+\tau_{x_0}a)+\mathsf{b}(a-\tau_{x_0}a)=4\,b(f)+12\,\mathsf{b}(a,a,\tau_{x_0}a,\tau_{x_0}a),
\]
where all the numbers are real by the same lemma. Consequently
\[
\mathrm{Ext}^\mu_{x_0}(f)\ge e^4\cdot\left\{b(f)-3\,\mathsf{b}(a,a,\tau_{x_0}a,\tau_{x_0}a)\right\}-C_*\cdot e^{4+\sigma}.
\]
Let $B\subset\mathbb{R}^3$ be a ball of radius $r$ that contains $\operatorname{supp}\beta$. The field $a$ is formed from first derivatives of $\beta$, so that $\operatorname{supp}a\subset B$ and $|\operatorname{supp}a|\le\frac{4\pi r^3}{3}$. The support of $\tau_{x_0}a$ is $x_0+\operatorname{supp}a\subset x_0+B$; moreover $|\operatorname{supp}\tau_{x_0}a|=|\operatorname{supp}a|$, $\|\tau_{x_0}a\|_\infty=\|a\|_\infty$, and
\[
\operatorname{dist}\left(\operatorname{supp}a,\operatorname{supp}\tau_{x_0}a\right)\ge\operatorname{dist}(B,x_0+B)=|x_0|-2r=d_0>0.
\]
The right side of the inequality of Lemma~\ref{lem:polarisation-decay-quartic}(b) decreases as $\mathfrak{d}$ increases. Hence Lemma~\ref{lem:polarisation-decay-quartic}(b), applied to $a$ and $a'=\tau_{x_0}a$, and \eqref{eq:fourth-cumulants-11} give
\[
\left|\mathsf{b}(a,a,\tau_{x_0}a,\tau_{x_0}a)\right|\le\frac{b(f)}{12}.
\]
Together with the hypothesis $C_*\cdot e^{\sigma}\le\frac{b(f)}{4}$ this yields
\[
\mathrm{Ext}^\mu_{x_0}(f)\ge\frac{3}{4}e^4\cdot b(f)-\frac{1}{4}e^4\cdot b(f)=\frac{1}{2}e^4\cdot b(f),
\]
which is (i).

(ii), (iii) Let $\delta_{\mathrm{Ext}}=\mathrm{Ext}^\mu_{x_0}(f)$ be the number of Theorem~\ref{thm:non-gaussian-components} for $x=x_0$; by (i), $\delta_{\mathrm{Ext}}\ge\frac{1}{2}e^4\cdot b(f)>0$. Applied with $x=x_0$, that theorem gives the integral inequality
\[
\int_{\Xi_0}\mathrm{d}\nu(\xi)\,\mathrm{Ext}^\xi_{x_0}(f)\ge\delta_{\mathrm{Ext}},
\]
which implies (ii), and the upper bound $\mathrm{Ext}^\xi_{x_0}(f)\le(c_4+6)\cdot\|f\|_2^4$ for $\xi\in\Xi_0$. Apply Lemma~\ref{lem:large-values-variable} on the probability space $\Xi_0$, endowed with the restriction of $\nu$, to the random variable $Z(\xi)=\mathrm{Ext}^\xi_{x_0}(f)$, with $M=(c_4+6)\cdot\|f\|_2^4$ and $\delta=\frac{1}{2}e^4\cdot b(f)$. Its hypotheses hold, since $Z\le M$ on $\Xi_0$ and $\int_{\Xi_0}\mathrm{d}\nu\,Z\ge\delta_{\mathrm{Ext}}\ge\delta>0$. The lemma gives
\[
\nu\left(\left\{\xi\in\Xi_0:\mathrm{Ext}^\xi_{x_0}(f)\ge\frac{\delta}{2}\right\}\right)\ge\frac{\delta}{2M}=\frac{e^4\cdot b(f)}{4(c_4+6)\cdot\|f\|_2^4},
\]
which is (iii), as $\frac{\delta}{2}=\frac{1}{4}e^4\cdot b(f)$. Finally $\exp(2)<7.3891$, so that $c_4<11.0837$ and $c_4+6<17.09$.

(iv) Let $\xi\in\Xi_0$ satisfy $\mathrm{Ext}^\xi_{x_0}(f)\ge\frac{1}{4}e^4\cdot b(f)$. By Corollary~\ref{cor:properties-almost-every} the law $\mu_\xi$ is translation invariant and satisfies \eqref{eq:fourth-cumulants-H1}; by Lemma~\ref{lem:moment-bounds}(a),(c) it is centred and has finite fourth moments. If $\mu_\xi$ were Gaussian, Section~\ref{sec:cumulants-gaussian-laws} would give $\mathrm{Ext}^\xi_{x_0}(f)=0$, contrary to $\mathrm{Ext}^\xi_{x_0}(f)\ge\frac{1}{4}e^4\cdot b(f)>0$; hence $\mu_\xi$ is not Gaussian. Moreover, \eqref{eq:fourth-cumulants-1} for the translation-invariant law $\mu_\xi$ gives, as in the proof of Theorem~\ref{thm:non-gaussian-components},
\begin{align*}
\frac{1}{4}e^4\cdot b(f)\le\mathrm{Ext}^\xi_{x_0}(f)&=2\,\kappa_4^\xi(f)-\frac{1}{4}\left[\kappa_4^\xi(g_+)+\kappa_4^\xi(g_-)\right]\\
&\le2\left|\kappa_4^\xi(f)\right|+\frac{1}{4}\left|\kappa_4^\xi(g_+)\right|+\frac{1}{4}\left|\kappa_4^\xi(g_-)\right|\\
&\le\frac{5}{2}\max\left\{\left|\kappa_4^\xi(f)\right|,\left|\kappa_4^\xi(g_+)\right|,\left|\kappa_4^\xi(g_-)\right|\right\}.
\end{align*}
Since $g_\pm=f\pm\tau_{x_0}f$ and $\frac{2}{5}\cdot\frac{1}{4}=\frac{1}{10}$, this is the inequality of (iv).
\end{proof}

\begin{remark}
The data of Theorem~\ref{thm:non-gaussian-components-2} can be fixed successively. First $\beta$ is chosen, together with the radius $r$ of a ball that contains $\operatorname{supp}\beta$; the form $\beta$ determines $a$, $f$ and $b(f)$. Next $x_0$ is chosen: apart from $|x_0|$, the conditions $d_0>0$ and \eqref{eq:fourth-cumulants-11} involve only $m$, $r$, $\|a\|_\infty$ and $b(f)$, quantities attached to $m$, $\beta$ and $r$ that do not involve the law $\mu$. After $x_0$ come the constants $C_0$, $C_+$, $C_-$ of \eqref{eq:fourth-cumulants-4g} at the three tests $f$, $f+\tau_{x_0}f$, $f-\tau_{x_0}f$, and last the upper bound on $e$ expressed by $C_*\cdot e^\sigma\le\frac{b(f)}{4}$. Of each of the two-sided bounds \eqref{eq:fourth-cumulants-4g} at the three tests, the proof of Theorem~\ref{thm:non-gaussian-components-2} uses one side only: at the test $f$ the bound serves to estimate $\kappa_4^\mu(f)$ from below, and at the tests $f\pm\tau_{x_0}f$ it serves to estimate the sum $\kappa_4^\mu(f+\tau_{x_0}f)+\kappa_4^\mu(f-\tau_{x_0}f)$ from above.
\end{remark}

\subsection{Clustering and ergodicity}\label{sec:clustering-ergodicity}

\begin{lemma}[Density of polynomials]\label{lem:density-polynomials}
Let $\mu'$ be a probability law on $(\Omega,\mathcal F)$ with expectation $E'$, let $V$ be a linear subspace of $\mathcal S$, and let $\mathcal F_V=\sigma\left(X(f):f\in V\right)$, as in Lemma~\ref{lem:exponentials-determination-law}. A \emph{polynomial in the variables $X(f)$, $f\in V$,} is a complex linear combination of finitely many products $X(f_1)\cdots X(f_k)$ with $k\ge0$ and $f_1,\dots,f_k\in V$, the empty product being the constant function $1$. Assume that for every $f\in V$ there is a number $t(f)>0$ such that
\[
E'\exp\left(t(f)\cdot|X(f)|\right)<\infty .
\]
Then every polynomial in the variables $X(f)$, $f\in V$, belongs to $L^2(\Omega,\mathcal F_V,\mu')$, and these polynomials form a dense linear subspace of $L^2(\Omega,\mathcal F_V,\mu')$.
\end{lemma}

\begin{proof}
Let $f\in V$ and $t=t(f)$. For $u\in\mathbb{R}$ and an integer $k\ge0$ the number $\frac{1}{k!}\left(t|u|\right)^k$ is one of the non-negative terms of the series of $\exp(t|u|)$, so that $|u|^k\le k!\cdot t^{-k}\cdot\exp(t|u|)$. Hence
\[
E'\left[|X(f)|^k\right]\le k!\cdot t^{-k}\cdot E'\exp\left(t\cdot|X(f)|\right)<\infty ,
\]
and every $X(f)$, $f\in V$, has finite moments of all orders. Let $k\ge1$ and $f_1,\dots,f_k\in V$. H\"older's inequality with $k$ exponents equal to $k$ gives
\[
E'\left[\left|X(f_1)\cdots X(f_k)\right|^2\right]\le\prod_{i=1}^k\left(E'\left[|X(f_i)|^{2k}\right]\right)^{\frac{1}{k}}<\infty .
\]
The product $X(f_1)\cdots X(f_k)$ is $\mathcal F_V$-measurable, and therefore an element of $L^2(\Omega,\mathcal F_V,\mu')$; consequently every polynomial in the variables $X(f)$, $f\in V$, belongs to this space. The polynomials form a linear subspace of it, and it remains to show that the orthogonal complement of this subspace in $L^2(\Omega,\mathcal F_V,\mu')$ is $\{0\}$.

Let $H\in L^2(\Omega,\mathcal F_V,\mu')$ be orthogonal to every polynomial $Q$ in the variables $X(f)$, $f\in V$; that is, $E'\left[\overline{H}\cdot Q\right]=0$ for all such $Q$. The products $X(f_1)\cdots X(f_k)$ are real, so that the complex conjugate $\overline{Q}$ of such a polynomial is again one. Applying the orthogonality to $\overline{Q}$ and taking complex conjugates, one finds
\[
E'\left[H\cdot Q\right]=0\qquad\text{for all such polynomials $Q$.}
\]

Fix $f\in V$, write $X=X(f)$ and $t=t(f)$, and let $W=\left\{z\in\mathbb{C}:|\operatorname{Re}z|<\frac{t}{2}\right\}$. For $0\le s<\frac{t}{2}$ one has $2s<t$, and the Cauchy--Schwarz inequality gives
\[
E'\left[|H|\cdot\exp\left(s\cdot|X|\right)\right]\le\|H\|_2\cdot\left(E'\exp\left(2s\cdot|X|\right)\right)^{\frac{1}{2}}<\infty ,
\]
where $\|H\|_2$ is the norm of $H$ in $L^2(\mu')$. For $z\in W$ one has $\left|\exp(zX)\right|=\exp\left(\operatorname{Re}z\cdot X\right)\le\exp\left(|\operatorname{Re}z|\cdot|X|\right)$, and the estimate with $s=|\operatorname{Re}z|$ shows that $H\cdot\exp(zX)$ is integrable; so the function
\[
\varphi\colon W\to\mathbb{C},\qquad\varphi(z)=E'\left[H\cdot\exp(zX)\right],
\]
is well defined. Let $z_0\in W$, and let $\delta>0$ be such that $|\operatorname{Re}z_0|+\delta<\frac{t}{2}$. For $z\in\mathbb{C}$ with $|z-z_0|<\delta$ one has
\[
\exp(zX)=\sum_{n=0}^{\infty}\frac{(z-z_0)^n}{n!}\cdot X^n\cdot\exp(z_0X)
\]
and
\begin{align*}
\sum_{n=0}^{\infty}\frac{|z-z_0|^n}{n!}\cdot|X|^n\cdot\left|\exp(z_0X)\right|&=\exp\left(|z-z_0|\cdot|X|+\operatorname{Re}z_0\cdot X\right)\\
&\le\exp\left(\left(|\operatorname{Re}z_0|+\delta\right)\cdot|X|\right).
\end{align*}
By the preceding estimate with $s=|\operatorname{Re}z_0|+\delta$, the right side, multiplied by $|H|$, is integrable. Dominated convergence, applied to the partial sums of the series, therefore yields
\[
\varphi(z)=\sum_{n=0}^{\infty}\frac{(z-z_0)^n}{n!}\cdot E'\left[H\cdot X^n\cdot\exp(z_0X)\right]\qquad\left(|z-z_0|<\delta\right).
\]
Thus $\varphi$ is analytic on $W$. At $z_0=0$ the coefficients of the series are the numbers $\frac{1}{n!}E'\left[H\cdot X(f)^n\right]$, which vanish because $X(f)^n$ is a polynomial in the variables $X(g)$, $g\in V$; so $\varphi$ vanishes on a disc centred at $0$. The strip $W$ is connected, and by the identity theorem for analytic functions $\varphi=0$ on $W$. Since $i\in W$,
\[
E'\left[H\cdot\exp\left(iX(f)\right)\right]=\varphi(i)=0 .
\]

This identity holds for every $f\in V$. By complex conjugation, $E'\left[\exp\left(-iX(f)\right)\cdot\overline{H}\right]=0$ for every $f\in V$; in other words, the element $\overline{H}$ of $L^2(\Omega,\mathcal F_V,\mu')$ is orthogonal to all functions $\exp\left(iX(f)\right)$, $f\in V$. By Lemma~\ref{lem:exponentials-determination-law}(a), the complex linear span of these functions is dense in $L^2(\Omega,\mathcal F_V,\mu')$. Hence $\overline{H}=0$, and $H=0$.
\end{proof}

\begin{theorem}[Clustering, mixing and ergodicity]\label{thm:clustering-mixing-ergodicity}
Let $\mu$ be a translation-invariant law on $(\Omega,\mathcal F)$ with expectation $E$, and assume that for every $f\in\mathcal D$ there is a number $t=t(f)>0$ such that $E\exp\left(t\cdot|X(f)|\right)<\infty$. (If $\mu$ satisfies \eqref{eq:fourth-cumulants-H1}, then $E\exp\left(t\cdot|X(f)|\right)<\infty$ for all $f\in\mathcal S$ and all $t>0$, by Lemma~\ref{lem:moment-bounds}(d).) Let $\langle\cdot,\cdot\rangle$ and $\|\cdot\|$ be the scalar product and the norm of $L^2(\mu)$, as in Section~\ref{sec:test-forms-distributions-2}. For all integers $m,n\ge1$ and all $f_1,\dots,f_m,g_1,\dots,g_n\in\mathcal D$ the function
\[
\mathbb{R}^3\to\mathbb{R},\qquad a\mapsto E\left[\prod_{i=1}^mX(f_i)\cdot\prod_{l=1}^nX(\tau_ag_l)\right],
\]
is bounded and measurable. Consider the following statements.
\begin{multline}\tag{C1}\label{eq:fourth-cumulants-C1}
\lim_{|a|\to\infty}E\left[\prod_{i=1}^mX(f_i)\cdot\prod_{l=1}^nX(\tau_ag_l)\right]\\
=E\left[\prod_{i=1}^mX(f_i)\right]\cdot E\left[\prod_{l=1}^nX(g_l)\right]\\
\text{for all integers $m,n\ge1$ and all $f_1,\dots,f_m,g_1,\dots,g_n\in\mathcal D$.}
\end{multline}
\begin{multline}\tag{C2}\label{eq:fourth-cumulants-C2}
\text{There is a sequence $(a_k)_{k\ge1}$ in $\mathbb{R}^3$ such that}\\
\lim_{k\to\infty}E\left[\prod_{i=1}^mX(f_i)\cdot\prod_{l=1}^nX(\tau_{a_k}g_l)\right]\\
=E\left[\prod_{i=1}^mX(f_i)\right]\cdot E\left[\prod_{l=1}^nX(g_l)\right]\\
\text{for all integers $m,n\ge1$ and all $f_1,\dots,f_m,g_1,\dots,g_n\in\mathcal D$.}
\end{multline}
\begin{multline}\tag{C3}\label{eq:fourth-cumulants-C3}
\lim_{L\to\infty}L^{-3}\int_{[0,L]^3}\mathrm{d}a\,E\left[\prod_{i=1}^mX(f_i)\cdot\prod_{l=1}^nX(\tau_ag_l)\right]\\
=E\left[\prod_{i=1}^mX(f_i)\right]\cdot E\left[\prod_{l=1}^nX(g_l)\right]\\
\text{for all integers $m,n\ge1$ and all $f_1,\dots,f_m,g_1,\dots,g_n\in\mathcal D$.}
\end{multline}
\begin{equation}\tag{M}\label{eq:fourth-cumulants-M}
\lim_{|a|\to\infty}\langle G,\mathbb{U}_aH\rangle=\langle G,1\rangle\cdot\langle1,H\rangle\qquad\text{for all }G,H\in L^2(\mu).
\end{equation}
\begin{equation}\tag{Erg}\label{eq:fourth-cumulants-Erg}
\text{The law $\mu$ is ergodic in the sense of Definition~\ref{def:ergodicity-ergodic-decompositions}.}
\end{equation}
\begin{multline}\tag{T}\label{eq:fourth-cumulants-T}
\text{For every ergodic decomposition $(\Xi,B,\nu)$, $(\mu_\xi)_{\xi\in\Xi}$ of $\mu$,}\\
\text{one has $\mu_\xi=\mu$ for $\nu$-almost every $\xi\in\Xi$.}
\end{multline}
Then
\[
\eqref{eq:fourth-cumulants-C1}\Leftrightarrow\eqref{eq:fourth-cumulants-M}\Rightarrow\eqref{eq:fourth-cumulants-C2}\Rightarrow\eqref{eq:fourth-cumulants-Erg}\Leftrightarrow\eqref{eq:fourth-cumulants-C3}\Leftrightarrow\eqref{eq:fourth-cumulants-T}.
\]
\end{theorem}

\begin{proof}
The objects $I$, $H_{\mathrm{inv}}$, $P=E[\,\cdot\mid I]$ and $A_L$ are those of Lemma~\ref{lem:averages-translations-2}, formed with $\mu'=\mu$. Since $\mu$ is translation invariant, for every $a\in\mathbb{R}^3$ the operator $\mathbb{U}_a\colon G\mapsto G\circ T_a^{-1}=G\circ T_{-a}$ maps $L^2(\mu)$ onto itself and preserves the norm; moreover $\mathbb{U}_a1=1$ and $\mathbb{U}_a(G\cdot H)=(\mathbb{U}_aG)\cdot(\mathbb{U}_aH)$. As $\mathbb{U}_aX(g)=X(\tau_ag)$ for $g\in\mathcal S$ (Section~\ref{sec:test-forms-distributions-2}), it follows that
\[
\mathbb{U}_a\prod_{l=1}^nX(g_l)=\prod_{l=1}^nX(\tau_ag_l)\qquad(n\ge1,\ g_1,\dots,g_n\in\mathcal S).
\]

\emph{Step 1: polynomials.} Let $\mathrm{Pol}$ be the set of the polynomials in the variables $X(f)$, $f\in\mathcal D$, in the sense of Lemma~\ref{lem:density-polynomials}, and put $\mathcal F_{\mathcal D}=\sigma\left(X(f):f\in\mathcal D\right)$. The space $\mathcal D$ is a linear subspace of $\mathcal S$, and by hypothesis $E\exp\left(t(f)\cdot|X(f)|\right)<\infty$ for every $f\in\mathcal D$. By Lemma~\ref{lem:density-polynomials}, applied with $\mu'=\mu$ and $V=\mathcal D$, the set $\mathrm{Pol}$ is a dense linear subspace of $L^2(\Omega,\mathcal F_{\mathcal D},\mu)$. Further, $\mathcal D$ is sequentially dense in $\mathcal S$: if $\chi\in C_c^\infty(\mathbb{R}^3)$ equals $1$ in a neighbourhood of $0$, then for every $f\in\mathcal S$ the forms $x\mapsto\chi\left(\frac{x}{k}\right)\cdot f(x)$, $k\ge1$, belong to $\mathcal D$ and converge to $f$ in $\mathcal S$ as $k\to\infty$. Lemma~\ref{lem:exponentials-determination-law}(b), applied with $\mu'=\mu$ and $V_0=\mathcal D$, therefore gives $\mathcal F_{\mathcal D}=\mathcal F$. Consequently $\mathrm{Pol}$ is a dense linear subspace of $L^2(\mu)$. In particular, all products $X(h_1)\cdots X(h_k)$ with $k\ge1$ and $h_1,\dots,h_k\in\mathcal D$ lie in $L^2(\mu)$, and the expectations in \eqref{eq:fourth-cumulants-C1}--\eqref{eq:fourth-cumulants-C3} are finite.

\emph{Step 2: averages of correlations.} Let $G,H\in L^2(\mu)$. For every $a\in\mathbb{R}^3$ the Cauchy--Schwarz inequality and the invariance of the norm under $\mathbb{U}_a$ give
\[
E\left[|G|\cdot\left|H\circ T_{-a}\right|\right]\le\|G\|\cdot\|H\| ,
\]
and in particular $\left|\langle G,\mathbb{U}_aH\rangle\right|\le\|G\|\cdot\|H\|$. Fix measurable representatives of $G$ and $H$. The map $(x,\omega)\mapsto T_x\omega$ is jointly measurable by Lemma~\ref{lem:averages-translations-2}; hence the function $(a,\omega)\mapsto\overline{G(\omega)}\cdot H(T_{-a}\omega)$ on $\mathbb{R}^3\times\Omega$ is jointly measurable, and for each $a$ its integral with respect to $\mu$ is $\langle G,\mathbb{U}_aH\rangle$. The theorem of Tonelli, applied to the positive and the negative parts of its real and its imaginary part, whose integrals with respect to $\mu$ are finite by the bound just proved, shows that $a\mapsto\langle G,\mathbb{U}_aH\rangle$ is a measurable function on $\mathbb{R}^3$. For $L>0$ the same bound gives $\int_{[0,L]^3}\mathrm{d}a\,E\left[|G|\cdot\left|H\circ T_{-a}\right|\right]\le L^3\cdot\|G\|\cdot\|H\|$, so that Fubini's theorem and the definition of $A_L$ yield
\begin{align*}
L^{-3}\int_{[0,L]^3}\mathrm{d}a\,\langle G,\mathbb{U}_aH\rangle&=E\left[\overline{G}\cdot L^{-3}\int_{[0,L]^3}\mathrm{d}a\,H\circ T_{-a}\right]\\
&=\langle G,A_LH\rangle .
\end{align*}
Now let $m,n\ge1$ and $f_1,\dots,f_m,g_1,\dots,g_n\in\mathcal D$, and take $G=\prod_{i=1}^mX(f_i)$ and $H=\prod_{l=1}^nX(g_l)$. These functions are real and lie in $L^2(\mu)$ by Step~1, and by the identity for $\mathbb{U}_a$ displayed above, $\langle G,\mathbb{U}_aH\rangle=E\left[\prod_{i=1}^mX(f_i)\cdot\prod_{l=1}^nX(\tau_ag_l)\right]$. This proves that the function of $a$ in the statement is bounded and measurable.

\emph{Step 3: reduction to polynomials.} For $Q,Q'\in\mathrm{Pol}$ and $a\in\mathbb{R}^3$ put
\[
\Psi_a(Q,Q')=\langle Q,\mathbb{U}_aQ'\rangle-\langle Q,1\rangle\cdot\langle1,Q'\rangle .
\]
By Step~2 this is a bounded measurable function of $a$; it is conjugate linear in $Q$ and linear in $Q'$. Let $Q=\prod_{i=1}^mX(f_i)$ and $Q'=\prod_{l=1}^nX(g_l)$ with $m,n\ge1$ and $f_1,\dots,f_m,g_1,\dots,g_n\in\mathcal D$. These functions are real, and
\begin{align*}
\langle Q,\mathbb{U}_aQ'\rangle&=E\left[\prod_{i=1}^mX(f_i)\cdot\prod_{l=1}^nX(\tau_ag_l)\right],\\
\langle Q,1\rangle\cdot\langle1,Q'\rangle&=E\left[\prod_{i=1}^mX(f_i)\right]\cdot E\left[\prod_{l=1}^nX(g_l)\right].
\end{align*}
If instead $Q=1$, then $\Psi_a(1,Q')=E\left[\mathbb{U}_aQ'\right]-E[Q']=0$ for every $a$, by the translation invariance of $\mu$; and if $Q'=1$, then $\Psi_a(Q,1)=\langle Q,1\rangle-\langle Q,1\rangle=0$, because $\mathbb{U}_a1=1$. Every element of $\mathrm{Pol}$ is a finite linear combination of the constant $1$ and of products of the form just considered. By sesquilinearity it follows that \eqref{eq:fourth-cumulants-C1} holds if and only if
\[
\lim_{|a|\to\infty}\Psi_a(Q,Q')=0\qquad\text{for all }Q,Q'\in\mathrm{Pol} ;
\]
that \eqref{eq:fourth-cumulants-C2} holds if and only if there is a sequence $(a_k)_{k\ge1}$ in $\mathbb{R}^3$ such that $\lim_{k\to\infty}\Psi_{a_k}(Q,Q')=0$ for all $Q,Q'\in\mathrm{Pol}$; and that \eqref{eq:fourth-cumulants-C3} holds if and only if
\[
\lim_{L\to\infty}L^{-3}\int_{[0,L]^3}\mathrm{d}a\,\Psi_a(Q,Q')=0\qquad\text{for all }Q,Q'\in\mathrm{Pol} .
\]

\emph{\eqref{eq:fourth-cumulants-C1} implies \eqref{eq:fourth-cumulants-M}.} Let $G,H\in L^2(\mu)$ and $\varepsilon>0$. By Step~1 there are $Q,Q'\in\mathrm{Pol}$ with $\|G-Q\|<\varepsilon$ and $\|H-Q'\|<\varepsilon$; then $\|Q\|<\|G\|+\varepsilon$. For $a\in\mathbb{R}^3$ one has
\begin{align*}
\langle G,\mathbb{U}_aH\rangle-\langle Q,\mathbb{U}_aQ'\rangle&=\langle G-Q,\mathbb{U}_aH\rangle+\langle Q,\mathbb{U}_a(H-Q')\rangle,\\
\langle G,1\rangle\cdot\langle1,H\rangle-\langle Q,1\rangle\cdot\langle1,Q'\rangle&=\langle G-Q,1\rangle\cdot\langle1,H\rangle\\
&\qquad+\langle Q,1\rangle\cdot\langle1,H-Q'\rangle .
\end{align*}
By the Cauchy--Schwarz inequality, the invariance of the norm under $\mathbb{U}_a$ and $\|1\|=1$, the modulus of each of the two right sides is at most $\|G-Q\|\cdot\|H\|+\|Q\|\cdot\|H-Q'\|\le\varepsilon\cdot\left(\|H\|+\|Q\|\right)$, whatever $a$ is. Consequently, for every $a\in\mathbb{R}^3$,
\begin{align*}
\left|\langle G,\mathbb{U}_aH\rangle-\langle G,1\rangle\cdot\langle1,H\rangle\right|&\le\left|\Psi_a(Q,Q')\right|+2\varepsilon\cdot\left(\|H\|+\|Q\|\right)\\
&\le\left|\Psi_a(Q,Q')\right|+2\varepsilon\cdot\left(\|H\|+\|G\|+\varepsilon\right).
\end{align*}
By \eqref{eq:fourth-cumulants-C1} and Step~3, $\Psi_a(Q,Q')\to0$ as $|a|\to\infty$. Hence the upper limit of the left side as $|a|\to\infty$ is at most $2\varepsilon\cdot\left(\|H\|+\|G\|+\varepsilon\right)$; as $\varepsilon>0$ is arbitrary, \eqref{eq:fourth-cumulants-M} follows.

\emph{\eqref{eq:fourth-cumulants-M} implies \eqref{eq:fourth-cumulants-C1}.} Let $m,n\ge1$ and $f_1,\dots,f_m,g_1,\dots,g_n\in\mathcal D$. The functions $Q=\prod_{i=1}^mX(f_i)$ and $Q'=\prod_{l=1}^nX(g_l)$ lie in $L^2(\mu)$ by Step~1, and by the two identities of Step~3, statement \eqref{eq:fourth-cumulants-M} with $G=Q$ and $H=Q'$ is the convergence asserted in \eqref{eq:fourth-cumulants-C1} for these forms.

\emph{\eqref{eq:fourth-cumulants-C1} implies \eqref{eq:fourth-cumulants-C2}.} The points $a_k=(k,0,0)$, $k\ge1$, satisfy $|a_k|\to\infty$; so \eqref{eq:fourth-cumulants-C1} gives the convergence required in \eqref{eq:fourth-cumulants-C2} along this sequence, for all $m,n$ and all forms.

\emph{\eqref{eq:fourth-cumulants-C2} implies \eqref{eq:fourth-cumulants-Erg}.} Let $(a_k)_{k\ge1}$ be a sequence as in \eqref{eq:fourth-cumulants-C2}. Let $G\in H_{\mathrm{inv}}$ and $\varepsilon>0$, and choose $Q\in\mathrm{Pol}$ with $\|G-Q\|<\varepsilon$ (Step~1); then $\|Q\|<\|G\|+\varepsilon$. Since $\mathbb{U}_{a_k}G=G$, one has $\|G\|^2=\langle G,\mathbb{U}_{a_k}G\rangle$ for every $k$, and the bound obtained in the proof that \eqref{eq:fourth-cumulants-C1} implies \eqref{eq:fourth-cumulants-M}, with $H=G$ and $Q'=Q$, gives
\[
\left|\|G\|^2-\langle Q,\mathbb{U}_{a_k}Q\rangle\right|\le\varepsilon\cdot\left(\|G\|+\|Q\|\right)\qquad(k\ge1).
\]
By \eqref{eq:fourth-cumulants-C2} and Step~3, $\Psi_{a_k}(Q,Q)\to0$, that is, $\langle Q,\mathbb{U}_{a_k}Q\rangle\to\langle Q,1\rangle\cdot\langle1,Q\rangle=\left|E[Q]\right|^2$ as $k\to\infty$. Moreover $\left|E[Q]-E[G]\right|=\left|\langle1,Q-G\rangle\right|\le\|Q-G\|<\varepsilon$. Letting $k\to\infty$ one obtains
\[
\|G\|^2\le\left(\left|E[G]\right|+\varepsilon\right)^2+\varepsilon\cdot\left(2\|G\|+\varepsilon\right),
\]
and letting $\varepsilon\to0$, $\|G\|^2\le\left|E[G]\right|^2$. Since $\left\|G-E[G]\right\|^2=\|G\|^2-\left|E[G]\right|^2$, the function $G$ coincides almost surely with the constant $E[G]$. Thus $H_{\mathrm{inv}}$ consists of constant functions, and $\mu$ is ergodic.

\emph{\eqref{eq:fourth-cumulants-Erg} is equivalent to \eqref{eq:fourth-cumulants-C3}.} Let $Q,Q'\in\mathrm{Pol}$. By Lemma~\ref{lem:averages-translations-2}(c), $A_LQ'\to PQ'$ in $L^2(\mu)$ as $L\to\infty$; together with Step~2 this gives
\[
L^{-3}\int_{[0,L]^3}\mathrm{d}a\,\langle Q,\mathbb{U}_aQ'\rangle=\langle Q,A_LQ'\rangle\longrightarrow\langle Q,PQ'\rangle
\]
as $L\to\infty$. By Lemma~\ref{lem:averages-translations-2}(b), $P$ is the orthogonal projection of $L^2(\mu)$ onto $H_{\mathrm{inv}}$. If $\mu$ is ergodic, $H_{\mathrm{inv}}$ is the space of the constant functions, so that $PQ'=\langle1,Q'\rangle\cdot1$ and $\langle Q,PQ'\rangle=\langle Q,1\rangle\cdot\langle1,Q'\rangle$; by the display, $L^{-3}\int_{[0,L]^3}\mathrm{d}a\,\Psi_a(Q,Q')\to0$ as $L\to\infty$, and \eqref{eq:fourth-cumulants-C3} follows by Step~3. Conversely, assume \eqref{eq:fourth-cumulants-C3}. By Step~3 and the display, $\langle Q,PQ'\rangle=\langle Q,1\rangle\cdot\langle1,Q'\rangle$ for all $Q,Q'\in\mathrm{Pol}$. Both sides are continuous functions of $(Q,Q')\in L^2(\mu)\times L^2(\mu)$, and $\mathrm{Pol}$ is dense in $L^2(\mu)$ by Step~1; hence
\[
\langle G,PH\rangle=\langle G,1\rangle\cdot\langle1,H\rangle=\left\langle G,\langle1,H\rangle\cdot1\right\rangle\qquad\left(G,H\in L^2(\mu)\right),
\]
that is, $PH=\langle1,H\rangle\cdot1$ for every $H\in L^2(\mu)$. Every $H\in H_{\mathrm{inv}}$ satisfies $H=PH$ and is therefore constant: $\mu$ is ergodic.

\emph{\eqref{eq:fourth-cumulants-Erg} implies \eqref{eq:fourth-cumulants-T}.} Let $\mu$ be ergodic, and let $(\Xi,B,\nu)$, $(\mu_\xi)_{\xi\in\Xi}$ be an ergodic decomposition of $\mu$. Let $V_0$ and $G_0$ be the countable sets of Steps~2 and~3 of the proof of Proposition~\ref{prop:existence-conditional-expectations}(a): $V_0$ is the $\mathbb{Q}$-linear span of the forms $\tau_qh_n$, $q\in\mathbb{Q}^3$, $n\in\mathbb{N}$, where $(h_n)$ is the Hermite basis of $L^2(\mathbb{R}^3;\Lambda^2)$, and it is sequentially dense in $\mathcal S$; $G_0$ is the set of the functions $\sum_{l=1}^pc_l\cdot\exp\left(iX(f_l)\right)$ with $p\ge1$, $c_1,\dots,c_p\in\mathbb{Q}+i\mathbb{Q}$ and $f_1,\dots,f_p\in V_0$. Let $G\in G_0$. Then $G$ is bounded, hence an element of $L^2(\mu)$, and since $H_{\mathrm{inv}}$ consists of constant functions, $E[G\mid I]=PG=E[G]$ almost surely. By Proposition~\ref{prop:existence-conditional-expectations}(b), the law of the function $\xi\mapsto E_\xi[G]$ under $\nu$ is the law of $E[G\mid I]$ under $\mu$, that is, the point mass at $E[G]$; so $E_\xi[G]=E[G]$ for $\nu$-almost every $\xi$. As $G_0$ is countable, there is a $\nu$-null set $N\subset\Xi$ such that $E_\xi[G]=E[G]$ for all $\xi\in\Xi\setminus N$ and all $G\in G_0$. The set $G_0$ contains the functions $\exp\left(iX(f)\right)$, $f\in V_0$. Hence, for $\xi\in\Xi\setminus N$, the characteristic functionals of $\mu_\xi$ and of $\mu$ coincide on the sequentially dense set $V_0$, and Lemma~\ref{lem:exponentials-determination-law}(b), applied with $\mu'=\mu$ and $\mu''=\mu_\xi$, gives $\mu_\xi=\mu$.

\emph{\eqref{eq:fourth-cumulants-T} implies \eqref{eq:fourth-cumulants-Erg}.} By Proposition~\ref{prop:existence-conditional-expectations}(a), the law $\mu$ has an ergodic decomposition $(\Xi,B,\nu)$, $(\mu_\xi)_{\xi\in\Xi}$. By \eqref{eq:fourth-cumulants-T} and \eqref{eq:fourth-cumulants-ED2} there are $\nu$-null sets $N_1,N_2\subset\Xi$ such that $\mu_\xi=\mu$ for $\xi\in\Xi\setminus N_1$, and $\mu_\xi$ is translation invariant and ergodic for $\xi\in\Xi\setminus N_2$. Since $\nu(\Xi)=1$, the set $\Xi\setminus(N_1\cup N_2)$ is not empty, and for each of its elements $\xi$ the law $\mu=\mu_\xi$ is ergodic.

The implications \eqref{eq:fourth-cumulants-C1}$\Rightarrow$\eqref{eq:fourth-cumulants-M}, \eqref{eq:fourth-cumulants-M}$\Rightarrow$\eqref{eq:fourth-cumulants-C1}, \eqref{eq:fourth-cumulants-C1}$\Rightarrow$\eqref{eq:fourth-cumulants-C2}, \eqref{eq:fourth-cumulants-C2}$\Rightarrow$\eqref{eq:fourth-cumulants-Erg}, \eqref{eq:fourth-cumulants-Erg}$\Leftrightarrow$\eqref{eq:fourth-cumulants-C3} and \eqref{eq:fourth-cumulants-Erg}$\Leftrightarrow$\eqref{eq:fourth-cumulants-T} are thus established, and together they give the assertion.
\end{proof}

\begin{corollary}[The decomposition of a clustering law]\label{cor:decomposition-clustering-law}
Let $\mu$ be a law on $(\Omega,\mathcal F)$, with expectation $E$, which satisfies \eqref{eq:fourth-cumulants-H1}, \eqref{eq:fourth-cumulants-H2}, \eqref{eq:fourth-cumulants-H3} and \eqref{eq:fourth-cumulants-C1}, and let $(\Xi,B,\nu)$, $(\mu_\xi)_{\xi\in\Xi}$ be an ergodic decomposition of $\mu$. Then the following hold.
\begin{enumerate}
\item[(i)] The law $\mu$ satisfies \eqref{eq:fourth-cumulants-M}, that is, $\mu$ is \emph{mixing}; and $\mu$ is ergodic.
\item[(ii)] One has $\mu_\xi=\mu$ for $\nu$-almost every $\xi\in\Xi$. Let the set of all laws on $(\Omega,\mathcal F)$ carry the $\sigma$-algebra generated by the maps $\rho\mapsto\rho(A)$, $A\in\mathcal F$. Then the map $\xi\mapsto\mu_\xi$ is measurable, and the image of $\nu$ under this map is the point mass at $\mu$. In particular, if $\Xi$ is the set of all laws on $(\Omega,\mathcal F)$, $B$ is this $\sigma$-algebra and $\mu_\xi=\xi$ for every $\xi\in\Xi$, then $\nu$ is the point mass at $\mu$.
\item[(iii)] For $\nu$-almost every $\xi\in\Xi$ the law $\mu_\xi$ has every property of $\mu$. In particular, $\mu_\xi$ satisfies \eqref{eq:fourth-cumulants-H1}, \eqref{eq:fourth-cumulants-H2}, including the invariance under the motions $(0,r)$ with $r\in O_{\mathrm{cub}}$, and \eqref{eq:fourth-cumulants-H3}; and for every integer $n\ge1$ and all $f_1,\dots,f_n\in\mathcal S$
\[
\kappa_n^\xi(f_1,\dots,f_n)=\kappa_n^\mu(f_1,\dots,f_n),
\]
where $\kappa_n^\lambda(f_1,\dots,f_n)$ denotes the joint cumulant of order $n$ of the variables $X(f_1),\dots,X(f_n)$ under a law $\lambda$ under which they have finite moments of all orders, and $\kappa_n^\xi=\kappa_n^{\mu_\xi}$.
\item[(iv)] The numbers $\Gamma(f_1,f_2;f_3,f_4)$ and $V(f)$ of Theorem~\ref{thm:fourth-cumulant-mixture}, formed with this ergodic decomposition, vanish for all $f_1,f_2,f_3,f_4,f\in\mathcal S$. For $f\in\mathcal S$ and integers $k\ge1$, let $V_k(f)$ be the variance under $\nu$ of the function $\xi\mapsto E_\xi\left[X(f)^k\right]$ on the set $\Xi_0$ of Corollary~\ref{cor:properties-almost-every}, so that $V_2(f)=V(f)$, and put
\[
\delta_V(f)=\kappa_4^\mu(f)-3V(f),
\]
which equals $\int\mathrm{d}\nu(\xi)\,\kappa_4^\xi(f)$ by \eqref{eq:fourth-cumulants-4}. Then $V_k(f)=0$ for every $k\ge1$, and $\delta_V(f)=\kappa_4^\mu(f)$.
\item[(v)] Let $f\in\mathcal S$, and let $e>0$, $\sigma>0$, $C_f<\infty$ and $b(f)>0$ be numbers such that \eqref{eq:fourth-cumulants-4g} holds for the test $g=f$ and $C_f\cdot e^\sigma\le\frac{1}{2}b(f)$. Then
\[
\kappa_4^\xi(f)=\kappa_4^\mu(f)\ge\frac{1}{2}e^4\cdot b(f)\qquad\text{for $\nu$-almost every }\xi\in\Xi .
\]
\end{enumerate}
\end{corollary}

\begin{proof}
By \eqref{eq:fourth-cumulants-H2}, applied to the motions $(x,1)$, $x\in\mathbb{R}^3$, the law $\mu$ is translation invariant. By \eqref{eq:fourth-cumulants-H1} and Lemma~\ref{lem:moment-bounds}(d), for all $f\in\mathcal S$ and $t>0$,
\[
E\exp\left(t\cdot|X(f)|\right)\le2\exp\left(\frac{1}{2}t^2\cdot\|f\|_2^2\right)<\infty .
\]
Thus Theorem~\ref{thm:clustering-mixing-ergodicity} applies to $\mu$, and by that theorem \eqref{eq:fourth-cumulants-C1} implies \eqref{eq:fourth-cumulants-M}, \eqref{eq:fourth-cumulants-Erg} and \eqref{eq:fourth-cumulants-T}. This proves (i) and the first assertion of (ii).

(ii) Let $N\in B$ be a $\nu$-null set such that $\mu_\xi=\mu$ for every $\xi\in\Xi\setminus N$. For $A\in\mathcal F$ the composition of the map $\xi\mapsto\mu_\xi$ with the map $\rho\mapsto\rho(A)$ is the function $\xi\mapsto\mu_\xi(A)$, which is $B$-measurable by \eqref{eq:fourth-cumulants-ED1}; as the maps $\rho\mapsto\rho(A)$ generate the $\sigma$-algebra on the set of laws, the map $\xi\mapsto\mu_\xi$ is measurable. Let $Y$ be a measurable set of laws, and let $\Xi_Y=\left\{\xi\in\Xi:\mu_\xi\in Y\right\}$ be its preimage. If $\mu\in Y$, then $\Xi_Y\supset\Xi\setminus N$, and $\nu(\Xi_Y)=1$; if $\mu\notin Y$, then $\Xi_Y\subset N$, and $\nu(\Xi_Y)=0$. So the image of $\nu$ is the point mass at $\mu$. If $\Xi$ is the set of all laws, $B$ is this $\sigma$-algebra and $\mu_\xi=\xi$ for every $\xi$, then $\xi\mapsto\mu_\xi$ is the identity map of $\Xi$, and the image of $\nu$ is $\nu$ itself.

(iii) For $\xi\in\Xi\setminus N$ the law $\mu_\xi$ is $\mu$. It therefore satisfies \eqref{eq:fourth-cumulants-H1}--\eqref{eq:fourth-cumulants-H3}; and the joint moments of finitely many of the variables $X(f)$, $f\in\mathcal S$, which are finite by Lemma~\ref{lem:moment-bounds}(e), and with them the joint cumulants, are the same under $\mu_\xi$ as under $\mu$.

(iv) Let $\Xi_0$ be the set of Corollary~\ref{cor:properties-almost-every}, let $f\in\mathcal S$, and let $k\ge1$. For $\xi\in\Xi_0$ the law $\mu_\xi$ satisfies \eqref{eq:fourth-cumulants-H1}, and the Cauchy--Schwarz inequality together with Lemma~\ref{lem:moment-bounds}(c), applied to $\mu_\xi$, gives
\[
\left|E_\xi\left[X(f)^k\right]\right|\le\left(E_\xi\left[X(f)^{2k}\right]\right)^{\frac{1}{2}}\le c_{2k}^{\frac{1}{2}}\cdot\|f\|_2^k .
\]
By \eqref{eq:fourth-cumulants-ED1}, extended by linearity and monotone convergence to the positive and the negative part of $X(f)^k$, the function $\xi\mapsto E_\xi\left[X(f)^k\right]$ is measurable on $\Xi_0$; being bounded there, it has a variance under $\nu$, and $V_k(f)$ is defined. As $C_\xi(f,f)=E_\xi\left[X(f)^2\right]$, one has $V_2(f)=V(f)$. For $\xi\in\Xi_0\setminus N$,
\begin{gather*}
C_\xi(f_1,f_2)=S_2^\mu(f_1,f_2)\qquad(f_1,f_2\in\mathcal S),\\
E_\xi\left[X(f)^k\right]=E\left[X(f)^k\right].
\end{gather*}
Hence the functions $\xi\mapsto C_\xi(f_1,f_2)$ and $\xi\mapsto E_\xi\left[X(f)^k\right]$ are constant on a set of full $\nu$-measure, and their covariances and variances under $\nu$ vanish: $\Gamma(f_1,f_2;f_3,f_4)=0$, $V(f)=\Gamma(f,f;f,f)=0$ and $V_k(f)=0$. Consequently $\delta_V(f)=\kappa_4^\mu(f)-3V(f)=\kappa_4^\mu(f)$, and by \eqref{eq:fourth-cumulants-4} this number equals $\int\mathrm{d}\nu(\xi)\,\kappa_4^\xi(f)$.

(v) By \eqref{eq:fourth-cumulants-4g} for the test $f$ and by the inequality $C_f\cdot e^\sigma\le\frac{1}{2}b(f)$,
\[
\kappa_4^\mu(f)\ge e^4\cdot b(f)-C_f\cdot e^{4+\sigma}=e^4\cdot\left(b(f)-C_f\cdot e^\sigma\right)\ge\frac{1}{2}e^4\cdot b(f).
\]
For $\xi\in\Xi\setminus N$ one has $\mu_\xi=\mu$, and therefore $\kappa_4^\xi(f)=\kappa_4^\mu(f)$.
\end{proof}

\section{Partition functions of the all-flux law: trace formula, Gaussian domination, convexity}\label{sec:partition-functions-all}

For lattice tori with arbitrary sides $n_{\mu}\ge2$ and a lattice mass $\boldsymbol{\mu}>-1$, it is shown here that the sector of flux zero of the all-flux law is the harmonic-averaged law (Lemma~\ref{lem:flux-sectors}) and that the shifted partition functions $Z_n(h)$ are traces (Proposition~\ref{prop:trace-formula}); the trace formula then gives Gaussian domination (Theorem~\ref{thm:gaussian-domination}) with its consequences, among them Estimate~1 for the all-flux law (Corollary~\ref{cor:partition-functions-1}), and the convexity and the subadditivity of $\log Z_n[1]$ in each side length (Theorem~\ref{thm:partition-functions-1}), from which a lower bound along the diagonal follows (Lemma~\ref{lem:partition-functions-1}). Section~\ref{sec:all-flux-functional} returns to the all-flux law, with observables in place of partition functions.

\subsection{Cochains, fermion weight and flux sectors}\label{sec:cochains-fermion-weight}

\subsubsection{Tori and cochains}\label{sec:tori-cochains}

Let $n=(n_1,n_2,n_3)$ be a triple of integers $n_{\mu}\ge2$, let $e_1,e_2,e_3$ be the unit vectors of $\mathbb{Z}^3$, and put
\[
\Lambda_n=\prod_{\mu=1}^3\mathbb{Z}/n_{\mu}\mathbb{Z},\qquad|n|=n_1n_2n_3.
\]
For $x\in\Lambda_n$ and a set $I=\{i_1,\dots,i_j\}\subset\{1,2,3\}$ with $i_1<\dots<i_j$, the $j$-cell $(x;I)$ is the cell spanned at $x$ by $e_{i_1},\dots,e_{i_j}$; it is positively oriented by $e_{i_1}\wedge\dots\wedge e_{i_j}$. The $j$-cells with $j=0,1,2,3$ are the sites, the links, the plaquettes and the cubes. The sets of the links and of the plaquettes are denoted $\mathrm{Lk}_n$ and $\mathrm{Pl}_n$, and $p_{\mu\nu}(x)=(x;\{\mu,\nu\})$ for $\mu<\nu$. For $j=0,1,2,3$ let $C^j_n$ be the real vector space of the $j$-cochains, that is, of the real functions on the $j$-cells. It carries the counting inner product $\langle u,v\rangle$, the sum over all $j$-cells of the product of the values of $u$ and $v$, and the norm $\|u\|=\langle u,u\rangle^{\frac{1}{2}}$. The integer-valued $j$-cochains form the subgroup $C^j_n(\mathbb{Z})$. The value of $\lambda\in C^0_n$ at the site $x$ is $\lambda(x)$; the value of $A\in C^1_n$ on the link from $x$ to $x+e_{\mu}$ is $A_{\mu}(x)$; the value of $h\in C^2_n$ on a plaquette $p$ is $h(p)$, and $h_{\mu\nu}(x)=h\left(p_{\mu\nu}(x)\right)$; the value of a three-cochain on the cube $(x;\{1,2,3\})$ carries the index $123$. The coboundary $d\colon C^j_n\to C^{j+1}_n$, $j=0,1,2$, is given by
\begin{align*}
(d\lambda)_{\mu}(x)&=\lambda(x+e_{\mu})-\lambda(x),\\
(dA)_{\mu\nu}(x)&=A_{\mu}(x)+A_{\nu}(x+e_{\mu})-A_{\mu}(x+e_{\nu})-A_{\nu}(x)\qquad(\mu<\nu),\\
(dh)_{123}(x)&=\left[h_{23}(x+e_1)-h_{23}(x)\right]-\left[h_{13}(x+e_2)-h_{13}(x)\right]+\left[h_{12}(x+e_3)-h_{12}(x)\right],
\end{align*}
and $d^{*}$ is its adjoint with respect to the counting inner products. A cochain $u$ is called \emph{closed} if $du=0$, and \emph{constant} if, for every $I$, its value on the cell $(x;I)$ does not depend on $x$.

Let $\mathrm{Harm}^2_n\subset C^2_n$ be the space of the constant two-cochains; thus $h\in C^2_n$ belongs to $\mathrm{Harm}^2_n$ if and only if, for each pair $\mu<\nu$, the number $h_{\mu\nu}(x)$ is the same for all $x\in\Lambda_n$. In degree one put $V_{\mathrm{g}}=dC^0_n$, let $V_{\mathrm{h}}\subset C^1_n$ be the space of the constant one-cochains, the $A$ for which $A_{\mu}(x)$ does not depend on $x$ ($\mu=1,2,3$), and let $V_{\mathrm{t}}$ be the orthogonal complement of $V_{\mathrm{g}}+V_{\mathrm{h}}$ in $C^1_n$. Then, all direct sums being orthogonal,
\begin{gather}
\ker\left(d|_{C^2_n}\right)=dC^1_n\oplus\mathrm{Harm}^2_n,\notag\\
C^1_n=V_{\mathrm{g}}\oplus V_{\mathrm{h}}\oplus V_{\mathrm{t}},\qquad\ker\left(d|_{C^1_n}\right)=V_{\mathrm{g}}\oplus V_{\mathrm{h}}.\label{eq:partition-functions-1}
\end{gather}
By \eqref{eq:partition-functions-1}, $V_{\mathrm{g}}$, $V_{\mathrm{h}}$ and $V_{\mathrm{t}}$ are the three spaces of the same names in Section~\ref{sec:cochains-periodic-lattice}. The identities \eqref{eq:partition-functions-1} form the Hodge decomposition on the discrete torus (Section~\ref{sec:results-literature}); in degree one they are Lemma~\ref{lem:harmonic-cochains-coercivity} together with Remark~\ref{rem:tori-unequal-sides}.

\begin{proof}[Proof of \eqref{eq:partition-functions-1}]
It is to be shown, for $j=1$ and for $j=2$, that the coboundaries and the constant cochains of degree $j$ are closed, that every coboundary of degree $j$ is orthogonal to every constant $j$-cochain, and that every closed $u\in C^j_n$ is of the form $u=dv+u_0$ with $v\in C^{j-1}_n$ and a constant $u_0\in C^j_n$. For $j=2$ this is the first identity of \eqref{eq:partition-functions-1}. For $j=1$ it is the third one, as $dC^0_n=V_{\mathrm{g}}$ and the constant one-cochains form $V_{\mathrm{h}}$; the second identity then follows from the definition of $V_{\mathrm{t}}$.

In every degree $j\in\{0,1,2,3\}$ consider the complex cochains $C^j_n\otimes\mathbb{C}$, with the counting inner product extended to a Hermitian form, and extend $d$ to them complex-linearly. A frequency is a triple $k=(k_1,k_2,k_3)$ with $k_{\mu}\in\frac{2\pi}{n_{\mu}}\mathbb{Z}$, taken modulo $2\pi$; there are $|n|$ frequencies. For $j\in\{0,1,2,3\}$ let $I$ range over the subsets $\{i_1<\dots<i_j\}$ of $\{1,2,3\}$ with $j$ elements, and put $e_I=e_{i_1}\wedge\dots\wedge e_{i_j}$ in the exterior power $\bigwedge^j\mathbb{C}^3$; for $j=0$ the only such set is $I=\emptyset$, and $e_{\emptyset}=1\in\bigwedge^0\mathbb{C}^3=\mathbb{C}$. For a frequency $k$ and $c=\sum_Ic_I\cdot e_I\in\bigwedge^j\mathbb{C}^3$ let $\boldsymbol{\omega}_{k,c}\in C^j_n\otimes\mathbb{C}$ be the cochain with the values
\[
\boldsymbol{\omega}_{k,c}\left((x;I)\right)=c_I\cdot\exp\left(i\sum_{\mu=1}^3k_{\mu}x_{\mu}\right).
\]
In each degree these cochains span $C^j_n\otimes\mathbb{C}$ (finite Fourier transform), and two of them with different frequencies are orthogonal. The shift $x\mapsto x+e_{\mu}$ multiplies the exponential by $\exp(ik_{\mu})$; therefore the three formulas for $d$ give, for $j\in\{0,1,2\}$, every frequency $k$ and every $c\in\bigwedge^j\mathbb{C}^3$,
\[
d\,\boldsymbol{\omega}_{k,c}=\boldsymbol{\omega}_{k,\xi\wedge c}\in C^{j+1}_n\otimes\mathbb{C},\qquad\xi=\xi(k)=\sum_{\mu=1}^3\left(\exp(ik_{\mu})-1\right)\cdot e_{\mu}.
\]
On the cochains of frequency $k$ the coboundary is thus the exterior multiplication by $\xi$; in particular $dd=0$ on $C^0_n$ and on $C^1_n$, since $\xi\wedge\xi=0$. Let $\mathfrak{i}$ be the contraction with $\bar{\xi}$, the linear map of the exterior algebra $\bigwedge\mathbb{C}^3=\bigoplus_{j=0}^3\bigwedge^j\mathbb{C}^3$ that vanishes on $\bigwedge^0\mathbb{C}^3$ and maps $\bigwedge^j\mathbb{C}^3$ into $\bigwedge^{j-1}\mathbb{C}^3$ for $1\le j\le3$ by
\[
\mathfrak{i}\left(e_{i_1}\wedge\dots\wedge e_{i_j}\right)=\sum_{l=1}^j(-1)^{l-1}\cdot\overline{\xi_{i_l}}\cdot e_{i_1}\wedge\dots\wedge e_{i_{l-1}}\wedge e_{i_{l+1}}\wedge\dots\wedge e_{i_j}.
\]
Contraction and exterior multiplication satisfy the anticommutation relation
\[
\mathfrak{i}(\xi\wedge c)+\xi\wedge\mathfrak{i}(c)=|\xi|^2\cdot c\qquad\left(c\in\bigwedge\mathbb{C}^3\right),\qquad|\xi|^2=\sum_{\mu=1}^3|\xi_{\mu}|^2.
\]
Now $\xi(k)=0$ if and only if $k=0$. Let $j\in\{1,2\}$, let $k$ be a frequency and $c\in\bigwedge^j\mathbb{C}^3$, and let $\boldsymbol{\omega}_{k,c}$ be closed, that is, $\xi\wedge c=0$. If $k\ne0$, then $c=\xi\wedge b$ with $b=|\xi|^{-2}\cdot\mathfrak{i}(c)\in\bigwedge^{j-1}\mathbb{C}^3$, so that $\boldsymbol{\omega}_{k,c}=d\,\boldsymbol{\omega}_{k,b}$ is the coboundary of a cochain of degree $j-1$ with the same frequency. If $k=0$, then $\xi=0$: the cochains of frequency $0$ are the constant cochains, all of them are closed, and, $d$ preserving the frequency, no coboundary has a non-zero component of frequency $0$. Splitting a closed complex cochain of degree $j$ into its components of the different frequencies, each of which is closed, one obtains the three assertions for complex cochains, with $v$ complex; the closedness of the coboundaries is $dd=0$, and the orthogonality is that of different frequencies.

Complex conjugation maps the cochains of frequency $k$ onto those of frequency $-k$, commutes with $d$, whose coefficients are real, and maps constant cochains onto constant cochains. The first two assertions for real cochains are special cases of those for complex cochains. If $u\in C^j_n$ is closed and $u=dv+u_0$ with $v\in C^{j-1}_n\otimes\mathbb{C}$ and a constant $u_0\in C^j_n\otimes\mathbb{C}$, then also $u=d\bar{v}+\overline{u_0}$, and hence $u=d(\operatorname{Re}v)+\operatorname{Re}u_0$. This proves the third assertion for real cochains.
\end{proof}

For $\mu<\nu$ let $\kappa$ be the third direction. For $t\in\mathbb{Z}/n_{\kappa}\mathbb{Z}$ the $n_{\mu}n_{\nu}$ plaquettes $p_{\mu\nu}(x)$ with $x_{\kappa}=t$ form the \emph{coordinate two-torus} $\Pi_{\mu\nu}(t)$; the $n_{\kappa}$ coordinate two-tori $\Pi_{\mu\nu}(t)$ are parallel to each other. Let moreover $\Lambda_s$ be a two-dimensional torus, a product of two cyclic groups of orders at least $2$, with its cells, its cochains $C^j(\Lambda_s)$, its integer-valued cochains $C^j(\Lambda_s;\mathbb{Z})$ and its coboundary $d_s$ defined as for $\Lambda_n$. Then, for every closed $z\in C^2_n(\mathbb{Z})$ and every $z_s\in C^2(\Lambda_s;\mathbb{Z})$,
\begin{multline}\label{eq:partition-functions-2}
\sum_{p\in\Pi_{\mu\nu}(t)}z(p)=\sum_{p\in\Pi_{\mu\nu}(0)}z(p)=:w_{\mu\nu}(z)\qquad\left(\mu<\nu,\ t\in\mathbb{Z}/n_{\kappa}\mathbb{Z}\right),\\
w(z):=\left(w_{12}(z),w_{13}(z),w_{23}(z)\right)\in\mathbb{Z}^3,\qquad z\in d\,C^1_n(\mathbb{Z})\iff w(z)=0,\\
z_s\in d_s\,C^1(\Lambda_s;\mathbb{Z})\iff\sum_pz_s(p)=0,
\end{multline}
where the last sum is taken over all plaquettes $p$ of $\Lambda_s$. The vector $w(z)$ is called the \emph{flux} of $z$.

\begin{proof}[Proof of \eqref{eq:partition-functions-2}]
Let $\mu<\nu$ and $t\in\mathbb{Z}/n_{\kappa}\mathbb{Z}$. The cubes $(x;\{1,2,3\})$ with $x_{\kappa}=t$ form the slab between $\Pi_{\mu\nu}(t)$ and $\Pi_{\mu\nu}(t+1)$. Sum the formula for $(dh)_{123}$, with $h=z$, over these cubes. The two brackets that do not contain $z_{\mu\nu}$ contribute $0$, each of their values occurring once with each sign, and the remaining bracket contributes
\[
\pm\left(\sum_{p\in\Pi_{\mu\nu}(t+1)}z(p)-\sum_{p\in\Pi_{\mu\nu}(t)}z(p)\right).
\]
As $dz=0$, this difference vanishes for every $t$, which proves the first identity.

The cells of $\Lambda_n$ form a cell decomposition of the torus $\mathbb{T}^3=\mathbb{R}^3/\left(n_1\mathbb{Z}\times n_2\mathbb{Z}\times n_3\mathbb{Z}\right)$, the product of three circles divided into $n_1$, $n_2$ and $n_3$ arcs, and the groups $C^0_n(\mathbb{Z}),\dots,C^3_n(\mathbb{Z})$ with the maps $d$ form its cellular cochain complex with integer coefficients. Hence the closed elements of $C^2_n(\mathbb{Z})$ modulo $d\,C^1_n(\mathbb{Z})$ form the group $H^2(\mathbb{T}^3;\mathbb{Z})$. Since $H_1(\mathbb{T}^3;\mathbb{Z})$ is free, the universal coefficient theorem shows that the evaluation of cohomology classes on two-cycles is an isomorphism of $H^2(\mathbb{T}^3;\mathbb{Z})$ onto the group of the homomorphisms from $H_2(\mathbb{T}^3;\mathbb{Z})$ to $\mathbb{Z}$. The K\"unneth formula, applied to this product cell structure, shows that $H_2(\mathbb{T}^3;\mathbb{Z})$ is generated by the classes of the three coordinate two-tori $\Pi_{12}(0)$, $\Pi_{13}(0)$, $\Pi_{23}(0)$, each taken as the sum of its positively oriented plaquettes (see Section~\ref{sec:results-literature}). The class of $z$ takes the value $w_{\mu\nu}(z)$ on $\Pi_{\mu\nu}(0)$. Hence this class vanishes, that is, $z\in d\,C^1_n(\mathbb{Z})$, if and only if $w(z)=0$.

On the two-dimensional torus, every value of $m\in C^1(\Lambda_s;\mathbb{Z})$ occurs in $\sum_p(d_sm)(p)$ once with each sign, so that this sum vanishes. Conversely, every link $b$ of $\Lambda_s$ lies in exactly two plaquettes, and the value $m(b)$ enters $d_sm$ at one of them with the sign $+$ and at the other with the sign $-$. Let $\mathcal{G}$ be the graph whose vertices are the plaquettes of $\Lambda_s$ and which has, for every link $b$, an edge directed from the plaquette at which $m(b)$ enters with the sign $-$ to the plaquette at which it enters with the sign $+$. A one-cochain $m$ is then a flow on $\mathcal{G}$, and $(d_sm)(p)$ is the inflow of this flow into $p$ minus its outflow from $p$, that is, its divergence at $p$. The graph $\mathcal{G}$ is connected, since the plaquettes spanned at $x$ and at $x+e$, for a unit vector $e$ of $\Lambda_s$, have a link in common. Suppose now that $\sum_pz_s(p)=0$. An integer flow with the divergence $z_s$ is constructed along a spanning tree of $\mathcal{G}$, as follows. For a vertex $u$ let $\delta_u$ be the function on the vertices that equals $1$ at $u$ and $0$ elsewhere. Fix a vertex $r$ and a spanning tree. For a vertex $v\ne r$ let $\phi_v$ be the flow that vanishes off the path from $r$ to $v$ in the tree and equals $+1$ on the edges of this path that point from $r$ towards $v$ and $-1$ on those that point from $v$ towards $r$. Along the path the contributions cancel at the inner vertices, so that the divergence of $\phi_v$ is $\delta_v-\delta_r$. Hence the integer flow $\sum_{v\ne r}z_s(v)\cdot\phi_v$ has the divergence
\[
\sum_{v\ne r}z_s(v)\cdot\left(\delta_v-\delta_r\right)=z_s-\left(\sum_pz_s(p)\right)\cdot\delta_r=z_s;
\]
read as an integer one-cochain $m$, this flow satisfies $d_sm=z_s$.
\end{proof}

\subsubsection{The Wilson determinant}\label{sec:wilson-determinant}

From now on $q>0$ and $\boldsymbol{\mu}>-1$. (In Theorem~\ref{thm:main-theorem} one has $q=e\eta^{\frac{1}{2}}$, and $\boldsymbol{\mu}=\bar{m}\cdot\eta+m_{0,N}$ is the lattice mass, with the charge $e$, the lattice spacing $\eta$, the mass $\bar{m}$ and the numbers $m_{0,N}$ of that theorem.) Let $\sigma_1,\sigma_2,\sigma_3$ be the Pauli matrices, and for $\mu=1,2,3$ let
\[
\gamma_{\mu}=\operatorname{diag}\left(\sigma_{\mu},-\sigma_{\mu}\right),\qquad P^{\pm}_{\mu}=\frac{1}{2}\left(1\pm\gamma_{\mu}\right).
\]
The \emph{spinors} are the functions $\psi\colon\mathbb{Z}^3\to\mathbb{C}^4$ with $\psi(x+n_{\mu}e_{\mu})=-\psi(x)$ for $\mu=1,2,3$; they form a complex vector space of dimension $4|n|$. Through the projection of $\mathbb{Z}^3$ onto $\Lambda_n$, every $A\in C^1_n$ assigns a number $A_{\mu}(x)$ to every $x\in\mathbb{Z}^3$ and every $\mu$, and this number does not change when $x$ is translated by an element of $n_1\mathbb{Z}\times n_2\mathbb{Z}\times n_3\mathbb{Z}$. The Wilson operator $D_A$ is the linear map of the space of spinors into itself given by
\begin{multline*}
(D_A\psi)(x)=3\psi(x)-\sum_{\mu=1}^3\left[P^{-}_{\mu}\cdot\exp\left(iqA_{\mu}(x)\right)\cdot\psi(x+e_{\mu})\right.\\
\left.+P^{+}_{\mu}\cdot\exp\left(-iqA_{\mu}(x-e_{\mu})\right)\cdot\psi(x-e_{\mu})\right].
\end{multline*}
The Wilson parameter in this formula equals one, and the phase attached to the link from $x$ to $x+e_{\mu}$ is $\exp\left(iqA_{\mu}(x)\right)$. This is the operator $D[qA]$ of \eqref{eq:free-coefficient-5}, and the operator $D_A$ of \eqref{eq:gaussian-bounds-1} at the antiperiodic character. The \emph{fermion weight} is the function
\[
\mathfrak{w}(A)=\det\left(D_A+\boldsymbol{\mu}\right)\qquad\left(A\in C^1_n\right),
\]
the determinant being taken on the space of spinors.

The fermion weight has the following properties. For all $A\in C^1_n$,
\begin{gather}
\mathfrak{w}(A)\ge0,\label{eq:partition-functions-3}\\
\mathfrak{w}(A+d\lambda)=\mathfrak{w}(A)\qquad\left(\lambda\in C^0_n\right),\label{eq:partition-functions-4}\\
\mathfrak{w}\left(A+\frac{2\pi}{q}m\right)=\mathfrak{w}(A)\qquad\left(m\in C^1_n(\mathbb{Z})\right).\label{eq:partition-functions-5}
\end{gather}
Moreover,
\begin{multline}\label{eq:partition-functions-6}
\mathfrak{w}\ \text{is continuous and bounded on}\ C^1_n\ \text{and has an entire extension to}\ C^1_n\otimes\mathbb{C},\\
\text{also written}\ \mathfrak{w},\ \text{for which}\ \eqref{eq:partition-functions-4}\ \text{and}\ \eqref{eq:partition-functions-5}\ \text{remain valid with complex}\ A,\ \text{and}\\
\sup\left\{\left|\mathfrak{w}(A)\right|:\ A\in C^1_n\otimes\mathbb{C},\ \max_{x,\mu}\left|\operatorname{Im}A_{\mu}(x)\right|\le c\right\}<\infty\qquad(c\ge0).
\end{multline}
Properties \eqref{eq:partition-functions-3}--\eqref{eq:partition-functions-6} are Theorem~\ref{thm:free-coefficient-1}(d1),(d2) for $B=qA$. By Theorem~\ref{thm:free-coefficient-1}(d3),
\begin{equation}\label{eq:partition-functions-7}
\mathfrak{w}(-A)=\mathfrak{w}(A)\qquad\left(A\in C^1_n\right).
\end{equation}
Finally, let $r$ be a signed permutation matrix, let $c\in\mathbb{Z}^3$, let $n'=(n'_1,n'_2,n'_3)$ be a triple of integers $n'_{\mu}\ge2$ with $r\left(\prod_{\mu}n_{\mu}\mathbb{Z}\right)=\prod_{\mu}n'_{\mu}\mathbb{Z}$, and let $g$ be the map $gx=rx+c$ of $\mathbb{Z}^3$. It induces an isometry of $\Lambda_n$ onto $\Lambda_{n'}$, again denoted $g$. One-cochains are read as odd functions on the oriented links, $A(x,x+e_{\mu})=A_{\mu}(x)$ and $A(x',x)=-A(x,x')$, and $g_{*}A\in C^1_{n'}$ is given by $(g_{*}A)(gx,gx')=A(x,x')$. With $\mathfrak{w}_n$ and $\mathfrak{w}_{n'}$ the fermion weights of $\Lambda_n$ and of $\Lambda_{n'}$, Theorem~\ref{thm:free-coefficient-1}(d4) gives
\begin{equation}\label{eq:partition-functions-8}
\mathfrak{w}_{n'}(g_{*}A)=\mathfrak{w}_n(A)\qquad\left(A\in C^1_n\right).
\end{equation}

\subsubsection{Summation over all fluxes}\label{sec:summation-all-fluxes}

Let $\mathrm{Gr}$ be the class of the Borel functions $G\colon C^2_n\to\mathbb{C}$ for which there are constants $c,c'\ge0$ with $|G(F)|\le c\cdot\exp\left(c'\|F\|\right)$ for all $F\in C^2_n$.

\begin{definition}[Partition functions and the all-flux law]\label{def:partition-functions-all}
For $G\in\mathrm{Gr}$ let
\[
Z_n[G]=\sum_{\substack{\nu\in C^2_n(\mathbb{Z})\\ d\nu=0}}\ \int_{\left[0,\frac{2\pi}{q}\right)^{\mathrm{Lk}_n}}\mathrm{d}A\,\exp\left(-\frac{1}{2}\|F\|^2\right)\cdot\mathfrak{w}(A)\cdot G(F),
\]
where $F=dA+\frac{2\pi}{q}\nu$, and $E_n[G]=\frac{Z_n[G]}{Z_n[1]}$; the series converges absolutely, and $Z_n[1]>0$, as shown below. For $h\in C^2_n$ let
\[
Z_n(h)=Z_n\left[\exp\left(-\langle h,F\rangle-\frac{1}{2}\|h\|^2\right)\right];
\]
this is $Z_n[1]$ with $\|F\|^2$ replaced by $\|F+h\|^2$. The \emph{all-flux law}~\eqref{eq:setup-af} is the distribution of the closed real two-cochain $F$, regarded as a random variable with the expectation $E_n$.
\end{definition}

The series in Definition~\ref{def:partition-functions-all} converges absolutely. Let $c,c'\ge0$ be constants for $G$ as above, and let $M$ be the supremum of $\|dA\|$ over the box $\left[0,\frac{2\pi}{q}\right)^{\mathrm{Lk}_n}$, which is finite. For $A$ in the box, the triangle inequality gives $\|F\|\ge t_{\nu}$ with $t_{\nu}=\frac{2\pi}{q}\|\nu\|-M$; the function $\mathfrak{w}$ is bounded by \eqref{eq:partition-functions-6}; and $s\mapsto c's-\frac{1}{2}s^2$ decreases on $[c',\infty)$. Hence, if $t_{\nu}\ge c'$, the modulus of the term of $\nu$ is at most
\[
\left(\frac{2\pi}{q}\right)^{3|n|}\cdot c\cdot\sup_{C^1_n}\left|\mathfrak{w}\right|\cdot\exp\left(c't_{\nu}-\frac{1}{2}t_{\nu}^2\right),
\]
the first factor being the volume of the box. This bound decreases faster than every power of $\|\nu\|$, whereas the number of integer two-cochains in a ball grows only polynomially with the radius,
\[
\#\left|\left\{\nu\in C^2_n(\mathbb{Z}):\ \|\nu\|\le r\right\}\right|\le(2r+1)^{3|n|};
\]
so the series converges absolutely.

Moreover $Z_n[1]>0$. All terms of $Z_n[1]$ are non-negative by \eqref{eq:partition-functions-3}, and the term $\nu=0$ is positive: its integrand $\exp\left(-\frac{1}{2}\|dA\|^2\right)\cdot\mathfrak{w}(A)$ is continuous and non-negative, and it is positive at $A=0$, hence on a subset of the box of positive measure, because
\[
\mathfrak{w}(0)=\prod_{k\in\mathcal{K}_n}\left(\left(\boldsymbol{\mu}+\sum_{\mu=1}^3\left(1-\cos k_{\mu}\right)\right)^2+\sum_{\mu=1}^3\sin^2k_{\mu}\right)^2>0.
\]
Here $\mathcal{K}_n$ is the set of the $|n|$ antiperiodic momenta, the $k\in(\mathbb{R}/2\pi\mathbb{Z})^3$ with $\exp(in_{\mu}k_{\mu})=-1$ for $\mu=1,2,3$. The spinors $x\mapsto u\cdot\exp\left(i\sum_{\mu}k_{\mu}x_{\mu}\right)$ with $u\in\mathbb{C}^4$ and $k\in\mathcal{K}_n$ span the space of spinors, and on those of momentum $k$ the operator $D_0+\boldsymbol{\mu}$ acts as the matrix
\[
a_k+i\sum_{\mu=1}^3\sin k_{\mu}\cdot\gamma_{\mu},\qquad a_k=\boldsymbol{\mu}+\sum_{\mu=1}^3\left(1-\cos k_{\mu}\right),
\]
whose two diagonal blocks $a_k\pm i\sum_{\mu}\sin k_{\mu}\cdot\sigma_{\mu}$ both have the determinant $a_k^2+\sum_{\mu}\sin^2k_{\mu}$. Every factor of the product is positive. If $\sin k_{\mu}=0$ for every $\mu$, then $k=(\pi,\pi,\pi)$, since an antiperiodic momentum $k_{\mu}$ is never $0$ modulo $2\pi$, and the first square in the factor is $(\boldsymbol{\mu}+6)^2>0$; otherwise $\sum_{\mu}\sin^2k_{\mu}>0$.

For $G\in\mathrm{Gr}$ the function $G^{-}(F)=G(-F)$ again belongs to $\mathrm{Gr}$, and $Z_n[G^{-}]=Z_n[G]$; thus $F$ and $-F$ have the same distribution under $E_n$. This rests on \eqref{eq:partition-functions-7} and is obtained by the substitutions $A\mapsto-A$ modulo $\frac{2\pi}{q}$ and $\nu\mapsto-\nu$. On the open box the first substitution is $A\mapsto\frac{2\pi}{q}\mathbf{1}-A$, where $\mathbf{1}\in C^1_n$ takes the value $1$ on every link; it preserves the Lebesgue measure, and $d\mathbf{1}=0$, so that $d\left(\frac{2\pi}{q}\mathbf{1}-A\right)-\frac{2\pi}{q}\nu=-F$. Moreover $\mathfrak{w}\left(\frac{2\pi}{q}\mathbf{1}-A\right)=\mathfrak{w}(-A)=\mathfrak{w}(A)$ by \eqref{eq:partition-functions-5} and \eqref{eq:partition-functions-7}.

Let $g\colon\Lambda_n\to\Lambda_{n'}$ be an isometry as in \eqref{eq:partition-functions-8}. It maps oriented plaquettes onto oriented plaquettes. Reading two-cochains as odd functions of the oriented plaquettes, let $g^{*}\colon C^2_{n'}\to C^2_n$ be the pull-back, $(g^{*}h)(p)=h(gp)$. Since $g^{*}$ is orthogonal, $G\circ g^{*}$ belongs to the class $\mathrm{Gr}$ of $\Lambda_{n'}$ for every $G\in\mathrm{Gr}$, and
\[
Z_{n'}\left[G\circ g^{*}\right]=Z_n[G]\qquad\left(G\in\mathrm{Gr}\right);
\]
the same holds for the parts of the two series with flux zero (Section~\ref{sec:decomposition-flux}). This follows from \eqref{eq:partition-functions-8} (Theorem~\ref{thm:free-coefficient-1}(d5)). If $g$ is a translation of $\Lambda_n$, or if $n_1=n_2=n_3$ and $g$ is an element of the cubic group $O_{\mathrm{cub}}$, then $n'=n$, and the identity gives $E_n[G\circ g^{*}]=E_n[G]$ for all $G\in\mathrm{Gr}$: the distribution of $F$ under $E_n$ is not changed by these maps.

Let $U$ be a contractible union of closed cells of $\Lambda_n$, and let $\nu\in C^2_n(\mathbb{Z})$ be closed. The restriction of $\nu$ to the plaquettes of $U$ is closed on $U$; as $H^2(U;\mathbb{Z})=0$, it is the coboundary of an integer one-cochain on the links of $U$, which extends to some $m\in C^1_n(\mathbb{Z})$. Then $\nu-dm$ vanishes on the plaquettes of $U$, the pair $\left(A+\frac{2\pi}{q}m,\nu-dm\right)$ gives the same $F$ as the pair $(A,\nu)$, and $\mathfrak{w}\left(A+\frac{2\pi}{q}m\right)=\mathfrak{w}(A)$ by \eqref{eq:partition-functions-5}. On $U$ the pair $(A,\nu)$ is thus equivalent to $\left(A+\frac{2\pi}{q}m,0\right)$, and there the integrand of Definition~\ref{def:partition-functions-all} is the weight $\exp\left(-\frac{1}{2}\|dA\|^2\right)\cdot\mathfrak{w}(A)$ of the noncompact gauge field without fluxes.

\subsubsection{Decomposition by flux}\label{sec:decomposition-flux}

For $w=(w_{12},w_{13},w_{23})\in\mathbb{Z}^3$ and $G\in\mathrm{Gr}$ let $Z_{n,w}[G]$ be the sum of the terms of the series in Definition~\ref{def:partition-functions-all} with $w(\nu)=w$, and put
\[
P_n(w)=\frac{Z_{n,w}[1]}{Z_n[1]},\qquad E_{n,w}[G]=\frac{Z_{n,w}[G]}{Z_{n,w}[1]}.
\]
Here $Z_{n,w}[1]>0$ for every $w$. Indeed, let $z_w\in C^2_n(\mathbb{Z})$ be the two-cochain with $(z_w)_{\mu\nu}(x)=w_{\mu\nu}$ if $x_{\mu}=x_{\nu}=0$ and $(z_w)_{\mu\nu}(x)=0$ otherwise ($\mu<\nu$). Each component of $z_w$ is independent of the third coordinate, so that the three brackets in the formula for $(dz_w)_{123}$ vanish and $z_w$ is closed; its flux is $w(z_w)=w$; and its term in $Z_{n,w}[1]$ is positive by the argument given for the term $\nu=0$ of $Z_n[1]$, while the other terms are non-negative. Let $h_w\in\mathrm{Harm}^2_n$ be the constant two-cochain with the components
\[
(h_w)_{\mu\nu}(x)=\frac{2\pi w_{\mu\nu}}{qn_{\mu}n_{\nu}}\qquad(\mu<\nu),
\]
and let
\[
\Gamma_n=\frac{2\pi}{q}\left\{m\in C^1_n(\mathbb{Z}):\ dm=0\right\}.
\]

\begin{lemma}[The flux sectors]\label{lem:flux-sectors}
In (a) and (c), let $\nu\in C^2_n(\mathbb{Z})$ be closed, let $w=w(\nu)$, let $A\in C^1_n$, and let $F=dA+\frac{2\pi}{q}\nu$.
\begin{enumerate}
\item[(a)] There is $\beta_{\nu}\in C^1_n$ with $\frac{2\pi}{q}\nu=d\beta_{\nu}+h_w$. Hence $F=d(A+\beta_{\nu})+h_w$, and
\[
\|F+c\|^2=\left\|d(A+\beta_{\nu})\right\|^2+\|h_w+c\|^2\qquad\left(c\in\mathrm{Harm}^2_n\right).
\]
Moreover
\[
\|h_w\|^2=\frac{4\pi^2}{q^2}\left(\frac{n_3}{n_1n_2}\cdot w_{12}^2+\frac{n_2}{n_1n_3}\cdot w_{13}^2+\frac{n_1}{n_2n_3}\cdot w_{23}^2\right);
\]
in particular, if $n_1=n_2=n_3$, then $\|h_w\|^2=\frac{4\pi^2|w|^2}{q^2n_1}$ with $|w|^2=w_{12}^2+w_{13}^2+w_{23}^2$.
\item[(b)] For $G\in\mathrm{Gr}$,
\[
Z_{n,0}[G]=\int_{C^1_n/\Gamma_n}\mathrm{d}A\,\exp\left(-\frac{1}{2}\|dA\|^2\right)\cdot\mathfrak{w}(A)\cdot G(dA).
\]
Let $\sigma$ be Lebesgue measure on $V_{\mathrm{t}}$, for $\theta\in[-\pi,\pi]^3$ let $H_{\theta}\in V_{\mathrm{h}}$ be the constant one-cochain with the value $\frac{\theta_{\mu}}{qn_{\mu}}$ on the links of direction $\mu$ ($\mu=1,2,3$), and for $G\in\mathrm{Gr}$ put
\[
Z^{\mathrm{av}}_n[G]=\int_{[-\pi,\pi]^3}\mathrm{d}\theta\int_{V_{\mathrm{t}}}\mathrm{d}\sigma(A)\,\exp\left(-\frac{1}{2}\|dA\|^2\right)\cdot\mathfrak{w}\left(A+H_{\theta}\right)\cdot G(dA).
\]
Then $E_{n,0}[G]=\frac{Z^{\mathrm{av}}_n[G]}{Z^{\mathrm{av}}_n[1]}$ for $G\in\mathrm{Gr}$, and $E_{n,0}$ coincides with the harmonic-averaged law of \eqref{eq:gaussian-bounds-3}, taken for rectangular tori as in Remark~\ref{rem:tori-unequal-sides}: under $E_{n,0}$ the field $F$ is distributed as $dA$, where $A$ is a random element of $V_{\mathrm{t}}$ whose distribution has, with respect to $\sigma$, a density proportional to $\int_{[-\pi,\pi]^3}\mathrm{d}\theta\,\exp\left(-\frac{1}{2}\|dA\|^2\right)\cdot\mathfrak{w}\left(A+H_{\theta}\right)$.
\item[(c)] If $f=da$ with $a\in C^1_n$, then $\langle f,F\rangle=\left\langle da,d(A+\beta_{\nu})\right\rangle$, whatever the flux $w$.
\end{enumerate}
\end{lemma}

\begin{proof}
(a) The real two-cochain $\frac{2\pi}{q}\nu$ is closed. By \eqref{eq:partition-functions-1} there are $\beta\in C^1_n$ and $c'\in\mathrm{Harm}^2_n$ with
\[
\frac{2\pi}{q}\nu=d\beta+c'.
\]
Sum this identity over a coordinate two-torus. The left side gives $\frac{2\pi}{q}$ times the corresponding component of the flux $w$. The coboundary $d\beta$ gives $0$, since every value of $\beta$ occurs in the sum once with each sign. The constant cochain $c'$ gives its value on the plaquettes of the coordinate two-torus times their number. For the coordinate two-tori $\Pi_{\mu\nu}(t)$, $\mu<\nu$, which consist of $n_{\mu}n_{\nu}$ plaquettes, this reads $\frac{2\pi}{q}w_{\mu\nu}=n_{\mu}n_{\nu}\cdot c'_{\mu\nu}$, where $c'_{\mu\nu}$ is the value of $c'$ on these plaquettes; hence $c'=h_w$. Thus $\beta_{\nu}=\beta$ has the required property, and $F=dA+d\beta_{\nu}+h_w=d(A+\beta_{\nu})+h_w$. For $c\in\mathrm{Harm}^2_n$ the cochain $h_w+c$ lies in $\mathrm{Harm}^2_n$, which is orthogonal to the coboundaries by \eqref{eq:partition-functions-1}; Pythagoras' theorem gives the identity for $\|F+c\|^2$. Finally, for each pair $\mu<\nu$ there are $|n|$ plaquettes $p_{\mu\nu}(x)$, one for every $x\in\Lambda_n$, so that
\[
\|h_w\|^2=|n|\cdot\sum_{\mu<\nu}\left(\frac{2\pi w_{\mu\nu}}{qn_{\mu}n_{\nu}}\right)^2,
\]
and $\frac{|n|}{(n_{\mu}n_{\nu})^2}=\frac{n_{\kappa}}{n_{\mu}n_{\nu}}$ for the third direction $\kappa$; if $n_1=n_2=n_3$, each of these three coefficients equals $\frac{1}{n_1}$.

(b) By \eqref{eq:partition-functions-2}, the closed $\nu\in C^2_n(\mathbb{Z})$ with $w(\nu)=0$ are the cochains $dm$ with $m\in C^1_n(\mathbb{Z})$, and $dm=dm'$ if and only if $m-m'$ is closed. Let $R\subset C^1_n(\mathbb{Z})$ contain exactly one element of each class of $C^1_n(\mathbb{Z})$ modulo the closed integer one-cochains, and put
\[
\Psi(A)=\exp\left(-\frac{1}{2}\|dA\|^2\right)\cdot\mathfrak{w}(A)\cdot G(dA)\qquad\left(A\in C^1_n\right).
\]
As $dA+\frac{2\pi}{q}dm=d\left(A+\frac{2\pi}{q}m\right)$, \eqref{eq:partition-functions-5} gives
\begin{align*}
Z_{n,0}[G]&=\sum_{m\in R}\int_{\left[0,\frac{2\pi}{q}\right)^{\mathrm{Lk}_n}}\mathrm{d}A\,\Psi\left(A+\frac{2\pi}{q}m\right)\\
&=\sum_{m\in R}\int_{\left[0,\frac{2\pi}{q}\right)^{\mathrm{Lk}_n}+\frac{2\pi}{q}m}\mathrm{d}A\,\Psi(A).
\end{align*}
The box $\left[0,\frac{2\pi}{q}\right)^{\mathrm{Lk}_n}$ is a fundamental domain for the translations by $\frac{2\pi}{q}C^1_n(\mathbb{Z})$, and $\frac{2\pi}{q}R$ contains exactly one element of each class of $\frac{2\pi}{q}C^1_n(\mathbb{Z})$ modulo $\Gamma_n$. Hence the translates of the box by the vectors $\frac{2\pi}{q}m$, $m\in R$, are pairwise disjoint, and their union is a fundamental domain for the translations by $\Gamma_n$. The function $\Psi$ is invariant under these translations, by \eqref{eq:partition-functions-5} and because $d\gamma=0$ for $\gamma\in\Gamma_n$. This proves the first assertion.

The function $\Psi$ does not depend on the component of $A$ in $V_{\mathrm{g}}$, since $\Psi(A+d\lambda)=\Psi(A)$ for $\lambda\in C^0_n$ by \eqref{eq:partition-functions-4}. By \eqref{eq:partition-functions-1}, $\Gamma_n$ lies in $\ker\left(d|_{C^1_n}\right)=V_{\mathrm{g}}\oplus V_{\mathrm{h}}$. Its intersection with $V_{\mathrm{g}}$ is $\frac{2\pi}{q}dC^0_n(\mathbb{Z})$, a lattice of full rank in $V_{\mathrm{g}}$: if an integer one-cochain equals $d\lambda$ with $\lambda\in C^0_n$, then $\lambda-\lambda(0)$ takes integer values, its value at $x$ being the sum of the values of the cochain along a path from $0$ to $x$. Let $P_h$ be the orthogonal projection of $C^1_n$ onto $V_{\mathrm{h}}$; it maps a one-cochain to the constant one-cochain whose value on the links of direction $\mu$ is the mean of its values on these links. For a closed $m\in C^1_n(\mathbb{Z})$ the sum of the values $m_{\mu}(x)$ along a line in the direction $\mu$ is an integer $k_{\mu}$ that does not depend on the line, for the sums along two adjacent parallel lines agree by $dm=0$; so the mean of the values of $m$ on the links of direction $\mu$ is $\frac{k_{\mu}}{n_{\mu}}$. Every $k\in\mathbb{Z}^3$ occurs, for instance for the closed cochain with $m_{\mu}(x)=k_{\mu}$ if $x_{\mu}=0$ and $m_{\mu}(x)=0$ otherwise. Hence
\[
P_h(\Gamma_n)=\left\{c\in V_{\mathrm{h}}:\ c_{\mu}=\frac{2\pi k_{\mu}}{qn_{\mu}}\ (\mu=1,2,3)\ \text{for some}\ k\in\mathbb{Z}^3\right\},
\]
$c_{\mu}$ denoting the value of $c$ on the links of direction $\mu$; this is a lattice of full rank in $V_{\mathrm{h}}$. So $\Gamma_n$ is a lattice of full rank in $V_{\mathrm{g}}\oplus V_{\mathrm{h}}$, and the quotient $\left(V_{\mathrm{g}}\oplus V_{\mathrm{h}}\right)/\Gamma_n$ is fibred over $V_{\mathrm{h}}/P_h(\Gamma_n)$, with the fibre $V_{\mathrm{g}}/\left(\Gamma_n\cap V_{\mathrm{g}}\right)$. Write $A=A_g+A_h+A_t$ along $C^1_n=V_{\mathrm{g}}\oplus V_{\mathrm{h}}\oplus V_{\mathrm{t}}$; then $dA=dA_t$ and $\Psi(A)=\Psi(A_h+A_t)$. Fubini's theorem gives
\[
Z_{n,0}[G]=\operatorname{vol}\left(V_{\mathrm{g}}/\left(\Gamma_n\cap V_{\mathrm{g}}\right)\right)\cdot\int_{V_{\mathrm{h}}/P_h(\Gamma_n)}\mathrm{d}A_h\int_{V_{\mathrm{t}}}\mathrm{d}\sigma(A_t)\,\Psi(A_h+A_t).
\]
The cochains $H_{\theta}$, $\theta\in[-\pi,\pi)^3$, contain exactly one element of each class of $V_{\mathrm{h}}$ modulo $P_h(\Gamma_n)$, and in the coordinates $\theta$ the Lebesgue measure of $V_{\mathrm{h}}$ is a constant multiple of $\mathrm{d}\theta$. As $dH_{\theta}=0$, it follows that
\[
Z_{n,0}[G]=V_n\cdot Z^{\mathrm{av}}_n[G]
\]
with a constant $V_n>0$ that does not depend on $G$. The functional $Z^{\mathrm{av}}_n$ is a constant multiple of the functional $\overline{Z}$ of \eqref{eq:gaussian-bounds-3}, formed with $c^{\theta}=H_{\theta}$ as in Remark~\ref{rem:tori-unequal-sides}. Since $0<Z_{n,0}[1]<\infty$ (Sections~\ref{sec:summation-all-fluxes} and~\ref{sec:decomposition-flux}), dividing by the same identity for $G=1$ gives $E_{n,0}[G]=\frac{Z^{\mathrm{av}}_n[G]}{Z^{\mathrm{av}}_n[1]}$. The description by a density on $V_{\mathrm{t}}$ follows by Fubini's theorem, since $Z^{\mathrm{av}}_n[|G|]=V_n^{-1}\cdot Z_{n,0}[|G|]<\infty$.

(c) By (a), $\langle f,F\rangle=\left\langle da,d(A+\beta_{\nu})\right\rangle+\langle da,h_w\rangle$, and $\langle da,h_w\rangle=\langle a,d^{*}h_w\rangle=0$, because $d^{*}h_w=0$: by \eqref{eq:partition-functions-1}, $h_w$ is orthogonal to every coboundary.
\end{proof}

The fermion weight may be multiplied by a vacuum factor $\exp\left(-\varepsilon_0|n|\right)$ with a real constant $\varepsilon_0$; for fixed $n$ this factor is a positive number. Properties \eqref{eq:partition-functions-3}--\eqref{eq:partition-functions-8} hold for the product as well, in the case of \eqref{eq:partition-functions-8} because $r$ permutes the coordinate axes, so that $n'$ is a permutation of $n$ and $|n'|=|n|$. Each of the quantities $Z_n[G]$, $Z_n(h)$ and $Z_{n,w}[G]$ is then multiplied by $\exp\left(-\varepsilon_0|n|\right)$, and this number cancels in $E_n$, $E_{n,w}$ and $P_n(w)$. The product formula for $\mathfrak{w}(0)$ in Section~\ref{sec:summation-all-fluxes} gives the value of the determinant without the factor; with the factor, $\mathfrak{w}(0)$ is this value times $\exp\left(-\varepsilon_0|n|\right)$, and is again positive.

\subsection{The partition function as a trace}\label{sec:partition-function-trace}

One of the three coordinate directions is taken as time; it is the direction~$1$, and the other two directions are reduced to this case after Proposition~\ref{prop:trace-formula}. Let $n_0=n_1$, and let
\[
\Lambda_s=\mathbb{Z}/n_2\mathbb{Z}\times\mathbb{Z}/n_3\mathbb{Z}
\]
be the cross-section, the two-dimensional torus in the directions $2$ and $3$. Let $\mathrm{Lk}_s$ and $\mathrm{Pl}_s$ be the sets of its links and of its plaquettes, let $d_s$ be its coboundary, given by the formulas of Section~\ref{sec:tori-cochains} in the directions $2$ and $3$, and let
\[
v=\#|\Lambda_s|,\qquad l=\#|\mathrm{Lk}_s|=2v.
\]
The points of $\Lambda_n$ are written $(t,x)$ with $t\in\mathbb{Z}/n_0\mathbb{Z}$ and $x\in\Lambda_s$, and each slice $\{t\}\times\Lambda_s$ is identified with $\Lambda_s$. A one-cochain on $\Lambda_n$ is written as a family $(A_t,a_t)_{t\in\mathbb{Z}/n_0\mathbb{Z}}$ with $A_t\in\mathbb{R}^{\mathrm{Lk}_s}$ and $a_t\in\mathbb{R}^{\Lambda_s}$: the number $A_t(b)$ is its value on the link $b$ of the slice $t$, and $a_t(x)$ is its value on the link from $(t,x)$ to $(t+1,x)$. For $t\in\mathbb{Z}/n_0\mathbb{Z}$ and a link $b\in\mathrm{Lk}_s$ from $x$ to $x+e_{\mu}$, $\mu\in\{2,3\}$, the plaquette $p_{1\mu}((t,x))$ is called the temporal plaquette over $(t,b)$; for a plaquette $p=p_{23}(x)\in\mathrm{Pl}_s$, the plaquette $p_{23}((t,x))$ is called the spatial plaquette $(t,p)$. Accordingly, a two-cochain $h$ on $\Lambda_n$ is written as a family $(h^{\mathrm{el}}_t,h^{\mathrm{mag}}_t)_{t\in\mathbb{Z}/n_0\mathbb{Z}}$: its electric part $h^{\mathrm{el}}_t\in\mathbb{R}^{\mathrm{Lk}_s}$ takes at $b$ the value of $h$ on the temporal plaquette over $(t,b)$, and its magnetic part $h^{\mathrm{mag}}_t\in\mathbb{R}^{\mathrm{Pl}_s}$ takes at $p$ the value of $h$ on the spatial plaquette $(t,p)$. With the orientations of Section~\ref{sec:tori-cochains}, the curvature $dA$ of a one-cochain $A=(A_t,a_t)_t$ has the parts
\[
(dA)^{\mathrm{el}}_t=A_{t+1}-A_t-d_sa_t,\qquad(dA)^{\mathrm{mag}}_t=d_sA_t;
\]
that is, its value on the temporal plaquette over $(t,b)$ is $A_{t+1}(b)-A_t(b)-(d_sa_t)(b)$, and its value on the spatial plaquette $(t,p)$ is $(d_sA_t)(p)$.

\subsubsection{Factorisation of the determinant over time slices}\label{sec:factorisation-determinant-time}

Fix $q>0$ and $\boldsymbol{\mu}>-1$, and apply Lemma~\ref{lem:slice-factorisation-determinant-2} to the cross-section $\Lambda_s$. It furnishes a finite-dimensional Hilbert space $\mathcal{F}$, a function $s\colon\mathbb{R}^{\mathrm{Lk}_s}\to(0,\infty)$, a function $Q$ on $\mathbb{R}^{\mathrm{Lk}_s}$ with values in the space of linear operators on $\mathcal{F}$, and a family $(Q_x)_{x\in\Lambda_s}$ of self-adjoint operators on $\mathcal{F}$, called the local charges, which commute with one another and have their spectra in $\mathbb{Z}$. These data depend on nothing but $\Lambda_s$, $q$ and $\boldsymbol{\mu}$; in particular, they are the same for all $n_0$. For $a\in\mathbb{R}^{\Lambda_s}$ put
\[
W(a)=\exp\left(iq\sum_{x\in\Lambda_s}a_x\cdot Q_x\right).
\]
If $n_0\ge2$ is an integer, $n=(n_0,n_2,n_3)$, and $(A_t,a_t)_{t\in\mathbb{Z}/n_0\mathbb{Z}}$ is a real one-cochain on $\Lambda_n$, then
\begin{equation}\tag{FK}\label{eq:partition-functions-FK}
\det(D_A+\boldsymbol{\mu})=\operatorname{tr}_{\mathcal{F}}\prod_{t=n_0-1}^0\left[s(A_{t+1})\cdot Q(A_{t+1})^{*}\cdot W(a_t)\cdot Q(A_t)\cdot s(A_t)\right];
\end{equation}
here $A_{n_0}=A_0$, the index $t$ being read modulo $n_0$; the product is ordered with later times to the left, the factor $t=n_0-1$ standing leftmost and the factor $t=0$ rightmost; and $\operatorname{tr}_{\mathcal{F}}$ is the ordinary trace on $\mathcal{F}$, which yields the antiperiodic closure in time. Moreover,
\begin{gather}
s\ \text{and}\ Q\ \text{are bounded and continuous, and}\ c_0:=\sup_{A\in\mathbb{R}^{\mathrm{Lk}_s}}s(A)\cdot\|Q(A)\|<\infty;\tag{F1}\label{eq:partition-functions-F1}\\
s\ \text{and}\ Q\ \text{are functions of}\ \left(\exp(iqA(b))\right)_{b\in\mathrm{Lk}_s},\ \text{hence}\ \frac{2\pi}{q}\mathbb{Z}^{\mathrm{Lk}_s}\text{-periodic};\tag{F2}\label{eq:partition-functions-F2}\\
W(a)\ \text{is unitary},\ W(a)^{*}=W(-a),\ W(a+a')=W(a)\cdot W(a')\ \text{for}\ a,a'\in\mathbb{R}^{\Lambda_s},\notag\\
W(a)=1\ \text{for}\ a\in\frac{2\pi}{q}\mathbb{Z}^{\Lambda_s};\tag{F3}\label{eq:partition-functions-F3}
\end{gather}
here $\|Q(A)\|$ is the operator norm on $\mathcal{F}$.

In the notation of Lemma~\ref{lem:slice-factorisation-determinant-2} one has $Q=\mathfrak{Q}$, $Q_x=Q^{\mathrm{ch}}_x$ and $W=\mathsf{G}$; with these identifications, \eqref{eq:partition-functions-FK} is the clause~\eqref{eq:free-coefficient-b} of that lemma, and \eqref{eq:partition-functions-F1}--\eqref{eq:partition-functions-F3} together form its clause~\eqref{eq:free-coefficient-a}. The proof of Lemma~\ref{lem:slice-factorisation-determinant-2} rests on Lemmas~\ref{lem:slice-factorisation-determinant} and~\ref{lem:properties-second-quantisation}, on Proposition~\ref{prop:trace-formula-determinant} and, for rectangular tori, on Remark~\ref{rem:tori-unequal-sides}; that proof and these four statements are independent of Section~\ref{sec:partition-functions-all}, so that no circularity arises. Factorising the Wilson determinant over time slices as in \eqref{eq:partition-functions-FK} is an idea of L\"uscher~\cite{Luescher77}. Section~\ref{sec:partition-function-trace} draws on no property of the fermion weight $\mathfrak{w}$ other than \eqref{eq:partition-functions-FK}, \eqref{eq:partition-functions-F1}--\eqref{eq:partition-functions-F3} and \eqref{eq:partition-functions-5}, except in the paragraph following Proposition~\ref{prop:trace-formula}: there the time direction $1$ is exchanged for $2$ or $3$ by means of the invariance under lattice isometries established in Section~\ref{sec:summation-all-fluxes}, so that \eqref{eq:partition-functions-8} enters as well.

\subsubsection{Paths and spatial fluxes}\label{sec:paths-spatial-fluxes}

Let
\[
\Gamma_s=\frac{2\pi}{q}\left\{m\in\mathbb{Z}^{\mathrm{Lk}_s}:d_sm=0\right\};
\]
this is a lattice of full rank in the kernel $\ker d_s\subset\mathbb{R}^{\mathrm{Lk}_s}$ of $d_s$ on the one-cochains. Let $\mathfrak{C}=\mathbb{R}^{\mathrm{Lk}_s}/\Gamma_s$, with the measure induced by the Lebesgue measure of $\mathbb{R}^{\mathrm{Lk}_s}$. The class of $A\in\mathbb{R}^{\mathrm{Lk}_s}$ in $\mathfrak{C}$ is denoted by $[A]$, and functions on $\mathfrak{C}$ are identified with $\Gamma_s$-periodic functions on $\mathbb{R}^{\mathrm{Lk}_s}$. Let
\[
\mathrm{box}_s=\left[0,\frac{2\pi}{q}\right)^{\mathrm{Lk}_s},\qquad\mathrm{box}_0=\left[0,\frac{2\pi}{q}\right)^{\Lambda_s}.
\]
Fix a plaquette $p_0\in\mathrm{Pl}_s$, and for $w\in\mathbb{Z}$ let $\nu^{(w)}=w\cdot\mathbf{1}_{p_0}\in\mathbb{Z}^{\mathrm{Pl}_s}$, where $\mathbf{1}_{p_0}$ is the indicator function of $p_0$. In Lemma~\ref{lem:integral-paths} the summation index $w$ corresponds to the flux $w_{23}(\nu)$, in the notation of \eqref{eq:partition-functions-2}, of the closed integer two-cochains $\nu$ of Definition~\ref{def:partition-functions-all}.

Let $w\in\mathbb{Z}$. A path of spatial flux $w$ is a family
\[
\mathsf{p}=(A_0,\dots,A_{n_0};a_0,\dots,a_{n_0-1}),\qquad A_t\in\mathbb{R}^{\mathrm{Lk}_s},\quad A_{n_0}-A_0\in\Gamma_s,\quad a_t\in\mathrm{box}_0.
\]
Its curvature is the two-cochain $F(\mathsf{p})\in C^2_n$ whose electric and magnetic parts are, for $0\le t<n_0$,
\[
F(\mathsf{p})^{\mathrm{el}}_t=\mathrm{El}_t=A_{t+1}-A_t-d_sa_t,\qquad F(\mathsf{p})^{\mathrm{mag}}_t=B_t=d_sA_t+\frac{2\pi}{q}\nu^{(w)},
\]
and $\mathfrak{w}(\mathsf{p})$ is the fermion weight with the link phases $\exp(iqA_t)$ on the links of the slice $t$ and $\exp(iqa_t)$ on the links from the slice $t$ to the slice $t+1$, $0\le t<n_0$; by \eqref{eq:partition-functions-5}, $\mathfrak{w}$ is a function of the link phases. The integral over the paths of spatial flux $w$ is
\[
\int\mathrm{d}\mathsf{p}=\int_{\mathfrak{C}}\mathrm{d}[A_0]\,\sum_{\gamma\in\Gamma_s}\int\prod_{t=1}^{n_0-1}\mathrm{d}A_t\,\int\prod_{t=0}^{n_0-1}\mathrm{d}a_t,
\]
where $A_{n_0}=A_0+\gamma$, each $A_t$ with $1\le t<n_0$ ranges over $\mathbb{R}^{\mathrm{Lk}_s}$, each $a_t$ ranges over $\mathrm{box}_0$, and the integrals are taken with respect to Lebesgue measure.

\begin{lemma}[Integral over paths]\label{lem:integral-paths}
For every $G\in\mathrm{Gr}$,
\[
Z_n[G]=\sum_{w\in\mathbb{Z}}\int\mathrm{d}\mathsf{p}\,\exp\left(-\frac{1}{2}\|F(\mathsf{p})\|^2\right)\cdot\mathfrak{w}(\mathsf{p})\cdot G(F(\mathsf{p})),
\]
where the integral of index $w$ is over the paths of spatial flux $w$. Replacing every $A_t$, $0\le t\le n_0$, by $A_t+\gamma'$, with one and the same $\gamma'\in\Gamma_s$, leaves the integrand unchanged. Hence, once the integrand has been integrated over $A_1,\dots,A_{n_0-1}$ and $a_0,\dots,a_{n_0-1}$ and summed over $\gamma\in\Gamma_s$, the result depends on $A_0$ only through its class $[A_0]$.
\end{lemma}

\begin{proof}
An integer two-cochain $\nu$ on $\Lambda_n$ is written as $(k_t,\nu_t)_{t\in\mathbb{Z}/n_0\mathbb{Z}}$, with the temporal part $k_t=\nu^{\mathrm{el}}_t\in\mathbb{Z}^{\mathrm{Lk}_s}$ and the spatial part $\nu_t=\nu^{\mathrm{mag}}_t\in\mathbb{Z}^{\mathrm{Pl}_s}$. Let $z=(t,x)\in\Lambda_n$ and $p=p_{23}(x)\in\mathrm{Pl}_s$. On the cube with the corner $z$,
\begin{align*}
(d\nu)_{123}(z)&=\left[\nu_{23}(z+e_1)-\nu_{23}(z)\right]-\left[\nu_{13}(z+e_2)-\nu_{13}(z)\right]+\left[\nu_{12}(z+e_3)-\nu_{12}(z)\right]\\
&=\nu_{t+1}(p)-\nu_t(p)-(d_sk_t)(p).
\end{align*}
Hence $d\nu=0$ if and only if $\nu_{t+1}=\nu_t+d_sk_t$ for all $t\in\mathbb{Z}/n_0\mathbb{Z}$. Let $m_0=0$ and $m_{t+1}=m_t+k_t$ for $0\le t<n_0$. The closed $\nu$ are therefore exactly the cochains with $\nu_t=\nu_0+d_sm_t$ for $0\le t<n_0$, where $\nu_0\in\mathbb{Z}^{\mathrm{Pl}_s}$ and $k_0,\dots,k_{n_0-1}\in\mathbb{Z}^{\mathrm{Lk}_s}$ are subject only to the condition $d_sm_{n_0}=0$, which expresses the periodicity of $t\mapsto\nu_t$.

By Definition~\ref{def:partition-functions-all}, $Z_n[G]$ is the sum over these $\nu$ of the integral over the one-cochains $A=(A_t,a_t)_t$ with $A_t\in\mathrm{box}_s$ and $a_t\in\mathrm{box}_0$, $0\le t<n_0$. Put
\[
\hat{A}_t=A_t+\frac{2\pi}{q}m_t,\qquad0\le t\le n_0,\qquad\text{where }A_{n_0}=A_0.
\]
For fixed $\nu_0$, the map
\[
(A_t,k_t)_{0\le t<n_0}\longmapsto\left(\hat{A}_0;\ \hat{A}_1,\dots,\hat{A}_{n_0-1};\ \frac{2\pi}{q}m_{n_0}\right),
\]
defined for $A_t\in\mathrm{box}_s$ and for families $(k_t)_t$ with $d_sm_{n_0}=0$, is a bijection onto $\mathrm{box}_s\times\left(\mathbb{R}^{\mathrm{Lk}_s}\right)^{n_0-1}\times\Gamma_s$, and it carries the product of Lebesgue measure and counting measure to the product of Lebesgue measure and counting measure. In the new variables the electric part of $F=dA+\frac{2\pi}{q}\nu$ is
\[
A_{t+1}-A_t-d_sa_t+\frac{2\pi}{q}k_t=\hat{A}_{t+1}-\hat{A}_t-d_sa_t,
\]
its magnetic part is
\[
d_sA_t+\frac{2\pi}{q}\nu_t=d_s\hat{A}_t+\frac{2\pi}{q}\nu_0,
\]
and the link phases are unchanged, because $q\hat{A}_t-qA_t\in2\pi\mathbb{Z}^{\mathrm{Lk}_s}$; so the fermion weight is unchanged by \eqref{eq:partition-functions-5}. Finally, let $w=\sum_{p\in\mathrm{Pl}_s}\nu_0(p)$; this is the flux $w_{23}(\nu)$. By the two-dimensional part of \eqref{eq:partition-functions-2}, $\nu_0=\nu^{(w)}+d_sm'$ with $m'\in\mathbb{Z}^{\mathrm{Lk}_s}$, and $m'$ is unique modulo the closed integer one-cochains. The translation of all $\hat{A}_t$, $0\le t\le n_0$, by $\frac{2\pi}{q}m'$ replaces $\nu_0$ by $\nu^{(w)}$ in the magnetic part and changes neither the electric part nor the link phases. As $m'$ ranges over a system of representatives of $\mathbb{Z}^{\mathrm{Lk}_s}$ modulo the closed integer one-cochains, the sets $\mathrm{box}_s+\frac{2\pi}{q}m'$ tile $\mathfrak{C}$, every point of $\mathfrak{C}$ being covered exactly once. Since the sum and the integrals in Definition~\ref{def:partition-functions-all} converge absolutely (Section~\ref{sec:summation-all-fluxes}), the formula for $Z_n[G]$ follows.

Let $\gamma'\in\Gamma_s$, and translate all $A_t$, $0\le t\le n_0$, by $\gamma'$. The parts $\mathrm{El}_t$ do not change, since they involve the $A_t$ only through differences; the parts $B_t$ do not change, since $d_s\gamma'=0$; and $\mathfrak{w}(\mathsf{p})$ does not change, by \eqref{eq:partition-functions-5}, or equally by \eqref{eq:partition-functions-FK} and \eqref{eq:partition-functions-F2}, since $\gamma'\in\frac{2\pi}{q}\mathbb{Z}^{\mathrm{Lk}_s}$. As Lebesgue measure is invariant under translations and $\gamma=A_{n_0}-A_0$ is unchanged, the integral over $A_1,\dots,A_{n_0-1}$ and $a_0,\dots,a_{n_0-1}$, summed over $\gamma$, takes the same value at $A_0$ and at $A_0+\gamma'$.
\end{proof}

\subsubsection{Transfer operators}\label{sec:transfer-operators}

Let
\[
\mathscr{H}_1=L^2(\mathfrak{C};\mathcal{F}),\qquad\mathscr{H}=\bigoplus_{w\in\mathbb{Z}}\mathscr{H}_1.
\]
The operators $S_x$, $\exp\left(-\frac{\tau}{2}E^2\right)$, $U_{\mathrm{g}}(a)$, $P_G$, $K$ and $K^{\frac{1}{2}}$ defined below act on $\mathscr{H}_1$, and on $\mathscr{H}$ diagonally, that is, in the same way on each summand. An operator $Y$ on $\mathscr{H}_1$ has the kernel $Y([A'],[A])$, a function on $\mathfrak{C}\times\mathfrak{C}$ whose values are linear operators on $\mathcal{F}$, if $(Y\Psi)([A'])=\int_{\mathfrak{C}}\mathrm{d}[A]\,Y([A'],[A])\cdot\Psi([A])$ for all $\Psi\in\mathscr{H}_1$.
\begin{itemize}
\item For $x\in\mathbb{R}^{\mathrm{Lk}_s}$ the translation $S_x$, given by $(S_x\Psi)(A)=\Psi(A-x)$, is unitary.
\item For $\tau>0$, the operator $\exp\left(-\frac{\tau}{2}E^2\right)$ is the convolution with the function
\[
g_{\tau}([A])=\sum_{\gamma\in\Gamma_s}(2\pi\tau)^{-\frac{l}{2}}\cdot\exp\left(-\frac{|A+\gamma|^2}{2\tau}\right)
\]
on $\mathfrak{C}$, where $|\cdot|$ is the Euclidean norm of $\mathbb{R}^{\mathrm{Lk}_s}$; that is, $\exp\left(-\frac{\tau}{2}E^2\right)$ has the kernel $g_{\tau}([A'-A])$. These operators form the heat semigroup of the flat manifold $\mathfrak{C}$: they are non-negative, and
\[
\exp\left(-\frac{\tau}{2}E^2\right)\cdot\exp\left(-\frac{\tau'}{2}E^2\right)=\exp\left(-\frac{\tau+\tau'}{2}E^2\right),\qquad\tau,\tau'>0.
\]
\item For $a\in\mathbb{R}^{\Lambda_s}$ let $(U_{\mathrm{g}}(a)\Psi)(A)=W(a)\cdot\Psi(A-d_sa)$. By \eqref{eq:partition-functions-F3}, and since $d_s\left(\frac{2\pi}{q}\mathbb{Z}^{\Lambda_s}\right)\subset\Gamma_s$, the map $a\mapsto U_{\mathrm{g}}(a)$ is a strongly continuous unitary representation of the compact group $\left(\mathbb{R}/\frac{2\pi}{q}\mathbb{Z}\right)^{\Lambda_s}$. Hence
\[
\int_{\mathrm{box}_0}\mathrm{d}a\,U_{\mathrm{g}}(a)=\left(\frac{2\pi}{q}\right)^v\cdot P_G,
\]
where $P_G$ is the orthogonal projection onto the subspace of the vectors that are invariant under all $U_{\mathrm{g}}(a)$.
\item Let
\[
K=\left(\frac{2\pi}{q}\right)^v\cdot\exp\left(-\frac{1}{2}E^2\right)\cdot P_G,\qquad K^{\frac{1}{2}}=\left(\frac{2\pi}{q}\right)^{\frac{v}{2}}\cdot\exp\left(-\frac{1}{4}E^2\right)\cdot P_G.
\]
The operators $S_x$, $\exp\left(-\frac{\tau}{2}E^2\right)$, $U_{\mathrm{g}}(a)$ and $P_G$ commute with one another, because they are built from the translations of the abelian group $\mathfrak{C}$ and from the operators $W(a)$, which do not depend on $A$. Hence $K\ge0$, $K^{\frac{1}{2}}\ge0$, $\left(K^{\frac{1}{2}}\right)^2=K$, and $S_x$ commutes with $K^{\frac{1}{2}}$. The kernel of $K\cdot S_x$ is
\[
(K\cdot S_x)([A'],[A])=\sum_{\gamma\in\Gamma_s}\int_{\mathrm{box}_0}\mathrm{d}a\,(2\pi)^{-\frac{l}{2}}\cdot\exp\left(-\frac{1}{2}\left|A'+\gamma-A-x-d_sa\right|^2\right)\cdot W(a).
\]
\item For $w\in\mathbb{Z}$ and $\beta\in\mathbb{R}^{\mathrm{Pl}_s}$ let
\[
M_{w,\beta}(A)=s(A)\cdot\exp\left(-\frac{1}{4}\left\|d_sA+\frac{2\pi}{q}\nu^{(w)}+\beta\right\|^2\right)\cdot Q(A),\qquad A\in\mathbb{R}^{\mathrm{Lk}_s},
\]
where $\|\cdot\|$ is the Euclidean norm of $\mathbb{R}^{\mathrm{Pl}_s}$. By \eqref{eq:partition-functions-F1}, \eqref{eq:partition-functions-F2} and $d_s\gamma=0$ for $\gamma\in\Gamma_s$, the function $M_{w,\beta}$ is $\Gamma_s$-periodic, bounded and continuous, with values in the linear operators on $\mathcal{F}$. Let $M_{\beta}=\bigoplus_{w\in\mathbb{Z}}M_{w,\beta}$ be the operator on $\mathscr{H}$ that multiplies the component of index $w$ by $M_{w,\beta}$, and let
\[
X_{\beta}=K^{\frac{1}{2}}\cdot M_{\beta},\qquad T=X_0^{*}\cdot X_0=M_0^{*}\cdot K\cdot M_0\ge0,
\]
where $X_0$ and $M_0$ are $X_{\beta}$ and $M_{\beta}$ for $\beta=0$.
\end{itemize}

\begin{lemma}[Hilbert--Schmidt property]\label{lem:hilbert-schmidt-property}
For every $\beta\in\mathbb{R}^{\mathrm{Pl}_s}$ the operator $X_{\beta}$ on $\mathscr{H}$ is a Hilbert--Schmidt operator. It follows that each of the operators $T=X_0^{*}\cdot X_0$ and $X_{\beta}\cdot X_{\beta}^{*}$ is non-negative and belongs to the trace class.
\end{lemma}

\begin{proof}
The computation of the kernel of $K\cdot S_x$, with $\tau=\frac{1}{2}$ in place of $\tau=1$ and with $x=0$, shows that $K^{\frac{1}{2}}$ has the kernel
\[
K^{\frac{1}{2}}([A'],[A])=\left(\frac{2\pi}{q}\right)^{-\frac{v}{2}}\int_{\mathrm{box}_0}\mathrm{d}a\,g_{1/2}([A'-A-d_sa])\cdot W(a).
\]
As $W(a)$ is unitary by \eqref{eq:partition-functions-F3} and $g_{1/2}\ge0$, the operator norm of this kernel is at most $k([A'-A])$, where
\[
k([A])=\left(\frac{2\pi}{q}\right)^{-\frac{v}{2}}\int_{\mathrm{box}_0}\mathrm{d}a\,g_{1/2}([A-d_sa]).
\]
Like $g_{1/2}$, the function $k$ is bounded and integrable on $\mathfrak{C}$; so $k\in L^2(\mathfrak{C})$. On the summand of index $w$ the operator $X_{\beta}$ has the kernel $K^{\frac{1}{2}}([A'],[A])\cdot M_{w,\beta}(A)$, whose Hilbert--Schmidt norm on $\mathcal{F}$ is at most $(\dim\mathcal{F})^{\frac{1}{2}}\cdot k([A'-A])\cdot\|M_{w,\beta}(A)\|$. The square of the Hilbert--Schmidt norm of an integral operator is the integral of the square of the Hilbert--Schmidt norm of its kernel; therefore
\[
\|X_{\beta}\|_{\mathrm{HS}}^2\le\dim\mathcal{F}\cdot\|k\|_2^2\cdot\sum_{w\in\mathbb{Z}}\int_{\mathfrak{C}}\mathrm{d}A\,\|M_{w,\beta}(A)\|^2,
\]
and by \eqref{eq:partition-functions-F1} the sum on the right is at most
\[
c_0^2\cdot\sum_{w\in\mathbb{Z}}\int_{\mathfrak{C}}\mathrm{d}A\,\exp\left(-\frac{1}{2}\left\|d_sA+\frac{2\pi}{q}\nu^{(w)}+\beta\right\|^2\right).
\]
Since $\Gamma_s\subset\ker d_s$, the space $\mathfrak{C}$ is the product of the compact torus $\ker d_s/\Gamma_s$ and of $(\ker d_s)^{\perp}$, on which $d_s$ is injective; the integrand depends on $A$ only through its component in $(\ker d_s)^{\perp}$. Let $\mathbf{1}\in\mathbb{R}^{\mathrm{Pl}_s}$ be the two-cochain equal to $1$ on every plaquette; as $\Lambda_s$ has $v$ plaquettes, $v^{-\frac{1}{2}}\mathbf{1}$ is a unit vector. The sum of $d_sA$ over all plaquettes of $\Lambda_s$ vanishes, so $d_sA$ is orthogonal to $\mathbf{1}$, and the component of $d_sA+\frac{2\pi}{q}\nu^{(w)}+\beta$ along $v^{-\frac{1}{2}}\mathbf{1}$ does not depend on $A$: it is
\[
\rho_w=v^{-\frac{1}{2}}\left(\frac{2\pi w}{q}+\sum_{p\in\mathrm{Pl}_s}\beta_p\right).
\]
Hence
\[
\left\|d_sA+\frac{2\pi}{q}\nu^{(w)}+\beta\right\|^2=\rho_w^2+\|d_sA+y_w\|^2,\qquad y_w=\frac{2\pi}{q}\nu^{(w)}+\beta-\rho_w\,v^{-\frac{1}{2}}\mathbf{1}.
\]
For every $y\in\mathbb{R}^{\mathrm{Pl}_s}$ the integral of $\exp\left(-\frac{1}{2}\|d_sA+y\|^2\right)$ over $A\in(\ker d_s)^{\perp}$ is at most its value for $y=0$ (write $y=d_s\tilde{A}+y'$ with $\tilde{A}\in(\ker d_s)^{\perp}$ and $y'$ orthogonal to the range of $d_s$, and translate $A$ by $\tilde{A}$), and this value is finite because $d_s$ is injective on $(\ker d_s)^{\perp}$. Therefore the integral of index $w$ is at most
\[
C\cdot\exp\left(-\frac{1}{2v}\left(\frac{2\pi w}{q}+\sum_{p\in\mathrm{Pl}_s}\beta_p\right)^2\right),
\]
where $C$, the volume of $\ker d_s/\Gamma_s$ times the value for $y=0$, does not depend on $w$; and this bound is summable over $w\in\mathbb{Z}$. Thus $X_{\beta}$ is a Hilbert--Schmidt operator. The operators $T=X_0^{*}\cdot X_0$ and $X_{\beta}\cdot X_{\beta}^{*}$ are products of two Hilbert--Schmidt operators, hence of trace class, and they are non-negative.
\end{proof}

\begin{proposition}[Trace formula]\label{prop:trace-formula}
Let $h\in C^2_n$, and for $t\in\mathbb{Z}/n_0\mathbb{Z}$ let $X_{(t)}:=X_{h^{\mathrm{mag}}_t}$, so that $X_{(n_0)}=X_{(0)}$. Then
\[
Z_n(h)=(2\pi)^{\frac{l\cdot n_0}{2}}\cdot\operatorname{Tr}_{\mathscr{H}}\prod_{t=n_0-1,\dots,0}\left(X_{(t+1)}^{*}\cdot S_{-h^{\mathrm{el}}_t}\cdot X_{(t)}\right),
\]
where the factors are ordered by decreasing $t$ from left to right. In particular,
\[
Z_n[1]=(2\pi)^{\frac{l\cdot n_0}{2}}\cdot\operatorname{Tr}_{\mathscr{H}}T^{n_0}.
\]
\end{proposition}

\begin{proof}
For $w\in\mathbb{Z}$, $A,A'\in\mathbb{R}^{\mathrm{Lk}_s}$ and $a\in\mathbb{R}^{\Lambda_s}$ let $B_w(A)\in\mathbb{R}^{\mathrm{Pl}_s}$ and the operator $\mathrm{Tk}_w(A',A;a)$ on $\mathcal{F}$ be given by
\begin{align*}
B_w(A)&:=d_sA+\frac{2\pi}{q}\,\nu^{(w)},\\
\mathrm{Tk}_w(A',A;a)&:=(2\pi)^{-\frac{l}{2}}\cdot\exp\left(-\frac{1}{2}|A'-A-d_sa|^2\right)\cdot M_{w,0}(A')^{*}\cdot W(a)\cdot M_{w,0}(A).
\end{align*}
Since $M_{w,0}(A)=s(A)\cdot\exp\left(-\frac{1}{4}\|B_w(A)\|^2\right)\cdot Q(A)$ and $W(a)$ is unitary, \eqref{eq:partition-functions-F1} and \eqref{eq:partition-functions-F3} give
\begin{multline*}
\|\mathrm{Tk}_w(A',A;a)\|\le c_0^2\cdot(2\pi)^{-\frac{l}{2}}\cdot\exp\left(-\frac{1}{2}|A'-A-d_sa|^2\right.\\
\left.-\frac{1}{4}\|B_w(A')\|^2-\frac{1}{4}\|B_w(A)\|^2\right).
\end{multline*}

Let $G_h\colon C^2_n\to(0,\infty)$ be the function $G_h(F):=\exp\left(-\langle h,F\rangle-\frac{1}{2}\|h\|^2\right)$. Then $G_h\in\mathrm{Gr}$, because $G_h(F)\le\exp(\|h\|\cdot\|F\|)$, and $Z_n(h)=Z_n[G_h]$ by Definition~\ref{def:partition-functions-all}. As $\|F\|^2+2\langle h,F\rangle+\|h\|^2=\|F+h\|^2$, Lemma~\ref{lem:integral-paths} yields
\[
Z_n(h)=\sum_{w\in\mathbb{Z}}\int\mathrm{d}\mathsf{p}\,\exp\left(-\frac{1}{2}\|F(\mathsf{p})+h\|^2\right)\cdot\mathfrak{w}(\mathsf{p}).
\]
Let $\mathsf{p}=(A_0,\dots,A_{n_0};a_0,\dots,a_{n_0-1})$ be a path of spatial flux $w$, with $A_{n_0}=A_0+\gamma$, $\gamma\in\Gamma_s$. The parts of its curvature are $\mathrm{El}_t=A_{t+1}-A_t-d_sa_t$ and $B_t=B_w(A_t)$, $0\le t<n_0$, so that
\[
\|F(\mathsf{p})+h\|^2=\sum_{t=0}^{n_0-1}\left(|\mathrm{El}_t+h^{\mathrm{el}}_t|^2+\|B_t+h^{\mathrm{mag}}_t\|^2\right).
\]
The index $t$ of $h^{\mathrm{el}}_t$, $h^{\mathrm{mag}}_t$ and $B_t$ is read modulo $n_0$; for $B_t$ this is consistent, because $d_s\gamma=0$ gives $B_w(A_{n_0})=B_w(A_0)$. The weight $\mathfrak{w}(\mathsf{p})$ is expressed by \eqref{eq:partition-functions-FK}, and in the factor with the index $t=n_0-1$ the variable $A_0$ may be replaced by $A_{n_0}$: by \eqref{eq:partition-functions-F2} the functions $s$ and $Q$ are $\frac{2\pi}{q}\mathbb{Z}^{\mathrm{Lk}_s}$-periodic, and $\Gamma_s\subset\frac{2\pi}{q}\mathbb{Z}^{\mathrm{Lk}_s}$. The variable $A_t$ enters exactly two factors of \eqref{eq:partition-functions-FK}, those with the indices $t$ and $t-1$, and each of them contains $s(A_t)$ once. Write the Gaussian factor $\exp\left(-\frac{1}{2}\|B_t+h^{\mathrm{mag}}_t\|^2\right)$ of the slice $t$ as the square of $\exp\left(-\frac{1}{4}\|B_t+h^{\mathrm{mag}}_t\|^2\right)$, attach one of the two equal factors to each of these two factors of \eqref{eq:partition-functions-FK}, and attach $\exp\left(-\frac{1}{2}|\mathrm{El}_t+h^{\mathrm{el}}_t|^2\right)$ to the factor with the index $t$. By the definition of $M_{w,\beta}$ in Section~\ref{sec:transfer-operators} this gives
\[
\exp\left(-\frac{1}{2}\|F(\mathsf{p})+h\|^2\right)\cdot\mathfrak{w}(\mathsf{p})=(2\pi)^{\frac{l\cdot n_0}{2}}\cdot\operatorname{tr}_{\mathcal{F}}\prod_{t=n_0-1,\dots,0}\Phi_t,
\]
with the factors ordered as in the statement, where
\[
\Phi_t:=M_{w,h^{\mathrm{mag}}_{t+1}}(A_{t+1})^{*}\cdot(2\pi)^{-\frac{l}{2}}\cdot\exp\left(-\frac{1}{2}|\mathrm{El}_t+h^{\mathrm{el}}_t|^2\right)\cdot W(a_t)\cdot M_{w,h^{\mathrm{mag}}_t}(A_t).
\]

One has $M_{w,\beta}(A)=\exp\left(\frac{1}{4}\|B_w(A)\|^2-\frac{1}{4}\|B_w(A)+\beta\|^2\right)\cdot M_{w,0}(A)$ for $\beta\in\mathbb{R}^{\mathrm{Pl}_s}$. Hence $\Phi_t$ is $\mathrm{Tk}_w(A_{t+1},A_t;a_t)$ multiplied by a positive number, and the product of these numbers over $t$ equals $\exp\left(\frac{1}{2}\|F(\mathsf{p})\|^2-\frac{1}{2}\|F(\mathsf{p})+h\|^2\right)=G_h(F(\mathsf{p}))$. Together with the bound on $\mathrm{Tk}_w$ and with $B_w(A_{n_0})=B_w(A_0)$ this gives
\[
(2\pi)^{\frac{l\cdot n_0}{2}}\cdot\dim\mathcal{F}\cdot\prod_{t=0}^{n_0-1}\|\Phi_t\|\le c_0^{2n_0}\cdot\dim\mathcal{F}\cdot\exp\left(-\frac{1}{2}\|F(\mathsf{p})\|^2\right)\cdot G_h(F(\mathsf{p})).
\]
In what follows, sums and integrals are applied to the product of the $\Phi_t$, to kernels obtained from it by summation and integration over some of the variables, and to their traces over $\mathcal{F}$. Since $|\operatorname{tr}_{\mathcal{F}}Y|\le\dim\mathcal{F}\cdot\|Y\|$ for every operator $Y$ on $\mathcal{F}$, the last inequality shows that each of these sums and integrals is dominated by the right side of the identity of Lemma~\ref{lem:integral-paths} with $\mathfrak{w}(\mathsf{p})$ replaced by the constant $c_0^{2n_0}\cdot\dim\mathcal{F}$ and $G$ replaced by $|G_h|=G_h$. This quantity is finite. Indeed, the identity of Lemma~\ref{lem:integral-paths} holds, with the same proof, when $\mathfrak{w}$ is replaced by a positive constant, which has the periodicity \eqref{eq:partition-functions-5} of $\mathfrak{w}$; so the dominating quantity equals the series of Definition~\ref{def:partition-functions-all} with $\mathfrak{w}$ replaced by $c_0^{2n_0}\cdot\dim\mathcal{F}$ and $G$ by $G_h$, and this series converges by the argument given in Section~\ref{sec:summation-all-fluxes} for the series of Definition~\ref{def:partition-functions-all}. Consequently all sums and integrals below converge absolutely, and by Fubini's theorem they may be taken in any order.

Let $D\subset\mathbb{R}^{\mathrm{Lk}_s}$ be a Borel set that contains exactly one point of every class of $\mathfrak{C}=\mathbb{R}^{\mathrm{Lk}_s}/\Gamma_s$. With Lebesgue measure, $D$ is identified with $\mathfrak{C}$, and
\[
\int_{\mathbb{R}^{\mathrm{Lk}_s}}\mathrm{d}A\,f(A)=\int_D\mathrm{d}A\sum_{\gamma'\in\Gamma_s}f(A+\gamma')
\]
for every integrable function $f\colon\mathbb{R}^{\mathrm{Lk}_s}\to\mathbb{C}$. Take $A_0\in D$ in the integral over $[A_0]\in\mathfrak{C}$, write $A_t=\tilde{A}_t+\gamma_t$ with $\tilde{A}_t\in D$ and $\gamma_t\in\Gamma_s$ for $1\le t<n_0$, and put $\tilde{A}_0=\tilde{A}_{n_0}:=A_0$, $\gamma_0:=0$ and $\gamma_{n_0}:=\gamma$; then $A_t=\tilde{A}_t+\gamma_t$ for $0\le t\le n_0$. The functions $M_{w,\beta}$ are $\Gamma_s$-periodic (Section~\ref{sec:transfer-operators}). Hence $\Phi_t$ depends on $\gamma_0,\dots,\gamma_{n_0}$ only through $\delta_t:=\gamma_{t+1}-\gamma_t$, which enters through $\mathrm{El}_t=\tilde{A}_{t+1}+\delta_t-\tilde{A}_t-d_sa_t$; moreover, $a_t$ occurs in $\Phi_t$ only. The map $(\gamma_1,\dots,\gamma_{n_0-1},\gamma)\mapsto(\delta_0,\dots,\delta_{n_0-1})$ is a bijection of $\Gamma_s^{n_0}$ onto itself, with the inverse $\gamma_t=\delta_0+\dots+\delta_{t-1}$. Therefore the sums over $\gamma_1,\dots,\gamma_{n_0-1}$ and over $\gamma$ turn into independent sums over $\delta_0,\dots,\delta_{n_0-1}\in\Gamma_s$, one in each factor. The formula for the kernel of $K\cdot S_x$ in Section~\ref{sec:transfer-operators}, at $x=-h^{\mathrm{el}}_t$, reads
\begin{multline*}
\sum_{\delta\in\Gamma_s}\int_{\mathrm{box}_0}\mathrm{d}a\,(2\pi)^{-\frac{l}{2}}\cdot\exp\left(-\frac{1}{2}|\tilde{A}_{t+1}+\delta-\tilde{A}_t-d_sa+h^{\mathrm{el}}_t|^2\right)\cdot W(a)\\
=\left(K\cdot S_{-h^{\mathrm{el}}_t}\right)([\tilde{A}_{t+1}],[\tilde{A}_t]).
\end{multline*}
Thus the summation over $\delta_t$ and the integration over $a_t\in\mathrm{box}_0$ turn $\Phi_t$ into
\[
Y^w_t([\tilde{A}_{t+1}],[\tilde{A}_t]):=M_{w,h^{\mathrm{mag}}_{t+1}}(\tilde{A}_{t+1})^{*}\cdot\left(K\cdot S_{-h^{\mathrm{el}}_t}\right)([\tilde{A}_{t+1}],[\tilde{A}_t])\cdot M_{w,h^{\mathrm{mag}}_t}(\tilde{A}_t),
\]
which is the kernel of the component in the summand $w$ of $\mathscr{H}=\bigoplus_{w\in\mathbb{Z}}\mathscr{H}_1$ of the operator
\[
Y_t:=M_{(t+1)}^{*}\cdot K\cdot S_{-h^{\mathrm{el}}_t}\cdot M_{(t)},\qquad M_{(t)}:=M_{h^{\mathrm{mag}}_t}.
\]
As $X_{(t)}=K^{\frac{1}{2}}\cdot M_{(t)}$, and as $K^{\frac{1}{2}}$ is self-adjoint, satisfies $(K^{\frac{1}{2}})^2=K$ and commutes with $S_x$ (Section~\ref{sec:transfer-operators}), one has $Y_t=X_{(t+1)}^{*}\cdot S_{-h^{\mathrm{el}}_t}\cdot X_{(t)}$. After renaming $\tilde{A}_t$ as $A_t$ it follows that
\begin{multline*}
Z_n(h)=(2\pi)^{\frac{l\cdot n_0}{2}}\sum_{w\in\mathbb{Z}}\int_{\mathfrak{C}^{n_0}}\mathrm{d}[A_0]\cdots\mathrm{d}[A_{n_0-1}]\\
\operatorname{tr}_{\mathcal{F}}\left(Y^w_{n_0-1}([A_0],[A_{n_0-1}])\cdot Y^w_{n_0-2}([A_{n_0-1}],[A_{n_0-2}])\cdots Y^w_0([A_1],[A_0])\right).
\end{multline*}
Each $Y_t$ is of trace class, as the product of the Hilbert--Schmidt operators $X_{(t+1)}^{*}$ and $S_{-h^{\mathrm{el}}_t}\cdot X_{(t)}$ (Lemma~\ref{lem:hilbert-schmidt-property}). Moreover, $Y_t$ maps every summand of $\mathscr{H}=\bigoplus_w\mathscr{H}_1$ into itself: $K$ and $S_{-h^{\mathrm{el}}_t}$ act in the same way on each summand, and $M_{(t)}$ and $M_{(t+1)}^{*}$ multiply the component of index $w$ by $M_{w,h^{\mathrm{mag}}_t}(A)$ and by $M_{w,h^{\mathrm{mag}}_{t+1}}(A)^{*}$, respectively. Hence the product $Y_{n_0-1}\cdots Y_0$, an operator of trace class, also maps every summand into itself, its component in the summand $w$ is the product of the components of the factors, and its trace on $\mathscr{H}$ is the sum over $w$ of the traces on $\mathscr{H}_1$ of these components. For Hilbert--Schmidt integral operators $Y$ and $Y'$ on $\mathscr{H}_1=L^2(\mathfrak{C};\mathcal{F})$ with the kernels $Y(\cdot,\cdot)$ and $Y'(\cdot,\cdot)$ one has
\[
\operatorname{Tr}(Y\cdot Y')=\int_{\mathfrak{C}}\mathrm{d}[A]\int_{\mathfrak{C}}\mathrm{d}[A']\,\operatorname{tr}_{\mathcal{F}}\left(Y([A],[A'])\cdot Y'([A'],[A])\right).
\]
Since $n_0\ge2$, this formula applies, in each summand $w$, to the components of the first factor $Y_{n_0-1}$ and of the product $Y_{n_0-2}\cdots Y_0$ of the others; the kernel of the latter is obtained by integrating $Y^w_{n_0-2}([A_{n_0-1}],[A_{n_0-2}])\cdots Y^w_0([A_1],[A_0])$ over $[A_1],\dots,[A_{n_0-2}]$ (no integration occurs if $n_0=2$). The formula for $Z_n(h)$ follows. For $h=0$ one has $Z_n(0)=Z_n[1]$ and $X_{(t)}=X_0$ for every $t$, the translation $S_0$ is the identity, and $X_0^{*}\cdot X_0=T$.
\end{proof}

The direction $1$ may be replaced by either of the other two: Proposition~\ref{prop:trace-formula} holds with the direction $\mu\in\{2,3\}$ as time, when $n_0$, the cross-section $\Lambda_s$ and the objects of Sections~\ref{sec:factorisation-determinant-time}--\ref{sec:transfer-operators} are formed for that direction. To see this, let $n'$ arise from $n$ by the exchange of $n_1$ and $n_{\mu}$, and let $g\colon\Lambda_n\to\Lambda_{n'}$ be the lattice isometry, as in \eqref{eq:partition-functions-8}, that exchanges the coordinates $x_1$ and $x_{\mu}$. It maps the plaquettes of $\Lambda_n$ bijectively onto those of $\Lambda_{n'}$, preserving or reversing their orientations; so the pull-back $g^{*}\colon C^2_{n'}\to C^2_n$ is orthogonal, and with $g_{*}:=(g^{*})^{-1}$
\[
\exp\left(-\langle h,g^{*}F'\rangle-\frac{1}{2}\|h\|^2\right)=\exp\left(-\langle g_{*}h,F'\rangle-\frac{1}{2}\|g_{*}h\|^2\right),\qquad F'\in C^2_{n'}.
\]
The invariance $Z_{n'}[G\circ g^{*}]=Z_n[G]$ of Section~\ref{sec:summation-all-fluxes} therefore gives $Z_n(h)=Z_{n'}(g_{*}h)$, and Proposition~\ref{prop:trace-formula} for $\Lambda_{n'}$ and $g_{*}h$, transported by $g$, is the assertion for $\Lambda_n$ with the direction $\mu$ as time.

\subsection{Inequalities from the trace formula}\label{sec:inequalities-trace-formula}

\subsubsection{Schatten classes}\label{sec:schatten-classes}

Let $1\le p<\infty$. The Schatten class $\mathscr{J}_p$ of a separable Hilbert space is the set of the compact operators $Y$ on it with
\[
\|Y\|_p:=\left(\operatorname{Tr}|Y|^p\right)^{\frac{1}{p}}<\infty,\qquad|Y|:=(Y^{*}\cdot Y)^{\frac{1}{2}}.
\]
Thus $\mathscr{J}_1$ is the class of the operators of trace class, and $\mathscr{J}_2$ the class of the Hilbert--Schmidt operators.

\begin{lemma}[Schatten norms and traces]\label{lem:schatten-norms-traces}
\begin{enumerate}
\item[(a)] Let $k\ge1$, let $p_1,\dots,p_k\in[1,\infty)$ with $\sum_{i=1}^k\frac{1}{p_i}=1$, and let $Y_i\in\mathscr{J}_{p_i}$ for $i=1,\dots,k$. Then $Y_1\cdots Y_k$ is of trace class, and
\[
\left|\operatorname{Tr}(Y_1\cdots Y_k)\right|\le\prod_{i=1}^k\|Y_i\|_{p_i}.
\]
\item[(b)] Let $1\le p'\le p<\infty$. Every $Y\in\mathscr{J}_{p'}$ belongs to $\mathscr{J}_p$, and $\|Y\|_p\le\|Y\|_{p'}$. For $Y\in\mathscr{J}_p$ and every unitary operator $U$ one has $\|U\cdot Y\|_p=\|Y\cdot U\|_p=\|Y\|_p$ and $\|Y^{*}\|_p=\|Y\|_p$; if $p$ is even, then $\|Y\|_p^p=\operatorname{Tr}(Y^{*}\cdot Y)^{\frac{p}{2}}$.
\item[(c)] Let $B\ge0$ be of trace class. Then $\operatorname{Tr}B^k\le(\operatorname{Tr}B)^k$ for every integer $k\ge1$. If in addition $B\ne0$, the function $s\mapsto\log\operatorname{Tr}B^s$ is convex on $[1,\infty)$.
\end{enumerate}
\end{lemma}

\begin{proof}
(a) By H\"older's inequality for Schatten norms, $Y\cdot Y'\in\mathscr{J}_r$ and $\|Y\cdot Y'\|_r\le\|Y\|_p\cdot\|Y'\|_{\tilde{p}}$ whenever $Y\in\mathscr{J}_p$, $Y'\in\mathscr{J}_{\tilde{p}}$ and $\frac{1}{r}=\frac{1}{p}+\frac{1}{\tilde{p}}\le1$. Induction on $j$ gives $\|Y_1\cdots Y_j\|_{r_j}\le\prod_{i=1}^j\|Y_i\|_{p_i}$ for $j=1,\dots,k$, where $\frac{1}{r_j}=\sum_{i=1}^j\frac{1}{p_i}\le1$; thus only exponents $r_j\ge1$ occur. For $j=k$ one has $r_k=1$, and $|\operatorname{Tr}Y|\le\|Y\|_1$ for every operator $Y$ of trace class.

(b) For a compact operator $Y$, $\|Y\|_p$ is the $\ell^p$-norm of the sequence of the non-zero singular values of $Y$, that is, of the non-zero eigenvalues of $|Y|$, repeated according to multiplicity. These sequences coincide for $Y$, $U\cdot Y$, $Y\cdot U$ and $Y^{*}$, and $\|x\|_{\ell^p}\le\|x\|_{\ell^{p'}}$ for every sequence $x$ when $p\ge p'$. For even $p$ one has $|Y|^p=(Y^{*}\cdot Y)^{\frac{p}{2}}$.

(c) Let $b_1,b_2,\dots>0$ be the non-zero eigenvalues of $B$, repeated according to multiplicity. For $s\ge1$ one has $\operatorname{Tr}B^s=\sum_ib_i^s<\infty$, since $b_i^s\le\|B\|^{s-1}\cdot b_i$. The expansion of $\left(\sum_ib_i\right)^k$ is a sum of non-negative terms among which the $b_i^k$ occur; hence $\sum_ib_i^k\le\left(\sum_ib_i\right)^k$, and for $B=0$ both sides vanish. For $s,s'\ge1$ and $0<u<1$, H\"older's inequality for sums with the exponents $\frac{1}{u}$ and $\frac{1}{1-u}$ gives
\[
\sum_ib_i^{u\cdot s+(1-u)\cdot s'}\le\left(\sum_ib_i^s\right)^u\cdot\left(\sum_ib_i^{s'}\right)^{1-u}.
\]
If $B\ne0$, all these sums are positive, and taking logarithms gives the convexity.
\end{proof}

\subsubsection{Gaussian domination and its consequences}\label{sec:gaussian-domination-consequences}

\begin{theorem}[Gaussian domination]\label{thm:gaussian-domination}
Let $n=(n_1,n_2,n_3)$ with integers $n_{\mu}\ge2$, each of them even or odd, and let $q>0$ and $\boldsymbol{\mu}>-1$. Then
\[
0<Z_n(h)\le Z_n(0)\qquad\text{for every }h\in C^2_n.
\]
\end{theorem}

\begin{proof}
Take the direction $1$ as time, so that $n_0=n_1$; let $X_{(t)}$, $t\in\mathbb{Z}/n_0\mathbb{Z}$, be the operators of Proposition~\ref{prop:trace-formula}, and let $U_t:=S_{-h^{\mathrm{el}}_t}$, a unitary operator on $\mathscr{H}$. All traces are taken on $\mathscr{H}$. Every $X_{(t)}$ is a Hilbert--Schmidt operator by Lemma~\ref{lem:hilbert-schmidt-property}. The leftmost factor of the product in Proposition~\ref{prop:trace-formula} is $X_{(n_0)}^{*}=X_{(0)}^{*}$. As $\operatorname{Tr}(Y\cdot Y')=\operatorname{Tr}(Y'\cdot Y)$ for Hilbert--Schmidt operators $Y$ and $Y'$, and as both $X_{(0)}^{*}$ and the product of the remaining factors are Hilbert--Schmidt operators, this factor may be moved to the right end; regrouping then gives
\begin{align*}
(2\pi)^{-\frac{l\cdot n_0}{2}}\cdot Z_n(h)&=\operatorname{Tr}\prod_{t=n_0-1,\dots,0}\left(X_{(t+1)}^{*}\cdot U_t\cdot X_{(t)}\right)\\
&=\operatorname{Tr}\prod_{t=n_0-1,\dots,0}\left(U_t\cdot X_{(t)}\cdot X_{(t)}^{*}\right).
\end{align*}
The operators $\hat{T}_t:=X_{(t)}\cdot X_{(t)}^{*}$ are non-negative and of trace class (Lemma~\ref{lem:hilbert-schmidt-property}), and hence belong to $\mathscr{J}_{n_0}$ by Lemma~\ref{lem:schatten-norms-traces}(b). Lemma~\ref{lem:schatten-norms-traces}(a), with $k=n_0$ and $p_1=\dots=p_k=n_0$, and Lemma~\ref{lem:schatten-norms-traces}(b), applied to the unitary operators $U_t$, give
\begin{align*}
\left|\operatorname{Tr}\prod_{t=n_0-1,\dots,0}\left(U_t\cdot\hat{T}_t\right)\right|&\le\prod_t\|U_t\cdot\hat{T}_t\|_{n_0}=\prod_t\|\hat{T}_t\|_{n_0}\\
&=\prod_t\left(\operatorname{Tr}\left(X_{(t)}\cdot X_{(t)}^{*}\right)^{n_0}\right)^{\frac{1}{n_0}}\\
&=\prod_t\left(\operatorname{Tr}\left(X_{(t)}^{*}\cdot X_{(t)}\right)^{n_0}\right)^{\frac{1}{n_0}}.
\end{align*}
Here $\|\hat{T}_t\|_{n_0}^{n_0}=\operatorname{Tr}\hat{T}_t^{n_0}$ because $\hat{T}_t\ge0$, and $\operatorname{Tr}(X\cdot X^{*})^{n_0}=\operatorname{Tr}\left(X\cdot(X^{*}\cdot X)^{n_0-1}\cdot X^{*}\right)=\operatorname{Tr}(X^{*}\cdot X)^{n_0}$ for $X=X_{(t)}$ by the cyclicity of the trace. For $t\in\mathbb{Z}/n_1\mathbb{Z}$ let $h^{(t)}\in C^2_n$ be the two-cochain whose electric part vanishes and whose magnetic part equals $h^{\mathrm{mag}}_t$ on every slice. With $h^{(t)}$ in place of $h$, every operator $X_{(t')}$ of Proposition~\ref{prop:trace-formula} is $X_{h^{\mathrm{mag}}_t}=X_{(t)}$, and every translation in it is $S_0$, the identity; so Proposition~\ref{prop:trace-formula} gives
\[
(2\pi)^{\frac{l\cdot n_0}{2}}\cdot\operatorname{Tr}\left(X_{(t)}^{*}\cdot X_{(t)}\right)^{n_0}=Z_n(h^{(t)}).
\]
Since $Z_n(h)\le|Z_n(h)|$, $n_0=n_1$ and $(2\pi)^{\frac{l\cdot n_0}{2}}=\prod_t\left((2\pi)^{\frac{l\cdot n_0}{2}}\right)^{\frac{1}{n_0}}$, the last three displays give
\[
Z_n(h)\le\prod_{t\in\mathbb{Z}/n_1\mathbb{Z}}Z_n(h^{(t)})^{\frac{1}{n_1}}.
\]
The two-cochain $h^{(t)}$ vanishes on all plaquettes except the $(23)$-plaquettes. Now take the direction $2$ as time, which is admissible by the paragraph following Proposition~\ref{prop:trace-formula}. Then the $(12)$- and the $(23)$-plaquettes carry the electric parts and the $(13)$-plaquettes the magnetic parts, so that the magnetic part of $h^{(t)}$ vanishes; consequently the two-cochains formed from $h^{(t)}$ in the way $h^{(t)}$ was formed from $h$, now with the direction $2$ as time, are all zero. The inequality just proved, with the direction $2$ as time and $h^{(t)}$ in place of $h$, therefore gives
\[
Z_n(h^{(t)})\le\prod_{t'\in\mathbb{Z}/n_2\mathbb{Z}}Z_n(0)^{\frac{1}{n_2}}=Z_n(0),
\]
and hence $Z_n(h)\le Z_n(0)$.

It remains to prove $Z_n(h)>0$. By Definition~\ref{def:partition-functions-all}, $Z_n(h)$ is the sum over the closed $\nu\in C^2_n(\mathbb{Z})$ of the integrals over $\left[0,\frac{2\pi}{q}\right)^{\mathrm{Lk}_n}$ of $\exp\left(-\frac{1}{2}\|F+h\|^2\right)\cdot\mathfrak{w}(A)$, where $F=dA+\frac{2\pi}{q}\nu$. These integrands are non-negative by \eqref{eq:partition-functions-3}. The integrand of the term $\nu=0$ is continuous by \eqref{eq:partition-functions-6}, and at $A=0$ it equals $\exp\left(-\frac{1}{2}\|h\|^2\right)\cdot\mathfrak{w}(0)>0$, since $\mathfrak{w}(0)>0$ (Section~\ref{sec:summation-all-fluxes}). It is therefore positive on a subset of $\left[0,\frac{2\pi}{q}\right)^{\mathrm{Lk}_n}$ of positive measure, and $Z_n(h)>0$.
\end{proof}

For even as well as for odd side lengths, Theorem~\ref{thm:gaussian-domination-positive} is the counterpart of Theorem~\ref{thm:gaussian-domination} for the Proca parent of \eqref{eq:gaussian-bounds-2}, and Theorem~\ref{thm:limit-vanishing-proca} is its counterpart for the harmonic-averaged law, which is $E_{n,0}$ by Lemma~\ref{lem:flux-sectors}(b).

\begin{corollary}\label{cor:partition-functions-1}
\begin{enumerate}
\item[(a)] For every $g\in C^2_n$,
\[
E_n\left[\exp\langle g,F\rangle\right]\le\exp\left(\frac{1}{2}\|g\|^2\right),
\]
and for $g_1,\dots,g_k\in C^2_n$ and $z_1,\dots,z_k\in\mathbb{C}$,
\[
\left|E_n\left[\exp\sum_{i=1}^kz_i\cdot\langle g_i,F\rangle\right]\right|\le\exp\left(\frac{1}{2}\left\|\sum_{i=1}^k(\operatorname{Re}z_i)\cdot g_i\right\|^2\right).
\]
\item[(b)] For every $g\in C^2_n$ and every integer $r\ge1$,
\begin{gather*}
E_n[\langle g,F\rangle]=0,\qquad E_n[\langle g,F\rangle^2]\le\|g\|^2,\\
E_n[\langle g,F\rangle^{2r}]\le\left(2r\cdot\exp(1)\cdot\|g\|^2\right)^r,\qquad E_n[\langle g,F\rangle^4]\le\frac{3}{2}\cdot\exp(2)\cdot\|g\|^4;
\end{gather*}
and for $g_1,\dots,g_4\in C^2_n$ the joint fourth cumulant of $\langle g_1,F\rangle,\dots,\langle g_4,F\rangle$ under $E_n$ satisfies
\[
\left|\kappa_4\left(\langle g_1,F\rangle,\dots,\langle g_4,F\rangle\right)\right|\le15\prod_{i=1}^4\|g_i\|.
\]
\item[(c)] \emph{Peierls bound.} Let $g_1,\dots,g_m\in C^2_n$ be non-zero with pairwise disjoint supports, and let $p_1,\dots,p_m\ge0$. Then
\[
P_n\left(|\langle g_i,F\rangle|\ge p_i\cdot\|g_i\|\ \text{for}\ i=1,\dots,m\right)\le\prod_{i=1}^m2\exp\left(-\frac{p_i^2}{2}\right),
\]
where $P_n(S):=E_n[\mathbf{1}_S]$ for a Borel set $S\subset C^2_n$.
\item[(d)] \emph{Flux sectors.} $P_n(w)\le\exp\left(-\frac{1}{2}\|h_w\|^2\right)$ for every $w\in\mathbb{Z}^3$. Let
\[
\varsigma_n:=\sum_{w\in\mathbb{Z}^3\setminus\{0\}}\exp\left(-\frac{1}{2}\|h_w\|^2\right).
\]
Then $1-\varsigma_n\le P_n(0)\le1$; and if $\varsigma_n\le\frac{1}{4}$, then for every $G\in\mathrm{Gr}$
\[
\left|E_{n,0}[G]-E_n[G]\right|\le3\,\varsigma_n^{\frac{1}{2}}\cdot\left(E_n\left[|G|^2\right]\right)^{\frac{1}{2}}.
\]
\end{enumerate}
\end{corollary}

\begin{proof}
The functions of $F$ to which $E_n$ is applied below belong to $\mathrm{Gr}$. By Definition~\ref{def:partition-functions-all} and \eqref{eq:partition-functions-3}, $E_n$ is the integral with respect to a probability measure on the pairs $(\nu,A)$.

(a) By Definition~\ref{def:partition-functions-all},
\[
Z_n(-g)=Z_n\left[\exp\left(\langle g,F\rangle-\frac{1}{2}\|g\|^2\right)\right]=\exp\left(-\frac{1}{2}\|g\|^2\right)\cdot Z_n\left[\exp\langle g,F\rangle\right],
\]
and Theorem~\ref{thm:gaussian-domination} gives $Z_n(-g)\le Z_n(0)=Z_n[1]$. Division by $Z_n[1]$ yields the first inequality. Since $E_n$ is the integral with respect to a positive measure and $|\exp z|=\exp(\operatorname{Re}z)$ for $z\in\mathbb{C}$,
\[
\left|E_n\left[\exp\sum_iz_i\cdot\langle g_i,F\rangle\right]\right|\le E_n\left[\exp\left\langle\sum_i(\operatorname{Re}z_i)\cdot g_i,F\right\rangle\right],
\]
and the first inequality, applied to $\sum_i(\operatorname{Re}z_i)\cdot g_i$, gives the second.

(b) The law of $F$ under $E_n$ is invariant under $F\mapsto-F$ (Section~\ref{sec:summation-all-fluxes}); hence $E_n[\langle g,F\rangle]=-E_n[\langle g,F\rangle]=0$. For $s>0$, part (a) applied to $s\cdot g$ and to $-s\cdot g$ gives
\[
E_n\left[\cosh(s\cdot\langle g,F\rangle)\right]\le\exp\left(\frac{s^2}{2}\cdot\|g\|^2\right).
\]
The second claim is obtained by comparing the terms of second order in $s$: as $\cosh x\ge1+\frac{x^2}{2}$, one has $\frac{s^2}{2}\cdot E_n[\langle g,F\rangle^2]\le\exp\left(\frac{s^2}{2}\cdot\|g\|^2\right)-1$, and division by $\frac{s^2}{2}$ and the limit $s\to0$ give $E_n[\langle g,F\rangle^2]\le\|g\|^2$. Since $\frac{x^{2r}}{(2r)!}\le\cosh x$ for real $x$,
\[
E_n[\langle g,F\rangle^{2r}]\le(2r)!\cdot s^{-2r}\cdot\exp\left(\frac{s^2}{2}\cdot\|g\|^2\right),\qquad s>0.
\]
For $g=0$ there is nothing to prove. Otherwise the choice $s^2=\frac{2r}{\|g\|^2}$ and $(2r)!\le(2r)^{2r}$ give
\[
E_n[\langle g,F\rangle^{2r}]\le(2r)^{2r}\cdot\left(\frac{\|g\|^2}{2r}\right)^r\cdot\exp(r)=\left(2r\cdot\exp(1)\cdot\|g\|^2\right)^r,
\]
and for $r=2$ the same choice, now with $4!=24$, gives
\[
E_n[\langle g,F\rangle^4]\le24\cdot\frac{\|g\|^4}{16}\cdot\exp(2)=\frac{3}{2}\cdot\exp(2)\cdot\|g\|^4.
\]
Let $Y_i:=\langle g_i,F\rangle$ for $i=1,\dots,4$. By the first claim these random variables are centred, and for centred random variables the joint fourth cumulant is the mixed fourth moment minus the sum, over the three partitions of $\{1,2,3,4\}$ into two pairs $\{i,j\}$ and $\{i',j'\}$, of the products $E_n[Y_i\cdot Y_j]\cdot E_n[Y_{i'}\cdot Y_{j'}]$:
\begin{align*}
\kappa_4(Y_1,\dots,Y_4)&=E_n[Y_1\cdot Y_2\cdot Y_3\cdot Y_4]-E_n[Y_1\cdot Y_2]\cdot E_n[Y_3\cdot Y_4]\\
&\quad-E_n[Y_1\cdot Y_3]\cdot E_n[Y_2\cdot Y_4]-E_n[Y_1\cdot Y_4]\cdot E_n[Y_2\cdot Y_3].
\end{align*}
H\"older's inequality and the bound on the fourth moment give $|E_n[Y_1\cdot Y_2\cdot Y_3\cdot Y_4]|\le\prod_iE_n[Y_i^4]^{\frac{1}{4}}\le\frac{3}{2}\cdot\exp(2)\cdot\prod_i\|g_i\|$, and the Cauchy--Schwarz inequality together with the second claim gives $|E_n[Y_i\cdot Y_j]|\le\|g_i\|\cdot\|g_j\|$. Hence
\[
|\kappa_4(Y_1,\dots,Y_4)|\le\left(\frac{3}{2}\cdot\exp(2)+3\right)\prod_{i=1}^4\|g_i\|\le15\prod_{i=1}^4\|g_i\|,
\]
because $\frac{3}{2}\cdot\exp(2)<12$.

(c) For signs $\boldsymbol{\sigma}_1,\dots,\boldsymbol{\sigma}_m\in\{-1,1\}$ put $g:=\sum_i\boldsymbol{\sigma}_i\cdot p_i\cdot\frac{g_i}{\|g_i\|}$. The $g_i$ have pairwise disjoint supports and are therefore orthogonal, so that $\|g\|^2=\sum_ip_i^2$; and $\langle g,F\rangle\ge\sum_ip_i^2$ whenever $\boldsymbol{\sigma}_i\cdot\langle g_i,F\rangle\ge p_i\cdot\|g_i\|$ for all $i$. Chernoff's bound, that is, Markov's inequality applied to $\exp\langle g,F\rangle$, and part (a) give
\begin{multline*}
P_n\left(\boldsymbol{\sigma}_i\cdot\langle g_i,F\rangle\ge p_i\cdot\|g_i\|\ \text{for}\ i=1,\dots,m\right)\\
\le\exp\left(-\sum_ip_i^2\right)\cdot E_n\left[\exp\langle g,F\rangle\right]\le\exp\left(-\frac{1}{2}\sum_ip_i^2\right).
\end{multline*}
The event in (c) is contained in the union of these events over the $2^m$ choices of the signs, and $2^m\cdot\exp\left(-\frac{1}{2}\sum_ip_i^2\right)=\prod_i2\exp\left(-\frac{p_i^2}{2}\right)$.

(d) Let $c\in\mathrm{Harm}^2_n$. For a closed $\nu\in C^2_n(\mathbb{Z})$ with $w(\nu)=w$ and for $A\in C^1_n$, Lemma~\ref{lem:flux-sectors}(a) gives $\|F+c\|^2=\|d(A+\beta_{\nu})\|^2+\|h_w+c\|^2$, where $F=dA+\frac{2\pi}{q}\nu$. Grouping the terms of the series of Definition~\ref{def:partition-functions-all} for $Z_n(c)$ according to the flux, one obtains
\begin{gather*}
Z_n(c)=\sum_{w\in\mathbb{Z}^3}\exp\left(-\frac{1}{2}\|h_w+c\|^2\right)\cdot y_w,\\
y_w:=\sum_{\nu}\int_{\left[0,\frac{2\pi}{q}\right)^{\mathrm{Lk}_n}}\mathrm{d}A\,\exp\left(-\frac{1}{2}\|d(A+\beta_{\nu})\|^2\right)\cdot\mathfrak{w}(A),
\end{gather*}
where $\nu$ ranges over the closed $\nu\in C^2_n(\mathbb{Z})$ with $w(\nu)=w$. The numbers $y_w$ are non-negative by \eqref{eq:partition-functions-3} and do not depend on $c$, and $Z_{n,w}[1]=\exp\left(-\frac{1}{2}\|h_w\|^2\right)\cdot y_w$ (the case $c=0$, restricted to the flux $w$). Let $w_0\in\mathbb{Z}^3$. Since $-h_{w_0}\in\mathrm{Harm}^2_n$, the choice $c=-h_{w_0}$, the non-negativity of all terms of the sum and Theorem~\ref{thm:gaussian-domination} give
\[
y_{w_0}\le Z_n(-h_{w_0})\le Z_n(0)=Z_n[1],
\]
so that $P_n(w_0)=\frac{Z_{n,w_0}[1]}{Z_n[1]}\le\exp\left(-\frac{1}{2}\|h_{w_0}\|^2\right)$. This is the first claim. The series $\varsigma_n$ converges, since $\|h_w\|^2$ is a positive definite quadratic form in $w$ (Lemma~\ref{lem:flux-sectors}(a)). As $\sum_wP_n(w)=1$, summation of the first claim over $w\ne0$ gives $1-P_n(0)\le\varsigma_n$, and $P_n(0)\le1$. For the third claim let $G\in\mathrm{Gr}$; then also $|G|^2\in\mathrm{Gr}$. Splitting the series of Definition~\ref{def:partition-functions-all} into the flux $0$ and the other fluxes gives
\[
E_n[G]=P_n(0)\cdot E_{n,0}[G]+\frac{1}{Z_n[1]}\sum_{w\ne0}Z_{n,w}[G].
\]
The Cauchy--Schwarz inequality for the positive measure of Definition~\ref{def:partition-functions-all}, restricted to the configurations with non-zero flux, gives
\begin{align*}
\left|\sum_{w\ne0}Z_{n,w}[G]\right|&\le\left(\sum_{w\ne0}Z_{n,w}[1]\right)^{\frac{1}{2}}\cdot\left(\sum_{w\ne0}Z_{n,w}\left[|G|^2\right]\right)^{\frac{1}{2}}\\
&\le\left((1-P_n(0))\cdot Z_n[1]\right)^{\frac{1}{2}}\cdot Z_n\left[|G|^2\right]^{\frac{1}{2}}.
\end{align*}
Hence
\[
\left|E_{n,0}[G]-E_n[G]\right|\le(1-P_n(0))\cdot\left|E_{n,0}[G]\right|+(1-P_n(0))^{\frac{1}{2}}\cdot\left(E_n\left[|G|^2\right]\right)^{\frac{1}{2}}.
\]
Moreover, since $Z_{n,0}[|G|]\le Z_n[|G|]$ and $Z_{n,0}[1]=P_n(0)\cdot Z_n[1]$,
\[
\left|E_{n,0}[G]\right|\le E_{n,0}[|G|]=\frac{Z_{n,0}[|G|]}{P_n(0)\cdot Z_n[1]}\le\frac{E_n[|G|]}{P_n(0)}.
\]
If $\varsigma_n\le\frac{1}{4}$, then $P_n(0)\ge\frac{3}{4}$ and $1-P_n(0)\le\varsigma_n\le\frac{1}{2}\varsigma_n^{\frac{1}{2}}$. Together with $E_n[|G|]\le\left(E_n\left[|G|^2\right]\right)^{\frac{1}{2}}$ this shows that the right side of the inequality for $\left|E_{n,0}[G]-E_n[G]\right|$ is at most
\[
\left(\frac{4}{3}\cdot\frac{1}{2}+1\right)\cdot\varsigma_n^{\frac{1}{2}}\cdot\left(E_n\left[|G|^2\right]\right)^{\frac{1}{2}}\le3\,\varsigma_n^{\frac{1}{2}}\cdot\left(E_n\left[|G|^2\right]\right)^{\frac{1}{2}}.
\]
\end{proof}

\subsubsection{Convexity in the side lengths}\label{sec:convexity-side-lengths}

For an integer $\boldsymbol{\lambda}\ge2$ and integers $N_1,N_2,N_3\ge1$ put $N=(N_1,N_2,N_3)$ and $\boldsymbol{\lambda} N=(\boldsymbol{\lambda} N_1,\boldsymbol{\lambda} N_2,\boldsymbol{\lambda} N_3)$. Every side $\boldsymbol{\lambda} N_{\mu}$ is at least $2$, so $\Lambda_{\boldsymbol{\lambda} N}$ is a torus as in Section~\ref{sec:tori-cochains}, and $Z_{\boldsymbol{\lambda} N}[1]$ is given by Definition~\ref{def:partition-functions-all}. A function $\psi\colon\{1,2,\dots\}\to\mathbb{R}$ is called \emph{convex} if, for all integers $1\le a<b<c$,
\[
(c-a)\cdot\psi(b)\le(c-b)\cdot\psi(a)+(b-a)\cdot\psi(c),
\]
and a function $\phi\colon\{1,2,\dots\}^3\to\mathbb{R}$ is called \emph{convex in each variable separately} if, for all integers $a,b,c\ge1$, the functions $x\mapsto\phi(x,b,c)$, $x\mapsto\phi(a,x,c)$ and $x\mapsto\phi(a,b,x)$ on $\{1,2,\dots\}$ are convex.

\begin{theorem}\label{thm:partition-functions-1}
Let $\boldsymbol{\lambda}\ge2$ be an integer, and let $N=(N_1,N_2,N_3)$ with integers $N_{\mu}\ge1$, even or odd. As a function of $N_1\in\{1,2,\dots\}$, with $N_2$ and $N_3$ fixed, $\log Z_{\boldsymbol{\lambda} N}[1]$ is convex, and
\[
Z_{\boldsymbol{\lambda}(kN_1,N_2,N_3)}[1]\le Z_{\boldsymbol{\lambda} N}[1]^k\qquad\text{for every integer }k\ge1.
\]
The same holds in the directions $2$ and $3$: $\log Z_{\boldsymbol{\lambda} N}[1]$ is a convex function of $N_2$ for fixed $N_1,N_3$ and a convex function of $N_3$ for fixed $N_1,N_2$, and for every integer $k\ge1$
\[
Z_{\boldsymbol{\lambda}(N_1,kN_2,N_3)}[1]\le Z_{\boldsymbol{\lambda} N}[1]^k,\qquad Z_{\boldsymbol{\lambda}(N_1,N_2,kN_3)}[1]\le Z_{\boldsymbol{\lambda} N}[1]^k.
\]
\end{theorem}

\begin{proof}
Take the direction $1$ as time for the torus $\Lambda_{\boldsymbol{\lambda} N}$, as at the beginning of Section~\ref{sec:partition-function-trace}. Then $n_0=\boldsymbol{\lambda} N_1$, and the cross-section $\Lambda_s$ is the two-dimensional torus with the sides $\boldsymbol{\lambda} N_2$ and $\boldsymbol{\lambda} N_3$. Let $l$ be the number of links of $\Lambda_s$, and let $T=X_0^{*}\cdot X_0$ be the operator on $\mathscr{H}$ of Section~\ref{sec:transfer-operators}. For the given $q>0$ and $\boldsymbol{\mu}>-1$, the number $l$ and the operator $T$ are determined by $\Lambda_s$ alone; in particular they do not change when $N_1$ varies, since $T$ is formed from the data of Section~\ref{sec:factorisation-determinant-time}, which do not depend on $n_0$, and from the objects on $\Lambda_s$ of Section~\ref{sec:paths-spatial-fluxes}. By Lemma~\ref{lem:hilbert-schmidt-property}, $T$ is a non-negative trace class operator. Hence $B:=T^{\boldsymbol{\lambda}}$ and its powers $B^m$, $m\ge1$, are non-negative trace class operators (the trace class is an ideal in the algebra of bounded operators). Proposition~\ref{prop:trace-formula}, applied with $n_0=\boldsymbol{\lambda} N_1$, gives
\[
Z_{\boldsymbol{\lambda} N}[1]=(2\pi)^{\frac{l\boldsymbol{\lambda} N_1}{2}}\cdot\operatorname{Tr}T^{\boldsymbol{\lambda} N_1}=(2\pi)^{\frac{l\boldsymbol{\lambda} N_1}{2}}\cdot\operatorname{Tr}B^{N_1}.
\]
Since $Z_{\boldsymbol{\lambda} N}[1]>0$ (Section~\ref{sec:summation-all-fluxes}), it follows that $\operatorname{Tr}B^{N_1}>0$; in particular $B\ne0$.

\emph{Convexity.} By Lemma~\ref{lem:schatten-norms-traces}(c), applied to $B$, the function $s\mapsto\log\operatorname{Tr}B^s$ is convex on $[1,\infty)$. Its restriction to the positive integers is therefore convex (for integers $1\le a<b<c$, the number $b$ is a convex combination of $a$ and $c$). Taking logarithms in the last identity, one has
\[
\log Z_{\boldsymbol{\lambda}(N_1,N_2,N_3)}[1]=\frac{l\boldsymbol{\lambda}\cdot\log(2\pi)}{2}\cdot N_1+\log\operatorname{Tr}B^{N_1},
\]
where neither $l$ nor $B$ depends on $N_1$. The first term is a linear function of $N_1$ and the second a convex one, so the sum is convex.

\emph{Subadditivity.} Let $k\ge1$ be an integer. The torus $\Lambda_{\boldsymbol{\lambda}(kN_1,N_2,N_3)}$ has the same cross-section, hence the same $l$ and $B$, and the identity above with $kN_1$ in place of $N_1$ reads
\[
Z_{\boldsymbol{\lambda}(kN_1,N_2,N_3)}[1]=\left((2\pi)^{\frac{l\boldsymbol{\lambda} N_1}{2}}\right)^k\cdot\operatorname{Tr}\left(B^{N_1}\right)^k.
\]
By Lemma~\ref{lem:schatten-norms-traces}(c), applied to the non-negative trace class operator $B^{N_1}$, one has $\operatorname{Tr}\left(B^{N_1}\right)^k\le\left(\operatorname{Tr}B^{N_1}\right)^k$. It follows that
\[
Z_{\boldsymbol{\lambda}(kN_1,N_2,N_3)}[1]\le\left((2\pi)^{\frac{l\boldsymbol{\lambda} N_1}{2}}\cdot\operatorname{Tr}B^{N_1}\right)^k=Z_{\boldsymbol{\lambda} N}[1]^k.
\]

\emph{The directions $2$ and $3$.} As noted after Proposition~\ref{prop:trace-formula}, that proposition holds with each of the three coordinate directions as time, by the invariance under lattice isometries of Section~\ref{sec:summation-all-fluxes}. With the direction $2$ (resp.\ $3$) as time, the cross-section is the torus with the sides $\boldsymbol{\lambda} N_1$ and $\boldsymbol{\lambda} N_3$ (resp.\ $\boldsymbol{\lambda} N_1$ and $\boldsymbol{\lambda} N_2$), and $l$ and $T$ are formed on it. The argument above, with $N_2$ (resp.\ $N_3$) in place of $N_1$, proves the assertions in the directions $2$ and $3$.
\end{proof}

\subsubsection{A lower bound along the diagonal}\label{sec:lower-bound-along}

\begin{lemma}\label{lem:partition-functions-1}
Let $\phi\colon\{1,2,\dots\}^3\to\mathbb{R}$ be convex in each variable separately, and suppose that
\begin{align*}
\phi(ka,b,c)&\le k\cdot\phi(a,b,c),\\
\phi(a,kb,c)&\le k\cdot\phi(a,b,c),\\
\phi(a,b,kc)&\le k\cdot\phi(a,b,c)
\end{align*}
for all integers $a,b,c,k\ge1$. Let $j\ge2$ be an integer, and put $D_j=j^3\cdot\phi(1,1,1)-\phi(j,j,j)$. Then $D_j\ge0$, and for every integer $k\ge j$
\[
\phi(k,k,k)\ge k^3\cdot\phi(1,1,1)-\left(\frac{k-1}{j-1}\right)^3\cdot D_j.
\]
\end{lemma}

\begin{proof}
Three applications of the hypothesis $\phi(ka,b,c)\le k\cdot\phi(a,b,c)$ and of its analogues in the second and the third variable give
\[
\phi(j,j,j)\le j\cdot\phi(1,j,j)\le j^2\cdot\phi(1,1,j)\le j^3\cdot\phi(1,1,1),
\]
that is, $D_j\ge0$.

Let $k\ge j$ be an integer, and put $\boldsymbol{\alpha}=\frac{k-1}{j-1}$. Then $\boldsymbol{\alpha}\ge1$ and $\boldsymbol{\alpha}(j-1)=k-1$, so that
\[
j=\frac{1}{\boldsymbol{\alpha}}\cdot k+\left(1-\frac{1}{\boldsymbol{\alpha}}\right)\cdot1
\]
is a convex combination of $k$ and $1$. Let $\psi$ be any of the functions on $\{1,2,\dots\}$ obtained from $\phi$ by fixing two of its three arguments. If $j<k$, the convexity of $\psi$ gives
\[
\psi(j)\le\frac{1}{\boldsymbol{\alpha}}\cdot\psi(k)+\left(1-\frac{1}{\boldsymbol{\alpha}}\right)\cdot\psi(1),
\]
and if $j=k$, then $\boldsymbol{\alpha}=1$ and this holds with equality. Multiplication by $\boldsymbol{\alpha}>0$ gives the lower bound
\[
\psi(k)\ge\boldsymbol{\alpha}\cdot\psi(j)-(\boldsymbol{\alpha}-1)\cdot\psi(1).
\]
The same hypothesis as above, with $a=1$, $b=1$ or $c=1$ according to the variable of $\psi$, gives the upper bounds $\psi(m)\le m\cdot\psi(1)$ for all integers $m\ge1$.

The lower bound, used in the first, the second and the third variable, gives
\begin{align*}
\phi(k,k,k)&\ge\boldsymbol{\alpha}\cdot\phi(j,k,k)-(\boldsymbol{\alpha}-1)\cdot\phi(1,k,k),\\
\phi(j,k,k)&\ge\boldsymbol{\alpha}\cdot\phi(j,j,k)-(\boldsymbol{\alpha}-1)\cdot\phi(j,1,k),\\
\phi(j,j,k)&\ge\boldsymbol{\alpha}\cdot\phi(j,j,j)-(\boldsymbol{\alpha}-1)\cdot\phi(j,j,1),
\end{align*}
and the upper bounds give
\begin{align*}
\phi(1,k,k)&\le k^2\cdot\phi(1,1,1),\\
\phi(j,1,k)&\le k\cdot\phi(j,1,1)\le kj\cdot\phi(1,1,1),\\
\phi(j,j,1)&\le j^2\cdot\phi(1,1,1).
\end{align*}
Insert the second lower bound into the first, and then the third into the result; the second and the third lower bound are thereby multiplied by the positive coefficients $\boldsymbol{\alpha}$ and $\boldsymbol{\alpha}^2$. In the inequality so obtained the terms $\phi(1,k,k)$, $\phi(j,1,k)$ and $\phi(j,j,1)$ carry the coefficients $-(\boldsymbol{\alpha}-1)$, $-\boldsymbol{\alpha}(\boldsymbol{\alpha}-1)$ and $-\boldsymbol{\alpha}^2(\boldsymbol{\alpha}-1)$, which are not positive; therefore each of these terms may be replaced by its upper bound. The result is
\begin{align*}
\phi(k,k,k)&\ge\boldsymbol{\alpha}^3\cdot\phi(j,j,j)\\
&\quad-\left[\boldsymbol{\alpha}^2(\boldsymbol{\alpha}-1)\cdot j^2+\boldsymbol{\alpha}(\boldsymbol{\alpha}-1)\cdot jk+(\boldsymbol{\alpha}-1)\cdot k^2\right]\cdot\phi(1,1,1).
\end{align*}
Now insert $\phi(j,j,j)=j^3\cdot\phi(1,1,1)-D_j$. Since $\boldsymbol{\alpha} j-(\boldsymbol{\alpha}-1)=\boldsymbol{\alpha}(j-1)+1=k$, three applications of this identity give
\begin{align*}
\boldsymbol{\alpha}^3j^3-\boldsymbol{\alpha}^2(\boldsymbol{\alpha}-1)\cdot j^2&=\boldsymbol{\alpha}^2j^2k,\\
\boldsymbol{\alpha}^2j^2k-\boldsymbol{\alpha}(\boldsymbol{\alpha}-1)\cdot jk&=\boldsymbol{\alpha} jk^2,\\
\boldsymbol{\alpha} jk^2-(\boldsymbol{\alpha}-1)\cdot k^2&=k^3.
\end{align*}
Hence the coefficient of $\phi(1,1,1)$ equals $k^3$, and
\[
\phi(k,k,k)\ge k^3\cdot\phi(1,1,1)-\boldsymbol{\alpha}^3\cdot D_j=k^3\cdot\phi(1,1,1)-\left(\frac{k-1}{j-1}\right)^3\cdot D_j.
\]
The proof makes no use of the signs of the values of $\phi$, since every inequality entering it was multiplied only by a factor whose sign had been determined beforehand.
\end{proof}

\section{The all-flux functional: slab functions, flux sectors and external fields}\label{sec:all-flux-functional}

For lattice masses $\boldsymbol{\mu}>-1/2$, the functional $\omega_n$ of Definition~\ref{def:all-flux-functional} is the functional $Z_n[\cdot]$ of Definition~\ref{def:partition-functions-all}, written with the directions numbered $0,1,2$ instead of $1,2,3$, with $C^k(n)$ in place of $C^k_n$ and with $\Delta(qA)$ in place of the fermion weight $\mathfrak{w}(A)$; the lattice setting and the slice factorisation are formulated anew in this notation, and of Section~\ref{sec:partition-functions-all} only \eqref{eq:partition-functions-2} is quoted. The functional $\omega_n$ is treated a second time in order to obtain a trace formula for products of slab functions which, through Bloch angles, resolves the flux sectors and holds for arbitrary slice data (Proposition~\ref{prop:trace-formula-bloch}); this formula gives Gaussian domination once more (Theorem~\ref{thm:all-flux-1}), and the section ends with the shift identities for an external field (Lemma~\ref{lem:shift-identities}) and with a bound for a single flux sector, deduced from Gaussian domination (Lemma~\ref{lem:one-flux-sector}).

Throughout this section lattice units are used: the lattice spacing is $1$, the lattice charge is $q>0$, $r:=2\pi/q$, and the lattice mass is $\boldsymbol{\mu}>-1/2$.

\subsection{The lattice setting}\label{sec:lattice-setting}

\subsubsection{Tori, cochains and lattice isometries}\label{sec:tori-cochains-lattice}

Let $n=(n_0,n_1,n_2)$ be a triple of integers $n_\mu\ge2$. Put
\[
\Lambda_n:=\prod_{\mu=0}^2\mathbb Z/n_\mu\mathbb Z,\qquad |n|:=n_0n_1n_2,
\]
and let $\hat\mu\in\mathbb Z^3$ be the unit vector of the direction $\mu\in\{0,1,2\}$. The cells of $\Lambda_n$ are the sites $x\in\Lambda_n$; the links $(x,\mu)$, joining $x$ to $x+\hat\mu$; the plaquettes $(x,\mu\nu)$ with $\mu<\nu$, whose corners are $x$, $x+\hat\mu$, $x+\hat\nu$, $x+\hat\mu+\hat\nu$; and the cubes, each labelled by the corner $x$ from which its edges $(x,0)$, $(x,1)$, $(x,2)$ issue. Every cell is oriented by its directions taken in increasing order. For $k=0,1,2,3$ let $C^k(n)$ be the real vector space of the functions on the sites, the links, the plaquettes, the cubes, respectively (the $k$-cochains). The value of $A\in C^1(n)$ at $(x,\mu)$ is denoted $A_\mu(x)$, and the value of $F\in C^2(n)$ at $(x,\mu\nu)$ is denoted $F_{\mu\nu}(x)$. The coboundary $d:C^k(n)\to C^{k+1}(n)$ is given by
\begin{align*}
(d\lambda)_\mu(x)&=\lambda(x+\hat\mu)-\lambda(x),\\
(dA)_{\mu\nu}(x)&=A_\mu(x)+A_\nu(x+\hat\mu)-A_\mu(x+\hat\nu)-A_\nu(x),\\
(dF)(x)&=\left[F_{12}(x+\hat0)-F_{12}(x)\right]-\left[F_{02}(x+\hat1)-F_{02}(x)\right]+\left[F_{01}(x+\hat2)-F_{01}(x)\right].
\end{align*}
The space $C^k(n)$ carries the counting pairing $\langle u,v\rangle:=\sum_\tau u(\tau)\,v(\tau)$, the sum over the $k$-cells $\tau$. It is extended to $C^k(n)\otimes\mathbb C$ as a complex \emph{bilinear} form, and $\|g\|^2:=\langle g,g\rangle$ for real $g$.

Every cochain is extended to the cells with the opposite orientation by a change of sign. The $k$-cell of $\Lambda_n$ labelled by $\bar x\in\Lambda_n$ with the set $I$ of its $k$ directions is identified with the class, modulo translations by the period lattice $\prod_\mu n_\mu\mathbb Z\subset\mathbb Z^3$, of the unit cube $\left\{x+\sum_{\mu\in I}t_\mu\hat\mu:\ t_\mu\in[0,1]\right\}\subset\mathbb R^3$, where $x\in\mathbb Z^3$ is a representative of $\bar x$. Let $n'$ be a second triple of integers $\ge2$. A \emph{lattice isometry} from $\Lambda_n$ onto $\Lambda_{n'}$ is an affine map
\[
\gamma:\mathbb Z^3\to\mathbb Z^3,\qquad\gamma x=\rho x+a,
\]
where $a\in\mathbb Z^3$, and $\rho$ is a signed permutation matrix such that $\rho\left(\prod_\mu n_\mu\mathbb Z\right)=\prod_\mu n'_\mu\mathbb Z$. One writes $n'=\gamma n$; the condition on $\rho$ means that $(\gamma n)_{\mu'}=n_\mu$ whenever $\rho\hat\mu=\pm\hat\mu'$. Equivalently, $\gamma$ is a composition of translations, of reflections $x_\mu\mapsto c-x_\mu$ with $c\in\mathbb Z$, and, between tori whose sides are permuted, of permutations of the axes, each of them regarded as an affine map of $\mathbb Z^3$ (a signed permutation matrix is the product of a permutation matrix and a diagonal matrix with diagonal entries $\pm1$). The matrix $\rho$ is the linear part of $\gamma$, and the bijection of $\Lambda_n$ onto $\Lambda_{\gamma n}$ induced by $\gamma$ is denoted $\gamma:\Lambda_n\to\Lambda_{\gamma n}$ as well. Extended to $\mathbb R^3$ by the same formula, $\gamma$ is an isometry which maps the unit cubes above onto unit cubes of the same kind, and the cosets of $\prod_\mu n_\mu\mathbb Z$ onto those of $\prod_\mu(\gamma n)_\mu\mathbb Z$. With the orientations transported by $\rho$, it therefore maps the oriented $k$-cells of $\Lambda_n$ bijectively onto the oriented $k$-cells of $\Lambda_{\gamma n}$; the image of an oriented cell $\tau$ is denoted $\gamma\tau$. The pullback $\gamma^*:C^k(\gamma n)\to C^k(n)$ is defined by $(\gamma^*u)(\tau):=u(\gamma\tau)$, and $\gamma_*:=(\gamma^{-1})^*:C^k(n)\to C^k(\gamma n)$; both maps are extended complex-linearly. Since $\gamma^*$ permutes the cells up to signs, it is orthogonal, so that
\[
\langle g,\gamma^*F\rangle=\langle\gamma_*g,F\rangle\qquad\left(g\in C^k(n)\otimes\mathbb C,\ F\in C^k(\gamma n)\otimes\mathbb C\right).
\]
Moreover $\gamma^*$ commutes with $d$. Indeed, for positively oriented cells $\tau$ of dimension $k+1$ and $\sigma$ of dimension $k$ let $[\tau:\sigma]$ be the coefficient of $u(\sigma)$ in $(du)(\tau)$, and let $o(\tau)$, $o(\sigma)$ be their orientations, the exterior products of the unit vectors of their directions in increasing order. By the three formulas for $d$, $[\tau:\sigma]=0$ unless $\sigma$ is a face of $\tau$, and $N\wedge o(\sigma)=[\tau:\sigma]\cdot o(\tau)$ if $\sigma$ is a face of $\tau$ and $N$ is the outward unit normal of $\tau$ at $\sigma$. Let $\tau'$, $\sigma'$ be the positively oriented cells underlying $\gamma\tau$, $\gamma\sigma$, and let $\varepsilon_\tau,\varepsilon_\sigma\in\{1,-1\}$ be the signs with $\gamma\tau=\varepsilon_\tau\cdot\tau'$ and $\gamma\sigma=\varepsilon_\sigma\cdot\sigma'$. The linear part of $\gamma$ carries $o(\tau)$ and $o(\sigma)$ to $\varepsilon_\tau\cdot o(\tau')$ and $\varepsilon_\sigma\cdot o(\sigma')$; as $\gamma$ is an isometry mapping $\tau$ onto $\tau'$ and $\sigma$ onto $\sigma'$, it also carries the outward unit normal of $\tau$ at $\sigma$ to that of $\tau'$ at $\sigma'$. Applying the linear part of $\gamma$, extended to exterior products, to the identity $N\wedge o(\sigma)=[\tau:\sigma]\cdot o(\tau)$ one finds $[\tau':\sigma']=\varepsilon_\tau\cdot\varepsilon_\sigma\cdot[\tau:\sigma]$, and for $u\in C^k(\gamma n)$
\[
(\gamma^*du)(\tau)=\varepsilon_\tau\cdot\sum_\sigma[\tau':\sigma']\cdot u(\sigma')=\sum_\sigma[\tau:\sigma]\cdot\varepsilon_\sigma\cdot u(\sigma')=(d\gamma^*u)(\tau),
\]
the sums being taken over the faces $\sigma$ of $\tau$. Let $c\in\mathbb Z$, and let $\mathsf r$ be the lattice isometry of $\Lambda_n$ onto itself given by the reflection $x_0\mapsto c-x_0$, that is, $\mathsf rx=(c-x_0,x_1,x_2)$. It maps the plaquette $(x,12)$ onto $(\mathsf r x,12)$, and the plaquette $(x,0i)$ onto $(\mathsf r x-\hat0,0i)$ with the opposite orientation; hence
\[
(\mathsf r^*F)_{12}(x)=F_{12}(\mathsf r x),\qquad(\mathsf r^*F)_{0i}(x)=-F_{0i}(\mathsf r x-\hat0)\qquad(i=1,2).
\]
Thus $\mathsf r^*$ transports the component $F_{12}$ without a change of sign and reverses the sign of the components $F_{01}$ and $F_{02}$, those with an index $0$; this is the transformation rule of an antisymmetric tensor under the reflection of the coordinate $x_0$.

\subsubsection{The determinant on antiperiodic spinors}\label{sec:determinant-antiperiodic-spinors}

Let $\sigma_0,\sigma_1,\sigma_2$ be the three Pauli matrices, $\gamma_\mu:=\operatorname{diag}(\sigma_\mu,-\sigma_\mu)$ for $\mu=0,1,2$, and $P^\pm_\mu:=\frac12(1\pm\gamma_\mu)$. Let $\mathcal S_n$ be the space of the spinors that are antiperiodic in every direction, that is, of the functions $\psi:\mathbb Z^3\to\mathbb C^4$ with $\psi(x+n_\mu\hat\mu)=-\psi(x)$ for $\mu=0,1,2$; its dimension is $4|n|$. Let $B\in C^1(n)\otimes\mathbb C$ be a complex link field. For $x\in\mathbb Z^3$ and a direction $\mu$, $B_\mu(x)$ denotes the value of $B$ at the link $(\bar x,\mu)$ of $\Lambda_n$, where $\bar x\in\Lambda_n$ is the class of $x$; in this way $B$ becomes a function on all links of $\mathbb Z^3$, invariant under the translations by $\prod_\mu n_\mu\mathbb Z$. The Wilson operator $D[B]$ on $\mathcal S_n$ and the determinant $\Delta(B)$, taken on $\mathcal S_n$, are
\begin{multline}\label{eq:all-flux-1}
(D[B]\psi)(x)=3\psi(x)-\sum_\mu\left[P^-_\mu\cdot\exp\left(iB_\mu(x)\right)\cdot\psi(x+\hat\mu)+P^+_\mu\cdot\exp\left(-iB_\mu(x-\hat\mu)\right)\cdot\psi(x-\hat\mu)\right],\\
\Delta(B)=\det\left(D[B]+\boldsymbol{\mu}\right).
\end{multline}
When the torus has to be named, $\Delta_n$ is written for $\Delta$. Let $A\in C^1(n)$ be real. After a renumbering of the directions, two identifications hold: the Wilson operator $D_A$ of \eqref{eq:gaussian-bounds-1}, written in the four-component representation and taken at the antiperiodic character $\chi^{\mathrm{ap}}$, is the operator $D[qA]$ (on rectangular tori in the sense of Remark~\ref{rem:tori-unequal-sides}); and the operator of \eqref{eq:free-coefficient-5} is $D[B]$. The arguments proving Lemma~\ref{lem:gaussian-bounds-1} and Theorem~\ref{thm:free-coefficient-1}(d1)--(d4) rest on nothing but the formula for the operator and the algebraic relations among the matrices $\gamma_\mu$, and a renumbering of the directions affects neither; applied to $D[B]$, they yield the following three properties of $\Delta$.

First,
\begin{multline}\label{eq:all-flux-2}
\Delta\text{ is an entire function on }C^1(n)\otimes\mathbb C,\\
\sup\left\{|\Delta(B)|:\ B\in C^1(n)\otimes\mathbb C,\ \max_b\left|\operatorname{Im}B_b\right|\le c\right\}<\infty,\\
\Delta(B+2\pi\delta_b)=\Delta(B),\qquad\Delta(B+d\lambda)=\Delta(B),
\end{multline}
for every $c\ge0$, every $B\in C^1(n)\otimes\mathbb C$, every link $b$ and every $\lambda\in C^0(n)\otimes\mathbb C$; here $B_b$ is the value of $B$ at the link $b$, the maximum is taken over the links, and $\delta_b\in C^1(n)$ is the indicator of $b$ (Theorem~\ref{thm:free-coefficient-1}(d1)). The last identity holds because $\exp(-i\lambda)\cdot D[B]\cdot\exp(i\lambda)=D[B+d\lambda]$, the functions $\exp(\pm i\lambda)$ acting on $\mathcal S_n$ by multiplication.

Secondly, on real arguments
\begin{equation}\label{eq:all-flux-3}
\Delta(B)\ge0\quad\left(B\in C^1(n)\right),\qquad\sup_{B\in C^1(n)}\Delta(B)<\infty
\end{equation}
(the sign by Lemma~\ref{lem:gaussian-bounds-1}(v) and Theorem~\ref{thm:free-coefficient-1}(d2); the bound because, for real $B$, $D[B]+\boldsymbol{\mu}$ is a matrix of fixed finite size with bounded entries).

Thirdly, for every lattice isometry $\gamma:\Lambda_n\to\Lambda_{\gamma n}$,
\begin{equation}\label{eq:all-flux-4}
\Delta_n(\gamma^*B)=\Delta_{\gamma n}(B)\qquad\left(B\in C^1(\gamma n)\otimes\mathbb C\right)
\end{equation}
(Theorem~\ref{thm:free-coefficient-1}(d4)).

\subsubsection{The all-flux functional}\label{sec:all-flux-functional-2}

Let $\mathfrak F(n)$ be the set of the Borel functions $\Phi:C^2(n)\to\mathbb C$ for which there is a constant $c$ with $|\Phi(F)|\le c\cdot\exp(c\|F\|)$ for all $F\in C^2(n)$.

\begin{definition}[All-flux functional]\label{def:all-flux-functional}
Let $Z^2_{\mathrm{cl}}(n)$ be the set of the integer valued $m\in C^2(n)$ with $dm=0$. For $m\in Z^2_{\mathrm{cl}}(n)$ and $\mu<\nu$, let $\iota$ be the third direction and
\[
\mathfrak n_{\mu\nu}(m):=\sum_{x\in\Lambda_n:\ x_\iota=0}m_{\mu\nu}(x)\in\mathbb Z,
\]
the flux of $m$ through a coordinate two-torus in the directions $\mu,\nu$. Since $dm=0$, the flux does not change when the condition $x_\iota=0$ is replaced by $x_\iota=t$: in the sum of $dm$ over the cubes $x$ with $x_\iota=t$, which vanishes, the terms that do not contain $m_{\mu\nu}$ cancel, and the remaining terms give, up to sign, the difference of the sums of $m_{\mu\nu}$ at the heights $t+1$ and $t$. Put $\mathfrak n(m):=\left(\mathfrak n_{12}(m),\mathfrak n_{02}(m),\mathfrak n_{01}(m)\right)\in\mathbb Z^3$. Let $\mathrm dA$ be the Lebesgue measure of $C^1(n)$, and let $[0,r)^{\mathrm{links}}$ be the set of the $A\in C^1(n)$ with $0\le A_\mu(x)<r$ for every link $(x,\mu)$. For $\Phi\in\mathfrak F(n)$ put
\begin{multline}\label{eq:all-flux-5}
\omega^{\mathrm{af}}_n(\Phi):=\int_{[0,r)^{\mathrm{links}}}\mathrm dA\sum_{m\in Z^2_{\mathrm{cl}}(n)}\Phi(dA+rm)\cdot\exp\left(-\frac12\|dA+rm\|^2\right)\cdot\Delta(qA),\\
\omega^{\mathrm{af}}_n=\sum_{\mathfrak n\in\mathbb Z^3}\omega^{(\mathfrak n)}_n,
\end{multline}
where $\omega^{(\mathfrak n)}_n(\Phi)$, the sector of flux $\mathfrak n$, is the same expression with the sum restricted to the $m$ with $\mathfrak n(m)=\mathfrak n$. Put $\omega_n:=\omega^{\mathrm{af}}_n$.
\end{definition}

The two-cochain $F:=dA+rm$ is the curvature; $dF=0$, because $dd=0$ and $dm=0$. By \eqref{eq:all-flux-3}, $\omega_n$ is the integral with respect to a positive Borel measure on $C^2(n)$, the image of the measure $\exp\left(-\frac12\|dA+rm\|^2\right)\cdot\Delta(qA)\,\mathrm dA$ on $[0,r)^{\mathrm{links}}\times Z^2_{\mathrm{cl}}(n)$ under the map $(A,m)\mapsto dA+rm$. This measure is finite on $\mathfrak F(n)$, that is, $\omega_n(|\Phi|)<\infty$ for every $\Phi\in\mathfrak F(n)$: the function $\Delta$ is bounded on real arguments by \eqref{eq:all-flux-3}; for $A\in[0,r)^{\mathrm{links}}$ the component of $F$ at every plaquette $p$ satisfies $|F_p|\ge r|m_p|-4r$; and $c\|F\|-\frac12\|F\|^2\le c^2-\frac14\|F\|^2$, so that the Gaussian weight makes the series over $m$ of $\sup_A c\cdot\exp(c\|F\|)\cdot\exp\left(-\frac12\|F\|^2\right)$ converge.

Moreover $\omega_n(1)>0$: already the summand $m=0$ of \eqref{eq:all-flux-5} with $\Phi=1$ is positive, since $\Delta\ge0$ is continuous and $\Delta(0)>0$. For the last inequality, let $p$ be an antiperiodic momentum, that is, $p_\mu\in\frac{2\pi}{n_\mu}\left(\mathbb Z+\frac12\right)$ for every $\mu$, and let $u\in\mathbb C^4$. The operator $D[0]+\boldsymbol{\mu}$ maps the function $x\mapsto\exp(ip\cdot x)\cdot u$ to $x\mapsto\exp(ip\cdot x)\cdot\widehat D(p)u$, where
\[
\widehat D(p):=\boldsymbol{\mu}+\sum_\mu\left(1-\cos p_\mu\right)+i\gamma\cdot\sin p,\qquad\gamma\cdot\sin p:=\sum_\mu\gamma_\mu\sin p_\mu,
\]
is the free symbol. These functions span $\mathcal S_n$, so $\Delta(0)$ is the product of the numbers $\det\widehat D(p)$ over the antiperiodic momenta modulo $2\pi$. One has
\[
\det\widehat D(p)=\left(\left(\boldsymbol{\mu}+\sum_\mu\left(1-\cos p_\mu\right)\right)^2+\sum_\mu\sin^2p_\mu\right)^2,
\]
which vanishes only if $\sin p_\mu=0$ for every $\mu$, which forces $p=(\pi,\pi,\pi)$, and if moreover $\boldsymbol{\mu}+6=0$. As $\boldsymbol{\mu}>-1/2$, the number $\boldsymbol{\mu}+6$ is positive; hence every factor of the product is positive, and $\Delta(0)>0$.

For $g\in C^2(n)\otimes\mathbb C$ the function $F\mapsto\exp\langle g,F\rangle$ lies in $\mathfrak F(n)$, because $|\exp\langle g,F\rangle|\le\exp\left(\|\operatorname{Re}g\|\cdot\|F\|\right)$. Put
\[
Z_n(g):=\omega_n\left(\exp\langle g,F\rangle\right)\quad\left(g\in C^2(n)\otimes\mathbb C\right),\qquad\langle\Phi\rangle_n:=\frac{\omega_n(\Phi)}{\omega_n(1)}\quad\left(\Phi\in\mathfrak F(n)\right).
\]

Let $\gamma:\Lambda_n\to\Lambda_{\gamma n}$ be a lattice isometry and $\Phi\in\mathfrak F(n)$. Then $\Phi\circ\gamma^*\in\mathfrak F(\gamma n)$, as $\gamma^*$ is orthogonal, and
\[
\omega_{\gamma n}(\Phi\circ\gamma^*)=\omega_n(\Phi),\qquad\omega^{(0)}_{\gamma n}(\Phi\circ\gamma^*)=\omega^{(0)}_n(\Phi)
\]
(Theorem~\ref{thm:free-coefficient-1}(d5)). For the proof, let $\mathbb Z^{\mathrm{links}}(n')\subset C^1(n')$, for a triple $n'$ of integers $\ge2$, be the group of the integer valued one-cochains, and for $A\in C^1(n)$ put
\[
\Xi_{n,\Phi}(A):=\sum_{m\in Z^2_{\mathrm{cl}}(n)}\Phi(dA+rm)\cdot\exp\left(-\frac12\|dA+rm\|^2\right)\cdot\Delta(qA);
\]
let $\Xi^{(0)}_{n,\Phi}(A)$ be the same sum restricted to the $m$ of flux $0$. Let $c\in\mathbb Z^{\mathrm{links}}(n)$. The cochain $dc$ lies in $Z^2_{\mathrm{cl}}(n)$ and has flux $0$. Indeed, for $\mu<\nu$ with third direction $\iota$,
\[
(dc)_{\mu\nu}(x)=\left[c_\mu(x)-c_\mu(x+\hat\nu)\right]+\left[c_\nu(x+\hat\mu)-c_\nu(x)\right],
\]
and the sum of each bracket over the sites $x$ with $x_\iota=t$ vanishes, since this set of sites is mapped onto itself by the translations by $\hat\mu$ and by $\hat\nu$. Hence $m\mapsto m+dc$ is a bijection of $Z^2_{\mathrm{cl}}(n)$ which preserves the flux. As $d(A+rc)+rm=dA+r(m+dc)$, and $\Delta(q(A+rc))=\Delta(qA+2\pi c)=\Delta(qA)$ by \eqref{eq:all-flux-2}, the series defining $\Xi_{n,\Phi}(A+rc)$ is a rearrangement of the series defining $\Xi_{n,\Phi}(A)$, and the same holds for $\Xi^{(0)}_{n,\Phi}$. Every $A\in C^1(n)$ is of the form $A_0+rc$ with $A_0\in[0,r)^{\mathrm{links}}$ and $c\in\mathbb Z^{\mathrm{links}}(n)$, and for $A_0\in[0,r)^{\mathrm{links}}$ both series converge absolutely, by the bound that gave the finiteness of $\omega_n$ on $\mathfrak F(n)$. Consequently both series converge absolutely for every $A\in C^1(n)$, and
\[
\Xi_{n,\Phi}(A+rc)=\Xi_{n,\Phi}(A),\qquad\Xi^{(0)}_{n,\Phi}(A+rc)=\Xi^{(0)}_{n,\Phi}(A)\qquad\left(A\in C^1(n),\ c\in\mathbb Z^{\mathrm{links}}(n)\right).
\]
By \eqref{eq:all-flux-5}, $\omega_n(\Phi)$ and $\omega^{(0)}_n(\Phi)$ are the integrals of $\Xi_{n,\Phi}$ and $\Xi^{(0)}_{n,\Phi}$ over $[0,r)^{\mathrm{links}}$, a fundamental domain of the lattice $r\mathbb Z^{\mathrm{links}}(n)$. Since these two functions are invariant under the translations by this lattice and integrable over $[0,r)^{\mathrm{links}}$ (as $\omega_n(|\Phi|)<\infty$), their integrals over every measurable fundamental domain of $r\mathbb Z^{\mathrm{links}}(n)$ take the same values.

As $\gamma$ maps the links of $\Lambda_n$ bijectively onto those of $\Lambda_{\gamma n}$, up to orientation, the matrix of the linear map $\gamma_*:C^1(n)\to C^1(\gamma n)$ in the bases formed by the indicators of the positively oriented links is a signed permutation matrix. Hence $\gamma_*$ preserves the Lebesgue measure and maps $r\mathbb Z^{\mathrm{links}}(n)$ onto $r\mathbb Z^{\mathrm{links}}(\gamma n)$, and $\gamma_*\left([0,r)^{\mathrm{links}}\right)$ is a measurable fundamental domain of $r\mathbb Z^{\mathrm{links}}(\gamma n)$. In every degree, $\gamma_*$ maps the integer valued cochains onto the integer valued cochains and commutes with $d$, since it is the pullback by the lattice isometry $\gamma^{-1}:x\mapsto\rho^{-1}(x-a)$ of $\Lambda_{\gamma n}$ onto $\Lambda_n$. Therefore $\gamma_*$ maps $Z^2_{\mathrm{cl}}(n)$ onto $Z^2_{\mathrm{cl}}(\gamma n)$, and the set of the cochains $dc$ with $c\in\mathbb Z^{\mathrm{links}}(n)$ onto the set of the cochains $dc'$ with $c'\in\mathbb Z^{\mathrm{links}}(\gamma n)$; by \eqref{eq:partition-functions-2}, the first of these sets is the set of the elements of flux $0$ of $Z^2_{\mathrm{cl}}(n)$, and the second is that of $Z^2_{\mathrm{cl}}(\gamma n)$. For $A\in C^1(n)$ and $m\in Z^2_{\mathrm{cl}}(n)$ one has $d\gamma_*A+r\gamma_*m=\gamma_*(dA+rm)$, hence
\[
\gamma^*\left(d\gamma_*A+r\gamma_*m\right)=dA+rm,\qquad\left\|d\gamma_*A+r\gamma_*m\right\|=\|dA+rm\|,
\]
and $\Delta_{\gamma n}(q\gamma_*A)=\Delta_n(qA)$ by \eqref{eq:all-flux-4} with $B=q\gamma_*A$. The substitution $m'=\gamma_*m$ in the series that define $\Xi_{\gamma n,\Phi\circ\gamma^*}(\gamma_*A)$ and $\Xi^{(0)}_{\gamma n,\Phi\circ\gamma^*}(\gamma_*A)$ gives
\[
\Xi_{\gamma n,\Phi\circ\gamma^*}(\gamma_*A)=\Xi_{n,\Phi}(A),\qquad\Xi^{(0)}_{\gamma n,\Phi\circ\gamma^*}(\gamma_*A)=\Xi^{(0)}_{n,\Phi}(A)\qquad\left(A\in C^1(n)\right).
\]
Integrating over $A\in[0,r)^{\mathrm{links}}$ and substituting $A'=\gamma_*A$, one finds that $\omega_n(\Phi)$ and $\omega^{(0)}_n(\Phi)$ are the integrals of $\Xi_{\gamma n,\Phi\circ\gamma^*}$ and $\Xi^{(0)}_{\gamma n,\Phi\circ\gamma^*}$ over the fundamental domain $\gamma_*\left([0,r)^{\mathrm{links}}\right)$ of $r\mathbb Z^{\mathrm{links}}(\gamma n)$. Since $\Phi\circ\gamma^*\in\mathfrak F(\gamma n)$, what was shown above for $n$ and $\Phi$ holds for $\gamma n$ and $\Phi\circ\gamma^*$: these two functions are invariant under the translations by $r\mathbb Z^{\mathrm{links}}(\gamma n)$ and integrable over its fundamental domains, so that $\gamma_*\left([0,r)^{\mathrm{links}}\right)$ may be replaced by the box $[0,r)^{\mathrm{links}}$ of $\gamma n$. Hence the two integrals are $\omega_{\gamma n}(\Phi\circ\gamma^*)$ and $\omega^{(0)}_{\gamma n}(\Phi\circ\gamma^*)$.

The single sectors are permuted: by the clause of Theorem~\ref{thm:free-coefficient-1}(d5) on the single sectors,
\[
\omega^{(\gamma\cdot\mathfrak n)}_{\gamma n}(\Phi\circ\gamma^*)=\omega^{(\mathfrak n)}_n(\Phi)\qquad\left(\mathfrak n\in\mathbb Z^3,\ \Phi\in\mathfrak F(n)\right),
\]
where $\gamma\cdot\mathfrak n\in\mathbb Z^3$ is defined as follows: if $\gamma$ maps the plaquettes in the directions $\mu<\nu$ onto the plaquettes in the directions $\mu'<\nu'$ and multiplies their orientations by $s\in\{1,-1\}$, then $(\gamma\cdot\mathfrak n)_{\mu'\nu'}:=s\cdot\mathfrak n_{\mu\nu}$. The numbers
\[
q_{\mathfrak n}(n)=\exp\left(-\frac{r^2}{2|n|}\sum_{\mu<\nu}n_\iota^2\,\mathfrak n_{\mu\nu}^2\right)\qquad(\iota\text{ the third direction of }\mu,\nu)
\]
of Section~\ref{sec:bound-single-flux} satisfy $q_{\gamma\cdot\mathfrak n}(\gamma n)=q_{\mathfrak n}(n)$, since the sides of $\gamma n$ are those of $n$, permuted as the directions. For the exchange $\gamma$ of the axes $0$ and $1$, which maps $\hat1\wedge\hat2$, $\hat0\wedge\hat2$, $\hat0\wedge\hat1$ to $\hat0\wedge\hat2$, $\hat1\wedge\hat2$, $-\hat0\wedge\hat1$,
\[
\gamma\cdot(\mathfrak n_{12},\mathfrak n_{02},\mathfrak n_{01})=(\mathfrak n_{02},\mathfrak n_{12},-\mathfrak n_{01});
\]
for the exchange $\gamma$ of the axes $0$ and $2$, which maps them to $-\hat0\wedge\hat1$, $-\hat0\wedge\hat2$, $-\hat1\wedge\hat2$, one has $\gamma\cdot(\mathfrak n_{12},\mathfrak n_{02},\mathfrak n_{01})=(-\mathfrak n_{01},-\mathfrak n_{02},-\mathfrak n_{12})$.

\begin{lemma}[The sector of zero flux]\label{lem:sector-zero-flux}
Let $V_{\mathrm{g}}:=dC^0(n)$, $V_{\mathrm{h}}:=\ker d\cap{V_{\mathrm{g}}}^\perp$ and $V_{\mathrm{t}}:=(V_{\mathrm{g}}\oplus V_{\mathrm{h}})^\perp$, subspaces of $C^1(n)$, where $d$ acts on $C^1(n)$ in $\ker d$. Let $\sigma_{V_{\mathrm{t}}}$ be the Lebesgue measure of the Euclidean space $V_{\mathrm{t}}$. For $\theta\in\mathbb R^3$ let $H_\theta\in C^1(n)$ be the constant cochain with $(H_\theta)_\mu(x)=\theta_\mu/(qn_\mu)$. Let $d^*:C^1(n)\to C^0(n)$ be the adjoint of $d:C^0(n)\to C^1(n)$ with respect to the counting pairings, let $\Delta_0:=d^*d$, an operator on $C^0(n)$, and let $\det{}'\Delta_0$ be the product of its nonzero eigenvalues. Then for every $\Phi\in\mathfrak F(n)$
\begin{align}
\omega^{(0)}_n(\Phi)&=\mathcal N_n\cdot\overline Z_n(\Phi),\qquad\mathcal N_n:=r^{|n|+2}\cdot\left(\det{}'\Delta_0\right)^{\frac12},\label{eq:all-flux-6}\\
\overline Z_n(\Phi)&:=\frac1{(2\pi)^3}\int_{[-\pi,\pi]^3}\mathrm d\theta\int_{V_{\mathrm{t}}}\mathrm d\sigma_{V_{\mathrm{t}}}(A)\,\Phi(dA)\cdot\exp\left(-\frac12\|dA\|^2\right)\cdot\Delta(qA+qH_\theta).\notag
\end{align}
\end{lemma}

The functional $\overline Z_n$ is the functional $\overline Z$ of \eqref{eq:gaussian-bounds-3} with $c^\theta=H_\theta$ (on rectangular tori as in Remark~\ref{rem:tori-unequal-sides}). As the summand $m=0$ of $\omega^{(0)}_n(1)$ is positive, $\overline Z_n(1)>0$; the normalised functional $\Phi\mapsto\overline Z_n(\Phi)/\overline Z_n(1)$ is called the \emph{harmonic-averaged law}.

\begin{proof}
Let $\Phi\in\mathfrak F(n)$, and put $\Psi(F):=\Phi(F)\cdot\exp\left(-\frac12\|F\|^2\right)$ for $F\in C^2(n)$. Write $\mathbb Z^{\mathrm{links}}:=\mathbb Z^{\mathrm{links}}(n)$, let $\mathbb Z^{\mathrm{sites}}$ be the group of the integer valued zero-cochains, and let $Z^1_{\mathrm{cl}}(n):=\{c\in\mathbb Z^{\mathrm{links}}:dc=0\}$ be the group of the closed integer one-cochains. By \eqref{eq:partition-functions-2}, which expresses that $H^2(\mathbb T^3;\mathbb Z)\cong\mathbb Z^3$ is detected by the three fluxes, the elements of $Z^2_{\mathrm{cl}}(n)$ of flux $0$ are the cochains $dc$ with $c\in\mathbb Z^{\mathrm{links}}$. The map $c\mapsto dc$ has the kernel $Z^1_{\mathrm{cl}}(n)$, and $\Delta(q(A+rc))=\Delta(qA+2\pi c)=\Delta(qA)$ for $c\in\mathbb Z^{\mathrm{links}}$ by \eqref{eq:all-flux-2}. Let $\mathcal R\subset\mathbb Z^{\mathrm{links}}$ contain exactly one element of each coset of $Z^1_{\mathrm{cl}}(n)$. The sets $[0,r)^{\mathrm{links}}+rc$, $c\in\mathcal R$, are pairwise disjoint, and their union is a fundamental domain for the translations of $C^1(n)$ by $rZ^1_{\mathrm{cl}}(n)$, because $[0,r)^{\mathrm{links}}$ is a fundamental domain for the translations by $r\mathbb Z^{\mathrm{links}}$; the function $A\mapsto\Psi(dA)\cdot\Delta(qA)$ is invariant under the translations by $rZ^1_{\mathrm{cl}}(n)$. Since $\omega_n(|\Phi|)<\infty$, sum and integral may be interchanged, and
\[
\omega^{(0)}_n(\Phi)=\sum_{c\in\mathcal R}\int_{[0,r)^{\mathrm{links}}}\mathrm dA\,\Psi\left(d(A+rc)\right)\cdot\Delta\left(q(A+rc)\right)=\int_{C^1(n)/rZ^1_{\mathrm{cl}}(n)}\mathrm dA\,\Psi(dA)\cdot\Delta(qA).
\]

The group $Z^1_{\mathrm{cl}}(n)$ is a lattice of full rank in $\ker d=V_{\mathrm{g}}\oplus V_{\mathrm{h}}$, since $\ker d$ is the kernel of a matrix with integer entries. Decompose $A=A_{V_{\mathrm{g}}}+A_{V_{\mathrm{h}}}+A_{V_{\mathrm{t}}}$ along $C^1(n)=V_{\mathrm{g}}\oplus V_{\mathrm{h}}\oplus V_{\mathrm{t}}$. Then $dA=dA_{V_{\mathrm{t}}}$, and $\Delta(qA)$ does not depend on $A_{V_{\mathrm{g}}}$ by \eqref{eq:all-flux-2}; so the integrand depends on $A_{V_{\mathrm{h}}}$ and $A_{V_{\mathrm{t}}}$ only. Let $P_{V_{\mathrm{h}}}$ be the orthogonal projection onto $V_{\mathrm{h}}$, let $\mathrm dh$ be the Lebesgue measure of $V_{\mathrm{h}}$, and for a lattice $L$ let $\operatorname{covol}(L)$ be the covolume of $L$ in its linear span. The torus $\ker d/rZ^1_{\mathrm{cl}}(n)$ is fibred over the torus $V_{\mathrm{h}}/P_{V_{\mathrm{h}}}\left(rZ^1_{\mathrm{cl}}(n)\right)$, each fibre being a copy of $V_{\mathrm{g}}/\left(rZ^1_{\mathrm{cl}}(n)\cap V_{\mathrm{g}}\right)$; integrating first along the fibres and then over the base (Fubini's theorem), one obtains
\[
\omega^{(0)}_n(\Phi)=\operatorname{covol}\left(rZ^1_{\mathrm{cl}}(n)\cap V_{\mathrm{g}}\right)\cdot\int_{V_{\mathrm{h}}/P_{V_{\mathrm{h}}}\left(rZ^1_{\mathrm{cl}}(n)\right)}\mathrm dh\int_{V_{\mathrm{t}}}\mathrm d\sigma_{V_{\mathrm{t}}}(A)\,\Psi(dA)\cdot\Delta(qA+qh);
\]
here the lattice $rZ^1_{\mathrm{cl}}(n)\cap V_{\mathrm{g}}$ spans $V_{\mathrm{g}}$, as is shown next.

If $\lambda\in C^0(n)$ and $d\lambda$ is integer valued, then $\lambda(x+\hat\mu)-\lambda(x)\in\mathbb Z$ for all $x$ and $\mu$, so that $\lambda$ is the sum of a constant and an element of $\mathbb Z^{\mathrm{sites}}$, the graph $\Lambda_n$ being connected. Hence $Z^1_{\mathrm{cl}}(n)\cap V_{\mathrm{g}}=d\mathbb Z^{\mathrm{sites}}$. Fix a site $x_0$, and let $\delta_x\in C^0(n)$ be the indicator of the site $x$. The kernel of $d$ on $C^0(n)$ consists of the constants; therefore the cochains $d\delta_x$, $x\ne x_0$, form a basis of $d\mathbb Z^{\mathrm{sites}}$, which thus spans $V_{\mathrm{g}}$, a space of dimension $|n|-1$. Since $\langle d\delta_x,d\delta_y\rangle=\langle\delta_x,\Delta_0\delta_y\rangle$, the Gram matrix $\left(\langle d\delta_x,d\delta_y\rangle\right)_{x,y\ne x_0}$ is the reduced Laplacian, the matrix of $\Delta_0$ in the basis $(\delta_x)_{x\in\Lambda_n}$ with the row and the column of $x_0$ removed, and by Kirchhoff's matrix-tree theorem its determinant is $\det{}'\Delta_0/|n|$. Consequently
\[
\operatorname{covol}\left(rZ^1_{\mathrm{cl}}(n)\cap V_{\mathrm{g}}\right)=r^{|n|-1}\cdot\left(\frac{\det{}'\Delta_0}{|n|}\right)^{\frac12}.
\]

For $\mu=0,1,2$ let $c^{(\mu)}\in\mathbb Z^{\mathrm{links}}$ be the indicator of the links $(x,\mu)$ with $x_\mu=0$. It is closed, since on a plaquette that contains the direction $\mu$ its two $\mu$-links have the same coordinate $x_\mu$ and enter $d$ with opposite signs, and its sum along a coordinate loop in the direction $\nu$ is $\delta_{\mu\nu}$. A closed integer one-cochain whose sums along the three coordinate loops vanish is $d\lambda$ with $\lambda\in\mathbb Z^{\mathrm{sites}}$, because $H^1(\mathbb T^3;\mathbb Z)\cong\mathbb Z^3$ is detected by these sums. Therefore $Z^1_{\mathrm{cl}}(n)$ is generated by $d\mathbb Z^{\mathrm{sites}}$ and $c^{(0)},c^{(1)},c^{(2)}$. By the Hodge decomposition on the discrete torus, $V_{\mathrm{h}}$ consists of the cochains that are constant on the links of each direction. Hence $P_{V_{\mathrm{h}}}$ vanishes on $d\mathbb Z^{\mathrm{sites}}\subset V_{\mathrm{g}}$, and $P_{V_{\mathrm{h}}}c^{(\mu)}$ is the cochain equal to $1/n_\mu$ on the $\mu$-links and to $0$ on the other links; its norm is $|n|^{\frac12}/n_\mu$, and the three cochains $P_{V_{\mathrm{h}}}c^{(\mu)}$ are orthogonal. Thus $P_{V_{\mathrm{h}}}\left(rZ^1_{\mathrm{cl}}(n)\right)$ is the lattice with the orthogonal basis $rP_{V_{\mathrm{h}}}c^{(\mu)}$, $\mu=0,1,2$, and the quotient $V_{\mathrm{h}}/P_{V_{\mathrm{h}}}\left(rZ^1_{\mathrm{cl}}(n)\right)$ is a torus of volume $r^3|n|^{\frac32}/|n|$. Since $H_\theta=\sum_\mu\frac{\theta_\mu}{2\pi}\cdot rP_{V_{\mathrm{h}}}c^{(\mu)}$, the map $\theta\mapsto H_\theta$ is a bijection of $[-\pi,\pi)^3$ onto a fundamental domain of this lattice, and it carries the measure $\mathrm d\theta/(2\pi)^3$ to the Lebesgue measure of $V_{\mathrm{h}}$ divided by that volume. So
\[
\int_{V_{\mathrm{h}}/P_{V_{\mathrm{h}}}\left(rZ^1_{\mathrm{cl}}(n)\right)}\mathrm dh\,f(h)=r^3\cdot|n|^{\frac12}\cdot\frac1{(2\pi)^3}\int_{[-\pi,\pi]^3}\mathrm d\theta\,f(H_\theta)
\]
for every integrable function $f$ on this torus. Multiplying the two factors,
\[
r^{|n|-1}\cdot\left(\frac{\det{}'\Delta_0}{|n|}\right)^{\frac12}\cdot r^3\cdot|n|^{\frac12}=r^{|n|+2}\cdot\left(\det{}'\Delta_0\right)^{\frac12}=\mathcal N_n,
\]
which proves \eqref{eq:all-flux-6}.
\end{proof}

Let $\varepsilon_0\in\mathbb R$. The function $\Delta$ of \eqref{eq:all-flux-1} may be replaced by its product $\exp(-\varepsilon_0|n|)\cdot\Delta$ with a constant vacuum factor. The factor then multiplies $\omega_n$, every sector $\omega^{(\mathfrak n)}_n$ and $\overline Z_n$, and it enters the function $s$ of the slice data of Section~\ref{sec:slice-factorisation-determinant} as $\exp\left(-\frac12\varepsilon_0n_1n_2\right)$. The logarithm $-\varepsilon_0|n|$ of the vacuum factor is exactly proportional to the volume $|n|$ of the torus. Since the factor is a positive constant with the same value on $\Lambda_n$ and on $\Lambda_{\gamma n}$ (as $|\gamma n|=|n|$), the replacement preserves the three properties of $\Delta$ and the inequality $\Delta(0)>0$, and with them the finiteness of $\omega_n$ on $\mathfrak F(n)$, the inequality $\omega_n(1)>0$ and the statements on lattice isometries. Both sides of \eqref{eq:all-flux-6} acquire the factor, so that Lemma~\ref{lem:sector-zero-flux} holds with the same constant $\mathcal N_n$, and no statement that follows is changed.

\subsection{A trace formula}\label{sec:trace-formula}

In this section the direction $0$ is distinguished and called time. All constructions of the section can be made with any other direction $\mu$ in place of $0$, and all statements remain true. Indeed, let $\gamma\colon\Lambda_n\to\Lambda_{\gamma n}$ be the exchange of the axes $0$ and $\mu$. By the invariance under lattice isometries and the covariance $\omega^{(\gamma\cdot\mathfrak{n})}_{\gamma n}(\Phi\circ\gamma^{*})=\omega^{(\mathfrak{n})}_n(\Phi)$ of the single flux sectors (Section~\ref{sec:all-flux-functional-2}), a statement with time direction $\mu$ on $\Lambda_n$ is the statement with time direction $0$ on $\Lambda_{\gamma n}$, transported by $\gamma$, with the fluxes relabelled by $\mathfrak{n}\mapsto\gamma\cdot\mathfrak{n}$. Accordingly, the Fourier variables in \eqref{eq:all-flux-11} are, up to sign, the two fluxes through the coordinate planes that contain the time direction, and the summands of the decomposition in Proposition~\ref{prop:trace-formula-bloch}(iii) are labelled by the third flux, again up to sign.

Let $n=(n_0,n_1,n_2)$, and let $\Sigma=\mathbb{Z}/n_1\mathbb{Z}\times\mathbb{Z}/n_2\mathbb{Z}$ be the cross-section of $\Lambda_n$; every slice $\{x\in\Lambda_n:x_0=j\}$, $j\in\mathbb{Z}/n_0\mathbb{Z}$, is identified with $\Sigma$. The torus $\Sigma$ has $v=n_1n_2$ sites, $l=2v$ links and $v$ plaquettes. Real functions on its sites and on its links are identified with the elements of $\mathbb{R}^v$ and of $\mathbb{R}^l$, and real and integer functions on its plaquettes with those of $\mathbb{R}^{\mathrm{plaq}(\Sigma)}$ and of $\mathbb{Z}^{\mathrm{plaq}(\Sigma)}$. The spaces $\mathbb{R}^l$ and $\mathbb{R}^{\mathrm{plaq}(\Sigma)}$ carry the Euclidean norm $\|\cdot\|$: the square of the norm of an element is the sum of the squares of its values on the links, respectively on the plaquettes, as for the cochains of $\Lambda_n$ in Section~\ref{sec:tori-cochains-lattice}. The coboundary $d_s$ of $\Sigma$ is given by the formulas of Section~\ref{sec:tori-cochains-lattice} in the directions $1$ and $2$. A one-cochain $A\in C^1(n)$ is written $(A_j,a_j)_{j\in\mathbb{Z}/n_0\mathbb{Z}}$, where $A_j\in\mathbb{R}^l$ consists of the values of $A$ on the links of directions $1$ and $2$ in the slice $j$, and $a_j\in\mathbb{R}^v$ of its values on the links of direction $0$ from the slice $j$ to the slice $j+1$.

\subsubsection{Slice factorisation of the determinant}\label{sec:slice-factorisation-determinant}

\begin{definition}[Slice data]
Let $\mathcal{F}$ be a Hilbert space of finite dimension, and let $Q^{\mathrm{ch}}_x$, one for each site $x$ of $\Sigma$, be self-adjoint operators on $\mathcal{F}$ which commute with one another and whose spectra are contained in $\mathbb{Z}$. Let $s\colon\mathbb{R}^l\to(0,\infty)$ be a function and $Q$ a map from $\mathbb{R}^l$ into the linear operators on $\mathcal{F}$; the argument $A\in\mathbb{R}^l$ of $s$ and of $Q$ is read as a real field on the links of $\Sigma$. These four objects are attached to $\Sigma$, $q$ and $\boldsymbol{\mu}$, and they do not change with $n_0$. Put
\[
\Gamma(a)=\exp\left(iq\sum_{x\in\Sigma}a_x\cdot Q^{\mathrm{ch}}_x\right)\qquad\left(a\in\mathbb{R}^v\right).
\]
The quadruple $\left(\mathcal{F},s,Q,(Q^{\mathrm{ch}}_x)_{x\in\Sigma}\right)$ is called slice data for $\Sigma$, $q$ and $\boldsymbol{\mu}$ if
\begin{gather}
\text{$s$ and $Q$ are continuous and bounded on $\mathbb{R}^l$;}\tag{SD1}\label{eq:all-flux-SD1}\\
\text{$s(A)=s(A')$ and $Q(A)=Q(A')$ for all $A,A'\in\mathbb{R}^l$ with}\notag\\
\exp(iqA_b)=\exp(iqA'_b)\ \text{for every link $b$ of $\Sigma$;}\notag\\
\Gamma(a)\ \text{is unitary},\quad\Gamma(a)\cdot\Gamma(a')=\Gamma(a+a'),\notag\\
\Gamma(a)^{*}=\Gamma(-a),\quad\Gamma(rk)=1\qquad(a,a'\in\mathbb{R}^v,\ k\in\mathbb{Z}^v);\notag
\end{gather}
and
\begin{gather}
\text{for all}\ n_0\ge2\ \text{and all}\ A=(A_j,a_j)_{j\in\mathbb{Z}/n_0\mathbb{Z}}\in C^1(n):\tag{SD2}\label{eq:all-flux-SD2}\\
\Delta(qA)=\operatorname{Tr}_{\mathcal{F}}\prod_{j=n_0-1}^0\left[s(A_{j+1})\cdot Q(A_{j+1})^{*}\cdot\Gamma(a_j)\cdot Q(A_j)\cdot s(A_j)\right].\label{eq:all-flux-7}
\end{gather}
Here $\Delta$ is the function \eqref{eq:all-flux-1} on $\Lambda_n$, $n=(n_0,n_1,n_2)$; the index $j$ is taken modulo $n_0$, so that $A_{n_0}=A_0$; and the factors of the product are ordered with later times to the left, the factor $j=n_0-1$ standing leftmost and the factor $j=0$ rightmost.
\end{definition}

The four conditions on $\Gamma$ in \eqref{eq:all-flux-SD1} are satisfied automatically. Indeed, the operators $Q^{\mathrm{ch}}_x$, being self-adjoint and pairwise commuting, possess a common orthonormal eigenbasis of $\mathcal{F}$, and their eigenvalues are integers; in such a basis $\Gamma(a)$ is diagonal, each diagonal entry being $\exp\left(iq\sum_{x\in\Sigma}a_x\cdot z_x\right)$, where $z_x\in\mathbb{Z}$ is the eigenvalue of $Q^{\mathrm{ch}}_x$ at the corresponding basis vector. As a function of $a$, such an entry is a unitary character of $\mathbb{R}^v$, and it is trivial on $r\mathbb{Z}^v$ because $qr=2\pi$. Lemma~\ref{lem:slice-factorisation-determinant-2} provides slice data for $\Sigma$, $q$ and $\boldsymbol{\mu}$ whenever $q>0$ and $\boldsymbol{\mu}>-1$, whatever $n_1$ and $n_2$ are. From now on slice data are fixed; every statement below holds for each choice of them, and they enter only through \eqref{eq:all-flux-SD1} and \eqref{eq:all-flux-SD2}.

\subsubsection{Transfer kernels}\label{sec:transfer-kernels}

Points of $\mathbb{R}^l\times\mathbb{Z}^{\mathrm{plaq}(\Sigma)}$ are written $x=(A,m)$ and $x'=(A',m')$. The group $\mathbb{Z}^l$ acts on this set by $(A,m)\mapsto(A+rc,\,m-d_sc)$, $c\in\mathbb{Z}^l$, and $\sim$ denotes the equivalence relation whose classes are the orbits:
\[
(A,m)\sim(A+rc,\,m-d_sc),\qquad c\in\mathbb{Z}^l.
\]
The slice space
\[
\mathcal{X}=[0,r)^l\times\mathbb{Z}^{\mathrm{plaq}(\Sigma)}
\]
carries the product of Lebesgue measure and counting measure, written $\mathrm{d}x$; it is a fundamental domain of $\sim$, that is, it meets every class in exactly one point. The temporal variables are $a\in[0,r)^v$, with the Lebesgue measure $\mathfrak{m}$ (so $\mathrm{d}\mathfrak{m}(a)=\mathrm{d}a$), and $c\in\mathbb{Z}^l$. For $x,x'\in\mathbb{R}^l\times\mathbb{Z}^{\mathrm{plaq}(\Sigma)}$, $a\in\mathbb{R}^v$ and $c\in\mathbb{Z}^l$, the magnetic curvature $\beta(x)\in\mathbb{R}^{\mathrm{plaq}(\Sigma)}$, the electric curvature $\mathcal{E}(x',x;a,c)\in\mathbb{R}^l$ and the constraint $\chi(x',x;c)$ are
\begin{align*}
\beta(x)&=d_sA+rm,\\
\mathcal{E}(x',x;a,c)&=A'-A-d_sa+rc,\\
\chi(x',x;c)&=\begin{cases}1&\text{if}\ m'=m+d_sc,\\0&\text{otherwise}.\end{cases}
\end{align*}

Let $\mathrm{row}_1=\{(x;1):x\in\Sigma,\ x_2=0\}$ and $\mathrm{row}_2=\{(x;2):x\in\Sigma,\ x_1=0\}$; thus $\mathrm{row}_1$ consists of the links of direction $1$ with $x_2=0$, and $\mathrm{row}_2$ of the links of direction $2$ with $x_1=0$. For a Bloch angle $\vartheta\in\mathbb{R}^2$ let $\mathcal{H}_{\vartheta}$ be the space of measurable maps $\Psi\colon\mathbb{R}^l\times\mathbb{Z}^{\mathrm{plaq}(\Sigma)}\to\mathcal{F}$ such that
\begin{multline}\label{eq:all-flux-8}
\Psi(A+rc,\,m-d_sc)=\exp(-2\pi i\,\vartheta\cdot w(c))\cdot\Psi(A,m),\\
w_i(c)=\sum_{b\in\mathrm{row}_i}c_b\quad(i=1,2),\qquad\|\Psi\|^2=\int_{\mathcal{X}}\mathrm{d}x\,|\Psi(x)|^2<\infty,
\end{multline}
for all $(A,m)\in\mathbb{R}^l\times\mathbb{Z}^{\mathrm{plaq}(\Sigma)}$ and all $c\in\mathbb{Z}^l$; here $\vartheta\cdot w(c)=\vartheta_1w_1(c)+\vartheta_2w_2(c)$, $|\cdot|$ is the norm of $\mathcal{F}$, and maps that agree almost everywhere are identified. With the norm of \eqref{eq:all-flux-8}, $\mathcal{H}_{\vartheta}$ is a Hilbert space. For $\vartheta=0$ the first identity of \eqref{eq:all-flux-8} says that $\Psi$ is constant on the classes of $\sim$, and restriction to $\mathcal{X}$ identifies $\mathcal{H}:=\mathcal{H}_0$ with $L^2(\mathcal{X})\otimes\mathcal{F}$.

For $x=(A,m)$ and $x'$ in $\mathbb{R}^l\times\mathbb{Z}^{\mathrm{plaq}(\Sigma)}$, $a\in\mathbb{R}^v$ and $c\in\mathbb{Z}^l$, linear operators $M(x)$ and $\mathsf{T}_{\vartheta}(x',x;a,c)$ on $\mathcal{F}$ are defined by
\begin{multline}\label{eq:all-flux-9}
M(x)=s(A)\cdot\exp\left(-\frac{1}{4}\|\beta(x)\|^2\right)\cdot Q(A),\\
\mathsf{T}_{\vartheta}(x',x;a,c)=\chi\cdot\exp(2\pi i\,\vartheta\cdot w(c))\cdot(2\pi)^{-\frac{l}{2}}\cdot\exp\left(-\frac{1}{2}\|\mathcal{E}\|^2\right)\cdot M(x')^{*}\cdot\Gamma(a)\cdot M(x),
\end{multline}
where $\chi=\chi(x',x;c)$ and $\mathcal{E}=\mathcal{E}(x',x;a,c)$.

\begin{definition}[Slab functions and mirrors]
Let $\mathsf{m}\ge1$ be an integer. A slab function of thickness $\mathsf{m}$ is a bounded complex-valued Borel function $P$ on $(\mathbb{R}^l)^{\mathsf{m}}\times(\mathbb{R}^{\mathrm{plaq}(\Sigma)})^{\mathsf{m}-1}$, whose arguments are written $(\mathcal{E}_0,\dots,\mathcal{E}_{\mathsf{m}-1};\beta_1,\dots,\beta_{\mathsf{m}-1})$. The mirror of $P$ is the slab function $P^{\theta}$ of thickness $\mathsf{m}$ given by
\[
P^{\theta}(\mathcal{E}_0,\dots,\mathcal{E}_{\mathsf{m}-1};\beta_1,\dots,\beta_{\mathsf{m}-1})=\overline{P(-\mathcal{E}_{\mathsf{m}-1},\dots,-\mathcal{E}_0;\beta_{\mathsf{m}-1},\dots,\beta_1)}.
\]
\end{definition}

A slab function is read as a function of the curvature $F\in C^2(n)$ by taking for $\mathcal{E}_j$ the restriction of $F$ to the plaquettes $(x;0i)$, $i\in\{1,2\}$, with $x_0=j$, which join the slice $j$ to the slice $j+1$, and for $\beta_j$ the restriction of $F$ to the plaquettes $(x;12)$ with $x_0=j$. Thus a slab function of thickness $\mathsf{m}$ is a function of the curvature on those plaquettes of the slab $\{0\le x_0\le\mathsf{m}\}$ that are contained in neither of its two boundary slices; on the slab $\{j_0\le x_0\le j_0+\mathsf{m}\}$ it is read in the same way, with all times shifted by $j_0$. Let $\mathsf{r}$ be the reflection of $\Lambda_n$ in one of the two boundary slices of $\{0\le x_0\le\mathsf{m}\}$; it maps the adjacent slab on the other side of that slice onto $\{0\le x_0\le\mathsf{m}\}$. If $P$ is read on $\{0\le x_0\le\mathsf{m}\}$ and $P^{\theta}$ on the adjacent slab, then $P^{\theta}(F)=\overline{P(\mathsf{r}^{*}F)}$ for every $F\in C^2(n)$, by the rule for $\mathsf{r}^{*}$ in Section~\ref{sec:tori-cochains-lattice}.

For $\vartheta\in\mathbb{R}^2$ and a slab function $P$ of thickness $\mathsf{m}$, the kernel $C_{\mathsf{m},\vartheta}[P]$ on $\mathcal{X}\times\mathcal{X}$, with values in the linear operators on $\mathcal{F}$, is
\begin{multline}\label{eq:all-flux-10}
C_{\mathsf{m},\vartheta}[P](x_{\mathsf{m}},x_0)=\int_{\mathcal{X}^{\mathsf{m}-1}}\prod_{j=1}^{\mathsf{m}-1}\mathrm{d}x_j\int_{([0,r)^v)^{\mathsf{m}}}\prod_{j=0}^{\mathsf{m}-1}\mathrm{d}\mathfrak{m}(a_j)\sum_{c_0,\dots,c_{\mathsf{m}-1}\in\mathbb{Z}^l}\\
P(\mathcal{E}_0,\dots,\mathcal{E}_{\mathsf{m}-1};\beta_1,\dots,\beta_{\mathsf{m}-1})\\
\cdot\mathsf{T}_{\vartheta}(x_{\mathsf{m}},x_{\mathsf{m}-1};a_{\mathsf{m}-1},c_{\mathsf{m}-1})\cdots\mathsf{T}_{\vartheta}(x_1,x_0;a_0,c_0),
\end{multline}
where $x_0,x_{\mathsf{m}}\in\mathcal{X}$, $\mathcal{E}_j=\mathcal{E}(x_{j+1},x_j;a_j,c_j)$ for $0\le j\le\mathsf{m}-1$, and $\beta_j=\beta(x_j)$ for $1\le j\le\mathsf{m}-1$. The operator on $\mathcal{H}_{\vartheta}$ with this kernel is denoted by the same symbol: for $\Psi\in\mathcal{H}_{\vartheta}$ and $x'\in\mathcal{X}$,
\[
(C_{\mathsf{m},\vartheta}[P]\Psi)(x')=\int_{\mathcal{X}}\mathrm{d}x\,C_{\mathsf{m},\vartheta}[P](x',x)\cdot\Psi(x),
\]
and the values of $C_{\mathsf{m},\vartheta}[P]\Psi$ off $\mathcal{X}$ are fixed by \eqref{eq:all-flux-8}.

\begin{proposition}[Trace formula with Bloch angles]\label{prop:trace-formula-bloch}
Let $\Sigma$ be a cross-section and $\vartheta\in\mathbb{R}^2$, put $\kappa=(2\pi)^{\frac{l}{2}}$, and let $M$ also denote the operator of multiplication by the function $M$ of \eqref{eq:all-flux-9}. For every choice of slice data in the sense of Section~\ref{sec:slice-factorisation-determinant}, with which the objects of Section~\ref{sec:transfer-kernels} are formed, assertions (i), (ii) and (iii) below hold, and there is a bounded operator $K_{\vartheta}\ge0$ on $\mathcal{H}_{\vartheta}$ such that assertions (iv) and (v) hold.
\begin{enumerate}
\item[(i)] For every integer $\mathsf{m}\ge1$ and every slab function $P$ of thickness $\mathsf{m}$, the kernel \eqref{eq:all-flux-10} defines a Hilbert--Schmidt operator $C_{\mathsf{m},\vartheta}[P]$ on $\mathcal{H}_{\vartheta}$, which depends linearly on $P$.
\item[(ii)] For every slab function $P$ of thickness $\mathsf{m}$,
\[
C_{\mathsf{m},\vartheta}[P^{\theta}]=C_{\mathsf{m},\vartheta}[P]^{*}.
\]
\item[(iii)] Let $R\ge2$ and $\mathsf{m}_1,\dots,\mathsf{m}_R\ge1$ be integers, let $n=(n_0,n_1,n_2)$ be the torus with cross-section $\Sigma$ and $n_0=\mathsf{m}_1+\dots+\mathsf{m}_R$, and put $t_{\alpha}=\mathsf{m}_1+\dots+\mathsf{m}_{\alpha}$ for $0\le\alpha\le R$. For $\alpha=1,\dots,R$ let $P_{\alpha}$ be a slab function of thickness $\mathsf{m}_{\alpha}$, read on the slab $\{t_{\alpha-1}\le x_0\le t_{\alpha}\}$ of $\Lambda_n$, that is, as a function of the curvature on those plaquettes of this slab that lie in neither of its two boundary slices; thus $\prod_{\alpha}P_{\alpha}$ is a bounded Borel function on $C^2(n)$. Then, with the functionals $\omega^{(\mathfrak{n})}_n$ of Definition~\ref{def:all-flux-functional},
\begin{equation}\label{eq:all-flux-11}
\kappa^{n_0}\cdot\operatorname{Tr}_{\mathcal{H}_{\vartheta}}\left(C_{\mathsf{m}_R,\vartheta}[P_R]\cdots C_{\mathsf{m}_1,\vartheta}[P_1]\right)=\sum_{\mathfrak{n}\in\mathbb{Z}^3}\exp\left(2\pi i\left(\vartheta_1\mathfrak{n}_{01}+\vartheta_2\mathfrak{n}_{02}\right)\right)\cdot\omega^{(\mathfrak{n})}_n\left(\prod_{\alpha=1}^RP_{\alpha}\right).
\end{equation}
For $\mathfrak{n}_{12}\in\mathbb{Z}$ let $\mathcal{H}^{(\mathfrak{n}_{12})}_{\vartheta}$ be the space of those $\Psi\in\mathcal{H}_{\vartheta}$ with $\Psi(A,m)=0$ whenever $\sum_pm_p\ne\mathfrak{n}_{12}$, the sum being taken over the plaquettes $p$ of $\Sigma$. Then $\mathcal{H}_{\vartheta}$ is the orthogonal direct sum of the spaces $\mathcal{H}^{(\mathfrak{n}_{12})}_{\vartheta}$, $\mathfrak{n}_{12}\in\mathbb{Z}$; every operator $C_{\mathsf{m},\vartheta}[P]$ maps each of them into itself; and \eqref{eq:all-flux-11} remains valid when the trace is taken over $\mathcal{H}^{(\mathfrak{n}_{12})}_{\vartheta}$ and the sum on the right is restricted to the $\mathfrak{n}$ with this value of $\mathfrak{n}_{12}$.
\item[(iv)] The operator $T_{\vartheta}:=C_{1,\vartheta}[1]$ equals $M^{*}K_{\vartheta}M$, and $K_{\vartheta}^{\frac{1}{2}}M$ is a Hilbert--Schmidt operator. The operator $T_{\vartheta}$ is of trace class and $T_{\vartheta}\ge0$, and $C_{\mathsf{m},\vartheta}[1]=T_{\vartheta}^{\mathsf{m}}$ for every integer $\mathsf{m}\ge1$.
\item[(v)] Let $h\in\mathbb{R}^l$. The function on $\mathcal{X}\times\mathcal{X}$ given by
\begin{multline*}
(x',x)\longmapsto\int_{[0,r)^v}\mathrm{d}a\sum_{c\in\mathbb{Z}^l}\chi(x',x;c)\cdot\exp\left(2\pi i\,\vartheta\cdot w(c)\right)\\
\cdot(2\pi)^{-\frac{l}{2}}\cdot\exp\left(-\frac{1}{2}\left\|\mathcal{E}(x',x;a,c)+h\right\|^2\right)\cdot\Gamma(a)
\end{multline*}
is the kernel of $K_{\vartheta}u_h$ for some unitary operator $u_h$ on $\mathcal{H}_{\vartheta}$ that commutes with $K_{\vartheta}$.
\end{enumerate}
\end{proposition}

\begin{proof}
\emph{Kernels.} Restriction to $\mathcal{X}$ is a unitary map of $\mathcal{H}_{\vartheta}$ onto $L^2(\mathcal{X})\otimes\mathcal{F}$: it preserves the norm by \eqref{eq:all-flux-8}, and since $\mathcal{X}$ meets every class of $\sim$ in exactly one point, every element of $L^2(\mathcal{X})\otimes\mathcal{F}$ extends in exactly one way to a map with the first property in \eqref{eq:all-flux-8}. As for \eqref{eq:all-flux-10}, a measurable function $f$ on $\mathcal{X}\times\mathcal{X}$ with values in the linear operators on $\mathcal{F}$ is called the kernel of an operator $X$ on $\mathcal{H}_{\vartheta}$ if, for every $\Psi\in\mathcal{H}_{\vartheta}$ and almost every $x'\in\mathcal{X}$, the function $x\mapsto f(x',x)\cdot\Psi(x)$ is integrable over $\mathcal{X}$ and $(X\Psi)(x')=\int_{\mathcal{X}}\mathrm{d}x\,f(x',x)\cdot\Psi(x)$.

\emph{The operators $\mathsf{H}$, $U(a)$ and $K_{\vartheta}$; assertion (v).} For $y\in\mathbb{R}^l$ let $(\tau_y\Psi)(A,m)=\Psi(A-y,m)$. The translations $\tau_y$ preserve the first identity in \eqref{eq:all-flux-8} and the norm (the function $|\Psi|$ is constant on the classes of $\sim$, and the translate of $\mathcal{X}$ by $y$ is again a fundamental domain of $\sim$), and they form a strongly continuous unitary representation of $\mathbb{R}^l$ on $\mathcal{H}_{\vartheta}$. By Stone's theorem, $\tau_y=\exp\left(-i\sum_by_bE_b\right)$, where $b$ ranges over the links of $\Sigma$ and the generators $E_b=-i\partial_{A_b}$ are commuting self-adjoint operators on $\mathcal{H}_{\vartheta}$. Put
\[
\mathsf{H}=\exp\left(-\frac{1}{2}\sum_bE_b^2\right)\ge0.
\]
For $z,\xi\in\mathbb{R}^l$ one has
\begin{gather*}
\int_{\mathbb{R}^l}\mathrm{d}y\,(2\pi)^{-\frac{l}{2}}\cdot\exp\left(-\frac{1}{2}|y-z|^2-i\sum_by_b\xi_b\right)=\exp\left(-i\sum_bz_b\xi_b-\frac{1}{2}|\xi|^2\right),\\
\int_{\mathbb{R}^l}\mathrm{d}y\,\pi^{-\frac{l}{2}}\cdot\exp\left(-|y|^2-i\sum_by_b\xi_b\right)=\exp\left(-\frac{1}{4}|\xi|^2\right).
\end{gather*}
By the spectral theorem for the commuting family $(E_b)_b$ it follows that, as strong integrals,
\begin{gather*}
\mathsf{H}\tau_z=\int_{\mathbb{R}^l}\mathrm{d}y\,(2\pi)^{-\frac{l}{2}}\cdot\exp\left(-\frac{1}{2}|y-z|^2\right)\tau_y,\\
\mathsf{H}^{\frac{1}{2}}=\int_{\mathbb{R}^l}\mathrm{d}y\,\pi^{-\frac{l}{2}}\cdot\exp\left(-|y|^2\right)\tau_y.
\end{gather*}
Let $\varphi$ be the weight of one of these two integrals, let $\Psi\in\mathcal{H}_{\vartheta}$ and $(A',m')\in\mathcal{X}$. The sets $[0,r)^l-rc$, $c\in\mathbb{Z}^l$, partition $\mathbb{R}^l$. Writing $A'-y=A-rc$ with $A\in[0,r)^l$ and $c\in\mathbb{Z}^l$, and using \eqref{eq:all-flux-8} in the form $\Psi(A-rc,m')=\exp\left(2\pi i\,\vartheta\cdot w(c)\right)\cdot\Psi(A,m'-d_sc)$, one obtains
\begin{multline*}
\left(\int_{\mathbb{R}^l}\mathrm{d}y\,\varphi(y)\tau_y\Psi\right)(A',m')\\
=\sum_{c\in\mathbb{Z}^l}\int_{[0,r)^l}\mathrm{d}A\,\varphi(A'-A+rc)\cdot\exp\left(2\pi i\,\vartheta\cdot w(c)\right)\cdot\Psi(A,m'-d_sc).
\end{multline*}
Since $\chi(x',x;c)=1$ exactly when $m=m'-d_sc$, the operators $\mathsf{H}\tau_z$ and $\mathsf{H}^{\frac{1}{2}}$ have the kernels
\begin{gather*}
\mathsf{h}_z(x',x)=\sum_{c\in\mathbb{Z}^l}\chi(x',x;c)\cdot\exp\left(2\pi i\,\vartheta\cdot w(c)\right)\cdot(2\pi)^{-\frac{l}{2}}\cdot\exp\left(-\frac{1}{2}|A'-A+rc-z|^2\right),\\
\mathsf{g}(x',x)=\sum_{c\in\mathbb{Z}^l}\chi(x',x;c)\cdot\exp\left(2\pi i\,\vartheta\cdot w(c)\right)\cdot\pi^{-\frac{l}{2}}\cdot\exp\left(-|A'-A+rc|^2\right),
\end{gather*}
and $\mathsf{h}:=\mathsf{h}_0$ is the kernel of $\mathsf{H}$.

For $a\in\mathbb{R}^v$ put $U(a)=\Gamma(a)\tau_{d_sa}$, where $\Gamma(a)$ acts on the values of $\Psi$. The two factors commute, since $\Gamma(a)$ acts on the values and $\tau_{d_sa}$ on the arguments; hence, by the unitarity of $\Gamma(a)$ and the group law in \eqref{eq:all-flux-SD1}, $U$ is a unitary representation of $\mathbb{R}^v$ on $\mathcal{H}_{\vartheta}$. This representation is strongly continuous, because $a\mapsto\tau_{d_sa}$ is strongly continuous and $a\mapsto\Gamma(a)$ is continuous in the norm of operators on $\mathcal{H}_{\vartheta}$: by the formula $\Gamma(a)=\exp\left(iq\sum_{x\in\Sigma}a_x\,Q^{\mathrm{ch}}_x\right)$ of Section~\ref{sec:slice-factorisation-determinant}, $\Gamma$ is continuous in the operator norm on the finite-dimensional space $\mathcal{F}$, and the norm on $\mathcal{H}_{\vartheta}$ of $\Gamma(a)-\Gamma(a')$, acting on the values, is at most its norm on $\mathcal{F}$. The representation $U$ commutes with every $\tau_y$, hence with $\mathsf{H}$ and with $\mathsf{H}^{\frac{1}{2}}$, which are averages of translations. Let $k\in\mathbb{Z}^v$. Then $\Gamma(rk)=1$ by \eqref{eq:all-flux-SD1}. The integer cochain $c=-d_sk$ satisfies $d_sc=0$, so that \eqref{eq:all-flux-8} with this $c$ gives
\[
\tau_{rd_sk}\Psi=\exp\left(2\pi i\,\vartheta\cdot w(d_sk)\right)\cdot\Psi=\Psi;
\]
indeed $w_i(d_sk)$ is the sum of the increments of $k$ along the closed row $\mathrm{row}_i$, and vanishes. Hence $U(rk)=1$ for every $k\in\mathbb{Z}^v$, so that $U$ descends to a strongly continuous unitary representation of the compact group $(\mathbb{R}/r\mathbb{Z})^v$, whose normalised Haar measure is $r^{-v}\,\mathrm{d}a$ on $[0,r)^v$. Let $\Pi$ be the orthogonal projection of $\mathcal{H}_{\vartheta}$ onto the closed subspace of the vectors fixed by every $U(a)$. Averaging $U$ against this Haar measure produces $\Pi$, that is,
\[
\int_{[0,r)^v}\mathrm{d}a\,U(a)=r^v\,\Pi;
\]
indeed, the operator $r^{-v}\int_{[0,r)^v}\mathrm{d}a\,U(a)$ is self-adjoint, because $U(a)^{*}=U(-a)$ and the Haar measure is invariant under $a\mapsto-a$, and it is left unchanged when multiplied by any $U(b)$, by the invariance of the Haar measure under translations; so it is idempotent, every vector in its range is fixed by $U$, and it acts as the identity on the fixed vectors. Put
\[
K_{\vartheta}=\mathsf{H}\int_{[0,r)^v}\mathrm{d}a\,U(a)=r^v\,\mathsf{H}\Pi.
\]
As $\Pi$ commutes with $\mathsf{H}^{\frac{1}{2}}$, one has $K_{\vartheta}=r^v\,\mathsf{H}^{\frac{1}{2}}\Pi\,\mathsf{H}^{\frac{1}{2}}$; so $K_{\vartheta}$ is bounded and $K_{\vartheta}\ge0$. Moreover $K_{\vartheta}=\int_{[0,r)^v}\mathrm{d}a\,\Gamma(a)\,\mathsf{H}\tau_{d_sa}$, because $\Gamma(a)$ commutes with $\mathsf{H}$, and $A'-A+rc-d_sa=\mathcal{E}(x',x;a,c)$. Therefore $K_{\vartheta}$ has the kernel
\begin{multline*}
\int_{[0,r)^v}\mathrm{d}a\,\Gamma(a)\cdot\mathsf{h}_{d_sa}(x',x)\\
=\int_{[0,r)^v}\mathrm{d}a\sum_{c\in\mathbb{Z}^l}\chi(x',x;c)\cdot\exp\left(2\pi i\,\vartheta\cdot w(c)\right)\\
\cdot(2\pi)^{-\frac{l}{2}}\cdot\exp\left(-\frac{1}{2}\left\|\mathcal{E}(x',x;a,c)\right\|^2\right)\cdot\Gamma(a).
\end{multline*}
For $h\in\mathbb{R}^l$ put $u_h=\tau_{-h}$. Then $K_{\vartheta}u_h=\int_{[0,r)^v}\mathrm{d}a\,\Gamma(a)\,\mathsf{H}\tau_{d_sa-h}$ has the kernel $\int_{[0,r)^v}\mathrm{d}a\,\Gamma(a)\cdot\mathsf{h}_{d_sa-h}(x',x)$, which is the kernel displayed in (v), since $A'-A+rc-d_sa+h=\mathcal{E}(x',x;a,c)+h$. The operator $u_h$ is unitary and commutes with $\mathsf{H}$ and with every $U(a)$, hence with $\Pi$ and with $K_{\vartheta}$. This proves (v).

By \eqref{eq:all-flux-SD1}, $s$ and $Q$ depend on $A$ only through $\left(\exp(iqA_b)\right)_b$, which does not change under $A\mapsto A+rc$, $c\in\mathbb{Z}^l$, because $qr=2\pi$; and $\beta(A+rc,m-d_sc)=d_sA+rd_sc+rm-rd_sc=\beta(A,m)$. So the function $M$ is constant on the classes of $\sim$, multiplication by $M$ preserves the first identity in \eqref{eq:all-flux-8}, and $M$ is a bounded operator on $\mathcal{H}_{\vartheta}$, because $s$ and $Q$ are bounded.

\emph{Assertion (i).} By \eqref{eq:all-flux-SD1} the number $\lambda_0=\sup_{A\in\mathbb{R}^l}s(A)\cdot\|Q(A)\|$ is finite, and $\|M(x)\|\le\lambda_0\cdot\exp\left(-\frac{1}{4}\|\beta(x)\|^2\right)$. Write $\beta=\beta(x)$ and $\beta'=\beta(x')$. As $\Gamma(a)$ is unitary, \eqref{eq:all-flux-9} gives
\begin{multline*}
\left\|\mathsf{T}_{\vartheta}(x',x;a,c)\right\|\le\chi(x',x;c)\cdot(2\pi)^{-\frac{l}{2}}\cdot\exp\left(-\frac{1}{2}\left\|\mathcal{E}(x',x;a,c)\right\|^2\right)\\
\cdot\lambda_0^2\cdot\exp\left(-\frac{1}{4}\left(\|\beta'\|^2+\|\beta\|^2\right)\right).
\end{multline*}
Let $\mathsf{k}(x',x)\in[0,\infty]$ be the integral over $a\in[0,r)^v$ of the sum over $c\in\mathbb{Z}^l$ of the right side. The number
\[
S=\sup_{y\in\mathbb{R}^l}\sum_{c\in\mathbb{Z}^l}(2\pi)^{-\frac{l}{2}}\cdot\exp\left(-\frac{1}{2}|y+rc|^2\right)
\]
is finite, the sum being a product of $l$ bounded sums over $\mathbb{Z}$. Since $\chi\le1$, and $\mathcal{E}(x',x;a,c)=y+rc$ with $y=A'-A-d_sa$, it follows that
\[
\mathsf{k}(x',x)\le r^v\cdot S\cdot\lambda_0^2\cdot\exp\left(-\frac{1}{4}\left(\|\beta'\|^2+\|\beta\|^2\right)\right).
\]
For $A\in[0,r)^l$ and every plaquette $p$ of $\Sigma$ one has $|(d_sA+rm)_p|\ge r|m_p|-4r$; hence
\[
\int_{\mathcal{X}}\mathrm{d}x\,\exp\left(-\frac{1}{2}\|\beta(x)\|^2\right)=\sum_{m\in\mathbb{Z}^{\mathrm{plaq}(\Sigma)}}\int_{[0,r)^l}\mathrm{d}A\,\exp\left(-\frac{1}{2}\|d_sA+rm\|^2\right)<\infty.
\]
Consequently $\mathsf{k}$ is square integrable on $\mathcal{X}\times\mathcal{X}$, and by the Cauchy--Schwarz inequality so is its $\mathsf{m}$-fold composition
\[
\mathsf{k}^{*\mathsf{m}}(x_{\mathsf{m}},x_0)=\int_{\mathcal{X}^{\mathsf{m}-1}}\mathrm{d}x_1\cdots\mathrm{d}x_{\mathsf{m}-1}\,\mathsf{k}(x_{\mathsf{m}},x_{\mathsf{m}-1})\cdots\mathsf{k}(x_1,x_0)\qquad\left(\mathsf{k}^{*1}=\mathsf{k}\right),
\]
with $\left\|\mathsf{k}^{*\mathsf{m}}\right\|_{L^2}\le\left\|\mathsf{k}\right\|_{L^2}^{\mathsf{m}}$. In \eqref{eq:all-flux-10} the norm of the integrand is at most $\|P\|_{\infty}$ times the product of the bounds above for its factors $\mathsf{T}_{\vartheta}$. Therefore the integrals and sums that define $C_{\mathsf{m},\vartheta}[P](x_{\mathsf{m}},x_0)$ converge absolutely wherever $\mathsf{k}^{*\mathsf{m}}(x_{\mathsf{m}},x_0)<\infty$, in particular almost everywhere, and
\[
\left\|C_{\mathsf{m},\vartheta}[P](x_{\mathsf{m}},x_0)\right\|\le\|P\|_{\infty}\cdot\mathsf{k}^{*\mathsf{m}}(x_{\mathsf{m}},x_0).
\]
Consequently
\begin{multline*}
\int_{\mathcal{X}\times\mathcal{X}}\mathrm{d}x_{\mathsf{m}}\,\mathrm{d}x_0\,\operatorname{Tr}_{\mathcal{F}}\left(C_{\mathsf{m},\vartheta}[P](x_{\mathsf{m}},x_0)^{*}\cdot C_{\mathsf{m},\vartheta}[P](x_{\mathsf{m}},x_0)\right)\\
\le\dim\mathcal{F}\cdot\|P\|_{\infty}^2\cdot\left\|\mathsf{k}^{*\mathsf{m}}\right\|_{L^2}^2<\infty.
\end{multline*}
Identify $\mathcal{H}_{\vartheta}$ with $L^2(\mathcal{X})\otimes\mathcal{F}$ by restriction to $\mathcal{X}$, and take an orthonormal basis of $\mathcal{H}_{\vartheta}$ whose members are products of members of an orthonormal basis of $L^2(\mathcal{X})$ and of one of $\mathcal{F}$. By Parseval's identity, applied in $L^2(\mathcal{X})$ and in $\mathcal{F}$, the sum of $\left\|C_{\mathsf{m},\vartheta}[P]\Psi\right\|^2$ over the members $\Psi$ of this basis equals the left side of the last display. So $C_{\mathsf{m},\vartheta}[P]$ is a Hilbert--Schmidt operator on $\mathcal{H}_{\vartheta}$; it depends linearly on $P$ because \eqref{eq:all-flux-10} does.

\emph{Assertion (ii).} By \eqref{eq:all-flux-SD1}, $\Gamma(a)^{*}=\Gamma(-a)$; moreover $w(-c)=-w(c)$, $\mathcal{E}(x,x';-a,-c)=-\mathcal{E}(x',x;a,c)$ and $\chi(x,x';-c)=\chi(x',x;c)$. Taking adjoints in \eqref{eq:all-flux-9}, one obtains
\[
\mathsf{T}_{\vartheta}(x',x;a,c)^{*}=\mathsf{T}_{\vartheta}(x,x';-a,-c).
\]
For $k\in\mathbb{Z}^v$ one has $\mathcal{E}(x',x;a+rk,c+d_sk)=\mathcal{E}(x',x;a,c)$, $\chi(x',x;c+d_sk)=\chi(x',x;c)$ (as $d_sd_s=0$), $w(c+d_sk)=w(c)$ (as $w(d_sk)=0$) and $\Gamma(a+rk)=\Gamma(a)$ (by \eqref{eq:all-flux-SD1}); hence $\mathsf{T}_{\vartheta}(x',x;a+rk,c+d_sk)=\mathsf{T}_{\vartheta}(x',x;a,c)$. Let $\mathcal{E}_j$ and $\beta_j$ be as in \eqref{eq:all-flux-10}. The integrand of \eqref{eq:all-flux-10} depends on $(a_j,c_j)$ only through $\mathsf{T}_{\vartheta}(x_{j+1},x_j;a_j,c_j)$ and $\mathcal{E}_j$; so its sum over $c_j\in\mathbb{Z}^l$ is invariant under $a_j\mapsto a_j+rk$ (replace $c_j$ by $c_j+d_sk$), and the integration over $a_j\in[0,r)^v$ may be replaced by the integration over $a_j\in(-r,0]^v$.

Write \eqref{eq:all-flux-10} for $P^{\theta}$ in this form, and substitute
\[
\tilde{x}_i=x_{\mathsf{m}-i}\quad(0\le i\le\mathsf{m}),\qquad\tilde{a}_i=-a_{\mathsf{m}-1-i},\quad\tilde{c}_i=-c_{\mathsf{m}-1-i}\quad(0\le i\le\mathsf{m}-1).
\]
Then the $\tilde{a}_i$ range over $[0,r)^v$, the $\tilde{c}_i$ over $\mathbb{Z}^l$, and $\tilde{x}_1,\dots,\tilde{x}_{\mathsf{m}-1}$ over $\mathcal{X}$. With $\tilde{\mathcal{E}}_i=\mathcal{E}(\tilde{x}_{i+1},\tilde{x}_i;\tilde{a}_i,\tilde{c}_i)$ one has $\tilde{\mathcal{E}}_i=-\mathcal{E}_{\mathsf{m}-1-i}$ and $\beta(\tilde{x}_i)=\beta_{\mathsf{m}-i}$, so that
\[
P^{\theta}(\mathcal{E}_0,\dots,\mathcal{E}_{\mathsf{m}-1};\beta_1,\dots,\beta_{\mathsf{m}-1})=\overline{P\left(\tilde{\mathcal{E}}_0,\dots,\tilde{\mathcal{E}}_{\mathsf{m}-1};\beta(\tilde{x}_1),\dots,\beta(\tilde{x}_{\mathsf{m}-1})\right)},
\]
and, by the identity for $\mathsf{T}_{\vartheta}^{*}$,
\begin{multline*}
\mathsf{T}_{\vartheta}(x_{\mathsf{m}},x_{\mathsf{m}-1};a_{\mathsf{m}-1},c_{\mathsf{m}-1})\cdots\mathsf{T}_{\vartheta}(x_1,x_0;a_0,c_0)\\
=\left[\mathsf{T}_{\vartheta}(\tilde{x}_{\mathsf{m}},\tilde{x}_{\mathsf{m}-1};\tilde{a}_{\mathsf{m}-1},\tilde{c}_{\mathsf{m}-1})\cdots\mathsf{T}_{\vartheta}(\tilde{x}_1,\tilde{x}_0;\tilde{a}_0,\tilde{c}_0)\right]^{*}.
\end{multline*}
As $\tilde{x}_0=x_{\mathsf{m}}$ and $\tilde{x}_{\mathsf{m}}=x_0$, it follows that $C_{\mathsf{m},\vartheta}[P^{\theta}](x_{\mathsf{m}},x_0)=C_{\mathsf{m},\vartheta}[P](x_0,x_{\mathsf{m}})^{*}$, and the right side is the kernel of the adjoint of $C_{\mathsf{m},\vartheta}[P]$. This proves (ii).

\emph{Assertion (iii).} For $\alpha=1,\dots,R$ put $\mathsf{F}_{\alpha}=C_{\mathsf{m}_{\alpha},\vartheta}[P_{\alpha}]$, and let $f_{\alpha}$ be its kernel \eqref{eq:all-flux-10}; by (i) and its proof, $\mathsf{F}_{\alpha}$ is a Hilbert--Schmidt operator on $\mathcal{H}_{\vartheta}$ and $f_{\alpha}$ is square integrable on $\mathcal{X}\times\mathcal{X}$. As $R\ge2$, the product $\mathsf{F}_R\cdots\mathsf{F}_1$ is of trace class, being the product of the two Hilbert--Schmidt operators $\mathsf{F}_R$ and $\mathsf{F}_{R-1}\cdots\mathsf{F}_1$, and its trace can be read off from the kernels, multiplied around the cycle $\zeta_0\to\zeta_1\to\dots\to\zeta_{R-1}\to\zeta_0$ of points of $\mathcal{X}$:
\begin{multline*}
\operatorname{Tr}_{\mathcal{H}_{\vartheta}}\left(\mathsf{F}_R\cdots\mathsf{F}_1\right)\\
=\int_{\mathcal{X}^R}\mathrm{d}\zeta_0\cdots\mathrm{d}\zeta_{R-1}\,\operatorname{Tr}_{\mathcal{F}}\left[f_R(\zeta_0,\zeta_{R-1})\cdot f_{R-1}(\zeta_{R-1},\zeta_{R-2})\cdots f_1(\zeta_1,\zeta_0)\right].
\end{multline*}
Indeed, the trace of $\mathsf{F}_R\cdot\left(\mathsf{F}_{R-1}\cdots\mathsf{F}_1\right)$ is the Hilbert--Schmidt inner product of $\mathsf{F}_R^{*}$ and $\mathsf{F}_{R-1}\cdots\mathsf{F}_1$, which Parseval's identity writes as an integral over $\mathcal{X}\times\mathcal{X}$ of their kernels; and the kernel of $\mathsf{F}_{R-1}\cdots\mathsf{F}_1$ is the composition of $f_{R-1},\dots,f_1$, all integrals converging absolutely by the Cauchy--Schwarz inequality. In this formula take $\zeta_{\alpha}=x_{t_{\alpha}}$, and insert \eqref{eq:all-flux-10} for every $f_{\alpha}$ and \eqref{eq:all-flux-9} for every factor $\mathsf{T}_{\vartheta}$. The variables are then $x_j=(A_j,m_j)\in\mathcal{X}$, $a_j\in[0,r)^v$ and $c_j\in\mathbb{Z}^l$ for $j\in\mathbb{Z}/n_0\mathbb{Z}$; put $\mathcal{E}_j=\mathcal{E}(x_{j+1},x_j;a_j,c_j)$ and $\beta_j=\beta(x_j)$. The kernel $f_{\alpha}$ contributes the factors $\mathsf{T}_{\vartheta}(x_{j+1},x_j;a_j,c_j)$, $t_{\alpha-1}\le j<t_{\alpha}$, and the value of $P_{\alpha}$ at $(\mathcal{E}_j)_{t_{\alpha-1}\le j<t_{\alpha}}$ and $(\beta_j)_{t_{\alpha-1}<j<t_{\alpha}}$. The absolute value of the integrand is at most $\dim\mathcal{F}\cdot\prod_{\alpha}\|P_{\alpha}\|_{\infty}$ times the product over $j\in\mathbb{Z}/n_0\mathbb{Z}$ of the bounds for $\|\mathsf{T}_{\vartheta}(x_{j+1},x_j;a_j,c_j)\|$ obtained in the proof of (i). Integrated over the $a_j$ and summed over the $c_j$, this product becomes $\prod_j\mathsf{k}(x_{j+1},x_j)$, whose integral over $x_0,\dots,x_{n_0-1}$ is finite by the Cauchy--Schwarz inequality, as $\mathsf{k}$ is square integrable and $n_0\ge2$. So all integrals and sums converge absolutely, and by Fubini's theorem they may be carried out in any order.

Each of the $n_0$ factors $\mathsf{T}_{\vartheta}(x_{j+1},x_j;a_j,c_j)$ contains the number $(2\pi)^{-\frac{l}{2}}$, and $(2\pi)^{-\frac{n_0l}{2}}=\kappa^{-n_0}$. The weight $\exp\left(-\frac{1}{4}\|\beta_j\|^2\right)$ of the slice $j$ occurs twice: in $M(x_j)$, within the factor of the step from $j$ to $j+1$, and in $M(x_j)^{*}$, within the factor of the step from $j-1$ to $j$. Hence the scalar factors of the integrand multiply to
\[
\kappa^{-n_0}\cdot\prod_{j\in\mathbb{Z}/n_0\mathbb{Z}}\chi(x_{j+1},x_j;c_j)\cdot\exp\left(2\pi i\,\vartheta\cdot\sum_jw(c_j)\right)\cdot\exp\left(-\frac{1}{2}\sum_j\left(\|\mathcal{E}_j\|^2+\|\beta_j\|^2\right)\right),
\]
while the factors $s$, $Q^{*}$, $\Gamma$, $Q$ form the product in \eqref{eq:all-flux-7}, whose trace over $\mathcal{F}$ is $\Delta(qA)$ by \eqref{eq:all-flux-SD2}, with $A=(A_j,a_j)_{j\in\mathbb{Z}/n_0\mathbb{Z}}\in C^1(n)$.

Let $m\in C^2(n)$ be the integer cochain whose restriction to the plaquettes of the directions $1,2$ in the slice $j$ is $m_j$, and whose restriction to the plaquettes of the directions $0,i$, $i\in\{1,2\}$, that join the slice $j$ to the slice $j+1$ is $c_j$, for every $j$. By the formula for $dA$ in Section~\ref{sec:tori-cochains-lattice}, the restriction of $dA$ to the plaquettes of the first kind is $d_sA_j$, and its restriction to those of the second kind is $A_{j+1}-A_j-d_sa_j$. Hence the curvature $F=dA+rm$ has the restrictions $\beta_j$ and $\mathcal{E}_j$ to these two sets of plaquettes; consequently $\sum_j\left(\|\mathcal{E}_j\|^2+\|\beta_j\|^2\right)=\|F\|^2$, and, by the way slab functions are read (Section~\ref{sec:transfer-kernels}), the product of the values of the $P_{\alpha}$ in the integrand is $\left(\prod_{\alpha}P_{\alpha}\right)(F)$. By the formula for $dF$ in Section~\ref{sec:tori-cochains-lattice}, the restriction of $dm$ to the cubes labelled by the sites of the slice $j$ is $m_{j+1}-m_j-d_sc_j$; thus the constraints $\chi(x_{j+1},x_j;c_j)=1$, $j\in\mathbb{Z}/n_0\mathbb{Z}$, say exactly that $dm=0$. As the $(A_j,a_j)$ range over $[0,r)^l\times[0,r)^v$ and the families $(m_j,c_j)_j$ over those that satisfy the constraints, the pair $(A,m)$ ranges over $[0,r)^{\mathrm{links}}\times Z^2_{\mathrm{cl}}(n)$, and every pair is attained once. Finally, $\mathrm{row}_1$ consists of the links of direction $1$ of $\Sigma$ whose second coordinate is $0$, and $\mathrm{row}_2$ of the links of direction $2$ whose first coordinate is $0$; so, writing $y$ for the sites of $\Lambda_n$, Definition~\ref{def:all-flux-functional} gives
\[
\sum_jw_1(c_j)=\sum_{y:\,y_2=0}m_{01}(y)=\mathfrak{n}_{01}(m),\qquad\sum_jw_2(c_j)=\sum_{y:\,y_1=0}m_{02}(y)=\mathfrak{n}_{02}(m).
\]
Altogether,
\begin{multline*}
\kappa^{n_0}\cdot\operatorname{Tr}_{\mathcal{H}_{\vartheta}}\left(C_{\mathsf{m}_R,\vartheta}[P_R]\cdots C_{\mathsf{m}_1,\vartheta}[P_1]\right)\\
=\int_{[0,r)^{\mathrm{links}}}\mathrm{d}A\sum_{m\in Z^2_{\mathrm{cl}}(n)}\exp\left(2\pi i\left(\vartheta_1\mathfrak{n}_{01}(m)+\vartheta_2\mathfrak{n}_{02}(m)\right)\right)\\
\cdot\left(\prod_{\alpha=1}^RP_{\alpha}\right)(dA+rm)\cdot\exp\left(-\frac{1}{2}\|dA+rm\|^2\right)\cdot\Delta(qA),
\end{multline*}
and grouping the cochains $m$ according to their fluxes $(\mathfrak{n}_{12},\mathfrak{n}_{02},\mathfrak{n}_{01})(m)=\mathfrak{n}$ yields \eqref{eq:all-flux-11}, by \eqref{eq:all-flux-5}.

For $c\in\mathbb{Z}^l$ the sum of $d_sc$ over the plaquettes of $\Sigma$ vanishes, every value $c_b$ entering twice, with opposite signs. Hence $\sum_pm_p$ is unchanged under $(A,m)\mapsto(A+rc,m-d_sc)$ and under $m\mapsto m+d_sc$. By the first invariance, the condition defining $\mathcal{H}^{(\mathfrak{n}_{12})}_{\vartheta}$ is compatible with \eqref{eq:all-flux-8}; the spaces $\mathcal{H}^{(\mathfrak{n}_{12})}_{\vartheta}$ are mutually orthogonal, and every $\Psi\in\mathcal{H}_{\vartheta}$ is the sum of its products with the indicator functions of the sets $\left\{\sum_pm_p=\mathfrak{n}_{12}\right\}$, $\mathfrak{n}_{12}\in\mathbb{Z}$; so $\mathcal{H}_{\vartheta}=\bigoplus_{\mathfrak{n}_{12}\in\mathbb{Z}}\mathcal{H}^{(\mathfrak{n}_{12})}_{\vartheta}$. By the second invariance, $\chi(x',x;c)=0$ unless $\sum_pm'_p=\sum_pm_p$; so the kernels \eqref{eq:all-flux-9} and \eqref{eq:all-flux-10} vanish unless $\sum_pm'_p=\sum_pm_p$, and $C_{\mathsf{m},\vartheta}[P]$ maps every $\mathcal{H}^{(\mathfrak{n}_{12})}_{\vartheta}$ into itself. The orthogonal projection onto $\mathcal{H}^{(\mathfrak{n}_{12})}_{\vartheta}$ is the multiplication by the indicator function of $\left\{\sum_pm_p=\mathfrak{n}_{12}\right\}$, and the trace over $\mathcal{H}^{(\mathfrak{n}_{12})}_{\vartheta}$ of the product in \eqref{eq:all-flux-11} is the trace over $\mathcal{H}_{\vartheta}$ of the product of this projection with it. Placed to the left of $\mathsf{F}_R$, the projection multiplies $f_R(\zeta_0,\zeta_{R-1})$ by the value of the indicator function at $\zeta_0$; in the computation above this inserts the factor $1$ if $\sum_pm_0(p)=\mathfrak{n}_{12}$ and the factor $0$ otherwise. Here $\sum_pm_0(p)$ is the sum of the values of $m$ on the plaquettes of the directions $1,2$ in the slice $0$, which is $\mathfrak{n}_{12}(m)$ by Definition~\ref{def:all-flux-functional}. The sum over $m$ is thereby restricted to the $m$ with $\mathfrak{n}_{12}(m)=\mathfrak{n}_{12}$, which proves the last assertion of (iii).

\emph{Assertion (iv).} The factors $M(x')^{*}$ and $M(x)$ in \eqref{eq:all-flux-9} do not depend on $(a,c)$. Hence, by \eqref{eq:all-flux-10} with $\mathsf{m}=1$ and $P=1$, the kernel of $T_{\vartheta}=C_{1,\vartheta}[1]$ is $M(x')^{*}$ times the kernel of $K_{\vartheta}$ times $M(x)$, that is, $T_{\vartheta}=M^{*}K_{\vartheta}M$. As $0\le\Pi\le1$ and $\Pi$ commutes with $\mathsf{H}^{\frac{1}{2}}$, one has $K_{\vartheta}=r^v\,\mathsf{H}^{\frac{1}{2}}\Pi\,\mathsf{H}^{\frac{1}{2}}\le r^v\,\mathsf{H}$, and therefore
\begin{multline*}
\left\|K_{\vartheta}^{\frac{1}{2}}M\right\|_{\mathrm{HS}}^2=\operatorname{Tr}_{\mathcal{H}_{\vartheta}}\left(M^{*}K_{\vartheta}M\right)\le r^v\cdot\operatorname{Tr}_{\mathcal{H}_{\vartheta}}\left(M^{*}\mathsf{H}M\right)=r^v\cdot\left\|\mathsf{H}^{\frac{1}{2}}M\right\|_{\mathrm{HS}}^2\\
=r^v\int_{\mathcal{X}}\mathrm{d}x\,\mathsf{h}(x,x)\cdot\operatorname{Tr}_{\mathcal{F}}\left(M(x)^{*}M(x)\right).
\end{multline*}
For the last equality: $\mathsf{H}^{\frac{1}{2}}M$ has the kernel $\mathsf{g}(x',x)\cdot M(x)$, so that $\left\|\mathsf{H}^{\frac{1}{2}}M\right\|_{\mathrm{HS}}^2$ is the integral over $x\in\mathcal{X}$ of $\int_{\mathcal{X}}\mathrm{d}x'\,|\mathsf{g}(x',x)|^2$ times $\operatorname{Tr}_{\mathcal{F}}\left(M(x)^{*}M(x)\right)$; and, unfolding the sums over $c$ as above and evaluating a Gaussian integral,
\begin{multline*}
\int_{\mathcal{X}}\mathrm{d}x'\,|\mathsf{g}(x',x)|^2=\sum_{c\in\mathbb{Z}^l:\,d_sc=0}\exp\left(2\pi i\,\vartheta\cdot w(c)\right)\\
\cdot\int_{\mathbb{R}^l}\mathrm{d}Y\,\pi^{-l}\cdot\exp\left(-|Y+rc|^2-|Y|^2\right)=\mathsf{h}(x,x).
\end{multline*}
The diagonal value
\[
\mathsf{h}(x,x)=(2\pi)^{-\frac{l}{2}}\sum_{c\in\mathbb{Z}^l:\,d_sc=0}\exp\left(2\pi i\,\vartheta\cdot w(c)\right)\cdot\exp\left(-\frac{r^2}{2}|c|^2\right)
\]
does not depend on $x$ and is finite, and $\operatorname{Tr}_{\mathcal{F}}\left(M(x)^{*}M(x)\right)\le\dim\mathcal{F}\cdot\lambda_0^2\cdot\exp\left(-\frac{1}{2}\|\beta(x)\|^2\right)$, whose integral over $\mathcal{X}$ is finite by the proof of (i). Hence $K_{\vartheta}^{\frac{1}{2}}M$ is a Hilbert--Schmidt operator, and $T_{\vartheta}=\left(K_{\vartheta}^{\frac{1}{2}}M\right)^{*}\left(K_{\vartheta}^{\frac{1}{2}}M\right)$ is of trace class and $T_{\vartheta}\ge0$. Finally, for $P=1$ the kernel \eqref{eq:all-flux-10} is the $\mathsf{m}$-fold composition of the kernel of $T_{\vartheta}$, the integrals converging absolutely by the proof of (i); hence $C_{\mathsf{m},\vartheta}[1]=T_{\vartheta}^{\mathsf{m}}$.
\end{proof}

\subsection{Gaussian domination for the all-flux functional}\label{sec:gaussian-domination-all}

Let $X$ be a compact operator on a Hilbert space and $|X|=\left(X^{*}X\right)^{\frac{1}{2}}$. For $1\le p<\infty$ the Schatten norm of order $p$ is
\[
\|X\|_p=\left(\operatorname{Tr}|X|^p\right)^{\frac{1}{p}}\in[0,\infty],
\]
and the Schatten class $\mathfrak S_p$ is the set of compact operators $X$ with $\|X\|_p<\infty$; thus $\mathfrak S_1$ is the trace class and $\mathfrak S_2$ is the class of Hilbert--Schmidt operators. H\"older's inequality for Schatten norms, applied repeatedly and combined with $|\operatorname{Tr}X|\le\|X\|_1$ for $X\in\mathfrak S_1$, gives, for every positive integer $R$ and all $X_1,\dots,X_R\in\mathfrak S_R$,
\[
\left|\operatorname{Tr}\left(X_1\cdots X_R\right)\right|\le\prod_{a=1}^R\|X_a\|_R.
\]
Every Hilbert--Schmidt operator belongs to $\mathfrak S_R$ for $R\ge2$. If $R$ is even, then $\|X\|_R^R=\operatorname{Tr}\left(\left(X^{*}X\right)^{\frac{R}{2}}\right)$. Finally, $\|uX\|_R=\|X\|_R$ for every unitary operator $u$.

\begin{theorem}\label{thm:all-flux-1}
Let $n=(n_0,n_1,n_2)$ be a triple of integers $n_{\mu}\ge2$, each of which may be even or odd, and let $\omega_n=\omega^{\mathrm{af}}_n$ be the functional \eqref{eq:all-flux-5} of Definition~\ref{def:all-flux-functional}. Then
\[
\omega_n\left(\exp\left(-\langle h,F\rangle-\frac{1}{2}\|h\|^2\right)\right)\le\omega_n(1)\qquad\text{for every real }h\in C^2(n).
\]
Equivalently, $\left\langle\exp\langle g,F\rangle\right\rangle_n\le\exp\left(\frac{1}{2}\|g\|^2\right)$ for every real $g\in C^2(n)$, and
\[
\left|\left\langle\exp\langle g,F\rangle\right\rangle_n\right|\le\exp\left(\frac{1}{2}\|\operatorname{Re}g\|^2\right)\qquad\text{for every }g\in C^2(n)\otimes\mathbb C.
\]
\end{theorem}

\begin{proof}
Direction $0$ is time, as in Section~\ref{sec:trace-formula}, and the index $j$ ranges over $\mathbb Z/n_0\mathbb Z$, so that $j+1=0$ for $j=n_0-1$. The Bloch angle is $\vartheta=0$; accordingly, traces are taken in $\mathcal H=\mathcal H_0$, and $K=K_0$ is the operator of Proposition~\ref{prop:trace-formula-bloch}(iv). As in Proposition~\ref{prop:trace-formula-bloch}, $\kappa=(2\pi)^{\frac{l}{2}}$.

For each $j$ let $h^{\mathrm{el}}_j\in\mathbb R^l$, the electric part of $h$ at $j$, be the restriction of $h$ to the plaquettes $(x;0i)$ with $x_0=j$ and $i\in\{1,2\}$, which join slice $j$ to slice $j+1$; and let $b_j\in\mathbb R^{\mathrm{plaq}(\Sigma)}$, the magnetic part of $h$ at $j$, be the restriction of $h$ to the plaquettes $(x;12)$ with $x_0=j$, which lie in slice $j$. Thus $h=\left(h^{\mathrm{el}}_j,b_j\right)_j$. Since $\|F+h\|^2=\|F\|^2+2\langle h,F\rangle+\|h\|^2$, one has
\[
\exp\left(-\frac{1}{2}\|F\|^2-\langle h,F\rangle-\frac{1}{2}\|h\|^2\right)=\exp\left(-\frac{1}{2}\|F+h\|^2\right).
\]

Let $s$, $Q$ be the slice data of Section~\ref{sec:slice-factorisation-determinant}. For $b\in\mathbb R^{\mathrm{plaq}(\Sigma)}$ and $x=(A,m)$ in the slice space $\mathcal X$ of Section~\ref{sec:transfer-kernels}, with magnetic curvature $\beta(x)=d_sA+rm$, put
\[
M^b(x)=s(A)\cdot\exp\left(-\frac{1}{4}\|\beta(x)+b\|^2\right)\cdot Q(A),
\]
so that $M^0=M$. Acting by multiplication, $M^b$ is a bounded operator on $\mathcal H$, and $K^{\frac{1}{2}}M^b$ is a Hilbert--Schmidt operator by the proof of Proposition~\ref{prop:trace-formula-bloch}(iv), because the function $x\mapsto\exp\left(-\frac{1}{2}\|\beta(x)+b\|^2\right)$ is integrable on $\mathcal X$, as in the proof of Proposition~\ref{prop:trace-formula-bloch}(i). Let $u_{h'}$, $h'\in\mathbb R^l$, be the unitary operators of Proposition~\ref{prop:trace-formula-bloch}(v), and put
\[
R_j=M^{b_j}\left(M^{b_j}\right)^{*},\qquad u_j=u_{h^{\mathrm{el}}_j},\qquad Y_j=K^{\frac{1}{2}}R_jK^{\frac{1}{2}}=\left(K^{\frac{1}{2}}M^{b_j}\right)\left(K^{\frac{1}{2}}M^{b_j}\right)^{*}.
\]
Then $R_j\ge0$ is a bounded multiplication operator, $Y_j\ge0$ is of trace class, and $u_j$ commutes with $K$, hence with $K^{\frac{1}{2}}$.

Let $b,b'\in\mathbb R^{\mathrm{plaq}(\Sigma)}$ and $h'\in\mathbb R^l$. The operator $\left(M^{b'}\right)^{*}Ku_{h'}M^b=\left(K^{\frac{1}{2}}M^{b'}\right)^{*}u_{h'}\left(K^{\frac{1}{2}}M^b\right)$ is of trace class, and by Proposition~\ref{prop:trace-formula-bloch}(v) its kernel is
\[
(x',x)\longmapsto\int_{[0,r)^v}\mathrm{d}a\sum_{c\in\mathbb Z^l}\chi(x',x;c)\cdot(2\pi)^{-\frac{l}{2}}\cdot\exp\left(-\frac{1}{2}\left\|\mathcal E(x',x;a,c)+h'\right\|^2\right)\cdot M^{b'}(x')^{*}\cdot\Gamma(a)\cdot M^b(x),
\]
that is, the kernel $\mathsf T_0$ of \eqref{eq:all-flux-9}, integrated in $a$ and summed over $c$, after $\mathcal E$ has been replaced by $\mathcal E+h'$, $M(x')$ by $M^{b'}(x')$ and $M(x)$ by $M^b(x)$. The computation in the proof of Proposition~\ref{prop:trace-formula-bloch}(iii) applies to the product of these operators with $(b',b,h')=\left(b_{j+1},b_j,h^{\mathrm{el}}_j\right)$, taken in the order $j=n_0-1,\dots,0$ from left to right, and all integrals converge absolutely, as there. In the notation of that proof, with $\mathcal E_j=\mathcal E\left(x_{j+1},x_j;a_j,c_j\right)$ and $\beta_j=\beta\left(x_j\right)$, the Gaussian factors now multiply to
\[
(2\pi)^{-\frac{n_0l}{2}}\cdot\exp\left(-\frac{1}{2}\sum_j\left(\left\|\mathcal E_j+h^{\mathrm{el}}_j\right\|^2+\left\|\beta_j+b_j\right\|^2\right)\right)=\kappa^{-n_0}\cdot\exp\left(-\frac{1}{2}\|F+h\|^2\right),
\]
the factor $\exp\left(-\frac{1}{4}\left\|\beta_j+b_j\right\|^2\right)$ of slice $j$ occurring in both kernels adjacent to it; by \eqref{eq:all-flux-7}, the trace of the product of the matrix factors built from $s$, $Q^{*}$, $\Gamma$, $Q$ equals $\Delta(qA)$; and the variables of integration and summation are the pairs $(A,m)$ with $m\in Z^2_{\mathrm{cl}}(n)$. By the identity of exponentials above and by \eqref{eq:all-flux-5}, this proves the first of the equalities
\begin{align*}
\omega_n\left(\exp\left(-\langle h,F\rangle-\frac{1}{2}\|h\|^2\right)\right)
&=\kappa^{n_0}\cdot\operatorname{Tr}\prod_{j=n_0-1}^0\left(\left(M^{b_{j+1}}\right)^{*}Ku_jM^{b_j}\right)\\
&=\kappa^{n_0}\cdot\operatorname{Tr}\prod_{j=n_0-1}^0\left(K^{\frac{1}{2}}u_jK^{\frac{1}{2}}R_j\right)\\
&=\kappa^{n_0}\cdot\operatorname{Tr}\prod_{j=n_0-1}^0\left(u_jY_j\right).
\end{align*}
In the second equality the bounded operator $\left(M^{b_0}\right)^{*}$ at the left end of the product is moved to its right end, and $Ku_j=K^{\frac{1}{2}}u_jK^{\frac{1}{2}}$; in the third, the bounded operator $K^{\frac{1}{2}}$ at the left end is moved to the right end. Both moves rest on the cyclicity of the trace, $\operatorname{Tr}(BZ)=\operatorname{Tr}(ZB)$ for bounded $B$ and trace-class $Z$, which applies because the remaining product $Z$ is of trace class: once $K$ is written as $K^{\frac{1}{2}}K^{\frac{1}{2}}$, it contains a factor $Y_j$, as $n_0\ge2$.

Each $u_jY_j$ is of trace class, hence lies in $\mathfrak S_{n_0}$. H\"older's inequality with $R=n_0$, the invariance of $\|\cdot\|_{n_0}$ under left multiplication by a unitary operator, and $Y_j\ge0$ give
\[
\left|\operatorname{Tr}\prod_{j=n_0-1}^0\left(u_jY_j\right)\right|\le\prod_j\|Y_j\|_{n_0}=\prod_j\left(\operatorname{Tr}Y_j^{n_0}\right)^{\frac{1}{n_0}}.
\]
Fix $j$, and let $h^{(j)}\in C^2(n)$ vanish on the plaquettes that contain direction $0$ and coincide with $b_j$ on the plaquettes $(x;12)$ of every slice. The electric parts of $h^{(j)}$ vanish, its magnetic part at every $j'\in\mathbb Z/n_0\mathbb Z$ is $b_j$, and $u_{h'}$ is the identity operator for $h'=0$ (in the proof of Proposition~\ref{prop:trace-formula-bloch}, $u_{h'}=\tau_{-h'}$). The three equalities above, written for $h^{(j)}$ in place of $h$, therefore give
\[
\omega_n\left(\exp\left(-\left\langle h^{(j)},F\right\rangle-\frac{1}{2}\left\|h^{(j)}\right\|^2\right)\right)=\kappa^{n_0}\cdot\operatorname{Tr}Y_j^{n_0}.
\]
The left side of the first equality is nonnegative, since $\omega_n$ is a positive measure (Section~\ref{sec:all-flux-functional-2}). Combining the last three displays one obtains, for every $n$ and every real $h\in C^2(n)$,
\[
\omega_n\left(\exp\left(-\langle h,F\rangle-\frac{1}{2}\|h\|^2\right)\right)\le\prod_{j\in\mathbb Z/n_0\mathbb Z}\omega_n\left(\exp\left(-\left\langle h^{(j)},F\right\rangle-\frac{1}{2}\left\|h^{(j)}\right\|^2\right)\right)^{\frac{1}{n_0}}.
\]

This inequality also holds with direction $1$ as time. Indeed, let $\gamma:\Lambda_n\to\Lambda_{\gamma n}$ be the lattice isometry that exchanges the directions $0$ and $1$. Since $\omega_{\gamma n}\left(\Phi\circ\gamma^{*}\right)=\omega_n(\Phi)$ (Section~\ref{sec:all-flux-functional-2}), $\left\langle h,\gamma^{*}F\right\rangle=\left\langle\gamma_{*}h,F\right\rangle$ and $\left\|\gamma_{*}h\right\|=\|h\|$, the inequality for $\Lambda_{\gamma n}$ and $\gamma_{*}h$, with direction $0$ as time there, is the inequality for $\Lambda_n$ and $h$ with direction $1$ as time. In the latter the slices are the sets $\left\{x_1=k\right\}$, $k\in\mathbb Z/n_1\mathbb Z$, the magnetic parts of a two-cochain are its restrictions to the plaquettes $(x;02)$ of these slices, and the product on the right has $n_1$ factors, each with the exponent $\frac{1}{n_1}$. Consider $h^{(j)}$ with direction $1$ as time: it vanishes outside the plaquettes $(x;12)$, which contain direction $1$, so all its magnetic parts vanish, and every two-cochain on the right side is zero. Hence
\[
\omega_n\left(\exp\left(-\left\langle h^{(j)},F\right\rangle-\frac{1}{2}\left\|h^{(j)}\right\|^2\right)\right)\le\prod_{k\in\mathbb Z/n_1\mathbb Z}\omega_n(1)^{\frac{1}{n_1}}=\omega_n(1)
\]
for every $j$, and the preceding inequality yields
\[
\omega_n\left(\exp\left(-\langle h,F\rangle-\frac{1}{2}\|h\|^2\right)\right)\le\prod_{j\in\mathbb Z/n_0\mathbb Z}\omega_n(1)^{\frac{1}{n_0}}=\omega_n(1).
\]

For real $g\in C^2(n)$ the substitution $h=-g$ turns this inequality into $\omega_n\left(\exp\left(\langle g,F\rangle-\frac{1}{2}\|g\|^2\right)\right)\le\omega_n(1)$, which is the first restatement because $\omega_n(1)>0$ (Section~\ref{sec:all-flux-functional-2}); conversely, the first restatement for $g=-h$ is the inequality. For $g\in C^2(n)\otimes\mathbb C$ the pairing is bilinear and $F$ is real, so that $\left|\exp\langle g,F\rangle\right|=\exp\langle\operatorname{Re}g,F\rangle$; as $\omega_n$ is a positive measure, the first restatement for $\operatorname{Re}g$ gives
\[
\left|\omega_n\left(\exp\langle g,F\rangle\right)\right|\le\omega_n\left(\exp\langle\operatorname{Re}g,F\rangle\right)\le\exp\left(\frac{1}{2}\|\operatorname{Re}g\|^2\right)\cdot\omega_n(1),
\]
which is the second restatement. For real $g$ the second restatement contains the first, since then $\left\langle\exp\langle g,F\rangle\right\rangle_n>0$.
\end{proof}

\subsection{External fields and flux sectors}\label{sec:external-fields-flux}

Throughout this subsection, $\omega_n=\omega^{\mathrm{af}}_n$ is the functional \eqref{eq:all-flux-5} of Definition~\ref{def:all-flux-functional}, and $\omega^{(\mathfrak{n})}_n$, $\mathfrak{n}\in\mathbb{Z}^3$, are its flux sectors.

\subsubsection{Shifting the connection}\label{sec:shifting-connection}

For $\mathfrak{n}\in\mathbb{Z}^3$ let $Z^2_{\mathrm{cl}}(n;\mathfrak{n})$ be the set of those $m\in Z^2_{\mathrm{cl}}(n)$ with $(\mathfrak{n}_{12},\mathfrak{n}_{02},\mathfrak{n}_{01})(m)=\mathfrak{n}$; the set $Z^2_{\mathrm{cl}}(n)$ is the disjoint union of these sets. For $b\in C^1(n)\otimes\mathbb{C}$ put
\begin{align*}
Z^{\mathrm{ext}}_n(b)&=\int_{[0,r)^{\mathrm{links}}}\mathrm{d}A\,\sum_{m\in Z^2_{\mathrm{cl}}(n)}\exp\left(-\frac{1}{2}\|dA+rm\|^2\right)\cdot\Delta(qA+b),\\
\overline{Z}^{\mathrm{ext}}_n(b)&=\frac{1}{(2\pi)^3}\int_{[-\pi,\pi]^3}\mathrm{d}\theta\int_{V_{\mathrm{t}}}\mathrm{d}\sigma_{V_{\mathrm{t}}}(A)\,\exp\left(-\frac{1}{2}\|dA\|^2\right)\cdot\Delta(qA+qH_{\theta}+b),
\end{align*}
with $V_{\mathrm{t}}$, $\sigma_{V_{\mathrm{t}}}$ and $H_{\theta}$ as in Lemma~\ref{lem:sector-zero-flux}, and let $Z^{\mathrm{ext},(\mathfrak{n})}_n(b)$ be defined as $Z^{\mathrm{ext}}_n(b)$, with the sum over $m$ restricted to $Z^2_{\mathrm{cl}}(n;\mathfrak{n})$. By the first of the three properties of $\Delta$ in Section~\ref{sec:determinant-antiperiodic-spinors}, the one stated with \eqref{eq:all-flux-2}, and by \eqref{eq:all-flux-3}, the function $\Delta$ is entire and $|\Delta(B)|$ is bounded on every set of complex link fields $B$ with bounded imaginary part; consequently $Z^{\mathrm{ext}}_n$ and $\overline{Z}^{\mathrm{ext}}_n$ are entire functions of $b$.

\begin{lemma}[Shift identities]\label{lem:shift-identities}
Let $n=(n_0,n_1,n_2)$ with $n_{\mu}\ge2$ for each $\mu$, of either parity, and let $a\in C^1(n)\otimes\mathbb{C}$. With $Z_n(g)=\omega_n(\exp\langle g,F\rangle)$ as in Section~\ref{sec:all-flux-functional-2} and with $\overline{Z}_n$ as in \eqref{eq:all-flux-6}, one has
\begin{align*}
Z_n(da)&=\exp\left(\frac{1}{2}\langle da,da\rangle\right)\cdot Z^{\mathrm{ext}}_n(qa),\\
\omega^{(\mathfrak{n})}_n\left(\exp\langle da,F\rangle\right)&=\exp\left(\frac{1}{2}\langle da,da\rangle\right)\cdot Z^{\mathrm{ext},(\mathfrak{n})}_n(qa)\qquad\text{for every }\mathfrak{n}\in\mathbb{Z}^3,\\
\overline{Z}_n\left(\exp\langle da,F\rangle\right)&=\exp\left(\frac{1}{2}\langle da,da\rangle\right)\cdot\overline{Z}^{\mathrm{ext}}_n(qa).
\end{align*}
\end{lemma}

\begin{proof}
Suppose first that $a$ is real, and fix $\mathfrak{n}\in\mathbb{Z}^3$. Define $\mathsf{G}\colon C^1(n)\to[0,\infty)$ by
\[
\mathsf{G}(A)=\sum_{m\in Z^2_{\mathrm{cl}}(n;\mathfrak{n})}\exp\left(-\frac{1}{2}\|dA+rm\|^2+\langle da,dA+rm\rangle\right)\cdot\Delta(qA),
\]
so that $\omega^{(\mathfrak{n})}_n(\exp\langle da,F\rangle)=\int_{[0,r)^{\mathrm{links}}}\mathrm{d}A\,\mathsf{G}(A)$ by \eqref{eq:all-flux-5}. Let $c\in C^1(n)$ take integer values. Then $d(A+rc)+rm=dA+r(m+dc)$, and $m\mapsto m+dc$ is a bijection of $Z^2_{\mathrm{cl}}(n;\mathfrak{n})$ onto itself: $dc$ takes integer values, $d(dc)=0$, and for $\mu<\nu$, with $\iota$ the third direction, the sum of $(dc)_{\mu\nu}(x)$ over the $x$ with $x_{\iota}=0$ vanishes, since each value of $c$ that occurs in it occurs once with each sign. Moreover $q(A+rc)=qA+2\pi c$, so that $\Delta(q(A+rc))=\Delta(qA)$ by \eqref{eq:all-flux-2}. Hence $\mathsf{G}(A+rc)=\mathsf{G}(A)$. Thus $\mathsf{G}$ is $r$-periodic in each coordinate, and its integral over the cube $[0,r)^{\mathrm{links}}$ is unchanged by a translation of the argument:
\[
\int_{[0,r)^{\mathrm{links}}}\mathrm{d}A\,\mathsf{G}(A)=\int_{[0,r)^{\mathrm{links}}}\mathrm{d}A\,\mathsf{G}(A+a).
\]
For $F=dA+rm$ one has $d(A+a)+rm=F+da$ and
\[
-\frac{1}{2}\|F+da\|^2+\langle da,F+da\rangle=-\frac{1}{2}\|F\|^2+\frac{1}{2}\|da\|^2,
\]
whence
\[
\mathsf{G}(A+a)=\exp\left(\frac{1}{2}\|da\|^2\right)\cdot\sum_{m\in Z^2_{\mathrm{cl}}(n;\mathfrak{n})}\exp\left(-\frac{1}{2}\|dA+rm\|^2\right)\cdot\Delta(qA+qa).
\]
Integration over $[0,r)^{\mathrm{links}}$, together with $\|da\|^2=\langle da,da\rangle$, gives the second identity for real $a$. All terms of these sums are non-negative by \eqref{eq:all-flux-3}, and $Z^2_{\mathrm{cl}}(n)$ is the disjoint union of the sets $Z^2_{\mathrm{cl}}(n;\mathfrak{n})$; the sum of the second identity over $\mathfrak{n}\in\mathbb{Z}^3$ is therefore the first identity, by \eqref{eq:all-flux-5}.

Both sides of the first two identities are entire functions of the coordinates of $a\in C^1(n)\otimes\mathbb{C}$. On the left, the integrand $\exp\langle da,F\rangle$ is entire in $a$ and bounded in modulus by $\exp\left(\|d\operatorname{Re}a\|\,\|F\|\right)$, and $\omega_n$ is finite on $\mathfrak{F}(n)$ (Section~\ref{sec:all-flux-functional-2}). On the right, $\langle da,da\rangle$ is a polynomial in the coordinates of $a$, and $Z^{\mathrm{ext},(\mathfrak{n})}_n$ is entire for the same reason as $Z^{\mathrm{ext}}_n$. Two entire functions on $C^1(n)\otimes\mathbb{C}$ that coincide on $C^1(n)$ have the same Taylor coefficients at $0$, and by the identity theorem they coincide. This proves the first two identities for every $a\in C^1(n)\otimes\mathbb{C}$.

For the third identity, Lemma~\ref{lem:sector-zero-flux} with $\Phi=\exp\langle da,F\rangle$, which belongs to $\mathfrak{F}(n)$, gives
\[
\omega^{(0)}_n\left(\exp\langle da,F\rangle\right)=\mathcal{N}_n\cdot\overline{Z}_n\left(\exp\langle da,F\rangle\right).
\]
The proof of Lemma~\ref{lem:sector-zero-flux} uses $\Delta$ through the invariances $\Delta(B+2\pi\delta_b)=\Delta(B)$ and $\Delta(B+d\lambda)=\Delta(B)$ of \eqref{eq:all-flux-2}, and otherwise only through the absolute convergence of the integral over $[0,r)^{\mathrm{links}}$ of the series over $m$, which justifies there the interchange of sum and integral and the integration along the fibres and over the base by Fubini's theorem; the positivity of $\Delta$ in \eqref{eq:all-flux-3} serves only to write this absolute convergence as $\omega_n(|\Phi|)<\infty$. The function $B\mapsto\Delta(B+qa)$ is entire and has the same two invariances, since $\Delta$ has them at the point $B+qa$. By \eqref{eq:all-flux-2}, applied with $c=\max_b\left|\operatorname{Im}(qa_b)\right|$, there is a finite number $M_a$ such that $|\Delta(B)|\le M_a$ for every $B\in C^1(n)\otimes\mathbb{C}$ with $\max_b\left|\operatorname{Im}B_b\right|\le\max_b\left|\operatorname{Im}(qa_b)\right|$. For real $A,h\in C^1(n)$ one has $\operatorname{Im}(qA+qh+qa)=\operatorname{Im}(qa)$, hence $|\Delta(qA+qh+qa)|\le M_a$; this applies in particular to $h=0$ and to $h=H_{\theta}$, $\theta\in\mathbb{R}^3$. Therefore, for every $\Phi\in\mathfrak{F}(n)$, the integral
\[
\int_{[0,r)^{\mathrm{links}}}\mathrm{d}A\,\sum_{m\in Z^2_{\mathrm{cl}}(n)}\left|\Phi(dA+rm)\right|\cdot\exp\left(-\frac{1}{2}\|dA+rm\|^2\right)\cdot\left|\Delta(qA+qa)\right|
\]
is at most $M_a$ times the same integral without the factor $\left|\Delta(qA+qa)\right|$, which is finite by the bound of Section~\ref{sec:all-flux-functional-2} on the series over $m$, a bound in which $\Delta$ does not enter; so the absolute convergence holds for the weight $\Delta(\cdot+qa)$. With it in place of $\omega_n(|\Phi|)<\infty$, the proof of Lemma~\ref{lem:sector-zero-flux} carries over to this weight: Lemma~\ref{lem:sector-zero-flux} holds with $\Delta(\cdot+qa)$ in place of $\Delta$, with the same constant $\mathcal{N}_n$, which does not involve the weight. For $\Phi=1$ the left side of \eqref{eq:all-flux-6}, formed with $\Delta(\cdot+qa)$, is $Z^{\mathrm{ext},(0)}_n(qa)$, and the identity reads
\[
Z^{\mathrm{ext},(0)}_n(qa)=\mathcal{N}_n\cdot\overline{Z}^{\mathrm{ext}}_n(qa).
\]
Inserting the displayed equalities for $\omega^{(0)}_n\left(\exp\langle da,F\rangle\right)$ and for $Z^{\mathrm{ext},(0)}_n(qa)$ into the second identity for $\mathfrak{n}=0$ and dividing by $\mathcal{N}_n>0$ gives the third identity.
\end{proof}

\subsubsection{A bound for a single flux sector}\label{sec:bound-single-flux}

For $\mathfrak{n}=(\mathfrak{n}_{12},\mathfrak{n}_{02},\mathfrak{n}_{01})\in\mathbb{Z}^3$ put
\[
q_{\mathfrak{n}}(n)=\exp\left(-\frac{r^2}{2|n|}\sum_{\mu<\nu}n_{\iota}^2\,\mathfrak{n}_{\mu\nu}^2\right),
\]
where, in the term of the pair $\mu<\nu$, $\iota$ denotes the third direction.

\begin{lemma}[One flux sector]\label{lem:one-flux-sector}
Let $n=(n_0,n_1,n_2)$ with $n_{\mu}\ge2$ for each $\mu$. Let $\mathfrak{n}\in\mathbb{Z}^3$, and let $g\in C^2(n)$ be real with
\[
\sum_{x\in\Lambda_n}g_{\mu\nu}(x)=0\qquad\text{for each pair }\mu<\nu
\]
(this holds in particular for every exact $g$). Then
\[
\omega^{(\mathfrak{n})}_n\left(\exp\langle g,F\rangle\right)\le q_{\mathfrak{n}}(n)\cdot\exp\left(\frac{1}{2}\|g\|^2\right)\cdot\omega_n(1).
\]
\end{lemma}

\begin{proof}
For $\mu<\nu$ let $\mathbf{1}_{\mu\nu}\in C^2(n)$ be the indicator function of the set of $\mu\nu$-plaquettes, and let $\iota$ be the third direction. Let $F=dA+rm$ with $A\in C^1(n)$ and $m\in Z^2_{\mathrm{cl}}(n;\mathfrak{n}')$, $\mathfrak{n}'\in\mathbb{Z}^3$. The sum of $(dA)_{\mu\nu}(x)$ over $x\in\Lambda_n$ vanishes, since each value of $A$ that occurs in it occurs once with each sign; and for each $t\in\mathbb{Z}/n_{\iota}\mathbb{Z}$ the sum of $m_{\mu\nu}(x)$ over the $x$ with $x_{\iota}=t$ equals $\mathfrak{n}'_{\mu\nu}$ (Definition~\ref{def:all-flux-functional}). Hence
\[
\langle\mathbf{1}_{\mu\nu},F\rangle=r\sum_{x\in\Lambda_n}m_{\mu\nu}(x)=r\,n_{\iota}\,\mathfrak{n}'_{\mu\nu}.
\]
Further, $\|\mathbf{1}_{\mu\nu}\|^2=|n|$, the three functions $\mathbf{1}_{\mu\nu}$ have disjoint supports, and $\langle g,\mathbf{1}_{\mu\nu}\rangle=0$ by hypothesis. For real numbers $s_{\mu\nu}$, $\mu<\nu$, the real cochain $h=-g-\sum_{\mu<\nu}s_{\mu\nu}\mathbf{1}_{\mu\nu}$ therefore satisfies $\|h\|^2=\|g\|^2+|n|\sum_{\mu<\nu}s_{\mu\nu}^2$. Theorem~\ref{thm:all-flux-1} for this $h$, together with the decomposition $\omega_n=\sum_{\mathfrak{n}'}\omega^{(\mathfrak{n}')}_n$ of \eqref{eq:all-flux-5}, gives
\begin{align*}
&\sum_{\mathfrak{n}'\in\mathbb{Z}^3}\omega^{(\mathfrak{n}')}_n\left(\exp\langle g,F\rangle\right)\cdot\exp\left(\sum_{\mu<\nu}s_{\mu\nu}\,r\,n_{\iota}\,\mathfrak{n}'_{\mu\nu}\right)\\
&\qquad\le\exp\left(\frac{1}{2}\|g\|^2+\frac{1}{2}|n|\sum_{\mu<\nu}s_{\mu\nu}^2\right)\cdot\omega_n(1).
\end{align*}
The functionals $\omega^{(\mathfrak{n}')}_n$ are positive by \eqref{eq:all-flux-3}, so every term on the left is non-negative, and the term with $\mathfrak{n}'=\mathfrak{n}$ alone obeys the same bound. Choose $s_{\mu\nu}=r\,n_{\iota}\,\mathfrak{n}_{\mu\nu}/|n|$. Then, for each pair $\mu<\nu$,
\[
\frac{1}{2}|n|\,s_{\mu\nu}^2-s_{\mu\nu}\,r\,n_{\iota}\,\mathfrak{n}_{\mu\nu}=-\frac{r^2n_{\iota}^2\,\mathfrak{n}_{\mu\nu}^2}{2|n|},
\]
and the bound for the term $\mathfrak{n}'=\mathfrak{n}$ becomes the assertion of the lemma.
\end{proof}

\section{Wilson operators with block penalties at zero field}\label{sec:wilson-operators-block}

This section is concerned with the inverses of the free Wilson operator with block penalties (Definition~\ref{def:blocked-wilson-problems}) on the lattice $\eta\mathbb{Z}^3$ and on its tori, for an arbitrary mesh $\eta>0$; exponential decay and H\"older bounds are proved for them, and neither a gauge field nor a charge occurs. The section is self-contained, its conventions being fixed in Section~\ref{sec:conventions} and its constants in Section~\ref{sec:constants-order}; the results of it that are applied later are Theorem~\ref{thm:sup-holder-bounds-2} and Proposition~\ref{prop:wilson-operators-2}, both in Section~\ref{sec:inverses-zero-field}.

\subsection{Conventions}\label{sec:conventions}

The following conventions are in force throughout this section.

\begin{enumerate}
\item\label{it:wilson-operators-1} \emph{Scaling.} The conventions and definitions of this section apply to an arbitrary mesh $\eta>0$ of the lattice $\eta\mathbb{Z}^3$. Let $t>0$. For a set $X\subset\eta\mathbb{Z}^3$ put $tX:=\{tx:\ x\in X\}$, and for a function $u$ on $X$ let $U_tu$ be the function on $tX$ given by $(U_tu)(x):=u(x/t)$. Then $U_t$ maps the space $\mathbb{F}(\eta,R,\boldsymbol{\theta})$ of Convention~\ref{it:wilson-operators-2} onto $\mathbb{F}(t\eta,tR,\boldsymbol{\theta})$, and $1_{tX}\,U_t=U_t\,1_X$ on the functions on $\eta\mathbb{Z}^3$. The dilation $x\mapsto tx$ acts on a blocked Wilson problem $\mathbb{P}=(\eta,R,\boldsymbol{\theta};\gamma;W,\mathbb{E},\operatorname{Pen},\mathfrak{m})$ (Definition~\ref{def:blocked-wilson-problems}) as follows: the mesh, the side $R$, the set $W$ and the blocks are multiplied by $t$, the character $\boldsymbol{\theta}$ and the Clifford data $\gamma$ are kept, the mass is divided by $t$, and the penalty is replaced by $t^{-1}$ times its transport under $U_t$. Thus $\mathbb{P}$ is carried into
\[
\tilde{\mathbb{P}}:=\left(t\eta,tR,\boldsymbol{\theta};\gamma;tW,t\mathbb{E},\tilde{\operatorname{Pen}},\mathfrak{m}/t\right),\qquad t\mathbb{E}:=\left\{tE:\ E\in\mathbb{E}\right\},\qquad\tilde{\operatorname{Pen}}:=t^{-1}\,U_t\,\operatorname{Pen}\,U_t^{-1}.
\]
If $E$ is the block of level $i$ with centre $c$ (Convention~\ref{it:wilson-operators-4}), then $tE$ is the block of level $i$ with centre $tc$ of the lattice $t\eta\mathbb{Z}^3$, and its side is $t\cdot\boldsymbol{\ell}_E$. The scale function of $\tilde{\mathbb{P}}$ (Convention~\ref{it:wilson-operators-6}) is $x\mapsto t\cdot\boldsymbol{\ell}(x/t)$, and $\tilde{\operatorname{Pen}}=\sum_{E\in\mathbb{E}}\tilde{\operatorname{Pen}}_{tE}$, where $\tilde{\operatorname{Pen}}_{tE}:=t^{-1}\,U_t\,\operatorname{Pen}_E\,U_t^{-1}$ is the operator~\eqref{eq:wilson-operators-B3} of the block $tE$ with the kernel
\[
\tilde{K}_{tE}(x,x'):=t^{-4}\,K_E\left(\frac{x}{t},\frac{x'}{t}\right)\qquad(x,x'\in tE).
\]
For $u,v\in L^2(E)$, with the scalar product $\langle\cdot,\cdot\rangle_E$ of Definition~\ref{def:blocked-wilson-problems} and the means $P_E$ of Section~\ref{sec:block-means-poincare}, one has $U_tu,U_tv\in L^2(tE)$, $\langle U_tu,U_tv\rangle_{tE}=t^3\,\langle u,v\rangle_E$ and $P_{tE}U_tu=U_tP_Eu$, and
\[
\tilde{\operatorname{Pen}}_{tE}\,U_tu=t^{-1}\,U_t\,\operatorname{Pen}_Eu,\qquad\tilde{\operatorname{Pen}}_{tE}^{*}\,U_tu=t^{-1}\,U_t\,\operatorname{Pen}_E^{*}u,
\]
the adjoints being taken in $L^2(tE)$ and in $L^2(E)$. Consequently each inequality in~\eqref{eq:wilson-operators-B2}, \eqref{eq:wilson-operators-Pen1}, \eqref{eq:wilson-operators-Pen2}, \eqref{eq:wilson-operators-Pen3} and~\eqref{eq:wilson-operators-B4}, written for $\tilde{\mathbb{P}}$ at the sites $tx$, $tz$, $tx'$ and at the function $U_tu$, arises from the same inequality for $\mathbb{P}$ at $x$, $z$, $x'$ and $u$ on multiplying both of its sides by one and the same power of $t$, namely by $t$, $t^2$, $t$, $t^{-4}$, and by $t^{-1}$ or $1$, respectively; and~\eqref{eq:wilson-operators-B1}, including the periodicity required there and in~\eqref{eq:wilson-operators-B3}, holds for $\tilde{\mathbb{P}}$ with $tR$ in place of $R$. Hence $\tilde{\mathbb{P}}$ is again a blocked Wilson problem, with the same separation $s$ and the same penalty constants $(c_1,C_1,C_2)$. Let $D^{(\eta)}$ denote the operator $D$ of~\eqref{eq:wilson-operators-7} on the functions on $\eta\mathbb{Z}^3$. Under $U_t$ the shift by $\eta e_{\mu}$ on $\eta\mathbb{Z}^3$ corresponds to the shift by $t\eta e_{\mu}$ on $t\eta\mathbb{Z}^3$, so that the second expression in~\eqref{eq:wilson-operators-7} gives
\[
D^{(t\eta)}\,U_t=t^{-1}\,U_t\,D^{(\eta)}.
\]
Together with $1_{tW}\,U_t=U_t\,1_W$ and the definition of $\tilde{\operatorname{Pen}}$, this shows that the operator~\eqref{eq:wilson-operators-10} of $\tilde{\mathbb{P}}$, which acts on $U_t\mathbb{H}_W$, is
\[
T_{\tilde{\mathbb{P}}}=t^{-1}\,U_t\,T_{\mathbb{P}}\,U_t^{-1}.
\]
By this covariance no unit of length is distinguished. Let $L$ be the odd integer of Convention~\ref{it:wilson-operators-4}. For an integer $k$, lengths are said to be measured in units of level $k$ if $\eta=L^{-k}$; in that case one writes
\[
\boldsymbol{\ell}_j:=L^{-(k-j)}=L^j\eta\qquad(j\in\mathbb{Z}).
\]
Every statement of this section that is formulated for an arbitrary mesh $\eta>0$ is invariant under the dilations, which act on blocked Wilson problems by $\mathbb{P}\mapsto\tilde{\mathbb{P}}$. Theorem~\ref{thm:sup-holder-bounds-2} is formulated at the normalised mesh $\eta=L^{-i}$, that is, in units of level $i$, where $i\ge1$ is the integer in its hypothesis. Lemma~\ref{lem:inverse-symbol-unit} is formulated at the mesh $\eta=1$ and for a mass $m\in[0,1]$; Proposition~\ref{prop:lattice-green-function}, which holds at every mesh, is reduced to it by a dilation, with $m=\mathfrak{m}\cdot\eta$. The identification of $D$ with the Wilson operator $D_A$ at $A=0$ in Section~\ref{sec:shifts-differences-free} is made at the mesh $\eta=1$, and the description of the penalties~\eqref{eq:wilson-operators-Std2} in units of level $k$ in Definition~\ref{def:standard-penalties} is given at the mesh $\eta=L^{-k}$.

\item\label{it:wilson-operators-2} \emph{Tori and characters.} Also for a torus the underlying lattice is $\eta\mathbb{Z}^3$; the torus is encoded by quasi-periodicity. For $\mu\in\{1,2,3\}$ let $e_\mu$ be the unit vector of $\mathbb{R}^3$ in the direction $\mu$. The torus of side $R\in\eta\mathbb{N}$ with character $\boldsymbol{\theta}\in\mathbb{R}^3$ is described by the space of quasi-periodic functions
\[
\mathbb{F}(\eta,R,\boldsymbol{\theta}):=\left\{u:\eta\mathbb{Z}^3\to\mathbb{C}^4 \;:\; u(x+Re_\mu)=\exp(i\boldsymbol{\theta}_\mu)\,u(x)\ \ (x\in\eta\mathbb{Z}^3,\ \mu=1,2,3)\right\}.
\]
For $R=\infty$, the bulk, no condition is imposed: $\mathbb{F}(\eta,\infty,\boldsymbol{\theta})$ is the space of all functions $\eta\mathbb{Z}^3\to\mathbb{C}^4$. Periodic spinors correspond to $\boldsymbol{\theta}=0$, antiperiodic spinors to $\boldsymbol{\theta}=(\pi,\pi,\pi)$. Sets of sites, partitions and operators are thus objects on the covering lattice $\eta\mathbb{Z}^3$. For $R<\infty$ the objects that describe a blocked Wilson problem (Definition~\ref{def:blocked-wilson-problems}) on the torus are invariant under the translations $x\mapsto x+Re_\mu$, $\mu=1,2,3$: the set $W$; the family $\mathbb{E}$ of blocks as a family, that is, $E\pm Re_\mu\in\mathbb{E}$ for every $E\in\mathbb{E}$; the functions $\boldsymbol{\ell}$, $\boldsymbol{\lambda}$ and $d$ built from them (Convention~\ref{it:wilson-operators-6}, Lemma~\ref{lem:lipschitz-minorant-scale}(a), Definition~\ref{def:path-distance-units}); and the operators $D$, $\operatorname{Pen}$ and $1_W(D+\mathfrak{m})1_W+\operatorname{Pen}$ on the functions $\eta\mathbb{Z}^3\to\mathbb{C}^4$, the last of which coincides on $\mathbb{H}_W$ with the operator $T$ of~\eqref{eq:wilson-operators-10}. These operators commute with the translations by $Re_\mu$ and therefore map $\mathbb{F}(\eta,R,\boldsymbol{\theta})$ into itself. A single block and a period cell (Convention~\ref{it:wilson-operators-5}), on the other hand, are not invariant under these translations, and the operators $\operatorname{Pen}_E$, $E\in\mathbb{E}$, and the means $P_E$ extended by zero off $E$, which are attached to single blocks, do not commute with them; nor are the sets $X$ of the norms $\|\cdot\|_{2,X}$ and $\|\cdot\|_{\infty,X}$ of Convention~\ref{it:wilson-operators-5}, with their indicator functions $1_X$, invariant in general. These are objects on the covering lattice $\eta\mathbb{Z}^3$.

Let $R<\infty$, and let $\boldsymbol{\ell}_{\max}$ be the largest side of the blocks (Convention~\ref{it:wilson-operators-4}) of a partition as in Definition~\ref{def:blocked-wilson-problems}; the partition being $R$-periodic, $R$ is an integer multiple of $\boldsymbol{\ell}_{\max}$. Consider a cube of side $R$ whose faces lie in faces of blocks of side $\boldsymbol{\ell}_{\max}$; examples are $[-\boldsymbol{\ell}_{\max}/2,\,R-\boldsymbol{\ell}_{\max}/2)^3$ and, when $R/\eta$ is odd, also the cube $[-R/2,R/2]^3$ centred at the origin. The sites of such a cube form a fundamental domain of $\eta\mathbb{Z}^3$ modulo $R\mathbb{Z}^3$, and restricting $u\in\mathbb{F}(\eta,R,\boldsymbol{\theta})$ to them is a bijection onto the functions on the torus $\eta\mathbb{Z}^3/R\mathbb{Z}^3$. This bijection transforms the free Wilson operator $D$ of~\eqref{eq:wilson-operators-7} into the operator on the torus in which the bonds that cross the boundary of the cube in the direction $\mu$ carry the phase $\exp(i\boldsymbol{\theta}_\mu)$. Each block of side at most $\boldsymbol{\ell}_{\max}$ lies either in the cube or in its complement, so that means over blocks remain plain means.

\item\label{it:wilson-operators-3} \emph{Clifford data.} $\gamma_1,\gamma_2,\gamma_3$ denote Hermitian $4\times4$ matrices with
\[
\gamma_\mu\gamma_\nu+\gamma_\nu\gamma_\mu=2\delta_{\mu\nu}\qquad(\mu,\nu=1,2,3),
\]
and $\hat{\Gamma}$ denotes a Hermitian $4\times4$ matrix with $\hat{\Gamma}^2=1$, subject to no condition that involves $\gamma_1,\gamma_2,\gamma_3$. Let
\[
\Pi_\mu^{\pm}:=\frac{1\pm\gamma_\mu}{2}\qquad(\mu=1,2,3).
\]
Since $\gamma_\mu$ is Hermitian with $\gamma_\mu^2=1$, these are orthogonal projections, and $\Pi_\mu^{+}\Pi_\mu^{-}=0$. The triples $(-\gamma_\mu)_{\mu}$ and $(-\gamma_\mu^T)_{\mu}$, and the matrix $\hat{\Gamma}^T$, satisfy the same conditions; for this reason the class of blocked Wilson problems (Definition~\ref{def:blocked-wilson-problems}) is stable under adjoints and transposes. The dimension of the spinor space $\mathbb{C}^4$ is immaterial.

\item\label{it:wilson-operators-4} \emph{Blocks.} $L$ is an odd integer with $L\ge5$, and $|z|_\infty:=\max_{\mu}|z_\mu|$ for $z\in\mathbb{R}^3$. For an integer $i\ge0$ and $c\in L^i\eta\mathbb{Z}^3$, the block of level $i$ with centre $c$ is
\[
B_i(c):=\left\{x\in\eta\mathbb{Z}^3 \;:\; |x-c|_\infty\le\frac{(L^i-1)\,\eta}{2}\right\}.
\]
Its side is $\boldsymbol{\ell}:=L^i\eta$; one has $\#|B_i(c)|=(\boldsymbol{\ell}/\eta)^3$, and the closed cube of the block is $c+[-\boldsymbol{\ell}/2,\boldsymbol{\ell}/2]^3$. For each $i$ the blocks of level $i$ form a partition of $\eta\mathbb{Z}^3$, and every block is a union of blocks of each lower level. Since $L^i$ is odd, every face of the closed cube of a block of level $i$ lies in a plane $\{x\in\mathbb{R}^3:\ x_\mu=(2m+1)\eta/2\}$ with $\mu\in\{1,2,3\}$ and $m\in\mathbb{Z}$, and such a plane contains no site.

\item\label{it:wilson-operators-5} \emph{Pairings, norms and kernels.} For $v,w\in\mathbb{C}^4$ let $v\cdot w:=\sum_{j=1}^4v_jw_j$. A period cell is a set $\Lambda\subset\eta\mathbb{Z}^3$ that contains exactly one point of each class of $\eta\mathbb{Z}^3$ modulo $R\mathbb{Z}^3$; for $R=\infty$, $\Lambda=\eta\mathbb{Z}^3$. For $u,v\in\mathbb{F}(\eta,R,\boldsymbol{\theta})$ let
\[
\langle u,v\rangle:=\eta^3\sum_{x\in\Lambda}\overline{u(x)}\cdot v(x),\qquad \|u\|^2:=\langle u,u\rangle\in[0,\infty],
\]
and for $u\in\mathbb{F}(\eta,R,\boldsymbol{\theta})$ and $v\in\mathbb{F}(\eta,R,-\boldsymbol{\theta})$ let
\[
(u,v):=\eta^3\sum_{x\in\Lambda}u(x)\cdot v(x);
\]
this form is bilinear in $(u,v)\in\mathbb{F}(\eta,R,\boldsymbol{\theta})\times\mathbb{F}(\eta,R,-\boldsymbol{\theta})$. In each of these sums the summand is $R$-periodic, so that the choice of $\Lambda$ is immaterial. Adjoints $A^{*}$ and transposes $A^T$ of operators $A$ are taken with respect to $\langle\cdot,\cdot\rangle$ and $(\cdot,\cdot)$, respectively. The kernel of an operator $A$ is the function $(x,x')\mapsto A(x,x')$ on $\eta\mathbb{Z}^3\times\eta\mathbb{Z}^3$, with values in the $4\times4$ matrices, for which
\[
(Af)(x)=\eta^3\sum_{x'}A(x,x')\,f(x'),
\]
the sum extending over one period cell, as above. If $R=\infty$, the period cell is infinite, and the sums defining $\langle u,v\rangle$ and $(u,v)$, as well as the sum in the kernel identity, are taken only for arguments for which the series converges absolutely. For the two pairings this is the case if $\|u\|$ and $\|v\|$ are finite or if one of $u$, $v$ has finite support; for the kernel identity it is the case if $f$ has finite support. The series $\langle u,u\rangle$ has nonnegative terms, so that $\|u\|\in[0,\infty]$ is defined for every $u$. On $\mathbb{C}^4$, $|\cdot|$ is the Euclidean norm, and on matrices it is the associated operator norm. For a set $X\subset\eta\mathbb{Z}^3$ of sites let
\[
\|u\|_{2,X}^2:=\eta^3\sum_{x\in X}|u(x)|^2,\qquad \|u\|_{\infty,X}:=\sup_{x\in X}|u(x)|,
\]
and let $1_X$ denote the indicator function of $X$.

\item\label{it:wilson-operators-6} \emph{The scale function.} Let $W\subset\eta\mathbb{Z}^3$ be partitioned into blocks as in Definition~\ref{def:blocked-wilson-problems}. The scale function $\boldsymbol{\ell}:\eta\mathbb{Z}^3\to(0,\infty)$ is
\[
\boldsymbol{\ell}(x):=\begin{cases}\boldsymbol{\ell}_E, & x\in W,\\ L\eta, & x\notin W,\end{cases}
\]
where $\boldsymbol{\ell}_E$ is the side of the block $E$ of the partition that contains $x$. Thus every site outside $W$ is assigned the side of a block of level one.

\item\label{it:wilson-operators-7} \emph{H\"older quotients.} Let $\boldsymbol{\alpha}\in(0,1)$ be the H\"older order of~\eqref{eq:wilson-operators-1}, and let $W$ and the scale function $\boldsymbol{\ell}$ be as in Convention~\ref{it:wilson-operators-6}. For a function $u$ on $W$ and $x\in W$ let
\[
[u]_{\boldsymbol{\alpha};x}:=\sup\left\{\frac{\boldsymbol{\ell}(x)^{\boldsymbol{\alpha}}\cdot|u(x)-u(\tilde{x})|}{|x-\tilde{x}|_\infty^{\boldsymbol{\alpha}}} \;:\; \tilde{x}\in W,\ 0<|x-\tilde{x}|_\infty\le\boldsymbol{\ell}(x)\right\}.
\]
Even when $R<\infty$, the points $x,\tilde{x}$ of this quotient, and of every later H\"older quotient, are sites of $\eta\mathbb{Z}^3$, and $|x-\tilde{x}|_\infty$ is the maximum norm of their difference in $\eta\mathbb{Z}^3$.
\end{enumerate}

\subsection{Constants and their order}\label{sec:constants-order}

The constants are chosen in the order of the displays \eqref{eq:wilson-operators-1}--\eqref{eq:wilson-operators-6}; the quantities in each display depend only on those in the displays before it. None of them depends on $\eta$, $R$, $\boldsymbol{\theta}$, the Clifford data of Convention~\ref{it:wilson-operators-3}, the number of block levels, $W$, the partition, the mass $\mathfrak{m}$, the resolvent parameter $y$ of Definition~\ref{def:standard-penalties}, or the separation $s$ of Definition~\ref{def:blocked-wilson-problems}.
\begin{align}
&b>0,\qquad \boldsymbol{\alpha}\in(0,1),\qquad L\ge5\ \text{an odd integer};\label{eq:wilson-operators-1}\\
&c_1:=b\left(1-\frac{1}{L}\right),\qquad C_1:=L\cdot b,\qquad C_2:=2L^4\cdot b;\label{eq:wilson-operators-2}\\
&K_1:=\sqrt{L}+\sqrt{L\cdot b}+\sqrt{6(L+1)},\qquad C_I:=\frac{2L}{c_1}+\frac{3}{4}\,L^2\cdot K_1^2+L;\label{eq:wilson-operators-3}\\
&c_3:=6\exp(1)+3\exp(1)\cdot L\cdot C_2,\qquad a_0:=\min\left\{\frac{1}{3L},\,\frac{1}{2C_I\cdot c_3}\right\},\qquad s_0:=\frac{48}{a_0};\label{eq:wilson-operators-4}\\
&C_G\ \text{an absolute constant},\qquad C_R:=\frac{7\cdot10^4\cdot C_G}{1-\boldsymbol{\alpha}};\label{eq:wilson-operators-5}\\
&n_L:=(L+2)^3,\quad \boldsymbol{\gamma}_3:=\frac{a_0}{2},\quad C_3,\quad C_{\mathrm{sh}}(\sigma),\quad \boldsymbol{\gamma}_4:=\frac{a_0}{4},\quad \tilde{C},\quad \boldsymbol{\gamma}_5,\quad C_5,\quad c_0.\label{eq:wilson-operators-6}
\end{align}
Here $c_1$, $C_1$, $C_2$ are the constants of Lemma~\ref{lem:bounds-standard-penalties}; $K_1$ and $C_I$ are those of Theorem~\ref{thm:weighted-bound}, and $c_3$ and $a_0$ those of Theorem~\ref{thm:exponentially-weighted-bound}. Since $a_0\le1/(3L)$, one has $s_0\ge144L$. Further, $C_G$ is the absolute constant of Proposition~\ref{prop:lattice-green-function}, and $C_R$ is the constant of Proposition~\ref{prop:interior-estimates}. In~\eqref{eq:wilson-operators-6}, $C_3$ is given by~\eqref{eq:wilson-operators-20}, and $C_{\mathrm{sh}}(\sigma)$, for $\sigma>0$, by~\eqref{eq:wilson-operators-11}; $\tilde{C}$ is the constant of Section~\ref{sec:bounded-right-hand}, $\boldsymbol{\gamma}_5$ and $C_5$ are the constants of Theorem~\ref{thm:sup-holder-bounds-2}, and $c_0$ is the constant of Proposition~\ref{prop:wilson-operators-2}.

The separation is subject to $s\ge6L$ in Theorems~\ref{thm:weighted-bound} and~\ref{thm:exponentially-weighted-bound}; Definition~\ref{def:blocked-wilson-problems} already requires $s\in[6L,\infty]$. It is subject to $s\ge s_0$ in Theorem~\ref{thm:sup-holder-bounds} and throughout Section~\ref{sec:bounded-right-hand}, hence in Corollary~\ref{cor:wilson-operators-1}. In Section~\ref{sec:homogeneous-blocked-wilson}, hence in Proposition~\ref{prop:wilson-operators-1} and in Theorem~\ref{thm:sup-holder-bounds-2}, the blocked Wilson problems have separation $s=\infty$, so that both thresholds are met; the separation is not among the data $i$, $\eta=L^{-i}$, $R$, $\boldsymbol{\theta}$, $\mathfrak{m}$, $y$ of Theorem~\ref{thm:sup-holder-bounds-2}.

\subsection{Difference calculus on the lattice}\label{sec:difference-calculus-lattice}

\subsubsection{Shifts, differences and the free Wilson operator}\label{sec:shifts-differences-free}

Fix $\eta>0$, $R\in\eta\mathbb{N}\cup\{\infty\}$ and $\boldsymbol{\theta}\in\mathbb{R}^3$. Let $\mathbb{F}(\eta,R,\boldsymbol{\theta})$ be the space of Convention~\ref{it:wilson-operators-2}, with the pairings $\langle\cdot,\cdot\rangle$, $(\cdot,\cdot)$ and the norm $\|\cdot\|$ of Convention~\ref{it:wilson-operators-5}, and put
\[
\mathbb{H}:=\mathbb{H}(\eta,R,\boldsymbol{\theta}):=\left\{u\in\mathbb{F}(\eta,R,\boldsymbol{\theta}):\|u\|<\infty\right\}.
\]
With the scalar product $\langle\cdot,\cdot\rangle$ this is a Hilbert space; for $R<\infty$ it is the whole of $\mathbb{F}(\eta,R,\boldsymbol{\theta})$. Indices $\mu,\nu$ of directions take the values $1,2,3$, and $\sum_{\mu}$ denotes the sum over $\mu\in\{1,2,3\}$.

The operators below are defined on $\mathbb{F}(\eta,R,\boldsymbol{\theta})$ and map this space into itself; for $R=\infty$ it consists of all functions $\eta\mathbb{Z}^3\to\mathbb{C}^4$. For $u\in\mathbb{F}(\eta,R,\boldsymbol{\theta})$ and $x\in\eta\mathbb{Z}^3$ let
\begin{align*}
(\tau_{\mu}u)(x)&:=u(x+\eta e_{\mu}),\\
\nabla_{\mu}^{+}&:=\frac{\tau_{\mu}-1}{\eta},\qquad
\nabla_{\mu}^{-}:=\frac{1-\tau_{\mu}^{-1}}{\eta},\qquad
\nabla_{\mu}^s:=\frac{\tau_{\mu}-\tau_{\mu}^{-1}}{2\eta},\\
\Delta_{\mu}&:=\frac{\tau_{\mu}+\tau_{\mu}^{-1}-2}{\eta^2},\qquad
\Delta:=\sum_{\mu}\Delta_{\mu},
\end{align*}
and $\|\nabla^{+}u\|^2:=\sum_{\mu}\|\nabla_{\mu}^{+}u\|^2$. The shifts $\tau_1,\tau_2,\tau_3$ commute with one another, with every matrix in $\mathbb{C}^{4\times4}$ acting on the values of a function, and with the translations $u\mapsto u(\cdot+a)$, $a\in\eta\mathbb{Z}^3$. They map $\mathbb{H}$ onto $\mathbb{H}$ and are unitary there: for $u\in\mathbb{F}(\eta,R,\boldsymbol{\theta})$ the function $|u|^2$ is $R$-periodic, so its sum over a period cell does not change when the cell is translated.

Let $D$ be the operator on $\mathbb{F}(\eta,R,\boldsymbol{\theta})$ given by the first of the three expressions
\begin{align}
D&:=\sum_{\mu}\gamma_{\mu}\nabla_{\mu}^s-\frac{\eta}{2}\,\Delta\notag\\
&=-\frac{1}{\eta}\sum_{\mu}\left(\Pi_{\mu}^{-}\tau_{\mu}+\Pi_{\mu}^{+}\tau_{\mu}^{-1}-1\right)\notag\\
&=\sum_{\mu}\left(\Pi_{\mu}^{+}\nabla_{\mu}^{-}-\Pi_{\mu}^{-}\nabla_{\mu}^{+}\right),\label{eq:wilson-operators-7}
\end{align}
where $\gamma_{\mu}$ and $\Pi_{\mu}^{\pm}=(1\pm\gamma_{\mu})/2$ are as in Convention~\ref{it:wilson-operators-3}; $D$ is the free Wilson operator, with Wilson parameter $1$. The three expressions coincide. Indeed, for each $\mu$,
\[
\gamma_{\mu}\,\frac{\tau_{\mu}-\tau_{\mu}^{-1}}{2\eta}-\frac{\tau_{\mu}+\tau_{\mu}^{-1}-2}{2\eta}
=-\frac{1}{\eta}\left(\frac{1-\gamma_{\mu}}{2}\,\tau_{\mu}+\frac{1+\gamma_{\mu}}{2}\,\tau_{\mu}^{-1}-1\right),
\]
which is the equality of the first two expressions, and the third results from the second on writing $1=\Pi_{\mu}^{+}+\Pi_{\mu}^{-}$.

For $\eta=1$, for $R$ equal to the odd integer $n$ of Section~\ref{sec:lattice-measures-laws} and for $\boldsymbol{\theta}=(\pi,\pi,\pi)$, the space $\mathbb{F}(1,R,\boldsymbol{\theta})$ consists of the functions $f\colon\mathbb{Z}^3\to\mathbb{C}^4$ with $f(x+R\,e_{\mu})=-f(x)$ for all $x\in\mathbb{Z}^3$ and all $\mu$. If in this case the direction $\mu$ is identified with the direction $\mu-1$ of Section~\ref{sec:dirac-matrices-lattice}, where the directions are numbered $0,1,2$, and $\gamma_{\mu}$ is the matrix $\gamma_{\mu-1}$ of that section ($\mu=1,2,3$), then $D$ is the Wilson operator $D_A$ of Section~\ref{sec:lattice-measures-laws} for $A=0$.

\begin{lemma}[Identities for the free Wilson operator]\label{lem:identities-free-wilson}
\begin{enumerate}
\item[(a)] For every $u\in\mathbb{H}$ one has $\operatorname{Re}\langle u,D u\rangle=\frac{\eta}{2}\,\|\nabla^{+}u\|^2$.
\item[(b)] For every $u\in\mathbb{H}$,
\[
\|D u\|^2=\|\nabla^{+}u\|^2+\frac{\eta^2}{4}\sum_{\mu\ne\nu}\|\nabla_{\mu}^{+}\nabla_{\nu}^{+}u\|^2\ge\|\nabla^{+}u\|^2.
\]
\item[(c)] Let $\psi\colon\eta\mathbb{Z}^3\to\mathbb{R}$ and $u\colon\eta\mathbb{Z}^3\to\mathbb{C}^4$ be arbitrary functions, no periodicity and no summability being assumed. Let $D$ act on all functions $\eta\mathbb{Z}^3\to\mathbb{C}^4$ by~\eqref{eq:wilson-operators-7}, let $\psi$ also denote the operator of multiplication by $\psi$, and let $[D,\psi]:=D\psi-\psi D$. Then, for every $x\in\eta\mathbb{Z}^3$,
\begin{multline}\label{eq:wilson-operators-8}
([D,\psi]u)(x)=-\frac{1}{\eta}\sum_{\mu}[\left(\psi(x+\eta e_{\mu})-\psi(x)\right)\Pi_{\mu}^{-}u(x+\eta e_{\mu})\\
+\left(\psi(x-\eta e_{\mu})-\psi(x)\right)\Pi_{\mu}^{+}u(x-\eta e_{\mu})].
\end{multline}
Moreover, for arbitrary vectors $a_{\mu},b_{\mu}\in\mathbb{C}^4$, $\mu=1,2,3$,
\begin{equation}\label{eq:wilson-operators-9}
\left|\sum_{\mu}\left(\Pi_{\mu}^{-}a_{\mu}+\Pi_{\mu}^{+}b_{\mu}\right)\right|^2\le3\sum_{\mu}\left(|a_{\mu}|^2+|b_{\mu}|^2\right).
\end{equation}
\item[(d)] The adjoint $D^{*}$ of $D$ on $\mathbb{H}$ is the operator~\eqref{eq:wilson-operators-7} built with the matrices $-\gamma_{\mu}$ instead of $\gamma_{\mu}$. The transpose $D^T$ of $D$ for the pairing $(\cdot,\cdot)$ is the operator~\eqref{eq:wilson-operators-7} built with $-\gamma_{\mu}^T$ instead of $\gamma_{\mu}$, acting on $\mathbb{F}(\eta,R,-\boldsymbol{\theta})$.
\end{enumerate}
\end{lemma}

\begin{proof}
As $\tau_{\mu}$ is unitary on $\mathbb{H}$, its adjoint is $\tau_{\mu}^{-1}$. Consequently
\[
\left(\nabla_{\mu}^{+}\right)^{*}=\frac{\tau_{\mu}^{-1}-1}{\eta}=-\nabla_{\mu}^{-},\qquad
-\Delta_{\mu}=\left(\nabla_{\mu}^{+}\right)^{*}\nabla_{\mu}^{+},\qquad
\left(\nabla_{\mu}^s\right)^{*}=-\nabla_{\mu}^s,\qquad
\Delta^{*}=\Delta.
\]
Let $u\in\mathbb{H}(\eta,R,\boldsymbol{\theta})$ and $v\in\mathbb{H}(\eta,R,-\boldsymbol{\theta})$; for $R<\infty$ these are arbitrary elements of $\mathbb{F}(\eta,R,\boldsymbol{\theta})$ and of $\mathbb{F}(\eta,R,-\boldsymbol{\theta})$. The pairings $(\tau_{\mu}u,v)$ and $(u,\tau_{\mu}^{-1}v)$ are then defined: for $R<\infty$ they are finite sums, and for $R=\infty$ the series converge absolutely by the Cauchy--Schwarz inequality, the shifts preserving the norm $\|\cdot\|$. Let $\Lambda$ be a period cell (Convention~\ref{it:wilson-operators-5}). Replacing $x$ by $x-\eta e_{\mu}$ in the sum one finds
\[
(\tau_{\mu}u,v)=\eta^3\sum_{x\in\Lambda}u(x+\eta e_{\mu})\cdot v(x)=\eta^3\sum_{x\in\Lambda+\eta e_{\mu}}u(x)\cdot v(x-\eta e_{\mu})=(u,\tau_{\mu}^{-1}v);
\]
the last equality holds for $R<\infty$ because the summand $x\mapsto u(x)\cdot v(x-\eta e_{\mu})$ is $R$-periodic and $\Lambda+\eta e_{\mu}$ is again a period cell, and for $R=\infty$ because then $\Lambda+\eta e_{\mu}=\Lambda=\eta\mathbb{Z}^3$. Thus $\tau_{\mu}^T=\tau_{\mu}^{-1}$ for the bilinear pairing, and the four identities above remain true with ${}^T$ in place of ${}^{*}$.

\emph{Proof of (d).} A matrix acting on the values of functions commutes with the shifts; for $\langle\cdot,\cdot\rangle$ its adjoint is the conjugate transposed matrix, and for $(\cdot,\cdot)$ its transpose is the transposed matrix. Since $\gamma_{\mu}$ is Hermitian,
\[
\left(\gamma_{\mu}\nabla_{\mu}^s\right)^{*}=\left(\nabla_{\mu}^s\right)^{*}\gamma_{\mu}=-\gamma_{\mu}\nabla_{\mu}^s,\qquad
\left(\gamma_{\mu}\nabla_{\mu}^s\right)^T=\left(\nabla_{\mu}^s\right)^T\gamma_{\mu}^T=-\gamma_{\mu}^T\nabla_{\mu}^s.
\]
Applied to the first expression in~\eqref{eq:wilson-operators-7}, together with $\Delta^{*}=\Delta^T=\Delta$, this gives (d).

\emph{Proof of (a).} Let $A:=\sum_{\mu}\gamma_{\mu}\nabla_{\mu}^s$ and $B:=-\frac{\eta}{2}\,\Delta$, so that $D=A+B$. By the last display, $A^{*}=-A$; hence $\langle u,A u\rangle$ is purely imaginary. On the other hand
\[
\langle u,B u\rangle=\frac{\eta}{2}\sum_{\mu}\left\langle u,\left(\nabla_{\mu}^{+}\right)^{*}\nabla_{\mu}^{+}u\right\rangle=\frac{\eta}{2}\,\|\nabla^{+}u\|^2,
\]
and (a) follows.

\emph{Proof of (b).} The operators $A$ and $B$ commute, since the shifts commute with one another and with the matrices $\gamma_{\mu}$; and $D^{*}=-A+B$, because $A^{*}=-A$ and $B^{*}=B$. Hence
\[
D^{*}D=(-A+B)(A+B)=-A^2+B^2.
\]
The relations $\gamma_{\mu}\gamma_{\nu}+\gamma_{\nu}\gamma_{\mu}=2\delta_{\mu\nu}$ of Convention~\ref{it:wilson-operators-3}, together with the commutativity of the shifts, give $A^2=\sum_{\mu}(\nabla_{\mu}^s)^2$. For fixed $\mu$ write $t:=\tau_{\mu}$; then
\begin{align*}
-\left(\nabla_{\mu}^s\right)^2+\frac{\eta^2}{4}\,\Delta_{\mu}^2
&=\frac{-\left(t-t^{-1}\right)^2+\left(t+t^{-1}-2\right)^2}{4\eta^2}\\
&=\frac{\left(2-t^2-t^{-2}\right)+\left(t^2+t^{-2}+6-4t-4t^{-1}\right)}{4\eta^2}\\
&=\frac{2-t-t^{-1}}{\eta^2}=-\Delta_{\mu}.
\end{align*}
As $B^2=\frac{\eta^2}{4}\left(\sum_{\mu}\Delta_{\mu}\right)^2$, summation over $\mu$ gives
\[
D^{*}D=-\Delta+\frac{\eta^2}{4}\sum_{\mu\ne\nu}\Delta_{\mu}\Delta_{\nu},\qquad
\Delta_{\mu}\Delta_{\nu}=\left(\nabla_{\mu}^{+}\nabla_{\nu}^{+}\right)^{*}\left(\nabla_{\mu}^{+}\nabla_{\nu}^{+}\right),
\]
the second identity because all operators involved are polynomials in the commuting shifts $\tau_1^{\pm1},\tau_2^{\pm1},\tau_3^{\pm1}$ and $-\Delta_{\mu}=(\nabla_{\mu}^{+})^{*}\nabla_{\mu}^{+}$. Now $\|D u\|^2=\langle u,D^{*}D u\rangle$, $\langle u,-\Delta u\rangle=\|\nabla^{+}u\|^2$ and $\langle u,\Delta_{\mu}\Delta_{\nu}u\rangle=\|\nabla_{\mu}^{+}\nabla_{\nu}^{+}u\|^2$, which is the identity in (b); the inequality holds because the added terms are non-negative.

\emph{Proof of (c).} By the second expression in~\eqref{eq:wilson-operators-7}, for every $x\in\eta\mathbb{Z}^3$,
\begin{multline*}
\left(D(\psi u)\right)(x)=-\frac{1}{\eta}\sum_{\mu}[\psi(x+\eta e_{\mu})\,\Pi_{\mu}^{-}u(x+\eta e_{\mu})\\
+\psi(x-\eta e_{\mu})\,\Pi_{\mu}^{+}u(x-\eta e_{\mu})-\psi(x)\,u(x)].
\end{multline*}
Subtracting $\psi(x)\,(D u)(x)$, the terms containing $u(x)$ cancel, and~\eqref{eq:wilson-operators-8} remains. For~\eqref{eq:wilson-operators-9}, recall from Convention~\ref{it:wilson-operators-3} that $\Pi_{\mu}^{-}$ and $\Pi_{\mu}^{+}$ are orthogonal projections with $\Pi_{\mu}^{+}\Pi_{\mu}^{-}=0$. Their ranges are therefore orthogonal, and for $a,b\in\mathbb{C}^4$
\[
\left|\Pi_{\mu}^{-}a+\Pi_{\mu}^{+}b\right|^2=\left|\Pi_{\mu}^{-}a\right|^2+\left|\Pi_{\mu}^{+}b\right|^2\le|a|^2+|b|^2.
\]
By the Cauchy--Schwarz inequality, $|v_1+v_2+v_3|^2\le3\sum_{\mu}|v_{\mu}|^2$ for $v_1,v_2,v_3\in\mathbb{C}^4$; with $v_{\mu}:=\Pi_{\mu}^{-}a_{\mu}+\Pi_{\mu}^{+}b_{\mu}$ this gives~\eqref{eq:wilson-operators-9}.
\end{proof}

\begin{remark}\label{rem:wilson-operators-1}
For $p\in\mathbb{R}^3$ put $t_{\mu}:=1-\cos p_{\mu}$. Under the Fourier transform, part~(b) of Lemma~\ref{lem:identities-free-wilson} amounts to the trigonometric identity
\[
\sum_{\mu}\sin^2p_{\mu}+\left(\sum_{\mu}\left(1-\cos p_{\mu}\right)\right)^2=2\sum_{\mu}t_{\mu}+\sum_{\mu\ne\nu}t_{\mu}\cdot t_{\nu}\qquad\left(p\in\mathbb{R}^3\right),
\]
which follows from $\sin^2p_{\mu}=t_{\mu}\cdot\left(2-t_{\mu}\right)$.
\end{remark}

\subsubsection{Block means and a Poincar\'e inequality}\label{sec:block-means-poincare}

Let $E$ be a block (Convention~\ref{it:wilson-operators-4}) of side $\boldsymbol{\ell}$, and let $n:=\boldsymbol{\ell}/\eta$. Each line of $E$ parallel to a coordinate axis contains $n$ sites, and $E$ contains $n^3$ sites. Denote by $L^2(E)$ the Hilbert space of functions $E\to\mathbb{C}^4$ with the scalar product $\eta^3\sum_{x\in E}\overline{u(x)}\cdot v(x)$, whose norm is $\|\cdot\|_{2,E}$. Let $X\subset\eta\mathbb{Z}^3$ be a set that contains $E$, and let $u\colon X\to\mathbb{C}^4$. For $\mu\in\{1,2,3\}$ and $x\in X$ with $x+\eta e_{\mu}\in X$ put
\[
\nabla_{\mu}^{+}u(x):=\frac{u(x+\eta e_{\mu})-u(x)}{\eta};
\]
for $X=\eta\mathbb{Z}^3$ this is the value at $x$ of the function $\nabla_{\mu}^{+}u=(\tau_{\mu}u-u)/\eta$ defined above. Let $P_Eu$ be the mean of $u$ over $E$, taken as a constant function on $E$, and let $\|\nabla u\|_E\ge0$ measure the differences of $u$ along the bonds with both end points in $E$:
\[
(P_Eu)(x):=\frac{1}{n^3}\sum_{x'\in E}u(x')\quad(x\in E),\qquad
\|\nabla u\|_E^2:=\eta^3\sum_{\mu}\sum_{x\in E,\ x+\eta e_{\mu}\in E}\left|\nabla_{\mu}^{+}u(x)\right|^2.
\]

\begin{lemma}[Poincar\'e inequality for blocks]\label{lem:poincare-inequality-blocks}
Let $E$ be a block of side $\boldsymbol{\ell}$, let $X\subset\eta\mathbb{Z}^3$ be a set that contains $E$, and let $u\colon X\to\mathbb{C}^4$. Then
\[
\|u-P_Eu\|_{2,E}^2\le\frac{\boldsymbol{\ell}^2}{8}\,\|\nabla u\|_E^2.
\]
\end{lemma}

\begin{proof}
\emph{One dimension.} Let $g_1,\dots,g_n\in\mathbb{C}^4$, let $\bar{g}:=\frac{1}{n}\sum_{i=1}^ng_i$ be their mean, and let $\boldsymbol{\delta}_l:=g_{l+1}-g_l$ for $1\le l\le n-1$. Then
\begin{align*}
\sum_{i=1}^n|g_i-\bar{g}|^2
&=\frac{1}{2n}\sum_{i,j=1}^n|g_i-g_j|^2
=\frac{1}{n}\sum_{1\le i<j\le n}|g_i-g_j|^2\\
&\le\frac{1}{n}\sum_{1\le i<j\le n}(j-i)\sum_{l=i}^{j-1}|\boldsymbol{\delta}_l|^2
=\frac{1}{n}\sum_{l=1}^{n-1}|\boldsymbol{\delta}_l|^2\sum_{i=1}^l\sum_{j=l+1}^n(j-i).
\end{align*}
Here the first equality holds because
\[
\sum_{i,j=1}^n|g_i-g_j|^2=2n\sum_{i=1}^n|g_i|^2-2\left|\sum_{i=1}^ng_i\right|^2=2n\sum_{i=1}^n|g_i-\bar{g}|^2;
\]
the inequality is the Cauchy--Schwarz inequality for $g_j-g_i=\sum_{l=i}^{j-1}\boldsymbol{\delta}_l$, a sum of $j-i$ terms; and the last equality exchanges the order of summation. Since $l\cdot(n-l)\le n^2/4$,
\[
\sum_{i=1}^l\sum_{j=l+1}^n(j-i)=l\sum_{j=l+1}^nj-(n-l)\sum_{i=1}^li=\frac{l\cdot(n-l)\cdot n}{2}\le\frac{n^3}{8},
\]
and therefore
\[
\sum_{i=1}^n|g_i-\bar{g}|^2\le\frac{n^2}{8}\sum_{l=1}^{n-1}|\boldsymbol{\delta}_l|^2.
\]

\emph{Three dimensions.} Only the values of $u$ on $E$ enter, so $u$ may be regarded as an element of $L^2(E)$. For $\mu\in\{1,2,3\}$ let $A_{\mu}\colon L^2(E)\to L^2(E)$ replace a function by its means over the lines of $E$ parallel to $e_{\mu}$:
\[
(A_{\mu}u)(x):=\frac{1}{n}\sum_{k\in\mathbb{Z}:\ x+k\eta e_{\mu}\in E}u(x+k\eta e_{\mu})\qquad(x\in E).
\]
Each $A_{\mu}$ is the orthogonal projection of $L^2(E)$ onto the functions that are constant on every line of $E$ in the direction $\mu$. Since $E$ is a cube, the projections $A_1,A_2,A_3$ commute and $P_E=A_1A_2A_3$. Hence
\[
1-P_E=(1-A_1)+A_1(1-A_2)+A_1A_2(1-A_3),
\]
and the three summands have mutually orthogonal ranges: the first range lies in the kernel of $A_1$, the second in the range of $A_1$ and in the kernel of $A_2$, the third in the ranges of $A_1$ and of $A_2$. By Pythagoras and $\|A_{\mu}\|\le1$ it follows that
\[
\|u-P_Eu\|_{2,E}^2\le\sum_{\mu}\|(1-A_{\mu})u\|_{2,E}^2.
\]
The one-dimensional inequality, applied to the values of $u$ along each line of $E$ in the direction $\mu$ and summed over these lines, gives
\[
\|(1-A_{\mu})u\|_{2,E}^2\le\frac{n^2}{8}\cdot\eta^3\sum_{x\in E,\ x+\eta e_{\mu}\in E}|u(x+\eta e_{\mu})-u(x)|^2.
\]
Summing over $\mu$, writing $u(x+\eta e_{\mu})-u(x)=\eta\cdot\nabla_{\mu}^{+}u(x)$ by the definition of $\nabla_{\mu}^{+}u(x)$, and using $n\cdot\eta=\boldsymbol{\ell}$, one obtains
\[
\|u-P_Eu\|_{2,E}^2\le\frac{n^2\cdot\eta^2}{8}\,\|\nabla u\|_E^2=\frac{\boldsymbol{\ell}^2}{8}\,\|\nabla u\|_E^2.
\]
\end{proof}

\subsection{The Green function of the free Wilson operator}\label{sec:green-function-free}

For $p\in\mathbb{R}^3$, $|p|$ is the Euclidean norm, and for $x\in\mathbb{R}^3$, $|x|_{\infty}:=\max_{\mu}|x_{\mu}|$; a matrix carries its operator norm (see Convention~\ref{it:wilson-operators-5}). A multi-index is a triple $\boldsymbol{\beta}=(\boldsymbol{\beta}_1,\boldsymbol{\beta}_2,\boldsymbol{\beta}_3)$ of non-negative integers, with $|\boldsymbol{\beta}|:=\boldsymbol{\beta}_1+\boldsymbol{\beta}_2+\boldsymbol{\beta}_3$ and $\partial^{\boldsymbol{\beta}}:=\partial_1^{\boldsymbol{\beta}_1}\partial_2^{\boldsymbol{\beta}_2}\partial_3^{\boldsymbol{\beta}_3}$, where $\partial_{\mu}$ is the derivative with respect to $p_{\mu}$. One writes $\boldsymbol{\beta}'\le\boldsymbol{\beta}$ if $\boldsymbol{\beta}'_{\mu}\le\boldsymbol{\beta}_{\mu}$ for every $\mu$, and then $\binom{\boldsymbol{\beta}}{\boldsymbol{\beta}'}:=\prod_{\mu}\binom{\boldsymbol{\beta}_{\mu}}{\boldsymbol{\beta}'_{\mu}}$.

\begin{lemma}[Decay of Fourier coefficients]\label{lem:decay-fourier-coefficients}
Let $\mathsf{k}<3$ be real, let $N$ be an integer with $N>3-\mathsf{k}$, and let $A\ge0$. Let $a$ be a matrix valued function on $\mathbb{R}^3\setminus2\pi\mathbb{Z}^3$ which is $2\pi\mathbb{Z}^3$-periodic and of class $C^N$, and which satisfies
\[
\left|\partial^{\boldsymbol{\beta}}a(p)\right|\le A\cdot|p|^{-\mathsf{k}-|\boldsymbol{\beta}|}
\]
for all $p\in[-\pi,\pi]^3\setminus\{0\}$ and all multi-indices $\boldsymbol{\beta}$ with $|\boldsymbol{\beta}|\le N$. Then $a$ is integrable over $[-\pi,\pi]^3$, and its Fourier coefficients
\[
\hat a(x):=(2\pi)^{-3}\int_{[-\pi,\pi]^3}\mathrm{d}p\,\exp(i\,p\cdot x)\,a(p)\qquad(x\in\mathbb{Z}^3)
\]
satisfy
\[
|\hat a(x)|\le c(\mathsf{k},N)\cdot A\cdot(1+|x|_{\infty})^{-(3-\mathsf{k})}\qquad(x\in\mathbb{Z}^3),
\]
with a constant $c(\mathsf{k},N)$ that depends only on $\mathsf{k}$ and $N$.
\end{lemma}

\begin{proof}
The cube $[-\pi,\pi]^3$ lies in the ball of radius $\pi\sqrt{3}$ about the origin, and $\mathsf{k}<3$; therefore
\[
\int_{[-\pi,\pi]^3}\mathrm{d}p\,|p|^{-\mathsf{k}}\le\int_{\{|p|\le\pi\sqrt{3}\}}\mathrm{d}p\,|p|^{-\mathsf{k}}=\frac{4\pi\,(\pi\sqrt{3})^{3-\mathsf{k}}}{3-\mathsf{k}}.
\]
Since $|a(p)|\le A\cdot|p|^{-\mathsf{k}}$ on the cube, $a$ is integrable there, and $|\hat a(0)|$ is at most $(2\pi)^{-3}$ times the right side times $A$. This is the assertion for $x=0$.

Let now $x\in\mathbb{Z}^3\setminus\{0\}$. Choose a direction $\mu$ with $|x_{\mu}|=|x|_{\infty}$ and put $t:=1/|x|_{\infty}$, so that $0<t\le1$. Fix once and for all a function $\chi_0\in C^{\infty}(\mathbb{R})$ with $0\le\chi_0\le1$, $\chi_0=1$ on $(-\infty,1]$ and $\chi_0=0$ on $[2,\infty)$. On $[-\pi,\pi]^3$ let $\chi_t(p):=\chi_0(|p|/t)$, and extend $\chi_t$ to $\mathbb{R}^3$ by $2\pi\mathbb{Z}^3$-periodicity; the extension is smooth, because $\chi_t$ vanishes for $|p|\ge2t$ and $2t\le2<\pi$. The function $\chi_1=\chi_0(|\cdot|)$ is smooth on $\mathbb{R}^3$, as it is identically $1$ near the origin, and $\chi_t(p)=\chi_1(p/t)$ on the cube. Consequently there are constants $c_l$, depending only on $l$, such that
\[
\left|\partial_{\mu}^l\chi_t\right|\le c_l\cdot t^{-l}\qquad(l\ge0),
\]
and for $l\ge1$ the function $\partial_{\mu}^l\chi_t$ vanishes on the cube outside the shell $\{t\le|p|\le2t\}$.

One has $\hat a(x)=(2\pi)^{-3}(J_1+J_2)$ with
\[
J_1:=\int_{[-\pi,\pi]^3}\mathrm{d}p\,\exp(i\,p\cdot x)\,a(p)\,\chi_t(p),\qquad J_2:=\int_{[-\pi,\pi]^3}\mathrm{d}p\,\exp(i\,p\cdot x)\,a(p)\left(1-\chi_t(p)\right).
\]
As $0\le\chi_t\le1$ and $\chi_t$ vanishes on the cube outside the ball $\{|p|\le2t\}$,
\[
|J_1|\le A\int_{\{|p|\le2t\}}\mathrm{d}p\,|p|^{-\mathsf{k}}=\frac{4\pi A\cdot(2t)^{3-\mathsf{k}}}{3-\mathsf{k}}.
\]

The function $F:=a\cdot(1-\chi_t)$ is $2\pi\mathbb{Z}^3$-periodic and of class $C^N$ on all of $\mathbb{R}^3$, since $1-\chi_t$ vanishes in the balls of radius $t$ about the points of $2\pi\mathbb{Z}^3$. As $x\in\mathbb{Z}^3$, the function $p\mapsto\exp(i\,p\cdot x)$ is $2\pi\mathbb{Z}^3$-periodic as well, and $\exp(i\,p\cdot x)=(i x_{\mu})^{-1}\,\partial_{\mu}\exp(i\,p\cdot x)$. Integrating by parts $N$ times in the variable $p_{\mu}$ over one period, where the boundary terms cancel by periodicity, one obtains
\[
|J_2|=|x_{\mu}|^{-N}\left|\int_{[-\pi,\pi]^3}\mathrm{d}p\,\exp(i\,p\cdot x)\,\partial_{\mu}^NF(p)\right|\le|x_{\mu}|^{-N}\int_{[-\pi,\pi]^3}\mathrm{d}p\,\left|\partial_{\mu}^NF(p)\right|.
\]
By the Leibniz rule,
\[
\partial_{\mu}^NF=\left(\partial_{\mu}^Na\right)(1-\chi_t)-\sum_{l=1}^N\binom{N}{l}\left(\partial_{\mu}^{N-l}a\right)\partial_{\mu}^l\chi_t.
\]
On the cube the first term vanishes where $|p|<t$ and is bounded by $A\cdot|p|^{-\mathsf{k}-N}$ elsewhere; the term of index $l$ vanishes outside the shell $\{t\le|p|\le2t\}$ and is bounded there by $\binom{N}{l}c_l\cdot t^{-l}\cdot A\cdot|p|^{-\mathsf{k}-N+l}$. Hence
\[
\int_{[-\pi,\pi]^3}\mathrm{d}p\,\left|\partial_{\mu}^NF\right|\le A\int_{\{p\in\mathbb{R}^3:\,|p|\ge t\}}\mathrm{d}p\,|p|^{-\mathsf{k}-N}+A\sum_{l=1}^N\binom{N}{l}c_l\cdot t^{-l}\int_{\{t\le|p|\le2t\}}\mathrm{d}p\,|p|^{-\mathsf{k}-N+l}.
\]
The first integral extends over the exterior of a ball in $\mathbb{R}^3$; since $N+\mathsf{k}>3$, it equals $4\pi\,t^{3-\mathsf{k}-N}/(N+\mathsf{k}-3)$. With $k_l:=2-\mathsf{k}-N+l$, the integral of index $l$ equals
\[
4\pi\int_t^{2t}\mathrm{d}r\,r^{k_l}\le4\pi\cdot2^{|k_l|}\cdot t^{3-\mathsf{k}-N+l}.
\]
It follows that
\[
\int_{[-\pi,\pi]^3}\mathrm{d}p\,\left|\partial_{\mu}^NF\right|\le c'(\mathsf{k},N)\cdot A\cdot t^{3-\mathsf{k}-N}
\]
with a constant $c'(\mathsf{k},N)$ depending only on $\mathsf{k}$ and $N$. Since $|x_{\mu}|^{-N}=t^N$, this gives $|J_2|\le c'(\mathsf{k},N)\cdot A\cdot t^{3-\mathsf{k}}$.

Finally $|x|_{\infty}\ge1$, so that $t=1/|x|_{\infty}\le2/(1+|x|_{\infty})$, and $3-\mathsf{k}>0$. Therefore
\[
|\hat a(x)|\le(2\pi)^{-3}\left(|J_1|+|J_2|\right)\le(2\pi)^{-3}\left(\frac{4\pi\cdot2^{3-\mathsf{k}}}{3-\mathsf{k}}+c'(\mathsf{k},N)\right)A\cdot t^{3-\mathsf{k}},
\]
and $t^{3-\mathsf{k}}\le2^{3-\mathsf{k}}\,(1+|x|_{\infty})^{-(3-\mathsf{k})}$. Together with the case $x=0$ this proves the lemma, $c(\mathsf{k},N)$ being the larger of the two constants obtained.
\end{proof}

\begin{lemma}[Inverse of the symbol at unit mesh]\label{lem:inverse-symbol-unit}
Let $0\le m\le1$. For $p\in\mathbb{R}^3$ put $s_{\mu}:=\sin p_{\mu}$, $w:=\sum_{\mu}(1-\cos p_{\mu})$,
\[
\hat{D}(p):=i\sum_{\mu}\gamma_{\mu}s_{\mu}+w,
\]
and
\[
n(p):=-i\sum_{\mu}\gamma_{\mu}s_{\mu}+w+m,\qquad N(p):=\sum_{\mu}s_{\mu}^2+(w+m)^2,\qquad a_m(p):=\frac{n(p)}{N(p)}.
\]
At mesh $\eta=1$ the operator $D$ of~\eqref{eq:wilson-operators-7} maps the function $x\mapsto\exp(i\,p\cdot x)\,v$, $v\in\mathbb{C}^4$, to $x\mapsto\exp(i\,p\cdot x)\cdot\hat{D}(p)\cdot v$. Wherever $N(p)\ne0$,
\[
\left(\hat{D}(p)+m\right)a_m(p)=a_m(p)\left(\hat{D}(p)+m\right)=1.
\]
On $[-\pi,\pi]^3$ one has
\[
N(p)\ge\frac{4}{\pi^2}\left(|p|^2+m^2\right).
\]
Moreover, there is a family $(A_{\boldsymbol{\beta}})_{\boldsymbol{\beta}}$ of absolute constants, indexed by the multi-indices and depending neither on $m\in[0,1]$ nor on the choice of the matrices $\gamma_{\mu}$ subject to Convention~\ref{it:wilson-operators-3}, such that for all multi-indices $\boldsymbol{\beta}$
\[
\left|\partial^{\boldsymbol{\beta}}a_m(p)\right|\le A_{\boldsymbol{\beta}}\cdot|p|^{-1-|\boldsymbol{\beta}|}\qquad\left(p\in[-\pi,\pi]^3\setminus\{0\}\right).
\]
\end{lemma}

\begin{proof}
Write $\gamma\cdot s:=\sum_{\mu}\gamma_{\mu}s_{\mu}$, so that $\hat{D}(p)=i\,\gamma\cdot s+w$ and $n(p)=-i\,\gamma\cdot s+w+m$.

At mesh one, $\nabla^s_{\mu}$ multiplies the function $x\mapsto\exp(i\,p\cdot x)$ by $\frac{1}{2}\left(\exp(i\,p_{\mu})-\exp(-i\,p_{\mu})\right)=i\,s_{\mu}$, and $\Delta_{\mu}$ multiplies it by $2\cos p_{\mu}-2$; so $-\frac{1}{2}\Delta$ multiplies it by $w$. The first form of~\eqref{eq:wilson-operators-7} therefore maps $x\mapsto\exp(i\,p\cdot x)\,v$ to $x\mapsto\exp(i\,p\cdot x)\,(i\,\gamma\cdot s+w)\,v$.

By the anticommutation relations of Convention~\ref{it:wilson-operators-3}, $(\gamma\cdot s)^2=\sum_{\mu}s_{\mu}^2$ times the unit matrix. The matrices $i\,\gamma\cdot s+w+m$ and $-i\,\gamma\cdot s+w+m$ commute, and
\[
\left(i\,\gamma\cdot s+w+m\right)\left(-i\,\gamma\cdot s+w+m\right)=(\gamma\cdot s)^2+(w+m)^2=N(p).
\]
Thus $(\hat{D}(p)+m)\,n(p)=n(p)\,(\hat{D}(p)+m)=N(p)$, and division by $N(p)$, where it does not vanish, gives the two identities for $a_m(p)$.

Let $t_{\mu}:=1-\cos p_{\mu}\ge0$. By Remark~\ref{rem:wilson-operators-1},
\[
\sum_{\mu}s_{\mu}^2+w^2=2w+\sum_{\mu\ne\nu}t_{\mu}t_{\nu}\ge2w,
\]
so that $N(p)=\sum_{\mu}s_{\mu}^2+w^2+2wm+m^2\ge2w+m^2$. For $|\theta|\le\pi$ one has $1-\cos\theta=2\sin^2(\theta/2)\ge2\theta^2/\pi^2$, because $\sin y\ge2y/\pi$ for $0\le y\le\pi/2$ (concavity of the sine). Hence $2w\ge4|p|^2/\pi^2$ on $[-\pi,\pi]^3$, and together with $m^2\ge4m^2/\pi^2$ this gives the lower bound for $N$. In particular $N>0$ on $[-\pi,\pi]^3\setminus\{0\}$.

Put $\mathsf{r}:=\left(|p|^2+m^2\right)^{\frac{1}{2}}$. On the cube $\mathsf{r}\le\mathsf{r}_{\mathrm{max}}:=\left(3\pi^2+1\right)^{\frac{1}{2}}$. In the rest of the proof $c$ denotes an absolute constant, independent of $p$, of $m\in[0,1]$ and of the matrices $\gamma_{\mu}$, and $c_{\boldsymbol{\beta}}$ such a constant that may depend on $\boldsymbol{\beta}$; their values may change from one occurrence to the next. All bounds below hold on $[-\pi,\pi]^3\setminus\{0\}$.

\emph{The numerator.} The matrix $\gamma\cdot s$ is Hermitian with square $|s|^2$, so its norm is $|s|\le|p|$; moreover $0\le w\le|p|^2/2$. Hence $|n|\le|p|+|p|^2/2+m\le c\cdot\mathsf{r}$, as $|p|\le\pi\sqrt{3}$. The function $n-m$ is the sum over $\mu$ of $-i\gamma_{\mu}\sin p_{\mu}+1-\cos p_{\mu}$, a function of $p_{\mu}$ alone; therefore a derivative of $n$ of order $k\ge1$ either vanishes or equals $-i\gamma_{\mu}\,\partial_{\mu}^k\sin p_{\mu}-\partial_{\mu}^k\cos p_{\mu}$ for some $\mu$, and its norm is at most $2$. For $|\boldsymbol{\beta}|\ge1$ one has $2\le2\,\mathsf{r}_{\mathrm{max}}^{|\boldsymbol{\beta}|-1}\cdot\mathsf{r}^{1-|\boldsymbol{\beta}|}$. Consequently $\left|\partial^{\boldsymbol{\beta}}n\right|\le c_{\boldsymbol{\beta}}\cdot\mathsf{r}^{1-|\boldsymbol{\beta}|}$ for every $\boldsymbol{\beta}$.

\emph{The denominator.} One has $N\le|p|^2+(|p|^2/2+m)^2\le c\cdot\mathsf{r}^2$. Next,
\[
\partial_{\mu}N=2s_{\mu}\cos p_{\mu}+2(w+m)\,s_{\mu},
\]
and $|s_{\mu}|\le|p|$, $0\le w+m\le7$; so $|\partial_{\mu}N|\le c\cdot|p|\le c\cdot\mathsf{r}$. Finally, $N$ is a polynomial in $m$ and in the sines and cosines of $p_1,p_2,p_3$; as $0\le m\le1$, each of its derivatives of order at least two is bounded by an absolute constant $c_{\boldsymbol{\beta}}$, and $c_{\boldsymbol{\beta}}\le c_{\boldsymbol{\beta}}\cdot\mathsf{r}_{\mathrm{max}}^{|\boldsymbol{\beta}|-2}\cdot\mathsf{r}^{2-|\boldsymbol{\beta}|}$ for $|\boldsymbol{\beta}|\ge2$. Consequently $\left|\partial^{\boldsymbol{\beta}}N\right|\le c_{\boldsymbol{\beta}}\cdot\mathsf{r}^{2-|\boldsymbol{\beta}|}$ for every $\boldsymbol{\beta}$.

\emph{The reciprocal.} It is shown by induction on $|\boldsymbol{\beta}|$ that $\left|\partial^{\boldsymbol{\beta}}(N^{-1})\right|\le c_{\boldsymbol{\beta}}\cdot\mathsf{r}^{-2-|\boldsymbol{\beta}|}$. The case $\boldsymbol{\beta}=0$, with which the induction begins, follows from the inequality $N\ge(4/\pi^2)\cdot\mathsf{r}^2$ shown above: it gives $N^{-1}\le(\pi^2/4)\cdot\mathsf{r}^{-2}$. For $|\boldsymbol{\beta}|\ge1$, applying $\partial^{\boldsymbol{\beta}}$ to the identity $N\cdot N^{-1}=1$ gives
\[
\partial^{\boldsymbol{\beta}}(N^{-1})=-N^{-1}\sum_{0<\boldsymbol{\beta}'\le\boldsymbol{\beta}}\binom{\boldsymbol{\beta}}{\boldsymbol{\beta}'}\,\partial^{\boldsymbol{\beta}'}N\cdot\partial^{\boldsymbol{\beta}-\boldsymbol{\beta}'}(N^{-1}),
\]
and by the bounds for $N$ and the induction hypothesis the term of index $\boldsymbol{\beta}'$ is bounded by
\[
c\cdot\mathsf{r}^{-2}\cdot\mathsf{r}^{2-|\boldsymbol{\beta}'|}\cdot\mathsf{r}^{-2-|\boldsymbol{\beta}|+|\boldsymbol{\beta}'|}=c\cdot\mathsf{r}^{-2-|\boldsymbol{\beta}|}.
\]

\emph{The quotient.} By the Leibniz rule,
\[
\left|\partial^{\boldsymbol{\beta}}a_m\right|\le\sum_{\boldsymbol{\beta}'\le\boldsymbol{\beta}}\binom{\boldsymbol{\beta}}{\boldsymbol{\beta}'}\left|\partial^{\boldsymbol{\beta}'}n\right|\cdot\left|\partial^{\boldsymbol{\beta}-\boldsymbol{\beta}'}(N^{-1})\right|\le c_{\boldsymbol{\beta}}\cdot\mathsf{r}^{-1-|\boldsymbol{\beta}|}\le c_{\boldsymbol{\beta}}\cdot|p|^{-1-|\boldsymbol{\beta}|},
\]
the last step because $\mathsf{r}\ge|p|$. This is the assertion, with $A_{\boldsymbol{\beta}}:=c_{\boldsymbol{\beta}}$.
\end{proof}

\begin{proposition}[Lattice Green function]\label{prop:lattice-green-function}
One can choose an absolute constant $C_G$ such that the assertions below hold for every mesh $\eta>0$ and every mass $\mathfrak{m}$ with $0\le\mathfrak{m}\le1/\eta$. For such $\eta$ and $\mathfrak{m}$, let $a_{\mathfrak{m}\cdot\eta}$ be the function $a_m$ of Lemma~\ref{lem:inverse-symbol-unit} with $m=\mathfrak{m}\cdot\eta\in[0,1]$, and let $\widehat{a_{\mathfrak{m}\cdot\eta}}$ denote its Fourier coefficients in the sense of Lemma~\ref{lem:decay-fourier-coefficients}. Define $G=G^{\eta}_{\mathfrak{m}}$ on $\eta\mathbb{Z}^3$, with values in the complex $4\times4$ matrices, by
\[
G(x):=G^{\eta}_{\mathfrak{m}}(x):=\eta^{-2}\,\widehat{a_{\mathfrak{m}\cdot\eta}}(x/\eta)\qquad(x\in\eta\mathbb{Z}^3).
\]
\begin{enumerate}
\item[(a)] Let $D$ act on each column of $G$. Then $\left((D+\mathfrak{m})G\right)(x)$ is $\eta^{-3}$ times the unit matrix for $x=0$, and vanishes for $x\in\eta\mathbb{Z}^3\setminus\{0\}$.
\item[(b)] For every $x\in\eta\mathbb{Z}^3$ and every direction $\mu$,
\[
|G(x)|\le C_G\cdot(\eta+|x|_{\infty})^{-2},\qquad\left|\nabla^{\pm}_{\mu}G(x)\right|\le C_G\cdot(\eta+|x|_{\infty})^{-3}.
\]
\item[(c)] Every finitely supported $v\colon\eta\mathbb{Z}^3\to\mathbb{C}^4$ satisfies
\[
v(x)=\eta^3\sum_{x'\in\eta\mathbb{Z}^3}G(x-x')\left((D+\mathfrak{m})v\right)(x')\qquad(x\in\eta\mathbb{Z}^3).
\]
\end{enumerate}
\end{proposition}

\begin{proof}
\emph{Reduction to unit mesh.} For a mesh $\eta>0$ denote by $D^{(\eta)}$ and $\nabla^{\pm,(\eta)}_{\mu}$ the operators $D$ of~\eqref{eq:wilson-operators-7} and $\nabla^{\pm}_{\mu}$ of Section~\ref{sec:shifts-differences-free}, acting on functions on $\eta\mathbb{Z}^3$. For a function $g$ on $\mathbb{Z}^3$ let $g_{\eta}(x):=g(x/\eta)$, $x\in\eta\mathbb{Z}^3$. Under $x\mapsto x/\eta$ the shift by $\eta e_{\mu}$ on $\eta\mathbb{Z}^3$ becomes the shift by $e_{\mu}$ on $\mathbb{Z}^3$; therefore the second form of~\eqref{eq:wilson-operators-7} and the definitions of $\nabla^{\pm}_{\mu}$ give
\[
D^{(\eta)}g_{\eta}=\eta^{-1}\left(D^{(1)}g\right)_{\eta},\qquad\nabla^{\pm,(\eta)}_{\mu}g_{\eta}=\eta^{-1}\left(\nabla^{\pm,(1)}_{\mu}g\right)_{\eta}
\]
(see Convention~\ref{it:wilson-operators-1}). Let $m:=\mathfrak{m}\cdot\eta\in[0,1]$ and $G_1:=\widehat{a_m}$ on $\mathbb{Z}^3$, so that $G^{\eta}_{\mathfrak{m}}=\eta^{-2}\,(G_1)_{\eta}$, and suppose that (a)--(c) hold for the mesh $1$, the mass $m$ and the function $G_1$, with the constant $C_G$. Since $\mathfrak{m}=\eta^{-1}m$,
\[
\left(D^{(\eta)}+\mathfrak{m}\right)G^{\eta}_{\mathfrak{m}}=\eta^{-3}\left(\left(D^{(1)}+m\right)G_1\right)_{\eta},
\]
and (a) follows. Next, $|G^{\eta}_{\mathfrak{m}}(x)|=\eta^{-2}\,|G_1(x/\eta)|$ and $\nabla^{\pm,(\eta)}_{\mu}G^{\eta}_{\mathfrak{m}}=\eta^{-3}\left(\nabla^{\pm,(1)}_{\mu}G_1\right)_{\eta}$; as $\eta^{-k}\left(1+|x|_{\infty}/\eta\right)^{-k}=(\eta+|x|_{\infty})^{-k}$ for $k=2,3$, (b) follows with the same constant $C_G$. Finally, let $v$ be finitely supported on $\eta\mathbb{Z}^3$ and put $v_1(y):=v(\eta y)$ for $y\in\mathbb{Z}^3$; then $v=(v_1)_{\eta}$ and $(D^{(\eta)}+\mathfrak{m})v=\eta^{-1}\left((D^{(1)}+m)v_1\right)_{\eta}$. Substituting $x'=\eta y'$ one obtains
\begin{align*}
\eta^3\sum_{x'\in\eta\mathbb{Z}^3}G^{\eta}_{\mathfrak{m}}(x-x')\left(\left(D^{(\eta)}+\mathfrak{m}\right)v\right)(x')&=\sum_{y'\in\mathbb{Z}^3}G_1(x/\eta-y')\left(\left(D^{(1)}+m\right)v_1\right)(y')\\
&=v_1(x/\eta)=v(x),
\end{align*}
which is (c). It remains to prove (a)--(c) for $\eta=1$, $0\le m\le1$ and $G=\widehat{a_m}$. In the rest of the proof $D$ and $\nabla^{\pm}_{\mu}$ are taken at mesh one.

\emph{(a)} The functions $n$ and $N$ of Lemma~\ref{lem:inverse-symbol-unit} are trigonometric polynomials, and $N>0$ on $[-\pi,\pi]^3\setminus\{0\}$ by the lower bound of that lemma, hence on $\mathbb{R}^3\setminus2\pi\mathbb{Z}^3$ by periodicity. Thus $a_m$ is $2\pi\mathbb{Z}^3$-periodic and smooth on $\mathbb{R}^3\setminus2\pi\mathbb{Z}^3$, and by Lemma~\ref{lem:inverse-symbol-unit}
\[
\left|\partial^{\boldsymbol{\beta}}a_m(p)\right|\le A\cdot|p|^{-1-|\boldsymbol{\beta}|}\quad\text{for }p\in[-\pi,\pi]^3\setminus\{0\}\text{ and }|\boldsymbol{\beta}|\le3,\qquad\text{where }A:=\max_{|\boldsymbol{\beta}|\le3}A_{\boldsymbol{\beta}}.
\]
So Lemma~\ref{lem:decay-fourier-coefficients} applies to $a_m$ with $\mathsf{k}=1$ and with the integer $3$ in the place of its $N$ (note that $3>3-1$); in particular $a_m$ is integrable, and $G$ is well defined. The operator $D+m$ is a finite linear combination of shifts with constant matrix coefficients; hence it may be applied, column by column, under the integral sign. By Lemma~\ref{lem:inverse-symbol-unit} it maps $x\mapsto\exp(i\,p\cdot x)\,a_m(p)$ to $x\mapsto\exp(i\,p\cdot x)\left(\hat{D}(p)+m\right)a_m(p)$, and this is $\exp(i\,p\cdot x)$ times the unit matrix for $p\ne0$. Therefore, for $x\in\mathbb{Z}^3$,
\[
\left((D+m)G\right)(x)=(2\pi)^{-3}\int_{[-\pi,\pi]^3}\mathrm{d}p\,\exp(i\,p\cdot x)=1_{\{0\}}(x)
\]
times the unit matrix. This is (a).

\emph{(b)} Lemma~\ref{lem:decay-fourier-coefficients}, applied as in (a), gives $|G(x)|\le c(1,3)\cdot A\cdot(1+|x|_{\infty})^{-2}$. Next, for a direction $\mu$,
\[
\nabla^{+}_{\mu}G(x)=G(x+e_{\mu})-G(x),\qquad\nabla^{-}_{\mu}G(x)=G(x)-G(x-e_{\mu}),
\]
so that $\nabla^{\pm}_{\mu}G=\widehat{b_{\pm}}$, where
\[
b_{\pm}(p):=\pm f_{\pm}(p)\,a_m(p),\qquad f_{\pm}(p):=\exp(\pm i\,p_{\mu})-1.
\]
On $[-\pi,\pi]^3$ one has $|f_{\pm}(p)|=2\,|\sin(p_{\mu}/2)|\le|p|$; for $l\ge1$ the derivative $\partial_{\mu}^lf_{\pm}=(\pm i)^l\exp(\pm i\,p_{\mu})$ has modulus $1\le(\pi\sqrt{3})^{l-1}|p|^{1-l}$; and all other derivatives of $f_{\pm}$ vanish. Hence $\left|\partial^{\boldsymbol{\beta}'}f_{\pm}(p)\right|\le(\pi\sqrt{3})^{|\boldsymbol{\beta}'|}\cdot|p|^{1-|\boldsymbol{\beta}'|}$ for every multi-index $\boldsymbol{\beta}'$, and the Leibniz rule together with Lemma~\ref{lem:inverse-symbol-unit} gives, for $p\in[-\pi,\pi]^3\setminus\{0\}$ and $|\boldsymbol{\beta}|\le4$,
\[
\left|\partial^{\boldsymbol{\beta}}b_{\pm}(p)\right|\le\sum_{\boldsymbol{\beta}'\le\boldsymbol{\beta}}\binom{\boldsymbol{\beta}}{\boldsymbol{\beta}'}(\pi\sqrt{3})^{|\boldsymbol{\beta}'|}A_{\boldsymbol{\beta}-\boldsymbol{\beta}'}\cdot|p|^{-|\boldsymbol{\beta}|}\le A'\cdot|p|^{-|\boldsymbol{\beta}|},
\]
with $A':=\max_{|\boldsymbol{\beta}|\le4}\sum_{\boldsymbol{\beta}'\le\boldsymbol{\beta}}\binom{\boldsymbol{\beta}}{\boldsymbol{\beta}'}(\pi\sqrt{3})^{|\boldsymbol{\beta}'|}A_{\boldsymbol{\beta}-\boldsymbol{\beta}'}$. As $b_{\pm}$ is $2\pi\mathbb{Z}^3$-periodic and smooth on $\mathbb{R}^3\setminus2\pi\mathbb{Z}^3$, Lemma~\ref{lem:decay-fourier-coefficients} applies with $\mathsf{k}=0$ and the integer $4$ (note that $4>3$), and gives $\left|\nabla^{\pm}_{\mu}G(x)\right|\le c(0,4)\cdot A'\cdot(1+|x|_{\infty})^{-3}$. Thus (b) holds with $C_G:=\max\{c(1,3)\cdot A,\ c(0,4)\cdot A'\}$, which is an absolute constant because the $A_{\boldsymbol{\beta}}$ are.

\emph{(c)} Let $v\colon\mathbb{Z}^3\to\mathbb{C}^4$ be finitely supported, and let $g:=(D+m)v$; it is finitely supported as well. Put
\[
\hat{v}(p):=\sum_{x\in\mathbb{Z}^3}\exp(-i\,p\cdot x)\,v(x),\qquad\hat{g}(p):=\sum_{x\in\mathbb{Z}^3}\exp(-i\,p\cdot x)\,g(x).
\]
These are trigonometric polynomials. Replacing $v$ by $v(\cdot\pm e_{\mu})$ multiplies $\hat{v}$ by $\exp(\pm i\,p_{\mu})$, which is also the factor by which the same shift multiplies the function $x\mapsto\exp(i\,p\cdot x)$. Since $D+m$ is a finite linear combination of shifts with constant matrix coefficients, Lemma~\ref{lem:inverse-symbol-unit} gives $\hat{g}=\left(\hat{D}+m\right)\cdot\hat{v}$, and the identity $a_m\left(\hat{D}+m\right)=1$ of that lemma gives $a_m\cdot\hat{g}=\hat{v}$ for $p\ne0$. As the sum over $x'$ is finite, it follows that
\begin{align*}
\sum_{x'\in\mathbb{Z}^3}G(x-x')\,g(x')&=(2\pi)^{-3}\int_{[-\pi,\pi]^3}\mathrm{d}p\,\exp(i\,p\cdot x)\,a_m(p)\cdot\hat{g}(p)\\
&=(2\pi)^{-3}\int_{[-\pi,\pi]^3}\mathrm{d}p\,\exp(i\,p\cdot x)\cdot\hat{v}(p)=v(x),
\end{align*}
the last equality by the orthogonality relations $(2\pi)^{-3}\int_{[-\pi,\pi]^3}\mathrm{d}p\,\exp(i\,p\cdot y)=1_{\{0\}}(y)$, $y\in\mathbb{Z}^3$.
\end{proof}

\begin{lemma}[Two lattice sums; a H\"older bound for $G$]\label{lem:two-lattice-sums}
Let $\eta>0$, let $\boldsymbol{\alpha}\in(0,1)$ be the H\"older order fixed in~\eqref{eq:wilson-operators-1}, and let $\varrho\ge\eta$.
\begin{enumerate}
\item[(a)] With sums over $z\in\eta\mathbb{Z}^3$,
\[
\sum_{|z|_{\infty}\le\varrho}\eta^3(\eta+|z|_{\infty})^{-2}\le27\varrho,\qquad\sum_{|z|_{\infty}\le\varrho}\eta^3(\eta+|z|_{\infty})^{-2-\boldsymbol{\alpha}}\le\frac{27\,\varrho^{1-\boldsymbol{\alpha}}}{1-\boldsymbol{\alpha}}.
\]
\item[(b)] Let $0\le\mathfrak{m}\le1/\eta$ and $G=G^{\eta}_{\mathfrak{m}}$ as in Proposition~\ref{prop:lattice-green-function}. Let $x,\tilde x,x'\in\eta\mathbb{Z}^3$ with $h:=|x-\tilde x|_{\infty}>0$. Then
\[
\left|G(x-x')-G(\tilde x-x')\right|\le24\,C_G\cdot h^{\boldsymbol{\alpha}}\left((\eta+|x-x'|_{\infty})^{-2-\boldsymbol{\alpha}}+(\eta+|\tilde x-x'|_{\infty})^{-2-\boldsymbol{\alpha}}\right).
\]
\end{enumerate}
\end{lemma}

\begin{proof}
\emph{(a)} For every integer $n\ge1$,
\[
\#\left|\left\{z\in\eta\mathbb{Z}^3:|z|_{\infty}=n\eta\right\}\right|=(2n+1)^3-(2n-1)^3=24n^2+2\le26n^2.
\]
Let $N:=\lfloor\varrho/\eta\rfloor$ be the integer part of $\varrho/\eta$. Then $N\ge1$, and the sites $z$ with $|z|_{\infty}\le\varrho$ are those with $|z|_{\infty}=n\eta$ for an integer $n$ with $0\le n\le N$. Let $q\in\{2,2+\boldsymbol{\alpha}\}$. The site $z=0$ contributes $\eta^{3-q}$, and $(\eta+n\eta)^{-q}\le(n\eta)^{-q}$ for $n\ge1$; hence
\[
\sum_{|z|_{\infty}\le\varrho}\eta^3(\eta+|z|_{\infty})^{-q}\le\eta^{3-q}\left(1+26\sum_{n=1}^Nn^{2-q}\right).
\]
For $q=2$ the right side equals $\eta\,(1+26N)\le27\,\eta N\le27\varrho$. For $q=2+\boldsymbol{\alpha}$, the inequality $n^{-\boldsymbol{\alpha}}\le\int_{n-1}^n\mathrm{d}r\,r^{-\boldsymbol{\alpha}}$ gives $\sum_{n=1}^Nn^{-\boldsymbol{\alpha}}\le N^{1-\boldsymbol{\alpha}}/(1-\boldsymbol{\alpha})$; since moreover $1\le N^{1-\boldsymbol{\alpha}}/(1-\boldsymbol{\alpha})$, the right side is at most $27\,(\eta N)^{1-\boldsymbol{\alpha}}/(1-\boldsymbol{\alpha})\le27\,\varrho^{1-\boldsymbol{\alpha}}/(1-\boldsymbol{\alpha})$.

\emph{(b)} Put $\mathsf{r}:=\eta+\min\{|x-x'|_{\infty},|\tilde x-x'|_{\infty}\}$. Then $\eta+|y-x'|_{\infty}\ge\mathsf{r}$ for $y\in\{x,\tilde x\}$, with equality for at least one of the two points.

Let first $h\ge\mathsf{r}/2$. By the triangle inequality the larger of the two numbers $\eta+|x-x'|_{\infty}$ and $\eta+|\tilde x-x'|_{\infty}$ is at most $\mathsf{r}+h$; so both are at most $\mathsf{r}+h\le3h$. For $y\in\{x,\tilde x\}$, Proposition~\ref{prop:lattice-green-function}(b) gives
\[
|G(y-x')|\le C_G\cdot(\eta+|y-x'|_{\infty})^{-2-\boldsymbol{\alpha}}\cdot(\eta+|y-x'|_{\infty})^{\boldsymbol{\alpha}}\le3^{\boldsymbol{\alpha}}\,C_G\cdot h^{\boldsymbol{\alpha}}\cdot(\eta+|y-x'|_{\infty})^{-2-\boldsymbol{\alpha}}.
\]
Since $|G(x-x')-G(\tilde x-x')|\le|G(x-x')|+|G(\tilde x-x')|$ and $3^{\boldsymbol{\alpha}}\le24$, the assertion follows.

Let now $h<\mathsf{r}/2$. A monotone coordinate path from $x$ to $\tilde x$ is a finite sequence of sites of $\eta\mathbb{Z}^3$, beginning at $x$ and ending at $\tilde x$, in which consecutive sites differ by $\pm\eta e_{\mu}$ for some $\mu$, and along which each coordinate changes monotonically. Choose such a path. It has $\sum_{\mu}|x_{\mu}-\tilde x_{\mu}|/\eta\le3h/\eta$ bonds, and every coordinate of each of its sites $z$ lies between the corresponding coordinates of $x$ and $\tilde x$; hence $|z-x|_{\infty}\le h$, $|z-\tilde x|_{\infty}\le h$, and
\[
\eta+|z-x'|_{\infty}\ge\eta+|x-x'|_{\infty}-|z-x|_{\infty}\ge\mathsf{r}-h>\frac{\mathsf{r}}{2}.
\]
For a bond of the path with end points $z$ and $z+\eta e_{\mu}$, Proposition~\ref{prop:lattice-green-function}(b) gives
\[
\left|G(z+\eta e_{\mu}-x')-G(z-x')\right|=\eta\cdot\left|\nabla^{+}_{\mu}G(z-x')\right|\le\eta\,C_G\cdot(\eta+|z-x'|_{\infty})^{-3}\le\eta\,C_G\cdot\left(\frac{\mathsf{r}}{2}\right)^{-3}.
\]
Writing $G(x-x')-G(\tilde x-x')$ as a telescoping sum over the bonds of the path, one obtains
\[
\left|G(x-x')-G(\tilde x-x')\right|\le\frac{3h}{\eta}\cdot\eta\,C_G\cdot\left(\frac{\mathsf{r}}{2}\right)^{-3}=24\,C_G\cdot h\cdot\mathsf{r}^{-3}=24\,C_G\cdot h^{\boldsymbol{\alpha}}\cdot\mathsf{r}^{-2-\boldsymbol{\alpha}}\cdot\left(\frac{h}{\mathsf{r}}\right)^{1-\boldsymbol{\alpha}}.
\]
Here $(h/\mathsf{r})^{1-\boldsymbol{\alpha}}\le1$, because $h<\mathsf{r}/2$ and $\boldsymbol{\alpha}<1$; and $\mathsf{r}^{-2-\boldsymbol{\alpha}}$ is one of the two terms $(\eta+|x-x'|_{\infty})^{-2-\boldsymbol{\alpha}}$ and $(\eta+|\tilde x-x'|_{\infty})^{-2-\boldsymbol{\alpha}}$. This proves (b).
\end{proof}

\subsection{The class of blocked Wilson problems}\label{sec:class-blocked-wilson}

\subsubsection{Definition and standard penalties}\label{sec:definition-standard-penalties}

\begin{definition}[Blocked Wilson problems]\label{def:blocked-wilson-problems}
Let $s\in[6L,\infty]$ and $c_1,C_1,C_2>0$. Consider a tuple
\[
\mathbb{P}=(\eta,R,\boldsymbol{\theta};\gamma;W,\mathbb{E},\operatorname{Pen},\mathfrak{m})
\]
formed by a mesh $\eta>0$, a side $R\in\eta\mathbb{N}\cup\{\infty\}$ and a character $\boldsymbol{\theta}\in\mathbb{R}^3$ as in Convention~\ref{it:wilson-operators-2}, Clifford data $\gamma=(\gamma_1,\gamma_2,\gamma_3)$ as in Convention~\ref{it:wilson-operators-3}, a set $W\subset\eta\mathbb{Z}^3$, a set $\mathbb{E}$ of blocks (Convention~\ref{it:wilson-operators-4}), an operator $\operatorname{Pen}$ on the functions $\eta\mathbb{Z}^3\to\mathbb{C}^4$, and a real number $\mathfrak{m}$. It is called a \emph{blocked Wilson problem}, with \emph{separation} $s$ and \emph{penalty constants} $(c_1,C_1,C_2)$, if the following conditions hold. First,
\begin{equation}\tag{B1}\label{eq:wilson-operators-B1}
W\ne\emptyset,\qquad W=\bigcup_{E\in\mathbb{E}}E,\qquad \sup_{E\in\mathbb{E}}\boldsymbol{\ell}_E<\infty,
\end{equation}
where $W$ and $\mathbb{E}$ are $R$-periodic and the elements of $\mathbb{E}$ are pairwise disjoint blocks of levels $\ge1$ whose sides $\boldsymbol{\ell}_E$ satisfy $\boldsymbol{\ell}_E\le R$. Second, the scale function $\boldsymbol{\ell}\colon\eta\mathbb{Z}^3\to(0,\infty)$ of Convention~\ref{it:wilson-operators-6}, for which $\boldsymbol{\ell}(x)=\boldsymbol{\ell}_E$ if $x\in E\in\mathbb{E}$ and $\boldsymbol{\ell}(x)=L\eta$ if $x\notin W$, satisfies
\begin{equation}\tag{B2}\label{eq:wilson-operators-B2}
\boldsymbol{\ell}(x)\le L\left(\boldsymbol{\ell}(z)+\frac{|x-z|_{\infty}}{s}\right)\qquad\text{for all }x,z\in\eta\mathbb{Z}^3;
\end{equation}
for $s=\infty$ this means $\boldsymbol{\ell}(x)\le L\cdot\boldsymbol{\ell}(z)$. Third, $\operatorname{Pen}=\sum_{E\in\mathbb{E}}\operatorname{Pen}_E$, where the operator $\operatorname{Pen}_E$ acts on functions $u\colon\eta\mathbb{Z}^3\to\mathbb{C}^4$ by
\begin{equation}\tag{B3}\label{eq:wilson-operators-B3}
(\operatorname{Pen}_Eu)(x)=\begin{cases}\eta^3\sum_{x'\in E}K_E(x,x')\,u(x'), & x\in E,\\ 0, & x\notin E,\end{cases}
\end{equation}
with a kernel $K_E\colon E\times E\to\mathbb{C}^{4\times4}$; the operator $\operatorname{Pen}$ commutes with the translations by $R\,e_{\mu}$, $\mu=1,2,3$; and the inequalities~\eqref{eq:wilson-operators-Pen1}--\eqref{eq:wilson-operators-Pen3} below hold for every $E\in\mathbb{E}$ and every $u\in L^2(E)$. Here $L^2(E)$ is the space of the functions $E\to\mathbb{C}^4$ of Section~\ref{sec:block-means-poincare}, with the scalar product $\langle u,v\rangle_E:=\eta^3\sum_{x\in E}\overline{u(x)}\cdot v(x)$ and the norm $\|\cdot\|_{2,E}$, and $(u,v)_E:=\eta^3\sum_{x\in E}u(x)\cdot v(x)$ denotes the bilinear pairing of $L^2(E)$ with itself. The values of $\operatorname{Pen}_Eu$ on $E$ involve only the values of $u$ on $E$, and $\operatorname{Pen}_Eu$ vanishes off $E$; hence $\operatorname{Pen}_E$ is also an operator on $L^2(E)$. Conversely, $L^2(E)$ being of finite dimension, every operator on $L^2(E)$ is of the form~\eqref{eq:wilson-operators-B3} with exactly one kernel, and it is identified with the operator~\eqref{eq:wilson-operators-B3} on the functions $\eta\mathbb{Z}^3\to\mathbb{C}^4$ that has this kernel. As an operator on $L^2(E)$, $\operatorname{Pen}_E$ has the adjoint $\operatorname{Pen}_E^{*}$ for $\langle\cdot,\cdot\rangle_E$ and the transpose $\operatorname{Pen}_E^T$ for $(\cdot,\cdot)_E$; these are the operators~\eqref{eq:wilson-operators-B3} with the kernels $(x,x')\mapsto K_E(x',x)^{*}$ and $(x,x')\mapsto K_E(x',x)^T$. With the mean $P_E$ over $E$ (Section~\ref{sec:block-means-poincare}), the inequalities read
\begin{equation}\tag{Pen1}\label{eq:wilson-operators-Pen1}
\operatorname{Re}\langle u,\operatorname{Pen}_Eu\rangle_E\ge\frac{c_1}{\boldsymbol{\ell}_E}\,\|P_Eu\|_{2,E}^2,
\end{equation}
\begin{multline}\tag{Pen2}\label{eq:wilson-operators-Pen2}
\|\operatorname{Pen}_Eu\|_{2,E}^2\le\frac{C_1}{\boldsymbol{\ell}_E}\operatorname{Re}\langle u,\operatorname{Pen}_Eu\rangle_E\quad\text{and}\\
\|\operatorname{Pen}_E^{*}u\|_{2,E}^2\le\frac{C_1}{\boldsymbol{\ell}_E}\operatorname{Re}\langle u,\operatorname{Pen}_Eu\rangle_E,
\end{multline}
\begin{equation}\tag{Pen3}\label{eq:wilson-operators-Pen3}
|K_E(x,x')|\le\frac{C_2}{\boldsymbol{\ell}_E^4}\qquad(x,x'\in E).
\end{equation}
Fourth,
\begin{equation}\tag{B4}\label{eq:wilson-operators-B4}
\mathfrak{m}\ge0\qquad\text{and}\qquad\mathfrak{m}\cdot\sup_W\boldsymbol{\ell}\le L.
\end{equation}
Let $\mathbb{H}_W:=\{u\in\mathbb{H}:\ u=0\text{ off }W\}$, let $[X]_W:=1_WX1_W$ for an operator $X$ on $\mathbb{H}$, and let
\begin{equation}\label{eq:wilson-operators-10}
T:=T_{\mathbb{P}}:=[D+\mathfrak{m}]_W+\operatorname{Pen}\qquad\text{on }\mathbb{H}_W.
\end{equation}
The blocks modulo $R$ are the elements of a fixed set of representatives of $\mathbb{E}$ modulo the translations by $R\mathbb{Z}^3$; for $R=\infty$ they are all the blocks of $\mathbb{E}$. Let $R<\infty$. Two distinct sites of one block are not congruent modulo $R\mathbb{Z}^3$, because the sides are at most $R$ by~\eqref{eq:wilson-operators-B1}; and two blocks containing congruent sites are congruent, because $\mathbb{E}$ is an $R$-periodic partition of the $R$-periodic set $W$. Hence the union of the blocks modulo $R$ and of a set of representatives of the sites outside $W$ modulo $R\mathbb{Z}^3$ is a period cell (Convention~\ref{it:wilson-operators-5}); for $R=\infty$ the same holds trivially. Consequently, if $w\in\mathbb{F}(\eta,R,\boldsymbol{\theta})$ vanishes off $W$, then $\|w\|^2=\sum_E\|w\|_{2,E}^2$ in $[0,\infty]$, and if moreover $w\in\mathbb{H}(\eta,R,\boldsymbol{\theta})$, then
\[
\langle u,w\rangle=\sum_E\langle u,w\rangle_E\quad\left(u\in\mathbb{H}(\eta,R,\boldsymbol{\theta})\right),\qquad
(u,w)=\sum_E(u,w)_E\quad\left(u\in\mathbb{H}(\eta,R,-\boldsymbol{\theta})\right);
\]
here $E$ ranges over the blocks modulo $R$, the functions are restricted to $E$, and for $R=\infty$ the series converge absolutely by the Cauchy--Schwarz inequality. The same holds with $-\boldsymbol{\theta}$ in place of $\boldsymbol{\theta}$, and $\|u\|^2\ge\sum_E\|u\|_{2,E}^2$ for every $u\in\mathbb{F}(\eta,R,\boldsymbol{\theta})$.

The operator $T$ is bounded. Indeed, $\|D\|\le12/\eta$, and~\eqref{eq:wilson-operators-Pen3} gives $\|\operatorname{Pen}_E\|\le C_2/\boldsymbol{\ell}_E\le C_2/(L\eta)$ on $L^2(E)$ for every $E\in\mathbb{E}$. As $\operatorname{Pen}$ commutes with the translations by $R\,e_{\mu}$, it maps $\mathbb{F}(\eta,R,\boldsymbol{\theta})$ into itself; moreover $\operatorname{Pen} v$ vanishes off $W$ and coincides with $\operatorname{Pen}_Ev$ on every $E\in\mathbb{E}$. Therefore
\[
\|\operatorname{Pen} v\|^2=\sum_E\|\operatorname{Pen}_Ev\|_{2,E}^2\le\frac{C_2^2}{(L\eta)^2}\sum_E\|v\|_{2,E}^2\le\frac{C_2^2}{(L\eta)^2}\,\|v\|^2\qquad(v\in\mathbb{H}).
\]
When $T$ is bijective, its inverse is denoted by $S:=S_{\mathbb{P}}:=T^{-1}$.

The \emph{adjoint problem} and the \emph{transposed problem} are
\[
\mathbb{P}^{*}:=(\eta,R,\boldsymbol{\theta};-\gamma;W,\mathbb{E},\operatorname{Pen}^{*},\mathfrak{m}),\qquad
\mathbb{P}^T:=(\eta,R,-\boldsymbol{\theta};-\gamma^T;W,\mathbb{E},\operatorname{Pen}^T,\mathfrak{m}),
\]
where $-\gamma^T:=(-\gamma_1^T,-\gamma_2^T,-\gamma_3^T)$, $\operatorname{Pen}^{*}:=\sum_{E\in\mathbb{E}}\operatorname{Pen}_E^{*}$ and $\operatorname{Pen}^T:=\sum_{E\in\mathbb{E}}\operatorname{Pen}_E^T$. Both are again blocked Wilson problems, with separation $s$ and penalty constants $(c_1,C_1,C_2)$; the operators $\operatorname{Pen}^{*}$ and $\operatorname{Pen}^T$ are, on $\mathbb{H}(\eta,R,\boldsymbol{\theta})$ and on $\mathbb{H}(\eta,R,-\boldsymbol{\theta})$ respectively, the adjoint and the transpose of $\operatorname{Pen}$ for the pairings of Convention~\ref{it:wilson-operators-5}; and
\[
T_{\mathbb{P}^{*}}=T^{*},\qquad T_{\mathbb{P}^T}=T^T.
\]
Indeed, $-\gamma$ and $-\gamma^T$ are again Clifford data (Convention~\ref{it:wilson-operators-3}), and the conditions~\eqref{eq:wilson-operators-B1}, \eqref{eq:wilson-operators-B2}, \eqref{eq:wilson-operators-B4} involve neither $\gamma$ nor $\boldsymbol{\theta}$ nor $\operatorname{Pen}$. For $E\in\mathbb{E}$ and $u\in L^2(E)$ one has
\[
\operatorname{Re}\langle u,\operatorname{Pen}_E^{*}u\rangle_E=\operatorname{Re}\langle u,\operatorname{Pen}_Eu\rangle_E,\qquad
\left(\operatorname{Pen}_E^{*}\right)^{*}=\operatorname{Pen}_E,\qquad
\left|K_E(x',x)^{*}\right|=\left|K_E(x',x)\right|;
\]
hence $\operatorname{Pen}_E^{*}$ satisfies~\eqref{eq:wilson-operators-Pen1} and~\eqref{eq:wilson-operators-Pen3}, and the two inequalities of~\eqref{eq:wilson-operators-Pen2} for $\operatorname{Pen}_E^{*}$ are those for $\operatorname{Pen}_E$, interchanged. The kernels show that $\operatorname{Pen}_E^Tu=\overline{\operatorname{Pen}_E^{*}\,\overline{u}}$ and $\left(\operatorname{Pen}_E^T\right)^{*}u=\overline{\operatorname{Pen}_E\,\overline{u}}$; as complex conjugation preserves $\|\cdot\|_{2,E}$ and commutes with $P_E$, it carries~\eqref{eq:wilson-operators-Pen1}--\eqref{eq:wilson-operators-Pen3} for $\operatorname{Pen}_E^{*}$ over to the same conditions for $\operatorname{Pen}_E^T$. The translations by $R\,e_{\mu}$ permute the blocks of $\mathbb{E}$, and the commutation of $\operatorname{Pen}$ with them amounts to $K_{E+R\,e_{\mu}}(x+R\,e_{\mu},x'+R\,e_{\mu})=K_E(x,x')$ for $E\in\mathbb{E}$ and $x,x'\in E$. The kernels $K_E(x',x)^{*}$ and $K_E(x',x)^T$ inherit this invariance; so $\operatorname{Pen}^{*}$ and $\operatorname{Pen}^T$ commute with the same translations and map $\mathbb{F}(\eta,R,\boldsymbol{\theta})$ and $\mathbb{F}(\eta,R,-\boldsymbol{\theta})$ into themselves, and, their blocks satisfying~\eqref{eq:wilson-operators-Pen3}, they are bounded on $\mathbb{H}(\eta,R,\boldsymbol{\theta})$ and on $\mathbb{H}(\eta,R,-\boldsymbol{\theta})$ by the argument given above for $\operatorname{Pen}$. The identities for the sums over the blocks modulo $R$, applied to $w=\operatorname{Pen} v$, $w=\operatorname{Pen}^{*}u$ and $w=\operatorname{Pen}^Tu'$, now give
\begin{align*}
\langle u,\operatorname{Pen} v\rangle&=\sum_E\langle u,\operatorname{Pen}_Ev\rangle_E=\sum_E\langle\operatorname{Pen}_E^{*}u,v\rangle_E=\langle\operatorname{Pen}^{*}u,v\rangle,\\
(u',\operatorname{Pen} v)&=\sum_E(u',\operatorname{Pen}_Ev)_E=\sum_E(\operatorname{Pen}_E^Tu',v)_E=(\operatorname{Pen}^Tu',v)
\end{align*}
for $u,v\in\mathbb{H}(\eta,R,\boldsymbol{\theta})$ and $u'\in\mathbb{H}(\eta,R,-\boldsymbol{\theta})$, the sums ranging over the blocks modulo $R$. Finally, by Lemma~\ref{lem:identities-free-wilson}(d), $D^{*}$ is the operator~\eqref{eq:wilson-operators-7} built with $-\gamma$, and $D^T$ is the operator~\eqref{eq:wilson-operators-7} built with $-\gamma^T$, acting on $\mathbb{F}(\eta,R,-\boldsymbol{\theta})$; as $\mathfrak{m}$ is real and the multiplication by $1_W$ is its own adjoint and its own transpose, the two identities for $T$ follow.
\end{definition}

\begin{definition}[Standard penalties]\label{def:standard-penalties}
Let $b>0$ and $c:=b/L$. The \emph{standard penalties} on a block $E$ are the operators $\operatorname{Pen}_E$ on $L^2(E)$ of the following two types:
\begin{equation}\tag{Std1}\label{eq:wilson-operators-Std1}
\operatorname{Pen}_E:=\frac{\boldsymbol{\beta}}{\boldsymbol{\ell}_E}\,P_E\qquad\text{with a number }\boldsymbol{\beta}\in\left[b\left(1-\frac{1}{L}\right),b\right];
\end{equation}
and, if $E$ has level $\ge2$,
\begin{multline}\tag{Std2}\label{eq:wilson-operators-Std2}
\operatorname{Pen}_E:=\frac{1}{\boldsymbol{\ell}'}\left[\boldsymbol{\alpha}_y\left(P'_E-P_E\right)+\left(\boldsymbol{\alpha}_y+\boldsymbol{\beta}_y\right)P_E\right]\\
=\frac{1}{\boldsymbol{\ell}'}\left[\boldsymbol{\alpha}_yP'_E+\boldsymbol{\beta}_yP_E\right],
\end{multline}
where $\boldsymbol{\ell}':=\boldsymbol{\ell}_E/L$, $P'_E$ is the orthogonal projection of $L^2(E)$ onto the subspace of functions constant on every one of the $L^3$ blocks of side $\boldsymbol{\ell}'$ into which $E$ decomposes, $b'\in[b(1-1/L),b]$, and $y\in\mathbb{R}\cup\{\infty\}$. Here $\boldsymbol{\alpha}_{\infty}:=b'$ and $\boldsymbol{\beta}_{\infty}:=0$, and for $y\in\mathbb{R}$ the $4\times4$ matrices $\boldsymbol{\alpha}_y$, $\boldsymbol{\beta}_y$ are defined by
\begin{gather*}
\boldsymbol{\alpha}_y:=b'\,i\hat{\Gamma} y\left(b'+i\hat{\Gamma} y\right)^{-1},\\
\boldsymbol{\alpha}_y+\boldsymbol{\beta}_y:=b'-b'^2\left(b'+c+i\hat{\Gamma} y\right)^{-1},
\end{gather*}
with the matrix $\hat{\Gamma}$ of Convention~\ref{it:wilson-operators-3}; the inverses exist because $\hat{\Gamma}$ is Hermitian and $b'>0$. In units of level $k$ (Convention~\ref{it:wilson-operators-1}), for a block $E$ of side $L$ one has $\boldsymbol{\ell}'=1$, and with $b'=b_k$ the operator~\eqref{eq:wilson-operators-Std2} is the restriction to $E$ of $\boldsymbol{\alpha}_yP_k+\boldsymbol{\beta}_yP_{k+1}$, where $P_i$ denotes the projection of Proposition~\ref{prop:wilson-operators-1}. The numbers
\[
b_j=b\left(1-L^{-1}\right)\left(1-L^{-j}\right)^{-1},\qquad j\ge1,
\]
of Section~\ref{sec:running-quantities-tori} lie in $[b(1-1/L),b]$, since $1-L^{-1}\le1-L^{-j}<1$.
\end{definition}

\begin{lemma}[Bounds for the standard penalties]\label{lem:bounds-standard-penalties}
Let $y\in\mathbb{R}\cup\{\infty\}$. Every standard penalty, as well as its adjoint and its transpose, satisfies~\eqref{eq:wilson-operators-Pen1}--\eqref{eq:wilson-operators-Pen3} with the constants
\[
c_1=b\left(1-\frac{1}{L}\right),\qquad C_1=L\cdot b,\qquad C_2=2L^4\cdot b.
\]
\end{lemma}

\begin{proof}
Let $u\in L^2(E)$, and write $\|\cdot\|:=\|\cdot\|_{2,E}$ and $\langle\cdot,\cdot\rangle:=\langle\cdot,\cdot\rangle_E$.

Let $\operatorname{Pen}_E$ be of type~\eqref{eq:wilson-operators-Std1}. Since $E$ has $(\boldsymbol{\ell}_E/\eta)^3$ sites, the kernel of $P_E$ is the constant $\boldsymbol{\ell}_E^{-3}$ times the unit matrix; thus $K_E$ is $\boldsymbol{\beta}\cdot\boldsymbol{\ell}_E^{-4}$ times the unit matrix. This kernel is real and symmetric, and $\operatorname{Pen}_E^{*}$, $\operatorname{Pen}_E^T$ have the kernels $K_E(x',x)^{*}$, $K_E(x',x)^T$ (Definition~\ref{def:blocked-wilson-problems}); hence $\operatorname{Pen}_E^{*}=\operatorname{Pen}_E^T=\operatorname{Pen}_E$. Moreover $|K_E|\le b\cdot\boldsymbol{\ell}_E^{-4}\le2L^4\cdot b\cdot\boldsymbol{\ell}_E^{-4}$, which is~\eqref{eq:wilson-operators-Pen3}. As $\operatorname{Pen}_E$ is $\boldsymbol{\beta}/\boldsymbol{\ell}_E$ times the orthogonal projection $P_E$, one has $\operatorname{Re}\langle u,\operatorname{Pen}_Eu\rangle=(\boldsymbol{\beta}/\boldsymbol{\ell}_E)\,\|P_Eu\|^2$, which gives~\eqref{eq:wilson-operators-Pen1} because $\boldsymbol{\beta}\ge b(1-1/L)=c_1$; and
\[
\|\operatorname{Pen}_Eu\|^2=\frac{\boldsymbol{\beta}}{\boldsymbol{\ell}_E}\operatorname{Re}\langle u,\operatorname{Pen}_Eu\rangle\le\frac{L\cdot b}{\boldsymbol{\ell}_E}\operatorname{Re}\langle u,\operatorname{Pen}_Eu\rangle,
\]
because $\boldsymbol{\beta}\le b$. This is~\eqref{eq:wilson-operators-Pen2}, whose two inequalities coincide because $\operatorname{Pen}_E^{*}=\operatorname{Pen}_E$.

Let now $\operatorname{Pen}_E$ be of type~\eqref{eq:wilson-operators-Std2} with $y\in\mathbb{R}$, and write $P:=P_E$, $P':=P'_E$. A function that is constant on $E$ is constant on each block of side $\boldsymbol{\ell}'$ in $E$, so $P'P=PP'=P$; hence $P'-P$ and $P$ are orthogonal projections with orthogonal ranges. Both act on the variable $x\in E$ only, and therefore commute with every matrix acting on $\mathbb{C}^4$. As $\hat{\Gamma}$ is Hermitian with $\hat{\Gamma}^2=1$, the space $\mathbb{C}^4$ is the orthogonal sum of the eigenspaces $\{\hat{\Gamma}=\boldsymbol{\sigma}\}$, $\boldsymbol{\sigma}\in\{1,-1\}$. The matrices $\boldsymbol{\alpha}_y$ and $\boldsymbol{\alpha}_y+\boldsymbol{\beta}_y$ are functions of $\hat{\Gamma}$, hence normal, and on $\{\hat{\Gamma}=\boldsymbol{\sigma}\}$ they act as multiplication by the numbers, again denoted $\boldsymbol{\alpha}$ and $\boldsymbol{\alpha}+\boldsymbol{\beta}$,
\begin{gather*}
\boldsymbol{\alpha}=\frac{i\,b'y'}{b'+iy'}=\frac{b'y'^2+i\,b'^2y'}{b'^2+y'^2},\\
\boldsymbol{\alpha}+\boldsymbol{\beta}=\frac{b'(c+iy')}{b'+c+iy'},\qquad\text{where }y':=\boldsymbol{\sigma}\,y.
\end{gather*}
Consequently
\[
|\boldsymbol{\alpha}|^2=\frac{b'^2y'^2}{b'^2+y'^2}=b'\cdot\operatorname{Re}\boldsymbol{\alpha},
\]
further
\[
\operatorname{Re}(\boldsymbol{\alpha}+\boldsymbol{\beta})=b'\,\frac{c(b'+c)+y'^2}{(b'+c)^2+y'^2}\ge\frac{b'c}{b'+c},
\]
since the quotient is a non-decreasing function of $y'^2$, as $c(b'+c)\le(b'+c)^2$, and
\[
|\boldsymbol{\alpha}+\boldsymbol{\beta}|^2=\frac{b'^2(c^2+y'^2)}{(b'+c)^2+y'^2}\le b'\cdot\operatorname{Re}(\boldsymbol{\alpha}+\boldsymbol{\beta}),
\]
since $c^2\le c(b'+c)$. For a matrix $A$ let $\operatorname{Re}A:=(A+A^{*})/2$. Because $P'-P$ and $P$ are orthogonal projections commuting with matrices,
\begin{multline*}
\operatorname{Re}\langle u,\operatorname{Pen}_Eu\rangle\\
=\frac{1}{\boldsymbol{\ell}'}\left[\left\langle(P'-P)u,\left(\operatorname{Re}\boldsymbol{\alpha}_y\right)(P'-P)u\right\rangle+\left\langle Pu,\operatorname{Re}\left(\boldsymbol{\alpha}_y+\boldsymbol{\beta}_y\right)Pu\right\rangle\right].
\end{multline*}
The ranges of $P'-P$ and $P$ are orthogonal, and the normal matrices $\boldsymbol{\alpha}_y$, $\boldsymbol{\alpha}_y+\boldsymbol{\beta}_y$ satisfy the bounds above on each eigenspace of $\hat{\Gamma}$; hence
\begin{align*}
\|\operatorname{Pen}_Eu\|^2=\|\operatorname{Pen}_E^{*}u\|^2&=\frac{1}{\boldsymbol{\ell}'^2}\left(\left\|\boldsymbol{\alpha}_y(P'-P)u\right\|^2+\left\|\left(\boldsymbol{\alpha}_y+\boldsymbol{\beta}_y\right)Pu\right\|^2\right)\\
&\le\frac{b'}{\boldsymbol{\ell}'}\operatorname{Re}\langle u,\operatorname{Pen}_Eu\rangle\le\frac{L\cdot b}{\boldsymbol{\ell}_E}\operatorname{Re}\langle u,\operatorname{Pen}_Eu\rangle,
\end{align*}
where $\boldsymbol{\ell}'=\boldsymbol{\ell}_E/L$ and $b'\le b$ were used in the last step. This is~\eqref{eq:wilson-operators-Pen2}. The identity $|\boldsymbol{\alpha}|^2=b'\cdot\operatorname{Re}\boldsymbol{\alpha}$ shows that $\operatorname{Re}\boldsymbol{\alpha}_y$ is positive semidefinite, so the expression for $\operatorname{Re}\langle u,\operatorname{Pen}_Eu\rangle$ gives
\[
\operatorname{Re}\langle u,\operatorname{Pen}_Eu\rangle\ge\frac{L}{\boldsymbol{\ell}_E}\cdot\frac{b'c}{b'+c}\,\|P_Eu\|^2,
\]
and, as $c=b/L$ and $b'\ge b(1-1/L)$,
\[
\frac{L\,b'c}{b'+c}=\frac{b}{1+b\,(Lb')^{-1}}\ge\frac{b}{1+(L-1)^{-1}}=b\left(1-\frac{1}{L}\right).
\]
This is~\eqref{eq:wilson-operators-Pen1}. The kernel of $P'$ equals $\boldsymbol{\ell}'^{-3}$ at the pairs $(x,x')$ that lie in one block of side $\boldsymbol{\ell}'$ and vanishes at the other pairs, and the kernel of $P$ is $\boldsymbol{\ell}_E^{-3}$. Hence $K_E(x,x')$ equals $\boldsymbol{\ell}'^{-1}\left(\boldsymbol{\alpha}_y\cdot\boldsymbol{\ell}'^{-3}+\boldsymbol{\beta}_y\cdot\boldsymbol{\ell}_E^{-3}\right)$ if $x$ and $x'$ lie in one block of side $\boldsymbol{\ell}'$, and $\boldsymbol{\ell}'^{-1}\cdot\boldsymbol{\beta}_y\cdot\boldsymbol{\ell}_E^{-3}$ otherwise. The formula for $|\boldsymbol{\alpha}|^2$ gives $|\boldsymbol{\alpha}_y|\le b'$, and on $\{\hat{\Gamma}=\boldsymbol{\sigma}\}$ the matrix $\boldsymbol{\beta}_y$ acts as multiplication by
\[
\boldsymbol{\beta}=(\boldsymbol{\alpha}+\boldsymbol{\beta})-\boldsymbol{\alpha}=\frac{b'^2c}{(b'+c+iy')(b'+iy')},
\]
so that $|\boldsymbol{\beta}_y|\le b'^2c/((b'+c)\,b')\le c$. With $\boldsymbol{\ell}'=\boldsymbol{\ell}_E/L$, $b'\le b$ and $L\cdot c=b$ it follows that
\[
|K_E(x,x')|\le\boldsymbol{\ell}_E^{-4}\left(L^4\cdot b+L\cdot c\right)\le2L^4\cdot b\cdot\boldsymbol{\ell}_E^{-4},
\]
which is~\eqref{eq:wilson-operators-Pen3}.

For $y=\infty$ the operator~\eqref{eq:wilson-operators-Std2} is the limit of the operators with $y\in\mathbb{R}$ as $|y|\to\infty$, because $\boldsymbol{\alpha}_y\to b'=\boldsymbol{\alpha}_{\infty}$ and $\boldsymbol{\beta}_y\to0=\boldsymbol{\beta}_{\infty}$; the non-strict inequalities~\eqref{eq:wilson-operators-Pen1}--\eqref{eq:wilson-operators-Pen3} are preserved in this limit. Since $\hat{\Gamma}$ is Hermitian, $\boldsymbol{\alpha}_y^{*}=\boldsymbol{\alpha}_{-y}$ and $(\boldsymbol{\alpha}_y+\boldsymbol{\beta}_y)^{*}=\boldsymbol{\alpha}_{-y}+\boldsymbol{\beta}_{-y}$ for $y\in\mathbb{R}$; so the adjoint of the penalty~\eqref{eq:wilson-operators-Std2} with parameter $y$ is the penalty with parameter $-y$, and the penalty with $y=\infty$ is self-adjoint. The kernels of $P$ and $P'$ are real and symmetric; so the transpose of the penalty~\eqref{eq:wilson-operators-Std2} is the penalty~\eqref{eq:wilson-operators-Std2} built with $\hat{\Gamma}^T$ in place of $\hat{\Gamma}$, and $\hat{\Gamma}^T$ is again Hermitian with square $1$ (Convention~\ref{it:wilson-operators-3}). The bounds proved above hold for every $y$ and for every Hermitian matrix $\hat{\Gamma}$ with $\hat{\Gamma}^2=1$; therefore they hold for the adjoints and for the transposes.
\end{proof}

\subsubsection{Regularisation of the scale function}\label{sec:regularisation-scale-function}

\begin{lemma}[A Lipschitz minorant of the scale function]\label{lem:lipschitz-minorant-scale}
Let $s\in[6L,\infty]$, let $\eta$, $R$, $W$, $\mathbb{E}$ satisfy~\eqref{eq:wilson-operators-B1} and~\eqref{eq:wilson-operators-B2}, and define $\boldsymbol{\lambda}\colon\eta\mathbb{Z}^3\to\mathbb{R}$ by
\[
\boldsymbol{\lambda}(x):=\inf_{z\in\eta\mathbb{Z}^3}\left(\boldsymbol{\ell}(z)+\frac{|x-z|_{\infty}}{s}\right)\quad\text{if }s<\infty,\qquad
\boldsymbol{\lambda}:=\inf_{\eta\mathbb{Z}^3}\boldsymbol{\ell}\quad\text{if }s=\infty.
\]
Here and below $t/s:=0$ for $t\in\mathbb{R}$ if $s=\infty$. Then the following holds for all $x,x',z\in\eta\mathbb{Z}^3$.
\begin{enumerate}
\item[(a)] $\boldsymbol{\lambda}$ is $R$-periodic, $\boldsymbol{\ell}/L\le\boldsymbol{\lambda}\le\boldsymbol{\ell}$, and $\boldsymbol{\lambda}\ge\eta$.
\item[(b)] $|\boldsymbol{\lambda}(x)-\boldsymbol{\lambda}(x')|\le|x-x'|_{\infty}/s$.
\item[(c)] If $x$ has a nearest neighbour outside $W$, then $\boldsymbol{\lambda}(x)\le(L+1)\eta$.
\item[(d)] If $|x-z|_{\infty}\le\boldsymbol{\ell}(x)$, then $\frac{5}{6}\,\boldsymbol{\lambda}(x)\le\boldsymbol{\lambda}(z)\le\frac{7}{6}\,\boldsymbol{\lambda}(x)$ and $\boldsymbol{\ell}(x)/L\le\boldsymbol{\ell}(z)\le L\cdot\boldsymbol{\ell}(x)$.
\item[(e)] If $x\in W$ and $\boldsymbol{\ell}(x)\ge L^3\eta$, then $|x-z|_{\infty}\ge4\,\boldsymbol{\ell}(x)$ for every $z\notin W$.
\end{enumerate}
\end{lemma}

\begin{proof}
(a) By~\eqref{eq:wilson-operators-B1} the sets $W$ and $\mathbb{E}$ are $R$-periodic, hence so are $\boldsymbol{\ell}$ and $\boldsymbol{\lambda}$. Let $s<\infty$. The choice $z=x$ in the infimum gives $\boldsymbol{\lambda}\le\boldsymbol{\ell}$, and~\eqref{eq:wilson-operators-B2} states that $\boldsymbol{\ell}(z)+|x-z|_{\infty}/s\ge\boldsymbol{\ell}(x)/L$ for every $z$, so that $\boldsymbol{\lambda}\ge\boldsymbol{\ell}/L$. For $s=\infty$,~\eqref{eq:wilson-operators-B2} gives $\boldsymbol{\ell}(x)\le L\cdot\boldsymbol{\ell}(z)$ for all $x,z$, and again $\boldsymbol{\ell}/L\le\boldsymbol{\lambda}\le\boldsymbol{\ell}$. All blocks of $\mathbb{E}$ have level $\ge1$, and $\boldsymbol{\ell}=L\eta$ off $W$; thus $\boldsymbol{\ell}\ge L\eta$ everywhere, and $\boldsymbol{\lambda}\ge\boldsymbol{\ell}/L\ge\eta$.

(b) For $s=\infty$ the function $\boldsymbol{\lambda}$ is constant. For $s<\infty$ and fixed $z$, the function $x\mapsto\boldsymbol{\ell}(z)+|x-z|_{\infty}/s$ is Lipschitz with constant $1/s$; so is the infimum $\boldsymbol{\lambda}$ of these functions, which is finite by (a).

(c) Let $z\notin W$ be a nearest neighbour of $x$. Then $\boldsymbol{\ell}(z)=L\eta$ and $|x-z|_{\infty}=\eta$, whence $\boldsymbol{\lambda}(x)\le L\eta+\eta/s\le(L+1)\eta$.

(d) By (b), (a) and $s\ge6L$,
\[
\left|\boldsymbol{\lambda}(z)-\boldsymbol{\lambda}(x)\right|\le\frac{\boldsymbol{\ell}(x)}{s}\le\frac{L\cdot\boldsymbol{\lambda}(x)}{s}\le\frac{\boldsymbol{\lambda}(x)}{6},
\]
which is the first assertion. Using (a) again,
\begin{gather*}
\boldsymbol{\ell}(z)\le L\cdot\boldsymbol{\lambda}(z)\le\frac{7}{6}\,L\cdot\boldsymbol{\lambda}(x)\le\frac{7}{6}\,L\cdot\boldsymbol{\ell}(x)<L^2\cdot\boldsymbol{\ell}(x),\\
\boldsymbol{\ell}(z)\ge\boldsymbol{\lambda}(z)\ge\frac{5}{6}\,\boldsymbol{\lambda}(x)\ge\frac{5}{6}\cdot\frac{\boldsymbol{\ell}(x)}{L}>\frac{\boldsymbol{\ell}(x)}{L^2}.
\end{gather*}
Every value of $\boldsymbol{\ell}$ is of the form $L^i\eta$ with an integer $i\ge1$ (Convention~\ref{it:wilson-operators-4}); hence $\boldsymbol{\ell}(x)/L\le\boldsymbol{\ell}(z)\le L\cdot\boldsymbol{\ell}(x)$.

(e) Let $z\notin W$, so that $\boldsymbol{\ell}(z)=L\eta$. For $s=\infty$,~\eqref{eq:wilson-operators-B2} would give $\boldsymbol{\ell}(x)\le L^2\eta$, contrary to $\boldsymbol{\ell}(x)\ge L^3\eta$; so for $s=\infty$ no such $z$ exists. Let $s<\infty$. By~\eqref{eq:wilson-operators-B2}, and since $L\eta\le\boldsymbol{\ell}(x)/L^2$, $s\ge6L$ and $L\ge5$ (Convention~\ref{it:wilson-operators-4}),
\[
|x-z|_{\infty}\ge s\left(\frac{\boldsymbol{\ell}(x)}{L}-L\eta\right)\ge s\cdot\frac{\boldsymbol{\ell}(x)}{L}\left(1-\frac{1}{L}\right)\ge6\cdot\frac{4}{5}\,\boldsymbol{\ell}(x)\ge4\,\boldsymbol{\ell}(x).
\]
\end{proof}

\subsubsection{A distance adapted to the block sizes}\label{sec:distance-adapted-block}

\begin{definition}[Path distance in units of the local scale]\label{def:path-distance-units}
Let $s$, $\eta$, $R$, $W$, $\mathbb{E}$ be as in Lemma~\ref{lem:lipschitz-minorant-scale}, and let $\boldsymbol{\lambda}$ be the function defined there. A \emph{bond} is a pair $\{x,x'\}$ of nearest neighbours in $\eta\mathbb{Z}^3$, and its \emph{cost} is $\eta/\min\{\boldsymbol{\lambda}(x),\boldsymbol{\lambda}(x')\}$; the cost of a nearest neighbour path in $\eta\mathbb{Z}^3$ is the sum of the costs of its bonds. For $x,x'\in\eta\mathbb{Z}^3$ let $d_{\infty}(x,x')$ be the least cost of a nearest neighbour path in $\eta\mathbb{Z}^3$ from $x$ to $x'$, the paths not being required to stay in $W$, and let
\[
d(x,x'):=\min_{n\in\mathbb{Z}^3}d_{\infty}(x,x'+Rn)\quad\text{if }R<\infty,\qquad d:=d_{\infty}\quad\text{if }R=\infty.
\]
The function $d$ is a metric on the torus $\eta\mathbb{Z}^3/R\mathbb{Z}^3$, and on $\eta\mathbb{Z}^3$ if $R=\infty$: by Lemma~\ref{lem:lipschitz-minorant-scale}(a) and~\eqref{eq:wilson-operators-B1}, $\boldsymbol{\lambda}$ is $R$-periodic and bounded, so that $d_{\infty}$ is invariant under the translations by $R\mathbb{Z}^3$ and every bond costs at least $\eta/\sup\boldsymbol{\lambda}>0$. For sets $A,A'\subset\eta\mathbb{Z}^3$ let $d(A,A'):=\inf\{d(x,x'):x\in A,\ x'\in A'\}$, define $d_{\infty}(A,A')$ in the same way, and write $d(x,A'):=d(\{x\},A')$ and $d_{\infty}(x,A'):=d_{\infty}(\{x\},A')$. If $R<\infty$, then $|x-x'|_{\infty}$ denotes the periodic sup distance $\min_{n\in\mathbb{Z}^3}|x-x'-Rn|_{\infty}$ of sites $x,x'$, the norm on the right being that of Convention~\ref{it:wilson-operators-4}; the exceptions are the H\"older quotients of Convention~\ref{it:wilson-operators-7} and the arguments carried out on the covering lattice $\eta\mathbb{Z}^3$, in which $|x-x'|_{\infty}$ is the sup distance of $\eta\mathbb{Z}^3$.
\end{definition}

\begin{lemma}[Properties of the path distance]\label{lem:properties-path-distance}
In the setting of Definition~\ref{def:path-distance-units} the following holds for all $x,x',z\in\eta\mathbb{Z}^3$.
\begin{enumerate}
\item[(a)] $|\log(\boldsymbol{\lambda}(x)/\boldsymbol{\lambda}(x'))|\le d(x,x')/s$; consequently $\boldsymbol{\ell}(x)/\boldsymbol{\ell}(x')\le L\cdot\exp\left(\frac{d(x,x')}{s}\right)$.
\item[(b)] One has
\[
|x-x'|_{\infty}\le d(x,x')\cdot\boldsymbol{\lambda}(x)\cdot\exp\left(\frac{d(x,x')}{s}\right),
\]
and $d(x,x')\ge|x-x'|_{\infty}/\sup\boldsymbol{\lambda}$.
\item[(c)] If $|x-z|_{\infty}\le\boldsymbol{\ell}(x)$, then $d(x,z)\le4L$. If $x$ and $z$ lie in one block of $\mathbb{E}$, then $d(x,z)\le3L$.
\item[(d)] (Shell sums.) Let $\sigma>0$ and $s\ge12/\sigma$. Then for every site $x$, the sum on the left being taken over the blocks $E'$ of $\mathbb{E}$ modulo $R$,
\begin{multline}\label{eq:wilson-operators-11}
\sum_{E'}\exp\left(-\sigma\,d(x,E')\right)\\
\le C_{\mathrm{sh}}(\sigma):=8\exp\left(\frac{\sigma}{2}\right)\sum_{n\ge0}(n+2)^3\exp\left(-\frac{\sigma n}{2}\right).
\end{multline}
\end{enumerate}
\end{lemma}

\begin{proof}
(a) Let $\{w,w'\}$ be a bond, and let $m:=\min\{\boldsymbol{\lambda}(w),\boldsymbol{\lambda}(w')\}$. By Lemma~\ref{lem:lipschitz-minorant-scale}(b), $|\boldsymbol{\lambda}(w)-\boldsymbol{\lambda}(w')|\le\eta/s$, hence
\[
\frac{\max\{\boldsymbol{\lambda}(w),\boldsymbol{\lambda}(w')\}}{m}\le1+\frac{\eta}{s\,m}\le\exp\left(\frac{\eta}{s\,m}\right);
\]
that is, $|\log(\boldsymbol{\lambda}(w)/\boldsymbol{\lambda}(w'))|$ is at most the cost of the bond divided by $s$. Adding these inequalities along a nearest neighbour path from $x$ to a point $x''$, one finds that $|\log(\boldsymbol{\lambda}(x)/\boldsymbol{\lambda}(x''))|$ is at most the cost of the path divided by $s$. For $x''=x'+Rn$ one has $\boldsymbol{\lambda}(x'')=\boldsymbol{\lambda}(x')$ by Lemma~\ref{lem:lipschitz-minorant-scale}(a); the minimum over the paths and over $n\in\mathbb{Z}^3$, only $n=0$ being admitted if $R=\infty$, gives the first inequality. The second follows from the first, since $\boldsymbol{\ell}(x)\le L\cdot\boldsymbol{\lambda}(x)$ and $\boldsymbol{\lambda}(x')\le\boldsymbol{\ell}(x')$ by Lemma~\ref{lem:lipschitz-minorant-scale}(a).

(b) Choose $x''=x'+Rn$, with $x''=x'$ if $R=\infty$, and a nearest neighbour path from $x$ to $x''$ of cost $d(x,x')$. A point $w$ of the path is joined to $x$ by a part of the path, whose cost is at most $d(x,x')$; hence $d(x,w)\le d(x,x')$, and (a) gives
\[
\boldsymbol{\lambda}(w)\le\boldsymbol{\lambda}(x)\cdot\exp\left(\frac{d(x,x')}{s}\right).
\]
A bond $\{w,w'\}$ of the path has length $\eta$, which is its cost times $\min\{\boldsymbol{\lambda}(w),\boldsymbol{\lambda}(w')\}$ and therefore at most its cost times $\boldsymbol{\lambda}(x)\cdot\exp\left(\frac{d(x,x')}{s}\right)$. Adding over the bonds, one obtains that the sup distance in $\eta\mathbb{Z}^3$ of $x$ and $x''$ is at most $d(x,x')\cdot\boldsymbol{\lambda}(x)\cdot\exp\left(\frac{d(x,x')}{s}\right)$; and $|x-x'|_{\infty}$ does not exceed this sup distance. For the second inequality, observe that every bond costs at least $\eta/\sup\boldsymbol{\lambda}$, and that a nearest neighbour path from $x$ to a point $x'+Rn$ has at least $|x-x'|_{\infty}/\eta$ bonds.

(c) Let $|x-z|_{\infty}\le\boldsymbol{\ell}(x)$. Replacing $z$ by a suitable point $z+Rn$, which does not change $d(x,z)$, one may assume that the sup distance in $\eta\mathbb{Z}^3$ of $x$ and $z$ is at most $\boldsymbol{\ell}(x)$. A monotone coordinate path from $x$ to $z$, that is, a nearest neighbour path along which each coordinate is monotone, has at most $3\,\boldsymbol{\ell}(x)/\eta$ bonds, and all its points $w$ satisfy $|w-x|_{\infty}\le\boldsymbol{\ell}(x)$. At these points $\boldsymbol{\lambda}(w)\ge\frac{5}{6}\,\boldsymbol{\lambda}(x)\ge\frac{5}{6}\,\boldsymbol{\ell}(x)/L$, by Lemma~\ref{lem:lipschitz-minorant-scale}(d) and~(a). Each bond of the path therefore costs at most $\frac{6}{5}\,L\eta/\boldsymbol{\ell}(x)$, and
\[
d(x,z)\le3\cdot\frac{6}{5}\,L=\frac{18}{5}\,L\le4L.
\]
If $x$ and $z$ lie in one block $E$ of $\mathbb{E}$, a monotone coordinate path from $x$ to $z$ stays in $E$ and has at most $3\,\boldsymbol{\ell}_E/\eta$ bonds, and $\boldsymbol{\lambda}\ge\boldsymbol{\ell}/L=\boldsymbol{\ell}_E/L$ on $E$ by Lemma~\ref{lem:lipschitz-minorant-scale}(a); so each bond costs at most $L\eta/\boldsymbol{\ell}_E$, and $d(x,z)\le3L$.

(d) Let $R<\infty$, and let $E'$ be a block of $\mathbb{E}$. By Definition~\ref{def:path-distance-units}, $d(x,E')=\min_{n\in\mathbb{Z}^3}d_{\infty}(x,E'+Rn)$, so that
\[
\exp\left(-\sigma\,d(x,E')\right)\le\sum_{n\in\mathbb{Z}^3}\exp\left(-\sigma\,d_{\infty}(x,E'+Rn)\right);
\]
and when $E'$ ranges over a set of representatives of the blocks modulo $R$ and $n$ over $\mathbb{Z}^3$, the blocks $E'+Rn$ are pairwise distinct elements of $\mathbb{E}$. For $R=\infty$ one has $d=d_{\infty}$. In either case it suffices to prove
\[
\sum_{E''\in\mathbb{E}}\exp\left(-\sigma\,d_{\infty}(x,E'')\right)\le C_{\mathrm{sh}}(\sigma),
\]
the sum being taken over all blocks of $\mathbb{E}$, as subsets of the covering lattice $\eta\mathbb{Z}^3$; in this argument $|\cdot|_{\infty}$ is the sup distance of $\eta\mathbb{Z}^3$. Let $n\ge0$ be an integer, let $E''\in\mathbb{E}$ satisfy $n\le d_{\infty}(x,E'')<n+1$, and choose $x'\in E''$ with $d_{\infty}(x,x')<n+1$. The argument given for (b), applied to a nearest neighbour path from $x$ to $x'$ of cost less than $n+1$, yields
\[
|x-x'|_{\infty}\le(n+1)\,\boldsymbol{\lambda}(x)\cdot\exp\left(\frac{n+1}{s}\right)=:\mathsf{r}_n.
\]
As $d(x,x')\le d_{\infty}(x,x')<n+1$, part (a) and Lemma~\ref{lem:lipschitz-minorant-scale}(a) yield
\[
\boldsymbol{\ell}_{E''}=\boldsymbol{\ell}(x')\ge\boldsymbol{\lambda}(x')\ge\boldsymbol{\lambda}(x)\cdot\exp\left(-\frac{n+1}{s}\right)=:a_n.
\]
The closed cube of $E''$ (Convention~\ref{it:wilson-operators-4}) contains $x'$ and has side $\boldsymbol{\ell}_{E''}\ge a_n$; hence it contains a closed cube of side $a_n$ that contains $x'$, and the latter lies in the closed cube of half side $\mathsf{r}_n+a_n$ centred at $x$. Distinct blocks of $\mathbb{E}$ are disjoint, so their closed cubes have disjoint interiors (Convention~\ref{it:wilson-operators-4}). Comparing volumes, one finds that the number of the blocks $E''$ with $n\le d_{\infty}(x,E'')<n+1$ is at most
\begin{align*}
\left(\frac{2(\mathsf{r}_n+a_n)}{a_n}\right)^3&=8\left((n+1)\exp\left(\frac{2(n+1)}{s}\right)+1\right)^3\\
&\le8\,(n+2)^3\exp\left(\frac{6(n+1)}{s}\right)\le8\,(n+2)^3\exp\left(\frac{\sigma(n+1)}{2}\right),
\end{align*}
where $6/s\le\sigma/2$ was used in the last step. Each of these blocks contributes at most $\exp(-\sigma n)$ to the sum; summing over $n$,
\[
\sum_{E''\in\mathbb{E}}\exp\left(-\sigma\,d_{\infty}(x,E'')\right)\le\sum_{n\ge0}8\,(n+2)^3\exp\left(\frac{\sigma(n+1)}{2}-\sigma n\right)=C_{\mathrm{sh}}(\sigma).
\]
\end{proof}

\subsection{Invertibility and a weighted $L^2$ bound}\label{sec:invertibility-weighted-bound}

\begin{theorem}[Weighted $L^2$ bound]\label{thm:weighted-bound}
Let $\mathbb{P}$ be a blocked Wilson problem with separation $s \ge 6L$ and penalty constants $(c_1, C_1, C_2)$, and define
\[
K_1 := L^{\frac{1}{2}} + C_1^{\frac{1}{2}} + (6(L+1))^{\frac{1}{2}} , \qquad C_I := \frac{2L}{c_1} + \frac{3}{4} L^2 K_1^2 + L ,
\]
the expressions of~\eqref{eq:wilson-operators-3} with $c_1$ and $C_1$ in the roles of $b \left( 1 - \frac{1}{L} \right)$ and $L \cdot b$. Then for every $u \in \mathbb{H}_W$ the function $f := T u$ satisfies
\begin{equation}\label{eq:wilson-operators-12}
\left\| \boldsymbol{\lambda}^{-\frac{1}{2}} u \right\| \le C_I \left\| \boldsymbol{\lambda}^{\frac{1}{2}} f \right\| .
\end{equation}
The operator $T$ is a bijection of $\mathbb{H}_W$ onto itself, and its inverse $S = T^{-1}$ satisfies
\[
\left\| \boldsymbol{\lambda}^{-\frac{1}{2}} S \boldsymbol{\lambda}^{-\frac{1}{2}} \right\| \le C_I ,
\]
where $\|A\| := \sup \left\{ \|A w\| : w \in \mathbb{H}_W,\ \|w\| \le 1 \right\}$ denotes the operator norm of a linear operator $A$ on $\mathbb{H}_W$; in particular $\|S\| \le C_I \sup_W \boldsymbol{\lambda} \le C_I \sup_W \boldsymbol{\ell}$. The same holds with $T$ replaced by $T^{*}$ or by $T^T$.
\end{theorem}

\begin{proof}
Let $u \in \mathbb{H}_W$ and $f := T u$, and put
\[
N := \left\| \boldsymbol{\lambda}^{-\frac{1}{2}} u \right\| , \qquad F := \left\| \boldsymbol{\lambda}^{\frac{1}{2}} f \right\| , \qquad X := \operatorname{Re} \langle u, f \rangle , \qquad \varepsilon_0 := \frac{1}{s} \le \frac{1}{6L} .
\]
By Lemma~\ref{lem:lipschitz-minorant-scale}(a) and condition~\eqref{eq:wilson-operators-B1} of Definition~\ref{def:blocked-wilson-problems} one has $\eta \le \boldsymbol{\lambda} \le \sup \boldsymbol{\ell} < \infty$. Hence the multiplications by $\boldsymbol{\lambda}^{\frac{1}{2}}$ and by $\boldsymbol{\lambda}^{-\frac{1}{2}}$ are bounded operators on $\mathbb{H}$; as $D$ and $\operatorname{Pen}$ are bounded as well (Definition~\ref{def:blocked-wilson-problems}), all norms and sums below are finite. Sums over sites extend over a period cell (over $\eta\mathbb{Z}^3$ if $R = \infty$), and sums over $E$ over the blocks of $\mathbb{E}$ modulo $R$ (over all blocks if $R = \infty$). By condition~\eqref{eq:wilson-operators-B1} every block has side at most $R$, so that its sites are pairwise distinct modulo $R$; the blocks modulo $R$ therefore partition a period cell of $W$, and for every $w \in \mathbb{H}_W$
\[
\|w\|^2 = \sum_E \|w\|_{2,E}^2 , \qquad \langle w, \operatorname{Pen} w \rangle = \sum_E \langle w, \operatorname{Pen}_E w \rangle_E .
\]

\emph{Step 1: the real part of the form.} Since $u = 1_W u$, the definition~\eqref{eq:wilson-operators-10} of $T$ gives $\langle u, f \rangle = \langle u, (D + \mathfrak{m}) u \rangle + \langle u, \operatorname{Pen} u \rangle$. The real part of $\langle u, D u \rangle$ is given by Lemma~\ref{lem:identities-free-wilson}(a), and the Cauchy--Schwarz inequality applies to $X = \operatorname{Re} \left\langle \boldsymbol{\lambda}^{-\frac{1}{2}} u, \boldsymbol{\lambda}^{\frac{1}{2}} f \right\rangle$. Thus
\begin{equation}\label{eq:wilson-operators-13}
X = \frac{\eta}{2} \left\| \nabla^{+} u \right\|^2 + \mathfrak{m} \|u\|^2 + \operatorname{Re} \langle u, \operatorname{Pen} u \rangle , \qquad X \le \left| \left\langle \boldsymbol{\lambda}^{-\frac{1}{2}} u, \boldsymbol{\lambda}^{\frac{1}{2}} f \right\rangle \right| \le N F .
\end{equation}
The three terms of the identity are non-negative: the first evidently, the second since $\mathfrak{m} \ge 0$ by condition~\eqref{eq:wilson-operators-B4}, the third since $\operatorname{Re} \langle u, \operatorname{Pen} u \rangle = \sum_E \operatorname{Re} \langle u, \operatorname{Pen}_E u \rangle_E$ with summands that are non-negative by condition~\eqref{eq:wilson-operators-Pen1}. Each of the three terms is therefore at most $X$, and four bounds follow.
\begin{enumerate}
\item[(i)] By condition~\eqref{eq:wilson-operators-Pen1}, $c_1 \sum_E \boldsymbol{\ell}_E^{-1} \|P_E u\|_{2,E}^2 \le \sum_E \operatorname{Re} \langle u, \operatorname{Pen}_E u \rangle_E \le X$, that is,
\[
\sum_E \boldsymbol{\ell}_E^{-1} \|P_E u\|_{2,E}^2 \le \frac{X}{c_1} .
\]
\item[(ii)] On a block $E$ one has $\boldsymbol{\lambda} \le \boldsymbol{\ell} = \boldsymbol{\ell}_E$ (Lemma~\ref{lem:lipschitz-minorant-scale}(a)) and $\operatorname{Pen} u = \operatorname{Pen}_E u$. With condition~\eqref{eq:wilson-operators-Pen2},
\[
\left\| \boldsymbol{\lambda}^{\frac{1}{2}} \cdot \operatorname{Pen} u \right\|^2 \le \sum_E \boldsymbol{\ell}_E \|\operatorname{Pen}_E u\|_{2,E}^2 \le C_1 \sum_E \operatorname{Re} \langle u, \operatorname{Pen}_E u \rangle_E \le C_1 X .
\]
\item[(iii)] As $u$ vanishes off $W$, $\boldsymbol{\lambda} \le \boldsymbol{\ell}$, and $\mathfrak{m} \cdot \sup_W \boldsymbol{\ell} \le L$ by condition~\eqref{eq:wilson-operators-B4},
\[
\left\| \mathfrak{m} \cdot \boldsymbol{\lambda}^{\frac{1}{2}} u \right\|^2 \le \mathfrak{m} \cdot \left( \sup_W \boldsymbol{\ell} \right) \cdot \mathfrak{m} \|u\|^2 \le L X .
\]
\item[(iv)] Let $\mathrm{bd}$ consist of the pairs $(x, x')$ whose first entry is a site $x \notin W$ and whose second entry is a nearest neighbour $x' \in W$ of $x$; if $R < \infty$, the pairs $(x, x')$ and $(x + R n, x' + R n)$ with $n \in \mathbb{Z}^3$ are identified. One has $\left\| \nabla^{+} u \right\|^2 = \eta \sum_{\mu = 1}^3 \sum_y |u(y + \eta e_{\mu}) - u(y)|^2$, so that each bond contributes $\eta$ times the squared modulus of the difference of the values of $u$ at its two ends. For $(x, x') \in \mathrm{bd}$ this difference has modulus $|u(x')|$, because $u(x) = 0$; and distinct elements of $\mathrm{bd}$ give distinct bonds, because exactly one end of the bond lies in $W$. Consequently
\[
\frac{\eta^2}{2} \sum_{(x, x') \in \mathrm{bd}} |u(x')|^2 \le \frac{\eta}{2} \left\| \nabla^{+} u \right\|^2 \le X .
\]
\end{enumerate}

\emph{Step 2: the free Wilson operator applied to $v := \boldsymbol{\lambda}^{\frac{1}{2}} u$.} The function $\boldsymbol{\lambda}$ is $R$-periodic and bounded (Lemma~\ref{lem:lipschitz-minorant-scale}(a)), so $v$ belongs to $\mathbb{H}$; it vanishes off $W$. Let $W^c := \eta\mathbb{Z}^3 \setminus W$. Lemma~\ref{lem:identities-free-wilson}(b) and the triangle inequality give
\[
\left\| \nabla^{+} v \right\| \le \|D v\| \le \left\| 1_W D v \right\| + \left\| 1_{W^c} D v \right\| .
\]
Consider first the sites of $W$. One has $D v = \boldsymbol{\lambda}^{\frac{1}{2}} D u + \left[ D, \boldsymbol{\lambda}^{\frac{1}{2}} \right] u$, and $1_W D u = f - \mathfrak{m} u - \operatorname{Pen} u$ by~\eqref{eq:wilson-operators-10}, because $u = 1_W u$ and $\operatorname{Pen} u$ vanishes off $W$. Hence
\[
1_W D v = \boldsymbol{\lambda}^{\frac{1}{2}} \left( f - \mathfrak{m} u - \operatorname{Pen} u \right) + 1_W \left[ D, \boldsymbol{\lambda}^{\frac{1}{2}} \right] u .
\]
For nearest neighbours $x$, $x'$ Lemma~\ref{lem:lipschitz-minorant-scale}(b) gives $|\boldsymbol{\lambda}(x') - \boldsymbol{\lambda}(x)| \le \varepsilon_0 \eta$, and therefore
\[
\left| \sqrt{\boldsymbol{\lambda}(x')} - \sqrt{\boldsymbol{\lambda}(x)} \right| = \frac{|\boldsymbol{\lambda}(x') - \boldsymbol{\lambda}(x)|}{\sqrt{\boldsymbol{\lambda}(x')} + \sqrt{\boldsymbol{\lambda}(x)}} \le \varepsilon_0 \eta \cdot \boldsymbol{\lambda}(x')^{-\frac{1}{2}} .
\]
Write $x' \sim x$ if $x'$ is one of the six nearest neighbours of $x$. The commutator formula~\eqref{eq:wilson-operators-8} with $\psi = \boldsymbol{\lambda}^{\frac{1}{2}}$, the inequality~\eqref{eq:wilson-operators-9} and the last bound give $\left| \left( \left[ D, \boldsymbol{\lambda}^{\frac{1}{2}} \right] u \right)(x) \right|^2 \le 3 \varepsilon_0^2 \sum_{x' \sim x} \boldsymbol{\lambda}(x')^{-1} |u(x')|^2$ at every site $x$. Every site is a nearest neighbour of six sites, so
\[
\left\| \left[ D, \boldsymbol{\lambda}^{\frac{1}{2}} \right] u \right\|^2 \le 3 \varepsilon_0^2 \eta^3 \sum_x \sum_{x' \sim x} \boldsymbol{\lambda}(x')^{-1} |u(x')|^2 = 18 \varepsilon_0^2 N^2 .
\]
Together with (ii) and (iii) this yields
\[
\left\| 1_W D v \right\| \le F + \left( L^{\frac{1}{2}} + C_1^{\frac{1}{2}} \right) X^{\frac{1}{2}} + \sqrt{18} \, \varepsilon_0 N .
\]
Consider next the sites off $W$. For $x \notin W$ one has $v(x) = 0$, and the second expression for $D$ in~\eqref{eq:wilson-operators-7} gives
\[
(D v)(x) = - \eta^{-1} \sum_{\mu = 1}^3 \left[ \Pi_{\mu}^{-} v(x + \eta e_{\mu}) + \Pi_{\mu}^{+} v(x - \eta e_{\mu}) \right] .
\]
Only the neighbours $x'$ of $x$ that lie in $W$ contribute, that is, the pairs $(x, x') \in \mathrm{bd}$, and for these $\boldsymbol{\lambda}(x') \le (L + 1) \eta$ by Lemma~\ref{lem:lipschitz-minorant-scale}(c), since $x'$ has the neighbour $x$ outside $W$. By~\eqref{eq:wilson-operators-9} and (iv),
\begin{align*}
\left\| 1_{W^c} D v \right\|^2 & \le 3 \eta^{-2} \eta^3 \sum_{(x, x') \in \mathrm{bd}} \boldsymbol{\lambda}(x') |u(x')|^2 \\
& \le 3 (L + 1) \eta^2 \sum_{(x, x') \in \mathrm{bd}} |u(x')|^2 \le 6 (L + 1) X .
\end{align*}
Adding the bounds for the two parts and recalling $K_1 = L^{\frac{1}{2}} + C_1^{\frac{1}{2}} + (6(L+1))^{\frac{1}{2}}$, one finds
\[
\left\| \nabla^{+} v \right\| \le F + K_1 X^{\frac{1}{2}} + \sqrt{18} \, \varepsilon_0 N .
\]

\emph{Step 3: from $v$ back to $u$.} Regard $u$ as a function on $\eta\mathbb{Z}^3$ that vanishes off $W$, put $\left( \boldsymbol{\lambda}^{\frac{1}{2}} \nabla^{+} u \right)_{\mu}(x) := \boldsymbol{\lambda}(x)^{\frac{1}{2}} \cdot \nabla_{\mu}^{+} u(x)$ for every site $x$ and every direction $\mu$, and
\[
\left\| \boldsymbol{\lambda}^{\frac{1}{2}} \nabla^{+} u \right\|^2 := \sum_{\mu = 1}^3 \left\| \left( \boldsymbol{\lambda}^{\frac{1}{2}} \nabla^{+} u \right)_{\mu} \right\|^2 .
\]
At every site $x$ and for every direction $\mu$,
\[
\nabla_{\mu}^{+} v(x) = \boldsymbol{\lambda}(x)^{\frac{1}{2}} \cdot \nabla_{\mu}^{+} u(x) + \left( \nabla_{\mu}^{+} \boldsymbol{\lambda}^{\frac{1}{2}} \right)(x) \, u(x + \eta e_{\mu}) , \qquad \left| \left( \nabla_{\mu}^{+} \boldsymbol{\lambda}^{\frac{1}{2}} \right)(x) \right| \le \varepsilon_0 \cdot \boldsymbol{\lambda}(x + \eta e_{\mu})^{-\frac{1}{2}} ,
\]
the bound by the estimate of Step~2. For each $\mu$ the second term therefore has norm at most $\varepsilon_0 N$, and the three directions together give $\left\| \boldsymbol{\lambda}^{\frac{1}{2}} \nabla^{+} u \right\| \le \left\| \nabla^{+} v \right\| + \sqrt{3} \, \varepsilon_0 N$. By Step~2, by the bound $X \le N F$ of~\eqref{eq:wilson-operators-13}, and since $\sqrt{18} + \sqrt{3} < 6$,
\begin{equation}\label{eq:wilson-operators-14}
\left\| \boldsymbol{\lambda}^{\frac{1}{2}} \nabla^{+} u \right\| \le F + K_1 (N F)^{\frac{1}{2}} + 6 \varepsilon_0 N .
\end{equation}

\emph{Step 4: the Poincar\'e inequality on the blocks.} On a block $E$ one has $\boldsymbol{\lambda} \ge \boldsymbol{\ell}_E/L$ (Lemma~\ref{lem:lipschitz-minorant-scale}(a)), and $\|u\|_{2,E}^2 = \|P_E u\|_{2,E}^2 + \|u - P_E u\|_{2,E}^2$ because $P_E$ is an orthogonal projection. Lemma~\ref{lem:poincare-inequality-blocks} and (i) give
\[
N^2 = \sum_E \left\| \boldsymbol{\lambda}^{-\frac{1}{2}} u \right\|_{2,E}^2 \le \sum_E \frac{L}{\boldsymbol{\ell}_E} \left( \|P_E u\|_{2,E}^2 + \|u - P_E u\|_{2,E}^2 \right) \le \frac{L}{c_1} X + \frac{L}{8} \sum_E \boldsymbol{\ell}_E \|\nabla u\|_E^2 .
\]
If both ends of a bond $\{x, x + \eta e_{\mu}\}$ lie in $E$, then $\boldsymbol{\ell}_E = \boldsymbol{\ell}(x) \le L \cdot \boldsymbol{\lambda}(x)$ by Lemma~\ref{lem:lipschitz-minorant-scale}(a). The bonds inside the blocks modulo $R$ are distinct bonds of a period cell, because every block has side at most $R$ (condition~\eqref{eq:wilson-operators-B1}). Consequently $\sum_E \boldsymbol{\ell}_E \|\nabla u\|_E^2 \le L \left\| \boldsymbol{\lambda}^{\frac{1}{2}} \nabla^{+} u \right\|^2$, and with $X \le N F$
\[
N^2 \le \frac{L}{c_1} N F + \frac{L^2}{8} \left\| \boldsymbol{\lambda}^{\frac{1}{2}} \nabla^{+} u \right\|^2 .
\]
Inserting~\eqref{eq:wilson-operators-14} and using $(r_1 + r_2 + r_3)^2 \le 3 \left( r_1^2 + r_2^2 + r_3^2 \right)$ for real numbers $r_1, r_2, r_3$, one obtains
\[
N^2 \le \frac{L}{c_1} N F + \frac{3 L^2}{8} \left( F^2 + K_1^2 N F \right) + \frac{27}{2} L^2 \varepsilon_0^2 N^2 .
\]
Since $\varepsilon_0 \le 1/(6L)$, the coefficient of the last term satisfies $\frac{27}{2} L^2 \varepsilon_0^2 \le \frac{27}{72} < \frac{1}{2}$. As $N$ is finite, this term may be absorbed into the left side, and
\[
N^2 \le 2 \left( \frac{L}{c_1} + \frac{3}{8} L^2 K_1^2 \right) N F + \frac{3}{4} L^2 F^2 .
\]
For non-negative numbers $N$, $F$, $A$, $B$ the inequality $N^2 \le A N F + B F^2$ implies $N \le \left( A + B^{\frac{1}{2}} \right) F$. Here $A = \frac{2L}{c_1} + \frac{3}{4} L^2 K_1^2$ and $B = \frac{3}{4} L^2$, so that $B^{\frac{1}{2}} = \frac{\sqrt{3}}{2} L \le L$, hence
\[
N \le \left( \frac{2L}{c_1} + \frac{3}{4} L^2 K_1^2 + L \right) F = C_I F ,
\]
which is~\eqref{eq:wilson-operators-12}.

\emph{Step 5: invertibility.} The functions $u$ and $f$ vanish off $W$, and $\sup_W \boldsymbol{\lambda} \le \sup \boldsymbol{\ell} < \infty$ by Lemma~\ref{lem:lipschitz-minorant-scale}(a) and condition~\eqref{eq:wilson-operators-B1}. Therefore~\eqref{eq:wilson-operators-12} gives
\[
\|u\| \le \left( \sup_W \boldsymbol{\lambda} \right)^{\frac{1}{2}} N \le C_I \left( \sup_W \boldsymbol{\lambda} \right)^{\frac{1}{2}} F \le C_I \left( \sup_W \boldsymbol{\lambda} \right) \|T u\| \le C_I \left( \sup \boldsymbol{\ell} \right) \|T u\| .
\]
Thus $T$ is injective and its range is closed. The adjoint problem $\mathbb{P}^{*}$ is a blocked Wilson problem with the same separation and the same penalty constants, and $T_{\mathbb{P}^{*}} = T^{*}$ (Definition~\ref{def:blocked-wilson-problems}); the same bound therefore holds for $T^{*}$, which is consequently injective. The closure of the range of $T$ is the orthogonal complement of the kernel of $T^{*}$, so the range of $T$ is dense; being closed, it is all of $\mathbb{H}_W$, and $T$ is a bijection of $\mathbb{H}_W$. For $R < \infty$ the set $W$ contains only finitely many sites modulo $R\mathbb{Z}^3$, so that $\mathbb{H}_W$ has finite dimension; in that case injectivity of $T$ by itself already implies bijectivity. For $g \in \mathbb{H}_W$ let $u := S \boldsymbol{\lambda}^{-\frac{1}{2}} g$; then $\boldsymbol{\lambda}^{\frac{1}{2}} T u = g$, and~\eqref{eq:wilson-operators-12} becomes
\[
\left\| \boldsymbol{\lambda}^{-\frac{1}{2}} S \boldsymbol{\lambda}^{-\frac{1}{2}} g \right\| \le C_I \|g\| .
\]
For $h \in \mathbb{H}_W$ and $g := \boldsymbol{\lambda}^{\frac{1}{2}} h$ this gives $\|S h\| \le \left( \sup_W \boldsymbol{\lambda} \right)^{\frac{1}{2}} C_I \left\| \boldsymbol{\lambda}^{\frac{1}{2}} h \right\| \le C_I \left( \sup_W \boldsymbol{\lambda} \right) \|h\|$, and $\sup_W \boldsymbol{\lambda} \le \sup_W \boldsymbol{\ell}$ by Lemma~\ref{lem:lipschitz-minorant-scale}(a). Finally, the adjoint problem $\mathbb{P}^{*}$ and the transposed problem $\mathbb{P}^T$ are blocked Wilson problems with the same separation, the same penalty constants and the same $W$ and $\mathbb{E}$, hence with the same function $\boldsymbol{\lambda}$ and the same constant $C_I$, and their operators are $T^{*}$ and $T^T$ (Definition~\ref{def:blocked-wilson-problems}). What has been proved for $T$ therefore holds for $T^{*}$ and for $T^T$.
\end{proof}

\subsection{Bounds with exponential weights}\label{sec:bounds-exponential-weights}

\begin{theorem}[Exponentially weighted bound]\label{thm:exponentially-weighted-bound}
Let $\mathbb{P}$ be a blocked Wilson problem with separation $s \ge 6L$ and penalty constants $(c_1, C_1, C_2)$, let $C_I$ be the constant of Theorem~\ref{thm:weighted-bound}, and put
\[
c_3 := 6 \exp(1) + 3 \exp(1) \cdot L \cdot C_2 , \qquad a_0 := \min \left\{ \frac{1}{3L}, \frac{1}{2 C_I \cdot c_3} \right\} .
\]
Let $\phi \colon \eta\mathbb{Z}^3 \to \mathbb{R}$ be bounded and $R$-periodic, with
\[
|\phi(x) - \phi(x')| \le a_0 \, d(x, x') \qquad \text{for all } x, x' \in \eta\mathbb{Z}^3 ,
\]
where $d$ is the distance of Definition~\ref{def:path-distance-units}. Then the inverse $S = T^{-1}$ (Theorem~\ref{thm:weighted-bound}) satisfies, for every $f \in \mathbb{H}_W$,
\begin{equation}\label{eq:wilson-operators-15}
\left\| \exp(\phi) \, \boldsymbol{\lambda}^{-\frac{1}{2}} S f \right\| \le 2 C_I \left\| \exp(\phi) \, \boldsymbol{\lambda}^{\frac{1}{2}} f \right\| .
\end{equation}
The same holds with $S$ replaced by $S^{*}$ or by $S^T$.
\end{theorem}

\begin{proof}
By Theorem~\ref{thm:weighted-bound} the operator $T$ is a bijection of $\mathbb{H}_W$, so that $S = T^{-1}$ exists. Write $a := a_0$; then $a \le 1/(3L) \le 1$ and $C_I \cdot c_3 \cdot a \le \frac{1}{2}$. Since $\phi$ is real, bounded and $R$-periodic, multiplication by $\exp(\phi)$ and by $\exp(-\phi)$ maps $\mathbb{H}_W$ boundedly onto itself. Define operators on $\mathbb{H}_W$ by
\[
T_{\phi} := \exp(\phi) \, T \exp(-\phi) , \qquad \operatorname{Rm} := T_{\phi} - T .
\]
Multiplication operators commute with $\exp(\pm \phi)$, so the mass term contributes nothing to $\operatorname{Rm}$, and
\[
\operatorname{Rm} = 1_W \left( \exp(\phi) \, D \exp(-\phi) - D \right) 1_W + \exp(\phi) \, \operatorname{Pen} \exp(-\phi) - \operatorname{Pen} .
\]
For sites $x$, $y$ put $q(x, y) := \exp(\phi(x) - \phi(y)) - 1$. For $x \in W$ let $E(x)$ be the block of $\mathbb{E}$ that contains $x$, and for a direction $\mu$ let $x_{\pm} := x \pm \eta e_{\mu}$. The second expression for $D$ in~\eqref{eq:wilson-operators-7} and the kernel of $\operatorname{Pen}_E$ in condition~\eqref{eq:wilson-operators-B3} show that, for $w \in \mathbb{H}_W$, $\operatorname{Rm} w = \operatorname{Rm}_{\mathrm{hop}} w + \operatorname{Rm}_{\mathrm{pen}} w$, where both terms vanish off $W$ and, for $x \in W$,
\begin{align*}
(\operatorname{Rm}_{\mathrm{hop}} w)(x) & = - \eta^{-1} \sum_{\mu = 1}^3 \left[ q(x, x_{+}) \, \Pi_{\mu}^{-} w(x_{+}) + q(x, x_{-}) \, \Pi_{\mu}^{+} w(x_{-}) \right] , \\
(\operatorname{Rm}_{\mathrm{pen}} w)(x) & = \eta^3 \sum_{x' \in E(x)} K_{E(x)}(x, x') \, q(x, x') \, w(x') .
\end{align*}
As $w$ vanishes off $W$, only the neighbours $x_{\pm}$ that lie in $W$ contribute to the first sum.

\emph{The nearest neighbour terms.} Let $x, x' \in W$ be nearest neighbours. The bond $\{x, x'\}$ is a path from $x$ to $x'$ of cost $\eta / \min(\boldsymbol{\lambda}(x), \boldsymbol{\lambda}(x'))$ (Definition~\ref{def:path-distance-units}), so that
\[
|\phi(x) - \phi(x')| \le a \, d(x, x') \le \frac{a \eta}{\min(\boldsymbol{\lambda}(x), \boldsymbol{\lambda}(x'))} =: t ,
\]
and $t \le a \le 1$ because $\boldsymbol{\lambda} \ge \eta$ (Lemma~\ref{lem:lipschitz-minorant-scale}(a)). As $|\exp(r) - 1| \le |r| \exp(|r|)$ for real $r$, it follows that $|q(x, x')| \le t \exp(a)$. By Lemma~\ref{lem:lipschitz-minorant-scale}(a), (b), the larger of the two numbers $\boldsymbol{\lambda}(x)$, $\boldsymbol{\lambda}(x')$ exceeds the smaller one by at most $\eta / s$, so their quotient is at most $1 + \eta / (s \min(\boldsymbol{\lambda}(x), \boldsymbol{\lambda}(x'))) \le 2$; hence
\[
\boldsymbol{\lambda}(x)^{\frac{1}{2}} \cdot \eta^{-1} t = \frac{a \, (\boldsymbol{\lambda}(x) \cdot \boldsymbol{\lambda}(x'))^{\frac{1}{2}}}{\min(\boldsymbol{\lambda}(x), \boldsymbol{\lambda}(x'))} \cdot \boldsymbol{\lambda}(x')^{-\frac{1}{2}} \le \sqrt{2} \, a \cdot \boldsymbol{\lambda}(x')^{-\frac{1}{2}} .
\]
At every $x \in W$ the inequality~\eqref{eq:wilson-operators-9} therefore gives, with $x' \sim x$ meaning that $x'$ is one of the six nearest neighbours of $x$,
\[
\boldsymbol{\lambda}(x) \left| (\operatorname{Rm}_{\mathrm{hop}} w)(x) \right|^2 \le 3 \cdot 2 a^2 \exp(2a) \sum_{x' \sim x, \, x' \in W} \boldsymbol{\lambda}(x')^{-1} |w(x')|^2 ,
\]
and summation over $x$ (every site is a neighbour of six sites) yields
\[
\left\| \boldsymbol{\lambda}^{\frac{1}{2}} \operatorname{Rm}_{\mathrm{hop}} w \right\|^2 \le 3 \cdot 2 a^2 \exp(2a) \cdot 6 \left\| \boldsymbol{\lambda}^{-\frac{1}{2}} w \right\|^2 ,
\]
that is, $\left\| \boldsymbol{\lambda}^{\frac{1}{2}} \operatorname{Rm}_{\mathrm{hop}} w \right\| \le 6 a \exp(a) \left\| \boldsymbol{\lambda}^{-\frac{1}{2}} w \right\|$.

\emph{The penalty terms.} Let $E$ be a block and $x, x' \in E$. By Lemma~\ref{lem:properties-path-distance}(c), $d(x, x') \le 3L$, hence $|\phi(x) - \phi(x')| \le 3 L a \le 1$ and $|q(x, x')| \le 3 L a \exp(1)$. Condition~\eqref{eq:wilson-operators-Pen3}, the Cauchy--Schwarz inequality and $\eta^3 \, \#|E| = \boldsymbol{\ell}_E^3$ give, for $x \in E$,
\[
\left| (\operatorname{Rm}_{\mathrm{pen}} w)(x) \right| \le C_2 \cdot \boldsymbol{\ell}_E^{-4} \cdot 3 L a \exp(1) \cdot \boldsymbol{\ell}_E^{\frac{3}{2}} \|w\|_{2,E} .
\]
Since $\boldsymbol{\lambda} \le \boldsymbol{\ell}_E$ on $E$ (Lemma~\ref{lem:lipschitz-minorant-scale}(a)), and since $\|g\|_{2,E} \le \boldsymbol{\ell}_E^{\frac{3}{2}} \sup_E |g|$ for every function $g$ on $E$, it follows that
\begin{align*}
\left\| \boldsymbol{\lambda}^{\frac{1}{2}} \operatorname{Rm}_{\mathrm{pen}} w \right\|_{2,E} & \le \boldsymbol{\ell}_E^{\frac{1}{2}} \cdot \boldsymbol{\ell}_E^{\frac{3}{2}} \cdot C_2 \cdot \boldsymbol{\ell}_E^{-\frac{5}{2}} \cdot 3 \exp(1) \cdot L a \, \|w\|_{2,E} \\
& = 3 \exp(1) \cdot L \cdot C_2 \cdot a \cdot \boldsymbol{\ell}_E^{-\frac{1}{2}} \|w\|_{2,E} \le 3 \exp(1) \cdot L \cdot C_2 \cdot a \left\| \boldsymbol{\lambda}^{-\frac{1}{2}} w \right\|_{2,E} .
\end{align*}
On $E$ the function $\operatorname{Rm}_{\mathrm{pen}} w$ depends only on the values of $w$ on $E$; summing the squares of these bounds over the blocks modulo $R$ (condition~\eqref{eq:wilson-operators-B1}) gives $\left\| \boldsymbol{\lambda}^{\frac{1}{2}} \operatorname{Rm}_{\mathrm{pen}} w \right\| \le 3 \exp(1) \cdot L \cdot C_2 \cdot a \left\| \boldsymbol{\lambda}^{-\frac{1}{2}} w \right\|$. Together with the bound for the nearest neighbour terms and $\exp(a) \le \exp(1)$,
\[
\left\| \boldsymbol{\lambda}^{\frac{1}{2}} \operatorname{Rm} w \right\| \le \left( 6 \exp(1) + 3 \exp(1) \cdot L \cdot C_2 \right) a \left\| \boldsymbol{\lambda}^{-\frac{1}{2}} w \right\| = c_3 \cdot a \left\| \boldsymbol{\lambda}^{-\frac{1}{2}} w \right\| \qquad (w \in \mathbb{H}_W) .
\]

\emph{Conclusion.} Let $f \in \mathbb{H}_W$, $u := S f$ and $w := \exp(\phi) u$, an element of $\mathbb{H}_W$. Then $T_{\phi} w = \exp(\phi) T u = \exp(\phi) f$, so that $T w = \exp(\phi) f - \operatorname{Rm} w$. The bound~\eqref{eq:wilson-operators-12}, applied to $w$ in place of $u$, yields
\[
\left\| \boldsymbol{\lambda}^{-\frac{1}{2}} w \right\| \le C_I \left\| \boldsymbol{\lambda}^{\frac{1}{2}} T w \right\| \le C_I \left\| \boldsymbol{\lambda}^{\frac{1}{2}} \exp(\phi) f \right\| + C_I \cdot c_3 \cdot a \left\| \boldsymbol{\lambda}^{-\frac{1}{2}} w \right\| ,
\]
where $C_I \cdot c_3 \cdot a \le \frac{1}{2}$. As $\left\| \boldsymbol{\lambda}^{-\frac{1}{2}} w \right\|$ is finite, it follows that $\left\| \boldsymbol{\lambda}^{-\frac{1}{2}} w \right\| \le 2 C_I \left\| \exp(\phi) \, \boldsymbol{\lambda}^{\frac{1}{2}} f \right\|$, and this is~\eqref{eq:wilson-operators-15}, because $\boldsymbol{\lambda}^{-\frac{1}{2}} w = \exp(\phi) \, \boldsymbol{\lambda}^{-\frac{1}{2}} S f$. The adjoint problem $\mathbb{P}^{*}$ and the transposed problem $\mathbb{P}^T$ are blocked Wilson problems with the same separation, the same penalty constants and the same $W$ and $\mathbb{E}$ (Definition~\ref{def:blocked-wilson-problems}); hence they have the same $\boldsymbol{\lambda}$, $d$, $C_I$, $c_3$ and $a_0$, and by Theorem~\ref{thm:weighted-bound} their operators $T^{*}$ and $T^T$ have the inverses $S^{*}$ and $S^T$. The argument above, applied to these two problems, proves~\eqref{eq:wilson-operators-15} for $S^{*}$ and for $S^T$.
\end{proof}

\subsection{Interior regularity and pointwise bounds for the inverse}\label{sec:interior-regularity-pointwise}

\subsubsection{Interior estimates uniform in the mesh}\label{sec:interior-estimates-uniform}

\begin{proposition}[Interior estimates]\label{prop:interior-estimates}
Let $x_0\in\eta\mathbb{Z}^3$, let $r\ge8\eta$ and $0\le\mathfrak{m}\le1/\eta$, and let $B$, $B^{+}$ and $B'$ be the sets of the sites $x\in\eta\mathbb{Z}^3$ with $|x-x_0|_{\infty}<r$, with $|x-x_0|_{\infty}<r+\eta$ and with $|x-x_0|_{\infty}\le r/4$, respectively. Let $u\colon B^{+}\to\mathbb{C}^4$ and $g\colon B\to\mathbb{C}^4$ satisfy
\[
\left((D+\mathfrak{m})u\right)(x)=g(x)\qquad\text{for every }x\in B;
\]
the left side is defined, because the nearest neighbours of a site of $B$ belong to $B^{+}$. Let $\boldsymbol{\alpha}\in(0,1)$ be the H\"older order of~\eqref{eq:wilson-operators-1}, and let $C_R=7\cdot10^4\,C_G/(1-\boldsymbol{\alpha})$, where $C_G$ is the constant of Proposition~\ref{prop:lattice-green-function}. Then for all $x,\tilde x\in B'$ the following holds.
\begin{enumerate}
\item[(a)] The value of $u$ at $x$ satisfies
\[
|u(x)|\le C_R\left(r\,\|g\|_{\infty,B}+r^{-\frac{3}{2}}\,\|u\|_{2,B^{+}}\right).
\]
\item[(b)] The values of $u$ at $x$ and at $\tilde x$ satisfy
\[
|u(x)-u(\tilde x)|\le C_R\left(\frac{|x-\tilde x|_{\infty}}{r}\right)^{\boldsymbol{\alpha}}\left(r\,\|g\|_{\infty,B}+r^{-\frac{3}{2}}\,\|u\|_{2,B^{+}}\right).
\]
\end{enumerate}
\end{proposition}

\begin{proof}
Let $\chi\colon\eta\mathbb{Z}^3\to[0,1]$ be the function
\[
\chi(x):=\min\left\{1,\max\left\{0,\,2-\frac{2\,|x-x_0|_{\infty}}{r}\right\}\right\}.
\]
It equals $1$ where $|x-x_0|_{\infty}\le r/2$ and $0$ outside $B$, and $|\chi(x)-\chi(x')|\le2\eta/r$ for any two nearest neighbours $x,x'$. Extend $u$ by zero to $\eta\mathbb{Z}^3\setminus B^{+}$ and put $v:=\chi u$; this function has finite support. At every site,
\[
(D+\mathfrak{m})v=\chi\cdot(D+\mathfrak{m})u+[D,\chi]\,u.
\]
If $x\in B$, the nearest neighbours of $x$ lie in $B^{+}$, so that $\left((D+\mathfrak{m})u\right)(x)=g(x)$; if $x\notin B$, then $\chi(x)=0$. With $\chi g$ extended by zero outside $B$, it follows that
\[
(D+\mathfrak{m})v=\chi g+k,\qquad k:=[D,\chi]\,u.
\]
For a finitely supported function $q\colon\eta\mathbb{Z}^3\to\mathbb{C}^4$ put
\[
(G*q)(x):=\eta^3\sum_{x'\in\eta\mathbb{Z}^3}G(x-x')\,q(x')\qquad(x\in\eta\mathbb{Z}^3),
\]
where $G=G^{\eta}_{\mathfrak{m}}$ is the function of Proposition~\ref{prop:lattice-green-function}. The functions $\chi g$ and $k=(D+\mathfrak{m})v-\chi g$ have finite support, and since $v$ has finite support, Proposition~\ref{prop:lattice-green-function}(c) gives
\[
v=G*\left((D+\mathfrak{m})v\right)=G*(\chi g)+G*k.
\]

\emph{The term $G*k$.} By~\eqref{eq:wilson-operators-8} of Lemma~\ref{lem:identities-free-wilson}(c), $k(x)$ is the sum, over the directions $\mu$ and over the two signs, of the vectors
\[
-\eta^{-1}\left(\chi(x\pm\eta e_{\mu})-\chi(x)\right)\Pi_{\mu}^{\mp}\,u(x\pm\eta e_{\mu}).
\]
For nearest neighbours $x,x'$ the difference $\chi(x')-\chi(x)$ vanishes unless the larger of the numbers $|x-x_0|_{\infty}$, $|x'-x_0|_{\infty}$ exceeds $r/2$ and the smaller one is less than $r$. Therefore $k$ is supported in the set $\left\{x:r/2-\eta<|x-x_0|_{\infty}<r+\eta\right\}$, and only the values of $u$ on $B^{+}$ enter $k$. The matrices $\Pi_{\mu}^{\pm}$ are orthogonal projections (Convention~\ref{it:wilson-operators-3}), and every site has six nearest neighbours. With the Cauchy--Schwarz inequality and the bound
\[
\eta^3\,\#|B^{+}|\le(2r+3\eta)^3\le(2.4\,r)^3\qquad(\text{recall }\eta\le r/8)
\]
one obtains
\begin{align*}
\|k\|_{L^1}:=\eta^3\sum_{x\in\eta\mathbb{Z}^3}|k(x)|
&\le\eta^{-1}\cdot\frac{2\eta}{r}\cdot6\,\eta^3\sum_{x\in B^{+}}|u(x)|\\
&\le\frac{12}{r}\left(\eta^3\,\#|B^{+}|\right)^{\frac{1}{2}}\|u\|_{2,B^{+}}\le45\,r^{\frac{1}{2}}\,\|u\|_{2,B^{+}}.
\end{align*}
If $z\in B'$ and $x''\in\operatorname{supp}k$, then $|z-x''|_{\infty}>r/4-\eta\ge r/8$. Hence Proposition~\ref{prop:lattice-green-function}(b) gives, for every $x\in B'$ and every direction $\mu$,
\begin{align*}
|(G*k)(x)|&\le C_G\cdot\left(\frac{r}{8}\right)^{-2}\|k\|_{L^1}\le2880\,C_G\cdot r^{-\frac{3}{2}}\,\|u\|_{2,B^{+}},\\
\left|\nabla_{\mu}^{\pm}(G*k)(x)\right|&\le C_G\cdot\left(\frac{r}{8}\right)^{-3}\|k\|_{L^1}\le23040\,C_G\cdot r^{-\frac{5}{2}}\,\|u\|_{2,B^{+}}.
\end{align*}
Now let $x,\tilde x\in B'$ and $h:=|x-\tilde x|_{\infty}$; then $h\le r/2$. A monotone coordinate path from $x$ to $\tilde x$ lies in the cube $B'$ and has at most $3h/\eta$ bonds, and along a bond $\{z,z\pm\eta e_{\mu}\}$ with $z\in B'$ the function $G*k$ changes by $\pm\eta\cdot\nabla_{\mu}^{\pm}(G*k)(z)$. Summing the second of the last two bounds over the bonds of the path, and using $h/r\le1$, one finds
\begin{align*}
|(G*k)(x)-(G*k)(\tilde x)|&\le69120\,C_G\cdot h\cdot r^{-\frac{5}{2}}\,\|u\|_{2,B^{+}}\\
&\le69120\,C_G\cdot\left(\frac{h}{r}\right)^{\boldsymbol{\alpha}}r^{-\frac{3}{2}}\,\|u\|_{2,B^{+}}.
\end{align*}

\emph{The term $G*(\chi g)$.} One has $|\chi g|\le\|g\|_{\infty,B}$, and $|x-x'|_{\infty}\le5r/4$ for $x\in B'$ and $x'\in B$. By Proposition~\ref{prop:lattice-green-function}(b) and the first bound of Lemma~\ref{lem:two-lattice-sums}(a) with $\varrho=5r/4$, every $x\in B'$ satisfies
\[
|(G*(\chi g))(x)|\le27\cdot\frac{5r}{4}\,C_G\cdot\|g\|_{\infty,B}\le34\,C_G\cdot r\,\|g\|_{\infty,B}.
\]
For $x,\tilde x\in B'$ with $h=|x-\tilde x|_{\infty}>0$, Lemma~\ref{lem:two-lattice-sums}(b), followed by the second bound of Lemma~\ref{lem:two-lattice-sums}(a) with $\varrho=5r/4$ (once about $x$ and once about $\tilde x$), gives
\begin{align*}
|(G*(\chi g))(x)-(G*(\chi g))(\tilde x)|&\le24\,C_G\cdot h^{\boldsymbol{\alpha}}\cdot2\cdot27\left(\frac{5r}{4}\right)^{1-\boldsymbol{\alpha}}\frac{\|g\|_{\infty,B}}{1-\boldsymbol{\alpha}}\\
&\le\frac{1620\,C_G}{1-\boldsymbol{\alpha}}\left(\frac{h}{r}\right)^{\boldsymbol{\alpha}}r\,\|g\|_{\infty,B}.
\end{align*}

\emph{Conclusion.} As $\chi=1$ on the set $\left\{x:|x-x_0|_{\infty}\le r/2\right\}$, which contains $B'$, one has $v=u$ on $B'$. Each of the numbers $34$, $2880$, $1620/(1-\boldsymbol{\alpha})$ and $69120$ is at most $7\cdot10^4/(1-\boldsymbol{\alpha})=C_R/C_G$. Adding the bounds obtained for $G*(\chi g)$ and for $G*k$ therefore gives (a) and (b).
\end{proof}

\subsubsection{Pointwise bounds from block to block}\label{sec:pointwise-bounds-block}

\begin{theorem}[Sup and H\"older bounds for the inverse]\label{thm:sup-holder-bounds}
Let $\mathbb{P}$ be a blocked Wilson problem with separation $s\ge s_0=48/a_0$, and let $E,E'\in\mathbb{E}$. Let $f\in\mathbb{F}(\eta,R,\boldsymbol{\theta})$ be bounded, that is, $\|f\|_{\infty}:=\sup_{z\in\eta\mathbb{Z}^3}|f(z)|<\infty$, and supported in the orbit $\bigcup_{n\in\mathbb{Z}^3}(E'+Rn)$ of $E'$ (the set $E'$ itself if $R=\infty$); then $f\in\mathbb{H}_W$. Let $u:=Sf$. Then every $x\in E$ satisfies
\begin{equation}\label{eq:wilson-operators-16}
|u(x)|+[u]_{\boldsymbol{\alpha};x}\le C_3\cdot\min\left\{\boldsymbol{\ell}_E,\boldsymbol{\ell}_{E'}\right\}\cdot\exp\left(-\boldsymbol{\gamma}_3\cdot d(x,E')\right)\cdot\|f\|_{\infty},
\end{equation}
where $[u]_{\boldsymbol{\alpha};x}$ is the H\"older seminorm of Convention~\ref{it:wilson-operators-7}, $d$ is the distance of Definition~\ref{def:path-distance-units}, $\boldsymbol{\gamma}_3=a_0/2$, and $C_3$ is the number defined in~\eqref{eq:wilson-operators-20}. In the definitions of $a_0$, $s_0$ and $C_3$ the penalty constants $(c_1,C_1,C_2)$ are those of $\mathbb{P}$; when they are the constants~\eqref{eq:wilson-operators-2}, these three numbers are functions of $(L,b,\boldsymbol{\alpha})$ alone. The same holds with $S^{*}$ or $S^T$ in place of $S$ (for $S^T$, with $f\in\mathbb{F}(\eta,R,-\boldsymbol{\theta})$).
\end{theorem}

\begin{proof}
By~\eqref{eq:wilson-operators-4} one has $a_0\le1/(3L)$; hence $s\ge s_0\ge144L\ge6L$, and Theorem~\ref{thm:exponentially-weighted-bound} and Lemmas~\ref{lem:lipschitz-minorant-scale} and~\ref{lem:properties-path-distance} apply. By linearity it suffices to treat the case $\|f\|_{\infty}=1$. The argument is carried out on the covering lattice $\eta\mathbb{Z}^3$, in the sense of Definition~\ref{def:path-distance-units}: sites and blocks are those of $\eta\mathbb{Z}^3$, a block being a member of the $R$-periodic partition $\mathbb{E}$ of $W$, and $|z-z'|_{\infty}$ denotes the sup distance in $\eta\mathbb{Z}^3$ of sites $z,z'$, also when $R<\infty$. As $d(x,E'+Rn)=d(x,E')$ for $n\in\mathbb{Z}^3$, the number $d(x,E')$ is the distance from $x$ to the orbit of $E'$.

\emph{Step 0: weighted square sums.} By~\eqref{eq:wilson-operators-B1} every block has side at most $R$, so that two distinct sites of one block are never congruent modulo $R\mathbb{Z}^3$. Let $q\ge0$ be an $R$-periodic function on $\eta\mathbb{Z}^3$ and $E''$ a block. Then the sum of $q$ over $E''$ is at most the sum of $q$ over a period cell, and the sum of $q$ over the sites of a period cell that lie in the orbit of $E''$ is at most the sum of $q$ over $E''$. The functions $|f|$ and $|u|$ are $R$-periodic, because $f$ and $u$ belong to $\mathbb{F}(\eta,R,\boldsymbol{\theta})$, and so are $\boldsymbol{\lambda}$ (Lemma~\ref{lem:lipschitz-minorant-scale}(a)) and $d(\cdot,E')$. Since $f$ vanishes outside the orbit of $E'$, the second statement, applied to $q=\boldsymbol{\lambda}\cdot|f|^2$ and $E''=E'$, together with $\boldsymbol{\lambda}\le\boldsymbol{\ell}=\boldsymbol{\ell}_{E'}$ on $E'$ (Lemma~\ref{lem:lipschitz-minorant-scale}(a)), with $|f|\le1$ and with the number $\left(\boldsymbol{\ell}_{E'}/\eta\right)^3$ of the sites of $E'$ (Convention~\ref{it:wilson-operators-4}), gives
\[
\left\|\boldsymbol{\lambda}^{\frac{1}{2}}\cdot f\right\|\le\boldsymbol{\ell}_{E'}^{\frac{1}{2}}\cdot\boldsymbol{\ell}_{E'}^{\frac{3}{2}}=\boldsymbol{\ell}_{E'}^2.
\]
For $M>0$ let $\phi_M\colon\eta\mathbb{Z}^3\to\mathbb{R}$ be the function $\phi_M:=a_0\min\left\{d(\cdot,E'),M\right\}$. It is bounded and $R$-periodic, it vanishes on the support of $f$, and $|\phi_M(x)-\phi_M(x')|\le a_0\,d(x,x')$ by the triangle inequality for $d$. Theorem~\ref{thm:exponentially-weighted-bound}, that is, inequality~\eqref{eq:wilson-operators-15}, yields
\begin{align*}
\left\|\exp(\phi_M)\cdot\boldsymbol{\lambda}^{-\frac{1}{2}}\cdot u\right\|&\le2C_I\cdot\left\|\exp(\phi_M)\cdot\boldsymbol{\lambda}^{\frac{1}{2}}\cdot f\right\|\\
&=2C_I\cdot\left\|\boldsymbol{\lambda}^{\frac{1}{2}}\cdot f\right\|\le2C_I\cdot\boldsymbol{\ell}_{E'}^2.
\end{align*}
As $M\to\infty$, the functions $\exp(\phi_M)$ increase to $\exp\left(a_0\,d(\cdot,E')\right)$, and the monotone convergence theorem gives the first line of
\begin{align}
\left\|\exp\left(a_0\,d(\cdot,E')\right)\cdot\boldsymbol{\lambda}^{-\frac{1}{2}}\cdot u\right\|&\le U_0:=2C_I\cdot\boldsymbol{\ell}_{E'}^2,\notag\\
\|u\|_{2,E''}&\le\boldsymbol{\ell}_{E''}^{\frac{1}{2}}\cdot\exp\left(-a_0\,d(E'',E')\right)U_0\qquad\text{for every block }E''.\label{eq:wilson-operators-17}
\end{align}
The second line follows from the first and from the first statement above, applied to $q=\exp\left(2a_0\,d(\cdot,E')\right)\cdot\boldsymbol{\lambda}^{-1}\cdot|u|^2$, since $\boldsymbol{\lambda}\le\boldsymbol{\ell}_{E''}$ and $d(\cdot,E')\ge d(E'',E')$ on $E''$.

\emph{Step 1: low blocks.} A block $E$ is called low if $\boldsymbol{\ell}_E\le L^2\eta$, and high if $\boldsymbol{\ell}_E\ge L^3\eta$. By~\eqref{eq:wilson-operators-B1} and Convention~\ref{it:wilson-operators-4}, $\boldsymbol{\ell}_E=L^i\eta$ with an integer $i\ge1$, so every block is low or high. Let $E$ be low and $x\in E$. The square of the norm in the first line of~\eqref{eq:wilson-operators-17} is the sum, over a period cell, of the $R$-periodic function $\eta^3\exp\left(2a_0\,d(\cdot,E')\right)\cdot\boldsymbol{\lambda}^{-1}\cdot|u|^2$; hence the value of this function at $x$ is at most $U_0^2$. Together with $\boldsymbol{\lambda}(x)\le\boldsymbol{\ell}_E$ and $\eta\ge\boldsymbol{\ell}_E/L^2$ this gives
\begin{align*}
|u(x)|=\eta^{-\frac{3}{2}}\,\|u\|_{2,\{x\}}&\le\eta^{-\frac{3}{2}}\cdot\boldsymbol{\lambda}(x)^{\frac{1}{2}}\cdot\exp\left(-a_0\,d(x,E')\right)U_0\\
&\le L^3\cdot\boldsymbol{\ell}_E^{-1}\cdot\exp\left(-a_0\,d(x,E')\right)U_0.
\end{align*}

\emph{Step 2: high blocks.} Let $E$ be high and $x\in E$. Put $r:=\boldsymbol{\ell}_E/4$; then $r\ge31\eta$. Let $B$, $B^{+}$ and $B'$ be the sets of Proposition~\ref{prop:interior-estimates} with $x_0:=x$. Since $r+\eta<4\boldsymbol{\ell}_E$, Lemma~\ref{lem:lipschitz-minorant-scale}(e) gives $B^{+}\subset W$. As $u$ vanishes off $W$ and $Tu=f$, with $T=[D+\mathfrak{m}]_W+\operatorname{Pen}$ as in~\eqref{eq:wilson-operators-10}, every $z\in W$ satisfies
\[
\left((D+\mathfrak{m})u\right)(z)=\left([D+\mathfrak{m}]_Wu\right)(z)=g(z):=f(z)-(\operatorname{Pen}\,u)(z).
\]
Moreover $\mathfrak{m}\cdot\eta\le\mathfrak{m}\cdot\boldsymbol{\ell}_E/L\le1$ by~\eqref{eq:wilson-operators-B4}, because $\boldsymbol{\ell}_E\ge L\eta$. Thus Proposition~\ref{prop:interior-estimates} applies to the restrictions of $u$ to $B^{+}$ and of $g$ to $B$.

Let $E^{*}(x)$ be the union of the blocks $E''$ that meet $B^{+}$. For such a block and a site $z\in E''\cap B^{+}$ one has $|x-z|_{\infty}<r+\eta\le\boldsymbol{\ell}(x)$, and Lemma~\ref{lem:lipschitz-minorant-scale}(d) gives $\boldsymbol{\ell}_E/L\le\boldsymbol{\ell}_{E''}\le L\cdot\boldsymbol{\ell}_E$. The closed cubes of these blocks have pairwise disjoint interiors (Convention~\ref{it:wilson-operators-4}) and meet the closed cube of side $2(r+\eta)\le0.6\,\boldsymbol{\ell}_E$ centred at $x$; hence their number is at most $(0.6L+2)^3\le n_L=(L+2)^3$. By Lemma~\ref{lem:properties-path-distance}(c), applied first to $x$ and a site of $E''\cap B^{+}$ and then to two sites of $E''$, every $z\in E''$ satisfies $d(x,z)\le4L+3L$; moreover $7La_0\le7/3$. Consequently $d(E'',E')\ge d(x,E')-7L$, and~\eqref{eq:wilson-operators-17} gives
\[
\|u\|_{2,E^{*}(x)}\le n_L^{\frac{1}{2}}\cdot\left(L\cdot\boldsymbol{\ell}_E\right)^{\frac{1}{2}}\cdot\exp\left(\frac{7}{3}\right)\cdot\exp\left(-a_0\,d(x,E')\right)U_0.
\]
For $z\in E''$ one has $(\operatorname{Pen}\,u)(z)=\eta^3\sum_{x'\in E''}K_{E''}(z,x')\,u(x')$ by~\eqref{eq:wilson-operators-B3}. By~\eqref{eq:wilson-operators-Pen3}, the Cauchy--Schwarz inequality and $\boldsymbol{\ell}_{E''}\ge\boldsymbol{\ell}_E/L$,
\[
\|\operatorname{Pen}\,u\|_{\infty,E''}\le C_2\cdot\boldsymbol{\ell}_{E''}^{-4}\cdot\boldsymbol{\ell}_{E''}^{\frac{3}{2}}\cdot\|u\|_{2,E''}\le C_2\cdot L^{\frac{5}{2}}\cdot\boldsymbol{\ell}_E^{-\frac{5}{2}}\cdot\|u\|_{2,E^{*}(x)}.
\]
Let $A$ be the set of the sites $z\in\eta\mathbb{Z}^3$ with $d(z,E')\le4L$, and let $1_A$ be its indicator function (Convention~\ref{it:wilson-operators-5}); thus $1_A(x)=1$ if $d(x,E')\le4L$, and $1_A(x)=0$ otherwise. If no site of the orbit of $E'$ belongs to $B$, then $f=0$ on $B$. If a site $z$ of this orbit belongs to $B$, then $|x-z|_{\infty}<r\le\boldsymbol{\ell}(x)$, and Lemma~\ref{lem:properties-path-distance}(c) yields $d(x,E')\le d(x,z)\le4L$. As $|f|\le1$, in both cases $\|f\|_{\infty,B}\le1_A(x)$. Since $B\subset B^{+}\subset E^{*}(x)$ (recall $B^{+}\subset W$) and $r=\boldsymbol{\ell}_E/4$, the number
\[
\mathscr{N}(x):=r\,\|g\|_{\infty,B}+r^{-\frac{3}{2}}\,\|u\|_{2,B^{+}}
\]
satisfies
\[
\mathscr{N}(x)\le\frac{\boldsymbol{\ell}_E}{4}\cdot1_A(x)+\left(\frac{C_2\cdot L^{\frac{5}{2}}}{4}+8\right)\cdot\boldsymbol{\ell}_E^{-\frac{3}{2}}\cdot\|u\|_{2,E^{*}(x)},
\]
and the bound on $\|u\|_{2,E^{*}(x)}$, with $U_0=2C_I\cdot\boldsymbol{\ell}_{E'}^2$, turns this into
\begin{align}
\mathscr{N}(x)&\le\frac{\boldsymbol{\ell}_E}{4}\cdot1_A(x)+C_3^{\mathrm{b}}\cdot\boldsymbol{\ell}_E^{-1}\cdot\boldsymbol{\ell}_{E'}^2\cdot\exp\left(-a_0\,d(x,E')\right),\label{eq:wilson-operators-18}\\
C_3^{\mathrm{b}}&:=\left(\frac{C_2\cdot L^{\frac{5}{2}}}{4}+8\right)n_L^{\frac{1}{2}}\cdot L^{\frac{1}{2}}\cdot\exp\left(\frac{7}{3}\right)\cdot2C_I.\notag
\end{align}

\emph{Step 3: comparison of the scales.} Let $\mathfrak{d}:=d(x,E')$ and $m:=\min\left\{\boldsymbol{\ell}_E,\boldsymbol{\ell}_{E'}\right\}$. By Lemma~\ref{lem:properties-path-distance}(a), $\boldsymbol{\ell}_{E'}\le L\cdot\boldsymbol{\ell}_E\cdot\exp\left(\frac{\mathfrak{d}}{s}\right)$ and $\boldsymbol{\ell}_E\le L\cdot\boldsymbol{\ell}_{E'}\cdot\exp\left(\frac{\mathfrak{d}}{s}\right)$. Hence $\boldsymbol{\ell}_E^{-1}\cdot\boldsymbol{\ell}_{E'}^2\le L^2m\cdot\exp\left(\frac{2\mathfrak{d}}{s}\right)$, and $2/s\le a_0/24$. If $\mathfrak{d}\le4L$, then $\boldsymbol{\ell}_E\le2Lm$ (since $4L/s\le1/36$, so that $\exp\left(\frac{4L}{s}\right)\le2$) and $1\le\exp\left(\frac{4}{3}\right)\cdot\exp(-a_0\cdot\mathfrak{d})$. As $\boldsymbol{\gamma}_3=a_0/2\le a_0-a_0/24$, these facts and~\eqref{eq:wilson-operators-18} give the first line of
\begin{align}
\mathscr{N}(x)&\le C_3^{\mathrm{n}}\cdot m\cdot\exp\left(-\boldsymbol{\gamma}_3\cdot\mathfrak{d}\right)\qquad\text{if $E$ is high},\notag\\
C_3^{\mathrm{n}}&:=\frac{L}{2}\exp\left(\frac{4}{3}\right)+L^2C_3^{\mathrm{b}},\notag\\
|u(x)|&\le C_3''\cdot m\cdot\exp\left(-\boldsymbol{\gamma}_3\cdot\mathfrak{d}\right)\qquad\text{for every block $E$},\notag\\
C_3''&:=\max\left\{C_R\cdot C_3^{\mathrm{n}},\,2L^5C_I\right\}.\label{eq:wilson-operators-19}
\end{align}
The third line follows from the first line and Proposition~\ref{prop:interior-estimates}(a) if $E$ is high, and from Step~1 and the facts above if $E$ is low.

\emph{Step 4: H\"older quotients.} Let $\tilde x\in W$ with $0<h:=|x-\tilde x|_{\infty}\le\boldsymbol{\ell}_E$, and let $\tilde E$ be the block of $\tilde x$. Lemmas~\ref{lem:lipschitz-minorant-scale}(d) and~\ref{lem:properties-path-distance}(c) give $\boldsymbol{\ell}_{\tilde E}\le L\cdot\boldsymbol{\ell}_E$ and $d(\tilde x,E')\ge\mathfrak{d}-4L$; moreover $\exp(4L\cdot\boldsymbol{\gamma}_3)\le\exp\left(\frac{2}{3}\right)$. If $E$ is high and $h\le\boldsymbol{\ell}_E/16=r/4$, then $\tilde x\in B'$, and Proposition~\ref{prop:interior-estimates}(b), together with $\left(\boldsymbol{\ell}_E/r\right)^{\boldsymbol{\alpha}}=4^{\boldsymbol{\alpha}}\le4$, gives
\[
\frac{\boldsymbol{\ell}_E^{\boldsymbol{\alpha}}\cdot|u(x)-u(\tilde x)|}{h^{\boldsymbol{\alpha}}}\le4C_R\cdot\mathscr{N}(x).
\]
In the remaining cases $h\ge\boldsymbol{\ell}_E/16$, or $h\ge\eta\ge\boldsymbol{\ell}_E/L^2$; in both, $1\le\boldsymbol{\ell}_E/h\le\max\left\{16,L^2\right\}=L^2$, so that $\left(\boldsymbol{\ell}_E/h\right)^{\boldsymbol{\alpha}}\le L^2$. Applying~\eqref{eq:wilson-operators-19} at $x$ and at $\tilde x$, where $\min\left\{\boldsymbol{\ell}_{\tilde E},\boldsymbol{\ell}_{E'}\right\}\le Lm$, one obtains
\[
\frac{\boldsymbol{\ell}_E^{\boldsymbol{\alpha}}\cdot|u(x)-u(\tilde x)|}{h^{\boldsymbol{\alpha}}}\le L^2C_3''\left(1+L\exp\left(\frac{2}{3}\right)\right)m\cdot\exp\left(-\boldsymbol{\gamma}_3\cdot\mathfrak{d}\right).
\]

\emph{Conclusion.} Since $\boldsymbol{\ell}(x)=\boldsymbol{\ell}_E$, the seminorm $[u]_{\boldsymbol{\alpha};x}$ of Convention~\ref{it:wilson-operators-7} is the supremum of the quotients bounded in Step~4. Together with the bound on $\mathscr{N}(x)$ and the bound on $|u(x)|$ in~\eqref{eq:wilson-operators-19}, Steps~3 and~4 give~\eqref{eq:wilson-operators-16} with
\begin{equation}\label{eq:wilson-operators-20}
C_3:=C_3''+4C_R\cdot C_3^{\mathrm{n}}+L^2C_3''\left(1+L\exp\left(\frac{2}{3}\right)\right).
\end{equation}
The adjoint problem $\mathbb{P}^{*}$ and the transposed problem $\mathbb{P}^T$ of Definition~\ref{def:blocked-wilson-problems} are blocked Wilson problems with the same separation, the same penalty constants and the same $W$ and $\mathbb{E}$ as $\mathbb{P}$, hence with the same functions $\boldsymbol{\ell}$, $\boldsymbol{\lambda}$ and the same distance $d$; their operators are $T^{*}$ and $T^T$, whose inverses are $S^{*}$ and $S^T$ by Theorem~\ref{thm:weighted-bound}. The argument above, applied to these two problems, gives the bound for $S^{*}$ and for $S^T$.
\end{proof}

\subsection{Bounded right-hand sides of arbitrary support}\label{sec:bounded-right-hand}

Throughout this subsection $\mathbb{P}$ is a blocked Wilson problem (Definition~\ref{def:blocked-wilson-problems}) with penalty constants $(c_1,C_1,C_2)$, and $a_0$, $s_0=48/a_0$ and $C_3$ are the numbers of Theorem~\ref{thm:sup-holder-bounds}, namely those given by~\eqref{eq:wilson-operators-4} and~\eqref{eq:wilson-operators-20} when these are formed with the penalty constants of $\mathbb{P}$. The separation $s$ of $\mathbb{P}$ satisfies $s\ge s_0$. The mesh $\eta$, the side $R$, the character $\boldsymbol{\theta}$, the set $W$, the set of blocks $\mathbb{E}$, the penalty $\operatorname{Pen}$ and the mass $\mathfrak{m}$ are those of $\mathbb{P}$, and the operators $T$ and $S=T^{-1}$ are those of Definition~\ref{def:blocked-wilson-problems} and Theorem~\ref{thm:weighted-bound}. As in~\eqref{eq:wilson-operators-6}, $\boldsymbol{\gamma}_4=a_0/4$; moreover
\[
\tilde{C}:=2C_3\cdot C_{\mathrm{sh}}(a_0/4),
\]
where $C_{\mathrm{sh}}$ is the function defined in~\eqref{eq:wilson-operators-11}. Hence $\boldsymbol{\gamma}_4$ and $\tilde{C}$, like $a_0$, $s_0$ and $C_3$, are formed with the penalty constants of $\mathbb{P}$; for the penalty constants of~\eqref{eq:wilson-operators-2} these five numbers depend on $(L,b,\boldsymbol{\alpha})$ only.

The operator $T$ extends to the functions $v\in\mathbb{F}(\eta,R,\boldsymbol{\theta})$ that vanish off $W$: for such $v$ let $Tv\colon\eta\mathbb{Z}^3\to\mathbb{C}^4$ be the function
\[
(Tv)(x):=\begin{cases}\left((D+\mathfrak{m})v\right)(x)+(\operatorname{Pen}\,v)(x), & x\in W,\\ 0, & x\notin W,\end{cases}
\]
where $D$ is the free Wilson operator~\eqref{eq:wilson-operators-7}, which acts on $\mathbb{F}(\eta,R,\boldsymbol{\theta})$, and $\operatorname{Pen}=\sum_{E\in\mathbb{E}}\operatorname{Pen}_E$ acts on all functions $\eta\mathbb{Z}^3\to\mathbb{C}^4$ by~\eqref{eq:wilson-operators-B3}. For $v\in\mathbb{H}_W$ the function $Tv$ is the one given by~\eqref{eq:wilson-operators-10}, because $[D+\mathfrak{m}]_Wv$ coincides with $(D+\mathfrak{m})v$ on $W$ and vanishes off $W$, while $\operatorname{Pen}\,v$ vanishes off $W$ by~\eqref{eq:wilson-operators-B3}.

\begin{corollary}\label{cor:wilson-operators-1}
Let $A'$ be an $R$-periodic subset of $W$, let
\[
\boldsymbol{\ell}^{+}(A'):=\sup\left\{\boldsymbol{\ell}_{E'}:E'\in\mathbb{E},\ E'\cap A'\ne\emptyset\right\},
\]
and let $f\in\mathbb{F}(\eta,R,\boldsymbol{\theta})$ be bounded and vanish off $A'$. For $E'\in\mathbb{E}$ let $\mathbf{1}_{E'}$ be the indicator function of the orbit of $E'$, that is, of the union of the translates of $E'$ by the elements of $R\mathbb{Z}^3$ (of $E'$ itself if $R=\infty$), and let $E'$ range over a set of representatives of these orbits (the blocks of $\mathbb{E}$ modulo $R$). Then the series
\[
u:=\sum_{E'}S(\mathbf{1}_{E'}\cdot f)
\]
converges absolutely at every site of $\eta\mathbb{Z}^3$, its sum $u$ lies in $\mathbb{F}(\eta,R,\boldsymbol{\theta})$ and vanishes off $W$, one has $Tu=f$, and $u=Sf$ if $f\in\mathbb{H}$. Moreover, for every block $E\in\mathbb{E}$ and every site $x\in E$,
\begin{equation}\label{eq:wilson-operators-21}
|u(x)|+[u]_{\boldsymbol{\alpha};x}\le\tilde{C}\cdot\min\left\{\boldsymbol{\ell}_E,\boldsymbol{\ell}^{+}(A')\right\}\cdot\exp(-\boldsymbol{\gamma}_4\cdot d(x,A'))\cdot\|f\|_{\infty},
\end{equation}
where $[\,\cdot\,]_{\boldsymbol{\alpha};x}$ is the seminorm of Convention~\ref{it:wilson-operators-7} and $d$ is the distance of Definition~\ref{def:path-distance-units}.
\end{corollary}

\begin{proof}
For a block $E'$ put $f_{E'}:=\mathbf{1}_{E'}\cdot f$ and $u_{E'}:=Sf_{E'}$; all sums over $E'$ below extend over the blocks of $\mathbb{E}$ modulo $R$.

The orbit of $E'$ is invariant under the translations by $R\mathbb{Z}^3$; hence $f_{E'}\in\mathbb{F}(\eta,R,\boldsymbol{\theta})$. Moreover $\|f_{E'}\|_{\infty}\le\|f\|_{\infty}$, and $f_{E'}$ vanishes off the orbit of $E'$. It lies in $\mathbb{H}_W$, because it vanishes off $W$, its support is finite if $R=\infty$, and $\mathbb{H}=\mathbb{F}(\eta,R,\boldsymbol{\theta})$ if $R<\infty$. Hence $u_{E'}\in\mathbb{H}_W$ is defined, and Theorem~\ref{thm:sup-holder-bounds} applies to $f_{E'}$. Since $A'$ is $R$-periodic and $f$ vanishes off $A'$, one has $f_{E'}=0$, and therefore $u_{E'}=0$, unless $E'$ meets $A'$; if $E'$ meets $A'$, then $\boldsymbol{\ell}_{E'}\le\boldsymbol{\ell}^{+}(A')$.

Let $E\in\mathbb{E}$, $x\in E$, and let $E'$ meet $A'$. Inequality~\eqref{eq:wilson-operators-16}, applied to $f_{E'}$ with the rate $\boldsymbol{\gamma}_3=a_0/2$ of Theorem~\ref{thm:sup-holder-bounds}, gives
\begin{align*}
|u_{E'}(x)|+[u_{E'}]_{\boldsymbol{\alpha};x}
&\le C_3\cdot\min\left\{\boldsymbol{\ell}_E,\boldsymbol{\ell}_{E'}\right\}\cdot\exp(-\boldsymbol{\gamma}_3\cdot d(x,E'))\cdot\|f_{E'}\|_{\infty}\\
&\le C_3\cdot\min\left\{\boldsymbol{\ell}_E,\boldsymbol{\ell}^{+}(A')\right\}\cdot\exp(-\boldsymbol{\gamma}_3\cdot d(x,E'))\cdot\|f\|_{\infty}.
\end{align*}
Fix $a'\in E'\cap A'$. For $x'\in E'$ the sites $x'$ and $a'$ lie in the same block $E'$, so that $d(x',a')\le 3L$ by Lemma~\ref{lem:properties-path-distance}(c), and the triangle inequality gives $d(x,A')\le d(x,a')\le d(x,x')+3L$. The infimum over $x'\in E'$ yields $d(x,E')\ge d(x,A')-3L$. Since $\boldsymbol{\gamma}_3=\boldsymbol{\gamma}_4+a_0/4$, it follows that
\begin{align*}
\exp(-\boldsymbol{\gamma}_3\cdot d(x,E'))
&=\exp(-\boldsymbol{\gamma}_4\cdot d(x,E'))\cdot\exp\left(-\frac{a_0}{4}\cdot d(x,E')\right)\\
&\le\exp(3L\cdot\boldsymbol{\gamma}_4)\cdot\exp(-\boldsymbol{\gamma}_4\cdot d(x,A'))\cdot\exp\left(-\frac{a_0}{4}\cdot d(x,E')\right).
\end{align*}
Moreover $a_0\le 1/(3L)$, because~\eqref{eq:wilson-operators-4} defines $a_0$ as the minimum of $1/(3L)$ and a second positive number; therefore
\[
\exp(3L\cdot\boldsymbol{\gamma}_4)=\exp\left(\frac{3L\cdot a_0}{4}\right)\le\exp\left(\frac{1}{4}\right)\le 2.
\]
Consequently
\begin{align*}
|u_{E'}(x)|+[u_{E'}]_{\boldsymbol{\alpha};x}
&\le 2C_3\cdot\min\left\{\boldsymbol{\ell}_E,\boldsymbol{\ell}^{+}(A')\right\}\cdot\exp(-\boldsymbol{\gamma}_4\cdot d(x,A'))\\
&\qquad\cdot\exp\left(-\frac{a_0}{4}\cdot d(x,E')\right)\cdot\|f\|_{\infty}
\end{align*}
for every block $E'$ that meets $A'$, while the left side vanishes for the other blocks.

Lemma~\ref{lem:properties-path-distance}(d) with $\sigma=a_0/4$ is applicable, since its hypothesis $s\ge 12/\sigma=48/a_0$ is the assumption $s\ge s_0$; by~\eqref{eq:wilson-operators-11} it gives
\[
\sum_{E'}\exp\left(-\frac{a_0}{4}\cdot d(x,E')\right)\le C_{\mathrm{sh}}(a_0/4).
\]
Summing the bound for a single term over the blocks $E'$ that meet $A'$, the other terms being zero, and using the last display, one obtains, for every block $E\in\mathbb{E}$ and every $x\in E$,
\[
\sum_{E'}\left(|u_{E'}(x)|+[u_{E'}]_{\boldsymbol{\alpha};x}\right)\le\tilde{C}\cdot\min\left\{\boldsymbol{\ell}_E,\boldsymbol{\ell}^{+}(A')\right\}\cdot\exp(-\boldsymbol{\gamma}_4\cdot d(x,A'))\cdot\|f\|_{\infty}.
\]

Every site of $W$ lies in a block of $\mathbb{E}$, so by the last display the series $\sum_{E'}u_{E'}$ converges absolutely at every site of $W$; off $W$ all its terms vanish. Hence the series converges absolutely at every site of $\eta\mathbb{Z}^3$; its sum $u$ vanishes off $W$, and it belongs to $\mathbb{F}(\eta,R,\boldsymbol{\theta})$, since the relations that define $\mathbb{F}(\eta,R,\boldsymbol{\theta})$ are linear and are preserved under pointwise convergence. Let $E\in\mathbb{E}$ and $x\in E$. Then $|u(x)|\le\sum_{E'}|u_{E'}(x)|$. For $\tilde{x}\in W$ with $0<|x-\tilde{x}|_{\infty}\le\boldsymbol{\ell}(x)$ the series converges absolutely at $x$ and at $\tilde{x}$, whence
\[
\frac{\boldsymbol{\ell}(x)^{\boldsymbol{\alpha}}\cdot|u(x)-u(\tilde{x})|}{|x-\tilde{x}|_{\infty}^{\boldsymbol{\alpha}}}\le\sum_{E'}\frac{\boldsymbol{\ell}(x)^{\boldsymbol{\alpha}}\cdot|u_{E'}(x)-u_{E'}(\tilde{x})|}{|x-\tilde{x}|_{\infty}^{\boldsymbol{\alpha}}}\le\sum_{E'}[u_{E'}]_{\boldsymbol{\alpha};x}.
\]
The supremum over $\tilde{x}$ gives $[u]_{\boldsymbol{\alpha};x}\le\sum_{E'}[u_{E'}]_{\boldsymbol{\alpha};x}$. Thus $|u(x)|+[u]_{\boldsymbol{\alpha};x}$ is at most the sum estimated in the preceding paragraph, and~\eqref{eq:wilson-operators-21} follows.

Let $v\in\mathbb{F}(\eta,R,\boldsymbol{\theta})$ vanish off $W$. By definition $(Tv)(x)=0$ for $x\notin W$, while for $x\in W$, by~\eqref{eq:wilson-operators-7} and~\eqref{eq:wilson-operators-B3}, the value $(Tv)(x)$ is the image, under a linear map that does not depend on $v$, of the values of $v$ at $x$, at the six nearest neighbours of $x$ and at the sites of the block of $\mathbb{E}$ that contains $x$; these sites are finitely many. The function $u$ lies in $\mathbb{F}(\eta,R,\boldsymbol{\theta})$ and vanishes off $W$, the series $\sum_{E'}u_{E'}$ converges to $u$ at every site, and $Tu_{E'}=f_{E'}$, because $u_{E'}=Sf_{E'}\in\mathbb{H}_W$. It follows that
\[
(Tu)(x)=\sum_{E'}(Tu_{E'})(x)=\sum_{E'}f_{E'}(x)=f(x)
\]
for every site $x$; the last equality holds because the orbits of the blocks of $\mathbb{E}$ partition $W$ and $f$ vanishes off $W$.

Finally, let $f\in\mathbb{H}$; then $f\in\mathbb{H}_W$. The functions $f_{E'}$ have pairwise disjoint supports, so they are pairwise orthogonal in $\mathbb{H}$ and $\sum_{E'}\|f_{E'}\|^2=\|f\|^2$; hence $\sum_{E'}f_{E'}$ converges to $f$ in $\mathbb{H}$. The operator $S$ is bounded on $\mathbb{H}_W$ by Theorem~\ref{thm:weighted-bound}, so $\sum_{E'}u_{E'}$ converges to $Sf$ in $\mathbb{H}$. Convergence in $\mathbb{H}$ implies convergence at every site, because $|v(x)|\le\eta^{-\frac{3}{2}}\cdot\|v\|$ for every $v\in\mathbb{H}$ and every site $x$. Therefore $u=Sf$.
\end{proof}

If $R=\infty$, one writes $Sf:=u$ for every bounded $f$ as in Corollary~\ref{cor:wilson-operators-1}, also when $f\notin\mathbb{H}$; by the corollary this is consistent with $S=T^{-1}$ on $\mathbb{H}_W$.

\subsection{Homogeneous blocked Wilson problems}\label{sec:homogeneous-blocked-wilson}

\subsubsection{Lattices and tori with blocks of one level}\label{sec:lattices-tori-blocks}

\begin{proposition}\label{prop:wilson-operators-1}
Let $\eta>0$, let $\gamma$ be Clifford data as in Convention~\ref{it:wilson-operators-3}, let $i\ge 1$ be an integer, and let
\[
\boldsymbol{\ell}_i:=L^i\eta,\qquad \boldsymbol{\theta}\in\mathbb{R}^3,\qquad 0\le\mathfrak{m}\le\boldsymbol{\ell}_i^{-1},\qquad y\in\mathbb{R}\cup\{\infty\},\qquad b'\in[b(1-1/L),b].
\]
For every integer $j\ge0$ let $P_j$ be the operator on the functions $u\colon\eta\mathbb{Z}^3\to\mathbb{C}^4$ that replaces $u$, on each block of level $j$, by its mean over this block (Section~\ref{sec:block-means-poincare}): if the site $x$ lies in the block $F$ of level $j$, then
\[
(P_ju)(x):=(P_Fu)(x)=L^{-3j}\sum_{x'\in F}u(x').
\]
If $j\ge0$ and $R\in L^j\eta\mathbb{N}\cup\{\infty\}$, then $P_j$ maps the space $\mathbb{H}(\eta,R,\boldsymbol{\theta})$ of Section~\ref{sec:shifts-differences-free} into itself, and its restriction to this space is the orthogonal projection of $\mathbb{H}(\eta,R,\boldsymbol{\theta})$ onto the subspace of the functions constant on every block of level $j$. The tuples $(\eta,R,\boldsymbol{\theta};\gamma;W,\mathbb{E},\operatorname{Pen},\mathfrak{m})$ given by the two items below are blocked Wilson problems with separation $s=\infty$ and penalty constants $c_1=b(1-1/L)$, $C_1=L\cdot b$, $C_2=2L^4\cdot b$, the constants of~\eqref{eq:wilson-operators-2}.
\begin{enumerate}
\item\label{it:wilson-operators-1-1} $R\in\boldsymbol{\ell}_i\mathbb{N}\cup\{\infty\}$, $W=\eta\mathbb{Z}^3$, $\mathbb{E}$ is the set of the blocks of level $i$, and $\operatorname{Pen}=b'\cdot\boldsymbol{\ell}_i^{-1}\cdot P_i$. The operator~\eqref{eq:wilson-operators-10} is then $T=D+\mathfrak{m}+b'\cdot\boldsymbol{\ell}_i^{-1}\cdot P_i$.
\item\label{it:wilson-operators-1-2} $R\in L\cdot\boldsymbol{\ell}_i\mathbb{N}\cup\{\infty\}$, $W=\eta\mathbb{Z}^3$, $\mathbb{E}$ is the set of the blocks of level $i+1$, and $\operatorname{Pen}=\sum_{E\in\mathbb{E}}\operatorname{Pen}_E$, where for every $E\in\mathbb{E}$ the operator $\operatorname{Pen}_E$ is the penalty of type~\eqref{eq:wilson-operators-Std2} of Definition~\ref{def:standard-penalties} formed with the numbers $b'$ and $y$, extended by zero off $E$; here $\boldsymbol{\ell}'=\boldsymbol{\ell}_E/L=\boldsymbol{\ell}_i$. The operator~\eqref{eq:wilson-operators-10} is then
\[
T=D+\mathfrak{m}+\boldsymbol{\ell}_i^{-1}\left(\boldsymbol{\alpha}_yP_i+\boldsymbol{\beta}_yP_{i+1}\right),
\]
where $\boldsymbol{\alpha}_y$ and $\boldsymbol{\beta}_y$ are the matrices of Definition~\ref{def:standard-penalties}.
\end{enumerate}
\end{proposition}

\begin{proof}
Let $j\ge0$ and $R\in L^j\eta\mathbb{N}\cup\{\infty\}$, and write $\mathbb{H}:=\mathbb{H}(\eta,R,\boldsymbol{\theta})$. By Convention~\ref{it:wilson-operators-4} the blocks of level $j$ partition $\eta\mathbb{Z}^3$, and their centres form the lattice $L^j\eta\mathbb{Z}^3$. A translation by a vector $a$ of this lattice therefore permutes the blocks of level $j$, so that
\[
(P_ju)(x+a)=\left(P_j\,u(\cdot+a)\right)(x)\qquad\left(u\colon\eta\mathbb{Z}^3\to\mathbb{C}^4,\ x\in\eta\mathbb{Z}^3\right).
\]
If $R<\infty$, the vectors $R\cdot e_{\mu}$, $\mu=1,2,3$, belong to this lattice, and for $u\in\mathbb{F}(\eta,R,\boldsymbol{\theta})$ (Convention~\ref{it:wilson-operators-2}) the identity with $a=R\cdot e_{\mu}$ gives $(P_ju)(x+R\cdot e_{\mu})=\exp(i\boldsymbol{\theta}_{\mu})\cdot(P_ju)(x)$; thus $P_j$ maps $\mathbb{F}(\eta,R,\boldsymbol{\theta})$ into itself. Moreover, for $R<\infty$ the union of the blocks of level $j$ whose centres lie in $L^j\eta\cdot\{0,1,\dots,m-1\}^3$, where $m:=R/(L^j\eta)$, is a period cell (Convention~\ref{it:wilson-operators-5}), because these centres represent each class of $L^j\eta\mathbb{Z}^3$ modulo $R\mathbb{Z}^3$ exactly once; for $R=\infty$ the period cell $\eta\mathbb{Z}^3$ is the union of all blocks of level $j$. Now let $F$ be a block of level $j$ and $u\in\mathbb{F}(\eta,R,\boldsymbol{\theta})$. By the Cauchy--Schwarz inequality, $L^{3j}\cdot\left|(P_Fu)(x)\right|^2\le\sum_{x'\in F}|u(x')|^2$ for $x\in F$, that is, $\|P_ju\|_{2,F}\le\|u\|_{2,F}$; and every function $w$ that is constant on $F$ satisfies
\[
\sum_{x\in F}\overline{w(x)}\cdot\left(u(x)-(P_ju)(x)\right)=0.
\]
Summation over the blocks of level $j$ that make up the period cell, the sums converging absolutely for $u,w\in\mathbb{H}$ by the Cauchy--Schwarz inequality, gives $\|P_ju\|\le\|u\|$, and $\langle w,u-P_ju\rangle=0$ for all $u\in\mathbb{H}$ and all $w\in\mathbb{H}$ that are constant on every block of level $j$. Hence $P_j$ maps $\mathbb{H}$ into itself. As $P_ju$ is constant on every block of level $j$, and $P_jw=w$ for every $w$ of this kind, the restriction of $P_j$ to $\mathbb{H}$ is the orthogonal projection of $\mathbb{H}$ onto the subspace of its elements constant on every block of level $j$.

Since $W=\eta\mathbb{Z}^3$ in both cases, no site lies outside $W$. Hence $\mathbb{H}_W=\mathbb{H}$, $[D+\mathfrak{m}]_W=D+\mathfrak{m}$, and the scale function $\boldsymbol{\ell}$ of Convention~\ref{it:wilson-operators-6} assigns to each site the side of the block of $\mathbb{E}$ containing it; thus $\boldsymbol{\ell}\equiv\boldsymbol{\ell}_i$ in case~\ref{it:wilson-operators-1-1} and $\boldsymbol{\ell}\equiv L\cdot\boldsymbol{\ell}_i$ in case~\ref{it:wilson-operators-1-2}.

Condition~\eqref{eq:wilson-operators-B1} of Definition~\ref{def:blocked-wilson-problems} holds. Indeed, $W=\eta\mathbb{Z}^3$ is nonempty, and by Convention~\ref{it:wilson-operators-4} the blocks of level $i$ (in case~\ref{it:wilson-operators-1-2}: of level $i+1$) are pairwise disjoint and have union $W$; their side is $\boldsymbol{\ell}_i$ (respectively $L\cdot\boldsymbol{\ell}_i$). If $R<\infty$, then $R$ is a positive multiple of this side; as shown above, the translations by $R\cdot e_{\mu}$, $\mu=1,2,3$, then permute the blocks of $\mathbb{E}$, and they leave $W$ invariant. Thus $W$ and $\mathbb{E}$ are $R$-periodic, the elements of $\mathbb{E}$ are blocks of level $i\ge1$ (respectively $i+1\ge2$) of side $\boldsymbol{\ell}_E=\boldsymbol{\ell}_i\le R$ (respectively $\boldsymbol{\ell}_E=L\cdot\boldsymbol{\ell}_i\le R$), and $\sup_{E\in\mathbb{E}}\boldsymbol{\ell}_E<\infty$.

For $s=\infty$, condition~\eqref{eq:wilson-operators-B2} asks for $\boldsymbol{\ell}(x)\le L\cdot\boldsymbol{\ell}(z)$ at all sites $x$, $z$. Because $W$ is the whole lattice and all blocks of $\mathbb{E}$ have one and the same side, $\boldsymbol{\ell}$ is constant, with the value $\boldsymbol{\ell}_i$, respectively $L\cdot\boldsymbol{\ell}_i$; as $L>1$, one has $\boldsymbol{\ell}(x)=\boldsymbol{\ell}(z)\le L\cdot\boldsymbol{\ell}(z)$.

As to condition~\eqref{eq:wilson-operators-B3}, let the mean $P_E$ over a block $E$ be extended by zero off $E$, and let $1_E$ be the indicator of $E$ (Convention~\ref{it:wilson-operators-5}). Let $j\ge0$. At each site of a block $E$ of level $j$ the function $P_ju$ takes as its value the mean of $u$ over $E$; hence $P_E=1_EP_j$ for every block $E$ of level $j$, and since these blocks partition $\eta\mathbb{Z}^3$, the operator $P_j$ is the sum of the operators $P_E$ over the blocks $E$ of level $j$, the sum having exactly one nonzero term at each site. In case~\ref{it:wilson-operators-1-1}, where $\boldsymbol{\ell}_E=\boldsymbol{\ell}_i$ for $E\in\mathbb{E}$, it follows that
\[
\operatorname{Pen}=b'\cdot\boldsymbol{\ell}_i^{-1}\cdot P_i=\sum_{E\in\mathbb{E}}b'\cdot\boldsymbol{\ell}_E^{-1}\cdot P_E,
\]
and for every $E\in\mathbb{E}$ the operator $\operatorname{Pen}_E=b'\cdot\boldsymbol{\ell}_E^{-1}\cdot P_E$ is the penalty of type~\eqref{eq:wilson-operators-Std1} of Definition~\ref{def:standard-penalties} with the number $b'\in[b(1-1/L),b]$. In case~\ref{it:wilson-operators-1-2} every $E\in\mathbb{E}$ has level $i+1\ge 2$, as type~\eqref{eq:wilson-operators-Std2} requires. Such a block is the union of $L^3$ blocks of level $i$, whose side is $\boldsymbol{\ell}_i=\boldsymbol{\ell}_E/L=\boldsymbol{\ell}'$ (Convention~\ref{it:wilson-operators-4}); as all sites of $E$ carry the same weight $\eta^3$ in $L^2(E)$, the projection $P'_E$ of Definition~\ref{def:standard-penalties} replaces a function, on each of these blocks, by its mean over it. Extended by zero off $E$, it therefore equals $1_EP_i$, while $P_E=1_EP_{i+1}$. Summation over $E\in\mathbb{E}$, with $\sum_{E\in\mathbb{E}}1_E=1$ and with the matrices $\boldsymbol{\alpha}_y$, $\boldsymbol{\beta}_y$ commuting with the multiplication by $1_E$, gives
\[
\operatorname{Pen}=\sum_{E\in\mathbb{E}}\boldsymbol{\ell}_i^{-1}\left(\boldsymbol{\alpha}_yP'_E+\boldsymbol{\beta}_yP_E\right)=\boldsymbol{\ell}_i^{-1}\left(\boldsymbol{\alpha}_yP_i+\boldsymbol{\beta}_yP_{i+1}\right).
\]
In both cases each $\operatorname{Pen}_E$ is given by a kernel on $E\times E$ and vanishes off $E$, and Lemma~\ref{lem:bounds-standard-penalties}, which holds for every $y\in\mathbb{R}\cup\{\infty\}$, yields conditions~\eqref{eq:wilson-operators-Pen1}--\eqref{eq:wilson-operators-Pen3} with $c_1=b(1-1/L)$, $C_1=L\cdot b$ and $C_2=2L^4\cdot b$. Finally, the translations by $R\cdot e_{\mu}$ permute the blocks of $\mathbb{E}$, and the penalty of every block is given by one and the same formula; hence $\operatorname{Pen}$ commutes with these translations.

Condition~\eqref{eq:wilson-operators-B4} holds: $\mathfrak{m}\ge0$ and $\sup\boldsymbol{\ell}\le L\cdot\boldsymbol{\ell}_i$, so that $\mathfrak{m}\cdot\sup_W\boldsymbol{\ell}\le\boldsymbol{\ell}_i^{-1}\cdot L\cdot\boldsymbol{\ell}_i=L$.

Lastly, $T=[D+\mathfrak{m}]_W+\operatorname{Pen}=D+\mathfrak{m}+\operatorname{Pen}$ on $\mathbb{H}_W=\mathbb{H}$, and the two expressions for $\operatorname{Pen}$ obtained above give the stated formulas for $T$.
\end{proof}

For the problems~\ref{it:wilson-operators-1-1} and~\ref{it:wilson-operators-1-2} of Proposition~\ref{prop:wilson-operators-1} one has $s=\infty$; hence the function $\boldsymbol{\lambda}$ of Lemma~\ref{lem:lipschitz-minorant-scale} is the constant $\inf\boldsymbol{\ell}$, namely $\boldsymbol{\lambda}=\boldsymbol{\ell}_i$ for~\ref{it:wilson-operators-1-1} and $\boldsymbol{\lambda}=L\cdot\boldsymbol{\ell}_i$ for~\ref{it:wilson-operators-1-2}, and in Definition~\ref{def:path-distance-units} every bond costs $\eta/\boldsymbol{\lambda}$. For $z\in\mathbb{R}^3$ let $|z|_1:=|z_1|+|z_2|+|z_3|$, so that $|z|_1\ge|z|_{\infty}$. Let $x,x''\in\eta\mathbb{Z}^3$. The two sites of a bond differ in exactly one coordinate, and there by $\eta$; hence a nearest neighbour path in $\eta\mathbb{Z}^3$ from $x$ to $x''$ has at least $|x-x''|_1/\eta$ bonds, and a monotone coordinate path has exactly that many. Therefore $d_{\infty}(x,x'')=|x-x''|_1/\boldsymbol{\lambda}$, and if $R<\infty$, then $d(x,x')=\boldsymbol{\lambda}^{-1}\min_{n\in\mathbb{Z}^3}|x-x'-Rn|_1$ for $x'\in\eta\mathbb{Z}^3$. When $R<\infty$, let $|x-x'|_1$ denote this minimum and $|x-x'|_{\infty}$ the periodic sup distance of Definition~\ref{def:path-distance-units}; the former is at least the latter, since $|x-x'-Rn|_1\ge|x-x'-Rn|_{\infty}$ for every $n\in\mathbb{Z}^3$. As moreover $\boldsymbol{\lambda}\le L\cdot\boldsymbol{\ell}_i$, it follows that
\[
d(x,x')=\frac{|x-x'|_1}{\boldsymbol{\lambda}}\ge\frac{|x-x'|_{\infty}}{L\cdot\boldsymbol{\ell}_i}\qquad(x,x'\in\eta\mathbb{Z}^3).
\]

\subsubsection{Inverses on homogeneous lattices and tori}\label{sec:inverses-homogeneous-lattices}

\begin{theorem}[Sup and H\"older bounds]\label{thm:sup-holder-bounds-2}
Put
\[
\boldsymbol{\gamma}_5:=\frac{\boldsymbol{\gamma}_4}{L}=\frac{a_0}{4L},\qquad C_5:=\exp(1)\cdot\tilde{C}\cdot L\cdot\left(2+2L\cdot\exp(1)\right),
\]
where $a_0$ is defined in~\eqref{eq:wilson-operators-4}, $\boldsymbol{\gamma}_4=a_0/4$ in~\eqref{eq:wilson-operators-6}, and $\tilde{C}$ in Subsection~\ref{sec:bounded-right-hand}, all three being formed with the penalty constants of~\eqref{eq:wilson-operators-2}. These two numbers depend on $(L,b,\boldsymbol{\alpha})$ only (Subsection~\ref{sec:constants-order}).
Let $i\ge 1$ be an integer, $\eta=L^{-i}$, $R\in L\mathbb{N}\cup\{\infty\}$, $\boldsymbol{\theta}\in\mathbb{R}^3$, $0\le\mathfrak{m}\le 1$ and $y\in\mathbb{R}\cup\{\infty\}$. Let $D$ be the free Wilson operator~\eqref{eq:wilson-operators-7}, $b_i$ the number of Definition~\ref{def:standard-penalties}, $\boldsymbol{\alpha}_y$ and $\boldsymbol{\beta}_y$ the matrices of Definition~\ref{def:standard-penalties} formed with $b'=b_i$, and $P_i$, $P_{i+1}$ the projections of Proposition~\ref{prop:wilson-operators-1}. The operators
\[
T^o:=D+\mathfrak{m}+b_i\cdot P_i,\qquad T^{*}_y:=D+\mathfrak{m}+\boldsymbol{\alpha}_y\cdot P_i+\boldsymbol{\beta}_y\cdot P_{i+1}
\]
are bijections of $\mathbb{H}(\eta,R,\boldsymbol{\theta})$ onto itself. Let $S$ be $(T^o)^{-1}$ or $(T^{*}_y)^{-1}$. For $x,x'\in\mathbb{Z}^3$ let $\Delta_x$ and $\Delta_{x'}$ be the blocks of level $i$ with centres $x$ and $x'$ (Convention~\ref{it:wilson-operators-4}); their side is $L^i\cdot\eta=1$. Let $f\in\mathbb{F}(\eta,R,\boldsymbol{\theta})$ be bounded and supported in the orbit $\bigcup_{n\in\mathbb{Z}^3}(\Delta_{x'}+Rn)$ of $\Delta_{x'}$ (the set $\Delta_{x'}$ itself if $R=\infty$), and let $u:=Sf$. Then for every $z\in\Delta_x$ and every $\tilde z\in\eta\mathbb{Z}^3$ with $0<|z-\tilde z|_\infty\le 1$,
\[
|u(z)|+\frac{|u(z)-u(\tilde z)|}{|z-\tilde z|_\infty^{\boldsymbol{\alpha}}}\le C_5\cdot\exp\left(-\boldsymbol{\gamma}_5\cdot|x-x'|_\infty\right)\cdot\|f\|_\infty .
\]
Here $|z-\tilde z|_\infty$ is the sup distance in $\eta\mathbb{Z}^3$, while $|x-x'|_\infty$ is, if $R<\infty$, the periodic sup distance $\min_{n\in\mathbb{Z}^3}|x-x'-Rn|_\infty$ of Definition~\ref{def:path-distance-units}.
\end{theorem}

\begin{proof}
Since $\eta=L^{-i}$, the blocks of level $i$ have side $\boldsymbol{\ell}_i=L^i\cdot\eta=1$ and centres in $\mathbb{Z}^3$, and the blocks of level $i+1$ have side $L$. The hypotheses of Proposition~\ref{prop:wilson-operators-1} are satisfied with $\boldsymbol{\ell}_i=1$ and $b'=b_i$: one has $0\le\mathfrak{m}\le 1=1/\boldsymbol{\ell}_i$; $b_i\in[b(1-1/L),b]$ by Definition~\ref{def:standard-penalties}; and $L\mathbb{N}\cup\{\infty\}$ is contained in $\boldsymbol{\ell}_i\mathbb{N}\cup\{\infty\}=\mathbb{N}\cup\{\infty\}$ and equals $L\cdot\boldsymbol{\ell}_i\mathbb{N}\cup\{\infty\}$. By Proposition~\ref{prop:wilson-operators-1}, the tuple of case~\ref{it:wilson-operators-1-1}, with the blocks of level $i$ as partition and the penalty $b_i\cdot\boldsymbol{\ell}_i^{-1}\cdot P_i=b_i\cdot P_i$, and the tuple of case~\ref{it:wilson-operators-1-2}, with the blocks of level $i+1$ as partition and the penalty $\boldsymbol{\ell}_i^{-1}\cdot(\boldsymbol{\alpha}_y\cdot P_i+\boldsymbol{\beta}_y\cdot P_{i+1})=\boldsymbol{\alpha}_y\cdot P_i+\boldsymbol{\beta}_y\cdot P_{i+1}$, are blocked Wilson problems with $W=\eta\mathbb{Z}^3$, separation $s=\infty$ and the penalty constants of~\eqref{eq:wilson-operators-2}. As $W=\eta\mathbb{Z}^3$, one has $[D+\mathfrak{m}]_W=D+\mathfrak{m}$ and $\mathbb{H}_W=\mathbb{H}(\eta,R,\boldsymbol{\theta})$, so that the operator $T$ of~\eqref{eq:wilson-operators-10} is $T^o$ for the first problem and $T^{*}_y$ for the second. Since $s=\infty$, the threshold $s\ge 6L$ of Theorem~\ref{thm:weighted-bound} and the threshold $s\ge s_0$ of Theorem~\ref{thm:sup-holder-bounds} and of Corollary~\ref{cor:wilson-operators-1} are met, and the constants $a_0$, $\boldsymbol{\gamma}_3=a_0/2$, $C_3$, $\boldsymbol{\gamma}_4$ and $\tilde{C}$ of these statements are those formed with the penalty constants of~\eqref{eq:wilson-operators-2}. By Theorem~\ref{thm:weighted-bound}, $T^o$ and $T^{*}_y$ are bijections of $\mathbb{H}(\eta,R,\boldsymbol{\theta})$.

For both problems the scale function $\boldsymbol{\ell}$ is constant, equal to $1$ for the first problem and to $L$ for the second. As $s=\infty$, the function $\boldsymbol{\lambda}$ of Lemma~\ref{lem:lipschitz-minorant-scale} is the constant $\inf\boldsymbol{\ell}=\boldsymbol{\ell}$; every bond costs $\eta/\boldsymbol{\lambda}$, and a shortest nearest neighbour path in $\eta\mathbb{Z}^3$ from $z$ to $z'$ has $|z-z'|_1/\eta$ bonds, where $|v|_1:=|v_1|+|v_2|+|v_3|$. By Definition~\ref{def:path-distance-units} it follows that
\[
d(z,z')=\frac{|z-z'|_1}{\boldsymbol{\lambda}}\ge\frac{|z-z'|_\infty}{\boldsymbol{\lambda}}\qquad(z,z'\in\eta\mathbb{Z}^3),
\]
where, if $R<\infty$, $|z-z'|_1$ and $|z-z'|_\infty$ denote here the periodic distances $\min_{n\in\mathbb{Z}^3}|z-z'-Rn|_1$ and $\min_{n\in\mathbb{Z}^3}|z-z'-Rn|_\infty$. Let $z\in\Delta_x$ and let $z'$ be a site of the orbit of $\Delta_{x'}$. A block of level $i$ consists of the sites at sup distance at most $(L^i-1)\eta/2=(1-\eta)/2$ from its centre (Convention~\ref{it:wilson-operators-4}); hence $|z-x|_\infty\le(1-\eta)/2$ and $|z'-x'|_\infty\le(1-\eta)/2$, and the triangle inequality for the distance $|\cdot|_\infty$ just described gives $|z-z'|_\infty\ge|x-x'|_\infty-(1-\eta)\ge|x-x'|_\infty-1$. Consequently
\[
d(z,\Delta_{x'})\ge\frac{|x-x'|_\infty-1}{\boldsymbol{\lambda}},
\]
and the same lower bound holds for the distance from $z$ to the orbit of $\Delta_{x'}$. Moreover $f\in\mathbb{H}(\eta,R,\boldsymbol{\theta})$: for $R<\infty$ this space is $\mathbb{F}(\eta,R,\boldsymbol{\theta})$, and for $R=\infty$ the function $f$ vanishes off the finite set $\Delta_{x'}$. Finally, $a_0\le 1/(3L)$ by~\eqref{eq:wilson-operators-4}, so that $\boldsymbol{\gamma}_5\le\boldsymbol{\gamma}_3\le 1$; and $C_{\mathrm{sh}}(a_0/4)\ge 64$ by~\eqref{eq:wilson-operators-11} (the term with $n=0$), so that $C_3\le\tilde{C}=2C_3\cdot C_{\mathrm{sh}}(a_0/4)$.

Let $S=(T^o)^{-1}$. Then $\boldsymbol{\lambda}=1$, the blocks $\Delta_x$ and $\Delta_{x'}$ belong to the partition, and $\Delta_x$ is the block of $z$. Theorem~\ref{thm:sup-holder-bounds}, applied with $E=\Delta_x$ and $E'=\Delta_{x'}$ (so that $\min(\boldsymbol{\ell}_E,\boldsymbol{\ell}_{E'})=1$), yields by~\eqref{eq:wilson-operators-16}
\[
|u(z)|+[u]_{\boldsymbol{\alpha};z}\le C_3\cdot\exp\left(-\boldsymbol{\gamma}_3\cdot d(z,\Delta_{x'})\right)\cdot\|f\|_\infty\le C_3\cdot\exp(\boldsymbol{\gamma}_3)\cdot\exp\left(-\boldsymbol{\gamma}_3\cdot|x-x'|_\infty\right)\cdot\|f\|_\infty .
\]
Since $\boldsymbol{\ell}(z)=1$ and $W=\eta\mathbb{Z}^3$, the seminorm $[u]_{\boldsymbol{\alpha};z}$ of Convention~\ref{it:wilson-operators-7} is the supremum of $|u(z)-u(\tilde z)|/|z-\tilde z|_\infty^{\boldsymbol{\alpha}}$ over all $\tilde z\in\eta\mathbb{Z}^3$ with $0<|z-\tilde z|_\infty\le 1$, the distance being that of $\eta\mathbb{Z}^3$. Hence the left side of the asserted inequality is bounded by the right side of the last display. As $\boldsymbol{\gamma}_3\le 1$ and $\boldsymbol{\gamma}_5\le\boldsymbol{\gamma}_3$, the latter is at most $\exp(1)\cdot C_3\cdot\exp(-\boldsymbol{\gamma}_5\cdot|x-x'|_\infty)\cdot\|f\|_\infty$, and $\exp(1)\cdot C_3\le\exp(1)\cdot\tilde{C}\le C_5$.

Let $S=(T^{*}_y)^{-1}$. Then $\boldsymbol{\lambda}=L$. Let $A'$ be the orbit of $\Delta_{x'}$; it is an $R$-periodic subset of $W$ which contains the support of $f$. All blocks of the partition have side $L$; hence $\boldsymbol{\ell}^{+}(A')=L$ and, for the block $E$ of $z$, $\min(\boldsymbol{\ell}_E,\boldsymbol{\ell}^{+}(A'))=L$. As $f\in\mathbb{H}(\eta,R,\boldsymbol{\theta})$, the function $u=Sf$ is the function $u$ of Corollary~\ref{cor:wilson-operators-1}, and~\eqref{eq:wilson-operators-21}, together with $d(z,A')\ge(|x-x'|_\infty-1)/L$ and $\boldsymbol{\gamma}_4/L=\boldsymbol{\gamma}_5$, gives
\[
|u(z)|+[u]_{\boldsymbol{\alpha};z}\le\tilde{C}\cdot L\cdot\exp\left(-\boldsymbol{\gamma}_4\cdot\frac{|x-x'|_\infty-1}{L}\right)\cdot\|f\|_\infty=\tilde{C}\cdot L\cdot\exp(\boldsymbol{\gamma}_5)\cdot\exp\left(-\boldsymbol{\gamma}_5\cdot|x-x'|_\infty\right)\cdot\|f\|_\infty .
\]
Since $\boldsymbol{\ell}(z)=L$, the seminorm $[u]_{\boldsymbol{\alpha};z}$ is the supremum of $L^{\boldsymbol{\alpha}}\cdot|u(z)-u(\tilde z)|/|z-\tilde z|_\infty^{\boldsymbol{\alpha}}$ over the sites $\tilde z\in\eta\mathbb{Z}^3$ with $0<|z-\tilde z|_\infty\le L$, the distance being again that of $\eta\mathbb{Z}^3$. As $L^{\boldsymbol{\alpha}}\ge 1$ and $1<L$, it bounds $|u(z)-u(\tilde z)|/|z-\tilde z|_\infty^{\boldsymbol{\alpha}}$ for every $\tilde z\in\eta\mathbb{Z}^3$ with $0<|z-\tilde z|_\infty\le 1$. Finally $\exp(\boldsymbol{\gamma}_5)\le\exp(1)$ and $\exp(1)\cdot\tilde{C}\cdot L\le C_5$.
\end{proof}

\subsection{Coercivity on blocks of level one}\label{sec:coercivity-blocks-level}

\begin{proposition}\label{prop:wilson-operators-2}
Let $\eta > 0$ and $\boldsymbol{\ell}_1 := L\eta$. Let $R \in \boldsymbol{\ell}_1\mathbb{N} \cup \{\infty\}$ and $\boldsymbol{\theta} \in \mathbb{R}^3$, and let $W$ be a nonempty $R$-periodic union of blocks of level one. For every block $E$ of level one contained in $W$ let a number $\boldsymbol{\beta}_E \in [b(1 - 1/L), b]$ be given, such that $\boldsymbol{\beta}_{E + R\,e_{\mu}} = \boldsymbol{\beta}_E$ for $\mu = 1, 2, 3$ if $R < \infty$. Let $\mathfrak{m} \ge 0$ and $\upsilon \in \mathbb{R}$, and let $T^{[\upsilon]} \colon \mathbb{H}_W \to \mathbb{H}_W$ be the operator
\[
T^{[\upsilon]} := [D + \mathfrak{m}]_W + \sum_E \boldsymbol{\beta}_E \cdot \boldsymbol{\ell}_1^{-1} P_E + i\,\hat{\Gamma} \cdot \upsilon,
\]
where $E$ ranges over all blocks of level one contained in $W$. Put
\[
c_0 := \min\left\{ b\left(1 - \frac{1}{L}\right), \frac{4}{L} \right\}.
\]
Then one has, for every $u \in \mathbb{H}_W$,
\[
\operatorname{Re}\left\langle u, T^{[\upsilon]} u \right\rangle \ge c_0 \cdot \boldsymbol{\ell}_1^{-1} \|u\|^2.
\]
The operator $T^{[\upsilon]}$ is a bijection of $\mathbb{H}_W$, and
\[
\left\| \left( T^{[\upsilon]} \right)^{-1} \right\| \le \frac{\boldsymbol{\ell}_1}{c_0}.
\]
These bounds are uniform in $\upsilon$, $\mathfrak{m}$, $W$, $R$, $\boldsymbol{\theta}$ and $\eta$: the constant $c_0$ depends on $b$ and $L$ only.
\end{proposition}

\begin{proof}
If $R < \infty$, then $R$ is a multiple of $\boldsymbol{\ell}_1$ and $W$ is $R$-periodic, so the translations by $R\,e_{\mu}$ ($\mu = 1, 2, 3$) permute the blocks of level one contained in $W$; as $\boldsymbol{\beta}_{E + R\,e_{\mu}} = \boldsymbol{\beta}_E$, the operator $T^{[\upsilon]}$ commutes with these translations and maps $\mathbb{H}_W$ into itself. In the sums over $E$ below, $E$ ranges over a set of blocks of level one contained in $W$ that contains exactly one block of every class $\left\{ E' + Rn : n \in \mathbb{Z}^3 \right\}$, $E'$ a block of level one contained in $W$ (all blocks of level one contained in $W$ if $R = \infty$). By this choice, and since $\boldsymbol{\ell}_1 \le R$, the sites of these blocks are pairwise incongruent modulo $R\mathbb{Z}^3$, and every site of $W$ is congruent to one of them; in particular, no two bonds inside these blocks are congruent modulo $R\mathbb{Z}^3$. Hence, for every $u \in \mathbb{H}_W$, which vanishes off $W$,
\[
\|u\|^2 = \sum_E \|u\|_{2,E}^2, \qquad \left\| \nabla^{+} u \right\|^2 \ge \sum_E \|\nabla u\|_E^2.
\]

Let $u \in \mathbb{H}_W$. The matrix $\hat{\Gamma}$ is Hermitian (Convention~\ref{it:wilson-operators-3}) and $\upsilon$ is real, so $i\,\hat{\Gamma} \cdot \upsilon$ is anti-Hermitian and $\operatorname{Re}\left\langle u, i\,\hat{\Gamma} \cdot \upsilon\, u \right\rangle = 0$. As $u$ vanishes off $W$, one has $\left\langle u, [D + \mathfrak{m}]_W\, u \right\rangle = \left\langle u, (D + \mathfrak{m})\, u \right\rangle$, and by Lemma~\ref{lem:identities-free-wilson}(a) the real part of this number is $(\eta/2) \left\| \nabla^{+} u \right\|^2 + \mathfrak{m}\, \|u\|^2$. On each block $E$ the operator $P_E$ is the orthogonal projection of $L^2(E)$ onto the constant functions, so the penalty term of $T^{[\upsilon]}$ contributes $\sum_E \boldsymbol{\beta}_E \cdot \boldsymbol{\ell}_1^{-1} \|P_E u\|_{2,E}^2$ to $\left\langle u, T^{[\upsilon]} u \right\rangle$. Since $\mathfrak{m} \ge 0$, it follows that
\[
\operatorname{Re}\left\langle u, T^{[\upsilon]} u \right\rangle \ge \frac{\eta}{2} \left\| \nabla^{+} u \right\|^2 + \sum_E \boldsymbol{\beta}_E \cdot \boldsymbol{\ell}_1^{-1} \|P_E u\|_{2,E}^2.
\]
Lemma~\ref{lem:poincare-inequality-blocks}, applied to the blocks $E$, whose side is $\boldsymbol{\ell}_1$, gives $\|u - P_E u\|_{2,E}^2 \le \left( \boldsymbol{\ell}_1^2/8 \right) \|\nabla u\|_E^2$. Together with the lower bound for $\left\| \nabla^{+} u \right\|^2$ above and with $\boldsymbol{\beta}_E \ge b(1 - 1/L)$, this yields
\[
\operatorname{Re}\left\langle u, T^{[\upsilon]} u \right\rangle \ge \sum_E \left[ \frac{\eta}{2} \cdot \frac{8}{\boldsymbol{\ell}_1^2}\, \|u - P_E u\|_{2,E}^2 + b\left(1 - \frac{1}{L}\right) \cdot \boldsymbol{\ell}_1^{-1} \|P_E u\|_{2,E}^2 \right],
\]
and, because $\boldsymbol{\ell}_1 = L\eta$,
\[
\frac{\eta}{2} \cdot \frac{8}{\boldsymbol{\ell}_1^2} = \frac{4}{L \cdot \boldsymbol{\ell}_1}.
\]
Both coefficients in the brackets are therefore at least $c_0 \cdot \boldsymbol{\ell}_1^{-1}$. The functions $P_E u$ and $u - P_E u$ are orthogonal in $L^2(E)$, so that $\|u - P_E u\|_{2,E}^2 + \|P_E u\|_{2,E}^2 = \|u\|_{2,E}^2$ by the theorem of Pythagoras, and
\[
\operatorname{Re}\left\langle u, T^{[\upsilon]} u \right\rangle \ge c_0 \cdot \boldsymbol{\ell}_1^{-1} \sum_E \|u\|_{2,E}^2 = c_0 \cdot \boldsymbol{\ell}_1^{-1} \|u\|^2.
\]
This is the lower bound for the real part asserted in the proposition.

By the Cauchy--Schwarz inequality, $c_0 \cdot \boldsymbol{\ell}_1^{-1} \|u\|^2 \le \|u\| \left\| T^{[\upsilon]} u \right\|$, that is,
\[
\left\| T^{[\upsilon]} u \right\| \ge c_0 \cdot \boldsymbol{\ell}_1^{-1} \|u\| \qquad \left( u \in \mathbb{H}_W \right).
\]
For the adjoint one has $\operatorname{Re}\left\langle u, \left( T^{[\upsilon]} \right)^{*} u \right\rangle = \operatorname{Re}\left\langle T^{[\upsilon]} u, u \right\rangle = \operatorname{Re}\left\langle u, T^{[\upsilon]} u \right\rangle$, so $\left( T^{[\upsilon]} \right)^{*}$ satisfies the same two lower bounds. The operator $T^{[\upsilon]}$ is bounded, since $D$ is bounded, the blocks of level one are disjoint, $\boldsymbol{\beta}_E \le b$, and $i\,\hat{\Gamma} \cdot \upsilon$ is a constant matrix. The argument of Step~5 of the proof of Theorem~\ref{thm:weighted-bound} therefore applies: $T^{[\upsilon]}$ is injective with closed range, and its range is dense because $\left( T^{[\upsilon]} \right)^{*}$ is injective. Hence $T^{[\upsilon]}$ is a bijection of $\mathbb{H}_W$, and the lower bound for $\left\| T^{[\upsilon]} u \right\|$ gives $\left\| \left( T^{[\upsilon]} \right)^{-1} \right\| \le \boldsymbol{\ell}_1/c_0$.
\end{proof}

\section{Polymer activities: algebra and estimates}\label{sec:polymer-activities-algebra}

This section contains the proofs of the propositions of Section~\ref{sec:norms-polymers-operations}; its notation is that of Section~\ref{sec:renormalisation-group-scheme}, and its constants are those chosen in the stages of Definition~\ref{def:order-choice}. Theorems~\ref{thm:unit-lattice-bounds} and~\ref{thm:free-determinant-across} draw moreover on results of Section~\ref{sec:operators-fermion-field} on the operators of the fermion field; Theorem~\ref{thm:free-determinant-across} draws also on Propositions~\ref{prop:depth-stop-level} and~\ref{prop:consequences}.

\paragraph{Classical facts.}
Five classical results are recalled.

Let $U\subset\mathbb{C}$ be an open set, $V$ a finite-dimensional complex vector space, and $\gamma$ a cycle in $U$ whose winding number $n(\gamma,a)$ about every point $a\in\mathbb{C}\setminus U$ vanishes. Let $f\colon U\to V$ be holomorphic, and let $a\in U$ be a point that does not lie on $\gamma$. Let $a_1,\dots,a_m$ be pairwise distinct points of $U$, none of which lies on $\gamma$, and let $g\colon U\setminus\{a_1,\dots,a_m\}\to V$ be holomorphic, with the residue $r_j\in V$ at $a_j$, that is, $r_j$ is the coefficient of $(\zeta-a_j)^{-1}$ in the Laurent expansion of $g$ about $a_j$. Cauchy's integral formula for each coordinate of $f$ and the residue theorem for each coordinate of $g$, the coordinates being taken with respect to a basis of $V$, give
\begin{gather}
n(\gamma,a)\cdot f(a)=\frac{1}{2\pi i}\oint_{\gamma}\mathrm{d}\zeta\,\frac{f(\zeta)}{\zeta-a},\notag\\
\frac{1}{2\pi i}\oint_{\gamma}\mathrm{d}\zeta\,g(\zeta)=\sum_{j=1}^mn(\gamma,a_j)\cdot r_j.\tag{cl1}\label{eq:polymer-activities-cl1}
\end{gather}

Let $U\subset\mathbb{C}^n$ be a connected open set and $f\colon U\to\mathbb{C}$ a holomorphic function. Suppose that $f$ vanishes on a non-empty open subset of $U$, or that $f$ and all its partial derivatives of all orders vanish at one point of $U$. The identity theorem for holomorphic functions of several complex variables states that then
\begin{equation}\tag{cl2}\label{eq:polymer-activities-cl2}
f(z)=0\qquad\text{for every }z\in U.
\end{equation}

Let $n,m\ge 1$, and let $A_1$, $A_2$, $A_3$, $A_4$ be complex matrices of the sizes $n\times n$, $n\times m$, $m\times n$, $m\times m$, respectively. The determinant of the block matrix formed by them is expressed through Schur complements:
\begin{align}
\det\begin{pmatrix}A_1&A_2\\A_3&A_4\end{pmatrix}&=\det A_1\cdot\det\left(A_4-A_3\cdot A_1^{-1}\cdot A_2\right)&&\text{if }A_1\text{ is invertible},\notag\\
\det\begin{pmatrix}A_1&A_2\\A_3&A_4\end{pmatrix}&=\det A_4\cdot\det\left(A_1-A_2\cdot A_4^{-1}\cdot A_3\right)&&\text{if }A_4\text{ is invertible}.\tag{cl3}\label{eq:polymer-activities-cl3}
\end{align}
Let $A$ be an $n\times m$ matrix and $B$ an $m\times n$ matrix. For $A_1=1$, $A_2=-A$, $A_3=B$ and $A_4=1$, where $1$ denotes the unit matrix of the appropriate size, the first formula gives $\det\left(1+B\cdot A\right)$ and the second gives $\det\left(1+A\cdot B\right)$; this is Sylvester's identity
\[
\det\left(1+B\cdot A\right)=\det\left(1+A\cdot B\right).
\]

Euler's theorem: let $G$ be a finite connected multigraph in which every vertex has even degree. Then
\begin{equation}\tag{cl4}\label{eq:polymer-activities-cl4}
\text{$G$ has a closed walk that traverses every edge of $G$ exactly once.}
\end{equation}

The Leibniz formula in a Grassmann algebra: let $\boldsymbol{\theta}_1,\dots,\boldsymbol{\theta}_n$ be elements of a Grassmann algebra over $\mathbb{C}$ with $\boldsymbol{\theta}_i\cdot\boldsymbol{\theta}_j=-\boldsymbol{\theta}_j\cdot\boldsymbol{\theta}_i$ for all $i,j\in\{1,\dots,n\}$, and let $A=(A_{ij})$ be a complex $n\times n$ matrix. Then
\begin{equation}\tag{cl5}\label{eq:polymer-activities-cl5}
\prod_{i=1}^n\left(\sum_{j=1}^nA_{ij}\cdot\boldsymbol{\theta}_j\right)=\det A\cdot\boldsymbol{\theta}_1\cdots\boldsymbol{\theta}_n,
\end{equation}
where the factors of the product are ordered by increasing $i$. Indeed, the left side is the sum of the products $A_{1j_1}\cdots A_{nj_n}\cdot\boldsymbol{\theta}_{j_1}\cdots\boldsymbol{\theta}_{j_n}$ over all $j_1,\dots,j_n\in\{1,\dots,n\}$. A product in which some index occurs twice vanishes: at the cost of a sign its factors can be rearranged so that two equal factors become adjacent, and $\boldsymbol{\theta}_j^2=0$ by the relation with $i=j$. The remaining products correspond to the permutations $\pi$ of $\{1,\dots,n\}$, with $j_i=\pi(i)$, and $\boldsymbol{\theta}_{\pi(1)}\cdots\boldsymbol{\theta}_{\pi(n)}=\operatorname{sgn}(\pi)\cdot\boldsymbol{\theta}_1\cdots\boldsymbol{\theta}_n$. The Leibniz formula $\det A=\sum_{\pi}\operatorname{sgn}(\pi)\cdot A_{1\pi(1)}\cdots A_{n\pi(n)}$ for the determinant now gives~\eqref{eq:polymer-activities-cl5}.

\paragraph{Elementary inequalities for $L$.}
Since $L\ge 1$, the first inequality $L^{\frac{1}{8}-\frac{3\bar{\varepsilon}}{2}}\ge 2400$ of stage~\ref{it:order-choice-4} gives $L^{\frac{1}{8}}\ge 2400$, that is,
\[
L\ge 2400^8.
\]
Let $\varepsilon$ be the number chosen in stage~\ref{it:order-choice-8}; thus $\varepsilon\le\bar{\varepsilon}=\frac{1}{200}$ (stage~\ref{it:order-choice-1}). It follows that
\begin{gather*}
L\ge 5,\qquad L^{\frac{1}{4}}\ge 2,\qquad L^{\frac{3}{4}}\ge 110,\qquad L^3\cdot\exp(-L)\le L^{-\frac{1}{4}},\\
L^{\frac{1}{4}-\frac{3\varepsilon}{2}}\ge 13.6,\qquad L^{-\frac{7}{4}}+4L^{-\frac{7}{8}+\frac{\varepsilon}{2}}<1.
\end{gather*}
The first three inequalities are immediate. The fourth is equivalent to $L^{\frac{13}{4}}\le\exp(L)$, which holds because $\exp(L)\ge\frac{L^4}{24}$ and $L^{\frac{3}{4}}\ge 24$. Since $\varepsilon\le\frac{1}{200}$, one has $\frac{1}{4}-\frac{3\varepsilon}{2}\ge\frac{1}{8}$ and $-\frac{7}{8}+\frac{\varepsilon}{2}\le-\frac{1}{2}$; hence $L^{\frac{1}{4}-\frac{3\varepsilon}{2}}\ge L^{\frac{1}{8}}\ge 2400$ and
\[
L^{-\frac{7}{4}}+4L^{-\frac{7}{8}+\frac{\varepsilon}{2}}\le 5L^{-\frac{1}{2}}\le 5\cdot 2400^{-4}<1.
\]

\subsection{Hypotheses, notation and constants}\label{sec:hypotheses-notation-constants}

\subsubsection{Settings and basic notation}\label{sec:settings-basic-notation}

The statements of Sections~\ref{sec:norms-gaussian-integrals}--\ref{sec:free-fermion-determinant} are made in four settings.

The statements of Section~\ref{sec:norms-gaussian-integrals} are abstract. In Section~\ref{sec:polymers} a level $k$ and a side $l$ are given with
\begin{equation}\tag{Sa}\label{eq:polymer-activities-Sa}
k\ge 0,\qquad l\in\mathrm{Cl}.
\end{equation}

In Sections~\ref{sec:lattice-symmetries-local}--\ref{sec:contraction-reblocking} the following part of hypothesis~\ref{def:hypothesis-one-step} is assumed, the system of primitives of that hypothesis being left aside: the constants of stages~\ref{it:order-choice-1}--\ref{it:order-choice-9} are chosen as prescribed in Section~\ref{sec:order-choice} and satisfy conditions~\eqref{eq:renormalisation-group-A0}--\eqref{eq:renormalisation-group-A10}, \eqref{eq:renormalisation-group-19}, \eqref{eq:renormalisation-group-20}--\eqref{eq:renormalisation-group-24} and~\eqref{eq:renormalisation-group-25} of Section~\ref{sec:conditions}, and
\begin{equation}\tag{Sb}\label{eq:polymer-activities-Sb}
k\ge 0,\qquad l\in\mathrm{Cl},\qquad\frac{l}{L}\in\mathrm{Cl},\qquad 0<e_k,\qquad e_{k+1}=L^{\frac{1}{2}}\cdot e_k\le e_0.
\end{equation}
The objects of level $k$ live on $T_k=\mathbb{L}_k/l\mathbb{Z}^3$, those of level $k+1$ on $T_{k+1}=\mathbb{L}_{k+1}/(l/L)\mathbb{Z}^3$. No system of primitives and no mass occurs in these sections. Of these hypotheses, the statements of Section~\ref{sec:gamma-matrices-symmetry} and Proposition~\ref{prop:properties-actions}, its clause~(d) included, require only $l\in\mathrm{Cl}$; consequently they are true also under the hypotheses of the scale-free setting of Section~\ref{sec:properties-scheme}. As to clause~(d), $l\in\mathrm{Cl}$ gives $l/L=n\cdot M$ with an odd integer $n\ge13$, so that $l/L$ is a multiple of $M$ and $l/L>2M$; and the definitions of the objects of level $k+1$ that occur in~(d) use nothing else of the side $l/L$. These objects are the torus $T_{k+1}$; the group $\mathbb{G}_{k+1}$ (Section~\ref{sec:dirac-matrices-lattice}) and the automorphisms $a_{g'}$ (Section~\ref{sec:lattice-gauge-charge}); the set $\mathrm{Pz}_{k+1}$, the shortest lifts of its pairs and the generators $\psi$, $\chi$ of level $k+1$ (Definition~\ref{def:pairs-sites-holder}, Sections~\ref{sec:grassmann-algebras-norms} and~\ref{sec:antiperiodicity}); and $A_L$, $\psi_L$, $\chi_L$ (Section~\ref{sec:local-parts-adjustment}). The condition $l/L\in\mathrm{Cl}$ of~\eqref{eq:polymer-activities-Sb} is needed only for the polymers of $D_{k+1}$, the reblocked family $L\,E$ and the domains of level $k+1$, none of which occurs in Section~\ref{sec:gamma-matrices-symmetry} or in Proposition~\ref{prop:properties-actions}. For the assertions (a)--(d) of Proposition~\ref{prop:symmetries-extraction}, which follow from Section~\ref{sec:gamma-matrices-symmetry} and Proposition~\ref{prop:properties-actions}, this is assertion~(b) of Theorem~\ref{thm:renormalisation-group-1}. Similarly, of the hypotheses of these sections, the vocabulary introduced at the beginning of Section~\ref{sec:locality} (spectators, the algebra $\mathrm{Gr}$, legs, and the length and the kinds of a generator~$\chi$), Definition~\ref{def:local-maps}, the descriptions given in Section~\ref{sec:locality} of the property~\eqref{it:renormalisation-group-1-k2} of Definition~\ref{def:renormalisation-group-1} by local maps, and assertions~(a) and~(b) of Theorem~\ref{thm:closure-properties-local} require only~\eqref{eq:polymer-activities-Sa}, that is, $k\ge 0$ and $l\in\mathrm{Cl}$, together with a charge $e_k>0$ with which the domains $\tilde{R}_k$ and $\mathfrak{X}_k$ are formed (all of them concern the single torus $T_k$ and involve no object of another level, no upper bound on $e_k$ and no condition of Section~\ref{sec:conditions}); since~\eqref{eq:polymer-activities-Sb} contains $l/L\in\mathrm{Cl}$ and $e_{k+1}=L^{\frac{1}{2}}\cdot e_k>0$, they hold also for the data $(k+1,l/L,e_{k+1})$ in place of $(k,l,e_k)$, that is, on the torus $T_{k+1}$, and it is with these data that Section~\ref{sec:locality} takes them at the level $k+1$.

In Sections~\ref{sec:bounds-unit-lattice} and~\ref{sec:block-covariance-real} the hypotheses are those of the scale-free setting of Section~\ref{sec:properties-scheme} with $k\ge 1$, namely
\begin{equation}\tag{Sc}\label{eq:polymer-activities-Sc}
k\ge 1,\qquad l\in\mathrm{Cl},\qquad\mathfrak{m}\in\left(0,\mathfrak{m}_0\right],\qquad e_k\in\left(0,e_0\right].
\end{equation}

In Section~\ref{sec:free-fermion-determinant} the hypotheses are
\begin{gather}
\text{for Lemma~\ref{lem:block-determinants}: any level $k\ge 0$, any side $l\in\{L,2L,3L,\dots\}$, any complex mass;}\notag\\
\text{for Theorem~\ref{thm:free-determinant-across}: hypothesis~\ref{def:hypothesis-flow}.}\tag{Sd}\label{eq:polymer-activities-Sd}
\end{gather}
Lemma~\ref{lem:block-determinants} is algebraic: none of the conditions of Section~\ref{sec:conditions} is assumed in it, and its side $l$ need not belong to $\mathrm{Cl}$.

Throughout, $\eta:=\eta_k=L^{-k}$. The thresholds of the domains $R_k$ and $\mathfrak{X}_k$ of Section~\ref{sec:domains-bond-functions} are denoted by
\begin{align*}
&t_{1,k}:=e_k^{-\frac{3}{4}+\varepsilon}, && t_{2,k}:=e_k^{-\frac{3}{4}+2\varepsilon}, && t_{3,k}:=e_k^{-\frac{3}{4}+3\varepsilon} && (\text{for }R_k),\\
&x_{1,k}:=e_k^{-1+2\varepsilon}, && x_{2,k}:=e_k^{-1+3\varepsilon} && && (\text{for }\mathfrak{X}_k).
\end{align*}

A bond function $b$ on $T_k$ is a complex function on the oriented bonds $(x,x')$ of $T_k$, $x$ and $x'$ being nearest neighbours, with $b(x',x)=-b(x,x')$; its components are $b_{\mu}(x):=b(x,x+\eta\cdot e_{\mu})$. Its three norms of Section~\ref{sec:domains-bond-functions} are written
\[
|b|:=|b|_{\infty},\qquad|\partial b|:=|\partial b|_{\infty},\qquad[\partial b]_{\boldsymbol{\alpha}}:=|\delta_{\alpha}\partial b|_{\infty}.
\]
For a function $\boldsymbol{\lambda}$ on the sites of $T_k$ and nearest neighbours $x$, $x'$ put
\[
(d\boldsymbol{\lambda})(x,x'):=\frac{\boldsymbol{\lambda}(x')-\boldsymbol{\lambda}(x)}{\eta},\qquad\text{so that}\qquad e_k\cdot\eta\cdot(d\boldsymbol{\lambda})(x,x')=e_k\cdot\left(\boldsymbol{\lambda}(x')-\boldsymbol{\lambda}(x)\right).
\]

Let $y\in\mathbb{Z}^3/l\mathbb{Z}^3$ be a unit site and $x\in B_k(y)$. For $0\le j\le k$ let $y_j$ be the centre of the block of side $L^j\cdot\eta$ that contains $x$, the blocks of side $L^j\cdot\eta$ being the cubes of this side centred at the points of $\left(L^j\cdot\eta\right)\mathbb{Z}^3$. Then $y_0=x$, and $y_k=y$ because $L^k\cdot\eta=1$. Put
\[
(\tau_kb)(y,x):=\sum_{j=0}^{k-1}(\tau b)(y_{j+1},y_j),
\]
where $\tau b$ is formed by the formula of Definition~\ref{def:pairs-sites-holder}, with the lift of the pair $(y_{j+1},y_j)$ described next; every summand with $y_{j+1}=y_j$ is zero. Let $\tilde{x}\in\mathbb{L}_k$ represent $x$, and for $0\le j\le k$ let $\tilde{y}_j\in\left(L^j\cdot\eta\right)\mathbb{Z}^3$ be the representative of $y_j$ with $|\tilde{x}-\tilde{y}_j|<\frac{L^j\cdot\eta}{2}$. For $0\le j\le k-1$ the vector $\tilde{y}_j-\tilde{y}_{j+1}$ belongs to $\left(L^j\cdot\eta\right)\mathbb{Z}^3$ and has sup norm less than $\frac{L+1}{2}\cdot L^j\cdot\eta$. Since $\frac{L+1}{2}$ is an integer, $L$ being odd, and since $L^k\cdot\eta=1$ and $l\ge1$, it follows that
\[
\left|\tilde{y}_{j+1}-\tilde{y}_j\right|\le\frac{L-1}{2}\cdot L^j\cdot\eta=\frac{\left(L^{j+1}-L^j\right)\cdot\eta}{2}<\frac{L^{j+1}\cdot\eta}{2}\le\frac{1}{2}\le\frac{l}{2}.
\]
If a vector of $\mathbb{L}_k$ that represents the difference of two sites of $T_k$ has sup norm less than $\frac{l}{2}$, then every other vector that represents this difference differs from it by a non-zero vector of $l\mathbb{Z}^3$ and therefore has sup norm greater than $\frac{l}{2}$. Hence $\left(\tilde{y}_{j+1},\tilde{y}_j\right)$ is a shortest lift of the pair $(y_{j+1},y_j)$, and the shortest lift is unique up to a common translation by a vector of $l\mathbb{Z}^3$, which does not change the bonds of $T_k$ traversed by the paths below. With this lift the six rectilinear paths $\Gamma^{\pi}(y_{j+1},y_j)$ of Section~\ref{sec:averaging-fluctuations-gauge} are formed, and
\[
(\tau b)(y_{j+1},y_j)=\frac{1}{3!}\sum_{\pi}b\left(\Gamma^{\pi}(y_{j+1},y_j)\right)
\]
is the mean over the six orders $\pi$ of the plain sums of $b$ over the bonds of these paths. This is the formula of Definition~\ref{def:pairs-sites-holder}, and it is the axial functional $\tau$ of Section~\ref{sec:averaging-fluctuations-gauge}. Of the side, the construction uses only $l\ge1$; hence $\tau_kb$ is defined for every side admitted in~\eqref{eq:polymer-activities-Sa}--\eqref{eq:polymer-activities-Sd}, and, at the level $k+1$, for the side $l/L\ge1$ of Lemma~\ref{lem:block-determinants}(iv). In the same way, let the side $l$ be a multiple of $L$, let $y\in L\mathbb{Z}^3/l\mathbb{Z}^3$, and let $u$ be a unit site of the $L$-block $B(y)$ (Section~\ref{sec:lattice-geometry}). Then $y$ and $u$ have representatives $\tilde{y}\in L\mathbb{Z}^3$ and $\tilde{u}\in\mathbb{Z}^3$ with $|\tilde{u}-\tilde{y}|<\frac{L}{2}$, hence with $|\tilde{u}-\tilde{y}|\le\frac{L-1}{2}<\frac{l}{2}$, as $l\ge L$. So $(\tilde{y},\tilde{u})$ is a shortest lift of $(y,u)$, unique up to a common translation by a vector of $l\mathbb{Z}^3$, and for a bond function $b$ on $T_k$ the number $(\tau b)(y,u)$ of Section~\ref{sec:fermion-integrals-blocking} is formed with the six rectilinear paths of the unit lattice from $y$ to $u$ in this lift, $b$ being summed over their fine bonds. The function $\tau_kb$ is the one in the exponent of the operator $Q_k(\cdot)$ of Section~\ref{sec:exact-fermion-operators}.

\subsubsection{Balls in the torus and their lifts}\label{sec:balls-torus-lifts}

On $\mathbb{R}^3$, and in particular on $\mathbb{L}_k$, $|\cdot|$ denotes the sup norm and $|\cdot|_1$ the $\ell^1$ norm. The sup distance on $T_k$ is denoted by $d(\cdot,\cdot)$: $d(x,x')$ is the least sup norm of a vector of $\mathbb{L}_k$ that represents $x-x'$.

The boundaries of the $M$-cubes, of the $LM$-cubes, of the $L$-blocks and of the unit cubes $\Delta_y$ contain no site of $\mathbb{L}_k$. Indeed, every boundary point of such a cube has a coordinate of the form $(2c+1)\cdot\frac{s}{2}$ with $c\in\mathbb{Z}$ and $s=M$, $L\cdot M$, $L$ or $1$, respectively, whereas the coordinates of the sites lie in $L^{-k}\mathbb{Z}$; and $(2c+1)\cdot\frac{L^k\cdot s}{2}$ is not an integer, because $L$ and $M$ are odd. For a union $Z$ of $M$-cubes, the letter $Z$ denotes also the set of those sites of $T_k$ that belong to it. A path is said to \emph{lie in} $Z$ when each of its sites belongs to $Z$, and a bond is said to \emph{meet} $Z$ when $Z$ contains at least one end point of the bond.

Let $x_0$ be a site of $T_k$, and let $r$ be real with $0<2r<l$. The ball of radius $r$ about $x_0$ is
\[
K(x_0,r):=\left\{x\in T_k:d(x,x_0)\le r\right\}.
\]
For $z\in\mathbb{L}_k$ with $|z|\le r$ the site $x_0+z$ lies in $K(x_0,r)$, as $z$ represents $(x_0+z)-x_0$; and every $x\in K(x_0,r)$ is of the form $x_0+z$ for exactly one $z\in\mathbb{L}_k$ with $|z|\le r$. Indeed, such a $z$ exists by the definition of $d$; and if $z\ne z'$ were two of them, then $z-z'$ would be a non-zero vector of $l\mathbb{Z}^3$ of sup norm at most $2r<l$, whereas every non-zero vector of $l\mathbb{Z}^3$ has sup norm at least $l$. Thus $z\mapsto x_0+z$ is a bijection of $\{z\in\mathbb{L}_k:|z|\le r\}$ onto $K(x_0,r)$. For $x\in K(x_0,r)$, $x-x_0$ denotes the vector $z$ with $|z|\le r$ and $x_0+z=x$. Assume moreover $2r+\eta<l$, and let $\tilde{x}_0\in\mathbb{L}_k$ represent $x_0$. By the bijection just obtained, the quotient map $\mathbb{L}_k\to T_k$ maps the set $\tilde{x}_0+\{z\in\mathbb{L}_k:|z|\le r\}$ bijectively onto $K(x_0,r)$; the sup-diameter of this set is at most $2r$, and $2r+\eta<l$. So the set is a lift of $K(x_0,r)$ in the sense of Definition~\ref{def:terminology}(9), and $K(x_0,r)$ does not wrap. Moreover $K(x_0,r)$ is connected through nearest neighbours, because the cube $\{z\in\mathbb{L}_k:|z|\le r\}$ is connected through nearest neighbours of $\mathbb{L}_k$ and $z\mapsto x_0+z$ maps nearest neighbours to nearest neighbours. By Definition~\ref{def:terminology}(9) the lifts of $K(x_0,r)$ are therefore exactly the sets $\tilde{x}_0+\{z\in\mathbb{L}_k:|z|\le r\}$, $\tilde{x}_0$ ranging over the representatives of $x_0$.

Let $s\ge 0$ with $2r+s<l$. Then any two points $x,x'\in K(x_0,r)$ with $d(x,x')\le s$ satisfy
\[
\left|(x-x_0)-(x'-x_0)\right|=d(x,x').
\]
Indeed, $(x-x_0)-(x'-x_0)$ represents $x-x'$ and has sup norm at most $2r$, whereas every representative of $x-x'$ other than one of sup norm $d(x,x')\le s$ differs from the latter by a non-zero vector of $l\mathbb{Z}^3$ and therefore has sup norm at least $l-s>2r$. Consequently, when the cube $\{z\in\mathbb{L}_k:|z|\le r\}$ of $\mathbb{L}_k$ is identified with $K(x_0,r)$ by $z\mapsto x_0+z$, the two sets have the same nearest neighbours if $2r+\eta<l$ (take $s=\eta$), and the same pairs of distinct points at sup distance at most $M$, that is, the same pairs of $\mathrm{Pz}_k$, if $2r+M<l$ (take $s=M$); in these cases the rectilinear paths of such pairs, and all difference quotients, can be computed in the cube.

\subsubsection{Antiperiodicity}\label{sec:antiperiodicity}

A spinor function on $T_k$ is a function $f\colon\mathbb{L}_k\to\mathbb{C}^4$ such that
\begin{equation}\tag{AP}\label{eq:polymer-activities-AP}
f(x+l\cdot e_{\mu})=-f(x)\qquad\text{for every }x\in\mathbb{L}_k\text{ and every direction }\mu
\end{equation}
(Section~\ref{sec:exact-fermion-operators}). Accordingly, the generators $\psi(x,\boldsymbol{\beta},\boldsymbol{\omega})$ are indexed by $x\in\mathbb{L}_k$, a spinor index $\boldsymbol{\beta}$ and $\boldsymbol{\omega}\in\{0,1\}$, with $\boldsymbol{\omega}=0$ for the unbarred and $\boldsymbol{\omega}=1$ for the barred generators, and they satisfy $\psi(x+l\cdot e_{\mu},\boldsymbol{\beta},\boldsymbol{\omega})=-\psi(x,\boldsymbol{\beta},\boldsymbol{\omega})$. The generators $\chi(x,x',\boldsymbol{\beta},\boldsymbol{\omega})$ are indexed by the pairs $(x,x')\in\mathbb{L}_k\times\mathbb{L}_k$ with $0<|x-x'|\le M$ and by $\boldsymbol{\beta}$, $\boldsymbol{\omega}$, and they satisfy $\chi(x+l\cdot e_{\mu},x'+l\cdot e_{\mu},\boldsymbol{\beta},\boldsymbol{\omega})=-\chi(x,x',\boldsymbol{\beta},\boldsymbol{\omega})$. The generators $\Psi$ and $W$ are indexed by $\mathbb{Z}^3$ in the same way. Restricting $x$ to a fundamental domain of $l\mathbb{Z}^3$ yields independent generators; they are indexed by $T_k$, $\mathrm{Pz}_k$ and $T^0_k$, respectively, as in Section~\ref{sec:grassmann-algebras-norms}.

A kernel is a function of arguments in $\mathbb{L}_k$, namely sites, or pairs of sites for the arguments of $\chi$, that is antiperiodic in each argument. The product of a kernel with the generators of its arguments is then a function on the torus, and the integrals $\int\mathrm{d}\xi$ and $\int\mathrm{d}\boldsymbol{\zeta}$ of Section~\ref{sec:grassmann-algebras-norms} are sums over an arbitrary fundamental domain. Since a kernel changes at most its sign when one of its arguments is translated by a vector of $l\mathbb{Z}^3$, its modulus is a function on the torus, and the norms of kernels do not depend on the fundamental domain over which they are computed. An even element does not change when every generator is replaced by its negative.

A kernel is integrated without generators in the extraction of Section~\ref{sec:local-parts-adjustment}, for the kernel of bidegree $(2,0)$, and in the adjustment of Definition~\ref{def:transport-mean}. In both cases the polymer $X$ is small. Put $r:=M\cdot(L+2)$. Then $X\subset K(x_0,r)$ for every site $x_0$ of $X$, by Proposition~\ref{prop:polymers}\eqref{it:polymers-t7}; moreover $2r+M<l$, since $l\ge 13L\cdot M$. One fixes a lift of $K(x_0,r)$, and all arguments are represented by their preimages in this lift. If $\tilde{x}_0\in\mathbb{L}_k$ is the preimage of $x_0$, the preimage of a site $x$ of $X$ is $\tilde{x}_0+(x-x_0)$; these points form the lift of $X$ about $\tilde{x}_0$. Two sites $x_0$ and $x_0'$ of $X$ lead to lifts of $X$ that differ at most by a translation by an element of $l\mathbb{Z}^3$. Indeed, let $x$ be a site of $X$. As $X$ lies in both balls $K(x_0,r)$ and $K(x_0',r)$, the vectors
\[
u:=x-x_0,\qquad v:=x_0'-x_0,\qquad u':=x-x_0'
\]
are defined, $u$ and $v$ with respect to the centre $x_0$ and $u'$ with respect to the centre $x_0'$, and $|u|,|v|,|u'|\le r$. The vector $w:=v+u'-u$ of $\mathbb{L}_k$ projects to $(x_0'-x_0)+(x-x_0')-(x-x_0)=0$ in $T_k$, so $w\in l\mathbb{Z}^3$; on the other hand $|w|\le 3r=3M\cdot(L+2)<13L\cdot M\le l$. Since no non-zero vector of $l\mathbb{Z}^3$ has sup norm less than $l$, it follows that $w=0$, that is, $u=v+u'$. Now let $\tilde{x}_0':=\tilde{x}_0+v$, the preimage of $x_0'$ in the lift of $K(x_0,r)$ fixed above. Then $\tilde{x}_0'+u'=\tilde{x}_0+u$ for every site $x$ of $X$; thus the lift of $X$ about $\tilde{x}_0'$, formed with the base site $x_0'$, is the lift of $X$ about $\tilde{x}_0$. Any other representative of $x_0'$ differs from $\tilde{x}_0'$ by an element of $l\mathbb{Z}^3$, and the lift of $X$ about it is the translate of this lift by that element. This lift of $X$ is the lift $\tilde{X}$ of Proposition~\ref{prop:polymers}\eqref{it:polymers-t10}; precisely, it is the set of the sites of $\mathbb{L}_k$ in that translate of $\tilde{X}$ by a vector of $l\mathbb{Z}^3$ which contains $\tilde{x}_0$. Indeed, as $X$ is small, $2M\cdot\left(d_M(X)+2\right)\le 2M\cdot(L+2)<13L\cdot M\le l$, so that~\eqref{it:polymers-t10} applies to $X$. The unions of $M$-cubes of $\mathbb{R}^3$ that it provides are the translates of one of them by the vectors of $l\mathbb{Z}^3$; as the projection maps each of them bijectively onto $X$, exactly one of them contains $\tilde{x}_0$; let $\tilde{X}$ denote this one. Let $x$ be a site of $X$ and $\tilde{x}$ its preimage in $\tilde{X}$. The last assertion of Proposition~\ref{prop:polymers}\eqref{it:polymers-t7}, applied to $\tilde{X}$, and the equality $d_M(\tilde{X})=d_M(X)$ of~\eqref{it:polymers-t10} give
\[
|\tilde{x}-\tilde{x}_0|\le M\cdot\left(d_M(X)+2\right)\le M\cdot(L+2)=r.
\]
So $\tilde{x}-\tilde{x}_0$ is a vector of $\mathbb{L}_k$ of sup norm at most $r$ that represents $x-x_0$; by the uniqueness in Section~\ref{sec:balls-torus-lifts}, which holds as $2r<l$, it is the vector $x-x_0$. Since the projection maps $\tilde{X}$ bijectively onto $X$, the sites of $\mathbb{L}_k$ in $\tilde{X}$ are therefore exactly the points $\tilde{x}_0+(x-x_0)$, $x$ a site of $X$.

If all arguments of a kernel are translated by one and the same vector of $l\mathbb{Z}^3$, the kernel is multiplied by $\sigma^p$, where $p$ is the number of its arguments and the sign $\sigma\in\{1,-1\}$ depends only on the vector, by the antiperiodicity in each argument. For even $p$ one has $\sigma^p=1$; hence, if a kernel has an even number of arguments and these are all taken in one lift, its values remain the same when that lift is replaced by another one.

\subsubsection{Coefficients and kernels in Grassmann algebras}\label{sec:coefficients-kernels-grassmann}

Let $p\ge0$ be an integer, and for $1\le i\le p$ let $S_i$ be a finite set of indices carrying a total order. The complex Grassmann algebra in which the elements $\boldsymbol{\phi}_i(s)$, $1\le i\le p$, $s\in S_i$, are independent generators is denoted by $\mathrm{Gr}=\mathrm{Gr}[\boldsymbol{\phi}_1,\dots,\boldsymbol{\phi}_p]$. For a subset $I\subset S_i$ let $\boldsymbol{\phi}_i(I)$ be the product of the generators $\boldsymbol{\phi}_i(s)$, $s\in I$, taken in increasing order. Every $F\in\mathrm{Gr}$ has a unique expansion
\[
F=\sum_{I_1\subset S_1,\dots,I_p\subset S_p}c_F(I_1,\dots,I_p)\cdot\boldsymbol{\phi}_1(I_1)\cdots\boldsymbol{\phi}_p(I_p)
\]
with complex coefficients $c_F(I_1,\dots,I_p)$, the sum being taken over all $p$-tuples of subsets. For a vector $g=(g_1,\dots,g_p)$ of weights $g_i>0$ put
\begin{equation}\label{eq:polymer-activities-1}
\|F\|_g:=\sum_{I_1,\dots,I_p}\left|c_F(I_1,\dots,I_p)\right|\cdot g_1^{\#|I_1|}\cdots g_p^{\#|I_p|}.
\end{equation}
In particular a single generator $\boldsymbol{\phi}_i(s)$ has the norm $g_i$. For $p=0$ there are no generators; then $\mathrm{Gr}=\mathbb{C}$, and by~\eqref{eq:polymer-activities-1} the norm of a number is its modulus.

Let, moreover, positive weights $\nu_i(s)$, $s\in S_i$, be given, and for an argument $s$ with values in $S_i$ write $\int\mathrm{d}s\,f(s):=\sum_{s\in S_i}\nu_i(s)\cdot f(s)$. Suppose that $F$ is written with kernels $K_n$, $n=(n_1,\dots,n_p)$, in the form
\begin{multline}\label{eq:polymer-activities-2}
F=\sum_{n_1,\dots,n_p\ge0}\frac{1}{n_1!\cdots n_p!}\int\mathrm{d}s\,K_n\left(s^1_1,\dots,s^1_{n_1};\dots;s^p_1,\dots,s^p_{n_p}\right)\\
\cdot\boldsymbol{\phi}_1(s^1_1)\cdots\boldsymbol{\phi}_1(s^1_{n_1})\cdots\boldsymbol{\phi}_p(s^p_1)\cdots\boldsymbol{\phi}_p(s^p_{n_p}),
\end{multline}
every argument $s^i_j$ being integrated over $S_i$ in this way. Then, by Proposition~\ref{prop:weighted-coefficient-norms}\eqref{it:weighted-coefficient-norms-N0},
\[
\|F\|_g\le\sum_n\left(\prod_{i=1}^p\frac{g_i^{n_i}}{n_i!}\right)\cdot\int\mathrm{d}s\,\left|K_n\right|.
\]
Equality holds here if, for every $n$, the kernel $K_n$ is antisymmetric separately in the arguments of every group, that is, changes its sign when two arguments $s^i_j$, $s^i_{j'}$ with the same $i$ and $j\ne j'$ are interchanged. Moreover, two representations~\eqref{eq:polymer-activities-2} of one element by kernels with this property have the same kernels.

Let $p=2$, let $S_1=\{\xi\}$ be the set of the indices of the generators $\psi(\xi)$, with the weight $\nu_1=\eta^3$, and let $S_2=\{\boldsymbol{\zeta}\}$ be the set of the indices of the generators $\chi(\boldsymbol{\zeta})$, with the weight $\nu_2=\eta^6$ (Section~\ref{sec:grassmann-algebras-norms}). Then the norm~\eqref{eq:polymer-activities-1} is the norm $\|E\|_h$ of Section~\ref{sec:grassmann-algebras-norms}. With the counting weight on the indices $u$ of the generators $\Psi(u)$ and a single weight $a$, the norm~\eqref{eq:polymer-activities-1} is the norm $|F|_a$ of Section~\ref{sec:grassmann-algebras-norms}.

The part of an element $F$ of \emph{bidegree} $(n,m)$ in $(\psi,\chi)$, that is, the sum of the terms of the expansion of $F$ that contain exactly $n$ generators $\psi(\xi)$ and $m$ generators $\chi(\boldsymbol{\zeta})$, is denoted by $F^{[n,m]}$. By~\eqref{eq:polymer-activities-1}, $\|F\|_g=\sum_{n,m}\|F^{[n,m]}\|_g$.

For a site $x$ of $T_k$ let $\psi(x)$ be the column of the four generators $\psi_{\boldsymbol{\beta}}(x):=\psi(x,\boldsymbol{\beta},0)$, and $\bar{\psi}(x)$ the column of the four generators $\bar{\psi}_{\boldsymbol{\beta}}(x):=\psi(x,\boldsymbol{\beta},1)$; the columns $\chi(x,x')$ and $\bar{\chi}(x,x')$ are formed in the same way from the generators $\chi(x,x',\boldsymbol{\beta},\boldsymbol{\omega})$. Each family $\mathscr{E}$ of Definition~\ref{def:renormalisation-group-1} is viewed as consisting of the elements $\mathscr{E}(X;\mathscr{A},\mathscr{C})\in\mathrm{Gr}[\psi,\chi]$; accordingly, its kernels $\mathscr{E}_{nm}$ are the kernels, antisymmetric in the sense explained above, by which these elements are represented in the form~\eqref{eq:polymer-activities-2}. The kernel $\mathscr{E}_{20}$ is arranged in $4\times4$ matrices: for sites $x$, $y$ of $T_k$, the matrix $\mathscr{E}_{20}(X;\mathscr{A},\mathscr{C};x,y)$ has in row $\boldsymbol{\alpha}$ and column $\boldsymbol{\beta}$ the entry
\[
\mathscr{E}_{20}(X;\mathscr{A},\mathscr{C};x,y)_{\boldsymbol{\alpha}\boldsymbol{\beta}}:=\mathscr{E}_{20}\left(X;\mathscr{A},\mathscr{C};(x,\boldsymbol{\alpha},1),(y,\boldsymbol{\beta},0)\right).
\]
Since the kernels are antisymmetric, and since every term that occurs has as many barred legs as unbarred ones, one has, for every pair of weights $h=(h_1,h_2)$ and with the arguments $\mathscr{A}$, $\mathscr{C}$ suppressed,
\begin{multline}\label{eq:polymer-activities-3}
\mathscr{E}^{[2,0]}(X)=\int\mathrm{d}x\,\mathrm{d}y\,\bar{\psi}(x)^T\cdot\mathscr{E}_{20}(X;x,y)\cdot\psi(y),\\
\frac{h_1^2}{2}\cdot\left\|\mathscr{E}_{20}(X)\right\|=\left\|\mathscr{E}^{[2,0]}(X)\right\|_h=h_1^2\cdot\sum_{\boldsymbol{\alpha},\boldsymbol{\beta}}\int\mathrm{d}x\,\mathrm{d}y\,\left|\mathscr{E}_{20}(X;x,y)_{\boldsymbol{\alpha}\boldsymbol{\beta}}\right|;
\end{multline}
here $\int\mathrm{d}x:=\sum_x\eta^3$, the sum over the sites of $T_k$, and $\|\mathscr{E}_{20}(X)\|$ is the kernel norm of Section~\ref{sec:grassmann-algebras-norms}.

Finally, let $O$ be an operator from the functions on an index set $S$ with weights $\nu(s)$ to the functions on an index set $T$, written as $(Of)(t)=\sum_sO(t,s)\cdot f(s)\cdot\nu(s)$. Put
\[
|O|_{1,\infty}:=\sup_t\sum_s\left|O(t,s)\right|\cdot\nu(s);
\]
this number is equal to $\sup\left\{|Of|_{\infty}:|f|_{\infty}\le1\right\}$, where $f$ ranges over the complex functions on $S$ and $|\cdot|_{\infty}$ is the supremum of the modulus over $S$, respectively over $T$. Indeed, $\left|(Of)(t)\right|\le\sum_s\left|O(t,s)\right|\cdot\nu(s)\cdot|f|_{\infty}$ for every $t$, so that the supremum is at most $|O|_{1,\infty}$. Conversely, for a given $t$ the function
\[
f(s):=\frac{\overline{O(t,s)}}{\left|O(t,s)\right|}\quad\text{if }O(t,s)\ne0,\qquad f(s):=0\quad\text{if }O(t,s)=0,
\]
satisfies $|f|_{\infty}\le1$ and $(Of)(t)=\sum_s\left|O(t,s)\right|\cdot\nu(s)$; hence the supremum is at least $\sum_s\left|O(t,s)\right|\cdot\nu(s)$ for every $t$, and therefore at least $|O|_{1,\infty}$.

\subsubsection{Lattice, gauge, charge and reality transformations}\label{sec:lattice-gauge-charge}

Let $g=(S,a)\in\mathbb{G}_k$ (Section~\ref{sec:dirac-matrices-lattice} and Section~\ref{sec:gamma-matrices-symmetry}), let $\mathsf{C}$ and $\varkappa$ be the matrices of Section~\ref{sec:dirac-matrices-lattice}, and let $\Lambda$ be a real function on the sites of $T_k$. Spinors are written as columns, as in Section~\ref{sec:coefficients-kernels-grassmann}. The four maps below are defined on $\mathrm{Gr}[\psi,\chi]$; on the Grassmann algebras generated by $\Psi$ and by $W$ they are defined by the same formulas, with the sites restricted to $T^0_k$.
\begin{itemize}
\item $a_g$ is the automorphism with
\begin{align*}
a_g\psi(x)&=S\cdot\psi(g^{-1}x),&a_g\bar{\psi}(x)&=(S^{-1})^T\cdot\bar{\psi}(g^{-1}x),\\
a_g\chi(x,x')&=S\cdot\chi(g^{-1}x,g^{-1}x'),&a_g\bar{\chi}(x,x')&=(S^{-1})^T\cdot\bar{\chi}(g^{-1}x,g^{-1}x').
\end{align*}
It is compatible with the antiperiodicity~\eqref{eq:polymer-activities-AP}, because $g^{-1}(x+l\cdot n)=g^{-1}x+l\cdot r(S)\cdot n$ for $n\in\mathbb{Z}^3$.
\item $G_{\Lambda}$ is the automorphism with
\begin{align*}
\psi(x)&\mapsto\exp\left(i\,e_k\cdot\Lambda(x)\right)\cdot\psi(x),&\bar{\psi}(x)&\mapsto\exp\left(-i\,e_k\cdot\Lambda(x)\right)\cdot\bar{\psi}(x),\\
\chi(x,x')&\mapsto\exp\left(i\,e_k\cdot\Lambda(x)\right)\cdot\chi(x,x'),&\bar{\chi}(x,x')&\mapsto\exp\left(-i\,e_k\cdot\Lambda(x)\right)\cdot\bar{\chi}(x,x').
\end{align*}
\item $\mathbf{C}$ is the automorphism with
\begin{align*}
\psi(x)&\mapsto\mathsf{C}\cdot\bar{\psi}(x),&\bar{\psi}(x)&\mapsto-\mathsf{C}\cdot\psi(x),\\
\chi(x,x')&\mapsto\mathsf{C}\cdot\bar{\chi}(x,x'),&\bar{\chi}(x,x')&\mapsto-\mathsf{C}\cdot\chi(x,x').
\end{align*}
\item $\mathbf{K}$ is the antilinear automorphism of rings with
\begin{align*}
\psi(x)&\mapsto\varkappa\cdot\psi(x),&\bar{\psi}(x)&\mapsto\varkappa\cdot\bar{\psi}(x),\\
\chi(x,x')&\mapsto\varkappa\cdot\chi(x,x'),&\bar{\chi}(x,x')&\mapsto\varkappa\cdot\bar{\chi}(x,x').
\end{align*}
Antilinearity means $\mathbf{K}(c\cdot F)=\overline{c}\cdot\mathbf{K}(F)$ for $c\in\mathbb{C}$; thus, in the expansion of an element, every coefficient is conjugated, while the generators are transformed as displayed. Note that $(\varkappa^{-1})^T=\varkappa$ by~\eqref{eq:polymer-activities-13}.
\end{itemize}
With these maps the clauses~\eqref{it:renormalisation-group-1-k3}, \eqref{it:renormalisation-group-1-k4}, \eqref{it:renormalisation-group-1-k5} and~\eqref{it:renormalisation-group-1-k7} of Definition~\ref{def:renormalisation-group-1} take the following form, for every $X\in D_k$ and all $(\mathscr{A},\mathscr{C})\in\tilde{R}_k\times\mathfrak{X}_k$ (Definition~\ref{def:domains}):
\begin{align*}
\eqref{it:renormalisation-group-1-k3}&\qquad a_g\left[\mathscr{E}(gX;g\mathscr{A},g\mathscr{C})\right]=\mathscr{E}(X;\mathscr{A},\mathscr{C})\qquad(g\in\mathbb{G}_k);\\
\eqref{it:renormalisation-group-1-k4}&\qquad G_{\boldsymbol{\lambda}+\boldsymbol{\lambda}'}\left[\mathscr{E}(X;\mathscr{A}^1,\mathscr{C}^1)\right]=\mathscr{E}(X;\mathscr{A},\mathscr{C});\\
\eqref{it:renormalisation-group-1-k5}&\qquad\mathbf{C}\left[\mathscr{E}(X;-\mathscr{A},-\mathscr{C})\right]=\mathscr{E}(X;\mathscr{A},\mathscr{C});\\
\eqref{it:renormalisation-group-1-k7}&\qquad\mathbf{K}\left[\mathscr{E}\left(X;-\overline{\mathscr{A}},-\overline{\mathscr{C}}\right)\right]=\mathscr{E}(X;\mathscr{A},\mathscr{C}).
\end{align*}
In~\eqref{it:renormalisation-group-1-k4} the hypotheses of Definition~\ref{def:renormalisation-group-1} on $(\mathscr{A}^1,\mathscr{C}^1,\boldsymbol{\lambda},\boldsymbol{\lambda}')$ are made: $(\mathscr{A}^1,\mathscr{C}^1)$ lies in $\tilde{R}_k\times\mathfrak{X}_k$ as well, $\boldsymbol{\lambda}$ and $\boldsymbol{\lambda}'$ are real functions on the sites of $T_k$, and $\mathscr{A}^1=\mathscr{A}-d\boldsymbol{\lambda}$, $\mathscr{C}^1=\mathscr{C}-d\boldsymbol{\lambda}'$ on the bonds at distance at most $1$ from $X$.

The matrices $S$, $(S^{-1})^T$, $\mathsf{C}$ and $\varkappa$ are monomial (Section~\ref{sec:gamma-matrices-symmetry}): every row and every column contains exactly one non-zero entry, and this entry has modulus one. Consequently, each of $a_g$, $G_{\Lambda}$, $\mathbf{C}$, $\mathbf{K}$ maps any generator onto a single generator multiplied by a phase factor, that is, by a complex number of absolute value one; and if the phase factors are disregarded, the generators are permuted. By~\eqref{eq:polymer-activities-1} it follows that, for every element $F$ and in each of the norms~\eqref{eq:polymer-activities-1} with one weight for the generators $\psi$ and one weight for the generators $\chi$,
\begin{multline}\label{eq:polymer-activities-4}
\|a_gF\|=\|G_{\Lambda}F\|=\|\mathbf{C} F\|=\|\mathbf{K} F\|=\|F\|,\\
(TF)^{[n,m]}=T\left(F^{[n,m]}\right)\qquad\left(T=a_g,G_{\Lambda},\mathbf{C},\mathbf{K};\ n,m\ge0\right);
\end{multline}
that is, the four maps preserve these norms and the bidegree.

Throughout, $g_1,g_2,\dots$ are weights. Without an index, $g$ stands for a vector of weights in Sections~\ref{sec:coefficients-kernels-grassmann} and~\ref{sec:norms-gaussian-integrals}; for a holomorphic function with values in a finite-dimensional complex vector space in the classical facts that precede Section~\ref{sec:hypotheses-notation-constants}; for a spinor function on $T_k$, or a vector of a finite-dimensional complex Hilbert space, introduced where it occurs, in Sections~\ref{sec:block-covariance-real} and~\ref{sec:determinant-identities}; and for an element of $\mathbb{G}_k$ in all other places.

\subsubsection{Absolute constants}\label{sec:absolute-constants}

The absolute constants $c_X$, $\kappa_0$, $K_0$, $c_{\mathrm{rb}}$ and $c_{\mathrm{fl}}$ of stage~\ref{it:order-choice-1} are
\begin{multline}\label{eq:polymer-activities-5}
c_X:=54,\qquad\kappa_0:=54\log72\quad(=230.9\ldots),\qquad K_0:=\frac{\exp(\kappa_0)}{36},\\
c_{\mathrm{rb}}:=c_X\cdot K_0,\qquad c_{\mathrm{fl}}:=61\,c_X\cdot K_0.
\end{multline}
These are explicit numbers; $c_X\ge1$, $\kappa_0\ge c_X$ and $K_0\ge1$. Further, $c':=3c_X$, and for $\gamma>0$
\[
c(\gamma):=\sum_{n\in\mathbb{Z}^3}\exp(-\gamma\cdot|n|)=1+\sum_{j\ge1}\left(24j^2+2\right)\cdot\exp(-\gamma\cdot j),
\]
$|n|$ being the sup norm; there are $24j^2+2$ points $n\in\mathbb{Z}^3$ with $|n|=j$. On every torus $\mathbb{Z}^3/l\mathbb{Z}^3$, with the sup distance $d$,
\[
\sup_y\sum_{y'}\exp\left(-\gamma\cdot d(y,y')\right)\le c(\gamma),
\]
because to every site $y'$ there is an $n\in\mathbb{Z}^3$ with $y'=y+n$ on the torus and $|n|=d(y,y')$, and distinct sites $y'$ give distinct $n$. Moreover
\[
c(\gamma)\le\left(1+\frac{6}{\gamma}\right)^3,\qquad\text{hence}\qquad c(\gamma)\le\left(\frac{12}{\gamma}\right)^3\quad\text{if }\gamma\le6.
\]
Indeed, $|n|\ge\frac{1}{3}|n|_1$ gives
\[
c(\gamma)\le\left(\sum_{j\in\mathbb{Z}}\exp\left(-\frac{\gamma\cdot|j|}{3}\right)\right)^3=\coth^3\left(\frac{\gamma}{6}\right),
\]
and $\coth x=1+\frac{2}{\exp(2x)-1}\le1+\frac{1}{x}$ for $x>0$. Finally, for the bounds on the unit lattice in Section~\ref{sec:bounds-unit-lattice} put
\[
C_0^{\mathrm{unit}}:=4b_{+}+16\,C_F\cdot c(\gamma_F),
\]
where $b_{+}$ is the number of Section~\ref{sec:running-quantities-tori} and $C_F$, $\gamma_F$ are the constants of Propositions~\ref{prop:decay-fine-propagator}--\ref{prop:bounds}.

No constant of Sections~\ref{sec:hypotheses-notation-constants}--\ref{sec:free-fermion-determinant} depends on $N$, on the level $k$, on the side $l$, on the twist, on the mass $\mathfrak{m}$, on the real pair $(m_{\mathrm{in}},\boldsymbol{\varepsilon}_{\mathrm{in}})$, on the external connection $a$, on the connection $\mathscr{C}$ in the second slot or on $\bar{m}$. Here the external connection is the complex one-cochain $a\in C^1\otimes\mathbb{C}$ on $\Lambda_n$ of Definition~\ref{def:twist-external-connection}, and $\mathscr{C}$ is the connection that occupies the second slot (Section~\ref{sec:slots-connection-covering}), an element of $\mathfrak{X}_k$ at the level $k$ (Definition~\ref{def:domains}); along the flow $\mathscr{C}=a^{[k]}$ (Definition~\ref{def:twist-external-connection}), whereas in the statements of these sections on a single level or on a single step $\mathscr{C}$ is an arbitrary element of $\mathfrak{X}_k$.

By Definition~\ref{def:twist-external-connection}, the twist $c^{(k)}_{\theta}$ is a constant connection with real values, placed in the first of the two slots. Real constant connections belong to $c\cdot\tilde{R}_k$ for $0<c\le1$ and to $\tilde{G}_k$, and these sets are invariant under the addition of real constant connections (Proposition~\ref{prop:properties-actions}(c)). Hence a statement below that is asserted on $\tilde{R}_k$, on $\tilde{G}_k$ or for arbitrary real connections (as in Theorem~\ref{thm:block-covariance}) is valid uniformly in $\theta$.

Sections~\ref{sec:lattice-symmetries-local}--\ref{sec:contraction-reblocking} include the level $k=0$, the first step; there $\eta=1$, and the pairs $(x,x')$ with $0<|x-x'|\le1$ are the pairs of neighbours for the sup norm. The level $0$ is covered by Lemma~\ref{lem:block-determinants}, through the case $k=0$ of part~(iii) and through part~(iv), which, together with~\eqref{eq:polymer-activities-25}, is stated for every $k\ge0$, and by Theorem~\ref{thm:free-determinant-across}, whose levels $j$ start at $j=0$; for $k\ge1$ the level $0$ is among its levels $0\le j<k$.

\subsubsection{Standing conventions}\label{sec:standing-conventions}

The following conventions are in force. Let $X$ be a union of $M$-cubes of $T_k$. A \emph{closed straight segment} of $\mathbb{R}^3/l\mathbb{Z}^3$ is the image $\sigma$ of a closed straight segment $S$ of $\mathbb{R}^3$ under the projection of $\mathbb{R}^3$ onto $\mathbb{R}^3/l\mathbb{Z}^3$, where $S$ may reduce to a single point; $\sigma$ is taken together with $S$, and $\mathrm{length}(\sigma)$ is the Euclidean length of $S$. A \emph{net} is a connected set $\mathcal{N}\subset\mathbb{R}^3/l\mathbb{Z}^3$ taken together with finitely many closed straight segments $\sigma_1,\dots,\sigma_r$ of $\mathbb{R}^3/l\mathbb{Z}^3$ such that $\mathcal{N}=\sigma_1\cup\dots\cup\sigma_r$; the \emph{length} of $\mathcal{N}$ is $\mathrm{length}(\mathcal{N}):=\sum_{j=1}^r\mathrm{length}(\sigma_j)$. Put
\begin{equation}\label{eq:polymer-activities-6}
M\cdot d_M(X):=\inf_{\mathcal{N}}\,\mathrm{length}(\mathcal{N}),
\end{equation}
the infimum over the nets $\mathcal{N}$ that meet every closed $M$-cube of $X$. Since this infimum is taken over all finite families of closed straight segments, $M\cdot d_M(X)$ is also the infimum, over the connected sets $\mathcal{N}$ that meet every closed $M$-cube of $X$, of the infimum of $\sum_j\mathrm{length}(\sigma_j)$ over all finite families $(\sigma_j)$ of closed straight segments with union $\mathcal{N}$.

Every net $\mathcal{N}$ is the union of the vertices and the edges of a finite graph whose edges are closed straight segments, two edges meeting at most in common end points, such that the lengths of the edges add up to at most $\mathrm{length}(\mathcal{N})$. Indeed, divide the segments of $\mathbb{R}^3$ that belong to $\sigma_1,\dots,\sigma_r$ into finitely many closed segments of length less than $l$, the pieces; the projection is injective on each piece, as a non-zero element of $l\mathbb{Z}^3$ has Euclidean length at least $l$. For pieces $I$, $I'$ and $n\in\mathbb{Z}^3$ the set $I\cap(I'+l\cdot n)$ is empty, a point or a closed segment, and it is non-empty for only finitely many $n$. Let $P$ be the finite set of the images of the end points of all pieces and of all non-empty sets $I\cap(I'+l\cdot n)$, and divide each piece further at its points that are mapped into $P$. If a point $p$ of one of the resulting segments, $s\subset I$, other than its end points, is mapped into the image of another one, $s'\subset I'$, then $p\in s'+l\cdot n$ for some $n$; as the image of $p$ does not lie in $P$, the point $p$ is an interior point of the segment $I\cap(I'+l\cdot n)$, and both $s$ and $s'+l\cdot n$ are the segment between the two points of $I\cap(I'+l\cdot n)$ next to $p$ that are mapped into $P$. Hence the images of two of the resulting segments either coincide or meet at most in common end points. The points of $P$, as vertices, and the images of the resulting segments that do not reduce to a point, each taken once, as edges, form the required graph; the lengths of its edges add up to at most the sum of the lengths of the pieces, which is $\mathrm{length}(\mathcal{N})$.

The infimum over those nets that are trees, that is, that contain no cycle when they are subdivided into a finite graph, is the same number. Indeed, let $\mathcal{N}$ be a net that meets every closed $M$-cube of $X$, and choose a point of $\mathcal{N}$ in each closed $M$-cube of $X$; these are finitely many points. Write $\mathcal{N}$ as the union of the vertices and the edges of a finite graph as in the preceding paragraph, two edges meeting at most in common end points and the lengths of the edges adding up to at most $\mathrm{length}(\mathcal{N})$, and suppose that there is a cycle in this graph. Pick an edge lying on this cycle, with end points $a$ and $b$. As only finitely many points have been chosen, the interior of this edge contains a non-empty open segment $J$ that avoids all of them; take $J$ away from $\mathcal{N}$. The set $\mathcal{N}\setminus J$, taken together with the edges of the graph, the chosen edge being replaced by the two closed segments that remain of it, is again a net, since $a$ and $b$ are still connected through the remaining edges of the cycle; its length is the sum of the lengths of the edges diminished by the length of $J$, and is therefore smaller than $\mathrm{length}(\mathcal{N})$; it contains all chosen points and therefore meets every closed $M$-cube of $X$; and the number of independent cycles of its graph is smaller by one. After finitely many such steps a tree is obtained which meets every closed $M$-cube of $X$ and is not longer than $\mathcal{N}$.

Further,
\begin{equation}\label{eq:polymer-activities-7}
\bar{X}^L:=\text{union of all $LM$-cubes in which at least one $M$-cube of $X$ lies.}
\end{equation}
The $M$-cubes are nested in the $LM$-cubes; therefore an $LM$-cube contains an $M$-cube of $X$ if and only if it has a site in common with $X$.

Spinor functions, the generators and the kernels obey
\begin{equation}\label{eq:polymer-activities-8}
\text{the antiperiodicity~\eqref{eq:polymer-activities-AP} of Section~\ref{sec:antiperiodicity}.}
\end{equation}

For the families of Definition~\ref{def:renormalisation-group-1},
\begin{equation}\label{eq:polymer-activities-9}
\text{the clauses~\eqref{it:renormalisation-group-1-k3}--\eqref{it:renormalisation-group-1-k7} are used in the form of Section~\ref{sec:lattice-gauge-charge}.}
\end{equation}

A cut-off function is defined at the points of $\mathbb{R}^3/l\mathbb{Z}^3$, and, by Definition~\ref{def:barycentres}, its value at a bond is its value at the barycentre of the bond, which is the midpoint. For a cut-off function $f$ and the bond of $T_k$ joining $x$ to $x+\eta\cdot e_{\mu}$ this means
\begin{equation}\label{eq:polymer-activities-10}
f(x,x+\eta\cdot e_{\mu}):=f\left(x+\frac{\eta}{2}\cdot e_{\mu}\right).
\end{equation}

Let $O$ be an operator acting on spinors on $T^0_k$. The following convention concerns the bounds $|O(y,y')|\le c$ on the $4\times4$ blocks $O(y,y')$ of such operators that are taken over from the propositions of Section~\ref{sec:fermion-operators-bounds}, where bounds on blocks occur in Proposition~\ref{prop:kernel-bounds-2}. In that section either the Euclidean norm or the largest modulus of a component may be taken as the norm on $\mathbb{C}^4$, and a $4\times4$ block carries the operator norm belonging to this choice. In Sections~\ref{sec:norms-gaussian-integrals}--\ref{sec:free-fermion-determinant} exactly one bound of this kind occurs, the second hypothesis $\left|M_k(A)(y,y')\right|\le C_F\cdot\exp\left(-\gamma_F\cdot d(y,y')\right)$ of Theorem~\ref{thm:unit-lattice-bounds}; it is used only in the proof of part~(a) of that theorem, and there only through the following estimate of the sixteen entries of the block:
\begin{equation}\label{eq:polymer-activities-11}
\left|O(y,y')_{\boldsymbol{\alpha}\boldsymbol{\beta}}\right|\le c\qquad\text{for all spinor indices }\boldsymbol{\alpha},\boldsymbol{\beta}.
\end{equation}
For either norm, $|O(y,y')|\le c$ implies~\eqref{eq:polymer-activities-11}. Indeed, let $w$ be the $\boldsymbol{\beta}$-th vector of the standard basis of $\mathbb{C}^4$; in both norms $|w|=1$, and the modulus of each component of a vector is at most the norm of the vector. Hence every complex $4\times4$ matrix $B$ satisfies
\[
\left|B_{\boldsymbol{\alpha}\boldsymbol{\beta}}\right|=\left|(B\cdot w)_{\boldsymbol{\alpha}}\right|\le|B\cdot w|\le|B|\cdot|w|=|B|.
\]
The convention does not apply in Section~\ref{sec:block-covariance-real}. There the modulus $|B|$ of a $4\times4$ matrix $B$ is its operator norm for the Euclidean norm on $\mathbb{C}^4$, as fixed at the beginning of that section and in Proposition~\ref{prop:block-covariance-real}. The proof of Theorem~\ref{thm:block-covariance} establishes the bound on the blocks $D_k(A)^{-1}(y,y')$ in this norm and derives from it the bound on the rows and the columns of the matrix of $D_k(A)^{-1}$ through the inequalities $\sum_{\boldsymbol{\beta}'}\left|B_{\boldsymbol{\beta}\boldsymbol{\beta}'}\right|\le2|B|$ and $\sum_{\boldsymbol{\beta}'}\left|B_{\boldsymbol{\beta}'\boldsymbol{\beta}}\right|\le2|B|$, which hold for this norm; bounds $\left|B_{\boldsymbol{\alpha}\boldsymbol{\beta}}\right|\le c$ on the sixteen entries alone give only $\sum_{\boldsymbol{\beta}'}\left|B_{\boldsymbol{\beta}\boldsymbol{\beta}'}\right|\le4c$.

The notion of a pair admissible for a polymer (Definition~\ref{def:pairs-sites-holder}) is used, with the same words, for every union $Z$ of $M$-cubes: a pair $(x,x')\in\mathrm{Pz}_k$ is admissible for $Z$ if
\begin{multline}\label{eq:polymer-activities-12}
x,x'\in Z\ \text{and the six paths}\ \Gamma^{\pi}(x,x')\ \text{lie in}\ Z,\\
\text{or}\ x,x'\ \text{are nearest neighbours in}\ T_k\ \text{and}\ \{x,x'\}\cap Z\ne\emptyset.
\end{multline}

\subsection{Norms and Gaussian integrals on Grassmann algebras}\label{sec:norms-gaussian-integrals}

\paragraph{Gaussian integrals.}
Let $W$ be a family of generators $W(u)$ of the kind introduced in Section~\ref{sec:grassmann-algebras-norms}, with the index set $\{u\}$, $u=(y,\boldsymbol{\beta},\boldsymbol{\omega})$; here $y$ is a site of the unit lattice $T^0_k$, $\boldsymbol{\beta}$ is a spinor index and $\boldsymbol{\omega}\in\{0,1\}$. Let $\boldsymbol{\phi}=(\boldsymbol{\phi}_1,\dots,\boldsymbol{\phi}_p)$ be further families of generators, and for a $p$-tuple $I=(I_1,\dots,I_p)$ of subsets of their index sets put $\boldsymbol{\phi}(I):=\boldsymbol{\phi}_1(I_1)\cdots\boldsymbol{\phi}_p(I_p)$. For $J\subset\{u\}$ let $W(J)$ be the product of the generators $W(u)$, $u\in J$, in increasing order (Section~\ref{sec:coefficients-kernels-grassmann}). For $b=(y,\boldsymbol{\beta})$ write
\[
W_b=W_{\boldsymbol{\beta}}(y):=W(y,\boldsymbol{\beta},0),\qquad\bar{W}_b=\bar{W}_{\boldsymbol{\beta}}(y):=W(y,\boldsymbol{\beta},1).
\]
For a complex matrix $O=(O_{bc})$ with rows and columns indexed by $\{(y,\boldsymbol{\beta})\}$ let $O W$ and $O\bar{W}$ be the families $(O W)_b:=\sum_cO_{bc}\cdot W_c$ and $(O\bar{W})_b:=\sum_cO_{bc}\cdot\bar{W}_c$, and put
\[
\left\langle\bar{W},O W\right\rangle:=\sum_b\bar{W}_b\cdot(O W)_b,\qquad\left\langle\bar{W},W\right\rangle:=\sum_{y,\boldsymbol{\beta}}\bar{W}_{\boldsymbol{\beta}}(y)\cdot W_{\boldsymbol{\beta}}(y).
\]
The elements $W_b\cdot\bar{W}_b$ are even, hence commute with each other, and the monomial
\[
\mathrm{Top}:=\prod_{(y,\boldsymbol{\beta})}W_{\boldsymbol{\beta}}(y)\cdot\bar{W}_{\boldsymbol{\beta}}(y)
\]
does not depend on the order of its factors. Let $\int\mathrm{D}W$ be the linear map of $\mathrm{Gr}[\boldsymbol{\phi},W]$ to $\mathrm{Gr}[\boldsymbol{\phi}]$ that sends $\boldsymbol{\phi}(I)\cdot W(J)$, for every $p$-tuple $I$ and every $J\subset\{u\}$, to $\boldsymbol{\phi}(I)$ multiplied by the coefficient of $\mathrm{Top}$ in $W(J)$. The summands $\bar{W}_b\cdot W_b$ of $\left\langle\bar{W},W\right\rangle$ are even, commute with each other and have square zero, and $-\bar{W}_b\cdot W_b=W_b\cdot\bar{W}_b$; therefore
\[
\exp\left(-\left\langle\bar{W},W\right\rangle\right)=\prod_{(y,\boldsymbol{\beta})}\left(1+W_{\boldsymbol{\beta}}(y)\cdot\bar{W}_{\boldsymbol{\beta}}(y)\right).
\]
The coefficient of $\mathrm{Top}$ in this product is $1$, so that $\int\mathrm{D}W\,\exp\left(-\left\langle\bar{W},W\right\rangle\right)=1$; this is the normalization of Section~\ref{sec:fermion-integrals-blocking}. For a complex matrix $\operatorname{Cov}$ with rows and columns indexed by $\{(y,\boldsymbol{\beta})\}$ that has an inverse $\operatorname{Cov}^{-1}$ let $\mathrm{d}\mu_I(W)$ and $\mathrm{d}\mu_{\operatorname{Cov}}(W)$ denote the linear maps of $\mathrm{Gr}[\boldsymbol{\phi},W]$ to $\mathrm{Gr}[\boldsymbol{\phi}]$ given by
\begin{align*}
\int\mathrm{d}\mu_I(W)\,F&:=\int\mathrm{D}W\,\exp\left(-\left\langle\bar{W},W\right\rangle\right)\cdot F,\\
\int\mathrm{d}\mu_{\operatorname{Cov}}(W)\,F&:=\det(\operatorname{Cov})\cdot\int\mathrm{D}W\,\exp\left(-\left\langle\bar{W},\operatorname{Cov}^{-1}W\right\rangle\right)\cdot F.
\end{align*}
The same definitions apply to every family of this kind; for $\Psi$ they give $\int\mathrm{D}\Psi$, $\mathrm{d}\mu_I(\Psi)$ and $\mathrm{d}\mu_{\operatorname{Cov}}(\Psi)$.

\paragraph{Paired sets.}
A subset $J\subset\{u\}$ is called \emph{paired} if, for every $(y,\boldsymbol{\beta})$, it contains $(y,\boldsymbol{\beta},0)$ if and only if it contains $(y,\boldsymbol{\beta},1)$. For $J\subset\{u\}$ put $m(J):=\int\mathrm{d}\mu_I(W)\,W(J)$. Then
\begin{equation}\tag{m}\label{eq:polymer-activities-m}
\begin{cases}
m(J)=0&\text{if $J$ is not paired},\\
m(J)\in\{1,-1\}\ \text{and}\ W(J)=m(J)\cdot\prod_{(y,\boldsymbol{\beta},0)\in J}W_{\boldsymbol{\beta}}(y)\cdot\bar{W}_{\boldsymbol{\beta}}(y)&\text{if $J$ is paired}.
\end{cases}
\end{equation}

\begin{proof}[Proof of~\eqref{eq:polymer-activities-m}]
The element $\exp\left(-\left\langle\bar{W},W\right\rangle\right)$ is even and commutes with $W(J)$; so $m(J)$ is the coefficient of $\mathrm{Top}$ in
\[
W(J)\cdot\prod_{(y,\boldsymbol{\beta})}\left(1+W_{\boldsymbol{\beta}}(y)\cdot\bar{W}_{\boldsymbol{\beta}}(y)\right)=\sum_QW(J)\cdot\prod_{(y,\boldsymbol{\beta})\in Q}W_{\boldsymbol{\beta}}(y)\cdot\bar{W}_{\boldsymbol{\beta}}(y),
\]
the sum being over the subsets $Q$ of $\{(y,\boldsymbol{\beta})\}$. The term of $Q$ has a non-zero component along $\mathrm{Top}$ only if, for every $(y,\boldsymbol{\beta})$, each of the two generators $W_{\boldsymbol{\beta}}(y)$, $\bar{W}_{\boldsymbol{\beta}}(y)$ occurs in it exactly once, both of them supplied by $W(J)$ or both by the factor $W_{\boldsymbol{\beta}}(y)\cdot\bar{W}_{\boldsymbol{\beta}}(y)$ of the product over $Q$. In that case $J$ consists of the triples $(y,\boldsymbol{\beta},\boldsymbol{\omega})$ with $(y,\boldsymbol{\beta})\notin Q$, and $J$ is paired. Hence $m(J)=0$ if $J$ is not paired. Let $J$ be paired, and let $P_J:=\prod_{(y,\boldsymbol{\beta},0)\in J}W_{\boldsymbol{\beta}}(y)\cdot\bar{W}_{\boldsymbol{\beta}}(y)$, a product of commuting even factors. A reordering of the factors of $W(J)$ shows that $W(J)=P_J$ or $W(J)=-P_J$. Only the term with $Q=\{(y,\boldsymbol{\beta}):(y,\boldsymbol{\beta},0)\notin J\}$ contributes, and since $P_J\cdot\prod_{(y,\boldsymbol{\beta})\in Q}W_{\boldsymbol{\beta}}(y)\cdot\bar{W}_{\boldsymbol{\beta}}(y)=\mathrm{Top}$, this term is $\mathrm{Top}$ if $W(J)=P_J$ and $-\mathrm{Top}$ if $W(J)=-P_J$. Hence $m(J)\in\{1,-1\}$ and $W(J)=m(J)\cdot P_J$.
\end{proof}

\paragraph{Contour integrals.}
For $\rho>1$ and a continuous function $f$ on the circle $\{t\in\mathbb{C}:|t|=\rho\}$ with values in a complex normed space of finite dimension let
\[
\oint_{\rho}f:=\frac{1}{2\pi i}\oint_{|t|=\rho}\mathrm{d}t\,\frac{f(t)}{t(t-1)},\qquad\oint^{(2)}_{\rho}f:=\frac{1}{2\pi i}\oint_{|t|=\rho}\mathrm{d}t\,\frac{f(t)}{t^2(t-1)},
\]
the circle being oriented counterclockwise.

\begin{theorem}\label{thm:polymer-activities-1}
Let the index sets $S_1,\dots,S_p$ with the weights $\nu_i$, the algebra $\mathrm{Gr}[\boldsymbol{\phi}_1,\dots,\boldsymbol{\phi}_p]$, the coefficients $c_F$, the norms $\|\cdot\|_g$, $g=(g_1,\dots,g_p)$, and the row sums $|\cdot|_{1,\infty}$ be as in Section~\ref{sec:coefficients-kernels-grassmann}, and let $W$ be a family of generators as above.

\emph{Kernels.} Let $F$ be given by~\eqref{eq:polymer-activities-2} with arbitrary complex kernels $K_n$, $n=(n_1,\dots,n_p)$. Then
\[
\|F\|_g\le\sum_n\prod_{i=1}^p\frac{g_i^{n_i}}{n_i!}\int\mathrm{d}s\,|K_n|,
\]
and equality holds if every $K_n$ is antisymmetric, that is, changes sign whenever two of its arguments that belong to the same family are exchanged. Kernels with this antisymmetry are uniquely fixed by $F$: for them,
\[
K_n(I_1;\dots;I_p)\cdot\prod_{i=1}^p\nu_i(I_i)=c_F(I_1,\dots,I_p),\qquad\nu_i(I):=\prod_{s\in I}\nu_i(s),
\]
for all subsets $I_i\subset S_i$ with $\#|I_i|=n_i$, where $K_n(I_1;\dots;I_p)$ denotes the value of $K_n$ at the elements of $I_1,\dots,I_p$, each set listed in increasing order.

\emph{Products.} For $F,G\in\mathrm{Gr}[\boldsymbol{\phi}_1,\dots,\boldsymbol{\phi}_p]$,
\[
\|F\cdot G\|_g\le\|F\|_g\cdot\|G\|_g.
\]
The even part of $\mathrm{Gr}[\boldsymbol{\phi}_1,\dots,\boldsymbol{\phi}_p]$, normed by $\|\cdot\|_g$, is a commutative unital Banach algebra, and $\|\exp F-1\|_g\le\exp\left(\|F\|_g\right)-1$ for every even $F$.

\emph{Linear substitutions.} Let $\boldsymbol{\phi}_1,\dots,\boldsymbol{\phi}_p$ and $\boldsymbol{\theta}_1,\dots,\boldsymbol{\theta}_q$ be families of generators, $\boldsymbol{\phi}_i$ indexed by a finite ordered set $T_i$ and $\boldsymbol{\theta}_j$ by a finite ordered set $S_j$ with weights $\mu_j(s)>0$. For all $i$, $j$ let $O_{ij}$ be an operator from functions on $S_j$ to functions on $T_i$, $(O_{ij}f)(t)=\sum_{s\in S_j}O_{ij}(t,s)\cdot f(s)\cdot\mu_j(s)$, with $|O_{ij}|_{1,\infty}\le C_{ij}$, and let $F\mapsto F'$ be the homomorphism of $\mathrm{Gr}[\boldsymbol{\phi}]$ to $\mathrm{Gr}[\boldsymbol{\theta}]$ with
\[
\boldsymbol{\phi}_i(t)\mapsto\sum_{j=1}^q\left(O_{ij}\boldsymbol{\theta}_j\right)(t)=\sum_{j=1}^q\sum_{s\in S_j}O_{ij}(t,s)\cdot\mu_j(s)\cdot\boldsymbol{\theta}_j(s)\qquad(1\le i\le p,\ t\in T_i).
\]
If the weights $g=(g_1,\dots,g_p)$ and $g'=(g'_1,\dots,g'_q)$ satisfy $\sum_jC_{ij}\cdot g'_j\le g_i$ for all $i$, then
\[
\|F'\|_{g'}\le\|F\|_g\qquad\left(F\in\mathrm{Gr}[\boldsymbol{\phi}]\right).
\]

\emph{Integration.} For $F\in\mathrm{Gr}[\boldsymbol{\phi}_1,\dots,\boldsymbol{\phi}_p,W]$,
\[
\left\|\int\mathrm{d}\mu_I(W)\,F\right\|_{(g_1,\dots,g_p)}\le\|F\|_{(g_1,\dots,g_p,1)}.
\]

\emph{Factorization.} Let $F_1,F_2\in\mathrm{Gr}[\boldsymbol{\phi}_1,\dots,\boldsymbol{\phi}_p,W]$. Suppose that the set of the unit sites $y$ at which $F_1$ contains a generator $W(y,\cdot,\cdot)$, that is, at which some $W(y,\boldsymbol{\beta},\boldsymbol{\omega})$ is a factor of a basis monomial with non-zero coefficient in $F_1$, is disjoint from the corresponding set of $F_2$. Then
\[
\int\mathrm{d}\mu_I(W)\,\left(F_1\cdot F_2\right)=\left(\int\mathrm{d}\mu_I(W)\,F_1\right)\cdot\left(\int\mathrm{d}\mu_I(W)\,F_2\right).
\]

\emph{Invariance.} Let $U$ be a complex matrix with rows and columns indexed by $\{(y,\boldsymbol{\beta})\}$ that has an inverse $U^{-1}$, and let $s_U$ be the homomorphism of $\mathrm{Gr}[\boldsymbol{\phi}_1,\dots,\boldsymbol{\phi}_p,W]$ to itself with $W\mapsto U W$ and $\bar{W}\mapsto\left(U^{-1}\right)^T\bar{W}$ that fixes the generators of the families $\boldsymbol{\phi}_i$. Then
\[
\int\mathrm{d}\mu_I(W)\,s_U(F)=\int\mathrm{d}\mu_I(W)\,F\qquad\left(F\in\mathrm{Gr}[\boldsymbol{\phi}_1,\dots,\boldsymbol{\phi}_p,W]\right).
\]

\emph{Contour integrals.} Let $\rho>1$, and let $f$ be an analytic function on a neighbourhood of the closed disc $\{t\in\mathbb{C}:|t|\le\rho\}$ with values in a complex normed space of finite dimension, with the norm $\|\cdot\|$. Then
\begin{align*}
\oint_{\rho}f&=f(1)-f(0),&\left\|\oint_{\rho}f\right\|&\le\frac{1}{\rho-1}\cdot\sup_{|t|=\rho}\|f(t)\|,\\
\oint^{(2)}_{\rho}f&=f(1)-f(0)-f'(0),&\left\|\oint^{(2)}_{\rho}f\right\|&\le\frac{1}{\rho(\rho-1)}\cdot\sup_{|t|=\rho}\|f(t)\|.
\end{align*}
In particular, let $F\in\mathrm{Gr}[\boldsymbol{\phi}_1,\dots,\boldsymbol{\phi}_p]$, let $\boldsymbol{\theta}=(\boldsymbol{\theta}_1,\dots,\boldsymbol{\theta}_p)$ be families of generators with $\boldsymbol{\theta}_i$ indexed by $S_i$, and for $t\in\mathbb{C}$ let $F(\boldsymbol{\phi}+t\cdot\boldsymbol{\theta})\in\mathrm{Gr}[\boldsymbol{\phi},\boldsymbol{\theta}]$ be the image of $F$ under the homomorphism with $\boldsymbol{\phi}_i(s)\mapsto\boldsymbol{\phi}_i(s)+t\cdot\boldsymbol{\theta}_i(s)$. If the weights $g'=(g'_1,\dots,g'_p)$ of $\boldsymbol{\phi}$ and $g''=(g''_1,\dots,g''_p)$ of $\boldsymbol{\theta}$ satisfy $g'_i+\rho\cdot g''_i\le g_i$ for all $i$, then
\[
\left\|F(\boldsymbol{\phi}+\boldsymbol{\theta})-F(\boldsymbol{\phi})\right\|_{(g',g'')}\le\frac{\|F\|_g}{\rho-1}.
\]

\emph{Monotonicity in the weights.} If $0<g'_i\le g_i$ for all $i$, then
\[
\|F\|_{g'}\le\|F\|_g\qquad\left(F\in\mathrm{Gr}[\boldsymbol{\phi}_1,\dots,\boldsymbol{\phi}_p]\right).
\]
If moreover $r_1,\dots,r_p\ge0$ are integers and every basis monomial with non-zero coefficient in $F$ has degree $n_i\ge r_i$ in $\boldsymbol{\phi}_i$ for every $i$, then
\[
\|F\|_{g'}\le\prod_{i=1}^p\left(\frac{g'_i}{g_i}\right)^{r_i}\cdot\|F\|_g.
\]

\emph{Norms on $\mathrm{Gr}[\Psi]$.} For a number $a>0$ let $|F|_a$ be the norm~\eqref{eq:polymer-activities-1} of $F\in\mathrm{Gr}[\Psi]$ with the single weight $a$, that is,
\[
|F|_a:=\sum_{J\subset\{u\}}|c_F(J)|\cdot a^{\#|J|}.
\]
Then the following hold.
\begin{enumerate}
\item[(i)] $|F\cdot G|_a\le|F|_a\cdot|G|_a$ for $F,G\in\mathrm{Gr}[\Psi]$.
\item[(ii)] $\left|\int\mathrm{d}\mu_I(\Psi)\,F\right|\le|F|_1$ for $F\in\mathrm{Gr}[\Psi]$.
\item[(iii)] For every invertible matrix $\operatorname{Cov}$ on $\{(y,\boldsymbol{\beta})\}$ and every $f=f(\bar{\Psi},\Psi)\in\mathrm{Gr}[\Psi]$,
\[
\int\mathrm{d}\mu_{\operatorname{Cov}}(\Psi)\,f(\bar{\Psi},\Psi)=\int\mathrm{d}\mu_I(\Psi)\,f(\bar{\Psi},\operatorname{Cov}\Psi),
\]
where $f(\bar{\Psi},\operatorname{Cov}\Psi)$ is the image of $f$ under the homomorphism with $\Psi\mapsto\operatorname{Cov}\Psi$ and $\bar{\Psi}\mapsto\bar{\Psi}$. Moreover $\int\mathrm{D}\Psi\,\exp\left(-\left\langle\bar{\Psi},O\Psi\right\rangle\right)=\det O$ for every matrix $O$ on $\{(y,\boldsymbol{\beta})\}$.
\item[(iv)] Let $E\in\mathrm{Gr}[\psi,\chi]$, let $h_1,h_2>0$, let $C_0>0$ be an arbitrary number, and let $H^{\#}=(H^{(1)},H^{(2)})$ be a pair of operators from functions on $\{u\}$ to functions on $\{\xi\}$, resp.\ $\{\boldsymbol{\zeta}\}$ (the index sets of Section~\ref{sec:grassmann-algebras-norms}, the set $\{u\}$ with the weight $1$), with $\left|H^{(i)}\right|_{1,\infty}\le C_0$ for $i=1,2$. Let $E(H^{(1)}\Psi,H^{(2)}\Psi)\in\mathrm{Gr}[\Psi]$ be the image of $E$ under the homomorphism with $\psi(\xi)\mapsto\left(H^{(1)}\Psi\right)(\xi)$ and $\chi(\boldsymbol{\zeta})\mapsto\left(H^{(2)}\Psi\right)(\boldsymbol{\zeta})$. If $C_0\cdot a\le\min(h_1,h_2)$, then
\[
\left|E(H^{(1)}\Psi,H^{(2)}\Psi)\right|_a\le\|E\|_{(h_1,h_2)}.
\]
\end{enumerate}
\end{theorem}

The nine paragraphs of Theorem~\ref{thm:polymer-activities-1}, from \emph{Kernels} to \emph{Norms on $\mathrm{Gr}[\Psi]$}, are, in their order, the items~\eqref{it:weighted-coefficient-norms-N0}--\eqref{it:weighted-coefficient-norms-N7} of Proposition~\ref{prop:weighted-coefficient-norms}, and each paragraph is referred to by the mark of the corresponding item.

\begin{proof}
\emph{Proof of~\eqref{it:weighted-coefficient-norms-N0}.} The computation is written for $p=1$; for several families the same computation is done for every family separately, with one permutation for each family. Let $n\ge0$, and consider the sum
\[
\sum_{s_1,\dots,s_n\in S_1}K_n(s_1,\dots,s_n)\cdot\nu_1(s_1)\cdots\nu_1(s_n)\cdot\boldsymbol{\phi}_1(s_1)\cdots\boldsymbol{\phi}_1(s_n).
\]
Its terms in which two arguments coincide vanish, because $\boldsymbol{\phi}_1(s)^2=0$. The other terms are grouped according to the set $I=\{s_1,\dots,s_n\}$: if $I=\{t_1,\dots,t_n\}$ with $t_1<\dots<t_n$, then $\boldsymbol{\phi}_1(t_{\pi(1)})\cdots\boldsymbol{\phi}_1(t_{\pi(n)})=\operatorname{sgn}(\pi)\cdot\boldsymbol{\phi}_1(I)$ for every permutation $\pi$ of $\{1,\dots,n\}$. Hence
\[
c_F(I)=\frac{1}{n!}\sum_{\pi}\operatorname{sgn}(\pi)\cdot K_n\left(t_{\pi(1)},\dots,t_{\pi(n)}\right)\cdot\nu_1(I),
\]
and consequently
\[
\sum_{\#|I|=n}|c_F(I)|\cdot g_1^n\le\frac{g_1^n}{n!}\sum_{s_1,\dots,s_n\ \text{distinct}}|K_n(s_1,\dots,s_n)|\cdot\nu_1(s_1)\cdots\nu_1(s_n)\le\frac{g_1^n}{n!}\int\mathrm{d}s\,|K_n|.
\]
Summation over $n$ gives~\eqref{it:weighted-coefficient-norms-N0}. If $K_n$ is antisymmetric, then the $n!$ terms $\operatorname{sgn}(\pi)\cdot K_n(t_{\pi(1)},\dots,t_{\pi(n)})$ all equal $K_n(t_1,\dots,t_n)$, and $K_n$ vanishes where two arguments coincide; so both inequalities are equalities, and $c_F(I)=K_n(t_1,\dots,t_n)\cdot\nu_1(I)$.

\emph{Proof of~\eqref{it:weighted-coefficient-norms-N1}.} Let $I=(I_1,\dots,I_p)$ and $J=(J_1,\dots,J_p)$ be $p$-tuples of subsets. The product $\boldsymbol{\phi}(I)\cdot\boldsymbol{\phi}(J)$ vanishes if $I_i\cap J_i\ne\emptyset$ for some $i$; otherwise it equals $\boldsymbol{\phi}(I_1\cup J_1,\dots,I_p\cup J_p)$ or its negative, and the weight $\prod_ig_i^{\#|I_i|+\#|J_i|}$ of this basis monomial in~\eqref{eq:polymer-activities-1} equals the weight of $\boldsymbol{\phi}(I)$ times that of $\boldsymbol{\phi}(J)$. Expanding $F\cdot G=\sum_{I,J}c_F(I)\cdot c_G(J)\cdot\boldsymbol{\phi}(I)\cdot\boldsymbol{\phi}(J)$ one obtains
\[
\|F\cdot G\|_g\le\sum_{I,J}|c_F(I)|\cdot|c_G(J)|\cdot\prod_{i=1}^pg_i^{\#|I_i|+\#|J_i|}=\|F\|_g\cdot\|G\|_g.
\]
Even elements commute with all elements, products of even elements are even, $\|1\|_g=1$, and the algebra has finite dimension, hence is complete; so the even part is a commutative unital Banach algebra. For even $F$ the inequality just proved gives $\|F^n\|_g\le\|F\|_g^n$; hence the series of $\exp F$ converges absolutely for $\|\cdot\|_g$, and
\[
\|\exp F-1\|_g\le\sum_{n\ge1}\frac{\|F\|_g^n}{n!}=\exp\left(\|F\|_g\right)-1.
\]

\emph{Proof of~\eqref{it:weighted-coefficient-norms-N2}.} Let $1\le i\le p$ and $t\in T_i$. The image $\sum_j\sum_{s\in S_j}\left[O_{ij}(t,s)\cdot\mu_j(s)\right]\cdot\boldsymbol{\theta}_j(s)$ of $\boldsymbol{\phi}_i(t)$ is a linear combination of generators, and by~\eqref{eq:polymer-activities-1} its norm for the weights $g'$ is
\[
\sum_{j=1}^q\sum_{s\in S_j}|O_{ij}(t,s)|\cdot\mu_j(s)\cdot g'_j\le\sum_{j=1}^qC_{ij}\cdot g'_j\le g_i.
\]
These images are odd; they anticommute with each other and have square zero, so that the substitution extends to the homomorphism $F\mapsto F'$. By~\eqref{it:weighted-coefficient-norms-N1}, every $\boldsymbol{\phi}(I)$ is mapped to an element of norm at most $\prod_ig_i^{\#|I_i|}$. Multiplying by $|c_F(I)|$ and summing over $I$ gives $\|F'\|_{g'}\le\|F\|_g$.

\emph{Proof of~\eqref{it:weighted-coefficient-norms-N3}.} For $G_1\in\mathrm{Gr}[\boldsymbol{\phi}]$ and $G_2\in\mathrm{Gr}[W]$ the even element $\exp\left(-\left\langle\bar{W},W\right\rangle\right)$ commutes with $G_1$, and the definition of $\int\mathrm{D}W$ gives
\[
\int\mathrm{d}\mu_I(W)\,\left(G_1\cdot G_2\right)=G_1\cdot\int\mathrm{d}\mu_I(W)\,G_2;
\]
in particular $\int\mathrm{d}\mu_I(W)\,\boldsymbol{\phi}(I)\cdot W(J)=m(J)\cdot\boldsymbol{\phi}(I)$. Writing $F=\sum_{I,J}c_F(I,J)\cdot\boldsymbol{\phi}(I)\cdot W(J)$, with $p$-tuples $I$ and subsets $J\subset\{u\}$, one obtains $\int\mathrm{d}\mu_I(W)\,F=\sum_{I,J}c_F(I,J)\cdot m(J)\cdot\boldsymbol{\phi}(I)$. As $|m(J)|\le1$ by~\eqref{eq:polymer-activities-m},
\[
\left\|\int\mathrm{d}\mu_I(W)\,F\right\|_{(g_1,\dots,g_p)}\le\sum_{I,J}|c_F(I,J)|\cdot\prod_{i=1}^pg_i^{\#|I_i|}=\|F\|_{(g_1,\dots,g_p,1)}.
\]

\emph{Proof of~\eqref{it:weighted-coefficient-norms-N4}.} Both sides of~\eqref{it:weighted-coefficient-norms-N4} are bilinear in $(F_1,F_2)$, and the hypothesis holds for every pair formed by a basis monomial with non-zero coefficient in $F_1$ and one in $F_2$. So let $F_1=\boldsymbol{\phi}(I_1)\cdot W(J_1)$ and $F_2=\boldsymbol{\phi}(I_2)\cdot W(J_2)$, where no unit site carries both an element of $J_1$ and an element of $J_2$, and let $d_2$ be the number of factors of $\boldsymbol{\phi}(I_2)$. Then
\[
F_1\cdot F_2=(-1)^{\#|J_1|\cdot d_2}\cdot\boldsymbol{\phi}(I_1)\cdot\boldsymbol{\phi}(I_2)\cdot W(J_1)\cdot W(J_2).
\]
By the proof of~\eqref{it:weighted-coefficient-norms-N3}, $\int\mathrm{d}\mu_I(W)\,F_{\ell}=m(J_{\ell})\cdot\boldsymbol{\phi}(I_{\ell})$ for $\ell=1,2$, and the left side of~\eqref{it:weighted-coefficient-norms-N4} equals $(-1)^{\#|J_1|\cdot d_2}\cdot\boldsymbol{\phi}(I_1)\cdot\boldsymbol{\phi}(I_2)\cdot\int\mathrm{d}\mu_I(W)\,W(J_1)\cdot W(J_2)$. The sets $J_1$ and $J_2$ are disjoint, so that $W(J_1)\cdot W(J_2)=\pm W(J_1\cup J_2)$. The two triples $(y,\boldsymbol{\beta},0)$ and $(y,\boldsymbol{\beta},1)$ have the same site $y$, and no site carries elements of both $J_1$ and $J_2$; hence $J_1\cup J_2$ is paired only if $J_1$ and $J_2$ are paired. If one of $J_1$, $J_2$ fails to be paired, then so does $J_1\cup J_2$, and by~\eqref{eq:polymer-activities-m} the two sides of~\eqref{it:weighted-coefficient-norms-N4} are $0$. If both are paired, then $\#|J_1|$ is even, and by~\eqref{eq:polymer-activities-m}
\[
W(J_1)\cdot W(J_2)=m(J_1)\cdot m(J_2)\cdot P,\qquad P:=\prod_{(y,\boldsymbol{\beta},0)\in J_1\cup J_2}W_{\boldsymbol{\beta}}(y)\cdot\bar{W}_{\boldsymbol{\beta}}(y).
\]
By~\eqref{eq:polymer-activities-m} for the paired set $J_1\cup J_2$, $P=m(J_1\cup J_2)\cdot W(J_1\cup J_2)$, so that $\int\mathrm{d}\mu_I(W)\,P=m(J_1\cup J_2)^2=1$. Both sides of~\eqref{it:weighted-coefficient-norms-N4} therefore equal $m(J_1)\cdot m(J_2)\cdot\boldsymbol{\phi}(I_1)\cdot\boldsymbol{\phi}(I_2)$.

\emph{Proof of~\eqref{it:weighted-coefficient-norms-N4p}.} For every $b$ the element $s_U(W_b)=(U W)_b$ is a linear combination of the $W_c$, and $s_U(\bar{W}_b)=\left(\left(U^{-1}\right)^T\bar{W}\right)_b$ is a linear combination of the $\bar{W}_c$. Hence $s_U$ maps every monomial to a linear combination of monomials with the same degree in $W$ and the same degree in $\bar{W}$; and the only basis monomial of $\mathrm{Gr}[W]$ of full degree in $W$ and in $\bar{W}$ is $W(\{u\})=\pm\mathrm{Top}$. Write $\mathrm{Top}=\pm\prod_bW_b\cdot\prod_b\bar{W}_b$, both products being taken in one fixed order. By the Leibniz formula~\eqref{eq:polymer-activities-cl5},
\[
s_U\left(\prod_bW_b\right)=\det U\cdot\prod_bW_b,\qquad s_U\left(\prod_b\bar{W}_b\right)=\det\left(\left(U^{-1}\right)^T\right)\cdot\prod_b\bar{W}_b,
\]
so that $s_U(\mathrm{Top})=\det U\cdot\det\left(\left(U^{-1}\right)^T\right)\cdot\mathrm{Top}=\mathrm{Top}$. As $s_U$ fixes the generators of the families $\boldsymbol{\phi}_i$, the element $s_U\left(\boldsymbol{\phi}(I)\cdot W(J)\right)=\boldsymbol{\phi}(I)\cdot s_U\left(W(J)\right)$ has no component along $\boldsymbol{\phi}(I)\cdot\mathrm{Top}$ if $J\ne\{u\}$, and equals $\boldsymbol{\phi}(I)\cdot W(J)$ if $J=\{u\}$. Therefore $\int\mathrm{D}W\,s_U(G)=\int\mathrm{D}W\,G$ for every $G\in\mathrm{Gr}[\boldsymbol{\phi}_1,\dots,\boldsymbol{\phi}_p,W]$. Moreover
\[
s_U\left(\left\langle\bar{W},W\right\rangle\right)=\sum_b\left(\left(U^{-1}\right)^T\bar{W}\right)_b\cdot(U W)_b=\sum_{b,c,d}\bar{W}_c\cdot\left(U^{-1}\right)_{cb}\cdot U_{bd}\cdot W_d=\left\langle\bar{W},W\right\rangle,
\]
so that $s_U\left(\exp\left(-\left\langle\bar{W},W\right\rangle\right)\cdot F\right)=\exp\left(-\left\langle\bar{W},W\right\rangle\right)\cdot s_U(F)$. The choice $G=\exp\left(-\left\langle\bar{W},W\right\rangle\right)\cdot F$ gives~\eqref{it:weighted-coefficient-norms-N4p}.

\emph{Proof of~\eqref{it:weighted-coefficient-norms-N5}.} Since $\rho>1$, the points $0$ and $1$ lie inside the circle $|t|=\rho$. The function $\frac{f(t)}{t(t-1)}$ has the residue $f(1)$ at $t=1$ and the residue $-f(0)$ at $t=0$. The function $\frac{f(t)}{t^2(t-1)}$ has the residue $f(1)$ at $t=1$; near $t=0$ one has $\frac{1}{t-1}=-\left(1+t+t^2+\cdots\right)$, so that its residue at $t=0$ is $-f(0)-f'(0)$. The residue theorem~\eqref{eq:polymer-activities-cl1}, applied to the components of $f$, gives the two identities. On the circle $|t-1|\ge|t|-1=\rho-1$, hence $|t(t-1)|\ge\rho(\rho-1)$ and $|t^2(t-1)|\ge\rho^2(\rho-1)$; as the circle has length $2\pi\rho$, the two bounds follow.

For the last assertion, $t\mapsto F(\boldsymbol{\phi}+t\cdot\boldsymbol{\theta})$ is a polynomial in $t$ with values in $\mathrm{Gr}[\boldsymbol{\phi},\boldsymbol{\theta}]$, a complex vector space of finite dimension, normed by $\|\cdot\|_{(g',g'')}$; its values at $t=1$ and at $t=0$ are $F(\boldsymbol{\phi}+\boldsymbol{\theta})$ and $F(\boldsymbol{\phi})$. Let $|t|=\rho$. The substitution $\boldsymbol{\phi}_i\mapsto\boldsymbol{\phi}_i+t\cdot\boldsymbol{\theta}_i$ is of the form considered in~\eqref{it:weighted-coefficient-norms-N2}, with the target families $\boldsymbol{\phi}_1,\dots,\boldsymbol{\phi}_p,\boldsymbol{\theta}_1,\dots,\boldsymbol{\theta}_p$: for each $i$ the operator into the target family $\boldsymbol{\phi}_i$ is the identity, whose row sums equal $1$, the operator into the target family $\boldsymbol{\theta}_i$ is $t$ times the identity, whose row sums equal $|t|=\rho$, and all other operators vanish. Since $g'_i+\rho\cdot g''_i\le g_i$, \eqref{it:weighted-coefficient-norms-N2} gives $\|F(\boldsymbol{\phi}+t\cdot\boldsymbol{\theta})\|_{(g',g'')}\le\|F\|_g$, and the bound for $\oint_{\rho}$ yields the assertion.

\emph{Proof of~\eqref{it:weighted-coefficient-norms-N6}.} By~\eqref{eq:polymer-activities-1}, the coefficient $c_F(I_1,\dots,I_p)$ enters $\|F\|_{g'}$ with the weight $\prod_i(g'_i)^{\#|I_i|}$, which does not exceed its weight $\prod_ig_i^{\#|I_i|}$ in $\|F\|_g$; this gives the first assertion. If $\#|I_i|\ge r_i$ for all $i$, then, as $\frac{g'_i}{g_i}\le1$,
\[
\prod_{i=1}^p(g'_i)^{\#|I_i|}=\prod_{i=1}^p\left(\frac{g'_i}{g_i}\right)^{\#|I_i|}\cdot\prod_{i=1}^pg_i^{\#|I_i|}\le\prod_{i=1}^p\left(\frac{g'_i}{g_i}\right)^{r_i}\cdot\prod_{i=1}^pg_i^{\#|I_i|};
\]
summation over the basis monomials with non-zero coefficients gives the second assertion.

\emph{Proof of~\eqref{it:weighted-coefficient-norms-N7}.} Statement~(i) is~\eqref{it:weighted-coefficient-norms-N1} for the single family $\Psi$ with the weight $a$. Statement~(ii) is~\eqref{it:weighted-coefficient-norms-N3} with $p=0$ and $\Psi$ in place of $W$; for $p=0$ the algebra without generators is $\mathbb{C}$, and the norm~\eqref{eq:polymer-activities-1} of a number is its modulus.

(iii) For a matrix $A$ on $\{(y,\boldsymbol{\beta})\}$ let $s_A$ be the homomorphism of $\mathrm{Gr}[\Psi]$ to itself with $\Psi\mapsto A\Psi$ and $\bar{\Psi}\mapsto\bar{\Psi}$. It preserves the degree in $\Psi$ and the degree in $\bar{\Psi}$, and by the Leibniz formula~\eqref{eq:polymer-activities-cl5} it maps $\prod_{(y,\boldsymbol{\beta})}\Psi_{\boldsymbol{\beta}}(y)$ to $\det A\cdot\prod_{(y,\boldsymbol{\beta})}\Psi_{\boldsymbol{\beta}}(y)$. Writing $\mathrm{Top}=\pm\prod_{(y,\boldsymbol{\beta})}\Psi_{\boldsymbol{\beta}}(y)\cdot\prod_{(y,\boldsymbol{\beta})}\bar{\Psi}_{\boldsymbol{\beta}}(y)$, where $\mathrm{Top}$ is now formed with $\Psi$ and both products are taken in one fixed order, one obtains $s_A(\mathrm{Top})=\det A\cdot\mathrm{Top}$, while the images of the other basis monomials have no component along $\mathrm{Top}$; therefore $\int\mathrm{D}\Psi\,s_A(F)=\det A\cdot\int\mathrm{D}\Psi\,F$ for every $F\in\mathrm{Gr}[\Psi]$. The homomorphism $s_O$ maps $\left\langle\bar{\Psi},\Psi\right\rangle$ to $\left\langle\bar{\Psi},O\Psi\right\rangle$; so $A=O$ and $F=\exp\left(-\left\langle\bar{\Psi},\Psi\right\rangle\right)$ give
\[
\int\mathrm{D}\Psi\,\exp\left(-\left\langle\bar{\Psi},O\Psi\right\rangle\right)=\det O\cdot\int\mathrm{D}\Psi\,\exp\left(-\left\langle\bar{\Psi},\Psi\right\rangle\right)=\det O.
\]
The homomorphism $s_{\operatorname{Cov}}$ maps $\left\langle\bar{\Psi},\operatorname{Cov}^{-1}\Psi\right\rangle$ to $\left\langle\bar{\Psi},\Psi\right\rangle$; so for $A=\operatorname{Cov}$ and $F=\exp\left(-\left\langle\bar{\Psi},\operatorname{Cov}^{-1}\Psi\right\rangle\right)\cdot f$ one has $s_{\operatorname{Cov}}(F)=\exp\left(-\left\langle\bar{\Psi},\Psi\right\rangle\right)\cdot s_{\operatorname{Cov}}(f)$, and therefore
\[
\int\mathrm{d}\mu_I(\Psi)\,s_{\operatorname{Cov}}(f)=\det(\operatorname{Cov})\cdot\int\mathrm{D}\Psi\,\exp\left(-\left\langle\bar{\Psi},\operatorname{Cov}^{-1}\Psi\right\rangle\right)\cdot f=\int\mathrm{d}\mu_{\operatorname{Cov}}(\Psi)\,f.
\]
Since $s_{\operatorname{Cov}}(f)=f(\bar{\Psi},\operatorname{Cov}\Psi)$, this is the first identity of~(iii).

(iv) The map $E\mapsto E(H^{(1)}\Psi,H^{(2)}\Psi)$ is the homomorphism of~\eqref{it:weighted-coefficient-norms-N2} with $p=2$ (the families $\psi$ and $\chi$, with the weights $h_1$ and $h_2$), $q=1$ (the family $\Psi$ on $\{u\}$, with the index weight $1$ and the weight $a$), $O_{11}=H^{(1)}$, $O_{21}=H^{(2)}$ and $C_{11}=C_{21}=C_0$. Its hypothesis $C_0\cdot a\le h_i$, $i=1,2$, is the hypothesis $C_0\cdot a\le\min(h_1,h_2)$, and its conclusion is the assertion of~(iv), since the norm~\eqref{eq:polymer-activities-1} on $\mathrm{Gr}[\Psi]$ with the weight $a$ is $|\cdot|_a$.
\end{proof}

\paragraph{Index sets.}
The generators are indexed as in Section~\ref{sec:grassmann-algebras-norms}: $\psi$ by $\{\xi\}$ with the weight $\eta^3$, $\chi$ by $\{\boldsymbol{\zeta}\}$ with the weight $\eta^6$, where $\{\boldsymbol{\zeta}\}$ is formed with all pairs of $\mathrm{Pz}_k$ (Definition~\ref{def:pairs-sites-holder}), and $\Psi$, $W$ by $\{u\}$ with the weight $1$. Accordingly, for the operator $H^{(2)}$ of~\eqref{it:weighted-coefficient-norms-N7}(iv), and for the operator $\delta_{\alpha,B}H$ of a substitution $\chi\mapsto\delta_{\alpha,B}H\Psi$ in~\eqref{it:weighted-coefficient-norms-N2}, where $H$ is an operator from functions on $\{u\}$ to functions on $\{\xi\}$, $B$ is a bond function and $\delta_{\alpha,B}$ is the operator of Definition~\ref{def:pairs-sites-holder}, the supremum in the row sum $|\cdot|_{1,\infty}$ is taken over \emph{all} pairs of $\mathrm{Pz}_k$.

By Theorem~\ref{thm:polymer-activities-1}, Proposition~\ref{prop:weighted-coefficient-norms} follows.

\subsection{Polymers}\label{sec:polymers}

Let $k$ and $l$ be as in~\eqref{eq:polymer-activities-Sa}. As $l\in\mathrm{Cl}$, the side $l$ is a multiple of $L\cdot M$, and $l\ge13L\cdot M$ (Section~\ref{sec:lattice-geometry}). The $M$-cubes of $T_k$ are the $\frac{l^3}{M^3}$ cubes of side $M$ whose centres are the points of $M\mathbb{Z}^3/l\mathbb{Z}^3$; the closed $M$-cube with a given centre is the closed cube of side $M$ with this centre in the continuum torus $\mathbb{R}^3/l\mathbb{Z}^3$. Two cubes of equal side are \emph{adjacent} if they have a common face, and a union of such cubes is \emph{connected through faces} if any two of its cubes are joined by a finite sequence of its cubes in which consecutive cubes are adjacent. A \emph{polymer} is a non-empty union of $M$-cubes of $T_k$ that is connected through faces, and $D_k$ is the set of the polymers (Section~\ref{sec:lattice-geometry}). For a union $X$ of $M$-cubes of $T_k$ let $|X|_M$ be the number of the $M$-cubes of $X$, and let $d_M(X)$ be given by~\eqref{eq:polymer-activities-6}. A net that meets every closed $M$-cube of $X$ is called a \emph{net for} $X$; thus $M\cdot d_M(X)$ is the infimum of the lengths of the nets for $X$.

The $M$-cubes of $\mathbb{R}^3$ are the cubes of side $M$ centred at the points of $M\mathbb{Z}^3$. A \emph{net in} $\mathbb{R}^3$ is a connected subset of $\mathbb{R}^3$ that is a finite union of closed straight segments, a point counting as a segment; its length is the sum of the Euclidean lengths of the segments. For a finite union $\tilde{X}$ of $M$-cubes of $\mathbb{R}^3$ that is connected through faces, $d_M(\tilde{X})$ is formed as in~\eqref{eq:polymer-activities-6} with nets in $\mathbb{R}^3$: $M\cdot d_M(\tilde{X})$ is the infimum of the lengths of the nets in $\mathbb{R}^3$ that meet every closed cube of $\tilde{X}$.

Two polymers are \emph{incompatible} if their union is connected through faces, and \emph{compatible} otherwise; incompatibility is a symmetric and reflexive relation. For $n\ge1$ and $(X_1,\dots,X_n)\in D_k^n$ let $\mathbf{G}(X_1,\dots,X_n)$ be the graph with the vertex set $\{1,\dots,n\}$ whose edges are the pairs $\{i,j\}$, $i\ne j$, for which $X_i$ and $X_j$ are incompatible. Put $\rho^T(X_1):=1$ and, for $n\ge2$,
\[
\rho^T(X_1,\dots,X_n):=\sum_G(-1)^{\#|G|},
\]
where $G$ ranges over the connected spanning subgraphs of $\mathbf{G}(X_1,\dots,X_n)$, that is, over the connected graphs with the vertex set $\{1,\dots,n\}$ all of whose edges are edges of $\mathbf{G}(X_1,\dots,X_n)$, and $\#|G|$ is the number of the edges of $G$. If $\mathbf{G}(X_1,\dots,X_n)$ is not connected, there is no such $G$, and $\rho^T(X_1,\dots,X_n)=0$.

\begin{theorem}\label{thm:polymer-activities-2}
Let $k\ge0$ and $l\in\mathrm{Cl}$, and let $c_X$, $\kappa_0$, $K_0$, $c_{\mathrm{rb}}$ be the numbers~\eqref{eq:polymer-activities-5}. Then the following assertions hold. The paragraphs below contain, in their order, the assertions~\eqref{it:polymers-t0}--\eqref{it:polymers-CE2} of Proposition~\ref{prop:polymers} and are designated by the marks of these assertions; the paragraph on monotonicity comprises two of them, its first inequality being~\eqref{it:polymers-t3} and its second~\eqref{it:polymers-t3p}.

\emph{Single cubes.} Every $M$-cube $\square$ of $T_k$ and every $X\in D_k$ satisfy
\[
d_M(\square)=0,\qquad d_M(X)\le|X|_M-1.
\]

\emph{Number of cubes.} Every $X\in D_k$ satisfies
\[
|X|_M\le c_X\cdot\left(d_M(X)+1\right).
\]

\emph{Sums over polymers.} For every $M$-cube $\square$ of $T_k$ and every union $Y$ of $M$-cubes of $T_k$,
\begin{align*}
\sum_{X\in D_k:\,X\supset\square}\exp\left(-\kappa_0\cdot d_M(X)\right)&\le K_0,\\
\sum_{X\in D_k:\,X\subset Y}\exp\left(-\kappa_0\cdot d_M(X)\right)&\le K_0\cdot|Y|_M;
\end{align*}
the second inequality is a consequence of the first.

\emph{Monotonicity.} If $X,Y'\in D_k$ and $X\subset Y'$, then
\begin{align*}
d_M(Y')&\le d_M(X)+|Y'\setminus X|_M,\\
d_M(X)&\le d_M(Y').
\end{align*}

\emph{Unions.} If $X_1,\dots,X_n\in D_k$ and the graph $\mathbf{G}(X_1,\dots,X_n)$ is connected, then $X_1\cup\dots\cup X_n$ is a polymer, and
\[
d_M\left(X_1\cup\dots\cup X_n\right)\le\sum_{i=1}^nd_M(X_i)+3(n-1).
\]

\emph{Reblocking.} Assume in addition that $\frac{l}{L}\in\mathrm{Cl}$. Level $k+1$ is formed on the torus $T_{k+1}=\mathbb{L}_{k+1}/\frac{l}{L}\mathbb{Z}^3$, and $D_{k+1}$ is the set of its polymers; for $Y\in D_{k+1}$, $LY$ denotes the image of $Y$ under the map $x\mapsto Lx$ of $T_{k+1}$ onto $T_k$. For every $X\in D_k$ there is exactly one $Y\in D_{k+1}$ with $\bar{X}^L=LY$, the set $\bar{X}^L$ being defined by~\eqref{eq:polymer-activities-7}, and this $Y$ satisfies
\[
L\cdot d_M(Y)\le d_M(X).
\]

\emph{Two elementary inequalities.} Every real number $d\ge0$ satisfies
\[
d+1\le\exp(d),\qquad d+1\le2\exp\left(\frac{d}{2}\right).
\]

\emph{Diameter.} Let $X\in D_k$. Any two sites $x$, $x'$ of $X$ satisfy
\[
d(x,x')\le M\cdot\left(d_M(X)+2\right).
\]
The same inequality holds for any two points $x$, $x'$ of the closed cubes of $X$, $d$ being the sup distance of $\mathbb{R}^3/l\mathbb{Z}^3$; and, with $|x-x'|$ in place of $d(x,x')$ and with $d_M$ formed with nets in $\mathbb{R}^3$, it holds for any two points $x$, $x'$ of the closed cubes of a finite union of $M$-cubes of $\mathbb{R}^3$ that is connected through faces. In particular a small polymer $X$, that is, one with $d_M(X)\le L$ (Section~\ref{sec:lattice-geometry}), is contained in the ball $K\left(x_0,M\cdot(L+2)\right)$ of Section~\ref{sec:balls-torus-lifts} for each of its sites $x_0$, and this ball does not wrap.

\emph{Reblocked sums.} In the situation of~\eqref{it:polymers-t5}, that is, if moreover $\frac{l}{L}\in\mathrm{Cl}$, every $Y\in D_{k+1}$ and every $\kappa'\ge\kappa_0$ satisfy
\[
\sum_{X\in D_k:\,\bar{X}^L=LY}\exp\left(-\kappa'\cdot d_M(X)\right)\le c_{\mathrm{rb}}\cdot L^3\cdot\exp\left(-L\cdot\left(\kappa'-\kappa_0-1\right)\cdot d_M(Y)\right).
\]

\emph{Incompatible polymers.} For every $Y'\in D_k$ there is a set $\mathcal{C}$ of $M$-cubes of $T_k$ with
\[
\#|\mathcal{C}|\le7\cdot|Y'|_M
\]
such that every polymer that is incompatible with $Y'$ contains a cube of $\mathcal{C}$.

\emph{Lifts.} Let $\mathrm{pr}\colon\mathbb{R}^3\to\mathbb{R}^3/l\mathbb{Z}^3$ be the projection; for finite unions of $M$-cubes of $\mathbb{R}^3$ that are connected through faces, $d_M$ is formed with nets in $\mathbb{R}^3$. Let $X\in D_k$ satisfy $2M\cdot\left(d_M(X)+2\right)<l$. Then $X=\mathrm{pr}(\tilde{X})$ for a finite union $\tilde{X}$ of $M$-cubes of $\mathbb{R}^3$ that is connected through faces and on which $\mathrm{pr}$ is injective; such an $\tilde{X}$ is unique up to translations by elements of $l\mathbb{Z}^3$; two cubes of $\tilde{X}$ are adjacent if and only if their images under $\mathrm{pr}$ are adjacent; and
\[
d_M(\tilde{X})=d_M(X).
\]
Conversely, if $\tilde{X}$ is a non-empty finite union of $M$-cubes of $\mathbb{R}^3$ that is connected through faces and satisfies $2M\cdot\left(d_M(\tilde{X})+2\right)<l$, then $\mathrm{pr}$ is injective on $\tilde{X}$, $\mathrm{pr}(\tilde{X})$ is a polymer, and $d_M\left(\mathrm{pr}(\tilde{X})\right)=d_M(\tilde{X})$. In particular, let $\square$ be an $M$-cube of $T_k$, and fix a lift of $\square$, that is, an $M$-cube of $\mathbb{R}^3$ whose image under $\mathrm{pr}$ is $\square$. Then the small polymers that contain $\square$ correspond bijectively, through $\mathrm{pr}$, to the finite unions $\tilde{X}$ of $M$-cubes of $\mathbb{R}^3$ that are connected through faces, contain this lift and satisfy $d_M(\tilde{X})\le L$; this set of unions $\tilde{X}$ does not depend on $l$.

\emph{Mayer expansion.} In the ring of formal power series in the commuting variables $z_X$, $X\in D_k$,
\[
\sum_{\{X_i\}}\prod_iz_{X_i}=\exp\left(\sum_{n\ge1}\frac{1}{n!}\sum_{(X_1,\dots,X_n)\in D_k^n}\rho^T(X_1,\dots,X_n)\cdot\prod_{i=1}^nz_{X_i}\right),
\]
the sum on the left being taken over all sets $\{X_i\}$ of pairwise compatible polymers, the empty set included.

\emph{Tree bound.} For every $n\ge1$ and every $(X_1,\dots,X_n)\in D_k^n$,
\[
\left|\rho^T(X_1,\dots,X_n)\right|\le\#\left|\left\{\text{spanning trees of }\mathbf{G}(X_1,\dots,X_n)\right\}\right|.
\]
\end{theorem}

\begin{proof}[Proof of Theorem~\ref{thm:polymer-activities-2}\eqref{it:polymers-t0}--\eqref{it:polymers-t9}]
\emph{Proof of~\eqref{it:polymers-t0}.} A single point of the closed cube $\square$ is a net that meets $\square$, and its length is $0$; hence $d_M(\square)=0$. Let $X\in D_k$ and $n:=|X|_M$. Since $X$ is connected through faces, its cubes can be numbered $\square_1,\dots,\square_n$ in such a way that every $\square_i$ with $i\ge2$ is adjacent to one of $\square_1,\dots,\square_{i-1}$. Begin with a point of the closed cube $\square_1$ and apply, for $i=2,\dots,n$ in turn, the construction in the proof of~\eqref{it:polymers-t3}, which adds a segment of length at most $M$ ending in the closed cube $\square_i$. The outcome is a net for $X$ of length at most $(n-1)\cdot M$; hence $d_M(X)\le n-1$.

\emph{Proof of~\eqref{it:polymers-t3}.} Put $m:=|Y'\setminus X|_M$. In the graph of the cubes of $Y'$ with the adjacency relation every cube has a finite distance from the set of the cubes of $X$, because $Y'$ is connected through faces and $X$ is not empty; a cube at distance $s\ge1$ is adjacent to a cube at distance $s-1$. Number the cubes of $Y'$ that are not cubes of $X$ as $\square_1,\dots,\square_m$ in the order of increasing distance. Then each $\square_i$ is adjacent to a cube $\square$ of $X\cup\square_1\cup\dots\cup\square_{i-1}$. Let $N$ be a net that meets the closed cube $\square$ at a point $p$. The segment that starts at $p$, is perpendicular to the common face of $\square$ and $\square_i$ and ends on this face has length at most $M$, and its end point lies in the closed cube $\square_i$; the union of $N$ and of this segment is a net that meets $\square_i$ and every closed cube met by $N$. Now let $\boldsymbol{\delta}>0$, start from a net for $X$ of length at most $M\cdot d_M(X)+\boldsymbol{\delta}$, and add such segments for $i=1,\dots,m$. The outcome is a net for $Y'$ of length at most $M\cdot d_M(X)+\boldsymbol{\delta}+m\cdot M$. As $\boldsymbol{\delta}>0$ is arbitrary, $d_M(Y')\le d_M(X)+m$.

\emph{Proof of~\eqref{it:polymers-t3p}.} Every cube of $X$ is a cube of $Y'$. Hence every net for $Y'$ is a net for $X$, and the infimum that defines $M\cdot d_M(X)$ is at most the infimum that defines $M\cdot d_M(Y')$.

\emph{Proof of~\eqref{it:polymers-t1}.} Let $N$ be a net for $X$ of length $\boldsymbol{\ell}$. Subdivide the segments of $N$ at their intersection points, keeping overlapping pieces only once; the length does not increase, and $N$ becomes a finite connected graph with straight edges whose total length is at most $\boldsymbol{\ell}$. With every edge doubled, this graph is a connected multigraph in which every vertex has even degree; by Euler's theorem~\eqref{eq:polymer-activities-cl4} it has a closed walk that traverses every edge of $N$ exactly twice. This walk is a closed curve with image $N$, and its length is twice the total length of the graph, hence at most $2\boldsymbol{\ell}$. Put
\[
n_a:=\max\left(1,\left\lceil\frac{2\boldsymbol{\ell}}{M}\right\rceil\right),
\]
and cut the closed curve into $n_a$ arcs of length at most $M$. Any two points of one arc have Euclidean distance at most $M$, hence sup distance at most $M$. Each arc lifts to a curve in $\mathbb{R}^3$. The projection of this curve onto a coordinate axis is an interval of length at most $M$, and an interval of length at most $M$ meets at most three of the closed intervals $\left[\left(j-\frac{1}{2}\right)M,\left(j+\frac{1}{2}\right)M\right]$, $j\in\mathbb{Z}$. So the lifted curve meets at most $27$ closed $M$-cubes of $\mathbb{R}^3$; as every closed $M$-cube of the torus that meets the arc is the image of a closed $M$-cube of $\mathbb{R}^3$ that meets the lifted curve, the arc meets at most $27$ closed $M$-cubes. Every closed cube of $X$ meets $N$, hence one of the arcs, and therefore
\[
|X|_M\le27\cdot n_a\le27\cdot\left(\frac{2\boldsymbol{\ell}}{M}+1\right).
\]
For nets for $X$ whose lengths $\boldsymbol{\ell}$ tend to $M\cdot d_M(X)$ this gives
\[
|X|_M\le27\cdot\left(2d_M(X)+1\right)\le54\cdot\left(d_M(X)+1\right),
\]
which is~\eqref{it:polymers-t1}, since $c_X=54$ by~\eqref{eq:polymer-activities-5}.

\emph{Proof of~\eqref{it:polymers-t2}.} Every $M$-cube of $T_k$ is adjacent to exactly six $M$-cubes, because $l>2M$; so the adjacency graph of the $M$-cubes has degree $6$. Let $X\in D_k$ contain $\square$, and let $n:=|X|_M$. The cubes of $X$, with the adjacency relation, form a connected graph; fix a spanning tree of it. By Euler's theorem~\eqref{eq:polymer-activities-cl4}, applied to this tree with every edge doubled, there is a closed walk that traverses every edge of the tree exactly twice, and it may be taken to start at $\square$. This walk has $2(n-1)$ steps and visits exactly the cubes of $X$. Choosing one such walk for every $X$ defines an injective map, since $X$ is the set of the cubes visited by its walk. A walk of $2(n-1)$ steps that starts at $\square$ has at most six choices at each step; hence at most $6^{2(n-1)}=36^{n-1}$ polymers $X\in D_k$ with $|X|_M=n$ contain $\square$. By~\eqref{it:polymers-t1}, each of them satisfies $d_M(X)\ge\frac{n}{c_X}-1$. As $\frac{\kappa_0}{c_X}=\log72$ and $K_0=\frac{\exp(\kappa_0)}{36}$ by~\eqref{eq:polymer-activities-5}, it follows that
\begin{align*}
\sum_{X\in D_k:\,X\supset\square}\exp\left(-\kappa_0\cdot d_M(X)\right)&\le\exp(\kappa_0)\cdot\sum_{n\ge1}36^{n-1}\cdot\exp\left(-\frac{\kappa_0\cdot n}{c_X}\right)\\
&=\exp(\kappa_0)\cdot\sum_{n\ge1}36^{n-1}\cdot72^{-n}=\frac{\exp(\kappa_0)}{36}=K_0.
\end{align*}
Every $X\in D_k$ with $X\subset Y$ contains a cube of $Y$. Summing the first inequality of~\eqref{it:polymers-t2} over the $|Y|_M$ cubes of $Y$ therefore gives the second.

\emph{Proof of~\eqref{it:polymers-t4}.} Fix a spanning tree of $\mathbf{G}(X_1,\dots,X_n)$; it has $n-1$ edges. Let $\boldsymbol{\delta}>0$, and for $i=1,\dots,n$ let $N_i$ be a net for $X_i$ of length at most $M\cdot d_M(X_i)+\boldsymbol{\delta}$. Let $\{i,j\}$ be an edge of the tree. Then $X_i$ and $X_j$ are incompatible, that is, $X_i\cup X_j$ is connected through faces; hence some cube $\square$ of $X_i$ and some cube $\square'$ of $X_j$ coincide or are adjacent. Let $p$ be a point at which $N_i$ meets the closed cube $\square$, and $p'$ a point at which $N_j$ meets the closed cube $\square'$. Both points lie in a box with sides $M$, $M$, $2M$ that contains the two closed cubes and, as $2M<l$, does not wrap; so their Euclidean distance satisfies
\[
|p-p'|_2\le\sqrt{6}\cdot M<3M.
\]
Join $p$ to $p'$ by a segment of this length. The union of $N_1,\dots,N_n$ and of the $n-1$ segments obtained in this way is connected, because the tree is connected; so it is a net, and it meets every closed cube of $X_1\cup\dots\cup X_n$. Its length is at most $\sum_{i=1}^n\left(M\cdot d_M(X_i)+\boldsymbol{\delta}\right)+3M\cdot(n-1)$, and letting $\boldsymbol{\delta}\to0$ gives~\eqref{it:polymers-t4}. The union $X_1\cup\dots\cup X_n$ is a polymer: it is not empty, and it is connected through faces, since every $X_i$ is and so is $X_i\cup X_j$ for every edge $\{i,j\}$ of the tree.

\emph{Proof of~\eqref{it:polymers-t5}.} The $LM$-cubes of $T_k$ are the cubes of side $L\cdot M$ centred at the points of $LM\mathbb{Z}^3/l\mathbb{Z}^3$; this makes sense because $l$ is a multiple of $L\cdot M$. Since $L$ is odd, every $M$-cube lies in exactly one $LM$-cube, and every $LM$-cube is the union of the $L^3$ $M$-cubes that it contains. Two adjacent $M$-cubes lie in one $LM$-cube or in two $LM$-cubes with a common face. As $X$ is connected through faces, it follows that $\bar{X}^L$ is a union of $LM$-cubes that is connected through faces. The map $x\mapsto\frac{x}{L}$ carries $T_k$ onto $T_{k+1}$ and the $LM$-cubes of $T_k$ onto the $M$-cubes of $T_{k+1}$. Hence $\bar{X}^L=LY$ for exactly one $Y\in D_{k+1}$, namely the image of $\bar{X}^L$ under this map. Let $N$ be a net for $X$. Every closed $LM$-cube of $\bar{X}^L$ contains a cube of $X$, so $N$ meets it. The map $x\mapsto\frac{x}{L}$ of $\mathbb{R}^3/l\mathbb{Z}^3$ onto $\mathbb{R}^3/\frac{l}{L}\mathbb{Z}^3$ divides all lengths by $L$; so the image of $N$ is a net for $Y$ whose length is the length of $N$ divided by $L$. Consequently
\[
M\cdot d_M(Y)\le\frac{M\cdot d_M(X)}{L}.
\]

\emph{Proof of~\eqref{it:polymers-t6}.} As $\exp(t)\ge1+t$ for every real $t$, one has $\exp(d)\ge1+d$ and $2\exp\left(\frac{d}{2}\right)\ge2+d\ge1+d$.

\emph{Proof of~\eqref{it:polymers-t7}.} Let $p$, $q$ be sites of $X$, contained in the cubes $\square_p$, $\square_q$ of $X$. Let $\boldsymbol{\delta}>0$, and let $N$ be a net for $X$ of length at most $M\cdot d_M(X)+\boldsymbol{\delta}$; it meets the closed cube $\square_p$ at a point $p'$ and the closed cube $\square_q$ at a point $q'$. As $N$ is connected, $p'$ and $q'$ are joined in $N$ by a curve of length at most $M\cdot d_M(X)+\boldsymbol{\delta}$. Hence
\[
d(p,q)\le d(p,p')+d(p',q')+d(q',q)\le M+\left(M\cdot d_M(X)+\boldsymbol{\delta}\right)+M,
\]
and letting $\boldsymbol{\delta}\to0$ gives~\eqref{it:polymers-t7}. The argument rests on three facts only: two points of one closed $M$-cube have sup distance at most $M$; two points of a net are joined in the net by a curve whose length does not exceed that of the net; the sup distance of two points is at most the length of any curve joining them. None of them involves lattice sites or the torus. Therefore the same argument applies to arbitrary points $p$, $q$ of the closed cubes of $X$, and, with nets in $\mathbb{R}^3$ and with the sup norm, to a finite union of $M$-cubes of $\mathbb{R}^3$ that is connected through faces. If $X$ is small, then $d_M(X)\le L$, so every site $x$ of $X$ satisfies $d(x,x_0)\le M\cdot(L+2)$ for every site $x_0$ of $X$; moreover, since $\eta=L^{-k}\le1\le M$, one has $2M\cdot(L+2)+\eta\le2L\cdot M+5M<13L\cdot M\le l$, so the ball $K\left(x_0,M\cdot(L+2)\right)$ is defined and does not wrap (Section~\ref{sec:balls-torus-lifts}).

\emph{Proof of~\eqref{it:polymers-t8}.} Let $Y\in D_{k+1}$ and $\kappa'\ge\kappa_0$, and let $X\in D_k$ satisfy $\bar{X}^L=LY$. By~\eqref{it:polymers-t5}, $L\cdot d_M(Y)\le d_M(X)$; as $\kappa'-\kappa_0\ge0$, this gives
\[
\exp\left(-\kappa'\cdot d_M(X)\right)\le\exp\left(-\left(\kappa'-\kappa_0\right)\cdot L\cdot d_M(Y)\right)\cdot\exp\left(-\kappa_0\cdot d_M(X)\right).
\]
Every such $X$ is contained in $\bar{X}^L=LY$, which is a union of $L^3\cdot|Y|_M$ $M$-cubes of $T_k$. The second inequality of~\eqref{it:polymers-t2}, then~\eqref{it:polymers-t1} at level $k+1$ (with the side $\frac{l}{L}\in\mathrm{Cl}$), then~\eqref{it:polymers-t6}, and finally $L\ge1$ and $d_M(Y)\ge0$ give
\begin{align*}
\sum_{X\in D_k:\,\bar{X}^L=LY}\exp\left(-\kappa_0\cdot d_M(X)\right)&\le K_0\cdot L^3\cdot|Y|_M\le c_X\cdot K_0\cdot L^3\cdot\left(d_M(Y)+1\right)\\
&\le c_X\cdot K_0\cdot L^3\cdot\exp\left(d_M(Y)\right)\le c_X\cdot K_0\cdot L^3\cdot\exp\left(L\cdot d_M(Y)\right).
\end{align*}
With $c_{\mathrm{rb}}=c_X\cdot K_0$ from~\eqref{eq:polymer-activities-5}, the last two displays yield
\[
\sum_{X\in D_k:\,\bar{X}^L=LY}\exp\left(-\kappa'\cdot d_M(X)\right)\le c_{\mathrm{rb}}\cdot L^3\cdot\exp\left(-\left(\kappa'-\kappa_0\right)\cdot L\cdot d_M(Y)+L\cdot d_M(Y)\right),
\]
which is~\eqref{it:polymers-t8}.

\emph{Proof of~\eqref{it:polymers-t9}.} Let $\mathcal{C}$ be the set of the $M$-cubes of $T_k$ that belong to $Y'$ or are adjacent to a cube of $Y'$; it has at most $|Y'|_M+6\cdot|Y'|_M=7\cdot|Y'|_M$ elements. Let $X$ be a polymer that is incompatible with $Y'$, so that $X\cup Y'$ is connected through faces. If no cube of $X$ belonged to $\mathcal{C}$, then no cube of $X$ would coincide with or be adjacent to a cube of $Y'$, and $X\cup Y'$ would not be connected through faces. Hence $X$ contains a cube of $Y'$ or a cube adjacent to a cube of $Y'$, that is, a cube of $\mathcal{C}$.
\end{proof}

\begin{proof}[Proof of Theorem~\ref{thm:polymer-activities-2}\eqref{it:polymers-t10}]
In this proof a polymer is regarded as the union of its closed $M$-cubes in the continuum torus $\mathbb{R}^3/l\mathbb{Z}^3$, and the $M$-cubes of $\mathbb{R}^3$ are the closed cubes of side $M$ centred at the points of $M\mathbb{Z}^3$. On $\mathbb{R}^3$, $|\cdot|$ is the sup norm; on the continuum torus, $d$ denotes the sup distance,
\[
d\left(\mathrm{pr}(x),\mathrm{pr}(y)\right)=\min_{v\in l\mathbb{Z}^3}|x-y-v|\qquad\left(x,y\in\mathbb{R}^3\right).
\]
As $l\in\mathrm{Cl}$, the side $l$ is a multiple of $L\cdot M$, hence of $M$, and $l\ge 13L\cdot M$. Consequently the translations by the elements of $l\mathbb{Z}^3$ permute the $M$-cubes of $\mathbb{R}^3$, $\mathrm{pr}$ maps every $M$-cube of $\mathbb{R}^3$ onto an $M$-cube of the torus, and every $M$-cube of the torus is of this form; the faces of the image of an $M$-cube of $\mathbb{R}^3$ are the images of its faces. In $\mathbb{R}^3$ as in the torus, a closed $M$-cube that is contained in a finite union of closed $M$-cubes of the same space is one of them, since its interior meets the interior of one of them and distinct $M$-cubes have disjoint interiors.

\emph{A ball on which the projection is injective.} Let $X\in D_k$ satisfy $2M\cdot\left(d_M(X)+2\right)<l$, and put
\[
\mathrm{Dm}:=M\cdot\left(d_M(X)+2\right),\qquad\text{so that}\qquad 2\mathrm{Dm}<l.
\]
Let $x_0$ be a point of a closed cube $\square_0$ of $X$, let $\tilde{x}_0\in\mathbb{R}^3$ be a point with $\mathrm{pr}(\tilde{x}_0)=x_0$, and let
\[
\tilde{B}:=\left\{x\in\mathbb{R}^3:|x-\tilde{x}_0|\le\mathrm{Dm}\right\},\qquad\bar{B}:=\left\{p\in\mathbb{R}^3/l\mathbb{Z}^3:d(p,x_0)\le\mathrm{Dm}\right\}
\]
be the closed sup-balls of radius $\mathrm{Dm}$ around $\tilde{x}_0$ in $\mathbb{R}^3$ and around $x_0$ in the torus. Any two points of $\tilde{B}$ differ by at most $2\mathrm{Dm}<l$ in each coordinate, whereas a non-zero element of $l\mathbb{Z}^3$ has a coordinate of modulus at least $l$. Hence $\mathrm{pr}$ is injective on $\tilde{B}$. Moreover $\mathrm{pr}(\tilde{B})=\bar{B}$, because $d\left(\mathrm{pr}(x),x_0\right)\le|x-\tilde{x}_0|$ for $x\in\mathbb{R}^3$ and every point of $\bar{B}$ has a representative at sup distance at most $\mathrm{Dm}$ from $\tilde{x}_0$. Since $\tilde{B}$ is compact, the restriction of $\mathrm{pr}$ to $\tilde{B}$ is a homeomorphism onto $\bar{B}$; it is locally isometric, because $d\left(\mathrm{pr}(x),\mathrm{pr}(y)\right)=|x-y|$ whenever $|x-y|\le\frac{l}{2}$. Further,
\[
\mathrm{pr}^{-1}(\bar{B})=\bigcup_{v\in l\mathbb{Z}^3}\left(\tilde{B}+v\right),
\]
and the closed sets $\tilde{B}+v$, $v\in l\mathbb{Z}^3$, are pairwise disjoint: if $z$ lies in $\tilde{B}+v$ and in $\tilde{B}+v'$, then $z-v$ and $z-v'$ are points of $\tilde{B}$, and their difference $v'-v\in l\mathbb{Z}^3$ has all its coordinates of modulus less than $l$, so that $v=v'$.

Let $\varphi\colon\bar{B}\to\tilde{B}$ be the inverse of the restriction of $\mathrm{pr}$ to $\tilde{B}$. A bounded subset of $\mathbb{R}^3$ meets only finitely many of the sets $\tilde{B}+v$. Therefore a compact connected subset $S$ of $\mathrm{pr}^{-1}(\bar{B})$, which is the union of the finitely many pairwise disjoint closed sets $S\cap\left(\tilde{B}+v\right)$, lies in one of the sets $\tilde{B}+v$. If $S$ is a closed segment or an $M$-cube of $\mathbb{R}^3$ with $\mathrm{pr}(S)\subset\bar{B}$, and $S\subset\tilde{B}+v$, then $S-v\subset\tilde{B}$ and $\mathrm{pr}(S-v)=\mathrm{pr}(S)$; so $\varphi\left(\mathrm{pr}(S)\right)=S-v$ is a segment of the same length as $S$, respectively an $M$-cube of $\mathbb{R}^3$. Thus $\varphi$ maps the closed $M$-cubes of the torus that lie in $\bar{B}$ onto $M$-cubes of $\mathbb{R}^3$, and the segments that lie in $\bar{B}$ onto segments of the same lengths. It also maps common faces to common faces. Indeed, let $c,c'\subset\bar{B}$ be adjacent $M$-cubes of the torus with the common face $f$. As $\mathrm{pr}$ maps the $M$-cube $\varphi(c)$ of $\mathbb{R}^3$ onto $c$, the face $f$ of $c$ is the image $\mathrm{pr}(F)$ of a face $F$ of $\varphi(c)$; since $F\subset\varphi(c)\subset\tilde{B}$ and $\varphi$ is the inverse of the restriction of $\mathrm{pr}$ to $\tilde{B}$, one has $\varphi(f)=F$. In the same way $\varphi(f)$ is a face of $\varphi(c')$. The cubes $\varphi(c)$ and $\varphi(c')$ are different, because their images $c$ and $c'$ are different; hence they are adjacent, and $\varphi(f)$ is their common face.

\emph{Existence and adjacency.} By~\eqref{it:polymers-t7}, applied to the points of the closed cubes of $X$, every point of a closed cube of $X$ has sup distance at most $\mathrm{Dm}$ from $x_0$. So the closed cubes of $X$ lie in $\bar{B}$. Let
\[
\tilde{X}:=\varphi(X)=\tilde{B}\cap\mathrm{pr}^{-1}(X)
\]
be the preimage of $X$ in $\tilde{B}$. It is the union of the $M$-cubes $\varphi(c)$ of $\mathbb{R}^3$, $c$ a cube of $X$; so $\tilde{X}$ is a finite union of $M$-cubes of $\mathbb{R}^3$, $\mathrm{pr}(\tilde{X})=X$, and $\mathrm{pr}$ is injective on $\tilde{X}\subset\tilde{B}$. The projection maps common faces of cubes of $\tilde{X}$ to common faces of their images: the images of two adjacent cubes of $\tilde{X}$ are distinct, and the image of their common face is a face of both. Combined with the corresponding property of $\varphi$, this shows that adjacency of two cubes of $\tilde{X}$ is equivalent to adjacency of their images under $\mathrm{pr}$. As $X$ is connected through faces, so is $\tilde{X}$.

\emph{Uniqueness.} Let $\tilde{X}'$ have the properties that~\eqref{it:polymers-t10} requires of $\tilde{X}$: it is a finite union of $M$-cubes of $\mathbb{R}^3$, it is connected through faces, the restriction of $\mathrm{pr}$ to it is injective, and $\mathrm{pr}(\tilde{X}')=X$. Each cube of $\tilde{X}'$ is mapped onto a cube of $X$, and each cube of $X$ is the image of a cube of $\tilde{X}'$. In particular some cube of $\tilde{X}'$ lies above $\square_0$, and it is the translate of the cube $\varphi(\square_0)$ of $\tilde{X}$ by an element of $l\mathbb{Z}^3$. After the translation of $\tilde{X}'$ by the opposite element, which preserves the hypotheses on $\tilde{X}'$, the cube $\varphi(\square_0)$ belongs to $\tilde{X}'$. Now let $\tilde{c}$ be a cube of $\tilde{X}'$ that belongs to $\tilde{X}$, and let $\tilde{c}'$ be a cube of $\tilde{X}'$ adjacent to $\tilde{c}$. Then $\mathrm{pr}(\tilde{c}')$ is a cube of $X$ adjacent to $\mathrm{pr}(\tilde{c})$, because $\mathrm{pr}$ is injective on $\tilde{X}'$ and maps the common face of $\tilde{c}$ and $\tilde{c}'$ to a face of both images. By the adjacency statement for $\tilde{X}$, the cube $\varphi\left(\mathrm{pr}(\tilde{c}')\right)$ of $\tilde{X}$ above $\mathrm{pr}(\tilde{c}')$ is adjacent to $\tilde{c}=\varphi\left(\mathrm{pr}(\tilde{c})\right)$. Thus $\tilde{c}'$ and $\varphi\left(\mathrm{pr}(\tilde{c}')\right)$ are neighbours of $\tilde{c}$ with the same image under $\mathrm{pr}$. Two different neighbours of $\tilde{c}$ differ by a non-zero vector of sup norm at most $2M$, and $2M<l$; this vector does not lie in $l\mathbb{Z}^3$, so their images are different. Hence $\tilde{c}'=\varphi\left(\mathrm{pr}(\tilde{c}')\right)$, and $\tilde{c}'$ belongs to $\tilde{X}$. Since $\tilde{X}'$ is connected through faces, each of its cubes is joined to $\varphi(\square_0)$ by a chain of cubes of $\tilde{X}'$ in which consecutive cubes are adjacent; the preceding step, applied along such a chain, gives $\tilde{X}'\subset\tilde{X}$. Finally $\tilde{X}'=\tilde{X}$: for a cube $c$ of $\tilde{X}$, the cube $\mathrm{pr}(c)$ of $X$ is the image of a cube $c'$ of $\tilde{X}'\subset\tilde{X}$, and $c=c'$ because $\mathrm{pr}$ is injective on $\tilde{X}$. Undoing the translation, one finds that the given union $\tilde{X}'$ is the image of $\tilde{X}$ under a translation by an element of $l\mathbb{Z}^3$, which proves the uniqueness assertion.

\emph{The numbers $d_M$.} Let $N$ be a net in $\mathbb{R}^3$, in the sense of convention~\eqref{eq:polymer-activities-6} with $\mathbb{R}^3$ in place of the torus, that meets every closed cube of $\tilde{X}$. Then $\mathrm{pr}(N)$ is a net in the torus of the same length, since it is connected and $\mathrm{pr}$ maps the segments of $N$ onto segments of the same lengths; and it meets every closed cube of $X$, each of which is the image of a cube of $\tilde{X}$. Hence $M\cdot d_M(X)$ is at most the length of $N$, and therefore $d_M(X)\le d_M(\tilde{X})$. Conversely, let $N$ be a net in the torus that meets every closed cube of $X$ and has length at most $M\cdot d_M(X)+M$, and let $p_0$ be a point of $N$ in $\square_0$. Since $N$ is connected, each point $p$ of $N$ is joined to $p_0$ by a curve in $N$ whose length is at most the length of $N$. The sup distance of the end points of a curve is at most its length, and $p_0$ and $x_0$ lie in the closed cube $\square_0$ of side $M$; hence
\[
d(p,x_0)\le d(p,p_0)+d(p_0,x_0)\le M\cdot d_M(X)+M+M=\mathrm{Dm} .
\]
So $N\subset\bar{B}$. The set $\varphi(N)$ is connected, and $\varphi$ maps the segments of $N$ onto segments of the same lengths; so $\varphi(N)$ is a net in $\mathbb{R}^3$ of the same length as $N$, and it meets every closed cube $\varphi(c)$ of $\tilde{X}$, $c$ a cube of $X$. Hence $M\cdot d_M(\tilde{X})$ is at most the length of $N$. The infimum of the lengths of the nets $N$ of this kind is $M\cdot d_M(X)$; therefore $d_M(\tilde{X})\le d_M(X)$, and $d_M(\tilde{X})=d_M(X)$.

\emph{The converse.} Let $\tilde{X}$ be a non-empty finite union of $M$-cubes of $\mathbb{R}^3$ that is connected through faces and satisfies $2M\cdot\left(d_M(\tilde{X})+2\right)<l$. By~\eqref{it:polymers-t7}, in its form for such unions in $\mathbb{R}^3$, two points of $\tilde{X}$ have sup distance at most $M\cdot\left(d_M(\tilde{X})+2\right)<\frac{l}{2}$. Hence $\mathrm{pr}$ is injective on $\tilde{X}$, since two distinct points of $\tilde{X}$ with the same image would differ by a non-zero element of $l\mathbb{Z}^3$ of sup norm less than $l$; and, as above, $\mathrm{pr}$ maps common faces of cubes of $\tilde{X}$ to common faces of their images. So $X:=\mathrm{pr}(\tilde{X})$ is a non-empty union of $M$-cubes of the torus that is connected through faces, that is, a polymer, and the argument with $\mathrm{pr}(N)$ above gives $d_M(X)\le d_M(\tilde{X})$. In particular $2M\cdot\left(d_M(X)+2\right)<l$, so the first part applies to $X$; by the uniqueness, $\tilde{X}$ is a translate, by an element of $l\mathbb{Z}^3$, of the union constructed there for $X$. The number $d_M$ does not change under translations of $\mathbb{R}^3$; hence $d_M\left(\mathrm{pr}(\tilde{X})\right)=d_M(X)=d_M(\tilde{X})$.

\emph{Small polymers.} Let $\square$ be an $M$-cube of the torus, and fix an $M$-cube $\tilde{\square}$ of $\mathbb{R}^3$ with $\mathrm{pr}(\tilde{\square})=\square$. A small polymer $X$ satisfies $d_M(X)\le L$ (Section~\ref{sec:lattice-geometry}), and
\[
2M\cdot(L+2)\le 6L\cdot M<13L\cdot M\le l;
\]
so the first part applies to every small polymer. Let $X$ be a small polymer that contains $\square$, and let $\tilde{X}$ be a union given by the first part. As $\square\subset X=\mathrm{pr}(\tilde{X})$ and $\mathrm{pr}$ is injective on $\tilde{X}$, exactly one cube of $\tilde{X}$ lies above $\square$; it is $\tilde{\square}+v_0$ for some $v_0\in l\mathbb{Z}^3$. For $v\in l\mathbb{Z}^3$ the translate $\tilde{X}+v$ contains $\tilde{\square}$ if and only if $\tilde{\square}-v$ is a cube of $\tilde{X}$ above $\square$, that is, if and only if $v=-v_0$. So exactly one translate of $\tilde{X}$ by an element of $l\mathbb{Z}^3$ contains $\tilde{\square}$; it is connected through faces, $\mathrm{pr}$ is injective on it, and its number $d_M$ is $d_M(X)\le L$. Assign this translate to $X$. Conversely, let $\tilde{X}$ be a finite union of $M$-cubes of $\mathbb{R}^3$, connected through faces, with $d_M(\tilde{X})\le L$ and $\tilde{\square}\subset\tilde{X}$. Then $\tilde{X}\ne\emptyset$ and $2M\cdot\left(d_M(\tilde{X})+2\right)<l$; by the converse, $\mathrm{pr}(\tilde{X})$ is a polymer that contains $\square$, with $d_M\left(\mathrm{pr}(\tilde{X})\right)=d_M(\tilde{X})\le L$; so it is small. Assign $\mathrm{pr}(\tilde{X})$ to $\tilde{X}$. The image under $\mathrm{pr}$ of the translate assigned to $X$ is $X$. In the other direction, $\tilde{X}$ is, by the uniqueness, a translate of the union given by the first part for $\mathrm{pr}(\tilde{X})$, and it contains $\tilde{\square}$; so it is the translate assigned to $\mathrm{pr}(\tilde{X})$. The two assignments are therefore inverse to each other. The set of the unions $\tilde{X}$ that occur here is defined in $\mathbb{R}^3$ by $\tilde{\square}$ and $L$ alone; it does not depend on $l$.
\end{proof}

\begin{proof}[Proof of Theorem~\ref{thm:polymer-activities-2}\eqref{it:polymers-CE1}]
For polymers $X$, $X'$ put $\mathfrak{i}(X,X'):=1$ if $X$ and $X'$ are incompatible, and $\mathfrak{i}(X,X'):=0$ otherwise. Since incompatibility is reflexive, $\mathfrak{i}(X,X)=1$. Hence, for $n\ge 0$ and $(X_1,\dots,X_n)\in D_k^n$, the product
\[
\prod_{1\le i<j\le n}\left(1-\mathfrak{i}(X_i,X_j)\right)
\]
equals $1$ if $X_1,\dots,X_n$ are distinct and pairwise compatible, and $0$ otherwise; and every set of $n$ distinct pairwise compatible polymers is the set of the entries of exactly $n!$ tuples. The left side of~\eqref{it:polymers-CE1} is therefore
\[
\sum_{n\ge 0}\frac{1}{n!}\sum_{(X_1,\dots,X_n)\in D_k^n}\ \prod_{1\le i<j\le n}\left(1-\mathfrak{i}(X_i,X_j)\right)\cdot\prod_{i=1}^nz_{X_i},
\]
the term with $n=0$ being $1$.

Graphs on a finite vertex set are identified with their sets of edges. Expanding the product over the pairs gives
\[
\prod_{1\le i<j\le n}\left(1-\mathfrak{i}(X_i,X_j)\right)=\sum_G\ \prod_{\{i,j\}\in G}\left(-\mathfrak{i}(X_i,X_j)\right),
\]
where $G$ ranges over all graphs on $\{1,\dots,n\}$. A graph $G$ on $\{1,\dots,n\}$ determines the partition $\{B_1,\dots,B_r\}$ of $\{1,\dots,n\}$ into the vertex sets of its connected components, and it is the union of connected graphs on $B_1,\dots,B_r$; conversely, for every partition $\{B_1,\dots,B_r\}$ of $\{1,\dots,n\}$ into non-empty blocks and all connected graphs $G_a$ on $B_a$, $1\le a\le r$, the union of the $G_a$ is a graph whose components have the vertex sets $B_1,\dots,B_r$. For a non-empty set $B\subset\{1,\dots,n\}$ let $(X_i)_{i\in B}$ denote the tuple of the $X_i$, $i\in B$, in increasing order of $i$. Then
\[
\sum_G\ \prod_{\{i,j\}\in G}\left(-\mathfrak{i}(X_i,X_j)\right)=\rho^T\left((X_i)_{i\in B}\right),
\]
the sum on the left being over the connected graphs $G$ on $B$. Indeed, a term on the left vanishes unless every edge of $G$ joins two incompatible polymers, that is, unless $G$, transported to $\{1,\dots,\#|B|\}$ by the increasing bijection, is a connected spanning subgraph of $\mathbf{G}\left((X_i)_{i\in B}\right)$; and then the term equals $(-1)^{\#|G|}$. (For $\#|B|=1$ the only graph on $B$ is the one without edges, and both sides equal $1$.) Grouping the graphs $G$ according to the vertex sets of their components one obtains
\[
\prod_{1\le i<j\le n}\left(1-\mathfrak{i}(X_i,X_j)\right)=\sum_{\{B_1,\dots,B_r\}}\ \prod_{a=1}^r\rho^T\left((X_i)_{i\in B_a}\right),
\]
the sum over all partitions of $\{1,\dots,n\}$ into non-empty blocks.

For $m\ge 1$ let
\[
f_m:=\sum_{(X_1,\dots,X_m)\in D_k^m}\rho^T(X_1,\dots,X_m)\cdot\prod_{i=1}^mz_{X_i};
\]
since $D_k$ is finite, $f_m$ is a homogeneous polynomial of degree $m$. For a fixed partition $\{B_1,\dots,B_r\}$ of $\{1,\dots,n\}$ the sum over the tuples factorises over the blocks:
\[
\sum_{(X_1,\dots,X_n)\in D_k^n}\ \prod_{a=1}^r\left(\rho^T\left((X_i)_{i\in B_a}\right)\cdot\prod_{i\in B_a}z_{X_i}\right)=\prod_{a=1}^rf_{\#|B_a|}.
\]
Each partition into $r$ blocks arises from exactly $r!$ sequences $(B_1,\dots,B_r)$ of pairwise disjoint non-empty sets with union $\{1,\dots,n\}$, and the number of such sequences with $\#|B_a|=m_a$ for $1\le a\le r$ is $\frac{n!}{m_1!\cdots m_r!}$. Hence the left side of~\eqref{it:polymers-CE1} equals
\begin{align*}
\sum_{n\ge 0}\frac{1}{n!}\sum_{r\ge 0}\frac{1}{r!}\sum_{\substack{m_1,\dots,m_r\ge 1\\ m_1+\dots+m_r=n}}\frac{n!}{m_1!\cdots m_r!}\prod_{a=1}^rf_{m_a}
&=\sum_{r\ge 0}\frac{1}{r!}\sum_{m_1,\dots,m_r\ge 1}\ \prod_{a=1}^r\frac{f_{m_a}}{m_a!}\\
&=\sum_{r\ge 0}\frac{1}{r!}\left(\sum_{m\ge 1}\frac{f_m}{m!}\right)^r=\exp\left(\sum_{m\ge 1}\frac{f_m}{m!}\right),
\end{align*}
and the last expression is the right side of~\eqref{it:polymers-CE1}. Only finitely many terms of these sums contribute to the monomials of a given degree $n$, since such terms have $r\le n$ and $m_1,\dots,m_r\le n$.
\end{proof}

\begin{proof}[Proof of Theorem~\ref{thm:polymer-activities-2}\eqref{it:polymers-CE2}]
For $n=1$ both sides equal $1$. Let $n\ge 2$ and $\mathbf{G}:=\mathbf{G}(X_1,\dots,X_n)$. If $\mathbf{G}$ is not connected, then $\rho^T(X_1,\dots,X_n)=0$ by definition and $\mathbf{G}$ has no spanning tree, so that both sides vanish. Let $\mathbf{G}$ be connected. Graphs on $\{1,\dots,n\}$ are identified with their sets of edges; thus a spanning subgraph of $\mathbf{G}$ is a subset of $\mathbf{G}$, and $\#|G|$ is the number of the edges of $G$.

\emph{The tree of a connected subgraph.} Let $G$ be a connected spanning subgraph of $\mathbf{G}$. For $i\in\{1,\dots,n\}$ let $d_G(i)$ be the distance from $1$ to $i$ in $G$, that is, the least number of edges of a path in $G$ from $1$ to $i$. For $i\ne 1$ the set of the neighbours $j$ of $i$ in $G$ with $d_G(j)=d_G(i)-1$ is not empty, as it contains the vertex that precedes $i$ on a shortest path from $1$ to $i$; let $j_G(i)$ be its least element. Let $T(G)$ be the graph with the edges $\{i,j_G(i)\}$, $i\ne 1$. Then $T(G)\subset G$, and
\[
d_{T(G)}(i)\le d_G(i)\qquad(1\le i\le n)
\]
by induction on $d_G(i)$: if $d_G(i)\ge 1$, then $d_{T(G)}(i)\le d_{T(G)}\left(j_G(i)\right)+1\le d_G(i)$. In particular $T(G)$ is connected. Its $n-1$ edges $\{i,j_G(i)\}$, $i\ne 1$, are distinct, since $i$ is the end point with the larger value of $d_G$. A connected graph on $n$ vertices with $n-1$ edges is a tree; so $T(G)$ is a spanning tree of $\mathbf{G}$ contained in $G$. The reverse inequality $d_{T(G)}\ge d_G$ holds because $T(G)$ is contained in $G$; thus $d_{T(G)}=d_G$.

\emph{The graphs $R(\boldsymbol{\tau})$.} Let $\boldsymbol{\tau}$ be a spanning tree of $\mathbf{G}$. For $i\ne 1$ let $\operatorname{pa}(i)$ be the neighbour of $i$ in $\boldsymbol{\tau}$ with $d_{\boldsymbol{\tau}}\left(\operatorname{pa}(i)\right)=d_{\boldsymbol{\tau}}(i)-1$. It is the vertex that follows $i$ on the path from $i$ to $1$ in $\boldsymbol{\tau}$, and it is unique, since two such neighbours would give two different paths from $i$ to $1$ in the tree $\boldsymbol{\tau}$. Let $R(\boldsymbol{\tau})$ be the graph whose edges are the edges of $\boldsymbol{\tau}$, the edges $\{i,j\}$ of $\mathbf{G}$ with $d_{\boldsymbol{\tau}}(i)=d_{\boldsymbol{\tau}}(j)$, and the edges $\{i,j\}$ of $\mathbf{G}$ with $d_{\boldsymbol{\tau}}(j)=d_{\boldsymbol{\tau}}(i)-1$ and $j>\operatorname{pa}(i)$. Then $\boldsymbol{\tau}\subset R(\boldsymbol{\tau})\subset\mathbf{G}$, and every edge $\{i,j\}$ of $R(\boldsymbol{\tau})$ satisfies $\left|d_{\boldsymbol{\tau}}(i)-d_{\boldsymbol{\tau}}(j)\right|\le 1$ (for the edges of $\boldsymbol{\tau}$ by the triangle inequality).

\emph{The fibres of $T$.} For every connected spanning subgraph $G$ of $\mathbf{G}$ and every spanning tree $\boldsymbol{\tau}$ of $\mathbf{G}$,
\[
T(G)=\boldsymbol{\tau}\qquad\text{if and only if}\qquad\boldsymbol{\tau}\subset G\subset R(\boldsymbol{\tau}).
\]
Let $T(G)=\boldsymbol{\tau}$. Then $\boldsymbol{\tau}\subset G$ and $d_G=d_{\boldsymbol{\tau}}$; moreover $j_G(i)=\operatorname{pa}(i)$ for $i\ne 1$, since $j_G(i)$ is a neighbour of $i$ in $\boldsymbol{\tau}$ with $d_{\boldsymbol{\tau}}\left(j_G(i)\right)=d_{\boldsymbol{\tau}}(i)-1$. Let $\{i,j\}$ be an edge of $G$; then $\left|d_G(i)-d_G(j)\right|\le 1$. If $d_{\boldsymbol{\tau}}(i)=d_{\boldsymbol{\tau}}(j)$, the edge lies in $R(\boldsymbol{\tau})$. Otherwise its end points can be named so that $d_{\boldsymbol{\tau}}(j)=d_{\boldsymbol{\tau}}(i)-1$; then $j\ge j_G(i)=\operatorname{pa}(i)$ by the definition of $j_G(i)$, and the edge lies in $\boldsymbol{\tau}$ if $j=\operatorname{pa}(i)$, and in $R(\boldsymbol{\tau})$ if $j>\operatorname{pa}(i)$. So $G\subset R(\boldsymbol{\tau})$. Conversely, let $\boldsymbol{\tau}\subset G\subset R(\boldsymbol{\tau})$. Along a path in $G$ from $1$ to $i$ the function $d_{\boldsymbol{\tau}}$ changes by at most $1$ at each edge; so $d_{\boldsymbol{\tau}}(i)\le d_G(i)$. As $\boldsymbol{\tau}\subset G$, also $d_G(i)\le d_{\boldsymbol{\tau}}(i)$; hence $d_G=d_{\boldsymbol{\tau}}$. Let $i\ne 1$, and let $j$ be a neighbour of $i$ in $G$ with $d_G(j)=d_G(i)-1$. The edge $\{i,j\}$ lies in $R(\boldsymbol{\tau})$, and $d_{\boldsymbol{\tau}}(j)=d_{\boldsymbol{\tau}}(i)-1$; so either it is an edge of $\boldsymbol{\tau}$, and then $j=\operatorname{pa}(i)$, or $j>\operatorname{pa}(i)$. Since $\operatorname{pa}(i)$ is itself a neighbour of $i$ in $G$ with $d_G\left(\operatorname{pa}(i)\right)=d_G(i)-1$, it follows that $j_G(i)=\operatorname{pa}(i)$. Thus all the edges $\{i,\operatorname{pa}(i)\}$ of $T(G)$ lie in $\boldsymbol{\tau}$; as $T(G)$ and $\boldsymbol{\tau}$ are trees on $\{1,\dots,n\}$ with $n-1$ edges each, $T(G)=\boldsymbol{\tau}$.

\emph{Conclusion.} Every graph $G$ with $\boldsymbol{\tau}\subset G\subset R(\boldsymbol{\tau})$ is a connected spanning subgraph of $\mathbf{G}$, because it contains the spanning tree $\boldsymbol{\tau}$ and $R(\boldsymbol{\tau})\subset\mathbf{G}$. By the preceding paragraph, the connected spanning subgraphs $G$ of $\mathbf{G}$ with $T(G)=\boldsymbol{\tau}$ are exactly the graphs $\boldsymbol{\tau}\cup A$ with $A\subset R(\boldsymbol{\tau})\setminus\boldsymbol{\tau}$, and $\#|\boldsymbol{\tau}\cup A|=n-1+\#|A|$. The sets $\left\{G:T(G)=\boldsymbol{\tau}\right\}$, $\boldsymbol{\tau}$ a spanning tree of $\mathbf{G}$, form a partition of the set of the connected spanning subgraphs of $\mathbf{G}$. Hence, by the definition of $\rho^T$,
\begin{align*}
\rho^T(X_1,\dots,X_n)&=\sum_{\boldsymbol{\tau}}\ \sum_{G:\,T(G)=\boldsymbol{\tau}}(-1)^{\#|G|}=\sum_{\boldsymbol{\tau}}(-1)^{n-1}\sum_{A\subset R(\boldsymbol{\tau})\setminus\boldsymbol{\tau}}(-1)^{\#|A|}\\
&=(-1)^{n-1}\cdot\#\left|\left\{\boldsymbol{\tau}:R(\boldsymbol{\tau})=\boldsymbol{\tau}\right\}\right|,
\end{align*}
where $\boldsymbol{\tau}$ ranges over the spanning trees of $\mathbf{G}$; the inner sum equals $(1-1)^{\#|R(\boldsymbol{\tau})\setminus\boldsymbol{\tau}|}$, which is $1$ if $R(\boldsymbol{\tau})=\boldsymbol{\tau}$ and $0$ otherwise. Taking absolute values, and noting that the spanning trees $\boldsymbol{\tau}$ with $R(\boldsymbol{\tau})=\boldsymbol{\tau}$ form a subset of the set of all spanning trees of $\mathbf{G}$, one obtains~\eqref{it:polymers-CE2}.

This completes the proof of Theorem~\ref{thm:polymer-activities-2}, and Proposition~\ref{prop:polymers} follows.
\end{proof}

\subsection{Lattice symmetries and local terms of activities}\label{sec:lattice-symmetries-local}

\subsubsection{Gamma matrices and symmetry transformations}\label{sec:gamma-matrices-symmetry}

The matrices, groups and actions of Section~\ref{sec:dirac-matrices-lattice} are the following. Let $\sigma_1$, $\sigma_2$, $\sigma_3$ be the Pauli matrices and $\sigma_0:=1$ the unit $2\times2$ matrix. The symbol $\otimes$ denotes the Kronecker product of matrices; thus $(N_1\otimes N_2)(N_3\otimes N_4)=N_1N_3\otimes N_2N_4$ and $(N_1\otimes N_2)^T=N_1^T\otimes N_2^T$, where $N^T$ is the transpose of a matrix $N$. For a complex matrix $N$, $\overline{N}$ is the matrix of the complex conjugate entries; the same notation is used for complex functions. Put
\[
\gamma_{\mu}:=\sigma_3\otimes\sigma_{\mu+1}\quad(\mu=0,1,2),\qquad\gamma_3:=\sigma_1\otimes1,\qquad\mathsf{C}:=i\gamma_1,\qquad\varkappa:=\gamma_3\cdot\mathsf{C}.
\]
The Pauli matrices satisfy $\sigma_a\sigma_b+\sigma_b\sigma_a=2\delta_{ab}$ for $a,b\in\{1,2,3\}$, and $\sigma_1\sigma_2=i\sigma_3$, $\sigma_2\sigma_3=i\sigma_1$, $\sigma_3\sigma_1=i\sigma_2$; the matrices $\sigma_1$ and $\sigma_3$ are real and symmetric, and $\sigma_2$ is imaginary and antisymmetric. From these facts one obtains
\begin{align}
&\gamma_{\mu}\gamma_{\nu}+\gamma_{\nu}\gamma_{\mu}=2\delta_{\mu\nu}\quad(0\le\mu,\nu\le3),\qquad\text{the matrices }\gamma_{\mu}\text{ are Hermitian},\notag\\
&\gamma_0,\ \gamma_2,\ \gamma_3\text{ are real and symmetric, and }\gamma_1\text{ is imaginary and antisymmetric},\notag\\
&\mathsf{C}\text{ is real},\qquad\mathsf{C}^T=\mathsf{C}^{-1}=-\mathsf{C},\qquad-\gamma_{\mu}^T=\mathsf{C}^{-1}\cdot\gamma_{\mu}\cdot\mathsf{C}\quad(0\le\mu\le3),\notag\\
&\varkappa\text{ is real},\qquad\varkappa^T=\varkappa=\varkappa^{-1},\qquad\overline{\gamma_{\mu}}=\varkappa^{-1}\cdot\gamma_{\mu}\cdot\varkappa\quad(0\le\mu\le2),\notag\\
&\overline{\gamma_3}=-\varkappa^{-1}\cdot\gamma_3\cdot\varkappa,\qquad\varkappa\cdot\mathsf{C}=-\mathsf{C}\cdot\varkappa.\label{eq:polymer-activities-13}
\end{align}

For $0\le\mu,\nu\le2$ one has $\gamma_{\mu}\gamma_{\nu}=1\otimes\sigma_{\mu+1}\sigma_{\nu+1}$ and $\gamma_{\mu}\gamma_3=\sigma_3\sigma_1\otimes\sigma_{\mu+1}=-\gamma_3\gamma_{\mu}$; moreover $\gamma_3^2=1$. Kronecker products of Hermitian matrices are Hermitian, Kronecker products of real symmetric matrices are real and symmetric, and $\gamma_1$ is the Kronecker product of the real symmetric matrix $\sigma_3$ with the imaginary antisymmetric matrix $\sigma_2$. This proves the relations of~\eqref{eq:polymer-activities-13} that involve the matrices $\gamma_{\mu}$ only. As $\gamma_1$ is imaginary and antisymmetric with $\gamma_1^2=1$, the matrix $\mathsf{C}=i\gamma_1$ is real, $\mathsf{C}^T=-\mathsf{C}$ and $\mathsf{C}^2=-1$. Next,
\[
\mathsf{C}^{-1}\cdot\gamma_{\mu}\cdot\mathsf{C}=\gamma_1\gamma_{\mu}\gamma_1
\]
is equal to $-\gamma_{\mu}$ for $\mu\ne1$ and to $\gamma_1$ for $\mu=1$; in both cases this is $-\gamma_{\mu}^T$. The matrix $\varkappa=i\gamma_3\gamma_1$ is real, being the product of the real matrices $\gamma_3$ and $\mathsf{C}$. One has $(\gamma_3\gamma_1)^2=-1$, hence $\varkappa^2=1$; and, since $\mathsf{C}$ anticommutes with $\gamma_3$,
\[
\varkappa^T=\mathsf{C}^T\cdot\gamma_3=-\mathsf{C}\cdot\gamma_3=\gamma_3\cdot\mathsf{C}=\varkappa.
\]
Further,
\[
\varkappa\cdot\gamma_{\mu}\cdot\varkappa=-\gamma_3\gamma_1\gamma_{\mu}\gamma_3\gamma_1
\]
is equal to $\gamma_{\mu}$ for $\mu=0,2$, to $-\gamma_1$ for $\mu=1$ and to $-\gamma_3$ for $\mu=3$. Since $\gamma_0$, $\gamma_2$, $\gamma_3$ are real and $\gamma_1$ is imaginary, this yields the two relations for $\overline{\gamma_{\mu}}$. Finally, $\varkappa\cdot\mathsf{C}=\gamma_3\cdot\mathsf{C}^2=-\gamma_3$ and $\mathsf{C}\cdot\varkappa=\mathsf{C}\cdot\gamma_3\cdot\mathsf{C}=-\gamma_3\cdot\mathsf{C}^2=\gamma_3$.

A complex square matrix is called \emph{monomial} if each of its rows and each of its columns contains exactly one non-zero entry, and this entry has modulus one. The monomial matrices of a given size form a group, which contains the unit matrix multiplied by any number of modulus one and is closed under transposition; the Kronecker product of two monomial matrices is again monomial. The matrices $\sigma_0,\dots,\sigma_3$ are monomial; hence so are $\gamma_0,\dots,\gamma_3$, $\mathsf{C}$ and $\varkappa$.

Let $\mathrm{Sp}$ be the group generated by $\gamma_0,\gamma_1,\gamma_2,\gamma_3$. Since these matrices anticommute pairwise and have square $1$, the elements of $\mathrm{Sp}$ are the matrices
\[
\pm\gamma_0^{n_0}\gamma_1^{n_1}\gamma_2^{n_2}\gamma_3^{n_3},\qquad n_0,n_1,n_2,n_3\in\{0,1\}.
\]
In particular $S$ and $\left(S^{-1}\right)^T$ are monomial for every $S\in\mathrm{Sp}$. Up to factors of modulus one, the sixteen products $\gamma_0^{n_0}\gamma_1^{n_1}\gamma_2^{n_2}\gamma_3^{n_3}$ are the sixteen matrices $\sigma_a\otimes\sigma_b$, $a,b\in\{0,1,2,3\}$. Indeed, the products with $n_3=0$ are, up to such factors, the eight matrices $\sigma_a\otimes\sigma_b$ with $a\in\{0,3\}$, because $\gamma_{\mu}\gamma_{\nu}=1\otimes\sigma_{\mu+1}\sigma_{\nu+1}$ for $0\le\mu,\nu\le2$ and $\gamma_0\gamma_1\gamma_2=i\sigma_3\otimes1$; and multiplication from the right by $\gamma_3=\sigma_1\otimes1$ sends $\sigma_0\otimes\sigma_b$ to $\sigma_1\otimes\sigma_b$ and $\sigma_3\otimes\sigma_b$ to $i\sigma_2\otimes\sigma_b$. The matrices $\sigma_a\otimes\sigma_b$ form a basis of the algebra of all complex $4\times4$ matrices. Hence this algebra is spanned by $\mathrm{Sp}$, a matrix that commutes with $\gamma_0,\dots,\gamma_3$ commutes with every $4\times4$ matrix, and
\begin{equation}\tag{Sch}\label{eq:polymer-activities-Sch}
\text{every $4\times4$ matrix that commutes with $\gamma_0$, $\gamma_1$, $\gamma_2$, $\gamma_3$ is a multiple of $1$.}
\end{equation}

Let $S\in\mathrm{Sp}$ and $0\le\nu\le2$. Each of the generators $\gamma_0,\dots,\gamma_3$ commutes or anticommutes with $\gamma_{\nu}$; therefore
\[
S^{-1}\gamma_{\nu}S=r_{\nu}(S)\,\gamma_{\nu}\qquad\text{with}\qquad r_{\nu}(S)\in\{1,-1\}.
\]
Put $r(S):=\operatorname{diag}\left(r_0(S),r_1(S),r_2(S)\right)$. The diagonal $3\times3$ matrices with diagonal entries $\pm1$ form a group, the group of signs. Since $(SS')^{-1}\gamma_{\nu}SS'=r_{\nu}(S)\,r_{\nu}(S')\,\gamma_{\nu}$, the map $r$ is a homomorphism of $\mathrm{Sp}$ into the group of signs. For $0\le\mu\le2$ the matrix $r(\gamma_{\mu})$ has the entry $1$ at the place $\mu$ and the entry $-1$ at the two other places, and $r(\gamma_3)=-1$. Let
\[
\mathrm{Sp}_3:=\left\{S\in\mathrm{Sp}:\ S\gamma_3=\gamma_3S\right\}.
\]
The matrix $\gamma_3$ anticommutes with $\gamma_0$, $\gamma_1$ and $\gamma_2$, hence with $\gamma_0\gamma_1\gamma_2$, and it commutes with $\gamma_{\mu}\gamma_0\gamma_1\gamma_2$ for $0\le\mu\le2$. So $\mathrm{Sp}_3$ contains $\gamma_3$ and the three matrices $\gamma_{\mu}\gamma_0\gamma_1\gamma_2$, $0\le\mu\le2$. As $r(\gamma_0\gamma_1\gamma_2)=r(\gamma_0)\,r(\gamma_1)\,r(\gamma_2)=1$, one has
\[
r(\gamma_{\mu}\gamma_0\gamma_1\gamma_2)=r(\gamma_{\mu})\qquad(0\le\mu\le2).
\]
Consequently the group $r(\mathrm{Sp}_3)$ contains $-1$ and $r(\gamma_0)$, $r(\gamma_1)$, $r(\gamma_2)$, hence the three matrices $-r(\gamma_{\mu})$, each of which has a single entry $-1$; every element of the group of signs is $1$ or a product of some of these three matrices. Hence $r$ maps $\mathrm{Sp}_3$, and a fortiori $\mathrm{Sp}$, onto the group of signs.

The maps
\[
S\mapsto\left(S^{-1}\right)^T,\qquad S\mapsto\mathsf{C}^{-1}\cdot S\cdot\mathsf{C},\qquad S\mapsto\overline{S},\qquad S\mapsto\varkappa^{-1}\cdot S\cdot\varkappa
\]
are homomorphisms of $\mathrm{Sp}$ into the group of the invertible $4\times4$ matrices. As $\gamma_{\mu}^{-1}=\gamma_{\mu}$, the relations~\eqref{eq:polymer-activities-13} show that the first two maps agree up to a sign on each generator $\gamma_{\mu}$, and so do the last two. Every element of $\mathrm{Sp}$ is a product of generators; therefore, for every $S\in\mathrm{Sp}$,
\[
\left(S^{-1}\right)^T=\pm\,\mathsf{C}^{-1}\cdot S\cdot\mathsf{C},\qquad\overline{S}=\pm\,\varkappa^{-1}\cdot S\cdot\varkappa,
\]
with signs that depend on $S$.

Let $k\ge0$ and $l\in\mathrm{Cl}$ (Section~\ref{sec:lattice-geometry}). The group $\mathbb{G}_k$ is the set of the pairs $(S,a)$ with $S\in\mathrm{Sp}$ and $a\in M\mathbb{Z}^3/l\mathbb{Z}^3$, with the product
\[
(S,a)(S',a'):=\left(SS',\,a+r(S)a'\right).
\]
An element $g=(S,a)$ acts on $\mathbb{R}^3/l\mathbb{Z}^3$ by
\[
gx:=r(S)x+a,
\]
and $(gh)x=g(hx)$ for $g,h\in\mathbb{G}_k$. The subgroup $\mathbb{G}_k^L$ consists of the elements $(S,a)$ of $\mathbb{G}_k$ with $a\in LM\mathbb{Z}^3/l\mathbb{Z}^3$. Given $S\in\mathrm{Sp}$ and an $M$-cube with centre $c$, define
\[
g_c:=\left(S,\,(1-r(S))\,c\right).
\]
The diagonal entries of $1-r(S)$ are $0$ or $2$, and $c\in M\mathbb{Z}^3$; hence $(1-r(S))\,c\in2M\mathbb{Z}^3$, so that $g_c$ belongs to $\mathbb{G}_k$, and
\[
g_cx=r(S)(x-c)+c.
\]

An element $g=(S,a)\in\mathbb{G}_k$ acts on functions, families and sets as follows:
\begin{itemize}
\item on the bond functions $A$ of each of the lattices $L^{-i}\mathbb{Z}^3/l\mathbb{Z}^3$, $-1\le i\le k$, by $(gA)(x,x'):=A\left(g^{-1}x,g^{-1}x'\right)$;
\item on scalar functions $\boldsymbol{\lambda}$ of the sites by $(g\boldsymbol{\lambda})(x):=\boldsymbol{\lambda}\left(g^{-1}x\right)$;
\item on plaquette functions $F$, that is, on functions defined on the oriented plaquettes whose value changes sign when the orientation of the plaquette is reversed, by $(gF)(p):=F\left(g^{-1}p\right)$;
\item on spinor functions $f$ (Section~\ref{sec:antiperiodicity}) by $(gf)(x):=S\cdot f\left(g^{-1}x\right)$;
\item on the generators of the Grassmann algebras by the automorphism $a_g$ of Section~\ref{sec:lattice-gauge-charge};
\item on families of parameters $s=(s_c)$, indexed by the parameter cubes $c$ (Definition~\ref{def:parameter-cubes-polydisks}), by $(gs)_c:=s_{g^{-1}c}$; if cubes of $\pi''_k$ are among the indices, $g$ is taken in $\mathbb{G}_k^L$, so that $g^{-1}c$ is again a cube of $\pi''_k$ whenever $c$ is (Proposition~\ref{prop:properties-actions}(a));
\item on sets $X$ by $X\mapsto gX$.
\end{itemize}
For a function with the antiperiodicity~\eqref{eq:polymer-activities-AP} the formulas above require a vector $a\in M\mathbb{Z}^3$ in place of its class modulo $l\mathbb{Z}^3$. The translation by $l\cdot n$, $n\in\mathbb{Z}^3$, multiplies such a function by $(-1)^{n_0+n_1+n_2}$, while the element $(-1,0)$, formed with $S=-1$, multiplies it by $-1$; each of these elements also multiplies all generators of the Grassmann algebras by one common sign, since the generators are antiperiodic in the same way and $\left((-1)^{-1}\right)^T=-1$. Since the group acts on an operator by conjugating it, and since an even element of a Grassmann algebra is unchanged when every generator is replaced by its negative (Section~\ref{sec:antiperiodicity}), these elements act as the identity on operators and on the even elements of the Grassmann algebras.

\begin{proposition}[Properties of the actions]\label{prop:properties-actions}
Let $g=(S,a)\in\mathbb{G}_k$ and $r:=r(S)$.
\begin{enumerate}
\item[(a)] Let the class $a$ be represented by a vector of $M\mathbb{Z}^3$, denoted again by $a$, and let $gx:=rx+a$ for $x\in\mathbb{R}^3$. This affine map of $\mathbb{R}^3$ is an isometry for the sup norm, the $\ell^1$ norm and the Euclidean norm, and $g(x+l\cdot n)=gx+l\cdot rn$ for $n\in\mathbb{Z}^3$; hence it induces on $\mathbb{R}^3/l\mathbb{Z}^3$ the map $x\mapsto gx$ defined above, whichever representative of $a$ is chosen. It maps each lattice $L^{-i}\mathbb{Z}^3$, $-1\le i\le k$, onto itself, and, for $y\in L^{-i}\mathbb{Z}^3$, the block $\left\{x:|x-y|<L^{-i}/2\right\}$ with centre $y$ onto the block with centre $gy$. The induced map of the torus, denoted again by $g$, maps the rectilinear path $\Gamma^{\pi}(y,x)$ (Section~\ref{sec:averaging-fluctuations-gauge}) onto the path $\Gamma^{\pi}(gy,gx)$ with the same order $\pi$, segments onto segments, $M$-cubes onto $M$-cubes, the centre of an $M$-cube onto the centre of its image, adjacent cubes onto adjacent cubes, and nets onto nets of the same length. Consequently $g$ maps polymers onto polymers and small polymers onto small polymers, and
\[
|gX|_M=|X|_M,\qquad d_M(gX)=d_M(X),\qquad(gX)^{\sim i}=g\left(X^{\sim i}\right),\qquad(g\square)^{\dagger}=g\left(\square^{\dagger}\right).
\]
Further, $g$ maps the cubes of $\pi_k$, resp.\ of $\pi'_k$, onto cubes of the same kind and linked cubes onto linked cubes; $K_{gz}=gK_z$; $g$ maps $\mathrm{Pz}_k$ onto $\mathrm{Pz}_k$, a pair that is admissible for $X$ of kind (I), resp.\ (II), onto a pair that is admissible for $gX$ of the same kind, and the shortest chains of adjacent cubes of $X$ onto the shortest chains of adjacent cubes of $gX$. If $g\in\mathbb{G}_k^L$, then $g$ maps $LM$-cubes onto $LM$-cubes and the cubes of $\pi''_k$ onto cubes of $\pi''_k$, and $g\left(\bar{X}^L\right)$ is the union of the $LM$-cubes that contain an $M$-cube of $gX$.
\item[(b)] The maps $A\mapsto gA$ on bond functions and $F\mapsto gF$ on plaquette functions preserve the sup, $\ell^1$ and $\ell^2$ norms and the pairings (Section~\ref{sec:averaging-fluctuations-gauge}). The map $A\mapsto gA$ preserves the three norms of Section~\ref{sec:domains-bond-functions}, also on subsets: for a set $W$ of sites, each of the three norms of $gA$ on $gW$ coincides with the same norm of $A$ on $W$. Moreover
\[
(\tau\,gA)(gy,gx)=(\tau A)(y,x),\qquad g(d\boldsymbol{\lambda})=d\left(\boldsymbol{\lambda}\circ g^{-1}\right),
\]
where $\tau$ is the axial functional (Section~\ref{sec:averaging-fluctuations-gauge}). Let $N$ be an element of $\mathrm{Sp}$, or $N=\mathsf{C}$, or $N=\varkappa$. Then $|Nf|_{\infty}=|f|_{\infty}$ for every spinor function $f$, where $(Nf)(x):=N\cdot f(x)$ and $|f|_{\infty}$ is the supremum, over all sites $x$, of the moduli of the four components of $f(x)$. For every real function $\Lambda$ on the sites of $T_k$, the four maps $a_g$, $G_{\Lambda}$, $\mathbf{C}$, $\mathbf{K}$ of Section~\ref{sec:lattice-gauge-charge} satisfy~\eqref{eq:polymer-activities-4}.
\item[(c)] For every $c\in(0,1]$, each of the sets $R_k$, $c\cdot\tilde{R}_k$, $\tilde{R}_k^{-}$, $G_k$, $\tilde{G}_k$, $\mathfrak{X}_k$ of Definition~\ref{def:domains} is invariant under the maps $A\mapsto gA$, $A\mapsto-A$ and $A\mapsto\overline{A}$. The sets $c\cdot\tilde{R}_k$, $\tilde{R}_k^{-}$ and $\tilde{G}_k$ are moreover invariant under $A\mapsto A-d\boldsymbol{\lambda}$ for every real function $\boldsymbol{\lambda}$ of the sites, and under the addition of real constant bond functions.
\item[(d)] Let $g\in\mathbb{G}_k^L$ and $g':=(S,a/L)$. The side $l$ is subject only to the condition $l\in\mathrm{Cl}$; thus $l/L=n\cdot M$ with an odd integer $n\ge13$, so that $l/L$ is a multiple of $M$ and $l/L>2M$, and $l/L$ need not lie in $\mathrm{Cl}$. The objects of level $k+1$ in~(d) are formed with the side $l/L$. The torus $T_{k+1}$, the set $\mathrm{Pz}_{k+1}$, the generators and the Grassmann algebra $\mathrm{Gr}[\psi,\chi]$ of level $k+1$, the scalings $\psi_L$ and $\chi_L$, and the scaling $A_L$ of the bond functions $A$ on $T_{k+1}$ are those defined, for every $l\in\mathrm{Cl}$, in the paragraph on the change of scale in Section~\ref{sec:local-parts-adjustment}; $\mathbb{G}_{k+1}$ is the group of Section~\ref{sec:dirac-matrices-lattice} for the side $l/L$, and $a_{g'}$ is the automorphism of Section~\ref{sec:lattice-gauge-charge} formed at level $k+1$ with this side. The map $g\mapsto g'$ is an isomorphism of $\mathbb{G}_k^L$ onto $\mathbb{G}_{k+1}$, and
\[
(g'A)_L=g\left(A_L\right),\qquad L\left(g'Y\right)=g\left(LY\right),\qquad S_L\circ a_g=a_{g'}\circ S_L
\]
for every bond function $A$ on one of the lattices $L^{-i}\mathbb{Z}^3/(l/L)\mathbb{Z}^3$, $0\le i\le k+1$, and every set $Y$ of level $k+1$. Here $A_L$ is the bond function on $L^{-(i-1)}\mathbb{Z}^3/l\mathbb{Z}^3$, one of the lattices on which $g$ acts (Section~\ref{sec:gamma-matrices-symmetry}), given by $A_L(b):=L^{-\frac{1}{2}}A(b/L)$. For $i=k+1$ the lattice of $A$ is $T_{k+1}$, and $A_L$ is the scaling of Section~\ref{sec:local-parts-adjustment}; for $i=0$ the lattice of $A$ is the unit lattice $T^0_{k+1}$, and, for a real bond function $A'$ on it, $A'_L$ is the function of Section~\ref{sec:averaging-fluctuations-gauge}. Further, $LY:=\{Ly:\ y\in Y\}$, and $S_L$ is the homomorphism of the Grassmann algebra $\mathrm{Gr}[\psi,\chi]$ of level $k$ into that of level $k+1$ with
\[
\psi(x)\mapsto L^{-1}\psi(x/L),\qquad\chi(x,x')\mapsto L^{-1-\alpha}\chi(x/L,x'/L),
\]
for the barred and the unbarred generators alike; that is, $S_LF(\psi,\chi)=F(\psi_L,\chi_L)$, with $\psi_L$ and $\chi_L$ as in Section~\ref{sec:local-parts-adjustment}.
\end{enumerate}
\end{proposition}

\begin{proof}
(a) Here $r$ is a diagonal matrix with diagonal entries $\pm1$, and $a\in M\mathbb{Z}^3$; every assertion of (a) rests on these two facts. The affine map $x\mapsto rx+a$ preserves the sup, $\ell^1$ and Euclidean distances, carries segments onto segments of the same length, and satisfies $g(x+l\cdot n)=gx+l\cdot rn$. Since $r$ maps $l\mathbb{Z}^3$ onto itself, the vectors that represent the difference $gx-gx'$ of the images of two points $x$, $x'$ of the torus are the images under $r$ of the vectors that represent $x-x'$; so the induced map of the torus preserves the distance of two points, which is the least norm of such a vector. Another representative $a+l\cdot n$ of the class $a$ alters the affine map by the translation by $l\cdot n$, which is the identity on the torus. The matrix $r$ maps each lattice $L^{-i}\mathbb{Z}^3$ onto itself, and $M=L^m$ (Section~\ref{sec:lattice-geometry}) with $m\ge2$ (stage~\ref{it:order-choice-6}), so that $a\in M\mathbb{Z}^3\subset L\mathbb{Z}^3\subset L^{-i}\mathbb{Z}^3$ for $-1\le i\le k$. Hence $g$ maps each of the lattices $L^{-i}\mathbb{Z}^3$, $-1\le i\le k$, onto itself; in the same way it maps $M\mathbb{Z}^3$ onto itself, and $LM\mathbb{Z}^3$ onto itself if $g\in\mathbb{G}_k^L$, since then $a\in LM\mathbb{Z}^3$.

Write the order as $\pi=(\pi_1,\pi_2,\pi_3)$, put $w^{(0)}:=y$, and for $i=1,2,3$ let $w^{(i)}$ be the point whose coordinate in the direction $\pi_i$ is that of $x$ and whose two other coordinates are those of $w^{(i-1)}$. Then $w^{(3)}=x$, and $\Gamma^{\pi}(y,x)$ is composed of the straight paths from $w^{(i-1)}$ to $w^{(i)}$, $i=1,2,3$. As $r$ is diagonal, the coordinate $\mu$ of $gw$ is determined by the coordinate $\mu$ of $w$ alone. Consequently $gw^{(0)}=gy$, and for $i=1,2,3$ the point $gw^{(i)}$ has the coordinate of $gx$ in the direction $\pi_i$ and the two other coordinates of $gw^{(i-1)}$; thus $gw^{(0)},\dots,gw^{(3)}$ are the corners of $\Gamma^{\pi}(gy,gx)$. Since $g$ carries segments onto segments, it carries the straight path from $w^{(i-1)}$ to $w^{(i)}$ onto the straight path from $gw^{(i-1)}$ to $gw^{(i)}$. Therefore $g$ maps $\Gamma^{\pi}(y,x)$ onto $\Gamma^{\pi}(gy,gx)$, bond by bond.

Let $y$ be a point of one of the lattices above and $t>0$. As $g$ preserves distances, it maps the set $\{x:\ |x-y|<t\}$ onto the set $\{x:\ |x-gy|<t\}$, and likewise with $\le$ in place of $<$; moreover $gy$ lies in the same lattice as $y$. The remaining objects of (a) are given by conditions of this form, or are built from objects so given, from segments, from the paths $\Gamma^{\pi}$ and from lengths: the blocks; the $M$-cubes and their centres; adjacency of cubes (Section~\ref{sec:polymers}); the polymers, $|X|_M$, the sets $X^{\sim i}$ and the sets $K_z$ (Section~\ref{sec:lattice-geometry}); the sets $\square^{\dagger}$ (Definition~\ref{def:enlarged-cubes}); the cubes of $\pi_k$ and $\pi'_k$ and the relation of being linked (Definition~\ref{def:parameter-cubes-polydisks}); the set $\mathrm{Pz}_k$ and the admissible pairs of kinds (I) and (II) (Definition~\ref{def:pairs-sites-holder}); the shortest chains of adjacent cubes; the nets, their lengths, $d_M$ (convention~\eqref{eq:polymer-activities-6}) and the small polymers; and, if $g\in\mathbb{G}_k^L$, the $LM$-cubes, the cubes of $\pi''_k$ and the sets $\bar{X}^L$ (convention~\eqref{eq:polymer-activities-7}). All data that enter these definitions are preserved by $g$ and by $g^{-1}$; hence $g$ carries each of these objects onto its counterpart formed with the images of the data.

(b) By (a), $g$ permutes the oriented bonds and the oriented plaquettes of each of the lattices and commutes with the reversal of the orientation. Since $(gA)(b)=A(g^{-1}b)$ and $(gF)(p)=F(g^{-1}p)$, the sup, $\ell^1$ and $\ell^2$ norms and the pairings are preserved. Now let $A$ be a bond function on $T_k$, whose spacing is $\eta=L^{-k}$, let $x$ be a site, $z:=g^{-1}x$, and $r=\operatorname{diag}(r_0,r_1,r_2)$; put $s_{\mu}:=\eta e_{\mu}$ if $r_{\mu}=-1$ and $s_{\mu}:=0$ if $r_{\mu}=1$. Since $g^{-1}(x+\eta e_{\mu})=z+r_{\mu}\eta e_{\mu}$ and $A$ is odd under reversal of the bonds,
\[
(gA)_{\mu}(x)=A\left(z,z+r_{\mu}\eta e_{\mu}\right)=r_{\mu}A_{\mu}(z-s_{\mu}),
\]
that is, $(gA)_{\mu}(x)=A_{\mu}(z)$ if $r_{\mu}=1$ and $(gA)_{\mu}(x)=-A_{\mu}(z-\eta e_{\mu})$ if $r_{\mu}=-1$. Applying this at $x$ and at $x+\eta e_{\nu}$, whose image under $g^{-1}$ is $z+r_{\nu}\eta e_{\nu}$, one obtains
\[
\left(\partial_{\nu}(gA)_{\mu}\right)(x)=r_{\mu}r_{\nu}\left(\partial_{\nu}A_{\mu}\right)(z-s_{\mu}-s_{\nu}).
\]
Here the point $z-s_{\mu}-s_{\nu}$ is one of $z$, $z-\eta e_{\mu}$, $z-\eta e_{\nu}$, $z-\eta e_{\mu}-\eta e_{\nu}$, and the vector $s_{\mu}+s_{\nu}$ and the sign $r_{\mu}r_{\nu}$ are determined by $\mu$, $\nu$ and $r$; they do not vary with $x$. As $x\mapsto g^{-1}x-s_{\mu}-s_{\nu}$ is a bijection of the set of the sites of $T_k$, for fixed $(\mu,\nu)$ the function $|\partial_{\nu}(gA)_{\mu}|$ takes the values of $|\partial_{\nu}A_{\mu}|$, permuted; the same holds for the H\"older quotients $|x-x'|^{-\alpha}\left|\partial_{\nu}A_{\mu}(x')-\partial_{\nu}A_{\mu}(x)\right|$, because the shift cancels in the difference of the two points and $|g^{-1}x-g^{-1}x'|=|x-x'|$. This proves the invariance of the three norms; on a subset $W$ the same computation applies to the bonds and differences in $W$ and in $gW$. The identity for $\tau$ follows from (a), since $g$ maps each oriented path $\Gamma^{\pi}(y,x)$ onto the oriented path $\Gamma^{\pi}(gy,gx)$, bond by bond. By the definition of $d\boldsymbol{\lambda}$ (Section~\ref{sec:settings-basic-notation}),
\[
\left(g(d\boldsymbol{\lambda})\right)(x,x')=(d\boldsymbol{\lambda})\left(g^{-1}x,g^{-1}x'\right)=\eta^{-1}\left(\boldsymbol{\lambda}\left(g^{-1}x'\right)-\boldsymbol{\lambda}\left(g^{-1}x\right)\right)=\left(d\left(\boldsymbol{\lambda}\circ g^{-1}\right)\right)(x,x').
\]
Finally, the matrix $N$ is monomial, being an element of $\mathrm{Sp}$ or one of $\mathsf{C}$, $\varkappa$; hence the components of $N\cdot f(x)$ are the components of $f(x)$, permuted and multiplied by numbers of modulus one. This gives $|Nf|_{\infty}=|f|_{\infty}$, and~\eqref{eq:polymer-activities-4} is derived from the same fact in Section~\ref{sec:lattice-gauge-charge}.

(c) By Definition~\ref{def:domains}, each of the six sets is described by strict bounds on the three norms, or on some of them, of complex bond functions on $T_k$; the sets $c\cdot\tilde{R}_k$, $\tilde{R}_k^{-}$ and $\tilde{G}_k$ are described in addition by a splitting $A=A_0+A_1$ with real $A_0$ such that, for every $M$-cube $\square$, $A_0=A'+d\boldsymbol{\lambda}_{\square}$ on $\square^{\dagger}$, resp.\ on $\square^{\sim5}$, with a real function $\boldsymbol{\lambda}_{\square}$ and a real bond function $A'$ whose norms on that set obey bounds. These descriptions are preserved by $A\mapsto gA$: by (a), $\square\mapsto g\square$ is a bijection of the set of the $M$-cubes, with $g\left(\square^{\dagger}\right)=(g\square)^{\dagger}$ and $g\left(\square^{\sim5}\right)=(g\square)^{\sim5}$; by (b), each norm of $gA'$ on $g\left(\square^{\dagger}\right)$, resp.\ $g\left(\square^{\sim5}\right)$, equals the same norm of $A'$ on $\square^{\dagger}$, resp.\ $\square^{\sim5}$, and the norms on $T_k$ are preserved; and $g(A'+d\boldsymbol{\lambda})=gA'+d\left(\boldsymbol{\lambda}\circ g^{-1}\right)$, where $\boldsymbol{\lambda}\circ g^{-1}$ is real. They are preserved by $A\mapsto-A$, which replaces $A_0$, $A_1$, $A'$, $\boldsymbol{\lambda}_{\square}$ by their negatives, and by $A\mapsto\overline{A}$, which replaces $A_1$ by $\overline{A_1}$ and leaves the real functions unchanged; neither map changes a norm. If $\boldsymbol{\lambda}$ is a real function of the sites and $A_0=A'+d\boldsymbol{\lambda}_{\square}$ on $\square^{\dagger}$, then
\[
A_0-d\boldsymbol{\lambda}=A'+d\left(\boldsymbol{\lambda}_{\square}-\boldsymbol{\lambda}\right)\qquad\text{on }\square^{\dagger},
\]
with the same $A'$; the same holds on $\square^{\sim5}$, and $A-d\boldsymbol{\lambda}=(A_0-d\boldsymbol{\lambda})+A_1$. Finally, the cubes $\square^{\dagger}$ and $\square^{\sim5}$ have the sides $(2L+1)M$ and $11M$ (Definition~\ref{def:enlarged-cubes}), both smaller than $13LM\le l$, as $l\in\mathrm{Cl}$. So these cubes do not wrap, and by Section~\ref{sec:balls-torus-lifts} the nearest neighbours in them and the vectors $x-x_{\square}$, where $x_{\square}$ is the centre of $\square$, are those of the corresponding cube of $\mathbb{L}_k$. Hence a real constant bond function $v$, with $v_{\mu}(x)=v_{\mu}\in\mathbb{R}$ for all sites $x$ and directions $\mu$, coincides on these cubes with $d$ of the real function $x\mapsto v\cdot(x-x_{\square})$, and
\[
A_0+v=A'+d\left(\boldsymbol{\lambda}_{\square}+v\cdot(x-x_{\square})\right)\qquad\text{on }\square^{\dagger},
\]
with the same $A'$; the same holds on $\square^{\sim5}$, and $A+v=(A_0+v)+A_1$.

(d) By the definition of $\mathrm{Cl}$ (Section~\ref{sec:lattice-geometry}), $l/L=n\cdot M$ with an odd integer $n\ge13$. Thus $l/L$ is a multiple of $M$ and $l/L\ge13M>2M$, and the objects of level $k+1$ that occur in~(d) are defined (Section~\ref{sec:local-parts-adjustment}, paragraph on the change of scale). As $l/L$ is a multiple of $M$, the map $a\mapsto a/L$ is a bijection of $LM\mathbb{Z}^3/l\mathbb{Z}^3$ onto $M\mathbb{Z}^3/(l/L)\mathbb{Z}^3$; hence $g\mapsto g'$ is a bijection of $\mathbb{G}_k^L$ onto $\mathbb{G}_{k+1}$. It respects products: for $h=(S',a')\in\mathbb{G}_k^L$,
\[
(gh)'=\left(SS',\,(a+r(S)a')/L\right)=\left(S,a/L\right)\left(S',a'/L\right)=g'h'.
\]
The automorphism $a_{g'}$ is well defined. Indeed, $r=r(S)$ is diagonal with diagonal entries $\pm1$, and $a/L\in M\mathbb{Z}^3\subset\mathbb{L}_{k+1}$; so the map $y\mapsto g'y=ry+a/L$ is an isometry of $\mathbb{R}^3$ that maps $\mathbb{L}_{k+1}$ onto itself and satisfies $g'\left(y+(l/L)\cdot n\right)=g'y+(l/L)\cdot rn$ for $n\in\mathbb{Z}^3$. Hence it induces an isometry of $T_{k+1}$, it maps $\mathrm{Pz}_{k+1}$ onto itself and shortest lifts onto shortest lifts, and, since $|rn|_1=|n|_1$, the formulas of Section~\ref{sec:lattice-gauge-charge} define $a_{g'}$ compatibly with the antiperiodicity of level $k+1$. This is the argument of~(a) with $l/L$ in place of $l$; the side enters it only through the fact that $r$ maps the lattice of periods onto itself. The homomorphism $S_L$ is well defined as well: the map $x\mapsto x/L$ carries $\mathbb{L}_k$ onto $\mathbb{L}_{k+1}$ and $l\mathbb{Z}^3$ onto $(l/L)\mathbb{Z}^3$ and divides distances by $L$, so that $(x/L,x'/L)\in\mathrm{Pz}_{k+1}$ for $(x,x')\in\mathrm{Pz}_k$, a shortest lift of $(x,x')$ going to a shortest lift of $(x/L,x'/L)$; and $(x+l\cdot e_{\mu})/L=x/L+(l/L)\cdot e_{\mu}$, so that $S_L$ respects the antiperiodicity of the generators. For $x\in\mathbb{R}^3$ one has
\[
g'^{-1}(x/L)=r^{-1}\left(x/L-a/L\right)=r^{-1}(x-a)/L=\left(g^{-1}x\right)/L,
\]
and in the same way $g'y=g(Ly)/L$, which gives $L\left(g'Y\right)=g\left(LY\right)$. Let $0\le i\le k+1$. As $-1\le i-1\le k$, statement~(a), applied to the element $g^{-1}$ of $\mathbb{G}_k$, shows that $x\mapsto g^{-1}x$ maps the lattice $L^{-(i-1)}\mathbb{Z}^3$ onto itself; hence it permutes the bonds of $L^{-(i-1)}\mathbb{Z}^3/l\mathbb{Z}^3$. Therefore, for every bond function $A$ on $L^{-i}\mathbb{Z}^3/(l/L)\mathbb{Z}^3$ and every bond $b$ of $L^{-(i-1)}\mathbb{Z}^3/l\mathbb{Z}^3$,
\[
(g'A)_L(b)=L^{-\frac{1}{2}}A\left(g'^{-1}(b/L)\right)=L^{-\frac{1}{2}}A\left(\left(g^{-1}b\right)/L\right)=\left(g\left(A_L\right)\right)(b).
\]
The maps $S_L\circ a_g$ and $a_{g'}\circ S_L$ are homomorphisms of Grassmann algebras, and they agree on the generators:
\[
S_L\left(a_g\psi(x)\right)=L^{-1}S\cdot\psi\left(\left(g^{-1}x\right)/L\right)=L^{-1}S\cdot\psi\left(g'^{-1}(x/L)\right)=a_{g'}\left(S_L\psi(x)\right);
\]
the same computation applies to $\bar{\psi}$, with $\left(S^{-1}\right)^T$ in place of $S$, and to $\chi$ and $\bar{\chi}$, with $L^{-1-\alpha}$ in place of $L^{-1}$ and two arguments in place of one. So $S_L\circ a_g=a_{g'}\circ S_L$.

Together with the monomiality of the elements of $\mathrm{Sp}$ and of $\mathsf{C}$, $\varkappa$, statements (a)--(d) give the facts (a)--(d) of Proposition~\ref{prop:symmetries-extraction}.
\end{proof}

\subsubsection{Symmetries of the kernels of low degree}\label{sec:symmetries-kernels-low}

For a set $W$ of sites of $T_k$ and a complex $4\times4$ matrix $N$ put
\[
\int_W\mathrm{d}x\;\bar{\psi}(x)^T\cdot N\cdot\psi(x):=\sum_{x\in W}\eta^3\cdot\bar{\psi}(x)^T\cdot N\cdot\psi(x),
\]
with the columns $\psi(x)$, $\bar{\psi}(x)$ of Section~\ref{sec:coefficients-kernels-grassmann}. A point of $\mathbb{L}_k$ represents the site of $T_k$ that is its class modulo $l\mathbb{Z}^3$. Let $\delta_W$ be the real function on $\mathbb{L}_k\times\mathbb{L}_k$ defined by
\[
\delta_W(x,y):=\begin{cases}(-1)^{|n|_1}&\text{if }y=x+l\,n\text{ with }n\in\mathbb{Z}^3\text{, and }x\text{ represents a site of }W,\\0&\text{otherwise,}\end{cases}
\]
where $|\cdot|_1$ is the $\ell^1$ norm of Section~\ref{sec:balls-torus-lifts}. A translation of one of the two arguments by $l\cdot e_{\mu}$ does not change the site that this argument represents, and it replaces $n$ by $n\mp e_{\mu}$, which changes the parity of $|n|_1$; hence $\delta_W$ is antiperiodic in each argument. Moreover $\delta_W(y,x)=\delta_W(x,y)$, since $x=y+l\cdot(-n)$ and $x$, $y$ represent the same site whenever $y-x\in l\mathbb{Z}^3$. If $x$ and $y$ lie in a subset of $\mathbb{L}_k$ that contains at most one representative of each site, such as a fundamental domain of $l\mathbb{Z}^3$ or one lift of a small polymer (Section~\ref{sec:antiperiodicity}), then $y-x\in l\mathbb{Z}^3$ only for $x=y$; so $\delta_W(x,y)=1$ if $x=y$ represents a site of $W$, and $\delta_W(x,y)=0$ otherwise. The integrals over $x$ and over $y$ are sums over a fundamental domain, each with the weight $\eta^3$ (Section~\ref{sec:antiperiodicity}); it follows that
\[
\int\mathrm{d}x\,\mathrm{d}y\;\bar{\psi}(x)^T\cdot\eta^{-3}\cdot\delta_W(x,y)\cdot N\cdot\psi(y)=\int_W\mathrm{d}x\;\bar{\psi}(x)^T\cdot N\cdot\psi(x).
\]

\begin{lemma}[Symmetries of $\boldsymbol{\alpha}_0$ and $\boldsymbol{\alpha}_2$]\label{lem:symmetries}
Let $\mathscr{E}$ be a family of elements $\mathscr{E}(X';\mathscr{A},\mathscr{C})$ of $\mathrm{Gr}[\psi,\chi]$, indexed by the polymers $X'\in D_k$, with the following properties: each element has equally many barred and unbarred legs; the family is defined at $(\mathscr{A},\mathscr{C})=(0,0)$; and the kernels $\mathscr{E}_{20}(X';0,0;x,y)$, in the matrix notation of Section~\ref{sec:coefficients-kernels-grassmann}, vanish unless both $x$ and $y$ lie in $X'$ (this is the first clause of the locality property~\eqref{it:renormalisation-group-1-k2} of Definition~\ref{def:renormalisation-group-1}). Let $X\in D_k$ be small, and let $\boldsymbol{\alpha}_0(\mathscr{E},X)$ and $\boldsymbol{\alpha}_2(\mathscr{E},X)$ be the number and the $4\times4$ matrix of Section~\ref{sec:local-parts-adjustment}.
\begin{enumerate}
\item[(i)] Assume that~\eqref{it:renormalisation-group-1-k3} holds at $(0,0)$, that is, $a_g\left[\mathscr{E}(gX';0,0)\right]=\mathscr{E}(X';0,0)$ for all $g\in\mathbb{G}_k$ and all $X'\in D_k$. Then, for every $g=(S,a)\in\mathbb{G}_k$,
\[
\boldsymbol{\alpha}_0(\mathscr{E},gX)=\boldsymbol{\alpha}_0(\mathscr{E},X),\qquad\boldsymbol{\alpha}_2(\mathscr{E},gX)=S\cdot\boldsymbol{\alpha}_2(\mathscr{E},X)\cdot S^{-1}.
\]
\item[(ii)] Assume that~\eqref{it:renormalisation-group-1-k5} holds at $(0,0)$, that is, $\mathbf{C}\left[\mathscr{E}(X';0,0)\right]=\mathscr{E}(X';0,0)$ for all $X'\in D_k$. Then
\[
\boldsymbol{\alpha}_2(\mathscr{E},X)=\mathsf{C}^{-1}\cdot\boldsymbol{\alpha}_2(\mathscr{E},X)^T\cdot\mathsf{C}.
\]
\item[(iii)] Assume that~\eqref{it:renormalisation-group-1-k7} holds at $(0,0)$, that is, $\mathbf{K}\left[\mathscr{E}(X';0,0)\right]=\mathscr{E}(X';0,0)$ for all $X'\in D_k$. Then $\boldsymbol{\alpha}_0(\mathscr{E},X)$ is real and
\[
\boldsymbol{\alpha}_2(\mathscr{E},X)=\varkappa^{-1}\cdot\overline{\boldsymbol{\alpha}_2(\mathscr{E},X)}\cdot\varkappa.
\]
\end{enumerate}
Moreover, for every complex $4\times4$ matrix $N$, every $g=(S,a)\in\mathbb{G}_k$ and every real function $\Lambda$ on the sites of $T_k$, the maps of Section~\ref{sec:lattice-gauge-charge} satisfy
\begin{align*}
a_g\left[\int_{gX}\mathrm{d}x\;\bar{\psi}(x)^T\cdot\left(S\cdot N\cdot S^{-1}\right)\cdot\psi(x)\right]&=\int_X\mathrm{d}x\;\bar{\psi}(x)^T\cdot N\cdot\psi(x),\\
G_{\Lambda}\left[\int_X\mathrm{d}x\;\bar{\psi}(x)^T\cdot N\cdot\psi(x)\right]&=\int_X\mathrm{d}x\;\bar{\psi}(x)^T\cdot N\cdot\psi(x),\\
\mathbf{C}\left[\int_X\mathrm{d}x\;\bar{\psi}(x)^T\cdot N\cdot\psi(x)\right]&=\int_X\mathrm{d}x\;\bar{\psi}(x)^T\cdot\left(\mathsf{C}^{-1}\cdot N^T\cdot\mathsf{C}\right)\cdot\psi(x),\\
\mathbf{K}\left[\int_X\mathrm{d}x\;\bar{\psi}(x)^T\cdot N\cdot\psi(x)\right]&=\int_X\mathrm{d}x\;\bar{\psi}(x)^T\cdot\left(\varkappa^{-1}\cdot\overline{N}\cdot\varkappa\right)\cdot\psi(x).
\end{align*}
\end{lemma}

\begin{proof}
The four maps $a_g$, $G_{\Lambda}$, $\mathbf{C}$, $\mathbf{K}$ preserve the bidegree, by~\eqref{eq:polymer-activities-4}. Hence each of the identities~\eqref{it:renormalisation-group-1-k3}, \eqref{it:renormalisation-group-1-k5}, \eqref{it:renormalisation-group-1-k7} at $(0,0)$ holds separately for the parts of bidegree $(0,0)$ and for the parts of bidegree $(2,0)$.

Let $K$ be a function on $\mathbb{L}_k\times\mathbb{L}_k$ with values in the complex $4\times4$ matrices that is antiperiodic in each argument (Section~\ref{sec:antiperiodicity}), and let $g=(S,a)\in\mathbb{G}_k$. Since $a_g\psi(y)=S\cdot\psi(g^{-1}y)$ and $a_g\bar{\psi}(x)=(S^{-1})^T\cdot\bar{\psi}(g^{-1}x)$,
\begin{align*}
a_g\left[\int\mathrm{d}x\,\mathrm{d}y\;\bar{\psi}(x)^T\cdot K(x,y)\cdot\psi(y)\right]
&=\int\mathrm{d}x\,\mathrm{d}y\;\bar{\psi}(g^{-1}x)^T\cdot S^{-1}\cdot K(x,y)\cdot S\cdot\psi(g^{-1}y)\\
&=\int\mathrm{d}x\,\mathrm{d}y\;\bar{\psi}(x)^T\cdot S^{-1}\cdot K(gx,gy)\cdot S\cdot\psi(y);
\end{align*}
in the second step the variables are changed by $x\mapsto gx$, $y\mapsto gy$, which preserves the measure and carries a lift of the ball of Section~\ref{sec:antiperiodicity} around $X$ onto a lift of the corresponding ball around $gX$. Since $\psi$ and $\bar{\psi}$ anticommute,
\begin{align*}
\mathbf{C}\left[\int\mathrm{d}x\,\mathrm{d}y\;\bar{\psi}(x)^T\cdot K(x,y)\cdot\psi(y)\right]
&=-\int\mathrm{d}x\,\mathrm{d}y\;\psi(x)^T\cdot\mathsf{C}^T\cdot K(x,y)\cdot\mathsf{C}\cdot\bar{\psi}(y)\\
&=\int\mathrm{d}x\,\mathrm{d}y\;\bar{\psi}(y)^T\cdot\mathsf{C}^T\cdot K(x,y)^T\cdot\mathsf{C}\cdot\psi(x),
\end{align*}
and, since $\mathbf{K}$ replaces the coefficients by their complex conjugates,
\[
\mathbf{K}\left[\int\mathrm{d}x\,\mathrm{d}y\;\bar{\psi}(x)^T\cdot K(x,y)\cdot\psi(y)\right]=\int\mathrm{d}x\,\mathrm{d}y\;\bar{\psi}(x)^T\cdot\varkappa^T\cdot\overline{K(x,y)}\cdot\varkappa\cdot\psi(y).
\]

Now let $K(x,y):=\mathscr{E}_{20}(X;0,0;x,y)$. Since every element of the family has equally many barred and unbarred legs, the first identity of~\eqref{eq:polymer-activities-3} expresses the parts of bidegree $(2,0)$ through the matrix kernels: for every $X'\in D_k$, in particular for $X'=X$ and for $X'=gX$,
\[
\mathscr{E}^{[2,0]}(X';0,0)=\int\mathrm{d}x\,\mathrm{d}y\;\bar{\psi}(x)^T\cdot\mathscr{E}_{20}(X';0,0;x,y)\cdot\psi(y),
\]
where $\mathscr{E}_{20}(X';0,0;\cdot\,,\cdot)$ is antiperiodic in each argument (Section~\ref{sec:antiperiodicity}). Conversely, an element $\int\mathrm{d}x\,\mathrm{d}y\;\bar{\psi}(x)^T\cdot K'(x,y)\cdot\psi(y)$, with a function $K'$ of the kind considered above, determines $K'$: it is written in the form~\eqref{eq:polymer-activities-2} by the antisymmetric kernel that takes the value $K'(x,y)_{\boldsymbol{\alpha}\boldsymbol{\beta}}$ at $\left((x,\boldsymbol{\alpha},1),(y,\boldsymbol{\beta},0)\right)$, the value $-K'(x,y)_{\boldsymbol{\alpha}\boldsymbol{\beta}}$ at $\left((y,\boldsymbol{\beta},0),(x,\boldsymbol{\alpha},1)\right)$ and the value $0$ at the pairs of two barred and at the pairs of two unbarred arguments, and by Proposition~\ref{prop:weighted-coefficient-norms}\eqref{it:weighted-coefficient-norms-N0} the antisymmetric kernels of an element are determined by the element. By the first paragraph of this proof, \eqref{it:renormalisation-group-1-k3}, \eqref{it:renormalisation-group-1-k5} and~\eqref{it:renormalisation-group-1-k7} at $(0,0)$ give, respectively,
\begin{gather*}
a_g\left[\mathscr{E}^{[2,0]}(gX;0,0)\right]=\mathscr{E}^{[2,0]}(X;0,0),\\
\mathbf{C}\left[\mathscr{E}^{[2,0]}(X;0,0)\right]=\mathscr{E}^{[2,0]}(X;0,0),\\
\mathbf{K}\left[\mathscr{E}^{[2,0]}(X;0,0)\right]=\mathscr{E}^{[2,0]}(X;0,0).
\end{gather*}
Write $\mathscr{E}^{[2,0]}(gX;0,0)$ and $\mathscr{E}^{[2,0]}(X;0,0)$ through their matrix kernels $\mathscr{E}_{20}(gX;0,0;\cdot\,,\cdot)$ and $K$, evaluate the left sides by the formulas for $a_g$, $\mathbf{C}$ and $\mathbf{K}$ of the preceding paragraph (for $\mathbf{C}$ with the names of $x$ and $y$ exchanged), and use $\mathsf{C}^T=\mathsf{C}^{-1}$, $\varkappa^T=\varkappa^{-1}$, which hold by~\eqref{eq:polymer-activities-13}. The functions of $(x,y)$ obtained in this way are again antiperiodic in each argument, because $g(x+l\,n)=gx+l\cdot r(S)n$ (Proposition~\ref{prop:properties-actions}(a)). Comparing the matrix kernels of the two sides one obtains, respectively,
\[
S^{-1}\cdot\mathscr{E}_{20}(gX;0,0;gx,gy)\cdot S=K(x,y),\qquad K(x,y)=\mathsf{C}^{-1}\cdot K(y,x)^T\cdot\mathsf{C},\qquad K(x,y)=\varkappa^{-1}\cdot\overline{K(x,y)}\cdot\varkappa.
\]
Integrate these identities over $x,y\in X$, all points being taken in one lift of $X$. By the hypothesis on the kernels,
\[
\int_{X\times X}\mathrm{d}x\,\mathrm{d}y\;K(x,y)=\operatorname{Vol}(X)\cdot\boldsymbol{\alpha}_2(\mathscr{E},X),\qquad\int_{X\times X}\mathrm{d}x\,\mathrm{d}y\;\mathscr{E}_{20}(gX;0,0;gx,gy)=\operatorname{Vol}(gX)\cdot\boldsymbol{\alpha}_2(\mathscr{E},gX),
\]
the second identity after the substitution $x\mapsto g^{-1}x$, $y\mapsto g^{-1}y$. Dividing by $\operatorname{Vol}(gX)=\operatorname{Vol}(X)$ (Proposition~\ref{prop:properties-actions}(a)) one obtains the assertions on $\boldsymbol{\alpha}_2$ in (i), (ii) and (iii). In bidegree $(0,0)$ the identities~\eqref{it:renormalisation-group-1-k3} and~\eqref{it:renormalisation-group-1-k7} at $(0,0)$ give
\[
\mathscr{E}_{00}(gX;0,0)=\mathscr{E}_{00}(X;0,0),\qquad\overline{\mathscr{E}_{00}(X;0,0)}=\mathscr{E}_{00}(X;0,0),
\]
because $a_g$ leaves the scalars fixed and $\mathbf{K}$ replaces them by their complex conjugates; division by $\operatorname{Vol}(gX)=\operatorname{Vol}(X)$ gives the assertions on $\boldsymbol{\alpha}_0$.

It remains to prove the four final identities. Let $N$ be a complex $4\times4$ matrix and $g=(S,a)\in\mathbb{G}_k$. For all $x,y\in\mathbb{L}_k$
\[
\delta_{gX}(gx,gy)=\delta_X(x,y),
\]
because $g(x+l\,n)=gx+l\cdot r(S)n$ (Proposition~\ref{prop:properties-actions}(a)), $|r(S)n|_1=|n|_1$, and $gx$ represents a site of $gX$ if and only if $x$ represents a site of $X$. By the identity displayed before the lemma, the elements $\int_{gX}\mathrm{d}x\;\bar{\psi}(x)^T\cdot\left(S\cdot N\cdot S^{-1}\right)\cdot\psi(x)$ and $\int_X\mathrm{d}x\;\bar{\psi}(x)^T\cdot N\cdot\psi(x)$ are formed with the kernels $\eta^{-3}\cdot\delta_{gX}(x,y)\cdot S\cdot N\cdot S^{-1}$ and $\eta^{-3}\cdot\delta_X(x,y)\cdot N$, both antiperiodic in each argument. The formulas for $a_g$, $\mathbf{C}$ and $\mathbf{K}$ of the second paragraph of this proof, applied to the first kernel with $a_g$ and to the second kernel with $\mathbf{C}$ and with $\mathbf{K}$, give elements of the same form with the kernels
\[
S^{-1}\cdot\eta^{-3}\cdot\delta_{gX}(gx,gy)\cdot S\cdot N\cdot S^{-1}\cdot S=\eta^{-3}\cdot\delta_X(x,y)\cdot N,\qquad\eta^{-3}\cdot\delta_X(y,x)\cdot\mathsf{C}^T\cdot N^T\cdot\mathsf{C},\qquad\eta^{-3}\cdot\delta_X(x,y)\cdot\varkappa^T\cdot\overline{N}\cdot\varkappa,
\]
the second after the names of $x$ and $y$ have been exchanged. As $\delta_X(y,x)=\delta_X(x,y)$, $\mathsf{C}^T=\mathsf{C}^{-1}$ and $\varkappa^T=\varkappa^{-1}$, the identity displayed before the lemma identifies these three elements with the right sides of the first, the third and the fourth final identity. Finally, $G_{\Lambda}$ multiplies $\bar{\psi}(x)$ by $\exp\left(-i\,e_k\cdot\Lambda(x)\right)$ and $\psi(x)$ by $\exp\left(i\,e_k\cdot\Lambda(x)\right)$; the product of these two numbers is $1$, so that $G_{\Lambda}$ fixes every product $\bar{\psi}_{\boldsymbol{\alpha}}(x)\cdot\psi_{\boldsymbol{\beta}}(x)$, and the second final identity follows.
\end{proof}

\subsubsection{Extraction of the constant and of the mass term}\label{sec:extraction-constant-mass}

\begin{theorem}[Extraction]\label{thm:extraction}
Assume the setting~\eqref{eq:polymer-activities-Sb}, and let $\mathscr{E}\in K^{x,C}_k$. Let $R$ and $L$ be the extraction and the reblocking of Section~\ref{sec:local-parts-adjustment}.
\begin{enumerate}
\item[(a)] For an $M$-cube $\square$ let
\[
\varepsilon(\mathscr{E},\square):=-\sum_{X\ \text{small},\ X\supset\square}\boldsymbol{\alpha}_0(\mathscr{E},X),\qquad m(\mathscr{E},\square):=-\sum_{X\ \text{small},\ X\supset\square}\boldsymbol{\alpha}_2(\mathscr{E},X).
\]
Neither $\varepsilon(\mathscr{E},\square)$ nor $m(\mathscr{E},\square)$ changes when $\square$ is replaced by another $M$-cube, and $m(\mathscr{E},\square)=m(\mathscr{E})\cdot\mathbf{1}$ with a complex number $m(\mathscr{E})$, $\mathbf{1}$ denoting the unit $4\times4$ matrix; thus $\varepsilon(\mathscr{E})=\varepsilon(\mathscr{E},\square)$ and $m(\mathscr{E})$ are the numbers of Section~\ref{sec:local-parts-adjustment}. The maps $\mathscr{E}\mapsto\varepsilon(\mathscr{E})$ and $\mathscr{E}\mapsto m(\mathscr{E})$ are complex linear, and
\[
\varepsilon(\mathscr{E})=\varepsilon(\mathscr{E}|_0),\qquad m(\mathscr{E})=m(\mathscr{E}|_0).
\]
If $\mathscr{E}\in K^x_k$, then $\varepsilon(\mathscr{E})$ and $m(\mathscr{E})$ are real.
\item[(b)] One has
\begin{align*}
|\varepsilon(\mathscr{E})|&\le K_0\cdot M^{-3}\cdot\|\mathscr{E}\|^x_k\le c_{\mathrm{fl}}\cdot\|\mathscr{E}\|^x_k,\\
|m(\mathscr{E})|&\le K_0\cdot M^{-3}\cdot e_k^{\frac{1}{2}}\cdot\|\mathscr{E}\|^x_k\le c_{\mathrm{fl}}\cdot e_k^{\frac{1}{2}}\cdot\|\mathscr{E}\|^x_k.
\end{align*}
\item[(c)] For all $(\mathscr{A},\mathscr{C})\in\tilde{R}_k\times\mathfrak{X}_k$ one has, in $\mathrm{Gr}[\psi,\chi]$,
\[
\sum_X\mathscr{E}(X;\mathscr{A},\mathscr{C})=-\varepsilon(\mathscr{E})\cdot n_k-m(\mathscr{E})\cdot\left\langle\bar{\psi},\psi\right\rangle+\sum_X(R\,\mathscr{E})(X;\mathscr{A},\mathscr{C}),
\]
where $n_k=l^3$ and $\left\langle\bar{\psi},\psi\right\rangle$ is the pairing of Section~\ref{sec:fermion-integrals-blocking}.
\item[(d)] The extraction $R$ is a linear map with $R\left(K^{x,C}_k\right)\subset K^{x,C}_k$ and $R\left(K^x_k\right)\subset K^x_k$. For all $X\in D_k$ and $(\mathscr{A},\mathscr{C})\in\tilde{R}_k\times\mathfrak{X}_k$
\[
\left\|(R\,\mathscr{E})(X;\mathscr{A},\mathscr{C})\right\|_{h_k}\le\left\|\mathscr{E}(X;\mathscr{A},\mathscr{C})\right\|_{h_k}+\left\|\mathscr{E}(X;0,0)\right\|_{h_k},
\]
and consequently $\|R\,\mathscr{E}\|^x_k\le2\|\mathscr{E}\|^x_k\le c_{\mathrm{fl}}\cdot\|\mathscr{E}\|^x_k$. For every small $X$
\[
(R\,\mathscr{E})_{00}(X;0,0)=0,\qquad\int\mathrm{d}x\,\mathrm{d}y\;(R\,\mathscr{E})_{20}(X;0,0;x,y)=0.
\]
Finally, $(R\,\mathscr{E})|_0=R\,(\mathscr{E}|_0)$.
\item[(e)] Let $E'=\left\{E'(X;\mathscr{A},\mathscr{C}):X\in D_k,\ (\mathscr{A},\mathscr{C})\in\tilde{R}_k\times\mathfrak{X}_k\right\}$ be a family of elements of the algebra $\mathrm{Gr}[\psi,\chi]$ of the level $k$. The homomorphism $S_L$ of Proposition~\ref{prop:properties-actions}(d) satisfies $S_LF(\psi,\chi)=F(\psi_L,\chi_L)$; by the definition of $L$ in Section~\ref{sec:local-parts-adjustment} it follows that
\[
(L\,E')(Y;\mathscr{A},\mathscr{C})=\sum_{X:\ \bar{X}^L=LY}S_L\left[E'(X;\mathscr{A}_L,\mathscr{C}_L)\right]\qquad\left(Y\in D_{k+1}\right),
\]
and by Lemma~\ref{lem:scaling-field-domains} the right side is defined for $(\mathscr{A},\mathscr{C})\in\tilde{R}_{k+1}\times\mathfrak{X}_{k+1}$. Suppose that $E'$ satisfies~\eqref{it:renormalisation-group-1-k3} for every $g\in\mathbb{G}_k^L$; then $L\,E'$ satisfies~\eqref{it:renormalisation-group-1-k3} at the level $k+1$ for every element of the group $\mathbb{G}_{k+1}$. Each of the properties~\eqref{it:renormalisation-group-1-k4}, \eqref{it:renormalisation-group-1-k5} and~\eqref{it:renormalisation-group-1-k7}, if $E'$ has it, holds for $L\,E'$ at the level $k+1$.
\end{enumerate}
Assertions (a)--(e) remain true for the families without the second slot, with $K_k^C$, $K_k$ and $\|\cdot\|_k$ in place of $K^{x,C}_k$, $K^x_k$ and $\|\cdot\|^x_k$.
\end{theorem}

\begin{proof}
By Definition~\ref{def:renormalisation-group-1}, in the form of Section~\ref{sec:lattice-gauge-charge}, every $\mathscr{E}\in K^{x,C}_k$ satisfies the hypotheses of Lemma~\ref{lem:symmetries} and of its parts (i) and (ii): its elements have equally many barred and unbarred legs, it is defined at $(0,0)\in\tilde{R}_k\times\mathfrak{X}_k$, its kernels have the locality property~\eqref{it:renormalisation-group-1-k2}, and it has the properties~\eqref{it:renormalisation-group-1-k3} and~\eqref{it:renormalisation-group-1-k5}. If $\mathscr{E}\in K^x_k$, it has the property~\eqref{it:renormalisation-group-1-k7} as well, and part (iii) of the lemma applies.

\emph{(a)} Let $a\in M\mathbb{Z}^3$ and $g=(\mathbf{1},a)\in\mathbb{G}_k$. By Proposition~\ref{prop:properties-actions}(a) the map $X\mapsto X+a$ is a bijection of the set of the small polymers that contain $\square$ onto the set of the small polymers that contain $\square+a$, and by Lemma~\ref{lem:symmetries}(i) corresponding terms of the sums that define $\varepsilon(\mathscr{E},\cdot)$ and $m(\mathscr{E},\cdot)$ at $\square$ and at $\square+a$ are equal. Any two $M$-cubes differ by a translation by an element of $M\mathbb{Z}^3$; hence $\varepsilon(\mathscr{E},\square)$ and $m(\mathscr{E},\square)$ do not depend on $\square$. Let $c$ be the centre of $\square$ and $S\in\mathrm{Sp}$. The element $g_c\in\mathbb{G}_k$ of Section~\ref{sec:gamma-matrices-symmetry}, formed with $S$, satisfies $g_cx=r(S)(x-c)+c$; so it maps $\square$ onto itself, and therefore the set of the small polymers that contain $\square$ onto itself (Proposition~\ref{prop:properties-actions}(a)). Lemma~\ref{lem:symmetries}(i) now gives
\[
m(\mathscr{E},\square)=S\cdot m(\mathscr{E},\square)\cdot S^{-1}\qquad(S\in\mathrm{Sp}).
\]
In particular $m(\mathscr{E},\square)$ commutes with $\gamma_0,\dots,\gamma_3$, and $m(\mathscr{E},\square)=m(\mathscr{E})\cdot\mathbf{1}$ by~\eqref{eq:polymer-activities-Sch}. The numbers $\boldsymbol{\alpha}_0(\mathscr{E},X)$ and the matrices $\boldsymbol{\alpha}_2(\mathscr{E},X)$ depend linearly on $\mathscr{E}$, and they contain only the element $\mathscr{E}(X;0,0)$, which is also $\mathscr{E}|_0(X;0)$; this proves the linearity and the identities $\varepsilon(\mathscr{E})=\varepsilon(\mathscr{E}|_0)$, $m(\mathscr{E})=m(\mathscr{E}|_0)$. Reality: for $\mathscr{E}\in K^x_k$, Lemma~\ref{lem:symmetries}(iii) gives $\overline{\varepsilon(\mathscr{E})}=\varepsilon(\mathscr{E})$ and
\[
m(\mathscr{E})\cdot\mathbf{1}=\varkappa^{-1}\cdot\overline{m(\mathscr{E})}\cdot\varkappa=\overline{m(\mathscr{E})}\cdot\mathbf{1}.
\]

\emph{(b)} The pair $(0,0)$ lies in $\tilde{R}_k\times\mathfrak{X}_k$. For a small polymer $X$ one has $\operatorname{Vol}(X)\cdot|\boldsymbol{\alpha}_0(\mathscr{E},X)|=|\mathscr{E}_{00}(X;0,0)|\le\|\mathscr{E}(X;0,0)\|_{h_k}$. Next, the definition of $\boldsymbol{\alpha}_2(\mathscr{E},X)$ and~\eqref{eq:polymer-activities-3}, applied with the weights $h=h_k=(h_{k,1},h_{k,2})$ of Section~\ref{sec:grassmann-algebras-norms}, whose first entry is $h_{k,1}=e_k^{-\frac{1}{4}}$, so that $h_{k,1}^{-2}=e_k^{\frac{1}{2}}$, give
\begin{multline}\label{eq:polymer-activities-14}
\operatorname{Vol}(X)\cdot\sum_{\boldsymbol{\alpha},\boldsymbol{\beta}}\left|\boldsymbol{\alpha}_2(\mathscr{E},X)_{\boldsymbol{\alpha}\boldsymbol{\beta}}\right|\le\sum_{\boldsymbol{\alpha},\boldsymbol{\beta}}\int\mathrm{d}x\,\mathrm{d}y\;\left|\mathscr{E}_{20}(X;0,0;x,y)_{\boldsymbol{\alpha}\boldsymbol{\beta}}\right|\\
=e_k^{\frac{1}{2}}\cdot\left\|\mathscr{E}^{[2,0]}(X;0,0)\right\|_{h_k}.
\end{multline}
Further, $\operatorname{Vol}(X)\ge M^3$; $\|\mathscr{E}(X;0,0)\|_{h_k}\le\exp\left(-\kappa\cdot d_M(X)\right)\cdot\|\mathscr{E}\|^x_k$ by the definition of the norm $\|\cdot\|^x_k$ in Definition~\ref{def:renormalisation-group-1}; and, since $\kappa\ge\kappa_0$ (stage~\ref{it:order-choice-3}), Proposition~\ref{prop:polymers}\eqref{it:polymers-t2} gives $\sum_{X\supset\square}\exp\left(-\kappa\cdot d_M(X)\right)\le K_0$. The norm is additive over the bidegrees (Section~\ref{sec:coefficients-kernels-grassmann}); hence $\|\mathscr{E}^{[2,0]}(X;0,0)\|_{h_k}\le\|\mathscr{E}(X;0,0)\|_{h_k}$. Finally, by (a) the matrix $\sum_X\boldsymbol{\alpha}_2(\mathscr{E},X)$, summed over the small polymers $X$ that contain $\square$, equals $-m(\mathscr{E})\cdot\mathbf{1}$; each of its diagonal entries is therefore $-m(\mathscr{E})$, and $|m(\mathscr{E})|$ does not exceed the sum of the moduli of its sixteen entries. With all sums over the small polymers $X$ that contain $\square$ it follows that
\begin{align*}
|\varepsilon(\mathscr{E})|&\le\sum_X\operatorname{Vol}(X)^{-1}\cdot\left\|\mathscr{E}(X;0,0)\right\|_{h_k}\le M^{-3}\cdot\|\mathscr{E}\|^x_k\cdot\sum_X\exp\left(-\kappa\cdot d_M(X)\right)\le K_0\cdot M^{-3}\cdot\|\mathscr{E}\|^x_k,\\
|m(\mathscr{E})|&\le\sum_X\sum_{\boldsymbol{\alpha},\boldsymbol{\beta}}\left|\boldsymbol{\alpha}_2(\mathscr{E},X)_{\boldsymbol{\alpha}\boldsymbol{\beta}}\right|\le\sum_X\operatorname{Vol}(X)^{-1}\cdot e_k^{\frac{1}{2}}\cdot\left\|\mathscr{E}^{[2,0]}(X;0,0)\right\|_{h_k}\\
&\le M^{-3}\cdot e_k^{\frac{1}{2}}\cdot\|\mathscr{E}\|^x_k\cdot\sum_X\exp\left(-\kappa\cdot d_M(X)\right)\le K_0\cdot M^{-3}\cdot e_k^{\frac{1}{2}}\cdot\|\mathscr{E}\|^x_k.
\end{align*}
The remaining inequalities of (b) hold because $M\ge1$ and $c_{\mathrm{fl}}=61\,c_X\cdot K_0\ge K_0$ by~\eqref{eq:polymer-activities-5}.

\emph{(c)} Fix $(\mathscr{A},\mathscr{C})\in\tilde{R}_k\times\mathfrak{X}_k$ and write $\mathscr{E}(X)$ for $\mathscr{E}(X;\mathscr{A},\mathscr{C})$. By the definition of $R$ in Section~\ref{sec:local-parts-adjustment},
\[
\sum_X\mathscr{E}(X)-\sum_X(R\,\mathscr{E})(X)=\sum_{X\ \text{small}}\left(\boldsymbol{\alpha}_0(\mathscr{E},X)\cdot\operatorname{Vol}(X)+\int_X\mathrm{d}x\;\bar{\psi}(x)^T\cdot\boldsymbol{\alpha}_2(\mathscr{E},X)\cdot\psi(x)\right).
\]
No site lies on the boundary of an $M$-cube (Section~\ref{sec:balls-torus-lifts}). Hence the $M$-cubes partition the sites of $T_k$, those contained in $X$ partition the sites of $X$, and
\[
\operatorname{Vol}(X)=\sum_{\square\subset X}M^3,\qquad\int_X\mathrm{d}x\;\bar{\psi}(x)^T\cdot N\cdot\psi(x)=\sum_{\square\subset X}\int_{\square}\mathrm{d}x\;\bar{\psi}(x)^T\cdot N\cdot\psi(x)
\]
for every matrix $N$. Insert this into the right side and exchange the two sums. By the definitions of $\varepsilon(\mathscr{E},\square)$ and $m(\mathscr{E},\square)$ in (a), the right side becomes the first of the expressions
\[
\sum_{\square}\left(-\varepsilon(\mathscr{E},\square)\cdot M^3-\int_{\square}\mathrm{d}x\;\bar{\psi}(x)^T\cdot m(\mathscr{E},\square)\cdot\psi(x)\right)=-\varepsilon(\mathscr{E})\cdot l^3-m(\mathscr{E})\cdot\left\langle\bar{\psi},\psi\right\rangle,
\]
the sum being over all $M$-cubes of $T_k$, whose number is $l^3/M^3$ (Section~\ref{sec:polymers}); the equality holds by the assertions of (a), because the $M$-cubes partition the sites of $T_k$.

\emph{(d)} Let $X$ be small. By the identity displayed before Lemma~\ref{lem:symmetries}, the element $\int_X\mathrm{d}x\;\bar{\psi}(x)^T\cdot\boldsymbol{\alpha}_2(\mathscr{E},X)\cdot\psi(x)$ is formed with the kernel $\eta^{-3}\cdot\delta_X(x,y)\cdot\boldsymbol{\alpha}_2(\mathscr{E},X)$, which is antiperiodic in each argument; as in the proof of Lemma~\ref{lem:symmetries}, this function is its matrix kernel in the notation of~\eqref{eq:polymer-activities-3}. Since no site lies on the boundary of an $M$-cube (Section~\ref{sec:balls-torus-lifts}), $X$ has exactly $\operatorname{Vol}(X)\cdot\eta^{-3}$ sites. On one lift of $X$, and on a fundamental domain of $l\mathbb{Z}^3$, the function $\delta_X$ takes the value $1$ at the $\operatorname{Vol}(X)\cdot\eta^{-3}$ pairs $(x,x)$ with $x$ representing a site of $X$, and the value $0$ at all other pairs. Therefore the integral of this kernel over one lift of $X$ is
\[
\operatorname{Vol}(X)\cdot\boldsymbol{\alpha}_2(\mathscr{E},X)=\int\mathrm{d}x\,\mathrm{d}y\;\mathscr{E}_{20}(X;0,0;x,y),
\]
the equality being the definition of $\boldsymbol{\alpha}_2$ in Section~\ref{sec:local-parts-adjustment}; and, for all spinor indices $\boldsymbol{\alpha}$, $\boldsymbol{\beta}$, the integral of the modulus of the entry $(\boldsymbol{\alpha},\boldsymbol{\beta})$ of this kernel over a fundamental domain is $\operatorname{Vol}(X)\cdot\left|\boldsymbol{\alpha}_2(\mathscr{E},X)_{\boldsymbol{\alpha}\boldsymbol{\beta}}\right|$. By~\eqref{eq:polymer-activities-3}, applied with $h=h_k$, the norm of this element is thus the left side of
\[
h_{k,1}^2\cdot\operatorname{Vol}(X)\cdot\sum_{\boldsymbol{\alpha},\boldsymbol{\beta}}\left|\boldsymbol{\alpha}_2(\mathscr{E},X)_{\boldsymbol{\alpha}\boldsymbol{\beta}}\right|\le\left\|\mathscr{E}^{[2,0]}(X;0,0)\right\|_{h_k};
\]
the inequality is~\eqref{eq:polymer-activities-14}, multiplied by $h_{k,1}^2=e_k^{-\frac{1}{2}}$.
The other term subtracted in the definition of $(R\,\mathscr{E})(X)$, the number $\boldsymbol{\alpha}_0(\mathscr{E},X)\cdot\operatorname{Vol}(X)=\mathscr{E}_{00}(X;0,0)$, has the norm $|\mathscr{E}_{00}(X;0,0)|$. Since the norm is additive over the bidegrees, the sum of the norms of the two subtracted terms is at most $\|\mathscr{E}(X;0,0)\|_{h_k}$. For a large polymer $X$ one has $(R\,\mathscr{E})(X)=\mathscr{E}(X)$. This proves the bound on $(R\,\mathscr{E})(X;\mathscr{A},\mathscr{C})$; with the definition of $\|\cdot\|^x_k$ it gives $\|R\,\mathscr{E}\|^x_k\le2\|\mathscr{E}\|^x_k$, and $2\le c_{\mathrm{fl}}$ by~\eqref{eq:polymer-activities-5}. Moreover, for small $X$,
\[
(R\,\mathscr{E})_{00}(X;0,0)=\mathscr{E}_{00}(X;0,0)-\boldsymbol{\alpha}_0(\mathscr{E},X)\cdot\operatorname{Vol}(X)=0,
\]
and, the matrix kernel of $(R\,\mathscr{E})(X;0,0)$ being $\mathscr{E}_{20}(X;0,0;x,y)-\eta^{-3}\cdot\delta_X(x,y)\cdot\boldsymbol{\alpha}_2(\mathscr{E},X)$, the integral of $(R\,\mathscr{E})_{20}(X;0,0;x,y)$ over one lift of $X$ is $\int\mathrm{d}x\,\mathrm{d}y\;\mathscr{E}_{20}(X;0,0;x,y)-\operatorname{Vol}(X)\cdot\boldsymbol{\alpha}_2(\mathscr{E},X)=0$.

It remains to verify the properties of Definition~\ref{def:renormalisation-group-1} for $R\,\mathscr{E}$. The map $R$ is linear, because $\boldsymbol{\alpha}_0$ and $\boldsymbol{\alpha}_2$ are linear in $\mathscr{E}$. Neither of the two subtracted terms depends on $(\mathscr{A},\mathscr{C})$. The first is a number; the second is a linear combination of the products $\bar{\psi}_{\boldsymbol{\alpha}}(x)\cdot\psi_{\boldsymbol{\beta}}(x)$ with $x$ a site of $X$, each of which consists of one barred and one unbarred generator at a site of $X$. Hence the kernels of $R\,\mathscr{E}$ differ from those of $\mathscr{E}$ by terms that are constant in $(\mathscr{A},\mathscr{C})$ and vanish unless all their arguments lie in $X$, and $R\,\mathscr{E}$ inherits~\eqref{it:renormalisation-group-1-k1} and~\eqref{it:renormalisation-group-1-k2} from $\mathscr{E}$. The property~\eqref{it:renormalisation-group-1-k3} follows from Lemma~\ref{lem:symmetries}(i) and the first of its final identities with $N=\boldsymbol{\alpha}_2(\mathscr{E},X)$; here $gX$ is small if and only if $X$ is small (Proposition~\ref{prop:properties-actions}(a)). The property~\eqref{it:renormalisation-group-1-k4} follows from the second of these identities. The properties~\eqref{it:renormalisation-group-1-k5} and~\eqref{it:renormalisation-group-1-k7} follow from Lemma~\ref{lem:symmetries}(ii), resp.\ (iii), and from the third, resp.\ fourth, of these identities; in the case of~\eqref{it:renormalisation-group-1-k7} the number $\boldsymbol{\alpha}_0(\mathscr{E},X)$ is real by Lemma~\ref{lem:symmetries}(iii). With the bound proved above it follows that $R\,\mathscr{E}\in K^{x,C}_k$, and that $R\,\mathscr{E}\in K^x_k$ if $\mathscr{E}\in K^x_k$. Finally, $(R\,\mathscr{E})|_0=R\,(\mathscr{E}|_0)$, because $\boldsymbol{\alpha}_0$ and $\boldsymbol{\alpha}_2$ are formed at $(0,0)$.

\emph{(e)} Let $g'=(S,a')\in\mathbb{G}_{k+1}$, and let $g:=(S,La')\in\mathbb{G}_k^L$; then $g'$ is the image of $g$ under the isomorphism of $\mathbb{G}_k^L$ onto $\mathbb{G}_{k+1}$ of Proposition~\ref{prop:properties-actions}(d), and every element of $\mathbb{G}_{k+1}$ arises in this way. By Proposition~\ref{prop:properties-actions}(a) and (d), the map $X\mapsto gX$ is a bijection of the set $\left\{X:\bar{X}^L=LY\right\}$ onto the set $\left\{X:\bar{X}^L=L(g'Y)\right\}$, and
\[
(g'\mathscr{A})_L=g(\mathscr{A}_L),\qquad(g'\mathscr{C})_L=g(\mathscr{C}_L),\qquad a_{g'}\circ S_L=S_L\circ a_g.
\]
Hence, by the property~\eqref{it:renormalisation-group-1-k3} of $E'$ for $g$,
\begin{align*}
a_{g'}\left[(L\,E')(g'Y;g'\mathscr{A},g'\mathscr{C})\right]&=\sum_{X:\ \bar{X}^L=LY}a_{g'}S_L\left[E'\left(gX;g(\mathscr{A}_L),g(\mathscr{C}_L)\right)\right]\\
&=\sum_{X:\ \bar{X}^L=LY}S_L\,a_g\left[E'\left(gX;g(\mathscr{A}_L),g(\mathscr{C}_L)\right)\right]\\
&=\sum_{X:\ \bar{X}^L=LY}S_L\left[E'(X;\mathscr{A}_L,\mathscr{C}_L)\right]=(L\,E')(Y;\mathscr{A},\mathscr{C}),
\end{align*}
which is the property~\eqref{it:renormalisation-group-1-k3} of $L\,E'$ at the level $k+1$.

The property~\eqref{it:renormalisation-group-1-k4}. Every real function $\boldsymbol{\lambda}$ of the sites of $T_{k+1}$ defines a real function $\boldsymbol{\lambda}_L$ of the sites of $T_k$ by $\boldsymbol{\lambda}_L(x):=L^{\frac{1}{2}}\cdot\boldsymbol{\lambda}(x/L)$. Since $e_{k+1}=L^{\frac{1}{2}}e_k$ by~\eqref{eq:polymer-activities-Sb},
\[
e_k\cdot\boldsymbol{\lambda}_L(x)=e_{k+1}\cdot\boldsymbol{\lambda}(x/L),\qquad(d\boldsymbol{\lambda})_L=d(\boldsymbol{\lambda}_L),\qquad G_{\Lambda}\circ S_L=S_L\circ G_{\Lambda_L}
\]
for all real functions $\boldsymbol{\lambda}$, $\Lambda$ on the sites of $T_{k+1}$. Let $(\mathscr{A},\mathscr{C})$ and $(\mathscr{A}^1,\mathscr{C}^1)$ lie in $\tilde{R}_{k+1}\times\mathfrak{X}_{k+1}$, and let $\boldsymbol{\lambda}$, $\boldsymbol{\lambda}'$ be real functions on the sites of $T_{k+1}$ with $\mathscr{A}^1=\mathscr{A}-d\boldsymbol{\lambda}$ and $\mathscr{C}^1=\mathscr{C}-d\boldsymbol{\lambda}'$ on the bonds at distance at most $1$ from $Y$. Then
\[
\mathscr{A}^1_L=\mathscr{A}_L-d\boldsymbol{\lambda}_L,\qquad\mathscr{C}^1_L=\mathscr{C}_L-d\boldsymbol{\lambda}'_L
\]
on the bonds at distance at most $L$ from $LY$, and these bonds include the bonds at distance at most $1$ from every polymer $X$ contained in $LY$, in particular from every $X$ with $\bar{X}^L=LY$. The pairs $(\mathscr{A}_L,\mathscr{C}_L)$ and $(\mathscr{A}^1_L,\mathscr{C}^1_L)$ lie in $\tilde{R}_k\times\mathfrak{X}_k$ by Lemma~\ref{lem:scaling-field-domains}. So the property~\eqref{it:renormalisation-group-1-k4} of $E'$, applied to each of these $X$ with the functions $\boldsymbol{\lambda}_L$, $\boldsymbol{\lambda}'_L$, and the relation $G_{\boldsymbol{\lambda}+\boldsymbol{\lambda}'}\circ S_L=S_L\circ G_{\boldsymbol{\lambda}_L+\boldsymbol{\lambda}'_L}$ give
\[
G_{\boldsymbol{\lambda}+\boldsymbol{\lambda}'}\left[(L\,E')(Y;\mathscr{A}^1,\mathscr{C}^1)\right]=\sum_{X:\ \bar{X}^L=LY}S_L\,G_{\boldsymbol{\lambda}_L+\boldsymbol{\lambda}'_L}\left[E'(X;\mathscr{A}^1_L,\mathscr{C}^1_L)\right]=(L\,E')(Y;\mathscr{A},\mathscr{C}).
\]

The properties~\eqref{it:renormalisation-group-1-k5} and~\eqref{it:renormalisation-group-1-k7}. The maps $\mathbf{C}$ and $\mathbf{K}$ commute with $S_L$, whose factors $L^{-1}$ and $L^{-1-\alpha}$ are positive numbers, and
\[
(-\mathscr{A})_L=-\mathscr{A}_L,\qquad\left(\overline{\mathscr{A}}\right)_L=\overline{\mathscr{A}_L},
\]
and the same holds for $\mathscr{C}$. Therefore the properties~\eqref{it:renormalisation-group-1-k5} and~\eqref{it:renormalisation-group-1-k7} of $E'$ at the pairs $(\mathscr{A}_L,\mathscr{C}_L)$ give the same properties of $L\,E'$ at the level $k+1$.

Without the second slot, that is, with $\mathscr{C}=0$ throughout and with $K_k^C$, $K_k$, $\|\cdot\|_k$ in place of $K^{x,C}_k$, $K^x_k$, $\|\cdot\|^x_k$, the same arguments prove the same statements. With Proposition~\ref{prop:properties-actions} and Lemma~\ref{lem:symmetries}, Proposition~\ref{prop:symmetries-extraction} follows.
\end{proof}

\subsection{Averaging the bilinear part over a polymer}\label{sec:averaging-bilinear-part}

The setting~\eqref{eq:polymer-activities-Sb} is assumed in this subsection. Fix $(\mathscr{A},\mathscr{C})\in\tilde{R}_k\times\mathfrak{X}_k$ and put $\mathscr{T}:=\mathscr{A}+\mathscr{C}$. Let $\mathscr{A}=A_0+A_1$ be any decomposition as in Definition~\ref{def:domains}, so that $A_0$ is real and $A_1\in R_k$. Since $A_0$ is real, $\operatorname{Im}\mathscr{T}=\operatorname{Im}A_1+\operatorname{Im}\mathscr{C}$, and the definitions of $R_k$ and of $\mathfrak{X}_k$ give $|A_1|_{\infty}<e_k^{-\frac{3}{4}+\varepsilon}$ and $|\mathscr{C}|_{\infty}<e_k^{-1+2\varepsilon}$. As $e_k<e_{k+1}\le e_0$ by~\eqref{eq:polymer-activities-Sb}, condition~\eqref{eq:renormalisation-group-A8} yields $e_k^{\varepsilon}\le\frac{1}{4}$; thus $e_k<1$, and $e_k^{\frac{1}{4}}\le e_k^{\varepsilon}$ because $0<\varepsilon\le\bar{\varepsilon}=\frac{1}{200}$ (stages~\ref{it:order-choice-1} and~\ref{it:order-choice-8}). Consequently each of the two numbers $e_k^{\frac{1}{4}+\varepsilon}=e_k^{\frac{1}{4}}\cdot e_k^{\varepsilon}$ and $e_k^{2\varepsilon}$ is at most $\frac{1}{4}\,e_k^{\varepsilon}$, and
\begin{equation}\label{eq:polymer-activities-15}
e_k\cdot\left|\operatorname{Im}\mathscr{T}\right|_{\infty}<e_k^{\frac{1}{4}+\varepsilon}+e_k^{2\varepsilon}\le\frac{1}{2}\,e_k^{\varepsilon}.
\end{equation}

\paragraph{Transports along pairs.}
For a complex bond function $B$ on $T_k$ and a pair $(p,p')\in\mathrm{Pz}_k$ let
\[
u_B(p,p'):=\exp\left\{i\,e_k\cdot\eta_k\cdot(\tau B)(p,p')\right\},
\]
where $\tau$ and the six rectilinear paths $\Gamma^{\pi}(p,p')$ are those of Definition~\ref{def:pairs-sites-holder}, and let $u_B(p,p):=1$ for every site $p$ of $T_k$; put $u:=u_{\mathscr{T}}$. Fix $(p,p')\in\mathrm{Pz}_k$. The transport $u(p,p')$ and its inverse satisfy
\begin{equation}\tag{u1}\label{eq:polymer-activities-u1}
\left|u(p,p')^{\pm1}\right|\le\exp\left(e_k\cdot\left|\operatorname{Im}\mathscr{T}\right|_{\infty}\cdot\left|p-p'\right|_1\right).
\end{equation}
Indeed, $\left|u(p,p')^{\pm1}\right|=\exp\left(\mp e_k\cdot\eta_k\cdot(\tau\operatorname{Im}\mathscr{T})(p,p')\right)$; here $\eta_k\cdot(\tau\operatorname{Im}\mathscr{T})(p,p')$ is the average, over the six paths $\Gamma^{\pi}(p,p')$, of the sums of $\eta_k\cdot\operatorname{Im}\mathscr{T}$ along them, and every one of these paths consists of $\left|p-p'\right|_1/\eta_k$ bonds. Written out with the formula for $\tau$ in Definition~\ref{def:pairs-sites-holder},
\begin{equation}\tag{u2}\label{eq:polymer-activities-u2}
u(p,p')=\exp\left\{\frac{i\,e_k}{6}\sum_{\pi}\ \sum_{b\in\Gamma^{\pi}(p,p')}\eta_k\cdot\mathscr{T}(b)\right\};
\end{equation}
thus $u(p,p')$ is the exponential of a linear form in the numbers $\mathscr{T}(b)$, $b$ ranging over the bonds that belong to one of the paths $\Gamma^{\pi}(p,p')$. In particular $u(p,p')$ is entire in these numbers, and it does not change if $\mathscr{T}$ is altered at other bonds.

Next, let $g\in\mathbb{G}_k$, let $\Lambda$ be a real function on the sites of $T_k$, and let $\mathscr{T}^1$ be a complex bond function on $T_k$ such that $\mathscr{T}^1(b)=\mathscr{T}(b)-(d\Lambda)(b)$ for every bond $b$ of every path $\Gamma^{\pi}(p,p')$. Then
\begin{multline}\tag{u3}\label{eq:polymer-activities-u3}
u_{g\mathscr{T}}(gp,gp')=u_{\mathscr{T}}(p,p'),\\
u_{\mathscr{T}^1}(p,p')=u_{\mathscr{T}}(p,p')\cdot\exp\left\{-i\,e_k\cdot\left(\Lambda(p')-\Lambda(p)\right)\right\},\\
u_{-\mathscr{T}}=u_{\mathscr{T}}^{-1},\qquad\overline{\left(u_{-\overline{\mathscr{T}}}\right)}=u_{\mathscr{T}}.
\end{multline}
The first identity is a consequence of Proposition~\ref{prop:properties-actions}(b). For the second, the sum of $\eta_k\cdot d\Lambda$ along each path $\Gamma^{\pi}(p,p')$ telescopes to $\Lambda(p')-\Lambda(p)$. The last two hold because $\tau$ is linear with real coefficients.

Finally, let $\psi(p)$, $\psi(p')$, $\bar{\psi}(p)$, $\bar{\psi}(p')$, $\chi(p,p')$, $\bar{\chi}(p,p')$ be elements of an algebra over $\mathbb{C}$ that are related by
\begin{multline}\tag{u4}\label{eq:polymer-activities-u4}
\chi(p,p')=\left|p-p'\right|^{-\alpha}\cdot\left(u(p,p')\cdot\psi(p')-\psi(p)\right),\\
\bar{\chi}(p,p')=\left|p-p'\right|^{-\alpha}\cdot\left(u(p,p')^{-1}\cdot\bar{\psi}(p')-\bar{\psi}(p)\right);
\end{multline}
with the operators of Definition~\ref{def:pairs-sites-holder} these relations read $\chi(p,p')=\left(\delta_{\alpha,\mathscr{T}}\psi\right)(p,p')$ and $\bar{\chi}(p,p')=\left(\delta_{\alpha,-\mathscr{T}}\bar{\psi}\right)(p,p')$. Since $p\ne p'$, they can be solved for $\psi(p)$ and $\bar{\psi}(p)$:
\begin{multline}\label{eq:polymer-activities-16}
\psi(p)=u(p,p')\cdot\psi(p')-\left|p-p'\right|^{\alpha}\cdot\chi(p,p'),\\
\bar{\psi}(p)=u(p,p')^{-1}\cdot\bar{\psi}(p')-\left|p-p'\right|^{\alpha}\cdot\bar{\chi}(p,p').
\end{multline}

\paragraph{Chains of cubes.}
Let $X$ be a small polymer, so that $d_M(X)\le L$. By Proposition~\ref{prop:polymers}\eqref{it:polymers-t1},
\[
|X|_M\le c_X\cdot\left(d_M(X)+1\right)\le c_X\cdot(L+1).
\]
Every site of $T_k$ lies in exactly one $M$-cube (Section~\ref{sec:balls-torus-lifts}). Let $y$, $y_0$ be sites of $X$. A sequence $(\square_0,\dots,\square_n)$ of $M$-cubes of $X$ is called a chain of length $n$ from $y$ to $y_0$ if $y\in\square_0$, $y_0\in\square_n$, and $\square_{i-1}$, $\square_i$ are adjacent (Section~\ref{sec:polymers}) for $1\le i\le n$. Let $\mathrm{Ch}(y,y_0)$ be the set of the chains from $y$ to $y_0$ whose length is the least possible; its elements are denoted by $\boldsymbol{\chi}$. As $X$ is connected through faces, $\mathrm{Ch}(y,y_0)$ is not empty, and all its elements have the same length $n$. In each of them the cubes are pairwise distinct, since a repeated cube would allow a shorter chain; hence $n+1\le|X|_M$, and $\mathrm{Ch}(y,y_0)$ is finite.

For $\boldsymbol{\chi}\in\mathrm{Ch}(y,y_0)$ define the points
\[
p_0:=y,\qquad p_{i+1}:=\text{the centre of }\square_i\quad(0\le i\le n),\qquad p_{n+2}:=y_0,
\]
all of them taken in one lift of $X$, as in Section~\ref{sec:antiperiodicity}. The centres belong to $M\mathbb{Z}^3\subset\mathbb{L}_k$, so they are sites of $T_k$. For every $i\in\{0,\dots,n+1\}$ with $p_i\ne p_{i+1}$,
\begin{equation}\tag{ch1}\label{eq:polymer-activities-ch1}
(p_i,p_{i+1})\in\mathrm{Pz}_k\quad\text{is a pair admissible for $X$ of kind (I).}
\end{equation}
Indeed, $p_i$ and $p_{i+1}$ are sites of $X$, and $\left|p_i-p_{i+1}\right|\le M$. Along a rectilinear path each coordinate varies monotonically between its values at the two end points; hence the path stays in every product of three intervals that contains both end points. If $i=0$ or $i=n+1$, then $p_i$ and $p_{i+1}$ belong to a single $M$-cube of $X$, which is such a product, so the six paths $\Gamma^{\pi}(p_i,p_{i+1})$ stay in it. If $1\le i\le n$, then $p_i$ and $p_{i+1}$ are the centres of the adjacent cubes $\square_{i-1}$ and $\square_i$, so $p_{i+1}-p_i=\pm M\cdot e_{\mu}$ for a direction $\mu$; every one of the six paths then consists of the lattice sites of the segment from $p_i$ to $p_{i+1}$, and these sites lie in $\square_{i-1}\cup\square_i$.

By~\eqref{eq:polymer-activities-ch1} the transports $u(p_j,p_{j+1})$, $0\le j\le n+1$, are defined, and one may put
\[
U_i:=\prod_{0\le j<i}u(p_j,p_{j+1}),\qquad U^{-}_i:=\prod_{0\le j<i}u(p_j,p_{j+1})^{-1}\qquad(0\le i\le n+2),
\]
an empty product being $1$. With the constant $c'=3c_X$ of Section~\ref{sec:absolute-constants},
\begin{multline}\tag{ch2}\label{eq:polymer-activities-ch2}
\sum_{i=0}^{n+1}\left|p_i-p_{i+1}\right|_1\le\frac{3M}{2}+n\cdot M+\frac{3M}{2}=(n+3)\cdot M\\
\le\left(c_X\cdot(L+1)+2\right)\cdot M\le c'\cdot L\cdot M,\\
\left|U_i\right|,\ \left|U^{-}_i\right|\le\exp\left(\frac{c'\cdot L\cdot M\cdot e_k^{\varepsilon}}{2}\right)\le\exp(1)\qquad(0\le i\le n+2).
\end{multline}
In the chain of inequalities for the sum, the terms with $i=0$ and $i=n+1$ are at most $\frac{3M}{2}$, because a site and the centre of the $M$-cube that contains it differ by at most $\frac{M}{2}$ in each coordinate, and each of the remaining $n$ terms equals $M$; moreover $n+1\le|X|_M\le c_X\cdot(L+1)$, and $c_X+2\le 2c_X\cdot L$ since $c_X\ge1$ (Section~\ref{sec:absolute-constants}) and $L\ge5$ (stage~\ref{it:order-choice-4}). The bound on $\left|U_i\right|$ and $\left|U^{-}_i\right|$ follows from~\eqref{eq:polymer-activities-u1}, \eqref{eq:polymer-activities-15} and the bound on the sum, together with
\[
\frac{1}{2}\,c'\cdot L\cdot M\cdot e_k^{\varepsilon}=\frac{3}{2}\,c_X\cdot L\cdot M\cdot e_k^{\varepsilon}\le\frac{3}{2}\,c_X\cdot L\cdot M\cdot e_0^{\varepsilon}<1,
\]
which holds because condition~\eqref{eq:renormalisation-group-A9} contains the inequality $90\,c_X\cdot L\cdot M\cdot e_0^{\varepsilon}\le1$. Finally, $0<\alpha<1$ (stage~\ref{it:order-choice-2}) gives $\left|p_i-p_{i+1}\right|^{\alpha}\le\max\left(1,\left|p_i-p_{i+1}\right|\right)\le M$ for every $i$, whence
\begin{equation}\tag{ch3}\label{eq:polymer-activities-ch3}
\sum_{i=0}^{n+1}\left|p_i-p_{i+1}\right|^{\alpha}\le(n+2)\cdot M\le c'\cdot L\cdot M.
\end{equation}

For a function $F$ on $\mathrm{Ch}(y,y_0)$ with values in a complex vector space let
\[
\operatorname{Mn} F:=\frac{1}{\#\left|\mathrm{Ch}(y,y_0)\right|}\sum_{\boldsymbol{\chi}\in\mathrm{Ch}(y,y_0)}F(\boldsymbol{\chi})
\]
be its mean. In the sums below, terms with $p_i=p_{i+1}$ are left out, by the convention of Section~\ref{sec:local-parts-adjustment}. For sites $x$, $x_0$, $y$, $y_0$ of $X$ define
\begin{align}
\mathbf{T}(y,y_0)&:=\operatorname{Mn}\,U_{n+2},\notag\\
\operatorname{Rm}(y,y_0)[\chi]&:=\operatorname{Mn}\sum_{i=0}^{n+1}U_i\cdot\left|p_i-p_{i+1}\right|^{\alpha}\cdot\chi(p_i,p_{i+1}),\notag\\
\mathbf{T}^{-}(x,x_0)&:=\operatorname{Mn}\,U^{-}_{n+2},\notag\\
\operatorname{Rm}^{-}(x,x_0)[\bar{\chi}]&:=\operatorname{Mn}\sum_{i=0}^{n+1}U^{-}_i\cdot\left|p_i-p_{i+1}\right|^{\alpha}\cdot\bar{\chi}(p_i,p_{i+1});\label{eq:polymer-activities-17}
\end{align}
in the first two lines the mean is over $\mathrm{Ch}(y,y_0)$, and in the last two it is over $\mathrm{Ch}(x,x_0)$, the points $p_i$ and the products $U^{-}_i$ being formed with $p_0=x$ and $p_{n+2}=x_0$. Here $\chi(p_i,p_{i+1})$ and $\bar{\chi}(p_i,p_{i+1})$ are the columns of generators of Section~\ref{sec:coefficients-kernels-grassmann}, and the spinor index is carried through unchanged: $\operatorname{Rm}(y,y_0)[\chi]$ and $\operatorname{Rm}^{-}(x,x_0)[\bar{\chi}]$ are columns of four elements of $\mathrm{Gr}[\psi,\chi]$, whereas $\mathbf{T}(y,y_0)$ and $\mathbf{T}^{-}(x,x_0)$ are complex numbers.

\begin{definition}[Transport to the mean]\label{def:transport-mean}
Let $X$ be a small polymer. Let $\boldsymbol{\Theta}_{X,\mathscr{T}}$, briefly $\boldsymbol{\Theta}$, be the homomorphism from the subalgebra of $\mathrm{Gr}[\psi,\chi]$ spanned by the products of the generators $\psi(\xi)$, $\xi\in X$, into $\mathrm{Gr}[\psi,\chi]$ that acts on these generators, for sites $x$ and $y$ of $X$, by
\begin{align*}
\boldsymbol{\Theta}\,\psi(y)&:=\operatorname{Vol}(X)^{-1}\cdot\int_X\mathrm{d}y_0\,\left(\mathbf{T}(y,y_0)\cdot\psi(y_0)-\operatorname{Rm}(y,y_0)[\chi]\right),\\
\boldsymbol{\Theta}\,\bar{\psi}(x)&:=\operatorname{Vol}(X)^{-1}\cdot\int_X\mathrm{d}x_0\,\left(\mathbf{T}^{-}(x,x_0)\cdot\bar{\psi}(x_0)-\operatorname{Rm}^{-}(x,x_0)[\bar{\chi}]\right).
\end{align*}
Let $E$ be a family of elements $E(X;\mathscr{A},\mathscr{C})$ of $\mathrm{Gr}[\psi,\chi]$, indexed by $X\in D_k$ and $(\mathscr{A},\mathscr{C})\in\tilde{R}_k\times\mathfrak{X}_k$, that has property~\eqref{it:renormalisation-group-1-k2}, and let $E^{[2,0]}$ denote its part of bidegree $(2,0)$ (Section~\ref{sec:coefficients-kernels-grassmann}). For small $X$ put
\begin{multline*}
E^{\mathrm{nat}}(X;\mathscr{A},\mathscr{C}):=E(X;\mathscr{A},\mathscr{C})-E^{[2,0]}(X;\mathscr{A},\mathscr{C})\\
+\boldsymbol{\Theta}_{X,\mathscr{T}}\left(E^{[2,0]}(X;\mathscr{A},\mathscr{C})\right),
\end{multline*}
where $\mathscr{T}=\mathscr{A}+\mathscr{C}$, and for large $X$ put $E^{\mathrm{nat}}(X;\mathscr{A},\mathscr{C}):=E(X;\mathscr{A},\mathscr{C})$. The map $E\mapsto E^{\mathrm{nat}}$ is called the adjustment.
\end{definition}

The homomorphism $\boldsymbol{\Theta}$ exists and is unique: the prescribed images are linear combinations of generators, hence odd elements, and odd elements anticommute and have square zero. By property~\eqref{it:renormalisation-group-1-k2} only generators $\psi(\xi)$ with $\xi\in X$ occur in $E^{[2,0]}(X;\mathscr{A},\mathscr{C})$, so this element lies in the domain of $\boldsymbol{\Theta}_{X,\mathscr{T}}$.

Let $X$ be small, and let the elements of $E$ have, as in Definition~\ref{def:renormalisation-group-1}, equally many barred and unbarred legs; then~\eqref{eq:polymer-activities-3} applies. Write $E_{20}(X;x,y)$ for the matrix kernel $E_{20}(X;\mathscr{A},\mathscr{C};x,y)$ of~\eqref{eq:polymer-activities-3}, and $E^{[2,0]}(X)$ for $E^{[2,0]}(X;\mathscr{A},\mathscr{C})$. All integrals below extend over $X$. Since $\boldsymbol{\Theta}$ is linear and multiplicative, \eqref{eq:polymer-activities-3} gives
\[
\boldsymbol{\Theta}\left(E^{[2,0]}(X)\right)=\int\mathrm{d}x\,\mathrm{d}y\;\left(\boldsymbol{\Theta}\,\bar{\psi}(x)\right)^T\cdot E_{20}(X;x,y)\cdot\left(\boldsymbol{\Theta}\,\psi(y)\right),
\]
and after insertion of the definition of $\boldsymbol{\Theta}$ the product expands into four terms:
\begin{align*}
&\operatorname{Vol}(X)^2\cdot\boldsymbol{\Theta}\left(E^{[2,0]}(X)\right)\\
&\quad=\int\mathrm{d}x_0\,\mathrm{d}y_0\;\bar{\psi}(x_0)^T\cdot\left[\int\mathrm{d}x\,\mathrm{d}y\;\mathbf{T}^{-}(x,x_0)\cdot E_{20}(X;x,y)\cdot\mathbf{T}(y,y_0)\right]\cdot\psi(y_0)\\
&\qquad-\int\mathrm{d}x\,\mathrm{d}y\,\mathrm{d}x_0\,\mathrm{d}y_0\;\bar{\psi}(x_0)^T\cdot\mathbf{T}^{-}(x,x_0)\cdot E_{20}(X;x,y)\cdot\operatorname{Rm}(y,y_0)[\chi]\\
&\qquad-\int\mathrm{d}x\,\mathrm{d}y\,\mathrm{d}x_0\,\mathrm{d}y_0\;\operatorname{Rm}^{-}(x,x_0)[\bar{\chi}]^T\cdot E_{20}(X;x,y)\cdot\mathbf{T}(y,y_0)\cdot\psi(y_0)\\
&\qquad+\int\mathrm{d}x\,\mathrm{d}y\,\mathrm{d}x_0\,\mathrm{d}y_0\;\operatorname{Rm}^{-}(x,x_0)[\bar{\chi}]^T\cdot E_{20}(X;x,y)\cdot\operatorname{Rm}(y,y_0)[\chi].
\end{align*}

\begin{theorem}[Properties of the adjustment]\label{thm:properties-adjustment}
Assume the setting~\eqref{eq:polymer-activities-Sb}. Let $E\in K^{x,C}_k$, let $X\in D_k$ be a small polymer, and let $(\mathscr{A},\mathscr{C})\in\tilde{R}_k\times\mathfrak{X}_k$, with the total connection $\mathscr{T}=\mathscr{A}+\mathscr{C}$. Let $E^{\mathrm{nat}}$ be the family of Definition~\ref{def:transport-mean}, let $c'=3c_X$ be the constant of Section~\ref{sec:absolute-constants}, and put $x_k:=c'\cdot L\cdot M\cdot e_k^{\varepsilon}$; then $x_k\le\frac{1}{30}$. The norms $\|\cdot\|$ of kernels are those of Section~\ref{sec:grassmann-algebras-norms}.
\begin{enumerate}
\item[(a)] \emph{Equality after insertion.} Let $\mathrm{ev}$ be an algebra homomorphism defined on $\mathrm{Gr}[\psi,\chi]$, with values in some algebra over $\mathbb{C}$. Assume that, at every pair $(p,p')$ of kind~(I) for $X$ (Definition~\ref{def:pairs-sites-holder}), the images under $\mathrm{ev}$ of the generators obey the relations~\eqref{eq:polymer-activities-u4}, entry by entry in the spinor index:
\begin{align*}
\mathrm{ev}(\chi(p,p'))&=|p-p'|^{-\alpha}\cdot\left(u(p,p')\cdot\mathrm{ev}(\psi(p'))-\mathrm{ev}(\psi(p))\right),\\
\mathrm{ev}(\bar{\chi}(p,p'))&=|p-p'|^{-\alpha}\cdot\left(u(p,p')^{-1}\cdot\mathrm{ev}(\bar{\psi}(p'))-\mathrm{ev}(\bar{\psi}(p))\right),
\end{align*}
where $u=u_{\mathscr{T}}$ is the transport of Section~\ref{sec:averaging-bilinear-part}. Then
\[
\mathrm{ev}\left(E^{\mathrm{nat}}(X;\mathscr{A},\mathscr{C})\right)=\mathrm{ev}\left(E(X;\mathscr{A},\mathscr{C})\right).
\]
In particular $E^{\mathrm{nat}}(X;\mathscr{A},\mathscr{C},\psi,\delta_{\alpha,\mathscr{T}}\psi)=E(X;\mathscr{A},\mathscr{C},\psi,\delta_{\alpha,\mathscr{T}}\psi)$, where, as in Definition~\ref{def:pairs-sites-holder}, the barred generators $\bar{\chi}$ are replaced by $\delta_{\alpha,-\mathscr{T}}\bar{\psi}$.
\item[(b)] \emph{Low degrees.} In the matrix notation of~\eqref{eq:polymer-activities-3}, the vanishing of the matrix $\int\mathrm{d}x\,\mathrm{d}y\;E_{20}(X;0,0;x,y)$ implies $E^{\mathrm{nat}}_{20}(X;0,0;\cdot,\cdot)=0$. For every bidegree $(n,m)$ outside $\{(2,0),(1,1),(0,2)\}$ the kernels $E^{\mathrm{nat}}_{nm}$ and $E_{nm}$ coincide; this includes $E^{\mathrm{nat}}_{00}=E_{00}$.
\item[(c)] \emph{Locality.} The element $E^{\mathrm{nat}}(X;\mathscr{A},\mathscr{C})$ contains a generator $\psi(\xi)$ only if $\xi$ lies in $X$, and a generator $\chi(\boldsymbol{\zeta})$ only if $\boldsymbol{\zeta}$ is admissible for $X$. If a generator $\chi(\boldsymbol{\zeta})$, $\boldsymbol{\zeta}=(p,p',\boldsymbol{\beta},\boldsymbol{\omega})$, occurs in a monomial with non-zero coefficient of $E^{\mathrm{nat}}(X;\mathscr{A},\mathscr{C})$ (Section~\ref{sec:coefficients-kernels-grassmann}) and in no such monomial of $E(X;\mathscr{A},\mathscr{C})$, then $(p,p')=(p_i,p_{i+1})$ with $p_i\ne p_{i+1}$ for some sites $y,y_0$ of $X$, some chain $\boldsymbol{\chi}\in\mathrm{Ch}(y,y_0)$ and some index $i$ (Section~\ref{sec:averaging-bilinear-part}); in particular the pair $(p,p')$ is admissible for $X$ of kind~(I), by~\eqref{eq:polymer-activities-ch1}. The coefficients of $E^{\mathrm{nat}}(X;\mathscr{A},\mathscr{C})$ depend on $(\mathscr{A},\mathscr{C})$ only through the kernels $E_{nm}(X;\mathscr{A},\mathscr{C})$ and through the values of $\mathscr{T}$ on the bonds whose two end points lie in $X$. Consequently $E^{\mathrm{nat}}$ has the property~\eqref{it:renormalisation-group-1-k2} of Definition~\ref{def:renormalisation-group-1}.
\item[(d)] \emph{Bounds.} With all kernels taken at $(\mathscr{A},\mathscr{C})$,
\begin{align*}
\|E^{\mathrm{nat}}_{20}(X)\|&\le\exp(2)\cdot\|E_{20}(X)\|,\\
\|E^{\mathrm{nat}}_{11}(X)\|&\le\|E_{11}(X)\|+\exp(2)\cdot c'\cdot L\cdot M\cdot\|E_{20}(X)\|,\\
\|E^{\mathrm{nat}}_{02}(X)\|&\le\|E_{02}(X)\|+\exp(2)\cdot(c'\cdot L\cdot M)^2\cdot\|E_{20}(X)\|,
\end{align*}
and
\begin{align*}
\|E^{\mathrm{nat}}(X;\mathscr{A},\mathscr{C})\|_{h_k}&\le\exp(2)\cdot(1+x_k)^2\cdot\|E(X;\mathscr{A},\mathscr{C})\|_{h_k}\\
&\le10\cdot\|E(X;\mathscr{A},\mathscr{C})\|_{h_k}.
\end{align*}
\item[(e)] \emph{Symmetries.} The family $E^{\mathrm{nat}}$ belongs to $K^{x,C}_k$, and to $K^x_k$ if $E$ does, and $\|E^{\mathrm{nat}}\|^x_k\le10\cdot\|E\|^x_k$. The map $E\mapsto E^{\mathrm{nat}}$ is linear, and $E^{\mathrm{nat}}|_0=(E|_0)^{\mathrm{nat}}$.
\end{enumerate}
Assertions~(a)--(d), and those of~(e) except the identity $E^{\mathrm{nat}}|_0=(E|_0)^{\mathrm{nat}}$, remain true for the families without the second slot, in the following form. Let $E\in K^C_k$ and $\mathscr{A}\in\tilde{R}_k$, put $\mathscr{T}:=\mathscr{A}$ and $u:=u_{\mathscr{A}}$, and let $E^{\mathrm{nat}}$ be the family given by Definition~\ref{def:transport-mean} with the total connection $\mathscr{A}$, that is, $E^{\mathrm{nat}}(Y;\mathscr{A})=E(Y;\mathscr{A})$ for large $Y\in D_k$ and
\[
E^{\mathrm{nat}}(Y;\mathscr{A})=E(Y;\mathscr{A})-E^{[2,0]}(Y;\mathscr{A})+\boldsymbol{\Theta}_{Y,\mathscr{A}}\left(E^{[2,0]}(Y;\mathscr{A})\right)
\]
for small $Y\in D_k$. Then these assertions hold for the small polymer $X$, with $\mathscr{A}$, $E(X;\mathscr{A})$, $E_{nm}(X;\mathscr{A})$, $E_{20}(X;0;x,y)$, $K^C_k$, $K_k$ and $\|\cdot\|_k$ in place of $(\mathscr{A},\mathscr{C})$, $E(X;\mathscr{A},\mathscr{C})$, $E_{nm}(X;\mathscr{A},\mathscr{C})$, $E_{20}(X;0,0;x,y)$, $K^{x,C}_k$, $K^x_k$ and $\|\cdot\|^x_k$, and with the same replacements for $E^{\mathrm{nat}}$.
\end{theorem}

\begin{proof}
For sites $y,y_0$ of $X$ and a chain $\boldsymbol{\chi}\in\mathrm{Ch}(y,y_0)$, the points $p_0=y,p_1,\dots,p_{n+2}=y_0$ and the products $U_i$, $U^{-}_i$ are those of Section~\ref{sec:averaging-bilinear-part}; all points are taken in one lift of $X$ (Section~\ref{sec:antiperiodicity}). For a $4\times4$ matrix $N$, the expression $\operatorname{Rm}(y,y_0)[N\cdot\chi]$ is the expression~\eqref{eq:polymer-activities-17} with $N\cdot\chi(p_i,p_{i+1})$ in place of $\chi(p_i,p_{i+1})$, and $\operatorname{Rm}^{-}(y,y_0)[N\cdot\bar{\chi}]$ is defined in the same way; both are linear in the argument in brackets. Since no site lies on the boundary of an $M$-cube (Section~\ref{sec:balls-torus-lifts}), the polymer $X$ has exactly $\operatorname{Vol}(X)\cdot\eta^{-3}$ sites, and therefore
\[
\int_X\mathrm{d}y_0\;1=\operatorname{Vol}(X).
\]

\emph{(a)} Fix sites $y,y_0$ of $X$ and a chain $\boldsymbol{\chi}\in\mathrm{Ch}(y,y_0)$, and let $0\le i\le n+1$. If $p_i\ne p_{i+1}$, the pair $(p_i,p_{i+1})$ is of kind~(I) for $X$ by~\eqref{eq:polymer-activities-ch1}; by hypothesis the images satisfy the relations~\eqref{eq:polymer-activities-u4} at this pair, and hence~\eqref{eq:polymer-activities-16}:
\[
\mathrm{ev}(\psi(p_i))=u(p_i,p_{i+1})\cdot\mathrm{ev}(\psi(p_{i+1}))-|p_i-p_{i+1}|^{\alpha}\cdot\mathrm{ev}(\chi(p_i,p_{i+1})).
\]
If $p_i=p_{i+1}$, then $u(p_i,p_{i+1})=1$, and the same identity holds without its last term, which is omitted, as in~\eqref{eq:polymer-activities-17}. As $U_0=1$ and $U_i\cdot u(p_i,p_{i+1})=U_{i+1}$, induction on $i$ yields, for $0\le i\le n+2$,
\[
\mathrm{ev}(\psi(y))=U_i\cdot\mathrm{ev}(\psi(p_i))-\sum_{j<i}U_j\cdot|p_j-p_{j+1}|^{\alpha}\cdot\mathrm{ev}(\chi(p_j,p_{j+1})).
\]
For $i=n+2$ the point $p_i$ is $y_0$. The map $\mathrm{ev}$ is linear; the mean over $\boldsymbol{\chi}\in\mathrm{Ch}(y,y_0)$ therefore gives, by~\eqref{eq:polymer-activities-17},
\[
\mathrm{ev}(\psi(y))=\mathbf{T}(y,y_0)\cdot\mathrm{ev}(\psi(y_0))-\mathrm{ev}\left(\operatorname{Rm}(y,y_0)[\chi]\right).
\]
Apply $\operatorname{Vol}(X)^{-1}\int_X\mathrm{d}y_0$ to both sides. Since $\int_X\mathrm{d}y_0\,1=\operatorname{Vol}(X)$, Definition~\ref{def:transport-mean} gives $\mathrm{ev}(\boldsymbol{\Theta}_{X,\mathscr{T}}\psi(y))=\mathrm{ev}(\psi(y))$. The second relation of~\eqref{eq:polymer-activities-16}, with $u^{-1}$, $U^{-}_i$, $\mathbf{T}^{-}$ and $\operatorname{Rm}^{-}$ in place of $u$, $U_i$, $\mathbf{T}$ and $\operatorname{Rm}$, gives in the same way $\mathrm{ev}(\boldsymbol{\Theta}_{X,\mathscr{T}}\bar{\psi}(x))=\mathrm{ev}(\bar{\psi}(x))$ for every site $x$ of $X$. Thus the homomorphisms $\mathrm{ev}\circ\boldsymbol{\Theta}_{X,\mathscr{T}}$ and $\mathrm{ev}$ coincide on the generators of the domain of $\boldsymbol{\Theta}_{X,\mathscr{T}}$, hence on the whole domain. By the property~\eqref{it:renormalisation-group-1-k2} of $E$, the element $E^{[2,0]}(X;\mathscr{A},\mathscr{C})$ contains generators $\psi(\xi)$ with $\xi\in X$ only; so it belongs to this domain, and Definition~\ref{def:transport-mean} yields
\begin{align*}
\mathrm{ev}\left(E^{\mathrm{nat}}(X;\mathscr{A},\mathscr{C})\right)&=\mathrm{ev}\left(E(X;\mathscr{A},\mathscr{C})\right)-\mathrm{ev}\left(E^{[2,0]}(X;\mathscr{A},\mathscr{C})\right)\\
&\qquad+\mathrm{ev}\left(\boldsymbol{\Theta}_{X,\mathscr{T}}\,E^{[2,0]}(X;\mathscr{A},\mathscr{C})\right)\\
&=\mathrm{ev}\left(E(X;\mathscr{A},\mathscr{C})\right).
\end{align*}
For the last assertion of~(a), let $\mathrm{ev}$ be the homomorphism that fixes $\psi$ and $\bar{\psi}$ and replaces $\chi$ by $\delta_{\alpha,\mathscr{T}}\psi$ and $\bar{\chi}$ by $\delta_{\alpha,-\mathscr{T}}\bar{\psi}$ (it exists, since the inserted elements are linear combinations of the generators $\psi$, $\bar{\psi}$, hence odd). By Definition~\ref{def:pairs-sites-holder}, and since $u_{-\mathscr{T}}=u_{\mathscr{T}}^{-1}$ by~\eqref{eq:polymer-activities-u3}, the images satisfy the relations of~\eqref{eq:polymer-activities-u4} at every pair of $\mathrm{Pz}_k$, in particular at every pair of kind~(I) for $X$.

\emph{(b)} Let $(\mathscr{A},\mathscr{C})=(0,0)$. Then $\mathscr{T}=0$, every transport equals $1$, and so do all $U_i$, $U^{-}_i$, $\mathbf{T}$ and $\mathbf{T}^{-}$. Put
\[
\psi_{\mathrm{av}}:=\operatorname{Vol}(X)^{-1}\int_X\mathrm{d}y_0\;\psi(y_0),\qquad\bar{\psi}_{\mathrm{av}}:=\operatorname{Vol}(X)^{-1}\int_X\mathrm{d}x_0\;\bar{\psi}(x_0).
\]
By Definition~\ref{def:transport-mean}, for sites $x,y$ of $X$,
\begin{align*}
\boldsymbol{\Theta}_{X,0}\,\psi(y)&=\psi_{\mathrm{av}}-\operatorname{Vol}(X)^{-1}\int_X\mathrm{d}y_0\;\operatorname{Rm}(y,y_0)[\chi],\\
\boldsymbol{\Theta}_{X,0}\,\bar{\psi}(x)&=\bar{\psi}_{\mathrm{av}}-\operatorname{Vol}(X)^{-1}\int_X\mathrm{d}x_0\;\operatorname{Rm}^{-}(x,x_0)[\bar{\chi}],
\end{align*}
and the second terms have the bidegree $(0,1)$. By~\eqref{eq:polymer-activities-3}, $E^{[2,0]}(X;0,0)=\int\mathrm{d}x\,\mathrm{d}y\;\bar{\psi}(x)^T\cdot E_{20}(X;0,0;x,y)\cdot\psi(y)$. Since $\boldsymbol{\Theta}_{X,0}$ is a homomorphism, the part of bidegree $(2,0)$ of $\boldsymbol{\Theta}_{X,0}\left(E^{[2,0]}(X;0,0)\right)$ is
\[
\bar{\psi}_{\mathrm{av}}^T\cdot\left(\int\mathrm{d}x\,\mathrm{d}y\;E_{20}(X;0,0;x,y)\right)\cdot\psi_{\mathrm{av}},
\]
and it vanishes under the hypothesis of~(b). In Definition~\ref{def:transport-mean} the parts of bidegree $(2,0)$ of $E(X;0,0)$ and of $E^{[2,0]}(X;0,0)$ cancel; hence the part of bidegree $(2,0)$ of $E^{\mathrm{nat}}(X;0,0)$ is zero, and so is its antisymmetric kernel $E^{\mathrm{nat}}_{20}(X;0,0;\cdot,\cdot)$, by Proposition~\ref{prop:weighted-coefficient-norms}\eqref{it:weighted-coefficient-norms-N0}.

At every $(\mathscr{A},\mathscr{C})$ the homomorphism $\boldsymbol{\Theta}_{X,\mathscr{T}}$ maps each generator of its domain to the sum of an element of bidegree $(1,0)$ and an element of bidegree $(0,1)$. As $E^{[2,0]}(X;\mathscr{A},\mathscr{C})$ is a linear combination of products of two generators, the part of bidegree $(n,m)$ of $\boldsymbol{\Theta}_{X,\mathscr{T}}\left(E^{[2,0]}(X;\mathscr{A},\mathscr{C})\right)$ vanishes unless $(n,m)$ is one of $(2,0)$, $(1,1)$, $(0,2)$. Consequently $E^{\mathrm{nat}}(X;\mathscr{A},\mathscr{C})$ and $E(X;\mathscr{A},\mathscr{C})$ have the same part of every other bidegree and, by Proposition~\ref{prop:weighted-coefficient-norms}\eqref{it:weighted-coefficient-norms-N0}, the same antisymmetric kernels of these bidegrees.

\emph{(c)} By Definition~\ref{def:transport-mean},
\[
E^{\mathrm{nat}}(X;\mathscr{A},\mathscr{C})=\left(E-E^{[2,0]}\right)(X;\mathscr{A},\mathscr{C})+\boldsymbol{\Theta}_{X,\mathscr{T}}\left(E^{[2,0]}(X;\mathscr{A},\mathscr{C})\right).
\]
The first summand consists of terms of $E(X;\mathscr{A},\mathscr{C})$: its coefficients are given by the kernels $E_{nm}(X;\mathscr{A},\mathscr{C})$, and its generators and its dependence on $(\mathscr{A},\mathscr{C})$ are restricted by the property~\eqref{it:renormalisation-group-1-k2} of $E$. In the second summand, the elements $\boldsymbol{\Theta}_{X,\mathscr{T}}\psi(y)$ and $\boldsymbol{\Theta}_{X,\mathscr{T}}\bar{\psi}(x)$ contain the generators $\psi(y_0)$ and $\bar{\psi}(x_0)$ with $x_0,y_0\in X$, and the generators $\chi$ and $\bar{\chi}$ at the pairs $(p_i,p_{i+1})$, $p_i\ne p_{i+1}$, of the chains; these pairs are of kind~(I) for $X$ by~\eqref{eq:polymer-activities-ch1}. In $\boldsymbol{\Theta}_{X,\mathscr{T}}\psi(y)$ and in $\boldsymbol{\Theta}_{X,\mathscr{T}}\bar{\psi}(x)$ the coefficient of each of these generators is a finite sum, extended over the sites $y_0$, respectively $x_0$, over the chains and over the indices $i$ that contribute; each term of this sum is a number independent of $(\mathscr{A},\mathscr{C})$ multiplied by a product of transports $u(p_j,p_{j+1})^{\pm1}$ with $p_j\ne p_{j+1}$, since the transports with $p_j=p_{j+1}$ are equal to $1$. By~\eqref{eq:polymer-activities-u2}, each such transport $u(p_j,p_{j+1})$ depends on $\mathscr{T}$ only through its values on the bonds of the six paths $\Gamma^{\pi}(p_j,p_{j+1})$, and these paths lie in $X$ by~\eqref{eq:polymer-activities-ch1}. The kernels of the second summand are formed from these coefficients and from the kernel $E_{20}(X;\mathscr{A},\mathscr{C})$. By the property~\eqref{it:renormalisation-group-1-k2} of $E$, the element $E^{[2,0]}(X;\mathscr{A},\mathscr{C})$ is a linear combination of products of two generators $\psi(\xi)$ with $\xi\in X$, and $\boldsymbol{\Theta}_{X,\mathscr{T}}$ is a homomorphism; so the second summand is a linear combination of products of two generators, each of which occurs in one of the elements $\boldsymbol{\Theta}_{X,\mathscr{T}}\psi(y)$, $\boldsymbol{\Theta}_{X,\mathscr{T}}\bar{\psi}(x)$. Hence every generator $\chi(\boldsymbol{\zeta})$ that occurs in a monomial with non-zero coefficient of the second summand has as its pair of sites a pair $(p_i,p_{i+1})$, $p_i\ne p_{i+1}$, of a chain. The coefficient of a monomial in $E^{\mathrm{nat}}(X;\mathscr{A},\mathscr{C})$ is the sum of its coefficients in the two summands; so a monomial with non-zero coefficient in $E^{\mathrm{nat}}(X;\mathscr{A},\mathscr{C})$ has a non-zero coefficient in the first or in the second summand. In the first case it is a monomial of $E(X;\mathscr{A},\mathscr{C})$ with the same coefficient, since the first summand is the sum of the parts of $E(X;\mathscr{A},\mathscr{C})$ of the bidegrees other than $(2,0)$, and all its generators occur in $E(X;\mathscr{A},\mathscr{C})$. Consequently, if a generator $\chi(\boldsymbol{\zeta})$ occurs in a monomial with non-zero coefficient of $E^{\mathrm{nat}}(X;\mathscr{A},\mathscr{C})$ and in no such monomial of $E(X;\mathscr{A},\mathscr{C})$, then this monomial has a non-zero coefficient in the second summand, and the pair of sites of $\chi(\boldsymbol{\zeta})$ is a pair $(p_i,p_{i+1})$, $p_i\ne p_{i+1}$, of a chain, which is of kind~(I) for $X$ by~\eqref{eq:polymer-activities-ch1}. This proves the assertions of~(c) on the generators, on the new legs and on the coefficients. Since a bond whose two end points lie in $X$ meets $X$, and since $\mathscr{T}=\mathscr{A}+\mathscr{C}$, the property~\eqref{it:renormalisation-group-1-k2} of $E^{\mathrm{nat}}$ follows.

\emph{(d)} Let $h=(h_1,h_2)$ be positive weights, $h_1$ for the generators $\psi$ and $h_2$ for the generators $\chi$, and let $\|\cdot\|_h$ be the norm~\eqref{eq:polymer-activities-1}; a single generator $\psi$ has the norm $h_1$, a single generator $\chi$ the norm $h_2$ (Section~\ref{sec:coefficients-kernels-grassmann}). Let $y$ be a site of $X$ and $\boldsymbol{\beta}$ a spinor index. The entry $\boldsymbol{\beta}$ of $\boldsymbol{\Theta}_{X,\mathscr{T}}\psi(y)$ is $P_{\boldsymbol{\beta}}(y)+Q_{\boldsymbol{\beta}}(y)$, where
\begin{align*}
P_{\boldsymbol{\beta}}(y)&:=\operatorname{Vol}(X)^{-1}\int_X\mathrm{d}y_0\;\mathbf{T}(y,y_0)\cdot\psi_{\boldsymbol{\beta}}(y_0),\\
Q_{\boldsymbol{\beta}}(y)&:=-\operatorname{Vol}(X)^{-1}\int_X\mathrm{d}y_0\;\operatorname{Rm}(y,y_0)[\chi]_{\boldsymbol{\beta}}
\end{align*}
have the bidegrees $(1,0)$ and $(0,1)$. By~\eqref{eq:polymer-activities-ch2}, $|U_i|\le\exp(1)$ for every $i$; hence $|\mathbf{T}(y,y_0)|\le\exp(1)$, and $\|P_{\boldsymbol{\beta}}(y)\|_h\le\exp(1)\cdot h_1$ by~\eqref{eq:polymer-activities-1} and $\int_X\mathrm{d}y_0\,1=\operatorname{Vol}(X)$. By the triangle inequality, \eqref{eq:polymer-activities-ch2} and~\eqref{eq:polymer-activities-ch3},
\begin{align*}
\|Q_{\boldsymbol{\beta}}(y)\|_h&\le\operatorname{Vol}(X)^{-1}\int_X\mathrm{d}y_0\;\operatorname{Mn}\sum_{i=0}^{n+1}|U_i|\cdot|p_i-p_{i+1}|^{\alpha}\cdot h_2\\
&\le\exp(1)\cdot c'\cdot L\cdot M\cdot h_2.
\end{align*}
In the same way, with $U^{-}_i$ in place of $U_i$, for every site $x$ of $X$ and every spinor index $\boldsymbol{\alpha}$ the entry $\boldsymbol{\alpha}$ of $\boldsymbol{\Theta}_{X,\mathscr{T}}\bar{\psi}(x)$ is $P'_{\boldsymbol{\alpha}}(x)+Q'_{\boldsymbol{\alpha}}(x)$, with $P'_{\boldsymbol{\alpha}}(x)$ of bidegree $(1,0)$, $Q'_{\boldsymbol{\alpha}}(x)$ of bidegree $(0,1)$, $\|P'_{\boldsymbol{\alpha}}(x)\|_h\le\exp(1)\cdot h_1$ and $\|Q'_{\boldsymbol{\alpha}}(x)\|_h\le\exp(1)\cdot c'\cdot L\cdot M\cdot h_2$. By~\eqref{eq:polymer-activities-3}, and since $\boldsymbol{\Theta}_{X,\mathscr{T}}$ is a homomorphism,
\begin{align*}
&\boldsymbol{\Theta}_{X,\mathscr{T}}\left(E^{[2,0]}(X)\right)\\
&\quad=\sum_{\boldsymbol{\alpha},\boldsymbol{\beta}}\int\mathrm{d}x\,\mathrm{d}y\;E_{20}(X;x,y)_{\boldsymbol{\alpha}\boldsymbol{\beta}}\cdot\left(P'_{\boldsymbol{\alpha}}(x)+Q'_{\boldsymbol{\alpha}}(x)\right)\cdot\left(P_{\boldsymbol{\beta}}(y)+Q_{\boldsymbol{\beta}}(y)\right),
\end{align*}
where only sites $x,y$ of $X$ contribute, by~\eqref{it:renormalisation-group-1-k2}. Put
\[
s:=\sum_{\boldsymbol{\alpha},\boldsymbol{\beta}}\int\mathrm{d}x\,\mathrm{d}y\;\left|E_{20}(X;x,y)_{\boldsymbol{\alpha}\boldsymbol{\beta}}\right|;
\]
by~\eqref{eq:polymer-activities-3}, $s=\frac{1}{2}\|E_{20}(X)\|$ and $\|E^{[2,0]}(X)\|_h=h_1^2\cdot s$. Write $F:=\boldsymbol{\Theta}_{X,\mathscr{T}}\left(E^{[2,0]}(X)\right)$, and $F^{[n,m]}$ for its part of bidegree $(n,m)$ (Section~\ref{sec:coefficients-kernels-grassmann}). In the expansion above, the products $P'_{\boldsymbol{\alpha}}(x)\cdot P_{\boldsymbol{\beta}}(y)$ are of bidegree $(2,0)$, the products $P'_{\boldsymbol{\alpha}}(x)\cdot Q_{\boldsymbol{\beta}}(y)$ and $Q'_{\boldsymbol{\alpha}}(x)\cdot P_{\boldsymbol{\beta}}(y)$ of bidegree $(1,1)$, and the products $Q'_{\boldsymbol{\alpha}}(x)\cdot Q_{\boldsymbol{\beta}}(y)$ of bidegree $(0,2)$. The triangle inequality, Proposition~\ref{prop:weighted-coefficient-norms}\eqref{it:weighted-coefficient-norms-N1} and the bounds on $P_{\boldsymbol{\beta}}$, $Q_{\boldsymbol{\beta}}$, $P'_{\boldsymbol{\alpha}}$, $Q'_{\boldsymbol{\alpha}}$ therefore give
\begin{align*}
\|F^{[2,0]}\|_h&\le\exp(2)\cdot h_1^2\cdot s,\\
\|F^{[1,1]}\|_h&\le2\exp(2)\cdot c'\cdot L\cdot M\cdot h_1\cdot h_2\cdot s,\\
\|F^{[0,2]}\|_h&\le\exp(2)\cdot(c'\cdot L\cdot M)^2\cdot h_2^2\cdot s.
\end{align*}
Proposition~\ref{prop:weighted-coefficient-norms}\eqref{it:weighted-coefficient-norms-N0}, in its case of equality, expresses the norm of an element of pure bidegree $(n,m)$ through its antisymmetric kernel $K$: this norm is $\frac{h_1^n\cdot h_2^m}{n!\,m!}\cdot\|K\|$. Let $K_{20}$, $K_{11}$ and $K_{02}$ be the antisymmetric kernels of $F^{[2,0]}$, $F^{[1,1]}$ and $F^{[0,2]}$. Since $2s=\|E_{20}(X)\|$, it follows that
\begin{align*}
\|K_{20}\|&\le\exp(2)\cdot\|E_{20}(X)\|,\\
\|K_{11}\|&\le\exp(2)\cdot c'\cdot L\cdot M\cdot\|E_{20}(X)\|,\\
\|K_{02}\|&\le\exp(2)\cdot(c'\cdot L\cdot M)^2\cdot\|E_{20}(X)\|.
\end{align*}
By Definition~\ref{def:transport-mean} and Proposition~\ref{prop:weighted-coefficient-norms}\eqref{it:weighted-coefficient-norms-N0}, $E^{\mathrm{nat}}_{20}(X)=K_{20}$, $E^{\mathrm{nat}}_{11}(X)=E_{11}(X)+K_{11}$ and $E^{\mathrm{nat}}_{02}(X)=E_{02}(X)+K_{02}$; this proves the three bounds on the kernels.

Now let $h=h_k=(h_{k,1},h_{k,2})=\left(e_k^{-\frac{1}{4}},e_k^{-\frac{1}{4}+\varepsilon}\right)$ (Section~\ref{sec:grassmann-algebras-norms}). Then $h_{k,2}=e_k^{\varepsilon}\cdot h_{k,1}$, that is, $c'\cdot L\cdot M\cdot h_{k,2}=x_k\cdot h_{k,1}$. The norm $\|F\|_{h_k}$ is at most the sum of the three norms above, hence at most
\begin{align*}
\exp(2)\cdot\left(h_{k,1}+c'\cdot L\cdot M\cdot h_{k,2}\right)^2\cdot s&=\exp(2)\cdot(1+x_k)^2\cdot h_{k,1}^2\cdot s\\
&=\exp(2)\cdot(1+x_k)^2\cdot\|E^{[2,0]}(X)\|_{h_k}.
\end{align*}
The norm~\eqref{eq:polymer-activities-1} is additive over the bidegrees (Section~\ref{sec:coefficients-kernels-grassmann}); so $\|(E-E^{[2,0]})(X)\|_{h_k}=\|E(X)\|_{h_k}-\|E^{[2,0]}(X)\|_{h_k}$, and Definition~\ref{def:transport-mean} and the triangle inequality give
\begin{align*}
\|E^{\mathrm{nat}}(X)\|_{h_k}&\le\|E(X)\|_{h_k}-\|E^{[2,0]}(X)\|_{h_k}\\
&\qquad+\exp(2)\cdot(1+x_k)^2\cdot\|E^{[2,0]}(X)\|_{h_k}\\
&\le\exp(2)\cdot(1+x_k)^2\cdot\|E(X)\|_{h_k},
\end{align*}
because $\exp(2)\cdot(1+x_k)^2\ge1$ and $\|E^{[2,0]}(X)\|_{h_k}\le\|E(X)\|_{h_k}$. Finally, the setting~\eqref{eq:polymer-activities-Sb} gives $e_k=L^{-\frac{1}{2}}e_{k+1}\le e_0$, and the condition $90\,c_X\cdot L\cdot M\cdot e_0^{\varepsilon}\le1$ of~\eqref{eq:renormalisation-group-A9} gives $x_k=3c_X\cdot L\cdot M\cdot e_k^{\varepsilon}\le\frac{1}{30}$; moreover $\exp(2)\cdot\left(\frac{31}{30}\right)^2<7.9<10$.

\emph{(e)} If $X$ is large, then $E^{\mathrm{nat}}(X)=E(X)$ for all connections, by Definition~\ref{def:transport-mean}; and every $g\in\mathbb{G}_k$ maps small polymers onto small polymers and large polymers onto large ones (Proposition~\ref{prop:properties-actions}(a)). So the clauses of Definition~\ref{def:renormalisation-group-1} have to be verified at small polymers only, and $X$ is small in what follows.

Each kernel of $E^{\mathrm{nat}}(X;\mathscr{A},\mathscr{C})$ arises by finitely many additions and multiplications from three kinds of ingredients: the kernels of $E(X;\mathscr{A},\mathscr{C})$, analytic on $\tilde{R}_k\times\mathfrak{X}_k$ by the property~\eqref{it:renormalisation-group-1-k1} of $E$; the transports $u^{\pm1}$, entire in $(\mathscr{A},\mathscr{C})$ by~\eqref{eq:polymer-activities-u2}; and numbers independent of $(\mathscr{A},\mathscr{C})$. Consequently the kernels of $E^{\mathrm{nat}}(X;\mathscr{A},\mathscr{C})$ are analytic on $\tilde{R}_k\times\mathfrak{X}_k$, which is the property~\eqref{it:renormalisation-group-1-k1} of $E^{\mathrm{nat}}$. The property~\eqref{it:renormalisation-group-1-k2} is part~(c). The homomorphism $\boldsymbol{\Theta}_{X,\mathscr{T}}$ maps unbarred generators to linear combinations of unbarred generators and barred generators to linear combinations of barred ones; therefore $E^{\mathrm{nat}}(X;\mathscr{A},\mathscr{C})$ is even and has as many barred legs as unbarred legs. Part~(d) gives the property~\eqref{it:renormalisation-group-1-k6} together with $\|E^{\mathrm{nat}}\|^x_k\le10\cdot\|E\|^x_k$.

The properties~\eqref{it:renormalisation-group-1-k3}, \eqref{it:renormalisation-group-1-k4}, \eqref{it:renormalisation-group-1-k5} and~\eqref{it:renormalisation-group-1-k7} are used in the form of Section~\ref{sec:lattice-gauge-charge}. Let $g=(S,a)\in\mathbb{G}_k$, let $\Lambda$ be a real function on the sites of $T_k$, and let $\mathscr{T}^1$ be a complex bond function with $\mathscr{T}^1=\mathscr{T}-d\Lambda$ on the bonds whose two end points lie in $X$. For a small polymer $Y$, such as $X$ or $gX$ (Proposition~\ref{prop:properties-actions}(a)), and a complex bond function $B$ on $T_k$, let $\boldsymbol{\Theta}_{Y,B}$ denote the homomorphism given by Definition~\ref{def:transport-mean} and~\eqref{eq:polymer-activities-17} with $Y$ in place of $X$ and with the products $U_i$, $U^{-}_i$ formed with the transports $u_B$ in place of $u=u_{\mathscr{T}}$; the prescribed images of the generators are again odd, so that this homomorphism exists and is unique, and for $Y=X$, $B=\mathscr{T}$ it is $\boldsymbol{\Theta}_{X,\mathscr{T}}$. It suffices to prove
\begin{multline}\label{eq:polymer-activities-18}
a_g\circ\boldsymbol{\Theta}_{gX,g\mathscr{T}}=\boldsymbol{\Theta}_{X,\mathscr{T}}\circ a_g,\qquad G_{\Lambda}\circ\boldsymbol{\Theta}_{X,\mathscr{T}^1}=\boldsymbol{\Theta}_{X,\mathscr{T}}\circ G_{\Lambda},\\
\mathbf{C}\circ\boldsymbol{\Theta}_{X,-\mathscr{T}}=\boldsymbol{\Theta}_{X,\mathscr{T}}\circ\mathbf{C},\qquad\mathbf{K}\circ\boldsymbol{\Theta}_{X,-\overline{\mathscr{T}}}=\boldsymbol{\Theta}_{X,\mathscr{T}}\circ\mathbf{K}.
\end{multline}
Here the first relation is an identity of maps on the domain of $\boldsymbol{\Theta}_{gX,g\mathscr{T}}$, which $a_g$ maps into the domain of $\boldsymbol{\Theta}_{X,\mathscr{T}}$; the other three are identities of maps on the domain of $\boldsymbol{\Theta}_{X,\mathscr{T}}$, which does not depend on the connection and which $G_{\Lambda}$, $\mathbf{C}$ and $\mathbf{K}$ map into itself (Definition~\ref{def:transport-mean}, Section~\ref{sec:lattice-gauge-charge}).

Indeed, assume~\eqref{eq:polymer-activities-18}. The polymer $gX$ is small, the pair $(g\mathscr{A},g\mathscr{C})$ lies in $\tilde{R}_k\times\mathfrak{X}_k$, and its total connection is $g\mathscr{T}$ (Proposition~\ref{prop:properties-actions}(a),(c)). The map $a_g$ preserves the bidegree by~\eqref{eq:polymer-activities-4}; so the property~\eqref{it:renormalisation-group-1-k3} of $E$ gives $a_g\left[E(gX;g\mathscr{A},g\mathscr{C})\right]=E(X;\mathscr{A},\mathscr{C})$ and $a_g\left[E^{[2,0]}(gX;g\mathscr{A},g\mathscr{C})\right]=E^{[2,0]}(X;\mathscr{A},\mathscr{C})$. With Definition~\ref{def:transport-mean} and the first relation of~\eqref{eq:polymer-activities-18} it follows that
\begin{align*}
a_g\left[E^{\mathrm{nat}}(gX;g\mathscr{A},g\mathscr{C})\right]&=\left(E-E^{[2,0]}\right)(X;\mathscr{A},\mathscr{C})\\
&\qquad+\boldsymbol{\Theta}_{X,\mathscr{T}}\left(a_g\left[E^{[2,0]}(gX;g\mathscr{A},g\mathscr{C})\right]\right)\\
&=E^{\mathrm{nat}}(X;\mathscr{A},\mathscr{C}),
\end{align*}
which is the property~\eqref{it:renormalisation-group-1-k3} of $E^{\mathrm{nat}}$. The second relation gives the property~\eqref{it:renormalisation-group-1-k4} in the same way. Indeed, let $(\mathscr{A},\mathscr{C})$ and $(\mathscr{A}^1,\mathscr{C}^1)$ lie in $\tilde{R}_k\times\mathfrak{X}_k$, let $\boldsymbol{\lambda}$, $\boldsymbol{\lambda}'$ be real functions on the sites of $T_k$, and let $\mathscr{A}^1=\mathscr{A}-d\boldsymbol{\lambda}$ and $\mathscr{C}^1=\mathscr{C}-d\boldsymbol{\lambda}'$ on the bonds at distance at most $1$ from $X$. Then $\Lambda:=\boldsymbol{\lambda}+\boldsymbol{\lambda}'$ and the total connection $\mathscr{T}^1=\mathscr{A}^1+\mathscr{C}^1$ satisfy $\mathscr{T}^1=\mathscr{T}-d\Lambda$ on these bonds, among which are the bonds whose two end points lie in $X$. The map $G_{\Lambda}$ preserves the bidegree by~\eqref{eq:polymer-activities-4}, and the property~\eqref{it:renormalisation-group-1-k4} of $E$ together with the second relation of~\eqref{eq:polymer-activities-18} yields $G_{\Lambda}\left[E^{\mathrm{nat}}(X;\mathscr{A}^1,\mathscr{C}^1)\right]=E^{\mathrm{nat}}(X;\mathscr{A},\mathscr{C})$. The third and the fourth relation give the properties~\eqref{it:renormalisation-group-1-k5} and~\eqref{it:renormalisation-group-1-k7} in the same way: the total connections of $(-\mathscr{A},-\mathscr{C})$ and of $(-\overline{\mathscr{A}},-\overline{\mathscr{C}})$ are $-\mathscr{T}$ and $-\overline{\mathscr{T}}$, these pairs lie in $\tilde{R}_k\times\mathfrak{X}_k$ by Proposition~\ref{prop:properties-actions}(c), and $\mathbf{C}$ and $\mathbf{K}$ preserve the bidegree by~\eqref{eq:polymer-activities-4}. The property~\eqref{it:renormalisation-group-1-k7} of $E^{\mathrm{nat}}$ is obtained when $E$ has it, that is, when $E\in K^x_k$.

Every map in~\eqref{eq:polymer-activities-18} is a ring homomorphism; the two sides of the fourth relation are antilinear, those of the other three are linear (Section~\ref{sec:lattice-gauge-charge}). It therefore suffices to verify each relation on the generators $\psi(gy)$, $\bar{\psi}(gx)$, respectively $\psi(y)$, $\bar{\psi}(x)$, where $x$ and $y$ are sites of $X$.

\emph{The first relation.} By Proposition~\ref{prop:properties-actions}(a), $g$ maps $X$ onto $gX$, the $M$-cubes of $X$ and their centres onto those of $gX$, and the shortest chains of adjacent cubes of $X$ onto those of $gX$. Hence $\boldsymbol{\chi}\mapsto g\boldsymbol{\chi}$ is a bijection of $\mathrm{Ch}(y,y_0)$ onto $\mathrm{Ch}(gy,gy_0)$, and the points of $g\boldsymbol{\chi}$ are $gp_0,\dots,gp_{n+2}$. As $g$ is an isometry, $|gp_i-gp_{i+1}|=|p_i-p_{i+1}|$, and $u_{g\mathscr{T}}(gp_j,gp_{j+1})=u_{\mathscr{T}}(p_j,p_{j+1})$ by~\eqref{eq:polymer-activities-u3}; so the products $U_i$ formed with $(gX,g\mathscr{T},g\boldsymbol{\chi})$ coincide with those formed with $(X,\mathscr{T},\boldsymbol{\chi})$. Moreover $y_0\mapsto gy_0$ is a bijection of the sites of $X$ onto those of $gX$, and $\operatorname{Vol}(gX)=\operatorname{Vol}(X)$. Since $a_g\psi(gy_0)=S\cdot\psi(y_0)$ and $a_g\chi(gp,gp')=S\cdot\chi(p,p')$,
\begin{align*}
a_g\left(\boldsymbol{\Theta}_{gX,g\mathscr{T}}\psi(gy)\right)&=\operatorname{Vol}(X)^{-1}\int_X\mathrm{d}y_0\;\left(\mathbf{T}(y,y_0)\cdot S\cdot\psi(y_0)-\operatorname{Rm}(y,y_0)[S\cdot\chi]\right)\\
&=\boldsymbol{\Theta}_{X,\mathscr{T}}\left(S\cdot\psi(y)\right)=\boldsymbol{\Theta}_{X,\mathscr{T}}\left(a_g\psi(gy)\right).
\end{align*}
For the barred generators the same computation holds with $(S^{-1})^T$, $u^{-1}$, $U^{-}_i$, $\mathbf{T}^{-}$ and $\operatorname{Rm}^{-}$.

\emph{The second relation.} Let $U^1_i$ denote the products formed with $\mathscr{T}^1$ in place of $\mathscr{T}$. For every chain, the six paths $\Gamma^{\pi}(p_j,p_{j+1})$ lie in $X$ by~\eqref{eq:polymer-activities-ch1}, so that $\mathscr{T}^1=\mathscr{T}-d\Lambda$ on their bonds, and~\eqref{eq:polymer-activities-u3} gives
\[
u_{\mathscr{T}^1}(p_j,p_{j+1})=u_{\mathscr{T}}(p_j,p_{j+1})\cdot\exp\left(-i\,e_k\cdot\left(\Lambda(p_{j+1})-\Lambda(p_j)\right)\right).
\]
The product telescopes:
\[
U^1_i=U_i\cdot\exp\left(-i\,e_k\cdot\left(\Lambda(p_i)-\Lambda(p_0)\right)\right).
\]
The map $G_{\Lambda}$ multiplies $\psi(y_0)=\psi(p_{n+2})$ by $\exp\left(i\,e_k\cdot\Lambda(p_{n+2})\right)$, and $\chi(p_i,p_{i+1})$ by $\exp\left(i\,e_k\cdot\Lambda(p_i)\right)$, the phase of the first point of the pair (Section~\ref{sec:lattice-gauge-charge}). Therefore $U^1_i\cdot\exp\left(i\,e_k\cdot\Lambda(p_i)\right)=U_i\cdot\exp\left(i\,e_k\cdot\Lambda(y)\right)$ for $0\le i\le n+2$, and
\[
G_{\Lambda}\left(\boldsymbol{\Theta}_{X,\mathscr{T}^1}\psi(y)\right)=\exp\left(i\,e_k\cdot\Lambda(y)\right)\cdot\boldsymbol{\Theta}_{X,\mathscr{T}}\psi(y)=\boldsymbol{\Theta}_{X,\mathscr{T}}\left(G_{\Lambda}\psi(y)\right).
\]
For the barred generators the same holds with $U^{-}_i$ and with all phases of the opposite sign.

\emph{The third relation.} By~\eqref{eq:polymer-activities-u3}, $u_{-\mathscr{T}}=u_{\mathscr{T}}^{-1}$. Hence the replacement of $\mathscr{T}$ by $-\mathscr{T}$ turns $U_i$, $\mathbf{T}$ and $\operatorname{Rm}$ into $U^{-}_i$, $\mathbf{T}^{-}$ and $\operatorname{Rm}^{-}$, and conversely. The map $\mathbf{C}$ sends $\psi$ and $\chi$ to $\mathsf{C}\cdot\bar{\psi}$ and $\mathsf{C}\cdot\bar{\chi}$, and $\bar{\psi}$ and $\bar{\chi}$ to $-\mathsf{C}\cdot\psi$ and $-\mathsf{C}\cdot\chi$. Therefore
\begin{align*}
\mathbf{C}\left(\boldsymbol{\Theta}_{X,-\mathscr{T}}\psi(y)\right)&=\operatorname{Vol}(X)^{-1}\int_X\mathrm{d}y_0\;\mathbf{T}^{-}(y,y_0)\cdot\mathsf{C}\cdot\bar{\psi}(y_0)\\
&\qquad-\operatorname{Vol}(X)^{-1}\int_X\mathrm{d}y_0\;\operatorname{Rm}^{-}(y,y_0)[\mathsf{C}\cdot\bar{\chi}]\\
&=\boldsymbol{\Theta}_{X,\mathscr{T}}\left(\mathsf{C}\cdot\bar{\psi}(y)\right)=\boldsymbol{\Theta}_{X,\mathscr{T}}\left(\mathbf{C}\,\psi(y)\right),\\
\mathbf{C}\left(\boldsymbol{\Theta}_{X,-\mathscr{T}}\bar{\psi}(x)\right)&=-\operatorname{Vol}(X)^{-1}\int_X\mathrm{d}x_0\;\mathbf{T}(x,x_0)\cdot\mathsf{C}\cdot\psi(x_0)\\
&\qquad+\operatorname{Vol}(X)^{-1}\int_X\mathrm{d}x_0\;\operatorname{Rm}(x,x_0)[\mathsf{C}\cdot\chi]\\
&=\boldsymbol{\Theta}_{X,\mathscr{T}}\left(-\mathsf{C}\cdot\psi(x)\right)=\boldsymbol{\Theta}_{X,\mathscr{T}}\left(\mathbf{C}\,\bar{\psi}(x)\right).
\end{align*}

\emph{The fourth relation.} By~\eqref{eq:polymer-activities-u3}, $\overline{u_{-\overline{\mathscr{T}}}}=u_{\mathscr{T}}$. The numbers $\operatorname{Vol}(X)^{-1}$, $\eta^3$ and $|p_i-p_{i+1}|^{\alpha}$, and the weights of the means over the chains, are real. Hence the complex conjugates of $U_i$, $U^{-}_i$, $\mathbf{T}$, $\mathbf{T}^{-}$ and of the coefficients of $\operatorname{Rm}$ and $\operatorname{Rm}^{-}$, formed with $-\overline{\mathscr{T}}$, are the same quantities formed with $\mathscr{T}$. The map $\mathbf{K}$ conjugates the coefficients and replaces each of the generators $\psi$, $\bar{\psi}$, $\chi$, $\bar{\chi}$ by $\varkappa$ times it, and $\varkappa$ acts on the spinor index only. Therefore
\begin{align*}
\mathbf{K}\left(\boldsymbol{\Theta}_{X,-\overline{\mathscr{T}}}\psi(y)\right)&=\operatorname{Vol}(X)^{-1}\int_X\mathrm{d}y_0\;\left(\mathbf{T}(y,y_0)\cdot\varkappa\cdot\psi(y_0)-\operatorname{Rm}(y,y_0)[\varkappa\cdot\chi]\right)\\
&=\boldsymbol{\Theta}_{X,\mathscr{T}}\left(\mathbf{K}\,\psi(y)\right),
\end{align*}
and in the same way $\mathbf{K}\left(\boldsymbol{\Theta}_{X,-\overline{\mathscr{T}}}\bar{\psi}(x)\right)=\boldsymbol{\Theta}_{X,\mathscr{T}}\left(\mathbf{K}\,\bar{\psi}(x)\right)$. This proves~\eqref{eq:polymer-activities-18}.

\emph{Linearity and the restriction to $\mathscr{C}=0$.} For fixed $X$ and $(\mathscr{A},\mathscr{C})$ the maps $E\mapsto E^{[2,0]}(X;\mathscr{A},\mathscr{C})$ and $\boldsymbol{\Theta}_{X,\mathscr{T}}$ are linear; hence $E\mapsto E^{\mathrm{nat}}$ is linear. For a family without the second slot, Definition~\ref{def:transport-mean} is applied with the total connection $\mathscr{A}$. The total connection of $(\mathscr{A},0)$ is $\mathscr{A}$; so $E^{\mathrm{nat}}(X;\mathscr{A},0)=(E|_0)^{\mathrm{nat}}(X;\mathscr{A})$ for every $X$, that is, $E^{\mathrm{nat}}|_0=(E|_0)^{\mathrm{nat}}$.

\emph{The families without the second slot.} Each identity and each estimate of the proof above concerns the family at a single pair $(\mathscr{A},\mathscr{C})$ of $\tilde{R}_k\times\mathfrak{X}_k$, at $(0,0)$, or at the pairs that the properties~\eqref{it:renormalisation-group-1-k3}, \eqref{it:renormalisation-group-1-k4}, \eqref{it:renormalisation-group-1-k5} and~\eqref{it:renormalisation-group-1-k7} relate to $(\mathscr{A},\mathscr{C})$, namely $(g\mathscr{A},g\mathscr{C})$, $(\mathscr{A}^1,\mathscr{C}^1)$, $(-\mathscr{A},-\mathscr{C})$ and $(-\overline{\mathscr{A}},-\overline{\mathscr{C}})$; if $\mathscr{C}=0$ and $\mathscr{C}^1=0$, each of these pairs has the second entry $0$. In detail, let $E\in K^C_k$ and $\mathscr{A}\in\tilde{R}_k$, and let $E^{\mathrm{nat}}$ be the family described after~(e) in the statement. Both norms of the zero bond function vanish, and the bounds $e_k^{-1+2\varepsilon}$ and $e_k^{-1+3\varepsilon}$ of Definition~\ref{def:domains}(f) are positive; hence $0\in\mathfrak{X}_k$, the pair $(\mathscr{A},0)$ lies in $\tilde{R}_k\times\mathfrak{X}_k$, and its total connection is $\mathscr{A}$. Consequently the statements of Section~\ref{sec:averaging-bilinear-part} that involve no family hold with $\mathscr{T}=\mathscr{A}$ and $u=u_{\mathscr{A}}$: the bound~\eqref{eq:polymer-activities-15}, the properties~\eqref{eq:polymer-activities-u1}--\eqref{eq:polymer-activities-u4} and~\eqref{eq:polymer-activities-ch1}--\eqref{eq:polymer-activities-ch3}, the definitions~\eqref{eq:polymer-activities-17}, and the homomorphism $\boldsymbol{\Theta}_{X,\mathscr{A}}$ of Definition~\ref{def:transport-mean}. The four relations~\eqref{eq:polymer-activities-18} involve no family either; they hold for the total connection of every pair in $\tilde{R}_k\times\mathfrak{X}_k$, in particular for $\mathscr{T}=\mathscr{A}$.

The proofs of~(a)--(d) are carried out at the polymer $X$ and at one pair of connections, the first assertion of~(b) at $(0,0)$. Of the family they use only the property~\eqref{it:renormalisation-group-1-k2}, the equality of the numbers of barred and unbarred legs (through~\eqref{eq:polymer-activities-3}), and the element and its kernels at that pair; the second slot enters them only through $\mathscr{T}$. The families of $K^C_k$ have these properties (Definition~\ref{def:renormalisation-group-1}), and $0\in\tilde{R}_k$ by Definition~\ref{def:domains}(a),(b), since all norms of the zero bond function vanish and the thresholds of $R_k$ are positive. Hence the same proofs, with $\mathscr{C}=0$, $\mathscr{T}=\mathscr{A}$ and the replacements listed in the statement, give~(a)--(d) for $E$ at $\mathscr{A}$, respectively at $0$.

In~(e), as above, it suffices to consider small polymers $X$. Each kernel of $E^{\mathrm{nat}}(X;\mathscr{A})$ arises by finitely many additions and multiplications from the kernels of $E(X;\mathscr{A})$, which are analytic on the open set $\tilde{R}_k$ by the property~\eqref{it:renormalisation-group-1-k1} of $E$, from the transports $u_{\mathscr{A}}^{\pm1}$, which are entire in $\mathscr{A}$ by~\eqref{eq:polymer-activities-u2}, and from numbers independent of $\mathscr{A}$; this is the property~\eqref{it:renormalisation-group-1-k1} of $E^{\mathrm{nat}}$. The property~\eqref{it:renormalisation-group-1-k2} is part~(c), and evenness and the balance of barred and unbarred legs follow as above. Part~(d), together with $E^{\mathrm{nat}}(Y;\mathscr{A})=E(Y;\mathscr{A})$ for large $Y$, gives the property~\eqref{it:renormalisation-group-1-k6} and $\|E^{\mathrm{nat}}\|_k\le10\cdot\|E\|_k$, the supremum being taken over $\mathscr{A}\in\tilde{R}_k$. Without the second slot the properties~\eqref{it:renormalisation-group-1-k3}, \eqref{it:renormalisation-group-1-k4}, \eqref{it:renormalisation-group-1-k5} and~\eqref{it:renormalisation-group-1-k7} are those of Section~\ref{sec:lattice-gauge-charge} at $\mathscr{C}=\mathscr{C}^1=0$:
\begin{gather*}
a_g\left[E(gX;g\mathscr{A})\right]=E(X;\mathscr{A}),\qquad G_{\boldsymbol{\lambda}+\boldsymbol{\lambda}'}\left[E(X;\mathscr{A}^1)\right]=E(X;\mathscr{A}),\\
\mathbf{C}\left[E(X;-\mathscr{A})\right]=E(X;\mathscr{A}),\qquad\mathbf{K}\left[E\left(X;-\overline{\mathscr{A}}\right)\right]=E(X;\mathscr{A});
\end{gather*}
in the second relation $\mathscr{A}$ and $\mathscr{A}^1$ lie in $\tilde{R}_k$, $\boldsymbol{\lambda}$ and $\boldsymbol{\lambda}'$ are real functions on the sites of $T_k$, and $\mathscr{A}^1=\mathscr{A}-d\boldsymbol{\lambda}$ and $d\boldsymbol{\lambda}'=0$ on the bonds at distance at most $1$ from $X$, so that $\Lambda:=\boldsymbol{\lambda}+\boldsymbol{\lambda}'$ satisfies $\mathscr{A}^1=\mathscr{A}-d\Lambda$ on these bonds. The connections $g\mathscr{A}$, $-\mathscr{A}$ and $-\overline{\mathscr{A}}$ lie in $\tilde{R}_k$ by Proposition~\ref{prop:properties-actions}(c); they are $g\mathscr{T}$, $-\mathscr{T}$ and $-\overline{\mathscr{T}}$ for $\mathscr{T}=\mathscr{A}$, and $\mathscr{A}^1$ is the total connection of the pair $(\mathscr{A}^1,0)$. Therefore the deduction of these properties of $E^{\mathrm{nat}}$ from~\eqref{eq:polymer-activities-18} given above applies with $\mathscr{C}=\mathscr{C}^1=0$; the property~\eqref{it:renormalisation-group-1-k7} of $E^{\mathrm{nat}}$ is obtained when $E\in K_k$. Finally, for fixed $X$ and $\mathscr{A}$ the maps $E\mapsto E^{[2,0]}(X;\mathscr{A})$ and $\boldsymbol{\Theta}_{X,\mathscr{A}}$ are linear; hence $E\mapsto E^{\mathrm{nat}}$ is linear on $K^C_k$.

Proposition~\ref{prop:adjustment} follows.
\end{proof}

\subsection{Scaled connections near a small polymer}\label{sec:scaled-connections-near}

Throughout Section~\ref{sec:scaled-connections-near} the setting~\eqref{eq:polymer-activities-Sb} is assumed. Since $\eta_{k+1}=\eta_k/L$ and the sides of $T_k$ and $T_{k+1}$ are $l$ and $l/L$, the map $x\mapsto x/L$ is a bijection of the sites of $T_k$ onto the sites of $T_{k+1}$; it maps nearest neighbours onto nearest neighbours and multiplies the sup distance by $L^{-1}$. For a bond $b$ of $T_k$, $b/L$ denotes its image, a bond of $T_{k+1}$. For a set $W'$ of sites of $T_{k+1}$ put $LW':=\left\{Lx:x\in W'\right\}$, a set of sites of $T_k$. A complex bond function $A$ defined on the bonds of $T_{k+1}$ with both end points in $W'$ gives the bond function $A_L$ of Section~\ref{sec:local-parts-adjustment}, defined on the bonds of $T_k$ with both end points in $LW'$, and a real function $\boldsymbol{\lambda}$ on $W'$ gives the real function $\boldsymbol{\lambda}_L$ on $LW'$:
\[
A_L(b)=L^{-\frac{1}{2}}A(b/L),\qquad \boldsymbol{\lambda}_L(x):=L^{\frac{1}{2}}\boldsymbol{\lambda}(x/L).
\]
For two points $x$, $y$ of the continuum torus $\mathbb{R}^3/l\mathbb{Z}^3$, $|x-y|$ and $|x-y|_1$ denote the least sup norm and the least $\ell^1$ norm of a vector of $\mathbb{R}^3$ that represents $x-y$; for two sites, $|x-y|=d(x,y)$. The same notation is used on $T_{k+1}$. The thresholds $t_{1,k}$, $t_{2,k}$, $t_{3,k}$ of $R_k$ and $x_{1,k}$, $x_{2,k}$ of $\mathfrak{X}_k$, and the three norms $|a|$, $|\partial a|$, $[\partial a]_{\boldsymbol{\alpha}}$ of a bond function $a$, are those of Section~\ref{sec:settings-basic-notation}; the thresholds of the level $k+1$ are formed with $e_{k+1}$ in place of $e_k$. On a set of sites the three norms are formed with the bonds and the difference quotients inside this set (Section~\ref{sec:domains-bond-functions}), and they are below the thresholds of $R_k$ if $|a|<t_{1,k}$, $|\partial a|<t_{2,k}$ and $[\partial a]_{\boldsymbol{\alpha}}<t_{3,k}$.

\begin{lemma}[Scaling of the field domains]\label{lem:scaling-field-domains}
\begin{enumerate}
\item[(a)] Let $W'$ be a set of sites of $T_{k+1}$ and $A$ a complex bond function on the bonds of $T_{k+1}$ with both end points in $W'$. For all directions $\mu$, $\nu$ and all sites $x$, $x'$ of $LW'$ for which the quantities on the left are formed inside $LW'$,
\begin{align*}
&(A_L)_{\mu}(x)=L^{-\frac{1}{2}}A_{\mu}(x/L),\qquad \left(\partial_{\nu}(A_L)_{\mu}\right)(x)=L^{-\frac{3}{2}}\left(\partial_{\nu}A_{\mu}\right)(x/L),\\
&\frac{\left|\partial_{\nu}(A_L)_{\mu}(x')-\partial_{\nu}(A_L)_{\mu}(x)\right|}{|x-x'|^{\alpha}}=L^{-\frac{3}{2}-\alpha}\cdot\frac{\left|\partial_{\nu}A_{\mu}(x'/L)-\partial_{\nu}A_{\mu}(x/L)\right|}{|x/L-x'/L|^{\alpha}},
\end{align*}
the last identity for $0<|x-x'|\le 1$. Suppose, in addition, that the norms of $A$, formed on $W'$, satisfy
\[
|A|<t_{1,k+1},\qquad |\partial A|<t_{2,k+1},\qquad [\partial A]_{\boldsymbol{\alpha}}<t_{3,k+1}.
\]
Then the norms of $A_L$, formed on $LW'$, obey
\[
|A_L|<L^{-\frac{7}{8}+\frac{\varepsilon}{2}}\cdot t_{1,k},\qquad |\partial A_L|<v_1,\qquad [\partial A_L]_{\boldsymbol{\alpha}}<L^{-\alpha}\cdot w_1\le w_1,
\]
where
\[
v_1:=L^{-\frac{15}{8}+\varepsilon}\cdot t_{2,k},\qquad w_1:=L^{-\frac{15}{8}+\frac{3\varepsilon}{2}}\cdot t_{3,k}.
\]
\item[(b)] If $\mathscr{C}\in\mathfrak{X}_{k+1}$, then
\[
|\mathscr{C}_L|<L^{-1+\varepsilon}\cdot x_{1,k},\qquad |\partial\mathscr{C}_L|<L^{-2+\frac{3\varepsilon}{2}}\cdot x_{2,k};
\]
hence $\mathscr{C}_L\in L^{-\frac{3}{4}}\cdot\mathfrak{X}_k$, that is, the two norms of $\mathscr{C}_L$ are below $L^{-\frac{3}{4}}$ times the thresholds of $\mathfrak{X}_k$.
\item[(c)] Let $\mathscr{A}=A_0+A_1\in\tilde{R}_{k+1}$ be split as in Definition~\ref{def:domains}, let $\square'$ be an $M$-cube of level $k+1$, and let $A_0=A'+d\boldsymbol{\lambda}$ on $\square'^{\dagger}$, with a real function $\boldsymbol{\lambda}$ on the sites and a real bond function $A'$ on the bonds of $\square'^{\dagger}$ as in Definition~\ref{def:domains}. Let $z'$ be the centre of the $LM$-cube $L\square'$. Then
\begin{gather*}
L\square'^{\dagger}=\left\{x\in T_k:|x-z'|<\left(L+\frac{1}{2}\right)LM\right\},\\
(A_0)_L=(A')_L+d(\boldsymbol{\lambda}_L)\quad\text{on }L\square'^{\dagger};
\end{gather*}
the bond function $(A')_L$ satisfies the three bounds of~(a) on $L\square'^{\dagger}$, and $(A_1)_L$ satisfies them on $T_k$. For every $M$-cube $\square$ of level $k$ contained in $L\square'$, the set $\square^{\dagger}$ is contained in $L\square'^{\dagger}$. Consequently $\mathscr{A}_L\in L^{-\frac{3}{4}}\cdot\tilde{R}_k$. This set is contained in $\tilde{R}_k$, and it is contained in $\tilde{R}_k^{-}$, because $L^{-\frac{3}{4}}\le\frac{1}{2}\le 1-e_k^{\varepsilon}$.
\end{enumerate}
\end{lemma}

\begin{proof}
(a) Since $\eta_{k+1}=\eta_k/L$, the map $x\mapsto x/L$ sends the bond $(x,x+\eta_ke_{\mu})$ to the bond $(x/L,x/L+\eta_{k+1}e_{\mu})$. The definition of $A_L$ therefore gives the first formula, and with it
\begin{align*}
\left(\partial_{\nu}(A_L)_{\mu}\right)(x)&=\frac{(A_L)_{\mu}(x+\eta_ke_{\nu})-(A_L)_{\mu}(x)}{\eta_k}=L^{-\frac{1}{2}}\cdot\frac{A_{\mu}(x/L+\eta_{k+1}e_{\nu})-A_{\mu}(x/L)}{L\eta_{k+1}}\\
&=L^{-\frac{3}{2}}\left(\partial_{\nu}A_{\mu}\right)(x/L).
\end{align*}
The third formula follows from the second and from $|x-x'|=L\,|x/L-x'/L|$. If $0<|x-x'|\le 1$, then $0<|x/L-x'/L|\le L^{-1}\le 1$. The bonds and the difference quotients inside $LW'$ are the images of those inside $W'$; hence, the norms of $A_L$ being taken on $LW'$ and those of $A$ on $W'$,
\[
|A_L|=L^{-\frac{1}{2}}|A|,\qquad |\partial A_L|=L^{-\frac{3}{2}}|\partial A|,\qquad [\partial A_L]_{\boldsymbol{\alpha}}\le L^{-\frac{3}{2}-\alpha}[\partial A]_{\boldsymbol{\alpha}}.
\]
Since $e_{k+1}=L^{\frac{1}{2}}e_k$, the thresholds of the two levels are related by
\[
t_{1,k+1}=L^{-\frac{3}{8}+\frac{\varepsilon}{2}}\cdot t_{1,k},\qquad t_{2,k+1}=L^{-\frac{3}{8}+\varepsilon}\cdot t_{2,k},\qquad t_{3,k+1}=L^{-\frac{3}{8}+\frac{3\varepsilon}{2}}\cdot t_{3,k}.
\]
As $L^{-\frac{1}{2}}L^{-\frac{3}{8}+\frac{\varepsilon}{2}}=L^{-\frac{7}{8}+\frac{\varepsilon}{2}}$, $L^{-\frac{3}{2}}L^{-\frac{3}{8}+\varepsilon}=L^{-\frac{15}{8}+\varepsilon}$ and $L^{-\frac{3}{2}-\alpha}L^{-\frac{3}{8}+\frac{3\varepsilon}{2}}=L^{-\alpha}L^{-\frac{15}{8}+\frac{3\varepsilon}{2}}$, the three bounds follow. Finally $L^{-\alpha}\le 1$, because $\alpha>0$ (stage~\ref{it:order-choice-2}) and $L>1$.

(b) Since $e_{k+1}=L^{\frac{1}{2}}e_k$, one has $x_{1,k+1}=L^{-\frac{1}{2}+\varepsilon}\cdot x_{1,k}$ and $x_{2,k+1}=L^{-\frac{1}{2}+\frac{3\varepsilon}{2}}\cdot x_{2,k}$. The first two formulas of~(a), applied with $W'$ the set of all sites of $T_{k+1}$, give
\begin{align*}
|\mathscr{C}_L|&=L^{-\frac{1}{2}}|\mathscr{C}|<L^{-\frac{1}{2}}\cdot x_{1,k+1}=L^{-1+\varepsilon}\cdot x_{1,k},\\
|\partial\mathscr{C}_L|&=L^{-\frac{3}{2}}|\partial\mathscr{C}|<L^{-\frac{3}{2}}\cdot x_{2,k+1}=L^{-2+\frac{3\varepsilon}{2}}\cdot x_{2,k}.
\end{align*}
Since $\varepsilon\le\bar{\varepsilon}=\frac{1}{200}$ (stages~\ref{it:order-choice-1} and~\ref{it:order-choice-8}), the exponents $-1+\varepsilon$ and $-2+\frac{3\varepsilon}{2}$ are smaller than $-\frac{3}{4}$; as $L>1$, it follows that $\mathscr{C}_L\in L^{-\frac{3}{4}}\cdot\mathfrak{X}_k$.

(c) Let $\boldsymbol{\lambda}$ be a real function on a set $W'$ of sites of $T_{k+1}$, and let $x$, $x'$ be nearest neighbours in $LW'$. Since $L^{-\frac{1}{2}}\eta_{k+1}^{-1}=L^{\frac{1}{2}}\eta_k^{-1}$,
\[
(d\boldsymbol{\lambda})_L(x,x')=L^{-\frac{1}{2}}\cdot\frac{\boldsymbol{\lambda}(x'/L)-\boldsymbol{\lambda}(x/L)}{\eta_{k+1}}=\frac{\boldsymbol{\lambda}_L(x')-\boldsymbol{\lambda}_L(x)}{\eta_k}=\left(d(\boldsymbol{\lambda}_L)\right)(x,x').
\]
Thus $(d\boldsymbol{\lambda})_L=d(\boldsymbol{\lambda}_L)$, and the identity $A_0=A'+d\boldsymbol{\lambda}$ on $\square'^{\dagger}$ gives $(A_0)_L=(A')_L+d(\boldsymbol{\lambda}_L)$ on $L\square'^{\dagger}$. By Definition~\ref{def:enlarged-cubes}, $\square'^{\dagger}$ is the cube of side $(2L+1)M$ with the same centre $z'/L$ as $\square'$. No site lies on the boundary of an $M$-cube (Section~\ref{sec:balls-torus-lifts}), so the sites of $\square'^{\dagger}$ are the sites at sup distance less than $\left(L+\frac{1}{2}\right)M$ from $z'/L$; their images under $y\mapsto Ly$ form the set displayed in~(c). By Definition~\ref{def:domains}, $A'$ satisfies the three inequalities assumed in~(a) with $W'=\square'^{\dagger}$, and $A_1$, an element of $R_{k+1}$, satisfies them with $W'$ the set of all sites of $T_{k+1}$. Part~(a) therefore yields the bounds for $(A')_L$ on $L\square'^{\dagger}$ and for $(A_1)_L$ on $T_k$.

As $L$ is odd, the $LM$-cube $L\square'$ is the union of the $L^3$ $M$-cubes of level $k$ whose centres $z$ satisfy $|z-z'|\le\frac{(L-1)M}{2}$. If $\square$ is one of them, with the centre $z$, then
\[
\square^{\dagger}=\left\{x:|x-z|<\left(L+\frac{1}{2}\right)M\right\}\subset\left\{x:|x-z'|<\frac{3LM}{2}\right\}\subset L\square'^{\dagger},
\]
because $\left(L+\frac{1}{2}\right)M+\frac{(L-1)M}{2}=\frac{3LM}{2}\le\left(L+\frac{1}{2}\right)LM$.

Now let $\square$ be an arbitrary $M$-cube of level $k$. Since $l\in\mathrm{Cl}$, the side $l$ is a multiple of $LM$; the $LM$-cubes of $T_k$ are the images under $y\mapsto Ly$ of the $M$-cubes of level $k+1$, and every $M$-cube lies in exactly one $LM$-cube. Choose $\square'$ with $\square\subset L\square'$, and $\boldsymbol{\lambda}$, $A'$ as in~(c). On $\square^{\dagger}$, which is contained in $L\square'^{\dagger}$, one has $(A_0)_L=(A')_L+d(\boldsymbol{\lambda}_L)$ with real $(A')_L$ and $\boldsymbol{\lambda}_L$, and the three norms of $(A')_L$ on $\square^{\dagger}$ are at most those on $L\square'^{\dagger}$, hence below $L^{-\frac{7}{8}+\frac{\varepsilon}{2}}\cdot t_{1,k}$, $v_1$ and $w_1$. Since $\varepsilon\le\frac{1}{200}$, the exponents $-\frac{7}{8}+\frac{\varepsilon}{2}$, $-\frac{15}{8}+\varepsilon$ and $-\frac{15}{8}+\frac{3\varepsilon}{2}$ are smaller than $-\frac{3}{4}$, so these three numbers are smaller than $L^{-\frac{3}{4}}\cdot t_{1,k}$, $L^{-\frac{3}{4}}\cdot t_{2,k}$ and $L^{-\frac{3}{4}}\cdot t_{3,k}$. The three norms of $(A_1)_L$ on $T_k$ obey the same bounds, so that $(A_1)_L\in L^{-\frac{3}{4}}R_k$, and $(A_0)_L$ is real. By Definition~\ref{def:domains}, with the constant $L^{-\frac{3}{4}}\le 1$, the connection $\mathscr{A}_L=(A_0)_L+(A_1)_L$ lies in $L^{-\frac{3}{4}}\cdot\tilde{R}_k$. The conditions that define $c\cdot\tilde{R}_k$ become weaker as $c$ increases; hence $L^{-\frac{3}{4}}\cdot\tilde{R}_k\subset\tilde{R}_k$. The set $\tilde{R}_k^{-}$ is defined by the same conditions, with the constant $1$ in the bounds on the representatives $A'$ of the real summand and with the complex summand in $(1-e_k^{\varepsilon})R_k$. By stage~\ref{it:order-choice-4}, $L^{-\frac{3}{4}}\le\frac{1}{2}$. By the setting~\eqref{eq:polymer-activities-Sb}, $e_k\le L^{\frac{1}{2}}e_k=e_{k+1}\le e_0$, and $e_0^{\varepsilon}\le\frac{1}{4}$ by condition~\eqref{eq:renormalisation-group-A8}; hence $e_k^{\varepsilon}\le\frac{1}{4}$ and $\frac{1}{2}\le 1-e_k^{\varepsilon}$. Therefore $L^{-\frac{3}{4}}R_k\subset(1-e_k^{\varepsilon})R_k$, and $L^{-\frac{3}{4}}\cdot\tilde{R}_k\subset\tilde{R}_k^{-}$.

In the setting~\eqref{eq:polymer-activities-Sb}, Lemma~\ref{lem:rescaling} follows from~(b) and~(c).
\end{proof}

\paragraph{Cut-off functions.}
Let $\check{f}$ be the real function on $\mathbb{R}$ given by
\[
\check{f}(s):=\begin{cases}
1, & |s|\le 1,\\
1-2\left(|s|-1\right)^2, & 1\le|s|\le\frac{3}{2},\\
2\left(2-|s|\right)^2, & \frac{3}{2}\le|s|\le 2,\\
0, & |s|\ge 2;
\end{cases}
\]
the expressions agree where two cases overlap. Since $\check{f}(s)$ depends on $|s|$ only, $\check{f}$ is even. For $s\ge 0$ one finds $\check{f}'(s)=0$ on $[0,1]$ and on $[2,\infty)$, $\check{f}'(s)=-4(s-1)$ on $\left[1,\frac{3}{2}\right]$ and $\check{f}'(s)=-4(2-s)$ on $\left[\frac{3}{2},2\right]$; these expressions agree at $s=1$, $s=\frac{3}{2}$ and $s=2$, and $\check{f}$ is constant on $[-1,1]$. By evenness, $\check{f}$ is therefore of class $C^1$ on $\mathbb{R}$, with an odd derivative. On $[0,\infty)$ this derivative takes its values in $[-2,0]$ and is continuous and piecewise linear with the slopes $0$, $-4$ and $4$; by oddness, the same holds on $(-\infty,0]$, with values in $[0,2]$. Consequently $|\check{f}'|\le 2$, and $\check{f}'$ has the Lipschitz constant $4$. Finally, $\check{f}$ is non-increasing on $[0,\infty)$, with $\check{f}(0)=1$ and $\check{f}(s)=0$ for $s\ge 2$; together with evenness this gives $0\le\check{f}\le 1$.

Let $x_0$ be a site of $T_k$ and $r>0$ a real number with $4r+11<l$. For a point $x$ of the continuum torus with $|x-x_0|\le 2r+5$ there is exactly one $z\in\mathbb{R}^3$ with $|z|\le 2r+5$ that represents $x-x_0$, because $2(2r+5)<l$; its coordinates are denoted by $(x-x_0)_{\mu}$. Define the function $\boldsymbol{\chi}_r$ on the continuum torus $\mathbb{R}^3/l\mathbb{Z}^3$ by
\[
\boldsymbol{\chi}_r(x):=\begin{cases}
\displaystyle\prod_{\mu}\check{f}\left(\frac{(x-x_0)_{\mu}}{r}\right), & |x-x_0|\le 2r+5,\\
0 & \text{otherwise},
\end{cases}
\]
the product being taken over the three directions, and for a site $x$ of $T_k$ and a direction $\mu$ let
\[
f_{\mu}(x):=\boldsymbol{\chi}_r\left(x+\frac{\eta}{2}e_{\mu}\right)
\]
be the value of $\boldsymbol{\chi}_r$ on the bond $(x,x+\eta e_{\mu})$ in the sense of convention~\eqref{eq:polymer-activities-10}. Then, for all sites $x$, $y$ of $T_k$ and all directions $\mu$, $\nu$,
\begin{align}
&0\le f_{\mu}\le 1;\quad f_{\mu}(x)=1\ \text{ if }x,\ x+\eta e_{\mu}\in K(x_0,r);\quad f_{\mu}(x)=0\ \text{ if }d(x,x_0)\ge 2r+1;\tag{f1}\label{eq:polymer-activities-f1}\\
&\left|f_{\mu}(x)-f_{\mu}(y)\right|\le\frac{2}{r}|x-y|_1\le\frac{6}{r}|x-y|,\qquad \left|\partial_{\nu}f_{\mu}(x)\right|\le\frac{2}{r};\tag{f2}\label{eq:polymer-activities-f2}\\
&\left|\partial_{\nu}f_{\mu}(x)-\partial_{\nu}f_{\mu}(y)\right|\le\frac{4}{r^2}|x-y|_1\le\frac{12}{r^2}|x-y|.\tag{f3}\label{eq:polymer-activities-f3}
\end{align}

Since $\check{f}$ vanishes outside $[-2,2]$, the function $\boldsymbol{\chi}_r$ vanishes at every point $x$ with $|x-x_0|\le 2r+5$ at which some coordinate $(x-x_0)_{\mu}$ has modulus at least $2r$. Consequently, as a function of one coordinate of its argument, the other coordinates being fixed, $\boldsymbol{\chi}_r$ is continuously differentiable on the circle $\mathbb{R}/l\mathbb{Z}$; its derivative $D_{\nu}\boldsymbol{\chi}_r$ with respect to the coordinate $\nu$ is
\[
D_{\nu}\boldsymbol{\chi}_r(x)=\begin{cases}
\displaystyle\frac{1}{r}\,\check{f}'\left(\frac{(x-x_0)_{\nu}}{r}\right)\prod_{\mathfrak{l}\ne\nu}\check{f}\left(\frac{(x-x_0)_{\mathfrak{l}}}{r}\right), & |x-x_0|\le 2r+5,\\
0 & \text{otherwise},
\end{cases}
\]
the product being taken over the two directions $\mathfrak{l}\ne\nu$; it vanishes as well at every point at which some coordinate $(x-x_0)_{\mu}$ has modulus at least $2r$.

\eqref{eq:polymer-activities-f1}: the values of $\check{f}$ lie in $[0,1]$. If $x$ and $x+\eta e_{\mu}$ lie in $K(x_0,r)$, then $x-x_0$ and $x+\eta e_{\mu}-x_0$ are represented by $z$ and $z+\eta e_{\mu}$ with $|z|\le r$ and $|z+\eta e_{\mu}|\le r$ (Section~\ref{sec:balls-torus-lifts}, as $2r+1<l$); so every coordinate of the barycentre $x+\frac{\eta}{2}e_{\mu}$ relative to $x_0$ has modulus at most $r$, and every factor of the product equals $1$. If $d(x,x_0)\ge 2r+1$, the barycentre has the sup distance at least $2r+1-\frac{\eta}{2}>2r$ from $x_0$; either this distance exceeds $2r+5$, or one of the coordinates of the barycentre relative to $x_0$ has modulus greater than $2r$, and in both cases $\boldsymbol{\chi}_r$ vanishes there.

\eqref{eq:polymer-activities-f2}: each function $s\mapsto\check{f}\left(\frac{s}{r}\right)$ takes its values in $[0,1]$ and is Lipschitz continuous with the constant $\frac{2}{r}$. Hence a change of one coordinate of the argument by an amount $\delta$ changes $\boldsymbol{\chi}_r$ by at most $\frac{2}{r}|\delta|$, and changing the three coordinates one after the other one obtains $\left|\boldsymbol{\chi}_r(p)-\boldsymbol{\chi}_r(q)\right|\le\frac{2}{r}|p-q|_1$ for all points $p$, $q$ of the continuum torus. With $p=x+\frac{\eta}{2}e_{\mu}$, $q=y+\frac{\eta}{2}e_{\mu}$ and $|p-q|_1=|x-y|_1\le 3|x-y|$ this is the first assertion of~\eqref{eq:polymer-activities-f2}. Taking $y=x+\eta e_{\nu}$ one obtains $\left|\partial_{\nu}f_{\mu}(x)\right|=\eta^{-1}\left|f_{\mu}(x+\eta e_{\nu})-f_{\mu}(x)\right|\le\frac{2}{r}$.

\eqref{eq:polymer-activities-f3}: by the fundamental theorem of calculus, applied in the coordinate $\nu$,
\begin{align*}
\partial_{\nu}f_{\mu}(x)&=\frac{1}{\eta}\left(\boldsymbol{\chi}_r\left(x+\frac{\eta}{2}e_{\mu}+\eta e_{\nu}\right)-\boldsymbol{\chi}_r\left(x+\frac{\eta}{2}e_{\mu}\right)\right)\\
&=\int_0^1\mathrm{d}t\,\left(D_{\nu}\boldsymbol{\chi}_r\right)\left(x+\frac{\eta}{2}e_{\mu}+t\eta e_{\nu}\right).
\end{align*}
As a function of any one coordinate of its argument, the others being fixed, $D_{\nu}\boldsymbol{\chi}_r$ has the Lipschitz constant $\frac{4}{r^2}$: for the coordinate $\nu$ because $s\mapsto\frac{1}{r}\check{f}'\left(\frac{s}{r}\right)$ has the Lipschitz constant $\frac{4}{r^2}$ and the other factors lie in $[0,1]$; for a coordinate $\mathfrak{l}\ne\nu$ because $\left|\frac{1}{r}\check{f}'\right|\le\frac{2}{r}$ and $s\mapsto\check{f}\left(\frac{s}{r}\right)$ has the Lipschitz constant $\frac{2}{r}$. Changing the coordinates one after the other gives $\left|D_{\nu}\boldsymbol{\chi}_r(p)-D_{\nu}\boldsymbol{\chi}_r(q)\right|\le\frac{4}{r^2}|p-q|_1$, and integration over $t$, with $|x-y|_1\le 3|x-y|$, gives~\eqref{eq:polymer-activities-f3}.

\begin{lemma}[Truncation of a bond function near a site]\label{lem:truncation-bond-function}
Let $x_0$ be a site of $T_k$ and $r>0$ with $4r+11<l$, and let $f_{\mu}$ be the functions defined above with these $x_0$ and $r$. Let $b$ be a complex bond function on the bonds of $T_k$ with both end points in $K(x_0,2r+5)$, and let $B_0:=|b|$, $B_1:=|\partial b|$ and $B_{\boldsymbol{\alpha}}:=[\partial b]_{\boldsymbol{\alpha}}$ be its three norms on $K(x_0,2r+5)$. Define the bond function $\boldsymbol{\chi}_rb$ on $T_k$ by
\[
(\boldsymbol{\chi}_rb)_{\mu}(x):=\begin{cases}
f_{\mu}(x)\,b_{\mu}(x), & x\in K(x_0,2r+1),\\
0 & \text{for all other sites }x\text{ of }T_k.
\end{cases}
\]
Then $\boldsymbol{\chi}_rb$ coincides with $b$ on every bond whose end points both lie in $K(x_0,r)$, and the three norms of $\boldsymbol{\chi}_rb$, formed on $T_k$, obey
\[
|\boldsymbol{\chi}_rb|\le B_0,\qquad |\partial(\boldsymbol{\chi}_rb)|\le B_1+\frac{2}{r}B_0,\qquad [\partial(\boldsymbol{\chi}_rb)]_{\boldsymbol{\alpha}}\le B_{\boldsymbol{\alpha}}+\frac{12}{r}B_1+\frac{12}{r^2}B_0.
\]
\end{lemma}

\begin{proof}
Since $2(2r+5)+1=4r+11<l$, the ball $K(x_0,2r+5)$ does not wrap, and nearest neighbours, pairs of sites at distance at most $1$, rectilinear paths and difference quotients inside it are those of the cube $\left\{z\in\mathbb{L}_k:|z|\le 2r+5\right\}$ (Section~\ref{sec:balls-torus-lifts}). For every site $x$ of $K(x_0,2r+4)$ and every direction $\mu$ the number $b_{\mu}(x)$ is defined, because $x+\eta e_{\mu}\in K(x_0,2r+5)$, and $(\boldsymbol{\chi}_rb)_{\mu}(x)=f_{\mu}(x)\,b_{\mu}(x)$: for $x\in K(x_0,2r+1)$ this is the definition, and otherwise $d(x,x_0)>2r+1$ and both sides vanish by~\eqref{eq:polymer-activities-f1}.

A bond with both end points in $K(x_0,r)$ is, up to orientation, of the form $(x,x+\eta e_{\mu})$ with $x$, $x+\eta e_{\mu}\in K(x_0,r)$. Then $x\in K(x_0,2r+1)$ and $f_{\mu}(x)=1$ by~\eqref{eq:polymer-activities-f1}, so that $(\boldsymbol{\chi}_rb)_{\mu}(x)=b_{\mu}(x)$. The first bound follows from $0\le f_{\mu}\le 1$.

Fix directions $\mu$, $\nu$ and a site $x$, and put $x^{+}:=x+\eta e_{\nu}$. If $d(x,x_0)\ge 2r+2$, then $d(x^{+},x_0)\ge 2r+1$, and $(\boldsymbol{\chi}_rb)_{\mu}$ vanishes at $x$ and at $x^{+}$, by its definition or by~\eqref{eq:polymer-activities-f1}; hence $\partial_{\nu}(\boldsymbol{\chi}_rb)_{\mu}(x)=0$. If $d(x,x_0)<2r+3$, then $x$, $x^{+}\in K(x_0,2r+4)$, and the product rule for difference quotients gives
\[
\partial_{\nu}(\boldsymbol{\chi}_rb)_{\mu}(x)=f_{\mu}(x^{+})\,\partial_{\nu}b_{\mu}(x)+b_{\mu}(x)\,\partial_{\nu}f_{\mu}(x).
\]
On the right, $b_{\mu}$ is evaluated only on the bonds $(x,x+\eta e_{\mu})$ and $(x^{+},x^{+}+\eta e_{\mu})$; as $x$ and $x^{+}$ lie in $K(x_0,2r+4)$, the end points of these two bonds belong to $K(x_0,2r+5)$, so that $|b_{\mu}(x)|\le B_0$ and $|\partial_{\nu}b_{\mu}(x)|\le B_1$. Together with~\eqref{eq:polymer-activities-f1} and~\eqref{eq:polymer-activities-f2} this bounds the modulus of the right-hand side by $B_1+\frac{2}{r}B_0$. Every site satisfies $d(x,x_0)\ge 2r+2$ or $d(x,x_0)<2r+3$; this proves the second bound.

Let now $x$, $y$ be sites with $0<d(x,y)\le 1$, and put $x^{+}:=x+\eta e_{\nu}$, $y^{+}:=y+\eta e_{\nu}$. If $d(x,x_0)\ge 2r+2$ and $d(y,x_0)\ge 2r+2$, both difference quotients vanish. Otherwise one of the two distances is smaller than $2r+2$, and both are smaller than $2r+3$. Write $f=f_{\mu}$, $b=b_{\mu}$ and $\partial=\partial_{\nu}$; by the preceding paragraph, $\partial_{\nu}(\boldsymbol{\chi}_rb)_{\mu}=\partial(fb)$ at $x$ and at $y$, and
\begin{multline*}
\partial(fb)(x)-\partial(fb)(y)=f(x^{+})\left[\partial b(x)-\partial b(y)\right]+\left[f(x^{+})-f(y^{+})\right]\partial b(y)\\
+\left[b(x)-b(y)\right]\partial f(x)+b(y)\left[\partial f(x)-\partial f(y)\right].
\end{multline*}
The sites $x$ and $y$ lie in $K(x_0,2r+3)$, which is the image of a cube of $\mathbb{L}_k$. A rectilinear path from $y$ to $x$ inside it has $\eta^{-1}|x-y|_1$ steps, and each step changes $b_{\mu}$ by $\eta$ times a difference quotient inside $K(x_0,2r+5)$; hence
\[
|b(x)-b(y)|\le B_1|x-y|_1\le 3B_1|x-y|.
\]
Further $|x^{+}-y^{+}|=|x-y|$, and the difference quotients of $b$ at $x$ and at $y$ are formed inside $K(x_0,2r+5)$. With~\eqref{eq:polymer-activities-f1}, \eqref{eq:polymer-activities-f2} and~\eqref{eq:polymer-activities-f3}, the moduli of the four terms on the right are at most
\[
B_{\boldsymbol{\alpha}}|x-y|^{\alpha},\qquad\frac{6}{r}B_1|x-y|,\qquad\frac{6}{r}B_1|x-y|,\qquad\frac{12}{r^2}B_0|x-y|.
\]
Since $|x-y|\le 1$ and $\alpha\le 1$ (stage~\ref{it:order-choice-2}), one has $|x-y|\le|x-y|^{\alpha}$, and the third bound follows.
\end{proof}

\begin{lemma}[Functions with vanishing sum on the torus]\label{lem:functions-vanishing-sum}
Let $w$ be a real number, and let $u$ be a complex-valued function on the set of sites of $T_k$ such that
\[
\sum_xu(x)=0\qquad\text{and}\qquad|u(x)-u(y)|\le w\quad\text{whenever }d(x,y)\le 1.
\]
Then $|u(x)|\le l\,w$ at every site $x$. Moreover, $u=\partial_{\nu}a_{\mu}$ satisfies $\sum_xu(x)=0$, whatever the bond function $a$ on $T_k$ and the directions $\mu$, $\nu$.
\end{lemma}

\begin{proof}
Let $x$, $y$ be sites of $T_k$; then $d(x,y)\le\frac{l}{2}$. Let $z\in\mathbb{L}_k$ represent $x-y$ with $|z|=d(x,y)$, and put $n:=\left\lceil\frac{l}{2}\right\rceil$. Every coordinate of $z$ is a multiple of $\eta$ of modulus at most $n$, and $1$ is a multiple of $\eta$; so $z$ is the sum of $n$ vectors of $\mathbb{L}_k$ of sup norm at most $1$, and $y$ is joined to $x$ by a chain of $n$ steps, each between two sites at distance at most $1$. As $l\ge 1$, one has $n\le l$, and therefore $|u(x)-u(y)|\le n\,w\le l\,w$. Since $\sum_yu(y)=0$, the number $u(x)$ is the mean of $u(x)-u(y)$ over the sites $y$ of $T_k$; hence $|u(x)|\le l\,w$. Finally, $x\mapsto x+\eta e_{\nu}$ is a bijection of the sites of $T_k$, so that
\[
\sum_x\partial_{\nu}a_{\mu}(x)=\frac{1}{\eta}\left(\sum_xa_{\mu}(x+\eta e_{\nu})-\sum_xa_{\mu}(x)\right)=0.
\]
\end{proof}

\begin{theorem}\label{thm:polymer-activities-3}
In the setting~\eqref{eq:polymer-activities-Sb} of Section~\ref{sec:settings-basic-notation}, let $X\in D_k$ be a small polymer, let $x_0$ be a site of $X$, let $(\mathscr{A},\mathscr{C})\in\tilde{R}_{k+1}\times\mathfrak{X}_{k+1}$, and put $\rho:=2LM$. Then there are two complex bond functions $\hat{b}$, $\mathscr{C}^b$ and two constant complex bond functions $c_1$, $c_2$ on $T_k$, together with a real function $\mathring{m}$ on the sites of $T_k$, with the properties (i)--(iii) below. They can be chosen, independently of the families $E$ of~(iii), in such a way that $\hat{b}$, $c_1$ and $\mathring{m}$ depend only on $\mathscr{A}$ and $x_0$, and not on $\mathscr{C}$, that $c_2$ is the constant bond function with the components $(\mathscr{C}_L)_{\mu}(x_0)$, and that $\mathscr{C}^b=\boldsymbol{\chi}_{\rho}\left(\mathscr{C}_L-c_2\right)$, where $\boldsymbol{\chi}_{\rho}$ is the truncation of Lemma~\ref{lem:truncation-bond-function} formed with the site $x_0$ and $r=\rho$. In particular, $\mathscr{C}^b=0$ and $c_2=0$ if $\mathscr{C}=0$.
\begin{enumerate}
\item[(i)] The identities
\[
\mathscr{A}_L-d\mathring{m}=\hat{b}+c_1,\qquad\mathscr{C}_L=\mathscr{C}^b+c_2
\]
hold on every bond of $T_k$ whose two end points lie in $K(x_0,\rho)$. Every bond that has an end point at distance at most $1$ from $X$ is a bond of this kind.
\item[(ii)] The bond functions $\hat{b}$ and $\mathscr{C}^b$ satisfy
\begin{gather*}
|\hat{b}|\le1382\,L^{-\frac{15}{8}+\frac{3\varepsilon}{2}}\cdot t_{1,k},\qquad|\partial\hat{b}|\le1382\,L^{-\frac{15}{8}+\frac{3\varepsilon}{2}}\cdot t_{2,k},\\
[\partial\hat{b}]_{\boldsymbol{\alpha}}\le1382\,L^{-\frac{15}{8}+\frac{3\varepsilon}{2}}\cdot t_{3,k},\\
|\mathscr{C}^b|<L^{-\frac{7}{4}}\cdot x_{1,k},\qquad|\partial\mathscr{C}^b|<L^{-\frac{7}{4}}\cdot x_{2,k};
\end{gather*}
consequently $(t\cdot\hat{b},t\cdot\mathscr{C}^b)\in R_k\times\mathfrak{X}_k$ for every complex number $t$ with $|t|\le L^{\frac{7}{4}}$. Moreover
\[
|c_1|\le L^{-\frac{7}{8}+\frac{\varepsilon}{2}}\cdot t_{1,k},\qquad|c_2|\le L^{-1+\varepsilon}\cdot x_{1,k}.
\]
\item[(iii)] Let $E=\{E(Y;\cdot,\cdot):Y\in D_k\}$ be a family of elements of $\mathrm{Gr}[\psi,\chi]$, defined on $\tilde{R}_k\times\mathfrak{X}_k$ and written with antisymmetric kernels $E_{nm}$ in the form displayed in Definition~\ref{def:renormalisation-group-1}, which has the properties~\eqref{it:renormalisation-group-1-k1}, \eqref{it:renormalisation-group-1-k2} and~\eqref{it:renormalisation-group-1-k4} of that definition; every $E\in K^{x,C}_k$ is such a family. Then, for all $n,m\ge0$ and all arguments $\xi=(\xi_1,\dots,\xi_n)$ and $\boldsymbol{\zeta}=(\boldsymbol{\zeta}_1,\dots,\boldsymbol{\zeta}_m)$,
\[
\left|E_{nm}(X;\mathscr{A}_L,\mathscr{C}_L;\xi,\boldsymbol{\zeta})\right|\le\exp\left(2(n+m)\right)\cdot\left|E_{nm}(X;\hat{b},\mathscr{C}^b;\xi,\boldsymbol{\zeta})\right|,
\]
and $E_{00}(X;\mathscr{A}_L,\mathscr{C}_L)=E_{00}(X;\hat{b},\mathscr{C}^b)$. More precisely, the modulus on the left is the modulus on the right multiplied by a product of $2(n+m)$ positive numbers, each at most $\exp(1)$, one for each argument and each of the two slots. Consequently
\[
\left\|E_{nm}(X;\mathscr{A}_L,\mathscr{C}_L)\right\|\le\exp\left(2(n+m)\right)\cdot\left\|E_{nm}(X;\hat{b},\mathscr{C}^b)\right\|.
\]
Suppose, in addition, that $\mathscr{C}=0$; then $\mathscr{C}^b=0$ and $c_2=0$ for the choice described before~(i). Let $E=\{E(Y;\cdot):Y\in D_k\}$ be a family of elements of $\mathrm{Gr}[\psi,\chi]$ that is defined on $\tilde{R}_k$ alone and written, in the form displayed in Definition~\ref{def:renormalisation-group-1} with the second slot suppressed, with antisymmetric kernels $E_{nm}(Y;\mathscr{A}';\xi,\boldsymbol{\zeta})$, $\mathscr{A}'\in\tilde{R}_k$. Assume that $E$ has the properties~\eqref{it:renormalisation-group-1-k1}, \eqref{it:renormalisation-group-1-k2} and~\eqref{it:renormalisation-group-1-k4} without the second slot, that is: for fixed $n$, $m$ and fixed arguments the kernel $E_{nm}(X;\mathscr{A}';\cdot)$ is an analytic function of $\mathscr{A}'$ on the open set $\tilde{R}_k$; property~\eqref{it:renormalisation-group-1-k2} holds with the single connection $\mathscr{A}'$ in place of the pair $(\mathscr{A},\mathscr{C})$; and, with the automorphisms $G_{\boldsymbol{\lambda}}$ of Section~\ref{sec:lattice-gauge-charge},
\[
G_{\boldsymbol{\lambda}}\left[E(X;\mathscr{A}^1)\right]=E(X;\mathscr{A}')
\]
whenever $\mathscr{A}'$ and $\mathscr{A}^1$ lie in $\tilde{R}_k$ and $\boldsymbol{\lambda}$ is a real function on the sites of $T_k$ with $\mathscr{A}^1=\mathscr{A}'-d\boldsymbol{\lambda}$ on the bonds at distance at most $1$ from $X$, which is~\eqref{it:renormalisation-group-1-k4} with $\mathscr{C}=\mathscr{C}^1=0$ and $\boldsymbol{\lambda}'=0$. Every $E\in K^C_k$ is such a family. Then, for all $n,m\ge0$ and all arguments $\xi=(\xi_1,\dots,\xi_n)$ and $\boldsymbol{\zeta}=(\boldsymbol{\zeta}_1,\dots,\boldsymbol{\zeta}_m)$, the modulus $\left|E_{nm}(X;\mathscr{A}_L;\xi,\boldsymbol{\zeta})\right|$ is $\left|E_{nm}(X;\hat{b};\xi,\boldsymbol{\zeta})\right|$ multiplied by a product of $n+m$ positive numbers, one for each argument, each at most $\exp(1)$; hence
\[
\left|E_{nm}(X;\mathscr{A}_L;\xi,\boldsymbol{\zeta})\right|\le\exp(n+m)\cdot\left|E_{nm}(X;\hat{b};\xi,\boldsymbol{\zeta})\right|\le\exp\left(2(n+m)\right)\cdot\left|E_{nm}(X;\hat{b};\xi,\boldsymbol{\zeta})\right|.
\]
Moreover $E_{00}(X;\mathscr{A}_L)=E_{00}(X;\hat{b})$, and
\[
\left\|E_{nm}(X;\mathscr{A}_L)\right\|\le\exp\left(2(n+m)\right)\cdot\left\|E_{nm}(X;\hat{b})\right\|.
\]
By~(ii), $t\cdot\hat{b}$ lies in $R_k$, hence in $\tilde{R}_k$, for every complex number $t$ with $|t|\le L^{\frac{7}{4}}$, so that $E(X;t\cdot\hat{b})$ is defined.
\end{enumerate}
\end{theorem}

\begin{proof}
\emph{Step~0: a decomposition of $\mathscr{A}$ near $x_0$.} By Definition~\ref{def:domains} at level $k+1$, one can write $\mathscr{A}=A_0+A_1$ with $A_1\in R_{k+1}$ and a real bond function $A_0$ which, on the enlarged cube of every $M$-cube of level $k+1$, is the sum of a real bond function with three norms below the thresholds of $R_{k+1}$ and of the differential of a real function. Fix such a decomposition and put $a_0:=(A_0)_L$ and $a_1:=(A_1)_L$. Then $a_0$ is real and $\mathscr{A}_L=a_0+a_1$. Let $\square'$ be the $M$-cube of level $k+1$ for which the $LM$-cube $L\square'$ of level $k$ contains $x_0$, and put $U:=L\square'^{\dagger}$. Lemma~\ref{lem:scaling-field-domains}(c), applied to the representation of $A_0$ on $\square'^{\dagger}$, provides a real function $\mathfrak{l}$ on the sites of $U$ and a real bond function $a_0'$ on $U$ with
\[
a_0=a_0'+d\mathfrak{l}\qquad\text{on }U,
\]
and Lemma~\ref{lem:scaling-field-domains}(a), (c), with the numbers $v_1$ and $w_1$ defined there, gives
\begin{gather*}
|a_1|<L^{-\frac{7}{8}+\frac{\varepsilon}{2}}\cdot t_{1,k},\qquad|\partial a_1|<v_1,\qquad[\partial a_1]_{\boldsymbol{\alpha}}<w_1\qquad\text{on }T_k,\\
|\partial a_0'|<v_1,\qquad[\partial a_0']_{\boldsymbol{\alpha}}<w_1\qquad\text{on }U.
\end{gather*}
No site lies on the boundary of an $LM$-cube (Section~\ref{sec:balls-torus-lifts}); hence the distance of $x_0$ from the centre of $L\square'$ is less than $\frac{1}{2}LM$. By Lemma~\ref{lem:scaling-field-domains}(c), $U$ consists of the sites whose distance from this centre is less than $\left(L+\frac{1}{2}\right)LM$. It follows that $K(x_0,r)\subset U$ whenever $r\le L^2M$. Since $L\ge5$, one has $2\rho+5=4LM+5\le L^2M$, and therefore
\[
K_2:=K(x_0,2\rho+5)\subset U.
\]
By Theorem~\ref{thm:polymer-activities-2}\eqref{it:polymers-t7}, $X\subset K\left(x_0,M(L+2)\right)$. If a bond has an end point at distance at most $1$ from $X$, then both of its end points have distance at most $2$ from $X$ and lie in $K\left(x_0,M(L+2)+2\right)$; as $M(L+2)+2\le2LM=\rho$, this proves the last assertion of~(i).

\emph{Step~1: Taylor remainders.} Since $l\ge13LM$ (because $l\in\mathrm{Cl}$), one has $2(2\rho+5)+1<l$; so the vectors $x-x_0$, $x\in K_2$, of Section~\ref{sec:balls-torus-lifts} are defined, and nearest neighbours and difference quotients in $K_2$ are those of the cube $\{z\in\mathbb{L}_k:|z|\le2\rho+5\}$. A bond function $f$ is described by its components $f_{\mu}(x)=f(x,x+\eta e_{\mu})$; the bond functions $\mathrm{Lin}[b]$ and $\mathrm{Ln}$ below are defined by formulas for these components, that is, for their values on the positively oriented bonds, written in terms of the initial site $x$ of the bond, and the same applies to the formula for $d\check{s}$ in Step~2.

For $b\in\{a_0',a_1\}$ and directions $\nu$, $\mu$ let $V[b]_{\nu\mu}:=(\partial_{\nu}b_{\mu})(x_0)$, and define on the bonds of $K_2$
\[
\mathrm{Lin}[b]_{\mu}(x):=\sum_{\nu}V[b]_{\nu\mu}\cdot(x-x_0)_{\nu},\qquad\mathsf{r}[b]:=b-b(x_0)-\mathrm{Lin}[b],
\]
where $b(x_0)$ is the constant bond function with the components $b_{\mu}(x_0)$. A constant bond function has vanishing derivatives, and $\partial_{\nu}\mathrm{Lin}[b]_{\mu}=V[b]_{\nu\mu}$; hence
\[
\partial_{\nu}\mathsf{r}[b]_{\mu}(x)=\partial_{\nu}b_{\mu}(x)-\partial_{\nu}b_{\mu}(x_0).
\]
Thus $\partial\mathsf{r}[b]$ has the H\"older quotients of $\partial b$, so that $[\partial\mathsf{r}[b]]_{\boldsymbol{\alpha}}\le w_1$ on $K_2$ by Step~0, and $\partial\mathsf{r}[b]$ vanishes at $x_0$. From every site of $K_2$ one reaches $x_0$ in at most $2\rho+6$ steps, each between two sites of $K_2$ at sup distance at most $1$: move all coordinates simultaneously towards those of $x_0$. In each step $\partial\mathsf{r}[b]$ changes by at most $w_1$, by the bound on $[\partial\mathsf{r}[b]]_{\boldsymbol{\alpha}}$. It follows that $|\partial\mathsf{r}[b]|\le(2\rho+6)\cdot w_1\le3\rho\cdot w_1$ on $K_2$. Moreover $\mathsf{r}[b]_{\mu}(x_0)=0$ and $|x-x_0|_1\le3(2\rho+5)$ for $x\in K_2$; summing the differences of $\mathsf{r}[b]$ along a rectilinear path from $x_0$ to $x$ one obtains
\[
|\mathsf{r}[b]|\le3(2\rho+5)(2\rho+6)\cdot w_1\le18\rho^2\cdot w_1\qquad\text{on }K_2,
\]
the last inequality because $\rho=2LM\ge14$. Lemma~\ref{lem:truncation-bond-function} applies with $r=\rho$, as $4\rho+11<13LM\le l$. With the three bounds just obtained for the norms of $\mathsf{r}[b]$ on $K_2$ it yields, for $b\in\{a_0',a_1\}$,
\begin{multline}\label{eq:polymer-activities-19}
|\boldsymbol{\chi}_{\rho}\mathsf{r}[b]|\le18\rho^2\cdot w_1,\qquad\left|\partial\left(\boldsymbol{\chi}_{\rho}\mathsf{r}[b]\right)\right|\le39\rho\cdot w_1,\\
\left[\partial\left(\boldsymbol{\chi}_{\rho}\mathsf{r}[b]\right)\right]_{\boldsymbol{\alpha}}\le(1+36+216)\,w_1=253\,w_1.
\end{multline}

\emph{Step~2: removal of the symmetric slope.} The matrix $V[a_0']$ is real; write it as $V[a_0']=s+\frac{1}{2}\mathfrak{f}$ with $s$ symmetric and $\mathfrak{f}$ antisymmetric. For a bond function $a$ put $(da)_{\nu\mu}:=\partial_{\nu}a_{\mu}-\partial_{\mu}a_{\nu}$. Then
\[
\mathfrak{f}_{\nu\mu}=\left(\partial_{\nu}a'_{0,\mu}-\partial_{\mu}a'_{0,\nu}\right)(x_0)=(da_0)_{\nu\mu}(x_0),
\]
the second equality because $a_0-a_0'=d\mathfrak{l}$ on $U$ and $d\,d\mathfrak{l}=0$. On the sites of $K_2$ define the real function
\[
\check{s}(x):=\sum_{\mu}\left(a'_{0,\mu}(x_0)-\frac{\eta}{2}s_{\mu\mu}\right)\cdot(x-x_0)_{\mu}+\frac{1}{2}\sum_{\nu,\mu}s_{\nu\mu}\cdot(x-x_0)_{\nu}\cdot(x-x_0)_{\mu}.
\]
The quadratic form $q(y):=\frac{1}{2}\sum_{\nu,\mu}s_{\nu\mu}y_{\nu}y_{\mu}$ on $\mathbb{R}^3$ satisfies $q(y+\eta e_{\mu})-q(y)=\eta\sum_{\nu}s_{\nu\mu}y_{\nu}+\frac{\eta^2}{2}s_{\mu\mu}$, because $s$ is symmetric. Hence, on the bonds of $K_2$,
\[
(d\check{s})_{\mu}(x)=a'_{0,\mu}(x_0)+\sum_{\nu}s_{\nu\mu}\cdot(x-x_0)_{\nu}.
\]
On the bonds of $K_2$ one has $a_0=a_0'+d\mathfrak{l}$ (Step~0), $a_0'=a_0'(x_0)+\mathrm{Lin}[a_0']+\mathsf{r}[a_0']$ and $a_1=a_1(x_0)+\mathrm{Lin}[a_1]+\mathsf{r}[a_1]$ (Step~1), and the components of $a_0'(x_0)+\mathrm{Lin}[a_0']-d\check{s}$ are $\frac{1}{2}\sum_{\nu}\mathfrak{f}_{\nu\mu}\cdot(x-x_0)_{\nu}$. Consequently
\begin{multline}\label{eq:polymer-activities-20}
a_0+a_1=d(\mathfrak{l}+\check{s})+a_1(x_0)+\mathrm{Ln}+\mathsf{r}[a_0']+\mathsf{r}[a_1]\qquad\text{on the bonds of }K_2,\\
\text{where}\qquad\mathrm{Ln}_{\mu}(x):=\sum_{\nu}\left(\frac{1}{2}\mathfrak{f}_{\nu\mu}+V[a_1]_{\nu\mu}\right)\cdot(x-x_0)_{\nu}.
\end{multline}

\emph{Step~3: bounds on the slopes.} By Step~0, $|V[a_1]_{\nu\mu}|\le|\partial a_1|<v_1$. The function $u:=\partial_{\nu}a_{1,\mu}$ on the sites of $T_k$ has sum zero (Lemma~\ref{lem:functions-vanishing-sum}), and $|u(x)-u(y)|\le w_1$ whenever $d(x,y)\le1$, since $[\partial a_1]_{\boldsymbol{\alpha}}<w_1$ on $T_k$; so Lemma~\ref{lem:functions-vanishing-sum} gives $|V[a_1]_{\nu\mu}|\le l\cdot w_1$. Hence
\[
|V[a_1]_{\nu\mu}|\le\min(v_1,l\cdot w_1).
\]
Now let $u:=(da_0)_{\nu\mu}=\partial_{\nu}a_{0,\mu}-\partial_{\mu}a_{0,\nu}$, a function on the sites of $T_k$ with sum zero by Lemma~\ref{lem:functions-vanishing-sum}. Let $x$ and $y$ be sites with $d(x,y)\le1$, and let $\square''$ be the $M$-cube of level $k+1$ for which the $LM$-cube $L\square''$ contains $x$. As in Step~0, Lemma~\ref{lem:scaling-field-domains}(c) provides a real bond function $a_0''$ on $L\square''^{\dagger}$ such that $a_0-a_0''$ is the differential of a real function on $L\square''^{\dagger}$ and $|\partial a_0''|<v_1$, $[\partial a_0'']_{\boldsymbol{\alpha}}<w_1$ there; in particular $da_0=da_0''$ on $L\square''^{\dagger}$. The sites $x$ and $y$ and the plaquettes that enter $u(x)$ and $u(y)$ lie within distance $3$ of $x$, hence in $L\square''^{\dagger}$. Therefore $|u(x)-u(y)|\le2w_1$ and $|u(x)|\le2v_1$, and Lemma~\ref{lem:functions-vanishing-sum} gives $|u|\le2l\cdot w_1$. Since $\mathfrak{f}_{\nu\mu}=u(x_0)$ by Step~2,
\[
|\mathfrak{f}_{\nu\mu}|\le2\min(v_1,l\cdot w_1).
\]
Consequently the coefficients $\frac{1}{2}\mathfrak{f}_{\nu\mu}+V[a_1]_{\nu\mu}$ of $\mathrm{Ln}$ are bounded in modulus by
\[
v':=2\min(v_1,l\cdot w_1).
\]

\emph{Step~4: construction of $\hat{b}$, $c_1$ and $\mathring{m}$.} By~\eqref{eq:polymer-activities-Sb}, $e_k<e_{k+1}\le e_0$; therefore the inequalities $160L^{1+\frac{3\varepsilon}{2}}M\cdot e_0^{\varepsilon}\le1$ and $9M\cdot L^{\varepsilon}\cdot e_0^{2\varepsilon}\le1$, the first and the third inequality of~\eqref{eq:renormalisation-group-A9}, remain true with $e_k$ in place of $e_0$. Put
\[
\rho':=\min\left(e_k^{-\varepsilon},\frac{l}{5}\right).
\]
Then $\rho'\ge\rho\ge50$: indeed $2LM\cdot e_k^{\varepsilon}\le160L^{1+\frac{3\varepsilon}{2}}M\cdot e_k^{\varepsilon}\le1$ and $\frac{l}{5}\ge\frac{13}{5}LM$, and $\rho=2LM\ge50$ because $L\ge5$ and $M\ge L$. Moreover $4\rho'+11\le\frac{4}{5}l+11<l$, because $l\ge13LM>55$. In particular $2(2\rho'+5)+1<l$, so the formula for $\mathrm{Ln}$ in~\eqref{eq:polymer-activities-20} defines $\mathrm{Ln}$ on all bonds of the ball $K(x_0,2\rho'+5)$, with the vectors $x-x_0$ of Section~\ref{sec:balls-torus-lifts}. Define
\[
\hat{b}:=\boldsymbol{\chi}_{\rho}\left(\mathsf{r}[a_0']+\mathsf{r}[a_1]\right)+\boldsymbol{\chi}_{\rho'}\mathrm{Ln},\qquad c_1:=a_1(x_0),
\]
and let $\mathring{m}:=\mathfrak{l}+\check{s}$ on the sites of $K_2$ and $\mathring{m}:=0$ on all other sites of $T_k$. The function $\mathring{m}$ is real, because $\mathfrak{l}$ and $a_0'$, hence $s$ and $\check{s}$, are real. By Step~0, $|c_1|\le|a_1|<L^{-\frac{7}{8}+\frac{\varepsilon}{2}}\cdot t_{1,k}$. By Lemma~\ref{lem:truncation-bond-function}, with $r=\rho$ and with $r=\rho'$, the cut-offs $\boldsymbol{\chi}_{\rho}$ and $\boldsymbol{\chi}_{\rho'}$ do not change a bond function on the bonds whose end points lie in $K(x_0,\rho)\subset K(x_0,\rho')$. On these bonds, moreover, $d\mathring{m}=d(\mathfrak{l}+\check{s})$, both end points lying in $K_2$. So~\eqref{eq:polymer-activities-20} gives
\[
\mathscr{A}_L-d\mathring{m}=\hat{b}+c_1\qquad\text{on the bonds with both end points in }K(x_0,\rho),
\]
which is the first identity of~(i).

\emph{Step~5: estimates for $\hat{b}$.} For $x\in K(x_0,2\rho'+5)$ one has $|x-x_0|_1\le3(2\rho'+5)$; by Step~3 and since $\rho'\ge50$,
\[
|\mathrm{Ln}|\le3(2\rho'+5)\cdot v'\le6.3\rho'\cdot v',\qquad|\partial\mathrm{Ln}|\le v',\qquad[\partial\mathrm{Ln}]_{\boldsymbol{\alpha}}=0
\]
on $K(x_0,2\rho'+5)$, the derivatives of $\mathrm{Ln}$ being its constant coefficients. Lemma~\ref{lem:truncation-bond-function} with $r=\rho'$ yields
\[
|\boldsymbol{\chi}_{\rho'}\mathrm{Ln}|\le6.3\rho'\cdot v',\qquad\left|\partial\left(\boldsymbol{\chi}_{\rho'}\mathrm{Ln}\right)\right|\le(1+12.6)\,v'=13.6\,v',\qquad\left[\partial\left(\boldsymbol{\chi}_{\rho'}\mathrm{Ln}\right)\right]_{\boldsymbol{\alpha}}\le(12+75.6)\,\frac{v'}{\rho'}=87.6\,\frac{v'}{\rho'}.
\]
Recall from Lemma~\ref{lem:scaling-field-domains}(a) and Section~\ref{sec:settings-basic-notation} that
\[
v_1=L^{-\frac{15}{8}+\varepsilon}\cdot t_{2,k},\qquad w_1=L^{-\frac{15}{8}+\frac{3\varepsilon}{2}}\cdot t_{3,k},\qquad t_{2,k}=e_k^{\varepsilon}\cdot t_{1,k},\qquad t_{3,k}=e_k^{\varepsilon}\cdot t_{2,k}.
\]
If $\rho'=e_k^{-\varepsilon}$, then, with $v'\le2v_1$, the three bounds for $\boldsymbol{\chi}_{\rho'}\mathrm{Ln}$ are at most
\[
12.6\,L^{-\frac{15}{8}+\varepsilon}\cdot t_{1,k},\qquad27.2\,L^{-\frac{15}{8}+\varepsilon}\cdot t_{2,k},\qquad175.2\,L^{-\frac{15}{8}+\varepsilon}\cdot t_{3,k}.
\]
If $\rho'=\frac{l}{5}<e_k^{-\varepsilon}$, then, with $v'\le2l\cdot w_1$ and $l<5e_k^{-\varepsilon}$, they are at most
\begin{gather*}
2.52\,l^2\cdot w_1\le63\,L^{-\frac{15}{8}+\frac{3\varepsilon}{2}}\cdot t_{1,k},\qquad27.2\,l\cdot w_1\le136\,L^{-\frac{15}{8}+\frac{3\varepsilon}{2}}\cdot t_{2,k},\\
876\,w_1=876\,L^{-\frac{15}{8}+\frac{3\varepsilon}{2}}\cdot t_{3,k}.
\end{gather*}
By~\eqref{eq:polymer-activities-19}, taken for $b=a_0'$ and for $b=a_1$, the three norms of $\boldsymbol{\chi}_{\rho}\left(\mathsf{r}[a_0']+\mathsf{r}[a_1]\right)$ are at most
\begin{gather*}
36\rho^2\cdot w_1=144\left(LM\cdot e_k^{\varepsilon}\right)^2\cdot L^{-\frac{15}{8}+\frac{3\varepsilon}{2}}\cdot t_{1,k},\qquad78\rho\cdot w_1=156\,LM\cdot e_k^{\varepsilon}\cdot L^{-\frac{15}{8}+\frac{3\varepsilon}{2}}\cdot t_{2,k},\\
506\,w_1=506\,L^{-\frac{15}{8}+\frac{3\varepsilon}{2}}\cdot t_{3,k},
\end{gather*}
and $144\left(LM\cdot e_k^{\varepsilon}\right)^2\le1$, $156\,LM\cdot e_k^{\varepsilon}\le1$ by the first inequality of~\eqref{eq:renormalisation-group-A9}. Since $L^{-\frac{15}{8}+\varepsilon}\le L^{-\frac{15}{8}+\frac{3\varepsilon}{2}}$, the two summands of $\hat{b}$ together give, in either case,
\[
|\hat{b}|\le64\,L^{-\frac{15}{8}+\frac{3\varepsilon}{2}}\cdot t_{1,k},\qquad|\partial\hat{b}|\le137\,L^{-\frac{15}{8}+\frac{3\varepsilon}{2}}\cdot t_{2,k},\qquad[\partial\hat{b}]_{\boldsymbol{\alpha}}\le1382\,L^{-\frac{15}{8}+\frac{3\varepsilon}{2}}\cdot t_{3,k}.
\]
They imply the bounds for $\hat{b}$ asserted in~(ii). Furthermore
\[
1382\,L^{-\frac{15}{8}+\frac{3\varepsilon}{2}}<L^{-\frac{7}{4}},
\]
because $L^{\frac{1}{8}-\frac{3\varepsilon}{2}}\ge L^{\frac{1}{8}-\frac{3\bar{\varepsilon}}{2}}\ge2400$, by $\varepsilon\le\bar{\varepsilon}$ (stage~\ref{it:order-choice-8}) and by the first inequality of stage~\ref{it:order-choice-4}.

\emph{Step~6: the second slot.} Let $c_2:=\mathscr{C}_L(x_0)$ be the constant bond function with the components $(\mathscr{C}_L)_{\mu}(x_0)$, and put
\[
\mathscr{C}^b:=\boldsymbol{\chi}_{\rho}\left(\mathscr{C}_L-c_2\right).
\]
Thus $c_2$ and $\mathscr{C}^b$ are formed from $\mathscr{C}_L$ and $x_0$ alone. If $\mathscr{C}=0$, then $\mathscr{C}_L=0$, the scaling being linear, so that $c_2=\mathscr{C}_L(x_0)=0$ and $\mathscr{C}^b=\boldsymbol{\chi}_{\rho}(0)=0$; indeed, each component $(\boldsymbol{\chi}_{\rho}b)_{\mu}(x)$ of a truncated bond function is either $f_{\mu}(x)\,b_{\mu}(x)$ or $0$ (Lemma~\ref{lem:truncation-bond-function}). The connection $\mathscr{C}$ does not occur in Steps~0--5: the decomposition $\mathscr{A}=A_0+A_1$, the representation $a_0=a_0'+d\mathfrak{l}$ on $U$, the matrices $V[a_0']$, $V[a_1]$, $s$ and $\mathfrak{f}$, the bond functions $\mathrm{Lin}[a_0']$, $\mathrm{Lin}[a_1]$, $\mathsf{r}[a_0']$, $\mathsf{r}[a_1]$ and $\mathrm{Ln}$, the functions $\mathfrak{l}$ and $\check{s}$, and the radius $\rho'$ are formed from $\mathscr{A}$, $x_0$ and the data $k$, $l$, $e_k$ of~\eqref{eq:polymer-activities-Sb}. So $\hat{b}$, $c_1=a_1(x_0)$ and $\mathring{m}$ can be fixed as functions of $(\mathscr{A},x_0)$. No family $E$ enters the construction of the five objects in Steps~0--6. This proves the description of $\hat{b}$, $c_1$, $\mathring{m}$, $c_2$ and $\mathscr{C}^b$ given in the statement before~(i).

By Lemma~\ref{lem:scaling-field-domains}(b), $|c_2|\le|\mathscr{C}_L|<L^{-1+\varepsilon}\cdot x_{1,k}$. Summing along rectilinear paths from $x_0$, as in Step~1, one finds
\[
|\mathscr{C}_L-c_2|\le3(2\rho+5)\cdot|\partial\mathscr{C}_L|\le6.3\rho\cdot|\partial\mathscr{C}_L|\qquad\text{on }K_2,
\]
the second inequality because $\rho\ge50$; the derivatives of $\mathscr{C}_L-c_2$ are those of $\mathscr{C}_L$. Lemma~\ref{lem:truncation-bond-function} with $r=\rho$ and Lemma~\ref{lem:scaling-field-domains}(b), which gives $|\partial\mathscr{C}_L|<L^{-2+\frac{3\varepsilon}{2}}\cdot x_{2,k}$, together with $6.3\rho=12.6\,LM$ and $x_{2,k}=e_k^{\varepsilon}\cdot x_{1,k}$, yield
\[
|\mathscr{C}^b|\le6.3\rho\cdot|\partial\mathscr{C}_L|<12.6\,M\cdot L^{-1+\frac{3\varepsilon}{2}}\cdot e_k^{\varepsilon}\cdot x_{1,k}\le L^{-\frac{7}{4}}\cdot x_{1,k};
\]
the last inequality holds because, by the first inequality of~\eqref{eq:renormalisation-group-A9},
\[
12.6\,M\cdot L^{\frac{3}{4}+\frac{3\varepsilon}{2}}\cdot e_k^{\varepsilon}=\frac{12.6}{160}\,L^{-\frac{1}{4}}\cdot160L^{1+\frac{3\varepsilon}{2}}M\cdot e_k^{\varepsilon}\le\frac{12.6}{160}\,L^{-\frac{1}{4}}\le1.
\]
In the same way
\[
|\partial\mathscr{C}^b|\le\left(1+\frac{2}{\rho}\cdot6.3\rho\right)|\partial\mathscr{C}_L|=13.6\,|\partial\mathscr{C}_L|<13.6\,L^{-2+\frac{3\varepsilon}{2}}\cdot x_{2,k}\le L^{-\frac{7}{4}}\cdot x_{2,k},
\]
where the last step is the inequality $13.6\le L^{\frac{1}{4}-\frac{3\varepsilon}{2}}$, a consequence of stages~\ref{it:order-choice-4} and~\ref{it:order-choice-8}. Lemma~\ref{lem:truncation-bond-function} also shows that $\mathscr{C}^b=\mathscr{C}_L-c_2$ on the bonds with both end points in $K(x_0,\rho)$; this is the second identity of~(i), and~(i) is proved.

Let $t$ be a complex number with $|t|\le L^{\frac{7}{4}}$. By Step~5 and the last two bounds, the three norms of $t\cdot\hat{b}$ are less than $t_{1,k}$, $t_{2,k}$, $t_{3,k}$, and $|t\cdot\mathscr{C}^b|<x_{1,k}$, $|\partial(t\cdot\mathscr{C}^b)|<x_{2,k}$. By Definition~\ref{def:domains}, $(t\cdot\hat{b},t\cdot\mathscr{C}^b)\in R_k\times\mathfrak{X}_k$. Together with the bounds for $|c_1|$ in Step~4 and for $|c_2|$ above, this proves~(ii).

\emph{Step~7: comparison of the kernels.} Let $E$ be a family as in~(iii). By Lemma~\ref{lem:scaling-field-domains}(c), (b), $\mathscr{A}_L\in\tilde{R}_k$ and $\mathscr{C}_L\in\mathfrak{X}_k$; as $\mathring{m}$ is real, also $\mathscr{A}_L-d\mathring{m}\in\tilde{R}_k$ (Proposition~\ref{prop:properties-actions}(c)). By stages~\ref{it:order-choice-4} and~\ref{it:order-choice-8}, $L^{-\frac{7}{4}}+4L^{-\frac{7}{8}+\frac{\varepsilon}{2}}<1$, and $L^{-1+\varepsilon}\le L^{-\frac{7}{8}+\frac{\varepsilon}{2}}$. Hence $|\hat{b}+c_1|<\left(L^{-\frac{7}{4}}+L^{-\frac{7}{8}+\frac{\varepsilon}{2}}\right)t_{1,k}<t_{1,k}$, and since a constant bond function has vanishing derivatives, Step~5 shows $\hat{b}+c_1\in R_k\subset\tilde{R}_k$. Likewise $|\mathscr{C}^b+c_2|<\left(L^{-\frac{7}{4}}+L^{-1+\varepsilon}\right)x_{1,k}<x_{1,k}$ and Step~6 show $\mathscr{C}^b+c_2\in\mathfrak{X}_k$. By~(i), the identities $\mathscr{A}_L-d\mathring{m}=\hat{b}+c_1$ and $\mathscr{C}_L-d0=\mathscr{C}^b+c_2$ hold on the bonds at distance at most $1$ from $X$. Property~\eqref{it:renormalisation-group-1-k4}, applied to the pairs $(\mathscr{A}_L,\mathscr{C}_L)$ and $(\hat{b}+c_1,\mathscr{C}^b+c_2)$ with the real functions $\boldsymbol{\lambda}:=\mathring{m}$ and $\boldsymbol{\lambda}':=0$, therefore gives
\[
G_{\mathring{m}}\left[E(X;\hat{b}+c_1,\mathscr{C}^b+c_2)\right]=E(X;\mathscr{A}_L,\mathscr{C}_L),
\]
with the automorphism $G_{\mathring{m}}$ of Section~\ref{sec:lattice-gauge-charge}. Since $\mathring{m}$ is real, $G_{\mathring{m}}$ multiplies every generator by a number of modulus one, and the kernels of the transformed element are the old kernels times the product of these numbers over the arguments; this product is symmetric in the arguments, so the new kernels are again antisymmetric. As the antisymmetric kernels are determined by the element (Proposition~\ref{prop:weighted-coefficient-norms}\eqref{it:weighted-coefficient-norms-N0}), the kernels of $E(X;\mathscr{A}_L,\mathscr{C}_L)$ have the same moduli as those of $E(X;\hat{b}+c_1,\mathscr{C}^b+c_2)$, and the two parts of bidegree $(0,0)$ are equal.

It remains to compare the kernels at $(\hat{b}+c_1,\mathscr{C}^b+c_2)$ with those at $(\hat{b},\mathscr{C}^b)$. Let $c_j^R$ and $c_j^I$ be the real constant bond functions with $c_j=c_j^R+i\,c_j^I$, $j=1,2$; then $|c_j^R|\le|c_j|$ and $|c_j^I|\le|c_j|$. For a constant bond function $c$ and $y\in\mathbb{R}^3$ put $\langle c,y\rangle:=\sum_{\mu}c_{\mu}y_{\mu}$, so that $|\langle c,y\rangle|\le|c|\cdot|y|_1$. Let
\[
\mathscr{D}:=\left\{z\in\mathbb{C}^4:|\operatorname{Im}z_1|<2,\ |\operatorname{Im}z_2|<2,\ |z_3|<2,\ |z_4|<2\right\},
\]
and for $z\in\mathscr{D}$ define the bond functions
\[
A(z):=\hat{b}+z_1\cdot c_1^R+z_2\cdot c_1^I,\qquad C(z):=\mathscr{C}^b+z_3\cdot c_2^R+z_4\cdot c_2^I,
\]
and the function $\Lambda_z$ on the sites of $T_k$ by
\[
\Lambda_z(x):=-\left\langle z_1\cdot c_1^R+z_2\cdot c_1^I+z_3\cdot c_2^R+z_4\cdot c_2^I,\,x-x_0\right\rangle\quad\text{for }x\in K(x_0,\rho),\qquad\Lambda_z(x):=0\quad\text{otherwise}.
\]
For every $z\in\mathscr{D}$ the pair $(A(z),C(z))$ lies in $\tilde{R}_k\times\mathfrak{X}_k$. Indeed, $A(z)$ is the sum of the real constant bond function $\operatorname{Re}z_1\cdot c_1^R+\operatorname{Re}z_2\cdot c_1^I$ and of $\hat{b}+i\left(\operatorname{Im}z_1\cdot c_1^R+\operatorname{Im}z_2\cdot c_1^I\right)$. The latter lies in $R_k$, since its modulus is less than $\left(L^{-\frac{7}{4}}+4L^{-\frac{7}{8}+\frac{\varepsilon}{2}}\right)t_{1,k}<t_{1,k}$ and its derivatives are those of $\hat{b}$; and an element of $R_k\subset\tilde{R}_k$ plus a real constant bond function lies in $\tilde{R}_k$ by Proposition~\ref{prop:properties-actions}(c). Further, $|C(z)|<\left(L^{-\frac{7}{4}}+4L^{-1+\varepsilon}\right)x_{1,k}<x_{1,k}$ and $\partial C(z)=\partial\mathscr{C}^b$, so that $C(z)\in\mathfrak{X}_k$. The domain $\mathscr{D}$ is open and convex, hence connected, and the map $z\mapsto(A(z),C(z))$ is affine.

Fix $n,m\ge0$ and arguments $\xi=(\xi_1,\dots,\xi_n)$, $\boldsymbol{\zeta}=(\boldsymbol{\zeta}_1,\dots,\boldsymbol{\zeta}_m)$. For an argument $(x,\boldsymbol{\beta},\boldsymbol{\omega})$ of $\psi$ let $x$ be its site, and for an argument $(x,x',\boldsymbol{\beta},\boldsymbol{\omega})$ of $\chi$ let $x$ be the first site of its pair; in both cases $\boldsymbol{\omega}$ is the bar index. Define, for $z\in\mathscr{D}$,
\[
\mathscr{K}(z):=E_{nm}\left(X;A(z),C(z);\xi,\boldsymbol{\zeta}\right)\cdot\prod\exp\left((-1)^{\boldsymbol{\omega}}\,i\,e_k\cdot\Lambda_z(x)\right),
\]
the product being taken over the $n+m$ arguments, with $x$ and $\boldsymbol{\omega}$ attached to each argument as just described. The kernel is an analytic function on the open set $\tilde{R}_k\times\mathfrak{X}_k$ of Definition~\ref{def:domains}, by~\eqref{it:renormalisation-group-1-k1}; its composition with the affine map $z\mapsto(A(z),C(z))$ is holomorphic on $\mathscr{D}$, and the exponential factors are entire functions of $z$. So $\mathscr{K}$ is holomorphic on $\mathscr{D}$.

Let $z\in\mathscr{D}\cap\mathbb{R}^4$. Then $\Lambda_z$ is real, and $\Lambda_z=\boldsymbol{\lambda}+\boldsymbol{\lambda}'$ with the real functions
\[
\boldsymbol{\lambda}(x):=-\left\langle z_1\cdot c_1^R+z_2\cdot c_1^I,\,x-x_0\right\rangle,\qquad\boldsymbol{\lambda}'(x):=-\left\langle z_3\cdot c_2^R+z_4\cdot c_2^I,\,x-x_0\right\rangle
\]
for $x\in K(x_0,\rho)$ and $\boldsymbol{\lambda}(x)=\boldsymbol{\lambda}'(x):=0$ for the other sites. On the bonds with both end points in $K(x_0,\rho)$ the differences of $x-x_0$ are those of $\mathbb{L}_k$ (Section~\ref{sec:balls-torus-lifts}); there $d\boldsymbol{\lambda}=-(z_1\cdot c_1^R+z_2\cdot c_1^I)$ and $d\boldsymbol{\lambda}'=-(z_3\cdot c_2^R+z_4\cdot c_2^I)$, that is,
\[
A(z)=\hat{b}-d\boldsymbol{\lambda},\qquad C(z)=\mathscr{C}^b-d\boldsymbol{\lambda}'.
\]
By Step~0 this holds in particular on the bonds at distance at most $1$ from $X$. Both $(\hat{b},\mathscr{C}^b)$, by~(ii) with $t=1$, and $(A(z),C(z))$ lie in $\tilde{R}_k\times\mathfrak{X}_k$. Thus~\eqref{it:renormalisation-group-1-k4} gives
\[
G_{\Lambda_z}\left[E(X;A(z),C(z))\right]=E(X;\hat{b},\mathscr{C}^b).
\]
As functions of the arguments, the kernels of the left side are the antisymmetric functions $\mathscr{K}(z)$, and those of the right side are $E_{nm}(X;\hat{b},\mathscr{C}^b;\xi,\boldsymbol{\zeta})=\mathscr{K}(0)$, because $A(0)=\hat{b}$, $C(0)=\mathscr{C}^b$ and $\Lambda_0=0$. By Proposition~\ref{prop:weighted-coefficient-norms}\eqref{it:weighted-coefficient-norms-N0}, $\mathscr{K}(z)=\mathscr{K}(0)$ for every $z\in\mathscr{D}\cap\mathbb{R}^4$.

The function $\mathscr{K}-\mathscr{K}(0)$ is holomorphic on the connected open set $\mathscr{D}$ and is zero on $\mathscr{D}\cap\mathbb{R}^4=\mathbb{R}^2\times(-2,2)^2$, a non-empty open subset of $\mathbb{R}^4$. At a point of $\mathscr{D}\cap\mathbb{R}^4$ its complex partial derivatives of every positive order are the corresponding partial derivatives along $\mathbb{R}^4$, and these vanish. By the identity theorem for holomorphic functions of several variables~\eqref{eq:polymer-activities-cl2}, $\mathscr{K}(z)=\mathscr{K}(0)$ for all $z\in\mathscr{D}$.

The point $z=(1,i,1,i)$ lies in $\mathscr{D}$. For it, $A(z)=\hat{b}+c_1$, $C(z)=\mathscr{C}^b+c_2$, and $\Lambda_z(x)=-\langle c_1+c_2,x-x_0\rangle$, with imaginary part $-\langle c_1^I+c_2^I,x-x_0\rangle$, for $x\in K(x_0,\rho)$. Hence
\[
E_{nm}\left(X;\hat{b}+c_1,\mathscr{C}^b+c_2;\xi,\boldsymbol{\zeta}\right)\cdot\prod\exp\left((-1)^{\boldsymbol{\omega}}\,i\,e_k\cdot\Lambda_z(x)\right)=E_{nm}\left(X;\hat{b},\mathscr{C}^b;\xi,\boldsymbol{\zeta}\right).
\]
By~\eqref{it:renormalisation-group-1-k2}, both kernels vanish unless the site of every $\xi_i$ lies in $X$ and every pair of the arguments $\boldsymbol{\zeta}_j$ is admissible for $X$. In that case the site $x$ of every argument has distance at most $\eta\le1$ from $X$ (Definition~\ref{def:pairs-sites-holder}); by Step~0, $x\in K\left(x_0,M(L+2)+1\right)\subset K(x_0,\rho)$ and
\[
|x-x_0|_1\le3\left(M(L+2)+1\right)\le9LM.
\]
By~(ii), $e_k\cdot|c_1|\le L^{-\frac{7}{8}+\frac{\varepsilon}{2}}\cdot e_k^{\frac{1}{4}+\varepsilon}$ and $e_k\cdot|c_2|\le L^{-1+\varepsilon}\cdot e_k^{2\varepsilon}$. Therefore
\[
e_k\cdot|c_1^I|\cdot9LM\le9LM\cdot L^{-\frac{7}{8}+\frac{\varepsilon}{2}}\cdot e_k^{\frac{1}{4}+\varepsilon}\le1,\qquad e_k\cdot|c_2^I|\cdot9LM\le9M\cdot L^{\varepsilon}\cdot e_k^{2\varepsilon}\le1:
\]
the middle term of the first chain equals $\frac{9}{160}\,L^{-\frac{7}{8}-\varepsilon}\cdot e_k^{\frac{1}{4}}\cdot160L^{1+\frac{3\varepsilon}{2}}M\cdot e_k^{\varepsilon}$, which is at most $\frac{9}{160}$ by the first inequality of~\eqref{eq:renormalisation-group-A9}, since this inequality also gives $e_k<1$ and hence $L^{-\frac{7}{8}-\varepsilon}\cdot e_k^{\frac{1}{4}}\le1$; the second chain ends with the third inequality of~\eqref{eq:renormalisation-group-A9}. Now $\exp\left((-1)^{\boldsymbol{\omega}}\,i\,e_k\cdot\Lambda_z(x)\right)$ is the product of the two numbers $\exp\left(-(-1)^{\boldsymbol{\omega}}\,i\,e_k\cdot\langle c_j,x-x_0\rangle\right)$, $j=1,2$, of the moduli $\exp\left((-1)^{\boldsymbol{\omega}}\,e_k\cdot\langle c_j^I,x-x_0\rangle\right)$, and $e_k\cdot|\langle c_j^I,x-x_0\rangle|\le e_k\cdot|c_j^I|\cdot|x-x_0|_1\le1$. So each of the two numbers has modulus between $\exp(-1)$ and $\exp(1)$, one for each slot, and each of the $n+m$ factors of the product has modulus between $\exp(-2)$ and $\exp(2)$. It follows that
\[
\left|E_{nm}\left(X;\hat{b}+c_1,\mathscr{C}^b+c_2;\xi,\boldsymbol{\zeta}\right)\right|\le\exp\left(2(n+m)\right)\cdot\left|E_{nm}\left(X;\hat{b},\mathscr{C}^b;\xi,\boldsymbol{\zeta}\right)\right|
\]
for all arguments, both sides vanishing for the arguments excluded above. For $n=m=0$ the product is empty, and $E_{00}(X;\hat{b}+c_1,\mathscr{C}^b+c_2)=E_{00}(X;\hat{b},\mathscr{C}^b)$. Combined with the comparison of $E(X;\mathscr{A}_L,\mathscr{C}_L)$ and $E(X;\hat{b}+c_1,\mathscr{C}^b+c_2)$ above, this proves the assertions of~(iii) on the kernels. Integration over the arguments, with the kernel norm of Section~\ref{sec:grassmann-algebras-norms}, gives the bound for $\|E_{nm}(X;\mathscr{A}_L,\mathscr{C}_L)\|$.

It remains to prove the assertions of~(iii) on families without the second slot. Let $\mathscr{C}=0$, and let $E=\{E(Y;\cdot):Y\in D_k\}$ be a family on $\tilde{R}_k$ alone as in~(iii). Then $\mathscr{C}_L=0$, and $c_2=0$, $\mathscr{C}^b=0$ by Step~6. Hence $c_2^R=c_2^I=0$ and $C(z)=0$ for every $z\in\mathscr{D}$, the function $\boldsymbol{\lambda}'$ above vanishes for every $z\in\mathscr{D}\cap\mathbb{R}^4$, and
\[
\Lambda_z(x)=-\left\langle z_1\cdot c_1^R+z_2\cdot c_1^I,\,x-x_0\right\rangle\quad\text{for }x\in K(x_0,\rho),\qquad\Lambda_z(x)=0\quad\text{otherwise},
\]
for every $z\in\mathscr{D}$. Thus each of the pairs at which a family defined on $\tilde{R}_k\times\mathfrak{X}_k$ has been evaluated earlier in this step, namely $(\mathscr{A}_L,\mathscr{C}_L)$, $(\hat{b}+c_1,\mathscr{C}^b+c_2)$, $(A(z),C(z))$ for $z\in\mathscr{D}$, and $(\hat{b},\mathscr{C}^b)$, has the second entry $0$. The bond function $\mathscr{A}_L$ lies in $\tilde{R}_k$ by Lemma~\ref{lem:scaling-field-domains}(c), and $\hat{b}+c_1\in R_k\subset\tilde{R}_k$, as shown at the beginning of this step. By~(i), $\mathscr{A}_L-d\mathring{m}=\hat{b}+c_1$ on the bonds at distance at most $1$ from $X$. The property~\eqref{it:renormalisation-group-1-k4} without the second slot, applied with $\mathscr{A}'=\mathscr{A}_L$, $\mathscr{A}^1=\hat{b}+c_1$ and $\boldsymbol{\lambda}=\mathring{m}$, gives
\[
G_{\mathring{m}}\left[E(X;\hat{b}+c_1)\right]=E(X;\mathscr{A}_L).
\]
Since $\mathring{m}$ is real, $G_{\mathring{m}}$ multiplies every generator by a number of modulus one, and the kernels of the transformed element are the old kernels multiplied by the product of these numbers over the arguments, which is symmetric in the arguments. By Proposition~\ref{prop:weighted-coefficient-norms}\eqref{it:weighted-coefficient-norms-N0}, the kernels of $E(X;\mathscr{A}_L)$ therefore have the same moduli as those of $E(X;\hat{b}+c_1)$, and the two parts of bidegree $(0,0)$ are equal.

Fix $n,m\ge0$ and arguments $\xi=(\xi_1,\dots,\xi_n)$, $\boldsymbol{\zeta}=(\boldsymbol{\zeta}_1,\dots,\boldsymbol{\zeta}_m)$, with the sites $x$ and the bar indices $\boldsymbol{\omega}$ attached to them as above, and define now, for $z\in\mathscr{D}$,
\[
\mathscr{K}(z):=E_{nm}\left(X;A(z);\xi,\boldsymbol{\zeta}\right)\cdot\prod\exp\left((-1)^{\boldsymbol{\omega}}\,i\,e_k\cdot\Lambda_z(x)\right),
\]
the product being taken over the $n+m$ arguments. As shown above, $A(z)\in\tilde{R}_k$ for every $z\in\mathscr{D}$. By~\eqref{it:renormalisation-group-1-k1} without the second slot, the kernel is an analytic function on the open set $\tilde{R}_k$; its composition with the affine map $z\mapsto A(z)$ is holomorphic on $\mathscr{D}$, and the exponential factors are entire functions of $z$. So $\mathscr{K}$ is holomorphic on $\mathscr{D}$. Let $z\in\mathscr{D}\cap\mathbb{R}^4$. Then $\Lambda_z$ is real, and $\Lambda_z=\boldsymbol{\lambda}$ because $\boldsymbol{\lambda}'=0$; so, as shown above, $A(z)=\hat{b}-d\Lambda_z$ on the bonds with both end points in $K(x_0,\rho)$, in particular on the bonds at distance at most $1$ from $X$. Both $\hat{b}$, which lies in $R_k\subset\tilde{R}_k$ by~(ii) with $t=1$, and $A(z)$ lie in $\tilde{R}_k$. The property~\eqref{it:renormalisation-group-1-k4} without the second slot, applied with $\mathscr{A}'=\hat{b}$, $\mathscr{A}^1=A(z)$ and $\boldsymbol{\lambda}=\Lambda_z$, gives
\[
G_{\Lambda_z}\left[E(X;A(z))\right]=E(X;\hat{b}).
\]
The kernels of the left side are the antisymmetric functions $\mathscr{K}(z)$, and those of the right side are $E_{nm}(X;\hat{b};\xi,\boldsymbol{\zeta})=\mathscr{K}(0)$, because $A(0)=\hat{b}$ and $\Lambda_0=0$; by Proposition~\ref{prop:weighted-coefficient-norms}\eqref{it:weighted-coefficient-norms-N0}, $\mathscr{K}(z)=\mathscr{K}(0)$. The function $\mathscr{K}-\mathscr{K}(0)$ is therefore holomorphic on the connected open set $\mathscr{D}$ and zero on the non-empty open subset $\mathscr{D}\cap\mathbb{R}^4$ of $\mathbb{R}^4$; at the points of this subset its complex partial derivatives of every positive order are partial derivatives along $\mathbb{R}^4$ and vanish. By the identity theorem~\eqref{eq:polymer-activities-cl2}, $\mathscr{K}(z)=\mathscr{K}(0)$ for all $z\in\mathscr{D}$.

At the point $z=(1,i,1,i)$ of $\mathscr{D}$ one has $A(z)=\hat{b}+c_1$ and $\Lambda_z(x)=-\langle c_1,x-x_0\rangle$ for $x\in K(x_0,\rho)$. Hence
\[
E_{nm}\left(X;\hat{b}+c_1;\xi,\boldsymbol{\zeta}\right)\cdot\prod\exp\left((-1)^{\boldsymbol{\omega}}\,i\,e_k\cdot\Lambda_z(x)\right)=E_{nm}\left(X;\hat{b};\xi,\boldsymbol{\zeta}\right).
\]
By~\eqref{it:renormalisation-group-1-k2}, both kernels vanish unless the site of every $\xi_i$ lies in $X$ and every pair of the arguments $\boldsymbol{\zeta}_j$ is admissible for $X$; in that case, as shown above, the site $x$ of every argument lies in $K(x_0,\rho)$ and satisfies $|x-x_0|_1\le9LM$. For these arguments each of the $n+m$ factors of the product is therefore $\exp\left(-(-1)^{\boldsymbol{\omega}}\,i\,e_k\cdot\langle c_1,x-x_0\rangle\right)$, a number of modulus $\exp\left((-1)^{\boldsymbol{\omega}}\,e_k\cdot\langle c_1^I,x-x_0\rangle\right)$, and $e_k\cdot|\langle c_1^I,x-x_0\rangle|\le e_k\cdot|c_1^I|\cdot9LM\le1$ by the first of the two chains above; so this modulus lies between $\exp(-1)$ and $\exp(1)$. Consequently $\left|E_{nm}(X;\hat{b}+c_1;\xi,\boldsymbol{\zeta})\right|$ is $\left|E_{nm}(X;\hat{b};\xi,\boldsymbol{\zeta})\right|$ multiplied by the reciprocals of these $n+m$ moduli, which are positive numbers, one for each argument, each at most $\exp(1)$; for the arguments excluded above both sides vanish, and these numbers may all be taken equal to $1$. For $n=m=0$ the product is empty, and $E_{00}(X;\hat{b}+c_1)=E_{00}(X;\hat{b})$. Together with the comparison of $E(X;\mathscr{A}_L)$ and $E(X;\hat{b}+c_1)$ above, this proves that $\left|E_{nm}(X;\mathscr{A}_L;\xi,\boldsymbol{\zeta})\right|$ is $\left|E_{nm}(X;\hat{b};\xi,\boldsymbol{\zeta})\right|$ multiplied by $n+m$ positive numbers, one for each argument, each at most $\exp(1)$, and that $E_{00}(X;\mathscr{A}_L)=E_{00}(X;\hat{b})$. The two inequalities between the moduli asserted in~(iii) follow, as $\exp(n+m)\le\exp\left(2(n+m)\right)$, and integration over the arguments, with the kernel norm of Section~\ref{sec:grassmann-algebras-norms}, gives the bound for $\|E_{nm}(X;\mathscr{A}_L)\|$.

This completes the proof, and Proposition~\ref{prop:connections-near-small} follows.
\end{proof}

\subsection{Locality}\label{sec:locality}

In this section the setting~\eqref{eq:polymer-activities-Sb} is assumed. Of these hypotheses, the definitions of this section, the descriptions of the property~\eqref{it:renormalisation-group-1-k2} by local maps and assertions~(a) and~(b) of Theorem~\ref{thm:closure-properties-local} require only~\eqref{eq:polymer-activities-Sa}, that is, $k\ge 0$ and $l\in\mathrm{Cl}$, together with $e_k>0$, the charge with which the domains $\tilde{R}_k$ and $\mathfrak{X}_k$ are formed; they concern the single torus $T_k=\mathbb{L}_k/l\mathbb{Z}^3$. Since~\eqref{eq:polymer-activities-Sb} gives $l/L\in\mathrm{Cl}$ and $e_{k+1}=L^{\frac{1}{2}}\cdot e_k>0$, they hold also for the data $(k+1,l/L,e_{k+1})$ in place of $(k,l,e_k)$, that is, on the torus $T_{k+1}=\mathbb{L}_{k+1}/(l/L)\mathbb{Z}^3$, with the domains $\tilde{R}_{k+1}$ and $\mathfrak{X}_{k+1}$ formed with $e_{k+1}$. In this section, a definition or an assertion taken at the level $k+1$ is formed with these data. Besides the generators $\psi(\xi)$ and $\chi(\boldsymbol{\zeta})$, the algebras considered here may contain further families of independent generators, each indexed by a finite set as in Section~\ref{sec:coefficients-kernels-grassmann}, for instance the generators $\Psi(u)$ and $W(u)$ of Section~\ref{sec:grassmann-algebras-norms}, or those of one of these two families; these further generators are called \emph{spectators}. The integral $\int\mathrm{d}\mu_I(W)$ of Section~\ref{sec:norms-gaussian-integrals} is taken over a family $W$ of spectators of the kind considered there. Let $\mathrm{Gr}$ be the Grassmann algebra over $\mathbb{C}$ whose independent generators are the $\psi(\xi)$, the $\chi(\boldsymbol{\zeta})$ and the spectators that are present.

The \emph{legs} of an element of $\mathrm{Gr}$ are the legs, in the sense of Definition~\ref{def:terminology}(5), of the products of generators that carry a non-zero coefficient in its expansion (Section~\ref{sec:coefficients-kernels-grassmann}). A generator $\chi(x,x',\boldsymbol{\beta},\boldsymbol{\omega})$ has the \emph{length} $|x-x'|$. For a union $Z$ of $M$-cubes of $T_k$, the generator $\chi(x,x',\boldsymbol{\beta},\boldsymbol{\omega})$ is \emph{of kind}~(I), resp.\ \emph{of kind}~(II), \emph{for} $Z$ if the pair $(x,x')$ is admissible for $Z$ of kind~(I), resp.\ of kind~(II), in the sense of Definition~\ref{def:pairs-sites-holder} and~\eqref{eq:polymer-activities-12}.

\begin{definition}[Local maps]\label{def:local-maps}
Let $Z$ be a union of $M$-cubes of $T_k$. Let $\mathfrak{A}_k(Z)$ be the subalgebra of $\mathrm{Gr}$ spanned by $1$ and by the products of the following generators: the $\psi(x,\boldsymbol{\beta},\boldsymbol{\omega})$ with $x\in Z$, the $\chi(x,x',\boldsymbol{\beta},\boldsymbol{\omega})$ for which the pair $(x,x')$ is admissible for $Z$ in the sense of~\eqref{eq:polymer-activities-12}, and all spectators. These are called the generators of $\mathfrak{A}_k(Z)$. Let $\mathrm{Dm}$ be a set of pairs $(\mathscr{A},\mathscr{C})$ of connections, that is, of complex bond functions on $T_k$. A map $F\colon\mathrm{Dm}\to\mathrm{Gr}$ is \emph{local in} $Z$ if the following two conditions hold:
\begin{enumerate}
\item[(i)] $F(\mathscr{A},\mathscr{C})\in\mathfrak{A}_k(Z)$ for every $(\mathscr{A},\mathscr{C})\in\mathrm{Dm}$;
\item[(ii)] $F(\mathscr{A},\mathscr{C})=F(\mathscr{A}',\mathscr{C}')$ whenever $(\mathscr{A},\mathscr{C})$ and $(\mathscr{A}',\mathscr{C}')$ lie in $\mathrm{Dm}$ and satisfy $\mathscr{A}(b)=\mathscr{A}'(b)$ and $\mathscr{C}(b)=\mathscr{C}'(b)$ for every bond $b$ of $T_k$ that meets $Z$, that is, that has an end point in $Z$.
\end{enumerate}
\end{definition}

Let $\mathscr{E}=\left\{\mathscr{E}(X;\mathscr{A},\mathscr{C}):X\in D_k,\ (\mathscr{A},\mathscr{C})\in\tilde{R}_k\times\mathfrak{X}_k\right\}$ be a family of elements of $\mathrm{Gr}[\psi,\chi]$. In the terms of Definition~\ref{def:local-maps}, the property~\eqref{it:renormalisation-group-1-k2} of Definition~\ref{def:renormalisation-group-1} states the following: for every $X\in D_k$ the map $(\mathscr{A},\mathscr{C})\mapsto\mathscr{E}(X;\mathscr{A},\mathscr{C})$ is local in $X$, with $\mathrm{Dm}=\tilde{R}_k\times\mathfrak{X}_k$ and without spectators. Indeed, the antisymmetric kernels of an element are determined by the element (Proposition~\ref{prop:weighted-coefficient-norms}\eqref{it:weighted-coefficient-norms-N0}).

\begin{theorem}[Closure properties of local maps]\label{thm:closure-properties-local}
Let $Z$, $Z'$ and $Z_1,\dots,Z_n$ be unions of $M$-cubes of $T_k$.
\begin{enumerate}
\item[(a)] Let $Z\subset Z'$. Then $\mathfrak{A}_k(Z)\subset\mathfrak{A}_k(Z')$, and, for every set $\mathrm{Dm}$ of pairs of connections as in Definition~\ref{def:local-maps}, a map from $\mathrm{Dm}$ to $\mathrm{Gr}$ that is local in $Z$ is also local in $Z'$. A pair that is admissible for $Z$ of kind~(I), resp.\ of kind~(II), is admissible for $Z'$ of kind~(I), resp.\ of kind~(II); and a bond with an end point in $Z$ meets $Z'$.
\item[(b)] Let $\mathrm{Dm}$ be a set of pairs of connections as in Definition~\ref{def:local-maps}. With the pointwise operations, the maps $F\colon\mathrm{Dm}\to\mathrm{Gr}$ that are local in $Z$ form an algebra, and a pointwise limit of such maps is again local in $Z$. Each of the following operations, applied to maps local in $Z$, produces maps local in $Z$:
\begin{itemize}
\item the exponential map $F\mapsto\exp F$, and the summation of a series of maps local in $Z$ which converges in norm at every point of $\mathrm{Dm}$;
\item differentiation and integration in an auxiliary parameter: if $F_{\tau}$ is local in $Z$ for every value of a parameter $\tau$, then every derivative of $F_{\tau}$ in $\tau$ and every integral of $F_{\tau}$ over $\tau$ is local in $Z$. Examples of such parameters are $s$, $t$ and $\mathsf{Z}$ of Section~\ref{sec:definition-step-map}; examples of such operations are $\partial/\partial s_P$ and $\int_{[0,1]^P}\mathrm{d}s_P$ of~\eqref{eq:renormalisation-group-57}, the contour integral $\oint_{\rho}$ of Section~\ref{sec:norms-gaussian-integrals} and the integral $\int\mathrm{d}\hat{\mu}(\mathsf{Z})$ of Section~\ref{sec:averaging-fluctuations-gauge};
\item the integration $\int\mathrm{d}\mu_I(W)$ of Section~\ref{sec:norms-gaussian-integrals} over the spectators $W$.
\end{itemize}
If $F_i\colon\mathrm{Dm}\to\mathrm{Gr}$ is local in $Z_i$ for $i=1,\dots,n$, then the product $F_1\cdots F_n$ is local in $Z_1\cup\dots\cup Z_n$.
\item[(c)] If a family $\mathscr{E}$ as above has the property~\eqref{it:renormalisation-group-1-k2}, then so has the family $R\,\mathscr{E}$, where $R$ is the extraction of Section~\ref{sec:local-parts-adjustment}.
\item[(d)] If a family $E$ as above has the property~\eqref{it:renormalisation-group-1-k2}, then so has the family $E^{\mathrm{nat}}$ of Definition~\ref{def:transport-mean}. Moreover, for every $X\in D_k$ and every $(\mathscr{A},\mathscr{C})\in\tilde{R}_k\times\mathfrak{X}_k$, every leg $\chi$ of $E^{\mathrm{nat}}(X;\mathscr{A},\mathscr{C})$ that is not a leg of $E(X;\mathscr{A},\mathscr{C})$ is of kind~(I) for $X$, and its length is at most $M$.
\item[(e)] Let $Y\in D_{k+1}$, and let $LY$ and the homomorphism $S_L$ be those of Proposition~\ref{prop:properties-actions}(d). Let $X\in D_k$ with $X\subset LY$, and let $F$ be a map from $\tilde{R}_k\times\mathfrak{X}_k$ into $\mathrm{Gr}[\psi,\chi]$ which is local in $X$, the algebra $\mathrm{Gr}$ of Definition~\ref{def:local-maps} being formed with no spectators. Let $\mathscr{A}_L$, $\mathscr{C}_L$ be the scalings of Section~\ref{sec:local-parts-adjustment}. Then the map
\[
(\mathscr{A},\mathscr{C})\mapsto S_L\left[F(\mathscr{A}_L,\mathscr{C}_L)\right]
\]
is local in $Y$ on $\tilde{R}_{k+1}\times\mathfrak{X}_{k+1}$, in the sense of Definition~\ref{def:local-maps} at the level $k+1$. The homomorphism $S_L$ sends $\chi(x,x',\boldsymbol{\beta},\boldsymbol{\omega})$ to $L^{-1-\alpha}\cdot\chi(x/L,x'/L,\boldsymbol{\beta},\boldsymbol{\omega})$. If $\chi(x,x',\boldsymbol{\beta},\boldsymbol{\omega})$ is of kind~(I) for $X$, then $\chi(x/L,x'/L,\boldsymbol{\beta},\boldsymbol{\omega})$ is of kind~(I) for $Y$, and its length is that of $\chi(x,x',\boldsymbol{\beta},\boldsymbol{\omega})$ divided by $L$; if $\chi(x,x',\boldsymbol{\beta},\boldsymbol{\omega})$ is of kind~(II) for $X$, then $\chi(x/L,x'/L,\boldsymbol{\beta},\boldsymbol{\omega})$ is of kind~(II) for $Y$. Consequently, if a family $E'$ of the level $k$ has the property~\eqref{it:renormalisation-group-1-k2}, then the family $L\,E'$ of Theorem~\ref{thm:extraction}(e) has the property~\eqref{it:renormalisation-group-1-k2} at the level $k+1$.
\end{enumerate}
Assertions~(c), (d) and~(e) remain true for the families without the second slot, in the following sense. A family $\mathscr{E}=\left\{\mathscr{E}(X;\mathscr{A}):X\in D_k,\ \mathscr{A}\in\tilde{R}_k\right\}$ of elements of $\mathrm{Gr}[\psi,\chi]$ is regarded as the family of the maps $(\mathscr{A},0)\mapsto\mathscr{E}(X;\mathscr{A})$ on $\mathrm{Dm}=\tilde{R}_k\times\{0\}$. Its property~\eqref{it:renormalisation-group-1-k2} states that each of these maps is local in $X$, without spectators (Proposition~\ref{prop:weighted-coefficient-norms}\eqref{it:weighted-coefficient-norms-N0}); for these maps, condition~(ii) of Definition~\ref{def:local-maps} reduces to $\mathscr{A}(b)=\mathscr{A}'(b)$ for every bond $b$ that meets $X$. With $\tilde{R}_k\times\{0\}$ and $\tilde{R}_{k+1}\times\{0\}$ in place of $\tilde{R}_k\times\mathfrak{X}_k$ and $\tilde{R}_{k+1}\times\mathfrak{X}_{k+1}$, the three assertions read as follows. If $\mathscr{E}$ has the property~\eqref{it:renormalisation-group-1-k2}, then so has $R\,\mathscr{E}$. If $E$ has the property~\eqref{it:renormalisation-group-1-k2}, then so has the family $E^{\mathrm{nat}}$, formed by Definition~\ref{def:transport-mean} with the total connection $\mathscr{T}=\mathscr{A}$; moreover, for every $X\in D_k$ and every $\mathscr{A}\in\tilde{R}_k$, every leg $\chi$ of $E^{\mathrm{nat}}(X;\mathscr{A})$ that is not a leg of $E(X;\mathscr{A})$ is of kind~(I) for $X$, and its length is at most $M$. If $Y$, $LY$, $S_L$ and $X\subset LY$ are as in~(e), and if $F$ is a map from $\tilde{R}_k\times\{0\}$ into $\mathrm{Gr}[\psi,\chi]$ which is local in $X$, then the map $(\mathscr{A},0)\mapsto S_L\left[F(\mathscr{A}_L,0)\right]$ is local in $Y$ on $\tilde{R}_{k+1}\times\{0\}$; consequently, if a family $E'$ of the level $k$ has the property~\eqref{it:renormalisation-group-1-k2}, then $L\,E'$ has the property~\eqref{it:renormalisation-group-1-k2} at the level $k+1$.
\end{theorem}

\begin{proof}
\emph{(a)} By Definition~\ref{def:pairs-sites-holder} and~\eqref{eq:polymer-activities-12}, a pair $(x,x')\in\mathrm{Pz}_k$ is admissible for $Z$ of kind~(I) if $x$, $x'$ and the six paths $\Gamma^{\pi}(x,x')$ lie in $Z$, and of kind~(II) if $x$ and $x'$ are nearest neighbours and at least one of them lies in $Z$. Both conditions persist when $Z$ is replaced by the larger set $Z'$, and so does the condition that a bond has an end point in $Z$. Hence every generator of $\mathfrak{A}_k(Z)$ is a generator of $\mathfrak{A}_k(Z')$, and $\mathfrak{A}_k(Z)\subset\mathfrak{A}_k(Z')$. Two pairs in $\mathrm{Dm}$ that coincide on the bonds that meet $Z'$ coincide on the bonds that meet $Z$; so a map local in $Z$ takes the same value at them, and it is local in $Z'$.

\emph{(b)} The generators of $\mathrm{Gr}$ are indexed by finite sets; so $\mathrm{Gr}$ is a complex vector space of finite dimension, and its linear subspace $\mathfrak{A}_k(Z)$ is closed. Moreover, $\mathfrak{A}_k(Z)$ is a subalgebra, since the product of two products of generators of $\mathfrak{A}_k(Z)$ is, up to sign, again such a product, or zero. Condition~(ii) of Definition~\ref{def:local-maps} is preserved by sums, by products and by pointwise limits. Hence the maps local in $Z$ form an algebra which contains the pointwise limits of its elements. The exponential $\exp F$ is the pointwise limit of the partial sums $\sum_{j=0}^q\frac{1}{j!}\cdot F^j$, $q=0,1,2,\dots$, and the sum of a series is the pointwise limit of its partial sums; these partial sums are local in $Z$. A derivative or an integral with respect to a parameter is a limit of linear combinations of values of $F_{\tau}$ at several values of $\tau$; hence it is local in $Z$ if $F_{\tau}$ is local in $Z$ for every value of $\tau$.

Let $W$ be a family of spectators, let $\boldsymbol{\phi}$ denote the other generators of $\mathrm{Gr}$, and let $\boldsymbol{\phi}(I)$ and $W(J)$ be monomials as in Section~\ref{sec:norms-gaussian-integrals}. By the definition of $\int\mathrm{d}\mu_I(W)$ in Section~\ref{sec:norms-gaussian-integrals} and by~\eqref{eq:polymer-activities-m},
\[
\int\mathrm{d}\mu_I(W)\,\boldsymbol{\phi}(I)\cdot W(J)=m(J)\cdot\boldsymbol{\phi}(I).
\]
Since the generators anticommute, every element of $\mathfrak{A}_k(Z)$ is a linear combination of products $\boldsymbol{\phi}(I)\cdot W(J)$ in which every factor of $\boldsymbol{\phi}(I)$ is a generator of $\mathfrak{A}_k(Z)$. Therefore $\int\mathrm{d}\mu_I(W)$ maps $\mathfrak{A}_k(Z)$ into the algebra $\mathfrak{A}_k(Z)$ formed without the spectators $W$. Being linear, it preserves condition~(ii).

Finally, let $F_i$ be local in $Z_i$ for $i=1,\dots,n$. By~(a), each $F_i$ is local in $Z_1\cup\dots\cup Z_n$, and so is the product $F_1\cdots F_n$, because the maps local in $Z_1\cup\dots\cup Z_n$ form an algebra.

\emph{(c)} Let $\mathscr{E}$ have the property~\eqref{it:renormalisation-group-1-k2}, and let $X\in D_k$. If $X$ is large, then $(R\,\mathscr{E})(X)=\mathscr{E}(X)$. If $X$ is small, then, by the definition of $R$ in Section~\ref{sec:local-parts-adjustment} and in the notation of Section~\ref{sec:symmetries-kernels-low},
\[
(R\,\mathscr{E})(X;\mathscr{A},\mathscr{C})=\mathscr{E}(X;\mathscr{A},\mathscr{C})-\boldsymbol{\alpha}_0(\mathscr{E},X)\cdot\operatorname{Vol}(X)-\int_X\mathrm{d}x\;\bar{\psi}(x)^T\cdot\boldsymbol{\alpha}_2(\mathscr{E},X)\cdot\psi(x).
\]
The number $\boldsymbol{\alpha}_0(\mathscr{E},X)$ and the matrix $\boldsymbol{\alpha}_2(\mathscr{E},X)$ are formed with $\mathscr{E}(X;0,0)$; so the two subtracted terms do not depend on $(\mathscr{A},\mathscr{C})$. They lie in $\mathfrak{A}_k(X)$: the first is a number, and the second is a linear combination of the products $\bar{\psi}_{\boldsymbol{\alpha}}(x)\cdot\psi_{\boldsymbol{\beta}}(x)$ with $x\in X$. Hence the constant maps with these values are local in $X$, and by~(b) the map $(\mathscr{A},\mathscr{C})\mapsto(R\,\mathscr{E})(X;\mathscr{A},\mathscr{C})$ is local in $X$.

\emph{(d)} For large $X$ one has $E^{\mathrm{nat}}(X)=E(X)$ (Definition~\ref{def:transport-mean}). For small $X$ the assertions are those of Theorem~\ref{thm:properties-adjustment}(c). The proof of that statement rests on~\eqref{eq:polymer-activities-ch1}, on~\eqref{eq:polymer-activities-u2} and on the property~\eqref{it:renormalisation-group-1-k2} of $E$, and therefore applies to every family $E$ with the property~\eqref{it:renormalisation-group-1-k2}. The new legs are generators $\chi(p_i,p_{i+1},\boldsymbol{\beta},\boldsymbol{\omega})$ with $p_i\ne p_{i+1}$, the points $p_i$ being those of~\eqref{eq:polymer-activities-17}; the pair $(p_i,p_{i+1})$ lies in $\mathrm{Pz}_k$ by~\eqref{eq:polymer-activities-ch1}, so that $|p_i-p_{i+1}|\le M$.

\emph{(e)} Let $(\mathscr{A},\mathscr{C})\in\tilde{R}_{k+1}\times\mathfrak{X}_{k+1}$. By Lemma~\ref{lem:scaling-field-domains}, $(\mathscr{A}_L,\mathscr{C}_L)$ lies in $\tilde{R}_k\times\mathfrak{X}_k$; so $F(\mathscr{A}_L,\mathscr{C}_L)$ is defined.

The homomorphism $S_L$ maps $\psi(x,\boldsymbol{\beta},\boldsymbol{\omega})$ to $L^{-1}\cdot\psi(x/L,\boldsymbol{\beta},\boldsymbol{\omega})$ and $\chi(x,x',\boldsymbol{\beta},\boldsymbol{\omega})$ to $L^{-1-\alpha}\cdot\chi(x/L,x'/L,\boldsymbol{\beta},\boldsymbol{\omega})$; here $x/L$ and $x'/L$ are sites of $T_{k+1}=\mathbb{L}_{k+1}/(l/L)\mathbb{Z}^3$, whose mesh is $\eta_{k+1}=\eta_k/L$. If $x\in X$, then $x/L\in Y$, because $X\subset LY$. Let $(x,x')$ be admissible for $X$ of kind~(I). Then $x/L$ and $x'/L$ are distinct sites of $T_{k+1}$ in $Y$, one has
\[
\left|\frac{x}{L}-\frac{x'}{L}\right|=\frac{|x-x'|}{L}\le M,
\]
and the six paths $\Gamma^{\pi}(x/L,x'/L)$ are the images of the six paths $\Gamma^{\pi}(x,x')$ under the map $z\mapsto z/L$; so they lie in $\left\{z/L:z\in X\right\}$, hence in $Y$. Therefore $(x/L,x'/L)$ is admissible for $Y$ of kind~(I), and the generator $\chi(x/L,x'/L,\boldsymbol{\beta},\boldsymbol{\omega})$ has the length $|x-x'|/L$. Let $(x,x')$ be admissible for $X$ of kind~(II). The map $z\mapsto z/L$ is a bijection from the sites of $T_k$ onto the sites of $T_{k+1}$ (it maps $\mathbb{L}_k=\eta_k\mathbb{Z}^3$ onto $\mathbb{L}_{k+1}=\eta_{k+1}\mathbb{Z}^3$ and $l\mathbb{Z}^3$ onto $(l/L)\mathbb{Z}^3$), and $\eta_{k+1}=\eta_k/L$; hence it maps the set of pairs of nearest neighbours of $T_k$ onto the set of pairs of nearest neighbours of $T_{k+1}$. In particular, the sites $x/L$ and $x'/L$ of $T_{k+1}$ are nearest neighbours. One of the sites $x$, $x'$ belongs to $X$, and its image belongs to $Y$. Thus $(x/L,x'/L)$ is admissible for $Y$ of kind~(II). This proves the assertions of~(e) on the generators $\chi$. It shows, moreover, that $S_L$ maps every generator of $\mathfrak{A}_k(X)$ to a multiple of a generator of $\mathfrak{A}_{k+1}(Y)$. Since $S_L$ is a homomorphism, it follows that $S_L\left[F(\mathscr{A}_L,\mathscr{C}_L)\right]\in\mathfrak{A}_{k+1}(Y)$.

Let $(\mathscr{A}',\mathscr{C}')\in\tilde{R}_{k+1}\times\mathfrak{X}_{k+1}$ coincide with $(\mathscr{A},\mathscr{C})$ on the bonds of $T_{k+1}$ that meet $Y$. For a bond $b$ of $T_k$ with the end points $x$ and $x'$, let $b/L$ be the bond of $T_{k+1}$ with the end points $x/L$ and $x'/L$. Then
\[
\mathscr{A}_L(b)=L^{-\frac{1}{2}}\cdot\mathscr{A}(b/L),\qquad\mathscr{C}_L(b)=L^{-\frac{1}{2}}\cdot\mathscr{C}(b/L),
\]
and $b/L$ meets $Y$ if $b$ meets $X$. Hence $(\mathscr{A}_L,\mathscr{C}_L)$ and $(\mathscr{A}'_L,\mathscr{C}'_L)$ coincide on the bonds that meet $X$, and $F(\mathscr{A}_L,\mathscr{C}_L)=F(\mathscr{A}'_L,\mathscr{C}'_L)$, because $F$ is local in $X$. Applying $S_L$ to both sides gives condition~(ii) of Definition~\ref{def:local-maps} at the level $k+1$. So the map $(\mathscr{A},\mathscr{C})\mapsto S_L\left[F(\mathscr{A}_L,\mathscr{C}_L)\right]$ is local in $Y$.

Let now $E'$ have the property~\eqref{it:renormalisation-group-1-k2} at the level $k$, and let $Y\in D_{k+1}$. By Theorem~\ref{thm:extraction}(e),
\[
(L\,E')(Y;\mathscr{A},\mathscr{C})=\sum_{X:\ \bar{X}^L=LY}S_L\left[E'(X;\mathscr{A}_L,\mathscr{C}_L)\right].
\]
Every $X$ with $\bar{X}^L=LY$ is contained in $LY$, because each $M$-cube of $X$ lies in an $LM$-cube of $\bar{X}^L$ by~\eqref{eq:polymer-activities-7}; and the map $(\mathscr{A},\mathscr{C})\mapsto E'(X;\mathscr{A},\mathscr{C})$, from $\tilde{R}_k\times\mathfrak{X}_k$ into $\mathrm{Gr}[\psi,\chi]$, is local in $X$ by the property~\eqref{it:renormalisation-group-1-k2} of $E'$. By the first part of~(e) every summand is local in $Y$. The sum is finite, and by~(b), applied at the level $k+1$, it is local in $Y$. Thus $L\,E'$ has the property~\eqref{it:renormalisation-group-1-k2} at the level $k+1$.

\emph{Families without the second slot.} Let $\mathrm{Dm}=\tilde{R}_k\times\{0\}$. Let $\mathscr{E}$ be a family without the second slot with the property~\eqref{it:renormalisation-group-1-k2}, and let $X\in D_k$. If $X$ is large, then $(R\,\mathscr{E})(X)=\mathscr{E}(X)$. If $X$ is small, the number $\boldsymbol{\alpha}_0(\mathscr{E},X)$ and the matrix $\boldsymbol{\alpha}_2(\mathscr{E},X)$ are formed with $\mathscr{E}(X;0)$; so the two terms subtracted by $R$ do not depend on $\mathscr{A}$, and, as in the proof of~(c), they lie in $\mathfrak{A}_k(X)$. The constant maps on $\mathrm{Dm}$ with these values are local in $X$, and~(b), applied with this $\mathrm{Dm}$, shows that $(\mathscr{A},0)\mapsto(R\,\mathscr{E})(X;\mathscr{A})$ is local in $X$. Let $E$ be a family without the second slot with the property~\eqref{it:renormalisation-group-1-k2}. For large $X$ one has $E^{\mathrm{nat}}(X)=E(X)$. For small $X$ the assertions of~(d) are those of Theorem~\ref{thm:properties-adjustment}(c) with $\mathscr{C}=0$, so that $\mathscr{T}=\mathscr{A}$: as noted in the proof of~(d), the proof of that statement rests only on~\eqref{eq:polymer-activities-ch1}, on~\eqref{eq:polymer-activities-u2} and on the property~\eqref{it:renormalisation-group-1-k2} of $E$; the new legs have length at most $M$ by~\eqref{eq:polymer-activities-ch1}, as in the proof of~(d). Let now $Y$, $LY$, $S_L$ and $X\subset LY$ be as in~(e), and let $F$ be a map from $\mathrm{Dm}$ into $\mathrm{Gr}[\psi,\chi]$ which is local in $X$. The scaling of the bond function $0$ is $0$, and $\mathscr{A}_L\in\tilde{R}_k$ for $\mathscr{A}\in\tilde{R}_{k+1}$ by Lemma~\ref{lem:scaling-field-domains}(c); so $F(\mathscr{A}_L,0)$ is defined, and the proof of~(e) applies unchanged with every second entry equal to $0$. Finally, let $E'$ be a family of the level $k$ without the second slot with the property~\eqref{it:renormalisation-group-1-k2}, and let $Y\in D_{k+1}$. By Theorem~\ref{thm:extraction}(e), which remains true without the second slot, $(L\,E')(Y;\mathscr{A})$ is the sum of the elements $S_L\left[E'(X;\mathscr{A}_L)\right]$ over the finitely many $X$ with $\bar{X}^L=LY$. Each of these $X$ lies in $LY$, as shown above, and the map $(\mathscr{A},0)\mapsto E'(X;\mathscr{A})$ is local in $X$ on $\mathrm{Dm}$; hence every summand is local in $Y$ on $\tilde{R}_{k+1}\times\{0\}$, and~(b), applied at the level $k+1$ with $\tilde{R}_{k+1}\times\{0\}$ in place of $\mathrm{Dm}$, shows that $L\,E'$ has the property~\eqref{it:renormalisation-group-1-k2} at the level $k+1$.

With~(a)--(e), Proposition~\ref{prop:locality} follows.
\end{proof}

\subsection{Contraction under reblocking}\label{sec:contraction-reblocking}

Assume the setting~\eqref{eq:polymer-activities-Sb}. The weights $h_k=(h_{k,1},h_{k,2})=\left(e_k^{-\frac{1}{4}},e_k^{-\frac{1}{4}+\varepsilon}\right)$ of Section~\ref{sec:grassmann-algebras-norms} and the corresponding weights $h_{k+1}$ of the level $k+1$ are related as follows, because $e_{k+1}=L^{\frac{1}{2}}\cdot e_k$:
\[
h_{k+1,1}=L^{-\frac{1}{8}}\cdot h_{k,1},\qquad h_{k+1,2}=L^{-\frac{1}{8}+\frac{\varepsilon}{2}}\cdot h_{k,2}.
\]

\begin{lemma}[Norms under rescaling]\label{lem:norms-rescaling}
Let $S_L$ be the homomorphism of Proposition~\ref{prop:properties-actions}(d) from the algebra $\mathrm{Gr}[\psi,\chi]$ of the level $k$ to the algebra $\mathrm{Gr}[\psi,\chi]$ of the level $k+1$. It maps distinct basis monomials of the level $k$ to multiples of distinct basis monomials of the level $k+1$, and for a basis monomial of bidegree $(n,m)$ the factor is $\pm L^{-n-(1+\alpha)\cdot m}$. Consequently, for every element $F$ of the algebra of the level $k$, all weights $g_1,g_2>0$ and all $n,m\ge0$,
\begin{gather*}
\left\|S_LF\right\|_{(g_1,g_2)}=\left\|F\right\|_{\left(L^{-1}\cdot g_1,\,L^{-1-\alpha}\cdot g_2\right)},\\
\left\|S_LF^{[n,m]}\right\|_{h_{k+1}}=L^{-n-(1+\alpha)\cdot m}\cdot L^{-\frac{n+m}{8}+\frac{m\cdot\varepsilon}{2}}\cdot\left\|F^{[n,m]}\right\|_{h_k}.
\end{gather*}
For indices $\xi=(x,\boldsymbol{\beta},\boldsymbol{\omega})$ and $\boldsymbol{\zeta}=(x,x',\boldsymbol{\beta},\boldsymbol{\omega})$ of the level $k+1$ put $L\xi:=(Lx,\boldsymbol{\beta},\boldsymbol{\omega})$ and $L\boldsymbol{\zeta}:=(Lx,Lx',\boldsymbol{\beta},\boldsymbol{\omega})$, and apply $L$ to tuples $\xi=(\xi_1,\dots,\xi_n)$ and $\boldsymbol{\zeta}=(\boldsymbol{\zeta}_1,\dots,\boldsymbol{\zeta}_m)$ entry by entry. Let $F_{nm}$ be the antisymmetric kernels of $F$, formed with the weights $\eta_k^3$, $\eta_k^6$ of the level $k$, and $(S_LF)_{nm}$ those of $S_LF$, formed with the weights $\eta_{k+1}^3$, $\eta_{k+1}^6$ of the level $k+1$ (Section~\ref{sec:coefficients-kernels-grassmann}). Then, with $\boldsymbol{\zeta}_j=(x_j,x'_j,\boldsymbol{\beta}_j,\boldsymbol{\omega}_j)$,
\[
(S_LF)_{nm}(\xi;\boldsymbol{\zeta})=\begin{cases}L^{3n+6m}\cdot L^{-n-(1+\alpha)\cdot m}\cdot F_{nm}(L\xi;L\boldsymbol{\zeta})&\text{if }|x_j-x'_j|\le\frac{M}{L}\text{ for }j=1,\dots,m,\\0&\text{otherwise},\end{cases}
\]
and therefore $\left\|(S_LF)_{nm}\right\|=L^{-n-(1+\alpha)\cdot m}\cdot\left\|F_{nm}\right\|$.
\end{lemma}

\begin{proof}
Since $\eta_{k+1}=\frac{\eta_k}{L}$, division by $L$ maps $\mathbb{L}_k$ onto $\mathbb{L}_{k+1}$ and $l\mathbb{Z}^3$ onto $\frac{l}{L}\mathbb{Z}^3$. It therefore induces a one-to-one correspondence $x\mapsto x/L$ between the sites of $T_k=\mathbb{L}_k/l\mathbb{Z}^3$ and those of $T_{k+1}=\mathbb{L}_{k+1}/(l/L)\mathbb{Z}^3$, under which every sup distance is multiplied by $L^{-1}$. Accordingly, $(x,x')\mapsto(x/L,x'/L)$ sends distinct pairs of $\mathrm{Pz}_k$ to distinct pairs of $\mathrm{Pz}_{k+1}$, and the pairs so obtained are exactly the pairs $(y,y')\in\mathrm{Pz}_{k+1}$ with $|y-y'|\le\frac{M}{L}$. For an index $\xi=(x,\boldsymbol{\beta},\boldsymbol{\omega})$ or $\boldsymbol{\zeta}=(x,x',\boldsymbol{\beta},\boldsymbol{\omega})$ of the level $k$, where $x,x'\in\mathbb{L}_k$, let $\xi/L:=(x/L,\boldsymbol{\beta},\boldsymbol{\omega})$ and $\boldsymbol{\zeta}/L:=(x/L,x'/L,\boldsymbol{\beta},\boldsymbol{\omega})$ be the corresponding indices of the level $k+1$. By the definition of $S_L$,
\[
S_L\,\psi(\xi)=L^{-1}\cdot\psi(\xi/L),\qquad S_L\,\chi(\boldsymbol{\zeta})=L^{-1-\alpha}\cdot\chi(\boldsymbol{\zeta}/L);
\]
these formulas hold at all points of $\mathbb{L}_k$, and they respect the antiperiodicity~\eqref{eq:polymer-activities-AP} of both levels, because $(x+l\cdot e_{\mu})/L=x/L+\frac{l}{L}\cdot e_{\mu}$ for every direction $\mu$. Let $I$ be a set of $n$ indices $\xi$ and $J$ a set of $m$ indices $\boldsymbol{\zeta}$ of independent generators of the level $k$ (Section~\ref{sec:antiperiodicity}), and let $I/L$, $J/L$ be their images. Then $S_L\left(\psi(I)\cdot\chi(J)\right)$ is $L^{-n-(1+\alpha)\cdot m}$ times the product of the generators $\psi(\xi/L)$, $\xi\in I$, and $\chi(\boldsymbol{\zeta}/L)$, $\boldsymbol{\zeta}\in J$, taken in the order of $I$ and of $J$. By~\eqref{eq:polymer-activities-AP}, each of these generators is, up to a sign, the independent generator of the level $k+1$ at the same site of $T_{k+1}$, or at the same pair of $\mathrm{Pz}_{k+1}$, and, the two maps above being injective, different factors give different independent generators; putting the factors in increasing order changes the product by a further sign. Hence the product is $\pm\psi(I/L)\cdot\chi(J/L)$, where $I/L$ and $J/L$ are regarded as sets of indices of independent generators of the level $k+1$. Distinct pairs $(I,J)$ have distinct images $(I/L,J/L)$. Consequently
\[
c_{S_LF}(I/L,J/L)=\pm L^{-\#|I|-(1+\alpha)\cdot\#|J|}\cdot c_F(I,J)
\]
for all $I$ and $J$, and the coefficients of $S_LF$ at all other basis monomials vanish. By~\eqref{eq:polymer-activities-1},
\[
\left\|S_LF\right\|_{(g_1,g_2)}=\sum_{I,J}\left|c_F(I,J)\right|\cdot\left(L^{-1}\cdot g_1\right)^{\#|I|}\cdot\left(L^{-1-\alpha}\cdot g_2\right)^{\#|J|}=\left\|F\right\|_{\left(L^{-1}\cdot g_1,\,L^{-1-\alpha}\cdot g_2\right)}.
\]
All basis monomials that occur in $F^{[n,m]}$ have bidegree $(n,m)$; so~\eqref{eq:polymer-activities-1} gives $\|F^{[n,m]}\|_{(c_1\cdot g_1,\,c_2\cdot g_2)}=c_1^n\cdot c_2^m\cdot\|F^{[n,m]}\|_{(g_1,g_2)}$ for all $c_1,c_2>0$. Apply the first identity with $(g_1,g_2)=h_{k+1}$. By the relations between $h_{k+1}$ and $h_k$ stated at the beginning of Section~\ref{sec:contraction-reblocking},
\[
L^{-1}\cdot h_{k+1,1}=L^{-1-\frac{1}{8}}\cdot h_{k,1},\qquad L^{-1-\alpha}\cdot h_{k+1,2}=L^{-1-\alpha-\frac{1}{8}+\frac{\varepsilon}{2}}\cdot h_{k,2},
\]
and the second identity follows.

\emph{Kernels.} By~\eqref{eq:polymer-activities-2}, with the weights $\eta_k^3$ and $\eta_k^6$,
\[
F^{[n,m]}=\frac{1}{n!\,m!}\sum_{\xi,\boldsymbol{\zeta}}\eta_k^{3n+6m}\cdot F_{nm}(\xi;\boldsymbol{\zeta})\cdot\psi(\xi_1)\cdots\psi(\xi_n)\cdot\chi(\boldsymbol{\zeta}_1)\cdots\chi(\boldsymbol{\zeta}_m),
\]
the sum being taken over all $n$-tuples $\xi$ and all $m$-tuples $\boldsymbol{\zeta}$ of indices of the level $k$. Apply $S_L$ and write $\xi=L\xi'$, $\boldsymbol{\zeta}=L\boldsymbol{\zeta}'$; then $\xi'$ ranges over the $n$-tuples of indices of the level $k+1$, and $\boldsymbol{\zeta}'$ over the $m$-tuples of indices $(y,y',\boldsymbol{\beta},\boldsymbol{\omega})$ of the level $k+1$ with $|y-y'|\le\frac{M}{L}$. Since $\eta_k^{3n+6m}=L^{3n+6m}\cdot\eta_{k+1}^{3n+6m}$, one obtains
\[
S_LF^{[n,m]}=\frac{1}{n!\,m!}\sum_{\xi',\boldsymbol{\zeta}'}\eta_{k+1}^{3n+6m}\cdot K(\xi';\boldsymbol{\zeta}')\cdot\psi(\xi'_1)\cdots\psi(\xi'_n)\cdot\chi(\boldsymbol{\zeta}'_1)\cdots\chi(\boldsymbol{\zeta}'_m),
\]
where now the sum is taken over all tuples of indices of the level $k+1$, and $K(\xi';\boldsymbol{\zeta}')$ denotes the right side of the formula for $(S_LF)_{nm}(\xi';\boldsymbol{\zeta}')$ in the lemma. The function $K$ is antisymmetric in the entries of $\xi'$ and in the entries of $\boldsymbol{\zeta}'$, because $F_{nm}$ is. The homomorphism $S_L$ preserves the bidegree; so $S_LF^{[n,m]}$ is the part of $S_LF$ of bidegree $(n,m)$, and by Proposition~\ref{prop:weighted-coefficient-norms}\eqref{it:weighted-coefficient-norms-N0} its antisymmetric kernel is $K$. Finally,
\[
\left\|(S_LF)_{nm}\right\|=\sum_{\xi',\boldsymbol{\zeta}'}\eta_{k+1}^{3n+6m}\cdot\left|K(\xi';\boldsymbol{\zeta}')\right|=L^{-n-(1+\alpha)\cdot m}\cdot\sum_{\xi,\boldsymbol{\zeta}}\eta_k^{3n+6m}\cdot\left|F_{nm}(\xi;\boldsymbol{\zeta})\right|=L^{-n-(1+\alpha)\cdot m}\cdot\left\|F_{nm}\right\|.
\]
\end{proof}

\begin{theorem}[Contraction]\label{thm:contraction}
Assume the setting~\eqref{eq:polymer-activities-Sb}; in particular $1-\varepsilon<\alpha<1$, by stages~\ref{it:order-choice-2} and~\ref{it:order-choice-8} of Definition~\ref{def:order-choice}. Let $R$, the adjustment $E\mapsto E^{\mathrm{nat}}$ and the reblocking $L$ be those of Section~\ref{sec:local-parts-adjustment} (Definition~\ref{def:transport-mean}, Theorem~\ref{thm:extraction}(e)), and for $\mathscr{E}\in K^{x,C}_k$ put
\[
L\,R\,\mathscr{E}:=L\left((R\,\mathscr{E})^{\mathrm{nat}}\right).
\]
\begin{enumerate}
\item[(a)] The map $\mathscr{E}\mapsto L\,R\,\mathscr{E}$ is linear; it maps $K^{x,C}_k$ into $K^{x,C}_{k+1}$ and $K^x_k$ into $K^x_{k+1}$.
\item[(b)] Every $\mathscr{E}\in K^{x,C}_k$ satisfies
\[
\left\|L\,R\,\mathscr{E}\right\|^x_{k+1}\le c_{\mathrm{fl}}\cdot L^{-\frac{1}{4}+2\varepsilon}\cdot\|\mathscr{E}\|^x_k.
\]
\item[(c)] Every $\mathscr{E}\in K^{x,C}_k$ satisfies $(L\,R\,\mathscr{E})|_0=L\,R\,(\mathscr{E}|_0)$. If the second slot is suppressed, that is, if $K^C_j$, $K_j$ and $\|\cdot\|_j$ take the places of $K^{x,C}_j$, $K^x_j$ and $\|\cdot\|^x_j$ for $j=k,k+1$, every assertion of this theorem still holds.
\item[(d)] Let $\mathscr{E}\in K^{x,C}_k$ and $E:=(R\,\mathscr{E})^{\mathrm{nat}}$, let $X\in D_k$ be small, and let $(\mathscr{A},\mathscr{C})\in\tilde{R}_{k+1}\times\mathfrak{X}_{k+1}$. Then
\[
\left\|S_L\left[E(X;\mathscr{A}_L,\mathscr{C}_L)\right]\right\|_{h_{k+1}}\le3\,L^{-\frac{13}{4}+\frac{3\varepsilon}{2}}\cdot\sup_{(\mathscr{A}',\mathscr{C}')\in\tilde{R}_k\times\mathfrak{X}_k}\left\|E(X;\mathscr{A}',\mathscr{C}')\right\|_{h_k}.
\]
\end{enumerate}
\end{theorem}

\begin{proof}
Let $\mathscr{E}\in K^{x,C}_k$ and $E:=(R\,\mathscr{E})^{\mathrm{nat}}$. By Theorem~\ref{thm:extraction}(d) the family $R\,\mathscr{E}$ lies in $K^{x,C}_k$, and in $K^x_k$ if $\mathscr{E}$ does, and $\|R\,\mathscr{E}\|^x_k\le2\,\|\mathscr{E}\|^x_k$. By Theorem~\ref{thm:properties-adjustment}(e), applied to $R\,\mathscr{E}$, the family $E$ lies in $K^{x,C}_k$, and in $K^x_k$ if $\mathscr{E}$ does, and $\|E\|^x_k\le10\,\|R\,\mathscr{E}\|^x_k\le20\,\|\mathscr{E}\|^x_k$. The map $R$ and the adjustment are linear by the same two theorems; hence $E$ depends linearly on $\mathscr{E}$. For a large polymer $X$ one has $E(X)=(R\,\mathscr{E})(X)=\mathscr{E}(X)$, by the definitions of Section~\ref{sec:local-parts-adjustment}. Together with the definition of the norm $\|\cdot\|^x_k$ in Definition~\ref{def:renormalisation-group-1} this gives, for every $X\in D_k$,
\begin{equation}\label{eq:polymer-activities-21}
\sup_{(\mathscr{A},\mathscr{C})\in\tilde{R}_k\times\mathfrak{X}_k}\left\|E(X;\mathscr{A},\mathscr{C})\right\|_{h_k}\le\begin{cases}20\,\|\mathscr{E}\|^x_k\cdot\exp\left(-\kappa\cdot d_M(X)\right)&\text{if }X\text{ is small},\\\|\mathscr{E}\|^x_k\cdot\exp\left(-\kappa\cdot d_M(X)\right)&\text{if }X\text{ is large}.\end{cases}
\end{equation}
Let $X$ be small. By Theorem~\ref{thm:extraction}(d), $(R\,\mathscr{E})_{00}(X;0,0)=0$ and $\int\mathrm{d}x\,\mathrm{d}y\;(R\,\mathscr{E})_{20}(X;0,0;x,y)=0$; hence Theorem~\ref{thm:properties-adjustment}(b), applied to $R\,\mathscr{E}$, yields
\[
E_{00}(X;0,0)=0,\qquad E_{20}(X;0,0;\cdot,\cdot)=0.
\]

\emph{(a)} The reblocking $L$ is linear, since $S_L$ is linear; together with the linearity of $\mathscr{E}\mapsto E$ this shows that $\mathscr{E}\mapsto L\,R\,\mathscr{E}=L\,E$ is linear. It remains to verify the properties of Definition~\ref{def:renormalisation-group-1} for $L\,E$ at the level $k+1$. Let $Y\in D_{k+1}$. The map $(\mathscr{A},\mathscr{C})\mapsto(\mathscr{A}_L,\mathscr{C}_L)$ is linear, by the definition of the scaling in Section~\ref{sec:local-parts-adjustment}, and by Lemma~\ref{lem:scaling-field-domains} it carries $\tilde{R}_{k+1}\times\mathfrak{X}_{k+1}$ into $\tilde{R}_k\times\mathfrak{X}_k$; by Lemma~\ref{lem:norms-rescaling} each value of a kernel of $S_L\left[E(X;\mathscr{A}_L,\mathscr{C}_L)\right]$ either vanishes or is a constant multiple of a value of a kernel of $E(X;\mathscr{A}_L,\mathscr{C}_L)$; and only finitely many polymers $X$ satisfy $\bar{X}^L=LY$. Hence the kernels of $(L\,E)(Y;\mathscr{A},\mathscr{C})$ are analytic on $\tilde{R}_{k+1}\times\mathfrak{X}_{k+1}$; this is the property~\eqref{it:renormalisation-group-1-k1}. By Theorem~\ref{thm:closure-properties-local}(c), (d) and~(e), in turn, the families $R\,\mathscr{E}$, $E$ and $L\,E$ have the property~\eqref{it:renormalisation-group-1-k2}, the last one at the level $k+1$. The group $\mathbb{G}_k^L$ is a subgroup of $\mathbb{G}_k$; so $E$ has the property~\eqref{it:renormalisation-group-1-k3} for the elements of $\mathbb{G}_k^L$, and Theorem~\ref{thm:extraction}(e), applied to $E'=E$, gives the properties~\eqref{it:renormalisation-group-1-k3}, \eqref{it:renormalisation-group-1-k4} and~\eqref{it:renormalisation-group-1-k5} of $L\,E$ at the level $k+1$, and the property~\eqref{it:renormalisation-group-1-k7} if $\mathscr{E}\in K^x_k$. The homomorphism $S_L$ sends every generator to a positive multiple of a generator of the same kind, barred or unbarred; so the elements $(L\,E)(Y;\mathscr{A},\mathscr{C})$ are even and have equally many barred and unbarred legs. Finally, the property~\eqref{it:renormalisation-group-1-k6} at the level $k+1$, the finiteness of $\|L\,E\|^x_{k+1}$, follows from~(b), which is proved below.

\emph{(d)} Let $X$ be small, let $x_0$ be a site of $X$ and $(\mathscr{A},\mathscr{C})\in\tilde{R}_{k+1}\times\mathfrak{X}_{k+1}$, and let $\hat{b}$ and $\mathscr{C}^b$ be the complex bond functions on $T_k$ that Theorem~\ref{thm:polymer-activities-3} provides for these data. Put $\rho_t:=L^{\frac{7}{4}}$. In this part $\sup\|E(X;\cdot)\|_{h_k}$ and $\sup\|E_{20}(X;\cdot)\|$ denote the suprema of $\|E(X;\mathscr{A}',\mathscr{C}')\|_{h_k}$ and of $\|E_{20}(X;\mathscr{A}',\mathscr{C}')\|$ over all $(\mathscr{A}',\mathscr{C}')\in\tilde{R}_k\times\mathfrak{X}_k$, the second norm being the kernel norm of Section~\ref{sec:grassmann-algebras-norms}, as in~\eqref{eq:polymer-activities-3}.

By Theorem~\ref{thm:polymer-activities-3}(ii) the pair $(t\cdot\hat{b},t\cdot\mathscr{C}^b)$ lies in $R_k\times\mathfrak{X}_k$ for every $t\in\mathbb{C}$ with $|t|\le\rho_t$. The set $R_k\times\mathfrak{X}_k$ is open, and it is contained in $\tilde{R}_k\times\mathfrak{X}_k$ (Definition~\ref{def:domains}). Hence there is an open neighbourhood $U$ of the closed disc $\{t\in\mathbb{C}:|t|\le\rho_t\}$ such that $(t\cdot\hat{b},t\cdot\mathscr{C}^b)\in R_k\times\mathfrak{X}_k\subset\tilde{R}_k\times\mathfrak{X}_k$ for all $t\in U$. The element $E(X;\mathscr{A}_L,\mathscr{C}_L)$ is even; so only bidegrees $(n,m)$ with even $n+m$ occur in it, and the parts $S_LE^{[n,m]}(X;\mathscr{A}_L,\mathscr{C}_L)$ are estimated one bidegree at a time.

\emph{Bidegree $(0,0)$.} The function $f(t):=E_{00}(X;t\cdot\hat{b},t\cdot\mathscr{C}^b)$ is analytic on $U$ by~\eqref{it:renormalisation-group-1-k1}, and $f(0)=E_{00}(X;0,0)=0$. The map $\mathbf{C}$ preserves the bidegree by~\eqref{eq:polymer-activities-4} and, being linear, fixes the multiples of the unit. Therefore the property~\eqref{it:renormalisation-group-1-k5}, in the form $\mathbf{C}\left[E(X;-\mathscr{A}',-\mathscr{C}')\right]=E(X;\mathscr{A}',\mathscr{C}')$ of Section~\ref{sec:lattice-gauge-charge}, gives $f(-t)=f(t)$ whenever $t$ and $-t$ lie in $U$; hence $f'(0)=0$. By Theorem~\ref{thm:polymer-activities-3}(iii), $E_{00}(X;\mathscr{A}_L,\mathscr{C}_L)=E_{00}(X;\hat{b},\mathscr{C}^b)=f(1)$, and by Proposition~\ref{prop:weighted-coefficient-norms}\eqref{it:weighted-coefficient-norms-N5}, $\oint^{(2)}_{\rho_t}f=f(1)-f(0)-f'(0)=f(1)$. Moreover $|f(t)|\le\|E(X;t\cdot\hat{b},t\cdot\mathscr{C}^b)\|_{h_k}$ by~\eqref{eq:polymer-activities-1}, and $\rho_t-1\ge\frac{\rho_t}{2}$. It follows that
\[
\left|E_{00}(X;\mathscr{A}_L,\mathscr{C}_L)\right|=\left|\oint^{(2)}_{\rho_t}f\right|\le\frac{\sup_{|t|=\rho_t}|f(t)|}{\rho_t\cdot(\rho_t-1)}\le2\,L^{-\frac{7}{2}}\cdot\sup\left\|E(X;\cdot)\right\|_{h_k}.
\]
The homomorphism $S_L$ fixes the numbers; so the left side is $\left\|S_LE^{[0,0]}(X;\mathscr{A}_L,\mathscr{C}_L)\right\|_{h_{k+1}}$.

\emph{Bidegree $(2,0)$.} Let $f(t):=E_{20}(X;t\cdot\hat{b},t\cdot\mathscr{C}^b;\cdot,\cdot)$ for $t\in U$; its values lie in the space of the kernels of bidegree $(2,0)$, which is finite-dimensional and is normed by the kernel norm of Section~\ref{sec:grassmann-algebras-norms}, an $L^1$ norm. The function $f$ is analytic by~\eqref{it:renormalisation-group-1-k1}, and $f(0)=E_{20}(X;0,0;\cdot,\cdot)=0$. By Proposition~\ref{prop:weighted-coefficient-norms}\eqref{it:weighted-coefficient-norms-N5}, $\oint_{\rho_t}f=f(1)-f(0)=f(1)$, and
\[
\|f(1)\|\le\frac{\sup_{|t|=\rho_t}\|f(t)\|}{\rho_t-1}\le2\,L^{-\frac{7}{4}}\cdot\sup\left\|E_{20}(X;\cdot)\right\|.
\]
By Theorem~\ref{thm:polymer-activities-3}(iii), $\|E_{20}(X;\mathscr{A}_L,\mathscr{C}_L)\|\le\exp(4)\cdot\|f(1)\|$. Lemma~\ref{lem:norms-rescaling} for $(n,m)=(2,0)$ and~\eqref{eq:polymer-activities-3} give
\[
\left\|S_LE^{[2,0]}(X;\mathscr{A}_L,\mathscr{C}_L)\right\|_{h_{k+1}}=L^{-\frac{9}{4}}\cdot\frac{h_{k,1}^2}{2}\cdot\left\|E_{20}(X;\mathscr{A}_L,\mathscr{C}_L)\right\|\le2\exp(4)\cdot L^{-4}\cdot\sup\left\|E(X;\cdot)\right\|_{h_k};
\]
in the last step $\frac{h_{k,1}^2}{2}\cdot\|E_{20}(X;\mathscr{A}',\mathscr{C}')\|=\|E^{[2,0]}(X;\mathscr{A}',\mathscr{C}')\|_{h_k}\le\|E(X;\mathscr{A}',\mathscr{C}')\|_{h_k}$ for all $(\mathscr{A}',\mathscr{C}')\in\tilde{R}_k\times\mathfrak{X}_k$, by~\eqref{eq:polymer-activities-3} and the additivity of the norm over the bidegrees (Section~\ref{sec:coefficients-kernels-grassmann}).

\emph{The other bidegrees.} By Lemma~\ref{lem:norms-rescaling}, $\|S_LE^{[n,m]}(X;\mathscr{A}_L,\mathscr{C}_L)\|_{h_{k+1}}$ is the product of the factor $L^{-n-(1+\alpha)\cdot m}\cdot L^{-\frac{n+m}{8}+\frac{m\cdot\varepsilon}{2}}$ with $\|E^{[n,m]}(X;\mathscr{A}_L,\mathscr{C}_L)\|_{h_k}$. Since $\alpha>1-\varepsilon$ (stage~\ref{it:order-choice-8}), this factor satisfies
\begin{align*}
L^{-2-\alpha-\frac{1}{4}+\frac{\varepsilon}{2}}&\le L^{-\frac{13}{4}+\frac{3\varepsilon}{2}}&&\text{for }(n,m)=(1,1),\\
L^{-2-2\alpha-\frac{1}{4}+\varepsilon}&\le L^{-\frac{13}{4}+\frac{3\varepsilon}{2}}&&\text{for }(n,m)=(0,2),\\
L^{-n-(1+\alpha)\cdot m}\cdot L^{-\frac{n+m}{8}+\frac{m\cdot\varepsilon}{2}}&\le L^{-4}&&\text{for }n+m\ge4;
\end{align*}
in the last case the exponent is at most $-\frac{9}{8}\cdot(n+m)$, because $\alpha>\frac{\varepsilon}{2}$. Since $L^{-4}\le L^{-\frac{13}{4}+\frac{3\varepsilon}{2}}$, the norm is additive over the bidegrees and $(\mathscr{A}_L,\mathscr{C}_L)\in\tilde{R}_k\times\mathfrak{X}_k$ (Lemma~\ref{lem:scaling-field-domains}), it follows that
\begin{align*}
\sum_{(n,m)\notin\{(0,0),(2,0)\}}\left\|S_LE^{[n,m]}(X;\mathscr{A}_L,\mathscr{C}_L)\right\|_{h_{k+1}}&\le L^{-\frac{13}{4}+\frac{3\varepsilon}{2}}\cdot\left\|E(X;\mathscr{A}_L,\mathscr{C}_L)\right\|_{h_k}\\
&\le L^{-\frac{13}{4}+\frac{3\varepsilon}{2}}\cdot\sup\left\|E(X;\cdot)\right\|_{h_k}.
\end{align*}
Finally, $2\,L^{-\frac{7}{2}}\le L^{-\frac{13}{4}}$ and $2\exp(4)\cdot L^{-4}\le L^{-\frac{13}{4}}$, because $L^{\frac{1}{4}}\ge2$ and $L^{\frac{3}{4}}\ge110>2\exp(4)$ (stage~\ref{it:order-choice-4}). Adding the three contributions and using $L^{-\frac{13}{4}}\le L^{-\frac{13}{4}+\frac{3\varepsilon}{2}}$ one obtains~(d).

\emph{(b)} Let $Y\in D_{k+1}$ and $(\mathscr{A},\mathscr{C})\in\tilde{R}_{k+1}\times\mathfrak{X}_{k+1}$. By the definition of $L$,
\[
(L\,R\,\mathscr{E})(Y;\mathscr{A},\mathscr{C})=\sum_{X:\ \bar{X}^L=LY}S_L\left[E(X;\mathscr{A}_L,\mathscr{C}_L)\right],
\]
and each polymer $X$ in this sum is small or large. Recall that $\kappa\ge\kappa_0+1$ by stage~\ref{it:order-choice-3}, so that Proposition~\ref{prop:polymers}\eqref{it:polymers-t8} applies with $\kappa'=\kappa$ and with $\kappa'=\kappa-1$, and that
\[
L\cdot\left(\kappa-\kappa_0-1\right)\ge L\cdot\left(\kappa-\kappa_0-2\right)\ge\kappa,
\]
a consequence of the third condition of stage~\ref{it:order-choice-4}.

\emph{Small polymers.} For small $X$, part~(d) and~\eqref{eq:polymer-activities-21} give
\[
\left\|S_L\left[E(X;\mathscr{A}_L,\mathscr{C}_L)\right]\right\|_{h_{k+1}}\le60\,L^{-\frac{13}{4}+\frac{3\varepsilon}{2}}\cdot\|\mathscr{E}\|^x_k\cdot\exp\left(-\kappa\cdot d_M(X)\right).
\]
By Proposition~\ref{prop:polymers}\eqref{it:polymers-t8} with $\kappa'=\kappa$, and by $L\cdot(\kappa-\kappa_0-1)\ge\kappa$,
\[
\sum_{X\ \text{small}:\ \bar{X}^L=LY}\left\|S_L\left[E(X;\mathscr{A}_L,\mathscr{C}_L)\right]\right\|_{h_{k+1}}\le60\,L^{-\frac{13}{4}+\frac{3\varepsilon}{2}}\cdot\|\mathscr{E}\|^x_k\cdot c_{\mathrm{rb}}\cdot L^3\cdot\exp\left(-\kappa\cdot d_M(Y)\right).
\]

\emph{Large polymers.} If $X$ is large, then $d_M(X)>L$, and hence
\[
\exp\left(-\kappa\cdot d_M(X)\right)\le\exp(-L)\cdot\exp\left(-(\kappa-1)\cdot d_M(X)\right).
\]
Since $L^{-1}\cdot h_{k+1,1}\le h_{k,1}$ and $L^{-1-\alpha}\cdot h_{k+1,2}\le h_{k,2}$, Lemma~\ref{lem:norms-rescaling} and Proposition~\ref{prop:weighted-coefficient-norms}\eqref{it:weighted-coefficient-norms-N6} give $\|S_LF\|_{h_{k+1}}\le\|F\|_{h_k}$ for every element $F$ of the algebra of the level $k$. With~\eqref{eq:polymer-activities-21}, with Proposition~\ref{prop:polymers}\eqref{it:polymers-t8} for $\kappa'=\kappa-1$, and with $L\cdot(\kappa-\kappa_0-2)\ge\kappa$, the sum of $\|S_L[E(X;\mathscr{A}_L,\mathscr{C}_L)]\|_{h_{k+1}}$ over the large $X$ with $\bar{X}^L=LY$ is therefore at most
\[
\exp(-L)\cdot c_{\mathrm{rb}}\cdot L^3\cdot\|\mathscr{E}\|^x_k\cdot\exp\left(-\kappa\cdot d_M(Y)\right),
\]
and $L^3\cdot\exp(-L)\le L^{-\frac{1}{4}}$ (stage~\ref{it:order-choice-4}).

Adding the two bounds, and using $L^3\cdot L^{-\frac{13}{4}+\frac{3\varepsilon}{2}}=L^{-\frac{1}{4}+\frac{3\varepsilon}{2}}$ and $L^{-\frac{1}{4}}\le L^{-\frac{1}{4}+\frac{3\varepsilon}{2}}$, one obtains
\[
\left\|(L\,R\,\mathscr{E})(Y;\mathscr{A},\mathscr{C})\right\|_{h_{k+1}}\le61\,c_{\mathrm{rb}}\cdot L^{-\frac{1}{4}+\frac{3\varepsilon}{2}}\cdot\|\mathscr{E}\|^x_k\cdot\exp\left(-\kappa\cdot d_M(Y)\right),
\]
where $61\,c_{\mathrm{rb}}=c_{\mathrm{fl}}$ by~\eqref{eq:polymer-activities-5}. Since $L^{-\frac{1}{4}+\frac{3\varepsilon}{2}}\le L^{-\frac{1}{4}+2\varepsilon}$, the definition of $\|\cdot\|^x_{k+1}$ in Definition~\ref{def:renormalisation-group-1} yields~(b).

\emph{(c)} By Theorem~\ref{thm:extraction}(d), $(R\,\mathscr{E})|_0=R\,(\mathscr{E}|_0)$, and by Theorem~\ref{thm:properties-adjustment}(e), $(E')^{\mathrm{nat}}|_0=(E'|_0)^{\mathrm{nat}}$ for every $E'\in K^{x,C}_k$; furthermore, the scaled connection of $\mathscr{C}=0$ is $\mathscr{C}_L=0$. Hence, for $Y\in D_{k+1}$ and $\mathscr{A}\in\tilde{R}_{k+1}$,
\[
(L\,R\,\mathscr{E})|_0(Y;\mathscr{A})=(L\,R\,\mathscr{E})(Y;\mathscr{A},0)=\sum_{X:\ \bar{X}^L=LY}S_L\left[\left(R\,(\mathscr{E}|_0)\right)^{\mathrm{nat}}(X;\mathscr{A}_L)\right]=\left(L\,R\,(\mathscr{E}|_0)\right)(Y;\mathscr{A}).
\]
\emph{Families without the second slot.} Let $\mathscr{E}\in K^C_k$, and let $E:=(R\,\mathscr{E})^{\mathrm{nat}}$, the adjustment being formed by Definition~\ref{def:transport-mean} with the total connection $\mathscr{A}$, as at the end of Theorem~\ref{thm:properties-adjustment}. By Theorem~\ref{thm:extraction}(e), in the form for families without the second slot stated at the end of Theorem~\ref{thm:extraction}, and by Lemma~\ref{lem:scaling-field-domains}(c), according to which $\mathscr{A}_L\in\tilde{R}_k$ for every $\mathscr{A}\in\tilde{R}_{k+1}$,
\[
(L\,R\,\mathscr{E})(Y;\mathscr{A})=\sum_{X:\ \bar{X}^L=LY}S_L\left[E(X;\mathscr{A}_L)\right]\qquad\left(Y\in D_{k+1},\ \mathscr{A}\in\tilde{R}_{k+1}\right).
\]
The part of the proof that precedes~(a), and the proofs of~(a), (d) and~(b), apply to $\mathscr{E}$ with $\mathscr{A}$, $0$, $\tilde{R}_j$, $K^C_j$, $K_j$ and $\|\cdot\|_j$ in place of $(\mathscr{A},\mathscr{C})$, $(0,0)$, $\tilde{R}_j\times\mathfrak{X}_j$, $K^{x,C}_j$, $K^x_j$ and $\|\cdot\|^x_j$ for $j=k,k+1$, and with the following statements on families without the second slot in place of the statements on families with two slots cited there. By Theorem~\ref{thm:extraction}(d), in the form stated at the end of that theorem, the extraction $R$ is linear, the family $R\,\mathscr{E}$ lies in $K^C_k$, and in $K_k$ if $\mathscr{E}$ does, $\|R\,\mathscr{E}\|_k\le2\,\|\mathscr{E}\|_k$, and $(R\,\mathscr{E})_{00}(X;0)=0$ and $\int\mathrm{d}x\,\mathrm{d}y\;(R\,\mathscr{E})_{20}(X;0;x,y)=0$ for every small $X$. By Theorem~\ref{thm:properties-adjustment}(b) and~(e), in the form stated at the end of that theorem, applied to $R\,\mathscr{E}$, the family $E$ lies in $K^C_k$, and in $K_k$ if $\mathscr{E}$ does, $\|E\|_k\le10\,\|R\,\mathscr{E}\|_k\le20\,\|\mathscr{E}\|_k$, the map $\mathscr{E}\mapsto E$ is linear, and, for every small $X$,
\[
E_{00}(X;0)=0,\qquad E_{20}(X;0;\cdot,\cdot)=0.
\]
As $E(X)=(R\,\mathscr{E})(X)=\mathscr{E}(X)$ for every large $X$ (Section~\ref{sec:local-parts-adjustment}), the definition of $\|\cdot\|_k$ in Definition~\ref{def:renormalisation-group-1} gives~\eqref{eq:polymer-activities-21} with $\sup_{\mathscr{A}\in\tilde{R}_k}\|E(X;\mathscr{A})\|_{h_k}$ on the left and with $\|\mathscr{E}\|_k$ in place of $\|\mathscr{E}\|^x_k$. By Theorem~\ref{thm:closure-properties-local}(c), (d) and~(e), in the form stated at the end of that theorem, the families $R\,\mathscr{E}$ and $E$ have the property~\eqref{it:renormalisation-group-1-k2}, and $L\,E$ has it at the level $k+1$. By Theorem~\ref{thm:extraction}(e), in the form stated at the end of that theorem, applied to $E'=E$, the family $L\,E$ has the properties~\eqref{it:renormalisation-group-1-k3}, \eqref{it:renormalisation-group-1-k4} and~\eqref{it:renormalisation-group-1-k5} at the level $k+1$, and the property~\eqref{it:renormalisation-group-1-k7} if $\mathscr{E}\in K_k$.

In the proof of~(d), Theorem~\ref{thm:polymer-activities-3} is applied to the pair $(\mathscr{A},0)$ with $\mathscr{A}\in\tilde{R}_{k+1}$; then $\mathscr{C}^b=0$ and $c_2=0$, and Theorem~\ref{thm:polymer-activities-3}(ii) gives $t\cdot\hat{b}\in R_k$ for every $t\in\mathbb{C}$ with $|t|\le\rho_t=L^{\frac{7}{4}}$. The set $R_k$ is open, and it is contained in $\tilde{R}_k$ (Definition~\ref{def:domains}). Hence the set $U$ of the numbers $t\in\mathbb{C}$ with $t\cdot\hat{b}\in R_k$ is an open neighbourhood of the closed disc $\{t\in\mathbb{C}:|t|\le\rho_t\}$, and $t\cdot\hat{b}\in\tilde{R}_k$ for every $t\in U$. In bidegree $(0,0)$ the function $f(t):=E_{00}(X;t\cdot\hat{b})$ is analytic on $U$ by~\eqref{it:renormalisation-group-1-k1}, one has $f(0)=E_{00}(X;0)=0$, and $f(-t)=f(t)$ whenever $t$ and $-t$ lie in $U$, by the property~\eqref{it:renormalisation-group-1-k5} with $\mathscr{C}=0$, in the form $\mathbf{C}\left[E(X;-\mathscr{A}')\right]=E(X;\mathscr{A}')$, $\mathscr{A}'\in\tilde{R}_k$, of Section~\ref{sec:lattice-gauge-charge}. In bidegree $(2,0)$ the function $f(t):=E_{20}(X;t\cdot\hat{b};\cdot,\cdot)$ is analytic on $U$ by~\eqref{it:renormalisation-group-1-k1}, and $f(0)=E_{20}(X;0;\cdot,\cdot)=0$. The part of Theorem~\ref{thm:polymer-activities-3}(iii) on families without the second slot, applied to $E\in K^C_k$, gives $E_{00}(X;\mathscr{A}_L)=E_{00}(X;\hat{b})$ and, by integration over the arguments of its bound for the kernels, in which the factor is $\exp(n+m)=\exp(2)$ for $(n,m)=(2,0)$,
\[
\left\|E_{20}(X;\mathscr{A}_L)\right\|\le\exp(2)\cdot\left\|E_{20}(X;\hat{b})\right\|\le\exp(4)\cdot\left\|E_{20}(X;\hat{b})\right\|.
\]
The suprema in the proof of~(d) are taken over $\mathscr{A}'\in\tilde{R}_k$, and the bounds for the bidegrees other than $(0,0)$ and $(2,0)$ are those of Lemma~\ref{lem:norms-rescaling}, taken at $\mathscr{A}_L\in\tilde{R}_k$. Lemma~\ref{lem:norms-rescaling}, Proposition~\ref{prop:weighted-coefficient-norms}\eqref{it:weighted-coefficient-norms-N0}, \eqref{it:weighted-coefficient-norms-N5}, \eqref{it:weighted-coefficient-norms-N6} and Proposition~\ref{prop:polymers}\eqref{it:polymers-t8} do not involve the connections and are applied as above. The constants are therefore those of the proof above: $2$, $10$ and $20$ before~\eqref{eq:polymer-activities-21}, $2\,L^{-\frac{7}{2}}$, $2\exp(4)\cdot L^{-4}$ and the factor $3$ in~(d), and $60$ and $61\,c_{\mathrm{rb}}=c_{\mathrm{fl}}$ in~(b).

Proposition~\ref{prop:contraction} follows.
\end{proof}

\paragraph{Remarks.}
\begin{enumerate}
\item The proof of Theorem~\ref{thm:contraction}(b) involves two inequalities between $\kappa$ and $L$, namely
\[
L\cdot\left(\kappa-\kappa_0-1\right)\ge\kappa\qquad\text{and}\qquad L\cdot\left(\kappa-\kappa_0-2\right)\ge\kappa.
\]
Both are consequences of the third condition of stage~\ref{it:order-choice-4}, which reads
\[
L\cdot\left(\kappa-6\kappa_0-4\right)-1\ge\kappa.
\]
Indeed, $L>0$ and $\kappa_0>0$, so that $\kappa-\kappa_0-1>\kappa-\kappa_0-2>\kappa-6\kappa_0-4$; hence
\[
L\cdot\left(\kappa-\kappa_0-1\right)\ge L\cdot\left(\kappa-\kappa_0-2\right)\ge L\cdot\left(\kappa-6\kappa_0-4\right)\ge\kappa+1>\kappa.
\]
\item In every step of the proof of Theorem~\ref{thm:contraction} the supremum over the connections is taken outside the sum over the bidegrees $(n,m)$. In the proof of part~(d), let $E=(R\,\mathscr{E})^{\mathrm{nat}}$, let $X$ be a small polymer, $x_0$ a site of $X$ and $(\mathscr{A},\mathscr{C})\in\tilde{R}_{k+1}\times\mathfrak{X}_{k+1}$ a fixed pair, let $\hat{b}$ and $\mathscr{C}^b$ be the bond functions of Theorem~\ref{thm:polymer-activities-3} for these data, and let $\rho_t=L^{\frac{7}{4}}$. The parts of bidegree $(0,0)$ and $(2,0)$ of $S_L\left[E(X;\mathscr{A}_L,\mathscr{C}_L)\right]$ are bounded by constant multiples of the suprema of
\[
\left|E_{00}(X;\cdot\,,\cdot)\right|\qquad\text{and}\qquad\frac{h_{k,1}^2}{2}\cdot\left\|E_{20}(X;\cdot\,,\cdot)\right\|
\]
over the disc $\left\{\left(t\cdot\hat{b},t\cdot\mathscr{C}^b\right):t\in\mathbb{C},\ |t|\le\rho_t\right\}$. This disc lies in $R_k\times\mathfrak{X}_k$ by Theorem~\ref{thm:polymer-activities-3}(ii), hence in $\tilde{R}_k\times\mathfrak{X}_k$. At each of its points the two quantities are the norms $\|\cdot\|_{h_k}$ of the parts of $E(X;\cdot\,,\cdot)$ of bidegree $(0,0)$ and $(2,0)$, by the definition of the norm and by~\eqref{eq:polymer-activities-3}; and the norm $\|\cdot\|_{h_k}$ of an element is the sum of the norms of its parts of all bidegrees (Section~\ref{sec:coefficients-kernels-grassmann}). Therefore each of the two suprema is at most
\[
\sup_{\tilde{R}_k\times\mathfrak{X}_k}\left\|E(X;\cdot\,,\cdot)\right\|_{h_k}.
\]
The parts of the other bidegrees are bounded through $\left\|E(X;\mathscr{A}_L,\mathscr{C}_L)\right\|_{h_k}$, that is, at the single pair $(\mathscr{A}_L,\mathscr{C}_L)$.
\item Let $w^{\mathrm{E}}_k=e_k^{\frac{1}{4}-7\varepsilon}$ be the window of Section~\ref{sec:running-quantities-tori}. Since $e_{k+1}=L^{\frac{1}{2}}\cdot e_k$, one has
\[
w^{\mathrm{E}}_{k+1}=L^{\frac{1}{8}-\frac{7\varepsilon}{2}}\cdot w^{\mathrm{E}}_k,
\]
so that, in units of the windows, Theorem~\ref{thm:contraction}(b) gives
\[
\frac{\left\|L\,R\,\mathscr{E}\right\|^x_{k+1}}{w^{\mathrm{E}}_{k+1}}\le c_{\mathrm{fl}}\cdot L^{-\frac{3}{8}+\frac{11\varepsilon}{2}}\cdot\frac{\left\|\mathscr{E}\right\|^x_k}{w^{\mathrm{E}}_k}.
\]
\end{enumerate}

\begin{remark}[Extraction, adjustment and scaling on a lift]\label{rem:extraction-adjustment-scaling}
Let $\mathscr{E}\in K^{x,C}_k$, let $(\mathscr{A},\mathscr{C})\in\tilde{R}_k\times\mathfrak{X}_k$ with the total connection $\mathscr{T}=\mathscr{A}+\mathscr{C}$, and let $X\in D_k$ be a polymer with
\[
2M\cdot\left(d_M(X)+2\right)+2<l.
\]
Let $\tilde{X}$ be the lift of $X$ given by Proposition~\ref{prop:polymers}\eqref{it:polymers-t10}. By Definition~\ref{def:terminology}(9), which applies since $2M\cdot\left(d_M(X)+2\right)<l$, the set of the sites of $\tilde{X}$ is a lift of $X$ in the sense of that clause. Put $\rho:=1$. By Proposition~\ref{prop:polymers}\eqref{it:polymers-t7} and~\eqref{it:polymers-t10}, any two sites $z$, $z'$ of $\tilde{X}$ satisfy $|z-z'|\le M\cdot\left(d_M(\tilde{X})+2\right)=M\cdot\left(d_M(X)+2\right)$; as $d_M(X)\ge0$, one has $M\le M\cdot\left(d_M(X)+2\right)$, and therefore
\[
|z-z'|+2\rho+M\le M\cdot\left(d_M(X)+2\right)+2+M\le2M\cdot\left(d_M(X)+2\right)+2<l.
\]
This is the margin which the last paragraph of Definition~\ref{def:terminology}(2) requires of the lift $\tilde{X}$ for the range $\rho$. Consider the following objects:
\begin{enumerate}
\item[(i)] the restrictions of the kernels of $\mathscr{E}(X;\cdot\,,\cdot)$ to arguments within distance $1$ of $\tilde{X}$;
\item[(ii)] the $M$-cubes of which $\tilde{X}$ is the union, together with the centre of each of them and the faces common to two of them;
\item[(iii)] the values taken by $\mathscr{T}$ on the bonds of $\tilde{X}$.
\end{enumerate}
As in Definition~\ref{def:terminology}(2), let $X^{\rho}$ be the set of the sites of $T_k$ at sup distance at most $\rho$ from $X$, let $\tilde{X}^{\rho}$ be the set of the sites of $\mathbb{L}_k$ at sup distance at most $\rho$ from $\tilde{X}$, and let the datum $u$ consist of the kernels of $\mathscr{E}(X;\cdot\,,\cdot)$ at the arguments contained in $X^{\rho}$ and of the values of $\mathscr{T}$ on the bonds both of whose end points lie in $X$. The covering map $\mathbb{L}_k\to T_k$ is injective on $\tilde{X}^{\rho}$ and maps $\tilde{X}$ onto $X$ (Definition~\ref{def:terminology}(2)); hence the transported datum $\tilde{u}$, that is, the restriction of $u$ to the indices contained in $\tilde{X}^{\rho}$, consists of the objects~(i) and~(iii). The objects~(ii) are determined by $\tilde{X}$, which is the first argument of the map $\Phi$ of that definition. It follows from the definitions in Section~\ref{sec:local-parts-adjustment}, in~\eqref{eq:polymer-activities-17}, in Definition~\ref{def:transport-mean} and in Lemma~\ref{lem:norms-rescaling} that each of the following objects is given by one formula on $\mathbb{L}_k$, in the sense of Definition~\ref{def:terminology}(2), with the range $\rho=1$, in terms of the datum $u$, that is, of the data~(i) and~(iii): for small $X$, the numbers $\boldsymbol{\alpha}_0(\mathscr{E},X)$ and $\boldsymbol{\alpha}_2(\mathscr{E},X)$; and the kernels of
\[
(R\,\mathscr{E})(X),\qquad(R\,\mathscr{E})^{\mathrm{nat}}(X),\qquad S_L\left[(R\,\mathscr{E})^{\mathrm{nat}}(X)\right].
\]
Each of these objects is the value of $\Phi$ at any lift of $X$ and at the datum $u$ transported to that lift. Indeed, $\mathscr{E}(X;\cdot\,,\cdot)$ takes even values (Section~\ref{sec:spaces-activities}), so that its kernels have an even number of arguments; $(R\,\mathscr{E})(X)$ differs from $\mathscr{E}(X)$ by elements of degree $0$ and $2$ (Section~\ref{sec:local-parts-adjustment}); and the homomorphisms $\boldsymbol{\Theta}_{X,\mathscr{T}}$ of Definition~\ref{def:transport-mean} and $S_L$ of Proposition~\ref{prop:properties-actions}(d) map generators to odd elements, namely to linear combinations of generators, and hence even elements to even elements. Consequently the three elements displayed above are even and their kernels have an even number of arguments, and the assertion follows from Definition~\ref{def:terminology}(2) and Section~\ref{sec:antiperiodicity}. Likewise, the index sets of the sums defining $\varepsilon(\mathscr{E})$ and $m(\mathscr{E})$ are independent of $l$. Indeed, by Proposition~\ref{prop:polymers}\eqref{it:polymers-t10}, the small polymers containing a given $M$-cube are in one-to-one correspondence with the lifts $\tilde{X}'$ satisfying $d_M(\tilde{X}')\le L$ which contain a fixed lift of the cube, and the set of these lifts is the same for all $l$.
\end{remark}

\subsection{Bounds on the unit lattice}\label{sec:bounds-unit-lattice}

In this section the hypotheses~\eqref{eq:polymer-activities-Sc} are in force. Let $C_F$ and $\gamma_F$ be the constants of Propositions~\ref{prop:decay-fine-propagator}--\ref{prop:bounds} (Section~\ref{sec:fermion-operators-bounds}; Theorem~\ref{thm:fermion-operators}(i)), and let
\[
C_0^{\mathrm{unit}}=4\,b_{+}+16\,C_F\cdot c(\gamma_F)
\]
be the number of Section~\ref{sec:absolute-constants}. The operators $D_A$, $P_k(A)$, $Q_k(A)$, $Q_k^T(-A)$, $S_k(A)$, $D_k(A)$, $M_k(A)=b_k-D_k(A)$, $H_k(A)$ and $\bar{H}_k(A)=\mathsf{C}^{-1}\cdot H_k(-A)\cdot\mathsf{C}$ are those of Section~\ref{sec:exact-fermion-operators}. The operator $H_k^{\#}(A)=\left(H_k(A),\delta_{\alpha,A}H_k(A)\right)$, on the barred generators $\left(\bar{H}_k(A),\delta_{\alpha,-A}\bar{H}_k(A)\right)$, and the pairings $\left\langle\bar{\Psi},O\Psi\right\rangle$, $\langle\bar{\psi},\psi\rangle$ are those of Section~\ref{sec:fermion-integrals-blocking}, with $\delta_{\alpha,A}$ as in Definition~\ref{def:pairs-sites-holder}. The norm $|\cdot|_a$ on $\mathrm{Gr}[\Psi]$ is that of~\eqref{it:weighted-coefficient-norms-N7}, and
\[
h_k=\left(h_{k,1},h_{k,2}\right)=\left(e_k^{-\frac{1}{4}},e_k^{-\frac{1}{4}+\varepsilon}\right)
\]
are the weights of Section~\ref{sec:grassmann-algebras-norms}.

\begin{theorem}[Unit-lattice bounds]\label{thm:unit-lattice-bounds}
Assume that the following statements, which are contained in Propositions~\ref{prop:decay-fine-propagator}, \ref{prop:kernel-bounds-2} and~\ref{prop:bounds}, hold for every $A\in\tilde{G}_k$.
\begin{itemize}
\item The operator $D_A+\mathfrak{m}+b_k\cdot P_k(A)$ is invertible, and its inverse $S_k(A)$ is an analytic function of $A$ on $\tilde{G}_k$ (Theorem~\ref{thm:expansion-walks}(a) and~(b)).
\item For all unit sites $y,y'\in T^0_k$ the $4\times4$ block $M_k(A)(y,y')$ satisfies
\[
\left|M_k(A)(y,y')\right|\le C_F\cdot\exp\left(-\gamma_F\cdot d(y,y')\right)
\]
(Proposition~\ref{prop:kernels-unit-lattice}(d)).
\item For every spinor function $f$ on $T^0_k$,
\[
\left|H_k^{\#}(A)f\right|_{\infty}\le C_F\cdot|f|_{\infty},
\]
the supremum on the left being taken over the sites of $T_k$ and over all pairs of $\mathrm{Pz}_k$ (Proposition~\ref{prop:operators}(c)).
\end{itemize}
Let $C_0\ge\max\left(C_0^{\mathrm{unit}},C_F\right)$, and put
\[
a_1:=e_k^{-\frac{1}{4}},\qquad a_2:=C_0^{-1}\cdot e_k^{-\frac{1}{4}+\varepsilon}.
\]
\begin{enumerate}
\item[(a)] For every $A\in\tilde{G}_k$,
\[
\left|\left\langle\bar{\Psi},D_k(A)\Psi\right\rangle\right|_{a_1}\le C_0\cdot e_k^{-\frac{1}{2}}\cdot l^3,
\]
and the entries of $D_k(A)$ are analytic functions of $A$ on $\tilde{G}_k$.
\item[(b)] For every $A\in\tilde{G}_k$, the element $\langle\bar{\psi},\psi\rangle$ of $\mathrm{Gr}[\Psi]$ obtained by the substitution $\psi=H_k(A)\Psi$, $\bar{\psi}=\bar{H}_k(A)\bar{\Psi}$ satisfies
\[
\left|\langle\bar{\psi},\psi\rangle\right|_{a_2}\le4\,e_k^{-\frac{1}{2}+2\varepsilon}\cdot l^3\le4\,e_k^{-\frac{1}{2}}\cdot l^3.
\]
\item[(c)] Let $(\mathscr{A},\mathscr{C})\in\tilde{R}_k\times\mathfrak{X}_k$, $\mathscr{T}:=\mathscr{A}+\mathscr{C}$, $\mathscr{E}\in K^{x,C}_k$ and $X\in D_k$, and let $\mathscr{E}\left(X;\mathscr{A},\mathscr{C},H_k^{\#}(\mathscr{T})\Psi\right)\in\mathrm{Gr}[\Psi]$ be obtained from $\mathscr{E}(X;\mathscr{A},\mathscr{C})$ by the substitution $\psi^{\#}=H_k^{\#}(\mathscr{T})\Psi$ of Section~\ref{sec:fermion-integrals-blocking}. Then
\[
\left|\mathscr{E}\left(X;\mathscr{A},\mathscr{C},H_k^{\#}(\mathscr{T})\Psi\right)\right|_{a_2}\le\left\|\mathscr{E}(X;\mathscr{A},\mathscr{C})\right\|_{h_k}.
\]
\end{enumerate}
\end{theorem}

\begin{proof}
By~\eqref{eq:polymer-activities-Sc} and the first inequality $e_0^{\varepsilon}\le\frac{1}{4}$ of condition~\eqref{eq:renormalisation-group-A8}, one has $e_k\le e_0<1$.

(a) Let $O$ be a complex matrix whose rows and columns are indexed by the pairs $(y,\boldsymbol{\beta})$ of a unit site $y\in T^0_k$ and a spinor index $\boldsymbol{\beta}$. Then
\[
\left\langle\bar{\Psi},O\Psi\right\rangle=\sum_{y,\boldsymbol{\beta}}\ \sum_{y',\boldsymbol{\beta}'}O(y,\boldsymbol{\beta};y',\boldsymbol{\beta}')\cdot\bar{\Psi}_{\boldsymbol{\beta}}(y)\cdot\Psi_{\boldsymbol{\beta}'}(y'),
\]
and the products $\bar{\Psi}_{\boldsymbol{\beta}}(y)\cdot\Psi_{\boldsymbol{\beta}'}(y')$ are, up to signs, pairwise distinct basis monomials of degree two of $\mathrm{Gr}[\Psi]$. Hence, for every $a>0$,
\[
\left|\left\langle\bar{\Psi},O\Psi\right\rangle\right|_a=a^2\cdot\sum_{y,\boldsymbol{\beta}}\ \sum_{y',\boldsymbol{\beta}'}\left|O(y,\boldsymbol{\beta};y',\boldsymbol{\beta}')\right|.
\]
Apply this to $O=D_k(A)=b_k-M_k(A)$, and fix $y\in T^0_k$. By the second hypothesis, read through~\eqref{eq:polymer-activities-11}, each of the sixteen entries of the block $M_k(A)(y,y')$ has modulus at most $C_F\cdot\exp\left(-\gamma_F\cdot d(y,y')\right)$. Moreover $b_k\le b_{+}$ (Section~\ref{sec:running-quantities-tori}), and $\sum_{y'}\exp\left(-\gamma\cdot d(y,y')\right)\le c(\gamma)$ for every $\gamma>0$ (Section~\ref{sec:absolute-constants}). Therefore
\[
\sum_{\boldsymbol{\beta}}\ \sum_{y',\boldsymbol{\beta}'}\left|D_k(A)(y,\boldsymbol{\beta};y',\boldsymbol{\beta}')\right|\le4\,b_k+16\,C_F\cdot\sum_{y'}\exp\left(-\gamma_F\cdot d(y,y')\right)\le4\,b_{+}+16\,C_F\cdot c(\gamma_F)=C_0^{\mathrm{unit}}.
\]
The unit lattice $T^0_k=\mathbb{Z}^3/l\mathbb{Z}^3$ has $l^3$ sites, and $a_1^2=e_k^{-\frac{1}{2}}$. It follows that
\[
\left|\left\langle\bar{\Psi},D_k(A)\Psi\right\rangle\right|_{a_1}\le e_k^{-\frac{1}{2}}\cdot l^3\cdot C_0^{\mathrm{unit}}\le C_0\cdot e_k^{-\frac{1}{2}}\cdot l^3.
\]
By their displays in Section~\ref{sec:exact-fermion-operators}, the entries of $Q_k(A)$ and of $Q_k^T(-A)$ are entire functions of $A$, and $S_k(A)$ is analytic on $\tilde{G}_k$ by the first hypothesis. Hence the entries of
\[
D_k(A)=b_k-b_k^2\cdot Q_k(A)\cdot S_k(A)\cdot Q_k^T(-A)
\]
are analytic functions of $A$ on $\tilde{G}_k$.

(b) Let first $\langle\bar{\psi},\psi\rangle$ be regarded as an element of $\mathrm{Gr}[\psi]$, with a weight $g_1>0$ for the generators $\psi(\xi)$. In
\[
\langle\bar{\psi},\psi\rangle=\sum_{x\in T_k}\eta^3\cdot\sum_{\boldsymbol{\beta}}\bar{\psi}_{\boldsymbol{\beta}}(x)\cdot\psi_{\boldsymbol{\beta}}(x)
\]
the products $\bar{\psi}_{\boldsymbol{\beta}}(x)\cdot\psi_{\boldsymbol{\beta}}(x)$ are, up to signs, pairwise distinct basis monomials of degree two, each with a coefficient of modulus $\eta^3$. Since $\sum_{x\in T_k}\eta^3=l^3$, the definition~\eqref{eq:polymer-activities-1} of the norm gives
\[
\left\|\langle\bar{\psi},\psi\rangle\right\|_{g_1}=g_1^2\cdot4\sum_{x\in T_k}\eta^3=4\,l^3\cdot g_1^2.
\]
By the third hypothesis and Section~\ref{sec:coefficients-kernels-grassmann}, $\left|H_k(A)\right|_{1,\infty}\le C_F$. The set $\tilde{G}_k$ is invariant under $A\mapsto-A$ (Proposition~\ref{prop:properties-actions}(c)); hence also $\left|H_k(-A)\right|_{1,\infty}\le C_F$. The matrix $\mathsf{C}$ is monomial (Section~\ref{sec:gamma-matrices-symmetry}); so multiplication by $\mathsf{C}$ or by $\mathsf{C}^{-1}$ preserves the supremum of a spinor function over the sites and the components, and therefore
\[
\left|\bar{H}_k(A)\right|_{1,\infty}=\left|\mathsf{C}^{-1}\cdot H_k(-A)\cdot\mathsf{C}\right|_{1,\infty}=\left|H_k(-A)\right|_{1,\infty}\le C_F\le C_0.
\]
The substitution $\psi=H_k(A)\Psi$, $\bar{\psi}=\bar{H}_k(A)\bar{\Psi}$ is a homomorphism as in~\eqref{it:weighted-coefficient-norms-N2}, with the family $\psi$ indexed by $\{\xi\}$, the family $\Psi$ indexed by $\{u\}$ with the weight $1$, and the operator that acts as $H_k(A)$ on the unbarred and as $\bar{H}_k(A)$ on the barred generators; the row sums of this operator are at most $C_F$. Take the weight $a_2$ for $\Psi$ and the weight
\[
g_1:=C_0\cdot a_2=e_k^{-\frac{1}{4}+\varepsilon}
\]
for $\psi$. Then $C_F\cdot a_2\le g_1$, and~\eqref{it:weighted-coefficient-norms-N2} gives
\[
\left|\langle\bar{\psi},\psi\rangle\right|_{a_2}\le4\,l^3\cdot g_1^2=4\,e_k^{-\frac{1}{2}+2\varepsilon}\cdot l^3.
\]
Since $e_k<1$, one has $e_k^{2\varepsilon}\le1$, and the second inequality of~(b) follows.

(c) First, $\mathscr{T}\in\tilde{G}_k$. Indeed, let $\mathscr{A}=A_0+A_1$ be a decomposition as in Definition~\ref{def:domains}, with $c=1$: $A_1\in R_k$, $A_0$ is real, and for every $M$-cube $\square$ there are a real function $\boldsymbol{\lambda}$ on the sites and a real bond function $A'$ on $\square^{\dagger}$ with $A_0=A'+d\boldsymbol{\lambda}$ on $\square^{\dagger}$ and $|A'|<t_{1,k}$ on $\square^{\dagger}$. Since $L\ge5$ (stage~\ref{it:order-choice-4}), the cube $\square^{\sim5}$, of side $11M$, is contained in $\square^{\dagger}$, of side $(2L+1)\cdot M$ (Definition~\ref{def:enlarged-cubes}); and $t_{1,k}=e_k^{-\frac{3}{4}+\varepsilon}\le e_k^{-1+\varepsilon}$, because $e_k<1$. Thus $A_0=A'+d\boldsymbol{\lambda}$ on $\square^{\sim5}$, with real $\boldsymbol{\lambda}$ and $A'$ and with $|A'|<e_k^{-1+\varepsilon}$ on $\square^{\sim5}$. Moreover $\varepsilon\le\bar{\varepsilon}=\frac{1}{200}$ (stages~\ref{it:order-choice-8} and~\ref{it:order-choice-1}) and $e_k\le e_0<1$, so the first inequality of condition~\eqref{eq:renormalisation-group-A8} gives
\[
e_k^{\frac{1}{4}}\le e_k^{\varepsilon}\le e_0^{\varepsilon}\le\frac{1}{4}.
\]
Since $|A_1|<t_{1,k}$ and $|\mathscr{C}|<x_{1,k}$, it follows that
\[
\left|A_1+\mathscr{C}\right|<t_{1,k}+x_{1,k}=\left(e_k^{\frac{1}{4}}+e_k^{\varepsilon}\right)\cdot e_k^{-1+\varepsilon}\le e_k^{-1+\varepsilon},
\]
that is, $A_1+\mathscr{C}\in G_k$. By Definition~\ref{def:domains}, $\mathscr{T}=A_0+\left(A_1+\mathscr{C}\right)\in\tilde{G}_k$, and $-\mathscr{T}\in\tilde{G}_k$ by Proposition~\ref{prop:properties-actions}(c).

By Section~\ref{sec:fermion-integrals-blocking}, the substitution $\psi^{\#}=H_k^{\#}(\mathscr{T})\Psi$ inserts the fields
\[
\psi=H_k(\mathscr{T})\Psi,\qquad\chi=\delta_{\alpha,\mathscr{T}}H_k(\mathscr{T})\Psi,\qquad\bar{\psi}=\bar{H}_k(\mathscr{T})\bar{\Psi},\qquad\bar{\chi}=\delta_{\alpha,-\mathscr{T}}\bar{H}_k(\mathscr{T})\bar{\Psi}.
\]
The operator $\delta_{\alpha,-\mathscr{T}}$ acts on the sites only and commutes with the multiplication by the constant matrix $\mathsf{C}^{-1}$; hence
\[
\delta_{\alpha,-\mathscr{T}}\bar{H}_k(\mathscr{T})=\mathsf{C}^{-1}\cdot\left[\delta_{\alpha,B}H_k(B)\right]_{B=-\mathscr{T}}\cdot\mathsf{C}.
\]
The third hypothesis at $A=\mathscr{T}$ and at $A=-\mathscr{T}$, together with the invariance of the supremum norm under the multiplication by $\mathsf{C}$ and by $\mathsf{C}^{-1}$, shows that each of the four operators
\[
H_k(\mathscr{T}),\qquad\delta_{\alpha,\mathscr{T}}H_k(\mathscr{T}),\qquad\bar{H}_k(\mathscr{T}),\qquad\delta_{\alpha,-\mathscr{T}}\bar{H}_k(\mathscr{T})
\]
has row sums at most $C_F\le C_0$, the row sums of the second and of the fourth being taken over all pairs of $\mathrm{Pz}_k$. Let $H^{(1)}$ be the operator from functions on $\{u\}$ to functions on $\{\xi\}$ that acts as $H_k(\mathscr{T})$ on the unbarred and as $\bar{H}_k(\mathscr{T})$ on the barred generators, and let $H^{(2)}$ be the operator from functions on $\{u\}$ to functions on $\{\boldsymbol{\zeta}\}$ formed in the same way with $\delta_{\alpha,\mathscr{T}}H_k(\mathscr{T})$ and $\delta_{\alpha,-\mathscr{T}}\bar{H}_k(\mathscr{T})$. Then $\left|H^{(i)}\right|_{1,\infty}\le C_0$ for $i=1,2$, and
\[
\mathscr{E}\left(X;\mathscr{A},\mathscr{C},H_k^{\#}(\mathscr{T})\Psi\right)=E\left(H^{(1)}\Psi,H^{(2)}\Psi\right)\qquad\text{with}\qquad E:=\mathscr{E}(X;\mathscr{A},\mathscr{C}).
\]
Finally, since $e_k<1$,
\[
C_0\cdot a_2=e_k^{-\frac{1}{4}+\varepsilon}=h_{k,2}\le h_{k,1}=e_k^{-\frac{1}{4}},
\]
so that $C_0\cdot a_2\le\min\left(h_{k,1},h_{k,2}\right)$. Proposition~\ref{prop:weighted-coefficient-norms}\eqref{it:weighted-coefficient-norms-N7}(iv), applied with $a=a_2$, gives the assertion~(c).
\end{proof}

In the setting~\eqref{eq:polymer-activities-Sc} one has $e_k\le e_0\le e^F_0$, by condition~\eqref{eq:renormalisation-group-A10}, and Theorem~\ref{thm:fermion-operators}(i) gives Propositions~\ref{prop:decay-fine-propagator}, \ref{prop:kernel-bounds-2} and~\ref{prop:bounds} with $(C_F,\gamma_F)=(C^F,\gamma^F)$. Thus the three assumptions of Theorem~\ref{thm:unit-lattice-bounds} hold for every $A\in\tilde{G}_k$, and $C_0^{\mathrm{unit}}=4b_{+}+16\,C^F\cdot c(\gamma^F)$. The constant $C_0$ of stage~\ref{it:order-choice-7} is the largest of three numbers, two of which are $C^F$ and $4b_{+}+16\,C^F\cdot c(\gamma^F)$; hence $C_0\ge\max\left(C_0^{\mathrm{unit}},C_F\right)$, and Theorem~\ref{thm:unit-lattice-bounds} holds with this constant $C_0$. Moreover $C_0\ge C_F\ge1$ (Section~\ref{sec:fermion-operators-bounds}) and $e_k\le e_0<1$, so that
\[
a_2=C_0^{-1}\cdot h_{k,2}\le h_{k,2}\le h_{k,1}=a_1;
\]
and the element $\left\langle\bar{\Psi},D_k(A)\Psi\right\rangle$ of $\mathrm{Gr}[\Psi]$ is an analytic function of $A$ on $\tilde{G}_k$, its coefficients being, up to signs, the entries of $D_k(A)$. Proposition~\ref{prop:bounds-unit-lattice} follows.

\subsection{The block covariance at real connections}\label{sec:block-covariance-real}

The setting of this subsection is~\eqref{eq:polymer-activities-Sc}, except that the mass $\mathfrak{m}$ may be any real number $\mathfrak{m}>0$. In the setting~\eqref{eq:polymer-activities-Sc} itself one has $\mathfrak{m}\le\mathfrak{m}_0\le1$.

Let $\ell^2(T_k)$ be the space of the spinor functions on $T_k$ (Section~\ref{sec:antiperiodicity}) with the scalar product
\[
(f,g):=\sum_{x\in T_k}\eta^3\cdot\overline{f(x)}\cdot g(x),\qquad\text{where}\qquad\overline{u}\cdot w:=\sum_{\boldsymbol{\beta}}\overline{u_{\boldsymbol{\beta}}}\cdot w_{\boldsymbol{\beta}}\quad\left(u,w\in\mathbb{C}^4\right),
\]
and let $\ell^2(T^0_k)$ be the space of the spinor functions on the unit lattice $T^0_k$ with the scalar product $(v,w):=\sum_{y\in T^0_k}\overline{v(y)}\cdot w(y)$, formed with the counting measure. In both cases the summand is a periodic function, so that the sum over the sites of the torus is well defined. In both spaces, and in every finite-dimensional complex Hilbert space, $|f|:=(f,f)^{\frac{1}{2}}$; for a linear map $T$ between two such spaces, $|T|$ denotes its operator norm and $T^{*}$ its adjoint. If $T$ maps such a space into itself, $\operatorname{Re}T:=\frac{1}{2}\left(T+T^{*}\right)$, and $\operatorname{Re}T\ge c$ means that $\operatorname{Re}(f,Tf)\ge c\cdot|f|^2$ for every $f$; for a Hermitian $T$, $T\ge0$ means that $(f,Tf)\ge0$ for every $f$. A number standing for an operator denotes that number times the identity. On $\mathbb{C}^4$ put $|u|:=\left(\overline{u}\cdot u\right)^{\frac{1}{2}}$, and for a $4\times4$ matrix $B$ let $|B|$ be its operator norm with respect to $|\cdot|$. An operator $O$ on $\ell^2(T^0_k)$ is written $(Ov)(y)=\sum_{y'}O(y,y')\cdot v(y')$ with $4\times4$ blocks $O(y,y')$; the numbers $O(y,\boldsymbol{\beta};y',\boldsymbol{\beta}'):=O(y,y')_{\boldsymbol{\beta}\boldsymbol{\beta}'}$ are the entries of the matrix of $O$, whose rows and columns are indexed by the pairs $(y,\boldsymbol{\beta})$ of a unit site $y$ and a spinor index $\boldsymbol{\beta}$. Here, as for the kernels of Section~\ref{sec:antiperiodicity}, unit sites are represented by points of $\mathbb{Z}^3$; by~\eqref{eq:polymer-activities-AP}, another choice of the representatives multiplies a block, an entry, or the value of a spinor function at a unit site by $1$ or by $-1$, so that the moduli and the norms of these quantities depend on the unit sites alone. Finally, $\mathbf{1}[y=y']$ is $1$ if $y=y'$ and $0$ otherwise.

For a bond function $A$ on $T_k$, a unit site $y$ and a site $x\in B_k(y)$ put
\[
\sigma_A(y,x):=e_k\cdot\eta\cdot(\tau_kA)(y,x),
\]
with $\tau_k$ as in Section~\ref{sec:settings-basic-notation}. The blocks $B_k(y)=\Delta_y$, $y\in T^0_k$, partition the sites of $T_k$ (Section~\ref{sec:balls-torus-lifts}), and the operators of Section~\ref{sec:exact-fermion-operators} are given by
\begin{gather*}
\left(Q_k(A)f\right)(y)=\eta^3\sum_{x\in B_k(y)}\exp\left(i\,\sigma_A(y,x)\right)\cdot f(x),\qquad\left(Q_k^T(-A)v\right)(x)=\exp\left(-i\,\sigma_A(y,x)\right)\cdot v(y)\quad\left(x\in B_k(y)\right),\\
P_k(A)=Q_k^T(-A)\cdot Q_k(A),\qquad S_k(A)=\left(D_A+\mathfrak{m}+b_k\cdot P_k(A)\right)^{-1},\qquad D_k(A)=b_k-b_k^2\cdot Q_k(A)\cdot S_k(A)\cdot Q_k^T(-A),
\end{gather*}
the last two whenever the inverse exists.

Let $A$ be a bond function on $T_k$. For a direction $\mu\in\{0,1,2\}$, with the unit vector $e_{\mu}$, and for $\boldsymbol{\sigma}\in\{1,-1\}$ put
\[
x^{\boldsymbol{\sigma}}:=x+\boldsymbol{\sigma}\cdot\eta\cdot e_{\mu},\qquad P^{\boldsymbol{\sigma}}_{\mu}:=\frac{1}{2}\left(1-\boldsymbol{\sigma}\cdot\gamma_{\mu}\right),\qquad\left(T^{\boldsymbol{\sigma}}_{\mu}(A)f\right)(x):=\exp\left(i\,e_k\cdot\eta\cdot A(x,x^{\boldsymbol{\sigma}})\right)\cdot f(x^{\boldsymbol{\sigma}}),
\]
the dependence of $x^{\boldsymbol{\sigma}}$ on $\mu$ being suppressed. Write $P^{\pm}_{\mu}$ and $T^{\pm}_{\mu}(A)$ for $\boldsymbol{\sigma}=\pm1$, and $T_{\mu}:=T^{+}_{\mu}(A)$. Since $P^{+}_{\mu}+P^{-}_{\mu}=1$, the definition of $D_A$ in Section~\ref{sec:exact-fermion-operators} takes the form
\begin{equation}\label{eq:polymer-activities-22}
D_A=\eta^{-1}\sum_{\mu=0}^2\sum_{\boldsymbol{\sigma}=\pm1}P^{\boldsymbol{\sigma}}_{\mu}\cdot\left(1-T^{\boldsymbol{\sigma}}_{\mu}(A)\right).
\end{equation}
By~\eqref{eq:polymer-activities-13} the matrix $\gamma_{\mu}$ is Hermitian and $\gamma_{\mu}^2=1$. Hence every $P^{\boldsymbol{\sigma}}_{\mu}$ is Hermitian and satisfies $\left(P^{\boldsymbol{\sigma}}_{\mu}\right)^2=P^{\boldsymbol{\sigma}}_{\mu}$; that is, the $P^{\boldsymbol{\sigma}}_{\mu}$ are orthogonal projections of $\mathbb{C}^4$, and they act on spinor functions at every site.

\begin{theorem}[Block covariance]\label{thm:block-covariance}
Let $\mathfrak{m}>0$ be real and let $A$ be a real bond function on $T_k$. Then $\operatorname{Re}\left(D_A+\mathfrak{m}\right)\ge\mathfrak{m}$, the operator $D_A+\mathfrak{m}+b_k\cdot P_k(A)$ is invertible, and $D_k(A)$ is invertible, with
\begin{gather*}
D_k(A)^{-1}=b_k^{-1}+Q_k(A)\cdot\left(D_A+\mathfrak{m}\right)^{-1}\cdot Q_k(A)^{*},\qquad Q_k(A)^{*}=Q_k^T(-A),\\
\left|D_k(A)^{-1}\right|\le b_k^{-1}+\mathfrak{m}^{-1},\qquad\left|D_k(A)\right|\le b_k\qquad\left(\text{operator norms on }\ell^2(T^0_k)\right).
\end{gather*}
If moreover $\mathfrak{m}\le12\exp(1)$, then the $4\times4$ blocks of $D_k(A)^{-1}$ satisfy, for all unit sites $y$ and $y'$,
\[
\left|D_k(A)^{-1}(y,y')\right|\le b_k^{-1}\cdot\mathbf{1}[y=y']+\frac{2\exp(1)}{\mathfrak{m}}\cdot\exp\left(-\frac{\mathfrak{m}}{12\exp(1)}\cdot d(y,y')\right),
\]
and, for every unit site $y$ and every spinor index $\boldsymbol{\beta}$, the sums
\[
\sum_{y'\in T^0_k}\sum_{\boldsymbol{\beta}'}\left|D_k(A)^{-1}(y,\boldsymbol{\beta};y',\boldsymbol{\beta}')\right|\qquad\text{and}\qquad\sum_{y'\in T^0_k}\sum_{\boldsymbol{\beta}'}\left|D_k(A)^{-1}(y',\boldsymbol{\beta}';y,\boldsymbol{\beta})\right|,
\]
which are the $\ell^1$ norms of the row with the index $(y,\boldsymbol{\beta})$ and of the column with the same index in the matrix of $D_k(A)^{-1}$, do not exceed
\[
2b_k^{-1}+\frac{4\exp(1)}{\mathfrak{m}}\cdot\left(\frac{144\exp(1)}{\mathfrak{m}}\right)^3.
\]
\end{theorem}

\begin{proof}
Let $V$ be a finite-dimensional complex Hilbert space, $T$ a linear map of $V$ into itself and $c>0$ a number such that $\operatorname{Re}(f,Tf)\ge c\cdot|f|^2$ for every $f\in V$. Then
\begin{equation}\label{eq:polymer-activities-23}
T\ \text{is invertible},\qquad\left|T^{-1}\right|\le\frac{1}{c},\qquad\operatorname{Re}T^{-1}\ge0.
\end{equation}
Indeed, $|Tf|\cdot|f|\ge\left|(f,Tf)\right|\ge c\cdot|f|^2$, so that $|Tf|\ge c\cdot|f|$ for every $f$; thus $T$ is injective, hence invertible, and $\left|T^{-1}\right|\le c^{-1}$. For $g=Tf$ one has $\operatorname{Re}\left(g,T^{-1}g\right)=\operatorname{Re}(Tf,f)=\operatorname{Re}(f,Tf)\ge0$.

Let $V_1$ and $V_2$ be finite-dimensional complex vector spaces, $\mathbf{R}$ an invertible linear map of $V_1$ onto itself, $Q\colon V_1\to V_2$ and $Q'\colon V_2\to V_1$ linear maps, $b>0$ a number, and $c:=Q\cdot\mathbf{R}^{-1}\cdot Q'$, a linear map of $V_2$ into itself. Then either both $\mathbf{R}+b\cdot Q'\cdot Q$ and $b^{-1}+c$ are invertible or neither of them is; in the first case the inverse $S:=\left(\mathbf{R}+b\cdot Q'\cdot Q\right)^{-1}$ satisfies
\begin{gather}
S=\mathbf{R}^{-1}-\mathbf{R}^{-1}\cdot Q'\cdot\left(b^{-1}+c\right)^{-1}\cdot Q\cdot\mathbf{R}^{-1},\qquad Q\cdot S\cdot Q'=c-c\cdot\left(b^{-1}+c\right)^{-1}\cdot c=c\cdot\left(b^{-1}+c\right)^{-1}\cdot b^{-1},\notag\\
b-b^2\cdot Q\cdot S\cdot Q'=b\cdot\left[\left(b^{-1}+c\right)-c\right]\cdot\left(b^{-1}+c\right)^{-1}=\left(b^{-1}+c\right)^{-1}.\label{eq:polymer-activities-24}
\end{gather}
Indeed, $\mathbf{R}+b\cdot Q'\cdot Q=\mathbf{R}\cdot\left(1+b\cdot\mathbf{R}^{-1}\cdot Q'\cdot Q\right)$, and Sylvester's identity, a consequence of~\eqref{eq:polymer-activities-cl3}, gives
\[
\det\left(1+\left(b\cdot\mathbf{R}^{-1}\cdot Q'\right)\cdot Q\right)=\det\left(1+Q\cdot\left(b\cdot\mathbf{R}^{-1}\cdot Q'\right)\right)=\det\left(1+b\cdot c\right)=\det\left(b\cdot\left(b^{-1}+c\right)\right);
\]
this proves the first assertion. Now let $\mathbf{R}+b\cdot Q'\cdot Q$ be invertible. Since $Q\cdot\mathbf{R}^{-1}\cdot Q'=c$ and $1+b\cdot c=b\cdot\left(b^{-1}+c\right)$,
\begin{multline*}
\left(\mathbf{R}+b\cdot Q'\cdot Q\right)\cdot\left(\mathbf{R}^{-1}-\mathbf{R}^{-1}\cdot Q'\cdot\left(b^{-1}+c\right)^{-1}\cdot Q\cdot\mathbf{R}^{-1}\right)\\
=1+Q'\cdot\left[b-\left(1+b\cdot c\right)\cdot\left(b^{-1}+c\right)^{-1}\right]\cdot Q\cdot\mathbf{R}^{-1}=1,
\end{multline*}
so that the right side of the formula for $S$ is the inverse of $\mathbf{R}+b\cdot Q'\cdot Q$; this is the Woodbury identity. Multiplying the formula for $S$ by $Q$ from the left and by $Q'$ from the right gives $Q\cdot S\cdot Q'=c-c\cdot\left(b^{-1}+c\right)^{-1}\cdot c$, and this equals $c\cdot\left(b^{-1}+c\right)^{-1}\cdot\left[\left(b^{-1}+c\right)-c\right]=c\cdot\left(b^{-1}+c\right)^{-1}\cdot b^{-1}$. As $b^{-1}$ is a number, the expression for $b-b^2\cdot Q\cdot S\cdot Q'$ in~\eqref{eq:polymer-activities-24} follows from the one for $Q\cdot S\cdot Q'$.

\emph{The inverse of $D_k(A)$.} Let $\mathfrak{m}>0$, let $A$ be real, and put $\mathbf{R}:=D_A+\mathfrak{m}$. For every direction $\mu$ the operator $T_{\mu}$ is unitary on $\ell^2(T_k)$: the translation $f\mapsto f(\cdot+\eta e_{\mu})$ maps the antiperiodic functions bijectively onto themselves and preserves the scalar product, and $T_{\mu}$ is this translation followed by the multiplication with the function $x\mapsto\exp\left(i\,e_k\cdot\eta\cdot A(x,x+\eta e_{\mu})\right)$, whose modulus is $1$. Moreover $T^{-}_{\mu}(A)=T_{\mu}^{*}$, because $A(x,x-\eta e_{\mu})=-A(x-\eta e_{\mu},x)$ and $A$ is real: the substitution $x=x'-\eta e_{\mu}$ gives, for $f,g\in\ell^2(T_k)$,
\[
\left(g,T_{\mu}f\right)=\sum_{x'\in T_k}\eta^3\cdot\overline{\exp\left(i\,e_k\cdot\eta\cdot A(x',x'-\eta e_{\mu})\right)\cdot g(x'-\eta e_{\mu})}\cdot f(x')=\left(T^{-}_{\mu}(A)g,f\right).
\]
Since $P^{\pm}_{\mu}=\frac{1}{2}\left(1\mp\gamma_{\mu}\right)$,
\[
P^{+}_{\mu}\cdot T_{\mu}+P^{-}_{\mu}\cdot T_{\mu}^{*}=\frac{1}{2}\left(T_{\mu}+T_{\mu}^{*}\right)-\frac{1}{2}\gamma_{\mu}\cdot\left(T_{\mu}-T_{\mu}^{*}\right),
\]
and~\eqref{eq:polymer-activities-22}, together with $P^{+}_{\mu}+P^{-}_{\mu}=1$, gives
\[
D_A=\eta^{-1}\sum_{\mu=0}^2\left(1-\frac{1}{2}\left(T_{\mu}+T_{\mu}^{*}\right)\right)+\eta^{-1}\sum_{\mu=0}^2\frac{1}{2}\gamma_{\mu}\cdot\left(T_{\mu}-T_{\mu}^{*}\right).
\]
The operators of the first sum are Hermitian and $\ge0$, because $\operatorname{Re}\left(f,T_{\mu}f\right)\le|f|\cdot\left|T_{\mu}f\right|=|f|^2$. The operators of the second sum are anti-Hermitian: $\gamma_{\mu}$ is Hermitian by~\eqref{eq:polymer-activities-13}, and it commutes with $T_{\mu}$ and $T_{\mu}^{*}$, which act on the argument of a spinor function and not on its spinor index. Hence $\operatorname{Re}(f,D_Af)\ge0$ for every $f$, and $\operatorname{Re}\mathbf{R}\ge\mathfrak{m}$.

For $v\in\ell^2(T^0_k)$ and $f\in\ell^2(T_k)$ one has, $\sigma_A$ being real,
\[
\left(v,Q_k(A)f\right)=\sum_{y\in T^0_k}\overline{v(y)}\cdot\eta^3\sum_{x\in B_k(y)}\exp\left(i\,\sigma_A(y,x)\right)\cdot f(x)=\left(Q_k^T(-A)v,f\right);
\]
thus $Q_k(A)^{*}=Q_k^T(-A)$. The block $B_k(y)$ has $\#|B_k(y)|=\eta^{-3}$ sites; therefore
\[
\left(Q_k(A)\cdot Q_k(A)^{*}v\right)(y)=\eta^3\sum_{x\in B_k(y)}\exp\left(i\,\sigma_A(y,x)\right)\cdot\exp\left(-i\,\sigma_A(y,x)\right)\cdot v(y)=v(y),
\]
that is, $Q_k(A)\cdot Q_k(A)^{*}=1$. Consequently $Q_k(A)^{*}$ is an isometry, $\left|Q_k(A)\right|=\left|Q_k(A)^{*}\right|=1$, and $P_k(A)=Q_k(A)^{*}\cdot Q_k(A)$ is an orthogonal projection: it is Hermitian, and $P_k(A)^2=Q_k(A)^{*}\cdot\left(Q_k(A)\cdot Q_k(A)^{*}\right)\cdot Q_k(A)=P_k(A)$.

By~\eqref{eq:polymer-activities-23}, applied to $T=\mathbf{R}$ and $c=\mathfrak{m}$, the operator $\mathbf{R}$ is invertible, $\left|\mathbf{R}^{-1}\right|\le\mathfrak{m}^{-1}$ and $\operatorname{Re}\mathbf{R}^{-1}\ge0$. As $P_k(A)\ge0$ and $b_k>0$ (Section~\ref{sec:running-quantities-tori}), one has $\operatorname{Re}\left(\mathbf{R}+b_k\cdot P_k(A)\right)\ge\mathfrak{m}$; by~\eqref{eq:polymer-activities-23} once more, $D_A+\mathfrak{m}+b_k\cdot P_k(A)=\mathbf{R}+b_k\cdot P_k(A)$ is invertible. The operator
\[
c:=Q_k(A)\cdot\mathbf{R}^{-1}\cdot Q_k(A)^{*}
\]
on $\ell^2(T^0_k)$ satisfies $\operatorname{Re}c\ge0$, because $\operatorname{Re}(v,cv)=\operatorname{Re}\left(Q_k(A)^{*}v,\mathbf{R}^{-1}\cdot Q_k(A)^{*}v\right)\ge0$, and $|c|\le\left|Q_k(A)\right|\cdot\left|\mathbf{R}^{-1}\right|\cdot\left|Q_k(A)^{*}\right|\le\mathfrak{m}^{-1}$. Apply~\eqref{eq:polymer-activities-24} with $V_1=\ell^2(T_k)$, $V_2=\ell^2(T^0_k)$, $Q=Q_k(A)$, $Q'=Q_k(A)^{*}=Q_k^T(-A)$ and $b=b_k$. Here $\mathbf{R}+b_k\cdot Q'\cdot Q=D_A+\mathfrak{m}+b_k\cdot P_k(A)$ is invertible; so $b_k^{-1}+c$ is invertible, $S=S_k(A)$, and the identity for $b-b^2\cdot Q\cdot S\cdot Q'$ in~\eqref{eq:polymer-activities-24} reads
\[
D_k(A)=b_k-b_k^2\cdot Q_k(A)\cdot S_k(A)\cdot Q_k^T(-A)=\left(b_k^{-1}+c\right)^{-1}.
\]
Thus $D_k(A)$ is invertible, and $D_k(A)^{-1}=b_k^{-1}+c$; by the definitions of $c$ and of $\mathbf{R}$, this is the expression for $D_k(A)^{-1}$ stated in the theorem. Moreover $\left|D_k(A)^{-1}\right|\le b_k^{-1}+|c|\le b_k^{-1}+\mathfrak{m}^{-1}$. Finally $\operatorname{Re}\left(b_k^{-1}+c\right)\ge b_k^{-1}$, and~\eqref{eq:polymer-activities-23}, applied to $T=b_k^{-1}+c$ with the number $b_k^{-1}$ as the lower bound, gives $\left|D_k(A)\right|=\left|\left(b_k^{-1}+c\right)^{-1}\right|\le b_k$.

\emph{Decay of the blocks.} Let $\mathfrak{m}\le12\exp(1)$, fix unit sites $y$ and $y'$, and put $\boldsymbol{\theta}:=\frac{\mathfrak{m}}{12\exp(1)}$, so that $0<\boldsymbol{\theta}\le1$. Define
\[
\phi(x):=\boldsymbol{\theta}\cdot\operatorname{dist}\left(x,\Delta_{y'}\right),\qquad\operatorname{dist}\left(x,\Delta_{y'}\right):=\min_{z\in\Delta_{y'}}d(x,z),
\]
$d$ being the sup distance on the torus $T_k$. Then $\phi$ is a function on the torus, so that the multiplication by $\exp(\pm\phi)$ maps $\ell^2(T_k)$ into itself; it is Hermitian, $\phi$ being real. For nearest neighbours $x$ and $x'$ one has $\left|\phi(x)-\phi(x')\right|\le\boldsymbol{\theta}\cdot d(x,x')=\boldsymbol{\theta}\cdot\eta$. For every $\mu$ and $\boldsymbol{\sigma}$,
\[
\left(\exp(\phi)\cdot T^{\boldsymbol{\sigma}}_{\mu}(A)\cdot\exp(-\phi)f\right)(x)=\exp\left(\phi(x)-\phi(x^{\boldsymbol{\sigma}})\right)\cdot\left(T^{\boldsymbol{\sigma}}_{\mu}(A)f\right)(x);
\]
so $\exp(\phi)\cdot T^{\boldsymbol{\sigma}}_{\mu}(A)\cdot\exp(-\phi)-T^{\boldsymbol{\sigma}}_{\mu}(A)$ is the operator $T^{\boldsymbol{\sigma}}_{\mu}(A)$ followed by the multiplication with the function $x\mapsto\exp\left(\phi(x)-\phi(x^{\boldsymbol{\sigma}})\right)-1$, whose modulus is at most $\exp(\boldsymbol{\theta}\cdot\eta)-1$, because $x^{\boldsymbol{\sigma}}$ is a nearest neighbour of $x$ and $\left|\exp(t)-1\right|\le\exp(|t|)-1$ for real $t$. As $T^{\boldsymbol{\sigma}}_{\mu}(A)$ is unitary ($T^{-}_{\mu}(A)=T_{\mu}^{*}$), the norm of this difference is at most $\exp(\boldsymbol{\theta}\cdot\eta)-1$. The orthogonal projection $P^{\boldsymbol{\sigma}}_{\mu}$ has norm $1$ and commutes with the multiplication by $\exp(\pm\phi)$, and so does the number $\mathfrak{m}$; hence~\eqref{eq:polymer-activities-22} gives
\[
\exp(\phi)\cdot\mathbf{R}\cdot\exp(-\phi)-\mathbf{R}=\eta^{-1}\sum_{\mu=0}^2\sum_{\boldsymbol{\sigma}=\pm1}P^{\boldsymbol{\sigma}}_{\mu}\cdot\left(T^{\boldsymbol{\sigma}}_{\mu}(A)-\exp(\phi)\cdot T^{\boldsymbol{\sigma}}_{\mu}(A)\cdot\exp(-\phi)\right),
\]
a sum of six terms, each of norm at most $\eta^{-1}\cdot\left(\exp(\boldsymbol{\theta}\cdot\eta)-1\right)$. It follows that
\[
\left|\exp(\phi)\cdot\mathbf{R}\cdot\exp(-\phi)-\mathbf{R}\right|\le6\eta^{-1}\cdot\left(\exp(\boldsymbol{\theta}\cdot\eta)-1\right)\le6\boldsymbol{\theta}\cdot\exp(\boldsymbol{\theta}\cdot\eta)\le6\exp(1)\cdot\boldsymbol{\theta}=\frac{\mathfrak{m}}{2};
\]
here $\exp(s)-1\le s\cdot\exp(s)$ for $s\ge0$ and $\boldsymbol{\theta}\cdot\eta\le1$ were used. Consequently $\operatorname{Re}\left(\exp(\phi)\cdot\mathbf{R}\cdot\exp(-\phi)\right)\ge\mathfrak{m}-\frac{\mathfrak{m}}{2}=\frac{\mathfrak{m}}{2}$, and~\eqref{eq:polymer-activities-23} gives
\[
\left|\left(\exp(\phi)\cdot\mathbf{R}\cdot\exp(-\phi)\right)^{-1}\right|\le\frac{2}{\mathfrak{m}}.
\]
Let $v,v'\in\mathbb{C}^4$ be unit vectors, let $v\cdot\delta_y\in\ell^2(T^0_k)$ be the spinor function with the value $v$ at $y$ and the value $0$ at the other unit sites, and define $v'\cdot\delta_{y'}$ in the same way. The functions
\[
F:=Q_k(A)^{*}\left(v\cdot\delta_y\right),\qquad F':=Q_k(A)^{*}\left(v'\cdot\delta_{y'}\right)
\]
vanish outside $\Delta_y$ and outside $\Delta_{y'}$, respectively, and $|F|=|F'|=1$, because $Q_k(A)^{*}$ is an isometry. As $\left(\exp(\phi)\cdot\mathbf{R}\cdot\exp(-\phi)\right)^{-1}=\exp(\phi)\cdot\mathbf{R}^{-1}\cdot\exp(-\phi)$,
\[
\left(F,\mathbf{R}^{-1}F'\right)=\left(\exp(-\phi)\cdot F,\left(\exp(\phi)\cdot\mathbf{R}\cdot\exp(-\phi)\right)^{-1}\cdot\exp(\phi)\cdot F'\right).
\]
On $\Delta_{y'}$ one has $\phi=0$. On $\Delta_y$ one has $\phi\ge\boldsymbol{\theta}\cdot\left(d(y,y')-1\right)$, since $d(y,y')\le d(y,x)+d(x,z)+d(z,y')<d(x,z)+1$ for $x\in\Delta_y$ and $z\in\Delta_{y'}$. Hence $\exp(\phi)\cdot F'=F'$, $\left|\exp(-\phi)\cdot F\right|\le\exp\left(-\boldsymbol{\theta}\cdot\left(d(y,y')-1\right)\right)$, and
\[
\left|\left(F,\mathbf{R}^{-1}F'\right)\right|\le\frac{2}{\mathfrak{m}}\cdot\exp\left(-\boldsymbol{\theta}\cdot\left(d(y,y')-1\right)\right)\le\frac{2\exp(1)}{\mathfrak{m}}\cdot\exp\left(-\boldsymbol{\theta}\cdot d(y,y')\right),
\]
the last step by $\boldsymbol{\theta}\le1$. By the formula for $D_k(A)^{-1}$ already proved,
\[
\overline{v}\cdot D_k(A)^{-1}(y,y')\cdot v'=\left(v\cdot\delta_y,D_k(A)^{-1}\left(v'\cdot\delta_{y'}\right)\right)=b_k^{-1}\cdot\mathbf{1}[y=y']\cdot\overline{v}\cdot v'+\left(F,\mathbf{R}^{-1}F'\right).
\]
Since $\left|\overline{v}\cdot v'\right|\le1$, and since the norm of a $4\times4$ matrix $B$ is the supremum of $\left|\overline{v}\cdot B\cdot v'\right|$ over the unit vectors $v$ and $v'$, the bound on $\left|D_k(A)^{-1}(y,y')\right|$ follows.

\emph{Rows and columns.} Let $\mathfrak{m}\le12\exp(1)$, and let $\boldsymbol{\theta}=\frac{\mathfrak{m}}{12\exp(1)}$ be as above. For every $4\times4$ matrix $B$ and every spinor index $\boldsymbol{\beta}$,
\[
\sum_{\boldsymbol{\beta}'}\left|B_{\boldsymbol{\beta}\boldsymbol{\beta}'}\right|\le2|B|\qquad\text{and}\qquad\sum_{\boldsymbol{\beta}'}\left|B_{\boldsymbol{\beta}'\boldsymbol{\beta}}\right|\le2|B|.
\]
Indeed, the Cauchy--Schwarz inequality gives $\sum_{\boldsymbol{\beta}'}\left|u_{\boldsymbol{\beta}'}\right|\le\sqrt{4}\cdot|u|=2|u|$ for every $u\in\mathbb{C}^4$, and each of the two sums is of this form with a vector $u$ such that $|u|\le|B|$: in the second sum $u$ is the image under $B$ of the $\boldsymbol{\beta}$-th vector of the standard basis of $\mathbb{C}^4$, and in the first sum $u$ is the complex conjugate of the image of that basis vector under $B^{*}$, where $\left|B^{*}\right|=|B|$. The row $(y,\boldsymbol{\beta})$ of the matrix of $D_k(A)^{-1}$ consists of the rows $\boldsymbol{\beta}$ of the blocks $D_k(A)^{-1}(y,y')$, $y'\in T^0_k$, and the column $(y,\boldsymbol{\beta})$ consists of the columns $\boldsymbol{\beta}$ of the blocks $D_k(A)^{-1}(y',y)$. By these two inequalities, by the bound on the blocks, and as $d(y',y)=d(y,y')$, the $\ell^1$ norm of each of them is at most
\[
2\sum_{y'\in T^0_k}\left(b_k^{-1}\cdot\mathbf{1}[y=y']+\frac{2\exp(1)}{\mathfrak{m}}\cdot\exp\left(-\boldsymbol{\theta}\cdot d(y,y')\right)\right)\le2b_k^{-1}+\frac{4\exp(1)}{\mathfrak{m}}\cdot c(\boldsymbol{\theta}),
\]
where $c(\cdot)$ is the function of Section~\ref{sec:absolute-constants}, which dominates $\sum_{y'}\exp\left(-\boldsymbol{\theta}\cdot d(y,y')\right)$ on every torus. As $\boldsymbol{\theta}\le1$, Section~\ref{sec:absolute-constants} also gives
\[
c(\boldsymbol{\theta})\le\left(\frac{12}{\boldsymbol{\theta}}\right)^3=\left(\frac{144\exp(1)}{\mathfrak{m}}\right)^3,
\]
and the bound on the rows and the columns follows.

In the setting~\eqref{eq:polymer-activities-Sc} one has $\mathfrak{m}\le\mathfrak{m}_0\le1<12\exp(1)$; hence Proposition~\ref{prop:block-covariance-real} follows.
\end{proof}

\subsection{The free fermion determinant}\label{sec:free-fermion-determinant}

\subsubsection{Determinant identities}\label{sec:determinant-identities}

Let $k\ge1$ be an integer, let $l$ be a positive integer multiple of $L$, let $B$ be a complex bond function on $T_k=\mathbb{L}_k/l\mathbb{Z}^3$, and let $z\in\mathbb{C}$ be the mass. The spinor functions on $T_k$ are the functions $f\colon\mathbb{L}_k\to\mathbb{C}^4$ with $f(x+le_{\mu})=-f(x)$ for every $x\in\mathbb{L}_k$ and every direction $\mu$ (Section~\ref{sec:antiperiodicity}). The spinor functions on $T^0_k=\mathbb{Z}^3/l\mathbb{Z}^3$ and on $L\mathbb{Z}^3/l\mathbb{Z}^3$ are the functions on $\mathbb{Z}^3$, resp.\ on $L\mathbb{Z}^3$, with values in $\mathbb{C}^4$ and with the same antiperiodicity. The complex vector spaces of these functions are denoted by $\mathcal{F}(T_k)$, $\mathcal{F}(T^0_k)$ and $\mathcal{F}(L\mathbb{Z}^3/l\mathbb{Z}^3)$. Restriction to a fundamental domain of $l\mathbb{Z}^3$ shows that their dimensions are $4N_f$, $4n_k$ and $4n_{k+1}$, where
\[
N_f:=l^3\eta_k^{-3},\qquad n_k=l^3,\qquad n_{k+1}=\frac{l^3}{L^3}
\]
are the numbers of the sites of $T_k$, of $T^0_k$ and of $L\mathbb{Z}^3/l\mathbb{Z}^3$. For an endomorphism $T$ of a finite-dimensional complex vector space, $\det T$ and $\operatorname{Tr}T$ denote its determinant and its trace; a number standing for an operator denotes that number times the identity.

Put
\[
R_z:=D_B+z,
\]
where $D_B$ is the Wilson operator of level $k$ with the connection $B$ (Section~\ref{sec:exact-fermion-operators}); thus $R_z$ is an endomorphism of $\mathcal{F}(T_k)$. The averaging operators of Sections~\ref{sec:exact-fermion-operators} and~\ref{sec:fermion-integrals-blocking}, taken at $B$, are abbreviated by
\begin{align*}
Q_k&:=Q_k(B)\colon\mathcal{F}(T_k)\to\mathcal{F}(T^0_k),&Q_k'&:=Q_k^T(-B)\colon\mathcal{F}(T^0_k)\to\mathcal{F}(T_k),\\
Q&:=Q(B)\colon\mathcal{F}(T^0_k)\to\mathcal{F}(L\mathbb{Z}^3/l\mathbb{Z}^3),&Q'&:=Q^T(-B)\colon\mathcal{F}(L\mathbb{Z}^3/l\mathbb{Z}^3)\to\mathcal{F}(T^0_k).
\end{align*}
By their definitions, for $f\in\mathcal{F}(T_k)$, $v\in\mathcal{F}(T^0_k)$ and a unit site $y\in\mathbb{Z}^3$,
\begin{gather*}
\left(Q_kf\right)(y)=\eta_k^3\sum_{x\in B_k(y)}\exp\left(i\,e_k\cdot\eta_k\cdot(\tau_kB)(y,x)\right)\cdot f(x),\\
\left(Q_k'v\right)(x)=\exp\left(-i\,e_k\cdot\eta_k\cdot(\tau_kB)(y,x)\right)\cdot v(y)\qquad\left(x\in B_k(y)\right),
\end{gather*}
and for $v\in\mathcal{F}(T^0_k)$, $w\in\mathcal{F}(L\mathbb{Z}^3/l\mathbb{Z}^3)$ and $y\in L\mathbb{Z}^3$,
\begin{gather*}
\left(Qv\right)(y)=L^{-3}\sum_{u\in B(y)}\exp\left(i\,e_k\cdot\eta_k\cdot(\tau B)(y,u)\right)\cdot v(u),\\
\left(Q'w\right)(u)=\exp\left(-i\,e_k\cdot\eta_k\cdot(\tau B)(y,u)\right)\cdot w(y)\qquad\left(u\in B(y)\right).
\end{gather*}
Here $B_k(y)$ is the set of the sites $x\in\mathbb{L}_k$ with $|x-y|<\frac{1}{2}$, and $B(y)$ is the $L$-block of the unit sites $u$ with $|u-y|<\frac{L}{2}$ (Section~\ref{sec:lattice-geometry}); $\tau_kB$ is the function of Section~\ref{sec:settings-basic-notation}, and $(\tau B)(y,u)$ is formed as in Definition~\ref{def:pairs-sites-holder} along the rectilinear paths of the unit lattice from $y$ to $u$, $B$ being summed over their fine bonds (Section~\ref{sec:fermion-integrals-blocking}).

Put
\[
P_k:=Q_k'\cdot Q_k,\qquad P:=Q'\cdot Q,\qquad P_{k+1}:=Q_k'\cdot P\cdot Q_k;
\]
$P_k$ and $P_{k+1}$ are endomorphisms of $\mathcal{F}(T_k)$, and $P$ is an endomorphism of $\mathcal{F}(T^0_k)$. Let further
\[
\mathring{a}:=bL^{-1},\qquad\mathring{a}':=b_{k+1}L^{-1}=\frac{\mathring{a}\cdot b_k}{b_k+\mathring{a}};
\]
the second expression for $\mathring{a}'$ follows from the relation $b_{k+1}=\frac{b\cdot b_k}{b_k+bL^{-1}}$ of Section~\ref{sec:running-quantities-tori}. Since $b=1$ (stage~\ref{it:order-choice-1}) and $b_j\ge b\left(1-L^{-1}\right)$ for every $j\ge1$ (Section~\ref{sec:running-quantities-tori}), the numbers $b_k$, $b_{k+1}$, $\mathring{a}$ and $\mathring{a}'$ are positive.

The operator $Q_k\cdot Q_k'$ is the identity of $\mathcal{F}(T^0_k)$, and $Q\cdot Q'$ is the identity of $\mathcal{F}(L\mathbb{Z}^3/l\mathbb{Z}^3)$. Indeed, the two phase factors in each of these products are inverse to each other, for complex $B$ as well, and $\#|B_k(y)|=\eta_k^{-3}$, $\#|B(y)|=L^3$; hence
\[
\left(Q_k\cdot Q_k'v\right)(y)=\eta_k^3\cdot\#|B_k(y)|\cdot v(y)=v(y),\qquad\left(Q\cdot Q'w\right)(y)=L^{-3}\cdot\#|B(y)|\cdot w(y)=w(y).
\]

\begin{lemma}[Block determinants]\label{lem:block-determinants}
Whenever $R_z+b_k\cdot P_k$ is invertible, put
\[
S_k:=\left(R_z+b_k\cdot P_k\right)^{-1},\qquad D_k:=b_k-b_k^2\cdot Q_k\cdot S_k\cdot Q_k',\qquad\delta Z_k:=\det\left(D_k+\mathring{a}\cdot P\right);
\]
here $S_k$ is an endomorphism of $\mathcal{F}(T_k)$, $D_k$ is an endomorphism of $\mathcal{F}(T^0_k)$, and $\delta Z_k$ is a complex number. Whenever $R_z+\mathring{a}'\cdot P_{k+1}$ is invertible, let $D^0_{k+1}$ be the endomorphism of $\mathcal{F}(L\mathbb{Z}^3/l\mathbb{Z}^3)$ given by
\[
D^0_{k+1}:=b_{k+1}-b_{k+1}\cdot\mathring{a}'\cdot\left(Q\cdot Q_k\right)\cdot\left(R_z+\mathring{a}'\cdot P_{k+1}\right)^{-1}\cdot\left(Q_k'\cdot Q'\right).
\]
\begin{enumerate}
\item[(i)] If $R_z+b_k\cdot P_k$ is invertible, then
\[
\det D_k\cdot\det\left(R_z+b_k\cdot P_k\right)=b_k^{4n_k}\cdot\det R_z.
\]
\item[(ii)] If $R_z+b_k\cdot P_k$ is invertible, then
\[
\delta Z_k\cdot\det\left(R_z+b_k\cdot P_k\right)=b_k^{4(n_k-n_{k+1})}\cdot\left(b_k+\mathring{a}\right)^{4n_{k+1}}\cdot\det\left(R_z+\mathring{a}'\cdot P_{k+1}\right).
\]
\item[(iii)] If $R_z+b_k\cdot P_k$ is invertible and $\delta Z_k\ne0$, then $R_z+\mathring{a}'\cdot P_{k+1}$ is invertible, and
\[
\det D^0_{k+1}=b^{4n_{k+1}}\cdot\frac{\det D_k}{\delta Z_k}.
\]
For $k=0$, let $B$ be a complex bond function on $T_0=T^0_0=\mathbb{Z}^3/l\mathbb{Z}^3$, where $\eta_0=1$, let $z\in\mathbb{C}$, and let $R_z$, $Q$, $Q'$, $P$, $\mathring{a}$, $n_0=l^3$ and $n_1=l^3/L^3$ be defined as above with $k=0$. Let $Q_0$ and $Q_0'$ be the identity of $\mathcal{F}(T_0)$, so that $P_1=Q_0'\cdot P\cdot Q_0=P$, and put
\[
D_0:=R_z,\qquad\delta Z_0:=\det\left(R_z+\mathring{a}\cdot P\right),\qquad\mathring{a}':=b_1L^{-1}=bL^{-1},
\]
where $b_1=b$. Let $D^0_1$ be given by the formula for $D^0_{k+1}$ with $k=0$, that is,
\[
D^0_1=b-b\cdot\mathring{a}\cdot Q\cdot\left(R_z+\mathring{a}\cdot P\right)^{-1}\cdot Q'.
\]
If $\delta Z_0\ne0$, then $R_z+\mathring{a}\cdot P$ is invertible, and
\[
\det D^0_1=b^{4n_1}\cdot\frac{\det D_0}{\delta Z_0}.
\]
\item[(iv)] (Scaling.) Let $k\ge0$, let $B'$ be a complex bond function on $T_{k+1}=\mathbb{L}_{k+1}/(l/L)\mathbb{Z}^3$, and let $z'\in\mathbb{C}$. Form the objects above, at level $k$, with the mass $z:=z'/L$ and with $B:=B'_L$, that is, $B(x,\tilde{x})=L^{-\frac{1}{2}}\cdot B'\left(\frac{x}{L},\frac{\tilde{x}}{L}\right)$ for nearest neighbours $x$, $\tilde{x}$ of $\mathbb{L}_k$ (Section~\ref{sec:local-parts-adjustment}), for $k=0$ with the conventions of (iii), and write $S_k(B)$, $D_k(B)$ and $\delta Z_k(B)$ for $S_k$, $D_k$ and $\delta Z_k$. Let $\mathcal{F}(T_{k+1})$ and $\mathcal{F}(T^0_{k+1})$ be the spaces of the spinor functions on $T_{k+1}$ and on $T^0_{k+1}=\mathbb{Z}^3/(l/L)\mathbb{Z}^3$, defined with the side $l/L$ in place of $l$. The operators with the argument $B'$ are those of Section~\ref{sec:exact-fermion-operators} at level $k+1$, formed with $\eta_{k+1}$, $e_{k+1}$ and the mass $z'$; thus
\begin{gather*}
P_{k+1}(B')=Q_{k+1}^T(-B')\cdot Q_{k+1}(B'),\qquad S_{k+1}(B')=\left(D_{B'}+z'+b_{k+1}\cdot P_{k+1}(B')\right)^{-1},\\
D_{k+1}(B')=b_{k+1}-b_{k+1}^2\cdot Q_{k+1}(B')\cdot S_{k+1}(B')\cdot Q_{k+1}^T(-B').
\end{gather*}
Define $U\colon\mathcal{F}(T_{k+1})\to\mathcal{F}(T_k)$ and $U_0\colon\mathcal{F}(T^0_{k+1})\to\mathcal{F}(L\mathbb{Z}^3/l\mathbb{Z}^3)$ by
\[
(Uf)(x):=f\left(\frac{x}{L}\right)\quad\left(x\in\mathbb{L}_k\right),\qquad(U_0v)(y):=v\left(\frac{y}{L}\right)\quad\left(y\in L\mathbb{Z}^3\right);
\]
thus $U_0$ relabels the sites of $L\mathbb{Z}^3/l\mathbb{Z}^3$ by the sites of $T^0_{k+1}$. Then $U$ and $U_0$ are linear bijections,
\[
D_{B'}+z'=L\cdot U^{-1}\cdot\left(D_B+z\right)\cdot U,\qquad P_{k+1}(B')=U^{-1}\cdot P_{k+1}\cdot U,
\]
and $S_{k+1}(B')$ exists if and only if $R_z+\mathring{a}'\cdot P_{k+1}$ is invertible; in this case
\[
D_{k+1}(B')=U_0^{-1}\cdot D^0_{k+1}\cdot U_0.
\]
Consequently, if $\delta Z_k(B)\ne0$ and, in case $k\ge1$, $S_k(B)$ exists, then
\begin{equation}\label{eq:polymer-activities-25}
\det D_{k+1}(B')=b^{4n_{k+1}}\cdot\frac{\det D_k(B)}{\delta Z_k(B)},
\end{equation}
where the left side is formed at level $k+1$ with the mass $z'=Lz$, and the right side at level $k$ with the mass $z$.
\end{enumerate}
\end{lemma}

\begin{proof}
\emph{A determinant identity.} Let $V$ and $W$ be finite-dimensional complex vector spaces, let $R$ be an endomorphism of $V$, let $G\colon V\to W$ and $G'\colon W\to V$ be linear maps, and let $c\ne0$ be a complex number such that $R+c\cdot G'\cdot G$ is invertible. Then
\[
\det\left(R+c\cdot G'\cdot G\right)\cdot\det\left(c-c^2\cdot G\cdot\left(R+c\cdot G'\cdot G\right)^{-1}\cdot G'\right)=c^{\dim W}\cdot\det R.
\]
Indeed, let $M$ be the endomorphism of $V\oplus W$ given by
\[
M(v,w):=\left(\left(R+c\cdot G'\cdot G\right)v-c\cdot G'w,\;-c\cdot Gv+c\cdot w\right).
\]
With respect to the decomposition $V\oplus W$, its blocks are $A_1=R+c\cdot G'\cdot G$, $A_2=-c\cdot G'$, $A_3=-c\cdot G$ and $A_4=c$, and the blocks $A_1$ and $A_4$ are invertible. In bases of $V$ and of $W$, the first formula of~\eqref{eq:polymer-activities-cl3} gives
\[
\det M=\det\left(R+c\cdot G'\cdot G\right)\cdot\det\left(c-c^2\cdot G\cdot\left(R+c\cdot G'\cdot G\right)^{-1}\cdot G'\right),
\]
because $A_4-A_3\cdot A_1^{-1}\cdot A_2=c-c^2\cdot G\cdot\left(R+c\cdot G'\cdot G\right)^{-1}\cdot G'$; and the second formula of~\eqref{eq:polymer-activities-cl3} gives $\det M=c^{\dim W}\cdot\det R$, because $\det A_4=c^{\dim W}$ and
\[
A_1-A_2\cdot A_4^{-1}\cdot A_3=R+c\cdot G'\cdot G-c\cdot G'\cdot G=R.
\]

\emph{Proof of (i).} Apply the determinant identity with $V=\mathcal{F}(T_k)$, $W=\mathcal{F}(T^0_k)$, $R=R_z$, $G=Q_k$, $G'=Q_k'$ and $c=b_k$. Then $G'\cdot G=P_k$, so that $R+c\cdot G'\cdot G=R_z+b_k\cdot P_k$ is invertible by assumption, and
\[
c-c^2\cdot G\cdot\left(R+c\cdot G'\cdot G\right)^{-1}\cdot G'=b_k-b_k^2\cdot Q_k\cdot S_k\cdot Q_k'=D_k.
\]
As $\dim W=4n_k$, the identity becomes (i).

\emph{Proof of (ii).} The operator $P$ is idempotent, since $P^2=Q'\cdot\left(Q\cdot Q'\right)\cdot Q=Q'\cdot Q=P$, and
\[
\operatorname{Tr}P=\operatorname{Tr}\left(Q'\cdot Q\right)=\operatorname{Tr}\left(Q\cdot Q'\right)=4n_{k+1},
\]
because $Q\cdot Q'$ is the identity of $\mathcal{F}(L\mathbb{Z}^3/l\mathbb{Z}^3)$, a space of dimension $4n_{k+1}$. A finite-dimensional space is the direct sum of the image and of the kernel of an idempotent endomorphism, which is the identity on its image and vanishes on its kernel; so the trace of an idempotent endomorphism is the dimension of its image. Hence $\mathcal{F}(T^0_k)$ is the direct sum of the image of $P$, of dimension $4n_{k+1}$, on which $b_k+\mathring{a}\cdot P$ is the multiplication by $b_k+\mathring{a}$, and of the kernel of $P$, of dimension $4n_k-4n_{k+1}$, on which $b_k+\mathring{a}\cdot P$ is the multiplication by $b_k$. Therefore
\[
\det\left(b_k+\mathring{a}\cdot P\right)=b_k^{4(n_k-n_{k+1})}\cdot\left(b_k+\mathring{a}\right)^{4n_{k+1}}\ne0,\qquad\left(b_k+\mathring{a}\cdot P\right)^{-1}=b_k^{-1}-\frac{\mathring{a}}{b_k\left(b_k+\mathring{a}\right)}\cdot P;
\]
the formula for the inverse follows from $P^2=P$:
\[
\left(b_k+\mathring{a}\cdot P\right)\cdot\left(b_k^{-1}-\frac{\mathring{a}}{b_k\left(b_k+\mathring{a}\right)}\cdot P\right)=1+\frac{\mathring{a}\cdot\left(b_k+\mathring{a}\right)-\mathring{a}\cdot b_k-\mathring{a}^2}{b_k\left(b_k+\mathring{a}\right)}\cdot P=1.
\]
Let $M_2$ be the endomorphism of $\mathcal{F}(T_k)\oplus\mathcal{F}(T^0_k)$ given by
\[
M_2(f,v):=\left(\left(R_z+b_k\cdot P_k\right)f-b_k\cdot Q_k'v,\;-b_k\cdot Q_kf+\left(b_k+\mathring{a}\cdot P\right)v\right).
\]
Its blocks are $A_1=R_z+b_k\cdot P_k$, $A_2=-b_k\cdot Q_k'$, $A_3=-b_k\cdot Q_k$ and $A_4=b_k+\mathring{a}\cdot P$, and $A_1$ and $A_4$ are invertible. By the first formula of~\eqref{eq:polymer-activities-cl3},
\[
\det M_2=\det\left(R_z+b_k\cdot P_k\right)\cdot\det\left(b_k+\mathring{a}\cdot P-b_k^2\cdot Q_k\cdot S_k\cdot Q_k'\right)=\det\left(R_z+b_k\cdot P_k\right)\cdot\delta Z_k.
\]
By the second formula of~\eqref{eq:polymer-activities-cl3},
\[
\det M_2=\det\left(b_k+\mathring{a}\cdot P\right)\cdot\det\left(R_z+b_k\cdot P_k-b_k^2\cdot Q_k'\cdot\left(b_k+\mathring{a}\cdot P\right)^{-1}\cdot Q_k\right).
\]
The formula for $\left(b_k+\mathring{a}\cdot P\right)^{-1}$ and the relation $\mathring{a}'=\frac{\mathring{a}\cdot b_k}{b_k+\mathring{a}}$ give
\[
b_k^2\cdot Q_k'\cdot\left(b_k+\mathring{a}\cdot P\right)^{-1}\cdot Q_k=b_k\cdot Q_k'\cdot Q_k-\frac{\mathring{a}\cdot b_k}{b_k+\mathring{a}}\cdot Q_k'\cdot P\cdot Q_k=b_k\cdot P_k-\mathring{a}'\cdot P_{k+1},
\]
so that the second factor equals $\det\left(R_z+\mathring{a}'\cdot P_{k+1}\right)$. Comparing the two expressions for $\det M_2$, with the value of $\det\left(b_k+\mathring{a}\cdot P\right)$ found above, one obtains (ii).

\emph{Proof of (iii).} Let $\delta Z_k\ne0$. Since $\det\left(R_z+b_k\cdot P_k\right)\ne0$, the left side of (ii) does not vanish; hence $\det\left(R_z+\mathring{a}'\cdot P_{k+1}\right)\ne0$, and $R_z+\mathring{a}'\cdot P_{k+1}$ is invertible. The averaging maps $Q\cdot Q_k\colon\mathcal{F}(T_k)\to\mathcal{F}(L\mathbb{Z}^3/l\mathbb{Z}^3)$ and $Q_k'\cdot Q'\colon\mathcal{F}(L\mathbb{Z}^3/l\mathbb{Z}^3)\to\mathcal{F}(T_k)$ satisfy
\[
\left(Q\cdot Q_k\right)\cdot\left(Q_k'\cdot Q'\right)=Q\cdot\left(Q_k\cdot Q_k'\right)\cdot Q'=Q\cdot Q'=1,\qquad\left(Q_k'\cdot Q'\right)\cdot\left(Q\cdot Q_k\right)=Q_k'\cdot P\cdot Q_k=P_{k+1}.
\]
The determinant identity with $V=\mathcal{F}(T_k)$, $W=\mathcal{F}(L\mathbb{Z}^3/l\mathbb{Z}^3)$, $R=R_z$, $G=Q\cdot Q_k$, $G'=Q_k'\cdot Q'$ and $c=\mathring{a}'$ therefore applies, and since $\dim W=4n_{k+1}$ it gives
\begin{multline*}
\det\left(R_z+\mathring{a}'\cdot P_{k+1}\right)\cdot\det\left(\mathring{a}'-\left(\mathring{a}'\right)^2\cdot\left(Q\cdot Q_k\right)\cdot\left(R_z+\mathring{a}'\cdot P_{k+1}\right)^{-1}\cdot\left(Q_k'\cdot Q'\right)\right)\\
=\left(\mathring{a}'\right)^{4n_{k+1}}\cdot\det R_z.
\end{multline*}
The operator $D^0_{k+1}$ is $\frac{b_{k+1}}{\mathring{a}'}$ times the operator in the second determinant on the left. Hence
\[
\det D^0_{k+1}=b_{k+1}^{4n_{k+1}}\cdot\frac{\det R_z}{\det\left(R_z+\mathring{a}'\cdot P_{k+1}\right)}.
\]
On the other hand, division of (i) by (ii) gives
\[
\frac{\det D_k}{\delta Z_k}=\left(\frac{b_k}{b_k+\mathring{a}}\right)^{4n_{k+1}}\cdot\frac{\det R_z}{\det\left(R_z+\mathring{a}'\cdot P_{k+1}\right)}.
\]
By Section~\ref{sec:running-quantities-tori}, $b_{k+1}=b\cdot\frac{b_k}{b_k+\mathring{a}}$, that is, $b_{k+1}\cdot\left(1+\mathring{a}\cdot b_k^{-1}\right)=b$. The last two displays therefore give
\[
\det D^0_{k+1}=b^{4n_{k+1}}\cdot\left(\frac{b_k}{b_k+\mathring{a}}\right)^{4n_{k+1}}\cdot\frac{\det R_z}{\det\left(R_z+\mathring{a}'\cdot P_{k+1}\right)}=b^{4n_{k+1}}\cdot\frac{\det D_k}{\delta Z_k}.
\]

\emph{The case $k=0$ of (iii).} Here $Q\cdot Q_0=Q$, $Q_0'\cdot Q'=Q'$, $P_1=P$, $\mathring{a}'=\mathring{a}$ and $b_1=b$. Let $\delta Z_0=\det\left(R_z+\mathring{a}\cdot P\right)\ne0$; then $R_z+\mathring{a}\cdot P$ is invertible. The determinant identity with $V=\mathcal{F}(T_0)$, $W=\mathcal{F}(L\mathbb{Z}^3/l\mathbb{Z}^3)$, $R=R_z$, $G=Q$, $G'=Q'$ and $c=\mathring{a}$, for which $G'\cdot G=P$ and $\dim W=4n_1$, gives
\[
\det\left(R_z+\mathring{a}\cdot P\right)\cdot\det\left(\mathring{a}-\mathring{a}^2\cdot Q\cdot\left(R_z+\mathring{a}\cdot P\right)^{-1}\cdot Q'\right)=\mathring{a}^{4n_1}\cdot\det R_z.
\]
As $D^0_1$ is $\frac{b}{\mathring{a}}$ times the operator in the second determinant, it follows directly that
\[
\det D^0_1=b^{4n_1}\cdot\frac{\det R_z}{\det\left(R_z+\mathring{a}\cdot P\right)}=b^{4n_1}\cdot\frac{\det D_0}{\delta Z_0}.
\]

\emph{Proof of (iv): the maps $U$ and $U_0$.} The map $x\mapsto x/L$ is a bijection of $\mathbb{L}_k$ onto $\mathbb{L}_{k+1}$ that carries $x+le_{\mu}$ to $x/L+(l/L)e_{\mu}$. Hence $U$ maps $\mathcal{F}(T_{k+1})$ bijectively onto $\mathcal{F}(T_k)$, with $(U^{-1}g)(x')=g(Lx')$ for $g\in\mathcal{F}(T_k)$ and $x'\in\mathbb{L}_{k+1}$. The map $y\mapsto y/L$ is a bijection of $L\mathbb{Z}^3$ onto $\mathbb{Z}^3$ that carries $y+le_{\mu}$ to $y/L+(l/L)e_{\mu}$; hence $U_0$ is bijective, with $(U_0^{-1}w)(y')=w(Ly')$ for $w\in\mathcal{F}(L\mathbb{Z}^3/l\mathbb{Z}^3)$ and $y'\in\mathbb{Z}^3$.

\emph{Link phases.} Let $x'$ and $x''$ be nearest neighbours in $\mathbb{L}_{k+1}$. Then $Lx'$ and $Lx''$ are nearest neighbours in $\mathbb{L}_k$, and $B(Lx',Lx'')=L^{-\frac{1}{2}}\cdot B'(x',x'')$ by the definition of $B'_L$ (Section~\ref{sec:local-parts-adjustment}). Since $e_{k+1}=L^{\frac{1}{2}}\cdot e_k$ and $\eta_{k+1}=L^{-1}\cdot\eta_k$ (Section~\ref{sec:running-quantities-tori}), it follows that
\[
e_k\cdot\eta_k\cdot B(Lx',Lx'')=e_{k+1}\cdot\eta_{k+1}\cdot B'(x',x'').
\]

\emph{The Wilson operators.} By~\eqref{eq:polymer-activities-22}, at level $k$ and at level $k+1$,
\[
D_B=\eta_k^{-1}\sum_{\mu=0}^2\sum_{\boldsymbol{\sigma}=\pm1}P^{\boldsymbol{\sigma}}_{\mu}\cdot\left(1-T^{\boldsymbol{\sigma}}_{\mu}(B)\right),\qquad D_{B'}=\eta_{k+1}^{-1}\sum_{\mu=0}^2\sum_{\boldsymbol{\sigma}=\pm1}P^{\boldsymbol{\sigma}}_{\mu}\cdot\left(1-T^{\boldsymbol{\sigma}}_{\mu}(B')\right),
\]
where $P^{\boldsymbol{\sigma}}_{\mu}=\frac{1}{2}\left(1-\boldsymbol{\sigma}\cdot\gamma_{\mu}\right)$ acts on the spinor index and where, for $f\in\mathcal{F}(T_k)$, $f'\in\mathcal{F}(T_{k+1})$, $x\in\mathbb{L}_k$ and $x'\in\mathbb{L}_{k+1}$,
\begin{align*}
\left(T^{\boldsymbol{\sigma}}_{\mu}(B)f\right)(x)&=\exp\left(i\,e_k\cdot\eta_k\cdot B\left(x,x+\boldsymbol{\sigma}\cdot\eta_k\cdot e_{\mu}\right)\right)\cdot f\left(x+\boldsymbol{\sigma}\cdot\eta_k\cdot e_{\mu}\right),\\
\left(T^{\boldsymbol{\sigma}}_{\mu}(B')f'\right)(x')&=\exp\left(i\,e_{k+1}\cdot\eta_{k+1}\cdot B'\left(x',x'+\boldsymbol{\sigma}\cdot\eta_{k+1}\cdot e_{\mu}\right)\right)\cdot f'\left(x'+\boldsymbol{\sigma}\cdot\eta_{k+1}\cdot e_{\mu}\right).
\end{align*}
Since $Lx'+\boldsymbol{\sigma}\cdot\eta_k\cdot e_{\mu}=L\left(x'+\boldsymbol{\sigma}\cdot\eta_{k+1}\cdot e_{\mu}\right)$, the relation of the link phases gives
\begin{multline*}
\left(U^{-1}\cdot T^{\boldsymbol{\sigma}}_{\mu}(B)\cdot Uf'\right)(x')=\exp\left(i\,e_k\cdot\eta_k\cdot B\left(Lx',L\left(x'+\boldsymbol{\sigma}\cdot\eta_{k+1}\cdot e_{\mu}\right)\right)\right)\cdot f'\left(x'+\boldsymbol{\sigma}\cdot\eta_{k+1}\cdot e_{\mu}\right)\\
=\left(T^{\boldsymbol{\sigma}}_{\mu}(B')f'\right)(x').
\end{multline*}
The matrices $P^{\boldsymbol{\sigma}}_{\mu}$ act on the spinor index only and therefore commute with $U$. Hence
\[
U^{-1}\cdot D_B\cdot U=\eta_k^{-1}\sum_{\mu=0}^2\sum_{\boldsymbol{\sigma}=\pm1}P^{\boldsymbol{\sigma}}_{\mu}\cdot\left(1-T^{\boldsymbol{\sigma}}_{\mu}(B')\right)=\frac{\eta_{k+1}}{\eta_k}\cdot D_{B'}=L^{-1}\cdot D_{B'},
\]
and, as $z'=Lz$, one has $D_{B'}+z'=L\cdot U^{-1}\cdot\left(D_B+z\right)\cdot U$.

\emph{The averaging operators.} Let $y'\in\mathbb{Z}^3$, let $x'\in B_{k+1}(y')$, that is, $x'\in\mathbb{L}_{k+1}$ and $|x'-y'|<\frac{1}{2}$, and put $y:=Ly'\in L\mathbb{Z}^3$ and $x:=Lx'\in\mathbb{L}_k$. For $0\le j\le k+1$ let $y'_j$ be the centre of the block of side $L^j\eta_{k+1}$ that contains $x'$, the blocks of this side being the cubes of this side centred at the points of $L^j\eta_{k+1}\mathbb{Z}^3$; thus $y'_0=x'$ and $y'_{k+1}=y'$. Since $L\cdot\eta_{k+1}=\eta_k$, multiplication by $L$ maps the blocks of side $L^j\eta_{k+1}$ onto the blocks of side $L^j\eta_k$. Hence $y_j:=Ly'_j$ is the centre of the block of side $L^j\eta_k$ that contains $x$. In particular $y_0=x$, $y_k$ is the unit site with $x\in B_k(y_k)$, and $y_{k+1}=y$ is the centre of the cube of side $L$ that contains $x$; as $L$ is odd, this cube is the union of the unit cubes centred at the points of $B(y)$, so that $y_k\in B(y)$. Multiplication by $L$ maps each rectilinear path of $\mathbb{L}_{k+1}$ from $y'_{j+1}$ to $y'_j$, bond by bond, onto the rectilinear path of $\mathbb{L}_k$ from $y_{j+1}$ to $y_j$ with the same order of the directions. By the relation of the link phases and Definition~\ref{def:pairs-sites-holder},
\[
e_{k+1}\cdot\eta_{k+1}\cdot(\tau B')(y'_{j+1},y'_j)=e_k\cdot\eta_k\cdot(\tau B)(y_{j+1},y_j)\qquad(0\le j\le k).
\]
Summation over $j$ and the definitions of $\tau_{k+1}$ and of $\tau_k$ (Section~\ref{sec:settings-basic-notation}) give
\[
e_{k+1}\cdot\eta_{k+1}\cdot(\tau_{k+1}B')(y',x')=e_k\cdot\eta_k\cdot(\tau B)(y,y_k)+e_k\cdot\eta_k\cdot(\tau_kB)(y_k,x):
\]
the terms with $j<k$ add up to the exponent of $Q_k$ at $(y_k,x)$, and the term with $j=k$ is the exponent of $Q$ at $(y,y_k)$, which is formed along rectilinear paths of the unit lattice with $B$ summed over their fine bonds. For $k=0$ the sum over $j<k$ is empty and $B_0(u)=\{u\}$ for $u\in\mathbb{Z}^3$, so that the display of $Q_k$ with $k=0$ gives the identity, in accordance with the convention $Q_0=1$ of (iii).

As $x'$ ranges over $B_{k+1}(y')$, the site $x=Lx'$ ranges over the sites of $\mathbb{L}_k$ in the cube of side $L$ centred at $y$, that is, over the disjoint union of the blocks $B_k(u)$, $u\in B(y)$; and $y_k=u$ for $x\in B_k(u)$. Since $\eta_{k+1}^3=L^{-3}\cdot\eta_k^3$, it follows for $f\in\mathcal{F}(T_k)$ that
\begin{align*}
\left(U_0\cdot Q_{k+1}(B')\cdot U^{-1}f\right)(y)&=\eta_{k+1}^3\sum_{x'\in B_{k+1}(y')}\exp\left(i\,e_{k+1}\cdot\eta_{k+1}\cdot(\tau_{k+1}B')(y',x')\right)\cdot f(Lx')\\
&=L^{-3}\sum_{u\in B(y)}\exp\left(i\,e_k\cdot\eta_k\cdot(\tau B)(y,u)\right)\cdot\eta_k^3\sum_{x\in B_k(u)}\exp\left(i\,e_k\cdot\eta_k\cdot(\tau_kB)(u,x)\right)\cdot f(x)\\
&=\left(Q\cdot Q_kf\right)(y),
\end{align*}
that is, $U_0\cdot Q_{k+1}(B')\cdot U^{-1}=Q\cdot Q_k$. Let now $w\in\mathcal{F}(L\mathbb{Z}^3/l\mathbb{Z}^3)$ and $x\in\mathbb{L}_k$, let $y'$ be the unit site with $x/L\in B_{k+1}(y')$, and let $x'=x/L$, $y=Ly'$ and $u=y_k$ be as above, so that $x\in B_k(u)$ and $u\in B(y)$. The display of $Q_{k+1}^T(-B')$ in Section~\ref{sec:exact-fermion-operators}, at level $k+1$, and the decomposition of the exponent just obtained give
\begin{align*}
\left(U\cdot Q_{k+1}^T(-B')\cdot U_0^{-1}w\right)(x)&=\exp\left(-i\,e_{k+1}\cdot\eta_{k+1}\cdot(\tau_{k+1}B')(y',x')\right)\cdot w(y)\\
&=\exp\left(-i\,e_k\cdot\eta_k\cdot(\tau_kB)(u,x)\right)\cdot\exp\left(-i\,e_k\cdot\eta_k\cdot(\tau B)(y,u)\right)\cdot w(y)\\
&=\left(Q_k'\cdot Q'w\right)(x),
\end{align*}
that is, $U\cdot Q_{k+1}^T(-B')\cdot U_0^{-1}=Q_k'\cdot Q'$. Consequently
\[
P_{k+1}(B')=\left(U^{-1}\cdot Q_k'\cdot Q'\cdot U_0\right)\cdot\left(U_0^{-1}\cdot Q\cdot Q_k\cdot U\right)=U^{-1}\cdot Q_k'\cdot P\cdot Q_k\cdot U=U^{-1}\cdot P_{k+1}\cdot U.
\]

\emph{Conclusion.} By the relations obtained for $D_{B'}+z'$ and for $P_{k+1}(B')$,
\[
D_{B'}+z'+b_{k+1}\cdot P_{k+1}(B')=L\cdot U^{-1}\cdot\left(R_z+b_{k+1}L^{-1}\cdot P_{k+1}\right)\cdot U=L\cdot U^{-1}\cdot\left(R_z+\mathring{a}'\cdot P_{k+1}\right)\cdot U.
\]
As $U$ is bijective, $S_{k+1}(B')$ exists if and only if $R_z+\mathring{a}'\cdot P_{k+1}$ is invertible, and then
\[
S_{k+1}(B')=L^{-1}\cdot U^{-1}\cdot\left(R_z+\mathring{a}'\cdot P_{k+1}\right)^{-1}\cdot U.
\]
In this case, inserting this formula, together with $Q_{k+1}(B')=U_0^{-1}\cdot Q\cdot Q_k\cdot U$ and $Q_{k+1}^T(-B')=U^{-1}\cdot Q_k'\cdot Q'\cdot U_0$, into the definition of $D_{k+1}(B')$, and using $b_{k+1}^2\cdot L^{-1}=b_{k+1}\cdot\mathring{a}'$, one obtains
\begin{align*}
D_{k+1}(B')&=b_{k+1}-b_{k+1}^2\cdot L^{-1}\cdot U_0^{-1}\cdot Q\cdot Q_k\cdot\left(R_z+\mathring{a}'\cdot P_{k+1}\right)^{-1}\cdot Q_k'\cdot Q'\cdot U_0\\
&=U_0^{-1}\cdot\left(b_{k+1}-b_{k+1}\cdot\mathring{a}'\cdot\left(Q\cdot Q_k\right)\cdot\left(R_z+\mathring{a}'\cdot P_{k+1}\right)^{-1}\cdot\left(Q_k'\cdot Q'\right)\right)\cdot U_0\\
&=U_0^{-1}\cdot D^0_{k+1}\cdot U_0.
\end{align*}
Finally, let $\delta Z_k(B)\ne0$, and let $S_k(B)$ exist if $k\ge1$. By (iii) if $k\ge1$, and by its case $k=0$ if $k=0$, the operator $R_z+\mathring{a}'\cdot P_{k+1}$ is invertible and $\det D^0_{k+1}=b^{4n_{k+1}}\cdot\frac{\det D_k(B)}{\delta Z_k(B)}$. By what has been shown, $S_{k+1}(B')$ exists and $D_{k+1}(B')=U_0^{-1}\cdot D^0_{k+1}\cdot U_0$. The determinant is invariant under conjugation by the linear bijection $U_0$; therefore $\det D_{k+1}(B')=\det D^0_{k+1}$, and~\eqref{eq:polymer-activities-25} follows.
\end{proof}

\subsubsection{The determinant of the free operator}\label{sec:determinant-free-operator}

\begin{theorem}[Free determinant across the levels]\label{thm:free-determinant-across}
Assume hypothesis~\ref{def:hypothesis-flow}, and let $k$ be an integer with $0\le k\le k^{*}(R)$, where $k^{*}(R)$ is the stop level of Definition~\ref{def:tori-stop-level}. For $0\le j\le k$ the objects of level $j$ are formed on the torus $T_j$ of side $l_j$, with the charge of level $j$ and with the running mass $\bar{m}\cdot a_j$ of level $j$ (Section~\ref{sec:running-quantities-tori}). Following Sections~\ref{sec:running-quantities-tori} and~\ref{sec:exact-fermion-operators} and Definition~\ref{def:notation}(b), this mass is denoted $\mathfrak{m}_j$ if $j\ge1$, and $\bar{m}\cdot\eta$ if $j=0$, a product in which $\eta$ means the spacing $L^{-N}=a_0$ of Section~\ref{sec:lattice-measures-laws} rather than $\eta_k=L^{-k}$. In sums over polymers, $D_j$ is the set of the polymers of $T_j$ (Section~\ref{sec:lattice-geometry}); $D_j(A)$, $S_j(A)$ and $\delta Z_j(A)$ are the operators and the determinant of Section~\ref{sec:exact-fermion-operators} at level $j$; and $D_a$, $D_0$ are the operators $D_A$ of Section~\ref{sec:lattice-measures-laws} at $A=a$ and at $A=0$.

For every level $j<k$ let a family~\eqref{it:primitives-p5} of level $j$ on $T_j$ be given for which Proposition~\ref{prop:local-expansion-determinant}(d0) holds; the functions $E^{\mathrm{det}}_j$ of Definition~\ref{def:determinant-functions} and $\mathfrak{P}_j$ of Section~\ref{sec:flow} are formed with these families. Assume that Proposition~\ref{prop:decay-fine-propagator} holds at the levels $1\le j\le k$ in the following form: for every $A\in\tilde{G}_j$ the operator $D_A+\mathfrak{m}_j+b_j\cdot P_j(A)$ is invertible, that is, $S_j(A)$ exists (Theorem~\ref{thm:expansion-walks}(a)). Let $a\in C^1\otimes\mathbb{C}$, let $a^{[j]}(b)=L^{\frac{j}{2}}\cdot a\left(L^j\cdot b\right)$ be $a$ in the coordinates of level $j$ (Section~\ref{sec:lattice-measures-laws}), and assume that $a^{[j]}\in\mathfrak{X}_j$ for all $j\le k$. Then $D_k\left(a^{[k]}\right)$ is defined, $\det D_j(0)\ne0$ for $0\le j\le k$, and
\begin{align*}
\frac{\det\left(D_a+\bar{m}\cdot\eta\right)}{\det\left(D_0+\bar{m}\cdot\eta\right)}&=\exp\left\{\sum_{X\in D_k}\mathfrak{P}_k\left(X,a^{[k]}\right)\right\}\cdot\frac{\det D_k\left(a^{[k]}\right)}{\det D_k(0)},\\
\sum_{X\in D_k}\mathfrak{P}_k\left(X,a^{[k]}\right)&=\sum_{j<k}\;\sum_{X\in D_j}E^{\mathrm{det}}_j\left(X,a^{[j]}\right).
\end{align*}
\end{theorem}

\begin{proof}
Throughout, $j$ denotes an integer with $0\le j\le k$.

\emph{Charges.} The charge of level $j$ is $e_j=e\cdot a_j^{\frac{1}{2}}$, where $a_j=L^{j-N}$ (Section~\ref{sec:running-quantities-tori}); for $j=0$ it is the number $q=e\cdot a_0^{\frac{1}{2}}$ of Section~\ref{sec:lattice-measures-laws}, whereas $e_0$ and $e_1$ denote the constants of stages~\ref{it:order-choice-9} and~\ref{it:order-choice-10} (Definition~\ref{def:notation}(e) and~(g)). By this definition $e_j$ increases with $j$. Proposition~\ref{prop:depth-stop-level}(c), whose hypotheses hold under hypothesis~\ref{def:hypothesis-flow} by Proposition~\ref{prop:consequences}, gives $k^{*}(R)\le K$ (Section~\ref{sec:running-quantities-tori}). As $j\le k\le k^{*}(R)\le K$, it follows that $e_j\le e_K$. Hypothesis~\ref{def:hypothesis-flow} gives $e_K=e\cdot h'^{\frac{1}{2}}\le e_1$, stage~\ref{it:order-choice-10} gives $e_1\le e_0$, and the first inequality $e_0^{\varepsilon}\le\frac{1}{4}$ of condition~\eqref{eq:renormalisation-group-A8} gives $e_0<1$. Since $e>0$, one obtains
\[
0<e_j\le e_K\le e_1\le e_0<1\qquad(0\le j\le k).
\]

\emph{Sides and masses.} For every $i\le k^{*}(R)$ the side $l_i$ belongs to the class $\mathrm{Cl}$ of Section~\ref{sec:lattice-geometry}, by Proposition~\ref{prop:depth-stop-level}(f); the hypotheses of that proposition, condition~\eqref{eq:renormalisation-flow-3} among them, hold under hypothesis~\ref{def:hypothesis-flow} by Proposition~\ref{prop:consequences}. Thus $l_j\in\mathrm{Cl}$; in particular $l_j$ is a multiple of $L\cdot M$, hence of $L$. By Section~\ref{sec:running-quantities-tori} one has $a_{j+1}=L\cdot a_j$, $l_j=\frac{R}{a_j}$ and $\eta_j=L^{-j}$; therefore, for $j<k$,
\[
l_{j+1}=\frac{l_j}{L},\qquad\eta_{j+1}=\frac{\eta_j}{L},\qquad\bar{m}\cdot a_{j+1}=L\cdot\bar{m}\cdot a_j,
\]
and the charge of level $j+1$ is $L^{\frac{1}{2}}$ times the charge of level $j$. In the notation of the theorem, $\mathfrak{m}_{j+1}=L\cdot\mathfrak{m}_j$ for $j\ge1$, and $\mathfrak{m}_1=L\cdot\bar{m}\cdot\eta$. The running masses are positive, $\bar{m}\cdot a_j>0$, because $0<\bar{m}\le1$ by hypothesis~\ref{def:hypothesis-flow}.

\emph{Domains.} By Definition~\ref{def:domains}, every element of $\mathfrak{X}_j$ has sup norm less than $x_{1,j}=e_j^{-1+2\varepsilon}$ (Section~\ref{sec:settings-basic-notation}), and $G_j$ is the set of the complex bond functions of sup norm less than $e_j^{-1+\varepsilon}$. As $0<e_j<1$, one has $x_{1,j}\le e_j^{-1+\varepsilon}$, so that $\mathfrak{X}_j\subset G_j$; moreover $G_j\subset\tilde{G}_j$, as one sees by taking the real summand $A_0=0$ in the definition of $\tilde{G}_j$. Hence
\[
a^{[j]}\in\mathfrak{X}_j\subset G_j\subset\tilde{G}_j,\qquad 0\in\tilde{G}_j\qquad(0\le j\le k).
\]

\emph{The operators.} For $1\le j\le k$ the operators $S_j\left(a^{[j]}\right)$ and $S_j(0)$ exist, by the form of Proposition~\ref{prop:decay-fine-propagator} assumed in the theorem; so $D_j\left(a^{[j]}\right)$ and $D_j(0)$ are defined (Section~\ref{sec:exact-fermion-operators}). At level~$0$ one has $D_0(A)=D_A+\bar{m}\cdot\eta$ for every $A$ (Section~\ref{sec:exact-fermion-operators}), and $a^{[0]}=a$; hence
\[
D_0\left(a^{[0]}\right)=D_a+\bar{m}\cdot\eta,\qquad D_0(0)=D_0+\bar{m}\cdot\eta.
\]
In particular $D_k\left(a^{[k]}\right)$ is defined. Let $j<k$. By the hypothesis that Proposition~\ref{prop:local-expansion-determinant}(d0) holds at level $j$, for every $A\in\tilde{G}_j$ the number $\delta Z_j(A)$ is the exponential of a complex number, so that $\delta Z_j(A)\ne0$, and
\[
\frac{\delta Z_j(A)}{\delta Z_j(0)}=\exp\left\{\sum_{X\in D_j}E^{\mathrm{det}}_j(X,A)\right\}.
\]

\emph{Non-vanishing of $\det D_j(0)$.} Let $j=0$. For the real connection $A=0$ the operator $D_A$ is, by~\eqref{eq:polymer-activities-22}, the sum of a non-negative Hermitian operator and an anti-Hermitian one; this argument, with which the proof of Theorem~\ref{thm:block-covariance} begins, holds at level~$0$, where the lattice spacing is $\eta_0=1$, as well. Hence
\[
\operatorname{Re}\left(D_0+\bar{m}\cdot\eta\right)\ge\bar{m}\cdot\eta>0,
\]
and $D_0(0)=D_0+\bar{m}\cdot\eta$ is invertible by~\eqref{eq:polymer-activities-23}. Let $1\le j\le k$. Theorem~\ref{thm:block-covariance} holds for every real positive mass (Section~\ref{sec:block-covariance-real}); applied at level $j$, with the side $l_j\in\mathrm{Cl}$, the charge $e_j\in\left(0,e_0\right]$ and the mass $\mathfrak{m}_j>0$, to the real connection $A=0$, it shows that $D_j(0)$ is invertible. Therefore
\[
\det D_j(0)\ne0\qquad(0\le j\le k).
\]

\emph{One step.} Let $0\le j<k$. The scaling $A\mapsto A_L$ of Section~\ref{sec:local-parts-adjustment}, $A_L(b)=L^{-\frac{1}{2}}\cdot A\left(\frac{b}{L}\right)$, maps $a^{[j+1]}$ to $a^{[j]}$: for every bond $b$ of $T_j$,
\[
\left(a^{[j+1]}\right)_L(b)=L^{-\frac{1}{2}}\cdot a^{[j+1]}\left(\frac{b}{L}\right)=L^{-\frac{1}{2}}\cdot L^{\frac{j+1}{2}}\cdot a\left(L^{j+1}\cdot\frac{b}{L}\right)=L^{\frac{j}{2}}\cdot a\left(L^j\cdot b\right)=a^{[j]}(b).
\]
Apply Lemma~\ref{lem:block-determinants}(iv), that is,~\eqref{eq:polymer-activities-25}, with $j$ in place of $k$, with the side $l_j$, with the mass $z=\bar{m}\cdot a_j$ at level $j$ and with the mass $z'=L\cdot z=\bar{m}\cdot a_{j+1}$ at level $j+1$; by the relations between the levels $j$ and $j+1$ obtained above, the objects of level $j+1$ are those of the lemma. Take first $B'=a^{[j+1]}$, so that $B=B'_L=a^{[j]}$, and then $B'=0$, so that $B=0$. In both cases $B\in\tilde{G}_j$; hence $S_j(B)$ exists if $j\ge1$, and $\delta Z_j(B)\ne0$. With $n_{j+1}=l_{j+1}^3$, \eqref{eq:polymer-activities-25} gives
\[
\det D_{j+1}\left(a^{[j+1]}\right)=b^{4n_{j+1}}\cdot\frac{\det D_j\left(a^{[j]}\right)}{\delta Z_j\left(a^{[j]}\right)},\qquad\det D_{j+1}(0)=b^{4n_{j+1}}\cdot\frac{\det D_j(0)}{\delta Z_j(0)}.
\]
Both sides of the second identity are different from $0$, since $\det D_j(0)\ne0$. Dividing the first identity by the second and inserting the quotient $\delta Z_j\left(a^{[j]}\right)/\delta Z_j(0)$ obtained above at $A=a^{[j]}$, one finds
\[
\frac{\det D_{j+1}\left(a^{[j+1]}\right)}{\det D_{j+1}(0)}=\frac{\det D_j\left(a^{[j]}\right)}{\det D_j(0)}\cdot\exp\left\{-\sum_{X\in D_j}E^{\mathrm{det}}_j\left(X,a^{[j]}\right)\right\}\qquad(0\le j<k).
\]

\emph{The first identity.} The product of these identities over $0\le j<k$ telescopes:
\[
\frac{\det D_k\left(a^{[k]}\right)}{\det D_k(0)}=\frac{\det D_0\left(a^{[0]}\right)}{\det D_0(0)}\cdot\exp\left\{-\sum_{j<k}\;\sum_{X\in D_j}E^{\mathrm{det}}_j\left(X,a^{[j]}\right)\right\};
\]
for $k=0$ the product is empty and the exponent vanishes. Since $D_0\left(a^{[0]}\right)=D_a+\bar{m}\cdot\eta$ and $D_0(0)=D_0+\bar{m}\cdot\eta$, it follows that
\[
\frac{\det\left(D_a+\bar{m}\cdot\eta\right)}{\det\left(D_0+\bar{m}\cdot\eta\right)}=\exp\left\{\sum_{j<k}\;\sum_{X\in D_j}E^{\mathrm{det}}_j\left(X,a^{[j]}\right)\right\}\cdot\frac{\det D_k\left(a^{[k]}\right)}{\det D_k(0)}.
\]
By the second identity, proved below, the exponent equals $\sum_{X\in D_k}\mathfrak{P}_k\left(X,a^{[k]}\right)$; so this display is the first identity.

\emph{The second identity.} Let $0\le j<k$. The sides $l_j$ and $l_{j+1}=\frac{l_j}{L}$ belong to $\mathrm{Cl}$, by Proposition~\ref{prop:depth-stop-level}(f), whose hypotheses, condition~\eqref{eq:renormalisation-flow-3} among them, hold under hypothesis~\ref{def:hypothesis-flow} by Proposition~\ref{prop:consequences}. Hence clause~\eqref{it:polymers-t5} of Proposition~\ref{prop:polymers} applies at level $j$: each polymer $X\in D_j$ satisfies $\bar{X}^L=LY$ for precisely one polymer $Y\in D_{j+1}$. Consequently
\[
\sum_{Y\in D_{j+1}}\;\sum_{X\in D_j:\,\bar{X}^L=LY}f(X)=\sum_{X\in D_j}f(X)
\]
for every function $f$ on $D_j$. The recursion of Section~\ref{sec:flow},
\[
\mathfrak{P}_{j+1}(Y,\mathscr{C})=\sum_{X\in D_j:\,\bar{X}^L=LY}\left(\mathfrak{P}_j\left(X,\mathscr{C}_L\right)+E^{\mathrm{det}}_j\left(X,\mathscr{C}_L\right)\right)\qquad(Y\in D_{j+1}),
\]
taken at $\mathscr{C}=a^{[j+1]}$, for which $\mathscr{C}_L=a^{[j]}$ by the computation above, therefore gives
\[
\sum_{Y\in D_{j+1}}\mathfrak{P}_{j+1}\left(Y,a^{[j+1]}\right)=\sum_{X\in D_j}\left(\mathfrak{P}_j\left(X,a^{[j]}\right)+E^{\mathrm{det}}_j\left(X,a^{[j]}\right)\right).
\]
As $\mathfrak{P}_0=0$ (Section~\ref{sec:flow}), induction on $j$ yields
\[
\sum_{X\in D_j}\mathfrak{P}_j\left(X,a^{[j]}\right)=\sum_{i<j}\;\sum_{X\in D_i}E^{\mathrm{det}}_i\left(X,a^{[i]}\right)\qquad(0\le j\le k).
\]
For $j=k$ this is the second identity, and the first identity follows as shown above.

\emph{The data $\left(j,l_j,\bar{m}\cdot a_j,e_j\right)$ satisfy the hypotheses of the statements on $S_j$ and on $\delta Z_j$.} Let $0\le j\le k$. Under hypothesis~\ref{def:hypothesis-flow} the hypotheses of Proposition~\ref{prop:depth-stop-level}, condition~\eqref{eq:renormalisation-flow-3} included, are fulfilled, by Proposition~\ref{prop:consequences}. Hence the side $l_j$ lies in $\mathrm{Cl}$, by part~(f) of Proposition~\ref{prop:depth-stop-level}; and part~(e) of that proposition and Section~\ref{sec:running-quantities-tori}, together with $\mathfrak{m}_0\le1$ from Definition~\ref{def:mass-bound-stop}, give for the running mass of level $j$
\[
0<\bar{m}\cdot a_j\le\mathfrak{m}_K\le L^{-m_0}\le\mathfrak{m}_0\le1.
\]
Moreover $0<e_j\le e_0$, by the inequalities obtained for the charges. The constants $L$, $M$ and $e_0$ satisfy the conditions of Theorem~\ref{thm:expansion-walks}, by stages~\ref{it:order-choice-4} and~\ref{it:order-choice-9} and by stage~\ref{it:order-choice-6} together with condition~\eqref{eq:renormalisation-group-19}. Hence Theorem~\ref{thm:expansion-walks} applies to the data $\left(j,l_j,\bar{m}\cdot a_j,e_j\right)$. If $j\ge1$, its part~(a) with $y=\infty$, for which $F_{\infty}(A)=D_A+\mathfrak{m}_j+b_j\cdot P_j(A)$ (Section~\ref{sec:operators-2}), shows that $S_j(A)$ exists for every $A\in\tilde{G}_j$. If $j<k$, Proposition~\ref{prop:polymer-expansion-determinant} applies at level $j$ to every family~\eqref{it:primitives-p5} on $T_j$ with the properties named in it, and clause~\eqref{it:polymer-expansion-determinant-a} of this proposition holds for every such family. Thus hypothesis~\ref{def:hypothesis-flow} and the inequality $k\le k^{*}(R)$ imply the hypothesis of the theorem on the operators $S_j$. The hypotheses of Proposition~\ref{prop:free-determinant} are hypothesis~\ref{def:hypothesis-flow}, the inequality $k\le k^{*}(R)$ and families~\eqref{it:primitives-p5} of the levels $j<k$ for which clause~\eqref{it:local-expansion-determinant-d0} of Proposition~\ref{prop:local-expansion-determinant} holds (Theorem~\ref{thm:renormalisation-group-1}(d)), together with $a\in C^1\otimes\mathbb{C}$ and $a^{[j]}\in\mathfrak{X}_j$ for $0\le j\le k$; they therefore imply all hypotheses of the theorem, whose two identities are the assertions of Proposition~\ref{prop:free-determinant}. Proposition~\ref{prop:free-determinant} follows.
\end{proof}

\section{Algebraic and Gaussian identities for a renormalisation group step}\label{sec:algebraic-gaussian-identities}

This section contains the proofs of the propositions of Section~\ref{sec:identities-exact-operators}; the directions and the norms are those of Section~\ref{sec:renormalisation-group-scheme}, the stages are those of Definition~\ref{def:order-choice}, and lattices, cochains and blocks are given in Section~\ref{sec:preliminaries} a notation of their own, for an arbitrary odd integer $L\ge3$. Block averages, constrained minimisers, axial gauges and the basis of the fluctuation space are also treated in Section~\ref{sec:operators-photon-field}, there in another notation and with a view to bounds, and the resolvent identities of the fermionic block operators in Section~\ref{sec:operators-fermion-field}.

\subsection{Preliminaries}\label{sec:preliminaries}

\subsubsection{Tori, cochains and the Hodge decomposition}\label{sec:tori-cochains-hodge}

Let $L\ge 3$ be an odd integer. The directions are $\mu=0,1,2$, and $e_{\mu}$ is the unit vector of $\mathbb{R}^3$ in the direction $\mu$. For $x\in\mathbb{R}^3$, $|x|:=\max_{\mu}|x_{\mu}|$ is the sup norm (Definition~\ref{def:sup-norm-distance}). For $a>0$ and $l\in a\mathbb{N}$, where $\mathbb{N}=\{1,2,3,\dots\}$, put
\[
\mathbb{L}(a):=a\mathbb{Z}^3,\qquad \mathbb{T}(a,l):=a\mathbb{Z}^3/l\mathbb{Z}^3.
\]
A function on the cells of $\mathbb{T}(a,l)$, that is, on its sites, oriented bonds and oriented plaquettes, is identified with an $l$-periodic function on the cells of $\mathbb{L}(a)$. The oriented bond from a site $x$ to a neighbouring site $x'$ is written $(x,x')$, and $[x_0,x_1,x_2,x_3]$ denotes the plaquette with the corners $x_0,x_1,x_2,x_3$, oriented in this cyclic order. For $p=0,1,2$ let $C^p(a,l)$ be the space of the real functions on the sites, on the oriented bonds, on the oriented plaquettes of $\mathbb{T}(a,l)$, respectively, which for $p=1,2$ change sign when the orientation is reversed. The components of $A\in C^1(a,l)$ and of $F\in C^2(a,l)$ at a site $x$ are
\[
A_{\mu}(x):=A(x,x+ae_{\mu}),\qquad F_{\mu\nu}(x):=F([x,x+ae_{\mu},x+ae_{\mu}+ae_{\nu},x+ae_{\nu}])\qquad(\mu\ne\nu),
\]
so that $F_{\nu\mu}=-F_{\mu\nu}$. For $\boldsymbol{\lambda}\in C^0(a,l)$, $A\in C^1(a,l)$ and $f,g\in C^p(a,l)$ put
\begin{align}
(d\boldsymbol{\lambda})_{\mu}(x)&:=a^{-1}\left(\boldsymbol{\lambda}(x+ae_{\mu})-\boldsymbol{\lambda}(x)\right),\notag\\
(dA)_{\mu\nu}(x)&:=a^{-1}\left(A_{\mu}(x)+A_{\nu}(x+ae_{\mu})-A_{\mu}(x+ae_{\nu})-A_{\nu}(x)\right),\label{eq:algebraic-gaussian-1}\\
\langle f,g\rangle&:=\begin{cases}\sum_xa^3f(x)\cdot g(x)&(p=0),\\ \sum_x\sum_{\mu}a^3f_{\mu}(x)\cdot g_{\mu}(x)&(p=1),\\ \sum_x\sum_{\mu<\nu}a^3f_{\mu\nu}(x)\cdot g_{\mu\nu}(x)&(p=2),\end{cases}\notag\\
|f|_2^2&:=\langle f,f\rangle,\qquad |f|_{\infty}:=\sup|f|,\qquad \delta:=d^T,\qquad \operatorname{Lap}:=\delta d\ \text{ on } C^0(a,l),\notag
\end{align}
where $x$ ranges over the sites of $\mathbb{T}(a,l)$ and the transpose is taken with respect to these pairings. The second formula is antisymmetric in $\mu,\nu$; thus $d$ maps $C^0(a,l)$ into $C^1(a,l)$ and $C^1(a,l)$ into $C^2(a,l)$, and $\delta$ maps $C^2(a,l)$ into $C^1(a,l)$ and $C^1(a,l)$ into $C^0(a,l)$. The coboundary of a bond function, with the above convention for the components of a plaquette function, is that of \cite[(9)]{DimockUR}; the weight $a^3$ in the pairings and the norm $|\cdot|_2$ on $C^2(a,l)$, a sum over the pairs $\mu<\nu$ and the sites $x$, are those of \cite[(11)]{DimockUR}. For $a=1$ these pairings and norms are the counting ones.

A path from $x$ to $x'$ is a finite sequence of oriented bonds $(x_0,x_1),(x_1,x_2),\dots,(x_{m-1},x_m)$ with $x_0=x$ and $x_m=x'$; it is closed if $x'=x$. For a chain $\boldsymbol{\Gamma}$ of oriented bonds, for instance a path, and for $A\in C^1(a,l)$ put
\[
A(\boldsymbol{\Gamma}):=\sum_{b\in\boldsymbol{\Gamma}}A(b);
\]
no weight is attached to the bonds. The boundary $\partial p$ of the oriented plaquette $p=[x_0,x_1,x_2,x_3]$ is the closed path $(x_0,x_1),(x_1,x_2),(x_2,x_3),(x_3,x_0)$. For $\boldsymbol{\lambda}\in C^0(a,l)$ and $A\in C^1(a,l)$ one has
\begin{align}
(d\boldsymbol{\lambda})(\boldsymbol{\Gamma})&=a^{-1}\left(\boldsymbol{\lambda}(x')-\boldsymbol{\lambda}(x)\right)\qquad\text{for every path }\boldsymbol{\Gamma}\text{ from }x\text{ to }x',\tag{S1}\label{eq:algebraic-gaussian-S1}\\
A(\partial p)&=a\cdot(dA)(p)\qquad\text{for every oriented plaquette }p.\tag{S2}\label{eq:algebraic-gaussian-S2}
\end{align}
Indeed, by \eqref{eq:algebraic-gaussian-1} the value of $d\boldsymbol{\lambda}$ on a bond $(x_i,x_{i+1})$ is $a^{-1}(\boldsymbol{\lambda}(x_{i+1})-\boldsymbol{\lambda}(x_i))$, also when $x_{i+1}-x_i=-ae_{\mu}$, and the sum over the bonds of $\boldsymbol{\Gamma}$ telescopes. Every oriented plaquette has the form $p=[x,x+ae_{\mu},x+ae_{\mu}+ae_{\nu},x+ae_{\nu}]$ with a site $x$ and $\mu\ne\nu$; for this $p$ the sum $A(\partial p)$ equals $A_{\mu}(x)+A_{\nu}(x+ae_{\mu})-A_{\mu}(x+ae_{\nu})-A_{\nu}(x)$, which is $a\cdot(dA)_{\mu\nu}(x)=a\cdot(dA)(p)$.

A bond function $c\in C^1(a,l)$ is constant if there are real numbers $c_{\mu}$ with $c_{\mu}(x)=c_{\mu}$ for all sites $x$ and all $\mu$; it is identified with the vector $(c_{\mu})_{\mu}\in\mathbb{R}^3$, so that the constant bond functions form a space of dimension $3$. The mean of $A\in C^1(a,l)$ is the vector $\operatorname{mean}(A)\in\mathbb{R}^3$ with the components
\[
\operatorname{mean}(A)_{\mu}:=\left(\frac{a}{l}\right)^3\sum_xA_{\mu}(x),
\]
the sum extending over the $\left(\frac{l}{a}\right)^3$ sites of $\mathbb{T}(a,l)$. By periodicity $\operatorname{mean}(d\boldsymbol{\lambda})=0$ for every $\boldsymbol{\lambda}\in C^0(a,l)$, and $\operatorname{mean}(c)=c$ for every constant $c$. Since $\langle A,c\rangle=l^3\sum_{\mu}\operatorname{mean}(A)_{\mu}\cdot c_{\mu}$, a bond function has mean zero if and only if it is orthogonal to the constants. A bond function $A$ is called coclosed if $\delta A=0$.

The coboundary satisfies $dd=0$ on $C^0(a,l)$: for an oriented plaquette $p$, \eqref{eq:algebraic-gaussian-S2} and \eqref{eq:algebraic-gaussian-S1} give $a\cdot(dd\boldsymbol{\lambda})(p)=(d\boldsymbol{\lambda})(\partial p)=0$, as $\partial p$ is closed. Hence $dC^0(a,l)$ and $\delta C^2(a,l)$ are orthogonal subspaces of $C^1(a,l)$, and the orthogonal complement of their sum is the intersection of the kernels of $\delta$ and of $d$ on $C^1(a,l)$, which equals $\ker(d\delta+\delta d)$, because $\langle A,(d\delta+\delta d)A\rangle=|\delta A|_2^2+|dA|_2^2$. This is the Hodge decomposition
\[
C^1(a,l)=dC^0(a,l)\oplus\delta C^2(a,l)\oplus\ker(d\delta+\delta d)\qquad\text{(orthogonal sum)}
\]
of the finite cochain complex $C^0(a,l)\to C^1(a,l)\to C^2(a,l)$ (B.~Eckmann, \emph{Harmonische Funktionen und Randwertaufgaben in einem Komplex}, Comment.\ Math.\ Helv.\ 17 (1944/45), 240--255). On the torus the difference operators of the different directions commute, so that $d\delta+\delta d$ acts on each component $A_{\mu}$ as $\operatorname{Lap}$. If $\operatorname{Lap}\boldsymbol{\lambda}=0$, then $|d\boldsymbol{\lambda}|_2^2=\langle\boldsymbol{\lambda},\operatorname{Lap}\boldsymbol{\lambda}\rangle=0$, and $\boldsymbol{\lambda}$ is constant by \eqref{eq:algebraic-gaussian-S1}, since any two sites are joined by a path. Consequently $\ker(d\delta+\delta d)$ consists of the constant bond functions, and
\begin{gather}
C^1(a,l)=dC^0(a,l)\oplus\{\text{constants}\}\oplus\delta C^2(a,l)\qquad\text{(orthogonal sum)},\notag\\
\ker d=dC^0(a,l)\oplus\{\text{constants}\}\ \text{ on } C^1(a,l),\qquad \ker d=\{\text{constants}\}\ \text{ on } C^0(a,l),\tag{Hodge}\label{eq:algebraic-gaussian-Hodge}\\
\delta C^1(a,l)=\operatorname{Lap} C^0(a,l)=\left\{\boldsymbol{\lambda}\in C^0(a,l):\sum_x\boldsymbol{\lambda}(x)=0\right\}.\notag
\end{gather}
The assertions on the kernels hold because the kernel of $d$ on $C^1(a,l)$ is the orthogonal complement of $\delta C^2(a,l)$, while on $C^0(a,l)$ it coincides with the kernel of $\operatorname{Lap}$, as $\langle\boldsymbol{\lambda},\operatorname{Lap}\boldsymbol{\lambda}\rangle=|d\boldsymbol{\lambda}|_2^2$. The range of $\delta$ on $C^1(a,l)$ and the range of the symmetric map $\operatorname{Lap}$ are both the orthogonal complement of the kernel of $d$ on $C^0(a,l)$, that is, of the constants.

Put
\[
V_{\mathrm{t}}:=\delta C^2(a,l).
\]
By \eqref{eq:algebraic-gaussian-Hodge}, $V_{\mathrm{t}}$ is the space of the coclosed bond functions with mean zero: the kernel of $\delta$ on $C^1(a,l)$ is the orthogonal complement of $dC^0(a,l)$, which is $\{\text{constants}\}\oplus V_{\mathrm{t}}$, and mean zero is orthogonality to the constants. The map $d$ is injective on $V_{\mathrm{t}}$, since $V_{\mathrm{t}}$ is orthogonal to the kernel $dC^0(a,l)\oplus\{\text{constants}\}$ of $d$ on $C^1(a,l)$, and $dV_{\mathrm{t}}=dC^1(a,l)$. As $\mathbb{T}(a,l)$ has $\left(\frac{l}{a}\right)^3$ sites, $\dim C^1(a,l)=3\left(\frac{l}{a}\right)^3$ and $\dim dC^0(a,l)=\left(\frac{l}{a}\right)^3-1$; therefore
\[
\dim dC^1(a,l)=\dim V_{\mathrm{t}}=3\left(\frac{l}{a}\right)^3-\left(\left(\frac{l}{a}\right)^3-1\right)-3=2\left(\frac{l}{a}\right)^3-2.
\]

\subsubsection{Blocks, levels and rectilinear paths}\label{sec:blocks-levels-rectilinear}

For $j\ge 0$ and $y\in\mathbb{L}(L^ja)$ put
\[
B_j(y):=\left\{x\in\mathbb{L}(a):|x-y|<\frac{L^ja}{2}\right\},\qquad B(y):=B_1(y).
\]
Since $L^j$ is odd, $B_j(y)=y+a\left\{-\frac{L^j-1}{2},\dots,\frac{L^j-1}{2}\right\}^3$. Hence $\#|B_j(y)|=L^{3j}$, the blocks $B_j(y)$ with $y\in\mathbb{L}(L^ja)$ form a partition of $\mathbb{L}(a)$, and for every $z\in\mathbb{L}(L^{j+1}a)$
\begin{equation}\label{eq:algebraic-gaussian-2}
B_{j+1}(z)=\bigsqcup_{y\in B(z)}B_j(y)\qquad\text{(disjoint union)},
\end{equation}
where $B(z):=\left\{y\in\mathbb{L}(L^ja):|y-z|<\frac{L^{j+1}a}{2}\right\}$ is the block of $z$ in the lattice $\mathbb{L}(L^ja)$. To verify \eqref{eq:algebraic-gaussian-2}, write a site $x\in\mathbb{L}(a)$ as $x=z+a(L^jq+r)$ with $q,r\in\mathbb{Z}^3$ and $|r|\le\frac{L^j-1}{2}$. Then $x\in B_j(z+L^jaq)$, and, coordinate by coordinate, $|x-z|<\frac{L^{j+1}a}{2}$ holds if and only if $|q|\le\frac{L-1}{2}$, that is, if and only if $z+L^jaq\in B(z)$.

For a level $k\ge 0$ and a side $l\in\mathbb{N}$ (Section~\ref{sec:lattice-geometry}) put $\eta:=\eta_k:=L^{-k}$,
\[
T_k:=\mathbb{T}(\eta_k,l),\qquad T^0_k:=\mathbb{T}(1,l),\qquad T^1_k:=\mathbb{T}(L,l)\ \text{ if } l\in L\mathbb{N},\qquad n_k:=l^3,
\]
and write $C^p(T_k)$, $C^p(T^0_k)$, $C^p(T^1_k)$ for $C^p(\eta_k,l)$, $C^p(1,l)$, $C^p(L,l)$. The sites of $T_k$ are called fine sites, those of $T^0_k$ unit sites. For $y\in\mathbb{Z}^3$, $B_k(y)$ is the block of fine sites
\[
B_k(y)=\left\{x\in\eta_k\mathbb{Z}^3:|x-y|<\frac{1}{2}\right\},
\]
that is, the block $B_j(y)$ with $a=\eta_k$ and $j=k$. In this case the partition property above says that each fine site $x$ lies in a single block $B_k(y)$; its centre $y$, a unit site, is denoted by $[x]$, so that $x\in B_k([x])$.

Let $x,x'\in\mathbb{L}(a)$ with $x'-x\in a\mathbb{Z}e_{\mu}$ for some direction $\mu$. The straight path $\boldsymbol{\Gamma}_{x,x'}$ from $x$ to $x'$ consists of the bonds of $\mathbb{L}(a)$ on the segment from $x$ to $x'$, each oriented from $x$ towards $x'$; for instance, $\boldsymbol{\Gamma}_{x,x+Lae_{\mu}}$ consists of the $L$ bonds $(x+iae_{\mu},x+(i+1)ae_{\mu})$, $0\le i\le L-1$. Let now $x,y\in\mathbb{L}(a)$, let $\pi=(\pi_1,\pi_2,\pi_3)$ be a permutation of the directions, and put
\[
y':=y+(x_{\pi_1}-y_{\pi_1})e_{\pi_1},\qquad y'':=y'+(x_{\pi_2}-y_{\pi_2})e_{\pi_2}.
\]
The differences $y'-y$, $y''-y'$, $x-y''$ lie in $a\mathbb{Z}e_{\pi_1}$, $a\mathbb{Z}e_{\pi_2}$, $a\mathbb{Z}e_{\pi_3}$, respectively. The rectilinear path $\boldsymbol{\Gamma}^{\pi}(y,x)$ from $y$ to $x$ is the straight path $\boldsymbol{\Gamma}_{y,y'}$, followed by $\boldsymbol{\Gamma}_{y',y''}$ and then by $\boldsymbol{\Gamma}_{y'',x}$; along it the coordinates are changed one at a time, in the order $\pi_1,\pi_2,\pi_3$, each from its value at $y$ to its value at $x$. If $y\in\mathbb{L}(La)$ and $x\in B(y)$, then all sites of $\boldsymbol{\Gamma}^{\pi}(y,x)$ lie in $B(y)$, since each coordinate of each of these sites lies between the corresponding coordinates of $y$ and of $x$.

\subsubsection{Scaling}\label{sec:scaling}

Let $l\in La\mathbb{N}$. The map $x\mapsto\frac{x}{L}$ carries the cells of $\mathbb{T}(La,l)$ bijectively onto those of $\mathbb{T}\left(a,\frac{l}{L}\right)$. For functions $\boldsymbol{\lambda}$, $A$, $F$ on the sites, on the bonds, on the plaquettes of $\mathbb{T}\left(a,\frac{l}{L}\right)$, functions on the cells of $\mathbb{T}(La,l)$ are defined by
\begin{equation}\label{eq:algebraic-gaussian-3}
\boldsymbol{\lambda}_L(x):=L^{\frac{1}{2}}\boldsymbol{\lambda}\left(\frac{x}{L}\right),\qquad A_L(b):=L^{-\frac{1}{2}}A\left(\frac{b}{L}\right),\qquad F_L(p):=L^{-\frac{3}{2}}F\left(\frac{p}{L}\right).
\end{equation}
For $\boldsymbol{\lambda}\in C^0\left(a,\frac{l}{L}\right)$, $A\in C^1\left(a,\frac{l}{L}\right)$ and $F\in C^2\left(a,\frac{l}{L}\right)$ it follows from \eqref{eq:algebraic-gaussian-1} that
\[
d(\boldsymbol{\lambda}_L)=(d\boldsymbol{\lambda})_L,\qquad d(A_L)=(dA)_L,\qquad |F_L|_2=|F|_2,\qquad \operatorname{mean}(A_L)=L^{-\frac{1}{2}}\cdot\operatorname{mean}(A).
\]
Indeed, $(A_L)_{\mu}(x)=L^{-\frac{1}{2}}A_{\mu}\left(\frac{x}{L}\right)$ and $(F_L)_{\mu\nu}(x)=L^{-\frac{3}{2}}F_{\mu\nu}\left(\frac{x}{L}\right)$, while for a site $x$ of $\mathbb{T}(La,l)$
\begin{align*}
(d(\boldsymbol{\lambda}_L))_{\mu}(x)&=(La)^{-1}L^{\frac{1}{2}}\left(\boldsymbol{\lambda}\left(\frac{x}{L}+ae_{\mu}\right)-\boldsymbol{\lambda}\left(\frac{x}{L}\right)\right)=L^{-\frac{1}{2}}\cdot(d\boldsymbol{\lambda})_{\mu}\left(\frac{x}{L}\right),\\
(d(A_L))_{\mu\nu}(x)&=(La)^{-1}L^{-\frac{1}{2}}\cdot a\cdot(dA)_{\mu\nu}\left(\frac{x}{L}\right)=L^{-\frac{3}{2}}\cdot(dA)_{\mu\nu}\left(\frac{x}{L}\right).
\end{align*}
In the pairing on $\mathbb{T}(La,l)$ every cell carries the weight $(La)^3$, and $(La)^3\cdot L^{-3}=a^3$ is the weight on $\mathbb{T}\left(a,\frac{l}{L}\right)$; this gives $|F_L|_2=|F|_2$. Both tori have $\left(\frac{l}{La}\right)^3$ sites, so that both means carry the factor $\left(\frac{La}{l}\right)^3$, and $\operatorname{mean}(A_L)=L^{-\frac{1}{2}}\cdot\operatorname{mean}(A)$ follows from the formula for $(A_L)_{\mu}$.

For $j\ge 0$ and $l\in L^ja\mathbb{N}$, $(\cdot)_{L^j}$ denotes the $j$-fold scaling, which takes functions on the cells of $\mathbb{T}\left(a,\frac{l}{L^j}\right)$ to functions on the cells of $\mathbb{T}(L^ja,l)$; on bond functions $A_{L^j}(b)=L^{-\frac{j}{2}}A\left(\frac{b}{L^j}\right)$. Its inverse on bond functions is $B\mapsto B^{[j]}$, where for a bond function $B$ on $\mathbb{T}(L^ja,l)$ and a bond $b$ of $\mathbb{T}\left(a,\frac{l}{L^j}\right)$
\[
B^{[j]}(b):=L^{\frac{j}{2}}B(L^jb)
\]
(Section~\ref{sec:lattice-measures-laws}).

\subsubsection{Quadratic minimisation and block matrices}\label{sec:quadratic-minimisation-block}

\begin{lemma}[Minimisation on an affine subspace]
Let $T:E\to E'$ be a linear map between finite-dimensional real Euclidean spaces, let $W$ be a linear subspace of $E$, and let $x_0\in E$. The restriction of the function $x\mapsto|Tx|^2$ to the affine subspace $x_0+W$ has a smallest value; the points of $x_0+W$ at which it is attained are called minimisers. They are exactly the points $x_{*}\in x_0+W$ with $Tx_{*}\perp T(W)$, that is, with $\langle Tx_{*},Tw\rangle=0$ for every $w\in W$. If $x_{*}$ is a minimiser, then the set of all minimisers is $x_{*}+(W\cap\ker T)$, and for every $w\in W$
\begin{equation}\tag{LA1}\label{eq:algebraic-gaussian-LA1}
|T(x_{*}+w)|^2=|Tx_{*}|^2+|Tw|^2.
\end{equation}
\end{lemma}

\begin{proof}
Let $\Pi$ be the orthogonal projection of $E'$ onto the subspace $T(W)$. As $\Pi Tx_0\in T(W)$, there is $w_0\in W$ with $Tw_0=-\Pi Tx_0$; put $\hat{x}:=x_0+w_0$. Then $T\hat{x}=(1-\Pi)Tx_0$ is orthogonal to $T(W)$, so that
\[
|T(\hat{x}+w)|^2=|T\hat{x}|^2+|Tw|^2\qquad(w\in W).
\]
Since $x_0+W=\hat{x}+W$, it follows that $\hat{x}$ is a minimiser, that $\hat{x}+w$ with $w\in W$ is a minimiser if and only if $Tw=0$, and hence that the set of the minimisers is $\hat{x}+(W\cap\ker T)$. In particular every minimiser $x$ satisfies $Tx=T\hat{x}$, which is orthogonal to $T(W)$. Conversely, let $x=\hat{x}+w$ with $w\in W$ be such that $Tx$ is orthogonal to $T(W)$; then $0=\langle Tx,Tw\rangle=\langle T\hat{x},Tw\rangle+|Tw|^2=|Tw|^2$, so that $x$ is a minimiser. Finally, for a minimiser $x_{*}$ one has $x_{*}+(W\cap\ker T)=\hat{x}+(W\cap\ker T)$, and $Tx_{*}=T\hat{x}$ is orthogonal to $T(W)$, which gives \eqref{eq:algebraic-gaussian-LA1}.
\end{proof}

\begin{lemma}[Schur complements]
Let $U$, $V$ be finite-dimensional complex vector spaces, and let $\mathfrak{M}$ be a linear map of $U\oplus V$ with the blocks $a:U\to U$, $b:V\to U$, $c:U\to V$, $d:V\to V$, that is,
\begin{equation}\tag{LA2}\label{eq:algebraic-gaussian-LA2}
\mathfrak{M}(u,v)=(au+bv,\,cu+dv)\qquad(u\in U,\ v\in V).
\end{equation}
If $d$ is invertible, put $s:=a-bd^{-1}c$. Then $\det\mathfrak{M}=\det d\cdot\det s$, the map $\mathfrak{M}$ is invertible if and only if $s$ is, and in this case
\[
\mathfrak{M}^{-1}(u,v)=\left(s^{-1}u-s^{-1}bd^{-1}v,\ -d^{-1}cs^{-1}u+\left(d^{-1}+d^{-1}cs^{-1}bd^{-1}\right)v\right).
\]
If $a$ is invertible, put $t:=d-ca^{-1}b$. Then $\det\mathfrak{M}=\det a\cdot\det t$, the map $\mathfrak{M}$ is invertible if and only if $t$ is, and in this case
\[
\mathfrak{M}^{-1}(u,v)=\left(\left(a^{-1}+a^{-1}bt^{-1}ca^{-1}\right)u-a^{-1}bt^{-1}v,\ -t^{-1}ca^{-1}u+t^{-1}v\right).
\]
\end{lemma}

\begin{proof}
Let $d$ be invertible, and define linear maps of $U\oplus V$ by
\[
G_1(u,v):=(u+bd^{-1}v,\,v),\qquad G_2(u,v):=(u,\,d^{-1}cu+v),\qquad (s\oplus d)(u,v):=(su,\,dv).
\]
Then $G_1(s\oplus d)G_2(u,v)=G_1(su,\,cu+dv)=(su+bd^{-1}cu+bv,\,cu+dv)=\mathfrak{M}(u,v)$, because $s+bd^{-1}c=a$. The maps $G_1-1$ and $G_2-1$ have square zero; hence $G_1$ and $G_2$ have determinant $1$ and are invertible, with
\[
G_1^{-1}(u,v)=(u-bd^{-1}v,\,v),\qquad G_2^{-1}(u,v)=(u,\,-d^{-1}cu+v).
\]
It follows that $\det\mathfrak{M}=\det(s\oplus d)=\det s\cdot\det d$, that $\mathfrak{M}$ is invertible if and only if $s\oplus d$ is, that is, if and only if $s$ is, and that in this case
\[
\mathfrak{M}^{-1}(u,v)=G_2^{-1}\left(s^{-1}\oplus d^{-1}\right)G_1^{-1}(u,v)=G_2^{-1}\left(s^{-1}u-s^{-1}bd^{-1}v,\,d^{-1}v\right),
\]
which is the first formula for $\mathfrak{M}^{-1}$. Let now $a$ be invertible, and define
\[
G_3(u,v):=(u,\,ca^{-1}u+v),\qquad G_4(u,v):=(u+a^{-1}bv,\,v),\qquad (a\oplus t)(u,v):=(au,\,tv).
\]
Then $G_3(a\oplus t)G_4(u,v)=G_3(au+bv,\,tv)=(au+bv,\,cu+ca^{-1}bv+tv)=\mathfrak{M}(u,v)$, because $ca^{-1}b+t=d$. The maps $G_3-1$ and $G_4-1$ have square zero; hence $G_3$ and $G_4$ have determinant $1$ and are invertible, with
\[
G_3^{-1}(u,v)=(u,\,-ca^{-1}u+v),\qquad G_4^{-1}(u,v)=(u-a^{-1}bv,\,v).
\]
It follows that $\det\mathfrak{M}=\det a\cdot\det t$, that $\mathfrak{M}$ is invertible if and only if $t$ is, and that in this case
\[
\mathfrak{M}^{-1}(u,v)=G_4^{-1}\left(a^{-1}\oplus t^{-1}\right)G_3^{-1}(u,v)=G_4^{-1}\left(a^{-1}u,\,-t^{-1}ca^{-1}u+t^{-1}v\right),
\]
which is the second formula for $\mathfrak{M}^{-1}$.
\end{proof}

\begin{lemma}[Sylvester's determinant identity]
Let $U$, $V$ be finite-dimensional complex vector spaces, and let $X:V\to U$ and $Y:U\to V$ be linear maps. Then
\begin{equation}\tag{LA3}\label{eq:algebraic-gaussian-LA3}
\det(1_U-XY)=\det(1_V-YX).
\end{equation}
\end{lemma}

\begin{proof}
Apply \eqref{eq:algebraic-gaussian-LA2} to $\mathfrak{M}(u,v):=(u+Xv,\,Yu+v)$, that is, to the blocks $1_U$, $X$, $Y$, $1_V$. The lower right block $1_V$ is invertible, with $s=1_U-XY$, and the upper left block $1_U$ is invertible, with $t=1_V-YX$. Hence $\det(1_U-XY)=\det\mathfrak{M}=\det(1_V-YX)$.
\end{proof}

\subsection{Block averages}\label{sec:block-averages}

\subsubsection{The averaging operators}\label{sec:averaging-operators}

Let $a>0$ and $l\in La\mathbb{N}$. For $\boldsymbol{\lambda}\in C^0(a,l)$, $A\in C^1(a,l)$, $F\in C^2(a,l)$, $y\in\mathbb{L}(La)$ and directions $\mu\ne\nu$ put
\begin{align}
(Q\boldsymbol{\lambda})(y)&:=L^{-3}\sum_{x\in B(y)}\boldsymbol{\lambda}(x),\notag\\
(QA)_{\mu}(y)&:=L^{-4}\sum_{x\in B(y)}\sum_{i=0}^{L-1}A_{\mu}(x+iae_{\mu})=L^{-4}\sum_{x\in B(y)}A\left(\boldsymbol{\Gamma}_{x,x+Lae_{\mu}}\right),\label{eq:algebraic-gaussian-4}\\
\left(Q^{(2)}F\right)_{\mu\nu}(y)&:=L^{-5}\sum_{x\in B(y)}\sum_{i,i'=0}^{L-1}F_{\mu\nu}(x+iae_{\mu}+i'ae_{\nu}).\notag
\end{align}
These expressions are $l$-periodic in $y$, and $\left(Q^{(2)}F\right)_{\nu\mu}=-\left(Q^{(2)}F\right)_{\mu\nu}$, as one sees by exchanging $i$ and $i'$. Hence \eqref{eq:algebraic-gaussian-4} defines, through values and components at the sites of $\mathbb{T}(La,l)$, a function $Q\boldsymbol{\lambda}$, a bond function $QA$ and a plaquette function $Q^{(2)}F$ on $\mathbb{T}(La,l)$, and thus linear maps
\[
Q:C^p(a,l)\to C^p(La,l)\quad(p=0,1),\qquad Q^{(2)}:C^2(a,l)\to C^2(La,l).
\]
On bond functions $Q$ is the average of \cite[(161)]{DimockUR} and \cite[(189)]{DimockMB}, and $Q^{(2)}$ is the average of \cite[(194), (195)]{DimockMB}.

For $j\ge 0$ and $l\in L^ja\mathbb{N}$ let
\[
Q_j:C^p(a,l)\to C^p(L^ja,l)\quad(p=0,1),\qquad Q_j^{(2)}:C^2(a,l)\to C^2(L^ja,l)
\]
be the $j$-fold compositions $Q\circ\dots\circ Q$ and $Q^{(2)}\circ\dots\circ Q^{(2)}$, in which the $m$-th factor from the right acts on the cochains on $\mathbb{T}(L^{m-1}a,l)$; $Q_0$ and $Q_0^{(2)}$ are the identity. At level $k$, with $a=\eta_k$ and $j=k$, so that $L^ja=1$, this gives
\[
Q_k:C^p(T_k)\to C^p(T^0_k)\quad(p=0,1),\qquad Q_k^{(2)}:C^2(T_k)\to C^2(T^0_k),
\]
and, for $l\in L\mathbb{N}$ and $a=1$, the map $Q:C^p(T^0_k)\to C^p(T^1_k)$ for $p=0,1$.

\subsubsection{Properties of the averages}\label{sec:properties-averages}

\begin{proposition}[Averaging operators]\label{prop:averaging-operators}
Let $j\ge0$ be an integer, let $a>0$, and let $l\in L^ja\mathbb{N}$.
\begin{enumerate}
\item[(a)] For $\boldsymbol{\lambda}\in C^0(a,l)$, $A\in C^1(a,l)$, $F\in C^2(a,l)$ and $y\in\mathbb{L}\left(L^ja\right)$ one has
\begin{align}
(Q_j\boldsymbol{\lambda})(y)&=L^{-3j}\sum_{x\in B_j(y)}\boldsymbol{\lambda}(x),\notag\\
(Q_jA)_{\mu}(y)&=L^{-4j}\sum_{x\in B_j(y)}\ \sum_{i=0}^{L^j-1}A_{\mu}(x+ia\,e_{\mu}),\notag\\
\left(Q_j^{(2)}F\right)_{\mu\nu}(y)&=L^{-5j}\sum_{x\in B_j(y)}\ \sum_{i,i'=0}^{L^j-1}F_{\mu\nu}(x+ia\,e_{\mu}+i'a\,e_{\nu}).\label{eq:algebraic-gaussian-5}
\end{align}
For $a=\eta_k$ and $j=k$ the formulas~\eqref{eq:algebraic-gaussian-5}, in which $L^{-3k}=\eta_k^3$, $L^{-4k}=\eta_k^3\cdot\eta_k$ and $L^{-5k}=\eta_k^3\cdot\eta_k^2$, are the formulas by which $Q_k$ and $Q_k^{(2)}$ are defined in Section~\ref{sec:averaging-fluctuations-gauge}; for $a=1$ the second line of~\eqref{eq:algebraic-gaussian-4} is the formula by which $Q$ is defined there.
\item[(b)] For $\boldsymbol{\lambda}\in C^0(a,l)$ and $A\in C^1(a,l)$,
\[
Q_j\,d\boldsymbol{\lambda}=d\,Q_j\boldsymbol{\lambda},\qquad d\,Q_jA=Q_j^{(2)}\,dA,
\]
each $d$ being the operator~\eqref{eq:algebraic-gaussian-1} of the torus on which its argument is defined.
\item[(c)] $Q_jc=c$ for every constant scalar function $c$ and for every constant bond function $c$ on $\mathbb{T}(a,l)$, and $\operatorname{mean}(Q_jA)=\operatorname{mean}(A)$ for every $A\in C^1(a,l)$.
\item[(d)] Write $T$ for $Q_j$ on $C^0(a,l)$, for $Q_j$ on $C^1(a,l)$ or for $Q_j^{(2)}$ on $C^2(a,l)$, the operators given by the three lines of~\eqref{eq:algebraic-gaussian-5}, and take $f$ in the domain of $T$. The cells of a torus are indexed as in the pairings of Section~\ref{sec:tori-cochains-hodge}: by the sites, by the pairs $(x,\mu)$, or by the triples $(x,\mu,\nu)$ with $\mu<\nu$, according to the kind of $f$. One has $|Tf|_{\infty}\le|f|_{\infty}$. If $Y$ is any set of cells of $\mathbb{T}\left(L^ja,l\right)$, and if $N(Y)$ denotes the set of all cells $z$ of $\mathbb{T}(a,l)$ at which $f$ is evaluated in the sum that~\eqref{eq:algebraic-gaussian-5} gives for $(Tf)(y)$, for at least one $y\in Y$, then
\begin{equation}\label{eq:algebraic-gaussian-6}
\sum_{y\in Y}\left(L^ja\right)^3\cdot\left|(Tf)(y)\right|^2\le\sum_{z\in N(Y)}a^3\cdot\left|f(z)\right|^2.
\end{equation}
In particular, for $\boldsymbol{\lambda}\in C^0(a,l)$, $A\in C^1(a,l)$ and $F\in C^2(a,l)$,
\[
|Q_j\boldsymbol{\lambda}|_2\le|\boldsymbol{\lambda}|_2,\qquad|Q_jA|_2\le|A|_2,\qquad\left|Q_j^{(2)}F\right|_2\le|F|_2.
\]
At level $k$, that is, for $a=\eta_k$ and $j=k$, the last bound reads $\left|Q_k^{(2)}F\right|_{\ell^2}\le C_Q\cdot|F|_2$ with $C_Q=1$, where $|\cdot|_{\ell^2}$ is the counting norm on $T^0_k$.
\item[(e)] At level $k$, that is, for $a=\eta=\eta_k$ and $j=k$, the transpose $Q_k^T\colon C^1(T^0_k)\to C^1(T_k)$ of $Q_k$ with respect to the pairings of Section~\ref{sec:tori-cochains-hodge} is
\[
\left(Q_k^TA'\right)_{\mu}(z)=L^{-k}\sum_{i=0}^{L^k-1}A'_{\mu}\left([z-i\eta e_{\mu}]\right)\qquad\left(A'\in C^1(T^0_k),\ z\in T_k\right),
\]
where $[x]\in T^0_k$ denotes the unit site with $x\in B_k([x])$. Moreover $\left|Q_k^TA'\right|_{\infty}\le|A'|_{\infty}$.
\item[(f)] For $A'\in C^1\left(L^ja,l\right)$ let $\boldsymbol{\Sigma}_jA'\in C^1(a,l)$ be the bond function with the values
\[
(\boldsymbol{\Sigma}_jA')_{\mu}(z):=\begin{cases}L^j\cdot A'_{\mu}(y)&\text{if }z+a\,e_{\mu}\in B_j\left(y+L^ja\,e_{\mu}\right),\\0&\text{otherwise}\end{cases}
\]
on the bonds $(z,z+a\,e_{\mu})$, $z\in\mathbb{L}(a)$, where $y\in\mathbb{L}\left(L^ja\right)$ is the point with $z\in B_j(y)$. Then $Q_j\boldsymbol{\Sigma}_jA'=A'$ for every $A'\in C^1\left(L^ja,l\right)$; in particular $Q_j$ maps $C^1(a,l)$ onto $C^1\left(L^ja,l\right)$.
\item[(g)] Assume $l\in L^2a\mathbb{N}$. Let $f$ be a scalar function or a bond function and $F$ a plaquette function on the cells of $\mathbb{T}(a,l/L)$, and let $f_L$, $F_L$ be their scalings~\eqref{eq:algebraic-gaussian-3}, which are functions on the cells of $\mathbb{T}(La,l)$. Then
\[
Q(f_L)=(Qf)_L,\qquad Q^{(2)}(F_L)=\left(Q^{(2)}F\right)_L,
\]
both sides being functions on the cells of $\mathbb{T}\left(L^2a,l\right)$.
\end{enumerate}
\end{proposition}

\begin{proof}
(a) The formulas are proved by induction on $j$. For $j=0$ one has $Q_0=1$, $Q_0^{(2)}=1$ and $B_0(y)=\{y\}$, so that~\eqref{eq:algebraic-gaussian-5} holds. Let~\eqref{eq:algebraic-gaussian-5} hold for some $j\ge0$, every $a>0$ and every $l\in L^ja\mathbb{N}$. Let $l\in L^{j+1}a\mathbb{N}$, put $a_j:=L^ja$, let $z\in\mathbb{L}(a_{j+1})$, and let
\[
B(z):=\left\{y\in\mathbb{L}(a_j):|y-z|<\frac{La_j}{2}\right\}
\]
be the block of $z$ in $\mathbb{L}(a_j)$. The definition of the blocks gives $B_j(y+v)=B_j(y)+v$ for $y,v\in\mathbb{L}(a_j)$. Since $Q_{j+1}=Q\circ Q_j$, the second line of~\eqref{eq:algebraic-gaussian-4}, applied on $\mathbb{T}(a_j,l)$, and the induction hypothesis give
\begin{align*}
(Q_{j+1}A)_{\mu}(z)&=L^{-4}\sum_{y\in B(z)}\ \sum_{m=0}^{L-1}(Q_jA)_{\mu}(y+ma_je_{\mu})\\
&=L^{-4(j+1)}\sum_{y\in B(z)}\ \sum_{m=0}^{L-1}\ \sum_{x\in B_j(y)}\ \sum_{i=0}^{L^j-1}A_{\mu}\left(x+\left(mL^j+i\right)a\,e_{\mu}\right);
\end{align*}
in the second step a site of $B_j(y+ma_je_{\mu})$ is written as $x+ma_je_{\mu}$ with $x\in B_j(y)$, and $a_j=L^ja$ is used. As $(m,i)$ ranges over $\{0,\dots,L-1\}\times\{0,\dots,L^j-1\}$, the integer $mL^j+i$ takes every value in $\{0,\dots,L^{j+1}-1\}$ exactly once. By~\eqref{eq:algebraic-gaussian-2}, as $y$ ranges over $B(z)$ and $x$ over $B_j(y)$, the site $x$ takes every value in $B_{j+1}(z)$ exactly once. Hence $(Q_{j+1}A)_{\mu}(z)$ is given by the second line of~\eqref{eq:algebraic-gaussian-5} with $j+1$ in place of $j$. For scalars and for plaquette functions, \eqref{eq:algebraic-gaussian-4} and the induction hypothesis give
\begin{align*}
(Q_{j+1}\boldsymbol{\lambda})(z)&=L^{-3}\sum_{y\in B(z)}(Q_j\boldsymbol{\lambda})(y)=L^{-3(j+1)}\sum_{y\in B(z)}\ \sum_{x\in B_j(y)}\boldsymbol{\lambda}(x),\\
\left(Q_{j+1}^{(2)}F\right)_{\mu\nu}(z)&=L^{-5}\sum_{y\in B(z)}\ \sum_{m,m'=0}^{L-1}\left(Q_j^{(2)}F\right)_{\mu\nu}(y+ma_je_{\mu}+m'a_je_{\nu})\\
&=L^{-5(j+1)}\sum_{y\in B(z)}\ \sum_{m,m'=0}^{L-1}\ \sum_{x\in B_j(y)}\ \sum_{i,i'=0}^{L^j-1}F_{\mu\nu}\left(x+\left(mL^j+i\right)a\,e_{\mu}+\left(m'L^j+i'\right)a\,e_{\nu}\right).
\end{align*}
By~\eqref{eq:algebraic-gaussian-2} the first of these sums is the first line of~\eqref{eq:algebraic-gaussian-5} with $j+1$ in place of $j$; the two counting facts above, the first one applied to $mL^j+i$ and to $m'L^j+i'$, turn the second into the third line of~\eqref{eq:algebraic-gaussian-5} with $j+1$ in place of $j$.

Let $a=\eta_k$ and $j=k$. Then $L^ja=1$, and for $y\in\mathbb{Z}^3$ the set $B_k(y)=\{x\in\mathbb{L}(\eta_k):|x-y|<1/2\}$ is the block of fine sites of Section~\ref{sec:blocks-levels-rectilinear}. Inserting $L^{-3k}=\eta_k^3$, $L^{-4k}=\eta_k^3\cdot\eta_k$ and $L^{-5k}=\eta_k^3\cdot\eta_k^2$ in~\eqref{eq:algebraic-gaussian-5}, one obtains term by term the three formulas of Section~\ref{sec:averaging-fluctuations-gauge} that define $Q_k$ and $Q_k^{(2)}$. The second line of~\eqref{eq:algebraic-gaussian-4} with $a=1$, in the form $L^{-4}\sum_{x\in B(y)}A(\boldsymbol{\Gamma}_{x,x+Le_{\mu}})$, is the formula for $Q$ in Section~\ref{sec:averaging-fluctuations-gauge}.

(b) For $j=0$ there is nothing to prove. Suppose that the identities are proved for $j=1$, every $a>0$ and every $l\in La\mathbb{N}$, and that they hold for some $j\ge1$, every $a>0$ and every $l\in L^ja\mathbb{N}$. Let $l\in L^{j+1}a\mathbb{N}$. The identities for $j$, together with the case $j=1$ applied on $\mathbb{T}\left(L^ja,l\right)$, give
\[
Q_{j+1}\,d\boldsymbol{\lambda}=Q\left(Q_j\,d\boldsymbol{\lambda}\right)=Q\,d\left(Q_j\boldsymbol{\lambda}\right)=d\,Q\left(Q_j\boldsymbol{\lambda}\right)=d\,Q_{j+1}\boldsymbol{\lambda}
\]
and
\[
d\,Q_{j+1}A=d\,Q\left(Q_jA\right)=Q^{(2)}\,d\left(Q_jA\right)=Q^{(2)}Q_j^{(2)}\,dA=Q_{j+1}^{(2)}\,dA.
\]
By induction on $j$ it therefore suffices to prove the identities for $j=1$, where $Q_1=Q$ and $Q_1^{(2)}=Q^{(2)}$.

Let $y\in\mathbb{L}(La)$, and let $B(y)\subset\mathbb{L}(a)$ be its block. For $x\in\mathbb{L}(a)$, the straight path $\boldsymbol{\Gamma}_{x,x+La\,e_{\mu}}$ consists of the $L$ bonds $(x+ia\,e_{\mu},x+(i+1)a\,e_{\mu})$, $0\le i\le L-1$. By~\eqref{eq:algebraic-gaussian-S1},
\[
a\sum_{i=0}^{L-1}(d\boldsymbol{\lambda})_{\mu}(x+ia\,e_{\mu})=\boldsymbol{\lambda}(x+La\,e_{\mu})-\boldsymbol{\lambda}(x).
\]
Inserting this in~\eqref{eq:algebraic-gaussian-4} and using $B(y+La\,e_{\mu})=B(y)+La\,e_{\mu}$, one obtains
\begin{align*}
(Q\,d\boldsymbol{\lambda})_{\mu}(y)&=L^{-4}\cdot a^{-1}\sum_{x\in B(y)}\left(\boldsymbol{\lambda}(x+La\,e_{\mu})-\boldsymbol{\lambda}(x)\right)\\
&=(La)^{-1}\left((Q\boldsymbol{\lambda})(y+La\,e_{\mu})-(Q\boldsymbol{\lambda})(y)\right)=(d\,Q\boldsymbol{\lambda})_{\mu}(y),
\end{align*}
the last equality being~\eqref{eq:algebraic-gaussian-1} on $\mathbb{L}(La)$.

Let $\mu\ne\nu$. For $x\in\mathbb{L}(a)$ let $P_x$ be the square with the corners $x$, $x+La\,e_{\mu}$, $x+La\,e_{\mu}+La\,e_{\nu}$, $x+La\,e_{\nu}$, and let $\partial P_x$ be its boundary, the closed path of $4L$ bonds through these corners in this order. By~\eqref{eq:algebraic-gaussian-1} on $\mathbb{L}(La)$,
\[
(d\,QA)_{\mu\nu}(y)=(La)^{-1}\left((QA)_{\mu}(y)+(QA)_{\nu}(y+La\,e_{\mu})-(QA)_{\mu}(y+La\,e_{\nu})-(QA)_{\nu}(y)\right).
\]
By the second form of~\eqref{eq:algebraic-gaussian-4}, and since $B(y+La\,e)=B(y)+La\,e$ for $e\in\{e_{\mu},e_{\nu}\}$, each of the four terms in the bracket is $L^{-4}$ times the sum over $x\in B(y)$ of the values of $A$ on one side of $P_x$; the two minus signs belong to the two sides that $\partial P_x$ traverses against the direction of $e_{\mu}$, resp.\ $e_{\nu}$, because $A$ is odd under reversal of the orientation. Therefore
\[
(d\,QA)_{\mu\nu}(y)=(La)^{-1}\cdot L^{-4}\sum_{x\in B(y)}A(\partial P_x).
\]
The square $P_x$ is the union of the $L^2$ plaquettes $p_{i,i'}$, $0\le i,i'\le L-1$, with the corners $x+ia\,e_{\mu}+i'a\,e_{\nu}$, $x+(i+1)a\,e_{\mu}+i'a\,e_{\nu}$, $x+(i+1)a\,e_{\mu}+(i'+1)a\,e_{\nu}$, $x+ia\,e_{\mu}+(i'+1)a\,e_{\nu}$, oriented like $P_x$; by the convention for $F_{\mu\nu}$ in Section~\ref{sec:tori-cochains-hodge}, $(dA)(p_{i,i'})=(dA)_{\mu\nu}(x+ia\,e_{\mu}+i'a\,e_{\nu})$. In the sum of the boundaries of these plaquettes each bond inside $P_x$ occurs twice, with opposite orientations, and each bond of $\partial P_x$ once, with its orientation in $\partial P_x$. By~\eqref{eq:algebraic-gaussian-S2} it follows that
\[
A(\partial P_x)=\sum_{i,i'=0}^{L-1}A(\partial p_{i,i'})=a\sum_{i,i'=0}^{L-1}(dA)_{\mu\nu}(x+ia\,e_{\mu}+i'a\,e_{\nu}),
\]
and therefore, by~\eqref{eq:algebraic-gaussian-4},
\[
(d\,QA)_{\mu\nu}(y)=L^{-5}\sum_{x\in B(y)}\ \sum_{i,i'=0}^{L-1}(dA)_{\mu\nu}(x+ia\,e_{\mu}+i'a\,e_{\nu})=\left(Q^{(2)}dA\right)_{\mu\nu}(y).
\]

(c) By~\eqref{eq:algebraic-gaussian-5}, $(Q_j\boldsymbol{\lambda})(y)$ is $L^{-3j}$ times a sum of $L^{3j}$ terms, each a value of $\boldsymbol{\lambda}$; $(Q_jA)_{\mu}(y)$ is $L^{-4j}$ times a sum of $L^{3j}\cdot L^j$ terms, each a value of $A_{\mu}$; and $\left(Q_j^{(2)}F\right)_{\mu\nu}(y)$ is $L^{-5j}$ times a sum of $L^{3j}\cdot L^{2j}$ terms, each a value of $F_{\mu\nu}$. In particular $Q_jc=c$ for constant scalar functions and constant bond functions $c$. The translations by $l\mathbb{Z}^3\subset\mathbb{L}\left(L^ja\right)$ map blocks onto blocks; hence, by Section~\ref{sec:blocks-levels-rectilinear}, as $y$ ranges over the sites of $\mathbb{T}\left(L^ja,l\right)$ and $x$ over $B_j(y)$, the site $x$ takes every value in $\mathbb{T}(a,l)$ exactly once. Since $x\mapsto x+ia\,e_{\mu}$ is a bijection of the sites of $\mathbb{T}(a,l)$, it follows that
\[
\sum_{y\in\mathbb{T}\left(L^ja,l\right)}(Q_jA)_{\mu}(y)=L^{-4j}\sum_{i=0}^{L^j-1}\ \sum_{x\in\mathbb{T}(a,l)}A_{\mu}(x+ia\,e_{\mu})=L^{-3j}\sum_{x\in\mathbb{T}(a,l)}A_{\mu}(x).
\]
The numbers of sites of $\mathbb{T}(a,l)$ and of $\mathbb{T}\left(L^ja,l\right)$ are $\left(\frac{l}{a}\right)^3$ and $\left(\frac{l}{L^ja}\right)^3=L^{-3j}\cdot\left(\frac{l}{a}\right)^3$. Dividing the last identity by $\left(\frac{l}{L^ja}\right)^3$, one finds that the average of $(Q_jA)_{\mu}$ over the sites of $\mathbb{T}\left(L^ja,l\right)$ equals the average of $A_{\mu}$ over the sites of $\mathbb{T}(a,l)$; by the definition of the mean in Section~\ref{sec:tori-cochains-hodge}, this is $\operatorname{mean}(Q_jA)=\operatorname{mean}(A)$.

(d) By~\eqref{eq:algebraic-gaussian-5}, each of the three operators has the form
\[
(Tf)(y)=\sum_zw(y,z)\cdot f(z),
\]
the sum over the cells $z$ of $\mathbb{T}(a,l)$, where $w(y,z)\ge0$ is the coefficient $L^{-3j}$, $L^{-4j}$ or $L^{-5j}$ of~\eqref{eq:algebraic-gaussian-5} times the number of those terms of the sum~\eqref{eq:algebraic-gaussian-5} at $y$ which are the value of $f$ at the cell $z$. By the count in the proof of (c), $\sum_zw(y,z)=1$ for every cell $y$ of $\mathbb{T}\left(L^ja,l\right)$; for $Q_j^{(2)}$ this is $L^{-5j}\cdot L^{3j}\cdot L^{2j}=1$. Moreover $\sum_yw(y,z)=L^{-3j}$ for every cell $z$ of $\mathbb{T}(a,l)$, the sum being over all cells $y$ of $\mathbb{T}\left(L^ja,l\right)$. Indeed, as $y$ ranges over the sites of $\mathbb{T}\left(L^ja,l\right)$ and $x$ over $B_j(y)$, the site $x$ takes every value in $\mathbb{T}(a,l)$ exactly once (proof of (c)); so for $Q_j$ on bond functions and the cell $(z,\mu)$ the sum is
\[
L^{-4j}\cdot\#\left|\left\{(x,i):x\in\mathbb{T}(a,l),\ 0\le i\le L^j-1,\ x+ia\,e_{\mu}=z\right\}\right|=L^{-4j}\cdot L^j,
\]
since $x=z-ia\,e_{\mu}$ is determined by $i$; for $Q_j^{(2)}$ it is $L^{-5j}\cdot L^{2j}$, the pair $(i,i')$ determining $x$; and for $Q_j$ on scalars it is $L^{-3j}$. The first of these two properties gives
\[
\left|(Tf)(y)\right|\le\sum_zw(y,z)\cdot\left|f(z)\right|\le|f|_{\infty},
\]
that is, $|Tf|_{\infty}\le|f|_{\infty}$. The Cauchy--Schwarz inequality, applied to the families $\left(w(y,z)^{\frac{1}{2}}\right)_z$ and $\left(w(y,z)^{\frac{1}{2}}f(z)\right)_z$, gives
\[
\left|(Tf)(y)\right|^2\le\left(\sum_zw(y,z)\right)\cdot\left(\sum_zw(y,z)\cdot\left|f(z)\right|^2\right)=\sum_zw(y,z)\cdot\left|f(z)\right|^2.
\]
For $y\in Y$ and $z\notin N(Y)$ one has $w(y,z)=0$, and for every $z$ one has $\sum_{y\in Y}w(y,z)\le L^{-3j}$. Multiplying by $\left(L^ja\right)^3$ and summing over $y\in Y$ one obtains
\[
\sum_{y\in Y}\left(L^ja\right)^3\cdot\left|(Tf)(y)\right|^2\le\left(L^ja\right)^3\sum_{z\in N(Y)}\left|f(z)\right|^2\cdot\sum_{y\in Y}w(y,z)\le\left(L^ja\right)^3\cdot L^{-3j}\sum_{z\in N(Y)}\left|f(z)\right|^2,
\]
and $\left(L^ja\right)^3\cdot L^{-3j}=a^3$; this is~\eqref{eq:algebraic-gaussian-6}. If $Y$ is the set of all cells of $\mathbb{T}\left(L^ja,l\right)$, the left side of~\eqref{eq:algebraic-gaussian-6} is $|Tf|_2^2$ and its right side is at most $|f|_2^2$, by the definition of the norms in Section~\ref{sec:tori-cochains-hodge}; this gives the three bounds for the norms. For $a=\eta_k$ and $j=k$ the torus $\mathbb{T}\left(L^ja,l\right)$ is $T^0_k=\mathbb{T}(1,l)$, on which $|\cdot|_2$ is the counting norm $|\cdot|_{\ell^2}$ (Section~\ref{sec:tori-cochains-hodge}); so the bound for $Q_k^{(2)}$ holds with $C_Q=1$.

(e) The pairing on $T^0_k$ carries the weight $1$, and the pairing on $T_k$ the weight $\eta^3$. For $A\in C^1(T_k)$ and $A'\in C^1(T^0_k)$, \eqref{eq:algebraic-gaussian-5} gives
\[
\langle A',Q_kA\rangle=\sum_{\mu}\ \sum_{y\in T^0_k}A'_{\mu}(y)\cdot L^{-4k}\sum_{x\in B_k(y)}\ \sum_{i=0}^{L^k-1}A_{\mu}(x+i\eta e_{\mu}).
\]
For fixed $\mu$ and $i$, the map $(y,x)\mapsto z:=x+i\eta e_{\mu}$ is a bijection of the set of the pairs $(y,x)$ with $y\in T^0_k$ and $x\in B_k(y)$ onto $T_k$, and its inverse is given by $x=z-i\eta e_{\mu}$, $y=[x]=[z-i\eta e_{\mu}]$. With $L^{-4k}=\eta^3\cdot L^{-k}$ it follows that
\[
\langle A',Q_kA\rangle=\sum_{\mu}\ \sum_{z\in T_k}\eta^3\cdot A_{\mu}(z)\cdot L^{-k}\sum_{i=0}^{L^k-1}A'_{\mu}\left([z-i\eta e_{\mu}]\right).
\]
As this holds for every $A\in C^1(T_k)$, the transpose $Q_k^TA'$ is the bond function stated in (e). Its value at $(z,\mu)$ is the average of $L^k$ values of $A'_{\mu}$; hence $\left|Q_k^TA'\right|_{\infty}\le|A'|_{\infty}$.

(f) The function $\boldsymbol{\Sigma}_jA'$ is periodic, because $A'$ is periodic and the translations by $l\mathbb{Z}^3\subset\mathbb{L}\left(L^ja\right)$ map blocks onto blocks; so it belongs to $C^1(a,l)$. Let $y\in\mathbb{L}\left(L^ja\right)$ and $x\in B_j(y)$. Since $L^j$ is odd, a site $x'\in\mathbb{L}(a)$ lies in $B_j(y)$ if and only if $|x'_{\nu}-y_{\nu}|\le\left(L^j-1\right)a/2$ for every direction $\nu$. The sites $x+ia\,e_{\mu}$, $0\le i\le L^j$, of the straight path from $x$ to $x+L^ja\,e_{\mu}$ differ from $x$ only in the coordinate $\mu$, and the number $x_{\mu}+ia-y_{\mu}$ is a multiple of $a$ which increases with $i$ from at least $-\left(L^j-1\right)a/2$ to at most $\left(3L^j-1\right)a/2$. Hence such a site lies in $B_j(y)$ if $x_{\mu}+ia-y_{\mu}\le\left(L^j-1\right)a/2$, and in $B_j\left(y+L^ja\,e_{\mu}\right)$ otherwise. Consequently, among the $L^j$ bonds $(x+ia\,e_{\mu},x+(i+1)a\,e_{\mu})$, $0\le i\le L^j-1$, of this path exactly one, namely the one with $x_{\mu}+ia-y_{\mu}=\left(L^j-1\right)a/2$, leads from $B_j(y)$ into $B_j\left(y+L^ja\,e_{\mu}\right)$, and each of the others has both end points in one block. By the definition of $\boldsymbol{\Sigma}_j$, the bond function $\boldsymbol{\Sigma}_jA'$ takes the value $L^j\cdot A'_{\mu}(y)$ on that bond, and the value $0$ on each of the other bonds of the path, since the final point of such a bond lies in the block of its initial point and therefore not in the next block in the direction $e_{\mu}$. With~\eqref{eq:algebraic-gaussian-5} it follows that
\[
(Q_j\boldsymbol{\Sigma}_jA')_{\mu}(y)=L^{-4j}\sum_{x\in B_j(y)}L^j\cdot A'_{\mu}(y)=L^{-4j}\cdot L^{3j}\cdot L^j\cdot A'_{\mu}(y)=A'_{\mu}(y).
\]

(g) Let $y\in\mathbb{L}\left(L^2a\right)$ and $y':=y/L\in\mathbb{L}(La)$. The dilation $x\mapsto x/L$ maps $\mathbb{L}(La)$ onto $\mathbb{L}(a)$ and the cells of $\mathbb{L}(La)$ onto those of $\mathbb{L}(a)$, and it maps the block $B(y)=\left\{x\in\mathbb{L}(La):|x-y|<L^2a/2\right\}$ onto the block $B(y')=\left\{x'\in\mathbb{L}(a):|x'-y'|<La/2\right\}$. By~\eqref{eq:algebraic-gaussian-3}, for $x\in\mathbb{L}(La)$,
\[
f_L(x)=L^{\frac{1}{2}}\cdot f(x/L)\ \text{ if }f\text{ is a scalar function},\qquad(f_L)_{\mu}(x)=L^{-\frac{1}{2}}\cdot f_{\mu}(x/L)\ \text{ if }f\text{ is a bond function},
\]
and $(F_L)_{\mu\nu}(x)=L^{-\frac{3}{2}}\cdot F_{\mu\nu}(x/L)$; in particular $(f_L)_{\mu}(x+iLa\,e_{\mu})=L^{-\frac{1}{2}}\cdot f_{\mu}(x/L+ia\,e_{\mu})$ and $(F_L)_{\mu\nu}(x+iLa\,e_{\mu}+i'La\,e_{\nu})=L^{-\frac{3}{2}}\cdot F_{\mu\nu}(x/L+ia\,e_{\mu}+i'a\,e_{\nu})$. The coefficients $L^{-3}$, $L^{-4}$, $L^{-5}$ in~\eqref{eq:algebraic-gaussian-4} do not depend on the lattice spacing. Hence~\eqref{eq:algebraic-gaussian-4}, applied on $\mathbb{T}(La,l)$ and on $\mathbb{T}(a,l/L)$, gives for a scalar function $f$
\[
\left(Q(f_L)\right)(y)=L^{-3}\sum_{x\in B(y)}L^{\frac{1}{2}}\cdot f(x/L)=L^{\frac{1}{2}}\cdot(Qf)(y')=\left((Qf)_L\right)(y),
\]
for a bond function $f$
\[
\left(Q(f_L)\right)_{\mu}(y)=L^{-4}\sum_{x\in B(y)}\ \sum_{i=0}^{L-1}L^{-\frac{1}{2}}\cdot f_{\mu}(x/L+ia\,e_{\mu})=L^{-\frac{1}{2}}\cdot(Qf)_{\mu}(y')=\left((Qf)_L\right)_{\mu}(y),
\]
and for the plaquette function $F$
\begin{align*}
\left(Q^{(2)}(F_L)\right)_{\mu\nu}(y)&=L^{-5}\sum_{x\in B(y)}\ \sum_{i,i'=0}^{L-1}L^{-\frac{3}{2}}\cdot F_{\mu\nu}(x/L+ia\,e_{\mu}+i'a\,e_{\nu})\\
&=L^{-\frac{3}{2}}\cdot\left(Q^{(2)}F\right)_{\mu\nu}(y')=\left(\left(Q^{(2)}F\right)_L\right)_{\mu\nu}(y).
\end{align*}

Let $a=\eta_k$ and $j=k$, so that the hypothesis $l\in L^ja\mathbb{N}$ reads $l\in\mathbb{N}$. By the last assertion of (a), the operators $Q_k$ and $Q_k^{(2)}$ of Section~\ref{sec:averaging-operators}, defined there as $k$-fold compositions, coincide with the operators of the same names that Section~\ref{sec:averaging-fluctuations-gauge} defines by closed formulas; likewise, for $a=1$ and $l\in L\mathbb{N}$, the operator $Q$ on bond functions coincides with the average of one step of Section~\ref{sec:averaging-fluctuations-gauge}. For these operators clauses (b), (c) and (d) give the identities and the bound of Proposition~\ref{prop:block-averages} for $Q_k$ and $Q_k^{(2)}$, with $C_Q=1$, and clause (c) with $a=1$, $j=1$ and $l\in L\mathbb{N}$ gives the equality of the means of $QA$ and $A$. Proposition~\ref{prop:block-averages} follows.
\end{proof}

\subsection{Gauge conditions and minimisers}\label{sec:gauge-conditions-minimisers}

For the gauge conditions and the minimisers of this section see \cite[Lemmas~13--17, (211), (272), (276)]{DimockMB} and \cite[(170), (187)--(192)]{DimockUR}.

\subsubsection{Towers, restricted gauge functions and axial functionals}\label{sec:towers-restricted-gauge}

A \emph{tower} is a triple $(a,J,l)$ consisting of a real number $a>0$, an integer $J\ge0$ and a number $l\in L^Ja\mathbb{N}$. For a tower one puts $a_j:=L^ja$ for $j\ge0$; the torus $\mathbb{T}(a,l)$ is the fine lattice of the tower, and $\mathbb{T}(a_J,l)$ is its coarse lattice.

Let $l\in\mathbb{N}$ be a side. For every $k\ge0$ the triple $(\eta_k,k,l)$ is a tower, the \emph{tower of level~$k$}; its fine lattice is $T_k$ and its coarse lattice is $T^0_k$. If $l\in L\mathbb{N}$, then $(\eta_k,k+1,l)$ is a tower with fine lattice $T_k$ and coarse lattice $T^1_k$. For $k=0$ and $l\in L\mathbb{N}$ the tower $(\eta_k,k+1,l)$ is the one-step tower $(1,1,l)$. Its fine lattice $\mathbb{T}(1,l)$ and its coarse lattice $\mathbb{T}(L,l)$ are, for every $k'\ge0$, the lattices $T^0_{k'}$ and $T^1_{k'}$ of level $k'$ with side $l$. If $l\in L\mathbb{N}$, the scaling~\eqref{eq:algebraic-gaussian-3} maps the cochains on the fine lattice and on the coarse lattice of the tower $(\eta_{k+1},k+1,l/L)$ of level $k+1$ onto the cochains on the fine lattice and on the coarse lattice, respectively, of the tower $(\eta_k,k+1,l)$.

Let $(a,J,l)$ be a tower. The elements of the subspace
\[
\mathcal{N}_J:=\left\{\boldsymbol{\lambda}\in C^0(a,l):Q_J\boldsymbol{\lambda}=0\right\}
\]
of $C^0(a,l)$ are called \emph{restricted gauge functions}, and $\mathscr{R}_J:C^0(a,l)\to C^0(a,l)$ denotes the orthogonal projection onto the subspace $\operatorname{Lap}(\mathcal{N}_J)$. For the tower of level $k$ the operator $Q_J=Q_k$ is the averaging operator of Section~\ref{sec:averaging-fluctuations-gauge} (Proposition~\ref{prop:averaging-operators}(a)), and $\mathscr{R}_J=\mathscr{R}_k$ is the projection of the Landau condition of Section~\ref{sec:averaging-fluctuations-gauge}.

Let $\mathscr{A}\in C^1(a,l)$ and $\boldsymbol{\lambda}\in\mathcal{N}_J$. Then $d(\mathscr{A}-d\boldsymbol{\lambda})=d\mathscr{A}$, because $d\,d\boldsymbol{\lambda}=0$, and $Q_J(\mathscr{A}-d\boldsymbol{\lambda})=Q_J\mathscr{A}-d\,Q_J\boldsymbol{\lambda}=Q_J\mathscr{A}$ by Proposition~\ref{prop:averaging-operators}(b). The \emph{restricted orbits} are the subsets $\mathscr{A}+d\mathcal{N}_J$ of $C^1(a,l)$, $\mathscr{A}\in C^1(a,l)$.

Let $0\le j<J$. For $A\in C^1(a_j,l)$, $y\in\mathbb{L}(a_{j+1})$ and $x$ in the block $B(y)$ of $\mathbb{L}(a_j)$ put
\begin{equation}\label{eq:algebraic-gaussian-7}
(\tau^{\pi}A)(y,x):=A\left(\boldsymbol{\Gamma}^{\pi}(y,x)\right),\qquad(\tau A)(y,x):=\frac{1}{3!}\sum_{\pi}(\tau^{\pi}A)(y,x);
\end{equation}
here $\pi$ ranges over the six permutations of the three directions, and $\boldsymbol{\Gamma}^{\pi}(y,x)$ is the rectilinear path of $\mathbb{L}(a_j)$ from $y$ to $x$ (Section~\ref{sec:blocks-levels-rectilinear}), all of whose sites lie in $B(y)$. The functional $\tau$ is that of \cite[(38)]{DimockUR}. In what follows $\tau^{\#}$ is one of these functionals, fixed once and for all: either $\tau$, the functional of the symmetrised axial gauge, which is the axial functional of Section~\ref{sec:averaging-fluctuations-gauge}, or $\tau^{\pi_0}$ for one fixed permutation $\pi_0$, the functional of the unsymmetrised axial gauge. For $\boldsymbol{\lambda}\in C^0(a_j,l)$ the rule~\eqref{eq:algebraic-gaussian-S1} on $\mathbb{L}(a_j)$, applied to the paths $\boldsymbol{\Gamma}^{\pi}(y,x)$ from $y$ to $x$, gives
\begin{equation}\label{eq:algebraic-gaussian-8}
(\tau^{\#}d\boldsymbol{\lambda})(y,x)=a_j^{-1}\left(\boldsymbol{\lambda}(x)-\boldsymbol{\lambda}(y)\right).
\end{equation}

A field $\mathscr{A}\in C^1(a,l)$ is called \emph{axial} if $(\tau^{\#}Q_j\mathscr{A})(y,x)=0$ for all $j$ with $0\le j<J$, all $y\in\mathbb{L}(a_{j+1})$ and all $x\in B(y)$ (compare \cite[(168)]{DimockUR} and \cite[(209), (260)]{DimockMB}); it is called \emph{Landau} if $\mathscr{R}_J\,\delta\mathscr{A}=0$ (compare \cite[(267)]{DimockMB}).

\subsubsection{Gauge fixing in restricted orbits}\label{sec:gauge-fixing-restricted}

\begin{lemma}[Closed fields]\label{lem:closed-fields}
Let $(a,J,l)$ be a tower, and let $B\in C^1(a,l)$ be closed, that is, $dB=0$.
\begin{enumerate}
\item[(a)] One has $d\,Q_JB=0$, and $Q_JB$ has the same mean as $B$: $\operatorname{mean}(Q_JB)=\operatorname{mean}(B)$.
\item[(b)] Suppose that $Q_JB=0$. Then $B$ is the coboundary of a restricted gauge function, and of only one: $B=d\boldsymbol{\lambda}$ holds for one and only one $\boldsymbol{\lambda}\in\mathcal{N}_J$.
\item[(c)] The restriction of $\operatorname{Lap}$ to $\mathcal{N}_J$ is injective.
\end{enumerate}
\end{lemma}

\begin{proof}
(a) By Proposition~\ref{prop:averaging-operators}(b), $d\,Q_JB=Q_J^{(2)}dB=0$, and Proposition~\ref{prop:averaging-operators}(c) gives $\operatorname{mean}(Q_JB)=\operatorname{mean}(B)$.

(b) As $B$ is closed, the Hodge decomposition~\eqref{eq:algebraic-gaussian-Hodge} gives $B=c+d\boldsymbol{\lambda}$ with a constant $c$ and some $\boldsymbol{\lambda}\in C^0(a,l)$. By Proposition~\ref{prop:averaging-operators}(b), (c), $Q_JB=c+d\,Q_J\boldsymbol{\lambda}$, and the mean of this field is $c$, since exact fields have mean zero (Section~\ref{sec:tori-cochains-hodge}). From $Q_JB=0$ it follows that $c=0$ and $d\,Q_J\boldsymbol{\lambda}=0$; by~\eqref{eq:algebraic-gaussian-Hodge} on $C^0(a_J,l)$, $Q_J\boldsymbol{\lambda}$ is then a constant $t$. By Proposition~\ref{prop:averaging-operators}(c), $Q_J(\boldsymbol{\lambda}-t)=t-t=0$; thus $\boldsymbol{\lambda}-t\in\mathcal{N}_J$, and $d(\boldsymbol{\lambda}-t)=d\boldsymbol{\lambda}=B$. If $\boldsymbol{\lambda},\boldsymbol{\lambda}'\in\mathcal{N}_J$ and $d\boldsymbol{\lambda}=d\boldsymbol{\lambda}'$, then $\boldsymbol{\lambda}-\boldsymbol{\lambda}'$ is a constant $s$ by~\eqref{eq:algebraic-gaussian-Hodge}, and $s=Q_Js=Q_J\boldsymbol{\lambda}-Q_J\boldsymbol{\lambda}'=0$.

(c) Let $\boldsymbol{\lambda}\in\mathcal{N}_J$ with $\operatorname{Lap}\boldsymbol{\lambda}=0$. Then $|d\boldsymbol{\lambda}|_2^2=\langle\boldsymbol{\lambda},\operatorname{Lap}\boldsymbol{\lambda}\rangle=0$, so $\boldsymbol{\lambda}$ is a constant by~\eqref{eq:algebraic-gaussian-Hodge}, and this constant equals $Q_J\boldsymbol{\lambda}=0$.
\end{proof}

\begin{lemma}[{Axial representatives \cite[Lemma~13]{DimockMB}}]\label{lem:axial-representatives}
Let $(a,J,l)$ be a tower. To each $\mathscr{A}\in C^1(a,l)$ corresponds a unique restricted gauge function $\boldsymbol{\lambda}\in\mathcal{N}_J$ for which the field $\mathscr{A}-d\boldsymbol{\lambda}$ is axial, and the map $\mathscr{A}\mapsto\boldsymbol{\lambda}$ so defined is linear. Each restricted orbit therefore meets the set of axial fields in a single point.
\end{lemma}

\begin{proof}
For $0\le j<J$ let $t_j:=\tau^{\#}Q_j\mathscr{A}$, a function of the pairs $(y,x)$ with $y\in\mathbb{L}(a_{j+1})$ and $x\in B(y)$; one has $t_j(y,y)=0$, since the path $\boldsymbol{\Gamma}^{\pi}(y,y)$ contains no bond. For $\boldsymbol{\lambda}\in C^0(a,l)$ put $\boldsymbol{\lambda}_j:=Q_j\boldsymbol{\lambda}\in C^0(a_j,l)$ for $0\le j\le J$; then $Q\boldsymbol{\lambda}_j=\boldsymbol{\lambda}_{j+1}$ for $j<J$. By Proposition~\ref{prop:averaging-operators}(b), $Q_j\,d\boldsymbol{\lambda}=d\boldsymbol{\lambda}_j$, and~\eqref{eq:algebraic-gaussian-8} yields
\[
\left(\tau^{\#}Q_j(\mathscr{A}-d\boldsymbol{\lambda})\right)(y,x)=t_j(y,x)-a_j^{-1}\left(\boldsymbol{\lambda}_j(x)-\boldsymbol{\lambda}_j(y)\right).
\]
Hence $\boldsymbol{\lambda}$ lies in $\mathcal{N}_J$ and $\mathscr{A}-d\boldsymbol{\lambda}$ is axial if and only if $\boldsymbol{\lambda}_J=0$ and
\[
\boldsymbol{\lambda}_j(x)=\boldsymbol{\lambda}_j(y)+a_j\cdot t_j(y,x)\qquad\text{for }0\le j<J,\ y\in\mathbb{L}(a_{j+1}),\ x\in B(y).
\]
One has $\#|B(y)|=L^3$ and, by~\eqref{eq:algebraic-gaussian-4}, $(Q\boldsymbol{\lambda}_j)(y)=L^{-3}\sum_{x\in B(y)}\boldsymbol{\lambda}_j(x)$. So the average of the last identity over $x\in B(y)$ reads
\[
\boldsymbol{\lambda}_{j+1}(y)=\boldsymbol{\lambda}_j(y)+a_j\cdot L^{-3}\sum_{x'\in B(y)}t_j(y,x').
\]
It follows that the two conditions are equivalent to $\boldsymbol{\lambda}_J=0$ together with
\begin{equation}\label{eq:algebraic-gaussian-9}
\boldsymbol{\lambda}_j(x)=\boldsymbol{\lambda}_{j+1}(y)+a_j\left(t_j(y,x)-L^{-3}\sum_{x'\in B(y)}t_j(y,x')\right)
\end{equation}
for $x\in B(y)$, $y\in\mathbb{L}(a_{j+1})$ and $j=J-1,\dots,0$. Indeed, \eqref{eq:algebraic-gaussian-9} is obtained by solving the averaged identity for $\boldsymbol{\lambda}_j(y)$ and inserting the result into the pointwise one; conversely, \eqref{eq:algebraic-gaussian-9} at $x=y$, where $t_j(y,y)=0$, is the averaged identity, and the difference of~\eqref{eq:algebraic-gaussian-9} at $x$ and at $y$ is the pointwise identity.

The blocks $B(y)$, $y\in\mathbb{L}(a_{j+1})$, form a partition of $\mathbb{L}(a_j)$. Therefore~\eqref{eq:algebraic-gaussian-9} defines, starting from $\boldsymbol{\lambda}_J:=0$, functions $\boldsymbol{\lambda}_{J-1},\dots,\boldsymbol{\lambda}_0$ one after the other; each $\boldsymbol{\lambda}_j$ is $l$-periodic, as $t_j$ and $\boldsymbol{\lambda}_{j+1}$ are, so that $\boldsymbol{\lambda}_j\in C^0(a_j,l)$. The bracket in~\eqref{eq:algebraic-gaussian-9} has average zero over $x\in B(y)$; hence $Q\boldsymbol{\lambda}_j=\boldsymbol{\lambda}_{j+1}$, and the function $\boldsymbol{\lambda}:=\boldsymbol{\lambda}_0$ satisfies $Q_j\boldsymbol{\lambda}=\boldsymbol{\lambda}_j$ for $0\le j\le J$. Thus $\boldsymbol{\lambda}$ has the two required properties. If $\boldsymbol{\lambda}'$ has them as well, then the functions $Q_j\boldsymbol{\lambda}'$ satisfy $Q_J\boldsymbol{\lambda}'=0$ and~\eqref{eq:algebraic-gaussian-9}, so they coincide with the functions $\boldsymbol{\lambda}_j$ just constructed, and $\boldsymbol{\lambda}'=Q_0\boldsymbol{\lambda}'=\boldsymbol{\lambda}_0=\boldsymbol{\lambda}$. The solution $\boldsymbol{\lambda}$ is a linear function of $(t_j)_{0\le j<J}$, hence of $\mathscr{A}$.

A restricted orbit $\mathscr{A}+d\mathcal{N}_J$ consists of the fields $\mathscr{A}-d\boldsymbol{\lambda}$ with $\boldsymbol{\lambda}\in\mathcal{N}_J$; by what has been proved, exactly one of them is axial.
\end{proof}

\begin{lemma}[{Landau representatives \cite[Lemma~14]{DimockMB}}]\label{lem:landau-representatives}
Let $(a,J,l)$ be a tower. Given $\mathscr{A}\in C^1(a,l)$, the equation
\[
\mathscr{R}_J\,\delta(\mathscr{A}-d\boldsymbol{\lambda})=0
\]
in the unknown $\boldsymbol{\lambda}\in\mathcal{N}_J$ has one and only one solution, and this solution is a linear function of $\mathscr{A}$. Each restricted orbit therefore meets the set of Landau fields in a single point.
\end{lemma}

\begin{proof}
For $\boldsymbol{\lambda}\in\mathcal{N}_J$ the function $\delta d\boldsymbol{\lambda}=\operatorname{Lap}\boldsymbol{\lambda}$ lies in $\operatorname{Lap}(\mathcal{N}_J)$, which is the range of $\mathscr{R}_J$; hence
\[
\mathscr{R}_J\,\delta(\mathscr{A}-d\boldsymbol{\lambda})=\mathscr{R}_J\,\delta\mathscr{A}-\operatorname{Lap}\boldsymbol{\lambda}.
\]
The range of $\mathscr{R}_J$ is $\operatorname{Lap}(\mathcal{N}_J)$ by definition, and $\operatorname{Lap}$ is injective on $\mathcal{N}_J$ by Lemma~\ref{lem:closed-fields}(c); so the restriction of $\operatorname{Lap}$ to $\mathcal{N}_J$ is a linear isomorphism of $\mathcal{N}_J$ onto the range of $\mathscr{R}_J$. As $\mathscr{R}_J\,\delta\mathscr{A}$ belongs to that range, the identity $\operatorname{Lap}\boldsymbol{\lambda}=\mathscr{R}_J\,\delta\mathscr{A}$ holds for one and only one $\boldsymbol{\lambda}\in\mathcal{N}_J$, the preimage of $\mathscr{R}_J\,\delta\mathscr{A}$ under this isomorphism; it is a linear function of $\mathscr{A}$. A restricted orbit $\mathscr{A}+d\mathcal{N}_J$ consists of the fields $\mathscr{A}-d\boldsymbol{\lambda}$ with $\boldsymbol{\lambda}\in\mathcal{N}_J$, and exactly one of them is Landau.
\end{proof}

\subsubsection{Minimisation under an averaging constraint}\label{sec:minimisation-averaging-constraint}

\begin{proposition}[{Constrained minimisers \cite[Lemma~15, (211), (231), (272), (276)]{DimockMB}, \cite[(170), (187)--(191)]{DimockUR}}]\label{prop:constrained-minimisers}
Let $(a,J,l)$ be a tower and $A\in C^1(a_J,l)$; the mean of $A$ is arbitrary.
\begin{enumerate}
\item[(i)] The function $\mathscr{A}\mapsto|d\mathscr{A}|_2^2$ attains its minimum on the affine space $\left\{\mathscr{A}\in C^1(a,l):Q_J\mathscr{A}=A\right\}$. A field $\mathscr{A}\in C^1(a,l)$ is a minimiser if and only if $Q_J\mathscr{A}=A$ and $\langle d\mathscr{A},dB\rangle=0$ for every $B\in C^1(a,l)$ with $Q_JB=0$. The set of all minimisers is a single restricted orbit. If $\mathscr{A}$ is a minimiser and $B\in C^1(a,l)$ satisfies $Q_JB=0$, then
\[
|d(\mathscr{A}+B)|_2^2=|d\mathscr{A}|_2^2+|dB|_2^2.
\]
\item[(ii)] Among the minimisers there is a unique Landau field, written $H_{a,J}A$, and a unique axial field, written $H^x_{a,J}A$. The field $H_{a,J}A$ is the unique solution of the problem of Section~\ref{sec:averaging-fluctuations-gauge}, posed for the tower $(a,J,l)$:
\[
\text{minimise }|d\mathscr{A}|_2^2\text{ over }\mathscr{A}\in C^1(a,l)\text{ subject to }Q_J\mathscr{A}=A\text{ and }\mathscr{R}_J\,\delta\mathscr{A}=0.
\]
Likewise, $H^x_{a,J}A$ is the unique solution of the problem with the constraints $Q_J\mathscr{A}=A$ and $\tau^{\#}Q_j\mathscr{A}=0$ for $0\le j<J$. The maps $H_{a,J},H^x_{a,J}:C^1(a_J,l)\to C^1(a,l)$ are linear with real coefficients and are right inverses of $Q_J$, that is, $Q_JH_{a,J}=Q_JH^x_{a,J}=1$. There is a linear map $O_{a,J}:C^1(a_J,l)\to\mathcal{N}_J$ such that
\[
H^x_{a,J}=H_{a,J}+d\,O_{a,J}.
\]
If $J=0$, then $H_{a,0}=H^x_{a,0}=1$.
\item[(iii)] $H_{a,J}c=c$ for every constant bond function $c$. One has $dA=0$ if and only if $dH_{a,J}A=0$. For every $\mu\in C^0(a_J,l)$ the field $H_{a,J}\,d\mu$ is exact. Moreover, $\operatorname{mean}(H_{a,J}A)=\operatorname{mean}(A)$.
\item[(iv)] Let $Q_J^T:C^1(a_J,l)\to C^1(a,l)$ be the transpose of $Q_J$ with respect to the inner products~\eqref{eq:algebraic-gaussian-1}, and fix a real number $a'>0$. Then
\[
\delta d+d\,\mathscr{R}_J\,\delta+a'\cdot Q_J^TQ_J
\]
is a symmetric positive definite operator of $C^1(a,l)$ into itself. Let $G$ be its inverse. The operator $Q_JGQ_J^T$ on $C^1(a_J,l)$ is positive definite as well, and
\[
H_{a,J}=GQ_J^T\left(Q_JGQ_J^T\right)^{-1}.
\]
\item[(v)] $|dA|_2\le|dH_{a,J}A|_2$.
\item[(vi)] Let $(a,J,l/L)$ be a tower; then $(La,J,l)$ is a tower as well. For every $A\in C^1(a_J,l/L)$
\[
H_{La,J}(A_L)=(H_{a,J}A)_L,\qquad H^x_{La,J}(A_L)=(H^x_{a,J}A)_L,
\]
where $H_{a,J}$, $H^x_{a,J}$ are the maps of the tower $(a,J,l/L)$, and $H_{La,J}$, $H^x_{La,J}$ those of the tower $(La,J,l)$.
\end{enumerate}
\end{proposition}

The maps $H^x_{a,J}$ and $O_{a,J}$ depend on the functional $\tau^{\#}$ that has been fixed; in Section~\ref{sec:one-averaging-step} and in all later sections $\tau^{\#}$ is the symmetrised functional $\tau$. At level $k$, that is, for the tower $(\eta_k,k,l)$, put
\[
H_k:=H_{\eta_k,k},\qquad H^x_k:=H^x_{\eta_k,k},\qquad\Delta_k:=(dH_k)^TdH_k,
\]
the transpose being taken with respect to the inner products~\eqref{eq:algebraic-gaussian-1} on $C^1(T^0_k)$ and on $C^2(T_k)$. By (ii), $dH^x_k=dH_k+d\,d\,O_{\eta_k,k}=dH_k$, so that also $\Delta_k=(dH^x_k)^TdH^x_k$; and by (i), (ii), for every $A\in C^1(T^0_k)$,
\[
\langle A,\Delta_kA\rangle=|dH_kA|_2^2=\min\left\{|d\mathscr{A}|_2^2:\mathscr{A}\in C^1(T_k),\ Q_k\mathscr{A}=A\right\}.
\]
By (ii), $H_0=H^x_0=1$, hence $\Delta_0=\delta d$. By (ii) and Proposition~\ref{prop:averaging-operators}(a), $H_k$ and $\Delta_k$ are the operators $H_k$ and $\Delta_k$ of Section~\ref{sec:averaging-fluctuations-gauge}.

\begin{proof}[Proof of Proposition~\ref{prop:constrained-minimisers}]
(i) Apply the lemma on minimisation on an affine subspace (Section~\ref{sec:quadratic-minimisation-block}) to the Euclidean spaces $E:=C^1(a,l)$ and $E':=C^2(a,l)$ with the inner products~\eqref{eq:algebraic-gaussian-1}, to the map $T:=d$, to the subspace $W:=\ker Q_J$ of $E$ and to the point $x_0:=\boldsymbol{\Sigma}_JA$. By Proposition~\ref{prop:averaging-operators}(f), $Q_Jx_0=A$, so that
\[
x_0+W=\left\{\mathscr{A}\in C^1(a,l):Q_J\mathscr{A}=A\right\}.
\]
The lemma gives the existence of the minimum, the characterisation of the minimisers, and the set of the minimisers in the form $\mathscr{A}_{*}+(W\cap\ker d)$, where $\mathscr{A}_{*}$ is any minimiser; its identity~\eqref{eq:algebraic-gaussian-LA1} is the identity for $|d(\mathscr{A}+B)|_2^2$. A field in $W\cap\ker d$ is closed and is annihilated by $Q_J$; by Lemma~\ref{lem:closed-fields}(b) it lies in $d\mathcal{N}_J$. Conversely, $d\mathcal{N}_J\subset W\cap\ker d$, as shown in Section~\ref{sec:towers-restricted-gauge}. Hence $W\cap\ker d=d\mathcal{N}_J$, and the minimisers form the restricted orbit $\mathscr{A}_{*}+d\mathcal{N}_J$.

(ii) By (i) the minimisers make up a single restricted orbit; this orbit contains one and only one Landau field by Lemma~\ref{lem:landau-representatives}, and one and only one axial field by Lemma~\ref{lem:axial-representatives}. Let $\mathscr{A}\in C^1(a,l)$ obey both constraints $Q_J\mathscr{A}=A$ and $\mathscr{R}_J\,\delta\mathscr{A}=0$. The first constraint makes $\mathscr{A}$ a competitor in the minimisation of (i); its value $|d\mathscr{A}|_2^2$ cannot fall below the least value of (i), and it reaches that value only if $\mathscr{A}$ is a minimiser, in which case the second constraint makes $\mathscr{A}$ the Landau minimiser $H_{a,J}A$. Conversely, $H_{a,J}A$ obeys both constraints and realises the least value of (i). Consequently $H_{a,J}A$ is the only solution of the constrained problem. For the axial problem, let $\mathscr{A}$ satisfy $Q_J\mathscr{A}=A$ and $\tau^{\#}Q_j\mathscr{A}=0$ for $0\le j<J$. Then $\mathscr{A}$ competes in (i), its value is not below the least value of (i), and it attains that value only if it is a minimiser, in which case, being axial, it is the axial minimiser $H^x_{a,J}A$; since $H^x_{a,J}A$ obeys these constraints and attains the least value, it is the only solution of the second problem.

Consider, for pairs $(\mathscr{A},A)\in C^1(a,l)\times C^1(a_J,l)$, the system
\[
Q_J\mathscr{A}=A;\qquad\langle d\mathscr{A},dB\rangle=0\ \text{ for all }B\in\ker Q_J;\qquad\mathscr{R}_J\,\delta\mathscr{A}=0.
\]
Each of its equations is a real linear equation without constant term in the pair; so the pairs that satisfy the system form a real linear subspace of $C^1(a,l)\times C^1(a_J,l)$. By the characterisation of the minimisers in (i), a pair satisfies the system exactly when $\mathscr{A}$ is a Landau minimiser for the datum $A$, that is, by the first paragraph, exactly when $\mathscr{A}=H_{a,J}A$. The subspace is therefore the set $\left\{(H_{a,J}A,A):A\in C^1(a_J,l)\right\}$; since it contains sums and real multiples of its elements, $H_{a,J}(A+A')=H_{a,J}A+H_{a,J}A'$ and $H_{a,J}(sA)=s\cdot H_{a,J}A$ for $A,A'\in C^1(a_J,l)$ and $s\in\mathbb{R}$. Thus $H_{a,J}$ is real and linear. The system
\[
Q_J\mathscr{A}=A;\qquad\langle d\mathscr{A},dB\rangle=0\ \text{ for all }B\in\ker Q_J;\qquad\tau^{\#}Q_j\mathscr{A}=0\ \text{ for }0\le j<J
\]
consists again of real linear equations without constant term in $(\mathscr{A},A)$, and by (i) and the first paragraph its solutions are exactly the pairs $(H^x_{a,J}A,A)$, $A\in C^1(a_J,l)$. These pairs therefore form a real linear subspace, whence $H^x_{a,J}(A+A')=H^x_{a,J}A+H^x_{a,J}A'$ and $H^x_{a,J}(sA)=s\cdot H^x_{a,J}A$ for $A,A'\in C^1(a_J,l)$ and $s\in\mathbb{R}$; so $H^x_{a,J}$ is real and linear. Applying $Q_J$ to $H_{a,J}A$ or to $H^x_{a,J}A$ gives back $A$, by the first constraint; this is the identity $Q_JH_{a,J}=Q_JH^x_{a,J}=1$. By (i), $B:=H^x_{a,J}A-H_{a,J}A$ lies in $d\mathcal{N}_J$; it is closed, and $Q_JB=0$, so by Lemma~\ref{lem:closed-fields}(b) there is exactly one element $O_{a,J}A$ of $\mathcal{N}_J$ with $B=d\,O_{a,J}A$. This uniqueness and the linearity of $H_{a,J}$ and $H^x_{a,J}$ show that $O_{a,J}$ is linear. For $J=0$ one has $Q_0=1$, and the constraint reads $\mathscr{A}=A$.

(iii) The constant field $c$ on $\mathbb{T}(a,l)$ satisfies $Q_Jc=c$ by Proposition~\ref{prop:averaging-operators}(c), and $dc=0$. So $c$ lies in the affine space of (i) for the datum $c$, where it attains the least possible value $0$; it is a minimiser. Moreover $\delta c=0$, because $\langle\delta c,\boldsymbol{\lambda}\rangle=\langle c,d\boldsymbol{\lambda}\rangle=0$ for every $\boldsymbol{\lambda}\in C^0(a,l)$ (an exact field has mean zero, and the bond functions with mean zero are those orthogonal to the constants). Thus $c$ is Landau, and $H_{a,J}c=c$ by (ii).

Let $dA=0$. By~\eqref{eq:algebraic-gaussian-Hodge} on $\mathbb{T}(a_J,l)$, $A=c+d\mu$ with a constant $c$ and some $\mu\in C^0(a_J,l)$. Let $\tilde{\mu}\in C^0(a,l)$ be the function with $\tilde{\mu}(x):=\mu(y)$ for $x\in B_J(y)$, $y\in\mathbb{L}(a_J)$. The fine field $c+d\tilde{\mu}$ is closed. Proposition~\ref{prop:averaging-operators}(a) gives $Q_J\tilde{\mu}=\mu$, since $\#|B_J(y)|=L^{3J}$, and with Proposition~\ref{prop:averaging-operators}(b), (c)
\[
Q_J(c+d\tilde{\mu})=c+d\,Q_J\tilde{\mu}=c+d\mu=A.
\]
So the minimum in (i) is $0$, and $dH_{a,J}A=0$. Conversely, if $dH_{a,J}A=0$, then $dA=d\,Q_JH_{a,J}A=Q_J^{(2)}dH_{a,J}A=0$ by Proposition~\ref{prop:averaging-operators}(b). By Proposition~\ref{prop:averaging-operators}(c), $\operatorname{mean}(H_{a,J}A)=\operatorname{mean}(Q_JH_{a,J}A)=\operatorname{mean}(A)$. Finally, let $\mu\in C^0(a_J,l)$. The field $d\mu$ is closed, hence so is $H_{a,J}\,d\mu$, and $\operatorname{mean}(H_{a,J}\,d\mu)=\operatorname{mean}(d\mu)=0$. By~\eqref{eq:algebraic-gaussian-Hodge} a closed bond function is the sum of an exact one and a constant one, and this constant is its mean; so $H_{a,J}\,d\mu$ is exact.

(iv) Since $\mathscr{R}_J$ is an orthogonal projection, $\mathscr{R}_J^T=\mathscr{R}_J=\mathscr{R}_J^2$. As $\delta=d^T$, the operator $\delta d+d\,\mathscr{R}_J\,\delta+a'\cdot Q_J^TQ_J$ is symmetric, and for $\mathscr{A}\in C^1(a,l)$
\[
\left\langle\mathscr{A},\left(\delta d+d\,\mathscr{R}_J\,\delta+a'\cdot Q_J^TQ_J\right)\mathscr{A}\right\rangle=|d\mathscr{A}|_2^2+|\mathscr{R}_J\,\delta\mathscr{A}|_2^2+a'\cdot|Q_J\mathscr{A}|_2^2.
\]
If this vanishes, then $d\mathscr{A}=0$ and $Q_J\mathscr{A}=0$, so that $\mathscr{A}=d\boldsymbol{\lambda}$ with $\boldsymbol{\lambda}\in\mathcal{N}_J$ by Lemma~\ref{lem:closed-fields}(b); moreover $0=\mathscr{R}_J\,\delta d\boldsymbol{\lambda}=\mathscr{R}_J\operatorname{Lap}\boldsymbol{\lambda}=\operatorname{Lap}\boldsymbol{\lambda}$, because $\operatorname{Lap}\boldsymbol{\lambda}\in\operatorname{Lap}(\mathcal{N}_J)$, and Lemma~\ref{lem:closed-fields}(c) gives $\boldsymbol{\lambda}=0$, hence $\mathscr{A}=0$. So the operator is positive definite, and so is its inverse $G$. The transpose $Q_J^T$ is injective, because $Q_J$ is onto by Proposition~\ref{prop:averaging-operators}(f). Hence $\langle v,Q_JGQ_J^Tv\rangle=\langle Q_J^Tv,GQ_J^Tv\rangle>0$ for every $v\ne0$ in $C^1(a_J,l)$, and $Q_JGQ_J^T$ is positive definite.

On the affine space $\left\{\mathscr{A}\in C^1(a,l):Q_J\mathscr{A}=A\right\}$ the form $\langle\mathscr{A},G^{-1}\mathscr{A}\rangle$ equals
\[
|d\mathscr{A}|_2^2+|\mathscr{R}_J\,\delta\mathscr{A}|_2^2+a'\cdot|A|_2^2\ \ge\ |dH_{a,J}A|_2^2+a'\cdot|A|_2^2
\]
by (i), and equality holds if and only if $\mathscr{A}$ is a minimiser of $|d\mathscr{A}|_2^2$ which is Landau, that is, by (ii), if and only if $\mathscr{A}=H_{a,J}A$. So $H_{a,J}A$ is the minimiser of this positive definite form on the affine space. On the other hand, for $B\in\ker Q_J$ and $s\in\mathbb{R}$,
\[
\left\langle\mathscr{A}+sB,G^{-1}(\mathscr{A}+sB)\right\rangle=\langle\mathscr{A},G^{-1}\mathscr{A}\rangle+2s\cdot\langle G^{-1}\mathscr{A},B\rangle+s^2\cdot\langle B,G^{-1}B\rangle;
\]
so the minimiser $\mathscr{A}_{*}$ of the form on the affine space is characterised by the condition that $G^{-1}\mathscr{A}_{*}$ be orthogonal to $\ker Q_J$, that is, that $G^{-1}\mathscr{A}_{*}=Q_J^Tv$ for some $v\in C^1(a_J,l)$, the orthogonal complement of $\ker Q_J$ being the range of $Q_J^T$. Then $\mathscr{A}_{*}=GQ_J^Tv$ and $A=Q_J\mathscr{A}_{*}=Q_JGQ_J^Tv$, so that $v=\left(Q_JGQ_J^T\right)^{-1}A$. With $\mathscr{A}_{*}=H_{a,J}A$ this is the asserted formula.

(v) By Proposition~\ref{prop:averaging-operators}(b) and $Q_JH_{a,J}=1$, one has $dA=d\,Q_JH_{a,J}A=Q_J^{(2)}dH_{a,J}A$, and Proposition~\ref{prop:averaging-operators}(d) gives $|dA|_2\le|dH_{a,J}A|_2$.

(vi) The condition $l/L\in L^Ja\mathbb{N}$ is the same as $l\in L^J(La)\mathbb{N}$; so $(a,J,l/L)$ is a tower if and only if $(La,J,l)$ is one, and $(La)_j=La_j$. For $0\le j\le J$ and $p=0,1,2$ the scaling $f\mapsto f_L$ of~\eqref{eq:algebraic-gaussian-3} is a linear bijection of $C^p(a_j,l/L)$ onto $C^p(La_j,l)$. By Section~\ref{sec:scaling}, $d(f_L)=(df)_L$ and $|F_L|_2=|F|_2$, so that $|d(\mathscr{A}_L)|_2=|d\mathscr{A}|_2$ for $\mathscr{A}\in C^1(a,l/L)$. Let $0\le j\le J$. The operator $Q_j$ of the tower $(a,J,l/L)$ is the composition of $j$ one-step operators $Q$; for $0\le i<j$ the factor number $i+1$ acts on the functions on $\mathbb{T}(a_i,l/L)$, and the corresponding factor of $Q_j$ for the tower $(La,J,l)$ acts on the functions on $\mathbb{T}(La_i,l)$. As $l/L\in L^Ja\mathbb{N}$ and $i+1\le J$,
\[
l\in L^{J+1}a\mathbb{N}\subset L^{i+2}a\mathbb{N}=L^2a_i\mathbb{N},
\]
which is the hypothesis of Proposition~\ref{prop:averaging-operators}(g) with $a_i$ in place of $a$. That clause therefore gives $Q(g_L)=(Qg)_L$ for the scalar functions and the bond functions $g$ on $\mathbb{T}(a_i,l/L)$, $0\le i<j$; applying these identities in turn for $i=0,1,\dots,j-1$ one obtains $Q_j(f_L)=(Q_jf)_L$ for the scalar functions and the bond functions $f$ on $\mathbb{T}(a,l/L)$. By~\eqref{eq:algebraic-gaussian-1} and~\eqref{eq:algebraic-gaussian-3},
\[
\langle\boldsymbol{\lambda}_L,\boldsymbol{\lambda}'_L\rangle=L^4\cdot\langle\boldsymbol{\lambda},\boldsymbol{\lambda}'\rangle,\qquad\langle A_L,A'_L\rangle=L^2\cdot\langle A,A'\rangle
\]
for $\boldsymbol{\lambda},\boldsymbol{\lambda}'\in C^0(a,l/L)$ and $A,A'\in C^1(a,l/L)$. Hence
\[
\langle\delta(A_L),\boldsymbol{\lambda}_L\rangle=\langle A_L,(d\boldsymbol{\lambda})_L\rangle=L^2\cdot\langle\delta A,\boldsymbol{\lambda}\rangle=L^{-2}\cdot\langle(\delta A)_L,\boldsymbol{\lambda}_L\rangle,
\]
that is, $\delta(A_L)=L^{-2}\cdot(\delta A)_L$, and $\operatorname{Lap}(\boldsymbol{\lambda}_L)=L^{-2}\cdot(\operatorname{Lap}\boldsymbol{\lambda})_L$. Since $Q_J(\boldsymbol{\lambda}_L)=(Q_J\boldsymbol{\lambda})_L$, the scaling maps the space $\mathcal{N}_J$ of the tower $(a,J,l/L)$ onto the space $\mathcal{N}_J$ of the tower $(La,J,l)$; by the formula for $\operatorname{Lap}(\boldsymbol{\lambda}_L)$ it maps $\operatorname{Lap}(\mathcal{N}_J)$ onto $\operatorname{Lap}(\mathcal{N}_J)$. On scalar functions it is $L^2$ times an isometry; so it maps the orthogonal complement of $\operatorname{Lap}(\mathcal{N}_J)$ onto the orthogonal complement of $\operatorname{Lap}(\mathcal{N}_J)$, whence $\mathscr{R}_J(f_L)=(\mathscr{R}_Jf)_L$ and
\[
\mathscr{R}_J\,\delta(\mathscr{A}_L)=L^{-2}\cdot(\mathscr{R}_J\,\delta\mathscr{A})_L.
\]
Let $0\le j<J$, $y\in\mathbb{L}(a_{j+1})$ and $x\in B(y)$. The path $\boldsymbol{\Gamma}^{\pi}(Ly,Lx)$ of $\mathbb{L}(La_j)$ consists of the bonds $Lb$ with $b$ a bond of $\boldsymbol{\Gamma}^{\pi}(y,x)$, and the pairs $(Ly,Lx)$ are exactly the pairs of~\eqref{eq:algebraic-gaussian-7} for the lattice $\mathbb{L}(La_j)$. Therefore
\[
\left(\tau^{\#}Q_j(\mathscr{A}_L)\right)(Ly,Lx)=L^{-\frac{1}{2}}\cdot\left(\tau^{\#}Q_j\mathscr{A}\right)(y,x).
\]
It follows that $\mathscr{A}\mapsto\mathscr{A}_L$ maps the affine space $\left\{\mathscr{A}:Q_J\mathscr{A}=A\right\}$ of the tower $(a,J,l/L)$ onto the affine space $\left\{\mathscr{A}:Q_J\mathscr{A}=A_L\right\}$ of the tower $(La,J,l)$, preserves $|d\mathscr{A}|_2$, and maps Landau fields onto Landau fields and axial fields onto axial fields. So it maps the minimisers of (i) for the first tower onto those for the second, and (ii) gives $H_{La,J}(A_L)=(H_{a,J}A)_L$ and $H^x_{La,J}(A_L)=(H^x_{a,J}A)_L$.

For the towers $(\eta_k,k,l)$, $k\ge0$, parts (ii) and (iii) contain all assertions of Proposition~\ref{prop:background-field} on $H_k$; Proposition~\ref{prop:background-field} follows.
\end{proof}

\subsubsection{Estimates in the axial gauge}\label{sec:estimates-axial-gauge}

\paragraph{Bond functions on $\mathbb{L}(a)$.}
Let $a>0$. A bond function on $\mathbb{L}(a)$ is a real function $Z$ on the oriented bonds of $\mathbb{L}(a)$ such that $Z(x',x)=-Z(x,x')$ for every oriented bond $(x,x')$; no periodicity is required, and by the identification of Section~\ref{sec:tori-cochains-hodge} the elements of $C^1(a,l)$ are the $l$-periodic bond functions on $\mathbb{L}(a)$. For such a function $Z$, the components $Z_{\mu}(x)$, the coboundary $dZ$, which is a function on the oriented plaquettes of $\mathbb{L}(a)$ changing sign with the orientation, and the numbers $Z(\boldsymbol{\Gamma})$, $\boldsymbol{\Gamma}$ a chain of oriented bonds, are defined by the formulas of Section~\ref{sec:tori-cochains-hodge}; in particular $(dZ)_{\nu\nu'}(x)$ is given by~\eqref{eq:algebraic-gaussian-1}. If the corners of a plaquette $p$ of $\mathbb{L}(a)$ are $x$, $x+ae_{\nu}$, $x+ae_{\nu}+ae_{\nu'}$, $x+ae_{\nu'}$ with $\nu<\nu'$, then $(dZ)(p):=(dZ)_{\nu\nu'}(x)$ is the value of $dZ$ on $p$ with the orientation $[x,x+ae_{\nu},x+ae_{\nu}+ae_{\nu'},x+ae_{\nu'}]$. For $y\in\mathbb{L}(La)$ and $x\in B(y)$ the numbers $(\tau^{\pi}Z)(y,x)$ and $(\tau Z)(y,x)$ are defined by~\eqref{eq:algebraic-gaussian-7}, with $a$ in the place of $a_j$, and $(\tau^{\#}Z)(y,x)$ is formed with the functional $\tau^{\#}$ of Section~\ref{sec:towers-restricted-gauge}; for a direction $\mu$ the number $(QZ)_{\mu}(y)$ is defined by~\eqref{eq:algebraic-gaussian-4}. The identity~\eqref{eq:algebraic-gaussian-S2} holds for $Z$ and every oriented plaquette $p$ of $\mathbb{L}(a)$: for $p=[x,x+ae_{\mu},x+ae_{\mu}+ae_{\nu},x+ae_{\nu}]$ with $\mu\ne\nu$, the oddness of $Z$ gives $Z(\partial p)=Z_{\mu}(x)+Z_{\nu}(x+ae_{\mu})-Z_{\mu}(x+ae_{\nu})-Z_{\nu}(x)$, which equals $a\cdot(dZ)_{\mu\nu}(x)=a\cdot(dZ)(p)$ by~\eqref{eq:algebraic-gaussian-1}.

The chain of a path is the sum of its bonds in the abelian group generated by the oriented bonds of $\mathbb{L}(a)$ subject to the relations $(x',x)=-(x,x')$. A bond function $Z$ extends to a homomorphism from this group to $\mathbb{R}$, and $Z(\boldsymbol{\Gamma})$ is the value of this homomorphism at the chain of $\boldsymbol{\Gamma}$; hence two paths with the same chain have the same value under $Z$.

\paragraph{Ribbons.}
Let $R=(w_0,w_1,\dots,w_n)$ be a sequence of points of $\mathbb{L}(a)$ in which any two consecutive points are nearest neighbours; it is regarded as the path with the bonds $(w_i,w_{i+1})$, $0\le i<n$. Let $\mu$ be a direction and $m\ge1$ an integer. For $0\le i<n$ the step $i$ of $R$ leads from $w_i$ to $w_{i+1}$; it is parallel to $e_{\mu}$ if $w_{i+1}-w_i=\pm ae_{\mu}$. Let $\mathrm{Rib}(R,me_{\mu})$ be the family of plaquettes indexed by the pairs $(i,j)$ such that $0\le j<m$ and the step $i$ is not parallel to $e_{\mu}$, whose member $(i,j)$ is the plaquette with the corners
\[
w_i+jae_{\mu},\qquad w_i+(j+1)ae_{\mu},\qquad w_{i+1}+(j+1)ae_{\mu},\qquad w_{i+1}+jae_{\mu}.
\]
For $v\in\mathbb{L}(a)$ let $R+v:=(w_0+v,\dots,w_n+v)$ be the translated path; for $w\in\mathbb{L}(a)$, $\boldsymbol{\Gamma}_{w,w+mae_{\mu}}$ is the straight path of Section~\ref{sec:blocks-levels-rectilinear} that joins $w$ to $w+mae_{\mu}$.

\begin{lemma}[Ribbon identity]\label{lem:ribbon-identity}
Let $R=(w_0,\dots,w_n)$, $\mu$ and $m$ be as above, and let $Z$ be a bond function on $\mathbb{L}(a)$. Then
\begin{equation}\label{eq:algebraic-gaussian-10}
Z\left(\boldsymbol{\Gamma}_{w_0,w_0+mae_{\mu}}\right)+Z(R+mae_{\mu})-Z\left(\boldsymbol{\Gamma}_{w_n,w_n+mae_{\mu}}\right)-Z(R)=a\sum_{p\in\mathrm{Rib}(R,me_{\mu})}\epsilon_p\cdot(dZ)(p)
\end{equation}
for suitable signs $\epsilon_p\in\{1,-1\}$ attached to the members $p=(i,j)$ of $\mathrm{Rib}(R,me_{\mu})$; the sign $\epsilon_p$ of the member $(i,j)$ is determined by $(R,\mu,m,i,j)$ and does not depend on $Z$. Here the sum extends over the members $p$ of the family, and $(dZ)(p)$ is the value of $dZ$ on the plaquette of the member $p$. The number of members of $\mathrm{Rib}(R,me_{\mu})$ is $m\cdot n_{\mathrm{perp}}$, where
\[
n_{\mathrm{perp}}:=\#\left|\left\{i:0\le i<n,\ w_{i+1}-w_i\ne\pm ae_{\mu}\right\}\right|.
\]
If each coordinate varies monotonically along $R$, then distinct members have distinct plaquettes. All corners of the members belong to the set $\left\{w_i+jae_{\mu}:0\le i\le n,\ 0\le j\le m\right\}$.
\end{lemma}

\begin{proof}
For a step $(w,w')$ of $R$ put
\[
c(w,w'):=Z\left(\boldsymbol{\Gamma}_{w,w+mae_{\mu}}\right)+Z(w+mae_{\mu},w'+mae_{\mu})-Z\left(\boldsymbol{\Gamma}_{w',w'+mae_{\mu}}\right)-Z(w,w').
\]
In the sum of the numbers $c(w_i,w_{i+1})$ over $0\le i<n$ the terms $Z(\boldsymbol{\Gamma}_{w_i,w_i+mae_{\mu}})$ with $0<i<n$ cancel, the second terms add up to $Z(R+mae_{\mu})$ and the fourth terms to $-Z(R)$; so this sum is the left side of~\eqref{eq:algebraic-gaussian-10}.

Let $(w,w')$ be a step of $R$. When $w'=w\pm ae_{\mu}$, both $\boldsymbol{\Gamma}_{w,w+mae_{\mu}}+(w+mae_{\mu},w'+mae_{\mu})$ and $(w,w')+\boldsymbol{\Gamma}_{w',w'+mae_{\mu}}$ are equal, as chains, to the straight path that joins $w$ to $w'+mae_{\mu}$; hence $c(w,w')=0$. Otherwise $c(w,w')=\sum_{j=0}^{m-1}s_j$, where
\begin{multline*}
s_j:=Z(w+jae_{\mu},w+(j+1)ae_{\mu})+Z(w+(j+1)ae_{\mu},w'+(j+1)ae_{\mu})\\
-Z(w'+jae_{\mu},w'+(j+1)ae_{\mu})-Z(w+jae_{\mu},w'+jae_{\mu}),
\end{multline*}
because in $\sum_js_j$ the terms $Z(w+jae_{\mu},w'+jae_{\mu})$ with $0<j<m$ cancel. The four bonds that occur in $s_j$ are the sides of the plaquette $p$ with the corners $w+jae_{\mu}$, $w+(j+1)ae_{\mu}$, $w'+(j+1)ae_{\mu}$, $w'+jae_{\mu}$; as $Z$ is odd under reversal of the bonds, $s_j=Z(\partial p')$ for the oriented plaquette
\[
p':=\left[w+jae_{\mu},\,w+(j+1)ae_{\mu},\,w'+(j+1)ae_{\mu},\,w'+jae_{\mu}\right].
\]
By~\eqref{eq:algebraic-gaussian-S2}, $s_j=a\cdot(dZ)(p')=\pm a\cdot(dZ)(p)$, with the sign $+$ if and only if $p'$ carries the orientation that enters the definition of $(dZ)(p)$; this sign is determined by $w$, $w'$, $j$ and $\mu$, and not by $Z$. For $(w,w')=(w_i,w_{i+1})$ the plaquette $p$ is the member $(i,j)$ of $\mathrm{Rib}(R,me_{\mu})$, and~\eqref{eq:algebraic-gaussian-10} follows.

The corners of the member $(i,j)$ are among the points $w_{i'}+j'ae_{\mu}$ with $i'\in\{i,i+1\}$ and $j'\in\{j,j+1\}$, and each step of $R$ that is not parallel to $e_{\mu}$ gives exactly $m$ members; this proves the assertions on the number of the members and on their corners. Assume now that each coordinate varies monotonically along $R$. The plaquette of the member $(i,j)$ lies in the plane spanned by $e_{\mu}$ and the direction of the step $i$, and its projection along $e_{\mu}$ is the projection of the bond $(w_i,w_{i+1})$. Two steps of $R$ in one direction $\nu\ne\mu$ change the coordinate $\nu$ on two different intervals, since this coordinate is monotone along $R$; hence, for two steps $i\ne i'$ of $R$, neither of them parallel to $e_{\mu}$, the bonds $(w_i,w_{i+1})$ and $(w_{i'},w_{i'+1})$ have different directions or different projections along $e_{\mu}$. Consequently two members $(i,j)$ and $(i',j')$ with $i\ne i'$ have distinct plaquettes; so have two members $(i,j)$ and $(i,j')$ with $j\ne j'$, whose plaquettes differ by the translation by $(j'-j)ae_{\mu}$.
\end{proof}

\paragraph{Three special cases.}
Fix a bond function $Z$ on $\mathbb{L}(a)$, a direction $\mu$, a point $y\in\mathbb{L}(La)$, a permutation $\pi$ of the directions and a point $x\in\mathbb{L}(a)$, and put $y':=y+Lae_{\mu}$. The path $\boldsymbol{\Gamma}^{\pi}(y,x)$ is composed of three straight legs, along which the coordinates $\pi_1$, $\pi_2$, $\pi_3$, in this order, are brought to their final values; the leg of the coordinate $\mu$ is called the $\mu$-leg, and it is empty if $x_{\mu}=y_{\mu}$.

\emph{First case.} Denote by $z$ the point at which the $\mu$-leg of $\boldsymbol{\Gamma}^{\pi}(y,x)$ ends; thus $z_{\nu}=x_{\nu}$ if $\nu=\mu$ or $\nu$ precedes $\mu$ in $\pi$, and $z_{\nu}=y_{\nu}$ if $\nu$ follows $\mu$ in $\pi$. Let $R$ denote the remaining piece of $\boldsymbol{\Gamma}^{\pi}(y,x)$, from $z$ to $x$. The part of $\boldsymbol{\Gamma}^{\pi}(y,x)$ from $y$ to $z$ is $\boldsymbol{\Gamma}^{\pi}(y,z)$. The legs of $\boldsymbol{\Gamma}^{\pi}(y,x+ae_{\mu})$ that precede its $\mu$-leg are those of $\boldsymbol{\Gamma}^{\pi}(y,x)$; its $\mu$-leg is, as a chain, the $\mu$-leg of $\boldsymbol{\Gamma}^{\pi}(y,x)$ plus the bond $(z,z+ae_{\mu})$; and its remaining legs form the path $R+ae_{\mu}$. Hence, as chains,
\[
\boldsymbol{\Gamma}^{\pi}(y,x)=\boldsymbol{\Gamma}^{\pi}(y,z)+R,\qquad\boldsymbol{\Gamma}^{\pi}(y,x+ae_{\mu})=\boldsymbol{\Gamma}^{\pi}(y,z)+(z,z+ae_{\mu})+(R+ae_{\mu}).
\]
The path $R$ consists of the legs of the directions $\nu$ that follow $\mu$ in $\pi$, the leg of $\nu$ having $a^{-1}|x_{\nu}-y_{\nu}|$ steps. The straight paths $\boldsymbol{\Gamma}_{z,z+ae_{\mu}}$ and $\boldsymbol{\Gamma}_{x,x+ae_{\mu}}$ are the single bonds $(z,z+ae_{\mu})$ and $(x,x+ae_{\mu})$. Lemma~\ref{lem:ribbon-identity} with $m=1$, applied to $R$, therefore gives
\begin{align}
Z(x,x+ae_{\mu})&=Z\left(\boldsymbol{\Gamma}^{\pi}(y,x+ae_{\mu})\right)-Z\left(\boldsymbol{\Gamma}^{\pi}(y,x)\right)-a\sum_{p\in\mathrm{Rib}(R,e_{\mu})}\epsilon_p\cdot(dZ)(p),\tag{R1}\label{eq:algebraic-gaussian-R1}\\
\#|\mathrm{Rib}(R,e_{\mu})|&=\sum_{\nu\ \text{after}\ \mu\ \text{in}\ \pi}a^{-1}|x_{\nu}-y_{\nu}|.\notag
\end{align}

\emph{Second case.} Let $R_1$ denote the initial piece of $\boldsymbol{\Gamma}^{\pi}(y,x)$ that ends where its $\mu$-leg begins; it is a path from $y$ to a point $y+t$, made up of the legs of the directions that precede $\mu$ in $\pi$. Since $y'$ differs from $y$ only in the coordinate $\mu$, the part of $\boldsymbol{\Gamma}^{\pi}(y',x)$ that precedes its $\mu$-leg is $R_1+Lae_{\mu}$, from $y'$ to $y'+t$. The legs that follow the $\mu$-legs coincide for the two paths, and the two $\mu$-legs lie on one line and end at the same point, the first beginning at $y+t$ and the second at $y'+t$. Hence, as chains,
\[
\boldsymbol{\Gamma}^{\pi}(y,x)-\boldsymbol{\Gamma}^{\pi}(y',x)=R_1-(R_1+Lae_{\mu})+\boldsymbol{\Gamma}_{y+t,y'+t}.
\]
Lemma~\ref{lem:ribbon-identity} with $m=L$, applied to $R_1$, gives, since $\boldsymbol{\Gamma}_{y,y+Lae_{\mu}}=\boldsymbol{\Gamma}_{y,y'}$ and $\boldsymbol{\Gamma}_{y+t,y+t+Lae_{\mu}}=\boldsymbol{\Gamma}_{y+t,y'+t}$,
\begin{align}
Z\left(\boldsymbol{\Gamma}^{\pi}(y,x)\right)&=Z\left(\boldsymbol{\Gamma}^{\pi}(y',x)\right)+Z\left(\boldsymbol{\Gamma}_{y,y'}\right)-a\sum_{p\in\mathrm{Rib}(R_1,Le_{\mu})}\epsilon_p\cdot(dZ)(p),\tag{R2}\label{eq:algebraic-gaussian-R2}\\
\#|\mathrm{Rib}(R_1,Le_{\mu})|&=L\sum_{\nu\ \text{before}\ \mu\ \text{in}\ \pi}a^{-1}|x_{\nu}-y_{\nu}|.\notag
\end{align}

\emph{Third case.} The translate of $\boldsymbol{\Gamma}^{\pi}(y,x)$ by $Lae_{\mu}$ is $\boldsymbol{\Gamma}^{\pi}(y',x+Lae_{\mu})$, and the steps of $\boldsymbol{\Gamma}^{\pi}(y,x)$ that are not parallel to $e_{\mu}$ are those of the legs of the directions $\nu\ne\mu$. Lemma~\ref{lem:ribbon-identity} with $R=\boldsymbol{\Gamma}^{\pi}(y,x)$ and $m=L$ gives
\begin{align}
Z\left(\boldsymbol{\Gamma}_{y,y'}\right)&=Z\left(\boldsymbol{\Gamma}_{x,x+Lae_{\mu}}\right)+Z\left(\boldsymbol{\Gamma}^{\pi}(y,x)\right)-Z\left(\boldsymbol{\Gamma}^{\pi}(y',x+Lae_{\mu})\right)\notag\\
&\quad+a\sum_{p\in\mathrm{Rib}(\boldsymbol{\Gamma}^{\pi}(y,x),Le_{\mu})}\epsilon_p\cdot(dZ)(p),\tag{R3}\label{eq:algebraic-gaussian-R3}\\
\#|\mathrm{Rib}(\boldsymbol{\Gamma}^{\pi}(y,x),Le_{\mu})|&=L\sum_{\nu\ne\mu}a^{-1}|x_{\nu}-y_{\nu}|.\notag
\end{align}

\paragraph{Plaquettes of blocks; means over the permutations.}
For $y\in\mathbb{L}(La)$ let $\mathrm{Pl}(y)$ consist of those plaquettes of $\mathbb{L}(a)$ all four corners of which belong to $B(y)$; for $y'=y+Lae_{\mu}$ let $\mathrm{Pl}(y,y')$ consist of those plaquettes of $\mathbb{L}(a)$ all four corners of which belong to $B(y)\cup B(y')$. For a quantity $g(\pi)$ that depends on the permutation $\pi$, the mean over $\pi$ is $\frac{1}{6}\sum_{\pi}g(\pi)$ if $\tau^{\#}=\tau$, and $g(\pi_0)$ if $\tau^{\#}=\tau^{\pi_0}$. By~\eqref{eq:algebraic-gaussian-7}, $(\tau^{\#}Z)(y,x)$ is the mean over $\pi$ of $Z(\boldsymbol{\Gamma}^{\pi}(y,x))$, for every $x\in B(y)$.

\begin{lemma}[Bounds inside a block and across a face]\label{lem:bounds-inside-block}
Let $y\in\mathbb{L}(La)$, and let $Z$ be a bond function on $\mathbb{L}(a)$ such that $(\tau^{\#}Z)(y,x)=0$ for every $x\in B(y)$.
\begin{enumerate}
\item[(a)] Let $b=(x,x+ae_{\mu})$ be a bond whose two ends lie in $B(y)$. For every $\pi$ there are a subset $S^{\pi}(b)$ of $\mathrm{Pl}(y)$ with at most $L-1$ elements and signs $\epsilon_p\in\{1,-1\}$, $p\in S^{\pi}(b)$, such that $Z(b)$ is the mean over $\pi$ of $-a\sum_{p\in S^{\pi}(b)}\epsilon_p\cdot(dZ)(p)$. Hence
\[
|Z(b)|\le(L-1)a\cdot\sup\left\{|(dZ)(p)|:p\in\mathrm{Pl}(y)\right\}.
\]
\item[(b)] Put $y':=y+Lae_{\mu}$, assume in addition that $(\tau^{\#}Z)(y',x)=0$ for every $x\in B(y')$, and let
\[
f:=a\cdot\sup\left\{|(dZ)(p)|:p\in\mathrm{Pl}(y,y')\right\}.
\]
Let $b=(x,x+ae_{\mu})$ be a bond that joins $B(y)$ and $B(y')$, that is, $x\in B(y)$ and $x+ae_{\mu}\in B(y')$. Then
\[
Z(b)=L\cdot(QZ)_{\mu}(y)+\mathrm{Tm}_1+\mathrm{Tm}_2+\mathrm{Tm}_3,
\]
where, for $i=1,2,3$, the number $a^{-1}\mathrm{Tm}_i$ is a mean of sums $\sum_{p\in S}\epsilon_p\cdot(dZ)(p)$, each formed with signs $\epsilon_p\in\{1,-1\}$ and with a set $S\subset\mathrm{Pl}(y,y')$ such that $\#|S|\le\acute{n}_i$; the mean is taken over $\pi$ for $i=1,2$, and over $\pi$ and over $x''\in B(y)$ for $i=3$; and $\acute{n}_1:=L-1$, $\acute{n}_2:=\acute{n}_3:=L(L-1)$. Hence
\begin{equation}\label{eq:algebraic-gaussian-11}
|Z(b)|\le L\cdot\left|(QZ)_{\mu}(y)\right|+(2L+1)(L-1)f\le L\cdot\left|(QZ)_{\mu}(y)\right|+2L^2f.
\end{equation}
\end{enumerate}
\end{lemma}

\begin{proof}
Each leg of a rectilinear path changes a single coordinate, always in the same sense, and leaves the other two coordinates fixed; so every coordinate varies monotonically along such a path and along any piece of it. The families $\mathrm{Rib}$ that occur below are formed with pieces of rectilinear paths; by Lemma~\ref{lem:ribbon-identity} they consist of pairwise distinct plaquettes, and they are treated as sets of plaquettes. By the description of the blocks in Section~\ref{sec:blocks-levels-rectilinear}, a point $z\in\mathbb{L}(a)$ belongs to $B(y)$ if and only if $|z_{\nu}-y_{\nu}|\le\frac{(L-1)a}{2}$ for every $\nu$.

(a) For every $\pi$, \eqref{eq:algebraic-gaussian-R1} gives
\[
Z(b)=Z\left(\boldsymbol{\Gamma}^{\pi}(y,x+ae_{\mu})\right)-Z\left(\boldsymbol{\Gamma}^{\pi}(y,x)\right)-a\sum_{p\in\mathrm{Rib}(R,e_{\mu})}\epsilon_p\cdot(dZ)(p),
\]
where $R$ is the part of $\boldsymbol{\Gamma}^{\pi}(y,x)$ that follows its $\mu$-leg. The mean over $\pi$ of the first two terms on the right is $(\tau^{\#}Z)(y,x+ae_{\mu})-(\tau^{\#}Z)(y,x)$, which vanishes because $x$ and $x+ae_{\mu}$ lie in $B(y)$. Put $S^{\pi}(b):=\mathrm{Rib}(R,e_{\mu})$; then $Z(b)$ is the mean over $\pi$ of $-a\sum_{p\in S^{\pi}(b)}\epsilon_p\cdot(dZ)(p)$. By~\eqref{eq:algebraic-gaussian-R1}, and as $x\in B(y)$,
\[
\#|S^{\pi}(b)|\le\sum_{\nu\ne\mu}a^{-1}|x_{\nu}-y_{\nu}|\le L-1.
\]
By Lemma~\ref{lem:ribbon-identity}, the corners of the plaquettes of $S^{\pi}(b)$ have the form $w+jae_{\mu}$ with $w$ a point of $R$ and $j\in\{0,1\}$. The points of $R$ have the coordinate $\mu$ equal to $x_{\mu}$, and each of their other coordinates lies in the closed interval bounded by the corresponding coordinates of $y$ and $x$. So every corner has its coordinate $\mu$ equal to $x_{\mu}$ or to $x_{\mu}+a$, and its other coordinates in these intervals; since $y$, $x$ and $x+ae_{\mu}$ lie in $B(y)$, the corner lies in $B(y)$. Thus $S^{\pi}(b)\subset\mathrm{Pl}(y)$. Each of the sums has at most $L-1$ terms, each of modulus at most $\sup\{|(dZ)(p)|:p\in\mathrm{Pl}(y)\}$, and the bound follows.

(b) Put $x':=x+ae_{\mu}$, so that $x\in B(y)$ and $x'\in B(y')$; the coordinates $\nu\ne\mu$ of $x'$ are those of $x$, and $x_{\mu}=y_{\mu}+\frac{(L-1)a}{2}$. For every $\pi$, \eqref{eq:algebraic-gaussian-R1}, and then \eqref{eq:algebraic-gaussian-R2} with $x'$ in the place of $x$, give
\begin{multline*}
Z(b)=Z\left(\boldsymbol{\Gamma}^{\pi}(y',x')\right)+Z\left(\boldsymbol{\Gamma}_{y,y'}\right)-Z\left(\boldsymbol{\Gamma}^{\pi}(y,x)\right)\\
-a\sum_{p\in\mathrm{Rib}(R_1,Le_{\mu})}\epsilon_p\cdot(dZ)(p)-a\sum_{p\in\mathrm{Rib}(R,e_{\mu})}\epsilon_p\cdot(dZ)(p),
\end{multline*}
where $R$ is the part of $\boldsymbol{\Gamma}^{\pi}(y,x)$ that follows its $\mu$-leg and $R_1$ is the part of $\boldsymbol{\Gamma}^{\pi}(y,x')$ that precedes its $\mu$-leg. The mean over $\pi$ of $Z(\boldsymbol{\Gamma}^{\pi}(y',x'))-Z(\boldsymbol{\Gamma}^{\pi}(y,x))$ is $(\tau^{\#}Z)(y',x')-(\tau^{\#}Z)(y,x)=0$. Hence
\[
Z(b)=Z\left(\boldsymbol{\Gamma}_{y,y'}\right)+\mathrm{Tm}_1+\mathrm{Tm}_2,
\]
where $\mathrm{Tm}_1$ and $\mathrm{Tm}_2$ are the means over $\pi$ of $-a\sum_{p\in\mathrm{Rib}(R,e_{\mu})}\epsilon_p\cdot(dZ)(p)$ and of $-a\sum_{p\in\mathrm{Rib}(R_1,Le_{\mu})}\epsilon_p\cdot(dZ)(p)$, respectively; they have the form required in (b), with the signs $-\epsilon_p$.

Let $x''\in B(y)$; then $x''+Lae_{\mu}\in B(y')$, so that $(\tau^{\#}Z)(y,x'')=0=(\tau^{\#}Z)(y',x''+Lae_{\mu})$. Write~\eqref{eq:algebraic-gaussian-R3} with $x''$ in the place of $x$, take the mean over $\pi$, and then take the mean over the $L^3$ points $x''\in B(y)$. The means over $\pi$ of $Z(\boldsymbol{\Gamma}^{\pi}(y,x''))$ and of $Z(\boldsymbol{\Gamma}^{\pi}(y',x''+Lae_{\mu}))$ are $(\tau^{\#}Z)(y,x'')$ and $(\tau^{\#}Z)(y',x''+Lae_{\mu})$, which vanish, and by~\eqref{eq:algebraic-gaussian-4}
\[
Z\left(\boldsymbol{\Gamma}_{y,y'}\right)=L^{-3}\sum_{x''\in B(y)}Z\left(\boldsymbol{\Gamma}_{x'',x''+Lae_{\mu}}\right)+\mathrm{Tm}_3=L\cdot(QZ)_{\mu}(y)+\mathrm{Tm}_3,
\]
where $\mathrm{Tm}_3$ is the mean over $\pi$ and over $x''\in B(y)$ of $a\sum_{p\in\mathrm{Rib}(\boldsymbol{\Gamma}^{\pi}(y,x''),Le_{\mu})}\epsilon_p\cdot(dZ)(p)$. This proves the formula for $Z(b)$.

It remains to bound the sizes of the three families and to locate their plaquettes. For $\nu\ne\mu$ one has $x'_{\nu}=x_{\nu}$, and $|x_{\nu}-y_{\nu}|\le\frac{(L-1)a}{2}$, $|x''_{\nu}-y_{\nu}|\le\frac{(L-1)a}{2}$ because $x,x''\in B(y)$. Inserted in the numbers of members recorded with~\eqref{eq:algebraic-gaussian-R1}, with~\eqref{eq:algebraic-gaussian-R2} at the point $x'$ and with~\eqref{eq:algebraic-gaussian-R3} at the point $x''$, these bounds show that the families entering $\mathrm{Tm}_1$, $\mathrm{Tm}_2$, $\mathrm{Tm}_3$ have at most $L-1=\acute{n}_1$, $L(L-1)=\acute{n}_2$ and $L(L-1)=\acute{n}_3$ members.

Lemma~\ref{lem:ribbon-identity} describes the corners: a corner of a plaquette of one of these families is a point $w+jae_{\mu}$, where either $w$ lies on $R$ and $0\le j\le1$, or $w$ lies on $R_1$ and $0\le j\le L$, or $w$ lies on $\boldsymbol{\Gamma}^{\pi}(y,x'')$ and $0\le j\le L$. In these three cases the coordinate $\mu$ of $w$ equals $x_{\mu}=y_{\mu}+\frac{(L-1)a}{2}$, equals $y_{\mu}$, and lies between $y_{\mu}$ and $x''_{\mu}$, respectively. Therefore each coordinate $\nu\ne\mu$ of the corner lies between the corresponding coordinates of $y$ and of $x$ (of $x''$ in the third case), and so differs from $y_{\nu}$ by at most $\frac{(L-1)a}{2}$, while its coordinate $\mu$ lies in the interval from $y_{\mu}-\frac{(L-1)a}{2}$ to $y_{\mu}+\frac{(L-1)a}{2}+La$. Every point of $\mathbb{L}(a)$ with these two properties belongs to $B(y)\cup B(y')$, since $y'_{\mu}-\frac{(L-1)a}{2}=y_{\mu}+\frac{(L-1)a}{2}+a$. So all three families are sets of plaquettes of $\mathrm{Pl}(y,y')$.

Finally, $a^{-1}\mathrm{Tm}_i$ is a mean of sums, each consisting of at most $\acute{n}_i$ terms whose moduli do not exceed $a^{-1}f$, so that $|\mathrm{Tm}_i|\le\acute{n}_i\cdot f$, and
\[
\acute{n}_1+\acute{n}_2+\acute{n}_3=(L-1)+2L(L-1)=(2L+1)(L-1)\le2L^2.
\]
This gives~\eqref{eq:algebraic-gaussian-11}.
\end{proof}

\begin{corollary}[{Axial gauge bounds \cite[Lemma~16]{DimockMB}, \cite[(192)]{DimockUR}}]\label{cor:axial-gauge-bounds}
Let $l\in L\mathbb{N}$, and let $Z\in C^1(1,l)$ satisfy $\tau^{\#}Z=0$ in all blocks, that is, $(\tau^{\#}Z)(y,x)=0$ for all $y\in L\mathbb{Z}^3$ and all $x\in B(y)$.
\begin{enumerate}
\item[(a)] $|Z|_{\infty}\le2L^2\cdot|dZ|_{\infty}+L\cdot|QZ|_{\infty}$. More precisely, if $y\in L\mathbb{Z}^3$, $\mu$ is a direction, $y':=y+Le_{\mu}$, and $b$ is a bond whose two ends lie in $B(y)\cup B(y')$, then
\[
|Z(b)|\le2L^2\cdot\sup\left\{|(dZ)(p)|:p\in\mathrm{Pl}(y,y')\right\}+L\cdot\left|(QZ)_{\mu}(y)\right|,
\]
where $(QZ)_{\mu}(y)$ is the value of $QZ$ on the bond $(y,y')$ of $L\mathbb{Z}^3$.
\item[(b)] If $QZ=0$ and $dZ=0$, then $Z=0$.
\item[(c)] If $QZ=0$, then $|Z|_2^2\le39L^4\cdot|dZ|_2^2$, with the same constant for every $l\in L\mathbb{N}$.
\item[(d)] If $QZ=0$, then for every $k\ge0$
\[
\langle Z,\Delta_kZ\rangle\ge|dZ|_2^2\ge\left(39L^4\right)^{-1}\cdot|Z|_2^2,
\]
and $|dZ|_2^2>0$ if $Z\ne0$.
\end{enumerate}
\end{corollary}

\begin{proof}
The bond function $Z$ is regarded as an $l$-periodic bond function on $\mathbb{L}(1)=\mathbb{Z}^3$; the hypotheses of Lemma~\ref{lem:bounds-inside-block} with $a=1$ hold at every $y\in L\mathbb{Z}^3$. For $a=1$ the number $|dZ|_2^2$ is the sum of $|(dZ)(p)|^2$ over the plaquettes $p$ of $\mathbb{T}(1,l)$.

(a) Let $y\in L\mathbb{Z}^3$, $y'=y+Le_{\mu}$, and let $b$ be a bond with both ends in $B(y)\cup B(y')$. If both ends lie in $B(y)$, or both lie in $B(y')$, Lemma~\ref{lem:bounds-inside-block}(a), applied with the centre of that block in the place of $y$, to $b$ or to its reversal, and with the direction of $b$ in the place of $\mu$, gives $|Z(b)|\le(L-1)\cdot\sup\{|(dZ)(p)|:p\in\mathrm{Pl}(y,y')\}$, because $\mathrm{Pl}(y)$ and $\mathrm{Pl}(y')$ are contained in $\mathrm{Pl}(y,y')$; and $L-1\le2L^2$. Otherwise $b$ is the bond $(x,x+e_{\mu})$ with $x\in B(y)$ and $x+e_{\mu}\in B(y')$, or its reversal, on which $Z$ has the same modulus; in this case~\eqref{eq:algebraic-gaussian-11} with $a=1$ is the asserted bound. Every bond of $\mathbb{Z}^3$ has both ends in a set $B(y)\cup B(y+Le_{\mu})$ with $y\in L\mathbb{Z}^3$ and some direction $\mu$; this gives the global bound.

(b) This follows from (a).

(c) Let $QZ=0$. If $S_1,\dots,S_N$ are subsets of a finite set $U$ with at most $n$ elements each, $\sigma_{r,p}\in\{1,-1\}$ for $1\le r\le N$ and $p\in U$, and $t_p$, $p\in U$, are real numbers, then the Cauchy--Schwarz inequality, applied twice, gives
\[
\left|\frac{1}{N}\sum_{r=1}^N\sum_{p\in S_r}\sigma_{r,p}t_p\right|^2\le\frac{1}{N}\sum_{r=1}^N\left|\sum_{p\in S_r}\sigma_{r,p}t_p\right|^2\le\frac{1}{N}\sum_{r=1}^N\#|S_r|\sum_{p\in S_r}t_p^2\le n\sum_{p\in U}t_p^2.
\]

Let $Y\subset L\mathbb{Z}^3$ contain exactly one point of each coset of $l\mathbb{Z}^3$ in $L\mathbb{Z}^3$; here $l\mathbb{Z}^3$ is a subgroup of $L\mathbb{Z}^3$ because $l\in L\mathbb{N}$. The blocks $B(y)$, $y\in L\mathbb{Z}^3$, partition $\mathbb{Z}^3$, and $B(y+lw)=B(y)+lw$ for $w\in\mathbb{Z}^3$; hence the union of the blocks $B(y)$, $y\in Y$, contains exactly one point of each coset of $l\mathbb{Z}^3$ in $\mathbb{Z}^3$. So $|Z|_2^2$ is the sum of $|Z(x,x+e_{\nu})|^2$ over the points $x$ of these blocks and the three directions $\nu$, and each of these bonds either has both ends in the block $B(y)$ that contains $x$, or joins $B(y)$ and $B(y+Le_{\nu})$.

Bonds inside one block. Let $y\in Y$ and let $b$ be a bond with both ends in $B(y)$. Lemma~\ref{lem:bounds-inside-block}(a) with $a=1$, and the inequality above with $U=\mathrm{Pl}(y)$, $n=L-1$ and $t_p=(dZ)(p)$, give
\[
|Z(b)|^2\le(L-1)\sum_{p\in\mathrm{Pl}(y)}|(dZ)(p)|^2.
\]
A block contains $3L^2(L-1)$ bonds $(x,x+e_{\nu})$ with both ends in it: for each of the three directions $\nu$, the block meets $L^2$ lines parallel to $e_{\nu}$, each in $L$ points joined by $L-1$ such bonds. Since distinct blocks are disjoint, a plaquette has its four corners in at most one block, so that no two of the sets $\mathrm{Pl}(y)$, $y\in L\mathbb{Z}^3$, meet; and $B(y+lw)=B(y)+lw$ implies $\mathrm{Pl}(y+lw)=\mathrm{Pl}(y)+lw$ for $w\in\mathbb{Z}^3$. Consequently a plaquette lies in at most one of the sets $\mathrm{Pl}(y)$, $y\in Y$, and no two distinct plaquettes in the union of these sets are translates of each other by a vector of $l\mathbb{Z}^3$. As $dZ$ is $l$-periodic, it follows that
\[
\sum_{y\in Y}\sum_{p\in\mathrm{Pl}(y)}|(dZ)(p)|^2\le|dZ|_2^2.
\]
Summing the bound for $|Z(b)|^2$ over all bonds $b=(x,x+e_{\nu})$ with both ends in one of the blocks $B(y)$, $y\in Y$, one finds that the part of $|Z|_2^2$ due to these bonds is at most $3L^2(L-1)^2\cdot|dZ|_2^2\le3L^4\cdot|dZ|_2^2$.

Bonds between two blocks. Let $y\in Y$, let $\mu$ be a direction, $y':=y+Le_{\mu}$, and let $b$ be a bond that joins $B(y)$ and $B(y')$. By Lemma~\ref{lem:bounds-inside-block}(b) with $a=1$ and $QZ=0$, $Z(b)=\mathrm{Tm}_1+\mathrm{Tm}_2+\mathrm{Tm}_3$. The Cauchy--Schwarz inequality, and the inequality above with $U=\mathrm{Pl}(y,y')$, $n=\acute{n}_i$ and $t_p=(dZ)(p)$, give
\[
|Z(b)|^2\le3\left(|\mathrm{Tm}_1|^2+|\mathrm{Tm}_2|^2+|\mathrm{Tm}_3|^2\right)\le3(2L+1)(L-1)\sum_{p\in\mathrm{Pl}(y,y')}|(dZ)(p)|^2\le6L^2\sum_{p\in\mathrm{Pl}(y,y')}|(dZ)(p)|^2;
\]
the bonds $b$ in question are the $L^2$ bonds $(x,x+e_{\mu})$ with $x\in B(y)$ and $x_{\mu}=y_{\mu}+\frac{L-1}{2}$. The corners of a plaquette of $\mathbb{Z}^3$ are distributed over one, two or four blocks; in the second case the two blocks are $B(y_1)$ and $B(y_1+Le_{\mu})$ for a single $y_1\in L\mathbb{Z}^3$ and a single direction $\mu$. Accordingly, the number of pairs $(y'',\mu)$ with $y''\in L\mathbb{Z}^3$, $\mu\in\{0,1,2\}$ and $p\in\mathrm{Pl}(y'',y''+Le_{\mu})$ is six if the corners of the plaquette $p$ lie in one block $B(y_0)$ (the pairs with $y''\in\{y_0,y_0-Le_{\mu}\}$, for each of the three directions $\mu$), one in the case of two blocks, and zero in the case of four; it never exceeds six. For a plaquette $p$ of $\mathbb{Z}^3$, the triples $(y,\mu,w)$ with $y\in Y$, $w\in\mathbb{Z}^3$ and $p+lw\in\mathrm{Pl}(y,y+Le_{\mu})$ correspond one to one, through $y''=y-lw$, to the pairs $(y'',\mu)$ with $y''\in L\mathbb{Z}^3$ and $p\in\mathrm{Pl}(y'',y''+Le_{\mu})$, because $\mathrm{Pl}(y+lw,y+lw+Le_{\mu})=\mathrm{Pl}(y,y+Le_{\mu})+lw$ and every $y''\in L\mathbb{Z}^3$ is $y-lw$ for exactly one pair $(y,w)\in Y\times\mathbb{Z}^3$. As $dZ$ is $l$-periodic, it follows that
\[
\sum_{y\in Y}\sum_{\mu}\sum_{p\in\mathrm{Pl}(y,y+Le_{\mu})}|(dZ)(p)|^2\le6\cdot|dZ|_2^2.
\]
Summing the bound for $|Z(b)|^2$ over the $L^2$ bonds $b$ that belong to a pair $(y,\mu)$ with $y\in Y$, and then over these pairs, one finds that the part of $|Z|_2^2$ due to the bonds with ends in two different blocks is at most $L^2\cdot6L^2\cdot6\cdot|dZ|_2^2=36L^4\cdot|dZ|_2^2$. Together with the bound for the bonds inside blocks this gives $|Z|_2^2\le(3+36)L^4\cdot|dZ|_2^2=39L^4\cdot|dZ|_2^2$.

(d) Let $QZ=0$ and $k\ge0$. The coarse lattice of the tower $(\eta_k,k,l)$ of level $k$ is $\mathbb{T}(1,l)$. By the definition of $\Delta_k$ and by Proposition~\ref{prop:constrained-minimisers}(v), applied to this tower,
\[
\langle Z,\Delta_kZ\rangle=|dH_kZ|_2^2\ge|dZ|_2^2.
\]
If $Z\ne0$, then $dZ\ne0$ by (b), so that $|dZ|_2^2>0$. The inequality $|dZ|_2^2\ge\left(39L^4\right)^{-1}\cdot|Z|_2^2$ is (c).
\end{proof}

\begin{remark}
Corollary~\ref{cor:axial-gauge-bounds}(c) holds with a constant smaller than $39L^4$, so that $39L^4$ is not optimal: the argument given for (c) yields the inequality with $3L^2(L-1)^2+18L^2(2L+1)(L-1)$ in place of $39L^4$, and this number is less than $3L^4+36L^4=39L^4$.
\end{remark}

\subsection{One averaging step of the photon field}\label{sec:one-averaging-step}

Throughout this section $k\ge0$ is an integer and $l\in L\mathbb{N}$. For the objects and identities of this section see \cite[(174)--(185), (191), (192)]{DimockUR}, \cite[(241)--(259)]{DimockMB} and \cite[(76)--(81), (132), (133), (139)]{DimockUVS}.

\subsubsection{Fluctuation spaces and the map $C$}\label{sec:fluctuation-spaces-map}

Every bond function occurring in this subsection is defined on the bonds of $T^0_k=\mathbb{T}(1,l)$ and is identified with the $l$-periodic bond function on $\mathbb{Z}^3$ that it induces. The \emph{inner} bonds of $\mathbb{Z}^3$ are those whose two ends belong to one and the same $L$-block $B(y)$, $y\in L\mathbb{Z}^3$. Every bond that is not inner joins $B(y)$ and $B(y+Le_{\mu})$ for exactly one pair $(y,\mu)$ of a point $y\in L\mathbb{Z}^3$ and a direction $\mu$; it is then called a \emph{face bond of $(y,\mu)$}. Every pair $(y,\mu)$ has $L^2$ face bonds, and the \emph{central} one among them is
\[
b(y,\mu):=\left(y+\frac{L-1}{2}\,e_{\mu},\ y+\frac{L+1}{2}\,e_{\mu}\right).
\]
For a bond function $Z$ the condition $\tau Z=0$ means that $(\tau Z)(y,x)=0$ for all $y\in L\mathbb{Z}^3$ and all $x\in B(y)$, where $\tau$ is the functional~\eqref{eq:algebraic-gaussian-7} with $a_j=1$. For $x\in B(y)$ all sites of the paths $\boldsymbol{\Gamma}^{\pi}(y,x)$ lie in $B(y)$ (Section~\ref{sec:blocks-levels-rectilinear}); hence $(\tau Z)(y,x)$ depends only on the values of $Z$ on the inner bonds of $B(y)$. Let
\begin{align*}
\mathring{Z}_k&:=\left\{Z\in C^1(T^0_k):\ QZ=0,\ \tau Z=0\right\},\\
V_k&:=\left\{\mathsf{Z}\in C^1(T^0_k):\ \tau\mathsf{Z}=0,\ \mathsf{Z}=0\text{ on all central bonds}\right\},
\end{align*}
both spaces being equipped with the inner product of $C^1(T^0_k)$.

The space $V_k$ is the fluctuation space of Section~\ref{sec:averaging-fluctuations-gauge}; there an element $\mathsf{Z}\in V_k$ is written as a pair $(\mathsf{Z}_1,\mathsf{Z}_2)$, whose first entry $\mathsf{Z}_1$ collects the values of $\mathsf{Z}$ on the inner bonds and whose second entry $\mathsf{Z}_2$ collects its values on the non-central face bonds. The family $\{\Upsilon\}$ of Section~\ref{sec:averaging-fluctuations-gauge}, given by~\eqref{eq:renormalisation-group-1} (see also Proposition~\ref{prop:fluctuation-space}(ii)), is made up of functions of two kinds. Those of the first kind are attached to blocks: the ones attached to $B(y)$, $y\in L\mathbb{Z}^3/l\mathbb{Z}^3$, arise by translation by $y$ from the members of one orthonormal basis, the same for all $y$, of the space of those functions on the inner bonds of $B(0)$ that are annihilated by $\tau$. Those of the second kind are attached to bonds: to each positively oriented non-central face bond $b$ belongs the function $\delta_b$, which takes the value $1$ at $b$ and the value $0$ at every other positively oriented bond. Two members of the family have disjoint supports unless both are attached to the same block, and the members attached to one block are orthonormal, since translation preserves the inner product; moreover, the condition $\tau\mathsf{Z}=0$ splits into separate conditions, one for each block $B(y)$, and the condition for $B(y)$ involves only the values of $\mathsf{Z}$ on the inner bonds of $B(y)$. It follows that $V_k$ has $\{\Upsilon\}$ as an orthonormal basis. Beyond this, \eqref{eq:renormalisation-group-1} demands that the orthonormal basis used for $B(0)$ consist of eigenvectors common to the three reflections; a basis with this property exists by Proposition~\ref{prop:fluctuation-space}(i), which is a part of Proposition~\ref{prop:basis-map-box}.

Let $(y,\mu)$ be a pair of a point $y\in L\mathbb{Z}^3$ and a direction $\mu$. The bonds of direction $\mu$ are oriented in the direction of $e_{\mu}$, and the $L$ sites of each line of direction $\mu$ in an $L$-block are numbered $1,\dots,L$ in the direction of $e_{\mu}$. The weights of the pair $(y,\mu)$ are the numbers
\[
w(b):=\begin{cases}i&\text{if $b$ leads from the $i$-th to the $(i+1)$-th site of its line in $B(y)$},\\ L-i&\text{if $b$ leads from the $i$-th to the $(i+1)$-th site of its line in $B(y+Le_{\mu})$},\end{cases}
\]
attached to the inner bonds $b$ of direction $\mu$ in $B(y)$ and in $B(y+Le_{\mu})$; here $1\le i\le L-1$.

\begin{lemma}[Face averages and the map $C$]\label{lem:face-averages-map}
\begin{enumerate}
\item[(a)] For every $Z\in C^1(T^0_k)$ and every pair $(y,\mu)$
\begin{equation}\label{eq:algebraic-gaussian-12}
L^4\cdot(QZ)_{\mu}(y)=L\cdot\sum_{\substack{b\text{ face bond}\\ \text{of }(y,\mu)}}Z(b)+\sum_bw(b)\cdot Z(b),
\end{equation}
where the second sum is over the inner bonds of direction $\mu$ in $B(y)$ and in $B(y+Le_{\mu})$, with the weights of the pair $(y,\mu)$. The sum of these weights is $\sum_bw(b)=L^3(L-1)$.
\item[(b)] For $\mathsf{Z}\in V_k$ let $C\mathsf{Z}$ be the bond function that coincides with $\mathsf{Z}$ on all bonds other than the central ones, and that takes at the central bond of every pair $(y,\mu)$ the value
\begin{equation}\label{eq:algebraic-gaussian-13}
(C\mathsf{Z})(b(y,\mu)):=-\sum_{\substack{b\text{ face bond of }(y,\mu)\\ b\ne b(y,\mu)}}\mathsf{Z}(b)-L^{-1}\sum_bw(b)\cdot\mathsf{Z}(b),
\end{equation}
the second sum being over the same bonds, with the same weights, as the second sum in~\eqref{eq:algebraic-gaussian-12}. Then $C$ maps $V_k$ linearly and bijectively onto $\mathring{Z}_k$, and for each $\mathsf{Z}\in V_k$ the field $C\mathsf{Z}$ is the unique element of $\mathring{Z}_k$ that agrees with $\mathsf{Z}$ on all bonds that are not central. Moreover
\[
|C\mathsf{Z}|_{\infty}\le L^3\cdot|\mathsf{Z}|_{\infty},\qquad|C\mathsf{Z}|_2\ge|\mathsf{Z}|_2,
\]
and, for every bond $b$, the number $(C\mathsf{Z})(b)$ is determined by the values that $\mathsf{Z}$ takes on the bonds having each of their two ends in one of the $L$-blocks that contain an end of $b$ (there are one or two such blocks).
\item[(c)] $\dim V_k=d_V\cdot n_k$, where $d_V=2-2L^{-3}$.
\end{enumerate}
\end{lemma}

\begin{proof}
(a) By~\eqref{eq:algebraic-gaussian-4} with $a=1$,
\[
L^4\cdot(QZ)_{\mu}(y)=\sum_{x\in B(y)}Z\left(\boldsymbol{\Gamma}_{x,x+Le_{\mu}}\right).
\]
Let $x$ be the $j$-th site of its line in $B(y)$. The straight path $\boldsymbol{\Gamma}_{x,x+Le_{\mu}}$ consists of the bonds of this line that lead from the $i$-th to the $(i+1)$-th site with $j\le i\le L-1$, of the face bond of $(y,\mu)$ on this line, and of the bonds of the continuation of the line in $B(y+Le_{\mu})$ that lead from the $i$-th to the $(i+1)$-th site with $1\le i\le j-1$. Of the $L$ paths $\boldsymbol{\Gamma}_{x,x+Le_{\mu}}$ with $x$ on a fixed line, every one contains the face bond of that line; the bond from the $i$-th to the $(i+1)$-th site in $B(y)$ lies on those with $j\le i$, which are $i$ in number; and the bond from the $i$-th to the $(i+1)$-th site in $B(y+Le_{\mu})$ lies on those with $j\ge i+1$, which are $L-i$ in number. Summation over the $L^2$ lines of direction $\mu$ in $B(y)$ gives~\eqref{eq:algebraic-gaussian-12}, and
\[
\sum_bw(b)=L^2\cdot\sum_{i=1}^{L-1}\left(i+(L-i)\right)=L^3(L-1).
\]

(b) By~\eqref{eq:algebraic-gaussian-12}, the value $Z(b(y,\mu))$ enters $(QZ)_{\mu}(y)$ with the coefficient $L\cdot L^{-4}=L^{-3}$. By the description of the paths in the proof of (a), a path $\boldsymbol{\Gamma}_{x,x+Le_{\mu'}}$ with $x\in B(y')$ contains no face bonds other than face bonds of $(y',\mu')$; so $Z(b(y,\mu))$ enters $(QZ)_{\mu'}(y')$ only for $(y',\mu')=(y,\mu)$. It enters none of the numbers $(\tau Z)(y',x)$, which depend on the inner bonds only. It follows that a field $Z\in C^1(T^0_k)$ with $\tau Z=0$ satisfies $QZ=0$ if and only if
\[
Z(b(y,\mu))=-\sum_{\substack{b\text{ face bond of }(y,\mu)\\ b\ne b(y,\mu)}}Z(b)-L^{-1}\sum_bw(b)\cdot Z(b)\qquad\text{for every pair }(y,\mu),
\]
that is, if and only if $Z$ fulfils~\eqref{eq:algebraic-gaussian-13} after $\mathsf{Z}$ and $C\mathsf{Z}$ have both been replaced by $Z$. No central bond occurs on the right side.

Let $\mathsf{Z}\in V_k$. The field $C\mathsf{Z}$ coincides with $\mathsf{Z}$ on the inner bonds; so $\tau C\mathsf{Z}=\tau\mathsf{Z}=0$, and $C\mathsf{Z}$ satisfies the equations just displayed. Hence $QC\mathsf{Z}=0$, and $C\mathsf{Z}\in\mathring{Z}_k$. If $Z\in\mathring{Z}_k$ coincides with $\mathsf{Z}$ off the central bonds, the same equations express the values of $Z$ at the central bonds through its values elsewhere, which are those of $\mathsf{Z}$; so $Z=C\mathsf{Z}$. Let now $Z\in\mathring{Z}_k$, and let $\mathsf{Z}$ be the bond function equal to $0$ on the central bonds and equal to $Z$ on all other bonds. Then $\tau\mathsf{Z}=\tau Z=0$, because $\mathsf{Z}$ and $Z$ agree on the inner bonds; thus $\mathsf{Z}\in V_k$, and $C\mathsf{Z}=Z$ by the uniqueness just proved. Conversely, for every $\mathsf{Z}\in V_k$, setting the values of $C\mathsf{Z}$ on the central bonds equal to $0$ gives back $\mathsf{Z}$. Hence the map $C$, which is linear by~\eqref{eq:algebraic-gaussian-13}, is a linear bijection of $V_k$ onto $\mathring{Z}_k$, and its inverse replaces the values of an element of $\mathring{Z}_k$ on the central bonds by $0$ and keeps all other values.

On the right side of~\eqref{eq:algebraic-gaussian-13} there are $L^2-1$ coefficients of modulus $1$ and the coefficients $L^{-1}w(b)$, whose sum is $L^{-1}\cdot L^3(L-1)$ by (a). The sum of the moduli of all coefficients is therefore
\[
L^2-1+L^{-1}\cdot L^3(L-1)=L^3-1,
\]
so that $|(C\mathsf{Z})(b(y,\mu))|\le(L^3-1)\cdot|\mathsf{Z}|_{\infty}$; as $C\mathsf{Z}=\mathsf{Z}$ off the central bonds, $|C\mathsf{Z}|_{\infty}\le L^3\cdot|\mathsf{Z}|_{\infty}$. Since $\mathsf{Z}$ vanishes on the central bonds and coincides with $C\mathsf{Z}$ on all other bonds,
\[
|C\mathsf{Z}|_2^2=|\mathsf{Z}|_2^2+\sum_{(y,\mu)}\left|(C\mathsf{Z})(b(y,\mu))\right|^2\ge|\mathsf{Z}|_2^2,
\]
the sum being over $y\in L\mathbb{Z}^3/l\mathbb{Z}^3$ and the three directions $\mu$. Finally, $(C\mathsf{Z})(b)=\mathsf{Z}(b)$ if $b$ is not central, and all bonds on the right side of~\eqref{eq:algebraic-gaussian-13} have both ends in $B(y)\cup B(y+Le_{\mu})$, which is the union of the blocks that contain the two ends of $b(y,\mu)$.

(c) Let $y\in L\mathbb{Z}^3$. By~\eqref{eq:algebraic-gaussian-8} with $a_j=1$ and $\tau^{\#}=\tau$, every function $\boldsymbol{\lambda}$ on the sites satisfies
\[
(\tau\,d\boldsymbol{\lambda})(y,x)=\boldsymbol{\lambda}(x)-\boldsymbol{\lambda}(y)\qquad(x\in B(y)),
\]
and only the values of $d\boldsymbol{\lambda}$ on the inner bonds of $B(y)$ enter the left side. Moreover $(\tau Z)(y,y)=0$ for every bond function $Z$, the paths $\boldsymbol{\Gamma}^{\pi}(y,y)$ being empty. Given a function $f$ on $B(y)\setminus\{y\}$, choose $\boldsymbol{\lambda}$ with $\boldsymbol{\lambda}(y)=0$ and $\boldsymbol{\lambda}=f$ on $B(y)\setminus\{y\}$; the restriction of $d\boldsymbol{\lambda}$ to the inner bonds of $B(y)$ is then mapped to $f$. Hence $Z\mapsto\left((\tau Z)(y,x)\right)_{x\in B(y)\setminus\{y\}}$ maps the space of the functions on the $3L^2(L-1)$ inner bonds of $B(y)$ onto the space of the functions on $B(y)\setminus\{y\}$, which has the dimension $L^3-1$. The equations $(\tau\mathsf{Z})(y',x)=0$, $x\in B(y')$, that belong to one block $B(y')$ involve only the values of $\mathsf{Z}$ on the inner bonds of this block; in particular, the condition $\tau\mathsf{Z}=0$ puts no restriction on the values on the face bonds. Every face bond belongs to exactly one pair $(y,\mu)$, and the three pairs $(y,\mu)$ with a fixed $y$ have together $3(L^2-1)$ face bonds that are not central. So every block contributes
\[
3L^2(L-1)-(L^3-1)+3L^2-3=2L^3-2
\]
to the dimension of $V_k$. The torus $T^0_k$ has $n_k=l^3$ sites and hence $n_k\cdot L^{-3}$ blocks; therefore $\dim V_k=(2L^3-2)\cdot n_k\cdot L^{-3}=d_V\cdot n_k$.
\end{proof}

\subsubsection{The operators $\tilde{\Delta}_k$, $C_k$ and $\mathrm{Sq}_k$}\label{sec:operators}

Let $C^T\colon C^1(T^0_k)\to V_k$ be the adjoint of the map $C\colon V_k\to C^1(T^0_k)$ of Lemma~\ref{lem:face-averages-map}(b), taken with respect to the inner products of $V_k$ and of $C^1(T^0_k)$, and put
\[
\tilde{\Delta}_k:=C^T\Delta_kC,
\]
where $\Delta_k=(dH_k)^TdH_k$ is the operator defined after Proposition~\ref{prop:constrained-minimisers}. This is the operator $\tilde{\Delta}_k$ of Section~\ref{sec:averaging-fluctuations-gauge}; it maps $V_k$ into itself.

\begin{proposition}[The fluctuation operator]\label{prop:fluctuation-operator}
\begin{enumerate}
\item[(a)] The operator $\tilde{\Delta}_k$ on $V_k$ is symmetric. For every $\mathsf{Z}\in V_k$
\[
\langle\mathsf{Z},\tilde{\Delta}_k\mathsf{Z}\rangle=|dH_kC\mathsf{Z}|_2^2\ge|dC\mathsf{Z}|_2^2,
\]
and $|dC\mathsf{Z}|_2^2>0$ if $\mathsf{Z}\ne0$. Moreover, for every $\mathsf{Z}\in V_k$,
\[
\langle\mathsf{Z},\tilde{\Delta}_k\mathsf{Z}\rangle\ge\left(39L^4\right)^{-1}\cdot|\mathsf{Z}|_2^2.
\]
Consequently $\tilde{\Delta}_k$ is invertible, $C_k:=\tilde{\Delta}_k^{-1}$ is an operator on $V_k$, and the operator
\[
C_k^{\frac{1}{2}}:=\pi^{-1}\int_0^{\infty}\mathrm{d}x\,x^{-\frac{1}{2}}\left(\tilde{\Delta}_k+x\right)^{-1}
\]
on $V_k$ is the positive definite symmetric square root of $C_k$.
\item[(b)] The operator $\mathrm{Sq}_k:=C\,C_k^{\frac{1}{2}}$ maps $V_k$ linearly and bijectively onto $\mathring{Z}_k$. Hence $Q\mathrm{Sq}_k=0$ and $\tau\mathrm{Sq}_k=0$, and the mean of $\mathrm{Sq}_k\zeta$ vanishes for every $\zeta\in V_k$.
\item[(c)] For all $\zeta,\zeta'\in V_k$
\[
\langle\mathrm{Sq}_k\zeta,\Delta_k\mathrm{Sq}_k\zeta'\rangle=\langle\zeta,\zeta'\rangle.
\]
\end{enumerate}
\end{proposition}

\begin{proof}
(a) As $\Delta_k=(dH_k)^TdH_k$ is symmetric, so is $\tilde{\Delta}_k=C^T\Delta_kC$, and for $\mathsf{Z}\in V_k$
\[
\langle\mathsf{Z},\tilde{\Delta}_k\mathsf{Z}\rangle=\langle C\mathsf{Z},\Delta_kC\mathsf{Z}\rangle=|dH_kC\mathsf{Z}|_2^2.
\]
By Lemma~\ref{lem:face-averages-map}(b) the field $Z:=C\mathsf{Z}$ lies in $\mathring{Z}_k$, that is, $QZ=0$ and $\tau Z=0$ in all $L$-blocks; and $Z\ne0$ if $\mathsf{Z}\ne0$, because $C$ is injective. Corollary~\ref{cor:axial-gauge-bounds}(d), applied to $Z$, gives
\[
|dH_kZ|_2^2=\langle Z,\Delta_kZ\rangle\ge|dZ|_2^2,
\]
with $|dZ|_2^2>0$ if $Z\ne0$, and $\langle Z,\Delta_kZ\rangle\ge(39L^4)^{-1}\cdot|Z|_2^2$. Since $|C\mathsf{Z}|_2\ge|\mathsf{Z}|_2$ by Lemma~\ref{lem:face-averages-map}(b), the lower bound for $\langle\mathsf{Z},\tilde{\Delta}_k\mathsf{Z}\rangle$ follows. Thus $\tilde{\Delta}_k$ is a positive definite symmetric operator on the finite-dimensional space $V_k$; in particular it is invertible.

By the spectral theorem, $V_k$ has an orthonormal basis $(v_i)_i$ with $\tilde{\Delta}_kv_i=t_iv_i$, and $t_i\ge(39L^4)^{-1}>0$ by the lower bound. For $x\ge0$ one has $(\tilde{\Delta}_k+x)^{-1}v_i=(t_i+x)^{-1}v_i$. For every number $t>0$ the substitution $x=tu^2$ gives
\[
\pi^{-1}\int_0^{\infty}\mathrm{d}x\,x^{-\frac{1}{2}}(t+x)^{-1}=\pi^{-1}\int_0^{\infty}\mathrm{d}u\,\frac{2}{1+u^2}\cdot t^{-\frac{1}{2}}=t^{-\frac{1}{2}}.
\]
Therefore $C_k^{\frac{1}{2}}v_i=t_i^{-\frac{1}{2}}v_i$ for every $i$. It follows that $C_k^{\frac{1}{2}}$ is symmetric and positive definite, and that its square maps $v_i$ to $t_i^{-1}v_i=C_kv_i$; so $C_k^{\frac{1}{2}}$ is the positive definite symmetric square root of $C_k$.

(b) The operator $C_k^{\frac{1}{2}}$ is invertible, being positive definite, and $C$ is a linear bijection of $V_k$ onto $\mathring{Z}_k$ by Lemma~\ref{lem:face-averages-map}(b). Hence $\mathrm{Sq}_k$ is a linear bijection of $V_k$ onto $\mathring{Z}_k$, and the definition of $\mathring{Z}_k$ gives $Q\mathrm{Sq}_k\zeta=0$ and $\tau\mathrm{Sq}_k\zeta=0$ for $\zeta\in V_k$. By Proposition~\ref{prop:averaging-operators}(c),
\[
\operatorname{mean}(\mathrm{Sq}_k\zeta)=\operatorname{mean}(Q\mathrm{Sq}_k\zeta)=0.
\]

(c) By the proof of (a), $C_k^{\frac{1}{2}}\tilde{\Delta}_kC_k^{\frac{1}{2}}v_i=t_i^{-\frac{1}{2}}\cdot t_i\cdot t_i^{-\frac{1}{2}}\,v_i=v_i$ for every $i$; thus $C_k^{\frac{1}{2}}\tilde{\Delta}_kC_k^{\frac{1}{2}}=1$. As $C_k^{\frac{1}{2}}$ is symmetric,
\[
\langle\mathrm{Sq}_k\zeta,\Delta_k\mathrm{Sq}_k\zeta'\rangle=\left\langle C_k^{\frac{1}{2}}\zeta,C^T\Delta_kC\,C_k^{\frac{1}{2}}\zeta'\right\rangle=\left\langle\zeta,C_k^{\frac{1}{2}}\tilde{\Delta}_kC_k^{\frac{1}{2}}\zeta'\right\rangle=\langle\zeta,\zeta'\rangle.
\]
By (a) and (b), Proposition~\ref{prop:fluctuation-fields} follows.
\end{proof}

\subsubsection{The minimiser $A^{\min}_k$ and the function $\boldsymbol{\omega}_k$}\label{sec:minimiser-function}

The lattices of level $k+1$ are taken with the side $l/L$: $T^0_{k+1}=\mathbb{T}(1,l/L)$ and $T_{k+1}=\mathbb{T}(\eta_{k+1},l/L)$, so that $H_{k+1}$ is the map of the tower $(\eta_{k+1},k+1,l/L)$. For $A'\in C^1(T^0_{k+1})$ the scaled field $A'_L$ of~\eqref{eq:algebraic-gaussian-3} lies in $C^1(T^1_k)$. Let
\[
Q_{k+1}:=QQ_k\colon C^p(T_k)\to C^p(T^1_k)\qquad(p=0,1);
\]
this is the averaging operator of the tower $(\eta_k,k+1,l)$, whose fine lattice is $T_k$ and whose coarse lattice is $T^1_k$.

\begin{theorem}[Minimisation in one step]\label{thm:minimisation-one-step}
Let $k\ge0$, $l\in L\mathbb{N}$ and $A'\in C^1(T^0_{k+1})$; the harmonic part of $A'$ is arbitrary.
\begin{enumerate}
\item[(a)] Put
\[
\mathrm{Sa}(A'):=\left\{A\in C^1(T^0_k):\ QA=A'_L,\ \tau A=0\right\}.
\]
The function $A\mapsto\langle A,\Delta_kA\rangle$ has exactly one minimiser on $\mathrm{Sa}(A')$; it is denoted by $A^{\min}_k(A')$. It is characterised by the two conditions
\[
A^{\min}_k(A')\in\mathrm{Sa}(A')\qquad\text{and}\qquad\langle A^{\min}_k(A'),\Delta_kZ\rangle=0\ \text{ for all }Z\in C^1(T^0_k)\text{ with }QZ=0,\ \tau Z=0.
\]
The map $A'\mapsto A^{\min}_k(A')$ is linear and real. For every $Z\in\mathring{Z}_k$
\[
\left\langle A^{\min}_k(A')+Z,\Delta_k\left(A^{\min}_k(A')+Z\right)\right\rangle=\langle A^{\min}_k(A'),\Delta_kA^{\min}_k(A')\rangle+\langle Z,\Delta_kZ\rangle.
\]
\item[(b)] There is a unique function $\boldsymbol{\omega}_k(A')\in C^0(T_k)$ with $Q_{k+1}\boldsymbol{\omega}_k(A')=0$ for which
\begin{equation}\label{eq:algebraic-gaussian-14}
H_kA^{\min}_k(A')=(H_{k+1}A')_L-d\,\boldsymbol{\omega}_k(A');
\end{equation}
the map $A'\mapsto\boldsymbol{\omega}_k(A')$ is real and linear.
\item[(c)] $|dH_kA^{\min}_k(A')|_2=|dH_{k+1}A'|_2$.
\item[(d)] $\operatorname{mean}(A^{\min}_k(A'))=L^{-\frac{1}{2}}\cdot\operatorname{mean}(A')$. If $A'$ is closed and has the mean $c'$, then $A^{\min}_k(A')$ is closed and has the mean $L^{-\frac{1}{2}}\cdot c'$.
\end{enumerate}
\end{theorem}

\begin{proof}
(a) The set $\mathrm{Sa}(A')$ is not empty. Indeed, let $\boldsymbol{\Sigma}_1$ be the map of Proposition~\ref{prop:averaging-operators}(f) with $j=1$ and $a=1$, that is, for the one-step tower $(1,1,l)$, whose fine and coarse lattices are $T^0_k$ and $T^1_k$. By its definition, the field $\boldsymbol{\Sigma}_1(A'_L)$ can differ from $0$ only on face bonds; hence it vanishes on the inner bonds, $\tau\boldsymbol{\Sigma}_1(A'_L)=0$, and Proposition~\ref{prop:averaging-operators}(f) gives $Q\boldsymbol{\Sigma}_1(A'_L)=A'_L$. As $Q$ and $\tau$ are linear,
\[
\mathrm{Sa}(A')=\boldsymbol{\Sigma}_1(A'_L)+\mathring{Z}_k.
\]
Apply the lemma on minimisation on an affine subspace (Section~\ref{sec:quadratic-minimisation-block}) to the Euclidean spaces $E:=C^1(T^0_k)$ and $E':=C^2(T_k)$, the map $T:=dH_k$, the subspace $W:=\mathring{Z}_k$ and the point $x_0:=\boldsymbol{\Sigma}_1(A'_L)$. By the definition of $\Delta_k$, one has $|TA|^2=\langle A,\Delta_kA\rangle$ and $\langle TA,TZ\rangle=\langle A,\Delta_kZ\rangle$ for $A,Z\in E$. By Corollary~\ref{cor:axial-gauge-bounds}(d), $\langle Z,\Delta_kZ\rangle>0$ for every $Z\in\mathring{Z}_k\setminus\{0\}$, that is, $W\cap\ker T=\{0\}$. The lemma therefore yields exactly one minimiser, its characterisation as stated, and, through its identity~\eqref{eq:algebraic-gaussian-LA1}, the identity at the end of (a). Finally, the two characterising conditions are homogeneous linear equations, with real coefficients, in the pair $(A^{\min}_k(A'),A')$, and for given $A'$ their only solution is $A^{\min}_k(A')$. Hence $A^{\min}_k(A')$ depends linearly on $A'$, with real coefficients: the map $A'\mapsto A^{\min}_k(A')$ is linear and real.

(b), (c) Consider the problem of minimising $|d\mathscr{A}|_2^2$ over the affine space
\[
\left\{\mathscr{A}\in C^1(T_k):\ Q_{k+1}\mathscr{A}=A'_L\right\}
\]
of the tower $(\eta_k,k+1,l)$.

\emph{Step 1.} Proposition~\ref{prop:constrained-minimisers}(vi), applied to the towers $(\eta_{k+1},k+1,l/L)$ and $(\eta_k,k+1,l)$, gives $(H_{k+1}A')_L=H_{\eta_k,k+1}(A'_L)$, and by Proposition~\ref{prop:constrained-minimisers}(ii) this field is a minimiser of the problem. By~\eqref{eq:algebraic-gaussian-3} and the identities that follow it,
\[
|d(H_{k+1}A')_L|_2=|(dH_{k+1}A')_L|_2=|dH_{k+1}A'|_2.
\]

\emph{Step 2.} The field $\mathscr{B}:=H_kA^{\min}_k(A')\in C^1(T_k)$ is a minimiser of the problem as well. First, $Q_kH_k=1$ by Proposition~\ref{prop:constrained-minimisers}(ii); so $Q_{k+1}\mathscr{B}=QA^{\min}_k(A')=A'_L$. Next, let $\mathscr{A}\in C^1(T_k)$ be any field with $Q_{k+1}\mathscr{A}=A'_L$, and put $A:=Q_k\mathscr{A}$; then $QA=A'_L$. Lemma~\ref{lem:axial-representatives}, applied to $A$ and to the tower $(1,1,l)$, gives a function $\mu\in C^0(T^0_k)$ with $Q\mu=0$ such that $A-d\mu$ is axial, that is, $\tau(A-d\mu)=0$. By Proposition~\ref{prop:averaging-operators}(b),
\[
Q(A-d\mu)=A'_L-d\,Q\mu=A'_L,
\]
so that $A-d\mu\in\mathrm{Sa}(A')$. Consequently
\[
|d\mathscr{A}|_2^2\ge|dH_kA|_2^2=|dH_k(A-d\mu)|_2^2\ge|dH_kA^{\min}_k(A')|_2^2=|d\mathscr{B}|_2^2.
\]
The first inequality holds because $H_kA$ minimises $|d\mathscr{A}|_2^2$ under the constraint $Q_k\mathscr{A}=A$ (Proposition~\ref{prop:constrained-minimisers}(i), (ii)). The equality holds because $H_k$ is linear and $dH_k\,d\mu=0$, the field $H_k\,d\mu$ being exact by Proposition~\ref{prop:constrained-minimisers}(iii). The last inequality follows from (a), because $|dH_kX|_2^2=\langle X,\Delta_kX\rangle$ for every $X\in C^1(T^0_k)$.

\emph{Step 3.} By Proposition~\ref{prop:constrained-minimisers}(i) for the tower $(\eta_k,k+1,l)$, all minimisers of the problem give the same value of $|d\mathscr{A}|_2$, and they form one restricted orbit. The first fact and Step 1 give
\[
|dH_kA^{\min}_k(A')|_2=|d(H_{k+1}A')_L|_2=|dH_{k+1}A'|_2,
\]
which is (c). The second fact gives
\[
(H_{k+1}A')_L-\mathscr{B}=d\,\boldsymbol{\omega}
\]
with some $\boldsymbol{\omega}\in\mathcal{N}_{k+1}$, that is, with $\boldsymbol{\omega}\in C^0(T_k)$ and $Q_{k+1}\boldsymbol{\omega}=0$; this is~\eqref{eq:algebraic-gaussian-14} with $\boldsymbol{\omega}_k(A'):=\boldsymbol{\omega}$. The field $d\,\boldsymbol{\omega}$ is closed, and $Q_{k+1}d\,\boldsymbol{\omega}=d\,Q_{k+1}\boldsymbol{\omega}=0$ by Proposition~\ref{prop:averaging-operators}(b); by Lemma~\ref{lem:closed-fields}(b), $\boldsymbol{\omega}$ is the only element of $\mathcal{N}_{k+1}$ whose coboundary is $d\,\boldsymbol{\omega}$. This proves the uniqueness in (b), and it shows that $d$ is injective on $\mathcal{N}_{k+1}$. Both minimisers $(H_{k+1}A')_L$ and $\mathscr{B}$ are linear and real in $A'$, by Proposition~\ref{prop:constrained-minimisers}(ii) and by (a). Hence $\boldsymbol{\omega}_k(A')$, the preimage of their difference under the injective linear map $d$ on $\mathcal{N}_{k+1}$, is real and linear in $A'$.

(d) By Proposition~\ref{prop:averaging-operators}(c) and by the formula for $\operatorname{mean}(A_L)$ that follows~\eqref{eq:algebraic-gaussian-3},
\[
\operatorname{mean}(A^{\min}_k(A'))=\operatorname{mean}(QA^{\min}_k(A'))=\operatorname{mean}(A'_L)=L^{-\frac{1}{2}}\cdot\operatorname{mean}(A').
\]
Let $dA'=0$. Then $dH_{k+1}A'=0$ by Proposition~\ref{prop:constrained-minimisers}(iii); hence $dH_kA^{\min}_k(A')=0$ by (c), and $dA^{\min}_k(A')=0$ by Proposition~\ref{prop:constrained-minimisers}(iii) again. Together with the formula for the mean this proves (d).

Theorem~\ref{thm:minimisation-one-step} contains all assertions of Proposition~\ref{prop:minimisers-prescribed-block}; Proposition~\ref{prop:minimisers-prescribed-block} follows.
\end{proof}

\subsection{The tower of block fields}\label{sec:tower-block-fields}

\subsubsection{Coordinates of axial fields}\label{sec:coordinates-axial-fields}

Let $k\ge1$ be an integer and let $n\in L^k\mathbb{N}$. For $0\le j\le k$ put
\[
l_j:=\frac{n}{L^j},\qquad n_j:=l_j^3,
\]
and let $\Lambda_n:=\mathbb{T}(1,n)$, so that $\Lambda_n=T^0_0$. The level $j$ is taken with the side $l_j$: then $T_j=\mathbb{T}(\eta_j,l_j)$ and $T^0_j=\mathbb{T}(1,l_j)$ (Section~\ref{sec:blocks-levels-rectilinear}), and $H_j$, $\Delta_j$ are the operators of Section~\ref{sec:minimisation-averaging-constraint} for the tower $(\eta_j,j,l_j)$. As $n\in L^k\mathbb{N}$, the side $l_j$ lies in $L\mathbb{N}$ for $j<k$; for such $j$ the spaces $V_j$, $\mathring{Z}_j$ and the maps $\mathrm{Sq}_j$, $A^{\min}_j$ of Section~\ref{sec:one-averaging-step} are those of the level $j$ and the side $l_j$, and $A^{\min}_j$ maps $C^1(T^0_{j+1})$ into $C^1(T^0_j)$, because $l_{j+1}=\frac{l_j}{L}$. The space $V_j$ carries the inner product of $C^1(T^0_j)$, and $|\zeta_j|$ denotes the norm of $\zeta_j\in V_j$. Under hypothesis~\ref{def:hypothesis-flow} one takes $n=\frac{R}{\eta}$, where $\eta=L^{-N}$ is the number of Section~\ref{sec:lattice-measures-laws}, and $k=k^{*}(R)$. By Section~\ref{sec:running-quantities-tori} and Definition~\ref{def:tori-stop-level} these values satisfy $k\ge1$ and $n\in L^k\mathbb{N}$: one has $k^{*}\ge1$, and the side $l_{*}=R\cdot L^{N-k^{*}}$ of the stop level is the positive integer $M\cdot r_{*}$, so that $n=R\cdot L^N=L^{k^{*}}\cdot l_{*}$.

For $A_0\in C^1(\Lambda_n)$ put
\begin{equation}\label{eq:algebraic-gaussian-15}
A_j:=\left(Q_jA_0\right)^{[j]}\in C^1(T^0_j)\quad(0\le j\le k);\qquad\text{then}\qquad\left(A_{j+1}\right)_L=Q\,A_j\quad(0\le j<k).
\end{equation}
Here $Q_jA_0\in C^1(L^j,n)$ is the average of Section~\ref{sec:averaging-operators}, and $B\mapsto B^{[j]}$, the inverse of the $j$-fold scaling (Section~\ref{sec:scaling}), maps $C^1(L^j,n)$ onto $C^1(1,l_j)=C^1(T^0_j)$. The identity in~\eqref{eq:algebraic-gaussian-15} is verified as follows. Let $j<k$. For a bond function $B$ on $\mathbb{T}(L^{j+1},n)$ and a bond $b$ of $\mathbb{T}(L,l_j)$ the scaling~\eqref{eq:algebraic-gaussian-3} gives
\[
\left(B^{[j+1]}\right)_L(b)=L^{-\frac{1}{2}}\cdot L^{\frac{j+1}{2}}\cdot B\left(L^jb\right)=B^{[j]}(b).
\]
Proposition~\ref{prop:averaging-operators}(g), applied $j$ times (its hypothesis is satisfied, since $n\in L^{j+1}\mathbb{N}$), shows that $Q\left(f^{[j]}\right)=(Qf)^{[j]}$ for every bond function $f$ on $\mathbb{T}(L^j,n)$. With $B=Q_{j+1}A_0=Q\,Q_jA_0$ it follows that
\[
\left(A_{j+1}\right)_L=\left(\left(Q_{j+1}A_0\right)^{[j+1]}\right)_L=\left(Q\,Q_jA_0\right)^{[j]}=Q\left(\left(Q_jA_0\right)^{[j]}\right)=Q\,A_j.
\]

Let
\[
\mathrm{Ax}_k:=\left\{A_0\in C^1(\Lambda_n):\tau A_j=0\ \text{ for }0\le j<k\right\},
\]
the condition $\tau A_j=0$ being understood on $T^0_j$ in the sense of Section~\ref{sec:fluctuation-spaces-map}. For $A\in C^1(T^0_k)$ and $\zeta=(\zeta_j)_{j<k}$ with $\zeta_j\in V_j$ define bond functions $A^{\downarrow}_j=A^{\downarrow}_j(A,\zeta)\in C^1(T^0_j)$ by the downward recursion
\[
A^{\downarrow}_k:=A,\qquad A^{\downarrow}_j:=A^{\min}_j\left(A^{\downarrow}_{j+1}\right)+\mathrm{Sq}_j\zeta_j\qquad(j=k-1,\dots,0),
\]
and let $A^{\downarrow}_0(A,\zeta)$ be the last of these fields. For $k=k^{*}$ this is the map $A^{\downarrow}_0$ of Definition~\ref{def:coordinates-slice-small} (Section~\ref{sec:averaging-fluctuations-gauge}).

\begin{theorem}[Hierarchical coordinates]\label{thm:hierarchical-coordinates}
\begin{enumerate}
\item[(a)] The map $(A,\zeta)\mapsto A^{\downarrow}_0(A,\zeta)$ is a linear bijection of $C^1(T^0_k)\times\prod_{j<k}V_j$ onto $\mathrm{Ax}_k$. If $A_0=A^{\downarrow}_0(A,\zeta)$, then the fields~\eqref{eq:algebraic-gaussian-15} formed with $A_0$ are $A_j=A^{\downarrow}_j(A,\zeta)$ for $0\le j\le k$; in particular $A=A_k$, and $\mathrm{Sq}_j\zeta_j=A_j-A^{\min}_j(A_{j+1})$ for $0\le j<k$.
\item[(b)] If $A_0=A^{\downarrow}_0(A,\zeta)$, then for $0\le j\le k$
\[
|dA_0|_2^2=|dH_jA_j|_2^2+\sum_{i<j}|\zeta_i|^2.
\]
\end{enumerate}
\end{theorem}

\begin{proof}
(a) The maps $A^{\min}_j$ and $\mathrm{Sq}_j$ are linear (Theorem~\ref{thm:minimisation-one-step}(a), Proposition~\ref{prop:fluctuation-operator}(b)); hence every $A^{\downarrow}_j(A,\zeta)$ depends linearly on $(A,\zeta)$. By Theorem~\ref{thm:minimisation-one-step}(a), for $A'\in C^1(T^0_{j+1})$ the field $A^{\min}_j(A')$ belongs to the set $\mathrm{Sa}(A')$ of that theorem, that is, $Q\,A^{\min}_j(A')=A'_L$ and $\tau\,A^{\min}_j(A')=0$; by Proposition~\ref{prop:fluctuation-operator}(b), $Q\,\mathrm{Sq}_j=0$ and $\tau\,\mathrm{Sq}_j=0$.

Let $A_0:=A^{\downarrow}_0(A,\zeta)$, and let $A_j$ be the fields~\eqref{eq:algebraic-gaussian-15} formed with it. The identities $A_j=A^{\downarrow}_j$ are proved by induction on $j$. For $j=0$ the identity is the definition of $A_0$, as $Q_0=1$ and $B^{[0]}=B$. Let $j<k$ and $A_j=A^{\downarrow}_j$. By~\eqref{eq:algebraic-gaussian-15}
\[
\left(A_{j+1}\right)_L=Q\,A_j=Q\,A^{\min}_j\left(A^{\downarrow}_{j+1}\right)+Q\,\mathrm{Sq}_j\zeta_j=\left(A^{\downarrow}_{j+1}\right)_L,
\]
and $A_{j+1}=A^{\downarrow}_{j+1}$, because the scaling~\eqref{eq:algebraic-gaussian-3} is injective. Moreover
\[
\tau A_j=\tau\,A^{\min}_j\left(A^{\downarrow}_{j+1}\right)+\tau\,\mathrm{Sq}_j\zeta_j=0.
\]
Consequently $A_0\in\mathrm{Ax}_k$, $A_k=A^{\downarrow}_k=A$, and $\mathrm{Sq}_j\zeta_j=A^{\downarrow}_j-A^{\min}_j(A^{\downarrow}_{j+1})=A_j-A^{\min}_j(A_{j+1})$ for $j<k$. Since $\mathrm{Sq}_j$ is injective (Proposition~\ref{prop:fluctuation-operator}(b)), the pair $(A,\zeta)$ is determined by $A_0$, and the map of (a) is injective.

It remains to show that every $A_0\in\mathrm{Ax}_k$ is attained. Let $A_j$ be the fields~\eqref{eq:algebraic-gaussian-15} formed with $A_0$, and for $j<k$ put $Z_j:=A_j-A^{\min}_j(A_{j+1})\in C^1(T^0_j)$. By~\eqref{eq:algebraic-gaussian-15} and Theorem~\ref{thm:minimisation-one-step}(a),
\[
Q\,Z_j=\left(A_{j+1}\right)_L-\left(A_{j+1}\right)_L=0,\qquad\tau Z_j=\tau A_j-\tau\,A^{\min}_j(A_{j+1})=0,
\]
the latter because $A_0\in\mathrm{Ax}_k$. Thus $Z_j\in\mathring{Z}_j$, and by Proposition~\ref{prop:fluctuation-operator}(b) there is $\zeta_j\in V_j$ with $Z_j=\mathrm{Sq}_j\zeta_j$. With $\zeta:=(\zeta_j)_{j<k}$ one has $A^{\downarrow}_j(A_k,\zeta)=A_j$ for $j=k,k-1,\dots,0$, by downward induction on $j$: $A^{\downarrow}_k(A_k,\zeta)=A_k$, and if $A^{\downarrow}_{j+1}(A_k,\zeta)=A_{j+1}$, then $A^{\downarrow}_j(A_k,\zeta)=A^{\min}_j(A_{j+1})+Z_j=A_j$. In particular $A^{\downarrow}_0(A_k,\zeta)=A_0$.

(b) Let $0\le i<k$. By (a), $A_i=A^{\min}_i(A_{i+1})+\mathrm{Sq}_i\zeta_i$ with $\mathrm{Sq}_i\zeta_i\in\mathring{Z}_i$. The identity $\langle A,\Delta_iA\rangle=|dH_iA|_2^2$ (Section~\ref{sec:minimisation-averaging-constraint}), the last assertion of Theorem~\ref{thm:minimisation-one-step}(a), Theorem~\ref{thm:minimisation-one-step}(c) and Proposition~\ref{prop:fluctuation-operator}(c) give
\begin{align*}
|dH_iA_i|_2^2&=\langle A_i,\Delta_iA_i\rangle=\left\langle A^{\min}_i(A_{i+1}),\Delta_i\,A^{\min}_i(A_{i+1})\right\rangle+\left\langle\mathrm{Sq}_i\zeta_i,\Delta_i\,\mathrm{Sq}_i\zeta_i\right\rangle\\
&=\left|dH_i\,A^{\min}_i(A_{i+1})\right|_2^2+|\zeta_i|^2=|dH_{i+1}A_{i+1}|_2^2+|\zeta_i|^2.
\end{align*}
For $0\le j\le k$, the sum of these identities over $i=0,\dots,j-1$ is $|dH_0A_0|_2^2=|dH_jA_j|_2^2+\sum_{i<j}|\zeta_i|^2$, and $H_0=1$ (Proposition~\ref{prop:constrained-minimisers}(ii)).
\end{proof}

\subsubsection{The slice and the space of curvatures}\label{sec:slice-space-curvatures}

Let $F_N:=dC^1(\Lambda_n)\subset C^2(\Lambda_n)$, with the inner product $\langle\cdot,\cdot\rangle$ and the norm $|\cdot|=|\cdot|_2$ of $C^2(\Lambda_n)$, and let $\gamma_N$ be the standard Gaussian measure of the Euclidean space $F_N$ (Section~\ref{sec:lattice-measures-laws}). Let $V_{\mathrm{t}}(T^0_k)$ be the space $V_{\mathrm{t}}$ of Section~\ref{sec:tori-cochains-hodge} for the torus $T^0_k=\mathbb{T}(1,l_k)$, that is, the space of the coclosed bond functions on $T^0_k$ with mean zero. Put
\[
\mathrm{Sl}:=\left\{A_0\in\mathrm{Ax}_k:A_k\in V_{\mathrm{t}}(T^0_k)\right\},
\]
and for $A_0\in\mathrm{Sl}$ let $A_{*}:=A_k$. Under hypothesis~\ref{def:hypothesis-flow} one takes $k=k^{*}$.

For $k=k^{*}$ the set $\mathrm{Sl}$ is the slice of Definition~\ref{def:coordinates-slice-small}, which is defined there by the conditions $\tau Q_jA_0=0$ for $0\le j<k^{*}$ together with the requirement that $Q_{k^{*}}A_0$ be coclosed with mean zero. The two descriptions agree. Let $0\le j<k$. By the definition of $B\mapsto B^{[j]}$ in Section~\ref{sec:scaling}, the value of $A_j=(Q_jA_0)^{[j]}$ on a bond $b$ of $T^0_j$ is $L^{\frac{j}{2}}$ times the value of $Q_jA_0$ on the dilated bond $L^jb$; and~\eqref{eq:algebraic-gaussian-7} forms sums of values on the bonds of paths, with coefficients that do not involve the lattice spacing. The dilation $x\mapsto L^jx$ carries the path $\boldsymbol{\Gamma}^{\pi}(y,x)$ of $\mathbb{L}(1)$ onto the path $\boldsymbol{\Gamma}^{\pi}(L^jy,L^jx)$ of $\mathbb{L}(L^j)$; hence
\[
(\tau A_j)(y,x)=L^{\frac{j}{2}}\cdot\left(\tau Q_jA_0\right)\left(L^jy,L^jx\right)\qquad\left(y\in L\mathbb{Z}^3,\ x\in B(y)\right).
\]
If $y$ ranges over $L\mathbb{Z}^3$ and $x$ over $B(y)$, then $(L^jy,L^jx)$ ranges over all the pairs at which $\tau Q_jA_0$ is defined; thus $\tau A_j=0$ holds exactly when $\tau Q_jA_0=0$. For the remaining condition, the definitions of Section~\ref{sec:tori-cochains-hodge} give
\[
\left(\delta B^{[k]}\right)(x)=L^{\frac{3k}{2}}\cdot(\delta B)\left(L^kx\right),\qquad\operatorname{mean}\left(B^{[k]}\right)=L^{\frac{k}{2}}\cdot\operatorname{mean}(B)
\]
for every bond function $B$ on $\mathbb{T}(L^k,n)$ and every site $x$ of $T^0_k$. The factors $L^{\frac{3k}{2}}$ and $L^{\frac{k}{2}}$ do not vanish, and $x\mapsto L^kx$ maps the sites of $T^0_k$ bijectively onto those of $\mathbb{T}(L^k,n)$; applied to $B=Q_kA_0$, this shows that $A_k$ is coclosed with mean zero if and only if $Q_kA_0$ is.

For $j<k$ let $\{\Upsilon\}$ be the basis of $V_j$ of Section~\ref{sec:averaging-fluctuations-gauge}; it is an orthonormal basis of $V_j$ (Section~\ref{sec:fluctuation-spaces-map}), and its existence is Proposition~\ref{prop:fluctuation-space}(i). For $\zeta_j\in V_j$ and an element $\Upsilon$ of this basis write $\zeta_j(\Upsilon):=\langle\zeta_j,\Upsilon\rangle$.

\begin{theorem}[Hierarchical coordinates of curvatures]\label{thm:hierarchical-coordinates-curvatures}
\begin{enumerate}
\item[(a)] The coboundary $d$ maps $\mathrm{Sl}$ isomorphically onto $F_N$.
\item[(b)] By (a), every $F\in F_N$ has exactly one preimage $A_0\in\mathrm{Sl}$ under $d$. By Theorem~\ref{thm:hierarchical-coordinates}(a), this $A_0$ can be written in exactly one way as $A_0=A^{\downarrow}_0(A_{*},\zeta)$ with $A_{*}\in C^1(T^0_k)$ and $\zeta=(\zeta_0,\dots,\zeta_{k-1})\in\prod_{j<k}V_j$, and then $A_{*}=A_k$. The map
\[
\mathsf{J}:F\mapsto(A_{*};\zeta)
\]
is a linear bijection of $F_N$ onto $V_{\mathrm{t}}(T^0_k)\times\prod_{j<k}V_j$, its inverse is $(A_{*};\zeta)\mapsto d\,A^{\downarrow}_0(A_{*},\zeta)$, and
\begin{equation}\label{eq:algebraic-gaussian-16}
|F|^2=|dH_kA_{*}|_2^2+\sum_{j<k}|\zeta_j|^2.
\end{equation}
\item[(c)] Through $\mathsf{J}$, the fields $A_{*}$ and $\zeta_j$ are linear functions of $F\in F_N$. There is a family $(\phi_{j,\Upsilon})$ in $F_N$, indexed by the pairs $(j,\Upsilon)$ with $j<k$ and $\Upsilon$ an element of the basis $\{\Upsilon\}$ of $V_j$, which is orthonormal and satisfies $\zeta_j(\Upsilon)=\langle\phi_{j,\Upsilon},F\rangle$ for all $F\in F_N$. For every linear functional $f$ on $V_{\mathrm{t}}(T^0_k)$ there is $\phi\in F_N$, orthogonal to all $\phi_{j,\Upsilon}$, with $f(A_{*})=\langle\phi,F\rangle$ for all $F\in F_N$.
\item[(d)] For every plaquette $p$ of $T^0_k$ there is $\phi_p\in F_N$ with $|\phi_p|\le1=C_Q$ and $(dA_{*})(p)=\langle\phi_p,F\rangle$ for all $F\in F_N$.
\item[(e)] Consider $A_{*}$ and the numbers $\zeta_j(\Upsilon)$, $j<k$, $\Upsilon$ an element of the basis $\{\Upsilon\}$ of $V_j$, as random variables on the probability space $(F_N,\gamma_N)$. The numbers $\zeta_j(\Upsilon)$ are mutually independent, each of them is standard normal, and the family that they form is independent of $A_{*}$. Let $\mathrm{d}\mathsf{m}_{*}$ be the distribution of $A_{*}$. It is a centred Gaussian measure on $V_{\mathrm{t}}(T^0_k)$, and its density with respect to the Lebesgue measure of $V_{\mathrm{t}}(T^0_k)$ is a constant multiple of the function $A\mapsto\exp\left(-\frac{1}{2}|dH_kA|_2^2\right)$.
\end{enumerate}
\end{theorem}

\begin{proof}
(a) By Theorem~\ref{thm:hierarchical-coordinates}(a), $A^{\downarrow}_0$ is a linear bijection of $C^1(T^0_k)\times\prod_{j<k}V_j$ onto $\mathrm{Ax}_k$, and $A_k=A$ for $A_0=A^{\downarrow}_0(A,\zeta)$. Hence $\mathrm{Sl}$ is the image of the subspace $V_{\mathrm{t}}(T^0_k)\times\prod_{j<k}V_j$ under $A^{\downarrow}_0$, and
\[
\dim\mathrm{Sl}=\dim V_{\mathrm{t}}(T^0_k)+\sum_{j<k}\dim V_j=2n_k-2+\sum_{j<k}\left(2n_j-2n_{j+1}\right)=2n_0-2=\dim F_N.
\]
Here $\dim V_{\mathrm{t}}(T^0_k)=2l_k^3-2=2n_k-2$ and $\dim F_N=\dim dC^1(1,n)=2n^3-2=2n_0-2$ by Section~\ref{sec:tori-cochains-hodge}; moreover $\dim V_j=d_V\cdot n_j=\left(2-2L^{-3}\right)\cdot n_j=2n_j-2n_{j+1}$ by Lemma~\ref{lem:face-averages-map}(c), as $n_{j+1}=L^{-3}\cdot n_j$, and the sum telescopes. The coboundary maps $\mathrm{Sl}\subset C^1(\Lambda_n)$ into $F_N$; since the dimensions agree, it suffices to show that $d$ is injective on $\mathrm{Sl}$. Let $A_0\in\mathrm{Sl}$ with $dA_0=0$, and write $A_0=A^{\downarrow}_0(A_k,\zeta)$ (Theorem~\ref{thm:hierarchical-coordinates}(a)). Theorem~\ref{thm:hierarchical-coordinates}(b) with $j=k$ gives
\[
0=|dH_kA_k|_2^2+\sum_{j<k}|\zeta_j|^2;
\]
hence $\zeta=0$ and $dH_kA_k=0$. By Proposition~\ref{prop:constrained-minimisers}(iii), $dA_k=0$, and $A_k=0$, because $A_k\in V_{\mathrm{t}}(T^0_k)$ and $d$ is injective on $V_{\mathrm{t}}$ (by~\eqref{eq:algebraic-gaussian-Hodge}; Section~\ref{sec:tori-cochains-hodge}). Then $A_0=A^{\downarrow}_0(0,0)=0$.

(b) By (a), the potential $A_0\in\mathrm{Sl}$ of $F$ exists, is unique and depends linearly on $F$. By Theorem~\ref{thm:hierarchical-coordinates}(a), $A_0=A^{\downarrow}_0(A_{*},\zeta)$ for exactly one pair $(A_{*},\zeta)\in C^1(T^0_k)\times\prod_{j<k}V_j$, which depends linearly on $A_0$, and $A_{*}=A_k\in V_{\mathrm{t}}(T^0_k)$, because $A_0\in\mathrm{Sl}$. So $\mathsf{J}$ is a linear map of $F_N$ into $V_{\mathrm{t}}(T^0_k)\times\prod_{j<k}V_j$. Conversely, let $(A_{*};\zeta)$ belong to this space, and put $A_0:=A^{\downarrow}_0(A_{*},\zeta)$. By Theorem~\ref{thm:hierarchical-coordinates}(a), $A_0\in\mathrm{Ax}_k$ and $A_k=A_{*}\in V_{\mathrm{t}}(T^0_k)$; so $A_0\in\mathrm{Sl}$ is the potential of $F:=dA_0\in F_N$, and $\mathsf{J} F=(A_{*};\zeta)$. Hence $\mathsf{J}$ is bijective, and its inverse is $(A_{*};\zeta)\mapsto d\,A^{\downarrow}_0(A_{*},\zeta)$. Finally,~\eqref{eq:algebraic-gaussian-16} is Theorem~\ref{thm:hierarchical-coordinates}(b) with $j=k$, since $|F|=|dA_0|_2$.

(c) On the space $\mathscr{W}:=V_{\mathrm{t}}(T^0_k)\times\prod_{j<k}V_j$ put
\[
\left\langle(A,\zeta),(A',\zeta')\right\rangle_W:=\langle A,\Delta_kA'\rangle+\sum_{j<k}\langle\zeta_j,\zeta'_j\rangle.
\]
This bilinear form is symmetric, since $\Delta_k=(dH_k)^TdH_k$, and positive definite, since for $0\ne A\in V_{\mathrm{t}}(T^0_k)$
\[
\langle A,\Delta_kA\rangle=|dH_kA|_2^2\ge|dA|_2^2>0
\]
by Proposition~\ref{prop:constrained-minimisers}(v) and the injectivity of $d$ on $V_{\mathrm{t}}$. By~\eqref{eq:algebraic-gaussian-16}, $|F|^2=\langle\mathsf{J} F,\mathsf{J} F\rangle_W$ for $F\in F_N$. As $\mathsf{J}$ is a linear bijection by (b), it is an isometry of $F_N$ onto the Euclidean space $(\mathscr{W},\langle\cdot,\cdot\rangle_W)$, and by polarisation
\[
\left\langle\mathsf{J} F,\mathsf{J} F'\right\rangle_W=\langle F,F'\rangle\qquad(F,F'\in F_N).
\]
For $j<k$ and an element $\Upsilon$ of the basis $\{\Upsilon\}$ of $V_j$ let $\mathbf{e}_{j,\Upsilon}\in\mathscr{W}$ be the element all of whose entries vanish, except the entry $\zeta_j$, which equals $\Upsilon$; put $\phi_{j,\Upsilon}:=\mathsf{J}^{-1}\mathbf{e}_{j,\Upsilon}\in F_N$. Since each basis $\{\Upsilon\}$ is orthonormal, the family $(\mathbf{e}_{j,\Upsilon})$ is orthonormal in $\mathscr{W}$, and therefore the family $(\phi_{j,\Upsilon})$ is orthonormal in $F_N$. For $F\in F_N$ with $\mathsf{J} F=(A_{*};\zeta)$ one has
\[
\zeta_j(\Upsilon)=\langle\zeta_j,\Upsilon\rangle=\left\langle\mathbf{e}_{j,\Upsilon},\mathsf{J} F\right\rangle_W=\left\langle\phi_{j,\Upsilon},F\right\rangle.
\]
Let $f$ be a linear functional on $V_{\mathrm{t}}(T^0_k)$. As $(A,A')\mapsto\langle A,\Delta_kA'\rangle$ is a scalar product on $V_{\mathrm{t}}(T^0_k)$, there is $B\in V_{\mathrm{t}}(T^0_k)$ with $f(A)=\langle B,\Delta_kA\rangle$ for all $A\in V_{\mathrm{t}}(T^0_k)$. Then
\[
f(A_{*})=\left\langle(B,0),\mathsf{J} F\right\rangle_W=\left\langle\mathsf{J}^{-1}(B,0),F\right\rangle\qquad(F\in F_N),
\]
and $\phi:=\mathsf{J}^{-1}(B,0)\in F_N$ is orthogonal to all $\phi_{j,\Upsilon}$, because $(B,0)$ is orthogonal to all $\mathbf{e}_{j,\Upsilon}$ in $\mathscr{W}$.

(d) Let $F\in F_N$. By Proposition~\ref{prop:constrained-minimisers}(ii), $Q_kH_k=1$, and by Proposition~\ref{prop:averaging-operators}(b), $d\,Q_k=Q_k^{(2)}d$; hence
\[
dA_{*}=d\,Q_kH_kA_{*}=Q_k^{(2)}dH_kA_{*}.
\]
On $T^0_k$ every plaquette carries the weight $1$ in the norm $|\cdot|_2$. For every plaquette $p$ of $T^0_k$, Proposition~\ref{prop:averaging-operators}(d) and~\eqref{eq:algebraic-gaussian-16} therefore give
\[
\left|(dA_{*})(p)\right|\le|dA_{*}|_2=\left|Q_k^{(2)}dH_kA_{*}\right|_2\le|dH_kA_{*}|_2\le|F|.
\]
So $F\mapsto(dA_{*})(p)$ is a linear functional of norm at most $1$ on the finite-dimensional Euclidean space $F_N$; it is $F\mapsto\langle\phi_p,F\rangle$ for exactly one $\phi_p\in F_N$, and $|\phi_p|\le1$. By Proposition~\ref{prop:averaging-operators}(d), $C_Q=1$.

(e) By the proof of (c), $\mathsf{J}$ is a linear isometry of $F_N$ onto the Euclidean space $(\mathscr{W},\langle\cdot,\cdot\rangle_W)$. On a Euclidean space $E$ the standard Gaussian measure has the density $v\mapsto(2\pi)^{-\frac{\dim E}{2}}\exp\left(-\frac{1}{2}|v|^2\right)$ with respect to the Lebesgue measure that assigns the volume $1$ to the cubes spanned by orthonormal bases of $E$. A linear isometry of $E$ onto a second Euclidean space preserves the norm and takes orthonormal bases to orthonormal bases, so that it also preserves this Lebesgue measure; consequently it transforms the standard Gaussian measure of $E$ into that of the second space. For $\mathsf{J}$ this says that $\gamma_N$ is carried to the standard Gaussian measure of $(\mathscr{W},\langle\cdot,\cdot\rangle_W)$. With this scalar product, $\mathscr{W}$ is the orthogonal sum of $V_{\mathrm{t}}(T^0_k)$, equipped with $\langle A,\Delta_kA'\rangle$, and of the spaces $V_j$, $j<k$. On an orthogonal sum one has $|w|^2=\sum_i|w_i|^2$ for every element $w$ with the components $w_i$, and the Lebesgue measure just described is the product of the corresponding measures of the summands; the density therefore factorises, and the standard Gaussian measure of $\mathscr{W}$ equals the product, over the summands, of their own standard Gaussian measures. It follows that $A_{*},\zeta_0,\dots,\zeta_{k-1}$ are independent random variables on $(F_N,\gamma_N)$. Each $\zeta_j$ is distributed according to the standard Gaussian measure of $V_j$; in the coordinates with respect to the orthonormal basis $\{\Upsilon\}$ this measure is the product of standard normal laws, one for each coordinate, so that the numbers $\zeta_j(\Upsilon)$ with fixed $j$ are independent and standard normal. Together with the independence of $A_{*},\zeta_0,\dots,\zeta_{k-1}$ this gives the assertions of (e) on the numbers $\zeta_j(\Upsilon)$. Finally, $A_{*}$ is distributed according to the standard Gaussian measure of $V_{\mathrm{t}}(T^0_k)$ for the scalar product $\langle A,\Delta_kA'\rangle$, that is, according to the centred Gaussian measure whose density with respect to the Lebesgue measure of $V_{\mathrm{t}}(T^0_k)$ is a constant multiple of
\[
\exp\left(-\frac{1}{2}\langle A,\Delta_kA\rangle\right)=\exp\left(-\frac{1}{2}|dH_kA|_2^2\right);
\]
indeed, the Lebesgue measures that two scalar products define on the same space differ by a constant factor.
\end{proof}

For $k=k^{*}$ under hypothesis~\ref{def:hypothesis-flow}, the set $\mathrm{Sl}$, the field $A_{*}$, the fields $\zeta_j$, which by Theorem~\ref{thm:hierarchical-coordinates}(a) satisfy $A_j=A^{\min}_j(A_{j+1})+\mathrm{Sq}_j\zeta_j$, and the law $\mathrm{d}\mathsf{m}_{*}$ are those of Definition~\ref{def:coordinates-slice-small}, and Proposition~\ref{prop:orthogonal-coordinates} follows.

\begin{corollary}[Adding a closed field at level $k$]\label{cor:adding-closed-field}
Let $A^c\in C^1(T^0_k)$ be closed with mean $c\in\mathbb{R}^3$. Then for every $A\in C^1(T^0_k)$ and every $\zeta=(\zeta_j)_{j<k}$ with $\zeta_j\in V_j$
\[
A^{\downarrow}_0(A+A^c,\zeta)=A^{\downarrow}_0(A,\zeta)+L^{-\frac{k}{2}}\cdot c+d\boldsymbol{\lambda}
\]
with a function $\boldsymbol{\lambda}$ on $\Lambda_n$; here $L^{-\frac{k}{2}}\cdot c$ is the constant bond function on $\Lambda_n$ with the components $L^{-\frac{k}{2}}\cdot c_{\mu}$. In particular, if hypothesis~\ref{def:hypothesis-flow} holds and $k=k^{*}$, then for every $\theta=(\theta_0,\theta_1,\theta_2)\in\mathbb{R}^3$ and all $A$, $\zeta$ as above
\[
A^{\downarrow}_0(A+c_{\theta},\zeta)=A^{\downarrow}_0(A,\zeta)+H_{\theta}+d\boldsymbol{\lambda}
\]
with a function $\boldsymbol{\lambda}$ on $\Lambda_n$; here $c_{\theta}:=c^{(k^{*})}_{\theta}$ is the twist of Definition~\ref{def:twist-external-connection} and $H_{\theta}$ is the constant bond function of Section~\ref{sec:lattice-measures-laws}, both formed with this $\theta$.
\end{corollary}

\begin{proof}
By Theorem~\ref{thm:hierarchical-coordinates}(a) the map $A^{\downarrow}_0$ is linear, and with $\zeta=0$ its recursion reads $A^{\downarrow}_j=A^{\min}_j(A^{\downarrow}_{j+1})$. Hence
\[
A^{\downarrow}_0(A+A^c,\zeta)-A^{\downarrow}_0(A,\zeta)=A^{\downarrow}_0(A^c,0)=A^{\min}_0\left(A^{\min}_1\left(\cdots A^{\min}_{k-1}(A^c)\cdots\right)\right).
\]
By Theorem~\ref{thm:minimisation-one-step}(d), if $A'\in C^1(T^0_{j+1})$ is closed with mean $c'$, then $A^{\min}_j(A')$ is closed with mean $L^{-\frac{1}{2}}\cdot c'$. Applying this $k$ times, at the levels $j=k-1,\dots,0$, one finds that $A^{\downarrow}_0(A^c,0)\in C^1(\Lambda_n)$ is closed with mean $L^{-\frac{k}{2}}\cdot c$. By~\eqref{eq:algebraic-gaussian-Hodge} the kernel of $d$ on $C^1(\Lambda_n)$ is $dC^0(\Lambda_n)\oplus\{\text{constants}\}$; as $\operatorname{mean}(d\boldsymbol{\lambda})=0$ for every $\boldsymbol{\lambda}\in C^0(\Lambda_n)$ and the mean of a constant bond function is the constant itself (Section~\ref{sec:tori-cochains-hodge}), it follows that $A^{\downarrow}_0(A^c,0)=L^{-\frac{k}{2}}\cdot c+d\boldsymbol{\lambda}$ with $\boldsymbol{\lambda}\in C^0(\Lambda_n)$. This proves the first assertion.

For the second assertion let $\theta\in\mathbb{R}^3$, with the components $\theta_{\mu}$, $\mu=0,1,2$, and recall the definitions of Sections~\ref{sec:lattice-measures-laws}, \ref{sec:running-quantities-tori} and of Definition~\ref{def:twist-external-connection}: $\eta=L^{-N}$, $n=\frac{R}{\eta}$ and $q=e\cdot\eta^{\frac{1}{2}}$; $a_j=L^{j-N}$, $e_j=e\cdot a_j^{\frac{1}{2}}$ and $l_j=\frac{R}{a_j}$, which agrees with $l_j=\frac{n}{L^j}$ of Section~\ref{sec:coordinates-axial-fields}; the twist $c^{(j)}_{\theta}$ is the constant bond function with the components $\frac{\theta_{\mu}}{e_j\cdot l_j}$, and $H_{\theta}$ is the constant bond function on $\Lambda_n$ with the components $\frac{\theta_{\mu}}{q\cdot n}$. Here $a_j=L^{j-N}$ is the scale of Section~\ref{sec:running-quantities-tori}, and not the number $L^ja$ of Section~\ref{sec:towers-restricted-gauge}. The field $c_{\theta}=c^{(k^{*})}_{\theta}$ is constant, hence closed, and its mean is $c^{(k^{*})}_{\theta}$; so the first assertion applies with $A^c=c_{\theta}$, and it remains to show that $L^{-\frac{k}{2}}\cdot c^{(k)}_{\theta}=H_{\theta}$. One has
\[
e_j\cdot l_j=e\cdot R\cdot a_j^{-\frac{1}{2}},\qquad\text{hence}\qquad\left(c^{(j)}_{\theta}\right)_{\mu}=\frac{\theta_{\mu}\cdot a_j^{\frac{1}{2}}}{e\cdot R},
\]
and $L^{-k}\cdot a_k=a_0=\eta$. Therefore
\[
L^{-\frac{k}{2}}\cdot\left(c^{(k)}_{\theta}\right)_{\mu}=\frac{\theta_{\mu}\cdot\left(L^{-k}\cdot a_k\right)^{\frac{1}{2}}}{e\cdot R}=\left(c^{(0)}_{\theta}\right)_{\mu}=\frac{\theta_{\mu}\cdot\eta^{\frac{1}{2}}}{e\cdot R}=\frac{\theta_{\mu}}{q\cdot n}=\left(H_{\theta}\right)_{\mu}.
\]
\end{proof}

With $k=k^{*}$ under hypothesis~\ref{def:hypothesis-flow}, Proposition~\ref{prop:closed-shifts-terminal} follows.

\subsection{Identities for the fermionic block operators}\label{sec:identities-fermionic-block}

The operators of this section and the identities proved for them occur in \cite[Lemmas~2--4, 18, 19, (23)--(25), (44), (57)--(86), (151), (152), (181), (182)]{DimockMB}, in \cite[(43), (44), (56)--(72), (78)--(116)]{DimockUR} and in \cite[(163), (164), (305)--(319), Section~3.15]{DimockUVS}.

\subsubsection{Spinor spaces and Berezin integration}\label{sec:spinor-spaces-berezin}

\paragraph{Spinors and pairings.}
For $a>0$ and $l\in a\mathbb{N}$ let $\mathfrak{S}(a,l)$ be the complex vector space of the maps $f\colon\mathbb{L}(a)\to\mathbb{C}^4$ with
\[
f(x+le_{\mu})=-f(x)\qquad\left(x\in\mathbb{L}(a),\ \mu=0,1,2\right).
\]
Its elements are the spinors on $\mathbb{T}(a,l)$, and $f_{\boldsymbol{\beta}}$ denotes the components of a spinor $f$. Let $\mathscr{F}$ be a fundamental domain of $\mathbb{L}(a)$ for the translations by $l\mathbb{Z}^3$, that is, a set which contains exactly one point of each coset of $l\mathbb{Z}^3$ in $\mathbb{L}(a)$. A spinor is determined by its values on the $\left(\frac{l}{a}\right)^3$ points of $\mathscr{F}$, and these values may be prescribed arbitrarily; hence $\dim\mathfrak{S}(a,l)=4\left(\frac{l}{a}\right)^3$. For $f,g\in\mathfrak{S}(a,l)$ put
\[
\langle g,f\rangle_a:=\sum_{x\in\mathscr{F}}a^3\sum_{\boldsymbol{\beta}}g_{\boldsymbol{\beta}}(x)\cdot f_{\boldsymbol{\beta}}(x),\qquad(g,f)_a:=\sum_{x\in\mathscr{F}}a^3\sum_{\boldsymbol{\beta}}\overline{g_{\boldsymbol{\beta}}(x)}\cdot f_{\boldsymbol{\beta}}(x).
\]
The first pairing is bilinear. The second is Hermitian; it is the first pairing of the complex conjugate of $g$ with $f$. The summands are $l$-periodic functions of $x$, being products of two antiperiodic factors, so that neither pairing depends on the choice of $\mathscr{F}$.

Let $a,a'>0$ with $l\in a\mathbb{N}\cap a'\mathbb{N}$, and let $O\colon\mathfrak{S}(a,l)\to\mathfrak{S}(a',l)$ be a linear map. Its transpose $O^T$ and its adjoint $O^{*}$ are the linear maps from $\mathfrak{S}(a',l)$ to $\mathfrak{S}(a,l)$ with
\[
\langle g,Of\rangle_{a'}=\langle O^Tg,f\rangle_a,\qquad(g,Of)_{a'}=(O^{*}g,f)_a\qquad\left(f\in\mathfrak{S}(a,l),\ g\in\mathfrak{S}(a',l)\right);
\]
they exist and are unique, since both pairings are non-degenerate. For an endomorphism $O$ of $\mathfrak{S}(a,l)$ put $\operatorname{Re}O:=\frac{1}{2}\left(O+O^{*}\right)$. The operator norm $|O|$ of a linear map $O\colon\mathfrak{S}(a,l)\to\mathfrak{S}(a',l)$ is taken with respect to the norms $f\mapsto(f,f)_a^{\frac{1}{2}}$ and $g\mapsto(g,g)_{a'}^{\frac{1}{2}}$. Thus $\operatorname{Re}O$ and $|O|$ refer to the Hermitian pairings, whereas the trace $\operatorname{Tr}O$ and the determinant $\det O$ of an endomorphism $O$ are those of a linear map of a finite-dimensional vector space. At level $k$ the lattices $\mathbb{L}(a)$ with $a=\eta_k$, $a=1$ and $a=L$ occur, and $\langle\cdot,\cdot\rangle_{\eta_k}$, $\langle\cdot,\cdot\rangle_1$, $\langle\cdot,\cdot\rangle_L$ are the three pairings of Section~\ref{sec:fermion-integrals-blocking}.

\paragraph{Clifford matrices.}
Let $\sigma_1$, $\sigma_2$, $\sigma_3$ be the Pauli matrices. They are Hermitian, $\sigma_1$ and $\sigma_3$ are real and symmetric, $\sigma_2$ is antisymmetric, $\sigma_i^2=1$, $\sigma_i\sigma_j=-\sigma_j\sigma_i$ for $i\ne j$, and $\operatorname{tr}\sigma_1=0$, $\det\sigma_1=-1$. As in Section~\ref{sec:dirac-matrices-lattice} put
\[
\gamma_0:=\sigma_3\otimes\sigma_1,\qquad\gamma_1:=\sigma_3\otimes\sigma_2,\qquad\gamma_2:=\sigma_3\otimes\sigma_3,\qquad\gamma_3:=\sigma_1\otimes1,\qquad\mathsf{C}:=i\gamma_1.
\]
From these properties and the rules $(A\otimes A')(C\otimes C')=AC\otimes A'C'$, $(A\otimes A')^T=A^T\otimes A'^T$, $(A\otimes A')^{*}=A^{*}\otimes A'^{*}$ it follows that the matrices $\gamma_{\mu}$, $0\le\mu\le3$, are Hermitian and satisfy
\[
\gamma_{\mu}\gamma_{\nu}+\gamma_{\nu}\gamma_{\mu}=2\delta_{\mu\nu}\qquad(0\le\mu,\nu\le3),
\]
that $\gamma_0$, $\gamma_2$, $\gamma_3$ are symmetric and anticommute with $\gamma_1$, and that $\gamma_1$ is antisymmetric. Moreover $\operatorname{tr}(A\otimes A')=\operatorname{tr}A\cdot\operatorname{tr}A'$ for square matrices $A$, $A'$, the diagonal entries of $A\otimes A'$ being the products of those of $A$ and of $A'$; and $\det(A\otimes1)=(\det A)^2$ for a $2\times2$ matrix $A$ and the $2\times2$ unit matrix $1$, since one and the same permutation of the rows and of the columns carries $A\otimes1$ into the block-diagonal matrix with the two diagonal blocks $A$. Hence $\operatorname{tr}\gamma_3=\operatorname{tr}\sigma_1\cdot\operatorname{tr}1=0$ and $\det\gamma_3=(\det\sigma_1)^2=1$. Since $\gamma_1^2=1$, one has $\mathsf{C}^{-1}=-i\gamma_1$, and therefore
\begin{equation}\label{eq:algebraic-gaussian-17}
\mathsf{C}^{-1}\cdot\gamma_{\mu}\cdot\mathsf{C}=\gamma_1\gamma_{\mu}\gamma_1=-\gamma_{\mu}^T\qquad(0\le\mu\le3):
\end{equation}
for $\mu\ne1$ one has $\gamma_1\gamma_{\mu}\gamma_1=-\gamma_{\mu}\gamma_1^2=-\gamma_{\mu}=-\gamma_{\mu}^T$, and $\gamma_1\gamma_1\gamma_1=\gamma_1=-\gamma_1^T$. Matrices of size $4\times4$ act on spinors, and on the fields introduced below, pointwise.

\paragraph{Grassmann fields.}
Let $\mathrm{Gr}$ be a Grassmann algebra over $\mathbb{C}$ with finitely many generators, and $\mathrm{Gr}_{\mathrm{odd}}$ its odd part. A field on $\mathbb{T}(a,l)$ is a map $g\colon\mathbb{L}(a)\to\left(\mathrm{Gr}_{\mathrm{odd}}\right)^4$ with $g(x+le_{\mu})=-g(x)$ for all $x\in\mathbb{L}(a)$ and $\mu=0,1,2$. For fields $g$, $f$ on $\mathbb{T}(a,l)$ the pairing $\langle g,f\rangle_a$ is defined by the same formula as for spinors, each product $g_{\boldsymbol{\beta}}(x)\cdot f_{\boldsymbol{\beta}}(x)$ being formed with its factors in this order. Let $(\theta_i)$ be a basis of $\mathrm{Gr}_{\mathrm{odd}}$. Every field $f$ on $\mathbb{T}(a,l)$ has a unique representation $f=\sum_i\theta_if_i$ with spinors $f_i\in\mathfrak{S}(a,l)$, and for a linear map $O\colon\mathfrak{S}(a,l)\to\mathfrak{S}(a',l)$ one puts
\[
Of:=\sum_i\theta_i\cdot Of_i,
\]
which does not depend on the choice of the basis. If $g=\sum_j\theta_jg_j$ is a field on $\mathbb{T}(a',l)$, then, by the definition of $O^T$, both sides of
\[
\langle g,Of\rangle_{a'}=\langle O^Tg,f\rangle_a
\]
are equal to $\sum_{i,j}\theta_j\theta_i\cdot\langle g_j,Of_i\rangle_{a'}$. Thus this rule holds for fields, and no Grassmann factors are exchanged in it.

Now let $\mathscr{F}$ be a fundamental domain of $\mathbb{Z}^3$ for the translations by $l\mathbb{Z}^3$, and let the elements $\Psi_{\boldsymbol{\beta}}(y)$, $\bar{\Psi}_{\boldsymbol{\beta}}(y)$, where $y\in\mathscr{F}$ and $\boldsymbol{\beta}$ ranges over the four components, be among the generators of $\mathrm{Gr}$. Extended to all $y\in\mathbb{Z}^3$ by $\Psi(y+le_{\mu}):=-\Psi(y)$ and $\bar{\Psi}(y+le_{\mu}):=-\bar{\Psi}(y)$, they form fields $\Psi$ and $\bar{\Psi}$ on $T^0_k=\mathbb{T}(1,l)$.

\paragraph{Integration over a family of generators.}
Let $\mathcal{M}_{\Psi}$ be the product, in a fixed order, of the $8n_k$ generators $\Psi_{\boldsymbol{\beta}}(y)$, $\bar{\Psi}_{\boldsymbol{\beta}}(y)$, $y\in\mathscr{F}$, where $n_k=l^3$. Every $F\in\mathrm{Gr}$ can be written in a unique way as $F=c\cdot\mathcal{M}_{\Psi}+R$, where $c$ lies in the subalgebra generated by the remaining generators of $\mathrm{Gr}$ and every term of $R$ has degree less than $8n_k$ in the generators $\Psi_{\boldsymbol{\beta}}(y)$, $\bar{\Psi}_{\boldsymbol{\beta}}(y)$. The Berezin integral over these generators is the linear map $F\mapsto\int\mathrm{D}\Psi\,F$ of $\mathrm{Gr}$ into itself which sends $F$ to the product of $c$ with a complex number that is the same for all $F$; this number is fixed by the normalisation
\[
\int\mathrm{D}\Psi\,\exp\left(-\langle\bar{\Psi},O\Psi\rangle_1\right)=\det O\qquad\left(O\text{ an endomorphism of }\mathfrak{S}(1,l)\right)
\]
of Section~\ref{sec:fermion-integrals-blocking}. Let $W_{\boldsymbol{\beta}}(y)$, $\bar{W}_{\boldsymbol{\beta}}(y)$ and $W'_{\boldsymbol{\beta}}(y)$, $\bar{W}'_{\boldsymbol{\beta}}(y)$, $y\in\mathscr{F}$, be two further families of generators of $\mathrm{Gr}$, extended antiperiodically to fields $W$, $\bar{W}$, $W'$, $\bar{W}'$ on $T^0_k$, let $\int\mathrm{D}W$ and $\int\mathrm{D}W'$ be the Berezin integrals over these families, defined and normalised in the same way, and put
\[
\int\mathrm{d}\mu_I(W)\,F:=\int\mathrm{D}W\,\exp\left(-\langle\bar{W},W\rangle_1\right)F\qquad(F\in\mathrm{Gr}),
\]
that is, $\mathrm{d}\mu_I(W):=\exp\left(-\langle\bar{W},W\rangle_1\right)\mathrm{D}W$.

For fields $h$, $h'$ on $T^0_k$ and $G\in\mathrm{Gr}$ let $G(h',h)$ be the image of $G$ under the homomorphism of algebras $\mathrm{Gr}\to\mathrm{Gr}$ which maps $\bar{\Psi}_{\boldsymbol{\beta}}(y)$ to $h'_{\boldsymbol{\beta}}(y)$ and $\Psi_{\boldsymbol{\beta}}(y)$ to $h_{\boldsymbol{\beta}}(y)$ for $y\in\mathscr{F}$ and fixes the remaining generators; this homomorphism exists and is unique, because the values of $h$ and $h'$ are odd. In particular $G(\bar{\Psi},\Psi)=G$. The same notation is used for the families $W$, $\bar{W}$ and $W'$, $\bar{W}'$; thus $G(\bar{W}',W')=G$.

\paragraph{Rules of Berezin integration.}
The following rules are classical (Section~\ref{sec:results-literature}; F.~A.~Berezin, \emph{The Method of Second Quantization}, Academic Press, 1966). In \eqref{eq:algebraic-gaussian-Ber1}--\eqref{eq:algebraic-gaussian-Ber3}, $g$, $\bar{g}$, $J$, $\bar{J}$ are fields on $T^0_k$ with values in the subalgebra of $\mathrm{Gr}$ generated by the generators other than the $\Psi_{\boldsymbol{\beta}}(y)$, $\bar{\Psi}_{\boldsymbol{\beta}}(y)$, and $G\in\mathrm{Gr}$. The rules and their proofs remain valid for every family of generators indexed in the same way by the sites of a torus $\mathbb{T}(1,l')$, in particular for generators $\Psi'$, $\bar{\Psi}'$ indexed by the sites of $T^0_{k+1}$, and for the families $W$, $\bar{W}$ and $W'$, $\bar{W}'$.

First,
\begin{equation}\tag{Ber1}\label{eq:algebraic-gaussian-Ber1}
\int\mathrm{D}\Psi\,G\left(\bar{\Psi}+\bar{g},\Psi+g\right)=\int\mathrm{D}\Psi\,G\left(\bar{\Psi},\Psi\right).
\end{equation}
Indeed, write $G=c\cdot\mathcal{M}_{\Psi}+R$ as above. The substitution fixes $c$, it maps $\mathcal{M}_{\Psi}$ to the sum of $\mathcal{M}_{\Psi}$ and of terms that contain a component of $g$ or of $\bar{g}$, and it maps every term of $R$ to terms of degree less than $8n_k$ in $\Psi$, $\bar{\Psi}$. The components of $g$ and $\bar{g}$ contain none of the generators $\Psi_{\boldsymbol{\beta}}(y)$, $\bar{\Psi}_{\boldsymbol{\beta}}(y)$; hence every term that contains such a component has degree less than $8n_k$ in $\Psi$, $\bar{\Psi}$, and the coefficient of the monomial $\mathcal{M}_{\Psi}$ of top degree in $G\left(\bar{\Psi}+\bar{g},\Psi+g\right)$ is again $c$.

Secondly, for every invertible endomorphism $O$ of $\mathfrak{S}(1,l)$,
\begin{equation}\tag{Ber2}\label{eq:algebraic-gaussian-Ber2}
\int\mathrm{D}\Psi\,\exp\left(-\langle\bar{\Psi},O\Psi\rangle_1+\langle\bar{\Psi},J\rangle_1+\langle\bar{J},\Psi\rangle_1\right)=\det O\cdot\exp\left\langle\bar{J},O^{-1}J\right\rangle_1.
\end{equation}
Indeed, apply \eqref{eq:algebraic-gaussian-Ber1} with $g:=O^{-1}J$ and $\bar{g}:=\left(O^{-1}\right)^T\bar{J}$. Then $Og=J$ and $O^T\bar{g}=\bar{J}$, and the transposition rule for fields gives $\langle\bar{g},O\Psi\rangle_1=\langle\bar{J},\Psi\rangle_1$ and $\langle\bar{g},Og\rangle_1=\langle\bar{g},J\rangle_1$. Of the eight terms of
\[
-\left\langle\bar{\Psi}+\bar{g},O(\Psi+g)\right\rangle_1+\left\langle\bar{\Psi}+\bar{g},J\right\rangle_1+\left\langle\bar{J},\Psi+g\right\rangle_1,
\]
the term $-\langle\bar{\Psi},Og\rangle_1$ cancels $\langle\bar{\Psi},J\rangle_1$, the term $-\langle\bar{g},O\Psi\rangle_1$ cancels $\langle\bar{J},\Psi\rangle_1$, and $-\langle\bar{g},Og\rangle_1$ cancels $\langle\bar{g},J\rangle_1$, while $-\langle\bar{\Psi},O\Psi\rangle_1$ and $\langle\bar{J},g\rangle_1=\left\langle\bar{J},O^{-1}J\right\rangle_1$ remain. Since the substitution is a homomorphism, it maps the exponential of the exponent to the exponential of its image, and the image of the exponent is
\[
-\langle\bar{\Psi},O\Psi\rangle_1+\left\langle\bar{J},O^{-1}J\right\rangle_1.
\]
These two terms are even and commute, so that the exponential of their sum is the product of their exponentials. The factor $\exp\left\langle\bar{J},O^{-1}J\right\rangle_1$ is an even element of the subalgebra generated by the generators other than $\Psi$, $\bar{\Psi}$; it can therefore be taken out of the integral, and the normalisation of $\int\mathrm{D}\Psi$ gives \eqref{eq:algebraic-gaussian-Ber2}.

Thirdly, let $u\colon\mathbb{Z}^3\to\mathbb{C}$ be an $l$-periodic function without zeros, and let $(u\Psi)(y):=u(y)\Psi(y)$ and $\left(u^{-1}\bar{\Psi}\right)(y):=u(y)^{-1}\bar{\Psi}(y)$. Then
\begin{equation}\tag{Ber3}\label{eq:algebraic-gaussian-Ber3}
\int\mathrm{D}\Psi\,G\left(\bar{\Psi},\Psi\right)=\int\mathrm{D}\Psi\,G\left(u^{-1}\bar{\Psi},u\Psi\right).
\end{equation}
Indeed, $u\Psi$ and $u^{-1}\bar{\Psi}$ are fields on $T^0_k$, because $u$ is periodic. The substitution fixes the coefficients from the remaining generators, multiplies every monomial in $\Psi$, $\bar{\Psi}$ by a number without changing its degree, and multiplies the monomial $\mathcal{M}_{\Psi}$ of top degree by
\[
\prod_{y\in\mathscr{F}}\left(u(y)\cdot u(y)^{-1}\right)^4=1;
\]
hence the coefficient of $\mathcal{M}_{\Psi}$ is the same on both sides.

Fourthly, for every invertible endomorphism $O$ of $\mathfrak{S}(1,l)$ and every $G$ in the subalgebra of $\mathrm{Gr}$ generated by the generators other than the $W_{\boldsymbol{\beta}}(y)$, $\bar{W}_{\boldsymbol{\beta}}(y)$,
\begin{equation}\tag{Ber4}\label{eq:algebraic-gaussian-Ber4}
\int\mathrm{D}W'\,\exp\left(-\langle\bar{W}',OW'\rangle_1\right)G\left(\bar{W}',W'\right)=\det O\cdot\int\mathrm{d}\mu_I(W)\,G\left(\bar{W},O^{-1}W\right).
\end{equation}
Indeed, let $c$ be an element of the subalgebra of $\mathrm{Gr}$ generated by the generators that belong to neither of the two families $W$, $\bar{W}$ and $W'$, $\bar{W}'$, that is, by the generators other than the components of $W$, $\bar{W}$, $W'$ and $\bar{W}'$. The Berezin integral over either family commutes with the multiplication by $c$ from the left: if $F=c'\cdot\mathcal{M}+R$ is the decomposition of $F\in\mathrm{Gr}$ by which that integral is defined, $\mathcal{M}$ being the product of the $8n_k$ generators of the family, then $c\cdot F=(c\cdot c')\cdot\mathcal{M}+c\cdot R$ is the decomposition of $c\cdot F$, since $c$ contains no generator of the family. Moreover, the substitution on the right of \eqref{eq:algebraic-gaussian-Ber4} fixes $c$, and the exponentials $\exp\left(-\langle\bar{W}',OW'\rangle_1\right)$ and $\exp\left(-\langle\bar{W},W\rangle_1\right)$ are even, hence commute with $c$. Consequently each side of \eqref{eq:algebraic-gaussian-Ber4}, formed with $c\cdot G$ in place of $G$, is $c$ times the same side formed with $G$. Each side is also additive in $G$, and every $G$ admitted in \eqref{eq:algebraic-gaussian-Ber4} is a finite sum of products $c\cdot m$ of elements $c$ of this subalgebra and monomials $m$ in the components of $W'$, $\bar{W}'$. It therefore suffices to prove \eqref{eq:algebraic-gaussian-Ber4} for these monomials. For $G=m$ the two sides are complex numbers, since the integrands lie in the subalgebras generated by the components of $W'$, $\bar{W}'$ and of $W$, $\bar{W}$, respectively; these numbers do not change when $\mathrm{Gr}$ is enlarged by further generators, because for an element $F$ of the subalgebra generated by one family the decomposition $F=c'\cdot\mathcal{M}+R$, and with it the number fixed by the normalisation, are the same in the enlarged algebra.

Enlarging $\mathrm{Gr}$ in this way, let $J$, $\bar{J}$ be fields on $T^0_k$ whose components $J_{\boldsymbol{\beta}}(y)$, $\bar{J}_{\boldsymbol{\beta}}(y)$, $y\in\mathscr{F}$, are distinct new generators. For
\[
G=\exp\left(\langle\bar{W}',J\rangle_1+\langle\bar{J},W'\rangle_1\right)
\]
the left side equals $\det O\cdot\exp\left\langle\bar{J},O^{-1}J\right\rangle_1$, by \eqref{eq:algebraic-gaussian-Ber2} for the family $W'$, $\bar{W}'$. On the right side
\[
G\left(\bar{W},O^{-1}W\right)=\exp\left(\langle\bar{W},J\rangle_1+\left\langle\left(O^{-1}\right)^T\bar{J},W\right\rangle_1\right),
\]
and \eqref{eq:algebraic-gaussian-Ber2} for the family $W$, $\bar{W}$, with the identity in place of $O$ and with $\left(O^{-1}\right)^T\bar{J}$ in place of $\bar{J}$, gives
\[
\int\mathrm{d}\mu_I(W)\,G\left(\bar{W},O^{-1}W\right)=\exp\left\langle\left(O^{-1}\right)^T\bar{J},J\right\rangle_1=\exp\left\langle\bar{J},O^{-1}J\right\rangle_1.
\]
So the two sides are equal for this $G$. On the other hand, the exponent of this $G$ is the sum of the even elements $\bar{W}'_{\boldsymbol{\beta}}(y)\cdot J_{\boldsymbol{\beta}}(y)$ and $\bar{J}_{\boldsymbol{\beta}}(y)\cdot W'_{\boldsymbol{\beta}}(y)$, $y\in\mathscr{F}$, whose squares vanish; hence $G$ is the product of the factors $1+\bar{W}'_{\boldsymbol{\beta}}(y)\cdot J_{\boldsymbol{\beta}}(y)$ and $1+\bar{J}_{\boldsymbol{\beta}}(y)\cdot W'_{\boldsymbol{\beta}}(y)$. Expanding this product and moving, in each term, the components of $J$, $\bar{J}$ to the left, one finds
\[
G=\sum_m\pm\,\tilde{m}\cdot m,
\]
where the sum extends over the monomials $m$ in the components of $W'$, $\bar{W}'$, each written with its factors in a fixed order, $\tilde{m}$ is the monomial in the components of $J$, $\bar{J}$ that contains $J_{\boldsymbol{\beta}}(y)$ if and only if $m$ contains $\bar{W}'_{\boldsymbol{\beta}}(y)$, and $\bar{J}_{\boldsymbol{\beta}}(y)$ if and only if $m$ contains $W'_{\boldsymbol{\beta}}(y)$, and the sign depends on $m$ only. The components of $J$, $\bar{J}$ belong to neither of the two families; so the first part of the proof, applied in the enlarged algebra with $\pm\,\tilde{m}$ in place of $c$, shows that each side of \eqref{eq:algebraic-gaussian-Ber4} at this $G$ is the sum over $m$ of $\pm\,\tilde{m}$ times the same side at $G=m$, which is a complex number. Distinct monomials in the components of $J$, $\bar{J}$ are linearly independent over $\mathbb{C}$. Hence the two sides of \eqref{eq:algebraic-gaussian-Ber4} agree at every monomial $m$, and \eqref{eq:algebraic-gaussian-Ber4} follows.

Fifthly, let $\int\mathcal{D}_1$ and $\int\mathcal{D}_2$ be the Berezin integrals over two disjoint families of generators of the kind considered here, let $F\in\mathrm{Gr}$ be even, let $\nu$ be a finite measure on a measurable space $\Omega$, and let $\Phi\colon\Omega\to\mathrm{Gr}$ be a bounded map, measurable with respect to the Borel sets of $\mathrm{Gr}$, its integral being taken coordinatewise with respect to a basis of $\mathrm{Gr}$. Then
\begin{equation}\tag{Ber5}\label{eq:algebraic-gaussian-Ber5}
\int\mathcal{D}_1\int\mathcal{D}_2\,F=\int\mathcal{D}_2\int\mathcal{D}_1\,F,\qquad\int\mathcal{D}_1\int\mathrm{d}\nu(t)\,\Phi(t)=\int\mathrm{d}\nu(t)\int\mathcal{D}_1\,\Phi(t).
\end{equation}
Indeed, by its definition a Berezin integral over a family of generators multiplies by a fixed constant the coefficient with which the product of all the generators of the family occurs in its argument; in particular it is a linear map of the finite-dimensional space $\mathrm{Gr}$ into itself. Both sides of the first identity are the same multiple of the coefficient of $F$ at the product of the two monomials of top degree of the two families, which have even degree; and the coordinatewise integral commutes with every linear map of $\mathrm{Gr}$ into itself.

\subsubsection{The block operators}\label{sec:block-operators}

\paragraph{Parameters.}
Let $k\ge0$ be an integer and $l\in L\mathbb{N}$. In every statement that involves the operators $\delta_{\alpha,B}$ or the set $\mathrm{Pz}_k$ of Definition~\ref{def:pairs-sites-holder}, $l$ is in addition required to satisfy the inequality $2M<l$ on which that definition rests, $M$ being the integer of Section~\ref{sec:lattice-geometry}. Under this requirement each pair of $\mathrm{Pz}_k$ has a shortest lift that is unique up to a common translation of its two points by a vector of $l\mathbb{Z}^3$. Division by $L$ turns the inequality into $\frac{2M}{L}<\frac{l}{L}$. The torus $T_{k+1}$ has the side $\frac{l}{L}$, and for $(x,x')\in\mathrm{Pz}_k$ the sites $\frac{x}{L}$, $\frac{x'}{L}$ of $T_{k+1}$ are at distance at most $\frac{M}{L}$ from each other, which by this inequality is less than half the side; hence the pair $\left(\frac{x}{L},\frac{x'}{L}\right)$, too, has a shortest lift that is unique up to a common translation by a vector of $\frac{l}{L}\mathbb{Z}^3$. Let $\mathfrak{m}\in\mathbb{C}$ (the mass), $e_k>0$ and $b>0$, and put $\tilde{b}:=bL^{-1}$ and
\[
b_1:=b,\qquad b_{j+1}:=\frac{b\cdot b_j}{b_j+\tilde{b}}\qquad(j\ge1).
\]
These are the numbers $b_j=\frac{b\left(1-L^{-1}\right)}{1-L^{-j}}$ of Section~\ref{sec:running-quantities-tori}: the formula holds for $j=1$, and if it holds for $j$, then $b_j+\tilde{b}=\frac{b\left(1-L^{-j-1}\right)}{1-L^{-j}}$, whence $b_{j+1}=\frac{b\left(1-L^{-1}\right)}{1-L^{-j-1}}$. Let $B$ be an $l$-periodic function from the set of the oriented bonds of $\mathbb{L}(\eta_k)$ to $\mathbb{C}$ that changes its sign when the orientation of a bond is reversed, and put $\eta:=\eta_k$.

\paragraph{Phase factors.}
For a chain $\boldsymbol{\Gamma}$ of bonds of $\mathbb{L}(\eta)$, called fine bonds, put
\[
u_B(\boldsymbol{\Gamma}):=\exp\left\{i\cdot e_k\cdot\eta\cdot B(\boldsymbol{\Gamma})\right\};
\]
for a fine bond $(x,x')$, $u_B(x,x')$ denotes the value at the chain formed by this bond. Let $y\in\mathbb{Z}^3$ and $x\in B_k(y)$. For $0\le j\le k$ denote by $y_j$ the point of $\mathbb{L}(L^j\eta)$ whose block $B_j(y_j)$ contains $x$, the blocks being those of $\mathbb{L}(\eta)$ (Section~\ref{sec:blocks-levels-rectilinear}); thus $y_0=x$, $y_k=y$, and $y_j\in B(y_{j+1})$ for $j<k$ by \eqref{eq:algebraic-gaussian-2}. Put
\begin{align*}
(\tau B)(y',y'')&:=\frac{1}{3!}\sum_{\pi}B\left(\boldsymbol{\Gamma}^{\pi}(y',y'')\right),\\
(\tau_kB)(y,x)&:=\sum_{j=0}^{k-1}(\tau B)(y_{j+1},y_j),\\
\boldsymbol{\phi}^B_k(y,x)&:=\exp\left\{i\cdot e_k\cdot\eta\cdot(\tau_kB)(y,x)\right\},
\end{align*}
where, for $j\ge0$, $y'\in\mathbb{L}(L^{j+1}\eta)$ and $y''\in B(y')\subset\mathbb{L}(L^j\eta)$, the sum extends over the six permutations $\pi$ of the directions, and the rectilinear path $\boldsymbol{\Gamma}^{\pi}(y',y'')$ of $\mathbb{L}(L^j\eta)$ (Section~\ref{sec:blocks-levels-rectilinear}) is regarded as the chain of the fine bonds that it covers. These phases occur in \cite[(58)]{DimockUR} and \cite[(173)]{DimockUVS}. For $z\in L\mathbb{Z}^3$ and $y\in B(z)$, the block of $z$ in $\mathbb{Z}^3$, put
\[
\boldsymbol{\phi}^B(z,y):=\exp\left\{i\cdot e_k\cdot\eta\cdot(\tau B)(z,y)\right\}
\]
as in Section~\ref{sec:fermion-integrals-blocking}, and for $x\in B_{k+1}(z)$ put
\[
(\tau_{k+1}B)(z,x):=(\tau B)(z,[x])+(\tau_kB)([x],x);
\]
here $[x]\in B(z)$ by \eqref{eq:algebraic-gaussian-2}.

\paragraph{The operators of level $k$.}
Put $P^{\mp}_{\mu}:=\frac{1}{2}\left(1\mp\gamma_{\mu}\right)$ for $\mu=0,1,2$. The following linear maps are defined as in Sections~\ref{sec:exact-fermion-operators} and~\ref{sec:fermion-integrals-blocking}. For $f\in\mathfrak{S}(\eta,l)$ and $x\in\mathbb{L}(\eta)$ let
\begin{align*}
(D_Bf)(x)&:=-\eta^{-1}\sum_{\mu=0}^2\left\{u_B(x,x+\eta e_{\mu})\cdot P^{-}_{\mu}f(x+\eta e_{\mu})\right.\\
&\qquad\qquad\qquad\left.+u_B(x,x-\eta e_{\mu})\cdot P^{+}_{\mu}f(x-\eta e_{\mu})-f(x)\right\},\\
K(B)&:=D_B+\mathfrak{m}.
\end{align*}
For $f\in\mathfrak{S}(\eta,l)$, $v\in\mathfrak{S}(1,l)$, $w\in\mathfrak{S}(L,l)$, $y\in\mathbb{Z}^3$, $x\in\mathbb{L}(\eta)$ and $z\in L\mathbb{Z}^3$ let
\begin{align*}
(Q_k(B)f)(y)&:=\eta^3\sum_{x\in B_k(y)}\boldsymbol{\phi}^B_k(y,x)\cdot f(x),&(Q_k^T(B)v)(x)&:=\boldsymbol{\phi}^B_k([x],x)\cdot v([x]),\\
(Q(B)v)(z)&:=L^{-3}\sum_{y\in B(z)}\boldsymbol{\phi}^B(z,y)\cdot v(y),&(Q^T(B)w)(y)&:=\boldsymbol{\phi}^B(z_y,y)\cdot w(z_y),
\end{align*}
where $z_y\in L\mathbb{Z}^3$ is the point with $y\in B(z_y)$, and let
\[
P_k(B):=Q_k^T(-B)\cdot Q_k(B),\qquad P(B):=Q^T(-B)\cdot Q(B).
\]
For $k\ge1$ let
\begin{align*}
X_k(B)&:=K(B)+b_k\cdot P_k(B),&S_k(B)&:=X_k(B)^{-1},\\
H_k(B)&:=b_k\cdot S_k(B)\cdot Q_k^T(-B),&\mathfrak{H}_k(B)&:=b_k\cdot Q_k(B)\cdot S_k(B),\\
D_k(B)&:=b_k-b_k^2\cdot Q_k(B)\cdot S_k(B)\cdot Q_k^T(-B),&M_k(B)&:=b_k-D_k(B),\\
\dot{M}_k(B)&:=\mathfrak{H}_k(B)\cdot H_k(B),&T_k(B)&:=b_k^{-1}\cdot Q_k(B)\cdot X_k(B),
\end{align*}
and for $k\ge0$ let
\begin{align*}
\Gamma_k(B)&:=\left(D_k(B)+\tilde{b}\cdot P(B)\right)^{-1},&\delta Z_k(B)&:=\det\left(D_k(B)+\tilde{b}\cdot P(B)\right),\\
\bar{H}_k(B)&:=\mathsf{C}^{-1}\cdot H_k(-B)\cdot\mathsf{C},&\bar{T}_k(B)&:=\mathsf{C}^{-1}\cdot T_k(-B)\cdot\mathsf{C},
\end{align*}
where $H_0$, $T_0$ and $D_0$ are defined in the paragraph on level zero below. Here $D_B$, $K(B)$, $P_k(B)$, $X_k(B)$ and $S_k(B)$ are endomorphisms of $\mathfrak{S}(\eta,l)$; $Q_k(B)$, $\mathfrak{H}_k(B)$, $T_k(B)$ and $\bar{T}_k(B)$ map $\mathfrak{S}(\eta,l)$ to $\mathfrak{S}(1,l)$; $Q_k^T(B)$, $H_k(B)$ and $\bar{H}_k(B)$ map $\mathfrak{S}(1,l)$ to $\mathfrak{S}(\eta,l)$; $Q(B)$ maps $\mathfrak{S}(1,l)$ to $\mathfrak{S}(L,l)$, and $Q^T(B)$ maps $\mathfrak{S}(L,l)$ to $\mathfrak{S}(1,l)$; $P(B)$, $D_k(B)$, $M_k(B)$, $\dot{M}_k(B)$ and $\Gamma_k(B)$ are endomorphisms of $\mathfrak{S}(1,l)$, a number such as $b_k$ standing for the corresponding multiple of the identity; $\delta Z_k(B)$ is a complex number. The operators formed with $S_k(B)$ are defined when $X_k(B)$ is invertible, and $\Gamma_k(B)$ is defined when $D_k(B)+\tilde{b}\cdot P(B)$ is invertible.

The right-hand sides of these formulas are antiperiodic, because $B$ is $l$-periodic and $B_k(y+le_{\mu})=B_k(y)+le_{\mu}$, $B(z+le_{\mu})=B(z)+le_{\mu}$, so that the phases do not change when all their arguments are translated by $le_{\mu}$. Hence all these maps preserve antiperiodicity, and they are well defined between the spaces just indicated. With respect to the bilinear pairings of Section~\ref{sec:spinor-spaces-berezin}, $Q_k^T(B)$ is the transpose of $Q_k(B)$, and $Q^T(B)$ is that of $Q(B)$. Indeed, let $\mathscr{F}$ be a fundamental domain of $\mathbb{Z}^3$ for $l\mathbb{Z}^3$, and $\mathscr{F}'$ one of $L\mathbb{Z}^3$ for $l\mathbb{Z}^3$; then the union of the blocks $B_k(y)$, $y\in\mathscr{F}$, is a fundamental domain of $\mathbb{L}(\eta)$, and the union of the blocks $B(z)$, $z\in\mathscr{F}'$, is one of $\mathbb{Z}^3$. Since $[x]=y$ for $x\in B_k(y)$, and $z_y=z$ for $y\in B(z)$, one obtains for $f\in\mathfrak{S}(\eta,l)$, $v\in\mathfrak{S}(1,l)$ and $w\in\mathfrak{S}(L,l)$
\begin{align*}
\langle v,Q_k(B)f\rangle_1&=\sum_{y\in\mathscr{F}}\ \sum_{x\in B_k(y)}\eta^3\cdot\boldsymbol{\phi}^B_k(y,x)\sum_{\boldsymbol{\beta}}v_{\boldsymbol{\beta}}(y)\cdot f_{\boldsymbol{\beta}}(x)=\langle Q_k^T(B)v,f\rangle_{\eta},\\
\langle w,Q(B)v\rangle_L&=\sum_{z\in\mathscr{F}'}\ \sum_{y\in B(z)}\boldsymbol{\phi}^B(z,y)\sum_{\boldsymbol{\beta}}w_{\boldsymbol{\beta}}(z)\cdot v_{\boldsymbol{\beta}}(y)=\langle Q^T(B)w,v\rangle_1;
\end{align*}
in the second line the factor $L^3$ of the pairing on $L\mathbb{Z}^3$ and the factor $L^{-3}$ of $Q(B)$ cancel.

\paragraph{Level zero.}
At level $k=0$ one has $\eta=1$, $B_0(y)=\{y\}$, and $\tau_0B=0$, the sum that defines it being empty; hence $Q_0(B)=1$. One puts
\[
H_0(B):=1,\qquad T_0(B):=1,\qquad D_0(B):=K(B),\qquad\dot{M}_0:=1,
\]
so that $\bar{H}_0(B)=\bar{T}_0(B)=1$ and $\Gamma_0(B)=\left(K(B)+\tilde{b}\cdot P(B)\right)^{-1}$ (Section~\ref{sec:exact-fermion-operators}). The number $b_0$ and the operators $S_0$ and $M_0$ are not defined.

\paragraph{Level $k+1$.}
At level $k+1$ the data are $\left(k+1,\frac{l}{L},L\cdot\mathfrak{m},e_{k+1},b\right)$, where $e_{k+1}:=L^{\frac{1}{2}}e_k$. The operators of level $k+1$ are formed from these data. They are distinguished from those of level $k$ by the index $k+1$ in place of $k$, where such an index is written, and by their argument $B'$, an $\frac{l}{L}$-periodic function from the set of the oriented bonds of $\mathbb{L}(\eta_{k+1})$ to $\mathbb{C}$ that changes its sign when the orientation of a bond is reversed.

\paragraph{Scaling by $L$.}
Let $a>0$ and $l\in La\mathbb{N}$. For $f\in\mathfrak{S}\left(a,\frac{l}{L}\right)$ put
\[
f_L(x):=L^{-1}f\left(\frac{x}{L}\right)\qquad\left(x\in\mathbb{L}(La)\right).
\]
Since $x\mapsto\frac{x}{L}$ maps $\mathbb{L}(La)$ bijectively onto $\mathbb{L}(a)$, and $x+le_{\mu}$ to $\frac{x}{L}+\frac{l}{L}e_{\mu}$, the map $f\mapsto f_L$ is a bijection of $\mathfrak{S}\left(a,\frac{l}{L}\right)$ onto $\mathfrak{S}(La,l)$. The same formula defines $f_L$ for fields $f$ on $\mathbb{T}\left(a,\frac{l}{L}\right)$. Let $\alpha$ be the H\"older order of stage~\ref{it:order-choice-2}. Let $\chi$ be a function, with values in a complex vector space, whose domain is a set of ordered pairs of points of $\mathbb{L}(a)$ which contains, together with each of its pairs $(x,x')$, every pair $(x+v,x'+v)$ with $v\in\frac{l}{L}\mathbb{Z}^3$; pairs of sites of $\mathbb{T}\left(a,\frac{l}{L}\right)$ are thus represented, in the manner of Section~\ref{sec:grassmann-algebras-norms}, by pairs of points of $\mathbb{L}(a)$ that project onto them. As in Section~\ref{sec:local-parts-adjustment}, $\chi_L$ denotes the function defined at those ordered pairs $(x,x')$ of points of $\mathbb{L}(La)$ for which $\left(\frac{x}{L},\frac{x'}{L}\right)$ lies in the domain of $\chi$, given by
\[
\chi_L(x,x'):=L^{-1-\alpha}\chi\left(\frac{x}{L},\frac{x'}{L}\right).
\]
For $v\in l\mathbb{Z}^3$ one has $\frac{x+v}{L}=\frac{x}{L}+\frac{v}{L}$ with $\frac{v}{L}\in\frac{l}{L}\mathbb{Z}^3$; hence the domain of $\chi_L$ contains, together with each of its pairs $(x,x')$, every pair $(x+v,x'+v)$ with $v\in l\mathbb{Z}^3$.

\begin{lemma}[Identities of the block operators]\label{lem:identities-block-operators}
Let the data and the operators be those of Section~\ref{sec:block-operators}, and assume that every inverse which is written below exists. The assertions in which $b_k$, $X_k$, $S_k$, $M_k$ or $\mathfrak{H}_k$ occur concern $k \ge 1$ only. For $k = 0$ the operators $T_0$, $\bar{T}_0$, $H_0$, $\bar{H}_0$ and $\dot{M}_0$ are the identity by their definitions at level~$0$ in Section~\ref{sec:block-operators}, and the assertions about them follow from these definitions; the assertions about $D_0 = K$ and $\Gamma_0$ are consequences of those about $D_B$ and $P(B)$.
\begin{enumerate}
\item[(a)] One has
\[
Q_k(B) \cdot Q_k^T(-B) = 1, \qquad Q(B) \cdot Q^T(-B) = 1.
\]
The operators $P_k(B)$ and $P(B)$ are idempotent, they commute with every complex $4 \times 4$ matrix acting on the spinor index, and
\[
P(B) \cdot Q^T(-B) = Q^T(-B), \qquad Q(B) \cdot P(B) = Q(B).
\]
For every complex $4 \times 4$ matrix $\mathrm{Mx}$, regarded as the operator of pointwise multiplication on $\mathfrak{S}(1,l)$,
\begin{align*}
\operatorname{Tr}\left[\mathrm{Mx} \cdot P(B)\right] &= n_{k+1} \cdot \operatorname{tr}\mathrm{Mx},\\
\operatorname{Tr}\left[\mathrm{Mx} \cdot \left(1 - P(B)\right)\right] &= \left(n_k - n_{k+1}\right) \cdot \operatorname{tr}\mathrm{Mx},
\end{align*}
where $n_k = l^3$ and $n_{k+1} = (l/L)^3$. Moreover $T_k(B) \cdot H_k(B) = 1$.
\item[(b)] Let $Q_{k+1}(B) := Q(B) \cdot Q_k(B)$, a linear map of $\mathfrak{S}(\eta,l)$ into $\mathfrak{S}(L,l)$. For $f \in \mathfrak{S}(\eta,l)$ and $z \in L\mathbb{Z}^3$,
\begin{align*}
&\left(Q_{k+1}(B) f\right)(z)\\
&\qquad = L^{-3} \cdot \eta^3 \cdot \sum_{x \in B_{k+1}(z)} \exp\left\{ i \cdot e_k \cdot \eta \cdot (\tau_{k+1} B)(z,x) \right\} \cdot f(x).
\end{align*}
With
\begin{align*}
Q_{k+1}^T(-B) &:= Q_k^T(-B) \cdot Q^T(-B),\\
P_{k+1}(B) &:= Q_{k+1}^T(-B) \cdot Q_{k+1}(B)
\end{align*}
one has $P_{k+1}(B) = Q_k^T(-B) \cdot P(B) \cdot Q_k(B)$.
\item[(c)] \emph{Transposes.} With respect to the bilinear pairings of Section~\ref{sec:spinor-spaces-berezin},
\[
D_B^T = \mathsf{C}^{-1} \cdot D_{-B} \cdot \mathsf{C}, \qquad P_k(B)^T = P_k(-B), \qquad P(B)^T = P(-B).
\]
Consequently $O(B)^T = \mathsf{C}^{-1} \cdot O(-B) \cdot \mathsf{C}$ for $O = X_k, S_k, D_k, M_k, \Gamma_k$, and
\begin{align*}
\bar{H}_k(B) &= \mathfrak{H}_k(B)^T,\\
\bar{T}_k(B) &= b_k^{-1} \cdot Q_k(-B) \cdot X_k(B)^T,\\
\bar{T}_k(B) \cdot \bar{H}_k(B) &= 1,\\
\left\langle \bar{H}_k(B) \bar{\Psi}, H_k(B) \Psi \right\rangle_{\eta} &= \left\langle \bar{\Psi}, \dot{M}_k(B) \Psi \right\rangle_1,
\end{align*}
where $\Psi$ and $\bar{\Psi}$ are the fields on $T^0_k$ formed from the generators of Section~\ref{sec:spinor-spaces-berezin}.
\item[(d)] \emph{Gauge transformations.} Let $\boldsymbol{\lambda} \colon \mathbb{L}(\eta) \to \mathbb{C}$, and let $d\boldsymbol{\lambda}$ be the function on the oriented bonds of $\mathbb{L}(\eta)$ given by
\[
(d\boldsymbol{\lambda})(x,x') := \eta^{-1} \cdot \left(\boldsymbol{\lambda}(x') - \boldsymbol{\lambda}(x)\right);
\]
it is odd under reversal, and its components $(d\boldsymbol{\lambda})_{\mu}(x) = (d\boldsymbol{\lambda})(x, x + \eta e_{\mu})$ are given by the formula of \eqref{eq:algebraic-gaussian-1} with $a = \eta$. Suppose that $d\boldsymbol{\lambda}$ and $\exp\left(i \cdot e_k \cdot \boldsymbol{\lambda}\right)$ are both $l$-periodic. These conditions are met if $\boldsymbol{\lambda}$ is $l$-periodic; they are also met if $\boldsymbol{\lambda}(x) = c \cdot x$ with a vector $c$ such that $e_k \cdot l \cdot c_{\mu} \in 2\pi\mathbb{Z}$ for $\mu = 0, 1, 2$, and then $d\boldsymbol{\lambda}$ is the constant bond function with the components $c_{\mu}$. For a function $f$ on one of the lattices $\mathbb{L}(\eta)$, $\mathbb{Z}^3$, $L\mathbb{Z}^3$, each of which is contained in $\mathbb{L}(\eta)$, put
\[
\left(\boldsymbol{\Upsilon}_{\boldsymbol{\lambda}} f\right)(x) := \exp\left(i \cdot e_k \cdot \boldsymbol{\lambda}(x)\right) \cdot f(x)
\]
at every point $x$ of that lattice. Then
\begin{align*}
D_{B - d\boldsymbol{\lambda}} &= \boldsymbol{\Upsilon}_{\boldsymbol{\lambda}} \cdot D_B \cdot \boldsymbol{\Upsilon}_{\boldsymbol{\lambda}}^{-1},\\
Q_k(B - d\boldsymbol{\lambda}) &= \boldsymbol{\Upsilon}_{\boldsymbol{\lambda}} \cdot Q_k(B) \cdot \boldsymbol{\Upsilon}_{\boldsymbol{\lambda}}^{-1},\\
Q_k^T\left(-(B - d\boldsymbol{\lambda})\right) &= \boldsymbol{\Upsilon}_{\boldsymbol{\lambda}} \cdot Q_k^T(-B) \cdot \boldsymbol{\Upsilon}_{\boldsymbol{\lambda}}^{-1},\\
Q(B - d\boldsymbol{\lambda}) &= \boldsymbol{\Upsilon}_{\boldsymbol{\lambda}} \cdot Q(B) \cdot \boldsymbol{\Upsilon}_{\boldsymbol{\lambda}}^{-1},\\
Q^T\left(-(B - d\boldsymbol{\lambda})\right) &= \boldsymbol{\Upsilon}_{\boldsymbol{\lambda}} \cdot Q^T(-B) \cdot \boldsymbol{\Upsilon}_{\boldsymbol{\lambda}}^{-1}.
\end{align*}
Consequently $O(B - d\boldsymbol{\lambda}) = \boldsymbol{\Upsilon}_{\boldsymbol{\lambda}} \cdot O(B) \cdot \boldsymbol{\Upsilon}_{\boldsymbol{\lambda}}^{-1}$ for the other operators of Section~\ref{sec:block-operators} without a bar, namely for $O = K, P_k, P, X_k, S_k, H_k, \mathfrak{H}_k, D_k, M_k, \dot{M}_k, T_k, \Gamma_k$, whereas the barred operators satisfy
\begin{align*}
\bar{H}_k(B - d\boldsymbol{\lambda}) &= \boldsymbol{\Upsilon}_{\boldsymbol{\lambda}}^{-1} \cdot \bar{H}_k(B) \cdot \boldsymbol{\Upsilon}_{\boldsymbol{\lambda}},\\
\bar{T}_k(B - d\boldsymbol{\lambda}) &= \boldsymbol{\Upsilon}_{\boldsymbol{\lambda}}^{-1} \cdot \bar{T}_k(B) \cdot \boldsymbol{\Upsilon}_{\boldsymbol{\lambda}}.
\end{align*}
Moreover $\delta Z_k(B - d\boldsymbol{\lambda}) = \delta Z_k(B)$. Finally, for every pair $(x,x') \in \mathrm{Pz}_k$ and every spinor function or field $f$ on $T_k$,
\[
\left(\delta_{\alpha, B - d\boldsymbol{\lambda}} \cdot \boldsymbol{\Upsilon}_{\boldsymbol{\lambda}} f\right)(x,x') = \exp\left(i \cdot e_k \cdot \boldsymbol{\lambda}(x)\right) \cdot \left(\delta_{\alpha,B} f\right)(x,x').
\]
\item[(e)] \emph{Adjoints for real $B$.} Let $B$ be real, and let the adjoints be taken with respect to the Hermitian pairings of Section~\ref{sec:spinor-spaces-berezin}. Then $Q_k(B)^{*} = Q_k^T(-B)$ and $Q(B)^{*} = Q^T(-B)$; in particular $P_k(B)$ and $P(B)$ are orthogonal projections. Moreover $D_B^{*} = \gamma_3 \cdot D_B \cdot \gamma_3$ and $\operatorname{Re} D_B \ge 0$.
\item[(f)] \emph{Change of scale.} Let $B'$ be a complex, $(l/L)$-periodic function on the oriented bonds of $\mathbb{L}(\eta_{k+1})$, odd under reversal, and suppose that $B = B'_L$. Then the identity
\[
e_k \cdot \eta_k \cdot B(b) = e_{k+1} \cdot \eta_{k+1} \cdot B'(b/L)
\]
holds at every bond $b$ of $\mathbb{L}(\eta_k)$. Let the operators with the argument $B'$ be formed at level $k+1$, let $Q_{k+1}(B)$ and $Q_{k+1}^T(-B)$ be the operators of~(b), and put $K_{k+1}(B') := D_{B'} + L \cdot \mathfrak{m}$. For $f \in \mathfrak{S}(\eta_{k+1}, l/L)$ and $v \in \mathfrak{S}(1, l/L)$,
\begin{align*}
K(B) f_L &= L^{-1} \cdot \left(K_{k+1}(B') f\right)_L,\\
Q_{k+1}(B) f_L &= \left(Q_{k+1}(B') f\right)_L,\\
Q_{k+1}^T(-B) v_L &= \left(Q_{k+1}^T(-B') v\right)_L,
\end{align*}
and $\left(\delta_{\alpha,B} f_L\right)(x,x') = \left(\delta_{\alpha,B'} f\right)_L(x,x')$ for every $(x,x') \in \mathrm{Pz}_k$. For $a > 0$ with $l/L \in a\mathbb{N}$ and $g, f \in \mathfrak{S}(a, l/L)$, and in the same way for fields on $\mathbb{T}(a, l/L)$,
\[
\left\langle g_L, f_L \right\rangle_{La} = L \cdot \langle g, f \rangle_a.
\]
\end{enumerate}
\end{lemma}

\begin{proof}
(a) Let $v \in \mathfrak{S}(1,l)$ and $y \in \mathbb{Z}^3$. For $x \in B_k(y)$ one has $[x] = y$ and $\boldsymbol{\phi}^B_k(y,x) \cdot \boldsymbol{\phi}^{-B}_k(y,x) = 1$, since $\tau_k$ is linear in $B$; moreover $\eta^3 \cdot \#|B_k(y)| = 1$, because $B_k(y)$ has $L^{3k} = \eta^{-3}$ points. Hence
\begin{align*}
\left(Q_k(B) \cdot Q_k^T(-B) v\right)(y) &= \eta^3 \cdot \sum_{x \in B_k(y)} \boldsymbol{\phi}^B_k(y,x) \cdot \boldsymbol{\phi}^{-B}_k(y,x) \cdot v(y)\\
&= \eta^3 \cdot \#|B_k(y)| \cdot v(y) = v(y).
\end{align*}
In the same way the relations $\boldsymbol{\phi}^B(z,y) \cdot \boldsymbol{\phi}^{-B}(z,y) = 1$ and $L^{-3} \cdot \#|B(z)| = 1$, for $z \in L\mathbb{Z}^3$ and $y$ in the block $B(z)$, give $Q(B) \cdot Q^T(-B) = 1$. In the formulas that define $Q_k(\pm B)$, $Q_k^T(\pm B)$, $Q(\pm B)$ and $Q^T(\pm B)$ the values of a spinor are only multiplied by numbers and added; hence these operators, and with them $P_k(B)$ and $P(B)$, commute with every $4 \times 4$ matrix acting on the spinor index. By the two identities just proved,
\begin{align*}
P_k(B)^2 &= Q_k^T(-B) \cdot \left(Q_k(B) \cdot Q_k^T(-B)\right) \cdot Q_k(B) = P_k(B),\\
P(B)^2 &= Q^T(-B) \cdot \left(Q(B) \cdot Q^T(-B)\right) \cdot Q(B) = P(B),\\
P(B) \cdot Q^T(-B) &= Q^T(-B) \cdot \left(Q(B) \cdot Q^T(-B)\right) = Q^T(-B),\\
Q(B) \cdot P(B) &= \left(Q(B) \cdot Q^T(-B)\right) \cdot Q(B) = Q(B).
\end{align*}
To compute the traces, let $z \in L\mathbb{Z}^3$ and $y \in B(z)$. As $\left(Q^T(-B) w\right)(y) = \boldsymbol{\phi}^{-B}(z,y) \cdot w(z)$ for $w \in \mathfrak{S}(L,l)$, the definition of $Q(B)$ gives
\[
\left(P(B) v\right)(y) = L^{-3} \cdot \sum_{y' \in B(z)} \boldsymbol{\phi}^{-B}(z,y) \cdot \boldsymbol{\phi}^B(z,y') \cdot v(y').
\]
Distinct points of $B(z)$ are not congruent modulo $l\mathbb{Z}^3$, their distance being at most $L - 1 < l$. Take as coordinates on $\mathfrak{S}(1,l)$ the values $v(y) \in \mathbb{C}^4$ at the $n_k$ points $y$ of a fundamental domain of $\mathbb{Z}^3$ for $l\mathbb{Z}^3$. In these coordinates the $4 \times 4$ block of $P(B)$ on the diagonal at such a point $y$ comes from the term $y' = y$ of the last sum alone, in which the two phases cancel; so this block is $L^{-3} \cdot 1_4$, where $1_4$ is the unit matrix of $\mathbb{C}^4$, and the corresponding block of $\mathrm{Mx} \cdot P(B)$ is $L^{-3} \cdot \mathrm{Mx}$. Adding the traces of these blocks over the $n_k$ points, one obtains
\[
\operatorname{Tr}\left[\mathrm{Mx} \cdot P(B)\right] = n_k \cdot L^{-3} \cdot \operatorname{tr}\mathrm{Mx} = n_{k+1} \cdot \operatorname{tr}\mathrm{Mx}.
\]
The second trace identity follows, as $\operatorname{Tr}\mathrm{Mx} = n_k \cdot \operatorname{tr}\mathrm{Mx}$. Finally, by the definitions of $T_k$, $H_k$ and $S_k$,
\begin{align*}
T_k(B) \cdot H_k(B) &= Q_k(B) \cdot X_k(B) \cdot S_k(B) \cdot Q_k^T(-B)\\
&= Q_k(B) \cdot Q_k^T(-B) = 1.
\end{align*}

(b) Let $z \in L\mathbb{Z}^3$ and $f \in \mathfrak{S}(\eta,l)$. Inserting the definition of $Q_k(B)$ into that of $Q(B)$, one finds
\[
\left(Q(B) \cdot Q_k(B) f\right)(z) = L^{-3} \cdot \eta^3 \cdot \sum_{y \in B(z)} \sum_{x \in B_k(y)} \boldsymbol{\phi}^B(z,y) \cdot \boldsymbol{\phi}^B_k(y,x) \cdot f(x).
\]
For $y \in B(z)$ and $x \in B_k(y)$ one has $y = [x]$, and therefore, by the definition of $\tau_{k+1}$,
\begin{align*}
\boldsymbol{\phi}^B(z,y) \cdot \boldsymbol{\phi}^B_k(y,x) &= \exp\left\{ i \cdot e_k \cdot \eta \cdot \left( (\tau B)(z,[x]) + (\tau_k B)([x],x) \right) \right\}\\
&= \exp\left\{ i \cdot e_k \cdot \eta \cdot (\tau_{k+1} B)(z,x) \right\}.
\end{align*}
By \eqref{eq:algebraic-gaussian-2}, applied in $\mathbb{L}(\eta)$, the blocks $B_k(y)$ with $y \in B(z)$ partition $B_{k+1}(z)$, so that the double sum is a single sum over $x \in B_{k+1}(z)$. This proves the formula for $Q_{k+1}(B)$. The identity for $P_{k+1}(B)$ follows from the definitions:
\begin{align*}
Q_{k+1}^T(-B) \cdot Q_{k+1}(B) &= Q_k^T(-B) \cdot Q^T(-B) \cdot Q(B) \cdot Q_k(B)\\
&= Q_k^T(-B) \cdot P(B) \cdot Q_k(B).
\end{align*}

(c) Let $g, f \in \mathfrak{S}(\eta,l)$, let $\mathscr{F}$ be a fundamental domain of $\mathbb{L}(\eta)$ for $l\mathbb{Z}^3$, and for $h, h' \in \mathbb{C}^4$ write $h^T \cdot h' := \sum_{\boldsymbol{\beta}} h_{\boldsymbol{\beta}} \cdot h'_{\boldsymbol{\beta}}$. For $\mu = 0, 1, 2$ and $x \in \mathbb{L}(\eta)$ put
\[
t^{\mp}_{\mu}(x) := g(x)^T \cdot P^{\mp}_{\mu} \cdot u_B(x, x \pm \eta e_{\mu}) \cdot f(x \pm \eta e_{\mu}).
\]
By the definition of $D_B$,
\[
\langle g, D_B f \rangle_{\eta} = -\eta^{-1} \cdot \sum_{x \in \mathscr{F}} \eta^3 \cdot \sum_{\mu} \left( t^{-}_{\mu}(x) + t^{+}_{\mu}(x) - g(x)^T \cdot f(x) \right).
\]
Each $t^{\mp}_{\mu}$ is an $l$-periodic function of $x$, as $g$ and $f$ are antiperiodic and $B$ is periodic; so its sum over $\mathscr{F}$ is unchanged if $x$ is replaced by $x \mp \eta e_{\mu}$. Because $B$ is odd under reversal, $u_B(x \mp \eta e_{\mu}, x) = u_{-B}(x, x \mp \eta e_{\mu})$, and therefore
\begin{align*}
\sum_{x \in \mathscr{F}} t^{\mp}_{\mu}(x) &= \sum_{x \in \mathscr{F}} s^{\mp}_{\mu}(x)^T \cdot f(x),\\
s^{\mp}_{\mu}(x) &:= \left(P^{\mp}_{\mu}\right)^T \cdot u_{-B}(x, x \mp \eta e_{\mu}) \cdot g(x \mp \eta e_{\mu}).
\end{align*}
It follows that $\langle g, D_B f \rangle_{\eta} = \langle D_B^T g, f \rangle_{\eta}$ with
\[
\left(D_B^T g\right)(x) = -\eta^{-1} \cdot \sum_{\mu} \left( s^{-}_{\mu}(x) + s^{+}_{\mu}(x) - g(x) \right).
\]
Since $P^{\mp}_{\mu} = \frac{1}{2}\left(1 \mp \gamma_{\mu}\right)$, the relation \eqref{eq:algebraic-gaussian-17} gives $\left(P^{\mp}_{\mu}\right)^T = \frac{1}{2}\left(1 \mp \gamma_{\mu}^T\right) = \mathsf{C}^{-1} \cdot P^{\pm}_{\mu} \cdot \mathsf{C}$. Hence
\begin{align*}
s^{-}_{\mu}(x) &= \mathsf{C}^{-1} \cdot u_{-B}(x, x - \eta e_{\mu}) \cdot P^{+}_{\mu} \cdot \left(\mathsf{C} g\right)(x - \eta e_{\mu}),\\
s^{+}_{\mu}(x) &= \mathsf{C}^{-1} \cdot u_{-B}(x, x + \eta e_{\mu}) \cdot P^{-}_{\mu} \cdot \left(\mathsf{C} g\right)(x + \eta e_{\mu}),
\end{align*}
and comparison with the definition of $D_{-B}$, applied to $\mathsf{C} g$, yields $D_B^T g = \mathsf{C}^{-1} \cdot D_{-B}\left(\mathsf{C} g\right)$. Next, $Q_k^T(\pm B)$ is the transpose of $Q_k(\pm B)$ (Section~\ref{sec:block-operators}); hence
\[
P_k(B)^T = Q_k(B)^T \cdot \left(Q_k^T(-B)\right)^T = Q_k^T(B) \cdot Q_k(-B) = P_k(-B),
\]
and in the same way $P(B)^T = P(-B)$.

The operators $X_k$, $S_k$, $D_k$, $M_k$ and $\Gamma_k$ are built from $K$, $P_k$, $P$ and from expressions $Q_k(B) \cdot O \cdot Q_k^T(-B)$ by sums, products and inverses. Transposition reverses the order of the factors of a product and commutes with sums and inverses; for the expressions just named
\[
\left(Q_k(B) \cdot O \cdot Q_k^T(-B)\right)^T = Q_k(-B) \cdot O^T \cdot Q_k^T(B);
\]
and $Q_k(\pm B)$, $Q_k^T(\pm B)$, $Q(\pm B)$, $Q^T(\pm B)$ commute with $\mathsf{C}$, by the proof of~(a). Explicitly: $K(B)^T = D_B^T + \mathfrak{m} = \mathsf{C}^{-1} \cdot K(-B) \cdot \mathsf{C}$; therefore $X_k(B)^T = K(B)^T + b_k \cdot P_k(-B) = \mathsf{C}^{-1} \cdot X_k(-B) \cdot \mathsf{C}$ and $S_k(B)^T = \left(X_k(B)^T\right)^{-1} = \mathsf{C}^{-1} \cdot S_k(-B) \cdot \mathsf{C}$; then
\begin{align*}
D_k(B)^T &= b_k - b_k^2 \cdot Q_k(-B) \cdot S_k(B)^T \cdot Q_k^T(B)\\
&= \mathsf{C}^{-1} \cdot \left( b_k - b_k^2 \cdot Q_k(-B) \cdot S_k(-B) \cdot Q_k^T\left(-(-B)\right) \right) \cdot \mathsf{C}\\
&= \mathsf{C}^{-1} \cdot D_k(-B) \cdot \mathsf{C},
\end{align*}
the same relation holds for $M_k = b_k - D_k$, and
\[
\Gamma_k(B)^T = \left(D_k(B)^T + \tilde{b} \cdot P(B)^T\right)^{-1} = \mathsf{C}^{-1} \cdot \Gamma_k(-B) \cdot \mathsf{C}.
\]
Further,
\begin{align*}
\mathfrak{H}_k(B)^T &= b_k \cdot S_k(B)^T \cdot Q_k^T(B) = \mathsf{C}^{-1} \cdot b_k \cdot S_k(-B) \cdot Q_k^T\left(-(-B)\right) \cdot \mathsf{C}\\
&= \mathsf{C}^{-1} \cdot H_k(-B) \cdot \mathsf{C} = \bar{H}_k(B),\\
\bar{T}_k(B) &= \mathsf{C}^{-1} \cdot T_k(-B) \cdot \mathsf{C} = b_k^{-1} \cdot Q_k(-B) \cdot \mathsf{C}^{-1} \cdot X_k(-B) \cdot \mathsf{C}\\
&= b_k^{-1} \cdot Q_k(-B) \cdot X_k(B)^T,
\end{align*}
and therefore, by~(a) with $-B$ in place of $B$,
\begin{align*}
\bar{T}_k(B) \cdot \bar{H}_k(B) &= Q_k(-B) \cdot X_k(B)^T \cdot \left(X_k(B)^T\right)^{-1} \cdot Q_k^T(B)\\
&= Q_k(-B) \cdot Q_k^T(B) = 1.
\end{align*}
Finally $\dot{M}_k(B) = \mathfrak{H}_k(B) \cdot H_k(B)$, and the rule $\langle g, O f\rangle_{a'} = \langle O^T g, f\rangle_a$, which holds for fields (Section~\ref{sec:spinor-spaces-berezin}), gives with $O = \mathfrak{H}_k(B)$
\begin{align*}
\left\langle \bar{\Psi}, \dot{M}_k(B) \Psi \right\rangle_1 &= \left\langle \mathfrak{H}_k(B)^T \bar{\Psi}, H_k(B) \Psi \right\rangle_{\eta}\\
&= \left\langle \bar{H}_k(B) \bar{\Psi}, H_k(B) \Psi \right\rangle_{\eta}.
\end{align*}

(d) Since $\exp\left(i \cdot e_k \cdot \boldsymbol{\lambda}\right)$ is $l$-periodic, $\boldsymbol{\Upsilon}_{\boldsymbol{\lambda}}$ maps antiperiodic functions to antiperiodic functions; since $d\boldsymbol{\lambda}$ is $l$-periodic, $B - d\boldsymbol{\lambda}$ is again a complex, $l$-periodic function on the oriented bonds of $\mathbb{L}(\eta)$, odd under reversal. Let $\boldsymbol{\Gamma}$ be a path in $\mathbb{L}(\eta)$ from $x$ to $x'$, formed by the bonds $(x_0,x_1), \dots, (x_{m-1},x_m)$, where $x_0 = x$ and $x_m = x'$. By the definition of $d\boldsymbol{\lambda}$ the contribution of each bond is a difference of two values of $\boldsymbol{\lambda}$, and these differences add up along the path, as in \eqref{eq:algebraic-gaussian-S1}; neither reality nor periodicity of $\boldsymbol{\lambda}$ enters:
\begin{align*}
\eta \cdot (d\boldsymbol{\lambda})(\boldsymbol{\Gamma}) &= \sum_{i=0}^{m-1} \eta \cdot (d\boldsymbol{\lambda})(x_i,x_{i+1})\\
&= \sum_{i=0}^{m-1} \left( \boldsymbol{\lambda}(x_{i+1}) - \boldsymbol{\lambda}(x_i) \right) = \boldsymbol{\lambda}(x') - \boldsymbol{\lambda}(x).
\end{align*}
For a single bond $(x,x')$ this reads $\eta \cdot (d\boldsymbol{\lambda})(x,x') = \boldsymbol{\lambda}(x') - \boldsymbol{\lambda}(x)$, so that $u_{B - d\boldsymbol{\lambda}}(x,x') = u_B(x,x') \cdot \exp\left\{ -i \cdot e_k \cdot \left(\boldsymbol{\lambda}(x') - \boldsymbol{\lambda}(x)\right) \right\}$, that is,
\[
u_{B - d\boldsymbol{\lambda}}(x,x') \cdot \exp\left(i \cdot e_k \cdot \boldsymbol{\lambda}(x')\right) = \exp\left(i \cdot e_k \cdot \boldsymbol{\lambda}(x)\right) \cdot u_B(x,x').
\]
Inserted into the definition of $D_B$, this yields $D_{B - d\boldsymbol{\lambda}} \cdot \boldsymbol{\Upsilon}_{\boldsymbol{\lambda}} f = \boldsymbol{\Upsilon}_{\boldsymbol{\lambda}} \cdot D_B f$ for every $f \in \mathfrak{S}(\eta,l)$.

Let $y \in \mathbb{Z}^3$, $x \in B_k(y)$, and let $y_0 = x, y_1, \dots, y_k = y$ be the points of Section~\ref{sec:block-operators}. For every permutation $\pi$ the rectilinear path $\boldsymbol{\Gamma}^{\pi}(y_{j+1}, y_j)$, regarded as a path in $\mathbb{L}(\eta)$, leads from $y_{j+1}$ to $y_j$. Averaging the path identity over $\pi$ gives $\eta \cdot (\tau d\boldsymbol{\lambda})(y_{j+1}, y_j) = \boldsymbol{\lambda}(y_j) - \boldsymbol{\lambda}(y_{j+1})$, and the sum over $j$ telescopes:
\[
\eta \cdot (\tau_k d\boldsymbol{\lambda})(y,x) = \sum_{j=0}^{k-1} \left( \boldsymbol{\lambda}(y_j) - \boldsymbol{\lambda}(y_{j+1}) \right) = \boldsymbol{\lambda}(x) - \boldsymbol{\lambda}(y).
\]
As $\tau_k$ is linear, it follows that
\[
\boldsymbol{\phi}^{B - d\boldsymbol{\lambda}}_k(y,x) \cdot \exp\left(i \cdot e_k \cdot \boldsymbol{\lambda}(x)\right) = \exp\left(i \cdot e_k \cdot \boldsymbol{\lambda}(y)\right) \cdot \boldsymbol{\phi}^B_k(y,x).
\]
For $z \in L\mathbb{Z}^3$ and $y \in B(z)$ the same argument, applied to the six paths that define $(\tau B)(z,y)$, gives
\[
\boldsymbol{\phi}^{B - d\boldsymbol{\lambda}}(z,y) \cdot \exp\left(i \cdot e_k \cdot \boldsymbol{\lambda}(y)\right) = \exp\left(i \cdot e_k \cdot \boldsymbol{\lambda}(z)\right) \cdot \boldsymbol{\phi}^B(z,y).
\]
Let $(x,x') \in \mathrm{Pz}_k$. The six rectilinear paths that define $(\tau B)(x,x')$ in Definition~\ref{def:pairs-sites-holder} lead from $x$ to $x'$, the two points being taken in the shortest lift of the pair; thus $\eta \cdot (\tau d\boldsymbol{\lambda})(x,x') = \boldsymbol{\lambda}(x') - \boldsymbol{\lambda}(x)$ at these points, and
\begin{align*}
&\exp\left\{ i \cdot e_k \cdot \eta \cdot \left(\tau (B - d\boldsymbol{\lambda})\right)(x,x') \right\} \cdot \exp\left(i \cdot e_k \cdot \boldsymbol{\lambda}(x')\right)\\
&\qquad = \exp\left(i \cdot e_k \cdot \boldsymbol{\lambda}(x)\right) \cdot \exp\left\{ i \cdot e_k \cdot \eta \cdot (\tau B)(x,x') \right\};
\end{align*}
both sides are unchanged when the lift is translated by a vector of $l\mathbb{Z}^3$, because $\exp\left(i \cdot e_k \cdot \boldsymbol{\lambda}\right)$ is $l$-periodic.

With these relations for the phases, the definitions of $Q_k$ and $Q$ give, for $f \in \mathfrak{S}(\eta,l)$,
\begin{align*}
\left(Q_k(B - d\boldsymbol{\lambda}) \cdot \boldsymbol{\Upsilon}_{\boldsymbol{\lambda}} f\right)(y) &= \eta^3 \cdot \sum_{x \in B_k(y)} \boldsymbol{\phi}^{B - d\boldsymbol{\lambda}}_k(y,x) \cdot \exp\left(i \cdot e_k \cdot \boldsymbol{\lambda}(x)\right) \cdot f(x)\\
&= \exp\left(i \cdot e_k \cdot \boldsymbol{\lambda}(y)\right) \cdot \left(Q_k(B) f\right)(y),
\end{align*}
and in the same way $Q(B - d\boldsymbol{\lambda}) \cdot \boldsymbol{\Upsilon}_{\boldsymbol{\lambda}} = \boldsymbol{\Upsilon}_{\boldsymbol{\lambda}} \cdot Q(B)$. Since $\tau$ and $\tau_k$ are linear, the phases formed with $-A$ are the reciprocals of those formed with $A$, for every bond function $A$. Taking reciprocals in the relation for $\boldsymbol{\phi}_k$ one therefore obtains
\[
\boldsymbol{\phi}^{-(B - d\boldsymbol{\lambda})}_k(y,x) \cdot \exp\left(i \cdot e_k \cdot \boldsymbol{\lambda}(y)\right) = \exp\left(i \cdot e_k \cdot \boldsymbol{\lambda}(x)\right) \cdot \boldsymbol{\phi}^{-B}_k(y,x),
\]
and the analogous relation for $\boldsymbol{\phi}$. Hence, for $v \in \mathfrak{S}(1,l)$ and $x \in \mathbb{L}(\eta)$,
\begin{align*}
\left(Q_k^T\left(-(B - d\boldsymbol{\lambda})\right) \cdot \boldsymbol{\Upsilon}_{\boldsymbol{\lambda}} v\right)(x) &= \boldsymbol{\phi}^{-(B - d\boldsymbol{\lambda})}_k([x],x) \cdot \exp\left(i \cdot e_k \cdot \boldsymbol{\lambda}([x])\right) \cdot v([x])\\
&= \exp\left(i \cdot e_k \cdot \boldsymbol{\lambda}(x)\right) \cdot \left(Q_k^T(-B) v\right)(x),
\end{align*}
and in the same way $Q^T\left(-(B - d\boldsymbol{\lambda})\right) \cdot \boldsymbol{\Upsilon}_{\boldsymbol{\lambda}} = \boldsymbol{\Upsilon}_{\boldsymbol{\lambda}} \cdot Q^T(-B)$. This proves the five relations displayed in~(d).

The other operators without a bar are formed from $D_B$, $Q_k(B)$, $Q_k^T(-B)$, $Q(B)$, $Q^T(-B)$ and numbers by sums, products and inverses, and conjugation by $\boldsymbol{\Upsilon}_{\boldsymbol{\lambda}}$ respects these operations; this gives $O(B - d\boldsymbol{\lambda}) = \boldsymbol{\Upsilon}_{\boldsymbol{\lambda}} \cdot O(B) \cdot \boldsymbol{\Upsilon}_{\boldsymbol{\lambda}}^{-1}$ for them, and
\[
\delta Z_k(B - d\boldsymbol{\lambda}) = \det\left( \boldsymbol{\Upsilon}_{\boldsymbol{\lambda}} \cdot \left( D_k(B) + \tilde{b} \cdot P(B) \right) \cdot \boldsymbol{\Upsilon}_{\boldsymbol{\lambda}}^{-1} \right) = \delta Z_k(B).
\]
The function $-\boldsymbol{\lambda}$ satisfies the hypotheses of~(d), $\boldsymbol{\Upsilon}_{-\boldsymbol{\lambda}} = \boldsymbol{\Upsilon}_{\boldsymbol{\lambda}}^{-1}$, $-(B - d\boldsymbol{\lambda}) = (-B) - d(-\boldsymbol{\lambda})$, and $\boldsymbol{\Upsilon}_{\boldsymbol{\lambda}}$ commutes with $\mathsf{C}$. Hence
\begin{align*}
\bar{H}_k(B - d\boldsymbol{\lambda}) &= \mathsf{C}^{-1} \cdot H_k\left((-B) - d(-\boldsymbol{\lambda})\right) \cdot \mathsf{C}\\
&= \mathsf{C}^{-1} \cdot \boldsymbol{\Upsilon}_{\boldsymbol{\lambda}}^{-1} \cdot H_k(-B) \cdot \boldsymbol{\Upsilon}_{\boldsymbol{\lambda}} \cdot \mathsf{C} = \boldsymbol{\Upsilon}_{\boldsymbol{\lambda}}^{-1} \cdot \bar{H}_k(B) \cdot \boldsymbol{\Upsilon}_{\boldsymbol{\lambda}},
\end{align*}
and in the same way $\bar{T}_k(B - d\boldsymbol{\lambda}) = \boldsymbol{\Upsilon}_{\boldsymbol{\lambda}}^{-1} \cdot \bar{T}_k(B) \cdot \boldsymbol{\Upsilon}_{\boldsymbol{\lambda}}$.

Finally, let $(x,x') \in \mathrm{Pz}_k$, and let $f$ be a spinor function or a field on $T_k$. Multiplying the relation for the phase of the pair by $f(x')$, one obtains
\begin{align*}
&\exp\left\{ i \cdot e_k \cdot \eta \cdot \left(\tau (B - d\boldsymbol{\lambda})\right)(x,x') \right\} \cdot \left(\boldsymbol{\Upsilon}_{\boldsymbol{\lambda}} f\right)(x')\\
&\qquad = \exp\left(i \cdot e_k \cdot \boldsymbol{\lambda}(x)\right) \cdot \exp\left\{ i \cdot e_k \cdot \eta \cdot (\tau B)(x,x') \right\} \cdot f(x'),
\end{align*}
while $\left(\boldsymbol{\Upsilon}_{\boldsymbol{\lambda}} f\right)(x) = \exp\left(i \cdot e_k \cdot \boldsymbol{\lambda}(x)\right) \cdot f(x)$. By the definition of $\delta_{\alpha,B}$ in Definition~\ref{def:pairs-sites-holder}, the difference of the two left-hand sides, multiplied by $|x - x'|^{-\alpha}$, is $\left(\delta_{\alpha, B - d\boldsymbol{\lambda}} \cdot \boldsymbol{\Upsilon}_{\boldsymbol{\lambda}} f\right)(x,x')$, and the difference of the two right-hand sides, multiplied by $|x - x'|^{-\alpha}$, is $\exp\left(i \cdot e_k \cdot \boldsymbol{\lambda}(x)\right) \cdot \left(\delta_{\alpha,B} f\right)(x,x')$.

(e) Let $B$ be real. Then $|u_B(x,x')| = 1$ and $\overline{u_B(x,x')} = u_B(x',x)$ for every bond $(x,x')$. For $\mu = 0, 1, 2$ define the operator $U_{\mu}$ on $\mathfrak{S}(\eta,l)$ by
\[
\left(U_{\mu} f\right)(x) := u_B(x, x + \eta e_{\mu}) \cdot f(x + \eta e_{\mu}).
\]
Substituting $x - \eta e_{\mu}$ for $x$ in the sum that defines $(g, U_{\mu} f)_{\eta}$ and using $\overline{u_B(x - \eta e_{\mu}, x)} = u_B(x, x - \eta e_{\mu})$, one obtains
\[
\left(U_{\mu}^{*} f\right)(x) = u_B(x, x - \eta e_{\mu}) \cdot f(x - \eta e_{\mu});
\]
moreover $U_{\mu} \cdot U_{\mu}^{*} = U_{\mu}^{*} \cdot U_{\mu} = 1$, as $u_B(x,x') \cdot u_B(x',x) = 1$. Thus $U_{\mu}$ is unitary. As $P^{\mp}_{\mu} = \frac{1}{2}\left(1 \mp \gamma_{\mu}\right)$, the definition of $D_B$ can be written
\[
D_B = \eta^{-1} \cdot \sum_{\mu} \left( 1 - \frac{1}{2} \cdot \left(U_{\mu} + U_{\mu}^{*}\right) + \frac{1}{2} \cdot \gamma_{\mu} \cdot \left(U_{\mu} - U_{\mu}^{*}\right) \right).
\]
The first part, $\eta^{-1} \cdot \sum_{\mu} \left(1 - \frac{1}{2} \cdot \left(U_{\mu} + U_{\mu}^{*}\right)\right)$, is self-adjoint and $\ge 0$, because $\operatorname{Re}(f, U_{\mu} f)_{\eta} \le (f,f)_{\eta}$ by the Cauchy--Schwarz inequality; it commutes with $\gamma_3$, since $U_{\mu}$ and $U_{\mu}^{*}$ multiply the values of a spinor by numbers. The second part, $\eta^{-1} \cdot \sum_{\mu} \frac{1}{2} \cdot \gamma_{\mu} \cdot \left(U_{\mu} - U_{\mu}^{*}\right)$, is anti-self-adjoint, because $\gamma_{\mu}$ is Hermitian and commutes with $U_{\mu}$; and it anticommutes with $\gamma_3$, because $\gamma_{\mu} \cdot \gamma_3 = -\gamma_3 \cdot \gamma_{\mu}$ for $\mu = 0, 1, 2$ (Section~\ref{sec:spinor-spaces-berezin}). Since $\gamma_3^2 = 1$, it follows that $D_B^{*} = \gamma_3 \cdot D_B \cdot \gamma_3$, and that $\operatorname{Re} D_B$ is the first part; so $\operatorname{Re} D_B \ge 0$.

For real $B$ one has $\overline{\boldsymbol{\phi}^B_k(y,x)} = \boldsymbol{\phi}^{-B}_k(y,x)$ and $\overline{\boldsymbol{\phi}^B(z,y)} = \boldsymbol{\phi}^{-B}(z,y)$. Let $v \in \mathfrak{S}(1,l)$ and $f \in \mathfrak{S}(\eta,l)$. The union of the blocks $B_k(y)$, with $y$ in a fundamental domain of $\mathbb{Z}^3$ for $l\mathbb{Z}^3$, is a fundamental domain of $\mathbb{L}(\eta)$, because $B_k(y + l e_{\mu}) = B_k(y) + l e_{\mu}$; with $y = [x]$ it follows that
\begin{align*}
\left(v, Q_k(B) f\right)_1 &= \sum_y \sum_{x \in B_k(y)} \eta^3 \cdot \sum_{\boldsymbol{\beta}} \overline{\boldsymbol{\phi}^{-B}_k([x],x) \cdot v_{\boldsymbol{\beta}}([x])} \cdot f_{\boldsymbol{\beta}}(x)\\
&= \left(Q_k^T(-B) v, f\right)_{\eta},
\end{align*}
that is, $Q_k(B)^{*} = Q_k^T(-B)$; in the same way $Q(B)^{*} = Q^T(-B)$. Hence $P_k(B) = Q_k(B)^{*} \cdot Q_k(B)$ and $P(B) = Q(B)^{*} \cdot Q(B)$ are self-adjoint; being idempotent by~(a), they are orthogonal projections.

(f) By \eqref{eq:algebraic-gaussian-3}, $B(b) = L^{-\frac{1}{2}} \cdot B'(b/L)$ for every bond $b$ of $\mathbb{L}(\eta_k)$; and since $e_{k+1} = L^{\frac{1}{2}} \cdot e_k$ and $\eta_{k+1} = \eta_k/L$, one has $e_k \cdot \eta_k \cdot L^{-\frac{1}{2}} = e_{k+1} \cdot \eta_{k+1}$. This is the first assertion. Division by $L$ maps $\mathbb{L}(L^j\eta_k)$ onto $\mathbb{L}(L^j\eta_{k+1})$ for every $j \ge 0$, the blocks $B_j(y)$ of $\mathbb{L}(\eta_k)$ onto the blocks $B_j(y/L)$ of $\mathbb{L}(\eta_{k+1})$, and rectilinear paths onto rectilinear paths. By the first assertion, therefore,
\[
u_B(x, x \pm \eta_k e_{\mu}) = u_{B'}(x/L, x/L \pm \eta_{k+1} e_{\mu}) \qquad \left(x \in \mathbb{L}(\eta_k)\right),
\]
where $u_{B'}$ is formed at level $k+1$. Let $z \in L\mathbb{Z}^3$ and $x \in B_{k+1}(z)$, let $y_0 = x, \dots, y_k = [x]$ be the points of Section~\ref{sec:block-operators}, and put $y_{k+1} := z$. Then $(\tau_{k+1} B)(z,x) = \sum_{j=0}^k (\tau B)(y_{j+1}, y_j)$, and at level $k+1$ the number $(\tau_{k+1} B')(z/L, x/L)$ is the same sum, formed with $B'$ and with the points $y_j/L$. Each term scales by the first assertion, so that
\[
e_k \cdot \eta_k \cdot (\tau_{k+1} B)(z,x) = e_{k+1} \cdot \eta_{k+1} \cdot (\tau_{k+1} B')(z/L, x/L);
\]
on taking exponentials,
\[
\exp\left\{ i \cdot e_k \cdot \eta_k \cdot (\tau_{k+1} B)(z,x) \right\} = \boldsymbol{\phi}^{B'}_{k+1}(z/L, x/L),
\]
the phase of level $k+1$ formed with $B'$. Likewise, for $(x,x') \in \mathrm{Pz}_k$,
\[
e_k \cdot \eta_k \cdot (\tau B)(x,x') = e_{k+1} \cdot \eta_{k+1} \cdot (\tau B')(x/L, x'/L);
\]
here $(x/L, x'/L)$ is a pair of distinct sites of $T_{k+1}$ at distance at most $M/L \le M$, hence an element of $\mathrm{Pz}_{k+1}$, and its shortest lift is the shortest lift of $(x,x')$ divided by $L$, as $2M/L < l/L$.

Let $f \in \mathfrak{S}(\eta_{k+1}, l/L)$, let $x \in \mathbb{L}(\eta_k)$, and put $\tilde{x} := x/L$. Then
\[
f_L(x) = L^{-1} \cdot f(\tilde{x}), \qquad f_L(x \pm \eta_k e_{\mu}) = L^{-1} \cdot f(\tilde{x} \pm \eta_{k+1} e_{\mu}).
\]
Together with the relation for $u_B$ this shows that each of the three terms in the braces of the definition of $\left(D_B f_L\right)(x)$ is $L^{-1}$ times the corresponding term of $\left(D_{B'} f\right)(\tilde{x})$, formed at level $k+1$. The factor $-\eta_k^{-1}$ in front of the sum is $L^{-1}$ times the corresponding factor $-\eta_{k+1}^{-1}$ of level $k+1$. Hence
\[
\left(D_B f_L\right)(x) = L^{-2} \cdot \left(D_{B'} f\right)(\tilde{x}) = L^{-1} \cdot \left(D_{B'} f\right)_L(x).
\]
Together with $\mathfrak{m} \cdot f_L = L^{-1} \cdot \left(L \cdot \mathfrak{m} \cdot f\right)_L$ this gives $K(B) f_L = L^{-1} \cdot \left(K_{k+1}(B') f\right)_L$.

For $z \in L\mathbb{Z}^3$ the map $x \mapsto x/L$ is a bijection of $B_{k+1}(z)$ onto the block $B_{k+1}(z/L)$ of level $k+1$. By~(b), by the relation for the phases of~(b) and by $L^{-3} \cdot \eta_k^3 = \eta_{k+1}^3$,
\begin{align*}
\left(Q_{k+1}(B) f_L\right)(z) &= L^{-3} \cdot \eta_k^3 \cdot \sum_{x \in B_{k+1}(z)} \exp\left\{ i \cdot e_k \cdot \eta_k \cdot (\tau_{k+1} B)(z,x) \right\} \cdot f_L(x)\\
&= L^{-1} \cdot \eta_{k+1}^3 \cdot \sum_{x \in B_{k+1}(z)} \boldsymbol{\phi}^{B'}_{k+1}(z/L, x/L) \cdot f(x/L)\\
&= L^{-1} \cdot \left(Q_{k+1}(B') f\right)(z/L) = \left(Q_{k+1}(B') f\right)_L(z).
\end{align*}
Let $w \in \mathfrak{S}(L,l)$, $z \in L\mathbb{Z}^3$ and $x \in B_{k+1}(z)$; then $[x] \in B(z)$, and the definitions give
\begin{align*}
\left(Q_{k+1}^T(-B) w\right)(x) &= \boldsymbol{\phi}^{-B}_k([x],x) \cdot \boldsymbol{\phi}^{-B}(z,[x]) \cdot w(z)\\
&= \exp\left\{ -i \cdot e_k \cdot \eta_k \cdot (\tau_{k+1} B)(z,x) \right\} \cdot w(z).
\end{align*}
For $w = v_L$ with $v \in \mathfrak{S}(1, l/L)$ one has $w(z) = L^{-1} \cdot v(z/L)$. Since $z/L$ is the centre $[x/L]$ of the block of level $k+1$ that contains $x/L$, the relation for the phases of~(b), applied with $-B$ and $-B'$, gives
\begin{align*}
\left(Q_{k+1}^T(-B) v_L\right)(x) &= L^{-1} \cdot \boldsymbol{\phi}^{-B'}_{k+1}(z/L, x/L) \cdot v(z/L)\\
&= L^{-1} \cdot \left(Q_{k+1}^T(-B') v\right)(x/L) = \left(Q_{k+1}^T(-B') v\right)_L(x).
\end{align*}
Let $(x,x') \in \mathrm{Pz}_k$. Then
\begin{align*}
&\left(\delta_{\alpha,B} f_L\right)(x,x')\\
&\quad = |x - x'|^{-\alpha} \cdot \left( \exp\left\{ i \cdot e_k \cdot \eta_k \cdot (\tau B)(x,x') \right\} \cdot f_L(x') - f_L(x) \right)\\
&\quad = |x - x'|^{-\alpha} \cdot L^{-1} \cdot |x/L - x'/L|^{\alpha} \cdot \left(\delta_{\alpha,B'} f\right)(x/L, x'/L)\\
&\quad = L^{-1-\alpha} \cdot \left(\delta_{\alpha,B'} f\right)(x/L, x'/L) = \left(\delta_{\alpha,B'} f\right)_L(x,x').
\end{align*}
The first equality is the definition of $\delta_{\alpha,B}$. The second follows from $f_L(x) = L^{-1} \cdot f(x/L)$, $f_L(x') = L^{-1} \cdot f(x'/L)$, the relation for the phases of pairs and the definition of $\delta_{\alpha,B'}$ at level $k+1$, applied to the pair $(x/L, x'/L)$ of $\mathrm{Pz}_{k+1}$. The third holds because $|x - x'|^{-\alpha} \cdot L^{-1} = L^{-1-\alpha} \cdot |x/L - x'/L|^{-\alpha}$, and the last by the scaling $\chi_L(x,x') = L^{-1-\alpha} \cdot \chi(x/L, x'/L)$ of functions of pairs (Section~\ref{sec:block-operators}). Finally, let $a > 0$ with $l/L \in a\mathbb{N}$. The map $x \mapsto x/L$ carries a fundamental domain of $\mathbb{L}(La)$ for $l\mathbb{Z}^3$ onto a fundamental domain of $\mathbb{L}(a)$ for $(l/L)\mathbb{Z}^3$; the values of $g_L$ and $f_L$ carry the factor $L^{-1}$ each, and $(La)^3 \cdot L^{-2} = L \cdot a^3$. This gives $\langle g_L, f_L\rangle_{La} = L \cdot \langle g, f\rangle_a$, for spinor functions and for fields alike.
\end{proof}

\subsubsection{Schur complement identities}\label{sec:schur-complement-identities}

Throughout this subsection let $k\ge1$; the data and the bond function $B$ are those of Section~\ref{sec:block-operators}, and $\eta=\eta_k$. Numbers and $4\times4$ matrices act on spinors pointwise: $(Mf)(x):=M\cdot f(x)$.

For a real number $y\ge0$ let (cf.~\cite[(151), (152), (182)]{DimockMB})
\[
a_y(B):=b_k+\tilde{b}\cdot P(B)+i\gamma_3\cdot y\qquad\text{on }\mathfrak{S}(1,l).
\]
For every real number $c>0$ the matrix $c+i\gamma_3\cdot y$ is invertible, and
\[
\left(c+i\gamma_3\cdot y\right)^{-1}=\frac{c-i\gamma_3\cdot y}{c^2+y^2},
\]
because $\gamma_3^2=1$ (Section~\ref{sec:spinor-spaces-berezin}). This applies to $c=b_k$ and to $c=b_k+\tilde{b}$; both numbers are positive, since $\tilde{b}>0$, $b_1=b>0$ and $b_{k+1}=b\cdot b_k\cdot\left(b_k+\tilde{b}\right)^{-1}$ (Section~\ref{sec:block-operators}). One has
\[
a_y(B)=\left(b_k+i\gamma_3\cdot y\right)\cdot\left(1-P(B)\right)+\left(b_k+\tilde{b}+i\gamma_3\cdot y\right)\cdot P(B).
\]
As $P(B)$ is idempotent and commutes with all spinor matrices (Lemma~\ref{lem:identities-block-operators}(a)), it follows that $a_y(B)$ is invertible, with the inverse
\[
B_{k,y}(B):=a_y(B)^{-1}=\left(b_k+i\gamma_3\cdot y\right)^{-1}\cdot\left(1-P(B)\right)+\left(b_k+\tilde{b}+i\gamma_3\cdot y\right)^{-1}\cdot P(B),\qquad B_k:=B_{k,0}.
\]
Further let $\boldsymbol{\alpha}_{k,y}$ and $\boldsymbol{\beta}_{k,y}$ be the $4\times4$ matrices
\[
\boldsymbol{\alpha}_{k,y}:=b_k\cdot i\gamma_3\cdot y\cdot\left(b_k+i\gamma_3\cdot y\right)^{-1},\qquad
\boldsymbol{\beta}_{k,y}:=b_k^2\cdot\tilde{b}\cdot\left(b_k+\tilde{b}+i\gamma_3\cdot y\right)^{-1}\cdot\left(b_k+i\gamma_3\cdot y\right)^{-1},
\]
and let $X_{k,y}(B)$ be the operator on $\mathfrak{S}(\eta,l)$ given by
\[
X_{k,y}(B):=K(B)+\boldsymbol{\alpha}_{k,y}\cdot P_k(B)+\boldsymbol{\beta}_{k,y}\cdot P_{k+1}(B).
\]
The matrices $b_k+i\gamma_3\cdot y$ and $b_k+\tilde{b}+i\gamma_3\cdot y$, their inverses, $\boldsymbol{\alpha}_{k,y}$ and $\boldsymbol{\beta}_{k,y}$ are functions of $\gamma_3$. Hence they commute with one another, and they commute with $P(B)$, $P_k(B)$, $P_{k+1}(B)$, $Q_k(B)$ and $Q(B)$, the kernels of these five operators being scalar multiples of the unit matrix of $\mathbb{C}^4$ (Section~\ref{sec:block-operators} and Lemma~\ref{lem:identities-block-operators}(b)). For $y=0$ one finds
\[
\boldsymbol{\alpha}_{k,0}=0,\qquad\boldsymbol{\beta}_{k,0}=\frac{b_k\cdot\tilde{b}}{b_k+\tilde{b}}=b_{k+1}\cdot L^{-1},
\]
the second equality being the definition of $b_{k+1}$ together with $\tilde{b}=b\cdot L^{-1}$; accordingly one puts
\[
X^0_{k+1}(B):=X_{k,0}(B)=K(B)+b_{k+1}\cdot L^{-1}\cdot P_{k+1}(B).
\]

\begin{proposition}[{cf.~\cite[(65), (87), (88), (115), (116), (379)]{DimockUR}, and \cite[(181), Lemmas~18, 19]{DimockMB} in the geometry without large-field regions}]\label{prop:cf-geometry-large}
Let $k\ge1$, and let $X_k(B)$ be invertible. Then for every real number $y\ge0$ one has
\begin{equation}\label{eq:algebraic-gaussian-18}
X_{k,y}(B)=X_k(B)-b_k^2\cdot Q_k^T(-B)\cdot B_{k,y}(B)\cdot Q_k(B),
\end{equation}
and the following two conditions are equivalent:
\begin{enumerate}
\item[(i)] the operator $D_k(B)+\tilde{b}\cdot P(B)+i\gamma_3\cdot y$ on $\mathfrak{S}(1,l)$ is invertible;
\item[(ii)] the operator $X_{k,y}(B)$ on $\mathfrak{S}(\eta,l)$ is invertible.
\end{enumerate}
If they hold, put
\[
\Gamma_{k,y}:=\left(D_k+\tilde{b}\cdot P+i\gamma_3\cdot y\right)^{-1},\qquad S_{k,y}:=X_{k,y}^{-1},
\]
all operators being taken at $B$. Then
\begin{align}
\Gamma_{k,y}&=B_{k,y}+b_k^2\cdot B_{k,y}\cdot Q_k\cdot S_{k,y}\cdot Q_k^T(-B)\cdot B_{k,y},\label{eq:algebraic-gaussian-19}\\
S_{k,y}&=S_k+H_k\cdot\Gamma_{k,y}\cdot\mathfrak{H}_k,\label{eq:algebraic-gaussian-20}\\
H_k\cdot\Gamma_{k,y}&=b_k\cdot S_{k,y}\cdot Q_k^T(-B)\cdot B_{k,y},\qquad\Gamma_{k,y}\cdot\mathfrak{H}_k=b_k\cdot B_{k,y}\cdot Q_k\cdot S_{k,y},\label{eq:algebraic-gaussian-21}\\
\det\left(D_k+\tilde{b}\cdot P+i\gamma_3\cdot y\right)\cdot\det X_k&=\det a_y\cdot\det X_{k,y}.\label{eq:algebraic-gaussian-22}
\end{align}
For $y=0$ one has $\Gamma_{k,0}=\Gamma_k$ and $S_{k,0}=S^0_{k+1}(B):=X^0_{k+1}(B)^{-1}$.
\end{proposition}

\begin{proof}
All operators are taken at $B$.

\emph{The identity \eqref{eq:algebraic-gaussian-18}.} Put $g:=i\gamma_3\cdot y$ and $\kappa:=b_k\cdot\left(\tilde{b}+g\right)\cdot\left(b_k+\tilde{b}+g\right)^{-1}$. These matrices and $\left(b_k+g\right)^{-1}$ are functions of $\gamma_3$ and commute with one another. One has
\[
b_k-b_k^2\cdot\left(b_k+g\right)^{-1}=b_k\cdot g\cdot\left(b_k+g\right)^{-1}=\boldsymbol{\alpha}_{k,y},\qquad b_k-b_k^2\cdot\left(b_k+\tilde{b}+g\right)^{-1}=\kappa,
\]
and
\[
\kappa-\boldsymbol{\alpha}_{k,y}=b_k\cdot\left[\left(\tilde{b}+g\right)\cdot\left(b_k+g\right)-g\cdot\left(b_k+\tilde{b}+g\right)\right]\cdot\left(b_k+\tilde{b}+g\right)^{-1}\cdot\left(b_k+g\right)^{-1}=\boldsymbol{\beta}_{k,y},
\]
since the expression in square brackets equals $b_k\cdot\tilde{b}$. One has
\begin{align*}
b_k\cdot P_k-b_k^2\cdot Q_k^T(-B)\cdot B_{k,y}\cdot Q_k
&=Q_k^T(-B)\cdot\left[b_k-b_k^2\cdot\left(b_k+g\right)^{-1}\cdot(1-P)-b_k^2\cdot\left(b_k+\tilde{b}+g\right)^{-1}\cdot P\right]\cdot Q_k\\
&=Q_k^T(-B)\cdot\left[\boldsymbol{\alpha}_{k,y}\cdot(1-P)+\kappa\cdot P\right]\cdot Q_k
=Q_k^T(-B)\cdot\left[\boldsymbol{\alpha}_{k,y}+\boldsymbol{\beta}_{k,y}\cdot P\right]\cdot Q_k\\
&=\boldsymbol{\alpha}_{k,y}\cdot P_k+\boldsymbol{\beta}_{k,y}\cdot P_{k+1}.
\end{align*}
The first equality follows from the formula for $B_{k,y}$ and from $b_k\cdot P_k=Q_k^T(-B)\cdot b_k\cdot Q_k$, which is the definition of $P_k$. The second follows from $b_k=b_k\cdot(1-P)+b_k\cdot P$ on $\mathfrak{S}(1,l)$ and from the expressions for $\boldsymbol{\alpha}_{k,y}$ and $\kappa$ obtained above; the third from $\boldsymbol{\alpha}_{k,y}\cdot(1-P)+\kappa\cdot P=\boldsymbol{\alpha}_{k,y}+\left(\kappa-\boldsymbol{\alpha}_{k,y}\right)\cdot P$ and $\kappa-\boldsymbol{\alpha}_{k,y}=\boldsymbol{\beta}_{k,y}$. For the last one: as $\boldsymbol{\alpha}_{k,y}$ and $\boldsymbol{\beta}_{k,y}$ commute with $P$ and $Q_k$, the operator in question equals
\[
Q_k^T(-B)\cdot Q_k\cdot\boldsymbol{\alpha}_{k,y}+Q_k^T(-B)\cdot P\cdot Q_k\cdot\boldsymbol{\beta}_{k,y}=P_k\cdot\boldsymbol{\alpha}_{k,y}+P_{k+1}\cdot\boldsymbol{\beta}_{k,y},
\]
because $Q_k^T(-B)\cdot Q_k=P_k$ by definition and $Q_k^T(-B)\cdot P\cdot Q_k=P_{k+1}$ by Lemma~\ref{lem:identities-block-operators}(b); and $\boldsymbol{\alpha}_{k,y}$, $\boldsymbol{\beta}_{k,y}$ commute with $P_k$ and $P_{k+1}$. Adding $K$ on both sides and using $X_k=K+b_k\cdot P_k$, one obtains \eqref{eq:algebraic-gaussian-18}.

\emph{Equivalence of \textup{(i)} and \textup{(ii)}; the identity \eqref{eq:algebraic-gaussian-22}.} Let $\mathfrak{M}_y$ be the linear map of $\mathfrak{S}(1,l)\oplus\mathfrak{S}(\eta,l)$ given by
\[
\mathfrak{M}_y(u,v):=\left(a_yu-b_k\cdot Q_kv,\;-b_k\cdot Q_k^T(-B)u+X_kv\right)\qquad\left(u\in\mathfrak{S}(1,l),\ v\in\mathfrak{S}(\eta,l)\right).
\]
This is a map of the form \eqref{eq:algebraic-gaussian-LA2}, with $U=\mathfrak{S}(1,l)$, $V=\mathfrak{S}(\eta,l)$ and the four blocks $a_y$, $-b_k\cdot Q_k$, $-b_k\cdot Q_k^T(-B)$, $X_k$; so the lemma on Schur complements applies to it. The lower right block $X_k$ is invertible by hypothesis, with the inverse $S_k$, and by the definition of $D_k$ in Section~\ref{sec:block-operators} the corresponding operator $s$ of that lemma is
\[
s=a_y-b_k^2\cdot Q_k\cdot S_k\cdot Q_k^T(-B)=D_k+\tilde{b}\cdot P+i\gamma_3\cdot y.
\]
The upper left block $a_y$ is invertible, with the inverse $B_{k,y}$, and by \eqref{eq:algebraic-gaussian-18} the corresponding operator $t$ is
\[
t=X_k-b_k^2\cdot Q_k^T(-B)\cdot B_{k,y}\cdot Q_k=X_{k,y}.
\]
By that lemma, $\mathfrak{M}_y$ is invertible if and only if $s$ is invertible, and also if and only if $t$ is invertible; this is the equivalence of (i) and (ii). The two expressions of that lemma for $\det\mathfrak{M}_y$, through $s$ and through $t$, give
\[
\det X_k\cdot\det\left(D_k+\tilde{b}\cdot P+i\gamma_3\cdot y\right)=\det\mathfrak{M}_y=\det a_y\cdot\det X_{k,y},
\]
which is \eqref{eq:algebraic-gaussian-22}.

\emph{The identities \eqref{eq:algebraic-gaussian-19}, \eqref{eq:algebraic-gaussian-20}, \eqref{eq:algebraic-gaussian-21}.} Let (i) and (ii) hold, and write
\[
\mathfrak{M}_y^{-1}(u,v)=\left(R_{11}u+R_{12}v,\;R_{21}u+R_{22}v\right)
\]
with linear maps $R_{11}$ of $\mathfrak{S}(1,l)$, $R_{12}\colon\mathfrak{S}(\eta,l)\to\mathfrak{S}(1,l)$, $R_{21}\colon\mathfrak{S}(1,l)\to\mathfrak{S}(\eta,l)$ and $R_{22}$ of $\mathfrak{S}(\eta,l)$. Here $s^{-1}=\Gamma_{k,y}$ and $t^{-1}=S_{k,y}$. The first formula for $\mathfrak{M}_y^{-1}$ in the lemma on Schur complements expresses the four blocks through $S_k$ and $\Gamma_{k,y}$, the second through $B_{k,y}$ and $S_{k,y}$; as both formulas describe the same map, one obtains
\begin{align*}
R_{11}&=\Gamma_{k,y}=B_{k,y}+b_k^2\cdot B_{k,y}\cdot Q_k\cdot S_{k,y}\cdot Q_k^T(-B)\cdot B_{k,y},\\
R_{12}&=b_k\cdot\Gamma_{k,y}\cdot Q_k\cdot S_k=b_k\cdot B_{k,y}\cdot Q_k\cdot S_{k,y},\\
R_{21}&=b_k\cdot S_k\cdot Q_k^T(-B)\cdot\Gamma_{k,y}=b_k\cdot S_{k,y}\cdot Q_k^T(-B)\cdot B_{k,y},\\
R_{22}&=S_k+b_k^2\cdot S_k\cdot Q_k^T(-B)\cdot\Gamma_{k,y}\cdot Q_k\cdot S_k=S_{k,y}.
\end{align*}
The first line is \eqref{eq:algebraic-gaussian-19}. By the definitions $H_k=b_k\cdot S_k\cdot Q_k^T(-B)$ and $\mathfrak{H}_k=b_k\cdot Q_k\cdot S_k$ of Section~\ref{sec:block-operators}, the middle members of the other three lines are $\Gamma_{k,y}\cdot\mathfrak{H}_k$, $H_k\cdot\Gamma_{k,y}$ and $S_k+H_k\cdot\Gamma_{k,y}\cdot\mathfrak{H}_k$. Hence the second and the third line are the two identities \eqref{eq:algebraic-gaussian-21}, and the fourth line is \eqref{eq:algebraic-gaussian-20}.

Finally, for $y=0$ the definitions give $\Gamma_{k,0}=\left(D_k+\tilde{b}\cdot P\right)^{-1}=\Gamma_k$ (Section~\ref{sec:block-operators}) and $S_{k,0}=X_{k,0}^{-1}=X^0_{k+1}(B)^{-1}$.
\end{proof}

\subsubsection{From one level to the next}\label{sec:one-level-next}

\begin{theorem}[Fermion operators of two consecutive levels]\label{thm:fermion-operators-two}
Let the data be those of Section~\ref{sec:block-operators}, with $k\ge 0$ and $l\in L\mathbb{N}$. Let $B'$ be a complex bond function on $T_{k+1}$, and let $B:=B'_L$ be its scaling~\eqref{eq:algebraic-gaussian-3}, a complex bond function on $T_k$. The operators of level $k+1$ are those of Section~\ref{sec:block-operators} for the data of level $k+1$, taken at $B'$; spinor functions, fields and functions on pairs of sites are scaled as in Section~\ref{sec:block-operators}. In (a)--(c) the function $B'$ is fixed, and it is supposed that $\Gamma_k(B)$ exists and, if $k\ge 1$, that $S_k(B)$ exists. Part~(d) is asserted without this supposition: in~(d), $B$ ranges over the complex bond functions on $T_k$ named there, and $B'$ denotes the complex bond function on $T_{k+1}$ with $B'_L=B$, that is, $B'(b')=L^{\frac{1}{2}}B(Lb')$ for every bond $b'$ of $T_{k+1}$. Such a $B'$ exists and is unique: as $b'\mapsto Lb'$ is a bijection of the bonds of $T_{k+1}$ onto those of $T_k$, the scaling $B'\mapsto B'_L$ is a bijection of the complex bond functions on $T_{k+1}$ onto those on $T_k$.
\begin{enumerate}
\item[(a)] \emph{Intertwining relations.} The inverses
\[
S^0_{k+1}(B):=\left(K(B)+b_{k+1}L^{-1}\cdot P_{k+1}(B)\right)^{-1}\quad\text{on }\mathfrak{S}(\eta_k,l),\qquad S_{k+1}(B')=X_{k+1}(B')^{-1}
\]
exist; consequently the operators $H_{k+1}(B')$ and $D_{k+1}(B')$ are defined. Let $D^0_{k+1}(B)$ be the operator on $\mathfrak{S}(L,l)$ given by
\[
D^0_{k+1}(B):=\tilde{b}-\tilde{b}^2\cdot Q(B)\Gamma_k(B)Q^T(-B).
\]
Then, for all $f\in\mathfrak{S}(\eta_{k+1},l/L)$ and $v\in\mathfrak{S}(1,l/L)$,
\begin{align}
S^0_{k+1}(B)f_L&=L\left(S_{k+1}(B')f\right)_L,\label{eq:algebraic-gaussian-23}\\
D^0_{k+1}(B)&=b_{k+1}L^{-1}-\left(b_{k+1}L^{-1}\right)^2\cdot Q_{k+1}(B)S^0_{k+1}(B)Q_{k+1}^T(-B),\notag\\
D^0_{k+1}(B)v_L&=L^{-1}\left(D_{k+1}(B')v\right)_L,\label{eq:algebraic-gaussian-24}\\
\tilde{b}\cdot H_k(B)\Gamma_k(B)Q^T(-B)&=b_{k+1}L^{-1}\cdot S^0_{k+1}(B)Q_{k+1}^T(-B),\notag\\
\tilde{b}\cdot H_k(B)\Gamma_k(B)Q^T(-B)v_L&=\left(H_{k+1}(B')v\right)_L.\label{eq:algebraic-gaussian-25}
\end{align}
\item[(b)] \emph{Completion of the square.} Let $\Psi'$, $\bar{\Psi}'$ be fields on $T^0_{k+1}$, let $W'$, $\bar{W}'$ be fields on $T^0_k$, and let
\begin{equation}\label{eq:algebraic-gaussian-26}
\Psi^c:=\tilde{b}\cdot\Gamma_k(B)Q^T(-B)\Psi'_L,\qquad\bar{\Psi}^c:=\tilde{b}\cdot\Gamma_k(B)^TQ^T(B)\bar{\Psi}'_L.
\end{equation}
If $\Psi=\Psi^c+W'$ and $\bar{\Psi}=\bar{\Psi}^c+\bar{W}'$, then
\begin{multline}\label{eq:algebraic-gaussian-27}
\left\langle\bar{\Psi},D_k(B)\Psi\right\rangle_1+\tilde{b}\cdot\left\langle\bar{\Psi}'_L-Q(-B)\bar{\Psi},\,\Psi'_L-Q(B)\Psi\right\rangle_L\\
=\left\langle\bar{\Psi}',D_{k+1}(B')\Psi'\right\rangle_1+\left\langle\bar{W}',\left(D_k(B)+\tilde{b}\cdot P(B)\right)W'\right\rangle_1.
\end{multline}
Moreover, the fields $\psi^0:=H_k(B)\Psi^c$ and $\bar{\psi}^0:=\bar{H}_k(B)\bar{\Psi}^c$ on $T_k$ satisfy
\begin{align*}
\psi^0&=\left(H_{k+1}(B')\Psi'\right)_L, & \Psi^c&=T_k(B)\psi^0,\\
\bar{\psi}^0&=\left(\bar{H}_{k+1}(B')\bar{\Psi}'\right)_L, & \bar{\Psi}^c&=\bar{T}_k(B)\bar{\psi}^0,
\end{align*}
and, on the pairs of $\mathrm{Pz}_k$,
\[
\delta_{\alpha,B}\psi^0=\left(\delta_{\alpha,B'}H_{k+1}(B')\Psi'\right)_L,\qquad\delta_{\alpha,-B}\bar{\psi}^0=\left(\delta_{\alpha,-B'}\bar{H}_{k+1}(B')\bar{\Psi}'\right)_L.
\]
If $k\ge 1$, then, as in \cite[(44), (63)]{DimockMB},
\[
\Psi^c=Q_k(B)\psi^0+\frac{\tilde{b}}{b_k+\tilde{b}}\cdot Q^T(-B)\left(\Psi'_L-Q_{k+1}(B)\psi^0\right).
\]
\item[(c)] \emph{Determinants.} With $n_{k+1}=(l/L)^3$,
\[
\det D_{k+1}(B')=c'_k\cdot\frac{\det D_k(B)}{\delta Z_k(B)},\qquad c'_k=b^{4n_{k+1}}.
\]
\item[(d)] \emph{Existence and analyticity.} Suppose that the hypotheses of Theorem~\ref{thm:expansion-walks} are satisfied, and let $B\in\tilde{G}_k$. Then $X_k(B)=D_B+\mathfrak{m}+b_kP_k(B)$ is invertible if $k\ge 1$, and $D_k(B)+\tilde{b}\cdot P(B)$ is invertible for every $k\ge 0$; hence $\Gamma_k(B)$ exists and, when $k\ge 1$, so does $S_k(B)$. Indeed, for $k\ge 1$ the operator $X_k(B)$ is $F_{\infty}(B)$, whose invertibility is asserted in Theorem~\ref{thm:expansion-walks}(a), while Proposition~\ref{prop:operators-fermion-1}(d) contains the invertibility of $D_k(B)+\tilde{b}\cdot P(B)$; for $k=0$ the operator $D_0(B)+\tilde{b}\cdot P(B)=D_B+\mathfrak{m}+\tilde{b}\cdot P(B)$ is $F_y(B)$ at $y=0$ on level $0$, whose invertibility Theorem~\ref{thm:expansion-walks}(a) asserts as well. For $k\ge 1$ both statements of invertibility are contained in Propositions~\ref{prop:decay-fine-propagator} and~\ref{prop:kernel-bounds-2}; for $k=0$ the invertibility of $D_0(B)+\tilde{b}\cdot P(B)$, that is, the existence of $\Gamma_0(B)$, which Section~\ref{sec:exact-fermion-operators} denotes also by $S_{0,0}(B)$, is stated in the last paragraph of that section. If $\mathfrak{m}>0$, then $\Gamma_k(B)$ and, when $k\ge 1$, $S_k(B)$ exist for every real-valued $B$ (Lemma~\ref{lem:algebraic-gaussian-1}). Finally, the set of those $B$ for which $\Gamma_k(B)$ exists and, when $k\ge 1$, $S_k(B)$ exists is an open subset of the space of complex bond functions on $T_k$; every operator occurring in (a)--(c) depends analytically on $B$ on this subset, and (a)--(c) hold on it as equalities of analytic functions of $B$. Under the hypotheses of Theorem~\ref{thm:expansion-walks} this subset contains $\tilde{G}_k$.
\end{enumerate}
\end{theorem}

\begin{proof}
Throughout the proof an operator written without argument is taken at $B$.

\emph{Step 1: the inverses in (a), and \eqref{eq:algebraic-gaussian-23}.} For every $k\ge 0$ let $X^0_{k+1}(B):=K(B)+b_{k+1}L^{-1}\cdot P_{k+1}(B)$, an operator on $\mathfrak{S}(\eta_k,l)$. Let $k\ge 1$. Then $X^0_{k+1}(B)$ is the operator $X_{k,0}(B)$ of Section~\ref{sec:schur-complement-identities}, and $S^0_{k+1}(B)$ as defined in (a) is the operator $S^0_{k+1}(B)=X^0_{k+1}(B)^{-1}$ of Proposition~\ref{prop:cf-geometry-large}, once it exists. Since $X_k(B)$ and $D_k+\tilde{b}\cdot P$ are invertible by assumption, the equivalence in Proposition~\ref{prop:cf-geometry-large}, taken at $y=0$, shows that $X^0_{k+1}(B)$ is invertible. Let $k=0$. Then $b_1L^{-1}=bL^{-1}=\tilde{b}$, and, as $Q_0=1$, Lemma~\ref{lem:identities-block-operators}(b) gives $Q_1(B)=Q(B)$, $Q_1^T(-B)=Q^T(-B)$ and hence $P_1(B)=P(B)$. So $X^0_1(B)=K+\tilde{b}\cdot P=\Gamma_0(B)^{-1}$ by the definition of level $0$ in Section~\ref{sec:block-operators}, and $S^0_1(B)=\Gamma_0(B)$. In both cases $S^0_{k+1}(B)$ exists.

Let $f\in\mathfrak{S}(\eta_{k+1},l/L)$. By Lemma~\ref{lem:identities-block-operators}(f) one has $K(B)f_L=L^{-1}\left(K_{k+1}(B')f\right)_L$ with $K_{k+1}(B')=D_{B'}+L\mathfrak{m}$, further $Q_{k+1}(B)f_L=\left(Q_{k+1}(B')f\right)_L$, and $Q_{k+1}^T(-B)v_L=\left(Q_{k+1}^T(-B')v\right)_L$ for $v\in\mathfrak{S}(1,l/L)$. Since $P_{k+1}(B)=Q_{k+1}^T(-B)Q_{k+1}(B)$ (Lemma~\ref{lem:identities-block-operators}(b)) and $P_{k+1}(B')=Q_{k+1}^T(-B')Q_{k+1}(B')$, the last two relations, the second one with $v=Q_{k+1}(B')f$, yield $P_{k+1}(B)f_L=\left(P_{k+1}(B')f\right)_L$. Therefore
\[
X^0_{k+1}(B)f_L=L^{-1}\left(\left(K_{k+1}(B')+b_{k+1}P_{k+1}(B')\right)f\right)_L=L^{-1}\left(X_{k+1}(B')f\right)_L.
\]
The map $f\mapsto f_L$ is a bijection of $\mathfrak{S}(\eta_{k+1},l/L)$ onto $\mathfrak{S}(\eta_k,l)$, and by the last identity it carries $L^{-1}X_{k+1}(B')$ into $X^0_{k+1}(B)$. As $X^0_{k+1}(B)$ is invertible, so is $X_{k+1}(B')$; thus $S_{k+1}(B')$ exists, and with it $H_{k+1}(B')$ and $D_{k+1}(B')$. The last identity, applied to $S_{k+1}(B')f$ in place of $f$, reads $X^0_{k+1}(B)\left(S_{k+1}(B')f\right)_L=L^{-1}f_L$; applying $S^0_{k+1}(B)$ to both sides gives \eqref{eq:algebraic-gaussian-23}.

\emph{Step 2: \eqref{eq:algebraic-gaussian-24} and \eqref{eq:algebraic-gaussian-25}.} Consider first the identities between operators in \eqref{eq:algebraic-gaussian-24} and \eqref{eq:algebraic-gaussian-25}, that is, the first line of each. Let $k\ge 1$, and let $a_0=b_k+\tilde{b}\cdot P$ and $B_k=B_{k,0}(B)=b_k^{-1}(1-P)+\left(b_k+\tilde{b}\right)^{-1}P$ be the operators of Section~\ref{sec:schur-complement-identities} at $y=0$, so that $B_k=a_0^{-1}$. By Lemma~\ref{lem:identities-block-operators}(a), $QP=Q$ and $PQ^T(-B)=Q^T(-B)$; hence
\[
QB_k=\frac{1}{b_k+\tilde{b}}\cdot Q,\qquad B_kQ^T(-B)=\frac{1}{b_k+\tilde{b}}\cdot Q^T(-B).
\]
By Proposition~\ref{prop:cf-geometry-large}, $\Gamma_{k,0}=\Gamma_k$ and $S_{k,0}=S^0_{k+1}(B)$, so that \eqref{eq:algebraic-gaussian-19} at $y=0$ reads $\Gamma_k=B_k+b_k^2\cdot B_kQ_kS^0_{k+1}Q_k^T(-B)B_k$. Multiply this identity by $Q$ on the left and by $Q^T(-B)$ on the right, and use $QQ^T(-B)=1$ (Lemma~\ref{lem:identities-block-operators}(a)) together with $QQ_k=Q_{k+1}$ and $Q_k^T(-B)Q^T(-B)=Q_{k+1}^T(-B)$ (Lemma~\ref{lem:identities-block-operators}(b)). One obtains
\[
Q\Gamma_kQ^T(-B)=\frac{1}{b_k+\tilde{b}}+\frac{b_k^2}{\left(b_k+\tilde{b}\right)^2}\cdot Q_{k+1}S^0_{k+1}Q_{k+1}^T(-B).
\]
The recursion $b_{k+1}=b\cdot b_k/\left(b_k+\tilde{b}\right)$ and $\tilde{b}=bL^{-1}$ (Section~\ref{sec:block-operators}) give
\[
\tilde{b}-\frac{\tilde{b}^2}{b_k+\tilde{b}}=\frac{\tilde{b}\cdot b_k}{b_k+\tilde{b}}=b_{k+1}L^{-1},\qquad\frac{\tilde{b}^2\cdot b_k^2}{\left(b_k+\tilde{b}\right)^2}=\left(b_{k+1}L^{-1}\right)^2;
\]
so multiplication of the preceding identity by $-\tilde{b}^2$ and addition of $\tilde{b}$ yields the operator identity in \eqref{eq:algebraic-gaussian-24}. Next, \eqref{eq:algebraic-gaussian-21} at $y=0$ states that $H_k\Gamma_k=b_k\cdot S^0_{k+1}Q_k^T(-B)B_k$. Consequently
\begin{align*}
\tilde{b}\cdot H_k\Gamma_kQ^T(-B)&=\tilde{b}\cdot b_k\cdot S^0_{k+1}Q_k^T(-B)B_kQ^T(-B)\\
&=\frac{\tilde{b}\cdot b_k}{b_k+\tilde{b}}\cdot S^0_{k+1}Q_{k+1}^T(-B)=b_{k+1}L^{-1}\cdot S^0_{k+1}Q_{k+1}^T(-B),
\end{align*}
which is the operator identity in \eqref{eq:algebraic-gaussian-25}. For $k=0$ both operator identities are the definitions, because $H_0=1$ (Section~\ref{sec:block-operators}), $Q_1=Q$, $S^0_1=\Gamma_0$ and $b_1L^{-1}=\tilde{b}$ (Step~1).

Now let $v\in\mathfrak{S}(1,l/L)$. Lemma~\ref{lem:identities-block-operators}(f) and \eqref{eq:algebraic-gaussian-23}, the latter applied to $f=Q_{k+1}^T(-B')v$, give
\begin{align*}
Q_{k+1}(B)S^0_{k+1}(B)Q_{k+1}^T(-B)v_L&=Q_{k+1}(B)S^0_{k+1}(B)\left(Q_{k+1}^T(-B')v\right)_L\\
&=L\left(Q_{k+1}(B')S_{k+1}(B')Q_{k+1}^T(-B')v\right)_L.
\end{align*}
Inserting this into the operator identity in \eqref{eq:algebraic-gaussian-24}, one finds
\begin{align*}
D^0_{k+1}(B)v_L&=L^{-1}\left(\left(b_{k+1}-b_{k+1}^2\cdot Q_{k+1}(B')S_{k+1}(B')Q_{k+1}^T(-B')\right)v\right)_L\\
&=L^{-1}\left(D_{k+1}(B')v\right)_L.
\end{align*}
In the same way, $b_{k+1}L^{-1}\cdot S^0_{k+1}(B)Q_{k+1}^T(-B)v_L=\left(b_{k+1}S_{k+1}(B')Q_{k+1}^T(-B')v\right)_L=\left(H_{k+1}(B')v\right)_L$; together with the operator identity in \eqref{eq:algebraic-gaussian-25} this is the second identity of \eqref{eq:algebraic-gaussian-25}.

\emph{Step 3: \eqref{eq:algebraic-gaussian-27}.} Let
\[
J:=\tilde{b}\cdot Q^T(-B)\Psi'_L,\qquad\bar{J}:=\tilde{b}\cdot Q^T(B)\bar{\Psi}'_L;
\]
these are fields on $T^0_k$, and $\Psi^c=\Gamma_kJ$, $\bar{\Psi}^c=\Gamma_k^T\bar{J}$ by \eqref{eq:algebraic-gaussian-26}. The relation between an operator and its transpose holds for fields, the factors keeping their order (Section~\ref{sec:spinor-spaces-berezin}), and $Q^T(\pm B)$ is the transpose of $Q(\pm B)$ (Section~\ref{sec:block-operators}). With $Q^T(-B)Q(B)=P(B)$ it follows that
\begin{align*}
\left\langle\bar{\Psi}'_L,Q(B)\Psi\right\rangle_L&=\left\langle Q^T(B)\bar{\Psi}'_L,\Psi\right\rangle_1,\\
\left\langle Q(-B)\bar{\Psi},\Psi'_L\right\rangle_L&=\left\langle\bar{\Psi},Q^T(-B)\Psi'_L\right\rangle_1,\\
\left\langle Q(-B)\bar{\Psi},Q(B)\Psi\right\rangle_L&=\left\langle\bar{\Psi},P\Psi\right\rangle_1.
\end{align*}
Expanding the second term on the left of \eqref{eq:algebraic-gaussian-27} by bilinearity, using these three relations and $D_k+\tilde{b}\cdot P=\Gamma_k^{-1}$, one finds that the left side of \eqref{eq:algebraic-gaussian-27} equals
\[
\left\langle\bar{\Psi},\Gamma_k^{-1}\Psi\right\rangle_1-\left\langle\bar{\Psi},J\right\rangle_1-\left\langle\bar{J},\Psi\right\rangle_1+\tilde{b}\cdot\left\langle\bar{\Psi}'_L,\Psi'_L\right\rangle_L.
\]
Insert $\Psi=\Gamma_kJ+W'$ and $\bar{\Psi}=\Gamma_k^T\bar{J}+\bar{W}'$. Since $\left\langle\Gamma_k^Tg,h\right\rangle_1=\left\langle g,\Gamma_kh\right\rangle_1$ for all fields $g$, $h$ on $T^0_k$,
\begin{align*}
\left\langle\bar{\Psi},\Gamma_k^{-1}\Psi\right\rangle_1&=\left\langle\bar{J},\Gamma_kJ\right\rangle_1+\left\langle\bar{J},W'\right\rangle_1+\left\langle\bar{W}',J\right\rangle_1+\left\langle\bar{W}',\Gamma_k^{-1}W'\right\rangle_1,\\
\left\langle\bar{\Psi},J\right\rangle_1&=\left\langle\bar{J},\Gamma_kJ\right\rangle_1+\left\langle\bar{W}',J\right\rangle_1,\\
\left\langle\bar{J},\Psi\right\rangle_1&=\left\langle\bar{J},\Gamma_kJ\right\rangle_1+\left\langle\bar{J},W'\right\rangle_1.
\end{align*}
Hence the left side of \eqref{eq:algebraic-gaussian-27} equals $\left\langle\bar{W}',\Gamma_k^{-1}W'\right\rangle_1-\left\langle\bar{J},\Gamma_kJ\right\rangle_1+\tilde{b}\cdot\left\langle\bar{\Psi}'_L,\Psi'_L\right\rangle_L$. The transpose relation for $Q(B)$ gives $\left\langle\bar{J},\Gamma_kJ\right\rangle_1=\tilde{b}^2\cdot\left\langle\bar{\Psi}'_L,Q\Gamma_kQ^T(-B)\Psi'_L\right\rangle_L$, so that, by the definition of $D^0_{k+1}(B)$ in (a), the left side of \eqref{eq:algebraic-gaussian-27} is
\[
\left\langle\bar{W}',\left(D_k+\tilde{b}\cdot P\right)W'\right\rangle_1+\left\langle\bar{\Psi}'_L,D^0_{k+1}(B)\Psi'_L\right\rangle_L.
\]
Finally, the second identity of \eqref{eq:algebraic-gaussian-24}, which extends to fields by linearity, and the relation $\left\langle g_L,f_L\right\rangle_L=L\left\langle g,f\right\rangle_1$ of Lemma~\ref{lem:identities-block-operators}(f) give
\[
\left\langle\bar{\Psi}'_L,D^0_{k+1}(B)\Psi'_L\right\rangle_L=L^{-1}\left\langle\bar{\Psi}'_L,\left(D_{k+1}(B')\Psi'\right)_L\right\rangle_L=\left\langle\bar{\Psi}',D_{k+1}(B')\Psi'\right\rangle_1,
\]
and \eqref{eq:algebraic-gaussian-27} follows.

\emph{Step 4: the fields $\psi^0$ and $\bar{\psi}^0$.} Since $H_k\Psi^c=\tilde{b}\cdot H_k\Gamma_kQ^T(-B)\Psi'_L$, the identity $\psi^0=\left(H_{k+1}(B')\Psi'\right)_L$ is the second identity of \eqref{eq:algebraic-gaussian-25}, extended to fields. By Lemma~\ref{lem:identities-block-operators}(a), $T_kH_k=1$ (for $k=0$ both operators are $1$, Section~\ref{sec:block-operators}); hence $T_k\psi^0=T_kH_k\Psi^c=\Psi^c$.

The assumption of the theorem is satisfied at $-B=(-B')_L$ as well. For $k\ge 1$, Lemma~\ref{lem:identities-block-operators}(c) gives $X_k(-B)=\mathsf{C} X_k(B)^T\mathsf{C}^{-1}$, so that $S_k(-B)$ and $D_k(-B)$ exist. For every $k\ge 0$ the same lemma gives (for $k=0$, where $D_0=K$, through its relations for $D_B$ and $P(B)$)
\[
\left(D_k+\tilde{b}\cdot P\right)(-B)=\mathsf{C}\left(\left(D_k+\tilde{b}\cdot P\right)(B)\right)^T\mathsf{C}^{-1};
\]
this operator is invertible, that is, $\Gamma_k(-B)$ exists, and $\Gamma_k(-B)^{-1}=\mathsf{C}\left(\Gamma_k(B)^{-1}\right)^T\mathsf{C}^{-1}$. So (a) holds at $-B'$ in place of $B'$, and $\Gamma_k(B)^T=\mathsf{C}^{-1}\Gamma_k(-B)\mathsf{C}$. As $\mathsf{C}$ acts on the spinor indices only, it commutes with $Q^T(B)$ and with the scaling; therefore
\[
\mathsf{C}\bar{\Psi}^c=\tilde{b}\cdot\Gamma_k(-B)Q^T\left(-(-B)\right)\left(\mathsf{C}\bar{\Psi}'\right)_L,
\]
which is the field $\Psi^c$ of \eqref{eq:algebraic-gaussian-26} formed at $-B$ with $\mathsf{C}\bar{\Psi}'$ in place of $\Psi'$. The second identity of \eqref{eq:algebraic-gaussian-25} at $-B$ thus gives $H_k(-B)\mathsf{C}\bar{\Psi}^c=\left(H_{k+1}(-B')\mathsf{C}\bar{\Psi}'\right)_L$. Multiplication by $\mathsf{C}^{-1}$, together with $\bar{H}_k(B)=\mathsf{C}^{-1}H_k(-B)\mathsf{C}$ and $\bar{H}_{k+1}(B')=\mathsf{C}^{-1}H_{k+1}(-B')\mathsf{C}$ (Section~\ref{sec:block-operators}), yields $\bar{\psi}^0=\bar{H}_k(B)\bar{\Psi}^c=\left(\bar{H}_{k+1}(B')\bar{\Psi}'\right)_L$. By Lemma~\ref{lem:identities-block-operators}(c), $\bar{T}_k\bar{H}_k=1$ (for $k=0$ both operators are $1$); hence $\bar{T}_k\bar{\psi}^0=\bar{\Psi}^c$.

For the H\"older derivatives, apply the relation $\delta_{\alpha,B}f_L=\left(\delta_{\alpha,B'}f\right)_L$ on the pairs of $\mathrm{Pz}_k$ (Lemma~\ref{lem:identities-block-operators}(f)), which extends to fields by linearity, to $f=H_{k+1}(B')\Psi'$, and the same relation at $(-B,-B')$ to $f=\bar{H}_{k+1}(B')\bar{\Psi}'$.

Let now $k\ge 1$, and let $\mathfrak{M}_0$ be the operator of the proof of Proposition~\ref{prop:cf-geometry-large} at $y=0$, that is, the linear map of $\mathfrak{S}(1,l)\oplus\mathfrak{S}(\eta_k,l)$ (and, by linearity, of the corresponding pairs of fields) given by
\[
\mathfrak{M}_0(g,h):=\left(a_0g-b_kQ_kh,\;-b_kQ_k^T(-B)g+X_kh\right).
\]
Its block $X_k$ is invertible, and the corresponding operator $s$ of \eqref{eq:algebraic-gaussian-LA2} is $a_0-b_k^2\cdot Q_kS_kQ_k^T(-B)=D_k+\tilde{b}\cdot P=\Gamma_k^{-1}$. By \eqref{eq:algebraic-gaussian-LA2} the map $\mathfrak{M}_0$ is invertible, and the formula for $\mathfrak{M}_0^{-1}$ there gives
\[
\mathfrak{M}_0^{-1}(J,0)=\left(\Gamma_kJ,\;b_kS_kQ_k^T(-B)\Gamma_kJ\right)=\left(\Psi^c,\psi^0\right),
\]
because $\Gamma_kJ=\Psi^c$ and $b_kS_kQ_k^T(-B)=H_k$. The first component of the identity $\mathfrak{M}_0(\Psi^c,\psi^0)=(J,0)$ reads $a_0\Psi^c-b_kQ_k\psi^0=J$. With $a_0^{-1}=B_k=b_k^{-1}(1-P)+\left(b_k+\tilde{b}\right)^{-1}P$ and the relation $B_kQ^T(-B)=\left(b_k+\tilde{b}\right)^{-1}Q^T(-B)$ of Step~2 it follows that
\[
\Psi^c=B_k\left(J+b_kQ_k\psi^0\right)=\frac{\tilde{b}}{b_k+\tilde{b}}\cdot Q^T(-B)\Psi'_L+\left(1-\frac{\tilde{b}}{b_k+\tilde{b}}\cdot P\right)Q_k\psi^0.
\]
Since $PQ_k=Q^T(-B)QQ_k=Q^T(-B)Q_{k+1}$ (Lemma~\ref{lem:identities-block-operators}(b)), this is the formula for $\Psi^c$ asserted in (b).

\emph{Step 5: (c).} Let $\mathscr{U}:\mathfrak{S}(1,l/L)\to\mathfrak{S}(L,l)$ be the bijection $v\mapsto v_L$. The second identity of \eqref{eq:algebraic-gaussian-24} states that $D^0_{k+1}(B)=L^{-1}\mathscr{U} D_{k+1}(B')\mathscr{U}^{-1}$. Both spaces have dimension $4n_{k+1}$; hence
\[
\det D^0_{k+1}(B)=L^{-4n_{k+1}}\det D_{k+1}(B').
\]
On the other hand $D^0_{k+1}(B)=\tilde{b}\left(1-\tilde{b}\cdot Q\Gamma_kQ^T(-B)\right)$. Apply \eqref{eq:algebraic-gaussian-LA3} with $U=\mathfrak{S}(L,l)$, $V=\mathfrak{S}(1,l)$, $X=Q:V\to U$ and $Y=\tilde{b}\cdot\Gamma_kQ^T(-B):U\to V$, for which $YX=\tilde{b}\cdot\Gamma_kP$. It gives
\begin{align*}
\det D^0_{k+1}(B)&=\tilde{b}^{4n_{k+1}}\det\left(1-\tilde{b}\cdot Q\Gamma_kQ^T(-B)\right)=\tilde{b}^{4n_{k+1}}\det\left(1-\tilde{b}\cdot\Gamma_kP\right)\\
&=\tilde{b}^{4n_{k+1}}\det\left(\Gamma_k\left(\Gamma_k^{-1}-\tilde{b}\cdot P\right)\right)=\tilde{b}^{4n_{k+1}}\cdot\frac{\det D_k(B)}{\delta Z_k(B)},
\end{align*}
since $\Gamma_k^{-1}-\tilde{b}\cdot P=D_k$ and $\det\Gamma_k=\delta Z_k(B)^{-1}$. As $L\cdot\tilde{b}=b$, comparison of the two expressions for $\det D^0_{k+1}(B)$ gives (c).

\emph{Step 6: (d).} For $B\in\tilde{G}_k$ the invertibility of $X_k(B)$ (if $k\ge 1$) and of $D_k(B)+\tilde{b}\cdot P(B)$ is the content of the clauses of Theorem~\ref{thm:expansion-walks}(a) and Proposition~\ref{prop:operators-fermion-1}(d) recalled in (d). For real-valued $B$ and $\mathfrak{m}>0$, Lemma~\ref{lem:algebraic-gaussian-1}(a) shows that $X_k(B)$ is invertible if $k\ge 1$, and Lemma~\ref{lem:algebraic-gaussian-1}(c) that $\Gamma_k(B)$ exists.

Every entry of each of the matrices $K(B)$, $Q_k(\pm B)$, $Q_k^T(\pm B)$, $Q(\pm B)$, $Q^T(\pm B)$, and of the operators $\delta_{\alpha,\pm B}$ on the pairs of $\mathrm{Pz}_k$, depends on $B$ only through its values $B(b)$ on the finitely many bonds $b$ of $T_k$, and is a finite sum of constant multiples of exponentials of linear combinations of these values; in particular it is an entire function of them. Since $B'(b')=L^{\frac{1}{2}}B(Lb')$ for every bond $b'$ of $T_{k+1}$, the operators of level $k+1$, at $B'$ and at $-B'$, have entries of the same kind.

Consider the set of those $B$ for which $\Gamma_k(B)$ exists and, if $k\ge 1$, $S_k(B)$ exists. For $k=0$ it is the set where $\det\left(K(B)+\tilde{b}\cdot P(B)\right)\ne 0$, which is open, the determinant being an entire function of $B$. Let $k\ge 1$. As $B\mapsto\det X_k(B)$ is entire, the set of those $B$ for which $S_k(B)$ exists is open, and on it every entry of $S_k(B)$ is a cofactor of $X_k(B)$ divided by $\det X_k(B)$ (Cramer's rule), hence analytic. Consequently, on this open set the operator
\[
D_k(B)+\tilde{b}\cdot P(B)=b_k-b_k^2\cdot Q_k(B)S_k(B)Q_k^T(-B)+\tilde{b}\cdot P(B)
\]
has analytic entries, and its determinant $\delta Z_k(B)$ is analytic as well. Since $\delta Z_k$ is continuous there, the points of this open set at which $\delta Z_k(B)\ne 0$ form an open set; they are exactly the $B$ for which both $S_k(B)$ and $\Gamma_k(B)$ exist. So the set under consideration is open for every $k\ge 0$, and $\Gamma_k(B)$, the inverse of a matrix whose entries are analytic on it, is analytic on it by Cramer's rule.

At each point of this set, $S^0_{k+1}(B)$ and $S_{k+1}(B')$ exist by (a), and $S_{k+1}(-B')$ and, if $k\ge 1$, $S_k(-B)$ exist by Step~4. These are the inverses of $X^0_{k+1}(B)$, $X_{k+1}(B')$, $X_{k+1}(-B')$ and $X_k(-B)$, whose entries are entire functions of $B$; by Cramer's rule they are analytic on the set as well. Every operator that occurs in (a)--(c) is obtained from the operators treated in this step by sums, products, transposes, conjugation with $\mathsf{C}$ and scaling; therefore it depends analytically on $B$ on the set, and so do the determinants in (c). Since Steps~1--5 prove (a)--(c) at each point of the set, (a)--(c) are equalities between functions of $B$ that are analytic on this open set.

Under the hypotheses of Theorem~\ref{thm:expansion-walks}, parts (a)--(d) hold for every $B\in\tilde{G}_k$, and Proposition~\ref{prop:comparison-fermion-operators} follows.
\end{proof}

\subsection{Real bond functions}\label{sec:real-bond-functions}

\subsubsection{Positivity and invertibility}\label{sec:positivity-invertibility}

In this subsection $(\cdot,\cdot)$ denotes, on each of the spaces $\mathfrak{S}(\eta,l)$ and $\mathfrak{S}(1,l)$, the Hermitian pairing of Section~\ref{sec:spinor-spaces-berezin}, that is, $(\cdot,\cdot)_{\eta}$ and $(\cdot,\cdot)_1$ respectively; $\lvert f\rvert:=(f,f)^{\frac{1}{2}}$, and $\lvert O\rvert$ is the corresponding operator norm of an operator $O$. As in Section~\ref{sec:spinor-spaces-berezin}, $\operatorname{Re}O:=\left(O+O^{*}\right)/2$. For a self-adjoint operator $O$ and a real number $c$, the inequality $O\ge c$ means $(f,Of)\ge c\cdot\lvert f\rvert^2$ for every $f$, and $O>0$ means $(f,Of)>0$ for every $f\ne0$. A complex $4\times4$ matrix, in particular $\gamma_3$, acts on spinor functions pointwise: $\left(\gamma_3f\right)(x):=\gamma_3\cdot f(x)$.

\begin{lemma}\label{lem:algebraic-gaussian-1}
Let the data be those of Section~\ref{sec:block-operators}, with $k\ge0$ and $\mathfrak{m}>0$, and let $B$ be real-valued. All operators of Section~\ref{sec:block-operators} are taken at $B$, and the argument $B$ is not written.
\begin{enumerate}
\item[(a)] One has $\operatorname{Re}K\ge\mathfrak{m}$, and $\operatorname{Re}X_k\ge\mathfrak{m}$ if $k\ge1$. In particular $K$ and, for $k\ge1$, $X_k$ are invertible.
\item[(b)] The operator $D_k$ is invertible, and $\operatorname{Re}D_k>0$. If $k\ge1$, then
\[
D_k^{-1}=b_k^{-1}+Q_k\cdot K^{-1}\cdot Q_k^{*}.
\]
\item[(c)] The operator $\mathfrak{K}:=D_k+\tilde{b}\cdot P$ on $\mathfrak{S}(1,l)$ satisfies
\[
\operatorname{Re}\mathfrak{K}>0,\qquad \mathfrak{K}^{*}=\gamma_3\cdot\mathfrak{K}\cdot\gamma_3.
\]
Consequently $\Gamma_k$ exists, the operator $\mathfrak{T}:=\mathfrak{K}\cdot\gamma_3$ on $\mathfrak{S}(1,l)$ is self-adjoint and invertible, and $\det\mathfrak{T}=\delta Z_k(B)$.
\item[(d)] For every $y\ge0$ the operator
\[
\mathfrak{K}+i\gamma_3\cdot y=\left(\mathfrak{T}+iy\right)\cdot\gamma_3
\]
on $\mathfrak{S}(1,l)$ is invertible, and its inverse
\[
\Gamma_{k,y}:=\left(\mathfrak{K}+i\gamma_3\cdot y\right)^{-1}=\gamma_3\cdot\left(\mathfrak{T}+iy\right)^{-1}
\quad\text{satisfies}\quad
\left\lvert\Gamma_{k,y}\right\rvert\le\min\left\{\frac{1}{y},\left\lvert\mathfrak{T}^{-1}\right\rvert\right\},
\]
with the convention $1/0:=+\infty$. If $k\ge1$, then $S_{k,y}(B)$ exists for every $y\ge0$. If $k=0$, then $\Gamma_{0,y}=S_{0,y}(B)$, where $S_{0,y}$ is the operator of Section~\ref{sec:exact-fermion-operators}.
\end{enumerate}
\end{lemma}

The self-adjoint operator $\mathfrak{T}=\left(D_k+\tilde{b}\cdot P\right)\cdot\gamma_3$ of (c), whose determinant is $\delta Z_k(B)$, and the operators $\Gamma_{k,y}$ of (d) occur in \cite[Lemma~17]{DimockUR} and in \cite[(179)--(186)]{DimockMB}.

\begin{proof}
Let $O$ be an operator on $\mathfrak{S}(\eta,l)$ or on $\mathfrak{S}(1,l)$ with $\operatorname{Re}O\ge c$ for some number $c>0$. For every $f$,
\[
\left\lvert(f,Of)\right\rvert\ge\operatorname{Re}(f,Of)=\left(f,(\operatorname{Re}O)f\right)\ge c\cdot\lvert f\rvert^2,
\]
so that $Of=0$ implies $f=0$. The same implication holds under the weaker hypothesis $\operatorname{Re}O>0$: if $Of=0$, then $\left(f,(\operatorname{Re}O)f\right)=\operatorname{Re}(f,Of)=0$, and hence $f=0$. In either case $O$ is injective, hence invertible, the space being of finite dimension. Moreover $\tilde{b}=b\cdot L^{-1}>0$, and $b_k>0$ for $k\ge1$ by the recursive definition of $b_k$ in Section~\ref{sec:block-operators}.

(a) Since $B$ is real-valued, Lemma~\ref{lem:identities-block-operators}(e) gives $\operatorname{Re}D_B\ge0$; since $\mathfrak{m}$ is real, $\operatorname{Re}K=\operatorname{Re}D_B+\mathfrak{m}\ge\mathfrak{m}$, and $K$ is invertible by the first paragraph of the proof, applied with $c=\mathfrak{m}$. Let $k\ge1$. By Lemma~\ref{lem:identities-block-operators}(e) one has $Q_k^{*}=Q_k^T(-B)$, so that $P_k=Q_k^{*}\cdot Q_k\ge0$, and therefore $\operatorname{Re}X_k=\operatorname{Re}K+b_k\cdot P_k\ge\mathfrak{m}$; by the first paragraph of the proof once more, $X_k$ is invertible.

(b) For $k=0$ one has $D_0=K$ by the definition of level $0$ in Section~\ref{sec:block-operators}, and the assertions follow from (a). Let $k\ge1$. Since $Q_k^T(-B)=Q_k^{*}$ (Lemma~\ref{lem:identities-block-operators}(e)), the definition of $D_k$ reads $D_k=b_k-b_k^2\cdot Q_k\cdot S_k\cdot Q_k^{*}$. Using $Q_k^{*}\cdot Q_k=P_k$ and $K^{-1}=K^{-1}\cdot X_k\cdot S_k$, one finds
\begin{align*}
\left(b_k^{-1}+Q_k\cdot K^{-1}\cdot Q_k^{*}\right)\cdot\left(b_k-b_k^2\cdot Q_k\cdot S_k\cdot Q_k^{*}\right)
&=1+b_k\cdot Q_k\cdot\left[-S_k+K^{-1}-K^{-1}\cdot b_k\cdot P_k\cdot S_k\right]\cdot Q_k^{*}\\
&=1+b_k\cdot Q_k\cdot K^{-1}\cdot\left[-K+X_k-b_k\cdot P_k\right]\cdot S_k\cdot Q_k^{*}=1,
\end{align*}
the last bracket being zero because $X_k=K+b_k\cdot P_k$. As $\mathfrak{S}(1,l)$ has finite dimension, a left inverse of $D_k$ is its inverse; this proves that $D_k$ is invertible and gives the formula for $D_k^{-1}$. Now let $h\in\mathfrak{S}(\eta,l)$ and $g:=Kh$. By (a),
\[
\operatorname{Re}\left(g,K^{-1}g\right)=\operatorname{Re}(Kh,h)=\operatorname{Re}(h,Kh)\ge\mathfrak{m}\cdot\lvert h\rvert^2\ge0.
\]
Since $K$ is invertible, every element of $\mathfrak{S}(\eta,l)$ is of the form $Kh$; hence $\operatorname{Re}K^{-1}\ge0$, and
\[
\operatorname{Re}D_k^{-1}=b_k^{-1}+Q_k\cdot\left(\operatorname{Re}K^{-1}\right)\cdot Q_k^{*}\ge b_k^{-1}.
\]
Let $v\ne0$ and $w:=D_kv$. Then $w\ne0$, and
\[
\operatorname{Re}\left(v,D_kv\right)=\operatorname{Re}\left(D_k^{-1}w,w\right)=\operatorname{Re}\left(w,D_k^{-1}w\right)\ge b_k^{-1}\cdot\lvert w\rvert^2>0.
\]
Thus $\operatorname{Re}D_k>0$.

(c) The matrix $\gamma_3$ is Hermitian and $\gamma_3^2=1$ (Section~\ref{sec:spinor-spaces-berezin}). By Lemma~\ref{lem:identities-block-operators}(e), $P$ is an orthogonal projection, so that $P^{*}=P\ge0$; together with (b) and $\tilde{b}>0$ this gives $\operatorname{Re}\mathfrak{K}=\operatorname{Re}D_k+\tilde{b}\cdot P>0$. The mass $\mathfrak{m}$ is real and $D_B^{*}=\gamma_3\cdot D_B\cdot\gamma_3$ (Lemma~\ref{lem:identities-block-operators}(e)); hence $K^{*}=\gamma_3\cdot K\cdot\gamma_3$, which for $k=0$ is the relation $D_0^{*}=\gamma_3\cdot D_0\cdot\gamma_3$. Let $k\ge1$. The operator $P_k$ is self-adjoint (Lemma~\ref{lem:identities-block-operators}(e)) and commutes with $\gamma_3$ (Lemma~\ref{lem:identities-block-operators}(a)), and $b_k$ is real; hence
\[
X_k^{*}=\gamma_3\cdot X_k\cdot\gamma_3,\qquad S_k^{*}=\left(X_k^{*}\right)^{-1}=\gamma_3\cdot S_k\cdot\gamma_3,\qquad D_k^{*}=b_k-b_k^2\cdot Q_k\cdot S_k^{*}\cdot Q_k^{*}=\gamma_3\cdot D_k\cdot\gamma_3.
\]
In the last equality $Q_k$ and $Q_k^{*}=Q_k^T(-B)$ commute with $\gamma_3$, since by their definition in Section~\ref{sec:block-operators} they act on the spinor components as multiples of the identity. So $D_k^{*}=\gamma_3\cdot D_k\cdot\gamma_3$ for every $k\ge0$. As $P$ is self-adjoint and commutes with $\gamma_3$ (Lemma~\ref{lem:identities-block-operators}(a)), it follows that
\[
\mathfrak{K}^{*}=D_k^{*}+\tilde{b}\cdot P^{*}=\gamma_3\cdot\left(D_k+\tilde{b}\cdot P\right)\cdot\gamma_3=\gamma_3\cdot\mathfrak{K}\cdot\gamma_3.
\]
Since $\operatorname{Re}\mathfrak{K}>0$, the first paragraph of the proof shows that $\mathfrak{K}$ is invertible; that is, $\Gamma_k=\mathfrak{K}^{-1}$ exists. Since $\gamma_3$ is invertible and Hermitian, $\mathfrak{T}=\mathfrak{K}\cdot\gamma_3$ is invertible, and
\[
\mathfrak{T}^{*}=\gamma_3\cdot\mathfrak{K}^{*}=\gamma_3^2\cdot\mathfrak{K}\cdot\gamma_3=\mathfrak{K}\cdot\gamma_3=\mathfrak{T}.
\]
Restriction to a fundamental domain of $\mathbb{Z}^3$ for $l\mathbb{Z}^3$, which has $n_k=l^3$ points, identifies $\mathfrak{S}(1,l)$ with the direct sum of $n_k$ copies of $\mathbb{C}^4$, and $\gamma_3$ acts on each copy separately. Hence the determinant of $\gamma_3$ on $\mathfrak{S}(1,l)$ is $\left(\det\gamma_3\right)^{n_k}=1$, because $\det\gamma_3=1$ (Section~\ref{sec:spinor-spaces-berezin}). It follows that $\det\mathfrak{T}=\det\mathfrak{K}=\delta Z_k(B)$, by the definition of $\delta Z_k$ in Section~\ref{sec:block-operators}.

(d) As $\gamma_3^2=1$, one has $\left(\mathfrak{T}+iy\right)\cdot\gamma_3=\mathfrak{K}\cdot\gamma_3^2+iy\cdot\gamma_3=\mathfrak{K}+i\gamma_3\cdot y$. By the spectral theorem, $\mathfrak{S}(1,l)$ has an orthonormal basis of eigenvectors of the self-adjoint operator $\mathfrak{T}$; the eigenvalues $t$ are real, and $t\ne0$ because $\mathfrak{T}$ is invertible. In this basis $\mathfrak{T}+iy$ is diagonal with the entries $t+iy$, none of which vanishes; so $\mathfrak{T}+iy$ is invertible, and $\left(\mathfrak{T}+iy\right)^{-1}$ is diagonal in the same basis. Since $\lvert t+iy\rvert\ge\lvert t\rvert$ and $\lvert t+iy\rvert\ge y$,
\[
\left\lvert\left(\mathfrak{T}+iy\right)^{-1}\right\rvert=\max_t\frac{1}{\lvert t+iy\rvert}\le\min\left\{\frac{1}{y},\max_t\frac{1}{\lvert t\rvert}\right\}=\min\left\{\frac{1}{y},\left\lvert\mathfrak{T}^{-1}\right\rvert\right\},
\]
the maxima being taken over the eigenvalues $t$ of $\mathfrak{T}$. Consequently $\mathfrak{K}+i\gamma_3\cdot y=\left(\mathfrak{T}+iy\right)\cdot\gamma_3$ is invertible, with the inverse $\gamma_3^{-1}\cdot\left(\mathfrak{T}+iy\right)^{-1}=\gamma_3\cdot\left(\mathfrak{T}+iy\right)^{-1}$; as $\gamma_3$ is unitary, $\left\lvert\Gamma_{k,y}\right\rvert=\left\lvert\left(\mathfrak{T}+iy\right)^{-1}\right\rvert$, and the bound follows. Let $k\ge1$. By (a) the operator $X_k$ is invertible, and $D_k+\tilde{b}\cdot P+i\gamma_3\cdot y=\mathfrak{K}+i\gamma_3\cdot y$ is invertible for every $y\ge0$. Proposition~\ref{prop:cf-geometry-large} therefore shows that $X_{k,y}(B)$ is invertible for every $y\ge0$, that is, $S_{k,y}(B)=X_{k,y}(B)^{-1}$ exists. Let $k=0$. Since $D_0=K=D_B+\mathfrak{m}$ and $\tilde{b}=b\cdot L^{-1}$, one has $\mathfrak{K}+i\gamma_3\cdot y=D_B+\mathfrak{m}+b\cdot L^{-1}\cdot P+i\gamma_3\cdot y$, whose inverse is the operator $S_{0,y}(B)$ of Section~\ref{sec:exact-fermion-operators}.
\end{proof}

\subsection{The formula for one renormalisation step}\label{sec:formula-one-renormalisation}

\subsubsection{Densities and the block transformation}\label{sec:densities-block-transformation}

Throughout this section the following conditions are assumed; they are part of the standing hypothesis~\ref{def:hypothesis-one-step} without its system of primitives. The constants are chosen as in Section~\ref{sec:order-choice}, subject to the conditions of Section~\ref{sec:conditions} that hypothesis~\ref{def:hypothesis-one-step} imposes; $k\ge 0$; the sides $l$ and $l/L$ belong to $\mathrm{Cl}$; the mass $\mathfrak{m}$ of Section~\ref{sec:block-operators} lies in $(0,\mathfrak{m}_0/L]$; and $e_{k+1}=L^{\frac{1}{2}}e_k\le e_0$.

The domains $\tilde{R}_k$, $\tilde{R}^{-}_k$, $\tilde{G}_k$, $\mathfrak{X}_k$ and $\mathrm{Dom}_k$ are those of Section~\ref{sec:domains-bond-functions}, and the spaces $K^{x,C}_k$ and $K^x_k$ of activities are those of Definition~\ref{def:renormalisation-group-1}. The numbers $\varepsilon(\mathscr{E})$ and $m(\mathscr{E})$, the operation $R$, the adjustment $E\mapsto E^{\mathrm{nat}}$ and the operation $L$ on families of activities are those of Section~\ref{sec:local-parts-adjustment}. The cube $\mathrm{Box}_k\subset V_k$, the standard Gaussian measure $\mu_I$ on $V_k$, the measure $\hat{\mu}$ and the number $d\varepsilon^0_k$ are those of Section~\ref{sec:averaging-fluctuations-gauge}; the cube $\mathrm{Box}_k$ is defined by means of the orthonormal basis $\{\Upsilon\}$ of $V_k$ of Proposition~\ref{prop:fluctuation-space}(i), (ii). The objects defined below are those of Section~\ref{sec:fermion-integrals-blocking}, written with the pairings $\langle\cdot,\cdot\rangle_{\eta_k}$, $\langle\cdot,\cdot\rangle_1$ and $\langle\cdot,\cdot\rangle_L$ of Section~\ref{sec:spinor-spaces-berezin}.

For a real bond function $A$ on $T^0_k$, $H_kA$ is the minimiser of Section~\ref{sec:averaging-fluctuations-gauge}; with a connection $B$ as argument, $H_k(B)$, $\bar{H}_k(B)$, $D_k(B)$ and $Q(B)$ are the operators of Section~\ref{sec:block-operators}, whereas $D_k$ without argument is the set of the polymers of Section~\ref{sec:lattice-geometry}. The generators $\Psi$, $\bar{\Psi}$ of Section~\ref{sec:spinor-spaces-berezin} are fields on $T^0_k$, and $\mathrm{Gr}[\Psi]$ is the Grassmann algebra of these generators (Section~\ref{sec:grassmann-algebras-norms}). The generators of level $k+1$ are denoted by $\Psi'$, $\bar{\Psi}'$; they are fields on $T^0_{k+1}$, and $\int\mathrm{D}\Psi'$ is the Berezin integral over them, normalised as in Section~\ref{sec:spinor-spaces-berezin} with the data of level $k+1$.

Let $(\mathscr{A},\mathscr{C})\in\tilde{R}_k\times\mathfrak{X}_k$ and $\mathscr{T}:=\mathscr{A}+\mathscr{C}$. The notation $\psi^{\#}=H^{\#}_k(\mathscr{T})\Psi$ stands for the four fields
\[
\psi=H_k(\mathscr{T})\Psi,\qquad \chi=\delta_{\alpha,\mathscr{T}}\psi,\qquad \bar{\psi}=\bar{H}_k(\mathscr{T})\bar{\Psi},\qquad \bar{\chi}=\delta_{\alpha,-\mathscr{T}}\bar{\psi},
\]
and, for $\mathscr{E}\in K^{x,C}_k$ and $X\in D_k$, the element $\mathscr{E}(X;\mathscr{A},\mathscr{C},\psi^{\#})$ of $\mathrm{Gr}[\Psi]$ is obtained from $\mathscr{E}(X;\mathscr{A},\mathscr{C},\psi,\chi)$ by inserting these fields for the generators $\psi$, $\chi$, $\bar{\psi}$, $\bar{\chi}$. For a complex number $m$ put
\begin{align*}
I^x_k[m,\mathscr{E}](\mathscr{A},\mathscr{C};\Psi)&:=-m\cdot\left\langle\bar{\psi},\psi\right\rangle_{\eta_k}+\sum_{X\in D_k}\mathscr{E}\left(X;\mathscr{A},\mathscr{C},\psi^{\#}\right),\\
\mathrm{Act}^x_k[m,\mathscr{E}](\mathscr{A},\mathscr{C};\Psi)&:=-\left\langle\bar{\Psi},D_k(\mathscr{T})\Psi\right\rangle_1+I^x_k[m,\mathscr{E}](\mathscr{A},\mathscr{C};\Psi).
\end{align*}
Let $\mathfrak{P}_k=\{\mathfrak{P}_k(X,\cdot):X\in D_k\}$ be an arbitrary family of complex functions on $\mathfrak{X}_k$; the families of Section~\ref{sec:flow} are of this kind. For a complex number $\boldsymbol{\varepsilon}$ and for $A\in\mathrm{Dom}_k$ put
\[
\mathscr{F}^x_k[\boldsymbol{\varepsilon},m,\mathscr{E}](A,\mathscr{C};\Psi):=\left(\det D_k(0)\right)^{-1}\cdot\exp\left\{\mathrm{Act}^x_k[m,\mathscr{E}](H_kA,\mathscr{C};\Psi)-\boldsymbol{\varepsilon}\cdot n_k+\sum_{X\in D_k}\mathfrak{P}_k(X,\mathscr{C})\right\}.
\]

Let $A'$ be a real bond function on $T^0_{k+1}$ and $\mathscr{C}'$ a complex bond function on $T_{k+1}$. From now on
\[
\mathscr{C}:=\mathscr{C}'_L,\qquad \mathscr{B}_k:=H_kA^{\min}_k(A'),\qquad c_k:=b^{-4n_{k+1}},
\]
and $\boldsymbol{\omega}_k$ denotes the real function $\boldsymbol{\omega}_k(A')$ of Theorem~\ref{thm:minimisation-one-step}(b), evaluated at the sites of $L\mathbb{Z}^3$. For a map $F$ with values in the Grassmann algebra of Section~\ref{sec:spinor-spaces-berezin}, defined at the arguments that occur on the right, put
\begin{multline*}
\left(\mathbf{T}^x_kF\right)(A',\mathscr{C}';\Psi'):=c_k\int_{\mathrm{Box}_k}\mathrm{d}\mu_I(\zeta)\int\mathrm{D}\Psi\,\exp\left\{-\tilde{b}\cdot\left\langle\exp\left(-ie_k\cdot\boldsymbol{\omega}_k\right)\bar{\Psi}'_L-Q\left(-\mathscr{B}_k-\mathscr{C}\right)\bar{\Psi},\right.\right.\\
\left.\left.\exp\left(ie_k\cdot\boldsymbol{\omega}_k\right)\Psi'_L-Q\left(\mathscr{B}_k+\mathscr{C}\right)\Psi\right\rangle_L\right\}\cdot F\left(A^{\min}_k(A')+\mathrm{Sq}_k\zeta,\mathscr{C};\Psi\right).
\end{multline*}

The elements $I^x_k[m,\mathscr{E}](\mathscr{A},\mathscr{C};\Psi)$ and $\mathrm{Act}^x_k[m,\mathscr{E}](\mathscr{A},\mathscr{C};\Psi)$ are even: the activities $\mathscr{E}(X;\mathscr{A},\mathscr{C},\psi,\chi)$ are even (Definition~\ref{def:renormalisation-group-1}), the inserted fields are linear in the generators, and the pairings are bilinear. Consequently the exponentials of even elements that occur in this section commute with one another.

\begin{lemma}[Three Berezin integrals]\label{lem:three-berezin-integrals}
\begin{enumerate}
\item[(a)] Let $\boldsymbol{\omega}$ be a real $l$-periodic function on $L\mathbb{Z}^3$, and let $f$, $\bar{f}$ be fields on $T^1_k$ whose values lie in the span of the monomials in the generators other than $\Psi'$, $\bar{\Psi}'$. Then
\[
c_k\int\mathrm{D}\Psi'\,\exp\left\{-\tilde{b}\cdot\left\langle\exp\left(-ie_k\cdot\boldsymbol{\omega}\right)\bar{\Psi}'_L-\bar{f},\ \exp\left(ie_k\cdot\boldsymbol{\omega}\right)\Psi'_L-f\right\rangle_L\right\}=1.
\]
\item[(b)] Assume the hypotheses of Proposition~\ref{prop:backgrounds-domains}(ii). Let $A'\in\mathrm{Dom}_{k+1}$ and $\mathscr{C}'\in\mathfrak{X}_{k+1}$, and let $F(\cdot,\mathscr{C}'_L;\Psi)$ be a bounded Borel map from $\mathrm{Dom}_k$ to the even part of $\mathrm{Gr}[\Psi]$. Then
\[
\int\mathrm{D}\Psi'\,\left(\mathbf{T}^x_kF\right)(A',\mathscr{C}';\Psi')=\int_{\mathrm{Box}_k}\mathrm{d}\mu_I(\zeta)\int\mathrm{D}\Psi\,F\left(A^{\min}_k(A')+\mathrm{Sq}_k\zeta,\mathscr{C}'_L;\Psi\right).
\]
\item[(c)] At level $0$ let the side be $l=n$ and the mass $\mathfrak{m}=\bar{m}\cdot\eta$, where $\eta=L^{-N}$, $\bar{m}$ and $n$ are those of Section~\ref{sec:lattice-measures-laws}, and let $(m_{\mathrm{in}},\boldsymbol{\varepsilon}_{\mathrm{in}})$ be a real pair. Then the function $y^x_0$ of Section~\ref{sec:fermion-integrals-blocking} satisfies, for $A_0\in\mathrm{Dom}_0$ and $a\in\mathfrak{X}_0$,
\[
y^x_0\left[\boldsymbol{\varepsilon}_{\mathrm{in}},m_{\mathrm{in}},0\right](A_0;a)=\exp\left(-\boldsymbol{\varepsilon}_{\mathrm{in}}\cdot n^3\right)\cdot\frac{\det\left(D_{A_0+a}+\bar{m}\cdot\eta+m_{\mathrm{in}}\right)}{\det\left(D_0+\bar{m}\cdot\eta\right)},
\]
where $D_0$ denotes the operator $D_B$ of Section~\ref{sec:block-operators} at $B=0$.
\end{enumerate}
\end{lemma}

\begin{proof}
(a) Define $u$ on $\mathbb{Z}^3$ by $u(y'):=\exp\left(-ie_k\cdot\boldsymbol{\omega}(Ly')\right)$. This function has no zeros, and it is $l/L$-periodic because $\boldsymbol{\omega}$ is $l$-periodic. By the definition of the spinor scaling in Section~\ref{sec:block-operators}, for $x\in L\mathbb{Z}^3$,
\[
\left(u\Psi'\right)_L(x)=\exp\left(-ie_k\cdot\boldsymbol{\omega}(x)\right)\Psi'_L(x),\qquad \left(u^{-1}\bar{\Psi}'\right)_L(x)=\exp\left(ie_k\cdot\boldsymbol{\omega}(x)\right)\bar{\Psi}'_L(x).
\]
Hence the substitution of $u^{-1}\bar{\Psi}'$, $u\Psi'$ for $\bar{\Psi}'$, $\Psi'$ turns the integrand into $\exp\left\{-\tilde{b}\cdot\left\langle\bar{\Psi}'_L-\bar{f},\Psi'_L-f\right\rangle_L\right\}$, and by~\eqref{eq:algebraic-gaussian-Ber3}, applied to the generators $\Psi'$, $\bar{\Psi}'$ on $T^0_{k+1}$, the integral is not changed by this substitution. Let $g$, $\bar{g}$ be the fields on $T^0_{k+1}$ with $g_L=f$ and $\bar{g}_L=\bar{f}$, that is, $g(y')=L\cdot f(Ly')$ and $\bar{g}(y')=L\cdot\bar{f}(Ly')$; their values lie in the same span as those of $f$ and $\bar{f}$. As $\Psi'_L-f=(\Psi'-g)_L$ and $\bar{\Psi}'_L-\bar{f}=(\bar{\Psi}'-\bar{g})_L$, the translation invariance~\eqref{eq:algebraic-gaussian-Ber1} gives
\[
c_k\int\mathrm{D}\Psi'\,\exp\left\{-\tilde{b}\cdot\left\langle\bar{\Psi}'_L-\bar{f},\Psi'_L-f\right\rangle_L\right\}=c_k\int\mathrm{D}\Psi'\,\exp\left\{-\tilde{b}\cdot\left\langle\bar{\Psi}'_L,\Psi'_L\right\rangle_L\right\}.
\]
By Lemma~\ref{lem:identities-block-operators}(f) with $a=1$, and since $\tilde{b}=bL^{-1}$ (Section~\ref{sec:block-operators}),
\[
\tilde{b}\cdot\left\langle\bar{\Psi}'_L,\Psi'_L\right\rangle_L=\tilde{b}\cdot L\cdot\left\langle\bar{\Psi}',\Psi'\right\rangle_1=b\cdot\left\langle\bar{\Psi}',\Psi'\right\rangle_1.
\]
The space $\mathfrak{S}(1,l/L)$ has the dimension $4(l/L)^3=4n_{k+1}$; so the normalisation of the Berezin integral, taken at level $k+1$, gives
\[
\int\mathrm{D}\Psi'\,\exp\left(-b\cdot\left\langle\bar{\Psi}',\Psi'\right\rangle_1\right)=\det(b\cdot 1)=b^{4n_{k+1}}.
\]
Since $c_k=b^{-4n_{k+1}}$, assertion (a) follows.

(b) Let $\zeta\in\mathrm{Box}_k$. By Proposition~\ref{prop:backgrounds-domains}(ii), $A^{\min}_k(A')+\mathrm{Sq}_k\zeta$ lies in $\mathrm{Dom}_k$, so that the integrand of $(\mathbf{T}^x_kF)(A',\mathscr{C}';\Psi')$ is defined. Let $G$ denote the exponential in the definition of $\mathbf{T}^x_kF$. It is even, and it does not depend on $\zeta$. The map $\zeta\mapsto A^{\min}_k(A')+\mathrm{Sq}_k\zeta$ is affine; hence $\zeta\mapsto F(A^{\min}_k(A')+\mathrm{Sq}_k\zeta,\mathscr{C}'_L;\Psi)$, and with it $\zeta\mapsto\int\mathrm{D}\Psi\,G\cdot F(A^{\min}_k(A')+\mathrm{Sq}_k\zeta,\mathscr{C}'_L;\Psi)$, is a bounded Borel function on $\mathrm{Box}_k$. The cube $\mathrm{Box}_k$, defined by means of the basis of Proposition~\ref{prop:fluctuation-space}(i), (ii), has finite measure $\mu_I(\mathrm{Box}_k)$. By~\eqref{eq:algebraic-gaussian-Ber5} the Berezin integral over $\Psi'$ therefore commutes with the integral over $\mathrm{Box}_k$, and, the element $G\cdot F$ being even, also with the Berezin integral over $\Psi$. The element $F(A^{\min}_k(A')+\mathrm{Sq}_k\zeta,\mathscr{C}'_L;\Psi)$ is even and contains none of the generators $\Psi'$, $\bar{\Psi}'$; since the integral over $\Psi'$ is the extraction of a coefficient, this element can be taken out of it. It follows that
\[
\int\mathrm{D}\Psi'\,\left(\mathbf{T}^x_kF\right)(A',\mathscr{C}';\Psi')=\int_{\mathrm{Box}_k}\mathrm{d}\mu_I(\zeta)\int\mathrm{D}\Psi\,\left(c_k\int\mathrm{D}\Psi'\,G\right)\cdot F\left(A^{\min}_k(A')+\mathrm{Sq}_k\zeta,\mathscr{C}'_L;\Psi\right).
\]
By Theorem~\ref{thm:minimisation-one-step}(b) the function $\boldsymbol{\omega}_k(A')$ is real; its restriction to $L\mathbb{Z}^3$ is $l$-periodic. The operators $Q(\mathscr{B}_k+\mathscr{C})$ and $Q(-\mathscr{B}_k-\mathscr{C})$ of Section~\ref{sec:block-operators} map fields on $T^0_k$ to fields on $T^1_k$; so $f:=Q(\mathscr{B}_k+\mathscr{C})\Psi$ and $\bar{f}:=Q(-\mathscr{B}_k-\mathscr{C})\bar{\Psi}$ are fields on $T^1_k$ with values in $\mathrm{Gr}[\Psi]$. Part~(a), applied with $\boldsymbol{\omega}=\boldsymbol{\omega}_k$ and these two fields, gives $c_k\int\mathrm{D}\Psi'\,G=1$. This proves (b); identity~\eqref{eq:renormalisation-group-3} follows.

(c) At level $0$ one has $\eta_0=1$ and $H_0(B)=\bar{H}_0(B)=1$, so that $\psi=\Psi$ and $\bar{\psi}=\bar{\Psi}$; moreover $D_0(B)=K(B)=D_B+\bar{m}\cdot\eta$ (Section~\ref{sec:block-operators}), $H_0A_0=A_0$, $n_0=n^3$, the activity is $0$, and $\mathfrak{P}_0=0$ (Section~\ref{sec:flow}). With $\mathscr{T}=A_0+a$ the definitions above give
\[
\mathscr{F}^x_0\left[\boldsymbol{\varepsilon}_{\mathrm{in}},m_{\mathrm{in}},0\right](A_0,a;\Psi)=\frac{\exp\left(-\boldsymbol{\varepsilon}_{\mathrm{in}}\cdot n^3\right)}{\det\left(D_0+\bar{m}\cdot\eta\right)}\cdot\exp\left\{-\left\langle\bar{\Psi},\left(D_{A_0+a}+\bar{m}\cdot\eta+m_{\mathrm{in}}\right)\Psi\right\rangle_1\right\}.
\]
By Section~\ref{sec:fermion-integrals-blocking}, $y^x_0\left[\boldsymbol{\varepsilon}_{\mathrm{in}},m_{\mathrm{in}},0\right](A_0;a)$ is the Berezin integral of this element over $\Psi$, and the normalisation of the Berezin integral (Section~\ref{sec:spinor-spaces-berezin}) turns the last exponential into $\det\left(D_{A_0+a}+\bar{m}\cdot\eta+m_{\mathrm{in}}\right)$. This proves (c).
\end{proof}

\subsubsection{The identity for one step with an external connection}\label{sec:identity-one-step}

\begin{theorem}[One-step formula with an external connection]\label{thm:one-step-formula}
Assume hypothesis~\ref{def:hypothesis-one-step} without its system of primitives, and assume the hypotheses of Lemmas~\ref{lem:inclusions-domains} and~\ref{lem:scaling-domains}, those of the two assertions of Lemma~\ref{lem:symmetries-domains} by which $\tfrac{1}{4}\tilde{R}_k$ contains $B-d\lambda$ along with each of its elements $B$, for every real function $\lambda$ on the sites of $T_k$, and $\tilde{G}_k$ contains $-B$ along with each of its elements $B$, namely $k\ge0$, $l\in\mathrm{Cl}$ and $0<e_k\le e_0$ (Section~\ref{sec:elementary-properties-field}), and those of Theorem~\ref{thm:expansion-walks} (together with Proposition~\ref{prop:operators-fermion-1}(d), this theorem yields the clauses of Propositions~\ref{prop:decay-fine-propagator} and~\ref{prop:kernel-bounds-2} recalled in Theorem~\ref{thm:fermion-operators-two}(d)), of Proposition~\ref{prop:backgrounds-domains}(ii) and of Theorem~\ref{thm:extraction}. These hypotheses follow from hypothesis~\ref{def:hypothesis-one-step} without its system of primitives, in which $0<e_k<e_{k+1}\le e_0$: those of the three lemmas, listed in Section~\ref{sec:elementary-properties-field}, are among its requirements; those of Theorem~\ref{thm:expansion-walks}, which for $k\ge1$ are also those of Proposition~\ref{prop:operators-fermion-1}(d), hold by stages~\ref{it:order-choice-4}, \ref{it:order-choice-6} and~\ref{it:order-choice-9}, by conditions~\eqref{eq:renormalisation-group-19}, \eqref{eq:renormalisation-group-20} and~\eqref{eq:renormalisation-group-21} and by $\mathfrak{m}\le\mathfrak{m}_0/L\le1$; those of Proposition~\ref{prop:backgrounds-domains}(ii) hold by conditions~\eqref{eq:renormalisation-group-A1}, \eqref{eq:renormalisation-group-A2} and~\eqref{eq:renormalisation-group-A10}, by $l/L\in\mathrm{Cl}$ and by Remark~\ref{rem:operators-photon-3}(ii); and the setting~\eqref{eq:polymer-activities-Sb} of Theorem~\ref{thm:extraction} is part of hypothesis~\ref{def:hypothesis-one-step}. Let $\boldsymbol{\varepsilon}$ and $m$ be complex numbers, let $\mathscr{E}\in K^{x,C}_k$, with $\mathscr{E}=0$ if $k=0$, and let $A'\in\mathrm{Dom}_{k+1}$ and $\mathscr{C}'\in\mathfrak{X}_{k+1}$. Define
\begin{gather*}
m':=m+m(\mathscr{E}),\qquad \mathscr{E}^r:=R\,\mathscr{E},\qquad m^{+}:=L\cdot m',\qquad \boldsymbol{\varepsilon}^{+}:=L^3\cdot\left(\boldsymbol{\varepsilon}+\varepsilon(\mathscr{E})+d\varepsilon^0_k\right),\\
\mathscr{A}':=H_{k+1}A',\qquad \mathscr{T}':=\mathscr{A}'+\mathscr{C}',\qquad \psi'^{\#}:=H^{\#}_{k+1}(\mathscr{T}')\Psi',\\
A:=\mathscr{A}'_L,\qquad \mathscr{C}:=\mathscr{C}'_L,\qquad \mathscr{T}:=A+\mathscr{C}=\mathscr{T}'_L,\qquad \psi^{0,\#}:=\left(\psi'^{\#}\right)_L.
\end{gather*}
Here $\psi'=H_{k+1}(\mathscr{T}')\Psi'$ and $\bar{\psi}'=\bar{H}_{k+1}(\mathscr{T}')\bar{\Psi}'$ are the unbarred and the barred spinor field contained in $\psi'^{\#}$, and $\psi^0=\psi'_L$ and $\bar{\psi}^0=\bar{\psi}'_L$ are those contained in $\psi^{0,\#}$, the scaling of Section~\ref{sec:block-operators} being applied to each component. Let $W$, $\bar{W}$ be the generators of the measure $\mathrm{d}\mu_I(W)$ of Section~\ref{sec:spinor-spaces-berezin}, which are fields on $T^0_k$, let $\mathsf{Z}\in\mathrm{Box}_k$, and define further
\begin{gather*}
\Psi^c:=T_k(\mathscr{T})\psi^0,\qquad \bar{\Psi}^c:=\bar{T}_k(\mathscr{T})\bar{\psi}^0,\qquad \mathscr{Z}:=H_k\,C\,C_k^{\frac{1}{2}}\mathsf{Z},\\
\Psi^f:=\Psi^c+\Gamma_k(\mathscr{T})W,\qquad \bar{\Psi}^f:=\bar{\Psi}^c+\bar{W},
\end{gather*}
\begin{multline*}
V^{\mathrm{ex},x}_k:=\mathrm{Act}^x_k[m',\mathscr{E}^r](A+\mathscr{Z},\mathscr{C};\Psi^f)-\mathrm{Act}^x_k[m',\mathscr{E}^r](A,\mathscr{C};\Psi^f)\\
+I^x_k[m',\mathscr{E}^r](A,\mathscr{C};\Psi^f)-I^x_k[m',\mathscr{E}^r](A,\mathscr{C};\Psi^c),
\end{multline*}
where a field in the last place of $I^x_k$ and of $\mathrm{Act}^x_k$ is understood as follows. Let $G$ be one of the elements $I^x_k[m,\mathscr{E}](\mathscr{A},\mathscr{C};\Psi)$ and $\mathrm{Act}^x_k[m,\mathscr{E}](\mathscr{A},\mathscr{C};\Psi)$ of Section~\ref{sec:densities-block-transformation}, and let $\Psi^s$, $\bar{\Psi}^s$ be two fields on $T^0_k$ that carry the same superscript $s$. The same symbol with $\Psi^s$ in the place of $\Psi$ denotes $G(\bar{\Psi}^s,\Psi^s)$, the image of $G$ under the homomorphism of Section~\ref{sec:spinor-spaces-berezin} that maps $\bar{\Psi}$ to $\bar{\Psi}^s$ and $\Psi$ to $\Psi^s$ and fixes the remaining generators. Thus the last argument $\Psi^f$, resp.\ $\Psi^c$, replaces the pair $(\Psi,\bar{\Psi})$ by $(\Psi^f,\bar{\Psi}^f)$, resp.\ by $(\Psi^c,\bar{\Psi}^c)$, whereas the operators $D_k$, $H_k$, $\bar{H}_k$ and the H\"older derivatives remain those of the total connection $\mathscr{A}+\mathscr{C}$ formed from the first two arguments; for instance,
\begin{multline*}
\mathrm{Act}^x_k[m',\mathscr{E}^r](A+\mathscr{Z},\mathscr{C};\Psi^f)=-\left\langle\bar{\Psi}^f,D_k(\mathscr{T}+\mathscr{Z})\Psi^f\right\rangle_1\\
-m'\cdot\left\langle\bar{H}_k(\mathscr{T}+\mathscr{Z})\bar{\Psi}^f,H_k(\mathscr{T}+\mathscr{Z})\Psi^f\right\rangle_{\eta_k}\\
+\sum_{X\in D_k}\mathscr{E}^r\left(X;A+\mathscr{Z},\mathscr{C},H^{\#}_k(\mathscr{T}+\mathscr{Z})\Psi^f\right),
\end{multline*}
where $H^{\#}_k(\mathscr{T}+\mathscr{Z})\Psi^f$ stands for the four fields of Section~\ref{sec:densities-block-transformation} formed at the connection $\mathscr{T}+\mathscr{Z}$ with $\Psi^f$, $\bar{\Psi}^f$ in the places of $\Psi$, $\bar{\Psi}$. Finally put
\[
\Xi^x_k:=\int\mathrm{d}\mu_I(W)\,\mathrm{d}\hat{\mu}(\mathsf{Z})\,\exp\left(V^{\mathrm{ex},x}_k\right).
\]
Then, for every $\mathsf{Z}\in\mathrm{Box}_k$, the pairs $(A,\mathscr{C})$ and $(A+\mathscr{Z},\mathscr{C})$ lie in $\tilde{R}_k\times\mathfrak{X}_k$ and the connections $\mathscr{T}$ and $\mathscr{T}+\mathscr{Z}$ lie in $\tilde{G}_k$; all objects introduced above are well defined, and the identity~\eqref{eq:renormalisation-group-9} holds, that is,
\begin{multline*}
\left(\mathbf{T}^x_k\mathscr{F}^x_k[\boldsymbol{\varepsilon},m,\mathscr{E}]\right)(A',\mathscr{C}';\Psi')=\left(\det D_{k+1}(0)\right)^{-1}\\
\cdot\exp\left\{-\left\langle\bar{\Psi}',D_{k+1}(\mathscr{T}')\Psi'\right\rangle_1-m^{+}\cdot\left\langle\bar{\psi}',\psi'\right\rangle_{\eta_{k+1}}-\boldsymbol{\varepsilon}^{+}\cdot n_{k+1}\right.\\
\left.+\sum_{X\in D_k}\mathscr{E}^r\left(X;A,\mathscr{C},\psi^{0,\#}\right)+\sum_{X\in D_k}\mathfrak{P}_k(X,\mathscr{C})\right\}\cdot\frac{\delta Z_k(\mathscr{T})}{\delta Z_k(0)}\cdot\Xi^x_k.
\end{multline*}
In particular, this identity holds for real $\boldsymbol{\varepsilon}$ and $m$ with $|m|\le w^{\mathrm{m}}_k$ and for $\mathscr{E}\in K^x_k$ with $\|\mathscr{E}\|^x_k\le w^{\mathrm{E}}_k$; the harmonic part of $A'$ is arbitrary.
\end{theorem}

\begin{proof}
\emph{Step~(0): domains.} By Theorem~\ref{thm:extraction}(a), $\varepsilon(\mathscr{E})$ and $m(\mathscr{E})$ are well defined complex numbers; so $m'$, $m^{+}$ and $\boldsymbol{\varepsilon}^{+}$ are defined. By the definition of $\mathrm{Dom}_{k+1}$ (Definition~\ref{def:domains}), the field $\mathscr{A}'=H_{k+1}A'$ lies in $\tfrac{1}{4}\tilde{R}_{k+1}$. Lemma~\ref{lem:scaling-domains} shows that $A=\mathscr{A}'_L$ lies in $\tilde{R}^{-}_k$, a subset of $\tilde{R}_k$, and that $\mathscr{C}=\mathscr{C}'_L$ lies in $\mathfrak{X}_k$. For $\zeta\in\mathrm{Box}_k$ let
\[
A_k:=A^{\min}_k(A')+\mathrm{Sq}_k\zeta .
\]
By Proposition~\ref{prop:backgrounds-domains}(ii), $A_k$ belongs to $\mathrm{Dom}_k$, which means that $H_kA_k$ lies in $\tfrac{1}{4}\tilde{R}_k$. The operator $H_k$ is linear (Proposition~\ref{prop:constrained-minimisers}(ii)), and $\mathrm{Sq}_k=C\,C_k^{\frac{1}{2}}$; therefore~\eqref{eq:algebraic-gaussian-14} gives
\begin{equation}\label{eq:algebraic-gaussian-28}
H_kA_k=A-d\boldsymbol{\omega}_k+\mathscr{Z},\qquad \mathscr{Z}=H_k\mathrm{Sq}_k\zeta=H_k\,C\,C_k^{\frac{1}{2}}\zeta .
\end{equation}
The function $\boldsymbol{\omega}_k=\boldsymbol{\omega}_k(A')$ is real (Theorem~\ref{thm:minimisation-one-step}(b)), and $\tfrac{1}{4}\tilde{R}_k$ is invariant under real gauge transformations (Lemma~\ref{lem:symmetries-domains}); hence $A+\mathscr{Z}=H_kA_k+d\boldsymbol{\omega}_k$ lies in $\tfrac{1}{4}\tilde{R}_k$, a subset of $\tilde{R}_k$. By Lemma~\ref{lem:inclusions-domains}(a), the set $\tilde{R}_k+\mathfrak{X}_k$ is contained in $\tilde{G}_k$. As $0$ belongs to $\tilde{R}_k$ and to $\mathfrak{X}_k$, it follows that the four connections
\[
\mathscr{T}=A+\mathscr{C},\qquad \mathscr{T}+\mathscr{Z}=(A+\mathscr{Z})+\mathscr{C},\qquad \mathscr{T}+\mathscr{Z}-d\boldsymbol{\omega}_k=H_kA_k+\mathscr{C},\qquad 0
\]
lie in $\tilde{G}_k$. For every $B\in\tilde{G}_k$, Theorem~\ref{thm:expansion-walks}(a) and Proposition~\ref{prop:operators-fermion-1}(d) show that $X_k(B)=D_B+\mathfrak{m}+b_kP_k(B)$ is invertible if $k\ge1$ and that $D_k(B)+\tilde{b}\cdot P(B)$ is invertible for every $k\ge0$, as recalled in Theorem~\ref{thm:fermion-operators-two}(d). Consequently the operators $S_k$ (for $k\ge1$), $H_k$, $D_k$ and $\Gamma_k$ exist at each of the four connections displayed above. Lemma~\ref{lem:algebraic-gaussian-1}(b),(c), applied at the real connection $0$ with the positive mass $\mathfrak{m}$, gives $\det D_k(0)\ne0$ and $\delta Z_k(0)\ne0$.

The cube $\mathrm{Box}_k$ of Section~\ref{sec:averaging-fluctuations-gauge} is defined by means of the basis $\{\Upsilon\}$ of $V_k$, which exists by Proposition~\ref{prop:fluctuation-space}(i),(ii); it is a compact subset of $V_k$. The integrand of $\mathbf{T}^x_k\mathscr{F}^x_k[\boldsymbol{\varepsilon},m,\mathscr{E}]$ is a continuous function of $\zeta$ on $\mathrm{Box}_k$, for the following reasons. The map $\zeta\mapsto H_kA_k$ is affine, with values in $\tfrac{1}{4}\tilde{R}_k$. The kernels of $\mathscr{E}$ are analytic functions of the connections, by property~\eqref{it:renormalisation-group-1-k1} of Definition~\ref{def:renormalisation-group-1}. The operators $D_k(B)$, $H_k(B)$ and $\bar{H}_k(B)=\mathsf{C}^{-1}H_k(-B)\mathsf{C}$ (Section~\ref{sec:block-operators}) are analytic functions of $B$ on $\tilde{G}_k$ by Theorem~\ref{thm:fermion-operators-two}(d), where $-B$ lies in $\tilde{G}_k$ together with $B$ by Lemma~\ref{lem:symmetries-domains}. The phases in $\delta_{\alpha,\pm B}$ (Definition~\ref{def:pairs-sites-holder}) are entire functions of $B$. For the same reasons $V^{\mathrm{ex},x}_k$ is a continuous function of $\mathsf{Z}$ on $\mathrm{Box}_k$. Finally, the operators of level $k+1$ exist at $\pm\mathscr{T}'$ and at $0$ by Theorem~\ref{thm:fermion-operators-two}(a), whose assumption holds at $B=\pm\mathscr{T}$ and at $B=0$ by Theorem~\ref{thm:fermion-operators-two}(d); and $\det D_{k+1}(0)\ne0$ by Theorem~\ref{thm:fermion-operators-two}(c), since $\det D_k(0)\ne0$.

\emph{Step~(a): extraction.} By Theorem~\ref{thm:extraction}(c),(d), for every $(\mathscr{A},\mathscr{C})\in\tilde{R}_k\times\mathfrak{X}_k$ one has
\[
\sum_{X\in D_k}\mathscr{E}(X;\mathscr{A},\mathscr{C},\psi,\chi)=-\varepsilon(\mathscr{E})\cdot n_k-m(\mathscr{E})\cdot\left\langle\bar{\psi},\psi\right\rangle_{\eta_k}+\sum_{X\in D_k}\mathscr{E}^r(X;\mathscr{A},\mathscr{C},\psi,\chi)
\]
as an identity between elements of the Grassmann algebra in the generators $\psi$, $\chi$, $\bar{\psi}$, $\bar{\chi}$, and $\mathscr{E}^r$ lies in $K^{x,C}_k$. Take $\mathscr{A}=H_kA_k$, which lies in $\tilde{R}_k$ by Step~(0), and insert the fields $H^{\#}_k(H_kA_k+\mathscr{C})\Psi$ for the generators. With the definitions of Section~\ref{sec:densities-block-transformation} this gives
\begin{multline*}
\mathscr{F}^x_k[\boldsymbol{\varepsilon},m,\mathscr{E}](A_k,\mathscr{C};\Psi)=\left(\det D_k(0)\right)^{-1}\\
\cdot\exp\left\{\mathrm{Act}^x_k[m',\mathscr{E}^r](H_kA_k,\mathscr{C};\Psi)-\left(\boldsymbol{\varepsilon}+\varepsilon(\mathscr{E})\right)\cdot n_k+\sum_{X\in D_k}\mathfrak{P}_k(X,\mathscr{C})\right\}.
\end{multline*}

\emph{Step~(b): phases.} Let $\boldsymbol{\Upsilon}_{\boldsymbol{\omega}_k}$ denote the multiplication by $\exp\left(ie_k\cdot\boldsymbol{\omega}_k\right)$ on the functions on each of the lattices $\mathbb{L}(\eta_k)$, $\mathbb{Z}^3$ and $L\mathbb{Z}^3$, as in Lemma~\ref{lem:identities-block-operators}(d). In the Berezin integral over $\Psi$ in the definition of $\mathbf{T}^x_k$ make the substitution
\[
\Psi=\boldsymbol{\Upsilon}_{\boldsymbol{\omega}_k}\Psi^n,\qquad \bar{\Psi}=\boldsymbol{\Upsilon}_{\boldsymbol{\omega}_k}^{-1}\bar{\Psi}^n;
\]
by~\eqref{eq:algebraic-gaussian-Ber3}, applied to the $l$-periodic function $\exp\left(ie_k\cdot\boldsymbol{\omega}_k\right)$ on $\mathbb{Z}^3$, which has no zeros, the integral over $\Psi$ equals the integral over the generators $\Psi^n$, $\bar{\Psi}^n$ of the transformed integrand. By~\eqref{eq:algebraic-gaussian-28}, the total connection at which $\mathscr{F}^x_k[\boldsymbol{\varepsilon},m,\mathscr{E}](A_k,\mathscr{C};\Psi)$ is formed is $H_kA_k+\mathscr{C}=\mathscr{T}+\mathscr{Z}-d\boldsymbol{\omega}_k$. Lemma~\ref{lem:identities-block-operators}(d), with the function $\boldsymbol{\omega}_k$ and the connection $\mathscr{T}+\mathscr{Z}$, gives
\begin{gather*}
\left\langle\bar{\Psi},D_k(\mathscr{T}+\mathscr{Z}-d\boldsymbol{\omega}_k)\Psi\right\rangle_1=\left\langle\bar{\Psi}^n,D_k(\mathscr{T}+\mathscr{Z})\Psi^n\right\rangle_1,\\
H_k(\mathscr{T}+\mathscr{Z}-d\boldsymbol{\omega}_k)\Psi=\boldsymbol{\Upsilon}_{\boldsymbol{\omega}_k}H_k(\mathscr{T}+\mathscr{Z})\Psi^n,\qquad \bar{H}_k(\mathscr{T}+\mathscr{Z}-d\boldsymbol{\omega}_k)\bar{\Psi}=\boldsymbol{\Upsilon}_{\boldsymbol{\omega}_k}^{-1}\bar{H}_k(\mathscr{T}+\mathscr{Z})\bar{\Psi}^n,
\end{gather*}
and, at every pair $(x,x')\in\mathrm{Pz}_k$, the H\"older derivative of $H_k(\mathscr{T}+\mathscr{Z}-d\boldsymbol{\omega}_k)\Psi$ formed with the connection $\mathscr{T}+\mathscr{Z}-d\boldsymbol{\omega}_k$ equals $\exp\left(ie_k\cdot\boldsymbol{\omega}_k(x)\right)$ times the H\"older derivative of $H_k(\mathscr{T}+\mathscr{Z})\Psi^n$ formed with the connection $\mathscr{T}+\mathscr{Z}$; for the barred fields $\bar{H}_k(\mathscr{T}+\mathscr{Z}-d\boldsymbol{\omega}_k)\bar{\Psi}$ and $\bar{H}_k(\mathscr{T}+\mathscr{Z})\bar{\Psi}^n$ the same holds with the factor $\exp\left(-ie_k\cdot\boldsymbol{\omega}_k(x)\right)$. Hence the fields $\psi^{\#}$ inserted in $\mathscr{F}^x_k[\boldsymbol{\varepsilon},m,\mathscr{E}](A_k,\mathscr{C};\Psi)$ are the gauge transforms, in the sense of property~\eqref{it:renormalisation-group-1-k4} of Definition~\ref{def:renormalisation-group-1} with $\boldsymbol{\lambda}=\boldsymbol{\omega}_k$ and $\boldsymbol{\lambda}'=0$, of the fields
\[
\psi^{n,\#}:=H^{\#}_k(\mathscr{T}+\mathscr{Z})\Psi^n.
\]
The pairing $\left\langle\bar{\psi},\psi\right\rangle_{\eta_k}$ is invariant under this gauge transformation. Both pairs $(A+\mathscr{Z}-d\boldsymbol{\omega}_k,\mathscr{C})$ and $(A+\mathscr{Z},\mathscr{C})$ lie in $\tilde{R}_k\times\mathfrak{X}_k$ (Step~(0)), and $\mathscr{E}^r\in K^{x,C}_k$ (Step~(a)); so property~\eqref{it:renormalisation-group-1-k4}, applied to $\mathscr{E}^r$, gives
\[
\mathscr{E}^r(X;A+\mathscr{Z}-d\boldsymbol{\omega}_k,\mathscr{C},\psi^{\#})=\mathscr{E}^r(X;A+\mathscr{Z},\mathscr{C},\psi^{n,\#})\qquad(X\in D_k).
\]
Together, these relations yield
\[
\mathrm{Act}^x_k[m',\mathscr{E}^r](H_kA_k,\mathscr{C};\Psi)=\mathrm{Act}^x_k[m',\mathscr{E}^r](A+\mathscr{Z},\mathscr{C};\Psi^n).
\]
In the Gaussian factor of $\mathbf{T}^x_k$ one has $\mathscr{B}_k+\mathscr{C}=\mathscr{T}-d\boldsymbol{\omega}_k$ by~\eqref{eq:algebraic-gaussian-14}, and Lemma~\ref{lem:identities-block-operators}(d) gives, on $L\mathbb{Z}^3$,
\[
Q(\mathscr{T}-d\boldsymbol{\omega}_k)\boldsymbol{\Upsilon}_{\boldsymbol{\omega}_k}\Psi^n=\boldsymbol{\Upsilon}_{\boldsymbol{\omega}_k}Q(\mathscr{T})\Psi^n,\qquad Q(-\mathscr{T}+d\boldsymbol{\omega}_k)\boldsymbol{\Upsilon}_{\boldsymbol{\omega}_k}^{-1}\bar{\Psi}^n=\boldsymbol{\Upsilon}_{\boldsymbol{\omega}_k}^{-1}Q(-\mathscr{T})\bar{\Psi}^n.
\]
So, at every site $z$ of $L\mathbb{Z}^3$, the first argument of the pairing $\langle\cdot,\cdot\rangle_L$ in the Gaussian factor takes the value $\exp\left(-ie_k\cdot\boldsymbol{\omega}_k(z)\right)\cdot\left(\bar{\Psi}'_L-Q(-\mathscr{T})\bar{\Psi}^n\right)(z)$, and the second argument the value $\exp\left(ie_k\cdot\boldsymbol{\omega}_k(z)\right)\cdot\left(\Psi'_L-Q(\mathscr{T})\Psi^n\right)(z)$. The pairing is a sum over these sites of products of the two values at the same site; the two phases cancel, and the quadratic form in the Gaussian factor becomes
\[
\tilde{b}\cdot\left\langle\bar{\Psi}'_L-Q(-\mathscr{T})\bar{\Psi}^n,\Psi'_L-Q(\mathscr{T})\Psi^n\right\rangle_L.
\]

\emph{Step~(c): splitting.} Since $A+\mathscr{C}=\mathscr{T}$, the definition of $\mathrm{Act}^x_k$ gives
\begin{multline*}
\mathrm{Act}^x_k[m',\mathscr{E}^r](A+\mathscr{Z},\mathscr{C};\Psi^n)=-\left\langle\bar{\Psi}^n,D_k(\mathscr{T})\Psi^n\right\rangle_1+I^x_k[m',\mathscr{E}^r](A,\mathscr{C};\Psi^n)\\
+\left\{\mathrm{Act}^x_k[m',\mathscr{E}^r](A+\mathscr{Z},\mathscr{C};\Psi^n)-\mathrm{Act}^x_k[m',\mathscr{E}^r](A,\mathscr{C};\Psi^n)\right\}.
\end{multline*}

\emph{Step~(d): completion of the square.} Apply Theorem~\ref{thm:fermion-operators-two} with $B=\mathscr{T}$ and $B'=\mathscr{T}'$. By~\eqref{eq:algebraic-gaussian-Ber1}, the Berezin integral over $\Psi^n$, $\bar{\Psi}^n$ does not change under the translation
\[
\Psi^n=\Psi^c+W',\qquad \bar{\Psi}^n=\bar{\Psi}^c+\bar{W}',
\]
where $\Psi^c$, $\bar{\Psi}^c$ are the fields~\eqref{eq:algebraic-gaussian-26}, whose values involve only the generators $\Psi'$, $\bar{\Psi}'$, and where $W'$, $\bar{W}'$ are the new generators of integration. By Theorem~\ref{thm:fermion-operators-two}(b), these fields are the fields $\Psi^c=T_k(\mathscr{T})\psi^0$ and $\bar{\Psi}^c=\bar{T}_k(\mathscr{T})\bar{\psi}^0$ of the statement, one has $\psi^{0,\#}=H^{\#}_k(\mathscr{T})\Psi^c$, and, by~\eqref{eq:algebraic-gaussian-27}, the quadratic forms of Steps~(c) and~(b) add up to
\begin{multline*}
\left\langle\bar{\Psi}^n,D_k(\mathscr{T})\Psi^n\right\rangle_1+\tilde{b}\cdot\left\langle\bar{\Psi}'_L-Q(-\mathscr{T})\bar{\Psi}^n,\Psi'_L-Q(\mathscr{T})\Psi^n\right\rangle_L\\
=\left\langle\bar{\Psi}',D_{k+1}(\mathscr{T}')\Psi'\right\rangle_1+\left\langle\bar{W}',\Gamma_k(\mathscr{T})^{-1}W'\right\rangle_1.
\end{multline*}
Now apply~\eqref{eq:algebraic-gaussian-Ber4} with $O=\Gamma_k(\mathscr{T})^{-1}=D_k(\mathscr{T})+\tilde{b}\cdot P(\mathscr{T})$, whose determinant is $\delta Z_k(\mathscr{T})$ by the definition of $\delta Z_k$ in Section~\ref{sec:block-operators}. It shows that the Berezin integral over $W'$, $\bar{W}'$ equals $\delta Z_k(\mathscr{T})$ times the integral with respect to $\mathrm{d}\mu_I(W)$ of the remaining integrand, evaluated at $W'=\Gamma_k(\mathscr{T})W$ and $\bar{W}'=\bar{W}$, that is, at $\Psi^n=\Psi^f$ and $\bar{\Psi}^n=\bar{\Psi}^f$.

\emph{Step~(e): collection.} Write $\mathsf{Z}$ in place of $\zeta$. Steps~(a)--(d) give
\begin{multline*}
\left(\mathbf{T}^x_k\mathscr{F}^x_k[\boldsymbol{\varepsilon},m,\mathscr{E}]\right)(A',\mathscr{C}';\Psi')=c_k\cdot\left(\det D_k(0)\right)^{-1}\cdot\delta Z_k(\mathscr{T})\\
\cdot\exp\left\{-\left\langle\bar{\Psi}',D_{k+1}(\mathscr{T}')\Psi'\right\rangle_1-\left(\boldsymbol{\varepsilon}+\varepsilon(\mathscr{E})\right)\cdot n_k+\sum_{X\in D_k}\mathfrak{P}_k(X,\mathscr{C})\right\}\\
\cdot\int_{\mathrm{Box}_k}\mathrm{d}\mu_I(\mathsf{Z})\int\mathrm{d}\mu_I(W)\,\exp\left\{I^x_k[m',\mathscr{E}^r](A,\mathscr{C};\Psi^f)\right.\\
\left.+\mathrm{Act}^x_k[m',\mathscr{E}^r](A+\mathscr{Z},\mathscr{C};\Psi^f)-\mathrm{Act}^x_k[m',\mathscr{E}^r](A,\mathscr{C};\Psi^f)\right\}.
\end{multline*}
By the definition of $V^{\mathrm{ex},x}_k$, the exponent of the last integrand equals $I^x_k[m',\mathscr{E}^r](A,\mathscr{C};\Psi^c)+V^{\mathrm{ex},x}_k$, and the first of these two terms depends neither on $W$ nor on $\mathsf{Z}$. As $A+\mathscr{C}=\mathscr{T}$ and $\psi^{0,\#}=H^{\#}_k(\mathscr{T})\Psi^c$ (Step~(d)),
\[
I^x_k[m',\mathscr{E}^r](A,\mathscr{C};\Psi^c)=-m'\cdot\left\langle\bar{\psi}^0,\psi^0\right\rangle_{\eta_k}+\sum_{X\in D_k}\mathscr{E}^r\left(X;A,\mathscr{C},\psi^{0,\#}\right),
\]
and Lemma~\ref{lem:identities-block-operators}(f) gives
\[
\left\langle\bar{\psi}^0,\psi^0\right\rangle_{\eta_k}=\left\langle\bar{\psi}'_L,\psi'_L\right\rangle_{\eta_k}=L\cdot\left\langle\bar{\psi}',\psi'\right\rangle_{\eta_{k+1}},
\]
so that $m'\cdot\left\langle\bar{\psi}^0,\psi^0\right\rangle_{\eta_k}=m^{+}\cdot\left\langle\bar{\psi}',\psi'\right\rangle_{\eta_{k+1}}$. By the definitions of $\hat{\mu}$ and of $d\varepsilon^0_k$ in Section~\ref{sec:averaging-fluctuations-gauge}, the restriction of $\mu_I$ to $\mathrm{Box}_k$ is the measure
\[
\mu_I(\mathrm{Box}_k)\cdot\hat{\mu}=\exp\left(-d\varepsilon^0_k\cdot n_k\right)\cdot\hat{\mu},
\]
and $n_k=L^3n_{k+1}$; hence $\left(\boldsymbol{\varepsilon}+\varepsilon(\mathscr{E})\right)\cdot n_k+d\varepsilon^0_k\cdot n_k=\boldsymbol{\varepsilon}^{+}\cdot n_{k+1}$. By~\eqref{eq:algebraic-gaussian-Ber5} the integrations over $W$ and over $\mathsf{Z}$ may be interchanged, the integrand being a continuous function of $\mathsf{Z}$ on the compact set $\mathrm{Box}_k$ (Step~(0)). Inserting these relations into the last display, one obtains the right side of the asserted identity, except that the factor $\left(\det D_{k+1}(0)\right)^{-1}\cdot\delta Z_k(\mathscr{T})/\delta Z_k(0)$ is replaced by $c_k\cdot\left(\det D_k(0)\right)^{-1}\cdot\delta Z_k(\mathscr{T})$.

\emph{Step~(f): the constant.} Theorem~\ref{thm:fermion-operators-two}(c) at $B'=0$ gives $\det D_{k+1}(0)=b^{4n_{k+1}}\cdot\det D_k(0)/\delta Z_k(0)$. Since $c_k=b^{-4n_{k+1}}$, it follows that
\[
c_k\cdot\left(\det D_k(0)\right)^{-1}\cdot\delta Z_k(\mathscr{T})=\left(\det D_{k+1}(0)\right)^{-1}\cdot\frac{\delta Z_k(\mathscr{T})}{\delta Z_k(0)},
\]
which completes the proof of the identity. If $\boldsymbol{\varepsilon}$ and $m$ are real with $|m|\le w^{\mathrm{m}}_k$, and if $\mathscr{E}\in K^x_k$ satisfies $\|\mathscr{E}\|^x_k\le w^{\mathrm{E}}_k$, then $\mathscr{E}$ belongs to $K^{x,C}_k$, and the theorem applies; thus Proposition~\ref{prop:step-external-connection} follows.
\end{proof}

\begin{remark}\label{rem:algebraic-gaussian-1}
\begin{enumerate}
\item[(i)] The additive constant of $\boldsymbol{\omega}_k$, which Theorem~\ref{thm:minimisation-one-step}(b) fixes by the condition $Q_{k+1}\boldsymbol{\omega}_k=0$, is immaterial: if $\boldsymbol{\omega}_k$ is replaced by $\boldsymbol{\omega}_k+t$ with a real number $t$, Steps~(b)--(f) of the proof of Theorem~\ref{thm:one-step-formula} remain unchanged.
\item[(ii)] For $k=0$ one has $T_0=H_0=1$, $\psi=\Psi$, $\mathscr{E}=0$, $m'=m$ and $D_0(\mathscr{T})=D_{\mathscr{T}}+\mathfrak{m}$; the proof then applies Theorem~\ref{thm:fermion-operators-two} with $k=0$, in which $b_0$ does not occur.
\item[(iii)] (The twist.) A constant $c'$ contained in $A'$ produces the constant $c'$ in $\mathscr{A}'$ (Proposition~\ref{prop:constrained-minimisers}(iii)) and the constant $L^{-\frac{1}{2}}c'$ in $A$; on the right side of the identity of Theorem~\ref{thm:one-step-formula} it occurs only through these two constants, since there $A'$ enters only through $\mathscr{A}'=H_{k+1}A'$: the connections $\mathscr{T}'$, $A$ and $\mathscr{T}$, the fields $\psi'^{\#}$, $\psi^{0,\#}$, $\Psi^c$ and $\bar{\Psi}^c$, the elements $V^{\mathrm{ex},x}_k$ and $\Xi^x_k$ and the number $\delta Z_k(\mathscr{T})$ are formed from $\mathscr{A}'$, $\mathscr{C}'$, $\boldsymbol{\varepsilon}$, $m$, $\mathscr{E}$, $\mathsf{Z}$ and the generators, and the remaining ingredients $\mathscr{E}^r=R\,\mathscr{E}$, $\mathscr{C}=\mathscr{C}'_L$, $\det D_{k+1}(0)$, $\delta Z_k(0)$, $m^{+}$, $\boldsymbol{\varepsilon}^{+}$, $n_{k+1}$ and $\mathfrak{P}_k(X,\mathscr{C})$ do not depend on $A'$ (the activity $\mathscr{E}^r$ is formed from $\mathscr{E}$ alone, and $\mathscr{C}$ from $\mathscr{C}'$ alone). On the left side, that is, in the definition of $\mathbf{T}^x_k$ in Section~\ref{sec:densities-block-transformation}, and in Steps~(0) and~(b) of the proof, $A'$ enters in addition through $A^{\min}_k(A')$, through $\mathscr{B}_k=H_kA^{\min}_k(A')=A-d\boldsymbol{\omega}_k(A')$ (see~\eqref{eq:algebraic-gaussian-14} and~\eqref{eq:algebraic-gaussian-28}) and through $\boldsymbol{\omega}_k(A')$, all three linear in $A'$ (Theorem~\ref{thm:minimisation-one-step}(a), (b) and Proposition~\ref{prop:constrained-minimisers}(ii)); so $c'$ occurs there also through $A^{\min}_k(c')$ and $\boldsymbol{\omega}_k(c')$. If $c'\ne0$, the function $\boldsymbol{\omega}_k(c')$ is not constant on $T_k$. Indeed, if $d\boldsymbol{\omega}_k(c')=0$, then~\eqref{eq:algebraic-gaussian-14}, Proposition~\ref{prop:constrained-minimisers}(iii) and~\eqref{eq:algebraic-gaussian-3} give
\[
H_kA^{\min}_k(c')=\left(H_{k+1}c'\right)_L=L^{-\frac{1}{2}}c'=H_k\left(L^{-\frac{1}{2}}c'\right),
\]
hence $A^{\min}_k(c')=L^{-\frac{1}{2}}c'$, as $Q_kH_k=1$ (Proposition~\ref{prop:constrained-minimisers}(ii)); this contradicts $\tau A^{\min}_k(c')=0$ (Theorem~\ref{thm:minimisation-one-step}(a)), because $(\tau c)(y,y+e_{\mu})=c_{\mu}$ for every constant bond function $c$ on $T^0_k$, every $y\in L\mathbb{Z}^3$ and every direction $\mu$, by~\eqref{eq:algebraic-gaussian-7}: the site $y+e_{\mu}$ lies in the block $B(y)$, as $L\ge3$, and each path $\boldsymbol{\Gamma}^{\pi}(y,y+e_{\mu})$ consists of the single bond $(y,y+e_{\mu})$. The constant $c'$ requires no separate argument: Theorem~\ref{thm:minimisation-one-step} holds for an arbitrary harmonic part of $A'$ and supplies the real function $\boldsymbol{\omega}_k(A')$, which Step~(b) of the proof removes by gauge covariance, whatever the harmonic part of $A'$ is.
\end{enumerate}
\end{remark}

\section{Operators of the photon field}\label{sec:operators-photon-field}

In this section the propositions of Section~\ref{sec:properties-exact-photon} and Theorem~\ref{thm:photon-operators-weakened} are proved, and the two photon families of a system of primitives are constructed, with the properties stated in Theorem~\ref{thm:properties-weakened-photon} and Proposition~\ref{prop:family-level-zero}; the notation and the constants are those of Section~\ref{sec:renormalisation-group-scheme}, the stages are the items of Definition~\ref{def:order-choice}, and a property cited by its number is one of the list in Section~\ref{sec:properties-systems-primitives}. Several facts that Section~\ref{sec:algebraic-gaussian-identities} derives algebraically, concerning the block averages, the operator $H_k$, the axial gauges, the constrained minimiser and the basis of the fluctuation space, are proved here again in the form required by the bounds of this section.

\paragraph{Standing conditions}
The ceiling of the charge satisfies condition~\eqref{eq:renormalisation-group-25}, that is,
\[
e_0\le\exp(-1).
\]
This condition is contained in the conditions~\eqref{eq:renormalisation-group-A9}: the first inequality of~\eqref{eq:renormalisation-group-A9} gives $e_0^{\varepsilon}\le\frac{1}{160}$, hence $e_0\le 160^{-200}$, because $\varepsilon\le\frac{1}{200}$ (stages~\ref{it:order-choice-1} and~\ref{it:order-choice-8}). For every $e_k\in(0,e_0]$ it follows that
\[
v_k=-\log e_k\ge 1,\qquad p_{0,k}=v_k^{p_0}\ge 1,
\]
and $p_{0,k}\ge 1$ is the hypothesis of the bound $d\varepsilon^0_k\le 12\exp\left(-\frac{p_{0,k}^2}{2}\right)$ in Proposition~\ref{prop:basis-map-box}.

The thresholds $L^P$, $\alpha^P$, $M^P$ that occur in stages~\ref{it:order-choice-4}, \ref{it:order-choice-5} and~\ref{it:order-choice-6} are
\[
L^P(\alpha):=5,\qquad \alpha^P:=\frac{1}{4},
\]
and the number $M^P(\alpha,L,\alpha_0)$ defined in \eqref{eq:operators-photon-20}. Of $L$ the proofs use only that $L$ is odd and $L\ge 3$, and, in Proposition~\ref{prop:fields-domains}(ii), that $L^{\frac{3}{4}}\ge\frac{8}{3}$; both requirements are contained in stage~\ref{it:order-choice-4}. No condition is imposed on $\kappa$, and none on $M$ except $M\ge M^P$ and $M\ge L^2$.

\subsection{Notation and elementary facts}\label{sec:notation-elementary-facts}

\subsubsection{Tori, cochains and kernels}\label{sec:tori-cochains-kernels}

Unless stated otherwise, the data of Section~\ref{sec:operators-photon-field} are those of the scale-free setting of Section~\ref{sec:properties-scheme}: an integer $k\ge 0$ and a side $l\in\mathrm{Cl}$. In Sections~\ref{sec:block-averages-minimiser}--\ref{sec:axial-gauges} the side $l$ is an arbitrary positive integer, not necessarily an element of $\mathrm{Cl}$, and wherever blocks occur in these sections, $l$ is a multiple of the side of the blocks. Section~\ref{sec:operators-photon-field} involves no mass. The charge $e_k$ appears explicitly in Section~\ref{sec:local-bounds-background}. In Sections~\ref{sec:notation-elementary-facts}--\ref{sec:constrained-minimisation-covariance} and~\ref{sec:weakened-photon-operators} it appears only through two objects. The first is the number $p_{0,k}=(-\log e_k)^{p_0}$, which enters the set $\mathrm{Box}_k$, the measure $\hat{\mu}$ and the number $d\varepsilon^0_k$ of Section~\ref{sec:averaging-fluctuations-gauge} (Propositions~\ref{prop:fluctuation-space}(iii), (v) and~\ref{prop:lattice-symmetries}), and the inequality $p_{0,k}\ge 1$ derived above from condition~\eqref{eq:renormalisation-group-25}. The second is the set $\tilde{G}_k$ of Definition~\ref{def:domains}(e), which occurs only in Section~\ref{sec:properties-two-families}, and there only as the domain of the connection $A$ in properties~\ref{it:renormalisation-group-2}, \ref{it:renormalisation-group-8}, \ref{it:renormalisation-group-9} and~\ref{it:renormalisation-group-14}; the families~\eqref{eq:operators-photon-21} and~\eqref{eq:operators-photon-22} do not depend on $A$. The letters $\mu,\nu,\mathfrak{l},\mathfrak{l}'$ denote directions in $\{0,1,2\}$, and $e_{\mu}$ is the unit vector in the direction $\mu$.

Put
\begin{align*}
\mathfrak{N}&:=L^k, & \eta&:=\eta_k=\frac{1}{\mathfrak{N}}, & \mathbf{h}&:=\frac{\mathfrak{N}-1}{2},\\
T_k&:=\eta\mathbb{Z}^3/l\mathbb{Z}^3, & T^0_k&:=\mathbb{Z}^3/l\mathbb{Z}^3, & B_k(y)&:=y+\eta\{-\mathbf{h},\dots,\mathbf{h}\}^3\quad(y\in T^0_k).
\end{align*}
The integer $\mathfrak{N}$ is odd. In Sections~\ref{sec:block-averages-minimiser}--\ref{sec:bounds-complex-strip} the numbers $L$ and $k$ enter only through $\mathfrak{N}$, and of $\mathfrak{N}$ only the fact that it is an odd positive integer is used.

For a vector $x$, $|x|$ is the sup norm, and $|x|_1$, $|x|_2$ are the $\ell^1$ and $\ell^2$ norms; $d(\cdot,\cdot)$ is the sup distance on the torus (Definition~\ref{def:sup-norm-distance}). The indices are the sites, the bonds and the basis elements $\Upsilon$ of Section~\ref{sec:blocks-fluctuation-space}. For an index $x$, $\operatorname{supp}x$ is the set of the end points of $x$: the site itself if $x$ is a site, the two ends if $x$ is a bond, and the end points of the bonds on which $\Upsilon$ does not vanish if $x=\Upsilon$. For two indices $x,x'$, $d(x,x')$ is the distance between $\operatorname{supp}x$ and $\operatorname{supp}x'$, and $x$ is contained in a set $W$ if $\operatorname{supp}x\subset W$ (Section~\ref{sec:lattice-geometry}). For real $t$ put $\langle t\rangle:=\max(1,|t|)$, and for $m\in\mathbb{Z}^3$ put $\langle m\rangle_{\infty}:=\max_{\nu}\langle m_{\nu}\rangle$.

\paragraph{Cochains}
A bond function $\mathscr{A}$ on $T_k$ changes sign under reversal of the bonds, $\mathscr{A}(x',x)=-\mathscr{A}(x,x')$, and its components are the functions $\mathscr{A}_{\mu}(x):=\mathscr{A}(x,x+\eta e_{\mu})$ of the initial point. For functions $f,g$ on the sites, a scalar $\lambda$ and a bond function $\mathscr{A}$ put
\begin{align*}
(\partial_{\mu}f)(x)&:=\eta^{-1}(f(x+\eta e_{\mu})-f(x)),\qquad (d\lambda)_{\mu}:=\partial_{\mu}\lambda,\\
(d\mathscr{A})_{\mu\nu}&:=\partial_{\mu}\mathscr{A}_{\nu}-\partial_{\nu}\mathscr{A}_{\mu}\quad(\mu<\nu),\\
\langle f,g\rangle&:=\sum_x\eta^3f(x)\cdot g(x),\qquad |f|_2^2:=\langle f,f\rangle,
\end{align*}
where for bond functions and plaquette functions the pairing is summed also over the components; further $\delta:=d^T$, and $\Delta:=\delta d$ on scalars. On $T^0_k$ the same definitions are used with $\eta=1$. It follows that
\[
(\delta\mathscr{A})(x)=\sum_{\mu}(\partial_{\mu}^T\mathscr{A}_{\mu})(x),\qquad (\partial_{\mu}^Tf)(x)=\eta^{-1}(f(x-\eta e_{\mu})-f(x)),
\]
and that on bond functions $\delta d=\Delta-d\delta$, where $\Delta:=\sum_{\nu}\partial_{\nu}^T\partial_{\nu}$ acts on each component. The vector $\left(l^{-3}\langle\mathscr{A}_{\mu},1\rangle\right)_{\mu=0,1,2}$ is called the mean of the bond function $\mathscr{A}$. For a path $\Gamma$ of bonds the line sum is
\[
\mathscr{A}(\Gamma):=\sum_{b\in\Gamma}\mathscr{A}(b),
\]
without weight, each bond $b$ of $\Gamma$ being oriented in the direction in which $\Gamma$ traverses it. A bond function $\mathscr{A}$ is closed if $d\mathscr{A}=0$, coclosed if $\delta\mathscr{A}=0$, and exact if $\mathscr{A}=d\lambda$ for a scalar $\lambda$.

The next lemma is the Hodge decomposition on the discrete torus; it is the discrete form of the classical fact that the first cohomology group of the three-dimensional torus with real coefficients is $\mathbb{R}^3$.

\begin{lemma}[Hodge decomposition on the discrete torus]
Let $\mathscr{A}$ be a bond function on $T_k$ or on $T^0_k$. If $\mathscr{A}$ is closed, then $\mathscr{A}=c+d\lambda$ with a constant bond function $c$ and a scalar $\lambda$, and $\lambda$ can be taken real if $\mathscr{A}$ is real. Every exact bond function has mean zero. Consequently, if $\mathscr{A}$ is closed and its mean vanishes, then $\mathscr{A}$ is exact; if $\mathscr{A}$ is both closed and coclosed, then $\mathscr{A}$ is constant.
\end{lemma}

\begin{proof}
The Fourier transform $f\mapsto\hat f$ and the functions $\hat{d}_{\nu}$, $\hat{d}^{*}_{\nu}$, $\hat{D}$ are those of Section~\ref{sec:elementary-symbols-fourier}; on $T^0_k$ they are taken with $\mathfrak{N}=1$. A function is recovered from its transform by the inversion formula
\[
f(x)=l^{-3}\sum_q\exp(iq\cdot x)\cdot\hat f(q).
\]
By \eqref{eq:operators-photon-4} the transform of $\partial_{\nu}f$ is $\hat{d}_{\nu}(q)\cdot\hat f(q)$. Hence $d\mathscr{A}=0$ reads
\[
\hat{d}_{\mu}(q)\cdot\hat{\mathscr{A}}_{\nu}(q)=\hat{d}_{\nu}(q)\cdot\hat{\mathscr{A}}_{\mu}(q)\qquad\text{for all }q\text{ and all }\mu,\nu.
\]
For $q\notin 2\pi\mathfrak{N}\mathbb{Z}^3$ one has $\hat{D}(q)>0$, because $\hat{D}(q)=\sum_{\nu}\operatorname{sn}(q_{\nu})^2$ and $\operatorname{sn}(q_{\nu})^2=4\mathfrak{N}^2\cdot\sin^2\left(\frac{q_{\nu}}{2\mathfrak{N}}\right)$ vanishes only for $q_{\nu}\in 2\pi\mathfrak{N}\mathbb{Z}$. For these $q$ put
\[
\hat\lambda(q):=\frac{1}{\hat{D}(q)}\sum_{\mu}\hat{d}^{*}_{\mu}(q)\cdot\hat{\mathscr{A}}_{\mu}(q),
\]
and put $\hat\lambda(0):=0$; the inversion formula, with $\hat\lambda$ in place of $\hat f$, defines a scalar $\lambda$. Since $\hat{D}=\sum_{\mu}\hat{d}^{*}_{\mu}\cdot\hat{d}_{\mu}$, the closedness relation gives, for $q\notin 2\pi\mathfrak{N}\mathbb{Z}^3$,
\[
\hat{d}_{\nu}(q)\cdot\hat\lambda(q)=\frac{1}{\hat{D}(q)}\sum_{\mu}\hat{d}^{*}_{\mu}(q)\cdot\hat{d}_{\mu}(q)\cdot\hat{\mathscr{A}}_{\nu}(q)=\hat{\mathscr{A}}_{\nu}(q),
\]
so that the transforms of $d\lambda$ and of $\mathscr{A}$ coincide at these $q$. At $q=0$ the transform of $d\lambda$ vanishes, since $\hat{d}_{\nu}(0)=0$, and the term $q=0$ of the inversion formula for $\mathscr{A}$ is its mean $c$. Hence $\mathscr{A}=c+d\lambda$. If $\mathscr{A}$ is real, then so is $c$, and the real parts give $\mathscr{A}=c+d\operatorname{Re}\lambda$.

The vanishing of the transform of $d\lambda$ at $q=0$ says that $d\lambda$ has mean zero, for every scalar $\lambda$. If $\mathscr{A}$ is closed and has mean zero, then $c=0$, and $\mathscr{A}=d\lambda$ is exact.

Let $\mathscr{A}$ be real, closed and coclosed, and take $\lambda$ real. Since $\delta c=0$ for the constant $c$, one has $\Delta\lambda=\delta d\lambda=\delta\mathscr{A}=0$; hence $|d\lambda|_2^2=\langle\lambda,\Delta\lambda\rangle=0$, and $\mathscr{A}=c$. If $\mathscr{A}$ is complex, closed and coclosed, then its real and imaginary parts are closed and coclosed, because $d$ and $\delta$ map real functions to real functions; so they are constant, and so is $\mathscr{A}$.
\end{proof}

\paragraph{Kernels}
Let $O$ be a linear map from the bond functions on $T^0_k$ to the bond functions on $T_k$. Its kernel is defined by
\[
(OA)(b)=\sum_{b'}O(b,b')\cdot A(b'),
\]
where $b$ is a positively oriented bond of $T_k$ and $b'$ ranges over the positively oriented bonds of $T^0_k$. For $b=(x,x+\eta e_{\mu})$ and $\tilde b=(\tilde x,\tilde x+\eta e_{\mu})$ with $0<|x-\tilde x|\le 1$ put
\begin{align*}
(\partial_{\mathfrak{l}}O)(b,b')&:=\eta^{-1}(O(b+\eta e_{\mathfrak{l}},b')-O(b,b')),\\
(\delta_{\alpha}\partial_{\mathfrak{l}}O)((b,\tilde b);b')&:=|x-\tilde x|^{-\alpha}((\partial_{\mathfrak{l}}O)(\tilde b,b')-(\partial_{\mathfrak{l}}O)(b,b')),
\end{align*}
where $b+\eta e_{\mathfrak{l}}$ is the translate of $b$ by $\eta e_{\mathfrak{l}}$. A bound on $|(\partial O)(b,b')|$ is meant for $\mathfrak{l}=0,1,2$. Since $(OA)_{\mu}(x)=(OA)(b)$ and $(OA)_{\mu}(x+\eta e_{\mathfrak{l}})=(OA)(b+\eta e_{\mathfrak{l}})$, and likewise with $\tilde x$, $\tilde b$ in place of $x$, $b$, the definition of the kernel gives, for every bond function $A$ on $T^0_k$,
\begin{align*}
(\partial_{\mathfrak{l}}(OA)_{\mu})(x)&=\sum_{b'}(\partial_{\mathfrak{l}}O)(b,b')\cdot A(b'),\\
|x-\tilde x|^{-\alpha}\left((\partial_{\mathfrak{l}}(OA)_{\mu})(\tilde x)-(\partial_{\mathfrak{l}}(OA)_{\mu})(x)\right)&=\sum_{b'}(\delta_{\alpha}\partial_{\mathfrak{l}}O)((b,\tilde b);b')\cdot A(b').
\end{align*}
Hence, if $|O(b,b')|$, $|(\partial O)(b,b')|$ and $|(\delta_{\alpha}\partial O)((b,\tilde b);b')|$ are at most $K(b,b')$ for all $b,\tilde b,b'$, then each of the three norms of $OA$ (Section~\ref{sec:domains-bond-functions}) is at most
\[
\sup_b\sum_{b'}K(b,b')\cdot|A(b')|.
\]

\paragraph{Lattice sums}
For $0<g\le 1$ put
\begin{multline}\label{eq:operators-photon-1}
\boldsymbol{\varsigma}(g):=\sum_{x\in\mathbb{Z}^3}\exp(-g|x|)\le\left(1+\frac{6}{g}\right)^3,\\
c_{\mathrm{bd}}(g):=3\exp(3g)\cdot\boldsymbol{\varsigma}(g),\qquad c_{\mathrm{bas}}(g):=2L^3\exp(2gL)\cdot\boldsymbol{\varsigma}(g).
\end{multline}
The inequality holds because $|x|\ge\frac{1}{3}|x|_1$, so that $\boldsymbol{\varsigma}(g)$ is at most the third power of
\[
\sum_{t\in\mathbb{Z}}\exp\left(-\frac{g|t|}{3}\right)=\coth\frac{g}{6}\le 1+\frac{6}{g}
\]
(for $s>0$ one has $\coth s=1+\frac{2}{\exp(2s)-1}\le 1+\frac{1}{s}$).

\begin{lemma}
Let $0<g\le 1$. On every torus $T_k$, $T^0_k$, for every bond $b$, fine or unit, and every element $\Upsilon$ of the basis of Section~\ref{sec:blocks-fluctuation-space},
\begin{equation}\label{eq:operators-photon-2}
\sum_{b'}\exp(-g\,d(b,b'))\le c_{\mathrm{bd}}(g),\qquad \sum_{\Upsilon'}\exp(-g\,d(\Upsilon,\Upsilon'))\le c_{\mathrm{bas}}(g),
\end{equation}
where $b'$ ranges over the positively oriented unit bonds and $\Upsilon'$ over the basis.
\end{lemma}

\begin{proof}
Let $y$ be a unit site. Every site $x'$ of $T^0_k$ has a lift to $\mathbb{Z}^3$ whose sup distance from a fixed lift of $y$ equals $d(x',y)$, and distinct sites have distinct lifts. Hence
\[
\sum_{x'\in T^0_k}\exp(-g\,d(y,x'))\le\boldsymbol{\varsigma}(g).
\]
Let $b$ be a bond, and let $y_b$ be a unit site at distance at most $\frac{1}{2}$ from an end of $b$; it exists, since the sets $B_k(y)$, $y\in T^0_k$, cover $T_k$. For a unit bond $b'$ let $x'(b')$ be its initial point, and let $z$ be an end of $b$ and $z'$ an end of $b'$ with $d(z,z')=d(b,b')$. A fine bond has length $\eta\le 1$ and a unit bond has length $1$; so the two ends of $b$ are at distance at most $1$ from each other, whether $b$ is fine or a unit bond, and the same holds for the two ends of $b'$. Hence $d(y_b,z)\le\frac{1}{2}+1$ and $d(z',x'(b'))\le 1$, and the triangle inequality gives
\[
d(y_b,x'(b'))\le\frac{1}{2}+1+d(b,b')+1\le d(b,b')+3.
\]
Three positively oriented unit bonds have the same initial point. It follows that
\[
\sum_{b'}\exp(-g\,d(b,b'))\le\exp(3g)\sum_{b'}\exp(-g\,d(y_b,x'(b')))\le 3\exp(3g)\cdot\boldsymbol{\varsigma}(g)=c_{\mathrm{bd}}(g).
\]
To each basis element $\Upsilon$ attach a unit site $x(\Upsilon)$: put $x(e_i^y):=y$, the centre of the block $B(y)$, and let $x(\delta_b)$ be the initial point of $b$. By Proposition~\ref{prop:fluctuation-space}(i), (ii), at most $2L^3-3L^2+4\le 2L^3$ basis elements have the same point $x(\Upsilon)$, namely the $2L^3-3L^2+1$ vectors $e_i^y$ of one block and at most three vectors $\delta_b$; moreover $\operatorname{supp}\Upsilon$ lies in one $L$-block or is one bond, so that every point of $\operatorname{supp}\Upsilon$ has distance at most $L$ from $x(\Upsilon)$, and
\[
d(x(\Upsilon),x(\Upsilon'))\le d(\Upsilon,\Upsilon')+2L.
\]
Hence
\[
\sum_{\Upsilon'}\exp(-g\,d(\Upsilon,\Upsilon'))\le\exp(2gL)\cdot 2L^3\sum_{x'\in T^0_k}\exp(-g\,d(x(\Upsilon),x'))\le c_{\mathrm{bas}}(g).
\]
\end{proof}

\subsubsection{Block averages and the axial functional}\label{sec:block-averages-axial}

For $y\in T^0_k$, a scalar $\lambda$, a bond function $\mathscr{A}$ and a plaquette function $F$ on $T_k$, and for a bond function $A$ on $T^0_k$, put
\begin{align*}
(Q_k\lambda)(y)&:=\sum_{x\in B_k(y)}\eta^3\lambda(x),\\
(Q_k\mathscr{A})_{\mu}(y)&:=\sum_{x\in B_k(y)}\eta^3\sum_{j=0}^{\mathfrak{N}-1}\eta\,\mathscr{A}_{\mu}(x+j\eta e_{\mu}),\\
(Q_k^{(2)}F)_{\mu\nu}(y)&:=\sum_{x\in B_k(y)}\eta^3\sum_{j,j'=0}^{\mathfrak{N}-1}\eta^2F_{\mu\nu}(x+j\eta e_{\mu}+j'\eta e_{\nu}),\\
(QA)(y,y+Le_{\mu})&:=L^{-4}\sum_{x\in B(y)}A(\Gamma_{x,x+Le_{\mu}})\qquad(y\in L\mathbb{Z}^3/l\mathbb{Z}^3),\\
(\tau A)(y,x)&:=\frac{1}{3!}\sum_{\pi}A(\Gamma^{\pi}(y,x))\qquad(y\in L\mathbb{Z}^3/l\mathbb{Z}^3,\ x\in B(y)).
\end{align*}
Thus $Q_k$ maps scalars and bond functions on $T_k$ to scalars and bond functions on $T^0_k$, respectively, $Q_k^{(2)}$ maps plaquette functions on $T_k$ to plaquette functions on $T^0_k$, and $QA$ is a bond function on $L\mathbb{Z}^3/l\mathbb{Z}^3$. Here $B(y)$ is the $L$-block of Section~\ref{sec:blocks-fluctuation-space}. If $x'-x=t\,e_{\mu}$ with a direction $\mu$ and a real number $t$, then $\Gamma_{x,x'}$ denotes the path that leads from $x$ to $x'$ through the consecutive bonds of the segment between them. Let $\pi=(\pi_1,\pi_2,\pi_3)$ be a permutation of the directions $0,1,2$. Put $w^{(0)}:=y$ and, for $i=1,2,3$, let $w^{(i)}$ be the site whose coordinate in the direction $\pi_i$ is $x_{\pi_i}$ and whose other two coordinates are those of $w^{(i-1)}$; thus $w^{(3)}=x$. Then $\Gamma^{\pi}(y,x)$ is the path $\Gamma_{w^{(0)},w^{(1)}}$, followed by $\Gamma_{w^{(1)},w^{(2)}}$ and then by $\Gamma_{w^{(2)},w^{(3)}}$; in the definition of $\tau$ the sum extends over the six permutations $\pi$. The averages $Q$, $Q_k$ and $Q_k^{(2)}$ are those of \eqref{eq:algebraic-gaussian-4} and \eqref{eq:algebraic-gaussian-5}. For a linear map $O$, the transpose $O^T$ is its adjoint for the pairings $\langle\cdot,\cdot\rangle$ of Section~\ref{sec:tori-cochains-kernels}. The projection $\mathscr{R}_k$ of the Landau condition of Section~\ref{sec:averaging-fluctuations-gauge}, in the notation of Definition~\ref{def:notation}(f), is the orthogonal projection, in the space of the scalars on $T_k$, onto $\Delta(\ker Q_k)$.

\subsubsection{Blocks and the fluctuation space}\label{sec:blocks-fluctuation-space}

The $L$-blocks of $T^0_k$ are the sets
\[
B(y):=y+\left\{-\frac{L-1}{2},\dots,\frac{L-1}{2}\right\}^3,\qquad y\in L\mathbb{Z}^3/l\mathbb{Z}^3.
\]
Every unit bond is either interior, that is, both its ends lie in one block, or a face bond, that is, its ends lie in $B(y)$ and in $B(y+Le_{\mu})$ for some $y$ and $\mu$. For each face $\langle y,y+Le_{\mu}\rangle$, that is, each pair of neighbouring blocks $B(y)$, $B(y+Le_{\mu})$, there are $L^2$ face bonds. The central bond of a face is the face bond on the straight line through the two centres $y$ and $y+Le_{\mu}$.

The fluctuation space $V_k$, its basis $\{\Upsilon\}$, the map $C$, the operators $\tilde{\Delta}_k$, $C_k$, $C_k^{\frac{1}{2}}$, $\mathrm{Sq}_k$, the map $A^{\min}_k$, the set $\mathrm{Box}_k$, the measures $\mu_I$ and $\hat{\mu}$, and the number $d\varepsilon^0_k$ are those of Section~\ref{sec:averaging-fluctuations-gauge}; $V_k$, $\{\Upsilon\}$ and $C$ are defined also in Section~\ref{sec:fluctuation-space-basis}.

\subsection{Block averages and the minimiser}\label{sec:block-averages-minimiser}

\begin{lemma}[Block averages]\label{lem:block-averages}
Let $l$ be any positive integer. The averages $Q_k$ and $Q_k^{(2)}$ of Section~\ref{sec:block-averages-axial}, which map the scalar functions, the bond functions and the plaquette functions on $T_k$ to the functions of the same kind on $T^0_k$, and the transpose $Q_k^T$ of $Q_k$ have the following properties.
\begin{enumerate}
\item[(a)] $Q_k1 = 1$ on scalar functions, and $Q_kc = c$ for every constant bond function $c$.
\item[(b)] $Q_k\,d\lambda = d\,Q_k\lambda$ for every scalar function $\lambda$ on $T_k$.
\item[(c)] For every bond function $\mathscr{A}$ on $T_k$, the bond function $Q_k\mathscr{A}$ has the same mean as $\mathscr{A}$.
\item[(d)] For every scalar function $\boldsymbol{\omega}$ on $T^0_k$ one has $\left(Q_k^T\boldsymbol{\omega}\right)(x) = \boldsymbol{\omega}(y)$ whenever $x \in B_k(y)$; moreover $Q_kQ_k^T = 1$ on scalar functions.
\item[(e)] $Q_k$ maps the space of the bond functions on $T_k$ onto the space of the bond functions on $T^0_k$.
\item[(f)] $d\,Q_k = Q_k^{(2)}\,d$ on bond functions, and $\left|Q_k^{(2)}F\right|_2 \le |F|_2$ for every plaquette function $F$ on $T_k$.
\end{enumerate}
\end{lemma}

\begin{proof}
Every block $B_k(y)$ has $\#|B_k(y)| = \mathfrak{N}^3$ fine sites, the blocks $B_k(y)$, $y \in T^0_k$, form a partition of $T_k$, and $B_k(y) + e_{\mu} = B_k(y + e_{\mu})$. For a fine site $x$, the straight path from $x$ to $x + e_{\mu}$ consists of the $\mathfrak{N}$ bonds whose initial points are $x + j\eta e_{\mu}$, $0 \le j \le \mathfrak{N} - 1$.

(a) In the definition of $Q_k$ the weights add up to $\mathfrak{N}^3\cdot\eta^3 = 1$ on scalar functions, and to $\mathfrak{N}^3\cdot\eta^3\cdot\mathfrak{N}\cdot\eta = 1$ on bond functions.

(b) Along the path from $x$ to $x + e_{\mu}$ the difference quotients telescope:
\[
\sum_{j=0}^{\mathfrak{N}-1}\eta\,(\partial_{\mu}\lambda)(x + j\eta e_{\mu}) = \lambda(x + e_{\mu}) - \lambda(x).
\]
Multiply by $\eta^3$ and sum over $x \in B_k(y)$; since the translation by $e_{\mu}$ maps $B_k(y)$ onto $B_k(y + e_{\mu})$, it follows that
\begin{align*}
(Q_k\,d\lambda)_{\mu}(y) &= \sum_{x\in B_k(y)}\eta^3\left(\lambda(x + e_{\mu}) - \lambda(x)\right)\\
&= (Q_k\lambda)(y + e_{\mu}) - (Q_k\lambda)(y) = (d\,Q_k\lambda)_{\mu}(y).
\end{align*}

(c) Since the blocks form a partition of $T_k$,
\[
\sum_{y\in T^0_k}(Q_k\mathscr{A})_{\mu}(y) = \sum_{j=0}^{\mathfrak{N}-1}\eta\sum_{x\in T_k}\eta^3\,\mathscr{A}_{\mu}(x + j\eta e_{\mu}) = \sum_{x\in T_k}\eta^3\,\mathscr{A}_{\mu}(x);
\]
in the second equality the sum over $x$ is invariant under the translation by $j\eta e_{\mu}$, and $\mathfrak{N}\cdot\eta = 1$. Multiplication by $l^{-3}$ gives the equality of the two means.

(d) For a scalar function $\lambda$ on $T_k$ the pairing on $T^0_k$ reads
\[
\langle Q_k\lambda, \boldsymbol{\omega}\rangle = \sum_{y\in T^0_k}\boldsymbol{\omega}(y)\sum_{x\in B_k(y)}\eta^3\lambda(x).
\]
The right side is the pairing on $T_k$ of $\lambda$ with the function that is equal to $\boldsymbol{\omega}(y)$ on $B_k(y)$ for every $y$. So $Q_k^T\boldsymbol{\omega}$ is this function, and $\left(Q_kQ_k^T\boldsymbol{\omega}\right)(y) = \boldsymbol{\omega}(y)\cdot(Q_k1)(y) = \boldsymbol{\omega}(y)$ by (a).

(e) Let $S$ be the linear map from the bond functions on $T^0_k$ to the bond functions on $T_k$ given, for a fine site $x \in B_k(y)$, by
\[
(SA)_{\mu}(x) := \begin{cases} \mathfrak{N}\cdot A_{\mu}(y) & \text{if } x_{\mu} = y_{\mu} + \mathbf{h}\cdot\eta,\\ 0 & \text{otherwise}. \end{cases}
\]
The first case is the case in which the bond $(x, x + \eta e_{\mu})$ leaves the cube $y + \eta\{-\mathbf{h}, \dots, \mathbf{h}\}^3$. Let $x = y + \eta t$ with $t \in \{-\mathbf{h}, \dots, \mathbf{h}\}^3$, and consider the path from $x$ to $x + e_{\mu}$. Its bond with $j = \mathbf{h} - t_{\mu}$ begins at a point of $B_k(y)$ whose $\mu$-th coordinate is $y_{\mu} + \mathbf{h}\cdot\eta$. For $j < \mathbf{h} - t_{\mu}$ the initial point is $y + \eta(t + je_{\mu}) \in B_k(y)$, and for $j > \mathbf{h} - t_{\mu}$ it is $y + e_{\mu} + \eta\left(t + (j - \mathfrak{N})e_{\mu}\right) \in B_k(y + e_{\mu})$; in both cases the $\mu$-th component of its offset from the centre of its block, namely $t_{\mu} + j$, resp.\ $t_{\mu} + j - \mathfrak{N}$, lies in $\{-\mathbf{h}, \dots, \mathbf{h} - 1\}$. So among the bonds of the path exactly one is of the first kind, and its initial point lies in $B_k(y)$. Consequently
\[
(Q_kSA)_{\mu}(y) = \sum_{x\in B_k(y)}\eta^3\cdot\eta\cdot\mathfrak{N}\cdot A_{\mu}(y) = A_{\mu}(y),
\]
that is, $Q_kS = 1$, and $Q_k$ is onto.

(f) As in (b), for every fine site $z$,
\[
\mathscr{A}_{\nu}(z + e_{\mu}) - \mathscr{A}_{\nu}(z) = \sum_{j=0}^{\mathfrak{N}-1}\eta\,(\partial_{\mu}\mathscr{A}_{\nu})(z + j\eta e_{\mu}).
\]
On $T^0_k$, where $\eta$ is replaced by $1$, the definition of $d$ gives
\[
(d\,Q_k\mathscr{A})_{\mu\nu}(y) = (Q_k\mathscr{A})_{\nu}(y + e_{\mu}) - (Q_k\mathscr{A})_{\nu}(y) - (Q_k\mathscr{A})_{\mu}(y + e_{\nu}) + (Q_k\mathscr{A})_{\mu}(y).
\]
By the definition of $Q_k$, the relation $B_k(y + e_{\mu}) = B_k(y) + e_{\mu}$ and the preceding identity at $z = x + j'\eta e_{\nu}$,
\begin{align*}
(Q_k\mathscr{A})_{\nu}(y + e_{\mu}) - (Q_k\mathscr{A})_{\nu}(y) &= \sum_{x\in B_k(y)}\eta^3\sum_{j'=0}^{\mathfrak{N}-1}\eta\left(\mathscr{A}_{\nu}(x + j'\eta e_{\nu} + e_{\mu}) - \mathscr{A}_{\nu}(x + j'\eta e_{\nu})\right)\\
&= \sum_{x\in B_k(y)}\eta^3\sum_{j,j'=0}^{\mathfrak{N}-1}\eta^2\,(\partial_{\mu}\mathscr{A}_{\nu})(x + j\eta e_{\mu} + j'\eta e_{\nu}),
\end{align*}
and in the same way $(Q_k\mathscr{A})_{\mu}(y + e_{\nu}) - (Q_k\mathscr{A})_{\mu}(y)$ is the same double sum with $\partial_{\nu}\mathscr{A}_{\mu}$ in place of $\partial_{\mu}\mathscr{A}_{\nu}$. The difference of the two double sums is $\left(Q_k^{(2)}d\mathscr{A}\right)_{\mu\nu}(y)$. For the inequality, fix $y$ and $\mu < \nu$. The number $\left(Q_k^{(2)}F\right)_{\mu\nu}(y)$ is a weighted mean of values of $F_{\mu\nu}$: all weights equal $\eta^5$, and their total is $\mathfrak{N}^3\cdot\eta^3\cdot\mathfrak{N}^2\cdot\eta^2 = 1$. Jensen's inequality, applied to the convex function $|\cdot|^2$, gives
\[
\left|\left(Q_k^{(2)}F\right)_{\mu\nu}(y)\right|^2 \le \sum_{x\in B_k(y)}\ \sum_{j,j'=0}^{\mathfrak{N}-1}\eta^5\left|F_{\mu\nu}(x + j\eta e_{\mu} + j'\eta e_{\nu})\right|^2.
\]
In the sum of the right side over $y \in T^0_k$, the value $F_{\mu\nu}(z)$ at a fine site $z$ occurs once for each of the $\mathfrak{N}^2$ pairs $(j, j')$, namely with $x = z - j\eta e_{\mu} - j'\eta e_{\nu}$, so with the total weight $\mathfrak{N}^2\cdot\eta^5 = \eta^3$. Summation over $y$ and over $\mu < \nu$ gives $\left|Q_k^{(2)}F\right|_2^2 \le |F|_2^2$.

Parts (a), (b), (c) and (f) are contained in Proposition~\ref{prop:averaging-operators}(b), (c), (d), taken at level $k$; they comprise the clauses of Proposition~\ref{prop:block-averages} on $Q_k$ and $Q_k^{(2)}$, and (f) gives the inequality there with $C_Q = 1$.
\end{proof}

\begin{lemma}[Characterisation of the minimiser]\label{lem:characterisation-minimiser}
Let $A$ be a real bond function on $T^0_k$, and consider the problem of Section~\ref{sec:averaging-fluctuations-gauge}: to minimise $|d\mathscr{A}|_2^2$ over the real bond functions $\mathscr{A}$ on $T_k$ subject to the two constraints
\[
Q_k\mathscr{A} = A, \qquad \mathscr{R}_k\,\delta\mathscr{A} = 0.
\]
A real bond function $\mathscr{A}$ on $T_k$ is a solution of this problem if and only if
\begin{align}
&Q_k\mathscr{A} = A, \tag{M1}\label{eq:operators-photon-M1}\\
&\delta d\mathscr{A} \ \text{lies in the range of } Q_k^T \text{ on bond functions}, \tag{M2}\label{eq:operators-photon-M2}\\
&\Delta\delta\mathscr{A} \ \text{lies in the range of } Q_k^T \text{ on scalar functions}. \tag{M3}\label{eq:operators-photon-M3}
\end{align}
The problem has exactly one solution.
\end{lemma}

\begin{proof}
On the scalar functions and on the bond functions on $T_k$ alike, the orthogonal complement of $\ker Q_k$ is the range of $Q_k^T$.

\emph{Step 1.} By Lemma~\ref{lem:block-averages}(e) the real bond functions $\mathscr{A}$ on $T_k$ with $Q_k\mathscr{A} = A$ form a non-empty affine space, namely $\mathscr{A}_0 + \ker Q_k$ for any one of them, $\mathscr{A}_0$. The map $\mathscr{A} \mapsto |d\mathscr{A}|_2^2$ is a non-negative quadratic form, and it attains its minimum on this affine space (as $\mathscr{A}$ ranges over it, $d\mathscr{A}$ ranges over the affine subspace $d\mathscr{A}_0 + d(\ker Q_k)$ of a finite-dimensional space, on which the square of the norm attains its minimum). For $B \in \ker Q_k$ and real $t$ one has
\[
\left|d(\mathscr{A} + tB)\right|_2^2 = |d\mathscr{A}|_2^2 + 2t\,\langle d\mathscr{A}, dB\rangle + t^2\,|dB|_2^2.
\]
Therefore $\mathscr{A}$ is a minimiser if and only if $\langle d\mathscr{A}, dB\rangle = \langle \delta d\mathscr{A}, B\rangle = 0$ for every $B \in \ker Q_k$, that is, if and only if $\delta d\mathscr{A}$ is orthogonal to $\ker Q_k$; this is \eqref{eq:operators-photon-M2}. If $\mathscr{A}$ and $\mathscr{A}'$ are two minimisers, then $\frac{1}{2}(\mathscr{A} + \mathscr{A}')$ lies in the affine space, and
\[
\left|\frac{1}{2}\,d(\mathscr{A} + \mathscr{A}')\right|_2^2 = \frac{1}{2}\,|d\mathscr{A}|_2^2 + \frac{1}{2}\,|d\mathscr{A}'|_2^2 - \frac{1}{4}\left|d(\mathscr{A} - \mathscr{A}')\right|_2^2
\]
forces $d\mathscr{A} = d\mathscr{A}'$; thus all minimisers have one and the same curvature.

\emph{Step 2.} Let $B$ be a real bond function on $T_k$ with $dB = 0$ and $Q_kB = 0$. By the Hodge decomposition on the discrete torus (Section~\ref{sec:tori-cochains-kernels}), $B = c + d\lambda$ with a constant $c$ and a real scalar function $\lambda$ on $T_k$. Lemma~\ref{lem:block-averages}(a), (b) gives
\[
0 = Q_kB = c + d\,Q_k\lambda.
\]
An exact bond function has mean zero (Section~\ref{sec:tori-cochains-kernels}); hence the mean of the right side is $c$, and $c = 0$. Then $d\,Q_k\lambda = 0$, so $Q_k\lambda$ is a constant $c_0$; by Lemma~\ref{lem:block-averages}(a), $\lambda - c_0 \in \ker Q_k$, and $B = d(\lambda - c_0) \in d(\ker Q_k)$. Conversely, for $\lambda \in \ker Q_k$ the field $d\lambda$ is closed, and $Q_kd\lambda = d\,Q_k\lambda = 0$ by Lemma~\ref{lem:block-averages}(b). Let $\mathscr{A}_1$ be one of the minimisers of Step 1. By Step 1 the minimisers are the fields $\mathscr{A}_1 + B$ with $B \in \ker Q_k$ and $dB = 0$; so they form the set $\mathscr{A}_1 + d(\ker Q_k)$.

\emph{Step 3.} By definition (Section~\ref{sec:block-averages-axial}), $\mathscr{R}_k$ is the orthogonal projection, in the scalar functions on $T_k$, onto $\Delta(\ker Q_k)$. For a scalar function $f$ on $T_k$, therefore, $\mathscr{R}_kf = 0$ holds if and only if $\langle f, \Delta\lambda\rangle = 0$ for every $\lambda \in \ker Q_k$. Since $\Delta$ is symmetric, this means that $\Delta f$ is orthogonal to $\ker Q_k$, that is, that $\Delta f$ lies in the range of $Q_k^T$. With $f = \delta\mathscr{A}$: the constraint $\mathscr{R}_k\,\delta\mathscr{A} = 0$ is equivalent to \eqref{eq:operators-photon-M3}.

\emph{Step 4.} Let $\lambda \in \ker Q_k$. Since $\delta d\lambda = \Delta\lambda$ and $\mathscr{R}_k\Delta\lambda = \Delta\lambda$,
\[
\mathscr{R}_k\,\delta(\mathscr{A}_1 + d\lambda) = \mathscr{R}_k\,\delta\mathscr{A}_1 + \Delta\lambda.
\]
The restriction of $\Delta$ to $\ker Q_k$ is injective: if $\Delta\lambda = 0$, then $|d\lambda|_2^2 = \langle\lambda, \Delta\lambda\rangle = 0$, so $\lambda$ is a constant, and this constant equals $Q_k\lambda = 0$ by Lemma~\ref{lem:block-averages}(a). As $-\mathscr{R}_k\,\delta\mathscr{A}_1$ lies in $\Delta(\ker Q_k)$, there is exactly one $\lambda \in \ker Q_k$ with $\Delta\lambda = -\mathscr{R}_k\,\delta\mathscr{A}_1$.

\emph{Step 5.} The argument of Step 4 does not use that $\mathscr{A}_1$ is a minimiser. Hence every real bond function $\mathscr{A}$ with $Q_k\mathscr{A} = A$ becomes, after the addition of $d\lambda$ for a suitable $\lambda \in \ker Q_k$, a field that satisfies both constraints (by Lemma~\ref{lem:block-averages}(b), $Q_kd\lambda = d\,Q_k\lambda = 0$) and has the same curvature. It follows that the minimum of $|d\mathscr{A}|_2^2$ under both constraints equals its minimum under the constraint $Q_k\mathscr{A} = A$ alone, and that the solutions of the problem of Section~\ref{sec:averaging-fluctuations-gauge} are those minimisers of Step 1 that satisfy $\mathscr{R}_k\,\delta\mathscr{A} = 0$, that is, by Step 3, \eqref{eq:operators-photon-M3}. Together with Step 1 this proves that the solutions are the fields with \eqref{eq:operators-photon-M1}, \eqref{eq:operators-photon-M2}, \eqref{eq:operators-photon-M3}. By Step 2 the minimisers of Step 1 are the fields $\mathscr{A}_1 + d\lambda$, $\lambda \in \ker Q_k$, and by Step 4 exactly one of them satisfies $\mathscr{R}_k\,\delta\mathscr{A} = 0$. So there is exactly one solution.
\end{proof}

\begin{proposition}[The minimiser $H_k$]\label{prop:minimiser}
For every real bond function $A$ on $T^0_k$ the problem of Section~\ref{sec:averaging-fluctuations-gauge} has exactly one solution, denoted $H_kA$. The map $H_k$ is linear and real, and $Q_kH_k = 1$. For every constant bond function $c$ one has $H_kc = c$, and $dA = 0$ implies $dH_kA = 0$. For $k = 0$ the solution is $A$, that is, $H_0 = 1$.
\end{proposition}

\begin{proof}
Lemma~\ref{lem:characterisation-minimiser} gives the existence and the uniqueness of the solution. Each of \eqref{eq:operators-photon-M1}, \eqref{eq:operators-photon-M2}, \eqref{eq:operators-photon-M3} is a linear condition on the pair $(\mathscr{A}, A)$; together with the uniqueness of the solution this shows that $H_k$ is linear. The problem is posed for real fields, so $H_k$ maps real fields to real fields. Condition \eqref{eq:operators-photon-M1} for $\mathscr{A} = H_kA$ is the identity $Q_kH_k = 1$. A constant field $c$ on $T_k$ satisfies \eqref{eq:operators-photon-M1} with $A = c$ by Lemma~\ref{lem:block-averages}(a), and $dc = 0$, $\delta c = 0$, so that \eqref{eq:operators-photon-M2} and \eqref{eq:operators-photon-M3} hold as well; hence $H_kc = c$.

Now let $dA = 0$. By the Hodge decomposition (Section~\ref{sec:tori-cochains-kernels}) on $T^0_k$, $A = c + d\lambda$ with a constant $c$ and a real scalar function $\lambda$ on $T^0_k$. The bond function $c + d\,Q_k^T\lambda$ on $T_k$ is closed, and by Lemma~\ref{lem:block-averages}(a), (b), (d)
\[
Q_k\left(c + d\,Q_k^T\lambda\right) = c + d\,Q_kQ_k^T\lambda = c + d\lambda = A.
\]
So the minimum of $|d\mathscr{A}|_2^2$ under the constraint $Q_k\mathscr{A} = A$ is zero. By Step 5 of the proof of Lemma~\ref{lem:characterisation-minimiser} it is also the minimum of the problem of Section~\ref{sec:averaging-fluctuations-gauge}, which is attained at $H_kA$; hence $dH_kA = 0$.

For $k = 0$ one has $\mathfrak{N} = 1$, hence $Q_0 = 1$, and \eqref{eq:operators-photon-M1} gives $H_0A = A$. Proposition~\ref{prop:background-field} follows.
\end{proof}

\subsection{The Fourier symbol of $H_k$}\label{sec:fourier-symbol}

Throughout this section and Section~\ref{sec:bounds-complex-strip}, the letter $p$ denotes a momentum, and not the exponent $p$ of stage~\ref{it:order-choice-8}; the vectors $q = p + 2\pi m$ with $m \in \mathbb{Z}^3$ are called the aliases of $p$.

\subsubsection{Elementary symbols and Fourier transforms}\label{sec:elementary-symbols-fourier}

For $s \in \mathbb{C}$ let
\[
\operatorname{sn}(s) := 2\mathfrak{N} \cdot \sin\left(\frac{s}{2\mathfrak{N}}\right), \qquad \operatorname{cs}(s) := 2\sin\left(\frac{s}{2}\right), \qquad u(s) := \frac{1}{\mathfrak{N}}\sum_{j=-\mathbf{h}}^{\mathbf{h}}\exp\left(\frac{isj}{\mathfrak{N}}\right).
\]
The function $u$ is entire, $u(0) = 1$, and
\[
\operatorname{cs}(s) = u(s) \cdot \operatorname{sn}(s) \qquad (s \in \mathbb{C}).
\]
To see the last identity, put $w := \exp\left(\frac{is}{2\mathfrak{N}}\right)$. Since $2\mathbf{h} + 1 = \mathfrak{N}$, the sum telescopes:
\[
\left(w - w^{-1}\right)\sum_{j=-\mathbf{h}}^{\mathbf{h}} w^{2j} = \sum_{j=-\mathbf{h}}^{\mathbf{h}}\left(w^{2j+1} - w^{2j-1}\right) = w^{\mathfrak{N}} - w^{-\mathfrak{N}},
\]
and here $\sum_j w^{2j} = \mathfrak{N} \cdot u(s)$, $\mathfrak{N} \cdot \left(w - w^{-1}\right) = i \cdot \operatorname{sn}(s)$ and $w^{\mathfrak{N}} - w^{-\mathfrak{N}} = i \cdot \operatorname{cs}(s)$; so the telescoped identity reads $i \cdot \operatorname{sn}(s) \cdot u(s) = i \cdot \operatorname{cs}(s)$.

For $q \in \mathbb{C}^3$, $p \in \mathbb{C}^3$ and $\nu \in \{0,1,2\}$ let
\begin{align*}
\hat{d}_{\nu}(q) &:= \eta^{-1}\left(\exp(i\eta q_{\nu}) - 1\right) = i\exp\left(\frac{i\eta q_{\nu}}{2}\right) \cdot \operatorname{sn}(q_{\nu}),\\
\hat{d}^{*}_{\nu}(q) &:= \eta^{-1}\left(\exp(-i\eta q_{\nu}) - 1\right) = -i\exp\left(-\frac{i\eta q_{\nu}}{2}\right) \cdot \operatorname{sn}(q_{\nu}),\\
\hat{D}(q) &:= \sum_{\nu}\operatorname{sn}(q_{\nu})^2 = \sum_{\nu}\hat{d}_{\nu}(q) \cdot \hat{d}^{*}_{\nu}(q),\\
\hat{U}(q) &:= u(q_0) \cdot u(q_1) \cdot u(q_2),\\
\hat{V}_{\nu}(q) &:= \eta\sum_{j=0}^{\mathfrak{N}-1}\exp(i\eta jq_{\nu}) = \exp\left(\frac{i(1-\eta)q_{\nu}}{2}\right) \cdot u(q_{\nu}),\\
\hat{V}^{*}_{\nu}(q) &:= \exp\left(-\frac{i(1-\eta)q_{\nu}}{2}\right) \cdot u(q_{\nu}),\\
\rho_{\nu}(p) &:= \exp(ip_{\nu}) - 1, \qquad \rho^{*}_{\nu}(p) := \exp(-ip_{\nu}) - 1.
\end{align*}
The second expressions for $\hat{d}_{\nu}$ and $\hat{d}^{*}_{\nu}$ follow from $\exp(ia) - 1 = 2i\exp\left(\frac{ia}{2}\right)\sin\left(\frac{a}{2}\right)$ and $\eta = \mathfrak{N}^{-1}$; the second expression for $\hat{V}_{\nu}$ follows from the substitution $j = j' + \mathbf{h}$, since $\eta \cdot \mathbf{h} = \frac{1-\eta}{2}$. All these functions are entire. The functions $\operatorname{sn}^2$, $u$, $\hat{d}_{\nu}$, $\hat{d}^{*}_{\nu}$, $\hat{D}$, $\hat{U}$, $\hat{V}_{\nu}$, $\hat{V}^{*}_{\nu}$ are periodic with period $2\pi\mathfrak{N}$ in each variable ($\mathfrak{N}$ is odd). For real arguments, $\hat{d}^{*}_{\nu}$, $\hat{V}^{*}_{\nu}$ and $\rho^{*}_{\nu}$ are the complex conjugates of $\hat{d}_{\nu}$, $\hat{V}_{\nu}$ and $\rho_{\nu}$; moreover $\hat{d}^{*}_{\nu}(q) = \hat{d}_{\nu}(-q)$ and $\hat{V}^{*}_{\nu}(q) = \hat{V}_{\nu}(-q)$ for all $q \in \mathbb{C}^3$, as $u$ is even.

Let $p \in \mathbb{C}^3$ and $q \in p + 2\pi\mathbb{Z}^3$. Then, for every $\nu$,
\begin{align}
\hat{V}_{\nu}(q) \cdot \hat{d}_{\nu}(q) &= \rho_{\nu}(p), \qquad \hat{V}^{*}_{\nu}(q) \cdot \hat{d}^{*}_{\nu}(q) = \rho^{*}_{\nu}(p), \notag\\
\hat{V}_{\nu}(q) \cdot \hat{V}^{*}_{\nu}(q) &= u(q_{\nu})^2, \qquad \rho_{\nu}(p) \cdot \rho^{*}_{\nu}(p) = \operatorname{cs}(p_{\nu})^2 = u(p_{\nu})^2 \cdot \operatorname{sn}(p_{\nu})^2. \label{eq:operators-photon-3}
\end{align}
Indeed, the sum
\[
\hat{V}_{\nu}(q) \cdot \hat{d}_{\nu}(q) = \sum_{j=0}^{\mathfrak{N}-1}\left(\exp\left(i\eta(j+1)q_{\nu}\right) - \exp\left(i\eta jq_{\nu}\right)\right) = \exp(iq_{\nu}) - 1
\]
telescopes, and $\exp(iq_{\nu}) = \exp(ip_{\nu})$. The second identity is the first one applied to $-q \in -p + 2\pi\mathbb{Z}^3$, by the relations $\hat{d}^{*}_{\nu}(q) = \hat{d}_{\nu}(-q)$, $\hat{V}^{*}_{\nu}(q) = \hat{V}_{\nu}(-q)$ and $\rho_{\nu}(-p) = \rho^{*}_{\nu}(p)$. The third identity is read off from the definitions, and the fourth follows from $\left(\exp(ia) - 1\right)\left(\exp(-ia) - 1\right) = 4\sin^2\left(\frac{a}{2}\right)$ and $\operatorname{cs} = u \cdot \operatorname{sn}$.

\paragraph{Fourier transforms.} For a function $f$ on $T_k$ and $q \in \frac{2\pi}{l}\mathbb{Z}^3$ let
\[
\hat{f}(q) := \sum_{x \in T_k}\eta^3\exp(-iq\cdot x) \cdot f(x),
\]
where $q \cdot x = \sum_{\nu} q_{\nu}x_{\nu}$ is formed with any representative of $x$ in $\eta\mathbb{Z}^3$. The value depends only on the class of $q$ modulo $2\pi\mathfrak{N}\mathbb{Z}^3$; these classes are called the frequencies of $T_k$. For a function $A$ on $T^0_k$ and $p \in \frac{2\pi}{l}\mathbb{Z}^3$ let
\[
\hat{A}(p) := \sum_{y \in T^0_k}\exp(-ip\cdot y) \cdot A(y),
\]
which depends only on the class of $p$ in
\[
P_l := \frac{2\pi}{l}\mathbb{Z}^3 \Big/ 2\pi\mathbb{Z}^3;
\]
the elements of $P_l$ are represented in $[-\pi,\pi)^3$. Functions are recovered from their transforms by
\[
f(x) = l^{-3}\sum_q\exp(iq\cdot x) \cdot \hat{f}(q), \qquad A(y) = l^{-3}\sum_{p \in P_l}\exp(ip\cdot y) \cdot \hat{A}(p),
\]
where $q$ ranges over the frequencies of $T_k$. Every frequency of $T_k$ contains exactly one vector of the form
\[
q_m := p + 2\pi m, \qquad p \in \frac{2\pi}{l}\mathbb{Z}^3 \cap [-\pi,\pi)^3, \quad m \in \mathfrak{M} := \{-\mathbf{h},\dots,\mathbf{h}\}^3.
\]
Bond functions are transformed componentwise, each bond being indexed by its initial point (Section~\ref{sec:tori-cochains-kernels}): for a bond function $\mathscr{A}$ on $T_k$, $\hat{\mathscr{A}}_{\mu}$ is the transform of the function $x \mapsto \mathscr{A}_{\mu}(x)$ on $T_k$, and the same convention is used on $T^0_k$.

Let $f$, $\lambda$ be functions and $\mathscr{A}$ a bond function on $T_k$, and let $\boldsymbol{\omega}$ be a function and $\boldsymbol{\Lambda}$ a bond function on $T^0_k$. The definitions of Sections~\ref{sec:tori-cochains-kernels} and~\ref{sec:block-averages-axial} give, for every frequency $q$ of $T_k$, every $p \in P_l$ and every $m \in \mathfrak{M}$,
\begin{align}
\widehat{\partial_{\nu}f}(q) &= \hat{d}_{\nu}(q) \cdot \hat{f}(q), \qquad \widehat{\partial^T_{\nu}f}(q) = \hat{d}^{*}_{\nu}(q) \cdot \hat{f}(q), \qquad \widehat{\Delta f}(q) = \hat{D}(q) \cdot \hat{f}(q), \notag\\
\widehat{Q_k\lambda}(p) &= \sum_{m \in \mathfrak{M}}\hat{U}(q_m) \cdot \hat{\lambda}(q_m), \notag\\
\widehat{(Q_k\mathscr{A})_{\mu}}(p) &= \sum_{m \in \mathfrak{M}}\hat{U}(q_m) \cdot \hat{V}_{\mu}(q_m) \cdot \hat{\mathscr{A}}_{\mu}(q_m), \notag\\
\widehat{Q_k^T\boldsymbol{\omega}}(q_m) &= \hat{U}(q_m) \cdot \hat{\boldsymbol{\omega}}(p), \notag\\
\widehat{(Q_k^T\boldsymbol{\Lambda})_{\mu}}(q_m) &= \hat{U}(q_m) \cdot \hat{V}^{*}_{\mu}(q_m) \cdot \hat{\boldsymbol{\Lambda}}_{\mu}(p). \label{eq:operators-photon-4}
\end{align}
The first line follows from the definitions of $\partial_{\nu}$, $\partial^T_{\nu}$ and $\Delta$ by a shift of the variable of summation. For the formulas for $Q_k$, insert $f(x) = l^{-3}\sum_q\exp(iq\cdot x)\cdot\hat{f}(q)$ into the definitions of $Q_k$ and use
\[
\sum_{t \in \eta\{-\mathbf{h},\dots,\mathbf{h}\}^3}\eta^3\exp(iq\cdot t) = \hat{U}(q), \qquad \eta\sum_{j=0}^{\mathfrak{N}-1}\exp(i\eta jq_{\mu}) = \hat{V}_{\mu}(q),
\]
together with the fact that $\sum_{y \in T^0_k}\exp(-ip\cdot y)\exp(iq\cdot y)$ equals $l^3$ if $q \in p + 2\pi\mathbb{Z}^3$ and vanishes otherwise. The formulas for $Q_k^T$ are the transposes of those for $Q_k$: they follow by Parseval's formula on $T_k$ and on $T^0_k$, because $\hat{U}$ is real and even on real arguments and $\hat{V}^{*}_{\mu}(q)$ is the complex conjugate of $\hat{V}_{\mu}(q)$ for real $q$.

\subsubsection{Sums over the aliases}\label{sec:sums-aliases}

For $p \in \mathbb{C}^3$ and $m \in \mathfrak{M}$ let $q_m := p + 2\pi m$, with components $q_{m,\mu}$, and write
\[
D_0 := \hat{D}(p), \qquad U_0 := \hat{U}(p), \qquad u_{\nu} := u(p_{\nu}), \qquad s_{\nu} := \operatorname{sn}(p_{\nu}).
\]
Let
\begin{align}
X_{\mu}(p) &:= \sum_{m \in \mathfrak{M}\setminus\{0\}}\frac{\hat{U}(q_m)^2 \cdot u(q_{m,\mu})^2}{\hat{D}(q_m)}, \notag\\
Y(p) &:= \sum_{m \in \mathfrak{M}\setminus\{0\}}\frac{\hat{U}(q_m)^2}{\hat{D}(q_m)^2}, \notag\\
N_{\mu}(p) &:= U_0^2u_{\mu}^2 + D_0X_{\mu}(p), \qquad N_G(p) := U_0^2 + D_0^2Y(p), \notag\\
\Pi(p) &:= N_0N_1N_2 - \sum_{\mathfrak{l}}s_{\mathfrak{l}}^2X_{\mathfrak{l}}\prod_{\mathfrak{l}' \ne \mathfrak{l}}N_{\mathfrak{l}'}, \label{eq:operators-photon-5}
\end{align}
where in the last line $N_{\mathfrak{l}} = N_{\mathfrak{l}}(p)$ and $X_{\mathfrak{l}} = X_{\mathfrak{l}}(p)$. These functions are defined and analytic on the open set of all $p \in \mathbb{C}^3$ with $\hat{D}(q_m) \ne 0$ for every $m \in \mathfrak{M}\setminus\{0\}$. At the points $p$ of this set at which, in addition, $N_0N_1N_2N_G\Pi \ne 0$, put $\mathcal{N} := N_0N_1N_2$, $\rho_{\mu} := \rho_{\mu}(p)$, $\rho^{*}_{\nu} := \rho^{*}_{\nu}(p)$, and
\begin{align}
\boldsymbol{\Lambda}_{\mu\nu}(p) &:= \delta_{\mu\nu}\frac{D_0}{N_{\mu}} - \frac{\rho_{\mu}\rho^{*}_{\nu}U_0^2 \cdot \mathcal{N}}{N_{\mu}N_{\nu}\Pi}, \notag\\
\boldsymbol{\omega}_{\nu}(p) &:= \frac{\rho^{*}_{\nu}D_0U_0^2 \cdot \mathcal{N}}{N_{\nu}N_G\Pi}, \notag\\
h_{\mu\nu}(p,m) &:= \frac{\hat{U}(q_m) \cdot \hat{V}^{*}_{\mu}(q_m) \cdot \boldsymbol{\Lambda}_{\mu\nu}(p)}{\hat{D}(q_m)} + \frac{\hat{d}_{\mu}(q_m) \cdot \hat{U}(q_m) \cdot \boldsymbol{\omega}_{\nu}(p)}{\hat{D}(q_m)^2} \qquad (m \ne 0), \notag\\
h_{\mu\nu}(p,0) &:= \delta_{\mu\nu}\frac{U_0 \cdot \hat{V}^{*}_{\mu}(p)}{N_{\mu}} + \frac{\hat{d}_{\mu}(p) \cdot \rho^{*}_{\nu} \cdot U_0^3\left(X_{\mu} - u_{\mu}^2D_0Y\right) \cdot \mathcal{N}}{N_{\mu}N_{\nu}N_G\Pi}. \label{eq:operators-photon-6}
\end{align}

\begin{lemma}[Algebraic identities]\label{lem:algebraic-identities}
Let $p \in \mathbb{C}^3$ be such that $\hat{D}(q_m) \ne 0$ for every $m \in \mathfrak{M}\setminus\{0\}$ and $N_0N_1N_2N_G\Pi \ne 0$ at $p$, so that the functions \eqref{eq:operators-photon-6} are defined at $p$. All functions below are taken at $p$.
\begin{enumerate}
\item[(a)] One has
\[
U_0^2\sum_{\mathfrak{l}}\frac{\rho_{\mathfrak{l}}\rho^{*}_{\mathfrak{l}}}{N_{\mathfrak{l}}} = \frac{D_0\Pi}{\mathcal{N}}.
\]
\item[(b)] For every $\nu$, $\sum_{\mu}\rho^{*}_{\mu}\boldsymbol{\Lambda}_{\mu\nu} = 0$.
\item[(c)] If $D_0 \ne 0$, then
\[
h_{\mu\nu}(p,0) = \frac{U_0 \cdot \hat{V}^{*}_{\mu}(p) \cdot \boldsymbol{\Lambda}_{\mu\nu}(p)}{D_0} + \frac{\hat{d}_{\mu}(p) \cdot U_0 \cdot \boldsymbol{\omega}_{\nu}(p)}{D_0^2};
\]
since $q_0 = p$, this says that the formula \eqref{eq:operators-photon-6} for $h_{\mu\nu}(p,m)$, $m \ne 0$, holds for $m = 0$ as well.
\item[(d)] If $D_0 \ne 0$, then
\[
\sum_{m \in \mathfrak{M}}\frac{\hat{U}(q_m)^2 \cdot u(q_{m,\mu})^2}{\hat{D}(q_m)} = \frac{N_{\mu}}{D_0}, \qquad \sum_{m \in \mathfrak{M}}\frac{\hat{U}(q_m)^2}{\hat{D}(q_m)^2} = \frac{N_G}{D_0^2}.
\]
\end{enumerate}
\end{lemma}

\begin{proof}
(d) The terms with $m = 0$ of the two sums are $\frac{U_0^2u_{\mu}^2}{D_0}$ and $\frac{U_0^2}{D_0^2}$, and the remaining terms add up to $X_{\mu}$ and to $Y$. Hence both identities are the definitions of $N_{\mu}$ and $N_G$ in \eqref{eq:operators-photon-5}.

(a) By the last identity of \eqref{eq:operators-photon-3} and the definition of $N_{\mathfrak{l}}$ in \eqref{eq:operators-photon-5},
\[
\frac{U_0^2\rho_{\mathfrak{l}}\rho^{*}_{\mathfrak{l}}}{N_{\mathfrak{l}}} = \frac{s_{\mathfrak{l}}^2U_0^2u_{\mathfrak{l}}^2}{N_{\mathfrak{l}}} = s_{\mathfrak{l}}^2 \cdot \frac{N_{\mathfrak{l}} - D_0X_{\mathfrak{l}}}{N_{\mathfrak{l}}}.
\]
Summation over $\mathfrak{l}$, with $\sum_{\mathfrak{l}}s_{\mathfrak{l}}^2 = D_0$, gives
\[
U_0^2\sum_{\mathfrak{l}}\frac{\rho_{\mathfrak{l}}\rho^{*}_{\mathfrak{l}}}{N_{\mathfrak{l}}} = D_0\left(1 - \sum_{\mathfrak{l}}\frac{s_{\mathfrak{l}}^2X_{\mathfrak{l}}}{N_{\mathfrak{l}}}\right) = \frac{D_0\Pi}{\mathcal{N}},
\]
where the last equality is the definition of $\Pi$ in \eqref{eq:operators-photon-5}, divided by $\mathcal{N}$.

(b) By \eqref{eq:operators-photon-6},
\[
\sum_{\mu}\rho^{*}_{\mu}\boldsymbol{\Lambda}_{\mu\nu} = \frac{\rho^{*}_{\nu}D_0}{N_{\nu}} - \left(U_0^2\sum_{\mu}\frac{\rho_{\mu}\rho^{*}_{\mu}}{N_{\mu}}\right)\frac{\rho^{*}_{\nu} \cdot \mathcal{N}}{N_{\nu}\Pi},
\]
and by (a) the second term on the right equals the first.

(c) By \eqref{eq:operators-photon-3} with $q = p$, one has $\hat{V}^{*}_{\mu}(p) \cdot \rho_{\mu}(p) = \hat{V}^{*}_{\mu}(p) \cdot \hat{V}_{\mu}(p) \cdot \hat{d}_{\mu}(p) = u_{\mu}^2 \cdot \hat{d}_{\mu}(p)$. Insert the definitions \eqref{eq:operators-photon-6} of $\boldsymbol{\Lambda}_{\mu\nu}$ and $\boldsymbol{\omega}_{\nu}$ into the right-hand side of (c); with this relation it becomes
\[
\delta_{\mu\nu}\frac{U_0 \cdot \hat{V}^{*}_{\mu}(p)}{N_{\mu}} + \hat{d}_{\mu}(p) \cdot \rho^{*}_{\nu}U_0^3 \cdot \frac{\mathcal{N}}{N_{\nu}\Pi} \cdot \frac{1}{D_0}\left(\frac{1}{N_G} - \frac{u_{\mu}^2}{N_{\mu}}\right).
\]
By \eqref{eq:operators-photon-5},
\[
N_{\mu} - u_{\mu}^2N_G = D_0\left(X_{\mu} - u_{\mu}^2D_0Y\right),
\]
so that
\[
\frac{1}{D_0}\left(\frac{1}{N_G} - \frac{u_{\mu}^2}{N_{\mu}}\right) = \frac{X_{\mu} - u_{\mu}^2D_0Y}{N_{\mu}N_G},
\]
and the expression above is $h_{\mu\nu}(p,0)$ as defined in \eqref{eq:operators-photon-6}.
\end{proof}

\subsubsection{The symbols of $H_k$ and of $\Delta_k$}\label{sec:symbols}

\begin{proposition}[Symbols of $H_k$ and $\Delta_k$]\label{prop:symbols}
Let $l$ be a positive integer and let $p \in P_l$, represented in $[-\pi,\pi)^3$. Then $\hat{D}(q_m) > 0$ for every $m \in \mathfrak{M}\setminus\{0\}$, and
\[
N_{\mu}(p) > 0, \qquad N_G(p) > 0, \qquad \Pi(p) > 0;
\]
in particular the functions \eqref{eq:operators-photon-6} are defined at $p$. Let $\Delta_k := (dH_k)^TdH_k$. For every real bond function $A$ on $T^0_k$,
\begin{equation}\label{eq:operators-photon-7}
\widehat{(H_kA)_{\mu}}(q_m) = \sum_{\nu}h_{\mu\nu}(p,m) \cdot \hat{A}_{\nu}(p) \quad (m \in \mathfrak{M}), \qquad \widehat{(\Delta_kA)_{\mu}}(p) = \sum_{\nu}\boldsymbol{\Lambda}_{\mu\nu}(p) \cdot \hat{A}_{\nu}(p).
\end{equation}
\end{proposition}

\begin{proof}
\emph{Positivity.} Let $p \in [-\pi,\pi]^3$ be real and let $m \in \mathfrak{M}\setminus\{0\}$. Choose $\nu$ with $m_{\nu} \ne 0$; then $1 \le |m_{\nu}| \le \mathbf{h}$. If $m_{\nu} \ge 1$, then $q_{m,\nu} = p_{\nu} + 2\pi m_{\nu}$ lies in $[\pi, 2\pi\mathbf{h} + \pi] = [\pi, \pi\mathfrak{N}]$; if $m_{\nu} \le -1$, then $q_{m,\nu} + 2\pi\mathfrak{N}$ lies in $[2\pi\mathfrak{N} - 2\pi\mathbf{h} - \pi, 2\pi\mathfrak{N} - \pi] = [\pi\mathfrak{N}, 2\pi\mathfrak{N} - \pi]$. In both cases $q_{m,\nu}$ is congruent modulo $2\pi\mathfrak{N}$ to a point of $[\pi, 2\pi\mathfrak{N} - \pi]$, hence $\sin\left(\frac{q_{m,\nu}}{2\mathfrak{N}}\right) \ne 0$ and $\operatorname{sn}(q_{m,\nu})^2 > 0$. The other summands of $\hat{D}(q_m)$ are squares of real numbers; therefore $\hat{D}(q_m) > 0$. As $\hat{U}$ and $u$ are real on real arguments, all terms of the sums $X_{\mu}$ and $Y$ are non-negative, so that $X_{\mu} \ge 0$ and $Y \ge 0$; also $D_0 \ge 0$. Moreover $u_{\nu} > 0$ for every $\nu$: one has $u(0) = 1$, and for $0 < |s| \le \pi$ the numbers $\operatorname{cs}(s) = 2\sin\left(\frac{s}{2}\right)$ and $\operatorname{sn}(s) = 2\mathfrak{N} \cdot \sin\left(\frac{s}{2\mathfrak{N}}\right)$ are different from zero and have the sign of $s$, so that $u(s) = \operatorname{cs}(s)/\operatorname{sn}(s) > 0$. Consequently $U_0 = u_0u_1u_2 > 0$, and
\[
N_{\mu} \ge U_0^2u_{\mu}^2 > 0, \qquad N_G \ge U_0^2 > 0.
\]
If $p = 0$, then $s_{\mathfrak{l}} = 0$ for every $\mathfrak{l}$, and $\Pi = \mathcal{N} > 0$. If $p \ne 0$, then $0 < |p_{\nu}| \le \pi$ for some $\nu$, whence $s_{\nu}^2 > 0$ and $D_0 > 0$. Dividing the definition of $\Pi$ in \eqref{eq:operators-photon-5} by $\mathcal{N}$ and using $\sum_{\mathfrak{l}}s_{\mathfrak{l}}^2 = D_0$ and $U_0^2u_{\mathfrak{l}}^2 = N_{\mathfrak{l}} - D_0X_{\mathfrak{l}}$, as in the proof of Lemma~\ref{lem:algebraic-identities}(a), one obtains
\[
\Pi = \mathcal{N}\sum_{\mathfrak{l}}\frac{s_{\mathfrak{l}}^2}{D_0}\left(1 - \frac{D_0X_{\mathfrak{l}}}{N_{\mathfrak{l}}}\right) = \mathcal{N}\sum_{\mathfrak{l}}\frac{s_{\mathfrak{l}}^2}{D_0} \cdot \frac{U_0^2u_{\mathfrak{l}}^2}{N_{\mathfrak{l}}} > 0.
\]
Since $p \in P_l$ was arbitrary, the functions \eqref{eq:operators-photon-6} are defined at every element of $P_l$.

\emph{A solution in Fourier variables.} Let $A$ be a real bond function on $T^0_k$. Let $\mathscr{A}$ be the complex bond function on $T_k$, $\boldsymbol{\Lambda}$ the complex bond function on $T^0_k$ and $\boldsymbol{\omega}$ the complex function on $T^0_k$ whose transforms are
\[
\hat{\mathscr{A}}_{\mu}(q_m) := \sum_{\nu}h_{\mu\nu}(p,m) \cdot \hat{A}_{\nu}(p), \qquad \hat{\boldsymbol{\Lambda}}_{\mu}(p) := \sum_{\nu}\boldsymbol{\Lambda}_{\mu\nu}(p) \cdot \hat{A}_{\nu}(p), \qquad \hat{\boldsymbol{\omega}}(p) := \sum_{\nu}\boldsymbol{\omega}_{\nu}(p) \cdot \hat{A}_{\nu}(p)
\]
for all $p \in P_l$ and $m \in \mathfrak{M}$; they exist and are unique, as functions are recovered from their transforms (Section~\ref{sec:elementary-symbols-fourier}). It is shown that
\[
Q_k\mathscr{A} = A, \qquad \delta d\mathscr{A} = Q_k^T\boldsymbol{\Lambda}, \qquad \Delta\delta\mathscr{A} = Q_k^T\boldsymbol{\omega},
\]
so that $\mathscr{A}$ satisfies the conditions \eqref{eq:operators-photon-M1}--\eqref{eq:operators-photon-M3} of Lemma~\ref{lem:characterisation-minimiser}. By \eqref{eq:operators-photon-4}, the transforms of $Q_k\mathscr{A}$, $\boldsymbol{\Lambda}$ and $\boldsymbol{\omega}$ at $p$, and those of $\delta d\mathscr{A}$, $\Delta\delta\mathscr{A}$, $Q_k^T\boldsymbol{\Lambda}$ and $Q_k^T\boldsymbol{\omega}$ at the aliases $q_m$, $m \in \mathfrak{M}$, involve only the transforms of $\mathscr{A}$ at these aliases and those of $A$, $\boldsymbol{\Lambda}$, $\boldsymbol{\omega}$ at $p$. Hence the three identities can be verified separately for each $p \in P_l$.

\emph{The frequencies with $p \ne 0$.} Let $p \in P_l$, $p \ne 0$. Then $D_0 > 0$, and by Lemma~\ref{lem:algebraic-identities}(c) the formula \eqref{eq:operators-photon-6} for $h_{\mu\nu}(p,m)$, $m \ne 0$, holds for $m = 0$ as well. Therefore, for every alias $q = q_m$ with $m \in \mathfrak{M}$,
\[
\hat{\mathscr{A}}_{\mu}(q) = \frac{\hat{U}(q) \cdot \hat{V}^{*}_{\mu}(q) \cdot \hat{\boldsymbol{\Lambda}}_{\mu}(p)}{\hat{D}(q)} + \frac{\hat{d}_{\mu}(q) \cdot \hat{U}(q) \cdot \hat{\boldsymbol{\omega}}(p)}{\hat{D}(q)^2}.
\]

\emph{Condition \eqref{eq:operators-photon-M1}.} By \eqref{eq:operators-photon-4}, the relations $\hat{V}_{\mu}(q) \cdot \hat{V}^{*}_{\mu}(q) = u(q_{\mu})^2$ and $\hat{V}_{\mu}(q) \cdot \hat{d}_{\mu}(q) = \rho_{\mu}(p)$ of \eqref{eq:operators-photon-3}, and Lemma~\ref{lem:algebraic-identities}(d),
\begin{align*}
\widehat{(Q_k\mathscr{A})_{\mu}}(p) &= \sum_{m \in \mathfrak{M}}\frac{\hat{U}(q_m)^2 \cdot u(q_{m,\mu})^2}{\hat{D}(q_m)} \cdot \hat{\boldsymbol{\Lambda}}_{\mu}(p) + \rho_{\mu}\sum_{m \in \mathfrak{M}}\frac{\hat{U}(q_m)^2}{\hat{D}(q_m)^2} \cdot \hat{\boldsymbol{\omega}}(p)\\
&= \frac{N_{\mu}}{D_0} \cdot \hat{\boldsymbol{\Lambda}}_{\mu}(p) + \frac{\rho_{\mu}N_G}{D_0^2} \cdot \hat{\boldsymbol{\omega}}(p).
\end{align*}
By \eqref{eq:operators-photon-6},
\begin{align*}
\frac{N_{\mu}}{D_0} \cdot \hat{\boldsymbol{\Lambda}}_{\mu}(p) &= \hat{A}_{\mu}(p) - \rho_{\mu}\sum_{\nu}\frac{\rho^{*}_{\nu}U_0^2 \cdot \mathcal{N}}{N_{\nu}\Pi D_0} \cdot \hat{A}_{\nu}(p),\\
\frac{\rho_{\mu}N_G}{D_0^2} \cdot \hat{\boldsymbol{\omega}}(p) &= \rho_{\mu}\sum_{\nu}\frac{\rho^{*}_{\nu}U_0^2 \cdot \mathcal{N}}{N_{\nu}\Pi D_0} \cdot \hat{A}_{\nu}(p).
\end{align*}
Adding, one finds $\widehat{(Q_k\mathscr{A})_{\mu}}(p) = \hat{A}_{\mu}(p)$.

\emph{The divergence.} One has $\delta\mathscr{A} = \sum_{\mu}\partial^T_{\mu}\mathscr{A}_{\mu}$ (Section~\ref{sec:tori-cochains-kernels}). By \eqref{eq:operators-photon-4}, the relation $\hat{V}^{*}_{\mu}(q) \cdot \hat{d}^{*}_{\mu}(q) = \rho^{*}_{\mu}(p)$ of \eqref{eq:operators-photon-3} and $\sum_{\mu}\hat{d}^{*}_{\mu} \cdot \hat{d}_{\mu} = \hat{D}$,
\[
\widehat{\delta\mathscr{A}}(q) = \sum_{\mu}\hat{d}^{*}_{\mu}(q) \cdot \hat{\mathscr{A}}_{\mu}(q) = \frac{\hat{U}(q)}{\hat{D}(q)}\left(\sum_{\mu}\rho^{*}_{\mu}\hat{\boldsymbol{\Lambda}}_{\mu}(p) + \hat{\boldsymbol{\omega}}(p)\right) = \frac{\hat{U}(q) \cdot \hat{\boldsymbol{\omega}}(p)}{\hat{D}(q)},
\]
because, by Lemma~\ref{lem:algebraic-identities}(b),
\[
\sum_{\mu}\rho^{*}_{\mu}\hat{\boldsymbol{\Lambda}}_{\mu}(p) = \sum_{\nu}\left(\sum_{\mu}\rho^{*}_{\mu}\boldsymbol{\Lambda}_{\mu\nu}(p)\right)\hat{A}_{\nu}(p) = 0.
\]

\emph{Condition \eqref{eq:operators-photon-M3}.} By \eqref{eq:operators-photon-4} and the formula for the divergence,
\[
\widehat{\Delta\delta\mathscr{A}}(q_m) = \hat{D}(q_m) \cdot \widehat{\delta\mathscr{A}}(q_m) = \hat{U}(q_m) \cdot \hat{\boldsymbol{\omega}}(p) = \widehat{Q_k^T\boldsymbol{\omega}}(q_m).
\]

\emph{Condition \eqref{eq:operators-photon-M2}.} On bond functions $\delta d = \Delta - d\delta$ (Section~\ref{sec:tori-cochains-kernels}), and $(d\delta\mathscr{A})_{\mu} = \partial_{\mu}\delta\mathscr{A}$. By \eqref{eq:operators-photon-4} and the formula for the divergence,
\[
\widehat{(\delta d\mathscr{A})_{\mu}}(q) = \hat{D}(q) \cdot \hat{\mathscr{A}}_{\mu}(q) - \hat{d}_{\mu}(q) \cdot \widehat{\delta\mathscr{A}}(q) = \hat{U}(q) \cdot \hat{V}^{*}_{\mu}(q) \cdot \hat{\boldsymbol{\Lambda}}_{\mu}(p) = \widehat{(Q_k^T\boldsymbol{\Lambda})_{\mu}}(q).
\]

\emph{The frequencies with $p = 0$.} Here $q_m = 2\pi m$. For $0 < |m_{\nu}| \le \mathbf{h}$ one has $\operatorname{cs}(2\pi m_{\nu}) = 2\sin(\pi m_{\nu}) = 0$ and $\operatorname{sn}(2\pi m_{\nu}) = 2\mathfrak{N} \cdot \sin\left(\frac{\pi m_{\nu}}{\mathfrak{N}}\right) \ne 0$; so $u(2\pi m_{\nu}) = 0$ by $\operatorname{cs} = u \cdot \operatorname{sn}$. Hence $\hat{U}(2\pi m) = 0$ for $m \in \mathfrak{M}\setminus\{0\}$, and $X_{\mu}(0) = Y(0) = 0$. Furthermore $u(0) = 1$, $\operatorname{sn}(0) = 0$ and $\rho_{\mu}(0) = \rho^{*}_{\mu}(0) = 0$; thus $U_0 = 1$, $D_0 = 0$ and
\[
N_{\mu}(0) = N_G(0) = \Pi(0) = 1.
\]
Since moreover $\hat{V}^{*}_{\mu}(0) = 1$ and $\hat{d}_{\mu}(0) = 0$, the definitions \eqref{eq:operators-photon-6} give
\[
h_{\mu\nu}(0,m) = \delta_{\mu\nu}\delta_{m,0} \quad (m \in \mathfrak{M}), \qquad \boldsymbol{\Lambda}_{\mu\nu}(0) = 0, \qquad \boldsymbol{\omega}_{\nu}(0) = 0.
\]
Consequently $\hat{\mathscr{A}}_{\mu}(2\pi m) = \delta_{m,0}\hat{A}_{\mu}(0)$ for $m \in \mathfrak{M}$, and $\hat{\boldsymbol{\Lambda}}_{\mu}(0) = 0$, $\hat{\boldsymbol{\omega}}(0) = 0$. At the frequencies $2\pi m$, $m \in \mathfrak{M}$, the transform of $\mathscr{A}$ thus coincides with the transform of the constant bond function $c$ on $T_k$ with components $c_{\mu} = l^{-3}\hat{A}_{\mu}(0)$. By Lemma~\ref{lem:block-averages}(a), $Q_kc$ is the constant bond function $c$ on $T^0_k$, whose transform at $0$ is $l^3c_{\mu} = \hat{A}_{\mu}(0)$; moreover $dc = 0$ and $\delta c = 0$. By \eqref{eq:operators-photon-4}, the transform of $Q_k\mathscr{A}$ at $0$ and the transforms of $\delta d\mathscr{A}$ and $\Delta\delta\mathscr{A}$ at the frequencies $2\pi m$ depend only on the transforms of $\mathscr{A}$ at these frequencies; so they coincide with the transform of $Q_kc$ at $0$, which is $\hat{A}_{\mu}(0)$, and with the transforms of $\delta dc = 0$ and $\Delta\delta c = 0$. On the other hand, the transforms of $Q_k^T\boldsymbol{\Lambda}$ and $Q_k^T\boldsymbol{\omega}$ vanish at the frequencies $2\pi m$ by \eqref{eq:operators-photon-4}, since $\hat{\boldsymbol{\Lambda}}_{\mu}(0) = 0$ and $\hat{\boldsymbol{\omega}}(0) = 0$. Therefore the three identities $Q_k\mathscr{A} = A$, $\delta d\mathscr{A} = Q_k^T\boldsymbol{\Lambda}$ and $\Delta\delta\mathscr{A} = Q_k^T\boldsymbol{\omega}$ hold at the frequency $0$ of $T^0_k$ and at the frequencies $2\pi m$, $m \in \mathfrak{M}$, of $T_k$.

\emph{Conclusion.} As functions are recovered from their transforms, $Q_k\mathscr{A} = A$, $\delta d\mathscr{A} = Q_k^T\boldsymbol{\Lambda}$ and $\Delta\delta\mathscr{A} = Q_k^T\boldsymbol{\omega}$. The operators $Q_k$, $d$, $\delta$, $\Delta$ and $Q_k^T$ map real functions to real functions, and $A$ is real. Hence $\operatorname{Re}\mathscr{A}$ satisfies \eqref{eq:operators-photon-M1}--\eqref{eq:operators-photon-M3}, with $\delta d\operatorname{Re}\mathscr{A} = Q_k^T\operatorname{Re}\boldsymbol{\Lambda}$ and $\Delta\delta\operatorname{Re}\mathscr{A} = Q_k^T\operatorname{Re}\boldsymbol{\omega}$, and $\operatorname{Im}\mathscr{A}$ satisfies \eqref{eq:operators-photon-M1}--\eqref{eq:operators-photon-M3} with $0$ in place of $A$, as does the zero field. By the uniqueness in Lemma~\ref{lem:characterisation-minimiser} and the definition of $H_k$ in Proposition~\ref{prop:minimiser}, $\operatorname{Re}\mathscr{A} = H_kA$ and $\operatorname{Im}\mathscr{A} = 0$. So $\mathscr{A} = H_kA$, and the first identity of \eqref{eq:operators-photon-7} is the definition of $\hat{\mathscr{A}}$. Finally, let $A'$ be a further real bond function on $T^0_k$. Since $\Delta_k = (dH_k)^TdH_k$, $\delta = d^T$, $\delta dH_kA = \delta d\mathscr{A} = Q_k^T\boldsymbol{\Lambda}$, and $Q_kH_k = 1$ by Proposition~\ref{prop:minimiser},
\[
\langle A', \Delta_kA\rangle = \langle dH_kA', dH_kA\rangle = \langle H_kA', \delta dH_kA\rangle = \langle H_kA', Q_k^T\boldsymbol{\Lambda}\rangle = \langle Q_kH_kA', \boldsymbol{\Lambda}\rangle = \langle A', \boldsymbol{\Lambda}\rangle.
\]
The pairing is bilinear and $A'$ is an arbitrary real bond function; hence $\Delta_kA = \boldsymbol{\Lambda}$, and the second identity of \eqref{eq:operators-photon-7} is the definition of $\hat{\boldsymbol{\Lambda}}$.
\end{proof}

\begin{remark}\label{rem:operators-photon-1}
Let $p$ be a point at which the functions \eqref{eq:operators-photon-6} are defined and $D_0 \ne 0$, and let
\[
E_{\mu} := \frac{N_{\mu}}{D_0}, \qquad G_0 := \frac{N_G}{D_0^2},
\]
which by Lemma~\ref{lem:algebraic-identities}(d) are sums over all $m \in \mathfrak{M}$. For such $p$ the solution of \eqref{eq:operators-photon-M1}--\eqref{eq:operators-photon-M3} constructed in the proof of Proposition~\ref{prop:symbols} is expressed through $E_{\mu}$ and $G_0$ in the display in the proof of Lemma~\ref{lem:periodicity}. At $p = 0$ one has $D_0 = 0$ and $N_{\mu} = N_G = 1$, so that $E_{\mu}$ and $G_0$ are singular there; the formulas \eqref{eq:operators-photon-6} give the same functions in a form from which this singularity has been divided out.
\end{remark}

\subsection{Bounds in a complex strip}\label{sec:bounds-complex-strip}

\subsubsection{Domains and numerical constants}\label{sec:domains-numerical-constants}

As in Section~\ref{sec:tori-cochains-kernels}, $\langle t\rangle=\max\{1,|t|\}$ for real $t$, and $\langle m\rangle_{\infty}=\max_{\nu}\langle m_{\nu}\rangle$ for $m\in\mathbb Z^3$. Define the open sets $\boldsymbol{\Omega}(a,g)\subset\mathbb C^3$ for $a,g>0$, the numbers $Z_a$ for $a>1$, and the constants $\mathrm{ar}_1,\dots,\gamma_H$ by
\begin{align}
\boldsymbol{\Omega}(a,g)&:=\left\{p\in\mathbb C^3:\ |\operatorname{Re}p_{\nu}|<a\ \text{and}\ |\operatorname{Im}p_{\nu}|<g\ \text{for}\ \nu=0,1,2\right\},\notag\\
\mathrm{ar}_1&:=\frac{5\pi}{4},\qquad g_1:=\frac{1}{4},\qquad\mathrm{ar}_2:=\mathrm{ar}_1-g_1>\pi+\frac{1}{2},\notag\\
a_{\mathrm{sn}}&:=\frac{8}{5\pi}\sin\left(\frac{5\pi}{8}\right)\in(0.4705,\,0.4706),\qquad c_{\mathrm{ch}}:=\cosh\left(\frac{1}{8}\right)<1.008,\notag\\
c_u&:=\frac{8c_{\mathrm{ch}}}{3\pi a_{\mathrm{sn}}}<1.82,\qquad c_{\rho}:=2c_{\mathrm{ch}}\cdot\exp\left(\frac{1}{8}\right)<2.29,\notag\\
Z_a&:=\sum_{t\in\mathbb Z}\langle t\rangle^{-a}\le3+\frac{2}{a-1}\quad(a>1),\qquad Z_2=1+\frac{\pi^2}{3},\qquad B_D:=48,\notag\\
B_X&:=c_u^8\cdot\left(Z_2^3-1\right),\qquad B_Y:=c_u^6\cdot\left(Z_2^3-1\right),\notag\\
B_N&:=c_{\mathrm{ch}}^8+B_D\cdot B_X,\qquad B_G:=c_{\mathrm{ch}}^6+B_D^2\cdot B_Y,\qquad B_{\Pi}:=B_N^3+B_D\cdot B_X\cdot B_N^2,\notag\\
c_N&:=a_{\mathrm{sn}}^8,\qquad c_G:=a_{\mathrm{sn}}^6,\qquad c_{\Pi}:=a_{\mathrm{sn}}^{24},\notag\\
g_2&:=g_1\cdot\min\left\{\frac{c_N}{4B_N},\ \frac{c_G}{4B_G},\ \frac{c_{\Pi}}{4B_{\Pi}}\right\},\qquad\gamma_H:=\frac{g_2}{2}.\label{eq:operators-photon-8}
\end{align}
(For the bound on $Z_a$ note that $Z_a=1+2\sum_{t\ge1}t^{-a}$ and $\sum_{t\ge2}t^{-a}\le\int_1^{\infty}\mathrm{d}x\,x^{-a}=\frac{1}{a-1}$; the value of $Z_2$ is Euler's evaluation $\sum_{t\ge1}t^{-2}=\frac{\pi^2}{6}$.) The numbers $Z_a$ and the constants $\mathrm{ar}_1,\dots,\gamma_H$ are absolute: none of them depends on $\mathfrak{N}$, $l$, $L$ or $\alpha$.

\begin{lemma}[One complex variable]\label{lem:one-complex-variable}
Let $\operatorname{sn}$, $\operatorname{cs}$ and $u$ be the entire functions of Section~\ref{sec:elementary-symbols-fourier}, and let $s=s'+is''$ with $s',s''\in\mathbb R$ and $|s''|\le\frac{1}{4}$.
\begin{enumerate}
\item[(a)] $|u(s)|\le c_{\mathrm{ch}}$ and $|\operatorname{cs}(s)|\le2c_{\mathrm{ch}}$.
\item[(b)] $|\operatorname{sn}(s)|\le c_{\mathrm{ch}}\cdot|s|$.
\item[(c)] If $|s'|\le\pi\mathfrak{N}+\frac{\pi}{4}$, then
\[
|\operatorname{sn}(s)|^2\ge a_{\mathrm{sn}}^2\cdot s'^2+s''^2\qquad\text{and}\qquad\operatorname{Re}\left(\operatorname{sn}(s)^2\right)\ge a_{\mathrm{sn}}^2\cdot s'^2-c_{\mathrm{ch}}^2\cdot s''^2.
\]
\item[(d)] If $s$ is real and $|s|\le\frac{5\pi}{4}$, then $a_{\mathrm{sn}}\le u(s)\le1$.
\end{enumerate}
\end{lemma}

\begin{proof}
For real $x$, $y$ one has $\sin(x+iy)=\sin x\cdot\cosh y+i\cos x\cdot\sinh y$ and $\cos(x+iy)=\cos x\cdot\cosh y-i\sin x\cdot\sinh y$. With $\cosh^2y=1+\sinh^2y$ it follows that
\[
|\sin(x+iy)|^2=\sin^2x+\sinh^2y,\qquad\operatorname{Re}\left(\sin^2(x+iy)\right)=\sin^2x-\sinh^2y+2\sin^2x\cdot\sinh^2y\ge\sin^2x-\sinh^2y,
\]
and that $|\cos(x+iy)|^2=\cos^2x+\sinh^2y\le\cosh^2y$.

(a) By definition $u(s)=\frac{1}{\mathfrak{N}}\sum_{j=-\mathbf{h}}^{\mathbf{h}}\exp\left(\frac{isj}{\mathfrak{N}}\right)$, and the term of index $j$ has modulus $\exp\left(-\frac{s''j}{\mathfrak{N}}\right)$. The set of indices is invariant under $j\mapsto-j$, and $|j|\le\mathbf{h}<\frac{\mathfrak{N}}{2}$; hence
\[
|u(s)|\le\frac{1}{\mathfrak{N}}\sum_{j=-\mathbf{h}}^{\mathbf{h}}\exp\left(-\frac{s''j}{\mathfrak{N}}\right)=\frac{1}{\mathfrak{N}}\sum_{j=-\mathbf{h}}^{\mathbf{h}}\cosh\left(\frac{s''j}{\mathfrak{N}}\right)\le\cosh\left(\frac{s''}{2}\right)\le c_{\mathrm{ch}}.
\]
Moreover $\operatorname{cs}(s)=2\sin\left(\frac{s}{2}\right)$, and the first of the identities above gives
\[
|\operatorname{cs}(s)|\le2\left(1+\sinh^2\left(\frac{s''}{2}\right)\right)^{\frac{1}{2}}=2\cosh\left(\frac{s''}{2}\right)\le2c_{\mathrm{ch}}.
\]

(b) For complex $w$ one has $\sin w=w\int_0^1\mathrm{d}t\,\cos(tw)$, and $|\cos(tw)|\le\cosh(t\operatorname{Im}w)\le\cosh(\operatorname{Im}w)$ for $0\le t\le1$. Take $w=\frac{s}{2\mathfrak{N}}$; then $|\operatorname{Im}w|=\frac{|s''|}{2\mathfrak{N}}\le\frac{1}{8}$, and therefore
\[
|\operatorname{sn}(s)|=2\mathfrak{N}\cdot|\sin w|\le|s|\cdot\cosh\left(\frac{1}{8}\right)=c_{\mathrm{ch}}\cdot|s|.
\]

(c) Put $x:=\frac{s'}{2\mathfrak{N}}$ and $y:=\frac{s''}{2\mathfrak{N}}$, so that $\operatorname{sn}(s)=2\mathfrak{N}\cdot\sin(x+iy)$ and $|y|\le\frac{1}{8}$. Since $\mathfrak{N}\ge1$,
\[
|x|\le\frac{\pi}{2}+\frac{\pi}{8\mathfrak{N}}\le\frac{5\pi}{8}.
\]
The function $r\mapsto\frac{\sin r}{r}$ is even, and decreasing on $(0,\pi)$; hence $\frac{\sin x}{x}\ge a_{\mathrm{sn}}>0$ if $x\ne0$, and so $\sin^2x\ge a_{\mathrm{sn}}^2\cdot x^2$, which is trivial for $x=0$. Moreover $y^2\le\sinh^2y\le y^2\cdot\cosh^2y\le c_{\mathrm{ch}}^2\cdot y^2$. The identities at the beginning of the proof now give
\[
|\operatorname{sn}(s)|^2=4\mathfrak{N}^2\cdot\left(\sin^2x+\sinh^2y\right)\ge4\mathfrak{N}^2\cdot\left(a_{\mathrm{sn}}^2\cdot x^2+y^2\right)=a_{\mathrm{sn}}^2\cdot s'^2+s''^2
\]
and
\[
\operatorname{Re}\left(\operatorname{sn}(s)^2\right)\ge4\mathfrak{N}^2\cdot\left(\sin^2x-\sinh^2y\right)\ge4\mathfrak{N}^2\cdot\left(a_{\mathrm{sn}}^2\cdot x^2-c_{\mathrm{ch}}^2\cdot y^2\right)=a_{\mathrm{sn}}^2\cdot s'^2-c_{\mathrm{ch}}^2\cdot s''^2.
\]

(d) Let $s$ be real. The terms of $u(s)$ with the indices $j$ and $-j$ are complex conjugate; so $u$ is real and even on $\mathbb R$, $u(0)=1$, and $|u(s)|\le1$, because $u(s)$ is the mean of $\mathfrak{N}$ numbers of modulus one. Let $0<s\le\frac{5\pi}{4}$ and $t:=\frac{s}{2}$. Then $\sin t>0$ and $0<\frac{t}{\mathfrak{N}}\le\frac{5\pi}{8}<\pi$, so that $0<\mathfrak{N}\cdot\sin\left(\frac{t}{\mathfrak{N}}\right)\le t$; in particular $\operatorname{sn}(s)=2\mathfrak{N}\cdot\sin\left(\frac{t}{\mathfrak{N}}\right)\ne0$. The identity $\operatorname{cs}=u\cdot\operatorname{sn}$ of Section~\ref{sec:elementary-symbols-fourier} and the monotonicity used in (c) give
\[
u(s)=\sin t\cdot\left(\mathfrak{N}\cdot\sin\left(\frac{t}{\mathfrak{N}}\right)\right)^{-1}\ge\frac{\sin t}{t}\ge a_{\mathrm{sn}}.
\]
As $u$ is even, the same lower bound holds for $-\frac{5\pi}{4}\le s<0$, and it holds at $s=0$ because $a_{\mathrm{sn}}<1$.
\end{proof}

\begin{lemma}[Estimates at the aliases]\label{lem:estimates-aliases}
Let $p\in\boldsymbol{\Omega}(\mathrm{ar}_1,g_1)$, $m\in\mathfrak{M}$ and $q=q_m=p+2\pi m$, and let $u_{\nu}$, $U_0$, $D_0$, $s_{\nu}$ be as in Section~\ref{sec:sums-aliases}.
\begin{enumerate}
\item[(a)] For $\nu=0,1,2$,
\[
|u(q_{\nu})|\le\frac{c_u}{\langle m_{\nu}\rangle},\qquad\left|\hat{V}^{*}_{\nu}(q)\right|\le\exp\left(\frac{1}{8}\right)\cdot\frac{c_u}{\langle m_{\nu}\rangle},\qquad|\hat{U}(q)|\le c_u^3\cdot\prod_{\sigma=0}^2\langle m_{\sigma}\rangle^{-1};
\]
for $m=0$, moreover, $|u_{\nu}|\le c_{\mathrm{ch}}$ and $|U_0|\le c_{\mathrm{ch}}^3$.
\item[(b)] If $m\ne0$, then $\operatorname{Re}\hat{D}(q_m)\ge|m|_2^2$.
\item[(c)] $\sum_{\nu}|\operatorname{sn}(p_{\nu})|^2\le B_D$; in particular $|D_0|\le B_D$.
\item[(d)] $|\rho_{\nu}(p)|\le c_{\rho}$ and $|\rho^{*}_{\nu}(p)|\le c_{\rho}$.
\item[(e)] $|\hat{d}_{\nu}(q_m)|\le12\langle m_{\nu}\rangle$, and $\sum_{\sigma=0}^2|q_{m,\sigma}|\le31.5\langle m\rangle_{\infty}$.
\item[(f)] The functions $X_{\mu}$, $Y$, $N_{\mu}$, $N_G$, $\Pi$ of \eqref{eq:operators-photon-5} are analytic on $\boldsymbol{\Omega}(\mathrm{ar}_1,g_1)$, and on this set
\[
|X_{\mu}|\le B_X,\qquad|Y|\le B_Y,\qquad|N_{\mu}|\le B_N,\qquad|N_G|\le B_G,\qquad|\Pi|\le B_{\Pi}.
\]
\item[(g)] If $p$ is real, $p\in(-\mathrm{ar}_1,\mathrm{ar}_1)^3$, then $N_{\mu}(p)\ge c_N$, $N_G(p)\ge c_G$ and $\Pi(p)\ge c_{\Pi}$.
\end{enumerate}
\end{lemma}

\begin{proof}
Since $|\operatorname{Re}p_{\nu}|<\frac{5\pi}{4}$, $|\operatorname{Im}p_{\nu}|<\frac{1}{4}$ and $|m_{\nu}|\le\mathbf{h}=\frac{\mathfrak{N}-1}{2}$, each component of $q$ satisfies
\[
|\operatorname{Re}q_{\nu}|\le\frac{5\pi}{4}+2\pi\mathbf{h}=\pi\mathfrak{N}+\frac{\pi}{4},\qquad|\operatorname{Im}q_{\nu}|=|\operatorname{Im}p_{\nu}|<\frac{1}{4}.
\]
Hence Lemma~\ref{lem:one-complex-variable}, clause (c) included, applies to $s=q_{\nu}$.

(a) If $m_{\nu}=0$, then $q_{\nu}=p_{\nu}$, and Lemma~\ref{lem:one-complex-variable}(a) gives $|u(q_{\nu})|\le c_{\mathrm{ch}}\le c_u=\frac{c_u}{\langle m_{\nu}\rangle}$. If $m_{\nu}\ne0$, then
\[
|\operatorname{Re}q_{\nu}|\ge2\pi|m_{\nu}|-\frac{5\pi}{4}\ge\frac{3\pi}{4}|m_{\nu}|,
\]
and Lemma~\ref{lem:one-complex-variable}(c) gives $|\operatorname{sn}(q_{\nu})|\ge a_{\mathrm{sn}}\cdot|\operatorname{Re}q_{\nu}|\ge\frac{3\pi}{4}a_{\mathrm{sn}}\cdot|m_{\nu}|>0$. With the identity $\operatorname{cs}=u\cdot\operatorname{sn}$ of Section~\ref{sec:elementary-symbols-fourier} and Lemma~\ref{lem:one-complex-variable}(a) it follows that
\[
|u(q_{\nu})|=\frac{|\operatorname{cs}(q_{\nu})|}{|\operatorname{sn}(q_{\nu})|}\le\frac{8c_{\mathrm{ch}}}{3\pi a_{\mathrm{sn}}\cdot|m_{\nu}|}=\frac{c_u}{\langle m_{\nu}\rangle}.
\]
Next, $\hat{V}^{*}_{\nu}(q)=\exp\left(-\frac{i}{2}(1-\eta)q_{\nu}\right)\cdot u(q_{\nu})$, and the exponential factor has the modulus $\exp\left(\frac{1}{2}(1-\eta)\operatorname{Im}q_{\nu}\right)\le\exp\left(\frac{1}{8}\right)$; this gives the bound on $\hat{V}^{*}_{\nu}(q)$. The bound on $\hat{U}(q)=u(q_0)\cdot u(q_1)\cdot u(q_2)$ is the product of the bounds on the three factors. For $m=0$, Lemma~\ref{lem:one-complex-variable}(a) gives $|u_{\nu}|=|u(p_{\nu})|\le c_{\mathrm{ch}}$, and hence $|U_0|=|u_0\cdot u_1\cdot u_2|\le c_{\mathrm{ch}}^3$.

(b) Let $m\ne0$. By Lemma~\ref{lem:one-complex-variable}(c), for each $\nu$,
\[
\operatorname{Re}\left(\operatorname{sn}(q_{m,\nu})^2\right)\ge a_{\mathrm{sn}}^2\cdot\left(\operatorname{Re}q_{m,\nu}\right)^2-c_{\mathrm{ch}}^2\cdot\left(\operatorname{Im}q_{m,\nu}\right)^2,
\]
where $\left(\operatorname{Re}q_{m,\nu}\right)^2\ge\left(\frac{3\pi}{4}\right)^2m_{\nu}^2$ (by the proof of (a) if $m_{\nu}\ne0$, trivially if $m_{\nu}=0$) and $\left(\operatorname{Im}q_{m,\nu}\right)^2<\frac{1}{16}$. As $\hat{D}(q)=\sum_{\nu}\operatorname{sn}(q_{\nu})^2$, summation over $\nu$ gives
\[
\operatorname{Re}\hat{D}(q_m)\ge a_{\mathrm{sn}}^2\cdot\left(\frac{3\pi}{4}\right)^2\cdot|m|_2^2-\frac{3c_{\mathrm{ch}}^2}{16}\ge1.228\,|m|_2^2-0.191\ge|m|_2^2,
\]
where the numerical bounds of \eqref{eq:operators-photon-8} were used, and the last inequality holds because $|m|_2^2\ge1$ for $m\ne0$.

(c) By Lemma~\ref{lem:one-complex-variable}(b) and $|p_{\nu}|^2=\left(\operatorname{Re}p_{\nu}\right)^2+\left(\operatorname{Im}p_{\nu}\right)^2<\mathrm{ar}_1^2+g_1^2$,
\[
\sum_{\nu}|\operatorname{sn}(p_{\nu})|^2\le3c_{\mathrm{ch}}^2\cdot\left(\mathrm{ar}_1^2+g_1^2\right)<47.2<B_D,
\]
and $|D_0|=\left|\sum_{\nu}\operatorname{sn}(p_{\nu})^2\right|\le\sum_{\nu}|\operatorname{sn}(p_{\nu})|^2$.

(d) By the definitions,
\[
\rho_{\nu}(p)=\exp(ip_{\nu})-1=i\exp\left(\frac{i}{2}p_{\nu}\right)\cdot\operatorname{cs}(p_{\nu}),\qquad\rho^{*}_{\nu}(p)=\exp(-ip_{\nu})-1=-i\exp\left(-\frac{i}{2}p_{\nu}\right)\cdot\operatorname{cs}(p_{\nu}).
\]
Both exponential factors have modulus at most $\exp\left(\frac{1}{2}|\operatorname{Im}p_{\nu}|\right)\le\exp\left(\frac{1}{8}\right)$, and $|\operatorname{cs}(p_{\nu})|\le2c_{\mathrm{ch}}$ by Lemma~\ref{lem:one-complex-variable}(a). So both moduli are at most $2c_{\mathrm{ch}}\cdot\exp\left(\frac{1}{8}\right)=c_{\rho}$.

(e) For each $\nu$,
\[
|q_{m,\nu}|\le|\operatorname{Re}q_{m,\nu}|+|\operatorname{Im}q_{m,\nu}|\le2\pi|m_{\nu}|+\frac{5\pi}{4}+\frac{1}{4}\le10.5\,\langle m_{\nu}\rangle,
\]
because $|m_{\nu}|\le\langle m_{\nu}\rangle$, $1\le\langle m_{\nu}\rangle$ and $2\pi+\frac{5\pi}{4}+\frac{1}{4}<10.5$. Summation over $\nu$, with $\langle m_{\nu}\rangle\le\langle m\rangle_{\infty}$, gives the second assertion. By Section~\ref{sec:elementary-symbols-fourier}, $\hat{d}_{\nu}(q)=i\exp\left(\frac{i}{2}\eta q_{\nu}\right)\cdot\operatorname{sn}(q_{\nu})$, and the exponential factor has the modulus $\exp\left(-\frac{1}{2}\eta\operatorname{Im}q_{\nu}\right)\le\exp\left(\frac{1}{8}\right)$. Lemma~\ref{lem:one-complex-variable}(b) and the bound just proved yield
\[
|\hat{d}_{\nu}(q_m)|\le\exp\left(\frac{1}{8}\right)\cdot c_{\mathrm{ch}}\cdot|q_{m,\nu}|\le10.5\exp\left(\frac{1}{8}\right)\cdot c_{\mathrm{ch}}\cdot\langle m_{\nu}\rangle\le12\langle m_{\nu}\rangle,
\]
since $10.5\exp\left(\frac{1}{8}\right)\cdot c_{\mathrm{ch}}<11.992$.

(f) For $m\ne0$, (b) gives $|\hat{D}(q_m)|\ge\operatorname{Re}\hat{D}(q_m)\ge|m|_2^2\ge1$; so no denominator in \eqref{eq:operators-photon-5} vanishes on $\boldsymbol{\Omega}(\mathrm{ar}_1,g_1)$. The sums in \eqref{eq:operators-photon-5} are finite; hence $X_{\mu}$ and $Y$ are analytic on $\boldsymbol{\Omega}(\mathrm{ar}_1,g_1)$, and so are $N_{\mu}$, $N_G$ and $\Pi$, which are polynomials in $X_0,X_1,X_2,Y$ and in entire functions of $p$. By (a), and since $|\hat{D}(q_m)|\ge|m|_2^2\ge1$, one has for $m\ne0$
\[
\left|\frac{\hat{U}(q_m)^2\cdot u(q_{m,\mu})^2}{\hat{D}(q_m)}\right|\le c_u^8\cdot\prod_{\sigma=0}^2\langle m_{\sigma}\rangle^{-2},\qquad\left|\frac{\hat{U}(q_m)^2}{\hat{D}(q_m)^2}\right|\le c_u^6\cdot\prod_{\sigma=0}^2\langle m_{\sigma}\rangle^{-2}.
\]
The sum of $\prod_{\sigma}\langle m_{\sigma}\rangle^{-2}$ over all $m\in\mathbb Z^3$ equals $Z_2^3$, and its term $m=0$ equals $1$. Summation of the first bound over $m\in\mathfrak{M}\setminus\{0\}\subset\mathbb Z^3\setminus\{0\}$ therefore gives
\[
|X_{\mu}|\le c_u^8\cdot\sum_{m\in\mathbb Z^3,\,m\ne0}\ \prod_{\sigma=0}^2\langle m_{\sigma}\rangle^{-2}=c_u^8\cdot\left(Z_2^3-1\right)=B_X,
\]
and the same summation of the second bound gives $|Y|\le c_u^6\cdot\left(Z_2^3-1\right)=B_Y$. From \eqref{eq:operators-photon-5}, (a) for $m=0$, and (c),
\begin{align*}
|N_{\mu}|&\le|U_0|^2\cdot|u_{\mu}|^2+|D_0|\cdot|X_{\mu}|\le c_{\mathrm{ch}}^8+B_D\cdot B_X=B_N,\\
|N_G|&\le|U_0|^2+|D_0|^2\cdot|Y|\le c_{\mathrm{ch}}^6+B_D^2\cdot B_Y=B_G,\\
|\Pi|&\le|N_0\cdot N_1\cdot N_2|+\sum_{\mathfrak{l}}|s_{\mathfrak{l}}|^2\cdot|X_{\mathfrak{l}}|\cdot\prod_{\mathfrak{l}'\ne\mathfrak{l}}|N_{\mathfrak{l}'}|\le B_N^3+B_D\cdot B_X\cdot B_N^2=B_{\Pi}.
\end{align*}

(g) Let $p\in(-\mathrm{ar}_1,\mathrm{ar}_1)^3$ be real. Then every $q_m$ is real, $u$ and $\operatorname{sn}$ are real on the real axis (for $u$ see the proof of Lemma~\ref{lem:one-complex-variable}(d)), and $\hat{D}(q_m)\ge|m|_2^2>0$ for $m\ne0$ by (b). Hence $X_{\mu}\ge0$, $Y\ge0$ and $D_0=\sum_{\mathfrak{l}}s_{\mathfrak{l}}^2\ge0$. As $|p_{\nu}|<\frac{5\pi}{4}$, Lemma~\ref{lem:one-complex-variable}(d) gives $u_{\nu}\ge a_{\mathrm{sn}}$, hence $U_0\ge a_{\mathrm{sn}}^3$, and \eqref{eq:operators-photon-5} yields
\[
N_{\mu}\ge U_0^2\cdot u_{\mu}^2\ge a_{\mathrm{sn}}^8=c_N,\qquad N_G\ge U_0^2\ge a_{\mathrm{sn}}^6=c_G.
\]
At $p=0$ one has $D_0=0$, $s_{\mathfrak{l}}=0$ and $u_{\nu}=u(0)=1$; so $N_{\mu}(0)=1$ and $\Pi(0)=1\ge c_{\Pi}$. Let $p\ne0$. For each $\mathfrak{l}$ one has $\left|\frac{p_{\mathfrak{l}}}{2\mathfrak{N}}\right|<\frac{5\pi}{8}<\pi$, so $s_{\mathfrak{l}}=2\mathfrak{N}\cdot\sin\left(\frac{p_{\mathfrak{l}}}{2\mathfrak{N}}\right)$ vanishes only if $p_{\mathfrak{l}}=0$; thus $D_0>0$. Since $D_0=\sum_{\mathfrak{l}}s_{\mathfrak{l}}^2$, one may write $N_0\cdot N_1\cdot N_2=\sum_{\mathfrak{l}}\frac{s_{\mathfrak{l}}^2}{D_0}\cdot N_{\mathfrak{l}}\cdot\prod_{\mathfrak{l}'\ne\mathfrak{l}}N_{\mathfrak{l}'}$; together with $N_{\mathfrak{l}}-D_0\cdot X_{\mathfrak{l}}=U_0^2\cdot u_{\mathfrak{l}}^2$, which is \eqref{eq:operators-photon-5}, the definition of $\Pi$ becomes, as in the proof of Proposition~\ref{prop:symbols},
\[
\Pi=\sum_{\mathfrak{l}}\frac{s_{\mathfrak{l}}^2}{D_0}\cdot\left(N_{\mathfrak{l}}-D_0\cdot X_{\mathfrak{l}}\right)\cdot\prod_{\mathfrak{l}'\ne\mathfrak{l}}N_{\mathfrak{l}'}=\sum_{\mathfrak{l}}\frac{s_{\mathfrak{l}}^2}{D_0}\cdot U_0^2\cdot u_{\mathfrak{l}}^2\cdot\prod_{\mathfrak{l}'\ne\mathfrak{l}}N_{\mathfrak{l}'}.
\]
The weights $\frac{s_{\mathfrak{l}}^2}{D_0}$ are non-negative with sum $1$, and $U_0^2\cdot u_{\mathfrak{l}}^2\cdot\prod_{\mathfrak{l}'\ne\mathfrak{l}}N_{\mathfrak{l}'}\ge a_{\mathrm{sn}}^8\cdot\left(a_{\mathrm{sn}}^8\right)^2=c_{\Pi}$ for each $\mathfrak{l}$; hence $\Pi(p)\ge c_{\Pi}$.
\end{proof}

\begin{lemma}[Lower bounds near the real cube]\label{lem:lower-bounds-near}
Let $B$ and $c>0$ be real numbers, and let $f\colon\boldsymbol{\Omega}(\mathrm{ar}_1,g_1)\to\mathbb C$ be an analytic function with $|f|\le B$ on $\boldsymbol{\Omega}(\mathrm{ar}_1,g_1)$, real on the cube $[-\mathrm{ar}_2,\mathrm{ar}_2]^3$ and satisfying $f\ge c$ there. Then
\[
|f(p)|\ge\frac{c}{2}\qquad\text{for all }p\in\boldsymbol{\Omega}\left(\mathrm{ar}_2,\frac{g_1c}{4B}\right).
\]
\end{lemma}

\begin{proof}
Put $g:=\frac{g_1c}{4B}$ and $R_0:=\frac{g_1}{g}=\frac{4B}{c}$. The origin belongs to $[-\mathrm{ar}_2,\mathrm{ar}_2]^3$, so $c\le f(0)\le B$ and $R_0\ge4$. Let $p\in\boldsymbol{\Omega}(\mathrm{ar}_2,g)$, and write $p=p'+ip''$ with $p',p''\in\mathbb R^3$; thus $p'\in(-\mathrm{ar}_2,\mathrm{ar}_2)^3$ and $|p''_{\nu}|<g$ for every $\nu$. Let $t=t_1+it_2$, with $t_1,t_2$ real, be a complex number with $|t|<R_0$. The point $p'+it\cdot p''$ has the real part $p'-t_2\cdot p''$ and the imaginary part $t_1\cdot p''$. Since $|t_1|<R_0$, $|t_2|<R_0$ and $R_0\cdot g=g_1$, its components satisfy
\[
\left|p'_{\nu}-t_2\cdot p''_{\nu}\right|<\mathrm{ar}_2+g_1=\mathrm{ar}_1,\qquad\left|t_1\cdot p''_{\nu}\right|<g_1\qquad(\nu=0,1,2),
\]
that is, $p'+it\cdot p''\in\boldsymbol{\Omega}(\mathrm{ar}_1,g_1)$. Consequently
\[
\varphi(t):=f\left(p'+it\cdot p''\right)-f(p')
\]
defines an analytic function on the disc $\{t\in\mathbb C:|t|<R_0\}$ with $|\varphi|\le2B$ and $\varphi(0)=0$. Schwarz's lemma, applied to $w\mapsto\frac{1}{2B}\cdot\varphi(R_0w)$ on the unit disc, gives $|\varphi(t)|\le\frac{2B}{R_0}\cdot|t|$ for $|t|<R_0$. The point $t=1$ lies in the disc, because $R_0\ge4$, and $\varphi(1)=f(p)-f(p')$; therefore
\[
|f(p)-f(p')|\le\frac{2B}{R_0}=\frac{c}{2}.
\]
As $p'\in[-\mathrm{ar}_2,\mathrm{ar}_2]^3$, one has $f(p')\ge c$, and it follows that $|f(p)|\ge c-\frac{c}{2}=\frac{c}{2}$.
\end{proof}

\subsubsection{Estimates for the symbol}\label{sec:estimates-symbol}

With the constants of \eqref{eq:operators-photon-8}, put
\begin{align}
\mathsf{A}_1&:=\frac{2B_D}{c_N}+\frac{8c_{\rho}^2\cdot c_{\mathrm{ch}}^6\cdot B_N^3}{c_N^2\cdot c_{\Pi}},\qquad\mathsf{A}_2:=\frac{8c_{\rho}\cdot B_D\cdot c_{\mathrm{ch}}^6\cdot B_N^3}{c_N\cdot c_G\cdot c_{\Pi}},\notag\\
\mathsf{A}_3&:=\frac{2\exp\left(\frac{1}{8}\right)\cdot c_{\mathrm{ch}}^4}{c_N}+\frac{192c_{\rho}\cdot c_{\mathrm{ch}}^9\cdot\left(B_X+c_{\mathrm{ch}}^2\cdot B_D\cdot B_Y\right)\cdot B_N^3}{c_N^2\cdot c_G\cdot c_{\Pi}},\notag\\
S_{*}(\alpha)&:=378\left(\mathsf{A}_3+\exp\left(\frac{1}{8}\right)\cdot c_u^4\cdot\mathsf{A}_1\cdot Z_2\cdot\left(3+\frac{4}{1-\alpha}\right)^2+12\cdot9^3c_u^3\cdot\mathsf{A}_2\right)\qquad(0\le\alpha<1).\label{eq:operators-photon-9}
\end{align}

\begin{proposition}[Analyticity and bounds of the symbol]\label{prop:analyticity-bounds-symbol}
Let $\mu,\nu\in\{0,1,2\}$, let $\boldsymbol{\Lambda}_{\mu\nu}$, $\boldsymbol{\omega}_{\nu}$ and $h_{\mu\nu}(\cdot,m)$, $m\in\mathfrak{M}$, be the functions defined in \eqref{eq:operators-photon-6}, and let $q_m=p+2\pi m$.
\begin{enumerate}
\item[(a)] On $\boldsymbol{\Omega}(\mathrm{ar}_2,g_2)$ one has $|N_{\mu}|\ge\frac{c_N}{2}$, $|N_G|\ge\frac{c_G}{2}$ and $|\Pi|\ge\frac{c_{\Pi}}{2}$. The functions $\boldsymbol{\Lambda}_{\mu\nu}$, $\boldsymbol{\omega}_{\nu}$ and $h_{\mu\nu}(\cdot,m)$, $m\in\mathfrak{M}$, are analytic on $\boldsymbol{\Omega}(\mathrm{ar}_2,g_2)$, and for every $p\in\boldsymbol{\Omega}(\mathrm{ar}_2,g_2)$
\begin{align}
|\boldsymbol{\Lambda}_{\mu\nu}(p)|&\le\mathsf{A}_1,\qquad|\boldsymbol{\omega}_{\nu}(p)|\le\mathsf{A}_2,\qquad|h_{\mu\nu}(p,0)|\le\mathsf{A}_3,\notag\\
|h_{\mu\nu}(p,m)|&\le\exp\left(\frac{1}{8}\right)\cdot c_u^4\cdot\mathsf{A}_1\cdot\langle m_{\mu}\rangle^{-1}\cdot\prod_{\sigma=0}^2\langle m_{\sigma}\rangle^{-1}\cdot|m|_2^{-2}\notag\\
&\qquad+12c_u^3\cdot\mathsf{A}_2\cdot\langle m_{\mu}\rangle\cdot\prod_{\sigma=0}^2\langle m_{\sigma}\rangle^{-1}\cdot|m|_2^{-4}\qquad(m\in\mathfrak{M},\ m\ne0).\label{eq:operators-photon-10}
\end{align}
\item[(b)] Let $0\le\alpha<1$. For every direction $\mathfrak{l}$ and every $p\in\boldsymbol{\Omega}(\mathrm{ar}_2,g_2)$,
\[
\sum_{m\in\mathfrak{M}}|h_{\mu\nu}(p,m)|\cdot\max\left\{1,\ |\hat{d}_{\mathfrak{l}}(q_m)|,\ |\hat{d}_{\mathfrak{l}}(q_m)|\cdot\left(\sum_{\sigma=0}^2|q_{m,\sigma}|\right)^{\alpha}\right\}\le S_{*}(\alpha).
\]
\end{enumerate}
\end{proposition}

\begin{proof}
(a) Let $(f,B,c)$ be one of the triples $(N_{\mu},B_N,c_N)$, $(N_G,B_G,c_G)$, $(\Pi,B_{\Pi},c_{\Pi})$. By Lemma~\ref{lem:estimates-aliases}(f), $f$ is analytic on $\boldsymbol{\Omega}(\mathrm{ar}_1,g_1)$ and $|f|\le B$ there; by Lemma~\ref{lem:estimates-aliases}(g), $f$ is real and $f\ge c$ on $[-\mathrm{ar}_2,\mathrm{ar}_2]^3\subset(-\mathrm{ar}_1,\mathrm{ar}_1)^3$. The definition of $g_2$ in \eqref{eq:operators-photon-8} gives $g_2\le\frac{g_1c}{4B}$, so that $\boldsymbol{\Omega}(\mathrm{ar}_2,g_2)\subset\boldsymbol{\Omega}\left(\mathrm{ar}_2,\frac{g_1c}{4B}\right)$, and Lemma~\ref{lem:lower-bounds-near} yields $|f|\ge\frac{c}{2}$ on $\boldsymbol{\Omega}(\mathrm{ar}_2,g_2)$. In particular $N_0$, $N_1$, $N_2$, $N_G$ and $\Pi$ have no zero in $\boldsymbol{\Omega}(\mathrm{ar}_2,g_2)$, and neither has $\hat{D}(q_m)$ for $m\ne0$, by Lemma~\ref{lem:estimates-aliases}(b). The functions in \eqref{eq:operators-photon-6} are quotients of analytic functions whose denominators are products of these functions; hence they are analytic on $\boldsymbol{\Omega}(\mathrm{ar}_2,g_2)$.

Let $p\in\boldsymbol{\Omega}(\mathrm{ar}_2,g_2)$ and $\mathcal{N}=N_0\cdot N_1\cdot N_2$. By $|N_{\mathfrak{l}}|\le B_N$ (Lemma~\ref{lem:estimates-aliases}(f)) and the lower bounds just proved,
\[
\left|\frac{\mathcal{N}}{N_{\mu}\cdot N_{\nu}\cdot\Pi}\right|\le B_N^3\cdot\left(\frac{2}{c_N}\right)^2\cdot\frac{2}{c_{\Pi}}=\frac{8B_N^3}{c_N^2\cdot c_{\Pi}},\qquad\left|\frac{\mathcal{N}}{N_{\nu}\cdot N_G\cdot\Pi}\right|\le\frac{8B_N^3}{c_N\cdot c_G\cdot c_{\Pi}},\qquad\left|\frac{\mathcal{N}}{N_{\mu}\cdot N_{\nu}\cdot N_G\cdot\Pi}\right|\le\frac{16B_N^3}{c_N^2\cdot c_G\cdot c_{\Pi}}.
\]
Lemma~\ref{lem:estimates-aliases} gives, in addition: $|D_0|\le B_D$ by (c); $|\rho_{\mu}(p)|\le c_{\rho}$ and $|\rho^{*}_{\nu}(p)|\le c_{\rho}$ by (d); $|U_0|\le c_{\mathrm{ch}}^3$ and $|u_{\mu}|\le c_{\mathrm{ch}}$ by (a) with $m=0$; $|\hat{d}_{\mu}(p)|\le12$ by (e) with $m=0$; $|X_{\mu}|\le B_X$ and $|Y|\le B_Y$ by (f). Finally, $\hat{V}^{*}_{\mu}(p)=\exp\left(-\frac{i}{2}(1-\eta)p_{\mu}\right)\cdot u_{\mu}$, where the exponential factor has modulus at most $\exp\left(\frac{1}{8}\right)$, as in the proof of Lemma~\ref{lem:estimates-aliases}(a); so $|\hat{V}^{*}_{\mu}(p)|\le\exp\left(\frac{1}{8}\right)\cdot c_{\mathrm{ch}}$. Inserting these bounds in the formulas \eqref{eq:operators-photon-6} for $\boldsymbol{\Lambda}_{\mu\nu}$, $\boldsymbol{\omega}_{\nu}$ and $h_{\mu\nu}(\cdot,0)$, one obtains
\begin{align*}
|\boldsymbol{\Lambda}_{\mu\nu}(p)|&\le\frac{|D_0|}{|N_{\mu}|}+|\rho_{\mu}|\cdot|\rho^{*}_{\nu}|\cdot|U_0|^2\cdot\left|\frac{\mathcal{N}}{N_{\mu}\cdot N_{\nu}\cdot\Pi}\right|\le\frac{2B_D}{c_N}+c_{\rho}^2\cdot c_{\mathrm{ch}}^6\cdot\frac{8B_N^3}{c_N^2\cdot c_{\Pi}}=\mathsf{A}_1,\\
|\boldsymbol{\omega}_{\nu}(p)|&\le|\rho^{*}_{\nu}|\cdot|D_0|\cdot|U_0|^2\cdot\left|\frac{\mathcal{N}}{N_{\nu}\cdot N_G\cdot\Pi}\right|\le c_{\rho}\cdot B_D\cdot c_{\mathrm{ch}}^6\cdot\frac{8B_N^3}{c_N\cdot c_G\cdot c_{\Pi}}=\mathsf{A}_2,\\
|h_{\mu\nu}(p,0)|&\le\frac{|U_0|\cdot|\hat{V}^{*}_{\mu}(p)|}{|N_{\mu}|}+|\hat{d}_{\mu}(p)|\cdot|\rho^{*}_{\nu}|\cdot|U_0|^3\cdot\left(|X_{\mu}|+|u_{\mu}|^2\cdot|D_0|\cdot|Y|\right)\cdot\left|\frac{\mathcal{N}}{N_{\mu}\cdot N_{\nu}\cdot N_G\cdot\Pi}\right|\\
&\le\frac{2\exp\left(\frac{1}{8}\right)\cdot c_{\mathrm{ch}}^4}{c_N}+12c_{\rho}\cdot c_{\mathrm{ch}}^9\cdot\left(B_X+c_{\mathrm{ch}}^2\cdot B_D\cdot B_Y\right)\cdot\frac{16B_N^3}{c_N^2\cdot c_G\cdot c_{\Pi}}=\mathsf{A}_3.
\end{align*}
Now let $m\ne0$. By Lemma~\ref{lem:estimates-aliases}(b), $|\hat{D}(q_m)|\ge\operatorname{Re}\hat{D}(q_m)\ge|m|_2^2$, and by Lemma~\ref{lem:estimates-aliases}(a), (e),
\[
|\hat{U}(q_m)|\le c_u^3\cdot\prod_{\sigma=0}^2\langle m_{\sigma}\rangle^{-1},\qquad|\hat{V}^{*}_{\mu}(q_m)|\le\exp\left(\frac{1}{8}\right)\cdot c_u\cdot\langle m_{\mu}\rangle^{-1},\qquad|\hat{d}_{\mu}(q_m)|\le12\langle m_{\mu}\rangle.
\]
The formula \eqref{eq:operators-photon-6} for $h_{\mu\nu}(p,m)$, together with these bounds and with $|\boldsymbol{\Lambda}_{\mu\nu}(p)|\le\mathsf{A}_1$, $|\boldsymbol{\omega}_{\nu}(p)|\le\mathsf{A}_2$, gives
\begin{align*}
|h_{\mu\nu}(p,m)|&\le\frac{|\hat{U}(q_m)|\cdot|\hat{V}^{*}_{\mu}(q_m)|\cdot|\boldsymbol{\Lambda}_{\mu\nu}(p)|}{|\hat{D}(q_m)|}+\frac{|\hat{d}_{\mu}(q_m)|\cdot|\hat{U}(q_m)|\cdot|\boldsymbol{\omega}_{\nu}(p)|}{|\hat{D}(q_m)|^2}\\
&\le c_u^3\cdot\prod_{\sigma=0}^2\langle m_{\sigma}\rangle^{-1}\cdot\exp\left(\frac{1}{8}\right)\cdot c_u\cdot\langle m_{\mu}\rangle^{-1}\cdot\mathsf{A}_1\cdot|m|_2^{-2}+12\langle m_{\mu}\rangle\cdot c_u^3\cdot\prod_{\sigma=0}^2\langle m_{\sigma}\rangle^{-1}\cdot\mathsf{A}_2\cdot|m|_2^{-4},
\end{align*}
which is \eqref{eq:operators-photon-10}.

(b) Let $m\in\mathfrak{M}$. By Lemma~\ref{lem:estimates-aliases}(e), $|\hat{d}_{\mathfrak{l}}(q_m)|\le12\langle m_{\mathfrak{l}}\rangle\le12\langle m\rangle_{\infty}$ and $\sum_{\sigma}|q_{m,\sigma}|\le31.5\langle m\rangle_{\infty}$. Since $\langle m\rangle_{\infty}\ge1$ and $0\le\alpha<1$, each of the three entries of the maximum is at most $12\langle m\rangle_{\infty}\cdot\left(31.5\langle m\rangle_{\infty}\right)^{\alpha}$, and
\[
\max\left\{1,\ |\hat{d}_{\mathfrak{l}}(q_m)|,\ |\hat{d}_{\mathfrak{l}}(q_m)|\cdot\left(\sum_{\sigma=0}^2|q_{m,\sigma}|\right)^{\alpha}\right\}\le12\langle m\rangle_{\infty}\cdot\left(31.5\langle m\rangle_{\infty}\right)^{\alpha}\le378\langle m\rangle_{\infty}^{1+\alpha}.
\]
For $m=0$ the right side equals $378$; by \eqref{eq:operators-photon-10}, the term $m=0$ of the sum in (b) is at most $378\mathsf{A}_3$.

Let $m\ne0$. Then $|m|_2\ge\max_{\sigma}|m_{\sigma}|\ge1$, and therefore $|m|_2\ge\langle m\rangle_{\infty}\ge\langle m_{\sigma}\rangle$ for $\sigma=0,1,2$. For the first summand on the right of \eqref{eq:operators-photon-10} this gives
\[
\langle m\rangle_{\infty}^{1+\alpha}\cdot|m|_2^{-2}\le\langle m\rangle_{\infty}^{-(1-\alpha)}\le\prod_{\sigma\ne\mu}\langle m_{\sigma}\rangle^{-\frac{1-\alpha}{2}},
\]
so that its product with $378\langle m\rangle_{\infty}^{1+\alpha}$ is at most
\[
378\exp\left(\frac{1}{8}\right)\cdot c_u^4\cdot\mathsf{A}_1\cdot\langle m_{\mu}\rangle^{-2}\cdot\prod_{\sigma\ne\mu}\langle m_{\sigma}\rangle^{-1-\frac{1-\alpha}{2}}.
\]
The sum of this over all $m\in\mathbb Z^3$ equals $378\exp\left(\frac{1}{8}\right)\cdot c_u^4\cdot\mathsf{A}_1\cdot Z_2\cdot Z_{1+\frac{1-\alpha}{2}}^2$, and $Z_{1+\frac{1-\alpha}{2}}\le3+\frac{4}{1-\alpha}$ by \eqref{eq:operators-photon-8}, applied with $a=1+\frac{1-\alpha}{2}>1$. For the second summand on the right of \eqref{eq:operators-photon-10},
\[
\langle m_{\mu}\rangle\cdot\langle m\rangle_{\infty}^{1+\alpha}\cdot|m|_2^{-4}\le\langle m\rangle_{\infty}^{-(2-\alpha)}\le\langle m\rangle_{\infty}^{-1}\le\prod_{\sigma=0}^2\langle m_{\sigma}\rangle^{-\frac{1}{3}},
\]
so that its product with $378\langle m\rangle_{\infty}^{1+\alpha}$ is at most $378\cdot12c_u^3\cdot\mathsf{A}_2\cdot\prod_{\sigma}\langle m_{\sigma}\rangle^{-\frac{4}{3}}$. The sum of this over all $m\in\mathbb Z^3$ equals $378\cdot12c_u^3\cdot\mathsf{A}_2\cdot Z_{\frac{4}{3}}^3$, and $Z_{\frac{4}{3}}\le3+6=9$ by \eqref{eq:operators-photon-8}. Since all terms are non-negative, the sum over $m\in\mathfrak{M}\setminus\{0\}$ is at most the sum over $\mathbb Z^3$. Adding the contribution of $m=0$ and the two contributions of the terms $m\ne0$, one obtains the bound
\[
378\left(\mathsf{A}_3+\exp\left(\frac{1}{8}\right)\cdot c_u^4\cdot\mathsf{A}_1\cdot Z_2\cdot\left(3+\frac{4}{1-\alpha}\right)^2+12\cdot9^3c_u^3\cdot\mathsf{A}_2\right)=S_{*}(\alpha)
\]
of \eqref{eq:operators-photon-9}.
\end{proof}

\subsubsection{Periodicity in the momentum variable}\label{sec:periodicity-momentum-variable}

For directions $\mu, \nu$, for $z \in \eta\mathbb{Z}^3$ and for $p \in \boldsymbol{\Omega}(\mathrm{ar}_2, g_2)$ let
\[
K_{\mu\nu}(p;z) := \sum_{m\in\mathfrak{M}}\exp\left(i\,q_m\cdot z\right)\cdot h_{\mu\nu}(p,m), \qquad q_m = p + 2\pi m,
\]
with the functions $h_{\mu\nu}(\cdot\,,m)$ of \eqref{eq:operators-photon-6}. Every $h_{\mu\nu}(\cdot\,,m)$, $m \in \mathfrak{M}$, is analytic on $\boldsymbol{\Omega}(\mathrm{ar}_2, g_2)$ (Proposition~\ref{prop:analyticity-bounds-symbol}(a)), and $p \mapsto \exp(i\,q_m\cdot z)$ is entire; hence $K_{\mu\nu}(\cdot\,;z)$ is an analytic function from $\boldsymbol{\Omega}(\mathrm{ar}_2, g_2)$ to $\mathbb{C}$.

\begin{lemma}[Periodicity]\label{lem:periodicity}
Let $\mathfrak{l} \in \{0,1,2\}$, and let $p \in \mathbb{C}^3$ be a point such that $p$ and $p + 2\pi e_{\mathfrak{l}}$ both lie in $\boldsymbol{\Omega}(\mathrm{ar}_2, g_2)$. Then for all directions $\mu, \nu$ and all $z \in \eta\mathbb{Z}^3$
\[
K_{\mu\nu}(p + 2\pi e_{\mathfrak{l}};z) = K_{\mu\nu}(p;z) \qquad\text{and}\qquad \boldsymbol{\Lambda}_{\mu\nu}(p + 2\pi e_{\mathfrak{l}}) = \boldsymbol{\Lambda}_{\mu\nu}(p).
\]
\end{lemma}

\begin{proof}
Fix $\mathfrak{l}$ and $z$, and let $G$ be the set of the points $p \in \mathbb{C}^3$ for which both $p$ and $p + 2\pi e_{\mathfrak{l}}$ belong to $\boldsymbol{\Omega}(\mathrm{ar}_2, g_2)$. As the intersection of two convex open subsets of $\mathbb{C}^3$, the set $G$ is convex and open, hence connected. A real point $p$ belongs to $G$ if and only if $-\mathrm{ar}_2 < p_{\mathfrak{l}} < \mathrm{ar}_2 - 2\pi$ and the two other coordinates of $p$ have modulus less than $\mathrm{ar}_2$. Since $\pi < \mathrm{ar}_2 < \frac{5\pi}{4}$ by \eqref{eq:operators-photon-8}, these real points form a non-empty open subset of $\mathbb{R}^3$.

Let $p$ be a real point of $G$, and let $p'$ stand for either of the two points $p$ and $p + 2\pi e_{\mathfrak{l}}$. The coordinates of $p'$ are real numbers of modulus less than $\mathrm{ar}_2$, and the coordinate of index $\mathfrak{l}$ is different from $0$, because $p_{\mathfrak{l}} < \mathrm{ar}_2 - 2\pi < 0$ and $p_{\mathfrak{l}} + 2\pi > 2\pi - \mathrm{ar}_2 > 0$. Consequently $D_0(p') > 0$: the number $D_0(p') = \hat{D}(p')$ is the sum of the three non-negative numbers $\operatorname{sn}(p'_0)^2$, $\operatorname{sn}(p'_1)^2$, $\operatorname{sn}(p'_2)^2$, and the one of index $\mathfrak{l}$ is positive, since $0 < |p'_{\mathfrak{l}}| < 2\pi \le 2\pi\mathfrak{N}$. Lemma~\ref{lem:one-complex-variable}(d), applied to the three coordinates of $p'$, gives $U_0(p') > 0$. For $m \in \mathfrak{M}\setminus\{0\}$ the alias $q_m = p' + 2\pi m$ satisfies $\operatorname{Re}\hat{D}(q_m) \ge |m|_2^2 > 0$ by Lemma~\ref{lem:estimates-aliases}(b) (note that $\boldsymbol{\Omega}(\mathrm{ar}_2, g_2) \subset \boldsymbol{\Omega}(\mathrm{ar}_1, g_1)$), so that $\hat{D}(q_m) \ne 0$; and $N_0\cdot N_1\cdot N_2\cdot N_G\cdot\Pi \ne 0$ at $p'$ by Proposition~\ref{prop:analyticity-bounds-symbol}(a). Thus the functions \eqref{eq:operators-photon-6} are defined at $p'$ (Section~\ref{sec:sums-aliases}), and Lemma~\ref{lem:algebraic-identities} applies at $p'$, where $D_0 \ne 0$.

At the point $p'$ define, for every direction $\mu$,
\[
E_{\mu} := \frac{N_{\mu}}{D_0}, \qquad G_0 := \frac{N_G}{D_0^2}, \qquad \mathrm{Sm} := \sum_{\mathfrak{l}'}\frac{\rho_{\mathfrak{l}'}\cdot\rho^{*}_{\mathfrak{l}'}}{E_{\mathfrak{l}'}}.
\]
Lemma~\ref{lem:algebraic-identities}(a) and $U_0 \ne 0$ give
\[
\mathrm{Sm} = D_0\cdot\sum_{\mathfrak{l}'}\frac{\rho_{\mathfrak{l}'}\cdot\rho^{*}_{\mathfrak{l}'}}{N_{\mathfrak{l}'}} = \frac{D_0^2\cdot\Pi}{U_0^2\cdot\mathcal{N}} \ne 0,
\]
and in these terms the definitions \eqref{eq:operators-photon-6} of $\boldsymbol{\Lambda}_{\mu\nu}$ and of $\boldsymbol{\omega}_{\nu}$ read
\[
\boldsymbol{\Lambda}_{\mu\nu} = \frac{\delta_{\mu\nu}}{E_{\mu}} - \frac{\rho_{\mu}\cdot\rho^{*}_{\nu}}{E_{\mu}\cdot E_{\nu}\cdot\mathrm{Sm}}, \qquad \boldsymbol{\omega}_{\nu} = \frac{\rho^{*}_{\nu}}{E_{\nu}\cdot G_0\cdot\mathrm{Sm}}.
\]
By Lemma~\ref{lem:algebraic-identities}(d), and by Lemma~\ref{lem:algebraic-identities}(c), according to which the formula \eqref{eq:operators-photon-6} for $h_{\mu\nu}(p',m)$, $m \ne 0$, holds for $m = 0$ as well, one has
\begin{align*}
E_{\mu} &= \sum_q\frac{\hat{U}(q)^2\cdot u(q_{\mu})^2}{\hat{D}(q)}, \qquad G_0 = \sum_q\frac{\hat{U}(q)^2}{\hat{D}(q)^2},\\
K_{\mu\nu}(p';z) &= \sum_q\exp\left(i\,q\cdot z\right)\cdot\left(\frac{\hat{U}(q)\cdot\hat{V}^{*}_{\mu}(q)\cdot\boldsymbol{\Lambda}_{\mu\nu}(p')}{\hat{D}(q)} + \frac{\hat{d}_{\mu}(q)\cdot\hat{U}(q)\cdot\boldsymbol{\omega}_{\nu}(p')}{\hat{D}(q)^2}\right),
\end{align*}
where $q$ ranges over the set $\{p' + 2\pi m : m \in \mathfrak{M}\}$. As $\mathfrak{M} = \{-\mathbf{h}, \dots, \mathbf{h}\}^3$ contains exactly one element of each class of $\mathbb{Z}^3$ modulo $\mathfrak{N}\mathbb{Z}^3$, this set contains exactly one element of each class of $p' + 2\pi\mathbb{Z}^3$ modulo $2\pi\mathfrak{N}\mathbb{Z}^3$. Every summand, as a function of $q$, has the period $2\pi\mathfrak{N}$ in each variable: the functions $\hat{U}$, $u$, $\hat{D}$, $\hat{V}^{*}_{\mu}$, $\hat{d}_{\mu}$ have this period (Section~\ref{sec:elementary-symbols-fourier}), and so has $q \mapsto \exp(i\,q\cdot z)$, because $\mathfrak{N}\cdot z \in \mathbb{Z}^3$. Therefore the three sums are sums over the classes of $p' + 2\pi\mathbb{Z}^3$ modulo $2\pi\mathfrak{N}\mathbb{Z}^3$; and these classes are the same for $p' = p$ and for $p' = p + 2\pi e_{\mathfrak{l}}$. The functions $\rho_{\nu}$, $\rho^{*}_{\nu}$ have the period $2\pi$ in each variable. It follows that $E_{\mu}$, $G_0$, $\rho_{\nu}$, $\rho^{*}_{\nu}$, $\mathrm{Sm}$, and with them $\boldsymbol{\Lambda}_{\mu\nu}$, $\boldsymbol{\omega}_{\nu}$ and $K_{\mu\nu}(\cdot\,;z)$, take the same values at $p$ and at $p + 2\pi e_{\mathfrak{l}}$. So both identities of the lemma hold at every real point of $G$.

The functions $p \mapsto K_{\mu\nu}(p + 2\pi e_{\mathfrak{l}};z) - K_{\mu\nu}(p;z)$ and $p \mapsto \boldsymbol{\Lambda}_{\mu\nu}(p + 2\pi e_{\mathfrak{l}}) - \boldsymbol{\Lambda}_{\mu\nu}(p)$ are analytic on $G$, since $K_{\mu\nu}(\cdot\,;z)$ and $\boldsymbol{\Lambda}_{\mu\nu}$ are analytic on $\boldsymbol{\Omega}(\mathrm{ar}_2, g_2)$ (Proposition~\ref{prop:analyticity-bounds-symbol}(a)). Let $f$ be either of them. By the preceding paragraph, $f = 0$ on $G \cap \mathbb{R}^3$, a non-empty open subset of $\mathbb{R}^3$. At a point of $G \cap \mathbb{R}^3$ every complex partial derivative of $f$, of any order, coincides with the derivative of the same order along the real coordinate axes, and the latter is determined by the values of $f$ on $G \cap \mathbb{R}^3$, which are zero; so $f$ and all its partial derivatives vanish at the real points of $G$. The set of the points of $G$ at which $f$ and all its partial derivatives vanish is therefore non-empty. It is closed in $G$, since these derivatives are continuous, and it is open, since near each point of $G$ the function $f$ is the sum of its Taylor series at that point. As $G$ is connected, this set is all of $G$; thus $f = 0$ on $G$ (identity theorem for analytic functions of several variables), and both identities of the lemma hold on $G$.
\end{proof}

\subsubsection{Exponential decay of the kernels}\label{sec:exponential-decay-kernels}

The operator $\Delta_k = (dH_k)^TdH_k$ of Proposition~\ref{prop:symbols} maps real bond functions on $T^0_k$ to real bond functions on $T^0_k$, since $H_k$ is real (Proposition~\ref{prop:minimiser}). Its kernel is defined by the formula of Section~\ref{sec:tori-cochains-kernels}, with positively oriented unit bonds in both places:
\[
(\Delta_kA)(b'') = \sum_{b'}\Delta_k(b'',b')\cdot A(b')
\]
for every real bond function $A$ on $T^0_k$ and every positively oriented unit bond $b''$, where $b'$ ranges over the positively oriented unit bonds.

\begin{theorem}[Decay of the kernels of $H_k$ and $\Delta_k$]\label{thm:decay-kernels}
Let $0 \le \alpha < 1$, and put
\[
C_{\mathrm{Hk}}(\alpha) := 2.5\cdot\left(3 + \frac{2}{\gamma_H}\right)^3\cdot S_{*}(\alpha), \qquad C_{\mathrm{De}} := \left(3 + \frac{2}{\gamma_H}\right)^3\cdot\mathsf{A}_1,
\]
with $\gamma_H$ of \eqref{eq:operators-photon-8} and $S_{*}(\alpha)$, $\mathsf{A}_1$ of \eqref{eq:operators-photon-9}. Let $\mathfrak{N} \ge 1$ be an odd integer (this includes $\mathfrak{N} = L^k$ for every integer $k \ge 0$ and every odd integer $L \ge 3$), and let $l$ be a positive integer. Let $b$, $\tilde{b}$ be fine bonds of $T_k$ of the same direction whose initial points $x$, $\tilde{x}$ satisfy $0 < |x - \tilde{x}| \le 1$, and let $b'$, $b''$ be unit bonds, all bonds being positively oriented. Then for $\mathfrak{l} = 0, 1, 2$ the kernels of $H_k$ (Section~\ref{sec:tori-cochains-kernels}) and the kernel of $\Delta_k$ satisfy
\begin{align}
&\left|H_k(b,b')\right|,\ \left|(\partial_{\mathfrak{l}}H_k)(b,b')\right|,\ \left|(\delta_{\alpha}\partial_{\mathfrak{l}}H_k)((b,\tilde{b});b')\right| \notag\\
&\qquad\qquad \le C_{\mathrm{Hk}}(\alpha)\cdot\exp\left(-\gamma_H\cdot d(b,b')\right), \notag\\
&\left|\Delta_k(b'',b')\right| \le C_{\mathrm{De}}\cdot\exp\left(-\gamma_H\cdot d(b'',b')\right). \label{eq:operators-photon-11}
\end{align}
Moreover, the norm of $\Delta_k$ as an operator on $\ell^2$, that is, on the real bond functions on $T^0_k$ with the norm $|\cdot|_2$, is at most $3\,\mathsf{A}_1$.
\end{theorem}

\begin{proof}
\emph{Step 1 (the kernel of $H_k$ as a sum over $P_l$).} Let $b = (x, x + \eta e_{\mu})$ and $b' = (y', y' + e_{\nu})$, and let $\delta_{b'}$ be the real bond function on $T^0_k$ that equals $1$ on $b'$ and $0$ on the other positively oriented unit bonds; then $H_k(b,b') = (H_k\delta_{b'})_{\mu}(x)$. The transform of a bond function is taken component by component, each component being a function of the initial point of the bond (Section~\ref{sec:elementary-symbols-fourier}); so the transform of $A := \delta_{b'}$ is
\[
\hat{A}_{\mathfrak{l}}(p) = \delta_{\mathfrak{l}\nu}\cdot\exp\left(-i\,p\cdot y'\right), \qquad p \in P_l,\ \mathfrak{l} = 0, 1, 2.
\]
Represent $x$ by a point of $\eta\mathbb{Z}^3$ and $y'$ by a point of $\mathbb{Z}^3$, denoted by the same letters, and let $z := x - y' \in \eta\mathbb{Z}^3$; such a $z$ is a lift of $x - y'$, and every lift of $x - y'$ is of the form $z + lw$, $w \in \mathbb{Z}^3$. For $p \in P_l$, represented in $[-\pi,\pi)^3$, and $m \in \mathfrak{M}$ one has $q_m - p = 2\pi m$ and $m\cdot y' \in \mathbb{Z}$, hence
\[
\exp\left(i\,q_m\cdot x\right)\cdot\exp\left(-i\,p\cdot y'\right) = \exp\left(i\,q_m\cdot z\right)\cdot\exp\left(2\pi i\,m\cdot y'\right) = \exp\left(i\,q_m\cdot z\right).
\]
The inverse Fourier transform on $T_k$ (Section~\ref{sec:elementary-symbols-fourier}) and \eqref{eq:operators-photon-7} therefore give
\begin{align*}
H_k(b,b') &= l^{-3}\sum_{p\in P_l}\sum_{m\in\mathfrak{M}}\exp\left(i\,q_m\cdot x\right)\cdot h_{\mu\nu}(p,m)\cdot\exp\left(-i\,p\cdot y'\right)\\
&= l^{-3}\sum_{p\in P_l}K_{\mu\nu}(p;z).
\end{align*}
The choice of the lift does not matter: for $w \in \mathbb{Z}^3$ and $p \in P_l$ one has $\exp(2\pi i\,l\,m\cdot w) = 1$ and $l\,p \in 2\pi\mathbb{Z}^3$, so that
\[
K_{\mu\nu}(p;z + lw) = \exp\left(i\,l\,p\cdot w\right)\cdot K_{\mu\nu}(p;z) = K_{\mu\nu}(p;z).
\]

\emph{Step 2 (the kernel in infinite volume).} The set of the points $p \in \mathbb{C}^3$ with $|\operatorname{Im}p_j| < g_2$ for $j = 0, 1, 2$ is called the strip. For every $z \in \eta\mathbb{Z}^3$ the function $K_{\mu\nu}(\cdot\,;z)$ has an extension to the strip which is analytic and has the period $2\pi$ in each variable; it is constructed as follows. Let $p$ lie in the strip. A vector $w \in \mathbb{Z}^3$ is called admissible at $p$ if $p - 2\pi w \in \boldsymbol{\Omega}(\mathrm{ar}_2, g_2)$. Such a vector exists: choose $w_j$ with $\operatorname{Re}p_j - 2\pi w_j \in [-\pi,\pi)$ and recall $\mathrm{ar}_2 > \pi$. Put $K_{\mu\nu}(p;z) := K_{\mu\nu}(p - 2\pi w;z)$ with a vector $w$ admissible at $p$. If $w'$ is a second vector admissible at $p$, then $2\pi|w_j - w'_j| < 2\mathrm{ar}_2 < 4\pi$, so that $|w_j - w'_j| \le 1$ for $j = 0, 1, 2$. Replace the coordinates of $w$ by those of $w'$ one at a time. The set $\boldsymbol{\Omega}(\mathrm{ar}_2, g_2)$ is the cartesian product of three copies of the rectangle $\{s \in \mathbb{C} : |\operatorname{Re}s| < \mathrm{ar}_2,\ |\operatorname{Im}s| < g_2\}$; hence every intermediate point $p - 2\pi w''$ lies in $\boldsymbol{\Omega}(\mathrm{ar}_2, g_2)$, and two consecutive such points coincide or differ by $\pm 2\pi e_{\mathfrak{l}}$ for some $\mathfrak{l}$. By Lemma~\ref{lem:periodicity} no replacement changes the value of $K_{\mu\nu}$, so the definition does not depend on the choice of $w$. At points of $\boldsymbol{\Omega}(\mathrm{ar}_2, g_2)$ the vector $w = 0$ is admissible, so the extension coincides with $K_{\mu\nu}(\cdot\,;z)$ there. A vector $w$ that is admissible at $p$ is admissible on a neighbourhood of $p$, where the extension equals the analytic function $\tilde{p} \mapsto K_{\mu\nu}(\tilde{p} - 2\pi w;z)$; so the extension is analytic. If $w$ is admissible at $p$, then $w + e_{\mathfrak{l}}$ is admissible at $p + 2\pi e_{\mathfrak{l}}$; so the extension has the period $2\pi$ in each variable. The extension is again denoted by $K_{\mu\nu}(\cdot\,;z)$.

For $v \in \mathbb{Z}^3$ one has $\exp(i\,q_m\cdot v) = \exp(i\,p\cdot v)$; hence $K_{\mu\nu}(p;z + v) = \exp(i\,p\cdot v)\cdot K_{\mu\nu}(p;z)$ on $\boldsymbol{\Omega}(\mathrm{ar}_2, g_2)$, and then, by the construction of the extension, on the strip. Define $H^{\infty}_{\mu\nu} : \eta\mathbb{Z}^3 \to \mathbb{C}$ by
\[
H^{\infty}_{\mu\nu}(z) := (2\pi)^{-3}\int_{[-\pi,\pi]^3}\mathrm{d}p\,K_{\mu\nu}(p;z).
\]
The restriction of $K_{\mu\nu}(\cdot\,;z)$ to $\mathbb{R}^3$ is smooth and has the period $2\pi$ in each variable, and its Fourier coefficients are, for $v \in \mathbb{Z}^3$,
\begin{align*}
(2\pi)^{-3}\int_{[-\pi,\pi]^3}\mathrm{d}p\,K_{\mu\nu}(p;z)\cdot\exp\left(i\,p\cdot v\right) &= (2\pi)^{-3}\int_{[-\pi,\pi]^3}\mathrm{d}p\,K_{\mu\nu}(p;z + v)\\
&= H^{\infty}_{\mu\nu}(z + v).
\end{align*}
The Fourier series of a smooth periodic function converges absolutely. Hence
\[
K_{\mu\nu}(p;z) = \sum_{v\in\mathbb{Z}^3}H^{\infty}_{\mu\nu}(z + v)\cdot\exp\left(-i\,p\cdot v\right) \qquad (p \in \mathbb{R}^3),
\]
the series being absolutely convergent. Insert this, for the lift $z$ of $x - y'$, into the formula of Step~1, and note that $l^{-3}\sum_{p\in P_l}\exp(-i\,p\cdot v)$ equals $1$ for $v \in l\mathbb{Z}^3$ and $0$ for the other $v \in \mathbb{Z}^3$ (orthogonality of the characters of $\mathbb{Z}^3/l\mathbb{Z}^3$). It follows that
\begin{equation}\label{eq:operators-photon-12}
H_k(b,b') = \sum_{w\in\mathbb{Z}^3}H^{\infty}_{\mu\nu}(z + lw).
\end{equation}

\emph{Step 3 (exponential decay of $H^{\infty}_{\mu\nu}$).} Let $\boldsymbol{\zeta} \in \{-1,1\}^3$ and $\hat{z} \in \eta\mathbb{Z}^3$. For $j = 0, 1, 2$, one after the other, apply Cauchy's theorem in the variable $p_j$, the other variables being kept fixed, to the function $K_{\mu\nu}(\cdot\,;\hat{z})$ and to the rectangle with the corners $-\pi$, $\pi$, $\pi + i\gamma_H\cdot\boldsymbol{\zeta}_j$, $-\pi + i\gamma_H\cdot\boldsymbol{\zeta}_j$. All points that occur have imaginary parts of modulus at most $\gamma_H = g_2/2 < g_2$, so they lie in the strip. The contributions of the two vertical sides cancel, since the integrand has the period $2\pi$ in $p_j$. Consequently, for all $\boldsymbol{\zeta} \in \{-1,1\}^3$ and all $\hat{z} \in \eta\mathbb{Z}^3$,
\[
H^{\infty}_{\mu\nu}(\hat{z}) = (2\pi)^{-3}\int_{[-\pi,\pi]^3}\mathrm{d}p\,K_{\mu\nu}\left(p + i\gamma_H\cdot\boldsymbol{\zeta};\hat{z}\right).
\]
For $p \in [-\pi,\pi]^3$ the point $p + i\gamma_H\cdot\boldsymbol{\zeta}$ lies in $\boldsymbol{\Omega}(\mathrm{ar}_2, g_2)$, where $K_{\mu\nu}$ is given by its defining sum; the aliases of this point are $q_m + i\gamma_H\cdot\boldsymbol{\zeta}$, $m \in \mathfrak{M}$, with the real vectors $q_m = p + 2\pi m$.

Now fix $z \in \eta\mathbb{Z}^3$ and let
\[
\boldsymbol{\zeta}_j := \begin{cases} 1 & \text{if } z_j \ge 0,\\ -1 & \text{if } z_j < 0, \end{cases} \qquad j = 0, 1, 2.
\]
In the formula for $H^{\infty}_{\mu\nu}(\hat{z})$ take this $\boldsymbol{\zeta}$ for each of the points $\hat{z} = z$, $z + \eta e_{\mathfrak{l}}$, $z + r$, $z + r + \eta e_{\mathfrak{l}}$, where $\mathfrak{l} \in \{0,1,2\}$ and $r \in \eta\mathbb{Z}^3$ with $0 < |r| \le 1$. In what follows $p \in [-\pi,\pi]^3$, $m \in \mathfrak{M}$, and $q := q_m + i\gamma_H\cdot\boldsymbol{\zeta}$ is the alias of index $m$ of the point $p + i\gamma_H\cdot\boldsymbol{\zeta}$. Then
\[
\left|\exp\left(i\,q\cdot z\right)\right| = \exp\left(-\gamma_H\cdot\boldsymbol{\zeta}\cdot z\right) = \exp\left(-\gamma_H\cdot|z|_1\right).
\]
The difference quotient $(\partial_{\mathfrak{l}}f)(\hat{z}) := \eta^{-1}\left(f(\hat{z} + \eta e_{\mathfrak{l}}) - f(\hat{z})\right)$ of a function $f$ on $\eta\mathbb{Z}^3$ (as in Section~\ref{sec:tori-cochains-kernels}) multiplies the function $\hat{z} \mapsto \exp(i\,q\cdot\hat{z})$ by $\hat{d}_{\mathfrak{l}}(q)$, and the translation by $r$ multiplies it by $\exp(i\,q\cdot r)$, where
\begin{align*}
\left|\exp\left(i\,q\cdot r\right) - 1\right| &\le \min\left\{1 + \exp\left(3\gamma_H\right),\ |q\cdot r|\cdot\exp\left(3\gamma_H\right)\right\}\\
&\le 2.5\cdot\min\left\{1,\ |r|\cdot\sum_j|q_j|\right\} \le 2.5\cdot|r|^{\alpha}\cdot\left(\sum_j|q_j|\right)^{\alpha}.
\end{align*}
(The imaginary part of $q\cdot r$ is $\gamma_H\cdot\boldsymbol{\zeta}\cdot r$, of modulus at most $3\gamma_H\cdot|r| \le 3\gamma_H$; for $s \in \mathbb{C}$ one has $|\exp(i\,s) - 1| \le |s|\cdot\exp(|\operatorname{Im}s|)$; by \eqref{eq:operators-photon-8}, $3\gamma_H \le \frac{3}{8}$, and $1 + \exp\left(\frac{3}{8}\right) < 2.5$; finally, $\min\{1,t\} \le t^{\alpha}$ for $t \ge 0$.) Accordingly,
\begin{align*}
H^{\infty}_{\mu\nu}(z) &= (2\pi)^{-3}\int_{[-\pi,\pi]^3}\mathrm{d}p\,\sum_{m\in\mathfrak{M}}\exp\left(i\,q\cdot z\right)\cdot h_{\mu\nu}\left(p + i\gamma_H\cdot\boldsymbol{\zeta}, m\right),\\
(\partial_{\mathfrak{l}}H^{\infty}_{\mu\nu})(z) &= (2\pi)^{-3}\int_{[-\pi,\pi]^3}\mathrm{d}p\,\sum_{m\in\mathfrak{M}}\exp\left(i\,q\cdot z\right)\cdot\hat{d}_{\mathfrak{l}}(q)\cdot h_{\mu\nu}\left(p + i\gamma_H\cdot\boldsymbol{\zeta}, m\right),
\end{align*}
and $(\partial_{\mathfrak{l}}H^{\infty}_{\mu\nu})(z + r) - (\partial_{\mathfrak{l}}H^{\infty}_{\mu\nu})(z)$ is the second integral with the additional factor $\exp(i\,q\cdot r) - 1$ in each summand. Proposition~\ref{prop:analyticity-bounds-symbol}(b), applied at the point $p + i\gamma_H\cdot\boldsymbol{\zeta} \in \boldsymbol{\Omega}(\mathrm{ar}_2, g_2)$, gives
\begin{align*}
&\sum_{m\in\mathfrak{M}}\left|h_{\mu\nu}\left(p + i\gamma_H\cdot\boldsymbol{\zeta}, m\right)\right|\cdot\max\left\{1,\ \left|\hat{d}_{\mathfrak{l}}(q)\right|,\ \left|\hat{d}_{\mathfrak{l}}(q)\right|\cdot\left(\sum_j|q_j|\right)^{\alpha}\right\}\\
&\qquad\qquad \le S_{*}(\alpha).
\end{align*}
Together with the two preceding bounds this yields, for all $z \in \eta\mathbb{Z}^3$, all $\mathfrak{l} \in \{0,1,2\}$ and all $r \in \eta\mathbb{Z}^3$ with $0 < |r| \le 1$,
\begin{multline}\label{eq:operators-photon-13}
\left|H^{\infty}_{\mu\nu}(z)\right|,\ \left|(\partial_{\mathfrak{l}}H^{\infty}_{\mu\nu})(z)\right|,\ |r|^{-\alpha}\cdot\left|(\partial_{\mathfrak{l}}H^{\infty}_{\mu\nu})(z + r) - (\partial_{\mathfrak{l}}H^{\infty}_{\mu\nu})(z)\right| \\
\le 2.5\cdot S_{*}(\alpha)\cdot\exp\left(-\gamma_H\cdot|z|_1\right).
\end{multline}

\emph{Step 4 (summation over the lifts).} Choose the lift $z$ of $x - y'$ with $|z_j| \le l/2$ for $j = 0, 1, 2$. For $w \in \mathbb{Z}\setminus\{0\}$ one has $|z_j + lw| \ge l|w| - |z_j| \ge |z_j| + l\left(|w| - 1\right)$. With the inequality $\frac{1}{1 - \exp(-t)} \le 1 + \frac{1}{t}$ for $t > 0$, and with $\gamma_H\cdot l \ge \gamma_H$, it follows that
\begin{align*}
\sum_{w\in\mathbb{Z}}\exp\left(-\gamma_H\cdot|z_j + lw|\right) &\le \exp\left(-\gamma_H\cdot|z_j|\right)\cdot\left(1 + \frac{2}{1 - \exp\left(-\gamma_H\cdot l\right)}\right)\\
&\le \left(3 + \frac{2}{\gamma_H}\right)\cdot\exp\left(-\gamma_H\cdot|z_j|\right),
\end{align*}
and the product over $j = 0, 1, 2$ gives
\[
\sum_{w\in\mathbb{Z}^3}\exp\left(-\gamma_H\cdot|z + lw|_1\right) \le \left(3 + \frac{2}{\gamma_H}\right)^3\cdot\exp\left(-\gamma_H\cdot|z|_1\right).
\]
Moreover $|z|_1 \ge |z| = d(x,y') \ge d(b,b')$, because the coordinates of $z$ have modulus at most $l/2$ and $x$, $y'$ are end points of $b$, $b'$. By \eqref{eq:operators-photon-12} and the first bound of \eqref{eq:operators-photon-13} at the points $z + lw$,
\begin{align*}
\left|H_k(b,b')\right| &\le 2.5\cdot S_{*}(\alpha)\cdot\sum_{w\in\mathbb{Z}^3}\exp\left(-\gamma_H\cdot|z + lw|_1\right)\\
&\le C_{\mathrm{Hk}}(\alpha)\cdot\exp\left(-\gamma_H\cdot d(b,b')\right).
\end{align*}
The bond $b + \eta e_{\mathfrak{l}}$ has the initial point $x + \eta e_{\mathfrak{l}}$, and $z + \eta e_{\mathfrak{l}}$ is a lift of $x + \eta e_{\mathfrak{l}} - y'$. Thus \eqref{eq:operators-photon-12}, applied to $b$ and to $b + \eta e_{\mathfrak{l}}$, gives $(\partial_{\mathfrak{l}}H_k)(b,b') = \sum_{w\in\mathbb{Z}^3}(\partial_{\mathfrak{l}}H^{\infty}_{\mu\nu})(z + lw)$, and the second bound of \eqref{eq:operators-photon-13}, summed in the same way, gives the bound for $(\partial_{\mathfrak{l}}H_k)(b,b')$. Let $r \in \eta\mathbb{Z}^3$ be a lift of $\tilde{x} - x$ with $|r| = |x - \tilde{x}|$, so that $0 < |r| \le 1$. Then $\tilde{z} := z + r$ is a lift of $\tilde{x} - y'$, and in the same way $(\partial_{\mathfrak{l}}H_k)(\tilde{b},b') = \sum_{w\in\mathbb{Z}^3}(\partial_{\mathfrak{l}}H^{\infty}_{\mu\nu})(z + lw + r)$. Hence
\begin{align*}
&(\delta_{\alpha}\partial_{\mathfrak{l}}H_k)((b,\tilde{b});b')\\
&\qquad\qquad = |r|^{-\alpha}\cdot\sum_{w\in\mathbb{Z}^3}\left((\partial_{\mathfrak{l}}H^{\infty}_{\mu\nu})(z + lw + r) - (\partial_{\mathfrak{l}}H^{\infty}_{\mu\nu})(z + lw)\right),
\end{align*}
and the third bound of \eqref{eq:operators-photon-13}, applied at the points $z + lw$ and summed in the same way, gives the bound \eqref{eq:operators-photon-11} for $(\delta_{\alpha}\partial_{\mathfrak{l}}H_k)((b,\tilde{b});b')$.

\emph{Step 5 (the operator $\Delta_k$).} Let $b'' = (y'', y'' + e_{\mu})$ and $b' = (y', y' + e_{\nu})$, and let $\delta_{b'}$ be the real bond function of Step~1. By the definition of the kernel of $\Delta_k$, $\Delta_k(b'',b') = (\Delta_k\delta_{b'})(b'') = (\Delta_k\delta_{b'})_{\mu}(y'')$. By \eqref{eq:operators-photon-7} the $\mu$-th component of the transform of $\Delta_k\delta_{b'}$ is the function $p \mapsto \boldsymbol{\Lambda}_{\mu\nu}(p)\cdot\exp(-i\,p\cdot y')$ on $P_l$, and the inverse Fourier transform on $T^0_k$ gives
\[
\Delta_k(b'',b') = l^{-3}\sum_{p\in P_l}\exp\left(i\,p\cdot(y'' - y')\right)\cdot\boldsymbol{\Lambda}_{\mu\nu}(p).
\]
By Proposition~\ref{prop:analyticity-bounds-symbol}(a) the function $\boldsymbol{\Lambda}_{\mu\nu}$ is analytic on $\boldsymbol{\Omega}(\mathrm{ar}_2, g_2)$ and bounded there by $\mathsf{A}_1$; by Lemma~\ref{lem:periodicity}, $\boldsymbol{\Lambda}_{\mu\nu}(p + 2\pi e_{\mathfrak{l}}) = \boldsymbol{\Lambda}_{\mu\nu}(p)$ whenever $p$ and $p + 2\pi e_{\mathfrak{l}}$ lie in $\boldsymbol{\Omega}(\mathrm{ar}_2, g_2)$. As in Step~2, $\boldsymbol{\Lambda}_{\mu\nu}$ extends to an analytic function on the strip with the period $2\pi$ in each variable; the extension is bounded by $\mathsf{A}_1$, because its values are values of $\boldsymbol{\Lambda}_{\mu\nu}$ on $\boldsymbol{\Omega}(\mathrm{ar}_2, g_2)$. Define $D^{\infty}_{\mu\nu} : \mathbb{Z}^3 \to \mathbb{C}$ by
\[
D^{\infty}_{\mu\nu}(z) := (2\pi)^{-3}\int_{[-\pi,\pi]^3}\mathrm{d}p\,\exp\left(i\,p\cdot z\right)\cdot\boldsymbol{\Lambda}_{\mu\nu}(p).
\]
These numbers are the Fourier coefficients of the restriction of $\boldsymbol{\Lambda}_{\mu\nu}$ to $\mathbb{R}^3$, a smooth periodic function; its Fourier series converges absolutely, and $\boldsymbol{\Lambda}_{\mu\nu}(p) = \sum_{v\in\mathbb{Z}^3}D^{\infty}_{\mu\nu}(v)\cdot\exp(-i\,p\cdot v)$ for $p \in \mathbb{R}^3$. As in Step~2 it follows that
\[
\Delta_k(b'',b') = \sum_{w\in\mathbb{Z}^3}D^{\infty}_{\mu\nu}(z + lw)
\]
for every lift $z \in \mathbb{Z}^3$ of $y'' - y'$. The shift of contours of Step~3, applied to the analytic periodic function $p \mapsto \exp(i\,p\cdot z)\cdot\boldsymbol{\Lambda}_{\mu\nu}(p)$ with $\boldsymbol{\zeta}_j := 1$ for $z_j \ge 0$ and $\boldsymbol{\zeta}_j := -1$ for $z_j < 0$, gives
\[
D^{\infty}_{\mu\nu}(z) = (2\pi)^{-3}\int_{[-\pi,\pi]^3}\mathrm{d}p\,\exp\left(i\,(p + i\gamma_H\cdot\boldsymbol{\zeta})\cdot z\right)\cdot\boldsymbol{\Lambda}_{\mu\nu}\left(p + i\gamma_H\cdot\boldsymbol{\zeta}\right),
\]
and hence $|D^{\infty}_{\mu\nu}(z)| \le \mathsf{A}_1\cdot\exp(-\gamma_H\cdot|z|_1)$. Take the lift $z$ with $|z_j| \le l/2$ for $j = 0, 1, 2$; then $|z|_1 \ge |z| = d(y'',y') \ge d(b'',b')$, and the summation of Step~4 gives
\begin{align*}
\left|\Delta_k(b'',b')\right| &\le \left(3 + \frac{2}{\gamma_H}\right)^3\cdot\mathsf{A}_1\cdot\exp\left(-\gamma_H\cdot d(b'',b')\right)\\
&= C_{\mathrm{De}}\cdot\exp\left(-\gamma_H\cdot d(b'',b')\right).
\end{align*}
This is the last line of \eqref{eq:operators-photon-11}.

Finally, let $A$ be a real bond function on $T^0_k$; for $p \in P_l$ let $\hat{A}(p)$ be the vector with the components $\hat{A}_{\nu}(p)$, and for $p \in \mathbb{R}^3$ let $\boldsymbol{\Lambda}(p)$ be the $3\times 3$ matrix with the entries $\boldsymbol{\Lambda}_{\mu\nu}(p)$ (the extension of $\boldsymbol{\Lambda}_{\mu\nu}$ being used where $p \notin \boldsymbol{\Omega}(\mathrm{ar}_2, g_2)$). By \eqref{eq:operators-photon-7} the transform of $\Delta_kA$ at $p \in P_l$ is the vector $\boldsymbol{\Lambda}(p)\hat{A}(p)$. Parseval's formula on $T^0_k$ gives $|A|_2^2 = l^{-3}\sum_{p\in P_l}\sum_{\nu}|\hat{A}_{\nu}(p)|^2$, and the same formula for the real bond function $\Delta_kA$; hence
\[
|\Delta_kA|_2 \le \sup_{p\in\mathbb{R}^3}\left\|\boldsymbol{\Lambda}(p)\right\|\cdot|A|_2,
\]
where $\|\cdot\|$ denotes the operator norm of $3\times 3$ matrices. For real $p$ every entry of $\boldsymbol{\Lambda}(p)$ has modulus at most $\mathsf{A}_1$, so that $\|\boldsymbol{\Lambda}(p)\|$ does not exceed the Hilbert--Schmidt norm $\left(\sum_{\mu,\nu}|\boldsymbol{\Lambda}_{\mu\nu}(p)|^2\right)^{\frac{1}{2}} \le 3\,\mathsf{A}_1$.
\end{proof}

\begin{remark}\label{rem:operators-photon-2}
\begin{enumerate}
\item[(i)] By \eqref{eq:operators-photon-9}, $C_{\mathrm{Hk}}(\alpha)$ is of the order of $(1-\alpha)^{-2}$ as $\alpha \to 1$, whereas $\gamma_H$ is a number which depends on none of $\alpha$, $\mathfrak{N}$, $l$, $L$.
\item[(ii)] Let $\mathfrak{N} = 1$. Then $\mathfrak{M} = \{0\}$, $u = 1$, $X_{\mu} = Y = 0$, $N_{\mu} = N_G = \Pi = 1$, $h_{\mu\nu}(p,0) = \delta_{\mu\nu}$ and $\boldsymbol{\Lambda}_{\mu\nu}(p) = \delta_{\mu\nu}\cdot D_0 - \rho_{\mu}\cdot\rho^{*}_{\nu}$ (see \eqref{eq:operators-photon-5}, \eqref{eq:operators-photon-6}). By \eqref{eq:operators-photon-7} it follows that $H_0 = 1$ and $\Delta_0 = \delta d$.
\item[(iii)] By \eqref{eq:operators-photon-12}, the kernel of $H_k$ on the torus of side $l$ is the periodisation, with respect to $l\mathbb{Z}^3$, of the kernel $(x,y') \mapsto H^{\infty}_{\mu\nu}(x - y')$ on $\eta\mathbb{Z}^3\times\mathbb{Z}^3$, and the function $H^{\infty}_{\mu\nu}$ does not depend on $l$.
\item[(iv)] Let $f$ be a bond function on $T^0_k$. Summation over $b'$ of the bounds \eqref{eq:operators-photon-11} for the three kernels of $H_k$, with \eqref{eq:operators-photon-2} at $g = \gamma_H \in (0,1]$ and the number $c_{\mathrm{bd}}(\gamma_H)$ of \eqref{eq:operators-photon-1}, gives
\[
|H_kf|_{\infty},\ |\partial H_kf|_{\infty},\ |\delta_{\alpha}\partial H_kf|_{\infty} \le C_{\mathrm{Hk}}(\alpha)\cdot c_{\mathrm{bd}}(\gamma_H)\cdot|f|_{\infty};
\]
here the three norms on the left are the suprema of $|(H_kf)(b)|$, of $|(\partial_{\mathfrak{l}}H_kf)(b)|$ and of $|x - \tilde{x}|^{-\alpha}\cdot|(\partial_{\mathfrak{l}}H_kf)(\tilde{b}) - (\partial_{\mathfrak{l}}H_kf)(b)|$ over $\mathfrak{l} \in \{0,1,2\}$ and over the fine bonds $b$, $\tilde{b}$ as in Theorem~\ref{thm:decay-kernels} (Section~\ref{sec:tori-cochains-kernels}).
\end{enumerate}
\end{remark}

\subsection{Axial gauges}\label{sec:axial-gauges}

\subsubsection{Blocks of odd side and weighted axial functionals}\label{sec:blocks-odd-side}

Let $n_b\ge 1$ be an odd integer and put $h_b:=\frac{n_b-1}{2}$. The underlying lattice is $\mathbb{Z}^3$, and a bond function is a function $A$ on its oriented bonds with $A(x',x)=-A(x,x')$. Let $l$ be a multiple of $n_b$, the case $l=n_b$ included. Through its periodic lift, every bond function on the torus $(\mathbb{Z}/l\mathbb{Z})^3$ is regarded as a bond function on $\mathbb{Z}^3$. Lifting commutes with the coboundary $d$, it preserves sup norms and line sums, and it is compatible with all objects introduced in this subsection. The blocks are the sets
\[
B(y):=y+\{-h_b,\dots,h_b\}^3,\qquad y\in n_b\mathbb{Z}^3.
\]
The interior bonds of a block, the faces $\langle y,y'\rangle$ with $y'=y+n_be_{\mu}$, and the face bonds of a face are defined as in Section~\ref{sec:blocks-fluctuation-space}, with $n_b$ in place of $L$; every face has $n_b^2$ face bonds.

Let $w=(w_{\pi})$ be a probability vector indexed by the six orders $\pi$ of the coordinates $0,1,2$. For a bond function $A$, a point $y\in n_b\mathbb{Z}^3$, a site $x\in B(y)$ and a direction $\mu$ put
\begin{align*}
(\tau_wA)(y,x)&:=\sum_{\pi}w_{\pi}\cdot A(\Gamma^{\pi}(y,x)),\\
(Q^{[n_b]}A)(y,y+n_be_{\mu})&:=n_b^{-4}\sum_{x'\in B(y)}A(\Gamma_{x',x'+n_be_{\mu}}),
\end{align*}
where $\Gamma^{\pi}(y,x)$ and $\Gamma_{x',x'+n_be_{\mu}}$ are the paths of Section~\ref{sec:block-averages-axial}; for a face $c=\langle y,y'\rangle$ write $(Q^{[n_b]}A)(c):=(Q^{[n_b]}A)(y,y')$. For $n_b=L$ the functional $Q^{[n_b]}$ is the functional $Q$ of Section~\ref{sec:block-averages-axial}, and for $n_b=L$ and the uniform vector, $w_{\pi}=\frac{1}{6}$ for every $\pi$, the functional $\tau_w$ is the functional $\tau$ of Section~\ref{sec:block-averages-axial}. The condition $\tau_wA=0$ is called the symmetrised axial gauge if $w$ is uniform, and the unsymmetrised axial gauge if $w$ is a point mass.

For $y\in n_b\mathbb{Z}^3$ let $\mathrm{plaq}(y)$ be the set of the plaquettes with all four corners in $B(y)$, and for a face $c=\langle y,y'\rangle$ let $\mathrm{plaq}(c)$ be the set of the plaquettes with two corners in $B(y)$ and two corners in $B(y')$. On the plaquette $P=[x,x+e_{\mu},x+e_{\mu}+e_{\nu},x+e_{\nu}]$, $\mu\ne\nu$, the coboundary of $A$ takes the value
\[
dA(P)=A(x,x+e_{\mu})+A(x+e_{\mu},x+e_{\mu}+e_{\nu})-A(x+e_{\nu},x+e_{\nu}+e_{\mu})-A(x,x+e_{\nu}).
\]
The box spanned by finitely many sites is the smallest set of the form $I_0\times I_1\times I_2$, with intervals of integers $I_0$, $I_1$, $I_2$, that contains them. A path, a bond or a plaquette lies in a box if all its sites, resp. both its ends, resp. all its corners, lie in the box.

\paragraph{Strips}
Let $z$ be a site, $t$ an integer, $S$ the straight segment of bonds from $z$ to $z+te_{\mathfrak{l}}$, and $\nu\ne\mathfrak{l}$. The strip between $S$ and its translate $S+e_{\nu}$ is the set of the $|t|$ plaquettes that have one side in $S$ and the opposite side in $S+e_{\nu}$. Then
\begin{equation}\label{eq:operators-photon-14}
A(S)+A(z+te_{\mathfrak{l}},z+te_{\mathfrak{l}}+e_{\nu})-A(S+e_{\nu})-A(z,z+e_{\nu})=\pm\sum_PdA(P),
\end{equation}
where $P$ ranges over the strip between $S$ and $S+e_{\nu}$, each of its plaquettes being written as $[x,x+e_{\mathfrak{l}},x+e_{\mathfrak{l}}+e_{\nu},x+e_{\nu}]$, and the sign is the sign of $t$. Indeed, when the expressions for $dA(P)$ are added over the strip, the bonds of direction $\mathfrak{l}$ contribute $A(S)-A(S+e_{\nu})$, and the bonds of direction $\nu$, which cancel in pairs except the two that join the end points of $S$ to those of $S+e_{\nu}$, contribute $A(z+te_{\mathfrak{l}},z+te_{\mathfrak{l}}+e_{\nu})-A(z,z+e_{\nu})$; both contributions carry the sign of $t$.

\begin{lemma}[Bonds in an axial gauge]\label{lem:bonds-axial-gauge}
Let $A$ be a real bond function such that $(\tau_wA)(y,x)=0$ for all $y\in n_b\mathbb{Z}^3$ and all $x\in B(y)$.
\begin{enumerate}
\item[(a)] Let $b=(x,x+e_{\nu})$ be an interior bond of $B(y)$. For every order $\pi$ there is a set $\boldsymbol{\Sigma}^{\pi}_b\subset\mathrm{plaq}(y)$ of plaquettes, contained in the box spanned by $y$, $x$ and $x+e_{\nu}$, with
\[
\#|\boldsymbol{\Sigma}^{\pi}_b|=\sum_{\mathfrak{l}}|x_{\mathfrak{l}}-y_{\mathfrak{l}}|,
\]
the sum being over the directions $\mathfrak{l}$ that come after $\nu$ in $\pi$, and there are signs, depending on $\pi$ and on $P$, such that
\[
A(b)=\sum_{\pi}w_{\pi}\sum_{P\in\boldsymbol{\Sigma}^{\pi}_b}\pm dA(P).
\]
Consequently, with $|dA|_y:=\sup_{P\in\mathrm{plaq}(y)}|dA(P)|$,
\[
|A(b)|\le\left(\sum_{\mathfrak{l}\ne\nu}|x_{\mathfrak{l}}-y_{\mathfrak{l}}|\right)\cdot|dA|_y\le(n_b-1)\cdot|dA|_y,
\]
and, if $w$ is uniform,
\[
|A(b)|\le\frac{1}{2}\left(\sum_{\mathfrak{l}\ne\nu}|x_{\mathfrak{l}}-y_{\mathfrak{l}}|\right)\cdot|dA|_y\le h_b\cdot|dA|_y.
\]
\item[(b)] Let $b$ be a face bond of the face $c=\langle y,y'\rangle$, and let $|dA|_c$ be the supremum of $|dA(P)|$ over the plaquettes $P$ of $\mathrm{plaq}(y)$, $\mathrm{plaq}(y')$ and $\mathrm{plaq}(c)$. Then
\[
|A(b)|\le c_{\mathrm{ax}}\cdot n_b^2\cdot|dA|_c+n_b\cdot\left|(Q^{[n_b]}A)(c)\right|,
\]
where $c_{\mathrm{ax}}:=3$, and $c_{\mathrm{ax}}:=\frac{3}{2}$ if $w$ is uniform.
\end{enumerate}
\end{lemma}

\begin{proof}
(a) Fix an order $\pi=(\pi_1,\pi_2,\pi_3)$, and let $j$ be the index with $\pi_j=\nu$. The path $\Gamma^{\pi}(y,x)$ is the sum $S_1+S_2+S_3$ of three straight segments, $S_i$ being the segment along which the coordinate $\pi_i$ changes from $y_{\pi_i}$ to $x_{\pi_i}$; let $S'_1$, $S'_2$, $S'_3$ be the segments of $\Gamma^{\pi}(y,x+e_{\nu})$ in the same sense. For $j\le i\le 3$ let $b_i$ be the translate of $b$ whose initial point is the end point of $S_i$; thus $b_3=b$, since $S_3$ ends at $x$. As chains of bonds,
\[
S'_i=S_i\quad(i<j),\qquad S'_j=S_j+b_j,\qquad S'_i=S_i+e_{\nu}\quad(i>j).
\]
Indeed, the two paths coincide up to the initial point of $S_j$. If $x_{\nu}\ge y_{\nu}$, then $S'_j$ is $S_j$ prolonged by the bond $b_j$; if $x_{\nu}<y_{\nu}$, then $S'_j$ is $S_j$ deprived of its last bond, which is $-b_j$. In both cases $S'_j$ ends at the end point of $S_j$ translated by $e_{\nu}$, and from there on the path to $x+e_{\nu}$ is the translate by $e_{\nu}$ of the path to $x$.

For $i>j$ the segment $S_i$ consists of $|x_{\pi_i}-y_{\pi_i}|$ bonds of the direction $\pi_i\ne\nu$; it begins at the initial point of $b_{i-1}$ and ends at the initial point of $b_i$. Hence \eqref{eq:operators-photon-14}, applied to $S_i$ and $\nu$, gives
\[
A(S_i)+A(b_i)-A(S_i+e_{\nu})-A(b_{i-1})=\pm\sum_PdA(P),
\]
the sum being over the strip between $S_i$ and $S_i+e_{\nu}$, which has $|x_{\pi_i}-y_{\pi_i}|$ elements. Adding these identities for $i>j$ and using the relations between the segments, one obtains
\begin{align*}
A(\Gamma^{\pi}(y,x))+A(b)-A(\Gamma^{\pi}(y,x+e_{\nu}))&=A(b)-A(b_j)+\sum_{i>j}\left(A(S_i)-A(S_i+e_{\nu})\right)\\
&=\sum_{P\in\boldsymbol{\Sigma}^{\pi}_b}\pm dA(P),
\end{align*}
where $\boldsymbol{\Sigma}^{\pi}_b$ is the union of the strips between $S_i$ and $S_i+e_{\nu}$, $i>j$. If $j=3$, the sums over $i>j$ are empty, $b_j=b$ and $\boldsymbol{\Sigma}^{\pi}_b=\emptyset$. The strips of distinct indices $i$ are disjoint, since their plaquettes are parallel to distinct coordinate planes, and the directions $\pi_i$, $i>j$, are the directions that come after $\nu$ in $\pi$; this gives $\#|\boldsymbol{\Sigma}^{\pi}_b|$. The path $\Gamma^{\pi}(y,x)$ lies in the box spanned by $y$ and $x$, the path $\Gamma^{\pi}(y,x+e_{\nu})$ in the box spanned by $y$ and $x+e_{\nu}$, and the bonds $b_i$ and the plaquettes of the strips have their ends, resp. corners, on these two paths. So all paths, bonds and plaquettes that occur lie in the box spanned by $y$, $x$ and $x+e_{\nu}$. This box is contained in $B(y)$, because $y$, $x$ and $x+e_{\nu}$ lie in $B(y)$; in particular $\boldsymbol{\Sigma}^{\pi}_b\subset\mathrm{plaq}(y)$.

Multiplying by $w_{\pi}$ and summing over $\pi$, one obtains, since $\sum_{\pi}w_{\pi}=1$,
\[
\sum_{\pi}w_{\pi}\sum_{P\in\boldsymbol{\Sigma}^{\pi}_b}\pm dA(P)=(\tau_wA)(y,x)+A(b)-(\tau_wA)(y,x+e_{\nu})=A(b),
\]
because $x$ and $x+e_{\nu}$ lie in $B(y)$, so that both values of $\tau_wA$ vanish. This is the asserted identity, and it follows that
\[
|A(b)|\le\sum_{\pi}w_{\pi}\cdot\#|\boldsymbol{\Sigma}^{\pi}_b|\cdot|dA|_y.
\]
For every $\pi$ one has $\#|\boldsymbol{\Sigma}^{\pi}_b|\le\sum_{\mathfrak{l}\ne\nu}|x_{\mathfrak{l}}-y_{\mathfrak{l}}|$, and this sum is at most $2h_b=n_b-1$, since $x\in B(y)$. If $w$ is uniform, every direction $\mathfrak{l}\ne\nu$ comes after $\nu$ in three of the six orders, so that
\[
\sum_{\pi}w_{\pi}\cdot\#|\boldsymbol{\Sigma}^{\pi}_b|=\frac{1}{2}\sum_{\mathfrak{l}\ne\nu}|x_{\mathfrak{l}}-y_{\mathfrak{l}}|\le h_b.
\]

(b) Put $c_i:=n_b-1$ in general and $c_i:=h_b$ if $w$ is uniform. Since $|dA|_y$ and $|dA|_{y'}$ are at most $|dA|_c$, part (a) gives
\[
|A(b'')|\le c_i\cdot|dA|_c
\]
for every interior bond $b''$ of $B(y)$ or of $B(y')$. Let $y'=y+n_be_{\mu}$. The face bonds of $c$ are the bonds $(v,v+e_{\mu})$, where $v$ ranges over the $n_b^2$ sites of $B(y)$ with $v_{\mu}=y_{\mu}+h_b$; for these sites put $a(v):=A(v,v+e_{\mu})$.

Let $v$ and $v+e_{\nu}$, $\nu\ne\mu$, be two of these sites, and $P:=[v,v+e_{\mu},v+e_{\mu}+e_{\nu},v+e_{\nu}]$. Then $P\in\mathrm{plaq}(c)$, and the formula for $dA(P)$ gives
\[
a(v+e_{\nu})-a(v)=-dA(P)+A(v+e_{\mu},v+e_{\mu}+e_{\nu})-A(v,v+e_{\nu}),
\]
where the last two bonds are interior bonds of $B(y')$ and of $B(y)$. Hence $|a(v+e_{\nu})-a(v)|\le(1+2c_i)\cdot|dA|_c$. Two of these sites $v$, $v'$ are joined, within the set of these sites, by a chain of $|v-v'|_1$ steps of this kind, so that
\[
|a(v)-a(v')|\le|v-v'|_1\,(1+2c_i)\cdot|dA|_c.
\]

Each of the $n_b^3$ straight paths $\Gamma_{x,x+n_be_{\mu}}$, $x\in B(y)$, has $n_b$ bonds: one face bond of $c$, and $n_b-1$ bonds interior to $B(y)$ or to $B(y')$. A given face bond of $c$ belongs to $n_b$ of these paths, namely to those with $x_{\mu}=y_{\mu}-h_b,\dots,y_{\mu}+h_b$ and the same other coordinates. Therefore
\[
n_b^4\,(Q^{[n_b]}A)(c)=n_b\sum_va(v)+E,
\]
where $E$ collects the remaining $n_b^3(n_b-1)$ terms, which are values of $A$ on interior bonds; hence $|E|\le n_b^3\cdot(n_b-1)\cdot c_i\cdot|dA|_c$. With the mean $\bar{a}:=n_b^{-2}\sum_va(v)$ it follows that
\[
|\bar{a}|\le n_b\cdot\left|(Q^{[n_b]}A)(c)\right|+(n_b-1)\cdot c_i\cdot|dA|_c.
\]

For $t\in\{-h_b,\dots,h_b\}$ one has
\[
\frac{1}{n_b}\sum_{t'=-h_b}^{h_b}|t-t'|\le h_b,
\]
because the left side is a convex function of $t$, whose maximum over $\{-h_b,\dots,h_b\}$ is attained at $t=\pm h_b$, where the sum equals $0+1+\dots+2h_b=h_b\cdot n_b$. The sites $v$ form a square of side $n_b$ in the two coordinates other than $\mu$; the last inequality, applied in each of these two coordinates, gives $n_b^{-2}\sum_{v'}|v-v'|_1\le 2h_b$. Since $a(v)-\bar{a}=n_b^{-2}\sum_{v'}(a(v)-a(v'))$, one obtains
\[
|a(v)-\bar{a}|\le\frac{1}{n_b^2}\sum_{v'}|a(v)-a(v')|\le 2h_b(1+2c_i)\cdot|dA|_c.
\]
Adding the bound for $|\bar{a}|$ and using $2h_b=n_b-1$, one finds
\[
|a(v)|\le n_b\cdot\left|(Q^{[n_b]}A)(c)\right|+(n_b-1)(c_i+1+2c_i)\cdot|dA|_c.
\]
The coefficient of $|dA|_c$ is $(n_b-1)(3n_b-2)\le 3n_b^2$ in general, and $\frac{1}{2}(n_b-1)(3n_b-1)\le\frac{3}{2}n_b^2$ if $w$ is uniform. Since $|A(b)|=|a(v)|$ for the end $v$ of $b$ in $B(y)$, assertion (b) follows.
\end{proof}

\subsubsection{The axial gauge on $L$-blocks}\label{sec:axial-gauge-blocks}

\begin{corollary}[Local bounds in the axial gauge]\label{cor:local-bounds-axial}
Let $l$ be a multiple of $L$, and let $A$ be a real bond function on $T^0_k$ with $\tau A=0$. Then every bond $b$ satisfies
\[
|A(b)|\le\frac{3}{2}L^2\cdot|dA|_{\infty}+L\cdot|QA|_{\infty},
\]
where the first supremum is taken over the plaquettes contained in the union of the blocks met by $b$, which are one or two in number, and the second over the face of $b$ if $b$ is a face bond; for an interior bond $b$ the second term is omitted. In particular,
\[
|Z|_{\infty}\le\frac{3}{2}L^2\cdot|dZ|_{\infty}
\]
for every real bond function $Z$ on $T^0_k$ with $QZ=0$ and $\tau Z=0$. If $w$ is a point mass, the same assertions hold with $\tau_w$, the functional of Section~\ref{sec:blocks-odd-side} for $n_b=L$, in place of $\tau$, and with $3L^2$ in place of $\frac{3}{2}L^2$.
\end{corollary}

\begin{proof}
Section~\ref{sec:blocks-odd-side} applies with $n_b=L$, since $l$ is a multiple of $L$: $A$ is identified with its periodic lift, and for uniform $w$ the functionals $\tau_w$ and $Q^{[L]}$ are $\tau$ and $Q$. If $b$ is an interior bond of $B(y)$, the plaquettes contained in $B(y)$ form $\mathrm{plaq}(y)$, and Lemma~\ref{lem:bonds-axial-gauge}(a) gives
\[
|A(b)|\le\frac{L-1}{2}\cdot|dA|_y\le\frac{3}{2}L^2\cdot|dA|_y.
\]
Let $b$ be a face bond of the face $c=\langle y,y'\rangle$, $y'=y+Le_{\mu}$. A site of $B(y)\cup B(y')$ belongs to $B(y)$ or to $B(y')$ according to its coordinate $\mu$, and each value of this coordinate is taken at two or at four corners of a plaquette; hence a plaquette contained in $B(y)\cup B(y')$ has four corners in one of the two blocks or two corners in each, and these plaquettes form $\mathrm{plaq}(y)\cup\mathrm{plaq}(y')\cup\mathrm{plaq}(c)$. Lemma~\ref{lem:bonds-axial-gauge}(b) gives
\[
|A(b)|\le\frac{3}{2}L^2\cdot|dA|_c+L\cdot|(QA)(c)|.
\]
If $QZ=0$ and $\tau Z=0$, the supremum over $b$ of these bounds gives $|Z|_{\infty}\le\frac{3}{2}L^2\cdot|dZ|_{\infty}$. If $w$ is a point mass, Lemma~\ref{lem:bonds-axial-gauge} gives in the same way $|A(b)|\le(L-1)\cdot|dA|_y\le 3L^2\cdot|dA|_y$ for an interior bond $b$ of $B(y)$, and $|A(b)|\le 3L^2\cdot|dA|_c+L\cdot|(QA)(c)|$ for a face bond $b$ of $c$. Proposition~\ref{prop:axial-gauge} follows.
\end{proof}

\begin{corollary}[An $\ell^2$ bound in the axial gauge]\label{cor:bound-axial-gauge}
Let $l$ be a multiple of $L$, and let $Z$ be a real bond function on $T^0_k$ with $QZ=0$ and $\tau Z=0$. Then
\[
|Z|_2^2\le 9L^7\cdot|dZ|_2^2.
\]
\end{corollary}

\begin{proof}
To a bond $b$ attach the set of plaquettes $\mathrm{Nb}(b):=\mathrm{plaq}(y)$ if $b$ is an interior bond of $B(y)$, and $\mathrm{Nb}(b):=\mathrm{plaq}(y)\cup\mathrm{plaq}(y')\cup\mathrm{plaq}(c)$ if $b$ is a face bond of the face $c=\langle y,y'\rangle$. By the proof of Corollary~\ref{cor:local-bounds-axial}, in either case $\mathrm{Nb}(b)$ consists of the plaquettes contained in the union of the blocks met by $b$. As $QZ=0$, that corollary therefore gives
\[
|Z(b)|^2\le\frac{9}{4}L^4\sup_{P\in\mathrm{Nb}(b)}|dZ(P)|^2\le\frac{9}{4}L^4\sum_{P\in\mathrm{Nb}(b)}|dZ(P)|^2.
\]
A plaquette of $\mathrm{plaq}(y)$ lies in $\mathrm{Nb}(b)$ only if $b$ is one of the $3L^2(L-1)$ interior bonds of $B(y)$ or one of the $6L^2$ face bonds of the six faces of $B(y)$, that is, for at most $3L^3+3L^2\le 4L^3$ bonds $b$, as $L\ge 3$; a plaquette of $\mathrm{plaq}(c)$ lies in $\mathrm{Nb}(b)$ only if $b$ is one of the $L^2$ face bonds of $c$. Summing over the bonds, one obtains
\[
|Z|_2^2=\sum_b|Z(b)|^2\le\frac{9}{4}L^4\cdot 4L^3\sum_P|dZ(P)|^2=9L^7\cdot|dZ|_2^2,
\]
where $b$ ranges over the positively oriented bonds and $P$ over the plaquettes of $T^0_k$.
\end{proof}

\subsubsection{Gauge fixing for bond functions of mean zero}\label{sec:gauge-fixing-bond}

\begin{lemma}[Line sums from a base point]\label{lem:line-sums-base}
Let $l\ge 2$ be an integer, let $A$ be a real bond function on $(\mathbb{Z}/l\mathbb{Z})^3$ with mean zero, and let $x_0$ be a site. Put
\[
W:=x_0+\left\{-\left\lfloor\frac{l-1}{2}\right\rfloor,\dots,\left\lfloor\frac{l}{2}\right\rfloor\right\}^3,
\]
so that every site of the torus occurs in $W$ exactly once. For $x\in W$ let
\[
\lambda(x):=A(\Gamma^{(0,1,2)}(x_0,x)),
\]
the path $\Gamma^{(0,1,2)}(x_0,x)$ of Section~\ref{sec:block-averages-axial} being taken inside $W$, and put
\[
A'':=A-d\lambda,\qquad p_A:=|dA|_{\infty}.
\]
A bond is specified by its initial point and its direction. For a bond $(x,x+e_{\nu})$ write $x=x_0+t$ with $t\in\left\{-\left\lfloor\frac{l-1}{2}\right\rfloor,\dots,\left\lfloor\frac{l}{2}\right\rfloor\right\}^3$; the bond is called inner if $t_{\nu}<\lfloor l/2\rfloor$, and every bond that is not inner is called a seam bond.
\begin{enumerate}
\item[(i)] For an inner bond $b'=(x,x+e_{\nu})$,
\[
|A''(b')|\le\left(\sum_{\mathfrak{l}>\nu}d(x_{\mathfrak{l}},x_{0,\mathfrak{l}})\right)\cdot p_A\le 2\,d(x,x_0)\cdot p_A\le l\cdot p_A,
\]
where $d(x_{\mathfrak{l}},x_{0,\mathfrak{l}})$ is the distance in $\mathbb{Z}/l\mathbb{Z}$.
\item[(ii)] For a seam bond $b'$,
\[
|A''(b')|\le\frac{3}{2}l^2\cdot p_A.
\]
\end{enumerate}
\end{lemma}

\begin{proof}
(i) For $x\in W$ one has
\[
A''(\Gamma^{(0,1,2)}(x_0,x))=\lambda(x)-(\lambda(x)-\lambda(x_0))=0,
\]
because the line sum of $d\lambda$ along a path is the difference of the values of $\lambda$ at its end points, and $\lambda(x_0)=0$; moreover $dA''=dA$. Let $b'=(x,x+e_{\nu})$ be an inner bond, $x=x_0+t$. Since $t_{\nu}+1\le\lfloor l/2\rfloor$, the site $x+e_{\nu}=x_0+t+e_{\nu}$ lies in $W$ as well, and so does the box spanned by $x_0$, $x$ and $x+e_{\nu}$. The proof of Lemma~\ref{lem:bonds-axial-gauge}(a), carried out for $A''$ with $x_0$ in place of $y$ and with $w$ the point mass at the order $(0,1,2)$, uses only the vanishing of $A''$ on the paths $\Gamma^{(0,1,2)}(x_0,x)$ and $\Gamma^{(0,1,2)}(x_0,x+e_{\nu})$, the identity \eqref{eq:operators-photon-14}, and paths, bonds and plaquettes in the box spanned by $x_0$, $x$ and $x+e_{\nu}$, which lies in $W$. Since the directions that come after $\nu$ in the order $(0,1,2)$ are the directions $\mathfrak{l}>\nu$, it yields
\[
A''(b')=\sum_P\pm dA''(P)=\sum_P\pm dA(P),
\]
the sums being over a set of $\sum_{\mathfrak{l}>\nu}|t_{\mathfrak{l}}|$ plaquettes. As $|t_{\mathfrak{l}}|\le\lfloor l/2\rfloor$, the number $|t_{\mathfrak{l}}|$ is the distance $d(x_{\mathfrak{l}},x_{0,\mathfrak{l}})$ in $\mathbb{Z}/l\mathbb{Z}$; this proves the first inequality. The sum over $\mathfrak{l}>\nu$ has at most two terms, each at most $d(x,x_0)$, and $d(x,x_0)\le\lfloor l/2\rfloor$; this proves the other two.

(ii) Let $b'$ be a seam bond of direction $\nu$. For a site $v$ put
\[
\mathrm{hol}(v):=\sum_{j\in\mathbb{Z}/l\mathbb{Z}}A_{\nu}(v+je_{\nu}),
\]
the line sum of $A$ along the closed straight line through $v$ in the direction $\nu$. It does not change when $A$ is replaced by $A''$, since the line sum of $d\lambda$ along a closed path vanishes. On every closed line in the direction $\nu$ exactly one bond has its initial point $x_0+t$ with $t_{\nu}=\lfloor l/2\rfloor$. So the closed line that contains $b'$ contains, besides $b'$, exactly $l-1$ bonds, all of them inner. Fix a site $u$ of this line; then $A''(b')$ is $\mathrm{hol}(u)$ minus the sum of $A''$ over these $l-1$ bonds, and (i) gives
\[
|A''(b')|\le|\mathrm{hol}(u)|+(l-1)\,l\cdot p_A.
\]
Now let $v$ be an arbitrary site and $\mathfrak{l}\ne\nu$ an arbitrary direction. The sum of $dA$ over the $l$ plaquettes $[v+je_{\nu},v+(j+1)e_{\nu},v+(j+1)e_{\nu}+e_{\mathfrak{l}},v+je_{\nu}+e_{\mathfrak{l}}]$, $j\in\mathbb{Z}/l\mathbb{Z}$, which lie between the lines through $v$ and through $v+e_{\mathfrak{l}}$, is $\mathrm{hol}(v)-\mathrm{hol}(v+e_{\mathfrak{l}})$, because every bond of direction $\mathfrak{l}$ occurs in two consecutive plaquettes of this family, with opposite signs. It follows that, for every site $v$ and every direction $\mathfrak{l}\ne\nu$,
\[
|\mathrm{hol}(v)-\mathrm{hol}(v+e_{\mathfrak{l}})|\le l\cdot p_A.
\]
The function $\mathrm{hol}$ depends only on the two coordinates of $v$ other than $\nu$, that is, on a point of $(\mathbb{Z}/l\mathbb{Z})^2$, and its mean over $(\mathbb{Z}/l\mathbb{Z})^2$ is $l^{-2}\sum_xA_{\nu}(x)=0$, since $A$ has mean zero. Two points of $(\mathbb{Z}/l\mathbb{Z})^2$ are joined by a chain of unit steps whose number is the sum of the distances in $\mathbb{Z}/l\mathbb{Z}$ of their two pairs of coordinates. Each unit step joins two points of the form $p$ and $p+e_{\mathfrak{l}}$, $\mathfrak{l}\ne\nu$, so that, by the last inequality, it changes $\mathrm{hol}$ by at most $l\cdot p_A$; consequently the values of $\mathrm{hol}$ at the two ends of the chain differ by at most $l\cdot p_A$ times the number of its steps. The mean of the distances in $\mathbb{Z}/l\mathbb{Z}$ from a fixed point to the points of $\mathbb{Z}/l\mathbb{Z}$ is at most $\frac{l}{4}$ (it equals $\frac{l}{4}$ for even $l$ and $\frac{l^2-1}{4l}$ for odd $l$). Hence, for every site $v$, with $v'$ ranging over $(\mathbb{Z}/l\mathbb{Z})^2$,
\[
|\mathrm{hol}(v)|\le\frac{1}{l^2}\sum_{v'}|\mathrm{hol}(v)-\mathrm{hol}(v')|\le l\cdot p_A\cdot 2\cdot\frac{l}{4}=\frac{l^2}{2}\cdot p_A.
\]
For $v=u$, together with the bound on $|A''(b')|$ above, this gives
\[
|A''(b')|\le\left(\frac{l^2}{2}+l(l-1)\right)\cdot p_A\le\frac{3}{2}l^2\cdot p_A.
\]
\end{proof}

\begin{corollary}[A bounded gauge on the torus]\label{cor:bounded-gauge-torus}
Let $l\ge 2$, and let $A$ be a real one-cochain on $(\mathbb{Z}/l\mathbb{Z})^3$ with mean zero. Then there is a function $\lambda$ on the sites of $(\mathbb{Z}/l\mathbb{Z})^3$ such that
\[
|A-d\lambda|_{\infty}\le 2l^2\cdot|dA|_{\infty}.
\]
\end{corollary}

\begin{proof}
Let $x_0$ be an arbitrary site, and let $\lambda$ be the function of Lemma~\ref{lem:line-sums-base}. Every bond is inner or a seam bond. By Lemma~\ref{lem:line-sums-base}(i) and (ii), $|(A-d\lambda)(b')|$ is at most $l\cdot|dA|_{\infty}$ in the first case and at most $\frac{3}{2}l^2\cdot|dA|_{\infty}$ in the second, and $l\le\frac{3}{2}l^2\le 2l^2$. Proposition~\ref{prop:gauge-fixing-controlled} follows.
\end{proof}

\subsection{The fluctuation space, its basis and the map $C$}\label{sec:fluctuation-space-basis}

Throughout this section $l$ is a multiple of $L$. In those assertions of Proposition~\ref{prop:fluctuation-space}(ii), (iii) that involve the group $\mathbb{G}_k$, the side $l$ belongs to $\mathrm{Cl}$, as in the setting of Section~\ref{sec:tori-cochains-kernels}. For $l=L$ the blocks, the faces and the bonds are taken on the periodic lift of the torus to $\mathbb{Z}^3$, as in Section~\ref{sec:blocks-odd-side}.

Let $y\in L\mathbb{Z}^3/l\mathbb{Z}^3$. The $L$-block $B(y)$ of Section~\ref{sec:blocks-fluctuation-space} contains $3L^2(L-1)$ positively oriented interior bonds (in each of the three directions, $L^2$ straight lines of $L-1$ bonds). Let $\mathrm{Bd}(y)$ be the space of the real functions on these bonds, with the counting scalar product
\[
\langle A,A'\rangle:=\sum_bA(b)\cdot A'(b),
\]
the sum being taken over the positively oriented interior bonds $b$ of $B(y)$. An element of $\mathrm{Bd}(y)$ is identified with the bond function on $T^0_k$ that coincides with it on these bonds and vanishes on all other positively oriented bonds. For $x\in B(y)$ the paths $\Gamma^{\pi}(y,x)$ consist of interior bonds of $B(y)$; hence the axial functional $\tau$ of Section~\ref{sec:block-averages-axial} is defined on $\mathrm{Bd}(y)$, and one puts
\[
\ker\tau_y:=\left\{\mathsf{Z}_1\in\mathrm{Bd}(y):(\tau\mathsf{Z}_1)(y,x)=0\text{ for all }x\in B(y)\right\}.
\]
For a positively oriented bond $b$ of $T^0_k$ let $\delta_b$ be the bond function that equals $1$ on $b$ and vanishes on every other positively oriented bond. The fluctuation space $V_k$ of Section~\ref{sec:averaging-fluctuations-gauge} is the orthogonal direct sum of the spaces $\ker\tau_y$, $y\in L\mathbb{Z}^3/l\mathbb{Z}^3$, and of the lines $\mathbb{R}\delta_b$, where $b$ ranges over the positively oriented face bonds that are not central. In particular every element of $V_k$ is a bond function on $T^0_k$ whose value at each central bond is zero. For $\mathsf{Z},\Upsilon\in V_k$ write $\mathsf{Z}(\Upsilon):=\langle\mathsf{Z},\Upsilon\rangle$.

\begin{proposition}[The fluctuation space]\label{prop:fluctuation-space}
\begin{enumerate}
\item[(i)] For $\mu=0,1,2$ let $\mathfrak{r}_{\mu}$ be the reflection of $\mathbb{Z}^3$ that changes the sign of the coordinate $x_{\mu}$, and let it act on $\mathrm{Bd}(0)$ by $(\mathfrak{r}_{\mu}A)(x,x'):=A(\mathfrak{r}_{\mu}x,\mathfrak{r}_{\mu}x')$. Then $\mathfrak{r}_0$, $\mathfrak{r}_1$, $\mathfrak{r}_2$ are commuting orthogonal involutions of $\mathrm{Bd}(0)$, and each of them maps $\ker\tau_0$ onto itself. Consequently $\ker\tau_0$ has an orthonormal basis $\{e_i^0\}$ of joint eigenvectors of the three reflections:
\[
\mathfrak{r}_{\mu}e_i^0=\pm e_i^0\qquad(\mu=0,1,2).
\]
The space $\mathrm{Bd}(0)$, the functional $\tau$ on $B(0)$ and the reflections $\mathfrak{r}_{\mu}$ depend on $L$ alone; for each $L$ one such basis is fixed, and the same basis is taken at every level $k$ and on every torus. Moreover
\[
\dim\ker\tau_0=2L^3-3L^2+1.
\]
\item[(ii)] For $y\in L\mathbb{Z}^3/l\mathbb{Z}^3$ let $e_i^y$ be the translate of $e_i^0$ by $y$, that is, $e_i^y(b):=e_i^0(b-y)$, where $b-y$ denotes the bond $b$ translated by $-y$. Then
\begin{equation}\label{eq:operators-photon-15}
\{\Upsilon\}:=\{e_i^y\}\cup\left\{\delta_b:b\text{ a positively oriented face bond that is not central}\right\}
\end{equation}
is an orthonormal basis of $V_k$ with $d_V\cdot n_k$ elements, where $d_V=2-2L^{-3}$ and $n_k=l^3$. For every $\Upsilon$ the set $\operatorname{supp}\Upsilon$ is contained in one $L$-block or is the set of the two ends of one bond. The action of each $g\in\mathbb{G}_k$ on bond functions (Section~\ref{sec:dirac-matrices-lattice}) restricts to an isometric bijection of $V_k$, and it maps the set $\{\pm\Upsilon\}$ of the basis vectors and their negatives bijectively onto itself.
\item[(iii)] The set $\mathrm{Box}_k$ and the measures $\mu_I$ and $\hat{\mu}$ of Section~\ref{sec:averaging-fluctuations-gauge} are invariant under $\mathbb{G}_k$ and under the map $\mathsf{Z}\mapsto-\mathsf{Z}$.
\item[(iv)] Let $\mathsf{Z}\in V_k$. Among the bond functions on $T^0_k$ that agree with $\mathsf{Z}$ outside the set of the central bonds, exactly one is annihilated by $Q$; it is denoted by $C\mathsf{Z}$. Its value at the central bond $b_c$ of the face $c=\langle y,y+L e_{\mu}\rangle$ is
\begin{equation}\label{eq:operators-photon-16}
(C\mathsf{Z})(b_c)=-\sum_{\substack{b\text{ face bond of }c\\ b\ne b_c}}\mathsf{Z}(b)-\sum_b\frac{n_c(b)}{L}\cdot\mathsf{Z}(b),
\end{equation}
where the first sum is over the positively oriented face bonds of $c$ other than $b_c$, the second over the positively oriented interior bonds $b$ of direction $\mu$ of $B(y)$ and of $B(y+L e_{\mu})$, and where $n_c(b)\in\{1,\dots,L-1\}$ is the number of the paths $\Gamma_{x,x+L e_{\mu}}$, $x\in B(y)$, that contain $b$. The map $C$ is a linear bijection of $V_k$ onto the space of the real bond functions $Z$ on $T^0_k$ with $Q Z=0$ and $\tau Z=0$. For every positively oriented bond $b$ the number $(C\mathsf{Z})(b)$ depends only on the coordinates $\mathsf{Z}(e_i^z)$ with $B(z)$ one of the (at most two) $L$-blocks that contain an end of $b$, and, if $b$ is a face bond, on the coordinates $\mathsf{Z}(\delta_{b'})$ with $b'$ a face bond of the same face as $b$. Moreover, with
\[
c_C:=3L^{\frac{9}{2}},\qquad c_1:=14L^{\frac{3}{2}},
\]
one has
\begin{align*}
|C\mathsf{Z}|_{\infty}&\le c_C\cdot\sup_{\Upsilon}|\mathsf{Z}(\Upsilon)|,\\
|\mathsf{Z}|_2\le|C\mathsf{Z}|_2&\le(5L^3)^{\frac{1}{2}}\cdot|\mathsf{Z}|_2,\\
\sum_b|(C\Upsilon)(b)|&\le c_1\qquad\text{for every }\Upsilon,
\end{align*}
the last sum being over the positively oriented bonds $b$ of $T^0_k$. Finally, for every $\Upsilon$ let $\operatorname{supp}C\Upsilon$ be the set of the end points of the bonds on which $C\Upsilon$ does not vanish, and let $d$ be the sup distance on the torus (Section~\ref{sec:tori-cochains-kernels}); then
\[
d(x,x')\le L\qquad\text{for all }x\in\operatorname{supp}C\Upsilon\text{ and }x'\in\operatorname{supp}\Upsilon.
\]
\item[(v)] Let $g$ be a standard normal random variable. Then
\[
d\varepsilon^0_k=-d_V\cdot\log P\left(|g|\le p_{0,k}\right),
\]
and if $p_{0,k}\ge1$, then
\[
0\le d\varepsilon^0_k\le12\exp\left(-\frac{p_{0,k}^2}{2}\right).
\]
The hypothesis $p_{0,k}\ge1$ is satisfied under condition~\eqref{eq:renormalisation-group-25}.
\end{enumerate}
\end{proposition}

\begin{proof}
(i) As $L$ is odd, the block $B(0)=\left\{-\frac{L-1}{2},\dots,\frac{L-1}{2}\right\}^3$ is centred at the site $0$, and each $\mathfrak{r}_{\mu}$ maps it onto itself. Consequently $\mathfrak{r}_{\mu}$ permutes the oriented interior bonds of $B(0)$: for $\nu\ne\mu$ the bond $(x,x+e_{\nu})$ goes to the positively oriented bond $(\mathfrak{r}_{\mu}x,\mathfrak{r}_{\mu}x+e_{\nu})$, whereas $(x,x+e_{\mu})$ goes to $(\mathfrak{r}_{\mu}x,\mathfrak{r}_{\mu}x-e_{\mu})$, the reversal of the positively oriented bond $(\mathfrak{r}_{\mu}x-e_{\mu},\mathfrak{r}_{\mu}x)$. Thus $A\mapsto\mathfrak{r}_{\mu}A$ permutes the coordinates $A(b)$, $b$ positively oriented, up to signs, and it is orthogonal. Since $\mathfrak{r}_{\mu}\circ\mathfrak{r}_{\mu}$ is the identity and $\mathfrak{r}_{\mu}\circ\mathfrak{r}_{\nu}=\mathfrak{r}_{\nu}\circ\mathfrak{r}_{\mu}$, the three maps are involutions and commute with each other. An orthogonal involution $O$ is symmetric, because $O^T=O^{-1}=O$.

The reflection $\mathfrak{r}_{\mu}$ is diagonal; therefore it maps the rectilinear path $\Gamma^{\pi}(0,x)$ onto the rectilinear path $\Gamma^{\pi}(0,\mathfrak{r}_{\mu}x)$ with the same order $\pi$ (Proposition~\ref{prop:properties-actions}(a)). For $x\in B(0)$ and every order $\pi$ this gives $(\mathfrak{r}_{\mu}A)(\Gamma^{\pi}(0,x))=A(\Gamma^{\pi}(0,\mathfrak{r}_{\mu}x))$, and the mean over the six orders gives
\[
(\tau\,\mathfrak{r}_{\mu}A)(0,x)=(\tau A)(0,\mathfrak{r}_{\mu}x).
\]
As $\mathfrak{r}_{\mu}x\in B(0)$, the map $\mathfrak{r}_{\mu}$ sends $\ker\tau_0$ into itself, and onto itself because it is an involution. The restrictions of $\mathfrak{r}_0$, $\mathfrak{r}_1$, $\mathfrak{r}_2$ to $\ker\tau_0$ are thus commuting symmetric involutions of a Euclidean space. By the simultaneous diagonalisation of commuting symmetric matrices, $\ker\tau_0$ has an orthonormal basis that consists of joint eigenvectors, and the eigenvalues of an involution are $\pm1$.

For every $A\in\mathrm{Bd}(0)$ one has $(\tau A)(0,0)=0$, the paths $\Gamma^{\pi}(0,0)$ being empty. The linear map that sends $A\in\mathrm{Bd}(0)$ to the function $x\mapsto(\tau A)(0,x)$ on $B(0)\setminus\{0\}$ has the kernel $\ker\tau_0$, and it is onto. Indeed, let $f$ be a function on $B(0)\setminus\{0\}$, and let $\lambda$ be the function on $B(0)$ with $\lambda(0):=0$ and $\lambda(x):=f(x)$ for $x\ne0$. The restriction of $d\lambda$ to the interior bonds of $B(0)$ lies in $\mathrm{Bd}(0)$, and its sum along every path in $B(0)$ from $0$ to $x$ telescopes to $\lambda(x)-\lambda(0)$; so
\[
(\tau\,d\lambda)(0,x)=\lambda(x)-\lambda(0)=f(x)\qquad(x\in B(0)\setminus\{0\}).
\]
Since $B(0)$ has $L^3$ sites, $\dim\ker\tau_0=3L^2(L-1)-(L^3-1)=2L^3-3L^2+1$.

(ii) Let $y\in L\mathbb{Z}^3/l\mathbb{Z}^3$. The translation by $y$ maps the interior bonds of $B(0)$ onto those of $B(y)$, and $\Gamma^{\pi}(y,x)=y+\Gamma^{\pi}(0,x-y)$ for $x\in B(y)$. Hence $A\mapsto A(\cdot-y)$ is an isometry of $\mathrm{Bd}(0)$ onto $\mathrm{Bd}(y)$ that maps $\ker\tau_0$ onto $\ker\tau_y$, and $\{e_i^y\}_i$ is an orthonormal basis of $\ker\tau_y$. The spaces $\ker\tau_y$ with different $y$ and the vectors $\delta_b$ with different face bonds $b$ are carried by pairwise disjoint sets of bonds, and are therefore mutually orthogonal; by the definition of $V_k$ the set \eqref{eq:operators-photon-15} is an orthonormal basis of $V_k$. Every $e_i^y$ vanishes off the interior bonds of $B(y)$, so that $\operatorname{supp}e_i^y\subset B(y)$, and $\operatorname{supp}\delta_b$ is the set of the two ends of $b$.

Every face is $\langle y,y+L e_{\mu}\rangle$ for exactly one $y\in L\mathbb{Z}^3/l\mathbb{Z}^3$ and exactly one $\mu$. To a given $y$ belong the $2L^3-3L^2+1$ vectors $e_i^y$, by (i), and the $3(L^2-1)$ vectors $\delta_b$ with $b$ a face bond, other than the central one, of one of the three faces $\langle y,y+L e_{\mu}\rangle$, $\mu=0,1,2$; together these are $2L^3-2$ basis elements. There are $l^3\cdot L^{-3}=n_k\cdot L^{-3}$ blocks, so that $\{\Upsilon\}$ has $(2L^3-2)\cdot n_k\cdot L^{-3}=d_V\cdot n_k$ elements.

Let $g=(S,a)\in\mathbb{G}_k$ and $r:=r(S)$, so that $g x=r x+a$, and $g$ acts on bond functions by $(g A)(x,x')=A(g^{-1}x,g^{-1}x')$ (Section~\ref{sec:dirac-matrices-lattice}). For a map $\varphi$ of the sites and a bond $b=(x,x')$ write $\varphi b:=(\varphi x,\varphi x')$, and put $(r A)(b):=A(r^{-1}b)$. The representative of $a$ can be taken in $M\mathbb{Z}^3$, which is contained in $L\mathbb{Z}^3$ because $M=L^m$ (Section~\ref{sec:lattice-geometry}). The matrix $r$ is diagonal with entries $\pm1$, hence a product of some of the reflections $\mathfrak{r}_{\mu}$, and (i) gives $r e_i^0=\pm e_i^0$. Since $g^{-1}b=r^{-1}(b-a)$ and $r^{-1}(b-a)-y=r^{-1}(b-a-r y)$,
\[
(g e_i^y)(b)=e_i^0\left(r^{-1}(b-a)-y\right)=(r e_i^0)(b-a-r y)=\pm e_i^{a+r y}(b),
\]
that is, $g e_i^y=\pm e_i^{a+r y}$, where $a+r y\in L\mathbb{Z}^3/l\mathbb{Z}^3$. For a face bond $b$ the value $(g\delta_b)(b'')=\delta_b(g^{-1}b'')$ vanishes unless $b''$ is $g b$ or its reversal; hence $g\delta_b=\pm\delta_{b'}$, where $b'$ is the positively oriented bond among $g b$ and its reversal. By Proposition~\ref{prop:properties-actions}(a), $g$ maps each $L$-block onto an $L$-block and its centre onto the centre of the image block. So $g$ maps interior bonds to interior bonds, faces to faces and face bonds to face bonds, and, since the central bond of a face is the face bond on the straight line through the two centres, central bonds to central bonds; in particular $b'$ is a face bond that is not central. Thus $g$ maps the set $\{\pm\Upsilon\}$ into itself, and since $g^{-1}$ does the same, $g$ permutes $\{\pm\Upsilon\}$. As $g$ permutes the oriented bonds, $A\mapsto g A$ preserves the counting scalar product (Proposition~\ref{prop:properties-actions}(b)). It maps the basis $\{\Upsilon\}$ into $V_k$, hence $V_k$ isometrically into $V_k$; the same holds for $g^{-1}$, so that the image is all of $V_k$.

(iii) Let $O$ be an orthogonal map of $V_k$ that permutes the set $\{\pm\Upsilon\}$; by (ii) every $g\in\mathbb{G}_k$ is such a map, and so is $\mathsf{Z}\mapsto-\mathsf{Z}$. Then $O^{-1}$ permutes $\{\pm\Upsilon\}$ as well: for every basis element $\Upsilon$ there are a basis element $\Upsilon'$ and a sign $s$ with $O^{-1}\Upsilon=s\cdot\Upsilon'$, and $\Upsilon\mapsto\Upsilon'$ is a bijection of the basis onto itself. For $\zeta\in V_k$ it follows that
\[
|(O\zeta)(\Upsilon)|=|\langle\zeta,O^{-1}\Upsilon\rangle|=|\zeta(\Upsilon')|,
\]
so that $\zeta$ lies in the cube $\mathrm{Box}_k=\left\{\zeta\in V_k:|\zeta(\Upsilon)|\le p_{0,k}\text{ for all }\Upsilon\right\}$ of Section~\ref{sec:averaging-fluctuations-gauge} if and only if $O\zeta$ does. In other words $O^{-1}\mathrm{Box}_k=\mathrm{Box}_k$, and hence $O\mathrm{Box}_k=\mathrm{Box}_k$. As $O$ is orthogonal, the standard Gaussian measure $\mu_I$ on $V_k$ is invariant under $O$, that is, $\mu_I(O^{-1}E)=\mu_I(E)$ for every Borel set $E\subset V_k$. Since $O^{-1}E\cap\mathrm{Box}_k=O^{-1}(E\cap\mathrm{Box}_k)$, this invariance, applied to the set $E\cap\mathrm{Box}_k$, gives
\[
\mu_I\left(O^{-1}E\cap\mathrm{Box}_k\right)=\mu_I\left(E\cap\mathrm{Box}_k\right)\qquad\text{for every Borel set }E\subset V_k.
\]
Dividing by $\mu_I(\mathrm{Box}_k)$, one obtains the invariance under $O$ of $\hat{\mu}$, the restriction of $\mu_I$ to $\mathrm{Box}_k$ divided by $\mu_I(\mathrm{Box}_k)$.

(iv) Let $c=\langle y,y+L e_{\mu}\rangle$ be a face, and for a bond function $Z$ on $T^0_k$ write $(Q Z)(c):=(Q Z)(y,y+L e_{\mu})$, with $Q$ as in Section~\ref{sec:block-averages-axial}. Each path $\Gamma_{x,x+L e_{\mu}}$, $x\in B(y)$, lies on a straight line of direction $\mu$ that meets $B(y)$ and $B(y+L e_{\mu})$; there are $L^2$ such lines. On one of them let $z$ be the site of $B(y)$ with $z_{\mu}=y_{\mu}-\frac{L-1}{2}$, and give the bond $(z+t e_{\mu},z+(t+1)e_{\mu})$ the number $t\in\{0,\dots,2L-2\}$. The sites $z+t e_{\mu}$ with $0\le t\le L-1$ lie in $B(y)$, and those with $L\le t\le 2L-1$ in $B(y+L e_{\mu})$; so the bond $t=L-1$ is the face bond of $c$ on this line, and the others are interior bonds of direction $\mu$. For $0\le j\le L-1$ the path $\Gamma_{z+j e_{\mu},z+(j+L)e_{\mu}}$ consists of the bonds $j,\dots,j+L-1$. Hence the bond $t$ lies on exactly $\min(t+1,2L-1-t)$ of the paths $\Gamma_{x,x+L e_{\mu}}$, $x\in B(y)$: on $L$ of them for $t=L-1$, and on a number of them in $\{1,\dots,L-1\}$ otherwise. Summing over the $L^2$ lines one finds, for every bond function $Z$ on $T^0_k$,
\[
L^4\cdot(Q Z)(c)=L\cdot\sum_{b\text{ face bond of }c}Z(b)+\sum_bn_c(b)\cdot Z(b),\qquad n_c(b)=\min(t+1,2L-1-t)\le L-1\quad(b\text{ interior}),
\]
where the second sum is over the same bonds as the second sum in \eqref{eq:operators-photon-16}, and $t$ is the number of $b$ on its line. If $Z$ coincides with $\mathsf{Z}$ off the central bonds, the equation $(Q Z)(c)=0$ is therefore the equation \eqref{eq:operators-photon-16} for the value $Z(b_c)$.

Central bonds of different faces are different. A path $\Gamma_{x,x+L e_{\nu}}$ with $x\in B(y')$ contains exactly one face bond, and this bond belongs to the face $\langle y',y'+L e_{\nu}\rangle$; the paths $\Gamma^{\pi}(y',x)$, $x\in B(y')$, consist of interior bonds of $B(y')$. Hence the central bond $b_c$ occurs in the equation $(Q Z)(c')=0$ only for $c'=c$, and in no path of $\tau$. The system $Q Z=0$ consists of one equation for each face $c$, and among the values of $Z$ at central bonds this equation contains only $Z(b_c)$, with the coefficient $L^{-3}\ne0$. Let $\mathsf{Z}\in V_k$. Once the values of $Z$ outside the central bonds are fixed to be those of $\mathsf{Z}$, the system $Q Z=0$ therefore splits into separate equations, the equation of the face $c$ being \eqref{eq:operators-photon-16} for the single unknown $Z(b_c)$, and it has exactly one solution. Thus $C\mathsf{Z}$ exists and is unique; by \eqref{eq:operators-photon-16} it is real, as $\mathsf{Z}$ is real, and it depends linearly on $\mathsf{Z}$. Since $\tau$ involves only interior bonds, $\tau C\mathsf{Z}=\tau\mathsf{Z}=0$. Conversely, let $Z$ be a real bond function on $T^0_k$ with $Q Z=0$ and $\tau Z=0$, and let $\mathsf{Z}$ coincide with $Z$ off the central bonds and vanish on them. Then $\mathsf{Z}$ is a real bond function that vanishes on the central bonds and satisfies $\tau\mathsf{Z}=\tau Z=0$, so that $\mathsf{Z}\in V_k$ by the definition of Section~\ref{sec:averaging-fluctuations-gauge}, and $Z=C\mathsf{Z}$ by the uniqueness. Finally $C$ is injective, because $\mathsf{Z}$ vanishes on the central bonds and coincides with $C\mathsf{Z}$ on all other bonds. So $C$ is a linear bijection of $V_k$ onto the space of the real bond functions $Z$ on $T^0_k$ with $Q Z=0$ and $\tau Z=0$.

Every $\mathsf{Z}\in V_k$ equals $\sum_{\Upsilon}\mathsf{Z}(\Upsilon)\cdot\Upsilon$. Let $b$ be a positively oriented bond. If $b$ is an interior bond of $B(z)$, then $(C\mathsf{Z})(b)=\mathsf{Z}(b)=\sum_i\mathsf{Z}(e_i^z)\cdot e_i^z(b)$. If $b$ is a face bond that is not central, then $(C\mathsf{Z})(b)=\mathsf{Z}(b)=\mathsf{Z}(\delta_b)$. If $b=b_c$, then by \eqref{eq:operators-photon-16} the number $(C\mathsf{Z})(b)$ is a linear combination of the numbers $\mathsf{Z}(\delta_{b'})$, $b'$ a face bond of $c$ other than $b_c$, and of the values of $\mathsf{Z}$ at interior bonds of $B(y)$ and of $B(y+L e_{\mu})$; these values are linear combinations of the numbers $\mathsf{Z}(e_i^y)$ and $\mathsf{Z}(e_i^{y+L e_{\mu}})$. In each case only the coordinates named in (iv) occur.

\emph{The bound in the supremum norm.} Both sides of the first inequality are homogeneous of degree one in $\mathsf{Z}$; it suffices therefore to consider $\mathsf{Z}$ with $|\mathsf{Z}(\Upsilon)|\le1$ for all $\Upsilon$. For an interior bond $b$ of $B(y)$ the Cauchy--Schwarz inequality and (i) give
\[
|\mathsf{Z}(b)|=\left|\sum_i\mathsf{Z}(e_i^y)\cdot e_i^y(b)\right|\le(\dim\ker\tau_0)^{\frac{1}{2}}\cdot\left(\sum_ie_i^y(b)^2\right)^{\frac{1}{2}}\le(2L^3)^{\frac{1}{2}},
\]
because $\sum_ie_i^y(b)^2=\sum_i\langle\delta_b,e_i^y\rangle^2$ is the squared norm of the orthogonal projection of $\delta_b$ onto $\ker\tau_y$, hence at most $1$. For a face bond $b$ that is not central, $|\mathsf{Z}(b)|=|\mathsf{Z}(\delta_b)|\le1$. The right side of \eqref{eq:operators-photon-16} has $L^2-1$ terms of the first kind, each of modulus at most $1$, and $2L^2(L-1)$ terms of the second kind, whose coefficients satisfy $\frac{n_c(b)}{L}<1$; hence
\[
|(C\mathsf{Z})(b_c)|\le L^2+2L^3\cdot(2L^3)^{\frac{1}{2}}=L^2+2\sqrt{2}\cdot L^{\frac{9}{2}}\le3L^{\frac{9}{2}},
\]
the last inequality because $L\ge3$ implies $L^{\frac{5}{2}}\ge3+2\sqrt{2}$. Since $(2L^3)^{\frac{1}{2}}$ and $1$ are also at most $3L^{\frac{9}{2}}$, it follows that $|C\mathsf{Z}|_{\infty}\le c_C$.

\emph{The bounds in the $\ell^2$ norm.} As $C\mathsf{Z}=\mathsf{Z}$ off the central bonds and $\mathsf{Z}$ vanishes on them,
\[
|C\mathsf{Z}|_2^2=|\mathsf{Z}|_2^2+\sum_c|(C\mathsf{Z})(b_c)|^2\ge|\mathsf{Z}|_2^2,
\]
the sum being over all faces $c$. The right side of \eqref{eq:operators-photon-16} is a linear combination of $(L^2-1)+2L^2(L-1)\le2L^3$ values of $\mathsf{Z}$, with coefficients of modulus at most $1$; by the Cauchy--Schwarz inequality
\[
|(C\mathsf{Z})(b_c)|^2\le2L^3\cdot\sum_b|\mathsf{Z}(b)|^2,
\]
the sum being over the bonds that occur in \eqref{eq:operators-photon-16} for the face $c$. A bond occurs there for at most two faces: an interior bond of direction $\mu$ of $B(z)$ for the faces $\langle z,z+L e_{\mu}\rangle$ and $\langle z-L e_{\mu},z\rangle$, and a face bond for its own face only. Summing over $c$,
\[
\sum_c|(C\mathsf{Z})(b_c)|^2\le4L^3\cdot|\mathsf{Z}|_2^2,\qquad\text{hence}\qquad|C\mathsf{Z}|_2^2\le(1+4L^3)\cdot|\mathsf{Z}|_2^2\le5L^3\cdot|\mathsf{Z}|_2^2.
\]

\emph{The bound in the $\ell^1$ norm.} Let $b$ be a face bond of the face $c$ other than $b_c$. By \eqref{eq:operators-photon-16}, $C\delta_b=\delta_b-\delta_{b_c}$, so that $\sum_{b'}|(C\delta_b)(b')|=2\le c_1$. Let $\Upsilon=e_i^y$. It is a unit vector carried by the $3L^2(L-1)\le3L^3$ interior bonds of $B(y)$, so the Cauchy--Schwarz inequality gives
\[
\sum_b|e_i^y(b)|\le(3L^3)^{\frac{1}{2}}\le2L^{\frac{3}{2}}.
\]
Off the central bonds $C e_i^y=e_i^y$. By \eqref{eq:operators-photon-16}, $C e_i^y$ vanishes at the central bond of every face of which $B(y)$ is not one of the two blocks, and at the central bonds of the six faces $\langle y,y+L e_{\mu}\rangle$, $\langle y-L e_{\mu},y\rangle$ ($\mu=0,1,2$) of $B(y)$ its modulus is at most $\sum_b|e_i^y(b)|\le2L^{\frac{3}{2}}$, the coefficients $\frac{n_c(b)}{L}$ being less than $1$. Hence
\[
\sum_b|(C e_i^y)(b)|\le2L^{\frac{3}{2}}+6\cdot2L^{\frac{3}{2}}=c_1.
\]

\emph{The supports.} Let $\Upsilon=e_i^y$. By the preceding paragraph, $C\Upsilon$ vanishes off the interior bonds of $B(y)$ and the central bonds of the six faces of $B(y)$. Each of these central bonds has one end in $B(y)$, and its other end has distance $1$ from that end. Since $\operatorname{supp}\Upsilon\subset B(y)$ and $B(y)$ has the diameter $L-1$, it follows that $d(x,x')\le(L-1)+1=L$ for $x\in\operatorname{supp}C\Upsilon$ and $x'\in\operatorname{supp}\Upsilon$. Next let $\Upsilon=\delta_b$, where $b\ne b_c$ is a face bond of the face $c=\langle y,y+L e_{\mu}\rangle$. As $C\delta_b=\delta_b-\delta_{b_c}$, the set $\operatorname{supp}C\Upsilon$ is formed by the end points of the two face bonds $b$ and $b_c$ of $c$. Write $b=(u,u+e_{\mu})$ and $b_c=(w,w+e_{\mu})$. The initial points $u$ and $w$ are sites of $B(y)$ with $u_{\mu}=w_{\mu}=y_{\mu}+\frac{L-1}{2}$, and the other two coordinates of $w$ are those of the centre $y$, because $b_c$ lies on the straight line through the two centres; hence $|u-w|\le\frac{L-1}{2}$. Consequently $d(x,x')\le\max\left\{\frac{L-1}{2},1\right\}\le L$ for any two points $x,x'$ among $u$, $u+e_{\mu}$, $w$, $w+e_{\mu}$.

(v) By Section~\ref{sec:averaging-fluctuations-gauge}, $d\varepsilon^0_k=-n_k^{-1}\cdot\log\mu_I(\mathrm{Box}_k)$. Under $\mu_I$ the coordinates $\zeta(\Upsilon)$ with respect to the orthonormal basis $\{\Upsilon\}$ are independent standard normal variables, and by (ii) there are $d_V\cdot n_k$ of them. Hence
\[
\mu_I(\mathrm{Box}_k)=P\left(|g|\le p_{0,k}\right)^{d_V\cdot n_k},
\]
and the formula for $d\varepsilon^0_k$ follows; as $P(|g|\le p_{0,k})\le1$, one has $d\varepsilon^0_k\ge0$. Let $p:=p_{0,k}\ge1$ and $t:=P(|g|>p)$. By the Gaussian tail bound, $t\le\exp\left(-\frac{p^2}{2}\right)$; moreover $t\le P(|g|>1)<0.32$. From the power series of the logarithm,
\[
-\log(1-t)=\sum_{n\ge1}\frac{t^n}{n}\le\frac{t}{1-t}\le1.5\,t.
\]
Since $d_V<2$,
\[
d\varepsilon^0_k=-d_V\cdot\log(1-t)\le3t\le3\exp\left(-\frac{p^2}{2}\right)\le12\exp\left(-\frac{p^2}{2}\right).
\]
Finally, for $e_k\in(0,e_0]$ condition~\eqref{eq:renormalisation-group-25} gives $v_k=-\log e_k\ge-\log e_0\ge1$, hence $p_{0,k}=v_k^{p_0}\ge1$ (Section~\ref{sec:running-quantities-tori}).
\end{proof}

For $l\in\mathrm{Cl}$, Proposition~\ref{prop:basis-map-box} follows from Proposition~\ref{prop:fluctuation-space}, with $c_C=3L^{\frac{9}{2}}$.

\subsection{The fluctuation form and its inverse square root}\label{sec:fluctuation-form-inverse}

Throughout this subsection $k\ge 0$, and the side $l$ is a multiple of $L$. The operator $\Delta_k=(dH_k)^TdH_k$ of Section~\ref{sec:averaging-fluctuations-gauge} acts on the bond functions on $T^0_k$, the transpose being taken with respect to the pairings of Section~\ref{sec:tori-cochains-kernels}. Let $C$ be the map of Proposition~\ref{prop:fluctuation-space}(iv), and let $C^T$ be its transpose with respect to the counting scalar products of $V_k$ and of the bond functions on $T^0_k$. The fluctuation form is the quadratic form $\mathsf{Z}\mapsto\langle\mathsf{Z},\tilde{\Delta}_k\mathsf{Z}\rangle$ on $V_k$, where
\[
\tilde{\Delta}_k:=C^T\Delta_kC .
\]
A linear operator $O$ on $V_k$ is identified with its matrix in the orthonormal basis $\{\Upsilon\}$ of $V_k$ (Proposition~\ref{prop:fluctuation-space}(ii)), namely $O(\Upsilon,\Upsilon'):=\langle\Upsilon,O\Upsilon'\rangle$. In particular
\[
\tilde{\Delta}_k(\Upsilon,\Upsilon')=\langle C\Upsilon,\Delta_kC\Upsilon'\rangle .
\]
Put
\begin{equation}\label{eq:operators-photon-17}
c_P:=\max\left\{9L^7\cdot C_Q^2,\ 15L^3\cdot\mathsf{A}_1\right\},\qquad C_{\mathrm{Dt}}:=c_1^2\cdot C_{\mathrm{De}}\cdot\exp(2\gamma_HL).
\end{equation}
Here $C_Q$ is the constant in the inequality $\left|Q_k^{(2)}F\right|_2\le C_Q|F|_2$ of Proposition~\ref{prop:block-averages}; one has $C_Q=1$ by Lemma~\ref{lem:block-averages}(f), and also by Proposition~\ref{prop:averaging-operators}(d). Further, $\gamma_H$ is the number of \eqref{eq:operators-photon-8}, $\mathsf{A}_1$ the number of \eqref{eq:operators-photon-9}, $C_{\mathrm{De}}$ the constant of Theorem~\ref{thm:decay-kernels}, and $c_1=14L^{\frac{3}{2}}$ the constant of Proposition~\ref{prop:fluctuation-space}(iv).

\begin{proposition}[Bounds on the fluctuation form]\label{prop:bounds-fluctuation-form}
For every $\mathsf{Z}\in V_k$
\[
c_P^{-1}|\mathsf{Z}|_2^2\le\langle\mathsf{Z},\tilde{\Delta}_k\mathsf{Z}\rangle\le c_P\,|\mathsf{Z}|_2^2,
\]
and for all elements $\Upsilon,\Upsilon'$ of the basis
\[
\left|\tilde{\Delta}_k(\Upsilon,\Upsilon')\right|\le C_{\mathrm{Dt}}\exp\left(-\gamma_H\,d(\Upsilon,\Upsilon')\right).
\]
\end{proposition}

\begin{proof}
Let $\mathsf{Z}\in V_k$ and $Z:=C\mathsf{Z}$. By Proposition~\ref{prop:fluctuation-space}(iv), $Z$ is a real bond function on $T^0_k$ with $QZ=0$ and $\tau Z=0$, and
\[
|\mathsf{Z}|_2\le|Z|_2\le\left(5L^3\right)^{\frac{1}{2}}|\mathsf{Z}|_2.
\]
The definitions of $\tilde{\Delta}_k$ and of $\Delta_k$ give
\[
\langle\mathsf{Z},\tilde{\Delta}_k\mathsf{Z}\rangle=\langle Z,\Delta_kZ\rangle=|dH_kZ|_2^2.
\]

\emph{The lower bound.} Since $Q_kH_k=1$ (Proposition~\ref{prop:minimiser}) and $d\,Q_k=Q_k^{(2)}d$ (Lemma~\ref{lem:block-averages}(f)), one has
\[
dZ=d\,Q_kH_kZ=Q_k^{(2)}\,dH_kZ,
\]
and the inequality of Proposition~\ref{prop:block-averages} (Lemma~\ref{lem:block-averages}(f)) yields $|dZ|_2\le C_Q\,|dH_kZ|_2$. Corollary~\ref{cor:bound-axial-gauge} applies to $Z$; together with $|\mathsf{Z}|_2\le|Z|_2$ it gives
\[
|dZ|_2^2\ge\frac{1}{9L^7}\,|Z|_2^2\ge\frac{1}{9L^7}\,|\mathsf{Z}|_2^2.
\]
Consequently
\[
\langle\mathsf{Z},\tilde{\Delta}_k\mathsf{Z}\rangle\ge C_Q^{-2}|dZ|_2^2\ge\frac{1}{9L^7\cdot C_Q^2}\,|\mathsf{Z}|_2^2\ge c_P^{-1}|\mathsf{Z}|_2^2.
\]

\emph{The upper bound.} By Theorem~\ref{thm:decay-kernels} the operator norm of $\Delta_k$ on $\ell^2$ is at most $3\mathsf{A}_1$. Therefore
\[
\langle\mathsf{Z},\tilde{\Delta}_k\mathsf{Z}\rangle=\langle Z,\Delta_kZ\rangle\le 3\mathsf{A}_1\cdot|Z|_2^2\le 15L^3\cdot\mathsf{A}_1\cdot|\mathsf{Z}|_2^2\le c_P\,|\mathsf{Z}|_2^2.
\]

\emph{The matrix.} Writing the pairing as a sum over the positively oriented unit bonds $b$, $b'$, one has
\[
\tilde{\Delta}_k(\Upsilon,\Upsilon')=\sum_{b,b'}(C\Upsilon)(b)\cdot\Delta_k(b,b')\cdot(C\Upsilon')(b').
\]
Let $b,b'$ be bonds with $(C\Upsilon)(b)\ne 0$ and $(C\Upsilon')(b')\ne 0$, and choose an end $p$ of $b$ and an end $p'$ of $b'$ with $d(p,p')=d(b,b')$. Then $p$ lies in $\operatorname{supp}C\Upsilon$ and $p'$ in $\operatorname{supp}C\Upsilon'$, the sets of the end points of the bonds on which $C\Upsilon$, resp.\ $C\Upsilon'$, does not vanish. By Proposition~\ref{prop:fluctuation-space}(iv) the point $p$ has distance at most $L$ from every point of $\operatorname{supp}\Upsilon$, and $p'$ has distance at most $L$ from every point of $\operatorname{supp}\Upsilon'$; by the triangle inequality,
\[
d(\Upsilon,\Upsilon')\le d(b,b')+2L .
\]
For all such pairs $(b,b')$, \eqref{eq:operators-photon-11} therefore gives
\[
|\Delta_k(b,b')|\le C_{\mathrm{De}}\exp(-\gamma_H\,d(b,b'))\le C_{\mathrm{De}}\exp(2\gamma_HL)\exp(-\gamma_H\,d(\Upsilon,\Upsilon')).
\]
Since $\sum_b|(C\Upsilon)(b)|\le c_1$ for every $\Upsilon$ (Proposition~\ref{prop:fluctuation-space}(iv)), it follows that
\[
\left|\tilde{\Delta}_k(\Upsilon,\Upsilon')\right|\le C_{\mathrm{De}}\exp(2\gamma_HL)\exp(-\gamma_H\,d(\Upsilon,\Upsilon'))\sum_b|(C\Upsilon)(b)|\sum_{b'}|(C\Upsilon')(b')|\le C_{\mathrm{Dt}}\exp(-\gamma_H\,d(\Upsilon,\Upsilon')).
\]
\end{proof}

\begin{lemma}[A Combes--Thomas estimate]\label{lem:combes-thomas-estimate}
Let $I$ be a finite set, and let $\mathbb{R}^I$ carry the Euclidean scalar product $\langle\cdot,\cdot\rangle$ and norm $|\cdot|_2$. Let $\boldsymbol{\varphi}\ge 0$, and let $d\colon I\times I\to[0,\infty)$ be a symmetric function with $d(a,a)=0$ and
\[
d(a,c)\le d(a,b)+d(b,c)+\boldsymbol{\varphi}\qquad(a,b,c\in I).
\]
Let $T$ be a real symmetric matrix on $I$, and let $t_0>0$, $C_T$ and $g>0$ be numbers such that $T\ge t_0$, that is, $\langle v,Tv\rangle\ge t_0|v|_2^2$ for all $v\in\mathbb{R}^I$, and
\[
|T(a,b)|\le C_T\exp(-g\,d(a,b))\qquad(a,b\in I).
\]
Let $S_1>0$ be a number with
\[
S_1\ge\exp\left(\frac{g\cdot\boldsymbol{\varphi}}{2}\right)\sup_{a\in I}\ \sum_{b\in I}\left(d(a,b)+\boldsymbol{\varphi}\right)\exp\left(-\frac{g\,d(a,b)}{2}\right),
\]
and put
\[
\boldsymbol{\theta}:=\min\left\{\frac{g}{2},\ \frac{t_0}{2C_T\cdot S_1}\right\}
\]
(here $C_T\ge T(a,a)\ge t_0$ for every $a$, so that $\boldsymbol{\theta}>0$). Then for every $x\ge 0$ and all $a,c\in I$
\[
\left|(T+x)^{-1}(a,c)\right|\le\left(\frac{t_0}{2}+x\right)^{-1}\exp\left(-\boldsymbol{\theta}\cdot d(a,c)\right)
\]
and
\[
\left|T^{-\frac{1}{2}}(a,c)\right|\le\left(\frac{2}{t_0}\right)^{\frac{1}{2}}\exp\left(-\boldsymbol{\theta}\cdot d(a,c)\right).
\]
Here $T+x$ stands for $T+x\cdot 1$, with the identity matrix $1$ on $I$, and $T^{-\frac{1}{2}}$ is the positive definite square root of $T^{-1}$.
\end{lemma}

\begin{proof}
For a matrix $M$ on $I$ let $\|M\|$ be its operator norm on $\mathbb{R}^I$; then $|M(a,b)|\le\|M\|$ for all $a,b\in I$.

Fix $c\in I$, and put $\mathbf{f}(a):=d(a,c)$ for $a\in I$. The hypotheses on $d$ give $\mathbf{f}(a)\le d(a,b)+\mathbf{f}(b)+\boldsymbol{\varphi}$ and $\mathbf{f}(b)\le d(a,b)+\mathbf{f}(a)+\boldsymbol{\varphi}$, that is,
\[
|\mathbf{f}(a)-\mathbf{f}(b)|\le d(a,b)+\boldsymbol{\varphi}\qquad(a,b\in I).
\]
Let $\tilde{W}$ be the diagonal matrix with the diagonal entries $\exp(\boldsymbol{\theta}\cdot\mathbf{f}(a))$, $a\in I$, and put $E:=\tilde{W} T\tilde{W}^{-1}-T$. The entries of $E$ are
\[
E(a,b)=T(a,b)\left(\exp\left(\boldsymbol{\theta}\cdot\left(\mathbf{f}(a)-\mathbf{f}(b)\right)\right)-1\right).
\]
For real numbers $u,t$ with $|u|\le t$ one has
\[
|\exp(u)-1|\le\exp(t)-1\le t\exp(t);
\]
indeed, for $u<0$ one has $1-\exp(u)\le 1-\exp(-t)\le\exp(t)-1$, and $\exp(t)-1\le t\exp(t)$ by the mean value theorem. This is applied with $u:=\boldsymbol{\theta}\cdot(\mathbf{f}(a)-\mathbf{f}(b))$ and $t:=\boldsymbol{\theta}\cdot(d(a,b)+\boldsymbol{\varphi})$; with $\boldsymbol{\theta}\le\frac{g}{2}$ it gives
\begin{align*}
|E(a,b)|&\le C_T\exp(-g\,d(a,b))\cdot\boldsymbol{\theta}\cdot\left(d(a,b)+\boldsymbol{\varphi}\right)\exp\left(\boldsymbol{\theta}\cdot\left(d(a,b)+\boldsymbol{\varphi}\right)\right)\\
&\le\boldsymbol{\theta}\cdot C_T\exp\left(\frac{g\cdot\boldsymbol{\varphi}}{2}\right)\left(d(a,b)+\boldsymbol{\varphi}\right)\exp\left(-\frac{g\,d(a,b)}{2}\right)=:K(a,b).
\end{align*}
The matrix $K$ has non-negative entries, and it is symmetric because $d$ is. By the Schur test, every matrix $M$ with $|M(a,b)|\le K(a,b)$ for all $a,b$ satisfies $\|M\|\le\sup_a\sum_bK(a,b)$ (by the Cauchy--Schwarz inequality, $|\langle v,Mw\rangle|\le\left(\sum_{a,b}K(a,b)v_a^2\right)^{\frac{1}{2}}\left(\sum_{a,b}K(a,b)w_b^2\right)^{\frac{1}{2}}$). With the choice of $S_1$ and $\boldsymbol{\theta}\le\frac{t_0}{2C_T\cdot S_1}$ it follows that
\[
\|E\|\le\sup_{a\in I}\ \sum_{b\in I}K(a,b)\le\boldsymbol{\theta}\cdot C_T\cdot S_1\le\frac{t_0}{2}.
\]

Let $x\ge 0$. The matrix $T+x$ is symmetric and $T+x\ge t_0+x$; hence it is invertible, $\left\|(T+x)^{-1}\right\|\le(t_0+x)^{-1}$, and
\[
\left\|(T+x)^{-1}E\right\|\le\frac{t_0}{2(t_0+x)}<1.
\]
By the Neumann series the matrix $1+(T+x)^{-1}E$ is invertible, and the norm of its inverse is at most $\left(1-\left\|(T+x)^{-1}E\right\|\right)^{-1}\le\frac{2(t_0+x)}{t_0+2x}$. So $T+x+E=(T+x)\left(1+(T+x)^{-1}E\right)$ is invertible, and
\[
\left\|(T+x+E)^{-1}\right\|\le\frac{2(t_0+x)}{t_0+2x}\cdot\frac{1}{t_0+x}=\left(\frac{t_0}{2}+x\right)^{-1}.
\]
Since $\tilde{W}$ is diagonal, $\tilde{W}(T+x)\tilde{W}^{-1}=T+x+E$, that is,
\[
(T+x)^{-1}=\tilde{W}^{-1}(T+x+E)^{-1}\tilde{W} .
\]
The entry $(a,c)$ of the right side equals $\exp(-\boldsymbol{\theta}\cdot\mathbf{f}(a))\,(T+x+E)^{-1}(a,c)\,\exp(\boldsymbol{\theta}\cdot\mathbf{f}(c))$, where $\mathbf{f}(c)=d(c,c)=0$ and $\mathbf{f}(a)=d(a,c)$. Hence
\[
\left|(T+x)^{-1}(a,c)\right|\le\left(\frac{t_0}{2}+x\right)^{-1}\exp\left(-\boldsymbol{\theta}\cdot d(a,c)\right).
\]
As $c$ was arbitrary, this proves the first bound.

For every $t>0$
\[
\frac{1}{\pi}\int_0^{\infty}\mathrm{d}x\,x^{-\frac{1}{2}}(t+x)^{-1}=t^{-\frac{1}{2}},
\]
since the substitution $x=ts^2$ turns the left side into $\frac{2}{\pi}\,t^{-\frac{1}{2}}\int_0^{\infty}\mathrm{d}s\,(1+s^2)^{-1}$. By the spectral theorem for symmetric matrices, $T=\sum_jt_jP_j$ with eigenvalues $t_j\ge t_0$ and orthogonal projections $P_j$ with $\sum_jP_j=1$; then $(T+x)^{-1}=\sum_j(t_j+x)^{-1}P_j$ and $T^{-\frac{1}{2}}=\sum_jt_j^{-\frac{1}{2}}P_j$, and the scalar identity, applied to each $t_j$, gives
\[
T^{-\frac{1}{2}}=\frac{1}{\pi}\int_0^{\infty}\mathrm{d}x\,x^{-\frac{1}{2}}(T+x)^{-1}.
\]
The first bound and the scalar identity at $t=\frac{t_0}{2}$ now yield
\[
\left|T^{-\frac{1}{2}}(a,c)\right|\le\frac{1}{\pi}\int_0^{\infty}\mathrm{d}x\,x^{-\frac{1}{2}}\left(\frac{t_0}{2}+x\right)^{-1}\exp\left(-\boldsymbol{\theta}\cdot d(a,c)\right)=\left(\frac{2}{t_0}\right)^{\frac{1}{2}}\exp\left(-\boldsymbol{\theta}\cdot d(a,c)\right).
\]
\end{proof}

Put
\begin{multline}\label{eq:operators-photon-18}
S_P:=\exp\left(\frac{\gamma_HL}{2}\right)\left(L+\frac{2}{\gamma_H}\right)\cdot c_{\mathrm{bas}}\left(\frac{\gamma_H}{4}\right),\qquad\gamma_P:=\min\left\{\frac{\gamma_H}{2},\ \frac{1}{2c_P\cdot C_{\mathrm{Dt}}\cdot S_P}\right\},\\
C_P:=\max\left\{C_{\mathrm{Hk}}(\alpha),\ C_{\mathrm{Dt}},\ (2c_P)^{\frac{1}{2}},\ 4\exp(2)\right\},
\end{multline}
where $c_{\mathrm{bas}}$ is the function of \eqref{eq:operators-photon-1}, and $C_{\mathrm{Hk}}(\alpha)$, $0\le\alpha<1$, is the constant of Theorem~\ref{thm:decay-kernels}.

\begin{proposition}[Decay of the resolvents and of the root]\label{prop:decay-resolvents-root}
The operator $\tilde{\Delta}_k$ is invertible, and the operator
\[
C_k^{\frac{1}{2}}=\frac{1}{\pi}\int_0^{\infty}\mathrm{d}x\,x^{-\frac{1}{2}}(\tilde{\Delta}_k+x)^{-1}
\]
of Section~\ref{sec:averaging-fluctuations-gauge} is the positive definite square root of $C_k=\tilde{\Delta}_k^{-1}$. For all elements $\Upsilon,\Upsilon'$ of the basis and every $x\ge 0$
\[
\left|(\tilde{\Delta}_k+x)^{-1}(\Upsilon,\Upsilon')\right|\le\left(\frac{1}{2c_P}+x\right)^{-1}\exp\left(-\gamma_P\,d(\Upsilon,\Upsilon')\right)
\]
and
\[
\left|C_k^{\frac{1}{2}}(\Upsilon,\Upsilon')\right|\le(2c_P)^{\frac{1}{2}}\exp\left(-\gamma_P\,d(\Upsilon,\Upsilon')\right).
\]
\end{proposition}

\begin{proof}
Let $I$ be the basis $\{\Upsilon\}$ of $V_k$, a finite set, and let $d(\Upsilon,\Upsilon')$ be the distance of the supports (Section~\ref{sec:tori-cochains-kernels}); this function is symmetric and vanishes on the diagonal. By Proposition~\ref{prop:fluctuation-space}(ii) the support of every $\Upsilon$ lies in one $L$-block or consists of the two ends of one bond, so that its diameter is at most $L-1$. For $\Upsilon_1,\Upsilon_2,\Upsilon_3\in I$ choose $p\in\operatorname{supp}\Upsilon_1$ and $q\in\operatorname{supp}\Upsilon_2$ with $d(p,q)=d(\Upsilon_1,\Upsilon_2)$, and $q'\in\operatorname{supp}\Upsilon_2$ and $r\in\operatorname{supp}\Upsilon_3$ with $d(q',r)=d(\Upsilon_2,\Upsilon_3)$; then
\[
d(\Upsilon_1,\Upsilon_3)\le d(p,q)+d(q,q')+d(q',r)\le d(\Upsilon_1,\Upsilon_2)+d(\Upsilon_2,\Upsilon_3)+L .
\]
Thus $d$ satisfies the hypotheses of Lemma~\ref{lem:combes-thomas-estimate} with $\boldsymbol{\varphi}:=L$.

Let $T$ be the matrix of $\tilde{\Delta}_k$ in the orthonormal basis $\{\Upsilon\}$ (Proposition~\ref{prop:fluctuation-space}(ii)). It is real and symmetric, because $\langle\mathsf{Z},\tilde{\Delta}_k\mathsf{Z}'\rangle=\langle dH_kC\mathsf{Z},dH_kC\mathsf{Z}'\rangle$ and $H_k$ and $C$ map real fields to real fields (Proposition~\ref{prop:minimiser} and \eqref{eq:operators-photon-16}). By Proposition~\ref{prop:bounds-fluctuation-form}, and since $|\mathsf{Z}|_2^2=\sum_{\Upsilon}\mathsf{Z}(\Upsilon)^2$, the hypotheses of Lemma~\ref{lem:combes-thomas-estimate} on $T$ hold with
\[
t_0:=\frac{1}{c_P},\qquad C_T:=C_{\mathrm{Dt}},\qquad g:=\gamma_H.
\]
For $t\ge 0$ one has $t\exp\left(-\frac{gt}{4}\right)\le\frac{4}{\exp(1)\,g}\le\frac{2}{g}$, the maximum of the left side being attained at $t=\frac{4}{g}$; hence
\[
(t+L)\exp\left(-\frac{gt}{2}\right)\le\left(L+\frac{2}{g}\right)\exp\left(-\frac{gt}{4}\right)\qquad(t\ge 0).
\]
By \eqref{eq:operators-photon-8}, $0<\gamma_H<1$; so \eqref{eq:operators-photon-2} applies with $\frac{g}{4}$ in place of $g$, and
\[
\exp\left(\frac{gL}{2}\right)\sup_{\Upsilon}\ \sum_{\Upsilon'}\left(d(\Upsilon,\Upsilon')+L\right)\exp\left(-\frac{g\,d(\Upsilon,\Upsilon')}{2}\right)\le\exp\left(\frac{gL}{2}\right)\left(L+\frac{2}{g}\right)\cdot c_{\mathrm{bas}}\left(\frac{g}{4}\right)=S_P.
\]
So the number $S_1:=S_P>0$ is admissible in Lemma~\ref{lem:combes-thomas-estimate}, and by \eqref{eq:operators-photon-18}
\[
\boldsymbol{\theta}=\min\left\{\frac{\gamma_H}{2},\ \frac{1}{2c_P\cdot C_{\mathrm{Dt}}\cdot S_P}\right\}=\gamma_P,\qquad\frac{t_0}{2}=\frac{1}{2c_P},\qquad\left(\frac{2}{t_0}\right)^{\frac{1}{2}}=(2c_P)^{\frac{1}{2}}.
\]
Since $T\ge c_P^{-1}>0$, the operator $\tilde{\Delta}_k$ is invertible. By the last part of the proof of Lemma~\ref{lem:combes-thomas-estimate}, applied to $T$, the integral $\frac{1}{\pi}\int_0^{\infty}\mathrm{d}x\,x^{-\frac{1}{2}}(\tilde{\Delta}_k+x)^{-1}$ equals $\tilde{\Delta}_k^{-\frac{1}{2}}$, the positive definite square root of $C_k=\tilde{\Delta}_k^{-1}$. The matrices of $(\tilde{\Delta}_k+x)^{-1}$ and of $C_k^{\frac{1}{2}}$ in the basis $\{\Upsilon\}$ are $(T+x)^{-1}$ and $T^{-\frac{1}{2}}$, and Lemma~\ref{lem:combes-thomas-estimate} gives the two bounds.
\end{proof}

\paragraph{The constants $C_P$, $\gamma_P$, $c_P$}
By \eqref{eq:operators-photon-18}, $\gamma_P\le\frac{\gamma_H}{2}\le\gamma_H$, and $C_P$ is at least each of the numbers $C_{\mathrm{Hk}}(\alpha)$, $C_{\mathrm{Dt}}$ and $(2c_P)^{\frac{1}{2}}$. Since $\exp(-\gamma_Hs)\le\exp(-\gamma_Ps)$ for $s\ge 0$, the bounds \eqref{eq:operators-photon-11} on the kernels of $H_k$, $\partial H_k$ and $\delta_{\alpha}\partial H_k$, the bound of Proposition~\ref{prop:bounds-fluctuation-form} on the matrix of $\tilde{\Delta}_k$ and the bound of Proposition~\ref{prop:decay-resolvents-root} on the matrix of $C_k^{\frac{1}{2}}$ remain valid when their right sides are replaced by $C_P\exp(-\gamma_P\,d(b,b'))$, resp.\ $C_P\exp(-\gamma_P\,d(\Upsilon,\Upsilon'))$. Moreover $\gamma_P\in(0,1]$, because $0<\gamma_P\le\frac{\gamma_H}{2}$ and $\gamma_H<1$ by \eqref{eq:operators-photon-8}, and $C_P\ge 4\exp(2)\ge 1$. The numbers $\gamma_P$ and $c_P$ depend on $L$ only, and $C_P$ depends on $(L,\alpha)$ only: $C_Q=1$; the quantities $\gamma_H$, $\mathsf{A}_1$ and $C_{\mathrm{De}}$ are absolute constants (\eqref{eq:operators-photon-8}, \eqref{eq:operators-photon-9}, Theorem~\ref{thm:decay-kernels}); $c_1=14L^{\frac{3}{2}}$ and $c_{\mathrm{bas}}\left(\frac{\gamma_H}{4}\right)$ depend on $L$ only; and $C_{\mathrm{Hk}}(\alpha)$ depends on $\alpha$ only. Propositions~\ref{prop:kernel-bounds}, \ref{prop:spectral-bounds-decay} and~\ref{prop:decay-exact-root} follow, in the form displayed in Section~\ref{sec:properties-exact-photon}, with the constant $c_P$ of \eqref{eq:operators-photon-17} and the constants $C_P$, $\gamma_P$ of \eqref{eq:operators-photon-18}.

At the level $k=0$ one has $\eta=1$, the fine bonds are the unit bonds, and $H_0=1$ (Proposition~\ref{prop:minimiser}); for bonds $b,\tilde b$ as in \eqref{eq:operators-photon-11} the initial points $x,\tilde x$ satisfy $|x-\tilde x|=1$, so that $|x-\tilde x|^{-\alpha}=1$. Hence the entries of the kernels of $H_0$, $\partial H_0$ and $\delta_{\alpha}\partial H_0$ have modulus at most $1$, $1$ and $2$, respectively, and they vanish unless $d(b,b')$ is at most $0$, $1$ and $2$, respectively. These entries are bounded by $C_P\exp(-\gamma_P\,d(b,b'))$ as well, because $C_P\exp(-\gamma_Ps)\ge 4\exp(2)\exp(-2)=4$ for $0\le s\le 2$.

\subsection{Constrained minimisation, covariance and coverings}\label{sec:constrained-minimisation-covariance}

\subsubsection{The constrained minimiser}\label{sec:constrained-minimiser}

\begin{lemma}[Existence, uniqueness and characterisation]\label{lem:existence-uniqueness}
Let $l$ be a multiple of $L$, and let $A'$ be a real bond function on $T^0_{k+1}$, the unit lattice of level $k+1$ and side $l/L$. Let $A'_L$ be the bond function on $L\mathbb{Z}^3/l\mathbb{Z}^3$ given by $A'_L(b):=L^{-\frac{1}{2}}A'(b/L)$, and put
\[
\mathrm{Adm}(A'):=\left\{A\ \text{a real bond function on}\ T^0_k:\ QA=A'_L,\ \tau A=0\right\}.
\]
Then $\mathrm{Adm}(A')$ is a non-empty affine space with the direction
\[
\left\{Z:\ QZ=0,\ \tau Z=0\right\}=C(V_k).
\]
On $\mathrm{Adm}(A')$ the function $A\mapsto\langle A,\Delta_kA\rangle$ has exactly one minimiser, denoted $A^{\min}_k(A')$. An element $A$ of $\mathrm{Adm}(A')$ is this minimiser if and only if
\[
\langle A,\Delta_kZ\rangle=0\qquad\text{for every}\ Z\in C(V_k).
\]
The map $A'\mapsto A^{\min}_k(A')$ is linear.
\end{lemma}

The existence and the uniqueness of the minimiser and the linearity of $A^{\min}_k$ form Proposition~\ref{prop:minimisers-prescribed-block}(i), and the relation $\langle A^{\min}_k(A'),\Delta_kZ\rangle=0$ for every $Z$ with $QZ=0$ and $\tau Z=0$ is Proposition~\ref{prop:minimisers-prescribed-block}(ii); Theorem~\ref{thm:minimisation-one-step}(a) contains these assertions together with the characterisation of the minimiser.

\begin{proof}
For a face $c=\langle y,y+Le_{\mu}\rangle$ let $A_0$ take the value $L\cdot A'_L(y,y+Le_{\mu})$ on each of the $L^2$ face bonds of $c$, oriented in the direction of $e_{\mu}$, and let $A_0$ vanish on the interior bonds; this defines $A_0$ on every unit bond, since a unit bond is either interior or a face bond of exactly one face. Let $x\in B(y)$. Among the $L$ bonds of the straight path $\Gamma_{x,x+Le_{\mu}}$ exactly one is a face bond of $c$; the others are interior bonds of $B(y)$ or of $B(y+Le_{\mu})$. Therefore $A_0(\Gamma_{x,x+Le_{\mu}})=L\cdot A'_L(y,y+Le_{\mu})$, and summation over the $L^3$ sites of $B(y)$ gives
\[
(QA_0)(y,y+Le_{\mu})=L^{-4}\cdot L^3\cdot L\cdot A'_L(y,y+Le_{\mu})=A'_L(y,y+Le_{\mu}).
\]
For $x\in B(y)$ the rectilinear paths $\Gamma^{\pi}(y,x)$ remain in the box $B(y)$ and consist of interior bonds; hence $\tau A_0=0$, and $A_0\in\mathrm{Adm}(A')$. As $Q$ and $\tau$ are linear,
\[
\mathrm{Adm}(A')=A_0+\left\{Z:\ QZ=0,\ \tau Z=0\right\},
\]
and by Proposition~\ref{prop:fluctuation-space}(iv) the map $C$ is a linear bijection of $V_k$ onto $\{Z:\ QZ=0,\ \tau Z=0\}$.

Let $Z=C\mathsf{Z}$ with $\mathsf{Z}\in V_k$. Since $\tilde{\Delta}_k=C^T\Delta_kC$, Proposition~\ref{prop:bounds-fluctuation-form} gives
\[
\langle Z,\Delta_kZ\rangle=\langle\mathsf{Z},\tilde{\Delta}_k\mathsf{Z}\rangle\ge c_P^{-1}\cdot|\mathsf{Z}|_2^2,
\]
and $\mathsf{Z}\ne0$ whenever $Z\ne0$, because $C$ is linear and therefore maps $0$ to $0$; so the form $Z\mapsto\langle Z,\Delta_kZ\rangle$ is positive definite on $C(V_k)$. The operator $\Delta_k=(dH_k)^TdH_k$ is symmetric. Hence, for $A\in\mathrm{Adm}(A')$, $Z\in C(V_k)$ and real $t$,
\[
\langle A+tZ,\Delta_k(A+tZ)\rangle=\langle A,\Delta_kA\rangle+2t\cdot\langle A,\Delta_kZ\rangle+t^2\cdot\langle Z,\Delta_kZ\rangle.
\]
Taken at $A=A_0$ and $t=1$, this identity shows that $Z\mapsto\langle A_0+Z,\Delta_k(A_0+Z)\rangle$ is a quadratic function on the finite-dimensional space $C(V_k)$ whose quadratic part is positive definite. Such a function is strictly convex and tends to infinity at infinity; so it has exactly one minimiser, and the form has exactly one minimiser on $\mathrm{Adm}(A')$. If $A$ is the minimiser, then for every $Z$ the right side is minimal at $t=0$, whence $\langle A,\Delta_kZ\rangle=0$. Conversely, if $\langle A,\Delta_kZ\rangle=0$ for every $Z\in C(V_k)$, the identity at $t=1$ gives $\langle A+Z,\Delta_k(A+Z)\rangle\ge\langle A,\Delta_kA\rangle$ for every $Z$, and $A$ is the minimiser. The characterisation is linear in the pair $(A,A')$: if $A_j=A^{\min}_k(A'_j)$ for $j=1,2$ and $a_1,a_2$ are real numbers, then $a_1A_1+a_2A_2$ lies in $\mathrm{Adm}(a_1A'_1+a_2A'_2)$ and satisfies $\langle a_1A_1+a_2A_2,\Delta_kZ\rangle=0$ for every $Z\in C(V_k)$; so it equals $A^{\min}_k(a_1A'_1+a_2A'_2)$.
\end{proof}

\subsubsection{Commutation with the lattice group}\label{sec:commutation-lattice-group}

\begin{proposition}[Lattice symmetries]\label{prop:lattice-symmetries}
Let $l\in\mathrm{Cl}$ and $g\in\mathbb{G}_k$. Then $g$ commutes with $d$, $\delta$ and $\Delta$, with the averages $Q_k$ on scalars and on bond functions and $Q_k^{(2)}$ on plaquette functions, with $Q$ and with $\mathscr{R}_k$; and
\[
(\tau gA)(gy,gx)=(\tau A)(y,x)
\]
for every bond function $A$ on $T^0_k$, every $y\in L\mathbb{Z}^3/l\mathbb{Z}^3$ and every $x\in B(y)$. For every real number $x\ge0$,
\begin{gather*}
gH_k=H_kg,\qquad g\Delta_k=\Delta_kg,\qquad gC=Cg,\qquad g\tilde{\Delta}_k=\tilde{\Delta}_kg,\\
g(\tilde{\Delta}_k+x)^{-1}=(\tilde{\Delta}_k+x)^{-1}g,\qquad gC_k^{\frac{1}{2}}=C_k^{\frac{1}{2}}g,\qquad g\mathrm{Sq}_k=\mathrm{Sq}_kg.
\end{gather*}
For $g=(S,a)\in\mathbb{G}_k^L$ let $g'=(S,a/L)$ be the element of $\mathbb{G}_{k+1}$, the group being formed with the side $l/L$, given by Proposition~\ref{prop:properties-actions}(d). Then
\[
A^{\min}_k(g'A')=g\,A^{\min}_k(A')
\]
for every real bond function $A'$ on $T^0_{k+1}$. All the operators named in this proposition are real. The set $\mathrm{Box}_k$ and the measures $\mu_I$ and $\hat{\mu}$ are invariant under $\mathbb{G}_k$ and under $\mathsf{Z}\mapsto-\mathsf{Z}$.
\end{proposition}

\begin{proof}
Write $g=(S,a)$ and $r:=r(S)$. By Section~\ref{sec:dirac-matrices-lattice}, $r$ is a diagonal matrix with diagonal entries $\pm1$, $a\in M\mathbb{Z}^3$, and $gx=rx+a$ for every point $x$. On functions, $g$ acts through composition with $g^{-1}$ (Section~\ref{sec:dirac-matrices-lattice}):
\[
(gA)(x,x')=A(g^{-1}x,g^{-1}x'),\qquad(g\boldsymbol{\lambda})(x)=\boldsymbol{\lambda}(g^{-1}x),\qquad(gF)(p)=F(g^{-1}p),
\]
for a bond function $A$, a scalar $\boldsymbol{\lambda}$ and a plaquette function $F$; here $p$ is an oriented plaquette, and $F(p)$ changes its sign when the orientation of $p$ is reversed. Proposition~\ref{prop:properties-actions}(a), (b) gives the following properties of $g$. It induces a permutation of the oriented cells of $T_k$, and one of the oriented cells of $T^0_k$. The cube $B_k(y)$ is carried onto $B_k(gy)$, the $L$-block $B(y)$ onto $B(gy)$, every segment onto a segment, and the path $\Gamma^{\pi}(y,x)$ onto the path $\Gamma^{\pi}(gy,gx)$ with the same order $\pi$. The pairings are invariant under $g$; consequently $g^T=g^{-1}$.

\emph{The operators $d$, $\delta$, $\Delta$.} Up to a constant factor, the value of $d$ at a cell is the sum of the values at its faces, and the value of $\delta$ at a cell is the sum of the values at its cofaces, the incidence numbers being the coefficients. The map $g$ preserves the incidence numbers. Hence $g$ commutes with $d$ and with $\delta$, and also with $\Delta$, which is $\delta d$ on scalars and $\delta d+d\delta$ on bond functions.

\emph{The averages.} Let $\mathscr{A}$ be a bond function on $T_k$ and $y$ a unit site. The straight path $\Gamma_{x,x+e_{\mu}}$ consists of the $\mathfrak{N}$ bonds with the initial points $x+j\eta e_{\mu}$, $0\le j\le\mathfrak{N}-1$, and $\#|B_k(y)|=\mathfrak{N}^3=\eta^{-3}$. By the definition of $Q_k$ in Section~\ref{sec:block-averages-axial}, therefore,
\[
(Q_k\mathscr{A})(y,y+e_{\mu})=\sum_{x\in B_k(y)}\eta^3\cdot\eta\cdot\mathscr{A}(\Gamma_{x,x+e_{\mu}})
\]
is the mean over $x\in B_k(y)$ of $\eta\cdot\mathscr{A}(\Gamma_{x,x+e_{\mu}})$. The value at the reversed bond, $(Q_k\mathscr{A})(y+e_{\mu},y)$, is the negative of this number. A line sum changes its sign with the orientation of the path, and $B_k(y)+e_{\mu}=B_k(y+e_{\mu})$; so this value is the mean over $x''=x+e_{\mu}\in B_k(y+e_{\mu})$ of $\eta\cdot\mathscr{A}(\Gamma_{x'',x''-e_{\mu}})$. Consequently, for both orientations of a unit bond $(y,y')$,
\[
(Q_k\mathscr{A})(y,y')=\sum_{x\in B_k(y)}\eta^4\cdot\mathscr{A}(\Gamma_{x,x+y'-y}).
\]
As $g(x+y'-y)=gx+r(y'-y)=gx+gy'-gy$, the map $g$ carries the path $\Gamma_{x,x+y'-y}$ onto $\Gamma_{gx,gx+gy'-gy}$, and it carries $B_k(y)$ onto $B_k(gy)$; so $(Q_kg\mathscr{A})(gy,gy')=(Q_k\mathscr{A})(y,y')$, that is, $Q_kg=gQ_k$ on bond functions. Each of the other averages is likewise a sum over the translates of one configuration:
\begin{align*}
(Q_k\lambda)(y)&=\sum_{x\in B_k(y)}\eta^3\cdot\lambda(x),\\
(QA)(y,y')&=L^{-4}\sum_{x\in B(y)}A(\Gamma_{x,x+y'-y})\qquad(y'-y=\pm Le_{\mu}),\\
\left(Q_k^{(2)}F\right)(p)&=\sum_{x\in B_k(y)}\eta^5\cdot F(p+x-y)\qquad(y\ \text{a corner of}\ p).
\end{align*}
In the second line the case $y'-y=-Le_{\mu}$ follows from the definition of $Q$ by the argument just given for $Q_k$, because $B(y)+Le_{\mu}=B(y+Le_{\mu})$. In the third line $p$ is an oriented unit plaquette, and $F(p+x-y)$ is the sum of the values of $F$ at the $\mathfrak{N}^2$ fine plaquettes that tile the translate $p+x-y$, each oriented as $p$; the set of the translates $p+x-y$, $x\in B_k(y)$, is the same for the four corners $y$ of $p$, because $B_k(y+w)=B_k(y)+w$ for $w\in\mathbb{Z}^3$. The map $g$ carries $B_k(y)$, $B(y)$, the paths and the translates of $p$ onto the corresponding objects formed with $gy$, $gy'$ and $gp$. Hence $g$ commutes with $Q_k$ on scalars, with $Q$ and with $Q_k^{(2)}$.

\emph{The functional $\tau$.} For $x\in B(y)$ one has $gx\in B(gy)$, and $g$ maps the six paths $\Gamma^{\pi}(y,x)$ onto the paths $\Gamma^{\pi}(gy,gx)$; hence
\[
(\tau gA)(gy,gx)=\frac{1}{3!}\sum_{\pi}(gA)(\Gamma^{\pi}(gy,gx))=\frac{1}{3!}\sum_{\pi}A(\Gamma^{\pi}(y,x))=(\tau A)(y,x).
\]

\emph{The projection $\mathscr{R}_k$.} As $g$ commutes with $Q_k$ and with $\Delta$ on scalars, it maps $\ker Q_k$ and $\Delta(\ker Q_k)$ onto themselves. Since $g$ is orthogonal, it commutes with the orthogonal projection $\mathscr{R}_k$ onto $\Delta(\ker Q_k)$.

\emph{The minimiser $H_k$.} Let $A$ be a real bond function on $T^0_k$, and let $\mathscr{A}$ be a real bond function on $T_k$ with $Q_k\mathscr{A}=A$ and $\mathscr{R}_k\delta\mathscr{A}=0$. Then
\[
Q_kg\mathscr{A}=gA,\qquad\mathscr{R}_k\delta g\mathscr{A}=g\mathscr{R}_k\delta\mathscr{A}=0,\qquad|dg\mathscr{A}|_2=|gd\mathscr{A}|_2=|d\mathscr{A}|_2,
\]
and the same holds with $g^{-1}$ in place of $g$. So $g$ maps the problem of Section~\ref{sec:averaging-fluctuations-gauge} with the datum $A$ onto the problem with the datum $gA$, and the uniqueness of the solution (Proposition~\ref{prop:minimiser}) gives $gH_kA=H_kgA$. By linearity, $gH_k=H_kg$.

\emph{The operator $\Delta_k$.} From $gdH_k=dH_kg$ and $g^T=g^{-1}$ it follows that $g(dH_k)^T=(dH_k)^Tg$, and therefore
\[
g\Delta_k=g(dH_k)^TdH_k=(dH_k)^TdH_kg=\Delta_kg.
\]

\emph{The map $C$.} Let $\mathsf{Z}\in V_k$. By Proposition~\ref{prop:fluctuation-space}(ii), $g\mathsf{Z}\in V_k$. Moreover, $g$ maps central bonds to central bonds. Indeed, $g$ carries each $L$-block $B(y)$ onto $B(gy)$, hence the centre of a block onto the centre of its image; since $g(y+Le_{\mu})=gy+r(Le_{\mu})=gy\pm Le_{\mu}$, it carries the two blocks of the face $\langle y,y+Le_{\mu}\rangle$ onto the two blocks of a face, and therefore the face bonds of the first face onto face bonds of the second. The central bond of a face is its face bond on the straight line through the centres of its two blocks (Section~\ref{sec:blocks-fluctuation-space}), and this bond lies on the segment that joins the two centres. As $g$ maps segments onto segments, the central bond of $\langle y,y+Le_{\mu}\rangle$ goes onto the face bond of the image face that lies on the segment joining $gy$ and $g(y+Le_{\mu})$, that is, onto the central bond of the image face. Being injective on the finite set of the bonds, $g$ thus permutes the central bonds. Hence $gC\mathsf{Z}$ coincides with $g\mathsf{Z}$ off the central bonds, and $QgC\mathsf{Z}=gQC\mathsf{Z}=0$. By the uniqueness in Proposition~\ref{prop:fluctuation-space}(iv), $gC\mathsf{Z}=Cg\mathsf{Z}$.

\emph{The operators $\tilde{\Delta}_k$, $(\tilde{\Delta}_k+x)^{-1}$, $C_k^{\frac{1}{2}}$, $\mathrm{Sq}_k$.} The map $g$ is orthogonal on $V_k$ (Proposition~\ref{prop:fluctuation-space}(ii)) and on the bond functions on $T^0_k$. Transposing $gC=Cg$ gives $gC^T=C^Tg$, and so
\[
g\tilde{\Delta}_k=gC^T\Delta_kC=C^T\Delta_kCg=\tilde{\Delta}_kg.
\]
By Propositions~\ref{prop:bounds-fluctuation-form} and~\ref{prop:decay-resolvents-root}, $\tilde{\Delta}_k+x$ is invertible for every $x\ge0$, and
\[
C_k^{\frac{1}{2}}=\pi^{-1}\int_0^{\infty}\mathrm{d}x\,x^{-\frac{1}{2}}\cdot(\tilde{\Delta}_k+x)^{-1}.
\]
Hence the resolvents $(\tilde{\Delta}_k+x)^{-1}$ and the root $C_k^{\frac{1}{2}}$ commute with $g$, and so does $\mathrm{Sq}_k=C\,C_k^{\frac{1}{2}}$.

\emph{The minimiser $A^{\min}_k$.} Let $g\in\mathbb{G}_k^L$, let $g'$ be as in the statement, and let $A'$ be a real bond function on $T^0_{k+1}$. For $A\in\mathrm{Adm}(A')$, Proposition~\ref{prop:properties-actions}(d) gives
\[
QgA=gQA=g(A'_L)=(g'A')_L,
\]
and $(\tau gA)(gy,gx)=(\tau A)(y,x)=0$ for all $y$ and $x\in B(y)$; so $gA\in\mathrm{Adm}(g'A')$. Since $g\mapsto g'$ is an isomorphism of $\mathbb{G}_k^L$ onto $\mathbb{G}_{k+1}$ (Proposition~\ref{prop:properties-actions}(d)), one has $(g^{-1})'=(g')^{-1}$, and the same argument shows that $g^{-1}$ maps $\mathrm{Adm}(g'A')$ into $\mathrm{Adm}(A')$; thus $g$ maps $\mathrm{Adm}(A')$ onto $\mathrm{Adm}(g'A')$. Moreover
\[
\langle gA,\Delta_kgA\rangle=\langle gA,g\Delta_kA\rangle=\langle A,\Delta_kA\rangle.
\]
By the uniqueness of the minimiser (Lemma~\ref{lem:existence-uniqueness}), $g\,A^{\min}_k(A')=A^{\min}_k(g'A')$.

\emph{Reality.} The operators $d$, $\delta$, $\Delta$, $Q_k$, $Q_k^{(2)}$, $Q$, $\tau$ and $C$ have real coefficients, and $\mathscr{R}_k$ is the orthogonal projection onto a space spanned by real functions. The problems that define $H_k$ and $A^{\min}_k$ have real data. The remaining operators are formed from these by transposition, products, inversion and an integral with real weights. Hence all of them are real.

\emph{The box and the measures.} The last assertion is Proposition~\ref{prop:fluctuation-space}(iii).

Proposition~\ref{prop:symmetries-reality} follows.
\end{proof}

\subsubsection{A torus and a covering torus}\label{sec:torus-covering-torus}

Throughout Section~\ref{sec:torus-covering-torus}, $l_1$ and $l_2$ are sides in $\mathrm{Cl}$ with $l_2=J\cdot l_1$ for an odd integer $J$, as in property~\ref{it:renormalisation-group-18}, and $T_i:=\mathbb{L}_k/l_i\mathbb{Z}^3$ for $i=1,2$, where $\mathbb{L}_k=\eta_k\mathbb{Z}^3$. Since $T_i$, $\mathbb{Z}^3/l_i\mathbb{Z}^3$ and $L\mathbb{Z}^3/l_i\mathbb{Z}^3$ are the quotients of $\mathbb{L}_k$, $\mathbb{Z}^3$ and $L\mathbb{Z}^3$ by $l_i\mathbb{Z}^3$, a function on the cells of one of these tori has as its lift the function on the cells of the corresponding infinite lattice that takes at each cell the value at the image of that cell in the torus, and the lifts are exactly the $l_i\mathbb{Z}^3$-periodic functions. A function and its lift are identified, and the elements of $V_k(T_i)$ are treated in the same way. As $l_2\mathbb{Z}^3$ is contained in $l_1\mathbb{Z}^3$, every $l_1\mathbb{Z}^3$-periodic function is $l_2\mathbb{Z}^3$-periodic; the inclusion of the $l_1\mathbb{Z}^3$-periodic functions into the $l_2\mathbb{Z}^3$-periodic functions is denoted $\mathrm{pj}^{*}$. For $v\in l_1\mathbb{Z}^3$ the translation $T_v$ is defined by $(T_vf)(x):=f(x-v)$, where $x-v$ is the translate of the cell or index $x$ by $-v$. Put
\[
\mathrm{Per}:=\left\{f\ \text{a function on}\ T_2:\ T_vf=f\ \text{for all}\ v\in l_1\mathbb{Z}^3\right\};
\]
this is the image of $\mathrm{pj}^{*}$. A period cell of $l_2\mathbb{Z}^3$ is the union of $J^3$ translates of a period cell of $l_1\mathbb{Z}^3$; hence, for functions $f$ and $h$ of the same kind on $T_1$,
\[
\left\langle\mathrm{pj}^{*}f,\mathrm{pj}^{*}h\right\rangle_{T_2}=J^3\cdot\langle f,h\rangle_{T_1}.
\]

\begin{lemma}[Three rules for periodic restrictions]\label{lem:three-rules-periodic}
\begin{enumerate}
\item[(a)] By an index set is meant one of the following sets of objects of $\mathbb{L}_k$, each of which is mapped onto itself by the translations by the vectors of $l_1\mathbb{Z}^3$: the set of the cells of one of the lattices $\mathbb{L}_k$, $\mathbb{Z}^3$, $L\mathbb{Z}^3$; the set of the pairs $(y,x)$ with $y\in L\mathbb{Z}^3$ and $x\in B(y)$; the set of the basis elements $\Upsilon$. Let $I$ and $I'$ be index sets, and let $O$ be a linear map from the functions on $I$ to the functions on $I'$ whose rows have finite support and which satisfies $OT_v=T_vO$ for all $v\in l_1\mathbb{Z}^3$. Then, for $i=1$ and for $i=2$, $O$ carries the $l_i\mathbb{Z}^3$-periodic functions on $I$ into the $l_i\mathbb{Z}^3$-periodic functions on $I'$, and the restrictions $O^{(i)}$ of $O$ to the $l_i\mathbb{Z}^3$-periodic functions on $I$ satisfy
\[
O^{(2)}\mathrm{pj}^{*}=\mathrm{pj}^{*}O^{(1)}.
\]
By their defining formulas, each of the operators $d$, $\delta$, $\Delta$, $Q$, $Q_k$, $Q_k^T$, $\tau$, $C$ on $T_i$ is the restriction $O^{(i)}$ of an operator $O$ of this kind, because the translations by $l_1\mathbb{Z}^3$ preserve the partition of $\mathbb{L}_k$ into the cubes $B_k(y)$, $y\in\mathbb{Z}^3$, and the partition of $\mathbb{Z}^3$ into the $L$-blocks. In the sense of Definition~\ref{def:terminology}(2), each of these eight operators is given by one formula on $\mathbb{L}_k$, and $L$ is a range of each of these formulas.
\item[(b)] Let $O^{(i)}$ and $P^{(i)}$, $i=1,2$, be linear maps between spaces of functions on $T_i$ with $O^{(2)}\mathrm{pj}^{*}=\mathrm{pj}^{*}O^{(1)}$ and $P^{(2)}\mathrm{pj}^{*}=\mathrm{pj}^{*}P^{(1)}$. Then $\left(O^{(2)}+P^{(2)}\right)\mathrm{pj}^{*}=\mathrm{pj}^{*}\left(O^{(1)}+P^{(1)}\right)$ and, whenever the products are defined, $O^{(2)}P^{(2)}\mathrm{pj}^{*}=\mathrm{pj}^{*}O^{(1)}P^{(1)}$; the relation is preserved under limits and under integrals of families of pairs $\left(O^{(1)},O^{(2)}\right)$ that satisfy it. Suppose, finally, that $O^{(1)}$ and $O^{(2)}$ are square matrices. If $O^{(2)}$ has an inverse, then so has $O^{(1)}$, and
\[
\left(O^{(2)}\right)^{-1}\mathrm{pj}^{*}=\mathrm{pj}^{*}\left(O^{(1)}\right)^{-1}.
\]
\item[(c)] Let $O^{(2)}$ be real, let it commute with the $T_v$, $v\in l_1\mathbb{Z}^3$, and let $O^{(2)}\mathrm{pj}^{*}=\mathrm{pj}^{*}O^{(1)}$. Then, the transposes being taken with respect to the pairings on $T_2$ and on $T_1$,
\[
\left(O^{(2)}\right)^T\mathrm{pj}^{*}=\mathrm{pj}^{*}\left(O^{(1)}\right)^T.
\]
\end{enumerate}
\end{lemma}

\begin{proof}
(a) Let $f$ be a function on $I$. Every row of $O$ has finite support; so each value of $Of$ is a finite linear combination of values of $f$, and $Of$ is defined whether $f$ is periodic or not. Let $i\in\{1,2\}$. Since $l_2\mathbb{Z}^3=J\cdot l_1\mathbb{Z}^3$ is contained in $l_1\mathbb{Z}^3$, the map $O$ commutes with $T_v$ for every $v\in l_i\mathbb{Z}^3$. If $f$ is $l_i\mathbb{Z}^3$-periodic, it follows that $T_v(Of)=O(T_vf)=Of$ for these $v$, that is, $Of$ is $l_i\mathbb{Z}^3$-periodic. Now let $f$ be an $l_1\mathbb{Z}^3$-periodic function on $I$, that is, a function on the corresponding objects of $T_1$. The inclusion $\mathrm{pj}^{*}$ leaves the underlying periodic function unchanged; hence $O^{(2)}\mathrm{pj}^{*}f$ and $\mathrm{pj}^{*}O^{(1)}f$ are both the function $Of$.

It remains to prove the last assertion of (a). Let $O$ be one of the eight operators named there, formed on $\mathbb{L}_k$ by its defining formula. No bond function and no union of $M$-cubes enters this formula. So one and the same operator $O$ corresponds to all bond functions and all unions of $M$-cubes of $\mathbb{L}_k$, its rows do not depend on these data, and condition~(ii) of Definition~\ref{def:terminology}(2) states that $O$ commutes with $T_v$ for every $v\in(L\cdot M)\mathbb{Z}^3$. This holds for the reason given in (a) for the vectors of $l_1\mathbb{Z}^3$: the vectors of $(L\cdot M)\mathbb{Z}^3$ lie in $L\mathbb{Z}^3$ and in $\mathbb{Z}^3$, so that the translations by them preserve the partition of $\mathbb{L}_k$ into the cubes $B_k(y)$, the partition of $\mathbb{Z}^3$ into the $L$-blocks together with their centres, and the family of the basis elements $\Upsilon$. Conditions~(i) and~(iii) hold, since the rows of $O$ have finite support and the operator on $T_i$ is $O^{(i)}$.

Let $x$ be a row index and $x'$ a column index of $O$ with $O(x,x')\ne0$. If $O$ is $d$ or $\delta$, the supports of $x$ and $x'$ meet: the value of $df$ at a cell $c$ is a combination of the values of $f$ at the cells in the boundary of $c$, whose supports lie in that of $c$, and $\delta=d^T$ has its non-zero entries at the same pairs of indices as $d$, with rows and columns exchanged. As $\Delta$ is $\delta d$ on scalars and $\delta d+d\delta$ on bond functions, a non-zero entry $\Delta(x,x')$ requires a site, a bond or a plaquette whose support meets the supports of $x$ and of $x'$; such a support has sup-diameter at most the spacing of the lattice, and so the supports of $x$ and $x'$ have sup distance at most this spacing, which is at most $1$. Let $y$ be a unit site. The columns of the row of $Q_k$ at $y$ are the sites of $B_k(y)$, which have sup distance less than $\frac{1}{2}$ from $y$. The columns of its row at a unit bond $(y,y+e_{\mu})$ are the fine bonds of the paths $\Gamma_{z,z+e_{\mu}}$, $z\in B_k(y)$, and every site $w$ of these paths has sup distance less than $1$ from $y$ or from $y+e_{\mu}$, because $|w_{\nu}-y_{\nu}|<\frac{1}{2}$ for $\nu\ne\mu$ and $-\frac{1}{2}<w_{\mu}-y_{\mu}<\frac{3}{2}$. The non-zero entries of $Q_k^T$ stand at the same pairs of indices as those of $Q_k$, with rows and columns exchanged. Hence the supports of $x$ and $x'$ have sup distance less than $1$ for $O=Q_k$ and for $O=Q_k^T$. Now let $y\in L\mathbb{Z}^3$. The columns of the row of $\tau$ at a pair $(y,z)$, $z\in B(y)$, are bonds of the six paths $\Gamma^{\pi}(y,z)$, which change each coordinate monotonically from $y_{\nu}$ to $z_{\nu}$; so every site $w$ of these paths satisfies $|w-y|\le|z-y|\le\frac{L-1}{2}$. The columns of the row of $Q$ at a bond $(y,y+Le_{\mu})$ are the unit bonds of the paths $\Gamma_{z,z+Le_{\mu}}$, $z\in B(y)$, and every site $w$ of these paths has sup distance at most $\frac{L-1}{2}$ from $y$ or from $y+Le_{\mu}$, because $|w_{\nu}-y_{\nu}|=|z_{\nu}-y_{\nu}|\le\frac{L-1}{2}$ for $\nu\ne\mu$, and $w_{\mu}-y_{\mu}$ is an integer with $-\frac{L-1}{2}\le w_{\mu}-y_{\mu}\le\frac{3L-1}{2}$, while $L$ is odd. Hence the supports of $x$ and $x'$ have sup distance at most $\frac{L-1}{2}$ for $O=\tau$ and for $O=Q$. Finally, the entry of $C$ at a unit bond $b$ and a basis element $\Upsilon$ is $(C\Upsilon)(b)$; if it does not vanish, the ends of $b$ belong to $\operatorname{supp}C\Upsilon$, and Proposition~\ref{prop:fluctuation-space}(iv) bounds the sup distance of the supports by $L$. All these bounds are at most $L$; so $L$ is a range of each of the eight formulas.

(b) The assertions on sums, limits and integrals follow from the linearity and the continuity of $\mathrm{pj}^{*}$, and the assertion on products from
\[
O^{(2)}P^{(2)}\mathrm{pj}^{*}=O^{(2)}\mathrm{pj}^{*}P^{(1)}=\mathrm{pj}^{*}O^{(1)}P^{(1)}.
\]
Let $O^{(1)}$ and $O^{(2)}$ be square matrices, $O^{(2)}$ invertible. If $O^{(1)}f=0$, then $O^{(2)}\mathrm{pj}^{*}f=\mathrm{pj}^{*}O^{(1)}f=0$, and $f=0$ because $O^{(2)}$ and $\mathrm{pj}^{*}$ are injective. So $O^{(1)}$ is injective, hence invertible, and
\[
\mathrm{pj}^{*}\left(O^{(1)}\right)^{-1}=\left(O^{(2)}\right)^{-1}O^{(2)}\mathrm{pj}^{*}\left(O^{(1)}\right)^{-1}=\left(O^{(2)}\right)^{-1}\mathrm{pj}^{*}O^{(1)}\left(O^{(1)}\right)^{-1}=\left(O^{(2)}\right)^{-1}\mathrm{pj}^{*}.
\]

(c) The translations $T_v$ are orthogonal for the pairing on $T_2$, and $T_v^T=T_{-v}$. Transposing $O^{(2)}T_v=T_vO^{(2)}$ for all $v\in l_1\mathbb{Z}^3$ shows that $\left(O^{(2)}\right)^T$ commutes with the $T_v$; hence it maps $\mathrm{Per}$ into $\mathrm{Per}$, and for every function $h$ on $T_1$ there is a function $h_1$ on $T_1$ with $\left(O^{(2)}\right)^T\mathrm{pj}^{*}h=\mathrm{pj}^{*}h_1$. For every function $f$ on $T_1$,
\begin{align*}
J^3\cdot\langle f,h_1\rangle_{T_1}&=\left\langle\mathrm{pj}^{*}f,\left(O^{(2)}\right)^T\mathrm{pj}^{*}h\right\rangle_{T_2}=\left\langle O^{(2)}\mathrm{pj}^{*}f,\mathrm{pj}^{*}h\right\rangle_{T_2}\\
&=\left\langle\mathrm{pj}^{*}O^{(1)}f,\mathrm{pj}^{*}h\right\rangle_{T_2}=J^3\cdot\left\langle f,\left(O^{(1)}\right)^Th\right\rangle_{T_1}.
\end{align*}
Hence $h_1=\left(O^{(1)}\right)^Th$, which is the assertion.
\end{proof}

\begin{proposition}[Periodisation of the photon operators]\label{prop:periodisation-photon-operators}
For each of the operators
\[
O=d,\ \delta,\ \Delta,\ Q,\ Q_k,\ Q_k^T,\ \tau,\ C,\ H_k,\ \Delta_k,\ \tilde{\Delta}_k,\ (\tilde{\Delta}_k+x)^{-1}\ (x\ge0),\ C_k^{\frac{1}{2}},\ \mathrm{Sq}_k
\]
let $O^{(i)}$ denote the operator $O$ formed on the torus $T_i$. Then
\[
O^{(2)}\mathrm{pj}^{*}=\mathrm{pj}^{*}O^{(1)}.
\]
Moreover, for every real bond function $A'$ on the unit lattice of level $k+1$ and side $l_1/L$,
\[
{A^{\min}}_k^{(2)}\left(\mathrm{pj}^{*}A'\right)=\mathrm{pj}^{*}{A^{\min}}_k^{(1)}(A'),
\]
where $\mathrm{pj}^{*}$ on the left is the inclusion at level $k+1$, for the sides $l_1/L$ and $l_2/L$.
\end{proposition}

\begin{proof}
For $O=d,\delta,\Delta,Q,Q_k,Q_k^T,\tau$ the relation is Lemma~\ref{lem:three-rules-periodic}(a). For $O=C$ it is Lemma~\ref{lem:three-rules-periodic}(a) as well, since $C$ is the identity off the central bonds and is given by the formula~\eqref{eq:operators-photon-16} on the central bonds.

\emph{The minimiser $H_k$.} By linearity it suffices to consider a real bond function $A$ on the unit lattice of side $l_1$. Let $\mathscr{B}:=H_k^{(1)}A$. By Lemma~\ref{lem:characterisation-minimiser} on $T_1$, $Q_k\mathscr{B}=A$, and there are a bond function $\boldsymbol{\Lambda}$ and a scalar $\boldsymbol{\omega}$ on the unit lattice of side $l_1$ with
\[
\delta d\mathscr{B}=Q_k^T\boldsymbol{\Lambda},\qquad\Delta\delta\mathscr{B}=Q_k^T\boldsymbol{\omega}.
\]
Applying $\mathrm{pj}^{*}$, and using the relation for $Q_k$, $d$, $\delta$, $\Delta$ and $Q_k^T$, one obtains
\[
Q_k\mathrm{pj}^{*}\mathscr{B}=\mathrm{pj}^{*}A,\qquad\delta d\,\mathrm{pj}^{*}\mathscr{B}=Q_k^T\mathrm{pj}^{*}\boldsymbol{\Lambda},\qquad\Delta\delta\,\mathrm{pj}^{*}\mathscr{B}=Q_k^T\mathrm{pj}^{*}\boldsymbol{\omega}.
\]
So the real bond function $\mathrm{pj}^{*}\mathscr{B}$ satisfies \eqref{eq:operators-photon-M1}, \eqref{eq:operators-photon-M2} and \eqref{eq:operators-photon-M3} on $T_2$ for the datum $\mathrm{pj}^{*}A$, and Lemma~\ref{lem:characterisation-minimiser} on $T_2$ gives $\mathrm{pj}^{*}\mathscr{B}=H_k^{(2)}\mathrm{pj}^{*}A$.

\emph{Commutation with the translations.} For $v\in l_1\mathbb{Z}^3$ the element $(1,v)$ of $\mathbb{G}_k(T_2)$, the group $\mathbb{G}_k$ formed with the side $l_2$, acts on functions as $T_v$, because $l_1\mathbb{Z}^3\subset M\mathbb{Z}^3$. As $l_2\in\mathrm{Cl}$, Proposition~\ref{prop:lattice-symmetries} on $T_2$ shows that $d^{(2)}$, $H_k^{(2)}$ and $C^{(2)}$ commute with the $T_v$.

\emph{The operators $\Delta_k$ and $\tilde{\Delta}_k$.} By Lemma~\ref{lem:three-rules-periodic}(b) the relation holds for the product $dH_k$. The operator $(dH_k)^{(2)}=d^{(2)}H_k^{(2)}$ is real (Proposition~\ref{prop:lattice-symmetries}) and commutes with the $T_v$; so Lemma~\ref{lem:three-rules-periodic}(c) gives the relation for $(dH_k)^T$, and Lemma~\ref{lem:three-rules-periodic}(b) gives it for $\Delta_k=(dH_k)^TdH_k$. In the same way $C^{(2)}$ is real and commutes with the $T_v$, so that Lemma~\ref{lem:three-rules-periodic}(c) gives the relation for $C^T$, and Lemma~\ref{lem:three-rules-periodic}(b) gives it for $\tilde{\Delta}_k=C^T\Delta_kC$.

\emph{The resolvents, the root and $\mathrm{Sq}_k$.} By Proposition~\ref{prop:bounds-fluctuation-form}, $\tilde{\Delta}_k$ is positive definite on $V_k(T_1)$ and on $V_k(T_2)$; so $\tilde{\Delta}_k+x$ is invertible for every $x\ge0$. The relation holds for $\tilde{\Delta}_k+x$, since it holds for the identity, and Lemma~\ref{lem:three-rules-periodic}(b) gives it for $(\tilde{\Delta}_k+x)^{-1}$, then for the integral
\[
C_k^{\frac{1}{2}}=\pi^{-1}\int_0^{\infty}\mathrm{d}x\,x^{-\frac{1}{2}}\cdot(\tilde{\Delta}_k+x)^{-1}
\]
of Proposition~\ref{prop:decay-resolvents-root}, and for the product $\mathrm{Sq}_k=C\,C_k^{\frac{1}{2}}$.

\emph{The minimiser $A^{\min}_k$.} Let $A:={A^{\min}}_k^{(1)}(A')$. The inclusions are the identity on the underlying periodic functions; hence $\left(\mathrm{pj}^{*}A'\right)_L=\mathrm{pj}^{*}(A'_L)$, and the relation for $Q$ and for $\tau$ gives
\[
Q\mathrm{pj}^{*}A=\mathrm{pj}^{*}QA=\left(\mathrm{pj}^{*}A'\right)_L,\qquad\tau\mathrm{pj}^{*}A=\mathrm{pj}^{*}\tau A=0,
\]
that is, $\mathrm{pj}^{*}A\in\mathrm{Adm}\left(\mathrm{pj}^{*}A'\right)$ on $T_2$. Let $\tilde Z$ be a real bond function on the unit lattice of side $l_2$ with $Q\tilde Z=0$ and $\tau\tilde Z=0$, and put
\[
\mathbf{Z}:=\sum_{v\in l_1\mathbb{Z}^3/l_2\mathbb{Z}^3}T_v\tilde Z,
\]
the sum being over the $J^3$ classes of $l_1\mathbb{Z}^3$ modulo $l_2\mathbb{Z}^3$; on functions on $T_2$ the translation $T_v$ depends only on the class of $v$. Then $\mathbf{Z}\in\mathrm{Per}$, since $T_w\mathbf{Z}$ is the same sum with the classes permuted, for every $w\in l_1\mathbb{Z}^3$; so $\mathbf{Z}=\mathrm{pj}^{*}Z$ with a real bond function $Z$ on the unit lattice of side $l_1$. Since $Q$ and $\tau$ commute with the $T_v$ (their formulas are invariant under the translations by $l_1\mathbb{Z}^3$), one has $\mathrm{pj}^{*}QZ=Q\mathbf{Z}=0$ and $\mathrm{pj}^{*}\tau Z=\tau\mathbf{Z}=0$, whence $QZ=0$ and $\tau Z=0$, that is, $Z\in C(V_k(T_1))$. The function $F:=\Delta_k^{(2)}\mathrm{pj}^{*}A=\mathrm{pj}^{*}\Delta_k^{(1)}A$ lies in $\mathrm{Per}$, and the $T_v$ are orthogonal; hence $\langle F,T_v\tilde Z\rangle_{T_2}=\langle T_{-v}F,\tilde Z\rangle_{T_2}=\langle F,\tilde Z\rangle_{T_2}$ for every $v$, and $\langle F,\mathbf{Z}\rangle_{T_2}=J^3\cdot\langle F,\tilde Z\rangle_{T_2}$. As $\Delta_k^{(2)}$ is symmetric,
\begin{align*}
\left\langle\mathrm{pj}^{*}A,\Delta_k^{(2)}\tilde Z\right\rangle_{T_2}&=J^{-3}\cdot\left\langle\Delta_k^{(2)}\mathrm{pj}^{*}A,\mathbf{Z}\right\rangle_{T_2}=J^{-3}\cdot\left\langle\mathrm{pj}^{*}\Delta_k^{(1)}A,\mathrm{pj}^{*}Z\right\rangle_{T_2}\\
&=\left\langle\Delta_k^{(1)}A,Z\right\rangle_{T_1}=\left\langle A,\Delta_k^{(1)}Z\right\rangle_{T_1}=0,
\end{align*}
where the last equality is the characterisation of Lemma~\ref{lem:existence-uniqueness} on $T_1$. By the same characterisation on $T_2$, $\mathrm{pj}^{*}A={A^{\min}}_k^{(2)}\left(\mathrm{pj}^{*}A'\right)$.

Proposition~\ref{prop:periodisation} follows.
\end{proof}

\subsection{Local bounds for the background field}\label{sec:local-bounds-background}

Throughout this section $l\in\mathrm{Cl}$. Put
\[
s_d:=(2L+1)M .
\]
This is the side of the cube $\square^{\dagger}$ of Definition~\ref{def:enlarged-cubes}, the enlargement of an $M$-cube $\square$ by $L$ layers of $M$-cubes. As $l/(LM)$ is an odd integer not smaller than $13$, one has $l\ge 13LM$, and therefore
\[
s_d\le 3LM\le\frac{3l}{13};
\]
in particular $s_d<l$. Moreover $l$ is odd, since $L$ and $M$ are odd. With the constants $C_P$, $\gamma_P$ of \eqref{eq:operators-photon-18} and the function $c_{\mathrm{bd}}$ of \eqref{eq:operators-photon-1} define
\begin{equation}\label{eq:operators-photon-19}
C_B:=C_P\left[(s_d+4)\cdot c_{\mathrm{bd}}(\gamma_P)+\left(\frac{2}{\gamma_P}+\frac{62}{\gamma_P^2}\right)\cdot c_{\mathrm{bd}}\left(\frac{\gamma_P}{2}\right)\right].
\end{equation}

\begin{proposition}[Regularity on enlarged cubes]\label{prop:regularity-enlarged-cubes}
Let $\square$ be an $M$-cube of $T_k$, and let $A$ be a real bond function on $T^0_k$ whose harmonic part is arbitrary. Then there exists a function $\nu$ from the set of the sites of $T_k$ lying in $\square^{\dagger}$ to $\mathbb{R}$ with the following property: the bond function $H_kA-d\nu$, given on the bonds of $T_k$ contained in $\square^{\dagger}$, has each of its three norms on $\square^{\dagger}$ bounded by $C_B\cdot|dA|_{\infty}$.
\end{proposition}

\begin{proof}
Put $p_A:=|dA|_{\infty}$. Let $c$ be the mean of $A$, regarded also as a constant bond function, and put $A^0:=A-c$; then $A^0$ is real, it has mean zero, and $dA^0=dA$. Let $x_0\in M\mathbb{Z}^3$ be the centre of $\square$. Since $l$ is odd, the box of Lemma~\ref{lem:line-sums-base} with centre $x_0$ is
\[
W=x_0+\left\{-\frac{l-1}{2},\dots,\frac{l-1}{2}\right\}^3.
\]
As in Lemma~\ref{lem:line-sums-base}, applied to $A^0$ and $x_0$, let $\lambda(x):=A^0(\Gamma^{(0,1,2)}(x_0,x))$, the path being taken inside $W$, and $A'':=A^0-d\lambda$. The inner bonds and the seam bonds are those of Lemma~\ref{lem:line-sums-base}, and $p_A$ is the number $|dA^0|_{\infty}$ of that lemma.

\emph{Step 1.} By Proposition~\ref{prop:background-field} the operator $H_k$ is linear and $H_kc=c$; since $A=c+d\lambda+A''$,
\[
H_kA=c+H_kd\lambda+H_kA''.
\]
The field $d\lambda$ is closed, hence so is $H_kd\lambda$ (Proposition~\ref{prop:background-field}). By Lemma~\ref{lem:block-averages}(c) (see also Proposition~\ref{prop:averaging-operators}(c), taken at the level $k$) and by $Q_kH_k=1$, the mean of $H_kd\lambda$ equals the mean of $Q_kH_kd\lambda=d\lambda$, which is zero because $d\lambda$ is exact (Section~\ref{sec:tori-cochains-kernels}). The field $H_kd\lambda$ is real. By the Hodge decomposition on the discrete torus (Section~\ref{sec:tori-cochains-kernels}) there is therefore a real function $\nu'$ on $T_k$ with $H_kd\lambda=d\nu'$. Since $s_d<l$, the cube $\square^{\dagger}$ does not wrap (Definition~\ref{def:terminology}(9)); let $x$ denote the coordinate on its lift with origin $x_0$. Then $d(c\cdot x)=c$ on the bonds of $\square^{\dagger}$, and the real function
\[
\nu:=c\cdot x+\nu'
\]
on the sites of $\square^{\dagger}$ satisfies $H_kA-d\nu=H_kA''$ on the bonds of $\square^{\dagger}$.

\emph{Step 2.} Consider a bond $b$ of $T_k$ that is contained in $\square^{\dagger}$, and choose one of its end points, $x_b$. As $\square^{\dagger}$ is the cube of side $s_d$ centred at $x_0$ (Definition~\ref{def:enlarged-cubes}) and $d$ is the sup distance, one has $d(x_b,x_0)\le\frac{s_d}{2}$. Consider further a positively oriented unit bond $b'$, with initial point $x'$ and direction $\mathfrak{l}$, that is, $b'=(x',x'+e_{\mathfrak{l}})$. The two ends of $b$ have the distance $\eta\le 1$, and the two ends of $b'$ the distance $1$; hence $d(x',x_b)\le d(b,b')+2$. If $b'$ is an inner bond, Lemma~\ref{lem:line-sums-base}(i) gives
\[
|A''(b')|\le 2\,d(x',x_0)\cdot p_A\le 2\left(d(x',x_b)+\frac{s_d}{2}\right)p_A\le\left(2\,d(b,b')+4+s_d\right)p_A .
\]
If $b'$ is a seam bond, then $x'=x_0+t$ with $t_{\mathfrak{l}}=\frac{l-1}{2}$, and the coordinate in the direction $\mathfrak{l}$ of $x'+e_{\mathfrak{l}}-x_0$ is $\frac{l+1}{2}\equiv-\frac{l-1}{2}$ modulo $l$. As $\frac{l-1}{2}$ is the largest distance in $\mathbb{Z}/l\mathbb{Z}$, both ends of $b'$ have the distance $\frac{l-1}{2}$ from $x_0$, and consequently
\[
d(b,b')\ge\frac{l-1-s_d}{2}.
\]
For a seam bond $b'$, Lemma~\ref{lem:line-sums-base}(ii) gives $|A''(b')|\le\frac{3}{2}l^2\cdot p_A$.

\emph{Step 3.} By Proposition~\ref{prop:kernel-bounds}, which holds with the constants $C_P$, $\gamma_P$ of \eqref{eq:operators-photon-18} (Theorem~\ref{thm:decay-kernels} and Section~\ref{sec:fluctuation-form-inverse}), the kernels of $H_k$, $\partial H_k$ and $\delta_{\alpha}\partial H_k$ are bounded by $C_P\exp(-\gamma_P\,d(b,b'))$. By the paragraph on kernels in Section~\ref{sec:tori-cochains-kernels}, each of the three norms of $H_kA''$ on $\square^{\dagger}$, formed with the bonds and the differences inside $\square^{\dagger}$, is therefore at most
\[
\sup_b\ \sum_{b'}C_P\exp(-\gamma_P\,d(b,b'))\cdot|A''(b')|,
\]
where $b$ ranges over the bonds of $T_k$ in $\square^{\dagger}$ and $b'$ over the unit bonds. Write $g:=\gamma_P$; then $0<g\le 1$ (Section~\ref{sec:fluctuation-form-inverse}). For $t\ge 0$ and $a>0$ one has $t\exp(-at)\le\frac{1}{\exp(1)\cdot a}$; with $a=\frac{g}{2}$ and $\frac{4}{\exp(1)}\le 2$ this gives
\[
2t\exp(-gt)\le\frac{2}{g}\exp\left(-\frac{gt}{2}\right)\qquad(t\ge 0).
\]
By Step~2 and \eqref{eq:operators-photon-2}, the inner bonds $b'$ contribute at most $C_P\cdot p_A$ times
\[
\sum_{b'}\exp(-g\,d(b,b'))\left(2\,d(b,b')+4+s_d\right)\le(s_d+4)\cdot c_{\mathrm{bd}}(g)+\frac{2}{g}\,c_{\mathrm{bd}}\left(\frac{g}{2}\right).
\]
For a seam bond $b'$ the lower bound on $d(b,b')$ in Step~2 gives
\[
\exp(-g\,d(b,b'))\le\exp\left(-\frac{g(l-1-s_d)}{4}\right)\exp\left(-\frac{g\,d(b,b')}{2}\right)\le\exp\left(\frac{g}{4}\right)\exp\left(-\frac{5gl}{26}\right)\exp\left(-\frac{g\,d(b,b')}{2}\right),
\]
because $\frac{l-s_d}{4}\ge\frac{5l}{26}$ by $s_d\le\frac{3l}{13}$. Moreover, by the same elementary inequality with $a=\frac{5g}{52}$,
\[
l^2\exp\left(-\frac{5gl}{26}\right)=\left(l\exp\left(-\frac{5gl}{52}\right)\right)^2\le\left(\frac{52}{5\exp(1)\cdot g}\right)^2\le\frac{15}{g^2}.
\]
Together with $|A''(b')|\le\frac{3}{2}l^2\cdot p_A$ and \eqref{eq:operators-photon-2}, the seam bonds $b'$ contribute at most $C_P\cdot p_A$ times
\[
\frac{3}{2}\cdot\frac{15}{g^2}\exp\left(\frac{g}{4}\right)c_{\mathrm{bd}}\left(\frac{g}{2}\right)\le\frac{62}{g^2}\,c_{\mathrm{bd}}\left(\frac{g}{2}\right),
\]
since $g\le 1$ and $\frac{45}{2}\exp\left(\frac{1}{4}\right)<62$. By \eqref{eq:operators-photon-19} the sum of the two contributions is $C_B\cdot p_A$. Hence each of the three norms of $H_kA''$ on $\square^{\dagger}$ is at most $C_B\cdot|dA|_{\infty}$, and by Step~1 the same holds for $H_kA-d\nu$, which coincides with $H_kA''$ on the bonds of $\square^{\dagger}$. Proposition~\ref{prop:local-regularity} follows.
\end{proof}

\begin{proposition}[Backgrounds in the domains]\label{prop:backgrounds-domains}
Let $\varepsilon\in(0,\bar{\varepsilon}]$, let $L^{\frac{3}{4}}\ge\frac{8}{3}$, and let $0<e_k<1$; the numbers $v_k=-\log e_k$, $p_k=v_k^p$ and $p_{0,k}=v_k^{p_0}$ are formed with $e_k$ as in Section~\ref{sec:running-quantities-tori}, the exponents $p$, $p_0$ being those of stage~\ref{it:order-choice-8} of Definition~\ref{def:order-choice}. Assume
\[
8\,C_B\cdot p_k\le e_k^{-\frac{3}{4}+3\varepsilon},\qquad 8\,e_k^{\frac{3}{4}-\frac{7\varepsilon}{2}}\le 1 .
\]
Let $C_w\ge C_P$ and $\gamma_w\le\gamma_P$ be numbers, let $c_{\gamma}$ be formed with $\gamma_w$ as in stage~\ref{it:order-choice-7}, and assume that conditions~\eqref{eq:renormalisation-group-A1} and~\eqref{eq:renormalisation-group-A2} of Section~\ref{sec:conditions} hold at $e_k$ with these numbers. In particular
\[
p_{0,k}\le e_k^{-\frac{\varepsilon}{4}}\qquad\text{and}\qquad 2\,C_w^2\cdot c_{\gamma}^2\cdot c_C\cdot e_k^{\frac{\varepsilon}{4}}\le 1,
\]
the first inequality being the first one of condition~\eqref{eq:renormalisation-group-A1} and the second being condition~\eqref{eq:renormalisation-group-A2}, where $c_C=3L^{\frac{9}{2}}$ is the constant of Proposition~\ref{prop:fluctuation-space}(iv). The domains are those of Definition~\ref{def:domains}. Then the following hold.
\begin{enumerate}
\item[(i)] The set $\{A\ \text{real on the bonds of}\ T^0_k : |dA|_{\infty}\le p_k\}$ lies in $\mathrm{Dom}_k$, and it contains $0$ and the constant bond functions.
\item[(ii)] Let moreover $l/L\in\mathrm{Cl}$, and let $\mathrm{Dom}_{k+1}$ be formed with the side $l/L$ and with $e_{k+1}=L^{\frac{1}{2}}e_k$. Then $A^{\min}_k(A')+\mathrm{Sq}_k\zeta\in\mathrm{Dom}_k$ for all $A'\in\mathrm{Dom}_{k+1}$ and all $\zeta\in\mathrm{Box}_k$.
\item[(iii)] For every $\mathsf{Z}\in\mathrm{Box}_k$ each of the three norms of $\mathscr{Z}:=H_kC\,C_k^{\frac{1}{2}}\mathsf{Z}$ is at most $\frac{1}{2}e_k^{-\frac{\varepsilon}{2}}$.
\end{enumerate}
\end{proposition}

\begin{proof}
Number the three norms of Section~\ref{sec:domains-bond-functions} by $i=1,2,3$ in the order $|\cdot|_{\infty}$, $|\partial\,\cdot|_{\infty}$, $|\delta_{\alpha}\partial\,\cdot|_{\infty}$; then the $i$-th threshold of $R_k$ is $e_k^{-\frac{3}{4}+i\varepsilon}$, and that of $R_{k+1}$ is $e_{k+1}^{-\frac{3}{4}+i\varepsilon}$. Since $e_k<1$,
\[
e_k^{-\frac{3}{4}+\varepsilon}\ge e_k^{-\frac{3}{4}+2\varepsilon}\ge e_k^{-\frac{3}{4}+3\varepsilon}=:\mathrm{thr} .
\]

\emph{Clause (i).} Let $A$ be real with $|dA|_{\infty}\le p_k$, and let $\square$ be an $M$-cube. By Proposition~\ref{prop:local-regularity} there is a real function $\nu$ on the sites of $\square^{\dagger}$ such that the real bond function $A^b:=H_kA-d\nu$ on the bonds of $\square^{\dagger}$ has its three norms on $\square^{\dagger}$ at most
\[
C_B\cdot|dA|_{\infty}\le C_B\cdot p_k\le\frac{\mathrm{thr}}{8}<\frac{\mathrm{thr}}{4},
\]
hence below $\frac{1}{4}$ of each threshold of $R_k$; and $H_kA=A^b+d\nu$ on $\square^{\dagger}$. The field $H_kA$ is real (Proposition~\ref{prop:background-field}). Therefore $H_kA=A_0+A_1$, with $A_0:=H_kA$ and $A_1:=0\in\frac{1}{4}R_k$, lies in $\frac{1}{4}\tilde{R}_k$; that is, $A\in\mathrm{Dom}_k$. The bond function $0$ and every constant bond function $c$ belong to the set, because $dc=0$.

\emph{Clause (iii).} Let $\mathsf{Z}\in\mathrm{Box}_k$. Then $|\mathsf{Z}(\Upsilon)|\le p_{0,k}\le e_k^{-\frac{\varepsilon}{4}}$ for every $\Upsilon$, by condition~\eqref{eq:renormalisation-group-A1}. By Proposition~\ref{prop:decay-exact-root}, which holds with the constants of \eqref{eq:operators-photon-18} (Proposition~\ref{prop:decay-resolvents-root} and Section~\ref{sec:fluctuation-form-inverse}), and by $C_P\le C_w$, $\gamma_w\le\gamma_P$,
\[
\left|C_k^{\frac{1}{2}}(\Upsilon,\Upsilon')\right|\le C_w\exp(-\gamma_w\,d(\Upsilon,\Upsilon')).
\]
The number $c_{\gamma}$ bounds the sums of \eqref{eq:operators-photon-2} at $g=\gamma_w$: one has $\sum_{\Upsilon'}\exp(-\gamma_w\,d(\Upsilon,\Upsilon'))\le c_{\gamma}$ and $\sum_{b'}\exp(-\gamma_w\,d(b,b'))\le c_{\gamma}$. Hence every coordinate of $C_k^{\frac{1}{2}}\mathsf{Z}$ satisfies
\[
\left|\left(C_k^{\frac{1}{2}}\mathsf{Z}\right)(\Upsilon)\right|\le C_w\sum_{\Upsilon'}\exp(-\gamma_w\,d(\Upsilon,\Upsilon'))\cdot|\mathsf{Z}(\Upsilon')|\le C_w\cdot c_{\gamma}\cdot e_k^{-\frac{\varepsilon}{4}},
\]
and Proposition~\ref{prop:basis-map-box} (Proposition~\ref{prop:fluctuation-space}(iv)) gives
\[
\left|C\,C_k^{\frac{1}{2}}\mathsf{Z}\right|_{\infty}\le c_C\cdot C_w\cdot c_{\gamma}\cdot e_k^{-\frac{\varepsilon}{4}}.
\]
By Proposition~\ref{prop:kernel-bounds}, again with $C_P\le C_w$ and $\gamma_w\le\gamma_P$, the kernels of $H_k$, $\partial H_k$ and $\delta_{\alpha}\partial H_k$ are bounded by $C_w\exp(-\gamma_w\,d(b,b'))$. Summing over $b'$ as in the paragraph on kernels in Section~\ref{sec:tori-cochains-kernels}, one finds that each of the three norms of $\mathscr{Z}=H_k\left(C\,C_k^{\frac{1}{2}}\mathsf{Z}\right)$ is at most
\[
C_w\cdot c_{\gamma}\cdot\left|C\,C_k^{\frac{1}{2}}\mathsf{Z}\right|_{\infty}\le C_w^2\cdot c_{\gamma}^2\cdot c_C\cdot e_k^{-\frac{\varepsilon}{4}}\le\frac{1}{2}e_k^{-\frac{\varepsilon}{2}},
\]
where the last inequality is condition~\eqref{eq:renormalisation-group-A2}.

\emph{Clause (ii).} Let $A'\in\mathrm{Dom}_{k+1}$ and $\zeta\in\mathrm{Box}_k$, and put $\mathscr{A}':=H_{k+1}A'$ and $\mathscr{Z}:=H_kC\,C_k^{\frac{1}{2}}\zeta$. As $\mathrm{Sq}_k=C\,C_k^{\frac{1}{2}}$ and $H_k$ is linear (Proposition~\ref{prop:background-field}), Theorem~\ref{thm:minimisation-one-step}(b) gives
\[
H_k\left(A^{\min}_k(A')+\mathrm{Sq}_k\zeta\right)=(\mathscr{A}')_L-d\,\boldsymbol{\omega}_k(A')+\mathscr{Z} ,
\]
with the real scalar $\boldsymbol{\omega}_k(A')$ on $T_k$ of that theorem and the scaling $(\mathscr{A}')_L(b)=L^{-\frac{1}{2}}\mathscr{A}'(b/L)$ of Section~\ref{sec:local-parts-adjustment}.

The field $\mathscr{A}'$ is real (Proposition~\ref{prop:background-field}), and by the definition of $\mathrm{Dom}_{k+1}$ it lies in $\frac{1}{4}\tilde{R}_{k+1}$. So $\mathscr{A}'=A_0'+A_1'$ with $A_1'\in\frac{1}{4}R_{k+1}$ and $A_0'$ real, so that $A_1'=\mathscr{A}'-A_0'$ is real; and for every $M$-cube $\square'$ of $T_{k+1}$ there are a real function $\lambda'$ on the sites and a real bond function $A^b$ on the bonds of $\square'^{\dagger}$ with $A_0'=A^b+d\lambda'$ on $\square'^{\dagger}$, the three norms of $A^b$ on $\square'^{\dagger}$ being below $\frac{1}{4}$ of the thresholds of $R_{k+1}$. Consequently, on $\square'^{\dagger}$,
\[
\mathscr{A}'=A^c+d\lambda',\qquad A^c:=A^b+A_1',
\]
and, by the triangle inequality for each of the three norms (those of $A_1'$ on $\square'^{\dagger}$ do not exceed those on $T_{k+1}$), the $i$-th norm of $A^c$ on $\square'^{\dagger}$ is below $\frac{1}{2}e_{k+1}^{-\frac{3}{4}+i\varepsilon}$, $i=1,2,3$.

Let $\square$ be an $M$-cube of $T_k$. The map $x\mapsto Lx$ carries $T_{k+1}$ onto $T_k$, and the $M$-cubes of $T_{k+1}$ onto the cubes of side $LM$ of $T_k$ centred at the points of $LM\mathbb{Z}^3$; as $L$ is odd, every $M$-cube of $T_k$ lies in one of these cubes. Let $\square'$ be the $M$-cube of $T_{k+1}$ with $\square\subset L\square'$. The cube $\square^{\dagger}$ is $\square$ enlarged by $LM$ on each side, and $L(\square'^{\dagger})$ is $L\square'$ enlarged by $L^2M$ on each side; hence $\square^{\dagger}\subset L(\square'^{\dagger})$, and $x\mapsto x/L$ maps the bonds and the pairs of points of $\square^{\dagger}$ to bonds and pairs of points of $\square'^{\dagger}$. By the definition of the scaling, $((A^c)_L)_{\mu}(x)=L^{-\frac{1}{2}}A^c_{\mu}(x/L)$. Since the mesh $\eta_k$ is $L$ times the mesh $\eta_{k+1}$, the difference quotients of $(A^c)_L$ at $x$ are $L^{-\frac{3}{2}}$ times those of $A^c$ at $x/L$; and for points $x,x'$ of $T_k$ with $0<|x-x'|\le 1$ the points $x/L$, $x'/L$ of $T_{k+1}$ satisfy $0<|x/L-x'/L|\le\frac{1}{L}$, while $|x-x'|^{-\alpha}=L^{-\alpha}\cdot|x/L-x'/L|^{-\alpha}$. For $i=1,2,3$ it follows that the $i$-th norm of $(A^c)_L$, formed on $\square^{\dagger}$, is bounded by the $i$-th norm of $A^c$, formed on $\square'^{\dagger}$, multiplied by $L^{-\frac{1}{2}}$ if $i=1$, by $L^{-\frac{3}{2}}$ if $i=2$ and by $L^{-\frac{3}{2}-\alpha}$ if $i=3$. Moreover
\[
e_{k+1}^{-\frac{3}{4}+i\varepsilon}=L^{-\frac{3}{8}+\frac{i\varepsilon}{2}}\,e_k^{-\frac{3}{4}+i\varepsilon}.
\]
Each of the three factors is at most $L^{-\frac{1}{2}}$, and $i\le 3$; therefore the $i$-th norm of $(A^c)_L$ on $\square^{\dagger}$ is below
\[
\frac{1}{2}L^{-\frac{7}{8}+\frac{3\varepsilon}{2}}\,e_k^{-\frac{3}{4}+i\varepsilon}\le\frac{1}{2}L^{-\frac{3}{4}}\,e_k^{-\frac{3}{4}+i\varepsilon},
\]
because $\varepsilon\le\bar{\varepsilon}=\frac{1}{200}\le\frac{1}{12}$ (stage~\ref{it:order-choice-1}).

Put $(\lambda')_L(x):=L^{\frac{1}{2}}\lambda'(x/L)$ for the sites $x$ of $\square^{\dagger}$. By Remark~\ref{rem:operators-photon-3}(iii), $d((\lambda')_L)=(d\lambda')_L$ on the bonds of $\square^{\dagger}$; so, on $\square^{\dagger}$,
\[
H_k\left(A^{\min}_k(A')+\mathrm{Sq}_k\zeta\right)=(A^c)_L+\mathscr{Z}+d\left((\lambda')_L-\boldsymbol{\omega}_k(A')\right).
\]
The bond function $A^{\min}_k(A')+\mathrm{Sq}_k\zeta$ is real, since $A^{\min}_k$ and $\mathrm{Sq}_k$ are real (Proposition~\ref{prop:symmetries-reality}); so the left side is a real field (Proposition~\ref{prop:background-field}), and the fields $(A^c)_L$, $\mathscr{Z}=H_k\mathrm{Sq}_k\zeta$, $(\lambda')_L$ and $\boldsymbol{\omega}_k(A')$ are real as well. By (iii), applied to $\mathsf{Z}=\zeta$, the three norms of $\mathscr{Z}$ are at most $\frac{1}{2}e_k^{-\frac{\varepsilon}{2}}$. Hence the $i$-th norm of $(A^c)_L+\mathscr{Z}$ on $\square^{\dagger}$ is below
\[
\frac{1}{2}L^{-\frac{3}{4}}\,e_k^{-\frac{3}{4}+i\varepsilon}+\frac{1}{2}e_k^{-\frac{\varepsilon}{2}}\le\left(\frac{3}{16}+\frac{1}{16}\right)e_k^{-\frac{3}{4}+i\varepsilon},
\]
because $L^{\frac{3}{4}}\ge\frac{8}{3}$, an inequality that holds for $L\ge 4$, and because
\[
\frac{1}{2}e_k^{-\frac{\varepsilon}{2}}\le\frac{1}{16}\mathrm{thr}\le\frac{1}{16}e_k^{-\frac{3}{4}+i\varepsilon},
\]
where the first inequality is equivalent to $8\,e_k^{\frac{3}{4}-\frac{7\varepsilon}{2}}\le 1$. As $\square$ was arbitrary, $H_k\left(A^{\min}_k(A')+\mathrm{Sq}_k\zeta\right)$ lies in $\frac{1}{4}\tilde{R}_k$, with the summand $0$ in $\frac{1}{4}R_k$; that is, $A^{\min}_k(A')+\mathrm{Sq}_k\zeta\in\mathrm{Dom}_k$.
\end{proof}

\begin{remark}\label{rem:operators-photon-3}
\begin{enumerate}
\item[(i)] Required for all $e_k\in(0,e_0]$, the two inequalities $8\,C_B\cdot p_k\le e_k^{-\frac{3}{4}+3\varepsilon}$ and $8\,e_k^{\frac{3}{4}-\frac{7\varepsilon}{2}}\le 1$ of Proposition~\ref{prop:fields-domains} are conditions on $e_0$, and $e_0$ is chosen after $M$ (condition~\eqref{eq:renormalisation-group-A10} and stage~\ref{it:order-choice-9}). Indeed $C_B$ depends on $M$ through $s_d$, and $p_k=v_k^p$ is a power of $v_k=-\log e_k$, so that both inequalities hold for all $e_k\in(0,e_0]$ as soon as $e_0$ is small enough.
\item[(ii)] For the constants of Theorem~\ref{thm:photon-operators}, and for $\gamma_w$, $C_w$ as defined in Sections~\ref{sec:primitive-operators} and~\ref{sec:photon-operators}, one has
\[
C_w\ge C^P\ge C_P,\qquad \gamma_w\le\gamma^P\le\gamma_P.
\]
Moreover stages~\ref{it:order-choice-8} and~\ref{it:order-choice-4} give $\varepsilon\in(0,\bar{\varepsilon}]$ and $L^{\frac{3}{4}}\ge\frac{8}{3}$, and $8\,e_k^{\frac{3}{4}-\frac{7\varepsilon}{2}}\le 1$ implies $e_k<1$, because $\frac{3}{4}-\frac{7\varepsilon}{2}>0$. So the constants of Section~\ref{sec:order-choice} satisfy the hypotheses of Proposition~\ref{prop:backgrounds-domains} on $\varepsilon$, $L$, $e_k$, $C_w$ and $\gamma_w$, and Proposition~\ref{prop:fields-domains} follows.
\item[(iii)] For a function $\lambda'$ on a set of sites of $T_{k+1}$ let $(\lambda')_L(x):=L^{\frac{1}{2}}\lambda'(x/L)$ at the sites $x$ of $T_k$ with $x/L$ in that set. Then $d((\lambda')_L)=(d\lambda')_L$, since the mesh $\eta_k$ of $T_k$ is $L$ times the mesh $\eta_{k+1}$ of $T_{k+1}$: for $b=(x,x+\eta_ke_{\mu})$,
\[
\left(d(\lambda')_L\right)(b)=\eta_k^{-1}L^{\frac{1}{2}}\left(\lambda'\left(\frac{x}{L}+\eta_{k+1}e_{\mu}\right)-\lambda'\left(\frac{x}{L}\right)\right)=L^{-\frac{1}{2}}(d\lambda')(b/L)=(d\lambda')_L(b).
\]
\end{enumerate}
\end{remark}

\subsection{Weakened photon operators}\label{sec:weakened-photon-operators}

Throughout this section $l\in\mathrm{Cl}$ (Section~\ref{sec:lattice-geometry}), that is, $l/(LM)$ is an odd integer not smaller than $13$. Since $L$ and $M=L^m$ are odd, it follows that $l$ is odd, that $l$ is a multiple of $LM$, and that $l\ge 13LM$. Moreover
\[
M=L^m\ge L^2\ge 9,\qquad 0<\alpha_0\le\alpha^P=\frac{1}{4},
\]
where $\alpha_0$ is the exponent in the radius $M^{\alpha_0}$ of the parameter polydisk $\mathrm{Pd}_k$ (Definition~\ref{def:parameter-cubes-polydisks}). The unit sites are the sites of $T^0_k$, regarded as sites of $T_k$; the fine bonds are the bonds of $T_k$, and the unit bonds are the bonds of $T^0_k$.

\subsubsection{Parameter cubes covering connected sets}\label{sec:parameter-cubes-covering}

Let $a\in\{M,LM\}$. The cubes of $\pi'_k$ and of $\pi''_k$ are closed (Definition~\ref{def:parameter-cubes-polydisks}): for $c\in a\mathbb{Z}^3/l\mathbb{Z}^3$ the cube with the centre $c$ is the set
\[
\{x\in T_k:\ d(x,c)\le a\},
\]
and these cubes form $\pi'_k$ if $a=M$, and $\pi''_k$ if $a=LM$. The number $a$ is called the \emph{half side} of these cubes, and a cube and its centre are denoted by the same letter. Two cubes of the same half side are said to be \emph{linked} if their intersection is not empty. A set of cubes of one kind is \emph{link-connected} if any two of its members are joined by a finite sequence of members of the set in which any two consecutive cubes are linked. If $\mathrm{St}$ is a set of unit sites, those cubes of half side $a$ which meet $\mathrm{St}$ form the set
\[
\mathrm{Cub}_a(\mathrm{St}):=\left\{c\in a\mathbb{Z}^3/l\mathbb{Z}^3:\ d(x,c)\le a\ \text{for some}\ x\in\mathrm{St}\right\}.
\]

\paragraph{Rectilinear paths}
Let $p_1$ and $p_2$ be unit sites. As $l$ is odd, the difference $p_2-p_1$ has exactly one representative $z\in\mathbb{Z}^3$ with $|z_{\nu}|\le\frac{l-1}{2}$ for $\nu=0,1,2$. Put $q_0:=p_1$ and $q_{\nu+1}:=q_{\nu}+z_{\nu}e_{\nu}$ for $\nu=0,1,2$; thus $q_3=p_2$. Let
\[
\mathrm{Pth}(p_1,p_2):=\bigcup_{\nu=0}^2\left\{q_{\nu}+t\,e_{\nu}:\ t\in\mathbb{Z},\ \min(0,z_{\nu})\le t\le\max(0,z_{\nu})\right\}.
\]
For $\nu=0,1,2$ the set with the index $\nu$ in this union consists of those unit sites that lie on the segment of the direction $\nu$ joining $q_{\nu}$ to $q_{\nu+1}$. This segment has the length $|z_{\nu}|$, consecutive sites on it are nearest neighbours, and for $\nu=0,1$ it shares the end $q_{\nu+1}$ with the segment of the direction $\nu+1$. Along the three segments, taken in this order, the coordinates $0$, $1$, $2$ are changed one after the other, from $q_0=p_1$ to $q_3=p_2$. In particular, $\mathrm{Pth}(p_1,p_2)$ contains $p_1$ and $p_2$ and is connected through nearest neighbours. Since $|z_{\nu}|\le\frac{l}{2}$ for every $\nu$, one has $|z|=d(p_1,p_2)$.

\begin{lemma}[Cubes along sets of unit sites]\label{lem:cubes-along-sets}
Let $a\in\{M,LM\}$.
\begin{enumerate}
\item[(a)] For every site $x$ of $T_k$, the number of cubes of half side $a$ containing $x$ is at least $1$ and at most $27$, and any two of these cubes are linked.
\item[(b)] Let $\mathrm{St}\ne\emptyset$ be a set of unit sites, connected through nearest neighbours. The set $\mathrm{Cub}_a(\mathrm{St})$ is then link-connected.
\item[(c)] For all unit sites $p_1$ and $p_2$,
\[
\#\left|\mathrm{Cub}_a\left(\mathrm{Pth}(p_1,p_2)\right)\right|\le 9\left(\frac{3\,d(p_1,p_2)}{a}+9\right).
\]
If $\mathrm{St}$ is an $L$-block $B(y)$ or the set of the two ends of a unit bond, then $\#|\mathrm{Cub}_a(\mathrm{St})|\le 27$.
\item[(d)] For all $g\in\mathbb{G}_k$ and all unit sites $p_1$, $p_2$ one has $g\,\mathrm{Pth}(p_1,p_2)=\mathrm{Pth}(gp_1,gp_2)$. For every set $\mathrm{St}$ of unit sites, $g\,\mathrm{Cub}_M(\mathrm{St})=\mathrm{Cub}_M(g\mathrm{St})$ for all $g\in\mathbb{G}_k$, and $g\,\mathrm{Cub}_{LM}(\mathrm{St})=\mathrm{Cub}_{LM}(g\mathrm{St})$ for all $g\in\mathbb{G}_k^L$.
\end{enumerate}
\end{lemma}

\begin{proof}
The number $a$ is odd, $L^2\le a\le LM$, and hence $l\ge 13a$. Consecutive multiples of $a$ are at the distance $a$ from each other; hence the number of multiples of $a$ in a closed interval of length $t$ does not exceed $\frac{t}{a}+1$. In each of the counts below, every coordinate of the centres in question has a representative in a fixed closed interval of $\mathbb{R}$ of length less than $\frac{l}{2}+2a$; as $\frac{l}{2}+2a<l$, distinct points of such an interval are distinct modulo $l$, and the counts may be made on these representatives.

(a) Let $x$ be a site of $T_k$, and fix representatives $x_{\nu}\in\mathbb{R}$ of its coordinates. The cube $c$ contains $x$ if and only if, for $\nu=0,1,2$, the coordinate $c_{\nu}$ has a representative in the interval $[x_{\nu}-a,x_{\nu}+a]$, whose length is $2a$. Such an interval contains at least one and at most three multiples of $a$. Consequently the number of cubes of half side $a$ containing $x$ lies between $1$ and $3^3=27$. The second assertion of (a) is the definition of linked cubes.

(b) Let $x$ and $x+e_{\nu}$ be sites of $\mathrm{St}$. For $\mathfrak{l}=0,1,2$ the interval $\left[x_{\mathfrak{l}}-\frac{a-1}{2},x_{\mathfrak{l}}+\frac{a-1}{2}\right]$ consists of $a$ consecutive integers, because $a$ is odd, and therefore contains exactly one multiple of $a$; so there is a centre $c$ with $|x_{\mathfrak{l}}-c_{\mathfrak{l}}|\le\frac{a-1}{2}$ for $\mathfrak{l}=0,1,2$. As $\frac{a-1}{2}+1\le a$, both $x$ and $x+e_{\nu}$ belong to the cube $c$; therefore, by (a), each cube that contains $x$ or $x+e_{\nu}$ is linked to $c$.

Now let $c'$ and $c''$ be cubes of $\mathrm{Cub}_a(\mathrm{St})$ that contain the sites $x'$ and $x''$ of $\mathrm{St}$. As $\mathrm{St}$ is connected through nearest neighbours, there are sites $x'=x^0,x^1,\dots,x^n=x''$ of $\mathrm{St}$ such that $x^{i-1}$ and $x^i$ are nearest neighbours for $1\le i\le n$. For each $i$ let $c^i$ be the cube constructed above for the pair $x^{i-1}$, $x^i$; it belongs to $\mathrm{Cub}_a(\mathrm{St})$. In the sequence $c',c^1,\dots,c^n,c''$ any two consecutive cubes contain a common site of $\mathrm{St}$, namely one of the sites $x^0,\dots,x^n$; by (a) they coincide or are linked. If $n=0$, the sequence is $c',c''$, and both cubes contain $x'$. So $\mathrm{Cub}_a(\mathrm{St})$ is link-connected.

(c) Let $x$ be a unit site, $\nu$ a direction and $t_1$ an integer with $0\le t_1\le\frac{l-1}{2}$. If a cube $c$ contains a site of the segment
\[
\{x+t\,e_{\nu}:\ t\in\mathbb{Z},\ 0\le t\le t_1\},
\]
then for $\mathfrak{l}\ne\nu$ the coordinate $c_{\mathfrak{l}}$ has a representative in the interval $[x_{\mathfrak{l}}-a,x_{\mathfrak{l}}+a]$ of length $2a$, and $c_{\nu}$ has a representative in the interval $[x_{\nu}-a,x_{\nu}+t_1+a]$ of length $t_1+2a$. Hence at most
\[
3\cdot 3\cdot\left(\frac{t_1+2a}{a}+1\right)=9\left(\frac{t_1}{a}+3\right)
\]
cubes contain a site of the segment. The set $\mathrm{Pth}(p_1,p_2)$ is the union of three segments of this form, with $x=q_{\nu}$ or $x=q_{\nu+1}$ and $t_1=|z_{\nu}|$ for $\nu=0,1,2$, and $|z|_1\le 3|z|=3\,d(p_1,p_2)$. Summation over the three segments gives
\[
\#\left|\mathrm{Cub}_a\left(\mathrm{Pth}(p_1,p_2)\right)\right|\le 9\left(\frac{|z|_1}{a}+9\right)\le 9\left(\frac{3\,d(p_1,p_2)}{a}+9\right).
\]
If $\mathrm{St}$ is $B(y)$ or the set of the two ends of a unit bond, then in each coordinate the sites of $\mathrm{St}$ lie in an interval of length at most $L-1$. So each coordinate $c_{\nu}$ of a cube $c\in\mathrm{Cub}_a(\mathrm{St})$ has a representative in an interval of length at most $2a+L-1$, and $2a+L-1<3a$ because $a\ge L^2$. Such an interval contains at most three multiples of $a$; hence $\#|\mathrm{Cub}_a(\mathrm{St})|\le 27$.

(d) Every $g\in\mathbb{G}_k$ acts by $gx=rx+a_g$, where $r$ is a diagonal matrix with the diagonal entries $r_0,r_1,r_2\in\{1,-1\}$ and $a_g\in M\mathbb{Z}^3/l\mathbb{Z}^3$; for $g\in\mathbb{G}_k^L$ moreover $a_g\in LM\mathbb{Z}^3/l\mathbb{Z}^3$ (Section~\ref{sec:dirac-matrices-lattice}). So $g$ maps unit sites to unit sites and preserves $d$. Let $z$ and $q_0,\dots,q_3$ be as in the definition of $\mathrm{Pth}(p_1,p_2)$. Then $gp_2-gp_1=r(p_2-p_1)$, and $rz$ is the representative of $gp_2-gp_1$ whose coordinates have modulus at most $\frac{l-1}{2}$. Since $gq_0=gp_1$ and $gq_{\nu+1}=gq_{\nu}+r_{\nu}z_{\nu}e_{\nu}$, the points $gq_0,\dots,gq_3$ are the corresponding points for $\mathrm{Pth}(gp_1,gp_2)$. Moreover
\[
g(q_{\nu}+t\,e_{\nu})=gq_{\nu}+r_{\nu}t\,e_{\nu},
\]
and $r_{\nu}t$ ranges over the integers between $0$ and $r_{\nu}z_{\nu}$ when $t$ ranges over the integers between $0$ and $z_{\nu}$. Thus $g$ maps the segment of the direction $\nu$ of $\mathrm{Pth}(p_1,p_2)$ onto the segment of the direction $\nu$ of $\mathrm{Pth}(gp_1,gp_2)$, in the same order, for $\nu=0,1,2$, and $g\,\mathrm{Pth}(p_1,p_2)=\mathrm{Pth}(gp_1,gp_2)$. If $a_g\in a\mathbb{Z}^3/l\mathbb{Z}^3$, then $g$ maps the set $a\mathbb{Z}^3/l\mathbb{Z}^3$ of the centres onto itself; as $g$ preserves $d$, it maps the cube $c$ onto the cube $gc$, and the cube $c$ contains a site of $\mathrm{St}$ if and only if the cube $gc$ contains a site of $g\mathrm{St}$. This applies to $a=M$ and all $g\in\mathbb{G}_k$, and to $a=LM$ and all $g\in\mathbb{G}_k^L$; the last two assertions follow.
\end{proof}

\subsubsection{Cut-off functions and a partition of unity}\label{sec:cut-off-functions}

Define functions $\vartheta$ and $\boldsymbol{\phi}$ on $\mathbb{R}$ by
\[
\vartheta(r):=\begin{cases}0, & r\le 0,\\ 3r^2-2r^3, & 0\le r\le 1,\\ 1, & r\ge 1,\end{cases}\qquad\qquad \boldsymbol{\phi}(t):=\vartheta\left(2\left(\frac{3}{4}-|t|\right)\right).
\]
The function $\vartheta$ is non-decreasing with values in $[0,1]$, and $\vartheta(r)+\vartheta(1-r)=1$ for $0\le r\le 1$, because $3(1-r)^2-2(1-r)^3=1-3r^2+2r^3$. It is continuously differentiable, with $\vartheta'(r)=6r(1-r)$ for $0\le r\le 1$ and $\vartheta'(r)=0$ otherwise; hence $0\le\vartheta'\le\frac{3}{2}$, and $\vartheta'$ is Lipschitz continuous with the constant $6$, since $|6-12r|\le 6$ for $0\le r\le 1$. It follows that $0\le\boldsymbol{\phi}\le 1$, that $\boldsymbol{\phi}$ is even, that $\boldsymbol{\phi}(t)=0$ for $|t|\ge\frac{3}{4}$, and that $\boldsymbol{\phi}(t)=1$ for $|t|\le\frac{1}{4}$. Moreover $\boldsymbol{\phi}$ is continuously differentiable, $|\boldsymbol{\phi}'|\le 3$, and $\boldsymbol{\phi}'$ is Lipschitz continuous with the constant $24$: indeed $\boldsymbol{\phi}'$ vanishes on $\left[-\frac{1}{4},\frac{1}{4}\right]$, and
\[
\boldsymbol{\phi}'(t)=-2\operatorname{sign}(t)\,\vartheta'\left(2\left(\frac{3}{4}-|t|\right)\right)\qquad(t\ne 0).
\]
Finally,
\[
\sum_{j\in\mathbb{Z}}\boldsymbol{\phi}(t-j)=1\qquad(t\in\mathbb{R}).
\]
The left side has the period $1$, so let $0\le t<1$; then only the terms with $j=0$ and $j=1$ can differ from zero. They are $1$ and $0$ if $t\le\frac{1}{4}$, and $0$ and $1$ if $t\ge\frac{3}{4}$. If $\frac{1}{4}\le t\le\frac{3}{4}$, put $r:=2\left(\frac{3}{4}-t\right)\in[0,1]$; then $2\left(\frac{3}{4}-|t-1|\right)=1-r$, and
\[
\boldsymbol{\phi}(t)+\boldsymbol{\phi}(t-1)=\vartheta(r)+\vartheta(1-r)=1.
\]

For $c\in M\mathbb{Z}^3/l\mathbb{Z}^3$ and a point $x$ of $\mathbb{R}^3/l\mathbb{Z}^3$ put
\[
\boldsymbol{\chi}_c(x):=\prod_{\nu=0}^2\boldsymbol{\phi}\left(\frac{x_{\nu}-c_{\nu}}{M}\right),
\]
where $x-c$ is represented by the vector with $|x_{\nu}-c_{\nu}|\le\frac{l}{2}$ for $\nu=0,1,2$; as $\boldsymbol{\phi}$ is even, the value does not depend on the choice of this vector when $|x_{\nu}-c_{\nu}|=\frac{l}{2}$. For a fine bond $b$ put $\boldsymbol{\chi}_c(b):=\boldsymbol{\chi}_c(x_b)$, where $x_b$ is the barycentre of $b$ (Definition~\ref{def:barycentres}). For a direction $\mathfrak{l}$ let
\[
(\partial_{\mathfrak{l}}\boldsymbol{\chi}_c)(b):=\eta^{-1}\left(\boldsymbol{\chi}_c(b+\eta e_{\mathfrak{l}})-\boldsymbol{\chi}_c(b)\right),
\]
where $b+\eta e_{\mathfrak{l}}$ is the translate of $b$ by $\eta e_{\mathfrak{l}}$, as in Section~\ref{sec:tori-cochains-kernels}.

\begin{lemma}[A partition of unity on the fine bonds]\label{lem:partition-unity-fine}
\begin{enumerate}
\item[(a)] $\boldsymbol{\chi}_c\ge 0$, and $\sum_c\boldsymbol{\chi}_c(x)=1$ for every point $x$ of $\mathbb{R}^3/l\mathbb{Z}^3$, the sum being over $c\in M\mathbb{Z}^3/l\mathbb{Z}^3$; in particular $\sum_c\boldsymbol{\chi}_c(b)=1$ for every fine bond $b$. Moreover $\boldsymbol{\chi}_{gc}(gb)=\boldsymbol{\chi}_c(b)$ for all $g\in\mathbb{G}_k$, all $c$ and all fine bonds $b$.
\item[(b)] Let $b$ be a fine bond such that $\boldsymbol{\chi}_c(b)\ne 0$. Then both ends $x$ of $b$ satisfy $d(x,c)\le\frac{3M}{4}+\frac{1}{2}$, and $b$ lies in the cube $c$ of $\pi'_k$.
\item[(c)] Let $b$ and $\tilde b$ be fine bonds of the same direction with the barycentres $x$ and $\tilde x=x+r$, where $|r|\le 1$. Then, for every $c$ and every direction $\mathfrak{l}$,
\[
\left|(\partial_{\mathfrak{l}}\boldsymbol{\chi}_c)(b)\right|\le\frac{3}{M},\qquad \left|\boldsymbol{\chi}_c(\tilde b)-\boldsymbol{\chi}_c(b)\right|\le\frac{9|r|}{M},\qquad \left|(\partial_{\mathfrak{l}}\boldsymbol{\chi}_c)(\tilde b)-(\partial_{\mathfrak{l}}\boldsymbol{\chi}_c)(b)\right|\le\frac{42|r|}{M^2}.
\]
\end{enumerate}
\end{lemma}

\begin{proof}
For $s\in\mathbb{R}$ put
\[
f(s):=\sum_{j\in\mathbb{Z}}\boldsymbol{\phi}\left(\frac{s-jl}{M}\right).
\]
Since $\boldsymbol{\phi}$ vanishes outside $\left(-\frac{3}{4},\frac{3}{4}\right)$ and $\frac{3M}{2}<l$, the terms of this sum have pairwise disjoint supports. By the properties of $\boldsymbol{\phi}$ it follows that $0\le f\le 1$, that $f$ is continuously differentiable with $|f'|\le\frac{3}{M}$, and that $f'$ is Lipschitz continuous with the constant $\frac{24}{M^2}$. Let $x$ be a point of $\mathbb{R}^3/l\mathbb{Z}^3$ and $c\in M\mathbb{Z}^3/l\mathbb{Z}^3$, and fix representatives $x_{\nu}\in\mathbb{R}$ and $c_{\nu}\in M\mathbb{Z}$ of their coordinates. If $v_{\nu}$ is the representative of $x_{\nu}-c_{\nu}$ with $|v_{\nu}|\le\frac{l}{2}$, then $v_{\nu}=x_{\nu}-c_{\nu}-j_0l$ for some $j_0\in\mathbb{Z}$, and the terms of $f(x_{\nu}-c_{\nu})$ with $j\ne j_0$ vanish, because $|v_{\nu}+(j_0-j)l|\ge\frac{l}{2}>\frac{3M}{4}$. Hence
\[
\boldsymbol{\chi}_c(x)=\prod_{\nu=0}^2f(x_{\nu}-c_{\nu})
\]
for every choice of the representatives.

(a) Since $\boldsymbol{\phi}\ge 0$, one has $\boldsymbol{\chi}_c\ge 0$. Let $c_{\nu}$ range over the representatives $0,M,2M,\dots,l-M$ of the elements of $M\mathbb{Z}/l\mathbb{Z}$. Then $c_{\nu}+jl$, $j\in\mathbb{Z}$, ranges over $M\mathbb{Z}$, each element being attained once; hence $\frac{c_{\nu}+jl}{M}$ ranges over $\mathbb{Z}$, each integer being attained once, and the identity for the integer translates of $\boldsymbol{\phi}$, with $t=\frac{x_{\nu}}{M}$, gives
\[
\sum_{c_{\nu}}f(x_{\nu}-c_{\nu})=\sum_{c_{\nu}}\sum_{j\in\mathbb{Z}}\boldsymbol{\phi}\left(\frac{x_{\nu}}{M}-\frac{c_{\nu}+jl}{M}\right)=1.
\]
The factor of index $\nu$ of $\boldsymbol{\chi}_c(x)$ depends on $c$ only through $c_{\nu}$; so $\sum_c\boldsymbol{\chi}_c(x)$ is the product over $\nu$ of these three sums, and equals $1$. Every $g\in\mathbb{G}_k$ is the composition of a diagonal linear map with diagonal entries $\pm 1$ and of the translation by an element of $M\mathbb{Z}^3/l\mathbb{Z}^3$ (Section~\ref{sec:dirac-matrices-lattice}). Hence $gc\in M\mathbb{Z}^3/l\mathbb{Z}^3$, $g$ maps the barycentre of a fine bond $b$ to the barycentre of $gb$, and if $v$ represents $x-c$ with $|v_{\nu}|\le\frac{l}{2}$, then $gx-gc$ is represented by a vector whose coordinates are $\pm v_{\nu}$. As $\boldsymbol{\phi}$ is even, $\boldsymbol{\chi}_{gc}(gx)=\boldsymbol{\chi}_c(x)$, and therefore $\boldsymbol{\chi}_{gc}(gb)=\boldsymbol{\chi}_c(b)$.

(b) Let $\boldsymbol{\chi}_c(b)\ne 0$, let $x_b$ be the barycentre of $b$, and let $v$ represent $x_b-c$ with $|v_{\nu}|\le\frac{l}{2}$. All three factors of $\boldsymbol{\chi}_c(x_b)$ differ from zero, so $|v_{\nu}|<\frac{3M}{4}$ for $\nu=0,1,2$. The ends of $b$ are $x_b\pm\frac{\eta}{2}e_{\mu}$, where $\mu$ is the direction of $b$, and $\frac{\eta}{2}\le\frac{1}{2}$. Hence $d(x,c)<\frac{3M}{4}+\frac{1}{2}\le M$ for both ends $x$ of $b$, because $M\ge 2$; so both ends lie in the cube $c$ of $\pi'_k$, that is, $b$ is contained in $c$.

(c) Fix $c$ and $\mathfrak{l}$, and representatives as above, with $\tilde x_{\nu}=x_{\nu}+r_{\nu}$. The translates $b+\eta e_{\mathfrak{l}}$ and $\tilde b+\eta e_{\mathfrak{l}}$ have the barycentres $x+\eta e_{\mathfrak{l}}$ and $\tilde x+\eta e_{\mathfrak{l}}$, so the translation by $\eta e_{\mathfrak{l}}$ changes only the factor of index $\mathfrak{l}$ of $\boldsymbol{\chi}_c$. Put
\[
D(s):=\eta^{-1}\left(f(s+\eta)-f(s)\right)=\eta^{-1}\int_s^{s+\eta}\mathrm{d}u\,f'(u)\qquad(s\in\mathbb{R});
\]
it is the mean of $f'$ over an interval of length $\eta$. Hence $|D|\le\frac{3}{M}$ and $|D(s+r_{\mathfrak{l}})-D(s)|\le\frac{24|r_{\mathfrak{l}}|}{M^2}\le\frac{24|r|}{M^2}$. Moreover
\begin{align*}
(\partial_{\mathfrak{l}}\boldsymbol{\chi}_c)(b)&=D(x_{\mathfrak{l}}-c_{\mathfrak{l}})\cdot\prod_{\nu\ne\mathfrak{l}}f(x_{\nu}-c_{\nu}),\\
(\partial_{\mathfrak{l}}\boldsymbol{\chi}_c)(\tilde b)&=D(\tilde x_{\mathfrak{l}}-c_{\mathfrak{l}})\cdot\prod_{\nu\ne\mathfrak{l}}f(\tilde x_{\nu}-c_{\nu}).
\end{align*}
As $0\le f\le 1$, the first bound of (c) follows. For complex numbers $u_{\nu}$, $\tilde u_{\nu}$, $\nu=0,1,2$,
\[
\tilde u_0\tilde u_1\tilde u_2-u_0u_1u_2=(\tilde u_0-u_0)u_1u_2+\tilde u_0(\tilde u_1-u_1)u_2+\tilde u_0\tilde u_1(\tilde u_2-u_2).
\]
With $u_{\nu}=f(x_{\nu}-c_{\nu})$ and $\tilde u_{\nu}=f(\tilde x_{\nu}-c_{\nu})$, every factor has modulus at most $1$ and every difference has modulus at most $\frac{3|r|}{M}$; this gives
\[
\left|\boldsymbol{\chi}_c(\tilde b)-\boldsymbol{\chi}_c(b)\right|\le 3\cdot\frac{3|r|}{M}=\frac{9|r|}{M}.
\]
Now replace $u_{\mathfrak{l}}$ and $\tilde u_{\mathfrak{l}}$ by $D(x_{\mathfrak{l}}-c_{\mathfrak{l}})$ and $D(\tilde x_{\mathfrak{l}}-c_{\mathfrak{l}})$, and keep the other factors. The term that contains $\tilde u_{\mathfrak{l}}-u_{\mathfrak{l}}$ has modulus at most $\frac{24|r|}{M^2}$, and each of the two other terms has modulus at most $\frac{3}{M}\cdot\frac{3|r|}{M}$. Hence
\[
\left|(\partial_{\mathfrak{l}}\boldsymbol{\chi}_c)(\tilde b)-(\partial_{\mathfrak{l}}\boldsymbol{\chi}_c)(b)\right|\le\frac{24|r|}{M^2}+2\cdot\frac{3}{M}\cdot\frac{3|r|}{M}=\frac{42|r|}{M^2}.
\]
\end{proof}

\subsubsection{The weakened kernels}\label{sec:weakened-kernels}

Let $\gamma_P$ and $C_P$ be the constants of \eqref{eq:operators-photon-18}. Put
\begin{align}
n_P&:=402,\qquad \gamma^P:=\frac{\gamma_P}{2},\qquad C^P:=116\,C_P\cdot M^{\alpha_0n_P},\notag\\
M^P(\alpha,L,\alpha_0)&:=\min\left\{L^j:\ j\in\mathbb{Z},\ j\ge 2,\ 2\alpha_0n_P\log M\le\gamma_PM\ \text{for all}\ M\ge L^j\right\}.\label{eq:operators-photon-20}
\end{align}
The function $t\mapsto\frac{\log t}{t}$ is decreasing on $[\exp(1),\infty)$ and tends to $0$, and $L^2\ge 9>\exp(1)$. Hence the minimum exists, and for $j\ge 2$ the inequality $2\alpha_0n_P\log M\le\gamma_PM$ holds for all $M\ge L^j$ as soon as it holds for $M=L^j$. From now on
\[
M\ge M^P(\alpha,L,\alpha_0),
\]
so that $2\alpha_0n_P\log M\le\gamma_PM$.

\paragraph{The weakened minimiser}
Let $k\ge 1$. Let $c\in\pi'_k$ be a cube, whose centre $c$ is a unit site, let $b'$ be a unit bond with the ends $p'_1$ and $p'_2$, let $b$ be a fine bond, and let $s'=(s'_c)_{c\in\pi'_k}$ be a family of complex numbers. Put
\begin{align}
\mathrm{Pt}'(c,b')&:=\mathrm{Cub}_M\left(\mathrm{Pth}(c,p'_1)\cup\mathrm{Pth}(c,p'_2)\right),\notag\\
\mathrm{Pr}'(c,b';s')&:=\prod_{c''\in\mathrm{Pt}'(c,b')}s'_{c''},\notag\\
w(b,b';s')&:=\sum_{c\in\pi'_k}\boldsymbol{\chi}_c(b)\cdot\mathrm{Pr}'(c,b';s'),\notag\\
H_k(s')(b,b')&:=H_k(b,b')\cdot w(b,b';s'),\label{eq:operators-photon-21}
\end{align}
where $H_k(b,b')$ is the kernel, in the sense of Section~\ref{sec:tori-cochains-kernels}, of the minimiser $H_k$ of Proposition~\ref{prop:minimiser}. Thus $\mathrm{Pt}'(c,b')$ is a set of cubes of $\pi'_k$. The operators $H_k(s')$, $s'\in\mathrm{Pd}_k$, from the bond functions on $T^0_k$ to the bond functions on $T_k$ with the kernels \eqref{eq:operators-photon-21} constitute, for $k\ge 1$, the weakened family~\eqref{it:primitives-p1}.

\paragraph{The weakened root}
Let $k\ge 0$. For an element $\Upsilon$ of the basis of Proposition~\ref{prop:fluctuation-space}(ii) define two sets $\mathrm{St}(\Upsilon)$ and $\mathbf{P}(\Upsilon)$ of unit sites by
\[
\mathrm{St}(e_i^y):=B(y),\qquad \mathbf{P}(e_i^y):=\{y\},\qquad \mathrm{St}(\delta_b):=\mathbf{P}(\delta_b):=\text{the set of the two ends of }b.
\]
For basis elements $\Upsilon$, $\Upsilon'$ and a family $s'=(s'_c)_{c\in\pi''_k}$ of complex numbers put
\begin{align}
\mathrm{Pt}''(\Upsilon,\Upsilon')&:=\mathrm{Cub}_{LM}\left(\mathrm{St}(\Upsilon)\cup\mathrm{St}(\Upsilon')\cup\bigcup_{p\in\mathbf{P}(\Upsilon)}\ \bigcup_{p'\in\mathbf{P}(\Upsilon')}\mathrm{Pth}(p,p')\right),\notag\\
\mathrm{Pr}''(\Upsilon,\Upsilon';s')&:=\prod_{c\in\mathrm{Pt}''(\Upsilon,\Upsilon')}s'_c,\notag\\
C_k^{\frac{1}{2}}(s')(\Upsilon,\Upsilon')&:=C_k^{\frac{1}{2}}(\Upsilon,\Upsilon')\cdot\mathrm{Pr}''(\Upsilon,\Upsilon';s'),\label{eq:operators-photon-22}
\end{align}
where $C_k^{\frac{1}{2}}(\Upsilon,\Upsilon')$ is the matrix of the positive root $C_k^{\frac{1}{2}}$ (Section~\ref{sec:averaging-fluctuations-gauge}) in the orthonormal basis $\{\Upsilon\}$ of $V_k$. Thus $\mathrm{Pt}''(\Upsilon,\Upsilon')$ is a set of cubes of $\pi''_k$. The operators $C_k^{\frac{1}{2}}(s')$, $s'\in\mathrm{Pd}_k$, on $V_k$ whose matrices in this basis are given by \eqref{eq:operators-photon-22} constitute, for $k\ge 0$, the weakened family~\eqref{it:primitives-p2}.

\begin{lemma}[The sets of cubes of the weakened kernels]\label{lem:sets-cubes-weakened}
\begin{enumerate}
\item[(a)] Let $c\in\pi'_k$, and let $b'$ be a unit bond. Then $\mathrm{Pt}'(c,b')$ is a link-connected set of cubes; among its elements are $c$ itself and every cube of $\pi'_k$ to which an end of $b'$ belongs. Likewise, for any basis elements $\Upsilon$ and $\Upsilon'$, the set $\mathrm{Pt}''(\Upsilon,\Upsilon')$ is link-connected, and each cube of $\pi''_k$ containing a point of $\operatorname{supp}\Upsilon\cup\operatorname{supp}\Upsilon'$ belongs to it.
\item[(b)] Let $b$ be a fine bond, and let $x$ be its initial point. Suppose that $c\in\pi'_k$ and that some fine bond $\hat b$ with the same direction as $b$ satisfies $\boldsymbol{\chi}_c(\hat b)\ne 0$ and $d(x,\hat x)\le 2$, where $\hat x$ is the initial point of $\hat b$. Then, for every unit bond $b'$,
\[
\#|\mathrm{Pt}'(c,b')|\le n_P\left(1+\frac{d(b,b')}{M}\right).
\]
For fixed $b$, at most $27$ cubes $c$ have the property just described. Moreover, for all basis elements $\Upsilon$, $\Upsilon'$,
\[
\#|\mathrm{Pt}''(\Upsilon,\Upsilon')|\le n_P\left(1+\frac{d(\Upsilon,\Upsilon')}{M}\right).
\]
\item[(c)] $g\,\mathrm{Pt}'(c,b')=\mathrm{Pt}'(gc,gb')$ for all $g\in\mathbb{G}_k$. For all $g\in\mathbb{G}_k^L$, $g\,\mathrm{Pt}''(\Upsilon,\Upsilon')=\mathrm{Pt}''(\Upsilon_g,\Upsilon'_g)$, where $\Upsilon_g$ denotes the basis element with $g\Upsilon=\pm\Upsilon_g$, and $\Upsilon'_g$ the one with $g\Upsilon'=\pm\Upsilon'_g$.
\item[(d)] Let $s'\in\mathrm{Pd}_k$, with the index set $\pi'_k$, resp.\ $\pi''_k$. For $c$ and $b$ as in (b) and every unit bond $b'$,
\[
\left|\mathrm{Pr}'(c,b';s')\right|\le M^{\alpha_0n_P}\cdot\exp\left(\frac{\gamma_P\,d(b,b')}{2}\right),
\]
and for all basis elements $\Upsilon$, $\Upsilon'$,
\[
\left|\mathrm{Pr}''(\Upsilon,\Upsilon';s')\right|\le M^{\alpha_0n_P}\cdot\exp\left(\frac{\gamma_P\,d(\Upsilon,\Upsilon')}{2}\right).
\]
\end{enumerate}
\end{lemma}

\begin{proof}
(a) The paths $\mathrm{Pth}(c,p'_1)$ and $\mathrm{Pth}(c,p'_2)$ are connected through nearest neighbours, and both contain $c$; so their union is connected through nearest neighbours. Let
\[
S:=\mathrm{St}(\Upsilon)\cup\mathrm{St}(\Upsilon')\cup\bigcup_{p\in\mathbf{P}(\Upsilon)}\ \bigcup_{p'\in\mathbf{P}(\Upsilon')}\mathrm{Pth}(p,p').
\]
The sets $\mathrm{St}(\Upsilon)$ and $\mathrm{St}(\Upsilon')$ are connected through nearest neighbours, being an $L$-block or the two ends of a unit bond, and they contain $\mathbf{P}(\Upsilon)$, resp.\ $\mathbf{P}(\Upsilon')$; each path $\mathrm{Pth}(p,p')$ is connected through nearest neighbours and contains $p\in\mathrm{St}(\Upsilon)$ and $p'\in\mathrm{St}(\Upsilon')$. So $S$ is connected through nearest neighbours. By Lemma~\ref{lem:cubes-along-sets}(b), $\mathrm{Pt}'(c,b')$ and $\mathrm{Pt}''(\Upsilon,\Upsilon')$ are link-connected. The centre $c$ lies in $\mathrm{Pth}(c,p'_1)$ and in the cube $c$, so that $c\in\mathrm{Pt}'(c,b')$. The ends $p'_1$, $p'_2$ of $b'$ lie in $\mathrm{Pth}(c,p'_1)\cup\mathrm{Pth}(c,p'_2)$, so that every cube of $\pi'_k$ to which an end of $b'$ belongs is an element of $\mathrm{Pt}'(c,b')$. Finally, $\operatorname{supp}\Upsilon\subset\mathrm{St}(\Upsilon)$: by Proposition~\ref{prop:fluctuation-space}(i), (ii), $e_i^y$ vanishes off the interior bonds of $B(y)$, and $\operatorname{supp}\delta_b$ is the set of the ends of $b$. Likewise $\operatorname{supp}\Upsilon'\subset\mathrm{St}(\Upsilon')$. Hence each cube of $\pi''_k$ containing a point of $\operatorname{supp}\Upsilon\cup\operatorname{supp}\Upsilon'$ belongs to $\mathrm{Pt}''(\Upsilon,\Upsilon')$.

(b) By Lemma~\ref{lem:partition-unity-fine}(b), applied to $\hat b$, one has $d(\hat x,c)\le\frac{3M}{4}+\frac{1}{2}$, and $d(x,\hat x)\le 2$; hence, as $M\ge 9$,
\[
d(c,x)\le\frac{3M}{4}+\frac{1}{2}+2\le M+1.
\]
Let $p'_1$ and $p'_2$ be the ends of $b'$. The two ends of $b$ have the distance $\eta\le 1$, and those of $b'$ the distance $1$; so $d(x,p'_i)\le d(b,b')+2$, and $d(c,p'_i)\le d(b,b')+M+3$ for $i=1,2$. The set $\mathrm{Cub}_M$ of a union is the union of the sets $\mathrm{Cub}_M$ of its parts; therefore Lemma~\ref{lem:cubes-along-sets}(c) with $a=M$, applied to the two paths, and $M\ge 9$ give
\begin{align*}
\#|\mathrm{Pt}'(c,b')|&\le 18\left(\frac{3\left(d(b,b')+M+3\right)}{M}+9\right)=\frac{54\,d(b,b')}{M}+216+\frac{162}{M}\\
&\le\frac{54\,d(b,b')}{M}+234\le n_P\left(1+\frac{d(b,b')}{M}\right).
\end{align*}
For every cube $c$ as in (b), and for representatives of the coordinates chosen with $|x_{\nu}-\hat x_{\nu}|\le 2$, the bound $d(\hat x,c)\le\frac{3M}{4}+\frac{1}{2}$ shows that each coordinate $c_{\nu}$ has a representative in the interval $\left[x_{\nu}-\frac{3M}{4}-\frac{5}{2},\,x_{\nu}+\frac{3M}{4}+\frac{5}{2}\right]$, whose length is $\frac{3M}{2}+5\le 3M-8$, as $M\ge 9$. This length is less than $l$, so that distinct multiples of $M$ in the interval are distinct modulo $l$, and the interval contains at most $\frac{3M-8}{M}+1<4$ multiples of $M$, hence at most three. So each coordinate of $c$ takes at most three values, and the number of these cubes is at most $3^3=27$.

Now let $p\in\mathbf{P}(\Upsilon)$ and $p'\in\mathbf{P}(\Upsilon')$. Every point of $\operatorname{supp}\Upsilon$ has the distance at most $L$ from $p$: if $\Upsilon=e_i^y$, then $p=y$ and the point lies in $B(y)$; if $\Upsilon=\delta_b$, then $p$ and the point are ends of $b$. The same holds for $\Upsilon'$ and $p'$; hence $d(p,p')\le d(\Upsilon,\Upsilon')+2L$. There are at most four such pairs $(p,p')$. By Lemma~\ref{lem:cubes-along-sets}(c) with $a=LM$, each of the sets $\mathrm{Cub}_{LM}(\mathrm{St}(\Upsilon))$ and $\mathrm{Cub}_{LM}(\mathrm{St}(\Upsilon'))$ has at most $27$ elements, and
\[
\#\left|\mathrm{Cub}_{LM}\left(\mathrm{Pth}(p,p')\right)\right|\le 9\left(\frac{3\left(d(\Upsilon,\Upsilon')+2L\right)}{LM}+9\right).
\]
Therefore, as $M\ge 9$ and $L\ge 3$,
\begin{align*}
\#|\mathrm{Pt}''(\Upsilon,\Upsilon')|&\le 54+36\left(\frac{3\left(d(\Upsilon,\Upsilon')+2L\right)}{LM}+9\right)=\frac{108\,d(\Upsilon,\Upsilon')}{LM}+378+\frac{216}{M}\\
&\le\frac{108\,d(\Upsilon,\Upsilon')}{LM}+402\le n_P\left(1+\frac{d(\Upsilon,\Upsilon')}{M}\right).
\end{align*}

(c) Let $g\in\mathbb{G}_k$. The images under $g$ of the ends $p'_1$, $p'_2$ of $b'$ are the ends of $gb'$; so Lemma~\ref{lem:cubes-along-sets}(d) gives
\[
g\,\mathrm{Pt}'(c,b')=\mathrm{Cub}_M\left(g\,\mathrm{Pth}(c,p'_1)\cup g\,\mathrm{Pth}(c,p'_2)\right)=\mathrm{Cub}_M\left(\mathrm{Pth}(gc,gp'_1)\cup\mathrm{Pth}(gc,gp'_2)\right)=\mathrm{Pt}'(gc,gb').
\]
Now let $g\in\mathbb{G}_k^L$. By Proposition~\ref{prop:fluctuation-space}(ii), $g$ permutes the set $\{\pm\Upsilon\}$, so that $\Upsilon_g$ is well defined. By Section~\ref{sec:dirac-matrices-lattice}, $gx=rx+a_g$ with a diagonal matrix $r$ with entries $\pm 1$ and $a_g\in M\mathbb{Z}^3/l\mathbb{Z}^3$, and $M\mathbb{Z}^3\subset L\mathbb{Z}^3$; hence $g$ maps $B(y)$ onto $B(gy)$ and the interior bonds of $B(y)$ onto those of $B(gy)$. If $\Upsilon=e_i^y$, then $g\Upsilon$ vanishes off the interior bonds of $B(gy)$; the only basis elements with this property are the vectors $e_j^{gy}$, so that $\Upsilon_g=e_j^{gy}$ for some $j$. If $\Upsilon=\delta_b$, then $g\Upsilon=\pm\delta_{b''}$, where $b''$ is the positively oriented bond whose ends are the images of the ends of $b$. In both cases
\[
\mathrm{St}(\Upsilon_g)=g\,\mathrm{St}(\Upsilon),\qquad \mathbf{P}(\Upsilon_g)=g\,\mathbf{P}(\Upsilon),
\]
and the same holds for $\Upsilon'$. Since $\mathbb{G}_k^L\subset\mathbb{G}_k$, Lemma~\ref{lem:cubes-along-sets}(d) gives
\[
g\,\mathrm{Pt}''(\Upsilon,\Upsilon')=\mathrm{Cub}_{LM}\left(g\,\mathrm{St}(\Upsilon)\cup g\,\mathrm{St}(\Upsilon')\cup\bigcup_{p\in\mathbf{P}(\Upsilon)}\ \bigcup_{p'\in\mathbf{P}(\Upsilon')}\mathrm{Pth}(gp,gp')\right)=\mathrm{Pt}''(\Upsilon_g,\Upsilon'_g).
\]

(d) For $s'\in\mathrm{Pd}_k$ one has $|s'_c|\le M^{\alpha_0}$ for every $c$ (Definition~\ref{def:parameter-cubes-polydisks}), and $M^{\alpha_0}\ge 1$. Either let $N$ be the number of elements of $\mathrm{Pt}'(c,b')$ and $D:=d(b,b')$, or let $N$ be the number of elements of $\mathrm{Pt}''(\Upsilon,\Upsilon')$ and $D:=d(\Upsilon,\Upsilon')$. By (b), $N\le n_P\left(1+\frac{D}{M}\right)$, and since $M\ge M^P(\alpha,L,\alpha_0)$, \eqref{eq:operators-photon-20} gives $\alpha_0n_P\cdot\frac{\log M}{M}\le\frac{\gamma_P}{2}$. Hence the modulus of $\mathrm{Pr}'(c,b';s')$, resp.\ of $\mathrm{Pr}''(\Upsilon,\Upsilon';s')$, is at most
\[
M^{\alpha_0N}\le M^{\alpha_0n_P}\cdot M^{\frac{\alpha_0n_PD}{M}}=M^{\alpha_0n_P}\cdot\exp\left(\alpha_0n_P\cdot\frac{\log M}{M}\cdot D\right)\le M^{\alpha_0n_P}\cdot\exp\left(\frac{\gamma_PD}{2}\right).
\]
\end{proof}

\subsubsection{Properties of the two families}\label{sec:properties-two-families}

\begin{theorem}[Properties of the weakened photon operators]\label{thm:properties-weakened-photon}
Let $0\le\alpha<1$, let $L\ge 3$ be an odd integer, let $0<\alpha_0\le\frac{1}{4}$, let $M=L^m\ge M^P(\alpha,L,\alpha_0)$, and let $l\in\mathrm{Cl}$. The following assertions hold for the parameters $(\alpha_0,\gamma^P,\cdot\,;C^P)$ of Section~\ref{sec:primitive-operators}, with $\gamma^P$ and $C^P$ as in~\eqref{eq:operators-photon-20}.

For every $k\ge 1$ the family $H_k(s')$ of~\eqref{eq:operators-photon-21} has the properties~\ref{it:renormalisation-group-1}, \ref{it:renormalisation-group-2}, \ref{it:renormalisation-group-3}, \ref{it:renormalisation-group-8}, \ref{it:renormalisation-group-9}, \ref{it:renormalisation-group-10} and~\ref{it:renormalisation-group-16}, and it has property~\ref{it:renormalisation-group-14} for every $g\in\mathbb{G}_k$.

For every $k\ge 0$ the family $C_k^{\frac{1}{2}}(s')$ of~\eqref{eq:operators-photon-22} has the properties~\ref{it:renormalisation-group-1}, \ref{it:renormalisation-group-2}, \ref{it:renormalisation-group-4}, \ref{it:renormalisation-group-8}, \ref{it:renormalisation-group-9}, \ref{it:renormalisation-group-10}, \ref{it:renormalisation-group-14} and~\ref{it:renormalisation-group-16}.
\end{theorem}

\begin{proof}
Throughout the proof $C_P\ge 1$ and $\gamma_P\in(0,1]$ are the constants of~\eqref{eq:operators-photon-18}. Propositions~\ref{prop:kernel-bounds} and~\ref{prop:decay-exact-root} hold with these constants, by Theorem~\ref{thm:decay-kernels}, Proposition~\ref{prop:decay-resolvents-root} and the paragraph on the constants at the end of Section~\ref{sec:fluctuation-form-inverse}. By~\eqref{eq:operators-photon-20} one has $\gamma^P=\frac{\gamma_P}{2}$ and $C^P=116\,C_P\cdot M^{\alpha_0\cdot n_P}$; moreover $M\ge L^2\ge 9$. Neither family depends on the connection $A\in\tilde{G}_k$; hence in the properties~\ref{it:renormalisation-group-2}, \ref{it:renormalisation-group-8}, \ref{it:renormalisation-group-9} and~\ref{it:renormalisation-group-14} only the dependence on the parameter family $s'$ has to be examined.

\emph{Property~\ref{it:renormalisation-group-1}.} Let $s'=1$, that is, $s'_c=1$ for every parameter cube $c$. Then every product $\mathrm{Pr}'(c,b';s')$ and every product $\mathrm{Pr}''(\Upsilon,\Upsilon';s')$ is equal to $1$, and $w(b,b';1)=\sum_c\boldsymbol{\chi}_c(b)=1$ by Lemma~\ref{lem:partition-unity-fine}(a). Hence $H_k(1)=H_k$ and $C_k^{\frac{1}{2}}(1)=C_k^{\frac{1}{2}}$.

\emph{Properties~\ref{it:renormalisation-group-2} and~\ref{it:renormalisation-group-16}.} By~\eqref{eq:operators-photon-21} and~\eqref{eq:operators-photon-22}, every kernel $H_k(s')(b,b')$ and every kernel $C_k^{\frac{1}{2}}(s')(\Upsilon,\Upsilon')$ is a polynomial in the variables $s'_c$. Its coefficients are real, since the operators $H_k$ and $C_k^{\frac{1}{2}}$ are real by Proposition~\ref{prop:symmetries-reality} (Proposition~\ref{prop:lattice-symmetries}) and the functions $\boldsymbol{\chi}_c$ are real. A polynomial in $s'$ that does not depend on $A$ is continuous on $\mathrm{Pd}_k\times\tilde{G}_k$ and analytic in $(s',A)$ in the interior, and for every fixed $s'\in\mathrm{Pd}_k$, boundary included, it is analytic in $A$: this is property~\ref{it:renormalisation-group-2}. A polynomial with real coefficients is real at real $s'$: this is property~\ref{it:renormalisation-group-16}.

\emph{Property~\ref{it:renormalisation-group-4}.} Let $k\ge 0$ and $s'\in\mathrm{Pd}_k$. Proposition~\ref{prop:decay-exact-root} and Lemma~\ref{lem:sets-cubes-weakened}(d) give
\[
\left|C_k^{\frac{1}{2}}(s')(\Upsilon,\Upsilon')\right|=\left|C_k^{\frac{1}{2}}(\Upsilon,\Upsilon')\right|\cdot\left|\mathrm{Pr}''(\Upsilon,\Upsilon';s')\right|\le C_P\cdot M^{\alpha_0\cdot n_P}\cdot\exp\left(-\frac{\gamma_P\cdot d(\Upsilon,\Upsilon')}{2}\right),
\]
and the right side is at most $C^P\cdot\exp\left(-\gamma^P\cdot d(\Upsilon,\Upsilon')\right)$.

\emph{Property~\ref{it:renormalisation-group-3}.} Let $k\ge 1$; then $\eta=L^{-k}\le\frac{1}{3}$. Fix $s'\in\mathrm{Pd}_k$, a unit bond $b'$, a fine bond $b$ with initial point $x$, a direction $\mathfrak{l}$, and a fine bond $\tilde{b}$ of the same direction as $b$ whose initial point is $\tilde{x}=x+r$, where $0<|r|\le 1$. For the fine bonds $\beta$ of this direction put
\[
H(\beta):=H_k(\beta,b'),\qquad w(\beta):=w(\beta,b';s'),\qquad \beta^{+}:=\beta+\eta e_{\mathfrak{l}},
\]
and, for every function $f$ on these bonds and every direction $\nu\in\{0,1,2\}$, let
\[
(\partial_{\nu}f)(\beta):=\eta^{-1}\cdot\left(f(\beta+\eta e_{\nu})-f(\beta)\right);
\]
thus $(\partial_{\mathfrak{l}}f)(\beta)=\eta^{-1}\cdot\left(f(\beta^{+})-f(\beta)\right)$. When $f(\beta)=O(\beta,b')$ for a kernel $O$, one has $(\partial_{\nu}f)(\beta)=(\partial_{\nu}O)(\beta,b')$, with the difference in the first argument of Section~\ref{sec:properties-systems-primitives}. In particular $(\partial_{\nu}H)(\beta)=(\partial_{\nu}H_k)(\beta,b')$, the kernel that Proposition~\ref{prop:kernel-bounds} bounds for $\nu=0,1,2$; for $\nu=\mathfrak{l}$ and $O=H_k(s')$ one obtains the difference that occurs in property~\ref{it:renormalisation-group-3}. Let $d:=d(b,b')$ and
\[
\mathrm{Pm}:=M^{\alpha_0\cdot n_P}\cdot\exp\left(\frac{\gamma_P\cdot d}{2}\right).
\]
The bonds $b$, $b^{+}$, $\tilde{b}$, $\tilde{b}^{+}$, and the translates of $b$ between $b$ and $\tilde{b}$ that are used below, have initial points at distance at most $2$ from $x$, because $|r|\le 1$ and $\eta\le\frac{1}{3}$. By Lemma~\ref{lem:sets-cubes-weakened}(b), at most $27$ cubes $c$ satisfy $\boldsymbol{\chi}_c(\beta)\ne 0$ for one of these bonds $\beta$, and by Lemma~\ref{lem:sets-cubes-weakened}(d) each of these cubes satisfies $\left|\mathrm{Pr}'(c,b';s')\right|\le\mathrm{Pm}$. Since $w(\beta)=\sum_c\boldsymbol{\chi}_c(\beta)\cdot\mathrm{Pr}'(c,b';s')$, Lemma~\ref{lem:partition-unity-fine}(a), (c) and $M\ge 9$ give
\begin{align*}
|w(\beta)|&\le\mathrm{Pm}\qquad\left(\beta\in\{b,b^{+},\tilde{b},\tilde{b}^{+}\}\right),\\
\left|(\partial_{\mathfrak{l}}w)(\beta)\right|&\le 27\cdot\frac{3}{M}\cdot\mathrm{Pm}\le 9\,\mathrm{Pm}\qquad\left(\beta\in\{b,\tilde{b}\}\right),\\
\left|w(\tilde{b}^{+})-w(b^{+})\right|&\le 27\cdot\frac{9|r|}{M}\cdot\mathrm{Pm}\le 27|r|\cdot\mathrm{Pm},\\
\left|(\partial_{\mathfrak{l}}w)(\tilde{b})-(\partial_{\mathfrak{l}}w)(b)\right|&\le 27\cdot\frac{42|r|}{M^2}\cdot\mathrm{Pm}\le 14|r|\cdot\mathrm{Pm}.
\end{align*}
The first line follows from $\boldsymbol{\chi}_c\ge 0$ and $\sum_c\boldsymbol{\chi}_c=1$, the other three lines follow from the three bounds of Lemma~\ref{lem:partition-unity-fine}(c); in the third line the barycentres of $b^{+}$ and $\tilde{b}^{+}$ differ by $r$.

By Proposition~\ref{prop:kernel-bounds},
\[
|H(b)|,\ \left|(\partial_{\mathfrak{l}}H)(b)\right|\le C_P\cdot\exp(-\gamma_P\cdot d),\qquad\left|(\partial_{\mathfrak{l}}H)(\tilde{b})-(\partial_{\mathfrak{l}}H)(b)\right|\le C_P\cdot|r|^{\alpha}\cdot\exp(-\gamma_P\cdot d).
\]
To bound $H(\tilde{b})-H(b)$, join $x$ to $\tilde{x}$ by a rectilinear path of $n:=\eta^{-1}\cdot|r|_1$ steps of length $\eta$, each parallel to a coordinate axis, and let $b=\beta_0,\beta_1,\dots,\beta_n=\tilde{b}$ be the translates of $b$ whose initial points are the successive points of this path. These initial points have distance at most $|r|\le 1$ from $x$, so that $d(\beta_j,b')\ge d-1$ for all $j$. If $\beta_{j+1}$ is the translate of $\beta_j$ by $\eta e_{\nu}$, resp.\ by $-\eta e_{\nu}$, then $H(\beta_{j+1})-H(\beta_j)$ is equal to $\eta\cdot(\partial_{\nu}H)(\beta_j)$, resp.\ to $-\eta\cdot(\partial_{\nu}H)(\beta_{j+1})$, and Proposition~\ref{prop:kernel-bounds} bounds its modulus by $\eta\cdot C_P\cdot\exp(\gamma_P)\cdot\exp(-\gamma_P\cdot d)$. Summation over the $n$ steps, with $|r|_1\le 3|r|$, yields
\[
\left|H(\tilde{b})-H(b)\right|\le 3|r|\cdot C_P\cdot\exp(\gamma_P)\cdot\exp(-\gamma_P\cdot d).
\]

By~\eqref{eq:operators-photon-21}, $H_k(s')(\beta,b')=H(\beta)\cdot w(\beta)$, and the difference quotient of this product obeys the discrete product rule
\[
\left(\partial_{\mathfrak{l}}(H\cdot w)\right)(\beta)=(\partial_{\mathfrak{l}}H)(\beta)\cdot w(\beta^{+})+H(\beta)\cdot(\partial_{\mathfrak{l}}w)(\beta).
\]
At $\beta=b$ the bounds above give $|H(b)\cdot w(b)|\le C_P\cdot\mathrm{Pm}\cdot\exp(-\gamma_P\cdot d)$, and they bound the moduli of the two terms on the right of the product rule by $C_P\cdot\mathrm{Pm}\cdot\exp(-\gamma_P\cdot d)$ and by $9\,C_P\cdot\mathrm{Pm}\cdot\exp(-\gamma_P\cdot d)$. Hence
\[
\left|H_k(s')(b,b')\right|\le C_P\cdot\mathrm{Pm}\cdot\exp(-\gamma_P\cdot d),\qquad\left|(\partial_{\mathfrak{l}}H_k(s'))(b,b')\right|\le 10\,C_P\cdot\mathrm{Pm}\cdot\exp(-\gamma_P\cdot d).
\]
The product rule at $\tilde{b}$ and at $b$ gives the identity
\begin{multline*}
\left(\partial_{\mathfrak{l}}(H\cdot w)\right)(\tilde{b})-\left(\partial_{\mathfrak{l}}(H\cdot w)\right)(b)\\
=\left[(\partial_{\mathfrak{l}}H)(\tilde{b})-(\partial_{\mathfrak{l}}H)(b)\right]\cdot w(\tilde{b}^{+})+(\partial_{\mathfrak{l}}H)(b)\cdot\left[w(\tilde{b}^{+})-w(b^{+})\right]\\
+\left[H(\tilde{b})-H(b)\right]\cdot(\partial_{\mathfrak{l}}w)(\tilde{b})+H(b)\cdot\left[(\partial_{\mathfrak{l}}w)(\tilde{b})-(\partial_{\mathfrak{l}}w)(b)\right].
\end{multline*}
One has $|r|\le|r|^{\alpha}$, because $0<|r|\le 1$ and, by hypothesis, $\alpha<1$; and $\exp(\gamma_P)\le\exp(1)$, because $\gamma_P\le 1$. With these two inequalities the bounds above show that the four terms on the right have moduli at most $1$, $27$, $27\exp(1)$ and $14$ times $|r|^{\alpha}\cdot C_P\cdot\mathrm{Pm}\cdot\exp(-\gamma_P\cdot d)$. As $1+27+27\exp(1)+14\le 116$, it follows that
\[
\left|(\partial_{\mathfrak{l}}H_k(s'))(\tilde{b},b')-(\partial_{\mathfrak{l}}H_k(s'))(b,b')\right|\le 116\,|r|^{\alpha}\cdot C_P\cdot\mathrm{Pm}\cdot\exp(-\gamma_P\cdot d).
\]
Finally $\mathrm{Pm}\cdot\exp(-\gamma_P\cdot d)=M^{\alpha_0\cdot n_P}\cdot\exp(-\gamma^P\cdot d)$ and $C^P=116\,C_P\cdot M^{\alpha_0\cdot n_P}$. The bounds on $H_k(s')(b,b')$ and on $(\partial_{\mathfrak{l}}H_k(s'))(b,b')$, and the last bound divided by $|r|^{\alpha}=|\tilde{x}-x|^{\alpha}$, which is a bound on $(\delta_{\alpha}\partial_{\mathfrak{l}}H_k(s'))((b,\tilde{b});b')$, are therefore, since $1\le 10\le 116$, the three bounds of property~\ref{it:renormalisation-group-3} with the rate $\gamma^P$ and the constant $C^P$.

\emph{Link-components.} Let $\mathscr{Y}$ be a set of parameter cubes of the kind of the family under consideration (Definition~\ref{def:terminology}(3)), that is, $\mathscr{Y}\subset\pi'_k$ for~\eqref{eq:operators-photon-21} and $\mathscr{Y}\subset\pi''_k$ for~\eqref{eq:operators-photon-22}, and let $Y$ be the union of the cubes of $\mathscr{Y}$. Let $s'\in\mathrm{Pd}_k$ satisfy $s'_c=0$ for every $c\notin\mathscr{Y}$. Let $Y_1$ be a link-component of $Y$ and $\mathscr{Y}_1$ the set of its cubes, that is, the link-component of $\mathscr{Y}$ whose union is $Y_1$ (Definition~\ref{def:parameter-cubes-polydisks}). Let $\tilde{s}\in\mathrm{Pd}_k$ satisfy $\tilde{s}_c=0$ for every $c\notin\mathscr{Y}$ and $\tilde{s}_c=s'_c$ for every $c\in\mathscr{Y}_1$. If $\mathrm{Pt}$ is a link-connected set of cubes of the same kind that contains a cube of $\mathscr{Y}_1$, then
\begin{equation}\label{eq:operators-photon-23}
\begin{cases}
\mathrm{Pt}\subset\mathscr{Y}_1 & \text{and}\quad\prod_{c\in\mathrm{Pt}}s'_c=\prod_{c\in\mathrm{Pt}}\tilde{s}_c,\qquad\text{or}\\
\mathrm{Pt}\not\subset\mathscr{Y} & \text{and}\quad\prod_{c\in\mathrm{Pt}}s'_c=\prod_{c\in\mathrm{Pt}}\tilde{s}_c=0.
\end{cases}
\end{equation}
Indeed, if $\mathrm{Pt}\subset\mathscr{Y}$, then $\mathrm{Pt}$ is a link-connected subset of $\mathscr{Y}$ that meets the link-component $\mathscr{Y}_1$, hence $\mathrm{Pt}\subset\mathscr{Y}_1$, and the two products coincide because $\tilde{s}_c=s'_c$ for $c\in\mathscr{Y}_1$; otherwise $\mathrm{Pt}$ contains a cube $c\notin\mathscr{Y}$, and $s'_c=\tilde{s}_c=0$. Moreover, for every site $p$ of $T_k$ and every cube $c$ of the same kind,
\begin{equation}\label{eq:operators-photon-24}
p\in Y_1,\quad c\in\mathscr{Y},\quad p\in c\qquad\Longrightarrow\qquad c\in\mathscr{Y}_1.
\end{equation}
Indeed, $p$ lies in a cube $c_1$ of $\mathscr{Y}_1$, and $c$ is linked to $c_1$, since both cubes contain the site $p$.

\emph{Property~\ref{it:renormalisation-group-8} for~\eqref{eq:operators-photon-21}.} Let $k\ge 1$, let $\mathscr{Y}\subset\pi'_k$, and let $Y$, $s'$, $Y_1$, $\mathscr{Y}_1$, $\tilde{s}$ be as in the preceding paragraph. Fix a fine bond $b$ that is contained in $Y_1$, and an arbitrary unit bond $b'$. In the sum that defines $w(b,b';s')$ in~\eqref{eq:operators-photon-21} only the cubes $c$ with $\boldsymbol{\chi}_c(b)\ne 0$ contribute. Such a cube contains both ends of $b$ (Lemma~\ref{lem:partition-unity-fine}(b)), and it belongs to $\mathrm{Pt}'(c,b')$ (Lemma~\ref{lem:sets-cubes-weakened}(a)). If $c\notin\mathscr{Y}$, then $s'_c=\tilde{s}_c=0$, and therefore $\mathrm{Pr}'(c,b';s')=\mathrm{Pr}'(c,b';\tilde{s})=0$. If $c\in\mathscr{Y}$, then $c\in\mathscr{Y}_1$ by~\eqref{eq:operators-photon-24}, applied to an end of $b$, and~\eqref{eq:operators-photon-23} applies to the link-connected set $\mathrm{Pt}'(c,b')$ (Lemma~\ref{lem:sets-cubes-weakened}(a)): the products $\mathrm{Pr}'(c,b';s')$ and $\mathrm{Pr}'(c,b';\tilde{s})$ coincide, and if they do not vanish, then $\mathrm{Pt}'(c,b')\subset\mathscr{Y}_1$. In that case each end of $b'$ lies in a cube of $\pi'_k$ (Lemma~\ref{lem:cubes-along-sets}(a)), every cube of $\pi'_k$ that contains an end of $b'$ belongs to $\mathrm{Pt}'(c,b')$ (Lemma~\ref{lem:sets-cubes-weakened}(a)) and hence to $\mathscr{Y}_1$, and so $b'$ is contained in $Y_1$. Consequently $H_k(s')(b,\cdot)=H_k(\tilde{s})(b,\cdot)$, and $H_k(s')(b,b')\ne 0$ implies that $b'$ is contained in $Y_1$.

\emph{Property~\ref{it:renormalisation-group-9} for~\eqref{eq:operators-photon-21}.} Suppose now that $b'$ is a unit bond contained in $Y_1$, and denote by $p'_1$ either of its two ends. Some cube $c_1\in\mathscr{Y}_1$ contains $p'_1$, and $c_1$ belongs to $\mathrm{Pt}'(c,b')$ for every $c\in\pi'_k$ (Lemma~\ref{lem:sets-cubes-weakened}(a)). By~\eqref{eq:operators-photon-23}, for every $c$ the products $\mathrm{Pr}'(c,b';s')$ and $\mathrm{Pr}'(c,b';\tilde{s})$ coincide, and $\mathrm{Pr}'(c,b';s')\ne 0$ implies $\mathrm{Pt}'(c,b')\subset\mathscr{Y}_1$, in particular $c\in\mathscr{Y}_1$. Hence $w(\cdot,b';s')=w(\cdot,b';\tilde{s})$, and the columns of $H_k(s')$ and of $H_k(\tilde{s})$ at $b'$ coincide. If $H_k(s')(b,b')\ne 0$, then $\boldsymbol{\chi}_c(b)\cdot\mathrm{Pr}'(c,b';s')\ne 0$ for some $c$, and such a $c$ lies in $\mathscr{Y}_1$; by Lemma~\ref{lem:partition-unity-fine}(b), $b$ is contained in $c$, hence in $Y_1$.

\emph{Property~\ref{it:renormalisation-group-10} for~\eqref{eq:operators-photon-21}.} If the row $H_k(s')(b,\cdot)$ does not vanish, there are a unit bond $b'$ and a cube $c$ with $\boldsymbol{\chi}_c(b)\ne 0$ and $\mathrm{Pr}'(c,b';s')\ne 0$. Since $c\in\mathrm{Pt}'(c,b')$, the factor $s'_c$ is not zero, so that $c\in\mathscr{Y}$; and $b$ is contained in $c$ (Lemma~\ref{lem:partition-unity-fine}(b)), hence in the union of the link-component of $\mathscr{Y}$ that contains $c$, which is a link-component of $Y$.

\emph{Property~\ref{it:renormalisation-group-8} for~\eqref{eq:operators-photon-22}.} Let $k\ge 0$, let $\mathscr{Y}\subset\pi''_k$, and let $Y$, $s'$, $Y_1$, $\mathscr{Y}_1$, $\tilde{s}$ be as above. Let the basis element $\Upsilon$ be contained in $Y_1$, let $p$ be a point of $\operatorname{supp}\Upsilon$, and let $c_1$ be a cube of $\mathscr{Y}_1$ that contains $p$. For every basis element $\Upsilon'$ the cube $c_1$ belongs to the link-connected set $\mathrm{Pt}''(\Upsilon,\Upsilon')$ (Lemma~\ref{lem:sets-cubes-weakened}(a)). By~\eqref{eq:operators-photon-23} the products $\mathrm{Pr}''(\Upsilon,\Upsilon';s')$ and $\mathrm{Pr}''(\Upsilon,\Upsilon';\tilde{s})$ coincide, and if they do not vanish, then $\mathrm{Pt}''(\Upsilon,\Upsilon')\subset\mathscr{Y}_1$. In that case every point of $\operatorname{supp}\Upsilon'$ lies in a cube of $\mathrm{Pt}''(\Upsilon,\Upsilon')$ (Lemmas~\ref{lem:sets-cubes-weakened}(a) and~\ref{lem:cubes-along-sets}(a)), and $\Upsilon'$ is contained in $Y_1$. So the rows of $C_k^{\frac{1}{2}}(s')$ and of $C_k^{\frac{1}{2}}(\tilde{s})$ at $\Upsilon$ coincide, and $C_k^{\frac{1}{2}}(s')(\Upsilon,\Upsilon')\ne 0$ implies that $\Upsilon'$ is contained in $Y_1$.

\emph{Property~\ref{it:renormalisation-group-9} for~\eqref{eq:operators-photon-22}.} The preceding argument uses only two properties of the set $\mathrm{Pt}''(\Upsilon,\Upsilon')$: it is link-connected, and it contains every cube of $\pi''_k$ that contains a point of $\operatorname{supp}\Upsilon$ or of $\operatorname{supp}\Upsilon'$ (Lemma~\ref{lem:sets-cubes-weakened}(a)). Both properties are symmetric in $\Upsilon$ and $\Upsilon'$, whereas the set itself need not be, since for sites $p\in\mathbf{P}(\Upsilon)$ and $p'\in\mathbf{P}(\Upsilon')$ the paths $\mathrm{Pth}(p,p')$ and $\mathrm{Pth}(p',p)$ that enter~\eqref{eq:operators-photon-22} differ in general. The argument therefore holds with the roles of $\Upsilon$ and $\Upsilon'$ exchanged. In detail, let $\Upsilon'$ be contained in $Y_1$, and let $c'_1\in\mathscr{Y}_1$ contain a point of $\operatorname{supp}\Upsilon'$. For every basis element $\Upsilon$ the cube $c'_1$ belongs to $\mathrm{Pt}''(\Upsilon,\Upsilon')$; by~\eqref{eq:operators-photon-23} the products $\mathrm{Pr}''(\Upsilon,\Upsilon';s')$ and $\mathrm{Pr}''(\Upsilon,\Upsilon';\tilde{s})$ coincide, and if they do not vanish, then $\mathrm{Pt}''(\Upsilon,\Upsilon')\subset\mathscr{Y}_1$; as $\mathrm{Pt}''(\Upsilon,\Upsilon')$ covers $\operatorname{supp}\Upsilon$ (Lemmas~\ref{lem:sets-cubes-weakened}(a) and~\ref{lem:cubes-along-sets}(a)), $\Upsilon$ is then contained in $Y_1$. So the columns of $C_k^{\frac{1}{2}}(s')$ and of $C_k^{\frac{1}{2}}(\tilde{s})$ at $\Upsilon'$ coincide, and $C_k^{\frac{1}{2}}(s')(\Upsilon,\Upsilon')\ne 0$ implies that $\Upsilon$ is contained in $Y_1$.

\emph{Property~\ref{it:renormalisation-group-10} for~\eqref{eq:operators-photon-22}.} If $C_k^{\frac{1}{2}}(s')(\Upsilon,\Upsilon')\ne 0$, then $\mathrm{Pr}''(\Upsilon,\Upsilon';s')\ne 0$, so $s'_c\ne 0$ for every $c\in\mathrm{Pt}''(\Upsilon,\Upsilon')$, and $\mathrm{Pt}''(\Upsilon,\Upsilon')$ is a link-connected subset of $\mathscr{Y}$. It is therefore contained in one link-component of $\mathscr{Y}$, and it covers $\operatorname{supp}\Upsilon$ (Lemmas~\ref{lem:sets-cubes-weakened}(a) and~\ref{lem:cubes-along-sets}(a)). Hence $\Upsilon$ is contained in a link-component of $Y$.

\emph{Property~\ref{it:renormalisation-group-14} for~\eqref{eq:operators-photon-21}.} Let $k\ge 1$ and $g\in\mathbb{G}_k$. Extend the kernels to oriented bonds so that they change sign under the reversal of either of the two bonds; the factor $w(b,b';s')$ is not affected by the orientations, because $\boldsymbol{\chi}_c(b)$ depends only on the barycentre of $b$ and $\mathrm{Pt}'(c,b')$ only on the two ends of $b'$. In terms of kernels Proposition~\ref{prop:symmetries-reality} states that $H_k(gb,gb')=H_k(b,b')$. The element $g$ acts on parameter families by $(gs')_c=s'_{g^{-1}c}$ (Section~\ref{sec:dirac-matrices-lattice}), and Lemma~\ref{lem:sets-cubes-weakened}(c) gives
\[
\mathrm{Pr}'(gc,gb';gs')=\prod_{c''\in g\mathrm{Pt}'(c,b')}s'_{g^{-1}c''}=\prod_{c''\in\mathrm{Pt}'(c,b')}s'_{c''}=\mathrm{Pr}'(c,b';s').
\]
Moreover $\boldsymbol{\chi}_{gc}(gb)=\boldsymbol{\chi}_c(b)$ by Lemma~\ref{lem:partition-unity-fine}(a), and $c\mapsto gc$ is a bijection of $\pi'_k$, since $g$ maps the set $M\mathbb{Z}^3/l\mathbb{Z}^3$ of the centres onto itself. It follows that $w(gb,gb';gs')=w(b,b';s')$ and $H_k(gs')(gb,gb')=H_k(s')(b,b')$, that is, $H_k(gs')\,g=g\,H_k(s')$ for every $g\in\mathbb{G}_k$.

\emph{Property~\ref{it:renormalisation-group-14} for~\eqref{eq:operators-photon-22}.} Let $k\ge 0$ and $g\in\mathbb{G}_k^L$. By Proposition~\ref{prop:basis-map-box}, $g$ maps $V_k$ isometrically onto itself and permutes the set $\{\pm\Upsilon\}$; write $g^{-1}\Upsilon=\operatorname{sgn}(\Upsilon)\cdot\Upsilon^{-}$ with a basis element $\Upsilon^{-}$ and $\operatorname{sgn}(\Upsilon)\in\{1,-1\}$. For a linear operator $O$ on $V_k$ the matrix of $gOg^{-1}$ is therefore
\[
\langle\Upsilon,gOg^{-1}\Upsilon'\rangle=\langle g^{-1}\Upsilon,O\,g^{-1}\Upsilon'\rangle=\operatorname{sgn}(\Upsilon)\cdot\operatorname{sgn}(\Upsilon')\cdot O(\Upsilon^{-},\Upsilon'^{-}).
\]
For $O=C_k^{\frac{1}{2}}$, Proposition~\ref{prop:symmetries-reality} states that $gC_k^{\frac{1}{2}}g^{-1}=C_k^{\frac{1}{2}}$, that is, $\operatorname{sgn}(\Upsilon)\cdot\operatorname{sgn}(\Upsilon')\cdot C_k^{\frac{1}{2}}(\Upsilon^{-},\Upsilon'^{-})=C_k^{\frac{1}{2}}(\Upsilon,\Upsilon')$. Lemma~\ref{lem:sets-cubes-weakened}(c), applied to $g^{-1}\in\mathbb{G}_k^L$, gives $g^{-1}\mathrm{Pt}''(\Upsilon,\Upsilon')=\mathrm{Pt}''(\Upsilon^{-},\Upsilon'^{-})$, and hence
\[
\mathrm{Pr}''(\Upsilon,\Upsilon';gs')=\prod_{c\in\mathrm{Pt}''(\Upsilon,\Upsilon')}s'_{g^{-1}c}=\prod_{c\in g^{-1}\mathrm{Pt}''(\Upsilon,\Upsilon')}s'_c=\mathrm{Pr}''(\Upsilon^{-},\Upsilon'^{-};s').
\]
With~\eqref{eq:operators-photon-22} these two identities give
\begin{align*}
\left(gC_k^{\frac{1}{2}}(s')g^{-1}\right)(\Upsilon,\Upsilon')&=\operatorname{sgn}(\Upsilon)\cdot\operatorname{sgn}(\Upsilon')\cdot C_k^{\frac{1}{2}}(\Upsilon^{-},\Upsilon'^{-})\cdot\mathrm{Pr}''(\Upsilon^{-},\Upsilon'^{-};s')\\
&=C_k^{\frac{1}{2}}(\Upsilon,\Upsilon')\cdot\mathrm{Pr}''(\Upsilon,\Upsilon';gs')=C_k^{\frac{1}{2}}(gs')(\Upsilon,\Upsilon'),
\end{align*}
that is, $C_k^{\frac{1}{2}}(gs')\,g=g\,C_k^{\frac{1}{2}}(s')$.
\end{proof}

\subsubsection{The diagonal family of level zero}\label{sec:diagonal-family-level}

\begin{proposition}[The family of level zero]\label{prop:family-level-zero}
For $s'=(s'_c)_{c\in\pi'_0}\in\mathrm{Pd}_0$ let $H_0(s')$ be the linear map on the bond functions on $T_0=T^0_0$ whose kernel is
\[
H_0(s')(b,b'):=\delta_{bb'}\prod_{c\in\pi'_0,\ \operatorname{supp}b\subset c}s'_c,
\]
the product being taken over the cubes $c$ of $\pi'_0$ that contain $b$. The family $H_0(s')$, $s'\in\mathrm{Pd}_0$, has properties~\ref{it:renormalisation-group-1}, \ref{it:renormalisation-group-2}, \ref{it:renormalisation-group-8}, \ref{it:renormalisation-group-9}, \ref{it:renormalisation-group-10}, \ref{it:renormalisation-group-14} (for every $g\in\mathbb{G}_0$) and~\ref{it:renormalisation-group-16}. For every $\gamma>0$ it has property~\ref{it:renormalisation-group-3} with the rate $\gamma$ and the constant $4\exp(2\gamma)\cdot M^{27\alpha_0}$.
\end{proposition}

\begin{proof}
At level $0$ one has $\eta_0=1$ and $T_0=T^0_0$, so that fine bonds and unit bonds are the same objects. For a bond $b$ denote by $\bar{C}(b)$ the set of the cubes of $\pi'_0$ that contain both ends of $b$. Then
\[
H_0(s')(b,b')=\delta_{bb'}\prod_{c\in\bar{C}(b)}s'_c.
\]
Every cube of $\bar{C}(b)$ contains the initial point of $b$, so $\#|\bar{C}(b)|\le 27$ by Lemma~\ref{lem:cubes-along-sets}(a). The set $\bar{C}(b)$ is not empty. Indeed, let $b=(x,x+e_{\mu})$, and let $c\in M\mathbb{Z}^3/l\mathbb{Z}^3$ be a centre with $|x_{\nu}-c_{\nu}|\le\frac{M-1}{2}$ for $\nu=0,1,2$; it exists because $M$ is odd. As in the proof of Lemma~\ref{lem:cubes-along-sets}(b), the cube with centre $c$ contains $x$ and $x+e_{\mu}$, since $\frac{M-1}{2}+1\le M$.

Every entry of $H_0(s')$ is a polynomial in $s'$ whose coefficients are $0$ or $1$: the entry at $(b,b)$ is the monomial $\prod_{c\in\bar{C}(b)}s'_c$, and the entries off the diagonal are $0$. At $s'=1$ each of these monomials equals $1$; hence $H_0(1)$ is the identity, which is the operator $H_0$ by Proposition~\ref{prop:minimiser}. This is property~\ref{it:renormalisation-group-1}. Polynomials are continuous on $\mathrm{Pd}_0$ and analytic in its interior, which gives property~\ref{it:renormalisation-group-2}. Since the coefficients are real, $H_0(s')$ is real for real $s'$, which is property~\ref{it:renormalisation-group-16}.

Let $\gamma>0$. For $s'\in\mathrm{Pd}_0$ every factor satisfies $|s'_c|\le M^{\alpha_0}$, and $M^{\alpha_0}\ge 1$; as $\bar{C}(b)$ has at most $27$ elements,
\[
|H_0(s')(b,b)|\le M^{27\alpha_0}\qquad\text{for every bond }b.
\]
Let $x$ be the initial point of $b$, and let $\mathfrak{l}\in\{0,1,2\}$. Because $\eta_0=1$, one has $(\partial_{\mathfrak{l}}H_0(s'))(b,b')=H_0(s')(b+e_{\mathfrak{l}},b')-H_0(s')(b,b')$; therefore the row $(\partial_{\mathfrak{l}}H_0(s'))(b,\cdot)$ has the entry $H_0(s')(b+e_{\mathfrak{l}},b+e_{\mathfrak{l}})$ at $b'=b+e_{\mathfrak{l}}$, the entry $-H_0(s')(b,b)$ at $b'=b$, and no further non-zero entry. Let $\tilde b$ be a bond of the direction of $b$ whose initial point $\tilde x$ satisfies $0<|x-\tilde x|\le 1$. Both $x$ and $\tilde x$ are sites of $\mathbb{Z}^3/l\mathbb{Z}^3$, so $|x-\tilde x|=1$ and
\[
(\delta_{\alpha}\partial_{\mathfrak{l}}H_0(s'))((b,\tilde b);b')=(\partial_{\mathfrak{l}}H_0(s'))(\tilde b,b')-(\partial_{\mathfrak{l}}H_0(s'))(b,b').
\]
This kernel is made of the four terms $H_0(s')(\tilde b+e_{\mathfrak{l}},\tilde b+e_{\mathfrak{l}})$ at $b'=\tilde b+e_{\mathfrak{l}}$, $-H_0(s')(\tilde b,\tilde b)$ at $b'=\tilde b$, $-H_0(s')(b+e_{\mathfrak{l}},b+e_{\mathfrak{l}})$ at $b'=b+e_{\mathfrak{l}}$ and $H_0(s')(b,b)$ at $b'=b$, and it vanishes at every other $b'$. Two of these four bonds may coincide, namely $\tilde b=b+e_{\mathfrak{l}}$ or $\tilde b+e_{\mathfrak{l}}=b$, but no three of them; so each entry is a sum of at most two of the four terms, and its modulus is at most $2M^{27\alpha_0}$. Whenever one of the kernels of $H_0(s')$, $\partial_{\mathfrak{l}}H_0(s')$, $\delta_{\alpha}\partial_{\mathfrak{l}}H_0(s')$ has a non-zero entry at $b'$, the initial point of $b'$ has sup distance at most $2$ from $x$, and hence $d(b,b')\le 2$. At such an entry the modulus is therefore at most
\[
2M^{27\alpha_0}\le 4\exp\left(\gamma\,(2-d(b,b'))\right)\cdot M^{27\alpha_0}=4\exp(2\gamma)\cdot M^{27\alpha_0}\cdot\exp(-\gamma\,d(b,b')),
\]
and at all other entries the three kernels vanish. So property~\ref{it:renormalisation-group-3} holds with the rate $\gamma$ and the constant $4\exp(2\gamma)\cdot M^{27\alpha_0}$.

For the locality properties, let $\mathscr{Y}$ be a set of cubes of $\pi'_0$ and $Y$ its union, let $s'\in\mathrm{Pd}_0$ satisfy $s'_c=0$ for $c\notin\mathscr{Y}$, let $Y_1$ be a link-component of $Y$ and $\mathscr{Y}_1$ the set of its cubes, and let $\tilde s\in\mathrm{Pd}_0$ satisfy $\tilde s_c=0$ for $c\notin\mathscr{Y}$ and $\tilde s_c=s'_c$ for $c\in\mathscr{Y}_1$. As $H_0(s')$ and $H_0(\tilde s)$ are diagonal, the row and the column of a bond $b$ contained in $Y_1$ can have a non-zero entry only at the place $(b,b)$, whose other index $b$ is again contained in $Y_1$. It remains to show that $H_0(s')(b,b)=H_0(\tilde s)(b,b)$ for such $b$. Let $c\in\bar{C}(b)$. If $c\notin\mathscr{Y}$, then $s'_c=\tilde s_c=0$. If $c\in\mathscr{Y}$, then $c\in\mathscr{Y}_1$, as in~\eqref{eq:operators-photon-24} in the proof of Theorem~\ref{thm:properties-weakened-photon}: an end $p$ of $b$ lies in $Y_1$, hence in a cube of $\mathscr{Y}_1$, and this cube and $c$ have the site $p$ in common, so they are linked. In this case $\tilde s_c=s'_c$. Hence the products over $\bar{C}(b)$ at $s'$ and at $\tilde s$ agree factor by factor, and properties~\ref{it:renormalisation-group-8} and~\ref{it:renormalisation-group-9} hold.

Suppose that the row of $b$ does not vanish, that is, $H_0(s')(b,b)\ne 0$. Then $s'_c\ne 0$ for all $c\in\bar{C}(b)$, so $\bar{C}(b)$ is a non-empty subset of $\mathscr{Y}$. The bond $b$ is contained in each cube of $\bar{C}(b)$, hence in the link-component of $Y$ containing any one of them. Thus the row of $b$ vanishes whenever $b$ is contained in no link-component of $Y$, which is property~\ref{it:renormalisation-group-10}.

Let $g\in\mathbb{G}_0$, and let $b=(x,x+e_{\mu})$. The set $\bar{C}(b)$ is the intersection of $\mathrm{Cub}_M(\{x\})$ and $\mathrm{Cub}_M(\{x+e_{\mu}\})$. Lemma~\ref{lem:cubes-along-sets}(d), applied to these two one-point sets, and the injectivity of $g$ give
\[
g\,\bar{C}(b)=\mathrm{Cub}_M(\{gx\})\cap\mathrm{Cub}_M(\{g(x+e_{\mu})\})=\bar{C}(gb),
\]
where $gb$ is the bond with the ends $gx$ and $g(x+e_{\mu})$. Put $m_{s'}(b):=\prod_{c\in\bar{C}(b)}s'_c$. This number depends only on the two ends of $b$, and $(H_0(s')A)(b)=m_{s'}(b)\cdot A(b)$ for every bond function $A$ and for both orientations of $b$. With $(gs')_c=s'_{g^{-1}c}$ and $(gA)(b)=A(g^{-1}b)$ one obtains
\[
m_{gs'}(b)=\prod_{c\in\bar{C}(b)}s'_{g^{-1}c}=\prod_{c\in g^{-1}\bar{C}(b)}s'_c=\prod_{c\in\bar{C}(g^{-1}b)}s'_c=m_{s'}(g^{-1}b),
\]
and therefore
\[
(H_0(gs')\,gA)(b)=m_{gs'}(b)\cdot A(g^{-1}b)=m_{s'}(g^{-1}b)\cdot A(g^{-1}b)=(g\,H_0(s')A)(b).
\]
So $H_0(gs')\,g=g\,H_0(s')$ for every $g\in\mathbb{G}_0$, in particular for every $g\in\mathbb{G}_0^L$; this is property~\ref{it:renormalisation-group-14}.

The assertions of property~\ref{it:renormalisation-group-17} on the family $H_0(s')$ follow.
\end{proof}

\subsection{The constants, the photon operators and their weakened families}\label{sec:constants-photon-operators}

\subsubsection{The constants}\label{sec:constants-2}

The list below records, for each constant of Section~\ref{sec:operators-photon-field}, the place where it is defined and the quantities on which it depends. The list follows the order of choice of Section~\ref{sec:order-choice}.
\begin{itemize}
\item The numbers $a_{\mathrm{sn}}$, $c_{\mathrm{ch}}$, $c_u$, $c_{\rho}$, $B_D$, $B_X$, $B_Y$, $B_N$, $B_G$, $B_{\Pi}$, $c_N$, $c_G$, $c_{\Pi}$, $g_2$ and $\gamma_H$ are defined in \eqref{eq:operators-photon-8}. They are absolute constants.
\item The numbers $\mathsf{A}_1$, $\mathsf{A}_2$, $\mathsf{A}_3$ are defined in \eqref{eq:operators-photon-9}, and $C_{\mathrm{De}}$ is defined in Theorem~\ref{thm:decay-kernels}. They are absolute constants.
\item The number $S_{*}(\alpha)$ is defined in \eqref{eq:operators-photon-9}, and $C_{\mathrm{Hk}}(\alpha)$ in Theorem~\ref{thm:decay-kernels}. Both are functions of $\alpha$ alone.
\item The number $c_{\mathrm{ax}}$ is defined in Lemma~\ref{lem:bonds-axial-gauge}; it equals $\frac{3}{2}$ when the weight vector is uniform and $3$ in general. It is an absolute constant.
\item The numbers $c_C=3L^{\frac{9}{2}}$, $c_1=14L^{\frac{3}{2}}$ and $d_V=2-2L^{-3}$ are defined in Proposition~\ref{prop:fluctuation-space}. They are functions of $L$ alone.
\item The numbers $c_P$ and $C_{\mathrm{Dt}}$ are defined in \eqref{eq:operators-photon-17}, and $S_P$ and $\gamma_P$ in \eqref{eq:operators-photon-18}. They are functions of $L$ alone; they contain, besides, the constant $C_Q$ of \eqref{eq:operators-photon-17}, whose value is $1$.
\item The number $C_P$ is defined in \eqref{eq:operators-photon-18}. It is a function of $L$ and $\alpha$.
\item The number $\alpha^P=\frac{1}{4}$ is fixed in the standing conditions at the beginning of Section~\ref{sec:operators-photon-field}, and again at the beginning of Section~\ref{sec:weakened-photon-operators}. It is an absolute constant.
\item The numbers $n_P=402$ and $\gamma^P=\frac{\gamma_P}{2}$ are defined in \eqref{eq:operators-photon-20}. The first is an absolute constant, the second a function of $L$ alone.
\item The number $M^P=M^P(\alpha,L,\alpha_0)$ is defined in \eqref{eq:operators-photon-20}. It is determined by $L$, $\alpha_0$, $\gamma_P$ and $n_P$, and is therefore a function of $L$ and $\alpha_0$ alone.
\item The number $C^P=116\,C_P\cdot M^{\alpha_0n_P}$ is defined in \eqref{eq:operators-photon-20}. It is a function of $\alpha$, $L$, $M$ and $\alpha_0$.
\item The number $C_B$ is defined in \eqref{eq:operators-photon-19}. It is a function of $C_P$, $\gamma_P$ and $s_d=(2L+1)M$.
\end{itemize}
The constants of this list are independent of $k$, of $l$, of the integer $N$ of Section~\ref{sec:lattice-measures-laws}, of the twist, of the harmonic part of a block field, of any mass, of $e_k$, of $\varepsilon$ and of $\kappa$. The constants $\gamma_P$, $\gamma^P$, $C_P$, $c_P$ and $c_C$ are chosen in stage~\ref{it:order-choice-5}, before $M$ is chosen. The definition of $M^P$ involves neither $C_P$ nor $\varepsilon$ nor the charge $e$. The constants $C^P$ and $C_B$ are chosen in stage~\ref{it:order-choice-7}, after $M$. Required for every $e_k\in(0,e_0]$, the two inequalities $8\,C_B\cdot p_k\le e_k^{-\frac{3}{4}+3\varepsilon}$ and $8\,e_k^{\frac{3}{4}-\frac{7\varepsilon}{2}}\le1$ displayed in Proposition~\ref{prop:fields-domains} are conditions on $e_0$; in this form they are imposed in stage~\ref{it:order-choice-9} as part of condition~\eqref{eq:renormalisation-group-A10}. The other conditions of Proposition~\ref{prop:fields-domains} on the charge are~\eqref{eq:renormalisation-group-A1} and~\eqref{eq:renormalisation-group-A2}, which stage~\ref{it:order-choice-9} imposes as well.

\subsubsection{Photon operators and their weakened families}\label{sec:photon-operators-weakened}

\begin{theorem}[Photon operators]\label{thm:photon-operators}
Let $\alpha\in(1-\bar{\varepsilon},1)$, and let $L\ge 5$ be odd. Let $\alpha^P=\frac{1}{4}$, and let $\gamma_P$, $\gamma^P$, $C_P$, $c_P$, $c_C$, $M^P$ and $C^P$ be the constants of Section~\ref{sec:constants-2}. Let $\alpha_0\in(0,\alpha^P]$, let $M=L^m\ge M^P(\alpha,L,\alpha_0)$, let $k\ge 0$ be an integer and let $l\in\mathrm{Cl}$. Then the following hold.
\begin{enumerate}
\item Propositions~\ref{prop:background-field}, \ref{prop:kernel-bounds}, \ref{prop:spectral-bounds-decay}, \ref{prop:decay-exact-root}, \ref{prop:basis-map-box}, \ref{prop:symmetries-reality}, \ref{prop:local-regularity}, \ref{prop:periodisation} and~\ref{prop:gauge-fixing-controlled} hold.
\item There are families $H_k(s')$ and $C_k^{\frac{1}{2}}(s')$ of level $k$ as in~\eqref{it:primitives-p1} and~\eqref{it:primitives-p2}, the first of which, for $k=0$, is the family $H_0(s')$ of property~\ref{it:renormalisation-group-17}, such that properties~\ref{it:renormalisation-group-1}, \ref{it:renormalisation-group-2}, \ref{it:renormalisation-group-3}, \ref{it:renormalisation-group-4}, \ref{it:renormalisation-group-8}, \ref{it:renormalisation-group-9}, \ref{it:renormalisation-group-10}, \ref{it:renormalisation-group-14} and~\ref{it:renormalisation-group-16} hold for the parameters $(\alpha_0,\gamma^P,\cdot\,;C^P)$. Moreover the assertions of property~\ref{it:renormalisation-group-17} on $H_0(s')$ hold.
\item For every $\varepsilon\in(0,\bar{\varepsilon}]$, Proposition~\ref{prop:fields-domains} holds under the conditions stated in it.
\item Proposition~\ref{prop:axial-gauge} holds.
\end{enumerate}
\end{theorem}

\begin{proof}
The propositions named in the statement are obtained as follows.
\begin{itemize}
\item Proposition~\ref{prop:background-field} is Proposition~\ref{prop:minimiser}.
\item Proposition~\ref{prop:kernel-bounds} is Theorem~\ref{thm:decay-kernels}, read with $C_P\ge C_{\mathrm{Hk}}(\alpha)$ and $\gamma_P\le\gamma_H$ from \eqref{eq:operators-photon-18}. Propositions~\ref{prop:spectral-bounds-decay} and~\ref{prop:decay-exact-root} are Propositions~\ref{prop:bounds-fluctuation-form} and~\ref{prop:decay-resolvents-root}, read with the constants of \eqref{eq:operators-photon-17} and \eqref{eq:operators-photon-18}. That the three propositions hold in this form, with $C_P\ge 1$ and $\gamma_P\in(0,1]$, is shown in the paragraph on the constants $C_P$, $\gamma_P$, $c_P$ at the end of Section~\ref{sec:fluctuation-form-inverse}.
\item Proposition~\ref{prop:basis-map-box} is Proposition~\ref{prop:fluctuation-space}.
\item Proposition~\ref{prop:symmetries-reality} is Proposition~\ref{prop:lattice-symmetries}.
\item Proposition~\ref{prop:local-regularity} is Proposition~\ref{prop:regularity-enlarged-cubes}.
\item Proposition~\ref{prop:fields-domains} is Proposition~\ref{prop:backgrounds-domains}, which holds for every $\varepsilon\in(0,\bar{\varepsilon}]$. Its hypothesis $L^{\frac{3}{4}}\ge\frac{8}{3}$ is satisfied, since $L^{\frac{3}{4}}\ge 5^{\frac{3}{4}}>3$. The constants $C_w$ and $\gamma_w$ with which condition~\eqref{eq:renormalisation-group-A2} is formed satisfy $C_w\ge C_P$ and $\gamma_w\le\gamma_P$ (Remark~\ref{rem:operators-photon-3}(ii)), and the condition $8e_k^{\frac{3}{4}-\frac{7\varepsilon}{2}}\le 1$ implies $e_k<1$.
\item Proposition~\ref{prop:periodisation} is Proposition~\ref{prop:periodisation-photon-operators}.
\item Proposition~\ref{prop:gauge-fixing-controlled} is Corollary~\ref{cor:bounded-gauge-torus}.
\item Proposition~\ref{prop:axial-gauge} is Corollary~\ref{cor:local-bounds-axial}.
\end{itemize}
For the families, note first that the hypotheses of Theorem~\ref{thm:properties-weakened-photon} are satisfied: $0\le\alpha<1$, $L$ is odd and $L\ge 3$, $0<\alpha_0\le\frac{1}{4}$, $M=L^m\ge M^P(\alpha,L,\alpha_0)$ and $l\in\mathrm{Cl}$. For $k\ge 1$ take the family $H_k(s')$ of \eqref{eq:operators-photon-21}, and for every $k\ge 0$ the family $C_k^{\frac{1}{2}}(s')$ of \eqref{eq:operators-photon-22}; by Theorem~\ref{thm:properties-weakened-photon} they have the properties listed in the statement for the parameters $(\alpha_0,\gamma^P,\cdot\,;C^P)$.

For $k=0$ the first family is $H_0(s')$. By Proposition~\ref{prop:family-level-zero} it has properties~\ref{it:renormalisation-group-1}, \ref{it:renormalisation-group-2}, \ref{it:renormalisation-group-8}, \ref{it:renormalisation-group-9}, \ref{it:renormalisation-group-10}, \ref{it:renormalisation-group-14} and~\ref{it:renormalisation-group-16}, the assertions of property~\ref{it:renormalisation-group-17} on $H_0(s')$ hold, and property~\ref{it:renormalisation-group-3} holds at every rate. At the rate $\gamma^P$ the constant of Proposition~\ref{prop:family-level-zero} satisfies
\[
4\exp(2\gamma^P)\cdot M^{27\alpha_0}\le C_P\cdot M^{\alpha_0n_P}\le C^P.
\]
Indeed, $\gamma^P=\frac{\gamma_P}{2}\le 1$ because $\gamma_P\in(0,1]$ (end of Section~\ref{sec:fluctuation-form-inverse}); $C_P\ge 4\exp(2)$ by \eqref{eq:operators-photon-18}; $n_P=402\ge 27$ and $M\ge 1$; and $C^P=116\,C_P\cdot M^{\alpha_0n_P}$ by \eqref{eq:operators-photon-20}. Hence $H_0(s')$ has property~\ref{it:renormalisation-group-3} for the parameters $(\alpha_0,\gamma^P,\cdot\,;C^P)$. At the rate $\gamma_w$ of Section~\ref{sec:photon-operators} the same proposition gives the constant $4\exp(2\gamma_w)\cdot M^{27\alpha_0}$, which is at most $C_w^1$, since $C_w^1$ is defined in Section~\ref{sec:photon-operators} as the maximum of $C_w^0$ and this number.

Since $\alpha^P=\frac{1}{4}$, $\gamma^P=\frac{\gamma_P}{2}$ with $\gamma_P\in(0,1]$, and the constants depend only on the quantities listed in Section~\ref{sec:constants-2}, Theorem~\ref{thm:photon-operators-weakened} follows.
\end{proof}

\section{Operators of the fermion field}\label{sec:operators-fermion-field}

This section contains the proofs of the propositions on the fermion operators in Section~\ref{sec:fermion-operators-bounds} and of Theorem~\ref{thm:fermion-operators-weakened}; the three fermion families of a system of primitives are constructed, together with their properties, in Theorem~\ref{thm:properties-weakened-families} for $k\ge1$ and in Theorem~\ref{thm:properties-level-zero} for $k=0$. The inverses at zero field are estimated in Section~\ref{sec:inverses-zero-field} by means of the results of Section~\ref{sec:wilson-operators-block}, and the resolvent identities of Section~\ref{sec:resolvent-identities} are also obtained, in different notation, in Proposition~\ref{prop:cf-geometry-large}.

\paragraph{Notation}
Definitions~\ref{def:constants}--\ref{def:terminology} and the definitions of Sections~\ref{sec:objects-scheme}--\ref{sec:constants} are in force. Euler's number is denoted by $\mathbf{e}:=\exp(1)$.

The charge of level $k$ is denoted by $e_k$, also for $k=0$ in a formula stated for a general level $k$. Written on its own, the number $e_0$ denotes the ceiling fixed in stage~\ref{it:order-choice-9} of the order of choice of the constants (Section~\ref{sec:order-choice}); the vector $e_0$ is the unit vector $e_{\mu}$ with $\mu=0$ (Definition~\ref{def:block-torus}).

In Lemma~\ref{lem:symmetries-elementary-operators}, in Proposition~\ref{prop:symmetries-fermion-operators-2} and in their proofs, and in the part of the proof of Theorem~\ref{thm:properties-weakened-families} that concerns gauge covariance, which is property~\ref{it:renormalisation-group-13} in the list of properties of systems of primitives (Section~\ref{sec:properties-systems-primitives}), the letter $\boldsymbol{\omega}$ denotes a real connection (Section~\ref{sec:tori-nested-blocks}), closed, that is, $d\boldsymbol{\omega}=0$, where this is stated. In Section~\ref{sec:local-inverses-walks} and in the rest of Sections~\ref{sec:inverses-torus}--\ref{sec:level-zero}, $\boldsymbol{\omega}$ denotes a walk (Definition~\ref{def:walks}).

The letters $y$, $y'$, $y''$, $\tilde{y}$, $y_1$, $y_2$ denote unit sites, that is, points of the unit lattice $\Lambda^0$, and $y'$ also points of the $L$-lattice $\Lambda_{k+1}$ (Section~\ref{sec:tori-nested-blocks}). The letter $y$ also denotes the resolvent parameter of Section~\ref{sec:operators-2}, and $y'=\pm y$ a parameter of the same kind; such a parameter occurs in two ways: as a lower index of operators, as in $F_y(A)$, $S_y(A)$ and $S_{k,y}(A)$, and of the matrices $\boldsymbol{\alpha}_y$ and $\boldsymbol{\beta}_y$; and as an element of the set $Y_k$ defined there, its range or value being stated where it enters, as in $y\in Y_k$, $y\in\mathbb{R}$, $y\ge0$ or $y=\infty$.

For a site $x\in\Lambda$ and $0\le j\le k$, let $y_j$ be the point of $\Lambda_j$ with $x\in\mathrm{Bl}_j(y_j)$; it exists and is unique, since for each $j$ the blocks $\mathrm{Bl}_j(z)$, $z\in\Lambda_j$, partition $\Lambda$ (Definition~\ref{def:block-torus}). Thus $y_j$ is the centre of the block of side $L^j\eta$ that contains $x$; here $\eta$ is the mesh of $\Lambda$, namely $\eta=\eta_k=L^{-k}$ (Section~\ref{sec:parameters-clifford-relations}), so that the sides $L^j\eta=L^{j-k}$, $0\le j\le k$, are $\eta,L\eta,\dots,1$. As the blocks of side $\eta$ are single sites and $\Delta_y=\mathrm{Bl}_k(y)$, one has $y_0=x$, and $y_k$ is the unit site $y$ with $x\in\Delta_y$. These points $y_0,y_1,\dots,y_k$ are used in Section~\ref{sec:phase-factors-lattice} and in the proofs of Lemmas~\ref{lem:bounds-phases} and~\ref{lem:symmetries-elementary-operators}.

\paragraph{Conditions on the constants}
The integer $L$ is odd and satisfies $L\ge L^F:=5$; this requirement is contained in stage~\ref{it:order-choice-4}. The exponent $\alpha_0$ satisfies $\alpha_0\le\alpha^F:=\frac{1}{54}$.

The condition on $M$ belongs to stage~\ref{it:order-choice-6} of the order of choice of the constants (Section~\ref{sec:order-choice}). It involves the constants $C_V$ and $a$ of~\eqref{eq:operators-fermion-4}, which are functions of $(b,\alpha,L)$. For $\kappa'>0$ put
\[
M^F(b,\alpha,L,\alpha_0,\kappa'):=\max\left\{4L+1,\ \left(108\pi\cdot C_V\cdot a\right)^2,\ 108\pi\cdot C_V\cdot a\cdot\exp(27\kappa')\right\}.
\]
Condition~\eqref{eq:renormalisation-group-19} is the inequality
\[
M=L^m\ge M^F(b,\alpha,L,\alpha_0,\kappa').
\]

The conditions on the ceiling $e_0$ of stage~\ref{it:order-choice-9} are the following five:
\begin{align*}
&6L\cdot e_0^{\varepsilon}\le 1 &&\text{(condition~\eqref{eq:renormalisation-group-20}),}\\
&4C_V\cdot a\cdot e_0^{\varepsilon}\le 1 &&\text{(condition~\eqref{eq:renormalisation-group-21}),}\\
&3M\cdot e_0^{\varepsilon}\le 1 &&\text{(condition~\eqref{eq:renormalisation-group-22}),}\\
&e_0^{\varepsilon}\le\frac{1}{4},\qquad e_0^{\frac{1}{4}}+e_0^{\varepsilon}\le 1,\qquad 4e_0^{\frac{1}{4}+2\varepsilon}+2e_0^{3\varepsilon}\le 1 &&\text{(condition~\eqref{eq:renormalisation-group-23}),}\\
&8c_{\mathrm{rb}}\cdot L^3\cdot C_d\cdot M^3\cdot e_0^{2\varepsilon}\le 1 &&\text{(condition~\eqref{eq:renormalisation-group-24}).}
\end{align*}
Here $c_{\mathrm{rb}}$ is the constant of stage~\ref{it:order-choice-1}, and $C_d$ is the constant~\eqref{eq:operators-fermion-10}, formed with the constant $C_w$ of the parameter set for which Proposition~\ref{prop:polymer-expansion-determinant} is applied; for the parameter set $\mathrm{par}$ of Section~\ref{sec:primitive-operators} it is the constant $C_d$ of Proposition~\ref{prop:local-expansion-determinant}. Condition~\eqref{eq:renormalisation-group-22} is the third inequality of condition~\eqref{eq:renormalisation-group-A8}. The first two inequalities of condition~\eqref{eq:renormalisation-group-23} are the first inequality of condition~\eqref{eq:renormalisation-group-A8} and its fourth inequality multiplied by $e_0^{1-\varepsilon}$. In each of these inequalities the left-hand side is an increasing function of $e_0$; consequently each of the five conditions remains valid when $e_0$ is replaced by any $e_k\in(0,e_0]$.

Let $e^F_0$ be the largest number $t\in(0,1]$ such that conditions~\eqref{eq:renormalisation-group-20}--\eqref{eq:renormalisation-group-23}, and condition~\eqref{eq:renormalisation-group-24} with $C_w=C^F$, hold with $t$ in place of $e_0$; here $C^F$ is the constant of Section~\ref{sec:fermion-operators-weakened}. This number exists, since every left-hand side is a continuous increasing function of $t$ that tends to $0$ as $t\to 0$. For the projections of systems of primitives of Definition~\ref{def:periodisation-projection}, whose constant is the constant $C_w$ of $\mathrm{par}$, condition~\eqref{eq:renormalisation-group-24} is a part of condition~\eqref{eq:renormalisation-group-A10}. Conditions~\eqref{eq:renormalisation-group-20} and~\eqref{eq:renormalisation-group-21} do not involve $M$.

\subsection{Local inverses and walks}\label{sec:local-inverses-walks}

Let $F:=F_y(A)=D_A+\mathfrak{m}+\operatorname{Pen}_y(A)$ be the operator~\eqref{eq:operators-fermion-2} on $T_k$, the sum of the Wilson operator $D_A$, the mass and the block penalty $\operatorname{Pen}_y(A)$ of Section~\ref{sec:operators-2}. It is inverted by means of the parametrix
\[
\sum_{\square}h_{\square}\,\mathrm{Sl}_{\square}\,h_{\square}
\]
and of a Neumann series. Here $\square$ ranges over the $M$-cubes of $T_k$, the functions $h_{\square}$ are those of Section~\ref{sec:quadratic-partition-unity}, $\mathrm{Sl}_{\square}$ is the local inverse attached to $\square$, and the terms of the Neumann series are indexed by the walks of Definition~\ref{def:walks}.

Fix an $M$-cube $\square$, and put $W:=\square^{\sim 1}$, the enlargement of $\square$ by one layer of $M$-cubes (Section~\ref{sec:lattice-geometry}); thus $W$ is a cube of side $3M$. Reflections in the faces of $W$ yield a torus of side $6M$, the mirror torus of Definition~\ref{def:mirror-torus-fold}, and the fold map sends this torus onto $W$. The field $A$ on $W$ is pulled back by the fold map to the mirror torus, the bonds that cross a mirror plane receiving the value $0$ (Definition~\ref{def:folded-fields-local}). This pull-back is a linear map, and it commutes with the coboundary $d$. Consequently, if some real gauge transformation makes a field small on $W$, then some real gauge transformation makes its pull-back small in the sup norm on the entire mirror torus. On the mirror torus the operator is then a small bounded perturbation of the operator at the field $0$, whose inverse is bounded as follows. For $k\ge1$, Theorem~\ref{thm:sup-holder-bounds-2}, which Proposition~\ref{prop:zero-field-levels} carries over to every block torus of level $k$, and so to the lattice $\frac{\eta}{2}(1,1,1)+\eta\mathbb{Z}^3$ of the mirror torus, gives a bound in the sup norm, a H\"older bound and exponential decay, with constants independent of the mesh. For $k=0$, Lemma~\ref{lem:zero-field-level}, obtained from the coercivity bound of Proposition~\ref{prop:wilson-operators-2} by conjugation with exponential weights, bounds the kernel of the inverse by $C_{00}\cdot\left(1+|y|\right)^{-1}\cdot\exp\left(-a_z\cdot d(x,x')\right)$, that is, with the factor $w_0(y)$ of~\eqref{eq:operators-fermion-2}; there is no H\"older bound at this level, since level $0$ has no short pairs. The two cases are combined in~\eqref{eq:operators-fermion-5}. The local inverse $\mathrm{Sl}_{\square}$ is the compression to $W$ of the inverse of the operator on the mirror torus. No choice enters this definition. The local inverse depends only on the values of $A$ in $W$; it is analytic in $A$ and covariant under gauge transformations and under the lattice symmetries; and the rows of $F\,\mathrm{Sl}_{\square}$ at the sites of $W$ all of whose nearest neighbours lie in $W$ coincide with the rows of the identity (Proposition~\ref{prop:properties-local-inverse}). Neither a boundary condition on the faces of $W$, nor a decomposition of $W$ into shells of several scales, nor a bound on the curvature of $A$ enters the construction.

The term of the Neumann series that belongs to a walk $\boldsymbol{\omega}$ has a kernel supported in $X_{\boldsymbol{\omega}}\times X_{\boldsymbol{\omega}}$, and it depends only on the values of $A$ in $X_{\boldsymbol{\omega}}$ (Proposition~\ref{prop:properties-weakened-sums}(d)). Multiplying the term of every walk $\boldsymbol{\omega}$ by the product of the parameters $s_c$ of the cubes $c\in\pi_k$ contained in $X_{\boldsymbol{\omega}}$ (Definition~\ref{def:weakened-sums-walks}) yields families that are exactly local in the rows and in the columns.

\subsection{Setting and notation}\label{sec:setting-notation}

\subsubsection{Parameters and Clifford relations}\label{sec:parameters-clifford-relations}

Throughout, the following data are given: an odd integer $L\ge 5$; numbers $b>0$ and $\alpha\in(0,1)$, and $b_L:=\frac{b}{L}$; a power $M=L^m>4L$; a number $\varepsilon\in(0,\bar{\varepsilon}]$, where $\bar{\varepsilon}$ is the number of stage~\ref{it:order-choice-1}; a level $k\ge 0$, with $\eta:=\eta_k=L^{-k}$; a side $l\in\mathrm{Cl}$; a mass $\mathfrak{m}\in(0,1]$; and a number $e_k\in(0,e_0]$. For $k\ge 1$ the number $b_k=\frac{b\left(1-L^{-1}\right)}{1-L^{-k}}$ of Section~\ref{sec:running-quantities-tori} satisfies
\[
b_{-}:=b\left(1-\frac{1}{L}\right)\le b_k\le b.
\]
The numbers $k$, $l$, $\mathfrak{m}$, $e_k$ are arbitrary in the ranges just given, and no relation between them is required. The word constant has the meaning of Definition~\ref{def:constants}: a constant is either defined by a formula from quantities that precede it in the order of choice of Section~\ref{sec:order-choice}, or chosen once and for all among the numbers that satisfy the conditions imposed on it at its stage; it depends on no quantity fixed at a later stage, and it does not vary with $k$, $l$, $\mathfrak{m}$ or $e_k$.

The matrices $\gamma_0,\dots,\gamma_3$, $\mathsf{C}$, $\varkappa$, the groups $\mathrm{Sp}$, $\mathrm{Sp}_3$, $\mathbb{G}_k$, $\mathbb{G}_k^L$, $\mathbb{G}_{k,3}$ and the map $r(\cdot)$ are those of Section~\ref{sec:dirac-matrices-lattice}. Given $S\in\mathrm{Sp}$ and $\nu\in\{0,1,2\}$, the relation $\gamma_{\nu}S=r_{\nu}(S)\,S\gamma_{\nu}$ determines the sign $r_{\nu}(S)$, and $r(S)$ is the diagonal $3\times3$ matrix whose diagonal entries are $r_0(S)$, $r_1(S)$, $r_2(S)$. The Clifford data satisfy the following relations, which are checked by multiplying the matrices $\sigma_a\otimes\sigma_b$ (Section~\ref{sec:gamma-matrices-symmetry}):
\begin{align}
&\gamma_{\mu}\ \text{is Hermitian},\qquad \gamma_{\mu}\gamma_{\nu}+\gamma_{\nu}\gamma_{\mu}=2\delta_{\mu\nu}\quad(0\le\mu,\nu\le 3);\notag\\
&\gamma_{\mu}^T=-\mathsf{C}^{-1}\cdot\gamma_{\mu}\cdot\mathsf{C}\quad(\mu\le 3),\qquad \mathsf{C}^T=\mathsf{C}^{-1}=-\mathsf{C},\qquad \mathsf{C}\ \text{is real};\notag\\
&\overline{\gamma_{\mu}}=\varkappa^{-1}\cdot\gamma_{\mu}\cdot\varkappa\quad(\mu\le 2),\qquad \overline{\gamma_3}=-\varkappa^{-1}\cdot\gamma_3\cdot\varkappa;\notag\\
&\varkappa\ \text{is real},\qquad \varkappa^T=\varkappa=\varkappa^{-1},\qquad \varkappa\cdot\mathsf{C}=-\mathsf{C}\cdot\varkappa;\notag\\
&\gamma_{\nu}S=r_{\nu}(S)\,S\gamma_{\nu}\quad(S\in\mathrm{Sp},\ \nu\le 2),\qquad r(\mathrm{Sp}_3)=\left\{\operatorname{diag}(\pm1,\pm1,\pm1)\right\}.\label{eq:operators-fermion-1}
\end{align}
For $\mu\le 2$ put
\[
\Pi^{-}_{\mu}:=\frac{1-\gamma_{\mu}}{2},\qquad \Pi^{+}_{\mu}:=\frac{1+\gamma_{\mu}}{2}.
\]
By~\eqref{eq:operators-fermion-1} these matrices are orthogonal projections, and $\Pi^{-}_{\mu}+\Pi^{+}_{\mu}=1$.

\subsubsection{Tori with nested blocks}\label{sec:tori-nested-blocks}

The operators of Section~\ref{sec:exact-fermion-operators} are considered on $T_k$ and on the mirror tori of Section~\ref{sec:local-inverses-mirror}. Both are block tori in the sense of the following definition, in which $|\cdot|$ is the sup norm on $\mathbb{R}^3$ (Definition~\ref{def:notation}(a)).

\begin{definition}[Block torus]\label{def:block-torus}
A \emph{block torus of level $k$} is a triple $\mathbb{T}=(t,R,\boldsymbol{\sigma})$, where $t\in\left\{0,\frac{\eta}{2}(1,1,1)\right\}$, $R$ is an integer multiple of $L$ with $R\ge 3L$, and $\boldsymbol{\sigma}\in\{1,-1\}$. Its lattice is $\Lambda:=t+\eta\mathbb{Z}^3$. For $0\le j\le k+1$ let $t_j:=0$ if $t=0$, and $t_j:=\frac{L^j\eta}{2}(1,1,1)$ otherwise; let $\Lambda_j:=t_j+L^j\eta\mathbb{Z}^3$, and for $z\in\Lambda_j$ let
\[
\mathrm{Bl}_j(z):=\left\{x\in\Lambda:\ |x-z|<\frac{L^j\eta}{2}\right\}
\]
be the block of side $L^j\eta$ with centre $z$.

Since $L$ is odd, $\frac{L-1}{2}$ is an integer, and $t_{j+1}-t_j$ is either $\frac{L-1}{2}L^j\eta\,(1,1,1)$ or $0$; hence $\Lambda_{j+1}\subset\Lambda_j$. Further, $\Lambda_0=\Lambda$. Every centre $z\in\Lambda_j$ lies in $\Lambda$, and $\mathrm{Bl}_j(z)$ consists of the $L^{3j}$ points $z+\eta n$ with $n\in\mathbb{Z}^3$ and $|n|\le\frac{L^j-1}{2}$, because $L^j$ is odd. The blocks of one side form a partition of $\Lambda$, and every block of side $L^{j+1}\eta$ is the union of the $L^3$ blocks of side $L^j\eta$ whose centres it contains. All these structures are invariant under the translations by $R\mathbb{Z}^3$, because $R$ is an integer multiple of $L^j\eta=L^{j-k}$ for $j\le k+1$.

The \emph{unit lattice} is $\Lambda^0:=\Lambda_k$; for $y\in\Lambda^0$ the unit block is $\Delta_y:=\mathrm{Bl}_k(y)$, which on $T_k$ is the block $B_k(y)$ of Section~\ref{sec:lattice-geometry}. The \emph{$L$-lattice} is $\Lambda_{k+1}$, and $B(y'):=\Lambda^0\cap\mathrm{Bl}_{k+1}(y')$ for $y'\in\Lambda_{k+1}$.

Let $j\in\{0,k,k+1\}$. A \emph{spinor} on $\Lambda_j$ is a function $f:\Lambda_j\to\mathbb{C}^4$ with $f(x+Re_{\mu})=\boldsymbol{\sigma}\cdot f(x)$ for all $x\in\Lambda_j$ and $\mu=0,1,2$, where $e_{\mu}$ denotes the unit vector in the direction $\mu$. The spinors on $\Lambda_0=\Lambda$ are also called spinors on $\mathbb{T}$, those on $\Lambda_k=\Lambda^0$ spinors on its unit lattice, and those on $\Lambda_{k+1}$ spinors on its $L$-lattice. A \emph{connection} is a complex function $A$ on the oriented nearest-neighbour bonds $(x,x')$ of $\Lambda$ that is invariant under the translations by $R\mathbb{Z}^3$ and satisfies $A(x',x)=-A(x,x')$. A \emph{scalar} is a function on $\Lambda$ that is invariant under the translations by $R\mathbb{Z}^3$. For a scalar $\boldsymbol{\lambda}$, a connection $A$, a bond $(x,x')$ and an oriented plaquette $p$ of $\Lambda$ put
\[
(d\boldsymbol{\lambda})(x,x'):=\frac{\boldsymbol{\lambda}(x')-\boldsymbol{\lambda}(x)}{\eta},\qquad (dA)(p):=\eta^{-1}\sum_{b\in\partial p}A(b),
\]
the bonds $b$ of the boundary $\partial p$ being oriented along $\partial p$. The distance of two sites is $d(x,x'):=\min_{n\in\mathbb{Z}^3}|x-x'-Rn|$, the sup distance modulo $R\mathbb{Z}^3$.

Let $\mathbb{T}'$ be a block torus of level $k$ with the same $R$ and $\boldsymbol{\sigma}$ as $\mathbb{T}$, and let $\Lambda'_j$ be its lattices. A \emph{morphism} $\mathbb{T}\to\mathbb{T}'$ is a map $\boldsymbol{\chi}:\mathbb{R}^3\to\mathbb{R}^3$ of the form $\boldsymbol{\chi}(x)=rx+a'$, with $r=\operatorname{diag}(\pm1,\pm1,\pm1)$ and $a'\in\mathbb{R}^3$, such that $\boldsymbol{\chi}(\Lambda_j)=\Lambda'_j$ for $0\le j\le k+1$. Since $|\boldsymbol{\chi} x-\boldsymbol{\chi} x'|=|r(x-x')|=|x-x'|$ for all $x,x'\in\mathbb{R}^3$ and $\boldsymbol{\chi}(\Lambda_j)=\Lambda'_j$, a morphism $\boldsymbol{\chi}$ sends, for $0\le j\le k+1$ and $z\in\Lambda_j$, the block of side $L^j\eta$ with centre $z$ onto the block of $\mathbb{T}'$ of the same side with centre $\boldsymbol{\chi} z\in\Lambda'_j$, and every nearest-neighbour bond of $\Lambda$ to a nearest-neighbour bond of $\Lambda'_0$. As $r$ is diagonal, $\boldsymbol{\chi}$ carries the rectilinear path $\Gamma^{\pi}(x,x')$ of Section~\ref{sec:phase-factors-lattice} onto $\Gamma^{\pi}(\boldsymbol{\chi} x,\boldsymbol{\chi} x')$, for every order $\pi$.
\end{definition}

The torus $T_k$ is the block torus $(0,l,-1)$; here $l$ is an integer multiple of $L$ with $l\ge 3L$, because $l\in\mathrm{Cl}$. Its blocks are those of Section~\ref{sec:lattice-geometry}. Every $g=(S,a')\in\mathbb{G}_k$, where $a'\in M\mathbb{Z}^3$, induces the morphism $x\mapsto r(S)x+a'$ of $T_k$ onto itself, since $M\mathbb{Z}^3\subset L^j\eta\mathbb{Z}^3$ for $j\le k+1$.

\subsubsection{Norms, pairings and kernels}\label{sec:norms-pairings-kernels}

For $v\in\mathbb{C}^4$, $|v|$ is the Euclidean norm, and for a complex $4\times4$ matrix $m$, $|m|$ is the corresponding operator norm; the entries of $m$ are bounded in modulus by $|m|$. For a function $f$ with values in $\mathbb{C}^4$ put $|f|_{\infty}:=\sup_x|f(x)|$, and let $\|O\|$ be the norm of a linear operator $O$ with respect to these sup norms. The maximum over the components, $|v|_c:=\max_{\boldsymbol{\beta}}|v_{\boldsymbol{\beta}}|$, satisfies
\[
|v|_c\le|v|\le 2|v|_c.
\]
Consequently, if $\|O\|\le C$, then the norm of $O$ is at most $2C$ when $|f|_{\infty}$ is replaced throughout by $\sup_x|f(x)|_c$, which is the sup norm of Proposition~\ref{prop:properties-actions}(b). The constants $C^F$ and $C_T$ of Section~\ref{sec:fermion-operators-weakened} contain this factor $2$; therefore every statement formulated with them holds for both norms on $\mathbb{C}^4$.

For spinors $f,g$ on $\Lambda$, on $\Lambda^0$ and on $\Lambda_{k+1}$ put, respectively,
\[
(f,g):=\sum_x\eta^3f(x)\cdot g(x),\qquad (f,g):=\sum_yf(y)\cdot g(y),\qquad (f,g):=\sum_{y'}L^3f(y')\cdot g(y'),
\]
where $v\cdot w:=\sum_{\boldsymbol{\beta}}v_{\boldsymbol{\beta}}w_{\boldsymbol{\beta}}$ is the bilinear product on $\mathbb{C}^4$ and each sum extends over one period cell, that is, over a set of sites of the lattice in question that contains exactly one site of each class modulo $R\mathbb{Z}^3$; the function $x\mapsto f(x)\cdot g(x)$ is invariant under the translations by $R\mathbb{Z}^3$, since $\boldsymbol{\sigma}^2=1$. The transpose $O^T$ of an operator $O$ is taken with respect to these pairings, and the adjoint $O^{*}$ with respect to the corresponding sesquilinear pairings, in which $f$ is replaced by $\overline{f}$. Let $j,j'\in\{0,k,k+1\}$, and let $O$ be a linear operator from the spinors on $\Lambda_{j'}$ to the spinors on $\Lambda_j$. The \emph{kernel} of $O$ is the function $(x,x')\mapsto O(x,x')$ on $\Lambda_j\times\Lambda_{j'}$, with values in the complex $4\times4$ matrices, such that
\[
(Of)(x)=\sum_{x'}O(x,x')\,f(x')\qquad\left(x\in\Lambda_j\right)
\]
for every spinor $f$ on $\Lambda_{j'}$, the sum extending over a period cell of $\Lambda_{j'}$, whichever period cell is chosen; the sum carries no weight. The kernel exists and is unique. Indeed, for $x'\in\Lambda_{j'}$ and $\zeta\in\mathbb{C}^4$ let $w_{x',\zeta}$ be the spinor on $\Lambda_{j'}$ that vanishes outside $x'+R\mathbb{Z}^3$ and takes the value $\zeta$ at $x'$. Every spinor $f$ on $\Lambda_{j'}$ is the sum of the spinors $w_{x',f(x')}$ over the sites $x'$ of any period cell; therefore the identity holds if and only if $O(x,x')\,\zeta=\left(O\,w_{x',\zeta}\right)(x)$ for all $x\in\Lambda_j$, $x'\in\Lambda_{j'}$ and $\zeta\in\mathbb{C}^4$. As $w_{x'+Re_{\mu},\zeta}=w_{x',\boldsymbol{\sigma}\zeta}$, and as $O\,w_{x',\zeta}$ is a spinor, it follows that
\[
O(x+Re_{\mu},x')=O(x,x'+Re_{\mu})=\boldsymbol{\sigma}\cdot O(x,x')\qquad(\mu=0,1,2).
\]
Hence $|O(x,x')|$ depends only on the classes of $x$ and $x'$ modulo $R\mathbb{Z}^3$, and a kernel that vanishes at a pair $(x,x')$ vanishes at every pair $(x+Rn,x'+Rn')$ with $n,n'\in\mathbb{Z}^3$. Accordingly, the statement that a kernel is supported in a set of pairs means that it vanishes at every pair $(x,x')$ such that no pair $(x+Rn,x'+Rn')$ with $n,n'\in\mathbb{Z}^3$ belongs to that set. For an operator on the spinors on the unit lattice, and for an operator from these spinors to the spinors on $\Lambda$, the identity reads $(Of)(y)=\sum_{y'}O(y,y')\,f(y')$ and $(Of)(x)=\sum_{y'}O(x,y')\,f(y')$, respectively. Since the identity defining the kernel is linear in $O$ and the kernel is unique, the kernel of a linear combination of operators from the spinors on $\Lambda_{j'}$ to the spinors on $\Lambda_j$ is the same linear combination of their kernels. Moreover, $O=0$ if and only if the kernel of $O$ vanishes identically: the kernel of the operator $0$ is the function $0$, and an operator whose kernel vanishes maps every spinor to $0$. The \emph{row} of $O$ at a site $x\in\Lambda_j$ is the function $x'\mapsto O(x,x')$ on $\Lambda_{j'}$; by the identity defining the kernel and by the formula $O(x,x')\,\zeta=\left(O\,w_{x',\zeta}\right)(x)$, the row of $O$ at $x$ determines the linear map $f\mapsto(Of)(x)$ from the spinors on $\Lambda_{j'}$ to $\mathbb{C}^4$, and it is determined by this map. The \emph{column} of $O$ at a site $x'\in\Lambda_{j'}$ is the function $x\mapsto O(x,x')$ on $\Lambda_j$, and the \emph{entries} of $O$ are the values $O(x,x')$ of its kernel. Let $h$ be a complex function on $\Lambda_j$ and $h'$ a complex function on $\Lambda_{j'}$, both invariant under the translations by $R\mathbb{Z}^3$; if $j=0$, then $h$ is a scalar, and if $j'=0$, then $h'$ is a scalar. These functions act on spinors by multiplication, $(hf)(x):=h(x)\cdot f(x)$, and, being invariant under the translations by $R\mathbb{Z}^3$, they map spinors to spinors. For every spinor $f$ on $\Lambda_{j'}$, every period cell of $\Lambda_{j'}$ and every $x\in\Lambda_j$ one has
\[
\left(\left(h\cdot O\cdot h'\right)f\right)(x)=h(x)\sum_{x'}O(x,x')\,h'(x')\,f(x')=\sum_{x'}\left(h(x)\cdot O(x,x')\cdot h'(x')\right)f(x');
\]
as the kernel is unique, the kernel of $h\cdot O\cdot h'$ is the function $(x,x')\mapsto h(x)\cdot O(x,x')\cdot h'(x')$. Let now $j'=j$. Taking for $h'$, respectively for $h$, the constant function $1$, one obtains the kernel $h(x)\cdot O(x,x')$ of $h\cdot O$ and the kernel $O(x,x')\cdot h'(x')$ of $O\cdot h'$; for $h'=h$ it follows in particular, by the linearity noted above, that the operator $h\cdot O-O\cdot h$ has the kernel
\[
\left(h\cdot O-O\cdot h\right)(x,x')=\left(h(x)-h(x')\right)\cdot O(x,x').
\]
The weights $\eta^3$ and $L^3$ of the pairings do not enter the definition, so that a factor $\eta^3$ or $L^{-3}$ in the formula for an operator is a factor of its kernel: by the formulas of Section~\ref{sec:phase-factors-lattice}, the kernel of $Q_k(A)$ is $\eta^3\cdot\boldsymbol{\phi}_A(y,x)$ at the pairs $(y,x)$ with $x\in\Delta_y$, and the kernel of $Q(A)$ is $L^{-3}\cdot U_A(y',y)$ at the pairs $(y',y)$ with $y\in B(y')$. In Convention~\ref{it:wilson-operators-5} of Section~\ref{sec:conventions}, where $(Of)(x)=\eta^3\sum_{x'}O(x,x')\,f(x')$, and in the representation $(Of)(t)=\sum_sO(t,s)\cdot f(s)\cdot\nu(s)$ of Proposition~\ref{prop:weighted-coefficient-norms} and of Section~\ref{sec:coefficients-kernels-grassmann}, with the weight $\nu=\eta^3$ at the sites of $\Lambda$, the factor $\eta^3$ stands outside the kernel; for an operator from the spinors on $\Lambda$, the kernel in that normalisation is $\eta^{-3}$ times the kernel defined here. On the torus $T_k=(0,l,-1)$ one has $R=l$ and $\boldsymbol{\sigma}=-1$; its lattice $\Lambda=\eta\mathbb{Z}^3$ is the lattice $\mathbb{L}_k$ of Section~\ref{sec:lattice-geometry}, and its unit lattice is $\mathbb{Z}^3$. For an operator between two of the spaces of spinors on these two lattices, the kernel defined here is the kernel of Section~\ref{sec:two-tori}, formed on a torus $T_i$ with $l_i=l$. Indeed, the character of the spinors in that section takes the value $-1$ at every unit vector $e_{\mu}$, so that the rule imposed there on kernels reads $O(x+l e_{\mu},x')=O(x,x'+l e_{\mu})=-O(x,x')$, which is the rule above with $R=l$ and $\boldsymbol{\sigma}=-1$; the sum there extends, without weight, over a set of representatives of the sites modulo $l\mathbb{Z}^3$, that is, over a period cell; and the factors $\eta^3$ are part of the kernel there as here. By this rule the summand $O(x,x')\,f(x')$ is invariant under the translations of $x'$ by $l\mathbb{Z}^3$, both factors acquiring the same sign; hence the sum there does not depend on the set of representatives, the kernel of Section~\ref{sec:two-tori} satisfies the identity that defines the kernel here, and the two kernels coincide, the kernel being unique.

\subsubsection{Phase factors and lattice operators}\label{sec:phase-factors-lattice}

Let $A$ be a connection on a block torus $(t,R,\boldsymbol{\sigma})$ of level $k$ (Definition~\ref{def:block-torus}), with lattice $\Lambda$, unit lattice $\Lambda^0$ and $L$-lattice $\Lambda_{k+1}$. For a lattice path $\Gamma$ put
\[
A(\Gamma):=\sum_{b\in\Gamma}A(b),
\]
each bond $b$ of $\Gamma$ being oriented in the direction in which $\Gamma$ traverses it. Let $x,x'$ be sites with $|x-x'|<R/2$, and let $\pi$ be an order of the coordinates $0,1,2$. Then $\Gamma^{\pi}(x,x')$ denotes the rectilinear path from $x$ to $x'$, taken on the shortest lift, that changes the coordinates in the order $\pi$ (Definition~\ref{def:pairs-sites-holder}), and
\[
(\tau A)(x,x'):=\frac{1}{6}\sum_{\pi}A\left(\Gamma^{\pi}(x,x')\right),
\]
the sum extending over the six orders $\pi$. Let $x$ be a site, and let $y\in\Lambda^0$ be the unit site with $x\in\Delta_y$. For $0\le j\le k$ let $y_j\in\Lambda_j$ be the centre of the block of side $L^j\eta$ that contains $x$, that is, $x\in\mathrm{Bl}_j(y_j)$; the sides of these blocks are $\eta,L\eta,\dots,1$, and $y_0=x$, $y_k=y$. Put
\[
(\tau_kA)(y,x):=\sum_{j=0}^{k-1}(\tau A)(y_{j+1},y_j).
\]
The phase factors of $A$ and the number $\mathfrak{i}(A)$ are
\begin{align*}
u_A(x,x')&:=\exp\left\{i\cdot e_k\cdot\eta\cdot A(x,x')\right\}&&\text{for a bond }(x,x'),\\
U_A(x,x')&:=\exp\left\{i\cdot e_k\cdot\eta\cdot(\tau A)(x,x')\right\}&&\text{for sites }x,x'\text{ with }|x-x'|<R/2,\\
\boldsymbol{\phi}_A(y,x)&:=\exp\left\{i\cdot e_k\cdot\eta\cdot(\tau_kA)(y,x)\right\}&&\text{for }x\in\Delta_y,\\
\boldsymbol{\phi}^{+}_A(y',x)&:=U_A(y',y)\cdot\boldsymbol{\phi}_A(y,x)&&\text{for }x\in\Delta_y\text{ and }y\in B(y'),\\
\mathfrak{i}(A)&:=e_k\cdot\left|\operatorname{Im}A\right|_{\infty}.
\end{align*}

Let $f$ be a spinor on the block torus, $v$ a spinor on its unit lattice and $w$ a spinor on its $L$-lattice. On every block torus the operators of Sections~\ref{sec:exact-fermion-operators} and~\ref{sec:fermion-integrals-blocking}, and the H\"older derivative of Definition~\ref{def:pairs-sites-holder}, are given by
\begin{align*}
\left(D_Af\right)(x)&=-\eta^{-1}\sum_{\mu=0}^2\left\{\Pi^{-}_{\mu}\cdot u_A(x,x+\eta e_{\mu})\cdot f(x+\eta e_{\mu})\right.\\
&\qquad\qquad\left.+\Pi^{+}_{\mu}\cdot u_A(x,x-\eta e_{\mu})\cdot f(x-\eta e_{\mu})-f(x)\right\},\\
\left(Q_k(A)f\right)(y)&=\eta^3\sum_{x\in\Delta_y}\boldsymbol{\phi}_A(y,x)\cdot f(x)\qquad(y\in\Lambda^0),\\
\left(Q_k^T(-A)v\right)(x)&=\boldsymbol{\phi}_A(y,x)^{-1}\cdot v(y)\qquad(x\in\Delta_y),\\
P_k(A)&:=Q_k^T(-A)\cdot Q_k(A),\\
\left(Q(A)v\right)(y')&=L^{-3}\sum_{y\in B(y')}U_A(y',y)\cdot v(y)\qquad(y'\in\Lambda_{k+1}),\\
\left(Q^T(-A)w\right)(y)&=U_A(y',y)^{-1}\cdot w(y')\qquad(y\in B(y')),\\
P(A)&:=Q^T(-A)\cdot Q(A),\qquad P_{k+1}(A):=Q_k^T(-A)\cdot P(A)\cdot Q_k(A),\\
\left(\delta_{\alpha,A}f\right)(x,x')&:=|x-x'|^{-\alpha}\cdot\left(U_A(x,x')\cdot f(x')-f(x)\right)\qquad\left(x\ne x',\ |x-x'|<R/2\right).
\end{align*}
Composing these formulas, and using that $\mathrm{Bl}_{k+1}(y')$ is the union of the blocks $\Delta_y$ with $y\in B(y')$ (Definition~\ref{def:block-torus}), one obtains for $y'\in\Lambda_{k+1}$ and $x\in\mathrm{Bl}_{k+1}(y')$
\[
\left(P_{k+1}(A)f\right)(x)=\boldsymbol{\phi}^{+}_A(y',x)^{-1}\cdot L^{-3}\cdot\eta^3\sum_{x'\in\mathrm{Bl}_{k+1}(y')}\boldsymbol{\phi}^{+}_A(y',x')\cdot f(x').
\]
For $k=0$ one has $\eta=1$, $\Delta_y=\{y\}$ and $\tau_0A=0$, the sum defining $\tau_0A$ being empty; hence $Q_0(A)=1$, $P_0(A)=1$ and $P_1(A)=P(A)$. An ordered pair $(x,x')$ of sites is called a \emph{short pair} if $0<|x-x'|<1$; for $k=0$ there are no short pairs. Throughout the proofs, inside a term that refers to a direction $\mu$, the neighbours of the sites $x$ and $x'$ in this direction are written
\[
x_{\pm}:=x\pm\eta e_{\mu},\qquad x'_{\pm}:=x'\pm\eta e_{\mu}.
\]

\subsubsection{The operators $F_y(A)$ and $S_y(A)$}\label{sec:operators-2}

Let $k\ge1$ and $y\in Y_k:=\mathbb{R}\cup\{\infty\}$. For real $y$ define the $4\times4$ matrices
\begin{align*}
\boldsymbol{\alpha}_y&:=b_k\cdot i\gamma_3y\cdot\left(b_k+i\gamma_3y\right)^{-1},\\
\boldsymbol{\beta}_y&:=b_k^2\cdot b_L\cdot\left(b_k+b_L+i\gamma_3y\right)^{-1}\cdot\left(b_k+i\gamma_3y\right)^{-1},
\end{align*}
which are functions of the matrix $\gamma_3$, and put $\boldsymbol{\alpha}_{\infty}:=b_k$, $\boldsymbol{\beta}_{\infty}:=0$. The inverses exist: $\gamma_3$ is Hermitian with $\gamma_3^2=1$, so that the eigenvalues of $b_k+i\gamma_3y$ and of $b_k+b_L+i\gamma_3y$ are $b_k\pm iy$ and $b_k+b_L\pm iy$, and $b_k>0$. Put
\[
\operatorname{Pen}_y(A):=\boldsymbol{\alpha}_y\cdot P_k(A)+\boldsymbol{\beta}_y\cdot P_{k+1}(A).
\]
For $k=0$ and $y\in Y_0:=\mathbb{R}$ put $\operatorname{Pen}_y(A):=b_L\cdot P(A)+i\gamma_3y$. In both cases let
\begin{gather}
F_y(A):=D_A+\mathfrak{m}+\operatorname{Pen}_y(A),\qquad S_y(A):=F_y(A)^{-1},\notag\\
w_k(y):=1\quad(k\ge1),\qquad w_0(y):=\left(1+|y|\right)^{-1},\label{eq:operators-fermion-2}
\end{gather}
the operator $S_y(A)$ being defined whenever $F_y(A)$ is invertible.

Now let $k\ge1$. Since $\gamma_3$ is Hermitian with $\gamma_3^2=1$, the space $\mathbb{C}^4$ is the orthogonal sum of the eigenspaces of $\gamma_3$ for the eigenvalues $1$ and $-1$. For real $y$ the matrices $\boldsymbol{\alpha}_y$ and $\boldsymbol{\beta}_y$ act on the eigenspace for $\pm1$ as the multiplications by
\[
\frac{b_k\cdot(\pm iy)}{b_k\pm iy}\qquad\text{and}\qquad\frac{b_k^2\cdot b_L}{\left(b_k+b_L\pm iy\right)\cdot\left(b_k\pm iy\right)},
\]
respectively, and the norm of such a matrix is the larger of the moduli of its two eigenvalues. As $|b_k\pm iy|\ge\max\{b_k,|y|\}$, $|b_k+b_L\pm iy|\ge b_k+b_L$ and $b_k\le b$ (Section~\ref{sec:parameters-clifford-relations}), it follows that for all $y\in Y_k$
\begin{gather}
|\boldsymbol{\alpha}_y|\le b_k\le b,\qquad|\boldsymbol{\beta}_y|\le\frac{b_k^2\cdot b_L}{\left(b_k+b_L\right)\cdot b_k}\le b_L,\notag\\
\boldsymbol{\alpha}_0=0,\qquad\boldsymbol{\beta}_0=\frac{b_k\cdot b_L}{b_k+b_L}=\frac{b_{k+1}}{L};\label{eq:operators-fermion-3}
\end{gather}
for $y=\infty$ the two bounds hold by definition, and the last equality is the recursion $b_{k+1}=b\cdot b_k\cdot\left(b_k+b_L\right)^{-1}$ of Section~\ref{sec:running-quantities-tori}.

For the block torus $T_k$ these definitions reproduce the operators of Section~\ref{sec:exact-fermion-operators}: for $k\ge1$
\[
S_k(A)=S_{\infty}(A),\qquad S_{k,y}(A)=S_y(A),\qquad S^0_{k+1}(A)=S_0(A),
\]
and at level $0$ one has $S_{0,y}(A)=S_y(A)$ and $\Gamma_0(A)=S_0(A)$. The parameter $y$ is called \emph{neutral} if $F_y$ does not contain $\gamma_3$, that is, if $y\in\{0,\infty\}$ for $k\ge1$, and if $y=0$ for $k=0$.

\subsubsection{Norms with exponential weights}\label{sec:norms-exponential-weights}

Let $\gamma\ge0$ and $y\in\Lambda^0$. On spinors on $\Lambda$, $\mathbf{1}_{\Delta_y}$ denotes the multiplication by the indicator function of the orbit $\Delta_y+R\mathbb{Z}^3$ of $\Delta_y$; on functions of short pairs $(x,x')$ it denotes the multiplication by the value of this indicator function at $x$; on spinors on $\Lambda^0$ it denotes the multiplication by the indicator function of $y+R\mathbb{Z}^3$. For a linear operator $O$ from a space of spinors on $\Lambda$ or on $\Lambda^0$ to a space of spinors on $\Lambda$ or on $\Lambda^0$, or to the space of functions of short pairs, put
\[
N_{\gamma}(O):=\sup_y\sum_{y'}\exp\left(\gamma\cdot d(y,y')\right)\cdot\left\|\mathbf{1}_{\Delta_y}O\mathbf{1}_{\Delta_{y'}}\right\|,
\]
where $y$ and $y'$ range over $\Lambda^0$ modulo $R\mathbb{Z}^3$, and $\|\cdot\|$ is the norm of operators for the sup norms of Section~\ref{sec:norms-pairings-kernels}. For all such operators $O$, $O_1$, $O_2$ and all $y,y'\in\Lambda^0$ one has
\begin{align}
&\|O\|\le N_{\gamma}(O),\qquad\left\|\mathbf{1}_{\Delta_y}O\mathbf{1}_{\Delta_{y'}}\right\|\le\exp\left(-\gamma\cdot d(y,y')\right)\cdot N_{\gamma}(O),\tag{EW1}\label{eq:operators-fermion-EW1}\\
&N_{\gamma}(O_1+O_2)\le N_{\gamma}(O_1)+N_{\gamma}(O_2),\qquad N_{\gamma}(O_1O_2)\le N_{\gamma}(O_1)\cdot N_{\gamma}(O_2),\tag{EW2}\label{eq:operators-fermion-EW2}\\
&N_{\gamma}(g)=|g|_{\infty}\quad\text{if }g\text{ is a bounded scalar function acting by multiplication}.\tag{EW3}\label{eq:operators-fermion-EW3}
\end{align}
The sets $\Delta_{y'}+R\mathbb{Z}^3$, with $y'$ ranging over $\Lambda^0$ modulo $R\mathbb{Z}^3$, partition $\Lambda$; hence $1=\sum_{y'}\mathbf{1}_{\Delta_{y'}}$ on each of the spaces considered, and $\|O\|=\sup_y\left\|\mathbf{1}_{\Delta_y}O\right\|$ for the sup norms. Since $\gamma\cdot d(y,y')\ge0$,
\[
\left\|\mathbf{1}_{\Delta_y}O\right\|\le\sum_{y'}\left\|\mathbf{1}_{\Delta_y}O\mathbf{1}_{\Delta_{y'}}\right\|\le N_{\gamma}(O)
\]
for every $y$, which is the first inequality of~\eqref{eq:operators-fermion-EW1}. The second one holds because $\exp\left(\gamma\cdot d(y,y')\right)\cdot\left\|\mathbf{1}_{\Delta_y}O\mathbf{1}_{\Delta_{y'}}\right\|$ is a term of the sum that defines $N_{\gamma}(O)$. The first inequality of~\eqref{eq:operators-fermion-EW2} follows from the triangle inequality for the norm. For the second one, the relations $1=\sum_{y''}\mathbf{1}_{\Delta_{y''}}$ and $\mathbf{1}_{\Delta_{y''}}^2=\mathbf{1}_{\Delta_{y''}}$ give
\[
\mathbf{1}_{\Delta_y}O_1O_2\mathbf{1}_{\Delta_{y'}}=\sum_{y''}\left(\mathbf{1}_{\Delta_y}O_1\mathbf{1}_{\Delta_{y''}}\right)\cdot\left(\mathbf{1}_{\Delta_{y''}}O_2\mathbf{1}_{\Delta_{y'}}\right),
\]
and $d(y,y')\le d(y,y'')+d(y'',y')$; as $\gamma\ge0$, it follows that
\begin{align*}
&\sum_{y'}\exp\left(\gamma\cdot d(y,y')\right)\cdot\left\|\mathbf{1}_{\Delta_y}O_1O_2\mathbf{1}_{\Delta_{y'}}\right\|\\
&\qquad\le\sum_{y''}\exp\left(\gamma\cdot d(y,y'')\right)\cdot\left\|\mathbf{1}_{\Delta_y}O_1\mathbf{1}_{\Delta_{y''}}\right\|\\
&\qquad\qquad\cdot\sum_{y'}\exp\left(\gamma\cdot d(y'',y')\right)\cdot\left\|\mathbf{1}_{\Delta_{y''}}O_2\mathbf{1}_{\Delta_{y'}}\right\|\\
&\qquad\le N_{\gamma}(O_1)\cdot N_{\gamma}(O_2).
\end{align*}
As to~\eqref{eq:operators-fermion-EW3}, the operator $\mathbf{1}_{\Delta_y}g\mathbf{1}_{\Delta_{y'}}$ vanishes unless $y=y'$ modulo $R\mathbb{Z}^3$, and $\mathbf{1}_{\Delta_y}g\mathbf{1}_{\Delta_y}$ is the multiplication by $g$ on $\Delta_y+R\mathbb{Z}^3$, whose norm is the supremum of $|g|$ on this set.

Let $g$ be a bounded scalar function on $\Lambda$. On functions $F$ of short pairs, $g$ acts, like $\mathbf{1}_{\Delta_y}$, through its value at the first site:
\[
(gF)(x,x'):=g(x)\cdot F(x,x').
\]
For every operator $O$ as above with values in the space of functions of short pairs, the operator $gO$ is again such an operator, and
\[
N_{\gamma}(gO)\le|g|_{\infty}\cdot N_{\gamma}(O).
\]
Indeed, on functions of short pairs the multiplications by $g$ and by $\mathbf{1}_{\Delta_y}$ commute, both being multiplications by functions of $x$, and $|gF|_{\infty}\le|g|_{\infty}\cdot|F|_{\infty}$; therefore
\[
\left\|\mathbf{1}_{\Delta_y}gO\mathbf{1}_{\Delta_{y'}}\right\|=\left\|g\mathbf{1}_{\Delta_y}O\mathbf{1}_{\Delta_{y'}}\right\|\le|g|_{\infty}\cdot\left\|\mathbf{1}_{\Delta_y}O\mathbf{1}_{\Delta_{y'}}\right\|
\]
for all $y,y'\in\Lambda^0$, and the bound follows upon multiplication by $\exp\left(\gamma\cdot d(y,y')\right)$, summation over $y'$ and the supremum over $y$. For an operator $O$ with values in a space of spinors on $\Lambda$ or on $\Lambda^0$, and a bounded scalar function $g$ on the same lattice acting by multiplication, the same bound follows from~\eqref{eq:operators-fermion-EW2} and~\eqref{eq:operators-fermion-EW3}.

For $g>0$ put
\[
c(g):=\sum_{n\in\mathbb{Z}^3}\exp\left(-g\cdot|n|\right).
\]
On every block torus and for every $y\in\Lambda^0$ one has $\sum_{y'}\exp\left(-g\cdot d(y,y')\right)\le c(g)$, the sum being taken over $\Lambda^0$ modulo $R\mathbb{Z}^3$. Indeed, choose in each class $y'$ an element at minimal distance from $y$; for this element $|y'-y|=d(y,y')$, and distinct classes yield distinct differences $y'-y$, all lying in $\mathbb{Z}^3$, so that the sum is bounded by the series defining $c(g)$.

\subsubsection{The domain of complex connections}\label{sec:domain-complex-connections}

The set $\tilde{G}_k=\tilde{G}_k(T_k)$ is that of Definition~\ref{def:domains} (Section~\ref{sec:domains-bond-functions}). Thus a connection $A$ on $T_k$ lies in $\tilde{G}_k$ if and only if it can be written as $A=A_0+A_1$, where $A_1\in G_k$, that is, $A_1$ is complex and $|A_1|_{\infty}<e_k^{-1+\varepsilon}$, and where $A_0$ is real and, for each $M$-cube $\square$, satisfies
\[
A_0=A'+d\boldsymbol{\lambda}\quad\text{and}\quad|A'|_{\infty}<e_k^{-1+\varepsilon}\qquad\text{on }\square^{\sim 5}
\]
for some maps $\boldsymbol{\lambda}\colon\left\{\text{sites of }\square^{\sim 5}\right\}\to\mathbb{R}$ and $A'\colon\left\{\text{bonds of }\square^{\sim 5}\right\}\to\mathbb{R}$, which depend on $\square$. Here $\square^{\sim 5}$, the enlargement of $\square$ by five layers of $M$-cubes (Section~\ref{sec:lattice-geometry}), is a cube of side $11M$ (Definition~\ref{def:enlarged-cubes}). As $l\in\mathrm{Cl}$, one has $l\ge13L\cdot M>11M$, so that no set $\square^{\sim 5}$ wraps around the torus.

Let $A\in\tilde{G}_k$ and let $\square$ be an $M$-cube. For suitable maps $\boldsymbol{\lambda}\colon\left\{\text{sites of }\square^{\sim 1}\right\}\to\mathbb{R}$ and $B\colon\left\{\text{bonds of }\square^{\sim 1}\right\}\to\mathbb{C}$ one has
\begin{gather}
A(b)=B(b)+(d\boldsymbol{\lambda})(b)\quad\text{for every bond }b\text{ of }\square^{\sim 1},\qquad e_k\cdot|B|_{\infty}<2e_k^{\varepsilon},\notag\\
\text{and moreover}\quad\mathfrak{i}(A)<e_k^{\varepsilon}.\tag{G1}\label{eq:operators-fermion-G1}
\end{gather}
Indeed, let $A=A_0+A_1$ and $A_0=A'+d\boldsymbol{\lambda}$ on $\square^{\sim 5}$ be as in the definition, and let $B:=A'+A_1$ on the bonds of $\square^{\sim 1}$, which is contained in $\square^{\sim 5}$. Restricting $\boldsymbol{\lambda}$ to the sites of $\square^{\sim 1}$, one obtains $A(b)=B(b)+(d\boldsymbol{\lambda})(b)$ for each bond $b$ of $\square^{\sim 1}$; moreover $e_k\cdot|B|_{\infty}\le e_k\cdot|A'|_{\infty}+e_k\cdot|A_1|_{\infty}<2e_k^{\varepsilon}$. As $A_0$ is real, $\operatorname{Im}A=\operatorname{Im}A_1$, so that $\mathfrak{i}(A)\le e_k\cdot|A_1|_{\infty}<e_k^{\varepsilon}$.

Moreover,
\begin{gather}
\tilde{G}_k\ \text{is an open subset of the space of connections};\notag\\
\tilde{G}_k\ \text{is star-shaped with respect to }0\text{, hence connected, and }G_k\subset\tilde{G}_k;\notag\\
A-d\boldsymbol{\lambda},\quad-A,\quad\overline{A},\quad gA,\quad A+\mathrm{cst}\ \in\tilde{G}_k\quad\text{for every }A\in\tilde{G}_k,\tag{G2}\label{eq:operators-fermion-G2}
\end{gather}
for every real scalar $\boldsymbol{\lambda}$, every $g\in\mathbb{G}_k$ and every real constant connection $\mathrm{cst}$, that is, $\mathrm{cst}(x,x+\eta e_{\mu})=\mathrm{cst}_{\mu}$ for all sites $x$ and directions $\mu$, with $(\mathrm{cst}_0,\mathrm{cst}_1,\mathrm{cst}_2)\in\mathbb{R}^3$.

To prove~\eqref{eq:operators-fermion-G2}, let $A=A_0+A_1\in\tilde{G}_k$, and for every $M$-cube $\square$ let $A_0=A'_{\square}+d\boldsymbol{\lambda}_{\square}$ on $\square^{\sim 5}$ be as in the definition. If $V$ is a connection with $|V|_{\infty}<e_k^{-1+\varepsilon}-|A_1|_{\infty}$, then $A+V=A_0+(A_1+V)$ with $A_1+V\in G_k$; thus a small complex perturbation, added to $A_1$, does not leave $\tilde{G}_k$, which is therefore open. For $0\le t\le1$ one has $tA=tA_0+tA_1$ with $tA_1\in G_k$, and $tA_0=tA'_{\square}+d(t\boldsymbol{\lambda}_{\square})$ on $\square^{\sim 5}$ with $|tA'_{\square}|_{\infty}\le|A'_{\square}|_{\infty}$; so $\tilde{G}_k$ is star-shaped with respect to $0$, hence connected. A connection $A\in G_k$ is of the form $0+A$, so $G_k\subset\tilde{G}_k$. For a real scalar $\boldsymbol{\lambda}$ one has $A-d\boldsymbol{\lambda}=(A_0-d\boldsymbol{\lambda})+A_1$ and $A_0-d\boldsymbol{\lambda}=A'_{\square}+d(\boldsymbol{\lambda}_{\square}-\boldsymbol{\lambda})$ on $\square^{\sim 5}$. Further, $-A=(-A_0)+(-A_1)$ with $-A_0=(-A'_{\square})+d(-\boldsymbol{\lambda}_{\square})$ on $\square^{\sim 5}$, and $\overline{A}=A_0+\overline{A_1}$, because $A_0$ is real. Let $g=(S,a')\in\mathbb{G}_k$. The morphism $x\mapsto r(S)x+a'$ of $T_k$ induced by $g$ (Section~\ref{sec:tori-nested-blocks}) permutes the $M$-cubes, as $a'\in M\mathbb{Z}^3$ and $r(S)$ is diagonal with entries $\pm1$, and it maps $\square^{\sim 5}$ onto $(g\square)^{\sim 5}$. Hence $gA_0=gA'_{\square}+d(g\boldsymbol{\lambda}_{\square})$ on $(g\square)^{\sim 5}$, where $(g\boldsymbol{\lambda}_{\square})(x):=\boldsymbol{\lambda}_{\square}(g^{-1}x)$, with $|gA'_{\square}|_{\infty}=|A'_{\square}|_{\infty}$ and $|gA_1|_{\infty}=|A_1|_{\infty}$; so $gA=gA_0+gA_1\in\tilde{G}_k$. For a real constant connection $\mathrm{cst}$, the function $\boldsymbol{\lambda}'(x):=\sum_{\mu}\mathrm{cst}_{\mu}\cdot x_{\mu}$, defined on a lift of $\square^{\sim 5}$ to $\mathbb{R}^3$, satisfies $d\boldsymbol{\lambda}'=\mathrm{cst}$ on $\square^{\sim 5}$; hence $A_0+\mathrm{cst}=A'_{\square}+d(\boldsymbol{\lambda}_{\square}+\boldsymbol{\lambda}')$ on $\square^{\sim 5}$, and $A+\mathrm{cst}=(A_0+\mathrm{cst})+A_1\in\tilde{G}_k$.

Assume now condition~\eqref{eq:renormalisation-group-23}, and let $\tilde{R}_k$ and $\mathfrak{X}_k$ be the sets of Definition~\ref{def:domains}. Then
\begin{gather}
\mathfrak{X}_k\subset\tilde{G}_k,\qquad\tilde{R}_k+\mathfrak{X}_k\subset\tilde{G}_k,\qquad\text{and for every }\mathscr{T}\in\tilde{R}_k+\mathfrak{X}_k:\notag\\
\mathfrak{i}(\mathscr{T})\le e_k^{\frac{1}{4}+\varepsilon}+e_k^{2\varepsilon}\le1,\qquad e_k\cdot|d\mathscr{T}|_{\infty}\le4e_k^{\frac{1}{4}+2\varepsilon}+2e_k^{3\varepsilon}\le1.\tag{G3}\label{eq:operators-fermion-G3}
\end{gather}
Since $e_k\le e_0$ and each inequality of condition~\eqref{eq:renormalisation-group-23} is monotone in $e_0$, these inequalities hold with $e_k$ in place of $e_0$; in particular $e_k\le1$, $e_k^{\frac{1}{4}}+e_k^{\varepsilon}\le1$ and $4e_k^{\frac{1}{4}+2\varepsilon}+2e_k^{3\varepsilon}\le1$. Let $\mathscr{T}=A_0+A_1+\mathscr{C}$ with $A_0+A_1\in\tilde{R}_k$, $A_1\in R_k$, and $\mathscr{C}\in\mathfrak{X}_k$, and for an $M$-cube $\square$ let $A_0=A'+d\boldsymbol{\lambda}$ on $\square^{\dagger}$ be as in Definition~\ref{def:domains}. Then $\mathscr{T}=A_0+(A_1+\mathscr{C})$, and
\[
|A_1+\mathscr{C}|_{\infty}<e_k^{-\frac{3}{4}+\varepsilon}+e_k^{-1+2\varepsilon}\le e_k^{-1+\varepsilon},
\]
the last inequality being $e_k^{\frac{1}{4}}+e_k^{\varepsilon}\le1$. The cube $\square^{\dagger}=\square^{\sim L}$ of Definition~\ref{def:enlarged-cubes} contains $\square^{\sim 5}$, because $L\ge5$, and on it $|A'|_{\infty}<e_k^{-\frac{3}{4}+\varepsilon}\le e_k^{-1+\varepsilon}$. Hence $\mathscr{T}\in\tilde{G}_k$; since $0\in\tilde{R}_k$, the case $A_0=A_1=0$ gives $\mathfrak{X}_k\subset\tilde{G}_k$. As $A_0$ is real, $\operatorname{Im}\mathscr{T}=\operatorname{Im}(A_1+\mathscr{C})$, and therefore
\[
\mathfrak{i}(\mathscr{T})\le e_k\cdot|A_1+\mathscr{C}|_{\infty}<e_k^{\frac{1}{4}+\varepsilon}+e_k^{2\varepsilon}\le e_k^{\frac{1}{4}}+e_k^{\varepsilon}\le1.
\]
Let $p$ be a plaquette with corner $x$ in the directions $\mu,\nu$. For every connection $A$ one has $(dA)(p)=\left(\partial_{\mu}A_{\nu}-\partial_{\nu}A_{\mu}\right)(x)$, where $A_{\mu}(x):=A(x,x+\eta e_{\mu})$ and $\partial_{\nu}$ is the difference quotient of Section~\ref{sec:domains-bond-functions}; hence $|dA|\le2|\partial A|$. If $\square$ is an $M$-cube that contains the corner $x$ of $p$, then all corners of $p$ lie in $\square^{\sim 1}\subset\square^{\dagger}$; as $A_0-A'=d\boldsymbol{\lambda}$ on $\square^{\dagger}$ and $d(d\boldsymbol{\lambda})=0$, it follows that $(dA_0)(p)=(dA')(p)$. The bounds of Definition~\ref{def:domains} on the derivatives of $A'$, $A_1$ and $\mathscr{C}$ now give
\[
e_k\cdot|d\mathscr{T}|_{\infty}\le2e_k\cdot\left(e_k^{-\frac{3}{4}+2\varepsilon}+e_k^{-\frac{3}{4}+2\varepsilon}+e_k^{-1+3\varepsilon}\right)=4e_k^{\frac{1}{4}+2\varepsilon}+2e_k^{3\varepsilon}\le1.
\]
The inclusion $\tilde{R}_k+\mathfrak{X}_k\subset\tilde{G}_k$ and these two bounds are the assertions of Lemma~\ref{lem:inclusions-bounds} that concern $\tilde{R}_k+\mathfrak{X}_k$.

\subsection{Properties of the elementary operators}\label{sec:properties-elementary-operators}

In this subsection $\mathbb{T}$ denotes a block torus of level $k$ (Definition~\ref{def:block-torus}), and the letters $A$, $B$ denote connections on $\mathbb{T}$.

\begin{lemma}[Bounds for the phases]\label{lem:bounds-phases}
In the following inequalities $y\in\Lambda^0$, $y'\in\Lambda_{k+1}$, and $x$, $x'$ are sites of $\mathbb{T}$ in the ranges indicated on the right.
\begin{enumerate}
\item[(a)] One has
\begin{align*}
\left|u_A(x,x')^{\pm1}\right|&\le\exp\left(\eta\cdot\mathfrak{i}(A)\right)&&\left((x,x')\ \text{a bond}\right),\\
\left|\boldsymbol{\phi}_A(y,x)^{\pm1}\right|&\le\exp\left(\frac{3}{2}\,\mathfrak{i}(A)\right)&&\left(x\in\Delta_y\right),\\
\left|\boldsymbol{\phi}^{+}_A(y',x)^{\pm1}\right|&\le\exp\left(\frac{3}{2}\,L\cdot\mathfrak{i}(A)\right)&&\left(x\in\mathrm{Bl}_{k+1}(y')\right),\\
\left|U_A(x,x')^{\pm1}\right|&\le\exp\left(3\,|x-x'|\cdot\mathfrak{i}(A)\right)&&\left(|x-x'|<R/2\right).
\end{align*}
\item[(b)] If $\boldsymbol{\delta}$ is a number with $e_k|B|_{\infty}\le\boldsymbol{\delta}$, then
\begin{align*}
\left|u_B(x,x')-1\right|&\le\eta\cdot\boldsymbol{\delta}\cdot\exp\left(\eta\cdot\boldsymbol{\delta}\right)&&\left((x,x')\ \text{a bond}\right),\\
\left|\boldsymbol{\phi}_B(y,x)^{-1}\,\boldsymbol{\phi}_B(y,x')-1\right|&\le3\boldsymbol{\delta}\cdot\exp(3\boldsymbol{\delta})&&\left(x,x'\in\Delta_y\right),\\
\left|\boldsymbol{\phi}^{+}_B(y',x)^{-1}\,\boldsymbol{\phi}^{+}_B(y',x')-1\right|&\le3L\cdot\boldsymbol{\delta}\cdot\exp\left(3L\cdot\boldsymbol{\delta}\right)&&\left(x,x'\in\mathrm{Bl}_{k+1}(y')\right),\\
\left|U_B(x,x')-1\right|&\le3\,|x-x'|\cdot\boldsymbol{\delta}\cdot\exp\left(3\,|x-x'|\cdot\boldsymbol{\delta}\right)&&\left(|x-x'|<R/2\right).
\end{align*}
\end{enumerate}
\end{lemma}

\begin{proof}
Let $\pi$ be an order of the coordinates, and let $x$, $x'$ be sites with $|x-x'|<R/2$. The path $\Gamma^{\pi}(x,x')$ of Section~\ref{sec:phase-factors-lattice} has $|x-x'|_1/\eta$ bonds, where $|\cdot|_1$ is the $\ell^1$ norm on $\mathbb{R}^3$, and $|x-x'|_1\le 3|x-x'|$. Hence, for every connection $B$,
\[
\eta\left|(\tau B)(x,x')\right|\le 3\,|x-x'|\cdot|B|_{\infty}.
\]
Let $y\in\Lambda^0$ and $x\in\Delta_y$, and let $y_0,\dots,y_k$ be the centres of the blocks of sides $\eta,L\eta,\dots,1$ that contain $x$ (Section~\ref{sec:phase-factors-lattice}). For $j<k$ the block of side $L^j\eta$ with centre $y_j$ is one of the $L^3$ blocks of which the block of side $L^{j+1}\eta$ with centre $y_{j+1}$ is composed (Section~\ref{sec:tori-nested-blocks}); consequently $|y_{j+1}-y_j|\le\left(L^{j+1}-L^j\right)\eta/2$. The preceding bound and $\eta=L^{-k}$ give
\[
\eta\left|(\tau_kB)(y,x)\right|\le\sum_{j<k}\frac{3}{2}\left(L^{j+1}-L^j\right)\eta\,|B|_{\infty}=\frac{3}{2}\left(L^k-1\right)\eta\,|B|_{\infty}<\frac{3}{2}\,|B|_{\infty}.
\]
In the same way, for $y'\in\Lambda_{k+1}$ and $y\in B(y')$ the unit block $\Delta_y$ is one of the $L^3$ unit blocks of which $\mathrm{Bl}_{k+1}(y')$ is composed, so that $|y'-y|\le(L-1)/2$ and
\[
\eta\left|(\tau B)(y',y)\right|\le\frac{3}{2}\,(L-1)\,|B|_{\infty}.
\]
Finally, for a complex number $z$ one has $|\exp(iz)|=\exp(-\operatorname{Im}z)$ and $|\exp(z)-1|\le|z|\exp|z|$ (compare the power series term by term), and the function $t\mapsto t\exp(t)$ is increasing on $[0,\infty)$.

(a) The maps $A\mapsto\tau A$ and $A\mapsto\tau_kA$ are linear with real coefficients; hence $\operatorname{Im}(\tau A)=\tau(\operatorname{Im}A)$ and $\operatorname{Im}(\tau_kA)=\tau_k(\operatorname{Im}A)$. The definitions of Section~\ref{sec:phase-factors-lattice}, the three bounds above applied to $\operatorname{Im}A$ in place of $B$, and $\mathfrak{i}(A)=e_k|\operatorname{Im}A|_{\infty}$ give
\begin{align*}
\left|u_A(x,x')^{\pm1}\right|&=\exp\left(\mp e_k\eta\operatorname{Im}A(x,x')\right)\le\exp\left(\eta\cdot\mathfrak{i}(A)\right),\\
\left|\boldsymbol{\phi}_A(y,x)^{\pm1}\right|&=\exp\left(\mp e_k\eta\,(\tau_k\operatorname{Im}A)(y,x)\right)\le\exp\left(\frac{3}{2}\,\mathfrak{i}(A)\right),\\
\left|U_A(x,x')^{\pm1}\right|&=\exp\left(\mp e_k\eta\,(\tau\operatorname{Im}A)(x,x')\right)\le\exp\left(3\,|x-x'|\cdot\mathfrak{i}(A)\right).
\end{align*}
A site $x\in\mathrm{Bl}_{k+1}(y')$ lies in $\Delta_y$ for some $y\in B(y')$, and $\boldsymbol{\phi}^{+}_A(y',x)=U_A(y',y)\,\boldsymbol{\phi}_A(y,x)$; therefore
\[
\left|\boldsymbol{\phi}^{+}_A(y',x)^{\pm1}\right|\le\exp\left(\frac{3}{2}\,(L-1)\cdot\mathfrak{i}(A)\right)\exp\left(\frac{3}{2}\,\mathfrak{i}(A)\right)=\exp\left(\frac{3}{2}\,L\cdot\mathfrak{i}(A)\right).
\]

(b) Each of the four quantities has the form $|\exp(z)-1|$. In the first, $z=ie_k\eta B(x,x')$, and $|z|\le\eta\cdot\boldsymbol{\delta}$. In the second, $z=ie_k\eta\left((\tau_kB)(y,x')-(\tau_kB)(y,x)\right)$, and $|z|\le2\cdot\frac{3}{2}\,\boldsymbol{\delta}=3\boldsymbol{\delta}$. In the third, let $x\in\Delta_y$ and $x'\in\Delta_{\tilde{y}}$ with $y,\tilde{y}\in B(y')$; then
\[
z=ie_k\eta\left((\tau B)(y',\tilde{y})+(\tau_kB)(\tilde{y},x')-(\tau B)(y',y)-(\tau_kB)(y,x)\right),
\]
and $|z|\le2\cdot\frac{3}{2}\,(L-1)\,\boldsymbol{\delta}+2\cdot\frac{3}{2}\,\boldsymbol{\delta}=3L\cdot\boldsymbol{\delta}$. In the fourth, $z=ie_k\eta\,(\tau B)(x,x')$, and $|z|\le3\,|x-x'|\cdot\boldsymbol{\delta}$. In each case the inequality $|\exp(z)-1|\le|z|\exp|z|$ and the monotonicity of $t\mapsto t\exp(t)$ give the assertion.
\end{proof}

\begin{lemma}[Block averages and block projections]\label{lem:block-averages-block}
Let $A$ be a connection on $\mathbb{T}$. Then
\[
Q_k(A)\,Q_k^T(-A)=1,\qquad Q(A)\,Q^T(-A)=1;
\]
the operators $P_k(A)$, $P(A)$ and $P_{k+1}(A)$ are projections, that is, each of them is equal to its square, and
\[
P_k(A)\,P_{k+1}(A)=P_{k+1}(A)\,P_k(A)=P_{k+1}(A).
\]
The operators $Q_k(A)$, $Q_k^T(-A)$, $Q(A)$, $Q^T(-A)$, $P_k(A)$, $P(A)$ and $P_{k+1}(A)$ commute with the multiplication by every constant complex $4\times4$ matrix, and the trace of $P(A)$ on the space of the spinors on the unit lattice is
\[
\operatorname{Tr}P(A)=4\,(R/L)^3.
\]
The kernel of $P_k(A)$ vanishes at every pair of sites that do not lie in one unit block, and the kernels of $P(A)$ and of $P_{k+1}(A)$ vanish at every pair of sites that do not lie in one block of side $L$. The row of $P_k(A)$ at a site $x$ depends on $A$ only through the values of $A$ on the bonds that join two sites of the unit block containing $x$; the row of $P(A)$ at a unit site $y$, and the row of $P_{k+1}(A)$ at a site $x$, depend on $A$ only through the values of $A$ on the bonds that join two sites of the block of side $L$ containing $y$, respectively $x$.
\end{lemma}

\begin{proof}
Let $v$ be a spinor on the unit lattice and $w$ a spinor on the $L$-lattice, and let $y\in\Lambda^0$ and $y'\in\Lambda_{k+1}$. By the formulas of Section~\ref{sec:phase-factors-lattice},
\begin{align*}
\left(Q_k(A)\,Q_k^T(-A)\,v\right)(y)&=\eta^3\sum_{x\in\Delta_y}\boldsymbol{\phi}_A(y,x)\,\boldsymbol{\phi}_A(y,x)^{-1}\,v(y)=\eta^3\,\#|\Delta_y|\,v(y)=v(y),\\
\left(Q(A)\,Q^T(-A)\,w\right)(y')&=L^{-3}\sum_{\tilde{y}\in B(y')}U_A(y',\tilde{y})\,U_A(y',\tilde{y})^{-1}\,w(y')=L^{-3}\,\#|B(y')|\,w(y')=w(y'),
\end{align*}
because $\#|\Delta_y|=L^{3k}=\eta^{-3}$ and $\#|B(y')|=L^3$ (Section~\ref{sec:tori-nested-blocks}). Inserting these two identities one obtains
\begin{align*}
P_k(A)^2&=Q_k^T(-A)\left(Q_k(A)\,Q_k^T(-A)\right)Q_k(A)=P_k(A),\\
P(A)^2&=Q^T(-A)\left(Q(A)\,Q^T(-A)\right)Q(A)=P(A),\\
P_{k+1}(A)^2&=Q_k^T(-A)\,P(A)\left(Q_k(A)\,Q_k^T(-A)\right)P(A)\,Q_k(A)=Q_k^T(-A)\,P(A)^2\,Q_k(A)=P_{k+1}(A),\\
P_k(A)\,P_{k+1}(A)&=Q_k^T(-A)\left(Q_k(A)\,Q_k^T(-A)\right)P(A)\,Q_k(A)=P_{k+1}(A),\\
P_{k+1}(A)\,P_k(A)&=Q_k^T(-A)\,P(A)\left(Q_k(A)\,Q_k^T(-A)\right)Q_k(A)=P_{k+1}(A).
\end{align*}
The kernels of $Q_k(A)$, $Q_k^T(-A)$, $Q(A)$ and $Q^T(-A)$ are scalar functions times the unit matrix of $\mathbb{C}^4$. Hence these four operators, and their products $P_k(A)$, $P(A)$, $P_{k+1}(A)$, commute with the constant $4\times4$ matrices. The trace of the product of two linear maps between finite-dimensional spaces does not depend on the order of the factors; therefore
\[
\operatorname{Tr}P(A)=\operatorname{Tr}\left(Q^T(-A)\,Q(A)\right)=\operatorname{Tr}\left(Q(A)\,Q^T(-A)\right)
\]
is the dimension of the space of the spinors on the $L$-lattice, namely $4$ times the number $(R/L)^3$ of the sites of $\Lambda_{k+1}=t_{k+1}+L\mathbb{Z}^3$ in a period cell.

Let $f$ be a spinor on $\Lambda$. The formulas of Section~\ref{sec:phase-factors-lattice} give, for $y\in\Lambda^0$ and $x\in\Delta_y$, for $y'\in\Lambda_{k+1}$ and $y\in B(y')$, and for $y'\in\Lambda_{k+1}$ and $x\in\mathrm{Bl}_{k+1}(y')$, respectively,
\begin{align*}
\left(P_k(A)f\right)(x)&=\boldsymbol{\phi}_A(y,x)^{-1}\,\eta^3\sum_{x'\in\Delta_y}\boldsymbol{\phi}_A(y,x')\,f(x'),\\
\left(P(A)v\right)(y)&=U_A(y',y)^{-1}\,L^{-3}\sum_{\tilde{y}\in B(y')}U_A(y',\tilde{y})\,v(\tilde{y}),\\
\left(P_{k+1}(A)f\right)(x)&=\boldsymbol{\phi}^{+}_A(y',x)^{-1}\,L^{-3}\eta^3\sum_{x'\in\mathrm{Bl}_{k+1}(y')}\boldsymbol{\phi}^{+}_A(y',x')\,f(x').
\end{align*}
The assertions on the kernels can be read off. The phase $\boldsymbol{\phi}_A(y,z)$, $z\in\Delta_y$, is formed with the sums of $A$ along the paths $\Gamma^{\pi}$ between the centres of two consecutive blocks containing $z$; both centres lie in the larger of these two blocks, which is contained in $\Delta_y$, since the blocks are nested. The phase $U_A(y',\tilde{y})$, $\tilde{y}\in B(y')$, is formed with the paths between $y'$ and $\tilde{y}$, two sites of $\mathrm{Bl}_{k+1}(y')$, and $\boldsymbol{\phi}^{+}_A(y',z)=U_A(y',\tilde{y})\,\boldsymbol{\phi}_A(\tilde{y},z)$ for $z\in\Delta_{\tilde{y}}\subset\mathrm{Bl}_{k+1}(y')$. All sites of a rectilinear path lie in the smallest box that contains its two end points, and a block contains this box as soon as it contains the two end points. Therefore the row of $P_k(A)$ at $x\in\Delta_y$ contains $A$ only through its values on bonds that join two sites of $\Delta_y$, and the rows of $P(A)$ at $y\in B(y')$ and of $P_{k+1}(A)$ at $x\in\mathrm{Bl}_{k+1}(y')$ contain $A$ only through its values on bonds that join two sites of $\mathrm{Bl}_{k+1}(y')$.
\end{proof}

\begin{lemma}[Perturbations, commutators and H\"older differences]\label{lem:perturbations-commutators-holder}
Let $\gamma\ge0$ satisfy $\gamma L\le\log 2$, put
\[
C_V:=12\mathbf{e}\cdot(2+b),
\]
and let $y\in Y_k$.
\begin{enumerate}
\item[(a)] If $\boldsymbol{\delta}$ is a number with $e_k|B|_{\infty}\le\boldsymbol{\delta}$ and $3L\cdot\boldsymbol{\delta}\le1$, then
\[
N_{\gamma}\left(F_y(B)-F_y(0)\right)\le C_V\cdot\boldsymbol{\delta}.
\]
\item[(b)] Let $3L\cdot\mathfrak{i}(A)\le1$, let $L_h$ be a number, and let $h$ be a real scalar on $\mathbb{T}$ such that
\[
|h(x)-h(x')|\le L_h\cdot\sum_{\mu}|x_{\mu}-x'_{\mu}|\qquad\text{for all sites }x,x'\text{ with }|x-x'|\le L.
\]
Then
\[
N_{\gamma}\left(\left[F_y(A),h\right]\right)\le C_V\cdot L_h.
\]
The kernel of $\left[F_y(A),h\right]$ vanishes at $(x,x')$ unless $h(x)\ne h(x')$ and, in addition, either $x$, $x'$ are nearest neighbours or both belong to one block of side $L$.
\item[(c)] For a spinor $f$ on $\Lambda$, a real scalar $h$ on $\mathbb{T}$ and a short pair $(x,x')$,
\begin{align*}
\left(\delta_{\alpha,A}(hf)\right)(x,x')&=h(x)\left(\delta_{\alpha,A}f\right)(x,x')+|x-x'|^{-\alpha}\left(h(x')-h(x)\right)U_A(x,x')\,f(x'),\\
\left(\delta_{\alpha,B}f-\delta_{\alpha,0}f\right)(x,x')&=|x-x'|^{-\alpha}\left(U_B(x,x')-1\right)f(x').
\end{align*}
Let $O$ be a linear operator from the spinors on $\Lambda$ or on $\Lambda^0$ to the spinors on $\Lambda$ with $N_{\gamma}(O)<\infty$, let $g$ be a function on the short pairs and $g_0$ a number with $|g|\le g_0$, and let $E_gO$ be the operator from the domain of $O$ to the functions of short pairs given by
\[
\left(E_gOf\right)(x,x'):=g(x,x')\,(Of)(x').
\]
Then
\[
N_{\gamma}(E_gO)\le 27\exp(\gamma)\,g_0\,N_{\gamma}(O).
\]
\end{enumerate}
\end{lemma}

\begin{proof}
(a) By~\eqref{eq:operators-fermion-2}, $F_y(A)=D_A+\mathfrak{m}+\operatorname{Pen}_y(A)$, where $\operatorname{Pen}_y(A)=\boldsymbol{\alpha}_y\cdot P_k(A)+\boldsymbol{\beta}_y\cdot P_{k+1}(A)$ for $k\ge1$ and $\operatorname{Pen}_y(A)=b_L\cdot P(A)+i\gamma_3y$ for $k=0$ (Section~\ref{sec:operators-2}). In $V:=F_y(B)-F_y(0)$ the mass $\mathfrak{m}$ and, for $k=0$, the term $i\gamma_3y$ cancel, and
\[
V=\begin{cases}\left(D_B-D_0\right)+\boldsymbol{\alpha}_y\cdot\left(P_k(B)-P_k(0)\right)+\boldsymbol{\beta}_y\cdot\left(P_{k+1}(B)-P_{k+1}(0)\right), & k\ge1,\\ \left(D_B-D_0\right)+b_L\cdot\left(P(B)-P(0)\right), & k=0.\end{cases}
\]
Fix $y_1\in\Lambda^0$. By the definition of $N_{\gamma}$ in Section~\ref{sec:norms-exponential-weights} and the triangle inequality, it suffices to bound the sum $\sum_{y_2}\exp\left(\gamma d(y_1,y_2)\right)\left\|\mathbf{1}_{\Delta_{y_1}}(\cdot)\,\mathbf{1}_{\Delta_{y_2}}\right\|$ for each of the three parts of $V$, uniformly in $y_1$, and to add the bounds.

\emph{The hopping part.} With $x_{\pm}=x\pm\eta e_{\mu}$ (Section~\ref{sec:phase-factors-lattice}), the definition of $D_A$ gives
\[
\left((D_B-D_0)f\right)(x)=-\eta^{-1}\sum_{\mu}\left\{\Pi^{-}_{\mu}\left(u_B(x,x_{+})-1\right)f(x_{+})+\Pi^{+}_{\mu}\left(u_B(x,x_{-})-1\right)f(x_{-})\right\},
\]
a sum of six terms. The matrices $\Pi^{\pm}_{\mu}$ are orthogonal projections (Section~\ref{sec:parameters-clifford-relations}), and $\eta\le1$; so, by Lemma~\ref{lem:bounds-phases}(b), each term has norm at most $\eta^{-1}\cdot\eta\cdot\boldsymbol{\delta}\cdot\exp(\eta\cdot\boldsymbol{\delta})\le\boldsymbol{\delta}\cdot\exp(\boldsymbol{\delta})$. Let $x\in\Delta_{y_1}$. A neighbour $x\pm\eta e_{\mu}$ of $x$ belongs either to $\Delta_{y_1}$ or to one of the six unit blocks $\Delta_{y_1\pm e_{\nu}}$ that share a face with $\Delta_{y_1}$; it belongs to $\Delta_{y_1+e_{\nu}}$ only if it is $x+\eta e_{\nu}$, and to $\Delta_{y_1-e_{\nu}}$ only if it is $x-\eta e_{\nu}$. Hence $\mathbf{1}_{\Delta_{y_1}}(D_B-D_0)\,\mathbf{1}_{\Delta_{y_2}}$ vanishes unless $y_2=y_1$ or $y_2=y_1\pm e_{\nu}$ for some $\nu$; in the second case only one of the six terms contributes, and $d(y_1,y_2)=1$. It follows that
\[
\sum_{y_2}\exp\left(\gamma d(y_1,y_2)\right)\left\|\mathbf{1}_{\Delta_{y_1}}(D_B-D_0)\,\mathbf{1}_{\Delta_{y_2}}\right\|\le\left(6+6\exp(\gamma)\right)\boldsymbol{\delta}\cdot\exp(\boldsymbol{\delta}).
\]

\emph{The part with $P_k$} ($k\ge1$). By Lemma~\ref{lem:block-averages-block}, $P_k(B)-P_k(0)$ couples only sites of one unit block, so that only $y_2=y_1$ contributes. By the formulas of Section~\ref{sec:phase-factors-lattice} and $\boldsymbol{\phi}_0=1$,
\[
\left((P_k(B)-P_k(0))f\right)(x)=\eta^3\sum_{x'\in\Delta_{y_1}}\left(\boldsymbol{\phi}_B(y_1,x)^{-1}\,\boldsymbol{\phi}_B(y_1,x')-1\right)f(x')\qquad\left(x\in\Delta_{y_1}\right),
\]
and Lemma~\ref{lem:bounds-phases}(b), together with $\eta^3\,\#|\Delta_{y_1}|=1$, bounds the norm of this operator by $3\boldsymbol{\delta}\cdot\exp(3\boldsymbol{\delta})$. By~\eqref{eq:operators-fermion-3}, $|\boldsymbol{\alpha}_y|\le b$.

\emph{The part with $P_{k+1}$, and for $k=0$ the part with $P=P_1$} (Section~\ref{sec:phase-factors-lattice}). By Lemma~\ref{lem:block-averages-block}, the operator $\mathbf{1}_{\Delta_{y_1}}\left(P_{k+1}(B)-P_{k+1}(0)\right)\mathbf{1}_{\Delta_{y_2}}$ vanishes unless $y_1$ and $y_2$ lie in one set $B(y')$, $y'\in\Lambda_{k+1}$. In that case the formula for $P_{k+1}$ in Section~\ref{sec:phase-factors-lattice} and $\boldsymbol{\phi}^{+}_0=1$ give
\[
\left(\mathbf{1}_{\Delta_{y_1}}\left(P_{k+1}(B)-P_{k+1}(0)\right)\mathbf{1}_{\Delta_{y_2}}f\right)(x)=L^{-3}\eta^3\sum_{x'\in\Delta_{y_2}}\left(\boldsymbol{\phi}^{+}_B(y',x)^{-1}\,\boldsymbol{\phi}^{+}_B(y',x')-1\right)f(x')\qquad\left(x\in\Delta_{y_1}\right),
\]
and Lemma~\ref{lem:bounds-phases}(b) bounds its norm by $L^{-3}\cdot3L\cdot\boldsymbol{\delta}\cdot\exp(3L\cdot\boldsymbol{\delta})$. For a given $y_1$ there are $L^3$ such $y_2$, and each of them satisfies $d(y_1,y_2)\le L-1$. The coefficient of this part is $\boldsymbol{\beta}_y$, with $|\boldsymbol{\beta}_y|\le b_L$ by~\eqref{eq:operators-fermion-3}, for $k\ge1$, and $b_L$ for $k=0$.

\emph{Summation.} The numbers $1$, $\exp(\gamma)$ and $\exp\left(\gamma(L-1)\right)$ are at most $\exp(\gamma L)$, the numbers $\exp(\boldsymbol{\delta})$ and $\exp(3\boldsymbol{\delta})$ are at most $\exp(3L\cdot\boldsymbol{\delta})$, and $L\cdot b_L=b$. Hence
\begin{align*}
N_{\gamma}(V)&\le\left(6+6\exp(\gamma)\right)\boldsymbol{\delta}\cdot\exp(\boldsymbol{\delta})+3b\cdot\boldsymbol{\delta}\cdot\exp(3\boldsymbol{\delta})+L^3\exp\left(\gamma(L-1)\right)b_L\cdot L^{-3}\cdot3L\cdot\boldsymbol{\delta}\cdot\exp(3L\cdot\boldsymbol{\delta})\\
&\le\exp(\gamma L)\exp(3L\cdot\boldsymbol{\delta})\left(12+3b+3L\cdot b_L\right)\boldsymbol{\delta}=\exp(\gamma L)\exp(3L\cdot\boldsymbol{\delta})\,(12+6b)\,\boldsymbol{\delta}\\
&\le2\mathbf{e}\cdot(12+6b)\,\boldsymbol{\delta}=C_V\cdot\boldsymbol{\delta},
\end{align*}
where the last inequality uses $\exp(\gamma L)\le2$ and $\exp(3L\cdot\boldsymbol{\delta})\le\mathbf{e}$. For $k=0$ the second term on the right of the first line is absent.

(b) The mass $\mathfrak{m}$, the matrix $i\gamma_3y$ and the constant matrices $\boldsymbol{\alpha}_y$, $\boldsymbol{\beta}_y$ commute with the multiplication by $h$. Hence
\[
\left[F_y(A),h\right]=\begin{cases}\left[D_A,h\right]+\boldsymbol{\alpha}_y\cdot\left[P_k(A),h\right]+\boldsymbol{\beta}_y\cdot\left[P_{k+1}(A),h\right], & k\ge1,\\ \left[D_A,h\right]+b_L\cdot\left[P(A),h\right], & k=0.\end{cases}
\]
By the definition of $D_A$,
\[
\left(\left[D_A,h\right]f\right)(x)=-\eta^{-1}\sum_{\mu}\left\{\Pi^{-}_{\mu}u_A(x,x_{+})\left(h(x_{+})-h(x)\right)f(x_{+})+\Pi^{+}_{\mu}u_A(x,x_{-})\left(h(x_{-})-h(x)\right)f(x_{-})\right\}.
\]
Since $|x_{\pm}-x|=\eta\le L$ and $\sum_{\nu}|(x_{\pm})_{\nu}-x_{\nu}|=\eta$, the hypothesis on $h$ gives $|h(x_{\pm})-h(x)|\le L_h\cdot\eta$; with Lemma~\ref{lem:bounds-phases}(a), each of the six terms has norm at most $L_h\cdot\exp\left(\eta\cdot\mathfrak{i}(A)\right)$. For $x\in\Delta_{y_1}$,
\[
\left(\left[P_k(A),h\right]f\right)(x)=\boldsymbol{\phi}_A(y_1,x)^{-1}\,\eta^3\sum_{x'\in\Delta_{y_1}}\boldsymbol{\phi}_A(y_1,x')\left(h(x')-h(x)\right)f(x').
\]
Two sites $x$, $x'$ of one unit block satisfy $|x-x'|<1\le L$ and $\sum_{\nu}|x_{\nu}-x'_{\nu}|<3$, so that $|h(x')-h(x)|\le3L_h$; by Lemma~\ref{lem:bounds-phases}(a) the norm of $\left[P_k(A),h\right]$ is at most $3L_h\cdot\exp\left(3\mathfrak{i}(A)\right)$, and this operator couples only sites of one unit block. In the same way, $\mathbf{1}_{\Delta_{y_1}}\left[P_{k+1}(A),h\right]\mathbf{1}_{\Delta_{y_2}}$ vanishes unless $y_1$ and $y_2$ lie in one set $B(y')$, and in that case
\[
\left(\mathbf{1}_{\Delta_{y_1}}\left[P_{k+1}(A),h\right]\mathbf{1}_{\Delta_{y_2}}f\right)(x)=\boldsymbol{\phi}^{+}_A(y',x)^{-1}\,L^{-3}\eta^3\sum_{x'\in\Delta_{y_2}}\boldsymbol{\phi}^{+}_A(y',x')\left(h(x')-h(x)\right)f(x')\qquad\left(x\in\Delta_{y_1}\right).
\]
Two sites $x$, $x'$ of $\mathrm{Bl}_{k+1}(y')$ satisfy $|x-x'|<L$ and $\sum_{\nu}|x_{\nu}-x'_{\nu}|<3L$, so that $|h(x')-h(x)|\le3L\cdot L_h$; by Lemma~\ref{lem:bounds-phases}(a) the norm of this operator is at most $L^{-3}\cdot3L\cdot L_h\cdot\exp\left(3L\cdot\mathfrak{i}(A)\right)$. For $k=0$ the same holds for $P(A)=P_1(A)$. The sum over $y_2$ is formed as in (a); since $\exp\left(\eta\cdot\mathfrak{i}(A)\right)$ and $\exp\left(3\mathfrak{i}(A)\right)$ are at most $\exp\left(3L\cdot\mathfrak{i}(A)\right)\le\mathbf{e}$,
\begin{align*}
N_{\gamma}\left(\left[F_y(A),h\right]\right)&\le\left(6+6\exp(\gamma)\right)L_h\cdot\exp\left(\eta\cdot\mathfrak{i}(A)\right)+3b\cdot L_h\cdot\exp\left(3\mathfrak{i}(A)\right)\\
&\qquad+L^3\exp\left(\gamma(L-1)\right)b_L\cdot L^{-3}\cdot3L\cdot L_h\cdot\exp\left(3L\cdot\mathfrak{i}(A)\right)\\
&\le\exp(\gamma L)\exp\left(3L\cdot\mathfrak{i}(A)\right)(12+6b)\,L_h\le2\mathbf{e}\cdot(12+6b)\,L_h=C_V\cdot L_h.
\end{align*}
The kernel of $\left[F_y(A),h\right]$ at $(x,x')$ is $h(x')-h(x)$ times the kernel of $F_y(A)$ at $(x,x')$. By the definition of $D_A$ and by Lemma~\ref{lem:block-averages-block}, the kernel of $F_y(A)$ vanishes at $(x,x')$ unless $x=x'$, or $x$, $x'$ are nearest neighbours, or both belong to one block of side $L$ (every unit block is contained in a block of side $L$). This proves the assertion on the kernel.

(c) By the definition of $\delta_{\alpha,A}$ in Section~\ref{sec:phase-factors-lattice},
\[
\left(\delta_{\alpha,A}(hf)\right)(x,x')=|x-x'|^{-\alpha}\left(U_A(x,x')\,h(x')\,f(x')-h(x)\,f(x)\right);
\]
adding and subtracting $h(x)\,U_A(x,x')\,f(x')$ inside the parentheses gives the first identity. The second follows from the same definition and $U_0=1$.

Let $y_1,y_2\in\Lambda^0$. If $(x,x')$ is a short pair with $x\in\Delta_{y_1}$, then $x'\in\Delta_{\tilde{y}}$ for a unit site $\tilde{y}$ with
\[
|\tilde{y}-y_1|\le|\tilde{y}-x'|+|x'-x|+|x-y_1|<\frac{1}{2}+1+\frac{1}{2}=2;
\]
as $\tilde{y}-y_1$ has integer coordinates, $|\tilde{y}-y_1|\le1$. There are $27$ unit sites $\tilde{y}$ with $|\tilde{y}-y_1|\le1$. Hence, for every spinor $f$ in the domain of $O$,
\[
\left|\left(\mathbf{1}_{\Delta_{y_1}}E_gO\,\mathbf{1}_{\Delta_{y_2}}f\right)(x,x')\right|\le g_0\left|\left(O\,\mathbf{1}_{\Delta_{y_2}}f\right)(x')\right|\le g_0\sum_{|\tilde{y}-y_1|\le1}\left\|\mathbf{1}_{\Delta_{\tilde{y}}}O\,\mathbf{1}_{\Delta_{y_2}}\right\|\,|f|_{\infty},
\]
that is,
\[
\left\|\mathbf{1}_{\Delta_{y_1}}E_gO\,\mathbf{1}_{\Delta_{y_2}}\right\|\le g_0\sum_{|\tilde{y}-y_1|\le1}\left\|\mathbf{1}_{\Delta_{\tilde{y}}}O\,\mathbf{1}_{\Delta_{y_2}}\right\|.
\]
For $|\tilde{y}-y_1|\le1$ one has $d(y_1,y_2)\le1+d(\tilde{y},y_2)$, hence $\exp\left(\gamma d(y_1,y_2)\right)\le\exp(\gamma)\exp\left(\gamma d(\tilde{y},y_2)\right)$, and therefore
\[
\sum_{y_2}\exp\left(\gamma d(y_1,y_2)\right)\left\|\mathbf{1}_{\Delta_{y_1}}E_gO\,\mathbf{1}_{\Delta_{y_2}}\right\|\le g_0\exp(\gamma)\sum_{|\tilde{y}-y_1|\le1}\ \sum_{y_2}\exp\left(\gamma d(\tilde{y},y_2)\right)\left\|\mathbf{1}_{\Delta_{\tilde{y}}}O\,\mathbf{1}_{\Delta_{y_2}}\right\|\le27\exp(\gamma)\,g_0\,N_{\gamma}(O).
\]
The supremum over $y_1$ gives the assertion.
\end{proof}

\begin{lemma}[Symmetries of the elementary operators]\label{lem:symmetries-elementary-operators}
Let $y\in Y_k$, and put $-\infty:=\infty$. Let $\Lambda$, $\Lambda^0$ and $\Lambda_{k+1}$ be the lattice, the unit lattice and the $L$-lattice of $\mathbb{T}$.
\begin{enumerate}
\item[(a)] \emph{Gauge transformations.} Let $\boldsymbol{\omega}$ be a real connection on $\mathbb{T}$ with $d\boldsymbol{\omega}=0$, and let $u\colon\Lambda\to\mathbb{C}$ be a function, invariant under the translations by $R\mathbb{Z}^3$, with $|u|=1$ and
\[
u(x')=u(x)\cdot\exp\left(i\,e_k\cdot\eta\cdot\boldsymbol{\omega}(x,x')\right)\qquad\text{for every bond }(x,x')\text{ of }\Lambda.
\]
Let $u^0$ and $u^1$ be the restrictions of $u$ to $\Lambda^0$ and to $\Lambda_{k+1}$; the letters $u$, $u^0$, $u^1$ also denote the operators of multiplication by these functions on the spinors on $\Lambda$, on $\Lambda^0$ and on $\Lambda_{k+1}$. (For instance, $\boldsymbol{\omega}=d\boldsymbol{\lambda}$ and $u=\exp\left(i\,e_k\cdot\boldsymbol{\lambda}\right)$, with a real scalar $\boldsymbol{\lambda}$, have these properties.) Then
\begin{gather*}
D_{A-\boldsymbol{\omega}}=u\cdot D_A\cdot u^{-1},\\
Q_k(A-\boldsymbol{\omega})=u^0\cdot Q_k(A)\cdot u^{-1},\qquad Q_k^T\left(-(A-\boldsymbol{\omega})\right)=u\cdot Q_k^T(-A)\cdot\left(u^0\right)^{-1},\\
Q(A-\boldsymbol{\omega})=u^1\cdot Q(A)\cdot\left(u^0\right)^{-1},\qquad Q^T\left(-(A-\boldsymbol{\omega})\right)=u^0\cdot Q^T(-A)\cdot\left(u^1\right)^{-1},\\
P(A-\boldsymbol{\omega})=u^0\cdot P(A)\cdot\left(u^0\right)^{-1},\\
O(A-\boldsymbol{\omega})=u\cdot O(A)\cdot u^{-1}\qquad\text{for }O=P_k,\ P_{k+1},\ \operatorname{Pen}_y,\ F_y,
\end{gather*}
and, for every spinor $f$ on $\Lambda$ and all sites $x\ne x'$ with $|x-x'|<R/2$,
\[
\left(\delta_{\alpha,A-\boldsymbol{\omega}}(u\cdot f)\right)(x,x')=u(x)\cdot\left(\delta_{\alpha,A}f\right)(x,x').
\]
\item[(b)] \emph{Morphisms of block tori.} Let $\boldsymbol{\chi}\colon\mathbb{T}\to\mathbb{T}'$ be a morphism of block tori of level $k$ (Definition~\ref{def:block-torus}) with linear part $r$, let $\Lambda'_j$ be the lattices of $\mathbb{T}'$, and let $S\in\mathrm{Sp}$ satisfy $r(S)=r$. For a spinor $f$ on $\mathbb{T}$, on its unit lattice or on its $L$-lattice, let $gf$ be the spinor on $\mathbb{T}'$, on its unit lattice or on its $L$-lattice, respectively, defined by
\[
(gf)(x):=S\cdot f\left(\boldsymbol{\chi}^{-1}x\right),
\]
and let $gA$ be the connection on $\mathbb{T}'$ defined by $(gA)(x,x'):=A\left(\boldsymbol{\chi}^{-1}x,\boldsymbol{\chi}^{-1}x'\right)$. Then, the operators on the left-hand sides being those of $\mathbb{T}'$,
\begin{gather*}
D_{gA}\cdot g=g\cdot D_A,\qquad Q_k(gA)\cdot g=g\cdot Q_k(A),\qquad Q_k^T(-gA)\cdot g=g\cdot Q_k^T(-A),\\
Q(gA)\cdot g=g\cdot Q(A),\qquad Q^T(-gA)\cdot g=g\cdot Q^T(-A),\\
P_k(gA)\cdot g=g\cdot P_k(A),\qquad P(gA)\cdot g=g\cdot P(A),\qquad P_{k+1}(gA)\cdot g=g\cdot P_{k+1}(A);
\end{gather*}
for every spinor $f$ on $\Lambda$ and all sites $x\ne x'$ with $|x-x'|<R/2$
\[
\left(\delta_{\alpha,gA}\,gf\right)(\boldsymbol{\chi} x,\boldsymbol{\chi} x')=S\cdot\left(\delta_{\alpha,A}f\right)(x,x');
\]
and
\[
F_{y'}(gA)\cdot g=g\cdot F_y(A),
\]
where
\[
y':=\begin{cases}y&\text{if }S\in\mathrm{Sp}_3\text{ or if }y\text{ is neutral},\\-y&\text{if }S\notin\mathrm{Sp}_3.\end{cases}
\]
Where both lines apply, they give the same value, since a neutral $y$ (Section~\ref{sec:operators-2}) is $0$ or $\infty$.
\item[(c)] \emph{Transposition.} For the transposes with respect to the pairings of Section~\ref{sec:norms-pairings-kernels},
\begin{gather*}
D_A^T=\mathsf{C}^{-1}\cdot D_{-A}\cdot\mathsf{C},\qquad Q_k(A)^T=Q_k^T(A),\qquad Q(A)^T=Q^T(A),\\
P_k(A)^T=P_k(-A),\qquad P(A)^T=P(-A),\qquad P_{k+1}(A)^T=P_{k+1}(-A),\\
F_y(A)^T=\mathsf{C}^{-1}\cdot F_{-y}(-A)\cdot\mathsf{C}.
\end{gather*}
\item[(d)] \emph{Complex conjugation.} For a linear operator $O$ between spaces of spinors let $\overline{O}$ be the operator whose kernel is the complex conjugate of the kernel of $O$, that is, $\overline{O}f:=\overline{O\overline{f}}$, and let $\overline{A}$ be the complex conjugate of the connection $A$. Then
\[
\overline{O(\overline{A})}=\varkappa^{-1}\cdot O(-A)\cdot\varkappa
\]
for each of the operator-valued functions $O$ of the connection given by $O(A)=D_A$, $Q_k(A)$, $Q_k^T(-A)$, $Q(A)$, $Q^T(-A)$, $P_k(A)$, $P(A)$, $P_{k+1}(A)$ and $F_y(A)$, the same $y$ standing on both sides.
\item[(e)] \emph{Adjoints.} Let $A$ be real, and let the adjoints be taken with respect to the sesquilinear pairings of Section~\ref{sec:norms-pairings-kernels}. Then
\[
D_A^{*}=\gamma_3\cdot D_A\cdot\gamma_3,\qquad Q_k^T(-A)=Q_k(A)^{*},\qquad Q^T(-A)=Q(A)^{*},
\]
and $P_k(A)$, $P(A)$, $P_{k+1}(A)$ are orthogonal projections.
\end{enumerate}
\end{lemma}

\begin{proof}
The symbols $x_{\pm}$ and $x'_{\pm}$ are those of Section~\ref{sec:phase-factors-lattice}, and sums over $\mu$ extend over the directions $\mu=0,1,2$.

\emph{(a)} Let $(x,x')$ be a bond. As $u_{A-\boldsymbol{\omega}}(x,x')=u_A(x,x')\cdot\exp\left(-i\,e_k\cdot\eta\cdot\boldsymbol{\omega}(x,x')\right)$, the hypothesis on $u$ gives
\[
u_{A-\boldsymbol{\omega}}(x,x')\cdot u(x')=u_A(x,x')\cdot u(x).
\]
Inserting this into the formula for $D_{A-\boldsymbol{\omega}}$ of Section~\ref{sec:phase-factors-lattice}, one obtains for every spinor $f$ on $\Lambda$ and every $x\in\Lambda$
\begin{align*}
\left(D_{A-\boldsymbol{\omega}}(u\cdot f)\right)(x)
&=-\eta^{-1}\sum_{\mu}\Big\{\Pi^{-}_{\mu}\cdot u_{A-\boldsymbol{\omega}}(x,x_{+})\cdot u(x_{+})\cdot f(x_{+})\\
&\qquad\qquad+\Pi^{+}_{\mu}\cdot u_{A-\boldsymbol{\omega}}(x,x_{-})\cdot u(x_{-})\cdot f(x_{-})-u(x)\cdot f(x)\Big\}\\
&=u(x)\cdot\left(D_Af\right)(x),
\end{align*}
that is, $D_{A-\boldsymbol{\omega}}\cdot u=u\cdot D_A$.

Next let $x$, $x'$ be sites with $|x-x'|<R/2$, taken on the shortest lift in $t+\eta\mathbb{Z}^3\subset\mathbb{R}^3$, as for the paths $\Gamma^{\pi}(x,x')$ of Section~\ref{sec:phase-factors-lattice}; here $t$ is the shift of $\mathbb{T}$ in Definition~\ref{def:block-torus}. Let $\mathbb{W}$ be the set of the points of $t+\eta\mathbb{Z}^3$ each coordinate of which lies between the corresponding coordinates of $x$ and of $x'$. Then $\mathbb{W}$ is a product of three intervals of $\eta\mathbb{Z}$ or of $\eta/2+\eta\mathbb{Z}$, and it contains the six paths $\Gamma^{\pi}(x,x')$, since along each of them every coordinate changes monotonically from its value at $x$ to its value at $x'$. On the bonds of $\mathbb{W}$ the connection $\boldsymbol{\omega}$ is real and antisymmetric, and $d\boldsymbol{\omega}$ vanishes on the plaquettes of $\mathbb{W}$. Let $s_W$ be the largest of the lengths of the three intervals. Lemma~\ref{lem:gauge-fixing-along} below, whose proof rests on the definitions of Section~\ref{sec:tori-nested-blocks} alone, applied to $\boldsymbol{\omega}$ on $\mathbb{W}$, provides a real function $\boldsymbol{\lambda}'$ on $\mathbb{W}$ such that
\[
\left|\boldsymbol{\omega}(b)-(d\boldsymbol{\lambda}')(b)\right|\le2s_W\cdot\sup_p\left|(d\boldsymbol{\omega})(p)\right|=0
\]
for every bond $b$ of $\mathbb{W}$, where $p$ ranges over the plaquettes of $\mathbb{W}$; thus $\boldsymbol{\omega}=d\boldsymbol{\lambda}'$ on the bonds of $\mathbb{W}$. Summing $d\boldsymbol{\lambda}'$ along a path one obtains
\[
\boldsymbol{\omega}\left(\Gamma^{\pi}(x,x')\right)=\eta^{-1}\cdot\left(\boldsymbol{\lambda}'(x')-\boldsymbol{\lambda}'(x)\right)
\]
for each of the six orders $\pi$, so that $(\tau\boldsymbol{\omega})(x,x')=\boldsymbol{\omega}\left(\Gamma^{\pi}(x,x')\right)$ for every $\pi$. The product of the relations $u(z')=u(z)\cdot\exp\left(i\,e_k\cdot\eta\cdot\boldsymbol{\omega}(z,z')\right)$ over the consecutive bonds $(z,z')$ of one of these paths is therefore
\[
u(x')=u(x)\cdot\exp\left(i\,e_k\cdot\eta\cdot(\tau\boldsymbol{\omega})(x,x')\right)=u(x)\cdot U_{\boldsymbol{\omega}}(x,x').
\]
Since $\tau$ is linear, $U_{A-\boldsymbol{\omega}}(x,x')=U_A(x,x')\cdot U_{\boldsymbol{\omega}}(x,x')^{-1}$, and consequently
\[
U_{A-\boldsymbol{\omega}}(x,x')\cdot u(x')=U_A(x,x')\cdot u(x).
\]

Let $y\in\Lambda^0$, $x\in\Delta_y$, and let $x=y_0,y_1,\dots,y_k=y$ be the centres of the blocks of sides $\eta,L\eta,\dots,1$ that contain $x$ (Section~\ref{sec:phase-factors-lattice}). By the definitions of $\tau_k$ and of the phases,
\[
\boldsymbol{\phi}_A(y,x)=\prod_{j<k}U_A(y_{j+1},y_j),\qquad\boldsymbol{\phi}_{A-\boldsymbol{\omega}}(y,x)=\prod_{j<k}U_{A-\boldsymbol{\omega}}(y_{j+1},y_j).
\]
By the last identity of the preceding paragraph, applied to the pairs $(y_{j+1},y_j)$, one has $U_{A-\boldsymbol{\omega}}(y_{j+1},y_j)=U_A(y_{j+1},y_j)\cdot u(y_{j+1})\cdot u(y_j)^{-1}$; the product of these factors telescopes, and
\[
\boldsymbol{\phi}_{A-\boldsymbol{\omega}}(y,x)\cdot u(x)=\boldsymbol{\phi}_A(y,x)\cdot u(y).
\]

Let $f$, $v$ and $w$ be spinors on $\Lambda$, on $\Lambda^0$ and on $\Lambda_{k+1}$, let $y\in\Lambda^0$, $x\in\Delta_y$, and let $y'$ be the site of $\Lambda_{k+1}$ with $y\in B(y')$. The identities for $\boldsymbol{\phi}_{A-\boldsymbol{\omega}}$ and $U_{A-\boldsymbol{\omega}}$ and the formulas of Section~\ref{sec:phase-factors-lattice} give
\begin{align*}
\left(Q_k(A-\boldsymbol{\omega})f\right)(y)
&=u(y)\cdot\eta^3\sum_{\tilde{x}\in\Delta_y}\boldsymbol{\phi}_A(y,\tilde{x})\cdot u(\tilde{x})^{-1}\cdot f(\tilde{x})
=\left(u^0\cdot Q_k(A)\cdot u^{-1}f\right)(y),\\
\left(Q_k^T\left(-(A-\boldsymbol{\omega})\right)v\right)(x)
&=u(x)\cdot\boldsymbol{\phi}_A(y,x)^{-1}\cdot u(y)^{-1}\cdot v(y)
=\left(u\cdot Q_k^T(-A)\cdot\left(u^0\right)^{-1}v\right)(x),\\
\left(Q(A-\boldsymbol{\omega})v\right)(y')
&=u(y')\cdot L^{-3}\sum_{\tilde{y}\in B(y')}U_A(y',\tilde{y})\cdot u(\tilde{y})^{-1}\cdot v(\tilde{y})
=\left(u^1\cdot Q(A)\cdot\left(u^0\right)^{-1}v\right)(y'),\\
\left(Q^T\left(-(A-\boldsymbol{\omega})\right)w\right)(y)
&=u(y)\cdot U_A(y',y)^{-1}\cdot u(y')^{-1}\cdot w(y')
=\left(u^0\cdot Q^T(-A)\cdot\left(u^1\right)^{-1}w\right)(y).
\end{align*}
These are the relations for $Q_k$, $Q_k^T$, $Q$ and $Q^T$. With the definitions $P_k(A)=Q_k^T(-A)\cdot Q_k(A)$, $P(A)=Q^T(-A)\cdot Q(A)$ and $P_{k+1}(A)=Q_k^T(-A)\cdot P(A)\cdot Q_k(A)$ of Section~\ref{sec:phase-factors-lattice} they give
\begin{align*}
P_k(A-\boldsymbol{\omega})&=u\cdot Q_k^T(-A)\cdot\left(u^0\right)^{-1}\cdot u^0\cdot Q_k(A)\cdot u^{-1}=u\cdot P_k(A)\cdot u^{-1},\\
P(A-\boldsymbol{\omega})&=u^0\cdot Q^T(-A)\cdot\left(u^1\right)^{-1}\cdot u^1\cdot Q(A)\cdot\left(u^0\right)^{-1}=u^0\cdot P(A)\cdot\left(u^0\right)^{-1},\\
P_{k+1}(A-\boldsymbol{\omega})&=u\cdot Q_k^T(-A)\cdot\left(u^0\right)^{-1}\cdot u^0\cdot P(A)\cdot\left(u^0\right)^{-1}\cdot u^0\cdot Q_k(A)\cdot u^{-1}\\
&=u\cdot P_{k+1}(A)\cdot u^{-1}.
\end{align*}

Now let $y\in Y_k$ be the parameter of the lemma. For $k\ge1$ the coefficients $\boldsymbol{\alpha}_y$ and $\boldsymbol{\beta}_y$ of Section~\ref{sec:operators-2} are constant matrices; they commute with the multiplication by the scalar function $u$, and
\[
\operatorname{Pen}_y(A-\boldsymbol{\omega})=\boldsymbol{\alpha}_y\cdot P_k(A-\boldsymbol{\omega})+\boldsymbol{\beta}_y\cdot P_{k+1}(A-\boldsymbol{\omega})=u\cdot\operatorname{Pen}_y(A)\cdot u^{-1}.
\]
For $k=0$ one has $\Lambda^0=\Lambda$ and $u^0=u$, so that $P(A-\boldsymbol{\omega})=u\cdot P(A)\cdot u^{-1}$; the constant matrix $i\gamma_3\cdot y$ commutes with $u$, and $\operatorname{Pen}_y(A-\boldsymbol{\omega})=b_L\cdot P(A-\boldsymbol{\omega})+i\gamma_3\cdot y=u\cdot\operatorname{Pen}_y(A)\cdot u^{-1}$. The number $\mathfrak{m}$ commutes with $u$ as well; hence, by~\eqref{eq:operators-fermion-2} and the relation for $D_{A-\boldsymbol{\omega}}$,
\[
F_y(A-\boldsymbol{\omega})=D_{A-\boldsymbol{\omega}}+\mathfrak{m}+\operatorname{Pen}_y(A-\boldsymbol{\omega})=u\cdot F_y(A)\cdot u^{-1}.
\]
Finally, for a spinor $f$ on $\Lambda$ and sites $x\ne x'$ with $|x-x'|<R/2$, the identity for $U_{A-\boldsymbol{\omega}}$ gives
\begin{align*}
\left(\delta_{\alpha,A-\boldsymbol{\omega}}(u\cdot f)\right)(x,x')
&=|x-x'|^{-\alpha}\cdot\left(U_A(x,x')\cdot u(x)\cdot f(x')-u(x)\cdot f(x)\right)\\
&=u(x)\cdot\left(\delta_{\alpha,A}f\right)(x,x').
\end{align*}
As to the example, let $\boldsymbol{\lambda}$ be a real scalar, $\boldsymbol{\omega}=d\boldsymbol{\lambda}$ and $u=\exp\left(i\,e_k\cdot\boldsymbol{\lambda}\right)$. Then $u$ is invariant under the translations by $R\mathbb{Z}^3$, $|u|=1$, and on every bond
\[
u(x')=u(x)\cdot\exp\left(i\,e_k\cdot\left(\boldsymbol{\lambda}(x')-\boldsymbol{\lambda}(x)\right)\right)=u(x)\cdot\exp\left(i\,e_k\cdot\eta\cdot(d\boldsymbol{\lambda})(x,x')\right);
\]
moreover $d(d\boldsymbol{\lambda})=0$, because the values of $d\boldsymbol{\lambda}$ on the four bonds of the boundary of a plaquette add up to zero.

\emph{(b)} Let $r_0,r_1,r_2\in\{1,-1\}$ be the diagonal entries of $r=r(S)$. By~\eqref{eq:operators-fermion-1}, $\gamma_{\mu}\cdot S=r_{\mu}\cdot S\cdot\gamma_{\mu}$ for $\mu\le2$; hence
\[
\Pi^{\mp}_{\mu}\cdot S=\frac{1}{2}\,S\cdot\left(1\mp r_{\mu}\cdot\gamma_{\mu}\right),
\]
which is $S\cdot\Pi^{\mp}_{\mu}$ if $r_{\mu}=1$ and $S\cdot\Pi^{\pm}_{\mu}$ if $r_{\mu}=-1$. The morphism has the form $\boldsymbol{\chi}(x)=rx+a'$ with $a'\in\mathbb{R}^3$ (Definition~\ref{def:block-torus}), and $re_{\mu}=r_{\mu}e_{\mu}$; therefore $\boldsymbol{\chi}(x\pm r_{\mu}\eta e_{\mu})=\boldsymbol{\chi}(x)\pm\eta e_{\mu}$ for every $x\in\Lambda$. By the definition of $gA$, $u_{gA}(\boldsymbol{\chi} x,\boldsymbol{\chi} z)=u_A(x,z)$ for every bond $(x,z)$ of $\Lambda$. For a spinor $f$ on $\Lambda$ and $x\in\Lambda$ the terms $\Pi^{-}_{\mu}\cdot u_A(x,x_{+})\cdot f(x_{+})$ and $\Pi^{+}_{\mu}\cdot u_A(x,x_{-})\cdot f(x_{-})$ of $\left(D_Af\right)(x)$ are called its forward and its backward term of direction $\mu$. Since $(gf)\left(\boldsymbol{\chi}(x+r_{\mu}\eta e_{\mu})\right)=S\cdot f(x+r_{\mu}\eta e_{\mu})$,
\[
\Pi^{-}_{\mu}\cdot u_{gA}(\boldsymbol{\chi} x,\boldsymbol{\chi} x+\eta e_{\mu})\cdot(gf)(\boldsymbol{\chi} x+\eta e_{\mu})
=S\cdot\frac{1}{2}\left(1-r_{\mu}\cdot\gamma_{\mu}\right)\cdot u_A(x,x+r_{\mu}\eta e_{\mu})\cdot f(x+r_{\mu}\eta e_{\mu}),
\]
which is $S$ times the forward term of direction $\mu$ in $\left(D_Af\right)(x)$ if $r_{\mu}=1$, and $S$ times the backward term if $r_{\mu}=-1$. In the same way
\[
\Pi^{+}_{\mu}\cdot u_{gA}(\boldsymbol{\chi} x,\boldsymbol{\chi} x-\eta e_{\mu})\cdot(gf)(\boldsymbol{\chi} x-\eta e_{\mu})
=S\cdot\frac{1}{2}\left(1+r_{\mu}\cdot\gamma_{\mu}\right)\cdot u_A(x,x-r_{\mu}\eta e_{\mu})\cdot f(x-r_{\mu}\eta e_{\mu})
\]
is $S$ times the backward term if $r_{\mu}=1$, and $S$ times the forward term if $r_{\mu}=-1$. Together with $(gf)(\boldsymbol{\chi} x)=S\cdot f(x)$ this gives $\left(D_{gA}\,gf\right)(\boldsymbol{\chi} x)=S\cdot\left(D_Af\right)(x)=\left(g\,D_Af\right)(\boldsymbol{\chi} x)$. Since $\boldsymbol{\chi}$ maps $\Lambda$ onto $\Lambda'_0$, it follows that $D_{gA}\cdot g=g\cdot D_A$.

By Definition~\ref{def:block-torus}, $\boldsymbol{\chi}$ is an isometry for the sup norm, it maps blocks onto blocks and centres to centres, and it maps every path $\Gamma^{\pi}(x,x')$ onto $\Gamma^{\pi}(\boldsymbol{\chi} x,\boldsymbol{\chi} x')$, bond by bond and with the orientation. Hence $(gA)\left(\Gamma^{\pi}(\boldsymbol{\chi} x,\boldsymbol{\chi} x')\right)=A\left(\Gamma^{\pi}(x,x')\right)$ for each $\pi$, so that
\[
(\tau\,gA)(\boldsymbol{\chi} x,\boldsymbol{\chi} x')=(\tau A)(x,x'),\qquad U_{gA}(\boldsymbol{\chi} x,\boldsymbol{\chi} x')=U_A(x,x')\qquad\left(|x-x'|<R/2\right).
\]
If $x\in\Delta_y$ and $x=y_0,\dots,y_k=y$ are the centres of the blocks of sides $\eta,L\eta,\dots,1$ that contain $x$, then $\boldsymbol{\chi} y_0,\dots,\boldsymbol{\chi} y_k$ are the centres of the blocks of $\mathbb{T}'$ of the same sides that contain $\boldsymbol{\chi} x$; therefore $(\tau_kgA)(\boldsymbol{\chi} y,\boldsymbol{\chi} x)=(\tau_kA)(y,x)$ and $\boldsymbol{\phi}_{gA}(\boldsymbol{\chi} y,\boldsymbol{\chi} x)=\boldsymbol{\phi}_A(y,x)$. Moreover $\boldsymbol{\chi}(\Delta_y)=\Delta'_{\boldsymbol{\chi} y}$ and $\boldsymbol{\chi}\left(B(z)\right)=B'(\boldsymbol{\chi} z)$ for every $z\in\Lambda_{k+1}$, where $\Delta'$ and $B'$ denote the corresponding sets of $\mathbb{T}'$. The phases are scalars and commute with $S$. Let $f$, $v$ and $w$ be spinors on $\Lambda$, on $\Lambda^0$ and on $\Lambda_{k+1}$, let $y\in\Lambda^0$, $x\in\Delta_y$, and let $z$ be the site of $\Lambda_{k+1}$ with $y\in B(z)$. The formulas of Section~\ref{sec:phase-factors-lattice} give
\begin{align*}
\left(Q_k(gA)\,gf\right)(\boldsymbol{\chi} y)&=\eta^3\sum_{\tilde{x}\in\Delta_y}\boldsymbol{\phi}_{gA}(\boldsymbol{\chi} y,\boldsymbol{\chi}\tilde{x})\cdot S\cdot f(\tilde{x})=S\cdot\left(Q_k(A)f\right)(y),\\
\left(Q_k^T(-gA)\,gv\right)(\boldsymbol{\chi} x)&=\boldsymbol{\phi}_{gA}(\boldsymbol{\chi} y,\boldsymbol{\chi} x)^{-1}\cdot S\cdot v(y)=S\cdot\left(Q_k^T(-A)v\right)(x),\\
\left(Q(gA)\,gv\right)(\boldsymbol{\chi} z)&=L^{-3}\sum_{\tilde{y}\in B(z)}U_{gA}(\boldsymbol{\chi} z,\boldsymbol{\chi}\tilde{y})\cdot S\cdot v(\tilde{y})=S\cdot\left(Q(A)v\right)(z),\\
\left(Q^T(-gA)\,gw\right)(\boldsymbol{\chi} y)&=U_{gA}(\boldsymbol{\chi} z,\boldsymbol{\chi} y)^{-1}\cdot S\cdot w(z)=S\cdot\left(Q^T(-A)w\right)(y).
\end{align*}
As $\boldsymbol{\chi}$ maps $\Lambda$, $\Lambda^0$ and $\Lambda_{k+1}$ onto $\Lambda'_0$, $\Lambda'_k$ and $\Lambda'_{k+1}$, these are the relations for $Q_k$, $Q_k^T$, $Q$ and $Q^T$. By the definitions of the projections in Section~\ref{sec:phase-factors-lattice},
\begin{align*}
P_k(gA)\cdot g&=Q_k^T(-gA)\cdot Q_k(gA)\cdot g=Q_k^T(-gA)\cdot g\cdot Q_k(A)=g\cdot Q_k^T(-A)\cdot Q_k(A)=g\cdot P_k(A),\\
P(gA)\cdot g&=Q^T(-gA)\cdot Q(gA)\cdot g=Q^T(-gA)\cdot g\cdot Q(A)=g\cdot Q^T(-A)\cdot Q(A)=g\cdot P(A),\\
P_{k+1}(gA)\cdot g&=Q_k^T(-gA)\cdot P(gA)\cdot Q_k(gA)\cdot g=Q_k^T(-gA)\cdot P(gA)\cdot g\cdot Q_k(A)\\
&=Q_k^T(-gA)\cdot g\cdot P(A)\cdot Q_k(A)=g\cdot Q_k^T(-A)\cdot P(A)\cdot Q_k(A)=g\cdot P_{k+1}(A).
\end{align*}
For a spinor $f$ on $\Lambda$ and sites $x\ne x'$ with $|x-x'|<R/2$ one has $|\boldsymbol{\chi} x-\boldsymbol{\chi} x'|=|x-x'|$, and
\[
\left(\delta_{\alpha,gA}\,gf\right)(\boldsymbol{\chi} x,\boldsymbol{\chi} x')=|x-x'|^{-\alpha}\cdot\left(U_{gA}(\boldsymbol{\chi} x,\boldsymbol{\chi} x')\cdot S\cdot f(x')-S\cdot f(x)\right)=S\cdot\left(\delta_{\alpha,A}f\right)(x,x').
\]

It remains to treat $F_y$, where $y\in Y_k$ is the parameter of the lemma. The group $\mathrm{Sp}$ is generated by $\gamma_0,\dots,\gamma_3$ (Section~\ref{sec:dirac-matrices-lattice}), and by~\eqref{eq:operators-fermion-1} each $\gamma_{\nu}$ satisfies $\gamma_{\nu}\cdot\gamma_3\cdot\gamma_{\nu}^{-1}=\pm\gamma_3$; hence $S\cdot\gamma_3\cdot S^{-1}=\pm\gamma_3$, and by the definition of $\mathrm{Sp}_3$ (Section~\ref{sec:dirac-matrices-lattice}) the sign is $+$ exactly when $S\in\mathrm{Sp}_3$. If $y$ is not neutral, then $y$ is real, $y'=y$ for $S\in\mathrm{Sp}_3$ and $y'=-y$ otherwise, so that $S\cdot\left(\gamma_3\cdot y\right)\cdot S^{-1}=\gamma_3\cdot y'$; consequently $S\cdot\phi(\gamma_3\cdot y)\cdot S^{-1}=\phi(\gamma_3\cdot y')$ for every rational function $\phi$ for which $\phi(\gamma_3\cdot y)$ is defined, and by Section~\ref{sec:operators-2} the coefficients $\boldsymbol{\alpha}_y$, $\boldsymbol{\beta}_y$ ($k\ge1$) and $i\gamma_3\cdot y$ ($k=0$) are such functions of $\gamma_3\cdot y$. If $y$ is neutral, then $y'=y$, and these coefficients are numbers, namely $\boldsymbol{\alpha}_{\infty}=b_k$, $\boldsymbol{\beta}_{\infty}=0$, $\boldsymbol{\alpha}_0=0$ and $\boldsymbol{\beta}_0=b_k\cdot b_L\cdot\left(b_k+b_L\right)^{-1}$ for $k\ge1$ (Section~\ref{sec:operators-2} and~\eqref{eq:operators-fermion-3}), and $i\gamma_3\cdot0=0$ for $k=0$; they commute with $S$. In all cases
\[
S\cdot\boldsymbol{\alpha}_y=\boldsymbol{\alpha}_{y'}\cdot S,\qquad S\cdot\boldsymbol{\beta}_y=\boldsymbol{\beta}_{y'}\cdot S\quad(k\ge1),\qquad S\cdot\left(i\gamma_3\cdot y\right)=\left(i\gamma_3\cdot y'\right)\cdot S\quad(k=0).
\]
For a constant matrix $\Theta$ and a spinor $h$ on $\Lambda$ one has $(\Theta\cdot gh)(\boldsymbol{\chi} x)=\Theta\cdot S\cdot h(x)$; hence $\boldsymbol{\alpha}_{y'}\cdot g=g\cdot\boldsymbol{\alpha}_y$, $\boldsymbol{\beta}_{y'}\cdot g=g\cdot\boldsymbol{\beta}_y$ and $\left(i\gamma_3\cdot y'\right)\cdot g=g\cdot\left(i\gamma_3\cdot y\right)$, while the number $\mathfrak{m}$ commutes with $g$. With~\eqref{eq:operators-fermion-2}, the definition of $\operatorname{Pen}_y$ in Section~\ref{sec:operators-2} and the relations for $D$, $P_k$, $P_{k+1}$ and $P$ proved above, one obtains for $k\ge1$
\begin{align*}
F_{y'}(gA)\cdot g&=D_{gA}\cdot g+\mathfrak{m}\cdot g+\boldsymbol{\alpha}_{y'}\cdot P_k(gA)\cdot g+\boldsymbol{\beta}_{y'}\cdot P_{k+1}(gA)\cdot g\\
&=g\cdot D_A+g\cdot\mathfrak{m}+\boldsymbol{\alpha}_{y'}\cdot g\cdot P_k(A)+\boldsymbol{\beta}_{y'}\cdot g\cdot P_{k+1}(A)\\
&=g\cdot\left(D_A+\mathfrak{m}+\boldsymbol{\alpha}_y\cdot P_k(A)+\boldsymbol{\beta}_y\cdot P_{k+1}(A)\right)=g\cdot F_y(A),
\end{align*}
and for $k=0$
\begin{align*}
F_{y'}(gA)\cdot g&=D_{gA}\cdot g+\mathfrak{m}\cdot g+b_L\cdot P(gA)\cdot g+\left(i\gamma_3\cdot y'\right)\cdot g\\
&=g\cdot\left(D_A+\mathfrak{m}+b_L\cdot P(A)+i\gamma_3\cdot y\right)=g\cdot F_y(A).
\end{align*}

\emph{(c)} Let $h$ and $f$ be spinors on $\Lambda$. By Section~\ref{sec:norms-pairings-kernels},
\[
(h,D_Af)=\sum_x\eta^3\cdot h(x)\cdot\left(D_Af\right)(x),
\]
the sum extending over a period cell. Into this sum the formula for $D_A$ is inserted; in the terms that contain $f(x_{+})$ one substitutes $x=x'_{-}$, in those that contain $f(x_{-})$ one substitutes $x=x'_{+}$, and, the summands being invariant under the translations by $R\mathbb{Z}^3$, the new variable $x'$ again ranges over a period cell. Collecting the terms that contain $f(x')$, and using $a\cdot(\Theta b)=(\Theta^Ta)\cdot b$ for $a,b\in\mathbb{C}^4$ and a $4\times4$ matrix $\Theta$, one finds $(h,D_Af)=(D_A^Th,f)$ with
\[
\left(D_A^Th\right)(x')=-\eta^{-1}\sum_{\mu}\Big\{\left(\Pi^{-}_{\mu}\right)^T\cdot u_A(x'_{-},x')\cdot h(x'_{-})+\left(\Pi^{+}_{\mu}\right)^T\cdot u_A(x'_{+},x')\cdot h(x'_{+})-h(x')\Big\}.
\]
Since $A$ is antisymmetric, $u_A(x'_{\mp},x')=u_{-A}(x',x'_{\mp})$, and by~\eqref{eq:operators-fermion-1}
\[
\left(\Pi^{\mp}_{\mu}\right)^T=\frac{1}{2}\left(1\mp\gamma_{\mu}^T\right)=\frac{1}{2}\left(1\pm\mathsf{C}^{-1}\cdot\gamma_{\mu}\cdot\mathsf{C}\right)=\mathsf{C}^{-1}\cdot\Pi^{\pm}_{\mu}\cdot\mathsf{C}.
\]
As $\mathsf{C}$ is a constant matrix and the phases are scalars, it follows that
\begin{align*}
\left(D_A^Th\right)(x')
&=-\eta^{-1}\cdot\mathsf{C}^{-1}\cdot\sum_{\mu}\Big\{\Pi^{+}_{\mu}\cdot u_{-A}(x',x'_{-})\cdot\left(\mathsf{C} h\right)(x'_{-})\\
&\qquad\qquad+\Pi^{-}_{\mu}\cdot u_{-A}(x',x'_{+})\cdot\left(\mathsf{C} h\right)(x'_{+})-\left(\mathsf{C} h\right)(x')\Big\}\\
&=\left(\mathsf{C}^{-1}\cdot D_{-A}\cdot\mathsf{C}\,h\right)(x'),
\end{align*}
the term with $\Pi^{-}_{\mu}$ and $x'_{+}$ being the forward term, and the term with $\Pi^{+}_{\mu}$ and $x'_{-}$ the backward term, of direction $\mu$ in $\left(D_{-A}\,\mathsf{C} h\right)(x')$. Thus $D_A^T=\mathsf{C}^{-1}\cdot D_{-A}\cdot\mathsf{C}$.

Let $f$ be a spinor on $\Lambda$ and $v$ a spinor on $\Lambda^0$. The unit blocks $\Delta_y$ partition $\Lambda$; hence
\[
\left(Q_k(A)f,v\right)=\sum_y\sum_{x\in\Delta_y}\eta^3\cdot\boldsymbol{\phi}_A(y,x)\cdot f(x)\cdot v(y)=\sum_x\eta^3\cdot f(x)\cdot\left(\boldsymbol{\phi}_A(y,x)\cdot v(y)\right),
\]
where in the last sum $y$ is the unit site with $x\in\Delta_y$. The formula for $Q_k^T(-A)$ in Section~\ref{sec:phase-factors-lattice}, with $-A$ in place of $A$, and the identity $\boldsymbol{\phi}_{-A}=\boldsymbol{\phi}_A^{-1}$, which holds because the exponent of $\boldsymbol{\phi}_A$ is linear in $A$, give $\left(Q_k^T(A)v\right)(x)=\boldsymbol{\phi}_{-A}(y,x)^{-1}\cdot v(y)=\boldsymbol{\phi}_A(y,x)\cdot v(y)$. Hence $\left(Q_k(A)f,v\right)=\left(f,Q_k^T(A)v\right)$, that is, $Q_k(A)^T=Q_k^T(A)$. Let $v$ be a spinor on $\Lambda^0$ and $w$ a spinor on $\Lambda_{k+1}$. The weight $L^3$ of the pairing on $\Lambda_{k+1}$ cancels the factor $L^{-3}$ in $Q(A)$, and the sets $B(y')$ partition $\Lambda^0$; so
\[
\left(Q(A)v,w\right)=\sum_{y'}\sum_{y\in B(y')}U_A(y',y)\cdot v(y)\cdot w(y')=\sum_yv(y)\cdot\left(U_A(y',y)\cdot w(y')\right)=\left(v,Q^T(A)w\right),
\]
where in the second and third expressions $y'$ is the site of $\Lambda_{k+1}$ with $y\in B(y')$, and where $\left(Q^T(A)w\right)(y)=U_{-A}(y',y)^{-1}\cdot w(y')=U_A(y',y)\cdot w(y')$ by the formula for $Q^T(-A)$ with $-A$ in place of $A$. Thus $Q(A)^T=Q^T(A)$.

Applied to $-A$ and transposed, these two relations give $Q_k^T(-A)^T=Q_k(-A)$ and $Q^T(-A)^T=Q(-A)$. Transposition reverses the order of the factors; therefore
\begin{align*}
P_k(A)^T&=Q_k(A)^T\cdot Q_k^T(-A)^T=Q_k^T(A)\cdot Q_k(-A)=P_k(-A),\\
P(A)^T&=Q(A)^T\cdot Q^T(-A)^T=Q^T(A)\cdot Q(-A)=P(-A),\\
P_{k+1}(A)^T&=Q_k(A)^T\cdot P(A)^T\cdot Q_k^T(-A)^T=Q_k^T(A)\cdot P(-A)\cdot Q_k(-A)=P_{k+1}(-A).
\end{align*}

A constant $4\times4$ matrix $\Theta$, acting at every site, has as its transpose the matrix $\Theta^T$, acting at every site. Let $\phi$ be a rational function and $\xi$ a real number such that $\phi(\gamma_3\cdot\xi)$ is defined. Then $\phi(\gamma_3\cdot\xi)^T=\phi(\gamma_3^T\cdot\xi)$, and $\gamma_3^T=-\mathsf{C}^{-1}\cdot\gamma_3\cdot\mathsf{C}$ by~\eqref{eq:operators-fermion-1}; hence
\[
\phi(\gamma_3\cdot\xi)^T=\mathsf{C}^{-1}\cdot\phi(-\gamma_3\cdot\xi)\cdot\mathsf{C}.
\]
Now let $y\in Y_k$ be the parameter of the lemma. By Section~\ref{sec:operators-2}, for real $y$ the coefficients $\boldsymbol{\alpha}_y$, $\boldsymbol{\beta}_y$ ($k\ge1$) and $i\gamma_3\cdot y$ ($k=0$) are such functions of $\gamma_3\cdot y$, so that
\[
\boldsymbol{\alpha}_y^T=\mathsf{C}^{-1}\cdot\boldsymbol{\alpha}_{-y}\cdot\mathsf{C},\qquad\boldsymbol{\beta}_y^T=\mathsf{C}^{-1}\cdot\boldsymbol{\beta}_{-y}\cdot\mathsf{C},\qquad\left(i\gamma_3\cdot y\right)^T=\mathsf{C}^{-1}\cdot\left(i\gamma_3\cdot(-y)\right)\cdot\mathsf{C};
\]
for $y=\infty=-y$ the coefficients $\boldsymbol{\alpha}_{\infty}=b_k$ and $\boldsymbol{\beta}_{\infty}=0$ are numbers, and the same relations hold. The operators $P_k(A)$, $P(A)$ and $P_{k+1}(A)$ commute with the constant matrices (Lemma~\ref{lem:block-averages-block}). For $k\ge1$ it follows that
\begin{align*}
\left(\boldsymbol{\alpha}_y\cdot P_k(A)\right)^T&=P_k(A)^T\cdot\boldsymbol{\alpha}_y^T=P_k(-A)\cdot\mathsf{C}^{-1}\cdot\boldsymbol{\alpha}_{-y}\cdot\mathsf{C}=\mathsf{C}^{-1}\cdot\boldsymbol{\alpha}_{-y}\cdot P_k(-A)\cdot\mathsf{C},\\
\left(\boldsymbol{\beta}_y\cdot P_{k+1}(A)\right)^T&=P_{k+1}(A)^T\cdot\boldsymbol{\beta}_y^T=P_{k+1}(-A)\cdot\mathsf{C}^{-1}\cdot\boldsymbol{\beta}_{-y}\cdot\mathsf{C}=\mathsf{C}^{-1}\cdot\boldsymbol{\beta}_{-y}\cdot P_{k+1}(-A)\cdot\mathsf{C},
\end{align*}
so that $\operatorname{Pen}_y(A)^T=\mathsf{C}^{-1}\cdot\operatorname{Pen}_{-y}(-A)\cdot\mathsf{C}$. For $k=0$ one has $\operatorname{Pen}_y(A)^T=b_L\cdot P(-A)+\mathsf{C}^{-1}\cdot\left(i\gamma_3\cdot(-y)\right)\cdot\mathsf{C}=\mathsf{C}^{-1}\cdot\operatorname{Pen}_{-y}(-A)\cdot\mathsf{C}$, again because $P(-A)$ commutes with $\mathsf{C}$. The number $\mathfrak{m}$ is its own transpose. With~\eqref{eq:operators-fermion-2} and the relation for $D_A^T$ one concludes that $F_y(A)^T=\mathsf{C}^{-1}\cdot F_{-y}(-A)\cdot\mathsf{C}$.

\emph{(d)} Complex conjugation of operators is additive and multiplicative, $\overline{O_1\cdot O_2}=\overline{O_1}\cdot\overline{O_2}$, and it fixes real numbers. Since $e_k$ and $\eta$ are real and $\tau$, $\tau_k$ are linear with real coefficients,
\[
\overline{u_{\overline{A}}}=u_{-A},\qquad\overline{U_{\overline{A}}}=U_{-A},\qquad\overline{\boldsymbol{\phi}_{\overline{A}}}=\boldsymbol{\phi}_{-A}.
\]
By~\eqref{eq:operators-fermion-1}, $\varkappa$ is real, $\varkappa^{-1}=\varkappa$, $\overline{\gamma_{\mu}}=\varkappa^{-1}\cdot\gamma_{\mu}\cdot\varkappa$ for $\mu\le2$, and $\overline{\gamma_3}=-\varkappa^{-1}\cdot\gamma_3\cdot\varkappa$. Hence $\overline{\Pi^{\mp}_{\mu}}=\varkappa^{-1}\cdot\Pi^{\mp}_{\mu}\cdot\varkappa$, and for real $\xi$
\[
\overline{i\gamma_3\cdot\xi}=-i\,\overline{\gamma_3}\cdot\xi=\varkappa^{-1}\cdot\left(i\gamma_3\cdot\xi\right)\cdot\varkappa;
\]
consequently $\overline{\phi(i\gamma_3\cdot\xi)}=\varkappa^{-1}\cdot\phi(i\gamma_3\cdot\xi)\cdot\varkappa$ for every rational function $\phi$ with real coefficients for which $\phi(i\gamma_3\cdot\xi)$ is defined. For real $y$ the coefficients $\boldsymbol{\alpha}_y$ and $\boldsymbol{\beta}_y$ of Section~\ref{sec:operators-2} are of this form, the coefficients of the rational functions being formed with the real numbers $b_k$ and $b_L$, and $\boldsymbol{\alpha}_{\infty}=b_k$, $\boldsymbol{\beta}_{\infty}=0$ are real; thus
\[
\overline{\boldsymbol{\alpha}_y}=\varkappa^{-1}\cdot\boldsymbol{\alpha}_y\cdot\varkappa,\qquad\overline{\boldsymbol{\beta}_y}=\varkappa^{-1}\cdot\boldsymbol{\beta}_y\cdot\varkappa\qquad\left(y\in Y_k\right).
\]
In the formula for $D_{\overline{A}}$ of Section~\ref{sec:phase-factors-lattice} the coefficients are the real number $\eta^{-1}$, the matrices $\Pi^{\mp}_{\mu}$, the unit matrix and the phases $u_{\overline{A}}$. Replacing each of them by its complex conjugate one obtains
\[
\overline{D_{\overline{A}}}=\varkappa^{-1}\cdot D_{-A}\cdot\varkappa,
\]
since $\varkappa$ is a constant matrix and commutes with the scalar phases and with the translations of the lattice. The kernels of $Q_k(\overline{A})$, $Q_k^T(-\overline{A})$, $Q(\overline{A})$ and $Q^T(-\overline{A})$ are the scalars $\eta^3\cdot\boldsymbol{\phi}_{\overline{A}}$, $\boldsymbol{\phi}_{\overline{A}}^{-1}$, $L^{-3}\cdot U_{\overline{A}}$ and $U_{\overline{A}}^{-1}$ (Section~\ref{sec:phase-factors-lattice}); their complex conjugates $\eta^3\cdot\boldsymbol{\phi}_{-A}$, $\boldsymbol{\phi}_{-A}^{-1}$, $L^{-3}\cdot U_{-A}$ and $U_{-A}^{-1}$ are the kernels of $Q_k(-A)$, $Q_k^T(A)$, $Q(-A)$ and $Q^T(A)$, and these operators commute with $\varkappa$ (Lemma~\ref{lem:block-averages-block}). This is the assertion for $O(A)=Q_k(A)$, $Q_k^T(-A)$, $Q(A)$, $Q^T(-A)$. For the projections the multiplicativity of the conjugation gives
\begin{align*}
\overline{P_k(\overline{A})}&=\overline{Q_k^T(-\overline{A})}\cdot\overline{Q_k(\overline{A})}=Q_k^T(A)\cdot Q_k(-A)=P_k(-A)=\varkappa^{-1}\cdot P_k(-A)\cdot\varkappa,\\
\overline{P(\overline{A})}&=\overline{Q^T(-\overline{A})}\cdot\overline{Q(\overline{A})}=Q^T(A)\cdot Q(-A)=P(-A)=\varkappa^{-1}\cdot P(-A)\cdot\varkappa,\\
\overline{P_{k+1}(\overline{A})}&=\overline{Q_k^T(-\overline{A})}\cdot\overline{P(\overline{A})}\cdot\overline{Q_k(\overline{A})}=Q_k^T(A)\cdot P(-A)\cdot Q_k(-A)\\
&=P_{k+1}(-A)=\varkappa^{-1}\cdot P_{k+1}(-A)\cdot\varkappa.
\end{align*}
Finally, $\mathfrak{m}$ and $b_L$ are real. For $k\ge1$ it follows that
\begin{align*}
\overline{F_y(\overline{A})}
&=\varkappa^{-1}\cdot D_{-A}\cdot\varkappa+\mathfrak{m}+\varkappa^{-1}\cdot\boldsymbol{\alpha}_y\cdot\varkappa\cdot P_k(-A)\\
&\qquad+\varkappa^{-1}\cdot\boldsymbol{\beta}_y\cdot\varkappa\cdot P_{k+1}(-A)\\
&=\varkappa^{-1}\cdot F_y(-A)\cdot\varkappa,
\end{align*}
and for $k=0$
\[
\overline{F_y(\overline{A})}=\varkappa^{-1}\cdot D_{-A}\cdot\varkappa+\mathfrak{m}+b_L\cdot P(-A)+\varkappa^{-1}\cdot\left(i\gamma_3\cdot y\right)\cdot\varkappa=\varkappa^{-1}\cdot F_y(-A)\cdot\varkappa,
\]
the operators $P_k(-A)$, $P_{k+1}(-A)$ and $P(-A)$ commuting with $\varkappa$.

\emph{(e)} Let $A$ be real. By Section~\ref{sec:norms-pairings-kernels}, the sesquilinear pairing of two spinors $h$, $f$ on $\Lambda$ is $(\overline{h},f)=\sum_x\eta^3\cdot\overline{h(x)}\cdot f(x)$, and likewise on $\Lambda^0$ and on $\Lambda_{k+1}$, with the weights $1$ and $L^3$; the adjoint $O^{*}$ of an operator $O$ is characterised by $(\overline{h},Of)=(\overline{O^{*}h},f)$. The numbers $e_k\cdot\eta\cdot A(x,x')$, $e_k\cdot\eta\cdot(\tau A)(x,x')$ and $e_k\cdot\eta\cdot(\tau_kA)(y,x)$ are real, so that the phases have modulus $1$:
\[
\overline{u_A}=u_A^{-1},\qquad\overline{U_A}=U_A^{-1},\qquad\overline{\boldsymbol{\phi}_A}=\boldsymbol{\phi}_A^{-1};
\]
moreover $u_A(x,x')^{-1}=u_A(x',x)$, because $A$ is antisymmetric. For $\mu\in\{0,1,2\}$ let $K^{\pm}_{\mu}$ be the operator on the spinors on $\Lambda$ defined by $\left(K^{\pm}_{\mu}f\right)(x):=u_A(x,x_{\pm})\cdot f(x_{\pm})$. For spinors $h$, $f$ on $\Lambda$ the substitution $x=x'_{-}$ in the periodic summand gives
\begin{align*}
\left(\overline{h},K^{+}_{\mu}f\right)
&=\sum_x\eta^3\cdot\overline{h(x)}\cdot u_A(x,x_{+})\cdot f(x_{+})\\
&=\sum_{x'}\eta^3\cdot\overline{u_A(x',x'_{-})\cdot h(x'_{-})}\cdot f(x')=\left(\overline{K^{-}_{\mu}h},f\right),
\end{align*}
since $\overline{u_A(x'_{-},x')}=u_A(x'_{-},x')^{-1}=u_A(x',x'_{-})$. Thus $\left(K^{+}_{\mu}\right)^{*}=K^{-}_{\mu}$, and therefore $\left(K^{-}_{\mu}\right)^{*}=K^{+}_{\mu}$. The matrices $\Pi^{\mp}_{\mu}$ are Hermitian, the $\gamma_{\mu}$ being Hermitian by~\eqref{eq:operators-fermion-1}, and they commute with $K^{\pm}_{\mu}$, whose kernels are scalar. Since
\[
D_A=-\eta^{-1}\sum_{\mu}\left(\Pi^{-}_{\mu}\cdot K^{+}_{\mu}+\Pi^{+}_{\mu}\cdot K^{-}_{\mu}-1\right),
\]
it follows that
\[
D_A^{*}=-\eta^{-1}\sum_{\mu}\left(\Pi^{-}_{\mu}\cdot K^{-}_{\mu}+\Pi^{+}_{\mu}\cdot K^{+}_{\mu}-1\right).
\]
By~\eqref{eq:operators-fermion-1}, $\gamma_3^2=1$ and $\gamma_3\cdot\gamma_{\mu}\cdot\gamma_3=-\gamma_{\mu}$ for $\mu\le2$, so that $\gamma_3\cdot\Pi^{\mp}_{\mu}\cdot\gamma_3=\Pi^{\pm}_{\mu}$; as $\gamma_3$ commutes with $K^{\pm}_{\mu}$,
\[
\gamma_3\cdot D_A\cdot\gamma_3=-\eta^{-1}\sum_{\mu}\left(\Pi^{+}_{\mu}\cdot K^{+}_{\mu}+\Pi^{-}_{\mu}\cdot K^{-}_{\mu}-1\right)=D_A^{*}.
\]

Let $f$ be a spinor on $\Lambda$ and $v$ a spinor on $\Lambda^0$. By $\overline{\boldsymbol{\phi}_A}=\boldsymbol{\phi}_A^{-1}$ and the formula for $Q_k^T(-A)$ in Section~\ref{sec:phase-factors-lattice},
\begin{align*}
\left(\overline{v},Q_k(A)f\right)
&=\sum_y\overline{v(y)}\cdot\eta^3\sum_{x\in\Delta_y}\boldsymbol{\phi}_A(y,x)\cdot f(x)\\
&=\sum_x\eta^3\cdot\overline{\boldsymbol{\phi}_A(y,x)^{-1}\cdot v(y)}\cdot f(x)=\left(\overline{Q_k^T(-A)v},f\right),
\end{align*}
where in the second and third expressions $y$ is the unit site with $x\in\Delta_y$; thus $Q_k(A)^{*}=Q_k^T(-A)$. Let $v$ be a spinor on $\Lambda^0$ and $w$ a spinor on $\Lambda_{k+1}$. By $\overline{U_A}=U_A^{-1}$,
\begin{align*}
\left(\overline{w},Q(A)v\right)
&=\sum_{y'}L^3\cdot\overline{w(y')}\cdot L^{-3}\sum_{y\in B(y')}U_A(y',y)\cdot v(y)\\
&=\sum_y\overline{U_A(y',y)^{-1}\cdot w(y')}\cdot v(y)=\left(\overline{Q^T(-A)w},v\right),
\end{align*}
where in the second and third expressions $y'$ is the site of $\Lambda_{k+1}$ with $y\in B(y')$; thus $Q(A)^{*}=Q^T(-A)$.

Taking adjoints in these two relations one obtains $Q_k^T(-A)^{*}=Q_k(A)$ and $Q^T(-A)^{*}=Q(A)$. Hence
\begin{align*}
P_k(A)^{*}&=Q_k(A)^{*}\cdot Q_k^T(-A)^{*}=Q_k^T(-A)\cdot Q_k(A)=P_k(A),\\
P(A)^{*}&=Q(A)^{*}\cdot Q^T(-A)^{*}=Q^T(-A)\cdot Q(A)=P(A),\\
P_{k+1}(A)^{*}&=Q_k(A)^{*}\cdot P(A)^{*}\cdot Q_k^T(-A)^{*}=Q_k^T(-A)\cdot P(A)\cdot Q_k(A)=P_{k+1}(A).
\end{align*}
By Lemma~\ref{lem:block-averages-block} these three operators are projections; being self-adjoint, they are orthogonal projections.
\end{proof}

\subsection{Inverses at zero field}\label{sec:inverses-zero-field}

\begin{proposition}[Zero field, levels $k\ge1$]\label{prop:zero-field-levels}
There are numbers $C_5\ge1$ and $\boldsymbol{\gamma}_5\in\left(0,\frac{1}{12L^2}\right]$, depending only on $L$, $b$ and $\alpha$, with the following property. Let $\mathbb{T}$ be a block torus of level $k\ge1$, let $0\le\mathfrak{m}\le1$ and let $y\in\mathbb{R}\cup\{\infty\}$. Then $F_y(0)$ is bijective. Further, let $y_1$ and $y_2$ be unit sites, let $f$ be a spinor on $\mathbb{T}$ supported in the orbit of $\Delta_{y_2}$, and put $v:=S_y(0)f$. Then for every $x\in\Delta_{y_1}$ and every site $\tilde{x}$ with $0<|x-\tilde{x}|\le1$
\[
|v(x)|+|x-\tilde{x}|^{-\alpha}\cdot|v(x)-v(\tilde{x})|\le C_5\cdot\exp\left(-\boldsymbol{\gamma}_5\cdot d(y_1,y_2)\right)\cdot|f|_{\infty}.
\]
\end{proposition}

\begin{proof}
Let $C_5$ and $\boldsymbol{\gamma}_5$ be the numbers of Theorem~\ref{thm:sup-holder-bounds-2}; by that theorem they depend only on $(L,b,\alpha)$. The proposition is a special case of Theorem~\ref{thm:sup-holder-bounds-2}. This theorem is applied with $i=k$, with the H\"older exponent $\alpha$ and with $b'=b_k$; its directions $1,2,3$ and the matrices $\gamma_1,\gamma_2,\gamma_3$ of Convention~\ref{it:wilson-operators-3} of Section~\ref{sec:conventions} are taken to be the directions $0,1,2$ and the matrices $\gamma_0,\gamma_1,\gamma_2$ of Section~\ref{sec:parameters-clifford-relations}, and the matrix $\hat{\Gamma}$ of the same convention is taken to be $\gamma_3$. By~\eqref{eq:operators-fermion-1}, these matrices satisfy the requirements of that convention.

With these choices the operator $T^{*}_y$ of Theorem~\ref{thm:sup-holder-bounds-2} is $D+\mathfrak{m}+\boldsymbol{\alpha}_yP_k+\boldsymbol{\beta}_yP_{k+1}$, where $\boldsymbol{\alpha}_y$ and $\boldsymbol{\beta}_y$ are the matrices of Definition~\ref{def:standard-penalties} formed with $b'=b_k$ and $\hat{\Gamma}=\gamma_3$; the number $b/L$ that enters that definition is $b_L$. These matrices coincide with those of Section~\ref{sec:operators-2}. For $y=\infty$ both definitions give $\boldsymbol{\alpha}_{\infty}=b_k$ and $\boldsymbol{\beta}_{\infty}=0$. For real $y$ both give
\[
\boldsymbol{\alpha}_y=b_k\,i\gamma_3y\left(b_k+i\gamma_3y\right)^{-1}=b_k-b_k^2\left(b_k+i\gamma_3y\right)^{-1},
\]
and the formula for $\boldsymbol{\alpha}_y+\boldsymbol{\beta}_y$ in Definition~\ref{def:standard-penalties} yields
\begin{multline*}
\boldsymbol{\beta}_y=\left(\boldsymbol{\alpha}_y+\boldsymbol{\beta}_y\right)-\boldsymbol{\alpha}_y=b_k^2\left[\left(b_k+i\gamma_3y\right)^{-1}-\left(b_k+b_L+i\gamma_3y\right)^{-1}\right]\\
=b_k^2\cdot b_L\cdot\left(b_k+b_L+i\gamma_3y\right)^{-1}\left(b_k+i\gamma_3y\right)^{-1},
\end{multline*}
which is the matrix $\boldsymbol{\beta}_y$ of Section~\ref{sec:operators-2}; in the last step the two inverses commute, both being functions of $\gamma_3$. The operator $D$ of~\eqref{eq:wilson-operators-7} is the operator $D_0$ of Section~\ref{sec:phase-factors-lattice}.

Theorem~\ref{thm:sup-holder-bounds-2} concerns two objects introduced in Section~\ref{sec:conventions}: the space $\mathbb{F}(\eta,R,\boldsymbol{\theta})$ of quasi-periodic functions on $\eta\mathbb{Z}^3$ of Convention~\ref{it:wilson-operators-2}, which coincides with $\mathbb{H}(\eta,R,\boldsymbol{\theta})$ because $R<\infty$, and the blocks of Convention~\ref{it:wilson-operators-4}, whose centres form the lattices $L^j\eta\mathbb{Z}^3$. Here $\eta=L^{-k}$; the side $R$ is a multiple of $L$ by Definition~\ref{def:block-torus}; and $\boldsymbol{\theta}:=0$ if $\boldsymbol{\sigma}=1$, $\boldsymbol{\theta}:=(\pi,\pi,\pi)$ if $\boldsymbol{\sigma}=-1$.

Let first $t=0$. Then the spinors on $\mathbb{T}=(0,R,\boldsymbol{\sigma})$ are the elements of $\mathbb{F}(\eta,R,\boldsymbol{\theta})$. For $0\le j\le k+1$ and $z\in\Lambda_j=L^j\eta\mathbb{Z}^3$, the block $\mathrm{Bl}_j(z)$ is the block with centre $z$ and side $L^j\eta$ in the sense of Convention~\ref{it:wilson-operators-4} of Section~\ref{sec:conventions}, since a point $x\in\eta\mathbb{Z}^3$ satisfies $|x-z|<L^j\eta/2$ if and only if $|x-z|\le(L^j-1)\eta/2$, $L^j$ being odd. By the formulas of Section~\ref{sec:phase-factors-lattice}, the operators $P_k(0)$ and $P_{k+1}(0)$ replace a spinor on every block of side $1$, respectively of side $L$, by its mean value over that block; so they are the projections $P_k$ and $P_{k+1}$ of Theorem~\ref{thm:sup-holder-bounds-2}, and $T^{*}_y=F_y(0)$. Further, $\Delta_{y_1}$ and $\Delta_{y_2}$ are the blocks of side one with centres $y_1$ and $y_2$, and $d(y_1,y_2)$ is the periodic sup distance of these centres, the distance in the exponent of the bound of Theorem~\ref{thm:sup-holder-bounds-2}. That theorem, applied to the function $f$, to the point $x$ of $\Delta_{y_1}$ and to the point $\tilde{x}$, gives the bijectivity of $F_y(0)$ and the asserted inequality.

Let now $t=\frac{\eta}{2}(1,1,1)$. As $L=L^{k+1}\eta$, one has for $0\le j\le k+1$
\[
\frac{L^j\eta}{2}-\frac{L}{2}=L^j\eta\cdot\frac{1-L^{k+1-j}}{2},
\]
and $1-L^{k+1-j}$ is even, $L$ being odd. Hence the translation $x\mapsto x-\frac{L}{2}(1,1,1)$ maps $\Lambda_j$ onto $L^j\eta\mathbb{Z}^3$ for $0\le j\le k+1$. Transport along this translation, $f\mapsto f\left(\cdot+\frac{L}{2}(1,1,1)\right)$, maps the spinors on $\mathbb{T}$ bijectively onto the spinors on the block torus $(0,R,\boldsymbol{\sigma})$ of level $k$, since the case $j=0$ gives $\Lambda-\frac{L}{2}(1,1,1)=\eta\mathbb{Z}^3$ and the rule $f(x+Re_{\mu})=\boldsymbol{\sigma}\cdot f(x)$ is kept; it preserves the sup norm, and the translation preserves the periodic distance $d$ as well as the distances $|x-\tilde{x}|$. The blocks of side $1$ and of side $L$ of $\mathbb{T}$ are carried onto those of $(0,R,\boldsymbol{\sigma})$, which by the case $t=0$ are blocks in the sense of Convention~\ref{it:wilson-operators-4} of Section~\ref{sec:conventions}; hence the transport conjugates the mean values over these blocks, $P_k(0)$ and $P_{k+1}(0)$, to the projections $P_k$ and $P_{k+1}$ of Theorem~\ref{thm:sup-holder-bounds-2}. The formula for $D_0$ in Section~\ref{sec:phase-factors-lattice} expresses $(D_0f)(x)$ through $f(x)$ and the values $f(x\pm\eta e_{\mu})$ alone, with matrices that do not depend on $x$; so the transport conjugates $D_0$ to the operator given by the same formula on $(0,R,\boldsymbol{\sigma})$. The multiplication by $\mathfrak{m}$ and by the matrices $\boldsymbol{\alpha}_y$ and $\boldsymbol{\beta}_y$, which are the same at every site, commutes with the transport. The assertions therefore follow from the case $t=0$.

Finally, $\boldsymbol{\gamma}_5=a_0/(4L)$, where $0<a_0\le1/(3L)$ by~\eqref{eq:wilson-operators-4}; hence $0<\boldsymbol{\gamma}_5\le1/(12L^2)$. Moreover
\[
C_5=\mathbf{e}\cdot\tilde{C}\cdot L\cdot\left(2+2L\cdot\mathbf{e}\right)\ge1,
\]
because $\tilde{C}=2C_3\cdot C_{\mathrm{sh}}(a_0/4)$ (Section~\ref{sec:bounded-right-hand}), $C_{\mathrm{sh}}(a_0/4)\ge64$ by~\eqref{eq:wilson-operators-11} (the term with $n=0$), $C_3\ge C_3''\ge2L^5C_I$ by~\eqref{eq:wilson-operators-20} and~\eqref{eq:wilson-operators-19}, and $C_I\ge L$ by~\eqref{eq:wilson-operators-3}.
\end{proof}

\begin{lemma}[Zero field, level $0$]\label{lem:zero-field-level}
Put
\begin{gather*}
c_{\mathrm{re}}:=\min\left\{b\left(1-\frac{1}{L}\right),\frac{4}{L}\right\},\qquad a_z:=\min\left\{\frac{1}{3L},\frac{c_{\mathrm{re}}}{2\mathbf{e}\cdot L\cdot(6+3b)}\right\},\\
k_1:=14+b,\qquad C_{00}:=\frac{2L}{c_{\mathrm{re}}}\cdot\left(1+2k_1\right).
\end{gather*}
Let $\mathbb{T}$ be a block torus of level $0$, let $0\le\mathfrak{m}\le1$ and let $y\in\mathbb{R}$. Then the operator
\[
F_y(0)=D_0+\mathfrak{m}+b_L\cdot P(0)+i\gamma_3y
\]
is bijective, and for all sites $x$ and $x'$
\[
\left|S_y(0)(x,x')\right|\le C_{00}\cdot\left(1+|y|\right)^{-1}\cdot\exp\left(-a_z\cdot d(x,x')\right).
\]
\end{lemma}

\begin{proof}
At level $0$ one has $\eta=1$ and $x_{\pm}=x\pm e_{\mu}$, and by Section~\ref{sec:operators-2} the operator $T:=F_y(0)$ is $D_0+\mathfrak{m}+b_L\cdot P(0)+i\gamma_3y$. For spinors $w$, $w'$ on $\mathbb{T}$ let $\langle w,w'\rangle:=\sum_x\overline{w(x)}\cdot w'(x)$, the sum extending over a period cell, let $\|w\|_2:=\langle w,w\rangle^{\frac{1}{2}}$, and let $\|O\|_2$ be the corresponding norm of a linear operator $O$.

Proposition~\ref{prop:wilson-operators-2} is applied with $\eta=1$, $\boldsymbol{\ell}_1=L$, $W$ the whole lattice, $\boldsymbol{\beta}_E=b$ for every block $E$ of level one, $\upsilon=y$ and $\hat{\Gamma}=\gamma_3$, the directions and the Clifford matrices being identified as in the proof of Proposition~\ref{prop:zero-field-levels}, and the character $\boldsymbol{\theta}$ being chosen as there; $R$ is a multiple of $L=\boldsymbol{\ell}_1$. If $t\ne0$, it is applied after the translation $x\mapsto x-\frac{L}{2}(1,1,1)$ of that proof. With these data the operator of Proposition~\ref{prop:wilson-operators-2} is $D_0+\mathfrak{m}+b_L\cdot P(0)+i\gamma_3y=T$, because the sum of the projections $P_E$ over the blocks $E$ of side $L$ is $P(0)$; and the number $c_0$ of Proposition~\ref{prop:wilson-operators-2} is $c_{\mathrm{re}}$. Therefore
\[
\operatorname{Re}\langle w,Tw\rangle\ge\frac{c_{\mathrm{re}}}{L}\cdot\|w\|_2^2
\]
for every spinor $w$ on $\mathbb{T}$.

Let $\phi$ be a real scalar on $\mathbb{T}$ such that $|\phi(x)-\phi(x')|\le a_z\cdot|x-x'|_1$ for all sites $x$ and $x'$, where $|\cdot|_1$ is the $\ell^1$ norm on $\mathbb{R}^3$. Put
\[
T_{\phi}:=\exp(\phi)\,T\exp(-\phi),\qquad\operatorname{Rm}:=T_{\phi}-T,
\]
where $\exp(\pm\phi)$ acts by multiplication. The mass $\mathfrak{m}$ and the matrix $i\gamma_3y$ commute with these multiplications, and so they do not contribute to $\operatorname{Rm}$. By the formula for $D_0$ in Section~\ref{sec:phase-factors-lattice}, the contribution of $D_0$ to $\operatorname{Rm}$ is the sum of the six operators that map $w$ to the spinors
\[
x\mapsto-\left(\exp\left(\phi(x)-\phi(x_{\pm})\right)-1\right)\Pi^{\mp}_{\mu}\,w(x_{\pm})\qquad(\mu=0,1,2),
\]
and each of these operators has norm at most $a_z\cdot\mathbf{e}$ for $\|\cdot\|_2$. Indeed, $|\phi(x)-\phi(x_{\pm})|\le a_z\le1$; $|\exp(s)-1|\le|s|\cdot\mathbf{e}$ for every real $s$ with $|s|\le1$; $\Pi^{\mp}_{\mu}$ is an orthogonal projection; and the map $w\mapsto w(\cdot\pm e_{\mu})$ is an isometry. The contribution of $b_L\cdot P(0)$ to $\operatorname{Rm}$ has the kernel
\[
b_L\cdot L^{-3}\left(\exp\left(\phi(x)-\phi(x')\right)-1\right)
\]
at the pairs $(x,x')$ of sites of one block of side $L$, and the kernel $0$ at all other pairs; for the former pairs $|\phi(x)-\phi(x')|\le3(L-1)\cdot a_z\le1$. Every row and every column of this kernel has at most $L^3$ nonzero entries, each of modulus at most $b_L\cdot L^{-3}\cdot3(L-1)\cdot a_z\cdot\mathbf{e}$. So the row sums and the column sums are at most $b_L\cdot3L\cdot a_z\cdot\mathbf{e}=3b\cdot a_z\cdot\mathbf{e}$, and by the Schur test the norm of this contribution is at most $3b\cdot a_z\cdot\mathbf{e}$. By the definition of $a_z$ it follows that
\[
\|\operatorname{Rm}\|_2\le(6+3b)\cdot\mathbf{e}\cdot a_z\le\frac{c_{\mathrm{re}}}{2L},
\]
and by the Cauchy--Schwarz inequality
\[
\operatorname{Re}\langle w,T_{\phi}w\rangle\ge\operatorname{Re}\langle w,Tw\rangle-\|\operatorname{Rm}\|_2\cdot\|w\|_2^2\ge\frac{c_{\mathrm{re}}}{2L}\cdot\|w\|_2^2.
\]
Again by the Cauchy--Schwarz inequality, $\|T_{\phi}w\|_2\ge\frac{c_{\mathrm{re}}}{2L}\cdot\|w\|_2$ for every $w$. Hence $T_{\phi}$ is injective; as the space of spinors on $\mathbb{T}$ has finite dimension, $T_{\phi}$ is bijective, and
\[
\left\|T_{\phi}^{-1}\right\|_2\le\frac{2L}{c_{\mathrm{re}}}.
\]
Since $T=\exp(-\phi)\,T_{\phi}\exp(\phi)$, the operator $T=F_y(0)$ is bijective.

By the formula of Section~\ref{sec:phase-factors-lattice}, $D_0$ is the sum of nine operators of norm at most $1$ (for each direction, the identity and two orthogonal projections composed with translations), so that $\|D_0\|_2\le9\le12$. Further, $\mathfrak{m}\le1$, and $P(0)$, which replaces a spinor on every block of side $L$ by its mean value over that block, is an orthogonal projection. As $T_{\phi}-i\gamma_3y=D_0+\mathfrak{m}+b_L\cdot P(0)+\operatorname{Rm}$, it follows that
\[
\left\|T_{\phi}-i\gamma_3y\right\|_2\le12+1+b_L+\frac{c_{\mathrm{re}}}{2L}\le k_1,
\]
because $b_L\le b$ and $c_{\mathrm{re}}/(2L)\le2/L^2\le1$. The matrix $\gamma_3$ is Hermitian with $\gamma_3^2=1$, so that $\|i\gamma_3y\,w\|_2=|y|\cdot\|w\|_2$; therefore
\[
\|T_{\phi}w\|_2\ge\left(|y|-k_1\right)\|w\|_2,
\]
and for $|y|\ge2k_1$ this gives $\left\|T_{\phi}^{-1}\right\|_2\le2/|y|$. The two bounds for $\left\|T_{\phi}^{-1}\right\|_2$ together yield
\[
\left\|T_{\phi}^{-1}\right\|_2\le\frac{C_{00}}{1+|y|}:
\]
for $|y|<2k_1$ because $2L/c_{\mathrm{re}}=C_{00}/(1+2k_1)$, and for $|y|\ge2k_1$ because $2(1+|y|)/|y|=2+2/|y|\le3\le C_{00}$, the last inequality following from $c_{\mathrm{re}}\le4/L$, which gives $C_{00}\ge\frac{L^2}{2}\cdot(1+2k_1)$.

Now fix a site $x'$ and let $\phi:=a_z\cdot d_1(\cdot,x')$, where $d_1$ is the $\ell^1$ distance modulo $R\mathbb{Z}^3$. This $\phi$ is a real scalar, and by the triangle inequality $|\phi(x)-\phi(x'')|\le a_z\cdot d_1(x,x'')\le a_z\cdot|x-x''|_1$ for all sites $x$ and $x''$. Since $T_{\phi}^{-1}=\exp(\phi)\,S_y(0)\exp(-\phi)$ and $\phi(x')=0$, one has for every site $x$
\[
S_y(0)(x,x')=\exp\left(-\phi(x)+\phi(x')\right)T_{\phi}^{-1}(x,x')=\exp\left(-a_z\cdot d_1(x,x')\right)T_{\phi}^{-1}(x,x').
\]
Let $\zeta\in\mathbb{C}^4$, and let $w_{\zeta}$ be the spinor that vanishes outside the orbit $x'+R\mathbb{Z}^3$ and satisfies $w_{\zeta}(x'+Rn)=\boldsymbol{\sigma}^{n_0+n_1+n_2}\cdot\zeta$ for $n=n_0e_0+n_1e_1+n_2e_2\in\mathbb{Z}^3$. A period cell contains exactly one site of this orbit. In the sums that define $\|w_{\zeta}\|_2^2$ and $\left(T_{\phi}^{-1}w_{\zeta}\right)(x)$, taken over a period cell containing $x'$, only the term at $x'$ is nonzero; thus $\|w_{\zeta}\|_2=|\zeta|$ and $\left(T_{\phi}^{-1}w_{\zeta}\right)(x)=T_{\phi}^{-1}(x,x')\,\zeta$. Every spinor $u$ satisfies $|u(x)|\le\|u\|_2$, since the sum that defines $\|u\|_2^2$ may be taken over a period cell containing $x$; therefore
\[
\left|T_{\phi}^{-1}(x,x')\,\zeta\right|\le\left\|T_{\phi}^{-1}w_{\zeta}\right\|_2\le\left\|T_{\phi}^{-1}\right\|_2\cdot|\zeta|.
\]
As $\zeta$ is arbitrary and $|\cdot|$ is the operator norm of $4\times4$ matrices, $\left|T_{\phi}^{-1}(x,x')\right|\le\left\|T_{\phi}^{-1}\right\|_2$. Moreover, $d_1(x,x')$ and $d(x,x')$ are the minima over $n\in\mathbb{Z}^3$ of $|x-x'-Rn|_1$ and of $|x-x'-Rn|$, respectively, and $|z|_1\ge|z|$ for $z\in\mathbb{R}^3$, the $\ell^1$ norm dominating the sup norm; so $d_1(x,x')\ge d(x,x')$. Together with the identity for $S_y(0)(x,x')$ above, it follows that
\[
\left|S_y(0)(x,x')\right|\le\left\|T_{\phi}^{-1}\right\|_2\cdot\exp\left(-a_z\cdot d(x,x')\right)\le C_{00}\cdot\left(1+|y|\right)^{-1}\cdot\exp\left(-a_z\cdot d(x,x')\right).
\]
\end{proof}

\paragraph{The constants $C_z$, $\gamma$ and $a$}
The numbers
\begin{multline}\label{eq:operators-fermion-4}
C_z:=\max\left\{C_5,C_{00}\right\},\qquad\gamma_z:=\min\left\{\boldsymbol{\gamma}_5,a_z\right\},\qquad\gamma:=\gamma^F:=\frac{\gamma_z}{2},\\
a:=C_z\cdot c(\gamma),\qquad C_V=12\mathbf{e}\cdot(2+b)
\end{multline}
depend only on $b$, $\alpha$ and $L$. Here $C_5$ and $\boldsymbol{\gamma}_5$ are the numbers of Proposition~\ref{prop:zero-field-levels}, $C_{00}$ and $a_z$ are those of Lemma~\ref{lem:zero-field-level}, $c(\cdot)$ is the function of Section~\ref{sec:norms-exponential-weights}, and $C_V$ is the constant of Lemma~\ref{lem:perturbations-commutators-holder}. Since $0<\gamma\le a_z/2\le1/(6L)$, one has
\[
0<\gamma L\le\frac{1}{6}<\log2.
\]

Let $\mathbb{T}$ be a block torus of level $k\ge0$, let $0\le\mathfrak{m}\le1$ and let $y\in Y_k$. Then
\begin{equation}\label{eq:operators-fermion-5}
N_{\gamma}\left(S_y(0)\right)\le a\cdot w_k(y),\qquad N_{\gamma}\left(\delta_{\alpha,0}S_y(0)\right)\le a.
\end{equation}
To see this, let first $k\ge1$, so that $Y_k=\mathbb{R}\cup\{\infty\}$ and $w_k(y)=1$, and let $y_1$ and $y_2$ be unit sites. For a spinor $f$, the spinor $\mathbf{1}_{\Delta_{y_2}}f$ is supported in the orbit of $\Delta_{y_2}$, and $|\mathbf{1}_{\Delta_{y_2}}f|_{\infty}\le|f|_{\infty}$. For a spinor $v$ and a short pair $(x,\tilde{x})$ one has $0<|x-\tilde{x}|<1$ and, as $U_0=1$,
\[
\left(\delta_{\alpha,0}v\right)(x,\tilde{x})=|x-\tilde{x}|^{-\alpha}\left(v(\tilde{x})-v(x)\right).
\]
The moduli $|v(x)|$ and $|v(\tilde{x})-v(x)|$ do not change when $x$ and $\tilde{x}$ are translated by one and the same element of $R\mathbb{Z}^3$. Hence, by Proposition~\ref{prop:zero-field-levels}, each of the two operators $\mathbf{1}_{\Delta_{y_1}}S_y(0)\mathbf{1}_{\Delta_{y_2}}$ and $\mathbf{1}_{\Delta_{y_1}}\delta_{\alpha,0}S_y(0)\mathbf{1}_{\Delta_{y_2}}$ has norm at most $C_5\cdot\exp\left(-\boldsymbol{\gamma}_5\cdot d(y_1,y_2)\right)$. As $\gamma\le\boldsymbol{\gamma}_5/2$,
\[
\exp\left(\gamma\cdot d(y_1,y_2)\right)\cdot\exp\left(-\boldsymbol{\gamma}_5\cdot d(y_1,y_2)\right)\le\exp\left(-\gamma\cdot d(y_1,y_2)\right);
\]
summing over $y_2$ and using Section~\ref{sec:norms-exponential-weights} and $C_5\le C_z$, one obtains for both operators the bound $C_5\cdot c(\gamma)\le a$ for the sum that defines $N_{\gamma}$, uniformly in $y_1$. This proves~\eqref{eq:operators-fermion-5} for $k\ge1$. Let now $k=0$, so that $Y_0=\mathbb{R}$. Then the unit blocks are single sites, and for sites $x$ and $x'$ the norm of $\mathbf{1}_{\Delta_x}S_y(0)\mathbf{1}_{\Delta_{x'}}$ is at most $\left|S_y(0)(x,x')\right|$. By Lemma~\ref{lem:zero-field-level}, by $\gamma\le a_z/2$ and by $C_{00}\le C_z$, the same summation gives
\[
N_{\gamma}\left(S_y(0)\right)\le C_{00}\cdot\left(1+|y|\right)^{-1}\cdot c(\gamma)\le a\cdot w_0(y).
\]
At level $0$ there are no short pairs, so that $\delta_{\alpha,0}S_y(0)=0$.

\subsection{Inverses at small fields}\label{sec:inverses-small-fields}

\begin{proposition}[Small fields]\label{prop:small-fields}
Let $\mathbb{T}$ be a block torus of level $k\ge0$, let $0<\mathfrak{m}\le1$ and $y\in Y_k$, and let $B$ be a connection on $\mathbb{T}$ with $e_k|B|_{\infty}\le\boldsymbol{\delta}$, where the number $\boldsymbol{\delta}$ satisfies
\begin{equation}\label{eq:operators-fermion-6}
3L\cdot\boldsymbol{\delta}\le1,\qquad2C_V\cdot a\cdot\boldsymbol{\delta}\le1.
\end{equation}
Put $V:=F_y(B)-F_y(0)$.
\begin{enumerate}
\item[(a)] The operator $F_y(B)$ is bijective, and
\[
S_y(B)=S_y(0)\sum_{n\ge0}\left(-V S_y(0)\right)^n,
\]
where the series converges with respect to $N_{\gamma}$, uniformly over all connections $B$ on $\mathbb{T}$ with $e_k|B|_{\infty}\le\boldsymbol{\delta}$.
\item[(b)] $N_{\gamma}\left(S_y(B)\right)\le2a\cdot w_k(y)$.
\item[(c)] $N_{\gamma}\left(\delta_{\alpha,B}S_y(B)\right)\le18a$.
\end{enumerate}
\end{proposition}

\begin{proof}
Lemma~\ref{lem:perturbations-commutators-holder}(a) applies, since $0<\gamma L\le\frac{1}{6}<\log2$ and $3L\cdot\boldsymbol{\delta}\le1$; it gives $N_{\gamma}(V)\le C_V\cdot\boldsymbol{\delta}$. Since $0<\mathfrak{m}\le1$, the bound~\eqref{eq:operators-fermion-5} and the inequality $w_k\le1$, which follows from the definition~\eqref{eq:operators-fermion-2}, give $N_{\gamma}\left(S_y(0)\right)\le a$. With~\eqref{eq:operators-fermion-EW2} and~\eqref{eq:operators-fermion-6} it follows that
\[
N_{\gamma}\left(V S_y(0)\right)\le C_V\cdot\boldsymbol{\delta}\cdot a\le\frac{1}{2}.
\]
The space of spinors on $\mathbb{T}$ has finite dimension; $N_{\gamma}$ is a norm on its linear operators (by~\eqref{eq:operators-fermion-EW1}, $N_{\gamma}(O)=0$ only for $O=0$), submultiplicative by~\eqref{eq:operators-fermion-EW2}, and $N_{\gamma}(1)=1$ by~\eqref{eq:operators-fermion-EW3}. By the Neumann series, the series
\[
\Xi:=\sum_{n\ge0}\left(-V S_y(0)\right)^n
\]
converges in $N_{\gamma}$, its terms satisfying $N_{\gamma}\left(\left(-V S_y(0)\right)^n\right)\le2^{-n}$ for every connection $B$ as in the statement; moreover $\Xi=\left(1+V S_y(0)\right)^{-1}$ and $N_{\gamma}(\Xi)\le2$. The operator $F_y(0)$ is bijective by Proposition~\ref{prop:zero-field-levels} and Lemma~\ref{lem:zero-field-level}, and
\[
F_y(B)=F_y(0)+V=\left(1+V S_y(0)\right)F_y(0).
\]
Hence $F_y(B)$ is bijective, with inverse $S_y(B)=S_y(0)\,\Xi$; this is (a). By~\eqref{eq:operators-fermion-EW2}, $N_{\gamma}\left(S_y(B)\right)\le N_{\gamma}\left(S_y(0)\right)\cdot N_{\gamma}(\Xi)\le2a\cdot w_k(y)$, which is (b).

(c) By~\eqref{eq:operators-fermion-EW2} and~\eqref{eq:operators-fermion-5},
\[
N_{\gamma}\left(\delta_{\alpha,0}S_y(B)\right)=N_{\gamma}\left(\left(\delta_{\alpha,0}S_y(0)\right)\Xi\right)\le a\cdot2=2a.
\]
By Lemma~\ref{lem:perturbations-commutators-holder}(c),
\[
\delta_{\alpha,B}S_y(B)-\delta_{\alpha,0}S_y(B)=E_gS_y(B),\qquad g(x,x'):=|x-x'|^{-\alpha}\left(U_B(x,x')-1\right),
\]
and by Lemma~\ref{lem:bounds-phases}(b), for every short pair $(x,x')$,
\[
|g(x,x')|\le3\boldsymbol{\delta}\cdot\exp(3\boldsymbol{\delta})\cdot|x-x'|^{1-\alpha}\le L^{-1}\exp\left(\frac{1}{L}\right)\le\frac{1}{4};
\]
here $|x-x'|<1$ and $\alpha<1$ were used, together with $3\boldsymbol{\delta}\le1/L$ (by~\eqref{eq:operators-fermion-6}), $L\ge5$ and the monotonicity of the function $s\mapsto s\cdot\exp(s)$ on $[0,\infty)$. As $\gamma\le1/(6L)\le1/30$, one has $27\exp(\gamma)/4\le8$. Lemma~\ref{lem:perturbations-commutators-holder}(c), applied with $O=S_y(B)$ and $g_0=\frac{1}{4}$, together with (b) and $w_k\le1$, gives
\[
N_{\gamma}\left(E_gS_y(B)\right)\le8\cdot N_{\gamma}\left(S_y(B)\right)\le8\cdot2a=16a.
\]
By~\eqref{eq:operators-fermion-EW2}, $N_{\gamma}\left(\delta_{\alpha,B}S_y(B)\right)\le2a+16a=18a$.
\end{proof}

For $\boldsymbol{\delta}=2e_k^{\varepsilon}$ the two conditions~\eqref{eq:operators-fermion-6} read $6L\cdot e_k^{\varepsilon}\le1$ and $4C_V\cdot a\cdot e_k^{\varepsilon}\le1$; these are conditions~\eqref{eq:renormalisation-group-20} and~\eqref{eq:renormalisation-group-21} with $e_k$ in place of $e_0$.

\subsection{Local inverses on mirror tori}\label{sec:local-inverses-mirror}

\subsubsection{Folding a torus onto a cube}\label{sec:folding-torus-onto}

Let $\square$ be an $M$-cube of $T_k$, and let $\mathsf{c}\in M\mathbb{Z}^3$ be a lift of its centre. Put
\begin{gather*}
\tilde{W}:=\left\{x\in\eta\mathbb{Z}^3:|x-\mathsf{c}|<\frac{3M}{2}\right\},\qquad W^{-}:=\left\{x\in\tilde{W}:|x-\mathsf{c}|<\frac{3M}{2}-\eta\right\},\\
o:=\mathsf{c}-\frac{3M}{2}(1,1,1).
\end{gather*}
The set $\tilde{W}$ is a lift to $\eta\mathbb{Z}^3$ of $W:=\square^{\sim 1}$, the enlargement of $\square$ by one layer of $M$-cubes (Section~\ref{sec:lattice-geometry}), and the point $o$ is a vertex of the closed cube $\left\{z\in\mathbb{R}^3:|z-\mathsf{c}|\le\frac{3M}{2}\right\}$. A bond of $\eta\mathbb{Z}^3$ is a \emph{bond of $\tilde{W}$} if both of its end points lie in $\tilde{W}$. A site $x\in\tilde{W}$ belongs to $W^{-}$ if and only if all six nearest neighbours of $x$ belong to $\tilde{W}$.

\begin{definition}[Mirror torus, fold map, chambers]\label{def:mirror-torus-fold}
Consider the triple
\[
\mathbb{T}^W:=\left(\frac{\eta}{2}(1,1,1),\,6M,\,+1\right).
\]
Since $6M$ is an integer multiple of $L$ and $6M\ge3L$, it is a block torus of level $k$ in the sense of Definition~\ref{def:block-torus}, called the \emph{mirror torus}. Its lattice is $\Lambda^W:=\frac{\eta}{2}(1,1,1)+\eta\mathbb{Z}^3$, and its lattices $\Lambda_j$ of Definition~\ref{def:block-torus} are denoted by
\[
\Lambda^W_j:=\frac{L^j\eta}{2}(1,1,1)+L^j\eta\mathbb{Z}^3\qquad(0\le j\le k+1).
\]
Let $\mathrm{fo}:\mathbb{R}\to[0,3M]$ be the function $\mathrm{fo}(t):=\operatorname{dist}(t,6M\mathbb{Z})$. The \emph{fold map} $\mathrm{Fo}:\Lambda^W\to\mathbb{R}^3$ and the map $\mathrm{emb}:\eta\mathbb{Z}^3\to\mathbb{R}^3/6M\mathbb{Z}^3$ are defined by
\[
\mathrm{Fo}(u):=o+\left(\mathrm{fo}(u_0),\mathrm{fo}(u_1),\mathrm{fo}(u_2)\right),\qquad\mathrm{emb}(x):=x-o\quad\text{modulo }6M\mathbb{Z}^3.
\]
A \emph{mirror plane} is a plane $\left\{u\in\mathbb{R}^3:u_{\mu}=3Mn\right\}$ with $\mu\in\{0,1,2\}$ and $n\in\mathbb{Z}$. The mirror planes divide $\mathbb{T}^W$ into eight \emph{chambers}, indexed by $s\in\{0,1\}^3$; the chamber of index $s$ is
\[
\left\{u\in\Lambda^W:u_{\mu}\in\left(3Ms_{\mu},3M(s_{\mu}+1)\right)+6M\mathbb{Z}\ \text{ for }\mu=0,1,2\right\}.
\]
Each chamber is invariant under the translations by $6M\mathbb{Z}^3$ and is regarded as a set of sites of $\mathbb{T}^W$. The chamber of index $(0,0,0)$,
\[
\mathrm{Ch}:=\left\{u\in\Lambda^W:u_{\mu}\in(0,3M)+6M\mathbb{Z}\ \text{ for }\mu=0,1,2\right\},
\]
is called \emph{the chamber}.
\end{definition}

\begin{lemma}[Fold map and embedding]\label{lem:fold-map-embedding}
\begin{enumerate}
\item[(a)] The number $3M$ is an odd multiple of $L^j\eta$ for each $j\in\{0,\dots,k+1\}$. It follows that the mirror planes contain no site of $\Lambda^W$, and that every block of $\mathbb{T}^W$ lies inside one chamber. Restricted to $\tilde{W}$, the map $\mathrm{emb}$ is a bijection $\tilde{W}\to\mathrm{Ch}$ which sends any two nearest neighbours to two nearest neighbours; it carries each block of $T_k$ lying in $\tilde{W}$ onto a block of $\mathbb{T}^W$ lying in $\mathrm{Ch}$, with the centre going to the centre, and each block of $\mathbb{T}^W$ lying in $\mathrm{Ch}$ is obtained in this manner. Finally, the blocks of side $L$ of $T_k$ that are contained in $\tilde{W}$ cover $\tilde{W}$.
\item[(b)] For all $u\in\Lambda^W$ and $n\in\mathbb{Z}^3$ one has $\mathrm{Fo}(u+6Mn)=\mathrm{Fo}(u)$; moreover $\mathrm{Fo}\left(\Lambda^W\right)=\tilde{W}$, and $\mathrm{Fo}\left(\mathrm{emb}(x)\right)=x$ for every $x\in\tilde{W}$. Let $u,u'$ be nearest neighbours in $\Lambda^W$. Either some chamber contains both of them, and then $\mathrm{Fo}(u)$ and $\mathrm{Fo}(u')$ are nearest neighbours in $\tilde{W}$; or else $u$ and $u'$ lie on opposite sides of a mirror plane, and then $\mathrm{Fo}(u)=\mathrm{Fo}(u')$.
\item[(c)] Let $y,y'$ be unit sites in $\tilde{W}$. On $T_k$ the distance of $y$ and $y'$ equals $|y-y'|$, and on $\mathbb{T}^W$ the distance of $\mathrm{emb}(y)$ and $\mathrm{emb}(y')$ equals $|y-y'|$ as well.
\end{enumerate}
\end{lemma}

\begin{proof}
(a) As $M=L^m$ and $\eta=L^{-k}$,
\[
\frac{3M}{L^j\eta}=3L^{m+k-j},
\]
and $m+k-j\ge m-1\ge0$ for $j\le k+1$; since $L$ is odd, this is an odd integer. It follows that $\frac{3M}{2}\in L^j\eta\left(\mathbb{Z}+\frac{1}{2}\right)$. Moreover $\mathsf{c}\in M\mathbb{Z}^3\subset L^j\eta\mathbb{Z}^3$, because $\frac{M}{L^j\eta}=L^{m+k-j}$ is an integer. Hence, for $0\le j\le k+1$, the translation
\[
x\mapsto x-o=x-\mathsf{c}+\frac{3M}{2}(1,1,1)
\]
maps $L^j\eta\mathbb{Z}^3$ onto $\frac{L^j\eta}{2}(1,1,1)+L^j\eta\mathbb{Z}^3=\Lambda^W_j$; for $j=0$ it maps $\eta\mathbb{Z}^3$ onto $\Lambda^W$. The blocks of a block torus are defined through the norm $|\cdot|$ (Definition~\ref{def:block-torus}), which is invariant under translations. Therefore this translation maps every block of side $L^j\eta$ of $T_k$, lifted to $\eta\mathbb{Z}^3$, onto a block of side $L^j\eta$ of $\mathbb{T}^W$, and the centre of the former onto the centre of the latter.

By Definition~\ref{def:block-torus}, a block of side $L^j\eta$ of $\mathbb{T}^W$ has a centre $z$ with $z_{\mu}=L^j\eta\left(n_{\mu}+\frac{1}{2}\right)$ for some $n\in\mathbb{Z}^3$, and it consists of the sites $u\in\Lambda^W$ with $L^j\eta\,n_{\mu}<u_{\mu}<L^j\eta\,(n_{\mu}+1)$ for $\mu=0,1,2$. Since $3M$ is an integer multiple of $L^j\eta$, none of these open intervals contains a multiple of $3M$. Hence every block lies in a single chamber; for $j=0$ the blocks are single sites, so that no site lies on a mirror plane.

A site $x\in\eta\mathbb{Z}^3$ belongs to $\tilde{W}$ if and only if $x-o\in(0,3M)^3$, and every class of $\mathrm{Ch}$ modulo $6M\mathbb{Z}^3$ has exactly one representative in $(0,3M)^3$. Therefore $\mathrm{emb}$ is a bijection of $\tilde{W}$ onto $\mathrm{Ch}$. Being induced by a translation, it maps nearest neighbours to nearest neighbours. As shown above, it maps the blocks of $T_k$ contained in $\tilde{W}$ onto blocks of $\mathbb{T}^W$ contained in $\mathrm{Ch}$, centres to centres; and every block of $\mathbb{T}^W$ contained in $\mathrm{Ch}$ is obtained in this way. Indeed, each of the three open intervals that define such a block lies between two consecutive multiples of $3M$, as shown above; as the block is contained in $\mathrm{Ch}$, it has a translate by an element of $6M\mathbb{Z}^3$ inside $(0,3M)^3$, and the preimage of this translate under the translation $x\mapsto x-o$ is a block of $T_k$ contained in $\tilde{W}$. Finally, the blocks of side $L=L^{k+1}\eta$ of $\mathbb{T}^W$ partition $\Lambda^W$ (Definition~\ref{def:block-torus}), and each of them lies in a single chamber; so $\mathrm{Ch}$ is the union of those among them that are contained in $\mathrm{Ch}$, and by what precedes $\tilde{W}$ is the union of the corresponding blocks of side $L$ of $T_k$.

(b) The function $\mathrm{fo}$ is even and $6M$-periodic, $\mathrm{fo}(t)=t$ for $0\le t\le3M$, and on each interval $[3Mn,3M(n+1)]$, $n\in\mathbb{Z}$, it is affine with slope $1$ or $-1$; in particular $|\mathrm{fo}(t)-\mathrm{fo}(t')|=|t-t'|$ for $t,t'$ in one such interval. Evenness and periodicity give
\[
\mathrm{fo}(3Mn+s)=\mathrm{fo}(-3Mn-s)=\mathrm{fo}(3Mn-s)\qquad(n\in\mathbb{Z},\ s\in\mathbb{R}).
\]
The periodicity of $\mathrm{fo}$ implies the invariance of $\mathrm{Fo}$ under $6M\mathbb{Z}^3$. Let $u\in\Lambda^W$ and $\mu\in\{0,1,2\}$. Then $u_{\mu}\in\eta\left(\mathbb{Z}+\frac{1}{2}\right)$, whereas $3M\mathbb{Z}\subset\eta\mathbb{Z}$ by (a); as $\mathrm{fo}(u_{\mu})=|u_{\mu}-6Mn|$ for a suitable $n\in\mathbb{Z}$, one obtains $\mathrm{fo}(u_{\mu})\in\eta\left(\mathbb{Z}+\frac{1}{2}\right)\cap(0,3M)$. Since $\mathsf{c}_{\mu}-\frac{3M}{2}\in\eta\left(\mathbb{Z}+\frac{1}{2}\right)$ by (a), the coordinate $\mathrm{Fo}(u)_{\mu}=\mathsf{c}_{\mu}-\frac{3M}{2}+\mathrm{fo}(u_{\mu})$ lies in $\eta\mathbb{Z}$ and satisfies $\left|\mathrm{Fo}(u)_{\mu}-\mathsf{c}_{\mu}\right|<\frac{3M}{2}$. Thus $\mathrm{Fo}(u)\in\tilde{W}$. If $x\in\tilde{W}$, the representative $x-o$ of $\mathrm{emb}(x)$ lies in $(0,3M)^3$, where $\mathrm{fo}$ acts as the identity on each coordinate; hence $\mathrm{Fo}\left(\mathrm{emb}(x)\right)=o+(x-o)=x$. In particular $\mathrm{Fo}\left(\Lambda^W\right)=\tilde{W}$.

Let $u'=u\pm\eta e_{\mu}$ be a nearest neighbour of $u$; the two other coordinates of $u$ and $u'$ agree. If no multiple of $3M$ lies between $u_{\mu}$ and $u'_{\mu}$, then $u$ and $u'$ lie in one chamber, $u_{\mu}$ and $u'_{\mu}$ lie in one interval $[3Mn,3M(n+1)]$, and $\left|\mathrm{fo}(u'_{\mu})-\mathrm{fo}(u_{\mu})\right|=\eta$; so $\mathrm{Fo}(u)$ and $\mathrm{Fo}(u')$ are nearest neighbours in $\tilde{W}$. Otherwise a mirror plane $\left\{u_{\mu}=3Mn\right\}$ separates $u$ from $u'$. As $u_{\mu}$ and $u'_{\mu}$ belong to $\eta\left(\mathbb{Z}+\frac{1}{2}\right)$ and differ by $\eta$, while $3Mn\in\eta\mathbb{Z}$, one has $\left\{u_{\mu},u'_{\mu}\right\}=\left\{3Mn-\frac{\eta}{2},3Mn+\frac{\eta}{2}\right\}$, and the symmetry of $\mathrm{fo}$ about $3Mn$ gives $\mathrm{Fo}(u)=\mathrm{Fo}(u')$.

(c) Both $y$ and $y'$ lie within the distance $\frac{3M}{2}$ of $\mathsf{c}$, so that $|y-y'|<3M$. The difference of $\mathrm{emb}(y)$ and $\mathrm{emb}(y')$ is the class of $y-y'$ modulo $6M\mathbb{Z}^3$, and for $n\in\mathbb{Z}^3\setminus\{0\}$
\[
|y-y'-6Mn|\ge6M-|y-y'|>|y-y'|.
\]
Hence the distance of $\mathrm{emb}(y)$ and $\mathrm{emb}(y')$ on $\mathbb{T}^W$, the minimum of $|y-y'-6Mn|$ over $n\in\mathbb{Z}^3$, is $|y-y'|$: the difference is shorter than half the period. As $l\in\mathrm{Cl}$, one has $l\ge13LM$, so that $|y-y'|<3M<\frac{l}{2}$; hence $|y-y'-ln|\ge l-|y-y'|>|y-y'|$ for $n\in\mathbb{Z}^3\setminus\{0\}$, and the distance of $y$ and $y'$ on $T_k$, the minimum of $|y-y'-ln|$ over $n\in\mathbb{Z}^3$, is $|y-y'|$ as well.
\end{proof}

\begin{definition}[Folded fields and local inverse]\label{def:folded-fields-local}
Let $A$ be a complex antisymmetric function on the bonds of $\tilde{W}$, and let $\boldsymbol{\lambda}$ be a function on $\tilde{W}$; as in Definition~\ref{def:block-torus}, $d\boldsymbol{\lambda}$ denotes the antisymmetric function on the bonds of $\tilde{W}$ given by
\[
(d\boldsymbol{\lambda})(x,x'):=\frac{\boldsymbol{\lambda}(x')-\boldsymbol{\lambda}(x)}{\eta}.
\]
The functions $A^W$ on the bonds of $\mathbb{T}^W$ and $\boldsymbol{\lambda}^W$ on $\Lambda^W$ are
\[
A^W(u,u'):=\begin{cases}A\left(\mathrm{Fo}(u),\mathrm{Fo}(u')\right)&\text{if }\mathrm{Fo}(u)\ne\mathrm{Fo}(u'),\\0&\text{if }\mathrm{Fo}(u)=\mathrm{Fo}(u'),\end{cases}\qquad\boldsymbol{\lambda}^W:=\boldsymbol{\lambda}\circ\mathrm{Fo};
\]
by Lemma~\ref{lem:fold-map-embedding}(b), $\left(\mathrm{Fo}(u),\mathrm{Fo}(u')\right)$ is a bond of $\tilde{W}$ whenever $\mathrm{Fo}(u)\ne\mathrm{Fo}(u')$. For a function $f:\tilde{W}\to\mathbb{C}^4$ and a spinor $v$ on $\mathbb{T}^W$ let
\[
(\operatorname{ext} f)(u):=\begin{cases}f\left(\mathrm{Fo}(u)\right)&(u\in\mathrm{Ch}),\\0&(u\in\Lambda^W\setminus\mathrm{Ch}),\end{cases}\qquad(\operatorname{rs} v)(x):=v\left(\mathrm{emb}(x)\right)\quad(x\in\tilde{W}).
\]
For $y\in Y_k$ let $F^W_y(A):=F_y(A^W)$ be the operator of Section~\ref{sec:operators-2} on the spinors of $\mathbb{T}^W$, formed with the connection $A^W$, and, whenever $F^W_y(A)$ is invertible, let
\[
\mathrm{Sl}_{\square,y}(A):=\operatorname{rs}{}\cdot F^W_y(A)^{-1}\cdot\operatorname{ext} .
\]
This operator on the functions $\tilde{W}\to\mathbb{C}^4$ is called the \emph{local inverse}.
\end{definition}

Since $\mathrm{Fo}$ and $\mathrm{Ch}$ are invariant under the translations by $6M\mathbb{Z}^3$ (Lemma~\ref{lem:fold-map-embedding}(b) and Definition~\ref{def:mirror-torus-fold}), the function $A^W$, which is antisymmetric, is a connection on $\mathbb{T}^W$, the function $\boldsymbol{\lambda}^W$ is a scalar on $\mathbb{T}^W$, and $\operatorname{ext} f$ is a spinor on $\mathbb{T}^W$. The map $A\mapsto A^W$ is complex linear, and $\left(\overline{A}\right)^W=\overline{A^W}$. Moreover $|A^W|_{\infty}=|A|_{\infty}$, the supremum on the right being taken over the bonds of $\tilde{W}$: each value of $A^W$ is $0$ or a value of $A$, and for every bond $(z,z')$ of $\tilde{W}$ the pair $\left(\mathrm{emb}(z),\mathrm{emb}(z')\right)$ is a bond inside $\mathrm{Ch}$ at which $A^W$ takes the value $A(z,z')$ (Lemma~\ref{lem:fold-map-embedding}(a), (b)). Furthermore
\[
(d\boldsymbol{\lambda})^W=d\left(\boldsymbol{\lambda}^W\right).
\]
Indeed, let $(u,u')$ be a bond of $\Lambda^W$. If one chamber contains $u$ and $u'$, then $\mathrm{Fo}(u)\ne\mathrm{Fo}(u')$ by Lemma~\ref{lem:fold-map-embedding}(b), and at $(u,u')$ both sides take the value $\eta^{-1}\left(\boldsymbol{\lambda}\left(\mathrm{Fo}(u')\right)-\boldsymbol{\lambda}\left(\mathrm{Fo}(u)\right)\right)$. Otherwise $u$ and $u'$ lie on opposite sides of a mirror plane, so that $\mathrm{Fo}(u)=\mathrm{Fo}(u')$ by the same lemma, and both sides are $0$ at $(u,u')$. Since $\mathrm{Fo}\left(\mathrm{emb}(x)\right)=x$ for $x\in\tilde{W}$ and $\mathrm{emb}(\tilde{W})=\mathrm{Ch}$, one has $\operatorname{rs}{}\cdot\operatorname{ext}=1$. Let the functions on $\tilde{W}$ be paired by $(f,g):=\sum_{x\in\tilde{W}}\eta^3f(x)\cdot g(x)$ and the spinors on $\mathbb{T}^W$ by the pairing of Section~\ref{sec:norms-pairings-kernels}, and let transposes of operators between these spaces be taken with respect to these pairings. As $\mathrm{emb}$ is a bijection of $\tilde{W}$ onto $\mathrm{Ch}$ and $\operatorname{ext} f$ vanishes off $\mathrm{Ch}$,
\[
(v,\operatorname{ext} f)=\sum_{x\in\tilde{W}}\eta^3v\left(\mathrm{emb}(x)\right)\cdot f(x)=(\operatorname{rs} v,f),
\]
that is, $\operatorname{ext}^T=\operatorname{rs}$.

For $\boldsymbol{\delta}>0$ let $G^W(\boldsymbol{\delta})$ be the set of the complex antisymmetric functions $A$ on the bonds of $\tilde{W}$ that admit a decomposition
\[
A=B+d\boldsymbol{\lambda},\qquad\boldsymbol{\lambda}:\tilde{W}\to\mathbb{R},\qquad e_k|B|_{\infty}<\boldsymbol{\delta},
\]
with a complex antisymmetric function $B$ on the bonds of $\tilde{W}$. Let $A=B+d\boldsymbol{\lambda}$ be such a decomposition. Every complex antisymmetric function $A'$ on the bonds of $\tilde{W}$ with $e_k|A'|_{\infty}<\boldsymbol{\delta}-e_k|B|_{\infty}$ satisfies $A+A'=(B+A')+d\boldsymbol{\lambda}$ and $e_k|B+A'|_{\infty}<\boldsymbol{\delta}$, whence $A+A'\in G^W(\boldsymbol{\delta})$. Thus each element of $G^W(\boldsymbol{\delta})$ has a neighbourhood contained in $G^W(\boldsymbol{\delta})$; that is, $G^W(\boldsymbol{\delta})$ is an open subset of the vector space of all complex antisymmetric functions on the bonds of $\tilde{W}$.

\begin{proposition}[Properties of the local inverse]\label{prop:properties-local-inverse}
Let the number $\boldsymbol{\delta}>0$ satisfy~\eqref{eq:operators-fermion-6}, let $0<\mathfrak{m}\le1$ and $y\in Y_k$, and let $A\in G^W(\boldsymbol{\delta})$. As in Lemma~\ref{lem:symmetries-elementary-operators}, $-\infty:=\infty$.
\begin{enumerate}
\item[(a)] The operator $F^W_y(A)$ is bijective. The map $A\mapsto\mathrm{Sl}_{\square,y}(A)$ is analytic on $G^W(\boldsymbol{\delta})$, and the map $(y,A)\mapsto\mathrm{Sl}_{\square,y}(A)$ is continuous on $\mathbb{R}\times G^W(\boldsymbol{\delta})$.
\item[(b)] For an operator $O$ on the functions $\tilde{W}\to\mathbb{C}^4$, with values in the functions on $\tilde{W}$ or in the functions on the short pairs in $\tilde{W}$, put
\[
N_{\gamma}(O):=\sup_{y_1}\sum_{y_2}\exp\left(\gamma|y_1-y_2|\right)\cdot\left\|\mathbf{1}_{\Delta_{y_1}}O\mathbf{1}_{\Delta_{y_2}}\right\|,
\]
where $y_1$ and $y_2$ range over the unit sites in $\tilde{W}$; the blocks $\Delta_{y_1}$, $\Delta_{y_2}$ are contained in $\tilde{W}$, since by Lemma~\ref{lem:fold-map-embedding}(a) each of the sites $y_1$, $y_2$ lies in a block of side $L$ contained in $\tilde{W}$, and by Definition~\ref{def:block-torus} a block of side $L$ is the union of the unit blocks whose centres it contains. Then
\[
N_{\gamma}\left(\mathrm{Sl}_{\square,y}(A)\right)\le2a\cdot w_k(y),\qquad N_{\gamma}\left(\delta_{\alpha,A}\mathrm{Sl}_{\square,y}(A)\right)\le18a,
\]
where in the second bound $\delta_{\alpha,A}$ is taken on the short pairs $(x,x')$ in $\tilde{W}$; for such a pair the six paths $\Gamma^{\pi}(x,x')$ lie in $\tilde{W}$, so that $U_A(x,x')$ is defined.
\item[(c)] Continue $A$ in an arbitrary way to a complex antisymmetric function on the bonds of $\eta\mathbb{Z}^3$, and let $F_y(A)$ act on the functions $\eta\mathbb{Z}^3\to\mathbb{C}^4$ by the formulas of Sections~\ref{sec:phase-factors-lattice} and~\ref{sec:operators-2}, formed with this continuation and with the blocks of $T_k$ lifted to $\eta\mathbb{Z}^3$. For $x\in W^{-}$ the value $\left(F_y(A)v\right)(x)$ depends only on the restriction of $v$ to $\tilde{W}$ and on the values of $A$ on the bonds of $\tilde{W}$. If $f$ is a function on $\tilde{W}$ and $v$ is the function that equals $\mathrm{Sl}_{\square,y}(A)f$ on $\tilde{W}$ and vanishes elsewhere, then
\[
\left(F_y(A)v\right)(x)=f(x)\qquad\text{for every }x\in W^{-}.
\]
If $g$ is a function on $\eta\mathbb{Z}^3$ that vanishes outside $W^{-}$, then $F_y(A)g$ vanishes outside $\tilde{W}$, and $\mathrm{Sl}_{\square,y}(A)\left(F_y(A)g\right)=g$ on $\tilde{W}$.
\item[(d)] For every real function $\boldsymbol{\lambda}$ on $\tilde{W}$,
\[
\mathrm{Sl}_{\square,y}(A-d\boldsymbol{\lambda})=\exp\left(ie_k\cdot\boldsymbol{\lambda}\right)\cdot\mathrm{Sl}_{\square,y}(A)\cdot\exp\left(-ie_k\cdot\boldsymbol{\lambda}\right).
\]
\item[(e)] Let $g=(S,a')\in\mathbb{G}_k$ with $a'\in M\mathbb{Z}^3$, let $gx:=r(S)x+a'$, and let
\[
y':=\begin{cases}y&\text{if }S\in\mathrm{Sp}_3\text{ or if }y\text{ is neutral},\\-y&\text{if }S\notin\mathrm{Sp}_3,\end{cases}
\]
as in Lemma~\ref{lem:symmetries-elementary-operators}(b). Form the local inverse of the cube $g\square$ with the lift $r(S)\mathsf{c}+a'$ of its centre; the corresponding lift of $(g\square)^{\sim 1}$ is $r(S)\tilde{W}+a'$. The element $g$ maps a function $f$ on $\tilde{W}$ to the function $(gf)(x):=S\cdot f\left(g^{-1}x\right)$ on $r(S)\tilde{W}+a'$, and $A$ to the function $(gA)(x,x'):=A\left(g^{-1}x,g^{-1}x'\right)$ on the bonds of $r(S)\tilde{W}+a'$. Then
\[
\mathrm{Sl}_{g\square,y'}(gA)\cdot g=g\cdot\mathrm{Sl}_{\square,y}(A).
\]
\item[(f)] For an operator $O$ on the functions $\tilde{W}\to\mathbb{C}^4$ let $\overline{O}$ be the operator $f\mapsto\overline{O\overline{f}}$, as in Lemma~\ref{lem:symmetries-elementary-operators}(d). Then
\[
\mathrm{Sl}_{\square,y}(A)^T=\mathsf{C}^{-1}\cdot\mathrm{Sl}_{\square,-y}(-A)\cdot\mathsf{C},\qquad\overline{\mathrm{Sl}_{\square,y}\left(\overline{A}\right)}=\varkappa^{-1}\cdot\mathrm{Sl}_{\square,y}(-A)\cdot\varkappa.
\]
\item[(g)] Suppose that $A$ is the restriction to the bonds of $\tilde{W}$ of a complex antisymmetric function $\tilde{A}$ on the bonds of $\eta\mathbb{Z}^3$ that is invariant under the translations by $l\mathbb{Z}^3$. Let $n\in\mathbb{Z}^3$; use the lift $\mathsf{c}+ln$ of the centre of $\square$ in the place of $\mathsf{c}$, so that $\tilde{W}+ln$ takes the place of $\tilde{W}$, and let $A_n$ be the restriction of $\tilde{A}$ to the bonds of $\tilde{W}+ln$. Then the local inverse $\mathrm{Sl}_{\square,y}(A_n)$ formed with the lift $\mathsf{c}+ln$ is the translate of $\mathrm{Sl}_{\square,y}(A)$ by $ln$: it sends a function $f$ on $\tilde{W}+ln$ to the function $x\mapsto\left(\mathrm{Sl}_{\square,y}(A)\,f(\cdot+ln)\right)(x-ln)$.
\end{enumerate}
\end{proposition}

\begin{proof}
In the whole proof, $A=B+d\boldsymbol{\lambda}_0$, with a real function $\boldsymbol{\lambda}_0$ on $\tilde{W}$, is a decomposition as in the definition of $G^W(\boldsymbol{\delta})$.

\emph{(a) and (b).} By the linearity of the fold and by $(d\boldsymbol{\lambda}_0)^W=d\left(\boldsymbol{\lambda}_0^W\right)$,
\[
A^W=B^W+d\left(\boldsymbol{\lambda}_0^W\right),
\]
where $\boldsymbol{\lambda}_0^W$ is a real scalar on $\mathbb{T}^W$ and $e_k|B^W|_{\infty}=e_k|B|_{\infty}<\boldsymbol{\delta}$. Apply Lemma~\ref{lem:symmetries-elementary-operators}(a) on the block torus $\mathbb{T}^W$ to the connection $B^W$, to the closed real connection $d\left(-\boldsymbol{\lambda}_0^W\right)$ and to the function $\exp\left(-ie_k\cdot\boldsymbol{\lambda}_0^W\right)$ on $\Lambda^W$; these have the properties required there, by the example given in that lemma with the real scalar $-\boldsymbol{\lambda}_0^W$. It gives
\[
F^W_y(A)=\exp\left(-ie_k\cdot\boldsymbol{\lambda}_0^W\right)\cdot F_y(B^W)\cdot\exp\left(ie_k\cdot\boldsymbol{\lambda}_0^W\right)
\]
and, for every spinor $v$ on $\mathbb{T}^W$ and every short pair $(u,u')$ of $\mathbb{T}^W$,
\[
\left(\delta_{\alpha,A^W}\left(\exp\left(-ie_k\cdot\boldsymbol{\lambda}_0^W\right)\cdot v\right)\right)(u,u')=\exp\left(-ie_k\cdot\boldsymbol{\lambda}_0^W(u)\right)\cdot\left(\delta_{\alpha,B^W}v\right)(u,u').
\]
By Proposition~\ref{prop:small-fields}(a), applied on $\mathbb{T}^W$ to the connection $B^W$, the operator $F_y(B^W)$ is bijective. Hence $F^W_y(A)$ is bijective, and
\[
F^W_y(A)^{-1}=\exp\left(-ie_k\cdot\boldsymbol{\lambda}_0^W\right)\cdot S_y(B^W)\cdot\exp\left(ie_k\cdot\boldsymbol{\lambda}_0^W\right).
\]

Each value of $A^W$ is $0$ or one of the finitely many numbers $A(b)$, $b$ a bond of $\tilde{W}$, and the phases of Section~\ref{sec:phase-factors-lattice} formed with $A^W$ are exponentials of linear functions of these numbers. The parameter $y$ enters $F_y$ only through the matrices $\boldsymbol{\alpha}_y$, $\boldsymbol{\beta}_y$ if $k\ge1$, among which $\boldsymbol{\alpha}_{\infty}=b_k$ and $\boldsymbol{\beta}_{\infty}=0$, and through $i\gamma_3y$ if $k=0$ (Section~\ref{sec:operators-2}); these matrices do not depend on $A$, and for $y\in\mathbb{R}$ they are continuous functions of $y$. Consider the matrix of $F^W_y(A)$ with respect to the sites of a period cell of $\Lambda^W$ and the four components. For every $y\in Y_k$, each entry of this matrix is a finite sum of products of a number that depends on $y$ alone and an entire function of the numbers $A(b)$; hence it is an entire function of these numbers, and for $y\in\mathbb{R}$ it is a continuous function of $y$ and of the numbers $A(b)$ jointly. By Cramer's rule the entries of the inverse of an invertible matrix are polynomials in its entries divided by its determinant; they are therefore analytic, and in particular continuous, functions of the entries on the open set where the determinant does not vanish. By what precedes, $F^W_y(A)$ is invertible for all $y\in Y_k$ and $A\in G^W(\boldsymbol{\delta})$; the maps $\operatorname{rs}$ and $\operatorname{ext}$ depend neither on $y$ nor on $A$; and $G^W(\boldsymbol{\delta})$ is open. It follows that, for every $y\in Y_k$, the map $A\mapsto\mathrm{Sl}_{\square,y}(A)$ is analytic on $G^W(\boldsymbol{\delta})$, and that $(y,A)\mapsto\mathrm{Sl}_{\square,y}(A)$ is continuous on $\mathbb{R}\times G^W(\boldsymbol{\delta})$. This proves (a).

Multiplication by a function of modulus one is an isometry for the sup norm and commutes with the multiplications $\mathbf{1}_{\Delta_z}$, $z$ a unit site; the multiplication of the functions of short pairs $(u,u')$ by $\exp\left(-ie_k\cdot\boldsymbol{\lambda}_0^W(u)\right)$ is an isometry for the sup norm as well, and it commutes with the operators $\mathbf{1}_{\Delta_z}$, which act on such functions as multiplications by functions of $u$. Therefore the formulas above and Proposition~\ref{prop:small-fields}(b), (c), applied on $\mathbb{T}^W$, give, for the norm $N_{\gamma}$ of Section~\ref{sec:norms-exponential-weights} on $\mathbb{T}^W$,
\begin{align*}
N_{\gamma}\left(F^W_y(A)^{-1}\right)&=N_{\gamma}\left(S_y(B^W)\right)\le2a\cdot w_k(y),\\
N_{\gamma}\left(\delta_{\alpha,A^W}F^W_y(A)^{-1}\right)&=N_{\gamma}\left(\delta_{\alpha,B^W}S_y(B^W)\right)\le18a.
\end{align*}
Let $(x,x')$ be a short pair of sites of $\tilde{W}$. Each coordinate of a site of a rectilinear path from $x$ to $x'$ lies between the corresponding coordinates of $x$ and of $x'$, and the set $\left\{z\in\mathbb{R}^3:|z-\mathsf{c}|<\frac{3M}{2}\right\}$, which contains $x$ and $x'$, is a product of three open intervals; hence all sites and bonds of the six paths $\Gamma^{\pi}(x,x')$ are sites and bonds of $\tilde{W}$. Since $\mathrm{emb}$ is induced by a translation, it maps these paths onto the paths $\Gamma^{\pi}\left(\mathrm{emb}(x),\mathrm{emb}(x')\right)$, which lie in $\mathrm{Ch}$; and $A^W\left(\mathrm{emb}(z),\mathrm{emb}(z')\right)=A(z,z')$ for every bond $(z,z')$ of $\tilde{W}$, because $\mathrm{Fo}\left(\mathrm{emb}(z)\right)=z$ (Lemma~\ref{lem:fold-map-embedding}(b)). Hence $U_{A^W}\left(\mathrm{emb}(x),\mathrm{emb}(x')\right)=U_A(x,x')$, and as $\left|\mathrm{emb}(x)-\mathrm{emb}(x')\right|=|x-x'|$,
\[
\left(\delta_{\alpha,A^W}v\right)\left(\mathrm{emb}(x),\mathrm{emb}(x')\right)=\left(\delta_{\alpha,A}(\operatorname{rs} v)\right)(x,x')
\]
for every spinor $v$ on $\mathbb{T}^W$. Let $\operatorname{rs}$ act also on the functions $w$ of the short pairs of $\mathbb{T}^W$, by $(\operatorname{rs} w)(x,x'):=w\left(\mathrm{emb}(x),\mathrm{emb}(x')\right)$ for the short pairs $(x,x')$ in $\tilde{W}$. Then
\[
\mathrm{Sl}_{\square,y}(A)=\operatorname{rs}{}\cdot F^W_y(A)^{-1}\cdot\operatorname{ext},\qquad\delta_{\alpha,A}\mathrm{Sl}_{\square,y}(A)=\operatorname{rs}{}\cdot\delta_{\alpha,A^W}F^W_y(A)^{-1}\cdot\operatorname{ext},
\]
and (b) follows from the two bounds on $\mathbb{T}^W$ once it is shown that
\[
N_{\gamma}\left(\operatorname{rs}{}\cdot O\cdot\operatorname{ext}\right)\le N_{\gamma}(O)
\]
for every operator $O$ on the spinors of $\mathbb{T}^W$ with values in the spinors or in the functions of short pairs. Let $y_1,y_2$ be unit sites in $\tilde{W}$. By Lemma~\ref{lem:fold-map-embedding}(a), $\mathrm{emb}$ maps $\Delta_{y_1}$ and $\Delta_{y_2}$ onto $\Delta_{\mathrm{emb}(y_1)}$ and $\Delta_{\mathrm{emb}(y_2)}$; hence
\[
\mathbf{1}_{\Delta_{y_1}}\cdot\operatorname{rs}=\operatorname{rs}{}\cdot\mathbf{1}_{\Delta_{\mathrm{emb}(y_1)}},\qquad\operatorname{ext}{}\cdot\mathbf{1}_{\Delta_{y_2}}=\mathbf{1}_{\Delta_{\mathrm{emb}(y_2)}}\cdot\operatorname{ext} .
\]
The operators $\operatorname{rs}$ and $\operatorname{ext}$ have norm at most $1$ for the sup norms, $|y_1-y_2|$ is the distance of $\mathrm{emb}(y_1)$ and $\mathrm{emb}(y_2)$ on $\mathbb{T}^W$ by Lemma~\ref{lem:fold-map-embedding}(c), and $\mathrm{emb}$ is injective on $\tilde{W}$. Therefore
\begin{align*}
&\sum_{y_2}\exp\left(\gamma|y_1-y_2|\right)\cdot\left\|\mathbf{1}_{\Delta_{y_1}}\cdot\operatorname{rs}{}\cdot O\cdot\operatorname{ext}{}\cdot\mathbf{1}_{\Delta_{y_2}}\right\|\\
&\qquad\le\sum_{y''}\exp\left(\gamma\,d\left(\mathrm{emb}(y_1),y''\right)\right)\cdot\left\|\mathbf{1}_{\Delta_{\mathrm{emb}(y_1)}}O\mathbf{1}_{\Delta_{y''}}\right\|\le N_{\gamma}(O),
\end{align*}
where $y_2$ ranges over the unit sites in $\tilde{W}$ and $y''$ over the unit sites of $\mathbb{T}^W$ modulo $6M\mathbb{Z}^3$. The supremum over $y_1$ gives the claim, and with it (b).

\emph{(c).} The row of $D_A$ at a site $x$ involves the values of the function at $x$ and at the six nearest neighbours of $x$, and the values of $A$ on the bonds from $x$ to these neighbours; the term $\mathfrak{m}$ acts diagonally. By Lemma~\ref{lem:block-averages-block}, the rows at $x$ of $P_k(A)$, $P(A)$ and $P_{k+1}(A)$, and hence the row of $\operatorname{Pen}_y(A)$ at $x$, involve only the values of the function on the block of side $L$ containing $x$ and the values of $A$ on the bonds inside this block; that lemma remains valid, with the same proof, for the blocks of $T_k$ lifted to $\eta\mathbb{Z}^3$ and for any continuation of $A$, since its proof rests only on the formulas of Section~\ref{sec:phase-factors-lattice}, on the nesting of the blocks and on the fact that a rectilinear path between two sites of a block lies in this block. For $k=0$ the term $i\gamma_3y$ of $\operatorname{Pen}_y(A)$ acts diagonally. If $x\in W^{-}$, the six neighbours of $x$ lie in $\tilde{W}$, and so does the block of side $L$ containing $x$, because $\tilde{W}$ is a union of such blocks (Lemma~\ref{lem:fold-map-embedding}(a)). This proves the first assertion.

Let $f$ be a function on $\tilde{W}$, and put $v':=F^W_y(A)^{-1}\cdot\operatorname{ext} f$, so that $v=\operatorname{rs} v'$ on $\tilde{W}$ and $v=0$ elsewhere. Let $x\in W^{-}$ and $u:=\mathrm{emb}(x)\in\mathrm{Ch}$. By Lemma~\ref{lem:fold-map-embedding}(a), the translation that induces $\mathrm{emb}$ maps the nearest neighbours of $x$ onto those of $u$, and it maps the block of side $L$ containing $x$, together with the blocks of the smaller sides contained in it, their centres and the rectilinear paths between their sites, onto the block of side $L$ containing $u$, together with the corresponding objects for $u$; moreover $A^W\left(\mathrm{emb}(z),\mathrm{emb}(z')\right)=A(z,z')$ for every bond $(z,z')$ of $\tilde{W}$. Consequently the row of $F_y(A^W)$ at $u$, applied to $v'$, coincides with the row of $F_y(A)$ at $x$, applied to $v$, and
\[
f(x)=(\operatorname{ext} f)(u)=\left(F_y(A^W)v'\right)(u)=\left(F_y(A)v\right)(x).
\]

Now let $g$ vanish outside $W^{-}$. If $x\notin\tilde{W}$, then no nearest neighbour of $x$ lies in $W^{-}$, since all nearest neighbours of a site of $W^{-}$ lie in $\tilde{W}$; and the block of side $L$ containing $x$ is disjoint from $\tilde{W}$, since $\tilde{W}$ is a union of blocks of side $L$. By the description of the rows given above, $\left(F_y(A)g\right)(x)=0$. Thus $F_y(A)g$ vanishes outside $\tilde{W}$. Next,
\[
F_y(A^W)\cdot\operatorname{ext} g=\operatorname{ext}\left(F_y(A)g\right),
\]
where $\operatorname{ext}$ is applied to the restrictions of $g$ and of $F_y(A)g$ to $\tilde{W}$. Indeed, let $u\in\Lambda^W$. If $u\in\mathrm{Ch}$, then $u=\mathrm{emb}(x)$ with $x:=\mathrm{Fo}(u)\in\tilde{W}$. The block of side $L$ containing $u$ lies in $\mathrm{Ch}$ and corresponds to the block of side $L$ containing $x$ as before; the neighbours of $u$ in $\mathrm{Ch}$ are the images of the neighbours of $x$ in $\tilde{W}$, with $A^W=A$ on the corresponding bonds; and a neighbour $u\pm\eta e_{\mu}$ of $u$ outside $\mathrm{Ch}$ corresponds to the neighbour $x\pm\eta e_{\mu}$ of $x$ outside $\tilde{W}$ (a site $z$ lies in $\tilde{W}$ if and only if $z-o\in(0,3M)^3$), so that the terms of the two rows that belong to it vanish, $\operatorname{ext} g$ vanishing at $u\pm\eta e_{\mu}$ and $g$ at $x\pm\eta e_{\mu}$. Hence $\left(F_y(A^W)\operatorname{ext} g\right)(u)=\left(F_y(A)g\right)(x)=\left(\operatorname{ext}\left(F_y(A)g\right)\right)(u)$. If $u\notin\mathrm{Ch}$, then $(\operatorname{ext} g)(u)=0$; the block of side $L$ containing $u$ lies in a single chamber (Lemma~\ref{lem:fold-map-embedding}(a)), which is not $\mathrm{Ch}$, so that $\operatorname{ext} g$ vanishes on it; and a neighbour $u'$ of $u$ that lies in $\mathrm{Ch}$ is separated from $u$ by a mirror plane, so that $u'_{\mu}\in\left\{\frac{\eta}{2},3M-\frac{\eta}{2}\right\}+6M\mathbb{Z}$ for the direction $\mu$ of the bond $(u,u')$. Then $u'=\mathrm{emb}(z)$ with $z:=\mathrm{Fo}(u')\in\tilde{W}$ and $|z_{\mu}-\mathsf{c}_{\mu}|=\frac{3M}{2}-\frac{\eta}{2}$, so that $z\notin W^{-}$ and $(\operatorname{ext} g)(u')=g(z)=0$. Hence $\left(F_y(A^W)\operatorname{ext} g\right)(u)=0=\left(\operatorname{ext}\left(F_y(A)g\right)\right)(u)$. Applying $\operatorname{rs}{}\cdot F^W_y(A)^{-1}$ to the identity and using $\operatorname{rs}{}\cdot\operatorname{ext}=1$, one obtains $g=\mathrm{Sl}_{\square,y}(A)\left(F_y(A)g\right)$ on $\tilde{W}$.

\emph{(d).} Let $\boldsymbol{\lambda}$ be a real function on $\tilde{W}$. Then $A-d\boldsymbol{\lambda}=B+d(\boldsymbol{\lambda}_0-\boldsymbol{\lambda})$ lies in $G^W(\boldsymbol{\delta})$, and $(A-d\boldsymbol{\lambda})^W=A^W-d\left(\boldsymbol{\lambda}^W\right)$. Lemma~\ref{lem:symmetries-elementary-operators}(a) on $\mathbb{T}^W$, applied to the connection $A^W$, to the closed real connection $d\left(\boldsymbol{\lambda}^W\right)$ and to the function $\exp\left(ie_k\cdot\boldsymbol{\lambda}^W\right)$, as in the example given in that lemma, yields
\[
F^W_y(A-d\boldsymbol{\lambda})=\exp\left(ie_k\cdot\boldsymbol{\lambda}^W\right)\cdot F^W_y(A)\cdot\exp\left(-ie_k\cdot\boldsymbol{\lambda}^W\right).
\]
Since $\mathrm{Fo}\left(\mathrm{emb}(x)\right)=x$ for $x\in\tilde{W}$, the definitions of $\operatorname{rs}$ and $\operatorname{ext}$ give
\[
\operatorname{rs}{}\cdot\exp\left(ie_k\cdot\boldsymbol{\lambda}^W\right)=\exp\left(ie_k\cdot\boldsymbol{\lambda}\right)\cdot\operatorname{rs},\qquad\exp\left(-ie_k\cdot\boldsymbol{\lambda}^W\right)\cdot\operatorname{ext}=\operatorname{ext}{}\cdot\exp\left(-ie_k\cdot\boldsymbol{\lambda}\right).
\]
Inserting these relations into $\mathrm{Sl}_{\square,y}(A-d\boldsymbol{\lambda})=\operatorname{rs}{}\cdot F^W_y(A-d\boldsymbol{\lambda})^{-1}\cdot\operatorname{ext}$ gives (d).

\emph{(e).} Put $r:=r(S)$, the diagonal matrix with the diagonal entries $r_{\mu}:=r_{\mu}(S)\in\{1,-1\}$, $\mu=0,1,2$ (Section~\ref{sec:parameters-clifford-relations}), and $\boldsymbol{\chi}(x):=rx+a'$, so that $gx=\boldsymbol{\chi}(x)$. The point $\boldsymbol{\chi}(\mathsf{c})\in M\mathbb{Z}^3$ is the chosen lift of the centre of $g\square$. The objects of Section~\ref{sec:folding-torus-onto} formed with this lift are marked by a prime: $\tilde{W}'=\boldsymbol{\chi}(\tilde{W})$, since $\boldsymbol{\chi}$ preserves the norm of differences, $o'=\boldsymbol{\chi}(\mathsf{c})-\frac{3M}{2}(1,1,1)$, and $\mathrm{Fo}'$, $\mathrm{emb}'$, $\operatorname{ext}'$, $\operatorname{rs}'$, the fold $E\mapsto E^{W'}$ of the complex antisymmetric functions $E$ on the bonds of $\tilde{W}'$ and $F^{W'}_{y'}(E):=F_{y'}\left(E^{W'}\right)$ are defined with $o'$ and $\tilde{W}'$ in place of $o$ and $\tilde{W}$; the torus $\mathbb{T}^W$ and the chamber $\mathrm{Ch}$ are the same as for $\square$. Let
\[
\boldsymbol{\chi}^W(u):=r\left(u-\frac{3M}{2}(1,1,1)\right)+\frac{3M}{2}(1,1,1)\qquad(u\in\mathbb{R}^3).
\]
In a coordinate with $r_{\mu}=-1$ this is the reflection $u_{\mu}\mapsto3M-u_{\mu}$, and $3M$ is an integer multiple of $L^j\eta$ for $0\le j\le k+1$ (Lemma~\ref{lem:fold-map-embedding}(a)). So $\boldsymbol{\chi}^W$ maps each lattice $\Lambda^W_j$ onto itself and is a morphism of $\mathbb{T}^W$ onto itself (Definition~\ref{def:block-torus}); it also maps $\mathrm{Ch}$ onto $\mathrm{Ch}$, since $t\mapsto3M-t$ maps $(0,3M)+6M\mathbb{Z}$ onto itself. The function $\mathrm{fo}$ satisfies
\[
\mathrm{fo}(3M-t)=3M-\mathrm{fo}(t)\qquad(t\in\mathbb{R}),
\]
because both sides are even and $6M$-periodic functions of $t$ that agree on $[0,3M]$. Hence
\[
\mathrm{Fo}'\circ\boldsymbol{\chi}^W=\boldsymbol{\chi}\circ\mathrm{Fo},\qquad\boldsymbol{\chi}^W\circ\mathrm{emb}=\mathrm{emb}'\circ\boldsymbol{\chi},
\]
as is seen coordinate by coordinate: at a point $u$, the $\mu$-th coordinates of the two sides of the first identity are both equal to $\boldsymbol{\chi}(\mathsf{c})_{\mu}-\frac{3M}{2}+\mathrm{fo}(u_{\mu})$ if $r_{\mu}=1$, and to $\boldsymbol{\chi}(\mathsf{c})_{\mu}+\frac{3M}{2}-\mathrm{fo}(u_{\mu})$ if $r_{\mu}=-1$; at a point $x$, the $\mu$-th coordinates of the two sides of the second identity are both equal, modulo $6M$, to $x_{\mu}-\mathsf{c}_{\mu}+\frac{3M}{2}$ if $r_{\mu}=1$, and to $\mathsf{c}_{\mu}-x_{\mu}+\frac{3M}{2}$ if $r_{\mu}=-1$.

Let $g^W:=(S,\boldsymbol{\chi}^W)$ act on the spinors and on the connections of $\mathbb{T}^W$ as in Lemma~\ref{lem:symmetries-elementary-operators}(b): $(g^Wv)(u):=S\cdot v\left((\boldsymbol{\chi}^W)^{-1}u\right)$ and $(g^WA'')(u,u'):=A''\left((\boldsymbol{\chi}^W)^{-1}u,(\boldsymbol{\chi}^W)^{-1}u'\right)$. The two identities, together with the definitions of the fold, of $\operatorname{ext}$ and of $\operatorname{rs}$, give
\[
(gA)^{W'}=g^W\left(A^W\right),\qquad\operatorname{ext}'\cdot g=g^W\cdot\operatorname{ext},\qquad\operatorname{rs}'\cdot g^W=g\cdot\operatorname{rs} .
\]
Indeed, the two identities say that $\boldsymbol{\chi}^{-1}\circ\mathrm{Fo}'=\mathrm{Fo}\circ(\boldsymbol{\chi}^W)^{-1}$ and $(\boldsymbol{\chi}^W)^{-1}\circ\mathrm{emb}'=\mathrm{emb}\circ\boldsymbol{\chi}^{-1}$. In particular $\mathrm{Fo}'(u)=\mathrm{Fo}'(u')$ holds if and only if $\mathrm{Fo}\left((\boldsymbol{\chi}^W)^{-1}u\right)=\mathrm{Fo}\left((\boldsymbol{\chi}^W)^{-1}u'\right)$, in which case $(gA)^{W'}(u,u')$ and $\left(g^WA^W\right)(u,u')$ both vanish. If $\mathrm{Fo}'(u)\ne\mathrm{Fo}'(u')$, the first of the following chains of equalities holds; the second holds for $u\in\mathrm{Ch}$, where $(\boldsymbol{\chi}^W)^{-1}u\in\mathrm{Ch}$, and the third for $x\in\tilde{W}'$:
\begin{align*}
(gA)^{W'}(u,u')&=A\left(\boldsymbol{\chi}^{-1}\mathrm{Fo}'(u),\boldsymbol{\chi}^{-1}\mathrm{Fo}'(u')\right)\\
&=A^W\left((\boldsymbol{\chi}^W)^{-1}u,(\boldsymbol{\chi}^W)^{-1}u'\right)=\left(g^WA^W\right)(u,u'),\\
\left(\operatorname{ext}'(gf)\right)(u)&=S\cdot f\left(\boldsymbol{\chi}^{-1}\mathrm{Fo}'(u)\right)\\
&=S\cdot f\left(\mathrm{Fo}\left((\boldsymbol{\chi}^W)^{-1}u\right)\right)=\left(g^W\operatorname{ext} f\right)(u),\\
\left(\operatorname{rs}'(g^Wv)\right)(x)&=S\cdot v\left((\boldsymbol{\chi}^W)^{-1}\mathrm{emb}'(x)\right)\\
&=S\cdot v\left(\mathrm{emb}\left(\boldsymbol{\chi}^{-1}x\right)\right)=\left(g\,\operatorname{rs} v\right)(x);
\end{align*}
and for $u\notin\mathrm{Ch}$ both $\left(\operatorname{ext}'(gf)\right)(u)$ and $\left(g^W\operatorname{ext} f\right)(u)$ vanish.
The function $gA$ belongs to the set $G^W(\boldsymbol{\delta})$ formed with $\tilde{W}'$, since $gA=gB+d\left(\boldsymbol{\lambda}_0\circ\boldsymbol{\chi}^{-1}\right)$ and $|gB|_{\infty}=|B|_{\infty}$. By Lemma~\ref{lem:symmetries-elementary-operators}(b), applied to the morphism $\boldsymbol{\chi}^W$ of $\mathbb{T}^W$ onto itself,
\[
F_{y'}\left(g^WA^W\right)\cdot g^W=g^W\cdot F_y\left(A^W\right),\qquad\text{that is,}\qquad F^{W'}_{y'}(gA)\cdot g^W=g^W\cdot F^W_y(A).
\]
Both $F^W_y(A)$ and $F^{W'}_{y'}(gA)$ are invertible by (a). Hence $F^{W'}_{y'}(gA)^{-1}\cdot g^W=g^W\cdot F^W_y(A)^{-1}$, and
\begin{align*}
\mathrm{Sl}_{g\square,y'}(gA)\cdot g&=\operatorname{rs}'\cdot F^{W'}_{y'}(gA)^{-1}\cdot g^W\cdot\operatorname{ext}=\operatorname{rs}'\cdot g^W\cdot F^W_y(A)^{-1}\cdot\operatorname{ext}\\
&=g\cdot\operatorname{rs}{}\cdot F^W_y(A)^{-1}\cdot\operatorname{ext}=g\cdot\mathrm{Sl}_{\square,y}(A).
\end{align*}

\emph{(f).} The functions $-A=(-B)+d(-\boldsymbol{\lambda}_0)$ and $\overline{A}=\overline{B}+d\boldsymbol{\lambda}_0$ belong to $G^W(\boldsymbol{\delta})$, and $(-A)^W=-A^W$, $\left(\overline{A}\right)^W=\overline{A^W}$. By Lemma~\ref{lem:symmetries-elementary-operators}(c) on $\mathbb{T}^W$, $F_y(A^W)^T=\mathsf{C}^{-1}\cdot F_{-y}(-A^W)\cdot\mathsf{C}$; hence
\[
\left(F^W_y(A)^{-1}\right)^T=\mathsf{C}^{-1}\cdot F^W_{-y}(-A)^{-1}\cdot\mathsf{C} .
\]
As $\operatorname{ext}^T=\operatorname{rs}$, $\operatorname{rs}^T=\operatorname{ext}$, and $\operatorname{rs}$, $\operatorname{ext}$ commute with constant matrices,
\begin{align*}
\mathrm{Sl}_{\square,y}(A)^T&=\operatorname{rs}{}\cdot\left(F^W_y(A)^{-1}\right)^T\cdot\operatorname{ext}=\mathsf{C}^{-1}\cdot\operatorname{rs}{}\cdot F^W_{-y}(-A)^{-1}\cdot\operatorname{ext}{}\cdot\mathsf{C}\\
&=\mathsf{C}^{-1}\cdot\mathrm{Sl}_{\square,-y}(-A)\cdot\mathsf{C} .
\end{align*}
By Lemma~\ref{lem:symmetries-elementary-operators}(d) on $\mathbb{T}^W$, $\overline{F_y\left(\overline{A^W}\right)}=\varkappa^{-1}\cdot F_y(-A^W)\cdot\varkappa$. The operators $\operatorname{rs}$ and $\operatorname{ext}$ map real functions to real functions and commute with $\varkappa$; therefore
\begin{align*}
\overline{\mathrm{Sl}_{\square,y}\left(\overline{A}\right)}&=\operatorname{rs}{}\cdot\left(\overline{F^W_y\left(\overline{A}\right)}\right)^{-1}\cdot\operatorname{ext}=\operatorname{rs}{}\cdot\varkappa^{-1}\cdot F^W_y(-A)^{-1}\cdot\varkappa\cdot\operatorname{ext}\\
&=\varkappa^{-1}\cdot\mathrm{Sl}_{\square,y}(-A)\cdot\varkappa .
\end{align*}

\emph{(g).} When $\mathsf{c}$ is replaced by $\mathsf{c}+ln$, the sets $\tilde{W}$, $W^{-}$ and the point $o$ are replaced by their translates by $ln$, the fold map $\mathrm{Fo}$ by $\mathrm{Fo}+ln$, and $\mathrm{emb}$ by the map $x\mapsto\mathrm{emb}(x-ln)$, while $\mathbb{T}^W$ and $\mathrm{Ch}$ do not change. By the invariance of $\tilde{A}$, one has $A_n\left(\mathrm{Fo}(u)+ln,\mathrm{Fo}(u')+ln\right)=\tilde{A}\left(\mathrm{Fo}(u),\mathrm{Fo}(u')\right)=A\left(\mathrm{Fo}(u),\mathrm{Fo}(u')\right)$ whenever $\mathrm{Fo}(u)\ne\mathrm{Fo}(u')$; so the fold of $A_n$ formed with the new lift is $A^W$, and the operator $F^W_y$ formed with the new lift takes at $A_n$ the value $F^W_y(A)$. The new map $\operatorname{ext}$ sends $f$ to $\operatorname{ext}\left(f(\cdot+ln)\right)$, and the new map $\operatorname{rs}$ sends $v$ to the function $x\mapsto(\operatorname{rs} v)(x-ln)$. Inserting this into the definition of the local inverse gives (g).
\end{proof}

\subsubsection{Operators on a cube as operators on the torus}\label{sec:operators-cube-operators}

Consider a linear operator $O$ that sends functions $\tilde{W}\to\mathbb{C}^4$ to functions on $\tilde{W}$, or to functions on the short pairs of sites of $\tilde{W}$. As $3M<l$, the translates $\tilde{W}+ln$, $n\in\mathbb{Z}^3$, are pairwise disjoint. A spinor $f$ of $T_k$ is an antiperiodic function on $\eta\mathbb{Z}^3$, and $f|_{\tilde{W}}$ denotes its restriction to $\tilde{W}$. The operator $O$ acts on the spinors of $T_k$ by
\begin{gather*}
(Of)(x):=\left(O\left(f|_{\tilde{W}}\right)\right)(x)\qquad(x\in\tilde{W}),\\
(Of)(x+ln):=(-1)^{n_0+n_1+n_2}\cdot(Of)(x)\qquad(x\in\tilde{W},\ n\in\mathbb{Z}^3),
\end{gather*}
and $(Of)(x):=0$ for the sites $x$ outside the orbit $\bigcup_{n\in\mathbb{Z}^3}(\tilde{W}+ln)$ of $\tilde{W}$; if $O$ takes values in the functions of short pairs, the same formulas define $Of$ on the short pairs $(x,x')$ in $\tilde{W}$ and on their translates $(x+ln,x'+ln)$, and $Of$ vanishes on all other short pairs. Let $n\in\mathbb{Z}^3$, and let $O_n$ be the translate of $O$ by $ln$, which sends a function $h$ on $\tilde{W}+ln$ to the function $x\mapsto\left(O\,h(\cdot+ln)\right)(x-ln)$. The restriction of a spinor $f$ to $\tilde{W}+ln$ is $(-1)^{n_0+n_1+n_2}$ times the translate by $ln$ of $f|_{\tilde{W}}$; hence $O_n$ maps it to $(-1)^{n_0+n_1+n_2}$ times the translate by $ln$ of $O\left(f|_{\tilde{W}}\right)$, and the formulas above, written with $\tilde{W}+ln$ and $O_n$, define the same operator on the spinors of $T_k$ as those written with $\tilde{W}$ and $O$. Now let $A$ be a connection on $T_k$ whose restriction to the bonds of $\tilde{W}$ lies in $G^W(\boldsymbol{\delta})$ for a number $\boldsymbol{\delta}>0$ satisfying~\eqref{eq:operators-fermion-6}, and let $y\in Y_k$. Two lifts of the centre of $\square$ differ by an element of $l\mathbb{Z}^3$, and $A$ is invariant under the translations by $l\mathbb{Z}^3$; hence, by Proposition~\ref{prop:properties-local-inverse}(g) and what precedes, the operator on the spinors of $T_k$ defined by the local inverse of the restriction of $A$ to the bonds of $\tilde{W}$ is determined by $\square$ and does not depend on the choice of the lift $\mathsf{c}$. By Lemma~\ref{lem:fold-map-embedding}(c), the localised norms $N_{\gamma}$ on $T_k$ (Section~\ref{sec:norms-exponential-weights}) of the operators so obtained are equal to the norms $N_{\gamma}$ of Proposition~\ref{prop:properties-local-inverse}(b). Let $A\in\tilde{G}_k$. By~\eqref{eq:operators-fermion-G1}, the restriction of $A$ to the bonds of $\tilde{W}$ is an element of $G^W\left(2e_k^{\varepsilon}\right)$; when $\boldsymbol{\delta}:=2e_k^{\varepsilon}$ satisfies~\eqref{eq:operators-fermion-6}, the operator on the spinors of $T_k$ defined by the local inverse of this restriction is again denoted by $\mathrm{Sl}_{\square,y}(A)$.

Let $O$ be an operator on the functions $\tilde{W}\to\mathbb{C}^4$ with values in the functions on $\tilde{W}$, acting on the spinors of $T_k$ by the formulas above. Let $u$ be a function on $\eta\mathbb{Z}^3$ without zeros that is invariant under the translations by $l\mathbb{Z}^3$, and let $u|_{\tilde{W}}$ be its restriction to $\tilde{W}$. The multiplication by $u$ maps the spinors of $T_k$ to spinors of $T_k$, the multiplication by $u|_{\tilde{W}}$ maps the functions on $\tilde{W}$ to functions on $\tilde{W}$, and
\[
\left(u|_{\tilde{W}}\cdot O\cdot\left(u|_{\tilde{W}}\right)^{-1}\right)f=u\cdot O\left(u^{-1}f\right)\qquad\text{for every spinor }f\text{ of }T_k,
\]
where on the left the operator $u|_{\tilde{W}}\cdot O\cdot\left(u|_{\tilde{W}}\right)^{-1}$ on the functions on $\tilde{W}$ acts on $f$ by the formulas above. Indeed, $\left(u^{-1}f\right)|_{\tilde{W}}=\left(u|_{\tilde{W}}\right)^{-1}\cdot f|_{\tilde{W}}$, so that the two sides coincide on $\tilde{W}$; at a site $x+ln$ with $x\in\tilde{W}$ and $n\in\mathbb{Z}^3$ each side takes $(-1)^{n_0+n_1+n_2}$ times its value at $x$, since $u(x+ln)=u(x)$; and both sides vanish outside the orbit of $\tilde{W}$. The same argument gives this identity with $u$ and $u|_{\tilde{W}}$ both replaced by one constant invertible complex $4\times4$ matrix.

Let $g=(S,a')\in\mathbb{G}_k$ with $a'\in M\mathbb{Z}^3$, and let $gx:=r(S)x+a'$. The element $g$ acts on the spinors of $T_k$ and on the connections on $T_k$ as in Lemma~\ref{lem:symmetries-elementary-operators}(b), applied to the morphism $x\mapsto gx$ of $T_k$ onto itself (Section~\ref{sec:tori-nested-blocks}), and it maps the functions on $\tilde{W}$ bijectively onto the functions on $r(S)\tilde{W}+a'$ as in Proposition~\ref{prop:properties-local-inverse}(e); the set $r(S)\tilde{W}+a'$ is the lift of $(g\square)^{\sim 1}$ that corresponds to the lift $r(S)\mathsf{c}+a'$ of the centre of $g\square$. Let $O^g$ be the operator $h\mapsto g\left(O\left(g^{-1}h\right)\right)$ on the functions on $r(S)\tilde{W}+a'$, acting on the spinors of $T_k$ by the formulas above with $r(S)\tilde{W}+a'$ in the place of $\tilde{W}$. Then
\[
O^gf=g\left(O\left(g^{-1}f\right)\right)\qquad\text{for every spinor }f\text{ of }T_k.
\]
Indeed, $\left(g^{-1}f\right)|_{\tilde{W}}=g^{-1}\left(f|_{r(S)\tilde{W}+a'}\right)$. Let $x\in\tilde{W}$ and $n\in\mathbb{Z}^3$. As $r(S)$ is diagonal with the diagonal entries $\pm1$, one has $g^{-1}\left(gx+l\,r(S)n\right)=x+ln$, and the sum of the coordinates of $r(S)n\in\mathbb{Z}^3$ has the parity of $n_0+n_1+n_2$. Hence
\begin{align*}
\left(g\left(O\left(g^{-1}f\right)\right)\right)\left(gx+l\,r(S)n\right)&=S\cdot\left(O\left(g^{-1}f\right)\right)(x+ln)\\
&=(-1)^{n_0+n_1+n_2}\cdot S\cdot\left(O\left(g^{-1}\left(f|_{r(S)\tilde{W}+a'}\right)\right)\right)(x)\\
&=(-1)^{n_0+n_1+n_2}\cdot\left(O^g\left(f|_{r(S)\tilde{W}+a'}\right)\right)(gx),
\end{align*}
and the last expression is the value of $O^gf$ at $gx+l\,r(S)n$. Both sides of the identity vanish outside the orbit of $r(S)\tilde{W}+a'$, which is the image of the orbit of $\tilde{W}$ under the map $x\mapsto gx$.

Let $O^T$ be the transpose of $O$ with respect to the pairing of the functions on $\tilde{W}$ of Section~\ref{sec:folding-torus-onto}. Two sites of $\tilde{W}$ differ by a vector of sup norm less than $3M<l$, so that no two of them are congruent modulo $l\mathbb{Z}^3$. Hence $\tilde{W}$, completed by one site of each class modulo $l\mathbb{Z}^3$ that contains no site of $\tilde{W}$, is a period cell of $T_k$, and the sites of this cell outside $\tilde{W}$ lie outside the orbit of $\tilde{W}$. Computing the pairing of Section~\ref{sec:norms-pairings-kernels} on this cell, one obtains for all spinors $f,f'$ of $T_k$
\[
(f',Of)=\sum_{x\in\tilde{W}}\eta^3f'(x)\cdot\left(O\left(f|_{\tilde{W}}\right)\right)(x)=\left(O^T\left(f'|_{\tilde{W}}\right),f|_{\tilde{W}}\right)=\left(O^Tf',f\right);
\]
that is, the operator on the spinors of $T_k$ defined by $O^T$ is the transpose of the operator on the spinors of $T_k$ defined by $O$. Let $\overline{O}$ be the operator $f\mapsto\overline{O\overline{f}}$ on the functions on $\tilde{W}$, as in Proposition~\ref{prop:properties-local-inverse}(f). The restriction to $\tilde{W}$ commutes with complex conjugation, and the signs $(-1)^{n_0+n_1+n_2}$ are real; therefore $\overline{O}f=\overline{O\overline{f}}$ for every spinor $f$ of $T_k$, the operators $\overline{O}$ and $O$ acting on the spinors of $T_k$ by the formulas above. Thus the operator on the spinors of $T_k$ defined by $\overline{O}$ is the complex conjugate, in the sense of Lemma~\ref{lem:symmetries-elementary-operators}(d), of the operator defined by $O$.

Let $A\in\tilde{G}_k$, let $\boldsymbol{\delta}:=2e_k^{\varepsilon}$ satisfy~\eqref{eq:operators-fermion-6}, let $y\in Y_k$, with $-\infty:=\infty$ as in Proposition~\ref{prop:properties-local-inverse}, let $\boldsymbol{\lambda}$ be a real scalar on $T_k$, let $g=(S,a')$ be as above, and let $y'$ be formed from $y$ and $S$ as in Proposition~\ref{prop:properties-local-inverse}(e). By~\eqref{eq:operators-fermion-G2}, the connections $A-d\boldsymbol{\lambda}$, $gA$, $-A$ and $\overline{A}$ belong to $\tilde{G}_k$. Let $A|_{\tilde{W}}$ be the restriction of $A$ to the bonds of $\tilde{W}$. The restrictions of $A-d\boldsymbol{\lambda}$, $-A$ and $\overline{A}$ to the bonds of $\tilde{W}$ are $A|_{\tilde{W}}-d\left(\boldsymbol{\lambda}|_{\tilde{W}}\right)$, with $d$ as in Definition~\ref{def:folded-fields-local}, $-A|_{\tilde{W}}$ and $\overline{A|_{\tilde{W}}}$; the restriction of $gA$ to the bonds of $r(S)\tilde{W}+a'$ is the function $(x,x')\mapsto A\left(g^{-1}x,g^{-1}x'\right)$, that is, the function $g\left(A|_{\tilde{W}}\right)$ of Proposition~\ref{prop:properties-local-inverse}(e). The operator $\mathrm{Sl}_{g\square,y'}(gA)$ on the spinors of $T_k$ does not depend on the lift of the centre of $g\square$; it is formed here with the lift $r(S)\mathsf{c}+a'$. Proposition~\ref{prop:properties-local-inverse}(d), (e), (f), applied to $A|_{\tilde{W}}$, and the four identities proved above, the first of them taken with $u:=\exp\left(ie_k\cdot\boldsymbol{\lambda}\right)$ and with the constant matrices $\mathsf{C}^{-1}$ and $\varkappa^{-1}$, give
\begin{gather*}
\mathrm{Sl}_{\square,y}(A-d\boldsymbol{\lambda})=\exp\left(ie_k\cdot\boldsymbol{\lambda}\right)\cdot\mathrm{Sl}_{\square,y}(A)\cdot\exp\left(-ie_k\cdot\boldsymbol{\lambda}\right),\qquad\mathrm{Sl}_{g\square,y'}(gA)\cdot g=g\cdot\mathrm{Sl}_{\square,y}(A),\\
\mathrm{Sl}_{\square,y}(A)^T=\mathsf{C}^{-1}\cdot\mathrm{Sl}_{\square,-y}(-A)\cdot\mathsf{C},\qquad\overline{\mathrm{Sl}_{\square,y}\left(\overline{A}\right)}=\varkappa^{-1}\cdot\mathrm{Sl}_{\square,y}(-A)\cdot\varkappa,
\end{gather*}
as identities between operators on the spinors of $T_k$; the transpose is taken with respect to the pairing of Section~\ref{sec:norms-pairings-kernels}, and the complex conjugate in the sense of Lemma~\ref{lem:symmetries-elementary-operators}(d).

\subsection{Inverses on the torus $T_k$}\label{sec:inverses-torus}

\subsubsection{A quadratic partition of unity}\label{sec:quadratic-partition-unity}

Let $\boldsymbol{\theta}\colon\mathbb{R}\to[0,1]$ be the function
\[
\boldsymbol{\theta}(t):=
\begin{cases}
1 & \text{if } |t|\le\frac{1}{4},\\
\cos\left(\frac{\pi}{2}\left(2|t|-\frac{1}{2}\right)\right) & \text{if } \frac{1}{4}\le|t|\le\frac{3}{4},\\
0 & \text{if } |t|\ge\frac{3}{4}.
\end{cases}
\]
Let $\square$ be an $M$-cube of $T_k$, and let $\mathsf{c}\in M\mathbb{Z}^3$ and $\tilde{W}\subset\eta\mathbb{Z}^3$ be as in Section~\ref{sec:folding-torus-onto}. Since $3M<l$, the sets $\tilde{W}+ln$, $n\in\mathbb{Z}^3$, are pairwise disjoint, and a function $h_{\square}$ on the sites of $T_k$, that is, an $l\mathbb{Z}^3$-periodic function on $\eta\mathbb{Z}^3$, with values in $[0,1]$ is defined by
\[
h_{\square}(x):=\prod_{\mu=0}^2\boldsymbol{\theta}\left(\frac{x_{\mu}-\mathsf{c}_{\mu}}{M}\right)\quad(x\in\tilde{W}),\qquad h_{\square}(x):=0\quad\left(x\notin\tilde{W}+l\mathbb{Z}^3\right).
\]
The operator of multiplication by $h_{\square}$ on the spinors on $T_k$ is denoted by $h_{\square}$ as well.

The function $\boldsymbol{\theta}$ is even; its values lie in $[0,1]$, since for $\frac{1}{4}\le|t|\le\frac{3}{4}$ the argument of the cosine lies in $\left[0,\frac{\pi}{2}\right]$. It is continuous, and outside the four points $\pm\frac{1}{4}$, $\pm\frac{3}{4}$ it is differentiable with a derivative of modulus at most $\pi$; hence $|\boldsymbol{\theta}(t)-\boldsymbol{\theta}(t')|\le\pi|t-t'|$ for all $t,t'\in\mathbb{R}$. Moreover
\[
\sum_{n\in\mathbb{Z}}\boldsymbol{\theta}(t-n)^2=1\qquad(t\in\mathbb{R}).
\]
Indeed, the left side has period $1$ in $t$. For $|t|\le\frac{1}{4}$ its only non-zero term is $\boldsymbol{\theta}(t)^2=1$. For $\frac{1}{4}\le t\le\frac{3}{4}$ its only non-zero terms are those with $n=0$ and $n=1$, and with $\phi:=\pi t-\frac{\pi}{4}$ one has $\boldsymbol{\theta}(t)=\cos\phi$ and $\boldsymbol{\theta}(t-1)=\cos\left(\frac{\pi}{2}-\phi\right)=\sin\phi$.

It follows that
\begin{multline}\label{eq:operators-fermion-7}
\sum_{\square}h_{\square}^2=1,\qquad\operatorname{supp}h_{\square}\subseteq\left\{x\in\eta\mathbb{Z}^3:|x-\mathsf{c}|<\frac{3M}{4}\right\}+l\mathbb{Z}^3,\\
\left|h_{\square}(x)-h_{\square}(x')\right|\le\frac{\pi}{M}\sum_{\mu}\left|x_{\mu}-x'_{\mu}\right|,\qquad h_{g\square}(gx)=h_{\square}(x)\quad(g\in\mathbb{G}_k),
\end{multline}
where the sum in the first relation is over all $M$-cubes of $T_k$, and the other three relations hold for every $M$-cube $\square$ of $T_k$, all $x,x'\in\eta\mathbb{Z}^3$ and all $g\in\mathbb{G}_k$. For the proof let $\square$ be an $M$-cube of $T_k$. For a lift $c\in M\mathbb{Z}^3$ of its centre let $h_{\square,c}\colon\eta\mathbb{Z}^3\to[0,1]$ be the function
\[
h_{\square,c}(x):=\prod_{\mu=0}^2\boldsymbol{\theta}\left(\frac{x_{\mu}-c_{\mu}}{M}\right);
\]
it vanishes at $x$ unless $|x-c|<\frac{3M}{4}$. Two distinct lifts of the centre differ by a non-zero element of $l\mathbb{Z}^3$, so that their distance is at least $l>3M$; hence for each $x\in\eta\mathbb{Z}^3$ at most one of the numbers $h_{\square,c}(x)$ is non-zero. Moreover
\[
h_{\square}(x)=\sum_ch_{\square,c}(x)\qquad\left(x\in\eta\mathbb{Z}^3\right),
\]
the sum being over the lifts $c$ of the centre of $\square$. Indeed, on $\tilde{W}+ln$, $n\in\mathbb{Z}^3$, the periodic function $h_{\square}$ coincides with $h_{\square,\mathsf{c}+ln}$, while every other lift $c$ has distance more than $l-\frac{3M}{2}>\frac{3M}{4}$ from the points of $\tilde{W}+ln$; outside $\tilde{W}+l\mathbb{Z}^3$ both sides vanish, because there $|x-c|\ge\frac{3M}{2}$ for every lift $c$. In particular, $h_{\square}$ does not depend on the choice of the lift $\mathsf{c}$.

For the first relation let $x\in\eta\mathbb{Z}^3$. Since $l$ is a multiple of $M$, every point of $M\mathbb{Z}^3$ is a lift of the centre of exactly one $M$-cube, and the square of a sum with at most one non-zero term is the sum of the squares of its terms. Therefore
\[
\sum_{\square}h_{\square}(x)^2=\sum_{c\in M\mathbb{Z}^3}\ \prod_{\mu=0}^2\boldsymbol{\theta}\left(\frac{x_{\mu}-c_{\mu}}{M}\right)^2=\prod_{\mu=0}^2\ \sum_{j\in\mathbb{Z}}\boldsymbol{\theta}\left(\frac{x_{\mu}}{M}-j\right)^2=1.
\]
The second relation follows from the representation of $h_{\square}$ as a sum over the lifts, since $h_{\square,c}$ vanishes outside $\left\{x:|x-c|<\frac{3M}{4}\right\}$ and the lifts of the centre of $\square$ are the points $\mathsf{c}+ln$, $n\in\mathbb{Z}^3$.

For the third relation let $x,x'\in\eta\mathbb{Z}^3$. If $h_{\square}(x)=h_{\square}(x')=0$, its left side is $0$; as both sides are symmetric in $x$ and $x'$, it suffices to treat the case $h_{\square}(x)\ne0$. Then $h_{\square}(x)=h_{\square,c}(x)$ for exactly one lift $c$. Suppose first that $h_{\square}(x')=h_{\square,c}(x')$ with the same $c$; this holds if $h_{\square}(x')=0$, and also if the non-zero term at $x'$ has the index $c$. The inequality $\left|\prod_{\mu}a_{\mu}-\prod_{\mu}b_{\mu}\right|\le\sum_{\mu}\left|a_{\mu}-b_{\mu}\right|$ for numbers $a_{\mu},b_{\mu}\in[0,1]$ and the Lipschitz bound for $\boldsymbol{\theta}$ give
\[
\left|h_{\square}(x)-h_{\square}(x')\right|\le\sum_{\mu=0}^2\left|\boldsymbol{\theta}\left(\frac{x_{\mu}-c_{\mu}}{M}\right)-\boldsymbol{\theta}\left(\frac{x'_{\mu}-c_{\mu}}{M}\right)\right|\le\frac{\pi}{M}\sum_{\mu}\left|x_{\mu}-x'_{\mu}\right|.
\]
Otherwise $h_{\square}(x')=h_{\square,c'}(x')\ne0$ for a lift $c'\ne c$. Then $|x-c|<\frac{3M}{4}$, $|x'-c'|<\frac{3M}{4}$ and $|c-c'|\ge l$, so that
\[
\sum_{\mu}\left|x_{\mu}-x'_{\mu}\right|\ge|x-x'|>l-\frac{3M}{2}>\frac{3M}{2};
\]
the right side of the third relation is then larger than $\frac{3\pi}{2}>1$, while its left side is at most $1$, both values of $h_{\square}$ lying in $[0,1]$.

The fourth relation holds because $\boldsymbol{\theta}$ is even. Indeed, an element $g=(S,a')$ of $\mathbb{G}_k$ maps $x$ to $r(S)x+a'$ and the centre of $\square$ to the centre of $g\square$; as $r(S)$ is diagonal with the entries $r_{\mu}(S)\in\{1,-1\}$, the map $c\mapsto r(S)c+a'$ is a bijection of the set of the lifts of the centre of $\square$ onto the set of the lifts of the centre of $g\square$, and
\[
\left(r(S)x+a'\right)_{\mu}-\left(r(S)c+a'\right)_{\mu}=r_{\mu}(S)\cdot\left(x_{\mu}-c_{\mu}\right),\qquad\text{so that}\qquad h_{g\square,r(S)c+a'}(gx)=h_{\square,c}(x);
\]
the sum over $c$ gives $h_{g\square}(gx)=h_{\square}(x)$.

Let $z$ be a unit site. If the unit block $\Delta_z$ meets $\operatorname{supp}h_{\square}$, then $|z_{\mu}-\mathsf{c}_{\mu}|<\frac{3M}{4}+\frac{1}{2}$ for $\mu=0,1,2$ and a suitable lift $\mathsf{c}$ of the centre of $\square$; an open interval of length $\frac{3M}{2}+1<2M$ contains at most two points of $M\mathbb{Z}$. Hence $\Delta_z$ meets $\operatorname{supp}h_{\square}$ for at most eight $M$-cubes $\square$.

\subsubsection{A parametrix from local inverses}\label{sec:parametrix-local-inverses}

In Sections~\ref{sec:parametrix-local-inverses} and~\ref{sec:expansion-walks}, conditions~\eqref{eq:renormalisation-group-20} and~\eqref{eq:renormalisation-group-21} hold and $M>4L$; moreover $A\in\tilde{G}_k$ and $y\in Y_k$ are fixed, and $F:=F_y(A)$ is the operator of~\eqref{eq:operators-fermion-2} on the spinors on $T_k$. For an $M$-cube $\square$ of $T_k$ let
\[
\mathrm{Sl}_{\square}:=\mathrm{Sl}_{\square,y}(A),\qquad\mathrm{Cm}_{\square}:=\mathrm{Cm}_{\square,y}(A):=-\left[F_y(A),h_{\square}\right],
\]
where $\mathrm{Sl}_{\square,y}(A)$ is the local inverse of Definition~\ref{def:folded-fields-local}, acting on the spinors on $T_k$ as in Section~\ref{sec:operators-cube-operators}. Both are linear operators on the space of spinors on $T_k$.

\begin{lemma}[Parametrix identities]\label{lem:parametrix-identities}
Let $\square$ and $\square'$ be $M$-cubes of $T_k$, and let $\mathsf{c}$, $\mathsf{c}'$ be lifts of their centres.
\begin{enumerate}
\item[(a)] $F\cdot h_{\square}\cdot\mathrm{Sl}_{\square}\cdot h_{\square}=h_{\square}^2-\mathrm{Cm}_{\square}\cdot\mathrm{Sl}_{\square}\cdot h_{\square}$.
\item[(b)] $h_{\square}\cdot\mathrm{Sl}_{\square}\cdot h_{\square}\cdot F=h_{\square}^2+h_{\square}\cdot\mathrm{Sl}_{\square}\cdot\mathrm{Cm}_{\square}$.
\item[(c)] The kernel $\mathrm{Cm}_{\square}(x,x')$ vanishes unless both $x$ and $x'$ belong to the orbit under $l\mathbb{Z}^3$ of the set $\left\{\tilde{x}\in\eta\mathbb{Z}^3:|\tilde{x}-\mathsf{c}|<\frac{3M}{4}+L\right\}$, and this set is contained in $W^{-}$. The operator $\mathrm{Cm}_{\square}$ depends on $A$ only through the values of $A$ on the bonds of $\tilde{W}$. Moreover $h_{\square'}\cdot\mathrm{Cm}_{\square}=0=\mathrm{Cm}_{\square}\cdot h_{\square'}$ unless $d(\mathsf{c},\mathsf{c}')\le M$.
\item[(d)] One has $N_{\gamma}(\mathrm{Cm}_{\square})\le\frac{\pi\cdot C_V}{M}$ and $N_{\gamma}\left(h_{\square}\cdot\mathrm{Sl}_{\square}\cdot h_{\square}\right)\le2a\cdot w_k(y)$; each of the numbers
\[
N_{\gamma}\left(\mathrm{Cm}_{\square}\cdot\mathrm{Sl}_{\square}\cdot h_{\square}\right),\qquad N_{\gamma}\left(h_{\square}\cdot\mathrm{Sl}_{\square}\cdot\mathrm{Cm}_{\square}\right),\qquad N_{\gamma}\left(\mathrm{Cm}_{\square}\cdot h_{\square'}\cdot\mathrm{Sl}_{\square'}\right)
\]
is at most $\frac{2\pi\cdot C_V\cdot a}{M}$; and if $M\ge1600$, then $N_{\gamma}\left(\delta_{\alpha,A}\left(h_{\square}\cdot\mathrm{Sl}_{\square}\right)\right)\le19a$.
\end{enumerate}
\end{lemma}

\begin{proof}
As noted after~\eqref{eq:operators-fermion-4}, $\gamma L\le\frac{1}{6}<\log2$, so that Lemma~\ref{lem:perturbations-commutators-holder} applies; in particular $\gamma\le\frac{1}{6L}\le\frac{1}{30}$, as $L\ge5$. By~\eqref{eq:operators-fermion-G1}, $\mathfrak{i}(A)<e_k^{\varepsilon}$, and $6L\cdot e_k^{\varepsilon}\le6L\cdot e_0^{\varepsilon}\le1$ by condition~\eqref{eq:renormalisation-group-20}; hence $3L\cdot\mathfrak{i}(A)\le1$, and in particular $3\mathfrak{i}(A)\le1$. For $\boldsymbol{\delta}:=2e_k^{\varepsilon}$ the conditions~\eqref{eq:operators-fermion-6} follow from conditions~\eqref{eq:renormalisation-group-20} and~\eqref{eq:renormalisation-group-21}, as $e_k\le e_0$, and the restriction of $A$ to the bonds of $\tilde{W}$ lies in $G^W(\boldsymbol{\delta})$ (Section~\ref{sec:operators-cube-operators}); so Proposition~\ref{prop:properties-local-inverse} applies to $\mathrm{Sl}_{\square}$ and to $\mathrm{Sl}_{\square'}$. Finally, by~\eqref{eq:operators-fermion-7} the real function $h_{\square}$ satisfies the hypothesis of Lemma~\ref{lem:perturbations-commutators-holder}(b) with $L_h:=\pi/M$.

(c) If $\mathrm{Cm}_{\square}(x,x')\ne0$, then Lemma~\ref{lem:perturbations-commutators-holder}(b) shows that $|x-x'|<L$ and $h_{\square}(x)\ne h_{\square}(x')$; in particular one of $x$, $x'$ lies in $\operatorname{supp}h_{\square}$. By~\eqref{eq:operators-fermion-7}, both sites then have lifts at distance less than $\frac{3M}{4}+L$ from $\mathsf{c}$. Since $M>4L$, one has $\frac{3M}{4}+L<\frac{3M}{2}-\eta$, so that these lifts lie in $W^{-}$. As $\mathrm{Cm}_{\square}=h_{\square}\cdot F_y(A)-F_y(A)\cdot h_{\square}$, the kernel of $\mathrm{Cm}_{\square}$ is
\[
\mathrm{Cm}_{\square}(x,x')=\left(h_{\square}(x)-h_{\square}(x')\right)\cdot F_y(A)(x,x').
\]
By what has just been shown, the rows of $\mathrm{Cm}_{\square}$ at the sites without a lift in $W^{-}$ are zero for every $A$. At a site $x$ with a lift in $W^{-}$, Proposition~\ref{prop:properties-local-inverse}(c) shows that the entries $F_y(A)(x,x')$, $x'$ arbitrary, are determined by the restriction of $A$ to the bonds of $\tilde{W}$, and the formula for the kernel shows the same for the entries $\mathrm{Cm}_{\square}(x,x')$. This proves the second assertion of (c). Let now $\mathrm{Cm}_{\square}(x,x')\ne0$, and suppose that $x$ or $x'$ lies in $\operatorname{supp}h_{\square'}$. Then $d(\mathsf{c},\mathsf{c}')<\frac{3M}{2}+L<2M$; being a multiple of $M$, the distance $d(\mathsf{c},\mathsf{c}')$ is then at most $M$. As $\left(h_{\square'}\cdot\mathrm{Cm}_{\square}\right)(x,x')=h_{\square'}(x)\cdot\mathrm{Cm}_{\square}(x,x')$ and $\left(\mathrm{Cm}_{\square}\cdot h_{\square'}\right)(x,x')=\mathrm{Cm}_{\square}(x,x')\cdot h_{\square'}(x')$, this proves the last assertion of (c).

(a) Let $f$ be a spinor on $T_k$ and $v:=\mathrm{Sl}_{\square}\cdot h_{\square}f$; by Section~\ref{sec:operators-cube-operators}, $v$ vanishes outside $\tilde{W}+l\mathbb{Z}^3$. One has
\[
F\left(h_{\square}v\right)=h_{\square}\cdot Fv+\left[F,h_{\square}\right]v=h_{\square}\cdot Fv-\mathrm{Cm}_{\square}v.
\]
By~\eqref{eq:operators-fermion-7} and (c), the support of $h_{\square}$ lies in the orbit of $W^{-}$. Proposition~\ref{prop:properties-local-inverse}(c) shows that for $x\in W^{-}$ the value $(Fv)(x)$ depends only on the restriction of $v$ to $\tilde{W}$; applied to the restriction of $h_{\square}f$ to $\tilde{W}$, the same proposition gives $(Fv)(x)=(h_{\square}f)(x)$ for $x\in W^{-}$. Hence $h_{\square}\cdot Fv=h_{\square}^2f$, and
\[
F\cdot h_{\square}\cdot\mathrm{Sl}_{\square}\cdot h_{\square}f=h_{\square}^2f-\mathrm{Cm}_{\square}\cdot\mathrm{Sl}_{\square}\cdot h_{\square}f.
\]

(b) Let $f$ be a spinor on $T_k$. Then $h_{\square}\cdot Ff=F\left(h_{\square}f\right)+\mathrm{Cm}_{\square}f$. Moreover $\mathrm{Sl}_{\square}\cdot F\left(h_{\square}f\right)=h_{\square}f$. For the proof of this identity let $F$ act also on all functions $\eta\mathbb{Z}^3\to\mathbb{C}^4$ as in Proposition~\ref{prop:properties-local-inverse}(c), the connection $A$ serving as the continuation of its restriction to the bonds of $\tilde{W}$; on the spinors on $T_k$, which are the antiperiodic functions on $\eta\mathbb{Z}^3$, this is the operator $F$, both being given by the formulas of Sections~\ref{sec:phase-factors-lattice} and~\ref{sec:operators-2}. As shown in the proof of Proposition~\ref{prop:properties-local-inverse}(c), by these formulas and Lemma~\ref{lem:block-averages-block}, the value of $Fu$ at a site $x$ depends only on the values of the function $u$ at the nearest neighbours of $x$ and on the block of side $L$ containing $x$. Let $g\colon\eta\mathbb{Z}^3\to\mathbb{C}^4$ be the function equal to $h_{\square}f$ on $W^{-}$ and to $0$ outside $W^{-}$. As the support of $h_{\square}$ lies in the orbit of $W^{-}$ (see (a)), the function $h_{\square}f-g$ vanishes outside the union of the sets $W^{-}+ln$ with $n\in\mathbb{Z}^3\setminus\{0\}$; these sets are contained in the sets $\tilde{W}+ln$, $n\ne0$, which are disjoint from $\tilde{W}$ (Section~\ref{sec:operators-cube-operators}). Let $x\in\tilde{W}$. The block of side $L$ containing $x$ lies in $\tilde{W}$, because $\tilde{W}$ is a union of blocks of side $L$ (Lemma~\ref{lem:fold-map-embedding}(a)). A nearest neighbour $x'$ of $x$ lies in no set $W^{-}+ln$ with $n\ne0$: otherwise
\[
l\le|ln|\le|x-\mathsf{c}|+|x'-x|+|x'-ln-\mathsf{c}|<\frac{3M}{2}+\eta+\left(\frac{3M}{2}-\eta\right)=3M,
\]
whereas $3M<l$. Hence $h_{\square}f$ and $g$ coincide on the block of side $L$ containing $x$ and at the nearest neighbours of $x$, so that $\left(F\left(h_{\square}f\right)\right)(x)=(Fg)(x)$. By Section~\ref{sec:operators-cube-operators}, the value of $\mathrm{Sl}_{\square}\cdot F\left(h_{\square}f\right)$ at $x$ is the value at $x$ of the local inverse of the restriction of $A$ to the bonds of $\tilde{W}$, applied to the restriction of $F\left(h_{\square}f\right)$ to $\tilde{W}$, that is, to the restriction of $Fg$ to $\tilde{W}$. Since $g$ vanishes outside $W^{-}$, the last assertion of Proposition~\ref{prop:properties-local-inverse}(c) shows that this value is $g(x)=\left(h_{\square}f\right)(x)$. Thus $\mathrm{Sl}_{\square}\cdot F\left(h_{\square}f\right)$ and $h_{\square}f$ coincide on $\tilde{W}$. Both are antiperiodic and vanish outside the orbit of $\tilde{W}$, by Section~\ref{sec:operators-cube-operators} and by~\eqref{eq:operators-fermion-7}; therefore they coincide. Applying $h_{\square}\cdot\mathrm{Sl}_{\square}$ to the first identity gives
\[
h_{\square}\cdot\mathrm{Sl}_{\square}\cdot h_{\square}\cdot Ff=h_{\square}^2f+h_{\square}\cdot\mathrm{Sl}_{\square}\cdot\mathrm{Cm}_{\square}f.
\]

(d) By Lemma~\ref{lem:perturbations-commutators-holder}(b) with $L_h=\pi/M$,
\[
N_{\gamma}(\mathrm{Cm}_{\square})=N_{\gamma}\left(\left[F,h_{\square}\right]\right)\le\frac{\pi\cdot C_V}{M}.
\]
By Proposition~\ref{prop:properties-local-inverse}(b) and Section~\ref{sec:operators-cube-operators}, $N_{\gamma}(\mathrm{Sl}_{\square})\le2a\cdot w_k(y)$ on $T_k$, and the same holds for $\mathrm{Sl}_{\square'}$; by~\eqref{eq:operators-fermion-EW3}, $N_{\gamma}(h_{\square})\le1$, since $0\le h_{\square}\le1$; and $w_k(y)\le1$ by~\eqref{eq:operators-fermion-2}. By~\eqref{eq:operators-fermion-EW2} it follows that $N_{\gamma}\left(h_{\square}\cdot\mathrm{Sl}_{\square}\cdot h_{\square}\right)\le2a\cdot w_k(y)$, and that each of $N_{\gamma}\left(\mathrm{Cm}_{\square}\cdot\mathrm{Sl}_{\square}\cdot h_{\square}\right)$, $N_{\gamma}\left(h_{\square}\cdot\mathrm{Sl}_{\square}\cdot\mathrm{Cm}_{\square}\right)$ and $N_{\gamma}\left(\mathrm{Cm}_{\square}\cdot h_{\square'}\cdot\mathrm{Sl}_{\square'}\right)$ is at most
\[
\frac{\pi\cdot C_V}{M}\cdot2a\cdot w_k(y)\le\frac{2\pi\cdot C_V\cdot a}{M}.
\]
For the last bound let $f$ be a spinor on $T_k$. The first identity of Lemma~\ref{lem:perturbations-commutators-holder}(c), applied to the spinor $\mathrm{Sl}_{\square}f$, gives
\[
\delta_{\alpha,A}\left(h_{\square}\cdot\mathrm{Sl}_{\square}\right)=h_{\square}(x)\,\delta_{\alpha,A}\mathrm{Sl}_{\square}+E_g\mathrm{Sl}_{\square},\qquad g(x,x'):=|x-x'|^{-\alpha}\left(h_{\square}(x')-h_{\square}(x)\right)U_A(x,x'),
\]
where $h_{\square}(x)\,\delta_{\alpha,A}\mathrm{Sl}_{\square}$ maps $f$ to the function $(x,x')\mapsto h_{\square}(x)\cdot\left(\delta_{\alpha,A}\mathrm{Sl}_{\square}f\right)(x,x')$ of the short pairs, and $E_g$ is the operator of Lemma~\ref{lem:perturbations-commutators-holder}(c). For a short pair one has $|x-x'|<1$ and $\sum_{\mu}|x_{\mu}-x'_{\mu}|\le3|x-x'|$; by~\eqref{eq:operators-fermion-7}, Lemma~\ref{lem:bounds-phases}(a) and $3\mathfrak{i}(A)\le1$,
\[
|g(x,x')|\le\frac{3\pi}{M}\cdot|x-x'|^{1-\alpha}\cdot\exp\left(3\mathfrak{i}(A)\right)\le\frac{3\pi\cdot\mathbf{e}}{M}.
\]
If $g(x,x')\ne0$ or $h_{\square}(x)\ne0$, then one of the sites $x$, $x'$ lies in $\operatorname{supp}h_{\square}$ and the other one at distance less than $1$ from it; by~\eqref{eq:operators-fermion-7} both lie in $\tilde{W}+l\mathbb{Z}^3$, as $\frac{3M}{4}+1<\frac{3M}{2}$. Hence only the values of $\delta_{\alpha,A}\mathrm{Sl}_{\square}f$ on short pairs in $\tilde{W}$ enter, and $h_{\square}(x)\,\delta_{\alpha,A}\mathrm{Sl}_{\square}=h_{\square}\cdot O$; here $O$ is the operator $\delta_{\alpha,A}\mathrm{Sl}_{\square,y}(A)$ of Proposition~\ref{prop:properties-local-inverse}(b), formed with the restriction of $A$ to the bonds of $\tilde{W}$ and acting on the spinors on $T_k$ as in Section~\ref{sec:operators-cube-operators}, and $h_{\square}$ acts on the functions of short pairs through its value at the first site. By Proposition~\ref{prop:properties-local-inverse}(b) and Section~\ref{sec:operators-cube-operators}, $N_{\gamma}(O)\le18a$. The bound $N_{\gamma}\left(h_{\square}\cdot O\right)\le|h_{\square}|_{\infty}\cdot N_{\gamma}(O)$ of Section~\ref{sec:norms-exponential-weights} for operators with values in the functions of short pairs, together with $0\le h_{\square}\le1$, gives $N_{\gamma}\left(h_{\square}\cdot O\right)\le18a$. Lemma~\ref{lem:perturbations-commutators-holder}(c) with $g_0:=\frac{3\pi\cdot\mathbf{e}}{M}$, together with $N_{\gamma}(\mathrm{Sl}_{\square})\le2a\cdot w_k(y)\le2a$, gives $N_{\gamma}\left(E_g\mathrm{Sl}_{\square}\right)\le27\exp(\gamma)\cdot g_0\cdot2a$. By~\eqref{eq:operators-fermion-EW2}, applied to the sum of these two operators, it follows that
\[
N_{\gamma}\left(\delta_{\alpha,A}\left(h_{\square}\cdot\mathrm{Sl}_{\square}\right)\right)\le18a+27\exp(\gamma)\cdot\frac{3\pi\cdot\mathbf{e}}{M}\cdot2a=18a+\frac{162\pi\cdot\mathbf{e}\cdot\exp(\gamma)}{M}\cdot a\le19a
\]
for $M\ge1600$, because $162\pi\cdot\mathbf{e}\cdot\exp(\gamma)\le162\pi\cdot\mathbf{e}\cdot\exp\left(\frac{1}{30}\right)\le1600$.
\end{proof}

\subsubsection{Expansion in walks}\label{sec:expansion-walks}

\begin{definition}[Walks]\label{def:walks}
Let $n\ge0$ be an integer. By a \emph{walk} of length $|\boldsymbol{\omega}|=n$ one means a sequence $\boldsymbol{\omega}=(\square_0,\dots,\square_n)$ of $M$-cubes of $T_k$ in which the centres $\mathsf{c}_0,\dots,\mathsf{c}_n$ of consecutive cubes satisfy $d(\mathsf{c}_{i-1},\mathsf{c}_i)\le M$ ($1\le i\le n$). The walks of length $n$ form the set $\boldsymbol{\Omega}_n$, and $\boldsymbol{\Omega}:=\bigcup_{n\ge0}\boldsymbol{\Omega}_n$. For $\boldsymbol{\omega}\in\boldsymbol{\Omega}_n$ let $h_i:=h_{\square_i}$, $\mathrm{Sl}_i:=\mathrm{Sl}_{\square_i,y}(A)$ and $\mathrm{Cm}_i:=\mathrm{Cm}_{\square_i,y}(A)$, and define the linear operators $S_{y,\boldsymbol{\omega}}(A)$, $S'_{y,\boldsymbol{\omega}}(A)$ on the spinors on $T_k$ and the set $X_{\boldsymbol{\omega}}$ by
\begin{align*}
S_{y,\boldsymbol{\omega}}(A)&:=\left(h_0\cdot\mathrm{Sl}_0\cdot h_0\right)\cdot\left(\mathrm{Cm}_1\cdot\mathrm{Sl}_1\cdot h_1\right)\cdot\left(\mathrm{Cm}_2\cdot\mathrm{Sl}_2\cdot h_2\right)\cdots\left(\mathrm{Cm}_n\cdot\mathrm{Sl}_n\cdot h_n\right),\\
S'_{y,\boldsymbol{\omega}}(A)&:=\left(h_n\cdot\mathrm{Sl}_n\right)\cdot\left(-\mathrm{Cm}_n\cdot h_{n-1}\cdot\mathrm{Sl}_{n-1}\right)\cdot\left(-\mathrm{Cm}_{n-1}\cdot h_{n-2}\cdot\mathrm{Sl}_{n-2}\right)\cdots\\
&\qquad\qquad\cdots\left(-\mathrm{Cm}_1\cdot h_0\cdot\mathrm{Sl}_0\right)\cdot h_0,\\
X_{\boldsymbol{\omega}}&:=\bigcup_{0\le i\le n}\square_i^{\sim1}.
\end{align*}
For $n=0$ the two products are $h_0\cdot\mathrm{Sl}_0\cdot h_0$ and $\left(h_0\cdot\mathrm{Sl}_0\right)\cdot h_0$.
\end{definition}

The set $X_{\boldsymbol{\omega}}$ belongs to $D_k$, that is, it is a polymer. Indeed, every enlarged cube $\square_i^{\sim1}$ is connected through faces, and two consecutive ones, $\square_{i-1}^{\sim1}$ and $\square_i^{\sim1}$, have the cube $\square_i$ in common. Since each $\square_i^{\sim1}$ consists of $27$ $M$-cubes,
\[
|X_{\boldsymbol{\omega}}|_M\le27(n+1),\qquad d_M(X_{\boldsymbol{\omega}})\le|X_{\boldsymbol{\omega}}|_M-1\le27(n+1).
\]
For the middle inequality let $X$ be a polymer, and consider the graph whose vertices are the $M$-cubes of $X$, two of them being adjacent when they have a common face; it is connected. Choose a spanning tree of this graph, and connect the centres of any two cubes that are adjacent in this tree by a straight segment of length $M$. The union of these segments, or the centre of the only cube of $X$ if $|X|_M=1$, is a net of length $\left(|X|_M-1\right)M$; it contains the centre of each $M$-cube of $X$, hence it meets every closed $M$-cube of $X$, and the definition of $d_M$ gives $M\cdot d_M(X)\le\left(|X|_M-1\right)M$.

\begin{lemma}[Walks of a given length]\label{lem:walks-given-length}
Let $w\ge1$ and assume
\[
q_w:=\frac{54\pi\cdot C_V\cdot a\cdot w}{M}\le\frac{1}{2}.
\]
Let $n\ge0$ be an integer, and let $S^{\#}_{\boldsymbol{\omega}}$ denote either $S_{y,\boldsymbol{\omega}}(A)$ for all $\boldsymbol{\omega}\in\boldsymbol{\Omega}_n$, or $S'_{y,\boldsymbol{\omega}}(A)$ for all $\boldsymbol{\omega}\in\boldsymbol{\Omega}_n$. Then, $z$ and $z'$ ranging over the unit sites of $T_k$,
\begin{align*}
\sup_z\sum_{z'}\exp\left(\gamma\cdot d(z,z')\right)\sum_{\boldsymbol{\omega}\in\boldsymbol{\Omega}_n}w^{n+1}\left\|\mathbf{1}_{\Delta_z}S^{\#}_{\boldsymbol{\omega}}\mathbf{1}_{\Delta_{z'}}\right\|&\le16a\cdot w\cdot w_k(y)\cdot q_w^n,\\
\sup_z\sum_{z'}\exp\left(\gamma\cdot d(z,z')\right)\sum_{\boldsymbol{\omega}\in\boldsymbol{\Omega}_n}w^{n+1}\left\|\mathbf{1}_{\Delta_z}\delta_{\alpha,A}S^{\#}_{\boldsymbol{\omega}}\mathbf{1}_{\Delta_{z'}}\right\|&\le152a\cdot w\cdot q_w^n.
\end{align*}
\end{lemma}

\begin{proof}
Let $\boldsymbol{\omega}=(\square_0,\dots,\square_n)\in\boldsymbol{\Omega}_n$, with the notation of Definition~\ref{def:walks}. If $S^{\#}_{\boldsymbol{\omega}}=S_{y,\boldsymbol{\omega}}(A)$, put $O_0:=h_0\cdot\mathrm{Sl}_0\cdot h_0$ and $O_i:=\mathrm{Cm}_i\cdot\mathrm{Sl}_i\cdot h_i$ for $1\le i\le n$, so that $S^{\#}_{\boldsymbol{\omega}}=O_0\cdot O_1\cdots O_n$. If $S^{\#}_{\boldsymbol{\omega}}=S'_{y,\boldsymbol{\omega}}(A)$, put $O_0:=h_n\cdot\mathrm{Sl}_n$ and $O_i:=-\mathrm{Cm}_{n+1-i}\cdot h_{n-i}\cdot\mathrm{Sl}_{n-i}$ for $1\le i\le n$, so that $S^{\#}_{\boldsymbol{\omega}}=O_0\cdot O_1\cdots O_n\cdot h_0$.

Fix a unit site $z$. Insert $1=\sum_{z_i}\mathbf{1}_{\Delta_{z_i}}$, with $z_i$ ranging over the unit sites, between $O_{i-1}$ and $O_i$ for $1\le i\le n$, and estimate the resulting sum term by term, using $d(z,z')\le d(z,z_1)+d(z_1,z_2)+\dots+d(z_n,z')$ as in the proof of~\eqref{eq:operators-fermion-EW2}; this gives
\[
\sum_{z'}\exp\left(\gamma\cdot d(z,z')\right)\left\|\mathbf{1}_{\Delta_z}S^{\#}_{\boldsymbol{\omega}}\mathbf{1}_{\Delta_{z'}}\right\|\le\prod_{i=0}^nN_{\gamma}(O_i),
\]
where in the primed case $N_{\gamma}(h_0)\le1$, by~\eqref{eq:operators-fermion-EW3}, has been used for the last factor. The left side vanishes unless $\mathbf{1}_{\Delta_z}O_0\ne0$, and this requires that $\Delta_z$ meet the support of the leftmost factor of $O_0$, which is $h_0$, resp.\ $h_n$. For $\delta_{\alpha,A}S^{\#}_{\boldsymbol{\omega}}=\left(\delta_{\alpha,A}O_0\right)\cdot O_1\cdots O_n$ (followed by $h_0$ in the primed case) the same holds with $\delta_{\alpha,A}O_0$ in place of $O_0$, and the left side vanishes unless $\Delta_z$ or a unit block $\Delta_{\tilde{z}}$ with $|\tilde{z}-z|\le1$ meets the support of this leftmost factor, because a short pair $(x,x')$ with $x\in\Delta_z$ has $x'$ in such a block. In both cases a lift $\mathsf{c}$ of the centre of the cube at the left end ($\square_0$, resp.\ $\square_n$) satisfies $|z_{\mu}-\mathsf{c}_{\mu}|<\frac{3M}{4}+\frac{3}{2}$ for $\mu=0,1,2$; since an open interval of length $\frac{3M}{2}+3<2M$ contains at most two points of $M\mathbb{Z}$, there are at most $8$ choices for this cube (for $S^{\#}_{\boldsymbol{\omega}}$ itself this is the count following~\eqref{eq:operators-fermion-7}). Once it is chosen, there are at most $27$ choices for each further cube of the walk ($\square_1,\dots,\square_n$, resp.\ $\square_{n-1},\dots,\square_0$), because consecutive cubes of a walk have centres at distance at most $M$. Thus for each $z$ at most $8\cdot27^n$ walks of $\boldsymbol{\Omega}_n$ contribute.

Lemma~\ref{lem:parametrix-identities}(d) bounds the factors; for the factor $O_0=h_n\cdot\mathrm{Sl}_n$ of the primed case one uses instead Proposition~\ref{prop:properties-local-inverse}(b), Section~\ref{sec:operators-cube-operators}, \eqref{eq:operators-fermion-EW2} and~\eqref{eq:operators-fermion-EW3}. If $S^{\#}_{\boldsymbol{\omega}}=S_{y,\boldsymbol{\omega}}(A)$, then $N_{\gamma}(O_0)\le2a\cdot w_k(y)$ and $N_{\gamma}(O_i)\le\frac{2\pi\cdot C_V\cdot a}{M}$ for $1\le i\le n$. If $S^{\#}_{\boldsymbol{\omega}}=S'_{y,\boldsymbol{\omega}}(A)$ and $n\ge1$, then $N_{\gamma}(O_0)=N_{\gamma}\left(h_n\cdot\mathrm{Sl}_n\right)\le2a\cdot w_k(y)$ and $N_{\gamma}(O_i)=N_{\gamma}\left(\mathrm{Cm}_{n+1-i}\cdot h_{n-i}\cdot\mathrm{Sl}_{n-i}\right)\le\frac{2\pi\cdot C_V\cdot a}{M}$ for $1\le i\le n$, while the last factor $h_0$ satisfies $N_{\gamma}(h_0)\le1$; for $n=0$ one has $S'_{y,\boldsymbol{\omega}}(A)=S_{y,\boldsymbol{\omega}}(A)$. Consequently the left side of the first inequality is at most
\[
8\cdot27^n\cdot w^{n+1}\cdot2a\cdot w_k(y)\cdot\left(\frac{2\pi\cdot C_V\cdot a}{M}\right)^n=16a\cdot w\cdot w_k(y)\cdot q_w^n.
\]
In the second inequality the left factor is $\delta_{\alpha,A}O_0$, and Lemma~\ref{lem:parametrix-identities}(d), together with~\eqref{eq:operators-fermion-EW2} and $N_{\gamma}(h_0)\le1$ when $O_0=h_0\cdot\mathrm{Sl}_0\cdot h_0$ (then $\delta_{\alpha,A}O_0=\delta_{\alpha,A}\left(h_0\cdot\mathrm{Sl}_0\right)\cdot h_0$), gives $N_{\gamma}\left(\delta_{\alpha,A}O_0\right)\le19a$ in place of $2a\cdot w_k(y)$; the factor in front becomes $8\cdot19a=152a$. Lemma~\ref{lem:parametrix-identities}(d) applies here, since $M\ge1600$: from $q_w\le\frac{1}{2}$ and $w\ge1$ one gets $M\ge108\pi\cdot C_V\cdot a\cdot w\ge108\pi\cdot C_V\cdot a$, and $108\pi\cdot C_V\cdot a\ge1600$, because $C_V=12\mathbf{e}\cdot(2+b)\ge24\mathbf{e}$ and $a\ge1$. Here $a=C_z\cdot c(\gamma)\ge1$ by~\eqref{eq:operators-fermion-4}, since $C_z\ge C_5\ge1$ by Proposition~\ref{prop:zero-field-levels}, and $c(\gamma)\ge1$, the summand at the origin in the definition of $c(\gamma)$ in Section~\ref{sec:norms-exponential-weights} being $1$.
\end{proof}

\begin{theorem}[Expansion of $S_y(A)$ in walks]\label{thm:expansion-walks}
Let $L\ge5$ be odd, let $M=L^m$ satisfy $M>4L$ and $M\ge108\pi\cdot C_V\cdot a$, and let conditions~\eqref{eq:renormalisation-group-20} and~\eqref{eq:renormalisation-group-21} hold. Let $k\ge0$, $l\in\mathrm{Cl}$, $\mathfrak{m}\in(0,1]$, $e_k\in(0,e_0]$, $y\in Y_k$ and $A\in\tilde{G}_k$. Then the following hold.
\begin{enumerate}
\item[(a)] $F_y(A)$ is a bijection of the space of spinors on $T_k$ onto itself, and
\[
S_y(A)=\sum_{\boldsymbol{\omega}\in\boldsymbol{\Omega}}S_{y,\boldsymbol{\omega}}(A)=\sum_{\boldsymbol{\omega}\in\boldsymbol{\Omega}}S'_{y,\boldsymbol{\omega}}(A),
\]
where each of the two series converges absolutely with respect to $N_{\gamma}$, uniformly in $A\in\tilde{G}_k$.
\item[(b)] The map $A\mapsto S_y(A)$ is analytic on $\tilde{G}_k$.
\item[(c)] $N_{\gamma}\left(S_y(A)\right)\le32a\cdot w_k(y)$ and $N_{\gamma}\left(\delta_{\alpha,A}S_y(A)\right)\le304a$. In particular, for all unit sites $z$, $z'$ and every spinor $f$ on $T_k$,
\begin{align*}
\left|\mathbf{1}_{\Delta_z}S_y(A)\mathbf{1}_{\Delta_{z'}}f\right|_{\infty}&\le32a\cdot w_k(y)\cdot\exp\left(-\gamma\cdot d(z,z')\right)\cdot|f|_{\infty},\\
\left|\mathbf{1}_{\Delta_z}\delta_{\alpha,A}S_y(A)\mathbf{1}_{\Delta_{z'}}f\right|_{\infty}&\le304a\cdot\exp\left(-\gamma\cdot d(z,z')\right)\cdot|f|_{\infty},
\end{align*}
where in the second inequality the supremum is taken over the short pairs $(x,x')$.
\end{enumerate}
\end{theorem}

\begin{proof}
The hypotheses of Sections~\ref{sec:parametrix-local-inverses} and~\ref{sec:expansion-walks} are satisfied, and $M\ge108\pi\cdot C_V\cdot a$ means $q_1\le\frac{1}{2}$, where $q_1=\frac{54\pi\cdot C_V\cdot a}{M}$ is the number $q_w$ of Lemma~\ref{lem:walks-given-length} for $w=1$. Let
\[
S^{*}:=\sum_{\square}h_{\square}\cdot\mathrm{Sl}_{\square}\cdot h_{\square},\qquad\mathfrak{R}:=\sum_{\square}\mathrm{Cm}_{\square}\cdot\mathrm{Sl}_{\square}\cdot h_{\square},
\]
the sums being over all $M$-cubes of $T_k$; these are linear operators on the spinors on $T_k$. By Lemma~\ref{lem:parametrix-identities}(a) and the first relation of~\eqref{eq:operators-fermion-7}, $F\cdot S^{*}=1-\mathfrak{R}$.

If $\mathbf{1}_{\Delta_z}\mathrm{Cm}_{\square}\ne0$ for a unit site $z$, then $|z_{\mu}-\mathsf{c}_{\mu}|<\frac{3M}{4}+L+\frac{1}{2}$ for $\mu=0,1,2$ and a suitable lift $\mathsf{c}$ of the centre of $\square$, by Lemma~\ref{lem:parametrix-identities}(c); an open interval of length $\frac{3M}{2}+2L+1<3M$ contains at most three points of $M\mathbb{Z}$. So for each $z$ at most $27$ cubes contribute to $\mathbf{1}_{\Delta_z}\mathfrak{R}$, and Lemma~\ref{lem:parametrix-identities}(d) gives
\[
N_{\gamma}(\mathfrak{R})\le27\cdot\frac{2\pi\cdot C_V\cdot a}{M}=q_1\le\frac{1}{2}.
\]
By~\eqref{eq:operators-fermion-EW1} and~\eqref{eq:operators-fermion-EW2}, $N_{\gamma}$ is a submultiplicative norm on the linear operators of the finite-dimensional space of spinors on $T_k$. By the Neumann series, $1-\mathfrak{R}$ is invertible and $(1-\mathfrak{R})^{-1}=\sum_{n\ge0}\mathfrak{R}^n$, the series converging in $N_{\gamma}$; moreover $F\cdot S^{*}\cdot(1-\mathfrak{R})^{-1}=1$. Since $F$ acts on a finite-dimensional space and has a right inverse, it is bijective, and
\[
S_y(A)=F^{-1}=S^{*}\cdot\sum_{n\ge0}\mathfrak{R}^n.
\]
When the sums defining $S^{*}$ and $\mathfrak{R}$ are multiplied out, $S^{*}\cdot\mathfrak{R}^n$ becomes the sum, over all sequences $(\square_0,\dots,\square_n)$ of $M$-cubes of $T_k$, of the products $\left(h_0\cdot\mathrm{Sl}_0\cdot h_0\right)\cdot\left(\mathrm{Cm}_1\cdot\mathrm{Sl}_1\cdot h_1\right)\cdots\left(\mathrm{Cm}_n\cdot\mathrm{Sl}_n\cdot h_n\right)$, where $h_i:=h_{\square_i}$, $\mathrm{Sl}_i:=\mathrm{Sl}_{\square_i}$ and $\mathrm{Cm}_i:=\mathrm{Cm}_{\square_i}$, and $\mathsf{c}_i$ denotes the centre of $\square_i$. A sequence that is not a walk has an index $i\in\{1,\dots,n\}$ with $d(\mathsf{c}_{i-1},\mathsf{c}_i)>M$; its product then contains the adjacent factors $h_{i-1}$ and $\mathrm{Cm}_i$, and $h_{i-1}\cdot\mathrm{Cm}_i=0$ by the last assertion of Lemma~\ref{lem:parametrix-identities}(c). Only the walks are left, and $S^{*}\cdot\mathfrak{R}^n=\sum_{\boldsymbol{\omega}\in\boldsymbol{\Omega}_n}S_{y,\boldsymbol{\omega}}(A)$.

Lemma~\ref{lem:walks-given-length} with $w=1$, summed over $n$ with $\sum_{n\ge0}q_1^n\le2$, gives
\begin{align*}
\sup_z\sum_{z'}\exp\left(\gamma\cdot d(z,z')\right)\sum_{\boldsymbol{\omega}\in\boldsymbol{\Omega}}\left\|\mathbf{1}_{\Delta_z}S_{y,\boldsymbol{\omega}}(A)\mathbf{1}_{\Delta_{z'}}\right\|&\le32a\cdot w_k(y),\\
\sup_z\sum_{z'}\exp\left(\gamma\cdot d(z,z')\right)\sum_{\boldsymbol{\omega}\in\boldsymbol{\Omega}}\left\|\mathbf{1}_{\Delta_z}\delta_{\alpha,A}S_{y,\boldsymbol{\omega}}(A)\mathbf{1}_{\Delta_{z'}}\right\|&\le304a,
\end{align*}
and, for every integer $N\ge0$ and every $A\in\tilde{G}_k$, the same left sides restricted to the walks of length at least $N$ are at most $32a\cdot q_1^N$ and $304a\cdot q_1^N$. As $T_k$ has finitely many unit sites, the sums $\sum_{\boldsymbol{\omega}\in\boldsymbol{\Omega}}N_{\gamma}\left(S_{y,\boldsymbol{\omega}}(A)\right)$ and $\sum_{\boldsymbol{\omega}\in\boldsymbol{\Omega}}N_{\gamma}\left(\delta_{\alpha,A}S_{y,\boldsymbol{\omega}}(A)\right)$ are finite, and their parts over the walks of length at least $N$ tend to $0$ as $N\to\infty$, uniformly in $A\in\tilde{G}_k$. So the series $\sum_{\boldsymbol{\omega}\in\boldsymbol{\Omega}}S_{y,\boldsymbol{\omega}}(A)$ converges absolutely in $N_{\gamma}$, uniformly on $\tilde{G}_k$; its sum is $S_y(A)$; and the two displayed bounds give $N_{\gamma}\left(S_y(A)\right)\le32a\cdot w_k(y)$ and $N_{\gamma}\left(\delta_{\alpha,A}S_y(A)\right)\le304a$. The second form of (c) follows from the first by~\eqref{eq:operators-fermion-EW1}.

Starting from Lemma~\ref{lem:parametrix-identities}(b) instead of (a), one obtains likewise from~\eqref{eq:operators-fermion-7} that $S^{*}\cdot F=1-\mathfrak{R}'$, where
\[
\mathfrak{R}':=-\sum_{\square}h_{\square}\cdot\mathrm{Sl}_{\square}\cdot\mathrm{Cm}_{\square}.
\]
For each unit site $z$ one has $\mathbf{1}_{\Delta_z}h_{\square}\ne0$ for at most $8$ cubes, by the count following~\eqref{eq:operators-fermion-7}; so $N_{\gamma}(\mathfrak{R}')\le8\cdot\frac{2\pi\cdot C_V\cdot a}{M}\le q_1$ by Lemma~\ref{lem:parametrix-identities}(d). Hence $1-\mathfrak{R}'$ is invertible, and $S_y(A)=(1-\mathfrak{R}')^{-1}\cdot S^{*}=\sum_{n\ge0}(\mathfrak{R}')^n\cdot S^{*}$. Multiplying out the sums, $(\mathfrak{R}')^n\cdot S^{*}$ becomes the sum, over all sequences $(\square_0,\dots,\square_n)$ of $M$-cubes, of the products
\begin{align*}
&\left(-h_n\cdot\mathrm{Sl}_n\cdot\mathrm{Cm}_n\right)\cdots\left(-h_1\cdot\mathrm{Sl}_1\cdot\mathrm{Cm}_1\right)\cdot\left(h_0\cdot\mathrm{Sl}_0\cdot h_0\right)\\
&\qquad=\left(h_n\cdot\mathrm{Sl}_n\right)\cdot\left(-\mathrm{Cm}_n\cdot h_{n-1}\cdot\mathrm{Sl}_{n-1}\right)\cdots\left(-\mathrm{Cm}_1\cdot h_0\cdot\mathrm{Sl}_0\right)\cdot h_0,
\end{align*}
with $h_i$, $\mathrm{Sl}_i$, $\mathrm{Cm}_i$ and $\mathsf{c}_i$ as before. If $d(\mathsf{c}_{i-1},\mathsf{c}_i)>M$ for some $i\in\{1,\dots,n\}$, such a product has the neighbouring factors $\mathrm{Cm}_i$ and $h_{i-1}$, and $\mathrm{Cm}_i\cdot h_{i-1}=0$ by Lemma~\ref{lem:parametrix-identities}(c); so again only the walks remain, and $(\mathfrak{R}')^n\cdot S^{*}=\sum_{\boldsymbol{\omega}\in\boldsymbol{\Omega}_n}S'_{y,\boldsymbol{\omega}}(A)$. Lemma~\ref{lem:walks-given-length} for the primed products then gives the absolute convergence of $\sum_{\boldsymbol{\omega}\in\boldsymbol{\Omega}}S'_{y,\boldsymbol{\omega}}(A)$ in $N_{\gamma}$, uniformly on $\tilde{G}_k$, exactly as above. This proves (a) and (c).

(b) Restricting $A$ to the bonds of $\tilde{W}$ is a linear map, which sends $\tilde{G}_k$ into $G^W\left(2e_k^{\varepsilon}\right)$ (Section~\ref{sec:operators-cube-operators}), and on $G^W\left(2e_k^{\varepsilon}\right)$ the local inverse $\mathrm{Sl}_{\square,y}$ is analytic by Proposition~\ref{prop:properties-local-inverse}(a). The entries of $\mathrm{Cm}_{\square}$ are constant multiples of entries of $F_y(A)$, hence entire functions of $A$. Therefore every $S_{y,\boldsymbol{\omega}}(A)$, a finite product of such operators and of multiplication operators, depends analytically on $A\in\tilde{G}_k$. By (a) the series $\sum_{\boldsymbol{\omega}\in\boldsymbol{\Omega}}S_{y,\boldsymbol{\omega}}(A)$ converges uniformly on $\tilde{G}_k$, and by the theorem of Weierstrass a uniformly convergent series of analytic functions has an analytic sum. Hence $A\mapsto S_y(A)$ is analytic on $\tilde{G}_k$.
\end{proof}

Let $k\ge1$. By Section~\ref{sec:operators-2}, $S_k(A)=S_{\infty}(A)$ and $S_{k,y}(A)=S_y(A)$, and $S^0_{k+1}(A)=S_0(A)$ by~\eqref{eq:operators-fermion-3}. Hence Proposition~\ref{prop:decay-fine-propagator} follows from Theorem~\ref{thm:expansion-walks} with $y=\infty$, and Proposition~\ref{prop:decay-resolvents} follows from it with $y\ge0$. Let $k=0$. Then every unit block consists of a single site, $w_0(y)=(1+|y|)^{-1}$ by~\eqref{eq:operators-fermion-2}, and $S_{0,y}(A)=S_y(A)$ (Section~\ref{sec:operators-2}); so the second form of Theorem~\ref{thm:expansion-walks}(c) gives
\[
\left|S_{0,y}(A)(x,x')\right|\le32a\cdot(1+|y|)^{-1}\cdot\exp\left(-\gamma\cdot d(x,x')\right)
\]
for all sites $x$, $x'$, in particular for $\Gamma_0(A)=S_{0,0}(A)$. The bounds of level $0$ in Proposition~\ref{prop:kernel-bounds-2} follow.

\subsection{Weakened walk expansions and polymer propagators}\label{sec:weakened-walk-expansions}

\begin{definition}[Weakened sums over walks]\label{def:weakened-sums-walks}
Let $s=(s_c)_{c\in\pi_k}\in\mathrm{Pd}_k$, $X\in D_k$, $y\in Y_k$ and $A\in\tilde{G}_k$. To every walk $\boldsymbol{\omega}\in\boldsymbol{\Omega}$ (Definition~\ref{def:walks}) attach the coefficient
\[
s_{\boldsymbol{\omega}}:=\prod\left\{s_c:\ c\in\pi_k,\ c\subset X_{\boldsymbol{\omega}}\right\}.
\]
Thus, for every walk $\boldsymbol{\omega}$, those of length zero included, the coefficient $s_{\boldsymbol{\omega}}$ contains the parameter $s_c$ of each $M$-cube $c$ of $X_{\boldsymbol{\omega}}$. Define the following operators on the spinors of $T_k$:
\begin{align*}
S_y(s,A)&:=\sum_{\boldsymbol{\omega}\in\boldsymbol{\Omega}}s_{\boldsymbol{\omega}}\cdot S_{y,\boldsymbol{\omega}}(A),\qquad&S_y(X,A)&:=\sum_{\boldsymbol{\omega}\in\boldsymbol{\Omega}:\ X_{\boldsymbol{\omega}}=X}S_{y,\boldsymbol{\omega}}(A),\\
S'_y(s,A)&:=\sum_{\boldsymbol{\omega}\in\boldsymbol{\Omega}}s_{\boldsymbol{\omega}}\cdot S'_{y,\boldsymbol{\omega}}(A),\qquad&S'_y(X,A)&:=\sum_{\boldsymbol{\omega}\in\boldsymbol{\Omega}:\ X_{\boldsymbol{\omega}}=X}S'_{y,\boldsymbol{\omega}}(A),
\end{align*}
and the means
\begin{align*}
S^{\mathrm{sy}}_y(s,A)&:=\frac{1}{2}\left(S_y(s,A)+S'_y(s,A)\right),\qquad S^{\mathrm{sy}}_y(X,A):=\frac{1}{2}\left(S_y(X,A)+S'_y(X,A)\right),\\
S^{\mathrm{sy}}_{y,\boldsymbol{\omega}}(A)&:=\frac{1}{2}\left(S_{y,\boldsymbol{\omega}}(A)+S'_{y,\boldsymbol{\omega}}(A)\right)\qquad(\boldsymbol{\omega}\in\boldsymbol{\Omega}).
\end{align*}
Under the hypotheses of Proposition~\ref{prop:properties-weakened-sums} these series converge absolutely (Proposition~\ref{prop:properties-weakened-sums}(a)).
\end{definition}

\begin{proposition}[Properties of the weakened sums]\label{prop:properties-weakened-sums}
Let $L$, $M$ and $e_0$ satisfy the hypotheses of Theorem~\ref{thm:expansion-walks}, let $0<\alpha_0\le\frac{1}{54}$ and $\kappa'>0$, and let $M$ satisfy condition~\eqref{eq:renormalisation-group-19}. Put
\[
w_s:=M^{27\alpha_0},\qquad w_X:=\exp(27\kappa'),
\]
and let $S^{\#}$ stand for any one of $S$, $S'$ and $S^{\mathrm{sy}}$. Let $k\ge0$, $l\in\mathrm{Cl}$, $\mathfrak{m}\in(0,1]$, $e_k\in(0,e_0]$, $y\in Y_k$, $A\in\tilde{G}_k$, $s\in\mathrm{Pd}_k$ and $X\in D_k$.
\begin{enumerate}
\item[(a)] \emph{Bounds.} The series of Definition~\ref{def:weakened-sums-walks} converge absolutely in the norm $N_{\gamma}$, uniformly on $\mathrm{Pd}_k\times\tilde{G}_k$. For all unit sites $z$ and $z'$
\begin{align*}
N_{\gamma}\left(S^{\#}_y(s,A)\right)&\le 32a\cdot w_s\cdot w_k(y),\qquad N_{\gamma}\left(\delta_{\alpha,A}S^{\#}_y(s,A)\right)\le 304a\cdot w_s,\\
\left\|\mathbf{1}_{\Delta_z}S^{\#}_y(X,A)\,\mathbf{1}_{\Delta_{z'}}\right\|&\le 32a\cdot w_X\cdot w_k(y)\cdot\exp\left(-\gamma\cdot d(z,z')-\kappa'\cdot d_M(X)\right).
\end{align*}
\item[(b)] \emph{Normalisation.} Let $1$ also denote the family $s\in\mathrm{Pd}_k$ with $s_c=1$ for every $c\in\pi_k$. Then
\[
S^{\#}_y(1,A)=S_y(A)\qquad\text{and}\qquad\sum_{X\in D_k}S^{\#}_y(X,A)=S_y(A).
\]
\item[(c)] \emph{Analyticity.} The map $(s,A)\mapsto S^{\#}_y(s,A)$ is continuous on $\mathrm{Pd}_k\times\tilde{G}_k$ and analytic on the interior of $\mathrm{Pd}_k\times\tilde{G}_k$; for every fixed $s\in\mathrm{Pd}_k$ the map $A\mapsto S^{\#}_y(s,A)$ is analytic on $\tilde{G}_k$; the map $A\mapsto S^{\#}_y(X,A)$ is analytic on $\tilde{G}_k$. For $y\in\mathbb{R}$ all these operators are continuous functions of $y$.
\item[(d)] \emph{Localisation.} Let $\boldsymbol{\omega}\in\boldsymbol{\Omega}$. The kernels of $S_{y,\boldsymbol{\omega}}(A)$ and of $S'_{y,\boldsymbol{\omega}}(A)$ vanish at $(x,x')$ unless both $x$ and $x'$ lie in $X_{\boldsymbol{\omega}}$. If $A,\tilde{A}\in\tilde{G}_k$ coincide on all bonds with both end points in $X_{\boldsymbol{\omega}}$, then $S^{\#}_{y,\boldsymbol{\omega}}(A)=S^{\#}_{y,\boldsymbol{\omega}}(\tilde{A})$. Let $x$ and $x'$ be sites with $S_{y,\boldsymbol{\omega}}(A)(x,x')\ne0$ or $S'_{y,\boldsymbol{\omega}}(A)(x,x')\ne0$, and let $K_x=\left\{\tilde{x}:\ |\tilde{x}-x|\le 2L\right\}$ and $K_{x'}$ be the sets of Section~\ref{sec:lattice-geometry}. Then
\[
c\subset X_{\boldsymbol{\omega}}\qquad\text{for every }c\in\pi_k\text{ with }c\cap K_x\ne\emptyset\text{ or }c\cap K_{x'}\ne\emptyset.
\]
\item[(e)] \emph{Gauge covariance.} Let $\boldsymbol{\lambda}$ be a real scalar on $T_k$. For every walk $\boldsymbol{\omega}$
\[
S^{\#}_{y,\boldsymbol{\omega}}(A-d\boldsymbol{\lambda})=\exp\left(i\,e_k\cdot\boldsymbol{\lambda}\right)\cdot S^{\#}_{y,\boldsymbol{\omega}}(A)\cdot\exp\left(-i\,e_k\cdot\boldsymbol{\lambda}\right),
\]
and consequently
\begin{align*}
S^{\#}_y(s,A-d\boldsymbol{\lambda})&=\exp\left(i\,e_k\cdot\boldsymbol{\lambda}\right)\cdot S^{\#}_y(s,A)\cdot\exp\left(-i\,e_k\cdot\boldsymbol{\lambda}\right),\\
S^{\#}_y(X,A-d\boldsymbol{\lambda})&=\exp\left(i\,e_k\cdot\boldsymbol{\lambda}\right)\cdot S^{\#}_y(X,A)\cdot\exp\left(-i\,e_k\cdot\boldsymbol{\lambda}\right).
\end{align*}
\item[(f)] \emph{Lattice symmetries.} Let $g=(S,a')\in\mathbb{G}_k$. For a walk $\boldsymbol{\omega}=(\square_0,\dots,\square_n)$ put $g\boldsymbol{\omega}:=(g\square_0,\dots,g\square_n)$; for $s\in\mathrm{Pd}_k$ put $(gs)_c:=s_{g^{-1}c}$ $(c\in\pi_k)$; let $y':=y$ if $S\in\mathrm{Sp}_3$ or if $y$ is neutral, and $y':=-y$ otherwise. Then
\[
X_{g\boldsymbol{\omega}}=gX_{\boldsymbol{\omega}},\qquad(gs)_{g\boldsymbol{\omega}}=s_{\boldsymbol{\omega}},\qquad S^{\#}_{y',g\boldsymbol{\omega}}(gA)\cdot g=g\cdot S^{\#}_{y,\boldsymbol{\omega}}(A),
\]
and consequently
\[
S^{\#}_{y'}(gs,gA)\cdot g=g\cdot S^{\#}_y(s,A),\qquad S^{\#}_{y'}(gX,gA)\cdot g=g\cdot S^{\#}_y(X,A).
\]
\item[(g)] \emph{Transposition.} Put $-\infty:=\infty$. For every walk $\boldsymbol{\omega}$
\[
S_{y,\boldsymbol{\omega}}(A)^T=\mathsf{C}^{-1}\cdot S'_{-y,\boldsymbol{\omega}}(-A)\cdot\mathsf{C},\qquad S'_{y,\boldsymbol{\omega}}(A)^T=\mathsf{C}^{-1}\cdot S_{-y,\boldsymbol{\omega}}(-A)\cdot\mathsf{C}.
\]
Consequently, if $y$ is neutral,
\[
S^{\mathrm{sy}}_y(s,A)^T=\mathsf{C}^{-1}\cdot S^{\mathrm{sy}}_y(s,-A)\cdot\mathsf{C},\qquad S^{\mathrm{sy}}_y(X,A)^T=\mathsf{C}^{-1}\cdot S^{\mathrm{sy}}_y(X,-A)\cdot\mathsf{C}.
\]
\item[(h)] \emph{Complex conjugation.} For every walk $\boldsymbol{\omega}$
\[
\overline{S^{\#}_{y,\boldsymbol{\omega}}(\overline{A})}=\varkappa^{-1}\cdot S^{\#}_{y,\boldsymbol{\omega}}(-A)\cdot\varkappa,
\]
and consequently
\[
\overline{S^{\#}_y(\overline{s},\overline{A})}=\varkappa^{-1}\cdot S^{\#}_y(s,-A)\cdot\varkappa,\qquad\overline{S^{\#}_y(X,\overline{A})}=\varkappa^{-1}\cdot S^{\#}_y(X,-A)\cdot\varkappa.
\]
\end{enumerate}
\end{proposition}

\begin{proof}
Throughout, $\boldsymbol{\omega}=(\square_0,\dots,\square_n)$ is a walk, so that $n=|\boldsymbol{\omega}|$; $\mathsf{c}_i$ is the centre of $\square_i$, and $h_i$, $\mathrm{Sl}_i$, $\mathrm{Cm}_i$ are the operators of Definition~\ref{def:walks}, with $\mathrm{Cm}_i=\mathrm{Cm}_{\square_i,y}(A)=-\left[F_y(A),h_{\square_i}\right]$ (Section~\ref{sec:parametrix-local-inverses}). Thus
\begin{align*}
S_{y,\boldsymbol{\omega}}(A)&=\left(h_0\cdot\mathrm{Sl}_0\cdot h_0\right)\cdot\left(\mathrm{Cm}_1\cdot\mathrm{Sl}_1\cdot h_1\right)\cdots\left(\mathrm{Cm}_n\cdot\mathrm{Sl}_n\cdot h_n\right),\\
S'_{y,\boldsymbol{\omega}}(A)&=\left(h_n\cdot\mathrm{Sl}_n\right)\cdot\left(-\mathrm{Cm}_n\cdot h_{n-1}\cdot\mathrm{Sl}_{n-1}\right)\cdots\left(-\mathrm{Cm}_1\cdot h_0\cdot\mathrm{Sl}_0\right)\cdot h_0.
\end{align*}

(a) Every parameter satisfies $|s_c|\le M^{\alpha_0}$ (Definition~\ref{def:parameter-cubes-polydisks}), the product $s_{\boldsymbol{\omega}}$ has $|X_{\boldsymbol{\omega}}|_M$ factors, and $|X_{\boldsymbol{\omega}}|_M\le 27(n+1)$ (Section~\ref{sec:expansion-walks}). Since $M\ge1$,
\[
|s_{\boldsymbol{\omega}}|\le M^{\alpha_0|X_{\boldsymbol{\omega}}|_M}\le M^{27\alpha_0(n+1)}=w_s^{|\boldsymbol{\omega}|+1}.
\]
Condition~\eqref{eq:renormalisation-group-19} gives $M\ge\left(108\pi\cdot C_V\cdot a\right)^2$, that is, $M^{\frac{1}{2}}\ge 108\pi\cdot C_V\cdot a$; and $27\alpha_0\le\frac{1}{2}$ gives $w_s\le M^{\frac{1}{2}}$. Hence
\[
M=M^{\frac{1}{2}}\cdot M^{\frac{1}{2}}\ge 108\pi\cdot C_V\cdot a\cdot M^{\frac{1}{2}}\ge 108\pi\cdot C_V\cdot a\cdot w_s,
\]
which says that $q_{w_s}\le\frac{1}{2}$ for the number $q_w=\frac{54\pi\cdot C_V\cdot a\cdot w}{M}$ of Lemma~\ref{lem:walks-given-length}. Condition~\eqref{eq:renormalisation-group-19} also gives $M\ge 108\pi\cdot C_V\cdot a\cdot\exp(27\kappa')=108\pi\cdot C_V\cdot a\cdot w_X$, that is, $q_{w_X}\le\frac{1}{2}$. As $w_s\ge1$ and $w_X\ge1$, Lemma~\ref{lem:walks-given-length} applies with $w=w_s$ and with $w=w_X$. Since $S^{\mathrm{sy}}_{y,\boldsymbol{\omega}}(A)$ is the mean of $S_{y,\boldsymbol{\omega}}(A)$ and $S'_{y,\boldsymbol{\omega}}(A)$, the two estimates of Lemma~\ref{lem:walks-given-length} hold, by the triangle inequality, also for $S^{\mathrm{sy}}_{y,\boldsymbol{\omega}}(A)$ in place of $S_{y,\boldsymbol{\omega}}(A)$ or $S'_{y,\boldsymbol{\omega}}(A)$.

Let $z$ be a unit site and $n\ge0$. The bound on $|s_{\boldsymbol{\omega}}|$ and Lemma~\ref{lem:walks-given-length} with $w=w_s$ give, for this $z$ and this $n$,
\begin{align*}
\sum_{z'}\exp\left(\gamma\cdot d(z,z')\right)\sum_{\boldsymbol{\omega}\in\boldsymbol{\Omega}_n}|s_{\boldsymbol{\omega}}|\cdot\left\|\mathbf{1}_{\Delta_z}S^{\#}_{y,\boldsymbol{\omega}}(A)\,\mathbf{1}_{\Delta_{z'}}\right\|&\le 16a\cdot w_s\cdot w_k(y)\cdot q_{w_s}^n,\\
\sum_{z'}\exp\left(\gamma\cdot d(z,z')\right)\sum_{\boldsymbol{\omega}\in\boldsymbol{\Omega}_n}|s_{\boldsymbol{\omega}}|\cdot\left\|\mathbf{1}_{\Delta_z}\delta_{\alpha,A}S^{\#}_{y,\boldsymbol{\omega}}(A)\,\mathbf{1}_{\Delta_{z'}}\right\|&\le 152a\cdot w_s\cdot q_{w_s}^n.
\end{align*}
The right-hand sides depend neither on $z$ nor on $(s,A)$, and $\sum_{n\ge N}q_{w_s}^n\le 2q_{w_s}^N$ for every integer $N\ge0$, because $q_{w_s}\le\frac{1}{2}$.

In the definition of $N_{\gamma}$ in Section~\ref{sec:norms-exponential-weights}, the supremum and the sum extend over the unit sites of $T_k$; these are finitely many, and $\nu$ denotes their number. A supremum of non-negative numbers does not exceed their sum; hence every linear operator $O$ on the spinors of $T_k$, with values in the spinors or in the functions of short pairs, satisfies
\[
N_{\gamma}(O)\le\sum_z\sum_{z'}\exp\left(\gamma\cdot d(z,z')\right)\cdot\left\|\mathbf{1}_{\Delta_z}O\,\mathbf{1}_{\Delta_{z'}}\right\|.
\]
Let $N\ge0$. Take $O=S^{\#}_{y,\boldsymbol{\omega}}(A)$ and $O=\delta_{\alpha,A}S^{\#}_{y,\boldsymbol{\omega}}(A)$ in this inequality, multiply by $|s_{\boldsymbol{\omega}}|$, and sum over the walks $\boldsymbol{\omega}$ of length at least $N$. All terms being non-negative, the order of the summations may be changed so that the sum over $z$ stands outside; then, for each of the $\nu$ unit sites $z$ and each $n\ge N$, the sum over $z'$ and over $\boldsymbol{\Omega}_n$ is estimated by the bounds for fixed $z$ and $n$. The result is
\begin{align*}
\sum_{n\ge N}\ \sum_{\boldsymbol{\omega}\in\boldsymbol{\Omega}_n}|s_{\boldsymbol{\omega}}|\cdot N_{\gamma}\left(S^{\#}_{y,\boldsymbol{\omega}}(A)\right)&\le\nu\cdot 32a\cdot w_s\cdot w_k(y)\cdot q_{w_s}^N,\\
\sum_{n\ge N}\ \sum_{\boldsymbol{\omega}\in\boldsymbol{\Omega}_n}|s_{\boldsymbol{\omega}}|\cdot N_{\gamma}\left(\delta_{\alpha,A}S^{\#}_{y,\boldsymbol{\omega}}(A)\right)&\le\nu\cdot 304a\cdot w_s\cdot q_{w_s}^N.
\end{align*}
For $N=0$ the left-hand sides are finite, and as $N\to\infty$ they tend to $0$, uniformly in $(s,A)\in\mathrm{Pd}_k\times\tilde{G}_k$. Hence the series $\sum_{\boldsymbol{\omega}}s_{\boldsymbol{\omega}}\cdot S^{\#}_{y,\boldsymbol{\omega}}(A)$ and $\sum_{\boldsymbol{\omega}}s_{\boldsymbol{\omega}}\cdot\delta_{\alpha,A}S^{\#}_{y,\boldsymbol{\omega}}(A)$, whose sums are $S^{\#}_y(s,A)$ and $\delta_{\alpha,A}S^{\#}_y(s,A)$, converge absolutely in $N_{\gamma}$, uniformly on $\mathrm{Pd}_k\times\tilde{G}_k$. Moreover, for every unit site $z$, the triangle inequality and the bounds for fixed $z$ and $n$, summed over all $n\ge0$, estimate the sum over $z'$ in the definition of $N_{\gamma}$ (Section~\ref{sec:norms-exponential-weights}) for $O=S^{\#}_y(s,A)$ and for $O=\delta_{\alpha,A}S^{\#}_y(s,A)$; the supremum over $z$ then gives
\begin{align*}
N_{\gamma}\left(S^{\#}_y(s,A)\right)&\le\sum_{n\ge0}16a\cdot w_s\cdot w_k(y)\cdot q_{w_s}^n\le 32a\cdot w_s\cdot w_k(y),\\
N_{\gamma}\left(\delta_{\alpha,A}S^{\#}_y(s,A)\right)&\le\sum_{n\ge0}152a\cdot w_s\cdot q_{w_s}^n\le 304a\cdot w_s.
\end{align*}
As $1\le w_s^{n+1}$, all these estimates hold with $1$ in place of $|s_{\boldsymbol{\omega}}|$. In particular the sum $\sum_{\boldsymbol{\omega}\in\boldsymbol{\Omega}}N_{\gamma}\left(S^{\#}_{y,\boldsymbol{\omega}}(A)\right)$ is finite, and its part over the walks of length at least $N$ is at most $\nu\cdot 32a\cdot w_s\cdot w_k(y)\cdot q_{w_s}^N$ for every $A\in\tilde{G}_k$. The terms of the series that define $S^{\#}_y(X,A)$ are terms of the series $\sum_{\boldsymbol{\omega}\in\boldsymbol{\Omega}}S^{\#}_{y,\boldsymbol{\omega}}(A)$; so these series converge absolutely in $N_{\gamma}$, uniformly on $\tilde{G}_k$.

Let $\boldsymbol{\omega}$ be a walk with $X_{\boldsymbol{\omega}}=X$. Then $d_M(X)=d_M(X_{\boldsymbol{\omega}})\le 27(|\boldsymbol{\omega}|+1)$ (Section~\ref{sec:expansion-walks}), and therefore
\[
1\le\exp\left(-\kappa'\cdot d_M(X)\right)\cdot\exp\left(27\kappa'(|\boldsymbol{\omega}|+1)\right)=\exp\left(-\kappa'\cdot d_M(X)\right)\cdot w_X^{|\boldsymbol{\omega}|+1}.
\]
Let $z$ and $z'$ be unit sites. By the definition of $S^{\#}_y(X,A)$, the triangle inequality and this bound, after the sum over the walks with $X_{\boldsymbol{\omega}}=X$ has been extended to all walks,
\begin{align*}
&\exp\left(\gamma\cdot d(z,z')+\kappa'\cdot d_M(X)\right)\cdot\left\|\mathbf{1}_{\Delta_z}S^{\#}_y(X,A)\,\mathbf{1}_{\Delta_{z'}}\right\|\\
&\qquad\le\sum_{n\ge0}\sum_{\boldsymbol{\omega}\in\boldsymbol{\Omega}_n}w_X^{n+1}\cdot\exp\left(\gamma\cdot d(z,z')\right)\cdot\left\|\mathbf{1}_{\Delta_z}S^{\#}_{y,\boldsymbol{\omega}}(A)\,\mathbf{1}_{\Delta_{z'}}\right\|\\
&\qquad\le\sum_{n\ge0}16a\cdot w_X\cdot w_k(y)\cdot q_{w_X}^n\le 32a\cdot w_X\cdot w_k(y);
\end{align*}
the second inequality is Lemma~\ref{lem:walks-given-length} with $w=w_X$, whose first estimate contains the given term $z'$ in its sum over the unit sites. This is the third bound of (a).

(b) By Theorem~\ref{thm:expansion-walks}(a),
\[
S_y(A)=\sum_{\boldsymbol{\omega}\in\boldsymbol{\Omega}}S_{y,\boldsymbol{\omega}}(A)=\sum_{\boldsymbol{\omega}\in\boldsymbol{\Omega}}S'_{y,\boldsymbol{\omega}}(A),
\]
both series converging absolutely. For the family $s=1$ one has $s_{\boldsymbol{\omega}}=1$ for every walk; hence $S_y(1,A)=S'_y(1,A)=S_y(A)$, and therefore $S^{\mathrm{sy}}_y(1,A)=S_y(A)$. Every walk $\boldsymbol{\omega}$ determines exactly one set $X_{\boldsymbol{\omega}}$, and this set is a polymer, that is, an element of $D_k$ (Section~\ref{sec:expansion-walks}). Grouping the terms of the absolutely convergent series according to $X_{\boldsymbol{\omega}}$ gives
\[
\sum_{X\in D_k}S^{\#}_y(X,A)=\sum_{\boldsymbol{\omega}\in\boldsymbol{\Omega}}S^{\#}_{y,\boldsymbol{\omega}}(A)=S_y(A).
\]

(c) The factors $h_i$ do not depend on $(s,A,y)$. For $A\in\tilde{G}_k$ the restriction of $A$ to the bonds of the lift of $\square_i^{\sim 1}$ depends linearly on $A$ and lies in $G^W(2e_k^{\varepsilon})$ (Section~\ref{sec:operators-cube-operators}); so, by Proposition~\ref{prop:properties-local-inverse}(a), the factor $\mathrm{Sl}_i=\mathrm{Sl}_{\square_i,y}(A)$ is analytic in $A\in\tilde{G}_k$ and continuous in $y\in\mathbb{R}$. The entries of $\mathrm{Cm}_i=-\left[F_y(A),h_{\square_i}\right]$ are entire functions of $A$. For $k\ge1$ the mass and the matrices $\boldsymbol{\alpha}_y$, $\boldsymbol{\beta}_y$ commute with $h_{\square_i}$, so that
\[
\mathrm{Cm}_i=-\left[D_A,h_{\square_i}\right]-\boldsymbol{\alpha}_y\cdot\left[P_k(A),h_{\square_i}\right]-\boldsymbol{\beta}_y\cdot\left[P_{k+1}(A),h_{\square_i}\right]
\]
contains $y$ only through $\boldsymbol{\alpha}_y$ and $\boldsymbol{\beta}_y$, which are continuous functions of $y\in\mathbb{R}$ (Section~\ref{sec:operators-2}); for $k=0$ the factor $\mathrm{Cm}_i$ does not contain $y$ at all, since $i\gamma_3y$ commutes with $h_{\square_i}$. Consequently each term $s_{\boldsymbol{\omega}}\cdot S^{\#}_{y,\boldsymbol{\omega}}(A)$ of the series of Definition~\ref{def:weakened-sums-walks} is a monomial in $s$ multiplied by an operator that is analytic in $A\in\tilde{G}_k$ and continuous in $y\in\mathbb{R}$. The estimates in the proof of (a) contain $y$ only through the factor $w_k(y)\le1$ of~\eqref{eq:operators-fermion-2}; hence these series converge uniformly on $\mathrm{Pd}_k\times\tilde{G}_k$, and, for $y\in\mathbb{R}$, uniformly on $\mathrm{Pd}_k\times\tilde{G}_k\times\mathbb{R}$. The set $\tilde{G}_k$ is open by~\eqref{eq:operators-fermion-G2}. A uniform limit of continuous functions is continuous, and a uniform limit of analytic functions of several complex variables is analytic (theorem of Weierstrass). This proves (c).

(d) The leftmost factor of the product $S_{y,\boldsymbol{\omega}}(A)$ is $h_0$ and its rightmost factor is $h_n$; for $S'_{y,\boldsymbol{\omega}}(A)$ these factors are $h_n$ and $h_0$. Denote the leftmost factor by $h_i$ and the rightmost one by $h_j$. The kernel of the product at $(x,x')$ vanishes unless $h_i(x)\ne0$ and $h_j(x')\ne0$. By~\eqref{eq:operators-fermion-7} the support of $h_{\square}$ consists of sites $x$ with $|x-\mathsf{c}|<\frac{3M}{4}$ on the lift, and these lie in $\square^{\sim 1}$. Hence a non-vanishing entry at $(x,x')$ has $x\in\square_i^{\sim 1}\subset X_{\boldsymbol{\omega}}$ and $x'\in\square_j^{\sim 1}\subset X_{\boldsymbol{\omega}}$.

The factors $h_i$ do not depend on $A$. The factor $\mathrm{Sl}_i$ depends only on the values of $A$ on the bonds of the lift of $\square_i^{\sim 1}$ (Section~\ref{sec:folding-torus-onto}, Definition~\ref{def:folded-fields-local}, Section~\ref{sec:operators-cube-operators}), and so does $\mathrm{Cm}_i$ (Lemma~\ref{lem:parametrix-identities}(c)). These bonds have both end points in $\square_i^{\sim 1}\subset X_{\boldsymbol{\omega}}$. Hence $S_{y,\boldsymbol{\omega}}(A)$ and $S'_{y,\boldsymbol{\omega}}(A)$, and with them their mean $S^{\mathrm{sy}}_{y,\boldsymbol{\omega}}(A)$, do not change when $A$ is replaced by an $\tilde{A}\in\tilde{G}_k$ that coincides with $A$ on the bonds with both end points in $X_{\boldsymbol{\omega}}$.

Finally, let $x$ and $x'$ be sites at which the kernel of $S_{y,\boldsymbol{\omega}}(A)$ or that of $S'_{y,\boldsymbol{\omega}}(A)$ is non-zero, and let $h_i$ and $h_j$ be the leftmost and the rightmost factor of that product. By the first paragraph, $h_i(x)\ne0$ and $h_j(x')\ne0$. From $h_i(x)\ne0$ and~\eqref{eq:operators-fermion-7} one gets $|x-\mathsf{c}_i|<\frac{3M}{4}$, so that every $\tilde{x}\in K_x$ satisfies $|\tilde{x}-\mathsf{c}_i|<\frac{3M}{4}+2L$. Since $M>4L$, one has $\frac{3M}{4}+2L<\frac{3M}{2}$; hence $K_x\subset\square_i^{\sim 1}$. The set $\square_i^{\sim 1}$ is a union of $M$-cubes, and distinct $M$-cubes are disjoint. Consequently a cube $c\in\pi_k$ with $c\cap K_x\ne\emptyset$ has a site in common with one of the $M$-cubes of $\square_i^{\sim 1}$, coincides with it, and satisfies $c\subset\square_i^{\sim 1}\subset X_{\boldsymbol{\omega}}$. The same argument, with $x'$, $h_j$, $\mathsf{c}_j$ and $\square_j$ in place of $x$, $h_i$, $\mathsf{c}_i$ and $\square_i$, gives $c\subset X_{\boldsymbol{\omega}}$ for every $c\in\pi_k$ with $c\cap K_{x'}\ne\emptyset$.

(e) Let $u:=\exp\left(i\,e_k\cdot\boldsymbol{\lambda}\right)$, also regarded as the operator of multiplication by this function. By~\eqref{eq:operators-fermion-G2}, $A-d\boldsymbol{\lambda}\in\tilde{G}_k$. Lemma~\ref{lem:symmetries-elementary-operators}(a), applied on $T_k$ to the closed real connection $d\boldsymbol{\lambda}$ and to the function $u$, gives $F_y(A-d\boldsymbol{\lambda})=u\cdot F_y(A)\cdot u^{-1}$. The multiplication operators $h_{\square}$ and $u$ commute; hence
\[
\mathrm{Cm}_{\square,y}(A-d\boldsymbol{\lambda})=-\left[u\cdot F_y(A)\cdot u^{-1},h_{\square}\right]=u\cdot\mathrm{Cm}_{\square,y}(A)\cdot u^{-1}.
\]
Proposition~\ref{prop:properties-local-inverse}(d), applied to the restriction of $\boldsymbol{\lambda}$ to the lift of $\square^{\sim 1}$ (Section~\ref{sec:operators-cube-operators}), gives $\mathrm{Sl}_{\square,y}(A-d\boldsymbol{\lambda})=u\cdot\mathrm{Sl}_{\square,y}(A)\cdot u^{-1}$. Thus every factor $O$ of the two products displayed at the beginning of the proof satisfies $O(A-d\boldsymbol{\lambda})=u\cdot O(A)\cdot u^{-1}$, and in each product the inner factors $u^{-1}\cdot u$ cancel. This gives the relation for $S_{y,\boldsymbol{\omega}}$ and for $S'_{y,\boldsymbol{\omega}}$, hence for their mean $S^{\mathrm{sy}}_{y,\boldsymbol{\omega}}$. The coefficients $s_{\boldsymbol{\omega}}$ and the sets $X_{\boldsymbol{\omega}}$ do not depend on $A$; so
\begin{align*}
S^{\#}_y(s,A-d\boldsymbol{\lambda})&=\sum_{\boldsymbol{\omega}\in\boldsymbol{\Omega}}s_{\boldsymbol{\omega}}\cdot u\cdot S^{\#}_{y,\boldsymbol{\omega}}(A)\cdot u^{-1}=u\cdot S^{\#}_y(s,A)\cdot u^{-1},\\
S^{\#}_y(X,A-d\boldsymbol{\lambda})&=\sum_{\boldsymbol{\omega}\in\boldsymbol{\Omega}:\ X_{\boldsymbol{\omega}}=X}u\cdot S^{\#}_{y,\boldsymbol{\omega}}(A)\cdot u^{-1}=u\cdot S^{\#}_y(X,A)\cdot u^{-1}.
\end{align*}

(f) The morphism $x\mapsto r(S)x+a'$ of $T_k$ induced by $g$ (Section~\ref{sec:tori-nested-blocks}) has $a'\in M\mathbb{Z}^3$ and a diagonal linear part with entries $\pm1$. Therefore it preserves the distance $d$, permutes the $M$-cubes, maps the centre of every $M$-cube $\square$ to the centre of $g\square$, and maps $\square^{\sim 1}$ onto $(g\square)^{\sim 1}$. Consequently $g\boldsymbol{\omega}$ is a walk, the map $\boldsymbol{\omega}\mapsto g\boldsymbol{\omega}$ is a bijection of $\boldsymbol{\Omega}_n$ onto itself for every $n$ (with inverse $\boldsymbol{\omega}\mapsto g^{-1}\boldsymbol{\omega}$), and
\[
X_{g\boldsymbol{\omega}}=\bigcup_{i=0}^n(g\square_i)^{\sim 1}=g\bigcup_{i=0}^n\square_i^{\sim 1}=gX_{\boldsymbol{\omega}}.
\]
The $M$-cubes contained in $gX_{\boldsymbol{\omega}}$ are the cubes $gc$ with $c\in\pi_k$, $c\subset X_{\boldsymbol{\omega}}$; hence
\[
(gs)_{g\boldsymbol{\omega}}=\prod_{c\in\pi_k,\ c\subset X_{\boldsymbol{\omega}}}(gs)_{gc}=\prod_{c\in\pi_k,\ c\subset X_{\boldsymbol{\omega}}}s_c=s_{\boldsymbol{\omega}}.
\]
Moreover $gs\in\mathrm{Pd}_k$, $gX\in D_k$, and $gA\in\tilde{G}_k$ by~\eqref{eq:operators-fermion-G2}. By~\eqref{eq:operators-fermion-7}, $h_{g\square}(gx)=h_{\square}(x)$; since the matrix $S$ commutes with scalar functions, this means $g\cdot h_{\square}\cdot g^{-1}=h_{g\square}$ for the multiplication operators on spinors. Lemma~\ref{lem:symmetries-elementary-operators}(b), applied on $T_k$ to the morphism induced by $g$, gives $F_{y'}(gA)\cdot g=g\cdot F_y(A)$; hence
\begin{align*}
\mathrm{Cm}_{g\square,y'}(gA)\cdot g&=-F_{y'}(gA)\cdot h_{g\square}\cdot g+h_{g\square}\cdot F_{y'}(gA)\cdot g\\
&=-g\cdot F_y(A)\cdot h_{\square}+g\cdot h_{\square}\cdot F_y(A)=g\cdot\mathrm{Cm}_{\square,y}(A).
\end{align*}
Proposition~\ref{prop:properties-local-inverse}(e) gives $\mathrm{Sl}_{g\square,y'}(gA)\cdot g=g\cdot\mathrm{Sl}_{\square,y}(A)$. The walk $g\boldsymbol{\omega}$ consists of the cubes $g\square_i$; by these three relations the operator $g$ can therefore be moved, factor by factor, from the right end to the left end of each of the products $S_{y',g\boldsymbol{\omega}}(gA)\cdot g$ and $S'_{y',g\boldsymbol{\omega}}(gA)\cdot g$ (Definition~\ref{def:walks}). This gives
\[
S_{y',g\boldsymbol{\omega}}(gA)\cdot g=g\cdot S_{y,\boldsymbol{\omega}}(A),\qquad S'_{y',g\boldsymbol{\omega}}(gA)\cdot g=g\cdot S'_{y,\boldsymbol{\omega}}(A),
\]
and half the sum of these two identities is $S^{\mathrm{sy}}_{y',g\boldsymbol{\omega}}(gA)\cdot g=g\cdot S^{\mathrm{sy}}_{y,\boldsymbol{\omega}}(A)$, by the definition of $S^{\mathrm{sy}}_{y,\boldsymbol{\omega}}$ in Definition~\ref{def:weakened-sums-walks}. Finally, as $\boldsymbol{\omega}\mapsto g\boldsymbol{\omega}$ is a bijection of $\boldsymbol{\Omega}$ and the series converge absolutely by (a),
\begin{align*}
S^{\#}_{y'}(gs,gA)\cdot g&=\sum_{\boldsymbol{\omega}\in\boldsymbol{\Omega}}(gs)_{g\boldsymbol{\omega}}\cdot S^{\#}_{y',g\boldsymbol{\omega}}(gA)\cdot g\\
&=\sum_{\boldsymbol{\omega}\in\boldsymbol{\Omega}}s_{\boldsymbol{\omega}}\cdot g\cdot S^{\#}_{y,\boldsymbol{\omega}}(A)=g\cdot S^{\#}_y(s,A),
\end{align*}
and, since $X_{g\boldsymbol{\omega}}=gX$ holds if and only if $X_{\boldsymbol{\omega}}=X$,
\begin{align*}
S^{\#}_{y'}(gX,gA)\cdot g&=\sum_{\boldsymbol{\omega}\in\boldsymbol{\Omega}:\ X_{\boldsymbol{\omega}}=X}S^{\#}_{y',g\boldsymbol{\omega}}(gA)\cdot g\\
&=\sum_{\boldsymbol{\omega}\in\boldsymbol{\Omega}:\ X_{\boldsymbol{\omega}}=X}g\cdot S^{\#}_{y,\boldsymbol{\omega}}(A)=g\cdot S^{\#}_y(X,A).
\end{align*}

(g) By~\eqref{eq:operators-fermion-G2}, $-A\in\tilde{G}_k$. For an operator $O$ on the spinors of $T_k$ and a real scalar function $h$, whose multiplication operator is its own transpose, one has $[O,h]^T=h\cdot O^T-O^T\cdot h=-\left[O^T,h\right]$. Lemma~\ref{lem:symmetries-elementary-operators}(c) gives $F_y(A)^T=\mathsf{C}^{-1}\cdot F_{-y}(-A)\cdot\mathsf{C}$, and the constant matrix $\mathsf{C}$ commutes with $h_{\square}$; hence
\[
\mathrm{Cm}_{\square,y}(A)^T=\left[F_y(A)^T,h_{\square}\right]=\mathsf{C}^{-1}\cdot\left[F_{-y}(-A),h_{\square}\right]\cdot\mathsf{C}=-\mathsf{C}^{-1}\cdot\mathrm{Cm}_{\square,-y}(-A)\cdot\mathsf{C}.
\]
By Proposition~\ref{prop:properties-local-inverse}(f), $\mathrm{Sl}_{\square,y}(A)^T=\mathsf{C}^{-1}\cdot\mathrm{Sl}_{\square,-y}(-A)\cdot\mathsf{C}$. Transposition reverses the order of the factors of $S_{y,\boldsymbol{\omega}}(A)$:
\begin{align*}
S_{y,\boldsymbol{\omega}}(A)^T&=h_n\cdot\mathrm{Sl}_n^T\cdot\mathrm{Cm}_n^T\cdot h_{n-1}\cdot\mathrm{Sl}_{n-1}^T\cdot\mathrm{Cm}_{n-1}^T\cdots\\
&\qquad\cdots h_1\cdot\mathrm{Sl}_1^T\cdot\mathrm{Cm}_1^T\cdot h_0\cdot\mathrm{Sl}_0^T\cdot h_0.
\end{align*}
Put $\mathrm{Sl}'_i:=\mathrm{Sl}_{\square_i,-y}(-A)$ and $\mathrm{Cm}'_i:=\mathrm{Cm}_{\square_i,-y}(-A)$. Inserting the two relations and using that $\mathsf{C}$ commutes with the functions $h_i$, one obtains
\begin{align*}
S_{y,\boldsymbol{\omega}}(A)^T&=\mathsf{C}^{-1}\cdot\left(h_n\cdot\mathrm{Sl}'_n\right)\cdot\left(-\mathrm{Cm}'_n\cdot h_{n-1}\cdot\mathrm{Sl}'_{n-1}\right)\cdots\left(-\mathrm{Cm}'_1\cdot h_0\cdot\mathrm{Sl}'_0\right)\cdot h_0\cdot\mathsf{C}\\
&=\mathsf{C}^{-1}\cdot S'_{-y,\boldsymbol{\omega}}(-A)\cdot\mathsf{C},
\end{align*}
which is the first relation. Applied with $(-y,-A)$ in place of $(y,A)$, it reads $S_{-y,\boldsymbol{\omega}}(-A)^T=\mathsf{C}^{-1}\cdot S'_{y,\boldsymbol{\omega}}(A)\cdot\mathsf{C}$. Transposing this identity and using $\mathsf{C}^T=\mathsf{C}^{-1}=-\mathsf{C}$ (see~\eqref{eq:operators-fermion-1}), one finds $S_{-y,\boldsymbol{\omega}}(-A)=\mathsf{C}^{-1}\cdot S'_{y,\boldsymbol{\omega}}(A)^T\cdot\mathsf{C}$, that is,
\[
S'_{y,\boldsymbol{\omega}}(A)^T=\mathsf{C}\cdot S_{-y,\boldsymbol{\omega}}(-A)\cdot\mathsf{C}^{-1}=\mathsf{C}^{-1}\cdot S_{-y,\boldsymbol{\omega}}(-A)\cdot\mathsf{C},
\]
the second relation. Let now $y$ be neutral; then $-y=y$, by the convention $-\infty=\infty$ of (g), which is that of Lemma~\ref{lem:symmetries-elementary-operators}. Adding the two relations and dividing by $2$ gives $S^{\mathrm{sy}}_{y,\boldsymbol{\omega}}(A)^T=\mathsf{C}^{-1}\cdot S^{\mathrm{sy}}_{y,\boldsymbol{\omega}}(-A)\cdot\mathsf{C}$. The coefficients $s_{\boldsymbol{\omega}}$ and the sets $X_{\boldsymbol{\omega}}$ are the same for the two families, and transposition is linear and continuous, so that it may be applied term by term:
\begin{align*}
S^{\mathrm{sy}}_y(s,A)^T&=\sum_{\boldsymbol{\omega}\in\boldsymbol{\Omega}}s_{\boldsymbol{\omega}}\cdot\mathsf{C}^{-1}\cdot S^{\mathrm{sy}}_{y,\boldsymbol{\omega}}(-A)\cdot\mathsf{C}=\mathsf{C}^{-1}\cdot S^{\mathrm{sy}}_y(s,-A)\cdot\mathsf{C},\\
S^{\mathrm{sy}}_y(X,A)^T&=\sum_{\boldsymbol{\omega}\in\boldsymbol{\Omega}:\ X_{\boldsymbol{\omega}}=X}\mathsf{C}^{-1}\cdot S^{\mathrm{sy}}_{y,\boldsymbol{\omega}}(-A)\cdot\mathsf{C}=\mathsf{C}^{-1}\cdot S^{\mathrm{sy}}_y(X,-A)\cdot\mathsf{C}.
\end{align*}

(h) By~\eqref{eq:operators-fermion-G2}, $\overline{A}\in\tilde{G}_k$ and $-A\in\tilde{G}_k$; also $\overline{s}\in\mathrm{Pd}_k$. Lemma~\ref{lem:symmetries-elementary-operators}(d) gives $\overline{F_y(\overline{A})}=\varkappa^{-1}\cdot F_y(-A)\cdot\varkappa$. The functions $h_{\square}$ are real and commute with the constant matrix $\varkappa$; hence
\[
\overline{\mathrm{Cm}_{\square,y}(\overline{A})}=-\left[\overline{F_y(\overline{A})},h_{\square}\right]=\varkappa^{-1}\cdot\mathrm{Cm}_{\square,y}(-A)\cdot\varkappa,
\]
and Proposition~\ref{prop:properties-local-inverse}(f) gives $\overline{\mathrm{Sl}_{\square,y}(\overline{A})}=\varkappa^{-1}\cdot\mathrm{Sl}_{\square,y}(-A)\cdot\varkappa$. Complex conjugation of the products that define $S_{y,\boldsymbol{\omega}}$ and $S'_{y,\boldsymbol{\omega}}$, factor by factor, with cancellation of the inner factors $\varkappa\cdot\varkappa^{-1}$, gives $\overline{S^{\#}_{y,\boldsymbol{\omega}}(\overline{A})}=\varkappa^{-1}\cdot S^{\#}_{y,\boldsymbol{\omega}}(-A)\cdot\varkappa$ for $S^{\#}=S$ and $S^{\#}=S'$, and hence for their mean. Moreover $(\overline{s})_{\boldsymbol{\omega}}=\overline{s_{\boldsymbol{\omega}}}$, that is, $\overline{(\overline{s})_{\boldsymbol{\omega}}}=s_{\boldsymbol{\omega}}$. Therefore
\begin{align*}
\overline{S^{\#}_y(\overline{s},\overline{A})}&=\sum_{\boldsymbol{\omega}\in\boldsymbol{\Omega}}\overline{(\overline{s})_{\boldsymbol{\omega}}}\cdot\overline{S^{\#}_{y,\boldsymbol{\omega}}(\overline{A})}\\
&=\sum_{\boldsymbol{\omega}\in\boldsymbol{\Omega}}s_{\boldsymbol{\omega}}\cdot\varkappa^{-1}\cdot S^{\#}_{y,\boldsymbol{\omega}}(-A)\cdot\varkappa=\varkappa^{-1}\cdot S^{\#}_y(s,-A)\cdot\varkappa,\\
\overline{S^{\#}_y(X,\overline{A})}&=\sum_{\boldsymbol{\omega}\in\boldsymbol{\Omega}:\ X_{\boldsymbol{\omega}}=X}\overline{S^{\#}_{y,\boldsymbol{\omega}}(\overline{A})}\\
&=\sum_{\boldsymbol{\omega}\in\boldsymbol{\Omega}:\ X_{\boldsymbol{\omega}}=X}\varkappa^{-1}\cdot S^{\#}_{y,\boldsymbol{\omega}}(-A)\cdot\varkappa=\varkappa^{-1}\cdot S^{\#}_y(X,-A)\cdot\varkappa.
\end{align*}
\end{proof}

\subsection{Operators on the unit lattice for $k\ge1$}\label{sec:operators-unit-lattice}

Throughout this section $k\ge1$, and the hypotheses of Theorem~\ref{thm:expansion-walks} are satisfied. For $A\in\tilde{G}_k$ let
\[
\mathsf{Q}:=Q_k(A),\qquad\tilde{Q}:=Q_k^T(-A),\qquad P:=P(A).
\]
For a unit site $z$, $\mathbf{1}_z$ denotes the operator $\mathbf{1}_{\Delta_z}$ of Section~\ref{sec:norms-exponential-weights} on the spinors on the unit lattice, that is, the multiplication by the indicator function of the orbit of $z$; for $v\in\mathbb{C}^4$, $\mathbf{1}_zv$ denotes the spinor on the unit lattice that takes the value $v$ at $z$ and vanishes outside the orbit of $z$. With $S_k(A)=S_{\infty}(A)$, the operators $H_k(A)$, $M_k(A)$, $D_k(A)$, $\dot{M}_k(A)$, $B_{k,y}(A)$, $B_k(A)$ and $\delta Z_k(A)$ of Section~\ref{sec:exact-fermion-operators} are given by the following display, which also defines the operator $\mathfrak{H}_k(A)$:
\begin{gather*}
H_k(A):=b_k\cdot S_k(A)\cdot\tilde{Q},\qquad\mathfrak{H}_k(A):=b_k\cdot\mathsf{Q}\cdot S_k(A),\qquad M_k(A):=b_k^2\cdot\mathsf{Q}\cdot S_k(A)\cdot\tilde{Q},\\
D_k(A):=b_k-M_k(A),\qquad\dot{M}_k(A):=b_k^2\cdot\mathsf{Q}\cdot S_k(A)^2\cdot\tilde{Q},\\
B_{k,y}(A):=\left(b_k+i\gamma_3\cdot y\right)^{-1}\cdot(1-P)+\left(b_k+b_L+i\gamma_3\cdot y\right)^{-1}\cdot P\quad(y\in\mathbb{R}),\qquad B_k(A):=B_{k,0}(A),\\
\delta Z_k(A):=\det\left(D_k(A)+b_L\cdot P(A)\right).
\end{gather*}
The operator $H_k(A)$ maps the spinors on the unit lattice to the spinors on $T_k$, the operator $\mathfrak{H}_k(A)$ maps the spinors on $T_k$ to the spinors on the unit lattice, and $M_k(A)$, $D_k(A)$, $\dot{M}_k(A)$, $B_{k,y}(A)$ act on the finite-dimensional space of the spinors on the unit lattice, on which the determinant $\delta Z_k(A)$ is taken.

\subsubsection{Resolvent identities}\label{sec:resolvent-identities}

\begin{proposition}\label{prop:operators-fermion-1}
Let $A\in\tilde{G}_k$ and $y\in\mathbb{R}$.
\begin{enumerate}
\item[(a)] One has $B_{k,y}(A)=\left(b_k+b_L\cdot P+i\gamma_3\cdot y\right)^{-1}$ and
\[
F_y(A)=F_{\infty}(A)-b_k^2\cdot\tilde{Q}\cdot B_{k,y}(A)\cdot\mathsf{Q}.
\]
\item[(b)] The operator $D_k(A)+b_L\cdot P(A)+i\gamma_3\cdot y$ is bijective, and its inverse is
\[
\Gamma_{k,y}(A):=B_{k,y}(A)+b_k^2\cdot B_{k,y}(A)\cdot\mathsf{Q}\cdot S_y(A)\cdot\tilde{Q}\cdot B_{k,y}(A).
\]
\item[(c)] $S_y(A)=S_k(A)+H_k(A)\cdot\Gamma_{k,y}(A)\cdot\mathfrak{H}_k(A)$.
\item[(d)] The operator $D_k(A)+b_L\cdot P(A)$ is invertible, its inverse $\Gamma_k(A):=\left(D_k(A)+b_L\cdot P(A)\right)^{-1}$ is given by
\[
\Gamma_k(A)=B_k(A)+b_k^2\cdot B_k(A)\cdot\mathsf{Q}\cdot S^0_{k+1}(A)\cdot\tilde{Q}\cdot B_k(A),
\]
and $\delta Z_k(A)\ne0$.
\end{enumerate}
\end{proposition}

\begin{proof}
Let $G:=i\gamma_3\cdot y$. By Lemma~\ref{lem:block-averages-block}, $P$ is a projection, and $P$, $\mathsf{Q}$, $\tilde{Q}$ commute with every constant matrix acting on the spin components; in particular $P$ commutes with $G$.

(a) The projections $1-P$ and $P$ add up to $1$ and commute with $b_k+G$ and with $b_k+b_L+G$. Hence the inverse of
\[
b_k+b_L\cdot P+G=\left(b_k+G\right)\cdot(1-P)+\left(b_k+b_L+G\right)\cdot P
\]
is $B_{k,y}(A)$. All factors in the following computation are functions of $\gamma_3$ and commute with each other; by the definitions of $\boldsymbol{\alpha}_y$ and $\boldsymbol{\beta}_y$ in Section~\ref{sec:operators-2},
\begin{align*}
b_k-b_k^2\cdot\left(b_k+G\right)^{-1}&=b_k\cdot G\cdot\left(b_k+G\right)^{-1}=\boldsymbol{\alpha}_y,\\
b_k-b_k^2\cdot\left(b_k+b_L+G\right)^{-1}-\boldsymbol{\alpha}_y&=b_k\cdot\left[\left(b_L+G\right)\cdot\left(b_k+G\right)-G\cdot\left(b_k+b_L+G\right)\right]\cdot\left(b_k+b_L+G\right)^{-1}\cdot\left(b_k+G\right)^{-1}\\
&=b_k^2\cdot b_L\cdot\left(b_k+b_L+G\right)^{-1}\cdot\left(b_k+G\right)^{-1}=\boldsymbol{\beta}_y.
\end{align*}
Consequently $b_k-b_k^2\cdot B_{k,y}(A)=\boldsymbol{\alpha}_y\cdot(1-P)+\left(\boldsymbol{\alpha}_y+\boldsymbol{\beta}_y\right)\cdot P=\boldsymbol{\alpha}_y+\boldsymbol{\beta}_y\cdot P$. Since $P_k(A)=\tilde{Q}\cdot\mathsf{Q}$ and $P_{k+1}(A)=\tilde{Q}\cdot P\cdot\mathsf{Q}$ (Section~\ref{sec:phase-factors-lattice}), it follows that
\[
\tilde{Q}\cdot\left(b_k-b_k^2\cdot B_{k,y}(A)\right)\cdot\mathsf{Q}=\boldsymbol{\alpha}_y\cdot P_k(A)+\boldsymbol{\beta}_y\cdot P_{k+1}(A).
\]
As $\boldsymbol{\alpha}_{\infty}=b_k$ and $\boldsymbol{\beta}_{\infty}=0$, one has $F_{\infty}(A)=D_A+\mathfrak{m}+b_k\cdot P_k(A)=D_A+\mathfrak{m}+\tilde{Q}\cdot b_k\cdot\mathsf{Q}$, and therefore, by~\eqref{eq:operators-fermion-2},
\[
F_{\infty}(A)-b_k^2\cdot\tilde{Q}\cdot B_{k,y}(A)\cdot\mathsf{Q}=D_A+\mathfrak{m}+\operatorname{Pen}_y(A)=F_y(A).
\]

(b) Let $S:=S_k(A)$ and $S':=S_y(A)$; both operators exist by Theorem~\ref{thm:expansion-walks}(a). Let further $\check{B}:=B_{k,y}(A)$ and $Z:=b_k^2\cdot\tilde{Q}\cdot\check{B}\cdot\mathsf{Q}$. By (a), $S^{-1}-S'^{-1}=Z$, and therefore
\[
S'-S=S\cdot Z\cdot S'=S'\cdot Z\cdot S.
\]
Let $\mathsf{F}:=D_k(A)+b_L\cdot P+G$. By the definition of $D_k(A)$ and by (a), $\mathsf{F}=\check{B}^{-1}-b_k^2\cdot\mathsf{Q}\cdot S\cdot\tilde{Q}$. Since $\Gamma_{k,y}(A)=\check{B}+b_k^2\cdot\check{B}\cdot\mathsf{Q}\cdot S'\cdot\tilde{Q}\cdot\check{B}$, multiplying out gives
\[
\mathsf{F}\cdot\Gamma_{k,y}(A)=1+b_k^2\cdot\mathsf{Q}\cdot\left[S'-S-S\cdot Z\cdot S'\right]\cdot\tilde{Q}\cdot\check{B}=1.
\]
The operator $\mathsf{F}$ acts on a finite-dimensional space; so it is bijective, with the inverse $\Gamma_{k,y}(A)$.

(c) By the definitions of $H_k(A)$, $\mathfrak{H}_k(A)$ and $\Gamma_{k,y}(A)$,
\begin{align*}
S+H_k(A)\cdot\Gamma_{k,y}(A)\cdot\mathfrak{H}_k(A)&=S+S\cdot\tilde{Q}\cdot\left(b_k^2\cdot\check{B}+b_k^4\cdot\check{B}\cdot\mathsf{Q}\cdot S'\cdot\tilde{Q}\cdot\check{B}\right)\cdot\mathsf{Q}\cdot S\\
&=S+S\cdot Z\cdot S+S\cdot Z\cdot S'\cdot Z\cdot S=S+S\cdot Z\cdot\left(S+S'\cdot Z\cdot S\right).
\end{align*}
By the two expressions for $S'-S$ in (b), one has $S+S'\cdot Z\cdot S=S'$ and $S+S\cdot Z\cdot S'=S'$; hence the right-hand side equals $S'$.

(d) One has $B_{k,0}(A)=B_k(A)$ and, by Section~\ref{sec:operators-2}, $S_0(A)=S^0_{k+1}(A)$. So the first two assertions are (b) at $y=0$, and $\delta Z_k(A)$, the determinant of an invertible operator, does not vanish.
\end{proof}

\subsubsection{Bounds}\label{sec:bounds}

Let $b_{-}=b\cdot\left(1-L^{-1}\right)$ be the number of Section~\ref{sec:parameters-clifford-relations}, so that $b_{-}\le b_k\le b$, and let $a$ and $\gamma$ be the constants~\eqref{eq:operators-fermion-4}. Put
\begin{equation}\label{eq:operators-fermion-8}
C_{\mathrm{By}}:=\sqrt{2}\cdot\left(1+2\mathbf{e}\right)\cdot\max\left(1,\frac{1}{b_{-}}\right),\qquad C_{\mathrm{Gam}}:=\frac{2\left(1+2\mathbf{e}\right)}{b_{-}}+128\,\mathbf{e}\cdot\left(1+2\mathbf{e}\right)^2\cdot a.
\end{equation}

\begin{proposition}[Kernels on the unit lattice]\label{prop:kernels-unit-lattice}
Let $A\in\tilde{G}_k$ and $y\in\mathbb{R}$, and let $x$, $x'$ be unit sites.
\begin{enumerate}
\item[(a)] If $x,x'\in B(y')$ for a site $y'$ of the $L$-lattice, then $P(A)(x,x')=L^{-3}\cdot U_A(y',x)^{-1}\cdot U_A(y',x')$, and $P(A)(x,x')=0$ if no set $B(y')$ contains both $x$ and $x'$. In both cases $\left|P(A)(x,x')\right|\le\mathbf{e}\cdot L^{-3}$.
\item[(b)] $B_{k,y}(A)(x,x')=0$ unless $x$ and $x'$ lie in one $L$-block $B(y')$, and
\[
\left|B_{k,y}(A)(x,x')\right|\le\left(b_k^2+y^2\right)^{-\frac{1}{2}}\cdot\left(\mathbf{1}[x=x']+2\mathbf{e}\cdot L^{-3}\right)\le C_{\mathrm{By}}\cdot\left(1+|y|\right)^{-1}.
\]
Here $\mathbf{1}[\cdot]$ is $1$ if the condition in the brackets holds and $0$ otherwise, and the equality $x=x'$, like the property of lying in one $L$-block, refers to $x$ and $x'$ as sites of the torus $T^0_k$, that is, it is understood modulo $l\mathbb{Z}^3$. For fixed $x$ the sum of $\left|B_{k,y}(A)(x,x')\right|$ over $x'$, and for fixed $x'$ the sum over $x$, is at most $\left(1+2\mathbf{e}\right)\cdot\left(b_k^2+y^2\right)^{-\frac{1}{2}}$.
\item[(c)] For every linear operator $O$ on the spinors on $T_k$,
\[
\left|\left(\mathsf{Q}\cdot O\cdot\tilde{Q}\right)(x,x')\right|\le\mathbf{e}\cdot\left\|\mathbf{1}_{\Delta_x}\cdot O\cdot\mathbf{1}_{\Delta_{x'}}\right\|.
\]
\item[(d)] One has
\begin{align*}
\left|M_k(A)(x,x')\right|&\le32\,\mathbf{e}\cdot b^2\cdot a\cdot\exp\left(-\gamma\cdot d(x,x')\right),\\
\left|\dot{M}_k(A)(x,x')\right|&\le\mathbf{e}\cdot b^2\cdot(32\,a)^2\cdot\exp\left(-\gamma\cdot d(x,x')\right),\\
\left|\left(\Gamma_{k,y}(A)-B_{k,y}(A)\right)(x,x')\right|&\le128\,\mathbf{e}\cdot\left(1+2\mathbf{e}\right)^2\cdot a\cdot b_k^2\cdot\left(b_k^2+y^2\right)^{-1}\cdot\exp\left(-\gamma\cdot d(x,x')\right),\\
\left|\Gamma_k(A)(x,x')\right|&\le C_{\mathrm{Gam}}\cdot\exp\left(-\gamma\cdot d(x,x')\right).
\end{align*}
\end{enumerate}
\end{proposition}

\begin{proof}
By~\eqref{eq:operators-fermion-G1}, by condition~\eqref{eq:renormalisation-group-20} and since $e_k\le e_0$,
\[
\mathfrak{i}(A)<e_k^{\varepsilon}\le\frac{1}{6L}.
\]

(a) The formula for the kernel follows from $P(A)=Q^T(-A)\cdot Q(A)$ and the formulas for $Q(A)$ and $Q^T(-A)$ in Section~\ref{sec:phase-factors-lattice}. If $x,x'\in B(y')$, then $|x-y'|\le\frac{L-1}{2}$ and $|x'-y'|\le\frac{L-1}{2}$, and Lemma~\ref{lem:bounds-phases}(a) bounds the modulus of the kernel by $L^{-3}\cdot\exp\left(3L\cdot\mathfrak{i}(A)\right)\le L^{-3}\cdot\exp\left(\frac{1}{2}\right)\le\mathbf{e}\cdot L^{-3}$.

(b) By the definition of $B_{k,y}(A)$,
\[
B_{k,y}(A)=\boldsymbol{\beta}_1+\left(\boldsymbol{\beta}_2-\boldsymbol{\beta}_1\right)\cdot P,\qquad\boldsymbol{\beta}_1:=\left(b_k+i\gamma_3\cdot y\right)^{-1},\qquad\boldsymbol{\beta}_2:=\left(b_k+b_L+i\gamma_3\cdot y\right)^{-1}.
\]
By~\eqref{eq:operators-fermion-1}, $\gamma_3$ is Hermitian and $\gamma_3^2=1$. So $\boldsymbol{\beta}_1$ and $\boldsymbol{\beta}_2$ are normal matrices, with the eigenvalues $\left(b_k\pm iy\right)^{-1}$ and $\left(b_k+b_L\pm iy\right)^{-1}$ respectively, and both have norm at most $\left(b_k^2+y^2\right)^{-\frac{1}{2}}$. The kernel of $B_{k,y}(A)$ at $(x,x')$ is $\mathbf{1}[x=x']\cdot\boldsymbol{\beta}_1+P(A)(x,x')\cdot\left(\boldsymbol{\beta}_2-\boldsymbol{\beta}_1\right)$. By (a) it vanishes unless $x$ and $x'$ lie in one $L$-block, and its norm is at most $\left(b_k^2+y^2\right)^{-\frac{1}{2}}\cdot\left(\mathbf{1}[x=x']+2\mathbf{e}\cdot L^{-3}\right)$. Moreover
\[
\left(1+|y|\right)^2\le2\left(1+y^2\right)\le2\max\left(1,b_k^{-2}\right)\cdot\left(b_k^2+y^2\right),
\]
so that $\left(b_k^2+y^2\right)^{-\frac{1}{2}}\le\sqrt{2}\cdot\max\left(1,b_{-}^{-1}\right)\cdot\left(1+|y|\right)^{-1}$ by $b_{-}\le b_k$. Together with $\mathbf{1}[x=x']+2\mathbf{e}\cdot L^{-3}\le1+2\mathbf{e}$ and~\eqref{eq:operators-fermion-8} this gives the second inequality. An $L$-block has $\#|B(y')|=L^3$ sites; summing the first bound over the $x'$ in the $L$-block of $x$, or over the $x$ in the $L$-block of $x'$, gives the bound on the sums.

(c) Let $v\in\mathbb{C}^4$, and let $g$ be the spinor on $T_k$ with $g(x'')=\boldsymbol{\phi}_A(x',x'')^{-1}v$ for $x''\in\Delta_{x'}$ that vanishes outside the orbit of $\Delta_{x'}$. By the formula for $Q_k^T(-A)$ in Section~\ref{sec:phase-factors-lattice}, $\tilde{Q}\cdot\mathbf{1}_{x'}v=g$, so that
\[
\left(\mathsf{Q}\cdot O\cdot\tilde{Q}\right)(x,x')\,v=\left(\mathsf{Q}\cdot O\,g\right)(x),\qquad g=\mathbf{1}_{\Delta_{x'}}g,\qquad|g|_{\infty}\le\exp\left(\frac{3}{2}\,\mathfrak{i}(A)\right)\cdot|v|,
\]
the last bound by Lemma~\ref{lem:bounds-phases}(a). For every spinor $h$ on $T_k$, the formula for $Q_k(A)$, Lemma~\ref{lem:bounds-phases}(a) and $\eta^3\cdot\#|\Delta_x|=1$ give $\left|\left(\mathsf{Q}\,h\right)(x)\right|\le\exp\left(\frac{3}{2}\,\mathfrak{i}(A)\right)\cdot\left|\mathbf{1}_{\Delta_x}h\right|_{\infty}$. With $h:=O\,g$ it follows that
\[
\left|\left(\mathsf{Q}\cdot O\cdot\tilde{Q}\right)(x,x')\,v\right|\le\exp\left(3\,\mathfrak{i}(A)\right)\cdot\left\|\mathbf{1}_{\Delta_x}\cdot O\cdot\mathbf{1}_{\Delta_{x'}}\right\|\cdot|v|,
\]
and $\exp\left(3\,\mathfrak{i}(A)\right)\le\exp\left(\frac{1}{2L}\right)\le\mathbf{e}$.

(d) Theorem~\ref{thm:expansion-walks}(c) is applied with the resolvent parameter $\infty$, for which $S_{\infty}(A)=S_k(A)$ and $w_k(\infty)=1$. With (c), \eqref{eq:operators-fermion-EW1} and $b_k\le b$ it gives
\begin{align*}
\left|M_k(A)(x,x')\right|&\le b_k^2\cdot\mathbf{e}\cdot\left\|\mathbf{1}_{\Delta_x}\cdot S_k(A)\cdot\mathbf{1}_{\Delta_{x'}}\right\|\le b^2\cdot\mathbf{e}\cdot\exp\left(-\gamma\cdot d(x,x')\right)\cdot N_{\gamma}\left(S_k(A)\right)\\
&\le32\,\mathbf{e}\cdot b^2\cdot a\cdot\exp\left(-\gamma\cdot d(x,x')\right).
\end{align*}
For $\dot{M}_k(A)$ the same argument applies to $S_k(A)^2$ in place of $S_k(A)$, because $N_{\gamma}\left(S_k(A)^2\right)\le N_{\gamma}\left(S_k(A)\right)^2\le(32\,a)^2$ by~\eqref{eq:operators-fermion-EW2}.

By Proposition~\ref{prop:operators-fermion-1}(b),
\[
\Gamma_{k,y}(A)-B_{k,y}(A)=b_k^2\cdot B_{k,y}(A)\cdot\left(\mathsf{Q}\cdot S_y(A)\cdot\tilde{Q}\right)\cdot B_{k,y}(A).
\]
The kernel of the right-hand side at $(x,x')$ is the sum over unit sites $x_1$, $x_2$ of
\[
b_k^2\cdot B_{k,y}(A)(x,x_1)\cdot\left(\mathsf{Q}\cdot S_y(A)\cdot\tilde{Q}\right)(x_1,x_2)\cdot B_{k,y}(A)(x_2,x'),
\]
and by (b) only the terms with $x_1$ in the $L$-block of $x$ and $x_2$ in the $L$-block of $x'$ contribute. For these $d(x_1,x_2)\ge d(x,x')-2L$. By (c), \eqref{eq:operators-fermion-EW1} and Theorem~\ref{thm:expansion-walks}(c), where $w_k(y)=1$ because $k\ge1$,
\begin{align*}
\left|\left(\mathsf{Q}\cdot S_y(A)\cdot\tilde{Q}\right)(x_1,x_2)\right|&\le32\,\mathbf{e}\cdot a\cdot\exp\left(-\gamma\cdot d(x_1,x_2)\right)\le32\,\mathbf{e}\cdot a\cdot\exp\left(2\gamma L\right)\cdot\exp\left(-\gamma\cdot d(x,x')\right)\\
&\le128\,\mathbf{e}\cdot a\cdot\exp\left(-\gamma\cdot d(x,x')\right),
\end{align*}
since $\gamma L\le\frac{1}{6}<\log2$ (see~\eqref{eq:operators-fermion-4}). The bound of (b) on the sums over $x_1$ of $\left|B_{k,y}(A)(x,x_1)\right|$ and over $x_2$ of $\left|B_{k,y}(A)(x_2,x')\right|$ now yields the third inequality.

Finally, $\Gamma_k(A)=\Gamma_{k,0}(A)$ by Proposition~\ref{prop:operators-fermion-1}(b) and~(d). By (b) at $y=0$ and $b_k\ge b_{-}$,
\[
\left|B_k(A)(x,x')\right|\le\frac{1+2\mathbf{e}}{b_{-}}\cdot\mathbf{1}[x,x'\text{ lie in one }L\text{-block}],
\]
and, as two sites of one $L$-block have distance less than $L$,
\[
\mathbf{1}[x,x'\text{ lie in one }L\text{-block}]\le\exp(\gamma L)\cdot\exp\left(-\gamma\cdot d(x,x')\right)\le2\exp\left(-\gamma\cdot d(x,x')\right).
\]
Together with the third inequality at $y=0$, where $b_k^2\cdot\left(b_k^2+y^2\right)^{-1}=1$, and with~\eqref{eq:operators-fermion-8} this gives the bound on $\left|\Gamma_k(A)(x,x')\right|$.
\end{proof}

For $k\ge1$, Proposition~\ref{prop:kernel-bounds-2} follows from Proposition~\ref{prop:operators-fermion-1}(d) and Proposition~\ref{prop:kernels-unit-lattice}(b), (d).

\begin{proposition}[The operators $H_k$ and $\bar{H}_k$]\label{prop:operators}
Assume moreover condition~\eqref{eq:renormalisation-group-22}. Let $A\in\tilde{G}_k$, let $y$, $y'$ be unit sites, $v\in\mathbb{C}^4$, $f$ a spinor on the unit lattice, and $(x,x')\in\mathrm{Pz}_k$.
\begin{enumerate}
\item[(a)] One has
\begin{align*}
\left\|\mathbf{1}_{\Delta_y}\cdot H_k(A)\cdot\mathbf{1}_{y'}\right\|&\le32\,\mathbf{e}\cdot b\cdot a\cdot\exp\left(-\gamma\cdot d(y,y')\right),\\
\left\|\mathbf{1}_{\Delta_y}\cdot\delta_{\alpha,A}H_k(A)\cdot\mathbf{1}_{y'}\right\|&\le304\,\mathbf{e}\cdot b\cdot a\cdot\exp\left(-\gamma\cdot d(y,y')\right),
\end{align*}
where in the second bound $\delta_{\alpha,A}H_k(A)$ is taken with values in the functions on the short pairs.
\item[(b)] If $|x-x'|\ge1$, then $\left|U_A(x,x')^{\pm1}\right|\le\mathbf{e}$, and for every spinor $g$ on $T_k$
\[
\left|\left(\delta_{\alpha,A}g\right)(x,x')\right|\le\mathbf{e}\cdot|g(x')|+|g(x)|\le4\,|g|_{\infty}.
\]
In particular, if moreover $x\in\Delta_z$ and $x'\in\Delta_{z'}$ for unit sites $z$, $z'$, then
\[
\left|\left(\delta_{\alpha,A}H_k(A)\cdot\mathbf{1}_{y'}v\right)(x,x')\right|\le128\,\mathbf{e}\cdot b\cdot a\cdot\exp\left(-\gamma\cdot\min\left(d(z,y'),d(z',y')\right)\right)\cdot|v|.
\]
\item[(c)] One has
\[
\left|H_k(A)f\right|_{\infty}\le32\,\mathbf{e}\cdot b\cdot a\cdot|f|_{\infty},\qquad\sup_{(x,x')\in\mathrm{Pz}_k}\left|\left(\delta_{\alpha,A}H_k(A)f\right)(x,x')\right|\le304\,\mathbf{e}\cdot b\cdot a\cdot|f|_{\infty}.
\]
\end{enumerate}
The bounds (a)--(c) remain valid with $\bar{H}_k(A)=\mathsf{C}^{-1}\cdot H_k(-A)\cdot\mathsf{C}$ in place of $H_k(A)$ and $\delta_{\alpha,-A}$ in place of $\delta_{\alpha,A}$. For $k=0$ let $A\in\tilde{G}_0$, let $y$, $y'$ be unit sites, that is, sites of $T_0=T^0_0$, let $v\in\mathbb{C}^4$, let $f$ be a spinor on the unit lattice, and let $(x,x')\in\mathrm{Pz}_0$. Then $H_0(A)=1$,
\[
\bar{H}_0(A):=\mathsf{C}^{-1}\cdot H_0(-A)\cdot\mathsf{C}=1,
\]
there is no short pair, and
\[
\left|H_0^{\#}(A)f\right|_{\infty}\le4\,|f|_{\infty}.
\]
Moreover $\Delta_y=\{y\}$, and the operator $\mathbf{1}_{\Delta_y}\cdot H_0(A)\cdot\mathbf{1}_{y'}$ is $\mathbf{1}_{y'}$ if $y$ and $y'$ have the same orbit, in which case $d(y,y')=0$, and $0$ otherwise; thus $\mathbf{1}_{\Delta_y}\cdot H_0(A)\cdot\mathbf{1}_{y'}v$ equals $\mathbf{1}_{y'}v$ in the first case and $0$ in the second, and
\[
\left\|\mathbf{1}_{\Delta_y}\cdot H_0(A)\cdot\mathbf{1}_{y'}\right\|\le\exp\left(-\gamma\cdot d(y,y')\right).
\]
As there is no short pair, the operator $\delta_{\alpha,A}H_0(A)$, taken with values in the functions on the short pairs, is $0$, so that the second bound of (a) holds for $k=0$. With $z:=x$ and $z':=x'$, which satisfy $x\in\Delta_z$ and $x'\in\Delta_{z'}$,
\[
\left|\left(\delta_{\alpha,A}H_0(A)\cdot\mathbf{1}_{y'}v\right)(x,x')\right|\le4\cdot\exp\left(-\gamma\cdot\min\left(d(z,y'),d(z',y')\right)\right)\cdot|v|,
\]
and the left-hand side vanishes unless $x$ or $x'$ lies in the orbit of $y'$. These assertions on $H_0(A)$, the bound on $\left|H_0^{\#}(A)f\right|_{\infty}$ included, remain valid with $\bar{H}_0(A)$ in place of $H_0(A)$ and $\delta_{\alpha,-A}$ in place of $\delta_{\alpha,A}$. Here $H_k^{\#}(A):=\left(H_k(A),\delta_{\alpha,A}H_k(A)\right)$, as in Section~\ref{sec:fermion-integrals-blocking}, and $\left|H_k^{\#}(A)f\right|_{\infty}$ is the larger of $\sup_x\left|\left(H_k(A)f\right)(x)\right|$, the supremum over the sites, and $\sup_{(x,x')\in\mathrm{Pz}_k}\left|\left(\delta_{\alpha,A}H_k(A)f\right)(x,x')\right|$.
\end{proposition}

\begin{proof}
By~\eqref{eq:operators-fermion-G1}, by condition~\eqref{eq:renormalisation-group-20} and since $e_k\le e_0$, one has $\mathfrak{i}(A)<e_k^{\varepsilon}\le\frac{1}{6L}$. Hence, by Lemma~\ref{lem:bounds-phases}(a), $\left|\boldsymbol{\phi}_A(z,x'')^{\pm1}\right|\le\exp\left(\frac{3}{2}\,\mathfrak{i}(A)\right)\le\mathbf{e}$ for every unit site $z$ and every $x''\in\Delta_z$.

(a) By the definition of $H_k(A)$ and the formula for $Q_k^T(-A)$ in Section~\ref{sec:phase-factors-lattice},
\[
H_k(A)\cdot\mathbf{1}_{y'}v=b_k\cdot S_k(A)\,g,
\]
where $g$ is the spinor on $T_k$ with $g(x'')=\boldsymbol{\phi}_A(y',x'')^{-1}v$ for $x''\in\Delta_{y'}$ that vanishes outside the orbit of $\Delta_{y'}$. Thus $g=\mathbf{1}_{\Delta_{y'}}g$ and $|g|_{\infty}\le\mathbf{e}\cdot|v|$. Theorem~\ref{thm:expansion-walks}(c) with the resolvent parameter $\infty$, for which $S_{\infty}(A)=S_k(A)$ and $w_k(\infty)=1$, applied to the spinor $g=\mathbf{1}_{\Delta_{y'}}g$, together with $b_k\le b$ (Section~\ref{sec:parameters-clifford-relations}), gives both bounds.

(b) As $(x,x')\in\mathrm{Pz}_k$, one has $|x-x'|\le M$ (Definition~\ref{def:pairs-sites-holder}), so that the paths $\Gamma^{\pi}(x,x')$ have length $|x-x'|_1\le3M$. Lemma~\ref{lem:bounds-phases}(a), \eqref{eq:operators-fermion-G1} and condition~\eqref{eq:renormalisation-group-22}, with $e_k\le e_0$, give
\[
\left|U_A(x,x')^{\pm1}\right|\le\exp\left(3M\cdot\mathfrak{i}(A)\right)\le\exp\left(3M\cdot e_k^{\varepsilon}\right)\le\mathbf{e}.
\]
If $|x-x'|\ge1$, then $|x-x'|^{-\alpha}\le1$, and the definition of $\delta_{\alpha,A}$ in Section~\ref{sec:phase-factors-lattice} gives
\[
\left|\left(\delta_{\alpha,A}g\right)(x,x')\right|\le\left|U_A(x,x')\right|\cdot|g(x')|+|g(x)|\le\mathbf{e}\cdot|g(x')|+|g(x)|\le\left(\mathbf{e}+1\right)\cdot|g|_{\infty}\le4\,|g|_{\infty}.
\]
In particular, Lemma~\ref{lem:phase-factors} holds for the pairs $(x,x')\in\mathrm{Pz}_k$ with $|x-x'|\ge1$. Now let $g:=H_k(A)\cdot\mathbf{1}_{y'}v$, $x\in\Delta_z$ and $x'\in\Delta_{z'}$. By the first bound of (a),
\[
|g(x)|\le32\,\mathbf{e}\cdot b\cdot a\cdot\exp\left(-\gamma\cdot d(z,y')\right)\cdot|v|,\qquad|g(x')|\le32\,\mathbf{e}\cdot b\cdot a\cdot\exp\left(-\gamma\cdot d(z',y')\right)\cdot|v|,
\]
and, since $\left(\mathbf{e}+1\right)\cdot32\le128$, the last assertion of (b) follows.

(c) One has $H_k(A)f=b_k\cdot S_k(A)\cdot Q_k^T(-A)f$, and by the formula for $Q_k^T(-A)$ and Lemma~\ref{lem:bounds-phases}(a),
\[
\left|Q_k^T(-A)f\right|_{\infty}\le\exp\left(\frac{3}{2}\,\mathfrak{i}(A)\right)\cdot|f|_{\infty}\le\mathbf{e}\cdot|f|_{\infty}.
\]
By~\eqref{eq:operators-fermion-EW1} and Theorem~\ref{thm:expansion-walks}(c) with the resolvent parameter $\infty$, one has $\left\|S_k(A)\right\|\le32\,a$, and $\left\|\delta_{\alpha,A}S_k(A)\right\|\le304\,a$ for $\delta_{\alpha,A}S_k(A)$ taken on the short pairs. With $b_k\le b$ this gives the bound on $\left|H_k(A)f\right|_{\infty}$, and the bound on $\left|\left(\delta_{\alpha,A}H_k(A)f\right)(x,x')\right|$ for the short pairs. For the pairs $(x,x')\in\mathrm{Pz}_k$ with $|x-x'|\ge1$, the first assertion of (b), applied to $g:=H_k(A)f$, and the bound on $\left|H_k(A)f\right|_{\infty}$ give
\[
\left|\left(\delta_{\alpha,A}H_k(A)f\right)(x,x')\right|\le4\cdot32\,\mathbf{e}\cdot b\cdot a\cdot|f|_{\infty}\le304\,\mathbf{e}\cdot b\cdot a\cdot|f|_{\infty}.
\]

The operator $\bar{H}_k(A)$. By Section~\ref{sec:dirac-matrices-lattice}, $\mathsf{C}=i\gamma_1$, and by~\eqref{eq:operators-fermion-1} the matrix $\gamma_1$ is Hermitian with $\gamma_1^2=1$; hence $\mathsf{C}$ is unitary. By~\eqref{eq:operators-fermion-G2}, $-A\in\tilde{G}_k$, so that (a)--(c) hold for $H_k(-A)$ with $\delta_{\alpha,-A}$. The multiplications by the constant matrices $\mathsf{C}$ and $\mathsf{C}^{-1}$ commute with $\mathbf{1}_{\Delta_y}$, $\mathbf{1}_{y'}$ and $\delta_{\alpha,-A}$ and preserve the norm of $\mathbb{C}^4$; therefore the bounds carry over to $\bar{H}_k(A)=\mathsf{C}^{-1}\cdot H_k(-A)\cdot\mathsf{C}$, with $\delta_{\alpha,-A}$.

The case $k=0$. By Section~\ref{sec:exact-fermion-operators}, $H_0(A)=1$, and likewise $H_0(-A)=1$; hence $\bar{H}_0(A)=\mathsf{C}^{-1}\cdot\mathsf{C}=1$. As $\eta_0=1$, distinct sites have distance at least $1$, so that there is no short pair and every pair of $\mathrm{Pz}_0$ satisfies $|x-x'|\ge1$. The argument for the first assertion of (b) holds without change for $k=0$: for every spinor $g$ on the unit lattice and every $(x,x')\in\mathrm{Pz}_0$ it gives
\[
\left|\left(\delta_{\alpha,A}g\right)(x,x')\right|\le\mathbf{e}\cdot|g(x')|+|g(x)|\le4\,|g|_{\infty}.
\]
In particular $\left|\left(\delta_{\alpha,A}f\right)(x,x')\right|\le4\,|f|_{\infty}$ for all $(x,x')\in\mathrm{Pz}_0$. Since moreover $\left|H_0(A)f\right|_{\infty}=|f|_{\infty}$, one obtains $\left|H_0^{\#}(A)f\right|_{\infty}\le4\,|f|_{\infty}$. As $\eta_0=1$, the block $\Delta_y$ is the single site $y$, so that $\mathbf{1}_{\Delta_y}$ is the multiplication by the indicator function of the orbit of $y$. Two orbits are equal or disjoint; hence $\mathbf{1}_{\Delta_y}\cdot H_0(A)\cdot\mathbf{1}_{y'}=\mathbf{1}_{\Delta_y}\cdot\mathbf{1}_{y'}$ is $\mathbf{1}_{y'}$ if $y$ and $y'$ have the same orbit, and then $d(y,y')=0$, the distance $d$ being taken modulo $l\mathbb{Z}^3$; it vanishes otherwise. As $\mathbf{1}_{y'}\cdot\mathbf{1}_{y'}v=\mathbf{1}_{y'}v$ and $\left\|\mathbf{1}_{y'}\right\|\le1$, the assertions on $\mathbf{1}_{\Delta_y}\cdot H_0(A)\cdot\mathbf{1}_{y'}$ follow, with the constant $1$. Since there is no short pair, the functions on the short pairs form the space $\{0\}$, and $\delta_{\alpha,A}H_0(A)$ with values in this space is the zero operator. Now let $(x,x')\in\mathrm{Pz}_0$ and $g:=\mathbf{1}_{y'}v=H_0(A)\cdot\mathbf{1}_{y'}v$. The spinor $g$ vanishes outside the orbit of $y'$, and $|g(x'')|\le|v|$ at every site $x''$, its values on that orbit being $\pm v$. If neither $x$ nor $x'$ lies in the orbit of $y'$, the right-hand side of $\left|\left(\delta_{\alpha,A}g\right)(x,x')\right|\le\mathbf{e}\cdot|g(x')|+|g(x)|$ vanishes, and with it $\left(\delta_{\alpha,A}g\right)(x,x')$. Otherwise $\min\left(d(x,y'),d(x',y')\right)=0$, and the right-hand side is at most $\left(\mathbf{e}+1\right)\cdot|v|\le4\,|v|$; it is even at most $\mathbf{e}\cdot|v|$, because $|x-x'|>0$, so that $x$ and $x'$ lie in distinct orbits and at most one of them lies in the orbit of $y'$. With $z=x$ and $z'=x'$ this is the asserted bound on $\left(\delta_{\alpha,A}H_0(A)\cdot\mathbf{1}_{y'}v\right)(x,x')$, together with its vanishing unless $x$ or $x'$ lies in the orbit of $y'$. Finally, $-A\in\tilde{G}_0$ by~\eqref{eq:operators-fermion-G2}, and $\bar{H}_0(A)=1=H_0(-A)$. So the assertions on $\bar{H}_0(A)$ with $\delta_{\alpha,-A}$ are the assertions on $H_0(-A)$ with $\delta_{\alpha,-A}$; these hold, since everything above has been proved for an arbitrary connection in $\tilde{G}_0$, in particular for $-A$.
\end{proof}

Proposition~\ref{prop:bounds} follows from Proposition~\ref{prop:operators}: for $k\ge1$ from (a)--(c) and the assertion on $\bar{H}_k(A)$, and for $k=0$ from the assertions on the case $k=0$. The constants of the case $k=0$ are $1$ and $4$. By~\eqref{eq:operators-fermion-4}, $a=C_z\cdot c(\gamma)$ with $C_z\ge C_5\ge1$ (Proposition~\ref{prop:zero-field-levels}) and $c(\gamma)\ge1$ (the term $n=0$ of the sum that defines $c(\gamma)$ in Section~\ref{sec:norms-exponential-weights}); so $a\ge1$, and
\[
1<4\le304\,\mathbf{e}\cdot\max(1,b)\cdot a.
\]
Hence both constants lie below the maximum in the definition of the constant $C^F$ of Theorem~\ref{thm:fermion-operators}.

\subsubsection{The weakened families}\label{sec:weakened-families}

\begin{definition}[Weakened families of level $k\ge 1$]\label{def:weakened-families-level}
Let $s=(s_c)_{c\in\pi_k}\in\mathrm{Pd}_k$, $A\in\tilde{G}_k$, $X\in D_k$ and $y\ge 0$, and let $S_{\infty}(s,A)$, $S^{\mathrm{sy}}_0(s,A)$, $S^{\mathrm{sy}}_{\infty}(X,A)$ and $S_y(X,A)$ be the operators of Definition~\ref{def:weakened-sums-walks}. The families~\eqref{it:primitives-p3}, \eqref{it:primitives-p4} and~\eqref{it:primitives-p5} of a system of primitives of level $k$ on $T_k$ (Section~\ref{sec:primitive-operators}) are defined as follows. The family~\eqref{it:primitives-p3} consists of the operators
\[
H_k(s,A):=b_k\cdot S_{\infty}(s,A)\cdot Q_k^T(-A),\qquad\bar{H}_k(s,A):=\mathsf{C}^{-1}\cdot H_k(s,-A)\cdot\mathsf{C},
\]
where $-A\in\tilde{G}_k$ by~\eqref{eq:operators-fermion-G2}; they map the spinors on the unit lattice $T^0_k$ to the spinors on $T_k$. The family~\eqref{it:primitives-p4} consists of the operators
\[
\Gamma_k(s,A):=B_k(A)+b_k^2\cdot B_k(A)\cdot Q_k(A)\cdot S^{\mathrm{sy}}_0(s,A)\cdot Q_k^T(-A)\cdot B_k(A)
\]
and
\[
M_k(X,A):=b_k^2\cdot Q_k(A)\cdot S^{\mathrm{sy}}_{\infty}(X,A)\cdot Q_k^T(-A),
\]
which act on the spinors on $T^0_k$. The family~\eqref{it:primitives-p5} consists of the operators
\[
S_{k,y}(X,A):=S_y(X,A)
\]
on the spinors on $T_k$.
\end{definition}

\begin{theorem}[Properties of the weakened families]\label{thm:properties-weakened-families}
Let $L\ge 5$ be odd, let $0<\alpha_0\le\frac{1}{54}$ and $\kappa'>0$, and let conditions~\eqref{eq:renormalisation-group-19}, \eqref{eq:renormalisation-group-20} and~\eqref{eq:renormalisation-group-21} hold. Let $\gamma$ and $a$ be the constants~\eqref{eq:operators-fermion-4}, let $C_{\mathrm{Gam}}$ be the constant~\eqref{eq:operators-fermion-8}, and put
\[
C_1:=\max\left\{304\mathbf{e}\cdot\max(1,b)\cdot a\cdot M^{27\alpha_0},\ C_{\mathrm{Gam}}\cdot M^{27\alpha_0},\ 32\mathbf{e}\cdot\max\left(1,b^2\right)\cdot a\cdot\exp(27\kappa')\right\}.
\]
Let $k\ge 1$, $l\in\mathrm{Cl}$, $\mathfrak{m}\in(0,1]$ and $e_k\in(0,e_0]$. For the Euclidean norm $|v|$ on $\mathbb{C}^4$, the corresponding operator norm $|m|$ of $4\times 4$ matrices, and the sup norms $|f|_{\infty}$ and operator norms $\|O\|$ formed with them (Section~\ref{sec:norms-pairings-kernels}), the three families of Definition~\ref{def:weakened-families-level} satisfy the bounds of the properties~\ref{it:renormalisation-group-5}, \ref{it:renormalisation-group-6} and~\ref{it:renormalisation-group-7} of Section~\ref{sec:properties-systems-primitives} with the rates $\gamma$, $\kappa'$ and with the constant $C_1$ in place of $C$. In the sense of Section~\ref{sec:properties-systems-primitives}, in which the bounds of the properties~\ref{it:renormalisation-group-5}, \ref{it:renormalisation-group-6} and~\ref{it:renormalisation-group-7} are required for both norms on $\mathbb{C}^4$, the Euclidean norm and $|v|_c=\max_{\boldsymbol{\beta}}|v_{\boldsymbol{\beta}}|$, each with the corresponding operator norm of $4\times 4$ blocks and the corresponding sup norms, the three families satisfy each of the properties~\ref{it:renormalisation-group-1}, \ref{it:renormalisation-group-2} and~\ref{it:renormalisation-group-5}--\ref{it:renormalisation-group-16} of that section with the parameter set $(\alpha_0,\gamma,\kappa';2C_1)$. The relations of property~\ref{it:renormalisation-group-14} hold for every $g\in\mathbb{G}_k$, and, for the family~\eqref{it:primitives-p5}, for every $g\in\mathbb{G}_{k,3}$.
\end{theorem}

\begin{proof}
Let $w_s:=M^{27\alpha_0}$ and $w_X:=\exp(27\kappa')$ be the weights of Proposition~\ref{prop:properties-weakened-sums}. By~\eqref{eq:operators-fermion-4}, $C_V=12\mathbf{e}\cdot(2+b)\ge 1$ and $a=C_z\cdot c(\gamma)\ge 1$, since $C_z\ge C_5\ge 1$ (Proposition~\ref{prop:zero-field-levels}) and $c(\gamma)\ge 1$ (Section~\ref{sec:norms-exponential-weights}). Hence condition~\eqref{eq:renormalisation-group-19} yields $M\ge 4L+1>4L$ and $M\ge\left(108\pi\cdot C_V\cdot a\right)^2\ge 108\pi\cdot C_V\cdot a$; together with the remaining hypotheses this shows that Theorem~\ref{thm:expansion-walks} and Proposition~\ref{prop:properties-weakened-sums} apply. By the definitions of Section~\ref{sec:operators-unit-lattice},
\[
B_k(A)=B_{k,0}(A)=b_k^{-1}\cdot\left(1-P(A)\right)+\left(b_k+b_L\right)^{-1}\cdot P(A).
\]

\emph{Normalisation} (property~\ref{it:renormalisation-group-1}). Let $1$ denote the family with $s_c=1$ for every $c\in\pi_k$. By Proposition~\ref{prop:properties-weakened-sums}(b),
\[
S_{\infty}(1,A)=S_{\infty}(A),\qquad S^{\mathrm{sy}}_0(1,A)=S_0(A),\qquad\sum_{X\in D_k}S^{\mathrm{sy}}_{\infty}(X,A)=S_{\infty}(A),\qquad\sum_{X\in D_k}S_y(X,A)=S_y(A),
\]
and on $T_k$ one has $S_{\infty}(A)=S_k(A)$, $S_0(A)=S^0_{k+1}(A)$ and $S_y(A)=S_{k,y}(A)$ (Section~\ref{sec:operators-2}). With the definitions of $H_k(A)$ and $M_k(A)$ in Section~\ref{sec:operators-unit-lattice}, and with Proposition~\ref{prop:operators-fermion-1}(d) for $\Gamma_k(A)$, it follows that
\begin{align*}
H_k(1,A)&=b_k\cdot S_k(A)\cdot Q_k^T(-A)=H_k(A),\\
\Gamma_k(1,A)&=B_k(A)+b_k^2\cdot B_k(A)\cdot Q_k(A)\cdot S^0_{k+1}(A)\cdot Q_k^T(-A)\cdot B_k(A)=\Gamma_k(A),\\
\sum_{X\in D_k}M_k(X,A)&=b_k^2\cdot Q_k(A)\cdot S_k(A)\cdot Q_k^T(-A)=M_k(A),\\
\sum_{X\in D_k}S_{k,y}(X,A)&=S_{k,y}(A).
\end{align*}

\emph{Analyticity} (property~\ref{it:renormalisation-group-2}). By Proposition~\ref{prop:properties-weakened-sums}(c), the operators $S_{\infty}(s,A)$ and $S^{\mathrm{sy}}_0(s,A)$ are continuous functions of $(s,A)$ on $\mathrm{Pd}_k\times\tilde{G}_k$, analytic on the interior of this set, and analytic functions of $A\in\tilde{G}_k$ for every fixed $s\in\mathrm{Pd}_k$; and $S^{\mathrm{sy}}_{\infty}(X,\cdot)$, $S_y(X,\cdot)$ are analytic on $\tilde{G}_k$. Every matrix element of $Q_k(A)$, of $Q_k^T(-A)$ and of $B_k(A)$ depends on $A$ as an entire function: by the formulas of Section~\ref{sec:phase-factors-lattice} and the expression of $B_k(A)$ above, each of them is a linear combination, with coefficients independent of $A$, of products of exponentials of linear functions of $A$. Sums and products of operators with the properties just listed, of operators with entire entries and of the numbers $b_k$, $b_k^2$ again have these properties. This gives property~\ref{it:renormalisation-group-2} for $H_k(s,A)$, $\Gamma_k(s,A)$, $M_k(X,\cdot)$ and $S_{k,y}(X,\cdot)$. Moreover, $S_y(X,A)$ is a continuous function of $y\in\mathbb{R}$ (Proposition~\ref{prop:properties-weakened-sums}(c)); so $S_{k,y}(X,A)$ is a Borel function of $y\ge 0$, as Section~\ref{sec:primitive-operators} requires of the family~\eqref{it:primitives-p5}.

\emph{The bound of property~\ref{it:renormalisation-group-5}.} Let $f$ be a spinor on the unit lattice. For a unit site $z$ and a site $x\in\Delta_z$ one has $\left(Q_k^T(-A)f\right)(x)=\boldsymbol{\phi}_A(z,x)^{-1}\cdot f(z)$, and $\left|\boldsymbol{\phi}_A(z,x)^{-1}\right|\le\exp\left(\frac{3}{2}\mathfrak{i}(A)\right)$ by Lemma~\ref{lem:bounds-phases}(a). By~\eqref{eq:operators-fermion-G1} and condition~\eqref{eq:renormalisation-group-20}, $\mathfrak{i}(A)<e_k^{\varepsilon}\le\frac{1}{6L}<\frac{2}{3}$. Hence
\[
\left|Q_k^T(-A)f\right|_{\infty}\le\exp\left(\frac{3}{2}\mathfrak{i}(A)\right)\cdot|f|_{\infty}\le\mathbf{e}\cdot|f|_{\infty}.
\]
By~\eqref{eq:operators-fermion-EW1} and Proposition~\ref{prop:properties-weakened-sums}(a), where $w_k(\infty)=1$ by~\eqref{eq:operators-fermion-2},
\[
\left\|S_{\infty}(s,A)\right\|\le N_{\gamma}\left(S_{\infty}(s,A)\right)\le 32a\cdot w_s,\qquad\left\|\delta_{\alpha,A}S_{\infty}(s,A)\right\|\le N_{\gamma}\left(\delta_{\alpha,A}S_{\infty}(s,A)\right)\le 304a\cdot w_s,
\]
the operator $\delta_{\alpha,A}S_{\infty}(s,A)$ taking its values in the functions of short pairs. As $H_k(s,A)f=b_k\cdot S_{\infty}(s,A)\left(Q_k^T(-A)f\right)$ and $b_k\le b$ (Section~\ref{sec:parameters-clifford-relations}), it follows that
\begin{align*}
\left|H_k(s,A)f\right|_{\infty}&\le 32\mathbf{e}\cdot b\cdot a\cdot w_s\cdot|f|_{\infty},\\
\sup_{0<|x-x'|<1}\left|\left(\delta_{\alpha,A}H_k(s,A)f\right)(x,x')\right|&\le 304\mathbf{e}\cdot b\cdot a\cdot w_s\cdot|f|_{\infty}.
\end{align*}
Both constants are at most $C_1$.

\emph{The bounds of property~\ref{it:renormalisation-group-6}.} Let $x$ and $x'$ be unit sites. The bound for $\Gamma_k(s,A)$ is obtained as in the proof of Proposition~\ref{prop:kernels-unit-lattice}(d), the estimate of Theorem~\ref{thm:expansion-walks}(c) being replaced there by that of Proposition~\ref{prop:properties-weakened-sums}(a). Put $O:=S^{\mathrm{sy}}_0(s,A)$. Proposition~\ref{prop:properties-weakened-sums}(a), with $w_k(0)=1$, and~\eqref{eq:operators-fermion-EW1} give
\[
\left\|\mathbf{1}_{\Delta_{x_1}}\cdot O\cdot\mathbf{1}_{\Delta_{x_2}}\right\|\le 32a\cdot w_s\cdot\exp\left(-\gamma\cdot d(x_1,x_2)\right)
\]
for all unit sites $x_1$, $x_2$; hence, by Proposition~\ref{prop:kernels-unit-lattice}(c),
\[
\left|\left(Q_k(A)\cdot O\cdot Q_k^T(-A)\right)(x_1,x_2)\right|\le 32\mathbf{e}\cdot a\cdot w_s\cdot\exp\left(-\gamma\cdot d(x_1,x_2)\right).
\]
By Proposition~\ref{prop:kernels-unit-lattice}(b) at $y=0$, an entry $B_k(A)(z,z')$ vanishes unless the unit sites $z$ and $z'$ lie in one $L$-block, and the sum of the moduli of the entries of $B_k(A)$ along a row, and along a column, is at most $\frac{1+2\mathbf{e}}{b_k}$. In the expansion
\begin{align*}
&\left(B_k(A)\cdot Q_k(A)\cdot O\cdot Q_k^T(-A)\cdot B_k(A)\right)(x,x')\\
&\qquad=\sum_{x_1,x_2}B_k(A)(x,x_1)\cdot\left(Q_k(A)\cdot O\cdot Q_k^T(-A)\right)(x_1,x_2)\cdot B_k(A)(x_2,x')
\end{align*}
only the sites $x_1$ of the $L$-block of $x$ and the sites $x_2$ of the $L$-block of $x'$ contribute. For these sites $d(x_1,x_2)\ge d(x,x')-2L$, and $\exp(2\gamma\cdot L)\le 4$ because $\gamma\cdot L\le\frac{1}{6}<\log 2$ (Section~\ref{sec:inverses-zero-field}). Bounding the middle factor by $4\cdot 32\mathbf{e}\cdot a\cdot w_s\cdot\exp\left(-\gamma\cdot d(x,x')\right)$, the sum over $x_1$ by the row sum at $x$ and the sum over $x_2$ by the column sum at $x'$, one obtains
\begin{align*}
b_k^2\cdot\left|\left(B_k(A)\cdot Q_k(A)\cdot O\cdot Q_k^T(-A)\cdot B_k(A)\right)(x,x')\right|&\le b_k^2\cdot\left(\frac{1+2\mathbf{e}}{b_k}\right)^2\cdot 4\cdot 32\mathbf{e}\cdot a\cdot w_s\cdot\exp\left(-\gamma\cdot d(x,x')\right)\\
&=128\mathbf{e}\cdot(1+2\mathbf{e})^2\cdot a\cdot w_s\cdot\exp\left(-\gamma\cdot d(x,x')\right).
\end{align*}
By Proposition~\ref{prop:kernels-unit-lattice}(b) and $b_{-}\le b_k$ (Section~\ref{sec:parameters-clifford-relations}), moreover, $\left|B_k(A)(x,x')\right|\le\frac{1+2\mathbf{e}}{b_{-}}$, and $B_k(A)(x,x')=0$ unless $x$ and $x'$ lie in one $L$-block; in that case $d(x,x')<L$, so that $1\le\exp(\gamma\cdot L)\cdot\exp\left(-\gamma\cdot d(x,x')\right)\le 2\exp\left(-\gamma\cdot d(x,x')\right)$. As $w_s\ge 1$, these bounds and~\eqref{eq:operators-fermion-8} yield
\begin{align*}
\left|\Gamma_k(s,A)(x,x')\right|&\le\left(\frac{2(1+2\mathbf{e})}{b_{-}}+128\mathbf{e}\cdot(1+2\mathbf{e})^2\cdot a\cdot w_s\right)\cdot\exp\left(-\gamma\cdot d(x,x')\right)\\
&\le C_{\mathrm{Gam}}\cdot w_s\cdot\exp\left(-\gamma\cdot d(x,x')\right),
\end{align*}
and $C_{\mathrm{Gam}}\cdot w_s\le C_1$. For $M_k(X,A)$, Proposition~\ref{prop:kernels-unit-lattice}(c), the inequality $b_k\le b$ and the third bound of Proposition~\ref{prop:properties-weakened-sums}(a), with $w_k(\infty)=1$, give
\[
\left|M_k(X,A)(x,x')\right|\le b^2\cdot\mathbf{e}\cdot\left\|\mathbf{1}_{\Delta_x}\cdot S^{\mathrm{sy}}_{\infty}(X,A)\cdot\mathbf{1}_{\Delta_{x'}}\right\|\le 32\mathbf{e}\cdot b^2\cdot a\cdot w_X\cdot\exp\left(-\gamma\cdot d(x,x')-\kappa'\cdot d_M(X)\right),
\]
and $32\mathbf{e}\cdot b^2\cdot a\cdot w_X\le C_1$.

\emph{The bound of property~\ref{it:renormalisation-group-7}.} Let $y\ge 0$, and let $x$, $x'$ be unit sites. The third bound of Proposition~\ref{prop:properties-weakened-sums}(a), with $w_k(y)=1$, gives
\[
\left\|\mathbf{1}_{\Delta_x}\cdot S_{k,y}(X,A)\cdot\mathbf{1}_{\Delta_{x'}}\right\|\le 32a\cdot w_X\cdot\exp\left(-\gamma\cdot d(x,x')-\kappa'\cdot d_M(X)\right),
\]
and $32a\cdot w_X\le C_1$.

\emph{The maximum norm on $\mathbb{C}^4$.} In the three preceding paragraphs the norms of vectors of $\mathbb{C}^4$, of $4\times 4$ matrices, of functions and of operators are the Euclidean norms of Section~\ref{sec:norms-pairings-kernels}: the norm $|v|$ on $\mathbb{C}^4$, the operator norm $|m|$ of $4\times 4$ matrices, and the norms $|f|_{\infty}$ and $\|O\|$ formed with them. These paragraphs therefore prove the bounds of the properties~\ref{it:renormalisation-group-5}, \ref{it:renormalisation-group-6} and~\ref{it:renormalisation-group-7} with the constant $C_1$ for the Euclidean norms, which is the first assertion of the theorem. For $v\in\mathbb{C}^4$ the maximum norm $|v|_c=\max_{\boldsymbol{\beta}}|v_{\boldsymbol{\beta}}|$ satisfies $|v|_c\le|v|\le 2|v|_c$ (Section~\ref{sec:norms-pairings-kernels}). For a function $g$ with values in $\mathbb{C}^4$ let $\sup|g|$ and $\sup|g|_c$ be the suprema of $|g|$ and of $|g|_c$ over the domain of $g$. Let $O$ be a linear map between spaces of functions with values in $\mathbb{C}^4$, and let $\theta\ge 0$ be a number such that $\sup|Of|\le\theta\cdot\sup|f|$ for every $f$. Then $\sup|Of|_c\le\sup|Of|\le\theta\cdot\sup|f|\le 2\theta\cdot\sup|f|_c$ for every $f$. Applied with $\theta=C_1$ to $O=H_k(s,A)$ and to $O=\delta_{\alpha,A}H_k(s,A)$, whose values are functions on the short pairs, this gives the two bounds of property~\ref{it:renormalisation-group-5} for $|\cdot|_c$ with the constant $2C_1$. Applied with $\theta=C_1\cdot\exp\left(-\gamma\cdot d(x,x')-\kappa'\cdot d_M(X)\right)$ to $O=\mathbf{1}_{\Delta_x}\cdot S_{k,y}(X,A)\cdot\mathbf{1}_{\Delta_{x'}}$, it gives the bound of property~\ref{it:renormalisation-group-7} for $|\cdot|_c$ with the constant $2C_1$. A $4\times 4$ matrix $m$ is the case of the functions on a single point: for $v\in\mathbb{C}^4$ one has $|mv|_c\le|mv|\le|m|\cdot|v|\le 2|m|\cdot|v|_c$, so that the operator norm of $m$ for $|\cdot|_c$ is at most $2|m|$. Applied to the blocks $\Gamma_k(s,A)(x,x')$ and $M_k(X,A)(x,x')$, this gives the bounds of property~\ref{it:renormalisation-group-6} for $|\cdot|_c$ with the constant $2C_1$. As $C_1\le 2C_1$, the bounds for the Euclidean norms hold with $2C_1$ as well. Hence the three families have the properties~\ref{it:renormalisation-group-5}, \ref{it:renormalisation-group-6} and~\ref{it:renormalisation-group-7} for the parameter set $(\alpha_0,\gamma,\kappa';2C_1)$ in the sense of Section~\ref{sec:properties-systems-primitives}. The remaining properties~\ref{it:renormalisation-group-1}, \ref{it:renormalisation-group-2} and~\ref{it:renormalisation-group-8}--\ref{it:renormalisation-group-16} contain neither the constant nor a norm on $\mathbb{C}^4$.

\emph{Two facts on locality.} For the properties~\ref{it:renormalisation-group-8}--\ref{it:renormalisation-group-11} let $\mathscr{Y}$ be a set of cubes of $\pi_k$, let $Y$ be the union of the cubes of $\mathscr{Y}$, let $s\in\mathrm{Pd}_k$ satisfy $s_c=0$ for every cube $c\in\pi_k$ that does not belong to $\mathscr{Y}$, and let $Y_1$ be a link-component of $Y$, that is, the union of the cubes of a link-component of $\mathscr{Y}$ (Definition~\ref{def:parameter-cubes-polydisks}).
\begin{enumerate}
\item\label{it:operators-fermion-1} If $\boldsymbol{\omega}$ is a walk with $s_{\boldsymbol{\omega}}\ne 0$, then $X_{\boldsymbol{\omega}}\subset Y$, and $X_{\boldsymbol{\omega}}$ lies in a single link-component of $Y$. Indeed, $X_{\boldsymbol{\omega}}$ is a union of cubes of $\pi_k$, and $s_{\boldsymbol{\omega}}$ is the product of the parameters of these cubes (Definition~\ref{def:weakened-sums-walks}); if $X_{\boldsymbol{\omega}}\not\subset Y$, one of these cubes does not belong to $\mathscr{Y}$, and its parameter vanishes. If $X_{\boldsymbol{\omega}}\subset Y$, every cube of $X_{\boldsymbol{\omega}}$ belongs to $\mathscr{Y}$. Moreover $X_{\boldsymbol{\omega}}$ (Definition~\ref{def:walks}) is a polymer, as shown in Section~\ref{sec:expansion-walks} after that definition, and is therefore connected through faces (Section~\ref{sec:lattice-geometry}): any two of its cubes are joined by a chain of cubes of $X_{\boldsymbol{\omega}}$ in which consecutive cubes have a common face, that is, are linked.
\item\label{it:operators-fermion-2} Each unit block is contained in one $M$-cube, and for every point $y'$ of the $L$-lattice the set $\bigcup_{z\in B(y')}\Delta_z$, which contains the $L$-block $B(y')$, is contained in one $M$-cube as well. Indeed, by fact~\eqref{it:renormalisation-group-g0} of Section~\ref{sec:cubes-parameter-cubes} (here $L$ is odd, $M=L^m>4L$ and $l\in\mathrm{Cl}$), every unit block and every $L$-block lies in one $M$-cube, and the $M$-cubes partition the sites; for $z\in B(y')$ the site $z$ belongs both to $\Delta_z$ and to $B(y')$, so that $\Delta_z$ lies in the $M$-cube that contains $B(y')$. Consequently, for a union $X$ of $M$-cubes and a unit site $x'$, the block $\Delta_{x'}$ meets $X$ exactly when $x'\in X$; in that case $X$ contains $\Delta_{x'}$, the $L$-block of $x'$ and the unit blocks of the sites of this $L$-block. Moreover, the row of $Q_k(A)$ and the column of $Q_k^T(-A)$ at a unit site $x'$ read $A$ only on bonds inside $\Delta_{x'}$, and the row and the column of $B_k(A)$ at $x'$ read $A$ only on bonds inside the block of side $L$ of the fine lattice that contains $x'$ (Lemma~\ref{lem:block-averages-block} and the formulas of Section~\ref{sec:phase-factors-lattice}).
\end{enumerate}

\emph{Rows of the family~\eqref{it:primitives-p3}} (property~\ref{it:renormalisation-group-8}). Let $x$ be a fine site contained in $Y_1$. By Definition~\ref{def:weakened-sums-walks},
\[
H_k(s,A)(x,x')=b_k\cdot\sum_{\boldsymbol{\omega}\in\boldsymbol{\Omega}}s_{\boldsymbol{\omega}}\cdot\left(S_{\infty,\boldsymbol{\omega}}(A)\cdot Q_k^T(-A)\right)(x,x').
\]
The entry of the term of a walk $\boldsymbol{\omega}$ at $(x,x')$ involves the kernel of $S_{\infty,\boldsymbol{\omega}}(A)$ only at the pairs $(x,x'')$ with $x''\in\Delta_{x'}$; by Proposition~\ref{prop:properties-weakened-sums}(d) and fact~\ref{it:operators-fermion-2}, it therefore vanishes, for every $A\in\tilde{G}_k$, unless $x\in X_{\boldsymbol{\omega}}$ and $\Delta_{x'}$ meets $X_{\boldsymbol{\omega}}$, that is, unless $x,x'\in X_{\boldsymbol{\omega}}$. Let $\boldsymbol{\omega}$ be a walk whose term has a non-vanishing entry at $(x,x')$. Then $s_{\boldsymbol{\omega}}\ne 0$ and $x,x'\in X_{\boldsymbol{\omega}}$; by fact~\ref{it:operators-fermion-1}, $X_{\boldsymbol{\omega}}$ lies in a link-component of $Y$, which contains $x$ and is therefore $Y_1$; so $x'\in X_{\boldsymbol{\omega}}\subset Y_1$. In particular $H_k(s,A)(x,x')=0$ unless $x'$ is contained in $Y_1$. The term of a walk $\boldsymbol{\omega}$ depends on $s$ only through the parameters of the cubes of $X_{\boldsymbol{\omega}}$, and on $A$ only through its values on the bonds with both end points in $X_{\boldsymbol{\omega}}$: its entries at the columns $x'\notin X_{\boldsymbol{\omega}}$ vanish for every $A$, and at a column $x'\in X_{\boldsymbol{\omega}}$ the factor $S_{\infty,\boldsymbol{\omega}}(A)$ reads $A$ only on these bonds by Proposition~\ref{prop:properties-weakened-sums}(d), while the factor $Q_k^T(-A)$ reads $A$ only on bonds inside $\Delta_{x'}\subset X_{\boldsymbol{\omega}}$ by fact~\ref{it:operators-fermion-2}. Now let $\tilde{s}\in\mathrm{Pd}_k$ vanish at the cubes that do not belong to $\mathscr{Y}$ and coincide with $s$ at the cubes of $Y_1$, and let $\tilde{A}\in\tilde{G}_k$ coincide with $A$ on the bonds that meet $Y_1$. What has been shown applies to $(\tilde{s},\tilde{A})$ as well. Thus, for both pairs, the term of a walk $\boldsymbol{\omega}$ in the row at $x$ vanishes unless $x\in X_{\boldsymbol{\omega}}\subset Y_1$; and for a walk with $x\in X_{\boldsymbol{\omega}}\subset Y_1$ the two terms coincide, because $\tilde{s}_{\boldsymbol{\omega}}=s_{\boldsymbol{\omega}}$ and $\tilde{A}=A$ on the bonds with both end points in $X_{\boldsymbol{\omega}}$. Hence $H_k(s,A)(x,\cdot)=H_k(\tilde{s},\tilde{A})(x,\cdot)$.

\emph{Columns of the family~\eqref{it:primitives-p3}} (property~\ref{it:renormalisation-group-9}). Let $x'$ be a unit site contained in $Y_1$, let $\tilde{s}$, $\tilde{A}$ be as in the preceding paragraph, and let $\boldsymbol{\omega}$ be a walk whose term has a non-vanishing entry at $(x,x')$ for some fine site $x$. As shown there, $s_{\boldsymbol{\omega}}\ne 0$, $x\in X_{\boldsymbol{\omega}}$, and $\Delta_{x'}$ meets $X_{\boldsymbol{\omega}}$. By fact~\ref{it:operators-fermion-2}, $\Delta_{x'}$ lies in the $M$-cube that contains $x'$, which is a cube of $Y_1$; so this cube is a cube of $X_{\boldsymbol{\omega}}$, the set $X_{\boldsymbol{\omega}}$ meets $Y_1$, and $X_{\boldsymbol{\omega}}\subset Y_1$ by fact~\ref{it:operators-fermion-1}. Hence the row $x$ lies in $Y_1$, that is, $H_k(s,A)(x,x')=0$ unless $x$ is contained in $Y_1$. The term of a walk depends on $(s,A)$ in the way described in the preceding paragraph, and the same conclusions hold for $(\tilde{s},\tilde{A})$. Consequently the columns at $x'$ of $H_k(s,A)$ and of $H_k(\tilde{s},\tilde{A})$ are sums of equal terms over the same walks, namely the walks $\boldsymbol{\omega}$ with $X_{\boldsymbol{\omega}}\subset Y_1$ for which $\Delta_{x'}$ meets $X_{\boldsymbol{\omega}}$; the column at $x'$ thus depends on $s$ only through its values at the cubes of $Y_1$ and on $A$ only through its values on the bonds that meet $Y_1$.

\emph{Vanishing rows of the family~\eqref{it:primitives-p3}} (properties~\ref{it:renormalisation-group-10} and~\ref{it:renormalisation-group-11}). Let $x$ be a fine site at which the row of $H_k(s,A)$ does not vanish, and let $\boldsymbol{\omega}$ be a walk whose term has a non-vanishing entry in this row. Then $x\in X_{\boldsymbol{\omega}}\subset Y$ by Proposition~\ref{prop:properties-weakened-sums}(d) and fact~\ref{it:operators-fermion-1}; so $x$ is contained in a link-component of $Y$, which is property~\ref{it:renormalisation-group-10}. By Proposition~\ref{prop:properties-weakened-sums}(d), every $M$-cube that meets $K_x$ is contained in $X_{\boldsymbol{\omega}}$, hence in $Y$, and is therefore a cube of $\mathscr{Y}$; this is property~\ref{it:renormalisation-group-11}.

\emph{Rows and columns of $\Gamma_k$} (properties~\ref{it:renormalisation-group-8} and~\ref{it:renormalisation-group-9}). Let $x$, $x'$ be unit sites, and let $\tilde{s}$, $\tilde{A}$ be as in the proof of property~\ref{it:renormalisation-group-8} for the family~\eqref{it:primitives-p3}. The summand $B_k(A)$ in the definition of $\Gamma_k(s,A)$ is independent of $s$. By Proposition~\ref{prop:kernels-unit-lattice}(b), $B_k(A)(x,x')=0$ unless $x$ and $x'$ lie in one $L$-block, and then $x$ and $x'$ lie in one $M$-cube by fact~\ref{it:operators-fermion-2}; if $x$ or $x'$ is contained in $Y_1$, this $M$-cube is a cube of $Y_1$, and the entry reads $A$ only on bonds inside this cube (fact~\ref{it:operators-fermion-2}). By Definition~\ref{def:weakened-sums-walks}, the second term of $\Gamma_k(s,A)$ is the sum over the walks $\boldsymbol{\omega}$ of the operators $s_{\boldsymbol{\omega}}\cdot b_k^2\cdot B_k(A)\cdot Q_k(A)\cdot S^{\mathrm{sy}}_{0,\boldsymbol{\omega}}(A)\cdot Q_k^T(-A)\cdot B_k(A)$, where
\begin{align*}
&\left(B_k(A)\cdot Q_k(A)\cdot S^{\mathrm{sy}}_{0,\boldsymbol{\omega}}(A)\cdot Q_k^T(-A)\cdot B_k(A)\right)(x,x')\\
&\qquad=\sum_{x_1,x_2}B_k(A)(x,x_1)\cdot\left(Q_k(A)\cdot S^{\mathrm{sy}}_{0,\boldsymbol{\omega}}(A)\cdot Q_k^T(-A)\right)(x_1,x_2)\cdot B_k(A)(x_2,x'),
\end{align*}
and only the sites $x_1$ of the $L$-block of $x$ and the sites $x_2$ of the $L$-block of $x'$ occur. If a summand does not vanish, then $\Delta_{x_1}$ and $\Delta_{x_2}$ meet $X_{\boldsymbol{\omega}}$ by Proposition~\ref{prop:properties-weakened-sums}(d); by fact~\ref{it:operators-fermion-2}, $x_1$ and $x_2$ lie in $X_{\boldsymbol{\omega}}$, and so do their $L$-blocks, which contain $x$ and $x'$. Let now $x$ (for property~\ref{it:renormalisation-group-8}) or $x'$ (for property~\ref{it:renormalisation-group-9}) be contained in $Y_1$, and let the term of a walk $\boldsymbol{\omega}$ have a non-vanishing entry at $(x,x')$. Then $s_{\boldsymbol{\omega}}\ne 0$ and $X_{\boldsymbol{\omega}}$ meets $Y_1$, so that $X_{\boldsymbol{\omega}}\subset Y_1$ by fact~\ref{it:operators-fermion-1}; thus $x$ and $x'$ are contained in $Y_1$, and the term depends on $s$ only through the parameters of the cubes of $X_{\boldsymbol{\omega}}$ and on $A$ only through its values on bonds inside $X_{\boldsymbol{\omega}}$ (Proposition~\ref{prop:properties-weakened-sums}(d) and fact~\ref{it:operators-fermion-2}). The same holds with $(\tilde{s},\tilde{A})$ in place of $(s,A)$. Consequently, if $x$ (resp.\ $x'$) is contained in $Y_1$, then $\Gamma_k(s,A)(x,x')=0$ unless $x'$ (resp.\ $x$) is contained in $Y_1$, and the row of $\Gamma_k(s,A)$ at $x$ (resp.\ its column at $x'$) coincides with that of $\Gamma_k(\tilde{s},\tilde{A})$: both consist of equal entries of $B_k(A)$ and $B_k(\tilde{A})$ and of equal terms of the same walks, namely those with $X_{\boldsymbol{\omega}}\subset Y_1$.

\emph{Polymer families} (property~\ref{it:renormalisation-group-12}). Let $X\in D_k$. By Definition~\ref{def:weakened-sums-walks}, $S_y(X,A)$ and $S^{\mathrm{sy}}_{\infty}(X,A)$ are sums over the walks $\boldsymbol{\omega}$ with $X_{\boldsymbol{\omega}}=X$. By Proposition~\ref{prop:properties-weakened-sums}(d), their kernels vanish at $(x,x')$ unless both $x$ and $x'$ lie in $X$, and they do not change if $A$ is replaced by an element $\tilde{A}\in\tilde{G}_k$ that coincides with $A$ on the bonds with both end points in $X$, in particular by one that coincides with $A$ on the bonds that meet $X$. This is property~\ref{it:renormalisation-group-12} for $S_{k,y}(X,A)$. The entry $M_k(X,A)(x,x')$ involves the kernel of $S^{\mathrm{sy}}_{\infty}(X,A)$ only at the pairs of sites in $\Delta_x\times\Delta_{x'}$; if it does not vanish, then $\Delta_x$ and $\Delta_{x'}$ meet $X$, and $x,x'\in X$ by fact~\ref{it:operators-fermion-2}. Hence the kernel of $M_k(X,A)$ vanishes at $(x,x')$ unless $x,x'\in X$; and for $x,x'\in X$ the entry depends on $A$ only through $S^{\mathrm{sy}}_{\infty}(X,A)$, through the row of $Q_k(A)$ at $x$ and through the column of $Q_k^T(-A)$ at $x'$, which read $A$ only on bonds inside $\Delta_x$ and $\Delta_{x'}$, and these blocks lie in $X$ (fact~\ref{it:operators-fermion-2}). So property~\ref{it:renormalisation-group-12} holds for $M_k(X,A)$ as well.

\emph{Gauge covariance} (property~\ref{it:renormalisation-group-13}). Let $\boldsymbol{\lambda}$ be a real scalar on $T_k$, and let $\boldsymbol{\lambda}^u$ be its restriction to the unit lattice $T^0_k$; then $A-d\boldsymbol{\lambda}\in\tilde{G}_k$ by~\eqref{eq:operators-fermion-G2}. Lemma~\ref{lem:symmetries-elementary-operators}(a), applied to the closed real connection $\boldsymbol{\omega}:=d\boldsymbol{\lambda}$ and to the function $u:=\exp\left(i\,e_k\cdot\boldsymbol{\lambda}\right)$, whose restriction to the unit lattice is $\exp\left(i\,e_k\cdot\boldsymbol{\lambda}^u\right)$, gives
\begin{align*}
Q_k^T\left(-(A-d\boldsymbol{\lambda})\right)&=\exp\left(i\,e_k\cdot\boldsymbol{\lambda}\right)\cdot Q_k^T(-A)\cdot\exp\left(-i\,e_k\cdot\boldsymbol{\lambda}^u\right),\\
Q_k(A-d\boldsymbol{\lambda})&=\exp\left(i\,e_k\cdot\boldsymbol{\lambda}^u\right)\cdot Q_k(A)\cdot\exp\left(-i\,e_k\cdot\boldsymbol{\lambda}\right),\\
P(A-d\boldsymbol{\lambda})&=\exp\left(i\,e_k\cdot\boldsymbol{\lambda}^u\right)\cdot P(A)\cdot\exp\left(-i\,e_k\cdot\boldsymbol{\lambda}^u\right),
\end{align*}
and hence $B_k(A-d\boldsymbol{\lambda})=\exp\left(i\,e_k\cdot\boldsymbol{\lambda}^u\right)\cdot B_k(A)\cdot\exp\left(-i\,e_k\cdot\boldsymbol{\lambda}^u\right)$. By Proposition~\ref{prop:properties-weakened-sums}(e), for $S^{\#}=S$ and for $S^{\#}=S^{\mathrm{sy}}$,
\begin{align*}
S^{\#}_y(s,A-d\boldsymbol{\lambda})&=\exp\left(i\,e_k\cdot\boldsymbol{\lambda}\right)\cdot S^{\#}_y(s,A)\cdot\exp\left(-i\,e_k\cdot\boldsymbol{\lambda}\right),\\
S^{\#}_y(X,A-d\boldsymbol{\lambda})&=\exp\left(i\,e_k\cdot\boldsymbol{\lambda}\right)\cdot S^{\#}_y(X,A)\cdot\exp\left(-i\,e_k\cdot\boldsymbol{\lambda}\right).
\end{align*}
When these relations are inserted into the products of Definition~\ref{def:weakened-families-level}, the phases between adjacent factors cancel, and one obtains
\begin{align*}
H_k(s,A-d\boldsymbol{\lambda})&=\exp\left(i\,e_k\cdot\boldsymbol{\lambda}\right)\cdot H_k(s,A)\cdot\exp\left(-i\,e_k\cdot\boldsymbol{\lambda}^u\right),\\
\Gamma_k(s,A-d\boldsymbol{\lambda})&=\exp\left(i\,e_k\cdot\boldsymbol{\lambda}^u\right)\cdot\Gamma_k(s,A)\cdot\exp\left(-i\,e_k\cdot\boldsymbol{\lambda}^u\right),\\
M_k(X,A-d\boldsymbol{\lambda})&=\exp\left(i\,e_k\cdot\boldsymbol{\lambda}^u\right)\cdot M_k(X,A)\cdot\exp\left(-i\,e_k\cdot\boldsymbol{\lambda}^u\right),\\
S_{k,y}(X,A-d\boldsymbol{\lambda})&=\exp\left(i\,e_k\cdot\boldsymbol{\lambda}\right)\cdot S_{k,y}(X,A)\cdot\exp\left(-i\,e_k\cdot\boldsymbol{\lambda}\right).
\end{align*}

\emph{Lattice symmetries} (property~\ref{it:renormalisation-group-14}). Let $g=(S,a')\in\mathbb{G}_k$; it acts on $T_k$ and on $T^0_k$ through the morphism $x\mapsto r(S)x+a'$ of Section~\ref{sec:tori-nested-blocks}, and $gA\in\tilde{G}_k$ by~\eqref{eq:operators-fermion-G2}. The parameters $y=\infty$ and $y=0$ are neutral (Section~\ref{sec:operators-2}); so Proposition~\ref{prop:properties-weakened-sums}(f) gives
\[
S_{\infty}(gs,gA)\cdot g=g\cdot S_{\infty}(s,A),\qquad S^{\mathrm{sy}}_0(gs,gA)\cdot g=g\cdot S^{\mathrm{sy}}_0(s,A),\qquad S^{\mathrm{sy}}_{\infty}(gX,gA)\cdot g=g\cdot S^{\mathrm{sy}}_{\infty}(X,A),
\]
and, if $S\in\mathrm{Sp}_3$, also $S_y(gX,gA)\cdot g=g\cdot S_y(X,A)$ for every $y\ge 0$. Lemma~\ref{lem:symmetries-elementary-operators}(b) on $T_k$ gives $Q_k(gA)\cdot g=g\cdot Q_k(A)$, $Q_k^T(-gA)\cdot g=g\cdot Q_k^T(-A)$ and $P(gA)\cdot g=g\cdot P(A)$, hence $B_k(gA)\cdot g=g\cdot B_k(A)$. Composing these relations according to Definition~\ref{def:weakened-families-level}, one obtains
\[
H_k(gs,gA)\cdot g=g\cdot H_k(s,A),\qquad\Gamma_k(gs,gA)\cdot g=g\cdot\Gamma_k(s,A),\qquad M_k(gX,gA)\cdot g=g\cdot M_k(X,A)
\]
for every $g\in\mathbb{G}_k$, and $S_{k,y}(gX,gA)\cdot g=g\cdot S_{k,y}(X,A)$ for every $g\in\mathbb{G}_{k,3}$. As $\mathbb{G}_k^L\subset\mathbb{G}_k$ and $\mathbb{G}^L_{k,3}\subset\mathbb{G}_{k,3}$, property~\ref{it:renormalisation-group-14} follows, together with the last assertion of the theorem.

\emph{Transposition} (property~\ref{it:renormalisation-group-15}). By Lemma~\ref{lem:symmetries-elementary-operators}(c),
\[
Q_k(A)^T=Q_k^T(A)=Q_k^T\left(-(-A)\right),\qquad\left(Q_k^T(-A)\right)^T=Q_k(-A),\qquad P(A)^T=P(-A),
\]
so that $B_k(A)^T=B_k(-A)$. These operators have scalar kernels and therefore commute with $\mathsf{C}$ (Lemma~\ref{lem:block-averages-block}). The parameters $y=\infty$ and $y=0$ are neutral, $-A\in\tilde{G}_k$ by~\eqref{eq:operators-fermion-G2}, and Proposition~\ref{prop:properties-weakened-sums}(g) gives
\[
S^{\mathrm{sy}}_{\infty}(X,A)^T=\mathsf{C}^{-1}\cdot S^{\mathrm{sy}}_{\infty}(X,-A)\cdot\mathsf{C},\qquad S^{\mathrm{sy}}_0(s,A)^T=\mathsf{C}^{-1}\cdot S^{\mathrm{sy}}_0(s,-A)\cdot\mathsf{C}.
\]
Transposing the products of Definition~\ref{def:weakened-families-level} factor by factor, in the reversed order, one obtains
\begin{align*}
M_k(X,A)^T&=b_k^2\cdot Q_k(-A)\cdot\mathsf{C}^{-1}\cdot S^{\mathrm{sy}}_{\infty}(X,-A)\cdot\mathsf{C}\cdot Q_k^T(A)=\mathsf{C}^{-1}\cdot M_k(X,-A)\cdot\mathsf{C},\\
\Gamma_k(s,A)^T&=B_k(-A)+b_k^2\cdot B_k(-A)\cdot Q_k(-A)\cdot\mathsf{C}^{-1}\cdot S^{\mathrm{sy}}_0(s,-A)\cdot\mathsf{C}\cdot Q_k^T(A)\cdot B_k(-A)\\
&=\mathsf{C}^{-1}\cdot\Gamma_k(s,-A)\cdot\mathsf{C}.
\end{align*}

\emph{Complex conjugation} (property~\ref{it:renormalisation-group-16}). Here $\overline{A}\in\tilde{G}_k$ by~\eqref{eq:operators-fermion-G2}. By Lemma~\ref{lem:symmetries-elementary-operators}(d),
\[
\overline{Q_k(\overline{A})}=\varkappa^{-1}\cdot Q_k(-A)\cdot\varkappa,\qquad\overline{Q_k^T(-\overline{A})}=\varkappa^{-1}\cdot Q_k^T(A)\cdot\varkappa,\qquad\overline{P(\overline{A})}=\varkappa^{-1}\cdot P(-A)\cdot\varkappa;
\]
as $b_k$ and $b_L$ are real, also $\overline{B_k(\overline{A})}=\varkappa^{-1}\cdot B_k(-A)\cdot\varkappa$. By Proposition~\ref{prop:properties-weakened-sums}(h),
\[
\overline{S^{\#}_y(\overline{s},\overline{A})}=\varkappa^{-1}\cdot S^{\#}_y(s,-A)\cdot\varkappa,\qquad\overline{S^{\#}_y(X,\overline{A})}=\varkappa^{-1}\cdot S^{\#}_y(X,-A)\cdot\varkappa.
\]
Since $b_k$ is real, the complex conjugate of each product of Definition~\ref{def:weakened-families-level} is the product of the complex conjugates of its factors; therefore
\begin{align*}
\overline{H_k(\overline{s},\overline{A})}&=\varkappa^{-1}\cdot H_k(s,-A)\cdot\varkappa,\qquad\overline{\Gamma_k(\overline{s},\overline{A})}=\varkappa^{-1}\cdot\Gamma_k(s,-A)\cdot\varkappa,\\
\overline{M_k(X,\overline{A})}&=\varkappa^{-1}\cdot M_k(X,-A)\cdot\varkappa.
\end{align*}
\end{proof}


\subsection{Level zero}\label{sec:level-zero}

\begin{definition}[Families of level zero]\label{def:families-level-zero}
Let $k=0$, $s=(s_c)_{c\in\pi_0}\in\mathrm{Pd}_0$, $A\in\tilde{G}_0$, $X\in D_0$ and $y\ge 0$. With the operators of Definition~\ref{def:weakened-sums-walks} put
\[
\Gamma_0(s,A):=S^{\mathrm{sy}}_0(s,A),\qquad S_{0,y}(X,A):=S_y(X,A),
\]
and let $H_0(s,A)$ be the operator on the spinors on $T_0$ with the kernel
\[
H_0(s,A)(x,x'):=\delta_{xx'}\cdot\prod\left\{s_c:\ c\in\pi_0,\ c\ \text{meets}\ K_x\right\},
\]
the right-hand side being understood as this number times the unit $4\times 4$ matrix; here $K_x=\left\{x'':\ |x''-x|\le 2L\right\}$ (Section~\ref{sec:lattice-geometry}).
\end{definition}

\begin{theorem}[Properties at level zero]\label{thm:properties-level-zero}
Let the hypotheses of Theorem~\ref{thm:properties-weakened-families} hold, let $k=0$, $l\in\mathrm{Cl}$, $\mathfrak{m}\in(0,1]$ and $e_k\in(0,e_0]$, and let $C_1$ be the constant of Theorem~\ref{thm:properties-weakened-families}.
\begin{enumerate}
\item[(a)] The families $\Gamma_0(s,A)$ and $S_{0,y}(X,A)$ of Definition~\ref{def:families-level-zero} have the properties~\ref{it:renormalisation-group-1}, \ref{it:renormalisation-group-2}, \ref{it:renormalisation-group-6}--\ref{it:renormalisation-group-9} and~\ref{it:renormalisation-group-12}--\ref{it:renormalisation-group-16} for the parameters $(\alpha_0,\gamma,\kappa';2C_1)$ in the sense of Section~\ref{sec:properties-systems-primitives}, that is, with the bounds valid for both norms of a $4\times 4$ block, the Euclidean operator norm and the operator norm belonging to $|v|_c=\max_{\boldsymbol{\beta}}|v_{\boldsymbol{\beta}}|$. Here property~\ref{it:renormalisation-group-7} takes the form of the kernel bound of property~\ref{it:renormalisation-group-17} with $C=2C_1$:
\[
\left|S_{0,y}(X,A)(x,x')\right|\le 2C_1\cdot(1+y)^{-1}\cdot\exp\left(-\gamma\cdot d(x,x')-\kappa'\cdot d_M(X)\right).
\]
For the Euclidean operator norm of Section~\ref{sec:norms-pairings-kernels}, the bound of property~\ref{it:renormalisation-group-6} for $\Gamma_0(s,A)$ and this kernel bound hold with $C_1$ in place of $2C_1$.
\item[(b)] The family $H_0(s,A)$ of Definition~\ref{def:families-level-zero} has the properties~\ref{it:renormalisation-group-1}, \ref{it:renormalisation-group-2}, \ref{it:renormalisation-group-8}--\ref{it:renormalisation-group-11}, \ref{it:renormalisation-group-13}, \ref{it:renormalisation-group-14} and~\ref{it:renormalisation-group-16}, and it has property~\ref{it:renormalisation-group-5} with the constant $M^{8\alpha_0}$, which satisfies $M^{8\alpha_0}\le 4\exp(2\gamma)\cdot M^{27\alpha_0}$.
\end{enumerate}
\end{theorem}

\begin{proof}
At the beginning of the proof of Theorem~\ref{thm:properties-weakened-families} it is shown that the hypotheses of Theorem~\ref{thm:properties-weakened-families} imply those of Theorem~\ref{thm:expansion-walks} and of Proposition~\ref{prop:properties-weakened-sums}, and that $M\ge 4L+1>4L$; these statements are applied here with $k=0$. Let $w_s=M^{27\alpha_0}$ and $w_X=\exp(27\kappa')$ be as there.

(a) At $k=0$ one has $\eta=1$, and the unit blocks are the sites: $\Delta_x=\{x\}$. Hence $\left\|\mathbf{1}_{\Delta_x}\cdot O\cdot\mathbf{1}_{\Delta_{x'}}\right\|=\left|O(x,x')\right|$ for every operator $O$ on the spinors on $T_0$ and all sites $x$, $x'$. By Section~\ref{sec:operators-2}, $F_0(A)=D_A+\mathfrak{m}+b_L\cdot P(A)$, $S_0(A)=\Gamma_0(A)$ and $S_y(A)=S_{0,y}(A)$, and $y=0$ is neutral.

Property~\ref{it:renormalisation-group-1}: by Proposition~\ref{prop:properties-weakened-sums}(b),
\[
\Gamma_0(1,A)=S^{\mathrm{sy}}_0(1,A)=S_0(A)=\Gamma_0(A),\qquad\sum_{X\in D_0}S_{0,y}(X,A)=\sum_{X\in D_0}S_y(X,A)=S_y(A)=S_{0,y}(A).
\]
Property~\ref{it:renormalisation-group-2} is Proposition~\ref{prop:properties-weakened-sums}(c). By the same proposition the map $y\mapsto S_{0,y}(X,A)$ is continuous; so it is Borel measurable on $[0,\infty)$.

Properties~\ref{it:renormalisation-group-6} and~\ref{it:renormalisation-group-7}: by~\eqref{eq:operators-fermion-EW1} and Proposition~\ref{prop:properties-weakened-sums}(a), with $w_0(0)=1$ and $w_0(y)=(1+y)^{-1}$ for $y\ge 0$ (see~\eqref{eq:operators-fermion-2}),
\begin{align*}
\left|\Gamma_0(s,A)(x,x')\right|&\le\exp\left(-\gamma\cdot d(x,x')\right)\cdot N_{\gamma}\left(S^{\mathrm{sy}}_0(s,A)\right)\le 32a\cdot w_s\cdot\exp\left(-\gamma\cdot d(x,x')\right),\\
\left|S_{0,y}(X,A)(x,x')\right|&\le 32a\cdot w_X\cdot(1+y)^{-1}\cdot\exp\left(-\gamma\cdot d(x,x')-\kappa'\cdot d_M(X)\right),
\end{align*}
and $32a\cdot w_s\le C_1$, $32a\cdot w_X\le C_1$. These are bounds for the Euclidean operator norm of the blocks $\Gamma_0(s,A)(x,x')$ and $S_{0,y}(X,A)(x,x')$, because the identity $\left\|\mathbf{1}_{\Delta_x}\cdot O\cdot\mathbf{1}_{\Delta_{x'}}\right\|=\left|O(x,x')\right|$ at $k=0$ holds with the norms of Section~\ref{sec:norms-pairings-kernels} on both sides. For a $4\times 4$ matrix $m$ and $v\in\mathbb{C}^4$ one has $|mv|_c\le|mv|\le|m|\cdot|v|\le 2|m|\cdot|v|_c$ (Section~\ref{sec:norms-pairings-kernels}), so that the operator norm of $m$ for $|\cdot|_c$ is at most $2|m|$, as in the proof of Theorem~\ref{thm:properties-weakened-families}. Hence the bound for $\Gamma_0(s,A)(x,x')$ and the bound for $S_{0,y}(X,A)(x,x')$ hold for the operator norm belonging to $|\cdot|_c$ with $2C_1$ in place of $C_1$; as $C_1\le 2C_1$, they hold with $2C_1$ for the Euclidean operator norm as well. This is the assertion of (a) on the properties~\ref{it:renormalisation-group-6} and~\ref{it:renormalisation-group-7} for the parameters $(\alpha_0,\gamma,\kappa';2C_1)$, together with the bounds with $C_1$ for the Euclidean operator norm.

Properties~\ref{it:renormalisation-group-8} and~\ref{it:renormalisation-group-9}: let $\mathscr{Y}$, $Y$, $s$ and $Y_1$ be as in the proof of Theorem~\ref{thm:properties-weakened-families}, now with $k=0$, and let $\tilde{s}\in\mathrm{Pd}_0$ and $\tilde{A}\in\tilde{G}_0$ be as in property~\ref{it:renormalisation-group-8}. By Definition~\ref{def:weakened-sums-walks}, $\Gamma_0(s,A)=\sum_{\boldsymbol{\omega}\in\boldsymbol{\Omega}}s_{\boldsymbol{\omega}}\cdot S^{\mathrm{sy}}_{0,\boldsymbol{\omega}}(A)$. If the entry of the term $s_{\boldsymbol{\omega}}\cdot S^{\mathrm{sy}}_{0,\boldsymbol{\omega}}(A)$ at $(x,x')$ does not vanish, then $x,x'\in X_{\boldsymbol{\omega}}$ by Proposition~\ref{prop:properties-weakened-sums}(d), and $X_{\boldsymbol{\omega}}\subset Y$ lies in a single link-component of $Y$ by fact~\ref{it:operators-fermion-1} in the proof of Theorem~\ref{thm:properties-weakened-families}, namely in the one that contains $x$ and $x'$. The term depends on $s$ only through the parameters of the cubes of $X_{\boldsymbol{\omega}}$, and on $A$ only through its values on the bonds with both end points in $X_{\boldsymbol{\omega}}$ (Proposition~\ref{prop:properties-weakened-sums}(d)); the same holds for $(\tilde{s},\tilde{A})$. Hence, if $x$ (resp.\ $x'$) is contained in $Y_1$, then $\Gamma_0(s,A)(x,x')=0$ unless $x'$ (resp.\ $x$) is contained in $Y_1$, and the row at $x$ (resp.\ the column at $x'$) of $\Gamma_0(s,A)$ and of $\Gamma_0(\tilde{s},\tilde{A})$ are sums of equal terms over the same walks, namely those with $X_{\boldsymbol{\omega}}\subset Y_1$.

Properties~\ref{it:renormalisation-group-12}--\ref{it:renormalisation-group-16}: by Definition~\ref{def:weakened-sums-walks}, $S_{0,y}(X,A)$ is the sum of the operators $S_{y,\boldsymbol{\omega}}(A)$ over the walks with $X_{\boldsymbol{\omega}}=X$; by Proposition~\ref{prop:properties-weakened-sums}(d), its kernel vanishes at $(x,x')$ unless $x,x'\in X$, and it depends on $A$ only through the values on the bonds with both end points in $X$; this is property~\ref{it:renormalisation-group-12}. Property~\ref{it:renormalisation-group-13} is Proposition~\ref{prop:properties-weakened-sums}(e), since at $k=0$ the unit lattice is the lattice of $T_0$ and $\boldsymbol{\lambda}^u=\boldsymbol{\lambda}$. Property~\ref{it:renormalisation-group-14} is Proposition~\ref{prop:properties-weakened-sums}(f): $y=0$ is neutral, and the elements $g=(S,a')$ of $\mathbb{G}^L_{0,3}$, which occur for $S_{0,y}(X,A)$, have $S\in\mathrm{Sp}_3$. Property~\ref{it:renormalisation-group-15} is Proposition~\ref{prop:properties-weakened-sums}(g) at the neutral parameter $y=0$, and property~\ref{it:renormalisation-group-16} is Proposition~\ref{prop:properties-weakened-sums}(h).

(b) For a site $x$ of $T_0$ let $\mathrm{Cu}(x)\subset\pi_0$ be the set of the $M$-cubes $c$ with $c\cap K_x\ne\emptyset$. The sup-ball $K_x$ of radius $2L$ about $x$ has edges of length $4L$, and $4L<M$; its projection to each coordinate axis consists of $4L+1\le M$ consecutive integers. The $M$-cubes cut every coordinate axis into intervals of $M$ consecutive integers, and a set of at most $M$ consecutive integers meets at most two of these intervals, which are then adjacent. An $M$-cube meets $K_x$ exactly when, in each of the three directions, its interval meets the projection of $K_x$; therefore $\mathrm{Cu}(x)$ is the Cartesian product of three sets of intervals, one for each direction, each of which consists of a single interval or of two adjacent ones. Consequently $\#\left|\mathrm{Cu}(x)\right|\le 8$, the union of the cubes of $\mathrm{Cu}(x)$ is connected through faces, and $\mathrm{Cu}(x)$ contains the $M$-cube that contains $x$. By Definition~\ref{def:families-level-zero}, $H_0(s,A)$ is diagonal, and its entry at $(x,x)$ is $\prod_{c\in\mathrm{Cu}(x)}s_c$ times the unit matrix.

Property~\ref{it:renormalisation-group-1}: for $s=1$ all factors equal $1$, so that $H_0(1,A)$ is the identity, which is $H_0(A)$ (Section~\ref{sec:exact-fermion-operators}).

Property~\ref{it:renormalisation-group-2}: the entries of $H_0(s,A)$ are monomials in $s$ and do not depend on $A$.

Property~\ref{it:renormalisation-group-5}: the operator $H_0(s,A)$ multiplies the value $f(x)$ of a spinor $f$ at each site $x$ by the number $\prod_{c\in\mathrm{Cu}(x)}s_c$. As $|s_c|\le M^{\alpha_0}$ (Definition~\ref{def:parameter-cubes-polydisks}), $M^{\alpha_0}\ge 1$ and $\#\left|\mathrm{Cu}(x)\right|\le 8$, it follows that, for the Euclidean norm and for the norm $|\cdot|_c$ of Section~\ref{sec:norms-pairings-kernels} alike,
\[
\left|\left(H_0(s,A)f\right)(x)\right|=\left|\prod_{c\in\mathrm{Cu}(x)}s_c\right|\cdot\left|f(x)\right|\le M^{8\alpha_0}\cdot\left|f(x)\right|
\]
at every site $x$. Hence $\left|H_0(s,A)f\right|_{\infty}\le M^{8\alpha_0}\cdot|f|_{\infty}$ for every spinor $f$, and the same bound, with the same constant $M^{8\alpha_0}$, holds for the sup norms formed with $|\cdot|_c$. In the second bound of property~\ref{it:renormalisation-group-5} the supremum extends over the pairs $(z,z')\in\mathrm{Pz}_0$ with $d(z,z')<1$, that is, over the short pairs of Section~\ref{sec:phase-factors-lattice}. The two sites of a pair of $\mathrm{Pz}_0$ are distinct (Definition~\ref{def:pairs-sites-holder}), and at $k=0$, where $\eta=1$, two distinct sites are at distance at least $1$. Hence there is no such pair, this supremum extends over the empty set, and the second bound holds trivially. Finally $M^{8\alpha_0}\le M^{27\alpha_0}\le 4\exp(2\gamma)\cdot M^{27\alpha_0}$.

Properties~\ref{it:renormalisation-group-10} and~\ref{it:renormalisation-group-11}: let $\mathscr{Y}$, $Y$, $s$ and $Y_1$ be as in property~\ref{it:renormalisation-group-8}. If the entry of $H_0(s,A)$ at $(x,x)$ does not vanish, then $s_c\ne 0$ for every $c\in\mathrm{Cu}(x)$, so that $\mathrm{Cu}(x)\subset\mathscr{Y}$; this is property~\ref{it:renormalisation-group-11}. In particular the $M$-cube that contains $x$ belongs to $\mathscr{Y}$, and $x$ is contained in $Y$, hence in a link-component of $Y$; this is property~\ref{it:renormalisation-group-10}. Whenever $\mathrm{Cu}(x)\subset\mathscr{Y}$, the cubes of $\mathrm{Cu}(x)$, whose union is connected through faces, belong to one link-component of $\mathscr{Y}$; so $\mathrm{Cu}(x)$ consists of cubes of the link-component of $Y$ that contains $x$.

Properties~\ref{it:renormalisation-group-8} and~\ref{it:renormalisation-group-9}: let $x$ be contained in $Y_1$, and let $\tilde{s}\in\mathrm{Pd}_0$ vanish at the cubes that do not belong to $\mathscr{Y}$ and coincide with $s$ at the cubes of $Y_1$. Since $H_0(s,A)$ is diagonal, its row and its column at $x$ contain only the entry at $(x,x)$, and $x$ is contained in $Y_1$. If $\mathrm{Cu}(x)\subset\mathscr{Y}$, then $\mathrm{Cu}(x)$ consists of cubes of $Y_1$ by the preceding paragraph, and $\prod_{c\in\mathrm{Cu}(x)}\tilde{s}_c=\prod_{c\in\mathrm{Cu}(x)}s_c$; otherwise $\mathrm{Cu}(x)$ contains a cube that does not belong to $\mathscr{Y}$, and both products vanish. As $H_0(s,A)$ does not depend on $A$, properties~\ref{it:renormalisation-group-8} and~\ref{it:renormalisation-group-9} follow.

Properties~\ref{it:renormalisation-group-13} and~\ref{it:renormalisation-group-16}: the operator $H_0(s,A)$ is diagonal, scalar in the spin indices and independent of $A$. Hence it commutes with the multiplication by $\exp\left(\pm i\,e_k\cdot\boldsymbol{\lambda}\right)$ and with every constant spin matrix, so that $H_0(s,A-d\boldsymbol{\lambda})=\exp\left(i\,e_k\cdot\boldsymbol{\lambda}\right)\cdot H_0(s,A)\cdot\exp\left(-i\,e_k\cdot\boldsymbol{\lambda}\right)$ for every real scalar $\boldsymbol{\lambda}$, which is property~\ref{it:renormalisation-group-13} because $\boldsymbol{\lambda}^u=\boldsymbol{\lambda}$ at $k=0$. Moreover $\overline{\prod_c\overline{s_c}}=\prod_cs_c$, so that $\overline{H_0(\overline{s},\overline{A})}=H_0(s,-A)=\varkappa^{-1}\cdot H_0(s,-A)\cdot\varkappa$, which is property~\ref{it:renormalisation-group-16}.

Property~\ref{it:renormalisation-group-14}: let $g=(S,a')\in\mathbb{G}_0$. The map $x\mapsto gx=r(S)x+a'$ preserves the sup distance and maps $M$-cubes onto $M$-cubes; so it maps $K_x$ onto $K_{gx}$, and $\mathrm{Cu}(gx)=g\,\mathrm{Cu}(x)$. With $(gs)_c=s_{g^{-1}c}$ it follows that
\[
\prod_{c\in\mathrm{Cu}(gx)}(gs)_c=\prod_{c\in\mathrm{Cu}(x)}s_c,
\]
that is, the entry of $H_0(gs,gA)$ at $(gx,gx)$ equals the entry of $H_0(s,A)$ at $(x,x)$. Since $(gf)(gx)=S\cdot f(x)$ for every spinor $f$, this gives $H_0(gs,gA)\cdot g=g\cdot H_0(s,A)$.

By (a) and (b), property~\ref{it:renormalisation-group-17} for the families $H_0(s,A)$, $\Gamma_0(s,A)$ and $S_{0,y}(X,A)$ follows, with the constant $2C_1$ in place of $C$, in the sense of Section~\ref{sec:properties-systems-primitives}, that is, for both norms on $\mathbb{C}^4$. With $C_1$ in place of $C$, the bounds of property~\ref{it:renormalisation-group-17} for $\Gamma_0(s,A)$ and $S_{0,y}(X,A)$ hold for the Euclidean norms of Section~\ref{sec:norms-pairings-kernels}. For the family $H_0(s,A)$, property~\ref{it:renormalisation-group-17} requires the bound of property~\ref{it:renormalisation-group-5} with the constant $4\exp(2\gamma)\cdot M^{27\alpha_0}$, which does not involve $C$; by (b), this bound holds for both norms.
\end{proof}

\subsection{The averaged operator $T_k$}\label{sec:averaged-operator}

\begin{lemma}[Gauge fixing along rectilinear paths]\label{lem:gauge-fixing-along}
Let $G$ be one of the two sets $\eta\mathbb{Z}$ and $\frac{\eta}{2}+\eta\mathbb{Z}$. An interval of $G$ is a set $\{a,a+\eta,\dots,a+m\eta\}$ with $a\in G$ and an integer $m\ge 0$; its length is $m\eta$. Let $\mathbb{W}=I_0\times I_1\times I_2$, where $I_0,I_1,I_2$ are intervals of $G$, each of length at most $s_W$. The bonds of $\mathbb{W}$ are the ordered pairs $(x,x\pm\eta e_{\nu})$, $\nu\in\{0,1,2\}$, of sites of $\mathbb{W}$, and the plaquettes of $\mathbb{W}$ are the plaquettes of the lattice $G^3$ whose four corners lie in $\mathbb{W}$. A lattice path in $\mathbb{W}$ is a finite sequence of bonds of $\mathbb{W}$, each of which begins at the end point of the preceding one; it is closed if it ends at its initial site.

Let $A$ be a complex function on the bonds of $\mathbb{W}$ with $A(x',x)=-A(x,x')$, and for a lattice path $\Gamma$ in $\mathbb{W}$ let $A(\Gamma)$ be the sum of the values of $A$ on the bonds of $\Gamma$, each bond being oriented along $\Gamma$. Points of $\mathbb{R}^3$ are written $x=(x_0,x_1,x_2)$. Fix a site $o=(o_0,o_1,o_2)$ of $\mathbb{W}$, and for a site $x$ of $\mathbb{W}$ let $\Gamma(o,x)$ be the rectilinear lattice path from $o$ to $x$ that changes the coordinates in the order $0,1,2$: it follows the straight segments from $o$ to $(x_0,o_1,o_2)$, from $(x_0,o_1,o_2)$ to $(x_0,x_1,o_2)$, and from $(x_0,x_1,o_2)$ to $x$, and it lies in $\mathbb{W}$ because $\mathbb{W}$ is a product of intervals. Put
\[
\boldsymbol{\lambda}(x):=\eta\cdot A\left(\Gamma(o,x)\right),\qquad A'':=A-d\boldsymbol{\lambda},
\]
where $d\boldsymbol{\lambda}$, a function on the bonds of $\mathbb{W}$, and $dA$, a function on the oriented plaquettes of $\mathbb{W}$, are given by the formulas of Section~\ref{sec:tori-nested-blocks}. Then $\boldsymbol{\lambda}$ is real whenever $A$ is real. For every bond $b$ of $\mathbb{W}$,
\[
\left|A''(b)\right|\le 2s_W\cdot\sup_p\left|(dA)(p)\right|,
\]
the supremum being taken over the plaquettes $p$ of $\mathbb{W}$. For every closed lattice path $\boldsymbol{\ell}$ in $\mathbb{W}$ that consists of $n$ bonds,
\[
A(\boldsymbol{\ell})=A''(\boldsymbol{\ell}),\qquad \left|A(\boldsymbol{\ell})\right|\le 2n\cdot s_W\cdot\sup_p\left|(dA)(p)\right|.
\]
\end{lemma}

\begin{proof}
If $A$ is real, every number $\boldsymbol{\lambda}(x)$ is $\eta$ times a sum of real numbers. Since $A$ and $d\boldsymbol{\lambda}$ are antisymmetric under the reversal of the bonds, so is $A''$, and it suffices to bound $A''(b)$ for $b=(x,x+\eta e_{\nu})$. By the definition of $d\boldsymbol{\lambda}$,
\[
A''(b)=A(b)-\frac{\boldsymbol{\lambda}(x+\eta e_{\nu})-\boldsymbol{\lambda}(x)}{\eta}=A\left(\Gamma(o,x)\right)+A(b)-A\left(\Gamma(o,x+\eta e_{\nu})\right).
\]
By the antisymmetry of $A$, the right side is the value of $A$ on the closed lattice path $\boldsymbol{\ell}_b$ that follows $\Gamma(o,x)$, then the bond $b$, and then $\Gamma(o,x+\eta e_{\nu})$ in the reverse direction; this value does not change if bonds that $\boldsymbol{\ell}_b$ traverses once in each direction are removed from it. There are three cases.

Let $\nu=2$. The paths $\Gamma(o,x)$ and $\Gamma(o,x+\eta e_2)$ have the same first two segments, and their last segments lie on one line parallel to $e_2$ and lead from the common point $(x_0,x_1,o_2)$ to $x$ and to $x+\eta e_2$, respectively. Hence the longer of the two paths arises from the shorter one by appending the bond that joins $x$ and $x+\eta e_2$, every bond of $\boldsymbol{\ell}_b$ is traversed once in each direction, and $A''(b)=0$.

Let $\nu=1$. The two paths have the same first segment. Their second segments lie on one line and differ by the bond that joins $(x_0,x_1,o_2)$ and $(x_0,x_1+\eta,o_2)$, and their third segments are parallel to $e_2$ at the distance $\eta$. After the removal of the bonds traversed in both directions, $\boldsymbol{\ell}_b$ is the boundary of the rectangle with the corners
\[
(x_0,x_1,o_2),\quad (x_0,x_1+\eta,o_2),\quad (x_0,x_1+\eta,x_2),\quad (x_0,x_1,x_2),
\]
which is the union of $\frac{|x_2-o_2|}{\eta}$ plaquettes parallel to $e_1$ and $e_2$.

Let $\nu=0$. The first segments of the two paths lie on one line and differ by one bond, and their second segments, as well as their third segments, are parallel at the distance $\eta$. After the removal of the bonds traversed in both directions, $\boldsymbol{\ell}_b$ is the boundary of the union of two rectangles with a common side: the rectangle with the corners
\[
(x_0,o_1,o_2),\quad (x_0+\eta,o_1,o_2),\quad (x_0+\eta,x_1,o_2),\quad (x_0,x_1,o_2),
\]
the union of $\frac{|x_1-o_1|}{\eta}$ plaquettes parallel to $e_0$ and $e_1$, and the rectangle with the corners
\[
(x_0,x_1,o_2),\quad (x_0+\eta,x_1,o_2),\quad (x_0+\eta,x_1,x_2),\quad (x_0,x_1,x_2),
\]
the union of $\frac{|x_2-o_2|}{\eta}$ plaquettes parallel to $e_0$ and $e_2$.

In each case every coordinate of a corner lies between the corresponding coordinates of $o$ and of $x$ or $x+\eta e_{\nu}$; since $\mathbb{W}$ is a product of intervals, the rectangles lie in $\mathbb{W}$. By the definition of $dA$, the sum of $A$ over the boundary of an oriented plaquette $p$ is $\eta\cdot(dA)(p)$. When the plaquettes of a rectangle, or of the union of the two rectangles above, are oriented coherently, the bonds inside cancel in the sum of these boundary sums; hence the sum of $A$ over the boundary is a sum of terms $\pm\eta\cdot(dA)(p)$, one for each plaquette $p$. The number of these plaquettes is at most
\[
\frac{|x_1-o_1|+|x_2-o_2|}{\eta}\le\frac{2s_W}{\eta},
\]
because $o$ and $x$ lie in $\mathbb{W}$. It follows that
\[
\left|A''(b)\right|\le\frac{2s_W}{\eta}\cdot\eta\cdot\sup_p\left|(dA)(p)\right|=2s_W\cdot\sup_p\left|(dA)(p)\right|.
\]
Finally, let $\boldsymbol{\ell}$ be a closed lattice path in $\mathbb{W}$ with $n$ bonds. The sum of $d\boldsymbol{\lambda}$ over $\boldsymbol{\ell}$ is $\eta^{-1}$ times a telescoping sum of differences of values of $\boldsymbol{\lambda}$, and it vanishes. Therefore $A(\boldsymbol{\ell})=A''(\boldsymbol{\ell})$, and the bound on $A''$, applied to each of the $n$ bonds of $\boldsymbol{\ell}$, gives $|A(\boldsymbol{\ell})|\le 2n\cdot s_W\cdot\sup_p|(dA)(p)|$.
\end{proof}

Let $k\ge 1$, and let $A$ be a connection on $T_k$. For $\nu\in\{0,1,2\}$ and $\boldsymbol{\sigma}\in\{+,-\}$, where in formulas $\boldsymbol{\sigma}$ stands for the number $\pm1$, put
\[
\mathrm{Pr}^{\boldsymbol{\sigma}}_{\nu}:=\frac{1-\boldsymbol{\sigma}\cdot\gamma_{\nu}}{2},\qquad x^{\boldsymbol{\sigma}}:=x+\boldsymbol{\sigma}\cdot\eta e_{\nu};
\]
the direction $\nu$ is not displayed in the notation $x^{\boldsymbol{\sigma}}$. Thus $\mathrm{Pr}^{+}_{\nu}=\Pi^{-}_{\nu}$ and $\mathrm{Pr}^{-}_{\nu}=\Pi^{+}_{\nu}$ (Section~\ref{sec:parameters-clifford-relations}), and $\mathrm{Pr}^{+}_{\nu}+\mathrm{Pr}^{-}_{\nu}=1$. Consequently the operator $D_A$ of Section~\ref{sec:phase-factors-lattice} takes the form
\[
(D_Af)(x)=-\eta^{-1}\sum_{\nu,\boldsymbol{\sigma}}\mathrm{Pr}^{\boldsymbol{\sigma}}_{\nu}\left(u_A(x,x^{\boldsymbol{\sigma}})\cdot f(x^{\boldsymbol{\sigma}})-f(x)\right),
\]
the sum extending over $\nu\in\{0,1,2\}$ and $\boldsymbol{\sigma}\in\{+,-\}$. By Lemma~\ref{lem:block-averages-block}, $Q_k(A)\cdot P_k(A)=Q_k(A)\cdot Q_k^T(-A)\cdot Q_k(A)=Q_k(A)$. Hence the operator $T_k(A)$ of Section~\ref{sec:exact-fermion-operators} satisfies
\[
T_k(A)=b_k^{-1}\cdot Q_k(A)\left(D_A+\mathfrak{m}+b_k\cdot P_k(A)\right)=b_k^{-1}\cdot Q_k(A)\cdot D_A+\left(1+\frac{\mathfrak{m}}{b_k}\right)Q_k(A),
\]
and the formula for $Q_k(A)$ in Section~\ref{sec:phase-factors-lattice} gives, for every spinor $\psi$ on $T_k$ and every unit site $y$,
\begin{multline}
\left(T_k(A)\psi\right)(y)=-b_k^{-1}\cdot\eta^2\sum_{x\in\Delta_y}\ \sum_{\nu,\boldsymbol{\sigma}}\boldsymbol{\phi}_A(y,x)\cdot\mathrm{Pr}^{\boldsymbol{\sigma}}_{\nu}\left(u_A(x,x^{\boldsymbol{\sigma}})\cdot\psi(x^{\boldsymbol{\sigma}})-\psi(x)\right)\\
+\left(1+\frac{\mathfrak{m}}{b_k}\right)\eta^3\sum_{x\in\Delta_y}\boldsymbol{\phi}_A(y,x)\cdot\psi(x).\label{eq:operators-fermion-9}
\end{multline}
The double sum over $x\in\Delta_y$ and $(\nu,\boldsymbol{\sigma})$ is called the first sum of~\eqref{eq:operators-fermion-9}.

Let $\mathbf{Y}$ be a polymer (Section~\ref{sec:lattice-geometry}) and $y$ a unit site of $\mathbf{Y}$. For a spinor $\psi$ on $T_k$ and a function $\chi$ with values in $\mathbb{C}^4$ on the set $\mathrm{Pz}_k$ of pairs of sites of Definition~\ref{def:pairs-sites-holder} define
\begin{align*}
\left(T^{(1)}(A)\psi\right)(y):={}&-b_k^{-1}\cdot\eta^2\sum_{x\in\Delta_y}\ \sum_{(\nu,\boldsymbol{\sigma}):\,x^{\boldsymbol{\sigma}}\in\mathbf{Y}}\boldsymbol{\phi}_A(y,x)\cdot\mathrm{Pr}^{\boldsymbol{\sigma}}_{\nu}\left(u_A(x,x^{\boldsymbol{\sigma}})\cdot\psi(x^{\boldsymbol{\sigma}})-\psi(x)\right)\\
&+\left(1+\frac{\mathfrak{m}}{b_k}\right)\eta^3\sum_{x\in\Delta_y}\boldsymbol{\phi}_A(y,x)\cdot\psi(x),\\
\left(T^{(2)}(A)\chi\right)(y):={}&-b_k^{-1}\cdot\eta^{2+\alpha}\sum_{x\in\Delta_y}\ \sum_{(\nu,\boldsymbol{\sigma}):\,x^{\boldsymbol{\sigma}}\notin\mathbf{Y}}\boldsymbol{\phi}_A(y,x)\cdot\mathrm{Pr}^{\boldsymbol{\sigma}}_{\nu}\cdot\chi(x,x^{\boldsymbol{\sigma}}).
\end{align*}
Every pair $(x,x^{\boldsymbol{\sigma}})$ at which $\chi$ is evaluated lies in $\mathrm{Pz}_k$, because $x\ne x^{\boldsymbol{\sigma}}$ and $|x-x^{\boldsymbol{\sigma}}|=\eta<M$ (Section~\ref{sec:parameters-clifford-relations}). So $T^{(1)}(A)$ maps the spinors on $T_k$, and $T^{(2)}(A)$ maps the $\mathbb{C}^4$-valued functions on $\mathrm{Pz}_k$, to the $\mathbb{C}^4$-valued functions on the unit sites of $\mathbf{Y}$. Comparison with~\eqref{eq:operators-fermion-9} shows that $\left(T^{(1)}(A)\psi\right)(y)$ is what remains of its right side when the terms of the first sum with $x^{\boldsymbol{\sigma}}\notin\mathbf{Y}$ are deleted, and that $\left(T^{(1)}(A)\psi\right)(y)+\left(T^{(2)}(A)\chi\right)(y)$ is obtained from that right side by replacing, in every term with $x^{\boldsymbol{\sigma}}\notin\mathbf{Y}$, the difference $u_A(x,x^{\boldsymbol{\sigma}})\cdot\psi(x^{\boldsymbol{\sigma}})-\psi(x)$ by $\eta^{\alpha}\cdot\chi(x,x^{\boldsymbol{\sigma}})$.

If $A\in\tilde{G}_k$, then $T^{(1)}(A)$ and $T^{(2)}(A)$ are the operators of the display~\eqref{eq:renormalisation-group-58} in Section~\ref{sec:localisation-exponent}, formed with the polymer $\mathbf{Y}$ and with $A$ in place of $B$; and since $-A\in\tilde{G}_k$ by~\eqref{eq:operators-fermion-G2}, the operators $\bar{T}^{(1)}(A)$ and $\bar{T}^{(2)}(A)$ of Proposition~\ref{prop:identities-bounds}(d) are the barred operators of that display. Indeed, on $T_k$ one has $\eta=\eta_k$ (Section~\ref{sec:parameters-clifford-relations}) and $\Delta_y=B_k(y)$ (Section~\ref{sec:tori-nested-blocks}); the site $x^{\boldsymbol{\sigma}}$, the matrix $P^{\boldsymbol{\sigma}}_{\nu}$, the factor $u_A(x,x^{\boldsymbol{\sigma}})$ and the phase $\phi_{A,y}(x)$ of Section~\ref{sec:notation} are given by the same formulas as the site $x^{\boldsymbol{\sigma}}$ and the matrix $\mathrm{Pr}^{\boldsymbol{\sigma}}_{\nu}$ above and the numbers $u_A(x,x^{\boldsymbol{\sigma}})$ and $\boldsymbol{\phi}_A(y,x)$ of Section~\ref{sec:phase-factors-lattice}; and $\eta^3\cdot(b_k\cdot\eta)^{-1}=b_k^{-1}\cdot\eta^2$. Once the braces in~\eqref{eq:renormalisation-group-58} are multiplied out, its formulas therefore coincide term by term with those above.

\begin{proposition}[Identities and bounds for $T_k$, $T^{(1)}$ and $T^{(2)}$]\label{prop:identities-bounds}
Let $k\ge 1$ and $0<\mathfrak{m}\le 1$.
\begin{enumerate}
\item[(a)] Under the hypotheses of Theorem~\ref{thm:expansion-walks}, for every $A\in\tilde{G}_k$,
\[
T_k(A)\cdot H_k(A)=1,\qquad \bar{T}_k(A)\cdot\bar{H}_k(A)=1.
\]
\item[(b)] Let $A$ be a connection on $T_k$, $\mathbf{Y}$ a polymer and $y$ a unit site of $\mathbf{Y}$. If $\psi$ is a spinor on $T_k$ and $\chi=\delta_{\alpha,A}\psi$ on $\mathrm{Pz}_k$, then
\[
\left(T^{(1)}(A)\psi\right)(y)+\left(T^{(2)}(A)\chi\right)(y)=\left(T_k(A)\psi\right)(y).
\]
For all $\psi$ and $\chi$, the value $\left(T^{(1)}(A)\psi\right)(y)$ depends only on the values of $\psi$ at the sites of $\mathbf{Y}$, and $\left(T^{(2)}(A)\chi\right)(y)$ depends only on the values of $\chi$ at the pairs $(x,x')$ of nearest neighbours with $x\in\Delta_y$ and $x'\notin\mathbf{Y}$. Both values depend on $A$ only through the values of $A$ on the bonds that have an end in $\Delta_y$.
\item[(c)] Let $A$ be a connection on $T_k$ and $\mathbf{Y}$ a polymer; abbreviate $\mathfrak{i}:=\mathfrak{i}(A)$ and put $q:=e_k\cdot|dA|_{\infty}$. For the norms from the sup norm over the sites (for $T^{(2)}(A)$: over $\mathrm{Pz}_k$) to the sup norm over the unit sites of $\mathbf{Y}$,
\begin{align*}
\left\|T^{(1)}(A)\right\|,\ \left\|T_k(A)\right\|&\le\exp\left(\frac{3}{2}\mathfrak{i}\right)\left[b_{-}^{-1}\left(54L\cdot q\cdot\exp(6q)+12\exp(\mathfrak{i})+\mathfrak{m}\right)+1\right],\\
\left\|T^{(2)}(A)\right\|&\le 6b_{-}^{-1}\cdot\exp\left(\frac{3}{2}\mathfrak{i}\right)\cdot\eta^{\alpha}.
\end{align*}
\item[(d)] Let condition~\eqref{eq:renormalisation-group-23} hold and let $A\in\tilde{R}_k+\mathfrak{X}_k$. For every polymer $\mathbf{Y}$ and for $i=1,2$,
\[
\left|T^{(i)}(A)\right|_{1,\infty}\le C_T:=2\mathbf{e}^{\frac{3}{2}}\left[b_{-}^{-1}\left(54L\cdot\mathbf{e}^6+12\mathbf{e}+1\right)+1\right],
\]
where $|\cdot|_{1,\infty}$ is the norm of Section~\ref{sec:coefficients-kernels-grassmann}. The same bound holds for $T_k(A)$, for $\bar{T}^{(i)}(A):=\mathsf{C}^{-1}\cdot T^{(i)}(-A)\cdot\mathsf{C}$ $(i=1,2)$ and for $\bar{T}_k(A)$.
\end{enumerate}
At the level $k=0$ one has $T_0(A)=H_0(A)=\bar{T}_0(A)=\bar{H}_0(A)=1$ for every connection $A$ on $T_0$ (Section~\ref{sec:exact-fermion-operators}); hence the two identities of (a) hold also for $k=0$, that is, $T_0(A)\cdot H_0(A)=1$ and $\bar{T}_0(A)\cdot\bar{H}_0(A)=1$ for every such $A$.
\end{proposition}

\begin{proof}
(a) By Theorem~\ref{thm:expansion-walks}(a) with $y=\infty$, the operator $F_{\infty}(A)=D_A+\mathfrak{m}+b_k\cdot P_k(A)$ of Section~\ref{sec:operators-2} is invertible, and its inverse is $S_k(A)$. By the definitions of $T_k(A)$ and of $H_k(A)=b_k\cdot S_k(A)\cdot Q_k^T(-A)$ in Section~\ref{sec:exact-fermion-operators}, and by Lemma~\ref{lem:block-averages-block},
\[
T_k(A)\cdot H_k(A)=b_k^{-1}\cdot Q_k(A)\cdot S_k(A)^{-1}\cdot b_k\cdot S_k(A)\cdot Q_k^T(-A)=Q_k(A)\cdot Q_k^T(-A)=1.
\]
By~\eqref{eq:operators-fermion-G2}, $-A$ lies in $\tilde{G}_k$ as well, so the same identity holds at $-A$. With the definitions of $\bar{T}_k$ and $\bar{H}_k$ in Section~\ref{sec:exact-fermion-operators} one obtains
\[
\bar{T}_k(A)\cdot\bar{H}_k(A)=\mathsf{C}^{-1}\cdot T_k(-A)\cdot\mathsf{C}\cdot\mathsf{C}^{-1}\cdot H_k(-A)\cdot\mathsf{C}=\mathsf{C}^{-1}\cdot T_k(-A)\cdot H_k(-A)\cdot\mathsf{C}=1.
\]

(b) If $x$ and $x'$ are nearest neighbours, each of the six paths $\Gamma^{\pi}(x,x')$ of Section~\ref{sec:phase-factors-lattice} is the bond $(x,x')$, so that $(\tau A)(x,x')=A(x,x')$ and $U_A(x,x')=u_A(x,x')$; moreover $|x-x'|=\eta$. The definition of $\delta_{\alpha,A}$ in Section~\ref{sec:phase-factors-lattice} therefore gives
\[
\eta^{\alpha}\left(\delta_{\alpha,A}\psi\right)(x,x^{\boldsymbol{\sigma}})=u_A(x,x^{\boldsymbol{\sigma}})\cdot\psi(x^{\boldsymbol{\sigma}})-\psi(x).
\]
Hence, for $\chi=\delta_{\alpha,A}\psi$, the value $\left(T^{(2)}(A)\chi\right)(y)$ is $-b_k^{-1}\cdot\eta^2$ times the sum of those terms of the first sum of~\eqref{eq:operators-fermion-9} for which $x^{\boldsymbol{\sigma}}\notin\mathbf{Y}$, and $\left(T^{(1)}(A)\psi\right)(y)$ is the remaining part of the right side of~\eqref{eq:operators-fermion-9}; the identity follows. In $\left(T^{(1)}(A)\psi\right)(y)$ the spinor $\psi$ is evaluated at sites $x\in\Delta_y$ and at sites $x^{\boldsymbol{\sigma}}\in\mathbf{Y}$, and $\Delta_y\subset\mathbf{Y}$, because the unit block $\Delta_y$ lies in the $M$-cube that contains $y$. In $\left(T^{(2)}(A)\chi\right)(y)$ the function $\chi$ is evaluated at the pairs $(x,x^{\boldsymbol{\sigma}})$ of nearest neighbours with $x\in\Delta_y$ and $x^{\boldsymbol{\sigma}}\notin\mathbf{Y}$. The connection enters both values through the numbers $\boldsymbol{\phi}_A(y,x)$, $x\in\Delta_y$, and $u_A(x,x^{\boldsymbol{\sigma}})$. The number $\boldsymbol{\phi}_A(y,x)$ depends only on the values of $A$ on the bonds inside $\Delta_y$: the paths that define $(\tau_kA)(y,x)$ in Section~\ref{sec:phase-factors-lattice} join centres of nested blocks that contain $x$, all contained in $\Delta_y$, and such a path does not leave the larger of the two blocks whose centres it joins, since along a rectilinear path each coordinate varies monotonically between its values at the end points. The number $u_A(x,x^{\boldsymbol{\sigma}})$ depends only on the value of $A$ on the bond $(x,x^{\boldsymbol{\sigma}})$, which has the end $x\in\Delta_y$.

(c) Fix a unit site $y$ of $\mathbf{Y}$ and a pair $(\nu,\boldsymbol{\sigma})$; put $E:=\Delta_y$. By Lemma~\ref{lem:bounds-phases}(a) and $\eta\le 1$,
\[
\left|\boldsymbol{\phi}_A(y,x)\right|\le\exp\left(\frac{3}{2}\mathfrak{i}\right),\qquad \left|u_A(x,x')\right|\le\exp(\mathfrak{i}).
\]
With $x^{-\boldsymbol{\sigma}}:=x-\boldsymbol{\sigma}\cdot\eta e_{\nu}$ let
\begin{align*}
E_1&:=\left\{x\in E:\ x^{\boldsymbol{\sigma}}\in E\right\},& E_2&:=\left\{x\in E:\ x^{\boldsymbol{\sigma}}\in\mathbf{Y}\setminus E\right\},\\
E_3&:=\left\{x\in E:\ x^{\boldsymbol{\sigma}}\notin\mathbf{Y}\right\},& E_4&:=\left\{x\in E:\ x^{-\boldsymbol{\sigma}}\notin E\right\}.
\end{align*}
The block $E$ is a cube of $\eta^{-3}$ sites, and $E\subset\mathbf{Y}$ (see the proof of (b)). Hence $E_2\cup E_3$ is the set of the $x\in E$ with $x^{\boldsymbol{\sigma}}\notin E$, that is, the face of $E$ on the side of $\boldsymbol{\sigma}\cdot e_{\nu}$, and $E_4$ is the opposite face; so $\#|E_2\cup E_3|=\#|E_4|=\eta^{-2}$, and $x\mapsto x^{\boldsymbol{\sigma}}$ is a bijection of $E_1$ onto $E\setminus E_4$. The terms of the first sum in the definition of $T^{(1)}(A)$ that belong to $(\nu,\boldsymbol{\sigma})$ are those with $x\in E_1\cup E_2$. Now
\[
-\sum_{x\in E_1\cup E_2}\boldsymbol{\phi}_A(y,x)\cdot\psi(x)=-\sum_{x\in E}\boldsymbol{\phi}_A(y,x)\cdot\psi(x)+\sum_{x\in E_3}\boldsymbol{\phi}_A(y,x)\cdot\psi(x),
\]
and, by the bijection just described,
\[
-\sum_{x\in E}\boldsymbol{\phi}_A(y,x)\cdot\psi(x)=-\sum_{x\in E_1}\boldsymbol{\phi}_A(y,x^{\boldsymbol{\sigma}})\cdot\psi(x^{\boldsymbol{\sigma}})-\sum_{x'\in E_4}\boldsymbol{\phi}_A(y,x')\cdot\psi(x').
\]
It follows that
\begin{align*}
\sum_{x\in E_1\cup E_2}\boldsymbol{\phi}_A(y,x)\left(u_A(x,x^{\boldsymbol{\sigma}})\cdot\psi(x^{\boldsymbol{\sigma}})-\psi(x)\right)={}&\sum_{x\in E_1}\left(\boldsymbol{\phi}_A(y,x)\cdot u_A(x,x^{\boldsymbol{\sigma}})-\boldsymbol{\phi}_A(y,x^{\boldsymbol{\sigma}})\right)\psi(x^{\boldsymbol{\sigma}})\\
&+\sum_{x\in E_2}\boldsymbol{\phi}_A(y,x)\cdot u_A(x,x^{\boldsymbol{\sigma}})\cdot\psi(x^{\boldsymbol{\sigma}})-\sum_{x'\in E_4}\boldsymbol{\phi}_A(y,x')\cdot\psi(x')\\
&+\sum_{x\in E_3}\boldsymbol{\phi}_A(y,x)\cdot\psi(x).
\end{align*}
The last three sums on the right have together $\#|E_2\cup E_3|+\#|E_4|=2\eta^{-2}$ terms, each of modulus at most $\exp\left(\frac{3}{2}\mathfrak{i}\right)\cdot\exp(\mathfrak{i})\cdot|\psi|_{\infty}$ (recall $\mathfrak{i}\ge 0$). For $T_k(A)$ the first sum of~\eqref{eq:operators-fermion-9} extends over all $x\in E$; the same computation then holds with $E_2\cup E_3$ in place of $E_2$ and without the sum over $E_3$, and the same bound applies to the remaining sums.

It remains to bound the first sum on the right. For $x\in E_1$ the definitions of $u_A$ and $\boldsymbol{\phi}_A$ give
\begin{gather*}
\boldsymbol{\phi}_A(y,x)\cdot u_A(x,x^{\boldsymbol{\sigma}})-\boldsymbol{\phi}_A(y,x^{\boldsymbol{\sigma}})=\boldsymbol{\phi}_A(y,x^{\boldsymbol{\sigma}})\left(\exp(z)-1\right),\\
z:=i\cdot e_k\cdot\eta\left[(\tau_kA)(y,x)+A(x,x^{\boldsymbol{\sigma}})-(\tau_kA)(y,x^{\boldsymbol{\sigma}})\right].
\end{gather*}
For $0\le i\le k$ denote by $x_i$ the centre of the block of side $L^i\eta$ to which $x$ belongs, and by $(x^{\boldsymbol{\sigma}})_i$ the centre of the block of the same side to which $x^{\boldsymbol{\sigma}}$ belongs; thus $x_0=x$, $(x^{\boldsymbol{\sigma}})_0=x^{\boldsymbol{\sigma}}$ and $x_k=(x^{\boldsymbol{\sigma}})_k=y$. Let $j=j(x)$ be the least index in $\{1,\dots,k\}$ with $x_j=(x^{\boldsymbol{\sigma}})_j$. Since the blocks are nested, $x_i=(x^{\boldsymbol{\sigma}})_i$ for all $i\ge j$, and the terms of index $i\ge j$ of
\[
(\tau_kA)(y,x)=\sum_{i<k}(\tau A)(x_{i+1},x_i)
\]
and of the corresponding sum for $x^{\boldsymbol{\sigma}}$ cancel in the bracket. Each of the remaining terms is the mean of the values of $A$ on the six paths $\Gamma^{\pi}$, and a sum of means over independent choices is the mean of the sums over all joint choices. The bracket is therefore the average of $A(\boldsymbol{\ell})$ over all choices of the orders $\pi_0,\dots,\pi_{j-1}$ and $\pi'_0,\dots,\pi'_{j-1}$, where $\boldsymbol{\ell}$ is the closed lattice path
\begin{align*}
\boldsymbol{\ell}:={}&\Gamma^{\pi_{j-1}}(x_j,x_{j-1})+\dots+\Gamma^{\pi_0}(x_1,x)+(x,x^{\boldsymbol{\sigma}})\\
&-\left[\Gamma^{\pi'_{j-1}}\left(x_j,(x^{\boldsymbol{\sigma}})_{j-1}\right)+\dots+\Gamma^{\pi'_0}\left((x^{\boldsymbol{\sigma}})_1,x^{\boldsymbol{\sigma}}\right)\right],
\end{align*}
where $+$ denotes the concatenation of paths and $-[\,\cdot\,]$ the reversed path. Each such path $\boldsymbol{\ell}$ lies in the block of side $L^j\eta$ with the centre $x_j$. Indeed, this block contains $x$ and $x^{\boldsymbol{\sigma}}$, hence the blocks of smaller side that contain one of them, and so all the points $x_i$ and $(x^{\boldsymbol{\sigma}})_i$, $0\le i\le j$; and a block that contains two sites contains every rectilinear path between them, as each coordinate varies monotonically along such a path. Since $|x_{i+1}-x_i|\le\frac{\left(L^{i+1}-L^i\right)\eta}{2}$, the path $\Gamma^{\pi_i}(x_{i+1},x_i)$ has at most $\frac{3}{2}\left(L^{i+1}-L^i\right)$ bonds, and the same holds for the paths that belong to $x^{\boldsymbol{\sigma}}$; hence $\boldsymbol{\ell}$ has at most
\[
2\cdot\frac{3}{2}\left(L^j-1\right)+1\le 3L^j
\]
bonds. Lifted to $\eta\mathbb{Z}^3$, the block of side $L^j\eta$ is a product of three intervals of $\eta\mathbb{Z}$ of length at most $L^j\eta$. Lemma~\ref{lem:gauge-fixing-along} with $s_W=L^j\eta$ yields
\[
\left|A(\boldsymbol{\ell})\right|\le 2\cdot 3L^j\cdot L^j\cdot\eta\cdot|dA|_{\infty}=6L^{2j}\cdot\eta\cdot|dA|_{\infty}.
\]
As $\eta=L^{-k}$, it follows that
\[
|z|\le 6L^{2(j-k)}\cdot q\le 6q,\qquad \left|\exp(z)-1\right|\le|z|\cdot\exp|z|\le 6L^{2(j-k)}\cdot q\cdot\exp(6q).
\]
Next, $\#|\{x\in E_1:\ j(x)=j\}|\le\eta^{-3}\cdot L^{1-j}$. For $j=1$ this is the bound $\#|E|=\eta^{-3}$. Let $j\ge 2$ and $j(x)=j$. As $j$ is the least index with $x_j=(x^{\boldsymbol{\sigma}})_j$, one has $x_{j-1}\ne(x^{\boldsymbol{\sigma}})_{j-1}$: the step from $x$ to $x^{\boldsymbol{\sigma}}$ leaves the block of side $L^{j-1}\eta$ with the centre $x_{j-1}$. This step alters only the coordinate in the direction $\nu$, by $\boldsymbol{\sigma}\cdot\eta$; so, among the sites of that block, $x$ has the largest coordinate in the direction $\nu$ if $\boldsymbol{\sigma}=+$, and the smallest one if $\boldsymbol{\sigma}=-$. In a block of side $L^{j-1}\eta$ this coordinate takes $L^{j-1}$ values, each of them at the same number of sites, and $E$ is a disjoint union of such blocks; consequently at most the fraction $L^{1-j}$ of the $\eta^{-3}$ sites of $E$ satisfy $j(x)=j$, which is the bound. With $\eta^2=L^{-2k}$ and $\eta^{-3}=L^{3k}$ one obtains
\begin{align*}
\eta^2\sum_{x\in E_1}\left|\boldsymbol{\phi}_A(y,x^{\boldsymbol{\sigma}})\right|\cdot\left|\exp(z)-1\right|&\le 6q\cdot\exp(6q)\cdot\exp\left(\frac{3}{2}\mathfrak{i}\right)\sum_{j=1}^kL^{-2k}\cdot L^{3k-j+1}\cdot L^{2j-2k}\\
&=6q\cdot\exp(6q)\cdot\exp\left(\frac{3}{2}\mathfrak{i}\right)\cdot L^{1-k}\sum_{j=1}^kL^j\\
&\le 9L\cdot q\cdot\exp(6q)\cdot\exp\left(\frac{3}{2}\mathfrak{i}\right),
\end{align*}
because $L^{1-k}\sum_{j=1}^kL^j\le\frac{L^2}{L-1}\le\frac{3L}{2}$ for $L\ge 5$. Consequently the modulus of the first sum on the right of the identity above is at most $\eta^{-2}\cdot 9L\cdot q\cdot\exp(6q)\cdot\exp\left(\frac{3}{2}\mathfrak{i}\right)\cdot|\psi|_{\infty}$.

The matrix $\mathrm{Pr}^{\boldsymbol{\sigma}}_{\nu}$ is a non-zero orthogonal projection (Section~\ref{sec:parameters-clifford-relations}), so that $|\mathrm{Pr}^{\boldsymbol{\sigma}}_{\nu}|=1$; moreover $b_k\ge b_{-}$ (Section~\ref{sec:parameters-clifford-relations}), and $\left|\left(Q_k(A)\psi\right)(y)\right|\le\exp\left(\frac{3}{2}\mathfrak{i}\right)\cdot|\psi|_{\infty}$, because $\eta^3\cdot\#|\Delta_y|=1$. Multiplying the bounds for the four sums by $b_k^{-1}\cdot\eta^2$, adding them over the six pairs $(\nu,\boldsymbol{\sigma})$, and adding the term with $Q_k(A)$, one finds
\[
\left|\left(T^{(1)}(A)\psi\right)(y)\right|\le\left[6b_{-}^{-1}\left(9L\cdot q\cdot\exp(6q)+2\exp(\mathfrak{i})\right)+1+\frac{\mathfrak{m}}{b_{-}}\right]\cdot\exp\left(\frac{3}{2}\mathfrak{i}\right)\cdot|\psi|_{\infty},
\]
which is the asserted bound for $T^{(1)}(A)$; the same estimate holds for $\left|\left(T_k(A)\psi\right)(y)\right|$, and the bound for $T_k(A)$ follows. Finally, for each pair $(\nu,\boldsymbol{\sigma})$ the sites $x\in\Delta_y$ with $x^{\boldsymbol{\sigma}}\notin\mathbf{Y}$ form the set $E_3$, which has at most $\eta^{-2}$ elements. So $\left(T^{(2)}(A)\chi\right)(y)$ is a sum of at most $6\eta^{-2}$ terms, each of modulus at most $b_{-}^{-1}\cdot\eta^{2+\alpha}\cdot\exp\left(\frac{3}{2}\mathfrak{i}\right)\cdot|\chi|_{\infty}$, and the bound for $T^{(2)}(A)$ follows.

(d) Let $A\in\tilde{R}_k+\mathfrak{X}_k$. Under condition~\eqref{eq:renormalisation-group-23}, \eqref{eq:operators-fermion-G3} gives $\mathfrak{i}\le 1$ and $q\le 1$; moreover $\mathfrak{m}\le 1$ and $\eta^{\alpha}\le 1$. By (c),
\begin{align*}
\left\|T^{(1)}(A)\right\|,\ \left\|T_k(A)\right\|&\le\mathbf{e}^{\frac{3}{2}}\left[b_{-}^{-1}\left(54L\cdot\mathbf{e}^6+12\mathbf{e}+1\right)+1\right]=\frac{C_T}{2},\\
\left\|T^{(2)}(A)\right\|&\le 6b_{-}^{-1}\cdot\mathbf{e}^{\frac{3}{2}}\le\frac{C_T}{2}.
\end{align*}
Every unit site lies in an $M$-cube, and an $M$-cube is a polymer (Section~\ref{sec:lattice-geometry}); since (c) holds for every polymer, the bound for $\left\|T_k(A)\right\|$ holds also for the norm into the sup norm over all unit sites. For an operator $O$, Section~\ref{sec:coefficients-kernels-grassmann} identifies $|O|_{1,\infty}$ with the least constant $C$ such that $\sup|Of|_c\le C\cdot\sup|f|_c$ for all $f$, where $|\cdot|_c$ is the largest modulus of a component; by Section~\ref{sec:norms-pairings-kernels} this constant is at most $2\|O\|$. This proves the bounds for $T^{(1)}(A)$, $T^{(2)}(A)$ and $T_k(A)$. The conditions that define $\tilde{R}_k$ and $\mathfrak{X}_k$ in Definition~\ref{def:domains} are invariant under $A\mapsto -A$; hence $-A\in\tilde{R}_k+\mathfrak{X}_k$, and the same bounds hold for $T^{(1)}(-A)$, $T^{(2)}(-A)$ and $T_k(-A)$. The matrix $\mathsf{C}=i\gamma_1$ (Section~\ref{sec:dirac-matrices-lattice}) is a monomial matrix whose non-zero entries are $\pm1$, and so is $\mathsf{C}^{-1}=-\mathsf{C}$; conjugation with $\mathsf{C}$ therefore leaves $|\cdot|_{1,\infty}$ unchanged. This gives the bounds for $\bar{T}^{(1)}(A)$, $\bar{T}^{(2)}(A)$ and $\bar{T}_k(A)=\mathsf{C}^{-1}\cdot T_k(-A)\cdot\mathsf{C}$ (Section~\ref{sec:exact-fermion-operators}).

As the operators of the display~\eqref{eq:renormalisation-group-58} are $T^{(1)}(A)$ and $T^{(2)}(A)$ for $A\in\tilde{G}_k$, Proposition~\ref{prop:left-inverse} follows from (a), (b) and (d), and, at the level $k=0$, from $T_0(A)=H_0(A)=\bar{T}_0(A)=\bar{H}_0(A)=1$ (Section~\ref{sec:exact-fermion-operators}).
\end{proof}

\subsection{Covariances and periodisation}\label{sec:covariances-periodisation}

\begin{proposition}[Symmetries of the fermion operators]\label{prop:symmetries-fermion-operators-2}
Assume the hypotheses of Theorem~\ref{thm:expansion-walks}. Let $k\ge1$, $A\in\tilde{G}_k$ and $y\in Y_k=\mathbb{R}\cup\{\infty\}$, where $y\in\mathbb{R}$ in every relation that contains $B_{k,y}$ or $\Gamma_{k,y}$.
\begin{enumerate}
\item[(a)] \emph{Gauge transformations.} Let $\boldsymbol{\omega}$ be a real connection on $T_k$ with $d\boldsymbol{\omega}=0$, and let $u$ be a function on the sites of $T_k$ such that $|u|=1$ and
\[
u(x')=u(x)\cdot\exp\left\{ie_k\eta\cdot\boldsymbol{\omega}(x,x')\right\}
\]
for every bond $(x,x')$; let $u^0$ be the restriction of $u$ to the unit lattice $T^0_k$. Then $A-\boldsymbol{\omega}\in\tilde{G}_k$, the relations of Lemma~\ref{lem:symmetries-elementary-operators}(a) hold, and
\begin{align*}
S_y(A-\boldsymbol{\omega})&=u\cdot S_y(A)\cdot u^{-1},\\
H_k(A-\boldsymbol{\omega})&=u\cdot H_k(A)\cdot\left(u^0\right)^{-1},\\
T_k(A-\boldsymbol{\omega})&=u^0\cdot T_k(A)\cdot u^{-1},\\
O(A-\boldsymbol{\omega})&=u^0\cdot O(A)\cdot\left(u^0\right)^{-1}\qquad\text{for }O=D_k,\ M_k,\ \dot{M}_k,\ B_{k,y},\ \Gamma_{k,y},\\
\delta Z_k(A-\boldsymbol{\omega})&=\delta Z_k(A).
\end{align*}
The hypotheses on $\boldsymbol{\omega}$ and $u$ are met, in particular, by $\boldsymbol{\omega}=d\boldsymbol{\lambda}$ and $u=\exp\left(ie_k\boldsymbol{\lambda}\right)$, where $\boldsymbol{\lambda}$ is a real scalar, and by a real constant connection $\boldsymbol{\omega}$, that is, $\boldsymbol{\omega}(x,x+\eta e_{\mu})=\boldsymbol{\omega}_{\mu}$ for all sites $x$ and $\mu=0,1,2$, such that $e_k\cdot l\cdot\boldsymbol{\omega}_{\mu}\in2\pi\mathbb{Z}$ for $\mu=0,1,2$, together with $u(x)=\exp\left(ie_k\cdot(\boldsymbol{\omega}\cdot x)\right)$, where $\boldsymbol{\omega}\cdot x:=\sum_{\mu}\boldsymbol{\omega}_{\mu}x_{\mu}$ for $x\in\eta\mathbb{Z}^3$.
\item[(b)] \emph{Lattice symmetries.} Let $g=(S,a')\in\mathbb{G}_k$, acting on connections and on spinors as in Section~\ref{sec:dirac-matrices-lattice}. Then
\[
O(gA)\cdot g=g\cdot O(A)
\]
for $O(A)=D_A$, $Q_k(A)$, $Q_k^T(-A)$, $Q(A)$, $Q^T(-A)$, $P_k(A)$, $P(A)$, $P_{k+1}(A)$, $S_k(A)$, $S^0_{k+1}(A)$, $H_k(A)$, $T_k(A)$, $D_k(A)$, $M_k(A)$, $\dot{M}_k(A)$, $B_k(A)$, $\Gamma_k(A)$. Put $y'=y$ if $S\in\mathrm{Sp}_3$ or $y\in\{0,\infty\}$, and $y'=-y$ otherwise. Then
\[
O_{y'}(gA)\cdot g=g\cdot O_y(A)\qquad\text{for }O_y=S_{k,y},\ B_{k,y},\ \Gamma_{k,y},
\]
and $\delta Z_k(gA)=\delta Z_k(A)$.
\item[(c)] \emph{Transposition.} One has
\[
O(A)^T=\mathsf{C}^{-1}\cdot O(-A)\cdot\mathsf{C}
\]
for $O(A)=D_A$, $P_k(A)$, $P_{k+1}(A)$, $S_k(A)$, $S^0_{k+1}(A)$, $P(A)$, $D_k(A)$, $M_k(A)$, $\dot{M}_k(A)$, $B_k(A)$, $\Gamma_k(A)$, and, for real $y$,
\[
O_y(A)^T=\mathsf{C}^{-1}\cdot O_{-y}(-A)\cdot\mathsf{C}\qquad\text{for }O_y=S_{k,y},\ B_{k,y},\ \Gamma_{k,y}.
\]
With $\mathfrak{H}_k(A)=b_k\cdot Q_k(A)\cdot S_k(A)$ (Section~\ref{sec:operators-unit-lattice}) one has $H_k(A)^T=\mathsf{C}^{-1}\cdot\mathfrak{H}_k(-A)\cdot\mathsf{C}$, and the operators $\bar{H}_k(A)=\mathsf{C}^{-1}\cdot H_k(-A)\cdot\mathsf{C}$ and $\bar{T}_k(A)=\mathsf{C}^{-1}\cdot T_k(-A)\cdot\mathsf{C}$ of Section~\ref{sec:exact-fermion-operators} are given by
\begin{align*}
\bar{H}_k(A)&=b_k\cdot S_k(A)^T\cdot Q_k^T(A),\\
\bar{T}_k(A)&=b_k^{-1}\cdot Q_k(-A)\cdot\left(D_A+\mathfrak{m}+b_k\cdot P_k(A)\right)^T.
\end{align*}
\item[(d)] \emph{Complex conjugation.} For an operator $O$ let $\overline{O}$ denote the operator $f\mapsto\overline{O\,\overline{f}}$. Then
\[
\overline{O\left(\overline{A}\right)}=\varkappa^{-1}\cdot O(-A)\cdot\varkappa
\]
for all operators $O(A)$ of the first list in (b), and $\overline{O_y\left(\overline{A}\right)}=\varkappa^{-1}\cdot O_y(-A)\cdot\varkappa$ for $O_y=S_{k,y}$, $B_{k,y}$, $\Gamma_{k,y}$, with the same $y$ on both sides.
\item[(e)] \emph{The determinant.} $\delta Z_k(-A)=\delta Z_k(A)$ and $\overline{\delta Z_k\left(\overline{A}\right)}=\delta Z_k(A)$.
\end{enumerate}
For $k=0$ and $A\in\tilde{G}_0$ the assertions (a)--(e) hold for the operators $D_A$, $Q(A)$, $P(A)$, $S_{0,y}(A)$ with $y\in\mathbb{R}$, and $\Gamma_0(A)$, and for $\delta Z_0(A):=\det F_0(A)$; at this level the unit lattice is $T_0$ itself, so that $u^0=u$.
\end{proposition}

\begin{proof}
Every operator that occurs in the statement is built, by sums, products and inverses, from the operators of Lemma~\ref{lem:symmetries-elementary-operators}, from real numbers and from functions of the matrix $i\gamma_3y$ (Section~\ref{sec:exact-fermion-operators} and the definitions at the head of Section~\ref{sec:operators-unit-lattice}). The inverses that are needed exist at each of the connections $A$, $A-\boldsymbol{\omega}$, $gA$, $-A$ and $\overline{A}$, by Theorem~\ref{thm:expansion-walks}(a) and Proposition~\ref{prop:operators-fermion-1}(b), (d), because these connections lie in $\tilde{G}_k$: for $gA$, $-A$ and $\overline{A}$ by~\eqref{eq:operators-fermion-G2}, and for $A-\boldsymbol{\omega}$ by the first assertion of (a), which is established at the beginning of the proof of (a).

\emph{(a)} Let $A=A_0+A_1$ be as in the definition of $\tilde{G}_k$ (Section~\ref{sec:domain-complex-connections}): $A_1\in G_k$, $A_0$ is real, and for every $M$-cube $\square$ there are real functions $\boldsymbol{\lambda}$ and $A'$ with $A_0=A'+d\boldsymbol{\lambda}$ and $|A'|_{\infty}<e_k^{-1+\varepsilon}$ on $\square^{\sim 5}$. Fix $\square$. The cube $\square^{\sim 5}$, of side $11M$, does not wrap around $T_k$ (Section~\ref{sec:domain-complex-connections}), and it lifts to a product of three intervals of $\eta\mathbb{Z}$, each of length at most $11M$. Choose a site $x_0$ of this product and put $\nu(x):=\eta\cdot\boldsymbol{\omega}\left(\Gamma(x_0,x)\right)$, with the path $\Gamma(x_0,x)$ of Lemma~\ref{lem:gauge-fixing-along}. By that lemma, applied to $\boldsymbol{\omega}$ with $s_W=11M$, the function $\nu$ is real and
\[
\left|\boldsymbol{\omega}(b)-(d\nu)(b)\right|\le2s_W\cdot\sup_p\left|(d\boldsymbol{\omega})(p)\right|=0
\]
for every bond $b$ of the product. Hence $\boldsymbol{\omega}=d\nu$ on $\square^{\sim 5}$, and
\[
A_0-\boldsymbol{\omega}=A'+d(\boldsymbol{\lambda}-\nu)\qquad\text{on }\square^{\sim 5},
\]
with $A'$ and $\boldsymbol{\lambda}-\nu$ real. As $A_0-\boldsymbol{\omega}$ is real and $A_1\in G_k$, it follows that $A-\boldsymbol{\omega}=(A_0-\boldsymbol{\omega})+A_1\in\tilde{G}_k$.

In the two particular cases the hypotheses are verified as follows. If $\boldsymbol{\omega}=d\boldsymbol{\lambda}$, then $d\boldsymbol{\omega}=0$, and $u=\exp\left(ie_k\boldsymbol{\lambda}\right)$ satisfies
\[
u(x')=u(x)\cdot\exp\left\{ie_k\left(\boldsymbol{\lambda}(x')-\boldsymbol{\lambda}(x)\right)\right\}=u(x)\cdot\exp\left\{ie_k\eta\cdot\boldsymbol{\omega}(x,x')\right\}.
\]
If $\boldsymbol{\omega}$ is a real constant connection with $e_k\cdot l\cdot\boldsymbol{\omega}_{\mu}\in2\pi\mathbb{Z}$, then $d\boldsymbol{\omega}=0$; the function $x\mapsto\exp\left(ie_k\cdot(\boldsymbol{\omega}\cdot x)\right)$ on $\eta\mathbb{Z}^3$ satisfies $u(x+le_{\mu})=u(x)$, so that it is a function on the sites of $T_k$, and $u(x+\eta e_{\mu})=u(x)\cdot\exp\left(ie_k\eta\cdot\boldsymbol{\omega}_{\mu}\right)$; in this case $A-\boldsymbol{\omega}\in\tilde{G}_k$ follows also from~\eqref{eq:operators-fermion-G2}, since $-\boldsymbol{\omega}$ is a real constant connection.

Lemma~\ref{lem:symmetries-elementary-operators}(a) gives
\begin{gather*}
F_y(A-\boldsymbol{\omega})=u\cdot F_y(A)\cdot u^{-1},\qquad P(A-\boldsymbol{\omega})=u^0\cdot P(A)\cdot\left(u^0\right)^{-1},\\
Q_k(A-\boldsymbol{\omega})=u^0\cdot Q_k(A)\cdot u^{-1},\qquad Q_k^T\left(-(A-\boldsymbol{\omega})\right)=u\cdot Q_k^T(-A)\cdot\left(u^0\right)^{-1}.
\end{gather*}
Inverting the first relation gives $S_y(A-\boldsymbol{\omega})=u\cdot S_y(A)\cdot u^{-1}$, in particular $S_k(A-\boldsymbol{\omega})=u\cdot S_k(A)\cdot u^{-1}$ for $y=\infty$. Recall the definitions $H_k(A)=b_k\cdot S_k(A)\cdot Q_k^T(-A)$, $M_k(A)=b_k^2\cdot Q_k(A)\cdot S_k(A)\cdot Q_k^T(-A)$, $D_k(A)=b_k-M_k(A)$ and $\dot{M}_k(A)=b_k^2\cdot Q_k(A)\cdot S_k(A)^2\cdot Q_k^T(-A)$ (Section~\ref{sec:operators-unit-lattice}), and $T_k(A)=b_k^{-1}\cdot Q_k(A)\cdot F_{\infty}(A)$ with $F_{\infty}(A)=D_A+\mathfrak{m}+b_k\cdot P_k(A)$ (Section~\ref{sec:exact-fermion-operators} and~\eqref{eq:operators-fermion-2}, as $\boldsymbol{\alpha}_{\infty}=b_k$ and $\boldsymbol{\beta}_{\infty}=0$). In the products the factors $u^{-1}\cdot u$ cancel:
\begin{align*}
H_k(A-\boldsymbol{\omega})&=b_k\cdot u\cdot S_k(A)\cdot u^{-1}\cdot u\cdot Q_k^T(-A)\cdot\left(u^0\right)^{-1}=u\cdot H_k(A)\cdot\left(u^0\right)^{-1},\\
M_k(A-\boldsymbol{\omega})&=b_k^2\cdot u^0\cdot Q_k(A)\cdot u^{-1}\cdot u\cdot S_k(A)\cdot u^{-1}\cdot u\cdot Q_k^T(-A)\cdot\left(u^0\right)^{-1}\\
&=u^0\cdot M_k(A)\cdot\left(u^0\right)^{-1},\\
\dot{M}_k(A-\boldsymbol{\omega})&=b_k^2\cdot u^0\cdot Q_k(A)\cdot u^{-1}\cdot\left(u\cdot S_k(A)\cdot u^{-1}\right)^2\cdot u\cdot Q_k^T(-A)\cdot\left(u^0\right)^{-1}\\
&=u^0\cdot\dot{M}_k(A)\cdot\left(u^0\right)^{-1},\\
T_k(A-\boldsymbol{\omega})&=b_k^{-1}\cdot u^0\cdot Q_k(A)\cdot u^{-1}\cdot u\cdot F_{\infty}(A)\cdot u^{-1}=u^0\cdot T_k(A)\cdot u^{-1},
\end{align*}
and $D_k(A-\boldsymbol{\omega})=b_k-u^0\cdot M_k(A)\cdot\left(u^0\right)^{-1}=u^0\cdot D_k(A)\cdot\left(u^0\right)^{-1}$. The constant matrices $\left(b_k+i\gamma_3y\right)^{-1}$ and $\left(b_k+b_L+i\gamma_3y\right)^{-1}$ commute with the multiplication by the scalar function $u^0$; together with the relation for $P$, the definition
\[
B_{k,y}(A)=\left(b_k+i\gamma_3y\right)^{-1}\cdot\left(1-P(A)\right)+\left(b_k+b_L+i\gamma_3y\right)^{-1}\cdot P(A)
\]
yields $B_{k,y}(A-\boldsymbol{\omega})=u^0\cdot B_{k,y}(A)\cdot\left(u^0\right)^{-1}$. Further,
\[
D_k(A-\boldsymbol{\omega})+b_L\cdot P(A-\boldsymbol{\omega})+i\gamma_3y=u^0\cdot\left(D_k(A)+b_L\cdot P(A)+i\gamma_3y\right)\cdot\left(u^0\right)^{-1};
\]
by Proposition~\ref{prop:operators-fermion-1}(b), applied at $A-\boldsymbol{\omega}$ and at $A$, both sides are invertible, with the inverses $\Gamma_{k,y}(A-\boldsymbol{\omega})$ and $u^0\cdot\Gamma_{k,y}(A)\cdot\left(u^0\right)^{-1}$, which are therefore equal. Finally, a determinant is invariant under conjugation by an invertible operator, so that
\[
\delta Z_k(A-\boldsymbol{\omega})=\det\left(u^0\cdot\left(D_k(A)+b_L\cdot P(A)\right)\cdot\left(u^0\right)^{-1}\right)=\delta Z_k(A).
\]

\emph{(b)} The element $g$ induces the morphism $x\mapsto r(S)x+a'$ of $T_k$ onto itself (Section~\ref{sec:tori-nested-blocks}), and its actions on connections and on spinors are those of Lemma~\ref{lem:symmetries-elementary-operators}(b) for this morphism. Lemma~\ref{lem:symmetries-elementary-operators}(b) gives $O(gA)\cdot g=g\cdot O(A)$ for $O(A)=D_A$, $Q_k(A)$, $Q_k^T(-A)$, $Q(A)$, $Q^T(-A)$, $P_k(A)$, $P(A)$, $P_{k+1}(A)$, and $F_{y'}(gA)\cdot g=g\cdot F_y(A)$; inverting the latter,
\[
S_{y'}(gA)\cdot g=g\cdot S_y(A).
\]
For $y=\infty$ and $y=0$, where $y'=y$, these are the relations for $S_k$ and $S^0_{k+1}$, and for real $y$ the relation for $S_{k,y}$. If $O_1(gA)\cdot g=g\cdot O_1(A)$ and $O_2(gA)\cdot g=g\cdot O_2(A)$, then
\[
O_1(gA)\cdot O_2(gA)\cdot g=O_1(gA)\cdot g\cdot O_2(A)=g\cdot O_1(A)\cdot O_2(A),
\]
and the relation holds for $O_1+O_2$ and for the real multiples of $O_1$. Since $H_k$, $M_k$, $D_k$, $\dot{M}_k$ and $T_k$ are formed from $S_k$, $Q_k$, $Q_k^T(-\,\cdot\,)$, $F_{\infty}$ and real numbers by the sums and products recalled in the proof of (a), the relations for these five operators follow. For a constant $4\times4$ matrix $m$ one has $m\cdot g=g\cdot\left(S^{-1}\cdot m\cdot S\right)$, because $(gf)(x)=S\cdot f\left(g^{-1}x\right)$. By~\eqref{eq:operators-fermion-1} each of the generators $\gamma_0,\dots,\gamma_3$ of $\mathrm{Sp}$ commutes or anticommutes with $\gamma_3$, and $\mathrm{Sp}_3$ consists of the elements of $\mathrm{Sp}$ that commute with $\gamma_3$ (Section~\ref{sec:dirac-matrices-lattice}); hence $S^{-1}\gamma_3S=\pm\gamma_3$, with the sign $+$ exactly for $S\in\mathrm{Sp}_3$. Therefore $S^{-1}\cdot\left(c+i\gamma_3y'\right)\cdot S=c+i\gamma_3y$, that is,
\[
\left(c+i\gamma_3y'\right)^{-1}\cdot S=S\cdot\left(c+i\gamma_3y\right)^{-1}\qquad\text{for real }y\text{ and }c\in\left\{b_k,\,b_k+b_L\right\}.
\]
With $P(gA)\cdot g=g\cdot P(A)$ and the definition of $B_{k,y}$ this gives $B_{k,y'}(gA)\cdot g=g\cdot B_{k,y}(A)$, and for $y=0$ the relation for $B_k=B_{k,0}$. The relations for $D_k$ and $P$, together with $i\gamma_3y'\cdot g=g\cdot S^{-1}\cdot\left(i\gamma_3y'\right)\cdot S=g\cdot i\gamma_3y$, give
\[
\left(D_k(gA)+b_L\cdot P(gA)+i\gamma_3y'\right)\cdot g=g\cdot\left(D_k(A)+b_L\cdot P(A)+i\gamma_3y\right);
\]
the operators in parentheses are invertible by Proposition~\ref{prop:operators-fermion-1}(b), applied at $gA$ and at $A$, with the inverses $\Gamma_{k,y'}(gA)$ and $\Gamma_{k,y}(A)$, so that $\Gamma_{k,y'}(gA)\cdot g=g\cdot\Gamma_{k,y}(A)$, and for $y=0$ this is the relation for $\Gamma_k=\Gamma_{k,0}$. Finally,
\[
\delta Z_k(gA)=\det\left(g\cdot\left(D_k(A)+b_L\cdot P(A)\right)\cdot g^{-1}\right)=\delta Z_k(A).
\]

\emph{(c)} Lemma~\ref{lem:symmetries-elementary-operators}(c) gives $D_A^T=\mathsf{C}^{-1}\cdot D_{-A}\cdot\mathsf{C}$; $Q_k(V)^T=Q_k^T(V)$ and $Q(V)^T=Q^T(V)$ for every connection $V$; $P_k(A)^T=P_k(-A)$, $P(A)^T=P(-A)$, $P_{k+1}(A)^T=P_{k+1}(-A)$; and $F_y(A)^T=\mathsf{C}^{-1}\cdot F_{-y}(-A)\cdot\mathsf{C}$ for real $y$. By Lemma~\ref{lem:block-averages-block}, the operators $Q_k(V)$, $Q_k^T(V)$, $P_k(V)$, $P(V)$ and $P_{k+1}(V)$ commute with the constant matrix $\mathsf{C}$; this gives the assertion for $P_k$, $P$ and $P_{k+1}$, and, with the relation for $D_A$, also
\[
F_{\infty}(A)^T=D_A^T+\mathfrak{m}+b_k\cdot P_k(A)^T=\mathsf{C}^{-1}\cdot F_{\infty}(-A)\cdot\mathsf{C}.
\]
Transposing $Q_k(-A)^T=Q_k^T(-A)$ gives $Q_k^T(-A)^T=Q_k(-A)$. Since the transpose of an inverse is the inverse of the transpose, the relations for $F_y$ give
\[
S_y(A)^T=\mathsf{C}^{-1}\cdot S_{-y}(-A)\cdot\mathsf{C}\quad(y\in\mathbb{R}),\qquad S_k(A)^T=\mathsf{C}^{-1}\cdot S_k(-A)\cdot\mathsf{C},
\]
which are the assertions for $S_{k,y}$, for $S^0_{k+1}=S_0$ and for $S_k=S_{\infty}$. With $M_k(-A)=b_k^2\cdot Q_k(-A)\cdot S_k(-A)\cdot Q_k^T(A)$ and $\dot{M}_k(-A)=b_k^2\cdot Q_k(-A)\cdot S_k(-A)^2\cdot Q_k^T(A)$ it follows that
\begin{align*}
M_k(A)^T&=b_k^2\cdot Q_k^T(-A)^T\cdot S_k(A)^T\cdot Q_k(A)^T=b_k^2\cdot Q_k(-A)\cdot\mathsf{C}^{-1}\cdot S_k(-A)\cdot\mathsf{C}\cdot Q_k^T(A)\\
&=\mathsf{C}^{-1}\cdot M_k(-A)\cdot\mathsf{C},\\
\dot{M}_k(A)^T&=b_k^2\cdot Q_k(-A)\cdot\left(\mathsf{C}^{-1}\cdot S_k(-A)\cdot\mathsf{C}\right)^2\cdot Q_k^T(A)=\mathsf{C}^{-1}\cdot\dot{M}_k(-A)\cdot\mathsf{C},\\
D_k(A)^T&=b_k-M_k(A)^T=\mathsf{C}^{-1}\cdot D_k(-A)\cdot\mathsf{C}.
\end{align*}
The transpose of the multiplication by a constant matrix $m$ is the multiplication by $m^T$, and $\gamma_3^T=-\mathsf{C}^{-1}\gamma_3\mathsf{C}$ by~\eqref{eq:operators-fermion-1}; hence, for real $y$ and $c\in\left\{b_k,\,b_k+b_L\right\}$,
\[
\left(\left(c+i\gamma_3y\right)^{-1}\right)^T=\left(c+i\gamma_3^Ty\right)^{-1}=\mathsf{C}^{-1}\cdot\left(c-i\gamma_3y\right)^{-1}\cdot\mathsf{C}.
\]
As $P(A)^T=P(-A)$ commutes with constant matrices,
\[
B_{k,y}(A)^T=\left(b_k+i\gamma_3^Ty\right)^{-1}\cdot\left(1-P(-A)\right)+\left(b_k+b_L+i\gamma_3^Ty\right)^{-1}\cdot P(-A)=\mathsf{C}^{-1}\cdot B_{k,-y}(-A)\cdot\mathsf{C},
\]
and $y=0$ gives the assertion for $B_k$. The relations for $D_k$, $P$ and $\gamma_3^T$ give
\[
\left(D_k(A)+b_L\cdot P(A)+i\gamma_3y\right)^T=\mathsf{C}^{-1}\cdot\left(D_k(-A)+b_L\cdot P(-A)-i\gamma_3y\right)\cdot\mathsf{C};
\]
by Proposition~\ref{prop:operators-fermion-1}(b), applied at $A$ and at $-A$, the operators in parentheses have the inverses $\Gamma_{k,y}(A)$ and $\Gamma_{k,-y}(-A)$, whence $\Gamma_{k,y}(A)^T=\mathsf{C}^{-1}\cdot\Gamma_{k,-y}(-A)\cdot\mathsf{C}$; the case $y=0$ is the assertion for $\Gamma_k$. Finally, since $H_k(-A)=b_k\cdot S_k(-A)\cdot Q_k^T(A)$ and $T_k(-A)=b_k^{-1}\cdot Q_k(-A)\cdot F_{\infty}(-A)$,
\begin{align*}
H_k(A)^T&=b_k\cdot Q_k^T(-A)^T\cdot S_k(A)^T=b_k\cdot Q_k(-A)\cdot\mathsf{C}^{-1}\cdot S_k(-A)\cdot\mathsf{C}=\mathsf{C}^{-1}\cdot\mathfrak{H}_k(-A)\cdot\mathsf{C},\\
\mathsf{C}^{-1}\cdot H_k(-A)\cdot\mathsf{C}&=b_k\cdot\mathsf{C}^{-1}\cdot S_k(-A)\cdot\mathsf{C}\cdot Q_k^T(A)=b_k\cdot S_k(A)^T\cdot Q_k^T(A),\\
\mathsf{C}^{-1}\cdot T_k(-A)\cdot\mathsf{C}&=b_k^{-1}\cdot Q_k(-A)\cdot\mathsf{C}^{-1}\cdot F_{\infty}(-A)\cdot\mathsf{C}=b_k^{-1}\cdot Q_k(-A)\cdot F_{\infty}(A)^T,
\end{align*}
and $F_{\infty}(A)=D_A+\mathfrak{m}+b_k\cdot P_k(A)$.

\emph{(d)} By Lemma~\ref{lem:symmetries-elementary-operators}(d), $\overline{O\left(\overline{A}\right)}=\varkappa^{-1}\cdot O(-A)\cdot\varkappa$ for $O(A)=D_A$, $Q_k(A)$, $Q_k^T(-A)$, $Q(A)$, $Q^T(-A)$, $P_k(A)$, $P(A)$, $P_{k+1}(A)$ and $F_y(A)$, with the same $y$. By~\eqref{eq:operators-fermion-1}, $\overline{\gamma_3}=-\varkappa^{-1}\gamma_3\varkappa$; hence, for real $y$, $\overline{i\gamma_3y}=-iy\cdot\overline{\gamma_3}=\varkappa^{-1}\cdot\left(i\gamma_3y\right)\cdot\varkappa$, and
\[
\overline{\left(c+i\gamma_3y\right)^{-1}}=\varkappa^{-1}\cdot\left(c+i\gamma_3y\right)^{-1}\cdot\varkappa\qquad\text{for }c\in\left\{b_k,\,b_k+b_L\right\}.
\]
The maps $O\mapsto\overline{O}$ and $O\mapsto\varkappa^{-1}\cdot O\cdot\varkappa$ are additive and multiplicative, leave the multiplication by a real number unchanged, and carry the inverse of an operator to the inverse of its image. The relation therefore holds for every operator obtained from those just named, from real numbers and from the matrices $\left(c+i\gamma_3y\right)^{-1}$ by sums, products and inverses, the inverses existing at $\overline{A}$ and at $-A$ by the first paragraph of this proof. By~\eqref{eq:operators-fermion-2} and the definitions of Sections~\ref{sec:exact-fermion-operators} and~\ref{sec:operators-unit-lattice}, each operator named in (b), with the same $y$, is of this kind: $S_y=F_y^{-1}$ gives $S_k$, $S^0_{k+1}$ and $S_{k,y}$; $H_k$, $M_k$, $D_k$, $\dot{M}_k$ and $T_k$ are the sums and products recalled in the proof of (a); $B_{k,y}$ and $B_k=B_{k,0}$ are formed from $P$ and the matrices $\left(c+i\gamma_3y\right)^{-1}$; and $\Gamma_{k,y}$ and $\Gamma_k=\Gamma_{k,0}$ are the inverses of $D_k+b_L\cdot P+i\gamma_3y$.

\emph{(e)} By (c), $D_k(-A)+b_L\cdot P(-A)=\mathsf{C}\cdot\left(D_k(A)+b_L\cdot P(A)\right)^T\cdot\mathsf{C}^{-1}$. A determinant is invariant under transposition and under conjugation by an invertible operator; hence $\delta Z_k(-A)=\delta Z_k(A)$. By (d),
\[
\overline{D_k\left(\overline{A}\right)+b_L\cdot P\left(\overline{A}\right)}=\varkappa^{-1}\cdot\left(D_k(-A)+b_L\cdot P(-A)\right)\cdot\varkappa,
\]
and the determinant of $\overline{O}$ is the complex conjugate of the determinant of $O$; therefore $\overline{\delta Z_k\left(\overline{A}\right)}=\delta Z_k(-A)=\delta Z_k(A)$.

\emph{The level $0$.} For $k=0$ one has $Y_0=\mathbb{R}$, $F_y(A)=D_A+\mathfrak{m}+b_L\cdot P(A)+i\gamma_3y$, $S_{0,y}(A)=S_y(A)$ and $\Gamma_0(A)=S_0(A)$ (Section~\ref{sec:operators-2}); Lemma~\ref{lem:symmetries-elementary-operators} holds on $T_0$, and $S_y(A)$ exists for $A\in\tilde{G}_0$ by Theorem~\ref{thm:expansion-walks}(a). The argument given at the beginning of the proof of (a), with $k=0$, shows $A-\boldsymbol{\omega}\in\tilde{G}_0$. Inverting the relations of Lemma~\ref{lem:symmetries-elementary-operators}(a)--(d) for $F_y$ gives
\begin{gather*}
S_y(A-\boldsymbol{\omega})=u\cdot S_y(A)\cdot u^{-1},\qquad S_{y'}(gA)\cdot g=g\cdot S_y(A),\\
S_y(A)^T=\mathsf{C}^{-1}\cdot S_{-y}(-A)\cdot\mathsf{C},\qquad\overline{S_y\left(\overline{A}\right)}=\varkappa^{-1}\cdot S_y(-A)\cdot\varkappa,
\end{gather*}
with $y'$ as in (b); these are the assertions for $S_{0,y}$ and for $\Gamma_0=S_{0,0}$. The assertions for $D_A$, $Q(A)$ and $P(A)$ are contained in Lemma~\ref{lem:symmetries-elementary-operators}, the relation $P(A)^T=\mathsf{C}^{-1}\cdot P(-A)\cdot\mathsf{C}$ following from $P(A)^T=P(-A)$ and Lemma~\ref{lem:block-averages-block}. For $\delta Z_0(A)=\det F_0(A)$ one has the relations
\[
F_0(A-\boldsymbol{\omega})=u\cdot F_0(A)\cdot u^{-1},\quad F_0(gA)=g\cdot F_0(A)\cdot g^{-1},\quad F_0(-A)=\mathsf{C}\cdot F_0(A)^T\cdot\mathsf{C}^{-1},\quad\overline{F_0\left(\overline{A}\right)}=\varkappa^{-1}\cdot F_0(-A)\cdot\varkappa;
\]
since a determinant is invariant under conjugation by an invertible operator and under transposition, and $\det\overline{O}$ is the complex conjugate of $\det O$, they give $\delta Z_0(A-\boldsymbol{\omega})=\delta Z_0(gA)=\delta Z_0(-A)=\delta Z_0(A)$ and $\overline{\delta Z_0\left(\overline{A}\right)}=\delta Z_0(A)$.
\end{proof}

Proposition~\ref{prop:symmetries-fermion-operators} follows from Lemma~\ref{lem:symmetries-elementary-operators} and Proposition~\ref{prop:symmetries-fermion-operators-2}.

\begin{proposition}[Operators on a covering torus]\label{prop:operators-covering-torus}
Let $k\ge0$, let $l_1,l_2\in\mathrm{Cl}$ with $l_2=J\cdot l_1$ for an odd integer $J$, and for $i=1,2$ let $T_i:=\eta\mathbb{Z}^3/l_i\mathbb{Z}^3$ be the torus of level $k$ and side $l_i$. A spinor on $T_1$ is a function on $\eta\mathbb{Z}^3$ that is antiperiodic under $l_1\mathbb{Z}^3$; a connection on $T_1$ is a function on the bonds of $\eta\mathbb{Z}^3$, and a scalar on $T_1$ a function on $\eta\mathbb{Z}^3$, that is periodic under $l_1\mathbb{Z}^3$. Since $J$ is odd, such a function is antiperiodic, respectively periodic, under $l_2\mathbb{Z}^3$, and so it is a spinor, a connection or a scalar on $T_2$. This inclusion, together with the corresponding inclusion of the spinors on the unit lattice and on the $L$-lattice, is denoted by $\mathrm{pj}^{*}$ (Definition~\ref{def:periodisation-projection}). For an operator $O$ of Section~\ref{sec:exact-fermion-operators} let $O^{(i)}$ denote the operator $O$ on $T_i$. Assume the hypotheses of Theorem~\ref{thm:expansion-walks}, and let $A\in\tilde{G}_k(T_1)$. Then $\mathrm{pj}^{*}A\in\tilde{G}_k(T_2)$, and for every real $y$ the relation
\[
O^{(2)}\left(\mathrm{pj}^{*}A\right)\cdot\mathrm{pj}^{*}=\mathrm{pj}^{*}\cdot O^{(1)}(A)
\]
holds, if $k\ge1$, for $O(A)=D_A$, $Q_k(A)$, $Q_k^T(-A)$, $P_k(A)$, $P_{k+1}(A)$, $Q(A)$, $Q^T(-A)$, $P(A)$, $S_k(A)$, $H_k(A)$, $\bar{H}_k(A)$, $D_k(A)$, $M_k(A)$, $\dot{M}_k(A)$, $\Gamma_k(A)$, $\Gamma_{k,y}(A)$, $B_{k,y}(A)$, $S_{k,y}(A)$, $S^0_{k+1}(A)$, $T_k(A)$, $\bar{T}_k(A)$, and, if $k=0$, for $O(A)=D_A$, $Q(A)$, $Q^T(-A)$, $P(A)$, $\Gamma_0(A)$, $S_{0,y}(A)$. Moreover, for every spinor $f$ on $T_1$,
\[
\left(\delta_{\alpha,\mathrm{pj}^{*}A}\,\mathrm{pj}^{*}f\right)(x,x')=\left(\delta_{\alpha,A}f\right)(x,x')
\]
at every pair $(x,x')$ of points of $\eta\mathbb{Z}^3$ with $0<|x-x'|<l_1/2$.
\end{proposition}

\begin{proof}
The definition of $\tilde{G}_k$ (Section~\ref{sec:domain-complex-connections}) consists of bounds on sup norms and of representations on the sets $\square^{\sim 5}$, which have side $11M<13LM\le l_1$ and therefore embed in both tori. In detail, let $A=A_0+A_1$ be as in that definition, on $T_1$. Then $\mathrm{pj}^{*}A=\mathrm{pj}^{*}A_0+\mathrm{pj}^{*}A_1$, where $\mathrm{pj}^{*}A_0$ is real and $\left|\mathrm{pj}^{*}A_1\right|_{\infty}=|A_1|_{\infty}<e_k^{-1+\varepsilon}$. Let $\square$ be an $M$-cube of $T_2$, and $\square_1$ its image under the covering map $T_2\to T_1$, an $M$-cube of $T_1$. As $11M<l_1$, the covering map restricts to a bijection of $\square^{\sim 5}$ onto $\square_1^{\sim 5}$ that maps bonds onto bonds. Composing with this bijection the real functions $\boldsymbol{\lambda}$ and $A'$ of a representation $A_0=A'+d\boldsymbol{\lambda}$ on $\square_1^{\sim 5}$ gives a representation of $\mathrm{pj}^{*}A_0$ on $\square^{\sim 5}$ with the same bound $|A'|_{\infty}<e_k^{-1+\varepsilon}$. Hence $\mathrm{pj}^{*}A\in\tilde{G}_k(T_2)$, and the assertion of Lemma~\ref{lem:periodic-extension} on $\tilde{G}_k$ follows.

The relations for the operators rest on three observations.

\emph{(i)} For every connection $V$ on $T_1$, the operators $D_V$, $Q_k(V)$, $Q_k^T(-V)$, $Q(V)$, $Q^T(-V)$, $P_k(V)$, $P(V)$, $P_{k+1}(V)$, $\operatorname{Pen}_y(V)$, $B_{k,y}(V)$ and $\delta_{\alpha,V}$ are given by the formulas of Sections~\ref{sec:phase-factors-lattice}, \ref{sec:operators-2} and~\ref{sec:operators-unit-lattice}, which express the image of a function at a site, or at a pair of sites, through the values of the function and of the connection on $\eta\mathbb{Z}^3$. These formulas involve the lattices $\Lambda_j=L^j\eta\mathbb{Z}^3$, the blocks $\mathrm{Bl}_j(z)$, and paths $\Gamma^{\pi}(x,x')$ between points of $\eta\mathbb{Z}^3$ at distance less than $l_1/2$; all of these are the same for $T_1$ and for $T_2$. Consequently $O^{(2)}\left(\mathrm{pj}^{*}V\right)\cdot\mathrm{pj}^{*}f$ and $\mathrm{pj}^{*}\cdot O^{(1)}(V)f$ are the same function, for each of these operators and every $f$ on $T_1$, the functions of pairs $\delta_{\alpha,\mathrm{pj}^{*}V}\,\mathrm{pj}^{*}f$ and $\delta_{\alpha,V}f$ being compared at the pairs $(x,x')$ with $0<|x-x'|<l_1/2$.

\emph{(ii)} If the relation holds for $O_1$ and $O_2$, then it holds for $O_1+O_2$ and, since
\[
O_1^{(2)}\left(\mathrm{pj}^{*}V\right)\cdot O_2^{(2)}\left(\mathrm{pj}^{*}V\right)\cdot\mathrm{pj}^{*}=O_1^{(2)}\left(\mathrm{pj}^{*}V\right)\cdot\mathrm{pj}^{*}\cdot O_2^{(1)}(V)=\mathrm{pj}^{*}\cdot O_1^{(1)}(V)\cdot O_2^{(1)}(V),
\]
for $O_1O_2$; it holds for the multiplication by a number and by a constant $4\times4$ matrix.

\emph{(iii)} If the relation holds for $O$, and $O^{(1)}(V)$ and $O^{(2)}\left(\mathrm{pj}^{*}V\right)$ are bijective, then it holds for their inverses: for every $f$ on $T_1$,
\[
O^{(2)}\left(\mathrm{pj}^{*}V\right)\cdot\mathrm{pj}^{*}\cdot O^{(1)}(V)^{-1}f=\mathrm{pj}^{*}\cdot O^{(1)}(V)\cdot O^{(1)}(V)^{-1}f=\mathrm{pj}^{*}f,
\]
and applying $O^{(2)}\left(\mathrm{pj}^{*}V\right)^{-1}$ gives $\mathrm{pj}^{*}\cdot O^{(1)}(V)^{-1}f=O^{(2)}\left(\mathrm{pj}^{*}V\right)^{-1}\cdot\mathrm{pj}^{*}f$.

Now let $V\in\{A,-A\}$. Then $V\in\tilde{G}_k(T_1)$ by~\eqref{eq:operators-fermion-G2}, and $\mathrm{pj}^{*}V\in\tilde{G}_k(T_2)$ by the first paragraph; the hypotheses of Theorem~\ref{thm:expansion-walks} hold on both tori, as $l_1,l_2\in\mathrm{Cl}$. By Theorem~\ref{thm:expansion-walks}(a), $F_y(V)$ and $F_y\left(\mathrm{pj}^{*}V\right)$ are bijective for $y\in Y_k$; for $k\ge1$ and real $y$, so are $D_k(V)+b_L\cdot P(V)+i\gamma_3y$ and the same operator at $\mathrm{pj}^{*}V$, by Proposition~\ref{prop:operators-fermion-1}(b). By (i) and (ii) the relation holds for $F_y(V)=D_V+\mathfrak{m}+\operatorname{Pen}_y(V)$, and by (iii) for $S_y$, that is, for $S_k=S_{\infty}$, $S^0_{k+1}=S_0$ and $S_{k,y}$ if $k\ge1$, and for $S_{0,y}=S_y$ and $\Gamma_0=S_0$ if $k=0$. By (ii) it holds for
\begin{gather*}
H_k(V)=b_k\cdot S_k(V)\cdot Q_k^T(-V),\qquad M_k(V)=b_k^2\cdot Q_k(V)\cdot S_k(V)\cdot Q_k^T(-V),\qquad D_k(V)=b_k-M_k(V),\\
\dot{M}_k(V)=b_k^2\cdot Q_k(V)\cdot S_k(V)^2\cdot Q_k^T(-V),\qquad T_k(V)=b_k^{-1}\cdot Q_k(V)\cdot F_{\infty}(V),
\end{gather*}
and by (ii) and (iii) for $\Gamma_{k,y}(V)=\left(D_k(V)+b_L\cdot P(V)+i\gamma_3y\right)^{-1}$ and for $\Gamma_k=\Gamma_{k,0}$. Finally, $\mathrm{pj}^{*}$ commutes with the multiplication by the constant matrix $\mathsf{C}$, and $\mathrm{pj}^{*}(-A)=-\mathrm{pj}^{*}A$; the relations for $H_k$ and $T_k$ at $V=-A$ give
\begin{align*}
\bar{H}_k^{(2)}\left(\mathrm{pj}^{*}A\right)\cdot\mathrm{pj}^{*}&=\mathsf{C}^{-1}\cdot H_k^{(2)}\left(\mathrm{pj}^{*}(-A)\right)\cdot\mathrm{pj}^{*}\cdot\mathsf{C}=\mathsf{C}^{-1}\cdot\mathrm{pj}^{*}\cdot H_k^{(1)}(-A)\cdot\mathsf{C}=\mathrm{pj}^{*}\cdot\bar{H}_k^{(1)}(A),\\
\bar{T}_k^{(2)}\left(\mathrm{pj}^{*}A\right)\cdot\mathrm{pj}^{*}&=\mathsf{C}^{-1}\cdot T_k^{(2)}\left(\mathrm{pj}^{*}(-A)\right)\cdot\mathrm{pj}^{*}\cdot\mathsf{C}=\mathsf{C}^{-1}\cdot\mathrm{pj}^{*}\cdot T_k^{(1)}(-A)\cdot\mathsf{C}=\mathrm{pj}^{*}\cdot\bar{T}_k^{(1)}(A).
\end{align*}
\end{proof}

Proposition~\ref{prop:coverings-tori} follows.

\subsection{The fermion determinant}\label{sec:fermion-determinant}

Put $n_k:=l^3$, which is the number of unit sites of $T_k$ modulo $l\mathbb{Z}^3$, and
\[
c^d_k:=4\left[\left(1-L^{-3}\right)\log b_k+L^{-3}\log\left(b_k+b_L\right)\right]\quad\text{for }k\ge1,\qquad c^d_0:=0.
\]
At level $0$ the determinant is defined by
\[
\delta Z_0(A):=\det F_0(A)=\det\left(D_A+\mathfrak{m}+b_L\cdot P(A)\right),
\]
where $F_0(A)$ is the operator~\eqref{eq:operators-fermion-2} of level $0$ at $y=0$.

\begin{lemma}[Logarithm of the determinant at real connections]\label{lem:logarithm-determinant-real}
Assume the hypotheses of Theorem~\ref{thm:expansion-walks}, and let $A\in\tilde{G}_k$ be real. If $k\ge1$, then
\[
\delta Z_k(A)=\exp\left\{c^d_k\cdot n_k-i\,b_k^2\int_0^{\infty}\mathrm{d}y\,\operatorname{Tr}\left[\gamma_3\cdot B_{k,y}(A)\cdot Q_k(A)\cdot S_{k,y}(A)\cdot Q_k^T(-A)\cdot B_{k,y}(A)\right]\right\}.
\]
If $k=0$, then
\[
\delta Z_0(A)=\exp\left\{\int_1^{\infty}\frac{\mathrm{d}y}{y}\,\operatorname{Tr}\left[F_0(A)\cdot S_{0,y}(A)\right]-i\int_0^1\mathrm{d}y\,\operatorname{Tr}\left[\gamma_3\cdot S_{0,y}(A)\right]\right\}.
\]
All integrals in these two formulas converge absolutely.
\end{lemma}

\begin{proof}
\emph{Two elementary integrals.} For a complex number $z\ne0$ with $\operatorname{Im}z\ge0$ let
\[
\operatorname{Log} z:=\log|z|+i\arg z,\qquad 0\le\arg z\le\pi .
\]
On the closed upper half plane with $0$ removed, this function coincides with the branch of the logarithm that is holomorphic on the complement of the half line $\left\{-it:\ t\ge0\right\}$ and whose imaginary part lies in $\left(-\frac{\pi}{2},\frac{3\pi}{2}\right)$. Consequently, for real $x\ne0$, the function $y\mapsto\operatorname{Log}(x+iy)$ is differentiable on $[0,\infty)$, with the derivative $i\,(x+iy)^{-1}$. Moreover $|x+iR_0|/R_0\to1$ and $\arg(x+iR_0)\to\frac{\pi}{2}$, so that
\[
\operatorname{Log}\left(x+iR_0\right)=\log R_0+\frac{i\pi}{2}+o(1)\qquad\left(R_0\to\infty\right).
\]
Let $x\ne0$ be real and $R_0>0$. The two identities
\begin{align*}
-i\int_0^{R_0}\frac{\mathrm{d}y}{x+iy}&=\operatorname{Log} x-\operatorname{Log}\left(x+iR_0\right),\\
\int_1^{\infty}\frac{\mathrm{d}y}{y}\,\frac{x}{x+iy}-i\int_0^1\frac{\mathrm{d}y}{x+iy}&=\operatorname{Log} x-\frac{i\pi}{2}
\end{align*}
hold. The first one is obtained by integrating the derivative of $y\mapsto\operatorname{Log}(x+iy)$ over $[0,R_0]$. For the second one, note that
\[
\frac{x}{y\,(x+iy)}=\frac{1}{y}-\frac{i}{x+iy}\qquad(y>0);
\]
integrating, one finds for $R>1$
\[
\int_1^R\frac{\mathrm{d}y}{y}\,\frac{x}{x+iy}=\log R-\operatorname{Log}(x+iR)+\operatorname{Log}(x+i),
\]
and the right-hand side tends to $\operatorname{Log}(x+i)-\frac{i\pi}{2}$ as $R\to\infty$. The integral over $[1,\infty)$ converges absolutely, since its integrand is bounded in modulus by $|x|\cdot y^{-2}$. The first identity with $R_0=1$ gives $-i\int_0^1\frac{\mathrm{d}y}{x+iy}=\operatorname{Log} x-\operatorname{Log}(x+i)$. Adding the last two results one obtains the second identity.

\emph{Reduction to a Hermitian matrix.} Let
\[
K:=D_k(A)+b_L\cdot P(A)\quad\text{if }k\ge1,\qquad K:=F_0(A)\quad\text{if }k=0,\qquad\text{and}\qquad T:=K\cdot\gamma_3.
\]
Both $K$ and $T$ are linear maps of the space of the spinors on the unit lattice of $T_k$ into itself; for $k=0$ the unit lattice is the lattice of $T_0$. This space has the dimension $4n_k$, and $K$, $T$ are regarded as matrices of size $4n_k$. The adjoints are taken with respect to the sesquilinear pairing of Section~\ref{sec:norms-pairings-kernels}, which on this space is the sum over the unit sites; thus $O^{*}$ is the conjugate transpose of a matrix $O$.

Since $A$ is real, Lemma~\ref{lem:symmetries-elementary-operators}(e) gives
\[
D_A^{*}=\gamma_3\cdot D_A\cdot\gamma_3,\qquad Q_k^T(-A)=Q_k(A)^{*},
\]
and shows that $P_k(A)$ and $P(A)$ are orthogonal projections, hence self-adjoint. By Lemma~\ref{lem:block-averages-block}, the operators $Q_k(A)$, $Q_k^T(-A)$, $P_k(A)$ and $P(A)$ commute with $\gamma_3$. The mass $\mathfrak{m}$, the number $b_L$ and, for $k\ge1$, the number $b_k$ are real, and by~\eqref{eq:operators-fermion-1} the matrix $\gamma_3$ is Hermitian with $\gamma_3^2=1$.

Let $k\ge1$. The operator $F_{\infty}(A)=D_A+\mathfrak{m}+b_k\cdot P_k(A)$ of Section~\ref{sec:operators-2} therefore satisfies $F_{\infty}(A)^{*}=\gamma_3\cdot F_{\infty}(A)\cdot\gamma_3$, and inverting both sides yields $S_k(A)^{*}=\gamma_3\cdot S_k(A)\cdot\gamma_3$. As $M_k(A)=b_k^2\cdot Q_k(A)\cdot S_k(A)\cdot Q_k^T(-A)$ (Section~\ref{sec:operators-unit-lattice}), it follows that
\[
M_k(A)^{*}=b_k^2\cdot Q_k(A)\cdot S_k(A)^{*}\cdot Q_k^T(-A)=\gamma_3\cdot M_k(A)\cdot\gamma_3,
\]
and, since $D_k(A)=b_k-M_k(A)$, that $K^{*}=\gamma_3\cdot K\cdot\gamma_3$. For $k=0$ the relation $K^{*}=\gamma_3\cdot K\cdot\gamma_3$ follows directly from $K=D_A+\mathfrak{m}+b_L\cdot P(A)$. In both cases
\[
T^{*}=\gamma_3\cdot K^{*}=\gamma_3^2\cdot K\cdot\gamma_3=T,
\]
so that $T$ is a Hermitian matrix.

The matrix $K$ is invertible: for $k\ge1$ its inverse is $\Gamma_k(A)$, by Proposition~\ref{prop:operators-fermion-1}(d), and for $k=0$ the operator $K=F_0(A)$ is bijective by Theorem~\ref{thm:expansion-walks}(a). As $\gamma_3$ is invertible, $T$ is invertible as well. The matrix $\gamma_3=\sigma_1\otimes1$ (Section~\ref{sec:dirac-matrices-lattice}) has the eigenvalues $1,1,-1,-1$, hence the determinant $1$; the multiplication by $\gamma_3$ on the spinors on the unit lattice is the direct sum of $n_k$ copies of this matrix, one for each unit site, and its determinant is $1$ too. Therefore
\[
\det T=\det K=\delta Z_k(A),
\]
the second equality being the definition of $\delta Z_k(A)$ in Section~\ref{sec:operators-unit-lattice} if $k\ge1$, and the definition of $\delta Z_0(A)$ given above if $k=0$.

By the spectral theorem for Hermitian matrices there is an orthonormal basis of eigenvectors of $T$. Let $x_1,\dots,x_{4n_k}$ be the corresponding eigenvalues; they are real, and they are different from $0$ because $T$ is invertible. Since $\exp\left(\operatorname{Log} z\right)=z$,
\[
\delta Z_k(A)=\prod_{j=1}^{4n_k}x_j=\exp\left\{\sum_{j=1}^{4n_k}\operatorname{Log} x_j\right\}.
\]
Let $y$ be real. As $\gamma_3^2=1$, one has $T+iy=\left(K+i\gamma_3y\right)\cdot\gamma_3$. The eigenvalues $x_j+iy$ of $T+iy$ do not vanish, so $T+iy$ and $K+i\gamma_3y$ are invertible, and
\[
(T+iy)^{-1}=\gamma_3\cdot\left(K+i\gamma_3y\right)^{-1}=
\begin{cases}
\gamma_3\cdot\Gamma_{k,y}(A), & k\ge1,\\
\gamma_3\cdot S_{0,y}(A), & k=0,
\end{cases}
\]
by Proposition~\ref{prop:operators-fermion-1}(b) if $k\ge1$, and by the definition of $S_{0,y}(A)=F_y(A)^{-1}$ in Section~\ref{sec:operators-2} if $k=0$. In the orthonormal basis of eigenvectors the matrices $(T+iy)^{-1}$ and $T\cdot(T+iy)^{-1}$ are diagonal, with the diagonal entries $(x_j+iy)^{-1}$ and $x_j\cdot(x_j+iy)^{-1}$. Consequently
\[
\sum_{j=1}^{4n_k}\frac{1}{x_j+iy}=\operatorname{Tr}\left[(T+iy)^{-1}\right],\qquad\sum_{j=1}^{4n_k}\frac{x_j}{x_j+iy}=\operatorname{Tr}\left[T\cdot(T+iy)^{-1}\right].
\]

\emph{The case $k=0$.} Here $T\cdot(T+iy)^{-1}=K\cdot\gamma_3^2\cdot S_{0,y}(A)=F_0(A)\cdot S_{0,y}(A)$ and $(T+iy)^{-1}=\gamma_3\cdot S_{0,y}(A)$. Apply the second of the two elementary identities to $x=x_j$ and add over $j$. By the two trace formulas just proved,
\begin{align*}
&\int_1^{\infty}\frac{\mathrm{d}y}{y}\,\operatorname{Tr}\left[F_0(A)\cdot S_{0,y}(A)\right]-i\int_0^1\mathrm{d}y\,\operatorname{Tr}\left[\gamma_3\cdot S_{0,y}(A)\right]\\
&\qquad=\sum_{j=1}^{4n_0}\operatorname{Log} x_j-4n_0\cdot\frac{i\pi}{2},
\end{align*}
and $4n_0\cdot\frac{i\pi}{2}=2\pi i\,n_0$. Since $\exp\left(2\pi i\,n_0\right)=1$, exponentiation gives the assertion for $k=0$. By Theorem~\ref{thm:expansion-walks}(c), with $w_0(y)=(1+|y|)^{-1}$ (see~\eqref{eq:operators-fermion-2}), the entries of the kernel of $S_{0,y}(A)$ are $O\left((1+y)^{-1}\right)$ for $y\ge0$, uniformly in the pair of sites; the trace being a finite sum, $\left|\operatorname{Tr}\left[F_0(A)\cdot S_{0,y}(A)\right]\right|=O\left((1+y)^{-1}\right)$, and the first integral converges absolutely. The integrand of the second integral equals $\sum_j(x_j+iy)^{-1}$ and is bounded in modulus by $\sum_j|x_j|^{-1}$.

\emph{The case $k\ge1$.} Apply the first of the two elementary identities to $x=x_j$, add over $j$, and use the asymptotic formula for $\operatorname{Log}\left(x_j+iR_0\right)$ and the expression of $(T+iy)^{-1}$ found above. This gives, as $R_0\to\infty$,
\[
\sum_{j=1}^{4n_k}\operatorname{Log} x_j=-i\int_0^{R_0}\mathrm{d}y\,\operatorname{Tr}\left[\gamma_3\cdot\Gamma_{k,y}(A)\right]+4n_k\cdot\left(\log R_0+\frac{i\pi}{2}\right)+o(1).
\]
The left-hand side does not depend on $R_0$; hence it is the limit, as $R_0\to\infty$, of the sum of the first two terms on the right. By Proposition~\ref{prop:operators-fermion-1}(b), in which $S_y(A)=S_{k,y}(A)$ (Section~\ref{sec:operators-2}),
\begin{align*}
\operatorname{Tr}\left[\gamma_3\cdot\Gamma_{k,y}(A)\right]&=\operatorname{Tr}\left[\gamma_3\cdot B_{k,y}(A)\right]\\
&\qquad+b_k^2\cdot\operatorname{Tr}\left[\gamma_3\cdot B_{k,y}(A)\cdot Q_k(A)\cdot S_{k,y}(A)\cdot Q_k^T(-A)\cdot B_{k,y}(A)\right].
\end{align*}
The second trace on the right is $O\left(\left(b_k^2+y^2\right)^{-1}\right)$, by Proposition~\ref{prop:kernels-unit-lattice}(b), (c) and Theorem~\ref{thm:expansion-walks}(c). Indeed, the kernel of $B_{k,y}(A)$ vanishes at the pairs of unit sites that do not lie in one $L$-block, and its entries are bounded by $\left(b_k^2+y^2\right)^{-\frac{1}{2}}$ times a number independent of $y$; the entries of the kernel of $Q_k(A)\cdot S_{k,y}(A)\cdot Q_k^T(-A)$ are bounded independently of $y$, because $w_k(y)=1$ for $k\ge1$; and the trace is a finite sum of products of such entries. Therefore the integral of the second trace over $[0,\infty)$ converges absolutely, and its integral over $[0,R_0]$ tends to the integral over $[0,\infty)$ as $R_0\to\infty$.

To compute the first trace, recall from Lemma~\ref{lem:block-averages-block} that $P(A)$ commutes with the constant $4\times4$ matrices, so that its kernel at every pair of unit sites is a multiple of the unit matrix, and that $\operatorname{Tr}P(A)=4\left(l/L\right)^3=4n_k\cdot L^{-3}$. Let $\operatorname{tr}$ denote the trace of $4\times4$ matrices, and let $\phi(\gamma_3)$ be a $4\times4$ matrix that is a function of $\gamma_3$, acting on the spinors at every unit site. Then $\operatorname{Tr}\phi(\gamma_3)=n_k\cdot\operatorname{tr}\phi(\gamma_3)$ and $\operatorname{Tr}\left[\phi(\gamma_3)\cdot P(A)\right]=\frac{1}{4}\operatorname{Tr}P(A)\cdot\operatorname{tr}\phi(\gamma_3)=n_k\cdot L^{-3}\cdot\operatorname{tr}\phi(\gamma_3)$, whence
\[
\operatorname{Tr}\left[\phi(\gamma_3)\cdot\left(1-P(A)\right)\right]=n_k\cdot\left(1-L^{-3}\right)\cdot\operatorname{tr}\phi(\gamma_3).
\]
The definition of $B_{k,y}(A)$ in Section~\ref{sec:operators-unit-lattice}, together with these two formulas for $\phi(\gamma_3)=\gamma_3\cdot\left(b_k+i\gamma_3y\right)^{-1}$ and for $\phi(\gamma_3)=\gamma_3\cdot\left(b_k+b_L+i\gamma_3y\right)^{-1}$, gives
\begin{align*}
\operatorname{Tr}\left[\gamma_3\cdot B_{k,y}(A)\right]&=n_k\cdot\left(1-L^{-3}\right)\cdot\operatorname{tr}\left[\gamma_3\cdot\left(b_k+i\gamma_3y\right)^{-1}\right]\\
&\qquad+n_k\cdot L^{-3}\cdot\operatorname{tr}\left[\gamma_3\cdot\left(b_k+b_L+i\gamma_3y\right)^{-1}\right].
\end{align*}
Let $\boldsymbol{\beta}>0$. Since $\gamma_3$ has the eigenvalues $1,1,-1,-1$,
\[
\operatorname{tr}\left[\gamma_3\cdot\left(\boldsymbol{\beta}+i\gamma_3y\right)^{-1}\right]=\frac{2}{\boldsymbol{\beta}+iy}-\frac{2}{\boldsymbol{\beta}-iy}=-\frac{4iy}{\boldsymbol{\beta}^2+y^2},
\]
and therefore, as $R_0\to\infty$,
\begin{align*}
-i\int_0^{R_0}\mathrm{d}y\,\left(-\frac{4iy}{\boldsymbol{\beta}^2+y^2}\right)&=-2\log\left(\boldsymbol{\beta}^2+R_0^2\right)+4\log\boldsymbol{\beta}\\
&=-4\log R_0+4\log\boldsymbol{\beta}+o(1).
\end{align*}
The numbers $b_k$ and $b_k+b_L$ are positive, as $b_k\ge b_{-}>0$ (Section~\ref{sec:parameters-clifford-relations}). Taking $\boldsymbol{\beta}=b_k$ with the weight $n_k\cdot\left(1-L^{-3}\right)$ and $\boldsymbol{\beta}=b_k+b_L$ with the weight $n_k\cdot L^{-3}$, and comparing with the definition of $c^d_k$, one obtains
\[
-i\int_0^{R_0}\mathrm{d}y\,\operatorname{Tr}\left[\gamma_3\cdot B_{k,y}(A)\right]=-4n_k\cdot\log R_0+c^d_k\cdot n_k+o(1)\qquad\left(R_0\to\infty\right).
\]
When these results on the two traces are inserted into the limit that represents $\sum_j\operatorname{Log} x_j$, the multiples of $\log R_0$ cancel each other; as $4n_k\cdot\frac{i\pi}{2}=2\pi i\,n_k$, it follows that
\begin{align*}
\sum_{j=1}^{4n_k}\operatorname{Log} x_j&=c^d_k\cdot n_k-i\,b_k^2\int_0^{\infty}\mathrm{d}y\,\operatorname{Tr}\left[\gamma_3\cdot B_{k,y}(A)\cdot Q_k(A)\cdot S_{k,y}(A)\cdot Q_k^T(-A)\cdot B_{k,y}(A)\right]\\
&\qquad+2\pi i\,n_k.
\end{align*}
Since $\exp\left(2\pi i\,n_k\right)=1$, exponentiation gives the assertion for $k\ge1$.
\end{proof}

Let a system of primitives of level $k\ge0$ on $T_k$ be given (Section~\ref{sec:primitive-operators}), and let $S_{k,y}(X,A)$, for $X\in D_k$ (the set of polymers of Section~\ref{sec:lattice-geometry}), $A\in\tilde{G}_k$ and $y\ge0$, be its family~\eqref{it:primitives-p5}; for $k=0$ this is the family $S_{0,y}(X,A)$ of the system. Let $\operatorname{Tr}$ denote the trace of linear operators on the spinors on the unit lattice $T^0_k$, which for $k=0$ is $T_0$, and let $\operatorname{tr}$ denote the trace of complex $4\times4$ matrices. For a complex number, or a complex function, $z$ let $\overline{z}$ be its complex conjugate. For $X\in D_k$ and $A\in\tilde{G}_k$ put, as in Definitions~\ref{def:contour-functional-functions} and~\ref{def:determinant-functions},
\begin{align*}
\acute{E}_k(X,A)&:=-i\,b_k^2\int_0^{\infty}\mathrm{d}y\,\operatorname{Tr}\left[\gamma_3\cdot B_{k,y}(A)\cdot Q_k(A)\cdot S_{k,y}(X,A)\cdot Q_k^T(-A)\cdot B_{k,y}(A)\right]\qquad(k\ge1),\\
\acute{E}_0(X,A)&:=\int_1^{\infty}\frac{\mathrm{d}y}{y}\,\operatorname{Tr}\left[F_0(A)\cdot S_{0,y}(X,A)\right]-i\int_0^1\mathrm{d}y\,\operatorname{Tr}\left[\gamma_3\cdot S_{0,y}(X,A)\right],
\end{align*}
where $B_{k,y}(A)$ is the operator of Section~\ref{sec:operators-unit-lattice} and $F_0(A)=D_A+\mathfrak{m}+b_L\cdot P(A)$ is the operator~\eqref{eq:operators-fermion-2} for $k=0$ and $y=0$. By~\eqref{eq:operators-fermion-G2}, the connections $-A$, $\overline{A}$ and $-\overline{A}$ lie in $\tilde{G}_k$ together with $A$, and $0\in G_k\subset\tilde{G}_k$. Put
\begin{align*}
E^d_k(X,A)&:=\frac{1}{4}\left[\acute{E}_k(X,A)+\acute{E}_k(X,-A)+\overline{\acute{E}_k(X,\overline{A})}+\overline{\acute{E}_k(X,-\overline{A})}\right],\\
E^{\mathrm{det}}_k(X,A)&:=E^d_k(X,A)-E^d_k(X,0),
\end{align*}
and, for connections $\mathscr{A}$, $\mathscr{C}$ on $T_k$ with $\mathscr{C}\in\tilde{G}_k$ and $\mathscr{A}+\mathscr{C}\in\tilde{G}_k$,
\[
\delta E^{\mathrm{det}}_k(X;\mathscr{A},\mathscr{C}):=E^d_k(X,\mathscr{A}+\mathscr{C})-E^d_k(X,\mathscr{C}).
\]
Finally, for a number $C_w>0$, with the number $c_X$ of stage~\ref{it:order-choice-1}, let
\begin{equation}
C_d:=16\,c_X\cdot C_w\cdot\max\left\{\frac{\pi}{2}\cdot b\cdot\mathbf{e}\cdot\left(1+2\mathbf{e}\right),\ 5+6\mathbf{e}+b\cdot\mathbf{e}\right\}.\label{eq:operators-fermion-10}
\end{equation}

\begin{proposition}[Polymer expansion of the determinant]\label{prop:polymer-expansion-determinant}
Assume the hypotheses of Theorem~\ref{thm:expansion-walks}. Let $\kappa$, $\kappa_0$, $c_X$, $c_{\mathrm{rb}}$, $r$ and $\bar{\varepsilon}$ be the numbers fixed in stages~\ref{it:order-choice-1}, \ref{it:order-choice-3} and~\ref{it:order-choice-8} of Definition~\ref{def:order-choice}, and let $\kappa_w=\kappa+\kappa_0+3$. Let $\gamma_w>0$ and $C_w>0$ be arbitrary numbers, and let $(\alpha_0,\gamma_w,\kappa_w;C_w)$ be the parameter set (Section~\ref{sec:primitive-operators}) formed with them; in this proposition and in its proof the letters $\gamma_w$ and $C_w$ denote these two numbers. The case used in the scheme is the parameter set $\mathrm{par}$ of Section~\ref{sec:primitive-operators}, in which $\gamma_w$ and $C_w$ are the numbers so denoted in that section; for $\mathrm{par}$ the proposition yields Proposition~\ref{prop:local-expansion-determinant}. Let $k\ge0$, and let a family~\eqref{it:primitives-p5} of level $k$ on $T_k$, with members $S_{k,y}(X,A)$, have the properties~\ref{it:renormalisation-group-1}, \ref{it:renormalisation-group-2}, \ref{it:renormalisation-group-7}, \ref{it:renormalisation-group-12}, \ref{it:renormalisation-group-13} and~\ref{it:renormalisation-group-14} for the parameter set $(\alpha_0,\gamma_w,\kappa_w;C_w)$, where for $k=0$ property~\ref{it:renormalisation-group-7} is taken in the form stated in property~\ref{it:renormalisation-group-17}. Let $\acute{E}_k$, $E^d_k$, $E^{\mathrm{det}}_k$ and $\delta E^{\mathrm{det}}_k$ be formed with this family, and let $C_d$ be the number~\eqref{eq:operators-fermion-10} formed with this $C_w$. Then the following hold.
\begin{enumerate}
\item[(a)]\itemmark{a}\label{it:polymer-expansion-determinant-a} For every $A\in\tilde{G}_k$, with the numbers $c^d_k$ and $n_k$ of Section~\ref{sec:fermion-determinant},
\[
\delta Z_k(A)=\exp\left\{c^d_k\cdot n_k+\sum_{X\in D_k}\acute{E}_k(X,A)\right\},\qquad\frac{\delta Z_k(A)}{\delta Z_k(0)}=\exp\left\{\sum_{X\in D_k}E^{\mathrm{det}}_k(X,A)\right\}.
\]
\item[(b)]\itemmark{b}\label{it:polymer-expansion-determinant-b} For every $X\in D_k$ the functions $\acute{E}_k(X,\cdot)$ and $E^d_k(X,\cdot)$ are analytic on $\tilde{G}_k$, and for every $A\in\tilde{G}_k$
\[
\left|\acute{E}_k(X,A)\right|\le C_d\cdot M^3\cdot\exp\left(-(\kappa+1)\cdot d_M(X)\right),\qquad\left|E^d_k(X,A)\right|\le C_d\cdot M^3\cdot\exp\left(-(\kappa+1)\cdot d_M(X)\right).
\]
\item[(c)]\itemmark{c}\label{it:polymer-expansion-determinant-c} Let $X\in D_k$, and let $A,\tilde{A}\in\tilde{G}_k$ coincide on the bonds that meet $X$. Then $\acute{E}_k(X,A)=\acute{E}_k(X,\tilde{A})$ and $E^d_k(X,A)=E^d_k(X,\tilde{A})$.
\item[(d)]\itemmark{d}\label{it:polymer-expansion-determinant-d} Let $X\in D_k$ and $A\in\tilde{G}_k$. For every real scalar $\boldsymbol{\lambda}$ on $T_k$,
\[
\acute{E}_k(X,A-d\boldsymbol{\lambda})=\acute{E}_k(X,A),\qquad E^d_k(X,A-d\boldsymbol{\lambda})=E^d_k(X,A);
\]
moreover $E^d_k(X,-A)=E^d_k(X,A)$, $\overline{E^d_k(X,\overline{A})}=E^d_k(X,A)$, and $E^d_k(gX,gA)=E^d_k(X,A)$ for every $g\in\mathbb{G}_k^L$.
\item[(e)]\itemmark{e}\label{it:polymer-expansion-determinant-e} Assume condition~\eqref{eq:renormalisation-group-23}, $l\ge M\cdot r_k$, $v_k\ge9$ and $8M\cdot\left(\frac{v_k}{4}+4\right)\cdot e_k^{\varepsilon}\le1$. Then for every $X\in D_k$ and every $(\mathscr{A},\mathscr{C})\in\tilde{R}_k\times\mathfrak{X}_k$
\[
\left|\delta E^{\mathrm{det}}_k(X;\mathscr{A},\mathscr{C})\right|\le8C_d\cdot M^3\cdot e_k^{\frac{1}{4}}\cdot\exp\left(-\kappa\cdot d_M(X)\right);
\]
in particular $\left|E^{\mathrm{det}}_k(X,\mathscr{A})\right|$ satisfies the same bound.
\item[(f)]\itemmark{f}\label{it:polymer-expansion-determinant-f} Assume, in addition to the hypotheses of~\eqref{it:polymer-expansion-determinant-e}, that $l/L\in\mathrm{Cl}$ and $e_{k+1}=L^{\frac{1}{2}}\cdot e_k\le e_0$. Then $L\,E^{\mathrm{det}}_k\in K_{k+1}$ and $L\,\delta E^{\mathrm{det}}_k\in K^x_{k+1}$; here $L$ is the operation of Section~\ref{sec:local-parts-adjustment}, and $K_{k+1}$, $K^x_{k+1}$ are the spaces of Definition~\ref{def:renormalisation-group-1} at level $k+1$ for the side $l/L$.
\item[(g)]\itemmark{g}\label{it:polymer-expansion-determinant-g} Assume, in addition to the hypotheses of~\eqref{it:polymer-expansion-determinant-f}, that $L(\kappa-\kappa_0-1)\ge\kappa$ and that condition~\eqref{eq:renormalisation-group-24} holds. Then
\[
\left\|L\,E^{\mathrm{det}}_k\right\|_{k+1}\le\left\|L\,\delta E^{\mathrm{det}}_k\right\|^x_{k+1}\le e_{k+1}^{\frac{1}{4}-2\varepsilon}.
\]
\end{enumerate}
\end{proposition}

\begin{proof}[Proof of~\eqref{it:polymer-expansion-determinant-a}--\eqref{it:polymer-expansion-determinant-d}]
By~\eqref{eq:operators-fermion-G1}, condition~\eqref{eq:renormalisation-group-20} and $e_k\le e_0$, every $A\in\tilde{G}_k$ satisfies
\[
\mathfrak{i}(A)<e_k^{\varepsilon}\le\frac{1}{6L};
\]
in particular $\exp\left(\mathfrak{i}(A)\right)\le\exp\left(3\,\mathfrak{i}(A)\right)\le\mathbf{e}$.

By fact~\eqref{it:renormalisation-group-g0} of Section~\ref{sec:cubes-parameter-cubes}, which applies here because $L$ is odd, $M=L^m>4L$ and $l\in\mathrm{Cl}$ (Theorem~\ref{thm:expansion-walks}), the $M$-cubes form a partition of the set of the sites of $T_k$, every unit block $\Delta_{x'}$ is contained in a single $M$-cube, and so is the $L$-block $B(y')$ of unit sites for every point $y'$ of the $L$-lattice. By Definition~\ref{def:block-torus}, the block of side $L$ with centre $y'$, a set of sites of $T_k$, is the union of the unit blocks $\Delta_z$ with $z\in B(y')$. Each such site $z$ belongs both to $\Delta_z$ and to $B(y')$; as the $M$-cubes are pairwise disjoint, $\Delta_z$ therefore lies in the $M$-cube that contains $B(y')$, and so does the whole block of side $L$ with centre $y'$. As a polymer $X\in D_k$ is a union of $M$-cubes, a unit site $x'$ belongs to $X$ if and only if $\Delta_{x'}$ meets $X$, and in that case the block of side $L$ that contains $x'$ is contained in $X$.

\emph{The bound in~\eqref{it:polymer-expansion-determinant-b} for $k\ge1$.} Let $k\ge1$, $X\in D_k$, $A\in\tilde{G}_k$ and $y\ge0$, and let
\[
\mathcal{N}:=Q_k(A)\cdot S_{k,y}(X,A)\cdot Q_k^T(-A),
\]
a linear operator on the spinors on $T^0_k$, with the kernel $\mathcal{N}(x_1,x_2)$ of Section~\ref{sec:norms-pairings-kernels}. By the formulas of Section~\ref{sec:phase-factors-lattice}, the operators $Q_k(A)$ and $Q_k^T(-A)$ couple a unit site $x'$ only with the sites of $\Delta_{x'}$, and by property~\ref{it:renormalisation-group-12} the kernel of $S_{k,y}(X,A)$ vanishes unless both of its arguments lie in $X$. With the preceding paragraph it follows that $\mathcal{N}(x_1,x_2)=0$ unless $x_1\in X$ and $x_2\in X$. Let $v$ be one of the standard unit vectors of $\mathbb{C}^4$. The column $\mathcal{N}(x_1,x_2)\cdot v$ of the $4\times4$ matrix $\mathcal{N}(x_1,x_2)$ is
\[
\mathcal{N}(x_1,x_2)\cdot v=\left(Q_k(A)\cdot S_{k,y}(X,A)\cdot f\right)(x_1),\qquad f:=\boldsymbol{\phi}_A(x_2,\cdot)^{-1}\cdot\mathbf{1}_{\Delta_{x_2}}\cdot v.
\]
By Lemma~\ref{lem:bounds-phases}(a), $|f|_{\infty}\le\exp\left(\frac{3}{2}\,\mathfrak{i}(A)\right)$, and $\left|\boldsymbol{\phi}_A(x_1,x)\right|\le\exp\left(\frac{3}{2}\,\mathfrak{i}(A)\right)$ for $x\in\Delta_{x_1}$. Since $\eta^3\cdot\#|\Delta_{x_1}|=1$ and $f=\mathbf{1}_{\Delta_{x_2}}\cdot f$, property~\ref{it:renormalisation-group-7} yields
\begin{align*}
\left|\mathcal{N}(x_1,x_2)\cdot v\right|&\le\exp\left(\frac{3}{2}\,\mathfrak{i}(A)\right)\cdot\left|\mathbf{1}_{\Delta_{x_1}}\cdot S_{k,y}(X,A)\cdot\mathbf{1}_{\Delta_{x_2}}\cdot f\right|_{\infty}\\
&\le\exp\left(3\,\mathfrak{i}(A)\right)\cdot C_w\cdot\exp\left(-\gamma_w\cdot d(x_1,x_2)-\kappa_w\cdot d_M(X)\right)\le\mathbf{e}\cdot C_w\cdot\exp\left(-\kappa_w\cdot d_M(X)\right).
\end{align*}
Property~\ref{it:renormalisation-group-7} is assumed for both norms on $\mathbb{C}^4$, the Euclidean norm and $|\cdot|_c$ (Section~\ref{sec:properties-systems-primitives}). In the chain of inequalities just displayed it is applied with the Euclidean norm of Section~\ref{sec:norms-pairings-kernels}, which is in force in this section. Since a standard unit vector $v$ has norm $1$ for both norms, the chain holds equally when $|\cdot|$ is replaced by $|\cdot|_c$ and the sup norms by those formed with $|\cdot|_c$. Every entry of a vector of $\mathbb{C}^4$ is bounded in modulus by the norm of the vector, for either norm; and as $v$ ranges over the standard unit vectors, $\mathcal{N}(x_1,x_2)\cdot v$ ranges over the columns of $\mathcal{N}(x_1,x_2)$. Therefore every entry of $\mathcal{N}(x_1,x_2)$ has modulus at most $\mathbf{e}\cdot C_w\cdot\exp\left(-\kappa_w\cdot d_M(X)\right)$.

Let $P:=P(A)$. By Lemma~\ref{lem:block-averages-block}, $P$ is a projection that commutes with the constant $4\times4$ matrices, and $\left(b_k+i\gamma_3y\right)^{-1}$, $\left(b_k+b_L+i\gamma_3y\right)^{-1}$ are functions of $\gamma_3$. The definition of $B_{k,y}(A)$ in Section~\ref{sec:operators-unit-lattice} gives therefore, as in the proof of Proposition~\ref{prop:kernels-unit-lattice}(b),
\[
B_{k,y}(A)\cdot\gamma_3\cdot B_{k,y}(A)=\gamma_3\cdot\left[\left(b_k+i\gamma_3y\right)^{-2}\cdot(1-P)+\left(b_k+b_L+i\gamma_3y\right)^{-2}\cdot P\right].
\]
Since $\gamma_3$ is Hermitian with $\gamma_3^2=1$ (Section~\ref{sec:parameters-clifford-relations}), the matrices $\left(b_k+i\gamma_3y\right)^{-2}$ and $\left(b_k+b_L+i\gamma_3y\right)^{-2}$ act on the eigenspaces of $\gamma_3$ for the eigenvalues $\pm1$, which are orthogonal, as the multiplications by $\left(b_k\pm iy\right)^{-2}$ and by $\left(b_k+b_L\pm iy\right)^{-2}$; these numbers have modulus at most $\left(b_k^2+y^2\right)^{-1}$, and hence so have the norms of the two matrices. Moreover $|\gamma_3|=1$, and by Proposition~\ref{prop:kernels-unit-lattice}(a) the kernel $P(x_2,x_1)$ vanishes unless the unit sites $x_1$, $x_2$ lie in one block of side $L$, and $\left|P(x_2,x_1)\right|\le\mathbf{e}\cdot L^{-3}$. Writing the right side of the last display as
\[
\gamma_3\cdot\left(b_k+i\gamma_3y\right)^{-2}+\gamma_3\cdot\left[\left(b_k+b_L+i\gamma_3y\right)^{-2}-\left(b_k+i\gamma_3y\right)^{-2}\right]\cdot P,
\]
one obtains for all unit sites $x_1$, $x_2$
\[
\left|\left(B_{k,y}(A)\cdot\gamma_3\cdot B_{k,y}(A)\right)(x_2,x_1)\right|\le\left(b_k^2+y^2\right)^{-1}\cdot\left(\mathbf{1}[x_1=x_2]+2\mathbf{e}\cdot L^{-3}\cdot\mathbf{1}[x_1,x_2\text{ lie in one block of side }L]\right),
\]
where $\mathbf{1}[\cdot]$ equals $1$ if the condition in the brackets holds and $0$ otherwise.

For complex $4\times4$ matrices $m$ and $n$, the trace $\operatorname{tr}(m\cdot n)$ is a sum of $16$ products of an entry of $m$ and an entry of $n$, and every entry of $m$ has modulus at most $|m|$; hence $\left|\operatorname{tr}(m\cdot n)\right|$ is at most $16\,|m|$ times the largest modulus of an entry of $n$. The trace $\operatorname{Tr}$ is invariant under cyclic permutations of the factors. With the kernels of Section~\ref{sec:norms-pairings-kernels}, and since only unit sites $x_1,x_2\in X$ contribute,
\[
\operatorname{Tr}\left[\gamma_3\cdot B_{k,y}(A)\cdot\mathcal{N}\cdot B_{k,y}(A)\right]=\operatorname{Tr}\left[B_{k,y}(A)\cdot\gamma_3\cdot B_{k,y}(A)\cdot\mathcal{N}\right]=\sum_{x_1,x_2\in X}\operatorname{tr}\left[\left(B_{k,y}(A)\cdot\gamma_3\cdot B_{k,y}(A)\right)(x_2,x_1)\cdot\mathcal{N}(x_1,x_2)\right].
\]
For fixed $x_1$, the block of side $L$ that contains $x_1$ has $L^3$ unit sites. The three preceding estimates therefore give
\[
\left|\operatorname{Tr}\left[\gamma_3\cdot B_{k,y}(A)\cdot\mathcal{N}\cdot B_{k,y}(A)\right]\right|\le\#|X\cap T^0_k|\cdot16\left(1+2\mathbf{e}\right)\cdot\left(b_k^2+y^2\right)^{-1}\cdot\mathbf{e}\cdot C_w\cdot\exp\left(-\kappa_w\cdot d_M(X)\right).
\]
Further,
\[
\int_0^{\infty}\frac{\mathrm{d}y}{b_k^2+y^2}=\frac{\pi}{2b_k},
\]
and, as every $M$-cube contains $M^3$ unit sites, Proposition~\ref{prop:polymers}\eqref{it:polymers-t1} and the inequality $d+1\le\exp(d)$ for $d\ge0$ give
\[
\#|X\cap T^0_k|=M^3\cdot|X|_M\le M^3\cdot c_X\cdot\left(d_M(X)+1\right)\le M^3\cdot c_X\cdot\exp\left(d_M(X)\right).
\]
Finally $\kappa_w-1=\kappa+\kappa_0+2\ge\kappa+1$, as $\kappa_0>0$, and $b_k\le b$ (Section~\ref{sec:parameters-clifford-relations}). The integrand in the definition of $\acute{E}_k(X,A)$ is $\operatorname{Tr}\left[\gamma_3\cdot B_{k,y}(A)\cdot\mathcal{N}\cdot B_{k,y}(A)\right]$. Hence the integral converges absolutely, and
\begin{align*}
\left|\acute{E}_k(X,A)\right|&\le b_k^2\cdot\frac{\pi}{2b_k}\cdot16\left(1+2\mathbf{e}\right)\cdot\mathbf{e}\cdot C_w\cdot M^3\cdot c_X\cdot\exp\left(-(\kappa_w-1)\cdot d_M(X)\right)\\
&\le8\pi\cdot b\cdot\mathbf{e}\cdot\left(1+2\mathbf{e}\right)\cdot c_X\cdot C_w\cdot M^3\cdot\exp\left(-(\kappa+1)\cdot d_M(X)\right)\le C_d\cdot M^3\cdot\exp\left(-(\kappa+1)\cdot d_M(X)\right),
\end{align*}
the last inequality by~\eqref{eq:operators-fermion-10}.

\emph{The bound in~\eqref{it:polymer-expansion-determinant-b} for $k=0$.} Let $k=0$, $X\in D_0$, $A\in\tilde{G}_0$ and $y\ge0$. At level $0$ one has $\eta=1$, the unit lattice is $T_0$, and $F_0(A)=D_A+\mathfrak{m}+b_L\cdot P(A)$. For every site $x$,
\[
\sum_{x'}\left|F_0(A)(x,x')\right|\le\left(3+\mathfrak{m}\right)+6\exp\left(\mathfrak{i}(A)\right)+b_L\cdot\mathbf{e}\le4+6\mathbf{e}+b\cdot\mathbf{e}.
\]
Indeed, the diagonal part of $D_A+\mathfrak{m}$ is $3+\mathfrak{m}$ times the unit matrix; each of the six terms of $D_A$ that couple $x$ with a nearest neighbour has norm at most $\exp\left(\mathfrak{i}(A)\right)$, by Lemma~\ref{lem:bounds-phases}(a) and because the matrices $\Pi^{\pm}_{\mu}$ are orthogonal projections; and by Proposition~\ref{prop:kernels-unit-lattice}(a), whose proof applies without change for $k=0$, the row of $P(A)$ at $x$ has at most $L^3$ non-vanishing entries, each of modulus at most $\mathbf{e}\cdot L^{-3}$. Moreover $\mathfrak{m}\le1$, $\mathfrak{i}(A)<1$ and $b_L\le b$.

By properties~\ref{it:renormalisation-group-12} and~\ref{it:renormalisation-group-7}, the latter in the form stated in property~\ref{it:renormalisation-group-17}, the block $S_{0,y}(X,A)(x',x)$ vanishes unless $x,x'\in X$, and its entries have modulus at most $C_w\cdot(1+y)^{-1}\cdot\exp\left(-\kappa_w\cdot d_M(X)\right)$. With the estimate for $\left|\operatorname{tr}(m\cdot n)\right|$ above it follows that
\begin{align*}
\left|\operatorname{Tr}\left[F_0(A)\cdot S_{0,y}(X,A)\right]\right|&\le\#|X\cap T^0_0|\cdot16\left(4+6\mathbf{e}+b\cdot\mathbf{e}\right)\cdot C_w\cdot(1+y)^{-1}\cdot\exp\left(-\kappa_w\cdot d_M(X)\right),\\
\left|\operatorname{Tr}\left[\gamma_3\cdot S_{0,y}(X,A)\right]\right|&\le\#|X\cap T^0_0|\cdot16\,C_w\cdot\exp\left(-\kappa_w\cdot d_M(X)\right).
\end{align*}
One has
\[
\int_1^{\infty}\frac{\mathrm{d}y}{y(1+y)}=\log2\le1,
\]
and $\#|X\cap T^0_0|\le M^3\cdot c_X\cdot\exp\left(d_M(X)\right)$, as for $k\ge1$. Hence both integrals in the definition of $\acute{E}_0(X,A)$ converge absolutely, and, as $\left(4+6\mathbf{e}+b\cdot\mathbf{e}\right)\cdot\log2+1\le5+6\mathbf{e}+b\cdot\mathbf{e}$,
\[
\left|\acute{E}_0(X,A)\right|\le16\left(5+6\mathbf{e}+b\cdot\mathbf{e}\right)\cdot c_X\cdot C_w\cdot M^3\cdot\exp\left(-(\kappa_w-1)\cdot d_M(X)\right)\le C_d\cdot M^3\cdot\exp\left(-(\kappa+1)\cdot d_M(X)\right),
\]
again by~\eqref{eq:operators-fermion-10}.

\emph{Analyticity in~\eqref{it:polymer-expansion-determinant-b}.} Let $X\in D_k$. For fixed $y$, the integrands in the definition of $\acute{E}_k(X,\cdot)$ are analytic functions on $\tilde{G}_k$: the function $S_{k,y}(X,\cdot)$ is analytic on $\tilde{G}_k$ by property~\ref{it:renormalisation-group-2}, and the entries of $B_{k,y}(A)$, $Q_k(A)$, $Q_k^T(-A)$ and $F_0(A)$ are entire functions of $A$ (Sections~\ref{sec:phase-factors-lattice} and~\ref{sec:operators-unit-lattice}). For fixed $A$, they are Borel functions of $y$, since the family~\eqref{it:primitives-p5} is Borel in $y$ (Section~\ref{sec:primitive-operators}) and the other factors are continuous in $y$. By the estimates just proved, whose right sides do not depend on $A$, they are dominated, uniformly in $A\in\tilde{G}_k$, by integrable functions of $y$. By the theorems of Fubini and Morera, $\acute{E}_k(X,\cdot)$ is analytic on $\tilde{G}_k$. The maps $A\mapsto-A$ and $A\mapsto\overline{A}$ map $\tilde{G}_k$ onto itself, by~\eqref{eq:operators-fermion-G2}, and $A\mapsto\overline{f(\overline{A})}$ is analytic whenever $f$ is. Hence $E^d_k(X,\cdot)$ is analytic on $\tilde{G}_k$; as the mean of four numbers of modulus at most $C_d\cdot M^3\cdot\exp\left(-(\kappa+1)\cdot d_M(X)\right)$, it satisfies the same bound. This proves~\eqref{it:polymer-expansion-determinant-b}.

\emph{Proof of~\eqref{it:polymer-expansion-determinant-c}.} Let $A,\tilde{A}\in\tilde{G}_k$ coincide on the bonds that meet $X$. Let first $k\ge1$. By the expression of the trace obtained in the proof of~\eqref{it:polymer-expansion-determinant-b}, at $A$ as well as at $\tilde{A}$, the integrand in the definition of $\acute{E}_k(X,\cdot)$ is a sum over unit sites $x_1,x_2\in X$. For a unit site $x'$ of $X$, the entries of $Q_k(A)$, $Q_k^T(-A)$ and $B_{k,y}(A)$ that involve $x'$ depend on $A$ only through its values on bonds inside the block of side $L$ that contains $x'$ (Section~\ref{sec:phase-factors-lattice} and Lemma~\ref{lem:block-averages-block}), and this block is contained in $X$; the kernel of $S_{k,y}(X,A)$ depends on $A$ only through its values on the bonds that meet $X$, by property~\ref{it:renormalisation-group-12}. Hence, for every $y$, the integrand in the definition of $\acute{E}_k(X,\cdot)$ has the same value at $A$ and at $\tilde{A}$. Let now $k=0$. Then
\[
\operatorname{Tr}\left[F_0(A)\cdot S_{0,y}(X,A)\right]=\sum_{x,x'\in X}\operatorname{tr}\left[F_0(A)(x,x')\cdot S_{0,y}(X,A)(x',x)\right],\qquad\operatorname{Tr}\left[\gamma_3\cdot S_{0,y}(X,A)\right]=\sum_{x\in X}\operatorname{tr}\left[\gamma_3\cdot S_{0,y}(X,A)(x,x)\right],
\]
and for $x,x'\in X$ the block $F_0(A)(x,x')$ depends on $A$ only through the value of $A$ on the bond $(x,x')$, if $x$ and $x'$ are nearest neighbours, and through the values of $A$ on bonds inside the block of side $L$ that contains $x$ (Lemma~\ref{lem:block-averages-block}), which is contained in $X$; the blocks of $S_{0,y}(X,\cdot)$ are covered by property~\ref{it:renormalisation-group-12} as before. In both cases $\acute{E}_k(X,A)=\acute{E}_k(X,\tilde{A})$. The pairs $\left(-A,-\tilde{A}\right)$, $\left(\overline{A},\overline{\tilde{A}}\right)$ and $\left(-\overline{A},-\overline{\tilde{A}}\right)$ consist again of elements of $\tilde{G}_k$, by~\eqref{eq:operators-fermion-G2}, that coincide on the bonds that meet $X$; so the definition of $E^d_k$ gives $E^d_k(X,A)=E^d_k(X,\tilde{A})$.

\emph{Proof of~\eqref{it:polymer-expansion-determinant-d}.} Let $\boldsymbol{\lambda}$ be a real scalar on $T_k$, let $\boldsymbol{\lambda}^u$ be its restriction to $T^0_k$, and put $u:=\exp\left(i\,e_k\cdot\boldsymbol{\lambda}\right)$ and $u^0:=\exp\left(i\,e_k\cdot\boldsymbol{\lambda}^u\right)$, the restriction of $u$ to $T^0_k$. Then $A-d\boldsymbol{\lambda}\in\tilde{G}_k$ by~\eqref{eq:operators-fermion-G2}, and the closed real connection $d\boldsymbol{\lambda}$ and the function $u$ satisfy the hypotheses of Lemma~\ref{lem:symmetries-elementary-operators}(a). Let first $k\ge1$. Lemma~\ref{lem:symmetries-elementary-operators}(a) gives
\[
Q_k(A-d\boldsymbol{\lambda})=u^0\cdot Q_k(A)\cdot u^{-1},\qquad Q_k^T\left(-(A-d\boldsymbol{\lambda})\right)=u\cdot Q_k^T(-A)\cdot(u^0)^{-1},\qquad P(A-d\boldsymbol{\lambda})=u^0\cdot P(A)\cdot(u^0)^{-1};
\]
as the matrices $\left(b_k+i\gamma_3y\right)^{-1}$ and $\left(b_k+b_L+i\gamma_3y\right)^{-1}$ commute with the multiplication by the scalar function $u^0$, it follows that $B_{k,y}(A-d\boldsymbol{\lambda})=u^0\cdot B_{k,y}(A)\cdot(u^0)^{-1}$. Property~\ref{it:renormalisation-group-13} gives $S_{k,y}(X,A-d\boldsymbol{\lambda})=u\cdot S_{k,y}(X,A)\cdot u^{-1}$. Since $\gamma_3$ commutes with $u^0$, the product of these relations yields
\[
\gamma_3\cdot B_{k,y}(A-d\boldsymbol{\lambda})\cdot Q_k(A-d\boldsymbol{\lambda})\cdot S_{k,y}(X,A-d\boldsymbol{\lambda})\cdot Q_k^T\left(-(A-d\boldsymbol{\lambda})\right)\cdot B_{k,y}(A-d\boldsymbol{\lambda})=u^0\cdot\left[\gamma_3\cdot B_{k,y}(A)\cdot Q_k(A)\cdot S_{k,y}(X,A)\cdot Q_k^T(-A)\cdot B_{k,y}(A)\right]\cdot(u^0)^{-1},
\]
and the two traces coincide for every $y$. Let now $k=0$; then $u^0=u$. Lemma~\ref{lem:symmetries-elementary-operators}(a) and property~\ref{it:renormalisation-group-13} give $F_0(A-d\boldsymbol{\lambda})=u\cdot F_0(A)\cdot u^{-1}$ and $S_{0,y}(X,A-d\boldsymbol{\lambda})=u\cdot S_{0,y}(X,A)\cdot u^{-1}$, hence
\[
F_0(A-d\boldsymbol{\lambda})\cdot S_{0,y}(X,A-d\boldsymbol{\lambda})=u\cdot F_0(A)\cdot S_{0,y}(X,A)\cdot u^{-1},\qquad\gamma_3\cdot S_{0,y}(X,A-d\boldsymbol{\lambda})=u\cdot\gamma_3\cdot S_{0,y}(X,A)\cdot u^{-1}.
\]
In both cases $\acute{E}_k(X,A-d\boldsymbol{\lambda})=\acute{E}_k(X,A)$. One has $-(A-d\boldsymbol{\lambda})=-A-d(-\boldsymbol{\lambda})$ and $\overline{A-d\boldsymbol{\lambda}}=\overline{A}-d\boldsymbol{\lambda}$, with the real scalars $-\boldsymbol{\lambda}$ and $\boldsymbol{\lambda}$; applying the identity just proved at the four arguments $A$, $-A$, $\overline{A}$, $-\overline{A}$ of $\acute{E}_k(X,\cdot)$ in the definition of $E^d_k$, one obtains $E^d_k(X,A-d\boldsymbol{\lambda})=E^d_k(X,A)$. The relations $E^d_k(X,-A)=E^d_k(X,A)$ and $\overline{E^d_k(X,\overline{A})}=E^d_k(X,A)$ hold because the substitution $A\mapsto-A$, and the substitution $A\mapsto\overline{A}$ followed by complex conjugation, permute the four terms in the definition of $E^d_k(X,A)$.

Let $g=(S,a')\in\mathbb{G}_k^L$. The polymer $gX$ and the connection $gA$ depend on $g$ only through $r(S)$ and $a'$ (Section~\ref{sec:dirac-matrices-lattice}), and by~\eqref{eq:operators-fermion-1} there is $S'\in\mathrm{Sp}_3$ with $r(S')=r(S)$. Replacing $S$ by $S'$, one may therefore assume $S\in\mathrm{Sp}_3$; then $g\in\mathbb{G}^L_{k,3}$, and $g$ commutes with $\gamma_3$. For $k\ge1$, property~\ref{it:renormalisation-group-14} and Proposition~\ref{prop:symmetries-fermion-operators-2}(b), the latter with the same parameter $y$ on both sides because $S\in\mathrm{Sp}_3$, give
\[
S_{k,y}(gX,gA)\cdot g=g\cdot S_{k,y}(X,A),\qquad Q_k(gA)\cdot g=g\cdot Q_k(A),\qquad Q_k^T(-gA)\cdot g=g\cdot Q_k^T(-A),\qquad B_{k,y}(gA)\cdot g=g\cdot B_{k,y}(A).
\]
For $k=0$, property~\ref{it:renormalisation-group-14} gives $S_{0,y}(gX,gA)\cdot g=g\cdot S_{0,y}(X,A)$, and Proposition~\ref{prop:symmetries-fermion-operators-2}(b), applied at level $0$ to $D_A$ and $P(A)$, gives $F_0(gA)\cdot g=g\cdot F_0(A)$. In both cases each operator under a trace in the definition of $\acute{E}_k(gX,gA)$ equals $g$ times the corresponding operator for $(X,A)$ times $g^{-1}$, so that $\acute{E}_k(gX,gA)=\acute{E}_k(X,A)$. Moreover $g(-A)=-gA$ and $g(\overline{A})=\overline{gA}$, and $-A$, $\overline{A}$, $-\overline{A}$ lie in $\tilde{G}_k$; the definition of $E^d_k$ gives $E^d_k(gX,gA)=E^d_k(X,A)$.

\emph{Proof of~\eqref{it:polymer-expansion-determinant-a}.} The set $D_k$ is finite. Put
\[
\mathrm{Lg}(A):=\sum_{X\in D_k}\acute{E}_k(X,A)\qquad\left(A\in\tilde{G}_k\right);
\]
by~\eqref{it:polymer-expansion-determinant-b}, $\mathrm{Lg}$ is analytic on $\tilde{G}_k$. Let $A\in\tilde{G}_k$ be real. By property~\ref{it:renormalisation-group-1}, $S_{k,y}(A)=\sum_{X\in D_k}S_{k,y}(X,A)$ for $y\ge0$. Inserting this into Lemma~\ref{lem:logarithm-determinant-real}, and interchanging the finite sum over $X$ with the integrals over $y$, which is permitted because each term is integrable by the proof of~\eqref{it:polymer-expansion-determinant-b}, one obtains
\[
\delta Z_k(A)=\exp\left\{c^d_k\cdot n_k+\mathrm{Lg}(A)\right\},
\]
where $c^d_0=0$. The function $\delta Z_k$ is analytic on $\tilde{G}_k$, being the determinant of a matrix whose entries are analytic functions of $A$: for $k\ge1$, parts (a) and (b) of Theorem~\ref{thm:expansion-walks} give the existence of $S_k(A)=S_{\infty}(A)$ and its analyticity on $\tilde{G}_k$, and for $k=0$ the entries of $F_0(A)$ are entire. By~\eqref{eq:operators-fermion-G2}, $\tilde{G}_k$ is open and connected, and its real elements form a non-empty open subset of the space of real connections, which contains $0$. At every real point of $\tilde{G}_k$, the two analytic functions $\delta Z_k$ and $\exp\left\{c^d_k\cdot n_k+\mathrm{Lg}\right\}$ coincide on a neighbourhood of that point in the space of real connections; hence their partial derivatives of all orders with respect to the real coordinates coincide at that point, and these are their complex partial derivatives, so that the two functions have the same Taylor coefficients there. By the identity theorem they coincide on $\tilde{G}_k$. This is the first identity of~\eqref{it:polymer-expansion-determinant-a}.

Let $A\in\tilde{G}_k$. Then $-A,\overline{A}\in\tilde{G}_k$ by~\eqref{eq:operators-fermion-G2}, and the number $c^d_k\cdot n_k$ is real. The first identity of~\eqref{it:polymer-expansion-determinant-a} and Proposition~\ref{prop:symmetries-fermion-operators-2}(e) give
\begin{align*}
\exp\mathrm{Lg}(-A)&=\exp\left(-c^d_k\cdot n_k\right)\cdot\delta Z_k(-A)=\exp\left(-c^d_k\cdot n_k\right)\cdot\delta Z_k(A)=\exp\mathrm{Lg}(A),\\
\exp\overline{\mathrm{Lg}(\overline{A})}&=\exp\left(-c^d_k\cdot n_k\right)\cdot\overline{\delta Z_k(\overline{A})}=\exp\left(-c^d_k\cdot n_k\right)\cdot\delta Z_k(A)=\exp\mathrm{Lg}(A).
\end{align*}
Consequently the functions $A\mapsto\mathrm{Lg}(-A)-\mathrm{Lg}(A)$ and $A\mapsto\overline{\mathrm{Lg}(\overline{A})}-\mathrm{Lg}(A)$ are continuous on $\tilde{G}_k$ and take their values in $2\pi i\,\mathbb{Z}$. Since $\tilde{G}_k$ is connected, both are constant; their values at $A=0$ are $0$ and $\overline{\mathrm{Lg}(0)}-\mathrm{Lg}(0)$. Applying the first of these facts at $A$ and at $\overline{A}$, and then the second one, one obtains for $A\in\tilde{G}_k$
\begin{align*}
\sum_{X\in D_k}E^d_k(X,A)&=\frac{1}{4}\left[\mathrm{Lg}(A)+\mathrm{Lg}(-A)+\overline{\mathrm{Lg}(\overline{A})}+\overline{\mathrm{Lg}(-\overline{A})}\right]=\frac{1}{2}\left[\mathrm{Lg}(A)+\overline{\mathrm{Lg}(\overline{A})}\right]\\
&=\mathrm{Lg}(A)+\frac{1}{2}\left(\overline{\mathrm{Lg}(0)}-\mathrm{Lg}(0)\right),
\end{align*}
and therefore
\[
\sum_{X\in D_k}E^{\mathrm{det}}_k(X,A)=\mathrm{Lg}(A)-\mathrm{Lg}(0).
\]
The first identity of~\eqref{it:polymer-expansion-determinant-a}, at $A$ and at $0$, gives $\delta Z_k(A)/\delta Z_k(0)=\exp\left\{\mathrm{Lg}(A)-\mathrm{Lg}(0)\right\}$, which is the second identity of~\eqref{it:polymer-expansion-determinant-a}.

Clauses (d0)--(d3) of Proposition~\ref{prop:local-expansion-determinant} follow.
\end{proof}

\begin{proof}[Proof of~\eqref{it:polymer-expansion-determinant-e}--\eqref{it:polymer-expansion-determinant-g}]
\emph{Clause~\eqref{it:polymer-expansion-determinant-e}.} Let the hypotheses of~\eqref{it:polymer-expansion-determinant-e} be satisfied, and let $X\in D_k$ and $(\mathscr{A},\mathscr{C})\in\tilde{R}_k\times\mathfrak{X}_k$. By~\eqref{eq:operators-fermion-G3}, which holds under condition~\eqref{eq:renormalisation-group-23}, the connections $\mathscr{C}$ and $\mathscr{A}+\mathscr{C}$ lie in $\tilde{G}_k$. Since $e_k\le e_0$, condition~\eqref{eq:renormalisation-group-23} gives $e_k^{\varepsilon}\le\frac{1}{4}$. Put $n:=d_M(X)$, and write $\mathscr{A}=A_0+A_1$ as in Definition~\ref{def:domains}. Thus $A_1\in R_k$, so that $|A_1|_{\infty}<e_k^{-\frac{3}{4}+\varepsilon}$; the connection $A_0$ is real; and for every $M$-cube $\square$ there are a real function $\boldsymbol{\lambda}_{\square}$ on the sites and a real connection $A'_{\square}$ on the bonds of $\square^{\dagger}$ such that $A_0=A'_{\square}+d\boldsymbol{\lambda}_{\square}$ on $\square^{\dagger}$ and the three norms of $A'_{\square}$, formed on $\square^{\dagger}$, are below the thresholds of $R_k$; in particular $|\partial A'_{\square}|_{\infty}<e_k^{-\frac{3}{4}+2\varepsilon}$. Moreover $|\mathscr{C}|_{\infty}<e_k^{-1+2\varepsilon}$, because $\mathscr{C}\in\mathfrak{X}_k$.

\emph{Case 1: $n\le\frac{v_k}{4}$.} By the definition of $d_M$ (Section~\ref{sec:lattice-geometry}; convention~\eqref{eq:polymer-activities-6}), $M\cdot n$ is the infimum of the lengths of the nets in the continuum torus $\mathbb{R}^3/l\mathbb{Z}^3$ that meet every closed $M$-cube of $X$. By the statement that follows convention~\eqref{eq:polymer-activities-6} in Section~\ref{sec:standing-conventions}, this infimum does not change when only nets that are trees are admitted; thus $M\cdot n$ is also the infimum of the lengths of the trees in $\mathbb{R}^3/l\mathbb{Z}^3$ that meet every closed $M$-cube of $X$. Choose such a tree, of length $M\cdot n$ if the infimum is attained, and of length at most $M\cdot\left(n+\frac{1}{2}\right)$ otherwise. Being contractible, the chosen tree lifts along the covering map $\mathbb{R}^3\to\mathbb{R}^3/l\mathbb{Z}^3$; as this map is a local isometry, the lift is a tree in $\mathbb{R}^3$ whose length equals that of the chosen tree. Any two points of the lift are joined by a path in the lift whose length does not exceed the length of the tree; hence the lift lies in a box of side at most $M\cdot n$, respectively $M\cdot\left(n+\frac{1}{2}\right)$, and the closed $M$-cubes of $\mathbb{R}^3$ that meet the lift lie in a box of side at most $M\cdot(n+2)$, respectively $M\cdot\left(n+\frac{5}{2}\right)$. Every $M$-cube of $X$ is the image of one of these cubes, since the tree meets it. A bond that meets $X$ has an end point in $X$, and its other end point lies at distance at most $\eta\le1\le\frac{M}{4}$ from $X$. Consequently there is a box $\mathbb{W}$ of $\eta\mathbb{Z}^3$, that is, a product of three intervals of $\eta\mathbb{Z}$, each of length at most
\[
M\cdot\left(n+\frac{5}{2}\right)+2\eta\le s_W:=M\cdot(n+3),
\]
such that $X$ and all bonds that meet $X$ lie in the image of $\mathbb{W}$ in $T_k$. By the hypotheses of~\eqref{it:polymer-expansion-determinant-e}, $v_k\ge9$ and $l\ge M\cdot r_k$. Moreover $r_k\ge v_k^r$ (Section~\ref{sec:running-quantities-tori}), where $r\ge2$ is the integer of stage~\ref{it:order-choice-8}; as $v_k\ge1$, this gives $r_k\ge v_k^2$. Hence
\[
s_W\le M\cdot\left(\frac{v_k}{4}+3\right)<M\cdot v_k^2\le M\cdot r_k\le l,
\]
so that the projection of $\mathbb{R}^3$ onto $\mathbb{R}^3/l\mathbb{Z}^3$ is injective on $\mathbb{W}$; in what follows $\mathbb{W}$ is identified with its image in $T_k$.

Every plaquette $p$ of $T_k$ lies in $\square^{\dagger}$ for some $M$-cube $\square$, and there $dA_0=dA'_{\square}$, because $d\,d\boldsymbol{\lambda}_{\square}=0$. Since $|(dB)(p)|\le2|\partial B|_{\infty}$ for every connection $B$ (Section~\ref{sec:domain-complex-connections}), it follows that
\[
|dA_0|_{\infty}<2e_k^{-\frac{3}{4}+2\varepsilon}.
\]
Let $\boldsymbol{\lambda}$ be the function of Lemma~\ref{lem:gauge-fixing-along} for the restriction of $A_0$ to the bonds of $\mathbb{W}$, formed with a site $o$ of $\mathbb{W}$; it is real, because $A_0$ is real. Continue $\boldsymbol{\lambda}$ by $0$ to a real scalar on $T_k$, and put $A'':=A_0-d\boldsymbol{\lambda}$. By Lemma~\ref{lem:gauge-fixing-along}, for every bond $b$ of $\mathbb{W}$,
\[
|A''(b)|\le2s_W\cdot\sup_p\left|(dA_0)(p)\right|\le4s_W\cdot e_k^{-\frac{3}{4}+2\varepsilon}\le\frac{1}{2}e_k^{-\frac{3}{4}+\varepsilon};
\]
the last inequality is $8s_W\cdot e_k^{\varepsilon}\le1$, which follows from $s_W\le M\cdot\left(\frac{v_k}{4}+3\right)$ and from the hypothesis $8M\cdot\left(\frac{v_k}{4}+4\right)\cdot e_k^{\varepsilon}\le1$.

Let $\mathfrak{B}$ be the connection on $T_k$ with $\mathfrak{B}:=A''+A_1$ on the bonds that meet $X$ and $\mathfrak{B}:=0$ on all other bonds. The bonds that meet $X$ are bonds of $\mathbb{W}$; hence
\[
|\mathfrak{B}|_{\infty}<\left(\frac{1}{2}+1\right)e_k^{-\frac{3}{4}+\varepsilon}<2e_k^{-\frac{3}{4}+\varepsilon}.
\]
Put $\rho:=\frac{1}{4}e_k^{-\frac{1}{4}}$. Since $v_k=-\log e_k$ (Section~\ref{sec:running-quantities-tori}) and $v_k\ge9>4\log8$, one has $\rho=\frac{1}{4}\exp\left(\frac{v_k}{4}\right)>\frac{1}{4}\exp(\log8)=2$. For every complex number $s$ with $|s|\le\rho$ the bound $e_k^{\varepsilon}\le\frac{1}{4}$ gives
\[
\left|s\cdot\mathfrak{B}+\mathscr{C}\right|_{\infty}<\frac{1}{2}e_k^{-1+\varepsilon}+e_k^{-1+2\varepsilon}\le\frac{3}{4}e_k^{-1+\varepsilon},
\]
so that $s\cdot\mathfrak{B}+\mathscr{C}\in G_k\subset\tilde{G}_k$ by~\eqref{eq:operators-fermion-G2}; since $1<\rho$, the case $s=1$ gives $\mathfrak{B}+\mathscr{C}\in\tilde{G}_k$. Also by~\eqref{eq:operators-fermion-G2}, the connection $\mathscr{A}+\mathscr{C}-d\boldsymbol{\lambda}$ lies in $\tilde{G}_k$; on the bonds that meet $X$ it coincides with $\mathfrak{B}+\mathscr{C}$, because $\mathscr{A}-d\boldsymbol{\lambda}=A''+A_1$ there. Therefore~\eqref{it:polymer-expansion-determinant-d} and~\eqref{it:polymer-expansion-determinant-c} give
\[
E^d_k(X,\mathscr{A}+\mathscr{C})=E^d_k(X,\mathscr{A}+\mathscr{C}-d\boldsymbol{\lambda})=E^d_k(X,\mathfrak{B}+\mathscr{C}).
\]
The set $U:=\left\{s\in\mathbb{C}:\ s\cdot\mathfrak{B}+\mathscr{C}\in G_k\right\}$ is open and contains the closed disc $\left\{s\in\mathbb{C}:\ |s|\le\rho\right\}$. The function
\[
f(s):=E^d_k\left(X,s\cdot\mathfrak{B}+\mathscr{C}\right)\qquad(s\in U)
\]
is the composition of an affine map with $E^d_k(X,\cdot)$, which is analytic on $\tilde{G}_k$ by~\eqref{it:polymer-expansion-determinant-b}; so $f$ is analytic on $U$, and~\eqref{it:polymer-expansion-determinant-b} gives
\[
|f(s)|\le C_d\cdot M^3\cdot\exp\left(-(\kappa+1)n\right)\qquad(s\in U).
\]
As $\rho>2$, the circle $|s|=\rho$ encloses the points $0$ and $1$. By Cauchy's integral formula,
\[
f(1)-f(0)=\frac{1}{2\pi i}\oint_{|s|=\rho}\mathrm{d}s\,f(s)\left(\frac{1}{s-1}-\frac{1}{s}\right)=\frac{1}{2\pi i}\oint_{|s|=\rho}\mathrm{d}s\,\frac{f(s)}{s(s-1)}.
\]
On this circle $|s-1|\ge\rho-1\ge\frac{\rho}{2}$. Hence
\begin{align*}
|f(1)-f(0)|&\le\frac{1}{\rho-1}\sup_{|s|=\rho}|f(s)|\le\frac{2}{\rho}\sup_{|s|=\rho}|f(s)|=8e_k^{\frac{1}{4}}\cdot\sup_{|s|=\rho}|f(s)|\\
&\le8C_d\cdot M^3\cdot e_k^{\frac{1}{4}}\cdot\exp\left(-(\kappa+1)n\right).
\end{align*}
As $f(1)-f(0)=E^d_k(X,\mathscr{A}+\mathscr{C})-E^d_k(X,\mathscr{C})=\delta E^{\mathrm{det}}_k(X;\mathscr{A},\mathscr{C})$ and $\exp\left(-(\kappa+1)n\right)\le\exp(-\kappa\cdot n)$, the inequality of~\eqref{it:polymer-expansion-determinant-e} holds in this case.

\emph{Case 2: $n>\frac{v_k}{4}$.} Then $\exp(-n)<\exp\left(-\frac{v_k}{4}\right)=e_k^{\frac{1}{4}}$. The two terms of
\[
\delta E^{\mathrm{det}}_k(X;\mathscr{A},\mathscr{C})=E^d_k(X,\mathscr{A}+\mathscr{C})-E^d_k(X,\mathscr{C})
\]
are bounded by~\eqref{it:polymer-expansion-determinant-b}, as $\mathscr{A}+\mathscr{C}$ and $\mathscr{C}$ lie in $\tilde{G}_k$; hence
\[
\left|\delta E^{\mathrm{det}}_k(X;\mathscr{A},\mathscr{C})\right|\le2C_d\cdot M^3\cdot\exp(-n)\cdot\exp(-\kappa\cdot n)\le2C_d\cdot M^3\cdot e_k^{\frac{1}{4}}\cdot\exp(-\kappa\cdot n).
\]
This proves the inequality of~\eqref{it:polymer-expansion-determinant-e} in both cases. As $0\in\mathfrak{X}_k$ and $\delta E^{\mathrm{det}}_k(X;\mathscr{A},0)=E^{\mathrm{det}}_k(X,\mathscr{A})$, the same bound holds for $\left|E^{\mathrm{det}}_k(X,\mathscr{A})\right|$.

\emph{Clauses~\eqref{it:polymer-expansion-determinant-f} and~\eqref{it:polymer-expansion-determinant-g}.} Let, in addition, $l/L\in\mathrm{Cl}$ and $e_{k+1}=L^{\frac{1}{2}}\cdot e_k\le e_0$. Let $T_{k+1}$ be the fine lattice of level $k+1$ and side $l/L$, let $D_{k+1}$ be the set of its polymers, and let $\tilde{R}_{k+1}$ and $\mathfrak{X}_{k+1}$ be formed on $T_{k+1}$ with the charge $e_{k+1}$. For $Y\in D_{k+1}$ and $(\mathscr{A},\mathscr{C})\in\tilde{R}_{k+1}\times\mathfrak{X}_{k+1}$ the definition of the map $L$ (Section~\ref{sec:local-parts-adjustment}) gives
\[
(L\,\delta E^{\mathrm{det}}_k)(Y;\mathscr{A},\mathscr{C})=\sum_{X\in D_k:\ \bar{X}^L=LY}\delta E^{\mathrm{det}}_k\left(X;\mathscr{A}_L,\mathscr{C}_L\right),
\]
where $\mathscr{A}_L(b):=L^{-\frac{1}{2}}\cdot\mathscr{A}(b/L)$ for the bonds $b$ of $T_k$, and $\mathscr{C}_L$ is defined in the same way. Thus $L\,\delta E^{\mathrm{det}}_k$ is a family of complex numbers, that is, of activities without fermion legs, and $L\,E^{\mathrm{det}}_k$ is its restriction to $\mathscr{C}=0$, since $\delta E^{\mathrm{det}}_k(X;\mathscr{A},0)=E^{\mathrm{det}}_k(X,\mathscr{A})$. The proofs of parts (a) and (c) of Lemma~\ref{lem:scaling-domains} use, apart from the definitions of the sets involved, only fact~\eqref{it:renormalisation-group-g0} of Section~\ref{sec:cubes-parameter-cubes}, the relations $l,l/L\in\mathrm{Cl}$ and $e_{k+1}=L^{\frac{1}{2}}\cdot e_k$, and the inequality $\varepsilon\le\frac{1}{4}$ (each exponent of $L$ that these two proofs bound by $-\frac{3}{4}$ is at most $-\frac{3}{4}$ as soon as $\varepsilon\le\frac{1}{4}$); condition~\eqref{eq:renormalisation-group-38} and the inequality $L^{-\frac{3}{4}}\le\frac{1}{2}$ of stage~\ref{it:order-choice-4} enter only the proof of part (b) of that lemma. Here fact~\eqref{it:renormalisation-group-g0} applies, since $L$ is odd, $M=L^m>4L$ and $l\in\mathrm{Cl}$ (Theorem~\ref{thm:expansion-walks}); the two relations hold by assumption; and $\varepsilon\le\bar{\varepsilon}=\frac{1}{200}<\frac{1}{4}$ (stage~\ref{it:order-choice-1}). Hence part (a) of that lemma, with $c=1$, and part (c) give
\[
\mathscr{A}_L\in L^{-\frac{3}{4}}\tilde{R}_k\subset\tilde{R}_k,\qquad\mathscr{C}_L\in\mathfrak{X}_k,
\]
the inclusion because $L^{-\frac{3}{4}}\le1$. Consequently~\eqref{it:polymer-expansion-determinant-b}--\eqref{it:polymer-expansion-determinant-e} apply to every term of the sum, and $\mathscr{A}_L+\mathscr{C}_L$ and $\mathscr{C}_L$ lie in $\tilde{G}_k$ by~\eqref{eq:operators-fermion-G3}. The properties~\eqref{it:renormalisation-group-1-k1}--\eqref{it:renormalisation-group-1-k7} of Definition~\ref{def:renormalisation-group-1}, at the level $k+1$, are verified as follows.
\begin{itemize}
\item Property~\eqref{it:renormalisation-group-1-k1}. The maps $(\mathscr{A},\mathscr{C})\mapsto\mathscr{A}_L+\mathscr{C}_L$ and $\mathscr{C}\mapsto\mathscr{C}_L$ are linear, with values in $\tilde{G}_k$, and $E^d_k(X,\cdot)$ is analytic on $\tilde{G}_k$ by~\eqref{it:polymer-expansion-determinant-b}. So every term, and therefore the finite sum, is analytic in $(\mathscr{A},\mathscr{C})$ on $\tilde{R}_{k+1}\times\mathfrak{X}_{k+1}$.
\item Property~\eqref{it:renormalisation-group-1-k2}. A bond $b$ of $T_k$ that meets $X$ has an end point in $X\subset\bar{X}^L=LY$, so that the bond $b/L$ of $T_{k+1}$ meets $Y$. By~\eqref{it:polymer-expansion-determinant-c} the term of $X$ depends only on the numbers $\mathscr{A}_L(b)=L^{-\frac{1}{2}}\cdot\mathscr{A}(b/L)$ and $\mathscr{C}_L(b)=L^{-\frac{1}{2}}\cdot\mathscr{C}(b/L)$ at the bonds $b$ that meet $X$; hence the sum depends only on the values of $\mathscr{A}$ and $\mathscr{C}$ at the bonds that meet $Y$. The condition on the fermion legs is void.
\item Property~\eqref{it:renormalisation-group-1-k3}. Let $g=(S,a')\in\mathbb{G}_k^L$ (Section~\ref{sec:dirac-matrices-lattice}), with $a'\in LM\mathbb{Z}^3$. By~\eqref{it:polymer-expansion-determinant-d}, $E^d_k(gX,gB)=E^d_k(X,B)$ for $B\in\tilde{G}_k$; hence $\delta E^{\mathrm{det}}_k(gX;gB,gB')=\delta E^{\mathrm{det}}_k(X;B,B')$ whenever $B+B'$ and $B'$ lie in $\tilde{G}_k$. The map $x\mapsto r(S)x+a'$ maps $LM\mathbb{Z}^3$ and $M\mathbb{Z}^3$ onto themselves; so it maps the $LM$-cubes centred at $LM\mathbb{Z}^3$ onto such cubes and the $M$-cubes onto $M$-cubes, and $g$ commutes with the map $X\mapsto\bar{X}^L$. Therefore $X\mapsto gX$ maps the set of the $X\in D_k$ with $\bar{X}^L=LY$ bijectively onto the set of the $X\in D_k$ with $\bar{X}^L=g(LY)$. Every element of $\mathbb{G}_{k+1}$ is of the form $g'=(S,a'/L)$ with such a $g$, and $g'^{-1}(x/L)=g^{-1}(x)/L$ for $x\in\mathbb{R}^3$; hence
\[
(g'\mathscr{A})_L=g(\mathscr{A}_L),\qquad(g'\mathscr{C})_L=g(\mathscr{C}_L),\qquad L(g'Y)=g(LY).
\]
It follows that
\begin{align*}
(L\,\delta E^{\mathrm{det}}_k)(g'Y;g'\mathscr{A},g'\mathscr{C})&=\sum_{X:\ \bar{X}^L=g(LY)}\delta E^{\mathrm{det}}_k\left(X;g(\mathscr{A}_L),g(\mathscr{C}_L)\right)\\
&=\sum_{X:\ \bar{X}^L=LY}\delta E^{\mathrm{det}}_k\left(gX;g(\mathscr{A}_L),g(\mathscr{C}_L)\right)=(L\,\delta E^{\mathrm{det}}_k)(Y;\mathscr{A},\mathscr{C}).
\end{align*}
\item Property~\eqref{it:renormalisation-group-1-k4}. Let also $(\mathscr{A}^1,\mathscr{C}^1)\in\tilde{R}_{k+1}\times\mathfrak{X}_{k+1}$, and let $\boldsymbol{\lambda}$ and $\boldsymbol{\lambda}'$ be real functions on the sites of $T_{k+1}$ such that $\mathscr{A}^1=\mathscr{A}-d\boldsymbol{\lambda}$ and $\mathscr{C}^1=\mathscr{C}-d\boldsymbol{\lambda}'$ on the bonds at distance at most $1$ from $Y$; functions given only near $Y$ are first continued in any way to real functions on all sites of $T_{k+1}$. For the sites $x$ of $T_k$ put $\boldsymbol{\lambda}_L(x):=L^{\frac{1}{2}}\cdot\boldsymbol{\lambda}(x/L)$ and $\boldsymbol{\lambda}'_L(x):=L^{\frac{1}{2}}\cdot\boldsymbol{\lambda}'(x/L)$. As $\eta_k=L\cdot\eta_{k+1}$, one has $(d\boldsymbol{\lambda})_L=d(\boldsymbol{\lambda}_L)$ and $(d\boldsymbol{\lambda}')_L=d(\boldsymbol{\lambda}'_L)$. A bond $b$ of $T_k$ that meets $X$ has an end point in $LY$; so $b/L$ has an end point in $Y$, its other end point lies at distance at most $\eta_{k+1}\le1$ from $Y$, and $b/L$ is at distance at most $1$ from $Y$. Hence the connection $\mathscr{A}_L+\mathscr{C}_L-d(\boldsymbol{\lambda}_L+\boldsymbol{\lambda}'_L)$, which lies in $\tilde{G}_k$ by~\eqref{eq:operators-fermion-G2}, coincides with $\mathscr{A}^1_L+\mathscr{C}^1_L$ on the bonds that meet $X$, and $\mathscr{C}_L-d\boldsymbol{\lambda}'_L\in\tilde{G}_k$ coincides with $\mathscr{C}^1_L$ there. By~\eqref{it:polymer-expansion-determinant-c} and~\eqref{it:polymer-expansion-determinant-d},
\[
E^d_k\left(X,\mathscr{A}^1_L+\mathscr{C}^1_L\right)=E^d_k\left(X,\mathscr{A}_L+\mathscr{C}_L-d(\boldsymbol{\lambda}_L+\boldsymbol{\lambda}'_L)\right)=E^d_k\left(X,\mathscr{A}_L+\mathscr{C}_L\right),
\]
and in the same way $E^d_k(X,\mathscr{C}^1_L)=E^d_k(X,\mathscr{C}_L)$. So the terms of the sum at $(\mathscr{A}^1,\mathscr{C}^1)$ and at $(\mathscr{A},\mathscr{C})$ coincide; as there are no fermion legs, this is property~\eqref{it:renormalisation-group-1-k4}.
\item Properties~\eqref{it:renormalisation-group-1-k5} and~\eqref{it:renormalisation-group-1-k7}. By~\eqref{it:polymer-expansion-determinant-d}, $E^d_k(X,-B)=E^d_k(X,B)$ and $\overline{E^d_k(X,\overline{B})}=E^d_k(X,B)$ for $B\in\tilde{G}_k$; together they give $\overline{E^d_k(X,-\overline{B})}=E^d_k(X,B)$. Since $(-\mathscr{A})_L=-\mathscr{A}_L$ and $(\overline{\mathscr{A}})_L=\overline{\mathscr{A}_L}$, and likewise for $\mathscr{C}$, every term of the sum satisfies
\begin{align*}
\delta E^{\mathrm{det}}_k\left(X;-\mathscr{A}_L,-\mathscr{C}_L\right)&=\delta E^{\mathrm{det}}_k\left(X;\mathscr{A}_L,\mathscr{C}_L\right),\\
\overline{\delta E^{\mathrm{det}}_k\left(X;-\overline{\mathscr{A}_L},-\overline{\mathscr{C}_L}\right)}&=\delta E^{\mathrm{det}}_k\left(X;\mathscr{A}_L,\mathscr{C}_L\right).
\end{align*}
For a family without fermion legs these two relations are properties~\eqref{it:renormalisation-group-1-k5} and~\eqref{it:renormalisation-group-1-k7}.
\item Property~\eqref{it:renormalisation-group-1-k6}. For every $X$ with $\bar{X}^L=LY$, \eqref{it:polymer-expansion-determinant-e} gives
\[
\left|\delta E^{\mathrm{det}}_k(X;\mathscr{A}_L,\mathscr{C}_L)\right|\le8C_d\cdot M^3\cdot e_k^{\frac{1}{4}}\cdot\exp\left(-\kappa\cdot d_M(X)\right).
\]
Since $\kappa\ge7\kappa_0+7\ge\kappa_0$ (stage~\ref{it:order-choice-3}) and $l,l/L\in\mathrm{Cl}$, inequality~\eqref{it:polymers-t8} of Proposition~\ref{prop:polymers}, with $\kappa'=\kappa$, gives
\[
\sum_{X:\ \bar{X}^L=LY}\exp\left(-\kappa\cdot d_M(X)\right)\le c_{\mathrm{rb}}\cdot L^3\cdot\exp\left(-L(\kappa-\kappa_0-1)\cdot d_M(Y)\right).
\]
Hence, for all $Y\in D_{k+1}$ and $(\mathscr{A},\mathscr{C})\in\tilde{R}_{k+1}\times\mathfrak{X}_{k+1}$,
\begin{align*}
&\exp\left(\kappa\cdot d_M(Y)\right)\cdot\left|(L\,\delta E^{\mathrm{det}}_k)(Y;\mathscr{A},\mathscr{C})\right|\\
&\qquad\le8C_d\cdot M^3\cdot e_k^{\frac{1}{4}}\cdot c_{\mathrm{rb}}\cdot L^3\cdot\exp\left(-\left(L(\kappa-\kappa_0-1)-\kappa\right)\cdot d_M(Y)\right).
\end{align*}
For a family of numbers the norm $\|\cdot\|_{h_{k+1}}$ of Section~\ref{sec:grassmann-algebras-norms} is the modulus. The right-hand side does not depend on $(\mathscr{A},\mathscr{C})$, and $D_{k+1}$ is a finite set; so the supremum in property~\eqref{it:renormalisation-group-1-k6} is finite.
\end{itemize}
Thus $L\,\delta E^{\mathrm{det}}_k$ lies in $K^x_{k+1}$, and therefore its restriction $L\,E^{\mathrm{det}}_k$ to $\mathscr{C}=0$ lies in $K_{k+1}$ (Section~\ref{sec:spaces-activities}). This proves~\eqref{it:polymer-expansion-determinant-f}.

Assume moreover that $L(\kappa-\kappa_0-1)\ge\kappa$ and that condition~\eqref{eq:renormalisation-group-24} holds. Then the exponential factor on the right-hand side of the last inequality is at most $1$. Condition~\eqref{eq:renormalisation-group-24} and $e_k\le e_0$ give $8c_{\mathrm{rb}}\cdot L^3\cdot C_d\cdot M^3\cdot e_k^{2\varepsilon}\le1$; moreover $e_k\le e_{k+1}\le e_0\le1$, the last inequality by condition~\eqref{eq:renormalisation-group-23}, and $\frac{1}{4}-2\varepsilon>0$. It follows that
\[
\exp\left(\kappa\cdot d_M(Y)\right)\cdot\left|(L\,\delta E^{\mathrm{det}}_k)(Y;\mathscr{A},\mathscr{C})\right|\le8c_{\mathrm{rb}}\cdot L^3\cdot C_d\cdot M^3\cdot e_k^{\frac{1}{4}}\le e_k^{\frac{1}{4}-2\varepsilon}\le e_{k+1}^{\frac{1}{4}-2\varepsilon}.
\]
Taking the supremum over $(\mathscr{A},\mathscr{C})$ and over $Y$ gives $\left\|L\,\delta E^{\mathrm{det}}_k\right\|^x_{k+1}\le e_{k+1}^{\frac{1}{4}-2\varepsilon}$, and $\left\|L\,E^{\mathrm{det}}_k\right\|_{k+1}\le\left\|L\,\delta E^{\mathrm{det}}_k\right\|^x_{k+1}$, since the norm of the restriction to $\mathscr{C}=0$ does not exceed the norm of the family (Section~\ref{sec:spaces-activities}). This proves~\eqref{it:polymer-expansion-determinant-g}.

Clauses (d4)--(d6) of Proposition~\ref{prop:local-expansion-determinant} follow.
\end{proof}

\subsection{Fermion operators and their weakened families}\label{sec:fermion-operators-weakened}

\begin{theorem}[Fermion operators]\label{thm:fermion-operators}
Let $b>0$ and $\alpha\in(0,1)$, and let $L\ge5$ be odd. Put $\mathfrak{m}_0:=1$ and $\alpha^F:=\frac{1}{54}$, and let $\gamma^F:=\gamma$ be the number defined in~\eqref{eq:operators-fermion-4}; it is a function of $(b,\alpha,L)$, and $0<\gamma^F\le\frac{1}{6L}$. Let $\alpha_0\in(0,\alpha^F]$ and $\kappa'>0$, let $M^F=M^F(b,\alpha,L,\alpha_0,\kappa')$ be the number of condition~\eqref{eq:renormalisation-group-19}, and let $M=L^m\ge M^F$. Put
\[
C^F:=2\max\left\{C_1,\ C_{\mathrm{By}},\ 304\mathbf{e}\cdot\max(1,b)\cdot a,\ C_{\mathrm{Gam}},\ \mathbf{e}\cdot b^2\cdot(32a)^2\right\},
\]
where $C_1$ is the constant of Theorem~\ref{thm:properties-weakened-families}, $C_{\mathrm{By}}$ and $C_{\mathrm{Gam}}$ are the constants~\eqref{eq:operators-fermion-8}, and $a$ is the constant of~\eqref{eq:operators-fermion-4}. Let $\varepsilon\in(0,\bar{\varepsilon}]$, and let $e^F_0$ be the largest number $t\in(0,1]$ such that conditions~\eqref{eq:renormalisation-group-20}--\eqref{eq:renormalisation-group-23}, and condition~\eqref{eq:renormalisation-group-24} with $C_w=C^F$, hold with $t$ in place of $e_0$. Then for every $k\ge0$, $l\in\mathrm{Cl}$, $\mathfrak{m}\in(0,\mathfrak{m}_0]$ and $e_k\in(0,e^F_0]$ the following assertions hold.
\begin{enumerate}
\item[(i)] Propositions~\ref{prop:decay-fine-propagator}--\ref{prop:bounds} hold with $(C_F,\gamma_F)=(C^F,\gamma^F)$; Proposition~\ref{prop:left-inverse} holds, its constant being the number $C_T$ of Proposition~\ref{prop:identities-bounds}(d); and Propositions~\ref{prop:symmetries-fermion-operators} and~\ref{prop:coverings-tori} hold.
\item[(ii)] The families of Definition~\ref{def:weakened-families-level} if $k\ge1$, and those of Definition~\ref{def:families-level-zero} if $k=0$, form the families~\eqref{it:primitives-p3}, \eqref{it:primitives-p4} and~\eqref{it:primitives-p5} of a system of primitives of level $k$ on $T_k$, with the properties~\ref{it:renormalisation-group-1}, \ref{it:renormalisation-group-2} and~\ref{it:renormalisation-group-5}--\ref{it:renormalisation-group-17} for the parameters $(\alpha_0,\gamma^F,\kappa';C^F)$.
\item[(iii)] Proposition~\ref{prop:polymer-expansion-determinant} holds with $e^F_0$ in place of $e_0$ throughout, that is, in the hypotheses of Theorem~\ref{thm:expansion-walks}, in conditions~\eqref{eq:renormalisation-group-23} and~\eqref{eq:renormalisation-group-24}, and in the hypothesis $e_{k+1}\le e_0$ of clause~\eqref{it:polymer-expansion-determinant-f}; with this replacement, the hypotheses of Theorem~\ref{thm:expansion-walks} and condition~\eqref{eq:renormalisation-group-23} are satisfied, and so is the inequality $L(\kappa-\kappa_0-1)\ge\kappa$ required in clause~\eqref{it:polymer-expansion-determinant-g}. Explicitly, let $\gamma_w>0$ and $C_w>0$ be arbitrary, let a family~\eqref{it:primitives-p5} of level $k$ on $T_k$ have the properties required in Proposition~\ref{prop:polymer-expansion-determinant} for the parameter set $(\alpha_0,\gamma_w,\kappa_w;C_w)$, where $\kappa_w=\kappa+\kappa_0+3$, and let $C_d$ be the number~\eqref{eq:operators-fermion-10} formed with this $C_w$. Then clauses~\eqref{it:polymer-expansion-determinant-a}--\eqref{it:polymer-expansion-determinant-d} of Proposition~\ref{prop:polymer-expansion-determinant} hold; clause~\eqref{it:polymer-expansion-determinant-e} holds if $l\ge M\cdot r_k$, $v_k\ge9$ and $8M\cdot\left(\frac{v_k}{4}+4\right)\cdot e_k^{\varepsilon}\le1$; clause~\eqref{it:polymer-expansion-determinant-f} holds if moreover $l/L\in\mathrm{Cl}$ and $e_{k+1}=L^{\frac{1}{2}}\cdot e_k\le e^F_0$; and clause~\eqref{it:polymer-expansion-determinant-g} holds if moreover
\[
8c_{\mathrm{rb}}\cdot L^3\cdot C_d\cdot M^3\cdot\left(e^F_0\right)^{2\varepsilon}\le1.
\]
If $\kappa'\ge\kappa_w$, the family~\eqref{it:primitives-p5} of the system of~(ii) has the properties required in Proposition~\ref{prop:polymer-expansion-determinant} for the parameter set $(\alpha_0,\gamma^F,\kappa_w;C^F)$, and for $C_w=C^F$ the last inequality holds by the definition of $e^F_0$. Hence, for this system, with $C_d$ formed with $C_w=C^F$, clauses~\eqref{it:polymer-expansion-determinant-a}--\eqref{it:polymer-expansion-determinant-d} hold, clause~\eqref{it:polymer-expansion-determinant-e} holds under the three inequalities listed for it, and clauses~\eqref{it:polymer-expansion-determinant-f} and~\eqref{it:polymer-expansion-determinant-g} hold if moreover $l/L\in\mathrm{Cl}$ and $e_{k+1}=L^{\frac{1}{2}}\cdot e_k\le e^F_0$.
\end{enumerate}
All constants that occur in the theorem are independent of $k$, $l$, $\mathfrak{m}$, $e_k$ and $A$. The sets $\tilde{G}_k$ and $\tilde{R}_k$ are invariant under the addition of real constant connections (twists), and $\tilde{R}_k+\mathfrak{X}_k\subset\tilde{G}_k$ by~\eqref{eq:operators-fermion-G3}. Consequently all statements hold at $\mathscr{T}=\mathscr{A}+\mathscr{C}+\mathrm{cst}$, where $\mathscr{A}\in\tilde{R}_k$, $\mathscr{C}\in\mathfrak{X}_k$ and $\mathrm{cst}$ is a real constant connection, the twist being put into the first slot: the connection $\mathscr{A}+\mathrm{cst}\in\tilde{R}_k$ takes the place of $\mathscr{A}$.
\end{theorem}

\begin{proof}
The hypotheses of the results cited below are satisfied. Indeed, $L\ge5$ is odd, $\alpha_0\le\frac{1}{54}$ and $\mathfrak{m}\le1$. Condition~\eqref{eq:renormalisation-group-19} holds; since $C_V=12\mathbf{e}\cdot(2+b)\ge1$, and $a=C_z\cdot c(\gamma)\ge1$ by Proposition~\ref{prop:zero-field-levels} and Section~\ref{sec:norms-exponential-weights}, it gives $M\ge4L+1>4L$ and $M\ge\left(108\pi\cdot C_V\cdot a\right)^2\ge108\pi\cdot C_V\cdot a$. By the definition of $e^F_0$, conditions~\eqref{eq:renormalisation-group-20}--\eqref{eq:renormalisation-group-23} hold with $e^F_0$ in place of $e_0$.

(i) Theorem~\ref{thm:expansion-walks} gives Propositions~\ref{prop:decay-fine-propagator} and~\ref{prop:decay-resolvents}, for $k\ge1$, and the bounds of level $0$ in Proposition~\ref{prop:kernel-bounds-2}. For $k\ge1$, Proposition~\ref{prop:kernel-bounds-2} follows from Propositions~\ref{prop:operators-fermion-1} and~\ref{prop:kernels-unit-lattice}. Proposition~\ref{prop:bounds} follows from Proposition~\ref{prop:operators}; Proposition~\ref{prop:left-inverse}, with the constant $C_T$, from Proposition~\ref{prop:identities-bounds}; Proposition~\ref{prop:symmetries-fermion-operators} from Proposition~\ref{prop:symmetries-fermion-operators-2} together with Lemma~\ref{lem:symmetries-elementary-operators}(c), which gives the relations $Q_k(A)^T=Q_k^T(A)$ and $Q(A)^T=Q^T(A)$, the latter also at the level $0$; and Proposition~\ref{prop:coverings-tori} from Proposition~\ref{prop:operators-covering-torus}. The constants with which Theorem~\ref{thm:expansion-walks} and Propositions~\ref{prop:kernels-unit-lattice} and~\ref{prop:operators} give the bounds of Propositions~\ref{prop:decay-fine-propagator}--\ref{prop:bounds} do not exceed the maximum in the definition of $C^F$, and the factor $2$ in $C^F$ accounts for the sup norms formed with $|\cdot|_c$ (Section~\ref{sec:norms-pairings-kernels}).

(ii) For $k\ge1$ Theorem~\ref{thm:properties-weakened-families}, and for $k=0$ Theorem~\ref{thm:properties-level-zero}, give the properties~\ref{it:renormalisation-group-1}, \ref{it:renormalisation-group-2} and~\ref{it:renormalisation-group-5}--\ref{it:renormalisation-group-17} for the parameter set $(\alpha_0,\gamma,\kappa';2C_1)$ in the sense of Section~\ref{sec:properties-systems-primitives}, in which the bounds are required for both norms on $\mathbb{C}^4$; moreover $\gamma^F=\gamma$, and $2C_1\le C^F$ by the definition of $C^F$. The bounds of the properties~\ref{it:renormalisation-group-5}, \ref{it:renormalisation-group-6} and~\ref{it:renormalisation-group-7}, and the kernel bounds of property~\ref{it:renormalisation-group-17} for $\Gamma_0$ and $S_{0,y}$, remain valid when the constant in them is replaced by a larger number; the other properties do not contain the constant of the parameter set. At $k=0$ the bound of property~\ref{it:renormalisation-group-5} for $H_0(s,A)$, with the constant $4\exp(2\gamma)\cdot M^{27\alpha_0}$ that property~\ref{it:renormalisation-group-17} prescribes, is Theorem~\ref{thm:properties-level-zero}(b); it holds for both norms, since $H_0(s,A)$ is diagonal and its entries are numbers times the unit matrix (Definition~\ref{def:families-level-zero}). Hence the families have the properties of~(ii) for $(\alpha_0,\gamma^F,\kappa';C^F)$. As in~(i), the factor $2$ in the definition of $C^F$ accounts for the norm $|\cdot|_c$ (Section~\ref{sec:norms-pairings-kernels}).

(iii) In the statements and proofs of Section~\ref{sec:operators-fermion-field}, the ceiling $e_0$ enters in three ways only: through conditions~\eqref{eq:renormalisation-group-20}--\eqref{eq:renormalisation-group-24}; through the requirement $e_k\in(0,e_0]$, which is used in the form $e_k\le e_0$; and, in clause~\eqref{it:polymer-expansion-determinant-f} of Proposition~\ref{prop:polymer-expansion-determinant} and in the proofs of clauses~\eqref{it:polymer-expansion-determinant-f} and~\eqref{it:polymer-expansion-determinant-g}, through the inequality $e_{k+1}\le e_0$. Five statements from outside Section~\ref{sec:operators-fermion-field} enter the proof of Proposition~\ref{prop:polymer-expansion-determinant}: fact~\eqref{it:renormalisation-group-g0} of Section~\ref{sec:cubes-parameter-cubes}, on the sites, blocks and cubes of $T_k$, which involves only $L$, $M=L^m$ and $l$ and contains no charge; inequalities~\eqref{it:polymers-t1} and~\eqref{it:polymers-t8} of Proposition~\ref{prop:polymers}, which do not involve the charge; parts (a) and (c) of Lemma~\ref{lem:scaling-domains}, whose proofs use, apart from definitions, only fact~\eqref{it:renormalisation-group-g0}, the relations $l,l/L\in\mathrm{Cl}$ and $e_{k+1}=L^{\frac{1}{2}}\cdot e_k$, and the inequality $\varepsilon\le\frac{1}{4}$ (proof of clause~\eqref{it:polymer-expansion-determinant-f}); the statement of Section~\ref{sec:standing-conventions}, following convention~\eqref{eq:polymer-activities-6}, that the infimum in this convention does not change when only nets that are trees are admitted (Case~1 of the proof of clause~\eqref{it:polymer-expansion-determinant-e}), which concerns the torus $\mathbb{R}^3/l\mathbb{Z}^3$ and the closed $M$-cubes of $X$, involves only $l$ and $M$ and contains no charge; and the statement of Section~\ref{sec:spaces-activities} that the restriction to $\mathscr{C}=0$ of a family in $K^x_{k+1}$ lies in $K_{k+1}$ and that its norm does not exceed the norm of the family (end of the proofs of clauses~\eqref{it:polymer-expansion-determinant-f} and~\eqref{it:polymer-expansion-determinant-g}), whose proof uses only $0\in\mathfrak{X}_{k+1}$ and the definitions of the spaces. The number $e_0$ occurs in neither of the last two statements. Apart from these five statements and from classical theorems of analysis and topology, the proof takes from outside Section~\ref{sec:operators-fermion-field} only definitions; hypotheses of Proposition~\ref{prop:polymer-expansion-determinant}, namely conditions~\eqref{eq:renormalisation-group-20}, \eqref{eq:renormalisation-group-21}, \eqref{eq:renormalisation-group-23} and~\eqref{eq:renormalisation-group-24}, the properties~\ref{it:renormalisation-group-1}, \ref{it:renormalisation-group-2}, \ref{it:renormalisation-group-7}, \ref{it:renormalisation-group-12}, \ref{it:renormalisation-group-13}, \ref{it:renormalisation-group-14} and~\ref{it:renormalisation-group-17} of the family~\eqref{it:primitives-p5}, and the Borel measurability of this family in $y$ (Section~\ref{sec:primitive-operators}); and the numbers $\bar{\varepsilon}$, $c_X$, $\kappa_0$, $c_{\mathrm{rb}}$, $\kappa$ and $r$ of stages~\ref{it:order-choice-1}, \ref{it:order-choice-3} and~\ref{it:order-choice-8} of Definition~\ref{def:order-choice}, with the relations $\bar{\varepsilon}=\frac{1}{200}$, $\kappa_0=54\log72>0$, $\kappa\ge7\kappa_0+7$ and $r\ge2$ prescribed there. These stages precede stage~\ref{it:order-choice-9}, in which $e_0$ is chosen, so that their numbers do not depend on $e_0$; and in the definitions and hypotheses so taken, $e_0$ occurs only in conditions~\eqref{eq:renormalisation-group-20}, \eqref{eq:renormalisation-group-21}, \eqref{eq:renormalisation-group-23} and~\eqref{eq:renormalisation-group-24} and in the requirement $e_k\in(0,e_0]$ of the scale-free setting of Section~\ref{sec:properties-scheme}, in which the definitions of Sections~\ref{sec:fermion-integrals-blocking}, \ref{sec:primitive-operators} and~\ref{sec:properties-systems-primitives} are made. Hence Proposition~\ref{prop:polymer-expansion-determinant} and its proof remain valid when $e_0$ is replaced throughout by a number $t\in(0,1]$ such that conditions~\eqref{eq:renormalisation-group-20}, \eqref{eq:renormalisation-group-21} and~\eqref{eq:renormalisation-group-23} hold with $t$ in place of $e_0$, condition~\eqref{eq:renormalisation-group-24} in clause~\eqref{it:polymer-expansion-determinant-g} being taken with $t$ in place of $e_0$ as well. As shown at the beginning of this proof, $t=e^F_0$ is such a number, and the other hypotheses of Theorem~\ref{thm:expansion-walks} are satisfied. Moreover $\kappa\ge7\kappa_0+7$ (stage~\ref{it:order-choice-3}) and $L\ge5$ give
\[
L(\kappa-\kappa_0-1)-\kappa=(L-1)\cdot\kappa-L\cdot(\kappa_0+1)\ge7(L-1)\cdot(\kappa_0+1)-L\cdot(\kappa_0+1)=(6L-7)\cdot(\kappa_0+1)>0,
\]
so that the inequality $L(\kappa-\kappa_0-1)\ge\kappa$ required in clause~\eqref{it:polymer-expansion-determinant-g} holds. This proves the assertions of~(iii) on an arbitrary family~\eqref{it:primitives-p5}.

Now let $\kappa'\ge\kappa_w$. By~(ii), the family~\eqref{it:primitives-p5} of the system of~(ii) has the properties~\ref{it:renormalisation-group-1}, \ref{it:renormalisation-group-2}, \ref{it:renormalisation-group-7}, \ref{it:renormalisation-group-12}, \ref{it:renormalisation-group-13} and~\ref{it:renormalisation-group-14} for $(\alpha_0,\gamma^F,\kappa';C^F)$, at $k=0$ with property~\ref{it:renormalisation-group-7} in the form stated in property~\ref{it:renormalisation-group-17}. For the family~\eqref{it:primitives-p5}, the properties~\ref{it:renormalisation-group-1}, \ref{it:renormalisation-group-2}, \ref{it:renormalisation-group-12}, \ref{it:renormalisation-group-13} and~\ref{it:renormalisation-group-14} contain none of the numbers $\gamma$, $\kappa'$, $C$ of a parameter set. In property~\ref{it:renormalisation-group-7}, and in its form at $k=0$, one has $\exp\left(-\kappa'\cdot d_M(X)\right)\le\exp\left(-\kappa_w\cdot d_M(X)\right)$, because $\kappa'\ge\kappa_w$ and $d_M(X)\ge0$. So the family has these six properties for $(\alpha_0,\gamma^F,\kappa_w;C^F)$. For the system of (ii) and $C_w=C^F$, condition~\eqref{eq:renormalisation-group-24} holds with $e^F_0$ in place of $e_0$, by the definition of $e^F_0$. The assertions on an arbitrary family, applied with $(\gamma_w,C_w)=(\gamma^F,C^F)$, now give the assertions of~(iii) on the system of~(ii).

\emph{Real constants.} For $\tilde{G}_k$ the invariance under the addition of a real constant connection is part of~\eqref{eq:operators-fermion-G2}. Let $\mathscr{A}=A_0+A_1\in\tilde{R}_k$ be written as in Definition~\ref{def:domains}, let $\mathrm{cst}$ be a real constant connection, and let $\square$ be an $M$-cube. Let $A_0=A'+d\boldsymbol{\lambda}$ on $\square^{\dagger}$, where $\boldsymbol{\lambda}$ is a real function on the sites and $A'$ a real connection on the bonds of $\square^{\dagger}$ whose three norms on $\square^{\dagger}$ are below the thresholds of $R_k$. The cube $\square^{\dagger}$ has the side $(2L+1)\cdot M$ (Definition~\ref{def:enlarged-cubes}), and $(2L+1)\cdot M<13L\cdot M\le l$, because $l\in\mathrm{Cl}$ (Section~\ref{sec:lattice-geometry}). Hence $\square^{\dagger}$ lifts to a box of $\eta\mathbb{Z}^3$, and on this box $\mathrm{cst}=d\boldsymbol{\lambda}_1$ with $\boldsymbol{\lambda}_1(x):=\sum_{\mu}\mathrm{cst}_{\mu}x_{\mu}$. Adding $\mathrm{cst}=d\boldsymbol{\lambda}_1$ to $A_0=A'+d\boldsymbol{\lambda}$ yields $A_0+\mathrm{cst}=A'+d(\boldsymbol{\lambda}+\boldsymbol{\lambda}_1)$ on $\square^{\dagger}$; the real connection $A'$, and with it the bounds on its three norms, is unchanged. As $A_0+\mathrm{cst}$ is real and $\square$ is arbitrary, the connection $\mathscr{A}+\mathrm{cst}=(A_0+\mathrm{cst})+A_1$ lies in $\tilde{R}_k$; this is the assertion of Lemma~\ref{lem:invariances} on $\tilde{R}_k$ and real constants. Consequently $\mathscr{T}=(\mathscr{A}+\mathrm{cst})+\mathscr{C}\in\tilde{R}_k+\mathfrak{X}_k\subset\tilde{G}_k$ by~\eqref{eq:operators-fermion-G3}, and the statements whose domain is $\tilde{G}_k$, $\tilde{R}_k+\mathfrak{X}_k$ or $\tilde{R}_k\times\mathfrak{X}_k$ apply at $\mathscr{T}$, with $\mathscr{A}+\mathrm{cst}$ in place of $\mathscr{A}$.
\end{proof}

Theorem~\ref{thm:fermion-operators-weakened} follows from Theorem~\ref{thm:fermion-operators}.

\paragraph{Order of the constants}
The constants of Section~\ref{sec:operators-fermion-field} are defined in the following order; each item depends only on the quantities named in it and on the preceding items.
\begin{itemize}
\item $b_L=\frac{b}{L}$ and $b_{-}=b\left(1-\frac{1}{L}\right)$: functions of $(b,L)$.
\item $C_5$ and $\boldsymbol{\gamma}_5$ of Proposition~\ref{prop:zero-field-levels}, which are the constants of Theorem~\ref{thm:sup-holder-bounds-2}: functions of $(b,\alpha,L)$.
\item $c_{\mathrm{re}}$, $a_z$, $k_1$ and $C_{00}$ of Lemma~\ref{lem:zero-field-level}: functions of $(b,L)$.
\item $C_z$, $\gamma_z$, $\gamma=\gamma^F$, $a$ and $C_V$ of~\eqref{eq:operators-fermion-4}: functions of $(b,\alpha,L)$.
\item $\mathfrak{m}_0=1$, $\alpha^F=\frac{1}{54}$ and $L^F=5$: absolute numbers.
\item $C_{\mathrm{By}}$ and $C_{\mathrm{Gam}}$ of~\eqref{eq:operators-fermion-8}: functions of $(b,\alpha,L)$.
\item $C_T$ of Proposition~\ref{prop:identities-bounds}(d): a function of $(b,L)$.
\item $M^F$ of condition~\eqref{eq:renormalisation-group-19}: a function of $(b,\alpha,L,\kappa')$; it does not depend on $\alpha_0$.
\item $C_1$ of Theorem~\ref{thm:properties-weakened-families} and $C^F$ of Theorem~\ref{thm:fermion-operators}: functions of $(b,\alpha,L,M,\alpha_0,\kappa')$.
\item $C_d$ of~\eqref{eq:operators-fermion-10}: a function of $(b,c_X,C_w)$.
\item $e^F_0$, defined through conditions~\eqref{eq:renormalisation-group-20}--\eqref{eq:renormalisation-group-24}: a function of $(b,\alpha,L,M,\varepsilon,C_d,c_{\mathrm{rb}})$.
\end{itemize}
This is the order of stages~\ref{it:order-choice-4}, \ref{it:order-choice-5}, \ref{it:order-choice-6}, \ref{it:order-choice-7} and~\ref{it:order-choice-9} of Section~\ref{sec:order-choice}. Accordingly, $\gamma^F$, $\alpha^F$ and $\mathfrak{m}_0$ precede $M$, and none of $C^F$, $\varepsilon$, $e$ enters a condition imposed on $M$.

\section{The renormalisation group step without large fields, with an external connection}\label{sec:renormalisation-group-step}

This section contains the construction of the step map of Definition~\ref{def:renormalisation-step} (Section~\ref{sec:definition-step-map}), the proof of Theorem~\ref{thm:step-small-fields}, and the proofs of Lemmas~\ref{lem:inclusions-bounds}--\ref{lem:periodic-extension} and~\ref{lem:phase-factors} (Section~\ref{sec:elementary-properties-field}). The statements of Section~\ref{sec:renormalisation-group-scheme} on which these proofs rest are recalled, some of them in abridged form, in Section~\ref{sec:recalled-auxiliary-results}; the notation is that of Section~\ref{sec:renormalisation-group-scheme} and of Section~\ref{sec:notation}, and a property of a system of primitives that is cited by its number is the item with that number in the list of Section~\ref{sec:properties-systems-primitives}.

\subsection{Setting, notation and recalled results}\label{sec:setting-notation-recalled}

\subsubsection{Hypotheses and conditions}\label{sec:hypotheses-conditions}

Throughout Section~\ref{sec:renormalisation-group-step} the standing hypothesis~\ref{def:hypothesis-one-step} of Section~\ref{sec:one-step-operators} is in force. Written out, it consists of the following requirements.

The constants of stages~\ref{it:order-choice-1}--\ref{it:order-choice-9} are chosen as in Section~\ref{sec:order-choice}, subject to the conditions~\eqref{eq:renormalisation-group-A0}--\eqref{eq:renormalisation-group-A10}, \eqref{eq:renormalisation-group-19}, \eqref{eq:renormalisation-group-20}--\eqref{eq:renormalisation-group-24} and~\eqref{eq:renormalisation-group-25} of Section~\ref{sec:conditions}. The data are an integer $k \ge 0$, a number $e_k > 0$, a side $l$ and a mass $\mathfrak{m}$, subject to the following requirements. The number $e_k$ satisfies
\[
e_{k+1} := L^{\frac{1}{2}} \cdot e_k \le e_0;
\]
here $e_0$ is the number fixed in stage~\ref{it:order-choice-9}, an upper bound for the charges. In the case $k = 0$ the letter $e_k$ stands for the running charge $q$ of level $0$ (Definition~\ref{def:notation}(e)), never for this bound. With $e_k$ and with the exponents $r$, $p_0$, $p$, $\varepsilon$ of stage~\ref{it:order-choice-8} one forms, as in Section~\ref{sec:running-quantities-tori}, the numbers
\begin{align*}
& v_k := -\log e_k, \qquad r_k := \min\left\{ L^i : i \ge 0 \text{ an integer},\ L^i \ge v_k^r \right\}, \qquad p_k := v_k^p,\\
& p_{0,k} := v_k^{p_0}, \qquad w^{\mathrm{m}}_k := e_k^{\frac{3}{4} - 8\varepsilon}, \qquad w^{\mathrm{E}}_k := e_k^{\frac{1}{4} - 7\varepsilon}.
\end{align*}
The side satisfies $l \in \mathrm{Cl}$, $l/L \in \mathrm{Cl}$ and $l \ge M \cdot r_k$, where the class $\mathrm{Cl}$ of Section~\ref{sec:lattice-geometry} consists of the numbers $l$ for which $l/(L \cdot M)$ is an odd integer $\ge 13$. The mass satisfies $0 < \mathfrak{m} \le \mathfrak{m}_0/L$, with $\mathfrak{m}_0$ as in Definition~\ref{def:mass-bound-stop}. Put $\eta_k := L^{-k}$. The lattices $\mathbb{L}_k = \eta_k\mathbb{Z}^3$, $T_k = \mathbb{L}_k/l\mathbb{Z}^3$ and $T^0_k = \mathbb{Z}^3/l\mathbb{Z}^3$ of Section~\ref{sec:lattice-geometry}, the operators of Sections~\ref{sec:averaging-fluctuations-gauge} and~\ref{sec:exact-fermion-operators}, and the domains of Section~\ref{sec:domains-bond-functions} are formed with $(k, l, \mathfrak{m}, e_k)$. The objects of level $k+1$ are formed with the triple $\left(l/L,\, L \cdot \mathfrak{m},\, e_{k+1}\right)$ of side, mass and charge. The unit lattice $T^0_k$ has $n_k = l^3$ points, and the unit lattice of level $k+1$, whose side is $l/L$, has $n_{k+1} = (l/L)^3$ points. On $T_k$ a system of primitives of level $k$ in the sense of Section~\ref{sec:primitive-operators} is given, and it has the properties~\ref{it:renormalisation-group-1}--\ref{it:renormalisation-group-17} for the parameters $\mathrm{par} = (\alpha_0, \gamma_w, \kappa_w; C_w)$. No further relation among $k$, $l$, $\mathfrak{m}$ and $e_k$ is assumed.

The exponential function is written $\exp$ (Definition~\ref{def:notation}(a)).

The inequalities \eqref{eq:renormalisation-group-30}--\eqref{eq:renormalisation-group-38} below are the conditions~\eqref{eq:renormalisation-group-A0}--\eqref{eq:renormalisation-group-A8} of Section~\ref{sec:conditions}, or parts of them, written for the number $e_k$. Each of them is required in Section~\ref{sec:conditions} for every value of $e_k$ in $(0, e_0]$, or is required there at $e_0$ and remains true when $e_0$ is replaced by a smaller positive number. Since $e_k \le e_{k+1} \le e_0$, they hold as stated, and they remain valid when $e_k$ is replaced by $e_{k+1}$. The constants in them are those of Section~\ref{sec:order-choice}; the numbers $c_{\gamma}$, $C_w$, $c_C$ and $C_T$ are described in Section~\ref{sec:notation}.
\begin{align}
& L \cdot \left(\kappa - 6\kappa_0 - 4\right) - 1 \ge \kappa, \qquad \kappa_w = \kappa + \kappa_0 + 3, \qquad M > 4L, \notag\\
& M^{\alpha_0} \ge 1 + \exp(\hat{\kappa}), \qquad \text{where } \hat{\kappa} = N_c \cdot \kappa + \log\left(6c_X\right) \text{ and } N_c = (2L+1)^3; \label{eq:renormalisation-group-30}\\
& p_{0,k} \le e_k^{-\frac{\varepsilon}{4}}, \qquad 12 \exp\left(-\frac{p_{0,k}^2}{2}\right) \le e_k^7, \qquad e_k^{\frac{5\varepsilon}{2}} \le \frac{1}{2}, \notag\\
& c_{\mathrm{fl}} \cdot e_k^{\varepsilon} \le 1, \qquad 3M \cdot e_k^{\varepsilon} \le 1; \label{eq:renormalisation-group-31}\\
& 2\, C_w^2 \cdot c_{\gamma}^2 \cdot c_C \cdot e_k^{\frac{\varepsilon}{4}} \le 1; \label{eq:renormalisation-group-32}\\
& 16\, c_X \cdot c_{\gamma} \cdot C_w \cdot C_T^2 \cdot M^3 \cdot e_k^{\frac{\varepsilon}{2}} \le 1, \qquad 96\, M^3 \cdot e_k^{\frac{1}{4}} \le 1; \label{eq:renormalisation-group-33}\\
& 32\, C_w \cdot C_T \cdot e_k^{\frac{3\varepsilon}{2}} \le 1; \label{eq:renormalisation-group-34}\\
& 88\, M^3 \cdot e_k^{\frac{1}{4} - 6\varepsilon} \le 1; \label{eq:renormalisation-group-35}\\
& 16\, C_w^2 \cdot c_{\gamma} \cdot e_k^{2\varepsilon} \le 3; \label{eq:renormalisation-group-36}\\
& c_S \cdot e_k^{\frac{1}{4} - 5\varepsilon} \le 1, \qquad c_S \cdot L^3 \cdot e_k^{\varepsilon} \le 1, \qquad \mathbf{H}_0 \le 1, \qquad \text{where} \notag\\
& c_S := 10^5\, c_X^5 \cdot K_0^5 \cdot \exp(6\kappa) \quad \text{and} \quad \mathbf{H}_0 := 36\, c_X^2 \cdot K_0^2 \cdot e_k^{\frac{1}{4} - 5\varepsilon}; \label{eq:renormalisation-group-37}\\
& e_k^{\varepsilon} \le \frac{1}{4}, \qquad e_k^{-\frac{3}{4} + \varepsilon} + e_k^{-1 + 2\varepsilon} \le e_k^{-1 + \varepsilon}. \label{eq:renormalisation-group-38}
\end{align}
The requirements in~\eqref{eq:renormalisation-group-30} are contained in stages~\ref{it:order-choice-4}--\ref{it:order-choice-6}. Multiplication by $e_k^{1 - \varepsilon} > 0$ shows that the second inequality in~\eqref{eq:renormalisation-group-38} is equivalent to
\[
e_k^{\frac{1}{4}} + e_k^{\varepsilon} \le 1.
\]

Besides the third of its conditions, which is the first inequality in~\eqref{eq:renormalisation-group-30}, stage~\ref{it:order-choice-4} yields the following facts: $L$ is odd; $L \ge 4^8$, so that in particular $L \ge 5$; and $L^{-\frac{3}{4}} \le \frac{1}{2}$.

The conditions~\eqref{eq:renormalisation-group-A9} and~\eqref{eq:renormalisation-group-A10} of Section~\ref{sec:conditions} are the conditions on $e_0$ which Propositions~\ref{prop:adjustment}, \ref{prop:contraction}, \ref{prop:fields-domains}, \ref{prop:left-inverse}, \ref{prop:local-expansion-determinant} and~\ref{prop:comparison-fermion-operators}, in the form recalled in Section~\ref{sec:recalled-auxiliary-results}, and Proposition~\ref{prop:step-external-connection}, in the form recalled at the beginning of Section~\ref{sec:proof-clause-theorem}, require in addition to~\eqref{eq:renormalisation-group-31}--\eqref{eq:renormalisation-group-38}; in the arguments of Sections~\ref{sec:elementary-properties-field}--\ref{sec:data-level-locality} these two conditions enter only through the seven propositions. For the other constants, the hypotheses of these propositions are described in the first paragraph of Section~\ref{sec:recalled-auxiliary-results}, and the conditions on which their statements and proofs rest, among them conditions on $\alpha$, $\varepsilon$, $L$, $M$ and $r$ and relations satisfied by $C_w$ and $\gamma_w$, are listed in Section~\ref{sec:uniformity}.

The number $e_{k+1}$ lies in $(0, e_0]$, and the side $l/L$ of level $k+1$ belongs to $\mathrm{Cl}$. Hence~\eqref{eq:renormalisation-group-38} and the inequality $3M \cdot e_k^{\varepsilon} \le 1$ of~\eqref{eq:renormalisation-group-31} remain true when $e_k$ is replaced by $e_{k+1}$, and Lemmas~\ref{lem:openness}, \ref{lem:inclusions-domains} and~\ref{lem:phases-differences-along} are valid at the level $k+1$. Lemma~\ref{lem:symmetries-domains} is valid at the level $k+1$ for each of its domains other than $\mathrm{Dom}_{k+1}$; its assertion on $\mathrm{Dom}_k$ is derived from the family~\eqref{it:primitives-p1} of the given system of primitives, which is a system of level $k$.

\subsubsection{Notation}\label{sec:notation}

\paragraph{Constants.}
Indices, their supports, and the distance $d(x,x')$ of two indices, namely the sup-distance of $\operatorname{supp} x$ and $\operatorname{supp} x'$ on the torus, are those of Section~\ref{sec:lattice-geometry}. On $\mathbb{R}^3$ put $d(\tilde{y},\tilde{y}') := |\tilde{y} - \tilde{y}'|$, the sup-distance of Definition~\ref{def:sup-norm-distance}; with it the distance $d(\tilde{x},\tilde{x}')$ of two indices of $\mathbb{L}_k$ or of $\mathbb{Z}^3$ is formed from their supports as on the torus, sup-diameters are taken with respect to $d$, and for a point $z$ and a finite set $B$ of points, both on the torus or both in $\mathbb{R}^3$, $\operatorname{dist}(z,B) := \min_{y \in B} d(z,y)$. For a level $k \ge 0$ and a side $l \in \mathrm{Cl}$ consider the following three sets of indices: the pairs formed by a unit site and a spinor index, the support of such a pair being the site; the unit bonds; the elements of the basis $\{\Upsilon\}$ of $V_k$ (Section~\ref{sec:averaging-fluctuations-gauge}). Let $c_{\gamma}$ be the supremum of the sums
\[
\sum_{x'} \exp\left(-\gamma_w \cdot d(x,x')\right),
\]
taken over all levels $k \ge 0$, all sides $l \in \mathrm{Cl}$, all indices $x$ of level $k$ on the torus of side $l$, and the three sets, the index $x'$ ranging over the set in question.

The number $c_{\gamma}$ is finite, satisfies $c_{\gamma} \ge 1$, and depends on $(\gamma_w, L)$ only. The inequality $c_{\gamma} \ge 1$ holds because, for an index $x$ of the set over which $x'$ ranges, the term with $x' = x$ equals $1$. Finiteness rests on three facts, valid on the torus and on $\mathbb{Z}^3$.
\begin{enumerate}
\item[(i)] The support of every index has sup-diameter at most $L$; for the elements of $\{\Upsilon\}$ this is contained in Proposition~\ref{prop:basis-map-box}.
\item[(ii)] Let $B$ be an $L$-block. For each of the three sets,
\[
\#\left|\left\{ x' : \operatorname{supp} x' \cap B \ne \emptyset \right\}\right| \le 6L^3,
\]
the index $x'$ ranging over the set in question.
\item[(iii)] Let $z$ be a point of the torus and $\tilde{z} \in \mathbb{R}^3$ a lift of $z$. The sum of $\exp\left(-\gamma_w \cdot \operatorname{dist}(z, B)\right)$ over the $L$-blocks $B$ of the torus is at most the sum of $\exp\left(-\gamma_w \cdot \operatorname{dist}(\tilde{z}, \tilde{B})\right)$ over the $L$-blocks $\tilde{B}$ of $\mathbb{Z}^3$, and the latter sum is bounded, uniformly in $\tilde{z}$, by a number that depends on $\gamma_w$ and $L$ only.
\end{enumerate}
Indeed, let $x$ be an index and $z$ a point of $\operatorname{supp} x$. If $\operatorname{supp} x'$ meets an $L$-block $B$, then $\operatorname{dist}(z, B) \le d(x,x') + 2L$, by (i). Assign to each $x'$ one $L$-block that $\operatorname{supp} x'$ meets. Then (ii) and (iii), with a lift $\tilde{z}$ of $z$, give
\[
\sum_{x'} \exp\left(-\gamma_w \cdot d(x,x')\right) \le 6L^3 \cdot \exp\left(2\gamma_w \cdot L\right) \cdot \sum_{\tilde{B}} \exp\left(-\gamma_w \cdot \operatorname{dist}(\tilde{z}, \tilde{B})\right),
\]
and the right side is bounded in terms of $\gamma_w$ and $L$ alone.

Although $M$ occurs in the definition of $c_{\gamma}$, namely through the class $\mathrm{Cl}$ and in no other way, the value of $c_{\gamma}$ is independent of $M$. To begin with, a single family of bond functions on $\mathbb{Z}^3$, which is mapped onto itself by every translation by a vector of $L\mathbb{Z}^3$ and is determined by $L$ alone, is carried onto the basis $\{\Upsilon\}$ of the torus of side $l$ by the covering map $\mathbb{Z}^3 \to \mathbb{Z}^3/l\mathbb{Z}^3$, for every $l \in \mathrm{Cl}$ (Proposition~\ref{prop:basis-map-box}; the vectors $e^0_i$ of Proposition~\ref{prop:fluctuation-space}(i) depend on $L$ only, and not on $k$ or $l$). Now fix a level $k \ge 0$ and one of the three sets, let $\tilde{x}$ be an index of $\mathbb{L}_k$ or of $\mathbb{Z}^3$, and, for $l \in \mathrm{Cl}$, let $x$ denote the image of $\tilde{x}$ on the torus of side $l$. For each index $x'$ of the torus that belongs to the fixed set, select a lift $\hat{x}'$ of $x'$ for which $d(\tilde{x},\hat{x}')$ is smallest among all lifts of $x'$. By the definition of the distance on the torus (Definition~\ref{def:sup-norm-distance}), $d(x,x') = d(\tilde{x},\hat{x}')$; moreover, different indices $x'$ have different lifts $\hat{x}'$. Hence
\[
\sum_{x'} \exp\left(-\gamma_w \cdot d(x,x')\right) = \sum_{x'} \exp\left(-\gamma_w \cdot d(\tilde{x},\hat{x}')\right) \le S(\tilde{x}) := \sum_{\tilde{x}'} \exp\left(-\gamma_w \cdot d(\tilde{x},\tilde{x}')\right),
\]
where $\tilde{x}'$ ranges over all indices of $\mathbb{Z}^3$ of the fixed kind: the pairs of a site of $\mathbb{Z}^3$ and a spinor index, the bonds of $\mathbb{Z}^3$, or the members of the family just described. Every lift $\tilde{x}''$ of $x'$ other than $\hat{x}'$ satisfies $d(\tilde{x},\tilde{x}'') \ge l/2 - L$: it differs from $\hat{x}'$, which is at least as near to $\tilde{x}$, by a non-zero vector of $l\mathbb{Z}^3$, whose sup-norm is at least $l$, and all supports have sup-diameter at most $L$. Therefore $S(\tilde{x})$ exceeds the left side at most by the sum of those terms of $S(\tilde{x})$ for which $d(\tilde{x},\tilde{x}') \ge l/2 - L$. The series $S(\tilde{x})$ converges by (i)--(iii), and the class $\mathrm{Cl}$ is unbounded; hence the supremum of the left side over $l \in \mathrm{Cl}$ equals $S(\tilde{x})$. Thus $c_{\gamma}$ is the supremum of the numbers $S(\tilde{x})$ over all levels $k \ge 0$, all indices $\tilde{x}$ of $\mathbb{L}_k$ and of $\mathbb{Z}^3$, and the three sets, and $M$ does not enter.

The constant $C_w$ of $\mathrm{par}$ satisfies $C_w \ge 1$. Indeed, by Sections~\ref{sec:primitive-operators} and~\ref{sec:photon-operators},
\[
C_w = 64\, K_0 \cdot C_w^1, \qquad C_w^1 \ge 4 \exp\left(2\gamma_w\right) \cdot M^{27\alpha_0} \ge 4,
\]
and $K_0 \ge 1$ (Section~\ref{sec:absolute-constants}); in particular $C_w \ge C_w^1 \ge 4 \exp\left(2\gamma_w\right) \cdot M^{27\alpha_0}$. The number $c_C$ is the constant of Proposition~\ref{prop:basis-map-box}, and $C_T$ is the constant of Proposition~\ref{prop:left-inverse}. By Theorem~\ref{thm:fermion-operators}(i), Proposition~\ref{prop:left-inverse} holds with the constant of Proposition~\ref{prop:identities-bounds}(d),
\[
C_T = 2 \exp\left(\frac{3}{2}\right) \cdot \left[ b_{-}^{-1} \cdot \left(54L \cdot \exp(6) + 12 \exp(1) + 1\right) + 1 \right],
\]
where $b_{-} = b \cdot \left(1 - L^{-1}\right) > 0$ is the number of Section~\ref{sec:running-quantities-tori}; hence $C_T \ge 2 \exp\left(\frac{3}{2}\right) \ge 1$. Put
\[
C_W := 4\, C_w^2 \cdot c_{\gamma}.
\]
The rates $\kappa_1, \dots, \kappa_4$ are the numbers
\[
\kappa_1 := \kappa - \kappa_0 - 1, \qquad \kappa_2 := \kappa - 2\kappa_0 - 2, \qquad \kappa_4 := \kappa - 3\kappa_0 - 3, \qquad \kappa_3 := \kappa - 5\kappa_0 - 4,
\]
and the weights $h_k$ of Section~\ref{sec:grassmann-algebras-norms} are written with their two components as
\[
h_k = \left(h_{k,1}, h_{k,2}\right) = \left(e_k^{-\frac{1}{4}}, e_k^{-\frac{1}{4} + \varepsilon}\right).
\]

\paragraph{Cubes and polymers.}
The following notions are those of Section~\ref{sec:lattice-geometry}: the $M$-cubes; the enlargement $X^{\sim i}$ of a union $X$ of $M$-cubes by $i$ layers of $M$-cubes; the polymers, that is, the non-empty unions of $M$-cubes that are connected through faces, and the set $D_k$ of the polymers; the numbers $d_M(X)$ and $|X|_M$, the latter being the number of the $M$-cubes of $X$; the small polymers; the union $\bar{X}^L$ of the $LM$-cubes that contain an $M$-cube of $X$; the parameter cubes of the three kinds $\pi_k$, $\pi'_k$, $\pi''_k$, the relation of being linked, and the link-components; the polydisk $\mathrm{Pd}_k$; the number $N_c=(2L+1)^3$. A lattice site is a site of $T_k$. The cubes of $\pi'_k$ and of $\pi''_k$ are closed: they are the cubes $\{x:|x-c|\le M\}$ with $c\in M\mathbb{Z}^3$, and the cubes $\{x:|x-c|\le LM\}$ with $c\in LM\mathbb{Z}^3$, respectively. Let $Y$ be a union of $M$-cubes. A \emph{component} of $Y$ is a union of $M$-cubes of $Y$ that is connected through faces and is maximal with this property.

\paragraph{Indices.}
Each of the two arguments $x$, $x'$ of the kernel $\mathrm{Op}(x,x')$ of an operator is a site, a bond or an element $\Upsilon$ of the basis of Section~\ref{sec:averaging-fluctuations-gauge}; objects of these three kinds are called \emph{indices}. The support $\operatorname{supp}x$ of an index $x$ is the set of its end points: $\operatorname{supp}x=\{x\}$ for a site $x$, the set of the two end points for a bond, and, for $x=\Upsilon$, the set of the end points of the bonds in the support of $\Upsilon$. Let $W$ be a set of sites. An index $x$ \emph{meets} $W$ if $\operatorname{supp}x\cap W\ne\emptyset$; it \emph{is contained in} $W$, or \emph{lies in} $W$, if $\operatorname{supp}x\subset W$. A parameter cube is contained in $W$ if each of its lattice sites lies in $W$, and it meets an $M$-cube if the two have a lattice site in common. A bond is \emph{at distance at most $1$ from} $W$ if each of its two end points has sup-distance at most $1$ from $W$. For indices $x$, $x'$ let $d(x,x')$ be the sup-distance on the torus between $\operatorname{supp}x$ and $\operatorname{supp}x'$. For a site $z$ of $T_k$ put
\[
K_z:=\left\{x\in T_k:|x-z|\le 2L\right\}.
\]

\paragraph{Antiperiodicity.}
Directions are denoted by $\nu\in\{0,1,2\}$, and $e_{\nu}$ is the unit vector in the direction $\nu$. Spinor functions and the generators $\psi$, $\chi$, $\Psi$, $W$ are indexed by the points of $\mathbb{L}_k$ (spinor functions on $T_k$, and $\psi$, $\chi$) or of $\mathbb{Z}^3$ (spinor functions on $T^0_k$, and $\Psi$, $W$), and they are antiperiodic:
\[
f(x+l\cdot e_{\nu})=-f(x),\qquad \chi(x+l\cdot e_{\nu},x'+l\cdot e_{\nu})=-\chi(x,x')\qquad(\nu=0,1,2),
\]
where $f$ is a spinor function or one of the generators $\psi$, $\Psi$, $W$, and the remaining indices of a generator are left unchanged. As in Section~\ref{sec:grassmann-algebras-norms}, a pair of $\mathrm{Pz}_k$ enters through its shortest lift, and a fundamental domain of $l\mathbb{Z}^3$ is fixed; the generators whose point $x$ in the display lies in this domain are independent. For $\chi$ this point is the first point of the pair, and the second point is then determined by the shortest lift. Another choice of the fundamental domain changes the signs of some generators. The kernels of an element of one of the Grassmann algebras of Section~\ref{sec:grassmann-algebras-norms} are the coefficient functions of its expansion in products of the generators $\psi(\xi)$, $\chi(\boldsymbol{\zeta})$, $\Psi(u)$, $W(u)$; their arguments are the arguments $\xi$, $\boldsymbol{\zeta}$, $u$ of these generators, with points in $\mathbb{L}_k$ for $\xi$ and $\boldsymbol{\zeta}$, and in $\mathbb{Z}^3$ for $u$. Each such kernel is antiperiodic in each of its arguments: it changes its sign when the point, or the pair of points, of one argument is translated by one of the vectors $l\cdot e_{\nu}$. Consequently the product of such a kernel with the generators of its arguments is a function on the torus, and neither the moduli of these kernels nor any of the norms depends on the choice of the fundamental domain. An even element is left fixed by the simultaneous change of sign of all the generators $\psi$, $\chi$, $\Psi$, $W$, barred and unbarred (Section~\ref{sec:antiperiodicity}).

\paragraph{Covariant differences and block averages.}
As in Section~\ref{sec:hypotheses-conditions}, $\mathfrak{m}$ denotes the mass. For a site $x$ of $T_k$ and $\boldsymbol{\sigma}\in\{1,-1\}$, the direction $\nu$ being understood, put
\[
x^{\boldsymbol{\sigma}}:=x+\boldsymbol{\sigma}\cdot\eta_k\cdot e_{\nu},\qquad P^{\boldsymbol{\sigma}}_{\nu}:=\frac{1}{2}\left(1-\boldsymbol{\sigma}\cdot\gamma_{\nu}\right),
\]
where $\gamma_{\nu}$ are the matrices of Section~\ref{sec:dirac-matrices-lattice}. Let $B$ be a complex bond function on $T_k$. For a bond $(x,x')$ of $T_k$ put $u_B(x,x'):=\exp\left(i\,e_k\cdot\eta_k\cdot B(x,x')\right)$. For a unit site $y$ and a site $x\in B_k(y)$ put
\[
\phi_{B,y}(x):=\exp\left(i\,e_k\cdot\eta_k\cdot(\tau_kB)(y,x)\right),\qquad (\tau_kB)(y,x):=\sum_{0\le i<k}(\tau B)(x_{i+1},x_i),
\]
where $x_i$ is the centre of the block of side $L^i\eta_k$ that contains $x$, so that $x_0=x$ and $x_k=y$, and where $\tau$ is the mean over the six rectilinear paths (Definition~\ref{def:pairs-sites-holder}). One has $\sum_{\boldsymbol{\sigma}\in\{1,-1\}}P^{\boldsymbol{\sigma}}_{\nu}=1$, and $Q_k(B)\cdot Q_k^T(-B)=1$: the phases of $Q_k(B)$ and of $Q_k^T(-B)$ cancel, and $\#|B_k(y)|=\eta_k^{-3}$. By these two relations the definitions of Section~\ref{sec:exact-fermion-operators} take the following form, for spinor functions $f$, sites $x$ of $T_k$ and unit sites $y$:
\begin{align}
\left(D_Bf\right)(x)&=-\eta_k^{-1}\sum_{\nu,\boldsymbol{\sigma}}P^{\boldsymbol{\sigma}}_{\nu}\cdot\left[u_B(x,x^{\boldsymbol{\sigma}})\cdot f(x^{\boldsymbol{\sigma}})-f(x)\right],\notag\\
\left(Q_k(B)f\right)(y)&=\eta_k^3\sum_{x\in B_k(y)}\phi_{B,y}(x)\cdot f(x),\label{eq:renormalisation-group-39}\\
T_k(B)&=b_k^{-1}\cdot Q_k(B)\cdot\left(D_B+\mathfrak{m}\right)+Q_k(B)\quad(k\ge1),\qquad T_0(B)=1.\notag
\end{align}

\paragraph{A contour functional.}
Let
\[
\rho_1:=\frac{1}{2}\,e_k^{-\frac{3}{4}+\frac{9\varepsilon}{2}},\qquad \rho_2:=e_k^{-\frac{3}{4}+7\varepsilon},\qquad \rho_4:=e_k^{-\frac{1}{4}+3\varepsilon}.
\]
Let $V$ be a finite-dimensional complex vector space with a norm $|\cdot|$; the case $V=\mathbb{C}$, with the modulus as norm, is included. For $\rho>0$ with $\rho\ne1$ and a continuous map $f$ from the circle $\{t\in\mathbb{C}:|t|=\rho\}$ to $V$ define the element
\[
\oint_{\rho}f:=\frac{1}{2\pi i}\oint_{|t|=\rho}\mathrm{d}t\,\frac{f(t)}{t(t-1)}
\]
of $V$, the circle being oriented counterclockwise. Let $\rho>1$, and let $f$ be a continuous map from the closed disc $\{t\in\mathbb{C}:|t|\le\rho\}$ to $V$ that is analytic in the open disc $\{t\in\mathbb{C}:|t|<\rho\}$. Then
\[
\oint_{\rho}f=f(1)-f(0),\qquad \left|\oint_{\rho}f\right|\le\frac{1}{\rho-1}\cdot\sup_{|t|=\rho}|f(t)|.
\]
Indeed, with respect to a basis of $V$ each coordinate of $f$ is a complex function that is continuous on the closed disc and analytic in the open disc, and the residues of the corresponding coordinate of $t\mapsto f(t)/(t(t-1))$ at $0$ and at $1$ are the coordinates of $-f(0)$ and of $f(1)$; the residue theorem is applied to each coordinate on the circles of radius $r\in(1,\rho)$, and $r\to\rho$. For the bound: the norm of the integral of a continuous map from the circle to $V$ is at most the integral of its norm with respect to arc length, the circle has length $2\pi\rho$, and $|t(t-1)|\ge\rho(\rho-1)$ for $|t|=\rho$. By \eqref{eq:renormalisation-group-31} and \eqref{eq:renormalisation-group-35} one has $\rho_2\le\rho_1$ and $\rho_1,\rho_2,\rho_4\ge2$, so that
\[
\frac{1}{\rho_i-1}\le\frac{2}{\rho_i}\qquad(i=1,2,4).
\]
In detail: $\rho_2/\rho_1=2e_k^{\frac{5\varepsilon}{2}}\le1$ by the condition $e_k^{\frac{5\varepsilon}{2}}\le\frac{1}{2}$ of \eqref{eq:renormalisation-group-31}, which also gives $e_k<1$. The condition \eqref{eq:renormalisation-group-35} gives $e_k^{\frac{1}{4}-6\varepsilon}\le\frac{1}{88}$; hence $\rho_4^{-1}=e_k^{\frac{1}{4}-3\varepsilon}\le e_k^{\frac{1}{4}-6\varepsilon}\le\frac{1}{88}$, and $\rho_2^{-1}=e_k^{\frac{3}{4}-7\varepsilon}\le\rho_4^{-1}$, because $\frac{3}{4}-7\varepsilon\ge\frac{1}{4}-3\varepsilon$ for $\varepsilon\le\frac{1}{200}$ (stages~\ref{it:order-choice-1} and~\ref{it:order-choice-8}). Thus $\rho_1\ge\rho_2\ge\rho_4\ge88$.

\paragraph{Pairings.}
As in Section~\ref{sec:fermion-integrals-blocking},
\[
\langle\bar{\psi},\psi\rangle:=\sum_{x\in T_k}\eta_k^3\sum_{\boldsymbol{\beta}}\bar{\psi}_{\boldsymbol{\beta}}(x)\cdot\psi_{\boldsymbol{\beta}}(x),\qquad \langle\bar{\Psi},O\Psi\rangle:=\sum_{y,y'\in T^0_k}\bar{\Psi}(y)\cdot O(y,y')\cdot\Psi(y'),
\]
where the inner sum in the first expression is taken over the spinor index $\boldsymbol{\beta}$, and $O$ is an operator with the kernel $O(y,y')$, $y,y'\in T^0_k$. For an $M$-cube $\square$, $\langle\bar{\psi},\psi\rangle_{\square}$ is the first sum with $x$ restricted to the sites of $\square$.

\subsubsection{Activities with and without slot}\label{sec:activities-slot}

For $j=k,k+1$ the spaces $K^{x,C}_j$, $K^x_j$ (with slot) and $K^C_j$, $K_j$ (without slot) are those of Definition~\ref{def:renormalisation-group-1} in Section~\ref{sec:spaces-activities}, defined by the clauses~\eqref{it:renormalisation-group-1-k1}--\eqref{it:renormalisation-group-1-k7}. The definitions and the arguments of Sections~\ref{sec:definition-step-map}--\ref{sec:data-level-locality} are given for activities with slot.

The \emph{case without slot} is the specialisation in which, for $j=k,k+1$, the set $\mathfrak{X}_j$ is replaced by $\{0\}$. In this case, throughout, $\mathscr{C}=0$, the total connection is $\mathscr{T}=A+\mathscr{C}=A$, and the gauge function $\lambda'$ of the second slot in clause~\eqref{it:renormalisation-group-1-k4} is $0$; the spaces $K^{x,C}_j$, $K^x_j$, the norms $\|\cdot\|^x_j$ and the ball $\mathrm{Ball}^x_k$ are replaced by $K^C_j$, $K_j$, $\|\cdot\|_j$ and $\mathrm{Ball}_k$, and $\delta E^{\mathrm{det}}_k$ is replaced by $E^{\mathrm{det}}_k$.

Propositions~\ref{prop:symmetries-extraction}, \ref{prop:adjustment}, \ref{prop:contraction} and~\ref{prop:local-expansion-determinant}, in the form stated in Section~\ref{sec:recalled-auxiliary-results}, include the case without slot. For Propositions~\ref{prop:symmetries-extraction}, \ref{prop:contraction} and~\ref{prop:local-expansion-determinant} this case consists of the parts without slot of Proposition~\ref{prop:symmetries-extraction}, of Proposition~\ref{prop:contraction} and of Proposition~\ref{prop:local-expansion-determinant}(d6); these parts are contained in Theorem~\ref{thm:extraction}, whose assertions hold also without the second slot, in Theorem~\ref{thm:contraction}(c), and in Proposition~\ref{prop:polymer-expansion-determinant}(f), (g). For Proposition~\ref{prop:adjustment} the case without slot is its last assertion. This assertion is contained in the last assertion of Theorem~\ref{thm:properties-adjustment}, which is proved for all the parts (a)--(e) of that theorem at the end of its proof; the adjustment is formed pointwise in the connections, and $\mathscr{T}=\mathscr{A}$ if $\mathscr{C}=0$. The part without slot of Theorem~\ref{thm:contraction} rests on the forms without the second slot of Theorems~\ref{thm:extraction}, \ref{thm:properties-adjustment}, \ref{thm:polymer-activities-3} and~\ref{thm:closure-properties-local}(c)--(e), each of which is stated in the theorem concerned. For Theorems~\ref{thm:extraction} and~\ref{thm:properties-adjustment}, and for parts~(c)--(e) of Theorem~\ref{thm:closure-properties-local}, this form is the last assertion of the theorem. For Theorem~\ref{thm:polymer-activities-3} it consists of the assertion preceding part~(i) that, for the choice described there, $\mathscr{C}^b=0$ and $c_2=0$ if $\mathscr{C}=0$, and of the part of~(iii) that concerns families defined on $\tilde{R}_k$ alone.

In Sections~\ref{sec:definition-step-map}--\ref{sec:symmetries-step-map}, and in Proposition~\ref{prop:local-constituents-step} and its proof, the sets $\mathfrak{X}_k$ and $\mathfrak{X}_{k+1}$ enter in two ways only: by way of the following five properties of a pair $(\mathrm{Sd}_k,\mathrm{Sd}_{k+1})$ in which $\mathrm{Sd}_j$ is a set of complex bond functions on $T_j$ ($j=k,k+1$), and by way of their openness in the space of all complex bond functions on $T_j$ (Lemma~\ref{lem:openness}, valid at the level $k+1$ by Section~\ref{sec:hypotheses-conditions}), whose role is described after the list. Hence the statements and proofs there remain valid when $(\mathfrak{X}_k,\mathfrak{X}_{k+1})$ is replaced by a pair of open sets with these five properties. For such a pair, analyticity on $\tilde{R}_j\times\mathrm{Sd}_j$ and on $\tilde{R}_k^{-}\times\mathrm{Sd}_k$ is holomorphy on an open set, as fixed in Section~\ref{sec:domains-bond-functions}, and the openness of $\mathrm{Sd}_k$ takes the place of the two appeals to Lemma~\ref{lem:openness} for $\mathfrak{X}_k$. For a pair that is not open, clause~\eqref{it:renormalisation-group-1-k1}, which property~\eqref{it:renormalisation-group-e} imposes, has no meaning; an example is the pair formed by the real elements of $\mathfrak{X}_k$ and of $\mathfrak{X}_{k+1}$, which has the properties~\eqref{it:renormalisation-group-a}--\eqref{it:renormalisation-group-d}. The pair $(\{0\},\{0\})$ of the case without slot is not open either; it is not an instance of this principle and is treated after the list.
\begin{itemize}
\item[(a)]\itemmark{a}\label{it:renormalisation-group-a} $0\in\mathrm{Sd}_k$ and $0\in\mathrm{Sd}_{k+1}$ (Section~\ref{sec:local-parts-adjustment}; Proposition~\ref{prop:bounds-regularity-localised}(a)).
\item[(b)]\itemmark{b}\label{it:renormalisation-group-b} The set $\tilde{R}_k+\mathrm{Sd}_k$ is contained in $\tilde{G}_k$ and in the set on which the bound of Proposition~\ref{prop:left-inverse} holds (Lemma~\ref{lem:bounds-analyticity}; Proposition~\ref{prop:bounds-regularity-localised}(b)).
\item[(c)]\itemmark{c}\label{it:renormalisation-group-c} For $j=k$ and for $j=k+1$ the set $\mathrm{Sd}_j$ is invariant under $\mathbb{G}_j$ and under the maps $\mathscr{C}\mapsto-\mathscr{C}$ and $\mathscr{C}\mapsto\overline{\mathscr{C}}$ (Proposition~\ref{prop:bounds-regularity-localised}(b); Section~\ref{sec:invariance-covariance}; at the level $k+1$, Proposition~\ref{prop:symmetries-step-map}, whose proof evaluates the values of the step map at $(g'\mathscr{A},g'\mathscr{C})$, $(-\mathscr{A},-\mathscr{C})$ and $(-\overline{\mathscr{A}},-\overline{\mathscr{C}})$ for $g'\in\mathbb{G}_{k+1}$ and $(\mathscr{A},\mathscr{C})\in\tilde{R}_{k+1}\times\mathrm{Sd}_{k+1}$). For $\mathfrak{X}_k$ and $\mathfrak{X}_{k+1}$ this invariance is Lemma~\ref{lem:symmetries-domains}, at the level $k+1$ by Section~\ref{sec:hypotheses-conditions}.
\item[(d)]\itemmark{d}\label{it:renormalisation-group-d} If $\mathscr{C}'\in\mathrm{Sd}_{k+1}$, then $\mathscr{C}'_L\in\mathrm{Sd}_k$ (Section~\ref{sec:bound-analyticity-after}).
\item[(e)]\itemmark{e}\label{it:renormalisation-group-e} The data of the step map are families with the clauses~\eqref{it:renormalisation-group-1-k1}--\eqref{it:renormalisation-group-1-k6}, and real data also with~\eqref{it:renormalisation-group-1-k7}, in which $\mathrm{Sd}_k$ is the domain of the second slot; the values of the step map are families of the same kind at level $k+1$, in which $\mathrm{Sd}_{k+1}$ is the domain of the second slot; and Proposition~\ref{prop:symmetries-extraction} holds for such data.
\end{itemize}
In these proofs no derivative and no integral with respect to $\mathscr{C}$ is taken, no contour integral in the variable $\mathscr{C}$ occurs, and the identity theorem is not applied in the variable $\mathscr{C}$: the contour integrals and the Cauchy formulas are taken in $t$ and in the parameters $s$, the derivatives and the integrals in $s$, $W$ and $\mathsf{Z}$, and the identity theorem is applied once, in the variable $s$, in the proof of Lemma~\ref{lem:renormalisation-group-2}(d). The openness of $\mathfrak{X}_k$ is invoked through Lemma~\ref{lem:openness} at two places. The first is the proof of part~\eqref{it:bounds-analyticity-q5} of Lemma~\ref{lem:bounds-analyticity}; there it is not needed for $\mathfrak{X}_k$, since the set of the points $(t,A,\mathscr{C})$ considered there is the preimage of the open set $\tilde{G}_k$ under a continuous map. The second is assertion~(3) of Section~\ref{sec:activity-before-reblocking}. Openness is presupposed wherever analyticity in the pair $(A,\mathscr{C})$ is asserted, in particular in clause~\eqref{it:renormalisation-group-1-k1} for the data and for the values; in part~\eqref{it:bounds-analyticity-q5} of Lemma~\ref{lem:bounds-analyticity}; in the paragraphs on continuity and analyticity in the proof of Proposition~\ref{prop:bounds-regularity-localised}, and in part~(e) of that proposition; in assertion~(3) of Section~\ref{sec:activity-before-reblocking}, with its compact set $Q$; in Section~\ref{sec:bound-analyticity-after}, in clause~\eqref{it:renormalisation-group-1-k1} at the level $k+1$ and in the Banach space of the paragraph on the analytic dependence on the data; and in the completeness of the spaces of Definition~\ref{def:renormalisation-group-1}. Its only role is to give a meaning to analyticity in $(A,\mathscr{C})$, which Section~\ref{sec:domains-bond-functions} defines on open sets only, and to carry this analyticity from the data to the values; by the first sentence of this paragraph, no argument of the proofs uses it otherwise. Each step that compares $\mathscr{C}$ with a second value of the second slot, or that applies a transformation to $\mathscr{C}$, requires of the pair only the properties~\eqref{it:renormalisation-group-b}--\eqref{it:renormalisation-group-e}; the comparison of rows at two connections uses at both points that $\tilde{R}_k+\mathrm{Sd}_k$ is contained in $\tilde{G}_k$ (Step~(2) of the proof of Lemma~\ref{lem:renormalisation-group-1}; beginning of the proof of Lemma~\ref{lem:renormalisation-group-2}(b)).

The pair $(\{0\},\{0\})$ has the properties~\eqref{it:renormalisation-group-a}--\eqref{it:renormalisation-group-e}. Property~\eqref{it:renormalisation-group-a} is evident, and~\eqref{it:renormalisation-group-c}, at both levels, and~\eqref{it:renormalisation-group-d} hold because $0$ is fixed by $\mathbb{G}_k$ and by $\mathbb{G}_{k+1}$, by $\mathscr{C}\mapsto-\mathscr{C}$, by $\mathscr{C}\mapsto\overline{\mathscr{C}}$ and by $\mathscr{C}'\mapsto\mathscr{C}'_L$. Property~\eqref{it:renormalisation-group-b} holds because $0\in\mathfrak{X}_k$: the set $\tilde{R}_k+\{0\}$ is contained in $\tilde{R}_k+\mathfrak{X}_k$, which lies in $\tilde{G}_k$ (Lemma~\ref{lem:inclusions-bounds}) and on which the bound of Proposition~\ref{prop:left-inverse} holds. The pair is not open, and the case without slot does not follow from the principle for pairs of open sets; the place of openness is taken by the following reading, which belongs to the definition of the case without slot. The second slot is suppressed, and $\mathscr{C}$ is not a variable. Each of the products $\tilde{R}_k\times\mathfrak{X}_k$, $\tilde{R}_k^{-}\times\mathfrak{X}_k$ and $\tilde{R}_{k+1}\times\mathfrak{X}_{k+1}$ is read as its first factor $D$, also in the presence of further factors $\mathrm{Pd}_k$, a disc in $t$ or $\mathrm{Box}_k$. Each assertion of continuity or of analyticity in $(A,\mathscr{C})$, alone or jointly with $s$, $t$ or $\mathsf{Z}$, is read as the same assertion in $A$ alone on $D$; the set $D$ is open by Lemma~\ref{lem:openness} at the levels $k$ and $k+1$. Assertions of locality in $\mathscr{C}$ and comparisons of two values of the second slot are void, and clause~\eqref{it:renormalisation-group-1-k4} is read with $\lambda'=0$. In particular, this reading of continuity and analyticity applies to clause~\eqref{it:renormalisation-group-1-k1} for the data and for the values; to part~\eqref{it:bounds-analyticity-q5} of Lemma~\ref{lem:bounds-analyticity}; to the paragraphs on continuity and analyticity in the proof of Proposition~\ref{prop:bounds-regularity-localised} and to part~(e) of that proposition; to assertion~(3) of Section~\ref{sec:activity-before-reblocking}, with $Q$ a compact subset of $\tilde{R}_k^{-}$; to clause~\eqref{it:renormalisation-group-1-k1} at the level $k+1$ and to the Banach space of Section~\ref{sec:bound-analyticity-after}; and to the Banach spaces $K^C_j$ of Definition~\ref{def:renormalisation-group-1}. With this reading the pair has property~\eqref{it:renormalisation-group-e}, Proposition~\ref{prop:symmetries-extraction} for such data being its form without slot. The proofs of Sections~\ref{sec:definition-step-map}--\ref{sec:symmetries-step-map} and of Proposition~\ref{prop:local-constituents-step} take from $\mathfrak{X}_k$ and $\mathfrak{X}_{k+1}$ only the properties~\eqref{it:renormalisation-group-a}--\eqref{it:renormalisation-group-e} and openness, and openness serves them only as described in the preceding paragraph; hence each of their steps remains valid under this reading. It follows that all statements of Sections~\ref{sec:definition-step-map}--\ref{sec:symmetries-step-map} and their proofs, and Proposition~\ref{prop:local-constituents-step} and its proof, are valid in the case without slot as well. Section~\ref{sec:proof-clause-theorem} concerns data with slot: it rests on Proposition~\ref{prop:step-external-connection}, which is formulated for such data.

In the case without slot the objects $E^{(i),x}$, $E^{\#,x}_k(m,\mathscr{E})$ and $E^{*,x}_k(m,\mathscr{E})$ of Section~\ref{sec:definition-step-map} are written without the index $x$: $E^{(i)}$, $E^{\#}_k(m,E)$ and $E^{*}_k(m,E)$.

\subsubsection{Recalled auxiliary results}\label{sec:recalled-auxiliary-results}

The hypotheses of the statements recalled in this subsection, and those of Proposition~\ref{prop:step-external-connection}, which is recalled at the beginning of Section~\ref{sec:proof-clause-theorem}, are contained in hypothesis~\ref{def:hypothesis-one-step}. Proposition~\ref{prop:weighted-coefficient-norms} is a statement on Grassmann algebras and needs no hypothesis on the data; Proposition~\ref{prop:polymers} holds for $l\in\mathrm{Cl}$, and its clause~\eqref{it:polymers-t5} under the additional hypothesis $l/L\in\mathrm{Cl}$; Propositions~\ref{prop:symmetries-extraction}, \ref{prop:adjustment}, \ref{prop:contraction}, \ref{prop:comparison-fermion-operators} and~\ref{prop:step-external-connection} hold under hypothesis~\ref{def:hypothesis-one-step} without its system of primitives, the last two by Theorem~\ref{thm:renormalisation-group-2}(b); the conditions on $L$, $M$ and $e_0$ under which Propositions~\ref{prop:basis-map-box}, \ref{prop:symmetries-reality}, \ref{prop:fields-domains}(ii), \ref{prop:left-inverse}, \ref{prop:symmetries-fermion-operators} and~\ref{prop:local-expansion-determinant} hold are contained in stages~\ref{it:order-choice-4} and~\ref{it:order-choice-6} and in the conditions~\eqref{eq:renormalisation-group-31}, \eqref{eq:renormalisation-group-32}, \eqref{eq:renormalisation-group-38}, \eqref{eq:renormalisation-group-A9} and~\eqref{eq:renormalisation-group-A10}; and identity~\eqref{eq:renormalisation-group-3} holds under hypothesis~\ref{def:hypothesis-one-step} without its system of primitives together with the hypotheses of Proposition~\ref{prop:fields-domains}(ii). Proposition~\ref{prop:comparison-fermion-operators} follows from Theorem~\ref{thm:fermion-operators-two} under the hypotheses of Theorem~\ref{thm:expansion-walks}. As conditions on the constants, these hypotheses require that $L\ge5$ be odd, that $M=L^m$ satisfy $M>4L$ and $M\ge108\pi\cdot C_V\cdot a$, where $C_V$ is the constant of Lemma~\ref{lem:perturbations-commutators-holder} and $a$ is the constant defined in~\eqref{eq:operators-fermion-4}, and that conditions~\eqref{eq:renormalisation-group-20} and~\eqref{eq:renormalisation-group-21} hold. The requirements on $L$ and $M$ follow from stages~\ref{it:order-choice-4} and~\ref{it:order-choice-6}, the inequality $M\ge108\pi\cdot C_V\cdot a$ through the bound $M\ge M^F\left(b,\alpha,L,\alpha_0,\kappa_w+\kappa_0\right)$ of stage~\ref{it:order-choice-6} and the definition of $M^F$ in condition~\eqref{eq:renormalisation-group-19}; conditions~\eqref{eq:renormalisation-group-20} and~\eqref{eq:renormalisation-group-21} are named in hypothesis~\ref{def:hypothesis-one-step}, and they follow also from the inequality $e_0\le e^F_0$ of~\eqref{eq:renormalisation-group-A10}, $e^F_0$ being the number of Theorem~\ref{thm:fermion-operators-weakened}. Finally, the assertions of Proposition~\ref{prop:local-expansion-determinant} hold for every family~\eqref{it:primitives-p5} that has the properties~\ref{it:renormalisation-group-1}, \ref{it:renormalisation-group-2}, \ref{it:renormalisation-group-7} and~\ref{it:renormalisation-group-12}--\ref{it:renormalisation-group-14} for the parameters $\mathrm{par}=(\alpha_0,\gamma_w,\kappa_w;C_w)$ of Section~\ref{sec:hypotheses-conditions}, by Proposition~\ref{prop:polymer-expansion-determinant}, which is formulated for a single family~\eqref{it:primitives-p5}, together with Theorem~\ref{thm:fermion-operators}(iii). Here, for $k=0$, property~\ref{it:renormalisation-group-7} is taken in the form stated in property~\ref{it:renormalisation-group-17}, that is, as the kernel bound
\[
\left|S_{0,y}(X,A)(x,x')\right|\le C_w\cdot(1+y)^{-1}\cdot\exp\left(-\gamma_w\cdot d(x,x')-\kappa_w\cdot d_M(X)\right)
\]
for all $y\ge0$, $A\in\tilde{G}_0$, $X\in D_0$ and all sites $x$, $x'$ of $T^0_0$, the modulus of the $4\times4$ block on the left being understood as in Section~\ref{sec:properties-systems-primitives}.

\paragraph{Norms on Grassmann algebras.}
Let $\mathrm{Gr}[\psi,\chi,\Psi,W]$ be the Grassmann algebra over $\mathbb{C}$ with the independent generators of Section~\ref{sec:grassmann-algebras-norms}. These are the generators $\psi(\xi)$, indexed by $\xi=(x,\boldsymbol{\beta},\boldsymbol{\omega})$ with a site $x\in T_k$, a spinor index $\boldsymbol{\beta}$ and $\boldsymbol{\omega}\in\{0,1\}$; the generators $\chi(\boldsymbol{\zeta})$, indexed by $\boldsymbol{\zeta}=(x,x',\boldsymbol{\beta},\boldsymbol{\omega})$ with $(x,x')\in\mathrm{Pz}_k$ (Definition~\ref{def:pairs-sites-holder}); and the generators $\Psi(u)$ and $W(u)$, indexed by $u=(y,\boldsymbol{\beta},\boldsymbol{\omega})$ with a unit site $y\in T^0_k$. The value $\boldsymbol{\omega}=0$ marks the unbarred and the value $\boldsymbol{\omega}=1$ the barred generators. Every index $s$ carries a weight $\mathrm{msw}(s)$, namely
\[
\mathrm{msw}(\xi)=\eta_k^3,\qquad\mathrm{msw}(\boldsymbol{\zeta})=\eta_k^6,\qquad\mathrm{msw}(u)=1,
\]
and $\mathrm{msw}(I):=\prod_{s\in I}\mathrm{msw}(s)$ for a finite set $I$ of indices. For a selection of the four families, the subalgebra spanned by the monomials in the generators of the selected families is denoted by listing these families, as in $\mathrm{Gr}[\psi,\chi]$, $\mathrm{Gr}[\Psi]$ or $\mathrm{Gr}[\psi,\chi,\Psi]$.

Each of the four index sets is totally ordered, and for a subset $I$ of one of them $\psi(I)$, $\chi(I)$, $\Psi(I)$, $W(I)$ denote the products of the corresponding generators, taken in increasing order of their indices. Every $F\in\mathrm{Gr}[\psi,\chi,\Psi,W]$ has a unique representation
\[
F=\sum_{I_1,I_2,I_3,I_4}F(I_1,I_2,I_3,I_4)\cdot\mathrm{msw}(I_1)\cdot\mathrm{msw}(I_2)\cdot\psi(I_1)\,\chi(I_2)\,\Psi(I_3)\,W(I_4),
\]
the sum extending over the subsets $I_1,I_2,I_3,I_4$ of the four index sets. For $\mathsf{w}=(\mathsf{w}_1,\mathsf{w}_2,\mathsf{w}_3,\mathsf{w}_4)$ with $\mathsf{w}_i>0$ put
\[
\|F\|_{\mathsf{w}}:=\sum_{I_1,I_2,I_3,I_4}\mathsf{w}_1^{\#|I_1|}\cdot\mathsf{w}_2^{\#|I_2|}\cdot\mathsf{w}_3^{\#|I_3|}\cdot\mathsf{w}_4^{\#|I_4|}\cdot\left|F(I_1,I_2,I_3,I_4)\right|\cdot\mathrm{msw}(I_1)\cdot\mathrm{msw}(I_2).
\]
Thus $\|F\|_{\mathsf{w}}$ is the sum, over the monomials $\psi(I_1)\,\chi(I_2)\,\Psi(I_3)\,W(I_4)$, of the modulus of the coefficient of the monomial in $F$ multiplied by the product of the numbers $\mathsf{w}_i$ attached to its generators. On one of the subalgebras the norm is written with the components of $\mathsf{w}$ that belong to the families named: $\|F\|_{(\mathsf{w}_1,\mathsf{w}_2)}$ for $F\in\mathrm{Gr}[\psi,\chi]$, $\|F\|_{(\mathsf{w}_1,\mathsf{w}_2,\mathsf{w}_3)}$ for $F\in\mathrm{Gr}[\psi,\chi,\Psi]$. A generator of the $i$-th family has the norm $\mathsf{w}_i$.

Put $\int\mathrm{d}\xi:=\eta_k^3\sum_{\xi}$ and $\int\mathrm{d}\boldsymbol{\zeta}:=\eta_k^6\sum_{\boldsymbol{\zeta}}$. Let $F\in\mathrm{Gr}[\psi,\chi]$ be written in the form
\begin{multline*}
F=\sum_{n,m\ge0}\frac{1}{n!\,m!}\int\mathrm{d}\xi_1\cdots\mathrm{d}\xi_n\,\mathrm{d}\boldsymbol{\zeta}_1\cdots\mathrm{d}\boldsymbol{\zeta}_m\\
E_{nm}(\xi_1,\dots,\xi_n;\boldsymbol{\zeta}_1,\dots,\boldsymbol{\zeta}_m)\,\psi(\xi_1)\cdots\psi(\xi_n)\,\chi(\boldsymbol{\zeta}_1)\cdots\chi(\boldsymbol{\zeta}_m),
\end{multline*}
with kernels $E_{nm}$ that are antisymmetric in $\xi_1,\dots,\xi_n$ and antisymmetric in $\boldsymbol{\zeta}_1,\dots,\boldsymbol{\zeta}_m$. Then $F$ determines these kernels, and
\[
\|F\|_{(\mathsf{w}_1,\mathsf{w}_2)}=\sum_{n,m\ge0}\frac{\mathsf{w}_1^n\cdot\mathsf{w}_2^m}{n!\,m!}\cdot\|E_{nm}\|,\qquad\|E_{nm}\|:=\int\mathrm{d}\xi_1\cdots\mathrm{d}\xi_n\,\mathrm{d}\boldsymbol{\zeta}_1\cdots\mathrm{d}\boldsymbol{\zeta}_m\,\left|E_{nm}(\xi_1,\dots,\xi_n;\boldsymbol{\zeta}_1,\dots,\boldsymbol{\zeta}_m)\right|;
\]
this is the norm $\|\cdot\|_h$ of Section~\ref{sec:grassmann-algebras-norms} with $h=(\mathsf{w}_1,\mathsf{w}_2)$.

Let $(S,\mathrm{msw})$ be an index set with weights, $T$ an index set, and $O$ an operator from the functions on $S$ to the functions on $T$, acting by $(Of)(t)=\sum_{s\in S}O(t,s)\,f(s)\cdot\mathrm{msw}(s)$. Put
\[
|O|_{1,\infty}:=\sup_{t\in T}\sum_{s\in S}\left|O(t,s)\right|\cdot\mathrm{msw}(s);
\]
then $|O|_{1,\infty}=\sup\left\{|Of|_{\infty}:|f|_{\infty}\le1\right\}$, as shown in Section~\ref{sec:coefficients-kernels-grassmann}.

The notation introduced above and the other facts just stated are those of Proposition~\ref{prop:weighted-coefficient-norms}. By the same proposition, for all choices of the weights, the assertions~\eqref{eq:renormalisation-group-40}--\eqref{eq:renormalisation-group-43} below hold, together with the assertion on Banach algebras that accompanies~\eqref{eq:renormalisation-group-40}.

For $F,G\in\mathrm{Gr}[\psi,\chi,\Psi,W]$,
\begin{equation}\label{eq:renormalisation-group-40}
\|FG\|_{\mathsf{w}}\le\|F\|_{\mathsf{w}}\cdot\|G\|_{\mathsf{w}},
\end{equation}
and with each of these norms the even elements form a commutative Banach algebra with unit.

Let $\boldsymbol{\phi}_1,\dots,\boldsymbol{\phi}_p$ and $\boldsymbol{\theta}_1,\dots,\boldsymbol{\theta}_q$ be families of generators, indexed by index sets with weights, and let $O_{ij}$ be operators from the functions on the index set of $\boldsymbol{\theta}_j$ to the functions on the index set of $\boldsymbol{\phi}_i$ with $|O_{ij}|_{1,\infty}\le C_{ij}$. Let $F\mapsto F'$ be the homomorphism from $\mathrm{Gr}[\boldsymbol{\phi}_1,\dots,\boldsymbol{\phi}_p]$ to $\mathrm{Gr}[\boldsymbol{\theta}_1,\dots,\boldsymbol{\theta}_q]$ that replaces every generator $\boldsymbol{\phi}_i(t)$ by
\[
\sum_{j=1}^q(O_{ij}\boldsymbol{\theta}_j)(t)=\sum_{j=1}^q\sum_sO_{ij}(t,s)\,\boldsymbol{\theta}_j(s)\cdot\mathrm{msw}(s),
\]
the inner sum extending over the index set of $\boldsymbol{\theta}_j$. The norms of these two algebras are defined as above, with the weight $\mathsf{w}_i$ for the family $\boldsymbol{\phi}_i$ and the weight $\mathsf{w}'_j$ for the family $\boldsymbol{\theta}_j$. If $\sum_jC_{ij}\cdot\mathsf{w}'_j\le\mathsf{w}_i$ for every $i$, then
\begin{equation}\label{eq:renormalisation-group-41}
\|F'\|_{\mathsf{w}'}\le\|F\|_{\mathsf{w}}\qquad\left(F\in\mathrm{Gr}[\boldsymbol{\phi}_1,\dots,\boldsymbol{\phi}_p]\right).
\end{equation}

For $F\in\mathrm{Gr}[\psi,\chi,\Psi,W]$,
\begin{equation}\label{eq:renormalisation-group-42}
\left\|\int\mathrm{d}\mu_I(W)\,F\right\|_{(\mathsf{w}_1,\mathsf{w}_2,\mathsf{w}_3)}\le\|F\|_{(\mathsf{w}_1,\mathsf{w}_2,\mathsf{w}_3,1)}.
\end{equation}
Here $\int\mathrm{d}\mu_I(W)$ is the Grassmann Gaussian integral over the generators $W$ with identity covariance in $(\bar{W},W)$, where $W_{\boldsymbol{\beta}}(y):=W(y,\boldsymbol{\beta},0)$ and $\bar{W}_{\boldsymbol{\beta}}(y):=W(y,\boldsymbol{\beta},1)$; it is the integral of Proposition~\ref{prop:weighted-coefficient-norms}, which coincides with the integral $\int\mathrm{d}\mu_I(W)$ of Section~\ref{sec:norms-gaussian-integrals}. A set $J$ of indices $u$ is called \emph{paired} if, for every pair $(y,\boldsymbol{\beta})$, it contains $(y,\boldsymbol{\beta},0)$ if and only if it contains $(y,\boldsymbol{\beta},1)$. Put $m(J):=\int\mathrm{d}\mu_I(W)\,W(J)$. By the display~\eqref{eq:polymer-activities-m} of Section~\ref{sec:norms-gaussian-integrals}, $m(J)=0$ if $J$ is not paired, while for a paired set $J$ one has $m(J)\in\{1,-1\}$ and
\[
W(J)=m(J)\cdot\prod_{(y,\boldsymbol{\beta},0)\in J}W_{\boldsymbol{\beta}}(y)\cdot\bar{W}_{\boldsymbol{\beta}}(y).
\]
Hence the integral of a product of generators $W(u)$ is $0$, $1$ or $-1$; in particular all moments are real, and $\int\mathrm{d}\mu_I(W)\,1=m(\emptyset)=1$, since the empty set is paired and $W(\emptyset)$, like the empty product, equals $1$. The two-point functions are
\[
\int\mathrm{d}\mu_I(W)\,W_{\boldsymbol{\beta}}(y)\,W_{\boldsymbol{\beta}'}(y')=0,\qquad\int\mathrm{d}\mu_I(W)\,\bar{W}_{\boldsymbol{\beta}}(y)\,\bar{W}_{\boldsymbol{\beta}'}(y')=0,\qquad\int\mathrm{d}\mu_I(W)\,W_{\boldsymbol{\beta}}(y)\,\bar{W}_{\boldsymbol{\beta}'}(y')=\boldsymbol{\varsigma}\cdot\delta_{yy'}\cdot\delta_{\boldsymbol{\beta}\boldsymbol{\beta}'},
\]
where $\boldsymbol{\varsigma}=1$, as follows from the display~\eqref{eq:polymer-activities-m}. Indeed, $W(u)\cdot W(u)=0$; for distinct indices $u$, $u'$ the product $W(u)\cdot W(u')$ is $W(\{u,u'\})$ or $-W(\{u,u'\})$, and the set $\{u,u'\}$ is paired only if it is $\{(y,\boldsymbol{\beta},0),(y,\boldsymbol{\beta},1)\}$ for a pair $(y,\boldsymbol{\beta})$, the integral of $W(u)\cdot W(u')$ vanishing otherwise; for this set $J$ the identity for $W(J)$ displayed above gives $W_{\boldsymbol{\beta}}(y)\cdot\bar{W}_{\boldsymbol{\beta}}(y)=m(J)\cdot W(J)$, whence $\int\mathrm{d}\mu_I(W)\,W_{\boldsymbol{\beta}}(y)\cdot\bar{W}_{\boldsymbol{\beta}}(y)=m(J)^2=1$.

All moments of the integral are given by Wick's theorem in terms of its two-point functions: if $v_1,\dots,v_r$ are linear combinations of the generators $W(u)$, then $\int\mathrm{d}\mu_I(W)\,v_1\cdots v_r$ is the Pfaffian of the $r\times r$ matrix whose entry at $(j,j')$ is $\int\mathrm{d}\mu_I(W)\,v_j\cdot v_{j'}$. Wick's theorem is one of the classical facts listed at the end of this subsection; there the Pfaffian is defined, and the theorem is derived from~\eqref{eq:polymer-activities-m} for this integral.

Moreover, by Proposition~\ref{prop:weighted-coefficient-norms}\eqref{it:weighted-coefficient-norms-N4p} the integral is invariant under the substitution $W\mapsto UW$, $\bar{W}\mapsto\left(U^{-1}\right)^T\bar{W}$, for every invertible matrix $U$ indexed by the pairs $(y,\boldsymbol{\beta})$; that is, $\int\mathrm{d}\mu_I(W)\,F$ does not change if $W$ and $\bar{W}$ are replaced in $F$ by $UW$ and $\left(U^{-1}\right)^T\bar{W}$. This invariance also follows from Wick's theorem, since the substitution preserves the two-point functions: $U\cdot\left(\boldsymbol{\varsigma}\,\mathbf{1}\right)\cdot U^{-1}=\boldsymbol{\varsigma}\,\mathbf{1}$.

An element $F\in\mathrm{Gr}[\psi,\chi,\Psi,W]$ is said to contain the generator $W(u)$ if $F(I_1,I_2,I_3,I_4)\ne0$ for some subsets with $u\in I_4$. Let $F_1,F_2\in\mathrm{Gr}[\psi,\chi,\Psi,W]$, and assume that the set of the unit sites $y$ for which $F_1$ contains a generator $W(y,\boldsymbol{\beta},\boldsymbol{\omega})$ is disjoint from the set of the unit sites $y$ for which $F_2$ contains a generator $W(y,\boldsymbol{\beta},\boldsymbol{\omega})$. Then
\begin{equation}\label{eq:renormalisation-group-43}
\int\mathrm{d}\mu_I(W)\,F_1F_2=\left(\int\mathrm{d}\mu_I(W)\,F_1\right)\cdot\left(\int\mathrm{d}\mu_I(W)\,F_2\right).
\end{equation}

\paragraph{Polymers.}
Let $\kappa_0\ge c_X\ge1$ and $K_0\ge1$ be the absolute constants of stage~\ref{it:order-choice-1}. By Proposition~\ref{prop:polymers}, the following assertions hold on every torus $T_k$ whose side $l$ lies in $\mathrm{Cl}$; for~\eqref{eq:renormalisation-group-49} it is assumed in addition that $l/L\in\mathrm{Cl}$.

Every $M$-cube $\square$ satisfies
\begin{equation}\label{eq:renormalisation-group-44}
d_M(\square)=0.
\end{equation}
Every polymer $X\in D_k$ satisfies
\begin{equation}\label{eq:renormalisation-group-45}
|X|_M\le c_X\cdot\left(d_M(X)+1\right).
\end{equation}
For every $M$-cube $\square$ and every union $Y$ of $M$-cubes,
\begin{equation}\label{eq:renormalisation-group-46}
\sum_{\substack{X\in D_k\\X\supseteq\square}}\exp\left(-\kappa_0\cdot d_M(X)\right)\le K_0,\qquad\sum_{\substack{X\in D_k\\X\subseteq Y}}\exp\left(-\kappa_0\cdot d_M(X)\right)\le K_0\cdot|Y|_M;
\end{equation}
the second inequality follows from the first, since every polymer contained in $Y$ contains an $M$-cube of $Y$, and $Y$ consists of $|Y|_M$ such cubes.
If $X$ and $Y'$ are polymers with $X\subseteq Y'$, then
\begin{equation}\label{eq:renormalisation-group-47}
d_M(Y')\le d_M(X)+|Y'\setminus X|_M.
\end{equation}
Let $X_1,\dots,X_n$ be polymers, and assume that the graph with the vertex set $\{1,\dots,n\}$ in which two vertices $i\ne j$ are joined whenever $X_i\cup X_j$ is connected through faces is connected. Then
\begin{equation}\label{eq:renormalisation-group-48}
X_1\cup\dots\cup X_n\in D_k,\qquad d_M\left(X_1\cup\dots\cup X_n\right)\le\sum_{i=1}^nd_M(X_i)+3\,(n-1).
\end{equation}
Let also $l/L\in\mathrm{Cl}$. If $X\in D_k$ and $Y\in D_{k+1}$ satisfy $\bar{X}^L=LY$, then
\begin{equation}\label{eq:renormalisation-group-49}
L\cdot d_M(Y)\le d_M(X).
\end{equation}
For every real number $d\ge0$,
\begin{equation}\label{eq:renormalisation-group-50}
d+1\le\exp(d),\qquad d+1\le2\exp\left(\frac{d}{2}\right).
\end{equation}

Two polymers are incompatible when their union is connected through faces; otherwise they are compatible. For polymers $X_1,\dots,X_n$ let $\mathcal{G}(X_1,\dots,X_n)$ be the set of the connected graphs with the vertex set $\{1,\dots,n\}$ each of whose edges $\{i,j\}$ joins two incompatible polymers $X_i$, $X_j$; a graph is identified with the set of its edges. Put $\rho^T(X_1):=1$ and, for $n\ge2$,
\[
\rho^T(X_1,\dots,X_n):=\sum_{G\in\mathcal{G}(X_1,\dots,X_n)}(-1)^{\#|G|}.
\]
As an identity of formal power series in commuting indeterminates $z_X$, $X\in D_k$, one has
\begin{equation}\label{eq:renormalisation-group-51}
\sum_{\mathcal{X}}\,\prod_{X\in\mathcal{X}}z_X=\exp\left(\sum_{n\ge1}\frac{1}{n!}\sum_{(X_1,\dots,X_n)\in D_k^n}\rho^T(X_1,\dots,X_n)\cdot\prod_{i=1}^nz_{X_i}\right),
\end{equation}
where $\mathcal{X}$ ranges over the finite sets of distinct, pairwise compatible polymers, the empty set included. For all polymers $X_1,\dots,X_n$,
\begin{equation}\label{eq:renormalisation-group-52}
\left|\rho^T(X_1,\dots,X_n)\right|\le\#\left|\left\{G\in\mathcal{G}(X_1,\dots,X_n):G\text{ is a tree}\right\}\right|,
\end{equation}
that is, $\left|\rho^T(X_1,\dots,X_n)\right|$ is at most the number of the trees with the vertex set $\{1,\dots,n\}$ all of whose edges join incompatible polymers.

\paragraph{Lattice symmetries.}
The clauses (a), (b) and (d) of Proposition~\ref{prop:symmetries-extraction} state the following properties of the groups and actions of Section~\ref{sec:dirac-matrices-lattice}.
\begin{enumerate}
\item[(a)] Every $g=(S,a)\in\mathbb{G}_k$ maps each of the lattices $L^{-i}\mathbb{Z}^3/l\mathbb{Z}^3$, $-1\le i\le k$, onto itself, blocks onto blocks and centres of blocks onto centres of blocks, the path $\Gamma^{\pi}(y,x)$ onto the path $\Gamma^{\pi}(gy,gx)$, $M$-cubes onto $M$-cubes, and every polymer $X$ onto a polymer $gX$ with $d_M(gX)=d_M(X)$. Every $g\in\mathbb{G}^L_k$ maps, moreover, $LM$-cubes onto $LM$-cubes, and $g(\bar{X}^L)$ is the union of the $LM$-cubes that contain an $M$-cube of $gX$.
\item[(b)] The action of $g$ preserves the sup norms and the pairings. The elements of $\mathrm{Sp}$ and the matrices $\mathsf{C}$ and $\varkappa$ are monomial matrices (each row and each column contains exactly one non-zero entry) whose non-zero entries have modulus one.
\item[(d)] The map $g=(S,a)\mapsto g':=(S,a/L)$ is an isomorphism of $\mathbb{G}^L_k$ onto $\mathbb{G}_{k+1}$, and
\[
(g'A)_L=g(A_L),\qquad(g'\psi)_L=g(\psi_L),\qquad(g'\chi)_L=g(\chi_L),\qquad L(g'Y)=g(LY)
\]
for the bond functions $A$, the generators $\psi$, $\chi$ and the sets $Y$ of level $k+1$, with the scalings $A_L(b)=L^{-\frac{1}{2}}A(b/L)$, $\psi_L(x)=L^{-1}\psi(x/L)$ and $\chi_L(x,x')=L^{-1-\alpha}\chi(x/L,x'/L)$ of Section~\ref{sec:local-parts-adjustment}.
\end{enumerate}

\paragraph{Extraction.}
Let $\mathscr{E}\in K^{x,C}_k$ (Definition~\ref{def:renormalisation-group-1}). By Proposition~\ref{prop:symmetries-extraction}, the number $\varepsilon(\mathscr{E})$ and the matrix that defines $m(\mathscr{E})$ in Section~\ref{sec:local-parts-adjustment}, both of which are formed with an $M$-cube $\square$, do not depend on $\square$, and this matrix is a multiple $m(\mathscr{E})\cdot\mathbf{1}$ of the unit matrix $\mathbf{1}$. The numbers $\varepsilon(\mathscr{E})$ and $m(\mathscr{E})$ are complex linear in $\mathscr{E}$ and real for $\mathscr{E}\in K^x_k$, and
\[
\left|\varepsilon(\mathscr{E})\right|\le c_{\mathrm{fl}}\cdot\|\mathscr{E}\|^x_k,\qquad\left|m(\mathscr{E})\right|\le c_{\mathrm{fl}}\cdot e_k^{\frac{1}{2}}\cdot\|\mathscr{E}\|^x_k.
\]
The operation $R$ of Section~\ref{sec:local-parts-adjustment} maps $K^{x,C}_k$ into itself and $K^x_k$ into itself.

\begin{lemma}[Clifford relations]
Let $\gamma_0,\dots,\gamma_3$, $\mathsf{C}$, $\varkappa$, the group $\mathrm{Sp}$ and the map $r$, $r(S)=\operatorname{diag}\left(r_0(S),r_1(S),r_2(S)\right)$, be those of Section~\ref{sec:dirac-matrices-lattice}, and let $P^{\boldsymbol{\sigma}}_{\nu}=\frac{1}{2}\left(1-\boldsymbol{\sigma}\cdot\gamma_{\nu}\right)$, $\nu\in\{0,1,2\}$, $\boldsymbol{\sigma}\in\{1,-1\}$, be the matrices of Section~\ref{sec:notation}. For a matrix $N$ let $N^T$ denote its transpose and $\overline{N}$ the matrix of the complex conjugates of its entries. Then the following hold.
\begin{enumerate}
\item[(a)] $\mathsf{C}^T=\mathsf{C}^{-1}=-\mathsf{C}$, and $-\gamma_{\nu}^T=\mathsf{C}^{-1}\cdot\gamma_{\nu}\cdot\mathsf{C}$ for $\nu\le3$.
\item[(b)] The matrices $\mathsf{C}$ and $\varkappa$ are real; $\varkappa^T=\varkappa=\varkappa^{-1}$, $\varkappa\cdot\mathsf{C}=-\mathsf{C}\cdot\varkappa$, and $\overline{\gamma_{\nu}}=\varkappa^{-1}\cdot\gamma_{\nu}\cdot\varkappa$ for $\nu\le2$.
\item[(c)] Let $S\in\mathrm{Sp}$. For $\nu\le2$ the matrix $S^{-1}\gamma_{\nu}S$ is $\gamma_{\nu}$ or $-\gamma_{\nu}$, so that $S^{-1}\gamma_{\nu}S=r_{\nu}(S)\cdot\gamma_{\nu}$ with $r_{\nu}(S)\in\{1,-1\}$. Moreover $\left(S^{-1}\right)^T=\pm\mathsf{C}^{-1}\cdot S\cdot\mathsf{C}$ and $\overline{S}=\pm\varkappa^{-1}\cdot S\cdot\varkappa$, with signs that depend on $S$.
\item[(d)] For $\nu\le2$, $\boldsymbol{\sigma}\in\{1,-1\}$ and $S\in\mathrm{Sp}$,
\begin{gather*}
\mathsf{C}^{-1}\cdot\left(P^{\boldsymbol{\sigma}}_{\nu}\right)^T\cdot\mathsf{C}=P^{-\boldsymbol{\sigma}}_{\nu},\\
\overline{P^{\boldsymbol{\sigma}}_{\nu}}=\varkappa^{-1}\cdot P^{\boldsymbol{\sigma}}_{\nu}\cdot\varkappa,\\
S^{-1}P^{\boldsymbol{\sigma}}_{\nu}S=P^{r_{\nu}(S)\boldsymbol{\sigma}}_{\nu}.
\end{gather*}
\end{enumerate}
\end{lemma}

\begin{proof}
By the definitions of Section~\ref{sec:dirac-matrices-lattice}, $\gamma_{\nu}=\sigma_3\otimes\sigma_{\nu+1}$ for $\nu\le2$ and $\gamma_3=\sigma_1\otimes\mathbf{1}$, where $\sigma_1$, $\sigma_2$, $\sigma_3$ are the Pauli matrices. Since the Pauli matrices are Hermitian, anticommute pairwise and have square $\mathbf{1}$, the four matrices $\gamma_{\nu}$ are Hermitian and
\[
\gamma_{\nu}\gamma_{\nu'}+\gamma_{\nu'}\gamma_{\nu}=2\delta_{\nu\nu'}\cdot\mathbf{1}\qquad(\nu,\nu'\le3);
\]
in particular any two of them anticommute and $\gamma_{\nu}^2=\mathbf{1}$. As $\sigma_1$ and $\sigma_3$ are real and symmetric and $\sigma_2$ is purely imaginary and antisymmetric, the matrices $\gamma_0$, $\gamma_2$, $\gamma_3$ are real and symmetric, while $\gamma_1$ is purely imaginary and antisymmetric.

(a) The matrix $\mathsf{C}=i\gamma_1$ is real, $\mathsf{C}^T=i\gamma_1^T=-\mathsf{C}$, and $\mathsf{C}^2=-\gamma_1^2=-\mathbf{1}$; consequently $\mathsf{C}^{-1}=-\mathsf{C}=\mathsf{C}^T$. For $\nu\le3$,
\[
\mathsf{C}^{-1}\cdot\gamma_{\nu}\cdot\mathsf{C}=\left(-i\gamma_1\right)\gamma_{\nu}\left(i\gamma_1\right)=\gamma_1\gamma_{\nu}\gamma_1.
\]
If $\nu\ne1$, the right side equals $-\gamma_{\nu}\gamma_1^2=-\gamma_{\nu}=-\gamma_{\nu}^T$, by anticommutation and since $\gamma_{\nu}$ is symmetric; if $\nu=1$, it equals $\gamma_1=-\gamma_1^T$. This proves (a). Transposing the second relation of (a), and using $\mathsf{C}^T=\mathsf{C}^{-1}$ and $\left(\mathsf{C}^{-1}\right)^T=\mathsf{C}$, one obtains
\[
\mathsf{C}^{-1}\cdot\gamma_{\nu}^T\cdot\mathsf{C}=-\gamma_{\nu}\qquad(\nu\le3).
\]

(b) The matrix $\varkappa=\gamma_3\cdot\mathsf{C}=i\gamma_3\gamma_1$ is real, being the product of the real matrices $\gamma_3$ and $\mathsf{C}$. Further,
\begin{gather*}
\varkappa^T=\mathsf{C}^T\cdot\gamma_3^T=-\mathsf{C}\cdot\gamma_3=-i\gamma_1\gamma_3=i\gamma_3\gamma_1=\varkappa,\\
\varkappa^2=-\gamma_3\gamma_1\gamma_3\gamma_1=\gamma_3^2\gamma_1^2=\mathbf{1},
\end{gather*}
so that $\varkappa^T=\varkappa=\varkappa^{-1}$. One has $\varkappa\cdot\mathsf{C}=\gamma_3\cdot\mathsf{C}^2=-\gamma_3$. On the other hand $\mathsf{C}$, a multiple of $\gamma_1$, anticommutes with $\gamma_3$, whence $\mathsf{C}\cdot\varkappa=\mathsf{C}\cdot\gamma_3\cdot\mathsf{C}=-\gamma_3\cdot\mathsf{C}^2=\gamma_3$; therefore $\varkappa\cdot\mathsf{C}=-\mathsf{C}\cdot\varkappa$. The complex conjugate of a Hermitian matrix is its transpose. For $\nu\le2$ one has $\gamma_3\gamma_{\nu}\gamma_3=-\gamma_{\nu}$, and $\varkappa^{-1}=\left(\gamma_3\cdot\mathsf{C}\right)^{-1}=\mathsf{C}^{-1}\cdot\gamma_3$; together with (a) this gives
\[
\overline{\gamma_{\nu}}=\gamma_{\nu}^T=-\mathsf{C}^{-1}\cdot\gamma_{\nu}\cdot\mathsf{C}=\mathsf{C}^{-1}\cdot\gamma_3\gamma_{\nu}\gamma_3\cdot\mathsf{C}=\varkappa^{-1}\cdot\gamma_{\nu}\cdot\varkappa\qquad(\nu\le2).
\]
In the same way $\varkappa^{-1}\cdot\gamma_3\cdot\varkappa=\mathsf{C}^{-1}\cdot\gamma_3^3\cdot\mathsf{C}=\mathsf{C}^{-1}\cdot\gamma_3\cdot\mathsf{C}=-\gamma_3^T=-\gamma_3$; since $\gamma_3$ is real, it follows that $\overline{\gamma_3}=\gamma_3=-\varkappa^{-1}\cdot\gamma_3\cdot\varkappa$.

(c) Let first $S=\gamma_{\nu'}$ with $\nu'\le3$. Then $S^{-1}=S$, and $S^{-1}\gamma_{\nu}S=\gamma_{\nu'}\gamma_{\nu}\gamma_{\nu'}$ is $\gamma_{\nu}$ for $\nu=\nu'$ and $-\gamma_{\nu}$ for $\nu\ne\nu'$. By (a), $\left(S^{-1}\right)^T=\gamma_{\nu'}^T=-\mathsf{C}^{-1}\cdot S\cdot\mathsf{C}$. By (b) and by the relation for $\gamma_3$ just proved, $\overline{S}=\varkappa^{-1}\cdot S\cdot\varkappa$ if $\nu'\le2$, and $\overline{S}=-\varkappa^{-1}\cdot S\cdot\varkappa$ if $\nu'=3$. Every element of $\mathrm{Sp}$ is a finite product of the matrices $\gamma_0,\dots,\gamma_3$, since each of them is its own inverse. The four maps
\[
S\mapsto\left(S^{-1}\right)^T,\qquad S\mapsto\mathsf{C}^{-1}\cdot S\cdot\mathsf{C},\qquad S\mapsto\overline{S},\qquad S\mapsto\varkappa^{-1}\cdot S\cdot\varkappa
\]
are multiplicative. Moreover, if $S_1^{-1}\gamma_{\nu}S_1$ and $S_2^{-1}\gamma_{\nu}S_2$ are $\pm\gamma_{\nu}$, then $\left(S_1S_2\right)^{-1}\gamma_{\nu}\left(S_1S_2\right)=S_2^{-1}\left(S_1^{-1}\gamma_{\nu}S_1\right)S_2$ is $\gamma_{\nu}$ multiplied by the product of the two signs. Consequently each of the three relations of (c), which holds for the matrices $\gamma_0,\dots,\gamma_3$ with a sign in $\{1,-1\}$, holds for every $S\in\mathrm{Sp}$, with the product of the signs of the factors. In the first relation this sign is $r_{\nu}(S)$, by the definition of $r$.

(d) The last display in the proof of (a) gives
\[
\mathsf{C}^{-1}\cdot\left(P^{\boldsymbol{\sigma}}_{\nu}\right)^T\cdot\mathsf{C}=\frac{1}{2}\left(1-\boldsymbol{\sigma}\cdot\mathsf{C}^{-1}\cdot\gamma_{\nu}^T\cdot\mathsf{C}\right)=\frac{1}{2}\left(1+\boldsymbol{\sigma}\cdot\gamma_{\nu}\right)=P^{-\boldsymbol{\sigma}}_{\nu};
\]
since $\boldsymbol{\sigma}$ is real, (b) gives $\overline{P^{\boldsymbol{\sigma}}_{\nu}}=\frac{1}{2}\left(1-\boldsymbol{\sigma}\cdot\overline{\gamma_{\nu}}\right)=\varkappa^{-1}\cdot P^{\boldsymbol{\sigma}}_{\nu}\cdot\varkappa$; and by (c), $S^{-1}P^{\boldsymbol{\sigma}}_{\nu}S=\frac{1}{2}\left(1-\boldsymbol{\sigma}\cdot r_{\nu}(S)\cdot\gamma_{\nu}\right)=P^{r_{\nu}(S)\boldsymbol{\sigma}}_{\nu}$.
\end{proof}

\paragraph{Adjustment.}
Let $E\in K^{x,C}_k$ be of the form $E=R\,\mathscr{E}$, let $(\mathscr{A},\mathscr{C})\in\tilde{R}_k\times\mathfrak{X}_k$, put $\mathscr{T}:=\mathscr{A}+\mathscr{C}$, and let $X$ be a polymer. By Proposition~\ref{prop:adjustment}, the elements $E^{\mathrm{nat}}(X;\mathscr{A},\mathscr{C},\psi,\chi)$ and $E(X;\mathscr{A},\mathscr{C},\psi,\chi)$ become equal after the substitution
\[
\chi=\delta_{\alpha,\mathscr{T}}\psi,\qquad\bar{\chi}=\delta_{\alpha,-\mathscr{T}}\bar{\psi}.
\]
The kernels of the difference $E^{\mathrm{nat}}(X)-E(X)$ vanish in every bidegree $(n,m)$ other than $(2,0)$, $(1,1)$ and $(0,2)$; in these three bidegrees they are expressions in the kernel $E_{20}(X;\mathscr{A},\mathscr{C};\cdot)$ and in the phases $u_{\mathscr{T}}$ along paths in $X$. Here $E^{\mathrm{nat}}$ is the adjustment of Section~\ref{sec:local-parts-adjustment}, given by Definition~\ref{def:transport-mean}.

\paragraph{Contraction.}
By Proposition~\ref{prop:contraction}, the map $L R$ sends $K^{x,C}_k$ into $K^{x,C}_{k+1}$ and $K^x_k$ into $K^x_{k+1}$, and every $\mathscr{E}\in K^{x,C}_k$ satisfies
\[
\left\|L R\,\mathscr{E}\right\|^x_{k+1}\le c_{\mathrm{fl}}\cdot L^{-\frac{1}{4}+2\varepsilon}\cdot\|\mathscr{E}\|^x_k,\qquad\left(L R\,\mathscr{E}\right)|_0=L R\left(\mathscr{E}|_0\right).
\]

\paragraph{The basis of the fluctuation space.}
By Proposition~\ref{prop:basis-map-box}, $\{\Upsilon\}$ is an orthonormal basis of $V_k$, and every $g\in\mathbb{G}_k$ maps $V_k$ onto itself and the set $\{\pm\Upsilon\}$ onto itself. The support of each $\Upsilon$ lies in one $L$-block, or it consists of one bond that joins two $L$-blocks. Let $C(b',\Upsilon)$ denote the kernel of $C$, so that $(C\mathsf{Z})(b')=\sum_{\Upsilon}C(b',\Upsilon)\cdot\mathsf{Z}(\Upsilon)$ for every bond $b'$ of $T^0_k$. Then $C(b',\Upsilon)=0$ unless $\operatorname{supp}\Upsilon$ is contained in the union of the one or two $L$-blocks that contain the end points of $b'$. Moreover
\[
\sup_{b'}\left|(C\mathsf{Z})(b')\right|\le c_C\cdot\sup_{\Upsilon}\left|\mathsf{Z}(\Upsilon)\right|\qquad(\mathsf{Z}\in V_k),
\]
where the supremum on the left is taken over the bonds $b'$ of $T^0_k$, so that the left side is the norm $|C\mathsf{Z}|_{\infty}$ of the bond function $C\mathsf{Z}$, while $\sup_{\Upsilon}|\mathsf{Z}(\Upsilon)|$ is the largest modulus of a coordinate of $\mathsf{Z}$; the constant $c_C=c_C(L)$ depends on $L$ only. This bound holds, with the same constant, also for the vectors $\mathsf{Z}=\sum_{\Upsilon}\mathsf{Z}(\Upsilon)\cdot\Upsilon$ with complex coordinates. Indeed, $C$ is a real matrix (Proposition~\ref{prop:symmetries-reality}). Fix a bond $b'$, and let $v_0\in V_k$ be the vector with the coordinates $v_0(\Upsilon):=\operatorname{sign}C(b',\Upsilon)$. Then $\sup_{\Upsilon}|v_0(\Upsilon)|\le1$ and $(C\,v_0)(b')=\sum_{\Upsilon}|C(b',\Upsilon)|$, so that the bound at $v_0$ gives $\sum_{\Upsilon}|C(b',\Upsilon)|\le c_C$; hence
\[
\left|(C\mathsf{Z})(b')\right|\le\sum_{\Upsilon}\left|C(b',\Upsilon)\right|\cdot\left|\mathsf{Z}(\Upsilon)\right|\le c_C\cdot\sup_{\Upsilon}\left|\mathsf{Z}(\Upsilon)\right|,
\]
and the supremum over the bonds $b'$ of $T^0_k$ gives the bound for $\mathsf{Z}$.
In the coordinates $\mathsf{Z}(\Upsilon)$ the measures $\mu_I$ and $\hat{\mu}$ are products, over the basis $\{\Upsilon\}$, of copies of one probability measure on $\mathbb{R}$: for $\mu_I$ of the standard normal law, and for $\hat{\mu}$ of its restriction to the interval $[-p_{0,k},p_{0,k}]$, normalised to total mass one. Finally,
\[
0\le d\varepsilon^0_k\le12\exp\left(-\frac{p_{0,k}^2}{2}\right).
\]
Here the hypothesis $p_{0,k}\ge1$ of Proposition~\ref{prop:fluctuation-space}(v) is satisfied: the first inequality $e_k^{\varepsilon}\le\frac{1}{4}$ of~\eqref{eq:renormalisation-group-38} and $\varepsilon\le\frac{1}{200}$ (stages~\ref{it:order-choice-1} and~\ref{it:order-choice-8}) give $e_k\le4^{-\frac{1}{\varepsilon}}\le4^{-200}$, so that $v_k=-\log e_k\ge1$ and $p_{0,k}=v_k^{p_0}\ge1$.

\paragraph{Symmetries of the photon objects.}
By Proposition~\ref{prop:symmetries-reality}, the operator $C$ is real and commutes with $\mathbb{G}_k$; the set $\mathrm{Box}_k$ and the measures $\mu_I$ and $\hat{\mu}$ are invariant under $\mathbb{G}_k$ and under the map $\mathsf{Z}\mapsto-\mathsf{Z}$; and $H_k$ commutes with $\mathbb{G}_k$.

\paragraph{Domains under the fluctuation step.}
By Proposition~\ref{prop:fields-domains}(ii), if $A'\in\mathrm{Dom}_{k+1}$, this domain being formed on the torus of side $l/L$, and if $\zeta\in\mathrm{Box}_k$, then
\[
A^{\min}_k(A')+\mathrm{Sq}_k\zeta\in\mathrm{Dom}_k.
\]

\paragraph{The operators $T_k$ and $T^{(i)}(B)$.}
By Proposition~\ref{prop:left-inverse}, $T_k(B)\cdot H_k(B)=1$ and $\bar{T}_k(B)\cdot\bar{H}_k(B)=1$ for $B\in\tilde{G}_k$. Let now $k\ge1$ and $B\in\tilde{R}_k+\mathfrak{X}_k$, and let $\mathbf{Y}$ be a polymer. Then the operators $T^{(1)}(B)$ and $T^{(2)}(B)$ of~\eqref{eq:renormalisation-group-58} satisfy
\[
\left|T^{(i)}(B)\right|_{1,\infty}\le C_T\qquad(i=1,2),
\]
where $C_T$ is a constant that depends on $(L,b_{-})$ only. In the row of $T^{(i)}(B)$ at a unit site $y$ the connection $B$ appears through the phases $\phi_{B,y}(x)$ and $u_B(x,x^{\boldsymbol{\sigma}})$, $x\in B_k(y)$, and this row reads $B$ on the bonds that have an end point in $B_k(y)$, in the sense of Definition~\ref{def:terminology}(1).

\paragraph{Symmetries of the block phases.}
Let $B\in\tilde{G}_k$ and $g\in\mathbb{G}_k$, and let $\lambda$ be a real function on the sites of $T_k$. By Proposition~\ref{prop:symmetries-fermion-operators}, in its part on $Q_k$, the phases $\phi_{B,y}(x)$ of Section~\ref{sec:notation} satisfy, for every unit site $y$ and every $x\in B_k(y)$,
\begin{gather*}
\phi_{g B,g y}(g x)=\phi_{B,y}(x),\\
\phi_{B-d\lambda,y}(x)=\exp\left(i\,e_k\cdot\left(\lambda(y)-\lambda(x)\right)\right)\cdot\phi_{B,y}(x),\\
\overline{\phi_{\overline{B},y}(x)}=\phi_{-B,y}(x).
\end{gather*}
These identities are the relations
\begin{gather*}
Q_k(g B)\cdot g=g\cdot Q_k(B),\\
Q_k(B-d\lambda)=\exp\left(i\,e_k\cdot\lambda^u\right)\cdot Q_k(B)\cdot\exp\left(-i\,e_k\cdot\lambda\right),\\
\overline{Q_k(\overline{B})}=\varkappa^{-1}\cdot Q_k(-B)\cdot\varkappa
\end{gather*}
of Lemma~\ref{lem:symmetries-elementary-operators}(b), (a) and (d), read off at the kernel $\eta_k^3\,\phi_{B,y}(x)\cdot\mathbf{1}$ of $Q_k(B)$, where $\mathbf{1}$ is the unit matrix in the spinor indices, $\lambda^u$ is the restriction of $\lambda$ to $T^0_k$, and the exponentials act as multiplication operators. All these relations hold for every connection $B$.

\paragraph{The fermion determinant.}
Let $E^d_k$, $E^{\mathrm{det}}_k$ and $\delta E^{\mathrm{det}}_k$ be as in Definition~\ref{def:determinant-functions}, formed with the family~\eqref{it:primitives-p5} of the given system of primitives. By Proposition~\ref{prop:local-expansion-determinant}, the following hold.
\begin{enumerate}
\item[(d0)] For $B\in\tilde{G}_k$ the number $\delta Z_k(B)$ is the exponential of a complex number; in particular $\delta Z_k(B)\ne0$. Moreover
\[
\frac{\delta Z_k(B)}{\delta Z_k(0)}=\exp\left\{\sum_{X\in D_k}E^{\mathrm{det}}_k(X,B)\right\}.
\]
\item[(d5)] $L\,E^{\mathrm{det}}_k\in K_{k+1}$ and $L\,\delta E^{\mathrm{det}}_k\in K^x_{k+1}$.
\item[(d6)] $\left\|L\,E^{\mathrm{det}}_k\right\|_{k+1}\le e_{k+1}^{\frac{1}{4}-2\varepsilon}$ and $\left\|L\,\delta E^{\mathrm{det}}_k\right\|^x_{k+1}\le e_{k+1}^{\frac{1}{4}-2\varepsilon}$.
\end{enumerate}

\paragraph{The rescaled fields of level $k+1$.}
Let $\mathscr{T}'$ be a complex bond function on $T_{k+1}$ such that $\mathscr{T}:=\mathscr{T}'_L$ lies in $\tilde{G}_k$. With the generators $\Psi'$, $\bar{\Psi}'$ of level $k+1$ put
\begin{gather*}
\psi^0:=\left(H_{k+1}(\mathscr{T}')\,\Psi'\right)_L,\qquad\bar{\psi}^0:=\left(\bar{H}_{k+1}(\mathscr{T}')\,\bar{\Psi}'\right)_L,\\
\Psi^c:=T_k(\mathscr{T})\,\psi^0,\qquad\bar{\Psi}^c:=\bar{T}_k(\mathscr{T})\,\bar{\psi}^0.
\end{gather*}
By Proposition~\ref{prop:comparison-fermion-operators}(b),
\begin{equation}\label{eq:renormalisation-group-53}
\psi^0=H_k(\mathscr{T})\,\Psi^c,\qquad\bar{\psi}^0=\bar{H}_k(\mathscr{T})\,\bar{\Psi}^c.
\end{equation}
Since $\Psi^c=T_k(\mathscr{T})\,\psi^0$, the first identity of~\eqref{eq:renormalisation-group-53} reads $H_k(\mathscr{T})\,T_k(\mathscr{T})\,\psi^0=\psi^0$; thus $\psi^0$ belongs to the range of $H_k(\mathscr{T})$. As $\mathscr{T}\in\tilde{G}_k$, Proposition~\ref{prop:left-inverse} gives $T_k(\mathscr{T})\cdot H_k(\mathscr{T})=1$, so that $H_k(\mathscr{T})\cdot T_k(\mathscr{T})$ is a projection onto the range of $H_k(\mathscr{T})$ (it is idempotent, and it fixes $H_k(\mathscr{T})\,w$ for every $w$ in the domain of $H_k(\mathscr{T})$); identity~\eqref{eq:renormalisation-group-53} is the additional fact that $\psi^0$ is fixed by this projection.

\paragraph{Integration of the block transformation.}
Let $A'\in\mathrm{Dom}_{k+1}$ and $\mathscr{C}'\in\mathfrak{X}_{k+1}$, and let $F(\cdot,\mathscr{C}'_L;\Psi)$ be a bounded Borel map from $\mathrm{Dom}_k$ to the even part of $\mathrm{Gr}[\Psi]$. By identity~\eqref{eq:renormalisation-group-3}, the block transformation $\mathbf{T}^x_k$ of Section~\ref{sec:fermion-integrals-blocking} satisfies
\[
\int\mathrm{D}\Psi'\,\left(\mathbf{T}^x_kF\right)(A',\mathscr{C}';\Psi')=\int_{\mathrm{Box}_k}\mathrm{d}\mu_I(\zeta)\int\mathrm{D}\Psi\,F\left(A^{\min}_k(A')+\mathrm{Sq}_k\zeta,\mathscr{C}'_L;\Psi\right).
\]

\paragraph{Classical facts.}
\begin{itemize}
\item Cauchy's integral formula and Cauchy's estimates, for analytic functions of one and of several complex variables; the identity theorem; the analyticity of uniform limits of analytic functions.
\item Fubini's theorem, which remains valid when Berezin integrals take part. In particular, a Berezin integral may be interchanged with the integral, against a finite measure, of a bounded Borel map with values in a Grassmann algebra: the Berezin integral is a linear map from one finite-dimensional space to another, and a linear map of this kind may be taken inside the integral (see~\eqref{eq:algebraic-gaussian-Ber5} in Section~\ref{sec:spinor-spaces-berezin}).
\item Wick's theorem for Grassmann Gaussian integrals, in the following form.

The Pfaffian of an antisymmetric complex $r\times r$ matrix $N=(N_{jj'})$, $r\ge0$, is $1$ if $r=0$ and $0$ if $r$ is odd; if $r=2s$ with $s\ge1$, it is the number
\[
\frac{1}{2^s\cdot s!}\sum_{\pi}\operatorname{sgn}(\pi)\cdot\prod_{a=1}^sN_{\pi(2a-1),\pi(2a)},
\]
the sum being taken over all permutations $\pi$ of $\{1,\dots,r\}$. The exchange of $\pi(2a-1)$ and $\pi(2a)$ changes the sign of $\pi$ and, by the antisymmetry of $N$, the sign of the factor $N_{\pi(2a-1),\pi(2a)}$; the exchange of two of the pairs $\left(\pi(2a-1),\pi(2a)\right)$, $1\le a\le s$, is an even permutation and only reorders the factors. Fix a partition of $\{1,\dots,r\}$ into $s$ pairs. There are $2^s\cdot s!$ permutations $\pi$ for which the sets $\{\pi(2a-1),\pi(2a)\}$, $1\le a\le s$, are the pairs of this partition; any two of them differ by exchanges of the two kinds just described, and therefore their terms in the sum are equal. Consequently the Pfaffian of $N$ equals
\[
\sum_P\operatorname{sgn}(\pi_P)\cdot\prod_{a=1}^sN_{i_a,j_a},
\]
where $P$ ranges over the partitions of $\{1,\dots,r\}$ into $s$ pairs, the pairs of $P$ are written $\{i_a,j_a\}$ with $i_a<j_a$ and $i_1<\dots<i_s$, and $\pi_P$ is the permutation with $\pi_P(2a-1)=i_a$ and $\pi_P(2a)=j_a$ for $1\le a\le s$.

Let $\int\mathrm{d}\mu_I(W)$ be the integral that occurs in~\eqref{eq:renormalisation-group-42}, that is, the integral of Proposition~\ref{prop:weighted-coefficient-norms} and of Section~\ref{sec:norms-gaussian-integrals}; for this integral the sign $\boldsymbol{\varsigma}$ in the two-point functions stated after~\eqref{eq:renormalisation-group-42} equals $1$ (Step~4 below). Let $r\ge0$, and let $v_1,\dots,v_r$ be complex linear combinations of the generators $W(u)$, $u=(y,\boldsymbol{\beta},\boldsymbol{\omega})$, of the unbarred ones ($\boldsymbol{\omega}=0$) as well as of the barred ones ($\boldsymbol{\omega}=1$). Any two such elements $v$, $v'$ satisfy $v\cdot v'=-v'\cdot v$; hence the $r\times r$ matrix $N$ with the entries
\[
N_{jj'}:=\int\mathrm{d}\mu_I(W)\,v_j\cdot v_{j'}\qquad(1\le j,j'\le r)
\]
is antisymmetric. Then $\int\mathrm{d}\mu_I(W)\,v_1\cdots v_r$ is the Pfaffian of $N$.

\emph{Proof.} \emph{Step~1.} Both sides are multilinear functions of $(v_1,\dots,v_r)$. For the left side this holds because the product $v_1\cdots v_r$ is multilinear and the integral is linear. If $r$ is odd, the right side is $0$, and if $r=0$, there are no arguments; for even $r\ge2$, every entry $N_{jj'}$ is linear in $v_j$ and in $v_{j'}$, and every term of the sum defining the Pfaffian of $N$ contains each index $j\in\{1,\dots,r\}$ in exactly one factor.

\emph{Step~2.} Let $\sigma$ be a permutation of $\{1,\dots,r\}$, and replace $(v_1,\dots,v_r)$ by $(v_{\sigma(1)},\dots,v_{\sigma(r)})$. The left side is multiplied by $\operatorname{sgn}(\sigma)$, because the elements $v_j$ anticommute with each other. So is the right side. If $r$ is odd, it is $0$, and if $r=0$, then $\sigma$ is the identity. For even $r\ge2$, the matrix $N$ is replaced by the matrix with the entries $N_{\sigma(j)\sigma(j')}$; in the sum defining the Pfaffian of the latter, the term of $\pi$ is $\operatorname{sgn}(\sigma)$ times the term of $\sigma\pi$ in the sum defining the Pfaffian of $N$, since $\operatorname{sgn}(\pi)=\operatorname{sgn}(\sigma)\cdot\operatorname{sgn}(\sigma\pi)$, and $\sigma\pi$ ranges over all permutations together with $\pi$. In particular, both sides vanish if $v_j=v_{j'}$ for some $j\ne j'$, since the transposition of $j$ and $j'$ then leaves the arguments unchanged and multiplies both sides by $-1$.

\emph{Step~3.} Each $v_j$ is a linear combination of generators. By Step~1 it therefore suffices to prove the identity for $v_j=W(u_j)$, $1\le j\le r$, where $u_1,\dots,u_r$ are indices of generators. If two of these indices coincide, both sides vanish by Step~2, and a change of their order multiplies both sides by the same sign; so the indices may be assumed to be distinct, and they may be listed in any order. Put $J:=\{u_1,\dots,u_r\}$. As in Section~\ref{sec:norms-gaussian-integrals}, $W(J)$ denotes the product of the generators $W(u)$, $u\in J$, in increasing order, and $J$ is called paired if, for every $(y,\boldsymbol{\beta})$, it contains $(y,\boldsymbol{\beta},0)$ if and only if it contains $(y,\boldsymbol{\beta},1)$. By~\eqref{eq:polymer-activities-m}, the number $m(J):=\int\mathrm{d}\mu_I(W)\,W(J)$ is $0$ if $J$ is not paired; if $J$ is paired, then $m(J)\in\{1,-1\}$ and $W(J)=m(J)\cdot P_J$, where $P_J$ is the product of the even elements $W_{\boldsymbol{\beta}}(y)\cdot\bar{W}_{\boldsymbol{\beta}}(y)$ over the indices $(y,\boldsymbol{\beta},0)\in J$. For a paired set $J$ it follows that $P_J=m(J)\cdot W(J)$, as $m(J)^2=1$, and hence
\[
\int\mathrm{d}\mu_I(W)\,P_J=m(J)^2=1.
\]

\emph{Step~4.} For an index $u=(y,\boldsymbol{\beta},\boldsymbol{\omega})$ put $\tilde{u}:=(y,\boldsymbol{\beta},1-\boldsymbol{\omega})$. Then $\tilde{u}\ne u$ and $\tilde{\tilde{u}}=u$, and a set $J$ of indices is paired if and only if $\tilde{u}\in J$ for every $u\in J$. For $b=(y,\boldsymbol{\beta})$ let $W_b:=W_{\boldsymbol{\beta}}(y)$ and $\bar{W}_b:=\bar{W}_{\boldsymbol{\beta}}(y)$, as in Section~\ref{sec:norms-gaussian-integrals}. The set $J_b:=\{(y,\boldsymbol{\beta},0),(y,\boldsymbol{\beta},1)\}$ is paired, and $P_{J_b}=W_b\cdot\bar{W}_b$; as $\bar{W}_b\cdot W_b=-W_b\cdot\bar{W}_b$, Step~3 gives
\[
\int\mathrm{d}\mu_I(W)\,W_b\cdot\bar{W}_b=1,\qquad\int\mathrm{d}\mu_I(W)\,\bar{W}_b\cdot W_b=-1.
\]
A comparison of the first of these identities with the two-point functions stated after~\eqref{eq:renormalisation-group-42} shows that $\boldsymbol{\varsigma}=1$. Moreover, $\int\mathrm{d}\mu_I(W)\,W(u)\cdot W(u')=0$ for all indices $u$, $u'$ with $u'\ne\tilde{u}$, that is, with $u\ne\widetilde{u'}$: if $u'=u$, because $W(u)^2=0$; if $u'\ne u$, because $W(u)\cdot W(u')=\pm W(\{u,u'\})$, while the set $\{u,u'\}$, which contains $u$ but not $\tilde{u}$, is not paired.

\emph{Step~5.} Let $J$ be paired. Then $J$ is the union of the two-element sets $\{u,\tilde{u}\}$, $u\in J$, any two of which are equal or disjoint; so $r$ is even, $r=2s$. If $r=0$, both sides equal $1$, as $\int\mathrm{d}\mu_I(W)\,1=1$. Let $s\ge1$. Then $J$ consists of the indices $(y_a,\boldsymbol{\beta}_a,0)$ and $(y_a,\boldsymbol{\beta}_a,1)$, $1\le a\le s$, for $s$ distinct pairs $b_a=(y_a,\boldsymbol{\beta}_a)$, and the generators may be listed as $v_{2a-1}=W_{b_a}$ and $v_{2a}=\bar{W}_{b_a}$, $1\le a\le s$. The elements $W_{b_a}\cdot\bar{W}_{b_a}$ are even and commute with each other, so that $v_1\cdots v_r=P_J$, and the left side equals $1$ by Step~3. By Step~4, $N_{2a-1,2a}=1$ and $N_{2a,2a-1}=-1$ for $1\le a\le s$, and all other entries of $N$ vanish, because $u_{j'}=\tilde{u}_j$ holds only if $\{j,j'\}$ is one of the sets $\{2a-1,2a\}$. Hence the term of a permutation $\pi$ in the sum defining the Pfaffian of $N$ vanishes unless each of the sets $\{\pi(2a-1),\pi(2a)\}$, $1\le a\le s$, is one of the sets $\{2c-1,2c\}$, $1\le c\le s$. These are the $2^s\cdot s!$ permutations that give the partition of $\{1,\dots,r\}$ into the pairs $\{2c-1,2c\}$, $1\le c\le s$; by what was shown after the definition of the Pfaffian, the term of each of them is the term of the identity permutation, which is $\prod_{a=1}^sN_{2a-1,2a}=1$. Therefore the Pfaffian of $N$ equals $\left(2^s\cdot s!\right)^{-1}\cdot2^s\cdot s!=1$.

\emph{Step~6.} Let $J$ be not paired. Then $v_1\cdots v_r=\pm W(J)$, so that the left side is $\pm m(J)=0$. If $r$ is odd, the Pfaffian of $N$ is $0$ by definition. Let $r$ be even; then $r\ge2$, since the empty set is paired. As $J$ is not paired, there is an index $j$ with $\tilde{u}_j\notin J$. In the term of every permutation in the sum defining the Pfaffian of $N$, the index $j$ occurs in exactly one factor, which is $N_{jj'}$ or $N_{j'j}$ with some $j'\ne j$. As $u_{j'}\in J$, one has $u_{j'}\ne\tilde{u}_j$, and this factor vanishes by Step~4. Hence every term vanishes, and the Pfaffian of $N$ is $0$. This completes the proof.
\item M\"obius inversion on the subsets of a finite set.
\item Cayley's formula for the number of trees with prescribed degrees, which follows from the Pr\"ufer code: if $d_1,\dots,d_n\ge1$ are integers with $\sum_jd_j=2n-2$, the number of the trees with the vertex set $\{1,\dots,n\}$ in which the vertex $j$ has the degree $d_j$ for every $j$ is
\[
\frac{(n-2)!}{\prod_j\left(d_j-1\right)!}.
\]
\end{itemize}

\subsection{The one-step theorem}\label{sec:one-step-theorem}

\begin{theorem}[The renormalisation step on small fields, with an external slot]\label{thm:renormalisation-step-small}
Assume hypothesis~\ref{def:hypothesis-one-step}. The step map is that of Definition~\ref{def:renormalisation-step}, constructed in Section~\ref{sec:definition-step-map} from the system of primitives which this hypothesis provides.
\begin{enumerate}
\item[(a)] \emph{Convergence, bound and analyticity.} Let $(m,\mathscr{E})\in\mathrm{Ball}^x_k(4)$, the ball of Section~\ref{sec:spaces-activities}; for $k=0$ one has $\mathrm{Ball}^x_0(4)=\left\{(m,0):|m|\le 4w^{\mathrm{m}}_0\right\}$, so that $\mathscr{E}=0$. For all $A\in\tilde{R}^{-}_k$, $\mathscr{C}\in\mathfrak{X}_k$ and $X\in D_k$, the series~\eqref{eq:renormalisation-group-61} defining $E^{\#,x}_k(m,\mathscr{E})(X;A,\mathscr{C},\cdot\,)$ converges absolutely in the norm $\|\cdot\|_{h_k/4}$ of Section~\ref{sec:grassmann-algebras-norms}, which is formed with the weights $h_k/4=\left(h_{k,1}/4,\,h_{k,2}/4\right)$. The family $E^{*,x}_k(m,\mathscr{E})$ belongs to $K^{x,C}_{k+1}$, and
\[
\left\|E^{*,x}_k(m,\mathscr{E})\right\|^x_{k+1}\le e_k^{\frac{1}{4}-6\varepsilon}.
\]
If $m$ is real and $\mathscr{E}\in K^x_k$, then $E^{*,x}_k(m,\mathscr{E})$ lies in $K^x_{k+1}$. Let $(m_a,\mathscr{E}_a)$ and $(m_b,\mathscr{E}_b)$ be elements of $\mathbb{C}\times K^{x,C}_k$. The set of those $t\in\mathbb{C}$ for which $(m_a+t\cdot m_b,\mathscr{E}_a+t\cdot\mathscr{E}_b)$ lies in $\mathrm{Ball}^x_k(4)$ is closed, and the map
\[
t\longmapsto E^{*,x}_k\left(m_a+t\cdot m_b,\,\mathscr{E}_a+t\cdot\mathscr{E}_b\right)\in K^{x,C}_{k+1}
\]
is continuous on this set and analytic in its interior. Finally, the number $d\varepsilon^0_k$ of Section~\ref{sec:averaging-fluctuations-gauge} satisfies $0\le d\varepsilon^0_k\le e_k^7$.
\item[(b)] \emph{The case without slot.} For $(m,\mathscr{E})\in\mathrm{Ball}^x_k(4)$ the pair $(m,\mathscr{E}|_0)$ lies in $\mathrm{Ball}_k(4)$, because $\left\|\mathscr{E}|_0\right\|_k\le\|\mathscr{E}\|^x_k$, and
\[
E^{*,x}_k(m,\mathscr{E})|_0=E^{*}_k\left(m,\mathscr{E}|_0\right).
\]
The map $E^{*}_k$ sends $\mathrm{Ball}_k(4)$ into $K^C_{k+1}$, and it sends the pairs $(m,E)\in\mathrm{Ball}_k(4)$ with real $m$ and $E\in K_k$ into $K_{k+1}$. For every $(m,E)\in\mathrm{Ball}_k(4)$,
\[
\left\|E^{*}_k(m,E)\right\|_{k+1}\le e_k^{\frac{1}{4}-6\varepsilon}.
\]
For $(m_a,E_a)$ and $(m_b,E_b)$ in $\mathbb{C}\times K^C_k$, the map $t\longmapsto E^{*}_k\left(m_a+t\cdot m_b,\,E_a+t\cdot E_b\right)\in K^C_{k+1}$ is continuous on the closed set of those $t\in\mathbb{C}$ for which $(m_a+t\cdot m_b,E_a+t\cdot E_b)$ lies in $\mathrm{Ball}_k(4)$, and analytic in the interior of this set.
\item[(c)] \emph{One-step identity.} Let $\boldsymbol{\varepsilon}$ and $m$ be real numbers with $|m|\le w^{\mathrm{m}}_k$, let $\mathscr{E}\in K^x_k$ with $\|\mathscr{E}\|^x_k\le w^{\mathrm{E}}_k$, where $\mathscr{E}=0$ if $k=0$, and let $\mathfrak{P}_k=\left\{\mathfrak{P}_k(X,\cdot\,):X\in D_k\right\}$ be an arbitrary family of functions from $\mathfrak{X}_k$ to $\mathbb{C}$. Put
\begin{gather*}
\boldsymbol{\varepsilon}^{+}:=L^3\cdot\left(\boldsymbol{\varepsilon}+\varepsilon(\mathscr{E})+d\varepsilon^0_k\right),\qquad m^{+}:=L\cdot\left(m+m(\mathscr{E})\right),\\
\mathscr{E}^{+}:=L R\,\mathscr{E}+L\,\delta E^{\mathrm{det}}_k+E^{*,x}_k(m,\mathscr{E}),
\end{gather*}
where $\varepsilon(\mathscr{E})$, $m(\mathscr{E})$ and the maps $L R$ and $L$ on families of polymer activities are those of Section~\ref{sec:local-parts-adjustment}, and $\delta E^{\mathrm{det}}_k$, $E^{\mathrm{det}}_k$ are the families of Definition~\ref{def:determinant-functions}; let $\mathfrak{P}_{k+1}=\left\{\mathfrak{P}_{k+1}(Y,\cdot\,):Y\in D_{k+1}\right\}$ be the family of functions on $\mathfrak{X}_{k+1}$ given by
\[
\mathfrak{P}_{k+1}(Y,\mathscr{C}'):=\sum_{X\in D_k:\ \bar{X}^L=LY}\left(\mathfrak{P}_k(X,\mathscr{C}'_L)+E^{\mathrm{det}}_k(X,\mathscr{C}'_L)\right)\qquad\left(Y\in D_{k+1},\ \mathscr{C}'\in\mathfrak{X}_{k+1}\right);
\]
here $\mathscr{C}'_L$ lies in $\mathfrak{X}_k$ by Lemma~\ref{lem:scaling-domains}(c). Let $\mathscr{F}^x_k$ be the density of Section~\ref{sec:fermion-integrals-blocking} formed with the family $\mathfrak{P}_k$, and let $\mathscr{F}^x_{k+1}$ be the density of level $k+1$ formed with the family $\mathfrak{P}_{k+1}$. Then, for all $A'\in\mathrm{Dom}_{k+1}$ and $\mathscr{C}'\in\mathfrak{X}_{k+1}$,
\begin{align}
\left(\mathbf{T}^x_k\mathscr{F}^x_k[\boldsymbol{\varepsilon},m,\mathscr{E}]\right)(A',\mathscr{C}';\,\cdot\,)&=\mathscr{F}^x_{k+1}[\boldsymbol{\varepsilon}^{+},m^{+},\mathscr{E}^{+}](A',\mathscr{C}';\,\cdot\,),\label{eq:renormalisation-group-54}\\
y^x_{k+1}[\boldsymbol{\varepsilon}^{+},m^{+},\mathscr{E}^{+}](A';\mathscr{C}')&=\int_{\mathrm{Box}_k}\mathrm{d}\mu_I(\zeta)\,y^x_k[\boldsymbol{\varepsilon},m,\mathscr{E}]\left(A^{\min}_k(A')+\mathrm{Sq}_k\zeta;\,\mathscr{C}'_L\right),\label{eq:renormalisation-group-55}
\end{align}
where the harmonic part of $A'$ is arbitrary. If $\mathfrak{P}_k(X,0)=0$ for every $X\in D_k$, then $\mathfrak{P}_{k+1}(Y,0)=0$ for every $Y\in D_{k+1}$, and at $\mathscr{C}'=0$ the identities~\eqref{eq:renormalisation-group-54} and~\eqref{eq:renormalisation-group-55} are the two identities of the same form for the objects $\mathscr{F}_k$, $y_k$, $\mathbf{T}_k$ without slot of Section~\ref{sec:fermion-integrals-blocking}, with the data $(\boldsymbol{\varepsilon},m,\mathscr{E}|_0)$ and their image $(\boldsymbol{\varepsilon}^{+},m^{+},\mathscr{E}^{+}|_0)$; the family $\mathscr{E}^{+}|_0$ is described in~(d).
\item[(d)] \emph{The data of the next level.} Under the hypotheses of~(c), and with its notation, $\mathscr{E}^{+}$ lies in $K^x_{k+1}$,
\[
\mathscr{E}^{+}|_0=L R\,(\mathscr{E}|_0)+L\,E^{\mathrm{det}}_k+E^{*}_k(m,\mathscr{E}|_0),
\]
and
\begin{gather*}
\left\|\mathscr{E}^{+}\right\|^x_{k+1}\le c_{\mathrm{fl}}\cdot L^{-\frac{1}{4}+2\varepsilon}\cdot\|\mathscr{E}\|^x_k+e_{k+1}^{\frac{1}{4}-2\varepsilon}+e_k^{\frac{1}{4}-6\varepsilon},\\
\left|m^{+}-L\cdot m\right|\le L\cdot c_{\mathrm{fl}}\cdot e_k^{\frac{1}{2}}\cdot\|\mathscr{E}\|^x_k,\qquad\left|\boldsymbol{\varepsilon}^{+}-L^3\cdot\boldsymbol{\varepsilon}\right|\le L^3\cdot\left(c_{\mathrm{fl}}\cdot\|\mathscr{E}\|^x_k+e_k^7\right).
\end{gather*}
\item[(e)] \emph{Lipschitz dependence on the couplings.} For $(m_1,\mathscr{E}_1)$ and $(m_2,\mathscr{E}_2)$ in $\mathrm{Ball}^x_k(1)$,
\[
\left\|E^{*,x}_k(m_1,\mathscr{E}_1)-E^{*,x}_k(m_2,\mathscr{E}_2)\right\|^x_{k+1}\le e_k^{\frac{1}{4}-6\varepsilon}\cdot\left(\frac{|m_1-m_2|}{w^{\mathrm{m}}_k}+\frac{\|\mathscr{E}_1-\mathscr{E}_2\|^x_k}{w^{\mathrm{E}}_k}\right).
\]
For $(m_1,E_1)$ and $(m_2,E_2)$ in $\mathrm{Ball}_k(1)$,
\[
\left\|E^{*}_k(m_1,E_1)-E^{*}_k(m_2,E_2)\right\|_{k+1}\le e_k^{\frac{1}{4}-6\varepsilon}\cdot\left(\frac{|m_1-m_2|}{w^{\mathrm{m}}_k}+\frac{\|E_1-E_2\|_k}{w^{\mathrm{E}}_k}\right).
\]
\item[(f)] \emph{Symmetries.} Let $(m,\mathscr{E})\in\mathrm{Ball}^x_k(4)$. The family $E^{*,x}_k(m,\mathscr{E})$ has property~\eqref{it:renormalisation-group-1-k3} at level $k+1$, that is, with respect to the group $\mathbb{G}_{k+1}$; it has property~\eqref{it:renormalisation-group-1-k4} with respect to both connections, and property~\eqref{it:renormalisation-group-1-k5}; and every monomial in it contains as many barred legs as unbarred legs. If $m$ is real and $\mathscr{E}\in K^x_k$, the family $E^{*,x}_k(m,\mathscr{E})$ has in addition property~\eqref{it:renormalisation-group-1-k7}.
\item[(g)] \emph{Locality and analyticity.} Let $(m,\mathscr{E})\in\mathrm{Ball}^x_k(4)$. The family $E^{*,x}_k(m,\mathscr{E})$ has properties~\eqref{it:renormalisation-group-1-k1} and~\eqref{it:renormalisation-group-1-k2} at level $k+1$: for every $Y\in D_{k+1}$, the kernels of $E^{*,x}_k(m,\mathscr{E})(Y;\cdot\,)$ are analytic on $\tilde{R}_{k+1}\times\mathfrak{X}_{k+1}$; they vanish unless all their sites lie in $Y$ and all their pairs are admissible for $Y$, that is, of kind~(I) or of kind~(II) in the sense of Definition~\ref{def:pairs-sites-holder}; and they depend on the two connections only through the values of these connections on the bonds that meet $Y$. Moreover, every assertion of Proposition~\ref{prop:local-constituents-step} holds.
\item[(h)] \emph{Uniformity.} The conditions on the constants on which (a)--(g) rest are conditions of Sections~\ref{sec:order-choice} and~\ref{sec:conditions}. No constant, and none of these conditions, depends on $k$, on $l$, on $\mathfrak{m}$, on the harmonic part of $A'$, on $\mathscr{C}'$, on the system of primitives other than through $\mathrm{par}$, or on the data $(m,\mathscr{E})$.
\end{enumerate}
\end{theorem}

\subsection{Elementary properties of the field domains}\label{sec:elementary-properties-field}

\paragraph{Conventions.}
Let $A$ be a bond function on $T_k$. Then $A(x',x)=-A(x,x')$ for every bond $(x,x')$, and the components of $A$ are $A_{\nu}(x):=A(x,x+\eta_k\cdot e_{\nu})$, $\nu=0,1,2$. For a function $\lambda$ on the sites of $T_k$ the bond function $d\lambda$ is
\[
(d\lambda)(x,x'):=\frac{\lambda(x')-\lambda(x)}{\eta_k}.
\]
For an oriented plaquette $p$ let $(dA)(p)$ be $\eta_k^{-1}$ times the sum of the values of $A$ on the four bonds of the boundary of $p$, each of these bonds being oriented as the boundary of $p$. For the plaquettes oriented from $e_{\nu}$ to $e_{\nu'}$ this means $(dA)_{\nu\nu'}=\partial_{\nu}A_{\nu'}-\partial_{\nu'}A_{\nu}$. Hence $|dA|_{\infty}\le 2|\partial A|_{\infty}$, and $dd\lambda=0$. For a bond function $B$ and a pair $(x,x')\in\mathrm{Pz}_k$ the number $\eta_k\cdot(\tau B)(x,x')$ is, by Definition~\ref{def:pairs-sites-holder}, the mean over the six rectilinear paths $\Gamma^{\pi}(x,x')$ of the sums $\sum_{b\in\Gamma^{\pi}(x,x')}\eta_k\cdot B(b)$. For $B=d\lambda$ each of these six sums telescopes to $\lambda(x')-\lambda(x)$; therefore
\[
\eta_k\cdot(\tau\,d\lambda)(x,x')=\lambda(x')-\lambda(x).
\]

The three norms of a complex bond function $A$ are $|A|_{\infty}$, $|\partial A|_{\infty}$ and $|\delta_{\alpha}\partial A|_{\infty}$ (Section~\ref{sec:domains-bond-functions}); on a set of sites they are formed with the bonds and the differences inside that set. The thresholds of $R_k$ in Definition~\ref{def:domains} are
\[
\boldsymbol{\theta}_1:=e_k^{-\frac{3}{4}+\varepsilon},\qquad \boldsymbol{\theta}_2:=e_k^{-\frac{3}{4}+2\varepsilon},\qquad \boldsymbol{\theta}_3:=e_k^{-\frac{3}{4}+3\varepsilon},
\]
and $\boldsymbol{\theta}_3\le\boldsymbol{\theta}_2\le\boldsymbol{\theta}_1$ whenever $e_k\le 1$, in particular whenever $e_k\le e_0$, by~\eqref{eq:renormalisation-group-38}. For a real number $c$ let $c\cdot R_k$ be the set of the complex bond functions on $T_k$ whose three norms are smaller than $c\cdot\boldsymbol{\theta}_1$, $c\cdot\boldsymbol{\theta}_2$ and $c\cdot\boldsymbol{\theta}_3$, respectively. If $c>0$, this is the set $\{c\cdot A:\ A\in R_k\}$ of Definition~\ref{def:domains}(a), since the three norms are positively homogeneous; if $c\le 0$, the set $c\cdot R_k$ is empty, since the three norms are non-negative. Consequently, if $e_k\ge 1$, the factor $1-e_k^{\varepsilon}$ is not positive, the set $(1-e_k^{\varepsilon})\cdot R_k$ is empty, and so is the set $\tilde{R}^{-}_k$ of Definition~\ref{def:domains}(c), each element of which has a splitting $A_0+A_1$ with $A_1$ in that set.

By Definition~\ref{def:barycentres}, $A_{\nu}(x)$ is a function of the unoriented bond $\{x,x+\eta_k\cdot e_{\nu}\}$, and $(\partial_{\nu'}A_{\nu})(x)$ is a function of the cell spanned by the two parallel bonds $(x,x+\eta_k\cdot e_{\nu})$ and $(x+\eta_k\cdot e_{\nu'},x+\eta_k\cdot e_{\nu'}+\eta_k\cdot e_{\nu})$. This cell is said to be of type $(\nu',\nu)$ and to have the base point $x$. Two cells of the same type, with the base points $x$ and $x'$, have the distance $|x-x'|$, which is the sup-distance of their barycentres.

Let $A_0$ be a real bond function on $T_k$, let $\square$ be an $M$-cube, and let $\square^{\dagger}:=\square^{\sim L}$ (Definition~\ref{def:enlarged-cubes}). A \emph{representative of $A_0$ on $\square^{\dagger}$} is a pair $(A',\lambda)$ of a real bond function $A'$ on $\square^{\dagger}$ and a real function $\lambda$ on the sites of $\square^{\dagger}$ such that $A_0=A'+d\lambda$ on the bonds of $\square^{\dagger}$. Representatives on $\square^{\sim 5}$ are defined in the same way. As $l\in\mathrm{Cl}$, the number $l/(L\cdot M)$ is an integer not smaller than $13$, so that $(2L+1)\cdot M<13\cdot L\cdot M\le l$. Consequently the cube $\square^{\dagger}$, of side $(2L+1)\cdot M$, and the cube $\square^{\sim 5}$, of side $11M$, embed in $T_k$, in the following sense. Let $W$ be one of these two cubes and $s$ its side, let $\tilde{c}\in\mathbb{L}_k$ represent the centre of $\square$, and let $\tilde{W}$ be the set of the points $y\in\mathbb{L}_k$ with $|y-\tilde{c}|<\frac{s}{2}$. The projection of $\mathbb{L}_k$ onto $T_k$ maps $\tilde{W}$ onto the set of the sites of $W$. Two points of $\tilde{W}$ have coordinates that differ by multiples of $\eta_k$ of modulus less than $s$; as $s$ is a multiple of $\eta_k$, the sup-diameter of $\tilde{W}$ is at most $s-\eta_k$. Since $s$ and $13\cdot L\cdot M$ are integers with $s<13\cdot L\cdot M$,
\[
(s-\eta_k)+1<s+1\le 13\cdot L\cdot M\le l.
\]
As $\eta_k\le 1$, it follows that $W$ does not wrap and that $\tilde{W}$ is a lift of $W$, in the sense of Definition~\ref{def:terminology}(9). Let $y,y'\in\tilde{W}$, and let $w\in\mathbb{L}_k$ with $|w|\le 1$ represent the difference of the images of $y$ and $y'$ in $T_k$, that is, $w-(y-y')\in l\mathbb{Z}^3$. Then $w=y-y'$: the vector $w-(y-y')$ of $l\mathbb{Z}^3$ has sup norm at most $1+(s-\eta_k)<l$, whereas every non-zero vector of $l\mathbb{Z}^3$ has sup norm at least $l$. Consequently the following hold.
(i) The choice $w=0$ shows that the projection maps $\tilde{W}$ bijectively onto the set of the sites of $W$. For a site $x$ of $W$ let $\tilde{x}$ denote its preimage in $\tilde{W}$.
(ii) Let $x$ and $x+\eta_k\cdot e_{\nu}$ be sites of $W$. The choice $w=\eta_k\cdot e_{\nu}$ shows that the preimage of $x+\eta_k\cdot e_{\nu}$ is $\tilde{x}+\eta_k\cdot e_{\nu}$. Hence the bonds of $T_k$ with both end points in $W$ are exactly the images of the bonds of $\mathbb{L}_k$ with both end points in $\tilde{W}$. The same holds for the cells of type $(\nu',\nu)$ all of whose sites lie in $W$ and the cells of the same type of $\mathbb{L}_k$ all of whose sites lie in $\tilde{W}$, since the sites $x$, $x+\eta_k\cdot e_{\nu}$, $x+\eta_k\cdot e_{\nu'}$ and $x+\eta_k\cdot e_{\nu'}+\eta_k\cdot e_{\nu}$ of a cell of type $(\nu',\nu)$ with the base point $x$ are joined by the bonds $(x,x+\eta_k\cdot e_{\nu})$, $(x,x+\eta_k\cdot e_{\nu'})$ and $(x+\eta_k\cdot e_{\nu'},x+\eta_k\cdot e_{\nu'}+\eta_k\cdot e_{\nu})$.
(iii) Let $x$ and $x'$ be sites of $W$. Their distance $|x-x'|$ on $T_k$ is the least sup norm of a vector of $\mathbb{L}_k$ that represents $x-x'$, and $\tilde{x}-\tilde{x}'$ is such a vector. If one of the numbers $|x-x'|$ and $|\tilde{x}-\tilde{x}'|$ is at most $1$, then a representative $w$ of $x-x'$ of least sup norm satisfies $|w|\le 1$, hence $w=\tilde{x}-\tilde{x}'$, and $|x-x'|=|\tilde{x}-\tilde{x}'|$.
Thus, when every site $x$ of $W$ is identified with $\tilde{x}$, the bonds of $W$, its cells of each type and the distances at most $1$ between its sites are those of $\tilde{W}$; in particular the bond function $d\lambda$, for a function $\lambda$ on the sites of $W$, and the three norms on $W$ may be computed on $\tilde{W}$.

\paragraph{Hypotheses.}
The constants are those of Section~\ref{sec:hypotheses-conditions}; by that section the conditions~\eqref{eq:renormalisation-group-31} and~\eqref{eq:renormalisation-group-38} hold for every $e_k\in(0,e_0]$. In Lemmas~\ref{lem:openness}, \ref{lem:inclusions-domains} and~\ref{lem:phases-differences-along}, and in Lemma~\ref{lem:symmetries-domains} except for its assertion on $\mathrm{Dom}_k$, the level $k\ge 0$ is arbitrary, $l$ is any element of $\mathrm{Cl}$, and $0<e_k\le e_0$. In Lemma~\ref{lem:scaling-domains} both $l$ and $l/L$ belong to $\mathrm{Cl}$, $0<e_k\le e_0$, and $e_{k+1}=L^{\frac{1}{2}}\cdot e_k$. In Lemma~\ref{lem:periodisation} the level $k\ge 0$ and the number $e_k>0$ are arbitrary, and $l_1$, $l_2$ belong to $\mathrm{Cl}$. These statements involve neither a mass nor a system of primitives. The assertion of Lemma~\ref{lem:symmetries-domains} on $\mathrm{Dom}_k$, and only this assertion, is made under hypothesis~\ref{def:hypothesis-one-step}.

\begin{lemma}[Openness]\label{lem:openness}
Each of the sets $R_k$, $G_k$, $\mathfrak{X}_k$, $c\cdot\tilde{R}_k$ for $0<c\le 1$, $\tilde{R}^{-}_k$ and $\tilde{G}_k$ is open in the finite-dimensional complex vector space formed by all complex bond functions on $T_k$. If $A=A_0+A_1$ is a splitting as in the definition of one of the sets $c\cdot\tilde{R}_k$, $\tilde{R}^{-}_k$, $\tilde{G}_k$, then $\operatorname{Im}A=\operatorname{Im}A_1$.
\end{lemma}

\begin{proof}
The complex bond functions on $T_k$ form a finite-dimensional space, on which the three norms are continuous. The sets $R_k$, $G_k$ and $\mathfrak{X}_k$ of Definition~\ref{def:domains}, and likewise the sets $c\cdot R_k$ with $c\in\mathbb{R}$, are defined by finitely many strict inequalities between these norms and numbers; hence they are open.

Let $A=A_0+A_1$ be an element of $c\cdot\tilde{R}_k$, split as in Definition~\ref{def:domains}: $A_1\in c\cdot R_k$, and $A_0$ is real and has on every cube $\square^{\dagger}$ a representative with the bounds required there. Since $c\cdot R_k$ is open, there is a neighbourhood $U$ of $0$ in the space of the complex bond functions such that $A_1+D\in c\cdot R_k$ for every $D\in U$. For $D\in U$ the function $A+D$ has the splitting $A_0+(A_1+D)$: the summand that the definition requires to be real is again $A_0$, with the same representatives on the cubes $\square^{\dagger}$, and the other summand $A_1+D$ lies in $c\cdot R_k$. So $A+D\in c\cdot\tilde{R}_k$, and $c\cdot\tilde{R}_k$ is open. The same argument, with the open sets $(1-e_k^{\varepsilon})\cdot R_k$ (note $1-e_k^{\varepsilon}>0$ by~\eqref{eq:renormalisation-group-38}) and $G_k$ in place of $c\cdot R_k$, shows that $\tilde{R}^{-}_k$ and $\tilde{G}_k$ are open.

In each of the three definitions the summand $A_0$ is real. Therefore $\operatorname{Im}A=\operatorname{Im}A_1$.
\end{proof}

\begin{lemma}[Inclusions between the domains]\label{lem:inclusions-domains}
\begin{enumerate}
\item[(a)] The sets $\tilde{R}_k$ and $\tilde{R}_k+\mathfrak{X}_k$ are contained in $\tilde{G}_k$.
\item[(b)] The set $\frac{1}{4}R_k$ is contained in $\mathfrak{X}_k$.
\item[(c)] Let $B$ be a complex bond function on $T_k$ with $|B|_{\infty}\le 2e_k^{-\frac{3}{4}+\varepsilon}$, let $s\in\mathbb{C}$ with $|s|\le\frac{1}{4}e_k^{-\frac{1}{4}}$, and let $\mathscr{C}\in\mathfrak{X}_k$. Then $s\cdot B+\mathscr{C}\in G_k$.
\item[(d)] Every $\mathscr{T}\in\tilde{R}_k+\mathfrak{X}_k$ satisfies
\[
e_k\cdot|\operatorname{Im}\mathscr{T}|_{\infty}\le e_k^{\frac{1}{4}+\varepsilon}+e_k^{2\varepsilon},\qquad e_k\cdot|d\mathscr{T}|_{\infty}\le 4e_k^{\frac{1}{4}+2\varepsilon}+2e_k^{3\varepsilon}.
\]
\end{enumerate}
\end{lemma}

\begin{proof}
Let $\mathscr{A}\in\tilde{R}_k$ and $\mathscr{C}\in\mathfrak{X}_k$. By Definition~\ref{def:domains} with $c=1$, $\mathscr{A}=A_0+A_1$, where $A_1\in R_k$ and $A_0$ is real, and for every $M$-cube $\square$ there is a representative $(A'_{\square},\lambda_{\square})$ of $A_0$ on $\square^{\dagger}$ whose three norms on $\square^{\dagger}$ are smaller than $\boldsymbol{\theta}_1$, $\boldsymbol{\theta}_2$, $\boldsymbol{\theta}_3$, respectively. By~\eqref{eq:renormalisation-group-38}, $e_k^{\varepsilon}\le\frac{1}{4}$; in particular $e_k<1$.

\emph{(a)} Let $\square$ be an $M$-cube. Since $L\ge 5$ (stage~\ref{it:order-choice-4}), the cube $\square^{\sim 5}$ is contained in $\square^{\sim L}=\square^{\dagger}$. On $\square^{\sim 5}$ the restriction of $(A'_{\square},\lambda_{\square})$ satisfies $A_0=A'_{\square}+d\lambda_{\square}$ and
\[
|A'_{\square}|_{\infty}<\boldsymbol{\theta}_1=e_k^{-\frac{3}{4}+\varepsilon}\le e_k^{-1+\varepsilon},
\]
the last inequality because $e_k<1$. Moreover, by~\eqref{eq:renormalisation-group-38},
\[
|A_1+\mathscr{C}|_{\infty}<e_k^{-\frac{3}{4}+\varepsilon}+e_k^{-1+2\varepsilon}\le e_k^{-1+\varepsilon},
\]
that is, $A_1+\mathscr{C}\in G_k$. Hence the splitting $\mathscr{A}+\mathscr{C}=A_0+(A_1+\mathscr{C})$ has the properties required in the definition of $\tilde{G}_k$, and $\mathscr{A}+\mathscr{C}\in\tilde{G}_k$. Since $0\in\mathfrak{X}_k$, the choice $\mathscr{C}=0$ gives $\tilde{R}_k\subset\tilde{G}_k$.

\emph{(b)} Let $A\in\frac{1}{4}R_k$. By stage~\ref{it:order-choice-8}, $\varepsilon\le\frac{1}{200}$, so that the exponent $\frac{1}{4}-\varepsilon$ is positive; together with $e_k<1$ this gives $e_k^{\frac{1}{4}-\varepsilon}\le 1\le 4$. It follows that
\[
|A|_{\infty}<\tfrac{1}{4}e_k^{-\frac{3}{4}+\varepsilon}\le e_k^{-1+2\varepsilon},\qquad |\partial A|_{\infty}<\tfrac{1}{4}e_k^{-\frac{3}{4}+2\varepsilon}\le e_k^{-1+3\varepsilon},
\]
and therefore $A\in\mathfrak{X}_k$.

\emph{(c)} One has $|s\cdot B|_{\infty}\le\frac{1}{4}e_k^{-\frac{1}{4}}\cdot 2e_k^{-\frac{3}{4}+\varepsilon}=\frac{1}{2}e_k^{-1+\varepsilon}$ and $|\mathscr{C}|_{\infty}<e_k^{-1+2\varepsilon}$. As $e_k^{\varepsilon}\le\frac{1}{2}$,
\[
|s\cdot B+\mathscr{C}|_{\infty}<\left(\tfrac{1}{2}+e_k^{\varepsilon}\right)\cdot e_k^{-1+\varepsilon}\le e_k^{-1+\varepsilon},
\]
so that $s\cdot B+\mathscr{C}\in G_k$.

\emph{(d)} Let $\mathscr{T}=\mathscr{A}+\mathscr{C}$. By Lemma~\ref{lem:openness}, $\operatorname{Im}\mathscr{A}=\operatorname{Im}A_1$, and so $\operatorname{Im}\mathscr{T}=\operatorname{Im}A_1+\operatorname{Im}\mathscr{C}$. Hence
\[
e_k\cdot|\operatorname{Im}\mathscr{T}|_{\infty}\le e_k\cdot\left(|A_1|_{\infty}+|\mathscr{C}|_{\infty}\right)<e_k\cdot\left(e_k^{-\frac{3}{4}+\varepsilon}+e_k^{-1+2\varepsilon}\right)=e_k^{\frac{1}{4}+\varepsilon}+e_k^{2\varepsilon}.
\]
For the second bound let $p$ be a plaquette of $T_k$. By fact~\eqref{it:renormalisation-group-g0} of Section~\ref{sec:cubes-parameter-cubes}, a corner of $p$ lies in an $M$-cube $\square$. The other corners of $p$ have sup-distance at most $\eta_k\le 1<M$ from this corner; hence $p$ lies in $\square^{\sim 1}$, and so in $\square^{\dagger}$. As $dd\lambda_{\square}=0$, one has $(dA_0)(p)=(dA'_{\square})(p)$, and $(dA'_{\square})(p)$ is formed with differences of $A'_{\square}$ inside $\square^{\dagger}$. Therefore
\[
|(dA_0)(p)|\le 2|\partial A'_{\square}|_{\infty}<2\boldsymbol{\theta}_2,
\]
the norm of $A'_{\square}$ being taken on $\square^{\dagger}$. Since $p$ was arbitrary, $|dA_0|_{\infty}\le 2\boldsymbol{\theta}_2$. Further, $|dA_1|_{\infty}\le 2|\partial A_1|_{\infty}<2\boldsymbol{\theta}_2$ and $|d\mathscr{C}|_{\infty}\le 2|\partial\mathscr{C}|_{\infty}<2e_k^{-1+3\varepsilon}$. Adding the three bounds and multiplying by $e_k$ one obtains
\[
e_k\cdot|d\mathscr{T}|_{\infty}\le e_k\cdot\left(4e_k^{-\frac{3}{4}+2\varepsilon}+2e_k^{-1+3\varepsilon}\right)=4e_k^{\frac{1}{4}+2\varepsilon}+2e_k^{3\varepsilon}.
\]
Lemma~\ref{lem:inclusions-bounds} follows.
\end{proof}

\begin{lemma}[Scaling of the domains]\label{lem:scaling-domains}
Let $l$ and $l/L$ lie in $\mathrm{Cl}$, and let $0<c\le 1$. The domains of level $k+1$ are those of the torus $T_{k+1}=\mathbb{L}_{k+1}/(l/L)\mathbb{Z}^3$, formed with the charge $e_{k+1}=L^{\frac{1}{2}}\cdot e_k$. For a bond function $A$ on $T_{k+1}$ let $A_L$ be the bond function on $T_k$ given by $A_L(b):=L^{-\frac{1}{2}}\cdot A(b/L)$ (Section~\ref{sec:local-parts-adjustment}), and put $c':=c\cdot L^{-\frac{3}{4}}$.
\begin{enumerate}
\item[(a)] If $A\in c\cdot R_{k+1}$, then $A_L\in c'\cdot R_k$. If $\mathscr{A}\in c\cdot\tilde{R}_{k+1}$, then $\mathscr{A}_L\in c'\cdot\tilde{R}_k$.
\item[(b)] The set $L^{-\frac{3}{4}}\cdot\tilde{R}_k$ is contained in $\tilde{R}^{-}_k$.
\item[(c)] If $\mathscr{C}\in\mathfrak{X}_{k+1}$, then
\[
|\mathscr{C}_L|_{\infty}<L^{-\frac{3}{4}}\cdot e_k^{-1+2\varepsilon},\qquad |\partial\mathscr{C}_L|_{\infty}<L^{-\frac{3}{4}}\cdot e_k^{-1+3\varepsilon},
\]
that is, the two norms of $\mathscr{C}_L$ that define $\mathfrak{X}_k$ are smaller than $L^{-\frac{3}{4}}$ times the corresponding thresholds of $\mathfrak{X}_k$; in particular $\mathscr{C}_L\in\mathfrak{X}_k$.
\end{enumerate}
\end{lemma}

\begin{proof}
One has $\eta_{k+1}=\eta_k/L$, and $x\mapsto x/L$ maps the sites and the bonds of $T_k$ bijectively onto those of $T_{k+1}$. Let $A$ be a complex bond function on $T_{k+1}$. Then $|A_L|_{\infty}=L^{-\frac{1}{2}}\cdot|A|_{\infty}$ and $(A_L)_{\nu}(x)=L^{-\frac{1}{2}}\cdot A_{\nu}(x/L)$, whence
\[
\left(\partial_{\nu'}(A_L)_{\nu}\right)(x)=\eta_k^{-1}\cdot L^{-\frac{1}{2}}\cdot\left(A_{\nu}\left(\tfrac{x}{L}+\eta_{k+1}\cdot e_{\nu'}\right)-A_{\nu}\left(\tfrac{x}{L}\right)\right)=L^{-\frac{3}{2}}\cdot\left(\partial_{\nu'}A_{\nu}\right)\left(\tfrac{x}{L}\right),
\]
and $|\partial A_L|_{\infty}=L^{-\frac{3}{2}}\cdot|\partial A|_{\infty}$. Let two cells of $T_k$ of the same type have the base points $x$ and $x'$, with $0<|x-x'|\le 1$. The cells of $T_{k+1}$ of the same type with the base points $x/L$ and $x'/L$ have the distance $|x-x'|/L\le 1/L\le 1$, and $|x-x'|^{-\alpha}=L^{-\alpha}\cdot|x/L-x'/L|^{-\alpha}$. Together with the formula for $\partial_{\nu'}(A_L)_{\nu}$ this gives
\[
|\delta_{\alpha}\partial A_L|_{\infty}\le L^{-\frac{3}{2}-\alpha}\cdot|\delta_{\alpha}\partial A|_{\infty}.
\]
The same three relations hold between the norms of $A_L$ on the set $LW$ and the norms of $A$ on $W$, for every set $W$ of sites of $T_{k+1}$.

Since $e_{k+1}=L^{\frac{1}{2}}\cdot e_k$, the thresholds of $R_{k+1}$ are
\[
e_{k+1}^{-\frac{3}{4}+\varepsilon}=L^{-\frac{3}{8}+\frac{\varepsilon}{2}}\cdot\boldsymbol{\theta}_1,\qquad e_{k+1}^{-\frac{3}{4}+2\varepsilon}=L^{-\frac{3}{8}+\varepsilon}\cdot\boldsymbol{\theta}_2,\qquad e_{k+1}^{-\frac{3}{4}+3\varepsilon}=L^{-\frac{3}{8}+\frac{3\varepsilon}{2}}\cdot\boldsymbol{\theta}_3.
\]
Let $A\in c\cdot R_{k+1}$. By the three relations, the three norms of $A_L$ are smaller than
\[
c\cdot L^{-\frac{7}{8}+\frac{\varepsilon}{2}}\cdot\boldsymbol{\theta}_1,\qquad c\cdot L^{-\frac{15}{8}+\varepsilon}\cdot\boldsymbol{\theta}_2,\qquad c\cdot L^{-\frac{15}{8}-\alpha+\frac{3\varepsilon}{2}}\cdot\boldsymbol{\theta}_3,
\]
respectively. As $\varepsilon\le\frac{1}{200}$ (stage~\ref{it:order-choice-8}) and $\alpha>0$, each of the three exponents of $L$ is at most $-\frac{3}{4}$. Hence $A_L\in c'\cdot R_k$, which is the first assertion of (a).

Let now $\mathscr{A}=A_0+A_1\in c\cdot\tilde{R}_{k+1}$, split as in Definition~\ref{def:domains}, so that $A_1\in c\cdot R_{k+1}$ and $A_0$ is real. Then $\mathscr{A}_L=A_{0,L}+A_{1,L}$, where $A_{0,L}:=(A_0)_L$ is real and $A_{1,L}:=(A_1)_L$ lies in $c'\cdot R_k$ by the first assertion. Let $\square$ be an $M$-cube of $T_k$. By fact~\eqref{it:renormalisation-group-g0} of Section~\ref{sec:cubes-parameter-cubes}, $\square$ lies in one $LM$-cube of $T_k$. Since $l/(L\cdot M)$ is an integer, the map $x\mapsto x/L$ carries the $LM$-cubes of $T_k$ onto the $M$-cubes of $T_{k+1}$; so there is an $M$-cube $\square'$ of $T_{k+1}$ with $\square\subset L\square'$. The set $L\square'^{\dagger}$ is a cube with the same centre as the $LM$-cube $L\square'$ and with the side $(2L+1)\cdot L\cdot M$; it contains every site whose sup-distance from $L\square'$ is at most $L\cdot M$. Every site of $\square^{\dagger}$ has sup-distance at most $L\cdot M$ from $\square$, hence from $L\square'$. Therefore $\square^{\dagger}\subset L\square'^{\dagger}$.

By Definition~\ref{def:domains}, $A_0$ has a representative $(A',\lambda)$ on $\square'^{\dagger}$ whose three norms on $\square'^{\dagger}$ are smaller than $c$ times the thresholds of $R_{k+1}$. For the sites $x$ of $L\square'^{\dagger}$ put $\lambda_L(x):=L^{\frac{1}{2}}\cdot\lambda(x/L)$. Then $(d\lambda)_L=d\lambda_L$ on the bonds of $L\square'^{\dagger}$, because
\[
L^{-\frac{1}{2}}\cdot\eta_{k+1}^{-1}=L^{\frac{1}{2}}\cdot\eta_k^{-1}.
\]
So $A_{0,L}=A'_L+d\lambda_L$ on the bonds of $L\square'^{\dagger}$, and in particular on those of $\square^{\dagger}$. By the scaling relations with $W=\square'^{\dagger}$ and by the expressions for the thresholds of $R_{k+1}$, exactly as in the proof of the first assertion, the three norms of $A'_L$ on $L\square'^{\dagger}$, and a fortiori on $\square^{\dagger}$, are smaller than $c'\cdot\boldsymbol{\theta}_1$, $c'\cdot\boldsymbol{\theta}_2$, $c'\cdot\boldsymbol{\theta}_3$, respectively. Thus the restriction of $(A'_L,\lambda_L)$ to $\square^{\dagger}$ is a representative of $A_{0,L}$ on $\square^{\dagger}$ with the bounds required in the definition of $c'\cdot\tilde{R}_k$. As $\square$ was arbitrary and $0<c'\le 1$, it follows that $\mathscr{A}_L\in c'\cdot\tilde{R}_k$. This proves (a).

\emph{(b)} An element of $L^{-\frac{3}{4}}\cdot\tilde{R}_k$ has a splitting $A_0+A_1$ with $A_1\in L^{-\frac{3}{4}}\cdot R_k$, in which $A_0$ has on every cube $\square^{\dagger}$ a representative whose three norms are smaller than $L^{-\frac{3}{4}}\cdot\boldsymbol{\theta}_i\le\boldsymbol{\theta}_i$, $i=1,2,3$. By stage~\ref{it:order-choice-4}, $L^{-\frac{3}{4}}\le\frac{1}{2}$, and by~\eqref{eq:renormalisation-group-38}, $e_k^{\varepsilon}\le\frac{1}{4}$; hence $L^{-\frac{3}{4}}\le\frac{1}{2}\le 1-e_k^{\varepsilon}$, and $A_1\in(1-e_k^{\varepsilon})\cdot R_k$. These are the conditions that define $\tilde{R}^{-}_k$.

\emph{(c)} Let $\mathscr{C}\in\mathfrak{X}_{k+1}$. By the scaling relations and by $e_{k+1}=L^{\frac{1}{2}}\cdot e_k$,
\[
|\mathscr{C}_L|_{\infty}<L^{-\frac{1}{2}}\cdot e_{k+1}^{-1+2\varepsilon}=L^{-1+\varepsilon}\cdot e_k^{-1+2\varepsilon},\qquad |\partial\mathscr{C}_L|_{\infty}<L^{-\frac{3}{2}}\cdot e_{k+1}^{-1+3\varepsilon}=L^{-2+\frac{3\varepsilon}{2}}\cdot e_k^{-1+3\varepsilon}.
\]
Since $\varepsilon\le\frac{1}{200}$, both powers of $L$ are at most $L^{-\frac{3}{4}}$, which proves the two bounds of (c); and $L^{-\frac{3}{4}}\le 1$ gives $\mathscr{C}_L\in\mathfrak{X}_k$.

Lemma~\ref{lem:rescaling} follows, by (a) with $c=1$, (b) and (c).
\end{proof}

\begin{lemma}[Symmetries of the domains]\label{lem:symmetries-domains}
The sets $R_k$, $c\cdot\tilde{R}_k$ with $0<c\le 1$, $\tilde{R}^{-}_k$, $G_k$, $\tilde{G}_k$, $\mathfrak{X}_k$ and $\mathrm{Dom}_k$ are invariant under the action of $\mathbb{G}_k$, under the map $A\mapsto-A$ and under the map $A\mapsto\overline{A}$. Moreover, each of the sets $c\cdot\tilde{R}_k$ with $0<c\le 1$, $\tilde{R}^{-}_k$ and $\tilde{G}_k$ contains, along with any of its elements $A$, every function $A-d\lambda$, where $\lambda$ is a real function on the sites of $T_k$, and every sum of $A$ with a real constant bond function.
\end{lemma}

\begin{proof}
\emph{The action of $\mathbb{G}_k$.} Let $g=(S,a)\in\mathbb{G}_k$, and let $r:=r(S)=\operatorname{diag}(r_0,r_1,r_2)$, so that $gx=r\cdot x+a$ (Section~\ref{sec:dirac-matrices-lattice}). As $r$ is diagonal with entries $\pm 1$ and $a\in M\mathbb{Z}^3/l\mathbb{Z}^3$, the map $g$ is an isometry of $T_k$ for the sup-distance; it maps bonds onto bonds and $M$-cubes onto $M$-cubes, and $g(\square^{\sim i})=(g\square)^{\sim i}$ for every $M$-cube $\square$ and every $i\ge 0$. On bond functions $g$ acts by $(gA)(x,x')=A(g^{-1}x,g^{-1}x')$. Regard $A_{\nu}$ as a function of the unoriented bonds of direction $\nu$ and $\partial_{\nu'}A_{\nu}$ as a function of the cells of type $(\nu',\nu)$, as in the conventions above (Definition~\ref{def:barycentres}). Let $b$ be a bond of direction $\nu$. Then $gb$ is a bond of direction $\nu$, and the value of $(gA)_{\nu}$ at $gb$ is $r_{\nu}$ times the value of $A_{\nu}$ at $b$: for $r_{\nu}=-1$ the bond $gb$ is traversed against $e_{\nu}$, and $A$ is odd under reversal. Further, $g$ maps the cell of type $(\nu',\nu)$ spanned by the bonds $b$ and $b+\eta_k\cdot e_{\nu'}$ onto a cell of the same type, and the value of $\partial_{\nu'}(gA)_{\nu}$ at the image cell is $r_{\nu}\cdot r_{\nu'}$ times the value of $\partial_{\nu'}A_{\nu}$ at the given cell: the factor $r_{\nu'}$ appears because for $r_{\nu'}=-1$ the order of the two bonds in the direction $e_{\nu'}$ is reversed. As $g$ preserves the distances of the barycentres, the three norms of $gA$ on $T_k$ are those of $A$, and the three norms of $gA$ on $gW$ are those of $A$ on $W$, for every set $W$ of sites. Moreover $g(d\lambda)=d(\lambda\circ g^{-1})$ for every function $\lambda$ on the sites.

The sets $R_k$, $G_k$ and $\mathfrak{X}_k$, and the sets $c\cdot R_k$ and $(1-e_k^{\varepsilon})\cdot R_k$, are defined through the three norms alone, and so they are invariant under $g$. Let $A=A_0+A_1$ be split as in the definition of $c\cdot\tilde{R}_k$. Then $gA=gA_0+gA_1$, where $gA_0$ is real and $gA_1\in c\cdot R_k$. If $(A',\lambda)$ is a representative of $A_0$ on $\square^{\dagger}$, then $(gA',\lambda\circ g^{-1})$ is a representative of $gA_0$ on $g(\square^{\dagger})=(g\square)^{\dagger}$, with the same three norms. As $\square$ ranges over the $M$-cubes, so does $g\square$; hence $gA\in c\cdot\tilde{R}_k$. The same reasoning applies to $\tilde{R}^{-}_k$, and to $\tilde{G}_k$ with the cubes $\square^{\sim 5}$ and $(g\square)^{\sim 5}$ in place of $\square^{\dagger}$ and $(g\square)^{\dagger}$.

\emph{Sign and complex conjugation.} The maps $A\mapsto-A$ and $A\mapsto\overline{A}$ preserve the three norms on every set of sites and map real bond functions to real bond functions. Under them a splitting $A_0+A_1$ becomes $(-A_0)+(-A_1)$, respectively $A_0+\overline{A_1}$; a representative $(A',\lambda)$ of $A_0$ becomes the representative $(-A',-\lambda)$ of $-A_0$, respectively remains a representative of $A_0$. This gives the invariance of $R_k$, $c\cdot\tilde{R}_k$, $\tilde{R}^{-}_k$, $G_k$, $\tilde{G}_k$ and $\mathfrak{X}_k$ under the two maps.

\emph{The set $\mathrm{Dom}_k$.} By Definition~\ref{def:domains}, $\mathrm{Dom}_k$ is the set of the real bond functions $A$ on $T^0_k$ with $H_kA\in\frac{1}{4}\tilde{R}_k$. By property~\ref{it:renormalisation-group-1}, $H_k=H_k(1)$ is a member of the family~\eqref{it:primitives-p1}, which consists of linear operators; at level $0$ one has $H_0=1$ (Section~\ref{sec:averaging-fluctuations-gauge}). Thus $H_k$ is linear, and it commutes with the action of $\mathbb{G}_k$ by Proposition~\ref{prop:symmetries-reality}. Let $A\in\mathrm{Dom}_k$ and $g\in\mathbb{G}_k$. The functions $gA$ and $-A$ are real, and $H_k(gA)=g(H_kA)$ and $H_k(-A)=-H_kA$ lie in $\frac{1}{4}\tilde{R}_k$, by the invariance of $\frac{1}{4}\tilde{R}_k$ proved above. Finally $\overline{A}=A$, since $A$ is real.

\emph{Gauge transformations and constants.} Let $A=A_0+A_1$ be split as in the definition of $c\cdot\tilde{R}_k$, of $\tilde{R}^{-}_k$ or of $\tilde{G}_k$, and for every $M$-cube $\square$ let $(A'_{\square},\lambda_{\square})$ be the representative of $A_0$ on $\square^{\dagger}$ (on $\square^{\sim 5}$, in the case of $\tilde{G}_k$) that the definition provides. Let $\lambda$ be a real function on the sites of $T_k$. Then $A-d\lambda=(A_0-d\lambda)+A_1$, the summand $A_0-d\lambda$ is real, and
\[
A_0-d\lambda=A'_{\square}+d(\lambda_{\square}-\lambda)
\]
on the bonds of $\square^{\dagger}$ (of $\square^{\sim 5}$). Thus $(A'_{\square},\lambda_{\square}-\lambda)$ is a representative of $A_0-d\lambda$ with the same norms, and $A-d\lambda$ lies in the same set as $A$. Next let $\Omega$ be a real constant bond function: $\Omega(x,x+\eta_k\cdot e_{\nu})=\Omega_{\nu}$ for all sites $x$ and all $\nu$, with $(\Omega_0,\Omega_1,\Omega_2)\in\mathbb{R}^3$. The cube $\square^{\dagger}$ embeds in $T_k$; let $\tilde{W}$ be a lift of $\square^{\dagger}$ as in the conventions above, and for a site $x$ of $\square^{\dagger}$, with the preimage $\tilde{x}\in\tilde{W}$, let $\Omega\cdot x:=\sum_{\nu}\Omega_{\nu}\cdot\tilde{x}_{\nu}$. By property~(ii) of the conventions, every bond of $T_k$ with both end points in $\square^{\dagger}$ is, for some $\nu$, a bond $(x,x')$ or $(x',x)$ with $\tilde{x}'=\tilde{x}+\eta_k\cdot e_{\nu}$, that is, the image of the bond $(\tilde{x},\tilde{x}+\eta_k\cdot e_{\nu})$ of $\mathbb{L}_k$ or of its reversal, the end points lying in $\tilde{W}$. For such sites one has $x'=x+\eta_k\cdot e_{\nu}$, and the value of $d(\Omega\cdot x)$ at the bond $(x,x')$ is
\[
\eta_k^{-1}\cdot\sum_{\mu}\Omega_{\mu}\cdot\left(\tilde{x}'_{\mu}-\tilde{x}_{\mu}\right)=\Omega_{\nu}=\Omega(x,x').
\]
Both $\Omega$ and $d(\Omega\cdot x)$ change sign when a bond is reversed. Therefore $\Omega=d(\Omega\cdot x)$ on the bonds of $\square^{\dagger}$. Hence $A+\Omega=(A_0+\Omega)+A_1$, where $A_0+\Omega$ is real and
\[
A_0+\Omega=A'_{\square}+d(\lambda_{\square}+\Omega\cdot x)
\]
on the bonds of $\square^{\dagger}$; the same holds on $\square^{\sim 5}$, which embeds as well, with $\Omega\cdot x$ formed with a lift of $\square^{\sim 5}$ as in the conventions above. So $(A'_{\square},\lambda_{\square}+\Omega\cdot x)$ is a representative of $A_0+\Omega$ with the same norms, and $A+\Omega$ lies in the same set as $A$. Lemma~\ref{lem:invariances} follows.
\end{proof}

\begin{lemma}[Periodisation]\label{lem:periodisation}
Let $l_1,l_2\in\mathrm{Cl}$ with $l_2=J\cdot l_1$ for an odd integer $J$, let $T_i:=\mathbb{L}_k/l_i\mathbb{Z}^3$ for $i=1,2$, and let $\mathrm{pj}\colon T_2\to T_1$ be the covering map. To a bond function $A$ on $T_1$ assign the bond function $\mathrm{pj}^{*}A$ on $T_2$ given by $(\mathrm{pj}^{*}A)(x,x'):=A(\mathrm{pj} x,\mathrm{pj} x')$; in this way $\mathrm{pj}^{*}$ is a bijection from the bond functions on $T_1$ onto those bond functions on $T_2$ that are invariant under the translations by $l_1\mathbb{Z}^3$. Then $\mathrm{pj}^{*}$ maps each of the sets $R_k$, $c\cdot\tilde{R}_k$ with $0<c\le 1$, $\tilde{R}^{-}_k$, $G_k$, $\tilde{G}_k$, $\mathfrak{X}_k$ of $T_1$ into the set of the same name of $T_2$.
\end{lemma}

\begin{proof}
A bond function on $T_i$ is an $l_i\mathbb{Z}^3$-periodic bond function on $\mathbb{L}_k$, and $\mathrm{pj}^{*}A$ is the same function on $\mathbb{L}_k$ as $A$. The norms $|A|_{\infty}$ and $|\partial A|_{\infty}$ on $T_i$ are the suprema over $\mathbb{L}_k$ of the periodic functions $|A_{\nu}|$ and $|\partial_{\nu'}A_{\nu}|$ of one cell. The norm $|\delta_{\alpha}\partial A|_{\infty}$ is a supremum over pairs of cells of the same type at distance at most $1$. Since $1<l_1/2\le l_2/2$, for two such cells the distance on $T_1$, the distance on $T_2$ and the distance of suitable lifts in $\mathbb{L}_k$ coincide; so this norm, too, is a supremum over $\mathbb{L}_k$ that is the same for $i=1$ and for $i=2$. Hence $\mathrm{pj}^{*}$ preserves the three norms, and it maps $R_k$, $G_k$, $\mathfrak{X}_k$ and the sets $c\cdot R_k$, $(1-e_k^{\varepsilon})\cdot R_k$ of $T_1$ into the sets of the same name of $T_2$. If $e_k\ge 1$, then $1-e_k^{\varepsilon}\le 0$, so that the set $\tilde{R}^{-}_k$ of $T_1$ is empty, and for this set there is nothing to prove.

As $M$ divides $l_1$, the covering map $\mathrm{pj}$ maps every $M$-cube $\square_2$ of $T_2$ onto an $M$-cube $\square_1$ of $T_1$. The map $\mathrm{pj}$ sends the image in $T_2$ of a point of $\mathbb{L}_k$ to the image in $T_1$ of the same point; so a point $\tilde{c}\in\mathbb{L}_k$ that represents the centre of $\square_2$ represents also the centre of $\square_1$. Let $\tilde{W}$ be the set of the points $y\in\mathbb{L}_k$ with $|y-\tilde{c}|<\frac{1}{2}(2L+1)\cdot M$. Its sup-diameter is at most $(2L+1)\cdot M-\eta_k$, and
\[
(2L+1)\cdot M-\eta_k+1<(2L+1)\cdot M+1\le 13\cdot L\cdot M\le l_1\le l_2.
\]
As $l_1$ and $l_2$ belong to $\mathrm{Cl}$, the conventions at the head of this section apply on $T_2$ and on $T_1$: $\tilde{W}$ is a lift both of $\square_2^{\dagger}$ and of $\square_1^{\dagger}$, and the properties (i)--(iii) stated there hold on both tori. For a site $x$ of $\square_2^{\dagger}$ with the preimage $\tilde{x}\in\tilde{W}$, the site $\mathrm{pj} x$ is the image of $\tilde{x}$ in $T_1$. By (i), applied on both tori, the map $\mathrm{pj}$ is therefore injective on $\square_2^{\dagger}$, and its image is $\square_1^{\dagger}$. By (ii), this bijection and its inverse map bonds to bonds and cells of a type to cells of the same type, since on either torus these are the images of the bonds and of the cells of $\mathbb{L}_k$ that lie in $\tilde{W}$. By (iii), if $x$ and $x'$ are sites of $\square_2^{\dagger}$ and one of the two distances, of $x$ and $x'$ on $T_2$ and of $\mathrm{pj} x$ and $\mathrm{pj} x'$ on $T_1$, is at most $1$, then both distances are equal to $|\tilde{x}-\tilde{x}'|$. Thus $\mathrm{pj}$ preserves on $\square_2^{\dagger}$ all lattice relations: bonds, cells and their distances. With the side $11M$ in place of $(2L+1)\cdot M$, the same holds for $\square_2^{\sim 5}$ and its image $\square_1^{\sim 5}$. Let $A=A_0+A_1$ be split as in the definition of one of the sets $c\cdot\tilde{R}_k$, $\tilde{R}^{-}_k$, $\tilde{G}_k$ of $T_1$. Then $\mathrm{pj}^{*}A=\mathrm{pj}^{*}A_0+\mathrm{pj}^{*}A_1$, where $\mathrm{pj}^{*}A_0$ is real and $\mathrm{pj}^{*}A_1$ satisfies the same bounds on its norms as $A_1$. If $(A',\lambda)$ is a representative of $A_0$ on $\square_1^{\dagger}$, then $(A'\circ\mathrm{pj},\lambda\circ\mathrm{pj})$ is a representative of $\mathrm{pj}^{*}A_0$ on $\square_2^{\dagger}$ with the same three norms. The same holds for the cubes $\square_2^{\sim 5}$ and $\square_1^{\sim 5}$. Therefore $\mathrm{pj}^{*}A$ lies in the set of the same name of $T_2$. Lemma~\ref{lem:periodic-extension} follows.
\end{proof}

\begin{lemma}[Phases and differences along pairs]\label{lem:phases-differences-along}
Let $B\in\tilde{G}_k$, let $(x,x')\in\mathrm{Pz}_k$, and let $f$ be a spinor function. Then
\[
\left|\exp\left(\pm i\cdot e_k\cdot\eta_k\cdot(\tau B)(x,x')\right)\right|\le\exp\left(3M\cdot e_k\cdot|\operatorname{Im}B|_{\infty}\right)\le\exp(1).
\]
If moreover $|x-x'|\ge 1$, then $\left|(\delta_{\alpha,B}f)(x,x')\right|\le 4|f|_{\infty}$.
\end{lemma}

\begin{proof}
By Definition~\ref{def:pairs-sites-holder}, $\eta_k\cdot(\tau B)(x,x')$ is the mean, over the six rectilinear paths $\Gamma^{\pi}(x,x')$ of the shortest lift of $(x,x')$, of the sums $\sum_{b\in\Gamma^{\pi}(x,x')}\eta_k\cdot B(b)$. Each of these paths consists of $|x-x'|_1/\eta_k$ bonds, where $|x-x'|_1$ denotes the sum of the moduli of the three coordinate differences of this lift; moreover $|x-x'|_1\le 3|x-x'|\le 3M$, since $|x-x'|\le M$ for the pairs of $\mathrm{Pz}_k$. Hence
\[
\left|\operatorname{Im}\left(\eta_k\cdot(\tau B)(x,x')\right)\right|\le 3M\cdot|\operatorname{Im}B|_{\infty}.
\]
The modulus of $\exp\left(\pm i\cdot e_k\cdot\eta_k\cdot(\tau B)(x,x')\right)$ is $\exp\left(\mp e_k\cdot\operatorname{Im}\left(\eta_k\cdot(\tau B)(x,x')\right)\right)$, and so it is at most $\exp\left(3M\cdot e_k\cdot|\operatorname{Im}B|_{\infty}\right)$. Write $B=A_0+A_1$ as in the definition of $\tilde{G}_k$. By Lemma~\ref{lem:openness}, $\operatorname{Im}B=\operatorname{Im}A_1$, and $|A_1|_{\infty}<e_k^{-1+\varepsilon}$; therefore
\[
3M\cdot e_k\cdot|\operatorname{Im}B|_{\infty}<3M\cdot e_k^{\varepsilon}\le 1,
\]
by~\eqref{eq:renormalisation-group-31}. This proves the first assertion. Let now $|x-x'|\ge 1$. Then $|x-x'|^{-\alpha}\le 1$, and by the definition of $\delta_{\alpha,B}$ in Definition~\ref{def:pairs-sites-holder} and by the first assertion,
\[
\left|(\delta_{\alpha,B}f)(x,x')\right|\le\exp(1)\cdot|f(x')|+|f(x)|\le\left(\exp(1)+1\right)\cdot|f|_{\infty}\le 4|f|_{\infty}.
\]
Lemma~\ref{lem:phase-factors} follows.
\end{proof}

\subsection{Definition of the step map}\label{sec:definition-step-map}

The displays~\eqref{eq:renormalisation-group-56}, \eqref{eq:renormalisation-group-57}, \eqref{eq:renormalisation-group-59}, \eqref{eq:renormalisation-group-58}, \eqref{eq:renormalisation-group-60}, \eqref{eq:renormalisation-group-61} and~\eqref{eq:renormalisation-group-62} of this section are definitions. Together, these seven displays define, in the sense of Definition~\ref{def:renormalisation-step} (Section~\ref{sec:renormalisation-step}), the step map of level~$k$ associated with the system of primitives of Section~\ref{sec:hypotheses-conditions}; the two connections enter by the rule~\eqref{eq:renormalisation-group-4}, and the signs are those of Definition~\ref{def:contour-functional-functions}.

\subsubsection{Cubes and parameter cubes}\label{sec:cubes-parameter-cubes}

The facts~\eqref{it:renormalisation-group-g0}--\eqref{it:renormalisation-group-g4} below are consequences of the notation of Section~\ref{sec:notation}, of the definitions of Section~\ref{sec:lattice-geometry}, in particular of Definition~\ref{def:parameter-cubes-polydisks}, and of the following properties of $L$, $M$ and $l$: $L$ is odd, by Section~\ref{sec:hypotheses-conditions}; $M=L^m$; $M>4L$, by~\eqref{eq:renormalisation-group-30}; and $l\in\mathrm{Cl}$. As $l/(LM)$ is an integer not smaller than $13$, the side $l$ is a multiple of $LM$, hence of $M$, and $(2L+1)M<13LM\le l$.

\begin{enumerate}
\item[(g0)]\itemmark{g0}\label{it:renormalisation-group-g0} Both $L$ and $M=L^m$ are odd. The faces of the $M$-cubes contain no site of $\mathbb{L}_k$. Indeed, in each coordinate the faces of the $M$-cubes lie in $M\left(\frac{1}{2}+\mathbb{Z}\right)$, and the coordinates of the sites are the numbers $j\eta_k=jL^{-k}$ with $j\in\mathbb{Z}$; an equality $jL^{-k}=M\left(n+\frac{1}{2}\right)$ with $n\in\mathbb{Z}$ would give $2j=ML^k(2n+1)$, whose right side is odd. Consequently the $M$-cubes $\left\{x:|x-c|<\frac{M}{2}\right\}$, $c\in M\mathbb{Z}^3$, which are defined on the torus because $M$ divides $l$, form a partition of the set of the sites of $T_k$, and also of the set of the sites of $T^0_k$. Each block $B_k(y)$ and each $L$-block is contained in one $M$-cube, and each $M$-cube is contained in one $LM$-cube (the $LM$-cubes are defined on the torus, as $LM$ divides $l$). For the proof, note that $B_k(y)$ lies in the open cube of side $1$ around the unit site $y$, that an $L$-block lies in the open cube of side $L$ around its centre, a point of $L\mathbb{Z}^3$, and that an $M$-cube is the open cube of side $M$ around its centre, a point of $M\mathbb{Z}^3$. Since $L$ and $M$ are odd, in each coordinate the points of $M\left(\frac{1}{2}+\mathbb{Z}\right)$ have distance $\frac{1}{2}$ from $\mathbb{Z}$ and distance at least $\frac{L}{2}$ from $L\mathbb{Z}$, and the points of $LM\left(\frac{1}{2}+\mathbb{Z}\right)$ have distance at least $\frac{M}{2}$ from $M\mathbb{Z}$. So each of the open cubes of the sides $1$ and $L$ just described contains no point of a face of an $M$-cube and therefore lies in a single $M$-cube, and an $M$-cube contains no point of a face of an $LM$-cube and therefore lies in a single $LM$-cube.

\item[(g1)]\itemmark{g1}\label{it:renormalisation-group-g1} Let $c'$ be a cube of $\pi'_k$, resp.\ of $\pi''_k$, and let $c$ be its centre, a point of $M\mathbb{Z}^3$, resp.\ of $LM\mathbb{Z}^3$. Denote by $c'^{+}$ the union of all $M$-cubes with centre at sup-distance at most $M$, resp.\ at most $LM$, from $c$. These centres are the points $c+Mv$ with $v\in\{-1,0,1\}^3$, resp.\ with $v\in\{-L,\dots,L\}^3$, and they are distinct on the torus, because $(2L+1)M<l$. Hence $c'^{+}$ is the cube
\[
c'^{+}=\left\{x:|x-c|<\frac{3M}{2}\right\},\qquad\text{resp.}\qquad c'^{+}=\left\{x:|x-c|<LM+\frac{M}{2}\right\}.
\]
If a point has sup-distance at most $2L$ from some point of $c'$, then it belongs to $c'^{+}$, since its sup-distance from $c$ is at most $M+2L<\frac{3M}{2}$, resp.\ at most $LM+2L<LM+\frac{M}{2}$, by $M>4L$. An $M$-cube meets $c'$ if and only if it is one of the $M$-cubes of $c'^{+}$. Indeed, a lattice site of $c'$ lies in $c'^{+}$, so that by~\eqref{it:renormalisation-group-g0} the $M$-cube containing it is one of the $M$-cubes of $c'^{+}$; conversely, the centre of an $M$-cube of $c'^{+}$ lies in $M\mathbb{Z}^3$ and is therefore a lattice site; it belongs to that $M$-cube and, its sup-distance from $c$ being at most $M$, resp.\ at most $LM$, also to the closed cube $c'$. Consequently no parameter cube, of whatever kind, meets more than $N_c=(2L+1)^3$ of the $M$-cubes; for a cube of $\pi_k$, which is itself an $M$-cube, this holds because it meets no other $M$-cube, by~\eqref{it:renormalisation-group-g0}.

\item[(g2)]\itemmark{g2}\label{it:renormalisation-group-g2} Let $\mathcal{L}$ be a non-empty link-connected set whose elements are parameter cubes of a single kind, and write $\mathrm{pol}(\mathcal{L})$ for the union of all those $M$-cubes which meet at least one cube of $\mathcal{L}$. Then $\mathrm{pol}(\mathcal{L})$ is a polymer containing every cube of $\mathcal{L}$. For the kind $\pi_k$, the set $\mathrm{pol}(\mathcal{L})$ is the union of the cubes of $\mathcal{L}$ by~\eqref{it:renormalisation-group-g0}, and two cubes of $\pi_k$ are linked if they have a common face; so this union is connected through faces. For the kinds $\pi'_k$ and $\pi''_k$, the set $\mathrm{pol}(\mathcal{L})$ is the union of the cubes $c'^{+}$, $c'\in\mathcal{L}$, by~\eqref{it:renormalisation-group-g1}. Every $c'^{+}$ contains $c'$ and, being a cube built from $M$-cubes, is connected through faces. Two linked cubes $c'_1,c'_2$ of $\mathcal{L}$ share a lattice site, and the $M$-cube containing this site is an $M$-cube of $c'^{+}_1$ and of $c'^{+}_2$. As $\mathcal{L}$ is link-connected, the union of the cubes $c'^{+}$, $c'\in\mathcal{L}$, is connected through faces.

\item[(g3)]\itemmark{g3}\label{it:renormalisation-group-g3} Let a set of parameter cubes of one kind be given. For two different link-components of this set, the unions of their cubes have no lattice site in common: two cubes of $\pi'_k$, or of $\pi''_k$, with a common lattice site are linked, and two different cubes of $\pi_k$ have no common site, by~\eqref{it:renormalisation-group-g0}. Moreover, a union $\mathbf{Y}$ of $M$-cubes contains at most $|\mathbf{Y}|_M$ parameter cubes of each kind. Indeed, the centre of a parameter cube contained in $\mathbf{Y}$ is a lattice site of this cube, hence of $\mathbf{Y}$; by~\eqref{it:renormalisation-group-g0} the $M$-cube with the same centre, which contains this site, is one of the $M$-cubes of $\mathbf{Y}$; and different cubes of one kind have different centres.

\item[(g4)]\itemmark{g4}\label{it:renormalisation-group-g4} Let $z$ be a site of $T_k$. The set of the cubes of $\pi_k$ that meet $K_z$ is link-connected. Indeed, $K_z$ is a cube of side $4L<M$. Hence, for each of the three coordinates, the set of the values of this coordinate at the sites of $K_z$ meets at most two of the intervals $\left(u-\frac{M}{2},u+\frac{M}{2}\right)$, $u\in M\mathbb{Z}$, and two consecutive ones if it meets two. The cubes of $\pi_k$ that meet $K_z$ are therefore the $M$-cubes whose centres have each of their three coordinates in such a set of one or two consecutive multiples of $M$, and any two of them are joined by a sequence of these cubes in which any two consecutive cubes have a common face.
\end{enumerate}

\subsubsection{Data, variables and fields}\label{sec:data-variables-fields}

The datum is a pair $(m,\mathscr{E})\in\mathrm{Ball}^x_k(4)$ (Section~\ref{sec:spaces-activities}), with $\mathscr{E}=0$ if $k=0$. With the number $m(\mathscr{E})$ and the map $R$ of Section~\ref{sec:local-parts-adjustment} put
\[
m':=m+m(\mathscr{E}),\qquad \mathscr{E}^r:=R\,\mathscr{E} .
\]
The background consists of a connection $A\in\tilde{R}_k^{-}$ in the first slot and a connection $\mathscr{C}\in\mathfrak{X}_k$ in the second slot (Definition~\ref{def:domains}), and $\mathscr{T}:=A+\mathscr{C}$. The fluctuation variables are $\mathsf{Z}\in\mathrm{Box}_k$ (Section~\ref{sec:averaging-fluctuations-gauge}) and the generators $W$ of Section~\ref{sec:grassmann-algebras-norms}.

\paragraph{Parameters.} The parameters form a family $s=(s',s'')$. The family $s'=(s'_c)$ is indexed by the cubes $c$ of $\pi'_k$ and by the cubes $c$ of $\pi''_k$: its values at the cubes of $\pi'_k$ are the parameters of $H_k(s')$, the family~\eqref{it:primitives-p1}, and its values at the cubes of $\pi''_k$ are the parameters of $C_k^{\frac{1}{2}}(s')$, the family~\eqref{it:primitives-p2}. The family $s''=(s''_c)$ is indexed by the cubes $c$ of $\pi_k$, and the same variables $s''$ stand in $H_k(s'',\cdot)$ and $\bar{H}_k(s'',\cdot)$ (family~\eqref{it:primitives-p3}) and in $\Gamma_k(s'',\cdot)$ (family~\eqref{it:primitives-p4}). The parameters at cubes of different kinds are separate, independent variables, also when the cubes have the same centre; each parameter $s_c$ ranges over the disk $|s_c|\le M^{\alpha_0}$, so that $s$ ranges over the polydisk $\mathrm{Pd}_k$ of Definition~\ref{def:parameter-cubes-polydisks}. For a union $U$ of $M$-cubes let $1_U$ be the family of parameters with $(1_U)_c=1$ if the parameter cube $c$ is contained in $U$, and $(1_U)_c=0$ otherwise. When $s'$ is complex, $H_k(s')$, $C$ and $C_k^{\frac{1}{2}}(s')$ are regarded as operators between the complexifications of their spaces.

\paragraph{Fields.} Let $B\in\tilde{G}_k$, and let $\mathfrak{Z}$ be a complex bond function on $T_k$ such that $B+\mathfrak{Z}\in\tilde{G}_k$. Define
\begin{align}
\mathscr{Z}(s')&:=H_k(s')\cdot C\cdot C_k^{\frac{1}{2}}(s')\cdot\mathsf{Z},\notag\\
H_k^{\#}(s'',B)&:=\left(H_k(s'',B),\ \delta_{\alpha,B}H_k(s'',B)\right)\quad\text{on unbarred fields},\notag\\
H_k^{\#}(s'',B)&:=\left(\bar{H}_k(s'',B),\ \delta_{\alpha,-B}\bar{H}_k(s'',B)\right)\quad\text{on barred fields},\notag\\
\bar{H}_k(s'',B)&:=\mathsf{C}^{-1}\cdot H_k(s'',-B)\cdot\mathsf{C},\notag\\
J_k^{\#}(s'',B,\mathfrak{Z})&:=H_k^{\#}(s'',B+\mathfrak{Z})-H_k^{\#}(s'',B),\notag\\
\tilde{\Gamma}_k(s'',B)\,W&:=\Gamma_k(s'',B)\cdot W,\qquad \tilde{\Gamma}_k(s'',B)\,\bar{W}:=\bar{W},\notag\\
W_k^{\#}(s'',B)&:=H_k^{\#}(s'',B)\cdot\tilde{\Gamma}_k(s'',B)\cdot W.\label{eq:renormalisation-group-56}
\end{align}
The operator $\delta_{\alpha,B}$ is the H\"older derivative of Definition~\ref{def:pairs-sites-holder}; it is evaluated at every pair of $\mathrm{Pz}_k$. The operators $\bar{H}_k(s'',B)$ and $\bar{H}_k(s'',B+\mathfrak{Z})$ are defined, because $-B$ and $-(B+\mathfrak{Z})$ lie in $\tilde{G}_k$ (Lemma~\ref{lem:symmetries-domains}). Thus $\mathscr{Z}(s')$ is a complex bond function on $T_k$, and $H_k^{\#}(s'',B)$ maps a spinor field on $T^0_k$ to a pair formed by a spinor field on $T_k$ and a function on $\mathrm{Pz}_k$. Written out, the unbarred components of $W_k^{\#}(s'',B)$ are
\[
H_k(s'',B)\cdot\Gamma_k(s'',B)\cdot W\qquad\text{and}\qquad\delta_{\alpha,B}H_k(s'',B)\cdot\Gamma_k(s'',B)\cdot W,
\]
and its barred components are $\bar{H}_k(s'',B)\cdot\bar{W}$ and $\delta_{\alpha,-B}\bar{H}_k(s'',B)\cdot\bar{W}$. For $k=0$, $H_0(s')$ and $H_0(s'',B)$ denote the diagonal families of property~\ref{it:renormalisation-group-17}.

\subsubsection{The localisation map}\label{sec:localisation-map}

Throughout this subsection a non-empty subset of the three kinds $\pi_k$, $\pi'_k$, $\pi''_k$ is fixed; a parameter cube is a cube of one of these kinds, and the parameters $s=(s_c)$ are indexed by the parameter cubes. For a polymer $X$ let $P_0(X)$ denote the set of those parameter cubes which are contained in $X$. If $P$ is a set of parameter cubes with $P\cap P_0(X)=\emptyset$, then $Y^{\mathrm{p}}(X,P)$ denotes the union of all cubes in $P_0(X)\cup P$, and $Y'(X,P)$ the union of all $M$-cubes having a lattice site in common with $X\cup Y^{\mathrm{p}}(X,P)$. Let $G(s,X)$ be a function of the parameters $s$ and of a polymer $X$. For $\mathbf{Y}\in D_k$ put
\begin{equation}\label{eq:renormalisation-group-57}
\mathrm{Loc}[G](\mathbf{Y}):=\sum_{(X,P):\,Y'(X,P)=\mathbf{Y}}\ \int_{[0,1]^P}\mathrm{d}s_P\,\frac{\partial}{\partial s_P}\left[G(s,X)\right]_{s=1\text{ on }P_0(X),\ s=0\text{ outside }P_0(X)\cup P}.
\end{equation}
The sum is taken over the pairs $(X,P)$ formed by a polymer $X$ and a set $P$ of parameter cubes with $P\cap P_0(X)=\emptyset$ and $Y'(X,P)=\mathbf{Y}$. Here $\mathrm{d}s_P:=\prod_{c\in P}\mathrm{d}s_c$ and $\frac{\partial}{\partial s_P}:=\prod_{c\in P}\frac{\partial}{\partial s_c}$; the variables $s_c$, $c\in P$, are integrated over $[0,1]$, and the subscript of the bracket fixes the remaining parameters: $s_c=1$ for $c\in P_0(X)$, and $s_c=0$ for every parameter cube $c$ outside $P_0(X)\cup P$. For $P=\emptyset$ the term of $(X,\emptyset)$ is the value of $G(\cdot,X)$ at the family that equals $1$ on $P_0(X)$ and $0$ at all other parameter cubes.

\subsubsection{Localisation of the exponent}\label{sec:localisation-exponent}

For a $4\times4$ matrix $O$ and sites $x$, $x'$ put
\begin{align*}
\bar{\psi}(x)\,O\,\psi(x'):=\sum_{a,b}\bar{\psi}_a(x)\cdot O_{ab}\cdot\psi_b(x').
\end{align*}
A function $G(\cdot,\square)$ that is defined for the single $M$-cubes $\square$ is extended to all polymers by the value $0$ on the polymers that are not single $M$-cubes.

Let $k\ge1$, let $\mathbf{Y}$ be a polymer and $B\in\tilde{G}_k$. For a spinor-valued function $\psi$ on the sites of $T_k$, a spinor-valued function $\chi$ on $\mathrm{Pz}_k$ and a unit site $y$ of $\mathbf{Y}$ put
\begin{align}
\left(T^{(1)}(B)\psi\right)(y):={}&\eta_k^3\sum_{x\in B_k(y)}\phi_{B,y}(x)\cdot\left\{\left(1+\frac{\mathfrak{m}}{b_k}\right)\psi(x)\right.\notag\\
&\left.{}-\left(b_k\cdot\eta_k\right)^{-1}\sum_{(\nu,\boldsymbol{\sigma}):\,x^{\boldsymbol{\sigma}}\in\mathbf{Y}}P^{\boldsymbol{\sigma}}_{\nu}\cdot\left[u_B\left(x,x^{\boldsymbol{\sigma}}\right)\cdot\psi\left(x^{\boldsymbol{\sigma}}\right)-\psi(x)\right]\right\},\notag\\
\left(T^{(2)}(B)\chi\right)(y):={}&-b_k^{-1}\cdot\eta_k^{2+\alpha}\sum_{x\in B_k(y)}\phi_{B,y}(x)\sum_{(\nu,\boldsymbol{\sigma}):\,x^{\boldsymbol{\sigma}}\notin\mathbf{Y}}P^{\boldsymbol{\sigma}}_{\nu}\cdot\chi\left(x,x^{\boldsymbol{\sigma}}\right),\notag\\
\bar{T}^{(i)}(B):={}&\mathsf{C}^{-1}\cdot T^{(i)}(-B)\cdot\mathsf{C}\qquad(i=1,2).\label{eq:renormalisation-group-58}
\end{align}
Let $F$ be an element of the Grassmann algebra $\mathrm{Gr}[\psi,\chi,\Psi]$ (Section~\ref{sec:grassmann-algebras-norms}), or a function of $(A,\mathscr{C},\mathsf{Z})\in\tilde{R}_k^{-}\times\mathfrak{X}_k\times\mathrm{Box}_k$ (Section~\ref{sec:data-variables-fields}) with values in $\mathrm{Gr}[\psi,\chi,\Psi]$, such as the function $E^{(i),*}(\mathbf{Y})$, $i=1,2,3$, of~\eqref{eq:renormalisation-group-59}; the substitution~\eqref{eq:renormalisation-group-58} in $F$ is the one formed with the polymer $\mathbf{Y}$ fixed above. The substitution~\eqref{eq:renormalisation-group-58} turns $F$ into the element obtained by inserting, at every unit site $y$ of $\mathbf{Y}$, the element $\left(T^{(1)}(\mathscr{T})\psi\right)(y)+\left(T^{(2)}(\mathscr{T})\chi\right)(y)$ in place of the generators $\Psi(y)$, and the element $\left(\bar{T}^{(1)}(\mathscr{T})\bar{\psi}\right)(y)+\left(\bar{T}^{(2)}(\mathscr{T})\bar{\chi}\right)(y)$ in place of the generators $\bar{\Psi}(y)$; the four operators are formed with the polymer $\mathbf{Y}$ and are applied to the generators $\psi$, $\chi$, $\bar{\psi}$, $\bar{\chi}$. Thus the substitution is the homomorphism of algebras that maps $\Psi(y)$ and $\bar{\Psi}(y)$, for every unit site $y$ of $\mathbf{Y}$, to these two elements and fixes all other generators, namely $\psi$, $\chi$, $\bar{\psi}$, $\bar{\chi}$, and $\Psi(y)$, $\bar{\Psi}(y)$ for the unit sites $y$ outside $\mathbf{Y}$; there is exactly one such homomorphism, since the images of all generators are linear combinations of generators, and such linear combinations anticommute with one another and have square $0$. The operators are taken at $\mathscr{T}=A+\mathscr{C}$; if $F$ is a function, the substitution is made at every point $(A,\mathscr{C},\mathsf{Z})$ of its domain, with $\mathscr{T}$ formed from the $A$ and the $\mathscr{C}$ of that point. If $F$ contains the generators $\Psi(y)$, $\bar{\Psi}(y)$ only at unit sites $y$ of $\mathbf{Y}$, then every generator occurring in $F$ is mapped into $\mathrm{Gr}[\psi,\chi]$, and the image of $F$ lies in $\mathrm{Gr}[\psi,\chi]$, respectively takes its values in $\mathrm{Gr}[\psi,\chi]$. By Lemma~\ref{lem:renormalisation-group-1}(a) this is the case for $F=E^{(i),*}(\mathbf{Y})$, $i=1,2,3$: every term $(X,P)$ of~\eqref{eq:renormalisation-group-57} has $Y'(X,P)=\mathbf{Y}$, which is connected, so that $Y'_1=\mathbf{Y}$ in that lemma, and the derivatives and integrals in~\eqref{eq:renormalisation-group-57} keep this property.

Let $k\ge1$. Let $J$ and $\bar{J}$ denote the components without H\"older derivative of $J^{\#}_k\left(s'',\mathscr{T},\mathscr{Z}(s')\right)$ on the unbarred and on the barred fields, and let $\mathcal{W}$ and $\bar{\mathcal{W}}$ denote those of $W^{\#}_k(s'',\mathscr{T})$. Here $J^{\#}_k$ is the operator and $W^{\#}_k$ the field defined in~\eqref{eq:renormalisation-group-56}; thus $J$ and $\bar{J}$ are operators, whereas $\mathcal{W}$ and $\bar{\mathcal{W}}$ are fields, linear in the generators $W$ and in the generators $\bar{W}$, respectively. The localisation~\eqref{eq:renormalisation-group-57} is formed with the three kinds $\pi_k$, $\pi'_k$, $\pi''_k$ of parameter cubes in the definition of $E^{(i),*}$, and with the kind $\pi_k$ alone in the definition of $E^{(\le5),x}$. The contour integrals $\oint_{\rho}$ of Section~\ref{sec:notation} are taken in the variable $t$. Put
\begin{align}
G^{(1)}(s,X)&:=\oint_{\rho_1}\left\langle\bar{\Psi},M_k\left(X,\mathscr{T}+t\cdot\mathscr{Z}(s')\right)\Psi\right\rangle,\notag\\
G^{(2)}(s,X)&:=\oint_{\rho_2}\mathscr{E}^r\left(X;A+t\cdot\mathscr{Z}(s'),\mathscr{C},\psi^{\#}+t\cdot J^{\#}_k\left(s'',\mathscr{T},\mathscr{Z}(s')\right)\Psi\right),\notag\\
G^{(3)}(s,\square)&:=\oint_{\rho_2}\left(-m'\right)\left\langle\bar{\psi}+t\cdot\bar{J}\,\bar{\Psi},\ \psi+t\cdot J\,\Psi\right\rangle_{\square},\notag\\
E^{(i),*}&:=\mathrm{Loc}\left[G^{(i)}\right]\qquad(i=1,2,3),\notag\\
E^{(i),x}(\mathbf{Y})&:=\text{the image of }E^{(i),*}(\mathbf{Y})\text{ under the substitution~\eqref{eq:renormalisation-group-58}}\qquad(i=1,2,3),\notag\\
E^{(\le3),x}&:=E^{(1),x}+E^{(2),x}+E^{(3),x},\notag\\
G^{(\le3)}(s'',X)&:=E^{(\le3),x}\left(X;A,\mathscr{C},\mathsf{Z},\psi^{\#}+W^{\#}_k(s'',\mathscr{T})\right),\notag\\
G^{(4)}(s'',X)&:=\oint_{\rho_4}\mathscr{E}^r\left(X;A,\mathscr{C},\psi^{\#}+t\cdot W^{\#}_k(s'',\mathscr{T})\right),\notag\\
G^{(5)}(s'',\square)&:=\oint_{\rho_4}\left(-m'\right)\left\langle\bar{\psi}+t\cdot\bar{\mathcal{W}},\ \psi+t\cdot\mathcal{W}\right\rangle_{\square},\notag\\
E^{(\le5),x}&:=\mathrm{Loc}\left[G^{(\le3)}+G^{(4)}+G^{(5)}\right].\label{eq:renormalisation-group-59}
\end{align}

Let $k=0$. Then the datum is $(m,0)$, hence $m'=m$ and $\mathscr{E}^r=0$; the field $\Psi$ does not occur, and the substitution~\eqref{eq:renormalisation-group-58} is not made. In this case $G^{(2)}$, $G^{(3)}$ and $G^{(4)}$ are set equal to $0$, $E^{(1),x}:=\mathrm{Loc}\left[G^{(1)}\right]$, and $G^{(1)}$ is the following function of the single $M$-cubes $\square$:
\begin{multline}
G^{(1)}(s,\square):=\frac{1}{2}\sum_{x\in\square}\sum_{\nu,\boldsymbol{\sigma}}\left\{\bar{\psi}(x)\,P^{\boldsymbol{\sigma}}_{\nu}\left(\exp\left(i\,e_k\cdot\mathscr{Z}(s')\left(x,x^{\boldsymbol{\sigma}}\right)\right)-1\right)\left(\psi(x)+\chi\left(x,x^{\boldsymbol{\sigma}}\right)\right)\right.\\
\left.{}+\left(\bar{\psi}(x)+\bar{\chi}\left(x,x^{\boldsymbol{\sigma}}\right)\right)P^{-\boldsymbol{\sigma}}_{\nu}\left(\exp\left(-i\,e_k\cdot\mathscr{Z}(s')\left(x,x^{\boldsymbol{\sigma}}\right)\right)-1\right)\psi(x)\right\}.\label{eq:renormalisation-group-60}
\end{multline}
All other definitions in~\eqref{eq:renormalisation-group-59} are kept at level $0$, with the following conventions: $\mathcal{W}$ and $\bar{\mathcal{W}}$ are the components without H\"older derivative of $W^{\#}_0(s'',\mathscr{T})$ on the unbarred and on the barred fields; the localisation~\eqref{eq:renormalisation-group-57} is formed with the three kinds $\pi_0$, $\pi'_0$, $\pi''_0$ of parameter cubes in the definitions of $E^{(1),x}$ and of $E^{(i),*}$, and with the kind $\pi_0$ alone in the definition of $E^{(\le5),x}$; and the contour integrals are taken in the variable $t$.

\begin{lemma}[Interior and boundary parts of $T_k$]\label{lem:interior-boundary-parts}
Let $k\ge1$, let $\mathbf{Y}$ be a polymer, $y$ a unit site of $\mathbf{Y}$ and $B\in\tilde{G}_k$.
\begin{enumerate}
\item[(a)] If $\chi=\delta_{\alpha,B}\psi$, then $\left(T^{(1)}(B)\psi\right)(y)+\left(T^{(2)}(B)\chi\right)(y)=\left(T_k(B)\psi\right)(y)$. If $\bar{\chi}=\delta_{\alpha,-B}\bar{\psi}$, then $\left(\bar{T}^{(1)}(B)\bar{\psi}\right)(y)+\left(\bar{T}^{(2)}(B)\bar{\chi}\right)(y)=\left(\bar{T}_k(B)\bar{\psi}\right)(y)$.
\item[(b)] In the row of $T^{(1)}(B)$ at $y$ the field $\psi$ occurs only at sites of $\mathbf{Y}$. In the row of $T^{(2)}(B)$ at $y$ the field $\chi$ occurs only at pairs $\left(x,x^{\boldsymbol{\sigma}}\right)$ with $x\in\mathbf{Y}$ and $x^{\boldsymbol{\sigma}}\notin\mathbf{Y}$; all of them are of kind~(II) for $\mathbf{Y}$. Both rows depend on $B$ through its values on the bonds with an end in $B_k(y)$ only, and $B_k(y)\subset\mathbf{Y}$. Every entry of the two rows is an entire function of $B$.
\end{enumerate}
\end{lemma}

\begin{proof}
(a) By~\eqref{eq:renormalisation-group-39},
\begin{align*}
\left(T_k(B)\psi\right)(y)={}&\eta_k^3\sum_{x\in B_k(y)}\phi_{B,y}(x)\cdot\left\{\left(1+\frac{\mathfrak{m}}{b_k}\right)\psi(x)\right.\\
&\left.{}-\left(b_k\cdot\eta_k\right)^{-1}\sum_{\nu,\boldsymbol{\sigma}}P^{\boldsymbol{\sigma}}_{\nu}\cdot\left[u_B\left(x,x^{\boldsymbol{\sigma}}\right)\cdot\psi\left(x^{\boldsymbol{\sigma}}\right)-\psi(x)\right]\right\}.
\end{align*}
The right side coincides with the formula for $\left(T^{(1)}(B)\psi\right)(y)$ in~\eqref{eq:renormalisation-group-58}, except that its inner sum extends over all pairs $(\nu,\boldsymbol{\sigma})$ and not only over those with $x^{\boldsymbol{\sigma}}\in\mathbf{Y}$. Hence $\left(T_k(B)\psi\right)(y)-\left(T^{(1)}(B)\psi\right)(y)$ is the sum over $x\in B_k(y)$ of $-\eta_k^3\cdot\left(b_k\cdot\eta_k\right)^{-1}\cdot\phi_{B,y}(x)$ times the terms of the inner sum with $x^{\boldsymbol{\sigma}}\notin\mathbf{Y}$. For every pair $(\nu,\boldsymbol{\sigma})$ the sites $x$ and $x^{\boldsymbol{\sigma}}$ are nearest neighbours; each of the six rectilinear paths from $x$ to $x^{\boldsymbol{\sigma}}$ consists of the single bond $\left(x,x^{\boldsymbol{\sigma}}\right)$, and $\left|x-x^{\boldsymbol{\sigma}}\right|=\eta_k$. Hence, by Definition~\ref{def:pairs-sites-holder},
\begin{align*}
\chi\left(x,x^{\boldsymbol{\sigma}}\right)=\eta_k^{-\alpha}\cdot\left(u_B\left(x,x^{\boldsymbol{\sigma}}\right)\cdot\psi\left(x^{\boldsymbol{\sigma}}\right)-\psi(x)\right).
\end{align*}
Since $-\eta_k^3\cdot\left(b_k\cdot\eta_k\right)^{-1}\cdot\eta_k^{\alpha}=-b_k^{-1}\cdot\eta_k^{2+\alpha}$, the difference is $\left(T^{(2)}(B)\chi\right)(y)$. This proves the first identity.

The matrix $\mathsf{C}$ acts on the spinor index only; therefore $\mathsf{C}\,\bar{\chi}=\delta_{\alpha,-B}\left(\mathsf{C}\,\bar{\psi}\right)$. By the definition of $\bar{T}^{(i)}$ in~\eqref{eq:renormalisation-group-58} and that of $\bar{T}_k$ in Section~\ref{sec:exact-fermion-operators}, the left side of the second identity is
\begin{align*}
\mathsf{C}^{-1}\cdot\left[T^{(1)}(-B)\cdot\mathsf{C}\,\bar{\psi}+T^{(2)}(-B)\cdot\delta_{\alpha,-B}\left(\mathsf{C}\,\bar{\psi}\right)\right](y)=\left(\mathsf{C}^{-1}\cdot T_k(-B)\cdot\mathsf{C}\,\bar{\psi}\right)(y)=\left(\bar{T}_k(B)\bar{\psi}\right)(y);
\end{align*}
the first equality is the computation just made, with $-B$ in place of $B$ and $\mathsf{C}\,\bar{\psi}$ in place of $\psi$.

(b) By~\eqref{it:renormalisation-group-g0}, $B_k(y)$ lies in one $M$-cube. This cube contains the site $y$ of $\mathbf{Y}$; hence it is one of the $M$-cubes of $\mathbf{Y}$, and $B_k(y)\subset\mathbf{Y}$. Therefore every site $x$ over which the sums in~\eqref{eq:renormalisation-group-58} extend lies in $\mathbf{Y}$. In the row of $T^{(1)}(B)$ the sites $x^{\boldsymbol{\sigma}}$ lie in $\mathbf{Y}$ by the condition on the inner sum. In the row of $T^{(2)}(B)$ each pair $\left(x,x^{\boldsymbol{\sigma}}\right)$ consists of nearest neighbours, with $x\in\mathbf{Y}$ and $x^{\boldsymbol{\sigma}}\notin\mathbf{Y}$; by Definition~\ref{def:pairs-sites-holder} such a pair is of kind~(II) for $\mathbf{Y}$. Every entry of the two rows is a polynomial, with constant coefficients, in finitely many of the numbers $\phi_{B,y}(x)$ and $u_B\left(x,x^{\boldsymbol{\sigma}}\right)$, and each of these numbers is the exponential of a linear form in $B$; hence the entries are entire functions of $B$. Let $x\in B_k(y)$. The number $u_B\left(x,x^{\boldsymbol{\sigma}}\right)$ depends on $B$ through its value on the bond $\left(x,x^{\boldsymbol{\sigma}}\right)$ only, and this bond has the end $x$ in $B_k(y)$. By the definition in Section~\ref{sec:notation}, the number $\phi_{B,y}(x)$ depends on $B$ only through its values on the bonds of the rectilinear paths from $x_{i+1}$ to $x_i$, $0\le i<k$, where $x_i$ is the centre of the block of side $L^i\cdot\eta_k$ that contains $x$. The blocks of the successive sides are nested (Section~\ref{sec:tori-nested-blocks}); hence $x_i$ and $x_{i+1}$ lie in the block of side $L^{i+1}\cdot\eta_k$ that contains $x$, and this block is contained in $B_k(y)$, the block of side $L^k\cdot\eta_k=1$ that contains $x$. The sites of a rectilinear path lie in the smallest box containing its two end points, and a block containing both end points contains that box; therefore every bond of these paths has both ends in $B_k(y)$. Consequently both rows depend on $B$ through its values on the bonds with an end in $B_k(y)$ only.
\end{proof}

\subsubsection{The fluctuation integral and the new activity}\label{sec:fluctuation-integral-new}

Let $\rho^T$ be as in Proposition~\ref{prop:polymers}. For a union $X$ of $M$-cubes the elements $\mathsf{K}(X)$ and $\mathsf{H}(X)$ are defined by the first two lines of the following display, where $\{X_i\}$ ranges over the finite sets of distinct polymers with union $X$, and $\mathsf{K}(\emptyset):=1$. For $X\in D_k$ the element $E^{\#,x}_k(X)$ is defined by its last line, where $(X_1,\dots,X_n)$ ranges over the $n$-tuples of polymers with union $X$:
\begin{align}
\mathsf{K}(X)&:=\sum_{\{X_i\}:\,\bigcup_iX_i=X}\ \prod_i\left(\exp\left(E^{(\le5),x}(X_i)\right)-1\right),\notag\\
\mathsf{H}(X)&:=\int\mathrm{d}\mu_I(W)\,\mathrm{d}\hat{\mu}(\mathsf{Z})\,\mathsf{K}(X),\notag\\
E^{\#,x}_k(X)=E^{\#,x}_k(m,\mathscr{E})(X;A,\mathscr{C},\psi,\chi)&:=\sum_{n\ge1}\frac{1}{n!}\sum_{(X_1,\dots,X_n):\,\bigcup_iX_i=X}\rho^T(X_1,\dots,X_n)\prod_i\mathsf{H}(X_i).\label{eq:renormalisation-group-61}
\end{align}
For $Y\in D_{k+1}$ and $(\mathscr{A},\mathscr{C})\in\tilde{R}_{k+1}\times\mathfrak{X}_{k+1}$ put
\begin{align}
E^{*,x}_k(m,\mathscr{E})&:=L\,E^{\#,x}_k(m,\mathscr{E}),\quad\text{that is,}\notag\\
E^{*,x}_k(m,\mathscr{E})(Y;\mathscr{A},\mathscr{C},\psi,\chi)&:=\sum_{X:\,\bar{X}^L=LY}E^{\#,x}_k(m,\mathscr{E})\left(X;\mathscr{A}_L,\mathscr{C}_L,\psi_L,\chi_L\right),\notag\\
\mathscr{E}^{+}(Y;\mathscr{A},\mathscr{C},\psi,\chi)&:=\sum_{X:\,\bar{X}^L=LY}\left[\left(R\,\mathscr{E}\right)^{\mathrm{nat}}+\delta E^{\mathrm{det}}_k+E^{\#,x}_k(m,\mathscr{E})\right]\left(X;\mathscr{A}_L,\mathscr{C}_L,\psi_L,\chi_L\right)\notag\\
&=\left(L\,R\,\mathscr{E}+L\,\delta E^{\mathrm{det}}_k+E^{*,x}_k(m,\mathscr{E})\right)(Y;\mathscr{A},\mathscr{C},\psi,\chi).\label{eq:renormalisation-group-62}
\end{align}

The functions $E^d_k$,
\begin{align*}
E^{\mathrm{det}}_k(X,B)=E^d_k(X,B)-E^d_k(X,0)\qquad\text{and}\qquad\delta E^{\mathrm{det}}_k(X;\mathscr{A},\mathscr{C})=E^d_k(X,\mathscr{A}+\mathscr{C})-E^d_k(X,\mathscr{C})
\end{align*}
are those of Definition~\ref{def:determinant-functions}. That definition expresses them through the functions $\acute{E}_k$ of convention~\eqref{eq:renormalisation-group-7} in Definition~\ref{def:contour-functional-functions}, and so through the family~\eqref{it:primitives-p5}; both for $k=0$ and for $k\ge1$ the functions $\acute{E}_k$ are also written out in Section~\ref{sec:fermion-determinant}, ahead of Proposition~\ref{prop:polymer-expansion-determinant}. The operators $L$ of reblocking and $R$ of extraction, the adjustment $(\cdot)^{\mathrm{nat}}$ and the scalings $A_L$, $\psi_L$, $\chi_L$ are defined in Section~\ref{sec:local-parts-adjustment}. For $(x,x')\in\mathrm{Pz}_k$ the points $x/L$ and $x'/L$ are distinct sites of $T_{k+1}$ whose distance does not exceed $M/L\le M$; thus $(x/L,x'/L)\in\mathrm{Pz}_{k+1}$, and $\chi_L(x,x')=L^{-1-\alpha}\chi(x/L,x'/L)$ is defined on all of $\mathrm{Pz}_k$.

\begin{remark}\label{rem:renormalisation-group-1}
\begin{enumerate}
\item[(i)] In the displays~\eqref{eq:renormalisation-group-56}--\eqref{eq:renormalisation-group-61} the connections enter in two ways: through the fermionic operators, exact or weakened, whose connection argument is $\mathscr{T}$ or $\mathscr{T}+t\cdot\mathscr{Z}(s')$; and through $\mathscr{E}^r$, which is taken with $A$ or $A+t\cdot\mathscr{Z}(s')$ in the first slot and with $\mathscr{C}$ in the second slot. The photon operators $H_k(s')$, $C$, $C_k^{\frac{1}{2}}(s')$, the measures, the contours and $\mathrm{Loc}$ do not depend on the connections.
\item[(ii)] The sign of $G^{(1)}$ is plus, as convention~\eqref{eq:renormalisation-group-5} of Definition~\ref{def:contour-functional-functions} requires; $G^{(3)}$ and $G^{(5)}$ are built from the form $-m'\left\langle\bar{\psi},\psi\right\rangle_{\square}$, as convention~\eqref{eq:renormalisation-group-6} requires.
\item[(iii)] Neither $A$ nor $\mathscr{C}$ occurs in~\eqref{eq:renormalisation-group-60}. The link phases of $\mathscr{T}$ enter through the generators $\chi$ and $\bar{\chi}$: in the proof of clause~\eqref{it:step-small-fields-c} of Theorem~\ref{thm:step-small-fields} in Section~\ref{sec:proof-clause-theorem} these generators are evaluated at $\delta_{\alpha,\mathscr{T}}\psi$ and $\delta_{\alpha,-\mathscr{T}}\bar{\psi}$.
\end{enumerate}
\end{remark}

\subsection{Estimates for the operators of the step}\label{sec:estimates-operators-step}

In this section the three norms of a complex bond function $B$ on $T_k$ are the norms $|B|_{\infty}$, $|\partial B|_{\infty}$ and $|\delta_{\alpha}\partial B|_{\infty}$ of Section~\ref{sec:domains-bond-functions}.

\begin{lemma}[Bounds and analyticity]\label{lem:bounds-analyticity}
Let $s=(s',s'')\in\mathrm{Pd}_k$, $\mathsf{Z}\in\mathrm{Box}_k$, $A\in\tilde{R}_k^{-}$ and $\mathscr{C}\in\mathfrak{X}_k$, and put $\mathscr{T}:=A+\mathscr{C}$. The field $\mathscr{Z}(s')$ and the operators $H_k^{\#}$, $J_k^{\#}$ and $\tilde{\Gamma}_k$ are those of~\eqref{eq:renormalisation-group-56}.
\begin{enumerate}
\item[(q1)]\itemmark{q1}\label{it:bounds-analyticity-q1} Each of the three norms of $\mathscr{Z}(s')$ is at most $\frac{1}{2}e_k^{-\frac{\varepsilon}{2}}$. Let $\mathfrak{Z}$ be a complex bond function on $T_k$ each of whose three norms is at most $\frac{1}{2}e_k^{-\frac{\varepsilon}{2}}$, and let $t\in\mathbb{C}$ with $|t|\le\rho_1$. Then
\[
t\cdot\mathfrak{Z}\in e_k^{\varepsilon}R_k,\qquad A+t\cdot\mathfrak{Z}\in\tilde{R}_k,\qquad \mathscr{T}+t\cdot\mathfrak{Z}\in\tilde{R}_k+\mathfrak{X}_k,
\]
and $\tilde{R}_k+\mathfrak{X}_k\subseteq\tilde{G}_k$.
\item[(q2)]\itemmark{q2}\label{it:bounds-analyticity-q2} Let $k\ge1$, $B\in\tilde{G}_k$, $X\in D_k$ and $\mathsf{w}_3>0$. Then
\[
\left\|\langle\bar{\Psi},M_k(X,B)\Psi\rangle\right\|_{\mathsf{w}_3}\le4c_X\cdot c_{\gamma}\cdot C_w\cdot M^3\cdot\mathsf{w}_3^2\cdot\exp\left(-(\kappa_w-1)\cdot d_M(X)\right),
\]
where $\|\cdot\|_{\mathsf{w}_3}$ is the norm of Proposition~\ref{prop:weighted-coefficient-norms} on $\mathrm{Gr}[\Psi]$, formed with the weight $\mathsf{w}_3$ of the generators $\Psi(u)$.
\item[(q3)]\itemmark{q3}\label{it:bounds-analyticity-q3} Let $B\in\tilde{G}_k$. On the unbarred and on the barred fields,
\[
\left|H_k^{\#}(s'',B)\right|_{1,\infty}\le4C_w,
\]
where in this norm the supremum is taken over the sites and over all pairs of $\mathrm{Pz}_k$. For $\mathfrak{Z}$ as in~\eqref{it:bounds-analyticity-q1},
\[
\left|J_k^{\#}(s'',\mathscr{T},\mathfrak{Z})\right|_{1,\infty}\le16C_w\cdot e_k^{\frac{3}{4}-\frac{9\varepsilon}{2}}.
\]
\item[(q4)]\itemmark{q4}\label{it:bounds-analyticity-q4} One has
\[
\left|H_k^{\#}(s'',\mathscr{T})\cdot\tilde{\Gamma}_k(s'',\mathscr{T})\right|_{1,\infty}\le C_W=4C_w^2\cdot c_{\gamma}.
\]
\item[(q5)]\itemmark{q5}\label{it:bounds-analyticity-q5} The field $\mathscr{Z}(s')$ is linear in $\mathsf{Z}$, continuous in $s'$ on $\mathrm{Pd}_k$, and analytic in $s'$ in the interior of $\mathrm{Pd}_k$. Let $X\in D_k$. The entries of $M_k(X,B)$, of $H_k^{\#}(s'',B)$ and of $\Gamma_k(s'',B)$ at
\[
B=\mathscr{T}+t\cdot\mathscr{Z}(s')
\]
are continuous functions of $(s,t,A,\mathscr{C},\mathsf{Z})$ on the set
\[
\mathrm{Pd}_k\times\left\{t\in\mathbb{C}:|t|\le\rho_1\right\}\times\tilde{R}_k^{-}\times\mathfrak{X}_k\times\mathrm{Box}_k;
\]
they are analytic in $(s,t,A,\mathscr{C})$ at the points at which $s$ is an interior point of $\mathrm{Pd}_k$ and $|t|<\rho_1$; and for fixed $s\in\mathrm{Pd}_k$ and $\mathsf{Z}\in\mathrm{Box}_k$ they are analytic in $(t,A,\mathscr{C})$ on a neighbourhood of $\left\{t\in\mathbb{C}:|t|\le\rho_1\right\}\times\tilde{R}_k^{-}\times\mathfrak{X}_k$.
\end{enumerate}
\end{lemma}

\begin{proof}
\emph{Part~\eqref{it:bounds-analyticity-q1}.} Since $\mathsf{Z}\in\mathrm{Box}_k$ (Section~\ref{sec:averaging-fluctuations-gauge}), condition~\eqref{eq:renormalisation-group-31} gives
\[
\left|\mathsf{Z}(\Upsilon)\right|\le p_{0,k}\le e_k^{-\frac{\varepsilon}{4}}\qquad\text{for every }\Upsilon .
\]
By property~\ref{it:renormalisation-group-4} and the definition of $c_{\gamma}$ in Section~\ref{sec:notation}, for every $\Upsilon$,
\[
\left|\left(C_k^{\frac{1}{2}}(s')\,\mathsf{Z}\right)(\Upsilon)\right|\le\sum_{\Upsilon'}C_w\cdot\exp\left(-\gamma_w\cdot d(\Upsilon,\Upsilon')\right)\cdot\left|\mathsf{Z}(\Upsilon')\right|\le C_w\cdot c_{\gamma}\cdot e_k^{-\frac{\varepsilon}{4}}.
\]
For complex $s'$ the coordinates of $C_k^{\frac{1}{2}}(s')\,\mathsf{Z}$ are complex numbers. By Proposition~\ref{prop:basis-map-box}, in the form for vectors with complex coordinates recalled in Section~\ref{sec:recalled-auxiliary-results}, the map $C$ multiplies the largest modulus of a coordinate by at most $c_C$; hence
\[
\left|C\,C_k^{\frac{1}{2}}(s')\,\mathsf{Z}\right|_{\infty}\le c_C\cdot C_w\cdot c_{\gamma}\cdot e_k^{-\frac{\varepsilon}{4}}.
\]
By property~\ref{it:renormalisation-group-3}, the moduli of the kernels $H_k(s')(b,b')$, $\left(\partial H_k(s')\right)(b,b')$ and $\left(\delta_{\alpha}\partial H_k(s')\right)((b,\tilde{b});b')$ are at most $C_w\cdot\exp\left(-\gamma_w\cdot d(b,b')\right)$, and by the definition of $c_{\gamma}$ the sum of this bound over the unit bonds $b'$ is at most $C_w\cdot c_{\gamma}$. The same holds for $k=0$: by property~\ref{it:renormalisation-group-17}, for the parameters $\mathrm{par}=(\alpha_0,\gamma_w,\kappa_w;C_w)$ of the system, the family $H_0(s')$ has property~\ref{it:renormalisation-group-3} with the constant $4\exp(2\gamma_w)\cdot M^{27\alpha_0}$, and this constant is at most $C_w$ (Section~\ref{sec:notation}). Since $\mathscr{Z}(s')=H_k(s')\,C\,C_k^{\frac{1}{2}}(s')\,\mathsf{Z}$, each of the three norms of $\mathscr{Z}(s')$ is at most
\[
C_w^2\cdot c_{\gamma}^2\cdot c_C\cdot e_k^{-\frac{\varepsilon}{4}}=\left(C_w^2\cdot c_{\gamma}^2\cdot c_C\cdot e_k^{\frac{\varepsilon}{4}}\right)\cdot e_k^{-\frac{\varepsilon}{2}}\le\frac{1}{2}e_k^{-\frac{\varepsilon}{2}},
\]
the last inequality being condition~\eqref{eq:renormalisation-group-32}.

Let now $\mathfrak{Z}$ and $t$ be as in~\eqref{it:bounds-analyticity-q1}. Each of the three norms of $t\cdot\mathfrak{Z}$ is at most
\[
\rho_1\cdot\frac{1}{2}e_k^{-\frac{\varepsilon}{2}}=\frac{1}{4}e_k^{-\frac{3}{4}+4\varepsilon}<e_k^{\varepsilon}\cdot e_k^{-\frac{3}{4}+3\varepsilon}=e_k^{\varepsilon}\cdot\boldsymbol{\theta}_3.
\]
The threshold $\boldsymbol{\theta}_3$ is the smallest of the thresholds $\boldsymbol{\theta}_1$, $\boldsymbol{\theta}_2$, $\boldsymbol{\theta}_3$ of $R_k$ (Section~\ref{sec:elementary-properties-field}). Hence each of the three norms of $t\cdot\mathfrak{Z}$ is smaller than $e_k^{\varepsilon}$ times the corresponding threshold, that is, $t\cdot\mathfrak{Z}\in e_k^{\varepsilon}R_k$. By Definition~\ref{def:domains}, $A=A_0+A_1$, where $A_1\in\left(1-e_k^{\varepsilon}\right)R_k$, and where $A_0$ is real and has, for every $M$-cube $\square$, a representative on $\square^{\dagger}$ whose three norms are below the thresholds of $R_k$, as in the definition of $\tilde{R}_k$. Each of the three norms is a seminorm. So, for $i=1,2,3$, the $i$-th norm of $A_1+t\cdot\mathfrak{Z}$ is smaller than $\left(1-e_k^{\varepsilon}\right)\cdot\boldsymbol{\theta}_i+e_k^{\varepsilon}\cdot\boldsymbol{\theta}_i=\boldsymbol{\theta}_i$; that is, $A_1+t\cdot\mathfrak{Z}\in R_k$, and therefore
\[
A+t\cdot\mathfrak{Z}=A_0+\left(A_1+t\cdot\mathfrak{Z}\right)\in\tilde{R}_k.
\]
Consequently $\mathscr{T}+t\cdot\mathfrak{Z}=\left(A+t\cdot\mathfrak{Z}\right)+\mathscr{C}\in\tilde{R}_k+\mathfrak{X}_k$, and $\tilde{R}_k+\mathfrak{X}_k\subseteq\tilde{G}_k$ by Lemma~\ref{lem:inclusions-domains}(a).

\emph{Part~\eqref{it:bounds-analyticity-q2}.} Let $k\ge1$. With spinor indices $a$ and $b$,
\[
\langle\bar{\Psi},M_k(X,B)\Psi\rangle=\sum_{y,y'\in T^0_k}\sum_{a,b}\bar{\Psi}_a(y)\cdot M_k(X,B)(y,a;y',b)\cdot\Psi_b(y').
\]
Up to sign, each term is one of the monomials of degree two in the generators $\Psi(u)$, and different terms give different monomials. By the definition of the norm in Proposition~\ref{prop:weighted-coefficient-norms}, the norm of this element is therefore $\mathsf{w}_3^2$ times the sum of the moduli of the entries $M_k(X,B)(y,a;y',b)$. By property~\ref{it:renormalisation-group-12}, an entry vanishes unless $y$ and $y'$ lie in $X$. The $M$-cubes partition the unit sites (fact~\eqref{it:renormalisation-group-g0}), and each of them contains $M^3$ unit sites; hence at most $4M^3\cdot|X|_M$ rows $(y,a)$ are different from zero. Every entry of a $4\times4$ block is bounded by the modulus of the block, in either reading of Section~\ref{sec:fermion-operators-bounds}; so property~\ref{it:renormalisation-group-6} gives
\[
\left|M_k(X,B)(y,a;y',b)\right|\le C_w\cdot\exp\left(-\gamma_w\cdot d(y,y')-\kappa_w\cdot d_M(X)\right).
\]
In the definition of $c_{\gamma}$ (Section~\ref{sec:notation}) one of the three sets over which the index $x'$ ranges is the set of the pairs formed by a unit site and a spinor index; consequently the moduli of the entries of one row add up to at most $c_{\gamma}\cdot C_w\cdot\exp\left(-\kappa_w\cdot d_M(X)\right)$. Finally, by~\eqref{eq:renormalisation-group-45} and~\eqref{eq:renormalisation-group-50},
\[
|X|_M\le c_X\cdot\left(d_M(X)+1\right)\le c_X\cdot\exp\left(d_M(X)\right).
\]
Multiplying the number of the rows, the bound for one row and $\mathsf{w}_3^2$, one obtains
\[
\left\|\langle\bar{\Psi},M_k(X,B)\Psi\rangle\right\|_{\mathsf{w}_3}\le\mathsf{w}_3^2\cdot4M^3\cdot c_X\cdot\exp\left(d_M(X)\right)\cdot c_{\gamma}\cdot C_w\cdot\exp\left(-\kappa_w\cdot d_M(X)\right),
\]
which is the assertion.

\emph{Part~\eqref{it:bounds-analyticity-q3}.} In the norm $|\cdot|_{1,\infty}$ of Proposition~\ref{prop:weighted-coefficient-norms} the supremum is taken over the sites, or pairs, and over the spinor components. Accordingly, property~\ref{it:renormalisation-group-5} is used with the modulus of a vector of $\mathbb{C}^4$ read as the largest modulus of its components; both readings of this modulus, the Euclidean one and this one, are part of the hypothesis~\ref{it:renormalisation-group-5} (Sections~\ref{sec:properties-systems-primitives} and~\ref{sec:fermion-operators-bounds}). Let $v$ be a spinor field on $T^0_k$, and put $f:=H_k(s'',B)\,v$. By property~\ref{it:renormalisation-group-5},
\[
|f|_{\infty}\le C_w\cdot|v|_{\infty},\qquad\left|\left(\delta_{\alpha,B}f\right)(x,x')\right|\le C_w\cdot|v|_{\infty}\quad\text{for }(x,x')\in\mathrm{Pz}_k\text{ with }|x-x'|<1.
\]
For $k=0$, property~\ref{it:renormalisation-group-17} gives the first of these bounds with the constant $4\exp(2\gamma_w)\cdot M^{27\alpha_0}\le C_w$, and no pair $(x,x')\in\mathrm{Pz}_0$ satisfies $|x-x'|<1$, because $\eta_0=1$. At a pair $(x,x')\in\mathrm{Pz}_k$ with $|x-x'|\ge1$, Lemma~\ref{lem:phases-differences-along} gives
\[
\left|\left(\delta_{\alpha,B}f\right)(x,x')\right|\le4|f|_{\infty}\le4C_w\cdot|v|_{\infty}.
\]
Hence $\left|H_k^{\#}(s'',B)\,v\right|_{\infty}\le4C_w\cdot|v|_{\infty}$ for every $v$, and by Proposition~\ref{prop:weighted-coefficient-norms} this is the first bound on the unbarred fields. On the barred fields,
\[
\bar{H}_k(s'',B)=\mathsf{C}^{-1}\cdot H_k(s'',-B)\cdot\mathsf{C},\qquad\delta_{\alpha,-B}\bar{H}_k(s'',B)=\mathsf{C}^{-1}\cdot\left[\delta_{\alpha,-B}H_k(s'',-B)\right]\cdot\mathsf{C},
\]
since $\delta_{\alpha,-B}$ acts on the arguments in $T_k$ and commutes with a constant matrix acting on the spinor components. By Proposition~\ref{prop:symmetries-extraction}(b), $\mathsf{C}$ is a monomial matrix whose non-zero entries have modulus one, so that $\mathsf{C}$ and $\mathsf{C}^{-1}$ preserve the largest modulus of the components; and $-B\in\tilde{G}_k$ by Lemma~\ref{lem:symmetries-domains}. The bound on the unbarred fields, applied at $-B$, therefore gives the same bound on the barred fields.

Let now $\mathfrak{Z}$ be as in~\eqref{it:bounds-analyticity-q1}. By~\eqref{it:bounds-analyticity-q1}, $\mathscr{T}+t\cdot\mathfrak{Z}\in\tilde{G}_k$ for $|t|\le\rho_1$. As $\tilde{G}_k$ is open (Lemma~\ref{lem:openness}), this holds for all $t$ in a neighbourhood of the closed disc $\left\{t\in\mathbb{C}:|t|\le\rho_1\right\}$, and by Lemma~\ref{lem:symmetries-domains} also $-\left(\mathscr{T}+t\cdot\mathfrak{Z}\right)\in\tilde{G}_k$ for these $t$. By property~\ref{it:renormalisation-group-2}, for fixed $s''\in\mathrm{Pd}_k$ the operator $H_k(s'',B)$ is analytic in $B\in\tilde{G}_k$; the phases $\exp\left(\pm i\,e_k\eta_k\cdot(\tau B)(x,x')\right)$ of $\delta_{\alpha,\pm B}$ (Definition~\ref{def:pairs-sites-holder}) are entire functions of $B$. Hence, for every spinor field $v$ on $T^0_k$, the function
\[
g_v(t):=H_k^{\#}\left(s'',\mathscr{T}+t\cdot\mathfrak{Z}\right)v
\]
is analytic on a neighbourhood of $\left\{t\in\mathbb{C}:|t|\le\rho_1\right\}$, on the unbarred and on the barred fields. Since $\rho_1\ge2$, the properties of the functional $\oint_{\rho}$ stated in Section~\ref{sec:notation}, applied to each component, give
\[
J_k^{\#}(s'',\mathscr{T},\mathfrak{Z})\,v=g_v(1)-g_v(0)=\oint_{\rho_1}g_v,\qquad\left|J_k^{\#}(s'',\mathscr{T},\mathfrak{Z})\,v\right|_{\infty}\le\frac{1}{\rho_1-1}\cdot\sup_{|t|=\rho_1}\left|g_v(t)\right|_{\infty}.
\]
By the first bound, applied at $B=\mathscr{T}+t\cdot\mathfrak{Z}$, one has $\left|g_v(t)\right|_{\infty}\le4C_w\cdot|v|_{\infty}$ for $|t|=\rho_1$; moreover $\frac{1}{\rho_1-1}\le\frac{2}{\rho_1}=4e_k^{\frac{3}{4}-\frac{9\varepsilon}{2}}$ (Section~\ref{sec:notation}). Therefore
\[
\left|J_k^{\#}(s'',\mathscr{T},\mathfrak{Z})\,v\right|_{\infty}\le4C_w\cdot|v|_{\infty}\cdot\frac{2}{\rho_1}=16C_w\cdot e_k^{\frac{3}{4}-\frac{9\varepsilon}{2}}\cdot|v|_{\infty}.
\]

\emph{Part~\eqref{it:bounds-analyticity-q4}.} By~\eqref{it:bounds-analyticity-q1} with $t=0$, $\mathscr{T}\in\tilde{G}_k$, so that~\eqref{it:bounds-analyticity-q3} applies at $B=\mathscr{T}$. The norm $|\cdot|_{1,\infty}$ is submultiplicative. On the generators $W$ the second factor is $\Gamma_k(s'',\mathscr{T})$. By property~\ref{it:renormalisation-group-6} and the definition of $c_{\gamma}$, as in the proof of~\eqref{it:bounds-analyticity-q2}, the moduli of the entries of a row of $\Gamma_k(s'',\mathscr{T})$ add up to at most $C_w\cdot c_{\gamma}$, that is, $\left|\Gamma_k(s'',\mathscr{T})\right|_{1,\infty}\le C_w\cdot c_{\gamma}$. Hence
\[
\left|H_k^{\#}(s'',\mathscr{T})\cdot\Gamma_k(s'',\mathscr{T})\right|_{1,\infty}\le4C_w\cdot C_w\cdot c_{\gamma}=C_W.
\]
On the generators $\bar{W}$ the second factor is the identity, and $\left|H_k^{\#}(s'',\mathscr{T})\right|_{1,\infty}\le4C_w\le C_W$, because $C_w\ge1$ and $c_{\gamma}\ge1$ (Section~\ref{sec:notation}).

\emph{Part~\eqref{it:bounds-analyticity-q5}.} By property~\ref{it:renormalisation-group-2} for the families~\eqref{it:primitives-p1} and~\eqref{it:primitives-p2}, the operators $H_k(s')$ and $C_k^{\frac{1}{2}}(s')$ are continuous in $s'$ on $\mathrm{Pd}_k$ and analytic in $s'$ in the interior, and $C$ is a fixed linear map. This gives the assertions on $\mathscr{Z}(s')=H_k(s')\,C\,C_k^{\frac{1}{2}}(s')\,\mathsf{Z}$. Consequently the map
\[
(s',t,A,\mathscr{C},\mathsf{Z})\longmapsto\mathscr{T}+t\cdot\mathscr{Z}(s')=A+\mathscr{C}+t\cdot\mathscr{Z}(s')
\]
is continuous on the set displayed in~\eqref{it:bounds-analyticity-q5}, analytic in $(s',t,A,\mathscr{C})$ for $s'$ in the interior of $\mathrm{Pd}_k$, and affine in $(t,A,\mathscr{C})$ for fixed $(s',\mathsf{Z})$. By~\eqref{it:bounds-analyticity-q1}, applied with $\mathfrak{Z}=\mathscr{Z}(s')$, it takes its values in $\tilde{G}_k$. The sets $\tilde{R}_k^{-}$, $\mathfrak{X}_k$ and $\tilde{G}_k$ are open (Lemma~\ref{lem:openness}); hence, for fixed $(s',\mathsf{Z})$, the set of the points $(t,A,\mathscr{C})$ at which $A+\mathscr{C}+t\cdot\mathscr{Z}(s')$ lies in $\tilde{G}_k$ is an open neighbourhood of $\left\{t\in\mathbb{C}:|t|\le\rho_1\right\}\times\tilde{R}_k^{-}\times\mathfrak{X}_k$. By property~\ref{it:renormalisation-group-2}, $M_k(X,\cdot)$ is analytic on $\tilde{G}_k$, while $H_k(s'',\cdot)$ and $\Gamma_k(s'',\cdot)$ are continuous on $\mathrm{Pd}_k\times\tilde{G}_k$, analytic in $(s'',B)$ for $s''$ in the interior, and analytic in $B\in\tilde{G}_k$ for every fixed $s''\in\mathrm{Pd}_k$. The entries of $\delta_{\alpha,B}H_k(s'',B)$ are finite sums of products of entries of $H_k(s'',B)$ with the numbers $|x-x'|^{-\alpha}$ and with the phases of Definition~\ref{def:pairs-sites-holder}, which are entire functions of $B$; and the operators on the barred fields are those at $-B$, which lies in $\tilde{G}_k$ by Lemma~\ref{lem:symmetries-domains}, conjugated by $\mathsf{C}$. Composing these maps with the map displayed above yields the continuity, the analyticity in $(s,t,A,\mathscr{C})$ in the interior, and, for fixed $s$ and $\mathsf{Z}$, the analyticity in $(t,A,\mathscr{C})$ on the neighbourhood just described.
\end{proof}

\subsection{Expansion in the parameters and resummation}\label{sec:expansion-parameters-resummation}

\begin{lemma}[Expansion in the parameters]\label{lem:expansion-parameters}
Let $\mathsf{P}$ be a finite set, $P_0\subseteq\mathsf{P}$ and $\rho>1$. Let $f$ be a continuous map from the closed polydisk
\[
\left\{s=(s_c)_{c\in\mathsf{P}}\in\mathbb{C}^{\mathsf{P}}:|s_c|\le\rho\ \text{for every }c\in\mathsf{P}\right\}
\]
to a finite-dimensional normed space, analytic in the interior of the polydisk, such that $\|f(s)\|\le B_f$ for every $s$ in the polydisk. For $P\subseteq\mathsf{P}\setminus P_0$ put
\[
f_P:=\int_{[0,1]^P}\mathrm{d}s_P\,\left[\frac{\partial}{\partial s_P}f(s)\right]_{s=1\text{ on }P_0,\ s=0\text{ outside }P_0\cup P},
\]
where $\mathrm{d}s_P:=\prod_{c\in P}\mathrm{d}s_c$ and $\frac{\partial}{\partial s_P}:=\prod_{c\in P}\frac{\partial}{\partial s_c}$; the variables $s_c$, $c\in P$, are integrated over $[0,1]$, and the remaining variables are fixed: $s_c=1$ for $c\in P_0$, and $s_c=0$ for $c\notin P_0\cup P$. Let $f(1)$ denote the value of $f$ at the point all of whose coordinates are $1$. Then
\[
f(1)=\sum_{P\subseteq\mathsf{P}\setminus P_0}f_P,\qquad\|f_P\|\le B_f\cdot(\rho-1)^{-\#|P|}.
\]
\end{lemma}

\begin{proof}
Since $\rho>1$, the cube $[0,1]^{\mathsf{P}}$ lies in the interior of the polydisk, where $f$ is analytic; so all partial derivatives of $f$ exist and are continuous on a neighbourhood of $[0,1]^{\mathsf{P}}$. For $c\in\mathsf{P}\setminus P_0$ and a function $g$ of the variable $s_c$ that is continuously differentiable on $[0,1]$, the fundamental theorem of calculus gives
\[
g(1)=g(0)+\int_0^1\mathrm{d}s_c\,g'(s_c).
\]
Fix $s_c=1$ for $c\in P_0$, and apply this identity to the variables $s_c$, $c\in\mathsf{P}\setminus P_0$, one after the other: first to $f$, then to each of the terms obtained, as a function of the next variable. The functions to which the identity is applied are continuously differentiable, and their derivatives are obtained by differentiation under the integral sign, because the partial derivatives of $f$ are continuous on a neighbourhood of $[0,1]^{\mathsf{P}}$. After the last variable, $f(1)$ is written as a sum of terms indexed by the subsets $P\subseteq\mathsf{P}\setminus P_0$, where $P$ is the set of the variables for which the second term of the identity was taken; the term indexed by $P$ is $f_P$. This proves the identity.

Let $s^0\in\mathbb{C}^{\mathsf{P}}$ be a point with $s^0_c\in[0,1]$ for $c\in P$, $s^0_c=1$ for $c\in P_0$ and $s^0_c=0$ for $c\notin P_0\cup P$. The closed polydisk of the points $s$ with $|s_c-s^0_c|\le\rho-1$ for $c\in P$ and $s_c=s^0_c$ for $c\notin P$ lies in the domain of $f$, since $|s_c|\le\rho$ for all its points. For $0<r'<\rho-1$, the closed polydisk of the radii $r'$ about $s^0$ in the variables $s_c$, $c\in P$, the other variables being fixed at $s^0$, lies in the interior of the domain, where $f$ is analytic. Cauchy's estimates on this polydisk give
\[
\left\|\frac{\partial}{\partial s_P}f(s^0)\right\|\le B_f\cdot(r')^{-\#|P|},
\]
and letting $r'$ tend to $\rho-1$ one obtains $\left\|\frac{\partial}{\partial s_P}f(s^0)\right\|\le B_f\cdot(\rho-1)^{-\#|P|}$. Integration over $[0,1]^P$, a set of volume $1$, gives the bound on $\|f_P\|$.
\end{proof}

\begin{lemma}[Resummation over polymers]\label{lem:resummation-polymers}
Let $\kappa_0\le\kappa'\le\kappa$ and $B_0\ge0$. Let the parameter cubes be those of $\pi_k$, or those of $\pi_k$, $\pi'_k$ and $\pi''_k$; let $\mathsf{P}$ be the set of all parameter cubes, and let $P_0(X)$ and $Y'(X,P)$ be as in Section~\ref{sec:localisation-map}. For every $X\in D_k$ let $f^X$ be a map as in Lemma~\ref{lem:expansion-parameters}, with values in a finite-dimensional normed space that is the same for all $X$, with
\[
\rho=M^{\alpha_0},\qquad P_0=P_0(X),\qquad B_f=B_0\cdot\exp\left(-\kappa'\cdot d_M(X)\right),
\]
and such that $f^X_P=0$ unless $Y'(X,P)$ is a polymer. For $\mathbf{Y}\in D_k$ put
\[
F(\mathbf{Y}):=\sum_{(X,P):\,Y'(X,P)=\mathbf{Y}}f^X_P,
\]
the sum being taken over the pairs formed by a polymer $X\in D_k$ and a set $P\subseteq\mathsf{P}\setminus P_0(X)$. Then
\[
\sum_{\mathbf{Y}\in D_k}F(\mathbf{Y})=\sum_{X\in D_k}f^X(1),\qquad\left\|F(\mathbf{Y})\right\|\le4c_X\cdot K_0\cdot B_0\cdot\exp\left(-(\kappa'-\kappa_0-1)\cdot d_M(\mathbf{Y})\right).
\]
\end{lemma}

\begin{proof}
By Lemma~\ref{lem:expansion-parameters}, $f^X(1)=\sum_Pf^X_P$ for every $X\in D_k$, the sum being taken over the sets $P\subseteq\mathsf{P}\setminus P_0(X)$. Summing over $X\in D_k$, and omitting the terms for which $Y'(X,P)$ is not a polymer, which vanish, one obtains the identity: every remaining term $f^X_P$ occurs in $F(\mathbf{Y})$ for exactly one $\mathbf{Y}\in D_k$, namely for $\mathbf{Y}=Y'(X,P)$.

Let $\mathbf{Y}\in D_k$, and let $(X,P)$ be a pair with $Y'(X,P)=\mathbf{Y}$. Every $M$-cube of $X$ has lattice sites in common with $X$; so $X\subseteq\mathbf{Y}$, and $d_M(X)\le d_M(\mathbf{Y})$, because a net that meets every closed $M$-cube of $\mathbf{Y}$ meets every closed $M$-cube of $X$ (Section~\ref{sec:lattice-geometry}). Let $\square$ be an $M$-cube of $\mathbf{Y}$ that is not contained in $X$. Different $M$-cubes have no lattice site in common (fact~\eqref{it:renormalisation-group-g0}); hence $\square$ meets neither $X$ nor a cube of $P_0(X)$, the cubes of $P_0(X)$ being contained in $X$. By the definition of $Y'(X,P)$, the cube $\square$ has a lattice site in common with $X\cup Y^{\mathrm{p}}(X,P)$; therefore it meets a cube of $P$. A parameter cube meets at most $N_c$ of the $M$-cubes (fact~\eqref{it:renormalisation-group-g1}). Hence $|\mathbf{Y}\setminus X|_M\le N_c\cdot\#|P|$, and~\eqref{eq:renormalisation-group-47} yields
\[
0\le d_M(\mathbf{Y})-d_M(X)\le N_c\cdot\#|P|.
\]
Put $a_X:=\log(6c_X)$. By condition~\eqref{eq:renormalisation-group-30}, $\rho-1=M^{\alpha_0}-1\ge\exp(\hat{\kappa})$, where $\hat{\kappa}=N_c\cdot\kappa+a_X$. Since $\kappa_0\ge1$ (Section~\ref{sec:recalled-auxiliary-results}) and $\kappa_0\le\kappa'\le\kappa$, one has $0\le\kappa'-\kappa_0\le\kappa$, and so
\begin{align*}
(\rho-1)^{-\#|P|}&\le\exp\left(-\hat{\kappa}\cdot\#|P|\right)=\exp\left(-N_c\cdot\kappa\cdot\#|P|\right)\cdot\exp\left(-a_X\cdot\#|P|\right)\\
&\le\exp\left(-(\kappa'-\kappa_0)\cdot\left(d_M(\mathbf{Y})-d_M(X)\right)\right)\cdot\exp\left(-a_X\cdot\#|P|\right).
\end{align*}
Together with the bound of Lemma~\ref{lem:expansion-parameters} and with $B_f=B_0\cdot\exp\left(-\kappa'\cdot d_M(X)\right)$ this gives
\begin{multline*}
\left\|f^X_P\right\|\le B_0\cdot\exp\left(-\kappa'\cdot d_M(X)\right)\cdot(\rho-1)^{-\#|P|}\\
\le B_0\cdot\exp\left(-(\kappa'-\kappa_0)\cdot d_M(\mathbf{Y})\right)\cdot\exp\left(-\kappa_0\cdot d_M(X)\right)\cdot\exp\left(-a_X\cdot\#|P|\right).
\end{multline*}

The cubes of $P$ are contained in $\mathbf{Y}$: every lattice site of a cube of $P$ lies in an $M$-cube (fact~\eqref{it:renormalisation-group-g0}), this $M$-cube meets the cube of $P$, and so it is one of the $M$-cubes of $Y'(X,P)=\mathbf{Y}$. By fact~\eqref{it:renormalisation-group-g3}, at most $3|\mathbf{Y}|_M$ parameter cubes are contained in $\mathbf{Y}$. Hence, the sum being taken over all sets $P$ of parameter cubes contained in $\mathbf{Y}$, and by~\eqref{eq:renormalisation-group-45} in the last step,
\begin{align*}
\sum_P\exp\left(-a_X\cdot\#|P|\right)&\le\left(1+\exp(-a_X)\right)^{3|\mathbf{Y}|_M}=\left(1+\frac{1}{6c_X}\right)^{3|\mathbf{Y}|_M}\\
&\le\exp\left(\frac{|\mathbf{Y}|_M}{2c_X}\right)\le\exp\left(\frac{d_M(\mathbf{Y})+1}{2}\right).
\end{align*}
By~\eqref{eq:renormalisation-group-46}, \eqref{eq:renormalisation-group-45} and~\eqref{eq:renormalisation-group-50},
\[
\sum_{X\in D_k,\,X\subseteq\mathbf{Y}}\exp\left(-\kappa_0\cdot d_M(X)\right)\le K_0\cdot|\mathbf{Y}|_M\le c_X\cdot K_0\cdot\left(d_M(\mathbf{Y})+1\right)\le2c_X\cdot K_0\cdot\exp\left(\frac{d_M(\mathbf{Y})}{2}\right).
\]
The sum of $\exp\left(-\kappa_0\cdot d_M(X)\right)\cdot\exp\left(-a_X\cdot\#|P|\right)$ over the pairs $(X,P)$ with $Y'(X,P)=\mathbf{Y}$ is at most the product of these two sums, hence at most
\[
2\exp\left(\frac{1}{2}\right)\cdot c_X\cdot K_0\cdot\exp\left(d_M(\mathbf{Y})\right)\le4c_X\cdot K_0\cdot\exp\left(d_M(\mathbf{Y})\right).
\]
With the bound on $\left\|f^X_P\right\|$ this gives
\begin{multline*}
\left\|F(\mathbf{Y})\right\|\le B_0\cdot\exp\left(-(\kappa'-\kappa_0)\cdot d_M(\mathbf{Y})\right)\cdot4c_X\cdot K_0\cdot\exp\left(d_M(\mathbf{Y})\right)\\
=4c_X\cdot K_0\cdot B_0\cdot\exp\left(-(\kappa'-\kappa_0-1)\cdot d_M(\mathbf{Y})\right).
\end{multline*}
\end{proof}

\subsection{Locality of the localised functions}\label{sec:locality-localised-functions}

The pairs $(x,x')\in\mathrm{Pz}_k$ that are admissible for a polymer $X$ are, by Definition~\ref{def:pairs-sites-holder}, those of kind~(I) and those of kind~(II) for $X$. Let $(x,x')$ be admissible for $X$, and let $X'\in D_k$ satisfy $X\subseteq X'$. Then
\begin{multline}\tag{adm}\label{eq:renormalisation-group-adm}
(x,x')\text{ is admissible for }X'\text{, of the same kind for }X'\text{ as for }X;\\
\text{for kind (I), each bond of the six paths }\Gamma^{\pi}(x,x')\text{ lies in }X;\\
\text{for kind (II), the bond }(x,x')\text{ is the only path, and it meets }X;\\
\text{each of }x,x'\text{ is within sup-distance }\eta_k\le1\text{ of a lattice site of }X.
\end{multline}
These statements follow from Definition~\ref{def:pairs-sites-holder}: for kind~(I), both sites and the six paths lie in $X$, hence in $X'$; for kind~(II), $x$ and $x'$ are nearest neighbours in $T_k$, so that $|x-x'|=\eta_k$, at least one of them lies in $X\subseteq X'$, and each of the six rectilinear paths between nearest neighbours is the single bond joining them; and $\eta_k=L^{-k}\le1$.

Let $F$ be a function of $(A,\mathscr{C},\mathsf{Z},s)$, for example a coefficient of an element of the Grassmann algebra, and let a set of parameter cubes be prescribed outside of which $s$ has to vanish. The statement that $F$ \emph{depends on $(A,\mathscr{C},\mathsf{Z},s)$ through certain data only} means: whenever two points of $\tilde{R}_k^{-}\times\mathfrak{X}_k\times\mathrm{Box}_k\times\mathrm{Pd}_k$ have parameter families vanishing outside the prescribed set and agree in these data, $F$ takes the same value at both.

\begin{lemma}\label{lem:renormalisation-group-1}
Let $G$ denote one of the six functions defined in~\eqref{eq:renormalisation-group-59}, namely $G^{(1)}$, $G^{(2)}$, $G^{(3)}$, which belong to the first localisation, or $G^{(\le3)}$, $G^{(4)}$, $G^{(5)}$, which belong to the second; for $k=0$ these definitions are taken with the modifications of~\eqref{eq:renormalisation-group-60}. In the first localisation the parameter cubes are the cubes of the kinds $\pi_k$, $\pi'_k$ and $\pi''_k$; in the second they are the cubes of kind $\pi_k$, and $s=s''$. Fix $X\in D_k$ and a set $P$ of parameter cubes disjoint from $P_0(X)$, and let $s\in\mathrm{Pd}_k$ vanish at every parameter cube outside $P_0(X)\cup P$. Write $Y':=Y'(X,P)$ (Section~\ref{sec:localisation-map}), and denote by $Y'_1$ the component of $Y'$ containing $X$.
\begin{enumerate}
\item[(a)] The generators occurring in $G(s,X)$ are among the following: $\psi(\xi)$ with $\xi$ in $X$; $\chi(\boldsymbol{\zeta})$ with $\boldsymbol{\zeta}$ admissible for $X$; $\Psi(u)$ and $W(u)$ with $u$ in $Y'_1$ (that is, the site of $\xi$ belongs to $X$, and the unit site of $u$ belongs to $Y'_1$). The coefficients of $G(s,X)$ depend on $(A,\mathscr{C},\mathsf{Z},s)$ through the following data only: the coordinates $\mathsf{Z}(\Upsilon)$ of the basis elements $\Upsilon$ contained in $Y'_1$; the values of $A$ and of $\mathscr{C}$ at the bonds meeting $Y'_1$; the parameters $s_c$ of the cubes $c$ contained in $Y'_1$.
\item[(b)] If $Y'$ is not connected, then $\frac{\partial}{\partial s_P}G(s,X)=0$.
\end{enumerate}
\end{lemma}

\begin{proof}
For $k=0$, property~\ref{it:renormalisation-group-17} contains properties~\ref{it:renormalisation-group-8}, \ref{it:renormalisation-group-10} and~\ref{it:renormalisation-group-11} of the diagonal families $H_0(s')$ and $H_0(s'',\cdot)$, property~\ref{it:renormalisation-group-8} of the family~\eqref{it:primitives-p2} and of $\Gamma_0$, and property~\ref{it:renormalisation-group-10} of the family~\eqref{it:primitives-p2}. So each of these properties holds at every level $k\ge0$ for the families to which Steps~(1), (2) and~(4) apply it.

\emph{Step~(0): a dichotomy for link-components.} For each of the three kinds $\pi'_k$, $\pi''_k$, $\pi_k$, take the cubes of that kind in $P_0(X)\cup P$ and form their union; let $\mathrm{LC}$ be the collection of all link-components of these three unions. An element $\mathcal{L}\in\mathrm{LC}$ is regarded as a set of parameter cubes of a single kind, as in fact~\eqref{it:renormalisation-group-g2}, and, as in property~\ref{it:renormalisation-group-8}, the same letter $\mathcal{L}$ also denotes the union of these cubes. Fix $\mathcal{L}\in\mathrm{LC}$. Being a link-component, $\mathcal{L}$ is non-empty and link-connected, and fact~\eqref{it:renormalisation-group-g2} shows that $\mathrm{pol}(\mathcal{L})$ is a polymer containing every cube of $\mathcal{L}$. Each $M$-cube of $\mathrm{pol}(\mathcal{L})$ meets a cube of $\mathcal{L}$, and this cube is a subset of $Y^{\mathrm{p}}(X,P)$; by the definition of $Y'(X,P)$, each such $M$-cube belongs to $Y'$. As a polymer is connected through faces, $\mathrm{pol}(\mathcal{L})$ is contained in a single component of $Y'$. Hence one of two alternatives holds: $\mathrm{pol}(\mathcal{L})\subseteq Y'_1$, or $\mathrm{pol}(\mathcal{L})$ and $Y'_1$ have no $M$-cube in common. In the first alternative $\mathcal{L}$ is said to be \emph{attached}; every cube of an attached $\mathcal{L}$ is contained in $\mathrm{pol}(\mathcal{L})$, hence in $Y'_1$. A sufficient condition for $\mathcal{L}$ to be attached is the existence of a lattice site $p$ of $Y'_1$ and of a cube $c$ of $\mathcal{L}$ such that $p\in c$, or such that $c$ is of kind $\pi'_k$ or $\pi''_k$ and $p\in c^{+}$. Indeed, by fact~\eqref{it:renormalisation-group-g0} exactly one $M$-cube contains $p$, and as $p$ is a lattice site of the union $Y'_1$ of $M$-cubes, this $M$-cube is one of the $M$-cubes of $Y'_1$. This $M$-cube meets $c$. If $p\in c$, the lattice site $p$ is common to the two cubes. If $p\in c^{+}$, the same uniqueness shows that the $M$-cube of $p$ is one of the $M$-cubes of which $c^{+}$ is the union, and by fact~\eqref{it:renormalisation-group-g1} these are exactly the $M$-cubes meeting $c$. In both cases the $M$-cube of $p$ meets a cube of $\mathcal{L}$; by the definition of $\mathrm{pol}(\mathcal{L})$ it is therefore an $M$-cube of $\mathrm{pol}(\mathcal{L})$, and it is one of $Y'_1$, which excludes the second alternative.

\emph{Step~(1): the field $\mathscr{Z}(s')$ near $Y'_1$.} Let $b$ be a bond of $T_k$ with an end point $p$ in $Y'_1$. By~\eqref{eq:renormalisation-group-56},
\[
\mathscr{Z}(s')(b)=\sum_{b'}\sum_{\Upsilon,\Upsilon'}H_k(s')(b,b')\cdot C(b',\Upsilon)\cdot C_k^{\frac{1}{2}}(s')(\Upsilon,\Upsilon')\cdot\mathsf{Z}(\Upsilon'),
\]
the sums extending over the unit bonds $b'$ and over the elements $\Upsilon$, $\Upsilon'$ of the basis of Proposition~\ref{prop:basis-map-box}. The parameters of $H_k(s')$ sit at the cubes of kind $\pi'_k$ and, as $s$ does, vanish outside $P_0(X)\cup P$. If no element of $\mathrm{LC}$ of kind $\pi'_k$ contains $b$, the row $H_k(s')(b,\cdot)$ is zero, by property~\ref{it:renormalisation-group-10}, and so is $\mathscr{Z}(s')(b)$. Otherwise let $\mathcal{L}_1\in\mathrm{LC}$ be an element of kind $\pi'_k$ that contains $b$. By property~\ref{it:renormalisation-group-8}, $H_k(s')(b,b')=0$ for every $b'$ not contained in $\mathcal{L}_1$, and the row at $b$ is determined by the parameters at the cubes of $\mathcal{L}_1$. The end point $p$ of $b$ is a lattice site of $Y'_1$ and lies in the union $\mathcal{L}_1$, that is, in a cube of $\mathcal{L}_1$; by Step~(0), $\mathcal{L}_1$ is attached.

Proposition~\ref{prop:basis-map-box} allows $C(b',\Upsilon)\ne0$ only if $\operatorname{supp}\Upsilon$ is contained in the union of the $L$-blocks that contain the end points of $b'$; this union is a box with at most $2L$ unit sites along each edge, so that two of its points have sup-distance less than $2L$. The parameters of $C_k^{\frac{1}{2}}(s')$ sit at the cubes of kind $\pi''_k$ and also vanish outside $P_0(X)\cup P$. If no element of $\mathrm{LC}$ of kind $\pi''_k$ contains $\Upsilon$, the row $C_k^{\frac{1}{2}}(s')(\Upsilon,\cdot)$ is zero, by property~\ref{it:renormalisation-group-10}. If $\Upsilon$ is contained in an element $\mathcal{L}_2\in\mathrm{LC}$ of kind $\pi''_k$, property~\ref{it:renormalisation-group-8} gives $C_k^{\frac{1}{2}}(s')(\Upsilon,\Upsilon')=0$ for every $\Upsilon'$ not contained in $\mathcal{L}_2$, and the row at $\Upsilon$ is determined by the parameters at the cubes of $\mathcal{L}_2$.

Thus only the terms with $b'$ contained in $\mathcal{L}_1$, with $C(b',\Upsilon)\ne0$, and with $\Upsilon$ contained in some $\mathcal{L}_2\in\mathrm{LC}$ of kind $\pi''_k$ can be non-zero. In such a term choose $q\in\operatorname{supp}\Upsilon$; then $q$ lies in a cube $c'$ of $\mathcal{L}_2$. As $q$ and the end points of $b'$ lie in the box just described, $|p'-q|<2L$ for an end point $p'$ of $b'$, and fact~\eqref{it:renormalisation-group-g1} yields $p'\in c'^{+}$. Being an end point of $b'$, the lattice site $p'$ lies in the union $\mathcal{L}_1$, which is contained in $Y'_1$ since $\mathcal{L}_1$ is attached. Step~(0) now shows that $\mathcal{L}_2$ is attached. Hence the cubes of $\mathcal{L}_1$ and of $\mathcal{L}_2$ are contained in $Y'_1$, and so is every basis element $\Upsilon'$ contained in $\mathcal{L}_2$. Consequently $\mathscr{Z}(s')(b)$ depends on the arguments only through the coordinates $\mathsf{Z}(\Upsilon')$ with $\Upsilon'$ contained in $Y'_1$ and the parameters at cubes contained in $Y'_1$.

\emph{Step~(2): rows of $H_k(s'',B)$ and $\bar{H}_k(s'',B)$ near $Y'_1$.} Let $B\in\tilde{G}_k$, let $p$ be a lattice site of $Y'_1$, and let $z$ be a site of $T_k$ with sup-distance $|z-p|\le1$. The family $s''$ of the parameters at the cubes of kind $\pi_k$ vanishes outside $P_0(X)\cup P$. If some cube of kind $\pi_k$ that meets $K_z$ (Section~\ref{sec:notation}) does not belong to $P_0(X)\cup P$, then $H_k(s'',B)(z,\cdot)=0$, by property~\ref{it:renormalisation-group-11}. Otherwise all these cubes belong to $P_0(X)\cup P$; they form a link-connected set, by fact~\eqref{it:renormalisation-group-g4}, and hence belong to a single element $\mathcal{L}_3\in\mathrm{LC}$. The $M$-cube containing $z$ meets $K_z$ and is therefore one of them, so that $z$ lies in the union $\mathcal{L}_3$. Since $|p-z|\le1<2L$, the site $p$ belongs to $K_z$ as well; the $M$-cube containing $p$ therefore meets $K_z$ and is a cube of $\mathcal{L}_3$ that contains $p$, and Step~(0) shows that $\mathcal{L}_3$ is attached. By property~\ref{it:renormalisation-group-8}, the non-zero entries of the row $H_k(s'',B)(z,\cdot)$ stand in columns contained in $\mathcal{L}_3$, and this row is determined by the parameters at the cubes of $\mathcal{L}_3$ together with the values of $B$ at the bonds meeting $\mathcal{L}_3$. The row of $\bar{H}_k(s'',B)=\mathsf{C}^{-1}\cdot H_k(s'',-B)\cdot\mathsf{C}$ (Section~\ref{sec:data-variables-fields}) has the same properties: Lemma~\ref{lem:symmetries-domains} gives $-B\in\tilde{G}_k$, so that what has been said applies to $H_k(s'',-B)$, and replacing $B$ by $-B$ does not change the bonds at which the connection is read. Because $\mathcal{L}_3$ is attached, every bond meeting $\mathcal{L}_3$ has an end point in $Y'_1$.

Let now $B=\mathscr{T}+t\cdot\mathscr{Z}(s')$ with $|t|\le\rho_1$. At a bond with an end point in $Y'_1$, the value of $\mathscr{T}=A+\mathscr{C}$ is the sum of the values of $A$ and of $\mathscr{C}$ at a bond meeting $Y'_1$, and the value of $\mathscr{Z}(s')$ is the subject of Step~(1). The values of $B$ at the bonds meeting $\mathcal{L}_3$ are therefore determined by the values of $A$ and of $\mathscr{C}$ at the bonds meeting $Y'_1$, by the coordinates $\mathsf{Z}(\Upsilon')$ with $\Upsilon'$ contained in $Y'_1$, and by the parameters at cubes contained in $Y'_1$. Property~\ref{it:renormalisation-group-8} compares rows at two connections in $\tilde{G}_k$. This hypothesis is met at both points of comparison: part~\eqref{it:bounds-analyticity-q1} of Lemma~\ref{lem:bounds-analyticity} gives $B\in\tilde{G}_k$, and Lemma~\ref{lem:symmetries-domains} then gives $-B\in\tilde{G}_k$.

\emph{Step~(3): H\"older derivatives at admissible pairs.} Let $B\in\tilde{G}_k$ and let $(x,x')$ be a pair admissible for $X$. By Definition~\ref{def:pairs-sites-holder}, for every spinor $f$ on $T^0_k$,
\begin{multline*}
\left(\delta_{\alpha,B}H_k(s'',B)f\right)(x,x')\\
=|x-x'|^{-\alpha}\cdot\left(\exp\left(i\cdot e_k\cdot\eta_k\cdot(\tau B)(x,x')\right)\cdot\left(H_k(s'',B)f\right)(x')-\left(H_k(s'',B)f\right)(x)\right),
\end{multline*}
and the same formula, with $\bar{H}_k(s'',B)$ in place of $H_k(s'',B)$ and with the phase $\exp\left(-i\cdot e_k\cdot\eta_k\cdot(\tau B)(x,x')\right)$, holds for $\delta_{\alpha,-B}\bar{H}_k(s'',B)$. So the row at $(x,x')$ is built from the phase and from the rows of the same operator at $x$ and at $x'$. By~\eqref{eq:renormalisation-group-adm}, each of $x$, $x'$ is within sup-distance $1$ of a lattice site of $X$, and $X\subseteq Y'_1$; hence Step~(2) describes the rows at $x$ and at $x'$. The phase depends on $B$ through the values at the bonds of the six paths $\Gamma^{\pi}(x,x')$ only. By~\eqref{eq:renormalisation-group-adm} these bonds meet $X$, so that they have an end point in $Y'_1$, and Step~(1) describes the values of $\mathscr{Z}(s')$ at them.

\emph{Step~(4): rows of $\Gamma_k(s'',\mathscr{T})$.} Let $y$ be a column at which some row treated in Step~(2) has a non-zero entry. Step~(2) places $y$ in $\mathcal{L}_3$, which is attached and is a link-component of the union of the cubes of kind $\pi_k$ in $P_0(X)\cup P$. The family $\Gamma_k(s'',\cdot)$ has its parameters at the cubes of this same kind $\pi_k$, and they are the same variables $s''$ (Section~\ref{sec:data-variables-fields}). Applied to $\Gamma_k$, property~\ref{it:renormalisation-group-8} shows that $\Gamma_k(s'',\mathscr{T})(y,y')=0$ for every $y'$ not contained in $\mathcal{L}_3$, and that the row at $y$ is determined by the parameters at the cubes of $\mathcal{L}_3$ and by the values of $\mathscr{T}$ at the bonds meeting $\mathcal{L}_3$. The hypothesis $\mathscr{T}\in\tilde{G}_k$ of property~\ref{it:renormalisation-group-8} is part~\eqref{it:bounds-analyticity-q1} of Lemma~\ref{lem:bounds-analyticity} at $t=0$.

\emph{Step~(5): $\mathscr{E}^r$, $M_k$ and the forms on $M$-cubes.} Proposition~\ref{prop:symmetries-extraction} shows that $R$ maps $K^{x,C}_k$ into itself, so $\mathscr{E}^r=R\,\mathscr{E}$ belongs to $K^{x,C}_k$ and has property~\eqref{it:renormalisation-group-1-k2} of Definition~\ref{def:renormalisation-group-1}, whose families are defined on $\tilde{R}_k\times\mathfrak{X}_k$. Therefore, for $(B_1,\mathscr{C})\in\tilde{R}_k\times\mathfrak{X}_k$, the only generators of $\mathscr{E}^r(X;B_1,\mathscr{C},\cdot)$ are $\psi(\xi)$ with $\xi$ in $X$ and $\chi(\boldsymbol{\zeta})$ with $\boldsymbol{\zeta}$ admissible for $X$, and, as $(B_1,\mathscr{C})$ ranges over $\tilde{R}_k\times\mathfrak{X}_k$, this element depends on $B_1$ and $\mathscr{C}$ through their values at the bonds meeting $X$ only. For $B\in\tilde{G}_k$, property~\ref{it:renormalisation-group-12} gives $M_k(X,B)(x,x')=0$ unless both $x$ and $x'$ lie in $X$, and $M_k(X,B)=M_k(X,\tilde B)$ for every $\tilde B\in\tilde{G}_k$ that coincides with $B$ at the bonds meeting $X$. For an $M$-cube $\square$, every generator $\psi(\xi)$ occurring in $\langle\bar{\psi},\psi\rangle_{\square}$ has its site in $\square$. Finally, at $k=0$ the function~\eqref{eq:renormalisation-group-60} of an $M$-cube $\square$ involves $\psi$ at the sites $x\in\square$, $\chi$ at the pairs $\left(x,x^{\boldsymbol{\sigma}}\right)$ with $x\in\square$, and $\mathscr{Z}(s')$ at the bonds $\left(x,x^{\boldsymbol{\sigma}}\right)$. Each such pair consists of nearest neighbours, the first of which lies in $\square$, so it is of kind~(II), hence admissible, for $\square$; and each such bond has its end point $x$ in $\square$.

\emph{Proof of~(a) for the first localisation.} Let $k\ge1$ first. For $|t|\le\rho_1$, part~\eqref{it:bounds-analyticity-q1} of Lemma~\ref{lem:bounds-analyticity} gives
\[
A+t\cdot\mathscr{Z}(s')\in\tilde{R}_k,\qquad\mathscr{T}+t\cdot\mathscr{Z}(s')\in\tilde{R}_k+\mathfrak{X}_k\subseteq\tilde{G}_k,
\]
and $\mathscr{C}\in\mathfrak{X}_k$. Since $1<\rho_2\le\rho_1$ (Section~\ref{sec:notation}), this applies on the circles $|t|=\rho_1$ and $|t|=\rho_2$ of the contour integrals and at $t=0$ and $t=1$, the values at which $J_k^{\#}\left(s'',\mathscr{T},\mathscr{Z}(s')\right)$ is formed; so the connections in the definitions of $G^{(1)}$, $G^{(2)}$, $G^{(3)}$ in~\eqref{eq:renormalisation-group-59} lie in the domains named in Steps~(2), (3) and~(5). In $G^{(1)}(s,X)$ the kernel of $M_k\left(X,\mathscr{T}+t\cdot\mathscr{Z}(s')\right)$ vanishes unless both indices lie in $X\subseteq Y'_1$, and only the values of $\mathscr{T}+t\cdot\mathscr{Z}(s')$ at the bonds meeting $X$ enter (Step~(5)); at these bonds, each of which has an end point in $Y'_1$, Step~(1) applies to $\mathscr{Z}(s')$. In $G^{(2)}(s,X)$ the element $\mathscr{E}^r\left(X;A+t\cdot\mathscr{Z}(s'),\mathscr{C},\cdot\right)$ contains $\psi(\xi)$ only with $\xi$ in $X$ and $\chi(\boldsymbol{\zeta})$ only with $\boldsymbol{\zeta}$ admissible for $X$, and only the values of $A+t\cdot\mathscr{Z}(s')$ and of $\mathscr{C}$ at the bonds meeting $X$ enter (Step~(5)); Step~(1) applies to $\mathscr{Z}(s')$ at these bonds once more. The function $G^{(3)}(s,X)$ is zero unless $X$ is a single $M$-cube $\square$ (Section~\ref{sec:localisation-exponent}), and in that case the generators $\psi(\xi)$ of $\langle\bar{\psi},\psi\rangle_{\square}$ have their sites in $\square$ (Step~(5)). In $G^{(2)}$, and in $G^{(3)}$ through its components without H\"older derivative, the field $t\cdot J_k^{\#}\left(s'',\mathscr{T},\mathscr{Z}(s')\right)\Psi$ is added to the generators; by what has just been shown, only those of its components enter that are indexed by a site of $X$ or by a pair admissible for $X$. Each of these components is obtained by applying to $\Psi$ the difference of the rows of $H_k^{\#}(s'',B)$ at $B=\mathscr{T}+\mathscr{Z}(s')$ and at $B=\mathscr{T}$, and these rows are described by Step~(2), with $p=z$, and by Step~(3): their non-zero columns are unit sites in attached link-components, hence in $Y'_1$, and their entries depend on the arguments only through data specified in~(a). Therefore $\Psi(u)$ occurs only with $u$ in $Y'_1$, and the coefficients depend on the arguments only through the data specified in~(a); the contour integrals in $t$ preserve both properties. For $k=0$ one has $G^{(2)}=G^{(3)}=0$, and Steps~(5) and~(1) give the assertion for the function $G^{(1)}$ of~\eqref{eq:renormalisation-group-60}.

\emph{Proof of~(b).} Assume that $Y'$ is not connected. Then some $M$-cube $\square$ of $Y'$ does not belong to $Y'_1$, and, by the definition of $Y'(X,P)$ (Section~\ref{sec:localisation-map}), $\square$ has a lattice site in common with $X\cup Y^{\mathrm{p}}(X,P)$. It has none in common with $X$: by fact~\eqref{it:renormalisation-group-g0}, an $M$-cube sharing a lattice site with $X$ is an $M$-cube of $X$, and the $M$-cubes of $X$ belong to $Y'_1$. So $\square$ shares a lattice site $q$ with a cube $c\in P_0(X)\cup P$, and since $q\notin X$, the cube $c$ is not contained in $X$. Every cube of $P_0(X)$ is contained in $X$; hence $c\in P$. By fact~\eqref{it:renormalisation-group-g0}, $\square$ is the only $M$-cube containing $q$, so $q\notin Y'_1$, and $c$ is not contained in $Y'_1$. A change of $s_c$, all other parameters being kept fixed, produces families that still vanish outside $P_0(X)\cup P$ and agree with $s$ at the cubes contained in $Y'_1$. By~(a), $G(s,X)$ is therefore independent of $s_c$ on the set of the families $s\in\mathrm{Pd}_k$ that vanish outside $P_0(X)\cup P$, and its derivative with respect to $s_c$ vanishes there. As $c\in P$, the operator $\frac{\partial}{\partial s_P}$ contains the factor $\frac{\partial}{\partial s_c}$, and $\frac{\partial}{\partial s_P}G(s,X)=0$.

\emph{Proof of~(a) for the second localisation.} Let $i\le3$. By~\eqref{eq:renormalisation-group-57}, $E^{(i),*}(X)=\mathrm{Loc}\left[G^{(i)}\right](X)$ is a sum of terms indexed by the pairs $(X_1,P_1)$ with $Y'(X_1,P_1)=X$. For each of them $X_1\subseteq X$, as every $M$-cube of $X_1$ belongs to $Y'(X_1,P_1)$, and, $X$ being connected, the component of $Y'(X_1,P_1)$ that contains $X_1$ is $X$ itself. Parts~(a) and~(b) for the first localisation, applied to these terms, give the following. The generators of $E^{(i),*}(X)$ are $\psi(\xi)$ with $\xi$ in $X$, $\chi(\boldsymbol{\zeta})$ with $\boldsymbol{\zeta}$ admissible for a polymer $X_1\subseteq X$, hence for $X$ by~\eqref{eq:renormalisation-group-adm}, and $\Psi(u)$ with $u$ in $X$; and its coefficients depend on the arguments only through the coordinates $\mathsf{Z}(\Upsilon)$ with $\Upsilon$ contained in $X$ and the values of $A$ and of $\mathscr{C}$ at the bonds meeting $X$. The derivatives and the integrals with respect to the parameters in~\eqref{eq:renormalisation-group-57} keep these properties. If $k\ge1$, the substitution~\eqref{eq:renormalisation-group-58} puts in the place of each $\Psi(u)$, $u$ in $X$, an element that contains, by Lemma~\ref{lem:interior-boundary-parts}(b), $\psi$ only at sites of $X$ and $\chi$ only at pairs of kind~(II) for $X$, and whose coefficients depend on $\mathscr{T}$ through its values at bonds meeting $X$ only; if $k=0$, no substitution is made. Hence $E^{(\le3),x}(X;A,\mathscr{C},\mathsf{Z},\cdot)$ has the locality established in Step~(5) for $\mathscr{E}^r$, and it depends in addition on the coordinates $\mathsf{Z}(\Upsilon)$ with $\Upsilon$ contained in $X$. Insert now Steps~(2)--(5) into the definitions of $G^{(\le3)}$, $G^{(4)}$, $G^{(5)}$ in~\eqref{eq:renormalisation-group-59}; part~\eqref{it:bounds-analyticity-q1} of Lemma~\ref{lem:bounds-analyticity} at $t=0$ gives $A\in\tilde{R}_k$ and $\mathscr{T}\in\tilde{G}_k$, so that, with $\mathscr{C}\in\mathfrak{X}_k$, the connections occurring there lie in the domains named in these steps. The field $W_k^{\#}(s'',\mathscr{T})$ is added to the generators in $G^{(\le3)}$, the field $t\cdot W_k^{\#}(s'',\mathscr{T})$ in $G^{(4)}$, and the components of the latter without H\"older derivative in $G^{(5)}$; by what precedes, only those components enter that are indexed by a site of $X$ or by a pair admissible for $X$. By~\eqref{eq:renormalisation-group-56}, these components are formed from the rows of $H_k^{\#}(s'',\mathscr{T})$ at those sites and pairs, which Steps~(2) and~(3) describe for $B=\mathscr{T}$, and from the rows of $\Gamma_k(s'',\mathscr{T})$ at the columns of the former, which Step~(4) describes. So the further generators are $W(u)$ with $u$ in $Y'_1$, and the coefficients depend on the arguments only through the data specified in~(a); the contour integrals in $t$ preserve this. Once~(a) is known for the second localisation, the proof of~(b) applies to it without change.
\end{proof}

\begin{corollary}\label{cor:renormalisation-group-1}
Let $\mathbf{Y}\in D_k$. The generators occurring in $E^{(i),x}(\mathbf{Y})$, $i\le3$, are of the forms $\psi(\xi)$ with $\xi$ in $\mathbf{Y}$ and $\chi(\boldsymbol{\zeta})$ with $\boldsymbol{\zeta}$ admissible for $\mathbf{Y}$; the generators occurring in $E^{(\le5),x}(\mathbf{Y})$ are of these two forms or of the form $W(u)$ with $u$ in $\mathbf{Y}$. In all these functions the coefficients depend on $(A,\mathscr{C},\mathsf{Z})$ only through the coordinates $\mathsf{Z}(\Upsilon)$ with $\Upsilon$ contained in $\mathbf{Y}$ and the values of $A$ and of $\mathscr{C}$ at the bonds meeting $\mathbf{Y}$. If two polymers have no $M$-cube in common, then they have no unit site in common, and no basis element $\Upsilon$ is contained in both.
\end{corollary}

\begin{proof}
The sum~\eqref{eq:renormalisation-group-57} defining $E^{(\le5),x}(\mathbf{Y})$ extends over the pairs $(X,P)$ with $Y'(X,P)=\mathbf{Y}$, and for each of them $Y'_1=Y'=\mathbf{Y}$, since $\mathbf{Y}$ is connected. To each term apply Lemma~\ref{lem:renormalisation-group-1}(a) for the second localisation; a pair admissible for $X$ is admissible for $\mathbf{Y}\supseteq X$, by~\eqref{eq:renormalisation-group-adm}; and the derivatives and integrals with respect to the parameters keep the properties so obtained. This proves the assertions on $E^{(\le5),x}(\mathbf{Y})$. Those on $E^{(i),x}(\mathbf{Y})$, $i\le3$, were established in the proof of Lemma~\ref{lem:renormalisation-group-1}(a) for the second localisation, with $\mathbf{Y}$ in place of $X$. Finally, by fact~\eqref{it:renormalisation-group-g0} the $M$-cubes partition the sites of $T_k$ and those of $T^0_k$; so a unit site common to two polymers would lie in an $M$-cube common to both, and a basis element contained in both would provide such a unit site, namely any point of its support, which is not empty.
\end{proof}

\subsection{Estimates for the localised functions}\label{sec:estimates-localised-functions}

The norms of Section~\ref{sec:recalled-auxiliary-results} are used below with the following weights. Put
\[
h_k/2:=\left(\frac{1}{2}h_{k,1},\,\frac{1}{2}h_{k,2}\right),\qquad h_k/4:=\left(\frac{1}{4}h_{k,1},\,\frac{1}{4}h_{k,2}\right);
\]
in each of these pairs the first entry is the weight of the generators $\psi$ and the second entry is the weight of the generators $\chi$. For a number $\mathsf{w}_3>0$, the norm $\|\cdot\|_{(h_k/2;\mathsf{w}_3)}$ on $\mathrm{Gr}[\psi,\chi,\Psi]$ is formed with the weights $h_k/2$ of $\psi$, $\chi$ and with the weight $\mathsf{w}_3$ of $\Psi$, and the norm $\|\cdot\|_{(h_k/4;1)}$ on $\mathrm{Gr}[\psi,\chi,W]$ is formed with the weights $h_k/4$ of $\psi$, $\chi$ and with the weight $1$ of $W$. By the definition of $h_k=(h_{k,1},h_{k,2})$ in Section~\ref{sec:notation},
\[
h_{k,2}=e_k^{\varepsilon}\cdot h_{k,1},\qquad h_{k,1}^2=e_k^{-\frac{1}{2}}.
\]

\begin{proposition}[Bounds and regularity of the localised functions]\label{prop:bounds-regularity-localised}
Let $(m,\mathscr{E})\in\mathrm{Ball}^x_k(4)$, $A\in\tilde{R}_k^{-}$, $\mathscr{C}\in\mathfrak{X}_k$ and $\mathsf{Z}\in\mathrm{Box}_k$. Let $m'=m+m(\mathscr{E})$ and $\mathscr{E}^r=R\,\mathscr{E}$ be as in Section~\ref{sec:data-variables-fields}, and let $G^{(i)}$, $E^{(i),*}$, $E^{(\le3),x}$, $G^{(\le3)}$ and $E^{(\le5),x}$ be the functions of Section~\ref{sec:localisation-exponent} formed with these data.
\begin{enumerate}
\item[(a)] For every $X\in D_k$, every $B_1\in\tilde{R}_k$ and every $\mathscr{C}\in\mathfrak{X}_k$,
\[
\left\|\mathscr{E}^r(X;B_1,\mathscr{C})\right\|_{h_k}\le3\,\|\mathscr{E}\|^x_k\cdot\exp\left(-\kappa\cdot d_M(X)\right).
\]
Moreover $|m'|\le8\,w^{\mathrm{m}}_k$.
\item[(b)] For every $\mathbf{Y}\in D_k$,
\[
\left\|E^{(\le3),x}(\mathbf{Y})\right\|_{h_k/2}\le8\,c_X\cdot K_0\cdot e_k^{\frac{1}{4}-5\varepsilon}\cdot\exp\left(-\kappa_1\cdot d_M(\mathbf{Y})\right).
\]
\item[(c)] For every $\mathbf{Y}\in D_k$,
\[
\left\|E^{(\le5),x}(\mathbf{Y})\right\|_{(h_k/4;1)}\le\mathbf{H}_0\cdot\exp\left(-\kappa_2\cdot d_M(\mathbf{Y})\right),\qquad\text{where}\quad\mathbf{H}_0=36\,c_X^2\cdot K_0^2\cdot e_k^{\frac{1}{4}-5\varepsilon}.
\]
\item[(d)] Let $1$ denote the family of parameters all of whose entries are equal to $1$. Then
\[
\sum_{\mathbf{Y}\in D_k}E^{(i),*}(\mathbf{Y})=\sum_{X\in D_k}G^{(i)}(1,X)\qquad(i=1,2,3),
\]
and
\[
\sum_{\mathbf{Y}\in D_k}E^{(\le5),x}(\mathbf{Y})=\sum_{X\in D_k}\left(G^{(\le3)}+G^{(4)}+G^{(5)}\right)(1,X).
\]
\item[(e)] For every $\mathbf{Y}\in D_k$ the elements $E^{(\le3),x}(\mathbf{Y})$ and $E^{(\le5),x}(\mathbf{Y})$ are even. They depend continuously on $(A,\mathscr{C},\mathsf{Z})\in\tilde{R}_k^{-}\times\mathfrak{X}_k\times\mathrm{Box}_k$ and analytically on $(A,\mathscr{C})\in\tilde{R}_k^{-}\times\mathfrak{X}_k$, and they are affine functions of $(m,\mathscr{E})$.
\end{enumerate}
\end{proposition}

\begin{proof}
\emph{(a)} If the polymer $X$ is not small, then $(R\,\mathscr{E})(X)=\mathscr{E}(X)$ by the definition of $R$ in Section~\ref{sec:local-parts-adjustment}. Let $X$ be small, $B_1\in\tilde{R}_k$ and $\mathscr{C}\in\mathfrak{X}_k$. As in Section~\ref{sec:local-parts-adjustment}, $\int_X\mathrm{d}x\,f(x)=\sum_{x\in X}\eta_k^3\cdot f(x)$, and $\operatorname{Vol}(X)=\int_X\mathrm{d}x\,1$. Since $\boldsymbol{\alpha}_0(\mathscr{E},X)\cdot\operatorname{Vol}(X)=\mathscr{E}_{00}(X;0,0)$, the definition of $R$ gives
\[
(R\,\mathscr{E})(X;B_1,\mathscr{C})=\mathscr{E}(X;B_1,\mathscr{C})-\mathscr{E}_{00}(X;0,0)-\int_X\mathrm{d}x\,\bar{\psi}(x)\,\boldsymbol{\alpha}_2(\mathscr{E},X)\,\psi(x),
\]
with the notation $\bar{\psi}(x)\,O\,\psi(x)$ of Section~\ref{sec:localisation-exponent}. The second term on the right is a number, and $\left|\mathscr{E}_{00}(X;0,0)\right|\le\left\|\mathscr{E}(X;0,0)\right\|_{h_k}$, because $\left|\mathscr{E}_{00}(X;0,0)\right|$ is the summand of bidegree $(0,0)$ in the expression of the norm through the kernels (Proposition~\ref{prop:weighted-coefficient-norms}). The third term is the sum, over the sites $x\in X$ and the spinor indices $a$, $b$, of $\eta_k^3\cdot\boldsymbol{\alpha}_2(\mathscr{E},X)_{ab}$ times the product $\bar{\psi}_a(x)\,\psi_b(x)$; these products are pairwise distinct monomials in two generators of the norm $h_{k,1}$ each. Its norm is therefore $h_{k,1}^2\cdot\operatorname{Vol}(X)\cdot\sum_{a,b}\left|\boldsymbol{\alpha}_2(\mathscr{E},X)_{ab}\right|$, and the definition of $\boldsymbol{\alpha}_2(\mathscr{E},X)$ in Section~\ref{sec:local-parts-adjustment} gives
\begin{align*}
h_{k,1}^2\cdot\operatorname{Vol}(X)\cdot\sum_{a,b}\left|\boldsymbol{\alpha}_2(\mathscr{E},X)_{ab}\right|&\le h_{k,1}^2\sum_{a,b}\int\mathrm{d}x\,\mathrm{d}y\,\left|\mathscr{E}_{20}\left(X;0,0;(x,a,1),(y,b,0)\right)\right|\\
&\le\frac{h_{k,1}^2}{2}\cdot\left\|\mathscr{E}_{20}(X;0,0)\right\|\le\left\|\mathscr{E}(X;0,0)\right\|_{h_k}.
\end{align*}
In the middle inequality, $\left\|\mathscr{E}_{20}(X;0,0)\right\|$ is an integral over ordered pairs of arguments; the pair $\left((x,a,1),(y,b,0)\right)$ and its reversal $\left((y,b,0),(x,a,1)\right)$ both enter this integral, and the antisymmetric kernel has equal moduli at the two. The last inequality holds because $\frac{1}{2}h_{k,1}^2\cdot\left\|\mathscr{E}_{20}(X;0,0)\right\|$ is the summand of bidegree $(2,0)$ of $\left\|\mathscr{E}(X;0,0)\right\|_{h_k}$ (Proposition~\ref{prop:weighted-coefficient-norms}). The pairs $(B_1,\mathscr{C})$ and $(0,0)$ lie in $\tilde{R}_k\times\mathfrak{X}_k$ (Definition~\ref{def:domains}). Hence clause~\eqref{it:renormalisation-group-1-k6} gives, for every polymer $X$,
\[
\left\|\mathscr{E}(X;B_1,\mathscr{C})\right\|_{h_k}\le\|\mathscr{E}\|^x_k\cdot\exp\left(-\kappa\cdot d_M(X)\right),\qquad\left\|\mathscr{E}(X;0,0)\right\|_{h_k}\le\|\mathscr{E}\|^x_k\cdot\exp\left(-\kappa\cdot d_M(X)\right).
\]
The first assertion of (a) follows: for small polymers from the three bounds on the terms of $(R\,\mathscr{E})(X;B_1,\mathscr{C})$, and for the other polymers from the first inequality of the last display.

By the definition of $\mathrm{Ball}^x_k(4)$ in Section~\ref{sec:spaces-activities}, $|m|\le4\,w^{\mathrm{m}}_k$ and $\|\mathscr{E}\|^x_k\le4\,w^{\mathrm{E}}_k$. The extraction bound of Proposition~\ref{prop:symmetries-extraction} gives $|m(\mathscr{E})|\le c_{\mathrm{fl}}\cdot e_k^{\frac{1}{2}}\cdot\|\mathscr{E}\|^x_k$. Since $w^{\mathrm{m}}_k=e_k^{\frac{3}{4}-8\varepsilon}$ and $w^{\mathrm{E}}_k=e_k^{\frac{1}{4}-7\varepsilon}$ (Section~\ref{sec:running-quantities-tori}), one has $e_k^{\frac{1}{2}}\cdot w^{\mathrm{E}}_k=e_k^{\varepsilon}\cdot w^{\mathrm{m}}_k$, and hence
\[
|m'|\le4\,w^{\mathrm{m}}_k+4c_{\mathrm{fl}}\cdot e_k^{\frac{1}{2}}\cdot w^{\mathrm{E}}_k=4\left(1+c_{\mathrm{fl}}\cdot e_k^{\varepsilon}\right)\cdot w^{\mathrm{m}}_k\le8\,w^{\mathrm{m}}_k,
\]
the last step by the inequality $c_{\mathrm{fl}}\cdot e_k^{\varepsilon}\le1$ of~\eqref{eq:renormalisation-group-31}.

\emph{(b)} Throughout the proof of (b), $\|\cdot\|$ denotes the norm $\|\cdot\|_{(h_k/2;\mathsf{w}_3)}$ with
\[
\mathsf{w}_3:=C_T\cdot h_{k,1};
\]
that is, the weights $h_k/2$ are given to $\psi$, $\chi$ and the weight $\mathsf{w}_3$ to $\Psi$. For $k=0$, where no generator $\Psi$ occurs, $\|\cdot\|$ is the norm $\|\cdot\|_{h_k/2}$.

\emph{Continuity and analyticity.} By part~\eqref{it:bounds-analyticity-q5} of Lemma~\ref{lem:bounds-analyticity}, by clause~\eqref{it:renormalisation-group-1-k1} and by the analyticity of $M_k(X,\cdot\,)$ (property~\ref{it:renormalisation-group-2}), each $G^{(i)}(s,X)$, $i\le3$, is a continuous function of $(s,A,\mathscr{C},\mathsf{Z})$ on $\mathrm{Pd}_k\times\tilde{R}_k^{-}\times\mathfrak{X}_k\times\mathrm{Box}_k$, analytic in $(s,A,\mathscr{C})$ for $s$ in the interior of $\mathrm{Pd}_k$. For $k\ge1$, Section~\ref{sec:localisation-exponent} defines the function $G^{(1)}$ by a contour integral over the circle $|t|=\rho_1$, and the functions $G^{(2)}$ and $G^{(3)}$ by contour integrals over the circle $|t|=\rho_2$. The algebras $\mathrm{Gr}[\psi,\chi,\Psi]$ and $\mathrm{Gr}[\psi,\chi,W]$ have finitely many generators and are therefore finite-dimensional complex vector spaces; the identity and the bound for $\oint_{\rho}$ of Section~\ref{sec:notation} are used below with $V=\mathrm{Gr}[\psi,\chi,\Psi]$ and the norm $\|\cdot\|_{(h_k/2;\mathsf{w}_3)}$, and in the proof of (c) with $V=\mathrm{Gr}[\psi,\chi,W]$ and the norm $\|\cdot\|_{(h_k/4;1)}$. In each of the three integrals above the integrand is continuous in $t$ on the closed disc bounded by the circle of integration and analytic in $t$ in the corresponding open disc; by Section~\ref{sec:notation}, $G^{(i)}(s,X)$ is therefore the value of its integrand at $t=1$ minus its value at $t=0$ ($i\le3$).

\emph{Bounds.} Let $s=(s',s'')\in\mathrm{Pd}_k$ and $X\in D_k$. For $k=0$ one has $G^{(2)}=G^{(3)}=0$ (Section~\ref{sec:localisation-exponent}), so that the bounds for these two functions obtained below hold trivially.

\emph{The function $G^{(1)}$ for $k\ge1$.} For $|t|=\rho_1$ the bond function $\mathscr{T}+t\cdot\mathscr{Z}(s')$ lies in $\tilde{G}_k$, by part~\eqref{it:bounds-analyticity-q1} of Lemma~\ref{lem:bounds-analyticity}. One has $\mathsf{w}_3^2=C_T^2\cdot e_k^{-\frac{1}{2}}$ and, by Section~\ref{sec:notation}, $\frac{1}{\rho_1-1}\le\frac{2}{\rho_1}=4e_k^{\frac{3}{4}-\frac{9\varepsilon}{2}}$. Part~\eqref{it:bounds-analyticity-q2} of Lemma~\ref{lem:bounds-analyticity} and the bound for contour integrals of Section~\ref{sec:notation} give
\begin{align*}
\left\|G^{(1)}(s,X)\right\|&\le4c_X\cdot c_{\gamma}\cdot C_w\cdot M^3\cdot C_T^2\cdot e_k^{-\frac{1}{2}}\cdot4e_k^{\frac{3}{4}-\frac{9\varepsilon}{2}}\cdot\exp\left(-(\kappa_w-1)\cdot d_M(X)\right)\\
&\le e_k^{\frac{1}{4}-5\varepsilon}\cdot\exp\left(-\kappa\cdot d_M(X)\right);
\end{align*}
the second inequality follows from~\eqref{eq:renormalisation-group-33} and from $\kappa_w-1=\kappa+\kappa_0+2\ge\kappa$, where $\kappa_w=\kappa+\kappa_0+3$ by~\eqref{eq:renormalisation-group-30}.

\emph{The function $G^{(1)}$ for $k=0$.} Let $b$ be a bond. By part~\eqref{it:bounds-analyticity-q1} of Lemma~\ref{lem:bounds-analyticity}, $\left|\mathscr{Z}(s')(b)\right|\le\frac{1}{2}e_k^{-\frac{\varepsilon}{2}}$; hence $e_k\cdot\left|\mathscr{Z}(s')(b)\right|\le\frac{1}{2}e_k^{1-\frac{\varepsilon}{2}}\le\frac{1}{2}$. Since $|\exp(z)-1|\le|z|\cdot\exp(|z|)\le2|z|$ for complex $z$ with $|z|\le\frac{1}{2}$,
\[
\left|\exp\left(\pm i\,e_k\cdot\mathscr{Z}(s')(b)\right)-1\right|\le2e_k\cdot\left|\mathscr{Z}(s')(b)\right|\le e_k^{1-\frac{\varepsilon}{2}}.
\]
The matrix $\gamma_{\nu}$ lies in the group $\mathrm{Sp}$ generated by $\gamma_0,\dots,\gamma_3$ (Section~\ref{sec:dirac-matrices-lattice}); by Proposition~\ref{prop:symmetries-extraction}(b) it is therefore monomial, that is, each of its four rows contains exactly one non-zero entry, and this entry has modulus one. The unit matrix $1$ has the same property. As $P^{\boldsymbol{\sigma}}_{\nu}=\frac{1}{2}\left(1-\boldsymbol{\sigma}\,\gamma_{\nu}\right)$ (Section~\ref{sec:notation}), the sum of the moduli of the sixteen entries of $P^{\boldsymbol{\sigma}}_{\nu}$ is at most $\frac{1}{2}\cdot(4+4)=4$. The right side of~\eqref{eq:renormalisation-group-60} is $\frac{1}{2}$ times the sum of two sums, each of which has $6M^3$ terms, indexed by the triples $(x,\nu,\boldsymbol{\sigma})$ with $x\in\square$; at level $0$ the $M$-cube $\square$ has $M^3$ sites. A generator $\psi$ or $\bar{\psi}$ has the norm $\frac{1}{2}h_{k,1}$, a generator $\chi$ or $\bar{\chi}$ has the norm $\frac{1}{2}h_{k,2}$, and $h_{k,2}\le h_{k,1}$. By~\eqref{eq:renormalisation-group-40}, each of the terms has norm at most
\[
4\cdot\frac{h_{k,1}}{2}\cdot\left(\frac{h_{k,1}}{2}+\frac{h_{k,2}}{2}\right)\cdot e_k^{1-\frac{\varepsilon}{2}}\le2h_{k,1}^2\cdot e_k^{1-\frac{\varepsilon}{2}}.
\]
Taking the factor $\frac{1}{2}$ into account, and using $h_{k,1}^2=e_k^{-\frac{1}{2}}$ and~\eqref{eq:renormalisation-group-33}, one obtains
\[
\left\|G^{(1)}(s,\square)\right\|\le12M^3\cdot e_k^{\frac{1}{2}-\frac{\varepsilon}{2}}\le e_k^{\frac{1}{4}-5\varepsilon}.
\]

\emph{The function $G^{(2)}$ for $k\ge1$.} Since $\rho_2\le\rho_1$, part~\eqref{it:bounds-analyticity-q1} of Lemma~\ref{lem:bounds-analyticity} gives $A+t\cdot\mathscr{Z}(s')\in\tilde{R}_k$ for $|t|\le\rho_2$. By part~\eqref{it:bounds-analyticity-q3} of Lemma~\ref{lem:bounds-analyticity}, applied with $\mathfrak{Z}=\mathscr{Z}(s')$, and by~\eqref{eq:renormalisation-group-34},
\[
\rho_2\cdot\left|J^{\#}_k\left(s'',\mathscr{T},\mathscr{Z}(s')\right)\right|_{1,\infty}\cdot\mathsf{w}_3\le16C_w\cdot C_T\cdot e_k^{\frac{5\varepsilon}{2}}\cdot h_{k,1}\le\frac{1}{2}h_{k,2}.
\]
For $|t|=\rho_2$ the integrand of $G^{(2)}(s,X)$ is the image of $\mathscr{E}^r\left(X;A+t\cdot\mathscr{Z}(s'),\mathscr{C}\right)\in\mathrm{Gr}[\psi,\chi]$ under the homomorphism that replaces $\psi^{\#}$ by $\psi^{\#}+t\cdot J^{\#}_k\left(s'',\mathscr{T},\mathscr{Z}(s')\right)\Psi$; in the notation of~\eqref{eq:renormalisation-group-41}, the operator from $\psi^{\#}$ to $\psi^{\#}$ is the identity, and the operator from $\Psi$ to $\psi^{\#}$ is $t\cdot J^{\#}_k\left(s'',\mathscr{T},\mathscr{Z}(s')\right)$. Since
\[
\frac{1}{2}h_{k,1}+\frac{1}{2}h_{k,2}\le h_{k,1},\qquad\frac{1}{2}h_{k,2}+\frac{1}{2}h_{k,2}=h_{k,2},
\]
the inequality~\eqref{eq:renormalisation-group-41} bounds the norm of this integrand by $\left\|\mathscr{E}^r\left(X;A+t\cdot\mathscr{Z}(s'),\mathscr{C}\right)\right\|_{h_k}$. By (a), with $\|\mathscr{E}\|^x_k\le4\,w^{\mathrm{E}}_k=4e_k^{\frac{1}{4}-7\varepsilon}$, and since $\frac{1}{\rho_2-1}\le\frac{2}{\rho_2}=2e_k^{\frac{3}{4}-7\varepsilon}$,
\[
\left\|G^{(2)}(s,X)\right\|\le2e_k^{\frac{3}{4}-7\varepsilon}\cdot12e_k^{\frac{1}{4}-7\varepsilon}\cdot\exp\left(-\kappa\cdot d_M(X)\right)=24e_k^{1-14\varepsilon}\cdot\exp\left(-\kappa\cdot d_M(X)\right).
\]

\emph{The function $G^{(3)}$ for $k\ge1$.} The $M$-cube $\square$ contains $M^3\cdot\eta_k^{-3}$ sites. The pairing $\langle\bar{\psi},\psi\rangle_{\square}$ is the sum, over these sites $x$ and over the four values of the spinor index $\boldsymbol{\beta}$, of the monomials $\eta_k^3\cdot\bar{\psi}_{\boldsymbol{\beta}}(x)\,\psi_{\boldsymbol{\beta}}(x)$, each of which has the norm $\eta_k^3\cdot h_{k,1}^2$ for the weights $h_k$; hence
\[
\left\|\langle\bar{\psi},\psi\rangle_{\square}\right\|_{h_k}=4M^3\cdot h_{k,1}^2.
\]
For $|t|=\rho_2$ the integrand of $G^{(3)}(s,\square)$ is the image of $-m'\cdot\langle\bar{\psi},\psi\rangle_{\square}$ under the homomorphism that replaces $\psi$ by $\psi+t\cdot J\,\Psi$ and $\bar{\psi}$ by $\bar{\psi}+t\cdot\bar{J}\,\bar{\Psi}$. The rows of $J$ and of $\bar{J}$ are rows of $J^{\#}_k\left(s'',\mathscr{T},\mathscr{Z}(s')\right)$, so that $\rho_2\cdot|J|_{1,\infty}\cdot\mathsf{w}_3$ and $\rho_2\cdot|\bar{J}|_{1,\infty}\cdot\mathsf{w}_3$ are at most $\frac{1}{2}h_{k,2}$. By~\eqref{eq:renormalisation-group-41}, applied with $\frac{1}{2}h_{k,1}+\frac{1}{2}h_{k,2}\le h_{k,1}$, the norm of the integrand is at most $|m'|\cdot\left\|\langle\bar{\psi},\psi\rangle_{\square}\right\|_{h_k}=|m'|\cdot4M^3\cdot h_{k,1}^2$. With $|m'|\le8\,w^{\mathrm{m}}_k=8e_k^{\frac{3}{4}-8\varepsilon}$ from (a), with $h_{k,1}^2=e_k^{-\frac{1}{2}}$ and with $\frac{1}{\rho_2-1}\le2e_k^{\frac{3}{4}-7\varepsilon}$,
\[
\left\|G^{(3)}(s,\square)\right\|\le8e_k^{\frac{3}{4}-8\varepsilon}\cdot4M^3\cdot e_k^{-\frac{1}{2}}\cdot2e_k^{\frac{3}{4}-7\varepsilon}=64M^3\cdot e_k^{1-15\varepsilon}.
\]

\emph{The sum of the three bounds.} The functions $G^{(3)}(s,\cdot\,)$ and, for $k=0$, $G^{(1)}(s,\cdot\,)$ vanish on the polymers that are not single $M$-cubes (Section~\ref{sec:localisation-exponent}), and $d_M(\square)=0$ for every $M$-cube $\square$ by~\eqref{eq:renormalisation-group-44}. Hence the bounds obtained for these functions remain valid when their right sides are multiplied by $\exp\left(-\kappa\cdot d_M(X)\right)$. By~\eqref{eq:renormalisation-group-35},
\[
24e_k^{\frac{3}{4}-9\varepsilon}+64M^3\cdot e_k^{\frac{3}{4}-10\varepsilon}\le88M^3\cdot e_k^{\frac{1}{4}-6\varepsilon}\le1,
\]
because $e_k<1$, $M\ge1$, and both exponents on the left are at least $\frac{1}{4}-6\varepsilon$. As $24e_k^{1-14\varepsilon}=24e_k^{\frac{3}{4}-9\varepsilon}\cdot e_k^{\frac{1}{4}-5\varepsilon}$ and $64M^3\cdot e_k^{1-15\varepsilon}=64M^3\cdot e_k^{\frac{3}{4}-10\varepsilon}\cdot e_k^{\frac{1}{4}-5\varepsilon}$, the bounds for $G^{(1)}$, $G^{(2)}$ and $G^{(3)}$ add up to
\[
\sum_{i=1}^3\left\|G^{(i)}(s,X)\right\|\le2e_k^{\frac{1}{4}-5\varepsilon}\cdot\exp\left(-\kappa\cdot d_M(X)\right)\qquad\left(s\in\mathrm{Pd}_k,\ X\in D_k\right).
\]

\emph{Resummation.} In the first localisation the parameter cubes are those of $\pi_k$, $\pi'_k$ and $\pi''_k$. For $X\in D_k$ let
\[
f^X(s):=\sum_{i\le3}G^{(i)}(s,X)\qquad\left(s\in\mathrm{Pd}_k\right),
\]
a map into $\mathrm{Gr}[\psi,\chi,\Psi]$, which carries the norm $\|\cdot\|$. By the continuity and analyticity shown above and by the last bound, $f^X$ is a map as in Lemma~\ref{lem:expansion-parameters}, with $\rho=M^{\alpha_0}$, $P_0=P_0(X)$ and $B_f=2e_k^{\frac{1}{4}-5\varepsilon}\cdot\exp\left(-\kappa\cdot d_M(X)\right)$. If $Y'(X,P)$ is not a polymer, that is, not connected, then by Lemma~\ref{lem:renormalisation-group-1}(b) for the first localisation the derivative $\frac{\partial}{\partial s_P}f^X(s)$ vanishes at every $s\in\mathrm{Pd}_k$ that vanishes outside $P_0(X)\cup P$; hence $f^X_P=0$. Since $\kappa_0\le\kappa$ (stage~\ref{it:order-choice-3}), the hypotheses of Lemma~\ref{lem:resummation-polymers} hold with $\kappa'=\kappa$ and $B_0=2e_k^{\frac{1}{4}-5\varepsilon}$. By~\eqref{eq:renormalisation-group-57} and~\eqref{eq:renormalisation-group-59}, the function $F$ of Lemma~\ref{lem:resummation-polymers} is $F(\mathbf{Y})=\sum_{i\le3}E^{(i),*}(\mathbf{Y})$, and $\kappa-\kappa_0-1=\kappa_1$ (Section~\ref{sec:notation}). Lemma~\ref{lem:resummation-polymers} therefore gives, for every $\mathbf{Y}\in D_k$,
\[
\left\|\sum_{i\le3}E^{(i),*}(\mathbf{Y})\right\|_{(h_k/2;\,C_T\cdot h_{k,1})}\le8c_X\cdot K_0\cdot e_k^{\frac{1}{4}-5\varepsilon}\cdot\exp\left(-\kappa_1\cdot d_M(\mathbf{Y})\right).
\]

\emph{The substitution.} For $k=0$ there is no substitution~\eqref{eq:renormalisation-group-58}, the functions $G^{(2)}$ and $G^{(3)}$ vanish, and $E^{(\le3),x}=\mathrm{Loc}\left[G^{(1)}\right]=E^{(1),*}=\sum_{i\le3}E^{(i),*}$ by~\eqref{eq:renormalisation-group-60}; in this case the last bound is (b). Let $k\ge1$. Then $E^{(\le3),x}(\mathbf{Y})$ is the image of $\sum_{i\le3}E^{(i),*}(\mathbf{Y})$ under the substitution~\eqref{eq:renormalisation-group-58}, formed with the polymer $\mathbf{Y}$ and with $B=\mathscr{T}$. Since $\tilde{R}_k^{-}\subset\tilde{R}_k$ (Definition~\ref{def:domains}), the field $\mathscr{T}=A+\mathscr{C}$ lies in $\tilde{R}_k+\mathfrak{X}_k$, and Proposition~\ref{prop:left-inverse} gives $\left|T^{(i)}(\mathscr{T})\right|_{1,\infty}\le C_T$ for $i=1,2$. On the barred fields the operators are $\bar{T}^{(i)}(\mathscr{T})=\mathsf{C}^{-1}\cdot T^{(i)}(-\mathscr{T})\cdot\mathsf{C}$; here $\mathsf{C}$ is a monomial matrix whose non-zero entries have modulus one (Proposition~\ref{prop:symmetries-extraction}(b)), and $-\mathscr{T}$ lies in $\tilde{R}_k+\mathfrak{X}_k$ (Lemma~\ref{lem:symmetries-domains}), so that $\left|\bar{T}^{(i)}(\mathscr{T})\right|_{1,\infty}\le C_T$ as well. The substitution leaves $\psi$, $\chi$, $\bar{\psi}$, $\bar{\chi}$ unchanged, and it replaces the generators $\Psi$, $\bar{\Psi}$, of the weight $\mathsf{w}_3=C_T\cdot h_{k,1}$, by combinations of $\psi$, $\chi$, resp.\ of $\bar{\psi}$, $\bar{\chi}$, of the weights $\frac{1}{2}h_{k,1}$, $\frac{1}{2}h_{k,2}$. Since
\[
C_T\cdot\frac{h_{k,1}}{2}+C_T\cdot\frac{h_{k,2}}{2}\le C_T\cdot h_{k,1},
\]
the inequality~\eqref{eq:renormalisation-group-41} yields
\[
\left\|E^{(\le3),x}(\mathbf{Y})\right\|_{h_k/2}\le\left\|\sum_{i\le3}E^{(i),*}(\mathbf{Y})\right\|_{(h_k/2;\,C_T\cdot h_{k,1})},
\]
and (b) follows from the preceding bound.

\emph{(c)} By~\eqref{eq:renormalisation-group-36}, $C_W=4C_w^2\cdot c_{\gamma}\le\frac{3}{4}e_k^{-2\varepsilon}$, and by~\eqref{eq:renormalisation-group-35}, $3e_k^{\frac{1}{4}-3\varepsilon}\le1$. As $h_{k,2}=e_k^{-\frac{1}{4}+\varepsilon}$ and $\rho_4=e_k^{-\frac{1}{4}+3\varepsilon}$, it follows that
\[
C_W\le\frac{1}{4}h_{k,2},\qquad\rho_4\cdot C_W\le\frac{3}{4}h_{k,2}.
\]
Let $s''\in\mathrm{Pd}_k$ and $X\in D_k$. For $k=0$ one has $G^{(4)}=0$ (Section~\ref{sec:localisation-exponent}), so that the bound for $G^{(4)}$ obtained below holds trivially. By~\eqref{eq:renormalisation-group-59}, the element $G^{(\le3)}(s'',X)$, the integrand of $G^{(4)}(s'',X)$ for $k\ge1$, and the integrand of $G^{(5)}(s'',\square)$ are the images of $E^{(\le3),x}(X)$, of $\mathscr{E}^r(X;A,\mathscr{C})$ and of $-m'\cdot\langle\bar{\psi},\psi\rangle_{\square}$, respectively, under the homomorphism that replaces $\psi^{\#}$ by $\psi^{\#}+t\cdot W^{\#}_k(s'',\mathscr{T})$, with $t=1$ in the first case and $|t|=\rho_4$ in the other two. Here $W^{\#}_k(s'',\mathscr{T})=H^{\#}_k(s'',\mathscr{T})\cdot\tilde{\Gamma}_k(s'',\mathscr{T})\cdot W$ by~\eqref{eq:renormalisation-group-56}, and
\[
\left|H^{\#}_k(s'',\mathscr{T})\cdot\tilde{\Gamma}_k(s'',\mathscr{T})\right|_{1,\infty}\le C_W
\]
by part~\eqref{it:bounds-analyticity-q4} of Lemma~\ref{lem:bounds-analyticity}. For $i=1,2$ one has, because $h_{k,2}\le h_{k,1}$,
\[
\frac{1}{4}h_{k,i}+\frac{1}{4}h_{k,2}\le\frac{1}{2}h_{k,i},\qquad\frac{1}{4}h_{k,i}+\frac{3}{4}h_{k,2}\le h_{k,i}.
\]
The inequality~\eqref{eq:renormalisation-group-41}, applied with the first of these inequalities for $t=1$ and with the second for $|t|=\rho_4$, together with (a), with the identity $\left\|\langle\bar{\psi},\psi\rangle_{\square}\right\|_{h_k}=4M^3\cdot h_{k,1}^2$ and with $\frac{1}{\rho_4-1}\le\frac{2}{\rho_4}=2e_k^{\frac{1}{4}-3\varepsilon}$, gives
\begin{align*}
\left\|G^{(\le3)}(s'',X)\right\|_{(h_k/4;1)}&\le\left\|E^{(\le3),x}(X)\right\|_{h_k/2},\\
\left\|G^{(4)}(s'',X)\right\|_{(h_k/4;1)}&\le2e_k^{\frac{1}{4}-3\varepsilon}\cdot12e_k^{\frac{1}{4}-7\varepsilon}\cdot\exp\left(-\kappa\cdot d_M(X)\right),\\
\left\|G^{(5)}(s'',\square)\right\|_{(h_k/4;1)}&\le8e_k^{\frac{3}{4}-8\varepsilon}\cdot4M^3\cdot e_k^{-\frac{1}{2}}\cdot2e_k^{\frac{1}{4}-3\varepsilon}.
\end{align*}
By~\eqref{eq:renormalisation-group-35},
\[
24e_k^{\frac{1}{2}-10\varepsilon}+64M^3\cdot e_k^{\frac{1}{2}-11\varepsilon}\le88M^3\cdot e_k^{\frac{1}{4}-6\varepsilon}\cdot e_k^{\frac{1}{4}-5\varepsilon}\le e_k^{\frac{1}{4}-5\varepsilon}.
\]
The function $G^{(5)}(s'',\cdot\,)$ vanishes on the polymers that are not single $M$-cubes, and $d_M(\square)=0$ by~\eqref{eq:renormalisation-group-44}; so the bound for $G^{(5)}$ remains valid when its right side is multiplied by $\exp\left(-\kappa\cdot d_M(X)\right)$. Let
\[
f^X(s''):=G^{(\le3)}(s'',X)+G^{(4)}(s'',X)+G^{(5)}(s'',X)\qquad\left(s''\in\mathrm{Pd}_k\right).
\]
By (b) for $\mathbf{Y}=X$, by the bounds for $G^{(\le3)}$, $G^{(4)}$, $G^{(5)}$ obtained above together with the consequence of~\eqref{eq:renormalisation-group-35} just displayed, and since $\kappa_1\le\kappa$, $c_X\ge1$ and $K_0\ge1$,
\[
\left\|f^X(s'')\right\|_{(h_k/4;1)}\le\left(8c_X\cdot K_0+1\right)\cdot e_k^{\frac{1}{4}-5\varepsilon}\cdot\exp\left(-\kappa_1\cdot d_M(X)\right)\le9c_X\cdot K_0\cdot e_k^{\frac{1}{4}-5\varepsilon}\cdot\exp\left(-\kappa_1\cdot d_M(X)\right).
\]
The map $f^X$ is continuous on $\mathrm{Pd}_k$ and analytic in its interior, by part~\eqref{it:bounds-analyticity-q5} of Lemma~\ref{lem:bounds-analyticity}. In the second localisation the parameter cubes are those of $\pi_k$. If $Y'(X,P)$ is not connected, then by Lemma~\ref{lem:renormalisation-group-1}(b) for the second localisation the derivative $\frac{\partial}{\partial s_P}f^X(s'')$ vanishes at every $s''$ that vanishes outside $P_0(X)\cup P$, so that $f^X_P=0$. Hence the hypotheses of Lemma~\ref{lem:resummation-polymers} hold with
\[
\kappa'=\kappa_1,\qquad B_0=9c_X\cdot K_0\cdot e_k^{\frac{1}{4}-5\varepsilon};
\]
indeed $\kappa_0\le\kappa_1\le\kappa$, because $\kappa\ge7\kappa_0+7$ by stage~\ref{it:order-choice-3}. By~\eqref{eq:renormalisation-group-57} and~\eqref{eq:renormalisation-group-59} the function $F$ of Lemma~\ref{lem:resummation-polymers} is $E^{(\le5),x}$, and $\kappa_1-\kappa_0-1=\kappa_2$ (Section~\ref{sec:notation}). Lemma~\ref{lem:resummation-polymers} therefore gives, for every $\mathbf{Y}\in D_k$,
\[
\left\|E^{(\le5),x}(\mathbf{Y})\right\|_{(h_k/4;1)}\le4c_X\cdot K_0\cdot9c_X\cdot K_0\cdot e_k^{\frac{1}{4}-5\varepsilon}\cdot\exp\left(-\kappa_2\cdot d_M(\mathbf{Y})\right)=\mathbf{H}_0\cdot\exp\left(-\kappa_2\cdot d_M(\mathbf{Y})\right).
\]

\emph{(d)} Let $i\in\{1,2,3\}$ and $X\in D_k$. By the proof of (b), the map $G^{(i)}(\cdot\,,X)$ is continuous on $\mathrm{Pd}_k$, analytic in its interior, and its norm is at most $2e_k^{\frac{1}{4}-5\varepsilon}\cdot\exp\left(-\kappa\cdot d_M(X)\right)$ on $\mathrm{Pd}_k$. Together with Lemma~\ref{lem:renormalisation-group-1}(b), this shows that the maps $G^{(i)}(\cdot\,,X)$, $X\in D_k$, satisfy the hypotheses of Lemma~\ref{lem:resummation-polymers} for the first localisation, with $\kappa'=\kappa$ and $B_0=2e_k^{\frac{1}{4}-5\varepsilon}$; by~\eqref{eq:renormalisation-group-57} and~\eqref{eq:renormalisation-group-59} the function $F$ of that lemma is $E^{(i),*}$. The first identity of (d) is therefore the identity of Lemma~\ref{lem:resummation-polymers}. The second identity of (d) is the identity of the same lemma for the maps $f^X$ of the proof of (c).

\emph{(e)} \emph{Evenness.} The element $\mathscr{E}^r(X)$ is even, because $R$ maps $K^{x,C}_k$ into itself (Proposition~\ref{prop:symmetries-extraction}). In the displays~\eqref{eq:renormalisation-group-59} and~\eqref{eq:renormalisation-group-60} the barred and the unbarred fields are coupled two by two, each barred field with an unbarred one. Consequently $E^{(\le3),x}(\mathbf{Y})$ and $E^{(\le5),x}(\mathbf{Y})$ are even.

\emph{Continuity and analyticity.} Let $G(s,X)$ be one of the functions to which $\mathrm{Loc}$ is applied in~\eqref{eq:renormalisation-group-59}, and let $P$ be a set of parameter cubes outside $P_0(X)$. Let $s^0$ be a point with $s^0_c\in[0,1]$ for $c\in P$, the other parameters being fixed at $0$ or $1$, and take a polycircle about $s^0$ in the variables $s_c$, $c\in P$, that is contained in $\mathrm{Pd}_k$. On this polycircle Cauchy's integral formula writes $\frac{\partial}{\partial s_P}G(s^0,X)$ as the integral of $G(\cdot\,,X)$ multiplied by a bounded continuous weight. Integrating over $s^0_P\in[0,1]^P$, one sees that the operation $\int_{[0,1]^P}\mathrm{d}s_P\,\frac{\partial}{\partial s_P}$ of~\eqref{eq:renormalisation-group-57}, which coincides with the operation $f\mapsto f_P$ of Lemma~\ref{lem:expansion-parameters}, amounts to integrating $G(\cdot\,,X)$ against bounded continuous weights. An integral of this type is continuous in $(A,\mathscr{C},\mathsf{Z})$ and analytic in $(A,\mathscr{C})$ whenever its integrand $G$ is. The functions $G^{(i)}$, $i\le3$, have these two properties by the proof of (b). The entries of $T^{(i)}(\mathscr{T})$ and of $\bar{T}^{(i)}(\mathscr{T})$ are entire functions of $\mathscr{T}=A+\mathscr{C}$ (Lemma~\ref{lem:interior-boundary-parts}(b)); hence $E^{(\le3),x}(\mathbf{Y})$ is continuous in $(A,\mathscr{C},\mathsf{Z})$ and analytic in $(A,\mathscr{C})$. By part~\eqref{it:bounds-analyticity-q5} of Lemma~\ref{lem:bounds-analyticity} and clause~\eqref{it:renormalisation-group-1-k1}, the functions $G^{(\le3)}$, $G^{(4)}$ and $G^{(5)}$ then have the two properties, and so has $E^{(\le5),x}(\mathbf{Y})$.

\emph{Dependence on $(m,\mathscr{E})$.} Neither $m$ nor $\mathscr{E}$ enters the definition of $G^{(1)}$. In $G^{(2)}$ and $G^{(4)}$ the pair $(m,\mathscr{E})$ occurs only through $\mathscr{E}^r=R\,\mathscr{E}$, on which both functions depend linearly, and $R\,\mathscr{E}$ depends linearly on $\mathscr{E}$. In $G^{(3)}$ and $G^{(5)}$ it occurs only through the factor $m'=m+m(\mathscr{E})$, on which both functions depend linearly; as the map $\mathscr{E}\mapsto m(\mathscr{E})$ is linear (Proposition~\ref{prop:symmetries-extraction}), $m'$ is an affine function of $(m,\mathscr{E})$. The operation $\mathrm{Loc}$, the substitution~\eqref{eq:renormalisation-group-58}, and the substitution by which $G^{(\le3)}$ is formed from $E^{(\le3),x}$ are linear. Hence $E^{(\le3),x}(\mathbf{Y})$ and $E^{(\le5),x}(\mathbf{Y})$ are affine in $(m,\mathscr{E})$.
\end{proof}

\subsection{The logarithm of the fluctuation integral}\label{sec:logarithm-fluctuation-integral}

\begin{lemma}[Integration of a polymer exponential]\label{lem:integration-polymer-exponential}
Assume~\eqref{eq:renormalisation-group-37}, and let $\mathbf{H}_0$ be the number defined there. Let $\mathsf{w}=(\mathsf{w}_1,\mathsf{w}_2)$ with $\mathsf{w}_1,\mathsf{w}_2>0$. Denote by $\|\cdot\|_{\mathsf{w}}$ the norm on $\mathrm{Gr}[\psi,\chi]$ formed with the weight $\mathsf{w}_1$ of $\psi$ and the weight $\mathsf{w}_2$ of $\chi$, and by $\|\cdot\|_{(\mathsf{w};1)}$ the norm on $\mathrm{Gr}[\psi,\chi,W]$ formed with the same weights of $\psi$, $\chi$ and with the weight $1$ of $W$ (Section~\ref{sec:recalled-auxiliary-results}). Let $\mathrm{Box}_k$ be the cube and $\hat{\mu}$ the probability measure on $\mathrm{Box}_k$ of Section~\ref{sec:averaging-fluctuations-gauge}, and let $\mathsf{Z}(\Upsilon)$ denote the coordinates of $\mathsf{Z}\in\mathrm{Box}_k$ with respect to the basis $\{\Upsilon\}$ of Proposition~\ref{prop:basis-map-box}. For every $X\in D_k$ let $V(X)$ be an even element of $\mathrm{Gr}[\psi,\chi,W]$ that depends continuously on $\mathsf{Z}\in\mathrm{Box}_k$, satisfies
\[
\left\|V(X)\right\|_{(\mathsf{w};1)}\le\mathbf{H}_0\cdot\exp\left(-\kappa_2\cdot d_M(X)\right)\qquad\left(\mathsf{Z}\in\mathrm{Box}_k\right),
\]
with the number $\kappa_2$ of Section~\ref{sec:notation}, contains no generator $W(u)$ whose unit site lies outside $X$, and depends on $\mathsf{Z}$ only through the coordinates $\mathsf{Z}(\Upsilon)$ of the basis elements $\Upsilon$ that are contained in $X$. Let $\mathsf{K}$, $\mathsf{H}$ and $E^{\#}$ be formed from the family $V=\left(V(X)\right)_{X\in D_k}$ as in~\eqref{eq:renormalisation-group-61}, with $V$ in place of $E^{(\le5),x}$: for a union $X$ of $M$-cubes,
\[
\mathsf{K}(X):=\sum_{\mathcal{X}}\prod_{X'\in\mathcal{X}}\left(\exp\left(V(X')\right)-1\right),\qquad\mathsf{H}(X):=\int\mathrm{d}\mu_I(W)\,\mathrm{d}\hat{\mu}(\mathsf{Z})\,\mathsf{K}(X),
\]
where $\mathcal{X}$ ranges over the sets of distinct polymers with union $X$, and $\mathsf{K}(\emptyset):=1$; and, for $X\in D_k$,
\[
E^{\#}(X):=\sum_{n\ge1}\frac{1}{n!}\sum_{\substack{(X_1,\dots,X_n)\in D_k^n\\X_1\cup\dots\cup X_n=X}}\rho^T(X_1,\dots,X_n)\cdot\prod_{i=1}^n\mathsf{H}(X_i).
\]
Then for every $X\in D_k$ the series that defines $E^{\#}(X)$ converges absolutely in the norm $\|\cdot\|_{\mathsf{w}}$, and
\begin{gather*}
\int\mathrm{d}\mu_I(W)\,\mathrm{d}\hat{\mu}(\mathsf{Z})\,\exp\left\{\sum_{X\in D_k}V(X)\right\}=\exp\left\{\sum_{X\in D_k}E^{\#}(X)\right\},\\
\left\|E^{\#}(X)\right\|_{\mathsf{w}}\le c'_S\cdot\mathbf{H}_0\cdot\exp\left(-\kappa_3\cdot d_M(X)\right)\quad\left(X\in D_k\right),\qquad c'_S:=14\,c_X^2\cdot K_0^2\cdot\exp(3\kappa),
\end{gather*}
with the number $\kappa_3$ of Section~\ref{sec:notation}. For every integer $\mathsf{n}\ge0$ the norms $\|\cdot\|_{\mathsf{w}}$ of the terms with $n>\mathsf{n}$ of the series for $E^{\#}(X)$ add up to at most $2^{-\mathsf{n}}\cdot c'_S\cdot\mathbf{H}_0\cdot\exp\left(-\kappa_3\cdot d_M(X)\right)$. For an integer $\mathsf{n}\ge1$ let $E^{\#}_{\mathsf{n}}(X)$ be the sum of the terms with $n\le\mathsf{n}$ of the series for $E^{\#}(X)$, in which $\exp\left(V(X')\right)-1$ is replaced, in the definition of $\mathsf{K}$, by the partial sum
\[
\mathrm{pl}_{\mathsf{n}}\left(V(X')\right):=\sum_{j=1}^{\mathsf{n}}\frac{1}{j!}\,V(X')^j
\]
of the exponential series. Then $E^{\#}_{\mathsf{n}}(X)$ is a polynomial in integrals $\int\mathrm{d}\mu_I(W)\,\mathrm{d}\hat{\mu}(\mathsf{Z})$ of products of the elements $V(X')$, $X'\in D_k$, $X'\subseteq X$; and $E^{\#}_{\mathsf{n}}(X)$ converges to $E^{\#}(X)$ in the norm $\|\cdot\|_{\mathsf{w}}$ as $\mathsf{n}\to\infty$, uniformly on every set of families $V$ that satisfy the hypotheses above.
\end{lemma}

\begin{proof}
Put
\begin{gather*}
\boldsymbol{\beta}_1:=\kappa_2-\kappa_0=\kappa-3\kappa_0-2,\qquad\boldsymbol{\beta}_2:=\kappa_4-2\kappa_0=\kappa-5\kappa_0-3,\\
\mathbf{H}_1:=7\,c_X\cdot K_0\cdot\exp(3\boldsymbol{\beta}_1)\cdot\mathbf{H}_0,\qquad z_S:=28\,\mathbf{H}_1\cdot\exp(\kappa_0)\cdot K_0\cdot\exp(3\boldsymbol{\beta}_2);
\end{gather*}
the equalities in the first line follow from the definitions of $\kappa_2$ and $\kappa_4$ in Section~\ref{sec:notation}. Since $\kappa\ge7\kappa_0+7$ by stage~\ref{it:order-choice-3}, one has $\boldsymbol{\beta}_1\ge\boldsymbol{\beta}_2>0$. The condition~\eqref{eq:renormalisation-group-37} gives
\begin{multline}\label{eq:renormalisation-group-63}
\mathbf{H}_0\le1,\qquad4\,\mathbf{H}_0\cdot K_0\cdot c_X\cdot\exp(3\boldsymbol{\beta}_1)\le144\,c_X^3\cdot K_0^3\cdot\exp(3\kappa)\cdot e_k^{\frac{1}{4}-5\varepsilon}\le1,\\
2z_S=14112\,c_X^3\cdot K_0^4\cdot\exp\left(3\boldsymbol{\beta}_1+3\boldsymbol{\beta}_2+\kappa_0\right)\cdot e_k^{\frac{1}{4}-5\varepsilon}\\
\le14112\,c_X^3\cdot K_0^4\cdot\exp(6\kappa)\cdot e_k^{\frac{1}{4}-5\varepsilon}\le1.
\end{multline}
Indeed, $\mathbf{H}_0=36\,c_X^2\cdot K_0^2\cdot e_k^{\frac{1}{4}-5\varepsilon}$, $\boldsymbol{\beta}_1\le\kappa$ and $3\boldsymbol{\beta}_1+3\boldsymbol{\beta}_2+\kappa_0=6\kappa-23\kappa_0-15\le6\kappa$; and the bounds by $1$ follow from the inequality $c_S\cdot e_k^{\frac{1}{4}-5\varepsilon}\le1$ of~\eqref{eq:renormalisation-group-37}, because $c_S=10^5\,c_X^5\cdot K_0^5\cdot\exp(6\kappa)$ is at least $144\,c_X^3\cdot K_0^3\cdot\exp(3\kappa)$ and at least $14112\,c_X^3\cdot K_0^4\cdot\exp(6\kappa)$.

For polymers $X_1,\dots,X_n$, the \emph{incompatibility graph} of $(X_1,\dots,X_n)$ is the graph with the vertex set $\{1,\dots,n\}$ in which two vertices $i\ne j$ are joined whenever $X_i$ and $X_j$ are incompatible, that is, whenever $X_i\cup X_j$ is connected through faces. If this graph is connected, then $X_1\cup\dots\cup X_n$ is a polymer and
\[
d_M\left(X_1\cup\dots\cup X_n\right)\le\sum_{i=1}^nd_M(X_i)+3\,(n-1),
\]
by~\eqref{eq:renormalisation-group-48}.

\emph{Step~1 (Mayer expansion).} Fix $\mathsf{Z}\in\mathrm{Box}_k$. Even elements of $\mathrm{Gr}[\psi,\chi,W]$ commute with each other; hence
\[
\exp\left\{\sum_{X\in D_k}V(X)\right\}=\prod_{X\in D_k}\left(1+\left(\exp\left(V(X)\right)-1\right)\right)=\sum_X\mathsf{K}(X),
\]
where $X$ ranges over the unions of $M$-cubes, the empty union included; the last equality is obtained by expanding the finite product and collecting the terms according to the union of the polymers that occur in them. Let $X$ be a union of $M$-cubes, and let $X_1,\dots,X_r$ be its components. A polymer that belongs to a set of polymers whose union is $X$ is connected through faces and contained in $X$, so that it lies in one component of $X$. Consequently the sets of distinct polymers whose union is $X$ are exactly the sets $\mathcal{X}_1\cup\dots\cup\mathcal{X}_r$ in which, for $a=1,\dots,r$, the polymers of $\mathcal{X}_a$ are distinct and have the union $X_a$; and
\[
\mathsf{K}(X)=\prod_{a=1}^r\mathsf{K}(X_a).
\]

Let $Y$ be a polymer, and let $X_1,\dots,X_n$ be $n$ distinct polymers whose union is $Y$. The incompatibility graph of $(X_1,\dots,X_n)$ is connected. Otherwise $\{1,\dots,n\}$ is the union of two disjoint non-empty sets $J_1$, $J_2$ such that no edge joins a vertex of $J_1$ to a vertex of $J_2$. The union $U_1$ of the $X_i$ with $i\in J_1$ and the union $U_2$ of the $X_j$ with $j\in J_2$ cover the polymer $Y$, which is connected through faces; so some $M$-cube $\square_1\subseteq U_1$ and some $M$-cube $\square_2\subseteq U_2$ are equal or share a face. Then a polymer $X_i\supseteq\square_1$ with $i\in J_1$ and a polymer $X_j\supseteq\square_2$ with $j\in J_2$ are incompatible, which is impossible. Hence, by~\eqref{eq:renormalisation-group-48},
\[
d_M(Y)\le\sum_{i=1}^nd_M(X_i)+3\,(n-1).
\]
By~\eqref{eq:renormalisation-group-40}, by the bound $\|V(X)\|_{(\mathsf{w};1)}\le\mathbf{H}_0\le1$, which follows from the hypothesis and~\eqref{eq:renormalisation-group-63}, and by the inequality $\exp(r)-1\le2r$ for $0\le r\le1$,
\[
\left\|\exp\left(V(X)\right)-1\right\|_{(\mathsf{w};1)}\le\exp\left(\|V(X)\|_{(\mathsf{w};1)}\right)-1\le2\,\|V(X)\|_{(\mathsf{w};1)}\qquad\left(X\in D_k\right).
\]
With the hypothesis on $\|V(X)\|_{(\mathsf{w};1)}$, with $\kappa_2=\boldsymbol{\beta}_1+\kappa_0$ and $\boldsymbol{\beta}_1>0$, and with the bound on $d_M(Y)$ just proved, it follows that
\[
\prod_{i=1}^n\left\|\exp\left(V(X_i)\right)-1\right\|_{(\mathsf{w};1)}\le\exp\left(-\boldsymbol{\beta}_1\cdot d_M(Y)\right)\cdot\left(2\,\mathbf{H}_0\cdot\exp(3\boldsymbol{\beta}_1)\right)^n\cdot\prod_{i=1}^n\exp\left(-\kappa_0\cdot d_M(X_i)\right).
\]
A set of $n$ distinct polymers is the set of the entries of exactly $n!$ ordered $n$-tuples. Drop the condition that the union is $Y$, and keep the condition that every entry is contained in $Y$. By~\eqref{eq:renormalisation-group-40} and~\eqref{eq:renormalisation-group-46},
\begin{multline*}
\left\|\mathsf{K}(Y)\right\|_{(\mathsf{w};1)}\le\exp\left(-\boldsymbol{\beta}_1\cdot d_M(Y)\right)\cdot\sum_{n\ge1}\frac{1}{n!}\left(2\,\mathbf{H}_0\cdot\exp(3\boldsymbol{\beta}_1)\cdot\sum_{\substack{X\in D_k\\X\subseteq Y}}\exp\left(-\kappa_0\cdot d_M(X)\right)\right)^n\\
\le\exp\left(-\boldsymbol{\beta}_1\cdot d_M(Y)\right)\cdot x\cdot\exp(x),\qquad x:=2\,\mathbf{H}_0\cdot K_0\cdot\exp(3\boldsymbol{\beta}_1)\cdot|Y|_M;
\end{multline*}
in the last step $\exp(x)-1\le x\cdot\exp(x)$ is used. By~\eqref{eq:renormalisation-group-63} and~\eqref{eq:renormalisation-group-45}, $x\le|Y|_M/(2c_X)\le\left(d_M(Y)+1\right)/2$; by~\eqref{eq:renormalisation-group-45} and~\eqref{eq:renormalisation-group-50},
\[
x\le4\,\mathbf{H}_0\cdot K_0\cdot c_X\cdot\exp(3\boldsymbol{\beta}_1)\cdot\exp\left(\frac{d_M(Y)}{2}\right).
\]
Therefore $x\cdot\exp(x)\le4\exp\left(\frac{1}{2}\right)\cdot c_X\cdot K_0\cdot\exp(3\boldsymbol{\beta}_1)\cdot\mathbf{H}_0\cdot\exp\left(d_M(Y)\right)$. Since $4\exp\left(\frac{1}{2}\right)\le7$ and $\boldsymbol{\beta}_1-1=\kappa_4$, one obtains
\[
\left\|\mathsf{K}(Y)\right\|_{(\mathsf{w};1)}\le\mathbf{H}_1\cdot\exp\left(-\kappa_4\cdot d_M(Y)\right)\qquad\left(Y\in D_k,\ \mathsf{Z}\in\mathrm{Box}_k\right).
\]

\emph{Step~2 (integration).} The measure $\hat{\mu}$ is a probability measure, and $\mathsf{K}(Y)$ depends continuously on $\mathsf{Z}$. By~\eqref{eq:renormalisation-group-42} and Step~1,
\[
\left\|\mathsf{H}(Y)\right\|_{\mathsf{w}}\le\mathbf{H}_1\cdot\exp\left(-\kappa_4\cdot d_M(Y)\right)\qquad\left(Y\in D_k\right).
\]
Every polymer of a set with union $Y$ is contained in $Y$. By the hypotheses on the generators $W$ and on the dependence on $\mathsf{Z}$, the element $\mathsf{K}(Y)$ therefore contains no generator $W(u)$ whose unit site lies outside $Y$, and it depends on $\mathsf{Z}$ only through the coordinates $\mathsf{Z}(\Upsilon)$ with $\Upsilon$ contained in $Y$. Let $Y_1$, $Y_2$ be compatible polymers. They have no $M$-cube in common, for otherwise their union would be connected through faces. By the last assertion of Corollary~\ref{cor:renormalisation-group-1} they have no unit site in common, and no basis element $\Upsilon$ is contained in both of them. Hence no unit site carries a generator $W$ of $\mathsf{K}(Y_1)$ and a generator $W$ of $\mathsf{K}(Y_2)$, and no coordinate $\mathsf{Z}(\Upsilon)$ enters both $\mathsf{K}(Y_1)$ and $\mathsf{K}(Y_2)$. Under $\hat{\mu}$ the coordinates $\mathsf{Z}(\Upsilon)$ are independent: $\hat{\mu}$ is a product, over the basis $\{\Upsilon\}$, of probability measures on $\mathbb{R}$ (Proposition~\ref{prop:basis-map-box}). Let $Y_1,\dots,Y_m$ be pairwise compatible polymers. By~\eqref{eq:renormalisation-group-43}, applied $m-1$ times at fixed $\mathsf{Z}$, and by Fubini's theorem,
\[
\int\mathrm{d}\mu_I(W)\,\mathrm{d}\hat{\mu}(\mathsf{Z})\,\prod_{i=1}^m\mathsf{K}(Y_i)=\prod_{i=1}^m\mathsf{H}(Y_i).
\]
The map that assigns to a union of $M$-cubes the set of its components is a bijection from the unions of $M$-cubes onto the finite sets of pairwise compatible polymers, the empty union corresponding to the empty set; the inverse map assigns to a set of polymers their union. Integrating the expansion of Step~1 and using its factorisation over components, one obtains
\begin{equation}\label{eq:renormalisation-group-64}
\int\mathrm{d}\mu_I(W)\,\mathrm{d}\hat{\mu}(\mathsf{Z})\,\exp\left\{\sum_{X\in D_k}V(X)\right\}=\sum_{\mathcal{Y}}\prod_{Y\in\mathcal{Y}}\mathsf{H}(Y),
\end{equation}
where $\mathcal{Y}$ ranges over the finite sets of distinct, pairwise compatible polymers, the empty set included.

\emph{Step~3 (bound on the cluster series).} For $Y\in D_k$ put $a(Y):=\mathbf{H}_1\cdot\exp\left(-\kappa_4\cdot d_M(Y)\right)$ and, for $n\ge1$,
\[
\boldsymbol{\Sigma}_n(Y):=\sum_{\substack{(Y_1,\dots,Y_n)\in D_k^n\\Y_1\cup\dots\cup Y_n=Y}}\left|\rho^T(Y_1,\dots,Y_n)\right|\cdot\prod_{i=1}^na(Y_i).
\]
Then $\boldsymbol{\Sigma}_1(Y)=a(Y)$. Let $n\ge2$. By~\eqref{eq:renormalisation-group-52}, $\boldsymbol{\Sigma}_n(Y)$ is at most the sum of $\prod_ia(Y_i)$ over the trees $\mathfrak{t}$ with the vertex set $\{1,\dots,n\}$ and over the $n$-tuples $(Y_1,\dots,Y_n)\in D_k^n$ with union $Y$ such that $Y_i$ and $Y_j$ are incompatible for every edge $\{i,j\}$ of $\mathfrak{t}$. For such an $n$-tuple the incompatibility graph contains $\mathfrak{t}$ and is therefore connected. Since $\kappa_4=\boldsymbol{\beta}_2+2\kappa_0$ and $\boldsymbol{\beta}_2>0$, the bound of~\eqref{eq:renormalisation-group-48} gives
\[
\prod_{i=1}^n\exp\left(-\kappa_4\cdot d_M(Y_i)\right)\le\exp\left(-\boldsymbol{\beta}_2\cdot d_M(Y)\right)\cdot\exp\left(3\boldsymbol{\beta}_2\cdot(n-1)\right)\cdot\prod_{i=1}^n\exp\left(-2\kappa_0\cdot d_M(Y_i)\right).
\]
From here on the condition that the union is $Y$ is replaced by the weaker condition $Y_i\subseteq Y$, $i=1,\dots,n$. Regard $\mathfrak{t}$ as a tree with the root $1$, and let $c_j$ count the children of the vertex $j$. If a polymer is incompatible with a polymer $Y'$, then it contains an $M$-cube of $Y'$ or an $M$-cube adjacent through a face to an $M$-cube of $Y'$; there are at most $7\,|Y'|_M$ such $M$-cubes. For every polymer $X$ and every integer $c\ge0$ one has $|X|_M^c\le c!\cdot\exp\left(|X|_M\right)$ and, by~\eqref{eq:renormalisation-group-45} and $\kappa_0\ge c_X$ (Proposition~\ref{prop:polymers}), $|X|_M\le c_X\cdot\left(d_M(X)+1\right)\le\kappa_0\cdot\left(d_M(X)+1\right)$; hence
\[
\exp\left(-2\kappa_0\cdot d_M(X)\right)\cdot|X|_M^c\le c!\cdot\exp(\kappa_0)\cdot\exp\left(-\kappa_0\cdot d_M(X)\right).
\]
With~\eqref{eq:renormalisation-group-46} it follows that, for every polymer $Y'$,
\begin{gather*}
\sum_{\substack{Y_j\in D_k\\Y_j\text{ incompatible with }Y'}}\exp\left(-2\kappa_0\cdot d_M(Y_j)\right)\cdot|Y_j|_M^{c_j}\le7\,c_j!\cdot\exp(\kappa_0)\cdot K_0\cdot|Y'|_M,\\
\sum_{\substack{Y_1\in D_k\\Y_1\subseteq Y}}\exp\left(-2\kappa_0\cdot d_M(Y_1)\right)\cdot|Y_1|_M^{c_1}\le c_1!\cdot\exp(\kappa_0)\cdot K_0\cdot|Y|_M.
\end{gather*}
Sum over the polymers $Y_j$ from the leaves of $\mathfrak{t}$ towards the root, each vertex after all of its children. When $Y_j$ is summed, the sums over the polymers of its $c_j$ children have contributed constant factors and the factor $|Y_j|_M^{c_j}$, and the first inequality applies, with $Y'$ the polymer of the parent of $j$, after the condition $Y_j\subseteq Y$ has been dropped; at the root the second inequality applies. Hence the sum of $\prod_i\exp\left(-2\kappa_0\cdot d_M(Y_i)\right)$ over the $n$-tuples $(Y_1,\dots,Y_n)$ with $Y_i\subseteq Y$ for every $i$ and with incompatible $Y_i$, $Y_j$ for every edge $\{i,j\}$ of $\mathfrak{t}$ is at most
\[
\left(7\exp(\kappa_0)\cdot K_0\right)^n\cdot|Y|_M\cdot\prod_{j=1}^nc_j!.
\]
In $\mathfrak{t}$ the root has the degree $c_1$, and every other vertex $j$ has the degree $c_j+1$; moreover $\sum_jc_j=n-1$ and $c_1\ge1$. By Cayley's formula for the number of trees with prescribed degrees, the number of the trees with the vertex set $\{1,\dots,n\}$ that have a given family $(c_j)$ is
\[
\frac{(n-2)!}{(c_1-1)!\cdot\prod_{j\ge2}c_j!}.
\]
Hence
\[
\sum_{\mathfrak{t}}\prod_{j=1}^nc_j!=(n-2)!\cdot\sum_{(c_j)}c_1\le(n-2)!\cdot(n-1)\cdot\binom{2n-2}{n-1}\le(n-1)!\cdot4^{n-1},
\]
where $(c_j)$ ranges over the families of integers $c_1\ge1$, $c_2,\dots,c_n\ge0$ with sum $n-1$; there are at most $\binom{2n-2}{n-1}$ such families, and $c_1\le n-1$. Altogether, by the definitions of $a$ and of $z_S$,
\[
\frac{\boldsymbol{\Sigma}_n(Y)}{n!}\le7\,\mathbf{H}_1\cdot\exp(\kappa_0)\cdot K_0\cdot|Y|_M\cdot\exp\left(-\boldsymbol{\beta}_2\cdot d_M(Y)\right)\cdot z_S^{n-1};
\]
this holds also for $n=1$, because $\kappa_4\ge\boldsymbol{\beta}_2$ and $7\exp(\kappa_0)\cdot K_0\cdot|Y|_M\ge1$. By~\eqref{eq:renormalisation-group-63}, $z_S\le\frac{1}{2}$. By~\eqref{eq:renormalisation-group-45} and~\eqref{eq:renormalisation-group-50}, $|Y|_M\le c_X\cdot\exp\left(d_M(Y)\right)$; moreover $\boldsymbol{\beta}_2-1=\kappa_3$. Therefore
\begin{multline*}
\sum_{n\ge1}\frac{\boldsymbol{\Sigma}_n(Y)}{n!}\le14\,c_X\cdot\exp(\kappa_0)\cdot K_0\cdot\mathbf{H}_1\cdot\exp\left(-\kappa_3\cdot d_M(Y)\right)\\
=98\,c_X^2\cdot K_0^2\cdot\exp\left(\kappa_0+3\boldsymbol{\beta}_1\right)\cdot\mathbf{H}_0\cdot\exp\left(-\kappa_3\cdot d_M(Y)\right)\le c'_S\cdot\mathbf{H}_0\cdot\exp\left(-\kappa_3\cdot d_M(Y)\right),
\end{multline*}
because $\kappa_0+3\boldsymbol{\beta}_1=3\kappa-8\kappa_0-6$ and $98\exp\left(-8\kappa_0-6\right)\le14$. For an integer $\mathsf{n}\ge0$ one has $\sum_{n>\mathsf{n}}z_S^{n-1}\le2\cdot2^{-\mathsf{n}}$, and the same computation gives
\[
\sum_{n>\mathsf{n}}\frac{\boldsymbol{\Sigma}_n(Y)}{n!}\le2^{-\mathsf{n}}\cdot c'_S\cdot\mathbf{H}_0\cdot\exp\left(-\kappa_3\cdot d_M(Y)\right).
\]

\emph{Step~4 (the logarithm).} The set $D_k$ is finite. Let $X_1,\dots,X_n\in D_k$. The number $\rho^T(X_1,\dots,X_n)$ vanishes unless the incompatibility graph of $(X_1,\dots,X_n)$ is connected, since otherwise no graph enters its definition; and if this graph is connected, the union $X_1\cup\dots\cup X_n$ is a polymer, by~\eqref{eq:renormalisation-group-48}. Therefore, by Step~3, the power series
\[
\boldsymbol{\Phi}(z):=\sum_{n\ge1}\frac{1}{n!}\sum_{(X_1,\dots,X_n)\in D_k^n}\rho^T(X_1,\dots,X_n)\cdot\prod_{i=1}^nz_{X_i}
\]
in the exponent of~\eqref{eq:renormalisation-group-51} converges absolutely for $|z_Y|\le a(Y)$, $Y\in D_k$: on this polydisk the moduli of its terms add up to at most $\sum_{Y\in D_k}\sum_{n\ge1}\boldsymbol{\Sigma}_n(Y)/n!$. Hence the power series obtained by expanding $\exp\boldsymbol{\Phi}(z)$ converges absolutely on the same polydisk: the moduli of its coefficients do not exceed the coefficients of the exponential of the series obtained from $\boldsymbol{\Phi}$ by replacing every number $\rho^T(X_1,\dots,X_n)$ by its modulus. Let two formal power series be equal, and let both converge absolutely on a polydisk. If one substitutes for the variables elements of a commutative unital Banach algebra whose norms do not exceed the radii of the polydisk, the two series turn into absolutely convergent series that differ only by a rearrangement of their terms, and so have the same sum. Put $z_Y:=\mathsf{H}(Y)$ for $Y\in D_k$. These are elements of the even part of $\mathrm{Gr}[\psi,\chi]$, because $\mathsf{K}(Y)$ is even and, by Wick's theorem, the integral over $W$ of a monomial that contains an odd number of generators $W$ vanishes. By~\eqref{eq:renormalisation-group-40} this even part is a commutative Banach algebra with unit for the norm $\|\cdot\|_{\mathsf{w}}$, and $\|\mathsf{H}(Y)\|_{\mathsf{w}}\le a(Y)$ by Step~2. After this substitution the left side of~\eqref{eq:renormalisation-group-51} is the right side of~\eqref{eq:renormalisation-group-64}. The terms of the exponent are grouped according to the union $X$ of $X_1,\dots,X_n$; every such union with $\rho^T(X_1,\dots,X_n)\ne0$ is a polymer, and the terms whose union is a given $X\in D_k$ make up the series for $E^{\#}(X)$. By~\eqref{eq:renormalisation-group-40}, $\left\|\prod_i\mathsf{H}(X_i)\right\|_{\mathsf{w}}\le\prod_ia(X_i)$; so the norms of the terms of the series for $E^{\#}(X)$ add up to at most $\sum_{n\ge1}\boldsymbol{\Sigma}_n(X)/n!$, and the norms of its terms with $n>\mathsf{n}$ add up to at most $\sum_{n>\mathsf{n}}\boldsymbol{\Sigma}_n(X)/n!$. By Step~3 this proves the absolute convergence and the two bounds of the lemma; the identity of the lemma is~\eqref{eq:renormalisation-group-64} combined with~\eqref{eq:renormalisation-group-51}.

\emph{Approximation by polynomials.} Let $\mathsf{n}\ge1$ be an integer. Let $\mathsf{K}_{\mathsf{n}}$, $\mathsf{H}_{\mathsf{n}}$ and $E^{\#}_{\mathsf{n}}(X)$ be formed as $\mathsf{K}$, $\mathsf{H}$ and $E^{\#}(X)$, with $\mathrm{pl}_{\mathsf{n}}\left(V(X')\right)$ in place of $\exp\left(V(X')\right)-1$ and with the series~\eqref{eq:renormalisation-group-61} stopped at $n=\mathsf{n}$; thus $E^{\#}_{\mathsf{n}}(X)$ is the partial sum of the statement. By construction, $\mathsf{K}_{\mathsf{n}}(Y)$ is a polynomial in the elements $V(X')$ with $X'\subseteq Y$, the element $\mathsf{H}_{\mathsf{n}}(Y)$ is a linear combination of integrals $\int\mathrm{d}\mu_I(W)\,\mathrm{d}\hat{\mu}(\mathsf{Z})$ of products of these elements, and $E^{\#}_{\mathsf{n}}(X)$ is a polynomial in the elements $\mathsf{H}_{\mathsf{n}}(X_i)$ with $X_i\subseteq X$. By~\eqref{eq:renormalisation-group-40}, $\left\|\mathrm{pl}_{\mathsf{n}}\left(V(X)\right)\right\|_{(\mathsf{w};1)}\le\exp\left(\|V(X)\|_{(\mathsf{w};1)}\right)-1$. The bounds of Steps~1--3 therefore apply to $\mathsf{K}_{\mathsf{n}}$, $\mathsf{H}_{\mathsf{n}}$ and $E^{\#}_{\mathsf{n}}$ without any change of the constants; in particular $\left\|\mathsf{H}_{\mathsf{n}}(Y)\right\|_{\mathsf{w}}\le a(Y)$. Since $\|V(X)\|_{(\mathsf{w};1)}\le\mathbf{H}_0\le1$,
\[
\left\|\exp\left(V(X)\right)-1-\mathrm{pl}_{\mathsf{n}}\left(V(X)\right)\right\|_{(\mathsf{w};1)}\le\sum_{j>\mathsf{n}}\frac{1}{j!}\le\frac{\exp(1)}{(\mathsf{n}+1)!}.
\]
Expanding the products, one writes $\mathsf{K}(Y)-\mathsf{K}_{\mathsf{n}}(Y)$ as a sum, finite because $D_k$ is finite, of products in which one factor is a difference of this kind and every other factor is $\exp\left(V(X')\right)-1$ or $\mathrm{pl}_{\mathsf{n}}\left(V(X')\right)$, hence of norm at most $\exp(1)-1\le2$. The number of these products is independent of $V$. By~\eqref{eq:renormalisation-group-42}, and since $\hat{\mu}$ is a probability measure, it follows that
\[
\boldsymbol{\delta}_{\mathsf{n}}:=\max_{Y\in D_k}\frac{\left\|\mathsf{H}(Y)-\mathsf{H}_{\mathsf{n}}(Y)\right\|_{\mathsf{w}}}{a(Y)}
\]
tends to $0$ as $\mathsf{n}\to\infty$, uniformly in $V$. In a term with $n\le\mathsf{n}$ of the series for $E^{\#}(X)$, the replacement of the elements $\mathsf{H}(X_i)$ by the elements $\mathsf{H}_{\mathsf{n}}(X_i)$ changes the product $\prod_i\mathsf{H}(X_i)$ by at most $n\cdot\boldsymbol{\delta}_{\mathsf{n}}\cdot\prod_ia(X_i)$ in the norm $\|\cdot\|_{\mathsf{w}}$: the difference of the two products is a sum of $n$ products, each of which has one factor $\mathsf{H}(X_i)-\mathsf{H}_{\mathsf{n}}(X_i)$, while its other factors have norms at most $a(X_i)$. By Step~3 the sum of these changes is at most
\[
\boldsymbol{\delta}_{\mathsf{n}}\cdot\sum_{n\ge1}n\cdot\frac{\boldsymbol{\Sigma}_n(X)}{n!}\le28\,\boldsymbol{\delta}_{\mathsf{n}}\cdot\mathbf{H}_1\cdot\exp(\kappa_0)\cdot K_0\cdot|X|_M,
\]
because $\sum_{n\ge1}n\cdot z_S^{n-1}\le4$. The norms of the terms with $n>\mathsf{n}$ of the series for $E^{\#}(X)$ add up to at most $2^{-\mathsf{n}}\cdot c'_S\cdot\mathbf{H}_0\cdot\exp\left(-\kappa_3\cdot d_M(X)\right)$, by Step~3. None of these bounds depends on $V$ otherwise than through the hypotheses of the lemma. Consequently
\[
\left\|E^{\#}(X)-E^{\#}_{\mathsf{n}}(X)\right\|_{\mathsf{w}}\le28\,\boldsymbol{\delta}_{\mathsf{n}}\cdot\mathbf{H}_1\cdot\exp(\kappa_0)\cdot K_0\cdot|X|_M+2^{-\mathsf{n}}\cdot c'_S\cdot\mathbf{H}_0\cdot\exp\left(-\kappa_3\cdot d_M(X)\right),
\]
and the right side tends to $0$ as $\mathsf{n}\to\infty$, uniformly on every set of families $V$ that satisfy the hypotheses of the lemma.
\end{proof}

\subsection{Bounds, analyticity and Lipschitz dependence of the step map}\label{sec:bounds-analyticity-lipschitz}

\subsubsection{The activity before reblocking}\label{sec:activity-before-reblocking}

Let $(m,\mathscr{E})\in\mathrm{Ball}^x_k(4)$, $A\in\tilde{R}^{-}_k$ and $\mathscr{C}\in\mathfrak{X}_k$. For $X\in D_k$ write $E^{\#,x}_k(X;A,\mathscr{C})$ for the element $E^{\#,x}_k(m,\mathscr{E})(X;A,\mathscr{C},\psi,\chi)$ of $\mathrm{Gr}[\psi,\chi]$ defined by~\eqref{eq:renormalisation-group-61}, and $E^{(\le5),x}(X)$ for the element defined by~\eqref{eq:renormalisation-group-59} with the same data. The following five assertions hold.

\begin{enumerate}
\item[(1)] \emph{Convergence and bound.} For every $X\in D_k$ the series~\eqref{eq:renormalisation-group-61} that defines $E^{\#,x}_k(X;A,\mathscr{C})$ converges absolutely in the norm $\|\cdot\|_{h_k/4}$, and
\[
\left\|E^{\#,x}_k(X;A,\mathscr{C})\right\|_{h_k/4}\le c_S'\cdot\mathbf{H}_0\cdot\exp\left(-\kappa_3\cdot d_M(X)\right),
\]
where $c_S'$ is the constant of Lemma~\ref{lem:integration-polymer-exponential} and $\mathbf{H}_0$ is the number defined in~\eqref{eq:renormalisation-group-37}.

Indeed, put $V(X):=E^{(\le5),x}(X)$ for $X\in D_k$. By Corollary~\ref{cor:renormalisation-group-1}, $V(X)$ lies in $\mathrm{Gr}[\psi,\chi,W]$, contains $W(u)$ only for $u$ in $X$, and depends on $\mathsf{Z}$ only through the coordinates $\mathsf{Z}(\Upsilon)$ with $\Upsilon$ contained in $X$. By Proposition~\ref{prop:bounds-regularity-localised}(e), $V(X)$ is even and continuous in $\mathsf{Z}\in\mathrm{Box}_k$, and by Proposition~\ref{prop:bounds-regularity-localised}(c),
\[
\left\|V(X)\right\|_{\left(h_k/4;1\right)}\le\mathbf{H}_0\cdot\exp\left(-\kappa_2\cdot d_M(X)\right).
\]
Together with condition~\eqref{eq:renormalisation-group-37}, these are the hypotheses of Lemma~\ref{lem:integration-polymer-exponential} for the weights $\mathsf{w}=h_k/4$. Since $\mathsf{K}$, $\mathsf{H}$ and $E^{\#,x}_k(X;A,\mathscr{C})$ arise from the family $\left\{V(X):X\in D_k\right\}$ through~\eqref{eq:renormalisation-group-61}, that lemma yields the absolute convergence and the bound.

\item[(2)] \emph{Locality.} For every $X\in D_k$ the element $E^{\#,x}_k(X;A,\mathscr{C})$ contains $\psi(\xi)$ only for indices $\xi=(x,\boldsymbol{\beta},\boldsymbol{\omega})$ with $x\in X$, and $\chi(\boldsymbol{\zeta})$ only for indices $\boldsymbol{\zeta}=(x,x',\boldsymbol{\beta},\boldsymbol{\omega})$ whose pair $(x,x')$ is admissible for $X$; and its coefficients take the same values at any two pairs of connections in $\tilde{R}^{-}_k\times\mathfrak{X}_k$ that coincide on the bonds that meet $X$.

Indeed, by Corollary~\ref{cor:renormalisation-group-1}, for every polymer $Z$ the element $E^{(\le5),x}(Z)$ contains $\psi$ only at sites of $Z$, $\chi$ only at pairs admissible for $Z$, and $W(u)$ only for $u$ in $Z$; its coefficients depend on $\mathsf{Z}$ only through the coordinates $\mathsf{Z}(\Upsilon)$ with $\Upsilon$ contained in $Z$, and on $(A,\mathscr{C})$ only through the values on the bonds that meet $Z$. Let $Z'\in D_k$. The element $\mathsf{K}(Z')$ is a sum of products of elements $\exp\left(E^{(\le5),x}(Z)\right)-1$ with polymers $Z\subseteq Z'$. A pair that is admissible for such a $Z$ is admissible for $Z'$, by~\eqref{eq:renormalisation-group-adm}, and a bond that meets $Z$ meets $Z'$. The integration with respect to $\mathrm{d}\mu_I(W)\,\mathrm{d}\hat{\mu}(\mathsf{Z})$, which turns $\mathsf{K}(Z')$ into $\mathsf{H}(Z')$, eliminates the generators $W$ and the variables $\mathsf{Z}$. Consequently $\mathsf{H}(Z')$ contains $\psi$ only at sites of $Z'$ and $\chi$ only at pairs admissible for $Z'$, and its coefficients depend on $(A,\mathscr{C})$ only through the values on the bonds that meet $Z'$. By~\eqref{eq:renormalisation-group-61}, $E^{\#,x}_k(X;A,\mathscr{C})$ is the sum of a series each term of which is a multiple of a product $\mathsf{H}(X_1)\cdots\mathsf{H}(X_n)$ with polymers $X_i\subseteq X$; by~\eqref{eq:renormalisation-group-adm} once more, each term has the two properties stated for $X$. Both properties concern single coefficients, and by~(1) each coefficient of the sum is the sum of the corresponding coefficients of the terms; so the sum has them as well.

\item[(3)] \emph{Analyticity.} The coefficients of $E^{\#,x}_k(X;A,\mathscr{C})$ are analytic functions of $(A,\mathscr{C})$ on the set $\tilde{R}^{-}_k\times\mathfrak{X}_k$, which is open by Lemma~\ref{lem:openness}. Moreover, let $(m_a,\mathscr{E}_a)$ and $(m_b,\mathscr{E}_b)$ be elements of $\mathbb{C}\times K^{x,C}_k$, and let $\Theta$ be the set of the $t\in\mathbb{C}$ for which $\left(m_a+t\cdot m_b,\,\mathscr{E}_a+t\cdot\mathscr{E}_b\right)\in\mathrm{Ball}^x_k(4)$; this set is closed, because the ball is closed. Then the coefficients of $E^{\#,x}_k\left(m_a+t\cdot m_b,\,\mathscr{E}_a+t\cdot\mathscr{E}_b\right)(X;A,\mathscr{C},\psi,\chi)$ are, as functions of $t$, uniform limits on $\Theta$ of polynomials in $t$; in particular they are continuous on $\Theta$ and analytic in the interior of $\Theta$.

Indeed, by the last assertion of Lemma~\ref{lem:integration-polymer-exponential}, applied to the families $V$ of~(1), the element $E^{\#,x}_k(X;A,\mathscr{C})$ is the limit in $\|\cdot\|_{h_k/4}$, uniformly in $(A,\mathscr{C})\in\tilde{R}^{-}_k\times\mathfrak{X}_k$ and in $(m,\mathscr{E})\in\mathrm{Ball}^x_k(4)$, of polynomials in integrals, with respect to $\mathrm{d}\mu_I(W)\,\mathrm{d}\hat{\mu}(\mathsf{Z})$, of products of the elements $E^{(\le5),x}(X_i)$, $X_i\subseteq X$. By Proposition~\ref{prop:bounds-regularity-localised}(e) these elements are continuous in $(A,\mathscr{C},\mathsf{Z})$, analytic in $(A,\mathscr{C})$ and affine in $(m,\mathscr{E})$. The integral over $W$ is a linear map between finite-dimensional spaces. Let $\Phi$ be a complex function on $\tilde{R}^{-}_k\times\mathfrak{X}_k\times\mathrm{Box}_k$ that is continuous and, for each fixed $\mathsf{Z}$, analytic in $(A,\mathscr{C})$. Then $(A,\mathscr{C})\mapsto\int\mathrm{d}\hat{\mu}(\mathsf{Z})\,\Phi(A,\mathscr{C},\mathsf{Z})$ is analytic, since it is approximated by finite sums of analytic functions, uniformly on each compact set. To see this, let $Q\subseteq\tilde{R}^{-}_k\times\mathfrak{X}_k$ be compact. As $\mathrm{Box}_k$ is compact, $\Phi$ is uniformly continuous on $Q\times\mathrm{Box}_k$. Choose finite partitions of $\mathrm{Box}_k$ into Borel sets $U_j$, with diameters tending to $0$, and points $\mathsf{Z}_j\in U_j$; then the sums $\sum_j\hat{\mu}(U_j)\cdot\Phi(\cdot\,,\mathsf{Z}_j)$, each of them analytic in $(A,\mathscr{C})$, converge to the integral uniformly on $Q$. So the approximating elements have coefficients that are analytic in $(A,\mathscr{C})$. For $(m,\mathscr{E})=\left(m_a+t\cdot m_b,\,\mathscr{E}_a+t\cdot\mathscr{E}_b\right)$ with $t\in\Theta$ they are polynomials in $t$, since products of affine functions of $t$ are polynomials in $t$ and the integrals are linear. A uniform limit of analytic functions is analytic, and a uniform limit of continuous functions is continuous; both assertions follow.

\item[(4)] \emph{Scaling of pairs.} Let $Y\in D_{k+1}$, and let $X$ be a union of $M$-cubes of $T_k$ with $X\subseteq LY$. The map $x\mapsto x/L$ is a bijection of the sites of $T_k$ onto the sites of $T_{k+1}$; it divides distances by $L$, maps nearest neighbours of $T_k$ to nearest neighbours of $T_{k+1}$, maps the path $\Gamma^{\pi}(x,x')$ of Definition~\ref{def:pairs-sites-holder} onto the path $\Gamma^{\pi}(x/L,x'/L)$, and maps the sites of $X$ into $Y$. As a consequence, if $(x,x')\in\mathrm{Pz}_k$ is admissible for $X$, then $(x/L,x'/L)$ is an element of $\mathrm{Pz}_{k+1}$ admissible for $Y$, namely of kind~(I) for $Y$ when $(x,x')$ is of kind~(I) for $X$, and of kind~(II) for $Y$ when $(x,x')$ is of kind~(II) for $X$; and if a bond $b$ of $T_k$ meets $X$, that is, has an end point in $X$, then its image $b/L$ is a bond of $T_{k+1}$ that meets $Y$. Indeed, two distinct sites at distance at most $M$ go to two distinct sites at distance at most $M/L\le M$; the two sites and the six paths of a pair of kind~(I) for $X$ lie in $X$, so their images lie in $Y$; the images of two nearest neighbours one of which lies in $X$ are nearest neighbours one of which lies in $Y$; and an end point in $X$ of a bond goes to an end point in $Y$ of the image bond.

\item[(5)] \emph{Reblocking.} For every $X\in D_k$ there is exactly one $Y\in D_{k+1}$ with $\bar{X}^L=LY$, and $X\subseteq LY$.

Indeed, $\bar{X}^L$ is a union of $LM$-cubes. By fact~\eqref{it:renormalisation-group-g0} of Section~\ref{sec:cubes-parameter-cubes}, every $M$-cube lies in a single $LM$-cube, and when two $M$-cubes share a face, the $LM$-cubes in which they lie are equal or share a face; as $X$ is connected through faces, so is $\bar{X}^L$. Since $l\in\mathrm{Cl}$, the quotient $l/(LM)$ is an integer; accordingly, under $x\mapsto x/L$ every $LM$-cube of $T_k$ becomes an $M$-cube of $T_{k+1}$, and every $M$-cube of $T_{k+1}$ arises in this way. Hence the image $Y$ of $\bar{X}^L$ under this map is a union of $M$-cubes of $T_{k+1}$ that is connected through faces, that is, $Y\in D_{k+1}$, and $LY=\bar{X}^L$; as the map is bijective, no other $Y'\in D_{k+1}$ satisfies $LY'=\bar{X}^L$. Finally, each $M$-cube of $X$ lies in an $LM$-cube, by fact~\eqref{it:renormalisation-group-g0}, and this $LM$-cube belongs to $\bar{X}^L$; so $X\subseteq LY$.
\end{enumerate}

\subsubsection{Bound and analyticity after reblocking}\label{sec:bound-analyticity-after}

\begin{proof}[Proof of clause~\eqref{it:step-small-fields-a} of Theorem~\ref{thm:step-small-fields}]
Let $(m,\mathscr{E})\in\mathrm{Ball}^x_k(4)$. For all $A\in\tilde{R}^{-}_k$ and $\mathscr{C}\in\mathfrak{X}_k$ the series~\eqref{eq:renormalisation-group-61} converge absolutely in $\|\cdot\|_{h_k/4}$, by assertion~(1) of Section~\ref{sec:activity-before-reblocking}.

\emph{Definition; the clauses~\eqref{it:renormalisation-group-1-k1} and~\eqref{it:renormalisation-group-1-k2}.} Let $(\mathscr{A},\mathscr{C})\in\tilde{R}_{k+1}\times\mathfrak{X}_{k+1}$. By Lemma~\ref{lem:scaling-domains}(a), (b), (c), applied with $c=1$, one has $\mathscr{A}_L\in L^{-\frac{3}{4}}\cdot\tilde{R}_k\subseteq\tilde{R}^{-}_k$ and $\mathscr{C}_L\in\mathfrak{X}_k$. So the elements $E^{\#,x}_k(m,\mathscr{E})\left(X;\mathscr{A}_L,\mathscr{C}_L,\cdot\,\right)$, $X\in D_k$, are defined, and with them the elements $E^{*,x}_k(m,\mathscr{E})(Y;\mathscr{A},\mathscr{C},\cdot\,)$, $Y\in D_{k+1}$, of Section~\ref{sec:fluctuation-integral-new}. The map $(\mathscr{A},\mathscr{C})\mapsto\left(\mathscr{A}_L,\mathscr{C}_L\right)$ is linear and carries $\tilde{R}_{k+1}\times\mathfrak{X}_{k+1}$ into $\tilde{R}^{-}_k\times\mathfrak{X}_k$, and the substitution $\psi\mapsto\psi_L$, $\chi\mapsto\chi_L$ does not involve the connections; hence assertion~(3) of Section~\ref{sec:activity-before-reblocking} shows that the kernels of $E^{*,x}_k(m,\mathscr{E})(Y;\mathscr{A},\mathscr{C})$ are analytic in $(\mathscr{A},\mathscr{C})$ on $\tilde{R}_{k+1}\times\mathfrak{X}_{k+1}$. This is clause~\eqref{it:renormalisation-group-1-k1} at level $k+1$.

Let $Y\in D_{k+1}$, and let $X\in D_k$ satisfy $\bar{X}^L=LY$; then $X\subseteq LY$, by assertion~(5) of Section~\ref{sec:activity-before-reblocking}. By assertion~(2), the element $E^{\#,x}_k(m,\mathscr{E})\left(X;\mathscr{A}_L,\mathscr{C}_L,\psi,\chi\right)$ contains $\psi$ only at sites of $X$ and $\chi$ only at pairs admissible for $X$, and it depends on $\left(\mathscr{A}_L,\mathscr{C}_L\right)$ only through the values on the bonds that meet $X$. The substitution $\psi\mapsto\psi_L$, $\chi\mapsto\chi_L$ replaces a generator at a site $x$, resp.\ at a pair $(x,x')$, by a multiple of the generator of level $k+1$ at $x/L$, resp.\ at $(x/L,x'/L)$; by assertion~(4) these sites lie in $Y$ and these pairs are admissible for $Y$. Moreover $\mathscr{A}_L(b)=L^{-\frac{1}{2}}\cdot\mathscr{A}(b/L)$ and $\mathscr{C}_L(b)=L^{-\frac{1}{2}}\cdot\mathscr{C}(b/L)$, and $b/L$ meets $Y$ whenever $b$ meets $X$, again by assertion~(4). Hence the kernels of $E^{*,x}_k(m,\mathscr{E})(Y;\mathscr{A},\mathscr{C})$ vanish unless all their sites lie in $Y$ and all their pairs are admissible for $Y$, and they depend on $\mathscr{A}$, $\mathscr{C}$ only through the values on the bonds that meet $Y$. This is clause~\eqref{it:renormalisation-group-1-k2} at level $k+1$.

\emph{Scaling of the norms.} Let $F$ be an element of $\mathrm{Gr}[\psi,\chi]$ of level $k$ with the kernels $F_{nm}$ (Proposition~\ref{prop:weighted-coefficient-norms}), and let $F'$ be the element of level $k+1$ obtained from $F$ by the substitution $\psi\mapsto\psi_L$, $\chi\mapsto\chi_L$ of Section~\ref{sec:local-parts-adjustment}, where $\psi_L(x)=L^{-1}\cdot\psi(x/L)$ and $\chi_L(x,x')=L^{-1-\alpha}\cdot\chi(x/L,x'/L)$. Every site $x$ of $T_k$ is $x=Lx'$ for exactly one site $x'$ of $T_{k+1}$, and $\eta_k^3=L^3\cdot\eta_{k+1}^3$, $\eta_k^6=L^6\cdot\eta_{k+1}^6$; thus $\sum_x\eta_k^3\cdot g(x)=L^3\cdot\sum_{x'}\eta_{k+1}^3\cdot g(Lx')$ for every function $g$ on the sites, and the analogous identity with the sixth powers holds for the pairs. For an index $\xi'=(x',\boldsymbol{\beta},\boldsymbol{\omega})$ of level $k+1$ put $L\xi':=(Lx',\boldsymbol{\beta},\boldsymbol{\omega})$, and for an index $\boldsymbol{\zeta}'=(x',x'',\boldsymbol{\beta},\boldsymbol{\omega})$ put $L\boldsymbol{\zeta}':=(Lx',Lx'',\boldsymbol{\beta},\boldsymbol{\omega})$. Every index $\boldsymbol{\zeta}$ of level $k$ is $L\boldsymbol{\zeta}'$ for exactly one index $\boldsymbol{\zeta}'$ of level $k+1$, since $(x/L,x'/L)\in\mathrm{Pz}_{k+1}$ for $(x,x')\in\mathrm{Pz}_k$ (Section~\ref{sec:fluctuation-integral-new}). It follows that the kernels of $F'$ are
\[
F'_{nm}\left(\xi'_1,\dots,\xi'_n;\boldsymbol{\zeta}'_1,\dots,\boldsymbol{\zeta}'_m\right)=L^{3n+6m}\cdot L^{-n}\cdot L^{-(1+\alpha)m}\cdot F_{nm}\left(L\xi'_1,\dots,L\xi'_n;L\boldsymbol{\zeta}'_1,\dots,L\boldsymbol{\zeta}'_m\right),
\]
where the right side is read as $0$ if the pair of one of the indices $L\boldsymbol{\zeta}'_j$ does not belong to $\mathrm{Pz}_k$, and that
\[
\left\|F'_{nm}\right\|=L^{-n}\cdot L^{-(1+\alpha)m}\cdot\left\|F_{nm}\right\|.
\]
By the expression of the norms through the kernels (Proposition~\ref{prop:weighted-coefficient-norms}),
\[
\left\|F'\right\|_{h_{k+1}}=\left\|F\right\|_{\left(L^{-1}h_{k+1,1},\,L^{-1-\alpha}h_{k+1,2}\right)}\le\left\|F\right\|_{h_k/4}.
\]
The inequality holds because the norm is non-decreasing in each weight, because $h_{k+1,i}\le h_{k,i}$ for $i=1,2$ (the weights $h_k=\left(e_k^{-\frac{1}{4}},e_k^{-\frac{1}{4}+\varepsilon}\right)$ of Section~\ref{sec:grassmann-algebras-norms} are negative powers of $e_k$, and $e_k\le e_{k+1}=L^{\frac{1}{2}}\cdot e_k$), and because $L^{-1-\alpha}\le L^{-1}\le\frac{1}{4}$, as $L\ge4$ by stage~\ref{it:order-choice-4}.

\emph{The sum over the polymers of level $k$.} By stage~\ref{it:order-choice-3}, $\kappa\ge7\kappa_0+7$; in particular $\kappa>0$ and $\kappa_3-\kappa_0=\kappa-6\kappa_0-4>0$. Let $Y\in D_{k+1}$. Every $X\in D_k$ with $\bar{X}^L=LY$ satisfies $L\cdot d_M(Y)\le d_M(X)$, by~\eqref{eq:renormalisation-group-49}, and $X\subseteq LY$, by assertion~(5) of Section~\ref{sec:activity-before-reblocking}; and $LY$ is the union of $L^3\cdot|Y|_M$ of the $M$-cubes of $T_k$. Hence, by~\eqref{eq:renormalisation-group-49} and~\eqref{eq:renormalisation-group-46}, then by~\eqref{eq:renormalisation-group-45} and~\eqref{eq:renormalisation-group-50}, and finally by~\eqref{eq:renormalisation-group-30} in the form $L\cdot\left(\kappa_3-\kappa_0\right)-1=L\cdot\left(\kappa-6\kappa_0-4\right)-1\ge\kappa$,
\begin{align*}
\sum_{X:\,\bar{X}^L=LY}\exp\left(-\kappa_3\cdot d_M(X)\right)&\le\exp\left(-L\cdot\left(\kappa_3-\kappa_0\right)\cdot d_M(Y)\right)\cdot\sum_{X\in D_k:\,X\subseteq LY}\exp\left(-\kappa_0\cdot d_M(X)\right)\\
&\le\exp\left(-L\cdot\left(\kappa_3-\kappa_0\right)\cdot d_M(Y)\right)\cdot K_0\cdot L^3\cdot|Y|_M\\
&\le c_X\cdot K_0\cdot L^3\cdot\exp\left(-\left(L\cdot\left(\kappa_3-\kappa_0\right)-1\right)\cdot d_M(Y)\right)\\
&\le c_X\cdot K_0\cdot L^3\cdot\exp\left(-\kappa\cdot d_M(Y)\right).
\end{align*}

\emph{The bound; clause~\eqref{it:renormalisation-group-1-k6}.} Let $Y\in D_{k+1}$ and $(\mathscr{A},\mathscr{C})\in\tilde{R}_{k+1}\times\mathfrak{X}_{k+1}$. The scaling of the norms, applied to $F=E^{\#,x}_k(m,\mathscr{E})\left(X;\mathscr{A}_L,\mathscr{C}_L\right)$, then assertion~(1) of Section~\ref{sec:activity-before-reblocking} at the connections $\left(\mathscr{A}_L,\mathscr{C}_L\right)\in\tilde{R}^{-}_k\times\mathfrak{X}_k$, and the sum over the polymers give
\begin{align*}
\left\|E^{*,x}_k(m,\mathscr{E})(Y;\mathscr{A},\mathscr{C})\right\|_{h_{k+1}}&\le\sum_{X:\,\bar{X}^L=LY}\left\|E^{\#,x}_k(m,\mathscr{E})\left(X;\mathscr{A}_L,\mathscr{C}_L\right)\right\|_{h_k/4}\\
&\le c_S'\cdot\mathbf{H}_0\cdot\sum_{X:\,\bar{X}^L=LY}\exp\left(-\kappa_3\cdot d_M(X)\right)\\
&\le c_X\cdot K_0\cdot L^3\cdot c_S'\cdot\mathbf{H}_0\cdot\exp\left(-\kappa\cdot d_M(Y)\right).
\end{align*}
With $c_S'=14\cdot c_X^2\cdot K_0^2\cdot\exp(3\kappa)$ (Lemma~\ref{lem:integration-polymer-exponential}) and $\mathbf{H}_0=36\cdot c_X^2\cdot K_0^2\cdot e_k^{\frac{1}{4}-5\varepsilon}$ (see~\eqref{eq:renormalisation-group-37}), the definition of the norm in clause~\eqref{it:renormalisation-group-1-k6} gives
\[
\left\|E^{*,x}_k(m,\mathscr{E})\right\|^x_{k+1}\le c_X\cdot K_0\cdot L^3\cdot c_S'\cdot\mathbf{H}_0=504\cdot c_X^5\cdot K_0^5\cdot\exp(3\kappa)\cdot L^3\cdot e_k^{\frac{1}{4}-5\varepsilon}\le e_k^{\frac{1}{4}-6\varepsilon}.
\]
For the last inequality note that $504\cdot\exp(3\kappa)\le10^5\cdot\exp(6\kappa)$, so that the middle term is at most $c_S\cdot L^3\cdot e_k^{\varepsilon}\cdot e_k^{\frac{1}{4}-6\varepsilon}$, and that $c_S\cdot L^3\cdot e_k^{\varepsilon}\le1$ by~\eqref{eq:renormalisation-group-37}. In particular $E^{*,x}_k(m,\mathscr{E})$ has clause~\eqref{it:renormalisation-group-1-k6} at level $k+1$, and it satisfies the bound of clause~\eqref{it:step-small-fields-a}.

\emph{Symmetries.} By Proposition~\ref{prop:symmetries-step-map}, $E^{*,x}_k(m,\mathscr{E})$ has clause~\eqref{it:renormalisation-group-1-k3} for the group $\mathbb{G}_{k+1}$ and the clauses~\eqref{it:renormalisation-group-1-k4} and~\eqref{it:renormalisation-group-1-k5}; every monomial in it has as many barred as unbarred legs, so that it is even; and if $m$ is real and $\mathscr{E}\in K^x_k$, it has clause~\eqref{it:renormalisation-group-1-k7}. Together with the clauses~\eqref{it:renormalisation-group-1-k1}, \eqref{it:renormalisation-group-1-k2} and~\eqref{it:renormalisation-group-1-k6} proved above, it follows that $E^{*,x}_k(m,\mathscr{E})\in K^{x,C}_{k+1}$ (Definition~\ref{def:renormalisation-group-1}), and that $E^{*,x}_k(m,\mathscr{E})\in K^x_{k+1}$ if $m$ is real and $\mathscr{E}\in K^x_k$.

\emph{Analytic dependence on the data.} Let $(m_a,\mathscr{E}_a)$ and $(m_b,\mathscr{E}_b)$ be elements of $\mathbb{C}\times K^{x,C}_k$, and let $\Theta$ be the set of the $t\in\mathbb{C}$ with $\left(m_a+t\cdot m_b,\,\mathscr{E}_a+t\cdot\mathscr{E}_b\right)\in\mathrm{Ball}^x_k(4)$. Since $\mathrm{Ball}^x_k(4)$ is closed and convex (Section~\ref{sec:spaces-activities}) and $t\mapsto\left(m_a+t\cdot m_b,\,\mathscr{E}_a+t\cdot\mathscr{E}_b\right)$ is affine and continuous, $\Theta$ is closed and convex. If $\Theta$ has at most one point, there is nothing to prove. Let $\Theta$ have more than one point; being convex, it is then infinite. For $t\in\Theta$ put $E(t):=E^{*,x}_k\left(m_a+t\cdot m_b,\,\mathscr{E}_a+t\cdot\mathscr{E}_b\right)$. Let $\mathcal{B}$ be the complex Banach space of the families $\left\{F(Y;\mathscr{A},\mathscr{C},\psi,\chi):Y\in D_{k+1}\right\}$ of elements of $\mathrm{Gr}[\psi,\chi]$ of level $k+1$ that have the clauses~\eqref{it:renormalisation-group-1-k1}, \eqref{it:renormalisation-group-1-k2} and~\eqref{it:renormalisation-group-1-k6} at level $k+1$, with the norm $\|\cdot\|^x_{k+1}$.

By assertion~(3) of Section~\ref{sec:activity-before-reblocking} and its proof, which rests on the last assertion of Lemma~\ref{lem:integration-polymer-exponential}, for each $X\in D_k$ the element $E^{\#,x}_k\left(m_a+t\cdot m_b,\,\mathscr{E}_a+t\cdot\mathscr{E}_b\right)(X;A,\mathscr{C})$ is the limit in $\|\cdot\|_{h_k/4}$ of polynomials in $t$, uniformly in $t\in\Theta$ and in $(A,\mathscr{C})\in\tilde{R}^{-}_k\times\mathfrak{X}_k$. Insert $(A,\mathscr{C})=\left(\mathscr{A}_L,\mathscr{C}_L\right)$, substitute $\psi\mapsto\psi_L$, $\chi\mapsto\chi_L$, and sum over the polymers $X$ with $\bar{X}^L=LY$. By the scaling of the norms, and since $D_k$ and $D_{k+1}$ are finite sets, this produces polynomials $\mathcal{P}_j$ in $t$, $j=1,2,\dots$, whose coefficients are families indexed by $D_{k+1}$, such that
\[
\lim_{j\to\infty}\ \sup_{t\in\Theta}\left\|E(t)-\mathcal{P}_j(t)\right\|^x_{k+1}=0.
\]
For $t\in\Theta$ the value $\mathcal{P}_j(t)$ has the clauses~\eqref{it:renormalisation-group-1-k1} and~\eqref{it:renormalisation-group-1-k2} at level $k+1$, since the arguments of~(2) and~(3) of Section~\ref{sec:activity-before-reblocking}, and those given above for these two clauses of $E^{*,x}_k(m,\mathscr{E})$, apply to it; and the norms $\left\|\mathcal{P}_j(t)\right\|^x_{k+1}$ are bounded on $\Theta$, since the terms of $\mathcal{P}_j(t)$ are majorised, term by term, by those of the series estimated in Lemma~\ref{lem:integration-polymer-exponential} and above. Thus $\mathcal{P}_j(t)\in\mathcal{B}$ for $t\in\Theta$. Let $d$ be the degree of $\mathcal{P}_j$. As $\Theta$ is infinite, it contains $d+1$ distinct points $t_0,\dots,t_d$, and the Vandermonde matrix $\left(t_i^q\right)_{0\le i,q\le d}$ is invertible; by interpolation, each coefficient of $\mathcal{P}_j$ is therefore a linear combination, with complex scalars, of the values $\mathcal{P}_j(t_0),\dots,\mathcal{P}_j(t_d)$, and so lies in $\mathcal{B}$. Consequently, on $\Theta$, the map $t\mapsto E(t)$ is the uniform limit in $\mathcal{B}$ of polynomials in $t$ with coefficients in $\mathcal{B}$. These polynomials are continuous on $\mathbb{C}$ and analytic as maps into $\mathcal{B}$; by the analyticity of uniform limits of analytic functions, applied to functions with values in the Banach space $\mathcal{B}$, the map $t\mapsto E(t)$ is continuous on $\Theta$ and analytic in the interior of $\Theta$ as a map into $\mathcal{B}$.

Within $\mathcal{B}$, the space $K^{x,C}_{k+1}$ is cut out by the clauses~\eqref{it:renormalisation-group-1-k3}, \eqref{it:renormalisation-group-1-k4} and~\eqref{it:renormalisation-group-1-k5}, together with the requirement that each monomial have as many barred legs as unbarred ones. Each of these requirements asks that certain finite linear combinations of values of kernels, every value taken at fixed arguments, vanish; and each such value depends continuously on the family, because a sequence converging in $\|\cdot\|^x_{k+1}$ has every kernel converging at every fixed choice of its arguments. Therefore $K^{x,C}_{k+1}$ is closed in $\mathcal{B}$. The values $E(t)$, $t\in\Theta$, lie in $K^{x,C}_{k+1}$, by the paragraph on symmetries. Hence $t\mapsto E(t)$ is continuous on $\Theta$ and analytic in the interior of $\Theta$ also as a map into $K^{x,C}_{k+1}$.

\emph{The number $d\varepsilon^0_k$.} Proposition~\ref{prop:basis-map-box} and~\eqref{eq:renormalisation-group-31} give
\[
0\le d\varepsilon^0_k\le12\cdot\exp\left(-\frac{p_{0,k}^2}{2}\right)\le e_k^7.
\]
\end{proof}

\subsubsection{The second slot at zero}\label{sec:second-slot-zero}

\begin{proof}[Proof of clause~\eqref{it:step-small-fields-b} of Theorem~\ref{thm:step-small-fields}]
By Section~\ref{sec:activities-slot}, the statements of Sections~\ref{sec:definition-step-map}--\ref{sec:symmetries-step-map} and their proofs remain valid in the case without slot, in which the sets $\mathfrak{X}_k$ and $\mathfrak{X}_{k+1}$ are replaced by $\{0\}$. In this case the proof given in Section~\ref{sec:bound-analyticity-after} shows that $E^{*}_k$ maps $\mathrm{Ball}_k(4)$ into $K^C_{k+1}$, and maps the pairs $(m,E)\in\mathrm{Ball}_k(4)$ with real $m$ and $E\in K_k$ into $K_{k+1}$; that $\left\|E^{*}_k(m,E)\right\|_{k+1}\le e_k^{\frac{1}{4}-6\varepsilon}$ for $(m,E)\in\mathrm{Ball}_k(4)$; and that, for $(m_a,E_a)$ and $(m_b,E_b)$ in $\mathbb{C}\times K^C_k$, the map $t\mapsto E^{*}_k\left(m_a+t\cdot m_b,\,E_a+t\cdot E_b\right)$ is continuous on the closed set of the $t\in\mathbb{C}$ with $\left(m_a+t\cdot m_b,\,E_a+t\cdot E_b\right)\in\mathrm{Ball}_k(4)$, and analytic in its interior.

Let $(m,\mathscr{E})\in\mathrm{Ball}^x_k(4)$. As shown in Section~\ref{sec:spaces-activities}, the family $\mathscr{E}|_0$ lies in $K^C_k$ and satisfies $\left\|\mathscr{E}|_0\right\|_k\le\|\mathscr{E}\|^x_k$; hence $\left(m,\mathscr{E}|_0\right)\in\mathrm{Ball}_k(4)$. The quantities $\boldsymbol{\alpha}_0(\mathscr{E},X)$ and $\boldsymbol{\alpha}_2(\mathscr{E},X)$ of Section~\ref{sec:local-parts-adjustment}, from which $m(\mathscr{E})$ and $R\,\mathscr{E}$ are formed, are formed with the kernels of $\mathscr{E}(X)$ at $(\mathscr{A},\mathscr{C})=(0,0)$; therefore
\[
m(\mathscr{E})=m\left(\mathscr{E}|_0\right),\qquad\left(R\,\mathscr{E}\right)|_0=R\left(\mathscr{E}|_0\right).
\]
At $\mathscr{C}=0$ the total connection is $\mathscr{T}=A$. In the displays~\eqref{eq:renormalisation-group-56}--\eqref{eq:renormalisation-group-61} the datum $\mathscr{E}$ occurs only through $m'=m+m(\mathscr{E})$ and through $\mathscr{E}^r=R\,\mathscr{E}$, whose second slot is $\mathscr{C}$; consequently, at $\mathscr{C}=0$, the objects that these displays define from the datum $(m,\mathscr{E})$ agree with the objects that the corresponding displays define, in the case without slot, from the datum $\left(m,\mathscr{E}|_0\right)$. Finally, $(0)_L=0$. Therefore, for $Y\in D_{k+1}$ and $\mathscr{A}\in\tilde{R}_{k+1}$,
\begin{align*}
E^{*,x}_k(m,\mathscr{E})(Y;\mathscr{A},0,\cdot\,)&=\sum_{X:\,\bar{X}^L=LY}E^{\#,x}_k(m,\mathscr{E})\left(X;\mathscr{A}_L,0,\psi_L,\chi_L\right)\\
&=\sum_{X:\,\bar{X}^L=LY}E^{\#}_k\left(m,\mathscr{E}|_0\right)\left(X;\mathscr{A}_L,\psi_L,\chi_L\right)=E^{*}_k\left(m,\mathscr{E}|_0\right)(Y;\mathscr{A},\cdot\,),
\end{align*}
that is, $E^{*,x}_k(m,\mathscr{E})|_0=E^{*}_k\left(m,\mathscr{E}|_0\right)$.
\end{proof}

\subsubsection{Lipschitz bound in the data}\label{sec:lipschitz-bound-data}

\begin{proof}[Proof of clause~\eqref{it:step-small-fields-e} of Theorem~\ref{thm:step-small-fields}]
\emph{An estimate by Cauchy's formula.} Let $\mathrm{rad}\ge\frac{3}{2}$, and let $f$ be a map from the closed disc $\left\{t\in\mathbb{C}:|t|\le\mathrm{rad}\right\}$ into $K^{x,C}_{k+1}$ which is continuous, analytic in the open disc, and satisfies $\left\|f(t)\right\|^x_{k+1}\le e_k^{\frac{1}{4}-6\varepsilon}$ for $|t|\le\mathrm{rad}$. Then
\[
\left\|f(1)-f(0)\right\|^x_{k+1}\le\frac{3\cdot e_k^{\frac{1}{4}-6\varepsilon}}{\mathrm{rad}}.
\]
To see this, fix $Y\in D_{k+1}$ and $(\mathscr{A},\mathscr{C})\in\tilde{R}_{k+1}\times\mathfrak{X}_{k+1}$. The evaluation of a kernel at fixed arguments is a continuous linear functional on $K^{x,C}_{k+1}$; so each kernel of $f(t)(Y;\mathscr{A},\mathscr{C})$, at fixed arguments, is continuous in $t$ on the closed disc and analytic in the open disc. Since $\mathrm{rad}>1$, each of these kernels obeys the identity of Section~\ref{sec:notation}, a consequence of Cauchy's integral formula, by which the contour functional $\oint_{\mathrm{rad}}$ returns the difference of the values at $t=1$ and at $t=0$; in other words,
\[
f(1)(Y;\mathscr{A},\mathscr{C})-f(0)(Y;\mathscr{A},\mathscr{C})=\oint_{\mathrm{rad}}f(\cdot\,)(Y;\mathscr{A},\mathscr{C}).
\]
By the triangle inequality for integrals, applied to the norm $\|\cdot\|_{h_{k+1}}$, and since $\left|t\cdot(t-1)\right|\ge\mathrm{rad}\cdot(\mathrm{rad}-1)$ on the circle $|t|=\mathrm{rad}$, whose length is $2\pi\cdot\mathrm{rad}$,
\[
\left\|f(1)(Y;\mathscr{A},\mathscr{C})-f(0)(Y;\mathscr{A},\mathscr{C})\right\|_{h_{k+1}}\le\frac{1}{\mathrm{rad}-1}\cdot\sup_{|t|=\mathrm{rad}}\left\|f(t)(Y;\mathscr{A},\mathscr{C})\right\|_{h_{k+1}}\le\frac{e_k^{\frac{1}{4}-6\varepsilon}}{\mathrm{rad}-1}\cdot\exp\left(-\kappa\cdot d_M(Y)\right).
\]
Multiplying by $\exp\left(\kappa\cdot d_M(Y)\right)$, taking the suprema of clause~\eqref{it:renormalisation-group-1-k6}, and using $\mathrm{rad}-1\ge\mathrm{rad}/3$, which holds because $\mathrm{rad}\ge\frac{3}{2}$, one obtains the estimate.

Let now $(m_1,\mathscr{E}_1)$ and $(m_2,\mathscr{E}_2)$ be elements of $\mathrm{Ball}^x_k(1)$, so that $|m_i|\le w^{\mathrm{m}}_k$ and $\left\|\mathscr{E}_i\right\|^x_k\le w^{\mathrm{E}}_k$ for $i=1,2$ (Section~\ref{sec:spaces-activities}). Decompose
\[
E^{*,x}_k(m_1,\mathscr{E}_1)-E^{*,x}_k(m_2,\mathscr{E}_2)=\left[E^{*,x}_k(m_1,\mathscr{E}_1)-E^{*,x}_k(m_2,\mathscr{E}_1)\right]+\left[E^{*,x}_k(m_2,\mathscr{E}_1)-E^{*,x}_k(m_2,\mathscr{E}_2)\right].
\]

\emph{The first bracket.} If $m_1=m_2$, the first bracket vanishes. Let $m_1\ne m_2$, and put
\[
f(t):=E^{*,x}_k\left(m_2+t\cdot(m_1-m_2),\,\mathscr{E}_1\right),\qquad\mathrm{rad}:=\frac{3\cdot w^{\mathrm{m}}_k}{\left|m_1-m_2\right|}.
\]
As $\left|m_1-m_2\right|\le2\cdot w^{\mathrm{m}}_k$, one has $\mathrm{rad}\ge\frac{3}{2}$. For $|t|\le\mathrm{rad}$,
\[
\left|m_2+t\cdot(m_1-m_2)\right|\le\left|m_2\right|+\mathrm{rad}\cdot\left|m_1-m_2\right|\le w^{\mathrm{m}}_k+3\cdot w^{\mathrm{m}}_k=4\cdot w^{\mathrm{m}}_k,
\]
and $\left\|\mathscr{E}_1\right\|^x_k\le w^{\mathrm{E}}_k\le4\cdot w^{\mathrm{E}}_k$; so the argument of $f$ lies in $\mathrm{Ball}^x_k(4)$. Clause~\eqref{it:step-small-fields-a} of Theorem~\ref{thm:step-small-fields}, applied with $(m_a,\mathscr{E}_a)=(m_2,\mathscr{E}_1)$ and $(m_b,\mathscr{E}_b)=(m_1-m_2,0)$, shows that $f$ is continuous on the closed disc $\left\{|t|\le\mathrm{rad}\right\}$ and analytic in the open disc as a map into $K^{x,C}_{k+1}$, and that $\left\|f(t)\right\|^x_{k+1}\le e_k^{\frac{1}{4}-6\varepsilon}$ there. Since $f(1)-f(0)$ is the first bracket, the estimate gives
\[
\left\|E^{*,x}_k(m_1,\mathscr{E}_1)-E^{*,x}_k(m_2,\mathscr{E}_1)\right\|^x_{k+1}\le\frac{3\cdot e_k^{\frac{1}{4}-6\varepsilon}}{\mathrm{rad}}=e_k^{\frac{1}{4}-6\varepsilon}\cdot\frac{\left|m_1-m_2\right|}{w^{\mathrm{m}}_k}.
\]

\emph{The second bracket.} If $\mathscr{E}_1=\mathscr{E}_2$, the second bracket vanishes. Let $\mathscr{E}_1\ne\mathscr{E}_2$, and put
\[
f(t):=E^{*,x}_k\left(m_2,\,\mathscr{E}_2+t\cdot(\mathscr{E}_1-\mathscr{E}_2)\right),\qquad\mathrm{rad}:=\frac{3\cdot w^{\mathrm{E}}_k}{\left\|\mathscr{E}_1-\mathscr{E}_2\right\|^x_k}.
\]
As $\left\|\mathscr{E}_1-\mathscr{E}_2\right\|^x_k\le2\cdot w^{\mathrm{E}}_k$, one has $\mathrm{rad}\ge\frac{3}{2}$. For $|t|\le\mathrm{rad}$,
\[
\left\|\mathscr{E}_2+t\cdot(\mathscr{E}_1-\mathscr{E}_2)\right\|^x_k\le\left\|\mathscr{E}_2\right\|^x_k+\mathrm{rad}\cdot\left\|\mathscr{E}_1-\mathscr{E}_2\right\|^x_k\le4\cdot w^{\mathrm{E}}_k,
\]
and $\left|m_2\right|\le w^{\mathrm{m}}_k\le4\cdot w^{\mathrm{m}}_k$; so the argument of $f$ lies in $\mathrm{Ball}^x_k(4)$. Clause~\eqref{it:step-small-fields-a} of Theorem~\ref{thm:step-small-fields}, applied with $(m_a,\mathscr{E}_a)=(m_2,\mathscr{E}_2)$ and $(m_b,\mathscr{E}_b)=(0,\mathscr{E}_1-\mathscr{E}_2)$, shows that $f$ satisfies the hypotheses of the estimate. Since $f(1)-f(0)$ is the second bracket,
\[
\left\|E^{*,x}_k(m_2,\mathscr{E}_1)-E^{*,x}_k(m_2,\mathscr{E}_2)\right\|^x_{k+1}\le\frac{3\cdot e_k^{\frac{1}{4}-6\varepsilon}}{\mathrm{rad}}=e_k^{\frac{1}{4}-6\varepsilon}\cdot\frac{\left\|\mathscr{E}_1-\mathscr{E}_2\right\|^x_k}{w^{\mathrm{E}}_k}.
\]
The sum of the bounds for the two brackets is the inequality of clause~\eqref{it:step-small-fields-e} for $E^{*,x}_k$.

\emph{Without slot.} Let $(m_1,E_1)$ and $(m_2,E_2)$ be elements of $\mathrm{Ball}_k(1)$. In the case without slot (Section~\ref{sec:activities-slot}) the preceding argument applies with $E^{*}_k$, $K^C_{k+1}$, $\mathrm{Ball}_k(\cdot)$, $\|\cdot\|_k$ and $\|\cdot\|_{k+1}$ in place of $E^{*,x}_k$, $K^{x,C}_{k+1}$, $\mathrm{Ball}^x_k(\cdot)$, $\|\cdot\|^x_k$ and $\|\cdot\|^x_{k+1}$, and with clause~\eqref{it:step-small-fields-b} of Theorem~\ref{thm:step-small-fields}, which provides for $E^{*}_k$ on $\mathrm{Ball}_k(4)$ the bound $e_k^{\frac{1}{4}-6\varepsilon}$ and the continuity and analyticity in $t$, in place of clause~\eqref{it:step-small-fields-a}. It gives
\[
\left\|E^{*}_k(m_1,E_1)-E^{*}_k(m_2,E_2)\right\|_{k+1}\le e_k^{\frac{1}{4}-6\varepsilon}\cdot\left(\frac{\left|m_1-m_2\right|}{w^{\mathrm{m}}_k}+\frac{\left\|E_1-E_2\right\|_k}{w^{\mathrm{E}}_k}\right).
\]
\end{proof}

\subsection{Symmetries of the step map}\label{sec:symmetries-step-map}

\subsubsection{Invariance and covariance}\label{sec:invariance-covariance}

Let $F$ be a function of some of the arguments $(X,P,s,t;A,\mathscr{C},\mathsf{Z})$ with values in a Grassmann algebra; here $X$ is a polymer, $P$ a set of parameter cubes, $s$ the family of the parameters and $t$ the contour variable; $A$ and $\mathscr{C}$ are the connections in the first and in the second slot, and $\mathsf{Z}$ is the fluctuation variable. The functions defined in Section~\ref{sec:definition-step-map} by~\eqref{eq:renormalisation-group-59}, \eqref{eq:renormalisation-group-60} and~\eqref{eq:renormalisation-group-61} are of this kind. Suppose that each generator is to be replaced by a linear combination of generators, and let $\boldsymbol{\theta}\psi^{\#}$, $\boldsymbol{\theta}\Psi$, $\boldsymbol{\theta} W$ be the families that take the place of $\psi^{\#}=(\psi,\chi)$, $\Psi$, $W$ and of their barred counterparts. The element obtained from $F$ by the complex linear ring homomorphism effecting this replacement is written $F\left(\,\cdot\,;\boldsymbol{\theta}\psi^{\#},\boldsymbol{\theta}\Psi,\boldsymbol{\theta} W\right)$.

\begin{definition}[Lattice symmetries, gauge transformations, charge conjugation, reality]
\begin{enumerate}
\item[(T-g)]\itemmark{T-g}\label{it:lattice-symmetries-gauge-Tg} Let $g=(S,a)\in\mathbb{G}^L_k$. It moves the sites by $x\mapsto gx:=r(S)x+a$, and with them the bonds, the pairs, the polymers and the parameter cubes. A bond function $A$ is carried to the bond function $gA$ given by $(gA)(b):=A\left(g^{-1}b\right)$; the analogous formulas define $g\mathscr{C}$ and $g\mathsf{Z}$, the fluctuation variable $\mathsf{Z}\in V_k$ being itself a bond function. The parameters are permuted according to $(gs)_c:=s_{g^{-1}c}$, and $t$ stays fixed. The generators are transformed by
\begin{align*}
(g\psi)(x)&:=S\cdot\psi\left(g^{-1}x\right),&(g\bar{\psi})(x)&:=\left(S^{-1}\right)^T\cdot\bar{\psi}\left(g^{-1}x\right),\\
(g\chi)(x,x')&:=S\cdot\chi\left(g^{-1}x,g^{-1}x'\right),&(g\bar{\chi})(x,x')&:=\left(S^{-1}\right)^T\cdot\bar{\chi}\left(g^{-1}x,g^{-1}x'\right),
\end{align*}
and by the same formulas on $(\Psi,\bar{\Psi})$ and on $(W,\bar{W})$. The function $F$ is called \emph{invariant} under $g$ if the identity
\[
F\left(gX,gP,gs,t;gA,g\mathscr{C},g\mathsf{Z};g\psi^{\#},g\Psi,gW\right)=F\left(X,P,s,t;A,\mathscr{C},\mathsf{Z};\psi^{\#},\Psi,W\right)
\]
holds for all values of the arguments.
\item[(T-G)]\itemmark{T-G}\label{it:lattice-symmetries-gauge-TG} Assign to every site $x$ of $T_k$ two real numbers $\lambda(x)$ and $\lambda'(x)$, and put $\boldsymbol{\Lambda}:=\lambda+\lambda'$. The generators $\psi$, $\chi$ and their barred counterparts are multiplied by phases,
\begin{align*}
\psi^{\boldsymbol{\Lambda}}(x)&:=\exp\left(i\,e_k\cdot\boldsymbol{\Lambda}(x)\right)\cdot\psi(x),&\bar{\psi}^{\boldsymbol{\Lambda}}(x)&:=\exp\left(-i\,e_k\cdot\boldsymbol{\Lambda}(x)\right)\cdot\bar{\psi}(x),\\
\chi^{\boldsymbol{\Lambda}}(x,x')&:=\exp\left(i\,e_k\cdot\boldsymbol{\Lambda}(x)\right)\cdot\chi(x,x'),&\bar{\chi}^{\boldsymbol{\Lambda}}(x,x')&:=\exp\left(-i\,e_k\cdot\boldsymbol{\Lambda}(x)\right)\cdot\bar{\chi}(x,x').
\end{align*}
The family $(\psi^{\boldsymbol{\Lambda}},\chi^{\boldsymbol{\Lambda}})$, together with $(\bar{\psi}^{\boldsymbol{\Lambda}},\bar{\chi}^{\boldsymbol{\Lambda}})$, is denoted $\psi^{\#,\boldsymbol{\Lambda}}$. On the generators $\Psi$, $\bar{\Psi}$, $W$, $\bar{W}$ the phases are formed with $\boldsymbol{\Lambda}^u:=\boldsymbol{\Lambda}|_{T^0_k}$ (notation of property~\ref{it:renormalisation-group-13}):
\[
\Psi^{\boldsymbol{\Lambda}}(y):=\exp\left(i\,e_k\cdot\boldsymbol{\Lambda}^u(y)\right)\cdot\Psi(y),\qquad\bar{\Psi}^{\boldsymbol{\Lambda}}(y):=\exp\left(-i\,e_k\cdot\boldsymbol{\Lambda}^u(y)\right)\cdot\bar{\Psi}(y),
\]
and $W^{\boldsymbol{\Lambda}}$, $\bar{W}^{\boldsymbol{\Lambda}}$ are defined by the same formulas. The transformation leaves $X$, $P$, $s$, $t$ and $\mathsf{Z}$ unchanged. Let $\mathbf{Y}$ be a polymer. Two pairs $(A,\mathscr{C})$ and $(A^1,\mathscr{C}^1)$ of connections are said to be \emph{related near} $\mathbf{Y}$ if $A^1=A-d\lambda$ and $\mathscr{C}^1=\mathscr{C}-d\lambda'$ on every bond at distance at most $1$ from $\mathbf{Y}$, that is, on every bond each of whose two end points has sup-distance at most $1$ from $\mathbf{Y}$ (Section~\ref{sec:notation}). The pairs of connections at which a function is evaluated range over its domain: this domain is $\tilde{R}^{-}_k\times\mathfrak{X}_k$ for the functions defined in Section~\ref{sec:definition-step-map} by~\eqref{eq:renormalisation-group-56}--\eqref{eq:renormalisation-group-61}, and it is $\tilde{R}_k\times\mathfrak{X}_k$ for the data $\mathscr{E}(X)$, $\mathscr{E}^r(X)$ of Definition~\ref{def:renormalisation-group-1} and for the forms $(A',\mathscr{C})\mapsto\left\langle\bar{\Psi},M_k\left(X,A'+\mathscr{C}\right)\Psi\right\rangle$. A function $F$ is called \emph{invariant} under the gauge transformations \emph{with respect to} $\mathbf{Y}$ if the identity
\[
F\left(\,\cdot\,;A^1,\mathscr{C}^1,\mathsf{Z};\psi^{\#,\boldsymbol{\Lambda}},\Psi^{\boldsymbol{\Lambda}},W^{\boldsymbol{\Lambda}}\right)=F\left(\,\cdot\,;A,\mathscr{C},\mathsf{Z};\psi^{\#},\Psi,W\right)
\]
holds for every such choice of $\lambda$ and $\lambda'$, for all pairs $(A,\mathscr{C})$ and $(A^1,\mathscr{C}^1)$ in the domain of $F$ that are related near $\mathbf{Y}$, and for all values of the remaining arguments $X$, $P$, $s$, $t$, $\mathsf{Z}$ that are admitted for $\mathbf{Y}$, namely those of the set $\mathrm{Ar}(\mathbf{Y})$ defined below. If the only argument of $F$ besides $t$, the connections and $\mathsf{Z}$ is a polymer, as for $\mathscr{E}(X)$, $\mathscr{E}^r(X)$, $\langle\bar{\psi},\psi\rangle_{\square}$ with the polymer $\square$, $E^{(i),*}$, $E^{(i),x}$, $E^{(\le3),x}$, $E^{(\le5),x}$, $\mathsf{K}$, $\mathsf{H}$ and $E^{\#,x}_k$, then $\mathbf{Y}$ is this polymer, and $F$ is called invariant under the gauge transformations if, for every polymer $X$, the function $F(X;\,\cdot\,)$ is invariant with respect to $X$. For a function $G(s,X)$ in the term of a pair $(X,P)$ of~\eqref{eq:renormalisation-group-57}, for the integrand of the contour integral that defines $G(s,X)$, and for this term itself, $\mathbf{Y}$ is the polymer $Y'(X,P)$. If $\mathbf{Y}$ is contained in a polymer $\mathbf{Y}^{+}$, then every bond at distance at most $1$ from $\mathbf{Y}$ is at distance at most $1$ from $\mathbf{Y}^{+}$, since each of its two end points has sup-distance at most $1$ from $\mathbf{Y}$ and therefore from the larger set $\mathbf{Y}^{+}$; hence pairs related near $\mathbf{Y}^{+}$ are related near $\mathbf{Y}$, and a function that is invariant with respect to $\mathbf{Y}$ is invariant with respect to $\mathbf{Y}^{+}$, at the same values of the other arguments (monotonicity).
\item[(T-C)]\itemmark{T-C}\label{it:lattice-symmetries-gauge-TC} The connections and the fluctuation variable change sign, $(A,\mathscr{C},\mathsf{Z})\mapsto(-A,-\mathscr{C},-\mathsf{Z})$, while $X$, $P$, $s$, $t$ stay fixed; the generators are transformed by
\[
\mathbf{C}\,\psi:=\mathsf{C}\cdot\bar{\psi},\qquad\mathbf{C}\,\bar{\psi}:=-\mathsf{C}\cdot\psi,
\]
and by the same formulas on $(\chi,\bar{\chi})$ and on $(\Psi,\bar{\Psi})$. The generators $W$ lie outside the domain of this transformation. A $\mathrm{Gr}[\psi,\chi,\Psi]$-valued function $F$ is called \emph{invariant} under charge conjugation when
\[
F\left(\,\cdot\,;-A,-\mathscr{C},-\mathsf{Z};\mathbf{C}\,\psi^{\#},\mathbf{C}\,\Psi\right)=F\left(\,\cdot\,;A,\mathscr{C},\mathsf{Z};\psi^{\#},\Psi\right).
\]
\item[(T-K)]\itemmark{T-K}\label{it:lattice-symmetries-gauge-TK} For a function $F$ and arguments $(X,P,s,t;A,\mathscr{C},\mathsf{Z})$ consider the element
\[
F\left(X,P,\overline{s},\overline{t};-\overline{A},-\overline{\mathscr{C}},-\mathsf{Z}\right).
\]
Replace each coefficient of this element in the basis of monomials of the generators by its complex conjugate, and apply to the result the complex linear ring homomorphism under which each of $\psi$, $\bar{\psi}$, $\chi$, $\bar{\chi}$, $\Psi$, $\bar{\Psi}$, $W$, $\bar{W}$ is replaced by its product with the matrix $\varkappa$. The element so obtained is denoted $(\mathbf{K} F)(X,P,s,t;A,\mathscr{C},\mathsf{Z})$, and $F$ is called \emph{invariant} under this transformation if $\mathbf{K} F=F$.
\end{enumerate}
\end{definition}

For a family $\mathscr{E}$ as in Definition~\ref{def:renormalisation-group-1}, the clauses~\eqref{it:renormalisation-group-1-k3}, \eqref{it:renormalisation-group-1-k4}, \eqref{it:renormalisation-group-1-k5} and~\eqref{it:renormalisation-group-1-k7}, written as in Section~\ref{sec:lattice-gauge-charge}, are statements of invariance in the present sense. By clause~\eqref{it:renormalisation-group-1-k3}, the maps of~\eqref{it:lattice-symmetries-gauge-Tg} formed with any $g\in\mathbb{G}_k$ leave $\mathscr{E}(X)$ invariant, and by clause~\eqref{it:renormalisation-group-1-k5} so does~\eqref{it:lattice-symmetries-gauge-TC}; clause~\eqref{it:renormalisation-group-1-k4} is the invariance of $\mathscr{E}(X)$ under~\eqref{it:lattice-symmetries-gauge-TG} for the polymer $\mathbf{Y}=X$, with $\tilde{R}_k$ taking the role of the domain $\tilde{R}^{-}_k$; and clause~\eqref{it:renormalisation-group-1-k7} amounts to the identity $\mathbf{K}\,\mathscr{E}(X)=\mathscr{E}(X)$, the matrix $\varkappa$ being real and equal to $\left(\varkappa^{-1}\right)^T$. For every complex number $c$ one has $\mathbf{K}(cF)=\overline{c}\cdot\mathbf{K} F$; moreover, $\mathbf{K}$ respects sums and products and preserves analyticity with respect to $(s,t,A,\mathscr{C})$.

By Lemma~\ref{lem:symmetries-domains}, each of $\tilde{R}^{-}_k$, $\tilde{R}_k$, $\mathfrak{X}_k$, $\tilde{G}_k$ is invariant under $B\mapsto gB$, $B\mapsto-B$ and $B\mapsto-\overline{B}$, which are the actions of~\eqref{it:lattice-symmetries-gauge-Tg}, \eqref{it:lattice-symmetries-gauge-TC} and~\eqref{it:lattice-symmetries-gauge-TK} on connections. By Proposition~\ref{prop:symmetries-reality}, $\mathrm{Box}_k$ is invariant under $\mathsf{Z}\mapsto g\mathsf{Z}$ and under $\mathsf{Z}\mapsto-\mathsf{Z}$. Finally, $s\mapsto gs$ and $s\mapsto\overline{s}$ map the polydisk $\mathrm{Pd}_k$ onto itself, and $t\mapsto\overline{t}$ maps each circle $\{|t|=\rho\}$ onto itself.

With every one of the four transformations comes a ring homomorphism $\boldsymbol{\theta}$ from the Grassmann algebra to itself: $\boldsymbol{\theta}$ sends every generator to a linear combination of generators belonging to the same family, namely to the combination prescribed in~\eqref{it:lattice-symmetries-gauge-Tg}, \eqref{it:lattice-symmetries-gauge-TG}, \eqref{it:lattice-symmetries-gauge-TC}, resp.~\eqref{it:lattice-symmetries-gauge-TK}. For~\eqref{it:lattice-symmetries-gauge-Tg}, \eqref{it:lattice-symmetries-gauge-TG} and~\eqref{it:lattice-symmetries-gauge-TC} the map $\boldsymbol{\theta}$ is complex linear. For~\eqref{it:lattice-symmetries-gauge-TK} it is antilinear: it first conjugates all coefficients in the expansion of an element in monomials of the generators, and then performs the prescribed replacement. Since $W$ lies outside the domain of~\eqref{it:lattice-symmetries-gauge-TC}, the homomorphism $\boldsymbol{\theta}$ of~\eqref{it:lattice-symmetries-gauge-TC} is defined on $\mathrm{Gr}[\psi,\chi,\Psi]$; in the proof of Lemma~\ref{lem:renormalisation-group-2}(c), Step~2, it acts on the algebra obtained from it by adjoining further independent generators $\boldsymbol{\phi}^{\#}$. Write $\mathrm{ar}=(X,P,s,t;A,\mathscr{C},\mathsf{Z})$ for the arguments and $\mathrm{ar}^{\boldsymbol{\theta}}$ for their transforms:
\begin{align*}
\mathrm{ar}^{\boldsymbol{\theta}}&:=\left(gX,gP,gs,t;gA,g\mathscr{C},g\mathsf{Z}\right)&&\text{for~\eqref{it:lattice-symmetries-gauge-Tg}},\\
\mathrm{ar}^{\boldsymbol{\theta}}&:=\left(X,P,s,t;-A,-\mathscr{C},-\mathsf{Z}\right)&&\text{for~\eqref{it:lattice-symmetries-gauge-TC}},\\
\mathrm{ar}^{\boldsymbol{\theta}}&:=\left(X,P,\overline{s},\overline{t};-\overline{A},-\overline{\mathscr{C}},-\mathsf{Z}\right)&&\text{for~\eqref{it:lattice-symmetries-gauge-TK}};
\end{align*}
for~\eqref{it:lattice-symmetries-gauge-TG} with respect to a polymer $\mathbf{Y}$, any $\left(X,P,s,t;A^1,\mathscr{C}^1,\mathsf{Z}\right)$ whose pair $(A^1,\mathscr{C}^1)$ lies in the domain of the function considered and is related near $\mathbf{Y}$ to the pair $(A,\mathscr{C})$ of $\mathrm{ar}$ may serve as $\mathrm{ar}^{\boldsymbol{\theta}}$. Let $\mathrm{Ar}$ denote a set of arguments that contains, together with each of its elements $\mathrm{ar}$, every $\mathrm{ar}^{\boldsymbol{\theta}}$. A function $F$ is called \emph{invariant on} $\mathrm{Ar}$ if
\[
\boldsymbol{\theta}\left(F\left(\mathrm{ar}^{\boldsymbol{\theta}}\right)\right)=F(\mathrm{ar})
\]
for every $\mathrm{ar}\in\mathrm{Ar}$ and every $\mathrm{ar}^{\boldsymbol{\theta}}$; in this form the four definitions of invariance become one. For~\eqref{it:lattice-symmetries-gauge-Tg}, \eqref{it:lattice-symmetries-gauge-TC} and~\eqref{it:lattice-symmetries-gauge-TK}, $\mathrm{Ar}$ is the set of all arguments: $X$ a polymer, $P$ a set of parameter cubes, $s\in\mathrm{Pd}_k$, $t\in\mathbb{C}$ with $|t|\le\rho_1$, $(A,\mathscr{C})$ in the domain of $F$ and $\mathsf{Z}\in\mathrm{Box}_k$. By the preceding paragraph, $\mathrm{ar}\mapsto\mathrm{ar}^{\boldsymbol{\theta}}$ maps this set into itself, and invariance on it is invariance as defined in~\eqref{it:lattice-symmetries-gauge-Tg}, \eqref{it:lattice-symmetries-gauge-TC}, resp.~\eqref{it:lattice-symmetries-gauge-TK}. For~\eqref{it:lattice-symmetries-gauge-TG} with respect to a polymer $\mathbf{Y}$, let $\mathrm{Ar}(\mathbf{Y})$ be the set of the arguments $\mathrm{ar}$ with $P\cap P_0(X)=\emptyset$ and $Y'(X,P)=\mathbf{Y}$, with $s\in\mathrm{Pd}_k$ and $s_c=0$ at every parameter cube $c$ outside $P_0(X)\cup P$, with $|t|\le\rho_1$, $\mathsf{Z}\in\mathrm{Box}_k$ and $(A,\mathscr{C})\in\tilde{R}^{-}_k\times\mathfrak{X}_k$. Here the parameter cubes are those of the localisation at hand: the cubes of the three kinds $\pi_k$, $\pi'_k$, $\pi''_k$ in the definition of $E^{(i),*}$, and the cubes of the kind $\pi_k$, with $s=s''$, in that of $E^{(\le5),x}$ (Section~\ref{sec:localisation-exponent}). On $\mathrm{Ar}(\mathbf{Y})$ the transforms $\mathrm{ar}^{\boldsymbol{\theta}}$ are formed with pairs $(A^1,\mathscr{C}^1)\in\tilde{R}^{-}_k\times\mathfrak{X}_k$. Since~\eqref{it:lattice-symmetries-gauge-TG} changes the connections only, $\mathrm{Ar}(\mathbf{Y})$ contains every $\mathrm{ar}^{\boldsymbol{\theta}}$ of each of its elements, and for a function with the domain $\tilde{R}^{-}_k\times\mathfrak{X}_k$ invariance on $\mathrm{Ar}(\mathbf{Y})$ is invariance with respect to $\mathbf{Y}$. Conditions that concern only arguments on which a function does not depend are void for this function. For a function of a polymer alone, the elements of $\mathrm{Ar}(\mathbf{Y})$ with $P=\emptyset$ are those with $X=\mathbf{Y}$, because $Y'(X,\emptyset)=X$; at them the invariance on $\mathrm{Ar}(\mathbf{Y})$ is the invariance with respect to $X$ required in~\eqref{it:lattice-symmetries-gauge-TG} for such functions, taken at the pairs in $\tilde{R}^{-}_k\times\mathfrak{X}_k$. As $X\subset Y'(X,P)$, a function of a polymer alone that is invariant under~\eqref{it:lattice-symmetries-gauge-TG} is, by the monotonicity stated in~\eqref{it:lattice-symmetries-gauge-TG}, invariant on $\mathrm{Ar}(\mathbf{Y})$ for every polymer $\mathbf{Y}$, provided that its domain contains the pairs at which it is evaluated.

Let $v$ be a family of odd elements, depending on the arguments, whose members are inserted at generators: each element of $v$ is either substituted for a generator or added to it. It is assumed that $\boldsymbol{\theta}$ maps the generators at which $v\left(\mathrm{ar}^{\boldsymbol{\theta}}\right)$ is inserted to linear combinations of the generators at which $v(\mathrm{ar})$ is inserted, and every other generator to a linear combination of other generators. In Section~\ref{sec:definition-step-map}, the families $W^{\#}_k(s'',\mathscr{T})$ and $J^{\#}_k\left(s'',\mathscr{T},\mathscr{Z}(s')\right)\Psi$ are inserted at the generators $\psi^{\#}$, and $T^{(1)}(\mathscr{T})\psi+T^{(2)}(\mathscr{T})\chi$ at the generators $\Psi$. For an index $i$ let $\Phi(i)$ be the generator at which the element $v\left(\mathrm{ar}^{\boldsymbol{\theta}}\right)(i)$ is inserted, for an index $j$ let $\Phi(j)$ be the generator at which $v(\mathrm{ar})(j)$ is inserted, and define the numbers $O_{ij}$ by
\[
\boldsymbol{\theta}\left(\Phi(i)\right)=\sum_jO_{ij}\cdot\Phi(j).
\]
The family $v$ is called \emph{covariant at the index} $i$ on $\mathrm{Ar}$ when
\[
\boldsymbol{\theta}\left(v\left(\mathrm{ar}^{\boldsymbol{\theta}}\right)(i)\right)=\sum_jO_{ij}\cdot v(\mathrm{ar})(j)
\]
for every $\mathrm{ar}\in\mathrm{Ar}$ and every $\mathrm{ar}^{\boldsymbol{\theta}}$, and \emph{covariant} on $\mathrm{Ar}$ if this holds at every index. For~\eqref{it:lattice-symmetries-gauge-Tg}, \eqref{it:lattice-symmetries-gauge-TC} and~\eqref{it:lattice-symmetries-gauge-TK}, $\mathrm{Ar}$ is the set of all arguments; for~\eqref{it:lattice-symmetries-gauge-TG}, covariance is, like invariance, taken with respect to a polymer $\mathbf{Y}$, on the set $\mathrm{Ar}(\mathbf{Y})$. Thus $\boldsymbol{\theta}$ acts on a covariant family exactly as on the generators at which the family is inserted.

The bond function $\mathscr{Z}$ of~\eqref{eq:renormalisation-group-56} is not inserted at generators; it enters through the connection argument. Its value $\mathscr{Z}(s')$ depends on the fluctuation variable and is written $\mathscr{Z}(s')[\mathsf{Z}]$ when this dependence is to be shown. Let $A^{\boldsymbol{\theta}}$, $t^{\boldsymbol{\theta}}$, $s'^{\boldsymbol{\theta}}$, $\mathsf{Z}^{\boldsymbol{\theta}}$ denote the entries of $\mathrm{ar}^{\boldsymbol{\theta}}$ in the positions of $A$, $t$, $s'$, $\mathsf{Z}$, and let $B\mapsto B^{\boldsymbol{\theta}}$ be the action on connections, that is, $B^{\boldsymbol{\theta}}:=gB$, $-B$, $-\overline{B}$ for~\eqref{it:lattice-symmetries-gauge-Tg}, \eqref{it:lattice-symmetries-gauge-TC}, \eqref{it:lattice-symmetries-gauge-TK}. For~\eqref{it:lattice-symmetries-gauge-TG}, let $\left(A+t\cdot\mathscr{Z}(s')[\mathsf{Z}]\right)^{\boldsymbol{\theta}}$ stand for any connection that coincides with $A+t\cdot\mathscr{Z}(s')[\mathsf{Z}]-d\lambda$ on every bond at distance at most $1$ from $\mathbf{Y}$. The function $\mathscr{Z}$ is called \emph{covariant} if
\[
A^{\boldsymbol{\theta}}+t^{\boldsymbol{\theta}}\cdot\mathscr{Z}\left(s'^{\boldsymbol{\theta}}\right)\left[\mathsf{Z}^{\boldsymbol{\theta}}\right]=\left(A+t\cdot\mathscr{Z}(s')[\mathsf{Z}]\right)^{\boldsymbol{\theta}}
\]
for all arguments, where for~\eqref{it:lattice-symmetries-gauge-TG} the identity means that the left side is one of the connections just described. For~\eqref{it:lattice-symmetries-gauge-TG} this is always the case: $s$, $t$ and $\mathsf{Z}$ are not changed, so neither is $\mathscr{Z}(s')[\mathsf{Z}]$, and since $A^1=A-d\lambda$ on the bonds in question, the left side $A^1+t\cdot\mathscr{Z}(s')[\mathsf{Z}]$ coincides there with $A+t\cdot\mathscr{Z}(s')[\mathsf{Z}]-d\lambda$.

\begin{lemma}[Operations that preserve invariance]
Fix one of the transformations~\eqref{it:lattice-symmetries-gauge-Tg}, \eqref{it:lattice-symmetries-gauge-TG}, \eqref{it:lattice-symmetries-gauge-TC}, \eqref{it:lattice-symmetries-gauge-TK} and a set $\mathrm{Ar}$ as in the text before the lemma; for~\eqref{it:lattice-symmetries-gauge-TG} fix in addition a polymer $\mathbf{Y}$, and let $\mathrm{Ar}=\mathrm{Ar}(\mathbf{Y})$. Invariant and covariant mean invariant and covariant on $\mathrm{Ar}$. The following operations, applied to invariant functions, yield functions that are invariant on $\mathrm{Ar}$, resp., after an integration or a summation, on the set of the arguments that remain.
\begin{enumerate}
\item[(Pr1)]\itemmark{Pr1}\label{it:operations-preserve-invariance-Pr1} Finite sums, finite products and exponentials, and sums of series that converge in norm; the substitution of a covariant family $v$ into an invariant function $F$; and the substitution of $A+t\cdot\mathscr{Z}(s')[\mathsf{Z}]$, with a covariant $\mathscr{Z}$, for the connection argument of a function $F$ that is invariant, for~\eqref{it:lattice-symmetries-gauge-TG} with respect to $\mathbf{Y}$, when its pairs of connections range over a domain containing the pair $\left(A+t\cdot\mathscr{Z}(s')[\mathsf{Z}],\mathscr{C}\right)$ for every $\mathrm{ar}\in\mathrm{Ar}$. For $F=\mathscr{E}^r(X)$ this domain is $\tilde{R}_k\times\mathfrak{X}_k$; if the pairs of the elements of $\mathrm{Ar}$ lie in $\tilde{R}^{-}_k\times\mathfrak{X}_k$, as for the functions of Section~\ref{sec:definition-step-map}, it contains these pairs by part~\eqref{it:bounds-analyticity-q1} of Lemma~\ref{lem:bounds-analyticity}, since $|t|\le\rho_1$. The rule holds at every single pair $\left(\mathrm{ar},\mathrm{ar}^{\boldsymbol{\theta}}\right)$: it uses only the invariance of the functions at this pair and the covariance of $v$, resp.\ of $\mathscr{Z}$, at this pair, that of $v$ at the indices of the generators that occur in $F\left(\mathrm{ar}^{\boldsymbol{\theta}}\right)$.
\item[(Pr2)]\itemmark{Pr2}\label{it:operations-preserve-invariance-Pr2} The contour integral $\oint_{\rho}$ in the variable $t$, provided that, for every $\mathrm{ar}\in\mathrm{Ar}$, all the arguments that differ from $\mathrm{ar}$ in the entry $t$ only and have this entry on the circle $\{|t|=\rho\}$ belong to $\mathrm{Ar}$. This holds for $\rho=\rho_1,\rho_2,\rho_4$, since the sets $\mathrm{Ar}$ of the text before the lemma restrict $t$ only by $|t|\le\rho_1$, and $\rho_4\le\rho_2\le\rho_1$ (Section~\ref{sec:notation}).
\item[(Pr3)]\itemmark{Pr3}\label{it:operations-preserve-invariance-Pr3} The operation of~\eqref{eq:renormalisation-group-57} that assigns to a function $G(s,X)$ and to a pair $(X,P)$ the element
\[
\int_{[0,1]^P}\mathrm{d}s_P\,\frac{\partial}{\partial s_P}\left[G(s,X)\right]_{s=1\text{ on }P_0(X),\ s=0\text{ outside }P_0(X)\cup P}.
\]
It reads $G(\cdot,X)$ only at the families $s$ with $s_c=1$ on $P_0(X)$, $s_c=0$ outside $P_0(X)\cup P$ and $(s_c)_{c\in P}$ in the polydisk $\left\{|s_c|\le M^{\alpha_0},\ c\in P\right\}$, a neighbourhood of $[0,1]^P$, and it differentiates in the variables $s_c$, $c\in P$, only. All these families vanish outside $P_0(X)\cup P$; so for~\eqref{it:lattice-symmetries-gauge-TG} and a pair $(X,P)$ with $Y'(X,P)=\mathbf{Y}$ the arguments at which $G(\cdot,X)$ is read belong to $\mathrm{Ar}(\mathbf{Y})$.
\item[(Pr4)]\itemmark{Pr4}\label{it:operations-preserve-invariance-Pr4} Summation over polymers and over sets of parameter cubes subject to inclusions, to connectedness, to the condition $Y'(X,P)=\mathbf{Y}$ or to the condition $\bigcup_iX_i=X$, possibly weighted by the coefficients $\rho^T$; these coefficients are real numbers determined by the incidence relations among the polymers alone. The rule holds term by term. For~\eqref{it:lattice-symmetries-gauge-TG}, the sum over the pairs $(X,P)$ with $Y'(X,P)=\mathbf{Y}$ of terms that are invariant with respect to $\mathbf{Y}$ is a function of the polymer $\mathbf{Y}$ that is invariant with respect to $\mathbf{Y}$; and in the sums over polymers $X_i$ contained in $\mathbf{Y}$, the summands and their factors that are invariant with respect to these smaller polymers are invariant with respect to $\mathbf{Y}$, by the monotonicity stated in~\eqref{it:lattice-symmetries-gauge-TG}.
\item[(Pr5)]\itemmark{Pr5}\label{it:operations-preserve-invariance-Pr5} The integration $\int\mathrm{d}\hat{\mu}(\mathsf{Z})\,\mathrm{d}\mu_I(W)$ of a function with values in $\mathrm{Gr}[\psi,\chi,W]$, provided that the action of the transformation on $\mathsf{Z}$ leaves $\hat{\mu}$ invariant and its action on the generators $W$, $\bar{W}$ has the form $W\mapsto UW$, $\bar{W}\mapsto\left(U^{-1}\right)^T\bar{W}$ with an invertible matrix $U$, and provided that, for every $\mathrm{ar}\in\mathrm{Ar}$, all the arguments that differ from $\mathrm{ar}$ in the entry $\mathsf{Z}$ only, with $\mathsf{Z}\in\mathrm{Box}_k$, belong to $\mathrm{Ar}$.
\end{enumerate}
\end{lemma}

\begin{proof}
Throughout the proof, $\mathrm{ar}$ denotes an element of $\mathrm{Ar}$ and $\mathrm{ar}^{\boldsymbol{\theta}}$ one of its transforms.

\emph{Proof of~\eqref{it:operations-preserve-invariance-Pr1}.} The proof is carried out at one such pair $\left(\mathrm{ar},\mathrm{ar}^{\boldsymbol{\theta}}\right)$: every identity in it involves only the values at $\mathrm{ar}$ and at $\mathrm{ar}^{\boldsymbol{\theta}}$, so that~\eqref{it:operations-preserve-invariance-Pr1} holds on every set $\mathrm{Ar}$. The algebra is finite-dimensional, and $\boldsymbol{\theta}$ is a ring homomorphism that is linear or antilinear over $\mathbb{C}$; hence $\boldsymbol{\theta}$ is continuous. As a ring homomorphism it respects sums and products; by continuity it carries limits to limits, and in particular $\boldsymbol{\theta}\left(\exp G\right)=\exp\left(\boldsymbol{\theta} G\right)$, the exponential series having rational coefficients. For invariant functions $F_1,F_2,\dots$ whose series converges in norm one obtains
\[
\boldsymbol{\theta}\left(\sum_nF_n\left(\mathrm{ar}^{\boldsymbol{\theta}}\right)\right)=\sum_n\boldsymbol{\theta}\left(F_n\left(\mathrm{ar}^{\boldsymbol{\theta}}\right)\right)=\sum_nF_n(\mathrm{ar}),
\]
and finite sums, finite products and exponentials are dealt with in the same manner.

Let $v$ be covariant, and denote by $\mathrm{sb}_v$ the complex linear ring homomorphism that substitutes each element of $v$ for the generator at which it is inserted, or adds it to that generator, and fixes every other generator. If $\Phi(i)$ is a generator at which $v\left(\mathrm{ar}^{\boldsymbol{\theta}}\right)$ is inserted, covariance at $i$ and the linearity of $\mathrm{sb}_{v(\mathrm{ar})}$ give
\[
\boldsymbol{\theta}\left(\mathrm{sb}_{v(\mathrm{ar}^{\boldsymbol{\theta}})}\Phi(i)\right)=\sum_jO_{ij}\cdot\mathrm{sb}_{v(\mathrm{ar})}\Phi(j)=\mathrm{sb}_{v(\mathrm{ar})}\left(\boldsymbol{\theta}\left(\Phi(i)\right)\right);
\]
on every other generator, $\boldsymbol{\theta}\circ\mathrm{sb}_{v(\mathrm{ar}^{\boldsymbol{\theta}})}$ and $\mathrm{sb}_{v(\mathrm{ar})}\circ\boldsymbol{\theta}$ both coincide with $\boldsymbol{\theta}$. The two composites are ring homomorphisms of the same kind, linear when $\boldsymbol{\theta}$ is linear and antilinear when $\boldsymbol{\theta}$ is antilinear; as they agree on the generators, they are equal. Applied to an invariant function $F$, this yields
\[
\boldsymbol{\theta}\left(\mathrm{sb}_{v(\mathrm{ar}^{\boldsymbol{\theta}})}F\left(\mathrm{ar}^{\boldsymbol{\theta}}\right)\right)=\mathrm{sb}_{v(\mathrm{ar})}\,\boldsymbol{\theta}\left(F\left(\mathrm{ar}^{\boldsymbol{\theta}}\right)\right)=\mathrm{sb}_{v(\mathrm{ar})}F(\mathrm{ar}),
\]
so that $\mathrm{ar}\mapsto\mathrm{sb}_{v(\mathrm{ar})}F(\mathrm{ar})$ is invariant.

Covariance of $v$ is needed only at those indices $i$ for which the generator $\Phi(i)$ occurs in $F\left(\mathrm{ar}^{\boldsymbol{\theta}}\right)$. Indeed, $S$, $\mathsf{C}$ and $\varkappa$ are monomial matrices (Proposition~\ref{prop:symmetries-extraction}(b)); so is $\left(S^{-1}\right)^T$, and the factors in~\eqref{it:lattice-symmetries-gauge-TG} are phases. All these matrices are invertible and $x\mapsto g^{-1}x$ is a bijection; hence $\boldsymbol{\theta}$ sends each generator to a non-zero multiple of a single generator, and distinct generators to multiples of distinct generators. Consequently $\boldsymbol{\theta}$ maps every monomial to a non-zero multiple of a monomial, and distinct monomials to multiples of distinct ones. No cancellation can therefore occur in $F(\mathrm{ar})=\boldsymbol{\theta}\left(F\left(\mathrm{ar}^{\boldsymbol{\theta}}\right)\right)$, and the image under $\boldsymbol{\theta}$ of the span of the monomials formed from the generators occurring in $F\left(\mathrm{ar}^{\boldsymbol{\theta}}\right)$ is the span of the monomials formed from those occurring in $F(\mathrm{ar})$; as for the indices, under~\eqref{it:lattice-symmetries-gauge-TG}, \eqref{it:lattice-symmetries-gauge-TC} and~\eqref{it:lattice-symmetries-gauge-TK} the sites and pairs stay in place, and under~\eqref{it:lattice-symmetries-gauge-Tg} those of $gX$ go to those of $X$. The homomorphisms $\boldsymbol{\theta}\circ\mathrm{sb}_{v(\mathrm{ar}^{\boldsymbol{\theta}})}$ and $\mathrm{sb}_{v(\mathrm{ar})}\circ\boldsymbol{\theta}$ need to agree only on the first of these two spans, which is a subalgebra and contains $F\left(\mathrm{ar}^{\boldsymbol{\theta}}\right)$. Being ring homomorphisms of the same kind, they agree there as soon as they agree on the generators occurring in $F\left(\mathrm{ar}^{\boldsymbol{\theta}}\right)$. This is the case if covariance of $v$ holds at the indices of those among these generators at which $v\left(\mathrm{ar}^{\boldsymbol{\theta}}\right)$ is inserted, since on each of the other generators both homomorphisms act as $\boldsymbol{\theta}$. For instance, $\mathscr{E}^r(X)$, $\langle\bar{\psi},\psi\rangle_{\square}$ with an $M$-cube $\square$ of $X$, and $E^{(\le3),x}(X)$ involve the generators $\psi$ only at sites of $X$ and the generators $\chi$ only at pairs that are admissible for $X$ (clause~\eqref{it:renormalisation-group-1-k2} of Definition~\ref{def:renormalisation-group-1}; the final paragraph in the proof of Lemma~\ref{lem:renormalisation-group-1}). For a substitution into one of these functions it therefore suffices that $v$ be covariant at the indices belonging to the sites of the polymer in $\mathrm{ar}^{\boldsymbol{\theta}}$ and to the pairs admissible for it.

Now let $\mathscr{Z}$ be covariant, and let $F$ be invariant, for~\eqref{it:lattice-symmetries-gauge-TG} with respect to $\mathbf{Y}$, when its pairs of connections range over a domain $D$ that contains the pair $\left(A+t\cdot\mathscr{Z}(s')[\mathsf{Z}],\mathscr{C}\right)$ for every $\mathrm{ar}\in\mathrm{Ar}$. For $F=\mathscr{E}^r(X)$, and also for $F=\mathscr{E}(X)$, the domain $D$ is $\tilde{R}_k\times\mathfrak{X}_k$ (clauses~\eqref{it:renormalisation-group-1-k3}, \eqref{it:renormalisation-group-1-k4}, \eqref{it:renormalisation-group-1-k5} and~\eqref{it:renormalisation-group-1-k7} of Definition~\ref{def:renormalisation-group-1}), which the actions of~\eqref{it:lattice-symmetries-gauge-Tg}, \eqref{it:lattice-symmetries-gauge-TC} and~\eqref{it:lattice-symmetries-gauge-TK} on connections map into itself (Lemma~\ref{lem:symmetries-domains}). If the pairs of the elements of $\mathrm{Ar}$ lie in $\tilde{R}^{-}_k\times\mathfrak{X}_k$, as for the functions of Section~\ref{sec:definition-step-map}, then every $\mathrm{ar}\in\mathrm{Ar}$ has $(A,\mathscr{C})\in\tilde{R}^{-}_k\times\mathfrak{X}_k$, $s\in\mathrm{Pd}_k$, $\mathsf{Z}\in\mathrm{Box}_k$ and $|t|\le\rho_1$, and part~\eqref{it:bounds-analyticity-q1} of Lemma~\ref{lem:bounds-analyticity} shows that $\left(A+t\cdot\mathscr{Z}(s')[\mathsf{Z}],\mathscr{C}\right)$ lies in $\tilde{R}_k\times\mathfrak{X}_k$; the same holds at $\mathrm{ar}^{\boldsymbol{\theta}}$, which lies in $\mathrm{Ar}$. Substitute $A+t\cdot\mathscr{Z}(s')[\mathsf{Z}]$ for the connection argument of $F$. The transformation of the arguments then acts on this connection as a whole:
\[
F\left(A^{\boldsymbol{\theta}}+t^{\boldsymbol{\theta}}\cdot\mathscr{Z}\left(s'^{\boldsymbol{\theta}}\right)\left[\mathsf{Z}^{\boldsymbol{\theta}}\right]\right)=F\left(\left(A+t\cdot\mathscr{Z}(s')[\mathsf{Z}]\right)^{\boldsymbol{\theta}}\right),
\]
where, for~\eqref{it:lattice-symmetries-gauge-TG}, the arguments of $F$ on the right are related to those of $F\left(A+t\cdot\mathscr{Z}(s')[\mathsf{Z}]\right)$ as in~\eqref{it:lattice-symmetries-gauge-TG}: the pairs $\left(A+t\cdot\mathscr{Z}(s')[\mathsf{Z}],\mathscr{C}\right)$ and $\left(A^1+t\cdot\mathscr{Z}(s')[\mathsf{Z}],\mathscr{C}^1\right)$ are related near $\mathbf{Y}$ by the same $\lambda$ and $\lambda'$, because $A^1=A-d\lambda$ and $\mathscr{C}^1=\mathscr{C}-d\lambda'$ on the bonds at distance at most $1$ from $\mathbf{Y}$, and therefore they are related near every polymer contained in $\mathbf{Y}$. In each case both pairs at which $F$ is evaluated lie in $D$, because $\mathrm{ar}$ and $\mathrm{ar}^{\boldsymbol{\theta}}$ lie in $\mathrm{Ar}$. By the invariance of $F$, the map $\boldsymbol{\theta}$ carries the value of the substituted function at $\mathrm{ar}^{\boldsymbol{\theta}}$ to its value at $\mathrm{ar}$, and so this function is invariant on $\mathrm{Ar}$. For $D=\tilde{R}_k\times\mathfrak{X}_k$, only the properties $(A,\mathscr{C})\in\tilde{R}^{-}_k\times\mathfrak{X}_k$, $s\in\mathrm{Pd}_k$, $\mathsf{Z}\in\mathrm{Box}_k$ and $|t|\le\rho_1$ of $\mathrm{ar}$ have been used. In particular, for~\eqref{it:lattice-symmetries-gauge-TG}, if $F$ is invariant with respect to $\mathbf{Y}$ at all pairs in $\tilde{R}_k\times\mathfrak{X}_k$, then the identity of~\eqref{it:lattice-symmetries-gauge-TG} holds for the function $(s,t;A,\mathscr{C},\mathsf{Z})\mapsto F\left(A+t\cdot\mathscr{Z}(s')[\mathsf{Z}],\mathscr{C}\right)$ at all pairs in $\tilde{R}^{-}_k\times\mathfrak{X}_k$ that are related near $\mathbf{Y}$, for all $s\in\mathrm{Pd}_k$, all $t$ with $|t|\le\rho_1$ and all $\mathsf{Z}\in\mathrm{Box}_k$.

\emph{Proof of~\eqref{it:operations-preserve-invariance-Pr2}.} For~\eqref{it:lattice-symmetries-gauge-Tg}, \eqref{it:lattice-symmetries-gauge-TG} and~\eqref{it:lattice-symmetries-gauge-TC} the variable $t$ is not changed and $\boldsymbol{\theta}$ is linear and continuous; so $\boldsymbol{\theta}$ can be taken under the contour integral, and the assertion follows from the invariance of the integrand at the arguments with $t$ on the circle $\{|t|=\rho\}$, which belong to $\mathrm{Ar}$ by the hypothesis of the rule. Consider~\eqref{it:lattice-symmetries-gauge-TK}, for which $\boldsymbol{\theta}$ is antilinear. Parametrising the circle by $t=\rho\exp(i\mathfrak{a})$, $0\le\mathfrak{a}\le2\pi$, one has
\[
\oint_{\rho}f=\frac{1}{2\pi}\int_0^{2\pi}\mathrm{d}\mathfrak{a}\,f\left(\rho\exp(i\mathfrak{a})\right)\cdot\left(\rho\exp(i\mathfrak{a})-1\right)^{-1}.
\]
Being real-linear and continuous, $\boldsymbol{\theta}$ can be interchanged with the integral in $\mathfrak{a}$; on the scalar factor it acts by complex conjugation, which gives $\left(\rho\exp(-i\mathfrak{a})-1\right)^{-1}$. Let $F$ be invariant, split its arguments as $\mathrm{ar}=(\mathrm{ar}_0;t)$ with $\mathrm{ar}_0$ the arguments other than $t$, and let $\mathrm{ar}_0^{\boldsymbol{\theta}}$ be the part of $\mathrm{ar}^{\boldsymbol{\theta}}$ in the same positions. By the hypothesis of the rule, all the points $(\mathrm{ar}_0;t')$ and $(\mathrm{ar}_0^{\boldsymbol{\theta}};t')$ with $|t'|=\rho$ belong to $\mathrm{Ar}$, the latter because $\mathrm{ar}^{\boldsymbol{\theta}}\in\mathrm{Ar}$. Apply the invariance of $F$ at the point $t=\rho\exp(-i\mathfrak{a})$, whose conjugate is $\rho\exp(i\mathfrak{a})$, and then replace $\mathfrak{a}$ by $-\mathfrak{a}$, which does not change the integral of a $2\pi$-periodic function:
\begin{align*}
\boldsymbol{\theta}\left(\oint_{\rho}F\left(\mathrm{ar}_0^{\boldsymbol{\theta}};\cdot\,\right)\right)&=\frac{1}{2\pi}\int_0^{2\pi}\mathrm{d}\mathfrak{a}\,\boldsymbol{\theta}\left(F\left(\mathrm{ar}_0^{\boldsymbol{\theta}};\rho\exp(i\mathfrak{a})\right)\right)\cdot\left(\rho\exp(-i\mathfrak{a})-1\right)^{-1}\\
&=\frac{1}{2\pi}\int_0^{2\pi}\mathrm{d}\mathfrak{a}\,F\left(\mathrm{ar}_0;\rho\exp(-i\mathfrak{a})\right)\cdot\left(\rho\exp(-i\mathfrak{a})-1\right)^{-1}=\oint_{\rho}F\left(\mathrm{ar}_0;\cdot\,\right).
\end{align*}
For a scalar function $f$ the same computation shows $\overline{\oint_{\rho}f}=\oint_{\rho}h$ with $h(t):=\overline{f(\overline{t})}$: the reflection $t\mapsto\overline{t}$ maps the circle $\{|t|=\rho\}$ onto itself and reverses its orientation.

\emph{Proof of~\eqref{it:operations-preserve-invariance-Pr3}.} Consider~\eqref{it:lattice-symmetries-gauge-Tg}. By Proposition~\ref{prop:symmetries-extraction}(a), $g$ maps the parameter cubes contained in $X$ onto those contained in $gX$, and $P$ onto $gP$; hence the change of variables $s\mapsto gs$ turns the operation belonging to $(X,P)$ into the operation belonging to $(gX,gP)$, and the linear map $\boldsymbol{\theta}$ can be interchanged with the derivatives and with the integrals in $s$. Under~\eqref{it:lattice-symmetries-gauge-TC} the map $\boldsymbol{\theta}$ is linear and $s$ is not changed. The same holds under~\eqref{it:lattice-symmetries-gauge-TG}; there let $Y'(X,P)=\mathbf{Y}$, let $(A,\mathscr{C})$ and $(A^1,\mathscr{C}^1)$ be pairs in $\tilde{R}^{-}_k\times\mathfrak{X}_k$ that are related near $\mathbf{Y}$, and fix $\mathsf{Z}\in\mathrm{Box}_k$ and $t$ with $|t|\le\rho_1$. The bracket in~\eqref{eq:renormalisation-group-57} sets $s_c=1$ for $c\in P_0(X)$ and $s_c=0$ for every parameter cube $c$ outside $P_0(X)\cup P$, and only the variables $s_c$ with $c\in P$ vary. For $(s_c)_{c\in P}$ in the polydisk $\left\{|s_c|\le M^{\alpha_0},\ c\in P\right\}$, which contains $[0,1]^P$ in its interior because $M^{\alpha_0}>1$ by~\eqref{eq:renormalisation-group-30}, the family $s$ vanishes outside $P_0(X)\cup P$. So the argument lies in $\mathrm{Ar}(\mathbf{Y})$, and the invariance of $G(\cdot,X)$ on $\mathrm{Ar}(\mathbf{Y})$ shows that $\boldsymbol{\theta}$ carries the value of $G(s,X)$ at the pair $(A^1,\mathscr{C}^1)$ to its value at $(A,\mathscr{C})$. This identity holds for all $(s_c)_{c\in P}$ in the polydisk; hence the derivatives $\partial/\partial s_P$ of its two sides coincide on $[0,1]^P$, and the linear and continuous map $\boldsymbol{\theta}$ commutes with these derivatives and with the integral over $[0,1]^P$. Under~\eqref{it:lattice-symmetries-gauge-TK}, differentiation in $s$ turns $s\mapsto\overline{f(\overline{s})}$ into $s\mapsto\overline{(\partial_sf)(\overline{s})}$; moreover the cube $[0,1]^P$ over which one integrates is real, and so are the values $0$ and $1$ given to the remaining parameters.

\emph{Proof of~\eqref{it:operations-preserve-invariance-Pr4}.} The coefficients $\rho^T$ are real, so that $\boldsymbol{\theta}\left(\rho^T\cdot F\right)=\rho^T\cdot\boldsymbol{\theta}(F)$ also for the antilinear map $\boldsymbol{\theta}$ of~\eqref{it:lattice-symmetries-gauge-TK}; and they are determined by the incidence relations among the polymers. Under~\eqref{it:lattice-symmetries-gauge-TG}, \eqref{it:lattice-symmetries-gauge-TC} and~\eqref{it:lattice-symmetries-gauge-TK} neither polymers nor parameter cubes are moved. Consider~\eqref{it:lattice-symmetries-gauge-Tg}. By Proposition~\ref{prop:symmetries-extraction}(a), $g$ maps polymers onto polymers and parameter cubes onto parameter cubes of the same kind, and it preserves inclusions and connectedness, hence also the incidence relations; consequently it carries the pairs $(X,P)$ with $Y'(X,P)=\mathbf{Y}$ onto the pairs with $Y'(X,P)=g\mathbf{Y}$, and the families $(X_i)$ with union $X$ onto the families with union $gX$. Consider~\eqref{it:lattice-symmetries-gauge-TG}. By the monotonicity stated in~\eqref{it:lattice-symmetries-gauge-TG}, a function that is invariant with respect to a polymer contained in $\mathbf{Y}$ is invariant with respect to $\mathbf{Y}$. So in the sums of~\eqref{eq:renormalisation-group-61} over polymers $X_i$ contained in $\mathbf{Y}$, every factor of a summand that is invariant with respect to one of the polymers $X_i$ is invariant with respect to $\mathbf{Y}$; by~\eqref{it:operations-preserve-invariance-Pr1} so is every summand, and hence so is the sum. Moreover, the terms of~\eqref{eq:renormalisation-group-57} with $Y'(X,P)=\mathbf{Y}$ are finite in number; if each of them is invariant with respect to $\mathbf{Y}$, then their sum, by the additivity of $\boldsymbol{\theta}$, is a function of the polymer $\mathbf{Y}$ that is invariant with respect to $\mathbf{Y}$.

\emph{Proof of~\eqref{it:operations-preserve-invariance-Pr5}.} The rule, like Lemma~\ref{lem:substitutions-integral}(b), concerns~\eqref{it:lattice-symmetries-gauge-Tg}, \eqref{it:lattice-symmetries-gauge-TG} and~\eqref{it:lattice-symmetries-gauge-TK} only, since $W$ lies outside the domain of~\eqref{it:lattice-symmetries-gauge-TC}. For these three transformations and for $F$ with values in $\mathrm{Gr}[\psi,\chi,W]$, \eqref{eq:renormalisation-group-42} and Lemma~\ref{lem:substitutions-integral}(b) yield
\[
\int\mathrm{d}\mu_I(W)\,\boldsymbol{\theta}(F)=\boldsymbol{\theta}\left(\int\mathrm{d}\mu_I(W)\,F\right).
\]
As $\hat{\mu}$ is a real measure and $\boldsymbol{\theta}$ is real-linear and continuous, $\boldsymbol{\theta}$ may also be taken inside the integral over $\mathsf{Z}$. For the antilinear $\boldsymbol{\theta}$ of~\eqref{it:lattice-symmetries-gauge-TK} these two interchanges rest on the fact that $\mu_I$ has real moments and that $\hat{\mu}$ is real. Split the arguments as $\mathrm{ar}=(\mathrm{ar}_0;\mathsf{Z})$ with $\mathrm{ar}_0$ the arguments other than $\mathsf{Z}$, let $\mathrm{ar}_0^{\boldsymbol{\theta}}$ be the part of $\mathrm{ar}^{\boldsymbol{\theta}}$ in the same positions, and let $\mathsf{Z}\mapsto\mathsf{Z}^{\boldsymbol{\theta}}$ be the action on the fluctuation variable. By the hypothesis of the rule, all the points $(\mathrm{ar}_0;\mathsf{Z})$ and $(\mathrm{ar}_0^{\boldsymbol{\theta}};\mathsf{Z})$ with $\mathsf{Z}\in\mathrm{Box}_k$ belong to $\mathrm{Ar}$, the latter because $\mathrm{ar}^{\boldsymbol{\theta}}\in\mathrm{Ar}$. The action $\mathsf{Z}\mapsto\mathsf{Z}^{\boldsymbol{\theta}}$ preserves $\hat{\mu}$: for~\eqref{it:lattice-symmetries-gauge-Tg} and~\eqref{it:lattice-symmetries-gauge-TK} by Proposition~\ref{prop:symmetries-reality}, while~\eqref{it:lattice-symmetries-gauge-TG} leaves $\mathsf{Z}$ unchanged. For invariant $F$, the change of variables $\mathsf{Z}\mapsto\mathsf{Z}^{\boldsymbol{\theta}}$ in the integral gives
\begin{align*}
\boldsymbol{\theta}\left(\int\mathrm{d}\hat{\mu}(\mathsf{Z})\,\mathrm{d}\mu_I(W)\,F\left(\mathrm{ar}_0^{\boldsymbol{\theta}};\mathsf{Z}\right)\right)&=\int\mathrm{d}\hat{\mu}(\mathsf{Z})\,\mathrm{d}\mu_I(W)\,\boldsymbol{\theta}\left(F\left(\mathrm{ar}_0^{\boldsymbol{\theta}};\mathsf{Z}\right)\right)\\
&=\int\mathrm{d}\hat{\mu}(\mathsf{Z})\,\mathrm{d}\mu_I(W)\,\boldsymbol{\theta}\left(F\left(\mathrm{ar}_0^{\boldsymbol{\theta}};\mathsf{Z}^{\boldsymbol{\theta}}\right)\right)\\
&=\int\mathrm{d}\hat{\mu}(\mathsf{Z})\,\mathrm{d}\mu_I(W)\,F\left(\mathrm{ar}_0;\mathsf{Z}\right).
\end{align*}
\end{proof}

\begin{lemma}[Substitutions in the integral over $W$]\label{lem:substitutions-integral}
\begin{enumerate}
\item[(a)] Let $\mathrm{Gr}[\psi,\chi,\boldsymbol{\phi}]$ be the Grassmann algebra over $\mathbb{C}$ whose generators are the $\psi(\xi)$, the $\chi(\boldsymbol{\zeta})$ and further independent generators $\boldsymbol{\phi}(i)$, where the index $i$ ranges over a finite set. Let $\mathrm{sb}_1$ and $\mathrm{sb}_2$ be homomorphisms of algebras from $\mathrm{Gr}[\psi,\chi,\boldsymbol{\phi}]$ to $\mathrm{Gr}[\psi,\chi,W]$ that leave every generator $\psi(\xi)$ and every generator $\chi(\boldsymbol{\zeta})$ fixed and send every $\boldsymbol{\phi}(i)$ into the linear span of the generators $W(u)$, the barred as well as the unbarred ones. Assume that
\[
\int\mathrm{d}\mu_I(W)\,\mathrm{sb}_1\left(\boldsymbol{\phi}(i)\right)\cdot\mathrm{sb}_1\left(\boldsymbol{\phi}(i')\right)=\int\mathrm{d}\mu_I(W)\,\mathrm{sb}_2\left(\boldsymbol{\phi}(i)\right)\cdot\mathrm{sb}_2\left(\boldsymbol{\phi}(i')\right)
\]
for all $i$ and $i'$. Then
\[
\int\mathrm{d}\mu_I(W)\,\mathrm{sb}_1(F)=\int\mathrm{d}\mu_I(W)\,\mathrm{sb}_2(F)\qquad\text{for every }F\in\mathrm{Gr}[\psi,\chi,\boldsymbol{\phi}].
\]
\item[(b)] Let $\boldsymbol{\theta}$ be a ring endomorphism of $\mathrm{Gr}[\psi,\chi,W]$ which is either linear or antilinear over $\mathbb{C}$ and which sends every generator of each of the three families $\psi$, $\chi$, $W$ into the linear span of the generators of that family. Assume that
\[
\int\mathrm{d}\mu_I(W)\,\boldsymbol{\theta}\left(W(u)\right)\cdot\boldsymbol{\theta}\left(W(u')\right)=\int\mathrm{d}\mu_I(W)\,W(u)\cdot W(u')
\]
for all indices $u$, $u'$ of the generators $W$. Then
\[
\int\mathrm{d}\mu_I(W)\,\boldsymbol{\theta}(F)=\boldsymbol{\theta}\left(\int\mathrm{d}\mu_I(W)\,F\right)\qquad\text{for every }F\in\mathrm{Gr}[\psi,\chi,W].
\]
The assumption is satisfied if there is an invertible matrix $U$, indexed by the pairs $(y,\boldsymbol{\beta})$, such that $\boldsymbol{\theta}$ sends the vectors $W=\left(W_{\boldsymbol{\beta}}(y)\right)$ and $\bar{W}=\left(\bar{W}_{\boldsymbol{\beta}}(y)\right)$ of~\eqref{eq:renormalisation-group-42} to $UW$ and to $\left(U^{-1}\right)^T\bar{W}$, respectively. This is the case for the homomorphisms $\boldsymbol{\theta}$ of Section~\ref{sec:invariance-covariance} that belong to the transformations~\eqref{it:lattice-symmetries-gauge-Tg}, \eqref{it:lattice-symmetries-gauge-TG} and~\eqref{it:lattice-symmetries-gauge-TK}.
\end{enumerate}
\end{lemma}

\begin{proof}
(a) Every $F\in\mathrm{Gr}[\psi,\chi,\boldsymbol{\phi}]$ is a linear combination of products $c\cdot\boldsymbol{\phi}(i_1)\cdots\boldsymbol{\phi}(i_r)$, $r\ge0$, in which $c$ is a monomial in the generators $\psi(\xi)$ and $\chi(\boldsymbol{\zeta})$. As $\mathrm{sb}_1$, $\mathrm{sb}_2$ and the integral are linear, it suffices to consider one such product. For $1\le j\le r$ put
\[
v_j:=\mathrm{sb}_1\left(\boldsymbol{\phi}(i_j)\right),\qquad w_j:=\mathrm{sb}_2\left(\boldsymbol{\phi}(i_j)\right).
\]
These are linear combinations of the generators $W(u)$, and the two homomorphisms send the product to $c\cdot v_1\cdots v_r$ and to $c\cdot w_1\cdots w_r$. The monomial $c$ contains no generator $W$; hence it can be taken out of the integral:
\begin{align*}
\int\mathrm{d}\mu_I(W)\,c\cdot v_1\cdots v_r&=c\cdot\int\mathrm{d}\mu_I(W)\,v_1\cdots v_r,\\
\int\mathrm{d}\mu_I(W)\,c\cdot w_1\cdots w_r&=c\cdot\int\mathrm{d}\mu_I(W)\,w_1\cdots w_r.
\end{align*}
Wick's theorem, the third of the classical facts listed in Section~\ref{sec:recalled-auxiliary-results}, holds for arbitrary elements of the linear span of the generators $W(u)$, barred and unbarred, and so applies to $v_1,\dots,v_r$ and to $w_1,\dots,w_r$; for $\int\mathrm{d}\mu_I(W)$ it follows from the formula~\eqref{eq:polymer-activities-m} of Section~\ref{sec:norms-gaussian-integrals} for the moments. By this theorem, $\int\mathrm{d}\mu_I(W)\,v_1\cdots v_r$ is the Pfaffian of the $r\times r$ matrix whose entry in the place $(j,j')$ is $\int\mathrm{d}\mu_I(W)\,v_j\cdot v_{j'}$, and it vanishes if $r$ is odd; the same holds with $w_j$ in place of $v_j$. With $i=i_j$ and $i'=i_{j'}$ the hypothesis states that $\int\mathrm{d}\mu_I(W)\,v_j\cdot v_{j'}=\int\mathrm{d}\mu_I(W)\,w_j\cdot w_{j'}$. So the two matrices coincide, and so do the two integrals.

(b) Let $F\in\mathrm{Gr}[\psi,\chi,W]$. For a set $J$ of indices of the generators $W$ let $W(J)$ be the product of the generators $W(u)$, $u\in J$, taken in increasing order of the indices. Then
\[
F=\sum_Jc_J\cdot W(J)
\]
with elements $c_J\in\mathrm{Gr}[\psi,\chi]$, the sum extending over all such sets $J$. By its hypothesis on the families $\psi$ and $\chi$, the homomorphism $\boldsymbol{\theta}$ maps $\mathrm{Gr}[\psi,\chi]$ into itself. Since $\boldsymbol{\theta}$ is additive and multiplicative, and factors from $\mathrm{Gr}[\psi,\chi]$ can be taken out of the integral,
\[
\int\mathrm{d}\mu_I(W)\,\boldsymbol{\theta}(F)=\sum_J\boldsymbol{\theta}(c_J)\cdot\int\mathrm{d}\mu_I(W)\,\boldsymbol{\theta}\left(W(J)\right).
\]
Apply (a) to further generators $\boldsymbol{\phi}(u)$, one for each index $u$ of the generators $W$, and to the homomorphisms $\mathrm{sb}_1$, $\mathrm{sb}_2$ that leave $\psi$ and $\chi$ fixed and satisfy
\[
\mathrm{sb}_1\left(\boldsymbol{\phi}(u)\right)=\boldsymbol{\theta}\left(W(u)\right),\qquad\mathrm{sb}_2\left(\boldsymbol{\phi}(u)\right)=W(u).
\]
Both values are linear combinations of the generators $W$, and the hypothesis of (a) is then the hypothesis of (b). Let $\boldsymbol{\phi}(J)$ be the product of the $\boldsymbol{\phi}(u)$, $u\in J$, in the order used for $W(J)$. Then $\mathrm{sb}_1\left(\boldsymbol{\phi}(J)\right)=\boldsymbol{\theta}\left(W(J)\right)$ and $\mathrm{sb}_2\left(\boldsymbol{\phi}(J)\right)=W(J)$, and (a) gives
\[
\int\mathrm{d}\mu_I(W)\,\boldsymbol{\theta}\left(W(J)\right)=\int\mathrm{d}\mu_I(W)\,W(J).
\]
By~\eqref{eq:renormalisation-group-42} these moments are real numbers, and $\boldsymbol{\theta}(c\cdot r)=\boldsymbol{\theta}(c)\cdot r$ for $c\in\mathrm{Gr}[\psi,\chi]$ and real $r$, whether $\boldsymbol{\theta}$ is linear or antilinear. Consequently
\begin{align*}
\int\mathrm{d}\mu_I(W)\,\boldsymbol{\theta}(F)&=\sum_J\boldsymbol{\theta}(c_J)\cdot\int\mathrm{d}\mu_I(W)\,W(J)\\
&=\boldsymbol{\theta}\left(\sum_Jc_J\cdot\int\mathrm{d}\mu_I(W)\,W(J)\right)=\boldsymbol{\theta}\left(\int\mathrm{d}\mu_I(W)\,F\right).
\end{align*}

It remains to prove the last two sentences of (b). From here on, each of the letters $u$, $u'$, $a$, $a'$ stands for a pair $(y,\boldsymbol{\beta})$ consisting of a unit site $y\in T^0_k$ and a spinor index $\boldsymbol{\beta}$. For such a pair $u=(y,\boldsymbol{\beta})$, the generators $W_{\boldsymbol{\beta}}(y)$ and $\bar{W}_{\boldsymbol{\beta}}(y)$ of~\eqref{eq:renormalisation-group-42}, the unbarred and the barred one with these indices, are denoted by $W(u)$ and $\bar{W}(u)$. In this notation~\eqref{eq:renormalisation-group-42} gives the pair expectations
\begin{gather*}
\int\mathrm{d}\mu_I(W)\,W(u)\cdot W(u')=0,\qquad\int\mathrm{d}\mu_I(W)\,\bar{W}(u)\cdot\bar{W}(u')=0,\\
\int\mathrm{d}\mu_I(W)\,W(u)\cdot\bar{W}(u')=\boldsymbol{\varsigma}\cdot\delta_{uu'},
\end{gather*}
and, by antisymmetry, $\int\mathrm{d}\mu_I(W)\,\bar{W}(u')\cdot W(u)=-\boldsymbol{\varsigma}\cdot\delta_{uu'}$. Suppose now that $\boldsymbol{\theta}$ sends $W$ to $UW$ and $\bar{W}$ to $\left(U^{-1}\right)^T\bar{W}$, that is,
\[
\boldsymbol{\theta}\left(W(u)\right)=\sum_aU_{ua}\cdot W(a),\qquad\boldsymbol{\theta}\left(\bar{W}(u')\right)=\sum_{a'}\left(\left(U^{-1}\right)^T\right)_{u'a'}\cdot\bar{W}(a').
\]
The images of two unbarred generators, and those of two barred generators, are linear combinations of generators of one kind; so their pair expectations vanish. Moreover
\begin{align*}
\int\mathrm{d}\mu_I(W)\,\boldsymbol{\theta}\left(W(u)\right)\cdot\boldsymbol{\theta}\left(\bar{W}(u')\right)&=\sum_{a,a'}U_{ua}\cdot\left(\left(U^{-1}\right)^T\right)_{u'a'}\cdot\boldsymbol{\varsigma}\cdot\delta_{aa'}\\
&=\boldsymbol{\varsigma}\cdot\sum_aU_{ua}\cdot\left(U^{-1}\right)_{au'}=\boldsymbol{\varsigma}\cdot\left(UU^{-1}\right)_{uu'}=\boldsymbol{\varsigma}\cdot\delta_{uu'},
\end{align*}
and by antisymmetry the pair expectation of $\boldsymbol{\theta}\left(\bar{W}(u')\right)$ and $\boldsymbol{\theta}\left(W(u)\right)$, in this order, equals $-\boldsymbol{\varsigma}\cdot\delta_{uu'}$. Thus $\boldsymbol{\theta}$ preserves the pair expectation of any two generators, which is the assumption of (b).

The transformation~\eqref{it:lattice-symmetries-gauge-Tg}. Let $g\in\mathbb{G}^L_k$, and let $S\in\mathrm{Sp}$ be its first component. Let $\Pi_g$ be the matrix, indexed by the unit sites of a fundamental domain, of the map $f\mapsto f\circ g^{-1}$ on the antiperiodic functions $f$ on $\mathbb{Z}^3$, each such function being described by its values on the fundamental domain. By the antiperiodicity (Section~\ref{sec:notation}), every row and every column of $\Pi_g$ contains exactly one non-zero entry, and this entry is $1$ or $-1$; hence $\left(\Pi_g^{-1}\right)^T=\Pi_g$. By the definition of~\eqref{it:lattice-symmetries-gauge-Tg}, $\boldsymbol{\theta}$ sends $W$ to $\left(S\otimes\Pi_g\right)W$ and $\bar{W}$ to $\left(\left(S^{-1}\right)^T\otimes\Pi_g\right)\bar{W}$: the same matrix $\Pi_g$, with the same signs, appears in both, because $W$ and $\bar{W}$ are antiperiodic in the same way. For $U:=S\otimes\Pi_g$ one has
\[
\left(U^{-1}\right)^T=\left(S^{-1}\right)^T\otimes\left(\Pi_g^{-1}\right)^T=\left(S^{-1}\right)^T\otimes\Pi_g.
\]
The transformation~\eqref{it:lattice-symmetries-gauge-TG}. Here $\boldsymbol{\theta}$ multiplies $W_{\boldsymbol{\beta}}(y)$ by $\exp\left(i\,e_k\cdot\boldsymbol{\Lambda}(y)\right)$ and $\bar{W}_{\boldsymbol{\beta}}(y)$ by $\exp\left(-i\,e_k\cdot\boldsymbol{\Lambda}(y)\right)$, $y\in T^0_k$. So $U$ is the diagonal matrix $\exp\left(i\,e_k\cdot\boldsymbol{\Lambda}^u\right)$, and $\left(U^{-1}\right)^T=\exp\left(-i\,e_k\cdot\boldsymbol{\Lambda}^u\right)$; in the pair expectation of $\boldsymbol{\theta}\left(W(u)\right)$ and $\boldsymbol{\theta}\left(\bar{W}(u)\right)$ the two phases cancel. The transformation~\eqref{it:lattice-symmetries-gauge-TK}. Here $\boldsymbol{\theta}$ multiplies $W$ and $\bar{W}$ by the matrix $\varkappa$, which acts on the spinor index. So $U=\varkappa$, and $\left(U^{-1}\right)^T=\left(\varkappa^{-1}\right)^T=\varkappa$, since $\varkappa^T=\varkappa=\varkappa^{-1}$ by the Clifford relations of Section~\ref{sec:recalled-auxiliary-results}; equivalently, $\varkappa\cdot\varkappa^T=\mathbf{1}$.
\end{proof}

\subsubsection{Symmetries of $E^{\#,x}_k$}\label{sec:symmetries}

\begin{lemma}\label{lem:renormalisation-group-2}
Let $(m,\mathscr{E})\in\mathrm{Ball}^x_k(4)$. Then $E^{\#,x}_k(m,\mathscr{E})$ is invariant
\begin{enumerate}
\item[(a)] under the transformation~\eqref{it:lattice-symmetries-gauge-Tg}, for every $g\in\mathbb{G}^L_k$;
\item[(b)] under the transformation~\eqref{it:lattice-symmetries-gauge-TG};
\item[(c)] under the transformation~\eqref{it:lattice-symmetries-gauge-TC};
\item[(d)] under the transformation~\eqref{it:lattice-symmetries-gauge-TK}, provided that $m$ is real and $\mathscr{E}\in K^x_k$.
\end{enumerate}
\end{lemma}

\begin{proof}[Proof of Lemma~\ref{lem:renormalisation-group-2}(a)]
Let $g=(S,a)\in\mathbb{G}^L_k$, and put $r_{\nu}:=r_{\nu}(S)$; thus $gx=r(S)x+a$ and $r(S)e_{\nu}=r_{\nu}e_{\nu}$ for $\nu=0,1,2$. Throughout the proof, $X$ is an arbitrary polymer, $s=(s',s'')$ an arbitrary point of the polydisk $\mathrm{Pd}_k$, and $B$ an arbitrary connection in $\tilde{G}_k$. The point $gs$ lies in $\mathrm{Pd}_k$ (Section~\ref{sec:invariance-covariance}). Since the group $\mathbb{G}_k$ contains $\mathbb{G}^L_k$ and $g$ acts linearly on bond functions, Lemma~\ref{lem:symmetries-domains} shows that the connections $gB$, $-B$ and $-gB=g(-B)$ lie in $\tilde{G}_k$ as well. Consequently property~\ref{it:renormalisation-group-14} and the definitions~\eqref{eq:renormalisation-group-56} and~\eqref{eq:renormalisation-group-58} apply at each of the connections $B$, $gB$, $-B$ and $g(-B)$.

\emph{Cubes, blocks and polymers.} By Proposition~\ref{prop:symmetries-extraction}(a) the map $g$ sends $M$-cubes onto $M$-cubes and $LM$-cubes onto $LM$-cubes, the centre of each such cube onto the centre of its image, and it is an isometry of the torus. Consequently $g$ sends the parameter cubes of each of the kinds $\pi_k$, $\pi'_k$, $\pi''_k$ of Definition~\ref{def:parameter-cubes-polydisks} onto parameter cubes of the same kind, linked cubes onto linked cubes and nearest neighbours onto nearest neighbours, and for every site $z$, every unit site $y$, every polymer $X$ and every set $P$ of parameter cubes outside $P_0(X)$
\[
gK_z=K_{gz},\qquad gB_k(y)=B_k(gy),\qquad gP_0(X)=P_0(gX),\qquad gY'(X,P)=Y'(gX,gP);
\]
moreover a pair of kind~(I), resp.\ of kind~(II), for $X$ in the sense of Definition~\ref{def:pairs-sites-holder} is sent to a pair of the same kind for $gX$.

\emph{The field $\mathscr{Z}$.} The map $g$ acts on $\mathsf{Z}\in V_k$, and it sends $V_k$ onto itself (Proposition~\ref{prop:basis-map-box}). By property~\ref{it:renormalisation-group-14}, $H_k(gs')\cdot g=g\cdot H_k(s')$ and $C_k^{\frac{1}{2}}(gs')\cdot g=g\cdot C_k^{\frac{1}{2}}(s')$, and $C\cdot g=g\cdot C$ by Proposition~\ref{prop:symmetries-reality}. Therefore
\begin{align*}
\mathscr{Z}(gs')[g\mathsf{Z}]&=H_k(gs')\cdot C\cdot C_k^{\frac{1}{2}}(gs')\cdot g\mathsf{Z}=H_k(gs')\cdot C\cdot g\cdot C_k^{\frac{1}{2}}(s')\cdot\mathsf{Z}\\
&=H_k(gs')\cdot g\cdot C\cdot C_k^{\frac{1}{2}}(s')\cdot\mathsf{Z}=g\cdot H_k(s')\cdot C\cdot C_k^{\frac{1}{2}}(s')\cdot\mathsf{Z}\\
&=g\,\mathscr{Z}(s')[\mathsf{Z}],
\end{align*}
that is, $\mathscr{Z}$ is covariant. As $t$ is not changed by~\eqref{it:lattice-symmetries-gauge-Tg}, the connection $\mathscr{T}+t\cdot\mathscr{Z}(s')[\mathsf{Z}]$ is replaced by $g\mathscr{T}+t\cdot\mathscr{Z}(gs')[g\mathsf{Z}]=g\left(\mathscr{T}+t\cdot\mathscr{Z}(s')[\mathsf{Z}]\right)$.

\emph{Pairings of barred and unbarred fields.} Let $N$ be an operator on the spinor functions on $T^0_k$, with the kernel $N(y,y')$. Since $(g\bar{\Psi})(y)=\left(S^{-1}\right)^T\bar{\Psi}\left(g^{-1}y\right)$ and $(g\Psi)(y')=S\cdot\Psi\left(g^{-1}y'\right)$,
\[
\left\langle g\bar{\Psi},N\,g\Psi\right\rangle=\sum_{y,y'}\bar{\Psi}\left(g^{-1}y\right)^T\cdot S^{-1}\cdot N(y,y')\cdot S\cdot\Psi\left(g^{-1}y'\right)=\left\langle\bar{\Psi},\left(g^{-1}\cdot N\cdot g\right)\Psi\right\rangle.
\]
For $N=M_k(gX,gB)=g\cdot M_k(X,B)\cdot g^{-1}$, which is property~\ref{it:renormalisation-group-14}, the right side is $\left\langle\bar{\Psi},M_k(X,B)\Psi\right\rangle$. Moreover
\[
(g\bar{\psi})(x)^T\cdot(g\psi)(x)=\bar{\psi}\left(g^{-1}x\right)^T\cdot S^{-1}\cdot S\cdot\psi\left(g^{-1}x\right)=\bar{\psi}\left(g^{-1}x\right)^T\cdot\psi\left(g^{-1}x\right),
\]
and summation over the sites $x$ of $g\square$ gives $\left\langle g\bar{\psi},g\psi\right\rangle_{g\square}=\langle\bar{\psi},\psi\rangle_{\square}$: the forms $\langle\bar{\psi},\psi\rangle_{\square}$ are invariant.

\emph{The fields $J_k^{\#}\Psi$ and $W_k^{\#}$.} By property~\ref{it:renormalisation-group-14},
\[
H_k(gs'',gB)\cdot g=g\cdot H_k(s'',B),\qquad\Gamma_k(gs'',gB)\cdot g=g\cdot\Gamma_k(s'',B).
\]
On the barred fields $g$ acts by the operator $\underline{g}$ given by $\left(\underline{g} f\right)(x):=\left(S^{-1}\right)^Tf\left(g^{-1}x\right)$. By the Clifford relations of Section~\ref{sec:recalled-auxiliary-results}, $\left(S^{-1}\right)^T=\pm\mathsf{C}^{-1}\cdot S\cdot\mathsf{C}$; hence $\underline{g}=\pm\mathsf{C}^{-1}\cdot g\cdot\mathsf{C}$, with the same sign, and $\mathsf{C}^{-1}\cdot g=\pm\underline{g}\cdot\mathsf{C}^{-1}$, $g^{-1}\cdot\mathsf{C}=\pm\mathsf{C}\cdot\underline{g}^{-1}$. Since $-gB=g(-B)$, the definition of $\bar{H}_k$ in~\eqref{eq:renormalisation-group-56} and property~\ref{it:renormalisation-group-14}, applied at the connection $-B$ of $\tilde{G}_k$, give
\begin{align*}
\bar{H}_k(gs'',gB)&=\mathsf{C}^{-1}\cdot H_k(gs'',-gB)\cdot\mathsf{C}=\mathsf{C}^{-1}\cdot g\cdot H_k(s'',-B)\cdot g^{-1}\cdot\mathsf{C}\\
&=\underline{g}\cdot\mathsf{C}^{-1}\cdot H_k(s'',-B)\cdot\mathsf{C}\cdot\underline{g}^{-1}=\underline{g}\cdot\bar{H}_k(s'',B)\cdot\underline{g}^{-1},
\end{align*}
the two signs in the third equality cancelling. By Proposition~\ref{prop:symmetries-extraction}(a) the map $g$ sends each path $\Gamma^{\pi}(x,x')$ onto $\Gamma^{\pi}(gx,gx')$, and $(gB)(gx,gx')=B(x,x')$ for every bond $(x,x')$; hence $(\tau\,gB)(gx,gx')=(\tau B)(x,x')$. As also $|gx-gx'|=|x-x'|$ and $(gf)(gx)=S\cdot f(x)$, the definition of $\delta_{\alpha,B}$ in Definition~\ref{def:pairs-sites-holder} gives, for every spinor function $f$ and every pair $(x,x')\in\mathrm{Pz}_k$,
\begin{align*}
\left(\delta_{\alpha,gB}\,gf\right)(gx,gx')&=|gx-gx'|^{-\alpha}\cdot\left(\exp\left(i\,e_k\eta_k\cdot(\tau\,gB)(gx,gx')\right)\cdot(gf)(gx')\right.\\
&\qquad\qquad\left.{}-(gf)(gx)\right)\\
&=|x-x'|^{-\alpha}\cdot\left(\exp\left(i\,e_k\eta_k\cdot(\tau B)(x,x')\right)\cdot S\cdot f(x')-S\cdot f(x)\right)\\
&=S\cdot\left(\delta_{\alpha,B}f\right)(x,x').
\end{align*}
The same equalities hold with $-B$, $\underline{g}$ and $\left(S^{-1}\right)^T$ in place of $B$, $g$ and $S$, because $\left(\tau\,(-gB)\right)(gx,gx')=\left(\tau\,(-B)\right)(x,x')$ and $(\underline{g} f)(gx)=\left(S^{-1}\right)^Tf(x)$; thus $\left(\delta_{\alpha,-gB}\,\underline{g} f\right)(gx,gx')=\left(S^{-1}\right)^T\cdot\left(\delta_{\alpha,-B}f\right)(x,x')$. Consequently the fields $J_k^{\#}\Psi$ and $W_k^{\#}$ of~\eqref{eq:renormalisation-group-56} are covariant.

\emph{The substitution for $\Psi$ and $\bar{\Psi}$.} For a polymer $\mathbf{Y}$ and a connection $B'\in\tilde{G}_k$ let $T^{(i)}_{\mathbf{Y}}(B')$ and $\bar{T}^{(i)}_{\mathbf{Y}}(B')$, $i=1,2$, denote the operators of~\eqref{eq:renormalisation-group-58} formed with $\mathbf{Y}$ and $B'$. As $g\mathbf{Y}$ is again a polymer, these operators are defined, with $\mathbf{Y}$ and with $g\mathbf{Y}$, at each of the connections $B$, $gB$, $-B$ and $g(-B)$ of $\tilde{G}_k$. For every site $x$, direction $\nu$ and sign $\boldsymbol{\sigma}\in\{1,-1\}$, and for the unit site $y$ with $x\in B_k(y)$,
\begin{gather*}
(gx)^{\boldsymbol{\sigma}}=g\left(x^{r_{\nu}\boldsymbol{\sigma}}\right),\qquad P^{\boldsymbol{\sigma}}_{\nu}\cdot S=S\cdot P^{r_{\nu}\boldsymbol{\sigma}}_{\nu},\\
u_{gB}\left(gx,g\left(x^{r_{\nu}\boldsymbol{\sigma}}\right)\right)=u_B\left(x,x^{r_{\nu}\boldsymbol{\sigma}}\right),\qquad\phi_{gB,gy}(gx)=\phi_{B,y}(x).
\end{gather*}
The first identity holds because $r(S)e_{\nu}=r_{\nu}e_{\nu}$ and $r_{\nu}^2=1$; the second is equivalent to $S^{-1}\cdot P^{\boldsymbol{\sigma}}_{\nu}\cdot S=P^{r_{\nu}\boldsymbol{\sigma}}_{\nu}$, which follows from $S^{-1}\cdot\gamma_{\nu}\cdot S=r_{\nu}\gamma_{\nu}$ (Clifford relations of Section~\ref{sec:recalled-auxiliary-results}); the third follows from $(gB)(gx,gx')=B(x,x')$; the fourth is Proposition~\ref{prop:symmetries-fermion-operators}. Moreover $(gx)^{\boldsymbol{\sigma}}\in g\mathbf{Y}$ if and only if $x^{r_{\nu}\boldsymbol{\sigma}}\in\mathbf{Y}$, and $B_k(gy)=gB_k(y)$. Let $y$ be a unit site of $\mathbf{Y}$. In the two formulas~\eqref{eq:renormalisation-group-58} for the operators formed with $(g\mathbf{Y},gB)$, applied to $g\psi$ and $g\chi$ at the unit site $gy$, write the sites of $B_k(gy)$ as $gx$ with $x\in B_k(y)$, and use $(g\psi)(gx)=S\cdot\psi(x)$, $(g\chi)(gx,gx')=S\cdot\chi(x,x')$ and the identities above. This gives
\begin{align*}
\left(T^{(1)}_{g\mathbf{Y}}(gB)\,g\psi\right)(gy)={}&\eta_k^3\sum_{x\in B_k(y)}\phi_{B,y}(x)\cdot S\cdot\left\{\left(1+\frac{\mathfrak{m}}{b_k}\right)\psi(x)\right.\\
&\left.{}-\left(b_k\cdot\eta_k\right)^{-1}\sum_{(\nu,\boldsymbol{\sigma}):\,x^{r_{\nu}\boldsymbol{\sigma}}\in\mathbf{Y}}P^{r_{\nu}\boldsymbol{\sigma}}_{\nu}\cdot\left[u_B\left(x,x^{r_{\nu}\boldsymbol{\sigma}}\right)\cdot\psi\left(x^{r_{\nu}\boldsymbol{\sigma}}\right)-\psi(x)\right]\right\},\\
\left(T^{(2)}_{g\mathbf{Y}}(gB)\,g\chi\right)(gy)={}&-b_k^{-1}\cdot\eta_k^{2+\alpha}\sum_{x\in B_k(y)}\phi_{B,y}(x)\cdot S\cdot\sum_{(\nu,\boldsymbol{\sigma}):\,x^{r_{\nu}\boldsymbol{\sigma}}\notin\mathbf{Y}}P^{r_{\nu}\boldsymbol{\sigma}}_{\nu}\cdot\chi\left(x,x^{r_{\nu}\boldsymbol{\sigma}}\right).
\end{align*}
For fixed $\nu$ the map $\boldsymbol{\sigma}\mapsto r_{\nu}\boldsymbol{\sigma}$ is a bijection of $\{1,-1\}$. Writing $\boldsymbol{\sigma}$ for $r_{\nu}\boldsymbol{\sigma}$ in the inner sums and comparing with~\eqref{eq:renormalisation-group-58}, one finds
\begin{align*}
\left(T^{(1)}_{g\mathbf{Y}}(gB)\,g\psi\right)(gy)&=S\cdot\left(T^{(1)}_{\mathbf{Y}}(B)\psi\right)(y),\\
\left(T^{(2)}_{g\mathbf{Y}}(gB)\,g\chi\right)(gy)&=S\cdot\left(T^{(2)}_{\mathbf{Y}}(B)\chi\right)(y).
\end{align*}
For the barred operators, $\bar{T}^{(i)}_{g\mathbf{Y}}(gB)=\mathsf{C}^{-1}\cdot T^{(i)}_{g\mathbf{Y}}(-gB)\cdot\mathsf{C}$ by~\eqref{eq:renormalisation-group-58}, and $-gB=g(-B)$. Let $f$ be $\bar{\psi}$ if $i=1$ and $\bar{\chi}$ if $i=2$. Using $\underline{g}=\pm\mathsf{C}^{-1}\cdot g\cdot\mathsf{C}$, then the two identities just proved, with the connection $-B$ and the field $\mathsf{C} f$ in place of $B$ and of $\psi$, resp.\ $\chi$, and finally the definition of $\bar{T}^{(i)}_{\mathbf{Y}}(B)$, one obtains
\begin{align*}
\left(\bar{T}^{(i)}_{g\mathbf{Y}}(gB)\,\underline{g} f\right)(gy)&=\pm\mathsf{C}^{-1}\cdot\left(T^{(i)}_{g\mathbf{Y}}\left(g(-B)\right)\,g\left(\mathsf{C} f\right)\right)(gy)\\
&=\pm\mathsf{C}^{-1}\cdot S\cdot\left(T^{(i)}_{\mathbf{Y}}(-B)\left(\mathsf{C} f\right)\right)(y)\\
&=\pm\mathsf{C}^{-1}\cdot S\cdot\mathsf{C}\cdot\left(\bar{T}^{(i)}_{\mathbf{Y}}(B)f\right)(y)\\
&=\left(S^{-1}\right)^T\cdot\left(\bar{T}^{(i)}_{\mathbf{Y}}(B)f\right)(y),
\end{align*}
where the last equality holds because $\left(S^{-1}\right)^T=\pm\mathsf{C}^{-1}\cdot S\cdot\mathsf{C}$ with the same sign, whose square is $1$. Hence the substitution~\eqref{eq:renormalisation-group-58}, formed with $\mathscr{T}$, is covariant: $\left(T^{(1)}(\mathscr{T})\psi\right)(y)+\left(T^{(2)}(\mathscr{T})\chi\right)(y)$ is transformed like $\Psi(y)$, and $\left(\bar{T}^{(1)}(\mathscr{T})\bar{\psi}\right)(y)+\left(\bar{T}^{(2)}(\mathscr{T})\bar{\chi}\right)(y)$ like $\bar{\Psi}(y)$.

\emph{The case $k=0$.} Let $k=0$. Then the substitution~\eqref{eq:renormalisation-group-58} is not made, and $G^{(1)}$ is the function~\eqref{eq:renormalisation-group-60} of the single $M$-cubes. For $x\in\square$, a direction $\nu$ and $\boldsymbol{\sigma}\in\{1,-1\}$, call the summand with these indices in the first, resp.\ in the second, of the two sums in~\eqref{eq:renormalisation-group-60} the term $(x,\nu,\boldsymbol{\sigma})$ of that sum. By the covariance of $\mathscr{Z}$ and by $(gx)^{\boldsymbol{\sigma}}=g\left(x^{r_{\nu}\boldsymbol{\sigma}}\right)$,
\[
\mathscr{Z}(gs')[g\mathsf{Z}]\left(gx,(gx)^{\boldsymbol{\sigma}}\right)=\mathscr{Z}(s')[\mathsf{Z}]\left(x,x^{r_{\nu}\boldsymbol{\sigma}}\right).
\]
Hence the term $(gx,\nu,\boldsymbol{\sigma})$ of the first sum of $G^{(1)}(gs,g\square)$, taken at the transformed fields and parameters, is
\[
\bar{\psi}(x)\,S^{-1}\cdot P^{\boldsymbol{\sigma}}_{\nu}\cdot S\left(\exp\left(i\,e_k\cdot\mathscr{Z}(s')\left(x,x^{r_{\nu}\boldsymbol{\sigma}}\right)\right)-1\right)\left(\psi(x)+\chi\left(x,x^{r_{\nu}\boldsymbol{\sigma}}\right)\right),
\]
and $S^{-1}\cdot P^{\boldsymbol{\sigma}}_{\nu}\cdot S=P^{r_{\nu}\boldsymbol{\sigma}}_{\nu}$ by the Clifford relations; so it is the term $(x,\nu,r_{\nu}\boldsymbol{\sigma})$ of the first sum of $G^{(1)}(s,\square)$. The term $(gx,\nu,\boldsymbol{\sigma})$ of the second sum, taken at the transformed fields and parameters, is
\[
\left(\bar{\psi}(x)+\bar{\chi}\left(x,x^{r_{\nu}\boldsymbol{\sigma}}\right)\right)S^{-1}\cdot P^{-\boldsymbol{\sigma}}_{\nu}\cdot S\left(\exp\left(-i\,e_k\cdot\mathscr{Z}(s')\left(x,x^{r_{\nu}\boldsymbol{\sigma}}\right)\right)-1\right)\psi(x),
\]
and $S^{-1}\cdot P^{-\boldsymbol{\sigma}}_{\nu}\cdot S=P^{-r_{\nu}\boldsymbol{\sigma}}_{\nu}$; so it is the term $(x,\nu,r_{\nu}\boldsymbol{\sigma})$ of the second sum. As $(x,\boldsymbol{\sigma})$ ranges over $\square\times\{1,-1\}$, so does $(x,r_{\nu}\boldsymbol{\sigma})$. Therefore $G^{(1)}(gs,g\square)$, taken at the transformed arguments, equals $G^{(1)}(s,\square)$. At level $0$ the families of property~\ref{it:renormalisation-group-17} have property~\ref{it:renormalisation-group-14}; so the relations of property~\ref{it:renormalisation-group-14} used above hold for $k=0$ as well.

\emph{The datum.} Since $\mathscr{E}\in K^{x,C}_k$ and $R$ maps $K^{x,C}_k$ into itself (Proposition~\ref{prop:symmetries-extraction}), the function $\mathscr{E}^r=R\,\mathscr{E}$ satisfies clause~\eqref{it:renormalisation-group-1-k3}: it is invariant under the action of every element of $\mathbb{G}_k$, in particular under~\eqref{it:lattice-symmetries-gauge-Tg}, as $\mathbb{G}_k$ contains $\mathbb{G}^L_k$. The number $m'$ is not changed.

\emph{The integrations.} The set $\mathrm{Box}_k$ and the measure $\hat{\mu}$ are invariant under $g$, by Proposition~\ref{prop:symmetries-reality}. On the generators $W$ and $\bar{W}$ the map $g$ acts by $W\mapsto UW$ and $\bar{W}\mapsto\left(U^{-1}\right)^T\bar{W}$ with $U=S\otimes\Pi_g$, as was shown in the proof of Lemma~\ref{lem:substitutions-integral}(b).

\emph{Conclusion.} Apply the rules~\eqref{it:operations-preserve-invariance-Pr1}--\eqref{it:operations-preserve-invariance-Pr5} to the displays~\eqref{eq:renormalisation-group-56}--\eqref{eq:renormalisation-group-61}, in their order. Let first $k\ge1$. The functions $G^{(1)}$, $G^{(2)}$, $G^{(3)}$ of~\eqref{eq:renormalisation-group-59} arise from the invariant functions $\left\langle\bar{\Psi},M_k(X,\cdot)\Psi\right\rangle$, $\mathscr{E}^r$, $\langle\bar{\psi},\psi\rangle_{\square}$ and from the number $m'$ by inserting the covariant fields $\mathscr{Z}$ and $J_k^{\#}\Psi$ and by contour integrals in $t$; so they are invariant by~\eqref{it:operations-preserve-invariance-Pr1} and~\eqref{it:operations-preserve-invariance-Pr2}. For $k=0$ one has $G^{(2)}=G^{(3)}=G^{(4)}=0$ by~\eqref{eq:renormalisation-group-60}, and the invariance of the function $G^{(1)}$ of~\eqref{eq:renormalisation-group-60} was proved in the paragraph on the case $k=0$. For every $k$, the localisation~\eqref{eq:renormalisation-group-57} preserves invariance by~\eqref{it:operations-preserve-invariance-Pr3} and~\eqref{it:operations-preserve-invariance-Pr4}, and so does the covariant substitution~\eqref{eq:renormalisation-group-58}, made for $k\ge1$ only, by~\eqref{it:operations-preserve-invariance-Pr1}; hence $E^{(i),*}$ and $E^{(i),x}$, $i\le3$, are invariant. By~\eqref{it:operations-preserve-invariance-Pr1} and~\eqref{it:operations-preserve-invariance-Pr2} the functions $G^{(\le3)}$, $G^{(4)}$, $G^{(5)}$, which are obtained from $E^{(\le3),x}$, $\mathscr{E}^r$, $\langle\bar{\psi},\psi\rangle_{\square}$ and $m'$ by inserting the covariant field $W_k^{\#}$ and by contour integrals in $t$, are invariant, and by~\eqref{it:operations-preserve-invariance-Pr3} and~\eqref{it:operations-preserve-invariance-Pr4} so is $E^{(\le5),x}$. In~\eqref{eq:renormalisation-group-61} the function $\mathsf{K}$ is invariant by~\eqref{it:operations-preserve-invariance-Pr1} and~\eqref{it:operations-preserve-invariance-Pr4}, the function $\mathsf{H}$ by~\eqref{it:operations-preserve-invariance-Pr5}, which applies by the paragraph on the integrations, and $E^{\#,x}_k$ by~\eqref{it:operations-preserve-invariance-Pr4} and~\eqref{it:operations-preserve-invariance-Pr1}; for the last step note that the series in~\eqref{eq:renormalisation-group-61} converge in norm, by assertion~(1) of Section~\ref{sec:activity-before-reblocking}. Thus $E^{(i),x}$, $E^{(\le5),x}$, $\mathsf{K}$, $\mathsf{H}$ and $E^{\#,x}_k$ are invariant under~\eqref{it:lattice-symmetries-gauge-Tg}.
\end{proof}

\begin{proof}[Proof of Lemma~\ref{lem:renormalisation-group-2}(b) and~(d)]
(b) Let $\mathbf{Y}$ be a polymer, let $\lambda$ and $\lambda'$ be real functions on the sites of $T_k$, put $\boldsymbol{\Lambda}:=\lambda+\lambda'$, and let $\boldsymbol{\Lambda}^u$ be the restriction of $\boldsymbol{\Lambda}$ to $T^0_k$. Let $(A,\mathscr{C})$ and $(A^1,\mathscr{C}^1)$ be elements of $\tilde{R}^{-}_k\times\mathfrak{X}_k$ such that $A^1=A-d\lambda$ and $\mathscr{C}^1=\mathscr{C}-d\lambda'$ on the bonds at distance at most $1$ from $\mathbf{Y}$, as in~\eqref{it:lattice-symmetries-gauge-TG}, and put
\[
\mathscr{T}:=A+\mathscr{C},\qquad\mathscr{T}^1:=A^1+\mathscr{C}^1.
\]
The transformation~\eqref{it:lattice-symmetries-gauge-TG} replaces the generators by $\psi^{\boldsymbol{\Lambda}}$, $\bar{\psi}^{\boldsymbol{\Lambda}}$, $\chi^{\boldsymbol{\Lambda}}$, $\bar{\chi}^{\boldsymbol{\Lambda}}$, $\Psi^{\boldsymbol{\Lambda}}$, $\bar{\Psi}^{\boldsymbol{\Lambda}}$, $W^{\boldsymbol{\Lambda}}$, $\bar{W}^{\boldsymbol{\Lambda}}$. For a real function $f$ on a set of sites, $\exp\left(i\,e_k\cdot f\right)$ denotes also the operator of multiplication by this function. Let $\mathsf{Z}\in\mathrm{Box}_k$, let $s\in\mathrm{Pd}_k$ vanish at every parameter cube that is not contained in $\mathbf{Y}$, let $t\in\mathbb{C}$ with $|t|\le\rho_1$, and put
\[
B:=\mathscr{T}+t\cdot\mathscr{Z}(s'),\qquad B^1:=\mathscr{T}^1+t\cdot\mathscr{Z}(s').
\]
By part~\eqref{it:bounds-analyticity-q1} of Lemma~\ref{lem:bounds-analyticity}, the connections $B$ and $B^1$ lie in $\tilde{G}_k$; by Lemma~\ref{lem:symmetries-domains}, so does $B-d\boldsymbol{\Lambda}$. On every bond at distance at most $1$ from $\mathbf{Y}$ one has
\[
B^1-B=\mathscr{T}^1-\mathscr{T}=-d\lambda-d\lambda'=-d\boldsymbol{\Lambda} .
\]
Every bond that meets $\mathbf{Y}$ is at distance at most $1$ from $\mathbf{Y}$: one of its end points lies in $\mathbf{Y}$, and the other one has sup-distance $\eta_k\le1$ from the first. Hence
\[
B^1=B-d\boldsymbol{\Lambda}\qquad\text{on every bond that meets }\mathbf{Y}.
\]

\emph{Primitives.} Let $(X,P)$ be a term of~\eqref{eq:renormalisation-group-57} with $Y'(X,P)=\mathbf{Y}$. In this term, the function $G(s,X)$ to which the localisation~\eqref{eq:renormalisation-group-57} is applied enters only through its values at families $s\in\mathrm{Pd}_k$ that vanish at every parameter cube outside $P_0(X)\cup P$: the bracket in~\eqref{eq:renormalisation-group-57} sets $s_c=0$ for every such cube $c$, and only the variables $s_c$ with $c\in P$ are differentiated and integrated. For such a family $s$, the polymer $X$, the set $P$ and $s$ satisfy the hypotheses of Lemma~\ref{lem:renormalisation-group-1}. Every cube of $P_0(X)\cup P$ is contained in $\mathbf{Y}$: by fact~\eqref{it:renormalisation-group-g0}, a lattice site of such a cube lies in an $M$-cube, this $M$-cube has that site in common with $Y^{\mathrm{p}}(X,P)$, and so it is one of the $M$-cubes of $Y'(X,P)=\mathbf{Y}$. Consequently a family that vanishes outside $P_0(X)\cup P$ vanishes at every parameter cube not contained in $\mathbf{Y}$, that is, it is one of the families $s$ fixed above. Let now $s\in\mathrm{Pd}_k$ vanish outside $P_0(X)\cup P$. Since $\mathbf{Y}$ is connected, the component $Y'_1$ of $Y'(X,P)$ that contains $X$ is $\mathbf{Y}$ itself. By Steps~(2) and~(4) of the proof of Lemma~\ref{lem:renormalisation-group-1}, every row of $H_k(s'',\cdot)$, of $\bar{H}_k(s'',\cdot)$ or of $\Gamma_k(s'',\cdot)$ that occurs in this term either vanishes for every connection, or has its index contained in a link-component $\mathcal{L}_3$ of the union of the cubes of kind $\pi_k$ in $P_0(X)\cup P$, and then $\mathcal{L}_3\subset\mathbf{Y}$. Let $z$ be the index of such a row of $H_k(s'',\cdot)$, contained in $\mathcal{L}_3$. The connections $B^1$ and $B-d\boldsymbol{\Lambda}$ lie in $\tilde{G}_k$, and they coincide on the bonds that meet $\mathcal{L}_3$, because these bonds meet $\mathbf{Y}$. Property~\ref{it:renormalisation-group-8}, applied to the set of the cubes of kind $\pi_k$ in $P_0(X)\cup P$ (the family $s''$ vanishes at every other cube of kind $\pi_k$) and to the link-component $\mathcal{L}_3$ of its union, and then property~\ref{it:renormalisation-group-13} give
\[
H_k(s'',B^1)(z,\cdot)=H_k(s'',B-d\boldsymbol{\Lambda})(z,\cdot)=\exp\left(i\,e_k\cdot\boldsymbol{\Lambda}(z)\right)\cdot\left(H_k(s'',B)\cdot\exp\left(-i\,e_k\cdot\boldsymbol{\Lambda}^u\right)\right)(z,\cdot).
\]
The same two properties, applied to $\Gamma_k$ at $t=0$, where $B=\mathscr{T}$ and $B^1=\mathscr{T}^1$, give for the index $y$ of such a row of $\Gamma_k(s'',\cdot)$
\[
\Gamma_k(s'',\mathscr{T}^1)(y,\cdot)=\Gamma_k(s'',\mathscr{T}-d\boldsymbol{\Lambda})(y,\cdot)=\exp\left(i\,e_k\cdot\boldsymbol{\Lambda}(y)\right)\cdot\left(\Gamma_k(s'',\mathscr{T})\cdot\exp\left(-i\,e_k\cdot\boldsymbol{\Lambda}^u\right)\right)(y,\cdot).
\]
At level $0$ the properties~\ref{it:renormalisation-group-8} and~\ref{it:renormalisation-group-13} of the families $H_0(s'',\cdot)$ and $\Gamma_0(s'',\cdot)$ are contained in property~\ref{it:renormalisation-group-17}. For $\bar{H}_k$ the phases are the opposite ones:
\[
\bar{H}_k(s'',B^1)(z,\cdot)=\exp\left(-i\,e_k\cdot\boldsymbol{\Lambda}(z)\right)\cdot\left(\bar{H}_k(s'',B)\cdot\exp\left(i\,e_k\cdot\boldsymbol{\Lambda}^u\right)\right)(z,\cdot).
\]
Indeed, $\bar{H}_k(s'',B)=\mathsf{C}^{-1}\cdot H_k(s'',-B)\cdot\mathsf{C}$ by~\eqref{eq:renormalisation-group-56}; the connections $-B$, $-B^1$ and $-B+d\boldsymbol{\Lambda}$ lie in $\tilde{G}_k$ (Lemma~\ref{lem:symmetries-domains}); $-B^1=-B-d(-\boldsymbol{\Lambda})$ on the bonds that meet $\mathbf{Y}$; and multiplication by a phase commutes with $\mathsf{C}$. So the identity for $H_k$, with $-B$, $-B^1$ and $-\boldsymbol{\Lambda}$ in place of $B$, $B^1$ and $\boldsymbol{\Lambda}$, gives the assertion. Let now $k\ge1$, let $X\subset\mathbf{Y}$ be a polymer, and let $s\in\mathrm{Pd}_k$ again be an arbitrary family that vanishes at every parameter cube not contained in $\mathbf{Y}$. The bonds that meet $X$ meet $\mathbf{Y}$, so that properties~\ref{it:renormalisation-group-12} and~\ref{it:renormalisation-group-13} give
\[
M_k(X,B^1)=M_k(X,B-d\boldsymbol{\Lambda})=\exp\left(i\,e_k\cdot\boldsymbol{\Lambda}^u\right)\cdot M_k(X,B)\cdot\exp\left(-i\,e_k\cdot\boldsymbol{\Lambda}^u\right);
\]
consequently $\left\langle\bar{\Psi}^{\boldsymbol{\Lambda}},M_k(X,B^1)\Psi^{\boldsymbol{\Lambda}}\right\rangle=\left\langle\bar{\Psi},M_k(X,B)\Psi\right\rangle$.

\emph{H\"older derivatives.} Let $(x,x')$ be a pair admissible for a polymer $X\subset\mathbf{Y}$. By~\eqref{eq:renormalisation-group-adm}, the bonds of the six paths $\Gamma^{\pi}(x,x')$ meet $X$, hence $\mathbf{Y}$, and $B^1=B-d\boldsymbol{\Lambda}$ on them. The sum of $d\boldsymbol{\Lambda}$ along each of the six paths telescopes (Section~\ref{sec:elementary-properties-field}); therefore
\[
\eta_k\cdot(\tau B^1)(x,x')=\eta_k\cdot(\tau B)(x,x')-\boldsymbol{\Lambda}(x')+\boldsymbol{\Lambda}(x).
\]
By the definition of $\delta_{\alpha,B}$ in Definition~\ref{def:pairs-sites-holder} it follows that, for every spinor-valued function $f$ on the sites of $T_k$,
\begin{align*}
\left(\delta_{\alpha,B^1}\left(\exp\left(i\,e_k\cdot\boldsymbol{\Lambda}\right)f\right)\right)(x,x')&=\exp\left(i\,e_k\cdot\boldsymbol{\Lambda}(x)\right)\cdot\left(\delta_{\alpha,B}f\right)(x,x'),\\
\left(\delta_{\alpha,-B^1}\left(\exp\left(-i\,e_k\cdot\boldsymbol{\Lambda}\right)f\right)\right)(x,x')&=\exp\left(-i\,e_k\cdot\boldsymbol{\Lambda}(x)\right)\cdot\left(\delta_{\alpha,-B}f\right)(x,x');
\end{align*}
the second identity is the first one with $-B$, $-B^1$ and $-\boldsymbol{\Lambda}$ in place of $B$, $B^1$ and $\boldsymbol{\Lambda}$. The value of a H\"older derivative at $(x,x')$ depends on the values of $f$ at $x$ and at $x'$ only.

\emph{Fermion fields.} Let $(X,P)$ be a term of~\eqref{eq:renormalisation-group-57} with $Y'(X,P)=\mathbf{Y}$, let $s\in\mathrm{Pd}_k$ vanish at every parameter cube outside $P_0(X)\cup P$, and let $z$ be a site of $X$ or one of the two sites of a pair admissible for $X$; by Steps~(2) and~(3) of the proof of Lemma~\ref{lem:renormalisation-group-1}, the rows at $z$ of the primitives that occur in this term are among those treated above. Let $v$ be one of the spinor-valued functions $H_k(s'',B)\Psi$ and $H_k(s'',\mathscr{T})\cdot\Gamma_k(s'',\mathscr{T})\,W$, and let $v^1$ be the same function formed with $B^1$, $\mathscr{T}^1$, $\Psi^{\boldsymbol{\Lambda}}$, $W^{\boldsymbol{\Lambda}}$ in place of $B$, $\mathscr{T}$, $\Psi$, $W$. The identities for the primitives give
\[
v^1(z)=\exp\left(i\,e_k\cdot\boldsymbol{\Lambda}(z)\right)\cdot v(z):
\]
the factor $\exp\left(-i\,e_k\cdot\boldsymbol{\Lambda}^u\right)$ to the right of the row of $H_k$ cancels the phase of $\Psi^{\boldsymbol{\Lambda}}$, resp.\ the factor $\exp\left(i\,e_k\cdot\boldsymbol{\Lambda}(y)\right)$ of the row of $\Gamma_k$ at each column $y$, and the factor to the right of the row of $\Gamma_k$ cancels the phase of $W^{\boldsymbol{\Lambda}}$. With the first identity for the H\"older derivatives it follows that
\[
\left(\delta_{\alpha,B^1}v^1\right)(x,x')=\exp\left(i\,e_k\cdot\boldsymbol{\Lambda}(x)\right)\cdot\left(\delta_{\alpha,B}v\right)(x,x')
\]
for every pair $(x,x')$ admissible for $X$, where $B=\mathscr{T}$ and $B^1=\mathscr{T}^1$ for the second function. The identity for $\bar{H}_k$ and the second identity for the H\"older derivatives give the same two relations, with $\exp\left(-i\,e_k\cdot\boldsymbol{\Lambda}\right)$ in place of $\exp\left(i\,e_k\cdot\boldsymbol{\Lambda}\right)$ and $\delta_{\alpha,-B}$, $\delta_{\alpha,-B^1}$ in place of $\delta_{\alpha,B}$, $\delta_{\alpha,B^1}$, for the barred functions $\bar{H}_k(s'',B)\bar{\Psi}$ and $\bar{H}_k(s'',\mathscr{T})\bar{W}$; on $\bar{W}$ the operator $\tilde{\Gamma}_k(s'',\mathscr{T})$ is the identity. Now $J_k^{\#}\left(s'',\mathscr{T},\mathscr{Z}(s')\right)$ is the difference of $H_k^{\#}(s'',B)$ at $t=1$ and at $t=0$, and $W_k^{\#}(s'',\mathscr{T})=H_k^{\#}(s'',\mathscr{T})\cdot\tilde{\Gamma}_k(s'',\mathscr{T})\,W$, by~\eqref{eq:renormalisation-group-56}. Hence the fields $J_k^{\#}\Psi$ and $W_k^{\#}$ are covariant at the sites of $X$ and at the pairs admissible for $X$, for every term $(X,P)$ of~\eqref{eq:renormalisation-group-57} with $Y'(X,P)=\mathbf{Y}$ and every family $s$ that vanishes outside $P_0(X)\cup P$.

\emph{The operators $T^{(i)}$.} Let $k\ge1$, and let $y$ be a unit site of $\mathbf{Y}$. By Lemma~\ref{lem:interior-boundary-parts}(b), the rows at $y$ of $T^{(1)}(\cdot)$ and $T^{(2)}(\cdot)$ depend on the connection through its values on the bonds with an end point in $B_k(y)$ only, and $B_k(y)\subset\mathbf{Y}$; on these bonds $\mathscr{T}^1=\mathscr{T}-d\boldsymbol{\Lambda}$, by the case $t=0$ above. Proposition~\ref{prop:symmetries-fermion-operators}, applied to the connections $\mathscr{T}$ and $\mathscr{T}-d\boldsymbol{\Lambda}$ of $\tilde{G}_k$, and the definition of $u_B$ in Section~\ref{sec:notation} give, for $x\in B_k(y)$ and every bond $(x,x')$,
\[
\phi_{\mathscr{T}^1,y}(x)=\exp\left(i\,e_k\cdot\left(\boldsymbol{\Lambda}(y)-\boldsymbol{\Lambda}(x)\right)\right)\cdot\phi_{\mathscr{T},y}(x),\qquad u_{\mathscr{T}^1}(x,x')=\exp\left(-i\,e_k\cdot\left(\boldsymbol{\Lambda}(x')-\boldsymbol{\Lambda}(x)\right)\right)\cdot u_{\mathscr{T}}(x,x').
\]
Hence each of the three kinds of terms of~\eqref{eq:renormalisation-group-58} acquires the factor $\exp\left(i\,e_k\cdot\boldsymbol{\Lambda}(y)\right)$ when $(\mathscr{T},\psi,\chi)$ is replaced by $(\mathscr{T}^1,\psi^{\boldsymbol{\Lambda}},\chi^{\boldsymbol{\Lambda}})$: for $x\in B_k(y)$,
\begin{align*}
\phi_{\mathscr{T}^1,y}(x)\cdot\psi^{\boldsymbol{\Lambda}}(x)&=\exp\left(i\,e_k\cdot\boldsymbol{\Lambda}(y)\right)\cdot\phi_{\mathscr{T},y}(x)\cdot\psi(x),\\
\phi_{\mathscr{T}^1,y}(x)\cdot u_{\mathscr{T}^1}\left(x,x^{\boldsymbol{\sigma}}\right)\cdot\psi^{\boldsymbol{\Lambda}}\left(x^{\boldsymbol{\sigma}}\right)&=\exp\left(i\,e_k\cdot\boldsymbol{\Lambda}(y)\right)\cdot\phi_{\mathscr{T},y}(x)\cdot u_{\mathscr{T}}\left(x,x^{\boldsymbol{\sigma}}\right)\cdot\psi\left(x^{\boldsymbol{\sigma}}\right),\\
\phi_{\mathscr{T}^1,y}(x)\cdot\chi^{\boldsymbol{\Lambda}}\left(x,x^{\boldsymbol{\sigma}}\right)&=\exp\left(i\,e_k\cdot\boldsymbol{\Lambda}(y)\right)\cdot\phi_{\mathscr{T},y}(x)\cdot\chi\left(x,x^{\boldsymbol{\sigma}}\right).
\end{align*}
The remaining factors of~\eqref{eq:renormalisation-group-58}, and the conditions $x^{\boldsymbol{\sigma}}\in\mathbf{Y}$, $x^{\boldsymbol{\sigma}}\notin\mathbf{Y}$ on the inner sums, do not change. Therefore
\[
\left(T^{(1)}(\mathscr{T}^1)\psi^{\boldsymbol{\Lambda}}\right)(y)+\left(T^{(2)}(\mathscr{T}^1)\chi^{\boldsymbol{\Lambda}}\right)(y)=\exp\left(i\,e_k\cdot\boldsymbol{\Lambda}(y)\right)\cdot\left[\left(T^{(1)}(\mathscr{T})\psi\right)(y)+\left(T^{(2)}(\mathscr{T})\chi\right)(y)\right],
\]
that is, $T^{(1)}\psi+T^{(2)}\chi$ is transformed like $\Psi(y)$. For the barred fields the phase is the opposite one. Indeed, $\bar{T}^{(i)}(\mathscr{T})=\mathsf{C}^{-1}\cdot T^{(i)}(-\mathscr{T})\cdot\mathsf{C}$ by~\eqref{eq:renormalisation-group-58}; the connections $-\mathscr{T}$ and $-\mathscr{T}+d\boldsymbol{\Lambda}$ lie in $\tilde{G}_k$ (Lemma~\ref{lem:symmetries-domains}); $-\mathscr{T}^1=-\mathscr{T}-d(-\boldsymbol{\Lambda})$ on the bonds with an end point in $B_k(y)$; and the phases commute with $\mathsf{C}$. So the computation just made, with $-\mathscr{T}$, $-\mathscr{T}^1$, $-\boldsymbol{\Lambda}$ in place of $\mathscr{T}$, $\mathscr{T}^1$, $\boldsymbol{\Lambda}$ and with $\mathsf{C}\,\bar{\psi}$, $\mathsf{C}\,\bar{\chi}$ in place of $\psi$, $\chi$, shows that $\left(\bar{T}^{(1)}(\mathscr{T}^1)\bar{\psi}^{\boldsymbol{\Lambda}}\right)(y)+\left(\bar{T}^{(2)}(\mathscr{T}^1)\bar{\chi}^{\boldsymbol{\Lambda}}\right)(y)$ is $\exp\left(-i\,e_k\cdot\boldsymbol{\Lambda}(y)\right)$ times $\left(\bar{T}^{(1)}(\mathscr{T})\bar{\psi}\right)(y)+\left(\bar{T}^{(2)}(\mathscr{T})\bar{\chi}\right)(y)$, like $\bar{\Psi}(y)$.

\emph{Data and bilinear forms.} Let $X\subset\mathbf{Y}$ be a polymer, and let $s\in\mathrm{Pd}_k$ be an arbitrary family that vanishes at every parameter cube not contained in $\mathbf{Y}$. By part~\eqref{it:bounds-analyticity-q1} of Lemma~\ref{lem:bounds-analyticity}, the pairs $\left(A+t\cdot\mathscr{Z}(s'),\mathscr{C}\right)$ and $\left(A^1+t\cdot\mathscr{Z}(s'),\mathscr{C}^1\right)$ lie in $\tilde{R}_k\times\mathfrak{X}_k$. Every bond at distance at most $1$ from $X$ is at distance at most $1$ from $\mathbf{Y}$, and on such bonds
\[
A^1+t\cdot\mathscr{Z}(s')=\left(A+t\cdot\mathscr{Z}(s')\right)-d\lambda,\qquad\mathscr{C}^1=\mathscr{C}-d\lambda'.
\]
Since $R$ maps $K^{x,C}_k$ into itself (Proposition~\ref{prop:symmetries-extraction}), the function $\mathscr{E}^r=R\,\mathscr{E}$ has clause~\eqref{it:renormalisation-group-1-k4}; hence $\mathscr{E}^r(X;\cdot)$ is invariant. The arguments $A$ and $\mathscr{C}$ occur neither in the forms $\langle\bar{\psi},\psi\rangle_{\square}$ nor in the function~\eqref{eq:renormalisation-group-60}. Both expressions are sums of products of one barred and one unbarred generator, with coefficients that do not depend on $(A,\mathscr{C})$, and in each such product the two generators belong to one site $x$: each of them is taken either at $x$ itself or, like $\chi\left(x,x^{\boldsymbol{\sigma}}\right)$ and $\bar{\chi}\left(x,x^{\boldsymbol{\sigma}}\right)$, at a pair whose first point is $x$, and~\eqref{it:lattice-symmetries-gauge-TG} gives such a generator the phase of $x$. Hence~\eqref{it:lattice-symmetries-gauge-TG} multiplies every product by $\exp\left(-i\,e_k\cdot\boldsymbol{\Lambda}(x)\right)\cdot\exp\left(i\,e_k\cdot\boldsymbol{\Lambda}(x)\right)=1$, and the two expressions are invariant. The field $\mathscr{Z}(s')$ is not changed, since $s$, $t$ and $\mathsf{Z}$ are fixed.

\emph{Measures.} The transformation fixes $\mathsf{Z}$, and on $(W,\bar{W})$ it acts by $W\mapsto UW$, $\bar{W}\mapsto\left(U^{-1}\right)^T\bar{W}$ with $U=\exp\left(i\,e_k\cdot\boldsymbol{\Lambda}^u\right)$; this is the situation of rule~\eqref{it:operations-preserve-invariance-Pr5}.

\emph{Conclusion.} In the remainder of the proof of~(b), $\mathbf{Y}$ is the polymer fixed at its beginning, which is arbitrary, and \emph{invariant} and \emph{covariant} mean invariant and covariant under~\eqref{it:lattice-symmetries-gauge-TG} with respect to $\mathbf{Y}$, on the set $\mathrm{Ar}(\mathbf{Y})$ of Section~\ref{sec:invariance-covariance}. The functions $\lambda$ and $\lambda'$, the pairs $(A,\mathscr{C})$ and $(A^1,\mathscr{C}^1)$ related near $\mathbf{Y}$, the fluctuation variable $\mathsf{Z}$, the family $s$ and the number $t$ fixed at the beginning are arbitrary. For an element of $\mathrm{Ar}(\mathbf{Y})$, the pair $(X,P)$ is a term of~\eqref{eq:renormalisation-group-57} with $Y'(X,P)=\mathbf{Y}$, and the family $s$ vanishes outside $P_0(X)\cup P$, hence at every parameter cube not contained in $\mathbf{Y}$, so that it is one of the families fixed at the beginning (paragraph \emph{Primitives}). Therefore the preceding paragraphs prove the following. The fields $J_k^{\#}\Psi$ and $W_k^{\#}$ are covariant on $\mathrm{Ar}(\mathbf{Y})$ at the sites of $X$ and at the pairs admissible for $X$ (paragraph \emph{Fermion fields}). For $k\ge1$, the families $T^{(1)}\psi+T^{(2)}\chi$ and $\bar{T}^{(1)}\bar{\psi}+\bar{T}^{(2)}\bar{\chi}$ are covariant at the unit sites of $\mathbf{Y}$, for pairs related near $\mathbf{Y}$, and for every polymer $X\subset\mathbf{Y}$ the form $\left\langle\bar{\Psi},M_k\left(X,\mathscr{T}+t\cdot\mathscr{Z}(s')\right)\Psi\right\rangle$ is invariant. For every polymer $X$, the function $\mathscr{E}^r(X)$ is invariant with respect to $X$ at all pairs of its domain $\tilde{R}_k\times\mathfrak{X}_k$ that are related near $X$ (clause~\eqref{it:renormalisation-group-1-k4}). The forms $\langle\bar{\psi},\psi\rangle_{\square}$ and the function~\eqref{eq:renormalisation-group-60} are invariant. The functions into which these fields are substituted contain generators at such indices only: $\mathscr{E}^r(X)$ contains $\psi$ at the sites of $X$ and $\chi$ at the pairs admissible for $X$ (clause~\eqref{it:renormalisation-group-1-k2}, which $\mathscr{E}^r$ has by Proposition~\ref{prop:symmetries-extraction}); $\langle\bar{\psi},\psi\rangle_{\square}$ contains $\psi$ at the sites of $\square$; $E^{(\le3),x}(X)$ contains $\psi$ at the sites of $X$ and $\chi$ at the pairs admissible for $X$, and $E^{(i),*}(\mathbf{Y})$ contains $\Psi(u)$ with $u$ in $\mathbf{Y}$ only (proof of Lemma~\ref{lem:renormalisation-group-1}(a) for the second localisation). By the proof of rule~\eqref{it:operations-preserve-invariance-Pr1}, covariance at the indices of the generators that the function contains suffices for a substitution; so rule~\eqref{it:operations-preserve-invariance-Pr1} applies to these substitutions.

The remaining hypotheses of the rules are satisfied as well. The hypotheses of rule~\eqref{it:operations-preserve-invariance-Pr1} on the substitution of $A+t\cdot\mathscr{Z}(s')$ into the first slot of $\mathscr{E}^r(X)$ hold: the bond function $\mathscr{Z}$ is covariant, as $s$, $t$ and $\mathsf{Z}$ are not changed; and for every element of $\mathrm{Ar}(\mathbf{Y})$ and every transform of it, the pairs $\left(A+t\cdot\mathscr{Z}(s'),\mathscr{C}\right)$ and $\left(A^1+t\cdot\mathscr{Z}(s'),\mathscr{C}^1\right)$ lie in $\tilde{R}_k\times\mathfrak{X}_k$, by part~\eqref{it:bounds-analyticity-q1} of Lemma~\ref{lem:bounds-analyticity}, and they are related near $\mathbf{Y}$, hence near $X$ (paragraph \emph{Data and bilinear forms}). For $t=0$ these are the pairs $(A,\mathscr{C})$ and $(A^1,\mathscr{C}^1)$, at which $G^{(4)}$ evaluates $\mathscr{E}^r(X)$. Rule~\eqref{it:operations-preserve-invariance-Pr2} applies to the contour integrals over the circles of the radii $\rho_1$, $\rho_2$ and $\rho_4$ in~\eqref{eq:renormalisation-group-59}: as $\rho_4\le\rho_2\le\rho_1$ (Section~\ref{sec:notation}), these circles lie in the disk $|t|\le\rho_1$, and $|t|\le\rho_1$ is the only condition that $\mathrm{Ar}(\mathbf{Y})$ imposes on $t$. Rule~\eqref{it:operations-preserve-invariance-Pr3} applies to every term $(X,P)$ of~\eqref{eq:renormalisation-group-57} with $Y'(X,P)=\mathbf{Y}$, because the bracket in~\eqref{eq:renormalisation-group-57} reads $G(\cdot,X)$ only at families that vanish outside $P_0(X)\cup P$, as stated in the first sentences of the paragraph \emph{Primitives}. Rule~\eqref{it:operations-preserve-invariance-Pr5} applies by the paragraph \emph{Measures}, since the only condition that $\mathrm{Ar}(\mathbf{Y})$ imposes on $\mathsf{Z}$ is $\mathsf{Z}\in\mathrm{Box}_k$.

Apply now the rules~\eqref{it:operations-preserve-invariance-Pr1}--\eqref{it:operations-preserve-invariance-Pr5} to the displays~\eqref{eq:renormalisation-group-56}--\eqref{eq:renormalisation-group-61}, in their order. Each step is proved for the arbitrary polymer $\mathbf{Y}$ and therefore holds for every polymer in place of $\mathbf{Y}$. The monotonicity stated in~\eqref{it:lattice-symmetries-gauge-TG} is applied to functions of polymers contained in $\mathbf{Y}$ that are invariant with respect to their own polymer, $\mathscr{E}^r(X)$ by clause~\eqref{it:renormalisation-group-1-k4} and the others by a preceding step; it shows that they are invariant with respect to $\mathbf{Y}$. This concerns $\mathscr{E}^r(X)$ and $E^{(\le3),x}(X)$ in the terms $(X,P)$ of~\eqref{eq:renormalisation-group-57}, where $X\subset Y'(X,P)=\mathbf{Y}$; $E^{(\le5),x}(X_i)$ in $\mathsf{K}(\mathbf{Y})$, where $X_i\subset\mathbf{Y}$; and $\mathsf{H}(X_i)$ in $E^{\#,x}_k(\mathbf{Y})$, where $X_i\subset\mathbf{Y}$. The functions $G^{(1)}$, $G^{(2)}$, $G^{(3)}$ of~\eqref{eq:renormalisation-group-59}, and for $k=0$ the function $G^{(1)}$ of~\eqref{eq:renormalisation-group-60}, are invariant by~\eqref{it:operations-preserve-invariance-Pr1} and~\eqref{it:operations-preserve-invariance-Pr2}. By~\eqref{it:operations-preserve-invariance-Pr3} and~\eqref{it:operations-preserve-invariance-Pr4} the localisation~\eqref{eq:renormalisation-group-57} preserves invariance, and by~\eqref{it:operations-preserve-invariance-Pr1} so does the substitution~\eqref{eq:renormalisation-group-58}; so $E^{(i),*}$ and $E^{(i),x}$, $i\le3$, are invariant. By~\eqref{it:operations-preserve-invariance-Pr1} and~\eqref{it:operations-preserve-invariance-Pr2} the functions $G^{(\le3)}$, $G^{(4)}$, $G^{(5)}$ are invariant, and by~\eqref{it:operations-preserve-invariance-Pr3} and~\eqref{it:operations-preserve-invariance-Pr4} so is $E^{(\le5),x}$. In~\eqref{eq:renormalisation-group-61} the function $\mathsf{K}$ is invariant by~\eqref{it:operations-preserve-invariance-Pr1} and~\eqref{it:operations-preserve-invariance-Pr4}, the function $\mathsf{H}$ by~\eqref{it:operations-preserve-invariance-Pr5}, and $E^{\#,x}_k$ by~\eqref{it:operations-preserve-invariance-Pr4} and~\eqref{it:operations-preserve-invariance-Pr1}, the series~\eqref{eq:renormalisation-group-61} being convergent in norm by assertion~(1) of Section~\ref{sec:activity-before-reblocking}. As $\mathbf{Y}$ is arbitrary, these steps give the invariance of $E^{(i),*}$, $E^{(i),x}$, $E^{(\le5),x}$, $\mathsf{K}$, $\mathsf{H}$ and $E^{\#,x}_k$ with respect to their polymer argument. Thus $E^{(i),x}$, $E^{(\le5),x}$, $\mathsf{K}$, $\mathsf{H}$ and $E^{\#,x}_k$ are invariant under~\eqref{it:lattice-symmetries-gauge-TG}; for $E^{\#,x}_k$ this is assertion~(b).

(d) Let $m$ be real and $\mathscr{E}\in K^x_k$. Let $s\in\mathrm{Pd}_k$, $t\in\mathbb{C}$ with $|t|\le\rho_1$, $(A,\mathscr{C})\in\tilde{R}^{-}_k\times\mathfrak{X}_k$ and $\mathsf{Z}\in\mathrm{Box}_k$, and put $\mathscr{T}:=A+\mathscr{C}$ and $B:=\mathscr{T}+t\cdot\mathscr{Z}(s')$. By part~\eqref{it:bounds-analyticity-q1} of Lemma~\ref{lem:bounds-analyticity}, $B\in\tilde{G}_k$, and by Lemma~\ref{lem:symmetries-domains} the connections $-B$, $\overline{B}$ and $-\overline{B}$ lie in $\tilde{G}_k$ as well. The map $\mathbf{K}$ of~\eqref{it:lattice-symmetries-gauge-TK} evaluates a function at the arguments $\left(\overline{s},\overline{t};-\overline{A},-\overline{\mathscr{C}},-\mathsf{Z}\right)$, replaces its coefficients by their complex conjugates, and multiplies every field by the matrix $\varkappa$, which acts on the spinor index at each site.

\emph{Photon field.} Let $O(s)$ be one of the operators $H_k(s')$ and $C_k^{\frac{1}{2}}(s')$ of the families~\eqref{it:primitives-p1} and~\eqref{it:primitives-p2}. By property~\ref{it:renormalisation-group-16}, $O(s)$ is real for real $s$. For $s\in\mathrm{Pd}_k$ put $\check{O}(s):=\overline{O\left(\overline{s}\right)}$. By property~\ref{it:renormalisation-group-2} (at level $0$, by property~\ref{it:renormalisation-group-17}), $O$ is continuous on the closed polydisk $\mathrm{Pd}_k$ and analytic in its interior; since $\check{O}$ arises from $O$ by conjugating both the variable and the value, $\check{O}$ has the same two properties. If all coordinates of $s$ are real, then $\check{O}(s)=\overline{O(s)}=O(s)$. Thus the analytic function $\check{O}-O$ vanishes at the interior points of $\mathrm{Pd}_k$ with real coordinates, which form a non-empty open subset of the real subspace. The interior of $\mathrm{Pd}_k$ is connected, so the identity theorem gives $\check{O}=O$ in the interior, and by continuity $\check{O}=O$ on all of $\mathrm{Pd}_k$, which is the closure of its interior. The operator $C$ is real (Proposition~\ref{prop:symmetries-reality}), and so is $\mathsf{Z}$. Writing $\mathscr{Z}(s')[\mathsf{Z}]$ for the field $\mathscr{Z}(s')$ formed with $\mathsf{Z}$, one obtains
\[
\overline{\mathscr{Z}\left(\overline{s'}\right)[-\mathsf{Z}]}=-\mathscr{Z}(s')[\mathsf{Z}].
\]
Hence, at the arguments at which $\mathbf{K}$ evaluates, the connection $B=\mathscr{T}+t\cdot\mathscr{Z}(s')$ is replaced by
\[
-\overline{A}-\overline{\mathscr{C}}+\overline{t}\cdot\mathscr{Z}\left(\overline{s'}\right)[-\mathsf{Z}]=-\overline{\mathscr{T}}-\overline{t\cdot\mathscr{Z}(s')[\mathsf{Z}]}=-\overline{B}.
\]

\emph{Fermion operators.} Property~\ref{it:renormalisation-group-16}, read at the connection $-B$, states that
\[
\overline{O\left(\overline{s},-\overline{B}\right)}=\varkappa^{-1}\cdot O(s,B)\cdot\varkappa
\]
for every member $O(s,\cdot)$ of the families~\eqref{it:primitives-p3} and~\eqref{it:primitives-p4}, the parameter being absent for $M_k(X,\cdot)$. With $\varkappa^T=\varkappa=\varkappa^{-1}$ (Clifford relations of Section~\ref{sec:recalled-auxiliary-results}) and the definition of $\mathbf{K}$ it follows that, for $k\ge1$ and every polymer $X$,
\[
\mathbf{K}\left\langle\bar{\Psi},M_k(X,B)\Psi\right\rangle=\sum_{y,y'\in T^0_k}\left(\varkappa\,\bar{\Psi}(y)\right)^T\cdot\varkappa^{-1}\cdot M_k(X,B)(y,y')\cdot\varkappa\cdot\varkappa\,\Psi(y')=\left\langle\bar{\Psi},M_k(X,B)\Psi\right\rangle,
\]
and that
\begin{align*}
\mathbf{K}\left(H_k(s'',B)\Psi\right)&=\varkappa^{-1}\cdot H_k(s'',B)\cdot\varkappa\cdot\varkappa\,\Psi=\varkappa\cdot H_k(s'',B)\Psi,\\
\mathbf{K}\left(\Gamma_k(s'',\mathscr{T})\,W\right)&=\varkappa^{-1}\cdot\Gamma_k(s'',\mathscr{T})\cdot\varkappa\cdot\varkappa\,W=\varkappa\cdot\Gamma_k(s'',\mathscr{T})\,W;
\end{align*}
so $H_k(s'',B)\Psi$ is covariant, and $\mathbf{K}\left(H_k(s'',\mathscr{T})\cdot\Gamma_k(s'',\mathscr{T})\,W\right)=\varkappa^{-1}\cdot H_k(s'',\mathscr{T})\cdot\Gamma_k(s'',\mathscr{T})\,W=\varkappa\cdot H_k(s'',\mathscr{T})\cdot\Gamma_k(s'',\mathscr{T})\,W$. For $\bar{H}_k$: by its definition in~\eqref{eq:renormalisation-group-56}, since $\mathsf{C}$ is real, by property~\ref{it:renormalisation-group-16} at the connection $B$, and by $\varkappa\cdot\mathsf{C}=-\mathsf{C}\cdot\varkappa$ (Clifford relations),
\begin{align*}
\overline{\bar{H}_k\left(\overline{s''},-\overline{B}\right)}&=\mathsf{C}^{-1}\cdot\overline{H_k\left(\overline{s''},\overline{B}\right)}\cdot\mathsf{C}=\mathsf{C}^{-1}\cdot\varkappa^{-1}\cdot H_k(s'',-B)\cdot\varkappa\cdot\mathsf{C}\\
&=\varkappa^{-1}\cdot\mathsf{C}^{-1}\cdot H_k(s'',-B)\cdot\mathsf{C}\cdot\varkappa=\varkappa^{-1}\cdot\bar{H}_k(s'',B)\cdot\varkappa,
\end{align*}
whence
\[
\mathbf{K}\left(\bar{H}_k(s'',B)\bar{\Psi}\right)=\varkappa\cdot\bar{H}_k(s'',B)\bar{\Psi},\qquad\mathbf{K}\left(\bar{H}_k(s'',\mathscr{T})\bar{W}\right)=\varkappa\cdot\bar{H}_k(s'',\mathscr{T})\bar{W},
\]
so that these functions are covariant as well. Since $e_k$ and $\eta_k$ are real and $\tau$ is linear with real coefficients,
\[
\overline{\exp\left(i\,e_k\eta_k\cdot\left(\tau\left(-\overline{B}\right)\right)(x,x')\right)}=\exp\left(i\,e_k\eta_k\cdot(\tau B)(x,x')\right),
\]
and the same holds with $-B$ in place of $B$. As $|x-x'|^{-\alpha}$ is real, the phases of the H\"older derivatives, $\delta_{\alpha,B}$ on the unbarred and $\delta_{\alpha,-B}$ on the barred fields, are invariant. The operator $J_k^{\#}\left(s'',\mathscr{T},\mathscr{Z}(s')\right)$ is the difference of $H_k^{\#}(s'',B)$ at $t=1$ and at $t=0$, and $W_k^{\#}(s'',\mathscr{T})=H_k^{\#}(s'',\mathscr{T})\cdot\tilde{\Gamma}_k(s'',\mathscr{T})\,W$; hence $J_k^{\#}\Psi$ and $W_k^{\#}$ are covariant.

\emph{The operators $T^{(i)}$.} Let $k\ge1$. By Proposition~\ref{prop:symmetries-fermion-operators}, $\overline{\phi_{-\overline{B},y}(x)}=\phi_{B,y}(x)$; by the definition of $u_B$, $\overline{u_{-\overline{B}}(x,x')}=u_B(x,x')$; by the Clifford relations, $\overline{P^{\boldsymbol{\sigma}}_{\nu}}=\varkappa^{-1}\cdot P^{\boldsymbol{\sigma}}_{\nu}\cdot\varkappa$; and the numbers $b_k$, $\mathfrak{m}$ and $\eta_k>0$ are real. Inserting these relations into~\eqref{eq:renormalisation-group-58} one obtains
\[
\overline{T^{(i)}\left(-\overline{B}\right)}=\varkappa^{-1}\cdot T^{(i)}(B)\cdot\varkappa\qquad(i=1,2),
\]
and therefore, by the definition of $\bar{T}^{(i)}$, the reality of $\mathsf{C}$ and $\varkappa\cdot\mathsf{C}=-\mathsf{C}\cdot\varkappa$,
\begin{align*}
\overline{\bar{T}^{(i)}\left(-\overline{B}\right)}&=\mathsf{C}^{-1}\cdot\overline{T^{(i)}\left(\overline{B}\right)}\cdot\mathsf{C}=\mathsf{C}^{-1}\cdot\varkappa^{-1}\cdot T^{(i)}(-B)\cdot\varkappa\cdot\mathsf{C}\\
&=\varkappa^{-1}\cdot\mathsf{C}^{-1}\cdot T^{(i)}(-B)\cdot\mathsf{C}\cdot\varkappa=\varkappa^{-1}\cdot\bar{T}^{(i)}(B)\cdot\varkappa\qquad(i=1,2).
\end{align*}
With $B=\mathscr{T}$ this shows that the substitution~\eqref{eq:renormalisation-group-58} is covariant: $\left(T^{(1)}(\mathscr{T})\psi\right)(y)+\left(T^{(2)}(\mathscr{T})\chi\right)(y)$ is transformed like $\Psi(y)$, and $\left(\bar{T}^{(1)}(\mathscr{T})\bar{\psi}\right)(y)+\left(\bar{T}^{(2)}(\mathscr{T})\bar{\chi}\right)(y)$ like $\bar{\Psi}(y)$.

\emph{Level $0$.} Let $k=0$, and let $b=\left(x,x^{\boldsymbol{\sigma}}\right)$ with $x\in\square$. By the identity for $\mathscr{Z}$ above and the Clifford relations,
\[
\overline{P^{\pm\boldsymbol{\sigma}}_{\nu}\cdot\left(\exp\left(\pm i\,e_k\cdot\mathscr{Z}\left(\overline{s'}\right)[-\mathsf{Z}](b)\right)-1\right)}=\varkappa^{-1}\cdot P^{\pm\boldsymbol{\sigma}}_{\nu}\cdot\varkappa\cdot\left(\exp\left(\pm i\,e_k\cdot\mathscr{Z}(s')[\mathsf{Z}](b)\right)-1\right),
\]
with the same choice of the sign at all places. For all columns $a$, $a'$ of four elements of the Grassmann algebra and every $4\times4$ matrix $N$ one has $\left(\varkappa\,a\right)^T\cdot\varkappa^{-1}\cdot N\cdot\varkappa\cdot\left(\varkappa\,a'\right)=a^T\cdot N\cdot a'$, since $\varkappa^T=\varkappa=\varkappa^{-1}$; hence each term of the two sums of~\eqref{eq:renormalisation-group-60} is invariant. The families of property~\ref{it:renormalisation-group-17} have property~\ref{it:renormalisation-group-16}; so the relations for the primitives used above hold at level $0$ as well.

\emph{Data.} The number $m'=m+m(\mathscr{E})$ is real, because $m$ is real and $m(\mathscr{E})$ is real for $\mathscr{E}\in K^x_k$ (Proposition~\ref{prop:symmetries-extraction}). Since $R$ maps $K^x_k$ into itself (Proposition~\ref{prop:symmetries-extraction}), the function $\mathscr{E}^r$ has clause~\eqref{it:renormalisation-group-1-k7}, that is, $\mathbf{K}\,\mathscr{E}^r(X)=\mathscr{E}^r(X)$. Finally $\left(\varkappa\,\bar{\psi}\right)^T\cdot\left(\varkappa\,\psi\right)=\bar{\psi}^T\cdot\psi$, so that the forms $\langle\bar{\psi},\psi\rangle_{\square}$ are invariant.

\emph{Measures.} The map $\mathsf{Z}\mapsto-\mathsf{Z}$ preserves $\hat{\mu}$ (Proposition~\ref{prop:symmetries-reality}), and on $(W,\bar{W})$ the transformation acts by $W\mapsto UW$, $\bar{W}\mapsto\left(U^{-1}\right)^T\bar{W}$ with $U=\varkappa=\left(U^{-1}\right)^T$.

\emph{Conclusion.} Apply the rules~\eqref{it:operations-preserve-invariance-Pr1}--\eqref{it:operations-preserve-invariance-Pr5} to the displays~\eqref{eq:renormalisation-group-56}--\eqref{eq:renormalisation-group-61}, in their order: the functions $G^{(1)}$, $G^{(2)}$, $G^{(3)}$, for $k=0$ the function~\eqref{eq:renormalisation-group-60}, and $G^{(\le3)}$, $G^{(4)}$, $G^{(5)}$ are invariant by~\eqref{it:operations-preserve-invariance-Pr1} and~\eqref{it:operations-preserve-invariance-Pr2}; the localisation~\eqref{eq:renormalisation-group-57} preserves invariance by~\eqref{it:operations-preserve-invariance-Pr3} and~\eqref{it:operations-preserve-invariance-Pr4}, the functions $G(s,X)$ to which it is applied being analytic in $s$ in the interior of $\mathrm{Pd}_k$, by part~\eqref{it:bounds-analyticity-q5} of Lemma~\ref{lem:bounds-analyticity} (see the proof of Proposition~\ref{prop:bounds-regularity-localised}(b), (c)); the substitution~\eqref{eq:renormalisation-group-58} preserves invariance by~\eqref{it:operations-preserve-invariance-Pr1}; $\mathsf{K}$ is invariant by~\eqref{it:operations-preserve-invariance-Pr1} and~\eqref{it:operations-preserve-invariance-Pr4}, $\mathsf{H}$ by~\eqref{it:operations-preserve-invariance-Pr5}, and $E^{\#,x}_k$ by~\eqref{it:operations-preserve-invariance-Pr4} and~\eqref{it:operations-preserve-invariance-Pr1}, the series~\eqref{eq:renormalisation-group-61} being convergent in norm by assertion~(1) of Section~\ref{sec:activity-before-reblocking}. Thus $E^{(i),x}$, $E^{(\le5),x}$, $\mathsf{K}$, $\mathsf{H}$ and $E^{\#,x}_k$ are invariant under~\eqref{it:lattice-symmetries-gauge-TK}.
\end{proof}

\begin{proof}[Proof of Lemma~\ref{lem:renormalisation-group-2}(c)]
The transformation~\eqref{it:lattice-symmetries-gauge-TC} leaves $X$, $P$, $s$ and $t$ unchanged, replaces $(A,\mathscr{C},\mathsf{Z})$ by $(-A,-\mathscr{C},-\mathsf{Z})$, and acts on the generators by $\mathbf{C}$. The same formulas apply to families of odd elements. Let $v^{\#}$ be a family, indexed like $\psi^{\#}=(\psi,\chi)$, of odd elements of a Grassmann algebra that contains $\mathrm{Gr}[\psi,\chi]$, and let $v_0$, $v_1$, $w_0$, $w_1$ be its components in the places of $\psi$, $\bar{\psi}$, $\chi$, $\bar{\chi}$, respectively. Then $\mathbf{C}\,v^{\#}$ denotes the family whose components in these four places are
\[
\mathsf{C}\cdot v_1,\qquad-\mathsf{C}\cdot v_0,\qquad\mathsf{C}\cdot w_1,\qquad-\mathsf{C}\cdot w_0;
\]
for $v^{\#}=\psi^{\#}$ this is the action of~\eqref{it:lattice-symmetries-gauge-TC} on the generators. A family $(v_0,v_1)$ indexed like $(\Psi,\bar{\Psi})$ is transformed into $(\mathsf{C}\cdot v_1,-\mathsf{C}\cdot v_0)$.

\emph{Step~1.} For every polymer $X$ and every $i\le3$ the elements $E^{(i),x}(X;A,\mathscr{C},\mathsf{Z},\psi^{\#})$ and $\mathscr{E}^r(X)$, and the elements $m'\cdot\langle\bar{\psi},\psi\rangle_{\square}$ for the $M$-cubes $\square$, are invariant under~\eqref{it:lattice-symmetries-gauge-TC}.

\emph{Bilinear forms.} Let $N$ be an operator on the spinor functions on $T^0_k$ with the kernel $N_{ab}(y,y')$, and let $N^T$ be its transpose, $N^T_{ba}(y',y):=N_{ab}(y,y')$. At the transformed generators the element $\left\langle\bar{\Psi},N\cdot\Psi\right\rangle$ becomes $\left\langle-\mathsf{C}\cdot\Psi,\,N\cdot\mathsf{C}\cdot\bar{\Psi}\right\rangle$. The components of $\mathsf{C}\cdot\Psi$ and of $\mathsf{C}\cdot\bar{\Psi}$ are odd elements, so that any two of them anticommute. Together with $\mathsf{C}^T=\mathsf{C}^{-1}$ this gives
\begin{align*}
\left\langle-\mathsf{C}\cdot\Psi,\,N\cdot\mathsf{C}\cdot\bar{\Psi}\right\rangle&=-\sum_{y,y'\in T^0_k}\ \sum_{a,b}\left(\mathsf{C}\cdot\Psi\right)_a(y)\cdot N_{ab}(y,y')\cdot\left(\mathsf{C}\cdot\bar{\Psi}\right)_b(y')\\
&=\sum_{y,y'\in T^0_k}\ \sum_{a,b}\left(\mathsf{C}\cdot\bar{\Psi}\right)_b(y')\cdot N^T_{ba}(y',y)\cdot\left(\mathsf{C}\cdot\Psi\right)_a(y)=\left\langle\bar{\Psi},\,\mathsf{C}^{-1}\cdot N^T\cdot\mathsf{C}\cdot\Psi\right\rangle.
\end{align*}
Let $s\in\mathrm{Pd}_k$ and $t\in\mathbb{C}$ with $|t|\le\rho_1$, and put $B:=\mathscr{T}+t\cdot\mathscr{Z}(s')$. By part~\eqref{it:bounds-analyticity-q1} of Lemma~\ref{lem:bounds-analyticity} and by Lemma~\ref{lem:symmetries-domains}, both $B$ and $-B$ lie in $\tilde{G}_k$. Since $\mathscr{Z}(s')$ is linear in $\mathsf{Z}$, the transformation replaces $B$ by $-B$. Property~\ref{it:renormalisation-group-15}, read at the connection $-B$, gives
\[
M_k(X,-B)^T=\mathsf{C}^{-1}\cdot M_k(X,B)\cdot\mathsf{C} .
\]
With $N=M_k(X,-B)$ and $\mathsf{C}^2=\mathsf{C}^{-2}=-1$ the preceding identity becomes
\[
\left\langle-\mathsf{C}\cdot\Psi,\,M_k(X,-B)\cdot\mathsf{C}\cdot\bar{\Psi}\right\rangle=\left\langle\bar{\Psi},\,\mathsf{C}^{-2}\cdot M_k(X,B)\cdot\mathsf{C}^2\cdot\Psi\right\rangle=\left\langle\bar{\Psi},\,M_k(X,B)\cdot\Psi\right\rangle .
\]
So, for $k\ge1$, the integrand of $G^{(1)}(s,X)$ in~\eqref{eq:renormalisation-group-59} is invariant. The same computation, made for the pairing over an $M$-cube $\square$ and with the identity operator in the place of $N$, gives
\[
\left\langle-\mathsf{C}\cdot\psi,\,\mathsf{C}\cdot\bar{\psi}\right\rangle_{\square}=\left\langle\bar{\psi},\,\mathsf{C}^{-1}\cdot\mathsf{C}\cdot\psi\right\rangle_{\square}=\langle\bar{\psi},\psi\rangle_{\square};
\]
as $m'$ is a number, $m'\cdot\langle\bar{\psi},\psi\rangle_{\square}$ is invariant.

\emph{Fermion fields.} By the definition $\bar{H}_k(s'',B)=\mathsf{C}^{-1}\cdot H_k(s'',-B)\cdot\mathsf{C}$ in~\eqref{eq:renormalisation-group-56}, and since $\mathsf{C}^{-1}=-\mathsf{C}$,
\[
H_k(s'',-B)\cdot\mathsf{C}=\mathsf{C}\cdot\bar{H}_k(s'',B),\qquad\bar{H}_k(s'',-B)\cdot\mathsf{C}=\mathsf{C}^{-1}\cdot H_k(s'',B)\cdot\mathsf{C}^2=\mathsf{C}\cdot H_k(s'',B);
\]
consequently
\[
H_k(s'',-B)\cdot\left(\mathsf{C}\cdot\bar{\Psi}\right)=\mathsf{C}\cdot\left(\bar{H}_k(s'',B)\cdot\bar{\Psi}\right),\qquad\bar{H}_k(s'',-B)\cdot\left(-\mathsf{C}\cdot\Psi\right)=-\mathsf{C}\cdot\left(H_k(s'',B)\cdot\Psi\right).
\]
Let $H^{\#}_k(s'',B)\Psi$ denote the family, indexed like $\psi^{\#}$, with the components
\[
H_k(s'',B)\cdot\Psi,\qquad\bar{H}_k(s'',B)\cdot\bar{\Psi},\qquad\delta_{\alpha,B}H_k(s'',B)\cdot\Psi,\qquad\delta_{\alpha,-B}\bar{H}_k(s'',B)\cdot\bar{\Psi}
\]
in the places of $\psi$, $\bar{\psi}$, $\chi$, $\bar{\chi}$. The H\"older derivative is formed with the connection on the unbarred and with its negative on the barred components, so that the replacement of $B$ by $-B$ interchanges the two H\"older derivatives; and $\mathsf{C}$, which acts on the spinor index only, commutes with both of them. Together with the two identities for $H_k(s'',-B)\cdot\left(\mathsf{C}\cdot\bar{\Psi}\right)$ and $\bar{H}_k(s'',-B)\cdot\left(-\mathsf{C}\cdot\Psi\right)$ this gives
\[
H^{\#}_k(s'',-B)\left(\mathbf{C}\,\Psi\right)=\mathbf{C}\left(H^{\#}_k(s'',B)\Psi\right).
\]
This holds for $B=\mathscr{T}$ and for $B=\mathscr{T}+\mathscr{Z}(s')$ (the values $t=0$ and $t=1$), which the transformation replaces by $-\mathscr{T}$ and by $-\mathscr{T}-\mathscr{Z}(s')$. Since $J^{\#}_k\left(s'',\mathscr{T},\mathscr{Z}(s')\right)=H^{\#}_k\left(s'',\mathscr{T}+\mathscr{Z}(s')\right)-H^{\#}_k(s'',\mathscr{T})$ by~\eqref{eq:renormalisation-group-56}, the family $J^{\#}_k\left(s'',\mathscr{T},\mathscr{Z}(s')\right)\Psi$ is covariant in the sense of Section~\ref{sec:invariance-covariance}.

\emph{The substitution~\eqref{eq:renormalisation-group-58}.} Let $k\ge1$, and let the operators $T^{(i)}$, $\bar{T}^{(i)}$ be formed with a polymer $\mathbf{Y}$. Their definition $\bar{T}^{(i)}(B)=\mathsf{C}^{-1}\cdot T^{(i)}(-B)\cdot\mathsf{C}$ in~\eqref{eq:renormalisation-group-58} gives in the same way
\[
T^{(i)}(-B)\cdot\mathsf{C}=\mathsf{C}\cdot\bar{T}^{(i)}(B),\qquad\bar{T}^{(i)}(-B)\cdot\mathsf{C}=\mathsf{C}\cdot T^{(i)}(B)\qquad(i=1,2).
\]
Hence, at the transformed arguments, the elements that the substitution~\eqref{eq:renormalisation-group-58} inserts at a unit site $y$ of $\mathbf{Y}$ in the places of $\Psi(y)$ and of $\bar{\Psi}(y)$ are
\begin{gather*}
\left(T^{(1)}(-\mathscr{T})\cdot\mathsf{C}\cdot\bar{\psi}\right)(y)+\left(T^{(2)}(-\mathscr{T})\cdot\mathsf{C}\cdot\bar{\chi}\right)(y)=\mathsf{C}\cdot\left[\left(\bar{T}^{(1)}(\mathscr{T})\cdot\bar{\psi}\right)(y)+\left(\bar{T}^{(2)}(\mathscr{T})\cdot\bar{\chi}\right)(y)\right],\\
\left(\bar{T}^{(1)}(-\mathscr{T})\cdot(-\mathsf{C})\cdot\psi\right)(y)+\left(\bar{T}^{(2)}(-\mathscr{T})\cdot(-\mathsf{C})\cdot\chi\right)(y)=-\mathsf{C}\cdot\left[\left(T^{(1)}(\mathscr{T})\cdot\psi\right)(y)+\left(T^{(2)}(\mathscr{T})\cdot\chi\right)(y)\right].
\end{gather*}
These are the images under $\mathbf{C}$ of the elements inserted at the original arguments: the substitution~\eqref{eq:renormalisation-group-58} is covariant.

\emph{The datum.} The family $\mathscr{E}$ lies in $K^{x,C}_k$, and by Proposition~\ref{prop:symmetries-extraction} the map $R$ sends $K^{x,C}_k$ into itself; so $\mathscr{E}^r=R\,\mathscr{E}$ has clause~\eqref{it:renormalisation-group-1-k5} of Definition~\ref{def:renormalisation-group-1}, which states that $\mathscr{E}^r(X)$ is invariant under~\eqref{it:lattice-symmetries-gauge-TC}.

\emph{Level~$0$.} Let $k=0$. For $x\in\square$, a direction $\nu$ and $\boldsymbol{\sigma}\in\{1,-1\}$, call the summand with these indices in the first, resp.\ in the second, of the two sums in~\eqref{eq:renormalisation-group-60} the term $(x,\nu,\boldsymbol{\sigma})$ of that sum. For columns $q$, $q'$ of four odd elements and a $4\times4$ matrix $O$ put $q^T\cdot O\cdot q':=\sum_{a,b}q_a\cdot O_{ab}\cdot q'_b$, so that the expression $\bar{\psi}(x)\,O\,\psi(x')$ of Section~\ref{sec:localisation-exponent} is $\bar{\psi}(x)^T\cdot O\cdot\psi(x')$; since the entries anticommute, $q^T\cdot O\cdot q'=-{q'}^T\cdot O^T\cdot q$. The transformation replaces $\bar{\psi}(x)$, $\psi(x)$, $\bar{\chi}\left(x,x^{\boldsymbol{\sigma}}\right)$, $\chi\left(x,x^{\boldsymbol{\sigma}}\right)$ by $-\mathsf{C}\cdot\psi(x)$, $\mathsf{C}\cdot\bar{\psi}(x)$, $-\mathsf{C}\cdot\chi\left(x,x^{\boldsymbol{\sigma}}\right)$, $\mathsf{C}\cdot\bar{\chi}\left(x,x^{\boldsymbol{\sigma}}\right)$, and $\mathscr{Z}(s')$ by $-\mathscr{Z}(s')$. So the term $(x,\nu,\boldsymbol{\sigma})$ of the first sum, taken at the transformed arguments, is
\begin{multline*}
\left(-\mathsf{C}\cdot\psi(x)\right)^T\cdot P^{\boldsymbol{\sigma}}_{\nu}\cdot\left(\exp\left(-i\,e_k\cdot\mathscr{Z}(s')\left(x,x^{\boldsymbol{\sigma}}\right)\right)-1\right)\cdot\mathsf{C}\cdot\left(\bar{\psi}(x)+\bar{\chi}\left(x,x^{\boldsymbol{\sigma}}\right)\right)\\
=\left(\bar{\psi}(x)+\bar{\chi}\left(x,x^{\boldsymbol{\sigma}}\right)\right)^T\cdot\mathsf{C}^{-1}\cdot\left(P^{\boldsymbol{\sigma}}_{\nu}\right)^T\cdot\mathsf{C}\cdot\left(\exp\left(-i\,e_k\cdot\mathscr{Z}(s')\left(x,x^{\boldsymbol{\sigma}}\right)\right)-1\right)\cdot\psi(x),
\end{multline*}
by the rule just stated and $\mathsf{C}^T=\mathsf{C}^{-1}$. By the Clifford relations of Section~\ref{sec:recalled-auxiliary-results}, $\mathsf{C}^{-1}\cdot\left(P^{\boldsymbol{\sigma}}_{\nu}\right)^T\cdot\mathsf{C}=P^{-\boldsymbol{\sigma}}_{\nu}$, because $\mathsf{C}^{-1}\cdot\gamma_{\nu}^T\cdot\mathsf{C}=-\gamma_{\nu}$. Thus the term $(x,\nu,\boldsymbol{\sigma})$ of the first sum, at the transformed arguments, is the term $(x,\nu,\boldsymbol{\sigma})$ of the second sum at the original arguments. The same computation, with $\mathsf{C}^{-1}\cdot\left(P^{-\boldsymbol{\sigma}}_{\nu}\right)^T\cdot\mathsf{C}=P^{\boldsymbol{\sigma}}_{\nu}$, shows that the term $(x,\nu,\boldsymbol{\sigma})$ of the second sum, at the transformed arguments, is the term $(x,\nu,\boldsymbol{\sigma})$ of the first sum at the original arguments. Hence the function $G^{(1)}(s,\square)$ of~\eqref{eq:renormalisation-group-60} is invariant.

\emph{End of Step~1.} Let $k\ge1$. In the integrand of $G^{(2)}(s,X)$ the connection $A+t\cdot\mathscr{Z}(s')$ in the first slot of $\mathscr{E}^r$ is replaced by its negative, and the covariant family $\psi^{\#}+t\cdot J^{\#}_k\left(s'',\mathscr{T},\mathscr{Z}(s')\right)\Psi$ is substituted into the invariant function $\mathscr{E}^r(X)$. In the integrand of $G^{(3)}(s,\square)$ the covariant fields $\psi+t\cdot J\Psi$ and $\bar{\psi}+t\cdot\bar{J}\,\bar{\Psi}$, formed with the components $J$ and $\bar{J}$ of $J^{\#}_k\left(s'',\mathscr{T},\mathscr{Z}(s')\right)$ without H\"older derivative, are substituted into $m'\cdot\langle\bar{\psi},\psi\rangle_{\square}$. By~\eqref{it:operations-preserve-invariance-Pr1} these integrands are invariant, and so is the integrand of $G^{(1)}(s,X)$, as shown above; by~\eqref{it:operations-preserve-invariance-Pr2} the functions $G^{(1)}$, $G^{(2)}$, $G^{(3)}$ are invariant. By~\eqref{it:operations-preserve-invariance-Pr3} and~\eqref{it:operations-preserve-invariance-Pr4}, applied to~\eqref{eq:renormalisation-group-57}, the elements $E^{(i),*}(\mathbf{Y})$, $i\le3$, are invariant; as the substitution~\eqref{eq:renormalisation-group-58} is covariant, the elements $E^{(i),x}(\mathbf{Y})$, $i\le3$, are invariant by~\eqref{it:operations-preserve-invariance-Pr1}. For $k=0$ the same rules apply to the function $G^{(1)}(s,\square)$ of~\eqref{eq:renormalisation-group-60}, while $G^{(2)}=G^{(3)}=0$ and $\mathscr{E}^r=0$. This proves Step~1.

\emph{Step~2.} Adjoin to $\mathrm{Gr}[\psi,\chi]$ a family $\boldsymbol{\phi}^{\#}$ of further independent generators, indexed like $\psi^{\#}$. Its components in the places of $\psi$ and $\bar{\psi}$ are denoted $\boldsymbol{\phi}$ and $\bar{\boldsymbol{\phi}}$; its two other components, indexed like $\chi$ and $\bar{\chi}$, stand in the places of the H\"older derivatives of $\boldsymbol{\phi}$ and $\bar{\boldsymbol{\phi}}$. The transformation~\eqref{it:lattice-symmetries-gauge-TC} acts on $\boldsymbol{\phi}^{\#}$ by $\mathbf{C}$, as on $\psi^{\#}$. For $X\in D_k$ put
\begin{multline*}
\mathbf{F}\left(X;\psi^{\#},\boldsymbol{\phi}^{\#}\right):=E^{(\le3),x}\left(X;A,\mathscr{C},\mathsf{Z},\psi^{\#}+\boldsymbol{\phi}^{\#}\right)+\mathscr{E}^r\left(X;A,\mathscr{C},\psi^{\#}+\boldsymbol{\phi}^{\#}\right)-\mathscr{E}^r\left(X;A,\mathscr{C},\psi^{\#}\right)\\
-1_{X=\square}\cdot m'\cdot\left(\left\langle\bar{\psi}+\bar{\boldsymbol{\phi}},\,\psi+\boldsymbol{\phi}\right\rangle_{\square}-\langle\bar{\psi},\psi\rangle_{\square}\right),
\end{multline*}
where $1_{X=\square}$ equals $1$ if $X$ consists of a single $M$-cube, which is then denoted $\square$, and $0$ otherwise. Where the arguments are to be shown, this element is written $\mathbf{F}\left(X;A,\mathscr{C},\mathsf{Z};\psi^{\#},\boldsymbol{\phi}^{\#}\right)$. Let $v^{\#}$ be a family, indexed like $\psi^{\#}$, of odd elements of a Grassmann algebra that contains $\mathrm{Gr}[\psi,\chi]$. Odd elements anticommute with each other and with the generators $\psi$, $\chi$, and their squares vanish; hence there is exactly one homomorphism of algebras, linear over $\mathbb{C}$, from the enlarged algebra into this Grassmann algebra that leaves $\psi$ and $\chi$ fixed and maps every component of $\boldsymbol{\phi}^{\#}$ to the corresponding component of $v^{\#}$. The image of $\mathbf{F}\left(X;\psi^{\#},\boldsymbol{\phi}^{\#}\right)$ under this homomorphism is denoted $\mathbf{F}\left(X;\psi^{\#},v^{\#}\right)$.

The element $\mathbf{F}$ is invariant under~\eqref{it:lattice-symmetries-gauge-TC}, acting on $\psi^{\#}$ and on $\boldsymbol{\phi}^{\#}$ together:
\[
\mathbf{F}\left(X;-A,-\mathscr{C},-\mathsf{Z};\mathbf{C}\,\psi^{\#},\mathbf{C}\,\boldsymbol{\phi}^{\#}\right)=\mathbf{F}\left(X;A,\mathscr{C},\mathsf{Z};\psi^{\#},\boldsymbol{\phi}^{\#}\right).
\]
Indeed, $\mathbf{C}\,\psi^{\#}+\mathbf{C}\,\boldsymbol{\phi}^{\#}=\mathbf{C}\left(\psi^{\#}+\boldsymbol{\phi}^{\#}\right)$, so that $\psi^{\#}+\boldsymbol{\phi}^{\#}$ is a covariant family; by~\eqref{it:operations-preserve-invariance-Pr1}, its substitution into the elements $E^{(\le3),x}(X)$, $\mathscr{E}^r(X)$ and $m'\cdot\langle\bar{\psi},\psi\rangle_{\square}$, which are invariant by Step~1, gives invariant elements, and the two remaining terms $-\mathscr{E}^r\left(X;A,\mathscr{C},\psi^{\#}\right)$ and $1_{X=\square}\cdot m'\cdot\langle\bar{\psi},\psi\rangle_{\square}$ are invariant by Step~1 itself.

For fixed $s''$ the maps
\[
t\longmapsto\mathscr{E}^r\left(X;A,\mathscr{C},\psi^{\#}+t\cdot W^{\#}_k(s'',\mathscr{T})\right),\qquad t\longmapsto\left\langle\bar{\psi}+t\cdot\bar{\mathcal{W}},\,\psi+t\cdot\mathcal{W}\right\rangle_{\square}
\]
are polynomials in $t$ with values in the finite-dimensional algebra $\mathrm{Gr}[\psi,\chi,W]$. As $\rho_4>1$, the formula $\oint_{\rho_4}f=f(1)-f(0)$ of Section~\ref{sec:notation} applies to them, and the definitions of $G^{(\le3)}$, $G^{(4)}$ and $G^{(5)}$ in~\eqref{eq:renormalisation-group-59} give, for all parameters $s''$,
\[
\left(G^{(\le3)}+G^{(4)}+G^{(5)}\right)(s'',X)=\mathbf{F}\left(X;\psi^{\#},W^{\#}_k(s'',\mathscr{T})\right).
\]

In the localisation that defines $E^{(\le5),x}$ the parameter cubes are those of $\pi_k$, which are $M$-cubes. By~\eqref{it:renormalisation-group-g0}, an $M$-cube that has a lattice site in common with a union of $M$-cubes is one of the $M$-cubes of this union. Hence $P_0(X)$ is the set of the $M$-cubes of $X$, and $Y'(X,P)$ is the union of $X$ and of the cubes of $P$.

Let $U$ be a union of $M$-cubes and $X\in D_k$ with $X\subseteq U$. Let $f$ be the map $s''\mapsto\mathbf{F}\left(X;\psi^{\#},W^{\#}_k(s'',\mathscr{T})\right)$, regarded as a function of the parameters $s''_c$ of the $M$-cubes $c\subseteq U$, each of which ranges over the disk $|s''_c|\le M^{\alpha_0}$, the parameters of all other cubes being equal to $0$. Its values lie in the finite-dimensional space $\mathrm{Gr}[\psi,\chi,W]$, and its coefficients are polynomials in the entries of $H^{\#}_k(s'',\mathscr{T})$ and of $\Gamma_k(s'',\mathscr{T})$. By part~\eqref{it:bounds-analyticity-q5} of Lemma~\ref{lem:bounds-analyticity}, $f$ is continuous on this closed polydisk, hence bounded there, and analytic in its interior; moreover $M^{\alpha_0}>1$, by~\eqref{eq:renormalisation-group-30}. Lemma~\ref{lem:expansion-parameters}, applied to $f$, to the set of the $M$-cubes of $U$, to its subset $P_0(X)$ and to $\rho=M^{\alpha_0}$, gives
\[
\mathbf{F}\left(X;\psi^{\#},W^{\#}_k(1_U,\mathscr{T})\right)=\sum_Pf_P,\qquad f_P:=\int_{[0,1]^P}\mathrm{d}s''_P\,\frac{\partial}{\partial s''_P}\left[f(s'')\right]_{s''=1\text{ on }P_0(X),\ s''=0\text{ outside }P_0(X)\cup P},
\]
where $P$ ranges over the sets of $M$-cubes that are contained in $U$ and not in $X$. By the identity for $G^{(\le3)}+G^{(4)}+G^{(5)}$ just proved, $f_P$ is the term of the pair $(X,P)$ in~\eqref{eq:renormalisation-group-57} for this function. By Lemma~\ref{lem:renormalisation-group-1}(b), $f_P=0$ unless $Y'(X,P)$, the union of $X$ and of the cubes of $P$, is connected, that is, a polymer. Conversely, if a pair $(X,P)$ of~\eqref{eq:renormalisation-group-57} has $Y'(X,P)\subseteq U$, then $X\subseteq U$ and $P$ is a set of $M$-cubes contained in $U$ and not in $X$. Summation over the polymers $X\subseteq U$, followed by the collection of the pairs $(X,P)$ according to the polymer $\mathbf{Y}=Y'(X,P)$, therefore gives, by~\eqref{eq:renormalisation-group-57} and~\eqref{eq:renormalisation-group-59},
\[
\sum_{\mathbf{Y}\in D_k,\,\mathbf{Y}\subseteq U}E^{(\le5),x}(\mathbf{Y})=\sum_{X\in D_k,\,X\subseteq U}\mathbf{F}\left(X;\psi^{\#},W^{\#}_k(1_U,\mathscr{T})\right).
\]

\emph{Step~3.} Let $U$ be a union of $M$-cubes. The element
\[
\Xi_U:=\int\mathrm{d}\mu_I(W)\,\mathrm{d}\hat{\mu}(\mathsf{Z})\,\exp\left\{\sum_{\mathbf{Y}\in D_k,\,\mathbf{Y}\subseteq U}E^{(\le5),x}(\mathbf{Y})\right\}
\]
of $\mathrm{Gr}[\psi,\chi]$, which depends on $(A,\mathscr{C})$, is invariant under~\eqref{it:lattice-symmetries-gauge-TC}.

Put
\[
\boldsymbol{\phi}_{\pm}:=H_k(1_U,\pm\mathscr{T})\cdot\Gamma_k(1_U,\pm\mathscr{T})\cdot W,\qquad\bar{\boldsymbol{\phi}}_{\pm}:=\bar{H}_k(1_U,\pm\mathscr{T})\cdot\bar{W};
\]
these are spinor functions on $T_k$ whose values are linear combinations of the generators $W$. By~\eqref{eq:renormalisation-group-56}, the components of $W^{\#}_k(1_U,\pm\mathscr{T})$ in the places of $\psi$, $\bar{\psi}$, $\chi$, $\bar{\chi}$ are
\[
\boldsymbol{\phi}_{\pm},\qquad\bar{\boldsymbol{\phi}}_{\pm},\qquad\delta_{\alpha,\pm\mathscr{T}}\boldsymbol{\phi}_{\pm},\qquad\delta_{\alpha,\mp\mathscr{T}}\bar{\boldsymbol{\phi}}_{\pm}.
\]
By Step~2 the integrand of $\Xi_U$ is the image of $\exp\left\{\sum_X\mathbf{F}\left(X;\psi^{\#},\boldsymbol{\phi}^{\#}\right)\right\}$, the sum extending over the polymers $X\subseteq U$, under the homomorphism that leaves $\psi$ and $\chi$ fixed and replaces $\boldsymbol{\phi}^{\#}$ by $W^{\#}_k(1_U,\mathscr{T})$. This homomorphism sends every component of $\boldsymbol{\phi}^{\#}$ to a linear combination of the generators $W$, and the H\"older derivatives $\delta_{\alpha,\mathscr{T}}\boldsymbol{\phi}_{+}$ and $\delta_{\alpha,-\mathscr{T}}\bar{\boldsymbol{\phi}}_{+}$ are linear images of $\boldsymbol{\phi}_{+}$ and $\bar{\boldsymbol{\phi}}_{+}$. By Lemma~\ref{lem:substitutions-integral}(a), which rests on Wick's theorem, the integral over $W$ of such an image is therefore determined by the two-point functions of $\boldsymbol{\phi}_{+}$ and $\bar{\boldsymbol{\phi}}_{+}$. For two columns $v=\left(v_a(x)\right)$ and $v'=\left(v'_b(x')\right)$ of linear combinations of the generators $W$, each indexed by the pairs of a site and a spinor index, where the sites are those of $T_k$ or those of $T^0_k$, let $\left\langle v\,{v'}^T\right\rangle$ be the matrix with the entries $\int\mathrm{d}\mu_I(W)\,v_a(x)\cdot v'_b(x')$, its rows indexed by the pairs $(x,a)$ and its columns by the pairs $(x',b)$. The spinor functions $\boldsymbol{\phi}_{\pm}$, $\bar{\boldsymbol{\phi}}_{\pm}$ on $T_k$ are such columns; so are the vectors $W=\left(W_{\boldsymbol{\beta}}(y)\right)$ and $\bar{W}=\left(\bar{W}_{\boldsymbol{\beta}}(y)\right)$, indexed by the pairs $(y,\boldsymbol{\beta})$ of a unit site $y\in T^0_k$ and a spinor index $\boldsymbol{\beta}$, and for them~\eqref{eq:renormalisation-group-42} gives $\left\langle W\,W^T\right\rangle=0$, $\left\langle\bar{W}\,\bar{W}^T\right\rangle=0$ and $\left\langle W\,\bar{W}^T\right\rangle=\boldsymbol{\varsigma}\cdot\mathbf{1}$, where $\boldsymbol{\varsigma}$ is the sign of~\eqref{eq:renormalisation-group-42}. Hence
\[
\left\langle\boldsymbol{\phi}_{+}\,\boldsymbol{\phi}_{+}^T\right\rangle=0,\qquad\left\langle\bar{\boldsymbol{\phi}}_{+}\,\bar{\boldsymbol{\phi}}_{+}^T\right\rangle=0,\qquad\left\langle\boldsymbol{\phi}_{+}\,\bar{\boldsymbol{\phi}}_{+}^T\right\rangle=\boldsymbol{\varsigma}\cdot\mathring{G}(\mathscr{T}),
\]
where, for $B\in\{\mathscr{T},-\mathscr{T}\}$, $\mathring{G}(B)$ is the operator on the spinor functions on $T_k$ given by
\[
\mathring{G}(B):=H_k(1_U,B)\cdot\Gamma_k(1_U,B)\cdot\bar{H}_k(1_U,B)^T,
\]
the transpose being taken in the site and the spinor indices together.

The definition of $\bar{H}_k$ in~\eqref{eq:renormalisation-group-56}, property~\ref{it:renormalisation-group-15} (for $k=0$ property~\ref{it:renormalisation-group-17}, which contains property~\ref{it:renormalisation-group-15} for $\Gamma_0$), and the relations $\mathsf{C}^T=\mathsf{C}^{-1}$, $\left(\mathsf{C}^{-1}\right)^T=\mathsf{C}$ give
\begin{gather*}
\bar{H}_k(1_U,-\mathscr{T})=\mathsf{C}^{-1}\cdot H_k(1_U,\mathscr{T})\cdot\mathsf{C},\qquad\Gamma_k(1_U,-\mathscr{T})^T=\mathsf{C}^{-1}\cdot\Gamma_k(1_U,\mathscr{T})\cdot\mathsf{C},\\
H_k(1_U,-\mathscr{T})^T=\left(\mathsf{C}\cdot\bar{H}_k(1_U,\mathscr{T})\cdot\mathsf{C}^{-1}\right)^T=\mathsf{C}\cdot\bar{H}_k(1_U,\mathscr{T})^T\cdot\mathsf{C}^{-1}.
\end{gather*}
With $\mathsf{C}^2=\mathsf{C}^{-2}=-1$, and with the operators $H_k$, $\Gamma_k$, $\bar{H}_k$ in the middle term taken at $(1_U,\mathscr{T})$, it follows that
\[
\mathring{G}(-\mathscr{T})^T=\bar{H}_k(1_U,-\mathscr{T})\cdot\Gamma_k(1_U,-\mathscr{T})^T\cdot H_k(1_U,-\mathscr{T})^T=\mathsf{C}^{-1}\cdot H_k\cdot\Gamma_k\cdot\mathsf{C}^2\cdot\bar{H}_k^T\cdot\mathsf{C}^{-1}=-\mathsf{C}^{-1}\cdot\mathring{G}(\mathscr{T})\cdot\mathsf{C}^{-1},
\]
that is,
\[
\mathsf{C}\cdot\mathring{G}(-\mathscr{T})^T\cdot\mathsf{C}^{-1}=\mathring{G}(\mathscr{T}).
\]

Now let the transformation act on $\Xi_U$, that is, evaluate $\Xi_U$ at $(-A,-\mathscr{C})$ and at $\mathbf{C}\,\psi^{\#}$. The measure $\hat{\mu}$ is invariant under $\mathsf{Z}\mapsto-\mathsf{Z}$ (Proposition~\ref{prop:symmetries-reality}). After the substitution of $-\mathsf{Z}$ for the variable of integration, the integrand is taken at the arguments $(-A,-\mathscr{C},-\mathsf{Z})$ and $\mathbf{C}\,\psi^{\#}$, at which the total connection is $-\mathscr{T}$. By Step~2, applied at these arguments, it equals
\[
\exp\left\{\sum_{X\in D_k,\,X\subseteq U}\mathbf{F}\left(X;-A,-\mathscr{C},-\mathsf{Z};\mathbf{C}\,\psi^{\#},W^{\#}_k(1_U,-\mathscr{T})\right)\right\}.
\]
Put
\[
\boldsymbol{\phi}':=\mathsf{C}\cdot\bar{\boldsymbol{\phi}}_{-},\qquad\bar{\boldsymbol{\phi}}':=\mathsf{C}^{-1}\cdot\boldsymbol{\phi}_{-}.
\]
Then $\boldsymbol{\phi}_{-}=\mathsf{C}\cdot\bar{\boldsymbol{\phi}}'$ and $\bar{\boldsymbol{\phi}}_{-}=\mathsf{C}^{-1}\cdot\boldsymbol{\phi}'=-\mathsf{C}\cdot\boldsymbol{\phi}'$. As $\mathsf{C}$ commutes with the H\"older derivatives, also
\[
\delta_{\alpha,-\mathscr{T}}\boldsymbol{\phi}_{-}=\mathsf{C}\cdot\delta_{\alpha,-\mathscr{T}}\bar{\boldsymbol{\phi}}',\qquad\delta_{\alpha,\mathscr{T}}\bar{\boldsymbol{\phi}}_{-}=-\mathsf{C}\cdot\delta_{\alpha,\mathscr{T}}\boldsymbol{\phi}'.
\]
Let ${\boldsymbol{\phi}'}^{\#}$ be the family with the components $\boldsymbol{\phi}'$, $\bar{\boldsymbol{\phi}}'$, $\delta_{\alpha,\mathscr{T}}\boldsymbol{\phi}'$, $\delta_{\alpha,-\mathscr{T}}\bar{\boldsymbol{\phi}}'$ in the places of $\psi$, $\bar{\psi}$, $\chi$, $\bar{\chi}$. The four identities just stated say that
\[
W^{\#}_k(1_U,-\mathscr{T})=\mathbf{C}\,{\boldsymbol{\phi}'}^{\#}.
\]
The homomorphism that leaves $\psi$, $\chi$ fixed and replaces $\boldsymbol{\phi}^{\#}$ by ${\boldsymbol{\phi}'}^{\#}$ maps $\mathbf{C}\,\boldsymbol{\phi}^{\#}$ to $\mathbf{C}\,{\boldsymbol{\phi}'}^{\#}$. Applied to both sides of the identity of Step~2 that expresses the invariance of $\mathbf{F}$, it shows that the integrand at $\mathsf{Z}$ equals
\[
\exp\left\{\sum_{X\in D_k,\,X\subseteq U}\mathbf{F}\left(X;A,\mathscr{C},\mathsf{Z};\psi^{\#},{\boldsymbol{\phi}'}^{\#}\right)\right\}.
\]
The components of ${\boldsymbol{\phi}'}^{\#}$ are linear combinations of the generators $W$, and they arise from $\boldsymbol{\phi}'$, $\bar{\boldsymbol{\phi}}'$ by the same linear maps by which the components of $W^{\#}_k(1_U,\mathscr{T})$ arise from $\boldsymbol{\phi}_{+}$, $\bar{\boldsymbol{\phi}}_{+}$. The generators anticommute; so $\left\langle\bar{W}\,W^T\right\rangle=-\boldsymbol{\varsigma}\cdot\mathbf{1}$, that is,
\[
\int\mathrm{d}\mu_I(W)\,\bar{W}_{\boldsymbol{\beta}}(y)\cdot W_{\boldsymbol{\beta}'}(y')=-\boldsymbol{\varsigma}\cdot\delta_{yy'}\cdot\delta_{\boldsymbol{\beta}\boldsymbol{\beta}'},
\]
for the pairs $(y,\boldsymbol{\beta})$, $(y',\boldsymbol{\beta}')$ of~\eqref{eq:renormalisation-group-42}. Hence
\[
\left\langle\boldsymbol{\phi}'\,{\boldsymbol{\phi}'}^T\right\rangle=\mathsf{C}\cdot\left\langle\bar{\boldsymbol{\phi}}_{-}\,\bar{\boldsymbol{\phi}}_{-}^T\right\rangle\cdot\mathsf{C}^T=0,\qquad\left\langle\bar{\boldsymbol{\phi}}'\,{\bar{\boldsymbol{\phi}}'}^T\right\rangle=\mathsf{C}^{-1}\cdot\left\langle\boldsymbol{\phi}_{-}\,\boldsymbol{\phi}_{-}^T\right\rangle\cdot\left(\mathsf{C}^{-1}\right)^T=0,
\]
and
\begin{align*}
\left\langle\boldsymbol{\phi}'\,{\bar{\boldsymbol{\phi}}'}^T\right\rangle&=-\boldsymbol{\varsigma}\cdot\mathsf{C}\cdot\bar{H}_k(1_U,-\mathscr{T})\cdot\left(\mathsf{C}^{-1}\cdot H_k(1_U,-\mathscr{T})\cdot\Gamma_k(1_U,-\mathscr{T})\right)^T=-\boldsymbol{\varsigma}\cdot\mathsf{C}\cdot\mathring{G}(-\mathscr{T})^T\cdot\left(\mathsf{C}^{-1}\right)^T\\
&=-\boldsymbol{\varsigma}\cdot\mathsf{C}\cdot\mathring{G}(-\mathscr{T})^T\cdot\mathsf{C}=\boldsymbol{\varsigma}\cdot\mathsf{C}\cdot\mathring{G}(-\mathscr{T})^T\cdot\mathsf{C}^{-1}=\boldsymbol{\varsigma}\cdot\mathring{G}(\mathscr{T}).
\end{align*}
So $\boldsymbol{\phi}'$, $\bar{\boldsymbol{\phi}}'$ have the same two-point functions as $\boldsymbol{\phi}_{+}$, $\bar{\boldsymbol{\phi}}_{+}$, and every pair expectation of two components of ${\boldsymbol{\phi}'}^{\#}$ equals the pair expectation of the corresponding two components of $W^{\#}_k(1_U,\mathscr{T})$. Lemma~\ref{lem:substitutions-integral}(a) gives, for every $\mathsf{Z}\in\mathrm{Box}_k$,
\[
\int\mathrm{d}\mu_I(W)\,\exp\left\{\sum_X\mathbf{F}\left(X;A,\mathscr{C},\mathsf{Z};\psi^{\#},{\boldsymbol{\phi}'}^{\#}\right)\right\}=\int\mathrm{d}\mu_I(W)\,\exp\left\{\sum_X\mathbf{F}\left(X;A,\mathscr{C},\mathsf{Z};\psi^{\#},W^{\#}_k(1_U,\mathscr{T})\right)\right\},
\]
both sums extending over the polymers $X\subseteq U$. By Step~2 the right side is the integral over $W$ of the integrand of $\Xi_U$ at $(A,\mathscr{C},\mathsf{Z})$ and $\psi^{\#}$. Integration with respect to $\mathrm{d}\hat{\mu}(\mathsf{Z})$ yields
\[
\Xi_U\left(-A,-\mathscr{C};\mathbf{C}\,\psi^{\#}\right)=\Xi_U\left(A,\mathscr{C};\psi^{\#}\right).
\]

\emph{Step~4.} The elements $E^{(\le5),x}(\mathbf{Y})$ are even, by Proposition~\ref{prop:bounds-regularity-localised}(e), and therefore commute with each other. Let $U$ be a union of $M$-cubes. As in Step~1 of the proof of Lemma~\ref{lem:integration-polymer-exponential}, the expansion of the product and the collection of its terms according to the union $X$ of the polymers that contribute to them give
\[
\exp\left\{\sum_{\mathbf{Y}\in D_k,\,\mathbf{Y}\subseteq U}E^{(\le5),x}(\mathbf{Y})\right\}=\prod_{\mathbf{Y}\in D_k,\,\mathbf{Y}\subseteq U}\left(1+\left(\exp\left(E^{(\le5),x}(\mathbf{Y})\right)-1\right)\right)=\sum_{X\subseteq U}\mathsf{K}(X),
\]
where $X$ ranges over the unions of $M$-cubes contained in $U$, the empty union included. Integration over $(W,\mathsf{Z})$ gives
\[
\Xi_U=\sum_{X\subseteq U}\mathsf{H}(X)
\]
for every union $U$ of $M$-cubes. By M\"obius inversion on the subsets of the set of the $M$-cubes it follows that, for every union $X$ of $M$-cubes,
\[
\mathsf{H}(X)=\sum_{U\subseteq X}(-1)^{|X\setminus U|_M}\cdot\Xi_U,
\]
the sum extending over the unions $U$ of $M$-cubes contained in $X$. By Step~3 and~\eqref{it:operations-preserve-invariance-Pr1}, $\mathsf{H}(X)$ is invariant under~\eqref{it:lattice-symmetries-gauge-TC}. The series~\eqref{eq:renormalisation-group-61} converge absolutely in norm, by assertion~(1) of Section~\ref{sec:activity-before-reblocking}; by~\eqref{it:operations-preserve-invariance-Pr1} and~\eqref{it:operations-preserve-invariance-Pr4}, applied to~\eqref{eq:renormalisation-group-61}, the element $E^{\#,x}_k(X)$ is invariant under~\eqref{it:lattice-symmetries-gauge-TC} for every $X\in D_k$.
\end{proof}

\begin{proposition}[Symmetries of the step map]\label{prop:symmetries-step-map}
Let $(m,\mathscr{E})\in\mathrm{Ball}^x_k(4)$. At level $k+1$ the family $E^{*,x}_k(m,\mathscr{E})$ has clause~\eqref{it:renormalisation-group-1-k3} of Definition~\ref{def:renormalisation-group-1} for the group $\mathbb{G}_{k+1}$, and the clauses~\eqref{it:renormalisation-group-1-k4} and~\eqref{it:renormalisation-group-1-k5}. Every monomial that occurs in $E^{*,x}_k(m,\mathscr{E})$ has as many barred legs as unbarred legs; in particular $E^{*,x}_k(m,\mathscr{E})$ is even. If $m$ is real and $\mathscr{E}\in K^x_k$, then $E^{*,x}_k(m,\mathscr{E})$ has clause~\eqref{it:renormalisation-group-1-k7}.
\end{proposition}

\begin{proof}
The generators are independent; an identity between two elements is therefore the same as the identity between their kernels, and it suffices to prove identities between elements. In this proof $(m,\mathscr{E})$ is fixed and omitted from the notation. Let $Y\in D_{k+1}$ and $(\mathscr{A},\mathscr{C})\in\tilde{R}_{k+1}\times\mathfrak{X}_{k+1}$. By Lemma~\ref{lem:scaling-domains}, $\mathscr{A}_L\in\tilde{R}^{-}_k$ and $\mathscr{C}_L\in\mathfrak{X}_k$, and by the definition in Section~\ref{sec:fluctuation-integral-new}
\[
E^{*,x}_k(Y;\mathscr{A},\mathscr{C},\psi,\chi)=\sum_{X\in D_k:\,\bar{X}^L=LY}E^{\#,x}_k\left(X;\mathscr{A}_L,\mathscr{C}_L,\psi_L,\chi_L\right),
\]
with the scalings of Section~\ref{sec:local-parts-adjustment}.

\emph{Clause~\eqref{it:renormalisation-group-1-k3}.} Let $g'\in\mathbb{G}_{k+1}$, and let $g\in\mathbb{G}^L_k$ be the element that the isomorphism $g\mapsto g'$ of Proposition~\ref{prop:symmetries-extraction}(d) sends to $g'$. By Proposition~\ref{prop:symmetries-extraction}(d),
\[
\left(g'\mathscr{A}\right)_L=g\left(\mathscr{A}_L\right),\qquad\left(g'\mathscr{C}\right)_L=g\left(\mathscr{C}_L\right),\qquad\left(g'\psi\right)_L=g\left(\psi_L\right),\qquad\left(g'\chi\right)_L=g\left(\chi_L\right),\qquad L\left(g'Y\right)=g(LY).
\]
By Proposition~\ref{prop:symmetries-extraction}(a), $g$ maps polymers onto polymers, and $g(\bar{X}^L)$ is the union of the $LM$-cubes that contain an $M$-cube of $gX$. Hence $X\mapsto gX$ is a bijection of $\left\{X\in D_k:\bar{X}^L=LY\right\}$ onto $\left\{X\in D_k:\bar{X}^L=L\left(g'Y\right)\right\}$, and
\begin{align*}
E^{*,x}_k\left(g'Y;g'\mathscr{A},g'\mathscr{C},g'\psi,g'\chi\right)&=\sum_{X\in D_k:\,\bar{X}^L=LY}E^{\#,x}_k\left(gX;g\left(\mathscr{A}_L\right),g\left(\mathscr{C}_L\right),g\left(\psi_L\right),g\left(\chi_L\right)\right)\\
&=\sum_{X\in D_k:\,\bar{X}^L=LY}E^{\#,x}_k\left(X;\mathscr{A}_L,\mathscr{C}_L,\psi_L,\chi_L\right)=E^{*,x}_k(Y;\mathscr{A},\mathscr{C},\psi,\chi),
\end{align*}
the second equality by Lemma~\ref{lem:renormalisation-group-2}(a).

\emph{Clause~\eqref{it:renormalisation-group-1-k4}.} Let $(\mathscr{A},\mathscr{C})$ and $(\mathscr{A}^1,\mathscr{C}^1)$ lie in $\tilde{R}_{k+1}\times\mathfrak{X}_{k+1}$, let $\boldsymbol{\lambda}$ and $\boldsymbol{\lambda}'$ be real functions on the sites of $T_{k+1}$, and let $\mathscr{A}^1=\mathscr{A}-d\boldsymbol{\lambda}$ and $\mathscr{C}^1=\mathscr{C}-d\boldsymbol{\lambda}'$ on the bonds at distance at most $1$ from $Y$. For the sites $x$ of $T_k$ put $\boldsymbol{\lambda}_L(x):=L^{\frac{1}{2}}\cdot\boldsymbol{\lambda}(x/L)$, and define $\boldsymbol{\lambda}'_L$ in the same way. The relation $(d\boldsymbol{\lambda})_L=d\boldsymbol{\lambda}_L$, shown in the proof of Lemma~\ref{lem:scaling-domains}, and its analogue for $\boldsymbol{\lambda}'$ give
\[
\mathscr{A}^1_L=\mathscr{A}_L-d\boldsymbol{\lambda}_L,\qquad\mathscr{C}^1_L=\mathscr{C}_L-d\boldsymbol{\lambda}'_L
\]
on every bond $b$ of $T_k$ whose two end points have sup-distance at most $L$ from $LY$, because for such $b$ the bond $b/L$ of $T_{k+1}$ is at distance at most $1$ from $Y$. Every $X\in D_k$ with $\bar{X}^L=LY$ is contained in $LY$, by assertion~(5) of Section~\ref{sec:activity-before-reblocking}; so these relations hold on the bonds at distance at most $1$ from every such $X$. By Lemma~\ref{lem:scaling-domains}, $\mathscr{A}_L$ and $\mathscr{A}^1_L$ lie in $\tilde{R}^{-}_k$, and $\mathscr{C}_L$, $\mathscr{C}^1_L$ in $\mathfrak{X}_k$. Since $e_{k+1}=L^{\frac{1}{2}}\cdot e_k$ (Section~\ref{sec:hypotheses-conditions}), one has $e_{k+1}\cdot\boldsymbol{\lambda}(x/L)=e_k\cdot\boldsymbol{\lambda}_L(x)$, and the same for $\boldsymbol{\lambda}'$; thus the phases of clause~\eqref{it:renormalisation-group-1-k4} at level $k+1$ become, after the scaling, the phases of~\eqref{it:lattice-symmetries-gauge-TG} at level $k$:
\[
\left(\psi^{\boldsymbol{\lambda}+\boldsymbol{\lambda}'}\right)_L=\left(\psi_L\right)^{\boldsymbol{\lambda}_L+\boldsymbol{\lambda}'_L},\qquad\left(\chi^{\boldsymbol{\lambda}+\boldsymbol{\lambda}'}\right)_L=\left(\chi_L\right)^{\boldsymbol{\lambda}_L+\boldsymbol{\lambda}'_L}.
\]
Lemma~\ref{lem:renormalisation-group-2}(b), applied to $E^{\#,x}_k(X)$ for every $X\in D_k$ with $\bar{X}^L=LY$ and to the gauge functions $\boldsymbol{\lambda}_L$, $\boldsymbol{\lambda}'_L$, and summation over these $X$ give
\[
E^{*,x}_k\left(Y;\mathscr{A}^1,\mathscr{C}^1,\psi^{\boldsymbol{\lambda}+\boldsymbol{\lambda}'},\chi^{\boldsymbol{\lambda}+\boldsymbol{\lambda}'}\right)=E^{*,x}_k(Y;\mathscr{A},\mathscr{C},\psi,\chi).
\]

\emph{Clauses~\eqref{it:renormalisation-group-1-k5} and~\eqref{it:renormalisation-group-1-k7}.} The scaling of bond functions is real linear: $(-\mathscr{A})_L=-\mathscr{A}_L$ and $\left(\overline{\mathscr{A}}\right)_L=\overline{\mathscr{A}_L}$, and the same holds for $\mathscr{C}$. The scalings $\psi\mapsto\psi_L$ and $\chi\mapsto\chi_L$ multiply the generators by the real numbers $L^{-1}$ and $L^{-1-\alpha}$ and change their sites only; hence they commute with the multiplication by the constant real matrices $\mathsf{C}$ and $\varkappa$ (Section~\ref{sec:recalled-auxiliary-results}) and with the complex conjugation of the coefficients. Consequently
\[
E^{*,x}_k\left(Y;-\mathscr{A},-\mathscr{C},\mathbf{C}\,\psi,\mathbf{C}\,\chi\right)=\sum_{X\in D_k:\,\bar{X}^L=LY}E^{\#,x}_k\left(X;-\mathscr{A}_L,-\mathscr{C}_L,\mathbf{C}\,\psi_L,\mathbf{C}\,\chi_L\right),
\]
and by Lemma~\ref{lem:renormalisation-group-2}(c) the right side equals $E^{*,x}_k(Y;\mathscr{A},\mathscr{C},\psi,\chi)$; this is clause~\eqref{it:renormalisation-group-1-k5}. In the same way, the image under $\mathbf{K}$ of $E^{*,x}_k\left(Y;-\overline{\mathscr{A}},-\overline{\mathscr{C}}\right)$ is the sum, over the same polymers $X$, of the images under $\mathbf{K}$ of $E^{\#,x}_k\left(X;-\overline{\mathscr{A}_L},-\overline{\mathscr{C}_L}\right)$, taken with $\psi_L$, $\chi_L$ in the places of $\psi$, $\chi$. If $m$ is real and $\mathscr{E}\in K^x_k$, Lemma~\ref{lem:renormalisation-group-2}(d) identifies each of these images with $E^{\#,x}_k\left(X;\mathscr{A}_L,\mathscr{C}_L,\psi_L,\chi_L\right)$, and clause~\eqref{it:renormalisation-group-1-k7} follows.

\emph{Legs.} Apply clause~\eqref{it:renormalisation-group-1-k4}, just proved, with $\mathscr{A}^1=\mathscr{A}$, $\mathscr{C}^1=\mathscr{C}$, a real constant $\boldsymbol{\lambda}$ and $\boldsymbol{\lambda}'=0$; its hypothesis holds because $d\boldsymbol{\lambda}=0$. At level $k+1$ the substitution $\psi\mapsto\psi^{\boldsymbol{\lambda}}$, $\chi\mapsto\chi^{\boldsymbol{\lambda}}$ multiplies every unbarred leg by $\exp\left(i\,e_{k+1}\cdot\boldsymbol{\lambda}\right)$ and every barred leg by $\exp\left(-i\,e_{k+1}\cdot\boldsymbol{\lambda}\right)$; so it multiplies a monomial with $n_u$ unbarred and $n_b$ barred legs by $\exp\left(i\,e_{k+1}\cdot\boldsymbol{\lambda}\cdot(n_u-n_b)\right)$. The coefficient $c$ of such a monomial in $E^{*,x}_k(Y;\mathscr{A},\mathscr{C},\psi,\chi)$ therefore satisfies
\[
c\cdot\left(\exp\left(i\,e_{k+1}\cdot\boldsymbol{\lambda}\cdot(n_u-n_b)\right)-1\right)=0
\]
for every real $\boldsymbol{\lambda}$. If $n_u\ne n_b$, the choice $\boldsymbol{\lambda}=\pi/\left(e_{k+1}\cdot(n_u-n_b)\right)$, admissible since $e_{k+1}>0$, gives $c=0$. Hence every monomial that occurs has as many barred as unbarred legs, and its degree $n_u+n_b=2n_u$ is even.

Clause~\eqref{it:step-small-fields-f} of Theorem~\ref{thm:step-small-fields} follows.
\end{proof}

\subsection{Proof of clause~\eqref{it:step-small-fields-c} of Theorem~\ref{thm:step-small-fields}}\label{sec:proof-clause-theorem}

The proof rests on the following one-step formula, which is Proposition~\ref{prop:step-external-connection} (see Theorem~\ref{thm:one-step-formula}). Objects of the level $k+1$ are formed with the side $l/L$, the mass $L\cdot\mathfrak{m}$ and the charge $e_{k+1}$ of Section~\ref{sec:hypotheses-conditions}.

Let $\boldsymbol{\varepsilon}$ and $m$ be real numbers with $|m|\le w^{\mathrm{m}}_k$, let $\mathscr{E}\in K^x_k$ with $\|\mathscr{E}\|^x_k\le w^{\mathrm{E}}_k$, where $\mathscr{E}=0$ if $k=0$, and let $A'\in\mathrm{Dom}_{k+1}$ and $\mathscr{C}'\in\mathfrak{X}_{k+1}$. With the numbers $\varepsilon(\mathscr{E})$, $m(\mathscr{E})$ and the map $R$ of Section~\ref{sec:local-parts-adjustment}, and with the number $d\varepsilon^0_k$ of Section~\ref{sec:averaging-fluctuations-gauge}, define
\[
m':=m+m(\mathscr{E}),\qquad\mathscr{E}^r:=R\,\mathscr{E},\qquad m^{+}:=L\cdot m',\qquad\boldsymbol{\varepsilon}^{+}:=L^3\cdot\left(\boldsymbol{\varepsilon}+\varepsilon(\mathscr{E})+d\varepsilon^0_k\right).
\]

Let $H_{k+1}$ and $H_k$, written without a connection, be the minimisers of Section~\ref{sec:averaging-fluctuations-gauge} at the levels $k+1$ and $k$; let $H^{\#}_{k+1}$ be the operator of Section~\ref{sec:fermion-integrals-blocking} at the level $k+1$; and let $\Psi'$ and $\bar{\Psi}'$ be the unbarred and the barred generators of Section~\ref{sec:grassmann-algebras-norms} at the level $k+1$. Put
\begin{gather*}
\mathscr{A}':=H_{k+1}A',\qquad\mathscr{T}':=\mathscr{A}'+\mathscr{C}',\qquad\psi'^{\#}:=H^{\#}_{k+1}(\mathscr{T}')\,\Psi',\\
A:=\mathscr{A}'_L,\qquad\mathscr{C}:=\mathscr{C}'_L,\qquad\mathscr{T}:=A+\mathscr{C},\qquad\psi^{0,\#}:=\left(\psi'^{\#}\right)_L,
\end{gather*}
the scalings $(\cdot)_L$ being those of Section~\ref{sec:local-parts-adjustment}, applied to every component of $\psi'^{\#}$. The unbarred and the barred spinor fields contained in $\psi'^{\#}$ are
\[
\psi'=H_{k+1}(\mathscr{T}')\,\Psi',\qquad\bar{\psi}'=\bar{H}_{k+1}(\mathscr{T}')\,\bar{\Psi}',
\]
and those contained in $\psi^{0,\#}$ are $\psi^0=\psi'_L$ and $\bar{\psi}^0=\bar{\psi}'_L$.

Let $W$ and $\bar{W}$ be the unbarred and the barred generators $W(u)$ of Section~\ref{sec:grassmann-algebras-norms}, which are indexed by the unit sites of $T^0_k$ and the spinor indices, and let $\mathsf{Z}\in\mathrm{Box}_k$. With the operators $T_k(\cdot)$, $\bar{T}_k(\cdot)$ and $\Gamma_k(\cdot)$ of Section~\ref{sec:exact-fermion-operators}, where for $k=0$ the operator $\bar{T}_0(\cdot)$ is the one defined in Section~\ref{sec:block-operators}, namely the identity, and with the operators $C$ and $C_k^{\frac{1}{2}}$ of Section~\ref{sec:averaging-fluctuations-gauge}, put
\begin{gather*}
\Psi^c:=T_k(\mathscr{T})\,\psi^0,\qquad\bar{\Psi}^c:=\bar{T}_k(\mathscr{T})\,\bar{\psi}^0,\qquad\mathscr{Z}:=H_k\cdot C\cdot C_k^{\frac{1}{2}}\cdot\mathsf{Z},\\
\Psi^f:=\Psi^c+\Gamma_k(\mathscr{T})\,W,\qquad\bar{\Psi}^f:=\bar{\Psi}^c+\bar{W},
\end{gather*}
and
\begin{multline*}
V^{\mathrm{ex},x}_k:=\mathrm{Act}^x_k[m',\mathscr{E}^r](A+\mathscr{Z},\mathscr{C};\Psi^f)-\mathrm{Act}^x_k[m',\mathscr{E}^r](A,\mathscr{C};\Psi^f)\\
+I^x_k[m',\mathscr{E}^r](A,\mathscr{C};\Psi^f)-I^x_k[m',\mathscr{E}^r](A,\mathscr{C};\Psi^c),
\end{multline*}
\[
\Xi^x_k:=\int\mathrm{d}\mu_I(W)\,\mathrm{d}\hat{\mu}(\mathsf{Z})\,\exp\left(V^{\mathrm{ex},x}_k\right).
\]
Here $\mathrm{Act}^x_k$ and $I^x_k$ are the functionals of Section~\ref{sec:fermion-integrals-blocking}, and the last argument $\Psi^f$, resp.\ $\Psi^c$, indicates that the generators $\Psi$ and $\bar{\Psi}$ of the level $k$ are replaced by $\Psi^f$ and $\bar{\Psi}^f$, resp.\ by $\Psi^c$ and $\bar{\Psi}^c$. The integral $\int\mathrm{d}\mu_I(W)$ is the Grassmann Gaussian integral over the generators $W$, $\bar{W}$ of Section~\ref{sec:recalled-auxiliary-results}, and $\hat{\mu}$ is the probability measure on $V_k$ of Section~\ref{sec:averaging-fluctuations-gauge}, which is concentrated on $\mathrm{Box}_k$. The density $\mathscr{F}^x_k$ and the block transformation $\mathbf{T}^x_k$ are those of Section~\ref{sec:fermion-integrals-blocking}. Let $\mathfrak{P}_k=\left\{\mathfrak{P}_k(X,\cdot\,):X\in D_k\right\}$ be an arbitrary family of functions from $\mathfrak{X}_k$ to $\mathbb{C}$; the density $\mathscr{F}^x_k$ on the left side of~\eqref{eq:renormalisation-group-65} and the last sum in the exponent on its right side are formed with this family. Along a flow, the family of level $k$ of the trajectory with slot of Section~\ref{sec:flow} is such a family. The operator $D_{k+1}(\cdot)$ and the numbers $\delta Z_k(\cdot)$ are those of Section~\ref{sec:exact-fermion-operators}, and the pairings are those of Section~\ref{sec:notation}, formed on the lattices $T_{k+1}$ and $T^0_{k+1}$. Then
\begin{multline}\label{eq:renormalisation-group-65}
\left(\mathbf{T}^x_k\mathscr{F}^x_k[\boldsymbol{\varepsilon},m,\mathscr{E}]\right)(A',\mathscr{C}';\Psi')\\
=\left(\det D_{k+1}(0)\right)^{-1}\cdot\exp\left\{-\left\langle\bar{\Psi}',D_{k+1}(\mathscr{T}')\Psi'\right\rangle-m^{+}\cdot\left\langle\bar{\psi}',\psi'\right\rangle-\boldsymbol{\varepsilon}^{+}\cdot n_{k+1}\right.\\
\left.+\sum_{X\in D_k}\mathscr{E}^r\left(X;A,\mathscr{C},\psi^{0,\#}\right)+\sum_{X\in D_k}\mathfrak{P}_k(X,\mathscr{C})\right\}\cdot\frac{\delta Z_k(\mathscr{T})}{\delta Z_k(0)}\cdot\Xi^x_k.
\end{multline}

\begin{remark}
The family $\mathfrak{P}_k$ plays no role in~\eqref{eq:renormalisation-group-65}: its replacement by another family of functions from $\mathfrak{X}_k$ to $\mathbb{C}$ multiplies both sides by one and the same nonzero complex number. For each family the identity~\eqref{eq:renormalisation-group-65} is contained in Theorem~\ref{thm:one-step-formula}, which is stated in the notation of Section~\ref{sec:densities-block-transformation}, where the family is arbitrary; and the identity for one family, for instance for the family $\mathfrak{P}_k=0$, which is defined in the setting of Section~\ref{sec:hypotheses-conditions}, is equivalent to the identity for every family. Indeed, any two families of functions from $\mathfrak{X}_k$ to $\mathbb{C}$ differ by a family $\varphi=\left\{\varphi(X,\cdot\,):X\in D_k\right\}$ of such functions, so that the second is $\mathfrak{P}_k+\varphi$ if the first is $\mathfrak{P}_k$. By the definition in Section~\ref{sec:fermion-integrals-blocking}, the family enters $\mathscr{F}^x_k[\boldsymbol{\varepsilon},m,\mathscr{E}]$ only through the factor $\exp\left\{\sum_{X\in D_k}\mathfrak{P}_k(X,\mathscr{C})\right\}$. Hence the replacement of $\mathfrak{P}_k$ by $\mathfrak{P}_k+\varphi$ multiplies $\mathscr{F}^x_k[\boldsymbol{\varepsilon},m,\mathscr{E}]$ by the complex number
\[
c_{\varphi}(\mathscr{C}):=\exp\left\{\sum_{X\in D_k}\varphi(X,\mathscr{C})\right\},
\]
which depends on the second slot $\mathscr{C}$ only, and neither on the block field in the first argument nor on the generators $\Psi$, $\bar{\Psi}$. The map $F\mapsto\mathbf{T}^x_kF$ is linear, and it evaluates $F$ at the second slot $\mathscr{C}=\mathscr{C}'_L$; so the left side of~\eqref{eq:renormalisation-group-65} is multiplied by $c_{\varphi}(\mathscr{C}'_L)$. The right side is multiplied by the same number, because the family enters it only through the summand $\sum_{X\in D_k}\mathfrak{P}_k(X,\mathscr{C})$ of the exponent, which is a complex number, taken at $\mathscr{C}=\mathscr{C}'_L$. Since $c_{\varphi}(\mathscr{C}'_L)\ne0$, it follows that~\eqref{eq:renormalisation-group-65} holds for $\mathfrak{P}_k+\varphi$ if and only if it holds for $\mathfrak{P}_k$.
\end{remark}

\begin{proof}[Proof of clause~\eqref{it:step-small-fields-c} of Theorem~\ref{thm:step-small-fields}]
Let $\boldsymbol{\varepsilon}$, $m$, $\mathscr{E}$, $\mathfrak{P}_k$, $A'$ and $\mathscr{C}'$ be as in that clause. The numbers $m'$, $m^{+}$, $\boldsymbol{\varepsilon}^{+}$, the family $\mathscr{E}^r$, the connections $\mathscr{A}'$, $\mathscr{T}'$, $A$, $\mathscr{C}$, $\mathscr{T}$, the fields $\psi'^{\#}$, $\psi^{0,\#}$, $\Psi^c$, $\bar{\Psi}^c$, $\Psi^f$, $\bar{\Psi}^f$, the bond function $\mathscr{Z}$, the element $V^{\mathrm{ex},x}_k$ and the integral $\Xi^x_k$ are those of Proposition~\ref{prop:step-external-connection}, as recalled in Section~\ref{sec:proof-clause-theorem}; the numbers $m^{+}$ and $\boldsymbol{\varepsilon}^{+}$ coincide with those of clause~\eqref{it:step-small-fields-c}. On the unbarred fields the two components of $\psi'^{\#}=H^{\#}_{k+1}(\mathscr{T}')\,\Psi'$ are
\[
\psi':=H_{k+1}(\mathscr{T}')\,\Psi',\qquad\chi':=\delta_{\alpha,\mathscr{T}'}\psi',
\]
and on the barred fields they are $\bar{\psi}':=\bar{H}_{k+1}(\mathscr{T}')\,\bar{\Psi}'$ and $\bar{\chi}':=\delta_{\alpha,-\mathscr{T}'}\bar{\psi}'$; here $\delta_{\alpha,\pm\mathscr{T}'}$ is the H\"older derivative of Definition~\ref{def:pairs-sites-holder} at the level $k+1$, formed with $e_{k+1}$ and $\eta_{k+1}$. Accordingly $\psi^{0,\#}=\left(\psi'^{\#}\right)_L$ has the components
\[
\psi^0=\psi'_L,\qquad\chi^0=\chi'_L,\qquad\bar{\psi}^0=\bar{\psi}'_L,\qquad\bar{\chi}^0=\bar{\chi}'_L,
\]
the scalings being those of Section~\ref{sec:local-parts-adjustment}.

\emph{Step~0: domains.} By the definition of $\mathrm{Dom}_{k+1}$ (Definition~\ref{def:domains}), the connection $\mathscr{A}'=H_{k+1}A'$ belongs to $\frac{1}{4}\tilde{R}_{k+1}$, which is a subset of $\tilde{R}_{k+1}$. Lemma~\ref{lem:scaling-domains}(a), with $c=1$, together with Lemma~\ref{lem:scaling-domains}(b), gives $A=\mathscr{A}'_L\in\tilde{R}^{-}_k$, and Lemma~\ref{lem:scaling-domains}(c) gives $\mathscr{C}=\mathscr{C}'_L\in\mathfrak{X}_k$. Lemma~\ref{lem:inclusions-domains}(a) holds at the level $k+1$ as well (Section~\ref{sec:hypotheses-conditions}); applied at the level $k+1$ it shows that $\mathscr{T}'=\mathscr{A}'+\mathscr{C}'$ lies in $\tilde{G}_{k+1}$, and applied at the level $k$, together with the inclusion $\tilde{R}^{-}_k\subset\tilde{R}_k$, it shows that $\mathscr{T}=A+\mathscr{C}$ lies in $\tilde{G}_k$. Since $|m|\le w^{\mathrm{m}}_k$ and $\|\mathscr{E}\|^x_k\le w^{\mathrm{E}}_k$, the pair $(m,\mathscr{E})$ belongs to $\mathrm{Ball}^x_k(1)$, a subset of $\mathrm{Ball}^x_k(4)$ (Section~\ref{sec:spaces-activities}). The left side of~\eqref{eq:renormalisation-group-54} is defined, because the field $A^{\min}_k(A')+\mathrm{Sq}_k\zeta$ lies in $\mathrm{Dom}_k$ for every $\zeta\in\mathrm{Box}_k$, by Proposition~\ref{prop:fields-domains}(ii).

\emph{The H\"older components of $\psi^{0,\#}$.} Let $(x,x')\in\mathrm{Pz}_k$. Then $(x/L,x'/L)$ is a pair of $\mathrm{Pz}_{k+1}$ (Section~\ref{sec:fluctuation-integral-new}), and the definitions of the scaling $\chi'\mapsto\chi'_L$ and of $\delta_{\alpha,\mathscr{T}'}$ give
\begin{align*}
\chi^0(x,x')&=L^{-1-\alpha}\cdot\chi'(x/L,x'/L)\\
&=L^{-1-\alpha}\cdot|x/L-x'/L|^{-\alpha}\\
&\qquad\cdot\left(\exp\left\{i\,e_{k+1}\cdot\eta_{k+1}\cdot(\tau\mathscr{T}')(x/L,x'/L)\right\}\cdot\psi'(x'/L)-\psi'(x/L)\right).
\end{align*}
In this expression $|x/L-x'/L|^{-\alpha}=L^{\alpha}\cdot|x-x'|^{-\alpha}$, and $\psi'(z/L)=L\cdot\psi^0(z)$ for every site $z$ of $T_k$, because $\psi^0=\psi'_L$. To treat the phase, recall that $e_{k+1}=L^{\frac{1}{2}}\cdot e_k$ by hypothesis~\ref{def:hypothesis-one-step} (Section~\ref{sec:hypotheses-conditions}) and that $\eta_{k+1}=L^{-1}\cdot\eta_k$, so that $e_{k+1}\cdot\eta_{k+1}=L^{-\frac{1}{2}}\cdot e_k\cdot\eta_k$; that the map $z\mapsto z/L$ carries the path $\Gamma^{\pi}(x,x')$ onto the path $\Gamma^{\pi}(x/L,x'/L)$, for each of the six orders $\pi$ (item~(4) of Section~\ref{sec:activity-before-reblocking}); and that $\mathscr{T}=\mathscr{A}'_L+\mathscr{C}'_L=\mathscr{T}'_L$, that is, $\mathscr{T}(b)=L^{-\frac{1}{2}}\cdot\mathscr{T}'(b/L)$ for every bond $b$ of $T_k$. Consequently, for every order $\pi$,
\[
e_{k+1}\cdot\eta_{k+1}\sum_{b'\in\Gamma^{\pi}(x/L,x'/L)}\mathscr{T}'(b')=e_k\cdot\eta_k\sum_{b\in\Gamma^{\pi}(x,x')}L^{-\frac{1}{2}}\cdot\mathscr{T}'(b/L)=e_k\cdot\eta_k\sum_{b\in\Gamma^{\pi}(x,x')}\mathscr{T}(b),
\]
and the mean over the six orders gives $e_{k+1}\cdot\eta_{k+1}\cdot(\tau\mathscr{T}')(x/L,x'/L)=e_k\cdot\eta_k\cdot(\tau\mathscr{T})(x,x')$. Inserting these three relations one obtains
\[
\chi^0(x,x')=|x-x'|^{-\alpha}\cdot\left(\exp\left\{i\,e_k\cdot\eta_k\cdot(\tau\mathscr{T})(x,x')\right\}\cdot\psi^0(x')-\psi^0(x)\right)=\left(\delta_{\alpha,\mathscr{T}}\psi^0\right)(x,x').
\]
Thus $\chi^0=\delta_{\alpha,\mathscr{T}}\psi^0$, and the same computation, with $-\mathscr{T}'$ and $-\mathscr{T}$ in place of $\mathscr{T}'$ and $\mathscr{T}$, gives $\bar{\chi}^0=\delta_{\alpha,-\mathscr{T}}\bar{\psi}^0$. As $\mathscr{T}=\mathscr{T}'_L$ lies in $\tilde{G}_k$ (Step~0), Proposition~\ref{prop:comparison-fermion-operators}(b) yields identity~\eqref{eq:renormalisation-group-53}: $\psi^0=H_k(\mathscr{T})\,\Psi^c$ and $\bar{\psi}^0=\bar{H}_k(\mathscr{T})\,\bar{\Psi}^c$. Together with the formulas for $\chi^0$ and $\bar{\chi}^0$ and with the definition of $H^{\#}_k$ in Section~\ref{sec:fermion-integrals-blocking}, this gives
\begin{equation}\label{eq:renormalisation-group-66}
\psi^{0,\#}=H^{\#}_k(\mathscr{T})\,\Psi^c\qquad\text{(on the unbarred and on the barred fields).}
\end{equation}

\emph{The parameters at the value one.} Denote by $s=1$ the family of parameters all of whose entries are equal to $1$; it lies in $\mathrm{Pd}_k$. By property~\ref{it:renormalisation-group-1}, for $k=0$ by property~\ref{it:renormalisation-group-17}, one has $H_k(1)=H_k$ and $C_k^{\frac{1}{2}}(1)=C_k^{\frac{1}{2}}$, hence $\mathscr{Z}(1)=\mathscr{Z}$, and for every $B\in\tilde{G}_k$
\[
H_k(1,B)=H_k(B),\qquad\Gamma_k(1,B)=\Gamma_k(B),\qquad\sum_{X\in D_k}M_k(X,B)=M_k(B)=b_k-D_k(B)\quad(k\ge1),
\]
where the last equality is the definition of $M_k(B)$ in Section~\ref{sec:exact-fermion-operators}. As $-B$ lies in $\tilde{G}_k$ too (Lemma~\ref{lem:symmetries-domains}), the relation for $H_k(1,\cdot\,)$ at $-B$ gives
\[
\bar{H}_k(1,B)=\mathsf{C}^{-1}\cdot H_k(1,-B)\cdot\mathsf{C}=\mathsf{C}^{-1}\cdot H_k(-B)\cdot\mathsf{C}=\bar{H}_k(B).
\]
Therefore $H^{\#}_k(1,B)=H^{\#}_k(B)$ on the unbarred and on the barred fields, and $W^{\#}_k(1,\mathscr{T})$ is the image under $H^{\#}_k(\mathscr{T})$ of $\Gamma_k(\mathscr{T})\,W$ on the unbarred and of $\bar{W}$ on the barred fields. Since $\Psi^f=\Psi^c+\Gamma_k(\mathscr{T})\,W$, $\bar{\Psi}^f=\bar{\Psi}^c+\bar{W}$ and $H^{\#}_k(\mathscr{T})$ is linear, \eqref{eq:renormalisation-group-66} implies
\begin{equation}\label{eq:renormalisation-group-67}
H^{\#}_k(\mathscr{T})\,\Psi^f=\psi^{0,\#}+W^{\#}_k(1,\mathscr{T})\qquad\text{(on the unbarred and on the barred fields).}
\end{equation}

\emph{The evaluations $\mathrm{ev}$ and $\mathrm{ev}_f$.} Let $\mathrm{Gr}[\Psi',W]$ be the Grassmann algebra over $\mathbb{C}$ with the generators $\Psi'(u')$, $u'=(y',\boldsymbol{\beta},\boldsymbol{\omega})$, $y'\in T^0_{k+1}$, and $W(u)$, $u=(y,\boldsymbol{\beta},\boldsymbol{\omega})$, $y\in T^0_k$, barred and unbarred. Let
\[
\mathrm{ev}\colon\mathrm{Gr}[\psi,\chi,W]\longrightarrow\mathrm{Gr}[\Psi',W]
\]
be the homomorphism of algebras that sends each generator $\psi$, $\chi$, $\bar{\psi}$, $\bar{\chi}$ to the corresponding entry of $\psi^0$, $\chi^0$, $\bar{\psi}^0$, $\bar{\chi}^0$, and each generator $W(u)$ to itself. For $k\ge1$ let
\[
\mathrm{ev}_f\colon\mathrm{Gr}[\psi,\chi,\Psi]\longrightarrow\mathrm{Gr}[\Psi',W]
\]
be the homomorphism of algebras that sends $\Psi$ to $\Psi^f$, $\bar{\Psi}$ to $\bar{\Psi}^f$, and $\psi^{\#}$ to $H^{\#}_k(\mathscr{T})\,\Psi^f$, that is, the generators $\psi$, $\chi$, $\bar{\psi}$, $\bar{\chi}$ to the entries of
\[
H_k(\mathscr{T})\,\Psi^f,\qquad\delta_{\alpha,\mathscr{T}}H_k(\mathscr{T})\,\Psi^f,\qquad\bar{H}_k(\mathscr{T})\,\bar{\Psi}^f,\qquad\delta_{\alpha,-\mathscr{T}}\bar{H}_k(\mathscr{T})\,\bar{\Psi}^f.
\]
For $k=0$ there is no generator $\Psi$, and $\mathrm{ev}_f\colon\mathrm{Gr}[\psi,\chi]\to\mathrm{Gr}[\Psi',W]$ is the homomorphism of algebras that sends $\psi^{\#}$ to $H^{\#}_0(\mathscr{T})\,\Psi^f$, that is, $\psi$, $\chi$, $\bar{\psi}$, $\bar{\chi}$ to the entries of $\Psi^f$, $\delta_{\alpha,\mathscr{T}}\Psi^f$, $\bar{\Psi}^f$, $\delta_{\alpha,-\mathscr{T}}\bar{\Psi}^f$. These prescriptions define homomorphisms of algebras, because every generator is sent to an odd element. The Grassmann algebras in question are complex vector spaces of finite dimension, on which every linear map is continuous; consequently $\mathrm{ev}$ and $\mathrm{ev}_f$ may be interchanged with each of the integrations in the variables $t$, $s$ and $\mathsf{Z}$ that occur in~\eqref{eq:renormalisation-group-57}, \eqref{eq:renormalisation-group-59} and~\eqref{eq:renormalisation-group-61}. The integration over $W$ may be interchanged with $\mathrm{ev}$ as well: this homomorphism leaves every generator $W(u)$ fixed, and the images of the generators $\psi$, $\chi$, $\bar{\psi}$, $\bar{\chi}$ contain no generator $W$.

\emph{Step~1: the first localisation at $s=1$.} One has
\[
\mathrm{ev}_f\left(\sum_{X\in D_k}\sum_{i\le3}G^{(i)}(1,X)\right)=\mathrm{Act}^x_k[m',\mathscr{E}^r](A+\mathscr{Z},\mathscr{C};\Psi^f)-\mathrm{Act}^x_k[m',\mathscr{E}^r](A,\mathscr{C};\Psi^f).
\]
By Lemma~\ref{lem:bounds-analyticity}(q1), applied with $\mathfrak{Z}=\mathscr{Z}(1)=\mathscr{Z}$, for $|t|\le\rho_1$ the connection $A+t\cdot\mathscr{Z}$ lies in $\tilde{R}_k$ and $\mathscr{T}+t\cdot\mathscr{Z}$ lies in $\tilde{G}_k$; in particular $(A+\mathscr{Z},\mathscr{C})\in\tilde{R}_k\times\mathfrak{X}_k$, and the right side is defined (Section~\ref{sec:fermion-integrals-blocking}).

Let first $k\ge1$. The family $\mathscr{E}^r=R\,\mathscr{E}$ lies in $K^x_k$ by Proposition~\ref{prop:symmetries-extraction}, so that its kernels are analytic on $\tilde{R}_k\times\mathfrak{X}_k$ by property~\eqref{it:renormalisation-group-1-k1}. With Lemma~\ref{lem:bounds-analyticity}(q5) and $\rho_2\le\rho_1$ it follows that, at $s=1$, the integrands of the contour integrals defining $G^{(1)}$, $G^{(2)}$ and $G^{(3)}$ are continuous on the closed discs $|t|\le\rho_1$, resp.\ $|t|\le\rho_2$, and analytic in the open discs. By Section~\ref{sec:notation} each of these contour integrals is the value of its integrand at $t=1$ minus its value at $t=0$. For $i=1$, this and the relations at $s=1$ give
\begin{align*}
\sum_{X\in D_k}\mathrm{ev}_f\left(G^{(1)}(1,X)\right)&=\sum_{X\in D_k}\left(\left\langle\bar{\Psi}^f,M_k(X,\mathscr{T}+\mathscr{Z})\,\Psi^f\right\rangle-\left\langle\bar{\Psi}^f,M_k(X,\mathscr{T})\,\Psi^f\right\rangle\right)\\
&=\left\langle\bar{\Psi}^f,\left(M_k(\mathscr{T}+\mathscr{Z})-M_k(\mathscr{T})\right)\Psi^f\right\rangle=-\left\langle\bar{\Psi}^f,\left(D_k(\mathscr{T}+\mathscr{Z})-D_k(\mathscr{T})\right)\Psi^f\right\rangle.
\end{align*}
For $i=2,3$ note that $J^{\#}_k(1,\mathscr{T},\mathscr{Z})=H^{\#}_k(\mathscr{T}+\mathscr{Z})-H^{\#}_k(\mathscr{T})$, by~\eqref{eq:renormalisation-group-56} and the relations at $s=1$. Hence $\mathrm{ev}_f$ sends $\psi^{\#}+J^{\#}_k(1,\mathscr{T},\mathscr{Z})\,\Psi$ to
\[
H^{\#}_k(\mathscr{T})\,\Psi^f+\left(H^{\#}_k(\mathscr{T}+\mathscr{Z})-H^{\#}_k(\mathscr{T})\right)\Psi^f=H^{\#}_k(\mathscr{T}+\mathscr{Z})\,\Psi^f,
\]
on the unbarred and on the barred fields. The element $G^{(2)}(1,X)$ is the difference of $\mathscr{E}^r(X;A+\mathscr{Z},\mathscr{C},\cdot\,)$, evaluated at $\psi^{\#}+J^{\#}_k(1,\mathscr{T},\mathscr{Z})\,\Psi$, and of $\mathscr{E}^r(X;A,\mathscr{C},\cdot\,)$, evaluated at $\psi^{\#}$; therefore
\[
\sum_{X\in D_k}\mathrm{ev}_f\left(G^{(2)}(1,X)\right)=\sum_{X\in D_k}\left(\mathscr{E}^r\left(X;A+\mathscr{Z},\mathscr{C},H^{\#}_k(\mathscr{T}+\mathscr{Z})\,\Psi^f\right)-\mathscr{E}^r\left(X;A,\mathscr{C},H^{\#}_k(\mathscr{T})\,\Psi^f\right)\right).
\]
In the same way the components $J$ and $\bar{J}$ of $J^{\#}_k(1,\mathscr{T},\mathscr{Z})$ without H\"older derivative satisfy $\mathrm{ev}_f(\psi+J\,\Psi)=H_k(\mathscr{T}+\mathscr{Z})\,\Psi^f$ and $\mathrm{ev}_f(\bar{\psi}+\bar{J}\,\bar{\Psi})=\bar{H}_k(\mathscr{T}+\mathscr{Z})\,\bar{\Psi}^f$. Since the $M$-cubes partition the sites, by~\eqref{it:renormalisation-group-g0}, the sum over the $M$-cubes $\square$ is
\[
\sum_{\square}\mathrm{ev}_f\left(G^{(3)}(1,\square)\right)=-m'\cdot\left\langle\bar{H}_k(\mathscr{T}+\mathscr{Z})\,\bar{\Psi}^f,H_k(\mathscr{T}+\mathscr{Z})\,\Psi^f\right\rangle+m'\cdot\left\langle\bar{H}_k(\mathscr{T})\,\bar{\Psi}^f,H_k(\mathscr{T})\,\Psi^f\right\rangle.
\]
By the definitions of $I^x_k$ and $\mathrm{Act}^x_k$ in Section~\ref{sec:fermion-integrals-blocking}, in which the total connection of the pair $(A+\mathscr{Z},\mathscr{C})$ is $\mathscr{T}+\mathscr{Z}$ and that of the pair $(A,\mathscr{C})$ is $\mathscr{T}$, the sum of the three expressions obtained for $i=1$, $i=2$ and $i=3$ is the right side of the assertion of Step~1.

Let now $k=0$. Then $\eta_0=1$, $\psi=\Psi$ and $D_0(B)=D_B+\mathfrak{m}$ (Section~\ref{sec:exact-fermion-operators}); moreover $m'=m$ and $\mathscr{E}^r=0$, and $G^{(2)}=G^{(3)}=0$ by~\eqref{eq:renormalisation-group-60}. Let $x$ be a site, $\nu$ a direction and $\boldsymbol{\sigma}\in\{1,-1\}$. The sites $x$ and $x^{\boldsymbol{\sigma}}$ are nearest neighbours, $|x-x^{\boldsymbol{\sigma}}|=1$, and the six rectilinear paths from $x$ to $x^{\boldsymbol{\sigma}}$ coincide with the bond $(x,x^{\boldsymbol{\sigma}})$. By the definition of the H\"older derivative (Definition~\ref{def:pairs-sites-holder}) it follows that
\[
\mathrm{ev}_f\left(\psi(x)+\chi(x,x^{\boldsymbol{\sigma}})\right)=u_{\mathscr{T}}(x,x^{\boldsymbol{\sigma}})\cdot\Psi^f(x^{\boldsymbol{\sigma}}),\qquad\mathrm{ev}_f\left(\bar{\psi}(x)+\bar{\chi}(x,x^{\boldsymbol{\sigma}})\right)=u_{\mathscr{T}}(x,x^{\boldsymbol{\sigma}})^{-1}\cdot\bar{\Psi}^f(x^{\boldsymbol{\sigma}}).
\]
Let $a_1(x,\nu,\boldsymbol{\sigma})$ and $a_2(x,\nu,\boldsymbol{\sigma})$ be the first and the second summand inside the braces of~\eqref{eq:renormalisation-group-60} at $s'=1$. The exponentials in them carry the charge $e_k$ of the level $k=0$, and $\exp\left(\pm i\,e_k\cdot\mathscr{Z}(1)(x,x^{\boldsymbol{\sigma}})\right)=u_{\mathscr{Z}}(x,x^{\boldsymbol{\sigma}})^{\pm1}$, by the definition of $u_B$ in Section~\ref{sec:notation}, since $\eta_0=1$ and $\mathscr{Z}(1)=\mathscr{Z}$. Since $u_{\mathscr{T}}\cdot u_{\mathscr{Z}}=u_{\mathscr{T}+\mathscr{Z}}$,
\[
\mathrm{ev}_f\left(a_1(x,\nu,\boldsymbol{\sigma})\right)=\bar{\Psi}^f(x)\,P^{\boldsymbol{\sigma}}_{\nu}\left(u_{\mathscr{T}+\mathscr{Z}}(x,x^{\boldsymbol{\sigma}})-u_{\mathscr{T}}(x,x^{\boldsymbol{\sigma}})\right)\Psi^f(x^{\boldsymbol{\sigma}}),
\]
and by~\eqref{eq:renormalisation-group-39}, with $\eta_0=1$, the sum of these elements over all triples $(x,\nu,\boldsymbol{\sigma})$ equals $-\left\langle\bar{\Psi}^f,\left(D_{\mathscr{T}+\mathscr{Z}}-D_{\mathscr{T}}\right)\Psi^f\right\rangle$. Further,
\[
\mathrm{ev}_f\left(a_2(x,\nu,\boldsymbol{\sigma})\right)=u_{\mathscr{T}}(x,x^{\boldsymbol{\sigma}})^{-1}\cdot\bar{\Psi}^f(x^{\boldsymbol{\sigma}})\,P^{-\boldsymbol{\sigma}}_{\nu}\left(u_{\mathscr{Z}}(x,x^{\boldsymbol{\sigma}})^{-1}-1\right)\Psi^f(x).
\]
Put $y:=x^{\boldsymbol{\sigma}}$. Then $x=y^{-\boldsymbol{\sigma}}$, and $u_B(x,y)^{-1}=u_B(y,x)$ for every bond function $B$, as bond functions are odd under the reversal of the bonds (Section~\ref{sec:averaging-fluctuations-gauge}). Hence $\mathrm{ev}_f\left(a_2(x,\nu,\boldsymbol{\sigma})\right)=\mathrm{ev}_f\left(a_1(y,\nu,-\boldsymbol{\sigma})\right)$. The map $(x,\nu,\boldsymbol{\sigma})\mapsto(x^{\boldsymbol{\sigma}},\nu,-\boldsymbol{\sigma})$ is a bijection of the set of the triples onto itself. Summing the last identity over all triples and reindexing the right side by this bijection, one finds that $\mathrm{ev}_f$ takes one and the same value on $\sum a_1$ and on $\sum a_2$, both sums being over all triples; by linearity, $\mathrm{ev}_f$ takes this value also on $\frac{1}{2}\sum\left(a_1+a_2\right)$. By~\eqref{it:renormalisation-group-g0} the latter element is $\sum_{\square}G^{(1)}(1,\square)$, and therefore, as $D_0(B)=D_B+\mathfrak{m}$,
\[
\mathrm{ev}_f\left(\sum_{\square}G^{(1)}(1,\square)\right)=-\left\langle\bar{\Psi}^f,\left(D_0(\mathscr{T}+\mathscr{Z})-D_0(\mathscr{T})\right)\Psi^f\right\rangle.
\]
On the other side, $\mathscr{E}^r=0$, and $-m\cdot\langle\bar{\psi},\psi\rangle$ at $(\psi,\bar{\psi})=(\Psi^f,\bar{\Psi}^f)$ does not depend on the connection; by Section~\ref{sec:fermion-integrals-blocking} the right side of the assertion of Step~1 is the same expression.

\emph{Step~2: the substitution~\eqref{eq:renormalisation-group-58}.} Let $k\ge1$ and let $\mathbf{Y}$ be a polymer. Let $F\in\mathrm{Gr}[\psi,\chi,\Psi]$ be an element to which the substitution~\eqref{eq:renormalisation-group-58} is applied with the polymer $\mathbf{Y}$, and which contains the generators $\Psi(u)$, $\bar{\Psi}(u)$ only for $u$ at unit sites of $\mathbf{Y}$. Let $F'\in\mathrm{Gr}[\psi,\chi]$ be the image of $F$ under this substitution. Then the image of $F'$ under the substitution $\psi^{\#}\mapsto H^{\#}_k(\mathscr{T})\,\Psi^f$, which is $\mathrm{ev}_f(F')$, equals $\mathrm{ev}_f(F)$. Indeed, this substitution replaces $\psi$ by $H_k(\mathscr{T})\,\Psi^f$ and $\chi$ by $\delta_{\alpha,\mathscr{T}}H_k(\mathscr{T})\,\Psi^f$. For a unit site $y$ of $\mathbf{Y}$, Lemma~\ref{lem:interior-boundary-parts}(a), applied with $B=\mathscr{T}\in\tilde{G}_k$, and Proposition~\ref{prop:left-inverse} give
\[
\left(T^{(1)}(\mathscr{T})\cdot H_k(\mathscr{T})\,\Psi^f\right)(y)+\left(T^{(2)}(\mathscr{T})\cdot\delta_{\alpha,\mathscr{T}}H_k(\mathscr{T})\,\Psi^f\right)(y)=\left(T_k(\mathscr{T})\cdot H_k(\mathscr{T})\,\Psi^f\right)(y)=\Psi^f(y),
\]
and in the same way the element that~\eqref{eq:renormalisation-group-58} puts in the place of $\bar{\Psi}(y)$ is sent to $\bar{\Psi}^f(y)$. So the substitution~\eqref{eq:renormalisation-group-58} followed by $\psi^{\#}\mapsto H^{\#}_k(\mathscr{T})\,\Psi^f$ agrees with $\mathrm{ev}_f$ on the generators $\psi^{\#}$ and on the generators $\Psi(u)$, $\bar{\Psi}(u)$ with $u$ at unit sites of $\mathbf{Y}$; both being homomorphisms of algebras, they agree on the subalgebra with these generators, which contains $F$. For $k=0$ there is no substitution~\eqref{eq:renormalisation-group-58}, one has $E^{(1),x}=E^{(1),*}$, and on $\mathrm{Gr}[\psi,\chi]$ the homomorphism $\mathrm{ev}_f$ is the substitution $\psi^{\#}\mapsto H^{\#}_0(\mathscr{T})\,\Psi^f$.

\emph{Step~3: the exponent.} One has
\[
\mathrm{ev}\left(\sum_{\mathbf{Y}\in D_k}E^{(\le5),x}(\mathbf{Y})\right)=V^{\mathrm{ex},x}_k.
\]
In the second localisation only the parameters $s''$ occur, and $1$ stands for $s''=1$. By Proposition~\ref{prop:bounds-regularity-localised}(d) the left side equals $\mathrm{ev}\left(\sum_{X\in D_k}\left(G^{(\le3)}+G^{(4)}+G^{(5)}\right)(1,X)\right)$.

By~\eqref{eq:renormalisation-group-59}, $G^{(\le3)}(1,X)$ is the image of $E^{(\le3),x}(X)$ under $\psi^{\#}\mapsto\psi^{\#}+W^{\#}_k(1,\mathscr{T})$. Since $\mathrm{ev}$ fixes $W$, the element $\mathrm{ev}\left(G^{(\le3)}(1,X)\right)$ is the image of $E^{(\le3),x}(X)$ under $\psi^{\#}\mapsto\psi^{0,\#}+W^{\#}_k(1,\mathscr{T})$, that is, by~\eqref{eq:renormalisation-group-67}, under $\psi^{\#}\mapsto H^{\#}_k(\mathscr{T})\,\Psi^f$. One has
\[
\mathrm{ev}\left(G^{(\le3)}(1,X)\right)=\mathrm{ev}_f\left(\sum_{i\le3}E^{(i),*}(X)\right).
\]
Indeed, let first $k\ge1$. Then $E^{(\le3),x}(X)=\sum_{i\le3}E^{(i),x}(X)$, where $E^{(i),x}(X)$ is the image of $E^{(i),*}(X)$ under the substitution~\eqref{eq:renormalisation-group-58} with the polymer $X$; and $E^{(i),*}(X)$ contains the generators $\Psi(u)$, $\bar{\Psi}(u)$ only for $u$ at unit sites of $X$, by Lemma~\ref{lem:renormalisation-group-1}(a) and~\eqref{eq:renormalisation-group-57}, as shown in the proof of Lemma~\ref{lem:renormalisation-group-1}. Step~2, applied with $\mathbf{Y}=X$ to $F=E^{(i),*}(X)$ for $i=1,2,3$, gives the identity. Let now $k=0$. Then $G^{(2)}=G^{(3)}=0$ by~\eqref{eq:renormalisation-group-60}; as $\mathrm{Loc}$ is linear, $E^{(2),*}=E^{(3),*}=0$, hence $E^{(2),x}=E^{(3),x}=0$, and $E^{(\le3),x}(X)=E^{(1),x}(X)=E^{(1),*}(X)$ by Step~2. This is an element of $\mathrm{Gr}[\psi,\chi]$, on which $\mathrm{ev}_f$ is the substitution $\psi^{\#}\mapsto H^{\#}_0(\mathscr{T})\,\Psi^f$, by the last sentence of Step~2; so its image under this substitution is $\mathrm{ev}_f\left(E^{(1),*}(X)\right)=\mathrm{ev}_f\left(\sum_{i\le3}E^{(i),*}(X)\right)$.
Summing over $X$ and using Proposition~\ref{prop:bounds-regularity-localised}(d) and Step~1, one obtains
\begin{align*}
\sum_{X\in D_k}\mathrm{ev}\left(G^{(\le3)}(1,X)\right)&=\mathrm{ev}_f\left(\sum_{X\in D_k}\sum_{i\le3}G^{(i)}(1,X)\right)\\
&=\mathrm{Act}^x_k[m',\mathscr{E}^r](A+\mathscr{Z},\mathscr{C};\Psi^f)-\mathrm{Act}^x_k[m',\mathscr{E}^r](A,\mathscr{C};\Psi^f).
\end{align*}
The integrands of the contour integrals defining $G^{(4)}$ and $G^{(5)}$ are polynomials in $t$, and $\rho_4>1$; by Section~\ref{sec:notation} these contour integrals equal the value of the integrand at $t=1$ minus its value at $t=0$. Together with~\eqref{eq:renormalisation-group-67}, \eqref{eq:renormalisation-group-66}, the partition of the sites into $M$-cubes and the definition of $I^x_k$ in Section~\ref{sec:fermion-integrals-blocking}, this gives
\begin{align*}
\mathrm{ev}\left(\sum_{X\in D_k}\left(G^{(4)}+G^{(5)}\right)(1,X)\right)={}&\sum_{X\in D_k}\left(\mathscr{E}^r\left(X;A,\mathscr{C},H^{\#}_k(\mathscr{T})\,\Psi^f\right)-\mathscr{E}^r\left(X;A,\mathscr{C},H^{\#}_k(\mathscr{T})\,\Psi^c\right)\right)\\
&-m'\cdot\left\langle\bar{H}_k(\mathscr{T})\,\bar{\Psi}^f,H_k(\mathscr{T})\,\Psi^f\right\rangle+m'\cdot\left\langle\bar{H}_k(\mathscr{T})\,\bar{\Psi}^c,H_k(\mathscr{T})\,\Psi^c\right\rangle\\
={}&I^x_k[m',\mathscr{E}^r](A,\mathscr{C};\Psi^f)-I^x_k[m',\mathscr{E}^r](A,\mathscr{C};\Psi^c).
\end{align*}
Adding this identity to the preceding one for $\sum_X\mathrm{ev}\left(G^{(\le3)}(1,X)\right)$, one obtains $V^{\mathrm{ex},x}_k$, by its definition in Proposition~\ref{prop:step-external-connection}.

\emph{Step~4: the fluctuation integral.} By item~(1) of Section~\ref{sec:activity-before-reblocking}, the elements $V(X):=E^{(\le5),x}(X)$, $X\in D_k$, satisfy the hypotheses of Lemma~\ref{lem:integration-polymer-exponential} with the weights $h_k/4$, and the elements that this lemma forms from them by~\eqref{eq:renormalisation-group-61} are the $E^{\#,x}_k(X)$. Hence, in $\mathrm{Gr}[\psi,\chi]$,
\[
\int\mathrm{d}\mu_I(W)\,\mathrm{d}\hat{\mu}(\mathsf{Z})\,\exp\left\{\sum_{X\in D_k}E^{(\le5),x}(X)\right\}=\exp\left\{\sum_{X\in D_k}E^{\#,x}_k(X;A,\mathscr{C},\psi,\chi)\right\}.
\]
Apply $\mathrm{ev}$ to both sides. Being a continuous homomorphism of algebras, $\mathrm{ev}$ commutes with the exponential and with the integrations over $W$ and over $\mathsf{Z}$; by Step~3 the left side becomes $\int\mathrm{d}\mu_I(W)\,\mathrm{d}\hat{\mu}(\mathsf{Z})\,\exp\left(V^{\mathrm{ex},x}_k\right)=\Xi^x_k$. Therefore
\[
\Xi^x_k=\exp\left\{\sum_{X\in D_k}E^{\#,x}_k\left(X;A,\mathscr{C},\psi^{0,\#}\right)\right\}.
\]

\emph{Step~5: the determinant and the regrouping.} Proposition~\ref{prop:local-expansion-determinant}(d0) at $B=\mathscr{T}\in\tilde{G}_k$ and Definition~\ref{def:determinant-functions} give
\begin{align*}
\frac{\delta Z_k(\mathscr{T})}{\delta Z_k(0)}&=\exp\left\{\sum_{X\in D_k}\left(E^d_k(X,\mathscr{T})-E^d_k(X,0)\right)\right\}\\
&=\exp\left\{\sum_{X\in D_k}\delta E^{\mathrm{det}}_k(X;A,\mathscr{C})+\sum_{X\in D_k}E^{\mathrm{det}}_k(X,\mathscr{C})\right\},
\end{align*}
because $E^d_k(X,A+\mathscr{C})-E^d_k(X,0)=\delta E^{\mathrm{det}}_k(X;A,\mathscr{C})+E^{\mathrm{det}}_k(X,\mathscr{C})$. Next, $\mathscr{E}^r=R\,\mathscr{E}$ lies in $K^x_k$ (Proposition~\ref{prop:symmetries-extraction}), the pair $(A,\mathscr{C})$ lies in $\tilde{R}_k\times\mathfrak{X}_k$ (Step~0), and $\chi^0=\delta_{\alpha,\mathscr{T}}\psi^0$, $\bar{\chi}^0=\delta_{\alpha,-\mathscr{T}}\bar{\psi}^0$; so Proposition~\ref{prop:adjustment} yields
\[
\sum_{X\in D_k}\mathscr{E}^r\left(X;A,\mathscr{C},\psi^{0,\#}\right)=\sum_{X\in D_k}(R\,\mathscr{E})^{\mathrm{nat}}\left(X;A,\mathscr{C},\psi^{0,\#}\right).
\]
By item~(5) of Section~\ref{sec:activity-before-reblocking}, for every $X\in D_k$ there is exactly one $Y\in D_{k+1}$ with $\bar{X}^L=LY$; so a sum over $X\in D_k$ is the sum over $Y\in D_{k+1}$ of the sums over the polymers $X$ with $\bar{X}^L=LY$. As $A=\mathscr{A}'_L$, $\mathscr{C}=\mathscr{C}'_L$, $\psi^{0,\#}=\left(\psi'^{\#}\right)_L$ and $(\mathscr{A}',\mathscr{C}')\in\tilde{R}_{k+1}\times\mathfrak{X}_{k+1}$, the definition~\eqref{eq:renormalisation-group-62} of $\mathscr{E}^{+}$ and the definition of $\mathfrak{P}_{k+1}$ in clause~\eqref{it:step-small-fields-c} give
\begin{align*}
\sum_{X\in D_k}\left[(R\,\mathscr{E})^{\mathrm{nat}}+\delta E^{\mathrm{det}}_k+E^{\#,x}_k(m,\mathscr{E})\right]\left(X;\mathscr{A}'_L,\mathscr{C}'_L,\left(\psi'^{\#}\right)_L\right)&=\sum_{Y\in D_{k+1}}\mathscr{E}^{+}\left(Y;\mathscr{A}',\mathscr{C}',\psi'^{\#}\right),\\
\sum_{X\in D_k}\left(\mathfrak{P}_k(X,\mathscr{C})+E^{\mathrm{det}}_k(X,\mathscr{C})\right)&=\sum_{Y\in D_{k+1}}\mathfrak{P}_{k+1}(Y,\mathscr{C}').
\end{align*}

\emph{Proof of~\eqref{eq:renormalisation-group-54}.} Insert the results of Steps~4 and~5 into identity~\eqref{eq:renormalisation-group-65} of Proposition~\ref{prop:step-external-connection}, which holds for the family $\mathfrak{P}_k$. The exponents are even elements of the Grassmann algebra and commute with each other, so that the product of the exponentials is the exponential of the sum of the exponents; the right side of~\eqref{eq:renormalisation-group-65} thus becomes
\begin{multline*}
\left(\det D_{k+1}(0)\right)^{-1}\cdot\exp\left\{-\left\langle\bar{\Psi}',D_{k+1}(\mathscr{T}')\,\Psi'\right\rangle-m^{+}\cdot\left\langle\bar{\psi}',\psi'\right\rangle\vphantom{\sum_{Y\in D_{k+1}}}\right.\\
\left.{}+\sum_{Y\in D_{k+1}}\mathscr{E}^{+}\left(Y;\mathscr{A}',\mathscr{C}',\psi'^{\#}\right)-\boldsymbol{\varepsilon}^{+}\cdot n_{k+1}+\sum_{Y\in D_{k+1}}\mathfrak{P}_{k+1}(Y,\mathscr{C}')\right\}.
\end{multline*}
By clause~\eqref{it:step-small-fields-d} of Theorem~\ref{thm:step-small-fields} (Section~\ref{sec:data-next-level}), the family $\mathscr{E}^{+}$ lies in $K^x_{k+1}$. Since $\mathscr{A}'=H_{k+1}A'$ and $\mathscr{T}'=\mathscr{A}'+\mathscr{C}'$, the definitions of Section~\ref{sec:fermion-integrals-blocking} at the level $k+1$ identify the last expression with $\mathscr{F}^x_{k+1}[\boldsymbol{\varepsilon}^{+},m^{+},\mathscr{E}^{+}](A',\mathscr{C}';\Psi')$, the density being formed with the family $\mathfrak{P}_{k+1}$. The left side of~\eqref{eq:renormalisation-group-65} is $\left(\mathbf{T}^x_k\mathscr{F}^x_k[\boldsymbol{\varepsilon},m,\mathscr{E}]\right)(A',\mathscr{C}';\Psi')$. This proves~\eqref{eq:renormalisation-group-54}.

\emph{Proof of~\eqref{eq:renormalisation-group-55}.} Integrate both sides of~\eqref{eq:renormalisation-group-54} with $\int\mathrm{D}\Psi'$. On the right one obtains $y^x_{k+1}[\boldsymbol{\varepsilon}^{+},m^{+},\mathscr{E}^{+}](A';\mathscr{C}')$, by the definition of $y^x_{k+1}$ (Section~\ref{sec:fermion-integrals-blocking}). On the left, identity~\eqref{eq:renormalisation-group-3} gives
\begin{align*}
\int\mathrm{D}\Psi'\,\left(\mathbf{T}^x_k\mathscr{F}^x_k[\boldsymbol{\varepsilon},m,\mathscr{E}]\right)(A',\mathscr{C}';\Psi')&=\int_{\mathrm{Box}_k}\mathrm{d}\mu_I(\zeta)\int\mathrm{D}\Psi\,\mathscr{F}^x_k[\boldsymbol{\varepsilon},m,\mathscr{E}]\left(A^{\min}_k(A')+\mathrm{Sq}_k\zeta,\mathscr{C}'_L;\Psi\right)\\
&=\int_{\mathrm{Box}_k}\mathrm{d}\mu_I(\zeta)\,y^x_k[\boldsymbol{\varepsilon},m,\mathscr{E}]\left(A^{\min}_k(A')+\mathrm{Sq}_k\zeta;\mathscr{C}'_L\right),
\end{align*}
and this is~\eqref{eq:renormalisation-group-55}. Identity~\eqref{eq:renormalisation-group-3} is applicable to the map $F:=\mathscr{F}^x_k[\boldsymbol{\varepsilon},m,\mathscr{E}](\cdot\,,\mathscr{C}'_L;\cdot\,)$ from $\mathrm{Dom}_k$ to $\mathrm{Gr}[\Psi]$, for the following reasons. The values of $F$ are even. Both sides of~\eqref{eq:renormalisation-group-3} contain $F$ only at the fields $A^{\min}_k(A')+\mathrm{Sq}_k\zeta$, $\zeta\in\mathrm{Box}_k$; these fields form a compact set, the image of the compact set $\mathrm{Box}_k$ (Section~\ref{sec:averaging-fluctuations-gauge}) under a continuous map, and this set is contained in $\mathrm{Dom}_k$ by Proposition~\ref{prop:fields-domains}(ii). On this set $F$ is continuous (Step~(0) of the proof of Theorem~\ref{thm:one-step-formula}), hence bounded. The map that equals $F$ on this set and $0$ elsewhere on $\mathrm{Dom}_k$ is a bounded Borel map from $\mathrm{Dom}_k$ to the even part of $\mathrm{Gr}[\Psi]$, and it gives the same two sides of~\eqref{eq:renormalisation-group-3} as $F$.

\emph{The value $\mathscr{C}'=0$.} Let $\mathscr{C}'=0$, and let $\mathfrak{P}_k(X,0)=0$ for every $X\in D_k$. Then $\mathscr{C}=(0)_L=0$. By Definition~\ref{def:determinant-functions}, $E^{\mathrm{det}}_k(X,0)=E^d_k(X,0)-E^d_k(X,0)=0$ for every $X\in D_k$, so that $\mathfrak{P}_{k+1}(Y,0)=0$ for every $Y\in D_{k+1}$; hence the factors $\exp\left\{\sum_{X\in D_k}\mathfrak{P}_k(X,0)\right\}$ and $\exp\left\{\sum_{Y\in D_{k+1}}\mathfrak{P}_{k+1}(Y,0)\right\}$ of $\mathscr{F}^x_k$ and $\mathscr{F}^x_{k+1}$ are equal to $1$. By Section~\ref{sec:fermion-integrals-blocking} the objects $\mathscr{F}_j$, $y_j$ ($j=k,k+1$) and $\mathbf{T}_k$ without slot are the objects $\mathscr{F}^x_j$, $y^x_j$ and $\mathbf{T}^x_k$ at the value $0$ of the second slot, and these contain $\mathscr{E}$, resp.\ $\mathscr{E}^{+}$, only through $\mathscr{E}|_0$, resp.\ $\mathscr{E}^{+}|_0$. Moreover $\varepsilon(\mathscr{E})=\varepsilon(\mathscr{E}|_0)$ and $m(\mathscr{E})=m(\mathscr{E}|_0)$, because the numbers $\boldsymbol{\alpha}_0(\mathscr{E},X)$ and $\boldsymbol{\alpha}_2(\mathscr{E},X)$ of Section~\ref{sec:local-parts-adjustment} are formed at $(\mathscr{A},\mathscr{C})=(0,0)$; so $\boldsymbol{\varepsilon}^{+}$ and $m^{+}$ are also the numbers formed by the same rule from $(\boldsymbol{\varepsilon},m,\mathscr{E}|_0)$. Consequently, at $\mathscr{C}'=0$ the identities~\eqref{eq:renormalisation-group-54} and~\eqref{eq:renormalisation-group-55} are the two identities of the same form for the objects $\mathscr{F}_k$, $y_k$, $\mathbf{T}_k$ without slot, with the data $(\boldsymbol{\varepsilon},m,\mathscr{E}|_0)$ and their image $(\boldsymbol{\varepsilon}^{+},m^{+},\mathscr{E}^{+}|_0)$.

No step of the proof depends on the harmonic part of $A'$; the identities hold for every $A'\in\mathrm{Dom}_{k+1}$, whatever its harmonic part.
\end{proof}

\subsection{The data of level $k+1$, locality and uniformity: clauses~\eqref{it:step-small-fields-d}, \eqref{it:step-small-fields-g} and~\eqref{it:step-small-fields-h} of Theorem~\ref{thm:step-small-fields}}\label{sec:data-level-locality}

\subsubsection{The data of the next level}\label{sec:data-next-level}

\begin{proof}[Proof of clause~\eqref{it:step-small-fields-d} of Theorem~\ref{thm:step-small-fields}]
Let $\boldsymbol{\varepsilon}$, $m$ and $\mathscr{E}$ satisfy the hypotheses of clause~\eqref{it:step-small-fields-c} of Theorem~\ref{thm:step-small-fields}, and let $\boldsymbol{\varepsilon}^{+}$, $m^{+}$, $\mathscr{E}^{+}$ be the quantities defined in that clause. Thus $m$ is a real number with $|m|\le w^{\mathrm{m}}_k$, and $\mathscr{E}$ is an element of $K^x_k$ with $\|\mathscr{E}\|^x_k\le w^{\mathrm{E}}_k$, equal to $0$ if $k=0$. By the definition of the balls in Section~\ref{sec:spaces-activities}, the pair $(m,\mathscr{E})$ is a real datum in $\mathrm{Ball}^x_k(1)\subseteq\mathrm{Ball}^x_k(4)$.

Consider the three summands of
\[
\mathscr{E}^{+}=L R\,\mathscr{E}+L\,\delta E^{\mathrm{det}}_k+E^{*,x}_k(m,\mathscr{E}).
\]
Proposition~\ref{prop:contraction} places $L R\,\mathscr{E}$ in $K^x_{k+1}$ and bounds its norm by $c_{\mathrm{fl}}\cdot L^{-\frac{1}{4}+2\varepsilon}\cdot\|\mathscr{E}\|^x_k$. Proposition~\ref{prop:local-expansion-determinant}, parts (d5) and (d6), places $L\,\delta E^{\mathrm{det}}_k$ in $K^x_{k+1}$, with norm at most $e_{k+1}^{\frac{1}{4}-2\varepsilon}$. Clause~\eqref{it:step-small-fields-a} of Theorem~\ref{thm:step-small-fields}, applied to the real datum $(m,\mathscr{E})$, places $E^{*,x}_k(m,\mathscr{E})$ in $K^x_{k+1}$, with norm at most $e_k^{\frac{1}{4}-6\varepsilon}$. Since $K^x_{k+1}$ is a real vector space, $\mathscr{E}^{+}$ belongs to it, and the triangle inequality gives
\[
\left\|\mathscr{E}^{+}\right\|^x_{k+1}\le c_{\mathrm{fl}}\cdot L^{-\frac{1}{4}+2\varepsilon}\cdot\|\mathscr{E}\|^x_k+e_{k+1}^{\frac{1}{4}-2\varepsilon}+e_k^{\frac{1}{4}-6\varepsilon}.
\]

The restriction of a family to the second slot $\mathscr{C}=0$ (Definition~\ref{def:renormalisation-group-1}) is linear; so $\mathscr{E}^{+}|_0$ is the sum of the restrictions of the three summands. For the first summand, Proposition~\ref{prop:contraction} gives $\left(L R\,\mathscr{E}\right)|_0=L R\left(\mathscr{E}|_0\right)$. For the second, Definition~\ref{def:determinant-functions} gives
\[
\delta E^{\mathrm{det}}_k(X;\mathscr{A},0)=E^d_k(X,\mathscr{A})-E^d_k(X,0)=E^{\mathrm{det}}_k(X,\mathscr{A})\qquad(X\in D_k);
\]
since the reblocking $L$ of Section~\ref{sec:local-parts-adjustment} evaluates a family of level $k$ at the scaled connections, and the scaling sends the connection $0$ to $0$, one has $\left(L\,\delta E^{\mathrm{det}}_k\right)|_0=L\,E^{\mathrm{det}}_k$. For the third, clause~\eqref{it:step-small-fields-b} of Theorem~\ref{thm:step-small-fields} gives $E^{*,x}_k(m,\mathscr{E})|_0=E^{*}_k\left(m,\mathscr{E}|_0\right)$. Consequently
\[
\mathscr{E}^{+}|_0=L R\left(\mathscr{E}|_0\right)+L\,E^{\mathrm{det}}_k+E^{*}_k\left(m,\mathscr{E}|_0\right).
\]

By the definitions of $m^{+}$ and of $\boldsymbol{\varepsilon}^{+}$,
\[
m^{+}-L\cdot m=L\cdot m(\mathscr{E}),\qquad\boldsymbol{\varepsilon}^{+}-L^3\cdot\boldsymbol{\varepsilon}=L^3\cdot\left(\varepsilon(\mathscr{E})+d\varepsilon^0_k\right).
\]
Proposition~\ref{prop:symmetries-extraction} gives $|m(\mathscr{E})|\le c_{\mathrm{fl}}\cdot e_k^{\frac{1}{2}}\cdot\|\mathscr{E}\|^x_k$ and $|\varepsilon(\mathscr{E})|\le c_{\mathrm{fl}}\cdot\|\mathscr{E}\|^x_k$, and clause~\eqref{it:step-small-fields-a} of Theorem~\ref{thm:step-small-fields} gives $0\le d\varepsilon^0_k\le e_k^7$. It follows that
\[
\left|m^{+}-L\cdot m\right|\le L\cdot c_{\mathrm{fl}}\cdot e_k^{\frac{1}{2}}\cdot\|\mathscr{E}\|^x_k,\qquad\left|\boldsymbol{\varepsilon}^{+}-L^3\cdot\boldsymbol{\varepsilon}\right|\le L^3\cdot\left(c_{\mathrm{fl}}\cdot\|\mathscr{E}\|^x_k+e_k^7\right),
\]
and clause~\eqref{it:step-small-fields-d} is proved.
\end{proof}

\subsubsection{Locality of the step}\label{sec:locality-step}

Let $(m,\mathscr{E})\in\mathrm{Ball}^x_k(4)$. By Section~\ref{sec:bound-analyticity-after}, the family $E^{*,x}_k(m,\mathscr{E})$ has the properties~\eqref{it:renormalisation-group-1-k1} and~\eqref{it:renormalisation-group-1-k2} at level $k+1$.

The H\"older legs of $E^{*,x}_k(m,\mathscr{E})$ arise as follows. Let $\mathbf{Y}\in D_k$. Every H\"older leg of $E^{\#,x}_k(m,\mathscr{E})(\mathbf{Y})$ is either a leg of an element $\mathscr{E}^r(X)$ with $X\subseteq\mathbf{Y}$, or a leg introduced by the substitution~\eqref{eq:renormalisation-group-58} or by the function~\eqref{eq:renormalisation-group-60}. A leg of the first sort is a pair of kind~(I) or of kind~(II) for $X$, because $\mathscr{E}^r=R\,\mathscr{E}$ has the property~\eqref{it:renormalisation-group-1-k2} (Proposition~\ref{prop:symmetries-extraction}); by~\eqref{eq:renormalisation-group-adm}, it has the same kind with respect to $\mathbf{Y}$. A leg of the second sort is a pair $\left(x,x^{\boldsymbol{\sigma}}\right)$ of nearest neighbours whose first point $x$ belongs to the polymer with which the substitution is formed, resp.\ to the $M$-cube $\square$ of $G^{(1)}(s,\square)$ (Lemma~\ref{lem:interior-boundary-parts}(b) and the display~\eqref{eq:renormalisation-group-60}); for this polymer the pair is of kind~(II), and it is of kind~(I) as well when $x^{\boldsymbol{\sigma}}$ also belongs to it. For $Y\in D_{k+1}$, the kernels of $E^{*,x}_k(m,\mathscr{E})(Y)$ are obtained from those of the elements $E^{\#,x}_k(m,\mathscr{E})(X)$, $\bar{X}^L=LY$, through the scaling $x\mapsto x/L$ in~\eqref{eq:renormalisation-group-62}. By item~(4) of Section~\ref{sec:activity-before-reblocking}, this scaling carries a pair of kind~(I) for $X$ to a pair of kind~(I) for $Y$, whose length is at most $\frac{M}{L}$, and a pair of kind~(II) for $X$ to a pair of kind~(II) for $Y$.

\begin{proposition}[Local constituents of the step]\label{prop:local-constituents-step}
Let $(m,\mathscr{E})\in\mathrm{Ball}^x_k(4)$ and $\mathbf{Y}\in D_k$. The elements
\[
E^{(i),x}(\mathbf{Y})\quad(i\le3),\qquad E^{(\le5),x}(\mathbf{Y}),\qquad\mathsf{K}(\mathbf{Y}),\qquad\mathsf{H}(\mathbf{Y}),\qquad E^{\#,x}_k(\mathbf{Y})
\]
are called the \emph{objects attached to} $\mathbf{Y}$. By the displays~\eqref{eq:renormalisation-group-56}--\eqref{eq:renormalisation-group-61}, each of them depends on five kinds of arguments: the datum $(m,\mathscr{E})$, which is the subject of~(a); the system of primitives, which is assumed to have the properties~\ref{it:renormalisation-group-1}--\ref{it:renormalisation-group-17} for the parameters $\mathrm{par}$ and is the subject of~(c); the exact operators and forms, which are listed in~(b) and are the subject of that clause; the pair of connections $(A,\mathscr{C})$, which is the subject of the paragraph that follows the list; and the integration variables, which are the subject of~(d). The expressions ``reads'' and ``enters through'' have the meaning fixed in Definition~\ref{def:terminology}(1), and the expression ``given by one formula on $\mathbb{L}_k$'' has the meaning fixed in Definition~\ref{def:terminology}(2).
\begin{enumerate}
\item[(a)] \emph{Polymers and the datum.} If a term $(X,P)$ of~\eqref{eq:renormalisation-group-57} satisfies $Y'(X,P)=\mathbf{Y}$, then $X\subseteq\mathbf{Y}$, and each cube of $P_0(X)\cup P$ is contained in $\mathbf{Y}$. In~\eqref{eq:renormalisation-group-61}, each member of a set $\{X_i\}$ or of a family $(X_1,\dots,X_n)$ with union $\mathbf{Y}$ is a polymer contained in $\mathbf{Y}$. For $Y\in D_{k+1}$, each polymer $X$ with $\bar{X}^L=LY$, that is, each index of the sums in~\eqref{eq:renormalisation-group-62} that define $E^{*,x}_k(m,\mathscr{E})(Y)$ and $\mathscr{E}^{+}(Y)$, is contained in $LY$. The objects attached to $\mathbf{Y}$ depend on the datum only through $m'=m+m(\mathscr{E})$ and through $\mathscr{E}^r(X;\cdot)$, $X\subseteq\mathbf{Y}$. As for $\mathscr{E}$ itself, $m(\mathscr{E})$ and the differences $\mathscr{E}^r(X)-\mathscr{E}(X)$ depend on it only through $\mathscr{E}_{00}(X;0,0)$ and $\mathscr{E}_{20}(X;0,0;\cdot)$ with $X$ small.
\item[(b)] \emph{Exact operators.} Besides the primitives, the displays~\eqref{eq:renormalisation-group-56}--\eqref{eq:renormalisation-group-61} contain precisely the following exact operators and forms: the operator $C$; the H\"older derivatives $\delta_{\alpha,B}$; the operators $T^{(i)}(\mathscr{T})$ and $\bar{T}^{(i)}(\mathscr{T})$; the constant matrices $P^{\boldsymbol{\sigma}}_{\nu}$ and $\mathsf{C}$; and the forms $\langle\bar{\psi},\psi\rangle_{\square}$. They enter the objects attached to $\mathbf{Y}$ in the following way.
\begin{itemize}
\item The operator $\delta_{\alpha,B}$ enters only through the pairs that are admissible for polymers contained in $\mathbf{Y}$.
\item The operators $T^{(i)}(\mathscr{T})$ and $\bar{T}^{(i)}(\mathscr{T})$ enter only through their rows at unit sites of $\mathbf{Y}$; in $E^{(i),x}(\mathbf{Y})$ they are formed with the polymer $\mathbf{Y}$, and in the other objects, where they appear inside the functions $G^{(\le3)}(s'',X)$, with polymers $X\subseteq\mathbf{Y}$.
\item A form $\langle\bar{\psi},\psi\rangle_{\square}$ enters only if $\square$ is an $M$-cube of $\mathbf{Y}$.
\end{itemize}
All these operators and forms, the forms through their kernels, are given by one formula on $\mathbb{L}_k$, assembled from bonds, $L$-blocks, the unit blocks $B_k(y)$ together with the smaller blocks of the sides $L^j\eta_k$, $0\le j<k$, inside them, the six rectilinear paths joining the two points of a pair, the condition that a point belong to a given union of $M$-cubes, and the real numbers $b_k$, $\mathfrak{m}$, $e_k$, $\eta_k$, $\alpha$. The kernels of $C$ and of $T^{(i)}$ vanish at pairs of indices whose sup-distance is $2L$ or more, and $2L<\frac{M}{2}$. A row of $\delta_{\alpha,B}$ at one of the pairs just named, and a row of $T^{(i)}(\mathscr{T})$ or of $\bar{T}^{(i)}(\mathscr{T})$ at one of the unit sites just named, reads the connection only on bonds meeting $\mathbf{Y}$; no connection occurs in $C$, $P^{\boldsymbol{\sigma}}_{\nu}$, $\mathsf{C}$ and $\langle\bar{\psi},\psi\rangle_{\square}$. Let $\tilde{B}$ be a bond function on $\mathbb{L}_k$ and $\tilde{Y}$ a union of $M$-cubes of $\mathbb{L}_k$. The number $M$ is a range of each of these formulas in the sense of Definition~\ref{def:terminology}(2): the entry of the operator on $\mathbb{L}_k$ in a row $x$ and a column $x'$ vanishes if $\operatorname{supp}x$ and $\operatorname{supp}x'$ are at sup-distance greater than $M$, the support of a pair $(x,x')$ of points of $\mathbb{L}_k$ being $\{x,x'\}$; the row at $x$ reads $\tilde{B}$ only on bonds whose two end points lie within sup-distance $M$ of $\operatorname{supp}x$; and this row takes the same value at any two unions of $M$-cubes of $\mathbb{L}_k$ that contain the same sites within sup-distance $M$ of $\operatorname{supp}x$. In detail, the following holds. The operator $C$ contains neither $\tilde{B}$ nor $\tilde{Y}$, and its entries vanish if the supports of the row and of the column are at sup-distance $L$ or more. The row of $\delta_{\alpha,\tilde{B}}$ at a pair $(x,x')$ does not involve $\tilde{Y}$, has its columns at $x$ and $x'$, and reads $\tilde{B}$ on the bonds of the six paths $\Gamma^{\pi}(x,x')$, whose end points lie within sup-distance $|x-x'|\le M$ of $x$. The operators $T^{(i)}(\tilde{B})$ and $\bar{T}^{(i)}(\tilde{B})$ are formed with $\tilde{Y}$ by~\eqref{eq:renormalisation-group-58} at the unit sites of $\tilde{Y}$ and have zero rows at the other unit sites; at a unit site $y$ of $\tilde{Y}$ their rows have their columns at sup-distance less than $\frac{1}{2}+\eta_k\le\frac{3}{2}$ from $y$, read $\tilde{B}$ on the bonds with an end in $B_k(y)$, whose end points lie at sup-distance less than $\frac{3}{2}$ from $y$, and read $\tilde{Y}$ only through the membership of $y$ and of the sites $x^{\boldsymbol{\sigma}}$, $x\in B_k(y)$. The matrices $P^{\boldsymbol{\sigma}}_{\nu}$ and $\mathsf{C}$ act at each site and contain neither $\tilde{B}$ nor $\tilde{Y}$. The form $\langle\bar{\psi},\psi\rangle_{\square}$ enters through its kernel, which on $\mathbb{L}_k$ is $\eta_k^3\cdot\mathbf{1}_{\{x=x'\}}\cdot\mathbf{1}_{\tilde{Y}}(x)$ times the unit matrix, the cube $\square$ standing in the place of $\tilde{Y}$; the row of this kernel at $x$ has no column other than $x$, contains no bond function, and reads $\tilde{Y}$ only at $x$.
\item[(c)] \emph{Primitives.} The families~\eqref{it:primitives-p1}--\eqref{it:primitives-p3} and $\Gamma_k$ enter the objects attached to $\mathbf{Y}$ only through the values they take at parameter families vanishing off $\mathbf{Y}$ in the sense of Section~\ref{sec:properties-systems-primitives}, and, for such parameter families, only through the following rows:
\begin{itemize}
\item the rows of $H_k(s')$ indexed by fine bonds that meet $\mathbf{Y}$;
\item the rows of $C_k^{\frac{1}{2}}(s')$ indexed by basis elements contained in $\mathbf{Y}$;
\item the rows of $H_k(s'',\cdot)$ and of $\bar{H}_k(s'',\cdot)$ indexed by fine sites whose sup-distance from $\mathbf{Y}$ does not exceed $\eta_k$;
\item the rows of $\Gamma_k(s'',\cdot)$ indexed by unit sites of $\mathbf{Y}$.
\end{itemize}
The indices of all these rows belong to $\mathbf{Y}^{\sim1}$; and at every parameter family that vanishes off $\mathbf{Y}$, the non-zero entries of each row of the list stand in columns contained in $\mathbf{Y}$. Of the family $M_k$, only the operators $M_k(X,\cdot)$ with $X\subseteq\mathbf{Y}$ enter. The operators of these families that depend on a connection are evaluated at $\mathscr{T}+t\cdot\mathscr{Z}(s')$ with $|t|\le\rho_1$, or at $\mathscr{T}$; the listed rows and the operators $M_k(X,\cdot)$ read these connections only on bonds meeting $\mathbf{Y}$.
\item[(d)] \emph{Integrations.} The integrations that occur are: in $t$, along the circles $|t|=\rho_i$, $i\in\{1,2,4\}$; in $s_P$, over $[0,1]^P$; in the generators $W(u)$ with $u$ in $\mathbf{Y}$, against the Grassmann Gaussian integral $\mu_I$ with identity covariance, which by~\eqref{eq:renormalisation-group-43} splits into a product over the single unit sites; and in the coordinates $\mathsf{Z}(\Upsilon)$ with $\Upsilon$ contained in $\mathbf{Y}$, against a product of copies of a single probability measure on $\mathbb{R}$ that depends on $p_{0,k}$ alone (Section~\ref{sec:averaging-fluctuations-gauge} and Proposition~\ref{prop:basis-map-box}).
\end{enumerate}
Let $(A,\mathscr{C})$ and $(A^1,\mathscr{C}^1)$ be elements of $\tilde{R}^{-}_k\times\mathfrak{X}_k$ that coincide on the bonds that meet $\mathbf{Y}$. Then each of the objects attached to $\mathbf{Y}$ takes the same value at the two pairs; for $E^{(i),x}(\mathbf{Y})$, $i\le3$, for $E^{(\le5),x}(\mathbf{Y})$ and for $\mathsf{K}(\mathbf{Y})$, which depend on $\mathsf{Z}$ as well, this holds at every $\mathsf{Z}\in\mathrm{Box}_k$. In the terminology of Definition~\ref{def:terminology}(1), with $\tilde{R}^{-}_k\times\mathfrak{X}_k$ as the set of the admitted values, the objects attached to $\mathbf{Y}$ read $A$ and $\mathscr{C}$ on the bonds that meet $\mathbf{Y}$. Clauses (b) and (c) refer to the objects attached to $\mathbf{Y}$. In the summands $\left(R\,\mathscr{E}\right)^{\mathrm{nat}}(X)$ and $\delta E^{\mathrm{det}}_k(X)$ of $\mathscr{E}^{+}(Y)$ in~\eqref{eq:renormalisation-group-62} further objects occur: in the first, the homomorphism $\boldsymbol{\Theta}_{X,\mathscr{T}}$ of Definition~\ref{def:transport-mean}; in the second, the exact operators and the family~\eqref{it:primitives-p5} that appear in the displays of Definitions~\ref{def:contour-functional-functions} and~\ref{def:determinant-functions}.
\end{proposition}

\begin{proof}
Let $\mathbf{Y}'$ be a polymer. A term $(X,P)$ of~\eqref{eq:renormalisation-group-57}, in the first or in the second localisation, with $Y'(X,P)=\mathbf{Y}'$ is called \emph{attributed to} $\mathbf{Y}'$.

\emph{The first assertion of~(a).} Let $(X,P)$ be attributed to $\mathbf{Y}$, and let $p$ be a lattice site that lies in $X$ or in some cube of $P_0(X)\cup P$. By fact~\eqref{it:renormalisation-group-g0}, exactly one $M$-cube contains $p$. This $M$-cube has the site $p$ in common with $X\cup Y^{\mathrm{p}}(X,P)$; by the definition of $Y'(X,P)$ in Section~\ref{sec:localisation-map}, it is therefore one of the $M$-cubes of $Y'(X,P)=\mathbf{Y}$. Hence all lattice sites of $X$ and of the cubes of $P_0(X)\cup P$ lie in $\mathbf{Y}$, as asserted.

\emph{The terms that enter the objects.} By~\eqref{eq:renormalisation-group-58}, \eqref{eq:renormalisation-group-59} and~\eqref{eq:renormalisation-group-61}, the elements $E^{(i),x}(\mathbf{Y})$ and $E^{(\le5),x}(\mathbf{Y})$ are formed from the terms attributed to $\mathbf{Y}$; in a term $(X,P)$ of $E^{(\le5),x}(\mathbf{Y})$, the function $G^{(\le3)}(s'',X)=E^{(\le3),x}(X;\cdot)$ is formed from the terms attributed to $X$; and $\mathsf{K}(\mathbf{Y})$, $\mathsf{H}(\mathbf{Y})$, $E^{\#,x}_k(\mathbf{Y})$ are built out of elements $E^{(\le5),x}(X_i)$ with $X_i\subseteq\mathbf{Y}$. Applying the first assertion of~(a) at each of these stages, one finds that every term $(X,P)$ entering an object attached to $\mathbf{Y}$ is attributed to some polymer $\mathbf{Y}'\subseteq\mathbf{Y}$, and that $X$ and all cubes of $P_0(X)\cup P$ are contained in $\mathbf{Y}'$. The properties named in (b), (c) and (d) persist when a polymer $\mathbf{Y}'\subseteq\mathbf{Y}$ is replaced by $\mathbf{Y}$: a pair admissible for a polymer contained in $\mathbf{Y}'$ is admissible for a polymer contained in $\mathbf{Y}$; a site, a cube, a basis element or a polymer contained in $\mathbf{Y}'$ is contained in $\mathbf{Y}$; a bond meeting $\mathbf{Y}'$ meets $\mathbf{Y}$; and a site at sup-distance at most $\eta_k$ from $\mathbf{Y}'$ is at sup-distance at most $\eta_k$ from $\mathbf{Y}$.

\emph{The other assertions of~(a).} The members of a set $\{X_i\}$ or of a family $(X_1,\dots,X_n)$ with union $\mathbf{Y}$ are subsets of $\mathbf{Y}$. A polymer $X$ with $\bar{X}^L=LY$ is contained in $\bar{X}^L=LY$, by item~(5) of Section~\ref{sec:activity-before-reblocking}. The datum occurs in none of the displays~\eqref{eq:renormalisation-group-56}, \eqref{eq:renormalisation-group-57}, \eqref{eq:renormalisation-group-58} and~\eqref{eq:renormalisation-group-61}. In~\eqref{eq:renormalisation-group-59} and~\eqref{eq:renormalisation-group-60} it occurs only through the number $m'$, in $G^{(3)}$ and $G^{(5)}$, and through the elements $\mathscr{E}^r(X;\cdot)$, in $G^{(2)}$ and $G^{(4)}$, where $X$ is the polymer of the term; by the preceding paragraph, $X\subseteq\mathbf{Y}$ for every term that enters. Finally, by the definitions of Section~\ref{sec:local-parts-adjustment}, $(R\,\mathscr{E})(X)$ equals $\mathscr{E}(X)$ for large $X$, while $m(\mathscr{E})$ and, for small $X$, the difference $(R\,\mathscr{E})(X)-\mathscr{E}(X)$ are computed from $\mathscr{E}_{00}(X;0,0)$ and $\mathscr{E}_{20}(X;0,0;\cdot)$, with $X$ small.

\emph{Proof of~(b): the list.} Apart from the datum and the primitives, the displays contain the following operators and forms: $C$, in the definition of $\mathscr{Z}(s')$; $\delta_{\alpha,B}$ and $\mathsf{C}$, in $H^{\#}_k(s'',B)$ and in $\bar{H}_k(s'',B)$, see~\eqref{eq:renormalisation-group-56}; the operators $T^{(i)}$ and $\bar{T}^{(i)}$, together with the matrices $\mathsf{C}$ and $P^{\boldsymbol{\sigma}}_{\nu}$, in~\eqref{eq:renormalisation-group-58}; the matrices $P^{\boldsymbol{\sigma}}_{\nu}$, in~\eqref{eq:renormalisation-group-60}; and the forms $\langle\bar{\psi},\psi\rangle_{\square}$, in $G^{(3)}$ and $G^{(5)}$. The displays~\eqref{eq:renormalisation-group-57} and~\eqref{eq:renormalisation-group-61} contain no operator.

\emph{Proof of~(b): the places.} The operator $\delta_{\alpha,B}$ stands only in the second components of $H^{\#}_k(s'',B)$, hence only in the fields $J^{\#}_k\left(s'',\mathscr{T},\mathscr{Z}(s')\right)\Psi$ and $W^{\#}_k(s'',\mathscr{T})$. By the proof of Lemma~\ref{lem:renormalisation-group-1}(a), these fields are inserted only at sites of $X$ and at pairs that are admissible for $X$, where $X$ denotes the polymer of the term; in $G^{(3)}$ and $G^{(5)}$ only their components without H\"older derivative occur. Since $X\subseteq\mathbf{Y}$, the operator $\delta_{\alpha,B}$ enters only through pairs admissible for polymers contained in $\mathbf{Y}$; by~\eqref{eq:renormalisation-group-adm}, such pairs are admissible for $\mathbf{Y}$ as well. In $E^{(i),x}(\mathbf{Y})$ the substitution~\eqref{eq:renormalisation-group-58} is formed with $\mathbf{Y}$, and it acts only on the generators $\Psi(y)$, $\bar{\Psi}(y)$ whose unit site $y$ lies in $\mathbf{Y}$; inside $G^{(\le3)}(s'',X)=E^{(\le3),x}(X;\cdot)$ the polymer used for the substitution is the polymer $X$ of the term, and $X\subseteq\mathbf{Y}$. In both cases only the rows of $T^{(i)}(\mathscr{T})$ and $\bar{T}^{(i)}(\mathscr{T})$ at unit sites of $\mathbf{Y}$ occur. The functions $G^{(3)}(s,X)$ and $G^{(5)}(s'',X)$ vanish unless $X$ is a single $M$-cube $\square$ (Section~\ref{sec:localisation-exponent}), and then $\square=X\subseteq\mathbf{Y}$.

\emph{Proof of~(b): one formula.} The defining formulas of these operators and forms make sense on $\mathbb{L}_k$: Definition~\ref{def:pairs-sites-holder} for $\delta_{\alpha,B}$; the display~\eqref{eq:renormalisation-group-58}, with the notation of Section~\ref{sec:notation} and with~\eqref{eq:renormalisation-group-39}, for $T^{(i)}$ and $\bar{T}^{(i)}$; the definition of $\langle\bar{\psi},\psi\rangle_{\square}$ in Section~\ref{sec:notation}; and the definition of $C$ in Section~\ref{sec:averaging-fluctuations-gauge}. Read on $\mathbb{L}_k$, they define, for every bond function $\tilde{B}$ on $\mathbb{L}_k$ and every union $\tilde{Y}$ of $M$-cubes of $\mathbb{L}_k$, operators between spaces of functions on the index sets of Definition~\ref{def:terminology}(2), each index possibly carrying a spinor index. One of these index sets consists of the pairs $(x,x')$ of points of $\mathbb{L}_k$ with $0<|x-x'|\le M$, which are the shortest lifts of the pairs of $\mathrm{Pz}_k$: a translation by $v$ acts on both points of a pair, $(T_vF)(x,x')=F(x-v,x'-v)$, the support of $(x,x')$ is $\{x,x'\}$, and the spinor functions on $\mathrm{Pz}_k$ are the spinor functions $F$ on these pairs with $F(x+l\cdot e_{\nu},x'+l\cdot e_{\nu})=-F(x,x')$ for $\nu=0,1,2$ (Section~\ref{sec:notation}). The operator $\delta_{\alpha,\tilde{B}}$ maps the spinor functions on $\mathbb{L}_k$ to the spinor functions on these pairs; $T^{(1)}(\tilde{B})$ and $\bar{T}^{(1)}(\tilde{B})$ map the spinor functions on $\mathbb{L}_k$, and $T^{(2)}(\tilde{B})$ and $\bar{T}^{(2)}(\tilde{B})$ the spinor functions on these pairs, to the spinor functions on $\mathbb{Z}^3$; $C$ maps the functions on the basis elements to the functions on the positively oriented bonds of $\mathbb{Z}^3$; the matrices $P^{\boldsymbol{\sigma}}_{\nu}$ and $\mathsf{C}$ are operators that act at each site; and the form $\langle\bar{\psi},\psi\rangle_{\square}$ enters through the kernel $\eta_k^3\cdot\mathbf{1}_{\{x=x'\}}\cdot\mathbf{1}_{\tilde{Y}}(x)$ times the unit matrix, the cube $\square$ standing in the place of $\tilde{Y}$. The display~\eqref{eq:renormalisation-group-58} is read at the unit sites of $\tilde{Y}$, and the rows of $T^{(i)}(\tilde{B})$ and of $\bar{T}^{(i)}(\tilde{B})$ at the unit sites outside $\tilde{Y}$ are put equal to zero; in the objects attached to $\mathbf{Y}$ only rows at unit sites of the polymer with which these operators are formed occur. These operators do not depend on the side $l$, and their rows have finite support. A translation by a vector of $LM\mathbb{Z}^3$ maps each of the following families of cubes onto itself: the blocks of the sides $L^j\eta_k$, the unit cubes, the $L$-blocks and the $M$-cubes. Since the formulas involve only the ingredients listed in~(b), the translation relation of condition~(ii) in Definition~\ref{def:terminology}(2) holds for every $v\in LM\mathbb{Z}^3$, and on $T_k$ each operator is the restriction of the corresponding operator on $\mathbb{L}_k$ to the periodic, resp.\ antiperiodic, functions, formed with the periodic extensions of the bond function and of the union of $M$-cubes. In detail, let $v\in LM\mathbb{Z}^3$; for the bonds $b$ of $\mathbb{L}_k$ one has $(T_v\tilde{B})(b)=\tilde{B}(b-v)$. The six paths of a pair $(x,x')$ are the translates by $v$ of the six paths of $(x-v,x'-v)$; hence $(\tau T_v\tilde{B})(x,x')=(\tau\tilde{B})(x-v,x'-v)$, and Definition~\ref{def:pairs-sites-holder} gives condition~(ii) for $\delta_{\alpha,\tilde{B}}$. As $v\in\mathbb{Z}^3$, the translation by $v$ permutes the blocks of each of the sides $L^j\eta_k$, $0\le j\le k$, and carries the centre of each block to the centre of its image; hence $B_k(y)+v=B_k(y+v)$ and $\phi_{T_v\tilde{B},y}(x)=\phi_{\tilde{B},y-v}(x-v)$, and moreover $u_{T_v\tilde{B}}\left(x,x^{\boldsymbol{\sigma}}\right)=u_{\tilde{B}}\left(x-v,(x-v)^{\boldsymbol{\sigma}}\right)$. Since $x^{\boldsymbol{\sigma}}$ lies in $\tilde{Y}+v$ if and only if $(x-v)^{\boldsymbol{\sigma}}$ lies in $\tilde{Y}$, and $y$ lies in $\tilde{Y}+v$ if and only if $y-v$ lies in $\tilde{Y}$, these identities, inserted into~\eqref{eq:renormalisation-group-58}, give condition~(ii) for $T^{(i)}(\tilde{B})$ and $\bar{T}^{(i)}(\tilde{B})$. For $C$ condition~(ii) holds because the formula of $C$ in Section~\ref{sec:averaging-fluctuations-gauge} involves only bonds, $L$-blocks and their faces, which the translations by the vectors of $L\mathbb{Z}^3$ permute, and because these translations map onto itself the family of bond functions on $\mathbb{Z}^3$ that the covering map carries onto the basis $\{\Upsilon\}$ (Section~\ref{sec:notation}). For the matrices $P^{\boldsymbol{\sigma}}_{\nu}$ and $\mathsf{C}$ it is evident, and the operator of multiplication by $\eta_k^3\cdot\mathbf{1}_{\tilde{Y}}$ satisfies $T_v\circ\left(\eta_k^3\cdot\mathbf{1}_{\tilde{Y}}\right)=\left(\eta_k^3\cdot\mathbf{1}_{\tilde{Y}+v}\right)\circ T_v$.

The form $\langle\bar{\psi},\psi\rangle_{\square}$ is the sum, over the sites $x$, $x'$ of $T_k$ and over the spinor index $\boldsymbol{\beta}$, of $\bar{\psi}_{\boldsymbol{\beta}}(x)\cdot\eta_k^3\cdot\mathbf{1}_{\{x=x'\}}\cdot\mathbf{1}_{\square}(x)\cdot\psi_{\boldsymbol{\beta}}(x')$. Its kernel is, in the sense of condition~(iii), the restriction to $T_k$ of the operator of multiplication by $\eta_k^3\cdot\mathbf{1}_{\tilde{Y}}$ on $\mathbb{L}_k$, formed with the periodic extension $\tilde{Y}$ of $\square$; each row of this operator has at most one non-zero entry, so that condition~(i) holds, and condition~(ii) has just been shown. The sum that defines the form extends over the sites of $\square$ in $T_k$, which correspond one to one to the sites of a lift of $\square$ (a lift exists by Definition~\ref{def:terminology}(9), as $M<l$); the periodic extension $\tilde{Y}$ occurs only in the operator on $\mathbb{L}_k$ whose restriction to $T_k$ is the kernel.

\emph{Proof of~(b): the range.} By Proposition~\ref{prop:basis-map-box}, $C(b',\Upsilon)=0$ unless $\operatorname{supp}\Upsilon$ lies in the union of the one or two $L$-blocks that contain the end points of $b'$, a set whose sup-diameter is at most $2L-1$. The row of $T^{(i)}(B)$ at a unit site $y$ involves only the sites $x\in B_k(y)$, for which $|x-y|<\frac{1}{2}$, and their nearest neighbours $x^{\boldsymbol{\sigma}}$, for which $\left|x^{\boldsymbol{\sigma}}-x\right|=\eta_k\le1$. So the kernels of both operators vanish at pairs of indices at sup-distance $2L$ or more. Since $M>4L$ (stage~\ref{it:order-choice-6}), one has $2L<\frac{M}{2}$. Each point of the union of $L$-blocks just named lies in the $L$-block of an end point of $b'$, hence at sup-distance at most $L-1$ from that end point; so the entries of $C$ vanish if the supports of the row and of the column are at sup-distance $L$ or more.

Now let $\tilde{B}$ be a bond function on $\mathbb{L}_k$ and $\tilde{Y}$ a union of $M$-cubes of $\mathbb{L}_k$. The operator $C$ and the matrices $P^{\boldsymbol{\sigma}}_{\nu}$ and $\mathsf{C}$ contain neither $\tilde{B}$ nor $\tilde{Y}$, and the matrices act at each site. The row of $\delta_{\alpha,\tilde{B}}$ at a pair $(x,x')$ has its columns at $x$ and $x'$, which belong to the support $\{x,x'\}$ of the pair. It depends on $\tilde{B}$ only through the phase $\exp\left(i\,e_k\cdot\eta_k\cdot(\tau\tilde{B})(x,x')\right)$, that is, through the values of $\tilde{B}$ on the bonds of the six paths $\Gamma^{\pi}(x,x')$. Along a rectilinear path each coordinate varies monotonically between its values at $x$ and at $x'$; so both end points of each of these bonds lie within sup-distance $|x-x'|\le M$ of $x$. The set $\tilde{Y}$ does not occur in this row. Let $y$ be a unit site. If $y\notin\tilde{Y}$, the rows of $T^{(i)}(\tilde{B})$ and of $\bar{T}^{(i)}(\tilde{B})$ at $y$ vanish. If $y\in\tilde{Y}$, the columns of these rows are the sites $x\in B_k(y)$ and $x^{\boldsymbol{\sigma}}$, resp.\ the pairs $\left(x,x^{\boldsymbol{\sigma}}\right)$, all of them at sup-distance less than $\frac{1}{2}+\eta_k\le\frac{3}{2}$ from $y$. By the proof of Lemma~\ref{lem:interior-boundary-parts}(b), read on $\mathbb{L}_k$, the number $\phi_{\tilde{B},y}(x)$ involves only bonds with both end points in $B_k(y)$, and $u_{\tilde{B}}\left(x,x^{\boldsymbol{\sigma}}\right)$ involves only the bond $\left(x,x^{\boldsymbol{\sigma}}\right)$; so the end points of all bonds on which these rows read $\tilde{B}$ lie at sup-distance less than $\frac{3}{2}$ from $y$. The set $\tilde{Y}$ enters~\eqref{eq:renormalisation-group-58} only through the conditions $x^{\boldsymbol{\sigma}}\in\tilde{Y}$ and $x^{\boldsymbol{\sigma}}\notin\tilde{Y}$, and through the condition $y\in\tilde{Y}$, which decides whether the row at $y$ is put equal to zero. The barred operators $\bar{T}^{(i)}(\tilde{B})=\mathsf{C}^{-1}\cdot T^{(i)}(-\tilde{B})\cdot\mathsf{C}$ involve the same bonds and the same sites. Finally, the row at $x$ of the kernel of the form has no column other than $x$, contains no bond function, and reads $\tilde{Y}$ only at $x$. Since $\frac{3}{2}<2L<\frac{M}{2}$ and $L<M$, the number $M$ is a range of every one of these formulas, as asserted in~(b).

\emph{Proof of~(b): the connections.} Let $(x,x')$ be a pair admissible for a polymer $X\subseteq\mathbf{Y}$. The row of $\delta_{\alpha,B}$ at $(x,x')$ depends on $B$ only through the phase $\exp\left(i\,e_k\cdot\eta_k\cdot(\tau B)(x,x')\right)$, that is, through the values of $B$ on the bonds of the paths $\Gamma^{\pi}(x,x')$; the same bonds occur for $\delta_{\alpha,-B}$ on the barred fields. By~\eqref{eq:renormalisation-group-adm} these bonds are contained in $X$, or they reduce to the bond $(x,x')$, which meets $X$; in both cases they meet $\mathbf{Y}$. Let $y$ be a unit site of a polymer $X\subseteq\mathbf{Y}$. By Lemma~\ref{lem:interior-boundary-parts}(b), applied with the polymer $X$, the rows of $T^{(i)}(B)$ at $y$ read $B$ on the bonds with an end in $B_k(y)\subseteq X$; the rows of $\bar{T}^{(i)}(B)=\mathsf{C}^{-1}\cdot T^{(i)}(-B)\cdot\mathsf{C}$ read $B$ on the same bonds, since the change of sign of $B$ leaves the bonds unchanged. All these bonds meet $\mathbf{Y}$. The operator $C$, the matrices $P^{\boldsymbol{\sigma}}_{\nu}$ and $\mathsf{C}$, and the forms $\langle\bar{\psi},\psi\rangle_{\square}$ are defined without any connection.

\emph{Proof of~(c).} By the paragraph on the terms that enter the objects, it suffices, except for the assertion on the columns, to consider a term $(X,P)$ attributed to $\mathbf{Y}$. In it, $s$ vanishes outside $P_0(X)\cup P$, by~\eqref{eq:renormalisation-group-57}, and every cube of $P_0(X)\cup P$ is contained in $\mathbf{Y}$, by~(a); so $s$ vanishes off $\mathbf{Y}$, which proves the assertion on the parameter families. As $\mathbf{Y}$ is a polymer, hence connected, the component $Y'_1$ of $Y'(X,P)=\mathbf{Y}$ that contains $X$ is $\mathbf{Y}$ itself. Thus Lemma~\ref{lem:renormalisation-group-1} and the steps of its proof apply with $Y'_1=\mathbf{Y}$; at level $0$ the properties of the primitives on which they rest are contained in property~\ref{it:renormalisation-group-17}.

The field $\mathscr{Z}(s')$ enters the term only through $M_k\left(X,\mathscr{T}+t\cdot\mathscr{Z}(s')\right)$, through $\mathscr{E}^r\left(X;A+t\cdot\mathscr{Z}(s'),\mathscr{C},\cdot\right)$, through the rows of the family~\eqref{it:primitives-p3} at the connection $\mathscr{T}+t\cdot\mathscr{Z}(s')$, through the phases of $\delta_{\alpha,\mathscr{T}+t\cdot\mathscr{Z}(s')}$, and, at level $0$, through the exponentials in~\eqref{eq:renormalisation-group-60}. By Steps~(2), (3) and~(5) of the proof of Lemma~\ref{lem:renormalisation-group-1}, all of these read $\mathscr{Z}(s')$ on bonds that meet $\mathbf{Y}$. Let $b$ be a fine bond meeting $\mathbf{Y}$, and let
\[
H_k(s')(b,b')\cdot C(b',\Upsilon)\cdot C_k^{\frac{1}{2}}(s')(\Upsilon,\Upsilon')\cdot\mathsf{Z}(\Upsilon')
\]
be a non-vanishing term of the sum that defines $\mathscr{Z}(s')(b)$. By Step~(1) of that proof, $b'$ is contained in an attached link-component $\mathcal{L}_1$, whose union lies in $\mathbf{Y}$. By Proposition~\ref{prop:basis-map-box}, $\operatorname{supp}\Upsilon$ lies in the union of the $L$-blocks of the end points of $b'$. Each of these $L$-blocks lies in a single $M$-cube, by fact~\eqref{it:renormalisation-group-g0}; this $M$-cube contains an end point of $b'$, hence it is an $M$-cube of $\mathbf{Y}$. So $\Upsilon$ is contained in $\mathbf{Y}$. Consequently $H_k(s')$ enters only through its rows indexed by fine bonds meeting $\mathbf{Y}$, and $C_k^{\frac{1}{2}}(s')$ only through its rows indexed by basis elements contained in $\mathbf{Y}$.

The rows of $H_k(s'',\cdot)$ and of $\bar{H}_k(s'',\cdot)$ that occur are taken at sites of $X$ and at end points of pairs admissible for $X$ (proof of Lemma~\ref{lem:renormalisation-group-1}(a), with its Step~(3)); by~\eqref{eq:renormalisation-group-adm}, these points have sup-distance at most $\eta_k$ from $X\subseteq\mathbf{Y}$. By~\eqref{eq:renormalisation-group-56}, the rows of $\Gamma_k(s'',\cdot)$ that occur are taken at the columns of the former rows; by Step~(2) of that proof these columns are unit sites of an attached link-component $\mathcal{L}_3$, which is contained in $\mathbf{Y}$. Since $\eta_k\le1<M$, all row indices of the list in~(c) belong to $\mathbf{Y}^{\sim1}$. By Step~(5) of that proof, that is, by property~\ref{it:renormalisation-group-12}, the kernel of $M_k(X,B)$ vanishes unless both indices lie in $X\subseteq\mathbf{Y}$, and $M_k(X,B)$ reads $B$ on the bonds that meet $X$; so $M_k$ enters only through the operators $M_k(X,\cdot)$ with $X\subseteq\mathbf{Y}$.

By~\eqref{eq:renormalisation-group-56} and~\eqref{eq:renormalisation-group-59}, the family~\eqref{it:primitives-p3} is evaluated at $\mathscr{T}+\mathscr{Z}(s')$ and at $\mathscr{T}$, in $J^{\#}_k$ and $W^{\#}_k$; the family $\Gamma_k$ at $\mathscr{T}$, in $W^{\#}_k$; and $M_k(X,\cdot)$ at $\mathscr{T}+t\cdot\mathscr{Z}(s')$ with $|t|=\rho_1$, in $G^{(1)}$. By Steps~(2), (4) and~(5) of the proof of Lemma~\ref{lem:renormalisation-group-1}, the rows of the list and the operators $M_k(X,\cdot)$ read these connections on the bonds that meet $\mathcal{L}_3$, resp.\ $X$, and these bonds meet $\mathbf{Y}$.

Now let two systems of primitives with the properties~\ref{it:renormalisation-group-1}--\ref{it:renormalisation-group-17} for $\mathrm{par}$ be given, whose rows in the list of~(c) coincide for every connection in $\tilde{G}_k$ and every parameter family that vanishes off $\mathbf{Y}$, and whose operators $M_k(X,\cdot)$ coincide for $X\subseteq\mathbf{Y}$. By what has been shown, applied to all the terms that enter, and since $C$ is common to both systems, the two fields $\mathscr{Z}(s')$ formed with them agree on the bonds that meet $\mathbf{Y}$. Hence the two connections $\mathscr{T}+t\cdot\mathscr{Z}(s')$, $|t|\le\rho_1$, agree on these bonds, and both lie in $\tilde{G}_k$, by part~\eqref{it:bounds-analyticity-q1} of Lemma~\ref{lem:bounds-analyticity}. By property~\ref{it:renormalisation-group-8} of the first system, used as in Step~(2) of the proof of Lemma~\ref{lem:renormalisation-group-1}, its rows of the family~\eqref{it:primitives-p3} in the list do not change when its connection is replaced by the connection of the second system; by hypothesis, they are then equal to the corresponding rows of the second system. The rows of $\Gamma_k$ are taken at the common connection $\mathscr{T}$ and coincide by hypothesis. In the same way, $M_k(X,\cdot)$, by property~\ref{it:renormalisation-group-12}, the elements $\mathscr{E}^r(X;\cdot)$, by property~\eqref{it:renormalisation-group-1-k2}, the phases of $\delta_{\alpha,B}$ and, at level $0$, the exponentials in~\eqref{eq:renormalisation-group-60} take the same values for the two systems. Therefore the objects attached to $\mathbf{Y}$ are the same for the two systems; in the sense of Definition~\ref{def:terminology}(1), this is the assertion of~(c) on the families of primitives.

\emph{Proof of~(c): the columns.} Let $s\in\mathrm{Pd}_k$ vanish off $\mathbf{Y}$, and let $B\in\tilde{G}_k$. Let $\mathrm{Op}$ be one of the families~\eqref{it:primitives-p1}--\eqref{it:primitives-p3} and $\Gamma_k$, let $\mathscr{Y}$ be the set of those parameter cubes of the kind of $\mathrm{Op}$ (Section~\ref{sec:properties-systems-primitives}) that are contained in $\mathbf{Y}$, and let $Y$ be the union of the cubes of $\mathscr{Y}$. A parameter cube of this kind that does not belong to $\mathscr{Y}$ is not contained in $\mathbf{Y}$, so that $s$ vanishes at it; hence $\mathscr{Y}$, $Y$ and $s$ are as in property~\ref{it:renormalisation-group-8}. Every link-component of $Y$ is a union of cubes of $\mathscr{Y}$ and is therefore contained in $\mathbf{Y}$. Let $x$ be the index of a row of $\mathrm{Op}$ that belongs to the list in~(c). If $x$ is contained in a link-component $Y_1$ of $Y$, property~\ref{it:renormalisation-group-8} gives $\mathrm{Op}(s,B)(x,x')=0$ unless $x'$ is contained in $Y_1$; so the non-zero entries of the row stand in columns contained in $\mathbf{Y}$. For $\Gamma_k$ this is always the case: the kind of $\Gamma_k$ is $\pi_k$, so that $\mathscr{Y}$ consists of the $M$-cubes of $\mathbf{Y}$, and the unit site $x$ of $\mathbf{Y}$ lies in one of them (fact~\eqref{it:renormalisation-group-g0}), hence in the link-component of $Y$ that contains this $M$-cube. If $\mathrm{Op}$ is one of the families~\eqref{it:primitives-p1}--\eqref{it:primitives-p3} and $x$ is contained in no link-component of $Y$, the row of $\mathrm{Op}(s,B)$ at $x$ vanishes, by property~\ref{it:renormalisation-group-10}. The row of $\bar{H}_k(s'',B)=\mathsf{C}^{-1}\cdot H_k(s'',-B)\cdot\mathsf{C}$ at a site $z$ arises from the row of $H_k(s'',-B)$ at $z$ through multiplication by the invertible constant matrix $\mathsf{C}$ and its inverse, which act on the spinor indices only; so its non-zero entries stand in the same columns, and the preceding argument applies to $H_k(s'',-B)$, since $-B\in\tilde{G}_k$ (Lemma~\ref{lem:symmetries-domains}). At level $0$ the properties~\ref{it:renormalisation-group-8} and~\ref{it:renormalisation-group-10} of the families used here are contained in property~\ref{it:renormalisation-group-17}. This proves the assertion of~(c) on the columns.

\emph{Proof of~(d).} The integrations in the displays are the contour integrals $\oint_{\rho_i}$ of~\eqref{eq:renormalisation-group-59}, over $t$ on the circles $|t|=\rho_i$; the integrals over $s_P\in[0,1]^P$ in~\eqref{eq:renormalisation-group-57}; and the integral with respect to $\mathrm{d}\mu_I(W)\,\mathrm{d}\hat{\mu}(\mathsf{Z})$ that defines $\mathsf{H}$ in~\eqref{eq:renormalisation-group-61}. In $\mathsf{H}(\mathbf{Y})$, and in the factors $\mathsf{H}(X_i)$, $X_i\subseteq\mathbf{Y}$, of $E^{\#,x}_k(\mathbf{Y})$, the integrand is formed from elements $E^{(\le5),x}(X')$ with $X'\subseteq\mathbf{Y}$. By Corollary~\ref{cor:renormalisation-group-1}, these contain a generator $W(u)$ only if $u$ is in $\mathbf{Y}$, and they depend on $\mathsf{Z}$ only through the coordinates $\mathsf{Z}(\Upsilon)$ with $\Upsilon$ contained in $\mathbf{Y}$. The integral $\mu_I$ factorises over the unit sites, by~\eqref{eq:renormalisation-group-43}, and its value at the constant $1$ is $1$ (Section~\ref{sec:recalled-auxiliary-results}). By Proposition~\ref{prop:basis-map-box} the basis $\{\Upsilon\}$ is orthonormal, so that in the coordinates $\mathsf{Z}(\Upsilon)=\langle\mathsf{Z},\Upsilon\rangle$ the standard Gaussian measure on $V_k$ is the product of standard normal laws, and the set $\mathrm{Box}_k$ is the product of the intervals $[-p_{0,k},p_{0,k}]$; as $\hat{\mu}$ is, by its definition in Section~\ref{sec:averaging-fluctuations-gauge}, the restriction of that measure to $\mathrm{Box}_k$, normalised to mass one, the measure $\hat{\mu}$ is the product, over the basis $\{\Upsilon\}$, of copies of the standard normal law restricted to $[-p_{0,k},p_{0,k}]$ and normalised to mass one. So the integration over the other generators and coordinates contributes the factor $1$, and the integrations are those described in~(d).

\emph{The pair of connections.} Let $(A,\mathscr{C})$ and $(A^1,\mathscr{C}^1)$ be as in the paragraph that follows the list. By Corollary~\ref{cor:renormalisation-group-1}, read in the sense fixed before Lemma~\ref{lem:renormalisation-group-1}, the coefficients of $E^{(i),x}(\mathbf{Y})$, $i\le3$, and of $E^{(\le5),x}(\mathbf{Y})$ take the same values at $(A,\mathscr{C},\mathsf{Z})$ and at $(A^1,\mathscr{C}^1,\mathsf{Z})$, for every $\mathsf{Z}\in\mathrm{Box}_k$: these two points agree in the coordinates of $\mathsf{Z}$ and in the values of the connections on the bonds that meet $\mathbf{Y}$. By the first line of~\eqref{eq:renormalisation-group-61}, $\mathsf{K}(\mathbf{Y})$ is a finite sum, over the sets $\{X_i\}$ of distinct polymers with union $\mathbf{Y}$, of products of the elements $\exp\left(E^{(\le5),x}(X_i)\right)-1$. Each $X_i$ is contained in $\mathbf{Y}$, and a bond that meets $X_i$ meets $\mathbf{Y}$; so the two pairs coincide on the bonds that meet $X_i$, and what has just been shown, applied with the polymer $X_i$ in place of $\mathbf{Y}$, gives the equality of every factor at every $\mathsf{Z}\in\mathrm{Box}_k$. Hence $\mathsf{K}(\mathbf{Y})$ takes the same value at the two pairs, at every $\mathsf{Z}\in\mathrm{Box}_k$. The element $\mathsf{H}(\mathbf{Y})$ is the integral of $\mathsf{K}(\mathbf{Y})$ with respect to $\mathrm{d}\mu_I(W)\,\mathrm{d}\hat{\mu}(\mathsf{Z})$, and these measures do not depend on the connections (Remark~\ref{rem:renormalisation-group-1}(i)); so $\mathsf{H}(\mathbf{Y})$ takes the same value at the two pairs. For $E^{\#,x}_k(\mathbf{Y})$ the assertion is assertion~(2) of Section~\ref{sec:activity-before-reblocking}. This proves the assertion on the pair of connections.
\end{proof}

Clause~\eqref{it:step-small-fields-g} of Theorem~\ref{thm:step-small-fields} follows from Section~\ref{sec:bound-analyticity-after} and Proposition~\ref{prop:local-constituents-step}.

\subsubsection{Uniformity}\label{sec:uniformity}

\begin{proof}[Proof of Theorem~\ref{thm:renormalisation-step-small}(h)]
Recall from Section~\ref{sec:hypotheses-conditions} that $v_k=-\log e_k$ and $p_{0,k}=v_k^{p_0}$; by stages~\ref{it:order-choice-1}, \ref{it:order-choice-2} and~\ref{it:order-choice-8}, one has $0<1-\alpha<\varepsilon\le\frac{1}{200}$. Consider the inequalities that make up the conditions~\eqref{eq:renormalisation-group-31}--\eqref{eq:renormalisation-group-38}. Each of them, the first two inequalities in~\eqref{eq:renormalisation-group-31} and the second inequality in~\eqref{eq:renormalisation-group-38} excepted, compares a positive power of $e_k$ with a number formed from the constants
\[
C_w,\quad c_{\gamma},\quad c_C,\quad C_T,\quad c_{\mathrm{fl}},\quad c_X,\quad K_0,\quad \kappa,\quad L,\quad M;
\]
the exponents of $e_k$ are positive because $0<6\varepsilon<\frac{1}{4}$. The second inequality in~\eqref{eq:renormalisation-group-38} is equivalent to $e_k^{\frac{1}{4}}+e_k^{\varepsilon}\le 1$ (Section~\ref{sec:hypotheses-conditions}), which compares the number $1$ with a sum of two positive powers of $e_k$. The first inequality in~\eqref{eq:renormalisation-group-31} compares the power $v_k^{p_0}$ of $v_k$ with the power $e_k^{-\frac{\varepsilon}{4}}$ of $e_k$. Since $e_k=\exp\left(-v_k\right)$, the second inequality in~\eqref{eq:renormalisation-group-31} is equivalent to
\[
\frac{1}{2}\,v_k^{2p_0}\ge 7v_k+\log 12,
\]
and here $2p_0>1$, because $p_0>3r+2\ge 8$ by stage~\ref{it:order-choice-8}. As $e_k\to 0$ one has $v_k\to+\infty$; every positive power of $e_k$ and the number
\[
v_k^{p_0}\cdot e_k^{\frac{\varepsilon}{4}}=v_k^{p_0}\cdot\exp\left(-\frac{\varepsilon\cdot v_k}{4}\right)
\]
tend to $0$, while $\frac{1}{2}\,v_k^{2p_0}-7v_k-\log 12$ tends to $+\infty$. Each of the inequalities therefore holds for all sufficiently small $e_k>0$. It follows that~\eqref{eq:renormalisation-group-31}--\eqref{eq:renormalisation-group-38} hold for every $e_k\in(0,e_0]$ as soon as $e_0>0$ is sufficiently small in terms of $(\alpha,b,\kappa,L,M,\mathrm{par},\varepsilon,p_0)$. Stage~\ref{it:order-choice-9} requires $e_0$ to be this small.

As far as the constants are concerned, the proofs of the clauses (a)--(g) of Theorem~\ref{thm:renormalisation-step-small} rest on hypothesis~\ref{def:hypothesis-one-step} alone; for the propositions recalled in Section~\ref{sec:recalled-auxiliary-results} this is stated there. Since hypothesis~\ref{def:hypothesis-one-step} requires the constants to be chosen as in Section~\ref{sec:order-choice}, subject to conditions of Section~\ref{sec:conditions}, every condition on the constants on which these clauses rest is a condition of Section~\ref{sec:order-choice} or of Section~\ref{sec:conditions}, or a consequence of such conditions. Besides the requirements on the data $k$, $e_k$, $l$ and $\mathfrak{m}$ stated in Section~\ref{sec:hypotheses-conditions}, the arguments of Sections~\ref{sec:elementary-properties-field}--\ref{sec:data-level-locality} invoke conditions on the constants, among them the following: the conditions~\eqref{eq:renormalisation-group-30}--\eqref{eq:renormalisation-group-38}; the inequality $\kappa\ge 7\kappa_0+7$ of stage~\ref{it:order-choice-3}; the facts that $L$ is odd, that $L\ge 4^8$, so that $L\ge 5$, and that $L^{-\frac{3}{4}}\le\frac{1}{2}$, which stage~\ref{it:order-choice-4} provides; the relation $M=L^m$ of stage~\ref{it:order-choice-6}; the chain of inequalities $0<1-\alpha<\varepsilon\le\frac{1}{200}$ obtained above, of which these arguments use $0<\varepsilon\le\frac{1}{200}$ and $\alpha>0$, the latter in the proof of Lemma~\ref{lem:scaling-domains}(a) (the exponent $-\frac{15}{8}-\alpha+\frac{3\varepsilon}{2}$ of $L$ is at most $-\frac{3}{4}$), in the proof of Lemma~\ref{lem:phases-differences-along} ($|x-x'|^{-\alpha}\le 1$ if $|x-x'|\ge 1$) and in Section~\ref{sec:bound-analyticity-after} ($L^{-1-\alpha}\le L^{-1}$); the inequality
\[
C_w\ge 4\exp\left(2\gamma_w\right)\cdot M^{27\alpha_0}
\]
of Section~\ref{sec:notation}, which follows from the definition of $C_w^1$ in Section~\ref{sec:photon-operators}; the inequality $c_{\gamma}\ge 1$, which holds by the definition of $c_{\gamma}$ in stage~\ref{it:order-choice-7} (Section~\ref{sec:notation}) and is used in the proof of Lemma~\ref{lem:bounds-analyticity}\eqref{it:bounds-analyticity-q4}; and the relations $\kappa_0\ge c_X\ge 1$ and $K_0\ge 1$, which hold for the values $c_X=54$, $\kappa_0=54\log 72$ and $K_0=\frac{1}{36}\exp\left(\kappa_0\right)$ fixed in stage~\ref{it:order-choice-1}, because $\log 72>1$ and $\kappa_0>\log 36$, and which enter, for instance, the proofs of Lemma~\ref{lem:resummation-polymers}, Proposition~\ref{prop:bounds-regularity-localised} and Lemma~\ref{lem:integration-polymer-exponential}. The arguments of Sections~\ref{sec:elementary-properties-field}--\ref{sec:data-level-locality} state further relations among the constants at the places where they use them, for instance $L\ge 4$ in Section~\ref{sec:bound-analyticity-after}. The propositions recalled in Section~\ref{sec:recalled-auxiliary-results} and at the beginning of Section~\ref{sec:proof-clause-theorem} are stated under the following hypotheses: Propositions~\ref{prop:symmetries-extraction}, \ref{prop:adjustment} and~\ref{prop:contraction} under hypothesis~\ref{def:hypothesis-one-step} without primitives (Theorem~\ref{thm:renormalisation-group-1}(a)), and likewise Propositions~\ref{prop:comparison-fermion-operators} and~\ref{prop:step-external-connection} (Theorem~\ref{thm:renormalisation-group-2}(b)); Propositions~\ref{prop:left-inverse}, \ref{prop:symmetries-fermion-operators} and~\ref{prop:local-expansion-determinant} under the hypotheses of Theorem~\ref{thm:fermion-operators-weakened}; and Propositions~\ref{prop:basis-map-box}, \ref{prop:symmetries-reality} and~\ref{prop:fields-domains} under those of Theorem~\ref{thm:photon-operators-weakened}. Their statements and proofs use, beyond the conditions just named, the following conditions on the constants. First, the conditions~\eqref{eq:renormalisation-group-A9} and~\eqref{eq:renormalisation-group-A10} on $e_0$; they contain the conditions~\eqref{eq:renormalisation-group-20}--\eqref{eq:renormalisation-group-25} named in hypothesis~\ref{def:hypothesis-one-step}, because the inequality $e_0\le e^F_0$ of~\eqref{eq:renormalisation-group-A10} gives~\eqref{eq:renormalisation-group-20}--\eqref{eq:renormalisation-group-23} by the definition of $e^F_0$ in Section~\ref{sec:conditions}, because condition~\eqref{eq:renormalisation-group-24}, in which $C_d$ is formed with the constant $C_w$ of $\mathrm{par}$, is itself one of the inequalities of~\eqref{eq:renormalisation-group-A10}, namely the condition on $e_0$ on which clause~\eqref{it:local-expansion-determinant-d6} of Proposition~\ref{prop:local-expansion-determinant} rests, as stated in Section~\ref{sec:conditions} and after Proposition~\ref{prop:local-expansion-determinant}, and because~\eqref{eq:renormalisation-group-25} is contained in~\eqref{eq:renormalisation-group-A9} (Section~\ref{sec:conditions}). Second, the inequalities $1-\varepsilon<\alpha<1$, which are part of the chain above: they are a hypothesis of Proposition~\ref{prop:contraction}, and the proof of Theorem~\ref{thm:contraction}(d) uses $\alpha>1-\varepsilon$ for the bidegrees $(1,1)$ and $(0,2)$. Third, the first inequality $L^{\frac{1}{8}-\frac{3\bar{\varepsilon}}{2}}\ge 2400$ of stage~\ref{it:order-choice-4} itself. Step~5 of the proof of Theorem~\ref{thm:polymer-activities-3}, which contains Proposition~\ref{prop:connections-near-small}, requires $1382\,L^{-\frac{15}{8}+\frac{3\varepsilon}{2}}<L^{-\frac{7}{4}}$; Theorem~\ref{thm:polymer-activities-3}(ii) rests on this inequality, Theorem~\ref{thm:contraction}(d) and~(b), that is, Proposition~\ref{prop:contraction}, rest on Theorem~\ref{thm:polymer-activities-3}(ii), and clause~(d) of Theorem~\ref{thm:renormalisation-step-small} rests on Proposition~\ref{prop:contraction}. The facts on $L$ named above do not imply this inequality: the odd number $L=4^8+1$ satisfies them, and $L^{\frac{1}{8}-\frac{3\varepsilon}{2}}<L^{\frac{1}{8}}<5$ for it. Fourth, the inequality $M\ge M^F\left(b,\alpha,L,\alpha_0,\kappa_w+\kappa_0\right)$ of stage~\ref{it:order-choice-6}, that is, condition~\eqref{eq:renormalisation-group-19} for $\kappa'=\kappa_w+\kappa_0$, which is a hypothesis of Theorem~\ref{thm:fermion-operators-weakened}. It enters through the inequalities $M>4L$ and $M\ge 108\pi\cdot C_V\cdot a$, which follow from it because $108\pi\cdot C_V\cdot a\ge 1$, and which are the hypotheses on $M$ of Theorem~\ref{thm:expansion-walks}. On this theorem rest the existence of the inverses in the definitions of the exact fermion operators (Section~\ref{sec:exact-fermion-operators}), the identity $T_k(A)\cdot H_k(A)=1$ of Proposition~\ref{prop:left-inverse} (Proposition~\ref{prop:identities-bounds}(a); of the conditions on the constants, the bound by $C_T$ in Proposition~\ref{prop:identities-bounds}(d) requires only~\eqref{eq:renormalisation-group-23}), Proposition~\ref{prop:local-expansion-determinant} (Proposition~\ref{prop:polymer-expansion-determinant}), and Propositions~\ref{prop:comparison-fermion-operators} and~\ref{prop:step-external-connection} (Theorems~\ref{thm:fermion-operators-two} and~\ref{thm:one-step-formula}). Of the conditions on $M$, the requirements~\eqref{eq:renormalisation-group-30} contain only $M>4L$ and $M^{\alpha_0}\ge 1+\exp(\hat{\kappa})$. Fifth, the inequality $M\ge M^P\left(\alpha,L,\alpha_0\right)$ of stage~\ref{it:order-choice-6}, $M^P$ being the number~\eqref{eq:operators-photon-20}: it is a hypothesis of Theorem~\ref{thm:photon-operators-weakened}, which is proved as Theorem~\ref{thm:photon-operators}, and under it Propositions~\ref{prop:basis-map-box}, \ref{prop:symmetries-reality} and~\ref{prop:fields-domains} are asserted and the number $C^P$, which enters $C_w$, is defined. The inequality $M\ge L^2$ of stage~\ref{it:order-choice-6}, cited in the proof of Proposition~\ref{prop:properties-actions} and at the beginning of Section~\ref{sec:operators-photon-field}, follows from $M=L^m$ and $M>4L$. Sixth, the inequality $r\ge 2$ of stage~\ref{it:order-choice-8}: in the proof of Proposition~\ref{prop:polymer-expansion-determinant}\eqref{it:polymer-expansion-determinant-e}, which is Proposition~\ref{prop:local-expansion-determinant}\eqref{it:local-expansion-determinant-d4}, it gives $r_k\ge v_k^r\ge v_k^2$ and hence $M\cdot\left(\frac{v_k}{4}+3\right)<M\cdot v_k^2\le M\cdot r_k\le l$; the clauses~\eqref{it:local-expansion-determinant-d5} and~\eqref{it:local-expansion-determinant-d6} of Proposition~\ref{prop:local-expansion-determinant} rest on clause~\eqref{it:local-expansion-determinant-d4}. Seventh, relations that hold by the definitions of stages~\ref{it:order-choice-5} and~\ref{it:order-choice-7}: the inequalities $C_w\ge C_P$ and $\gamma_w\le\gamma_P$ are hypotheses of Proposition~\ref{prop:backgrounds-domains}, from which Proposition~\ref{prop:fields-domains} follows, and they hold because $C_w\ge C^P\ge C_P$ and $\gamma_w\le\gamma^P\le\gamma_P$ (Remark~\ref{rem:operators-photon-3}(ii)); likewise the inequalities $\alpha_0\le\alpha^F$ and $\alpha_0\le\alpha^P$, which are hypotheses of Theorems~\ref{thm:fermion-operators-weakened} and~\ref{thm:photon-operators-weakened}, hold because $\alpha_0=\min\left(\alpha^F,\alpha^P\right)$ (Section~\ref{sec:photon-operators}). Every condition on the constants named in this paragraph, other than the conditions on $e_0$, is a condition of Section~\ref{sec:order-choice} or of Section~\ref{sec:conditions}, or a consequence of such conditions, of the values fixed in stage~\ref{it:order-choice-1} and of the definitions by which stages~\ref{it:order-choice-5} and~\ref{it:order-choice-7} fix $\alpha_0$, $\gamma_w$, $C_w$ and $c_{\gamma}$; it is a relation among the constants of stages~\ref{it:order-choice-1}--\ref{it:order-choice-8} alone, and therefore it depends on none of $k$, $l$, $\mathfrak{m}$, the harmonic part of $A'$, $\mathscr{C}'$, the system of primitives and the data $(m,\mathscr{E})$.

Among the conditions named in the preceding paragraph, \eqref{eq:renormalisation-group-31}--\eqref{eq:renormalisation-group-38}, \eqref{eq:renormalisation-group-A9} and~\eqref{eq:renormalisation-group-A10}, together with the conditions~\eqref{eq:renormalisation-group-20}--\eqref{eq:renormalisation-group-25} contained in the last two, are conditions on $e_0$, imposed in stage~\ref{it:order-choice-9}; the others belong to stages~\ref{it:order-choice-1}--\ref{it:order-choice-8}. Of the latter, $\alpha\in\left(1-\bar{\varepsilon},1\right)$, whence $0<\alpha<1$, belongs to stage~\ref{it:order-choice-2}, with $\bar{\varepsilon}=\frac{1}{200}$ from stage~\ref{it:order-choice-1}; the first inequality of stage~\ref{it:order-choice-4} belongs to that stage; $\alpha_0\le\alpha^F$, $\alpha_0\le\alpha^P$ and $\gamma_w\le\gamma_P$ are relations among numbers fixed in stage~\ref{it:order-choice-5}; $M\ge M^F$, $M\ge M^P$ and $M\ge L^2$ belong to stage~\ref{it:order-choice-6}; in $C_w\ge C_P$ the constant $C_w$ is fixed in stage~\ref{it:order-choice-7} and $C_P$ in stage~\ref{it:order-choice-5}; and $r\ge 2$ and $1-\alpha<\varepsilon\le\bar{\varepsilon}$ belong to stage~\ref{it:order-choice-8}. Of the conditions named there that belong to stages~\ref{it:order-choice-1}--\ref{it:order-choice-8}, three contain one of the constants $C_w$, $c_{\gamma}$, $c_C$ and $C_T$, namely the inequalities $C_w\ge 4\exp\left(2\gamma_w\right)\cdot M^{27\alpha_0}$, $C_w\ge C_P$ and $c_{\gamma}\ge 1$; each of the first two contains $C_w$ and none of the other three, the third contains $c_{\gamma}$ and no other constant, and none of these four constants occurs in the requirements on the data. No condition imposed in stages~\ref{it:order-choice-1}--\ref{it:order-choice-8} of Definition~\ref{def:order-choice}, whether the proofs invoke it or not, contains one of the four constants: stage~\ref{it:order-choice-1} fixes absolute numbers; stage~\ref{it:order-choice-2} contains, besides $\alpha$, only $\bar{\varepsilon}$; stage~\ref{it:order-choice-3} contains, besides $\kappa$, only $\kappa_0$; stages~\ref{it:order-choice-4} and~\ref{it:order-choice-6} contain none of them, as shown below; stage~\ref{it:order-choice-8} contains, besides the numbers chosen or defined in it, only $\alpha$, $\bar{\varepsilon}$, $a_L$ and $N^{*}$; and stages~\ref{it:order-choice-5} and~\ref{it:order-choice-7} consist of definitions. Among the conditions named in the preceding paragraph, the constants $c_C$ and $C_T$ occur only in conditions on $e_0$, and $c_{\gamma}$ occurs, apart from conditions on $e_0$, only in the inequality $c_{\gamma}\ge 1$; conditions on $e_0$ that contain one of these three constants are, for instance, \eqref{eq:renormalisation-group-32}, \eqref{eq:renormalisation-group-33}, \eqref{eq:renormalisation-group-34} and~\eqref{eq:renormalisation-group-36}. The inequalities $C_w\ge 4\exp\left(2\gamma_w\right)\cdot M^{27\alpha_0}$ and $C_w\ge C_P$ are lower bounds for $C_w$ which hold by the definition of $C_w$ in stage~\ref{it:order-choice-7}: one has $C_w=64\,K_0\cdot C_w^1$ with $K_0\ge 1$ (Sections~\ref{sec:primitive-operators} and~\ref{sec:absolute-constants}), $C_w^1=\max\left(C_w^0,\,4\exp\left(2\gamma_w\right)\cdot M^{27\alpha_0}\right)$ with $C_w^0\ge C^P$ (Section~\ref{sec:photon-operators}), and $C^P=116\,C_P\cdot M^{\alpha_0\cdot n_P}\ge C_P$ by~\eqref{eq:operators-photon-20}, since $M\ge 1$. The constant $C_P$ is fixed in stage~\ref{it:order-choice-5}, before $M$, and $C_w$ in stage~\ref{it:order-choice-7}, after $M$; hence neither of these two inequalities imposes a restriction on $\kappa$, $L$ or $M$. The inequality $c_{\gamma}\ge 1$ holds by the definition of $c_{\gamma}$ in stage~\ref{it:order-choice-7}, whatever the values of the constants chosen before: for an index $x$ of the set over which $x'$ ranges, the sum in that definition contains the term with $x'=x$, which equals $1$, and its other terms are positive (Section~\ref{sec:notation}). Hence it imposes no restriction on $\kappa$, $L$ or $M$ either. In particular, $C_w$, $C_T$, $\varepsilon$ and $e_0$ are absent from every restriction that the proofs place on $\kappa$, $L$ and $M$. Indeed, by the first two sentences of the preceding paragraph the clauses (a)--(g) rest on hypothesis~\ref{def:hypothesis-one-step} alone, and in this hypothesis the conditions on $\kappa$, $L$ and $M$ are those of stages~\ref{it:order-choice-3}, \ref{it:order-choice-4} and~\ref{it:order-choice-6}: by Section~\ref{sec:conditions}, the requirements~\eqref{eq:renormalisation-group-A0} are contained in stages~\ref{it:order-choice-4}--\ref{it:order-choice-6}, condition~\eqref{eq:renormalisation-group-19} for $\kappa'=\kappa_w+\kappa_0$ belongs to stage~\ref{it:order-choice-6}, the conditions~\eqref{eq:renormalisation-group-A1}--\eqref{eq:renormalisation-group-A10} and~\eqref{eq:renormalisation-group-20}--\eqref{eq:renormalisation-group-25} are conditions on $e_0$, and none of the remaining conditions adds a condition on $\kappa$, $L$ or $M$ to these three stages. Of these restrictions the proofs invoke the inequality $\kappa\ge 7\kappa_0+7$; the facts that $L$ is odd and satisfies the first inequality $L^{\frac{1}{8}-\frac{3\bar{\varepsilon}}{2}}\ge 2400$ of stage~\ref{it:order-choice-4}, with its consequences, among them $L\ge 4^8$ and $L^{-\frac{3}{4}}\le\frac{1}{2}$, and the third inequality $L\cdot\left(\kappa-6\kappa_0-4\right)-1\ge\kappa$ of that stage, which stands in~\eqref{eq:renormalisation-group-30}; the relation $M=L^m$; the inequalities $M>4L$ and $M^{\alpha_0}\ge 1+\exp(\hat{\kappa})$, which stand in~\eqref{eq:renormalisation-group-30}, and the consequence $M\ge L^2$ of $M=L^m$ and $M>4L$; and the inequalities $M\ge M^F\left(b,\alpha,L,\alpha_0,\kappa_w+\kappa_0\right)$ and $M\ge M^P\left(\alpha,L,\alpha_0\right)$. No condition of stages~\ref{it:order-choice-3}, \ref{it:order-choice-4} and~\ref{it:order-choice-6}, whether the proofs invoke it or not, contains $C_w$, $C_T$, $\varepsilon$ or $e_0$. Stage~\ref{it:order-choice-3} contains, besides $\kappa$, only $\kappa_0$. The four displayed inequalities of stage~\ref{it:order-choice-4} contain, besides $L$, only $\bar{\varepsilon}=\frac{1}{200}$, which is fixed in stage~\ref{it:order-choice-1} and is not $\varepsilon$, and $c_{\mathrm{fl}}$, $c_{\mathrm{rb}}$, $\kappa$ and $\kappa_0$; the thresholds of that stage are $L^F=L^P=5$. Stage~\ref{it:order-choice-6} contains, besides $M$, only $L$, $\gamma_w$, $\alpha_0$, the number $\hat{\kappa}=N_c\cdot\kappa+\log\left(6c_X\right)$ with $N_c=(2L+1)^3$, the integer $m_0$, and the numbers $M^F$ and $M^P$ with their arguments $b$, $\alpha$, $L$, $\alpha_0$ and $\kappa_w+\kappa_0$. By~\eqref{eq:renormalisation-group-19}, $M^F$ is the largest of the numbers $4L+1$, $\left(108\pi\cdot C_V\cdot a\right)^2$ and $108\pi\cdot C_V\cdot a\cdot\exp\left(27\kappa'\right)$ at $\kappa'=\kappa_w+\kappa_0=\kappa+2\kappa_0+3$, where $C_V$ and $a$ are functions of $(b,\alpha,L)$ (Section~\ref{sec:fermion-operators-weakened}); and $M^P$ is determined by $L$, $\alpha_0$, $\gamma_P$ and $n_P=402$, hence by $L$ and $\alpha_0$ (Section~\ref{sec:constants-2}). None of these numbers contains $C_w$, $C_T$, $\varepsilon$ or $e_0$, and stage~\ref{it:order-choice-6} itself states that none of its conditions on $M$ involves $C_w$, $\varepsilon$ or $e$. Where a proof needs an inequality between $L$ and $\varepsilon$, as the proof of Theorem~\ref{thm:polymer-activities-3} needs $1382\,L^{-\frac{15}{8}+\frac{3\varepsilon}{2}}<L^{-\frac{7}{4}}$ in Step~5, $13.6\le L^{\frac{1}{4}-\frac{3\varepsilon}{2}}$ in Step~6 and $L^{-\frac{7}{4}}+4L^{-\frac{7}{8}+\frac{\varepsilon}{2}}<1$ in Step~7, this inequality is derived from the first inequality of stage~\ref{it:order-choice-4}, which is written with $\bar{\varepsilon}$, and from $\varepsilon\le\bar{\varepsilon}$ (in Step~5 itself, and in the paragraph on elementary inequalities for $L$ at the beginning of Section~\ref{sec:polymer-activities-algebra}); and $\varepsilon\le\bar{\varepsilon}$ is a restriction on $\varepsilon$, which is chosen in stage~\ref{it:order-choice-8}, after $L$. Hence the restrictions that the proofs place on $\kappa$, $L$ and $M$ belong to stages~\ref{it:order-choice-3}--\ref{it:order-choice-6}, which precede the choice of $C_w$ in stage~\ref{it:order-choice-7} and the choices of $\varepsilon$ and of $e_0$. The constant $C_T$ is fixed in stage~\ref{it:order-choice-5}, before $M$; that it does not occur in the conditions of stage~\ref{it:order-choice-6} follows from the inspection of that stage above, not from the order of choice.

The primitives enter the proofs of Section~\ref{sec:renormalisation-group-step} only through the properties~\ref{it:renormalisation-group-1}--\ref{it:renormalisation-group-8} and~\ref{it:renormalisation-group-10}--\ref{it:renormalisation-group-17}; algebraic relations among the primitives, their positivity and resolvent identities for them do not enter the arguments. Apart from $\varepsilon$, which enters the bounds of Theorem~\ref{thm:renormalisation-step-small} only through exponents, the only constants occurring in these bounds are $L$ and $c_{\mathrm{fl}}$; besides them, the bounds contain the number $e_k$, a datum of hypothesis~\ref{def:hypothesis-one-step}, also through $e_{k+1}=L^{\frac{1}{2}}\cdot e_k$, $w^{\mathrm{m}}_k$ and $w^{\mathrm{E}}_k$. Finally, each constant and each condition occurring in Sections~\ref{sec:elementary-properties-field}--\ref{sec:data-level-locality} is the same for all data $(m,\mathscr{E})$ in the balls of Theorem~\ref{thm:renormalisation-step-small}, for all sides $l$ and all masses $\mathfrak{m}$, for every $\mathscr{C}'$ and every harmonic part of $A'$, and for every system of primitives with the properties~\ref{it:renormalisation-group-1}--\ref{it:renormalisation-group-17} for $\mathrm{par}$; the level $k$ affects them only in that some displays treat the cases $k=0$ and $k\ge 1$ separately. This proves Theorem~\ref{thm:renormalisation-step-small}(h); clause~\eqref{it:step-small-fields-h} of Theorem~\ref{thm:step-small-fields} follows.
\end{proof}

\section{The range of the running masses}\label{sec:range-running-masses}

This section consists of one statement, Proposition~\ref{prop:bounds-running-masses}, which bounds the running masses and locates the stop level; Proposition~\ref{prop:bounds-masses} follows from its clauses~(a) and~(c)--(e), and so does Theorem~\ref{thm:range-running-masses}, whose hypotheses are those of Proposition~\ref{prop:bounds-masses}. The running quantities and the tori are those of Section~\ref{sec:running-quantities-tori}, and the stages cited are those of Definition~\ref{def:order-choice}.

\begin{proposition}[Bounds on the running masses]\label{prop:bounds-running-masses}
Let $0 < \bar{m} \le 1$ and $\mathfrak{m}_0 \in (0,1]$, let $m_0$ be the least integer $m_0 \ge 0$ with $L^{-m_0} \le \mathfrak{m}_0$ (Definition~\ref{def:mass-bound-stop}), and let $h$, $h'$, $a_k$ and $K$ be defined as in Section~\ref{sec:running-quantities-tori}. For every integer $k$ the running mass of level $k$ is the number $\bar{m} \cdot a_k$; it is denoted by $\mathfrak{m}_k$ when $k \ge 1$, and it equals $\bar{m} \cdot \eta$ when $k = 0$ (Definition~\ref{def:notation}(b) and Section~\ref{sec:running-quantities-tori}). The symbol $\mathfrak{m}_0$ is reserved for the upper end of the mass range $(0, \mathfrak{m}_0]$.
\begin{enumerate}
\item[(a)] One has
\[
\frac{1}{L} < \bar{m} \cdot h \le 1, \qquad L^{-m_0-1} < \bar{m} \cdot a_K = \bar{m} \cdot h' \le L^{-m_0} \le \mathfrak{m}_0.
\]
\item[(b)] For every integer $k$, $\bar{m} \cdot a_{k+1} = L \cdot \bar{m} \cdot a_k$ and $0 < \bar{m} \cdot a_k = L^{-(K-k)} \cdot \bar{m} \cdot a_K$.
\end{enumerate}
Let $e > 0$ be such that $v_K > 0$ and $(1 + \mathrm{lg}/v_K)^r \le L$, let $R = h' \cdot L^T \in \mathrm{Tor}(e)$, and suppose that the integer $N$, with $\eta = L^{-N}$ (Section~\ref{sec:lattice-measures-laws}) and $a_k = L^{k-N}$, satisfies $N \ge N_{\mathrm{min}}(R,e)$; the notation is that of Section~\ref{sec:running-quantities-tori} and Definition~\ref{def:tori-stop-level}. Then the stop level $k^{*} = k^{*}(R)$ of Definition~\ref{def:tori-stop-level} exists and satisfies $1 \le k^{*} \le K$; a subscript $*$ denotes the value at $k = k^{*}$. In particular $K \ge 1$, so that $\mathfrak{m}_K = \bar{m} \cdot a_K = \bar{m} \cdot h'$ and $\mathfrak{m}_{*} = \bar{m} \cdot a_{k^{*}}$. If, in addition, $(1 + 3\,\mathrm{lg}/v_K)^r < L$, then $k^{*}(R^{*}) \in \{K-1, K-2\}$ when $R = R^{*}$ (and $R^{*}$ belongs to $\mathrm{Tor}(e)$ as soon as $M \ge L$), and $k_{\mathrm{top}} \in \{K, K-1\}$ when $R = R_{\mathrm{top}}$.
\begin{enumerate}
\item[(c)] For every integer $k \le k^{*}$,
\[
0 < \bar{m} \cdot a_k \le \mathfrak{m}_{*} \le \mathfrak{m}_K \le L^{-m_0} \le \mathfrak{m}_0.
\]
For every integer $k < k^{*}$, $\bar{m} \cdot a_k \le \mathfrak{m}_{*}/L \le \mathfrak{m}_0/L$; thus, at every level $k$ with $0 \le k < k^{*}$, the running mass $\bar{m} \cdot a_k$ lies in the interval $(0, \mathfrak{m}_0/L]$ prescribed for the mass in hypothesis~\ref{def:hypothesis-one-step}.
\item[(d)] One has
\[
\mathfrak{m}_{*} \cdot l_{*} = \bar{m} \cdot R \ge \bar{m} \cdot h' > L^{-m_0-1},
\]
and consequently $\mathfrak{m}_{*}^{-1} < L^{m_0+1} \cdot l_{*}$.
\item[(e)] Assume $(1 + 3\,\mathrm{lg}/v_K)^r < L$. If $R = R^{*}$, then $\mathfrak{m}_{*} \in \{\mathfrak{m}_K/L,\, \mathfrak{m}_K/L^2\}$, so that $L^{-m_0-3} < \mathfrak{m}_{*} \le L^{-m_0-1}$. If $R = R_{\mathrm{top}}$, then $\mathfrak{m}_{*} \in \{\mathfrak{m}_K,\, \mathfrak{m}_K/L\}$, so that $L^{-m_0-2} < \mathfrak{m}_{*} \le L^{-m_0}$.
\item[(f)] One has $\bar{m} \cdot R \le 1$ if and only if $T \le m_0$, that is, if and only if $R \le h$. For the reference tori, $\bar{m} \cdot R_{\mathrm{top}} = L \cdot \bar{m} \cdot R^{*}$ and, by stage~\ref{it:order-choice-6} of Definition~\ref{def:order-choice},
\[
\bar{m} \cdot R^{*} = \frac{\mathfrak{m}_K \cdot M \cdot r_K}{L} > L^{-m_0-2} \cdot M \cdot r_K \ge r_K;
\]
moreover $r_K \ge v_K^r \ge 1$ as soon as $v_K \ge 1$, and under hypothesis~\ref{def:hypothesis-flow} one has $v_K \ge -\log e_0 \ge 1$ by condition~\eqref{eq:renormalisation-group-A9}.
\end{enumerate}
\end{proposition}

\begin{proof}
By Section~\ref{sec:running-quantities-tori}, $h = L^{\left\lfloor \log_L \frac{1}{\bar{m}} \right\rfloor}$, $h' = h \cdot L^{-m_0}$, $a_k = L^{k-N}$ and $K = N + \log_L h'$, so that $a_K = h'$; further, $e_k = e \cdot a_k^{\frac{1}{2}}$, $v_k = -\log e_k$ and $\mathrm{lg} = \frac{1}{2} \log L$, the number $r_k$ is the least power $L^i$, with $i \ge 0$ an integer, such that $L^i \ge v_k^r$, and
\[
l_k = \frac{R}{a_k}, \qquad \mathfrak{n}_k = \log_L\left( \frac{l_k}{M \cdot r_k} \right),
\]
where $M = L^m$ (Section~\ref{sec:lattice-geometry}). These formulas define $a_k$, $e_k$, $v_k$, $r_k$, $l_k$ and $\mathfrak{n}_k$ for every integer $k$. By Definition~\ref{def:tori-stop-level}, $R_{\mathrm{top}} = h' \cdot M \cdot r_K$, and $\mathrm{Tor}(e)$ is the set of the numbers $R = h' \cdot L^T$ with an integer $T \ge 0$ and $R \le R_{\mathrm{top}}$; moreover $R^{*} = R_{\mathrm{top}}/L$, the stop level $k^{*}(R)$ is the first level $k$ with $\mathfrak{n}_k = 0$, $k_{\mathrm{top}} = k^{*}(R_{\mathrm{top}})$, and $N_{\mathrm{min}}(R,e)$ has the following defining property: $R \cdot L^{N''} \ge L \cdot M \cdot r_0$ for every integer $N'' \ge N_{\mathrm{min}}(R,e)$ with $N'' + \log_L h' \ge 0$, where $r_0$ is formed with $N''$ in the place of $N$.

\emph{The integers $\mathfrak{n}_k$ and the stop level.} Since $l_k/(M \cdot r_k) = h' \cdot L^{T+N-k}/(M \cdot r_k)$ and $h'$, $M$, $r_k$ are integral powers of $L$, every $\mathfrak{n}_k$ is an integer. From $a_{k+1} = L \cdot a_k$ one obtains $e_{k+1} = L^{\frac{1}{2}} \cdot e_k$, that is, $v_k = v_{k+1} + \mathrm{lg}$, and $l_{k+1} = l_k/L$. Let $k$ be an integer with $k + 1 \le K$. Then $v_k > v_{k+1} \ge v_K > 0$; hence $r_k \ge r_{k+1}$, and
\[
v_k^r = v_{k+1}^r \cdot \left( 1 + \frac{\mathrm{lg}}{v_{k+1}} \right)^r \le v_{k+1}^r \cdot \left( 1 + \frac{\mathrm{lg}}{v_K} \right)^r \le L \cdot v_{k+1}^r \le L \cdot r_{k+1}.
\]
Since $L \cdot r_{k+1}$ is a power $L^i$ with an integer $i \ge 0$, it follows that $r_k \le L \cdot r_{k+1}$. As $r_k$ and $r_{k+1}$ are powers of $L$, one concludes that $r_k \in \{r_{k+1}, L \cdot r_{k+1}\}$ and
\[
\mathfrak{n}_{k+1} - \mathfrak{n}_k = -1 + \log_L\left( \frac{r_k}{r_{k+1}} \right) \in \{-1, 0\}.
\]
In particular, $k \mapsto r_k$ and $k \mapsto \mathfrak{n}_k$ are non-increasing on the integers $k \le K$.

Since $R \in \mathrm{Tor}(e)$, one has $R \le R_{\mathrm{top}} = h' \cdot M \cdot r_K$, that is, $R/h' \le M \cdot r_K$. As $l_K = R/a_K = R/h'$, this says that $\mathfrak{n}_K = \log_L(R/R_{\mathrm{top}}) \le 0$.

Next, $K \ge 1$. Suppose, to the contrary, that $K \le 0$, and put $N' := N - K$. Then $N' \ge N \ge N_{\mathrm{min}}(R,e)$ and $N' + \log_L h' = 0$, so that the defining property of $N_{\mathrm{min}}(R,e)$ applies to $N'' = N'$. Replacing $N$ by $N'$ turns $a_0$ into $L^{-N'} = h'$, and therefore the charge $e \cdot a_0^{\frac{1}{2}}$ of level $0$ into $e \cdot h'^{\frac{1}{2}} = e_K$; since $r_0$ is formed with this charge, it turns into $r_K$. Together with $R/h' \le M \cdot r_K$ this yields
\[
R \cdot L^{N'} = \frac{R}{h'} \le M \cdot r_K < L \cdot M \cdot r_0,
\]
where $r_0$ is formed with $N'$; this is incompatible with the defining property. Hence $K \ge 1$. In particular $N + \log_L h' \ge 0$, and the defining property for $N'' = N$ gives $l_0 = R \cdot L^N \ge L \cdot M \cdot r_0$, that is, $\mathfrak{n}_0 \ge 1$.

Consequently $\mathfrak{n}_0, \mathfrak{n}_1, \dots, \mathfrak{n}_K$ is a sequence of integers which begins with a value $\ge 1$, ends with a value $\le 0$, and decreases by $0$ or by $1$ at each step. It takes the value $0$, and the first index at which it does so, namely $k^{*}(R)$, satisfies $1 \le k^{*} \le K$.

Assume now that $(1 + 3\,\mathrm{lg}/v_K)^r < L$. For $i \in \{1, 2, 3\}$ one has $v_{K-i} = v_K + i \cdot \mathrm{lg}$ and
\[
v_{K-i}^r = v_K^r \cdot \left( 1 + \frac{i \cdot \mathrm{lg}}{v_K} \right)^r < L \cdot v_K^r \le L \cdot r_K,
\]
so that $r_{K-i} \le L \cdot r_K$; as $r_{K-i} \ge r_K$ by the monotonicity of $k \mapsto r_k$, one has $r_{K-i} \in \{r_K, L \cdot r_K\}$.

If $M \ge L$, that is, $m \ge 1$, then $R^{*} = h' \cdot L^{m-1} \cdot r_K$, where $L^{m-1} \cdot r_K$ is a power of $L$ with a non-negative integral exponent, and $R^{*} \le R_{\mathrm{top}}$; thus $R^{*} \in \mathrm{Tor}(e)$.

Let $R = R^{*}$. As $h'/a_k = L^{K-k}$, one has $l_k = M \cdot r_K \cdot L^{K-k-1}$, and therefore
\[
\mathfrak{n}_{K-i} = i - 1 - \log_L\left( \frac{r_{K-i}}{r_K} \right) \qquad (i = 0, 1, 2, 3).
\]
Hence $\mathfrak{n}_K = -1$, $\mathfrak{n}_{K-1} \in \{0, -1\}$, $\mathfrak{n}_{K-2} \in \{1, 0\}$ and $\mathfrak{n}_{K-3} \in \{2, 1\}$, and $\mathfrak{n}_k \ge \mathfrak{n}_{K-3} \ge 1$ for $k \le K-3$, because $k \mapsto \mathfrak{n}_k$ is non-increasing. If $\mathfrak{n}_{K-2} = 0$, then $k^{*} = K-2$. Otherwise $\mathfrak{n}_{K-2} = 1$; the step from $K-2$ to $K-1$ being $0$ or $-1$, one gets $\mathfrak{n}_{K-1} \in \{1, 0\}$, hence $\mathfrak{n}_{K-1} = 0$ and $k^{*} = K-1$. Replacing $R^{*}$ by $R_{\mathrm{top}} = L \cdot R^{*}$ multiplies every $l_k$ by $L$ and leaves every $r_k$ unchanged, so that every $\mathfrak{n}_k$ increases by one. On $R_{\mathrm{top}}$ one therefore has $\mathfrak{n}_K = 0$, $\mathfrak{n}_{K-1} \in \{1, 0\}$, and $\mathfrak{n}_k \ge \mathfrak{n}_{K-2} \ge 1$ for $k \le K-2$; whence $k_{\mathrm{top}} \in \{K-1, K\}$.

\emph{Clause (a).} Put $t := \log_L(1/\bar{m})$, a real number $\ge 0$. Then $\bar{m} = L^{-t}$ and $h = L^{\lfloor t \rfloor}$, so that $\bar{m} \cdot h = L^{\lfloor t \rfloor - t}$ with $\lfloor t \rfloor - t \in (-1, 0]$; that is, $1/L < \bar{m} \cdot h \le 1$. Moreover $\bar{m} \cdot a_K = \bar{m} \cdot h' = \bar{m} \cdot h \cdot L^{-m_0}$, which lies in $(L^{-m_0-1}, L^{-m_0}]$, and $L^{-m_0} \le \mathfrak{m}_0$ by the choice of $m_0$.

\emph{Clause (b).} Let $k$ be an integer. One has $\bar{m} > 0$, $a_k = L^{k-N} > 0$ and $a_{k+1} = L \cdot a_k$; hence $\bar{m} \cdot a_k > 0$, $\bar{m} \cdot a_{k+1} = L \cdot \bar{m} \cdot a_k$ and $\bar{m} \cdot a_k = L^{k-K} \cdot \bar{m} \cdot a_K$.

\emph{Clause (c).} By (b), $k \mapsto \bar{m} \cdot a_k$ is increasing on the integers. As $1 \le k^{*} \le K$, one has $\mathfrak{m}_{*} = \bar{m} \cdot a_{k^{*}}$ and $\mathfrak{m}_K = \bar{m} \cdot a_K$; hence every integer $k \le k^{*}$ satisfies $0 < \bar{m} \cdot a_k \le \mathfrak{m}_{*} \le \mathfrak{m}_K$, and $\mathfrak{m}_K \le L^{-m_0} \le \mathfrak{m}_0$ by (a). For $k < k^{*}$, (b) gives $\bar{m} \cdot a_k = L^{-(k^{*}-k)} \cdot \mathfrak{m}_{*} \le \mathfrak{m}_{*}/L$, and $\mathfrak{m}_{*} \le \mathfrak{m}_0$.

\emph{Clause (d).} Since $l_{*} = R/a_{*}$ and $\mathfrak{m}_{*} = \bar{m} \cdot a_{*}$, one has $\mathfrak{m}_{*} \cdot l_{*} = \bar{m} \cdot R = \bar{m} \cdot h' \cdot L^T$. As $T \ge 0$, this is at least $\bar{m} \cdot h' = \mathfrak{m}_K$, and $\mathfrak{m}_K > L^{-m_0-1}$ by (a). It follows that $\mathfrak{m}_{*}^{-1} = l_{*}/(\bar{m} \cdot R) < L^{m_0+1} \cdot l_{*}$.

\emph{Clause (e).} It was shown above that $k^{*}(R^{*}) \in \{K-1, K-2\}$ and $k_{\mathrm{top}} \in \{K, K-1\}$. By (b), $\mathfrak{m}_{*} = L^{-(K-k^{*})} \cdot \mathfrak{m}_K$; so $\mathfrak{m}_{*} \in \{\mathfrak{m}_K/L,\, \mathfrak{m}_K/L^2\}$ for $R = R^{*}$, and $\mathfrak{m}_{*} \in \{\mathfrak{m}_K,\, \mathfrak{m}_K/L\}$ for $R = R_{\mathrm{top}}$. Together with $L^{-m_0-1} < \mathfrak{m}_K \le L^{-m_0}$, from (a), this gives
\[
L^{-m_0-3} < \frac{\mathfrak{m}_K}{L^2} \le \mathfrak{m}_{*} \le \frac{\mathfrak{m}_K}{L} \le L^{-m_0-1} \qquad (R = R^{*}),
\]
\[
L^{-m_0-2} < \frac{\mathfrak{m}_K}{L} \le \mathfrak{m}_{*} \le \mathfrak{m}_K \le L^{-m_0} \qquad (R = R_{\mathrm{top}}).
\]

\emph{Clause (f).} One has $\bar{m} \cdot R = \bar{m} \cdot h' \cdot L^T = (\bar{m} \cdot h) \cdot L^{T-m_0}$, where $T$ and $m_0$ are integers. In the case $T \le m_0$ the factor $L^{T-m_0}$ is at most $1$, and the upper bound $\bar{m} \cdot h \le 1$ of (a) gives $\bar{m} \cdot R \le 1$. In the case $T \ge m_0 + 1$ the factor $L^{T-m_0}$ is at least $L$, and the lower bound $\bar{m} \cdot h > L^{-1}$ of (a) gives $\bar{m} \cdot R > L^{-1} \cdot L = 1$. Since $R = h \cdot L^{T-m_0}$, the inequality $T \le m_0$ is equivalent to $R \le h$. Next, $R_{\mathrm{top}} = L \cdot R^{*}$ and $R^{*} = h' \cdot M \cdot r_K/L$ by Definition~\ref{def:tori-stop-level}, so that $\bar{m} \cdot R_{\mathrm{top}} = L \cdot \bar{m} \cdot R^{*}$ and $\bar{m} \cdot R^{*} = \mathfrak{m}_K \cdot M \cdot r_K/L$. By (a), $\mathfrak{m}_K/L > L^{-m_0-2}$, and $M \ge L^{m_0+2}$ by stage~\ref{it:order-choice-6}; therefore $\bar{m} \cdot R^{*} > L^{-m_0-2} \cdot M \cdot r_K \ge r_K$. By the definition of $r_K$, one has $r_K \ge v_K^r$, and $v_K^r \ge 1$ when $v_K \ge 1$. Finally, under hypothesis~\ref{def:hypothesis-flow} one has $e_K \le e_1$; moreover $e_1 \le e_0$ by stage~\ref{it:order-choice-10}, and $e_0 < 1$ and $\log^2 e_0 \ge 10 L$ by condition~\eqref{eq:renormalisation-group-A9}. Hence $v_K = -\log e_K \ge -\log e_0 > 0$; since moreover $\log^2 e_0 \ge 10 L \ge 1$, the positive number $-\log e_0$ is at least $1$.
\end{proof}

\section{Renormalisation flow, the tuned pair and partition functions on the reference tori}\label{sec:renormalisation-flow-tuned}

This section contains the proofs of the theorems of Section~\ref{sec:trajectories-pair-bounds}, of Proposition~\ref{prop:rows-component} and of property~\ref{it:renormalisation-group-18} in the list of Section~\ref{sec:properties-systems-primitives}, the definition of the number $e_1$ of stage~\ref{it:order-choice-10}, and the construction of the pair $P_N(e)$ of Definition~\ref{def:tuned-pair}, whose entries are the bare mass shift $m_{0,N}$ and the vacuum energy $\boldsymbol{\varepsilon}_{0,N}$. The notation is that of Section~\ref{sec:renormalisation-group-scheme}.

\subsection{Setting and constants}\label{sec:setting-constants}

\subsubsection{Notation and the standing hypothesis}\label{sec:notation-standing-hypothesis}

All constants, and the number $\bar{m}$, the probe $\beta$, the point $x_0$, the charge $e$, the side $R$ and the integer $N$, are taken from Section~\ref{sec:order-choice}, where they are fixed one after another, stage~\ref{it:order-choice-1} first and stage~\ref{it:order-choice-12} last. On the constants of stages~\ref{it:order-choice-1}--\ref{it:order-choice-10} one imposes, from Section~\ref{sec:conditions}, the conditions \eqref{eq:renormalisation-group-A0}--\eqref{eq:renormalisation-group-A10}, \eqref{eq:renormalisation-group-19}, \eqref{eq:renormalisation-group-20}--\eqref{eq:renormalisation-group-24}, \eqref{eq:renormalisation-group-25} and \eqref{eq:renormalisation-group-26}. The conditions \eqref{eq:renormalisation-flow-3}, \eqref{eq:renormalisation-group-fl} and \eqref{eq:renormalisation-group-T1}--\eqref{eq:renormalisation-group-T6}, which concern the side $R$ of the torus, the charge $e$ and the cutoff $N$, are recalled in Section~\ref{sec:constants-conditions-torus}. No condition of stage~\ref{it:order-choice-11} is in force unless the statement at hand names it.

The exponential function is written $\exp$ (Definition~\ref{def:notation}(a)). The letter $e$ denotes the charge, and $e_0$ is the ceiling of the running charges, fixed in stage~\ref{it:order-choice-9}. As in Section~\ref{sec:notation}, $\mathfrak{m}$ denotes a running mass and $\nu \in \{0, 1, 2\}$ a direction, whose unit vector is $e_{\nu}$. For $x \in \mathbb{R}^3$, $|x|$ is the sup norm, and $d(\cdot,\cdot)$ is the sup distance on a torus (Definition~\ref{def:sup-norm-distance}).

\paragraph{Scales and tori.}
Let $L$ be the odd integer of stage~\ref{it:order-choice-4}, $M = L^m$ the integer of stage~\ref{it:order-choice-6}, $r$, $p_0$, $p$ and $\varepsilon$ the exponents of stage~\ref{it:order-choice-8}, and $m_0$ the integer of Definition~\ref{def:mass-bound-stop}. Let $0 < \bar{m} \le 1$, let $N \ge 0$ be an integer, and put $\eta := L^{-N}$. As in Section~\ref{sec:running-quantities-tori}, for the levels $k \ge 0$ let
\begin{align*}
& h := L^{\lfloor \log_L \bar{m}^{-1} \rfloor}, \qquad h' := h L^{-m_0}, \qquad a_k := L^{k-N}, \qquad e_k := e\, a_k^{\frac{1}{2}}, \qquad v_k := -\log e_k,\\
& \mathrm{lg} := \frac{1}{2} \log L, \qquad \mathfrak{m}_k := \bar{m} \cdot a_k, \qquad r_k := \min\left\{ L^i : i \ge 0 \text{ an integer},\ L^i \ge v_k^r \right\},\\
& p_k := v_k^p, \qquad p_{0,k} := v_k^{p_0}, \qquad w^{\mathrm{m}}_k := e_k^{\frac{3}{4} - 8\varepsilon}, \qquad w^{\mathrm{E}}_k := e_k^{\frac{1}{4} - 7\varepsilon},
\end{align*}
and let $K$ be the level with $a_K = h'$, that is, $K = N + \log_L h'$. At $k = 0$ one has $e_k = q := e\, \eta^{\frac{1}{2}}$, the running charge of level $0$; the numbers $v_0$, $r_0$, $p_{0,0}$, $w^{\mathrm{m}}_0$, $w^{\mathrm{E}}_0$ are formed with $q$, while the symbol $e_0$ keeps its meaning as the ceiling. In the same way, among the numbers of the renormalisation group scheme, $p_0$ is the exponent of stage~\ref{it:order-choice-8} in $p_{0,k} = v_k^{p_0}$ and never the value of $p_k$ at $k = 0$, which is written $v_0^p$; and the running mass of level $0$ is written $\bar{m} \cdot \eta$, since $\mathfrak{m}_0$ is the upper end of the mass range (Definition~\ref{def:mass-bound-stop}). Other uses of the symbol $p_0$ are declared where they occur; for instance, in Section~\ref{sec:determinants-constant-connections} it is the component of index $0$ of a momentum $p = (p_0, p_1, p_2)$.

The sides of the tori are the numbers $R = h' L^T$ with an integer $T \ge 0$. For such $R$ put
\[
l_k := \frac{R}{a_k}, \qquad n_k := l_k^3, \qquad \mathfrak{n}_k := \log_L\left(\frac{l_k}{M r_k}\right),
\]
and let $T_k(R)$ be the fine lattice of level $k$ and side $l_k$ (Section~\ref{sec:lattice-geometry}). Following Definition~\ref{def:tori-stop-level}, let
\[
R_{\mathrm{top}} := h' M r_K, \qquad \mathrm{Tor}(e) := \left\{ h' L^T : T \ge 0 \text{ an integer},\ h' L^T \le R_{\mathrm{top}} \right\}, \qquad R^{*} := \frac{R_{\mathrm{top}}}{L};
\]
the two reference tori are $R^{*}$ and $L R^{*} = R_{\mathrm{top}}$. For $R \in \mathrm{Tor}(e)$ let $N_{\mathrm{min}}(R,e)$ be the least integer $N_1$ such that $R L^N \ge L M r_0$ for every $N \ge N_1$, where $r_0$ is formed with $q = e\, L^{-\frac{N}{2}}$ and so depends on $N$; here $N$ and $N_1$ range over the integers $\ge 0$ with $K = N + \log_L h' \ge 0$, $K$ being the index of a level. The stop level $k^{*} = k^{*}(R)$ is the least level $k$ with $\mathfrak{n}_k = 0$, and $k_{\mathrm{top}} := k^{*}(R_{\mathrm{top}})$. For the running quantities and the other objects of the scheme that are indexed by the level, a subscript $*$ marks the value at $k = k^{*}(R)$, as in $r_{*}$, $v_{*}$, $e_{*}$, $l_{*}$, $p_{*}$, $p_{0*}$, $\mathfrak{m}_{*}$; among these quantities, $n_{*}$ is excepted, being defined by $n_{*} := R^{*}/\eta$ (Definition~\ref{def:notation}(e)). Other uses of a subscript $*$ are declared where they occur. Finally, for $\theta \in \mathbb{R}^3$ and a level $k$, $c^{(k)}_{\theta}$ is the constant real bond function whose value on each bond of the direction $\nu$, oriented in the positive sense, is $\theta_{\nu}/(e_k l_k)$, and $c_{\theta} := c^{(k^{*})}_{\theta}$ (Definition~\ref{def:twist-external-connection}).

\paragraph{Lattice measures with twist and external connection.}
For a side $R$ as above let $n := R/\eta$ and $\Lambda_n := (\mathbb{Z}/n\mathbb{Z})^3$ (Section~\ref{sec:lattice-measures-laws}). Let $C^j$ be the space of real $j$-cochains on $\Lambda_n$ with the counting $\ell^2$ norm, $d$ the coboundary, and $C^1 = V_{\mathrm{g}} \oplus V_{\mathrm{h}} \oplus V_{\mathrm{t}}$ the orthogonal decomposition into the exact one-cochains, the constant ones, and the coclosed ones with mean zero. Put $F_N := dC^1$, let $\gamma_N$ be the standard Gaussian measure on $F_N$, and for $F \in F_N$ let $A(F)$, the potential of $F$, be the element of $V_{\mathrm{t}}$ with $dA(F) = F$. For a complex function $A$ on the oriented bonds of $\Lambda_n$ with $A(x',x) = -A(x,x')$, let $D_A$ be the Wilson operator of Section~\ref{sec:lattice-measures-laws}; it acts on the functions $f \colon \mathbb{Z}^3 \to \mathbb{C}^4$ with $f(x + n e_{\nu}) = -f(x)$, and it is the operator of Section~\ref{sec:exact-fermion-operators} at level $0$, with the charge $q$. For a real pair $(m_{\mathrm{in}}, \boldsymbol{\varepsilon}_{\mathrm{in}}) \in \mathbb{R}^2$, a twist $\theta \in \mathbb{R}^3$, an external connection $a \in C^1 \otimes \mathbb{C}$ and a real number $\mathbf{m}$, set
\begin{align}
H_{\theta}(x, x + e_{\nu}) &:= \frac{\theta_{\nu}}{q n}, \notag\\
\rho^{\theta,a}_N(F) &:= \exp\left(-\boldsymbol{\varepsilon}_{\mathrm{in}} \cdot n^3\right) \frac{\det\left(D_{A(F) + H_{\theta} + a} + \bar{m} \cdot \eta + m_{\mathrm{in}}\right)}{\det\left(D_0 + \bar{m} \cdot \eta\right)}, \notag\\
\nu^{\theta,a} &:= \rho^{\theta,a}_N\, \gamma_N, \qquad \nu^{\theta} := \nu^{\theta,0}, \qquad \bar{\nu} := \int_{[-\pi,\pi]^3} \frac{\mathrm{d}\theta}{(2\pi)^3}\, \nu^{\theta}, \qquad \bar{\mu} := \frac{\bar{\nu}}{\bar{\nu}(1)}, \notag\\
\rho_n(q, \mathbf{m}) &:= \int_{[-\pi,\pi]^3} \frac{\mathrm{d}\theta}{(2\pi)^3} \int \mathrm{d}\gamma_N(F)\, \det\left(D_{A(F) + H_{\theta}} + \mathbf{m}\right), \notag\\
\rho^{\mathrm{free}}_n &:= \int_{[-\pi,\pi]^3} \frac{\mathrm{d}\theta}{(2\pi)^3}\, \det\left(D_{H_{\theta}} + \bar{m} \cdot \eta\right). \label{eq:renormalisation-flow-1}
\end{align}
The measure $\bar{\mu}$ is defined when $\bar{\nu}(1) > 0$. Since $\det\left(D_0 + \bar{m} \cdot \eta\right) > 0$ (Lemma~\ref{lem:product-formula-volume}, applied with $\theta = 0$ and the mass $\bar{m} \cdot \eta \in (0, 1]$), the identity~\eqref{eq:renormalisation-flow-2} below shows that $\bar{\nu}(1)$ is a positive multiple of $\rho_n\left(q, \bar{m} \cdot \eta + m_{\mathrm{in}}\right)$. If $n \ge 2$, $e > 0$ and the lattice mass $\bar{m} \cdot \eta + m_{\mathrm{in}}$ exceeds $-1$, this number is positive by Theorem~\ref{thm:limit-vanishing-proca}, with the identifications made in the proof of Lemma~\ref{lem:positivity-large-fields}; in particular, under the standing hypothesis~\ref{def:hypothesis-flow}, $\bar{\nu}(1) > 0$ for every pair that is controlled on $R$ (Lemma~\ref{lem:positivity-large-fields}). The measure $\gamma_N$ is the image under $d$ of the probability measure on $V_{\mathrm{t}}$ whose density with respect to the Lebesgue measure $\mathrm{d}A$ is proportional to $\exp\left(-\frac{1}{2}|dA|^2\right)$. Integrating $\rho^{\theta,0}_N$ against $\gamma_N$ and averaging over $\theta \in [-\pi,\pi]^3$, one obtains from \eqref{eq:renormalisation-flow-1}
\begin{equation}\label{eq:renormalisation-flow-2}
\rho_n\left(q, \bar{m} \cdot \eta + m_{\mathrm{in}}\right) = \exp\left(\boldsymbol{\varepsilon}_{\mathrm{in}} \cdot n^3\right) \det\left(D_0 + \bar{m} \cdot \eta\right) \cdot \bar{\nu}(1).
\end{equation}

\paragraph{The standing hypothesis.}
Besides the conditions on the constants stated at the beginning of Section~\ref{sec:notation-standing-hypothesis}, the standing hypothesis~\ref{def:hypothesis-flow} of Section~\ref{sec:one-step-operators} comprises the following requirements: $e > 0$ and $e_K = e\, h'^{\frac{1}{2}} \le e_1$; $R \in \mathrm{Tor}(e)$; $N \ge N_{\mathrm{min}}(R,e)$; and the primitives are those of Definition~\ref{def:periodisation-projection}. Thus, on $T_k(R_{\mathrm{top}})$ with $k < k_{\mathrm{top}}$, the primitives are the systems formed by the families of Theorem~\ref{thm:fermion-operators-weakened}, taken at $\kappa' = \kappa_w + \kappa_0$, together with the families of Theorem~\ref{thm:photon-operators-weakened} (Lemma~\ref{lem:assembled-systems-primitives}); on $T_k(R)$ with $R < R_{\mathrm{top}}$ and $k < k^{*}(R)$, they are the projections of these systems (Definition~\ref{def:projection}).

\paragraph{Trajectories and controlled pairs.}
The spaces, norms and domains, the step map and the flows are those of Sections~\ref{sec:domains-bond-functions}--\ref{sec:flow}, the step map being given by the displays of Section~\ref{sec:definition-step-map}. Let $(m_{\mathrm{in}}, \boldsymbol{\varepsilon}_{\mathrm{in}})$ be a real pair. Its trajectory on $R$ is the sequence $(\boldsymbol{\varepsilon}^t_k, m_k, E_k)$, and its trajectory with slot is the sequence $(\mathscr{E}_k, \mathfrak{P}_k)$; both are given by the recursion of Section~\ref{sec:flow} on the torus $R$, with its systems of primitives at the levels $k < k^{*}(R)$, starting from $\boldsymbol{\varepsilon}^t_0 = \boldsymbol{\varepsilon}_{\mathrm{in}}$, the mass $m_{\mathrm{in}}$ at level $0$, $E_0 = 0$, $\mathscr{E}_0 = 0$ and $\mathfrak{P}_0 = 0$; at $k = 0$ the symbol $m_k$ stands for $m_{\mathrm{in}}$. In the proofs of Propositions~\ref{prop:existence-trajectory} and~\ref{prop:pair-vanishing-terminal} in Section~\ref{sec:flow-fixed-torus}, the term of index $0$ of the sequence $(m_k)$ is written $m_0$ (Definition~\ref{def:notation}(b)), both for a trajectory, where this term is $m_{\mathrm{in}}$, and for the fixed point of the map $\mathrm{Fx}$ of the proof of Proposition~\ref{prop:pair-vanishing-terminal}; there $m_0$ is not the integer of Definition~\ref{def:mass-bound-stop}. In the proof of Lemma~\ref{lem:dependence-twist-uniform}, $m_0$ is the component of index $0$ of an integer vector $m = (m_0, m_1, m_2)$. Everywhere else in Section~\ref{sec:renormalisation-flow-tuned}, $m_0$ is the integer of Definition~\ref{def:mass-bound-stop}. Let $k_1$ be a level. One says that the trajectory exists up to $k_1$ when, at each level $k < k_1$, the point $(m_k, \mathscr{E}_k)$ belongs to the domain of the step map at that level, namely the ball $\mathrm{Ball}^x_k(4)$ of Section~\ref{sec:spaces-activities}. Definition~\ref{def:fluctuation-field-controlled} imposes two requirements for the pair to be controlled on $R$: first, its trajectory on $R$ must exist up to the stop level $k^{*}(R)$; second, the masses of this trajectory must satisfy $|m_k| \le w^{\mathrm{m}}_k$ at all levels $k = 0, 1, \dots, k^{*}(R)$.

\paragraph{The small-field set.}
Let $\mathrm{Sl}$ be the slice of Definition~\ref{def:coordinates-slice-small}. The coboundary $d$ maps $\mathrm{Sl}$ isomorphically onto $F_N$ (Proposition~\ref{prop:orthogonal-coordinates}); for $F \in F_N$ let $A_0(F) \in \mathrm{Sl}$ be its preimage. Formed with $A_0 = A_0(F)$, the block fields $A_j$ on $T^0_j$ ($0 \le j \le k^{*}$), the field $A_{*} := A_{k^{*}}$ and the fluctuation variables $\zeta_j \in V_j$ ($0 \le j < k^{*}$) of Definition~\ref{def:coordinates-slice-small} are functions of $F$. Here $\zeta_j$ is the element of the fluctuation space $V_j$ of Section~\ref{sec:averaging-fluctuations-gauge} determined by $A_j = A^{\min}_j\left(A_{j+1}\right) + \mathrm{Sq}_j\, \zeta_j$, with the minimiser $A^{\min}_j$ and the map $\mathrm{Sq}_j$ of that section, and $\mathrm{Sq}_j\, \zeta_j = A_j - A^{\min}_j\left(A_{j+1}\right)$ is the fluctuation field of level $j$ (Definition~\ref{def:fluctuation-field-controlled}). For $\Upsilon$ in the basis $\{\Upsilon\}$ of $V_j$ (Section~\ref{sec:averaging-fluctuations-gauge}) put $\zeta_j(\Upsilon) := \langle \zeta_j, \Upsilon \rangle$. Let $\mathrm{d}\mathsf{m}_{*}$ be the law of $A_{*}$ under $\gamma_N$, and for a real bond function $A$ on $T^0_{k^{*}}$ let $\chi_{*}(A) := \mathbf{1}\left\{ |dA|_{\infty} \le p_{*} \right\}$. The small-field set is
\[
\mathfrak{G} := \left\{ F \in F_N : |\zeta_j(\Upsilon)| \le p_{0,j} \text{ for every } \Upsilon \text{ and every } j < k^{*}, \text{ and } |dA_{*}|_{\infty} \le p_{*} \right\}.
\]
For a real bond function $A$ on $T^0_{k^{*}}$ and a family $\zeta = (\zeta_j)_{j < k^{*}}$ with $\zeta_j \in V_j$, define real bond functions $A^{\downarrow}_j$ on $T^0_j$ by the downward recursion
\[
A^{\downarrow}_{k^{*}} := A, \qquad A^{\downarrow}_j := A^{\min}_j\left(A^{\downarrow}_{j+1}\right) + \mathrm{Sq}_j\, \zeta_j \qquad (j = k^{*} - 1, \dots, 0),
\]
with $A^{\min}_j$ and $\mathrm{Sq}_j$ as in Section~\ref{sec:averaging-fluctuations-gauge}, and denote the last of them by $A^{\downarrow}_0(A, \zeta)$.

\subsubsection{Constants, conditions on the torus, and the number $e_1$}\label{sec:constants-conditions-torus}

For $g > 0$ let
\[
c(g) := \sum_{z \in \mathbb{Z}^3} \exp(-g |z|), \qquad c_{\mathrm{bd}}(g) := 3 \exp(3g) \cdot c(g);
\]
$c_{\mathrm{bd}}$ is the function of \eqref{eq:operators-photon-1}. Let $c_{\kappa}$ be the supremum, over all levels $k \ge 0$, all sides $l \in \mathrm{Cl}$ and all $M$-cubes $\square$ of $T_k$, of
\[
\sum_{X \in D_k,\ X \supseteq \square} \exp\left(-\kappa\, d_M(X)\right).
\]
Since $d_M \ge 0$ and $\kappa \ge \kappa_0$ (stage~\ref{it:order-choice-3}), the bound of Proposition~\ref{prop:polymers}(t2) gives $c_{\kappa} \le K_0$, with the constants $\kappa_0$, $K_0$ of stage~\ref{it:order-choice-1}. With $b_{+} = b$ (Section~\ref{sec:running-quantities-tori}) define
\[
C_0 := \max\left\{ C_P \cdot c_{\mathrm{bd}}(\gamma_P),\ C^F,\ 4 b_{+} + 16\, C^F \cdot c(\gamma^F) \right\},
\]
where $C_P$, $\gamma_P$ are the constants of \eqref{eq:operators-photon-18}, and $C^F$, $\gamma^F$ are the constants of Theorem~\ref{thm:fermion-operators-weakened} at $\kappa' = \kappa_w + \kappa_0$, for which that theorem gives $(C_F, \gamma_F) = (C^F, \gamma^F)$.

The number $C_0$ is the constant of stage~\ref{it:order-choice-7}. It dominates the following numbers: the number $C_P \cdot c_{\mathrm{bd}}(\gamma_P)$, which by \eqref{eq:operators-photon-2} (applicable since $0 < \gamma_P \le 1$) bounds the sum over the second index $b'$ of the kernel bound $C_P \cdot \exp\left(-\gamma_P \cdot d(b,b')\right)$ of Proposition~\ref{prop:kernel-bounds}; the constant $C_Q = 1$ of Proposition~\ref{prop:block-averages}, because $C_P \ge 1$ and $c_{\mathrm{bd}} \ge 3$; the constant $C_F$ of Proposition~\ref{prop:bounds}; and the number $C_0^{\mathrm{unit}} = 4 b_{+} + 16\, C_F \cdot c(\gamma_F)$ of Proposition~\ref{prop:bounds-unit-lattice} (Section~\ref{sec:absolute-constants}).

Let $\gamma_1 \in (0, 1/16]$ and $C_{\Sigma} < \infty$ be the constants of Lemma~\ref{lem:dependence-twist-uniform} for the interval $[b_{-}, b_{+}]$, where $b_{-} = b(1 - L^{-1})$. They are constants of stage~\ref{it:order-choice-7} and functions of $(b_{-}, b_{+})$ alone. Section~\ref{sec:determinants-constant-connections} rests on no statement of Sections~\ref{sec:scales-stop-level}--\ref{sec:representation-stop-level}; the numbers $\gamma_1$ and $C_{\Sigma}$ enter the definition of $e_1$ below only as numbers, and Proposition~\ref{prop:consequences} uses no assertion of Lemma~\ref{lem:dependence-twist-uniform}. Hence the definition of $e_1$ involves no circularity. Recall also $d_U = 2L + 1$ and $c_l = 2L$ (Definition~\ref{def:parameter-cubes-polydisks}).

With the constant $c_{\mathrm{fl}}$ of stage~\ref{it:order-choice-1}, the conditions of the second kind of Section~\ref{sec:conditions}, namely \eqref{eq:renormalisation-flow-3}, \eqref{eq:renormalisation-group-fl} and \eqref{eq:renormalisation-group-T1}--\eqref{eq:renormalisation-group-T6}, are conditions on the triple $(R, e, N)$, stated for $R \in \mathrm{Tor}(e)$ and $N \ge N_{\mathrm{min}}(R,e)$; in this order they read
\begin{align}
& r_{*} > 31 L + 16; \label{eq:renormalisation-flow-3}\\
& \sum_{k=0}^{k^{*}} \mathfrak{s}_k \le \frac{1}{16}, \quad \text{where } \mathfrak{s}_k := \left(c_{\mathrm{fl}} + 2\right) e_k^{\varepsilon}; \qquad v_{*} \ge 8; \notag\\
& v_{*} \ge 8, \qquad p_0\, v_{*}^{p_0-1} \ge 28, \qquad w^{\mathrm{m}}_0 < 1 \quad \left(w^{\mathrm{m}}_0 \text{ formed with } q\right); \notag\\
& 8\, C_0^2\, l_{*}^6 \left(b_{-}^{-1} + \mathfrak{m}_{*}^{-1}\right) \le e_{*}^{-\frac{\varepsilon}{2}}, \qquad 4 \left(2 C_0 + c_{\kappa} + 5 c_{\mathrm{fl}} + 6\right) l_{*}^3 \le e_{*}^{-\varepsilon}, \notag\\
& \qquad e_{*}^{\frac{1}{4} - 9\varepsilon} \le \frac{1}{8}; \notag\\
& 4\, C_0\, l_{*}^2\, p_{*} < e_{*}^{-\varepsilon}; \notag\\
& 12\, l_{*}^3 \left( \exp\left(-\frac{p_{0*}}{2}\right) + \exp\left(-\frac{p_{*}^2}{2 C_0^2}\right) \right) \le e_{*}; \notag\\
& C_{\Sigma} \exp\left(-\gamma_1 l_{*}\right) \le e_{*}; \notag\\
& L^2 e^2 R \le 1, \notag
\end{align}
and the error terms of Section~\ref{sec:conditions} are
\begin{align*}
& \varepsilon_1 := \frac{1}{2}\, e_{*}^{\frac{1}{4} - 9\varepsilon}, \qquad \varepsilon_2 := C_{\Sigma}\, l_{*}^{-1} \exp\left(-\gamma_1 l_{*}\right),\\
& \varepsilon_3 := 6\, l_{*}^3 \exp\left(-\frac{p_{0*}}{2}\right) + 6\, l_{*}^3 \exp\left(-\frac{p_{*}^2}{2 C_0^2}\right), \qquad \varepsilon_4 := 6\, l_{*}^3 \exp\left(-\frac{p_{*}^2}{2 C_0^2}\right),\\
& \varepsilon_5(x) := -312 \log\left(1 - \exp\left(-\frac{x}{4}\right)\right) \qquad (x > 0).
\end{align*}
Under the standing hypothesis~\ref{def:hypothesis-flow} the conditions \eqref{eq:renormalisation-flow-3}, \eqref{eq:renormalisation-group-fl} and \eqref{eq:renormalisation-group-T1}--\eqref{eq:renormalisation-group-T6} hold, by Proposition~\ref{prop:consequences}.

\paragraph{The number $e_1$.}
For $v > 0$ put $\hat{l}(v) := L M v^r$, and recall $N^{*} = 8 L^3$ (stage~\ref{it:order-choice-5}). Let $\mathcal{E}$ be the set of the numbers $e' \in \left(0, \min\left(e_0, \frac{1}{3} L^{-\frac{m_0}{2}}\right)\right]$ such that for every $v \ge -\log e'$ the following inequalities hold:
\begin{align}
& v \ge 8, \qquad v^r > 31 L + 16, \qquad \left(1 + \frac{3 \mathrm{lg}}{v}\right)^r < L, \qquad p_0\, v^{p_0-1} \ge 28; \tag{e1}\label{eq:renormalisation-flow-e1}\\
& \left(c_{\mathrm{fl}} + 2\right) \exp(-\varepsilon v) \le \frac{1}{16} \left(1 - L^{-\frac{\varepsilon}{2}}\right); \tag{e2}\label{eq:renormalisation-flow-e2}\\
& 8\, C_0^2 \cdot \hat{l}(v)^6 \cdot \left(b_{-}^{-1} + L^{m_0+1} \cdot \hat{l}(v)\right) \le \exp\left(\frac{\varepsilon v}{2}\right), \notag\\
& 4 \left(2 C_0 + c_{\kappa} + 5 c_{\mathrm{fl}} + 6\right) \cdot \hat{l}(v)^3 \le \exp(\varepsilon v), \qquad \exp\left(-\left(\frac{1}{4} - 9\varepsilon\right) v\right) \le \frac{1}{8}; \tag{e3}\label{eq:renormalisation-flow-e3}\\
& 4\, C_0 \cdot \hat{l}(v)^2 \cdot v^p < \exp(\varepsilon v); \tag{e4}\label{eq:renormalisation-flow-e4}\\
& 12 \cdot \hat{l}(v)^3 \cdot \left( \exp\left(-\frac{v^{p_0}}{2}\right) + \exp\left(-\frac{v^{2p}}{2 C_0^2}\right) \right) \le \exp(-v); \tag{e5}\label{eq:renormalisation-flow-e5}\\
& C_{\Sigma} \exp\left(-\gamma_1 M v^r\right) \le \exp(-v); \tag{e6}\label{eq:renormalisation-flow-e6}\\
& L^2 \exp(-2v) \cdot \hat{l}(v) \le 1; \tag{e7}\label{eq:renormalisation-flow-e7}\\
& 6 N^{*} M v^r \exp(-v) \left( v \exp(-v) + \frac{1}{2} \right) \le \pi^2. \tag{e8}\label{eq:renormalisation-flow-e8}
\end{align}
If $e' \in \mathcal{E}$ and $0 < e'' < e'$, then $e'' \in \mathcal{E}$, because every $v \ge -\log e''$ satisfies $v \ge -\log e'$.

Each of the inequalities \eqref{eq:renormalisation-flow-e1}--\eqref{eq:renormalisation-flow-e8} holds for all sufficiently large $v$. In \eqref{eq:renormalisation-flow-e1}, the first inequality holds for $v \ge 8$; $v^r$ and $p_0\, v^{p_0-1}$ tend to infinity with $v$, as $r \ge 2$ and $p_0 \ge 2$; and $\left(1 + \frac{3 \mathrm{lg}}{v}\right)^r$ decreases to $1 < L$. In \eqref{eq:renormalisation-flow-e2}, $\exp(-\varepsilon v) \to 0$, whereas the right side is a positive number. In \eqref{eq:renormalisation-flow-e3}, \eqref{eq:renormalisation-flow-e4} and \eqref{eq:renormalisation-flow-e7}, the quotient of the left side by the right side is a polynomial in $v$ times $\exp(-c v)$ with a number $c > 0$; for the last inequality of \eqref{eq:renormalisation-flow-e3} one has $c = \frac{1}{4} - 9\varepsilon$, which is positive because $\varepsilon \le 1/200$ (stage~\ref{it:order-choice-8}). In \eqref{eq:renormalisation-flow-e8} this quotient is a sum of two terms of the same form, namely
\[
\frac{6 N^{*} M}{\pi^2} \left( v^{r+1} \exp(-2v) + \frac{1}{2}\, v^r \exp(-v) \right).
\]
All these quotients tend to $0$. In \eqref{eq:renormalisation-flow-e5},
\[
\hat{l}(v)^3 \exp\left(v - \frac{v^{p_0}}{2}\right) \to 0 \qquad \text{and} \qquad \hat{l}(v)^3 \exp\left(v - \frac{v^{2p}}{2 C_0^2}\right) \to 0,
\]
because $p_0 \ge 2$ and $2p \ge 2$. In \eqref{eq:renormalisation-flow-e6}, $\exp\left(v - \gamma_1 M v^r\right) \to 0$, because $r \ge 2$ and $\gamma_1 M > 0$.

Consequently there is a number $V > 0$ such that \eqref{eq:renormalisation-flow-e1}--\eqref{eq:renormalisation-flow-e8} hold for every $v \ge V$, and every $e'$ with $0 < e' \le \min\left(e_0, \frac{1}{3} L^{-\frac{m_0}{2}}, \exp(-V)\right)$ belongs to $\mathcal{E}$. So $\mathcal{E}$ is not empty, and the number
\[
e_1 := \frac{1}{2} \sup \mathcal{E}
\]
is positive. As $e_1 < \sup \mathcal{E}$, some element of $\mathcal{E}$ exceeds $e_1$, and therefore $e_1 \in \mathcal{E}$. Thus $0 < e_1 \le \min\left(e_0, \frac{1}{3} L^{-\frac{m_0}{2}}\right)$, and \eqref{eq:renormalisation-flow-e1}--\eqref{eq:renormalisation-flow-e8} hold for every $v \ge -\log e_1$.

In terms of the single variable $v$, the inequalities \eqref{eq:renormalisation-flow-e1}--\eqref{eq:renormalisation-flow-e7} express the conditions named in stage~\ref{it:order-choice-10}. The inequality \eqref{eq:renormalisation-flow-e8} serves to verify the condition~\eqref{eq:renormalisation-group-17}: for $e_K \le e_1$, Proposition~\ref{prop:consequences} deduces \eqref{eq:renormalisation-group-17} from \eqref{eq:renormalisation-flow-e8}, taken at $v = v_K$, under the further condition~\eqref{eq:renormalisation-group-18}, that is $e_K \log h' \le 1$. Apart from $v$ and absolute numbers, the inequalities \eqref{eq:renormalisation-flow-e1}--\eqref{eq:renormalisation-flow-e8} contain the constants $L$, $M$, $m_0$, $b_{-}$, $r$, $p_0$, $p$, $\varepsilon$, $c_{\mathrm{fl}}$, $c_{\kappa}$, $C_0$, $\gamma_1$, $C_{\Sigma}$, $N^{*}$ and no other quantity. The set $\mathcal{E}$, and with it $e_1$, is therefore determined by $e_0$ together with this list of constants; in particular $e_1$ takes the same value for every $\bar{m}$, $N$, $R$ and $\theta$. The number $e_1$ so defined is the one required by the condition~\eqref{eq:renormalisation-group-26} of Section~\ref{sec:conditions}.

\subsection{Scales and the stop level}\label{sec:scales-stop-level}

Let $e > 0$, and let $R = h' L^T$ with an integer $T \ge 0$. As in Section~\ref{sec:running-quantities-tori}, $e_K = e h'^{\frac{1}{2}}$, $v_K = -\log e_K$ and $\mathrm{lg} = \frac{1}{2} \log L$. For every integer $j \ge 0$, called the depth, put
\begin{align*}
\hat{a}_j &:= h' L^{-j}, & \hat{v}_j &:= v_K + j \cdot \mathrm{lg},\\
\hat{r}_j &:= \min\left\{ L^i : i \in \{0, 1, 2, \dots\},\ L^i \ge \hat{v}_j^r \right\}, & \hat{\mathfrak{n}}_j &:= T + j - \log_L\left( M \cdot \hat{r}_j \right).
\end{align*}
These four numbers are functions of $e h'^{\frac{1}{2}}$, $L$, $M$, $r$, $T$ and $j$, the number $\hat{a}_j$ also of $h'$; none of them depends on $N$.

Let $0 \le k \le K$. Since $a_k = L^{k-N}$ and $a_K = h'$, one has $a_k = h' L^{-(K-k)} = \hat{a}_{K-k}$. Hence $e_k = e a_k^{\frac{1}{2}} = e_K L^{-\frac{K-k}{2}}$ (for $k = 0$ this is the running charge $q$ of level $0$), so that
\[
v_k = v_K + (K-k) \cdot \mathrm{lg} = \hat{v}_{K-k}, \qquad r_k = \hat{r}_{K-k}.
\]
Moreover $l_k = R/a_k = L^{T+K-k}$, and therefore $\mathfrak{n}_k = \log_L\left( l_k/(M r_k) \right) = \hat{\mathfrak{n}}_{K-k}$.

\begin{proposition}[Depth and stop level]\label{prop:depth-stop-level}
Let $e > 0$ be such that $v_K \ge 1$ and $\left( 1 + \frac{\mathrm{lg}}{v_K} \right)^r \le L$, and let $R = h' L^T$ be an element of $\mathrm{Tor}(e)$.
\begin{enumerate}
\item[(a)] For every $j \ge 0$ one has $\hat{r}_{j+1} \in \left\{ \hat{r}_j, L \cdot \hat{r}_j \right\}$ and $\hat{\mathfrak{n}}_{j+1} - \hat{\mathfrak{n}}_j \in \{0, 1\}$. Consequently, for $0 \le k < K$, one has $r_{k+1} \in \left\{ r_k, r_k/L \right\}$ and $\mathfrak{n}_{k+1} - \mathfrak{n}_k \in \{-1, 0\}$.
\item[(b)] One has $\hat{\mathfrak{n}}_0 \le 0$, and $\hat{\mathfrak{n}}_j \to \infty$ as $j \to \infty$. Consequently the number
\[
j_{*} = j_{*}(R,e) := \max\left\{ j \ge 0 : \hat{\mathfrak{n}}_j = 0 \right\}
\]
exists; $\hat{\mathfrak{n}}_j \le 0$ for $0 \le j \le j_{*}$, and $\hat{\mathfrak{n}}_j \ge 1$ for $j > j_{*}$.
\item[(c)] One has $N_{\mathrm{min}}(R,e) = \max\left\{ 0, j_{*} + 1 - \log_L h' \right\}$. For every $N \ge N_{\mathrm{min}}(R,e)$ there is a first level $k$ with $\mathfrak{n}_k = 0$; it is $k^{*}(R) = K - j_{*}$, and $1 \le k^{*} \le K$. Hence the numbers
\begin{gather*}
K - k^{*} = j_{*}, \qquad a_{*} = h' L^{-j_{*}}, \qquad e_{*} = e a_{*}^{\frac{1}{2}},\\
l_{*} = M r_{*} = L^{T + j_{*}}, \qquad \mathfrak{m}_{*} = \bar{m} \cdot a_{*}
\end{gather*}
are the same for all $N \ge N_{\mathrm{min}}(R,e)$.
\item[(d)] Let $R' \in \mathrm{Tor}(e)$ with $R \le R'$. Then $j_{*}(R',e) \le j_{*}(R,e)$. Consequently $N_{\mathrm{min}}(R',e) \le N_{\mathrm{min}}(R,e)$, and $k^{*}(R) \le k^{*}(R')$ for every $N \ge N_{\mathrm{min}}(R,e)$.
\item[(e)] Let $N \ge N_{\mathrm{min}}(R,e)$. For $0 \le k \le k^{*}$ one has $\mathfrak{n}_k \ge 0$, that is $l_k \ge M r_k$, and
\[
r_k \ge r_{*}, \qquad v_k \ge v_{*} \ge v_K, \qquad e_k \le e_{*} \le e_K \le e_{*} \cdot l_{*}^{\frac{1}{2}};
\]
for $0 \le k < k^{*}$ one has $\mathfrak{n}_k \ge 1$. Further, the running mass $\bar{m} \cdot a_k$ of level $k$, which is $\mathfrak{m}_k$ for $k \ge 1$ and $\bar{m} \cdot \eta$ for $k = 0$, satisfies $\bar{m} \cdot a_{k+1} = L \cdot \bar{m} \cdot a_k$ for every level $k \ge 0$, and $0 < \bar{m} \cdot a_k \le \mathfrak{m}_{*} \le \mathfrak{m}_K$ for $0 \le k \le k^{*}$. With the upper end $\mathfrak{m}_0$ of the mass range (Definition~\ref{def:mass-bound-stop}) one has
\[
L^{-m_0-1} < \mathfrak{m}_K \le L^{-m_0} \le \mathfrak{m}_0, \qquad \mathfrak{m}_{*} \cdot l_{*} = \bar{m} \cdot R \ge \mathfrak{m}_K,
\]
hence $\mathfrak{m}_{*}^{-1} < L^{m_0+1} l_{*}$. Finally,
\[
M v_{*}^r \le l_{*} < L M v_{*}^r, \qquad e^2 R = e_{*}^2 \cdot l_{*}.
\]
\item[(f)] If condition~\eqref{eq:renormalisation-flow-3} holds, then $l_k \in \mathrm{Cl}$ for every $0 \le k \le k^{*}$.
\item[(g)] Assume that $\left( 1 + \frac{3 \mathrm{lg}}{v_K} \right)^r < L$. Then $j_{*}(R^{*},e) \in \{1, 2\}$ and $j_{*}(R_{\mathrm{top}},e) \in \{0, 1\}$. Consequently, for $N \ge N_{\mathrm{min}}(R^{*},e)$ one has $k^{*}(R^{*}) \in \{K-1, K-2\}$, and on the torus $R^{*}$ one has $a_{*} \in \left\{ h'/L, h'/L^2 \right\}$ and $r_{*} \in \left\{ r_K, L r_K \right\}$; for $N \ge N_{\mathrm{min}}(R_{\mathrm{top}},e)$ one has $k_{\mathrm{top}} \in \{K, K-1\}$.
\end{enumerate}
\end{proposition}

\begin{proof}
\emph{(a)} Since $\mathrm{lg} > 0$, one has $\hat{v}_{j+1} = \hat{v}_j + \mathrm{lg} > \hat{v}_j \ge v_K > 0$. Every power of $L$ that is at least $\hat{v}_{j+1}^r$ is therefore at least $\hat{v}_j^r$, and $\hat{r}_{j+1} \ge \hat{r}_j$. In the other direction, $\hat{v}_j \ge v_K$ gives
\[
\hat{v}_{j+1}^r = \hat{v}_j^r \left( 1 + \frac{\mathrm{lg}}{\hat{v}_j} \right)^r \le \hat{v}_j^r \left( 1 + \frac{\mathrm{lg}}{v_K} \right)^r \le L \cdot \hat{v}_j^r \le L \cdot \hat{r}_j.
\]
As $L \cdot \hat{r}_j$ is a power of $L$ with a non-negative integral exponent, the minimality in the definition of $\hat{r}_{j+1}$ yields $\hat{r}_{j+1} \le L \cdot \hat{r}_j$. Both $\hat{r}_j$ and $\hat{r}_{j+1}$ are powers of $L$; it follows that $\hat{r}_{j+1} \in \left\{ \hat{r}_j, L \cdot \hat{r}_j \right\}$ and
\[
\hat{\mathfrak{n}}_{j+1} - \hat{\mathfrak{n}}_j = 1 - \log_L \frac{\hat{r}_{j+1}}{\hat{r}_j} \in \{0, 1\}.
\]
For $0 \le k < K$ let $j := K - k - 1 \ge 0$. Then $r_k = \hat{r}_{j+1}$, $r_{k+1} = \hat{r}_j$, $\mathfrak{n}_k = \hat{\mathfrak{n}}_{j+1}$ and $\mathfrak{n}_{k+1} = \hat{\mathfrak{n}}_j$, and the assertions on $r_{k+1}$ and on $\mathfrak{n}_{k+1} - \mathfrak{n}_k$ follow.

\emph{(b)} By Definition~\ref{def:tori-stop-level}, $R \le R_{\mathrm{top}} = h' M r_K$. Since $\hat{v}_0 = v_K$, one has $\hat{r}_0 = r_K$, and the inequality $R \le R_{\mathrm{top}}$ reads $L^T \le M \cdot \hat{r}_0$, that is, $\hat{\mathfrak{n}}_0 \le 0$. Because $M = L^m$ (Section~\ref{sec:lattice-geometry}) and every $\hat{r}_j$ is a power of $L$ with a non-negative integral exponent, all the numbers $\hat{\mathfrak{n}}_j$ are integers. As $\hat{v}_j \ge v_K \ge 1$, the minimality in the definition of $\hat{r}_j$ gives $\hat{r}_j \le L \cdot \hat{v}_j^r$, a bound that grows polynomially in $j$. Consequently
\[
\hat{\mathfrak{n}}_j \ge T + j - \log_L\left( L M \cdot \hat{v}_j^r \right) \longrightarrow \infty \qquad (j \to \infty),
\]
since $\log_L\left( L M \cdot \hat{v}_j^r \right) = 1 + \log_L M + r \log_L\left( v_K + j \cdot \mathrm{lg} \right)$ grows only logarithmically in $j$. Thus $(\hat{\mathfrak{n}}_j)_{j \ge 0}$ is a sequence of integers which begins with a value $\le 0$, tends to infinity, and, by (a), increases by $0$ or by $1$ from each index to the next. Such a sequence attains the value $0$, and only at finitely many indices; it is non-decreasing, so it is $\le 0$ at every index up to the largest index $j_{*}$ at which it vanishes, and $\ge 1$ at every larger index.

\emph{(c)} One has $R L^N = R/a_0 = l_0$. So the inequality $R L^N \ge L M r_0$ of Definition~\ref{def:tori-stop-level} means $\mathfrak{n}_0 \ge 1$, that is, $\hat{\mathfrak{n}}_K \ge 1$; by (b) this is equivalent to $K > j_{*}$, i.e.\ to $N + \log_L h' \ge j_{*} + 1$. In Definition~\ref{def:tori-stop-level}, $N$ and $N_1$ are integers $\ge 0$ with $N + \log_L h' \ge 0$, respectively $N_1 + \log_L h' \ge 0$. Hence the inequality holds for all such $N \ge N_1$ if and only if $N_1 \ge 0$ and $N_1 + \log_L h' \ge j_{*} + 1$. Since $\log_L h'$ is an integer, the least such $N_1$ is $\max\left\{ 0, j_{*} + 1 - \log_L h' \right\}$, which is the asserted value of $N_{\mathrm{min}}(R,e)$.

Now let $N \ge N_{\mathrm{min}}(R,e)$, so that $K > j_{*}$. By (b), $\mathfrak{n}_k = \hat{\mathfrak{n}}_{K-k} \ge 1$ for $0 \le k < K - j_{*}$, and $\mathfrak{n}_{K-j_{*}} = \hat{\mathfrak{n}}_{j_{*}} = 0$. Hence the first level with $\mathfrak{n}_k = 0$ is $k^{*} = K - j_{*}$, and $1 \le k^{*} \le K$ because $0 \le j_{*} \le K - 1$. The equation $\mathfrak{n}_{k^{*}} = 0$ says $l_{*} = M r_{*}$; moreover $a_{*} = \hat{a}_{j_{*}} = h' L^{-j_{*}}$ and $l_{*} = L^{T + K - k^{*}} = L^{T + j_{*}}$. As $j_{*}$ does not depend on $N$, neither do $a_{*}$, $e_{*} = e a_{*}^{\frac{1}{2}}$, $l_{*}$ and $\mathfrak{m}_{*} = \bar{m} \cdot a_{*}$.

\emph{(d)} Write $R' = h' L^{T'}$, so that $T \le T'$, and let $\hat{\mathfrak{n}}_j(R)$, $\hat{\mathfrak{n}}_j(R')$ denote the numbers formed with $T$ and with $T'$. The numbers $\hat{r}_j$ are the same for both tori; hence
\[
\hat{\mathfrak{n}}_j(R') = \hat{\mathfrak{n}}_j(R) + T' - T = \hat{\mathfrak{n}}_j(R) + \log_L \frac{R'}{R} \ge \hat{\mathfrak{n}}_j(R).
\]
For $j > j_{*}(R,e)$, (b) gives $\hat{\mathfrak{n}}_j(R') \ge \hat{\mathfrak{n}}_j(R) \ge 1$; therefore $j_{*}(R',e) \le j_{*}(R,e)$. The two other assertions follow from the formulas $N_{\mathrm{min}}(R,e) = \max\left\{ 0, j_{*}(R,e) + 1 - \log_L h' \right\}$ and $k^{*}(R) = K - j_{*}(R,e)$ of (c), applied to $R$ and to $R'$, in which $h'$ and $K$ are the same for both tori.

\emph{(e)} Let $0 \le k \le k^{*}$. Then $K - k \ge K - k^{*} = j_{*}$, and (b), (c) give $\mathfrak{n}_k = \hat{\mathfrak{n}}_{K-k} \ge 0$, with $\mathfrak{n}_k \ge 1$ if $k < k^{*}$; the inequality $\mathfrak{n}_k \ge 0$ means $l_k \ge M r_k$. The sequences $(\hat{r}_j)$ and $(\hat{v}_j)$ are non-decreasing in $j$, by (a) and by the definition of $\hat{v}_j$, and $K - k^{*} = j_{*} \ge 0$ by (c). Hence
\[
r_k = \hat{r}_{K-k} \ge \hat{r}_{j_{*}} = r_{*}, \qquad v_k = \hat{v}_{K-k} \ge \hat{v}_{j_{*}} = v_{*} \ge \hat{v}_0 = v_K,
\]
and $e_k = \exp(-v_k) \le e_{*} \le e_K$. Further, $e_K/e_{*} = \left( a_K/a_{*} \right)^{\frac{1}{2}} = \left( h'/a_{*} \right)^{\frac{1}{2}}$, while $l_{*} = R/a_{*} \ge h'/a_{*}$ because $R = h' L^T \ge h'$; thus $e_K \le e_{*} \cdot l_{*}^{\frac{1}{2}}$.

The running mass of level $k$ is the number $\bar{m} \cdot a_k$ (Section~\ref{sec:running-quantities-tori}); as $1 \le k^{*} \le K$ by (c), one has $\mathfrak{m}_{*} = \bar{m} \cdot a_{k^{*}}$ and $\mathfrak{m}_K = \bar{m} \cdot a_K$. For every level $k \ge 0$ one has $a_{k+1} = L a_k > 0$, whence $\bar{m} \cdot a_{k+1} = L \cdot \bar{m} \cdot a_k$ and $\bar{m} \cdot a_k > 0$. The sequence $(a_k)$ is therefore increasing, and $k \le k^{*} \le K$ gives $0 < \bar{m} \cdot a_k \le \bar{m} \cdot a_{k^{*}} = \mathfrak{m}_{*} \le \bar{m} \cdot a_K = \mathfrak{m}_K$ for $0 \le k \le k^{*}$. Let $t := \log_L(1/\bar{m})$, a number $\ge 0$. Then $h = L^{\lfloor t \rfloor}$, so that $\bar{m} \cdot h = L^{\lfloor t \rfloor - t}$ lies in $(1/L, 1]$, because $\lfloor t \rfloor - t \in (-1, 0]$. Since $\mathfrak{m}_K = \bar{m} \cdot a_K = \bar{m} \cdot h' = \bar{m} \cdot h \cdot L^{-m_0}$, it follows that $L^{-m_0-1} < \mathfrak{m}_K \le L^{-m_0}$; and $L^{-m_0} \le \mathfrak{m}_0$ for the upper end $\mathfrak{m}_0$ of the mass range, by the choice of $m_0$ in Definition~\ref{def:mass-bound-stop} (the least integer $j \ge 0$ with $L^{-j} \le \mathfrak{m}_0$). Next,
\[
\mathfrak{m}_{*} \cdot l_{*} = \bar{m} \cdot a_{*} \cdot \frac{R}{a_{*}} = \bar{m} \cdot R \ge \bar{m} \cdot h' = \mathfrak{m}_K,
\]
as $R \ge h'$. Consequently $\mathfrak{m}_{*}^{-1} = l_{*}/(\bar{m} \cdot R) \le l_{*}/\mathfrak{m}_K < L^{m_0+1} l_{*}$.

By the definition of $r_{*}$ one has $r_{*} \ge v_{*}^r$; and $r_{*} < L v_{*}^r$ because $v_{*} \ge 1$: if $r_{*} = L^i$ with $i \ge 1$, the minimality in the definition gives $L^{i-1} < v_{*}^r$, and if $r_{*} = 1$, then $r_{*} < L \le L v_{*}^r$. Multiplying by $M$ and using $l_{*} = M r_{*}$ from (c), one obtains $M v_{*}^r \le l_{*} < L M v_{*}^r$. Finally, $e_{*}^2 \cdot l_{*} = e^2 a_{*} \cdot R/a_{*} = e^2 R$.

\emph{(f)} By stage~\ref{it:order-choice-4}, $L$ is odd and $L \ge 33$; hence $31 L + 16 < L^2$. Since also $31 L + 16 > L$, a power of $L$ with a non-negative integral exponent which exceeds $31 L + 16$ is at least $L^2$. Condition~\eqref{eq:renormalisation-flow-3} says $r_{*} > 31 L + 16$, and therefore $r_{*} \ge L^2$. Let $0 \le k \le k^{*}$. By the definition of $\mathfrak{n}_k$,
\[
\frac{l_k}{L M} = r_k L^{\mathfrak{n}_k - 1},
\]
where $r_k \ge r_{*} \ge L^2$ is a power of $L$ and $\mathfrak{n}_k \ge 0$ is an integer, by (e) and (b). Hence $l_k/(L M) = L^i$ with an integer $i \ge 1$; this is an odd integer, $L$ being odd, and it is at least $L \ge 13$. So $l_k \in \mathrm{Cl}$ (Section~\ref{sec:lattice-geometry}).

\emph{(g)} Write $M = L^m$ (Section~\ref{sec:lattice-geometry}) and $r_K = L^i$ with an integer $i \ge 0$. By Definition~\ref{def:tori-stop-level}, $R_{\mathrm{top}} = h' M r_K = h' L^{m+i}$ and $R^{*} = R_{\mathrm{top}}/L = h' L^{m+i-1}$. Stage~\ref{it:order-choice-6} gives $M \ge L^2$, that is $m \ge 2$; so the exponents $m + i$ and $m + i - 1$ are non-negative integers, and $R^{*} < R_{\mathrm{top}}$. Hence $R^{*}$ and $R_{\mathrm{top}}$ belong to $\mathrm{Tor}(e)$, and (a)--(c) apply to both of them.

For $j \in \{1, 2, 3\}$ one has $\left( 1 + \frac{j \cdot \mathrm{lg}}{v_K} \right)^r \le \left( 1 + \frac{3 \mathrm{lg}}{v_K} \right)^r < L$, and therefore
\[
\hat{v}_j^r = v_K^r \left( 1 + \frac{j \cdot \mathrm{lg}}{v_K} \right)^r < L v_K^r \le L \cdot \hat{r}_0.
\]
As $L \cdot \hat{r}_0$ is a power of $L$ and $\hat{r}_j \ge \hat{r}_0$ by (a), it follows that $\hat{r}_j \in \left\{ \hat{r}_0, L \cdot \hat{r}_0 \right\}$ for $j = 1, 2, 3$.

On the torus $R^{*}$ one has $L^T = R^{*}/h' = M \cdot \hat{r}_0/L$, since $\hat{r}_0 = r_K$; hence
\[
\hat{\mathfrak{n}}_j = j - 1 - \log_L \frac{\hat{r}_j}{\hat{r}_0} \qquad (j \ge 0).
\]
Consequently $\hat{\mathfrak{n}}_0 = -1$, $\hat{\mathfrak{n}}_1 \in \{0, -1\}$, $\hat{\mathfrak{n}}_2 \in \{1, 0\}$ and $\hat{\mathfrak{n}}_3 \in \{2, 1\}$. The sequence $(\hat{\mathfrak{n}}_j)$ is non-decreasing by (a); so $\hat{\mathfrak{n}}_j \ge 1$ for $j \ge 3$, and $j_{*}(R^{*},e) \le 2$. If $\hat{\mathfrak{n}}_2 = 0$, then $j_{*}(R^{*},e) = 2$. Otherwise $\hat{\mathfrak{n}}_2 = 1$; then (a) gives $\hat{\mathfrak{n}}_1 \in \{0, 1\}$, which together with $\hat{\mathfrak{n}}_1 \in \{0, -1\}$ forces $\hat{\mathfrak{n}}_1 = 0$, and $j_{*}(R^{*},e) = 1$.

On the torus $R_{\mathrm{top}} = L R^{*}$ every $\hat{\mathfrak{n}}_j$ is larger by one than on $R^{*}$, as in the proof of (d). Thus $\hat{\mathfrak{n}}_0 = 0$, $\hat{\mathfrak{n}}_1 \in \{1, 0\}$ and $\hat{\mathfrak{n}}_j \ge \hat{\mathfrak{n}}_2 \ge 1$ for $j \ge 2$, whence $j_{*}(R_{\mathrm{top}},e) \in \{0, 1\}$.

The remaining assertions follow from (c), where $r_{*} = r_{k^{*}} = \hat{r}_{K - k^{*}} = \hat{r}_{j_{*}}$. On $R^{*}$: $k^{*}(R^{*}) = K - j_{*}(R^{*},e) \in \{K-1, K-2\}$, $a_{*} = h' L^{-j_{*}} \in \left\{ h'/L, h'/L^2 \right\}$, and $r_{*} = \hat{r}_{j_{*}} \in \left\{ \hat{r}_0, L \cdot \hat{r}_0 \right\} = \left\{ r_K, L r_K \right\}$. On $R_{\mathrm{top}}$: $k_{\mathrm{top}} = K - j_{*}(R_{\mathrm{top}},e) \in \{K, K-1\}$.
\end{proof}

\begin{proposition}[Consequences of $e_K \le e_1$]\label{prop:consequences}
Let $e > 0$ satisfy $e_K = e h'^{\frac{1}{2}} \le e_1$, with the number $e_1$ of Section~\ref{sec:constants-conditions-torus}. Then
\[
e \le \frac{1}{3}, \qquad v_K \ge 8, \qquad \left( 1 + \frac{\mathrm{lg}}{v_K} \right)^r \le \left( 1 + \frac{3 \mathrm{lg}}{v_K} \right)^r < L,
\]
and for every $R \in \mathrm{Tor}(e)$ and every $N \ge N_{\mathrm{min}}(R,e)$ the conditions \eqref{eq:renormalisation-flow-3}, \eqref{eq:renormalisation-group-fl} and \eqref{eq:renormalisation-group-T1}--\eqref{eq:renormalisation-group-T6} hold. In terms of $e_K$ and $h'$, condition~\eqref{eq:renormalisation-group-17}, that is $6 N^{*} \cdot e^2 R^{*} \log \frac{1}{e} \le \pi^2$, is the inequality
\[
6 N^{*} \cdot \frac{M r_K}{L} \cdot e_K^2 \left( v_K + \frac{1}{2} \log h' \right) \le \pi^2.
\]
If moreover $e_K \log h' \le 1$, which is condition~\eqref{eq:renormalisation-group-18}, then condition~\eqref{eq:renormalisation-group-17} holds. In particular, \eqref{eq:renormalisation-group-17} holds if $e h' \le 1$.
\end{proposition}

\begin{proof}
Since $0 < \bar{m} \le 1$, the exponent $\lfloor \log_L(1/\bar{m}) \rfloor$ is $\ge 0$; so $h \ge 1$ and $h' = h L^{-m_0} \ge L^{-m_0}$. By its definition in Section~\ref{sec:constants-conditions-torus}, $e_1$ is one half of a supremum of numbers in the interval $\left( 0, \min\left( e_0, L^{-\frac{m_0}{2}}/3 \right) \right]$; in particular $e_1 \le L^{-\frac{m_0}{2}}/3$. Therefore
\[
e = e_K h'^{-\frac{1}{2}} \le e_1 L^{\frac{m_0}{2}} \le \frac{1}{3}.
\]
From $e_K \le e_1$ one gets $v_K = -\log e_K \ge -\log e_1$; as the inequalities \eqref{eq:renormalisation-flow-e1}--\eqref{eq:renormalisation-flow-e8} hold for every $v \ge -\log e_1$ (Section~\ref{sec:constants-conditions-torus}), they hold at $v = v_K$. By \eqref{eq:renormalisation-flow-e1} at $v = v_K$, $v_K \ge 8$ and $(1 + 3 \mathrm{lg}/v_K)^r < L$; and $(1 + \mathrm{lg}/v_K)^r \le (1 + 3 \mathrm{lg}/v_K)^r$ because $\mathrm{lg} > 0$. This proves the first three assertions, and it shows that every $R \in \mathrm{Tor}(e)$ satisfies the hypotheses of Proposition~\ref{prop:depth-stop-level}.

Let $R \in \mathrm{Tor}(e)$ and $N \ge N_{\mathrm{min}}(R,e)$. By Proposition~\ref{prop:depth-stop-level}(c), the stop level $k^{*} = k^{*}(R)$ exists and $1 \le k^{*} \le K$, so that the numbers with the index $*$ are defined. By Proposition~\ref{prop:depth-stop-level}(e), $v_{*} \ge v_K \ge -\log e_1$, so \eqref{eq:renormalisation-flow-e1}--\eqref{eq:renormalisation-flow-e8} hold also at $v = v_{*}$. With $\hat{l}(v) = L M v^r$ as in Section~\ref{sec:constants-conditions-torus}, Proposition~\ref{prop:depth-stop-level}(e) gives
\[
M v_{*}^r \le l_{*} < \hat{l}(v_{*}), \qquad \mathfrak{m}_{*}^{-1} < L^{m_0+1} \cdot \hat{l}(v_{*}),
\]
and by definition $e_{*} = \exp(-v_{*})$, $p_{*} = v_{*}^p$ and $p_{0*} = v_{*}^{p_0}$. The conditions are now verified one at a time.

\emph{Condition~\eqref{eq:renormalisation-flow-3}.} One has $r_{*} \ge v_{*}^r > 31 L + 16$, the second inequality by \eqref{eq:renormalisation-flow-e1} at $v = v_{*}$.

\emph{Condition~\eqref{eq:renormalisation-group-fl}.} For $0 \le k \le k^{*}$ one has $e_k = e_{*} L^{-\frac{k^{*} - k}{2}}$. Hence
\[
\sum_{k=0}^{k^{*}} e_k^{\varepsilon} = e_{*}^{\varepsilon} \sum_{i=0}^{k^{*}} L^{-\frac{i \varepsilon}{2}} \le e_{*}^{\varepsilon} \left( 1 - L^{-\frac{\varepsilon}{2}} \right)^{-1},
\]
and \eqref{eq:renormalisation-flow-e2} at $v = v_{*}$ gives, with $\mathfrak{s}_k = \left( c_{\mathrm{fl}} + 2 \right) e_k^{\varepsilon}$,
\[
\sum_{k=0}^{k^{*}} \mathfrak{s}_k \le \left( c_{\mathrm{fl}} + 2 \right) \exp(-\varepsilon v_{*}) \left( 1 - L^{-\frac{\varepsilon}{2}} \right)^{-1} \le \frac{1}{16}.
\]
The inequality $v_{*} \ge 8$ is contained in \eqref{eq:renormalisation-flow-e1} at $v = v_{*}$.

\emph{Condition~\eqref{eq:renormalisation-group-T1}.} The inequalities $v_{*} \ge 8$ and $p_0 v_{*}^{p_0-1} \ge 28$ are contained in \eqref{eq:renormalisation-flow-e1} at $v = v_{*}$. Proposition~\ref{prop:depth-stop-level}(e) at $k = 0$, where the running charge is $q$, gives $q \le e_{*} = \exp(-v_{*}) < 1$; since $\frac{3}{4} - 8 \varepsilon > 0$, it follows that $w^{\mathrm{m}}_0 = q^{\frac{3}{4} - 8 \varepsilon} < 1$.

\emph{Condition~\eqref{eq:renormalisation-group-T2}.} By \eqref{eq:renormalisation-flow-e3} at $v = v_{*}$, together with the bounds for $l_{*}$ and $\mathfrak{m}_{*}^{-1}$ displayed above,
\begin{align*}
8 C_0^2 \cdot l_{*}^6 \left( b_{-}^{-1} + \mathfrak{m}_{*}^{-1} \right) &\le 8 C_0^2 \cdot \hat{l}(v_{*})^6 \left( b_{-}^{-1} + L^{m_0+1} \cdot \hat{l}(v_{*}) \right)\\
&\le \exp\left( \frac{\varepsilon v_{*}}{2} \right) = e_{*}^{-\frac{\varepsilon}{2}},\\
4 \left( 2 C_0 + c_{\kappa} + 5 c_{\mathrm{fl}} + 6 \right) l_{*}^3 &\le 4 \left( 2 C_0 + c_{\kappa} + 5 c_{\mathrm{fl}} + 6 \right) \cdot \hat{l}(v_{*})^3\\
&\le \exp(\varepsilon v_{*}) = e_{*}^{-\varepsilon},\\
e_{*}^{\frac{1}{4} - 9 \varepsilon} &= \exp\left( -\left( \frac{1}{4} - 9 \varepsilon \right) v_{*} \right) \le \frac{1}{8}.
\end{align*}

\emph{Condition~\eqref{eq:renormalisation-group-T3}.} By \eqref{eq:renormalisation-flow-e4} at $v = v_{*}$, $4 C_0 \cdot l_{*}^2 \cdot p_{*} \le 4 C_0 \cdot \hat{l}(v_{*})^2 \cdot v_{*}^p < \exp(\varepsilon v_{*}) = e_{*}^{-\varepsilon}$.

\emph{Condition~\eqref{eq:renormalisation-group-T4}.} By \eqref{eq:renormalisation-flow-e5} at $v = v_{*}$,
\begin{align*}
& 12 l_{*}^3 \left( \exp\left( -\frac{p_{0*}}{2} \right) + \exp\left( -\frac{p_{*}^2}{2 C_0^2} \right) \right)\\
& \qquad \le 12 \cdot \hat{l}(v_{*})^3 \left( \exp\left( -\frac{v_{*}^{p_0}}{2} \right) + \exp\left( -\frac{v_{*}^{2p}}{2 C_0^2} \right) \right)\\
& \qquad \le \exp(-v_{*}) = e_{*}.
\end{align*}

\emph{Condition~\eqref{eq:renormalisation-group-T5}.} Since $l_{*} \ge M v_{*}^r$ and $\gamma_1 > 0$, \eqref{eq:renormalisation-flow-e6} at $v = v_{*}$ gives
\[
C_{\Sigma} \exp(-\gamma_1 l_{*}) \le C_{\Sigma} \exp\left( -\gamma_1 M v_{*}^r \right) \le \exp(-v_{*}) = e_{*}.
\]

\emph{Condition~\eqref{eq:renormalisation-group-T6}.} By Proposition~\ref{prop:depth-stop-level}(e) and by \eqref{eq:renormalisation-flow-e7} at $v = v_{*}$,
\[
L^2 e^2 R = L^2 e_{*}^2 \cdot l_{*} \le L^2 \exp(-2 v_{*}) \cdot \hat{l}(v_{*}) \le 1.
\]

\emph{Condition~\eqref{eq:renormalisation-group-17}.} By Definition~\ref{def:tori-stop-level}, $R^{*} = R_{\mathrm{top}}/L = h' M r_K/L$, so that $e^2 R^{*} = e_K^2 M r_K/L$; and
\[
\log \frac{1}{e} = -\log\left( e_K h'^{-\frac{1}{2}} \right) = v_K + \frac{1}{2} \log h',
\]
a positive number, as $e \le \frac{1}{3}$. The product $6 N^{*} \cdot e^2 R^{*} \log \frac{1}{e}$ is therefore the left side of the inequality displayed in the statement. Assume now that $e_K \log h' \le 1$. Since $v_K \ge 1$, one has $r_K < L v_K^r$, by the argument given for $r_{*}$ in the proof of Proposition~\ref{prop:depth-stop-level}(e); hence $e^2 R^{*} \le e_K^2 M v_K^r$. Recalling $e_K = \exp(-v_K)$, one obtains from \eqref{eq:renormalisation-flow-e8} at $v = v_K$
\begin{align*}
6 N^{*} \cdot e^2 R^{*} \log \frac{1}{e} &\le 6 N^{*} \cdot M v_K^r \cdot e_K \left( e_K v_K + \frac{1}{2} e_K \log h' \right)\\
&\le 6 N^{*} \cdot M v_K^r \cdot e_K \left( e_K v_K + \frac{1}{2} \right) \le \pi^2.
\end{align*}

Finally, let $e h' \le 1$. If $h' \le 1$, then $\log h' \le 0$, and $e_K \log h' \le 0 \le 1$. If $h' \ge 1$, then $\log h' \le h'^{\frac{1}{2}}$, because $\log x \le x^{\frac{1}{2}}$ for every $x \ge 1$; so $e_K \log h' = e h'^{\frac{1}{2}} \log h' \le e h' \le 1$. In both cases condition~\eqref{eq:renormalisation-group-18} holds, and condition~\eqref{eq:renormalisation-group-17} follows from what has just been shown.
\end{proof}

\subsection{Projection of primitives}\label{sec:projection-primitives}

\subsubsection{Two tori}\label{sec:two-tori}

Let $k \ge 0$ be a level, and let $l_1, l_2 \in \mathrm{Cl}$ with $l_2 = J \cdot l_1$, where $J$ is an odd integer. The constants satisfy
\[
M \ge L^2, \qquad M > 4L, \qquad \gamma_w \cdot M \ge 1;
\]
these inequalities follow from stage~\ref{it:order-choice-6}. For $i = 1, 2$ let
\[
T_i := \mathbb{L}_k/l_i\mathbb{Z}^3, \qquad T^0_i := \mathbb{Z}^3/l_i\mathbb{Z}^3
\]
be the fine lattice of level $k$ and side $l_i$ and its unit lattice, and let $\mathrm{pr}_i \colon \mathbb{L}_k \to T_i$ be the canonical projection. As $l_2\mathbb{Z}^3 \subset l_1\mathbb{Z}^3$, there is a covering map $\mathrm{pj} \colon T_2 \to T_1$ with $\mathrm{pj} \circ \mathrm{pr}_2 = \mathrm{pr}_1$. Both tori carry the same charge $e_k \in (0, e_0]$ and the same mass $\mathfrak{m} \in (0, \mathfrak{m}_0]$. Hence, for $i = 1, 2$, the quadruple $(k, l_i, \mathfrak{m}, e_k)$ consists of data of the scale free setting of Section~\ref{sec:properties-scheme}.

\paragraph{Periodic and antiperiodic objects.}
Under the identifications of Definition~\ref{def:periodisation-projection}, a function on the cells of $T_i$ (the sites and bonds of $T_i$ and of $T^0_i$), an element of $V_k(T_i)$ or a parameter family on $T_i$ is the same as an $l_i\mathbb{Z}^3$-periodic function on the corresponding cells or parameter cubes of $\mathbb{L}_k$, and a spinor on $T_i$ is the same as a function $f$ on the corresponding sites of $\mathbb{L}_k$ with $f(x + l_i \cdot e_{\nu}) = -f(x)$ for $\nu = 0, 1, 2$. Because $l_1$ is a multiple of $L \cdot M$, the unit lattice, the $L$-blocks, the $M$-cubes, the $LM$-cubes, the parameter cubes of each kind and the basis $\{\Upsilon\}$ are invariant under the translations by the elements of $l_1\mathbb{Z}^3$. Each object on $T_1$ is also an object on $T_2$, of the same type: an $l_1\mathbb{Z}^3$-periodic function is $l_2\mathbb{Z}^3$-periodic, and a function with $f(x + l_1 \cdot e_{\nu}) = -f(x)$ for all $\nu$ satisfies
\[
f(x + l_2 \cdot e_{\nu}) = (-1)^J f(x) = -f(x),
\]
$J$ being odd. The resulting inclusion map is denoted by $\mathrm{pj}^{*}$.

For $g \in l_1\mathbb{Z}^3$ and an object $f$ on $T_2$ let $(T_g f)(x) := f(x - g)$, where $x - g$ is the translate of the cell or parameter cube $x$ of $\mathbb{L}_k$ by $-g$. The character of a type of objects is the map $\boldsymbol{\chi} \colon \mathbb{Z}^3 \to \{1, -1\}$ given by $\boldsymbol{\chi}(w) := 1$ for scalar functions, bond functions, parameter families and elements of $V_k$, and by $\boldsymbol{\chi}(w) := (-1)^{w_0 + w_1 + w_2}$ for spinors. For spinors, the relations $f(x + l_1 \cdot e_{\nu}) = -f(x)$, $\nu = 0, 1, 2$, are equivalent to $f(x + l_1 \cdot w) = (-1)^{w_0 + w_1 + w_2} f(x)$ for all $w \in \mathbb{Z}^3$. Hence the image of $\mathrm{pj}^{*}$ is
\begin{align*}
\mathrm{Per} &:= \{ f \text{ on } T_2 : f(x + l_1 \cdot w) = \boldsymbol{\chi}(w) \cdot f(x) \text{ for all } w \in \mathbb{Z}^3 \} \\
&= \{ f \text{ on } T_2 : T_{-l_1 \cdot w} f = \boldsymbol{\chi}(w) \cdot f \text{ for all } w \in \mathbb{Z}^3 \}.
\end{align*}

The map $\mathrm{pj}^{*}$ is linear and injective, and $\lvert \mathrm{pj}^{*} f \rvert_{\infty} = \lvert f \rvert_{\infty}$, since $\mathrm{pj}^{*} f$ and $f$ are one and the same function on $\mathbb{L}_k$. On the fine lattice and on the unit lattice of a torus consider the bilinear pairings of Section~\ref{sec:fermion-integrals-blocking},
\[
\langle f, f' \rangle := \sum_x \eta_k^3 \sum_{\boldsymbol{\beta}} f_{\boldsymbol{\beta}}(x) \cdot f'_{\boldsymbol{\beta}}(x), \qquad \langle f, f' \rangle := \sum_y \sum_{\boldsymbol{\beta}} f_{\boldsymbol{\beta}}(y) \cdot f'_{\boldsymbol{\beta}}(y),
\]
where $x$ ranges over the sites of the fine lattice of the torus, $y$ over the sites of its unit lattice, and $\boldsymbol{\beta}$ over the components. For two objects $f$, $f'$ on $T_1$ with the same character one has
\[
\langle \mathrm{pj}^{*} f, \mathrm{pj}^{*} f' \rangle_{T_2} = J^3 \cdot \langle f, f' \rangle_{T_1}.
\]
Indeed, under the shift $x \mapsto x + l_1 \cdot w$, $w \in \mathbb{Z}^3$, both $f$ and $f'$ acquire the factor $\boldsymbol{\chi}(w) \in \{1, -1\}$; so every summand of the pairing is an $l_1\mathbb{Z}^3$-periodic function of the site. Moreover, a fundamental domain of $l_2\mathbb{Z}^3$ is the union of $J^3$ translates of a fundamental domain of $l_1\mathbb{Z}^3$ by elements of $l_1\mathbb{Z}^3$. For a parameter family $s$ on $T_1$ and a parameter cube $c$ of $T_2$ one has $(\mathrm{pj}^{*} s)_c = s_{\mathrm{pj}(c)}$. Consequently $\mathrm{pj}^{*}$ maps $\mathrm{Pd}_k(T_1)$ into $\mathrm{Pd}_k(T_2)$, and interior points of $\mathrm{Pd}_k(T_1)$ to interior points of $\mathrm{Pd}_k(T_2)$. Lemma~\ref{lem:periodic-extension} gives $\mathrm{pj}^{*}(\tilde{G}_k(T_1)) \subset \tilde{G}_k(T_2)$.

\paragraph{Indices and kernels.}
The indices (sites, bonds and basis elements), the support $\operatorname{supp} x$ of an index $x$, and the relations ``$x$ meets $W$'' and ``$x$ is contained in $W$'' for a set $W$ of sites are those of Section~\ref{sec:lattice-geometry} (see also Section~\ref{sec:notation}). Let $\square(x)$ denote the union of the $M$-cubes that $x$ meets. Every support has sup-diameter at most $L$, and $\square(x)$ is either one $M$-cube or the union of two $M$-cubes with a common face. (By fact~\eqref{it:renormalisation-group-g0} of Section~\ref{sec:cubes-parameter-cubes}, the $M$-cubes partition the sites, and every $L$-block lies in one $M$-cube. A site meets exactly one $M$-cube. The two end points of a bond differ in one coordinate only, and there by at most $1 < M$; so their $M$-cubes coincide or have a common face. By Proposition~\ref{prop:basis-map-box}, the support of a basis element either lies in one $L$-block, a set of sup-diameter $L - 1$ contained in one $M$-cube, or consists of the two end points of one bond.)

Let $O$ be an operator on $T_i$ from one space of functions to another, the two spaces having the same character $\boldsymbol{\chi}$. Its kernel is the function $O(x, x')$ of two indices $x$, $x'$ of $\mathbb{L}_k$ such that
\begin{gather*}
O(x + l_i \cdot e_{\nu}, x') = O(x, x' + l_i \cdot e_{\nu}) = \boldsymbol{\chi}(e_{\nu}) \cdot O(x, x') \quad (\nu = 0, 1, 2), \\
(O f)(x) = \sum_{x' \in \mathscr{F}_i} O(x, x') \cdot f(x'),
\end{gather*}
where $\mathscr{F}_i$ is a set of representatives of the column indices of $\mathbb{L}_k$ modulo $l_i\mathbb{Z}^3$; the factors $\eta_k^3$ that weight sums over the fine lattice are part of the kernel. For $g \in l_1\mathbb{Z}^3$ the operators $T_g O$ and $O T_g$ have the kernels $O(x - g, x')$ and $O(x, x' + g)$. Hence $T_g O = O T_g$ holds if and only if $O(x + g, x' + g) = O(x, x')$ for all indices $x$, $x'$.

\paragraph{Kinds.}
The family~\eqref{it:primitives-p1} has the kind $\pi'_k$ and bonds as indices: its row indices are the bonds of $T_k$, its column indices the bonds of $T^0_k$. The family~\eqref{it:primitives-p2} has the kind $\pi''_k$ and basis elements as indices. The family~\eqref{it:primitives-p3} and the family $\Gamma_k(s'', A)$ of~\eqref{it:primitives-p4} have the kind $\pi_k$ and sites as indices. Accordingly, an index belongs to the kind $\pi'_k$ if it is a bond, to the kind $\pi''_k$ if it is a basis element, and to the kind $\pi_k$ if it is a site. By Definition~\ref{def:parameter-cubes-polydisks}, the cubes of the kinds $\pi'_k$ and $\pi''_k$ are closed sets.

\subsubsection{Cover cubes}\label{sec:cover-cubes}

Let $l \in \mathrm{Cl}$ and $T_k = \mathbb{L}_k/l\mathbb{Z}^3$. Consider a set $\mathscr{Y}$ of parameter cubes on $T_k$, all of the same kind, and denote by $Y$ the union of its cubes. A link-component of $Y$ is the union of the cubes of a link-component of $\mathscr{Y}$ (Definition~\ref{def:parameter-cubes-polydisks}). For an index $x$, the \emph{cover cubes} of $x$ in $\mathscr{Y}$ are the cubes of $\mathscr{Y}$ that contain a point of $\operatorname{supp} x$.

\begin{lemma}[Cover cubes are linked]\label{lem:cover-cubes-linked}
Let $x$ be an index whose support is contained in $Y$ and which belongs to the kind of $\mathscr{Y}$ (Section~\ref{sec:two-tori}).
\begin{enumerate}
\item[(a)] Any two cover cubes of $x$ in $\mathscr{Y}$ are equal or linked.
\item[(b)] Exactly one link-component of $Y$ contains $x$.
\end{enumerate}
Consequently, among the indices that belong to the kind of $\mathscr{Y}$, those whose support lies in $Y$ are precisely those that are contained in some link-component of $Y$.
\end{lemma}

\begin{proof}
\emph{Reduction to $\mathbb{L}_k$.} Fix a point $q$ of $\operatorname{supp} x$ and a lift $\tilde{q} \in \mathbb{L}_k$ of $q$. Supports have sup-diameter at most $L$, and every parameter cube lies within sup-distance $L \cdot M$ of its centre; so every point of $\operatorname{supp} x$ and of the cover cubes of $x$ has sup-distance at most $2L \cdot M + L$ from $q$. As $l \in \mathrm{Cl}$, one has $l \ge 13L \cdot M$, hence $2L \cdot M + L < l/2$. Let $B$ be the sup-ball of radius $2L \cdot M + L$ in $\mathbb{L}_k$ around $\tilde{q}$. The projection of $\mathbb{L}_k$ onto $T_k$ is injective on $B$, because two points of $B$ differ by a vector of sup-norm less than $l$. Replace each point $p$ of $\operatorname{supp} x$ by its lift $\tilde{p}$ at sup-distance at most $L$ from $\tilde{q}$; this lift exists because $\operatorname{supp} x$ has sup-diameter at most $L$. The lifts keep the difference vectors of the points of $\operatorname{supp} x$: for $p, p' \in \operatorname{supp} x$ with the lifts $\tilde{p}$, $\tilde{p}'$, the vector $\tilde{p}' - \tilde{p}$ is the unique representative of sup-norm less than $l/2$ of the class $p' - p \in \mathbb{R}^3/l\mathbb{Z}^3$. Indeed, $\tilde{p}' - \tilde{p}$ represents this class, and $\lvert \tilde{p}' - \tilde{p} \rvert \le 2L < l/2$, because both lifts lie within sup-distance $L$ of $\tilde{q}$ and $l > 4L$; two representatives of sup-norm less than $l/2$ differ by a vector of $l\mathbb{Z}^3$ of sup-norm less than $l$, that is, by $0$. Every other representative of $p' - p$ has sup-norm at least $l - 2L > 2L$; as the sup-distance of $p$ and $p'$ is the least sup-norm of a representative of $p' - p$, it equals $\lvert \tilde{p}' - \tilde{p} \rvert$, and the lifted points have the same mutual sup-distances as the points themselves. In particular, if $x$ is a bond with the end points $p$ and $p + a \cdot e_{\nu}$, where $a \in \{\eta_k, 1\}$, then $a \cdot e_{\nu}$ represents the class $(p + a \cdot e_{\nu}) - p$ and has sup-norm $a \le 1 < l/2$; so the lifted end points are $\tilde{p}$ and $\tilde{p} + a \cdot e_{\nu}$, and they differ in the coordinate $\nu$ only. Lift each cover cube to the cube of $\mathbb{L}_k$ that contains the lift of one of its points of $\operatorname{supp} x$; this cube lies in $B$, and by the injectivity it contains the lift of every point of $\operatorname{supp} x$ that the cover cube contains. A lattice site common to two lifted cubes projects to a lattice site common to the two cover cubes. Hence, for the kinds $\pi'_k$ and $\pi''_k$, assertion (a) may be proved in $\mathbb{L}_k$ with these lifts; in that part of the proof all points and cubes are lifts.

\emph{Proof of (a).} Let $Q_1$, $Q_2$ be cover cubes of $x$.

If the kind is $\pi_k$, then $x$ is a site. The $M$-cubes partition the sites (fact~\eqref{it:renormalisation-group-g0} of Section~\ref{sec:cubes-parameter-cubes}); so exactly one $M$-cube contains $x$, and $Q_1 = Q_2$.

If the kind is $\pi'_k$ or $\pi''_k$, two cubes of this kind with a lattice site in common are equal or linked, by Definition~\ref{def:parameter-cubes-polydisks}. It therefore suffices to find a lattice site in $Q_1 \cap Q_2$.

Let the kind be $\pi'_k$. Then $x$ is a bond; let $p$ and $p' = p + a \cdot e_{\nu}$ be its end points, where $a = \eta_k$ for a bond of the fine lattice and $a = 1$ for a bond of the unit lattice. Both $p$ and $p'$ are lattice sites. Suppose that $Q_1$ and $Q_2$ have no end point of $x$ in common. Each of them contains an end point, so after renumbering $p \in Q_1$, $p' \notin Q_1$, $p' \in Q_2$ and $p \notin Q_2$. Let $c_1, c_2 \in M\mathbb{Z}^3$ be the centres, so that $Q_j = \{ y : \lvert y - c_j \rvert \le M \}$ for $j = 1, 2$. The points $p$ and $p'$ have the same coordinates except the coordinate $\nu$, and $p'_{\nu} = p_{\nu} + a > p_{\nu}$. From $p \in Q_1$ and $p' \notin Q_1$ one obtains $\lvert p'_{\nu} - c_{1,\nu} \rvert > M \ge \lvert p_{\nu} - c_{1,\nu} \rvert$; as $p'_{\nu} > p_{\nu} \ge c_{1,\nu} - M$, this gives $p'_{\nu} - c_{1,\nu} > M$. From $p' \in Q_2$ and $p \notin Q_2$ one obtains $\lvert p_{\nu} - c_{2,\nu} \rvert > M \ge \lvert p'_{\nu} - c_{2,\nu} \rvert$; as $p_{\nu} < p'_{\nu} \le c_{2,\nu} + M$, this gives $p_{\nu} - c_{2,\nu} < -M$. Thus
\[
p_{\nu} \le c_{1,\nu} + M < p_{\nu} + a, \qquad p_{\nu} < c_{2,\nu} - M \le p_{\nu} + a.
\]
All numbers in these inequalities are integral multiples of $a$, because $p \in a\mathbb{Z}^3$ and $M\mathbb{Z} \subset \mathbb{Z} \subset \eta_k\mathbb{Z}$. The only multiples of $a$ in the interval from $p_{\nu}$ to $p_{\nu} + a$ are its two end points; hence $c_{1,\nu} + M = p_{\nu}$ and $c_{2,\nu} - M = p_{\nu} + a$, and therefore
\[
c_{2,\nu} - c_{1,\nu} = 2M + a.
\]
Since $0 < a \le 1 < M$, this number is not a multiple of $M$, which contradicts $c_1, c_2 \in M\mathbb{Z}^3$. So $Q_1$ and $Q_2$ have the lattice site $p$ or the lattice site $p'$ in common.

Let the kind be $\pi''_k$. Then $x = \Upsilon$ is a basis element. By Proposition~\ref{prop:basis-map-box}, $\operatorname{supp} \Upsilon$ lies in one $L$-block or consists of the two end points of one bond; hence $\lvert p - p' \rvert \le L - 1$ for all $p, p' \in \operatorname{supp} \Upsilon$. Choose $p \in Q_1 \cap \operatorname{supp} \Upsilon$ and $p' \in Q_2 \cap \operatorname{supp} \Upsilon$, and let $c_1, c_2 \in LM\mathbb{Z}^3$ be the centres, so that $Q_j = \{ y : \lvert y - c_j \rvert \le L \cdot M \}$ for $j = 1, 2$. For every $\nu$,
\[
\lvert c_{1,\nu} - c_{2,\nu} \rvert \le \lvert c_{1,\nu} - p_{\nu} \rvert + \lvert p_{\nu} - p'_{\nu} \rvert + \lvert p'_{\nu} - c_{2,\nu} \rvert \le 2L \cdot M + L - 1 < 3L \cdot M,
\]
and $c_{1,\nu} - c_{2,\nu}$ is a multiple of $L \cdot M$; hence $c_{1,\nu} - c_{2,\nu} \in \{0, L \cdot M, -L \cdot M, 2L \cdot M, -2L \cdot M\}$. Let $z$ be the point with the coordinates
\[
z_{\nu} := \begin{cases} \frac{1}{2} (c_{1,\nu} + c_{2,\nu}) & \text{if } c_{1,\nu} - c_{2,\nu} \in \{0, 2L \cdot M, -2L \cdot M\}, \\ c_{1,\nu} & \text{if } c_{1,\nu} - c_{2,\nu} \in \{L \cdot M, -L \cdot M\}. \end{cases}
\]
In each case $z_{\nu}$ is a multiple of $L \cdot M$, and $\lvert z_{\nu} - c_{1,\nu} \rvert \le L \cdot M$, $\lvert z_{\nu} - c_{2,\nu} \rvert \le L \cdot M$. So $z \in \mathbb{Z}^3$ is a unit site, hence a lattice site, and $z$ lies in $Q_1 \cap Q_2$.

\emph{Proof of (b).} By (a), the cover cubes of $x$ form a link-connected subset of $\mathscr{Y}$; this subset is not empty, since $\operatorname{supp} x$ is a non-empty subset of $Y$. So the cover cubes of $x$ belong to one link-component of $\mathscr{Y}$; let $Y_1$ be the union of its cubes. Each point of $\operatorname{supp} x$ lies in $Y$, hence in a cube of $\mathscr{Y}$, and this cube is a cover cube of $x$. Therefore $\operatorname{supp} x \subset Y_1$, that is, $x$ is contained in $Y_1$. The points of $\operatorname{supp} x$ are lattice sites, and different link-components of $Y$ have no lattice site in common, by fact~\eqref{it:renormalisation-group-g3} of Section~\ref{sec:cubes-parameter-cubes}. Hence $Y_1$ is the only link-component of $Y$ that contains $x$.

\emph{The last assertion.} Let $u$ be an index that belongs to the kind of $\mathscr{Y}$. If $\operatorname{supp} u \subset Y$, then (b), applied to $u$, shows that $u$ is contained in a link-component of $Y$. Conversely, if $u$ is contained in a link-component $Y_1$ of $Y$, then $\operatorname{supp} u \subset Y_1 \subset Y$.
\end{proof}

\subsubsection{Projection of a system of primitives}\label{sec:projection-system-primitives}

\begin{definition}[Projection]\label{def:projection}
Let a system of primitives of level $k$ on $T_2$ (Section~\ref{sec:primitive-operators}) with property~\ref{it:renormalisation-group-14} be given. Each of its families is denoted by $\mathrm{Op}^{(2)}$, and its barred family of~\eqref{it:primitives-p3} by $\bar{H}_k^{(2)}$. Let $s \in \mathrm{Pd}_k(T_1)$, $A \in \tilde{G}_k(T_1)$, $X \in D_k(T_1)$ and $y \ge 0$, and let $f$ be a function on $T_1$ of the kind on which the family acts. The projection of the system to $T_1$ (Definition~\ref{def:periodisation-projection}) consists of the families $\mathrm{Op}^{\mathrm{pj}}$ given as follows. For each of the families~\eqref{it:primitives-p1}--\eqref{it:primitives-p3} and $\Gamma_k$,
\[
\mathrm{pj}^{*}\left( \mathrm{Op}^{\mathrm{pj}}(s,A) f \right) := \mathrm{Op}^{(2)}\left( \mathrm{pj}^{*} s, \mathrm{pj}^{*} A \right) \mathrm{pj}^{*} f,
\]
the argument $A$ being absent for the families~\eqref{it:primitives-p1} and~\eqref{it:primitives-p2}; for $M_k$ and $S_{k,y}$,
\[
\mathrm{pj}^{*}\left( \mathrm{Op}^{\mathrm{pj}}(X,A) f \right) := \sum_{X' \in D_k(T_2) :\ \mathrm{pj}(X') = X} \mathrm{Op}^{(2)}\left( X', \mathrm{pj}^{*} A \right) \mathrm{pj}^{*} f.
\]
The projected families are written $H_k^{\mathrm{pj}}$, $C_k^{\frac{1}{2},\mathrm{pj}}$, $\Gamma_k^{\mathrm{pj}}$, $M_k^{\mathrm{pj}}$, $S_{k,y}^{\mathrm{pj}}$. Moreover,
\[
\bar{H}_k^{\mathrm{pj}}(s,A) := \mathsf{C}^{-1} H_k^{\mathrm{pj}}(s,-A)\, \mathsf{C},
\]
where $-A \in \tilde{G}_k(T_1)$ by Lemma~\ref{lem:invariances}. Because $\mathrm{pj}^{*}$ is compatible with the change of sign $A \mapsto -A$ and with multiplication by constant spinor matrices,
\[
\mathrm{pj}^{*}\left( \bar{H}_k^{\mathrm{pj}}(s,A) f \right) = \bar{H}_k^{(2)}\left( \mathrm{pj}^{*} s, \mathrm{pj}^{*} A \right) \mathrm{pj}^{*} f;
\]
in other words, $\bar{H}_k^{\mathrm{pj}}$ is the projection of $\bar{H}_k^{(2)}$. A system of level $0$ has no family $M_0$ (property~\ref{it:renormalisation-group-17}). Finally, put
\[
\widetilde{\mathrm{Op}}(s,A) := \mathrm{Op}^{(2)}\left( \mathrm{pj}^{*} s, \mathrm{pj}^{*} A \right), \qquad \widetilde{\mathrm{Op}}(X,A) := \sum_{X' \in D_k(T_2) :\ \mathrm{pj}(X') = X} \mathrm{Op}^{(2)}\left( X', \mathrm{pj}^{*} A \right).
\]
\end{definition}

\begin{theorem}[Properties of the projection]\label{thm:properties-projection}
In the setting of Section~\ref{sec:two-tori}, in particular with $l_1, l_2 \in \mathrm{Cl}$, $l_2 = J \cdot l_1$, $J$ odd, and $\gamma_w \cdot M \ge 1$, let a system of primitives of level $k$ on $T_2$ have the properties~\ref{it:renormalisation-group-1}--\ref{it:renormalisation-group-17} for the parameters
\[
\mathrm{par}^0 = \left( \alpha_0, \gamma_w, \kappa_w + \kappa_0; C_w^1 \right)
\]
of Section~\ref{sec:primitive-operators}.
\begin{enumerate}
\item[(a)] The right-hand sides in Definition~\ref{def:projection} belong to $\mathrm{Per}$, so that the projected operators are well defined. Their kernels are given by
\begin{equation}\label{eq:renormalisation-flow-4}
\mathrm{Op}^{\mathrm{pj}}(\,\cdot\,)(x,x') = \sum_{w \in (\mathbb{Z}/J\mathbb{Z})^3} \boldsymbol{\chi}(w) \cdot \widetilde{\mathrm{Op}}(\,\cdot\,)(x, x' + l_1 \cdot w),
\end{equation}
where $(\,\cdot\,)$ stands for $(s,A)$, resp. for $(X,A)$, $x$ and $x'$ are indices of $\mathbb{L}_k$, and $\boldsymbol{\chi}$ is the character of the functions on which the family acts.
\item[(b)] The projected system satisfies the properties~\ref{it:renormalisation-group-1}--\ref{it:renormalisation-group-17} on $T_1$ for the parameters
\[
\mathrm{par} = \left( \alpha_0, \gamma_w, \kappa_w; C_w \right), \qquad C_w = 64 K_0 \cdot C_w^1.
\]
\item[(c)] At level $0$, projecting the two families $H_0(s')$ and $H_0(s'',A)$ of property~\ref{it:renormalisation-group-17} on $T_2$ yields the corresponding families of property~\ref{it:renormalisation-group-17} on $T_1$.
\end{enumerate}
\end{theorem}

\begin{proof}
Functions on the two tori, their characters, the space $\mathrm{Per}$, the translations $T_g$, the kernels and the sets $\mathscr{F}_i$ are those of Section~\ref{sec:two-tori}; indices, their supports and the distance of two indices are those of Section~\ref{sec:lattice-geometry}.

\emph{Part (a).} Let $g \in l_1\mathbb{Z}^3$. The objects $\mathrm{pj}^{*} s$ and $\mathrm{pj}^{*} A$ are $l_1\mathbb{Z}^3$-periodic, so that $T_g \mathrm{pj}^{*} s = \mathrm{pj}^{*} s$ and $T_g \mathrm{pj}^{*} A = \mathrm{pj}^{*} A$. Since $l_1 \in \mathrm{Cl}$ is a multiple of $L \cdot M$, the translation by $g$, with the spinor matrix $1$, is an element of $\mathbb{G}_k^L(T_2)$ and of $\mathbb{G}^L_{k,3}(T_2)$ (Section~\ref{sec:dirac-matrices-lattice}), and it acts by $T_g$. By property~\ref{it:renormalisation-group-14} on $T_2$, for $s_2 \in \mathrm{Pd}_k(T_2)$, $A_2 \in \tilde{G}_k(T_2)$ and $X' \in D_k(T_2)$,
\[
T_g \mathrm{Op}^{(2)}(s_2,A_2) T_g^{-1} = \mathrm{Op}^{(2)}(T_g s_2, T_g A_2), \qquad T_g \mathrm{Op}^{(2)}(X',A_2) T_g^{-1} = \mathrm{Op}^{(2)}(X' + g, T_g A_2).
\]
Hence $\widetilde{\mathrm{Op}}(s,A)$ commutes with $T_g$. So does $\widetilde{\mathrm{Op}}(X,A)$, since the translation $X' \mapsto X' + g$ is a bijection of the set $\{ X' \in D_k(T_2) : \mathrm{pj}(X') = X \}$ onto itself. Every family of the system preserves the character: it maps functions of character $\boldsymbol{\chi}$ to functions of the same character $\boldsymbol{\chi}$. Therefore, for $f \in \mathrm{Per}$ and $w \in \mathbb{Z}^3$,
\[
T_{-l_1 \cdot w} \widetilde{\mathrm{Op}} f = \widetilde{\mathrm{Op}} T_{-l_1 \cdot w} f = \boldsymbol{\chi}(w) \cdot \widetilde{\mathrm{Op}} f,
\]
that is, $\widetilde{\mathrm{Op}}$ maps $\mathrm{Per}$ into $\mathrm{Per}$. As $\mathrm{Per}$ is the image of the injective linear map $\mathrm{pj}^{*}$, the identities of Definition~\ref{def:projection} determine $\mathrm{Op}^{\mathrm{pj}}(\,\cdot\,) f$, and $\mathrm{Op}^{\mathrm{pj}}(\,\cdot\,)$ is linear.

For the kernels, take as $\mathscr{F}_2$ the union of the disjoint sets $\mathscr{F}_1 + l_1 \cdot w$, $w$ ranging over a set of representatives of $(\mathbb{Z}/J\mathbb{Z})^3$; recall that $f(x' + l_1 \cdot w) = \boldsymbol{\chi}(w) \cdot f(x')$ for $f \in \mathrm{Per}$. For a function $f$ on $T_1$ and an index $x$ of $\mathbb{L}_k$ it follows that
\[
\left( \mathrm{Op}^{\mathrm{pj}}(\,\cdot\,) f \right)(x) = \left( \widetilde{\mathrm{Op}}(\,\cdot\,) \mathrm{pj}^{*} f \right)(x) = \sum_{x' \in \mathscr{F}_1} \left( \sum_w \boldsymbol{\chi}(w) \cdot \widetilde{\mathrm{Op}}(\,\cdot\,)(x, x' + l_1 \cdot w) \right) \cdot f(x'),
\]
which is~\eqref{eq:renormalisation-flow-4}. The term of~\eqref{eq:renormalisation-flow-4} that belongs to a class is independent of the representative $w$ chosen. Indeed, if $w$ is replaced by $w + J \cdot e_{\nu}$, then $\widetilde{\mathrm{Op}}(\,\cdot\,)(x, x' + l_1 \cdot w)$ acquires the factor $\boldsymbol{\chi}(e_{\nu})$, as $J \cdot l_1 = l_2$, and $\boldsymbol{\chi}(w)$ acquires the factor $\boldsymbol{\chi}(J \cdot e_{\nu}) = \boldsymbol{\chi}(e_{\nu})$, as $J$ is odd; and $\boldsymbol{\chi}(e_{\nu})^2 = 1$.

\emph{Transposition and periodic extension.} Let $O^{(1)}$ and $O^{(2)}$ be operators on $T_1$ and on $T_2$ as in Section~\ref{sec:two-tori}, let $O^{(2)}$ commute with the translations $T_g$, $g \in l_1\mathbb{Z}^3$, and let $O^{(2)} \mathrm{pj}^{*} = \mathrm{pj}^{*} O^{(1)}$. Then
\begin{equation}\label{eq:renormalisation-flow-5}
\left( O^{(2)} \right)^T \mathrm{pj}^{*} = \mathrm{pj}^{*} \left( O^{(1)} \right)^T,
\end{equation}
the transposes being formed with respect to the bilinear pairings of Section~\ref{sec:two-tori}. Indeed, the pairings are invariant under the translations, $\langle T_g \varphi, T_g \varphi' \rangle_{T_2} = \langle \varphi, \varphi' \rangle_{T_2}$ for functions $\varphi$, $\varphi'$ on $T_2$ with the same character; hence $T_g^T = T_{-g}$, and transposing $O^{(2)} T_{-g} = T_{-g} O^{(2)}$ shows that $(O^{(2)})^T$ commutes with every $T_g$, $g \in l_1\mathbb{Z}^3$. As in part (a), it follows that $(O^{(2)})^T$ maps $\mathrm{Per}$ into $\mathrm{Per}$. So for every function $\varphi_1$ on $T_1$ there is a function $\varphi_2$ on $T_1$ with $(O^{(2)})^T \mathrm{pj}^{*} \varphi_1 = \mathrm{pj}^{*} \varphi_2$, and for every function $f$ on $T_1$
\begin{align*}
J^3 \cdot \langle f, \varphi_2 \rangle_{T_1} &= \left\langle \mathrm{pj}^{*} f, (O^{(2)})^T \mathrm{pj}^{*} \varphi_1 \right\rangle_{T_2} = \left\langle O^{(2)} \mathrm{pj}^{*} f, \mathrm{pj}^{*} \varphi_1 \right\rangle_{T_2} \\
&= \left\langle \mathrm{pj}^{*} O^{(1)} f, \mathrm{pj}^{*} \varphi_1 \right\rangle_{T_2} = J^3 \cdot \left\langle f, (O^{(1)})^T \varphi_1 \right\rangle_{T_1}.
\end{align*}
The pairing on $T_1$ being non-degenerate, $\varphi_2 = (O^{(1)})^T \varphi_1$.

\emph{Part (b), property~\ref{it:renormalisation-group-1}.} For $i = 1, 2$ the exact operators of Sections~\ref{sec:averaging-fluctuations-gauge} and~\ref{sec:exact-fermion-operators} on $T_i$ carry the superscript $(i)$; they are told apart from the families of the system by their arguments. Let $1$ also denote the parameter family all of whose entries are equal to $1$; then $\mathrm{pj}^{*} 1 = 1$. By Definition~\ref{def:projection}, property~\ref{it:renormalisation-group-1} on $T_2$ and Proposition~\ref{prop:periodisation},
\begin{align*}
\mathrm{pj}^{*}\left( H_k^{\mathrm{pj}}(1) f \right) &= H_k^{(2)}(\mathrm{pj}^{*} 1) \mathrm{pj}^{*} f = H_k^{(2)} \mathrm{pj}^{*} f = \mathrm{pj}^{*} H_k^{(1)} f, \\
\mathrm{pj}^{*}\left( C_k^{\frac{1}{2},\mathrm{pj}}(1) f \right) &= C_k^{\frac{1}{2},(2)}(\mathrm{pj}^{*} 1) \mathrm{pj}^{*} f = C_k^{\frac{1}{2},(2)} \mathrm{pj}^{*} f = \mathrm{pj}^{*} C_k^{\frac{1}{2},(1)} f.
\end{align*}
By Definition~\ref{def:projection}, property~\ref{it:renormalisation-group-1} on $T_2$ and Proposition~\ref{prop:coverings-tori},
\begin{align*}
\mathrm{pj}^{*}\left( H_k^{\mathrm{pj}}(1,A) f \right) &= H_k^{(2)}(\mathrm{pj}^{*} 1, \mathrm{pj}^{*} A) \mathrm{pj}^{*} f = H_k^{(2)}(\mathrm{pj}^{*} A) \mathrm{pj}^{*} f = \mathrm{pj}^{*} H_k^{(1)}(A) f, \\
\mathrm{pj}^{*}\left( \Gamma_k^{\mathrm{pj}}(1,A) f \right) &= \Gamma_k^{(2)}(\mathrm{pj}^{*} 1, \mathrm{pj}^{*} A) \mathrm{pj}^{*} f = \Gamma_k^{(2)}(\mathrm{pj}^{*} A) \mathrm{pj}^{*} f = \mathrm{pj}^{*} \Gamma_k^{(1)}(A) f.
\end{align*}
Since $\mathrm{pj}$ carries $M$-cubes onto $M$-cubes, and cubes with a common face onto cubes with a common face, it maps every polymer $X'$ of $T_2$ onto a polymer $\mathrm{pj}(X')$ of $T_1$. Hence the sets $\{ X' : \mathrm{pj}(X') = X \}$, $X \in D_k(T_1)$, form a partition of $D_k(T_2)$, and
\[
\sum_{X \in D_k(T_1)} \widetilde{\mathrm{Op}}(X,A) = \sum_{X' \in D_k(T_2)} \mathrm{Op}^{(2)}(X', \mathrm{pj}^{*} A),
\]
which is $M_k^{(2)}(\mathrm{pj}^{*} A)$, resp. $S_{k,y}^{(2)}(\mathrm{pj}^{*} A)$, by property~\ref{it:renormalisation-group-1} on $T_2$. By Proposition~\ref{prop:coverings-tori}, $M_k^{(2)}(\mathrm{pj}^{*} A) \mathrm{pj}^{*} = \mathrm{pj}^{*} M_k^{(1)}(A)$ and $S_{k,y}^{(2)}(\mathrm{pj}^{*} A) \mathrm{pj}^{*} = \mathrm{pj}^{*} S_{k,y}^{(1)}(A)$. Since $\mathrm{pj}^{*}$ is injective, the identities of property~\ref{it:renormalisation-group-1} hold on $T_1$.

\emph{Property~\ref{it:renormalisation-group-2}.} The map $(s,A) \mapsto (\mathrm{pj}^{*} s, \mathrm{pj}^{*} A)$ is linear. It maps $\mathrm{Pd}_k(T_1) \times \tilde{G}_k(T_1)$ into $\mathrm{Pd}_k(T_2) \times \tilde{G}_k(T_2)$, and interior points to interior points: for the first factor because $(\mathrm{pj}^{*} s)_c = s_{\mathrm{pj}(c)}$ (Section~\ref{sec:two-tori}), the polydisks being those of Definition~\ref{def:parameter-cubes-polydisks}; for the second factor because Lemma~\ref{lem:periodic-extension} gives $\mathrm{pj}^{*} A \in \tilde{G}_k(T_2)$, and every point of $\tilde{G}_k(T_2)$ is an interior point, this set being open (Section~\ref{sec:domains-bond-functions}; Lemma~\ref{lem:openness}). By~\eqref{eq:renormalisation-flow-4}, each entry of a projected family is a finite sum of entries of the corresponding family of the system on $T_2$ at these arguments, multiplied by the signs $\boldsymbol{\chi}(w)$. Consequently the projected families~\eqref{it:primitives-p1}--\eqref{it:primitives-p3} and $\Gamma_k^{\mathrm{pj}}$ are continuous on $\mathrm{Pd}_k(T_1) \times \tilde{G}_k(T_1)$ and analytic in $(s,A)$ in the interior, and $S^{\mathrm{pj}}_{k,y}(X,A)$ is Borel in $y$. Now fix $s \in \mathrm{Pd}_k(T_1)$, possibly on the boundary. The entries of $\mathrm{Op}^{(2)}(\mathrm{pj}^{*} s,\,\cdot\,)$ are analytic on $\tilde{G}_k(T_2)$, and $A \mapsto \mathrm{pj}^{*} A$ is a linear map of $\tilde{G}_k(T_1)$ into $\tilde{G}_k(T_2)$; hence the entries of $\mathrm{Op}^{\mathrm{pj}}(s,\,\cdot\,)$, and likewise those of $M_k^{\mathrm{pj}}(X,\,\cdot\,)$ and $S_{k,y}^{\mathrm{pj}}(X,\,\cdot\,)$, are analytic on $\tilde{G}_k(T_1)$.

\emph{Property~\ref{it:renormalisation-group-5}.} Let $s \in \mathrm{Pd}_k(T_1)$, $A \in \tilde{G}_k(T_1)$, and let $f$ be a spinor function on $T^0_1$. Since $\lvert \mathrm{pj}^{*} \varphi \rvert_{\infty} = \lvert \varphi \rvert_{\infty}$ for every function $\varphi$ on $T_1$, property~\ref{it:renormalisation-group-5} on $T_2$ gives
\[
\left\lvert H_k^{\mathrm{pj}}(s,A) f \right\rvert_{\infty} = \left\lvert H_k^{(2)}(\mathrm{pj}^{*} s, \mathrm{pj}^{*} A) \mathrm{pj}^{*} f \right\rvert_{\infty} \le C_w^1 \cdot \lvert f \rvert_{\infty}.
\]
Let $x$, $x'$ be distinct sites of $\mathbb{L}_k$ with $\lvert x - x' \rvert < 1$. As $\lvert x - x' \rvert < 1 < l_1/2$, their images in $T_1$ form a pair of $\mathrm{Pz}_k(T_1)$ and their images in $T_2$ a pair of $\mathrm{Pz}_k(T_2)$, both at sup-distance $\lvert x - x' \rvert$, and $(x,x')$ is the shortest lift of each of these two pairs (Definition~\ref{def:pairs-sites-holder}); conversely, every pair of $\mathrm{Pz}_k(T_1)$ at sup-distance less than $1$ is the image of its shortest lift, which is of this form. For a spinor function $h$ on $T_1$ the number $(\delta_{\alpha,\mathrm{pj}^{*} A} \mathrm{pj}^{*} h)(x,x')$, formed on $T_2$, is therefore equal to the number $(\delta_{\alpha,A} h)(x,x')$, formed on $T_1$, since the two computations use one lift, hence the same six paths, and the same values of the function and of the connection on $\mathbb{L}_k$. For $h := H_k^{\mathrm{pj}}(s,A) f$ one has $\mathrm{pj}^{*} h = H_k^{(2)}(\mathrm{pj}^{*} s, \mathrm{pj}^{*} A) \mathrm{pj}^{*} f$, and property~\ref{it:renormalisation-group-5} on $T_2$ yields
\[
\left\lvert \left( \delta_{\alpha,A} H_k^{\mathrm{pj}}(s,A) f \right)(x,x') \right\rvert \le C_w^1 \cdot \lvert f \rvert_{\infty}.
\]
Here $C_w^1 \le 64 K_0 \cdot C_w^1 = C_w$, since $K_0 = \exp(\kappa_0)/36 \ge 1$, $\kappa_0 = 54 \log 72$ being the number fixed in~\eqref{eq:polymer-activities-5} (Section~\ref{sec:absolute-constants}).

\emph{An estimate for sums over $(\mathbb{Z}/J\mathbb{Z})^3$.} For $i = 1, 2$ let $d_i$ be the sup distance on $T_i$; thus $d_i(x,x')$ is the sup distance of the images in $T_i$ of $\operatorname{supp} x$ and $\operatorname{supp} x'$. For subsets $a$, $a'$ of $\mathbb{L}_k$ let $\operatorname{dist}(a,a')$ be the least of the numbers $\lvert p - p' \rvert$, $p \in a$, $p' \in a'$. It is shown that, for all indices $x$, $x'$ of $\mathbb{L}_k$,
\begin{equation}\label{eq:renormalisation-flow-6}
\sum_{w \in (\mathbb{Z}/J\mathbb{Z})^3} \exp\left( -\gamma_w \cdot d_2(x, x' + l_1 \cdot w) \right) \le 64 \exp\left( -\gamma_w \cdot d_1(x,x') \right).
\end{equation}
Neither side changes if $x$ or $x'$ is replaced by its translate by an element of $l_1\mathbb{Z}^3$: on the left this permutes the classes $w$. One may therefore assume that $a := \operatorname{supp} x$ and $a' := \operatorname{supp} x'$ satisfy
\[
\operatorname{dist}(a,a') = d_1 := d_1(x,x') = \min_{v \in \mathbb{Z}^3} \operatorname{dist}(a, a' + l_1 \cdot v).
\]
In particular $\operatorname{dist}(a, a' + l_1 \cdot w') \ge d_1$ for every $w' \in \mathbb{Z}^3$. Moreover $d_1 \le l_1/2$, since every coordinate of a vector of $\mathbb{L}_k$ can be changed by a multiple of $l_1$ into a number of modulus at most $l_1/2$. For a class $w$ choose $w' \in w + J\mathbb{Z}^3$ with $d_2(x, x' + l_1 \cdot w) = \operatorname{dist}(a, a' + l_1 \cdot w')$; this is possible, because the distance on $T_2$ is the least of the numbers $\operatorname{dist}(a, a' + l_1 \cdot w + l_2 \cdot v)$, $v \in \mathbb{Z}^3$, and $l_2 = J \cdot l_1$. Different classes give different points $w'$.

Now let $w' \in \mathbb{Z}^3$ be arbitrary, let $p \in a$, $p' \in a'$ with $\lvert p - p' \rvert = d_1$, and let $u \in a$, $u' \in a'$ with $\lvert u - u' - l_1 \cdot w' \rvert = \operatorname{dist}(a, a' + l_1 \cdot w')$. Since every support has sup-diameter at most $L$ (Section~\ref{sec:two-tori}),
\[
l_1 \cdot \lvert w' \rvert \le \lvert u - u' - l_1 \cdot w' \rvert + \lvert u - p \rvert + \lvert p - p' \rvert + \lvert p' - u' \rvert \le \operatorname{dist}(a, a' + l_1 \cdot w') + d_1 + 2L.
\]
So, if $\lvert w' \rvert = t \ge 2$, then, by $d_1 \le l_1/2$ and $t - 1 \ge 1$,
\[
\operatorname{dist}(a, a' + l_1 \cdot w') \ge l_1 \cdot t - d_1 - 2L \ge d_1 + l_1 \cdot (t-1) - 2L \ge d_1 + (l_1 - 2L) \cdot (t-1).
\]
The cube $\lvert w' \rvert \le 1$ contains $27$ points of $\mathbb{Z}^3$, and the sphere $\lvert w' \rvert = t$ contains $(2t+1)^3 - (2t-1)^3 = 24 t^2 + 2 \le 26 t^2$ of them. As $l_1 \in \mathrm{Cl}$, one has $l_1 \ge 13 L \cdot M$, hence $l_1 - 2L \ge 6M$, and $\gamma_w \cdot (l_1 - 2L) \ge 6 \gamma_w \cdot M \ge 6$. Since different classes give different points $w'$, the left side of~\eqref{eq:renormalisation-flow-6} is at most the sum of $\exp(-\gamma_w \cdot \operatorname{dist}(a, a' + l_1 \cdot w'))$ over all $w' \in \mathbb{Z}^3$, and therefore at most
\[
\exp\left( -\gamma_w \cdot d_1 \right) \cdot \left( 27 + \sum_{t \ge 2} 26 t^2 \cdot \exp\left( -6(t-1) \right) \right).
\]
Here $\sum_{t \ge 2} 26 t^2 \cdot \exp(-6(t-1)) < 0.3$: the term $t = 2$ is $104 \exp(-6) < 0.26$, and the terms with $t \ge 3$ add up to less than $0.01$. Hence the left side of~\eqref{eq:renormalisation-flow-6} is less than $27.3 \exp(-\gamma_w \cdot d_1)$, which proves~\eqref{eq:renormalisation-flow-6}.

\emph{The properties~\ref{it:renormalisation-group-3}, \ref{it:renormalisation-group-4}, and property~\ref{it:renormalisation-group-6} for $\Gamma_k$.} In property~\ref{it:renormalisation-group-3} the difference quotients and the H\"older differences are taken in the first argument of the kernel. Hence, by~\eqref{eq:renormalisation-flow-4}, for every direction $\mathfrak{l} \in \{0,1,2\}$ and for fine bonds $b$, $\tilde{b}$ and a unit bond $b'$ as in property~\ref{it:renormalisation-group-3}, each of the kernels
\[
H_k^{\mathrm{pj}}(s)(b,b'), \qquad \left( \partial_{\mathfrak{l}} H_k^{\mathrm{pj}}(s) \right)(b,b'), \qquad \left( \delta_{\alpha} \partial_{\mathfrak{l}} H_k^{\mathrm{pj}}(s) \right)((b,\tilde{b});b')
\]
is the sum over $w$ of the corresponding kernel of $H_k^{(2)}(\mathrm{pj}^{*} s)$ with $b' + l_1 \cdot w$ in place of $b'$, the character of bond functions being $1$; and the distance $d_2(b, b' + l_1 \cdot w)$ is precisely the one that appears in property~\ref{it:renormalisation-group-3} on $T_2$. In the same way, by~\eqref{eq:renormalisation-flow-4}, for basis elements $\Upsilon$, $\Upsilon'$ the kernel $C_k^{\frac{1}{2},\mathrm{pj}}(s)(\Upsilon,\Upsilon')$ is the sum over $w$ of the kernels $C_k^{\frac{1}{2},(2)}(\mathrm{pj}^{*} s)(\Upsilon, \Upsilon' + l_1 \cdot w)$, and for unit sites $x$, $x'$ the kernel $\Gamma_k^{\mathrm{pj}}(s,A)(x,x')$ is the sum over $w$ of $\boldsymbol{\chi}(w) \cdot \Gamma_k^{(2)}(\mathrm{pj}^{*} s, \mathrm{pj}^{*} A)(x, x' + l_1 \cdot w)$, where $\lvert \boldsymbol{\chi}(w) \rvert = 1$. The properties~\ref{it:renormalisation-group-3}, \ref{it:renormalisation-group-4} and~\ref{it:renormalisation-group-6} on $T_2$, for the parameters $\mathrm{par}^0$, bound each summand by $C_w^1 \cdot \exp(-\gamma_w \cdot d_2)$, where $d_2$ is the distance on $T_2$ of the two indices of the summand, and~\eqref{eq:renormalisation-flow-6} bounds the sum by $64 C_w^1 \cdot \exp(-\gamma_w \cdot d_1)$, where $d_1$ is the distance on $T_1$ of the two indices on the left. So these bounds hold on $T_1$ with the rate $\gamma_w$ and the constant $64 C_w^1 \le 64 K_0 \cdot C_w^1 = C_w$; here $K_0 \ge 1$ by~\eqref{eq:polymer-activities-5} (Section~\ref{sec:absolute-constants}).

\emph{Property~\ref{it:renormalisation-group-6} for $M_k$, and property~\ref{it:renormalisation-group-7}.} For a unit site $x$ of $\mathbb{L}_k$ let $\square(x)$ be the $M$-cube of $T_2$ that contains the image of $x$. By fact~\eqref{it:renormalisation-group-g0} of Section~\ref{sec:cubes-parameter-cubes}, the $M$-cubes partition the sites of $T^0_2$ and of $T_2$, and the block $\Delta_x$ lies in one $M$-cube, which is $\square(x)$, because it contains $x$. A polymer $X' \in D_k(T_2)$ is a union of $M$-cubes; hence $x$, and any site of $\Delta_x$, lies in $X'$ if and only if $X'$ contains $\square(x)$. By property~\ref{it:renormalisation-group-12} on $T_2$, the kernel of $M_k^{(2)}(X',\,\cdot\,)$ at the row $x$, and the kernel of $S_{k,y}^{(2)}(X',\,\cdot\,)$ at the rows in $\Delta_x$, therefore vanish unless $X'$ contains $\square(x)$.

Next, $d_M(X') \ge d_M(\mathrm{pj}(X'))$ for every $X' \in D_k(T_2)$. Indeed, $M \cdot d_M$ of a polymer is the infimum of the lengths of the nets in the continuum torus that meet every closed $M$-cube of the polymer (Section~\ref{sec:lattice-geometry}; convention~\eqref{eq:polymer-activities-6}). The map of the continuum tori $\mathbb{R}^3/l_2\mathbb{Z}^3 \to \mathbb{R}^3/l_1\mathbb{Z}^3$ induced by the identity of $\mathbb{R}^3$ extends $\mathrm{pj}$ and is isometric on all sufficiently small balls. Under it, the image of a net meeting all closed $M$-cubes of $X'$ is a net meeting all closed $M$-cubes of $\mathrm{pj}(X')$, and its length does not exceed the length of the given net. Since $l_2 \in \mathrm{Cl}$, Proposition~\ref{prop:polymers}(t2) holds on $T_2$; together with $\kappa_w \ge 0$ it gives, for a unit site $x$ and $X \in D_k(T_1)$,
\begin{align*}
\sum_{\substack{X' \in D_k(T_2) :\ X' \supseteq \square(x), \\ \mathrm{pj}(X') = X}} \exp\left( -(\kappa_w + \kappa_0) \cdot d_M(X') \right) &\le \exp\left( -\kappa_w \cdot d_M(X) \right) \cdot \sum_{\substack{X' \in D_k(T_2) : \\ X' \supseteq \square(x)}} \exp\left( -\kappa_0 \cdot d_M(X') \right) \\
&\le K_0 \cdot \exp\left( -\kappa_w \cdot d_M(X) \right).
\end{align*}
By~\eqref{eq:renormalisation-flow-4}, the vanishing just shown, property~\ref{it:renormalisation-group-6} on $T_2$ for the parameters $\mathrm{par}^0$, \eqref{eq:renormalisation-flow-6} and this bound, for unit sites $x$, $x'$,
\begin{align*}
\left\lvert M_k^{\mathrm{pj}}(X,A)(x,x') \right\rvert &\le \sum_w \sum_{\substack{X' \supseteq \square(x) : \\ \mathrm{pj}(X') = X}} \left\lvert M_k^{(2)}(X', \mathrm{pj}^{*} A)(x, x' + l_1 \cdot w) \right\rvert \\
&\le C_w^1 \cdot \sum_w \exp\left( -\gamma_w \cdot d_2(x, x' + l_1 \cdot w) \right) \cdot \sum_{\substack{X' \supseteq \square(x) : \\ \mathrm{pj}(X') = X}} \exp\left( -(\kappa_w + \kappa_0) \cdot d_M(X') \right) \\
&\le 64 K_0 \cdot C_w^1 \cdot \exp\left( -\gamma_w \cdot d_1(x,x') - \kappa_w \cdot d_M(X) \right).
\end{align*}
Thus $M_k^{\mathrm{pj}}$ satisfies the bound of property~\ref{it:renormalisation-group-6} with the rates $\gamma_w$, $\kappa_w$ and the constant $64 K_0 \cdot C_w^1 = C_w$.

For property~\ref{it:renormalisation-group-7}, let $x$, $x'$ be unit sites and $f$ a spinor function on $T_1$. One has $\mathrm{pj}^{*}( \mathbf{1}_{\Delta_{x'}} f ) = \sum_w \mathbf{1}_{\Delta_{x' + l_1 \cdot w}} \mathrm{pj}^{*} f$; hence, by Definition~\ref{def:projection},
\[
\mathrm{pj}^{*}\left( S^{\mathrm{pj}}_{k,y}(X,A) \mathbf{1}_{\Delta_{x'}} f \right) = \sum_{X' :\ \mathrm{pj}(X') = X} \sum_w S^{(2)}_{k,y}(X', \mathrm{pj}^{*} A) \mathbf{1}_{\Delta_{x' + l_1 \cdot w}} \mathrm{pj}^{*} f.
\]
On the block $\Delta_x$ of a lift of $x$, by the vanishing shown above, property~\ref{it:renormalisation-group-7} on $T_2$ for the parameters $\mathrm{par}^0$, $\lvert \mathrm{pj}^{*} f \rvert_{\infty} = \lvert f \rvert_{\infty}$, \eqref{eq:renormalisation-flow-6} and the bound of the sum over $X'$, the modulus of this function is at most
\begin{align*}
&\sum_{\substack{X' \supseteq \square(x) : \\ \mathrm{pj}(X') = X}} \sum_w C_w^1 \cdot \exp\left( -\gamma_w \cdot d_2(x, x' + l_1 \cdot w) - (\kappa_w + \kappa_0) \cdot d_M(X') \right) \cdot \lvert f \rvert_{\infty} \\
&\qquad \le 64 K_0 \cdot C_w^1 \cdot \exp\left( -\gamma_w \cdot d_1(x,x') - \kappa_w \cdot d_M(X) \right) \cdot \lvert f \rvert_{\infty}.
\end{align*}
As $\mathrm{pj}^{*}$ does not change the values of a function, this is property~\ref{it:renormalisation-group-7} on $T_1$, for all $y \ge 0$, with the constant $64 K_0 \cdot C_w^1 = C_w$. At level $0$ property~\ref{it:renormalisation-group-7} is replaced by the kernel bound of property~\ref{it:renormalisation-group-17} for $S_{0,y}(X,A)$, which contains the factor $(1+y)^{-1}$; the same argument, applied to the kernels through~\eqref{eq:renormalisation-flow-4}, \eqref{eq:renormalisation-flow-6} and the bound of the sum over $X'$, gives this bound on $T_1$, the factor $(1+y)^{-1}$ being common to all terms.

\emph{Preimages of parameter cubes.} Let $\mathrm{Op}$ be one of the families of the system to which a kind is assigned in Section~\ref{sec:two-tori}, let $\mathscr{Y}$ be a set of parameter cubes of $T_1$ of the kind of $\mathrm{Op}$, and let $Y$ be the union of the cubes of $\mathscr{Y}$. Let $A\in\tilde{G}_k(T_1)$, and let $s\in\mathrm{Pd}_k(T_1)$ satisfy $s_c=0$ for every parameter cube $c$ that does not belong to $\mathscr{Y}$. Denote by $\mathrm{pj}^{-1}\mathscr{Y}$ the set of the parameter cubes $c$ of $T_2$ of the same kind for which $\mathrm{pj}(c)$ belongs to $\mathscr{Y}$. Since $\left(\mathrm{pj}^{*}s\right)_c=s_{\mathrm{pj}(c)}$, the family $\mathrm{pj}^{*}s$ vanishes at every cube that does not belong to $\mathrm{pj}^{-1}\mathscr{Y}$; the union of the cubes of $\mathrm{pj}^{-1}\mathscr{Y}$ is $\mathrm{pj}^{-1}(Y)$. The images under $\mathrm{pj}$ of two linked cubes of $T_2$ (Definition~\ref{def:parameter-cubes-polydisks}) are linked: a common lattice site is carried to a common lattice site, and a common face to a common face, because $2M<l_1$. Consequently the image of a link-component of $\mathrm{pj}^{-1}(Y)$ is link-connected and is contained in a single link-component of $Y$.

\emph{Property~\ref{it:renormalisation-group-8}.} Let $Y_1$ be a link-component of $Y$, and let $x$ be a row index of $\mathrm{Op}$ that is contained in $Y_1$. Let $\tilde s\in\mathrm{Pd}_k(T_1)$ vanish at every cube that does not belong to $\mathscr{Y}$ and agree with $s$ at the cubes of $Y_1$, and let $\tilde A\in\tilde{G}_k(T_1)$ agree with $A$ on the bonds that meet $Y_1$. A lift of $x$ is an index of $T_2$ whose image under $\mathrm{pj}$ is $x$. The map $\mathrm{pj}$ carries the support of a lift of $x$ onto $\operatorname{supp}x$, which lies in $Y_1\subseteq Y$; so every lift of $x$ is contained in $\mathrm{pj}^{-1}(Y)$. Choose a lift of $x$ and denote it again by $x$. As $l_2$ belongs to $\mathrm{Cl}$, Lemma~\ref{lem:cover-cubes-linked}(b) applies on $T_2$ to the set $\mathrm{pj}^{-1}\mathscr{Y}$: the lift is contained in exactly one link-component $\mathsf{Y}$ of $\mathrm{pj}^{-1}(Y)$. The image $\mathrm{pj}(\mathsf{Y})$ contains $\operatorname{supp}x$ and is contained in a link-component of $Y$; this link-component contains $x$, so it is $Y_1$, by Lemma~\ref{lem:cover-cubes-linked}(b) on $T_1$. Hence every cube of $\mathsf{Y}$ is mapped to a cube of $Y_1$, and every bond that meets $\mathsf{Y}$ to a bond that meets $Y_1$. It follows that $\mathrm{pj}^{*}\tilde s$ vanishes at every cube that does not belong to $\mathrm{pj}^{-1}\mathscr{Y}$ and agrees with $\mathrm{pj}^{*}s$ at the cubes of $\mathsf{Y}$, and that $\mathrm{pj}^{*}\tilde A$ agrees with $\mathrm{pj}^{*}A$ on the bonds that meet $\mathsf{Y}$. Property~\ref{it:renormalisation-group-8} of the system on $T_2$ therefore gives
\[
\widetilde{\mathrm{Op}}(s,A)(x,x'')=0\ \text{ unless $x''$ is contained in $\mathsf{Y}$},\qquad \widetilde{\mathrm{Op}}(s,A)(x,\cdot\,)=\widetilde{\mathrm{Op}}(\tilde s,\tilde A)(x,\cdot\,).
\]
By~\eqref{eq:renormalisation-flow-4}, every entry of the row of $\mathrm{Op}^{\mathrm{pj}}(s,A)$ at $x$ is a linear combination of entries of the row of $\widetilde{\mathrm{Op}}(s,A)$ at the lift, with the coefficients $\boldsymbol{\chi}(w)$, which do not depend on $(s,A)$. So $\mathrm{Op}^{\mathrm{pj}}(s,A)(x,\cdot\,)=\mathrm{Op}^{\mathrm{pj}}(\tilde s,\tilde A)(x,\cdot\,)$. If $\mathrm{Op}^{\mathrm{pj}}(s,A)(x,x')\ne0$, then by~\eqref{eq:renormalisation-flow-4} one has $\widetilde{\mathrm{Op}}(s,A)(x,x'+l_1\cdot w)\ne0$ for some $w$; by the last display the index $x'+l_1\cdot w$ is contained in $\mathsf{Y}$, and its image $x'$ is contained in $\mathrm{pj}(\mathsf{Y})\subseteq Y_1$.

\emph{Property~\ref{it:renormalisation-group-9}.} Let $x'$ be a column index of $\mathrm{Op}$ that is contained in $Y_1$, and let $\tilde s$, $\tilde A$ be as above. Every lift of $x'$ is contained in $\mathrm{pj}^{-1}(Y)$, since $\mathrm{pj}$ carries its support onto $\operatorname{supp}x'\subseteq Y_1$. Fix a lift of $x'$, denoted again by $x'$, and represent each row index $x$ by a lift, as in~\eqref{eq:renormalisation-flow-4}. Let $w\in\mathbb{Z}^3$. As $l_1\cdot w$ is a period of $T_1$, the index $x'+l_1\cdot w$ of $T_2$ is again a lift of the column index, hence contained in $\mathrm{pj}^{-1}(Y)$; since $l_2\in\mathrm{Cl}$, Lemma~\ref{lem:cover-cubes-linked}(b) on $T_2$ provides a unique link-component $\mathsf{Y}_w$ of $\mathrm{pj}^{-1}(Y)$ that contains this lift. The image $\mathrm{pj}(\mathsf{Y}_w)$ contains $\operatorname{supp}x'$ and lies in a link-component of $Y$, which contains $x'$ and is therefore $Y_1$, by Lemma~\ref{lem:cover-cubes-linked}(b) on $T_1$. Property~\ref{it:renormalisation-group-9} of the system on $T_2$, applied to the column index $x'+l_1\cdot w$, gives two facts. First, the entry $\widetilde{\mathrm{Op}}(s,A)(x,x'+l_1\cdot w)$ vanishes for every row index $x$ that is not contained in $\mathsf{Y}_w$. Secondly, the column $\widetilde{\mathrm{Op}}(s,A)(\cdot\,,x'+l_1\cdot w)$ is not altered if $\left(\mathrm{pj}^{*}s,\mathrm{pj}^{*}A\right)$ is replaced by $\left(\mathrm{pj}^{*}\tilde s,\mathrm{pj}^{*}\tilde A\right)$; indeed $\mathrm{pj}^{*}\tilde s$ vanishes at every cube that does not belong to $\mathrm{pj}^{-1}\mathscr{Y}$, since $\tilde s$ vanishes outside $\mathscr{Y}$, and, as $\mathrm{pj}$ maps the cubes of $\mathsf{Y}_w$ to cubes of $Y_1$ and the bonds that meet $\mathsf{Y}_w$ to bonds that meet $Y_1$, the family $\mathrm{pj}^{*}\tilde s$ agrees with $\mathrm{pj}^{*}s$ at the cubes of $\mathsf{Y}_w$ and $\mathrm{pj}^{*}\tilde A$ agrees with $\mathrm{pj}^{*}A$ on the bonds that meet $\mathsf{Y}_w$. By~\eqref{eq:renormalisation-flow-4}, the column $\mathrm{Op}^{\mathrm{pj}}(s,A)(\cdot\,,x')$ is therefore unchanged when $(s,A)$ is replaced by $(\tilde s,\tilde A)$. If $\mathrm{Op}^{\mathrm{pj}}(s,A)(x,x')\ne0$, then by~\eqref{eq:renormalisation-flow-4} one has $\widetilde{\mathrm{Op}}(s,A)(x,x'+l_1\cdot w)\ne0$ for some $w$; so the lift $x$ is contained in $\mathsf{Y}_w$, and its image is contained in $\mathrm{pj}(\mathsf{Y}_w)\subseteq Y_1$; that is, $x$ is contained in $Y_1$.

\emph{Property~\ref{it:renormalisation-group-10}.} Let $\mathrm{Op}$ be one of the families of property~\ref{it:renormalisation-group-10}, and let $x$ be a row index of $\mathrm{Op}$ that is contained in no link-component of $Y$. By Lemma~\ref{lem:cover-cubes-linked} on $T_1$, the support of $x$ is not contained in $Y$. Hence the support of a lift of $x$ is not contained in $\mathrm{pj}^{-1}(Y)$, and the lift is contained in no link-component of $\mathrm{pj}^{-1}(Y)$. By property~\ref{it:renormalisation-group-10} of the system on $T_2$ the row of $\widetilde{\mathrm{Op}}(s,A)$ at the lift vanishes, and by~\eqref{eq:renormalisation-flow-4} the row of $\mathrm{Op}^{\mathrm{pj}}(s,A)$ at $x$ vanishes as well.

\emph{Property~\ref{it:renormalisation-group-11}.} Now let $\mathrm{Op}$ be the family~\eqref{it:primitives-p3}, so that $\mathscr{Y}$ consists of cubes of $\pi_k$. Let $z$ be a fine site of $T_1$, let $\hat z$ be a fine site of $T_2$ with $\mathrm{pj}(\hat z)=z$, and let $K_z$ and $K_{\hat z}$ be the sets of the sites within the sup distance $2L$ of $z$, resp. of $\hat z$, as in Section~\ref{sec:lattice-geometry}. Because $4L+1<l_1$, the map $\mathrm{pj}$ is a bijection of $K_{\hat z}$ onto $K_z$; it is also a bijection of the set of the cubes of $\pi_k$ that meet $K_{\hat z}$ onto the set of the cubes of $\pi_k$ that meet $K_z$, all these cubes lying within the distance $2L+M<l_1/2$ of $\hat z$, resp. of $z$. Suppose that some cube of $\pi_k$ that meets $K_z$ does not belong to $\mathscr{Y}$. Then, by the definition of $\mathrm{pj}^{-1}\mathscr{Y}$, the cube of $\pi_k$ that meets $K_{\hat z}$ and is mapped onto it does not belong to $\mathrm{pj}^{-1}\mathscr{Y}$; so property~\ref{it:renormalisation-group-11} of the system on $T_2$ shows that the row of $H_k^{(2)}\left(\mathrm{pj}^{*}s,\mathrm{pj}^{*}A\right)$ at $\hat z$ vanishes, and by~\eqref{eq:renormalisation-flow-4} the row of $H_k^{\mathrm{pj}}(s,A)$ at $z$ vanishes.

\emph{Property~\ref{it:renormalisation-group-12}.} Let $\mathrm{Op}$ be $M_k$ or $S_{k,y}$, let $X\in D_k(T_1)$, and let $A,\tilde A\in\tilde{G}_k(T_1)$ agree on the bonds that meet $X$. Let $X'$ be a polymer of $T_2$ with $\mathrm{pj}(X')=X$. By property~\ref{it:renormalisation-group-12} of the system on $T_2$, the kernel of $\mathrm{Op}^{(2)}\left(X',\mathrm{pj}^{*}A\right)$ vanishes unless both of its indices lie in $X'$, and this operator depends on $\mathrm{pj}^{*}A$ only through its values on the bonds that meet $X'$. The image of a bond that meets $X'$ meets $\mathrm{pj}(X')=X$; so $\mathrm{pj}^{*}A$ and $\mathrm{pj}^{*}\tilde A$ agree on the bonds that meet $X'$, and, summing over these polymers $X'$, one obtains $\widetilde{\mathrm{Op}}(X,A)=\widetilde{\mathrm{Op}}(X,\tilde A)$. By~\eqref{eq:renormalisation-flow-4}, $\mathrm{Op}^{\mathrm{pj}}(X,A)=\mathrm{Op}^{\mathrm{pj}}(X,\tilde A)$. If $\mathrm{Op}^{\mathrm{pj}}(X,A)(x,x')\ne0$, then by~\eqref{eq:renormalisation-flow-4} there are a polymer $X'$ of $T_2$ with $\mathrm{pj}(X')=X$ and a vector $w\in\mathbb{Z}^3$ such that the kernel of $\mathrm{Op}^{(2)}\left(X',\mathrm{pj}^{*}A\right)$ does not vanish at the lift of $x$ and the index $x'+l_1\cdot w$; both of them lie in $X'$, and hence $x$ and $x'$ lie in $\mathrm{pj}(X')=X$.

\emph{Property~\ref{it:renormalisation-group-13}.} Let $\lambda$ be a real function on the sites of $T_1$, and let $\lambda^u$ be its restriction to the unit lattice $T^0_1$. The connection $A-d\lambda$ lies in $\tilde{G}_k(T_1)$ by Lemma~\ref{lem:invariances}, since this domain is invariant under real gauge transformations; so the projected families are defined at $A-d\lambda$. Each relation of property~\ref{it:renormalisation-group-13} has the form
\[
\mathrm{Op}(s,A-d\lambda)=\exp\left(i e_k\cdot\lambda'\right)\mathrm{Op}(s,A)\exp\left(-i e_k\cdot\lambda''\right),
\]
with $(\lambda',\lambda'')=\left(\lambda,\lambda^u\right)$ for the family~\eqref{it:primitives-p3} and $(\lambda',\lambda'')=\left(\lambda^u,\lambda^u\right)$ for $\Gamma_k$, or the same form with the polymer $X$ in place of $s$, with $(\lambda',\lambda'')=\left(\lambda^u,\lambda^u\right)$ for $M_k$ and $(\lambda',\lambda'')=(\lambda,\lambda)$ for $S_{k,y}$; the exponentials act as multiplication operators. One has $\mathrm{pj}^{*}(A-d\lambda)=\mathrm{pj}^{*}A-d\,\mathrm{pj}^{*}\lambda$; the function $\mathrm{pj}^{*}\lambda$ is real on the sites of $T_2$, and its restriction to the unit lattice $T^0_2$ is $\mathrm{pj}^{*}\lambda^u$. Let $h$ be a spinor on $T_1$ if $\lambda'=\lambda$, and on $T^0_1$ if $\lambda'=\lambda^u$. Then
\[
\exp\left(\pm i e_k\cdot\mathrm{pj}^{*}\lambda'\right)\mathrm{pj}^{*}h=\mathrm{pj}^{*}\left(\exp\left(\pm i e_k\cdot\lambda'\right)h\right),
\]
because under the identifications of Section~\ref{sec:two-tori} both sides are the function $\exp\left(\pm i e_k\cdot\lambda'\right)h$ on $\mathbb{L}_k$; the same holds with $\lambda''$ in place of $\lambda'$. Now let $f$ be an arbitrary function on $T_1$ on which the family under consideration acts, that is, a spinor on $T^0_1$ for the family~\eqref{it:primitives-p3}, for $\Gamma_k$ and for $M_k$, and a spinor on $T_1$ for $S_{k,y}$. Definition~\ref{def:projection}, property~\ref{it:renormalisation-group-13} of the system on $T_2$ at the data $\left(\mathrm{pj}^{*}s,\mathrm{pj}^{*}A,\mathrm{pj}^{*}\lambda\right)$, and the last display, used with $\lambda''$ for $h=f$ and with $\lambda'$ for $h=\mathrm{Op}^{\mathrm{pj}}(s,A)\exp\left(-i e_k\cdot\lambda''\right)f$, give
\begin{align*}
\mathrm{pj}^{*}\left(\mathrm{Op}^{\mathrm{pj}}(s,A-d\lambda)f\right)
&=\mathrm{Op}^{(2)}\left(\mathrm{pj}^{*}s,\mathrm{pj}^{*}A-d\,\mathrm{pj}^{*}\lambda\right)\mathrm{pj}^{*}f\\
&=\exp\left(i e_k\cdot\mathrm{pj}^{*}\lambda'\right)\mathrm{Op}^{(2)}\left(\mathrm{pj}^{*}s,\mathrm{pj}^{*}A\right)\exp\left(-i e_k\cdot\mathrm{pj}^{*}\lambda''\right)\mathrm{pj}^{*}f\\
&=\exp\left(i e_k\cdot\mathrm{pj}^{*}\lambda'\right)\mathrm{Op}^{(2)}\left(\mathrm{pj}^{*}s,\mathrm{pj}^{*}A\right)\mathrm{pj}^{*}\left(\exp\left(-i e_k\cdot\lambda''\right)f\right)\\
&=\mathrm{pj}^{*}\left(\exp\left(i e_k\cdot\lambda'\right)\mathrm{Op}^{\mathrm{pj}}(s,A)\exp\left(-i e_k\cdot\lambda''\right)f\right).
\end{align*}
For $M_k$ and $S_{k,y}$ the operator $\widetilde{\mathrm{Op}}(X,A)$ is the sum of the operators $\mathrm{Op}^{(2)}\left(X',\mathrm{pj}^{*}A\right)$ over the polymers $X'$ of $T_2$ with $\mathrm{pj}(X')=X$; property~\ref{it:renormalisation-group-13} on $T_2$ applies to each summand, and the four equalities above hold with $X$ in place of $s$ and of $\mathrm{pj}^{*}s$ and with this sum in place of $\mathrm{Op}^{(2)}$. Since $\mathrm{pj}^{*}$ is injective, the projected families satisfy the relations of property~\ref{it:renormalisation-group-13}.

\emph{Property~\ref{it:renormalisation-group-14}.} Let $g=(S,a)\in\mathbb{G}_k^L(T_1)$, acting as in Section~\ref{sec:dirac-matrices-lattice}, and let $a_1\in LM\mathbb{Z}^3$ be a representative of $a$. Then $g$ is induced by $S$ and by the map $x\mapsto r(S)\cdot x+a_1$ of $\mathbb{L}_k$. If $a_1$ is replaced by $a_1+l_1\cdot w$ with $w\in\mathbb{Z}^3$, the action on spinors is multiplied by the sign $\boldsymbol{\chi}(w)$; this sign occurs on both sides of each relation of property~\ref{it:renormalisation-group-14}. The matrix $S$ and the same map of $\mathbb{L}_k$ induce an element $g_2$ of $\mathbb{G}_k^L(T_2)$, and $g_2$ belongs to $\mathbb{G}^L_{k,3}(T_2)$ if $S\in\mathrm{Sp}_3$, that is, if $g\in\mathbb{G}^L_{k,3}(T_1)$; for $S_{k,y}$ the relations concern only such $g$. On functions, parameter families and connections one has $g_2\mathrm{pj}^{*}=\mathrm{pj}^{*}g$, and $\mathrm{pj}(g_2X')=g\,\mathrm{pj}(X')$ for every polymer $X'$ of $T_2$. The connection $gA$ lies in $\tilde{G}_k(T_1)$ by Lemma~\ref{lem:invariances}, the domain being invariant under $\mathbb{G}_k$. Let $f$ be an arbitrary function on $T_1$ on which the family under consideration acts: a bond function on $T^0_1$ for the family~\eqref{it:primitives-p1}, an element of $V_k(T_1)$ for the family~\eqref{it:primitives-p2}, a spinor on $T^0_1$ for the family~\eqref{it:primitives-p3}, for $\Gamma_k$ and for $M_k$, and a spinor on $T_1$ for $S_{k,y}$. By Definition~\ref{def:projection}, the relation $g_2\mathrm{pj}^{*}=\mathrm{pj}^{*}g$ and property~\ref{it:renormalisation-group-14} of the system on $T_2$,
\begin{align*}
\mathrm{pj}^{*}\left(\mathrm{Op}^{\mathrm{pj}}(gs,gA)\,gf\right)
&=\mathrm{Op}^{(2)}\left(g_2\mathrm{pj}^{*}s,g_2\mathrm{pj}^{*}A\right)g_2\mathrm{pj}^{*}f\\
&=g_2\,\mathrm{Op}^{(2)}\left(\mathrm{pj}^{*}s,\mathrm{pj}^{*}A\right)\mathrm{pj}^{*}f
=\mathrm{pj}^{*}\left(g\,\mathrm{Op}^{\mathrm{pj}}(s,A)f\right).
\end{align*}
The map $X'\mapsto g_2X'$ is a bijection of the set of the polymers $X'$ of $T_2$ with $\mathrm{pj}(X')=X$ onto the set of the polymers $X''$ of $T_2$ with $\mathrm{pj}(X'')=gX$. Hence
\begin{align*}
\mathrm{pj}^{*}\left(\mathrm{Op}^{\mathrm{pj}}(gX,gA)\,gf\right)
&=\sum_{X'\colon\mathrm{pj}(X')=X}\mathrm{Op}^{(2)}\left(g_2X',g_2\mathrm{pj}^{*}A\right)g_2\mathrm{pj}^{*}f\\
&=g_2\sum_{X'\colon\mathrm{pj}(X')=X}\mathrm{Op}^{(2)}\left(X',\mathrm{pj}^{*}A\right)\mathrm{pj}^{*}f
=\mathrm{pj}^{*}\left(g\,\mathrm{Op}^{\mathrm{pj}}(X,A)f\right).
\end{align*}
As $\mathrm{pj}^{*}$ is injective, the projected families satisfy the relations of property~\ref{it:renormalisation-group-14}.

\emph{Property~\ref{it:renormalisation-group-15}.} Let $\mathrm{Op}$ be $\Gamma_k$ or $M_k$. By the proof of part~(a), $\widetilde{\mathrm{Op}}(s,A)$ and $\widetilde{\mathrm{Op}}(X,A)$ commute with the translations $T_g$, $g\in l_1\mathbb{Z}^3$, and by Definition~\ref{def:projection} one has $\widetilde{\mathrm{Op}}(s,A)\,\mathrm{pj}^{*}=\mathrm{pj}^{*}\,\mathrm{Op}^{\mathrm{pj}}(s,A)$ and $\widetilde{\mathrm{Op}}(X,A)\,\mathrm{pj}^{*}=\mathrm{pj}^{*}\,\mathrm{Op}^{\mathrm{pj}}(X,A)$. So the rule~\eqref{eq:renormalisation-flow-5} applies to these operators. The connection $-A$ lies in $\tilde{G}_k(T_1)$ by Lemma~\ref{lem:invariances}, and $\mathrm{pj}^{*}(-A)=-\mathrm{pj}^{*}A$; the constant spinor matrix $\mathsf{C}$ commutes with $\mathrm{pj}^{*}$. With property~\ref{it:renormalisation-group-15} of the system on $T_2$, applied for $M_k$ to every polymer $X'$ with $\mathrm{pj}(X')=X$, one obtains
\begin{align*}
\mathrm{pj}^{*}\left(\Gamma_k^{\mathrm{pj}}(s,A)\right)^T
&=\left(\widetilde{\mathrm{Op}}(s,A)\right)^T\mathrm{pj}^{*}
=\mathsf{C}^{-1}\,\widetilde{\mathrm{Op}}(s,-A)\,\mathsf{C}\,\mathrm{pj}^{*}
=\mathrm{pj}^{*}\,\mathsf{C}^{-1}\,\Gamma_k^{\mathrm{pj}}(s,-A)\,\mathsf{C},\\
\mathrm{pj}^{*}\left(M_k^{\mathrm{pj}}(X,A)\right)^T
&=\left(\widetilde{\mathrm{Op}}(X,A)\right)^T\mathrm{pj}^{*}
=\mathsf{C}^{-1}\,\widetilde{\mathrm{Op}}(X,-A)\,\mathsf{C}\,\mathrm{pj}^{*}
=\mathrm{pj}^{*}\,\mathsf{C}^{-1}\,M_k^{\mathrm{pj}}(X,-A)\,\mathsf{C}.
\end{align*}
Since $\mathrm{pj}^{*}$ is injective, the projected families satisfy the relations of property~\ref{it:renormalisation-group-15}.

\emph{Property~\ref{it:renormalisation-group-16}.} Complex conjugation, of functions as well as of parameter families and of connections, is taken value by value, and the constant matrix $\varkappa$ acts on a spinor value by value as well; since under the identifications of Section~\ref{sec:two-tori} the map $\mathrm{pj}^{*}$ does not change values, these operations commute with $\mathrm{pj}^{*}$. The connections $\overline{A}$ and $-A$ lie in $\tilde{G}_k(T_1)$ by Lemma~\ref{lem:invariances}. Let $\mathrm{Op}$ be the family~\eqref{it:primitives-p3} or $\Gamma_k$, and let $f$ be an arbitrary spinor on $T^0_1$. Property~\ref{it:renormalisation-group-16} of the system on $T_2$, applied to $\mathrm{pj}^{*}f$, gives
\[
\mathrm{pj}^{*}\left(\overline{\mathrm{Op}^{\mathrm{pj}}\left(\overline{s},\overline{A}\right)}\,f\right)
=\overline{\widetilde{\mathrm{Op}}\left(\overline{s},\overline{A}\right)}\,\mathrm{pj}^{*}f
=\varkappa^{-1}\,\widetilde{\mathrm{Op}}(s,-A)\,\varkappa\,\mathrm{pj}^{*}f
=\mathrm{pj}^{*}\left(\varkappa^{-1}\,\mathrm{Op}^{\mathrm{pj}}(s,-A)\,\varkappa\,f\right).
\]
For $M_k$, which acts on the same spinors, the same equalities hold with $X$ in place of $s$ and of $\overline{s}$, the operator $\widetilde{\mathrm{Op}}\left(X,\overline{A}\right)$ being the sum of the operators $M_k^{(2)}\left(X',\mathrm{pj}^{*}\overline{A}\right)$ over the polymers $X'$ with $\mathrm{pj}(X')=X$, to each of which property~\ref{it:renormalisation-group-16} on $T_2$ applies. Since $\mathrm{pj}^{*}$ is injective, the projected families~\eqref{it:primitives-p3} and~\eqref{it:primitives-p4} satisfy the relation of property~\ref{it:renormalisation-group-16}. For the families~\eqref{it:primitives-p1} and~\eqref{it:primitives-p2} the coefficients $\boldsymbol{\chi}(w)$ in~\eqref{eq:renormalisation-flow-4} are equal to $1$, so that the entries~\eqref{eq:renormalisation-flow-4} of the projections are sums of entries of $\mathrm{Op}^{(2)}\left(\mathrm{pj}^{*}s\right)$. For real $s$ the family $\mathrm{pj}^{*}s$ is real; then these entries are real, and so are the entries of the projections.

\emph{Property~\ref{it:renormalisation-group-17}.} Let $k=0$. The assertions of property~\ref{it:renormalisation-group-17} on the family~\eqref{it:primitives-p2} and on the families $\Gamma_0$ and $S_{0,y}$ are contained in what has been proved above, taken at $k=0$, with property~\ref{it:renormalisation-group-7} in the form that it has at this level: the kernel bound for $S_{0,y}(X,A)$ with the factor $(1+y)^{-1}$, which was derived on $T_1$ together with property~\ref{it:renormalisation-group-7}. The assertions on the two diagonal families $H_0(s')$ and $H_0(s'',A)$ are obtained together with part~(c).

\emph{Part~(c).} In property~\ref{it:renormalisation-group-17}, the diagonal entry of $H_0(s')$ at a bond $b$ is the product of the numbers $s'_c$ over the cubes $c$ of $\pi'_0$ that contain $b$, and the diagonal entry of $H_0(s'',A)$ at a site $x$ is the product of the numbers $s''_c$ over the cubes $c$ of $\pi_0$ that meet $K_x$. These cubes lie within the distance $2M+2L<l_1/2$ of $b$, resp. of $x$. Hence $\mathrm{pj}$ maps the cubes of $\pi'_0$ that contain a lift of $b$ bijectively onto the cubes of $\pi'_0$ that contain $b$, and the cubes of $\pi_0$ that meet $K_{\hat x}$, for a lift $\hat x$ of $x$, bijectively onto the cubes of $\pi_0$ that meet $K_x$. Together with $\left(\mathrm{pj}^{*}s\right)_c=s_{\mathrm{pj}(c)}$ this gives, for every bond function $f$ on $T^0_1$ in the first identity and every spinor $f$ on $T^0_1$ in the second,
\[
H_0^{(2)}\left(\mathrm{pj}^{*}s'\right)\mathrm{pj}^{*}f=\mathrm{pj}^{*}\left(H_0^{(1)}(s')f\right),\qquad H_0^{(2)}\left(\mathrm{pj}^{*}s'',\mathrm{pj}^{*}A\right)\mathrm{pj}^{*}f=\mathrm{pj}^{*}\left(H_0^{(1)}(s'',A)f\right),
\]
where the superscripts $(1)$ and $(2)$ mark the families of property~\ref{it:renormalisation-group-17} on $T_1$ and on $T_2$. By Definition~\ref{def:projection} and the injectivity of $\mathrm{pj}^{*}$, the projections of the two families on $T_2$ are the two families on $T_1$. This proves~(c).

On $T_1$ the two families have the properties that property~\ref{it:renormalisation-group-17} lists for them: $H_0(s')$ by Proposition~\ref{prop:family-level-zero}, and $H_0(s'',A)$ by Theorem~\ref{thm:properties-level-zero}(b). In particular, the bound of property~\ref{it:renormalisation-group-3} for $H_0(s')$ and the bounds of property~\ref{it:renormalisation-group-5} for $H_0(s'',A)$ hold with the constant $4\exp(2\gamma_w)\cdot M^{27\alpha_0}$ of property~\ref{it:renormalisation-group-17}; the factor $64$ of~\eqref{eq:renormalisation-flow-6} does not occur for these two families. By the definition of $C_w^1$ in Section~\ref{sec:photon-operators}, this constant is at most $C_w^1$, and $C_w^1\le C_w=64 K_0\cdot C_w^1$ because $K_0\ge1$ (Section~\ref{sec:absolute-constants}). Thus the projected system satisfies property~\ref{it:renormalisation-group-17} for the parameters $\mathrm{par}$ of part~(b), and the proof of~(b) is complete.

This completes the proof of Theorem~\ref{thm:properties-projection}. Property~\ref{it:renormalisation-group-18} follows.
\end{proof}

\subsubsection{Rows of the primitives inside a component}\label{sec:rows-primitives-inside}

Let $W$ be a union of $M$-cubes of $T_k$. A parameter family $s = (s_c)$ is said to \emph{vanish off} $W$ if $s_c = 0$ for each parameter cube $c$ that is not contained in $W$. For every parameter family $s$ the family $s \cdot \mathbf{1}_W$ is defined by
\[
(s \cdot \mathbf{1}_W)_c := \begin{cases} s_c & \text{if the parameter cube } c \text{ is contained in } W, \\ 0 & \text{otherwise.} \end{cases}
\]
A \emph{component} of $W$ is a union of $M$-cubes of $W$ that is connected through faces and is maximal with this property.

\begin{lemma}[Locality of the rows]\label{lem:locality-rows}
Assume $l \in \mathrm{Cl}$ and $M > 4L$. Consider a system of primitives of level $k$ on $T_k$ with the properties~\ref{it:renormalisation-group-8} and~\ref{it:renormalisation-group-12}, and denote by $\mathrm{Op}$ one of its families~\eqref{it:primitives-p1}--\eqref{it:primitives-p3} and $\Gamma_k$. Fix a union $W$ of $M$-cubes together with a component $W_1$ of $W$. Let $s \in \mathrm{Pd}_k$, a parameter family of the kind of $\mathrm{Op}$, vanish off $W$, and let the connections $A, A_1 \in \tilde{G}_k$ agree on every bond meeting $W_1$; let $x$ be a row index of $\mathrm{Op}$ with $\square(x)^{\sim c_l} \subset W_1$. Then
\[
\mathrm{Op}(s,A)(x,x') = 0 \quad \text{unless } x' \text{ is contained in } W_1, \qquad \mathrm{Op}(s,A)(x,\cdot\,) = \mathrm{Op}\left(s \cdot \mathbf{1}_{W_1}, A_1\right)(x,\cdot\,).
\]
Moreover, for $X \in D_k$ and $y \ge 0$, the kernels of $M_k(X,A)$ and of $S_{k,y}(X,A)$ vanish at every pair of indices that do not both lie in $X$; and $M_k(X,A) = M_k(X,A_1)$, $S_{k,y}(X,A) = S_{k,y}(X,A_1)$ whenever $A$ and $A_1$ agree on every bond meeting $X$.
\end{lemma}

\begin{proof}
Write $\mathscr{Y}$ for the set of those parameter cubes of the kind of $\mathrm{Op}$ (Section~\ref{sec:two-tori}) which lie in $W$, and $Y$ for the union of the cubes of $\mathscr{Y}$. Since $s$ vanishes off $W$, the entry $s_c$ is zero at every parameter cube $c$ of this kind outside $\mathscr{Y}$.

\emph{A cube of $\mathscr{Y}$ around the support of $x$.} It is first shown that $\operatorname{supp} x$ lies in a parameter cube $c_x$ of the kind of $\mathrm{Op}$ that is contained in $\square(x)^{\sim c_l}$; here $c_l = 2L$ (Definition~\ref{def:parameter-cubes-polydisks}). For a point $p$ of $\operatorname{supp} x$ let $\square(p)$ denote the $M$-cube that contains $p$. Then $\square(p) \subset \square(x)$, and therefore $\square(p)^{\sim i} \subset \square(x)^{\sim c_l}$ for $0 \le i \le c_l$. By fact~\eqref{it:renormalisation-group-g0} of Section~\ref{sec:cubes-parameter-cubes}, no site lies on a face of an $M$-cube, the $M$-cubes partition the sites, and every $M$-cube lies in one $LM$-cube. Since $L$ is odd, every face of an $LM$-cube is a union of faces of $M$-cubes. Hence $\lvert p - u \rvert < M/2$ for the centre $u \in M\mathbb{Z}^3$ of $\square(p)$, and $\lvert p - u' \rvert < L \cdot M/2$ for the centre $u' \in LM\mathbb{Z}^3$ of the $LM$-cube that contains $p$.

If $x$ is a site, the kind is $\pi_k$, and one puts $c_x := \square(x)$, the $M$-cube of $x$.

If $x$ is a bond, the kind is $\pi'_k$. Let $p$ and $p'$ be the end points of $x$, let $u$ be the centre of $\square(p)$, and let $c_x := \{ z : \lvert z - u \rvert \le M \}$, a cube of $\pi'_k$. As $\lvert p' - p \rvert \le 1$, one has $\lvert p' - u \rvert < 1 + M/2 \le M$, so that both end points of $x$ lie in $c_x$. The cube $c_x$ lies in the union of the $M$-cubes whose centres have sup-distance at most $M$ from $u$ (fact~\eqref{it:renormalisation-group-g1} of Section~\ref{sec:cubes-parameter-cubes}), and this union is $\square(p)^{\sim 1}$. Hence $c_x \subset \square(p)^{\sim 1} \subset \square(x)^{\sim c_l}$.

If $x = \Upsilon$ is a basis element, the kind is $\pi''_k$. Let $p \in \operatorname{supp} \Upsilon$, let $u'$ be the centre of the $LM$-cube that contains $p$, and let $c_x := \{ z : \lvert z - u' \rvert \le L \cdot M \}$, a cube of $\pi''_k$. By Proposition~\ref{prop:basis-map-box}, $\operatorname{supp} \Upsilon$ lies in one $L$-block or consists of the two end points of one bond; hence $\lvert q - p \rvert < L$ for every $q \in \operatorname{supp} \Upsilon$, and
\[
\lvert q - u' \rvert \le \lvert q - p \rvert + \lvert p - u' \rvert < L + \frac{L \cdot M}{2} \le L \cdot M .
\]
So $\operatorname{supp} \Upsilon \subset c_x$. Let $z$ be a lattice site of $c_x$, and let $u_z$ and $u_p$ be the centres of the $M$-cubes that contain $z$ and $p$. One has $\lvert z - p \rvert \le \lvert z - u' \rvert + \lvert u' - p \rvert < 3L \cdot M/2$, and therefore
\[
\lvert u_z - u_p \rvert \le \lvert u_z - z \rvert + \lvert z - p \rvert + \lvert p - u_p \rvert < \left( \frac{3L}{2} + 1 \right) M .
\]
Each coordinate of $u_z - u_p$ equals $j \cdot M$ with an integer $j$ such that $\lvert j \rvert < (3L + 2)/2$. As $L$ is odd, $3L + 2$ is odd, and so $\lvert j \rvert \le (3L + 1)/2 \le 2L$. Thus the $M$-cube of $z$ lies in $\square(p)^{\sim 2L}$. So every lattice site of $c_x$ lies in $\square(p)^{\sim 2L} \subset \square(x)^{\sim c_l}$.

\emph{The link-component of $x$.} In each of the three cases every lattice site of $c_x$ lies in $\square(x)^{\sim c_l} \subset W_1 \subset W$. So $c_x$ is contained in $W$, that is, $c_x \in \mathscr{Y}$. Consequently $\operatorname{supp} x \subset Y$, and $x$ is contained in the link-component $Y_1$ of $Y$ that contains $c_x$.

\emph{The inclusion $Y_1 \subset W_1$.} Let $\mathrm{pol}(Y_1)$ be the union of the $M$-cubes that meet a cube of $\mathscr{Y}$ contained in $Y_1$. By fact~\eqref{it:renormalisation-group-g2} of Section~\ref{sec:cubes-parameter-cubes}, $\mathrm{pol}(Y_1)$ is a polymer, that is, a union of $M$-cubes connected through faces, and $Y_1 \subset \mathrm{pol}(Y_1)$. An $M$-cube that meets a cube $c \in \mathscr{Y}$ has a lattice site in common with $c$ (Section~\ref{sec:lattice-geometry}; fact~\eqref{it:renormalisation-group-g1} of Section~\ref{sec:cubes-parameter-cubes}). This site lies in $W$, because $c$ is contained in $W$; since $W$ is a union of $M$-cubes and the $M$-cubes partition the sites, the $M$-cube lies in $W$. Hence $\mathrm{pol}(Y_1)$ is a union of $M$-cubes of $W$ that is connected through faces. It has an $M$-cube in common with $W_1$, namely the $M$-cube of any point of $\operatorname{supp} x$, because
\[
\operatorname{supp} x \subset Y_1 \subset \mathrm{pol}(Y_1), \qquad \operatorname{supp} x \subset \square(x) \subset W_1 .
\]
So $\mathrm{pol}(Y_1) \cup W_1$ is a union of $M$-cubes of $W$ that is connected through faces, and the maximality of $W_1$ gives $\mathrm{pol}(Y_1) \subset W_1$. In particular $Y_1 \subset W_1$.

\emph{Conclusion.} The family $s$ vanishes at the parameter cubes outside $\mathscr{Y}$, and the row index $x$ is contained in the link-component $Y_1$ of $Y$. By property~\ref{it:renormalisation-group-8}, $\mathrm{Op}(s,A)(x,x') = 0$ unless $x'$ is contained in $Y_1$; as $Y_1 \subset W_1$, the first assertion follows. The family $s \cdot \mathbf{1}_{W_1}$ lies in $\mathrm{Pd}_k$, each of its entries being an entry of $s$ or zero. It vanishes at every parameter cube outside $\mathscr{Y}$, since such a cube is not contained in $W$ and hence not in $W_1$; and it coincides with $s$ at every cube of $\mathscr{Y}$ contained in $Y_1$, since such a cube is contained in $W_1$. The connection $A_1$ coincides with $A$ on the bonds that meet $Y_1$, because these bonds meet $W_1$. The facts just established, together with $A_1 \in \tilde{G}_k$, say that $\tilde{s} := s \cdot \mathbf{1}_{W_1}$ and $\tilde{A} := A_1$ satisfy the requirements that property~\ref{it:renormalisation-group-8} imposes on the parameter family and on the connection with which the pair $(s,A)$ is compared. Property~\ref{it:renormalisation-group-8}, applied to the pair $(s,A)$ and to the comparison pair $(\tilde{s},\tilde{A})$, therefore gives
\[
\mathrm{Op}(s,A)(x,\cdot\,) = \mathrm{Op}(\tilde{s},\tilde{A})(x,\cdot\,) = \mathrm{Op}\left(s \cdot \mathbf{1}_{W_1}, A_1\right)(x,\cdot\,).
\]
The assertions on $M_k(X,A)$ and $S_{k,y}(X,A)$ are property~\ref{it:renormalisation-group-12}.
\end{proof}

Proposition~\ref{prop:rows-component} follows from Lemma~\ref{lem:locality-rows}.

\subsubsection{The one-step hypothesis below the stop level}\label{sec:one-step-hypothesis}

\begin{lemma}[The assembled systems of primitives]\label{lem:assembled-systems-primitives}
Let $k \ge 0$ be an integer, $l \in \mathrm{Cl}$, $\mathfrak{m} \in (0, \mathfrak{m}_0]$ and $e_k \in (0, e_0]$. Consider on $T_k$ the families of Theorem~\ref{thm:photon-operators-weakened}, namely \eqref{eq:operators-photon-21} if $k \ge 1$, the family $H_0(s')$ of property~\ref{it:renormalisation-group-17} if $k = 0$, and \eqref{eq:operators-photon-22}; and the families of Theorem~\ref{thm:fermion-operators-weakened} at $\kappa' = \kappa_w + \kappa_0$, namely those of Definition~\ref{def:weakened-families-level} if $k \ge 1$ and those of Definition~\ref{def:families-level-zero} if $k = 0$. These families form a system of primitives of level $k$ on $T_k$ with the properties~\ref{it:renormalisation-group-1}--\ref{it:renormalisation-group-17} for $\mathrm{par}^0$, and hence also for $\mathrm{par}$.
\end{lemma}

\begin{proof}
By stage~\ref{it:order-choice-6}, $M \ge M^F(b, \alpha, L, \alpha_0, \kappa_w + \kappa_0)$ and $M \ge M^P(\alpha, L, \alpha_0)$. By Section~\ref{sec:photon-operators}, $\alpha_0 = \min(\alpha^F, \alpha^P)$, so that $\alpha_0 \le \alpha^F$ and $\alpha_0 \le \alpha^P$. By condition~\eqref{eq:renormalisation-group-A10}, $e_0 \le e^F_0$, so that $e_k \in (0, e^F_0]$. Hence the hypotheses of Theorem~\ref{thm:photon-operators-weakened}, and those of Theorem~\ref{thm:fermion-operators-weakened} with $\kappa' = \kappa_w + \kappa_0$, hold for the data $(k, l, \mathfrak{m}, e_k)$. They imply the hypotheses of the statements on the families named in the lemma: Theorem~\ref{thm:properties-weakened-photon}, Proposition~\ref{prop:family-level-zero}, and Theorem~\ref{thm:fermion-operators} with $\kappa' = \kappa_w + \kappa_0$.

The family~\eqref{eq:operators-photon-21}, for $k \ge 1$, and the family~\eqref{eq:operators-photon-22}, for $k \ge 0$, are families~\eqref{it:primitives-p1} and~\eqref{it:primitives-p2} of level $k$; by Theorem~\ref{thm:properties-weakened-photon}, for the parameter set $(\alpha_0, \gamma^P, \cdot\,; C^P)$, the first has the properties~\ref{it:renormalisation-group-1}--\ref{it:renormalisation-group-3}, \ref{it:renormalisation-group-8}--\ref{it:renormalisation-group-10}, \ref{it:renormalisation-group-14} and~\ref{it:renormalisation-group-16}, and the second has the properties~\ref{it:renormalisation-group-1}, \ref{it:renormalisation-group-2}, \ref{it:renormalisation-group-4}, \ref{it:renormalisation-group-8}--\ref{it:renormalisation-group-10}, \ref{it:renormalisation-group-14} and~\ref{it:renormalisation-group-16}. By Proposition~\ref{prop:family-level-zero}, the family $H_0(s')$ of property~\ref{it:renormalisation-group-17}, which is the family~\eqref{it:primitives-p1} at $k = 0$, has the properties~\ref{it:renormalisation-group-1}, \ref{it:renormalisation-group-2}, \ref{it:renormalisation-group-8}--\ref{it:renormalisation-group-10}, \ref{it:renormalisation-group-14} and~\ref{it:renormalisation-group-16}, and for every $\gamma > 0$ it has property~\ref{it:renormalisation-group-3} with the rate $\gamma$ and the constant $4 \exp(2 \gamma) \cdot M^{27 \alpha_0}$. By Theorem~\ref{thm:fermion-operators}(ii), the families of Definition~\ref{def:weakened-families-level} if $k \ge 1$, and those of Definition~\ref{def:families-level-zero} if $k = 0$, are families~\eqref{it:primitives-p3}, \eqref{it:primitives-p4} and~\eqref{it:primitives-p5} of level $k$ (for $k = 0$: $H_0(s'',A)$, $\Gamma_0$ and $S_{0,y}$) with the properties~\ref{it:renormalisation-group-1}, \ref{it:renormalisation-group-2} and~\ref{it:renormalisation-group-5}--\ref{it:renormalisation-group-17} for $(\alpha_0, \gamma^F, \kappa_w + \kappa_0; C^F)$. Together these families form a system of primitives of level $k$ on $T_k$ (Section~\ref{sec:primitive-operators}), and each of the properties~\ref{it:renormalisation-group-1}--\ref{it:renormalisation-group-17} is, for every family to which it refers, among the properties just listed.

By Section~\ref{sec:photon-operators}, $\gamma_w = \min(\gamma^F, \gamma^P)$ and $C_w^1 \ge C_w^0 = \max(C^F, C^P)$. The bounds in the properties~\ref{it:renormalisation-group-3}--\ref{it:renormalisation-group-7} remain valid when the constant is replaced by a larger one and the rate $\gamma$ by a smaller one; so each of these bounds, obtained with $(C^F, \gamma^F)$ or with $(C^P, \gamma^P)$, holds with $(C_w^1, \gamma_w)$. For the two diagonal families of level $0$, property~\ref{it:renormalisation-group-17} prescribes the constant $4 \exp(2 \gamma) \cdot M^{27 \alpha_0}$, where $\gamma$ is the rate of the parameter set; for $\mathrm{par}^0$ this constant is $4 \exp(2 \gamma_w) \cdot M^{27 \alpha_0}$, and $\gamma_w > 0$ because $\gamma^F$ and $\gamma^P$ are positive. Proposition~\ref{prop:family-level-zero}, applied with $\gamma = \gamma_w$, gives property~\ref{it:renormalisation-group-3} for $H_0(s')$ with the rate $\gamma_w$ and the constant $4 \exp(2 \gamma_w) \cdot M^{27 \alpha_0}$. Theorem~\ref{thm:properties-level-zero}(b), which holds at $k = 0$ under the hypotheses of Theorem~\ref{thm:fermion-operators} (proof of Theorem~\ref{thm:fermion-operators}(ii)), gives property~\ref{it:renormalisation-group-5} for $H_0(s'',A)$ with the constant $M^{8 \alpha_0}$; property~\ref{it:renormalisation-group-5} contains no rate, and
\[
M^{8 \alpha_0} \le M^{27 \alpha_0} \le 4 \exp(2 \gamma_w) \cdot M^{27 \alpha_0},
\]
since $M \ge 1$, $\alpha_0 > 0$ and $\gamma_w > 0$. So the two diagonal families satisfy the bounds that property~\ref{it:renormalisation-group-17} requires for $\mathrm{par}^0$. By the definition of $C_w^1$ in Section~\ref{sec:photon-operators}, $4 \exp(2 \gamma_w) \cdot M^{27 \alpha_0} \le C_w^1$; hence these two bounds also hold with the constant $C_w^1$, that is, $H_0(s')$ has property~\ref{it:renormalisation-group-3} and $H_0(s'',A)$ has property~\ref{it:renormalisation-group-5} with $(C_w^1, \gamma_w)$. This proves the properties~\ref{it:renormalisation-group-1}--\ref{it:renormalisation-group-17} for $\mathrm{par}^0 = (\alpha_0, \gamma_w, \kappa_w + \kappa_0; C_w^1)$ (Section~\ref{sec:primitive-operators}). Finally, $\kappa_w \le \kappa_w + \kappa_0$, and $C_w^1 \le C_w = 64 K_0 \cdot C_w^1$ since $K_0 \ge 1$ (Section~\ref{sec:order-choice}); the bounds in the properties~\ref{it:renormalisation-group-6} and~\ref{it:renormalisation-group-7} remain valid when the rate $\kappa'$ is replaced by a smaller one. So the properties~\ref{it:renormalisation-group-1}--\ref{it:renormalisation-group-17} hold for $\mathrm{par} = (\alpha_0, \gamma_w, \kappa_w; C_w)$.
\end{proof}

\begin{proposition}[The data along the flow]\label{prop:data-along-flow}
Assume hypothesis~\ref{def:hypothesis-flow}.
\begin{enumerate}
\item[(a)] Hypothesis~\ref{def:hypothesis-flow} holds with $R_{\mathrm{top}}$ in place of $R$ and with the same $e$ and $N$.
\item[(b)] Let $0 \le k < k^{*}(R)$. Hypothesis~\ref{def:hypothesis-one-step} is fulfilled by the data $(k, l_k, \bar{m} \cdot a_k, e_k)$ together with the system of primitives that Definition~\ref{def:periodisation-projection} provides on $T_k(R)$. Here $\bar{m} \cdot a_k$ is the running mass of level $k$, equal to $\mathfrak{m}_k$ if $k \ge 1$ and to $\bar{m} \cdot \eta$ if $k = 0$, and $e_k$ is the running charge of level $k$, equal to $q$ if $k = 0$ (Section~\ref{sec:running-quantities-tori}).
\item[(c)] The stop level satisfies $k^{*} = k^{*}(R) \ge 1$, and $l_{*} \in \mathrm{Cl}$, $\mathfrak{m}_{*} \in (0, \mathfrak{m}_0]$, $e_{*} \in (0, e_0]$; thus the quadruple $(k^{*}, l_{*}, \mathfrak{m}_{*}, e_{*})$ consists of data of the scale free setting of Section~\ref{sec:properties-scheme}.
\end{enumerate}
\end{proposition}

\begin{proof}
By hypothesis~\ref{def:hypothesis-flow}, $e_K = e \cdot h'^{\frac{1}{2}} \le e_1$, and $e_1 \le e_0$ by the definition of $e_1$ in Section~\ref{sec:constants-conditions-torus}. By Proposition~\ref{prop:consequences}, $v_K \ge 8$ and $(1 + \mathrm{lg}/v_K)^r < L$; so Proposition~\ref{prop:depth-stop-level} applies to every element of $\mathrm{Tor}(e)$.

\emph{(a)} The proof of Proposition~\ref{prop:depth-stop-level}(g) shows that $R_{\mathrm{top}}$ belongs to $\mathrm{Tor}(e)$, and $R \le R_{\mathrm{top}}$. By Proposition~\ref{prop:depth-stop-level}(d),
\[
N_{\mathrm{min}}(R_{\mathrm{top}},e) \le N_{\mathrm{min}}(R,e) \le N .
\]
The other requirements of hypothesis~\ref{def:hypothesis-flow} do not refer to the torus, except the one on the primitives, and Definition~\ref{def:periodisation-projection} fixes the primitives on every element of $\mathrm{Tor}(e)$.

\emph{(b)} Let $0 \le k < k^{*}(R)$, so that $1 \le k + 1 \le k^{*}(R)$. By Proposition~\ref{prop:consequences}, condition~\eqref{eq:renormalisation-flow-3} holds on $R$, and, by (a), on $R_{\mathrm{top}}$. By Proposition~\ref{prop:depth-stop-level}(e), (f), applied on $R$, the sides $l_k$ and $l_{k+1} = l_k/L$ lie in $\mathrm{Cl}$, and $l_k \ge M \cdot r_k$. Since $k + 1 \le k^{*}$, Proposition~\ref{prop:depth-stop-level}(e) gives moreover
\[
\bar{m} \cdot a_{k+1} = L \cdot \bar{m} \cdot a_k, \qquad 0 < \bar{m} \cdot a_{k+1} \le \mathfrak{m}_{*} \le \mathfrak{m}_K \le L^{-m_0} \le \mathfrak{m}_0, \qquad e_{k+1} \le e_{*} \le e_K,
\]
and $\bar{m} \cdot a_{k+1} = \mathfrak{m}_{k+1}$, because $k + 1 \ge 1$. Together with the inequalities $e_K \le e_1 \le e_0$ obtained at the beginning of the proof, this yields
\[
0 < \bar{m} \cdot a_k = \frac{\mathfrak{m}_{k+1}}{L} \le \frac{\mathfrak{m}_K}{L} \le \frac{L^{-m_0}}{L} \le \frac{\mathfrak{m}_0}{L}, \qquad e_{k+1} \le e_{*} \le e_K \le e_1 \le e_0 .
\]
Since $a_{k+1} = L \cdot a_k$ (Section~\ref{sec:running-quantities-tori}), the side, the mass and the charge of the level $k + 1$ arise from those of the level $k$ as hypothesis~\ref{def:hypothesis-one-step} prescribes:
\[
l_{k+1} = \frac{l_k}{L}, \qquad \mathfrak{m}_{k+1} = \bar{m} \cdot a_{k+1} = L \cdot \bar{m} \cdot a_k, \qquad e_{k+1} = L^{\frac{1}{2}} \cdot e_k.
\]
These relations, together with the facts established before them, are the requirements of hypothesis~\ref{def:hypothesis-one-step} on the data, and the requirements on the constants are contained in hypothesis~\ref{def:hypothesis-flow}. It remains to treat the primitives. By Proposition~\ref{prop:depth-stop-level}(d), $k < k^{*}(R) \le k_{\mathrm{top}}$, and by Proposition~\ref{prop:depth-stop-level}(f), applied on $R_{\mathrm{top}}$, the side $l_k(R_{\mathrm{top}})$ lies in $\mathrm{Cl}$. By Definition~\ref{def:periodisation-projection}, the primitives on $T_k(R_{\mathrm{top}})$ form the system of Lemma~\ref{lem:assembled-systems-primitives}; this lemma applies with the mass $\bar{m} \cdot a_k$ and the charge $e_k$, because $\bar{m} \cdot a_k \in (0, \mathfrak{m}_0/L] \subset (0, \mathfrak{m}_0]$ and $0 < e_k \le e_{k+1} \le e_0$, and the system has the properties~\ref{it:renormalisation-group-1}--\ref{it:renormalisation-group-17} for $\mathrm{par}^0$ and for $\mathrm{par}$. This proves (b) for $R = R_{\mathrm{top}}$. Let now $R < R_{\mathrm{top}}$. Both $R$ and $R_{\mathrm{top}}$ belong to $\mathrm{Tor}(e)$ (see~(a)), so the quotient $J := R_{\mathrm{top}}/R$ is an integral power of $L$; as $L$ is odd, so is $J$. Moreover $l_k(R_{\mathrm{top}}) = J \cdot l_k(R)$. By Definition~\ref{def:periodisation-projection}, the primitives on $T_k(R)$ are the projections (Definition~\ref{def:projection}) of the system on $T_k(R_{\mathrm{top}})$. By Theorem~\ref{thm:properties-projection}, applied with $l_1 = l_k(R)$, $l_2 = l_k(R_{\mathrm{top}})$, the mass $\bar{m} \cdot a_k$ and the charge $e_k$, the projected system has the properties~\ref{it:renormalisation-group-1}--\ref{it:renormalisation-group-17} for $\mathrm{par}$.

\emph{(c)} By Proposition~\ref{prop:depth-stop-level}(c), $k^{*} \ge 1$. Since condition~\eqref{eq:renormalisation-flow-3} holds on $R$ (Proposition~\ref{prop:consequences}), Proposition~\ref{prop:depth-stop-level}(f) gives $l_{*} \in \mathrm{Cl}$. By Proposition~\ref{prop:depth-stop-level}(e), $0 < \mathfrak{m}_{*} \le \mathfrak{m}_K \le L^{-m_0} \le \mathfrak{m}_0$ and $0 < e_{*} \le e_K \le e_1 \le e_0$.
\end{proof}

By Propositions~\ref{prop:depth-stop-level}, \ref{prop:consequences} and~\ref{prop:data-along-flow}, Theorem~\ref{thm:scales-stop-levels} follows.

Let hypothesis~\ref{def:hypothesis-flow} hold. By Proposition~\ref{prop:data-along-flow}(b), hypothesis~\ref{def:hypothesis-one-step} holds at every level $k < k^{*}(R)$ on $R$; by Proposition~\ref{prop:data-along-flow}(a) and~(b), applied with $R_{\mathrm{top}}$ in place of $R$, and by the inequality $k^{*}(R) \le k_{\mathrm{top}}$ of Proposition~\ref{prop:depth-stop-level}(d), it holds at these levels on $R_{\mathrm{top}}$ as well. Hence the assertions of Theorem~\ref{thm:step-small-fields} hold at every level $k < k^{*}(R)$, on $R$ and on $R_{\mathrm{top}}$. At the level $k = k^{*}(R)$, which satisfies $k^{*}(R) \ge 1$, the data belong to the scale free setting of Section~\ref{sec:properties-scheme}, by Proposition~\ref{prop:data-along-flow}(c). Therefore Propositions~\ref{prop:bounds-unit-lattice}, \ref{prop:block-covariance-real}, \ref{prop:background-field}, \ref{prop:kernel-bounds}, \ref{prop:fields-domains}(i) and~\ref{prop:gauge-fixing-controlled}, which hold in that setting (Theorems~\ref{thm:photon-operators-weakened} and~\ref{thm:renormalisation-group-1}, and condition~\eqref{eq:renormalisation-group-A10} for Proposition~\ref{prop:fields-domains}), hold at $k = k^{*}(R)$. Propositions~\ref{prop:decay-fine-propagator}, \ref{prop:kernel-bounds-2} and~\ref{prop:bounds}, from which Proposition~\ref{prop:bounds-unit-lattice} is derived, hold at this level by Theorem~\ref{thm:fermion-operators}(i), whose hypotheses are met: $l_{*} \in \mathrm{Cl}$, $\mathfrak{m}_{*} \in (0, \mathfrak{m}_0]$, and $e_{*} \le e_0 \le e^F_0$ by condition~\eqref{eq:renormalisation-group-A10}. Lemmas~\ref{lem:inclusions-bounds} and~\ref{lem:invariances}, the latter for the domains $c \cdot \tilde{R}_k$ with $0 < c \le 1$, $\tilde{R}_k^{-}$ and $\tilde{G}_k$, hold at $k = k^{*}(R)$ too. Indeed, $l_{*} \in \mathrm{Cl}$ and $e_{*} \le e_0$ by Proposition~\ref{prop:data-along-flow}(c), and Lemma~\ref{lem:inclusions-domains} and the assertions of Lemma~\ref{lem:symmetries-domains} on these domains hold at every level $k \ge 0$, for every side in $\mathrm{Cl}$ and every charge in $(0, e_0]$ (Section~\ref{sec:elementary-properties-field}). The same follows from Section~\ref{sec:hypotheses-conditions}: by Proposition~\ref{prop:data-along-flow}(b), hypothesis~\ref{def:hypothesis-one-step} holds at the level $k^{*}(R) - 1 \ge 0$, and under this hypothesis Lemma~\ref{lem:inclusions-domains} and the assertions of Lemma~\ref{lem:symmetries-domains} on these domains hold at the next level as well.

\subsection{Comparison of two tori on flat polymers}\label{sec:comparison-two-tori}

\subsubsection{Admissible sets and local families}\label{sec:admissible-sets-local}

In this subsection hypothesis~\ref{def:hypothesis-flow} is assumed, and $R < R_{\mathrm{top}}$. Put $k_m := k^{*}(R)$; by Proposition~\ref{prop:depth-stop-level}(d), $k_m \le k_{\mathrm{top}}$. Let $k \le k_m$ be a fixed level, and define
\begin{gather*}
T_1 := T_k(R), \qquad T_2 := T_k(R_{\mathrm{top}}), \\
l_1 := l_k(R), \qquad l_2 := l_k(R_{\mathrm{top}}) = J\cdot l_1,
\end{gather*}
where $J := R_{\mathrm{top}}/R$ is a power of $L$. These data satisfy the hypotheses of Section~\ref{sec:two-tori}, with the mass $\mathfrak{m} = \bar{m} \cdot a_k$ and the charge $e_k$ of level $k$. Indeed, by Proposition~\ref{prop:data-along-flow}(a), hypothesis~\ref{def:hypothesis-flow} holds also with $R_{\mathrm{top}}$ in place of $R$, with the same $e$ and $N$; in particular $R$ and $R_{\mathrm{top}}$ belong to $\mathrm{Tor}(e)$, $N \ge N_{\mathrm{min}}(R,e)$ and $N \ge N_{\mathrm{min}}(R_{\mathrm{top}},e)$. By Proposition~\ref{prop:consequences}, the conditions on $v_K$ in Proposition~\ref{prop:depth-stop-level} hold, and condition~\eqref{eq:renormalisation-flow-3} holds on $R$ and on $R_{\mathrm{top}}$. So Proposition~\ref{prop:depth-stop-level}(f), applied on $R$ with $k \le k_m$ and on $R_{\mathrm{top}}$ with $k \le k_m \le k_{\mathrm{top}}$, gives $l_1 \in \mathrm{Cl}$ and $l_2 \in \mathrm{Cl}$. As $R < R_{\mathrm{top}}$, the exponent of the power $J$ is a positive integer; and $J$ is odd, because $L$ is odd. Finally, Proposition~\ref{prop:depth-stop-level}(e), applied on $R$ with $k \le k_m$, the inequality $e_K \le e_1$ of hypothesis~\ref{def:hypothesis-flow} and the inequality $e_1 \le e_0$, which holds by the definition of $e_1$ in Section~\ref{sec:constants-conditions-torus}, give
\[
0 < \bar{m} \cdot a_k \le \mathfrak{m}_K \le \mathfrak{m}_0, \qquad 0 < e_k \le e_K \le e_1 \le e_0,
\]
and neither of the two numbers $\bar{m} \cdot a_k$, $e_k$ depends on $R$. All symbols of Section~\ref{sec:two-tori} keep their meaning: $\mathrm{pr}_i \colon \mathbb{L}_k \to T_i$ denotes the projection and $\mathrm{pj} \colon T_2 \to T_1$ the covering map, with $\mathrm{pj} \circ \mathrm{pr}_2 = \mathrm{pr}_1$, and $\mathrm{pj}^{*}$ embeds the objects on $T_1$ into those on $T_2$. The sup norm of $\mathbb{R}^3$ is denoted by $|\cdot|$, and the sup distance on $T_i$ by $d_i$. The definition of $k^{*}(R)$ (Definition~\ref{def:tori-stop-level}), Proposition~\ref{prop:depth-stop-level}(e), and Proposition~\ref{prop:consequences}, which gives~\eqref{eq:renormalisation-flow-3}, yield
\[
l_1 \ge l_{k_m}(R) = M\cdot r_{*} > (31L + 16)\cdot M .
\]
None of the numbers $e_k$, $\mathfrak{m}_k$, $b_k$, $\eta_k$, $p_{0,k}$, the contours of the step map (Definition~\ref{def:renormalisation-step}) and the windows $w^{\mathrm{m}}_k$, $w^{\mathrm{E}}_k$ depends on $R$; hence $T_1$ and $T_2$ share them. When $k < k_m$, Lemma~\ref{lem:assembled-systems-primitives} provides a system of primitives on $T_2$, denoted by $\mathrm{Op}^{(2)}$, and $\mathrm{Op}^{(1)}$ denotes the system on $T_1$ obtained from $\mathrm{Op}^{(2)}$ by projection (Definition~\ref{def:periodisation-projection}; see also Definition~\ref{def:projection}). The upper index $(i)$ marks objects formed on $T_i$. A union of $M$-cubes is called connected if it is connected through faces.

\begin{definition}[Admissibility, flatness and agreement on two tori]\label{def:admissibility-flatness-agreement}
\begin{enumerate}
\item[(a)] Let $\tilde{X}$ be a non-empty finite union of $M$-cubes of $\mathbb{L}_k$. Let $\mathrm{box}(\tilde{X})$ be the smallest product of three intervals that contains $\tilde{X}$ and is itself a union of $M$-cubes, let $\operatorname{diam}(\tilde{X})$ be the largest side of $\mathrm{box}(\tilde{X})$, and let
\[
N(\tilde{X}) := \mathrm{box}(\tilde{X})^{\sim 6 d_U}
\]
be the enlargement of $\mathrm{box}(\tilde{X})$ by $6 d_U$ layers of $M$-cubes. The set $\tilde{X}$ is called \emph{admissible} (at level $k$) when
\[
\operatorname{diam}(\tilde{X}) \le l_1 - 14 d_U\cdot M .
\]
If $\tilde{X}$ is admissible, then $N(\tilde{X})$ is a box all of whose sides are at most $l_1 - 2 d_U\cdot M$, the enlargement adding $12 d_U\cdot M$ to each side; moreover, for $i = 1, 2$, the following two properties hold. First,
\begin{equation}\label{eq:renormalisation-flow-a1}\tag{a1}
N(\tilde{X}) \cap \left( N(\tilde{X}) + g \right) = \emptyset \qquad \text{for every } g \in l_1\mathbb{Z}^3 \setminus \{0\} ,
\end{equation}
because such a $g$ has a coordinate of modulus at least $l_1$, which exceeds every side of $N(\tilde{X})$. Since $l_i\mathbb{Z}^3 \subset l_1\mathbb{Z}^3$, it follows that $\mathrm{pr}_i$ is injective on $N(\tilde{X})$. Second, for any two sites $u$, $u'$ of $N(\tilde{X})$,
\begin{equation}\label{eq:renormalisation-flow-a2}\tag{a2}
d_i(\mathrm{pr}_i u, \mathrm{pr}_i u') < 2 d_U\cdot M \quad \Longrightarrow \quad |u - u'| = d_i(\mathrm{pr}_i u, \mathrm{pr}_i u') .
\end{equation}
Indeed, $d_i(\mathrm{pr}_i u, \mathrm{pr}_i u') = |u - u' - l_i\cdot w|$ for some $w \in \mathbb{Z}^3$. As $|u - u'| \le l_1 - 2 d_U\cdot M \le l_i - 2 d_U\cdot M$, every $w \ne 0$ satisfies
\[
|u - u' - l_i\cdot w| \ge l_i - |u - u'| \ge 2 d_U\cdot M ,
\]
and since $d_i(\mathrm{pr}_i u, \mathrm{pr}_i u') < 2 d_U\cdot M$, one has $w = 0$, that is, $d_i(\mathrm{pr}_i u, \mathrm{pr}_i u') = |u - u'|$. Let $p \ge 2$, and consider a relation among $p$ objects of $\mathbb{L}_k$, each of which is a cell, an $M$-cube, an $LM$-cube, a parameter cube, an $L$-block, a block $B_k(y)$ or a basis element $\Upsilon$; the points of $\Upsilon$ are those of $\operatorname{supp}\Upsilon$. The relation is said to have \emph{range less than} $2 d_U\cdot M = (4L+2)\cdot M$ if, for every $g \in l_1\mathbb{Z}^3$, it holds for $o_1, \dots, o_p$ exactly when it holds for the translates $o_1 + g, \dots, o_p + g$, and if, whenever it holds for $o_1, \dots, o_p$, any two points of $o_1 \cup \dots \cup o_p$ have sup distance less than $2 d_U\cdot M$. Objects of $T_i$ stand in such a relation if they are the images under $\mathrm{pr}_i$ of objects of $\mathbb{L}_k$ that stand in it. Objects $o_1, \dots, o_p$ of $\mathbb{L}_k$ that are contained in $N(\tilde{X})$ then stand in the relation if and only if their images $\mathrm{pr}_i(o_1), \dots, \mathrm{pr}_i(o_p)$ do. Indeed, the objects of $\mathbb{L}_k$ of the kind of $o_j$ whose image under $\mathrm{pr}_i$ is $\mathrm{pr}_i(o_j)$ are the translates of $o_j$ by the elements of $l_i\mathbb{Z}^3$, since $T_i = \mathbb{L}_k/l_i\mathbb{Z}^3$, the objects of each kind are invariant under the translations by $l_1\mathbb{Z}^3 \supset l_i\mathbb{Z}^3$ (Section~\ref{sec:two-tori}), and each of them has sup diameter at most $2L\cdot M < l_i/2$. Hence, if the images stand in the relation, there are $g_1, \dots, g_p \in l_i\mathbb{Z}^3$ such that $o_1 + g_1, \dots, o_p + g_p$ stand in it, and, after the common translation by $-g_1 \in l_1\mathbb{Z}^3$, so do $o_1, o_2 + g_2 - g_1, \dots, o_p + g_p - g_1$. Let $2 \le j \le p$, let $u$ be a point of $o_1$ and $u'$ a point of $o_j$, and write $g_j - g_1 = l_i\cdot w$ with $w \in \mathbb{Z}^3$. Then $|u - u' - l_i\cdot w| < 2 d_U\cdot M$, while $|u - u'| \le l_1 - 2 d_U\cdot M$, because $u$ and $u'$ lie in the box $N(\tilde{X})$; the argument given for~\eqref{eq:renormalisation-flow-a2} yields $w = 0$, that is, $g_j = g_1$. Therefore $o_1, \dots, o_p$ stand in the relation. Conversely, if $o_1, \dots, o_p$ stand in the relation, then so do their images, by the definition of the relation on $T_i$. Relations of range less than $2 d_U\cdot M$ are, for instance: two $M$-cubes have a common face; two parameter cubes are linked (Definition~\ref{def:parameter-cubes-polydisks}); a cell lies in a parameter cube, an $M$-cube, an $LM$-cube, an $L$-block or a block $B_k(y)$; a cell lies in the support of a basis element $\Upsilon$ (Section~\ref{sec:averaging-fluctuations-gauge}); two sites are nearest neighbours; two sites $x$, $x'$ form a pair $(x,x')$ of $\mathrm{Pz}_k$, and a bond lies on one of the six paths $\Gamma^{\pi}(x,x')$ of such a pair (Definition~\ref{def:pairs-sites-holder}).
\item[(b)] A polymer of $T_i$ is called \emph{flat} if it is the image under $\mathrm{pr}_i$ of a connected admissible set; each such set is called a \emph{lift} of the polymer.
\item[(c)] On $T_i$ the generators $\psi(\xi)$, $\chi(\boldsymbol{\zeta})$, $\Psi(u)$, $W(u)$ of Section~\ref{sec:grassmann-algebras-norms} are indexed by points of $\mathbb{L}_k$, resp.\ of $\mathbb{Z}^3$, with the antiperiodicity under the translations by $l_i\mathbb{Z}^3$ described there (see also Section~\ref{sec:notation}). For a site $x$ let $\psi(x,\cdot)$ denote the generators $\psi(\xi)$ with $\xi = (x,\boldsymbol{\beta},\boldsymbol{\omega})$, $\boldsymbol{\beta}$ a spinor index and $\boldsymbol{\omega} \in \{0,1\}$; the notations $\chi(x,x',\cdot)$, $\Psi(y,\cdot)$, $W(y,\cdot)$ are formed in the same way. If $\tilde{X}$ is admissible, $\mathrm{Gr}_i(\tilde{X})$ denotes the smallest subalgebra of the Grassmann algebra over $\mathbb{C}$ in all generators of $T_i$ that contains
\begin{itemize}
\item the generators $\psi(x,\cdot)$, $x$ a site of $N(\tilde{X})$;
\item the generators $\chi(x,x',\cdot)$, $x$ and $x'$ sites of $N(\tilde{X})$ with $0 < |x - x'| \le M$;
\item the generators $\Psi(y,\cdot)$ and $W(y,\cdot)$, $y$ a unit site of $N(\tilde{X})$.
\end{itemize}
By~\eqref{eq:renormalisation-flow-a1}, distinct indices in $N(\tilde{X})$ project to distinct indices of $T_i$; so the listed generators are independent, and each of them is labelled by exactly one index in $N(\tilde{X})$. Let
\[
\mathfrak{i} \colon \mathrm{Gr}_1(\tilde{X}) \to \mathrm{Gr}_2(\tilde{X})
\]
denote the algebra isomorphism that sends every listed generator of $T_1$ to the listed generator of $T_2$ labelled by the same index in $N(\tilde{X})$. Let $\mathrm{St}$ be a set of cells of $\mathbb{L}_k$, let $f_1$ be a connection, a family $\mathsf{Z}$ (an element of the fluctuation space $V_k$ of Section~\ref{sec:averaging-fluctuations-gauge}) or a parameter family on $T_1$, and let $f_2$ be an object of the same type on $T_2$. One says that $f_1$ and $f_2$ \emph{have equal lifts on} $\mathrm{St}$ if
\[
f_1(\mathrm{pr}_1(c)) = f_2(\mathrm{pr}_2(c)) \qquad \text{for every } c \in \mathrm{St} .
\]
If $\tilde{X}$ is contained in an admissible set $\tilde{X}'$, then $\tilde{X}$ is admissible, $N(\tilde{X}) \subset N(\tilde{X}')$, $\mathrm{Gr}_i(\tilde{X}) \subset \mathrm{Gr}_i(\tilde{X}')$, and the isomorphism $\mathfrak{i}$ that belongs to $\tilde{X}$ is the restriction of the one that belongs to $\tilde{X}'$.
\item[(d)] Let $\mathrm{Dm}$ denote either $\tilde{R}_k \times \mathfrak{X}_k$ or $\tilde{R}_k^{-} \times \mathfrak{X}_k$ (Definition~\ref{def:domains}), and let $\mathrm{Dm}(T_i)$ be the corresponding set of pairs of bond functions on $T_i$. Consider a collection
\[
F = \left\{ F(X;\mathscr{A},\mathscr{C}) : X \in D_k(T_i),\ (\mathscr{A},\mathscr{C}) \in \mathrm{Dm}(T_i) \right\}
\]
whose members are even elements of $\mathrm{Gr}[\psi,\chi]$. It is called a \emph{local family} on $T_i$ if it has the locality property~\eqref{it:renormalisation-group-1-k2}, that is, if for all $X$ and $(\mathscr{A},\mathscr{C})$ no generator $\psi(\xi)$ with $\xi$ outside $X$ and no generator $\chi(\boldsymbol{\zeta})$ with $\boldsymbol{\zeta}$ not admissible for $X$ occurs in $F(X;\mathscr{A},\mathscr{C})$, and if $F(X;\mathscr{A},\mathscr{C}) = F(X;\mathscr{A}',\mathscr{C}')$ whenever $(\mathscr{A},\mathscr{C})$ and $(\mathscr{A}',\mathscr{C}')$ in $\mathrm{Dm}(T_i)$ take the same values on all bonds that meet $X$. Let $X = \mathrm{pr}_i(\tilde{X})$ be flat. By~\eqref{eq:renormalisation-group-adm}, every site of a pair admissible for $X$ lies within distance $\eta_k$ of $X$; such a pair therefore has a lift with both sites in $\tilde{X}^{\sim 1}$, which is a subset of $N(\tilde{X})$. Consequently $F(X;\mathscr{A},\mathscr{C})$ belongs to $\mathrm{Gr}_i(\tilde{X})$.
\item[(e)] Let $F^{(1)}$ be a local family on $T_1$ and $F^{(2)}$ one on $T_2$. One writes $F^{(1)} \sim F^{(2)}$, and says that $F^{(1)}$ and $F^{(2)}$ \emph{agree on flat polymers}, if
\[
\mathfrak{i} F^{(1)}(\mathrm{pr}_1\tilde{X};\mathscr{A}_1,\mathscr{C}_1) = F^{(2)}(\mathrm{pr}_2\tilde{X};\mathscr{A}_2,\mathscr{C}_2)
\]
for each connected admissible $\tilde{X}$ and for all pairs $(\mathscr{A}_1,\mathscr{C}_1) \in \mathrm{Dm}(T_1)$ and $(\mathscr{A}_2,\mathscr{C}_2) \in \mathrm{Dm}(T_2)$ with equal lifts on $N(\tilde{X})$. For $i = 1, 2$ let $\mathfrak{P}^{(i)} = \{ \mathfrak{P}^{(i)}(X,\cdot) : X \in D_k(T_i) \}$ be a family of complex functions on $\mathfrak{X}_k(T_i)$. The families $\mathfrak{P}^{(1)}$ and $\mathfrak{P}^{(2)}$ \emph{agree}, written $\mathfrak{P}^{(1)} \sim \mathfrak{P}^{(2)}$, if
\[
\mathfrak{P}^{(2)}(\mathrm{pr}_2\tilde{X}, \mathrm{pj}^{*}\mathscr{C}) = \mathfrak{P}^{(1)}(\mathrm{pr}_1\tilde{X}, \mathscr{C})
\]
for every connected admissible $\tilde{X}$ and every $\mathscr{C} \in \mathfrak{X}_k(T_1)$.
\end{enumerate}
\end{definition}

For families without the second connection, such as the restrictions $\mathscr{E}|_0 := \mathscr{E}(\,\cdot\,;\,\cdot\,,0,\,\cdot\,)$ of Section~\ref{sec:spaces-activities}, agreement is defined by the condition of Definition~\ref{def:admissibility-flatness-agreement}(e) with $\mathscr{C}_1 = 0$ and $\mathscr{C}_2 = 0$. With this definition, $\mathscr{E}^{(1)} \sim \mathscr{E}^{(2)}$ implies $\mathscr{E}^{(1)}|_0 \sim \mathscr{E}^{(2)}|_0$.

\begin{lemma}[Agreement criterion; small polymers; objects inside an admissible set]\label{lem:agreement-criterion-small}
Let $i \in \{1, 2\}$.
\begin{enumerate}
\item[(a)] For local families $F^{(1)}$ on $T_1$ and $F^{(2)}$ on $T_2$ the relation $F^{(1)} \sim F^{(2)}$ holds if and only if
\[
\mathfrak{i} F^{(1)}(\mathrm{pr}_1\tilde{X};\mathscr{A},\mathscr{C}) = F^{(2)}(\mathrm{pr}_2\tilde{X};\mathrm{pj}^{*}\mathscr{A},\mathrm{pj}^{*}\mathscr{C})
\]
for every connected admissible $\tilde{X}$ and all $(\mathscr{A},\mathscr{C}) \in \mathrm{Dm}(T_1)$. If they agree, then
\[
\mathfrak{i} F^{(1)}(\mathrm{pr}_1\tilde{X};\mathscr{A}_1,\mathscr{C}_1) = F^{(2)}(\mathrm{pr}_2\tilde{X};\mathscr{A}_2,\mathscr{C}_2)
\]
for every connected admissible $\tilde{X}$ and all $(\mathscr{A}_1,\mathscr{C}_1) \in \mathrm{Dm}(T_1)$, $(\mathscr{A}_2,\mathscr{C}_2) \in \mathrm{Dm}(T_2)$ that have equal lifts on the bonds that meet $\tilde{X}$.
\item[(b)] Each small polymer of $T_i$, that is, each polymer $X$ of $T_i$ with $d_M(X) \le L$ (Section~\ref{sec:lattice-geometry}), is flat. If $\tilde{\square}$ is an $M$-cube of $\mathbb{L}_k$, then $\tilde{X} \mapsto \mathrm{pr}_i(\tilde{X})$ maps the connected sets $\tilde{X}$ with $\tilde{\square} \subset \tilde{X}$ and $d_M(\tilde{X}) \le L$ bijectively onto the small polymers of $T_i$ that contain $\mathrm{pr}_i(\tilde{\square})$. The following three conditions on a connected admissible $\tilde{X}$ are equivalent: $\mathrm{pr}_1(\tilde{X})$ is small; $d_M(\tilde{X}) \le L$; $\mathrm{pr}_2(\tilde{X})$ is small.
\item[(c)] Let $\tilde{V}$ be an admissible set. Then projection by $\mathrm{pr}_i$ gives a bijection
\begin{itemize}
\item from the connected unions $\tilde{X} \subset \tilde{V}$ of $M$-cubes to the polymers $X \subset \mathrm{pr}_i(\tilde{V})$ of $T_i$;
\item for each kind of parameter cube, from the cubes $\tilde{c} \subset \tilde{V}$ of this kind to the cubes $c \subset \mathrm{pr}_i(\tilde{V})$ of this kind;
\item from the unit sites in $\tilde{V}$ to the unit sites in $\mathrm{pr}_i(\tilde{V})$;
\item from the basis elements of $\mathbb{L}_k$ contained in $\tilde{V}$ to the basis elements of $T_i$ contained in $\mathrm{pr}_i(\tilde{V})$.
\end{itemize}
These bijections are compatible with inclusions and unions, with the relation ``meets'', with common faces and links, with the formation of the closure $\bar{X}^L$ (Section~\ref{sec:lattice-geometry}), of $P_0(X)$ and of $Y'(X,P)$ (Section~\ref{sec:localisation-map}), and with the incompatibility of polymers (Proposition~\ref{prop:polymers}); on small polymers they preserve the number $d_M$.
\end{enumerate}
\end{lemma}

\begin{proof}
(a) Let $(\mathscr{A},\mathscr{C}) \in \mathrm{Dm}(T_1)$. By Lemma~\ref{lem:periodic-extension} the pair $\mathrm{pj}^{*}(\mathscr{A},\mathscr{C}) := (\mathrm{pj}^{*}\mathscr{A}, \mathrm{pj}^{*}\mathscr{C})$ belongs to $\mathrm{Dm}(T_2)$, and, because $\mathrm{pj} \circ \mathrm{pr}_2 = \mathrm{pr}_1$, the two pairs $(\mathscr{A},\mathscr{C})$ and $\mathrm{pj}^{*}(\mathscr{A},\mathscr{C})$ have equal lifts on every set of cells of $\mathbb{L}_k$. Hence $F^{(1)} \sim F^{(2)}$ implies the identity stated in (a).

Conversely, assume that this identity holds. Let $\tilde{X}$ be connected and admissible, and let $(\mathscr{A}_1,\mathscr{C}_1) \in \mathrm{Dm}(T_1)$ and $(\mathscr{A}_2,\mathscr{C}_2) \in \mathrm{Dm}(T_2)$ have equal lifts on the bonds that meet $\tilde{X}$. Every bond of $T_2$ meeting $\mathrm{pr}_2(\tilde{X})$ is the image $\mathrm{pr}_2(b)$ of some bond $b$ of $\mathbb{L}_k$ meeting $\tilde{X}$; the value of $\mathrm{pj}^{*}\mathscr{A}_1$ at $\mathrm{pr}_2(b)$ is the value of $\mathscr{A}_1$ at $\mathrm{pr}_1(b)$, which equals the value of $\mathscr{A}_2$ at $\mathrm{pr}_2(b)$, and the same holds for $\mathscr{C}_1$ and $\mathscr{C}_2$. Consequently $\mathrm{pj}^{*}(\mathscr{A}_1,\mathscr{C}_1)$ takes the same values as $(\mathscr{A}_2,\mathscr{C}_2)$ at every bond of $T_2$ meeting $\mathrm{pr}_2(\tilde{X})$. The locality property~\eqref{it:renormalisation-group-1-k2} of $F^{(2)}$, followed by the assumed identity, gives
\[
F^{(2)}(\mathrm{pr}_2\tilde{X};\mathscr{A}_2,\mathscr{C}_2) = F^{(2)}(\mathrm{pr}_2\tilde{X};\mathrm{pj}^{*}\mathscr{A}_1,\mathrm{pj}^{*}\mathscr{C}_1) = \mathfrak{i} F^{(1)}(\mathrm{pr}_1\tilde{X};\mathscr{A}_1,\mathscr{C}_1) .
\]
This is the second assertion of (a). Every bond that meets $\tilde{X}$ lies in $\tilde{X}^{\sim 1} \subset N(\tilde{X})$; so pairs with equal lifts on $N(\tilde{X})$ have equal lifts on the bonds that meet $\tilde{X}$, and therefore $F^{(1)} \sim F^{(2)}$.

(b) From $l_1 > (31L+16)\cdot M$ one obtains $2M\cdot(L+2) < l_1 \le l_i$. Let $X$ be a small polymer of $T_i$. Then $2M\cdot\left(d_M(X)+2\right) < l_i$, and Proposition~\ref{prop:polymers}(t10) yields a connected finite union $\tilde{X}$ of $M$-cubes on which $\mathrm{pr}_i$ is injective, with $\mathrm{pr}_i(\tilde{X}) = X$ and $d_M(\tilde{X}) = d_M(X) \le L$; the same clause contains the bijection asserted in (b). Since $\tilde{X}$ is a finite union of $M$-cubes of $\mathbb{R}^3$ that is connected through faces, Proposition~\ref{prop:polymers}(t7), applied to the points of the closed cubes of $\tilde{X}$, shows that any two such points have sup distance at most $M\cdot\left(d_M(\tilde{X})+2\right) \le M\cdot(L+2)$. The largest side of $\mathrm{box}(\tilde{X})$ is the largest difference between corresponding coordinates of two points of the closed cubes of $\tilde{X}$; hence
\[
\operatorname{diam}(\tilde{X}) \le M\cdot(L+2) \le l_1 - 14 d_U\cdot M ,
\]
where the second inequality holds because $14 d_U + L + 2 = 29L + 16 < 31L + 16$. Thus $\tilde{X}$ is admissible, and $X$ is flat.

Next, let $\tilde{X}$ be a connected admissible set such that $\mathrm{pr}_i(\tilde{X})$ is small. By Proposition~\ref{prop:polymers}(t10) there is a connected finite union $\tilde{X}'$ of $M$-cubes with $\mathrm{pr}_i(\tilde{X}') = \mathrm{pr}_i(\tilde{X})$ and $d_M(\tilde{X}') \le L$ on which $\mathrm{pr}_i$ is injective. By~\eqref{eq:renormalisation-flow-a1}, $\mathrm{pr}_i$ is injective on $\tilde{X}$ as well. The uniqueness assertion of Proposition~\ref{prop:polymers}(t10) shows that $\tilde{X}'$ is a translate of $\tilde{X}$; therefore $d_M(\tilde{X}) = d_M(\tilde{X}') \le L$. Conversely, if $d_M(\tilde{X}) \le L$, then $2M\cdot\left(d_M(\tilde{X})+2\right) < l_i$, and the converse part of Proposition~\ref{prop:polymers}(t10) gives $d_M(\mathrm{pr}_i\tilde{X}) = d_M(\tilde{X}) \le L$, so that $\mathrm{pr}_i(\tilde{X})$ is small. Applied with $i = 1$ and with $i = 2$, this proves the equivalence of the three conditions.

(c) Since $\tilde{V} \subset N(\tilde{V})$, the projection $\mathrm{pr}_i$ is injective on $\tilde{V}$ by~\eqref{eq:renormalisation-flow-a1}, and it maps an object of each of the four types that is contained in $\tilde{V}$ to an object of the same type contained in $\mathrm{pr}_i(\tilde{V})$. It remains to show that every object contained in $\mathrm{pr}_i(\tilde{V})$ is obtained in this way.

Given a polymer $X$ of $T_i$ with $X \subset \mathrm{pr}_i(\tilde{V})$, define $\tilde{X}$ as the union of those $M$-cubes of $\tilde{V}$ that are mapped into $X$. Every $M$-cube of $X$ lies in $\mathrm{pr}_i(\tilde{V})$ and is therefore the image of an $M$-cube of $\tilde{V}$; so $\mathrm{pr}_i(\tilde{X}) = X$. By~\eqref{eq:renormalisation-flow-a1} and~\eqref{eq:renormalisation-flow-a2}, two $M$-cubes of $\tilde{X}$ whose images have a common face have a common face themselves; as $X$ is connected, so is $\tilde{X}$.

Consider a parameter cube $c$ contained in $\mathrm{pr}_i(\tilde{V})$. The centre of $c$ is one of its lattice sites and therefore has a lift $\tilde{z} \in \tilde{V}$; denote by $\tilde{c}$ the cube of the kind of $c$ centred at $\tilde{z}$, so that $\mathrm{pr}_i$ maps $\tilde{c}$ onto $c$. A parameter cube has sides at most $2L\cdot M$ (Definition~\ref{def:parameter-cubes-polydisks}); hence every site of $\tilde{c}$ belongs to $\tilde{V}^{\sim (L+1)}$, which is contained in $N(\tilde{V})$ because $L + 1 \le 6 d_U$. If $u$ is a site of $\tilde{c}$, then $\mathrm{pr}_i(u) \in c \subset \mathrm{pr}_i(\tilde{V})$, so that $\mathrm{pr}_i(u) = \mathrm{pr}_i(u')$ for some $u' \in \tilde{V}$, and $u = u'$ by~\eqref{eq:renormalisation-flow-a1}. Thus $\tilde{c} \subset \tilde{V}$.

By~\eqref{eq:renormalisation-flow-a1}, each unit site of $T_i$ in $\mathrm{pr}_i(\tilde{V})$ has exactly one preimage $u$ in $\tilde{V}$; as $l_i$ is an integer, $u$ belongs to $\mathbb{Z}^3$.

Let $\Upsilon$ be a basis element contained in $\mathrm{pr}_i(\tilde{V})$. Choose a point of its support, let $\tilde{u} \in \tilde{V}$ be the lift of this point, and denote by $\tilde{\Upsilon}$ the basis element of $\mathbb{L}_k$ above $\Upsilon$ with $\tilde{u} \in \operatorname{supp}\tilde{\Upsilon}$. By Proposition~\ref{prop:basis-map-box}, $\operatorname{supp}\Upsilon$ lies in one $L$-block or is one bond; consequently every point of $\operatorname{supp}\tilde{\Upsilon}$ has sup distance at most $L$ from $\tilde{u}$, and $\operatorname{supp}\tilde{\Upsilon} \subset N(\tilde{V})$. For each point $u$ of $\operatorname{supp}\tilde{\Upsilon}$ there is $u' \in \tilde{V}$ with $\mathrm{pr}_i(u) = \mathrm{pr}_i(u')$, and $u = u'$ by~\eqref{eq:renormalisation-flow-a1}. Therefore $\tilde{\Upsilon}$ is contained in $\tilde{V}$.

Each relation and each set named in the last sentence of (c) is formed, by inclusions, intersections and unions, from four relations between the $M$-cubes, $LM$-cubes and parameter cubes that make up the objects concerned: two $M$-cubes coincide or have a common face; an $M$-cube and a parameter cube have a common lattice site; an $M$-cube lies in an $LM$-cube; two parameter cubes of one kind are linked. Indeed, $P_0(X)$ is the set of the parameter cubes contained in $X$; $Y'(X,P)$ is the union of the $M$-cubes that have a lattice site in common with $X$ or with a cube of $P_0(X) \cup P$; $\bar{X}^L$ is the union of the $LM$-cubes that contain an $M$-cube of $X$; and two polymers are incompatible if their union is connected through faces. The four relations have range less than $2 d_U\cdot M = (4L+2)\cdot M$ in the sense of Definition~\ref{def:admissibility-flatness-agreement}(a). They are invariant under the common translations by the elements of $l_1\mathbb{Z}^3$, because the $M$-cubes, the $LM$-cubes and the parameter cubes of each kind are (Section~\ref{sec:two-tori}); and two $M$-cubes that coincide or have a common face lie in a box of side $2M$; an $M$-cube and a parameter cube with a common lattice site lie in a box of side at most $2L\cdot M + M$; an $LM$-cube has side $L\cdot M$; two linked cubes of $\pi''_k$ lie in a box of side $4L\cdot M$, and two linked cubes of $\pi_k$ or of $\pi'_k$ in a box of side $4M$. Every cube that occurs here for objects contained in $\tilde{V}$ lies in $\tilde{V}^{\sim (L+1)} \subset N(\tilde{V})$: an $M$-cube that has a lattice site in common with a subset of $\tilde{V}$ is an $M$-cube of $\tilde{V}$, since every site lies in exactly one $M$-cube (fact~\eqref{it:renormalisation-group-g0}), and an $LM$-cube that contains an $M$-cube $\tilde{\square}$ of $\tilde{V}$ lies in $\tilde{\square}^{\sim (L-1)}$. As the $M$-cubes, the $LM$-cubes and the parameter cubes of $T_i$ are the images of those of $\mathbb{L}_k$, the same argument on $T_i$ shows that every cube of $T_i$ that occurs for objects contained in $\mathrm{pr}_i(\tilde{V})$ is the image of such a cube. Hence, by Definition~\ref{def:admissibility-flatness-agreement}(a), applied with $\tilde{V}$ in place of $\tilde{X}$, cubes contained in $N(\tilde{V})$ satisfy one of the four relations exactly when their images under $\mathrm{pr}_i$ satisfy it; and inclusions, intersections and unions of subsets of $N(\tilde{V})$ are preserved, since $\mathrm{pr}_i$ is injective on $N(\tilde{V})$ by~\eqref{eq:renormalisation-flow-a1}. Finally, $d_M$ is preserved on small polymers by Proposition~\ref{prop:polymers}(t10), as in the proof of (b).
\end{proof}

\subsubsection{Projected primitives near an admissible set}\label{sec:projected-primitives-near}

\begin{lemma}[Kernels near an admissible set]\label{lem:kernels-near-admissible}
Let $k<k_m$. As in Section~\ref{sec:admissible-sets-local}, $\mathrm{Op}^{(2)}$ denotes the system of primitives of Lemma~\ref{lem:assembled-systems-primitives} on $T_2$, and $\mathrm{Op}^{(1)}$ its projection to $T_1$ in the sense of Definition~\ref{def:projection}. Let $\mathrm{Op}$ be one of the families~\eqref{it:primitives-p1}--\eqref{it:primitives-p3}, $\Gamma_k$. Let $\tilde{Y}$ be an admissible union of $M$-cubes of $\mathbb{L}_k$ (Definition~\ref{def:admissibility-flatness-agreement}(a)), and put $Y_i:=\mathrm{pr}_i(\tilde{Y})$ for $i=1,2$. Let $s^{(1)}$ be an element of $\mathrm{Pd}_k(T_1)$ such that $s^{(1)}_c=0$ for every parameter cube $c$ of $T_1$ not contained in $Y_1$. A family $s^{(2)}$ indexed by the parameter cubes of $T_2$ is defined by
\[
s^{(2)}_{\mathrm{pr}_2(\tilde{c})}:=s^{(1)}_{\mathrm{pr}_1(\tilde{c})}
\]
for each parameter cube $\tilde{c}$ of $\mathbb{L}_k$ contained in $\tilde{Y}$, and by $s^{(2)}_c:=0$ for each parameter cube $c$ of $T_2$ that is not of the form $\mathrm{pr}_2(\tilde{c})$ with such a cube $\tilde{c}$. Let $j$ be an integer with $0\le j\le 5d_U$. For $i=1,2$, every index of $T_i$ that is contained in $Y_i^{\sim j}$ has exactly one lift to $\mathbb{L}_k$ (an index of $\mathbb{L}_k$ that $\mathrm{pr}_i$ maps to it) that is contained in $\tilde{Y}^{\sim j}$. Let $x$ be a row index of $\mathrm{Op}$ on $T_1$ that is contained in $Y_1^{\sim j}$ (that is, $\operatorname{supp}x$ lies in $Y_1^{\sim j}$, and hence so does $\square(x)$), and let $\tilde{x}$ be the lift of $x$ that is contained in $\tilde{Y}^{\sim j}$. Then the following holds for every $A^{(1)}\in\tilde{G}_k(T_1)$.
\begin{enumerate}
\item[(a)] One has $\mathrm{Op}^{(1)}\left(s^{(1)},A^{(1)}\right)(x,x')=0$ for every column index $x'$ that is not contained in $Y_1^{\sim(j+c_l)}$, and the row of $\mathrm{Op}^{(1)}\left(s^{(1)},A^{(1)}\right)$ at $x$ depends on $A^{(1)}$ only through the values of $A^{(1)}$ on the bonds contained in $Y_1^{\sim(j+d_U)}$.
\item[(b)] Let $A^{(2)}\in\tilde{G}_k(T_2)$ be such that $A^{(1)}$ and $A^{(2)}$ have equal lifts on the bonds contained in $\tilde{Y}^{\sim(j+d_U)}$. Then
\[
\mathrm{Op}^{(1)}\left(s^{(1)},A^{(1)}\right)(\tilde{x},\tilde{x}')=\mathrm{Op}^{(2)}\left(s^{(2)},A^{(2)}\right)(\tilde{x},\tilde{x}')
\]
for every column index $\tilde{x}'$ that is contained in $N(\tilde{Y})$, and $\mathrm{Op}^{(2)}\left(s^{(2)},A^{(2)}\right)(\tilde{x},x'')=0$ for every column index $x''$ of $T_2$ that is not contained in $Y_2^{\sim(j+c_l)}$.
\end{enumerate}
Moreover, let $\tilde{X}$ be a union of $M$-cubes of $\mathbb{L}_k$ that is connected through faces and contained in $\tilde{Y}$, let $\tilde{x}$ and $\tilde{x}'$ now be any two indices contained in $\tilde{X}$, and let $A^{(i)}\in\tilde{G}_k(T_i)$, $i=1,2$, have equal lifts on the bonds that meet $\tilde{X}$. Then
\[
M_k^{(1)}\left(\mathrm{pr}_1(\tilde{X}),A^{(1)}\right)(\tilde{x},\tilde{x}')=M_k^{(2)}\left(\mathrm{pr}_2(\tilde{X}),A^{(2)}\right)(\tilde{x},\tilde{x}'),
\]
and the same identity holds with $S_{k,y}$ in place of $M_k$, for every $y\ge0$.
\end{lemma}

\begin{proof}
\emph{Lifts.} Let $i\in\{1,2\}$. The $M$-cubes of $T_i$ are the images of those of $\mathbb{L}_k$, and the distance $d_i$ of two points of $T_i$ is the least sup distance between a lift of the one and a lift of the other; hence $\mathrm{pr}_i\left(\tilde{Y}^{\sim j}\right)=Y_i^{\sim j}$. Let $a$ be an index of $T_i$ that is contained in $Y_i^{\sim j}$. The lifts of $a$ are the translates of one of them by the elements of $l_i\mathbb{Z}^3$, and $\mathrm{pr}_i$ maps the support of each lift onto $\operatorname{supp}a$. Choose a point $p$ of $\operatorname{supp}a$ and a point $\tilde{p}$ of $\tilde{Y}^{\sim j}$ with $\mathrm{pr}_i(\tilde{p})=p$. If $\tilde{a}_0$ is a lift of $a$ and $\tilde{q}$ is a point of its support with $\mathrm{pr}_i(\tilde{q})=p$, then $\tilde{p}-\tilde{q}\in l_i\mathbb{Z}^3$, and $\tilde{a}:=\tilde{a}_0+(\tilde{p}-\tilde{q})$ is a lift of $a$ whose support contains $\tilde{p}$. Since the $M$-cubes partition the sites (fact~\eqref{it:renormalisation-group-g0} of Section~\ref{sec:cubes-parameter-cubes}), the $M$-cube that contains $\tilde{p}$ is an $M$-cube of $\tilde{Y}^{\sim j}$, and $\tilde{a}$ meets it. By Section~\ref{sec:two-tori}, $\square(\tilde{a})$ is either this cube or the union of this cube and a second $M$-cube with a common face, whose centre has sup distance $M$ from the centre of the first; by the definition of the enlargement, both cubes belong to $\tilde{Y}^{\sim(j+1)}$. As $\tilde{Y}\subseteq\mathrm{box}(\tilde{Y})$ and $j+1\le6d_U$, it follows that $\operatorname{supp}\tilde{a}\subseteq\square(\tilde{a})\subseteq\tilde{Y}^{\sim(j+1)}\subseteq N(\tilde{Y})$. Let $u$ be a point of $\operatorname{supp}\tilde{a}$. As $\mathrm{pr}_i(u)$ lies in $\operatorname{supp}a\subseteq\mathrm{pr}_i\left(\tilde{Y}^{\sim j}\right)$, there is a point $u'$ of $\tilde{Y}^{\sim j}$ with $\mathrm{pr}_i(u')=\mathrm{pr}_i(u)$; both $u$ and $u'$ lie in $N(\tilde{Y})$, on which $\mathrm{pr}_i$ is injective by~\eqref{eq:renormalisation-flow-a1}, and so $u=u'$. Thus $\tilde{a}$ is contained in $\tilde{Y}^{\sim j}$. Every other lift of $a$ is $\tilde{a}+g$ with $0\ne g\in l_i\mathbb{Z}^3\subseteq l_1\mathbb{Z}^3$, and it is not contained in $\tilde{Y}^{\sim j}$: otherwise, for $u\in\operatorname{supp}\tilde{a}$, the point $u+g$ would lie in $N(\tilde{Y})\cap\left(N(\tilde{Y})+g\right)$, which is empty by~\eqref{eq:renormalisation-flow-a1}. So $\tilde{a}$ is the only lift of $a$ that is contained in $\tilde{Y}^{\sim j}$.

By Lemma~\ref{lem:assembled-systems-primitives} the system $\mathrm{Op}^{(2)}$ has the properties~\ref{it:renormalisation-group-8} and~\ref{it:renormalisation-group-12} on $T_2$; by Theorem~\ref{thm:properties-projection}(b) the projected system $\mathrm{Op}^{(1)}$ has them on $T_1$, and by Theorem~\ref{thm:properties-projection}(a) its kernels are given by~\eqref{eq:renormalisation-flow-4}. Two parameter cubes that are contained in $\tilde{Y}$ and have the same image under $\mathrm{pr}_2$ differ by a translation by an element of $l_2\mathbb{Z}^3\subseteq l_1\mathbb{Z}^3$ and both lie in $N(\tilde{Y})$; by~\eqref{eq:renormalisation-flow-a1} they coincide. Hence $s^{(2)}$ is well defined, and it lies in $\mathrm{Pd}_k(T_2)$, each of its values being a value of $s^{(1)}$ or zero.

Put
\[
\tilde{W}:=\tilde{Y}^{\sim(j+c_l)},\qquad W_i:=\mathrm{pr}_i(\tilde{W})=Y_i^{\sim(j+c_l)}\quad(i=1,2).
\]
Since $c_l+1=d_U$ (Definition~\ref{def:parameter-cubes-polydisks}) and $j\le5d_U$, one has $j+c_l+1\le6d_U$; therefore the set $\tilde{W}^{\sim1}=\tilde{Y}^{\sim(j+d_U)}$ is contained in $N(\tilde{Y})$ (Definition~\ref{def:admissibility-flatness-agreement}(a)). The translates of $\tilde{W}^{\sim1}$ by the elements of $l_1\mathbb{Z}^3$ are pairwise disjoint, by~\eqref{eq:renormalisation-flow-a1}. It follows that, for $g\ne g'$ in $l_1\mathbb{Z}^3$, no $M$-cube of $\tilde{W}+g$ has a face in common with an $M$-cube of $\tilde{W}+g'$, for such a cube of $\tilde{W}+g'$ would lie in $\tilde{W}^{\sim1}+g$. As $\mathrm{pr}_1=\mathrm{pj}\circ\mathrm{pr}_2$, the set $\mathrm{pj}^{-1}W_1$ is the union of the sets $\mathrm{pr}_2(\tilde{W}+g)$, where $g$ ranges over a system of representatives of $l_1\mathbb{Z}^3/l_2\mathbb{Z}^3$.

The set $\square(\tilde{x})$ is one $M$-cube or the union of two $M$-cubes with a common face, and it lies in $\tilde{Y}^{\sim j}$, because $\tilde{x}$ is contained in $\tilde{Y}^{\sim j}$. Hence $\square(\tilde{x})^{\sim c_l}$ is connected and contained in $\tilde{W}$. Let $\tilde{W}_1$ be the component of $\tilde{W}$ that contains $\square(\tilde{x})^{\sim c_l}$. Then $\tilde{W}_1$ is also a component of
\[
\bigcup_{g\in l_1\mathbb{Z}^3}\left(\tilde{W}+g\right)=\mathrm{pr}_1^{-1}(W_1)=\mathrm{pr}_2^{-1}\left(\mathrm{pj}^{-1}W_1\right),
\]
since an $M$-cube of this union that has a face in common with an $M$-cube of $\tilde{W}_1$ belongs to $\tilde{W}$, by the preceding paragraph, and hence to $\tilde{W}_1$. If two $M$-cubes of $T_i$ have a common face, then every lift of the first has a common face with a lift of the second, $l_i$ being larger than $2M$. It follows that $W_1$ has $\mathrm{pr}_1(\tilde{W}_1)$ among its components and that $\mathrm{pj}^{-1}W_1$ has $\mathrm{pr}_2(\tilde{W}_1)$ among its components; since $\mathrm{pr}_2(\tilde{W}_1)\subseteq W_2\subseteq\mathrm{pj}^{-1}W_1$, the set $\mathrm{pr}_2(\tilde{W}_1)$ is also a component of $W_2$. Moreover, for $i=1,2$ the set
\[
\square\left(\mathrm{pr}_i(\tilde{x})\right)^{\sim c_l}=\mathrm{pr}_i\left(\square(\tilde{x})^{\sim c_l}\right)
\]
lies in $\mathrm{pr}_i(\tilde{W}_1)$; for $i=1$ it is the set $\square(x)^{\sim c_l}$.

\emph{Proof of~(a).} The family $s^{(1)}$ vanishes off $Y_1$, hence off $W_1\supseteq Y_1$. Apply Lemma~\ref{lem:locality-rows} on $T_1$ to the system $\mathrm{Op}^{(1)}$, the union $W_1$ and its component $\mathrm{pr}_1(\tilde{W}_1)$, the family $s^{(1)}$ and the row index $x$. First, it gives $\mathrm{Op}^{(1)}\left(s^{(1)},A^{(1)}\right)(x,x')=0$ unless $x'$ is contained in $\mathrm{pr}_1(\tilde{W}_1)$, which is a subset of $W_1=Y_1^{\sim(j+c_l)}$. Second, let $B,B'\in\tilde{G}_k(T_1)$ coincide on the bonds that meet $\mathrm{pr}_1(\tilde{W}_1)$. The second assertion of Lemma~\ref{lem:locality-rows}, applied to the pair of connections $(B,B')$ and to the pair $(B',B')$, shows that the rows at $x$ of $\mathrm{Op}^{(1)}\left(s^{(1)},B\right)$ and of $\mathrm{Op}^{(1)}\left(s^{(1)},B'\right)$ both coincide with the row at $x$ of $\mathrm{Op}^{(1)}\left(s^{(1)}\cdot\mathbf{1}_{\mathrm{pr}_1(\tilde{W}_1)},B'\right)$. Every bond that meets $\mathrm{pr}_1(\tilde{W}_1)$ is contained in $W_1^{\sim1}=Y_1^{\sim(j+c_l+1)}=Y_1^{\sim(j+d_U)}$. So two connections that coincide on the bonds contained in $Y_1^{\sim(j+d_U)}$ give the same row at $x$, and~(a) is proved.

\emph{Proof of~(b).} The family $\mathrm{pj}^{*}s^{(1)}$ lies in $\mathrm{Pd}_k(T_2)$ (Section~\ref{sec:two-tori}), and it vanishes off $\mathrm{pj}^{-1}W_1$: if a parameter cube $c$ of $T_2$ is not contained in $\mathrm{pj}^{-1}W_1$, then $\mathrm{pj}(c)$ is not contained in $W_1$, and $\left(\mathrm{pj}^{*}s^{(1)}\right)_c=s^{(1)}_{\mathrm{pj}(c)}=0$. The family $s^{(2)}$ vanishes off $W_2$, since a cube $\mathrm{pr}_2(\tilde{c})$ with $\tilde{c}$ contained in $\tilde{Y}$ is contained in $Y_2\subseteq W_2$.

Every parameter cube $c$ of $T_2$ that is contained in $\mathrm{pr}_2(\tilde{W}_1)$ is the image $\mathrm{pr}_2(\tilde{c})$ of a parameter cube $\tilde{c}$ of the same kind that is contained in $\tilde{W}_1$. Indeed, by Section~\ref{sec:lattice-geometry} and Definition~\ref{def:parameter-cubes-polydisks}, the centres of the parameter cubes lie in $M\mathbb{Z}^3$ or in $LM\mathbb{Z}^3$, both of which are contained in $\mathbb{L}_k$, and every parameter cube lies in the closed cube of sup-radius $L\cdot M$ about its centre. So the centre of $c$ is a lattice site of $c$, and it equals $\mathrm{pr}_2(\tilde{z})$ for a site $\tilde{z}$ of $\tilde{W}_1$. The site $\tilde{z}$ belongs to the grid of the centres, this grid being invariant under $l_2\mathbb{Z}^3$; let $\tilde{c}$ be the parameter cube of the kind of $c$ with the centre $\tilde{z}$. Then $\mathrm{pr}_2$ maps the lattice sites of $\tilde{c}$ onto those of $c$. Let $u$ be a lattice site of $\tilde{c}$. Since $\mathrm{pr}_2(u)$ lies in $c\subseteq\mathrm{pr}_2(\tilde{W}_1)$, there is a site $u'$ of $\tilde{W}_1$ with $\mathrm{pr}_2(u)=\mathrm{pr}_2(u')$, that is, with $u-u'\in l_2\mathbb{Z}^3$. The sites $\tilde{z}$ and $u'$ lie in the box $N(\tilde{Y})$, whose sides are at most $l_1-2d_U\cdot M$ (Definition~\ref{def:admissibility-flatness-agreement}(a)), and $|u-\tilde{z}|\le L\cdot M$, since $u$ lies in the parameter cube $\tilde{c}$ with the centre $\tilde{z}$. As $d_U=2L+1$ (Definition~\ref{def:parameter-cubes-polydisks}), this gives
\[
|u-u'|\le|u-\tilde{z}|+|\tilde{z}-u'|\le L\cdot M+l_1-2d_U\cdot M=l_1-(3L+2)\cdot M<l_1\le l_2,
\]
and therefore $u=u'$. Thus every lattice site of $\tilde{c}$ lies in $\tilde{W}_1$; that is, $\tilde{c}$ is contained in $\tilde{W}_1$, and $c=\mathrm{pr}_2(\tilde{c})$.

Now let $\tilde{c}$ be a parameter cube that is contained in $\tilde{W}_1$. Then
\[
\left(\mathrm{pj}^{*}s^{(1)}\right)_{\mathrm{pr}_2(\tilde{c})}=s^{(1)}_{\mathrm{pj}(\mathrm{pr}_2(\tilde{c}))}=s^{(1)}_{\mathrm{pr}_1(\tilde{c})}.
\]
If $\tilde{c}$ is contained in $\tilde{Y}$, this number equals $s^{(2)}_{\mathrm{pr}_2(\tilde{c})}$ by the definition of $s^{(2)}$. If $\tilde{c}$ is not contained in $\tilde{Y}$, both numbers are zero. On the one hand, $\mathrm{pr}_2(\tilde{c})$ is then not the image of a parameter cube contained in $\tilde{Y}$: such a cube would be $\tilde{c}+g$ with $0\ne g\in l_2\mathbb{Z}^3$, and $\tilde{c}$ and $\tilde{c}+g$ would both lie in $N(\tilde{Y})$, contrary to~\eqref{eq:renormalisation-flow-a1}; hence $s^{(2)}_{\mathrm{pr}_2(\tilde{c})}=0$. On the other hand, if $s^{(1)}_{\mathrm{pr}_1(\tilde{c})}$ were not zero, then $\mathrm{pr}_1(\tilde{c})$ would be contained in $Y_1=\mathrm{pr}_1(\tilde{Y})$, because $s^{(1)}$ vanishes off $Y_1$; as $\tilde{c}\subseteq\tilde{W}_1\subseteq N(\tilde{Y})$ and $\mathrm{pr}_1$ is injective on $N(\tilde{Y})$ by~\eqref{eq:renormalisation-flow-a1}, the cube $\tilde{c}$ would be contained in $\tilde{Y}$. Together with the preceding paragraph this shows
\[
\left(\mathrm{pj}^{*}s^{(1)}\right)\cdot\mathbf{1}_{\mathrm{pr}_2(\tilde{W}_1)}=s^{(2)}\cdot\mathbf{1}_{\mathrm{pr}_2(\tilde{W}_1)}.
\]

By Lemma~\ref{lem:periodic-extension} the connection $\mathrm{pj}^{*}A^{(1)}$ lies in $\tilde{G}_k(T_2)$. It coincides with $A^{(2)}$ on the bonds that meet $\mathrm{pr}_2(\tilde{W}_1)$. Indeed, such a bond is $\mathrm{pr}_2(b)$ for a bond $b$ of $\mathbb{L}_k$ that meets $\tilde{W}_1$; this bond $b$ is contained in $\tilde{W}^{\sim1}=\tilde{Y}^{\sim(j+d_U)}$, on whose bonds $A^{(1)}$ and $A^{(2)}$ have equal lifts, so that
\[
\left(\mathrm{pj}^{*}A^{(1)}\right)\left(\mathrm{pr}_2(b)\right)=A^{(1)}\left(\mathrm{pr}_1(b)\right)=A^{(2)}\left(\mathrm{pr}_2(b)\right).
\]

Lemma~\ref{lem:locality-rows} on $T_2$ is now applied twice, both times to the system $\mathrm{Op}^{(2)}$, the component $\mathrm{pr}_2(\tilde{W}_1)$ and the row index $\tilde{x}$: first to the union $\mathrm{pj}^{-1}W_1$, the family $\mathrm{pj}^{*}s^{(1)}$ and the pair of connections $\left(\mathrm{pj}^{*}A^{(1)},A^{(2)}\right)$; then to the union $W_2$, the family $s^{(2)}$ and the pair $\left(A^{(2)},A^{(2)}\right)$. Combined with the identity of the two restricted families, this yields
\begin{align*}
\mathrm{Op}^{(2)}\left(\mathrm{pj}^{*}s^{(1)},\mathrm{pj}^{*}A^{(1)}\right)(\tilde{x},\,\cdot\,)
&=\mathrm{Op}^{(2)}\left(s^{(2)}\cdot\mathbf{1}_{\mathrm{pr}_2(\tilde{W}_1)},A^{(2)}\right)(\tilde{x},\,\cdot\,)\\
&=\mathrm{Op}^{(2)}\left(s^{(2)},A^{(2)}\right)(\tilde{x},\,\cdot\,).
\end{align*}
By the first assertion of Lemma~\ref{lem:locality-rows}, a column index at which one of these rows has a nonzero entry is contained in $\mathrm{pr}_2(\tilde{W}_1)$. As $\mathrm{pr}_2(\tilde{W}_1)\subseteq W_2=Y_2^{\sim(j+c_l)}$, the second assertion of~(b) follows.

Let $\tilde{x}'$ be a column index that is contained in $N(\tilde{Y})$. By~\eqref{eq:renormalisation-flow-4},
\begin{multline*}
\mathrm{Op}^{(1)}\left(s^{(1)},A^{(1)}\right)(\tilde{x},\tilde{x}')\\
=\sum_{w\in(\mathbb{Z}/J\mathbb{Z})^3}\boldsymbol{\chi}(w)\cdot\mathrm{Op}^{(2)}\left(\mathrm{pj}^{*}s^{(1)},\mathrm{pj}^{*}A^{(1)}\right)\left(\tilde{x},\tilde{x}'+l_1\cdot w\right).
\end{multline*}
Let $w\in\mathbb{Z}^3$ represent a class whose term does not vanish. As the row of $\mathrm{Op}^{(2)}\left(\mathrm{pj}^{*}s^{(1)},\mathrm{pj}^{*}A^{(1)}\right)$ at $\tilde{x}$ vanishes at the column indices that are not contained in $\mathrm{pr}_2(\tilde{W}_1)$, the index $\tilde{x}'+l_1\cdot w$ of $T_2$ is contained in $\mathrm{pr}_2(\tilde{W}_1)$; that is, $\tilde{x}'+l_1\cdot w+l_2\cdot v$ is contained in $\tilde{W}_1$ for some $v\in\mathbb{Z}^3$. Since $l_2=J\cdot l_1$, this index equals $\tilde{x}'+l_1\cdot(w+J\cdot v)$; it is contained in $N(\tilde{Y})+l_1\cdot(w+J\cdot v)$, because $\tilde{x}'$ is contained in $N(\tilde{Y})$, and also in $\tilde{W}_1\subseteq N(\tilde{Y})$. So $N(\tilde{Y})$ and $N(\tilde{Y})+l_1\cdot(w+J\cdot v)$ meet, and~\eqref{eq:renormalisation-flow-a1} gives $w+J\cdot v=0$. Hence only the class of $0$ contributes, with $\boldsymbol{\chi}(0)=1$, and
\begin{align*}
\mathrm{Op}^{(1)}\left(s^{(1)},A^{(1)}\right)(\tilde{x},\tilde{x}')
&=\mathrm{Op}^{(2)}\left(\mathrm{pj}^{*}s^{(1)},\mathrm{pj}^{*}A^{(1)}\right)(\tilde{x},\tilde{x}')\\
&=\mathrm{Op}^{(2)}\left(s^{(2)},A^{(2)}\right)(\tilde{x},\tilde{x}'),
\end{align*}
which is the first assertion of~(b).

\emph{The polymer families.} In this paragraph $\mathrm{Op}$ stands for $M_k$ or for $S_{k,y}$ with some $y\ge0$, and $\tilde{X}$, $\tilde{x}$, $\tilde{x}'$, $A^{(1)}$, $A^{(2)}$ are as in the last assertion. Since $\tilde{X}^{\sim1}\subseteq\tilde{Y}^{\sim1}\subseteq N(\tilde{Y})$, the translates of $\tilde{X}^{\sim1}$ by the elements of $l_1\mathbb{Z}^3$ are pairwise disjoint, by~\eqref{eq:renormalisation-flow-a1}. Since $\mathrm{pr}_1=\mathrm{pj}\circ\mathrm{pr}_2$, the set $\mathrm{pj}^{-1}\mathrm{pr}_1(\tilde{X})$ is the union of the sets $\mathrm{pr}_2(\tilde{X}+g)$, where $g$ ranges over a system of representatives of $l_1\mathbb{Z}^3/l_2\mathbb{Z}^3$; by the disjointness just stated, no two of these sets contain $M$-cubes with a common face. Let $X'$ be a polymer of $T_2$ with $\mathrm{pj}(X')=\mathrm{pr}_1(\tilde{X})$. Being connected through faces, $X'$ lies in one of the sets $\mathrm{pr}_2(\tilde{X}+g)$. The covering map $\mathrm{pj}$ is injective on this set, because $\mathrm{pj}\circ\mathrm{pr}_2=\mathrm{pr}_1$ is injective on $\tilde{X}+g\subseteq N(\tilde{Y})+g$ by~\eqref{eq:renormalisation-flow-a1}; hence $X'=\mathrm{pr}_2(\tilde{X}+g)$. Conversely, $\mathrm{pr}_2(\tilde{X})$ is a polymer of $T_2$ with $\mathrm{pj}\left(\mathrm{pr}_2(\tilde{X})\right)=\mathrm{pr}_1(\tilde{X})$. By Definition~\ref{def:projection} and~\eqref{eq:renormalisation-flow-4},
\begin{multline*}
\mathrm{Op}^{(1)}\left(\mathrm{pr}_1(\tilde{X}),A^{(1)}\right)(\tilde{x},\tilde{x}')\\
=\sum_{w\in(\mathbb{Z}/J\mathbb{Z})^3}\boldsymbol{\chi}(w)\cdot\sum_{X'}\mathrm{Op}^{(2)}\left(X',\mathrm{pj}^{*}A^{(1)}\right)\left(\tilde{x},\tilde{x}'+l_1\cdot w\right),
\end{multline*}
where $X'$ ranges over the polymers of $T_2$ with $\mathrm{pj}(X')=\mathrm{pr}_1(\tilde{X})$. By property~\ref{it:renormalisation-group-12} on $T_2$, the kernel of $\mathrm{Op}^{(2)}\left(X',\,\cdot\,\right)$ vanishes unless both of its indices lie in $X'$. Let $X'=\mathrm{pr}_2(\tilde{X}+g)$. If the row index $\tilde{x}$ lies in $X'$, then $\tilde{x}+l_2\cdot v$ is contained in $\tilde{X}+g$ for some $v\in\mathbb{Z}^3$; as $\tilde{x}$ is contained in $\tilde{X}\subseteq N(\tilde{Y})$, the sets $N(\tilde{Y})$ and $N(\tilde{Y})+g-l_2\cdot v$ meet, so that $g=l_2\cdot v$ by~\eqref{eq:renormalisation-flow-a1} and $X'=\mathrm{pr}_2(\tilde{X})$. If the column index $\tilde{x}'+l_1\cdot w$ lies in $\mathrm{pr}_2(\tilde{X})$, then $\tilde{x}'+l_1\cdot(w+J\cdot v)$ is contained in $\tilde{X}$ for some $v\in\mathbb{Z}^3$; as $\tilde{x}'$ is contained in $\tilde{X}\subseteq N(\tilde{Y})$, the sets $N(\tilde{Y})$ and $N(\tilde{Y})+l_1\cdot(w+J\cdot v)$ meet, so that $w+J\cdot v=0$ by~\eqref{eq:renormalisation-flow-a1}. Thus only $X'=\mathrm{pr}_2(\tilde{X})$ and the class of $0$ contribute, and
\[
\mathrm{Op}^{(1)}\left(\mathrm{pr}_1(\tilde{X}),A^{(1)}\right)(\tilde{x},\tilde{x}')=\mathrm{Op}^{(2)}\left(\mathrm{pr}_2(\tilde{X}),\mathrm{pj}^{*}A^{(1)}\right)(\tilde{x},\tilde{x}').
\]
Finally, $\mathrm{pj}^{*}A^{(1)}$ lies in $\tilde{G}_k(T_2)$ by Lemma~\ref{lem:periodic-extension}, and it coincides with $A^{(2)}$ on the bonds that meet $\mathrm{pr}_2(\tilde{X})$: such a bond is $\mathrm{pr}_2(b)$ for a bond $b$ of $\mathbb{L}_k$ that meets $\tilde{X}$, and
\[
\left(\mathrm{pj}^{*}A^{(1)}\right)\left(\mathrm{pr}_2(b)\right)=A^{(1)}\left(\mathrm{pr}_1(b)\right)=A^{(2)}\left(\mathrm{pr}_2(b)\right),
\]
the connections $A^{(1)}$ and $A^{(2)}$ having equal lifts on the bonds that meet $\tilde{X}$. By property~\ref{it:renormalisation-group-12}, the connection $\mathrm{pj}^{*}A^{(1)}$ may therefore be replaced by $A^{(2)}$ in the kernel $\mathrm{Op}^{(2)}\left(\mathrm{pr}_2(\tilde{X}),\mathrm{pj}^{*}A^{(1)}\right)$, and the last assertion of the lemma follows.
\end{proof}

\subsubsection{Comparison of the steps on the two tori}\label{sec:comparison-steps-two}

\begin{proposition}[Agreement after one step]\label{prop:agreement-after-one}
Assume hypothesis~\ref{def:hypothesis-flow} and $R<R_{\mathrm{top}}$, let $0\le k<k_m$, and let $T_1$, $T_2$, $l_1$, $l_2$ be as in Section~\ref{sec:admissible-sets-local}. Let $m$ be a real number with $|m|\le w^{\mathrm{m}}_k$. For $i=1,2$ let $\mathscr{E}^{(i)}\in K^x_k(l_i)$ with $\|\mathscr{E}^{(i)}\|^x_k\le w^{\mathrm{E}}_k$, where $\mathscr{E}^{(i)}=0$ if $k=0$, and let $\mathfrak{P}^{(i)}_k$ be a family of complex functions on $\mathfrak{X}_k(T_i)$. Assume that $\mathscr{E}^{(1)}\sim\mathscr{E}^{(2)}$ and $\mathfrak{P}^{(1)}_k\sim\mathfrak{P}^{(2)}_k$ in the sense of Definition~\ref{def:admissibility-flatness-agreement}(e). Then the following hold.
\begin{enumerate}
\item[(a)] $\varepsilon(\mathscr{E}^{(1)})=\varepsilon(\mathscr{E}^{(2)})$ and $m(\mathscr{E}^{(1)})=m(\mathscr{E}^{(2)})$, and the number $d\varepsilon^0_k$ is the same on $T_1$ and on $T_2$.
\item[(b)] For $i=1,2$ let
\[
\mathscr{E}^{(i)+}:=L R\,\mathscr{E}^{(i)}+L\,{\delta E^{\mathrm{det}}}^{(i)}_k+E^{*,x,(i)}_k\left(m,\mathscr{E}^{(i)}\right),
\]
formed on $T_i$ with the operations of Section~\ref{sec:local-parts-adjustment}, the step map of Definition~\ref{def:renormalisation-step} and the determinant of Definition~\ref{def:determinant-functions}. Then $\mathscr{E}^{(1)+}\sim\mathscr{E}^{(2)+}$ at the level $k+1$, that is, for the tori $T_{k+1}(R)$ and $T_{k+1}(R_{\mathrm{top}})$.
\item[(c)] For $i=1,2$ let $\mathfrak{P}^{(i)}_{k+1}$ be formed from $\mathfrak{P}^{(i)}_k$ by the rule of clause~\eqref{it:step-small-fields-c} of Theorem~\ref{thm:step-small-fields}: for every polymer $Y$ of $T_{k+1}(R)$, resp.\ of $T_{k+1}(R_{\mathrm{top}})$, and every $\mathscr{C}'$ in $\mathfrak{X}_{k+1}$ of that torus,
\[
\mathfrak{P}^{(i)}_{k+1}(Y,\mathscr{C}')=\sum_{X:\ \bar{X}^L=LY}\left(\mathfrak{P}^{(i)}_k(X,\mathscr{C}'_L)+E^{\mathrm{det},(i)}_k(X,\mathscr{C}'_L)\right).
\]
Then $\mathfrak{P}^{(1)}_{k+1}\sim\mathfrak{P}^{(2)}_{k+1}$.
\end{enumerate}
In particular, let $E^{(i)}:=\mathscr{E}^{(i)}|_0$. By clause~\eqref{it:step-small-fields-d} of Theorem~\ref{thm:step-small-fields}, $\mathscr{E}^{(i)+}|_0=L R\,E^{(i)}+L\,E^{\mathrm{det},(i)}_k+E^{*,(i)}_k(m,E^{(i)})$, and agreement is preserved by the restriction to $\mathscr{C}=0$; hence
\[
L R\,E^{(1)}+L\,E^{\mathrm{det},(1)}_k+E^{*,(1)}_k\left(m,E^{(1)}\right)\sim L R\,E^{(2)}+L\,E^{\mathrm{det},(2)}_k+E^{*,(2)}_k\left(m,E^{(2)}\right).
\]
\end{proposition}

\begin{proof}
\emph{Part (a).} Let $\mathscr{E}$ stand for $\mathscr{E}^{(1)}$ or for $\mathscr{E}^{(2)}$, on the corresponding torus. By Proposition~\ref{prop:symmetries-extraction} the numbers $\varepsilon(\mathscr{E})$ and $m(\mathscr{E})$ do not depend on the $M$-cube $\square$ with which they are defined in Section~\ref{sec:local-parts-adjustment}. By that definition, $-\varepsilon(\mathscr{E})$ and $-m(\mathscr{E})\cdot\mathbf{1}$ (with $\mathbf{1}$ the unit matrix) are the sums, over the small polymers $X$ that contain $\square$, of $\boldsymbol{\alpha}_0(\mathscr{E},X)$ and of $\boldsymbol{\alpha}_2(\mathscr{E},X)$, respectively, where
\begin{gather*}
\boldsymbol{\alpha}_0(\mathscr{E},X)=\operatorname{Vol}(X)^{-1}\cdot\mathscr{E}_{00}(X;0,0),\\
\boldsymbol{\alpha}_2(\mathscr{E},X)=\operatorname{Vol}(X)^{-1}\int\mathrm{d}x\,\mathrm{d}y\;\mathscr{E}_{20}\left(X;0,0;(x,\cdot,1),(y,\cdot,0)\right),
\end{gather*}
all arguments being taken in one lift of $X$ (Section~\ref{sec:antiperiodicity}). For a small flat polymer $X=\mathrm{pr}_i(\tilde{X})$ these quantities are thus formed from the coefficients of $\mathscr{E}(X;0,0)$ at the generators whose indices lie in the lift $\tilde{X}$. The zero connections on $T_1$ and their images $\mathrm{pj}^{*}(0,0)=(0,0)$ on $T_2$ have equal lifts; so $\mathscr{E}^{(1)}\sim\mathscr{E}^{(2)}$ gives $\mathfrak{i}\,\mathscr{E}^{(1)}(\mathrm{pr}_1(\tilde{X});0,0)=\mathscr{E}^{(2)}(\mathrm{pr}_2(\tilde{X});0,0)$ for every connected admissible $\tilde{X}$, that is, the two elements have the same coefficients at generators with the same indices.

Fix an $M$-cube $\tilde{\square}$ of $\mathbb{L}_k$ and take $\square:=\mathrm{pr}_i(\tilde{\square})$ on $T_i$. By Lemma~\ref{lem:agreement-criterion-small}(b), for $i=1,2$ the map $\tilde{X}\mapsto\mathrm{pr}_i(\tilde{X})$ is a bijection of the set of the connected $\tilde{X}$ with $d_M(\tilde{X})\le L$ that contain $\tilde{\square}$ onto the set of the small polymers of $T_i$ that contain $\mathrm{pr}_i(\tilde{\square})$; these $\tilde{X}$ are admissible, and $\operatorname{Vol}(\mathrm{pr}_1(\tilde{X}))=\operatorname{Vol}(\mathrm{pr}_2(\tilde{X}))$, both numbers being $M^3$ times the number of the $M$-cubes of $\tilde{X}$. Consequently
\begin{gather*}
\boldsymbol{\alpha}_0\left(\mathscr{E}^{(1)},\mathrm{pr}_1(\tilde{X})\right)=\boldsymbol{\alpha}_0\left(\mathscr{E}^{(2)},\mathrm{pr}_2(\tilde{X})\right),\\
\boldsymbol{\alpha}_2\left(\mathscr{E}^{(1)},\mathrm{pr}_1(\tilde{X})\right)=\boldsymbol{\alpha}_2\left(\mathscr{E}^{(2)},\mathrm{pr}_2(\tilde{X})\right)
\end{gather*}
for each such $\tilde{X}$, and the sums that define $\varepsilon$ and $m$ on the two tori coincide term by term. Finally, $d\varepsilon^0_k=-d_V\cdot\log P\left(|g|\le p_{0,k}\right)$ with a standard normal variable $g$ (Proposition~\ref{prop:basis-map-box}); since $d_V$ depends on $L$ only and $p_{0,k}$ is the same on both tori (Section~\ref{sec:admissible-sets-local}), so is $d\varepsilon^0_k$.

\emph{Part (b), Step 0 (reduction to the level $k$).} Consider a union $\tilde{U}$ of $M$-cubes of $\mathbb{L}_{k+1}$, connected through faces, for which
\[
\operatorname{diam}(\tilde{U})\le l_1/L-14d_U\cdot M .
\]
As $l_{k+1}(R)=l_1/L$, this inequality is the admissibility of $\tilde{U}$ at the level $k+1$ (Definition~\ref{def:admissibility-flatness-agreement}(a)). The dilate $\tilde{V}:=L\,\tilde{U}$ is a union of $LM$-cubes of $\mathbb{L}_k$, and $\mathrm{box}(\tilde{V})=L\,\mathrm{box}(\tilde{U})$, so that
\[
\operatorname{diam}(\tilde{V})=L\cdot\operatorname{diam}(\tilde{U})\le l_1-14d_U\cdot LM\le l_1-14d_U\cdot M ;
\]
thus $\tilde{V}$ is admissible at the level $k$. If $\tilde{X}\subset\tilde{V}$ is a union of $M$-cubes of $\mathbb{L}_k$, then $\mathrm{box}(\tilde{X})\subset\mathrm{box}(\tilde{V})$; hence $\tilde{X}$ is admissible at the level $k$ too, and $N(\tilde{X})\subset N(\tilde{V})$. Finally, $N(\tilde{V})\subset L\,N(\tilde{U})$: the $6d_U$ layers of $M$-cubes around $\mathrm{box}(\tilde{V})$ lie within the $6d_U$ layers of $LM$-cubes around $L\,\mathrm{box}(\tilde{U})$.

By Proposition~\ref{prop:data-along-flow}, Theorem~\ref{thm:step-small-fields} applies at the level $k$ on both tori. By its clause~\eqref{it:step-small-fields-d}, $\mathscr{E}^{(i)+}$ lies in $K^x_{k+1}(l_i/L)$, whose elements have the locality property~\eqref{it:renormalisation-group-1-k2}; so $\mathscr{E}^{(i)+}$ is a local family of the level $k+1$ (Definition~\ref{def:admissibility-flatness-agreement}(d)). Applied to $\tilde{U}$ and to objects on $T_{k+1}(R)$, the symbols $\mathrm{pr}_i$ and $\mathrm{pj}^{*}$ stand for the maps of Section~\ref{sec:admissible-sets-local} formed at the level $k+1$: the projection of $\mathbb{L}_{k+1}$ onto $T_{k+1}(R)$ for $i=1$ and onto $T_{k+1}(R_{\mathrm{top}})$ for $i=2$, and the inclusion of the objects on $T_{k+1}(R)$ into those on $T_{k+1}(R_{\mathrm{top}})$. Applied to $\tilde{V}$, to its subsets and to objects on $T_1$, they keep their meaning at the level $k$, as do $T_1$, $T_2$, $l_1$, $l_2$. Let $Y_i:=\mathrm{pr}_i(\tilde{U})$, a polymer of $T_{k+1}(R)$ for $i=1$ and of $T_{k+1}(R_{\mathrm{top}})$ for $i=2$. By Lemma~\ref{lem:agreement-criterion-small}(a), applied at the level $k+1$, it suffices to show
\[
\mathfrak{i}\,\mathscr{E}^{(1)+}(Y_1;\mathscr{A},\mathscr{C})=\mathscr{E}^{(2)+}\left(Y_2;\mathrm{pj}^{*}\mathscr{A},\mathrm{pj}^{*}\mathscr{C}\right)
\]
for all such $\tilde{U}$ and all $(\mathscr{A},\mathscr{C})\in\tilde{R}_{k+1}\times\mathfrak{X}_{k+1}$ on $T_{k+1}(R)$, where $\mathfrak{i}$ is the isomorphism of Definition~\ref{def:admissibility-flatness-agreement}(c) for $\tilde{U}$ at the level $k+1$. Put $(A_1,\mathscr{C}_1):=(\mathscr{A}_L,\mathscr{C}_L)$, which lies in $\tilde{R}^{-}_k(T_1)\times\mathfrak{X}_k(T_1)$ by Lemma~\ref{lem:rescaling}, and $(A_2,\mathscr{C}_2):=\left((\mathrm{pj}^{*}\mathscr{A})_L,(\mathrm{pj}^{*}\mathscr{C})_L\right)=\mathrm{pj}^{*}(A_1,\mathscr{C}_1)$. Let $S_L$ be the homomorphism from the Grassmann algebra of the level $k$ on $T_i$ into that of the level $k+1$ with
\[
S_L\,\psi(x)=L^{-1}\cdot\psi(x/L),\qquad S_L\,\chi(x,x')=L^{-1-\alpha}\cdot\chi(x/L,x'/L).
\]
By the display~\eqref{eq:renormalisation-group-62}, for $i=1,2$,
\[
\mathscr{E}^{(i)+}(Y_i;\,\cdot\,)=\sum_{X:\ \bar{X}^L=LY_i}S_L\left[\left((R\,\mathscr{E}^{(i)})^{\mathrm{nat}}+{\delta E^{\mathrm{det}}}^{(i)}_k+E^{\#,x,(i)}_k\left(m,\mathscr{E}^{(i)}\right)\right)(X;A_i,\mathscr{C}_i)\right],
\]
where the dot stands for $(\mathscr{A},\mathscr{C})$ if $i=1$ and for $(\mathrm{pj}^{*}\mathscr{A},\mathrm{pj}^{*}\mathscr{C})$ if $i=2$. One has $LY_i=\mathrm{pr}_i(\tilde{V})$. By Lemma~\ref{lem:agreement-criterion-small}(c), the polymers $X$ of $T_i$ with $\bar{X}^L=\mathrm{pr}_i(\tilde{V})$ are the sets $\mathrm{pr}_i(\tilde{X})$, where $\tilde{X}$ is a connected union of $M$-cubes contained in $\tilde{V}$ such that the union of the $LM$-cubes that contain an $M$-cube of $\tilde{X}$ is $\tilde{V}$; so the two sums are indexed by the same sets $\tilde{X}$. The homomorphism $S_L$ intertwines the isomorphism $\mathfrak{i}$ of the level $k$ for $\tilde{V}$ with the isomorphism $\mathfrak{i}$ of the level $k+1$ for $\tilde{U}$. Indeed, by the display that defines $S_L$, on $T_i$, $i=1,2$, the generator $\psi(x,\cdot)$ labelled by a site $x$ of $N(\tilde{V})$ is mapped to $L^{-1}$ times the generator $\psi(x/L,\cdot)$ of the level $k+1$, and the generator $\chi(x,x',\cdot)$ labelled by two sites $x$, $x'$ of $N(\tilde{V})$ with $0<|x-x'|\le M$ is mapped to $L^{-1-\alpha}$ times $\chi(x/L,x'/L,\cdot)$; this rule does not contain $l_i$. As $N(\tilde{V})\subset L\,N(\tilde{U})$, the points $x/L$ and $x'/L$ are sites of $N(\tilde{U})$, and $0<|x/L-x'/L|\le M/L\le M$. So the images are multiples of generators listed for $\tilde{U}$ at the level $k+1$ in Definition~\ref{def:admissibility-flatness-agreement}(c), labelled by $x/L$, resp.\ by $(x/L,x'/L)$, on both tori. The isomorphism $\mathfrak{i}$ of the level $k$ for $\tilde{V}$ sends each listed generator of $T_1$ to the listed generator of $T_2$ with the same label, and the isomorphism $\mathfrak{i}$ of the level $k+1$ for $\tilde{U}$ sends each listed generator of $T_{k+1}(R)$ to the listed generator of $T_{k+1}(R_{\mathrm{top}})$ with the same label (Definition~\ref{def:admissibility-flatness-agreement}(c)). Consequently the homomorphisms $\mathfrak{i}\circ S_L$ and $S_L\circ\mathfrak{i}$, where the first $\mathfrak{i}$ is that of the level $k+1$ for $\tilde{U}$ and the second that of the level $k$ for $\tilde{V}$, agree on the listed generators $\psi(x,\cdot)$ and $\chi(x,x',\cdot)$ of $T_1$ for $\tilde{V}$, and hence on the smallest subalgebra that contains these generators. It therefore suffices to prove the following: for every connected union $\tilde{X}$ of $M$-cubes contained in $\tilde{V}$, every $(A_1,\mathscr{C}_1)\in\tilde{R}^{-}_k(T_1)\times\mathfrak{X}_k(T_1)$ and $(A_2,\mathscr{C}_2):=\mathrm{pj}^{*}(A_1,\mathscr{C}_1)$, each of the three families $F=(R\,\mathscr{E})^{\mathrm{nat}}$, $F=\delta E^{\mathrm{det}}_k$ and $F=E^{\#,x}_k(m,\mathscr{E})$ of the level $k$ satisfies
\[
\mathfrak{i}\,F^{(1)}\left(\mathrm{pr}_1(\tilde{X});A_1,\mathscr{C}_1\right)=F^{(2)}\left(\mathrm{pr}_2(\tilde{X});A_2,\mathscr{C}_2\right),
\]
where $F^{(i)}$ is the family formed on $T_i$, with $\mathscr{E}^{(i)}$ and the primitives on $T_i$. Put $\mathscr{T}_i:=A_i+\mathscr{C}_i$; then $\mathscr{T}_2=\mathrm{pj}^{*}\mathscr{T}_1$. In Steps 1--4, $\mathfrak{i}$ denotes the isomorphism of Definition~\ref{def:admissibility-flatness-agreement}(c) for $\tilde{V}$ at the level $k$.

\emph{Step 1 (the first family: extraction and adjustment).} By Lemma~\ref{lem:agreement-criterion-small}(b), $\mathrm{pr}_1(\tilde{X})$ is small if and only if $\mathrm{pr}_2(\tilde{X})$ is small. By Section~\ref{sec:local-parts-adjustment}, $(R\,\mathscr{E})(X)=\mathscr{E}(X)$ for a large polymer $X$, and
\[
(R\,\mathscr{E})(X)=\mathscr{E}(X)-\boldsymbol{\alpha}_0(\mathscr{E},X)\cdot\operatorname{Vol}(X)-\int_X\mathrm{d}x\;\bar{\psi}(x)\cdot\boldsymbol{\alpha}_2(\mathscr{E},X)\cdot\psi(x)
\]
for a small one. Together with $\mathscr{E}^{(1)}\sim\mathscr{E}^{(2)}$, part (a) and its proof, which give the equality of $\boldsymbol{\alpha}_0$, of $\boldsymbol{\alpha}_2$ and of the volumes at $\mathrm{pr}_1(\tilde{X})$ and at $\mathrm{pr}_2(\tilde{X})$, therefore yield $R\,\mathscr{E}^{(1)}\sim R\,\mathscr{E}^{(2)}$; here $R\,\mathscr{E}^{(i)}$ is a local family, because $R$ maps $K^x_k$ into itself (Proposition~\ref{prop:symmetries-extraction}). Next, $(R\,\mathscr{E})^{\mathrm{nat}}(X)=(R\,\mathscr{E})(X)$ for large $X$, and for small $X$
\[
(R\,\mathscr{E})^{\mathrm{nat}}(X)=(R\,\mathscr{E})(X)-(R\,\mathscr{E})^{[2,0]}(X)+\boldsymbol{\Theta}_{X,\mathscr{T}}\left((R\,\mathscr{E})^{[2,0]}(X)\right),
\]
where the upper index $[2,0]$ marks the part of bidegree $(2,0)$ and $\boldsymbol{\Theta}_{X,\mathscr{T}}$ is the homomorphism of Definition~\ref{def:transport-mean} (Section~\ref{sec:local-parts-adjustment}). In the sense of the last paragraph of Definition~\ref{def:terminology}(2), the homomorphism $\boldsymbol{\Theta}_{X,\mathscr{T}}$ is given by one formula on $\mathbb{L}_k$ with the range $0$, the datum being $\mathscr{T}$ on the bonds with both ends in a lift $\tilde{X}$ of $X$. The formula is the following map $\Phi$. For a finite union $\tilde{X}$ of $M$-cubes of $\mathbb{L}_k$ that is connected through faces, and for a bond function $\tilde{B}$ on the bonds with both ends in $\tilde{X}$, let $\Phi(\tilde{X},\tilde{B})$ be the homomorphism defined by the displays of Definition~\ref{def:transport-mean} and~\eqref{eq:polymer-activities-17} when they are read in $\mathbb{L}_k$, all points being taken in $\tilde{X}$ and the generators being labelled by their indices in $\tilde{X}$. It is built from five data. The first four are formed from $\tilde{X}$ alone: the number $\operatorname{Vol}(X)$, which is $M^3$ times the number of the $M$-cubes of $\tilde{X}$; the integration over $X$, which is the sum over the sites of $\tilde{X}$ with the weight $\eta_k^3$; for two $M$-cubes of $\tilde{X}$, the arithmetic mean over those chains of $M$-cubes of $\tilde{X}$ which lead from the one to the other, consecutive cubes having a common face, and whose length is minimal; and the centres of the cubes in such chains. The fifth consists of the transports at the pairs $(p,p')$ of kind (I) for $X$ (Definition~\ref{def:pairs-sites-holder}; see~\eqref{eq:polymer-activities-17}),
\[
u_{\tilde{B}}(p,p')=\exp\left\{i\cdot e_k\cdot\eta_k\cdot(\tau\tilde{B})(p,p')\right\},
\]
which read $\tilde{B}$ only on the bonds with both ends in $\tilde{X}$, because the six paths of each such pair lie in $\tilde{X}$ (see~\eqref{eq:polymer-activities-ch1}). None of the five data contains $l$, and a translation by a vector of $(L\cdot M)\mathbb{Z}^3$ maps $M$-cubes, their centres, their common faces, chains of $M$-cubes and the paths of pairs onto such; thus $\Phi$ does not depend on $l$ and commutes with these translations. The value of $\Phi$ at a lift $\tilde{X}$ of $X$ and at the values of $\mathscr{T}$, carried over by the covering map to the bonds with both ends in $\tilde{X}$, is $\boldsymbol{\Theta}_{X,\mathscr{T}}$. This value does not depend on the lift. Indeed, a change of the lift by a vector of $l\mathbb{Z}^3\subset(L\cdot M)\mathbb{Z}^3$ translates all the data, the connection being periodic; by the covariance of $\Phi$ it therefore only relabels the generators, and the relabelling multiplies every generator by one and the same sign, whereas $\Phi(\tilde{X},\tilde{B})$ maps each generator to a linear combination of generators. Now let $\mathrm{pr}_1(\tilde{X})$, and hence $\mathrm{pr}_2(\tilde{X})$, be small, and apply this to $X=\mathrm{pr}_i(\tilde{X})$ and $\mathscr{T}=\mathscr{T}_i$, with the lift $\tilde{X}$. For $i=1$ and for $i=2$, the projection $\mathrm{pr}_i$ maps the $M$-cubes of $\tilde{X}$ bijectively onto those of $\mathrm{pr}_i(\tilde{X})$ and the centre of each of them onto the centre of its image, and two $M$-cubes of $\tilde{X}$ have a common face if and only if their images have one (Lemma~\ref{lem:agreement-criterion-small}(c)). Moreover, as $\mathscr{T}_2=\mathrm{pj}^{*}\mathscr{T}_1$ and $\mathrm{pj}\circ\mathrm{pr}_2=\mathrm{pr}_1$,
\[
\mathscr{T}_1\left(\mathrm{pr}_1(b)\right)=\mathscr{T}_2\left(\mathrm{pr}_2(b)\right)\qquad\text{for every bond $b$ of $\mathbb{L}_k$.}
\]
So the five data, formed in $\tilde{X}$, are the same for $\left(\mathrm{pr}_1(\tilde{X}),\mathscr{T}_1\right)$ and for $\left(\mathrm{pr}_2(\tilde{X}),\mathscr{T}_2\right)$: both homomorphisms are values of $\Phi$ at $\tilde{X}$ and at one and the same bond function on the bonds with both ends in $\tilde{X}$. As $\tilde{X}$ is admissible, its sites have mutual sup distances less than $\operatorname{diam}(\tilde{X})\le l_1-14d_U\cdot M$; so these distances plus $M$ are less than $l_1\le l_i$, and the margin that Definition~\ref{def:terminology}(2) requires for the range $0$ holds on both tori. By the comparison of two tori in that definition, the two homomorphisms are intertwined by the isomorphism that sends each generator of $T_1$ labelled by an index in $\tilde{X}$ to the generator of $T_2$ with the same label, and this isomorphism is the restriction of $\mathfrak{i}$ (Definition~\ref{def:admissibility-flatness-agreement}(c)). Therefore
\[
\mathfrak{i}\circ\boldsymbol{\Theta}_{\mathrm{pr}_1(\tilde{X}),\mathscr{T}_1}=\boldsymbol{\Theta}_{\mathrm{pr}_2(\tilde{X}),\mathscr{T}_2}\circ\mathfrak{i} .
\]
Since $\mathfrak{i}$ preserves the bidegree, the first family agrees.

\emph{Step 2 (the second family: the determinant).} Let $X$ be a polymer of $T_1$ or of $T_2$, and $B$ an element of $\tilde{G}_k$ on the same torus. By property~\ref{it:renormalisation-group-12}, the traces in the displays of $\acute{E}_k$ (Definition~\ref{def:contour-functional-functions}; see also Section~\ref{sec:fermion-determinant}) are sums over the unit sites of $X$: for $k\ge1$ the trace is
\[
\sum_{x_1,x_2}\operatorname{tr}\left[\left(B_{k,y}\cdot\gamma_3\cdot B_{k,y}\right)(x_2,x_1)\cdot\left(Q_k\cdot S_{k,y}(X,B)\cdot Q^T_k\right)(x_1,x_2)\right],
\]
with $B_{k,y}=B_{k,y}(B)$, $Q_k=Q_k(B)$ and $Q^T_k=Q^T_k(-B)$, and for $k=0$ the two traces are
\begin{gather*}
\sum_{x,x'}\operatorname{tr}\left[F_0(B)(x,x')\cdot S_{0,y}(X,B)(x',x)\right],\\
\sum_x\operatorname{tr}\left[\gamma_3\cdot S_{0,y}(X,B)(x,x)\right],
\end{gather*}
where $F_0(B):=D_B+\mathfrak{m}+bL^{-1}\cdot P(B)$. The operators $Q_k(B)$, $Q^T_k(-B)$ and $B_{k,y}(B)$ of Section~\ref{sec:exact-fermion-operators}, and for $k=0$ the operator $F_0(B)$, are exact (Definition~\ref{def:terminology}(3)), and each of them is given by one formula on $\mathbb{L}_k$ in the sense of Definition~\ref{def:terminology}(2), with the range $L$. Indeed, let $\tilde{B}$ be a bond function on $\mathbb{L}_k$; in what follows a number stands for the corresponding multiple of the unit $4\times4$ matrix, and the weights are counted into the entries. By Sections~\ref{sec:exact-fermion-operators} and~\ref{sec:fermion-integrals-blocking}, for $k\ge1$, a unit site $x_1\in\mathbb{Z}^3$ and a site $x\in B_k(x_1)$ of $\mathbb{L}_k$,
\[
Q_k(\tilde{B})(x_1,x)=\eta_k^3\cdot\exp\left(i\cdot e_k\cdot\eta_k\cdot(\tau_k\tilde{B})(x_1,x)\right),\qquad Q^T_k(-\tilde{B})(x,x_1)=\exp\left(-i\cdot e_k\cdot\eta_k\cdot(\tau_k\tilde{B})(x_1,x)\right),
\]
and both entries vanish for the sites $x$ of $\mathbb{L}_k$ outside $B_k(x_1)$. For two unit sites $x_1$, $x_2$ the entry $P(\tilde{B})(x_1,x_2)$ vanishes unless $x_1$ and $x_2$ lie in one $L$-block $B(y')$, $y'\in L\mathbb{Z}^3$, and in that case, as $P(\tilde{B})=Q^T(-\tilde{B})\cdot Q(\tilde{B})$,
\[
P(\tilde{B})(x_1,x_2)=L^{-3}\cdot\exp\left(-i\cdot e_k\cdot\eta_k\cdot(\tau\tilde{B})(y',x_1)\right)\cdot\exp\left(i\cdot e_k\cdot\eta_k\cdot(\tau\tilde{B})(y',x_2)\right);
\]
moreover
\[
B_{k,y}(\tilde{B})=\left(b_k+i\gamma_3y\right)^{-1}\cdot\left(1-P(\tilde{B})\right)+\left(b_k+bL^{-1}+i\gamma_3y\right)^{-1}\cdot P(\tilde{B}).
\]
For $k=0$, where $\eta_0=1$, one has $F_0(\tilde{B})=D_{\tilde{B}}+\mathfrak{m}+bL^{-1}\cdot P(\tilde{B})$, with $D_{\tilde{B}}(x,x)=3$, with $D_{\tilde{B}}(x,x\pm e_{\mu})=-\frac{1\mp\gamma_{\mu}}{2}\cdot\exp\left(i\cdot e_k\cdot\tilde{B}(x,x\pm e_{\mu})\right)$ for $\mu=0,1,2$, and with $D_{\tilde{B}}(x,x')=0$ for all other pairs of sites. These entries do not contain $l$, no set occurs in them, and each row has finite support. The translations by the vectors of $(L\cdot M)\mathbb{Z}^3$ map the blocks of the sides $L^j\eta_k$, the blocks $B_k(y)$, the $L$-blocks, their centres and the paths $\Gamma^{\pi}$ onto such; and on $T_i$ the operators $Q_k(B)$, $Q^T_k(-B)$, $P(B)$, $B_{k,y}(B)$ and $D_B$ are these formulas, formed with the periodic extension $\tilde{B}$ of $B$ and applied to the antiperiodic functions (the proof of Proposition~\ref{prop:operators-covering-torus}, observation (i); compare Proposition~\ref{prop:coverings-tori}). Thus the conditions (i)--(iii) of Definition~\ref{def:terminology}(2) hold. The entry at $(x,x')$ vanishes unless $|x-x'|<\frac{1}{2}$ for $Q_k$ and $Q^T_k$, unless $|x-x'|\le L-1$ for $P$ and $B_{k,y}$ (Proposition~\ref{prop:kernel-bounds-2}), and unless $|x-x'|\le1$ for $D_{\tilde{B}}$. The row of $Q_k(\tilde{B})$ at $x_1$ and the entries $Q^T_k(-\tilde{B})(x,x_1)$ read $\tilde{B}$ only on the fine bonds with both ends in $B_k(x_1)$ (the proof of Lemma~\ref{lem:interior-boundary-parts}(b)); the entry $P(\tilde{B})(x_1,x_2)$, and with it $B_{k,y}(\tilde{B})(x_1,x_2)$, reads $\tilde{B}$ only on the fine bonds of which the unit-lattice paths $\Gamma^{\pi}(y',x_1)$ and $\Gamma^{\pi}(y',x_2)$ consist, and the sites of these paths lie within $\frac{L-1}{2}$ of $y'$; the entry $F_0(\tilde{B})(x,x')$ reads $\tilde{B}$, besides, on the bond $(x,x')$. All these bonds have both ends within $L-1$ of the index of the row; so $L$ is a range of these formulas. By~\eqref{it:renormalisation-group-g0} (Section~\ref{sec:cubes-parameter-cubes}), the open cube of side $1$ around $x_1$ and the open cube of side $L$ around $y'$ lie each in one $M$-cube. Hence, at indices in a polymer $X$, all the bonds just named have both ends in $X$, and the intermediate indices of the products $B_{k,y}\cdot\gamma_3\cdot B_{k,y}$ and $Q_k\cdot S_{k,y}(X,B)\cdot Q^T_k$ in the traces above lie in $X$.

Now let $B\in\tilde{G}_k(T_1)$. Then $\mathrm{pj}^{*}B\in\tilde{G}_k(T_2)$ (Lemma~\ref{lem:periodic-extension}), and $B$ and $\mathrm{pj}^{*}B$ have one and the same periodic extension $\tilde{B}$ to $\mathbb{L}_k$; in particular they have equal lifts on all bonds. Take all indices in the lift $\tilde{X}$ of $\mathrm{pr}_i(\tilde{X})$. As $\tilde{X}$ is contained in the admissible set $\tilde{V}$, the sites of $\tilde{X}$ have mutual sup distances less than $\operatorname{diam}(\tilde{X})\le l_1-14d_U\cdot M$; so these distances plus $L$ are less than $l_1\le l_i$. Let $O$ be one of the operators $Q_k$, $Q^T_k$, $P$, $B_{k,y}$, $F_0$, formed on $\mathbb{L}_k$ with $\tilde{B}$ as above, and let $O^{(i)}$ be the same operator on $T_i$, formed with $B$ if $i=1$ and with $\mathrm{pj}^{*}B$ if $i=2$. By condition (iii) of Definition~\ref{def:terminology}(2), the kernel of $O^{(i)}$ at two indices $x$, $x'$, taken in $\mathbb{L}_k$ or in $\mathbb{Z}^3$, is the finite sum
\[
O^{(i)}(x,x')=\sum_{n\in\mathbb{Z}^3}(-1)^{n_0+n_1+n_2}\cdot O(x,x'+l_in).
\]
If $x$ and $x'$ are contained in $\tilde{X}$ and $n\ne0$, then $|x-x'-l_in|\ge l_i-|x-x'|>L$, so that the term with the index $n$ vanishes, $L$ being a range. Hence $O^{(1)}(x,x')=O(x,x')=O^{(2)}(x,x')$ at any two indices $x$, $x'$ contained in $\tilde{X}$; this is the comparison of two tori in Definition~\ref{def:terminology}(2), applied with the range $L$ and the set $\tilde{X}$. So the entries of the exact operators just named coincide on the two tori, and so do the kernels of $S_{k,y}^{(1)}(\mathrm{pr}_1(\tilde{X}),B)$ and of $S_{k,y}^{(2)}(\mathrm{pr}_2(\tilde{X}),\mathrm{pj}^{*}B)$ at all pairs of indices contained in $\tilde{X}$, by the last assertion of Lemma~\ref{lem:kernels-near-admissible}, applied with $\tilde{Y}:=\tilde{V}$. Hence the integrands in the displays of $\acute{E}_k$ coincide for every $y$, and so do the integrals:
\[
\acute{E}^{(1)}_k\left(\mathrm{pr}_1(\tilde{X}),B\right)=\acute{E}^{(2)}_k\left(\mathrm{pr}_2(\tilde{X}),\mathrm{pj}^{*}B\right).
\]
By Lemma~\ref{lem:invariances} the fields $-B$, $\overline{B}$ and $-\overline{B}$ lie in $\tilde{G}_k(T_1)$ as well, so this identity holds at each of them. By Definition~\ref{def:determinant-functions}, $E^d_k(X,B)$ is the mean of $\acute{E}_k(X,B)$, $\acute{E}_k(X,-B)$, $\overline{\acute{E}_k(X,\overline{B})}$ and $\overline{\acute{E}_k(X,-\overline{B})}$; since $\mathrm{pj}^{*}(-B)=-\mathrm{pj}^{*}B$ and $\mathrm{pj}^{*}\overline{B}=\overline{\mathrm{pj}^{*}B}$, it follows that
\[
E^{d,(1)}_k\left(\mathrm{pr}_1(\tilde{X}),B\right)=E^{d,(2)}_k\left(\mathrm{pr}_2(\tilde{X}),\mathrm{pj}^{*}B\right).
\]
The fields $\mathscr{T}_1$, $\mathscr{C}_1$ and $0$ lie in $\tilde{G}_k(T_1)$ (Lemma~\ref{lem:inclusions-bounds}), and $\mathrm{pj}^{*}\mathscr{T}_1=\mathscr{T}_2$, $\mathrm{pj}^{*}\mathscr{C}_1=\mathscr{C}_2$. Since $\delta E^{\mathrm{det}}_k(X;A,\mathscr{C})=E^d_k(X,A+\mathscr{C})-E^d_k(X,\mathscr{C})$ and $E^{\mathrm{det}}_k(X,\mathscr{C})=E^d_k(X,\mathscr{C})-E^d_k(X,0)$ (Definition~\ref{def:determinant-functions}), the choices $B=\mathscr{T}_1$, $B=\mathscr{C}_1$ and $B=0$ give
\begin{gather*}
{\delta E^{\mathrm{det}}}^{(1)}_k\left(\mathrm{pr}_1(\tilde{X});A_1,\mathscr{C}_1\right)={\delta E^{\mathrm{det}}}^{(2)}_k\left(\mathrm{pr}_2(\tilde{X});A_2,\mathscr{C}_2\right),\\
E^{\mathrm{det},(1)}_k\left(\mathrm{pr}_1(\tilde{X}),\mathscr{C}_1\right)=E^{\mathrm{det},(2)}_k\left(\mathrm{pr}_2(\tilde{X}),\mathscr{C}_2\right).
\end{gather*}
Both sides of the first identity are numbers, which $\mathfrak{i}$ leaves unchanged; so the second family agrees.

\emph{Step 3 (localisation over a lifted polymer).} As in Step 0, $\mathfrak{i} \colon \mathrm{Gr}_1(\tilde{V}) \to \mathrm{Gr}_2(\tilde{V})$ denotes the isomorphism of Definition~\ref{def:admissibility-flatness-agreement}(c) attached to $\tilde{V}$. By Proposition~\ref{prop:data-along-flow}, hypothesis~\ref{def:hypothesis-one-step} holds at level $k$ on $T_1$ and on $T_2$; hence clause~\eqref{it:step-small-fields-a} of Theorem~\ref{thm:step-small-fields}, Lemma~\ref{lem:bounds-analyticity}, Corollary~\ref{cor:renormalisation-group-1} and Proposition~\ref{prop:local-constituents-step} hold on both tori.

Let $\tilde{\mathbf{Y}}$ be a union of $M$-cubes of $\mathbb{L}_k$ which is connected through faces and contained in $\tilde{V}$; by Step 0 it is admissible. Put $\mathbf{Y}_i := \mathrm{pr}_i(\tilde{\mathbf{Y}})$ for $i = 1, 2$, and let $\mathsf{B}$ be the set of the basis elements $\Upsilon$ of $\mathbb{L}_k$ that are contained in $\tilde{\mathbf{Y}}$. By~\eqref{eq:renormalisation-flow-a1} the map $\mathrm{pr}_i$ is injective on $N(\tilde{\mathbf{Y}})$, so that the following two definitions make sense. For
\[
z = (z_{\Upsilon})_{\Upsilon \in \mathsf{B}} \in [-p_{0,k}, p_{0,k}]^{\mathsf{B}}
\]
let $\mathsf{Z}^{(i)} \in \mathrm{Box}_k(T_i)$ be the element whose coordinate at $\mathrm{pr}_i(\Upsilon)$ is $z_{\Upsilon}$ for every $\Upsilon \in \mathsf{B}$, and whose coordinate at every other basis element of $T_i$ is $0$. Let $\boldsymbol{\sigma} = (\boldsymbol{\sigma}_{\tilde{c}})$ assign to each parameter cube $\tilde{c} \subset \tilde{\mathbf{Y}}$ of $\mathbb{L}_k$, of one of the three kinds $\pi_k$, $\pi'_k$, $\pi''_k$, a complex number $\boldsymbol{\sigma}_{\tilde{c}}$ with $|\boldsymbol{\sigma}_{\tilde{c}}| \le M^{\alpha_0}$, and let $s^{(i)} \in \mathrm{Pd}_k(T_i)$ be the parameter family which takes the value $\boldsymbol{\sigma}_{\tilde{c}}$ at $\mathrm{pr}_i(\tilde{c})$ and the value $0$ at every other parameter cube of $T_i$. Again by~\eqref{eq:renormalisation-flow-a1}, $\mathsf{Z}^{(1)}$ and $\mathsf{Z}^{(2)}$ have equal lifts on $N(\tilde{\mathbf{Y}})$. Every parameter cube of $T_1$ at which $s^{(1)}$ is not $0$ has the form $\mathrm{pr}_1(\tilde{c})$ with $\tilde{c} \subset \tilde{\mathbf{Y}}$, and is therefore contained in $\mathbf{Y}_1$; at each cube $\mathrm{pr}_2(\tilde{c})$ with $\tilde{c} \subset \tilde{\mathbf{Y}}$ the family $s^{(2)}$ takes the value of $s^{(1)}$ at $\mathrm{pr}_1(\tilde{c})$, and at all other parameter cubes of $T_2$ it is $0$. Thus $s^{(1)}$ and $s^{(2)}$ stand in the relation of Lemma~\ref{lem:kernels-near-admissible}, taken with $\tilde{Y} := \tilde{\mathbf{Y}}$. As in Section~\ref{sec:data-variables-fields}, write $s^{(i)} = \left(s'^{(i)}, s''^{(i)}\right)$, where $s'^{(i)}$ consists of the values at the cubes of $\pi'_k$ and of $\pi''_k$, and $s''^{(i)}$ of the values at the cubes of $\pi_k$. On $T_i$ let $\mathscr{Z}^{(i)}\left(s'^{(i)}\right)$ be the weakened fluctuation field $\mathscr{Z}(s')$ of~\eqref{eq:renormalisation-group-56}, formed with the primitives of $T_i$ and with $\mathsf{Z} = \mathsf{Z}^{(i)}$; it is a complex bond function on the fine lattice of $T_i$, and at a fine bond $b$
\[
\mathscr{Z}(s')(b) = \sum_{b'} \sum_{\Upsilon, \Upsilon'} H_k(s')(b,b') \cdot C(b',\Upsilon) \cdot C_k^{\frac{1}{2}}(s')(\Upsilon,\Upsilon') \cdot \mathsf{Z}(\Upsilon'),
\]
the sum extending over the positively oriented unit bonds $b'$ and over the basis elements $\Upsilon$, $\Upsilon'$.

\begin{enumerate}
\item\label{it:renormalisation-flow-1} \emph{The weakened fluctuation field.} Let $b$ be a fine bond of $\mathbb{L}_k$ contained in $\tilde{\mathbf{Y}}^{\sim(1+d_U)}$. Then
\[
\mathscr{Z}^{(1)}\left(s'^{(1)}\right)\left(\mathrm{pr}_1(b)\right) = \mathscr{Z}^{(2)}\left(s'^{(2)}\right)\left(\mathrm{pr}_2(b)\right).
\]
Indeed, by~\eqref{eq:renormalisation-flow-a1} the box $N(\tilde{\mathbf{Y}})$ is contained in a fundamental domain of $l_1\mathbb{Z}^3$ and in a fundamental domain of $l_2\mathbb{Z}^3$; the sums that define the two sides are taken over these domains, every index of $T_i$ being identified with its lift in the domain. Lemma~\ref{lem:kernels-near-admissible} with $j = 1 + d_U$, applied to the family~\eqref{it:primitives-p1}, shows that the rows of $H_k^{(1)}\left(s'^{(1)}\right)$ and of $H_k^{(2)}\left(s'^{(2)}\right)$ at $b$ vanish at every column that is not contained in $\tilde{\mathbf{Y}}^{\sim(1+2d_U)}$, and that the two rows are equal at the columns contained in this set. By Proposition~\ref{prop:basis-map-box}, the kernel $C(b',\Upsilon)$ is given on both tori by the formula~\eqref{eq:operators-photon-16} on $\mathbb{Z}^3$, and it vanishes unless $\operatorname{supp}\Upsilon$ lies in the $L$-blocks of the end points of $b'$; by fact~\eqref{it:renormalisation-group-g0} of Section~\ref{sec:cubes-parameter-cubes} these $L$-blocks lie in the $M$-cubes of the end points. Hence only basis elements $\Upsilon$ contained in $\tilde{\mathbf{Y}}^{\sim(1+2d_U)}$ contribute. For such $\Upsilon$, Lemma~\ref{lem:kernels-near-admissible} with $j = 1 + 2d_U \le 5d_U$, applied to the family~\eqref{it:primitives-p2}, shows in the same way that the rows of $C_k^{\frac{1}{2},(1)}\left(s'^{(1)}\right)$ and of $C_k^{\frac{1}{2},(2)}\left(s'^{(2)}\right)$ at $\Upsilon$ vanish at every column not contained in $\tilde{\mathbf{Y}}^{\sim(1+3d_U)}$, a subset of $N(\tilde{\mathbf{Y}})$, and are equal at the columns contained in it. Since $\mathsf{Z}^{(1)}$ and $\mathsf{Z}^{(2)}$ have equal lifts on $N(\tilde{\mathbf{Y}})$, the two sums agree term by term.

By Step 0, $(A_1, \mathscr{C}_1) \in \tilde{R}^{-}_k(T_1) \times \mathfrak{X}_k(T_1)$, and $(A_2, \mathscr{C}_2) = \mathrm{pj}^{*}(A_1, \mathscr{C}_1)$; hence $(A_2, \mathscr{C}_2) \in \tilde{R}^{-}_k(T_2) \times \mathfrak{X}_k(T_2)$ by Lemma~\ref{lem:periodic-extension}. Now let $|t| \le \rho_1$. By Lemma~\ref{lem:bounds-analyticity}(q1), applied on $T_i$ with $s = s^{(i)}$, $\mathsf{Z} = \mathsf{Z}^{(i)}$ and $(A, \mathscr{C}) = (A_i, \mathscr{C}_i)$, the fields
\[
B_i := \mathscr{T}_i + t \cdot \mathscr{Z}^{(i)}\left(s'^{(i)}\right) \qquad \text{and} \qquad A_i + t \cdot \mathscr{Z}^{(i)}\left(s'^{(i)}\right)
\]
lie in $\tilde{G}_k(T_i)$ and in $\tilde{R}_k(T_i)$, respectively. As $\mathscr{T}_2 = \mathrm{pj}^{*}(\mathscr{T}_1)$ and $A_2 = \mathrm{pj}^{*}(A_1)$, the identity just proved shows that $B_1$ and $B_2$ have equal lifts on the bonds of $\tilde{\mathbf{Y}}^{\sim(1+d_U)}$, and so have $A_1 + t \cdot \mathscr{Z}^{(1)}\left(s'^{(1)}\right)$ and $A_2 + t \cdot \mathscr{Z}^{(2)}\left(s'^{(2)}\right)$.

\item\label{it:renormalisation-flow-2} \emph{The fermionic primitives.} A term $(X,P)$ of~\eqref{eq:renormalisation-group-57} is attributed to $\mathbf{Y}_i$ if $Y'(X,P) = \mathbf{Y}_i$. By Proposition~\ref{prop:local-constituents-step}(c), and in the sense of Definition~\ref{def:terminology}(1), the families~\eqref{it:primitives-p3} and $\Gamma_k$ enter the terms attributed to $\mathbf{Y}_i$ only at parameter families that vanish outside the parameter cubes contained in $\mathbf{Y}_i$, and there only through the rows of $H_k(s'',\cdot)$ and of $\bar{H}_k(s'',\cdot)$ at the fine sites at sup-distance at most $\eta_k$ from $\mathbf{Y}_i$ and through the rows of $\Gamma_k(s'',\cdot)$ at the unit sites of $\mathbf{Y}_i$; the family $M_k$ enters only through the operators $M_k(X,\cdot)$ with $X \subset \mathbf{Y}_i$; and the connections at which these operators are evaluated are $B = \mathscr{T} + t \cdot \mathscr{Z}(s')$, $|t| \le \rho_1$, and $\mathscr{T}$. All the row indices just named are contained in $\mathbf{Y}_i^{\sim 1}$, and the parameter families just named are the families $s^{(i)}$, $\boldsymbol{\sigma}$ arbitrary, because the parameter cubes contained in $\mathbf{Y}_i$ are the cubes $\mathrm{pr}_i(\tilde{c})$ (Lemma~\ref{lem:agreement-criterion-small}(c)). Lemma~\ref{lem:kernels-near-admissible} is applied with $j \le 1$; its hypothesis on the connections concerns the bonds of $\tilde{\mathbf{Y}}^{\sim(j+d_U)} \subset \tilde{\mathbf{Y}}^{\sim(1+d_U)}$, and it holds for $B_1$, $B_2$ by item~\ref{it:renormalisation-flow-1}, the value $t = 0$, at which $B_i = \mathscr{T}_i$, included. Hence these rows of the primitives of $T_1$ and of $T_2$ have their columns in $N(\tilde{\mathbf{Y}})$ and coincide there. For the rows of $\bar{H}_k$ one uses
\[
\bar{H}_k(s'',B) = \mathsf{C}^{-1} \cdot H_k(s'',-B) \cdot \mathsf{C} ;
\]
here $-B_i$ lies in $\tilde{G}_k(T_i)$ by Lemma~\ref{lem:invariances}, and $-B_1$, $-B_2$ have equal lifts on the bonds of $\tilde{\mathbf{Y}}^{\sim(1+d_U)}$. Finally, let $\tilde{X}$ be a union of $M$-cubes connected through faces and contained in $\tilde{\mathbf{Y}}$. The bonds that meet $\tilde{X}$ are bonds of $\tilde{\mathbf{Y}}^{\sim(1+d_U)}$; so the last assertion of Lemma~\ref{lem:kernels-near-admissible} shows that the kernels of $M_k^{(1)}\left(\mathrm{pr}_1(\tilde{X}), B_1\right)$ and of $M_k^{(2)}\left(\mathrm{pr}_2(\tilde{X}), B_2\right)$ coincide.

\item\label{it:renormalisation-flow-3} \emph{Exact operators and forms.} By Proposition~\ref{prop:local-constituents-step}(b), the exact operators and forms in the displays~\eqref{eq:renormalisation-group-56}--\eqref{eq:renormalisation-group-61} are: $C$; the operators $\delta_{\alpha,B}$ at the pairs admissible for polymers contained in $\mathbf{Y}_i$; the operators $T^{(1)}(\mathscr{T})$, $T^{(2)}(\mathscr{T})$, $\bar{T}^{(1)}(\mathscr{T})$, $\bar{T}^{(2)}(\mathscr{T})$ at the unit sites of $\mathbf{Y}_i$; the matrices $P^{\boldsymbol{\sigma}}_{\nu}$ and $\mathsf{C}$; and the forms $\langle \bar{\psi}, \psi \rangle_{\square}$ with $\square \subset \mathbf{Y}_i$. Each of them is given by one formula on $\mathbb{L}_k$ in the sense of Definition~\ref{def:terminology}(2); this formula is built from bonds, $L$-blocks, the blocks $B_k(y)$, the paths of the pairs of $\mathrm{Pz}_k$, membership in $\mathbf{Y}_i$, and the numbers $b_k$, $\mathfrak{m}$, $e_k$, $\eta_k$, $\alpha$, which are the same on both tori. These operators and forms read the connection on the bonds that meet $\mathbf{Y}_i$: for $T^{(1)}$, $T^{(2)}$, $\bar{T}^{(1)}$, $\bar{T}^{(2)}$ this is Lemma~\ref{lem:interior-boundary-parts}(b), and for $\delta_{\alpha,B}$ it is~\eqref{eq:renormalisation-group-adm}. By~\eqref{eq:renormalisation-flow-a2}, a point of $\mathrm{pr}_i\left(N(\tilde{\mathbf{Y}})\right)$ belongs to $\mathbf{Y}_i$ if and only if its lift belongs to $\tilde{\mathbf{Y}}$.

From these properties the following comparison of the two tori is obtained. Let $\tilde{X}$ be a union of $M$-cubes connected through faces and contained in $\tilde{\mathbf{Y}}$; by Lemma~\ref{lem:agreement-criterion-small}(c), applied to the admissible set $\tilde{\mathbf{Y}}$, the polymers of $T_i$ contained in $\mathbf{Y}_i$ are the sets $\mathrm{pr}_i(\tilde{X})$ of this kind. Let $|t|\le\rho_1$, and let $B_i$ be the field of item~\ref{it:renormalisation-flow-1}, the value $t=0$, at which $B_i=\mathscr{T}_i$, included. As in item~\ref{it:renormalisation-flow-1}, every index of $T_i$ whose support lies in $\mathrm{pr}_i\left(N(\tilde{\mathbf{Y}})\right)$ is identified with its lift in $N(\tilde{\mathbf{Y}})$. For an index $x$ contained in $N(\tilde{\mathbf{Y}})$, the row at $x$ of an operator on $T_1$ and the row at $x$ of an operator on $T_2$ are said to \emph{coincide} if the non-zero entries of both rows stand in columns contained in $N(\tilde{\mathbf{Y}})$ and if the two kernels take the same value at $(x,x')$ for every index $x'$ contained in $N(\tilde{\mathbf{Y}})$. Then the following four assertions hold.

First, for $i=1,2$ the projection $\mathrm{pr}_i$ maps the pairs $(x,x')$ of sites of $\mathbb{L}_k$ with $0<|x-x'|\le M$ that are admissible for $\tilde{X}$ bijectively onto the pairs of $\mathrm{Pz}_k$ on $T_i$ that are admissible for $\mathrm{pr}_i(\tilde{X})$, and a pair is of kind~(I), resp.\ of kind~(II), if and only if its image is; here the two kinds of Definition~\ref{def:pairs-sites-holder} are read in $\mathbb{L}_k$, that is, $x,x'\in\tilde{X}$ and each of the six rectilinear paths of $\mathbb{L}_k$ from $x$ to $x'$ is contained in $\tilde{X}$, resp.\ $x'$ is a nearest neighbour of $x$ in $\mathbb{L}_k$ and $x\in\tilde{X}$ or $x'\in\tilde{X}$. At each such pair the rows of $\delta_{\alpha,B_1}$ and of $\delta_{\alpha,B_2}$ coincide, and so do the rows of $\delta_{\alpha,-B_1}$ and of $\delta_{\alpha,-B_2}$.

Second, let $k\ge1$ and $j\in\{1,2\}$. At each unit site $y$ of $\tilde{X}$ the rows of $T^{(j)}(\mathscr{T}_1)$ and of $\bar{T}^{(j)}(\mathscr{T}_1)$, formed with the polymer $\mathrm{pr}_1(\tilde{X})$, coincide with the rows of $T^{(j)}(\mathscr{T}_2)$ and of $\bar{T}^{(j)}(\mathscr{T}_2)$, respectively, formed with $\mathrm{pr}_2(\tilde{X})$; the non-zero entries of these rows stand in columns that are sites and pairs contained in $\tilde{X}^{\sim1}$. Hence $\mathfrak{i}$ maps the elements that the substitution~\eqref{eq:renormalisation-group-58} formed with $\mathrm{pr}_1(\tilde{X})$ puts in place of $\Psi(y)$ and of $\bar{\Psi}(y)$ on $T_1$ to the elements that the substitution formed with $\mathrm{pr}_2(\tilde{X})$ puts in place of $\Psi(y)$ and of $\bar{\Psi}(y)$ on $T_2$.

Third, the entries $C(b',\Upsilon)$ are the same on both tori, by item~\ref{it:renormalisation-flow-1}, and $P^{\boldsymbol{\sigma}}_{\nu}$ and $\mathsf{C}$ are constant matrices.

Fourth, if $\tilde{X}$ is a single $M$-cube, then $\mathfrak{i}$ maps the form $\langle\bar{\psi},\psi\rangle_{\mathrm{pr}_1(\tilde{X})}$ to the form $\langle\bar{\psi},\psi\rangle_{\mathrm{pr}_2(\tilde{X})}$; every $M$-cube of $\tilde{\mathbf{Y}}$ is such a set $\tilde{X}$.

For the proof, let $z$ be a site, let $x_0$ be a site of an $M$-cube with the centre $c$, and let $|z-x_0|\le s$. The $M$-cube that contains $z$ (fact~\eqref{it:renormalisation-group-g0}) has a centre $c'\in c+M\mathbb{Z}^3$ with $|c'-c|<\frac{M}{2}+s+\frac{M}{2}$; hence $|c'-c|\le M$ if $s<M$, and $|c'-c|\le2M$ if $s<2M$. As $\eta_k\le1<M$, it follows that the sites at sup-distance less than $M$ from a site of $\tilde{X}$ lie in $\tilde{X}^{\sim1}$, and that those at sup-distance at most $M+\eta_k$ from a site of $\tilde{X}$ lie in $\tilde{\mathbf{Y}}^{\sim2}$.

A pair admissible for $\tilde{X}$ consists of two sites of $\tilde{X}$, or of a site of $\tilde{X}$ and a nearest neighbour of it, as in~\eqref{eq:renormalisation-group-adm}; so its sites lie in $\tilde{X}^{\sim1}$, hence in the box $N(\tilde{\mathbf{Y}})$, which contains the six rectilinear paths between any two of its points. The relations ``two sites are nearest neighbours'', ``two sites form a pair of $\mathrm{Pz}_k$'' and ``a bond lies on one of the six paths of such a pair'' have range less than $2d_U\cdot M$; so they hold between cells of $N(\tilde{\mathbf{Y}})$ if and only if they hold between the images of these cells under $\mathrm{pr}_i$ (Definition~\ref{def:admissibility-flatness-agreement}(a)). Moreover, a site of $N(\tilde{\mathbf{Y}})$ lies in $\tilde{X}$ if and only if its image lies in $\mathrm{pr}_i(\tilde{X})$, by~\eqref{eq:renormalisation-flow-a1}. Therefore $\mathrm{pr}_i$ maps each pair admissible for $\tilde{X}$ to a pair of $\mathrm{Pz}_k$ on $T_i$ that is admissible for $\mathrm{pr}_i(\tilde{X})$, of the same kinds, and it is injective on these pairs, by~\eqref{eq:renormalisation-flow-a1}. Conversely, a pair of $\mathrm{Pz}_k$ on $T_i$ that is admissible for the flat polymer $\mathrm{pr}_i(\tilde{X})$ has a lift $(x,x')$ with $x,x'\in\tilde{X}^{\sim1}$ (Definition~\ref{def:admissibility-flatness-agreement}(d)); by the relations just named, $0<|x-x'|\le M$, the six paths of the pair are the images of the six rectilinear paths from $x$ to $x'$, and $(x,x')$ is admissible for $\tilde{X}$, of the same kinds. This proves the bijection in the first assertion.

The rows to be compared stand at the following indices $x$: in the first assertion at the pairs admissible for $\tilde{X}$, in the second at the unit sites of $\tilde{X}$, and in the fourth, for the kernels of the two forms, at the sites of $\tilde{X}$. Every point of $\operatorname{supp}x$ has sup-distance at most $\eta_k$ from a site of $\tilde{X}$; so the sites at sup-distance at most $M$ from $\operatorname{supp}x$, and the bonds joining two such sites, lie in $\tilde{\mathbf{Y}}^{\sim2}$, which is contained in $\tilde{\mathbf{Y}}^{\sim(1+d_U)}$ and in $N(\tilde{\mathbf{Y}})$. There the data of the two tori have the same periodic extensions to $\mathbb{L}_k$: those of $B_1$ and of $B_2$ coincide on the bonds of $\tilde{\mathbf{Y}}^{\sim(1+d_U)}$, where these fields have equal lifts (item~\ref{it:renormalisation-flow-1}), and so do those of $-B_1$ and of $-B_2$; those of $\mathscr{T}_1$ and of $\mathscr{T}_2=\mathrm{pj}^{*}\mathscr{T}_1$ coincide on all bonds; and the periodic extensions $\tilde{X}+l_1\mathbb{Z}^3$ and $\tilde{X}+l_2\mathbb{Z}^3$ of $\mathrm{pr}_1(\tilde{X})$ and of $\mathrm{pr}_2(\tilde{X})$ both meet $N(\tilde{\mathbf{Y}})$ in $\tilde{X}$, by~\eqref{eq:renormalisation-flow-a1}, as $\tilde{X}\subset N(\tilde{\mathbf{Y}})$ and $l_2\mathbb{Z}^3\subset l_1\mathbb{Z}^3$. Now apply the comparison of two tori in Definition~\ref{def:terminology}(2), with the range $\rho=M$, which Proposition~\ref{prop:local-constituents-step}(b) provides for each of the formulas concerned, and with $U$ the set of the sites of $N(\tilde{\mathbf{Y}})$. The sides of the box $N(\tilde{\mathbf{Y}})$ are at most $l_1-2d_U\cdot M$ (Definition~\ref{def:admissibility-flatness-agreement}(a)), and no site lies on a face of an $M$-cube (fact~\eqref{it:renormalisation-group-g0}); hence
\[
|z-z'|+M<l_1-2d_U\cdot M+M<l_1\le l_i\qquad\left(z,z'\in U,\ i=1,2\right).
\]
So the rows at $x$ of the operators on $\mathbb{L}_k$ that correspond to the data of $T_1$ and to those of $T_2$ are equal, and, for $i=1,2$ and every index $x'$ contained in $N(\tilde{\mathbf{Y}})$, the kernel on $T_i$ at $(x,x')$ is the entry at $(x,x')$ of the operator on $\mathbb{L}_k$. By Proposition~\ref{prop:local-constituents-step}(b), the row of $\delta_{\alpha,\pm B}$ at a pair has its columns at the two points of the pair, the rows of $T^{(j)}$ and of $\bar{T}^{(j)}$ at a unit site $y$ have their columns at sup-distance less than $\frac{3}{2}$ from $y$, and the row at $x$ of the kernel of a form has no column other than $x$; all these columns are contained in $\tilde{X}^{\sim1}$. By the formula for the kernel in Definition~\ref{def:terminology}(2), the kernel on $T_i$ vanishes at every column of which no representative in $\mathbb{L}_k$ is a column where the row on $\mathbb{L}_k$ is non-zero. Hence the rows to be compared coincide.

In the second assertion, the columns of the rows of $T^{(1)}$ and of $\bar{T}^{(1)}$ are sites, and those of $T^{(2)}$ and of $\bar{T}^{(2)}$ are pairs, all contained in $\tilde{X}^{\sim1}$, as has just been shown. On $T_i$ the substitution~\eqref{eq:renormalisation-group-58} puts $\left(T^{(1)}(\mathscr{T}_i)\psi\right)(y)+\left(T^{(2)}(\mathscr{T}_i)\chi\right)(y)$ in place of $\Psi(y)$ and $\left(\bar{T}^{(1)}(\mathscr{T}_i)\bar{\psi}\right)(y)+\left(\bar{T}^{(2)}(\mathscr{T}_i)\bar{\chi}\right)(y)$ in place of $\bar{\Psi}(y)$, that is, linear combinations of the generators labelled by these columns, with the entries of the rows as coefficients. These coefficients are the same for $i=1$ and for $i=2$, and $\mathfrak{i}$ sends each generator of $T_1$ labelled by an index contained in $\tilde{X}^{\sim1}\subset N(\tilde{\mathbf{Y}})\subset N(\tilde{V})$ (Step 0) to the generator of $T_2$ with the same label; this proves the second assertion. In the fourth assertion, the kernel of $\langle\bar{\psi},\psi\rangle_{\mathrm{pr}_i(\tilde{X})}$ is $\eta_k^3$ times the unit matrix at the arguments $(x,x)$ with $x$ a site of $\mathrm{pr}_i(\tilde{X})$, and $0$ at all other arguments; so this form is the sum, over the sites $x$ of $\tilde{X}$ and over the spinor index $\boldsymbol{\beta}$, of $\eta_k^3\cdot\bar{\psi}_{\boldsymbol{\beta}}(x)\cdot\psi_{\boldsymbol{\beta}}(x)$, formed with the generators of $T_i$ labelled by $x$, and $\mathfrak{i}$ maps the sum on $T_1$ to the sum on $T_2$.

\item\label{it:renormalisation-flow-4} \emph{The datum.} By part~(a), the number $m' = m + m\left(\mathscr{E}^{(i)}\right)$ is the same for $i = 1$ and for $i = 2$. Let $\tilde{X}$ be a union of $M$-cubes connected through faces and contained in $\tilde{\mathbf{Y}}$, and let $|t| \le \rho_1$. The pairs
\[
\left(A_i + t \cdot \mathscr{Z}^{(i)}\left(s'^{(i)}\right), \mathscr{C}_i\right) \in \tilde{R}_k(T_i) \times \mathfrak{X}_k(T_i)
\]
have equal lifts on the bonds that meet $\tilde{X}$, by item~\ref{it:renormalisation-flow-1} and because $\mathscr{C}_2 = \mathrm{pj}^{*}(\mathscr{C}_1)$ (these bonds are bonds of $\tilde{\mathbf{Y}}^{\sim(1+d_U)}$). Let $\mathscr{E}^{r,(i)}$ denote the activity $\mathscr{E}^r$ of Section~\ref{sec:data-variables-fields} formed with $\mathscr{E}^{(i)}$. By Step 1, $\mathscr{E}^{r,(1)} \sim \mathscr{E}^{r,(2)}$; so Lemma~\ref{lem:agreement-criterion-small}(a) gives
\[
\mathfrak{i}\left(\mathscr{E}^{r,(1)}\left(\mathrm{pr}_1(\tilde{X}); A_1 + t \cdot \mathscr{Z}^{(1)}\left(s'^{(1)}\right), \mathscr{C}_1\right)\right) = \mathscr{E}^{r,(2)}\left(\mathrm{pr}_2(\tilde{X}); A_2 + t \cdot \mathscr{Z}^{(2)}\left(s'^{(2)}\right), \mathscr{C}_2\right),
\]
and the same identity holds at $t = 0$.

\item\label{it:renormalisation-flow-5} \emph{The functions $E^{(j),x}$, $j \le 3$.} Let $\tilde{X}$ be as in item~\ref{it:renormalisation-flow-4}, and let $j \in \{1, 2, 3\}$. Inserting items~\ref{it:renormalisation-flow-1}--\ref{it:renormalisation-flow-4} into the displays~\eqref{eq:renormalisation-group-56}, \eqref{eq:renormalisation-group-59} and~\eqref{eq:renormalisation-group-60}, one finds that, for every $t$ on the contours, $\mathfrak{i}$ maps the integrand of $G^{(j),(1)}\left(s^{(1)}, \mathrm{pr}_1(\tilde{X})\right)$ to the integrand of $G^{(j),(2)}\left(s^{(2)}, \mathrm{pr}_2(\tilde{X})\right)$. Hence
\[
\mathfrak{i}\left(G^{(j),(1)}\left(s^{(1)}, \mathrm{pr}_1(\tilde{X})\right)\right) = G^{(j),(2)}\left(s^{(2)}, \mathrm{pr}_2(\tilde{X})\right), \qquad j = 1, 2, 3,
\]
for all families $\boldsymbol{\sigma}$ and all $z$. Since this identity holds throughout the polydisk of the families $\boldsymbol{\sigma}$, the derivatives of both sides with respect to the variables $\boldsymbol{\sigma}_{\tilde{c}}$ coincide, and so do their integrals over $[0,1]^P$. By Lemma~\ref{lem:agreement-criterion-small}(c), the pairs $(X,P)$ with $Y'(X,P) = \mathbf{Y}_1$ correspond to the pairs with $Y'(X,P) = \mathbf{Y}_2$. The substitution~\eqref{eq:renormalisation-group-58} is formed with the operators $T^{(1)}(\mathscr{T})$, $T^{(2)}(\mathscr{T})$, $\bar{T}^{(1)}(\mathscr{T})$, $\bar{T}^{(2)}(\mathscr{T})$ at the unit sites of $\mathbf{Y}_i$, which are among the exact operators treated in item~\ref{it:renormalisation-flow-3}. Therefore, by~\eqref{eq:renormalisation-group-57} and~\eqref{eq:renormalisation-group-59},
\[
\mathfrak{i}\left(E^{(j),x,(1)}(\mathbf{Y}_1)\right) = E^{(j),x,(2)}(\mathbf{Y}_2), \qquad j \le 3,
\]
at $\left(A_i, \mathscr{C}_i, \mathsf{Z}^{(i)}\right)$. This holds for every union $\tilde{\mathbf{Y}}$ of $M$-cubes connected through faces and contained in $\tilde{V}$. By Corollary~\ref{cor:renormalisation-group-1}, for every polymer $\mathbf{Y}$ the function $E^{(j),x}(\mathbf{Y})$ depends on $\mathsf{Z}$ only through the coordinates $\mathsf{Z}(\Upsilon)$ with $\Upsilon$ contained in $\mathbf{Y}$.

\item\label{it:renormalisation-flow-6} \emph{The functions $E^{(\le 5),x}$.} By~\eqref{eq:renormalisation-group-59},
\[
G^{(\le 3)}(s'',X) = E^{(\le 3),x}\left(X; A, \mathscr{C}, \mathsf{Z}, \psi^{\#} + W_k^{\#}(s'',\mathscr{T})\right),
\]
and for $X \subset \mathbf{Y}_i$ the functions $G^{(\le 3)}$, $G^{(4)}$, $G^{(5)}$ are formed from $E^{(\le 3),x}(X)$, from $\mathscr{E}^r(X; A, \mathscr{C}, \cdot)$, from the number $m'$, from the forms $\langle \bar{\psi}, \psi \rangle_{\square}$, and from the field
\[
W_k^{\#}(s'',\mathscr{T}) = H_k^{\#}(s'',\mathscr{T}) \cdot \tilde{\Gamma}_k(s'',\mathscr{T}) \cdot W,
\]
of which only the components indexed by the sites $x \in X$ and by the pairs $(x,x')$ admissible for $X$ occur. Let $\tilde{X}$ be as in item~\ref{it:renormalisation-flow-4}. Item~\ref{it:renormalisation-flow-5}, applied with $\tilde{X}$ in place of $\tilde{\mathbf{Y}}$, gives
\[
\mathfrak{i}\left(E^{(\le 3),x,(1)}\left(\mathrm{pr}_1(\tilde{X})\right)\right) = E^{(\le 3),x,(2)}\left(\mathrm{pr}_2(\tilde{X})\right)
\]
at $\left(A_i, \mathscr{C}_i, \mathsf{Z}^{(i)}\right)$: by Corollary~\ref{cor:renormalisation-group-1} only the coordinates of $\mathsf{Z}^{(i)}$ at the basis elements contained in $\mathrm{pr}_i(\tilde{X})$ enter, and these coordinates are the same whether $\mathsf{Z}^{(i)}$ is formed with $\tilde{X}$ or with $\tilde{\mathbf{Y}}$. Together with items~\ref{it:renormalisation-flow-2}--\ref{it:renormalisation-flow-4} it follows that
\[
\mathfrak{i}\left(\left(G^{(\le 3)} + G^{(4)} + G^{(5)}\right)^{(1)}\left(s''^{(1)}, \mathrm{pr}_1(\tilde{X})\right)\right) = \left(G^{(\le 3)} + G^{(4)} + G^{(5)}\right)^{(2)}\left(s''^{(2)}, \mathrm{pr}_2(\tilde{X})\right).
\]
The two sides of this identity agree for every $\boldsymbol{\sigma}$ in the polydisk; hence they still agree after differentiation in the variables $\boldsymbol{\sigma}_{\tilde{c}}$ that belong to cubes $\tilde{c}$ of the kind $\pi_k$, and after integration of these derivatives over $[0,1]^P$. Lemma~\ref{lem:agreement-criterion-small}(c) matches the pairs $(X,P)$ of $T_1$ with $Y'(X,P) = \mathbf{Y}_1$ bijectively with the pairs $(X,P)$ of $T_2$ with $Y'(X,P) = \mathbf{Y}_2$. By~\eqref{eq:renormalisation-group-57},
\[
\mathfrak{i}\left(E^{(\le 5),x,(1)}(\mathbf{Y}_1)\right) = E^{(\le 5),x,(2)}(\mathbf{Y}_2)
\]
at $\left(A_i, \mathscr{C}_i, \mathsf{Z}^{(i)}\right)$, for all $z$.
\end{enumerate}

\emph{Step 4 (the functions $\mathsf{K}$, $\mathsf{H}$ and $E^{\#,x}_k$).} Let $\tilde{X}$ be a union of $M$-cubes which is connected through faces and contained in $\tilde{V}$, and let $\mathsf{Z}^{(i)}$ be formed as in Step 3 with $\tilde{\mathbf{Y}} = \tilde{X}$. By~\eqref{eq:renormalisation-group-61}, $\mathsf{K}(X)$ is a finite sum, over the sets of distinct sub-polymers of $X$ whose union is $X$, of products of the elements $\exp\left(E^{(\le 5),x}(X')\right) - 1$, $X'$ a member of the set. By Lemma~\ref{lem:agreement-criterion-small}(c), these sets for $X = \mathrm{pr}_1(\tilde{X})$ and for $X = \mathrm{pr}_2(\tilde{X})$ correspond to each other through the lifts of their members, which are unions of $M$-cubes connected through faces and contained in $\tilde{X}$. For each member, item~\ref{it:renormalisation-flow-6} of Step 3, applied to its lift, gives the identity of the two functions $E^{(\le 5),x}$ under $\mathfrak{i}$, and by Corollary~\ref{cor:renormalisation-group-1} this identity also holds at the field $\mathsf{Z}^{(i)}$ formed with $\tilde{X}$. Since $\mathfrak{i}$ is a homomorphism of algebras,
\[
\mathfrak{i}\left(\mathsf{K}^{(1)}\left(\mathrm{pr}_1(\tilde{X})\right)\right) = \mathsf{K}^{(2)}\left(\mathrm{pr}_2(\tilde{X})\right)
\]
at $\mathsf{Z}^{(i)}$. The element $\mathsf{K}(X)$ contains the generators $W(u)$ with $u \in X$ only, and it depends on $\mathsf{Z}$ only through the coordinates $\mathsf{Z}(\Upsilon)$ with $\Upsilon$ contained in $X$ (Corollary~\ref{cor:renormalisation-group-1}). By~\eqref{eq:renormalisation-group-61},
\[
\mathsf{H}(X) = \int \mathrm{d}\mu_I(W)\, \mathrm{d}\hat{\mu}(\mathsf{Z})\, \mathsf{K}(X).
\]
By Proposition~\ref{prop:basis-map-box}, under $\hat{\mu}$ the coordinates $\mathsf{Z}(\Upsilon)$, $\Upsilon$ in the basis, are independent and identically distributed, with a law on $\mathbb{R}$ that is determined by $p_{0,k}$ alone. As $p_{0,k}$ is the same on $T_1$ and on $T_2$ (Section~\ref{sec:admissible-sets-local}), on both tori the integration over $\mathsf{Z}$ amounts to the integration over $z \in [-p_{0,k}, p_{0,k}]^{\mathsf{B}}$ against one and the same product measure. For a monomial in the generators $W(u)$, $u \in X$, the rule~\eqref{eq:polymer-activities-m} of Section~\ref{sec:norms-gaussian-integrals} gives the $\mu_I$-integral $1$ when the monomial is the product of elements $W_{\boldsymbol{\beta}}(y) \cdot \bar{W}_{\boldsymbol{\beta}}(y)$ with pairwise different indices $(y,\boldsymbol{\beta})$, and the $\mu_I$-integral $0$ when no permutation of its factors brings it into this shape. As $\mathfrak{i}$ maps each generator to the generator with the same index, it follows that
\[
\mathfrak{i}\left(\int \mathrm{d}\mu_I(W)\, F\right) = \int \mathrm{d}\mu_I(W)\, \mathfrak{i}(F) \qquad \text{for every } F \in \mathrm{Gr}_1(\tilde{V}),
\]
and therefore
\[
\mathfrak{i}\left(\mathsf{H}^{(1)}\left(\mathrm{pr}_1(\tilde{X})\right)\right) = \mathsf{H}^{(2)}\left(\mathrm{pr}_2(\tilde{X})\right).
\]
By~\eqref{eq:renormalisation-group-61},
\[
E^{\#,x}_k(X) = \sum_{n \ge 1} \frac{1}{n!} \sum_{(X_1,\dots,X_n)} \rho^T(X_1,\dots,X_n) \cdot \prod_{a=1}^n \mathsf{H}(X_a),
\]
where the inner sum extends over the $n$-tuples of polymers with union $X$, and $\rho^T$ is the coefficient of Proposition~\ref{prop:polymers}. On both tori this series converges in norm, by clause~\eqref{it:step-small-fields-a} of Theorem~\ref{thm:step-small-fields}, applied on $T_i$ at $(A_i, \mathscr{C}_i) \in \tilde{R}^{-}_k(T_i) \times \mathfrak{X}_k(T_i)$ (item~\ref{it:renormalisation-flow-1}). The number $\rho^T(X_1,\dots,X_n)$ is determined by the set of the pairs $\{a, a'\}$ for which $X_a$ and $X_{a'}$ are incompatible. By Lemma~\ref{lem:agreement-criterion-small}(c), for $i = 1$ and for $i = 2$ the map
\[
(\tilde{X}_1,\dots,\tilde{X}_n) \mapsto \left(\mathrm{pr}_i(\tilde{X}_1),\dots,\mathrm{pr}_i(\tilde{X}_n)\right)
\]
is a bijection from the $n$-tuples of unions of $M$-cubes connected through faces and contained in $\tilde{X}$ with $\tilde{X}_1 \cup \dots \cup \tilde{X}_n = \tilde{X}$ onto the $n$-tuples of polymers of $T_i$ with union $\mathrm{pr}_i(\tilde{X})$, and it preserves incompatibility; so the two images of a tuple $(\tilde{X}_1,\dots,\tilde{X}_n)$ carry the same coefficient $\rho^T$. The identity for $\mathsf{H}$ just obtained holds with each $\tilde{X}_a$ in place of $\tilde{X}$, and $\mathfrak{i}$ is multiplicative; hence $\mathfrak{i}$ maps the term of $\left(\mathrm{pr}_1(\tilde{X}_1),\dots,\mathrm{pr}_1(\tilde{X}_n)\right)$ in the series on $T_1$ to the term of $\left(\mathrm{pr}_2(\tilde{X}_1),\dots,\mathrm{pr}_2(\tilde{X}_n)\right)$ in the series on $T_2$. Finally, $\mathfrak{i}$ is a linear map between spaces of finite dimension, hence continuous, and may therefore be applied to the norm-convergent series term by term. Hence
\[
\mathfrak{i}\left(E^{\#,x,(1)}_k\left(\mathrm{pr}_1(\tilde{X}); A_1, \mathscr{C}_1\right)\right) = E^{\#,x,(2)}_k\left(\mathrm{pr}_2(\tilde{X}); A_2, \mathscr{C}_2\right).
\]
Together with Steps 1 and 2, each of the three families of level $k$ named in Step 0 satisfies the identity required there. This proves (b).

\emph{Part (c).} Let $\tilde{U}$, $\tilde{V}$ and $Y_1$, $Y_2$ be as in Step 0, and let $\mathscr{C}$ be an element of $\mathfrak{X}_{k+1}$ on the torus $T_{k+1}(R)$. By Lemma~\ref{lem:rescaling}, $\mathscr{C}_L$ lies in $\mathfrak{X}_k(T_1)$. Moreover $\left(\mathrm{pj}^{*}(\mathscr{C})\right)_L = \mathrm{pj}^{*}(\mathscr{C}_L)$: the scaling $\mathscr{C} \mapsto \mathscr{C}_L$ of Lemma~\ref{lem:rescaling} is given by the same formula on the periodic bond functions of both tori, and $\mathrm{pj}^{*}$ only regards a bond function on $T_{k+1}(R)$, resp.\ on $T_1$, as a periodic bond function on $T_{k+1}(R_{\mathrm{top}})$, resp.\ on $T_2$ (Definition~\ref{def:periodisation-projection}). By the rule of clause~\eqref{it:step-small-fields-c} of Theorem~\ref{thm:step-small-fields} that forms $\mathfrak{P}_{k+1}$,
\[
\mathfrak{P}_{k+1}^{(1)}(Y_1, \mathscr{C}) = \sum_{X \colon \bar{X}^L = \mathrm{pr}_1(\tilde{V})} \left( \mathfrak{P}_k^{(1)}(X, \mathscr{C}_L) + E_k^{\mathrm{det},(1)}(X, \mathscr{C}_L) \right),
\]
and $\mathfrak{P}_{k+1}^{(2)}\left(Y_2, \mathrm{pj}^{*}(\mathscr{C})\right)$ is the corresponding sum on $T_2$, with $\mathrm{pr}_2(\tilde{V})$ and $\mathrm{pj}^{*}(\mathscr{C}_L)$ in place of $\mathrm{pr}_1(\tilde{V})$ and $\mathscr{C}_L$. By Lemma~\ref{lem:agreement-criterion-small}(c), the polymers of the two sums correspond to each other through their lifts, which are unions of $M$-cubes connected through faces, admissible, and contained in $\tilde{V}$. For corresponding polymers the values of $\mathfrak{P}_k^{(1)}$ and $\mathfrak{P}_k^{(2)}$ coincide by the hypothesis $\mathfrak{P}_k^{(1)} \sim \mathfrak{P}_k^{(2)}$, and the values of $E_k^{\mathrm{det},(1)}$ and $E_k^{\mathrm{det},(2)}$ coincide by Step 2. Hence $\mathfrak{P}_{k+1}^{(2)}\left(Y_2, \mathrm{pj}^{*}(\mathscr{C})\right) = \mathfrak{P}_{k+1}^{(1)}(Y_1, \mathscr{C})$, that is, $\mathfrak{P}_{k+1}^{(1)} \sim \mathfrak{P}_{k+1}^{(2)}$, which is (c).

By Theorem~\ref{thm:properties-projection}, Lemmas~\ref{lem:locality-rows} and~\ref{lem:kernels-near-admissible}, and parts (a)--(c), Theorem~\ref{thm:projection-agreement-flat} follows.
\end{proof}

\subsection{The flow on a fixed torus}\label{sec:flow-fixed-torus}

Throughout this subsection hypothesis~\ref{def:hypothesis-flow} is assumed; every object below is formed on the torus $R$, with the primitives of Definition~\ref{def:periodisation-projection}. By Proposition~\ref{prop:data-along-flow}(b), the data of every level $k < k^{*}$ satisfy hypothesis~\ref{def:hypothesis-one-step}. Therefore Theorem~\ref{thm:step-small-fields} holds at every level $k < k^{*}$, and so do Propositions~\ref{prop:symmetries-extraction} and~\ref{prop:contraction} (Theorem~\ref{thm:renormalisation-group-1}) and Proposition~\ref{prop:local-expansion-determinant}.

Condition~\eqref{eq:renormalisation-group-fl} holds by Proposition~\ref{prop:consequences}. The terms of its sum are positive; hence
\[
c_{\mathrm{fl}}\, e_k^{\varepsilon} < \mathfrak{s}_k = \left( c_{\mathrm{fl}} + 2 \right) e_k^{\varepsilon} \le \frac{1}{16} \qquad (0 \le k \le k^{*}),
\]
and in particular $e_k < 1$ for $0 \le k \le k^{*}$.

Since $e_{k+1} = L^{\frac{1}{2}} e_k$ for every $k$ (Section~\ref{sec:running-quantities-tori}), the numbers $w^{\mathrm{m}}_k = e_k^{\frac{3}{4} - 8\varepsilon}$ and $w^{\mathrm{E}}_k = e_k^{\frac{1}{4} - 7\varepsilon}$ satisfy
\[
w^{\mathrm{m}}_{k+1} = L^{\frac{3}{8} - 4\varepsilon} \cdot w^{\mathrm{m}}_k, \qquad w^{\mathrm{E}}_{k+1} = L^{\frac{1}{8} - \frac{7}{2}\varepsilon} \cdot w^{\mathrm{E}}_k.
\]
Stage~\ref{it:order-choice-4} contains the inequalities $L^{-\frac{5}{8}} \le \frac{1}{4}$ and $c_{\mathrm{fl}}\, L^{-\frac{3}{8} + 6\bar{\varepsilon}} \le \frac{1}{16}$, and $\varepsilon \le \bar{\varepsilon}$ by stage~\ref{it:order-choice-8}. Consequently, for $0 \le k < k^{*}$,
\begin{align}
L^{-1} \cdot \frac{w^{\mathrm{m}}_{k+1}}{w^{\mathrm{m}}_k} &= L^{-\frac{5}{8} - 4\varepsilon} \le \frac{1}{4}, \label{eq:renormalisation-flow-7}\\
c_{\mathrm{fl}}\, e_k^{\frac{1}{2}} \cdot \frac{w^{\mathrm{E}}_k}{w^{\mathrm{m}}_k} &= c_{\mathrm{fl}}\, e_k^{\varepsilon}, \label{eq:renormalisation-flow-8}\\
c_{\mathrm{fl}}\, L^{-\frac{1}{4} + 2\varepsilon} \cdot \frac{w^{\mathrm{E}}_k}{w^{\mathrm{E}}_{k+1}} &= c_{\mathrm{fl}}\, L^{-\frac{3}{8} + \frac{11}{2}\varepsilon} \le \frac{1}{16}, \label{eq:renormalisation-flow-9}\\
\frac{e_{k+1}^{\frac{1}{4} - 2\varepsilon}}{w^{\mathrm{E}}_{k+1}} = e_{k+1}^{5\varepsilon} \le e_{k+1}^{\varepsilon}, \qquad \frac{e_k^{\frac{1}{4} - 6\varepsilon}}{w^{\mathrm{E}}_{k+1}} &= \left( \frac{e_k}{e_{k+1}} \right)^{\frac{1}{4} - 6\varepsilon} e_{k+1}^{\varepsilon} \le e_{k+1}^{\varepsilon}. \label{eq:renormalisation-flow-10}
\end{align}
The inequality in~\eqref{eq:renormalisation-flow-9} uses $\frac{11}{2}\varepsilon \le 6\bar{\varepsilon}$ and $L > 1$; the two inequalities in~\eqref{eq:renormalisation-flow-10} use $e_k < e_{k+1} < 1$ and $\frac{1}{4} - 6\varepsilon > 0$.

\begin{proposition}[Existence of the trajectory]\label{prop:existence-trajectory}
Let $(m_{\mathrm{in}}, \boldsymbol{\varepsilon}_{\mathrm{in}})$ be a real pair and let $0 \le k \le k^{*}$. Assume that, for $i = 0, 1, \dots, k-1$ in turn, the number $m_i$ is defined and satisfies $|m_i| \le w^{\mathrm{m}}_i$. Then the trajectory of the pair (Section~\ref{sec:flow}) exists up to level $k$, so that the data $\boldsymbol{\varepsilon}^t_k$, $m_k$, $E_k$, $\mathscr{E}_k$, $\mathfrak{P}_k$ of level $k$ are defined. Moreover, $\boldsymbol{\varepsilon}^t_k$ and $m_k$ are real, $\mathscr{E}_k \in K^x_k$, $E_k \in K_k$, $\mathscr{E}_k|_0 = E_k$, $\mathfrak{P}_k(X,0) = 0$ for every $X \in D_k$, and
\[
\|E_k\|_k \le \|\mathscr{E}_k\|^x_k \le 4\, e_k^{\varepsilon} \cdot w^{\mathrm{E}}_k \le w^{\mathrm{E}}_k.
\]
In particular, for a real pair the following are equivalent: the pair is controlled on $R$ (Definition~\ref{def:fluctuation-field-controlled}); for $k = 0, 1, \dots, k^{*}$ in turn, the number $m_k$ is defined and satisfies $|m_k| \le w^{\mathrm{m}}_k$.
\end{proposition}

\begin{proof}
The proof is by induction on $k$; besides the assertions, the induction carries the stronger bound
\[
\mathrm{Ax}_k := \frac{\|\mathscr{E}_k\|^x_k}{w^{\mathrm{E}}_k} \le \frac{32}{15}\, e_k^{\varepsilon}.
\]
The last inequality of the chain in the statement does not depend on the induction: for $0 \le k \le k^{*}$ one has $e_k \le e_{*} \le e_0$, where $e_0$ is the ceiling of the charge fixed in stage~\ref{it:order-choice-9} (Proposition~\ref{prop:data-along-flow}(c)), and therefore $4\, e_k^{\varepsilon} \le 1$ by condition~\eqref{eq:renormalisation-group-A8}.

For $k = 0$ the recursion of Section~\ref{sec:flow} starts from $\boldsymbol{\varepsilon}^t_0 = \boldsymbol{\varepsilon}_{\mathrm{in}}$, $m_0 = m_{\mathrm{in}}$, $E_0 = 0$, $\mathscr{E}_0 = 0$ and $\mathfrak{P}_0 = 0$. All assertions hold, and $\mathrm{Ax}_0 = 0$.

Let $0 \le k < k^{*}$, and assume that, for $i = 0, 1, \dots, k$ in turn, the number $m_i$ is defined and $|m_i| \le w^{\mathrm{m}}_i$. By the induction hypothesis the trajectory exists up to level $k$, the assertions hold at level $k$, and $\mathrm{Ax}_k \le \frac{32}{15}\, e_k^{\varepsilon}$. In particular $\|\mathscr{E}_k\|^x_k \le w^{\mathrm{E}}_k$; together with $|m_k| \le w^{\mathrm{m}}_k$ this gives $(m_k, \mathscr{E}_k) \in \mathrm{Ball}^x_k(1) \subset \mathrm{Ball}^x_k(4)$, so that the trajectory exists up to level $k+1$. The data $\boldsymbol{\varepsilon}^t_k$, $m_k$, $\mathscr{E}_k$ satisfy the hypotheses of clause~\eqref{it:step-small-fields-c} of Theorem~\ref{thm:step-small-fields} at level $k$ (for $k = 0$ note that $\mathscr{E}_0 = 0$). Since $\mathscr{E}_k|_0 = E_k$, Proposition~\ref{prop:symmetries-extraction} gives
\[
\varepsilon(\mathscr{E}_k) = \varepsilon(\mathscr{E}_k|_0) = \varepsilon(E_k), \qquad m(\mathscr{E}_k) = m(\mathscr{E}_k|_0) = m(E_k).
\]
Hence the recursion of Section~\ref{sec:flow} takes the form
\begin{align*}
\boldsymbol{\varepsilon}^t_{k+1} &= L^3 \left( \boldsymbol{\varepsilon}^t_k + \varepsilon(\mathscr{E}_k) + d\varepsilon^0_k \right), \qquad m_{k+1} = L \left( m_k + m(\mathscr{E}_k) \right),\\
\mathscr{E}_{k+1} &= L R\,\mathscr{E}_k + L\,\delta E^{\mathrm{det}}_k + E^{*,x}_k(m_k, \mathscr{E}_k);
\end{align*}
that is, $(\boldsymbol{\varepsilon}^t_{k+1}, m_{k+1}, \mathscr{E}_{k+1})$ is the triple $(\boldsymbol{\varepsilon}^{+}, m^{+}, \mathscr{E}^{+})$ that clause~\eqref{it:step-small-fields-c} of Theorem~\ref{thm:step-small-fields} forms from $(\boldsymbol{\varepsilon}^t_k, m_k, \mathscr{E}_k)$. The numbers $\varepsilon(E_k)$ and $m(E_k)$ are real by Proposition~\ref{prop:symmetries-extraction}, and $d\varepsilon^0_k$ is real by clause~\eqref{it:step-small-fields-a} of Theorem~\ref{thm:step-small-fields}; hence $\boldsymbol{\varepsilon}^t_{k+1}$ and $m_{k+1}$ are real.

By clause~\eqref{it:step-small-fields-d} of Theorem~\ref{thm:step-small-fields}, $\mathscr{E}_{k+1} \in K^x_{k+1}$,
\[
\mathscr{E}_{k+1}|_0 = L R\,E_k + L\,E^{\mathrm{det}}_k + E^{*}_k(m_k, E_k) = E_{k+1},
\]
the second equality being the recursion for $E_{k+1}$, and
\[
\|\mathscr{E}_{k+1}\|^x_{k+1} \le c_{\mathrm{fl}}\, L^{-\frac{1}{4} + 2\varepsilon} \|\mathscr{E}_k\|^x_k + e_{k+1}^{\frac{1}{4} - 2\varepsilon} + e_k^{\frac{1}{4} - 6\varepsilon}.
\]
Since $\mathscr{E}_{k+1} \in K^x_{k+1}$, its restriction $E_{k+1} = \mathscr{E}_{k+1}|_0$ belongs to $K_{k+1}$ and satisfies $\|E_{k+1}\|_{k+1} \le \|\mathscr{E}_{k+1}\|^x_{k+1}$ (Section~\ref{sec:spaces-activities}). Dividing the last bound by $w^{\mathrm{E}}_{k+1}$ and using~\eqref{eq:renormalisation-flow-9}, \eqref{eq:renormalisation-flow-10}, the bound on $\mathrm{Ax}_k$ and $e_k \le e_{k+1}$, one obtains
\[
\mathrm{Ax}_{k+1} \le \frac{1}{16}\, \mathrm{Ax}_k + 2\, e_{k+1}^{\varepsilon} \le \left( \frac{1}{16} \cdot \frac{32}{15} + 2 \right) e_{k+1}^{\varepsilon} = \frac{32}{15}\, e_{k+1}^{\varepsilon}.
\]
As $\frac{32}{15} < 4$, the chain of inequalities in the statement holds at level $k+1$. Finally, $E^{\mathrm{det}}_k(X,0) = 0$ for every $X \in D_k$ by Definition~\ref{def:determinant-functions}, and $\mathfrak{P}_k(X,0) = 0$. The recursion of Section~\ref{sec:flow} expresses $\mathfrak{P}_{k+1}(Y,0)$, for $Y \in D_{k+1}$, as a sum of numbers $\mathfrak{P}_k(X,0) + E^{\mathrm{det}}_k(X,0)$ with $X \in D_k$; hence $\mathfrak{P}_{k+1}(Y,0) = 0$. This completes the induction.

It remains to prove the last assertion. Definition~\ref{def:fluctuation-field-controlled} requires of a controlled pair exactly two things: its trajectory exists up to level $k^{*}$, and $|m_k| \le w^{\mathrm{m}}_k$ for $0 \le k \le k^{*}$. If these bounds hold for $k = 0, 1, \dots, k^{*}$ in turn, then the case $k = k^{*}$ of what has been proved shows that the trajectory exists up to level $k^{*}$; the converse implication is contained in the definition. Theorem~\ref{thm:trajectories-fixed-torus}(a) follows.
\end{proof}

\begin{proposition}[The pair with vanishing terminal values]\label{prop:pair-vanishing-terminal}
There is exactly one real pair $(m_{\mathrm{in}}, \boldsymbol{\varepsilon}_{\mathrm{in}})$, controlled on $R$, whose trajectory satisfies $m_{k^{*}} = 0$ and $\boldsymbol{\varepsilon}^t_{k^{*}} = 0$. For $0 \le k \le k^{*}$ its trajectory satisfies
\[
|m_k| \le \frac{w^{\mathrm{m}}_k}{8}, \qquad \boldsymbol{\varepsilon}^t_k = -\sum_{i=k}^{k^{*}-1} L^{-3(i-k)} \left( \varepsilon(E_i) + d\varepsilon^0_i \right), \qquad |\boldsymbol{\varepsilon}^t_k| \le 5\, c_{\mathrm{fl}}\, e_k^{\varepsilon} \cdot w^{\mathrm{E}}_k.
\]
\end{proposition}

\begin{proof}
\emph{A complete metric space.} Let $0 \le k \le k^{*}$. The space $K_k = K_k(l_k)$ of Definition~\ref{def:renormalisation-group-1} is a real Banach space. Indeed, the set $D_k$ of polymers of the torus is finite. For each $X \in D_k$, the bounded analytic maps from the set $\tilde{R}_k$, which is open (Definition~\ref{def:domains}, Lemma~\ref{lem:openness}), into the finite-dimensional space $\mathrm{Gr}[\psi,\chi]$ form a Banach space under the norm $F \mapsto \sup_{A \in \tilde{R}_k} \|F(A)\|_{h_k}$, since analyticity survives uniform convergence; as $D_k$ is finite, the norm $\|\cdot\|_k$ is equivalent to the maximum over $X \in D_k$ of these norms. The requirements on the form of the elements and the conditions~\eqref{it:renormalisation-group-1-k2}--\eqref{it:renormalisation-group-1-k5} and~\eqref{it:renormalisation-group-1-k7} are closed and real-linear.

Let $\mathsf{B}$ be the set of the sequences $\xi = (m_k, E_k)_{0 \le k \le k^{*}}$ with $E_0 = 0$ and
\[
m_k \in \mathbb{R}, \qquad |m_k| \le w^{\mathrm{m}}_k, \qquad E_k \in K_k, \qquad \|E_k\|_k \le w^{\mathrm{E}}_k \qquad (0 \le k \le k^{*}).
\]
For $\xi \in \mathsf{B}$ and $\tilde{\xi} = (\tilde{m}_k, \tilde{E}_k)_{0 \le k \le k^{*}} \in \mathsf{B}$ put
\[
a_k := \frac{\|E_k - \tilde{E}_k\|_k}{w^{\mathrm{E}}_k}, \qquad b_k := \frac{|m_k - \tilde{m}_k|}{w^{\mathrm{m}}_k}, \qquad \operatorname{dist}(\xi, \tilde{\xi}) := \max_{0 \le k \le k^{*}} \max\left( a_k, b_k \right).
\]
The function $\operatorname{dist}$ is the distance of the norm $\xi \mapsto \max_k \max\left( \|E_k\|_k/w^{\mathrm{E}}_k,\, |m_k|/w^{\mathrm{m}}_k \right)$ on the Banach space $\prod_{k=0}^{k^{*}} \left( \mathbb{R} \times K_k \right)$, and $\mathsf{B}$ is a closed subset of this space. Hence $(\mathsf{B}, \operatorname{dist})$ is a complete metric space.

\emph{The map.} For $\xi \in \mathsf{B}$ let $\mathrm{Fx}(\xi) := (m'_k, E'_k)_{0 \le k \le k^{*}}$, where
\begin{align*}
m'_k &:= L^{-1} m_{k+1} - m(E_k) \quad (0 \le k < k^{*}), & m'_{k^{*}} &:= 0,\\
E'_k &:= L R\,E_{k-1} + L\,E^{\mathrm{det}}_{k-1} + E^{*}_{k-1}(m_{k-1}, E_{k-1}) \quad (1 \le k \le k^{*}), & E'_0 &:= 0.
\end{align*}
The right sides are defined, since $(m_{k-1}, E_{k-1}) \in \mathrm{Ball}_{k-1}(1)$ for $1 \le k \le k^{*}$. An element $\xi$ of $\mathsf{B}$ is a fixed point of $\mathrm{Fx}$ if and only if $m_{k^{*}} = 0$ and
\[
m_{k+1} = L \left( m_k + m(E_k) \right), \qquad E_{k+1} = L R\,E_k + L\,E^{\mathrm{det}}_k + E^{*}_k(m_k, E_k) \qquad (0 \le k < k^{*});
\]
recall that $E_0 = 0$ for every element of $\mathsf{B}$.

\emph{$\mathrm{Fx}$ maps $\mathsf{B}$ into itself.} Let $\xi \in \mathsf{B}$. For $0 \le k < k^{*}$, Proposition~\ref{prop:symmetries-extraction} shows that $m(E_k)$ is real and that $|m(E_k)| \le c_{\mathrm{fl}}\, e_k^{\frac{1}{2}} \|E_k\|_k$; so $m'_k$ is real and, by~\eqref{eq:renormalisation-flow-7} and~\eqref{eq:renormalisation-flow-8},
\[
\frac{|m'_k|}{w^{\mathrm{m}}_k} \le L^{-1} \cdot \frac{w^{\mathrm{m}}_{k+1}}{w^{\mathrm{m}}_k} + c_{\mathrm{fl}}\, e_k^{\frac{1}{2}} \cdot \frac{w^{\mathrm{E}}_k}{w^{\mathrm{m}}_k} \le \frac{1}{4} + c_{\mathrm{fl}}\, e_k^{\varepsilon} \le 1.
\]
For $1 \le k \le k^{*}$, the three terms of $E'_k$ belong to $K_k$ by Proposition~\ref{prop:contraction}, Proposition~\ref{prop:local-expansion-determinant}(d5) and clause~\eqref{it:step-small-fields-b} of Theorem~\ref{thm:step-small-fields}, applied at level $k-1$. The same statements, with Proposition~\ref{prop:local-expansion-determinant}(d6) in place of~(d5), give
\[
\|E'_k\|_k \le c_{\mathrm{fl}}\, L^{-\frac{1}{4} + 2\varepsilon} \cdot w^{\mathrm{E}}_{k-1} + e_k^{\frac{1}{4} - 2\varepsilon} + e_{k-1}^{\frac{1}{4} - 6\varepsilon},
\]
and~\eqref{eq:renormalisation-flow-9}, \eqref{eq:renormalisation-flow-10} at level $k-1$ turn this into $\|E'_k\|_k/w^{\mathrm{E}}_k \le \frac{1}{16} + 2\, e_k^{\varepsilon} \le 1$. Thus $\mathrm{Fx}(\xi) \in \mathsf{B}$.

\emph{$\mathrm{Fx}$ is a contraction.} Let $\xi, \tilde{\xi} \in \mathsf{B}$, define $a_k$, $b_k$ as above, and let $a'_k$, $b'_k$ be the same quantities formed with $\mathrm{Fx}(\xi)$ and $\mathrm{Fx}(\tilde{\xi})$. Since $m(\cdot)$ is linear, Proposition~\ref{prop:symmetries-extraction} together with~\eqref{eq:renormalisation-flow-7} and~\eqref{eq:renormalisation-flow-8} gives
\[
b'_k \le \frac{1}{4}\, b_{k+1} + c_{\mathrm{fl}}\, e_k^{\varepsilon}\, a_k \quad (0 \le k < k^{*}), \qquad b'_{k^{*}} = 0.
\]
In $E'_k - \tilde{E}'_k$ the terms $L\,E^{\mathrm{det}}_{k-1}$ cancel, and $L R$ is linear; Proposition~\ref{prop:contraction}, clause~\eqref{it:step-small-fields-e} of Theorem~\ref{thm:step-small-fields} (applied to the pairs $(m_{k-1}, E_{k-1})$ and $(\tilde{m}_{k-1}, \tilde{E}_{k-1})$, which lie in $\mathrm{Ball}_{k-1}(1)$), \eqref{eq:renormalisation-flow-9} and~\eqref{eq:renormalisation-flow-10} give
\[
a'_k \le \frac{1}{16}\, a_{k-1} + e_k^{\varepsilon} \left( a_{k-1} + b_{k-1} \right) \quad (1 \le k \le k^{*}), \qquad a'_0 = 0.
\]
As $c_{\mathrm{fl}}\, e_k^{\varepsilon} \le \frac{1}{16}$ and $2\, e_k^{\varepsilon} \le \frac{1}{4}$, it follows that
\[
\operatorname{dist}\left( \mathrm{Fx}(\xi), \mathrm{Fx}(\tilde{\xi}) \right) \le \left( \frac{1}{4} + \frac{1}{16} \right) \operatorname{dist}(\xi, \tilde{\xi}).
\]
By Banach's fixed point theorem, $\mathrm{Fx}$ has exactly one fixed point $(m_k, E_k)_{0 \le k \le k^{*}}$ in $\mathsf{B}$.

\emph{Existence.} Let $(m_k, E_k)_{0 \le k \le k^{*}}$ be the fixed point of $\mathrm{Fx}$, and put
\[
m_{\mathrm{in}} := m_0, \qquad \boldsymbol{\varepsilon}_{\mathrm{in}} := -\sum_{i=0}^{k^{*}-1} L^{-3i} \left( \varepsilon(E_i) + d\varepsilon^0_i \right).
\]
This is a real pair, because $\varepsilon(E_i)$ is real (Proposition~\ref{prop:symmetries-extraction}) and $0 \le d\varepsilon^0_i \le e_i^7$ (clause~\eqref{it:step-small-fields-a} of Theorem~\ref{thm:step-small-fields}). For $0 \le k \le k^{*}$, the trajectory of this pair exists up to level $k$, and its numbers $m$ and activities $E$ of the levels $0, \dots, k$ are $m_0, \dots, m_k$ and $E_0, \dots, E_k$. For $k = 0$ this is the initial condition of the recursion of Section~\ref{sec:flow}. If it holds for some $k < k^{*}$, then $|m_i| \le w^{\mathrm{m}}_i$ for $0 \le i \le k$; so the trajectory exists up to level $k+1$ by Proposition~\ref{prop:existence-trajectory}, and by the recursion and the equations of a fixed point its terms of level $k+1$ are $L(m_k + m(E_k)) = m_{k+1}$ and $L R\,E_k + L\,E^{\mathrm{det}}_k + E^{*}_k(m_k, E_k) = E_{k+1}$. Hence the pair is controlled on $R$ (Definition~\ref{def:fluctuation-field-controlled}), and $m_{k^{*}} = 0$.

The relation for $\boldsymbol{\varepsilon}^t$ in the recursion of Section~\ref{sec:flow} is equivalent to $\boldsymbol{\varepsilon}^t_k = L^{-3} \boldsymbol{\varepsilon}^t_{k+1} - \varepsilon(E_k) - d\varepsilon^0_k$. By induction on $k$, every real pair whose trajectory exists up to level $k^{*}$ satisfies
\[
\boldsymbol{\varepsilon}^t_0 = L^{-3k}\, \boldsymbol{\varepsilon}^t_k - \sum_{i=0}^{k-1} L^{-3i} \left( \varepsilon(E_i) + d\varepsilon^0_i \right) \qquad (0 \le k \le k^{*}).
\]
For the pair constructed above one has $\boldsymbol{\varepsilon}^t_0 = \boldsymbol{\varepsilon}_{\mathrm{in}}$; at $k = k^{*}$ the identity gives $\boldsymbol{\varepsilon}^t_{k^{*}} = 0$, and for $0 \le k \le k^{*}$ it gives, after multiplication by $L^{3k}$, the formula for $\boldsymbol{\varepsilon}^t_k$ in the statement.

\emph{Uniqueness.} Consider any real pair, controlled on $R$, whose trajectory satisfies $m_{k^{*}} = 0$ and $\boldsymbol{\varepsilon}^t_{k^{*}} = 0$. By Proposition~\ref{prop:existence-trajectory}, applied for $k = 0, 1, \dots, k^{*}$, the sequence $(m_k, E_k)_{0 \le k \le k^{*}}$ of its trajectory lies in $\mathsf{B}$; by the recursion of Section~\ref{sec:flow} and $m_{k^{*}} = 0$, it is a fixed point of $\mathrm{Fx}$. Hence it is the fixed point found above. The first entry of the pair is the number $m_0$ of its trajectory, that is, the number $m_0$ of the fixed point; and the identity displayed above, taken at $k = k^{*}$ with $\boldsymbol{\varepsilon}^t_{k^{*}} = 0$, shows that the second entry equals $\boldsymbol{\varepsilon}_{\mathrm{in}}$. So the pair is the one constructed above.

\emph{Bounds.} For the trajectory of this pair put $B_k := |m_k|/w^{\mathrm{m}}_k$ and $A_k := \|E_k\|_k/w^{\mathrm{E}}_k$ for $0 \le k \le k^{*}$; then $A_k \le 1$ by Proposition~\ref{prop:existence-trajectory}, and $B_{k^{*}} = 0$ because $m_{k^{*}} = 0$. The equations of a fixed point give $m_k = L^{-1} m_{k+1} - m(E_k)$ for $k < k^{*}$, and the argument that bounded $|m'_k|$ gives
\[
B_k \le \frac{1}{4}\, B_{k+1} + c_{\mathrm{fl}}\, e_k^{\varepsilon}\, A_k \le \frac{1}{4}\, B_{k+1} + \frac{1}{16} \qquad (0 \le k < k^{*}).
\]
By descending induction on $k$, $B_k \le \frac{1}{16} \sum_{i \ge 0} 4^{-i} = \frac{1}{12} \le \frac{1}{8}$, which is the bound on $m_k$.

For $k \le i < k^{*}$, Proposition~\ref{prop:symmetries-extraction} and Proposition~\ref{prop:existence-trajectory} give, since $e_i = L^{\frac{i-k}{2}} e_k$,
\[
|\varepsilon(E_i)| \le c_{\mathrm{fl}} \|E_i\|_i \le 4\, c_{\mathrm{fl}}\, e_i^{\varepsilon} \cdot w^{\mathrm{E}}_i = 4\, c_{\mathrm{fl}}\, e_i^{\frac{1}{4} - 6\varepsilon} = 4\, c_{\mathrm{fl}}\, e_k^{\frac{1}{4} - 6\varepsilon} L^{(i-k)\left( \frac{1}{8} - 3\varepsilon \right)}.
\]
Since $L \ge 4$ (stage~\ref{it:order-choice-4}),
\[
\sum_{i=k}^{k^{*}-1} L^{-3(i-k)} |\varepsilon(E_i)| \le 4\, c_{\mathrm{fl}}\, e_k^{\frac{1}{4} - 6\varepsilon} \sum_{j \ge 0} L^{-j\left( \frac{23}{8} + 3\varepsilon \right)} \le 4\, c_{\mathrm{fl}}\, e_k^{\frac{1}{4} - 6\varepsilon} \left( 1 - 3^{-\frac{23}{8}} \right)^{-1} \le 4.2\, c_{\mathrm{fl}}\, e_k^{\frac{1}{4} - 6\varepsilon},
\]
because $\left( 1 - 3^{-\frac{23}{8}} \right)^{-1} \le 1.05$. Next, $0 \le d\varepsilon^0_i \le e_i^7$ by clause~\eqref{it:step-small-fields-a} of Theorem~\ref{thm:step-small-fields}, and $L^{-3(i-k)} = (e_k/e_i)^6$; hence
\[
\sum_{i=k}^{k^{*}-1} L^{-3(i-k)}\, d\varepsilon^0_i \le e_k^6 \sum_{i=0}^{k^{*}-1} e_i = e_k^6\, e_{*} \sum_{j=1}^{k^{*}} L^{-\frac{j}{2}} \le e_k^6\, e_{*},
\]
since $\sum_{j \ge 1} L^{-\frac{j}{2}} = \left( L^{\frac{1}{2}} - 1 \right)^{-1} \le 1$. The formula for $\boldsymbol{\varepsilon}^t_k$ now yields
\[
|\boldsymbol{\varepsilon}^t_k| \le 4.2\, c_{\mathrm{fl}}\, e_k^{\frac{1}{4} - 6\varepsilon} + e_k^6\, e_{*} \le 5\, c_{\mathrm{fl}}\, e_k^{\frac{1}{4} - 6\varepsilon} = 5\, c_{\mathrm{fl}}\, e_k^{\varepsilon} \cdot w^{\mathrm{E}}_k;
\]
for the second inequality note that $c_{\mathrm{fl}} \ge 1$ (stage~\ref{it:order-choice-1}), that $e_k^6 \le e_k^{\frac{1}{4} - 6\varepsilon}$, and that $e_{*} = \exp(-v_{*}) \le \exp(-8)$ by condition~\eqref{eq:renormalisation-group-fl}. Theorem~\ref{thm:trajectories-fixed-torus}(b) follows.
\end{proof}

\subsection{A single pair for all tori}\label{sec:single-pair-all}

\begin{definition}[The pair $P_N(e)$]
Let $e_K \le e_1$ and $N \ge N_{\mathrm{min}}(R_{\mathrm{top}}, e)$. As in Definition~\ref{def:tuned-pair}, $P_N(e) = (m_{0,N}, \boldsymbol{\varepsilon}_{0,N})$ is the real pair that Proposition~\ref{prop:pair-vanishing-terminal} provides on the torus $R_{\mathrm{top}}$. The symbols $(\boldsymbol{\varepsilon}^a_k, m^a_k, E^a_k)$ and $(\mathscr{E}^a_k, \mathfrak{P}^a_k)$, $0 \le k \le k_{\mathrm{top}}$, stand for the trajectory and the trajectory with slot of this pair on $R_{\mathrm{top}}$ (Section~\ref{sec:flow}). Thus $P_N(e)$ is the unique real pair that is controlled on $R_{\mathrm{top}}$ and satisfies $m^a_{k_{\mathrm{top}}} = 0$ and $\boldsymbol{\varepsilon}^a_{k_{\mathrm{top}}} = 0$.
\end{definition}

\begin{theorem}[Agreement of the trajectories]\label{thm:agreement-trajectories}
Assume hypothesis~\ref{def:hypothesis-flow}. Then the pair $P_N(e)$ is defined, and it is controlled on $R$. Its trajectory $(\boldsymbol{\varepsilon}^t_k, m_k, E_k)$ on $R$ and its trajectory with slot $(\mathscr{E}_k, \mathfrak{P}_k)$ on $R$ satisfy, for $0 \le k \le k^{*}(R)$,
\begin{equation}\label{eq:renormalisation-flow-11}
m_k = m^a_k, \qquad \boldsymbol{\varepsilon}^t_k = \boldsymbol{\varepsilon}^a_k, \qquad \mathscr{E}_k \sim \mathscr{E}^a_k, \qquad \mathfrak{P}_k \sim \mathfrak{P}^a_k, \qquad E_k \sim E^a_k;
\end{equation}
here $\sim$ is the agreement on flat polymers of Definition~\ref{def:admissibility-flatness-agreement}(e) (for the families $E_k$, $E^a_k$ without slot, taken with $\mathscr{C}_1 = \mathscr{C}_2 = 0$) if $R < R_{\mathrm{top}}$, and $\sim$ is equality if $R = R_{\mathrm{top}}$. Moreover, for $0 \le k \le k^{*}(R)$,
\[
|m_k| \le \frac{w^{\mathrm{m}}_k}{8}, \qquad \|E_k\|_k \le \|\mathscr{E}_k\|^x_k \le 4\, e_k^{\varepsilon} \cdot w^{\mathrm{E}}_k,
\]
and
\[
\left| \boldsymbol{\varepsilon}^t_{k^{*}(R)} \right| \le 5\, c_{\mathrm{fl}} \cdot e_{*}^{\varepsilon} \cdot w^{\mathrm{E}}_{*}.
\]
The pair $P_N(e)$ is real, and it is a function of $(N, e, \bar{m})$ and of the scheme, namely of the systems of primitives of Lemma~\ref{lem:assembled-systems-primitives} on the tori $T_k(R_{\mathrm{top}})$, $k < k_{\mathrm{top}}$, of the constants chosen in the stages~\ref{it:order-choice-1}--\ref{it:order-choice-10} of the order of choice of Section~\ref{sec:order-choice}, and of the basis $\{e_i^0\}$ of Section~\ref{sec:averaging-fluctuations-gauge}, which is fixed once for each $L$ (Proposition~\ref{prop:basis-map-box}); it is independent of the torus $R \in \mathrm{Tor}(e)$ and of the twist $\theta$.
\end{theorem}

\begin{proof}
By Proposition~\ref{prop:consequences} the hypotheses of Proposition~\ref{prop:depth-stop-level} are satisfied. As $R \le R_{\mathrm{top}}$, part~(d) of Proposition~\ref{prop:depth-stop-level} gives
\[
N \ge N_{\mathrm{min}}(R, e) \ge N_{\mathrm{min}}(R_{\mathrm{top}}, e),
\]
so that $P_N(e)$ is defined. By Proposition~\ref{prop:data-along-flow}(a), the torus $R_{\mathrm{top}}$ satisfies hypothesis~\ref{def:hypothesis-flow} with the same $(e, N)$; hence Propositions~\ref{prop:existence-trajectory} and~\ref{prop:pair-vanishing-terminal} hold on $R_{\mathrm{top}}$ as well as on $R$.

Let $R = R_{\mathrm{top}}$. Then the trajectory of $P_N(e)$ on $R$ is $(\boldsymbol{\varepsilon}^a_k, m^a_k, E^a_k)$, with slot $(\mathscr{E}^a_k, \mathfrak{P}^a_k)$; so the relations~\eqref{eq:renormalisation-flow-11}, with equality in place of $\sim$, are identities. That $P_N(e)$ is controlled on $R$, and the bounds for $|m_k|$, $\|E_k\|_k$, $\|\mathscr{E}_k\|^x_k$ and $|\boldsymbol{\varepsilon}^t_{k^{*}(R)}|$, are the assertions of Propositions~\ref{prop:existence-trajectory} and~\ref{prop:pair-vanishing-terminal} on $R_{\mathrm{top}}$.

Let $R < R_{\mathrm{top}}$, and put $k_m := k^{*}(R)$. By Proposition~\ref{prop:depth-stop-level}(d), $k_m \le k_{\mathrm{top}}$. It is shown by induction on $k \in \{0, 1, \dots, k_m\}$ that the trajectory of $P_N(e)$ on $R$ exists up to the level $k$ and that the relations~\eqref{eq:renormalisation-flow-11} hold at all levels $\le k$.

For $k = 0$ the initial values of Section~\ref{sec:flow} are the same on both tori: on $R$ and on $R_{\mathrm{top}}$ alike, the mass parameter of level $0$ is the first entry $m_{0,N}$ of $P_N(e)$, and $\boldsymbol{\varepsilon}^t_0 = \boldsymbol{\varepsilon}^a_0 = \boldsymbol{\varepsilon}_{0,N}$, $E_0 = E^a_0 = 0$, $\mathscr{E}_0 = \mathscr{E}^a_0 = 0$, $\mathfrak{P}_0 = \mathfrak{P}^a_0 = 0$. Families that vanish identically agree on flat polymers; so the relations~\eqref{eq:renormalisation-flow-11} hold at the level $0$.

Let $0 \le k < k_m$, and assume the assertion for $k$. For $i \le k$ one has $i < k_m \le k_{\mathrm{top}}$; so $m_i = m^a_i$ and Proposition~\ref{prop:pair-vanishing-terminal} on $R_{\mathrm{top}}$ give
\[
|m_i| = |m^a_i| \le \frac{w^{\mathrm{m}}_i}{8} \qquad (0 \le i \le k).
\]
Proposition~\ref{prop:existence-trajectory}, applied on $R$ and on $R_{\mathrm{top}}$ at the levels $i \le k + 1$ (note $k + 1 \le k_m \le k_{\mathrm{top}}$), yields the following: the trajectories on $R$ and on $R_{\mathrm{top}}$ both exist up to the level $k + 1$; the number $m_k$ is real; the families $\mathscr{E}_k$ and $\mathscr{E}^a_k$ lie in $K^x_k$, formed with the sides $l_k(R)$ and $l_k(R_{\mathrm{top}})$ respectively, and satisfy $\|\mathscr{E}_k\|^x_k \le w^{\mathrm{E}}_k$ and $\|\mathscr{E}^a_k\|^x_k \le w^{\mathrm{E}}_k$; and $\mathscr{E}_i|_0 = E_i$, $\mathscr{E}^a_i|_0 = E^a_i$ for $i \le k + 1$. Together with the induction hypothesis $\mathscr{E}_k \sim \mathscr{E}^a_k$, $\mathfrak{P}_k \sim \mathfrak{P}^a_k$, this shows that Proposition~\ref{prop:agreement-after-one} applies at the level $k$ with
\[
m = m_k = m^a_k, \qquad \mathscr{E}^{(1)} = \mathscr{E}_k, \qquad \mathscr{E}^{(2)} = \mathscr{E}^a_k, \qquad \mathfrak{P}^{(1)}_k = \mathfrak{P}_k, \qquad \mathfrak{P}^{(2)}_k = \mathfrak{P}^a_k.
\]
Part~(a) of that proposition gives $\varepsilon(\mathscr{E}_k) = \varepsilon(\mathscr{E}^a_k)$ and $m(\mathscr{E}_k) = m(\mathscr{E}^a_k)$, and shows that $d\varepsilon^0_k$ is the same number on both tori. Since $\varepsilon(\mathscr{E}) = \varepsilon(\mathscr{E}|_0)$ and $m(\mathscr{E}) = m(\mathscr{E}|_0)$ (Proposition~\ref{prop:symmetries-extraction}), and $\mathscr{E}_k|_0 = E_k$, $\mathscr{E}^a_k|_0 = E^a_k$, one has $\varepsilon(E_k) = \varepsilon(E^a_k)$ and $m(E_k) = m(E^a_k)$. The recursion of Section~\ref{sec:flow} and the induction hypothesis then give
\[
m_{k+1} = L \left( m_k + m(E_k) \right) = L \left( m^a_k + m(E^a_k) \right) = m^a_{k+1},
\]
\[
\boldsymbol{\varepsilon}^t_{k+1} = L^3 \left( \boldsymbol{\varepsilon}^t_k + \varepsilon(E_k) + d\varepsilon^0_k \right) = L^3 \left( \boldsymbol{\varepsilon}^a_k + \varepsilon(E^a_k) + d\varepsilon^0_k \right) = \boldsymbol{\varepsilon}^a_{k+1}.
\]
By the same recursion, $\mathscr{E}_{k+1}$ and $\mathscr{E}^a_{k+1}$ are the families $\mathscr{E}^{(1)+}$ and $\mathscr{E}^{(2)+}$ of Proposition~\ref{prop:agreement-after-one}(b), and $\mathfrak{P}_{k+1}$, $\mathfrak{P}^a_{k+1}$ are obtained from $\mathfrak{P}_k$, $\mathfrak{P}^a_k$ by the rule of clause~\eqref{it:step-small-fields-c} of Theorem~\ref{thm:step-small-fields}. Parts~(b) and~(c) of Proposition~\ref{prop:agreement-after-one} therefore give $\mathscr{E}_{k+1} \sim \mathscr{E}^a_{k+1}$ and $\mathfrak{P}_{k+1} \sim \mathfrak{P}^a_{k+1}$. Finally $E_{k+1} = \mathscr{E}_{k+1}|_0$ and $E^a_{k+1} = \mathscr{E}^a_{k+1}|_0$; taking $\mathscr{C}_1 = \mathscr{C}_2 = 0$ in Definition~\ref{def:admissibility-flatness-agreement}(e), one sees that the agreement of two families with slot implies the agreement of their restrictions to $\mathscr{C} = 0$. Hence $E_{k+1} \sim E^a_{k+1}$, and the induction is complete.

So the trajectory of $P_N(e)$ on $R$ exists up to the level $k_m = k^{*}(R)$, the relations~\eqref{eq:renormalisation-flow-11} hold for $0 \le k \le k_m$, and, by Proposition~\ref{prop:pair-vanishing-terminal} on $R_{\mathrm{top}}$ at the levels $k \le k_m \le k_{\mathrm{top}}$,
\[
|m_k| = |m^a_k| \le \frac{w^{\mathrm{m}}_k}{8} \le w^{\mathrm{m}}_k \qquad (0 \le k \le k_m).
\]
Thus $P_N(e)$ is controlled on $R$ (Definition~\ref{def:fluctuation-field-controlled}). The inequalities $\|E_k\|_k \le \|\mathscr{E}_k\|^x_k \le 4\, e_k^{\varepsilon} \cdot w^{\mathrm{E}}_k$ for $k \le k_m$ are Proposition~\ref{prop:existence-trajectory} on $R$. By~\eqref{eq:renormalisation-flow-11} at the level $k_m$ and by Proposition~\ref{prop:pair-vanishing-terminal} on $R_{\mathrm{top}}$ at the level $k_m \le k_{\mathrm{top}}$,
\[
\left| \boldsymbol{\varepsilon}^t_{k_m} \right| = \left| \boldsymbol{\varepsilon}^a_{k_m} \right| \le 5\, c_{\mathrm{fl}} \cdot e_{k_m}^{\varepsilon} \cdot w^{\mathrm{E}}_{k_m},
\]
and $e_{k_m}$, $w^{\mathrm{E}}_{k_m}$ are the numbers $e_{*}$, $w^{\mathrm{E}}_{*}$ of the torus $R$, since these numbers depend on the level and not on the torus.

In both cases $P_N(e)$ is real by Proposition~\ref{prop:pair-vanishing-terminal}. It is defined through its trajectory on $R_{\mathrm{top}}$ alone, which is formed from $(N, e, \bar{m})$, from the systems of primitives of Lemma~\ref{lem:assembled-systems-primitives} on the tori $T_k(R_{\mathrm{top}})$, $k < k_{\mathrm{top}}$, from the constants of the stages~\ref{it:order-choice-1}--\ref{it:order-choice-10} of Section~\ref{sec:order-choice}, and from the basis $\{e_i^0\}$ of Section~\ref{sec:averaging-fluctuations-gauge}. The basis enters the recursion of Section~\ref{sec:flow} through the family $\{\Upsilon\}$ of~\eqref{eq:renormalisation-group-1}: the set $\mathrm{Box}_k$ and the measure $\hat{\mu}$ are defined through the coordinates $\zeta(\Upsilon)$, the map $C$ enters through its matrix with respect to the basis $\{\Upsilon\}$ of $V_k$, and the kernels of the operators of the family~\eqref{it:primitives-p2} are indexed by the elements $\Upsilon$. For a given $L$ the basis $\{e_i^0\}$ is fixed once, the same vectors being taken at every level and for every side (Section~\ref{sec:averaging-fluctuations-gauge} and Proposition~\ref{prop:basis-map-box}), and it depends neither on $R$ nor on $\theta$. So the torus $R$ plays no role in this definition; nor does $\theta$, which appears nowhere in the recursion of Section~\ref{sec:flow} by which trajectories are defined. Theorem~\ref{thm:one-pair-all} follows.
\end{proof}

\subsection{The representation at the stop level}\label{sec:representation-stop-level}

Assume hypothesis~\ref{def:hypothesis-flow}, and let a real pair be controlled on $R$. The complex functions $Y^x_{*}$ on $\mathrm{Dom}_{k^{*}} \times \mathfrak{X}_{k^{*}}$ and $Y_{*}$ on $\mathrm{Dom}_{k^{*}}$ are defined by
\[
Y^x_{*}(A; \mathscr{C}) := y^x_{k^{*}}\left[ \boldsymbol{\varepsilon}^t_{*}, m_{*}, \mathscr{E}_{*} \right](A; \mathscr{C}), \qquad Y_{*}(A) := Y^x_{*}(A; 0),
\]
where $y^x_{k^{*}}$ is the function of Section~\ref{sec:fermion-integrals-blocking}, formed with the family $\mathfrak{P}_{*} := \mathfrak{P}_{k^{*}}$. Let $A \in \mathrm{Dom}_{k^{*}}$, and put $\mathscr{A} := H_{k^{*}} A$, $\psi^{\#} := H^{\#}_{k^{*}}(\mathscr{A}) \Psi$, with the first components $\psi = H_{k^{*}}(\mathscr{A}) \Psi$ and $\bar{\psi} = \bar{H}_{k^{*}}(\mathscr{A}) \bar{\Psi}$, and let $D_{*} := D_{k^{*}}$ be the operator of Section~\ref{sec:exact-fermion-operators} at the level $k^{*}$. By Proposition~\ref{prop:existence-trajectory}, $\mathscr{E}_{*}|_0 = E_{*}$ and $\mathfrak{P}_{*}(\cdot, 0) = 0$; moreover $n_{k^{*}} = l_{*}^3$. Hence the definitions of Section~\ref{sec:fermion-integrals-blocking} give
\begin{equation}\label{eq:renormalisation-flow-12}
Y_{*}(A) = \left( \det D_{*}(0) \right)^{-1} \int \mathrm{D}\Psi \, \exp\left\{ - \left\langle \bar{\Psi}, D_{*}(\mathscr{A}) \Psi \right\rangle - m_{*} \cdot \langle \bar{\psi}, \psi \rangle - \boldsymbol{\varepsilon}^t_{*} \cdot l_{*}^3 + \sum_X E_{*}\left( X, \mathscr{A}, \psi^{\#} \right) \right\}.
\end{equation}

\begin{theorem}[Representation with twist and external connection]\label{thm:representation-twist-external-2}
Assume hypothesis~\ref{def:hypothesis-flow}. Suppose that a real pair $(m_{\mathrm{in}}, \boldsymbol{\varepsilon}_{\mathrm{in}})$ is controlled on $R$, and that $a \in C^1 \otimes \mathbb{C}$ satisfies $a^{[j]} \in \mathfrak{X}_j$ for every $j \le k^{*}$; write $\mathscr{C}_{*} := a^{[k^{*}]}$. Let $A_{*}$, $\zeta = (\zeta_j)_{j<k^{*}}$ with $\zeta_j \in V_j$, $A^{\downarrow}_0$, $\mathrm{d}\mathsf{m}_{*}$, $\chi_{*}$ and $\mathfrak{G}$ be as in Definition~\ref{def:coordinates-slice-small}, and let $c_{\theta} = c^{(k^{*})}_{\theta}$ be the twist of Definition~\ref{def:twist-external-connection}.
\begin{enumerate}
\item[(a)] For every $A \in \mathrm{Dom}_{k^{*}}$,
\[
Y^x_{*}(A; \mathscr{C}_{*}) = \int \prod_{j<k^{*}} \left( \mathrm{d}\mu_I(\zeta_j) \, \mathbf{1}_{\mathrm{Box}_j}(\zeta_j) \right) y^x_0\left[ \boldsymbol{\varepsilon}_{\mathrm{in}}, m_{\mathrm{in}}, 0 \right]\left( A^{\downarrow}_0(A, \zeta); a \right),
\]
the integrand being read as $0$ where one of the factors $\mathbf{1}_{\mathrm{Box}_j}(\zeta_j)$ vanishes.
\item[(b)] Let $\lambda$ be a real function on $T^0_{k^{*}}$. If $A$ and $A + d\lambda$ lie in $\mathrm{Dom}_{k^{*}}$, then $Y^x_{*}(A + d\lambda; \mathscr{C}_{*}) = Y^x_{*}(A; \mathscr{C}_{*})$.
\item[(c)] Let $\theta \in \mathbb{R}^3$, and let $\Phi \colon V_{\mathrm{t}}(T^0_{k^{*}}) \to \mathbb{C}$ be a bounded Borel function; on the left side below, $\Phi$ is evaluated at the coordinate $A_{*}$ of $F \in F_N$. Then
\[
\nu^{\theta,a}\left( \mathbf{1}_{\mathfrak{G}} \cdot \Phi \right) = \int \mathrm{d}\mathsf{m}_{*}(A) \, \Phi(A) \cdot \chi_{*}(A) \cdot Y^x_{*}\left( A + c_{\theta}; \mathscr{C}_{*} \right),
\]
the integrand being read as $0$ where $\chi_{*}(A) = 0$.
\end{enumerate}
\end{theorem}

\begin{proof}
\emph{(0) Gauge invariance at the level $0$.} Let $B$ be a complex function on the oriented bonds of $\Lambda_n$ with $B(x',x) = -B(x,x')$, and let $\lambda$ be a real function on $\Lambda_n$. Since $(d\lambda)(x,x') = \lambda(x') - \lambda(x)$, every bond $(x,x')$ satisfies
\[
\exp\left( iq (B + d\lambda)(x,x') \right) \exp\left( -iq\lambda(x') \right) = \exp\left( -iq\lambda(x) \right) \exp\left( iq B(x,x') \right).
\]
Let $\exp(\pm iq\lambda)$ also denote the operators of multiplication by these functions; they map antiperiodic spinor functions to antiperiodic spinor functions. Let $f$ be antiperiodic and $x$ a site. By the identity just displayed, each term of $(D_{B + d\lambda} f)(x)$, formed with the formula for $D_A$ of Section~\ref{sec:lattice-measures-laws}, that carries a link phase,
\[
\frac{1 \mp \gamma_{\nu}}{2} \exp\left( iq (B + d\lambda)(x, x \pm e_{\nu}) \right) f(x \pm e_{\nu}),
\]
equals $\exp(-iq\lambda(x))$ times the corresponding term of $\left( D_B \exp(iq\lambda) f \right)(x)$, and the remaining term is $f(x) = \exp(-iq\lambda(x)) \cdot \exp(iq\lambda(x)) f(x)$. Hence
\[
D_{B + d\lambda} = \exp(-iq\lambda) \, D_B \, \exp(iq\lambda),
\]
and $\det\left( D_{B + d\lambda} + \mathbf{m} \right) = \det\left( D_B + \mathbf{m} \right)$ for every number $\mathbf{m}$. By Lemma~\ref{lem:three-berezin-integrals}(c), for $A_0 \in \mathrm{Dom}_0$,
\[
y^x_0\left[ \boldsymbol{\varepsilon}_{\mathrm{in}}, m_{\mathrm{in}}, 0 \right](A_0; a) = \exp\left( -\boldsymbol{\varepsilon}_{\mathrm{in}} \cdot n^3 \right) \cdot \frac{\det\left( D_{A_0 + a} + \bar{m} \cdot \eta + m_{\mathrm{in}} \right)}{\det\left( D_0 + \bar{m} \cdot \eta \right)}.
\]
The determinant identity with $B = A_0 + a$ and $\mathbf{m} = \bar{m} \cdot \eta + m_{\mathrm{in}}$ gives
\[
y^x_0\left[ \boldsymbol{\varepsilon}_{\mathrm{in}}, m_{\mathrm{in}}, 0 \right](A_0 + d\lambda; a) = y^x_0\left[ \boldsymbol{\varepsilon}_{\mathrm{in}}, m_{\mathrm{in}}, 0 \right](A_0; a)
\]
whenever $A_0$ and $A_0 + d\lambda$ lie in $\mathrm{Dom}_0$.

\emph{Proof of (a).} By the definition of $a^{[j]}$ in Section~\ref{sec:lattice-measures-laws} and that of the scaling $A_L(b) = L^{-\frac{1}{2}} A(b/L)$ in Section~\ref{sec:local-parts-adjustment}, for $j \ge 0$ and every bond $b$ of $T_j$,
\[
\left( a^{[j+1]} \right)_L(b) = L^{-\frac{1}{2}} \, a^{[j+1]}\left( \frac{b}{L} \right) = L^{-\frac{1}{2}} \cdot L^{\frac{j+1}{2}} \, a\left( L^{j+1} \cdot \frac{b}{L} \right) = L^{\frac{j}{2}} \, a\left( L^j b \right) = a^{[j]}(b);
\]
that is, $(a^{[j+1]})_L = a^{[j]}$.

By Proposition~\ref{prop:data-along-flow}(b), hypothesis~\ref{def:hypothesis-one-step} holds at every level $j < k^{*}$; hence Theorem~\ref{thm:step-small-fields} and Proposition~\ref{prop:fields-domains}(ii) apply at these levels. Let $j < k^{*}$. By Proposition~\ref{prop:existence-trajectory} the numbers $\boldsymbol{\varepsilon}^t_j$, $m_j$ are real, $\mathscr{E}_j \in K^x_j$ with $\|\mathscr{E}_j\|^x_j \le w^{\mathrm{E}}_j$, and $\mathscr{E}_j|_0 = E_j$; the initial value $\mathscr{E}_0 = 0$ is that of Section~\ref{sec:flow}; moreover $|m_j| \le w^{\mathrm{m}}_j$, the pair being controlled on $R$ (Definition~\ref{def:fluctuation-field-controlled}). So the data $(\boldsymbol{\varepsilon}^t_j, m_j, \mathscr{E}_j)$ satisfy the hypotheses of clause~\eqref{it:step-small-fields-c} of Theorem~\ref{thm:step-small-fields}. Since $\varepsilon(\mathscr{E}_j) = \varepsilon(E_j)$ and $m(\mathscr{E}_j) = m(E_j)$ (Proposition~\ref{prop:symmetries-extraction}), the recursion of Section~\ref{sec:flow} shows that the image of these data under the map of that clause is $(\boldsymbol{\varepsilon}^t_{j+1}, m_{j+1}, \mathscr{E}_{j+1})$, and that $\mathfrak{P}_{j+1}$ arises from $\mathfrak{P}_j$ by the rule of the same clause.

For $0 \le k \le k^{*}$, $A \in \mathrm{Dom}_k$ and $\zeta_j \in \mathrm{Box}_j$ ($j < k$) define
\[
A^{\downarrow}_k := A, \qquad A^{\downarrow}_j := A^{\min}_j(A^{\downarrow}_{j+1}) + \mathrm{Sq}_j \zeta_j \qquad (j = k - 1, \dots, 0).
\]
By Proposition~\ref{prop:fields-domains}(ii), applied successively for $j = k - 1, \dots, 0$, one has $A^{\downarrow}_j \in \mathrm{Dom}_j$ for all $j \le k$. It is shown by induction on $k$ that, for $0 \le k \le k^{*}$ and $A \in \mathrm{Dom}_k$,
\[
y^x_k\left[ \boldsymbol{\varepsilon}^t_k, m_k, \mathscr{E}_k \right]\left( A; a^{[k]} \right) = \int \prod_{j<k} \left( \mathrm{d}\mu_I(\zeta_j) \, \mathbf{1}_{\mathrm{Box}_j}(\zeta_j) \right) y^x_0\left[ \boldsymbol{\varepsilon}_{\mathrm{in}}, m_{\mathrm{in}}, 0 \right]\left( A^{\downarrow}_0; a \right),
\]
where $y^x_k$ is the function of Section~\ref{sec:fermion-integrals-blocking}, formed with the family $\mathfrak{P}_k$. For $k = 0$ no integration occurs, and both sides coincide: at the level $0$ the trajectory has $\boldsymbol{\varepsilon}^t_0 = \boldsymbol{\varepsilon}_{\mathrm{in}}$, the mass parameter $m_{\mathrm{in}}$, $\mathscr{E}_0 = 0$ and $\mathfrak{P}_0 = 0$ (Section~\ref{sec:flow}), and moreover $a^{[0]} = a$ and $A^{\downarrow}_0 = A$. Let $k < k^{*}$, assume the identity at the level $k$, and let $A \in \mathrm{Dom}_{k+1}$. Since $a^{[k+1]} \in \mathfrak{X}_{k+1}$ and $(a^{[k+1]})_L = a^{[k]}$, the identity~\eqref{eq:renormalisation-group-11} of clause~\eqref{it:step-small-fields-c} of Theorem~\ref{thm:step-small-fields}, taken with $A' = A$ and $\mathscr{C}' = a^{[k+1]}$, gives
\[
y^x_{k+1}\left[ \boldsymbol{\varepsilon}^t_{k+1}, m_{k+1}, \mathscr{E}_{k+1} \right]\left( A; a^{[k+1]} \right) = \int_{\mathrm{Box}_k} \mathrm{d}\mu_I(\zeta_k) \; y^x_k\left[ \boldsymbol{\varepsilon}^t_k, m_k, \mathscr{E}_k \right]\left( A^{\min}_k(A) + \mathrm{Sq}_k \zeta_k; a^{[k]} \right).
\]
For $\zeta_k \in \mathrm{Box}_k$ the field $A^{\min}_k(A) + \mathrm{Sq}_k \zeta_k$ lies in $\mathrm{Dom}_k$ (Proposition~\ref{prop:fields-domains}(ii)). The identity at the level $k$, applied at this field, writes the integrand as an integral over $(\zeta_j)_{j<k}$; and the field $A^{\downarrow}_0$ formed from $A^{\min}_k(A) + \mathrm{Sq}_k \zeta_k$ and $(\zeta_j)_{j<k}$ is the field $A^{\downarrow}_0$ formed from $A$ and $(\zeta_j)_{j \le k}$. Each integrand is continuous on the compact set over which it is integrated, a product of the sets $\mathrm{Box}_j$; hence Fubini's theorem applies, and it turns the integral over $\zeta_k$ of the inner integral over $(\zeta_j)_{j<k}$ into the integral over $(\zeta_j)_{j \le k}$ that stands in the identity at the level $k + 1$. For $k = k^{*}$ the left side of the identity is $Y^x_{*}(A; \mathscr{C}_{*})$, because $\mathscr{C}_{*} = a^{[k^{*}]}$ and $\mathfrak{P}_{*} = \mathfrak{P}_{k^{*}}$; this is (a).

\emph{Proof of (b).} The one-cochain $d\lambda$ on $T^0_{k^{*}}$ is closed and has mean zero. By Proposition~\ref{prop:closed-shifts-terminal}, applied with $A^c = d\lambda$ and mean $0$, for every $\zeta$ the difference $A^{\downarrow}_0(A + d\lambda, \zeta) - A^{\downarrow}_0(A, \zeta)$ is exact; both fields being real, it is the coboundary of a real function on $\Lambda_n$. If $\zeta_j \in \mathrm{Box}_j$ for all $j < k^{*}$, both fields lie in $\mathrm{Dom}_0$ (Proposition~\ref{prop:fields-domains}(ii)), and by (0) the integrands of (a) at $A + d\lambda$ and at $A$ coincide. As both integrands are read as $0$ whenever some $\zeta_j$ lies outside $\mathrm{Box}_j$, part (a) gives $Y^x_{*}(A + d\lambda; \mathscr{C}_{*}) = Y^x_{*}(A; \mathscr{C}_{*})$.

\emph{Proof of (c).} Let $F \in \mathfrak{G}$, let $A_0 := A_0(F)$ be the potential of $F$ in the slice $\mathrm{Sl}$ of Definition~\ref{def:coordinates-slice-small}, so that $A_0 \in \mathrm{Sl}$ and $dA_0 = F$, and let $A_j$, $0 \le j \le k^{*}$, be the fields formed from $A_0$ in that definition; for $j = 0$ this is $A_0$ itself, since $Q_0$ and the scaling of level $0$ are the identity. The coordinates $(A_{*}; \zeta)$, $\zeta = (\zeta_0, \dots, \zeta_{k^{*}-1})$, of $F$ in the sense of Proposition~\ref{prop:orthogonal-coordinates} are given by $A_{*} = A_{k^{*}}$ and by $A_j = A^{\min}_j(A_{j+1}) + \mathrm{Sq}_j \zeta_j$ for $j < k^{*}$ (Definition~\ref{def:coordinates-slice-small}). Thus the fields $A_j$ obey the recursion that defines the fields $A^{\downarrow}_j(A_{*}, \zeta)$, with the same top field $A_{*}$, and downward induction on $j$ gives $A^{\downarrow}_j(A_{*}, \zeta) = A_j$ for $j = k^{*}, k^{*} - 1, \dots, 0$. In particular
\[
A_0 = A^{\downarrow}_0(A_{*}, \zeta).
\]
By Definition~\ref{def:coordinates-slice-small}, $A_{*}$ is $Q_{k^{*}} A_0$ in level-$k^{*}$ scaling, and $Q_{k^{*}} A_0$ has mean zero; since $Q_{k^{*}}$ preserves means (Proposition~\ref{prop:averaging-operators}(c), taken for the operator $Q_{k^{*}}$ on the torus $\Lambda_n = \mathbb{T}(1, n)$ of spacing $1$ and side $n$; the hypothesis of that proposition, here $n \in L^{k^{*}} \mathbb{N}$, holds, as $n = L^{k^{*}} \cdot l_{*}$ with the positive integer $l_{*} = M \cdot r_{*}$, Section~\ref{sec:coordinates-axial-fields}), $A_0$ has mean zero. The potential $A(F) \in V_{\mathrm{t}}$ has mean zero as well, and $dA(F) = F = dA_0$. Therefore $A_0 - A(F)$ is a closed one-cochain with mean zero, hence exact by the Hodge decomposition on $\Lambda_n$. Combined with Proposition~\ref{prop:closed-shifts-terminal}, this gives real functions $\lambda$, $\lambda'$ on $\Lambda_n$ such that
\[
A^{\downarrow}_0(A_{*} + c_{\theta}, \zeta) = A_0 + H_{\theta} + d\lambda = A(F) + H_{\theta} + d\lambda'.
\]
As $F \in \mathfrak{G}$ and $c_{\theta}$ is constant, $|d(A_{*} + c_{\theta})|_{\infty} = |dA_{*}|_{\infty} \le p_{*}$ and $\zeta_j \in \mathrm{Box}_j$ for $j < k^{*}$. So $A_{*} + c_{\theta} \in \mathrm{Dom}_{k^{*}}$ by Proposition~\ref{prop:fields-domains}(i), which holds at the level $k^{*}$ (Section~\ref{sec:one-step-hypothesis}), and then $A^{\downarrow}_0(A_{*} + c_{\theta}, \zeta) \in \mathrm{Dom}_0$ by Proposition~\ref{prop:fields-domains}(ii). Lemma~\ref{lem:three-berezin-integrals}(c), followed by the determinant identity of (0) with $B = A(F) + H_{\theta} + a$ and the function $\lambda'$, gives
\[
y^x_0\left[ \boldsymbol{\varepsilon}_{\mathrm{in}}, m_{\mathrm{in}}, 0 \right]\left( A^{\downarrow}_0(A_{*} + c_{\theta}, \zeta); a \right) = \exp\left( -\boldsymbol{\varepsilon}_{\mathrm{in}} \cdot n^3 \right) \cdot \frac{\det\left( D_{A(F) + H_{\theta} + a} + \bar{m} \cdot \eta + m_{\mathrm{in}} \right)}{\det\left( D_0 + \bar{m} \cdot \eta \right)} = \rho^{\theta,a}_N(F)
\]
for every $F \in \mathfrak{G}$, the last equality being the definition of $\rho^{\theta,a}_N$ in Section~\ref{sec:lattice-measures-laws}. The function $\rho^{\theta,a}_N$ is bounded and continuous on $F_N$: the potential $A(F)$ depends linearly on $F$, and the link phases $\exp\left( iq (A(F) + H_{\theta} + a) \right)$ have the moduli $\exp\left( \mp q \cdot \operatorname{Im} a \right)$, which do not depend on $F$. Proposition~\ref{prop:orthogonal-coordinates} shows that $A_{*}, \zeta_0, \dots, \zeta_{k^{*}-1}$, as random variables on $(F_N, \gamma_N)$, are independent, with the law $\mu_I$ for each $\zeta_j$; the law of $A_{*}$ is $\mathrm{d}\mathsf{m}_{*}$ (Definition~\ref{def:coordinates-slice-small}). By Definition~\ref{def:coordinates-slice-small},
\[
\mathbf{1}_{\mathfrak{G}}(F) = \chi_{*}(A_{*}) \cdot \prod_{j<k^{*}} \mathbf{1}_{\mathrm{Box}_j}(\zeta_j).
\]
Let $A \in V_{\mathrm{t}}(T^0_{k^{*}})$. If $\chi_{*}(A) = 1$, then $|d(A + c_{\theta})|_{\infty} = |dA|_{\infty} \le p_{*}$; so $A + c_{\theta} \in \mathrm{Dom}_{k^{*}}$ by Proposition~\ref{prop:fields-domains}(i), and $A^{\downarrow}_0(A + c_{\theta}, \zeta) \in \mathrm{Dom}_0$ whenever $\zeta_j \in \mathrm{Box}_j$ for all $j < k^{*}$, by Proposition~\ref{prop:fields-domains}(ii). For such $A$ put
\[
J_{\theta}(A) := \int \prod_{j<k^{*}} \left( \mathrm{d}\mu_I(\zeta_j) \, \mathbf{1}_{\mathrm{Box}_j}(\zeta_j) \right) y^x_0\left[ \boldsymbol{\varepsilon}_{\mathrm{in}}, m_{\mathrm{in}}, 0 \right]\left( A^{\downarrow}_0(A + c_{\theta}, \zeta); a \right),
\]
and put $J_{\theta}(A) := 0$ if $\chi_{*}(A) = 0$. By (a), $J_{\theta}(A) = Y^x_{*}(A + c_{\theta}; \mathscr{C}_{*})$ whenever $\chi_{*}(A) = 1$. Since $\nu^{\theta,a} = \rho^{\theta,a}_N \cdot \gamma_N$ (Section~\ref{sec:lattice-measures-laws}),
\[
\nu^{\theta,a}\left( \mathbf{1}_{\mathfrak{G}} \cdot \Phi \right) = \int \mathrm{d}\gamma_N(F) \, \mathbf{1}_{\mathfrak{G}}(F) \cdot \Phi(A_{*}) \cdot \rho^{\theta,a}_N(F).
\]
Since $\mathbf{1}_{\mathfrak{G}}$ vanishes off $\mathfrak{G}$, the density $\rho^{\theta,a}_N(F)$ may here be replaced by $y^x_0\left[ \boldsymbol{\varepsilon}_{\mathrm{in}}, m_{\mathrm{in}}, 0 \right]\left( A^{\downarrow}_0(A_{*} + c_{\theta}, \zeta); a \right)$, by the identity proved above for $F \in \mathfrak{G}$. The resulting integrand is bounded, and Fubini's theorem, applied in the coordinates $(A_{*}; \zeta)$ with the product law just described, gives
\[
\nu^{\theta,a}\left( \mathbf{1}_{\mathfrak{G}} \cdot \Phi \right) = \int \mathrm{d}\mathsf{m}_{*}(A) \, \Phi(A) \cdot \chi_{*}(A) \cdot J_{\theta}(A) = \int \mathrm{d}\mathsf{m}_{*}(A) \, \Phi(A) \cdot \chi_{*}(A) \cdot Y^x_{*}\left( A + c_{\theta}; \mathscr{C}_{*} \right).
\]
This proves (c). Theorem~\ref{thm:representation-twist-external} follows.
\end{proof}

\subsection{Determinants at constant connections}\label{sec:determinants-constant-connections}

In the whole of this section the symbol $\mathbf{m}$ stands for a mass expressed in units of the lattice spacing. The letter $m$ in Section~\ref{sec:averaging-blocks-flat} is not a mass but an index, ranging over the set $\mathfrak{M}$ introduced there; in that subsection $q$ is a momentum variable, as $p$ is, whereas in Section~\ref{sec:constant-twists-free} the letter $q$ keeps the meaning it has in the link phases $\exp(i\,q\cdot A)$ of the operator $D_A$ of Section~\ref{sec:lattice-measures-laws}. For $k\in\{2,3\}$ and $v\in\mathbb{R}^k$, $|v|_2$ is the Euclidean norm of $v$.

\subsubsection{Constant twists of the free operator}\label{sec:constant-twists-free}

Let $\mathbf{m}>0$. The function $w_{\mathbf{m}}\colon\mathbb{R}^3\to\mathbb{R}$ is defined by
\[
w_{\mathbf{m}}(p):=\left(\mathbf{m}+\sum_{\nu=0}^2\left(1-\cos p_{\nu}\right)\right)^2+\sum_{\nu=0}^2\sin^2p_{\nu}.
\]
The first square is at least $\mathbf{m}^2$, so that $w_{\mathbf{m}}\ge\mathbf{m}^2>0$ on $\mathbb{R}^3$, and
\[
p^0(\mathbf{m}):=\frac{2}{(2\pi)^3}\int_{[-\pi,\pi]^3}\mathrm{d}^3p\,\log w_{\mathbf{m}}(p)
\]
is a real number. For an integer $n\ge1$ and $\theta\in\mathbb{R}^3$ let $P_{n,\theta}\subset\mathbb{R}^3$ be the set of the $n^3$ points
\[
p=\left(\frac{2\pi j_{\nu}+\pi+\theta_{\nu}}{n}\right)_{\nu=0,1,2},\qquad j\in\{0,\dots,n-1\}^3.
\]
As in Section~\ref{sec:lattice-measures-laws}, $H_{\theta}$ is the constant connection on the bonds of $\Lambda_n$ with $H_{\theta}(x,x+e_{\nu})=\theta_{\nu}/(q\cdot n)$, $D_A$ is the Wilson operator on the spinors $\phi\colon\mathbb{Z}^3\to\mathbb{C}^4$ with $\phi(x+n\cdot e_{\nu})=-\phi(x)$ for every $\nu$, and $D_0$ is $D_A$ at $A=0$.

\begin{lemma}[Product formula and volume asymptotics]\label{lem:product-formula-volume}
Let $n\ge1$ be an integer, $\theta\in\mathbb{R}^3$, $0<\mathbf{m}\le1$ and $q>0$. If $n\le2$, the bonds $(x,x+e_{\nu})$ and $(x,x-e_{\nu})$ of $\Lambda_n$ have the same end points; for such $n$ the connection $H_{\theta}$ is taken to be the function on the oriented links of $\mathbb{Z}^3$ that has period $n$ in each direction, equals $\theta_{\nu}/(q\cdot n)$ on every link from $x$ to $x+e_{\nu}$ and changes sign when the orientation is reversed, in the manner in which Section~\ref{sec:characters-twisted-spinors} treats one-cochains. Then
\begin{multline}\label{eq:renormalisation-flow-13}
\det\left(D_{H_{\theta}}+\mathbf{m}\right)=\prod_{p\in P_{n,\theta}}w_{\mathbf{m}}(p)^2>0,\\
\left|\log\det\left(D_{H_{\theta}}+\mathbf{m}\right)-p^0(\mathbf{m})\cdot n^3\right|\\
\le-312\log\left(1-\exp\left(-\frac{\mathbf{m}\cdot n}{4}\right)\right)=\varepsilon_5(\mathbf{m}\cdot n).
\end{multline}
\end{lemma}

\begin{proof}
\emph{(a) Plane waves.} For every $q>0$ one has $q\cdot H_{\theta}(x,x\pm e_{\nu})=\pm\theta_{\nu}/n$, the sign coming from the antisymmetry of $H_{\theta}$ under reversal of orientation (for $n\le2$ with $H_{\theta}$ as in the statement). Inserted into the formula for $D_A$ of Section~\ref{sec:lattice-measures-laws}, these values give $D_{H_{\theta}}$ the link phases $\exp\left(\pm\frac{i\theta_{\nu}}{n}\right)$:
\begin{align*}
\left(D_{H_{\theta}}\phi\right)(x)=3\phi(x)&-\sum_{\nu=0}^2\frac{1-\gamma_{\nu}}{2}\cdot\exp\left(\frac{i\theta_{\nu}}{n}\right)\cdot\phi(x+e_{\nu})\\
&-\sum_{\nu=0}^2\frac{1+\gamma_{\nu}}{2}\cdot\exp\left(-\frac{i\theta_{\nu}}{n}\right)\cdot\phi(x-e_{\nu}).
\end{align*}
For $j\in\{0,\dots,n-1\}^3$ let $p'\in\mathbb{R}^3$ have the coordinates $p'_{\nu}=(2\pi j_{\nu}+\pi)/n$. The $4n^3$ functions $x\mapsto\exp(i\,p'\cdot x)\cdot v$, where $j$ ranges over $\{0,\dots,n-1\}^3$ and $v$ over the standard basis of $\mathbb{C}^4$, are antiperiodic (as $\exp\left(i\,n\cdot p'_{\nu}\right)=-1$) and pairwise orthogonal for the scalar product $\sum_{x\in\{0,\dots,n-1\}^3}\phi(x)^{*}\cdot\phi'(x)$; the antiperiodic spinors form a space of dimension $4n^3$, so these functions are a basis of it. For real $p_{\nu}$,
\[
\frac{1-\gamma_{\nu}}{2}\cdot\exp(i\,p_{\nu})+\frac{1+\gamma_{\nu}}{2}\cdot\exp(-i\,p_{\nu})=\cos p_{\nu}-i\gamma_{\nu}\cdot\sin p_{\nu}.
\]
Hence, with $p:=p'+\theta/n$, a point of $P_{n,\theta}$,
\[
\left(D_{H_{\theta}}+\mathbf{m}\right)\left(\exp(i\,p'\cdot x)\cdot v\right)=\exp(i\,p'\cdot x)\cdot\left(a(p)+i\gamma\cdot s(p)\right)\cdot v,
\]
where
\[
a(p):=\mathbf{m}+\sum_{\nu=0}^2\left(1-\cos p_{\nu}\right),\qquad s_{\nu}(p):=\sin p_{\nu},\qquad\gamma\cdot s:=\sum_{\nu=0}^2\gamma_{\nu}\cdot s_{\nu}.
\]
The matrices $\gamma_{\nu}$ of Section~\ref{sec:dirac-matrices-lattice} are Hermitian and satisfy $\gamma_{\nu}\cdot\gamma_{\nu'}+\gamma_{\nu'}\cdot\gamma_{\nu}=2\delta_{\nu\nu'}$; they are traceless, since $\operatorname{tr}\gamma_{\nu}=\operatorname{tr}\left(\gamma_{\nu'}\cdot\gamma_{\nu'}\cdot\gamma_{\nu}\right)=-\operatorname{tr}\left(\gamma_{\nu'}\cdot\gamma_{\nu}\cdot\gamma_{\nu'}\right)=-\operatorname{tr}\gamma_{\nu}$ for $\nu'\ne\nu$. Thus $\gamma\cdot s$ is a traceless Hermitian matrix with $(\gamma\cdot s)^2=|s|_2^2$. Its eigenvalues are $|s|_2$ and $-|s|_2$, each of multiplicity two, and $a+i\gamma\cdot s$ has the eigenvalues $a+i|s|_2$ and $a-i|s|_2$, each of multiplicity two. It follows that
\[
\det\left(a(p)+i\gamma\cdot s(p)\right)=\left(a(p)^2+|s(p)|_2^2\right)^2=w_{\mathbf{m}}(p)^2.
\]
In the basis above $D_{H_{\theta}}+\mathbf{m}$ is the direct sum of these $4\times4$ matrices, and $p=p'+\theta/n$ ranges over $P_{n,\theta}$ as $j$ ranges over $\{0,\dots,n-1\}^3$. This proves the product formula in \eqref{eq:renormalisation-flow-13}; the product is positive, because $w_{\mathbf{m}}\ge\mathbf{m}^2>0$.

\emph{(b) Summation over $P_{n,\theta}$.} Let $f$ be a smooth complex valued function on $\mathbb{R}^3$ that is $2\pi$-periodic in each coordinate, and let
\[
\hat{f}(x):=\frac{1}{(2\pi)^3}\int_{[-\pi,\pi]^3}\mathrm{d}^3p\,f(p)\cdot\exp(-i\,p\cdot x)\qquad\left(x\in\mathbb{Z}^3\right)
\]
be its Fourier coefficients. As $f$ is smooth, its Fourier series $f(p)=\sum_{x\in\mathbb{Z}^3}\hat{f}(x)\cdot\exp(i\,p\cdot x)$ converges absolutely. Summing a geometric progression in each $j_{\nu}$, one finds for $x\in\mathbb{Z}^3$
\[
\sum_{p\in P_{n,\theta}}\exp(i\,p\cdot x)=\begin{cases}n^3\cdot\exp\left(\frac{i\,(\pi+\theta)\cdot x}{n}\right),&x\in n\mathbb{Z}^3,\\0,&x\notin n\mathbb{Z}^3,\end{cases}
\]
where $\pi+\theta$ is the vector $(\pi+\theta_{\nu})_{\nu}$. Inserting the Fourier series and interchanging the two sums, one obtains the identity in
\begin{multline}\label{eq:renormalisation-flow-14}
\sum_{p\in P_{n,\theta}}f(p)=n^3\cdot\sum_{w\in\mathbb{Z}^3}\hat{f}(n\cdot w)\cdot\exp\left(i\,(\pi+\theta)\cdot w\right),\\
\left|\sum_{p\in P_{n,\theta}}f(p)-n^3\cdot\hat{f}(0)\right|\le n^3\cdot\sum_{w\in\mathbb{Z}^3\setminus\{0\}}\left|\hat{f}(n\cdot w)\right|,
\end{multline}
and the inequality follows from it, since $\left|\exp\left(i\,(\pi+\theta)\cdot w\right)\right|=1$. The function $\log w_{\mathbf{m}}$ is smooth and $2\pi$-periodic in each coordinate, because $w_{\mathbf{m}}\ge\mathbf{m}^2>0$. From now on $f:=\log w_{\mathbf{m}}$, so that \eqref{eq:renormalisation-flow-14} holds for this $f$.

\emph{(c) Integration over one coordinate.} A permutation of the coordinates of $p$ leaves $w_{\mathbf{m}}(p)$ unchanged; hence $\hat{f}(x)$ is unchanged by every permutation of the coordinates of $x$. Let $x\in\mathbb{Z}^3\setminus\{0\}$; in (c)--(e) one may therefore assume $|x_0|=|x|\ge1$, where $|x|=\max_{\nu}|x_{\nu}|$. Write $\mathsf{p}:=(p_1,p_2)$ and
\[
\mathcal{W}:=\sum_{\nu=1}^2\left(1-\cos p_{\nu}\right),\qquad\mathcal{S}:=\sum_{\nu=1}^2\sin^2p_{\nu},\qquad\acute{M}:=1+\mathbf{m}+\mathcal{W}.
\]
Then
\begin{multline}\label{eq:renormalisation-flow-15}
w_{\mathbf{m}}(p)=\left(\acute{M}-\cos p_0\right)^2+\sin^2p_0+\mathcal{S}=\acute{M}^2+1+\mathcal{S}-2\acute{M}\cdot\cos p_0\\
=\acute{M}\cdot\left(2\cosh\mathrm{En}-2\cos p_0\right),\qquad2\cosh\mathrm{En}(\mathsf{p}):=\frac{\acute{M}^2+1+\mathcal{S}}{\acute{M}},
\end{multline}
and the number $\mathrm{En}=\mathrm{En}(\mathsf{p})>0$ is well defined, because $\acute{M}>0$ and $\acute{M}^2+1+\mathcal{S}-2\acute{M}=(\mathbf{m}+\mathcal{W})^2+\mathcal{S}>0$. For $E>0$ one has $2\cosh E-2\cos p_0=\exp(E)\cdot\left|1-\exp(-E+i\,p_0)\right|^2$, and the power series of the logarithm gives
\[
\log\left(2\cosh E-2\cos p_0\right)=E-\sum_{j\ge1}\frac{\exp(-j\cdot E)}{j}\cdot\left(\exp(i\,j\cdot p_0)+\exp(-i\,j\cdot p_0)\right),
\]
with uniform convergence in $p_0\in\mathbb{R}$. Consequently, for every integer $x_0\ne0$,
\[
\frac{1}{2\pi}\int_{-\pi}^{\pi}\mathrm{d}p_0\,\log\left(2\cosh E-2\cos p_0\right)\cdot\exp(-i\,p_0\cdot x_0)=-\frac{\exp\left(-|x_0|\cdot E\right)}{|x_0|}.
\]
By \eqref{eq:renormalisation-flow-15}, $f(p)=\log\acute{M}+\log\left(2\cosh\mathrm{En}(\mathsf{p})-2\cos p_0\right)$, where the first term is independent of $p_0$. Integrating first over $p_0$ (Fubini's theorem), one obtains
\[
\hat{f}(x)=-\frac{1}{|x_0|}\cdot\frac{1}{(2\pi)^2}\int_{[-\pi,\pi]^2}\mathrm{d}^2\mathsf{p}\,\exp\left(-|x_0|\cdot\mathrm{En}(\mathsf{p})\right)\cdot\exp\left(-i\,(p_1\cdot x_1+p_2\cdot x_2)\right).
\]

\emph{(d) A lower bound for $\mathrm{En}$.} By \eqref{eq:renormalisation-flow-15},
\[
4\sinh^2\frac{\mathrm{En}}{2}=2\cosh\mathrm{En}-2=\frac{(\mathbf{m}+\mathcal{W})^2+\mathcal{S}}{1+\mathbf{m}+\mathcal{W}}.
\]
For $\mathsf{p}\in[-\pi,\pi]^2$ the numerator satisfies
\begin{align*}
(\mathbf{m}+\mathcal{W})^2+\mathcal{S}&\ge\mathbf{m}^2+\mathcal{W}^2+\mathcal{S}\ge\mathbf{m}^2+\sum_{\nu=1}^2\left(\left(1-\cos p_{\nu}\right)^2+\sin^2p_{\nu}\right)\\
&=\mathbf{m}^2+\sum_{\nu=1}^24\sin^2\frac{p_{\nu}}{2}\ge\mathbf{m}^2+\frac{4}{\pi^2}\cdot|\mathsf{p}|_2^2,
\end{align*}
the last step by $|\sin(p_{\nu}/2)|\ge|p_{\nu}|/\pi$ for $|p_{\nu}|\le\pi$; the denominator is at most $6$, because $\mathbf{m}\le1$ and $\mathcal{W}\le4$. Hence
\[
\sinh\frac{\mathrm{En}}{2}\ge\mathsf{y}:=\frac{1}{2\sqrt{6}}\cdot\sqrt{\mathbf{m}^2+\frac{4}{\pi^2}\cdot|\mathsf{p}|_2^2},
\]
and $\mathsf{y}\le\mathsf{y}_0:=\frac{3}{2\sqrt{6}}\approx0.6124$, because $\mathbf{m}\le1$ and $|\mathsf{p}|_2^2\le2\pi^2$. As $\operatorname{arcsinh}$ is concave on $[0,\infty)$ with $\operatorname{arcsinh}(0)=0$, and $0\le\mathsf{y}\le\mathsf{y}_0$,
\[
\operatorname{arcsinh}(\mathsf{y})\ge\mathsf{y}\cdot\frac{\operatorname{arcsinh}(\mathsf{y}_0)}{\mathsf{y}_0}\ge0.9461\cdot\mathsf{y}.
\]
Since $\sinh$ is increasing, $\mathrm{En}\ge2\,\operatorname{arcsinh}(\mathsf{y})$. With $\sqrt{a^2+b^2}\ge(a+b)/\sqrt{2}$ for $a,b\ge0$, and with $0.9461/\sqrt{6}>0.3862$, $4/\pi^2>0.4052$, $0.3862/\sqrt{2}>0.2730$ and $0.3862\cdot\sqrt{0.4052/2}>0.1738$, it follows that
\begin{multline}\label{eq:renormalisation-flow-16}
\mathrm{En}(\mathsf{p})\ge2\cdot0.9461\cdot\mathsf{y}\ge0.3862\cdot\sqrt{\mathbf{m}^2+0.4052\cdot|\mathsf{p}|_2^2}\\
\ge0.2730\cdot\mathbf{m}+0.1738\cdot|\mathsf{p}|_2\ge\frac{\mathbf{m}}{4}+\frac{|\mathsf{p}|_2}{6}\qquad\left(\mathsf{p}\in[-\pi,\pi]^2\right).
\end{multline}

\emph{(e) Decay of the coefficients.} Let $t:=|x|=|x_0|$. By (c) and \eqref{eq:renormalisation-flow-16},
\begin{align*}
\left|\hat{f}(x)\right|&\le\frac{1}{t}\cdot\exp\left(-\frac{\mathbf{m}\cdot t}{4}\right)\cdot\frac{1}{(2\pi)^2}\int_{\mathbb{R}^2}\mathrm{d}^2\mathsf{p}\,\exp\left(-\frac{t\cdot|\mathsf{p}|_2}{6}\right)\\
&=\frac{1}{t}\cdot\exp\left(-\frac{\mathbf{m}\cdot t}{4}\right)\cdot\frac{36}{2\pi t^2}\le\frac{6}{|x|^3}\cdot\exp\left(-\frac{\mathbf{m}\cdot|x|}{4}\right).
\end{align*}
Since $\hat{f}(x)$ does not change when the coordinates of $x$ are permuted, this bound holds for every $x\in\mathbb{Z}^3\setminus\{0\}$.

\emph{(f) Conclusion.} For every integer $t\ge1$,
\[
\#\left|\left\{w\in\mathbb{Z}^3:|w|=t\right\}\right|=(2t+1)^3-(2t-1)^3=24t^2+2\le26t^2.
\]
Applying (e) at $x=n\cdot w$, where $|n\cdot w|=n\cdot t$ if $|w|=t$, one obtains
\begin{align*}
n^3\cdot\sum_{w\in\mathbb{Z}^3\setminus\{0\}}\left|\hat{f}(n\cdot w)\right|&\le n^3\cdot\sum_{t\ge1}26t^2\cdot\frac{6}{(n\cdot t)^3}\cdot\exp\left(-\frac{\mathbf{m}\cdot n\cdot t}{4}\right)\\
&=156\sum_{t\ge1}\frac{1}{t}\cdot\exp\left(-\frac{\mathbf{m}\cdot n\cdot t}{4}\right)=-156\log\left(1-\exp\left(-\frac{\mathbf{m}\cdot n}{4}\right)\right).
\end{align*}
By (a), $\log\det\left(D_{H_{\theta}}+\mathbf{m}\right)=2\sum_{p\in P_{n,\theta}}f(p)$, and $p^0(\mathbf{m})=2\,\hat{f}(0)$ by the definition of $p^0$. The inequality in \eqref{eq:renormalisation-flow-14} therefore gives
\begin{align*}
\left|\log\det\left(D_{H_{\theta}}+\mathbf{m}\right)-p^0(\mathbf{m})\cdot n^3\right|&=2\left|\sum_{p\in P_{n,\theta}}f(p)-n^3\cdot\hat{f}(0)\right|\\
&\le-312\log\left(1-\exp\left(-\frac{\mathbf{m}\cdot n}{4}\right)\right),
\end{align*}
which is the bound in \eqref{eq:renormalisation-flow-13}.
\end{proof}

\begin{corollary}[Twist ratios and the free partition function]\label{cor:twist-ratios-free}
Let $n\ge1$ be an integer, with $H_{\theta}$ understood as in Lemma~\ref{lem:product-formula-volume} when $n\le2$, let $0<\mathbf{m}\le1$ and $q>0$. For every $\theta\in\mathbb{R}^3$,
\[
\left|\log\frac{\det\left(D_{H_{\theta}}+\mathbf{m}\right)}{\det\left(D_0+\mathbf{m}\right)}\right|\le2\,\varepsilon_5(\mathbf{m}\cdot n).
\]
If, in addition, $\mathbf{m}=\bar{m}\cdot\eta$ and $n=R/\eta$, so that $\mathbf{m}\cdot n=\bar{m}\cdot R$, then the number $\rho^{\mathrm{free}}_n$ of \eqref{eq:renormalisation-flow-1} satisfies
\[
\left|\log\rho^{\mathrm{free}}_n-p^0(\bar{m}\cdot\eta)\cdot n^3\right|\le\varepsilon_5(\bar{m}\cdot R).
\]
\end{corollary}

\begin{proof}
At $\theta=0$ one has $H_0=0$, so that $D_{H_0}=D_0$. By \eqref{eq:renormalisation-flow-13}, applied at $\theta$ and at $0$, each of the numbers $\log\det\left(D_{H_{\theta}}+\mathbf{m}\right)$ and $\log\det\left(D_0+\mathbf{m}\right)$ differs from $p^0(\mathbf{m})\cdot n^3$ by at most $\varepsilon_5(\mathbf{m}\cdot n)$; this gives the first inequality. In the second case, \eqref{eq:renormalisation-flow-1} exhibits $\rho^{\mathrm{free}}_n$ as the mean, with respect to the probability measure $(2\pi)^{-3}\,\mathrm{d}\theta$ on $[-\pi,\pi]^3$, of the continuous function $\theta\mapsto\det\left(D_{H_{\theta}}+\bar{m}\cdot\eta\right)$. By \eqref{eq:renormalisation-flow-13} its values lie in the interval
\[
\left[\exp\left(p^0(\bar{m}\cdot\eta)\cdot n^3-\varepsilon_5(\bar{m}\cdot R)\right),\ \exp\left(p^0(\bar{m}\cdot\eta)\cdot n^3+\varepsilon_5(\bar{m}\cdot R)\right)\right],
\]
and so does their mean $\rho^{\mathrm{free}}_n$.
\end{proof}

\subsubsection{Averaging over blocks at a flat connection}\label{sec:averaging-blocks-flat}

Throughout this subsection $l\ge 1$ is an integer, $\mathfrak{N}\ge 3$ is an odd integer, $\boldsymbol{\eta}:=1/\mathfrak{N}$, $0<\mathfrak{m}\le 1$, $\mathsf{b}>0$ and $\mathsf{e}>0$. The fine lattice is $\boldsymbol{\eta}\mathbb{Z}^3/l\mathbb{Z}^3$ and the unit lattice is $\mathbb{Z}^3/l\mathbb{Z}^3$; the points of $\boldsymbol{\eta}\mathbb{Z}^3$ and of $\mathbb{Z}^3$ are called fine sites and unit sites. For $x\in\mathbb{R}^3$, $|x|$ is the sup norm and $|x|_2$ the Euclidean norm. A spinor on the fine lattice is a function $f\colon\boldsymbol{\eta}\mathbb{Z}^3\to\mathbb{C}^4$ with $f(x+l e_{\nu})=-f(x)$ for $\nu=0,1,2$; a spinor on the unit lattice is a function $g\colon\mathbb{Z}^3\to\mathbb{C}^4$ with $g(y+l e_{\nu})=-g(y)$ for $\nu=0,1,2$. The scalar product of two spinors on the fine lattice is $\boldsymbol{\eta}^3\sum_x$ of their pointwise scalar product in $\mathbb{C}^4$, and that of two spinors on the unit lattice is $\sum_y$ of it, where $x$ and $y$ range over the sites of the fine and of the unit lattice; adjoints refer to these scalar products. For a linear map $X$ between two such spaces, or between two spaces $\mathbb{C}^d$, $|X|$ denotes the operator norm; if $X$ maps a space into itself, then $\operatorname{Re}X:=\frac{1}{2}\left(X+X^{*}\right)$, and $\operatorname{Re}X\ge t$, for $t\in\mathbb{R}$, means that $\operatorname{Re}\langle f,Xf\rangle\ge t\,|f|^2$ for every $f$.

For $y\in\mathbb{Z}^3$ the block $B(y)$ is the set of the $\mathfrak{N}^3$ fine sites $x$ with $|x-y|<\frac{1}{2}$. Since $\mathfrak{N}$ is odd, every fine site $x$ lies in exactly one block; the centre of this block is denoted $y(x)$.

Let $\theta\in\mathbb{R}^3$, and let $c$ be the constant real connection on the fine lattice with
\[
\mathsf{e}\cdot c_{\nu}\,l=\theta_{\nu}\qquad(\nu=0,1,2),
\]
that is, $c(x,x+\boldsymbol{\eta}\cdot e_{\nu})=c_{\nu}$ for every fine site $x$, with these real numbers $c_{\nu}$. For $x\in\mathbb{R}^3$ put $c\cdot x:=\sum_{\nu}c_{\nu}x_{\nu}$. Let $D_c$ be the operator $D_A$ of Section~\ref{sec:exact-fermion-operators} at $A=c$, with $\boldsymbol{\eta}$ in place of $\eta_k$ and $\mathsf{e}$ in place of $e_k$; it acts on the spinors on the fine lattice. The operator $Q(c)$ from the spinors on the fine lattice to the spinors on the unit lattice, the operator $Q^T(-c)$ in the opposite direction, the operator $T_c$ on the spinors on the fine lattice, the operator $D(c)$ on the spinors on the unit lattice, and the number $R_{\theta}$ are defined by
\begin{gather*}
\left(Q(c)f\right)(y):=\boldsymbol{\eta}^3\sum_{x\in B(y)}\exp\left(i\,\mathsf{e}\cdot c\cdot(x-y)\right)f(x),\\
\left(Q^T(-c)g\right)(x):=\exp\left(-i\,\mathsf{e}\cdot c\cdot\left(x-y(x)\right)\right)g\left(y(x)\right),\\
T_c:=D_c+\mathfrak{m}+\mathsf{b}\cdot Q^T(-c)\,Q(c),\qquad D(c):=\mathsf{b}-\mathsf{b}^2\cdot Q(c)\,T_c^{-1}\,Q^T(-c),\qquad R_{\theta}:=\frac{\det D(c)}{\det D(0)}.
\end{gather*}
Here $T_0$ and $D(0)$ denote $T_c$ and $D(c)$ at $c=0$, that is, at $\theta=0$. It is shown below that $T_c$ and $D(c)$ are invertible for every $\theta$; so $D(c)$ and $R_{\theta}$ are well defined.

Let $\mathfrak{N}=L^k$ with an integer $k\ge 1$, so that $\boldsymbol{\eta}=\eta_k$, and let $\mathsf{e}=e_k$ and $\mathsf{b}=b_k$, with $e_k$ and $b_k$ as in Section~\ref{sec:running-quantities-tori}. Since $c$ is constant, the sum of $c$ over any path of bonds of the fine lattice from $y$ to $x$, in the sense of Section~\ref{sec:averaging-fluctuations-gauge}, equals $c\cdot(x-y)/\boldsymbol{\eta}$; hence $\boldsymbol{\eta}\cdot(\tau_kc)(y,x)=c\cdot(x-y)$, with $\tau_k$ as in Section~\ref{sec:exact-fermion-operators}. Therefore $Q(c)$, $Q^T(-c)$ and $D(c)$ are the operators $Q_k(c)$, $Q_k^T(-c)$ and $D_k(c)$ of Section~\ref{sec:exact-fermion-operators} on the torus of side $l$ with the mass $\mathfrak{m}$, and $S_k(c)=T_c^{-1}$. In particular, if
\[
\left(l,\mathfrak{N},\mathfrak{m},\mathsf{b},\mathsf{e}\right)=\left(l_{*},L^{k^{*}},\mathfrak{m}_{*},b_{k^{*}},e_{*}\right),
\]
the subscript $*$ referring to the stop level $k^{*}=k^{*}(R)$ of Definition~\ref{def:tori-stop-level}, then $c$ is the twist $c_{\theta}$ of Definition~\ref{def:twist-external-connection}, whose components are $\theta_{\nu}/(e_{*}l_{*})$, and $D(c)=D_{*}(c_{\theta})$.

\paragraph{Invertibility.} The operators $D_c+\mathfrak{m}$, $T_c$ and $D(c)$ are invertible, and
\begin{equation}\label{eq:renormalisation-flow-17}
D(c)^{-1}=\frac{1}{\mathsf{b}}+Q(c)\left(D_c+\mathfrak{m}\right)^{-1}Q^T(-c).
\end{equation}

\begin{proof}
For $\nu=0,1,2$ let $T_{\nu}$ be the covariant shift
\[
\left(T_{\nu}f\right)(x):=\exp\left(i\,\mathsf{e}\cdot\boldsymbol{\eta}\cdot c_{\nu}\right)f(x+\boldsymbol{\eta}\cdot e_{\nu})
\]
on the spinors on the fine lattice. It is unitary, and $\left(T_{\nu}^{*}f\right)(x)=\exp\left(-i\,\mathsf{e}\cdot\boldsymbol{\eta}\cdot c_{\nu}\right)f(x-\boldsymbol{\eta}\cdot e_{\nu})$. Since $c(x,x\pm\boldsymbol{\eta}\cdot e_{\nu})=\pm c_{\nu}$, the display of $D_A$ in Section~\ref{sec:exact-fermion-operators} becomes
\[
D_c=\frac{1}{\boldsymbol{\eta}}\sum_{\nu}\left[\frac{1}{2}\gamma_{\nu}\left(T_{\nu}-T_{\nu}^{*}\right)+\frac{1}{2}\left(2-T_{\nu}-T_{\nu}^{*}\right)\right].
\]
The matrices $\gamma_{\nu}$ are Hermitian (Section~\ref{sec:dirac-matrices-lattice}) and commute with $T_{\nu}$, which acts on the position only; hence the operators $\gamma_{\nu}\left(T_{\nu}-T_{\nu}^{*}\right)$ are anti-Hermitian, whereas $\left\langle f,\left(2-T_{\nu}-T_{\nu}^{*}\right)f\right\rangle=2|f|^2-2\operatorname{Re}\langle f,T_{\nu}f\rangle\ge 0$. Therefore $\operatorname{Re}\left(D_c+\mathfrak{m}\right)\ge\mathfrak{m}$. By the definitions, $Q^T(-c)=Q(c)^{*}$, and $Q(c)\,Q(c)^{*}=1$, because $y(x)=y$ for $x\in B(y)$ and $B(y)$ has $\mathfrak{N}^3=\boldsymbol{\eta}^{-3}$ sites; in particular $|Q(c)|=|Q(c)^{*}|=1$.

For every spinor $f$ on the fine lattice, $\mathfrak{m}\cdot|f|^2\le\operatorname{Re}\left\langle f,\left(D_c+\mathfrak{m}\right)f\right\rangle\le|f|\cdot\left|\left(D_c+\mathfrak{m}\right)f\right|$. Hence $G_c:=\left(D_c+\mathfrak{m}\right)^{-1}$ exists, and $|G_c|\le 1/\mathfrak{m}$. Multiplying the identity $\operatorname{Re}\left(D_c+\mathfrak{m}\right)=\frac{1}{2}\left(G_c^{-1}+\left(G_c^{-1}\right)^{*}\right)$ by $G_c^{*}$ from the left and by $G_c$ from the right gives
\[
\operatorname{Re}G_c=G_c^{*}\operatorname{Re}\left(D_c+\mathfrak{m}\right)G_c\ge 0.
\]
As $\mathsf{b}\cdot Q(c)^{*}Q(c)\ge 0$, one has $\operatorname{Re}T_c\ge\mathfrak{m}$, and $T_c$ is invertible by the same argument. The operator $S_c:=Q(c)\,G_c\,Q(c)^{*}$ on the spinors on the unit lattice satisfies $\operatorname{Re}S_c=Q(c)\left(\operatorname{Re}G_c\right)Q(c)^{*}\ge 0$ and $|S_c|\le 1/\mathfrak{m}$; in particular $\operatorname{Re}\left(\mathsf{b}^{-1}+S_c\right)\ge\mathsf{b}^{-1}$, so that $\mathsf{b}^{-1}+S_c$ is invertible.

Write $Q:=Q(c)$. Since $T_c=G_c^{-1}+\mathsf{b}\cdot Q^{*}Q$ and $Q\,G_c\,Q^{*}=S_c$,
\begin{align*}
\left(G_c^{-1}+\mathsf{b}\cdot Q^{*}Q\right)&\left(G_c-G_cQ^{*}\left(\mathsf{b}^{-1}+S_c\right)^{-1}QG_c\right)\\
&=1+Q^{*}\left[\mathsf{b}-\left(1+\mathsf{b}\cdot S_c\right)\left(\mathsf{b}^{-1}+S_c\right)^{-1}\right]QG_c=1,
\end{align*}
the bracket being zero because $1+\mathsf{b}\cdot S_c=\mathsf{b}\cdot\left(\mathsf{b}^{-1}+S_c\right)$. Consequently
\begin{gather*}
T_c^{-1}=G_c-G_cQ^{*}\left(\mathsf{b}^{-1}+S_c\right)^{-1}QG_c,\\
Q\,T_c^{-1}Q^{*}=S_c-S_c\left(\mathsf{b}^{-1}+S_c\right)^{-1}S_c=\mathsf{b}^{-1}\cdot S_c\left(\mathsf{b}^{-1}+S_c\right)^{-1},
\end{gather*}
where $S_c$ commutes with $\left(\mathsf{b}^{-1}+S_c\right)^{-1}$, and
\[
D(c)=\mathsf{b}-\mathsf{b}\cdot S_c\left(\mathsf{b}^{-1}+S_c\right)^{-1}=\left(\mathsf{b}^{-1}+S_c\right)^{-1}.
\]
So $D(c)$ is invertible, and $D(c)^{-1}=\mathsf{b}^{-1}+S_c$; this is \eqref{eq:renormalisation-flow-17}.
\end{proof}

\paragraph{Momentum space.} Let
\[
\mathfrak{M}:=\left\{-\frac{\mathfrak{N}-1}{2},\dots,\frac{\mathfrak{N}-1}{2}\right\}^3.
\]
For $z\in\mathbb{C}$ and $q=(q_0,q_1,q_2)\in\mathbb{C}^3$ put
\begin{gather*}
u_1(z):=\frac{1}{\mathfrak{N}}\sum_{g\in\mathbb{Z},\ |g|\le\frac{\mathfrak{N}-1}{2}}\exp\left(-\frac{izg}{\mathfrak{N}}\right),\qquad\hat{u}(q):=\prod_{\nu}u_1(q_{\nu}),\\
\hat{D}(q):=\frac{1}{\boldsymbol{\eta}}\sum_{\nu}\left[i\gamma_{\nu}\sin\left(\boldsymbol{\eta}\cdot q_{\nu}\right)+1-\cos\left(\boldsymbol{\eta}\cdot q_{\nu}\right)\right],\\
a(q):=\mathfrak{m}+\frac{1}{\boldsymbol{\eta}}\sum_{\nu}\left(1-\cos\left(\boldsymbol{\eta}\cdot q_{\nu}\right)\right),\qquad s_{\nu}(q):=\frac{\sin\left(\boldsymbol{\eta}\cdot q_{\nu}\right)}{\boldsymbol{\eta}},\qquad w(q):=a(q)^2+\sum_{\nu}s_{\nu}(q)^2.
\end{gather*}
Summation of the geometric series gives
\[
u_1(z)=\sin\left(\frac{z}{2}\right)\cdot\left(\mathfrak{N}\cdot\sin\left(\frac{z}{2\mathfrak{N}}\right)\right)^{-1}\qquad\text{whenever }\sin\left(\frac{z}{2\mathfrak{N}}\right)\ne 0,
\]
and $w(q)=\boldsymbol{\eta}^{-2}\cdot w_{\boldsymbol{\eta}\cdot\mathfrak{m}}(\boldsymbol{\eta}\cdot q)$, with the function $w_{\mathbf{m}}$ of Section~\ref{sec:constant-twists-free}, its defining formula being read for complex arguments. For $s\in\mathbb{C}^3$ put $\gamma\cdot s:=\sum_{\nu}\gamma_{\nu}s_{\nu}$, and let $s(q):=\left(s_{\nu}(q)\right)_{\nu}$; then $\hat{D}(q)+\mathfrak{m}=a(q)+i\gamma\cdot s(q)$. By the relations of Section~\ref{sec:dirac-matrices-lattice}, the matrix $\gamma\cdot s$ is traceless and $(\gamma\cdot s)^2=\sum_{\nu}s_{\nu}^2$. So its characteristic polynomial is $\left(\lambda^2-\sum_{\nu}s_{\nu}^2\right)^2$: if $\sum_{\nu}s_{\nu}^2\ne 0$, the matrix $\gamma\cdot s$ is diagonalisable with eigenvalues $\pm\left(\sum_{\nu}s_{\nu}^2\right)^{\frac{1}{2}}$, which occur equally often since the trace vanishes; otherwise it is nilpotent. With $a=a(q)$ and $s=s(q)$ it follows that
\[
(a+i\gamma\cdot s)(a-i\gamma\cdot s)=w(q),\qquad\det\left(\hat{D}(q)+\mathfrak{m}\right)=w(q)^2.
\]
For a momentum $p\in\mathbb{C}^3$ and for $m\in\mathfrak{M}$ put
\[
q_m:=p+2\pi m,\qquad R_m:=\hat{D}(q_m)+\mathfrak{m},\qquad\hat{u}_m:=\hat{u}(q_m);
\]
these quantities depend on $p$.

\begin{lemma}[Fourier decomposition]\label{lem:fourier-decomposition}
Let $\theta\in\mathbb{R}^3$, let $c$, $D(c)$, $T_c$ and $R_{\theta}$ be as above, and let $P_{l,\theta}$ be the set of momenta of Section~\ref{sec:constant-twists-free} with $n=l$.
\begin{enumerate}
\item[(a)] For $p\in\mathbb{R}^3$ the matrices $R_m$, $m\in\mathfrak{M}$, are invertible; put
\[
\mathrm{Ch}(p):=\frac{1}{\mathsf{b}}+G(p),\qquad G(p):=\sum_{m\in\mathfrak{M}}\hat{u}_m^2\cdot R_m^{-1}.
\]
Then $D(c)^{-1}$ is unitarily equivalent to the direct sum $\bigoplus_{p\in P_{l,\theta}}\mathrm{Ch}(p)$ of the $4\times 4$ matrices $\mathrm{Ch}(p)$. For $p\in\mathbb{R}^3$ one has $\mathrm{Ch}(p)=\boldsymbol{\alpha}-i\gamma\cdot\boldsymbol{\beta}$ with a number $\boldsymbol{\alpha}\ge 1/\mathsf{b}$ and a vector $\boldsymbol{\beta}\in\mathbb{R}^3$; moreover $\operatorname{Re}\mathrm{Ch}(p)\ge 1/\mathsf{b}$ and $\det\mathrm{Ch}(p)=\left(\boldsymbol{\alpha}^2+|\boldsymbol{\beta}|_2^2\right)^2>0$. Consequently
\[
\det D(c)=\prod_{p\in P_{l,\theta}}\det\mathrm{Ch}(p)^{-1}>0,
\]
and $R_{\theta}>0$.
\item[(b)] Let $T(p)$ be the operator
\[
T(p):=\operatorname{diag}_m\left(R_m\right)+\mathsf{b}\cdot|\hat{u}\rangle\langle\hat{u}|\otimes\mathbf{1}_4
\]
on $\mathbb{C}^{\mathfrak{M}}\otimes\mathbb{C}^4$, where $\operatorname{diag}_m(R_m)$ maps $\delta_m\otimes v$ to $\delta_m\otimes R_mv$, the vectors $\delta_m$ being the unit vectors of $\mathbb{C}^{\mathfrak{M}}$, the matrix $|\hat{u}\rangle\langle\hat{u}|$ is $\left(\hat{u}_m\cdot\hat{u}_{m'}\right)_{m,m'\in\mathfrak{M}}$, and $\mathbf{1}_4$ is the unit matrix of $\mathbb{C}^4$. Then $T_c$ is unitarily equivalent to the direct sum $\bigoplus_{p\in P_{l,\theta}}T(p)$, and for $p\in\mathbb{R}^3$
\[
\det T(p)=\mathsf{b}^4\cdot\det\mathrm{Ch}(p)\cdot\prod_{m\in\mathfrak{M}}w(q_m)^2>0.
\]
\item[(c)] Let $n:=l\cdot\mathfrak{N}$ and $\mathbf{m}:=\boldsymbol{\eta}\cdot\mathfrak{m}$, so that $\mathbf{m}\cdot n=\mathfrak{m}\cdot l$. Then
\[
R_{\theta}=\frac{\det\left(D_{H_{\theta}}+\mathbf{m}\right)}{\det\left(D_{H_0}+\mathbf{m}\right)}\cdot\frac{\det T_0}{\det T_c},
\]
where the first factor is the ratio of determinants of Section~\ref{sec:constant-twists-free}, formed on $\Lambda_n$ at the mass $\mathbf{m}$; here $H_0=0$.
\end{enumerate}
\end{lemma}

\begin{proof}
\emph{Conjugation by a phase.} Let $U$ be the multiplication by the function $x\mapsto\exp\left(i\,\mathsf{e}\cdot c\cdot x\right)$ on functions on $\boldsymbol{\eta}\mathbb{Z}^3$, and $U_0$ the multiplication by its restriction to $\mathbb{Z}^3$ on functions on $\mathbb{Z}^3$. Let $D_0$, $Q:=Q(0)$ and $Q^T:=Q^T(0)$ be the operators $D_c$, $Q(c)$ and $Q^T(-c)$ at $c=0$, read through their defining formulas as translation invariant operators: $D_0$ acting on all $\mathbb{C}^4$-valued functions on $\boldsymbol{\eta}\mathbb{Z}^3$, $Q$ mapping these to $\mathbb{C}^4$-valued functions on $\mathbb{Z}^3$, and $Q^T$ in the opposite direction. Inserting the definitions, one finds
\[
Q(c)=U_0^{-1}\,Q\,U,\qquad Q^T(-c)=U^{-1}\,Q^T\,U_0,\qquad D_c=U^{-1}\,D_0\,U;
\]
for the last identity note that $\exp\left(-i\,\mathsf{e}\cdot c\cdot x\right)\exp\left(i\,\mathsf{e}\cdot c\cdot(x\pm\boldsymbol{\eta}\cdot e_{\nu})\right)=\exp\left(\pm i\,\mathsf{e}\cdot\boldsymbol{\eta}\cdot c_{\nu}\right)$ is the link phase of $D_c$ on the bond from $x$ to $x\pm\boldsymbol{\eta}\cdot e_{\nu}$. Let $\mathscr{H}_{\theta}$ be the space of the $\mathbb{C}^4$-valued functions $g$ on $\boldsymbol{\eta}\mathbb{Z}^3$ (the fine part of $\mathscr{H}_{\theta}$) and on $\mathbb{Z}^3$ (its unit-lattice part) such that
\[
g(x+l e_{\nu})=-\exp\left(i\theta_{\nu}\right)g(x)\qquad(\nu=0,1,2),
\]
with the scalar products $\boldsymbol{\eta}^3\sum_x$, resp.\ $\sum_y$, of the pointwise scalar products, summed over the points of one cell of periods of $l\mathbb{Z}^3$. Since $\mathsf{e}\cdot c_{\nu}\,l=\theta_{\nu}$, the operator $U$ maps the spinors on the fine lattice unitarily onto the fine part of $\mathscr{H}_{\theta}$, and $U_0$ maps the spinors on the unit lattice unitarily onto its unit-lattice part. The operators $D_0$, $Q$ and $Q^T$ commute with the translations by the vectors of $\mathbb{Z}^3$; so they preserve the relations that define $\mathscr{H}_{\theta}$ and act between its parts. Hence
\begin{gather*}
T_c=U^{-1}\left(D_0+\mathfrak{m}+\mathsf{b}\cdot Q^TQ\right)U,\\
D(c)=U_0^{-1}\left(\mathsf{b}-\mathsf{b}^2\cdot Q\left(D_0+\mathfrak{m}+\mathsf{b}\cdot Q^TQ\right)^{-1}Q^T\right)U_0,
\end{gather*}
the operators in the outer parentheses acting in $\mathscr{H}_{\theta}$, and, by \eqref{eq:renormalisation-flow-17},
\[
U_0\,D(c)^{-1}\,U_0^{-1}=\frac{1}{\mathsf{b}}+Q\left(D_0+\mathfrak{m}\right)^{-1}Q^T
\]
on the unit-lattice part of $\mathscr{H}_{\theta}$.

\emph{Real momenta.} Let $q\in\mathbb{R}^3$ and $R:=\hat{D}(q)+\mathfrak{m}$. Then $a(q)\ge\mathfrak{m}>0$ and $w(q)\ge a(q)^2>0$, so that, by the identity $(a+i\gamma\cdot s)(a-i\gamma\cdot s)=w(q)$, the matrix $R$ is invertible and
\[
R^{-1}=\frac{a(q)-i\gamma\cdot s(q)}{w(q)}.
\]
The function $u_1$ is real on $\mathbb{R}$, its defining sum being invariant under $g\mapsto-g$. Hence, for $p\in\mathbb{R}^3$, the matrices $R_m$ are invertible and the numbers $\hat{u}_m$ are real.

\emph{Fourier bases.} For $q\in\mathbb{R}^3$ let $\varphi_q$ be the function $x\mapsto\exp(iq\cdot x)$, on $\boldsymbol{\eta}\mathbb{Z}^3$ as well as on $\mathbb{Z}^3$, and let $v$ range over an orthonormal basis of $\mathbb{C}^4$. The functions $l^{-\frac{3}{2}}\varphi_pv$ on $\mathbb{Z}^3$, $p\in P_{l,\theta}$, form an orthonormal basis of the unit-lattice part of $\mathscr{H}_{\theta}$, and the functions $l^{-\frac{3}{2}}\varphi_{q_m}v$ on $\boldsymbol{\eta}\mathbb{Z}^3$, $p\in P_{l,\theta}$, $m\in\mathfrak{M}$, form one of its fine part. Indeed, they lie in $\mathscr{H}_{\theta}$, because $\exp\left(ip_{\nu}l\right)=-\exp\left(i\theta_{\nu}\right)$ for $p\in P_{l,\theta}$ and $q_m-p\in 2\pi\mathbb{Z}^3$; they have norm one, because a cell of periods contains $l^3$ unit sites and $l^3\cdot\mathfrak{N}^3$ fine sites; and their numbers, $4l^3$ and $4l^3\cdot\mathfrak{N}^3$, are the dimensions of the two parts. As to orthogonality, let $p\in P_{l,\theta}$ be given, as in Section~\ref{sec:constant-twists-free}, by $p_{\nu}=\left(2\pi j_{\nu}+\pi+\theta_{\nu}\right)/l$ with $j\in\{0,\dots,l-1\}^3$. Then the components of $q_m$ are $q_{m,\nu}=\left(2\pi\left(j_{\nu}+l\cdot m_{\nu}\right)+\pi+\theta_{\nu}\right)/l$, and the integers $j_{\nu}+l\cdot m_{\nu}$ with $0\le j_{\nu}<l$ and $|m_{\nu}|\le\frac{\mathfrak{N}-1}{2}$ are $l\cdot\mathfrak{N}$ consecutive integers. Hence two distinct pairs $(p,m)$ give momenta $q_m$ whose difference lies in $\frac{2\pi}{l}\mathbb{Z}^3$ but not in $2\pi\cdot\mathfrak{N}\,\mathbb{Z}^3$, and two distinct elements of $P_{l,\theta}$ differ by an element of $\frac{2\pi}{l}\mathbb{Z}^3$ that does not lie in $2\pi\mathbb{Z}^3$. For $q\in\frac{2\pi}{l}\mathbb{Z}^3$ the function $\varphi_q$ is a character of $\boldsymbol{\eta}\mathbb{Z}^3/l\mathbb{Z}^3$ and of $\mathbb{Z}^3/l\mathbb{Z}^3$; its sum over the fine sites of a cell of periods vanishes unless $q\in 2\pi\cdot\mathfrak{N}\,\mathbb{Z}^3$, and its sum over the unit sites vanishes unless $q\in 2\pi\mathbb{Z}^3$. So two different functions of either family are orthogonal; if they have the same momentum, their vectors $v$ are orthogonal.

\emph{Proof of (a) and (b).} Let $p\in P_{l,\theta}$ and $v\in\mathbb{C}^4$. For $x$ in a block $B(y)$,
\[
\left(Q^T\varphi_pv\right)(x)=\exp(ip\cdot y)\,v=\varphi_p(x)\,g(x)\,v,\qquad g(x):=\exp\left(-ip\cdot\left(x-y(x)\right)\right).
\]
The function $g$ on $\boldsymbol{\eta}\mathbb{Z}^3$ has period one in each direction, because $y(x+z)=y(x)+z$ for $z\in\mathbb{Z}^3$. The characters of the finite group $\boldsymbol{\eta}\mathbb{Z}^3/\mathbb{Z}^3$ are the functions $\varphi_{2\pi m}$, $m\in\mathfrak{M}$, and $B(0)$ is a set of representatives of this group. Expanding $g$ in these characters (Fourier analysis on a finite abelian group), one has $g=\sum_{m\in\mathfrak{M}}\hat{g}(m)\,\varphi_{2\pi m}$ with
\[
\hat{g}(m)=\boldsymbol{\eta}^3\sum_{x\in B(0)}g(x)\exp\left(-2\pi i\,m\cdot x\right)=\boldsymbol{\eta}^3\sum_{x\in B(0)}\exp\left(-iq_m\cdot x\right)=\hat{u}_m;
\]
the last equality holds because $B(0)$ consists of the points $\boldsymbol{\eta}\cdot z$ with $z\in\mathbb{Z}^3$ and $|z_{\nu}|\le\frac{\mathfrak{N}-1}{2}$, so that the sum is the product over $\nu$ of the values of $u_1$ at the components of $q_m$. Therefore
\[
Q^T\varphi_pv=\sum_{m\in\mathfrak{M}}\hat{u}_m\cdot\varphi_{q_m}v.
\]
Moreover, for $q\in\mathbb{R}^3$ the display of $D_A$ in Section~\ref{sec:exact-fermion-operators}, at $A=0$, gives
\[
\left(D_0+\mathfrak{m}\right)\varphi_qv=\varphi_q\left(\hat{D}(q)+\mathfrak{m}\right)v,
\]
since
\[
\frac{1}{2}\left(1-\gamma_{\nu}\right)\exp\left(i\,\boldsymbol{\eta}\cdot q_{\nu}\right)+\frac{1}{2}\left(1+\gamma_{\nu}\right)\exp\left(-i\,\boldsymbol{\eta}\cdot q_{\nu}\right)=\cos\left(\boldsymbol{\eta}\cdot q_{\nu}\right)-i\gamma_{\nu}\sin\left(\boldsymbol{\eta}\cdot q_{\nu}\right);
\]
and
\[
Q\,\varphi_{q_m}v=\hat{u}_m\cdot\varphi_pv\qquad(m\in\mathfrak{M}),
\]
because $B(y)=y+B(0)$, the set $B(0)$ is symmetric under $x\mapsto-x$, and $\varphi_{q_m}=\varphi_p$ on $\mathbb{Z}^3$.

For $p\in P_{l,\theta}$ let $F_p$ be the span of the functions $\varphi_{q_m}v$, $m\in\mathfrak{M}$, $v\in\mathbb{C}^4$, and identify $l^{-\frac{3}{2}}\varphi_{q_m}v$ with $\delta_m\otimes v\in\mathbb{C}^{\mathfrak{M}}\otimes\mathbb{C}^4$. By the formulas displayed above for $Q^T\varphi_pv$, for $\left(D_0+\mathfrak{m}\right)\varphi_qv$ at $q=q_m$, and for $Q\,\varphi_{q_m}v$, the operators $D_0+\mathfrak{m}$ and $Q^TQ$ map $F_p$ into itself and act there as $\operatorname{diag}_m(R_m)$ and as $|\hat{u}\rangle\langle\hat{u}|\otimes\mathbf{1}_4$; indeed $Q^TQ\,\varphi_{q_m}v=\hat{u}_m\sum_{m'\in\mathfrak{M}}\hat{u}_{m'}\cdot\varphi_{q_{m'}}v$. The fine part of $\mathscr{H}_{\theta}$ is the orthogonal sum of the spaces $F_p$, $p\in P_{l,\theta}$; by the conjugation formula for $T_c$ above, $T_c$ is therefore unitarily equivalent to $\bigoplus_{p\in P_{l,\theta}}T(p)$. Since the matrices $R_m$ are invertible, $\left(D_0+\mathfrak{m}\right)^{-1}\varphi_{q_m}v=\varphi_{q_m}R_m^{-1}v$, and therefore
\[
\left(\frac{1}{\mathsf{b}}+Q\left(D_0+\mathfrak{m}\right)^{-1}Q^T\right)\varphi_pv=\varphi_p\left(\frac{1}{\mathsf{b}}+\sum_{m\in\mathfrak{M}}\hat{u}_m^2\cdot R_m^{-1}\right)v=\varphi_p\cdot\mathrm{Ch}(p)\,v.
\]
By the conjugation formula for $D(c)^{-1}$ above and the orthonormal basis of the unit-lattice part of $\mathscr{H}_{\theta}$, $D(c)^{-1}$ is unitarily equivalent to $\bigoplus_{p\in P_{l,\theta}}\mathrm{Ch}(p)$.

Now let $p\in\mathbb{R}^3$. By the formula for $R^{-1}$ and the reality of the numbers $\hat{u}_m$,
\[
\mathrm{Ch}(p)=\boldsymbol{\alpha}-i\gamma\cdot\boldsymbol{\beta},\qquad\boldsymbol{\alpha}:=\frac{1}{\mathsf{b}}+\sum_{m\in\mathfrak{M}}\hat{u}_m^2\cdot\frac{a(q_m)}{w(q_m)}\ge\frac{1}{\mathsf{b}},\qquad\boldsymbol{\beta}:=\sum_{m\in\mathfrak{M}}\hat{u}_m^2\cdot\frac{s(q_m)}{w(q_m)}\in\mathbb{R}^3.
\]
As $\gamma\cdot\boldsymbol{\beta}$ is Hermitian, $\operatorname{Re}\mathrm{Ch}(p)=\boldsymbol{\alpha}\ge 1/\mathsf{b}$. The computation of $\det\left(\hat{D}(q)+\mathfrak{m}\right)$ above, with $\boldsymbol{\alpha}$ in place of $a$ and $-\boldsymbol{\beta}$ in place of $s$, gives $\det\mathrm{Ch}(p)=\left(\boldsymbol{\alpha}^2+|\boldsymbol{\beta}|_2^2\right)^2>0$. By the unitary equivalence, $\det D(c)^{-1}=\prod_{p\in P_{l,\theta}}\det\mathrm{Ch}(p)$; so $\det D(c)>0$ for every $\theta$, in particular for $\theta=0$, and hence $R_{\theta}>0$. This proves (a).

The operator $\mathsf{b}\cdot|\hat{u}\rangle\langle\hat{u}|\otimes\mathbf{1}_4$ factors through $\mathbb{C}^4$: it is the map $\sum_m\delta_m\otimes v_m\mapsto\sum_m\hat{u}_m\cdot v_m$ from $\mathbb{C}^{\mathfrak{M}}\otimes\mathbb{C}^4$ to $\mathbb{C}^4$, followed by the map $v\mapsto\mathsf{b}\cdot\sum_m\hat{u}_m\cdot\delta_m\otimes v$. By the matrix determinant lemma, applied to this factorisation and to the invertible operator $\operatorname{diag}_m(R_m)$, and since $\det R_m=w(q_m)^2$,
\[
\det T(p)=\prod_{m\in\mathfrak{M}}\det R_m\cdot\det\left(1+\mathsf{b}\cdot\sum_{m\in\mathfrak{M}}\hat{u}_m^2\cdot R_m^{-1}\right)=\mathsf{b}^4\cdot\det\mathrm{Ch}(p)\cdot\prod_{m\in\mathfrak{M}}w(q_m)^2>0.
\]
This proves (b).

\emph{Proof of (c).} The spinors on the unit lattice form a space of dimension $4l^3$. By Sylvester's determinant identity $\det(1-XY)=\det(1-YX)$, applied with $X=\mathsf{b}\cdot Q(c)$ and $Y=T_c^{-1}Q^T(-c)$,
\begin{align*}
\det D(c)&=\mathsf{b}^{4l^3}\cdot\det\left(1-\mathsf{b}\cdot Q(c)\,T_c^{-1}Q^T(-c)\right)=\mathsf{b}^{4l^3}\cdot\det\left(T_c^{-1}\left(T_c-\mathsf{b}\cdot Q^T(-c)\,Q(c)\right)\right)\\
&=\mathsf{b}^{4l^3}\cdot\frac{\det\left(D_c+\mathfrak{m}\right)}{\det T_c}.
\end{align*}
Let $n:=l\cdot\mathfrak{N}$, and transport functions on $\boldsymbol{\eta}\mathbb{Z}^3$ to functions on $\mathbb{Z}^3$ by means of the bijection $x\mapsto x/\boldsymbol{\eta}$. This carries the spinors on the fine lattice onto the spinors on $\Lambda_n$ of Section~\ref{sec:constant-twists-free}, a space of dimension $4n^3$. Comparing the display of $D_A$ in Section~\ref{sec:exact-fermion-operators} with that in Section~\ref{sec:lattice-measures-laws}, one sees that the transport carries $\boldsymbol{\eta}\cdot D_c$ to the operator given by the formula for $D_A$ in Section~\ref{sec:lattice-measures-laws}, but with the phase factor $\exp\left(\pm i\,\mathsf{e}\cdot\boldsymbol{\eta}\cdot c_{\nu}\right)$ on the bond from $x$ to $x\pm e_{\nu}$. As $\mathsf{e}\cdot\boldsymbol{\eta}\cdot c_{\nu}=\theta_{\nu}/n$, these phase factors are $\exp\left(\pm i\theta_{\nu}/n\right)$, which by part (a) of the proof of Lemma~\ref{lem:product-formula-volume} are the link phases of $D_{H_{\theta}}$, with $H_{\theta}$ and $D_{H_{\theta}}$ as in Section~\ref{sec:constant-twists-free}. So that operator is $D_{H_{\theta}}$, the operator $\boldsymbol{\eta}\cdot\left(D_c+\mathfrak{m}\right)$ is carried to $D_{H_{\theta}}+\boldsymbol{\eta}\cdot\mathfrak{m}$, and
\[
\det\left(D_c+\mathfrak{m}\right)=\boldsymbol{\eta}^{-4n^3}\cdot\det\left(D_{H_{\theta}}+\boldsymbol{\eta}\cdot\mathfrak{m}\right).
\]
Dividing the resulting expression for $\det D(c)$ by the one for $c=0$, where $\theta=0$ and $H_0=0$, the factors $\mathsf{b}^{4l^3}$ and $\boldsymbol{\eta}^{-4n^3}$ cancel, and (c) follows.
\end{proof}

\begin{lemma}[Dependence on the twist, uniform in the mass]\label{lem:dependence-twist-uniform}
Let $0<b_{-}\le\mathsf{b}\le b_{+}$. There exist constants $\gamma_1\in(0,1/16]$ and $C_{\Sigma}<\infty$, depending only on $b_{-}$ and $b_{+}$, with the following property. For every integer $l\ge1$, every odd integer $\mathfrak{N}\ge3$, every $\mathfrak{m}\in(0,1]$ and every $\theta\in\mathbb{R}^3$, the operators $T_c$, $T_0$ and the number $R_{\theta}$ of Section~\ref{sec:averaging-blocks-flat}, formed with these parameters, satisfy
\begin{multline}\label{eq:renormalisation-flow-18}
\left|\log\frac{\det T_c}{\det T_0}\right|\le C_{\Sigma}\cdot l^{-1}\exp(-\gamma_1l),\\
\left|\log R_{\theta}\right|\le2\,\varepsilon_5(\mathfrak{m}\cdot l)+C_{\Sigma}\cdot l^{-1}\exp(-\gamma_1l).
\end{multline}
\end{lemma}

\begin{proof}
For a complex $4\times4$ matrix $X$, $|X|$ denotes the operator norm of $X$ on $\mathbb{C}^4$, where $\mathbb{C}^4$ carries the standard scalar product $\langle\cdot,\cdot\rangle$ and the norm $|v|=\sqrt{\langle v,v\rangle}$. For $v\in\mathbb{C}^3$, $|v|_1$ and $|v|_2$ are the $\ell^1$-norm and the Euclidean norm of $v$; for $x\in\mathbb{Z}^3$, $|x|:=\max_{\nu}|x_{\nu}|$. For every integer $t\ge1$,
\[
\#\left|\left\{x\in\mathbb{Z}^3:|x|=t\right\}\right|=(2t+1)^3-(2t-1)^3=24t^2+2\le26t^2.
\]

\emph{Step 1: the function $\Phi$ and its Fourier coefficients.} For $p\in\mathbb{C}^3$ let $T(p)$ be the operator on $\mathbb{C}^{\mathfrak{M}}\otimes\mathbb{C}^4$ given by the formula that defines it in Lemma~\ref{lem:fourier-decomposition}, with $q_m=p+2\pi m$, $R_m=\hat{D}(q_m)+\mathfrak{m}$ and $\hat{u}_m=\hat{u}(q_m)$ for $m\in\mathfrak{M}$. Its matrix elements are entire functions of $p$, because $u_1$ and $\hat{D}$ are entire and the entries of $|\hat{u}\rangle\langle\hat{u}|$ are the products $\hat{u}_m\cdot\hat{u}_{m'}$. Every $p\in\mathbb{R}^3$ belongs to $P_{l,\theta}$ for a suitable $\theta\in\mathbb{R}^3$; hence Lemma~\ref{lem:fourier-decomposition}(a), (b) apply at every real $p$. For real $p$ the numbers $\hat{u}_m$ are real, $w(q_m)>0$, and $R_m^{-1}=\left(a(q_m)-i\gamma\cdot s(q_m)\right)/w(q_m)$ by Section~\ref{sec:averaging-blocks-flat}; therefore $\mathrm{Ch}(p)=\boldsymbol{\alpha}-i\gamma\cdot\boldsymbol{\beta}$ with
\[
\boldsymbol{\alpha}=\frac{1}{\mathsf{b}}+\sum_{m\in\mathfrak{M}}\frac{\hat{u}_m^2\cdot a(q_m)}{w(q_m)},\qquad\boldsymbol{\beta}_{\nu}=\sum_{m\in\mathfrak{M}}\frac{\hat{u}_m^2\cdot s_{\nu}(q_m)}{w(q_m)},
\]
and Lemma~\ref{lem:fourier-decomposition}(a), (b) give
\[
\det T(p)=\mathsf{b}^4\det\mathrm{Ch}(p)\prod_{m\in\mathfrak{M}}w(q_m)^2=\mathsf{b}^4\left(\boldsymbol{\alpha}^2+|\boldsymbol{\beta}|_2^2\right)^2\prod_{m\in\mathfrak{M}}w(q_m)^2>0.
\]
Thus $\Phi(p):=\log\det T(p)$ is defined for $p\in\mathbb{R}^3$, and $\Phi$ is real analytic, being the logarithm of the restriction to $\mathbb{R}^3$ of an entire function that is positive there.

Since $\boldsymbol{\eta}=1/\mathfrak{N}$ and $u_1$ is a linear combination of the functions $z\mapsto\exp\left(-\frac{izg}{\mathfrak{N}}\right)$ with $g\in\mathbb{Z}$, the functions $\hat{u}^2$, $\hat{D}$, $a$, $s_{\nu}$ and $w$ have the period $2\pi\cdot\mathfrak{N}$ in each variable. As $\mathfrak{N}$ is odd, the set $\{-(\mathfrak{N}-1)/2,\dots,(\mathfrak{N}-1)/2\}$ is a system of representatives of $\mathbb{Z}/\mathfrak{N}\mathbb{Z}$. Hence for each direction $\nu$ there is a permutation $\sigma$ of $\mathfrak{M}$ with $m+e_{\nu}-\sigma(m)\in\mathfrak{N}\mathbb{Z}^3$ for all $m\in\mathfrak{M}$, and then $(p+2\pi e_{\nu})+2\pi m$ differs from $q_{\sigma(m)}=p+2\pi\sigma(m)$ by an element of $2\pi\cdot\mathfrak{N}\mathbb{Z}^3$. Replacing $p$ by $p+2\pi e_{\nu}$ therefore only permutes the terms of the sums and of the product above, and $\Phi$ is $2\pi$-periodic in each coordinate.

For a permutation $\tau$ of $\{0,1,2\}$ and $v\in\mathbb{C}^3$ write $v^{\tau}:=(v_{\tau(0)},v_{\tau(1)},v_{\tau(2)})$. The map $m\mapsto m^{\tau}$ is a bijection of $\mathfrak{M}$ onto itself, $p^{\tau}+2\pi m^{\tau}=(q_m)^{\tau}$, and
\[
\hat{u}(q^{\tau})=\hat{u}(q),\qquad a(q^{\tau})=a(q),\qquad w(q^{\tau})=w(q),\qquad s(q^{\tau})=s(q)^{\tau},
\]
where $s=(s_0,s_1,s_2)$. Replacing $p$ by $p^{\tau}$, and the index of summation $m$ by $m^{\tau}$, therefore leaves $\boldsymbol{\alpha}$ and $\prod_mw(q_m)$ unchanged and replaces $\boldsymbol{\beta}$ by $\boldsymbol{\beta}^{\tau}$; consequently $\Phi(p^{\tau})=\Phi(p)$.

By Lemma~\ref{lem:fourier-decomposition}(b), $T_c$ is unitarily equivalent to the direct sum of the operators $T(p)$, $p\in P_{l,\theta}$, and $T_0$, which is $T_c$ at $\theta=0$, to the direct sum of the operators $T(p)$, $p\in P_{l,0}$. Hence
\[
\log\frac{\det T_c}{\det T_0}=\sum_{p\in P_{l,\theta}}\Phi(p)-\sum_{p\in P_{l,0}}\Phi(p).
\]
Let
\[
\hat{\Phi}(x):=\frac{1}{(2\pi)^3}\int_{[-\pi,\pi]^3}\mathrm{d}^3p\,\Phi(p)\exp(-i\,p\cdot x)\qquad\left(x\in\mathbb{Z}^3\right)
\]
be the Fourier coefficients of $\Phi$. The identity in \eqref{eq:renormalisation-flow-14}, applied with $n=l$ to the smooth $2\pi$-periodic function $f=\Phi$, once at $\theta$ and once at $0$, gives
\begin{align*}
\log\frac{\det T_c}{\det T_0}&=l^3\sum_{w\in\mathbb{Z}^3\setminus\{0\}}\hat{\Phi}(lw)\\
&\qquad\cdot\left(\exp\left(i\,(\pi+\theta)\cdot w\right)-\exp\left(i\pi\,(w_0+w_1+w_2)\right)\right),
\end{align*}
since the two terms with $w=0$ cancel. It follows that
\[
\left|\log\frac{\det T_c}{\det T_0}\right|\le2\,l^3\sum_{w\in\mathbb{Z}^3\setminus\{0\}}\left|\hat{\Phi}(lw)\right|.
\]

\emph{Step 2: separation of the mode $m=0$.} Fix $\mathsf{p}=(p_1,p_2)\in[-\pi,\pi]^2$ and write $p=(\boldsymbol{\zeta},\mathsf{p})$, that is, $\boldsymbol{\zeta}=p_0$, where $\boldsymbol{\zeta}$ is now allowed to be complex; then $q_m$, $R_m$, $\hat{u}_m$ and $w(q_m)$ are functions of $\boldsymbol{\zeta}\in\mathbb{C}$. Let
\[
\mathsf{D}_2:=\left\{\boldsymbol{\zeta}\in\mathbb{C}:|\operatorname{Re}\boldsymbol{\zeta}|\le\pi+\frac{1}{2},\ |\operatorname{Im}\boldsymbol{\zeta}|\le\frac{1}{4}\right\}.
\]
For real $\boldsymbol{\zeta}$ put
\[
\mathrm{Reg}:=\frac{1}{\mathsf{b}}+\sum_{m\in\mathfrak{M}\setminus\{0\}}\hat{u}_m^2\cdot R_m^{-1},\qquad\mathsf{N}:=\hat{u}_0^2+R_0\cdot\mathrm{Reg},\qquad\boldsymbol{\Upsilon}:=\sum_{m\in\mathfrak{M}\setminus\{0\}}2\log w(q_m).
\]
By the definition of $\mathrm{Ch}(p)$ in Lemma~\ref{lem:fourier-decomposition}, $\mathrm{Ch}(p)=\hat{u}_0^2\cdot R_0^{-1}+\mathrm{Reg}$; hence $\mathsf{N}=R_0\cdot\mathrm{Ch}(p)$. Moreover $\det R_0=w(q_0)^2$ (Section~\ref{sec:averaging-blocks-flat}), $w(q_0)>0$ and $\det\mathrm{Ch}(p)>0$, so that $\det\mathsf{N}=w(q_0)^2\det\mathrm{Ch}(p)>0$. With $\boldsymbol{\Lambda}:=\log\det\mathsf{N}$ the formula for $\det T(p)$ in Step~1 becomes
\[
\Phi=\boldsymbol{\Lambda}+\boldsymbol{\Upsilon}+4\log\mathsf{b}\qquad(\boldsymbol{\zeta}\in\mathbb{R}).
\]

\emph{Step 3: the modes $m\ne0$ on $\mathsf{D}_2$.} Let $m\in\mathfrak{M}\setminus\{0\}$ and $\boldsymbol{\zeta}\in\mathsf{D}_2$, and put $\mathsf{q}:=\mathsf{p}+2\pi(m_1,m_2)\in\mathbb{R}^2$. Then $q_m=\left((q_m)_0,\mathsf{q}\right)$ with $(q_m)_0=\boldsymbol{\zeta}+2\pi m_0$, and $\boldsymbol{\eta}\cdot\mathsf{q}\in[-\pi,\pi]^2$, because $|p_{\nu}+2\pi m_{\nu}|\le\pi+\pi(\mathfrak{N}-1)=\pi\cdot\mathfrak{N}$ for $\nu=1,2$. By Section~\ref{sec:averaging-blocks-flat}, $w(q_m)=\boldsymbol{\eta}^{-2}\cdot w_{\boldsymbol{\eta}\cdot\mathfrak{m}}(\boldsymbol{\eta}\cdot q_m)$. Let $\mathcal{W}$, $\mathcal{S}$, $\acute{M}$ and $\mathrm{En}$ be the numbers of \eqref{eq:renormalisation-flow-15} and of the definitions preceding it in the proof of Lemma~\ref{lem:product-formula-volume}, formed for the mass $\mathbf{m}=\boldsymbol{\eta}\cdot\mathfrak{m}\in(0,1/3]$ and with $\boldsymbol{\eta}\cdot\mathsf{q}$ in place of $\mathsf{p}$; as $\mathsf{q}=\left((q_m)_1,(q_m)_2\right)$, this means
\[
\mathcal{W}=\sum_{\nu=1}^2\left(1-\cos\left(\boldsymbol{\eta}\cdot(q_m)_{\nu}\right)\right),\qquad\mathcal{S}=\sum_{\nu=1}^2\sin^2\left(\boldsymbol{\eta}\cdot(q_m)_{\nu}\right),\qquad\acute{M}=1+\mathbf{m}+\mathcal{W},
\]
and $2\cosh\mathrm{En}=(\acute{M}^2+1+\mathcal{S})/\acute{M}$ with $\mathrm{En}>0$. Since $0\le\mathcal{W}\le4$ and $0\le\mathcal{S}\le2$, one has $1\le\acute{M}\le1+1/3+4\le6$. The identity \eqref{eq:renormalisation-flow-15} remains valid for complex $p_0$, as it uses only $\cos^2+\sin^2=1$, and $2\cosh E-2\cos t=\exp(E)\left(1-\exp(-E+it)\right)\left(1-\exp(-E-it)\right)$ for $E\in\mathbb{R}$ and $t\in\mathbb{C}$. Therefore
\[
w(q_m)=\boldsymbol{\eta}^{-2}\cdot\acute{M}\cdot\exp(\mathrm{En})\left(1-z_{+}\right)\left(1-z_{-}\right),\qquad z_{\pm}:=\exp(-\mathrm{En})\exp\left(\pm i\boldsymbol{\eta}\cdot(q_m)_0\right).
\]
Moreover
\[
\frac{\boldsymbol{\eta}\cdot|\mathsf{q}|_2}{6}\le\mathrm{En}\le3.7:
\]
the lower bound is \eqref{eq:renormalisation-flow-16} at the point $\boldsymbol{\eta}\cdot\mathsf{q}\in[-\pi,\pi]^2$, and the upper bound follows from $2\cosh\mathrm{En}=(\acute{M}^2+1+\mathcal{S})/\acute{M}\le\acute{M}+3/\acute{M}\le9$, by $\mathcal{S}\le2$ and $1\le\acute{M}\le6$. Put
\[
\boldsymbol{\phi}:=\boldsymbol{\eta}\cdot\operatorname{Re}(q_m)_0,\qquad\boldsymbol{\kappa}:=\operatorname{Im}\boldsymbol{\zeta},
\]
so that $|z_{\pm}|=\exp(-\mathrm{En}\mp\boldsymbol{\eta}\cdot\boldsymbol{\kappa})$. As $|\operatorname{Re}(q_m)_0|\le\pi+1/2+\pi(\mathfrak{N}-1)=\pi\cdot\mathfrak{N}+1/2$ and $\boldsymbol{\eta}\le1/3$, one has $|\boldsymbol{\phi}|\le\pi+1/6$; since $\sin x\ge x/\sqrt{3}$ for $0\le x\le\pi/2+1/12$, it follows that $\sin^2(\boldsymbol{\phi}/2)\ge\boldsymbol{\phi}^2/12$.

\emph{A lower bound for $|1-z_{\pm}|$.} Neither $z_{+}$ nor $z_{-}$ lies in $[1,\infty)$, and
\[
\left|1-z_{\pm}\right|\ge c_1\cdot\boldsymbol{\eta}\cdot\left|\operatorname{Re}q_m\right|_2,\qquad c_1:=\frac{1}{50}.
\]
To see this, let $z$ be one of $z_{+}$, $z_{-}$ and let $\hat{r}:=|z|=\exp(-\mathrm{En}\mp\boldsymbol{\eta}\cdot\boldsymbol{\kappa})$. Then $z=\hat{r}\cdot\exp(\pm i\boldsymbol{\phi})$ and
\[
|1-z|^2=(1-\hat{r})^2+4\hat{r}\cdot\sin^2\frac{\boldsymbol{\phi}}{2}.
\]

\emph{Case A: $\mathrm{En}\ge2\boldsymbol{\eta}\cdot|\boldsymbol{\kappa}|$.} Then $\hat{r}\le\exp(-\mathrm{En}+\boldsymbol{\eta}\cdot|\boldsymbol{\kappa}|)\le\exp\left(-\frac{\mathrm{En}}{2}\right)<1$; in particular $z\notin[1,\infty)$. If $\hat{r}\ge1/4$, then $4\hat{r}\cdot\sin^2(\boldsymbol{\phi}/2)\ge\boldsymbol{\phi}^2/12$; moreover $1-\exp(-x)\ge0.4555\,x$ for $0\le x\le1.85$ and $\mathrm{En}/2\le1.85$, so that $1-\hat{r}\ge1-\exp\left(-\frac{\mathrm{En}}{2}\right)\ge0.2277\cdot\mathrm{En}$, and
\[
|1-z|^2\ge(0.2277\cdot\mathrm{En})^2+\frac{\boldsymbol{\phi}^2}{12}\ge0.022\left(\mathrm{En}^2+\boldsymbol{\phi}^2\right).
\]
If $\hat{r}<1/4$, then $|1-z|\ge1-\hat{r}>3/4$, and since $\mathrm{En}^2+\boldsymbol{\phi}^2\le3.7^2+(\pi+1/6)^2\le24.7$,
\[
|1-z|^2\ge\frac{9}{16}\ge0.022\left(\mathrm{En}^2+\boldsymbol{\phi}^2\right).
\]

\emph{Case B: $\mathrm{En}<2\boldsymbol{\eta}\cdot|\boldsymbol{\kappa}|$.} As $|\boldsymbol{\kappa}|\le1/4$, one has $\mathrm{En}<\boldsymbol{\eta}/2$, and the lower bound for $\mathrm{En}$ gives $|\mathsf{q}|_2\le6\mathrm{En}/\boldsymbol{\eta}<3<\pi$. If $(m_1,m_2)\ne0$, some coordinate of $\mathsf{q}=\mathsf{p}+2\pi(m_1,m_2)$ has modulus at least $2\pi-\pi=\pi$; hence $(m_1,m_2)=0$, and $m_0\ne0$ because $m\ne0$. It follows that $|\operatorname{Re}(q_m)_0|\ge2\pi-(\pi+1/2)=\pi-1/2$, so that $0<\boldsymbol{\eta}\cdot(\pi-1/2)\le|\boldsymbol{\phi}|\le\pi+1/6<2\pi$; thus $z=\hat{r}\cdot\exp(\pm i\boldsymbol{\phi})$ is not a positive real number, and $z\notin[1,\infty)$. Furthermore $\hat{r}\ge\exp(-\mathrm{En}-\boldsymbol{\eta}\cdot|\boldsymbol{\kappa}|)\ge\exp\left(-\frac{3\boldsymbol{\eta}}{4}\right)\ge0.77$, whence
\[
|1-z|^2\ge4\hat{r}\cdot\sin^2\frac{\boldsymbol{\phi}}{2}\ge\frac{4\cdot0.77}{12}\cdot\boldsymbol{\phi}^2\ge0.25\cdot\boldsymbol{\phi}^2,
\]
while $\mathrm{En}^2\le\boldsymbol{\eta}^2/4\le\boldsymbol{\phi}^2/27$, because $\boldsymbol{\phi}^2\ge(\pi-1/2)^2\cdot\boldsymbol{\eta}^2\ge27\boldsymbol{\eta}^2/4$. Hence $\mathrm{En}^2+\boldsymbol{\phi}^2\le\frac{28}{27}\cdot\boldsymbol{\phi}^2$, and again $|1-z|^2\ge0.022\left(\mathrm{En}^2+\boldsymbol{\phi}^2\right)$.

In both cases, since $\operatorname{Re}q_m=\left(\operatorname{Re}(q_m)_0,\mathsf{q}\right)$, $\mathrm{En}\ge\boldsymbol{\eta}\cdot|\mathsf{q}|_2/6$ and $\boldsymbol{\phi}=\boldsymbol{\eta}\cdot\operatorname{Re}(q_m)_0$,
\[
|1-z|^2\ge0.022\left(\mathrm{En}^2+\boldsymbol{\phi}^2\right)\ge\frac{0.022}{36}\cdot\boldsymbol{\eta}^2\cdot\left|\operatorname{Re}q_m\right|_2^2,
\]
and $\sqrt{0.022/36}\ge1/50=c_1$. This proves the lower bound for $|1-z_{\pm}|$. It has the following consequences.

\emph{(i)} Since $\acute{M}\ge1$ and $\exp(\mathrm{En})\ge1$,
\[
\left|w(q_m)\right|\ge c_1^2\cdot\left|\operatorname{Re}q_m\right|_2^2.
\]

\emph{(ii)} Let $\operatorname{Log}$ be the principal branch of the logarithm on $\mathbb{C}\setminus(-\infty,0]$, and put, for $\boldsymbol{\zeta}\in\mathsf{D}_2$,
\[
\log w(q_m):=-2\log\boldsymbol{\eta}+\log\acute{M}+\mathrm{En}+\operatorname{Log}(1-z_{+})+\operatorname{Log}(1-z_{-}).
\]
One has $1-z_{\pm}\notin(-\infty,0]$; as $z_{\pm}$ are entire functions of $\boldsymbol{\zeta}$ and $\acute{M}$, $\mathrm{En}$ do not depend on $\boldsymbol{\zeta}$, this is a holomorphic function of $\boldsymbol{\zeta}$ on $\mathsf{D}_2$, whose exponential is $w(q_m)$. For real $\boldsymbol{\zeta}$ the number $z_{-}$ is the complex conjugate of $z_{+}$, so that $\log w(q_m)$ is real and is the logarithm of the positive number $w(q_m)$. In particular $\boldsymbol{\Upsilon}=\sum_{m\ne0}2\log w(q_m)$ is a holomorphic function on $\mathsf{D}_2$ which coincides on the real axis with the function of Step~2. For $z\in\mathbb{C}\setminus[1,\infty)$,
\[
\left(z\frac{\mathrm{d}}{\mathrm{d}z}\right)^4\operatorname{Log}(1-z)=-\frac{z\left(1+4z+z^2\right)}{(1-z)^4},
\]
and $\partial_{\boldsymbol{\zeta}}z_{\pm}=\pm i\boldsymbol{\eta}\cdot z_{\pm}$; hence $\partial_{\boldsymbol{\zeta}}^4\operatorname{Log}(1-z_{\pm})=\boldsymbol{\eta}^4\cdot\left[\left(z\frac{\mathrm{d}}{\mathrm{d}z}\right)^4\operatorname{Log}(1-z)\right]_{z=z_{\pm}}$. Since $|z_{\pm}|\le\exp(\boldsymbol{\eta}\cdot|\boldsymbol{\kappa}|)\le\exp\left(\frac{1}{12}\right)$ and $\exp\left(\frac{1}{12}\right)\left(1+4\exp\left(\frac{1}{12}\right)+\exp\left(\frac{1}{6}\right)\right)\le7.1$, the lower bound for $|1-z_{\pm}|$ gives
\[
\left|\partial_{\boldsymbol{\zeta}}^4\,2\log w(q_m)\right|\le2\cdot2\cdot7.1\cdot\frac{\boldsymbol{\eta}^4}{\left(c_1\cdot\boldsymbol{\eta}\cdot\left|\operatorname{Re}q_m\right|_2\right)^4}=4\cdot7.1\cdot c_1^{-4}\left|\operatorname{Re}q_m\right|_2^{-4}.
\]
If $\nu$ is a direction with $|m_{\nu}|=|m|$, the $\nu$-th coordinate of $\operatorname{Re}q_m$ is $\operatorname{Re}p_{\nu}+2\pi m_{\nu}$ with $|\operatorname{Re}p_{\nu}|\le\pi+1/2$, hence of modulus at least $2\pi|m|-\pi-1/2$; therefore
\[
\left|\operatorname{Re}q_m\right|_2\ge\left(\pi-\frac{1}{2}\right)(2|m|-1).
\]
There are at most $26t^2$ points $m$ with $|m|=t$; as $\pi-1/2\ge2.64$ and
\[
\sum_{t\ge1}26t^2\left(2.64\,(2t-1)\right)^{-4}\le0.6,
\]
it follows that on $\mathsf{D}_2$
\[
\left|\partial_{\boldsymbol{\zeta}}^4\boldsymbol{\Upsilon}\right|\le4\cdot7.1\cdot0.6\cdot c_1^{-4}\le C_{\boldsymbol{\Upsilon}}:=18\,c_1^{-4}.
\]

\emph{(iii)} For $x\in\mathbb{C}$ one has $\sin x=x\int_0^1\mathrm{d}t\,\cos(tx)$ and $1-\cos x=x\int_0^1\mathrm{d}t\,\sin(tx)$. Since $|\cos(tx)|$ and $|\sin(tx)|$ are at most $\cosh(t\operatorname{Im}x)\le\cosh(\operatorname{Im}x)$ for $0\le t\le1$,
\[
|\sin x|\le|x|\cosh(\operatorname{Im}x),\qquad|1-\cos x|\le|x|\cosh(\operatorname{Im}x).
\]
Now let $m\in\mathfrak{M}$ be arbitrary, $m=0$ included, and $q:=q_m$. Then $\operatorname{Im}q=(\boldsymbol{\kappa},0,0)$, and these bounds at $x=\boldsymbol{\eta}\cdot q_{\nu}$, where $|\operatorname{Im}x|\le\boldsymbol{\eta}/4\le1/12$ and $2\cosh(1/12)\le2.01$, give
\[
|a(q)|+\sum_{\nu}|s_{\nu}(q)|\le\mathfrak{m}+2\cosh\frac{1}{12}\cdot|q|_1\le1+2.01\left(|\operatorname{Re}q|_1+\frac{1}{4}\right).
\]
For $m\ne0$ the right side is at most $5\left|\operatorname{Re}q_m\right|_2$, by $|\operatorname{Re}q_m|_1\le\sqrt{3}\,|\operatorname{Re}q_m|_2$ and $|\operatorname{Re}q_m|_2\ge\pi-1/2$; for $m=0$ it is at most $22$, by $|\operatorname{Re}q_0|_1\le3\pi+1/2$. By Section~\ref{sec:averaging-blocks-flat}, $R_m=a(q_m)+i\gamma\cdot s(q_m)$, and $\left(a+i\gamma\cdot s\right)\left(a-i\gamma\cdot s\right)=w$ also for complex arguments. For $m\ne0$ one has $w(q_m)\ne0$ by (i), so that $R_m$ is invertible with $R_m^{-1}=\left(a(q_m)-i\gamma\cdot s(q_m)\right)/w(q_m)$; as $|\gamma_{\nu}|=1$,
\[
\left|R_m^{-1}\right|\le\frac{|a(q_m)|+\sum_{\nu}|s_{\nu}(q_m)|}{|w(q_m)|}\le\frac{5c_1^{-2}}{\left|\operatorname{Re}q_m\right|_2}\le2c_1^{-2}\quad(m\ne0),\qquad|R_0|\le22.
\]

\emph{(iv)} The function $u_1(z)=\mathfrak{N}^{-1}\sum_g\exp\left(-\frac{izg}{\mathfrak{N}}\right)$, the sum being over the integers $g$ with $|g|\le(\mathfrak{N}-1)/2$, is the mean of $\mathfrak{N}$ exponentials $\exp(-izy)$ whose frequencies $y=g/\mathfrak{N}$ have modulus less than $1/2$; hence, for real $t$ and $\boldsymbol{\kappa}$,
\[
|u_1(t+i\boldsymbol{\kappa})|\le\exp\left(\frac{|\boldsymbol{\kappa}|}{2}\right).
\]
Summation of this geometric progression gives the closed form of Section~\ref{sec:averaging-blocks-flat},
\[
u_1(z)=\sin\frac{z}{2}\cdot\left(\mathfrak{N}\cdot\sin\frac{z}{2\mathfrak{N}}\right)^{-1}\qquad\text{whenever }\sin\frac{z}{2\mathfrak{N}}\ne0.
\]
Let $t$ be real with $0<|t|\le\pi\cdot\mathfrak{N}+1/2$. Then $|t|/(2\mathfrak{N})\le\pi/2+1/12\le1.66$, and $\sin x\ge x/2$ for $0\le x\le1.66$; together with $|\sin(x+iy)|^2=\sin^2x+\sinh^2y$ for real $x,y$ this gives
\[
\left|\mathfrak{N}\cdot\sin\frac{t+i\boldsymbol{\kappa}}{2\mathfrak{N}}\right|\ge\mathfrak{N}\cdot\left|\sin\frac{t}{2\mathfrak{N}}\right|\ge\frac{|t|}{4}>0.
\]
So the denominator of the closed form does not vanish at $z=t+i\boldsymbol{\kappa}$. As moreover $|\sin((t+i\boldsymbol{\kappa})/2)|\le\cosh(\boldsymbol{\kappa}/2)$, by the same identity, the closed form gives
\[
|u_1(t+i\boldsymbol{\kappa})|=\left|\sin\frac{t+i\boldsymbol{\kappa}}{2}\right|\cdot\left|\mathfrak{N}\cdot\sin\frac{t+i\boldsymbol{\kappa}}{2\mathfrak{N}}\right|^{-1}\le\frac{4\cosh(\boldsymbol{\kappa}/2)}{|t|}.
\]
Fix a direction $\nu$. As $m_{\nu}$ ranges over $\{-(\mathfrak{N}-1)/2,\dots,(\mathfrak{N}-1)/2\}$, the $\nu$-th coordinate of $q_m$ has the imaginary part $\boldsymbol{\kappa}$ if $\nu=0$ and $0$ if $\nu=1,2$, and its real part $t=\operatorname{Re}p_{\nu}+2\pi m_{\nu}$ satisfies $|t|\le\pi\cdot\mathfrak{N}+1/2$ and, if $m_{\nu}\ne0$, $|t|\ge(\pi-1/2)(2|m_{\nu}|-1)>0$. Using the first bound for $m_{\nu}=0$, the second for $m_{\nu}\ne0$, $\cosh^2(\boldsymbol{\kappa}/2)\le\exp(|\boldsymbol{\kappa}|)$ and $\sum_{j\ge1}(2j-1)^{-2}=\pi^2/8$, one obtains
\[
\sum_{m_{\nu}}\left|u_1\left((q_m)_{\nu}\right)\right|^2\le\exp(|\boldsymbol{\kappa}|)\left(1+32\cdot\frac{\pi^2}{8}\cdot\left(\pi-\frac{1}{2}\right)^{-2}\right)\le6.7\exp(|\boldsymbol{\kappa}|),
\]
where $\boldsymbol{\kappa}$ is to be replaced by $0$ if $\nu=1,2$. Since $|\hat{u}_m|^2=\prod_{\nu}|u_1((q_m)_{\nu})|^2$ and $|\boldsymbol{\kappa}|\le1/4$,
\[
\sum_{m\in\mathfrak{M}}\left|\hat{u}_m\right|^2\le6.7^3\cdot\exp\left(\frac{1}{4}\right)\le B_u:=400.
\]

\emph{(v)} By (i) and (iii), the matrices $R_m$, $m\ne0$, are invertible for $\boldsymbol{\zeta}\in\mathsf{D}_2$, and the formulas of Step~2 define $\mathrm{Reg}$ and $\mathsf{N}$ on $\mathsf{D}_2$. The matrix $\mathsf{N}$ is a holomorphic function of $\boldsymbol{\zeta}$ there, since $\hat{u}_m$, $a(q_m)$, $s_{\nu}(q_m)$, $w(q_m)$ and $R_m$ are entire functions of $\boldsymbol{\zeta}$ and $w(q_m)\ne0$ for $m\ne0$. By (iii) and (iv), on $\mathsf{D}_2$
\[
|\mathrm{Reg}|\le\frac{1}{\mathsf{b}}+\sum_{m\ne0}\left|\hat{u}_m\right|^2\cdot\left|R_m^{-1}\right|\le B_{\mathrm{Reg}}:=\frac{1}{b_{-}}+2B_u\,c_1^{-2},
\]
and, as $|\hat{u}_0|^2\le\exp(|\boldsymbol{\kappa}|)\le2$ by (iv),
\[
|\mathsf{N}|\le|\hat{u}_0|^2+|R_0|\cdot|\mathrm{Reg}|\le B_N:=2+22\,B_{\mathrm{Reg}}.
\]

\emph{Step 4: $\mathsf{N}$ on the real axis.} Let $\boldsymbol{\zeta}=p_0$ be real with $|p_0|\le\pi+1/4$. For $0<|t|<2\pi$ the numbers $\sin(t/2)$ and $\mathfrak{N}\cdot\sin(t/(2\mathfrak{N}))$ have the sign of $t$, and $|\mathfrak{N}\cdot\sin(t/(2\mathfrak{N}))|\le|t|/2$. In particular the second number is not zero, so that the closed form of $u_1$ recalled in (iv) of Step~3 applies at $z=t$ and gives
\[
u_1(t)=\sin\frac{t}{2}\cdot\left(\mathfrak{N}\cdot\sin\frac{t}{2\mathfrak{N}}\right)^{-1}\ge\frac{2}{t}\cdot\sin\frac{t}{2}\qquad(0<|t|<2\pi);
\]
moreover $u_1(0)=1$. As $\sin x/x$ is decreasing on $(0,\pi]$ and $\sin x/x\ge0.585$ at $x=(\pi+1/4)/2$, one has $u_1(p_0)\ge0.585$, and $u_1(p_{\nu})\ge2/\pi$ for $\nu=1,2$. Therefore
\[
\hat{u}_0^2\ge u_{-}:=\left(0.585\cdot\left(\frac{2}{\pi}\right)^2\right)^2\ge0.056.
\]
In $R_0=a(q_0)+i\gamma\cdot s(q_0)$ the number $a(q_0)\ge\mathfrak{m}>0$ and the vector $s(q_0)$ are real; as the matrices $\gamma_{\nu}$ are Hermitian, $R_0R_0^{*}=R_0^{*}R_0=w(q_0)$. Thus $R_0$ is normal, $R_0/\sqrt{w(q_0)}$ is unitary, $|R_0|=\sqrt{w(q_0)}$ and $|R_0^{-1}|=1/\sqrt{w(q_0)}$.

If $|R_0|\cdot B_{\mathrm{Reg}}\le u_{-}/2$, then $\left|\hat{u}_0^{-2}\cdot R_0\cdot\mathrm{Reg}\right|\le1/2$ by (v) of Step~3, and $\mathsf{N}=\hat{u}_0^2\left(1+\hat{u}_0^{-2}\cdot R_0\cdot\mathrm{Reg}\right)$ is invertible with $|\mathsf{N}^{-1}|\le2\hat{u}_0^{-2}\le2/u_{-}$.

Otherwise $|R_0^{-1}|=|R_0|^{-1}<2B_{\mathrm{Reg}}/u_{-}$. By Lemma~\ref{lem:fourier-decomposition}(a), $\operatorname{Re}\mathrm{Ch}(p)\ge1/\mathsf{b}$; so $|\mathrm{Ch}(p)v|\cdot|v|\ge\operatorname{Re}\langle v,\mathrm{Ch}(p)v\rangle\ge|v|^2/\mathsf{b}$ for $v\in\mathbb{C}^4$, and $\mathrm{Ch}(p)$ is invertible with $|\mathrm{Ch}(p)^{-1}|\le\mathsf{b}$. Since $\mathsf{N}=R_0\cdot\mathrm{Ch}(p)$,
\[
\left|\mathsf{N}^{-1}\right|=\left|\mathrm{Ch}(p)^{-1}\cdot R_0^{-1}\right|\le\frac{2\mathsf{b}\cdot B_{\mathrm{Reg}}}{u_{-}}.
\]
In either case, as $\mathsf{b}\le b_{+}$,
\[
\left|\mathsf{N}^{-1}\right|\le B'_N:=\frac{2}{u_{-}}\max\left(1,b_{+}B_{\mathrm{Reg}}\right).
\]
Finally $\det\mathsf{N}=w(q_0)^2\det\mathrm{Ch}(p)>0$, by Step~2.

\emph{Step 5: a logarithm of $\det\mathsf{N}$ in a strip.} Put
\[
\gamma_1:=\frac{1}{32\,B_NB'_N};
\]
since $B_N\ge2$ and $B'_N\ge1$, one has $\gamma_1\le1/16$. Let $p_0,\boldsymbol{\kappa}\in\mathbb{R}$ with $|p_0|\le\pi+1/4$ and $|\boldsymbol{\kappa}|\le2\gamma_1$, so that $|\boldsymbol{\kappa}|\le1/8$, and let $t$ lie in the closed interval with the end points $0$ and $\boldsymbol{\kappa}$. Every complex number at distance at most $1/8$ from $p_0+it$ has real part of modulus at most $\pi+1/4+1/8\le\pi+1/2$ and imaginary part of modulus at most $|\boldsymbol{\kappa}|+1/8\le1/4$, and hence belongs to $\mathsf{D}_2$. By (v) of Step~3, $\mathsf{N}$ is holomorphic on $\mathsf{D}_2$ and $|\mathsf{N}|\le B_N$ there; Cauchy's estimate for the first derivative, with the radius $1/8$, therefore yields
\[
\left|\partial_{\boldsymbol{\zeta}}\mathsf{N}(p_0+it)\right|\le8B_N.
\]
Integration along the vertical path $t\mapsto p_0+it$ shows that the matrix
\[
\Delta\mathsf{N}:=\mathsf{N}(p_0+i\boldsymbol{\kappa})-\mathsf{N}(p_0)=i\int_0^{\boldsymbol{\kappa}}\mathrm{d}t\,\partial_{\boldsymbol{\zeta}}\mathsf{N}(p_0+it)
\]
has norm $|\Delta\mathsf{N}|\le8B_N|\boldsymbol{\kappa}|\le16B_N\gamma_1=1/(2B'_N)$; together with the bound $\left|\mathsf{N}(p_0)^{-1}\right|\le B'_N$ of Step~4 this gives $\left|\mathsf{N}(p_0)^{-1}\cdot\Delta\mathsf{N}\right|\le1/2$. Consequently $\mathsf{N}(p_0+i\boldsymbol{\kappa})=\mathsf{N}(p_0)\left(1+\mathsf{N}(p_0)^{-1}\cdot\Delta\mathsf{N}\right)$ is invertible, that is, $\det\mathsf{N}\ne0$ on
\[
\mathsf{D}_1:=\left\{\boldsymbol{\zeta}\in\mathbb{C}:|\operatorname{Re}\boldsymbol{\zeta}|\le\pi+\frac{1}{4},\ |\operatorname{Im}\boldsymbol{\zeta}|\le2\gamma_1\right\}.
\]
For a complex $4\times4$ matrix $X$ with $|X|\le1/2$ put $\operatorname{Log}(1+X):=\sum_{j\ge1}(-1)^{j+1}X^j/j$; then $\exp\left(\operatorname{tr}\operatorname{Log}(1+X)\right)=\det(1+X)$ and $|\operatorname{tr}\operatorname{Log}(1+X)|\le4\sum_{j\ge1}2^{-j}/j=4\log2$. For $p_0+i\boldsymbol{\kappa}\in\mathsf{D}_1$ define
\[
\boldsymbol{\Lambda}(p_0+i\boldsymbol{\kappa}):=\log\det\mathsf{N}(p_0)+\operatorname{tr}\operatorname{Log}\left(1+\mathsf{N}(p_0)^{-1}\cdot\Delta\mathsf{N}\right),
\]
with the real logarithm of the positive number $\det\mathsf{N}(p_0)$. This is a continuous function on $\mathsf{D}_1$ whose exponential is the holomorphic function $\det\mathsf{N}$. A continuous logarithm of a holomorphic function is holomorphic; hence $\boldsymbol{\Lambda}$ is a holomorphic function of $\boldsymbol{\zeta}$ on $\mathsf{D}_1$. On the real axis $\Delta\mathsf{N}=0$, and there $\boldsymbol{\Lambda}$ is the real function of Step~2. Since $|\det X|\le|X|^4$ for every complex $4\times4$ matrix $X$, (v) of Step~3 and Step~4 give $\left(B'_N\right)^{-4}\le\det\mathsf{N}(p_0)\le B_N^4$, and therefore
\[
|\boldsymbol{\Lambda}|\le B_{\boldsymbol{\Lambda}}:=4\max\left(\log B_N,\log B'_N\right)+4\log2\qquad\left(\boldsymbol{\zeta}\in\mathsf{D}_1\right).
\]
If $|\operatorname{Re}\boldsymbol{\zeta}|\le\pi$ and $|\operatorname{Im}\boldsymbol{\zeta}|\le\gamma_1$, the closed disc of radius $\gamma_1$ centred at $\boldsymbol{\zeta}$ lies in $\mathsf{D}_1$, and Cauchy's estimates give
\[
\left|\partial_{\boldsymbol{\zeta}}^4\boldsymbol{\Lambda}(\boldsymbol{\zeta})\right|\le4!\cdot B_{\boldsymbol{\Lambda}}\gamma_1^{-4}=24\,B_{\boldsymbol{\Lambda}}\gamma_1^{-4}.
\]

\emph{Step 6: periodic holomorphic extension.} The function $\boldsymbol{\Lambda}+\boldsymbol{\Upsilon}+4\log\mathsf{b}$ is holomorphic on $\mathsf{D}_1\subset\mathsf{D}_2$ (Steps~3 and~5), and by Step~2 it coincides with $\Phi(\cdot,\mathsf{p})$ on the real segment $[-\pi-1/4,\pi+1/4]$. On the rectangle $\{\boldsymbol{\zeta}\in\mathsf{D}_1:\pi-1/4\le\operatorname{Re}\boldsymbol{\zeta}\le\pi+1/4\}$ the holomorphic functions $\boldsymbol{\zeta}\mapsto(\boldsymbol{\Lambda}+\boldsymbol{\Upsilon})(\boldsymbol{\zeta})$ and $\boldsymbol{\zeta}\mapsto(\boldsymbol{\Lambda}+\boldsymbol{\Upsilon})(\boldsymbol{\zeta}-2\pi)$ coincide on the real segment $[\pi-1/4,\pi+1/4]$, because $\Phi(\cdot,\mathsf{p})$ has the period $2\pi$ (Step~1); by the identity theorem they coincide on the whole rectangle. Therefore $\Phi(\cdot,\mathsf{p})$ extends to a $2\pi$-periodic holomorphic function on the strip $|\operatorname{Im}\boldsymbol{\zeta}|\le2\gamma_1$. By periodicity it suffices to estimate its derivatives at points with $|\operatorname{Re}\boldsymbol{\zeta}|\le\pi$, and there Step~5 and (ii) of Step~3 give
\[
\left|\partial_{\boldsymbol{\zeta}}^4\Phi(\boldsymbol{\zeta},\mathsf{p})\right|\le C_4:=24\,B_{\boldsymbol{\Lambda}}\gamma_1^{-4}+C_{\boldsymbol{\Upsilon}}\qquad\left(|\operatorname{Im}\boldsymbol{\zeta}|\le\gamma_1\right).
\]

\emph{Step 7: decay of $\hat{\Phi}$; conclusion.} Let $x\in\mathbb{Z}^3\setminus\{0\}$. By Step~1, $\Phi$ is invariant under permutations of the coordinates, hence so is $\hat{\Phi}$, and one may assume $|x_0|=|x|$. Fix $\mathsf{p}\in[-\pi,\pi]^2$. Four integrations by parts in $p_0$, without boundary terms because the integrand is $2\pi$-periodic, give
\[
\int_{-\pi}^{\pi}\mathrm{d}p_0\,\Phi(p)\exp(-ip_0x_0)=x_0^{-4}\int_{-\pi}^{\pi}\mathrm{d}p_0\,\partial_{p_0}^4\Phi(p)\exp(-ip_0x_0).
\]
The function $\boldsymbol{\zeta}\mapsto\partial_{\boldsymbol{\zeta}}^4\Phi(\boldsymbol{\zeta},\mathsf{p})\exp(-i\boldsymbol{\zeta}\cdot x_0)$ is holomorphic and $2\pi$-periodic on the strip $|\operatorname{Im}\boldsymbol{\zeta}|\le2\gamma_1$. By Cauchy's integral theorem on the rectangle with the vertices $-\pi$, $\pi$, $\pi-i\gamma_1x_0/|x_0|$ and $-\pi-i\gamma_1x_0/|x_0|$, whose vertical sides contribute amounts that cancel by periodicity, the integral of this function over $[-\pi,\pi]$ equals its integral over the segment $[-\pi,\pi]-i\gamma_1x_0/|x_0|$. On that segment $|\exp(-i\boldsymbol{\zeta}\cdot x_0)|=\exp(-\gamma_1|x_0|)$ and $|\partial_{\boldsymbol{\zeta}}^4\Phi|\le C_4$ by Step~6. Hence
\[
\left|\frac{1}{2\pi}\int_{-\pi}^{\pi}\mathrm{d}p_0\,\Phi(p)\exp(-ip_0x_0)\right|\le C_4\,|x_0|^{-4}\exp(-\gamma_1|x_0|)
\]
for every $\mathsf{p}\in[-\pi,\pi]^2$, and integration over $\mathsf{p}$ gives
\[
\left|\hat{\Phi}(x)\right|\le C_4\,|x|^{-4}\exp(-\gamma_1|x|).
\]
By the inequality at the end of Step~1, since there are at most $26t^2$ points $w\in\mathbb{Z}^3$ with $|w|=t$ and $|lw|=lt$ for such $w$,
\begin{align*}
\left|\log\frac{\det T_c}{\det T_0}\right|&\le2\,l^3\sum_{t\ge1}26t^2\cdot C_4\,(lt)^{-4}\exp(-\gamma_1lt)\\
&\le2\,C_4\,l^{-1}\exp(-\gamma_1l)\sum_{t\ge1}26t^{-2}\le86\,C_4\,l^{-1}\exp(-\gamma_1l),
\end{align*}
because $52\sum_{t\ge1}t^{-2}=52\pi^2/6\le86$. With
\[
C_{\Sigma}:=86\,C_4
\]
this is the first inequality in \eqref{eq:renormalisation-flow-18}.

By Lemma~\ref{lem:fourier-decomposition}(c),
\[
\log R_{\theta}=\log\frac{\det\left(D_{H_{\theta}}+\boldsymbol{\eta}\cdot\mathfrak{m}\right)}{\det\left(D_0+\boldsymbol{\eta}\cdot\mathfrak{m}\right)}-\log\frac{\det T_c}{\det T_0},
\]
where the first quotient is the quotient of Section~\ref{sec:constant-twists-free} on $\Lambda_n$, $n=l\cdot\mathfrak{N}$, at the mass $\mathbf{m}=\boldsymbol{\eta}\cdot\mathfrak{m}\in(0,1/3]$, for which $\mathbf{m}\cdot n=\mathfrak{m}\cdot l$. By Corollary~\ref{cor:twist-ratios-free} the modulus of the first logarithm is at most $2\,\varepsilon_5(\mathfrak{m}\cdot l)$, and the second inequality in \eqref{eq:renormalisation-flow-18} follows from the first.

The numbers $c_1$, $B_u$, $C_{\boldsymbol{\Upsilon}}$ and $u_{-}$ are absolute constants, and $B_{\mathrm{Reg}}$, $B_N$, $B'_N$, $\gamma_1$, $B_{\boldsymbol{\Lambda}}$, $C_4$ and $C_{\Sigma}$ depend only on $b_{-}$ and $b_{+}$.
\end{proof}

\subsection{Estimates on a single torus}\label{sec:estimates-single-torus}

Throughout this subsection hypothesis~\ref{def:hypothesis-flow} is assumed, and $(m_{\mathrm{in}},\boldsymbol{\varepsilon}_{\mathrm{in}})$ is a real pair that is controlled on $R$ (Definition~\ref{def:fluctuation-field-controlled}). As in Section~\ref{sec:notation-standing-hypothesis}, $(\boldsymbol{\varepsilon}^t_k,m_k,E_k)$ is its trajectory on $R$ and $(\mathscr{E}_k,\mathfrak{P}_k)$ its trajectory with slot, $0\le k\le k^{*}$, and the subscript $*$ marks the values at the level $k^{*}=k^{*}(R)$. Proposition~\ref{prop:existence-trajectory}, taken at $k=k^{*}$, and the definition of a controlled pair give
\[
E_{*}\in K_{k^{*}},\qquad\left\|E_{*}\right\|_{k^{*}}\le w^{\mathrm{E}}_{*},\qquad\left|m_{*}\right|\le w^{\mathrm{m}}_{*}.
\]

Let $D_{*}=D_{k^{*}}$ be as in~\eqref{eq:renormalisation-flow-12}, and for $\theta\in\mathbb{R}^3$ let $c_{\theta}$ be the twist of Definition~\ref{def:twist-external-connection}. By Proposition~\ref{prop:data-along-flow}(c), $k^{*}\ge1$ and $0<\mathfrak{m}_{*}\le\mathfrak{m}_0$, and $\mathfrak{m}_0=1$ by Definition~\ref{def:mass-bound-stop}; the integer $L$ is odd and at least $3$ (stage~\ref{it:order-choice-4}), so that $L^{k^{*}}$ is an odd integer $\ge3$; and $b_{-}\le b_{k^{*}}\le b_{+}$ (Section~\ref{sec:running-quantities-tori}). Consequently the data
\[
\left(l,\mathfrak{N},\mathfrak{m},\mathsf{b},\mathsf{e}\right):=\left(l_{*},L^{k^{*}},\mathfrak{m}_{*},b_{k^{*}},e_{*}\right)
\]
satisfy the requirements of Section~\ref{sec:averaging-blocks-flat}, and for them $D(c)=D_{*}(c_{\theta})$ and $D(0)=D_{*}(0)$. For every $\theta\in\mathbb{R}^3$ the number $R_{\theta}$ of Section~\ref{sec:averaging-blocks-flat}, formed with these data, is therefore
\[
R_{\theta}=\frac{\det D_{*}(c_{\theta})}{\det D_{*}(0)}.
\]

\subsubsection{The integral at the stop level}\label{sec:integral-stop-level}

\begin{proposition}[Estimate at the stop level]\label{prop:estimate-stop-level}
Let $\theta\in\mathbb{R}^3$, and let $A$ be a real bond function on $T^0_{k^{*}}$ with mean zero and $|dA|_{\infty}\le p_{*}$. Then $A+c_{\theta}\in\mathrm{Dom}_{k^{*}}$, and there is a complex number $\vartheta$ such that
\begin{equation}\label{eq:renormalisation-flow-19}
Y_{*}\left(A+c_{\theta}\right)=\exp\left(-\boldsymbol{\varepsilon}^t_{*}\cdot l_{*}^3\right)\cdot R_{\theta}\cdot\left(1+\vartheta\right),\qquad\left|\vartheta\right|\le\varepsilon_1=\frac{1}{2}\,e_{*}^{\frac{1}{4}-9\varepsilon}.
\end{equation}
\end{proposition}

\begin{proof}
Put $c:=c_{\theta}$. By Proposition~\ref{prop:data-along-flow}(c) and the last paragraph of Section~\ref{sec:one-step-hypothesis}, Propositions~\ref{prop:background-field}, \ref{prop:kernel-bounds}, \ref{prop:fields-domains}(i), \ref{prop:gauge-fixing-controlled}, \ref{prop:bounds-unit-lattice} and~\ref{prop:block-covariance-real} hold at the level $k^{*}$, that is, with the side $l_{*}$, the mass $\mathfrak{m}_{*}$ and the charge $e_{*}$. The last inequality in~\eqref{eq:renormalisation-group-T2} reads $e_{*}^{\frac{1}{4}-9\varepsilon}\le\frac{1}{8}$; as $\varepsilon\le\bar{\varepsilon}=\frac{1}{200}$ (stages~\ref{it:order-choice-1} and~\ref{it:order-choice-8}), the exponent $\frac{1}{4}-9\varepsilon$ is positive, and therefore $e_{*}<1$.

\emph{(1) Gauge fixing.} By Proposition~\ref{prop:depth-stop-level}(e), $l_{*}\ge M\cdot r_{*}\ge2$. Proposition~\ref{prop:gauge-fixing-controlled}, applied to $A$ on $T^0_{k^{*}}=\left(\mathbb{Z}/l_{*}\mathbb{Z}\right)^3$, provides a function $\lambda$ on the sites of $T^0_{k^{*}}$ with $\left|A-d\lambda\right|_{\infty}\le2\,l_{*}^2\cdot|dA|_{\infty}$. Since $A$ is real, replacing $\lambda$ by its real part does not increase $\left|A-d\lambda\right|_{\infty}$; so $\lambda$ may be taken real. Then $A':=A-d\lambda$ is a real bond function on $T^0_{k^{*}}$, and
\[
\left|A'\right|_{\infty}\le2\,l_{*}^2\cdot|dA|_{\infty}\le2\,l_{*}^2\cdot p_{*}.
\]
The constant bond function $c$ is closed, and $d(d\lambda)=0$; hence
\[
d\left(A'+c\right)=d\left(A+c\right)=dA.
\]
So $A+c$ and $A'+c$ are real bond functions on $T^0_{k^{*}}$ whose coboundaries have sup norm at most $p_{*}$, and both lie in $\mathrm{Dom}_{k^{*}}$ by Proposition~\ref{prop:fields-domains}(i). This is the first assertion. As $A+c=(A'+c)+d\lambda$ with $\lambda$ real, Theorem~\ref{thm:representation-twist-external-2}(b), applied with $a=0$ (so that $\mathscr{C}_{*}=0$ and $Y^x_{*}(\cdot\,;0)=Y_{*}$), gives
\[
Y_{*}(A+c)=Y_{*}(A'+c).
\]
Let $\mathscr{A}^0:=H_{k^{*}}A'$. By Proposition~\ref{prop:background-field}, $H_{k^{*}}$ is linear and real, and $H_{k^{*}}c=c$; hence $\mathscr{A}^0$ is a real bond function on $T_{k^{*}}$, and
\[
\mathscr{A}:=H_{k^{*}}\left(A'+c\right)=\mathscr{A}^0+c.
\]
The three norms of a bond function $B$ on $T_{k^{*}}$ are $|B|_{\infty}$, $|\partial B|_{\infty}$ and $|\delta_{\alpha}\partial B|_{\infty}$ (Section~\ref{sec:domains-bond-functions}). Summing the kernel bound of Proposition~\ref{prop:kernel-bounds} over the second argument by means of~\eqref{eq:operators-photon-2}, one finds that the three norms of $H_{k^{*}}B'$ are at most $C_P\cdot c_{\mathrm{bd}}(\gamma_P)\cdot|B'|_{\infty}\le C_0\cdot|B'|_{\infty}$ for every real bond function $B'$ on $T^0_{k^{*}}$, the last inequality by the definition of $C_0$ in Section~\ref{sec:constants-conditions-torus}. Taken at $B'=A'$, together with~\eqref{eq:renormalisation-group-T3}, this shows that each of the three norms of $\mathscr{A}^0$ is at most
\[
2\,C_0\cdot l_{*}^2\cdot p_{*}<\frac{1}{2}\,e_{*}^{-\varepsilon}.
\]

\emph{(2) A disc of connections.} Put $\mathrm{rad}:=e_{*}^{-\frac{3}{4}+4\varepsilon}$. Since $e_{*}<1$ and $\frac{3}{4}-4\varepsilon\ge\frac{1}{4}-9\varepsilon$, the last inequality in~\eqref{eq:renormalisation-group-T2} gives $\mathrm{rad}\ge e_{*}^{-\left(\frac{1}{4}-9\varepsilon\right)}\ge8$; in particular $\mathrm{rad}\ge2$. By Definition~\ref{def:domains}, $R_{k^{*}}$ is the set of the complex bond functions on $T_{k^{*}}$ whose three norms are smaller than $e_{*}^{-\frac{3}{4}+\varepsilon}$, $e_{*}^{-\frac{3}{4}+2\varepsilon}$ and $e_{*}^{-\frac{3}{4}+3\varepsilon}$, respectively; as $e_{*}<1$, the smallest of these three thresholds is $e_{*}^{-\frac{3}{4}+3\varepsilon}$. Let $s\in\mathbb{C}$ with $|s|<2\mathrm{rad}$. By (1), the three norms of $s\cdot\mathscr{A}^0$ are smaller than
\[
2\mathrm{rad}\cdot\frac{1}{2}\,e_{*}^{-\varepsilon}=e_{*}^{-\frac{3}{4}+3\varepsilon},
\]
so that $s\cdot\mathscr{A}^0\in R_{k^{*}}$. The set $\tilde{R}_{k^{*}}$ of Definition~\ref{def:domains} contains $R_{k^{*}}$: in its definition take the real summand equal to $0$. Lemmas~\ref{lem:inclusions-bounds} and~\ref{lem:invariances} hold at the level $k^{*}$, by the last paragraph of Section~\ref{sec:one-step-hypothesis}. By Lemma~\ref{lem:invariances}, $\tilde{R}_{k^{*}}$ is invariant under the addition of real constants, and by Lemma~\ref{lem:inclusions-bounds}, $\tilde{R}_{k^{*}}\subset\tilde{G}_{k^{*}}$. Thus
\[
s\cdot\mathscr{A}^0+c\in\tilde{R}_{k^{*}}\subset\tilde{G}_{k^{*}}\qquad\left(s\in\mathbb{C},\ |s|<2\mathrm{rad}\right),
\]
and in particular $\mathscr{A}=\mathscr{A}^0+c\in\tilde{R}_{k^{*}}$.

\emph{(3) Change of covariance.} Let $D_{\theta}:=D_{*}(c)$. As $c$ is real, Proposition~\ref{prop:block-covariance-real} shows that $D_{\theta}$ is invertible and that
\[
\operatorname{Cov}_{\theta}:=D_{\theta}^{-1}
\]
has operator norm at most $b_{k^{*}}^{-1}+\mathfrak{m}_{*}^{-1}\le b_{-}^{-1}+\mathfrak{m}_{*}^{-1}$. Let $\mathrm{Gr}[\Psi]$ be the Grassmann algebra of Section~\ref{sec:grassmann-algebras-norms} at the level $k^{*}$, formed with the components of $\Psi$ and $\bar{\Psi}$ at the sites of $T^0_{k^{*}}$, and let $\langle\bar{\Psi},O\Psi\rangle$ denote the pairing of Section~\ref{sec:fermion-integrals-blocking}. By (2), the element
\[
g(s):=\left\langle\bar{\Psi},D_{*}\left(s\cdot\mathscr{A}^0+c\right)\Psi\right\rangle\in\mathrm{Gr}[\Psi]
\]
is defined for $s\in\mathbb{C}$ with $|s|<2\mathrm{rad}$, and $g(0)=\langle\bar{\Psi},D_{\theta}\Psi\rangle$, $g(1)=\langle\bar{\Psi},D_{*}(\mathscr{A})\Psi\rangle$. By the definition of the Gaussian integral $\mu_{\operatorname{Cov}}$ in Proposition~\ref{prop:weighted-coefficient-norms}, taken with $\operatorname{Cov}=\operatorname{Cov}_{\theta}$, whose inverse is $D_{\theta}$ and whose determinant is $\left(\det D_{\theta}\right)^{-1}$,
\[
\int\mathrm{D}\Psi\,\exp\left(-\left\langle\bar{\Psi},D_{\theta}\Psi\right\rangle\right)F=\det D_{\theta}\cdot\int\mathrm{d}\mu_{\operatorname{Cov}_{\theta}}(\Psi)\,F\qquad\left(F\in\mathrm{Gr}[\Psi]\right).
\]
Now apply~\eqref{eq:renormalisation-flow-12} at $A'+c\in\mathrm{Dom}_{k^{*}}$, where the connection is $H_{k^{*}}(A'+c)=\mathscr{A}$, and write $-\langle\bar{\Psi},D_{*}(\mathscr{A})\Psi\rangle=-g(0)-\left(g(1)-g(0)\right)$ in the exponent. All terms of the exponent are even elements of $\mathrm{Gr}[\Psi]$, which commute with each other; so the exponential of their sum is the product of their exponentials. Since $\det D_{\theta}/\det D_{*}(0)=R_{\theta}$, it follows that
\begin{align*}
Y_{*}(A'+c)&=\exp\left(-\boldsymbol{\varepsilon}^t_{*}\cdot l_{*}^3\right)\cdot R_{\theta}\cdot\int\mathrm{d}\mu_{\operatorname{Cov}_{\theta}}(\Psi)\,\exp(\operatorname{Ev}),\\
\operatorname{Ev}&:=-\left(g(1)-g(0)\right)-m_{*}\cdot\langle\bar{\psi},\psi\rangle+\sum_XE_{*}\left(X,\mathscr{A},\psi^{\#}\right),\qquad\psi^{\#}=H^{\#}_{k^{*}}(\mathscr{A})\Psi,
\end{align*}
where, as in~\eqref{eq:renormalisation-flow-12}, $\psi=H_{k^{*}}(\mathscr{A})\Psi$ and $\bar{\psi}=\bar{H}_{k^{*}}(\mathscr{A})\bar{\Psi}$ are the first components of $\psi^{\#}$, and $X$ ranges over the polymers of $D_{k^{*}}$. The element $\operatorname{Ev}$ of $\mathrm{Gr}[\Psi]$ is even.

\emph{(4) The quadratic term.} For $a>0$ let $|\cdot|_a$ be the norm on $\mathrm{Gr}[\Psi]$ defined in Section~\ref{sec:grassmann-algebras-norms}. If an element $F$ of $\mathrm{Gr}[\Psi]$ is written as
\[
F=\sum_{j\ge0}\frac{1}{j!}\sum_{u_1,\dots,u_j}F_j(u_1,\dots,u_j)\cdot\Psi(u_1)\cdots\Psi(u_j),
\]
with kernels $F_j$ that are antisymmetric in their arguments, the indices $u_1,\dots,u_j$ ranging over the triples $u=(y,\boldsymbol{\beta},\boldsymbol{\omega})$ of a site $y$ of $T^0_{k^{*}}$, a spinor index $\boldsymbol{\beta}$ and $\boldsymbol{\omega}\in\{0,1\}$, then
\[
|F|_a=\sum_{j\ge0}\frac{a^j}{j!}\sum_{u_1,\dots,u_j}\left|F_j(u_1,\dots,u_j)\right|.
\]
Let
\[
a_1:=e_{*}^{-\frac{1}{4}},\qquad a_2:=C_0^{-1}\cdot e_{*}^{-\frac{1}{4}+\varepsilon}
\]
be the weights of Proposition~\ref{prop:bounds-unit-lattice} at the level $k^{*}$. By (2), the affine map $s\mapsto s\cdot\mathscr{A}^0+c$ maps the disc $\left\{s\in\mathbb{C}:|s|<2\mathrm{rad}\right\}$ into $\tilde{G}_{k^{*}}$. By Proposition~\ref{prop:bounds-unit-lattice}(a), the entries of $D_{*}(B)$ are analytic functions of $B\in\tilde{G}_{k^{*}}$, and $\left|\langle\bar{\Psi},D_{*}(B)\Psi\rangle\right|_{a_1}\le C_0\cdot e_{*}^{-\frac{1}{2}}\cdot l_{*}^3$ for every $B\in\tilde{G}_{k^{*}}$. Hence the coefficients of $g(s)$ are analytic functions of $s$ on this disc, and
\[
\left|g(s)\right|_{a_1}\le C_0\cdot e_{*}^{-\frac{1}{2}}\cdot l_{*}^3\qquad\left(|s|<2\mathrm{rad}\right).
\]
Since $\frac{1}{s(s-1)}=\frac{1}{s-1}-\frac{1}{s}$ and $1<\mathrm{rad}<2\mathrm{rad}$, Cauchy's formula, applied to each coefficient, gives
\[
g(1)-g(0)=\frac{1}{2\pi i}\oint_{|s|=\mathrm{rad}}\mathrm{d}s\,\frac{g(s)}{s(s-1)}.
\]
On the circle $|s|=\mathrm{rad}$, whose length is $2\pi\mathrm{rad}$, one has $|s(s-1)|\ge\mathrm{rad}\cdot(\mathrm{rad}-1)$; moreover $\mathrm{rad}-1\ge\frac{1}{2}\mathrm{rad}$, because $\mathrm{rad}\ge2$. Therefore
\[
\left|g(1)-g(0)\right|_{a_1}\le\frac{C_0\cdot e_{*}^{-\frac{1}{2}}\cdot l_{*}^3}{\mathrm{rad}-1}\le\frac{2\,C_0\cdot e_{*}^{-\frac{1}{2}}\cdot l_{*}^3}{\mathrm{rad}}=2\,C_0\cdot e_{*}^{\frac{1}{4}-4\varepsilon}\cdot l_{*}^3.
\]

\emph{(5) The remaining terms.} By (2), $\mathscr{A}\in\tilde{R}_{k^{*}}$; moreover $0\in\mathfrak{X}_{k^{*}}$, and $\mathscr{E}_{*}\in K^x_{k^{*}}$ with $\mathscr{E}_{*}|_0=E_{*}$ (Proposition~\ref{prop:existence-trajectory}). Proposition~\ref{prop:bounds-unit-lattice}(c), applied to $\mathscr{E}_{*}$ at the pair $(\mathscr{A},0)$, whose total connection is $\mathscr{A}$, therefore gives
\[
\left|E_{*}\left(X,\mathscr{A},\psi^{\#}\right)\right|_{a_2}\le\left\|E_{*}(X,\mathscr{A})\right\|_{h_{k^{*}}}\qquad\left(X\in D_{k^{*}}\right),
\]
with the weights $h_{k^{*}}=\left(e_{*}^{-\frac{1}{4}},e_{*}^{-\frac{1}{4}+\varepsilon}\right)$ of Section~\ref{sec:grassmann-algebras-norms}. By the definition of the norm of $K_{k^{*}}$ (Definition~\ref{def:renormalisation-group-1}) and by $\|E_{*}\|_{k^{*}}\le w^{\mathrm{E}}_{*}$, the right side is at most $\exp\left(-\kappa\cdot d_M(X)\right)\cdot w^{\mathrm{E}}_{*}$. Every polymer contains an $M$-cube, and the number of the $M$-cubes of $T_{k^{*}}$ is $(l_{*}/M)^3\le l_{*}^3$. So the sum of $\exp\left(-\kappa\cdot d_M(X)\right)$ over all polymers $X$ is at most the sum, over the $M$-cubes $\square$, of the sums over the polymers $X\supseteq\square$; each of the latter is at most $c_{\kappa}$, by the definition of $c_{\kappa}$ in Section~\ref{sec:constants-conditions-torus} and since $l_{*}\in\mathrm{Cl}$ (Proposition~\ref{prop:data-along-flow}(c)). Hence
\[
\left|\sum_XE_{*}\left(X,\mathscr{A},\psi^{\#}\right)\right|_{a_2}\le w^{\mathrm{E}}_{*}\cdot\sum_X\exp\left(-\kappa\cdot d_M(X)\right)\le c_{\kappa}\cdot w^{\mathrm{E}}_{*}\cdot l_{*}^3.
\]
By Proposition~\ref{prop:bounds-unit-lattice}(b), taken at $B=\mathscr{A}\in\tilde{G}_{k^{*}}$, where $\psi=H_{k^{*}}(\mathscr{A})\Psi$ and $\bar{\psi}=\bar{H}_{k^{*}}(\mathscr{A})\bar{\Psi}$, and by $|m_{*}|\le w^{\mathrm{m}}_{*}=e_{*}^{\frac{3}{4}-8\varepsilon}$,
\[
\left|m_{*}\cdot\langle\bar{\psi},\psi\rangle\right|_{a_2}\le4\,w^{\mathrm{m}}_{*}\cdot e_{*}^{-\frac{1}{2}}\cdot l_{*}^3=4\,e_{*}^{\frac{1}{4}-8\varepsilon}\cdot l_{*}^3.
\]
As $C_0\ge1$ and $e_{*}<1$, one has $a_2\le a_1$, and $|F|_a$ is non-decreasing in $a$ (Proposition~\ref{prop:weighted-coefficient-norms}(N6)); so (4) gives $\left|g(1)-g(0)\right|_{a_2}\le2\,C_0\cdot e_{*}^{\frac{1}{4}-4\varepsilon}\cdot l_{*}^3$. Finally $e_{*}^{\frac{1}{4}-4\varepsilon}\le e_{*}^{\frac{1}{4}-8\varepsilon}$ and $w^{\mathrm{E}}_{*}=e_{*}^{\frac{1}{4}-7\varepsilon}\le e_{*}^{\frac{1}{4}-8\varepsilon}$, because $e_{*}<1$. Adding the three bounds one obtains
\[
\left|\operatorname{Ev}\right|_{a_2}\le\left(2\,C_0+c_{\kappa}+4\right)\cdot e_{*}^{\frac{1}{4}-8\varepsilon}\cdot l_{*}^3.
\]

\emph{(6) Gaussian integration.} By Proposition~\ref{prop:weighted-coefficient-norms}\eqref{it:weighted-coefficient-norms-N7}(iii), applied to the invertible matrix $\operatorname{Cov}_{\theta}$,
\[
\int\mathrm{d}\mu_{\operatorname{Cov}_{\theta}}(\Psi)\,\exp(\operatorname{Ev})=\int\mathrm{d}\mu_I(\Psi)\,\exp(\operatorname{Ev}'),\qquad\operatorname{Ev}'\left(\bar{\Psi},\Psi\right):=\operatorname{Ev}\left(\bar{\Psi},\operatorname{Cov}_{\theta}\Psi\right);
\]
here $\operatorname{Ev}'$ is the image of $\operatorname{Ev}$ under the homomorphism of $\mathrm{Gr}[\Psi]$ that fixes the components of $\bar{\Psi}$ and replaces $\Psi$ by $\operatorname{Cov}_{\theta}\Psi$, and this homomorphism maps $\exp(\operatorname{Ev})$ to $\exp(\operatorname{Ev}')$. Let $|\operatorname{Cov}_{\theta}|_{1,\infty}$ be formed as in Proposition~\ref{prop:weighted-coefficient-norms}, the pairs $(y,\boldsymbol{\beta})$ that index the rows and the columns of $\operatorname{Cov}_{\theta}$ carrying the counting weight; thus
\[
|\operatorname{Cov}_{\theta}|_{1,\infty}=\max_{(y,\boldsymbol{\beta})}\sum_{(y',\boldsymbol{\beta}')}\left|\operatorname{Cov}_{\theta}\left((y,\boldsymbol{\beta}),(y',\boldsymbol{\beta}')\right)\right|,
\]
where $y,y'$ range over the sites of $T^0_{k^{*}}$ and $\boldsymbol{\beta},\boldsymbol{\beta}'$ over the four spinor indices. Put
\[
c':=\max\left(1,|\operatorname{Cov}_{\theta}|_{1,\infty}\right).
\]
The substitution has the row sum $1$ on the components of $\bar{\Psi}$, and row sums at most $|\operatorname{Cov}_{\theta}|_{1,\infty}$ on those of $\Psi$; both are at most $c'$. Proposition~\ref{prop:weighted-coefficient-norms}(N2), with the weight $c'$ for all components of $\Psi$ and $\bar{\Psi}$ in $\operatorname{Ev}$ and the weight $1$ for all of them in $\operatorname{Ev}'$, therefore gives $|\operatorname{Ev}'|_1\le|\operatorname{Ev}|_{c'}$. A row of $\operatorname{Cov}_{\theta}$ has $4\,l_{*}^3$ entries, one for each site of $T^0_{k^{*}}$ and each of the four spinor indices, and every entry has modulus at most the operator norm of $\operatorname{Cov}_{\theta}$. By (3), $|\operatorname{Cov}_{\theta}|_{1,\infty}\le4\,l_{*}^3\cdot\left(b_{-}^{-1}+\mathfrak{m}_{*}^{-1}\right)$, and this bound is at least $1$, since $\mathfrak{m}_{*}\le1$. With the first inequality in~\eqref{eq:renormalisation-group-T2} it follows that
\[
c'\le4\,l_{*}^3\cdot\left(b_{-}^{-1}+\mathfrak{m}_{*}^{-1}\right)\le\frac{e_{*}^{-\frac{\varepsilon}{2}}}{2\,C_0^2\cdot l_{*}^3}\le C_0^{-1}\cdot e_{*}^{-\frac{1}{4}+\varepsilon}=a_2;
\]
the last inequality amounts to $e_{*}^{\frac{1}{4}-\frac{3}{2}\varepsilon}\le2\,C_0\cdot l_{*}^3$, which holds because $e_{*}<1$, $\frac{3}{2}\varepsilon<\frac{1}{4}$, $C_0\ge1$ and $l_{*}\ge1$. The monotonicity of $|\cdot|_a$ in $a$ and (5) now give
\[
|\operatorname{Ev}'|_1\le|\operatorname{Ev}|_{c'}\le|\operatorname{Ev}|_{a_2}\le\left(2\,C_0+c_{\kappa}+4\right)\cdot e_{*}^{\frac{1}{4}-8\varepsilon}\cdot l_{*}^3\le\frac{1}{4}\,e_{*}^{\frac{1}{4}-9\varepsilon}\le1;
\]
the fourth inequality follows from the second inequality in~\eqref{eq:renormalisation-group-T2}, because $2\,C_0+c_{\kappa}+4\le2\,C_0+c_{\kappa}+5\,c_{\mathrm{fl}}+6$ (recall $c_{\mathrm{fl}}\ge1$, stage~\ref{it:order-choice-1}), and the last one from $e_{*}<1$. The element $\operatorname{Ev}'$ is even, since $\operatorname{Ev}$ is even. By Proposition~\ref{prop:weighted-coefficient-norms}\eqref{it:weighted-coefficient-norms-N7}(iii), taken with the identity matrix, $\int\mathrm{d}\mu_I(\Psi)\,1=1$. Hence
\[
\left|\int\mathrm{d}\mu_I(\Psi)\,\exp(\operatorname{Ev}')-1\right|=\left|\int\mathrm{d}\mu_I(\Psi)\left(\exp(\operatorname{Ev}')-1\right)\right|\le\left|\exp(\operatorname{Ev}')-1\right|_1\le\exp\left(|\operatorname{Ev}'|_1\right)-1\le2\,|\operatorname{Ev}'|_1.
\]
The first inequality is Proposition~\ref{prop:weighted-coefficient-norms}\eqref{it:weighted-coefficient-norms-N7}(ii); the second is the bound of Proposition~\ref{prop:weighted-coefficient-norms}(N1) for even elements, which rests on the submultiplicativity of $|\cdot|_1$ (Proposition~\ref{prop:weighted-coefficient-norms}\eqref{it:weighted-coefficient-norms-N7}(i)); the third holds because $0\le|\operatorname{Ev}'|_1\le1$ and, by convexity, $\exp(x)-1\le\left(\exp(1)-1\right)x\le2x$ for $0\le x\le1$.

\emph{Conclusion.} Let $\vartheta:=\int\mathrm{d}\mu_{\operatorname{Cov}_{\theta}}(\Psi)\,\exp(\operatorname{Ev})-1$. By (1), (3) and (6),
\[
Y_{*}(A+c)=Y_{*}(A'+c)=\exp\left(-\boldsymbol{\varepsilon}^t_{*}\cdot l_{*}^3\right)\cdot R_{\theta}\cdot\left(1+\vartheta\right),\qquad\left|\vartheta\right|\le2\,|\operatorname{Ev}'|_1\le\frac{1}{2}\,e_{*}^{\frac{1}{4}-9\varepsilon}=\varepsilon_1,
\]
which is~\eqref{eq:renormalisation-flow-19}.
\end{proof}

\subsubsection{Large fields and the averaged law}\label{sec:large-fields-averaged}

\begin{lemma}[Positivity and large fields]\label{lem:positivity-large-fields}
For every $\theta\in\mathbb{R}^3$ the measure $\nu^{\theta}$ on $F_N$ is non-negative. The number $\bar{\nu}(1)$ is positive and finite, $\bar{\mu}$ is a probability measure on $F_N$, and
\begin{equation}\label{eq:renormalisation-flow-20}
\bar{\mu}\left(F_N\setminus\mathfrak{G}\right)\le\varepsilon_3,\qquad\mathrm{d}\mathsf{m}_{*}\left(\chi_{*}=0\right)\le\varepsilon_4.
\end{equation}
\end{lemma}

\begin{proof}
\emph{Positivity.} Let $n=R/\eta=R\cdot L^N$ and $q=e\,\eta^{\frac{1}{2}}$ be as in Section~\ref{sec:notation-standing-hypothesis}. Since $N\ge N_{\mathrm{min}}(R,e)$, one has $n=R\cdot L^N\ge L\cdot M\cdot r_0\ge2$; and $q\ne0$, because $e>0$. Put
\[
\boldsymbol{\mu}:=\bar{m}\cdot\eta+m_{\mathrm{in}}.
\]
Since the pair is controlled on $R$, Definition~\ref{def:fluctuation-field-controlled} gives $|m_k|\le w^{\mathrm{m}}_k$ at the level $k=0$, at which the symbol $m_k$ stands for $m_{\mathrm{in}}$ (Section~\ref{sec:notation-standing-hypothesis}); hence $|m_{\mathrm{in}}|\le w^{\mathrm{m}}_0$, where $w^{\mathrm{m}}_0$ is formed with $q$. Moreover $w^{\mathrm{m}}_0<1$ by~\eqref{eq:renormalisation-group-T1}; as $\bar{m}\cdot\eta>0$, it follows that $\boldsymbol{\mu}>-1$. The matrices
\[
\gamma_{\nu}=\sigma_3\otimes\sigma_{\nu+1}=\operatorname{diag}\left(\sigma_{\nu+1},-\sigma_{\nu+1}\right)\qquad(\nu=0,1,2)
\]
of Section~\ref{sec:dirac-matrices-lattice} are the matrices of the four-component representation of Section~\ref{sec:clifford-triples}, the direction $\nu$ corresponding to the direction $\nu+1$ there; and the operator~\eqref{eq:gaussian-bounds-1} at the antiperiodic character $\chi^{\mathrm{ap}}$ is the operator $D_A$ of Section~\ref{sec:lattice-measures-laws}. So Theorems~\ref{thm:gaussian-domination-positive} and~\ref{thm:limit-vanishing-proca} and Corollary~\ref{cor:linear-functionals-curvature} apply on $\Lambda_n$, in the four-component representation, with $n\ge2$, the charge $q\ne0$ and the lattice mass $\boldsymbol{\mu}>-1$.

By Theorem~\ref{thm:gaussian-domination-positive}(b), $\det\left(D_A+\boldsymbol{\mu}\right)\ge0$ for every real one-cochain $A$ on $\Lambda_n$. Lemma~\ref{lem:product-formula-volume}, taken at $\theta=0$ with the mass $\bar{m}\cdot\eta\in(0,1]$, gives $\det\left(D_0+\bar{m}\cdot\eta\right)>0$. As $A(F)+H_{\theta}$ is real, it follows that
\[
\rho^{\theta,0}_N(F)=\exp\left(-\boldsymbol{\varepsilon}_{\mathrm{in}}\cdot n^3\right)\cdot\frac{\det\left(D_{A(F)+H_{\theta}}+\boldsymbol{\mu}\right)}{\det\left(D_0+\bar{m}\cdot\eta\right)}\ge0\qquad\left(F\in F_N,\ \theta\in\mathbb{R}^3\right).
\]
Moreover $\rho^{\theta,0}_N$ is bounded on $F_N$: the link phases of $D_{A(F)+H_{\theta}}$ have modulus one, so that the entries of the matrix $D_{A(F)+H_{\theta}}+\boldsymbol{\mu}$ are bounded independently of $F$. Hence $\nu^{\theta}=\rho^{\theta,0}_N\cdot\gamma_N$ is a non-negative finite measure on $F_N$.

\emph{The averaged law.} As in Section~\ref{sec:notation-standing-hypothesis}, $C^2$ is the space of real two-cochains on $\Lambda_n$, with the counting scalar product $\langle\cdot,\cdot\rangle$ and its norm $|\cdot|$; thus $F_N=dC^1\subset C^2$. For every Borel function $G$ on $C^2$ of at most exponential growth, that is, such that $|G(\omega)|\exp\left(-t\,|\omega|\right)$ is bounded on $C^2$ for some number $t\ge0$, put
\[
\overline{Z}[G]:=\int_{[-\pi,\pi]^3}\frac{\mathrm{d}\theta}{(2\pi)^3}\int_{V_{\mathrm{t}}}\mathrm{d}A\,G(dA)\exp\left(-\frac{1}{2}|dA|^2\right)\det\left(D_{A+H_{\theta}}+\boldsymbol{\mu}\right),
\]
where $\mathrm{d}A$ is the Lebesgue measure of the Euclidean space $V_{\mathrm{t}}$. This is the functional of~\eqref{eq:gaussian-bounds-3}, formed with the lattice mass $\boldsymbol{\mu}$. Indeed, the cochain $c^{\theta}$ of Section~\ref{sec:proca-parent-harmonic} is $H_{\theta}$; the coboundary $d_1$, the norm of Section~\ref{sec:cochains-periodic-lattice} and the measure $\sigma_{V_{\mathrm{t}}}$ are $d$, $|\cdot|$ and $\mathrm{d}A$; and, by the identification of the matrices $\gamma_{\nu}$ and of the Wilson operator made above, $\Delta^{\mathrm{ap}}\left(A+c^{\theta}\right)=\det\left(D_{A+H_{\theta}}+\boldsymbol{\mu}\right)$ for $A\in V_{\mathrm{t}}$. By Theorem~\ref{thm:limit-vanishing-proca}, $0<\overline{Z}[1]<\infty$, and the integral defining $\overline{Z}[G]$ converges absolutely for every such $G$; put $\overline{\mathsf{E}}[G]:=\overline{Z}[G]/\overline{Z}[1]$.

The measure $\gamma_N$ is the image under $d$ of the probability measure on $V_{\mathrm{t}}$ whose density with respect to $\mathrm{d}A$ is proportional to $\exp\left(-\frac{1}{2}|dA|^2\right)$, and $A(dA)=A$ for $A\in V_{\mathrm{t}}$. By the definitions of $\nu^{\theta}$ and $\bar{\nu}$ in Section~\ref{sec:notation-standing-hypothesis}, every bounded Borel function $G$ on $C^2$ therefore satisfies
\[
\bar{\nu}(G)=\exp\left(-\boldsymbol{\varepsilon}_{\mathrm{in}}\cdot n^3\right)\cdot\left(\det\left(D_0+\bar{m}\cdot\eta\right)\cdot\int_{V_{\mathrm{t}}}\mathrm{d}A\,\exp\left(-\frac{1}{2}|dA|^2\right)\right)^{-1}\cdot\overline{Z}[G],
\]
where $\bar{\nu}(G)$ is the integral of the restriction of $G$ to $F_N$; the factor in front of $\overline{Z}[G]$ is a positive number. With $G=1$ it follows that $\bar{\nu}(1)$ is positive and finite. Hence $\bar{\mu}=\bar{\nu}/\bar{\nu}(1)$ is a probability measure on $F_N$, and $\bar{\mu}(G)=\overline{\mathsf{E}}[G]$ for every bounded Borel function $G$ on $C^2$; that is, $\bar{\mu}$ is the harmonic-averaged law of Section~\ref{sec:proca-parent-harmonic}. The orthogonal projection of $C^2$ onto $F_N$ fixes every element of $F_N$. Therefore Corollary~\ref{cor:linear-functionals-curvature}(iii) gives the tail bound
\[
\bar{\mu}\left(\left\{F\in F_N:\langle\phi,F\rangle\ge y\right\}\right)\le\exp\left(-\frac{y^2}{2|\phi|^2}\right)\qquad\left(y\ge0,\ \phi\in F_N\setminus\{0\}\right).
\]

\emph{The first bound.} By Proposition~\ref{prop:orthogonal-coordinates}, for $j<k^{*}$ and every element $\Upsilon$ of the basis of $V_j$ one has $\zeta_j(\Upsilon)=\langle\phi_{j,\Upsilon},F\rangle$, where $\{\phi_{j,\Upsilon}\}$ is an orthonormal family in $F_N$; in particular $|\phi_{j,\Upsilon}|=1$. For every plaquette $p$ of $T^0_{k^{*}}$ one has $(dA_{*})(p)=\langle\phi_p,F\rangle$ with $\phi_p\in F_N$ and $|\phi_p|\le C_Q=1$, the constant $C_Q=1$ being that of Proposition~\ref{prop:block-averages}. Let $\phi\in F_N$ with $|\phi|\le1$, and let $y\ge0$. If $\phi\ne0$, the tail bound, applied to $\phi$ and to $-\phi$, gives
\[
\bar{\mu}\left(\left\{F\in F_N:\left|\langle\phi,F\rangle\right|>y\right\}\right)\le2\exp\left(-\frac{y^2}{2|\phi|^2}\right)\le2\exp\left(-\frac{y^2}{2}\right);
\]
if $\phi=0$, the set on the left is empty. By Definition~\ref{def:coordinates-slice-small}, $F\in F_N\setminus\mathfrak{G}$ if and only if $|\zeta_j(\Upsilon)|>p_{0,j}$ for some $j<k^{*}$ and some $\Upsilon$, or $|(dA_{*})(p)|>p_{*}$ for some plaquette $p$ of $T^0_{k^{*}}$. By Proposition~\ref{prop:basis-map-box}, the basis $\{\Upsilon\}$ of $V_j$ has $d_V\cdot n_j<2\,n_j$ elements, where $d_V=2-2L^{-3}$ (Section~\ref{sec:averaging-fluctuations-gauge}); and $T^0_{k^{*}}$ has $3\,l_{*}^3$ plaquettes. Applying the last bound to each of these functionals and adding, one obtains
\[
\bar{\mu}\left(F_N\setminus\mathfrak{G}\right)\le\sum_{j<k^{*}}4\,n_j\exp\left(-\frac{p_{0,j}^2}{2}\right)+6\,l_{*}^3\exp\left(-\frac{p_{*}^2}{2}\right).
\]
Let $j<k^{*}$ and $i:=k^{*}-j\ge1$. By the definitions of Section~\ref{sec:notation-standing-hypothesis}, $l_j=L^i\cdot l_{*}$ and $e_j=L^{-\frac{i}{2}}\cdot e_{*}$; hence
\[
n_j=l_{*}^3\cdot L^{3i}=l_{*}^3\exp\left(6\,\mathrm{lg}\cdot i\right),\qquad v_j=v_{*}+i\cdot\mathrm{lg}.
\]
As $v_j\ge v_{*}\ge8$ by~\eqref{eq:renormalisation-group-T1}, $p_{0,j}=v_j^{p_0}\ge1$, and therefore $p_{0,j}^2\ge p_{0,j}$. The function $x\mapsto x^{p_0}$ is convex on $(0,\infty)$, because $p_0\ge2$ (stage~\ref{it:order-choice-8}); so its graph lies above its tangent at $v_{*}$, and~\eqref{eq:renormalisation-group-T1} gives
\[
p_{0,j}=\left(v_{*}+i\cdot\mathrm{lg}\right)^{p_0}\ge v_{*}^{p_0}+p_0\,v_{*}^{p_0-1}\cdot i\cdot\mathrm{lg}\ge p_{0*}+28\,i\cdot\mathrm{lg}.
\]
Consequently
\[
n_j\exp\left(-\frac{p_{0,j}^2}{2}\right)\le l_{*}^3\exp\left(-\frac{p_{0*}}{2}\right)\exp\left((6-14)\cdot\mathrm{lg}\cdot i\right)=l_{*}^3\exp\left(-\frac{p_{0*}}{2}\right)L^{-4i},
\]
and $\sum_{i\ge1}L^{-4i}\le1$. Moreover $C_0\ge1$ (Section~\ref{sec:constants-conditions-torus}), so that $\exp\left(-\frac{p_{*}^2}{2}\right)\le\exp\left(-\frac{p_{*}^2}{2C_0^2}\right)$. It follows that
\[
\bar{\mu}\left(F_N\setminus\mathfrak{G}\right)\le4\,l_{*}^3\exp\left(-\frac{p_{0*}}{2}\right)+6\,l_{*}^3\exp\left(-\frac{p_{*}^2}{2}\right)\le\varepsilon_3.
\]

\emph{The second bound.} By Definition~\ref{def:coordinates-slice-small}, $\mathrm{d}\mathsf{m}_{*}$ is the law of $A_{*}$ under $\gamma_N$, and $\chi_{*}(A_{*})=0$ means $|dA_{*}|_{\infty}>p_{*}$. Hence
\[
\mathrm{d}\mathsf{m}_{*}\left(\chi_{*}=0\right)\le\sum_p\gamma_N\left(\left\{F\in F_N:\left|(dA_{*})(p)\right|>p_{*}\right\}\right),
\]
the sum extending over the $3\,l_{*}^3$ plaquettes $p$ of $T^0_{k^{*}}$. Under the standard Gaussian measure $\gamma_N$ on $F_N$ the variable $(dA_{*})(p)=\langle\phi_p,F\rangle$ is centred Gaussian with variance $|\phi_p|^2\le1$. If $\phi_p=0$, the corresponding term vanishes. Otherwise $(dA_{*})(p)=|\phi_p|\cdot g$ with a standard normal variable $g$; for $y\ge0$ the probability of $\{|g|>y\}$ is at most $2\exp\left(-\frac{y^2}{2}\right)$, so that the term is at most $2\exp\left(-\frac{p_{*}^2}{2|\phi_p|^2}\right)\le2\exp\left(-\frac{p_{*}^2}{2}\right)$. With $C_0\ge1$ it follows that
\[
\mathrm{d}\mathsf{m}_{*}\left(\chi_{*}=0\right)\le6\,l_{*}^3\exp\left(-\frac{p_{*}^2}{2}\right)\le\varepsilon_4.
\]
\end{proof}

\subsubsection{The average over the twist}\label{sec:average-twist}

\begin{theorem}[The twist average on one torus]\label{thm:twist-average-one}
Assume hypothesis~\ref{def:hypothesis-flow}. Let $(m_{\mathrm{in}}, \boldsymbol{\varepsilon}_{\mathrm{in}})$ be a real pair that is controlled on $R$ (Definition~\ref{def:fluctuation-field-controlled}), and let $\boldsymbol{\varepsilon}^t_{*} = \boldsymbol{\varepsilon}^t_{k^{*}}$ be the value at the stop level $k^{*} = k^{*}(R)$ of the first component of its trajectory on $R$. Let $R_{\theta}$ be the number defined in Section~\ref{sec:estimates-single-torus}, let $\varepsilon_1, \dots, \varepsilon_4$ be the numbers and $\varepsilon_5$ the function on $(0, \infty)$ defined in Section~\ref{sec:constants-conditions-torus}, let $p^0$ be the function of Section~\ref{sec:constant-twists-free}, and write $\rho_n$ for the number $\rho_n\left(q, \bar{m} \cdot \eta + m_{\mathrm{in}}\right)$ of~\eqref{eq:renormalisation-flow-1}. Then, for every $\theta \in \mathbb{R}^3$,
\begin{multline}\label{eq:renormalisation-flow-21}
\exp\left(-\boldsymbol{\varepsilon}^t_{*} \cdot l_{*}^3\right) \cdot R_{\theta} \cdot \left(1 - \varepsilon_1\right) \cdot \left(1 - \varepsilon_4\right) \le \nu^{\theta}(\mathfrak{G})\\
\le \exp\left(-\boldsymbol{\varepsilon}^t_{*} \cdot l_{*}^3\right) \cdot R_{\theta} \cdot \left(1 + \varepsilon_1\right).
\end{multline}
Moreover,
\begin{align}
& \bar{\nu}(\mathfrak{G}) \le \bar{\nu}(1) \le \frac{\bar{\nu}(\mathfrak{G})}{1 - \varepsilon_3}, \label{eq:renormalisation-flow-22}\\
& \left| \log \rho_n - \log \rho^{\mathrm{free}}_n - \boldsymbol{\varepsilon}_{\mathrm{in}} \cdot n^3 + \boldsymbol{\varepsilon}^t_{*} \cdot l_{*}^3 \right| \notag\\
& \qquad \le 2\,\varepsilon_1 + \varepsilon_2 + 2\,\varepsilon_3 + 2\,\varepsilon_4 \le 2\, e_{*}^{\frac{1}{4} - 9\varepsilon}, \label{eq:renormalisation-flow-23}\\
& \left| \log \rho_n - \left( p^0(\bar{m} \cdot \eta) + \boldsymbol{\varepsilon}_{\mathrm{in}} \right) \cdot n^3 + \boldsymbol{\varepsilon}^t_{*} \cdot l_{*}^3 \right| \notag\\
& \qquad \le \varepsilon_5(\bar{m} \cdot R) + 2\, e_{*}^{\frac{1}{4} - 9\varepsilon}. \label{eq:renormalisation-flow-24}
\end{align}
\end{theorem}

\begin{proof}
By Proposition~\ref{prop:consequences}, the conditions \eqref{eq:renormalisation-group-T2}, \eqref{eq:renormalisation-group-T4} and \eqref{eq:renormalisation-group-T5} are satisfied. The last inequality of \eqref{eq:renormalisation-group-T2} states $e_{*}^{\frac{1}{4} - 9\varepsilon} \le \frac{1}{8}$; as $\frac{1}{4} - 9\varepsilon > 0$, it implies $0 < e_{*} < 1$ and $\varepsilon_1 = \frac{1}{2}\, e_{*}^{\frac{1}{4} - 9\varepsilon} \le \frac{1}{16}$. From the definitions of $\varepsilon_3$ and $\varepsilon_4$ and from \eqref{eq:renormalisation-group-T4},
\[
0 \le \varepsilon_4 \le \varepsilon_3 = 6\, l_{*}^3 \cdot \left( \exp\left(-\frac{p_{0*}}{2}\right) + \exp\left(-\frac{p_{*}^2}{2 C_0^2}\right) \right) \le \frac{e_{*}}{2} < \frac{1}{2}.
\]

\emph{Proof of~\eqref{eq:renormalisation-flow-21}.} Fix $\theta \in \mathbb{R}^3$. Let $V_{\mathrm{t}}\left(T^0_{k^{*}}\right)$ be the space of Proposition~\ref{prop:orthogonal-coordinates}; its elements are the real bond functions on $T^0_{k^{*}}$ that are coclosed and have mean zero. By Proposition~\ref{prop:orthogonal-coordinates}(ii), the field $A_{*}$ of Definition~\ref{def:coordinates-slice-small} takes its values in this space; hence $\mathrm{d}\mathsf{m}_{*}$, the law of $A_{*}$ under $\gamma_N$, is a probability measure on $V_{\mathrm{t}}\left(T^0_{k^{*}}\right)$. The external connection $a = 0$ is admissible in Theorem~\ref{thm:representation-twist-external-2}, since $0 \in \mathfrak{X}_j$ for every $j$; for it one has $\mathscr{C}_{*} = 0$ and $Y^x_{*}(\cdot\,; 0) = Y_{*}$. Part~(c) of that theorem, taken with $a = 0$ and $\Phi = 1$, gives
\[
\nu^{\theta}(\mathfrak{G}) = \int \mathrm{d}\mathsf{m}_{*}(A)\, \chi_{*}(A) \cdot Y_{*}\left(A + c_{\theta}\right),
\]
the integrand being read as $0$ where $\chi_{*}(A) = 0$. Let $A \in V_{\mathrm{t}}\left(T^0_{k^{*}}\right)$ with $\chi_{*}(A) = 1$, that is, with $|dA|_{\infty} \le p_{*}$. As $A$ is a real bond function on $T^0_{k^{*}}$ with mean zero, Proposition~\ref{prop:estimate-stop-level} applies to $A$ and provides a complex number $\vartheta(A)$ such that
\[
Y_{*}\left(A + c_{\theta}\right) = \exp\left(-\boldsymbol{\varepsilon}^t_{*} \cdot l_{*}^3\right) \cdot R_{\theta} \cdot \left(1 + \vartheta(A)\right), \qquad \left|\vartheta(A)\right| \le \varepsilon_1.
\]
Consequently
\[
\nu^{\theta}(\mathfrak{G}) = \exp\left(-\boldsymbol{\varepsilon}^t_{*} \cdot l_{*}^3\right) \cdot R_{\theta} \cdot \int \mathrm{d}\mathsf{m}_{*}(A)\, \chi_{*}(A) \cdot \left(1 + \vartheta(A)\right).
\]
The factor in front of the integral is a positive real number: $\boldsymbol{\varepsilon}^t_{*}$ is real by Proposition~\ref{prop:existence-trajectory}, and $R_{\theta} > 0$ by Lemma~\ref{lem:fourier-decomposition}(a). The left side is real, because $\nu^{\theta} \ge 0$ (Lemma~\ref{lem:positivity-large-fields}). Hence the integral is a real number, and it equals its real part, in which $\vartheta(A)$ is replaced by $\operatorname{Re} \vartheta(A) \in \left[-\varepsilon_1, \varepsilon_1\right]$. It follows that
\[
\left(1 - \varepsilon_1\right) \cdot \mathrm{d}\mathsf{m}_{*}\left(\chi_{*} = 1\right) \le \int \mathrm{d}\mathsf{m}_{*}(A)\, \chi_{*}(A) \cdot \left(1 + \vartheta(A)\right) \le \left(1 + \varepsilon_1\right) \cdot \mathrm{d}\mathsf{m}_{*}\left(\chi_{*} = 1\right).
\]
As $\mathrm{d}\mathsf{m}_{*}$ is a probability measure, the second inequality of~\eqref{eq:renormalisation-flow-20} gives
\[
1 - \varepsilon_4 \le 1 - \mathrm{d}\mathsf{m}_{*}\left(\chi_{*} = 0\right) = \mathrm{d}\mathsf{m}_{*}\left(\chi_{*} = 1\right) \le 1.
\]
Since $1 - \varepsilon_1 > 0$, the last three displays together yield~\eqref{eq:renormalisation-flow-21}.

\emph{Proof of~\eqref{eq:renormalisation-flow-22}.} By Lemma~\ref{lem:positivity-large-fields}, $\nu^{\theta} \ge 0$ for every $\theta$, the number $\bar{\nu}(1)$ is positive and finite, and $\bar{\mu} = \bar{\nu}/\bar{\nu}(1)$ is a probability measure on $F_N$. Thus $\bar{\nu}$ is a positive measure, and $\bar{\nu}(\mathfrak{G}) \le \bar{\nu}(F_N) = \bar{\nu}(1)$. The first inequality of~\eqref{eq:renormalisation-flow-20} gives
\[
\bar{\nu}(1) - \bar{\nu}(\mathfrak{G}) = \bar{\nu}\left(F_N \setminus \mathfrak{G}\right) = \bar{\nu}(1) \cdot \bar{\mu}\left(F_N \setminus \mathfrak{G}\right) \le \varepsilon_3 \cdot \bar{\nu}(1),
\]
that is, $\left(1 - \varepsilon_3\right) \cdot \bar{\nu}(1) \le \bar{\nu}(\mathfrak{G})$. Since $\varepsilon_3 < 1$, this is the second inequality of~\eqref{eq:renormalisation-flow-22}.

\emph{Proof of~\eqref{eq:renormalisation-flow-23}.} As defined in Section~\ref{sec:estimates-single-torus}, $R_{\theta}$ is the number of Section~\ref{sec:averaging-blocks-flat} for the side $l_{*}$, the odd integer $\mathfrak{N} = L^{k^{*}} \ge 3$, the mass $\mathfrak{m}_{*} \in (0, 1]$, the parameter $b_{k^{*}} \in \left[b_{-}, b_{+}\right]$ and the charge $e_{*}$; for these data the constant connection $c$ of Section~\ref{sec:averaging-blocks-flat} is $c_{\theta}$, and the spacing of the fine lattice there is $\boldsymbol{\eta} = 1/\mathfrak{N} = L^{-k^{*}}$. Since $a_{*} = L^{k^{*} - N}$, $\mathfrak{m}_{*} = \bar{m} \cdot a_{*}$ and $\eta = L^{-N}$ (Sections~\ref{sec:lattice-measures-laws} and~\ref{sec:running-quantities-tori}),
\[
\boldsymbol{\eta} \cdot \mathfrak{m}_{*} = L^{-k^{*}} \cdot \bar{m} \cdot a_{*} = \bar{m} \cdot \eta, \qquad l_{*} \cdot L^{k^{*}} = \frac{R}{a_{*}} \cdot L^{k^{*}} = R \cdot L^N = n.
\]
So the lattice of Lemma~\ref{lem:fourier-decomposition}(c) is $\Lambda_n$, and its mass is $\bar{m} \cdot \eta$. Write $T_{c_{\theta}}$ and $T_0$ for the operator $T_c$ of Section~\ref{sec:averaging-blocks-flat} at the constant connections $c = c_{\theta}$ and $c = 0$. Lemma~\ref{lem:fourier-decomposition}(c) gives
\[
\det\left(D_0 + \bar{m} \cdot \eta\right) \cdot R_{\theta} = \det\left(D_{H_{\theta}} + \bar{m} \cdot \eta\right) \cdot \frac{\det T_0}{\det T_{c_{\theta}}},
\]
where $\det T_0$ and $\det T_{c_{\theta}}$ are positive by Lemma~\ref{lem:fourier-decomposition}(b). The data listed above satisfy the hypotheses of Lemma~\ref{lem:dependence-twist-uniform}, and the first inequality of~\eqref{eq:renormalisation-flow-18} gives
\[
\left| \log \frac{\det T_0}{\det T_{c_{\theta}}} \right| \le C_{\Sigma} \cdot l_{*}^{-1} \cdot \exp\left(-\gamma_1 \cdot l_{*}\right) = \varepsilon_2.
\]
As $0 < \bar{m} \cdot \eta \le 1$, Lemma~\ref{lem:product-formula-volume} shows that $\det\left(D_0 + \bar{m} \cdot \eta\right)$ and $\det\left(D_{H_{\theta}} + \bar{m} \cdot \eta\right)$ are positive. Multiplying~\eqref{eq:renormalisation-flow-21} by $\det\left(D_0 + \bar{m} \cdot \eta\right)$ and inserting the last two displays, one finds for every $\theta \in \mathbb{R}^3$
\begin{align*}
& \exp\left(-\varepsilon_2\right) \cdot \left(1 - \varepsilon_1\right) \cdot \left(1 - \varepsilon_4\right) \cdot \exp\left(-\boldsymbol{\varepsilon}^t_{*} \cdot l_{*}^3\right) \cdot \det\left(D_{H_{\theta}} + \bar{m} \cdot \eta\right)\\
& \qquad \le \det\left(D_0 + \bar{m} \cdot \eta\right) \cdot \nu^{\theta}(\mathfrak{G})\\
& \qquad \le \exp\left(\varepsilon_2\right) \cdot \left(1 + \varepsilon_1\right) \cdot \exp\left(-\boldsymbol{\varepsilon}^t_{*} \cdot l_{*}^3\right) \cdot \det\left(D_{H_{\theta}} + \bar{m} \cdot \eta\right).
\end{align*}
The function $(\theta, F) \mapsto \rho^{\theta,0}_N(F)$ is bounded and continuous on $\mathbb{R}^3 \times F_N$: up to a constant factor it is $\det\left(D_{A(F) + H_{\theta}} + \bar{m} \cdot \eta + m_{\mathrm{in}}\right)$, a polynomial in the entries of the operator $D_{A(F) + H_{\theta}}$, and these entries are bounded, the link phases of a real connection having modulus one, and depend continuously on $(\theta, F)$. By the dominated convergence theorem, the function
\[
\theta \mapsto \nu^{\theta}(\mathfrak{G}) = \int \mathrm{d}\gamma_N(F)\, \mathbf{1}_{\mathfrak{G}}(F) \cdot \rho^{\theta,0}_N(F)
\]
is continuous on $\mathbb{R}^3$; so are the outer members of the chain above. Integrating the chain over $\theta \in [-\pi, \pi]^3$ with respect to $(2\pi)^{-3}\, \mathrm{d}\theta$, and using the definitions of $\rho^{\mathrm{free}}_n$ and $\bar{\nu}$ in~\eqref{eq:renormalisation-flow-1}, one obtains
\begin{align*}
& \exp\left(-\varepsilon_2\right) \cdot \left(1 - \varepsilon_1\right) \cdot \left(1 - \varepsilon_4\right) \cdot \exp\left(-\boldsymbol{\varepsilon}^t_{*} \cdot l_{*}^3\right) \cdot \rho^{\mathrm{free}}_n\\
& \qquad \le \det\left(D_0 + \bar{m} \cdot \eta\right) \cdot \bar{\nu}(\mathfrak{G})\\
& \qquad \le \exp\left(\varepsilon_2\right) \cdot \left(1 + \varepsilon_1\right) \cdot \exp\left(-\boldsymbol{\varepsilon}^t_{*} \cdot l_{*}^3\right) \cdot \rho^{\mathrm{free}}_n.
\end{align*}
Here $\rho^{\mathrm{free}}_n > 0$, being the mean of a positive function, and $\rho_n > 0$, because~\eqref{eq:renormalisation-flow-2} gives
\[
\exp\left(-\boldsymbol{\varepsilon}_{\mathrm{in}} \cdot n^3\right) \cdot \rho_n = \det\left(D_0 + \bar{m} \cdot \eta\right) \cdot \bar{\nu}(1)
\]
and $\bar{\nu}(1) > 0$. Combined with~\eqref{eq:renormalisation-flow-22}, these relations show that the number
\begin{align*}
\Delta & := \log \rho_n - \log \rho^{\mathrm{free}}_n - \boldsymbol{\varepsilon}_{\mathrm{in}} \cdot n^3 + \boldsymbol{\varepsilon}^t_{*} \cdot l_{*}^3\\
& = \log \frac{\det\left(D_0 + \bar{m} \cdot \eta\right) \cdot \bar{\nu}(1)}{\exp\left(-\boldsymbol{\varepsilon}^t_{*} \cdot l_{*}^3\right) \cdot \rho^{\mathrm{free}}_n}
\end{align*}
satisfies
\[
-\varepsilon_2 + \log\left(1 - \varepsilon_1\right) + \log\left(1 - \varepsilon_4\right) \le \Delta \le \varepsilon_2 + \log\left(1 + \varepsilon_1\right) - \log\left(1 - \varepsilon_3\right).
\]
For $0 \le x \le \frac{1}{2}$ one has $\log(1 + x) \le x$ and $-\log(1 - x) \le 2x$; the latter holds because the convex function $x \mapsto -\log(1 - x)$ vanishes at $0$ and takes the value $\log 2 < 1$ at $\frac{1}{2}$. The numbers $\varepsilon_1$, $\varepsilon_3$, $\varepsilon_4$ lie in $\left[0, \frac{1}{2}\right]$; therefore
\[
-\varepsilon_2 - 2\,\varepsilon_1 - 2\,\varepsilon_4 \le \Delta \le \varepsilon_2 + \varepsilon_1 + 2\,\varepsilon_3,
\]
and $|\Delta| \le 2\,\varepsilon_1 + \varepsilon_2 + 2\,\varepsilon_3 + 2\,\varepsilon_4$, which is the first inequality of~\eqref{eq:renormalisation-flow-23}. For the second, note that $l_{*} = M \cdot r_{*} \ge 1$, so that~\eqref{eq:renormalisation-group-T5} gives $\varepsilon_2 \le C_{\Sigma} \cdot \exp\left(-\gamma_1 \cdot l_{*}\right) \le e_{*}$; together with $\varepsilon_4 \le \varepsilon_3 \le \frac{e_{*}}{2}$ this yields
\[
2\,\varepsilon_1 + \varepsilon_2 + 2\,\varepsilon_3 + 2\,\varepsilon_4 \le e_{*}^{\frac{1}{4} - 9\varepsilon} + 3\,e_{*}.
\]
Since $0 < e_{*} < 1$, one has $e_{*}^{\frac{1}{4}} \le e_{*}^{\frac{1}{4} - 9\varepsilon} \le \frac{1}{8}$, and hence
\[
3\,e_{*} = 3 \left(e_{*}^{\frac{1}{4}}\right)^3 \cdot e_{*}^{\frac{1}{4}} \le \frac{3}{512}\, e_{*}^{\frac{1}{4}} \le e_{*}^{\frac{1}{4} - 9\varepsilon},
\]
which gives the second inequality of~\eqref{eq:renormalisation-flow-23}.

\emph{Proof of~\eqref{eq:renormalisation-flow-24}.} Corollary~\ref{cor:twist-ratios-free}, applied with the mass $\bar{m} \cdot \eta \in (0, 1]$ and with $n = R/\eta$, gives
\[
\left| \log \rho^{\mathrm{free}}_n - p^0(\bar{m} \cdot \eta) \cdot n^3 \right| \le \varepsilon_5(\bar{m} \cdot R).
\]
By the triangle inequality, this bound and~\eqref{eq:renormalisation-flow-23} give~\eqref{eq:renormalisation-flow-24}.

With Lemma~\ref{lem:product-formula-volume}, Corollary~\ref{cor:twist-ratios-free}, Lemmas~\ref{lem:fourier-decomposition} and~\ref{lem:dependence-twist-uniform} for the data of Section~\ref{sec:estimates-single-torus}, Proposition~\ref{prop:estimate-stop-level} and Lemma~\ref{lem:positivity-large-fields}, Theorem~\ref{thm:estimates-one-torus} follows.
\end{proof}

\subsection{Partition functions on the tori of $\mathrm{Tor}(e)$}\label{sec:partition-functions-tori}

\begin{theorem}[The partition function with the pair $P_N(e)$]\label{thm:partition-function-pair}
Let $e > 0$ satisfy $e_K = e \cdot h'^{\frac{1}{2}} \le e_1$. Let $R \in \mathrm{Tor}(e)$, let $N \ge N_{\mathrm{min}}(R, e)$, put $n := R/\eta$, and let $P_N(e) = \left(m_{0,N}, \boldsymbol{\varepsilon}_{0,N}\right)$ be the pair of Definition~\ref{def:tuned-pair}. With the function $\rho_n$ of~\eqref{eq:renormalisation-flow-1} and the function $p^0$ of Section~\ref{sec:constant-twists-free}, put
\[
\rho_n := \rho_n\left(e \cdot \eta^{\frac{1}{2}}, \bar{m} \cdot \eta + m_{0,N}\right), \qquad p_F := p^0(\bar{m} \cdot \eta) + \boldsymbol{\varepsilon}_{0,N}.
\]
Then
\begin{align}
& \left| \log \rho_n - \log \rho^{\mathrm{free}}_n - \boldsymbol{\varepsilon}_{0,N} \cdot n^3 \right| \le 3\, e_{*}^{\frac{1}{4} - 9\varepsilon}, \label{eq:renormalisation-flow-25}\\
& \left| \log \rho_n - p_F \cdot n^3 \right| \le \varepsilon_P(R) := \varepsilon_5(\bar{m} \cdot R) + 3\, e_{*}^{\frac{1}{4} - 9\varepsilon} \notag\\
& \qquad \le \varepsilon_5\left(L^{-m_0-1}\right) + 3\, e_K^{\frac{1}{4} - 9\varepsilon}, \label{eq:renormalisation-flow-26}
\end{align}
where $e_{*} = e_{k^{*}(R)}$; for a torus $R'$ of $\mathrm{Tor}(e)$ with $N \ge N_{\mathrm{min}}(R', e)$, the number $\varepsilon_P(R')$ is defined by the same formula, with $\bar{m} \cdot R'$ in place of $\bar{m} \cdot R$ and with $e_{k^{*}(R')}$ in place of $e_{*}$. Furthermore, for every integer $N \ge N_{\mathrm{min}}\left(R^{*}, e\right)$, with $n_{*} = R^{*}/\eta$ and with $\rho_{n_{*}}$ and $\rho_{L \cdot n_{*}}$ formed with the same two arguments $\left(e \cdot \eta^{\frac{1}{2}}, \bar{m} \cdot \eta + m_{0,N}\right)$ as $\rho_n$, one has
\begin{multline}\label{eq:renormalisation-flow-27}
D^{\mathrm{mat}} := L^3 \cdot \log \rho_{n_{*}} - \log \rho_{L \cdot n_{*}} \le L^3 \cdot \varepsilon_P\left(R^{*}\right) + \varepsilon_P\left(L \cdot R^{*}\right)\\
\le \left(L^3 + 1\right) \cdot \left( \varepsilon_5\left(\bar{m} \cdot R^{*}\right) + 3\, e_K^{\frac{1}{4} - 9\varepsilon} \right),
\end{multline}
and the same two inequalities hold for $-D^{\mathrm{mat}}$. Neither $P_N(e)$ nor $p_F$ varies with the torus $R$; both are determined by $N$, $e$, $\bar{m}$ and the scheme of Theorem~\ref{thm:agreement-trajectories}.
\end{theorem}

\begin{proof}
The charge $e$, the side $R$ and the cutoff $N$, with the primitives of Definition~\ref{def:periodisation-projection}, satisfy hypothesis~\ref{def:hypothesis-flow}. By Proposition~\ref{prop:consequences}, the hypotheses of Proposition~\ref{prop:depth-stop-level} and the condition~\eqref{eq:renormalisation-group-T2} hold; the last inequality of~\eqref{eq:renormalisation-group-T2} gives $0 < e_{*} < 1$. By Theorem~\ref{thm:agreement-trajectories}, the pair $P_N(e)$ is defined, real and controlled on $R$, and the value $\boldsymbol{\varepsilon}^t_{*} = \boldsymbol{\varepsilon}^t_{k^{*}(R)}$ of the first component of its trajectory on $R$ satisfies
\[
\left| \boldsymbol{\varepsilon}^t_{*} \right| \le 5\, c_{\mathrm{fl}} \cdot e_{*}^{\varepsilon} \cdot w^{\mathrm{E}}_{*} = 5\, c_{\mathrm{fl}} \cdot e_{*}^{\frac{1}{4} - 6\varepsilon},
\]
because $w^{\mathrm{E}}_{*} = e_{*}^{\frac{1}{4} - 7\varepsilon}$. In the second inequality of~\eqref{eq:renormalisation-group-T2} the summands $2 C_0$, $c_{\kappa}$ and $6$ are non-negative; hence that inequality implies $5\, c_{\mathrm{fl}} \cdot l_{*}^3 \le \frac{1}{4}\, e_{*}^{-\varepsilon}$, and
\[
\left| \boldsymbol{\varepsilon}^t_{*} \right| \cdot l_{*}^3 \le \frac{1}{4}\, e_{*}^{\frac{1}{4} - 7\varepsilon} \le e_{*}^{\frac{1}{4} - 9\varepsilon},
\]
the last step because $0 < e_{*} < 1$. Theorem~\ref{thm:twist-average-one} applies to the real pair $(m_{\mathrm{in}}, \boldsymbol{\varepsilon}_{\mathrm{in}}) = P_N(e)$; as $q = e \cdot \eta^{\frac{1}{2}}$, the number $\rho_n$ of that theorem is the number $\rho_n$ defined above. By the triangle inequality, \eqref{eq:renormalisation-flow-23} and the last display give~\eqref{eq:renormalisation-flow-25}, and~\eqref{eq:renormalisation-flow-24} and the last display give the first inequality of~\eqref{eq:renormalisation-flow-26}. By Proposition~\ref{prop:depth-stop-level}(e),
\[
\bar{m} \cdot R = \mathfrak{m}_{*} \cdot l_{*} \ge \mathfrak{m}_K > L^{-m_0-1}, \qquad e_{*} \le e_K.
\]
The function $\varepsilon_5$ is decreasing on $(0, \infty)$, since $\exp\left(-\frac{x}{4}\right)$ decreases in $x$, and the exponent $\frac{1}{4} - 9\varepsilon$ is positive; so $\varepsilon_5(\bar{m} \cdot R) \le \varepsilon_5\left(L^{-m_0-1}\right)$ and $e_{*}^{\frac{1}{4} - 9\varepsilon} \le e_K^{\frac{1}{4} - 9\varepsilon}$. This is the second inequality of~\eqref{eq:renormalisation-flow-26}.

Let now $N \ge N_{\mathrm{min}}\left(R^{*}, e\right)$. One has $M = L^m$ (Section~\ref{sec:lattice-geometry}) and $M \ge L^2$ (stage~\ref{it:order-choice-6}), so that $m \ge 2$; by its definition, $r_K = L^i$ with an integer $i \ge 0$. With Definition~\ref{def:tori-stop-level} it follows that
\begin{align*}
R^{*} & = \frac{R_{\mathrm{top}}}{L} = \frac{h' \cdot M \cdot r_K}{L} = h' \cdot L^{m-1} \cdot r_K = h' \cdot L^{m-1+i},\\
L \cdot R^{*} & = R_{\mathrm{top}} = h' \cdot L^{m+i}.
\end{align*}
Both numbers have the form $h' \cdot L^T$ with an integer $T \ge 0$, and neither exceeds $R_{\mathrm{top}}$; hence $R^{*}$ and $L \cdot R^{*}$ belong to $\mathrm{Tor}(e)$. As $R^{*} \le R_{\mathrm{top}}$, Proposition~\ref{prop:depth-stop-level}(d) gives
\[
N \ge N_{\mathrm{min}}\left(R^{*}, e\right) \ge N_{\mathrm{min}}\left(R_{\mathrm{top}}, e\right).
\]
So~\eqref{eq:renormalisation-flow-26} holds on the torus $R^{*}$, where $n = n_{*}$, and on the torus $L \cdot R^{*}$, where $n = L \cdot n_{*}$. On both tori it is formed with the one pair $P_N(e)$ of Definition~\ref{def:tuned-pair}, which does not depend on the torus (Theorem~\ref{thm:agreement-trajectories}); thus $p_F$ is the same number on both. The upper bound of~\eqref{eq:renormalisation-flow-26} on $R^{*}$, multiplied by $L^3$, and the lower bound on $L \cdot R^{*}$ read
\begin{align*}
L^3 \cdot \log \rho_{n_{*}} & \le p_F \cdot L^3 \cdot n_{*}^3 + L^3 \cdot \varepsilon_P\left(R^{*}\right),\\
\log \rho_{L \cdot n_{*}} & \ge p_F \cdot \left(L \cdot n_{*}\right)^3 - \varepsilon_P\left(L \cdot R^{*}\right).
\end{align*}
Since $L^3 \cdot n_{*}^3 = \left(L \cdot n_{*}\right)^3$, the terms that contain $p_F$ cancel in the difference, and $D^{\mathrm{mat}} \le L^3 \cdot \varepsilon_P\left(R^{*}\right) + \varepsilon_P\left(L \cdot R^{*}\right)$. In the same way the lower bound on $R^{*}$ and the upper bound on $L \cdot R^{*}$,
\begin{align*}
L^3 \cdot \log \rho_{n_{*}} & \ge p_F \cdot L^3 \cdot n_{*}^3 - L^3 \cdot \varepsilon_P\left(R^{*}\right),\\
\log \rho_{L \cdot n_{*}} & \le p_F \cdot \left(L \cdot n_{*}\right)^3 + \varepsilon_P\left(L \cdot R^{*}\right),
\end{align*}
give $-D^{\mathrm{mat}} \le L^3 \cdot \varepsilon_P\left(R^{*}\right) + \varepsilon_P\left(L \cdot R^{*}\right)$. By Proposition~\ref{prop:depth-stop-level}(e), applied on each of the two tori, the number $e_{*}$ of either torus is at most $e_K$; and $\varepsilon_5\left(\bar{m} \cdot L \cdot R^{*}\right) \le \varepsilon_5\left(\bar{m} \cdot R^{*}\right)$, as $\varepsilon_5$ is decreasing. Hence each of $\varepsilon_P\left(R^{*}\right)$ and $\varepsilon_P\left(L \cdot R^{*}\right)$ is at most $\varepsilon_5\left(\bar{m} \cdot R^{*}\right) + 3\, e_K^{\frac{1}{4} - 9\varepsilon}$, which proves the second inequality of~\eqref{eq:renormalisation-flow-27}, for $D^{\mathrm{mat}}$ and for $-D^{\mathrm{mat}}$.

It remains to prove the last assertion of the theorem. By Theorem~\ref{thm:agreement-trajectories}, $P_N(e)$ is one and the same pair for all tori of $\mathrm{Tor}(e)$, and it is a function of $N$, $e$, $\bar{m}$ and the scheme. The number $p^0(\bar{m} \cdot \eta)$, where $\eta = L^{-N}$, is a function of $N$, $\bar{m}$ and the integer $L$ of stage~\ref{it:order-choice-4}, one of the constants of the scheme. Hence $p_F = p^0(\bar{m} \cdot \eta) + \boldsymbol{\varepsilon}_{0,N}$ has the same two properties, and Theorem~\ref{thm:partition-functions-tori} follows.
\end{proof}

\section{Four-jets in one probe per cell on the reference torus}\label{sec:four-jets-one}

This section contains the proofs of the theorems of Section~\ref{sec:probes-jets-reference}, from which Estimate~8 is deduced in Section~\ref{sec:proof-theorem}; it also contains, in the proof of Lemma~\ref{lem:scaling-external-connection}, the scaling computation on which the proof of Lemma~\ref{lem:external-connections} rests. Section~\ref{sec:standing-data-constants} states the definitions of Section~\ref{sec:probes-jets-reference} once more (probes, placements, mirror copies, the ring of jets, the constants), in greater detail and with the further constants that the proofs require; a shorter account, adapted to the proof of Theorem~\ref{thm:main-theorem}, is given in Section~\ref{sec:probes-discretisation}.

\subsection{Standing data, constants and smallness conditions}\label{sec:standing-data-constants}

\subsubsection{Conventions}\label{sec:conventions-2}

The notation fixed in Definition~\ref{def:notation}(a)--(h) is in force. In particular, $e$ denotes the charge and $e_j$, $e_{*}$ denote running charges; the exponential function is written $\exp(\cdot)$, for real and for imaginary arguments alike, and Euler's number is $\exp(1)$ (Definition~\ref{def:notation}(a), (e)). For $x \in \mathbb{R}^3$, $|x|$ is the sup norm and $d(\cdot,\cdot)$ is the sup distance on a torus (Definition~\ref{def:notation}(a)); $|x|_1$ denotes the $\ell^1$ norm. The directions are labelled $\nu = 0, 1, 2$ (Definition~\ref{def:notation}(d)).

Some letters are reserved. The jet radius of Section~\ref{sec:probes-mirror-copies} is denoted by $\mathbf{r}$. With the numbers $g$ and $c_1$ of Section~\ref{sec:constants-3}, one writes $h_G = 1 + g$ and $\tau_s = c_1\cdot l_{*}$. The set of the cells of Section~\ref{sec:cells-reflections} is $\mathrm{Cel}$. With $R^{*}$ and $\eta$ as in Section~\ref{sec:standing-hypothesis-scales}, $n_{*} := R^{*}/\eta$; the number of the unit sites at the stop level is not denoted by $n_{*}$, but by $l_{*}^3$.

On $\mathbb{C}^4$ the norm
\[
|v| := \max_{\boldsymbol{\beta}} |v_{\boldsymbol{\beta}}|
\]
is used, the maximum being taken over the four components $v_{\boldsymbol{\beta}}$ of $v$; for a function $f$ with values in $\mathbb{C}^4$, the norm $|f|_{\infty}$ is formed with it. The bounds of Section~\ref{sec:operators-fermion-field} that are applied in Section~\ref{sec:four-jets-one} remain valid when $|\cdot|_{\infty}$ is understood in this way. Indeed, the only ones among them in which the norm on $\mathbb{C}^4$ enters are the bounds of Propositions~\ref{prop:decay-fine-propagator}, \ref{prop:kernel-bounds-2} and~\ref{prop:bounds}, with the constants $(C_F,\gamma_F) = (C^F,\gamma^F)$ of Theorem~\ref{thm:fermion-operators-weakened}. As the Euclidean norm of every $v \in \mathbb{C}^4$ lies between $|v|$ and $2|v|$, a bound that holds with a constant $C$ when the Euclidean norm is used holds with the constant $2C$ when the norm $|\cdot|$ is used, and the constant $C^F$ contains this factor $2$ (Section~\ref{sec:norms-pairings-kernels}).

\subsubsection{The standing hypothesis and the scales at the stop level}\label{sec:standing-hypothesis-scales}

Throughout Section~\ref{sec:four-jets-one}, hypothesis~\ref{def:hypothesis-reference-torus} of Section~\ref{sec:probes-jets-reference} is in force. In its formulation and below, the numbers $h$, $h'$, the level $K$ with $a_K = h'$, the scales $a_k$, $e_k$, $v_k$, $r_k$, $p_k$, $p_{0,k}$, $l_k$, $\mathfrak{m}_k$, $b_k$, the bound $b_{-}$ and the number $\mathrm{lg} = \frac{1}{2}\log L$ are those of Section~\ref{sec:running-quantities-tori}, the sides $l_k$ being those of the torus $R^{*}$ introduced below; $R_{\mathrm{top}}$, $\mathrm{Tor}(e)$, $N_{\mathrm{min}}$ and the stop level $k^{*}(\cdot)$ are those of Definition~\ref{def:tori-stop-level}; and $\mathrm{Cl}$ is the class of sides of Section~\ref{sec:lattice-geometry}. Written out, hypothesis~\ref{def:hypothesis-reference-torus} consists of the following requirements.
\begin{itemize}
\item The constants of stages~\ref{it:order-choice-1}--\ref{it:order-choice-10} are chosen as in Section~\ref{sec:order-choice}, and they satisfy the conditions \eqref{eq:renormalisation-group-A0}--\eqref{eq:renormalisation-group-A10}, \eqref{eq:renormalisation-group-19}, \eqref{eq:renormalisation-group-20}--\eqref{eq:renormalisation-group-24}, \eqref{eq:renormalisation-group-25} and \eqref{eq:renormalisation-group-26} of Section~\ref{sec:conditions}.
\item $0 < \bar{m} \le 1$, and the charge $e > 0$ satisfies $e_K = e\cdot h'^{\frac{1}{2}} \le e_1$.
\item The side of the torus is
\[
R = R^{*} := \frac{h'\cdot M\cdot r_K(e)}{L}.
\]
\item The cutoff satisfies $N \ge N_{\mathrm{min}}(R^{*},e)$, and $\eta = L^{-N}$.
\item At the levels $k < k^{*}(R^{*})$ the systems of primitives are those of Definition~\ref{def:periodisation-projection}.
\item The pair is $\left(m_{\mathrm{in}}, \boldsymbol{\varepsilon}_{\mathrm{in}}\right) = P_N(e) = \left(m_{0,N}, \boldsymbol{\varepsilon}_{0,N}\right)$, the pair of Definition~\ref{def:tuned-pair}.
\end{itemize}
Thus hypothesis~\ref{def:hypothesis-reference-torus} is hypothesis~\ref{def:hypothesis-flow} of Section~\ref{sec:one-step-operators} for $R = R^{*}$, together with the choice $P_N(e)$ of the pair. (The side $R^{*} = h'\cdot L^{m-1}\cdot r_K(e)$ belongs to $\mathrm{Tor}(e)$: one has $M = L^m$ with $m \ge 2$ by stage~\ref{it:order-choice-6}, the number $r_K(e)$ is a power of $L$, and $R^{*} = R_{\mathrm{top}}/L < R_{\mathrm{top}}$.)

Let $k^{*} := k^{*}(R^{*})$. A subscript $*$ denotes the value at the level $k = k^{*}$. This applies to $a_{*}$, $e_{*}$, $v_{*}$, $r_{*}$, $p_{*}$, $p_{0*} := p_{0,k^{*}}$, $l_{*}$, $\mathfrak{m}_{*}$, $b_{*}$, $\eta_{*} = L^{-k^{*}}$, $T_{*} := T_{k^{*}}$ and $T^0_{*} := T^0_{k^{*}}$, and to $D_{*}$, which, taken as a set, is the set of the polymers of $T_{*}$. Finally, put
\[
\mathfrak{m}_{-} := L^{-m_0-3}.
\]

\begin{lemma}[Scales at the stop level]\label{lem:scales-stop-level}
\begin{enumerate}
\item[(a)] One has $k^{*} \in \{K-1, K-2\}$, $k^{*} \ge 1$, $l_{*} = M\cdot r_{*}$, $r_{*} \in \{r_K, L\cdot r_K\}$ and $a_{*} \in \left\{\frac{h'}{L}, \frac{h'}{L^2}\right\}$. Moreover
\[
\mathfrak{m}_{-} < \mathfrak{m}_{*} \le L^{-m_0-1} \le 1, \qquad e_{*} = e\cdot a_{*}^{\frac{1}{2}} \le e_K \le e_0, \qquad b_{-} = 1 - \frac{1}{L} \le b_{*} \le 1.
\]
\item[(b)] One has $v_K \le v_{*} \le v_K + \log L$ and $v_{*}^r \le r_{*} < L\cdot v_{*}^r$. The number $r_{*}$ is a power of $L$, and $r_{*} \ge L^2$.
\item[(c)] For $0 \le j \le k^{*}$ the side $l_j$ lies in $\mathrm{Cl}$. The data $\left(k^{*}, l_{*}, \mathfrak{m}_{*}, e_{*}\right)$ belong to the scale-free setting of Section~\ref{sec:properties-scheme} (thus $l_{*} \in \mathrm{Cl}$, $\mathfrak{m}_{*} \in (0, \mathfrak{m}_0]$ and $e_{*} \in (0, e_0]$). For $0 \le j < k^{*}$, hypothesis~\ref{def:hypothesis-one-step} holds for the data $\left(j, l_j, \bar{m}\cdot a_j, e_j\right)$ together with the primitives of Definition~\ref{def:periodisation-projection}. In these data, $\bar{m}\cdot a_j$ is the running mass of level $j$, that is, $\mathfrak{m}_j$ if $j \ge 1$ and $\bar{m}\cdot\eta$ if $j = 0$, and $e_j$ is the running charge of level $j$, that is, $q$ if $j = 0$ (Section~\ref{sec:running-quantities-tori} and Definition~\ref{def:notation}(b), (e)).
\end{enumerate}
\end{lemma}

\begin{proof}
Since $e_K \le e_1$, Proposition~\ref{prop:consequences} applies. It gives $v_K \ge 8$, in particular $v_K \ge 1$, and
\[
\left(1 + \frac{\mathrm{lg}}{v_K}\right)^r \le \left(1 + \frac{3\mathrm{lg}}{v_K}\right)^r < L;
\]
these are the hypotheses of Proposition~\ref{prop:depth-stop-level}, including the additional hypothesis of its part~(g). It also gives condition~\eqref{eq:renormalisation-group-16} on the torus $R^{*} \in \mathrm{Tor}(e)$, for every $N \ge N_{\mathrm{min}}(R^{*},e)$.

\emph{(a)} Proposition~\ref{prop:depth-stop-level}(g) yields $k^{*} \in \{K-1, K-2\}$, $a_{*} \in \left\{\frac{h'}{L}, \frac{h'}{L^2}\right\}$ and $r_{*} \in \{r_K, L\cdot r_K\}$. Part~(c) of the same proposition yields $k^{*} \ge 1$, $l_{*} = M\cdot r_{*}$ and $e_{*} = e\cdot a_{*}^{\frac{1}{2}}$, and part~(e) yields $e_{*} \le e_K$. One has $e_K \le e_1$ by hypothesis~\ref{def:hypothesis-reference-torus}, and $e_1 \le e_0$ by stage~\ref{it:order-choice-10}. For the torus $R^{*}$, Proposition~\ref{prop:bounds-running-masses}(e) gives $L^{-m_0-3} < \mathfrak{m}_{*} \le L^{-m_0-1}$, that is, $\mathfrak{m}_{-} < \mathfrak{m}_{*} \le L^{-m_0-1}$; and $L^{-m_0-1} \le 1$ because $m_0 \ge 0$. Finally, $b = 1$ by stage~\ref{it:order-choice-1}, and by Section~\ref{sec:running-quantities-tori} one has $b_{-} = b\cdot\left(1 - L^{-1}\right) \le b_k \le b$ for every level $k \ge 1$. For $k = k^{*} \ge 1$ this reads $1 - \frac{1}{L} \le b_{*} \le 1$.

\emph{(b)} By (a), $\frac{h'}{L^2} \le a_{*} \le \frac{h'}{L}$. Since $e_{*} = e\cdot a_{*}^{\frac{1}{2}}$ and $e_K = e\cdot h'^{\frac{1}{2}}$,
\[
v_{*} = -\log e_{*} = v_K + \frac{1}{2}\log\left(\frac{h'}{a_{*}}\right),
\]
and the logarithm on the right lies between $\log L$ and $2\log L$; hence $v_K \le v_{*} \le v_K + \log L$. In particular $v_{*} \ge v_K \ge 1$, so that $v_{*}^r \ge 1$. By Section~\ref{sec:running-quantities-tori}, $r_{*}$ is the least power $L^i$, with an integer $i \ge 0$, such that $L^i \ge v_{*}^r$. So $r_{*} \ge v_{*}^r$; if $r_{*} \ge L$, then $r_{*}/L < v_{*}^r$ by minimality, and if $r_{*} = 1$, then $r_{*} < L \le L\cdot v_{*}^r$. In either case $r_{*} < L\cdot v_{*}^r$. Finally, $r_{*}$ is a power of $L$, and condition~\eqref{eq:renormalisation-group-16} gives $r_{*} > 31L + 16 > L$; therefore $r_{*} \ge L^2$.

\emph{(c)} As condition~\eqref{eq:renormalisation-group-16} holds, Proposition~\ref{prop:depth-stop-level}(f) gives $l_j \in \mathrm{Cl}$ for $0 \le j \le k^{*}$. Hypothesis~\ref{def:hypothesis-reference-torus} contains hypothesis~\ref{def:hypothesis-flow} for $R = R^{*}$. Hence Proposition~\ref{prop:data-along-flow}(c) gives the assertion on $\left(k^{*}, l_{*}, \mathfrak{m}_{*}, e_{*}\right)$, and Proposition~\ref{prop:data-along-flow}(b) the assertion on the levels $j < k^{*}$.
\end{proof}

\subsubsection{Cells and their reflections}\label{sec:cells-reflections}

Unless another level is named, points of the torus are written in lattice units, that is, by their coordinates of level $0$ (Section~\ref{sec:lattice-measures-laws}). As $a_{*} = L^{k^{*}}\cdot\eta$ and $l_{*} = R^{*}/a_{*}$, the number of sites on each side of the torus is
\[
n_{*} = \frac{R^{*}}{\eta} = L^{k^{*}}\cdot l_{*}.
\]
Put $\mathrm{per} := n_{*}/L$. By Lemma~\ref{lem:scales-stop-level}(a), (b), $\mathrm{per} = L^{k^{*}-1}\cdot M\cdot r_{*}$ is a power of the odd integer $L$; let $m_{\mathrm{p}}$ be the integer with $\mathrm{per} = 2m_{\mathrm{p}} + 1$. For $\nu = 0, 1, 2$, the site planes are the planes $\{x_{\nu} = p\}$ with $p \in \mathrm{per}\cdot\mathbb{Z}$, and the link planes are the planes $\{x_{\nu} = p\}$ with $p \in \mathrm{per}\cdot\mathbb{Z} + m_{\mathrm{p}} + \frac{1}{2}$. Since $m_{\mathrm{p}} + \frac{1}{2} = \frac{\mathrm{per}}{2}$, the site planes and the link planes together are the planes $\{x_{\nu} = p\}$ with $2p \in \mathrm{per}\cdot\mathbb{Z}$. These planes cut the torus into $N^{*} = 8L^3$ cells $\mathfrak{c}_c$, indexed as in Section~\ref{sec:cells-reference-tori}; $\mathrm{Cel}$ is the set of the cells, a cell being identified with its index $c$. The cell $\mathfrak{c}_0$ is the one given by $0 \le x_{\nu} \le m_{\mathrm{p}} + \frac{1}{2}$ for $\nu = 0, 1, 2$. The physical side of a cell is
\[
l_c := \frac{R^{*}}{2L} = \left(m_{\mathrm{p}} + \frac{1}{2}\right)\cdot\eta .
\]
The reflections of the torus in the site planes and in the link planes generate a group; for $c \in \mathrm{Cel}$, $\mathsf{g}_c$ is the member of this group that maps $\mathfrak{c}_0$ onto $\mathfrak{c}_c$. For a lattice isometry $\tilde{g}$ and a bond function $A$, the push-forward $\tilde{g}_{*}A$ is the bond function
\[
\left(\tilde{g}_{*}A\right)(x,x') := A\left(\tilde{g}^{-1}x, \tilde{g}^{-1}x'\right).
\]

\begin{lemma}[The reflections as lattice symmetries]\label{lem:reflections-lattice-symmetries}
Let $c \in \mathrm{Cel}$.
\begin{enumerate}
\item[(a)] In the coordinates of level $0$, $\mathsf{g}_cx = r_cx + a_c$, with a matrix $r_c = \operatorname{diag}(\pm1, \pm1, \pm1)$ and a vector $a_c \in \mathrm{per}\cdot\mathbb{Z}^3$.
\item[(b)] Let $0 \le j \le k^{*}$. Then $L^{-j}\cdot a_c \in L\cdot M\cdot\mathbb{Z}^3$, and, written in the coordinates of level $j$, the map $\mathsf{g}_c$ takes the form
\[
x \mapsto r_cx + L^{-j}\cdot a_c.
\]
Hence, whenever $S \in \mathrm{Sp}$ satisfies $r(S) = r_c$, the step group $\mathbb{G}_j^L$ contains $\left(S, L^{-j}\cdot a_c\right)$, and this element acts on the points of the torus, in the coordinates of level $j$, as $\mathsf{g}_c$.
\end{enumerate}
\end{lemma}

\begin{proof}
\emph{(a)} Each plane is of the form $\{x_{\nu} = p\}$ with $\nu \in \{0, 1, 2\}$ and $2p \in \mathrm{per}\cdot\mathbb{Z}$, and the reflection in it replaces the coordinate $x_{\nu}$ by $2p - x_{\nu}$ and leaves the two other coordinates unchanged. This reflection is a map $x \mapsto rx + a$ with a diagonal sign matrix $r$ and the vector $a = 2p\cdot e_{\nu} \in \mathrm{per}\cdot\mathbb{Z}^3$. The maps of this form are closed under composition, because
\[
r\left(r'x + a'\right) + a = rr'x + \left(ra' + a\right)
\]
and $ra' \in \mathrm{per}\cdot\mathbb{Z}^3$. Hence every element of the group generated by the reflections, and in particular $\mathsf{g}_c$, is of the form stated in~(a).

\emph{(b)} The fine lattice $T_j$ of level $j$ has the side $l_j$ (Section~\ref{sec:lattice-geometry}) and the spacing $\eta_j = L^{-j}$ (Section~\ref{sec:running-quantities-tori}). By the definition of the coordinates of level $j$ in Section~\ref{sec:lattice-measures-laws}, a bond $b$ of $T_j$ corresponds to the bond $L^j\cdot b$ of the lattice $T_0 = \mathbb{Z}^3/n_{*}\mathbb{Z}^3$ of level $0$ (one has $L^j\cdot\eta_j = 1$ and $L^j\cdot l_j = l_0 = n_{*}$). Consequently a point with the coordinates $x$ of level $0$ has the coordinates $L^{-j}\cdot x$ of level $j$, and in the coordinates of level $j$ the map $\mathsf{g}_c$ reads
\[
x \mapsto L^{-j}\cdot\left(r_c\left(L^j\cdot x\right) + a_c\right) = r_cx + L^{-j}\cdot a_c.
\]
By the beginning of this subsection and by Lemma~\ref{lem:scales-stop-level}(a), $\mathrm{per} = L^{k^{*}-1}\cdot l_{*} = L^{k^{*}-1}\cdot M\cdot r_{*}$, so that
\[
L^{-j}\cdot\mathrm{per} = L^{k^{*}-1-j}\cdot M\cdot r_{*}.
\]
By Lemma~\ref{lem:scales-stop-level}(b), $r_{*}$ is a power of $L$ with $r_{*} \ge L^2$; hence $\frac{r_{*}}{L^2}$ is an integer, and $M\cdot r_{*}$ is an integral multiple of $L\cdot M$. If $j < k^{*}$, the exponent $k^{*}-1-j$ is non-negative, so that $L^{-j}\cdot\mathrm{per}$ is an integral multiple of $M\cdot r_{*}$, hence of $L\cdot M$. If $j = k^{*}$, then
\[
L^{-j}\cdot\mathrm{per} = \frac{M\cdot r_{*}}{L} = L\cdot M\cdot\frac{r_{*}}{L^2}.
\]
In either case $L^{-j}\cdot\mathrm{per} \in L\cdot M\cdot\mathbb{Z}$. As $a_c \in \mathrm{per}\cdot\mathbb{Z}^3$, it follows that $L^{-j}\cdot a_c \in L\cdot M\cdot\mathbb{Z}^3$.

The map $r$ of Section~\ref{sec:dirac-matrices-lattice} sends $\mathrm{Sp}$ onto the set of the eight diagonal sign matrices (Section~\ref{sec:dirac-matrices-lattice} and \eqref{eq:operators-fermion-1}); hence there are elements $S \in \mathrm{Sp}$ with $r(S) = r_c$. For every such $S$, the definition of $\mathbb{G}_j^L$ in Section~\ref{sec:dirac-matrices-lattice} and the inclusion $L^{-j}\cdot a_c \in L\cdot M\cdot\mathbb{Z}^3$ show that $\left(S, L^{-j}\cdot a_c\right)$ belongs to $\mathbb{G}_j^L$. It acts on the coordinates of level $j$ by
\[
x \mapsto r(S)x + L^{-j}\cdot a_c = r_cx + L^{-j}\cdot a_c,
\]
that is, as $\mathsf{g}_c$.
\end{proof}

\subsubsection{Probes and their mirror copies}\label{sec:probes-mirror-copies}

\paragraph{The two-form and its placement.}
Let $\beta\in C_c^{\infty}(\mathbb{R}^3;\Lambda^2)$ be a real two-form with $\beta\ne0$; its components satisfy $\beta_{\mu\nu}=-\beta_{\nu\mu}$. For a smooth function $\mathfrak{g}$ on $\mathbb{R}^3$ and an integer $i\ge0$ put
\[
\left|D^i\mathfrak{g}\right|_{\infty}:=\max_{\lambda_1,\dots,\lambda_i\in\{0,1,2\}}\ \sup_{x\in\mathbb{R}^3}\left|\left(\partial_{\lambda_1}\cdots\partial_{\lambda_i}\mathfrak{g}\right)(x)\right|\in[0,\infty],
\]
so that $\left|D^0\mathfrak{g}\right|_{\infty}=|\mathfrak{g}|_{\infty}$ is the sup norm of $\mathfrak{g}$; one writes $|D\mathfrak{g}|_{\infty}$ for $\left|D^1\mathfrak{g}\right|_{\infty}$. For an integer $i\ge0$ let
\[
|\beta|_{C^i}:=\max\left\{\left|D^j\beta_{\mu\nu}\right|_{\infty}:\ 0\le j\le i,\ \mu,\nu\in\{0,1,2\}\right\}.
\]
Let $\rho_{\beta}$ be the least number $\rho$ such that
\[
\operatorname{supp}\beta+[-2,2]^3\subset y+[-\rho,\rho]^3\qquad\text{for some }y\in\mathbb{R}^3
\]
(it exists, as $\operatorname{supp}\beta$ is compact and not empty), let $y_{\beta}\in\mathbb{R}^3$ be a point with $\operatorname{supp}\beta+[-2,2]^3\subset y_{\beta}+[-\rho_{\beta},\rho_{\beta}]^3$, and put $D_{\beta}:=2\rho_{\beta}$.

Multiplying the lattice coordinates of Section~\ref{sec:cells-reflections} by $\eta$ gives the \emph{physical coordinates}. In them $\mathfrak{c}_0=[0,l_c]^3$, the sites of $\Lambda_{n_{*}}$ are the classes modulo $R^{*}\mathbb{Z}^3$ of the points of $\eta\mathbb{Z}^3$, and a site is said to lie in a set $S\subset\mathbb{R}^3$ when $S$ contains one of its representatives in $\eta\mathbb{Z}^3$. The \emph{admitted placements} are the vectors $v\in\eta\mathbb{Z}^3$ for which
\begin{equation}\label{eq:four-jets-1}
\bar{Q}:=v+y_{\beta}+[-\rho_{\beta},\rho_{\beta}]^3\ \subset\ \left[\frac{l_c}{8},\frac{7l_c}{8}\right]^3.
\end{equation}
The centred placement of Section~\ref{sec:probes-discretisation}, that is, a vector $v\in\eta\mathbb{Z}^3$ with $\left|v+y_{\beta}-\frac{l_c}{2}(1,1,1)\right|\le\frac{\eta}{2}$, is admitted if $\rho_{\beta}+\frac{\eta}{2}\le\frac{3}{8}l_c$, because each coordinate of a point of $\bar{Q}$ then differs from $\frac{l_c}{2}$ by at most $\rho_{\beta}+\frac{\eta}{2}$.

\paragraph{The lattice probe.}
For a smooth function $\mathfrak{g}$ on $\mathbb{R}^3$, a point $x\in\mathbb{R}^3$ and $\mu\in\{0,1,2\}$ let
\[
\delta_{\eta,\mu}\mathfrak{g}(x):=\frac{\mathfrak{g}\left(x+\frac{\eta}{2}e_{\mu}\right)-\mathfrak{g}\left(x-\frac{\eta}{2}e_{\mu}\right)}{\eta},\qquad J_{\eta}\mathfrak{g}(x):=\eta^{-3}\int_{\left[-\frac{\eta}{2},\frac{\eta}{2}\right]^3}\mathrm{d}y\,\mathfrak{g}(x+y).
\]
For $P>0$, the \emph{periodisation} with the period $P$ of a function, a two-form or a bond function of compact support, given on $\mathbb{R}^3$ or on a lattice in $\mathbb{R}^3$ that is invariant under the translations by $P\mathbb{Z}^3$, is the sum of its translates by the vectors of $P\mathbb{Z}^3$; near each point only finitely many of these translates differ from zero, and the sum is invariant under the translations by $P\mathbb{Z}^3$. Now fix a vector $v\in\eta\mathbb{Z}^3$, and let $\tau_v\beta$ be the periodisation, with the period $R^{*}$, of the translate $\beta(\cdot-v)$:
\[
(\tau_v\beta)_{\mu\nu}(x)=\sum_{m\in\mathbb{Z}^3}\beta_{\mu\nu}\left(x-v-R^{*}\cdot m\right)\qquad\left(x\in\mathbb{R}^3,\ \mu,\nu\in\{0,1,2\}\right).
\]
For a site $x$ of $\Lambda_{n_{*}}$, in physical coordinates, and for $\mu,\nu\in\{0,1,2\}$ let
\[
u_{\mu\nu}(x):=\left(J_{\eta}(\tau_v\beta)_{\mu\nu}\right)\left(x+\frac{\eta}{2}\left(e_{\mu}+e_{\nu}\right)\right).
\]
Since $u_{\nu\mu}=-u_{\mu\nu}$, this defines a two-cochain $u$ on $\Lambda_{n_{*}}$: its value at the plaquette with the corner $x$ spanned by $\eta e_{\mu}$ and $\eta e_{\nu}$ is $u_{\mu\nu}(x)$, the value of $J_{\eta}(\tau_v\beta)_{\mu\nu}$ at the centre of the plaquette. Let $d$ be the coboundary, $d_{\eta}:=\eta^{-1}d$, and let $d_{\eta}^{*}$ be the adjoint of $d_{\eta}$ for the pairings $\eta^3\sum$, the sums extending over the sites and the components. Put
\[
a_{\eta}:=d_{\eta}^{*}u,\qquad\hat{\mathfrak{a}}:=\eta^{\frac{1}{2}}\cdot a_{\eta}\in C^1(\Lambda_{n_{*}}),\qquad\boldsymbol{\sigma}^2:=|d\hat{\mathfrak{a}}|_2^2,\qquad f:=dd^{*}\beta.
\]
Here $a_{\eta,\nu}(x)$ is the value of $a_{\eta}$ at the bond from $x$ to $x+\eta e_{\nu}$; $\hat{\mathfrak{a}}$ is read, in lattice units, as a real one-cochain on $\Lambda_{n_{*}}$; $|\cdot|_2$ is the counting norm of Section~\ref{sec:lattice-measures-laws}; and in the definition of $f$ the operators $d$ and $d^{*}$ are the exterior derivative and its formal adjoint on forms on $\mathbb{R}^3$.

\paragraph{Mirror copies.}
For $c\in\mathrm{Cel}$ let
\[
\hat{\mathfrak{a}}_c:=(\mathsf{g}_c)_{*}\hat{\mathfrak{a}},\qquad X_c:=\left\langle d\hat{\mathfrak{a}}_c,F\right\rangle\qquad\left(F\in F_N\right),
\]
where $(\mathsf{g}_c)_{*}$ is the push-forward of Section~\ref{sec:cells-reflections}, $F_N=dC^1(\Lambda_{n_{*}})$ is the space of Section~\ref{sec:lattice-measures-laws}, and $\langle\cdot,\cdot\rangle$ is the counting pairing. Put $X:=X_{c_0}$, where $c_0\in\mathrm{Cel}$ is the cell with $\mathfrak{c}_{c_0}=\mathfrak{c}_0$. For a complex bond function $a$ on $\Lambda_{n_{*}}$ let $|a|_{\mathrm{ext}}$ be the sum of the three norms of Section~\ref{sec:domains-bond-functions} of the bond function $a^{[k^{*}]}$ on $T_{*}$, where $a^{[k]}$ denotes $a$ in the coordinates of level $k$ (Section~\ref{sec:lattice-measures-laws}); this is the number $|a|_{\mathrm{ext}}$ of Lemma~\ref{lem:external-connections}.

\paragraph{Numbers and sets attached to the probe.}
Let
\[
C_{\beta}:=\sup|\hat{\mathfrak{a}}|_{\mathrm{ext}},\qquad\mathbf{D}:=\sup\boldsymbol{\sigma}^2,
\]
where in both suprema $N$, the number $a_{*}\in\left\{h'/L,h'/L^2\right\}$ and the admitted placement $v$ vary, and $\beta$ may be replaced by any of the forms $r_{*}\beta$. Here $r=\operatorname{diag}(r_0,r_1,r_2)$ ranges over the diagonal sign matrices, and $r_{*}\beta$ is the push-forward of $\beta$ under $r$, with the components $(r_{*}\beta)_{\mu\nu}(x)=r_{\mu}r_{\nu}\cdot\beta_{\mu\nu}(rx)$; the symbol $r_{*}$ standing alone is the number $r_{k^{*}}$. Further let
\begin{align*}
t_x&:=e_{*}^{3\varepsilon}\cdot\left(1+h'^{\frac{1}{2}}\cdot C_{\beta}\right)^{-1},\qquad t_x^{+}:=e_{*}^{3\varepsilon}\cdot\left(h'^{\frac{1}{2}}\cdot C_{\beta}\right)^{-1},\qquad\mathbf{r}:=\frac{t_x}{4},\\
\hat{c}_c&:=\frac{1}{e}\cdot\hat{\mathfrak{a}}_c^{[k^{*}]}\qquad\left(c\in\mathrm{Cel}\right),
\end{align*}
with $t_x^{+}:=\infty$ if $C_{\beta}=0$. Thus $t_x^{+}>t_x$, and $\hat{c}_c$ is a real bond function on $T_{*}$. For $c\in\mathrm{Cel}$ let
\begin{align*}
\mathrm{Bd}_c&:=\left\{b\ \text{a bond of}\ T_{*}:\ \hat{c}_c(b)\ne0\right\},\\
P_c&:=\left\{y\in T^0_{*}:\ \Delta_y\ \text{contains an end point of a bond of}\ \mathrm{Bd}_c\right\},\\
P_c^{+}&:=\left\{y\in T^0_{*}:\ d(y,P_c)\le M+1\right\},
\end{align*}
where $\Delta_y$ is the set of the sites $x$ of $T_{*}$ with $|x-y|<\frac{1}{2}$ (Section~\ref{sec:lattice-geometry}) and $d(\cdot,\cdot)$ is the sup distance on the torus; and put
\[
\mathsf{p}:=\left(\frac{D_{\beta}}{a_{*}}+2M+4\right)^3,\qquad d_{\mathrm{sep}}:=\min_{c\ne c'}d(P_c,P_{c'}),\qquad\epsilon_c:=e_{*}\cdot|\hat{c}_c|_{\infty},
\]
with $d_{\mathrm{sep}}:=\infty$ if $\hat{\mathfrak{a}}=0$.

\paragraph{The constants $C_{\beta}$ and $\mathbf{D}$.}
For $\nu\in\{0,1,2\}$ the function
\[
\mathring{g}^0_{\nu}:=-\sum_{\mu\ne\nu}\delta_{\eta,\mu}J_{\eta}\beta_{\mu\nu}
\]
is smooth, has compact support and is determined by $\beta$ and $\eta$. Let $w$ and $\acute{g}$ be the bond functions on $\eta\mathbb{Z}^3$ and on $\eta_{*}\mathbb{Z}^3$ whose values at the bonds from $x$ to $x+\eta e_{\nu}$ and from $y$ to $y+\eta_{*}e_{\nu}$ are
\[
w_{\nu}(x):=\eta^{\frac{1}{2}}\cdot\mathring{g}^0_{\nu}\left(x+\frac{\eta}{2}e_{\nu}\right)\quad\left(x\in\eta\mathbb{Z}^3\right),\qquad\acute{g}_{\nu}(y):=a_{*}^{\frac{1}{2}}\cdot\mathring{g}^0_{\nu}\left(a_{*}y+\frac{\eta}{2}e_{\nu}\right)\quad\left(y\in\eta_{*}\mathbb{Z}^3\right).
\]
For every $z\in\mathbb{R}^3$, applying $\delta_{\eta,\mu}$ or $J_{\eta}$ to $\mathfrak{g}(\cdot-z)$ gives the translate by $z$ of $\delta_{\eta,\mu}\mathfrak{g}$ or of $J_{\eta}\mathfrak{g}$. The values of $\delta_{\eta,\mu}\mathfrak{g}$ and of $J_{\eta}\mathfrak{g}$ at a point depend only on the values of $\mathfrak{g}$ on a bounded neighbourhood of the point, and on a bounded set a periodisation agrees with a finite sum of translates; therefore both operators commute with periodisations as well.

Let $v\in\eta\mathbb{Z}^3$. The computation that opens the proof of Lemma~\ref{lem:properties-probes}(b) is valid for every such $v$; it gives $a_{\eta,\nu}(x)=\mathring{g}_{\nu}\left(x+\frac{\eta}{2}e_{\nu}\right)$, so that $\hat{\mathfrak{a}}_{\nu}(x)=\eta^{\frac{1}{2}}\cdot\mathring{g}_{\nu}\left(x+\frac{\eta}{2}e_{\nu}\right)$, where $\mathring{g}_{\nu}$ is obtained from $\tau_v\beta$ in the way in which $\mathring{g}^0_{\nu}$ is obtained from $\beta$. Since $\tau_v\beta$ is the periodisation with the period $R^{*}$ of $\beta(\cdot-v)$, the two properties of $\delta_{\eta,\mu}$ and $J_{\eta}$ just stated give
\[
\mathring{g}_{\nu}(z)=\sum_{m\in\mathbb{Z}^3}\mathring{g}^0_{\nu}\left(z-v-R^{*}\cdot m\right)\qquad\left(z\in\mathbb{R}^3\right).
\]
A site of $T_{*}$ with the coordinates $y$ of level $k^{*}$ has the physical coordinates $a_{*}y$, and the definition of $a^{[k]}$ in Section~\ref{sec:lattice-measures-laws} gives $\hat{\mathfrak{a}}^{[k^{*}]}_{\nu}(y)=L^{\frac{k^{*}}{2}}\cdot\hat{\mathfrak{a}}_{\nu}(a_{*}y)$. With the relations $a_{*}=L^{k^{*}}\cdot\eta$, $a_{*}\cdot\eta_{*}=\eta$ and $R^{*}=a_{*}\cdot l_{*}$, which follow from the definitions of Section~\ref{sec:running-quantities-tori}, one obtains
\[
\hat{\mathfrak{a}}_{\nu}(x)=\sum_{m\in\mathbb{Z}^3}w_{\nu}\left(x-v-R^{*}\cdot m\right),\qquad\hat{\mathfrak{a}}^{[k^{*}]}_{\nu}(y)=\sum_{m\in\mathbb{Z}^3}\acute{g}_{\nu}\left(y-\frac{v}{a_{*}}-l_{*}\cdot m\right)
\]
for $\nu\in\{0,1,2\}$ and for all sites $x$ of $\Lambda_{n_{*}}$ and $y$ of $T_{*}$, represented by points of $\eta\mathbb{Z}^3$ and of $\eta_{*}\mathbb{Z}^3$; note that $\frac{v}{a_{*}}\in\eta_{*}\mathbb{Z}^3$. In either sum, the term of index $m$ is called the \emph{copy} of index $m$.

Now let $v$ be admitted. The reasoning in the proof of Lemma~\ref{lem:properties-probes}(b), carried out for the single translate $\beta(\cdot-v)$ in place of $\tau_v\beta$, shows that the end points of the bonds at which the copy of index $0$ of the first sum does not vanish lie in $\bar{Q}$. Since physical coordinates are $a_{*}$ times the coordinates of level $k^{*}$, the end points of the bonds at which the copy of index $0$ of the second sum does not vanish lie in the cube $a_{*}^{-1}\cdot\bar{Q}$. By \eqref{eq:four-jets-1}, by $l_c=\frac{R^{*}}{2L}$ (Section~\ref{sec:cells-reflections}) and by $R^{*}=a_{*}\cdot l_{*}$, the side of this cube satisfies
\[
\frac{D_{\beta}}{a_{*}}\le\frac{3l_c}{4a_{*}}=\frac{3l_{*}}{8L}.
\]
For $m\ne m'$, the copy of index $m'$ of the second sum is the translate of the copy of index $m$ by $l_{*}\cdot(m'-m)$, a vector of sup norm at least $l_{*}$. Hence an end point of a bond at which the copy of index $m$ does not vanish and an end point of a bond at which the copy of index $m'$ does not vanish have distance at least $l_{*}-\frac{3l_{*}}{8L}$, a number larger than $3$ because $l_{*}=M\cdot r_{*}\ge L^4$ (Lemma~\ref{lem:scales-stop-level}(a), (b), and $M\ge L^2$ by stage~\ref{it:order-choice-6}). For two copies of the first sum the corresponding physical distance is larger than $3a_{*}$, and $3a_{*}\ge3\eta$.

By Section~\ref{sec:domains-bond-functions}, each of the three norms of a bond function $A$ on $T_{*}$ is a supremum of numbers of the forms
\[
|A_{\mu}(x)|,\qquad|(\partial_{\nu}A_{\mu})(x)|,\qquad|x-x'|^{-\alpha}\cdot\left|(\partial_{\nu}A_{\mu})(x')-(\partial_{\nu}A_{\mu})(x)\right|\quad\left(0<|x-x'|\le1\right),
\]
and each of these numbers is computed from the values of $A$ at bonds whose end points have mutual distances at most $1+2\eta_{*}$; moreover $1+2\eta_{*}<3$, since $\eta_{*}=L^{-k^{*}}<1$ by $k^{*}\ge1$ (Lemma~\ref{lem:scales-stop-level}(a)). For $A=\hat{\mathfrak{a}}^{[k^{*}]}$, therefore, at most one copy of the second sum differs from zero at the bonds that enter one such number. Every copy of the second sum is a translate of $\acute{g}$ by a vector of $\eta_{*}\mathbb{Z}^3$. So each number of these forms computed with $\hat{\mathfrak{a}}^{[k^{*}]}$ on $T_{*}$ is the number of the same form computed with $\acute{g}$ on $\eta_{*}\mathbb{Z}^3$ at translated bonds, and every non-zero number of these forms computed with $\acute{g}$ arises in this way. Taking suprema one finds
\[
|\hat{\mathfrak{a}}|_{\mathrm{ext}}=|\acute{g}|_{\infty}+|\partial\acute{g}|_{\infty}+|\delta_{\alpha}\partial\acute{g}|_{\infty},
\]
where the norms on the right are given by the formulas of Section~\ref{sec:domains-bond-functions} on the lattice $\eta_{*}\mathbb{Z}^3$. The value of $d\hat{\mathfrak{a}}$ at a plaquette of $\Lambda_{n_{*}}$ is computed from the values of $\hat{\mathfrak{a}}$ at the four bonds of its boundary, whose end points are the vertices of the plaquette and have mutual physical distances at most $\eta$. Hence at most one copy of the first sum contributes to this value. As the plaquettes of $\Lambda_{n_{*}}$ are the classes modulo $R^{*}\mathbb{Z}^3$ of the plaquettes of $\eta\mathbb{Z}^3$, and as the translation by $v$ permutes the latter, one obtains
\[
\boldsymbol{\sigma}^2=|d\hat{\mathfrak{a}}|_2^2=\sum_p\left|(dw)(p)\right|^2,
\]
the sum being taken over all plaquettes $p$ of $\eta\mathbb{Z}^3$, with $d$ the coboundary on $\eta\mathbb{Z}^3$. Apart from the constants $L$ and $\alpha$, the right-hand side of the identity for $|\hat{\mathfrak{a}}|_{\mathrm{ext}}$ depends only on $\beta$, $N$ and $a_{*}$, and that of the identity for $\boldsymbol{\sigma}^2$ only on $\beta$ and $N$, since $\eta=L^{-N}$ and $\eta_{*}=\frac{\eta}{a_{*}}$.

Let $j_{*}=j_{*}(R^{*},e)$ be the number defined in Proposition~\ref{prop:depth-stop-level}(b) for the torus $R=R^{*}$. By Proposition~\ref{prop:depth-stop-level}(c), $N_{\mathrm{min}}(R^{*},e)=\max\left\{0,j_{*}+1-\log_Lh'\right\}$ and $a_{*}=h'\cdot L^{-j_{*}}$, that is, $\log_La_{*}=\log_Lh'-j_{*}$; so $N\ge N_{\mathrm{min}}(R^{*},e)$ holds if and only if $N\ge0$ and $N+\log_La_{*}\ge1$. Thus, in the definitions of $C_{\beta}$ and $\mathbf{D}$, the pair $(N,a_{*})$ ranges over all pairs of an integer $N\ge0$ and a number $a_{*}\in\left\{h'/L,h'/L^2\right\}$ with $N+\log_La_{*}\ge1$, and $r$ over the eight diagonal sign matrices; the quantity under the supremum is, for $C_{\beta}$, the number $|\acute{g}|_{\infty}+|\partial\acute{g}|_{\infty}+|\delta_{\alpha}\partial\acute{g}|_{\infty}$ and, for $\mathbf{D}$, the number $\sum_p\left|(dw)(p)\right|^2$, each formed with $r_{*}\beta$ in place of $\beta$. These data involve $\beta$, $h'$, $L$ and $\alpha$ only; hence $C_{\beta}$ and $\mathbf{D}$ are independent of the charge $e$, of the torus $R^{*}$ and of the placement. Both numbers are finite: at each such pair and for each $r$, the reasoning of the proof of Lemma~\ref{lem:properties-probes}(c) bounds these two quantities by the right-hand sides of the two inequalities of Lemma~\ref{lem:properties-probes}(c).

\paragraph{The external connection.}
For $t=(t_c)_{c\in\mathrm{Cel}}\in\mathbb{C}^{\mathrm{Cel}}$ let
\[
a(t):=\sum_{c\in\mathrm{Cel}}\frac{t_c}{e}\cdot\hat{\mathfrak{a}}_c,\qquad\mathscr{C}(t):=a(t)^{[k^{*}]}=\sum_{c\in\mathrm{Cel}}t_c\cdot\hat{c}_c,
\]
and for $\rho\in(0,\infty]$ let
\[
\mathbb{D}(\rho):=\left\{t\in\mathbb{C}^{\mathrm{Cel}}:\ |t_c|\le\rho\ \text{for all}\ c\right\},\qquad\mathbb{D}^{\circ}(\rho):=\left\{t\in\mathbb{C}^{\mathrm{Cel}}:\ |t_c|<\rho\ \text{for all}\ c\right\}.
\]
For $c\in\mathrm{Cel}$ let $e_c\in\mathbb{C}^{\mathrm{Cel}}$ be the unit vector of the cell $c$.

\begin{lemma}[Properties of the probes]\label{lem:properties-probes}
Let the placement $v$ be admitted. Then the following hold.
\begin{enumerate}
\item[(a)] $\hat{\mathfrak{a}}\in V_{\mathrm{t}}=(\ker d)^{\perp}$, and $\hat{\mathfrak{a}}_c\in V_{\mathrm{t}}$ for every $c\in\mathrm{Cel}$.
\item[(b)] For $\nu\in\{0,1,2\}$ let
\[
\mathring{g}_{\nu}:=-\sum_{\mu\ne\nu}\delta_{\eta,\mu}J_{\eta}(\tau_v\beta)_{\mu\nu}.
\]
These are smooth functions on $\mathbb{R}^3$; for every site $x$ one has $a_{\eta,\nu}(x)=\mathring{g}_{\nu}\left(x+\frac{\eta}{2}e_{\nu}\right)$, and $\left|D^i\mathring{g}_{\nu}\right|_{\infty}\le2|\beta|_{C^{i+1}}$ for every integer $i\ge0$. If $\hat{\mathfrak{a}}$ does not vanish at a bond $b$, then the two end points of $b$ lie in $\bar{Q}$.
\item[(c)] One has
\[
C_{\beta}\le2h'^{\frac{1}{2}}\cdot\left(1+h'+3h'^2\right)\cdot|\beta|_{C^3},\qquad\mathbf{D}\le48|\beta|_{C^2}^2\cdot\left(D_{\beta}+3\right)^3.
\]
\item[(d)] Let $c\in\mathrm{Cel}$. Each of the three norms of Section~\ref{sec:domains-bond-functions} has the same value at $\hat{\mathfrak{a}}_c^{[k^{*}]}$ as at $\hat{\mathfrak{a}}^{[k^{*}]}$, and $|d\hat{\mathfrak{a}}_c|_2^2=\boldsymbol{\sigma}^2$. If $c'\in\mathrm{Cel}$ and $c'\ne c$, there is no plaquette at which both $d\hat{\mathfrak{a}}_c$ and $d\hat{\mathfrak{a}}_{c'}$ differ from zero.
\item[(e)] One has $\#|P_c^{+}|\le\mathsf{p}$ for every $c\in\mathrm{Cel}$, and
\[
d_{\mathrm{sep}}\ge\frac{l_{*}}{8L}-1=\frac{M\cdot r_{*}}{8L}-1.
\]
For $c\ne c'$, every end point of a bond of $\mathrm{Bd}_c$ has distance at least $\frac{l_{*}}{8L}$ from every end point of a bond of $\mathrm{Bd}_{c'}$.
\item[(f)] Let $d_{\mathrm{sep}}\ge4$. For every $t\in\mathbb{C}^{\mathrm{Cel}}$ the three norms of Section~\ref{sec:domains-bond-functions} satisfy
\begin{gather*}
|\mathscr{C}(t)|_{\infty}=\max_{c\in\mathrm{Cel}}\left(|t_c|\cdot|\hat{c}_c|_{\infty}\right),\qquad|\partial\mathscr{C}(t)|_{\infty}=\max_{c\in\mathrm{Cel}}\left(|t_c|\cdot|\partial\hat{c}_c|_{\infty}\right),\\
|\delta_{\alpha}\partial\mathscr{C}(t)|_{\infty}=\max_{c\in\mathrm{Cel}}\left(|t_c|\cdot|\delta_{\alpha}\partial\hat{c}_c|_{\infty}\right).
\end{gather*}
Hence
\[
|a(t)|_{\mathrm{ext}}\le\max_{c\in\mathrm{Cel}}|t_c|\cdot\frac{C_{\beta}}{e}\qquad\left(t\in\mathbb{C}^{\mathrm{Cel}}\right),\qquad|a(t)|_{\mathrm{ext}}<e_{*}^{-1+3\varepsilon}\qquad\left(t\in\mathbb{D}^{\circ}(t_x^{+})\right).
\]
Moreover $t_x\cdot\epsilon_c<e_{*}^{3\varepsilon}$ for every $c\in\mathrm{Cel}$.
\end{enumerate}
\end{lemma}

\begin{proof}
\emph{(a)} Let $d^{*}$ be the adjoint of $d$ for the counting pairings. In every degree the pairing $\eta^3\sum$ is $\eta^3$ times the counting pairing, and $d_{\eta}=\eta^{-1}d$; hence $d_{\eta}^{*}=\eta^{-1}d^{*}$, and
\[
\hat{\mathfrak{a}}=\eta^{\frac{1}{2}}d_{\eta}^{*}u=d^{*}\left(\eta^{-\frac{1}{2}}u\right)
\]
is an element of the image of $d^{*}$. For a linear map between finite-dimensional real spaces with scalar products, the image of the adjoint is the orthogonal complement of the kernel; so $\hat{\mathfrak{a}}\in(\ker d)^{\perp}$. By the Hodge decomposition on the discrete torus, $\ker d=V_{\mathrm{g}}\oplus V_{\mathrm{h}}$, so that $(\ker d)^{\perp}=V_{\mathrm{t}}$ (Section~\ref{sec:lattice-measures-laws}). Further, $(\mathsf{g}_c)_{*}$ preserves the counting pairings and satisfies $d\circ(\mathsf{g}_c)_{*}=(\mathsf{g}_c)_{*}\circ d$. The second property gives $(\mathsf{g}_c)_{*}(\ker d)\subset\ker d$, with equality since $(\mathsf{g}_c)_{*}$ is injective and $\ker d$ is finite-dimensional; together with the first property this shows that $(\mathsf{g}_c)_{*}$ maps $(\ker d)^{\perp}$ onto itself. Hence $\hat{\mathfrak{a}}_c=(\mathsf{g}_c)_{*}\hat{\mathfrak{a}}\in V_{\mathrm{t}}$.

\emph{(b)} For $\mu<\nu$ and a one-cochain $A$ one has
\[
(d_{\eta}A)_{\mu\nu}(x)=\eta^{-1}\left[A_{\mu}(x)+A_{\nu}(x+\eta e_{\mu})-A_{\mu}(x+\eta e_{\nu})-A_{\nu}(x)\right].
\]
Summation by parts gives, with the convention $u_{\nu\mu}=-u_{\mu\nu}$,
\[
(d_{\eta}^{*}u)_{\nu}(x)=-\sum_{\mu\ne\nu}\eta^{-1}\left[u_{\mu\nu}(x)-u_{\mu\nu}(x-\eta e_{\mu})\right].
\]
Let $b$ be the bond from $x$ to $x+\eta e_{\nu}$, and let $c_b:=x+\frac{\eta}{2}e_{\nu}$ be its midpoint. The centres of the plaquettes with the corners $x$ and $x-\eta e_{\mu}$ spanned by $\eta e_{\mu}$ and $\eta e_{\nu}$ are $c_b+\frac{\eta}{2}e_{\mu}$ and $c_b-\frac{\eta}{2}e_{\mu}$. As $u_{\mu\nu}$ is $J_{\eta}(\tau_v\beta)_{\mu\nu}$ evaluated at the centre of the plaquette, the right side of the last display is
\[
-\sum_{\mu\ne\nu}\left(\delta_{\eta,\mu}J_{\eta}(\tau_v\beta)_{\mu\nu}\right)(c_b)=\mathring{g}_{\nu}(c_b).
\]
This is the asserted formula for $a_{\eta,\nu}(x)$.

By the choice of $y_{\beta}$ and $\rho_{\beta}$ one has $v+\operatorname{supp}\beta+[-2,2]^3\subset\bar{Q}$. Hence, for $m\in\mathbb{Z}^3$, the translate $\beta(\cdot-v-R^{*}\cdot m)$ vanishes outside the cube $\bar{Q}+R^{*}\cdot m$; and for $m\ne m'$ the cubes $\bar{Q}+R^{*}\cdot m$ and $\bar{Q}+R^{*}\cdot m'$ have positive distance, because the side $D_{\beta}$ of $\bar{Q}$ is at most $\frac{3}{4}l_c<R^{*}$ by \eqref{eq:four-jets-1}. So every point of $\mathbb{R}^3$ has a neighbourhood on which $\tau_v\beta$ either vanishes or coincides with one of the translates $\beta(\cdot-v-R^{*}\cdot m)$. In particular, at every point each partial derivative of $(\tau_v\beta)_{\mu\nu}$ has a modulus not larger than the supremum of the modulus of the same derivative of $\beta_{\mu\nu}$, and the periodisation plays no role in the estimates below.

Next,
\[
\delta_{\eta,\mu}\mathfrak{g}(y)=\int_{-\frac{1}{2}}^{\frac{1}{2}}\mathrm{d}s\,(\partial_{\mu}\mathfrak{g})(y+s\eta e_{\mu}),
\]
that is, $\delta_{\eta,\mu}\mathfrak{g}(y)$ is the mean value of $\partial_{\mu}\mathfrak{g}$ on the segment from $y-\frac{\eta}{2}e_{\mu}$ to $y+\frac{\eta}{2}e_{\mu}$, while $J_{\eta}\mathfrak{g}(x)$ is the mean value of $\mathfrak{g}$ on the cube $x+\left[-\frac{\eta}{2},\frac{\eta}{2}\right]^3$. Consequently $\delta_{\eta,\mu}$ and $J_{\eta}$ commute with partial derivatives, $|\delta_{\eta,\mu}\mathfrak{g}|_{\infty}\le|\partial_{\mu}\mathfrak{g}|_{\infty}$ and $|J_{\eta}\mathfrak{g}|_{\infty}\le|\mathfrak{g}|_{\infty}$. Applied to the partial derivatives of order $i$ of $(\tau_v\beta)_{\mu\nu}$, together with the bound for the derivatives of $\tau_v\beta$ obtained above, this gives $\left|D^i\delta_{\eta,\mu}J_{\eta}(\tau_v\beta)_{\mu\nu}\right|_{\infty}\le\left|D^{i+1}\beta_{\mu\nu}\right|_{\infty}\le|\beta|_{C^{i+1}}$, and the sum over the two indices $\mu\ne\nu$ gives $\left|D^i\mathring{g}_{\nu}\right|_{\infty}\le2|\beta|_{C^{i+1}}$.

The shifts by $\pm\frac{\eta}{2}e_{\mu}$ in $\delta_{\eta,\mu}$ and the averages over cubes of half-side $\frac{\eta}{2}$ in $J_{\eta}$ show that $\mathring{g}_{\nu}(y)$ is formed with the values of $\tau_v\beta$ on the cube $y+[-\eta,\eta]^3$ alone, whose half-side is the sum $\frac{\eta}{2}+\frac{\eta}{2}$ of the shift and of the half-side of the averaging cube. Hence, if $\mathring{g}_{\nu}(y)\ne0$, there are $m\in\mathbb{Z}^3$ and a point $z\in v+R^{*}\cdot m+\operatorname{supp}\beta$ with $|z-y|\le\eta$. Now let $\hat{\mathfrak{a}}$ be different from zero at the bond $b$. Then $\mathring{g}_{\nu}(c_b)\ne0$; take $m$ and $z$ as above for $y=c_b$. The end points $c_b\pm\frac{\eta}{2}e_{\nu}$ of $b$ have sup distance at most $\frac{3\eta}{2}\le2$ from $z$ (recall $\eta=L^{-N}\le1$). They therefore lie in $v+R^{*}\cdot m+\operatorname{supp}\beta+[-2,2]^3\subset\bar{Q}+R^{*}\cdot m$, and so, as sites, in $\bar{Q}$.

\emph{(c)} The definitions of Section~\ref{sec:running-quantities-tori} give $a_{*}=L^{k^{*}}\cdot\eta$ and $a_{*}\cdot\eta_{*}=\eta$. By the definition of $a^{[k]}$ (Section~\ref{sec:lattice-measures-laws}) and by (b), for a site $x$ of $T_{*}$, in the coordinates of level $k^{*}$, and for $\nu\in\{0,1,2\}$,
\[
\hat{\mathfrak{a}}^{[k^{*}]}_{\nu}(x)=L^{\frac{k^{*}}{2}}\cdot\eta^{\frac{1}{2}}\cdot a_{\eta,\nu}(a_{*}x)=a_{*}^{\frac{1}{2}}\cdot\mathring{g}_{\nu}\left(a_{*}x+\frac{\eta}{2}e_{\nu}\right).
\]
Let $|\mathring{g}|_{\infty}$, $|D\mathring{g}|_{\infty}$ and $\left|D^2\mathring{g}\right|_{\infty}$ be the maxima over $\nu$ of $|\mathring{g}_{\nu}|_{\infty}$, $|D\mathring{g}_{\nu}|_{\infty}$ and $\left|D^2\mathring{g}_{\nu}\right|_{\infty}$. Then $\left|\hat{\mathfrak{a}}^{[k^{*}]}\right|_{\infty}\le a_{*}^{\frac{1}{2}}\cdot|\mathring{g}|_{\infty}$. Moreover $\left(\partial_{\mu}\hat{\mathfrak{a}}^{[k^{*}]}_{\nu}\right)(x)=a_{*}^{\frac{3}{2}}\cdot G(a_{*}x)$ with
\[
G(y):=\int_0^1\mathrm{d}s\,(\partial_{\mu}\mathring{g}_{\nu})\left(y+s\eta e_{\mu}+\frac{\eta}{2}e_{\nu}\right),
\]
and $|G(y)|\le|D\mathring{g}|_{\infty}$. For sites $x,x'$ of $T_{*}$ with $0<|x-x'|\le1$ one has
\[
\left|G(a_{*}x')-G(a_{*}x)\right|\le3a_{*}\cdot|x-x'|\cdot\left|D^2\mathring{g}\right|_{\infty}\le3a_{*}\cdot|x-x'|^{\alpha}\cdot\left|D^2\mathring{g}\right|_{\infty},
\]
the factor $3$ coming from the three coordinate directions and the second inequality from $\alpha<1$ (stage~\ref{it:order-choice-2}). So the three norms of $\hat{\mathfrak{a}}^{[k^{*}]}$ are at most $a_{*}^{\frac{1}{2}}\cdot|\mathring{g}|_{\infty}$, $a_{*}^{\frac{3}{2}}\cdot|D\mathring{g}|_{\infty}$ and $3a_{*}^{\frac{5}{2}}\cdot\left|D^2\mathring{g}\right|_{\infty}$, and by (b)
\[
|\hat{\mathfrak{a}}|_{\mathrm{ext}}\le2a_{*}^{\frac{1}{2}}\cdot|\beta|_{C^1}+2a_{*}^{\frac{3}{2}}\cdot|\beta|_{C^2}+6a_{*}^{\frac{5}{2}}\cdot|\beta|_{C^3}\le2h'^{\frac{1}{2}}\cdot\left(1+h'+3h'^2\right)\cdot|\beta|_{C^3},
\]
because $a_{*}\le h'$ (Lemma~\ref{lem:scales-stop-level}(a)). The numbers $|\beta|_{C^i}$, and also $D_{\beta}$, do not change if $\beta$ is replaced by a translate or by $r_{*}\beta$ with a diagonal sign matrix $r$. So the bound holds for every number of which $C_{\beta}$ is the supremum, and the first inequality follows.

For the second, consider the plaquette with the corner $x$ spanned by $\eta e_{\mu}$ and $\eta e_{\nu}$, and let $b$ and $b'$ be the bonds from $x$ to $x+\eta e_{\nu}$ and from $x$ to $x+\eta e_{\mu}$, with the midpoints $c_b$ and $c_{b'}$. By (b),
\[
(d\hat{\mathfrak{a}})_{\mu\nu}(x)=\eta^{\frac{1}{2}}\left\{\left[\mathring{g}_{\nu}(c_b+\eta e_{\mu})-\mathring{g}_{\nu}(c_b)\right]-\left[\mathring{g}_{\mu}(c_{b'}+\eta e_{\nu})-\mathring{g}_{\mu}(c_{b'})\right]\right\},
\]
whose modulus is at most $2\eta^{\frac{3}{2}}\cdot|D\mathring{g}|_{\infty}\le4\eta^{\frac{3}{2}}\cdot|\beta|_{C^2}$. It vanishes if $\hat{\mathfrak{a}}$ vanishes at the four bonds of the plaquette; otherwise, by (b), a vertex of the plaquette lies in $\bar{Q}$. At most $3\left(\frac{D_{\beta}}{\eta}+3\right)^3$ plaquettes have a vertex in the cube $\bar{Q}$ of side $D_{\beta}$. Hence
\[
\boldsymbol{\sigma}^2\le3\left(\frac{D_{\beta}}{\eta}+3\right)^3\cdot16\eta^3\cdot|\beta|_{C^2}^2=48|\beta|_{C^2}^2\cdot\left(D_{\beta}+3\eta\right)^3\le48|\beta|_{C^2}^2\cdot\left(D_{\beta}+3\right)^3,
\]
as $\eta\le1$. With the invariance just noted, the second inequality follows.

\emph{(d)} Let $\tilde{g}$ be the map $\mathsf{g}_c$ expressed in the coordinates of level $k^{*}$. Lemma~\ref{lem:reflections-lattice-symmetries}(b), applied with $j=k^{*}$, gives $\tilde{g}(x)=rx+L^{-k^{*}}\cdot a_c$, where $r:=r_c=\operatorname{diag}(r_0,r_1,r_2)$ is a diagonal sign matrix; so $\tilde{g}$ preserves the sup distance on $T_{*}$, and $\hat{\mathfrak{a}}_c^{[k^{*}]}=\tilde{g}_{*}\left(\hat{\mathfrak{a}}^{[k^{*}]}\right)$. The map $\tilde{g}$ is a bijection of the set of the bonds of $T_{*}$ onto itself, which preserves the direction of each bond. Let $A$ be a bond function on $T_{*}$. If $b$ is the bond from $x$ to $x+\eta_{*}e_{\nu}$, the definition of the push-forward shows that the component $(\tilde{g}_{*}A)_{\nu}$ takes the value $r_{\nu}\cdot A_{\nu}(x)$ at the bond $\tilde{g} b$; and at the image under $\tilde{g}$ of a lattice cell, $\partial_{\mu}(\tilde{g}_{*}A)_{\nu}$ takes $r_{\mu}r_{\nu}$ times the value that $\partial_{\mu}A_{\nu}$ takes at that cell, functions on cells being evaluated at barycentres (Definition~\ref{def:barycentres}). Since $\tilde{g}$ preserves distances as well, the three norms of $\tilde{g}_{*}A$ coincide with those of $A$; for $A=\hat{\mathfrak{a}}^{[k^{*}]}$ this is the first assertion.

The four vertices of a plaquette lie within physical distance $\eta$ of each other. By (e), for $c\ne c'$ the end points of the bonds of the supports of $\hat{\mathfrak{a}}_c$ and of $\hat{\mathfrak{a}}_{c'}$ have distance at least $\frac{l_{*}}{8L}$ in units of level $k^{*}$, that is, physical distance at least $\frac{a_{*}\cdot l_{*}}{8L}=\frac{l_c}{4}$. Moreover $\frac{l_c}{4}\ge\frac{4}{3}>\eta$, because $\frac{3}{4}l_c\ge D_{\beta}=2\rho_{\beta}\ge4$ by \eqref{eq:four-jets-1}, and $\eta\le1$. Here $\rho_{\beta}\ge2$ because $\beta\ne0$, so that $\operatorname{supp}\beta+[-2,2]^3$ contains a cube of side $4$ (for $\beta=0$ all $\hat{\mathfrak{a}}_c$ would vanish). As in the proof of (c), $d\hat{\mathfrak{a}}_c$ vanishes at a plaquette unless one of its vertices is an end point of a bond of the support of $\hat{\mathfrak{a}}_c$. Hence no plaquette belongs to the supports of both $d\hat{\mathfrak{a}}_c$ and $d\hat{\mathfrak{a}}_{c'}$. Finally, $(\mathsf{g}_c)_{*}$ commutes with $d$ and preserves the counting norm (proof of (a)); so $|d\hat{\mathfrak{a}}_c|_2=\left|(\mathsf{g}_c)_{*}d\hat{\mathfrak{a}}\right|_2=|d\hat{\mathfrak{a}}|_2$.

\emph{(e)} Since $\hat{\mathfrak{a}}_c=(\mathsf{g}_c)_{*}\hat{\mathfrak{a}}$, every bond at which $\hat{\mathfrak{a}}_c$ does not vanish is the image under $\mathsf{g}_c$ of a bond at which $\hat{\mathfrak{a}}$ does not vanish, and the end points of the latter lie in $\bar{Q}$ by (b). By \eqref{eq:four-jets-1}, $\bar{Q}$ is a cube of side $D_{\beta}$ contained in $\mathfrak{c}_0=[0,l_c]^3$, each point of which has distance at least $\frac{l_c}{8}$ from the faces of $\mathfrak{c}_0$; and $\mathsf{g}_c$ maps $\mathfrak{c}_0$ isometrically onto $\mathfrak{c}_c$. So the end points of the bonds at which $\hat{\mathfrak{a}}_c$ does not vanish lie in the cube $\mathsf{g}_c\bar{Q}$ of side $D_{\beta}$, which is contained in $\mathfrak{c}_c$ and whose points have distance at least $\frac{l_c}{8}$ from the faces of $\mathfrak{c}_c$. Let $c\ne c'$, and let $\nu$ be a direction in which the two cells differ. The $\nu$-coordinates of the points of $\mathfrak{c}_c$ and of $\mathfrak{c}_{c'}$ fill two different arcs of length $l_c$ of the circle $\mathbb{R}/R^{*}\mathbb{Z}$, and the $\nu$-coordinate of an end point as above keeps the distance $\frac{l_c}{8}$ from the ends of its arc. Hence an end point for $c$ and an end point for $c'$ have physical distance at least $\frac{l_c}{4}$; in units of level $k^{*}$ it is at least $\frac{l_c}{4a_{*}}=\frac{l_{*}}{8L}$, as $l_c=\frac{R^{*}}{2L}$ (Section~\ref{sec:cells-reflections}) and $R^{*}=a_{*}\cdot l_{*}$ (Section~\ref{sec:running-quantities-tori}). This is the last assertion. If $y\in P_c$, the set $\Delta_y$ contains an end point $x$ of a bond of $\mathrm{Bd}_c$, and $|x-y|<\frac{1}{2}$ because $x\in\Delta_y$. Therefore, for $c\ne c'$, any two points $y\in P_c$ and $y'\in P_{c'}$ satisfy $d(y,y')>\frac{l_{*}}{8L}-\frac{1}{2}-\frac{1}{2}$, that is, $d(P_c,P_{c'})>\frac{l_{*}}{8L}-1$; here $l_{*}=M\cdot r_{*}$ by Lemma~\ref{lem:scales-stop-level}(a). If $\hat{\mathfrak{a}}=0$, then $d_{\mathrm{sep}}=\infty$. This is the bound for $d_{\mathrm{sep}}$. To count $P_c^{+}$, put $s:=\frac{D_{\beta}}{a_{*}}$ and fix a direction $\nu$. In the coordinates of level $k^{*}$ the end points of the bonds of $\mathrm{Bd}_c$ lie in the cube $a_{*}^{-1}\cdot\mathsf{g}_c\bar{Q}$ of side $s$, so that their $\nu$-coordinates lie in an arc $I$ of length $s$ of the circle $\mathbb{R}/l_{*}\mathbb{Z}$. The $\nu$-coordinates of the points of $P_c$ are points of $\mathbb{Z}/l_{*}\mathbb{Z}$ at distance less than $\frac{1}{2}$ from $I$; they lie in an open arc of length $s+1$ and therefore take at most $s+2$ values. Since $d(y,P_c)\le M+1$ for $y\in P_c^{+}$, the $\nu$-coordinates of the points of $P_c^{+}$ lie in an open arc of length $s+1+2(M+1)$ and take at most $s+2M+4$ values. As this holds for each of the three directions, $\#|P_c^{+}|\le(s+2M+4)^3=\mathsf{p}$.

\emph{(f)} The sets $\Delta_y$, $y\in T^0_{*}$, of Section~\ref{sec:lattice-geometry} cover the sites of $T_{*}$: as $L$ is odd, each coordinate of a site of $T_{*}$ differs from the nearest integer by at most $\frac{1-\eta_{*}}{2}<\frac{1}{2}$. Hence every end point of a bond of $\mathrm{Bd}_c$ lies in a set $\Delta_y$ with $y\in P_c$, at distance less than $\frac{1}{2}$ from $y$, and end points of bonds of $\mathrm{Bd}_c$ and of $\mathrm{Bd}_{c'}$, $c\ne c'$, have distance more than $d_{\mathrm{sep}}-1\ge3$. By Section~\ref{sec:domains-bond-functions}, the three norms of a bond function $A$ on $T_{*}$ are suprema of expressions of three kinds: the moduli $|A_{\nu}(x)|$; the moduli of the difference quotients $\eta_{*}^{-1}\cdot\left(A_{\nu}(x+\eta_{*}e_{\mu})-A_{\nu}(x)\right)$; and H\"older quotients, each of which contains the values of $A$ at four bonds. In every one of these expressions the end points of the bonds involved lie within distance $1+2\eta_{*}<3$ of each other ($\eta_{*}=L^{-k^{*}}<1$, as $k^{*}\ge1$ by Lemma~\ref{lem:scales-stop-level}(a)). So no expression reads $A$ on bonds of two different sets $\mathrm{Bd}_c$. As $\hat{c}_c$ vanishes outside $\mathrm{Bd}_c$, each expression at $A=\mathscr{C}(t)=\sum_ct_c\cdot\hat{c}_c$ either vanishes or equals the same expression at $A=t_c\cdot\hat{c}_c$ for one cell $c$. Taking suprema proves the first assertion. By (d), each of the three norms of $\hat{c}_c=e^{-1}\cdot\hat{\mathfrak{a}}_c^{[k^{*}]}$ is $e^{-1}$ times the same norm of $\hat{\mathfrak{a}}^{[k^{*}]}$. Summing the three maxima gives
\[
|a(t)|_{\mathrm{ext}}\le\max_c|t_c|\cdot\frac{|\hat{\mathfrak{a}}|_{\mathrm{ext}}}{e}\le\max_c|t_c|\cdot\frac{C_{\beta}}{e}.
\]
Since $e_{*}=e\cdot a_{*}^{\frac{1}{2}}$ and $a_{*}\le h'$ (Lemma~\ref{lem:scales-stop-level}(a)), one has $\frac{1}{e}=\frac{a_{*}^{\frac{1}{2}}}{e_{*}}\le\frac{h'^{\frac{1}{2}}}{e_{*}}$. For $t\in\mathbb{D}^{\circ}(t_x^{+})$ one has $\max_c|t_c|<t_x^{+}$, and it follows that
\[
|a(t)|_{\mathrm{ext}}\le\max_c|t_c|\cdot h'^{\frac{1}{2}}\cdot C_{\beta}\cdot e_{*}^{-1}<e_{*}^{3\varepsilon}\cdot e_{*}^{-1}=e_{*}^{-1+3\varepsilon};
\]
if $C_{\beta}=0$, then $|a(t)|_{\mathrm{ext}}=0$. For the last assertion note that $t_x<t_x^{+}$, so that $t_x\cdot e_c\in\mathbb{D}^{\circ}(t_x^{+})$, and that $a(t_x\cdot e_c)^{[k^{*}]}=t_x\cdot\hat{c}_c$. Hence
\[
t_x\cdot\epsilon_c=e_{*}\cdot\left|a(t_x\cdot e_c)^{[k^{*}]}\right|_{\infty}\le e_{*}\cdot\left|a(t_x\cdot e_c)\right|_{\mathrm{ext}}<e_{*}^{3\varepsilon}.
\]
\end{proof}

\subsubsection{The averaged law, the partition functions and the free ratio}\label{sec:averaged-law-partition}

The objects of Section~\ref{sec:lattice-measures-laws} are taken on the torus $R^{*}$, so that the lattice is $\Lambda_{n_{*}}=(\mathbb{Z}/n_{*}\mathbb{Z})^3$, and with the pair $\left(m_{\mathrm{in}},\boldsymbol{\varepsilon}_{\mathrm{in}}\right)=P_N(e)=\left(m_{0,N},\boldsymbol{\varepsilon}_{0,N}\right)$ of Definition~\ref{def:tuned-pair}. Thus $F_N=dC^1(\Lambda_{n_{*}})$ is the space of the exact two-cochains, $\gamma_N$ is its standard Gaussian measure, $A(F)$ is the element of $V_{\mathrm{t}}$ with $dA(F)=F$, and $q=e\cdot\eta^{\frac{1}{2}}$. For a complex bond function $A$ on $\Lambda_{n_{*}}$, $D_A$ is the Wilson operator of Section~\ref{sec:lattice-measures-laws}, acting on the antiperiodic spinor fields, with the link phases $\exp\left(i\,q\cdot A(x,x')\right)$; $D_0$ is this operator for $A=0$. For $\theta\in\mathbb{R}^3$, $a\in C^1(\Lambda_{n_{*}})\otimes\mathbb{C}$ and $F\in F_N$,
\begin{align*}
H_{\theta}(x,x+e_{\nu})&:=\frac{\theta_{\nu}}{q\cdot n_{*}},\\
\rho^{\theta,a}_N(F)&:=\exp\left(-\boldsymbol{\varepsilon}_{0,N}\cdot n_{*}^3\right)\cdot\frac{\det\left(D_{A(F)+H_{\theta}+a}+\bar{m}\cdot\eta+m_{0,N}\right)}{\det\left(D_0+\bar{m}\cdot\eta\right)},\\
\nu^{\theta,a}&:=\rho^{\theta,a}_N\cdot\gamma_N,\qquad\bar{\nu}^a:=\int_{[-\pi,\pi]^3}\frac{\mathrm{d}\theta}{(2\pi)^3}\,\nu^{\theta,a},\qquad\bar{\nu}:=\bar{\nu}^0,\qquad\bar{\mu}:=\frac{\bar{\nu}}{\bar{\nu}(F_N)}.
\end{align*}
The slice $\mathrm{Sl}$, the fields $A_{*}$ and $\zeta_j$, $0\le j<k^{*}$, the law $\mathrm{d}\mathsf{m}_{*}$ of $A_{*}$ under $\gamma_N$, the function $\chi_{*}$ and the sets $\mathfrak{G}$ and $\mathfrak{G}^{-}$ are those of Definition~\ref{def:coordinates-slice-small}, on the torus $R^{*}$. For $\theta\in\mathbb{R}^3$, $c_{\theta}:=c^{(k^{*})}_{\theta}$ denotes the bond function of Definition~\ref{def:twist-external-connection}, taken on $T_{*}$ as well as on $T^0_{*}$; it is constant, and its component in the direction $\nu$ equals $\theta_{\nu}/(e_{*}\cdot l_{*})$.

With the external connection $a(t)=\sum_{c\in\mathrm{Cel}}(t_c/e)\cdot\hat{\mathfrak{a}}_c$ of Section~\ref{sec:probes-mirror-copies}, define
\begin{align*}
\mathbf{Z}(t)&:=\bar{\nu}^{a(t)}(F_N),\qquad\mathring{\Phi}(t):=\frac{\mathbf{Z}(t)}{\mathbf{Z}(0)}\qquad\left(t\in\mathbb{C}^{\mathrm{Cel}}\right),\\
N_0(t)&:=\frac{\det\left(D_{(t/e)\cdot\hat{\mathfrak{a}}}+\bar{m}\cdot\eta\right)}{\det\left(D_0+\bar{m}\cdot\eta\right)}\qquad\left(t\in\mathbb{C}\right),\qquad b_{\eta,R^{*}}(0):=\left.\frac{\mathrm{d}^4}{\mathrm{d}t^4}\log N_0(t)\right|_{t=0}.
\end{align*}
The denominators are positive by Lemma~\ref{lem:signs-symmetries-wilson}(a) below.

Each entry of the matrix of $D_{A(F)+H_{\theta}+a(t)}$ is a polynomial in the numbers $\exp\left(\pm i\,q\cdot a(t)(b)\right)$, $b$ a bond of $\Lambda_{n_{*}}$, whose coefficients have modulus at most $3$ and depend continuously on $(F,\theta)$. Hence $\rho^{\theta,a(t)}_N(F)$ is an entire function of $t$, and it is bounded on every compact set of values of $t$, uniformly in $(F,\theta)$. By Fubini's theorem and Morera's theorem, $\mathbf{Z}$, $N_0$ and, for every Borel set $\operatorname{Ev}\subset F_N$, the function $t\mapsto\bar{\nu}^{a(t)}(\operatorname{Ev})$ are entire functions. Since moreover $N_0(0)=1$, the function $\log N_0$ is analytic in a neighbourhood of $t=0$, and $b_{\eta,R^{*}}(0)$ is defined.

\begin{lemma}[Signs and symmetries of the Wilson determinant]\label{lem:signs-symmetries-wilson}
\begin{enumerate}
\item[(a)] Let $a$ be a real bond function on $\Lambda_{n_{*}}$ and $z\in\mathbb{R}$. Then $\det\left(D_a+z\right)\ge0$, and $\operatorname{Re}\left(D_a+z\right)\ge z$, where $\operatorname{Re}X:=\frac{1}{2}\left(X+X^{*}\right)$ for an operator $X$ on the antiperiodic spinor fields with the counting scalar product. In particular $\det\left(D_0+z\right)>0$ for $z>0$. Consequently $\rho^{\theta,a}_N(F)\ge0$ for every real bond function $a$, every $\theta\in\mathbb{R}^3$ and every $F\in F_N$. Moreover $\mathbf{Z}(0)>0$.
\item[(b)] One has $N_0(0)=1$, $N_0(-t)=N_0(t)$ for every $t\in\mathbb{C}$, and $N_0(t)\in\mathbb{R}$ for every $t\in\mathbb{R}$. For every $c\in\mathrm{Cel}$ and every $t\in\mathbb{C}$,
\[
\det\left(D_{(t/e)\cdot\hat{\mathfrak{a}}_c}+\bar{m}\cdot\eta\right)=\det\left(D_{(t/e)\cdot\hat{\mathfrak{a}}}+\bar{m}\cdot\eta\right);
\]
thus the function $N_0$ stays the same when $\hat{\mathfrak{a}}_c$ is put in the place of $\hat{\mathfrak{a}}$.
\end{enumerate}
\end{lemma}

\begin{proof}
\emph{(a)} By Section~\ref{sec:dirac-matrices-lattice}, $\gamma_{\nu}=\sigma_3\otimes\sigma_{\nu+1}$; with respect to the decomposition $\mathbb{C}^4=\mathbb{C}^2\oplus\mathbb{C}^2$ defined by the first factor, $\gamma_{\nu}=\operatorname{diag}\left(\sigma_{\nu+1},-\sigma_{\nu+1}\right)$. Let $a$ be real, and define on the antiperiodic spinor fields the operators
\[
(T_{\nu}f)(x):=\exp\left(i\,q\cdot a(x,x+e_{\nu})\right)\cdot f(x+e_{\nu})\qquad(\nu=0,1,2).
\]
As $a$ is real and $a(x,x-e_{\nu})=-a(x-e_{\nu},x)$, each $T_{\nu}$ is unitary, with the adjoint $(T_{\nu}^{*}f)(x)=\exp\left(i\,q\cdot a(x,x-e_{\nu})\right)\cdot f(x-e_{\nu})$. The Wilson operator of Section~\ref{sec:lattice-measures-laws}, which is the operator of Section~\ref{sec:exact-fermion-operators} at level $0$, acts by
\begin{multline*}
(D_af)(x)=-\sum_{\nu}\left(\frac{1-\gamma_{\nu}}{2}\cdot\exp\left(i\,q\cdot a(x,x+e_{\nu})\right)\cdot f(x+e_{\nu})\right.\\
\left.+\frac{1+\gamma_{\nu}}{2}\cdot\exp\left(i\,q\cdot a(x,x-e_{\nu})\right)\cdot f(x-e_{\nu})-f(x)\right).
\end{multline*}
Therefore
\[
D_a+z=z+\sum_{\nu}\left(1-\frac{T_{\nu}+T_{\nu}^{*}}{2}\right)+\sum_{\nu}\gamma_{\nu}\cdot\frac{T_{\nu}-T_{\nu}^{*}}{2}=\operatorname{diag}\left(\mathbf{B},\mathbf{B}'\right),
\]
with
\begin{align*}
\mathbf{B}&:=z+\sum_{\nu}\left(1-\frac{T_{\nu}+T_{\nu}^{*}}{2}+\sigma_{\nu+1}\cdot\frac{T_{\nu}-T_{\nu}^{*}}{2}\right),\\
\mathbf{B}'&:=z+\sum_{\nu}\left(1-\frac{T_{\nu}+T_{\nu}^{*}}{2}-\sigma_{\nu+1}\cdot\frac{T_{\nu}-T_{\nu}^{*}}{2}\right).
\end{align*}
The matrices $\sigma_{\nu+1}$ are Hermitian and commute with $T_{\nu}$, and $z$ is real; hence $\mathbf{B}'=\mathbf{B}^{*}$, and
\[
\det\left(D_a+z\right)=\det\mathbf{B}\cdot\overline{\det\mathbf{B}}=\left|\det\mathbf{B}\right|^2\ge0.
\]
In the expression for $D_a+z$, the first sum is a non-negative Hermitian operator, because $\left\|T_{\nu}\right\|=1$, and the second sum is anti-Hermitian, because $\gamma_{\nu}$ is Hermitian and commutes with $T_{\nu}$. So $\operatorname{Re}\left(D_a+z\right)\ge z$. If $z>0$, then $\operatorname{Re}\left\langle f,\left(D_0+z\right)f\right\rangle\ge z\cdot|f|_2^2$ for every $f$; thus $D_0+z$ is injective, hence invertible, and $\det\left(D_0+z\right)>0$.

Now let $a$ be a real bond function, $\theta\in\mathbb{R}^3$ and $F\in F_N$. The bond functions $A(F)$ and $H_{\theta}$ take real values, and therefore so does the sum $A(F)+H_{\theta}+a$; by Theorem~\ref{thm:one-pair-all} the pair $\left(m_{0,N},\boldsymbol{\varepsilon}_{0,N}\right)$ is real, and $\bar{m}\cdot\eta>0$. Applying what has just been proved, once to the connection $A(F)+H_{\theta}+a$ with $z=\bar{m}\cdot\eta+m_{0,N}$ and once to the connection $0$ with $z=\bar{m}\cdot\eta$, one finds that the fraction in the definition of $\rho^{\theta,a}_N(F)$ has a non-negative numerator and a positive denominator; the exponential factor in front of it is positive. Therefore $\rho^{\theta,a}_N(F)\ge0$. Theorem~\ref{thm:one-pair-all} also gives $|m_{0,N}|\le w^{\mathrm{m}}_0/8$, where $w^{\mathrm{m}}_0=q^{\frac{3}{4}-8\varepsilon}$ is formed with $q$ (Section~\ref{sec:running-quantities-tori} and Definition~\ref{def:notation}(e)). One has $q=e\cdot\eta^{\frac{1}{2}}\le e_{*}<1$: indeed $a_{*}=L^{k^{*}}\cdot\eta\ge\eta$, $e_{*}\le e_0$ by Lemma~\ref{lem:scales-stop-level}(a), and $e_0<1$ by condition~\eqref{eq:renormalisation-group-A8}. Hence $w^{\mathrm{m}}_0<1$, so that $|m_{0,N}|<\frac{1}{2}$, and the lattice mass $\boldsymbol{\mu}:=\bar{m}\cdot\eta+m_{0,N}$ exceeds $-1$.

Finally, $\mathbf{Z}(0)=\bar{\nu}(F_N)$. The coboundary $d$ maps $V_{\mathrm{t}}$ bijectively onto $F_N$, with $A(dA)=A$ for $A\in V_{\mathrm{t}}$; let $\sigma_{V_{\mathrm{t}}}$ be the Lebesgue measure of the Euclidean space $V_{\mathrm{t}}$. The image under $d$ of the measure on $V_{\mathrm{t}}$ with the density $\exp\left(-\frac{1}{2}|dA|_2^2\right)$ with respect to $\sigma_{V_{\mathrm{t}}}$ is a positive multiple of $\gamma_N$. Since the factor $\exp\left(-\boldsymbol{\varepsilon}_{0,N}\cdot n_{*}^3\right)/\det\left(D_0+\bar{m}\cdot\eta\right)$ is positive, $\mathbf{Z}(0)$ is a positive multiple of the number $\overline{Z}[1]$ of~\eqref{eq:gaussian-bounds-3}, formed with $n=n_{*}$, the charge $q$, the lattice mass $\boldsymbol{\mu}$ and the constant cochains $c^{\theta}=H_{\theta}$. The hypotheses of Theorem~\ref{thm:limit-vanishing-proca} hold: $n_{*}\ge2$, $q\ne0$ and $\boldsymbol{\mu}>-1$. By that theorem $\overline{Z}[1]>0$, and therefore $\mathbf{Z}(0)>0$.

\emph{(b)} By Lemma~\ref{lem:scales-stop-level}(a) one has $k^{*}\ge1$, so that Lemma~\ref{lem:scales-stop-level}(c), taken with $j=0$, shows that hypothesis~\ref{def:hypothesis-one-step} holds for the data $\left(0,l_0,\bar{m}\cdot\eta,q\right)$. In particular the side $n_{*}=l_0$ of the torus lies in $\mathrm{Cl}$; the mass satisfies $0<\bar{m}\cdot\eta\le\mathfrak{m}_0/L\le\mathfrak{m}_0$; and the charge satisfies $q>0$ and $L^{\frac{1}{2}}\cdot q\le e_0$, hence $q\in(0,e_0]$. Thus the data $\left(0,n_{*},\bar{m}\cdot\eta,q\right)$ belong to the scale-free setting of Section~\ref{sec:properties-scheme}, and Proposition~\ref{prop:symmetries-fermion-operators} is available at the level $0$, for the connections $A$ that belong to the set $\tilde{G}_0$ of Definition~\ref{def:domains}, formed at the level $0$. Let $m\in\mathbb{C}$ and $A\in\tilde{G}_0$. Clause~(c) of Proposition~\ref{prop:symmetries-fermion-operators} gives $D_A^T=\mathsf{C}^{-1}\cdot D_{-A}\cdot\mathsf{C}$; adding $m$ on both sides and taking determinants, one obtains
\[
\det\left(D_{-A}+m\right)=\det\left(D_A+m\right).
\]
Let $c\in\mathrm{Cel}$. Lemma~\ref{lem:reflections-lattice-symmetries}(b), taken with $j=0$, provides an element $\tilde{g}$ of $\mathbb{G}_0^L\subset\mathbb{G}_0$ that induces the map $\mathsf{g}_c$. For this element, clause~(b) of Proposition~\ref{prop:symmetries-fermion-operators} yields $D_{\tilde{g}_{*}A}\cdot\tilde{g}=\tilde{g}\cdot D_A$, whence
\[
\det\left(D_{\tilde{g}_{*}A}+m\right)=\det\left(D_A+m\right).
\]
The matrix of $D_A$ depends on $A$ only through the link phases $\exp\left(\pm i\,q\cdot A(b)\right)$, $b$ a bond, and its entries are polynomial in them; as $\tilde{g}_{*}$ is linear, it follows that, for fixed $m$, each side of each of the two relations is an entire function of the complex bond function $A$. The two sides agree on $\tilde{G}_0$, which is a non-empty open subset of the space of the complex bond functions on $\Lambda_{n_{*}}$ (Section~\ref{sec:domains-bond-functions}; it contains $0$). By the identity theorem for holomorphic functions of several complex variables, both relations hold for every complex bond function $A$ and every $m\in\mathbb{C}$.

Now let $t\in\mathbb{C}$, and take $A=(t/e)\cdot\hat{\mathfrak{a}}$ and $m=\bar{m}\cdot\eta$. As $-A=(-t/e)\cdot\hat{\mathfrak{a}}$, the first relation gives $N_0(-t)=N_0(t)$. As $\tilde{g}_{*}\hat{\mathfrak{a}}=(\mathsf{g}_c)_{*}\hat{\mathfrak{a}}=\hat{\mathfrak{a}}_c$, one has $\tilde{g}_{*}A=(t/e)\cdot\hat{\mathfrak{a}}_c$, and the second relation becomes the equality of determinants stated in~(b). At $t=0$ the numerator and the denominator of $N_0(t)$ coincide, so that $N_0(0)=1$. Let $t$ be real. Then $(t/e)\cdot\hat{\mathfrak{a}}$ is a real bond function, and the link phases of $D_{(t/e)\cdot\hat{\mathfrak{a}}}$ are
\[
\exp\left(\pm i\,q\cdot\frac{t}{e}\cdot\hat{\mathfrak{a}}(b)\right)=\exp\left(\pm i\,t\cdot\eta\cdot a_{\eta}(b)\right),
\]
because $q=e\cdot\eta^{\frac{1}{2}}$ and $\hat{\mathfrak{a}}=\eta^{\frac{1}{2}}\cdot a_{\eta}$ (Section~\ref{sec:probes-mirror-copies}). By~(a), applied with $z=\bar{m}\cdot\eta$, the determinant $\det\left(D_{(t/e)\cdot\hat{\mathfrak{a}}}+\bar{m}\cdot\eta\right)$ is real, and the denominator of $N_0(t)$ is positive; so $N_0(t)$ is real.
\end{proof}

\subsubsection{The ring of jets}\label{sec:ring-jets}

Let
\[
\mathcal{J}:=\mathbb{C}\left[t_c:c\in\mathrm{Cel}\right]\big/\left(t_c^5:c\in\mathrm{Cel}\right).
\]
For a multi-index $\mathbf{a}=(\mathbf{a}_c)_{c\in\mathrm{Cel}}\in\{0,1,2,3,4\}^{\mathrm{Cel}}$ put $t^{\mathbf{a}}:=\prod_ct_c^{\mathbf{a}_c}$ and $|\mathbf{a}|:=\sum_c\mathbf{a}_c$. Every $P\in\mathcal{J}$ is the class of a unique polynomial $\sum_{\mathbf{a}}P_{\mathbf{a}}\cdot t^{\mathbf{a}}$, the sum extending over these multi-indices; for $\rho>0$ its norm is
\[
\|P\|_{\rho}:=\sum_{\mathbf{a}}|P_{\mathbf{a}}|\cdot\rho^{|\mathbf{a}|}.
\]
One has $\|P\cdot P'\|_{\rho}\le\|P\|_{\rho}\cdot\|P'\|_{\rho}$ for $P,P'\in\mathcal{J}$, since the coefficient of $t^{\mathbf{a}''}$ in $P\cdot P'$ is $\sum_{\mathbf{a}+\mathbf{a}'=\mathbf{a}''}P_{\mathbf{a}}\cdot P'_{\mathbf{a}'}$ and $|\mathbf{a}+\mathbf{a}'|=|\mathbf{a}|+|\mathbf{a}'|$.

Let $\mathfrak{g}$ be holomorphic in a neighbourhood of $t=0$ in $\mathbb{C}^{\mathrm{Cel}}$. Retaining from its Taylor series at $t=0$ the terms in which every variable $t_c$ occurs with an exponent at most $4$, one obtains a polynomial; its class in $\mathcal{J}$ is again denoted by $\mathfrak{g}$. On the ring of the functions holomorphic near $t=0$, two of them being identified when they agree near $t=0$, this assignment is a homomorphism of rings into $\mathcal{J}$. The element of $\mathcal{J}$ assigned to $\exp(\mathfrak{g})$ is $\sum_{n\ge0}\frac{1}{n!}\mathfrak{g}^n$, with the powers of the element $\mathfrak{g}$ formed in $\mathcal{J}$ (the series converges, as $\|\cdot\|_{\rho}$ is submultiplicative).

\subsubsection{Constants}\label{sec:constants-3}

The following constants of Section~\ref{sec:order-choice} are fixed: $K_0$, $\kappa_0$, $c_X$ and $c_{\mathrm{rb}}$; $b=1$ and $b_{-}=1-\frac{1}{L}$; $\kappa$; $L$; $N^{*}=8L^3$; $m_0$ and $h'$; $M$; the constants $C_F\ge1$ and $\gamma_F\in(0,1]$ of Propositions~\ref{prop:decay-fine-propagator}--\ref{prop:bounds}, which satisfy $\gamma_F\cdot M\ge16$, because $\gamma_w\le\gamma_F$ and by stage~\ref{it:order-choice-6}; $C_0\ge\max(1,C_F)$ (stage~\ref{it:order-choice-7}); the constant $C_d$ of Proposition~\ref{prop:local-expansion-determinant}; $C_Q=1$ (Proposition~\ref{prop:averaging-operators}(d)); and $r$, $p_0$, $p$, $\varepsilon$ and $s$.

The following numbers are functions of the constants of stages~\ref{it:order-choice-1}--\ref{it:order-choice-7}, with the exception of $\tau_s$, which contains $l_{*}$ and therefore depends on $e$:
\begin{align*}
\mathrm{ssum}(\tilde{\gamma})&:=\left(\frac{12}{\tilde{\gamma}}\right)^3\qquad(\tilde{\gamma}>0),\\
\mathbf{S}&:=C_F\cdot\mathrm{ssum}\left(\frac{\gamma_F}{2}\right),\qquad h_L:=6\,\mathbf{S},\\
C_K&:=100\,\mathbf{S}^2,\qquad C_K':=400\,\mathbf{S}^2,\qquad C_H:=100\,M\cdot C_F^2,\\
C_V&:=10^3\,C_H^2\cdot\mathrm{ssum}\left(\frac{\gamma_F}{2}\right)^4\cdot K_0,\\
\boldsymbol{\theta}_{\mathrm{CT}}&:=\frac{\mathfrak{m}_{-}}{12\exp(1)},\qquad\vartheta:=\min\left\{\frac{\boldsymbol{\theta}_{\mathrm{CT}}}{2},\frac{\gamma_F}{8},\frac{1}{M}\right\},\\
g&:=\left(b_{-}^{-1}+\mathfrak{m}_{-}^{-1}\right)^{\frac{1}{2}},\qquad h_G:=1+g,\qquad g_1:=b_{-}^{-1}+\frac{4}{\mathfrak{m}_{-}},\\
\boldsymbol{\alpha}&:=2\,b_{-}^{-1}+4\exp(1)\cdot\mathfrak{m}_{-}^{-1}\cdot\mathrm{ssum}\left(\frac{\boldsymbol{\theta}_{\mathrm{CT}}}{2}\right),\\
c_1&:=\min\left\{\frac{1}{8},\frac{\mathfrak{m}_{-}}{48},\frac{1}{80\,\mathbf{S}}\right\},\qquad\tau_s:=c_1\cdot l_{*},\\
C_W&:=3\,C_V\cdot h_G^2+2\,\boldsymbol{\alpha}\cdot C_V^2\cdot h_G^4+4\exp(6)\cdot C_d\cdot M^3\cdot K_0.
\end{align*}
For $0<\tilde{\gamma}\le6$, every integer $l\ge1$ and every $y\in(\mathbb{Z}/l\mathbb{Z})^3$ one has
\[
\sum_{y'\in(\mathbb{Z}/l\mathbb{Z})^3}\exp\left(-\tilde{\gamma}\cdot d(y,y')\right)\le\mathrm{ssum}(\tilde{\gamma}).
\]
Indeed, fix a lift $\tilde{y}\in\mathbb{Z}^3$ of $y$. Each class $y'$ can be lifted to a point $\tilde{y}'\in\mathbb{Z}^3$ with $|\tilde{y}'-\tilde{y}|=d(y,y')$, and distinct classes are lifted to distinct points. As $|n|\ge|n|_1/3$ for $n\in\mathbb{Z}^3$, the sum is at most
\[
\sum_{n\in\mathbb{Z}^3}\exp\left(-\frac{\tilde{\gamma}\cdot|n|_1}{3}\right)=\left(\sum_{j\in\mathbb{Z}}\exp\left(-\frac{\tilde{\gamma}\cdot|j|}{3}\right)\right)^3=\left(\coth\frac{\tilde{\gamma}}{6}\right)^3,
\]
and $\coth\frac{\tilde{\gamma}}{6}\le1+\frac{6}{\tilde{\gamma}}\le\frac{12}{\tilde{\gamma}}$, the last inequality because $\tilde{\gamma}\le6$.

The numbers $C_{\beta}$, $\mathbf{D}$ and $D_{\beta}$ of Section~\ref{sec:probes-mirror-copies} belong to stage~\ref{it:order-choice-11}; so does the constant
\[
C_{\mathrm{fr}}:=2592\,C_d\cdot K_0\cdot C_{\beta}^2\cdot\left(M^2\cdot D_{\beta}^3\cdot L^4\cdot h'^{-2}+8\,M^5\cdot h'\right).
\]
Lemma~\ref{lem:scales-stop-level}(a) allows the two values $h'/L$ and $h'/L^2$ for $a_{*}$; accordingly the number $\mathsf{p}=\left(D_{\beta}/a_{*}+2M+4\right)^3$ of Section~\ref{sec:probes-mirror-copies} takes two values. Since $a_{*}\ge h'/L^2$, each of them satisfies
\[
\mathsf{p}\le\left(\frac{D_{\beta}\cdot L^2}{h'}+2M+4\right)^3.
\]

The following quantities depend on $e$, with the exception of $z_2$ and $M_L$, which are constants of stage~\ref{it:order-choice-11}:
\begin{align*}
\bar{\kappa}&:=6\,e_{*}^{\frac{1}{4}-9\varepsilon},\qquad\bar{r}:=4\,h_G^2\cdot C_K\cdot e_{*}^{2\varepsilon}+650\,K_0\cdot e_{*}^{\frac{1}{4}-7\varepsilon},\\
B_1&:=16\,g_1\cdot C_K'\cdot\mathsf{p}\cdot e_{*}^{2\varepsilon},\\
B_{\mathrm{wr}}&:=16\exp(6)\cdot C_d\cdot M^3\cdot K_0\cdot\mathsf{p}\cdot\exp(-r_{*}),\\
B_0&:=C_{\mathrm{fr}}\cdot e_{*}^{4\varepsilon}+B_{\mathrm{wr}}+B_1,\\
\varpi_{\mathrm{sep}}&:=2^{N^{*}}\cdot C_W\cdot\mathsf{p}\cdot\exp\left(-\vartheta\cdot\left(d_{\mathrm{sep}}-4M-6\right)\right),\\
\varpi_0&:=18\,l_{*}^3\cdot\exp\left(-\frac{p_{0*}}{2}\right)+6\,l_{*}^3\cdot\exp\left(-\frac{(p_{*}-1)^2}{2C_0^2}\right),\\
z_2&:=\min\left\{\left(N^{*}\right)^{-\frac{1}{2}},\ \frac{1}{C_{\beta}},\ \frac{1}{2\sqrt{N^{*}\cdot\mathbf{D}}}\right\},\qquad t_2:=e\cdot z_2,\\
\varpi_L&:=\varpi_0^{\frac{1}{2}}\cdot\exp\left(\frac{\mathbf{D}}{2}\right),\qquad M_L:=2+\exp\left(\frac{\mathbf{D}}{2}\right),\\
\boldsymbol{\lambda}_L&:=\left(\varpi_L^{\frac{1}{3}}\cdot M_L^{\frac{2}{3}}\right)^{\frac{1}{N^{*}}}\cdot5\cdot\left(\frac{2N^{*}\cdot\mathbf{r}}{t_2}\right)^4.
\end{align*}
In the minimum that defines $z_2$, an entry whose denominator is $0$ is read as $+\infty$. As $C_Q=1$, the number $z_2$ is the number $z_2$ of Section~\ref{sec:probes-jets-reference}, in which the factor $2$ of the last entry is written $\max(2,C_Q)$.

\subsubsection{Smallness conditions}\label{sec:smallness-conditions}

The smallness conditions are the following.
\begin{align}
&5\,h_L\cdot h_G\le e_{*}^{-\frac{1}{4}+\varepsilon},\qquad\left(4+\frac{M^3}{2}\right)\cdot e_{*}^{\varepsilon}\le1,\qquad\boldsymbol{\alpha}\cdot\bar{r}\le\frac{1}{2};\tag{sm1}\label{eq:four-jets-sm1}\\
&d_{\mathrm{sep}}\ge6M+8,\qquad\text{and the placement is admitted, that is, \eqref{eq:four-jets-1} holds};\tag{sm2}\label{eq:four-jets-sm2}\\
&4\,C_0\cdot l_{*}^3\cdot g_1\le e_{*}^{-\frac{1}{4}+\varepsilon},\qquad4\,l_{*}^3\cdot\left(2C_K+K_0+4\right)\le e_{*}^{-2\varepsilon},\qquad e_{*}^{\frac{1}{4}-9\varepsilon}\le\frac{1}{8};\tag{sm3}\label{eq:four-jets-sm3}
\end{align}
\begin{multline}
48\,e_{*}^{2\varepsilon}\le\mathfrak{m}_{-},\qquad2\,g_1\cdot C_K'\cdot e_{*}^{2\varepsilon}\le\frac{1}{2},\qquad\tau_s\ge1,\\
52\,B_1\cdot\exp(-\tau_s)\le e_{*}^{\frac{1}{4}-9\varepsilon},\qquad B_0\le\frac{1}{4};\tag{sm4}\label{eq:four-jets-sm4}
\end{multline}
\begin{align}
&2\,\varpi_{\mathrm{sep}}\le\bar{\kappa}^{N^{*}};\tag{sm5}\label{eq:four-jets-sm5}\\
&N^{*}\cdot\bar{\kappa}^{1-s}\le\frac{1}{2};\tag{sm6}\label{eq:four-jets-sm6}\\
&\varpi_0\le\frac{1}{4},\qquad8\,N^{*}\cdot t_2\le t_x;\tag{sm7}\label{eq:four-jets-sm7}\\
&\boldsymbol{\lambda}_L\le\bar{\kappa}^s.\tag{sm8}\label{eq:four-jets-sm8}
\end{align}

Under hypothesis~\ref{def:hypothesis-reference-torus} the following conditions of Sections~\ref{sec:order-choice} and~\ref{sec:conditions} hold, by stages~\ref{it:order-choice-9} and~\ref{it:order-choice-10} and, for the conditions on $(R^{*},e)$, by Proposition~\ref{prop:consequences}.
\begin{itemize}
\item Condition~\eqref{eq:renormalisation-group-A8} contains $e_0^{\varepsilon}\le\frac{1}{4}$ and $3M\cdot e_0^{\varepsilon}\le1$. Since $e_{*}\le e_0$ (Lemma~\ref{lem:scales-stop-level}(a)), it follows that $e_{*}^{2\varepsilon}\le\frac{1}{16}$ and $3M\cdot e_{*}^{2\varepsilon}\le\frac{1}{4}$.
\item Condition~\eqref{eq:renormalisation-group-16} reads $r_{*}>31L+16$.
\item Condition~\eqref{eq:renormalisation-group-T1} contains $v_{*}\ge8$ and $p_0\cdot v_{*}^{p_0-1}\ge28$.
\item Condition~\eqref{eq:renormalisation-group-T3} reads $4\,C_0\cdot l_{*}^2\cdot p_{*}<e_{*}^{-\varepsilon}$.
\item Stages~\ref{it:order-choice-3} and~\ref{it:order-choice-4} give $\kappa\ge7\kappa_0+7$, $L\cdot\left(\kappa-\kappa_0-1\right)\ge\kappa$ (as $L\cdot\left(\kappa-6\kappa_0-4\right)-1\ge\kappa$, and $\kappa_0>0$ by stage~\ref{it:order-choice-1}), and $c_{\mathrm{rb}}\cdot L^3\cdot\exp\left(-\left(L\cdot\left(\kappa-\kappa_0-1\right)-\kappa\right)\right)\le\frac{1}{2}$.
\item Condition~\eqref{eq:renormalisation-group-A10} is a condition on the constants alone. It consists of the inequality $e_0\le e^F_0$, where $e^F_0$ is the number of Theorem~\ref{thm:fermion-operators-weakened} at $\kappa'=\kappa_w+\kappa_0$; of the four inequalities
\[
v_k\ge9,\qquad8M\cdot\left(\frac{v_k}{4}+4\right)\cdot e_k^{\varepsilon}\le1,\qquad8\,C_B\cdot p_k\le e_k^{-\frac{3}{4}+3\varepsilon},\qquad8\,e_k^{\frac{3}{4}-\frac{7\varepsilon}{2}}\le1,
\]
imposed for every $e_k\in(0,e_0]$, with $v_k=-\log e_k$ and $p_k=v_k^p$; and of condition~\eqref{eq:renormalisation-group-24}, in which $C_d$ is formed with the constant $C_w$ of $\mathrm{par}$. The first two of the four inequalities are the second and the third inequality of Proposition~\ref{prop:local-expansion-determinant}(d4), and the last two are the two displayed inequalities of Proposition~\ref{prop:fields-domains}, with the constant $C_B$ of Proposition~\ref{prop:local-regularity}; condition~\eqref{eq:renormalisation-group-24} is the condition on $e_0$ on which Proposition~\ref{prop:local-expansion-determinant}(d6) rests. Since $e_{*}\in(0,e_0]$ by Lemma~\ref{lem:scales-stop-level}(a), and since the running charge of every level $j<k^{*}$ lies in $(0,e_0]$ by Lemma~\ref{lem:scales-stop-level}(c), the four inequalities hold at $e_{*}$ and at the running charges of the levels $j<k^{*}$. The requirement $l\ge M\cdot r_k$ of Proposition~\ref{prop:local-expansion-determinant}(d4), the requirements $l/L\in\mathrm{Cl}$ and $e_{k+1}=L^{\frac{1}{2}}\cdot e_k\le e_0$ of Proposition~\ref{prop:local-expansion-determinant}(d5), (d6), and the requirement $l/L\in\mathrm{Cl}$ of Proposition~\ref{prop:fields-domains}(ii) are requirements on the data $(k,l,e_k)$; they are part of hypothesis~\ref{def:hypothesis-one-step}, which holds for the data of the levels $j<k^{*}$ by Lemma~\ref{lem:scales-stop-level}(c). Proposition~\ref{prop:fields-domains} assumes moreover the conditions~\eqref{eq:renormalisation-group-A1} and~\eqref{eq:renormalisation-group-A2}, which hold by stage~\ref{it:order-choice-9} as well.
\end{itemize}

\subsection{Gaussian domination and the stop level}\label{sec:gaussian-domination-stop}

\subsubsection{Gaussian domination}\label{sec:gaussian-domination}

\begin{proposition}[Sub-Gaussian bounds for the averaged law]\label{prop:sub-gaussian-bounds}
Let $\langle\cdot,\cdot\rangle$ and $|\cdot|_2$ be the counting scalar product and the counting norm on the real two-cochains on $\Lambda_{n_{*}}$, and let $P$ be the orthogonal projection onto $F_N$. Each real two-cochain $\mathbf{g}$ on $\Lambda_{n_{*}}$ satisfies
\[
\int\mathrm{d}\bar{\mu}(F)\,\exp\left\langle\mathbf{g},F\right\rangle\le\exp\left(\frac{1}{2}\left|P\mathbf{g}\right|_2^2\right).
\]
For every $\phi\in F_N$ with $|\phi|_2\le1$ and every $y\ge0$,
\[
\bar{\mu}\left(\left\{F\in F_N:\left|\langle\phi,F\rangle\right|>y\right\}\right)\le2\exp\left(-\frac{y^2}{2}\right).
\]
\end{proposition}

\begin{proof}
The coboundary $d$ restricts to a linear isomorphism of $V_{\mathrm{t}}$ onto $F_N$, and $A(dA)=A$ for $A\in V_{\mathrm{t}}$. Let $\sigma_{V_{\mathrm{t}}}$ be the Lebesgue measure of the Euclidean space $V_{\mathrm{t}}$; the image under $d$ of the measure on $V_{\mathrm{t}}$ with the density $\exp\left(-\frac{1}{2}|dA|_2^2\right)$ with respect to $\sigma_{V_{\mathrm{t}}}$ is a positive multiple of $\gamma_N$. One has $\rho^{\theta,0}_N\ge0$ and $\bar{\nu}(F_N)=\mathbf{Z}(0)>0$ by Lemma~\ref{lem:signs-symmetries-wilson}(a), and the positive factor $\exp\left(-\boldsymbol{\varepsilon}_{0,N}\cdot n_{*}^3\right)/\det\left(D_0+\bar{m}\cdot\eta\right)$ cancels in the normalisation. Hence, by the definition of $\bar{\mu}$ in Section~\ref{sec:lattice-measures-laws}, integration with respect to $\bar{\mu}$ is the functional $\overline{\mathsf{E}}$ of~\eqref{eq:gaussian-bounds-3}, that is, $\bar{\mu}$ is the harmonic-averaged law, formed with $n=n_{*}\ge2$, the charge $q=e\cdot\eta^{\frac{1}{2}}\ne0$, the constant cochains $c^{\theta}=H_{\theta}$ and the lattice mass $\boldsymbol{\mu}=\bar{m}\cdot\eta+m_{0,N}$, which exceeds $-1$ by the proof of Lemma~\ref{lem:signs-symmetries-wilson}(a). Therefore Theorem~\ref{thm:limit-vanishing-proca} applies, and the first assertion is the inequality~\eqref{eq:gaussian-bounds-5}.

Let $\phi\in F_N$ with $|\phi|_2\le1$, and let $y\ge0$. If $\phi=0$, the set on the left of the second assertion is empty. Let $\phi\ne0$. Then $P\phi=\phi$, and Corollary~\ref{cor:linear-functionals-curvature}(iii), applied to the two cochains $\phi$ and $-\phi$, for which $\sigma_{\pm\phi}=\left|P(\pm\phi)\right|_2=|\phi|_2\in(0,1]$, gives
\[
\bar{\mu}\left(\left\{F\in F_N:\pm\langle\phi,F\rangle\ge y\right\}\right)\le\exp\left(-\frac{y^2}{2|\phi|_2^2}\right)\le\exp\left(-\frac{y^2}{2}\right)
\]
for both signs. The set $\left\{F\in F_N:\left|\langle\phi,F\rangle\right|>y\right\}$ is contained in the union of these two sets, and the second assertion follows.
\end{proof}

\subsubsection{Fermionic operators at the stop level and majorants}\label{sec:fermionic-operators-stop}

Let $A$ be a complex bond function on $T_{*}$. Here $\eta_{*}=L^{-k^{*}}\le\frac{1}{L}$, because $k^{*}\ge1$ (Lemma~\ref{lem:scales-stop-level}(a)). The matrices $\gamma_{\nu}$ and $\mathsf{C}$ are those of Section~\ref{sec:dirac-matrices-lattice}; the axial functional $\tau$ is that of Section~\ref{sec:averaging-fluctuations-gauge}, and $\tau_{k^{*}}$ is formed from it as in Section~\ref{sec:exact-fermion-operators}; the set $\mathrm{Pz}_{*}$ of pairs and the six rectilinear paths $\Gamma^{\pi}(x,x')$ of a pair are those of Definition~\ref{def:pairs-sites-holder}; $\alpha$ is the H\"older order. Spinor functions on $T_{*}$ are antiperiodic, as in Section~\ref{sec:exact-fermion-operators}. The norm $|v|=\max_{\boldsymbol{\beta}}|v_{\boldsymbol{\beta}}|$ on $\mathbb{C}^4$, the norm $|f|_{\infty}$ formed with it, the sup norm $|x|$ and the $\ell^1$ norm $|x|_1$ of a vector $x$, and the sup distance $d$ are those of Section~\ref{sec:conventions-2}.

For $x\in T_{*}$, a direction $\nu\in\{0,1,2\}$ and a sign $\boldsymbol{\sigma}\in\{1,-1\}$ put
\[
x^{\boldsymbol{\sigma}}:=x+\boldsymbol{\sigma}\cdot\eta_{*}\cdot e_{\nu},\qquad P^{\boldsymbol{\sigma}}_{\nu}:=\frac{1}{2}\left(1-\boldsymbol{\sigma}\cdot\gamma_{\nu}\right),
\]
and define
\begin{align*}
u_A(x,x')&:=\exp\left(i\cdot e_{*}\cdot\eta_{*}\cdot A(x,x')\right)&&\text{for the bonds }(x,x')\text{ of }T_{*},\\
\sigma_A(y,x)&:=e_{*}\cdot\eta_{*}\cdot\left(\tau_{k^{*}}A\right)(y,x)&&\text{for }y\in T^0_{*}\text{ and }x\in\Delta_y,\\
\upsilon_A(x,x')&:=e_{*}\cdot\eta_{*}\cdot(\tau A)(x,x'),\qquad U_A(x,x'):=\exp\left(i\cdot\upsilon_A(x,x')\right)&&\text{for }(x,x')\in\mathrm{Pz}_{*}.
\end{align*}
In this notation the operators $D_A$, $Q_{*}(A)$ and $Q_{*}^T(-A)$ of Section~\ref{sec:exact-fermion-operators} and the H\"older difference quotient $\delta_{\alpha,A}$ of Definition~\ref{def:pairs-sites-holder}, formed at the level $k^{*}$, act on a spinor function $f$ on $T_{*}$ and on a spinor function $v$ on $T^0_{*}$ as follows:
\begin{align*}
\left(D_Af\right)(x)&=-\eta_{*}^{-1}\cdot\sum_{\nu=0}^2\sum_{\boldsymbol{\sigma}=\pm1}P^{\boldsymbol{\sigma}}_{\nu}\cdot\left[u_A(x,x^{\boldsymbol{\sigma}})\cdot f(x^{\boldsymbol{\sigma}})-f(x)\right],\\
\left(Q_{*}(A)f\right)(y)&=\eta_{*}^3\cdot\sum_{x\in\Delta_y}\exp\left(i\cdot\sigma_A(y,x)\right)\cdot f(x),\\
\left(Q_{*}^T(-A)v\right)(x)&=\exp\left(-i\cdot\sigma_A(y,x)\right)\cdot v(y)&&\text{for }x\in\Delta_y,\\
\left(\delta_{\alpha,A}f\right)(x,x')&=|x-x'|^{-\alpha}\cdot\left(U_A(x,x')\cdot f(x')-f(x)\right)&&\text{for }(x,x')\in\mathrm{Pz}_{*}.
\end{align*}
Further, let
\begin{align*}
P_{*}(A)&:=Q_{*}^T(-A)\cdot Q_{*}(A),\qquad\mathsf{D}(A):=D_A+\mathfrak{m}_{*}+b_{*}\cdot P_{*}(A),\\
S(A)&:=\mathsf{D}(A)^{-1}\qquad\text{whenever this inverse exists},\\
H_{*}(A)&:=b_{*}\cdot S(A)\cdot Q_{*}^T(-A),\qquad\bar{H}_{*}(A):=\mathsf{C}^{-1}\cdot H_{*}(-A)\cdot\mathsf{C},\\
D_{*}(A)&:=b_{*}-b_{*}^2\cdot Q_{*}(A)\cdot S(A)\cdot Q_{*}^T(-A)=b_{*}-b_{*}\cdot Q_{*}(A)\cdot H_{*}(A),
\end{align*}
and let $\mathrm{Sa}(A):=\delta_{\alpha,A}\cdot S(A)$ on the pairs $(x,x')\in\mathrm{Pz}_{*}$ with $|x-x'|<1$.

The maps $A\mapsto\sigma_A$ and $A\mapsto\upsilon_A$ are linear. For $y\in T^0_{*}$ the function $\sigma_A(y,\cdot)$ depends only on the values of $A$ on the bonds whose two end points lie in $\Delta_y$, and
\[
\left|\sigma_A(y,x)\right|\le\frac{3}{2}\,e_{*}\cdot|A|_{\infty}\qquad\text{for }x\in\Delta_y.
\]
Indeed, for $0\le j\le k^{*}$ let $y_j$ be the centre of the block of side $L^j\cdot\eta_{*}$ to which $x$ belongs. Then $y_0=x$, $y_{k^{*}}=y$, and $\left(\tau_{k^{*}}A\right)(y,x)=\sum_{j<k^{*}}(\tau A)(y_{j+1},y_j)$ by Section~\ref{sec:exact-fermion-operators}, where $(\tau A)(y_{j+1},y_j)$ is the mean of the sums of $A$ along the six paths $\Gamma^{\pi}(y_{j+1},y_j)$. These paths stay inside the block of side $L^{j+1}\cdot\eta_{*}$ centred at $y_{j+1}$, which is contained in $\Delta_y$, and each of them has at most $\frac{3}{2}\,(L-1)\cdot L^j$ bonds. This proves the first assertion, and it gives
\[
\left|\sigma_A(y,x)\right|\le e_{*}\cdot|A|_{\infty}\cdot\sum_{j<k^{*}}\frac{3}{2}\,(L-1)\cdot L^j\cdot\eta_{*},\qquad\sum_{j<k^{*}}\frac{3}{2}\,(L-1)\cdot L^j\cdot\eta_{*}=\frac{3}{2}\left(1-\eta_{*}\right)<\frac{3}{2}.
\]
For $(x,x')\in\mathrm{Pz}_{*}$ each of the six paths $\Gamma^{\pi}(x,x')$ consists of $\frac{|x-x'|_1}{\eta_{*}}\le\frac{3\,|x-x'|}{\eta_{*}}$ bonds, all of which lie in the box spanned by $x$ and $x'$. Consequently
\[
\left|\upsilon_A(x,x')\right|\le3\,|x-x'|\cdot e_{*}\cdot|A|_{\infty}.
\]
For $x\in T_{*}$, $[x]$ is the unit site whose cube $\Delta_{[x]}$ contains $x$. Put
\[
\mathfrak{i}(A):=e_{*}\cdot|\operatorname{Im}A|_{\infty}.
\]
The matrices $\gamma_{\nu}$ are monomial: every row and every column of $\gamma_{\nu}$ contains exactly one non-zero entry, and this entry has modulus one (Proposition~\ref{prop:symmetries-extraction}(b); Section~\ref{sec:gamma-matrices-symmetry}). By the triangle inequality every row of $P^{\boldsymbol{\sigma}}_{\nu}$ therefore has $\ell^1$ norm at most $\frac{1}{2}(1+1)=1$; equivalently, $\left|P^{\boldsymbol{\sigma}}_{\nu}\cdot v\right|\le|v|$ for every $v\in\mathbb{C}^4$.

\paragraph{Majorants.} Three kinds of functions occur: spinor functions on $T^0_{*}$, spinor functions on $T_{*}$, and functions on $\mathrm{Pz}_{*}$. Each argument of such a function is localised at a unit site: a site $y\in T^0_{*}$ at $y$; a site $x\in T_{*}$ at $[x]$, that is, in the cube $\Delta_{[x]}$; a pair $(x,x')\in\mathrm{Pz}_{*}$ at $[x]$, that is, in the cube of $x$. For $y\in T^0_{*}$ let $\mathbf{1}_y$ be the multiplication by the indicator function of the set of the arguments that are localised at $y$. Let $O$ be a linear operator that maps functions of one of the three kinds to functions of one of the three kinds. For $y,y'\in T^0_{*}$ and $\tilde{\gamma}\ge0$ define the majorant $|O|$ of $O$ and the norms $\|O\|_{\tilde{\gamma}}$ by
\begin{align*}
|O|(y,y')&:=\sup\left\{\left|\mathbf{1}_y\cdot O\cdot\mathbf{1}_{y'}\cdot f\right|_{\infty}:\ |f|_{\infty}\le1\right\},\\
\|O\|_{\tilde{\gamma}}&:=\max\left\{\sup_y\sum_{y'}|O|(y,y')\cdot\exp\left(\tilde{\gamma}\cdot d(y,y')\right),\right.\\
&\qquad\qquad\left.\sup_{y'}\sum_y|O|(y,y')\cdot\exp\left(\tilde{\gamma}\cdot d(y,y')\right)\right\},
\end{align*}
the sums and suprema being taken over $T^0_{*}$, and put
\[
E(y,y'):=\exp\left(-\gamma_F\cdot d(y,y')\right),
\]
with the constant $\gamma_F$ of Section~\ref{sec:constants-3}.

The majorants obey two rules. Let $O_1$ and $O_2$ be linear operators as above such that the product $O_1\cdot O_2$ is defined. Then, for all $y,y'\in T^0_{*}$ and $\tilde{\gamma}\ge0$,
\begin{equation}
|O_1\cdot O_2|(y,y')\le\sum_{z\in T^0_{*}}|O_1|(y,z)\cdot|O_2|(z,y'),\qquad\|O_1\cdot O_2\|_{\tilde{\gamma}}\le\|O_1\|_{\tilde{\gamma}}\cdot\|O_2\|_{\tilde{\gamma}}.\tag{Maj1}\label{eq:four-jets-Maj1}
\end{equation}
Let $O$ be a linear operator on the spinor functions on $T^0_{*}$, with the $4\times4$ blocks $O(y,y')=\left(O_{\boldsymbol{\beta}\boldsymbol{\beta}'}(y,y')\right)_{\boldsymbol{\beta},\boldsymbol{\beta}'}$, so that $(O\cdot v)(y)=\sum_{y'}O(y,y')\cdot v(y')$; let $|O|_{\ell^2}$ be its operator norm with respect to the $\ell^2$ norm, and $|O|_1$ its trace norm. Then, for all $y,y'\in T^0_{*}$ and every spinor index $\boldsymbol{\beta}$,
\begin{equation}
\sum_{\boldsymbol{\beta}'}|O_{\boldsymbol{\beta}\boldsymbol{\beta}'}(y,y')|\le|O|(y,y'),\qquad|O|_{\ell^2}\le2\,\|O\|_0,\qquad|O|_1\le4\sum_{z,z'\in T^0_{*}}|O|(z,z').\tag{Maj2}\label{eq:four-jets-Maj2}
\end{equation}

\begin{proof}[Proof of \eqref{eq:four-jets-Maj1} and \eqref{eq:four-jets-Maj2}]
Every argument is localised at exactly one unit site; hence $\sum_{z\in T^0_{*}}\mathbf{1}_z=1$ on each of the three spaces of functions. Let $|f|_{\infty}\le1$, and put $g_z:=\mathbf{1}_z\cdot O_2\cdot\mathbf{1}_{y'}\cdot f$ for $z\in T^0_{*}$. Then
\[
\mathbf{1}_y\cdot O_1\cdot O_2\cdot\mathbf{1}_{y'}\cdot f=\sum_{z\in T^0_{*}}\mathbf{1}_y\cdot O_1\cdot\mathbf{1}_z\cdot g_z,
\]
and $|g_z|_{\infty}\le|O_2|(z,y')$ by the definition of $|O_2|$. As $\mathbf{1}_z\cdot g_z=g_z$ and $O_1$ is linear, the definition of $|O_1|$ gives $\left|\mathbf{1}_y\cdot O_1\cdot\mathbf{1}_z\cdot g_z\right|_{\infty}\le|O_1|(y,z)\cdot|g_z|_{\infty}$. The triangle inequality for $|\cdot|_{\infty}$, followed by the supremum over $f$, now yields the first inequality of \eqref{eq:four-jets-Maj1}. For the second, let $w(y,y'):=\exp\left(\tilde{\gamma}\cdot d(y,y')\right)$. The triangle inequality for $d$ and $\tilde{\gamma}\ge0$ give $w(y,y')\le w(y,z)\cdot w(z,y')$ for all $y,z,y'$. Together with the first inequality this gives, for every $y$,
\[
\sum_{y'}|O_1\cdot O_2|(y,y')\cdot w(y,y')\le\sum_z\left(|O_1|(y,z)\cdot w(y,z)\cdot\sum_{y'}|O_2|(z,y')\cdot w(z,y')\right)\le\|O_1\|_{\tilde{\gamma}}\cdot\|O_2\|_{\tilde{\gamma}},
\]
and, for every $y'$,
\[
\sum_y|O_1\cdot O_2|(y,y')\cdot w(y,y')\le\sum_z\left(|O_2|(z,y')\cdot w(z,y')\cdot\sum_y|O_1|(y,z)\cdot w(y,z)\right)\le\|O_1\|_{\tilde{\gamma}}\cdot\|O_2\|_{\tilde{\gamma}}.
\]
This is the second inequality of \eqref{eq:four-jets-Maj1}.

For \eqref{eq:four-jets-Maj2}, let $f$ be a spinor function on $T^0_{*}$ with $f=\mathbf{1}_{y'}\cdot f$; then $|f|_{\infty}=\max_{\boldsymbol{\beta}'}|f_{\boldsymbol{\beta}'}(y')|$. The function $\mathbf{1}_y\cdot O\cdot\mathbf{1}_{y'}\cdot f$ vanishes outside $y$, and its value at $y$ is $O(y,y')\cdot f(y')$; so its norm is at least $\left|\sum_{\boldsymbol{\beta}'}O_{\boldsymbol{\beta}\boldsymbol{\beta}'}(y,y')\cdot f_{\boldsymbol{\beta}'}(y')\right|$ for every $\boldsymbol{\beta}$. Fix $\boldsymbol{\beta}$ and choose the numbers $f_{\boldsymbol{\beta}'}(y')$ of modulus one such that $O_{\boldsymbol{\beta}\boldsymbol{\beta}'}(y,y')\cdot f_{\boldsymbol{\beta}'}(y')=|O_{\boldsymbol{\beta}\boldsymbol{\beta}'}(y,y')|$ for every $\boldsymbol{\beta}'$; then $|f|_{\infty}=1$, and the last modulus equals $\sum_{\boldsymbol{\beta}'}|O_{\boldsymbol{\beta}\boldsymbol{\beta}'}(y,y')|$, the $\ell^1$ norm of the row $\boldsymbol{\beta}$ of $O(y,y')$. By the definition of $|O|(y,y')$ this gives the first inequality. Next, regard $O$ as a matrix whose rows and columns are indexed by the pairs $(y,\boldsymbol{\beta})$. By the first inequality, the row $(y,\boldsymbol{\beta})$ has $\ell^1$ norm
\[
\sum_{y'}\sum_{\boldsymbol{\beta}'}|O_{\boldsymbol{\beta}\boldsymbol{\beta}'}(y,y')|\le\sum_{y'}|O|(y,y')\le\|O\|_0,
\]
and, as there are four spinor indices, the column $(y',\boldsymbol{\beta}')$ has $\ell^1$ norm
\[
\sum_y\sum_{\boldsymbol{\beta}}|O_{\boldsymbol{\beta}\boldsymbol{\beta}'}(y,y')|\le4\sum_y|O|(y,y')\le4\,\|O\|_0.
\]
Schur's test gives $|O|_{\ell^2}\le\left(\|O\|_0\cdot4\,\|O\|_0\right)^{\frac{1}{2}}=2\,\|O\|_0$. Finally, a matrix is the sum of its entries times the matrix units, each of which has trace norm one; so its trace norm does not exceed the sum of the moduli of its entries. By the first inequality, $\sum_{\boldsymbol{\beta},\boldsymbol{\beta}'}|O_{\boldsymbol{\beta}\boldsymbol{\beta}'}(z,z')|\le4\,|O|(z,z')$ for all $z,z'\in T^0_{*}$, and summation over $z,z'$ gives the third inequality.
\end{proof}

For an analytic function $A\mapsto O(A)$ of the connection with values in linear operators, and for a complex bond function $B$ on $T_{*}$, put
\[
\left(\partial_BO\right)(A):=\left.\frac{\mathrm{d}}{\mathrm{d}u}O(A+u\cdot B)\right|_{u=0}.
\]

\subsubsection{The stop level}\label{sec:stop-level}

\begin{proposition}[Trajectory, fermionic operators, activities and determinant for $P_N(e)$ at the stop level of $R^{*}$]\label{prop:trajectory-fermionic-operators}
Assume hypothesis~\ref{def:hypothesis-reference-torus}. Here, for a real bond function $B$ on $T^0_{*}$, $H_{*}B:=H_{k^{*}}B$ is the minimiser of Section~\ref{sec:averaging-fluctuations-gauge}, and $\mathrm{Dom}_{*}:=\mathrm{Dom}_{k^{*}}$ is the set of Definition~\ref{def:domains}; for a bond function $A$ on $T_{*}$ the operators $\mathsf{D}(A)$, $S(A)$, $\mathrm{Sa}(A)$, $H_{*}(A)$, $\bar{H}_{*}(A)$ and $D_{*}(A)$, the majorants $|O|(y,y')$ of operators $O$ and the function $E(y,y')$ are those of Section~\ref{sec:fermionic-operators-stop}. Then the following statements hold.
\begin{enumerate}
\item[(SL1)]\itemmark{SL1}\label{it:trajectory-fermionic-operators-SL1} \emph{(The pair.)} The pair $P_N(e)$ is real and controlled on $R^{*}$ in the sense of Definition~\ref{def:fluctuation-field-controlled}. Let $\left(\boldsymbol{\varepsilon}^t_k,m_k,E_k\right)$ be its trajectory and $\left(\mathscr{E}_k,\mathfrak{P}_k\right)$ its trajectory with slot on $R^{*}$, $0\le k\le k^{*}$ (Section~\ref{sec:flow}), and put
\[
m_{*}:=m_{k^{*}},\qquad\boldsymbol{\varepsilon}^t_{*}:=\boldsymbol{\varepsilon}^t_{k^{*}},\qquad\mathscr{E}_{*}:=\mathscr{E}_{k^{*}},\qquad\mathfrak{P}_{*}:=\mathfrak{P}_{k^{*}}.
\]
Then $|m_{*}|\le\frac{1}{8}\,e_{*}^{\frac{3}{4}-8\varepsilon}$.
\item[(SL2)]\itemmark{SL2}\label{it:trajectory-fermionic-operators-SL2} \emph{(Domains and fermionic operators.)} For $0\le j\le k^{*}$ the sets $G_j$ and $\tilde{R}_j+\mathfrak{X}_j$, and in particular the set $\mathfrak{X}_j$, are contained in $\tilde{G}_j$. For $0\le j\le k^{*}$ and $0<c\le1$ the sets $c\cdot\tilde{R}_j$ and $\tilde{G}_j$ are invariant under the addition of real constant bond functions, under the map $A\mapsto-A$, and under the real gauge transformations $A\mapsto A-d\lambda$, $\lambda$ a real function on the sites of $T_j$. For every $A\in\tilde{G}_{*}$ the operator $\mathsf{D}(A)$ is invertible; the map $A\mapsto S(A)$ is analytic on $\tilde{G}_{*}$; and for $A\in\tilde{G}_{*}$ and $y,y'\in T^0_{*}$
\[
|S(A)|(y,y'),\quad|\mathrm{Sa}(A)|(y,y'),\quad|H_{*}(A)|(y,y'),\quad\left|\delta_{\alpha,A}H_{*}(A)\right|(y,y')\ \le\ C_F\cdot E(y,y'),
\]
where the second and the fourth majorant are formed with the pairs $(x,x')\in\mathrm{Pz}_{*}$ with $|x-x'|<1$ only. Moreover, for $A\in\tilde{G}_{*}$, every pair $(x,x')\in\mathrm{Pz}_{*}$ with $|x-x'|\ge1$ and every spinor function $f$ on $T_{*}$,
\[
\left|\left(\delta_{\alpha,A}f\right)(x,x')\right|\le\exp(1)\cdot|f(x')|+|f(x)|.
\]
\item[(SL3)]\itemmark{SL3}\label{it:trajectory-fermionic-operators-SL3} \emph{(The activity with slot.)} One has $\mathscr{E}_{*}\in K^x_{*}$, the space of Definition~\ref{def:renormalisation-group-1} at the level $k^{*}$. Let
\[
\mathscr{E}_{nm}\left(X;\mathscr{A},\mathscr{C};\xi_1,\dots,\xi_n;\boldsymbol{\zeta}_1,\dots,\boldsymbol{\zeta}_m\right)\qquad\left(X\in D_{*},\ n,m\ge0\right)
\]
be its kernels in the sense of Definition~\ref{def:renormalisation-group-1}. At fixed arguments, each kernel is an analytic function of $(\mathscr{A},\mathscr{C})$ on $\tilde{R}_{*}\times\mathfrak{X}_{*}$. A kernel is zero at every family of arguments in which the site of some $\xi_i$ lies outside $X$, or the pair of sites of some $\boldsymbol{\zeta}_j$ is not admissible for $X$, the admissible pairs being those of the kinds~(I) and~(II) of Definition~\ref{def:pairs-sites-holder}; it is zero as well at every family in which the number of barred arguments differs from the number of unbarred ones. Two elements $(\mathscr{A},\mathscr{C})$ and $(\mathscr{A}',\mathscr{C}')$ of $\tilde{R}_{*}\times\mathfrak{X}_{*}$ with $\mathscr{A}(b)=\mathscr{A}'(b)$ and $\mathscr{C}(b)=\mathscr{C}'(b)$ for every bond $b$ with an end point in $X$ give the same kernels. For $X\in D_{*}$ and $(\mathscr{A},\mathscr{C})\in\tilde{R}_{*}\times\mathfrak{X}_{*}$
\begin{equation}\label{eq:four-jets-2}
\sum_{n,m\ge0}\frac{e_{*}^{-\frac{n}{4}}\cdot e_{*}^{\left(-\frac{1}{4}+\varepsilon\right)m}}{n!\cdot m!}\cdot\left\|\mathscr{E}_{nm}(X;\mathscr{A},\mathscr{C})\right\|\le4\,e_{*}^{\varepsilon}\cdot e_{*}^{\frac{1}{4}-7\varepsilon}\cdot\exp\left(-\kappa\cdot d_M(X)\right),
\end{equation}
where $\|\cdot\|$ denotes the norm of kernels of Section~\ref{sec:grassmann-algebras-norms}.
\item[(SL4)]\itemmark{SL4}\label{it:trajectory-fermionic-operators-SL4} \emph{(Determinants and polymers.)} For $0\le j<k^{*}$ and $X\in D_j$ the functions $E^d_j(X,\cdot)$ of Definition~\ref{def:determinant-functions}, formed with the primitives of Definition~\ref{def:periodisation-projection}, are analytic on $\tilde{G}_j$, and $E^{\mathrm{det}}_j(X,A)=E^d_j(X,A)-E^d_j(X,0)$ for $A\in\tilde{G}_j$. The functions $\mathfrak{P}_j$ of the trajectory with slot satisfy $\mathfrak{P}_0=0$ and, for $0\le j<k^{*}$, $Y\in D_{j+1}$ and $\mathscr{C}\in\mathfrak{X}_{j+1}$,
\[
\mathfrak{P}_{j+1}(Y,\mathscr{C})=\sum_{X\in D_j:\ \bar{X}^L=LY}\left(\mathfrak{P}_j\left(X,\mathscr{C}_L\right)+E^{\mathrm{det}}_j\left(X,\mathscr{C}_L\right)\right),
\]
with the set $\bar{X}^L$ of Section~\ref{sec:lattice-geometry} and the scaled field $\mathscr{C}_L$ of Section~\ref{sec:local-parts-adjustment}. Let $0\le j<k^{*}$, $X\in D_j$ and $A,\tilde{A}\in\tilde{G}_j$. Then:
\begin{enumerate}
\item[(i)] $\left|E^d_j(X,A)\right|\le C_d\cdot M^3\cdot\exp\left(-(\kappa+1)\cdot d_M(X)\right)$;
\item[(ii)] $E^d_j(X,A)=E^d_j(X,\tilde{A})$ if $A$ and $\tilde{A}$ coincide on the bonds that meet $X$;
\item[(iii)] $E^d_j(X,A-d\boldsymbol{\lambda})=E^d_j(X,A)$ for every real function $\boldsymbol{\lambda}$ on the sites of $T_j$;
\item[(iv)] $E^d_j(X,-A)=E^d_j(X,A)$.
\end{enumerate}
Further:
\begin{enumerate}
\item[(v)] \emph{(Polymers.)} Let $0\le j\le k^{*}$. For every $M$-cube $\square$ of $T_j$
\[
\sum_{X\in D_j:\ X\supset\square}\exp\left(-\kappa_0\cdot d_M(X)\right)\le K_0,
\]
and any two points of a polymer $X\in D_j$ have distance at most $M\cdot\left(d_M(X)+2\right)$. If $j<k^{*}$, then every $X\in D_j$ has exactly one $Y\in D_{j+1}$ with $\bar{X}^L=LY$, and $d_M(X)\ge L\cdot d_M(Y)$ for this $Y$; and for every $Y\in D_{j+1}$
\[
\sum_{X\in D_j:\ \bar{X}^L=LY}\exp\left(-\kappa\cdot d_M(X)\right)\le c_{\mathrm{rb}}\cdot L^3\cdot\exp\left(-L\cdot\left(\kappa-\kappa_0-1\right)\cdot d_M(Y)\right).
\]
\item[(vi)] \emph{(The free determinant.)} For every $a\in C^1(\Lambda_{n_{*}})\otimes\mathbb{C}$ with $a^{[j]}\in\mathfrak{X}_j$ for $0\le j\le k^{*}$,
\begin{align}
\frac{\det\left(D_a+\bar{m}\cdot\eta\right)}{\det\left(D_0+\bar{m}\cdot\eta\right)}&=\exp\left\{\sum_{X\in D_{*}}\mathfrak{P}_{*}\left(X,a^{[k^{*}]}\right)\right\}\cdot\frac{\det D_{*}\left(a^{[k^{*}]}\right)}{\det D_{*}(0)},\notag\\
\sum_{X\in D_{*}}\mathfrak{P}_{*}\left(X,a^{[k^{*}]}\right)&=\sum_{j<k^{*}}\ \sum_{X\in D_j}E^{\mathrm{det}}_j\left(X,a^{[j]}\right).\label{eq:four-jets-3}
\end{align}
\end{enumerate}
\item[(SL5)]\itemmark{SL5}\label{it:trajectory-fermionic-operators-SL5} \emph{(Block fields and the small-field set.)} Let $A_{*}$, $\zeta_j$, $\mathrm{d}\mathsf{m}_{*}$ and $\chi_{*}$ be as in Definition~\ref{def:coordinates-slice-small}, and let $V_{\mathrm{t}}(T^0_{*})$ be the space of the real bond functions on $T^0_{*}$ that are coclosed and have mean zero (Proposition~\ref{prop:orthogonal-coordinates}). The map $F\mapsto\left(A_{*};\zeta_0,\dots,\zeta_{k^{*}-1}\right)$ is a linear bijection of $F_N$ onto $V_{\mathrm{t}}(T^0_{*})\times\prod_{j<k^{*}}V_j$. One has $\zeta_j(\Upsilon)=\left\langle\phi_{j,\Upsilon},F\right\rangle$ with an orthonormal family $\left\{\phi_{j,\Upsilon}\right\}$ in $F_N$; the space $V_j$ has $d_V\cdot n_j\le2\,n_j$ basis elements $\Upsilon$, where $n_j=l_{*}^3\cdot L^{3(k^{*}-j)}$; and for each of the $3\,l_{*}^3$ plaquettes $p$ of $T^0_{*}$ one has $(dA_{*})(p)=\left\langle\phi_p,F\right\rangle$ with an element $\phi_p$ of $F_N$ such that $|\phi_p|_2\le1$. On $V_{\mathrm{t}}(T^0_{*})$, the measure $\mathrm{d}\mathsf{m}_{*}$ is Gaussian, centred and non-degenerate. If $A\in V_{\mathrm{t}}(T^0_{*})$ and $\chi_{*}(A)=1$, then a real function $\boldsymbol{\lambda}$ on $T^0_{*}$ can be chosen so that the bond functions $A^g:=A-d\boldsymbol{\lambda}$ and $\mathscr{A}^0:=H_{*}A^g$ satisfy
\begin{align}
&|A^g|_{\infty}\le2\,l_{*}^2\cdot p_{*},\qquad\mathscr{A}^0\ \text{is real},\qquad|\mathscr{A}^0|_{\infty}<\frac{1}{2}\,e_{*}^{-\varepsilon},\notag\\
&\text{and for all }\theta\in\mathbb{R}^3\text{ and }u\in[0,1]:\notag\\
&A+c_{\theta}\in\mathrm{Dom}_{*},\qquad u\cdot A^g+c_{\theta}\in\mathrm{Dom}_{*},\notag\\
&H_{*}\left(u\cdot A^g+c_{\theta}\right)=u\cdot\mathscr{A}^0+c_{\theta}\in\frac{1}{4}\,\tilde{R}_{*}.\label{eq:four-jets-4}
\end{align}
\item[(SL6)]\itemmark{SL6}\label{it:trajectory-fermionic-operators-SL6} \emph{(Representation.)} Let $\Psi$, $\bar{\Psi}$ be the generators of Section~\ref{sec:grassmann-algebras-norms} on $T^0_{*}$, and let $\int\mathrm{D}\Psi$ and the pairings $\langle\cdot,\cdot\rangle$ be those of Section~\ref{sec:fermion-integrals-blocking}. For $A\in\mathrm{Dom}_{*}$ and $\mathscr{C}\in\mathfrak{X}_{*}$ put $\mathscr{A}:=H_{*}A$, $\mathscr{T}:=\mathscr{A}+\mathscr{C}$,
\[
\psi:=H_{*}(\mathscr{T})\Psi,\qquad\bar{\psi}:=\bar{H}_{*}(\mathscr{T})\bar{\Psi},\qquad\chi:=\delta_{\alpha,\mathscr{T}}\psi,\qquad\bar{\chi}:=\delta_{\alpha,-\mathscr{T}}\bar{\psi},\qquad\psi^{\#}:=(\psi,\chi),
\]
and let $Y^x_{*}(A;\mathscr{C}):=y^x_{k^{*}}\left[\boldsymbol{\varepsilon}^t_{*},m_{*},\mathscr{E}_{*}\right](A;\mathscr{C})$ be the function of Section~\ref{sec:fermion-integrals-blocking} at the level $k^{*}$, formed with $\mathfrak{P}_{*}$. Then
\begin{multline}\label{eq:four-jets-5}
Y^x_{*}(A;\mathscr{C})=\det D_{*}(0)^{-1}\cdot\exp\left\{-\boldsymbol{\varepsilon}^t_{*}\cdot l_{*}^3+\sum_{X\in D_{*}}\mathfrak{P}_{*}(X,\mathscr{C})\right\}\\
\cdot\int\mathrm{D}\Psi\,\exp\left\{-\left\langle\bar{\Psi},D_{*}(\mathscr{T})\Psi\right\rangle-m_{*}\cdot\left\langle\bar{\psi},\psi\right\rangle+\sum_{X\in D_{*}}\mathscr{E}_{*}\left(X;\mathscr{A},\mathscr{C},\psi^{\#}\right)\right\}.
\end{multline}
Let $a\in C^1(\Lambda_{n_{*}})\otimes\mathbb{C}$ with $a^{[j]}\in\mathfrak{X}_j$ for $0\le j\le k^{*}$. Then
\[
Y^x_{*}\left(A+d\boldsymbol{\lambda};a^{[k^{*}]}\right)=Y^x_{*}\left(A;a^{[k^{*}]}\right)
\]
for every real bond function $A$ and every real function $\boldsymbol{\lambda}$ on $T^0_{*}$ such that $A$ and $A+d\boldsymbol{\lambda}$ lie in $\mathrm{Dom}_{*}$; and for every $\theta\in\mathbb{R}^3$ and every bounded Borel function $\Phi$ on $V_{\mathrm{t}}(T^0_{*})$, evaluated on the left side at the coordinate $A_{*}$ of $F\in F_N$,
\begin{equation}\label{eq:four-jets-6}
\nu^{\theta,a}\left(\mathbf{1}_{\mathfrak{G}}\cdot\Phi\right)=\int\mathrm{d}\mathsf{m}_{*}(A)\,\Phi(A)\cdot\chi_{*}(A)\cdot Y^x_{*}\left(A+c_{\theta};a^{[k^{*}]}\right),
\end{equation}
the integrand being read as $0$ where $\chi_{*}(A)=0$.
\item[(SL7)]\itemmark{SL7}\label{it:trajectory-fermionic-operators-SL7} \emph{(The external connection.)} For every admitted placement and every $t\in\mathbb{D}^{\circ}(t_x^{+})$ one has $a(t)^{[j]}\in\mathfrak{X}_j$ for $0\le j\le k^{*}$, and $e_{*}\cdot|\mathscr{C}(t)|_{\infty}<e_{*}^{3\varepsilon}$.
\end{enumerate}
\end{proposition}

\begin{proof}
Hypothesis~\ref{def:hypothesis-reference-torus} is hypothesis~\ref{def:hypothesis-flow} for the torus $R=R^{*}$, with the pair $P_N(e)$; so every result that assumes hypothesis~\ref{def:hypothesis-flow} applies with $R=R^{*}$.

\emph{Proof of~\eqref{it:trajectory-fermionic-operators-SL1}.} By Theorem~\ref{thm:agreement-trajectories}, that is, Theorem~\ref{thm:one-pair-all}, the pair $P_N(e)$ is defined, real and controlled on $R^{*}$, and its trajectory on $R^{*}$ satisfies $|m_k|\le\frac{1}{8}\,w^{\mathrm{m}}_k$ for $0\le k\le k^{*}$. At the level $k^{*}$ one has $w^{\mathrm{m}}_{*}=e_{*}^{\frac{3}{4}-8\varepsilon}$ (Section~\ref{sec:running-quantities-tori}).

\emph{Proof of~\eqref{it:trajectory-fermionic-operators-SL2}.} Let $0\le j<k^{*}$. By Lemma~\ref{lem:scales-stop-level}(c) the data $\left(j,l_j,\bar{m}\cdot a_j,e_j\right)$, in which $\bar{m}\cdot a_j$ is the running mass of level $j$, satisfy, together with the primitives of Definition~\ref{def:periodisation-projection}, hypothesis~\ref{def:hypothesis-one-step}; hence Lemmas~\ref{lem:inclusions-domains} and~\ref{lem:symmetries-domains} hold at the level $j$. By Lemma~\ref{lem:scales-stop-level}(a), (c) one has $k^{*}\ge1$, and hypothesis~\ref{def:hypothesis-one-step} holds at the level $k^{*}-1\ge0$. By the last paragraph of Section~\ref{sec:hypotheses-conditions}, it follows that Lemmas~\ref{lem:openness}, \ref{lem:inclusions-domains} and~\ref{lem:phases-differences-along} hold at the level $k^{*}$, and so does Lemma~\ref{lem:symmetries-domains} for all its sets except $\mathrm{Dom}_{*}$. Thus, for $0\le j\le k^{*}$, Lemma~\ref{lem:inclusions-domains}(a), that is, Lemma~\ref{lem:inclusions-bounds}, gives $\tilde{R}_j+\mathfrak{X}_j\subset\tilde{G}_j$; since $0\in\tilde{R}_j$, also $\mathfrak{X}_j\subset\tilde{G}_j$. The inclusion $G_j\subset\tilde{G}_j$ holds by Definition~\ref{def:domains}: an element $A$ of $G_j$ is of the form $A_0+A_1$ required there, with $A_0=0$, $A_1=A$, and with $A'=0$, $\lambda=0$ on every cube $\square^{\sim5}$. Lemma~\ref{lem:symmetries-domains}, that is, Lemma~\ref{lem:invariances}, gives the invariance of the sets $c\cdot\tilde{R}_j$ and $\tilde{G}_j$ under the addition of real constant bond functions, under $A\mapsto-A$ and under the real gauge transformations.

The data $\left(k^{*},l_{*},\mathfrak{m}_{*},e_{*}\right)$ belong to the scale-free setting of Section~\ref{sec:properties-scheme}, and $k^{*}\ge1$ (Lemma~\ref{lem:scales-stop-level}(a), (c)). In this setting Propositions~\ref{prop:decay-fine-propagator} and~\ref{prop:bounds} hold, with the constants $C_F$, $\gamma_F$ of Section~\ref{sec:constants-3}, by Theorem~\ref{thm:expansion-walks} and Proposition~\ref{prop:operators}. At the level $k=k^{*}$, where $\mathfrak{m}_k=\mathfrak{m}_{*}$ and $b_k=b_{*}$, the operators $D_A+\mathfrak{m}_k+b_k\cdot P_k(A)$, $S_k(A)$ and $H_k(A)$ of Section~\ref{sec:exact-fermion-operators} are the operators $\mathsf{D}(A)$, $S(A)$ and $H_{*}(A)$ of Section~\ref{sec:fermionic-operators-stop}, and the norms of the blocks $\mathbf{1}_{\Delta_y}\cdot O\cdot\mathbf{1}_{\Delta_{y'}}$ and $\mathbf{1}_{\Delta_y}\cdot O\cdot\mathbf{1}_{y'}$ that are bounded in these two propositions are the majorants $|O|(y,y')$. Therefore Proposition~\ref{prop:decay-fine-propagator} shows that $\mathsf{D}(A)$ is invertible for every $A\in\tilde{G}_{*}$, that $S$ is analytic on $\tilde{G}_{*}$, and that $|S(A)|(y,y')$ and $|\mathrm{Sa}(A)|(y,y')$ are at most $C_F\cdot E(y,y')$; and Proposition~\ref{prop:bounds} gives the same bound for $|H_{*}(A)|(y,y')$ and, with the pairs $(x,x')$ such that $|x-x'|<1$, for $\left|\delta_{\alpha,A}H_{*}(A)\right|(y,y')$.

Finally let $A\in\tilde{G}_{*}$, let $(x,x')\in\mathrm{Pz}_{*}$ with $|x-x'|\ge1$, and let $f$ be a spinor function on $T_{*}$. Lemma~\ref{lem:phases-differences-along}, that is, Lemma~\ref{lem:phase-factors}, holds at the level $k^{*}$, as shown above, and it gives
\[
\left|U_A(x,x')\right|=\left|\exp\left(i\cdot e_{*}\cdot\eta_{*}\cdot(\tau A)(x,x')\right)\right|\le\exp(1).
\]
Moreover $|x-x'|^{-\alpha}\le1$, since $|x-x'|\ge1$ and $\alpha>0$. With the formula for $\delta_{\alpha,A}f$ in Section~\ref{sec:fermionic-operators-stop} this gives
\[
\left|\left(\delta_{\alpha,A}f\right)(x,x')\right|=|x-x'|^{-\alpha}\cdot\left|U_A(x,x')\cdot f(x')-f(x)\right|\le\exp(1)\cdot|f(x')|+|f(x)|.
\]

\emph{Proof of~\eqref{it:trajectory-fermionic-operators-SL3}.} The numbers $m_i$ of the trajectory of $P_N(e)$ on $R^{*}$ satisfy $|m_i|\le\frac{1}{8}\,w^{\mathrm{m}}_i\le w^{\mathrm{m}}_i$ for $0\le i\le k^{*}$, by Theorem~\ref{thm:agreement-trajectories}. Hence Proposition~\ref{prop:existence-trajectory}, that is, Theorem~\ref{thm:trajectories-fixed-torus}(a), applied on $R^{*}$ to the real pair $P_N(e)$ at the level $k=k^{*}$, shows that $\mathscr{E}_{*}\in K^x_{*}$ and
\[
\|\mathscr{E}_{*}\|^x_{*}\le4\,e_{*}^{\varepsilon}\cdot w^{\mathrm{E}}_{*},\qquad w^{\mathrm{E}}_{*}=e_{*}^{\frac{1}{4}-7\varepsilon}
\]
(Section~\ref{sec:running-quantities-tori}). The assertions on the kernels are the clauses~\eqref{it:renormalisation-group-1-k1} and~\eqref{it:renormalisation-group-1-k2} of Definition~\ref{def:renormalisation-group-1}, together with the requirement made there that the elements of $K^x_{*}$ have equally many barred and unbarred arguments. By Section~\ref{sec:grassmann-algebras-norms}, the left side of~\eqref{eq:four-jets-2} is the norm $\left\|\mathscr{E}_{*}(X;\mathscr{A},\mathscr{C})\right\|_{h_{k^{*}}}$ formed with the weights $h_{k^{*}}=\left(e_{*}^{-\frac{1}{4}},e_{*}^{-\frac{1}{4}+\varepsilon}\right)$. By the definition~\eqref{it:renormalisation-group-1-k6} of the norm $\|\cdot\|^x_{*}$, this norm is at most $\exp\left(-\kappa\cdot d_M(X)\right)\cdot\|\mathscr{E}_{*}\|^x_{*}$ for $X\in D_{*}$ and $(\mathscr{A},\mathscr{C})\in\tilde{R}_{*}\times\mathfrak{X}_{*}$. This gives~\eqref{eq:four-jets-2}.

\emph{Proof of~\eqref{it:trajectory-fermionic-operators-SL4}.} Let $0\le j<k^{*}$. By Definition~\ref{def:periodisation-projection} the system of primitives of the level $j$ on $T_j(R^{*})$ is the projection of the system of the level $j$ on $T_j(R_{\mathrm{top}})$, with $J=R_{\mathrm{top}}/R^{*}=L$, which is odd. By Lemma~\ref{lem:assembled-systems-primitives} the latter system has the properties~\ref{it:renormalisation-group-1}--\ref{it:renormalisation-group-17} for $\mathrm{par}^0$, and by Theorem~\ref{thm:properties-projection}(b), that is, by property~\ref{it:renormalisation-group-18}, its projection has the properties~\ref{it:renormalisation-group-1}--\ref{it:renormalisation-group-17} for $\mathrm{par}$; the two statements are combined in the proof of Proposition~\ref{prop:data-along-flow}(b). In particular the family~\eqref{it:primitives-p5} of the projected system has the properties~\ref{it:renormalisation-group-1}, \ref{it:renormalisation-group-2}, \ref{it:renormalisation-group-7}, \ref{it:renormalisation-group-12}, \ref{it:renormalisation-group-13} and~\ref{it:renormalisation-group-14} for $\mathrm{par}$, as Proposition~\ref{prop:polymer-expansion-determinant} requires. The remaining hypotheses of that proposition are contained in stages~\ref{it:order-choice-4}--\ref{it:order-choice-6} and in condition~\eqref{eq:renormalisation-group-A10}, and the data of the level $j$ satisfy hypothesis~\ref{def:hypothesis-one-step} (Lemma~\ref{lem:scales-stop-level}(c)). Hence the clauses~\eqref{it:polymer-expansion-determinant-b}, \eqref{it:polymer-expansion-determinant-c} and~\eqref{it:polymer-expansion-determinant-d} of Proposition~\ref{prop:polymer-expansion-determinant} hold for the functions $E^d_j(X,\cdot)$: the first gives their analyticity on $\tilde{G}_j$ and~(i), the second gives~(ii), and the third contains~(iii) and~(iv). The formula for $E^{\mathrm{det}}_j$ is Definition~\ref{def:determinant-functions}, and the relations for $\mathfrak{P}_j$ are the definition of the trajectory with slot in Section~\ref{sec:flow}.

\emph{(v)} By Lemma~\ref{lem:scales-stop-level}(c), $l_j\in\mathrm{Cl}$ for $0\le j\le k^{*}$. Hence the clauses~\eqref{it:polymers-t2} and~\eqref{it:polymers-t7} of Proposition~\ref{prop:polymers} hold at every level $j\le k^{*}$. For $j<k^{*}$ also $l_j/L=l_{j+1}\in\mathrm{Cl}$, so that the clauses~\eqref{it:polymers-t5} and~\eqref{it:polymers-t8} of Proposition~\ref{prop:polymers} hold as well; the latter is used with $\kappa'=\kappa$, which is allowed, since $\kappa\ge7\kappa_0+7\ge\kappa_0$ by stage~\ref{it:order-choice-3}.

\emph{(vi)} This is Theorem~\ref{thm:free-determinant-across}, that is, Proposition~\ref{prop:free-determinant}, on the torus $R^{*}$ with $k=k^{*}$. Its hypotheses hold: hypothesis~\ref{def:hypothesis-flow} holds for $R=R^{*}$; the systems~\eqref{it:primitives-p5} of the levels $j<k^{*}$ satisfy clause~\eqref{it:polymer-expansion-determinant-a} of Proposition~\ref{prop:polymer-expansion-determinant}, by the argument just given for the clauses~\eqref{it:polymer-expansion-determinant-b}--\eqref{it:polymer-expansion-determinant-d}; and Proposition~\ref{prop:decay-fine-propagator} holds at the levels $1\le j\le k^{*}$ by Theorem~\ref{thm:expansion-walks}, because the data of these levels belong to the scale-free setting of Section~\ref{sec:properties-scheme} (by hypothesis~\ref{def:hypothesis-one-step} for $j<k^{*}$, and by Lemma~\ref{lem:scales-stop-level}(c) for $j=k^{*}$).

\emph{Proof of~\eqref{it:trajectory-fermionic-operators-SL5}.} Proposition~\ref{prop:orthogonal-coordinates}, that is, Theorem~\ref{thm:hierarchical-coordinates-curvatures}(b)--(e) with $C_Q=1$, applied on $R^{*}$ at the level $k^{*}$, gives the following: the map $F\mapsto\left(A_{*};\zeta_0,\dots,\zeta_{k^{*}-1}\right)$ is a linear bijection of $F_N$ onto $V_{\mathrm{t}}(T^0_{*})\times\prod_{j<k^{*}}V_j$; there is an orthonormal family $\left\{\phi_{j,\Upsilon}\right\}$ in $F_N$ with $\zeta_j(\Upsilon)=\left\langle\phi_{j,\Upsilon},F\right\rangle$; for every plaquette $p$ of $T^0_{*}$ one has $(dA_{*})(p)=\left\langle\phi_p,F\right\rangle$ with $\phi_p\in F_N$ and $|\phi_p|_2\le1$; and $\mathrm{d}\mathsf{m}_{*}$ is the centred Gaussian measure on $V_{\mathrm{t}}(T^0_{*})$ whose density is proportional to $\exp\left(-\frac{1}{2}\left|dH_{*}A\right|_2^2\right)$, so that it is non-degenerate. By Section~\ref{sec:averaging-fluctuations-gauge} the basis $\{\Upsilon\}$ of $V_j$ has $d_V\cdot n_j$ elements, where $d_V=2-2L^{-3}\le2$. By Section~\ref{sec:running-quantities-tori}, $n_j=l_j^3$, and $l_j=R^{*}/a_j=l_{*}\cdot L^{k^{*}-j}$, since $a_j=L^{j-k^{*}}\cdot a_{*}$; hence $n_j=l_{*}^3\cdot L^{3(k^{*}-j)}$. The torus $T^0_{*}=\left(\mathbb{Z}/l_{*}\mathbb{Z}\right)^3$ has $l_{*}^3$ sites. Assigning to a plaquette the two directions $\mu<\nu$ of its plane and the corner $x$ from which it extends in the positive directions $\mu$ and $\nu$ defines a bijection of the set of the plaquettes onto the set of the pairs $\left(x,\{\mu,\nu\}\right)$ formed by a site $x$ and two directions $\mu<\nu$. As there are three choices of $\{\mu,\nu\}$, the torus $T^0_{*}$ has $3\,l_{*}^3$ plaquettes.

Now let $A\in V_{\mathrm{t}}(T^0_{*})$ with $\chi_{*}(A)=1$, that is, $|dA|_{\infty}\le p_{*}$ (Definition~\ref{def:coordinates-slice-small}). Then $A$ is real and has mean zero, and $l_{*}=M\cdot r_{*}\ge2$. By Proposition~\ref{prop:gauge-fixing-controlled}, that is, Corollary~\ref{cor:bounded-gauge-torus}, there is a function $\boldsymbol{\lambda}$ on $T^0_{*}$ with $|A-d\boldsymbol{\lambda}|_{\infty}\le2\,l_{*}^2\cdot|dA|_{\infty}$. It can be taken real: as $A$ is real, replacing $\boldsymbol{\lambda}$ by its real part does not increase $|A-d\boldsymbol{\lambda}|_{\infty}$. With $A^g=A-d\boldsymbol{\lambda}$ this gives $|A^g|_{\infty}\le2\,l_{*}^2\cdot p_{*}$.

Let $\theta\in\mathbb{R}^3$ and $u\in[0,1]$. The bond functions $u\cdot A^g+c_{\theta}$ and $A+c_{\theta}$ on $T^0_{*}$ are real, and
\[
d\left(u\cdot A^g+c_{\theta}\right)=u\cdot dA,\qquad d\left(A+c_{\theta}\right)=dA,
\]
because $d\,d\boldsymbol{\lambda}=0$ and the coboundary of the constant bond function $c_{\theta}$ vanishes. So both coboundaries have sup norm at most $p_{*}$. One has $0<e_{*}\le e_0$ by Lemma~\ref{lem:scales-stop-level}(a). Condition~\eqref{eq:renormalisation-group-A10} imposes, for every $e_k\in(0,e_0]$, the two inequalities $8\,C_B\cdot p_k\le e_k^{-\frac{3}{4}+3\varepsilon}$ and $8\,e_k^{\frac{3}{4}-\frac{7\varepsilon}{2}}\le1$ assumed in Proposition~\ref{prop:backgrounds-domains}, and conditions~\eqref{eq:renormalisation-group-A1} and~\eqref{eq:renormalisation-group-A2} are imposed for every $e_k\in(0,e_0]$ as well; hence these two inequalities and conditions~\eqref{eq:renormalisation-group-A1} and~\eqref{eq:renormalisation-group-A2} hold at $e_k=e_{*}$, where $p_k=p_{*}$. By Remark~\ref{rem:operators-photon-3}(ii), the hypotheses of Proposition~\ref{prop:backgrounds-domains} on $\varepsilon$, $L$, $C_w$ and $\gamma_w$ hold, and the second of the two inequalities gives $e_{*}<1$; moreover $l_{*}\in\mathrm{Cl}$ by Lemma~\ref{lem:scales-stop-level}(c). By Proposition~\ref{prop:fields-domains}(i), that is, Proposition~\ref{prop:backgrounds-domains}(i) at the level $k^{*}$, whose hypotheses therefore hold at $e_k=e_{*}$, both bond functions lie in $\mathrm{Dom}_{*}$; by the definition of $\mathrm{Dom}_{*}$, the operator $H_{*}$ maps them into $\frac{1}{4}\,\tilde{R}_{*}$. By Proposition~\ref{prop:background-field}, that is, Proposition~\ref{prop:minimiser}, the operator $H_{*}$ is linear and real, and $H_{*}c=c$ for every constant bond function $c$. Hence $\mathscr{A}^0=H_{*}A^g$ is real, and $H_{*}\left(u\cdot A^g+c_{\theta}\right)=u\cdot\mathscr{A}^0+c_{\theta}$.

It remains to bound $\mathscr{A}^0$. By Proposition~\ref{prop:kernel-bounds}, that is, by Theorem~\ref{thm:decay-kernels} with the constants $C_P$, $\gamma_P$ of~\eqref{eq:operators-photon-18}, the kernel of $H_{*}$ satisfies
\[
\left|H_{*}(b,b')\right|\le C_P\cdot\exp\left(-\gamma_P\cdot d(b,b')\right)
\]
for all bonds $b$ of $T_{*}$ and all bonds $b'$ of $T^0_{*}$. Summation over $b'$ by~\eqref{eq:operators-photon-2} gives $|H_{*}B|_{\infty}\le C_P\cdot c_{\mathrm{bd}}(\gamma_P)\cdot|B|_{\infty}$ for every real bond function $B$ on $T^0_{*}$, with the number $c_{\mathrm{bd}}(\gamma_P)$ of~\eqref{eq:operators-photon-1}, and $C_P\cdot c_{\mathrm{bd}}(\gamma_P)\le C_0$ by stage~\ref{it:order-choice-7}. Therefore
\[
|\mathscr{A}^0|_{\infty}\le C_0\cdot|A^g|_{\infty}\le2\,C_0\cdot l_{*}^2\cdot p_{*}<\frac{1}{2}\,e_{*}^{-\varepsilon},
\]
the last inequality by condition~\eqref{eq:renormalisation-group-T3}, which holds under hypothesis~\ref{def:hypothesis-reference-torus} (Section~\ref{sec:smallness-conditions}). This proves~\eqref{eq:four-jets-4}.

\emph{Proof of~\eqref{it:trajectory-fermionic-operators-SL6}.} Let $A\in\mathrm{Dom}_{*}$ and $\mathscr{C}\in\mathfrak{X}_{*}$. Then $\mathscr{A}=H_{*}A\in\frac{1}{4}\,\tilde{R}_{*}\subset\tilde{R}_{*}$, so that $\mathscr{T}$ and $-\mathscr{T}$ lie in $\tilde{G}_{*}$ by~\eqref{it:trajectory-fermionic-operators-SL2}; thus the operators $H_{*}(\mathscr{T})$, $\bar{H}_{*}(\mathscr{T})$ and $D_{*}(\mathscr{T})$ are defined, and so are the fields $\psi$, $\bar{\psi}$, $\chi$, $\bar{\chi}$. By the definitions of Section~\ref{sec:fermion-integrals-blocking} at the level $k^{*}$, with the data $\left(\boldsymbol{\varepsilon}^t_{*},m_{*},\mathscr{E}_{*}\right)$ and the functions $\mathfrak{P}_{*}$, the number $Y^x_{*}(A;\mathscr{C})$ is the integral $\int\mathrm{D}\Psi$ of
\[
\det D_{*}(0)^{-1}\cdot\exp\left\{-\left\langle\bar{\Psi},D_{*}(\mathscr{T})\Psi\right\rangle-m_{*}\cdot\left\langle\bar{\psi},\psi\right\rangle+\sum_{X\in D_{*}}\mathscr{E}_{*}\left(X;\mathscr{A},\mathscr{C},\psi^{\#}\right)-\boldsymbol{\varepsilon}^t_{*}\cdot n_{k^{*}}+\sum_{X\in D_{*}}\mathfrak{P}_{*}(X,\mathscr{C})\right\},
\]
where $\psi^{\#}=H^{\#}_{k^{*}}(\mathscr{T})\Psi=(\psi,\chi)$, the barred fields being $\bar{\psi}$ and $\bar{\chi}$. One has $n_{k^{*}}=l_{*}^3$ (Section~\ref{sec:running-quantities-tori}). The factor $\det D_{*}(0)^{-1}\cdot\exp\left\{-\boldsymbol{\varepsilon}^t_{*}\cdot l_{*}^3+\sum_{X\in D_{*}}\mathfrak{P}_{*}(X,\mathscr{C})\right\}$ is a number; taking it out of the integral gives~\eqref{eq:four-jets-5}. By~\eqref{it:trajectory-fermionic-operators-SL1} the pair $P_N(e)$ is real and controlled on $R^{*}$. Hence Theorem~\ref{thm:representation-twist-external-2}, that is, Theorem~\ref{thm:representation-twist-external}, applies on $R^{*}$ to every $a\in C^1(\Lambda_{n_{*}})\otimes\mathbb{C}$ with $a^{[j]}\in\mathfrak{X}_j$ for $0\le j\le k^{*}$: its part~(b) is the gauge invariance stated in~\eqref{it:trajectory-fermionic-operators-SL6}, and its part~(c) is~\eqref{eq:four-jets-6}.

\emph{Proof of~\eqref{it:trajectory-fermionic-operators-SL7}.} Let the placement be admitted. By Lemma~\ref{lem:properties-probes}(e),
\[
d_{\mathrm{sep}}\ge\frac{M\cdot r_{*}}{8L}-1\ge\frac{L^3}{8}-1\ge4,
\]
because $M\ge L^2$ (stage~\ref{it:order-choice-6}), $r_{*}\ge L^2$ (Lemma~\ref{lem:scales-stop-level}(b)) and $L\ge5$ (stage~\ref{it:order-choice-4}). So Lemma~\ref{lem:properties-probes}(f) applies; it gives $|a(t)|_{\mathrm{ext}}<e_{*}^{-1+3\varepsilon}$ for $t\in\mathbb{D}^{\circ}(t_x^{+})$, and $a(t)\in C^1(\Lambda_{n_{*}})\otimes\mathbb{C}$. By Lemma~\ref{lem:scaling-external-connection}, $a(t)^{[j]}\in\mathfrak{X}_j$ for $0\le j\le k^{*}$. Finally $\mathscr{C}(t)=a(t)^{[k^{*}]}$, and $|\mathscr{C}(t)|_{\infty}$ is one of the three norms whose sum is $|a(t)|_{\mathrm{ext}}$; hence
\[
e_{*}\cdot|\mathscr{C}(t)|_{\infty}\le e_{*}\cdot|a(t)|_{\mathrm{ext}}<e_{*}^{3\varepsilon}.
\]

With Lemmas~\ref{lem:wick-rule-gram}, \ref{lem:forest-formula-brydges} and~\ref{lem:tree-expansion}, Theorem~\ref{thm:stop-level-reference} follows.
\end{proof}

\begin{lemma}[Scaling of an external connection]\label{lem:scaling-external-connection}
Let $a\in C^1(\Lambda_{n_{*}})\otimes\mathbb{C}$ with $|a|_{\mathrm{ext}}<e_{*}^{-1+3\varepsilon}$. Then $a^{[k]}\in\mathfrak{X}_k$ for $0\le k\le k^{*}$.
\end{lemma}

\begin{proof}
Let $0\le k\le k^{*}$ and $i:=k^{*}-k\ge0$. By the definition of $a^{[k]}$ in Section~\ref{sec:lattice-measures-laws}, $a^{[k]}(b)=L^{\frac{k}{2}}\cdot a\left(L^k\cdot b\right)$ for a bond $b$ of $T_k$, written in the coordinates of the level $k$. Since $L^{k^{*}}\cdot\left(L^{-i}\cdot b\right)=L^k\cdot b$, it follows that
\[
a^{[k]}(b)=L^{-\frac{i}{2}}\cdot a^{[k^{*}]}\left(L^{-i}\cdot b\right),
\]
and $\eta_k=L^{-k}=L^i\cdot\eta_{*}$ (Section~\ref{sec:running-quantities-tori}). Hence $\left|a^{[k]}\right|_{\infty}=L^{-\frac{i}{2}}\cdot\left|a^{[k^{*}]}\right|_{\infty}$, and for every site $x$ of $T_k$ and all directions $\mu,\nu$
\[
\left(\partial_{\mu}a^{[k]}_{\nu}\right)(x)=\eta_k^{-1}\cdot L^{-\frac{i}{2}}\cdot\left[a^{[k^{*}]}_{\nu}\left(L^{-i}\cdot x+\eta_{*}\cdot e_{\mu}\right)-a^{[k^{*}]}_{\nu}\left(L^{-i}\cdot x\right)\right]=L^{-\frac{3i}{2}}\cdot\left(\partial_{\mu}a^{[k^{*}]}_{\nu}\right)\left(L^{-i}\cdot x\right),
\]
so that $\left|\partial a^{[k]}\right|_{\infty}=L^{-\frac{3i}{2}}\cdot\left|\partial a^{[k^{*}]}\right|_{\infty}$. The norms $\left|a^{[k^{*}]}\right|_{\infty}$ and $\left|\partial a^{[k^{*}]}\right|_{\infty}$ are two of the three norms of Section~\ref{sec:domains-bond-functions} whose sum is $|a|_{\mathrm{ext}}$; so both are smaller than $e_{*}^{-1+3\varepsilon}$. One has $e_{*}\le e_0$ by Lemma~\ref{lem:scales-stop-level}(a) and $e_0<1$ by condition~\eqref{eq:renormalisation-group-A8}; hence $e_{*}<1$ and $e_{*}^{-1+3\varepsilon}\le e_{*}^{-1+2\varepsilon}$. Moreover $e_k=e\cdot a_k^{\frac{1}{2}}\le e_{*}$, because $a_k=L^{k-N}$ increases with $k$. As the exponents $-1+2\varepsilon$ and $-1+3\varepsilon$ are negative, the thresholds of $\mathfrak{X}_k$ in Definition~\ref{def:domains} satisfy
\[
e_k^{-1+2\varepsilon}\ge e_{*}^{-1+2\varepsilon},\qquad e_k^{-1+3\varepsilon}\ge e_{*}^{-1+3\varepsilon}.
\]
Consequently
\[
\left|a^{[k]}\right|_{\infty}\le\left|a^{[k^{*}]}\right|_{\infty}<e_{*}^{-1+2\varepsilon}\le e_k^{-1+2\varepsilon},\qquad\left|\partial a^{[k]}\right|_{\infty}\le\left|\partial a^{[k^{*}]}\right|_{\infty}<e_{*}^{-1+3\varepsilon}\le e_k^{-1+3\varepsilon},
\]
that is, $a^{[k]}\in\mathfrak{X}_k$. Lemma~\ref{lem:external-connections} follows.
\end{proof}

\subsubsection{Grassmann norms, Gaussian integrals and the block covariance}\label{sec:grassmann-norms-gaussian}

The Grassmann algebra $\mathrm{Gr}[\Psi]$ over $\mathbb{C}$ has the generators $\Psi(u)$, indexed by the triples $u=(y,\boldsymbol{\beta},\boldsymbol{\omega})$ formed by a unit site $y\in T^0_{*}$, a spinor index $\boldsymbol{\beta}$ and an index $\boldsymbol{\omega}\in\{0,1\}$ (Section~\ref{sec:grassmann-algebras-norms}). One writes
\[
\Psi_{\boldsymbol{\beta}}(y):=\Psi(y,\boldsymbol{\beta},0),\qquad\bar{\Psi}_{\boldsymbol{\beta}}(y):=\Psi(y,\boldsymbol{\beta},1).
\]
Let $a>0$ be a number. For $F\in\mathrm{Gr}[\Psi]$ the norm $|F|_a$ is that of Section~\ref{sec:grassmann-algebras-norms}. Fix a total order $<$ on the finite set of the triples $u$, as in Section~\ref{sec:coefficients-kernels-grassmann}, and expand $F$ in the basis of the ordered monomials $\Psi(u_1)\cdots\Psi(u_p)$, $u_1<\dots<u_p$, $p\ge0$; then $|F|_a$ is the sum, over these monomials, of $a^p$ times the modulus of the coefficient of $\Psi(u_1)\cdots\Psi(u_p)$ in $F$. This sum does not depend on the chosen order: under a change of the order, the ordered monomial formed from a given set of triples is multiplied by $\pm1$, so that the coefficients of $F$ change by signs only.

Let $O$ be a complex matrix whose rows and columns are indexed by the pairs $(y,\boldsymbol{\beta})$. Its entries are written $O_{\boldsymbol{\beta}\boldsymbol{\beta}'}(y,y')$, and $O(y,y')=\left(O_{\boldsymbol{\beta}\boldsymbol{\beta}'}(y,y')\right)_{\boldsymbol{\beta},\boldsymbol{\beta}'}$ is its $4\times4$ block at $(y,y')$; in this way $O$ is also a linear operator on the spinor functions on $T^0_{*}$. Put
\[
\left\langle\bar{\Psi},O\Psi\right\rangle:=\sum_{y,y'\in T^0_{*}}\ \sum_{\boldsymbol{\beta},\boldsymbol{\beta}'}O_{\boldsymbol{\beta}\boldsymbol{\beta}'}(y,y')\cdot\bar{\Psi}_{\boldsymbol{\beta}}(y)\cdot\Psi_{\boldsymbol{\beta}'}(y'),
\]
and let $\left\langle\bar{\Psi},\Psi\right\rangle$ be this element for the unit matrix. The integral $\int\mathrm{D}\Psi$ is the linear form on $\mathrm{Gr}[\Psi]$ of Section~\ref{sec:fermion-integrals-blocking}; it is normalised so that $\int\mathrm{D}\Psi\,\exp\left(-\left\langle\bar{\Psi},O\Psi\right\rangle\right)=\det O$ for every such matrix $O$. The Gaussian integrals $\mathrm{d}\mu_I$ and $\mathrm{d}\mu_{\operatorname{Cov}}$ are those of Proposition~\ref{prop:weighted-coefficient-norms} (Section~\ref{sec:norms-gaussian-integrals}): for $F\in\mathrm{Gr}[\Psi]$ and an invertible matrix $\operatorname{Cov}$ with rows and columns indexed by the pairs $(y,\boldsymbol{\beta})$,
\[
\int\mathrm{d}\mu_I(\Psi)\,F:=\int\mathrm{D}\Psi\,\exp\left(-\left\langle\bar{\Psi},\Psi\right\rangle\right)\cdot F,\qquad\int\mathrm{d}\mu_{\operatorname{Cov}}(\Psi)\,F:=\det(\operatorname{Cov})\cdot\int\mathrm{D}\Psi\,\exp\left(-\left\langle\bar{\Psi},\operatorname{Cov}^{-1}\Psi\right\rangle\right)\cdot F.
\]
For a matrix $\operatorname{Cov}$ of this kind, invertible or not, $\operatorname{Cov}\Psi$ denotes the family with the members
\[
(\operatorname{Cov}\Psi)_{\boldsymbol{\beta}}(y):=\sum_{y'\in T^0_{*}}\ \sum_{\boldsymbol{\beta}'}\operatorname{Cov}_{\boldsymbol{\beta}\boldsymbol{\beta}'}(y,y')\cdot\Psi_{\boldsymbol{\beta}'}(y').
\]
An element $F$ of $\mathrm{Gr}[\Psi]$ is also written $F(\bar{\Psi},\Psi)$, and $F(\bar{\Psi},\operatorname{Cov}\Psi)$ denotes its image under the homomorphism of the algebra $\mathrm{Gr}[\Psi]$ into itself that maps every $\Psi_{\boldsymbol{\beta}}(y)$ to $(\operatorname{Cov}\Psi)_{\boldsymbol{\beta}}(y)$ and every $\bar{\Psi}_{\boldsymbol{\beta}}(y)$ to itself.

The following facts hold for every number $a>0$ and all $F,G\in\mathrm{Gr}[\Psi]$.

\emph{Products and exponentials.}
\begin{equation}\tag{Gr1}\label{eq:four-jets-Gr1}
|F\cdot G|_a\le|F|_a\cdot|G|_a,\qquad\left|\exp F-1\right|_a\le\exp\left(|F|_a\right)-1\quad\text{if $F$ is even.}
\end{equation}

\emph{Integration against $\mathrm{d}\mu_I$.}
\begin{equation}\tag{Gr2}\label{eq:four-jets-Gr2}
\left|\int\mathrm{d}\mu_I(\Psi)\,F\right|\le|F|_1.
\end{equation}

\emph{Change of covariance.} For every invertible matrix $\operatorname{Cov}$ with rows and columns indexed by the pairs $(y,\boldsymbol{\beta})$,
\begin{gather}
\int\mathrm{D}\Psi\,\exp\left(-\left\langle\bar{\Psi},\operatorname{Cov}^{-1}\Psi\right\rangle\right)\cdot F=\det\left(\operatorname{Cov}^{-1}\right)\cdot\int\mathrm{d}\mu_{\operatorname{Cov}}(\Psi)\,F,\notag\\
\int\mathrm{d}\mu_{\operatorname{Cov}}(\Psi)\,F(\bar{\Psi},\Psi)=\int\mathrm{d}\mu_I(\Psi)\,F(\bar{\Psi},\operatorname{Cov}\Psi).\tag{Gr3}\label{eq:four-jets-Gr3}
\end{gather}

\emph{Substitution.} Let $a\ge1$, and let $\operatorname{Cov}$ be a matrix with rows and columns indexed by the pairs $(y,\boldsymbol{\beta})$ each row of which has $\ell^1$ norm at most $a$, that is, $\sum_{y',\boldsymbol{\beta}'}\left|\operatorname{Cov}_{\boldsymbol{\beta}\boldsymbol{\beta}'}(y,y')\right|\le a$ for every pair $(y,\boldsymbol{\beta})$. Then
\begin{equation}\tag{Gr4}\label{eq:four-jets-Gr4}
\left|F(\bar{\Psi},\operatorname{Cov}\Psi)\right|_1\le|F|_a.
\end{equation}

\emph{Quadratic elements.} For every matrix $O$ with rows and columns indexed by the pairs $(y,\boldsymbol{\beta})$, with the majorant $|O|(y,y')$ of Section~\ref{sec:fermionic-operators-stop},
\begin{equation}\tag{Gr5}\label{eq:four-jets-Gr5}
\left|\left\langle\bar{\Psi},O\Psi\right\rangle\right|_a=a^2\cdot\sum_{y,y'\in T^0_{*}}\ \sum_{\boldsymbol{\beta},\boldsymbol{\beta}'}\left|O_{\boldsymbol{\beta}\boldsymbol{\beta}'}(y,y')\right|\le4a^2\cdot\sum_{y,y'\in T^0_{*}}|O|(y,y').
\end{equation}

\emph{Terminal bounds.} Put
\[
a_2:=C_0^{-1}\cdot e_{*}^{-\frac{1}{4}+\varepsilon},
\]
with the constant $C_0$ of stage~\ref{it:order-choice-7}; the domains $\tilde{G}_{*}$, $\tilde{R}_{*}$ and $\mathfrak{X}_{*}$ are those of Definition~\ref{def:domains} at the level $k^{*}$. Let $A\in\tilde{G}_{*}$, and let
\[
\langle\bar{\psi},\psi\rangle:=\sum_{x\in T_{*}}\eta_{*}^3\cdot\sum_{\boldsymbol{\beta}}\bar{\psi}_{\boldsymbol{\beta}}(x)\cdot\psi_{\boldsymbol{\beta}}(x)
\]
be the element of $\mathrm{Gr}[\Psi]$ formed, as in Section~\ref{sec:fermion-integrals-blocking}, with $\psi:=H_{*}(A)\Psi$ and $\bar{\psi}:=\bar{H}_{*}(A)\bar{\Psi}$, where $H_{*}(A)$ and $\bar{H}_{*}(A)$ are the operators of Section~\ref{sec:fermionic-operators-stop}. Let further $X\in D_{*}$, $(\mathscr{A},\mathscr{C})\in\tilde{R}_{*}\times\mathfrak{X}_{*}$ and $\mathscr{T}:=\mathscr{A}+\mathscr{C}$, and let $\mathscr{E}_{*}\left(X;\mathscr{A},\mathscr{C},H^{\#}_{*}(\mathscr{T})\Psi\right)\in\mathrm{Gr}[\Psi]$ be obtained from the activity $\mathscr{E}_{*}$ of~\eqref{it:trajectory-fermionic-operators-SL1} by the substitution $\psi^{\#}=H^{\#}_{*}(\mathscr{T})\Psi$ of Section~\ref{sec:fermion-integrals-blocking}. With the kernels $\mathscr{E}_{nm}$ of $\mathscr{E}_{*}$ of~\eqref{it:trajectory-fermionic-operators-SL3} and the kernel norm $\|\cdot\|$ of Section~\ref{sec:grassmann-algebras-norms},
\begin{gather}
\left|\langle\bar{\psi},\psi\rangle\right|_{a_2}\le4\,e_{*}^{-\frac{1}{2}+2\varepsilon}\cdot l_{*}^3,\notag\\
\left|\mathscr{E}_{*}\left(X;\mathscr{A},\mathscr{C},H^{\#}_{*}(\mathscr{T})\Psi\right)\right|_{a_2}\le\sum_{n,m\ge0}\frac{e_{*}^{-\frac{n}{4}}\cdot e_{*}^{\left(-\frac{1}{4}+\varepsilon\right)m}}{n!\cdot m!}\cdot\left\|\mathscr{E}_{nm}(X;\mathscr{A},\mathscr{C})\right\|.\tag{Gr6}\label{eq:four-jets-Gr6}
\end{gather}

\emph{The block covariance at real connections.} Let $A$ be a real bond function on $T_{*}$, and let $D_A$, $Q_{*}(A)$, $Q_{*}^T(-A)$ and $D_{*}(A)$ be the operators of Section~\ref{sec:fermionic-operators-stop}. Let $\ell^2(T_{*})$ and $\ell^2(T^0_{*})$ be the spaces of the spinor functions on $T_{*}$ and on $T^0_{*}$ with the scalar products $(\cdot,\cdot)$ of Section~\ref{sec:block-covariance-real}, formed at the level $k^{*}$. Then
\[
\operatorname{Re}\left(f,\left(D_A+\mathfrak{m}_{*}\right)f\right)\ge\mathfrak{m}_{*}\cdot(f,f)\qquad\left(f\in\ell^2(T_{*})\right),
\]
the adjoint of $Q_{*}(A)\colon\ell^2(T_{*})\to\ell^2(T^0_{*})$ is $Q_{*}(A)^{*}=Q_{*}^T(-A)$, the operator $D_{*}(A)$ is invertible, and for all $y,y'\in T^0_{*}$
\begin{gather}
\left|D_{*}(A)^{-1}\right|_{\ell^2}\le b_{*}^{-1}+\mathfrak{m}_{*}^{-1},\notag\\
\left|D_{*}(A)^{-1}(y,y')\right|\le b_{*}^{-1}\cdot\mathbf{1}[y=y']+\frac{2\exp(1)}{\mathfrak{m}_{*}}\cdot\exp\left(-\frac{\mathfrak{m}_{*}}{12\exp(1)}\cdot d(y,y')\right).\tag{Gr7}\label{eq:four-jets-Gr7}
\end{gather}
Here $\left|D_{*}(A)^{-1}\right|_{\ell^2}$ is the operator norm on $\ell^2(T^0_{*})$, the norm of the $4\times4$ block $D_{*}(A)^{-1}(y,y')$ is its operator norm with respect to the Euclidean norm on $\mathbb{C}^4$, as in Section~\ref{sec:block-covariance-real}, and $\mathbf{1}[y=y']$ equals $1$ if $y=y'$ and $0$ otherwise.

\begin{proof}[Proof of \eqref{eq:four-jets-Gr1}--\eqref{eq:four-jets-Gr7}]
The facts \eqref{eq:four-jets-Gr1}--\eqref{eq:four-jets-Gr4} are parts of Proposition~\ref{prop:weighted-coefficient-norms}, that is, of Theorem~\ref{thm:polymer-activities-1}. That theorem is applied with the generators of $\mathrm{Gr}[\Psi]$ grouped either as the single family $\Psi$, indexed by the triples $u$ with the counting weight, or as the two families of the $\Psi_{\boldsymbol{\beta}}(y)$ and of the $\bar{\Psi}_{\boldsymbol{\beta}}(y)$, each indexed by the pairs $(y,\boldsymbol{\beta})$ with the counting weight. The ordered monomials of the two groupings coincide up to signs. Hence for both groupings the norm of Theorem~\ref{thm:polymer-activities-1} in which every family carries the weight $a$ is $|\cdot|_a$.

\eqref{eq:four-jets-Gr1} is~\eqref{it:weighted-coefficient-norms-N1}, together with the bound for $\exp F-1$, $F$ even, which Theorem~\ref{thm:polymer-activities-1} states with it, for the single family $\Psi$ with the weight $a$.

\eqref{eq:four-jets-Gr2} is~\eqref{it:weighted-coefficient-norms-N3} for $p=0$, with $\Psi$ in the place of $W$; for $p=0$ the algebra without generators is $\mathbb{C}$, and its norm is the modulus.

In~\eqref{eq:four-jets-Gr3}, the first identity is the definition of $\mathrm{d}\mu_{\operatorname{Cov}}$, multiplied by $\det\left(\operatorname{Cov}^{-1}\right)=\det(\operatorname{Cov})^{-1}$, and the second identity is~\eqref{it:weighted-coefficient-norms-N7}(iii).

\eqref{eq:four-jets-Gr4} is~\eqref{it:weighted-coefficient-norms-N2} for the two families. The homomorphism $F\mapsto F(\bar{\Psi},\operatorname{Cov}\Psi)$ acts on the family of the $\Psi_{\boldsymbol{\beta}}(y)$ by the operator $O_{11}:=\operatorname{Cov}$, whose row sums are at most $C_{11}:=a$, and on the family of the $\bar{\Psi}_{\boldsymbol{\beta}}(y)$ by the identity $O_{22}$, whose row sums equal $C_{22}:=1$; the operators $O_{12}$ and $O_{21}$ vanish. Take the weights $g=(a,a)$ before and $g'=(1,1)$ after the substitution. The hypothesis of~\eqref{it:weighted-coefficient-norms-N2} then reads $a\cdot1\le a$ and $1\cdot1\le a$, and it holds because $a\ge1$. The conclusion $\left\|F(\bar{\Psi},\operatorname{Cov}\Psi)\right\|_{(1,1)}\le\|F\|_{(a,a)}$ is~\eqref{eq:four-jets-Gr4}.

\eqref{eq:four-jets-Gr5}: For distinct quadruples $(y,\boldsymbol{\beta},y',\boldsymbol{\beta}')$ the products $\bar{\Psi}_{\boldsymbol{\beta}}(y)\cdot\Psi_{\boldsymbol{\beta}'}(y')$ are, up to signs, distinct ordered monomials of degree $2$. By the definition of $|\cdot|_a$, the norm of $\left\langle\bar{\Psi},O\Psi\right\rangle$ is therefore $a^2$ times the sum of the moduli of all entries of $O$; this is the equality. By~\eqref{eq:four-jets-Maj2}, $\sum_{\boldsymbol{\beta}'}\left|O_{\boldsymbol{\beta}\boldsymbol{\beta}'}(y,y')\right|\le|O|(y,y')$ for each of the four spinor indices $\boldsymbol{\beta}$, and the sum over $\boldsymbol{\beta}$ gives the inequality.

\eqref{eq:four-jets-Gr6}, \eqref{eq:four-jets-Gr7}: Propositions~\ref{prop:bounds-unit-lattice} and~\ref{prop:block-covariance-real} are statements in the scale-free setting of Section~\ref{sec:properties-scheme}, at the levels $k\ge1$. By Lemma~\ref{lem:scales-stop-level}(a), (c), one has $k^{*}\ge1$, and the data $\left(k^{*},l_{*},\mathfrak{m}_{*},e_{*}\right)$ belong to this setting; so both propositions hold at $k=k^{*}$, with $l=l_{*}$.

\eqref{eq:four-jets-Gr6} is Theorem~\ref{thm:unit-lattice-bounds}(b), (c) (Proposition~\ref{prop:bounds-unit-lattice}) at $k=k^{*}$. The hypotheses of Theorem~\ref{thm:unit-lattice-bounds} that are taken from Propositions~\ref{prop:decay-fine-propagator}, \ref{prop:kernel-bounds-2} and~\ref{prop:bounds} hold by Theorem~\ref{thm:expansion-walks} and Propositions~\ref{prop:kernels-unit-lattice}, \ref{prop:operators}, and its hypothesis on $C_0$ holds by stage~\ref{it:order-choice-7}. Theorem~\ref{thm:unit-lattice-bounds}(b) contains two inequalities; the first bound of~\eqref{eq:four-jets-Gr6} is the first of them, the one with the exponent $-\frac{1}{2}+2\varepsilon$. For the second bound of~\eqref{eq:four-jets-Gr6}, Theorem~\ref{thm:unit-lattice-bounds}(c) is applied to $\mathscr{E}_{*}$, which lies in $K^x_{*}\subset K^{x,C}_{*}$ (Definition~\ref{def:renormalisation-group-1}) by~\eqref{it:trajectory-fermionic-operators-SL3}. The right side of the inequality of Theorem~\ref{thm:unit-lattice-bounds}(c) is $\left\|\mathscr{E}_{*}(X;\mathscr{A},\mathscr{C})\right\|_{h_{k^{*}}}$, and by the definition of this norm in Section~\ref{sec:grassmann-algebras-norms}, with the weights
\[
h_{k^{*}}=\left(e_{*}^{-\frac{1}{4}},e_{*}^{-\frac{1}{4}+\varepsilon}\right),
\]
it equals the sum on the right of~\eqref{eq:four-jets-Gr6}.

\eqref{eq:four-jets-Gr7} is Proposition~\ref{prop:block-covariance-real}, that is, Theorem~\ref{thm:block-covariance}, at $k=k^{*}$ and with the mass $\mathfrak{m}_{*}$. By Lemma~\ref{lem:scales-stop-level}(a), $0<\mathfrak{m}_{-}<\mathfrak{m}_{*}\le1$; so $\mathfrak{m}_{*}$ is a positive number, as Theorem~\ref{thm:block-covariance} requires, and $\mathfrak{m}_{*}\le12\exp(1)$, as its bound for the blocks requires.
\end{proof}

\subsection{Six auxiliary lemmas}\label{sec:six-auxiliary-lemmas}

\subsubsection{The M\"obius decomposition in the cell variables}\label{sec:mobius-decomposition-cell}

The polydiscs $\mathbb{D}(r)$ and $\mathbb{D}^{\circ}(r)$, $0<r\le\infty$, are those of Section~\ref{sec:probes-mirror-copies}. For a subset $T\subset\mathrm{Cel}$ and a point $t\in\mathbb{C}^{\mathrm{Cel}}$ let $t^T\in\mathbb{C}^{\mathrm{Cel}}$ be the point with the components
\[
\left(t^T\right)_c:=t_c\quad(c\in T),\qquad\left(t^T\right)_c:=0\quad(c\in\mathrm{Cel}\setminus T).
\]
The map $t\mapsto t^T$ of $\mathbb{C}^{\mathrm{Cel}}$ into itself is linear, and $\left|\left(t^T\right)_c\right|\le|t_c|$ for every $c\in\mathrm{Cel}$; hence it maps $\mathbb{D}(r)$ into $\mathbb{D}(r)$ and $\mathbb{D}^{\circ}(r)$ into $\mathbb{D}^{\circ}(r)$ for every $r\in(0,\infty]$. Let $U$ be one of these polydiscs, and let $\mathfrak{g}$ be a complex-valued function on $U$. For a subset $S\subset\mathrm{Cel}$ let $\mathfrak{g}_S$ be the function on $U$ given by
\[
\mathfrak{g}_S(t):=\sum_{T\subset S}(-1)^{\#|S\setminus T|}\cdot\mathfrak{g}\left(t^T\right)\qquad(t\in U).
\]
If $U'\subset U$ is another of these polydiscs, then the function formed in the same way from the restriction of $\mathfrak{g}$ to $U'$ is the restriction of $\mathfrak{g}_S$ to $U'$.

Let $\rho$ be a real number with $0<\rho<\infty$. A function $\mathfrak{g}$ on $\mathbb{D}(\rho)$ is said to be analytic in a neighbourhood of $\mathbb{D}(\rho)$ if it is the restriction to $\mathbb{D}(\rho)$ of a function that is analytic on an open subset $V$ of $\mathbb{C}^{\mathrm{Cel}}$ containing $\mathbb{D}(\rho)$. As $\mathbb{D}(\rho)$ is compact, $V$ then contains $\mathbb{D}^{\circ}(\rho_1)$ for some real number $\rho_1>\rho$. In this situation $\mathfrak{g}$ also denotes the restriction of that analytic function to $\mathbb{D}^{\circ}(\rho_1)$; by the identity theorem it is determined by the values of $\mathfrak{g}$ on $\mathbb{D}(\rho)$, because $\mathbb{D}^{\circ}(\rho_1)$ is connected and contains the non-empty open set $\mathbb{D}^{\circ}(\rho)$. Accordingly, $\mathfrak{g}_S$ also denotes the function on $\mathbb{D}^{\circ}(\rho_1)$ formed from it; its restriction to $\mathbb{D}(\rho)$ is the function $\mathfrak{g}_S$ formed from $\mathfrak{g}$ on $\mathbb{D}(\rho)$.

For $c\in\mathrm{Cel}$, $\partial_c$ denotes the partial derivative with respect to $t_c$, and $e_c$ is the unit vector of the cell $c$ (Section~\ref{sec:probes-mirror-copies}).

\begin{lemma}[Properties of the M\"obius components]\label{lem:properties-mobius-components}
Let $\rho$ be a real number with $0<\rho<\infty$, and let $\mathfrak{g}$ be a complex-valued function on $\mathbb{D}(\rho)$.
\begin{enumerate}
\item[(a)] One has $\mathfrak{g}=\sum_{S\subset\mathrm{Cel}}\mathfrak{g}_S$ on $\mathbb{D}(\rho)$, and $\mathfrak{g}_{\emptyset}$ is the constant function with the value $\mathfrak{g}(0)$. For every $S\subset\mathrm{Cel}$ the function $\mathfrak{g}_S$ depends on the variables $t_c$, $c\in S$, only, and $\mathfrak{g}_S(t)=0$ for every $t\in\mathbb{D}(\rho)$ such that $t_c=0$ for some $c\in S$. If $\mathfrak{g}$ does not depend on $t_c$, then $\mathfrak{g}_S=0$ for every $S\subset\mathrm{Cel}$ that contains $c$.
\item[(b)] For every $S\subset\mathrm{Cel}$,
\[
\sup_{\mathbb{D}(\rho)}\left|\mathfrak{g}_S\right|\le2^{\#|S|}\cdot\sup_{\mathbb{D}(\rho)}|\mathfrak{g}|.
\]
\item[(c)] Let $\mathfrak{g}$ be analytic in a neighbourhood of $\mathbb{D}(\rho)$, let $S\subset\mathrm{Cel}$, and let $c\ne c'$ be elements of $S$. Then
\[
\sup_{\mathbb{D}(\rho)}\left|\mathfrak{g}_S\right|\le2^{\#|S|-2}\cdot\rho^2\cdot\sup_{\mathbb{D}(\rho)}\left|\partial_c\partial_{c'}\mathfrak{g}\right|.
\]
\item[(d)] Let $\mathfrak{g}$ be analytic in a neighbourhood of $\mathbb{D}(\rho)$, let $S\subset\mathrm{Cel}$ be non-empty, and let $0<\rho'\le\frac{\rho}{2}$. Then every monomial that has a non-zero coefficient in the Taylor series of $\mathfrak{g}_S$ at $0$ contains every variable $t_c$ with $c\in S$ and no variable $t_c$ with $c\notin S$; and the image of $\mathfrak{g}_S$ in the ring $\mathcal{J}$ of Section~\ref{sec:ring-jets} satisfies
\[
\left\|\mathfrak{g}_S\right\|_{\rho'}\le\sup_{\mathbb{D}(\rho)}\left|\mathfrak{g}_S\right|.
\]
\end{enumerate}
\end{lemma}

\begin{proof}
(a) Let $t\in\mathbb{D}(\rho)$. After an exchange of the two summations,
\[
\sum_{S\subset\mathrm{Cel}}\mathfrak{g}_S(t)=\sum_{T\subset\mathrm{Cel}}\mathfrak{g}\left(t^T\right)\cdot\sum_{S:\ T\subset S\subset\mathrm{Cel}}(-1)^{\#|S\setminus T|}.
\]
The inner sum equals $\sum_{R\subset\mathrm{Cel}\setminus T}(-1)^{\#|R|}=(1-1)^{\#|\mathrm{Cel}\setminus T|}$, which is $1$ for $T=\mathrm{Cel}$ and $0$ for every other $T$ (M\"obius inversion on the subsets of the finite set $\mathrm{Cel}$). Since $t^{\mathrm{Cel}}=t$, the sum is $\mathfrak{g}(t)$. For $S=\emptyset$ the only subset $T$ is $\emptyset$, and $t^{\emptyset}=0$; so $\mathfrak{g}_{\emptyset}(t)=\mathfrak{g}(0)$. For $T\subset S$ the point $t^T$ depends on the $t_c$ with $c\in S$ only, and hence so does $\mathfrak{g}_S$.

Now let $c\in S$, and assume that $t_c=0$ or that $\mathfrak{g}$ does not depend on $t_c$. The subsets of $S$ are the sets $T$ and $T\cup\{c\}$ with $T\subset S\setminus\{c\}$. For such $T$ one has $\mathfrak{g}\left(t^{T\cup\{c\}}\right)=\mathfrak{g}\left(t^T\right)$: if $t_c=0$, the two points coincide, and otherwise they differ in the component $c$ only, on which $\mathfrak{g}$ does not depend. The signs $(-1)^{\#|S\setminus T|}$ and $(-1)^{\#|S\setminus(T\cup\{c\})|}$ are opposite. So the terms of $T$ and of $T\cup\{c\}$ cancel, and $\mathfrak{g}_S(t)=0$.

(b) The sum that defines $\mathfrak{g}_S(t)$ has $2^{\#|S|}$ terms, and each of them has modulus at most $\sup_{\mathbb{D}(\rho)}|\mathfrak{g}|$, because $t^T\in\mathbb{D}(\rho)$.

(c) Every subset of $S$ has exactly one of the forms $T$, $T\cup\{c\}$, $T\cup\{c'\}$, $T\cup\{c,c'\}$ with $T\subset S\setminus\{c,c'\}$, and there are $2^{\#|S|-2}$ such sets $T$. The sign of $T\cup\{c,c'\}$ in the definition of $\mathfrak{g}_S$ equals the sign $(-1)^{\#|S\setminus T|}$ of $T$, and the signs of $T\cup\{c\}$ and of $T\cup\{c'\}$ are opposite to it. Hence $\mathfrak{g}_S(t)$ is the sum, over the sets $T\subset S\setminus\{c,c'\}$, of $(-1)^{\#|S\setminus T|}$ times
\[
\mathfrak{g}\left(t^{T\cup\{c,c'\}}\right)-\mathfrak{g}\left(t^{T\cup\{c\}}\right)-\mathfrak{g}\left(t^{T\cup\{c'\}}\right)+\mathfrak{g}\left(t^T\right).
\]
Let $t\in\mathbb{D}(\rho)$, fix $T$, and put $\varphi(u,u'):=\mathfrak{g}\left(t^T+u\cdot t_c\cdot e_c+u'\cdot t_{c'}\cdot e_{c'}\right)$ for $u,u'\in[0,1]$; these points lie in $\mathbb{D}(\rho)$. Since $c,c'\notin T$, one has $t^{T\cup\{c\}}=t^T+t_c\cdot e_c$, $t^{T\cup\{c'\}}=t^T+t_{c'}\cdot e_{c'}$ and $t^{T\cup\{c,c'\}}=t^T+t_c\cdot e_c+t_{c'}\cdot e_{c'}$, so that the displayed expression is $\varphi(1,1)-\varphi(1,0)-\varphi(0,1)+\varphi(0,0)$. As $\mathfrak{g}$ is analytic in a neighbourhood of $\mathbb{D}(\rho)$, the chain rule gives $\partial_u\partial_{u'}\varphi(u,u')=t_c\cdot t_{c'}\cdot\left(\partial_c\partial_{c'}\mathfrak{g}\right)\left(t^T+u\cdot t_c\cdot e_c+u'\cdot t_{c'}\cdot e_{c'}\right)$, and the fundamental theorem of calculus, applied in $u$ and in $u'$, yields
\[
\varphi(1,1)-\varphi(1,0)-\varphi(0,1)+\varphi(0,0)=t_c\cdot t_{c'}\cdot\int_0^1\mathrm{d}u\int_0^1\mathrm{d}u'\,\left(\partial_c\partial_{c'}\mathfrak{g}\right)\left(t^T+u\cdot t_c\cdot e_c+u'\cdot t_{c'}\cdot e_{c'}\right).
\]
Since $|t_c|,|t_{c'}|\le\rho$, this has modulus at most $\rho^2\cdot\sup_{\mathbb{D}(\rho)}\left|\partial_c\partial_{c'}\mathfrak{g}\right|$. Summing over the $2^{\#|S|-2}$ sets $T$ gives~(c).

(d) Let $\rho_1>\rho$ be a real number such that $\mathfrak{g}$ is analytic on $\mathbb{D}^{\circ}(\rho_1)$, as above. For every $T\subset\mathrm{Cel}$ the map $t\mapsto t^T$ is linear and maps $\mathbb{D}^{\circ}(\rho_1)$ into itself; so the function $t\mapsto\mathfrak{g}\left(t^T\right)$ is analytic on $\mathbb{D}^{\circ}(\rho_1)$, and so is $\mathfrak{g}_S$, a linear combination of finitely many such functions. For the multi-indices $\mathbf{a}=(\mathbf{a}_c)_{c\in\mathrm{Cel}}$ with entries in $\{0,1,2,\dots\}$ put $t^{\mathbf{a}}:=\prod_ct_c^{\mathbf{a}_c}$, $|\mathbf{a}|:=\sum_c\mathbf{a}_c$ and $\partial^{\mathbf{a}}:=\prod_c\partial_c^{\mathbf{a}_c}$, and let
\[
\sum_{\mathbf{a}}g_{\mathbf{a}}\cdot t^{\mathbf{a}},\qquad g_{\mathbf{a}}:=\frac{1}{\prod_c\mathbf{a}_c!}\cdot\left(\partial^{\mathbf{a}}\mathfrak{g}_S\right)(0),
\]
be the Taylor series of $\mathfrak{g}_S$ at $0$. The open polydisc $\mathbb{D}^{\circ}(\rho)$ is a neighbourhood of $0$ contained in $\mathbb{D}(\rho)$. If $c\notin S$, then $\mathfrak{g}_S$ does not depend on $t_c$ by~(a); so $\partial_c\mathfrak{g}_S=0$ on $\mathbb{D}^{\circ}(\rho)$, and $g_{\mathbf{a}}=0$ whenever $\mathbf{a}_c\ge1$. If $c\in S$, then $\mathfrak{g}_S$ vanishes at the points of $\mathbb{D}^{\circ}(\rho)$ with $t_c=0$ by~(a). A coefficient $g_{\mathbf{a}}$ with $\mathbf{a}_c=0$ is formed with derivatives at $0$ in the variables $t_{c''}$, $c''\ne c$, only; these are derivatives of the restriction of $\mathfrak{g}_S$ to the points of $\mathbb{D}^{\circ}(\rho)$ with $t_c=0$, which vanishes, and so $g_{\mathbf{a}}=0$. This proves the first assertion: $g_{\mathbf{a}}\ne0$ only if $\mathbf{a}_c\ge1$ for all $c\in S$ and $\mathbf{a}_c=0$ for all $c\notin S$.

By Cauchy's estimates on the polydisc $\mathbb{D}(\rho)$, on a neighbourhood of which $\mathfrak{g}_S$ is analytic,
\[
\left|g_{\mathbf{a}}\right|\le\rho^{-|\mathbf{a}|}\cdot\sup_{\mathbb{D}(\rho)}\left|\mathfrak{g}_S\right|\qquad\text{for every }\mathbf{a}.
\]
The image of $\mathfrak{g}_S$ in $\mathcal{J}$ is the class of $\sum_{\mathbf{a}}g_{\mathbf{a}}\cdot t^{\mathbf{a}}$, the sum being restricted to the $\mathbf{a}$ with $\mathbf{a}_c\le4$ for all $c$ (Section~\ref{sec:ring-jets}). By the first assertion only the $\mathbf{a}$ with $1\le\mathbf{a}_c\le4$ for $c\in S$ and $\mathbf{a}_c=0$ for $c\notin S$ contribute to its norm. Hence, with $x:=\frac{\rho'}{\rho}\in\left(0,\frac{1}{2}\right]$,
\[
\left\|\mathfrak{g}_S\right\|_{\rho'}=\sum_{\mathbf{a}}\left|g_{\mathbf{a}}\right|\cdot\left(\rho'\right)^{|\mathbf{a}|}\le\sup_{\mathbb{D}(\rho)}\left|\mathfrak{g}_S\right|\cdot\sum_{\mathbf{a}}x^{|\mathbf{a}|}=\sup_{\mathbb{D}(\rho)}\left|\mathfrak{g}_S\right|\cdot\left(\sum_{i=1}^4x^i\right)^{\#|S|},
\]
the sums over $\mathbf{a}$ being over the multi-indices just described. Finally $\sum_{i=1}^4x^i\le\frac{1}{2}+\frac{1}{4}+\frac{1}{8}+\frac{1}{16}=\frac{15}{16}<1$, which gives~(d).
\end{proof}

\subsubsection{Gaussian integrals over a Grassmann algebra}\label{sec:gaussian-integrals-grassmann}

Let $\mathcal{I}$ be a finite set, and let $\mathrm{Gr}$ be the Grassmann algebra over $\mathbb{C}$ with the generators $\Psi(x)$, $\bar{\Psi}(x)$, $x\in\mathcal{I}$. A monomial is a product of distinct generators; since the generators anticommute, a product of generators in which some generator occurs twice is $0$. For a generator $\mathrm{gn}$, the left derivative $d_{\mathrm{gn}}$ is the linear map of $\mathrm{Gr}$ into itself such that
\[
d_{\mathrm{gn}}\left(\mathrm{gn}\cdot m\right)=m,\qquad d_{\mathrm{gn}}(m)=0
\]
for every monomial $m$ that does not contain $\mathrm{gn}$. Each $d_{\mathrm{gn}}$ is an anti-derivation, that is,
\[
d_{\mathrm{gn}}\left(F\cdot G\right)=\left(d_{\mathrm{gn}}F\right)\cdot G+(-1)^q\cdot F\cdot d_{\mathrm{gn}}G
\]
for every $F\in\mathrm{Gr}$ that is homogeneous of degree $q$ and every $G\in\mathrm{Gr}$; and any two left derivatives anticommute. Let $\epsilon_0\colon\mathrm{Gr}\to\mathbb{C}$ assign to each $F\in\mathrm{Gr}$ its coefficient of $1$, that is, its component of degree zero; $\epsilon_0$ is a homomorphism of algebras. For a complex matrix $C$ on $\mathcal{I}$, that is, a map $C\colon\mathcal{I}\times\mathcal{I}\to\mathbb{C}$, put
\[
\Delta_C:=\sum_{x,y\in\mathcal{I}}C(x,y)\cdot d_{\bar{\Psi}(y)}d_{\Psi(x)}.
\]
Any two products of two left derivatives commute with each other. The operator $\Delta_C$ lowers the degree by two; hence $\Delta_C^k=0$ for $k>\#|\mathcal{I}|$, and $\exp\left(\Delta_C\right)=\sum_{k\ge0}\frac{1}{k!}\Delta_C^k$ is a finite sum.

\paragraph{Copies of the generators.}
Let $n\ge1$, and let $\mathrm{Gr}_n$ be the Grassmann algebra over $\mathbb{C}$ with the generators $\Psi^{(i)}(x)$, $\bar{\Psi}^{(i)}(x)$, $1\le i\le n$, $x\in\mathcal{I}$. Left derivatives on $\mathrm{Gr}_n$ and the coefficient $\epsilon_0$ of $1$ on $\mathrm{Gr}_n$ are defined as on $\mathrm{Gr}$. For $F\in\mathrm{Gr}$ and $1\le i\le n$, the element $F(\psi^{(i)})\in\mathrm{Gr}_n$ is the image of $F$ under the homomorphism of $\mathrm{Gr}$ into $\mathrm{Gr}_n$ that sends $\Psi(x)$ to $\Psi^{(i)}(x)$ and $\bar{\Psi}(x)$ to $\bar{\Psi}^{(i)}(x)$ for every $x\in\mathcal{I}$. Let $\pi_n\colon\mathrm{Gr}_n\to\mathrm{Gr}$ be the homomorphism with
\[
\pi_n\left(\Psi^{(i)}(x)\right)=\Psi(x),\qquad\pi_n\left(\bar{\Psi}^{(i)}(x)\right)=\bar{\Psi}(x)\qquad\left(1\le i\le n,\ x\in\mathcal{I}\right).
\]
For a complex matrix $C$ on $\mathcal{I}$ and $i,j\in\{1,\dots,n\}$ define
\[
\Delta^{ij}:=\sum_{x,y\in\mathcal{I}}C(x,y)\cdot d_{\bar{\Psi}^{(j)}(y)}d_{\Psi^{(i)}(x)},\qquad\Delta_{\{i,j\}}:=\Delta^{ij}+\Delta^{ji}\quad(i\ne j),
\]
and, for a family $s=\left(s_{ij}\right)_{1\le i<j\le n}$ of complex numbers,
\[
\Delta[s]:=\sum_{i=1}^n\Delta^{ii}+\sum_{1\le i<j\le n}s_{ij}\cdot\Delta_{\{i,j\}}.
\]
Complete the family by $s_{ii}:=1$ and $s_{ji}:=s_{ij}$ for $i<j$, and let $C_s$ be the matrix on $\{1,\dots,n\}\times\mathcal{I}$ with the entries
\[
C_s\left((i,x),(j,y)\right):=s_{ij}\cdot C(x,y).
\]
If $\mathrm{Gr}_n$ is regarded as the Grassmann algebra over the index set $\{1,\dots,n\}\times\mathcal{I}$, the generators with the index $(i,x)$ being $\Psi^{(i)}(x)$ and $\bar{\Psi}^{(i)}(x)$, then $\Delta[s]=\Delta_{C_s}$. Every $\Delta^{ij}$ is a linear combination of products of two left derivatives; therefore the operators $\Delta^{ij}$, $i,j\in\{1,\dots,n\}$, commute with each other.

\begin{lemma}[Wick's rule and a Gram bound]\label{lem:wick-rule-gram}
Let $C$ be a complex matrix on $\mathcal{I}$, and let $n\ge1$.
\begin{enumerate}
\item[(a)] If $C$ is invertible, then
\[
\int\mathrm{d}\mu_C(\Psi)\,F=\epsilon_0\left(\exp\left(\Delta_C\right)F\right)\qquad\left(F\in\mathrm{Gr}\right),
\]
where $\mathrm{d}\mu_C$ is the Grassmann Gaussian integral with the covariance $C$, defined as in Section~\ref{sec:grassmann-norms-gaussian} with the generators of $\mathrm{Gr}$.
\item[(b)] Let $p,p'\ge0$ and $x_1,\dots,x_p,y_1,\dots,y_{p'}\in\mathcal{I}$. Then
\[
\epsilon_0\left(\exp\left(\Delta_C\right)\left[\Psi(x_1)\cdots\Psi(x_p)\cdot\bar{\Psi}(y_1)\cdots\bar{\Psi}(y_{p'})\right]\right)=
\begin{cases}
0&\text{if }p\ne p',\\
(-1)^{\frac{p(p-1)}{2}}\cdot\det\left[C(x_a,y_b)\right]_{a,b\le p}&\text{if }p=p'.
\end{cases}
\]
\item[(c)] Let $[1]$ denote the family with $s_{ij}=1$ for all $i<j$. Then
\[
\epsilon_0\left(\exp\left(\Delta[1]\right)F\right)=\epsilon_0\left(\exp\left(\Delta_C\right)\pi_n(F)\right)\qquad\left(F\in\mathrm{Gr}_n\right).
\]
\item[(d)] Let $\{1,\dots,n\}$ be the disjoint union of the sets $B_1,\dots,B_r$. For $1\le a\le r$ let $F_a$ be an even element of the subalgebra of $\mathrm{Gr}_n$ spanned by the monomials in the generators $\Psi^{(i)}(x)$, $\bar{\Psi}^{(i)}(x)$ with $i\in B_a$ and $x\in\mathcal{I}$, and let $O_a$ be a polynomial in the operators $\Delta^{ij}$ with $i,j\in B_a$. Then
\[
\epsilon_0\left(\left(\prod_{a=1}^rO_a\right)\prod_{a=1}^rF_a\right)=\prod_{a=1}^r\epsilon_0\left(O_aF_a\right).
\]
\item[(e)] Let $s=\left(s_{ij}\right)_{i<j}$, and let $\boldsymbol{\omega}_1,\dots,\boldsymbol{\omega}_n$ be unit vectors of a Hilbert space such that $s_{ij}=\left\langle\boldsymbol{\omega}_i,\boldsymbol{\omega}_j\right\rangle$ for all $i,j\le n$, with $s_{ii}=1$ and $s_{ji}=s_{ij}$ as above. Then, for every $p\ge0$ and every monomial $m$ of degree $2p$ in the generators of $\mathrm{Gr}_n$,
\[
\left|\epsilon_0\left(\exp\left(\Delta[s]\right)m\right)\right|\le\left|C\right|_{\ell^2\to\ell^2}^p,
\]
where $\left|C\right|_{\ell^2\to\ell^2}$ is the operator norm of $C$ on $\ell^2(\mathcal{I})$.
\end{enumerate}
\end{lemma}

\begin{proof}
For $F\in\mathrm{Gr}$ put $\mathfrak{J}(F):=\epsilon_0\left(\exp\left(\Delta_C\right)F\right)$.

\emph{A recursion for $\mathfrak{J}$.} Let $x\in\mathcal{I}$ and $F\in\mathrm{Gr}$. For $x'\in\mathcal{I}$ the anti-derivation rule gives $d_{\Psi(x')}\left(\Psi(x)\cdot F\right)=\delta_{xx'}\cdot F-\Psi(x)\cdot d_{\Psi(x')}F$; moreover $d_{\bar{\Psi}(y)}\Psi(x)=0$, so that $d_{\bar{\Psi}(y)}\left(\Psi(x)\cdot G\right)=-\Psi(x)\cdot d_{\bar{\Psi}(y)}G$ for all $G\in\mathrm{Gr}$. Therefore
\[
\Delta_C\left(\Psi(x)\cdot F\right)=\sum_{x',y\in\mathcal{I}}C(x',y)\cdot d_{\bar{\Psi}(y)}\left(\delta_{xx'}\cdot F-\Psi(x)\cdot d_{\Psi(x')}F\right)=T_xF+\Psi(x)\cdot\Delta_CF,
\]
with the operator $T_x:=\sum_{y\in\mathcal{I}}C(x,y)\cdot d_{\bar{\Psi}(y)}$. A left derivative anticommutes with each of the two factors of a product of two left derivatives, and hence commutes with the product. So each $d_{\bar{\Psi}(y)}$, and with it $T_x$, commutes with $\Delta_C$ and with $\exp\left(\Delta_C\right)$. By induction on $k$,
\[
\Delta_C^k\left(\Psi(x)\cdot F\right)=\Psi(x)\cdot\Delta_C^kF+k\cdot T_x\Delta_C^{k-1}F\qquad(k\ge1),
\]
and summation over $k$ with the weights $\frac{1}{k!}$ gives
\[
\exp\left(\Delta_C\right)\left(\Psi(x)\cdot F\right)=\Psi(x)\cdot\exp\left(\Delta_C\right)F+T_x\exp\left(\Delta_C\right)F.
\]
Every monomial of $\Psi(x)\cdot G$ contains $\Psi(x)$; hence $\epsilon_0\left(\Psi(x)\cdot G\right)=0$ for all $G\in\mathrm{Gr}$. Applying $\epsilon_0$ to the last identity and using that $d_{\bar{\Psi}(y)}$ commutes with $\exp\left(\Delta_C\right)$, one obtains
\begin{equation}\label{eq:four-jets-7}
\mathfrak{J}\left(\Psi(x)\cdot F\right)=\sum_{y\in\mathcal{I}}C(x,y)\cdot\mathfrak{J}\left(d_{\bar{\Psi}(y)}F\right)\qquad\left(x\in\mathcal{I},\ F\in\mathrm{Gr}\right).
\end{equation}
Further, $\Delta_C$ maps a monomial with $p$ factors $\Psi$ and $p'$ factors $\bar{\Psi}$ into the span of the monomials with $p-1$ factors $\Psi$ and $p'-1$ factors $\bar{\Psi}$. Consequently $\exp\left(\Delta_C\right)$ maps such a monomial into the span of the monomials with $p-k$ factors $\Psi$ and $p'-k$ factors $\bar{\Psi}$, $k\ge0$, and its coefficient of $1$ can differ from $0$ only if $p=p'$. Thus $\mathfrak{J}$ vanishes on every monomial with $p\ne p'$; and $\mathfrak{J}(1)=1$, since $\Delta_C1=0$.

\emph{(b)} The left side of~(b) is the value of $\mathfrak{J}$ at the product $\Psi(x_1)\cdots\Psi(x_p)\cdot\bar{\Psi}(y_1)\cdots\bar{\Psi}(y_{p'})$. If $p\ne p'$, this product is $0$ or a monomial with $p$ factors $\Psi$ and $p'$ factors $\bar{\Psi}$, and in both cases $\mathfrak{J}$ vanishes on it. For $p=p'$ the assertion is proved by induction on $p$; for $p=0$ it reads $\mathfrak{J}(1)=1$, the determinant of the empty matrix being $1$. Put $\boldsymbol{\sigma}_q:=(-1)^{\frac{q(q-1)}{2}}$ for $q\ge0$; then $\boldsymbol{\sigma}_p=(-1)^{p-1}\cdot\boldsymbol{\sigma}_{p-1}$ for $p\ge1$. Let $p\ge1$, and assume the assertion for $p-1$ and all points of $\mathcal{I}$. If $y_b=y_{b'}$ for some $b\ne b'$, then the product contains the generator $\bar{\Psi}(y_b)$ twice and is $0$, and the matrix $\left[C(x_a,y_b)\right]_{a,b\le p}$ has two equal columns, so that both sides are $0$. Let the points $y_1,\dots,y_p$ be distinct, and put
\[
G:=\Psi(x_2)\cdots\Psi(x_p)\cdot\bar{\Psi}(y_1)\cdots\bar{\Psi}(y_p),\qquad G_b:=\Psi(x_2)\cdots\Psi(x_p)\cdot\prod_{b'\ne b}\bar{\Psi}(y_{b'})\quad(1\le b\le p),
\]
the factors of the last product being taken in the increasing order of $b'$. In $G$ the factor $\bar{\Psi}(y_b)$ is preceded by $p-1+b-1$ factors; bringing it to the front, one finds
\[
d_{\bar{\Psi}(y)}G=\sum_{b=1}^p\delta_{yy_b}\cdot(-1)^{p-1+b-1}\cdot G_b\qquad\left(y\in\mathcal{I}\right).
\]
By~\eqref{eq:four-jets-7} with $x=x_1$ and $F=G$, and by the induction hypothesis applied to the products $G_b$,
\begin{align*}
\mathfrak{J}\left(\Psi(x_1)\cdot G\right)&=\sum_{b=1}^p(-1)^{p+b}\cdot C(x_1,y_b)\cdot\boldsymbol{\sigma}_{p-1}\cdot\det\left[C(x_a,y_{b'})\right]_{2\le a\le p,\ b'\ne b}\\
&=(-1)^{p-1}\cdot\boldsymbol{\sigma}_{p-1}\cdot\sum_{b=1}^p(-1)^{b-1}\cdot C(x_1,y_b)\cdot\det\left[C(x_a,y_{b'})\right]_{2\le a\le p,\ b'\ne b},
\end{align*}
the rows of the matrices being indexed by $a=2,\dots,p$ and their columns by the $b'\ne b$ in increasing order. The last sum is the expansion of $\det\left[C(x_a,y_b)\right]_{a,b\le p}$ along its first row, and $(-1)^{p-1}\cdot\boldsymbol{\sigma}_{p-1}=\boldsymbol{\sigma}_p$. This is the assertion for $p$.

\emph{(a)} Let $C$ be invertible, and let $C^{-1}(y,x)$, $x,y\in\mathcal{I}$, be the entries of the inverse matrix. For $y\in\mathcal{I}$ put
\[
\left(C^{-1}\Psi\right)(y):=\sum_{x\in\mathcal{I}}C^{-1}(y,x)\cdot\Psi(x),
\]
and let
\[
\mathfrak{Q}:=\left\langle\bar{\Psi},C^{-1}\Psi\right\rangle=\sum_{x,y\in\mathcal{I}}\bar{\Psi}(y)\cdot C^{-1}(y,x)\cdot\Psi(x).
\]
By the definition of the Gaussian integral and of $\int\mathrm{D}\Psi$ (Sections~\ref{sec:grassmann-norms-gaussian}, \ref{sec:fermion-integrals-blocking} and~\ref{sec:norms-gaussian-integrals}, formed with the generators of $\mathrm{Gr}$),
\[
\mathbf{I}(F):=\int\mathrm{d}\mu_C(\Psi)\,F=\det C\cdot\int\mathrm{D}\Psi\,\exp\left(-\mathfrak{Q}\right)\cdot F\qquad\left(F\in\mathrm{Gr}\right),
\]
where $\int\mathrm{D}\Psi\,G$ is the coefficient, in $G$, of the monomial $\prod_{x\in\mathcal{I}}\Psi(x)\cdot\bar{\Psi}(x)$ of top degree; the factors $\Psi(x)\cdot\bar{\Psi}(x)$ of this product are even and commute with each other. The monomial of top degree contains every generator, whereas no monomial of $d_{\bar{\Psi}(y)}G$ contains $\bar{\Psi}(y)$. Hence
\[
\int\mathrm{D}\Psi\,d_{\bar{\Psi}(y)}G=0\qquad\left(G\in\mathrm{Gr},\ y\in\mathcal{I}\right).
\]
The element $\mathfrak{Q}$ is even, and $d_{\bar{\Psi}(y)}\mathfrak{Q}=\left(C^{-1}\Psi\right)(y)$. Since an even element commutes with every element of $\mathrm{Gr}$, the anti-derivation rule gives $d_{\bar{\Psi}(y)}\left(\mathfrak{Q}^k\right)=k\cdot\left(C^{-1}\Psi\right)(y)\cdot\mathfrak{Q}^{k-1}$ for $k\ge1$, and therefore
\[
d_{\bar{\Psi}(y)}\exp\left(-\mathfrak{Q}\right)=-\left(C^{-1}\Psi\right)(y)\cdot\exp\left(-\mathfrak{Q}\right).
\]
Let $F\in\mathrm{Gr}$ and $G:=\exp\left(-\mathfrak{Q}\right)\cdot F$. As $\exp\left(-\mathfrak{Q}\right)$ is even,
\[
d_{\bar{\Psi}(y)}G=-\left(C^{-1}\Psi\right)(y)\cdot\exp\left(-\mathfrak{Q}\right)\cdot F+\exp\left(-\mathfrak{Q}\right)\cdot d_{\bar{\Psi}(y)}F,
\]
and $\left(C^{-1}\Psi\right)(y)\cdot\exp\left(-\mathfrak{Q}\right)=\exp\left(-\mathfrak{Q}\right)\cdot\left(C^{-1}\Psi\right)(y)$. Integrating with $\int\mathrm{D}\Psi$ and multiplying by $\det C$ one obtains
\[
\mathbf{I}\left(\left(C^{-1}\Psi\right)(y)\cdot F\right)=\mathbf{I}\left(d_{\bar{\Psi}(y)}F\right)\qquad\left(y\in\mathcal{I}\right).
\]
Multiply by $C(x,y)$ and sum over $y$. Since $\sum_{y\in\mathcal{I}}C(x,y)\cdot\left(C^{-1}\Psi\right)(y)=\Psi(x)$, the functional $\mathbf{I}$ satisfies~\eqref{eq:four-jets-7} in place of $\mathfrak{J}$.

By~\eqref{eq:four-jets-Gr3} with $F=1$, together with the normalization $\int\mathrm{D}\Psi\,\exp\left(-\left\langle\bar{\Psi},O\Psi\right\rangle\right)=\det O$ of Section~\ref{sec:fermion-integrals-blocking} for $O=C^{-1}$, one has $\mathbf{I}(1)=1$. Every monomial of $\exp\left(-\mathfrak{Q}\right)$ has as many factors $\Psi$ as factors $\bar{\Psi}$, and so has the monomial of top degree. If $m\ne1$ is a monomial without a factor $\Psi$, then every monomial of $\exp\left(-\mathfrak{Q}\right)\cdot m$ has more factors $\bar{\Psi}$ than factors $\Psi$; hence $\mathbf{I}(m)=0$. Such a monomial $m$ has no factor $\Psi$ and at least one factor $\bar{\Psi}$; as $\mathfrak{J}$ vanishes on every monomial whose numbers of factors $\Psi$ and $\bar{\Psi}$ differ, also $\mathfrak{J}(m)=0$. Moreover $\mathfrak{J}(1)=1$, because $\Delta_C1=0$.

Every monomial either has no factor $\Psi$, or is, up to sign, of the form $\Psi(x)\cdot m'$ with some $x\in\mathcal{I}$ and a monomial $m'$ with one factor $\Psi$ fewer; and $d_{\bar{\Psi}(y)}m'$ is $0$ or, up to sign, a monomial with as many factors $\Psi$ as $m'$. The linear functionals $\mathbf{I}$ and $\mathfrak{J}$ both satisfy~\eqref{eq:four-jets-7}, take the value $1$ at $1$, and vanish on the monomials $m\ne1$ without a factor $\Psi$. Induction on the number of factors $\Psi$ shows that they coincide on all monomials. Hence $\mathbf{I}=\mathfrak{J}$, which is (a).

\emph{(c)} Call a linear map $\mathfrak{D}\colon\mathrm{Gr}_n\to\mathrm{Gr}$ a $\pi_n$-anti-derivation if
\[
\mathfrak{D}\left(F\cdot G\right)=\left(\mathfrak{D} F\right)\cdot\pi_n(G)+(-1)^q\cdot\pi_n(F)\cdot\mathfrak{D} G
\]
for every $F\in\mathrm{Gr}_n$ homogeneous of degree $q$ and every $G\in\mathrm{Gr}_n$. The difference of two such maps is again one. If it vanishes on the generators, then it vanishes at $1$ (take $F=G=1$) and, by induction on the degree, on all monomials. Let $x\in\mathcal{I}$. As $\pi_n$ is a homomorphism of algebras that preserves degrees, $d_{\Psi(x)}\circ\pi_n$ and $\pi_n\circ\sum_{i=1}^nd_{\Psi^{(i)}(x)}$ are $\pi_n$-anti-derivations; both map $\Psi^{(i)}(x')$ to $\delta_{xx'}$ and $\bar{\Psi}^{(i)}(y)$ to $0$, for all $i$, $x'$, $y$. Hence they coincide, and in the same way for the generators $\bar{\Psi}$:
\[
d_{\Psi(x)}\circ\pi_n=\pi_n\circ\sum_{i=1}^nd_{\Psi^{(i)}(x)},\qquad d_{\bar{\Psi}(y)}\circ\pi_n=\pi_n\circ\sum_{j=1}^nd_{\bar{\Psi}^{(j)}(y)}\qquad\left(x,y\in\mathcal{I}\right).
\]
It follows that
\[
\Delta_C\circ\pi_n=\sum_{x,y\in\mathcal{I}}C(x,y)\cdot\pi_n\circ\sum_{i,j=1}^nd_{\bar{\Psi}^{(j)}(y)}d_{\Psi^{(i)}(x)}=\pi_n\circ\sum_{i,j=1}^n\Delta^{ij}=\pi_n\circ\Delta[1],
\]
then $\Delta_C^k\circ\pi_n=\pi_n\circ\Delta[1]^k$ for all $k\ge0$, and $\exp\left(\Delta_C\right)\circ\pi_n=\pi_n\circ\exp\left(\Delta[1]\right)$. Moreover $\epsilon_0\circ\pi_n=\epsilon_0$, because $\pi_n$ fixes $1$ and maps every monomial of positive degree to $0$ or to a monomial of the same degree. Applying $\epsilon_0$ to $\exp\left(\Delta_C\right)\pi_n(F)=\pi_n\left(\exp\left(\Delta[1]\right)F\right)$ gives (c).

\emph{(d)} For $1\le a'\le r$ call the subalgebra of $\mathrm{Gr}_n$ spanned by the monomials in the generators with an upper index $i\in B_{a'}$ the subalgebra of the block $B_{a'}$. Fix $a$, and let $d$ be the left derivative with respect to a generator $\Psi^{(i)}(x)$ or $\bar{\Psi}^{(i)}(x)$ with $i\in B_a$. Let $G_1,\dots,G_r\in\mathrm{Gr}_n$ be such that $G_{a'}$ lies in the subalgebra of the block $B_{a'}$ for every $a'$ and is even for $a'\ne a$. For $a'\ne a$ the element $G_{a'}$ contains no generator with an upper index in $B_a$; so $d\,G_{a'}=0$, the anti-derivation $d$ annihilates every product of such elements, and $d\left(G_{a'}\cdot G\right)=G_{a'}\cdot d\,G$ for all $G\in\mathrm{Gr}_n$, since $G_{a'}$ is even. Consequently
\[
d\left(G_1\cdots G_r\right)=G_1\cdots G_{a-1}\cdot\left(d\,G_a\right)\cdot G_{a+1}\cdots G_r.
\]
Applying this twice, one sees that every $\Delta^{ij}$ with $i,j\in B_a$ acts on such a product through the factor $G_a$ alone, and it maps even elements of the subalgebra of the block $B_a$ to even elements of the same subalgebra. By induction on the degree of the polynomial, the operator $O_a$ has the same two properties. Applying the commuting operators $O_r,O_{r-1},\dots,O_1$ in turn to $\prod_aF_a$ one obtains
\[
\left(\prod_{a=1}^rO_a\right)\prod_{a=1}^rF_a=\prod_{a=1}^r\left(O_aF_a\right),
\]
and (d) follows, as $\epsilon_0$ is multiplicative.

\emph{(e)} Regard $\mathrm{Gr}_n$ as the Grassmann algebra over $\{1,\dots,n\}\times\mathcal{I}$; then $\Delta[s]=\Delta_{C_s}$, and (b), applied to this index set and to the matrix $C_s$, gives the following. The number $\epsilon_0\left(\exp\left(\Delta[s]\right)m\right)$ is $0$, unless $m$ is, up to sign, of the form $\Psi^{(i_1)}(x_1)\cdots\Psi^{(i_p)}(x_p)\cdot\bar{\Psi}^{(j_1)}(y_1)\cdots\bar{\Psi}^{(j_p)}(y_p)$; in that case it equals, up to sign,
\[
\det\left[s_{i_aj_b}\cdot C(x_a,y_b)\right]_{a,b\le p}=\det\left[\left\langle\boldsymbol{\omega}_{i_a}\otimes\boldsymbol{\delta}_{x_a},\,\boldsymbol{\omega}_{j_b}\otimes C\boldsymbol{\delta}_{y_b}\right\rangle\right]_{a,b\le p}.
\]
Here $\boldsymbol{\delta}_x$, $x\in\mathcal{I}$, are the unit vectors of $\ell^2(\mathcal{I})$, the matrix $C$ acts on $\ell^2(\mathcal{I})$, and the scalar products are those of the tensor product of the given Hilbert space with $\ell^2(\mathcal{I})$, antilinear in the first argument; thus $\left\langle\boldsymbol{\omega}\otimes f,\boldsymbol{\omega}'\otimes g\right\rangle=\left\langle\boldsymbol{\omega},\boldsymbol{\omega}'\right\rangle\cdot\left\langle f,g\right\rangle$ and $\left\langle\boldsymbol{\delta}_x,C\boldsymbol{\delta}_y\right\rangle=C(x,y)$. Gram's inequality states that
\[
\left|\det\left[\left\langle f_a,\vec{g}_b\right\rangle\right]_{a,b\le p}\right|\le\prod_{a=1}^p\left|f_a\right|\cdot\prod_{b=1}^p\left|\vec{g}_b\right|
\]
for arbitrary vectors $f_1,\dots,f_p,\vec{g}_1,\dots,\vec{g}_p$ of a Hilbert space. For $f_a:=\boldsymbol{\omega}_{i_a}\otimes\boldsymbol{\delta}_{x_a}$ and $\vec{g}_b:=\boldsymbol{\omega}_{j_b}\otimes C\boldsymbol{\delta}_{y_b}$ one has $\left|f_a\right|=1$ and $\left|\vec{g}_b\right|=\left|C\boldsymbol{\delta}_{y_b}\right|\le\left|C\right|_{\ell^2\to\ell^2}$, which proves (e).
\end{proof}

\subsubsection{Interpolation over forests}\label{sec:interpolation-forests}

Let $n\ge1$, and let $P_n$ be the set of the pairs $\{i,j\}$ with $1\le i<j\le n$. A set of pairs is regarded as the set of the edges of a graph with the vertex set $\{1,\dots,n\}$, and its elements are called lines. A forest on $\{1,\dots,n\}$ is a subset $\mathrm{Fo}\subset P_n$ that contains no cycle; the empty set is a forest. A tree is a forest with $n-1$ elements. Two vertices $i\ne j$ are joined by $\mathrm{Fo}$ if they belong to the same connected component of the graph with the lines $\mathrm{Fo}$; as $\mathrm{Fo}$ has no cycle, there is then exactly one path from $i$ to $j$ in $\mathrm{Fo}$. For a forest $\mathrm{Fo}$ and $w=\left(w_l\right)_{l\in\mathrm{Fo}}\in[0,1]^{\mathrm{Fo}}$ let $s^{\mathrm{Fo}}(w)\in[0,1]^{P_n}$ be the point with the coordinates
\[
s^{\mathrm{Fo}}_{ij}(w):=
\begin{cases}
\min\left\{w_l:l\text{ is a line of the path from }i\text{ to }j\text{ in }\mathrm{Fo}\right\}&\text{if }i\text{ and }j\text{ are joined by }\mathrm{Fo},\\
0&\text{otherwise},
\end{cases}
\qquad\left(\{i,j\}\in P_n\right).
\]

\begin{lemma}[Forest formula of Brydges--Kennedy and Abdesselam--Rivasseau]\label{lem:forest-formula-brydges}
\begin{enumerate}
\item[(a)] For every polynomial $H$ in the variables $s_l$, $l\in P_n$,
\[
H(1,\dots,1)=\sum_{\mathrm{Fo}}\int_{[0,1]^{\mathrm{Fo}}}\mathrm{d}w\,\left(\left(\prod_{l\in\mathrm{Fo}}\partial_{s_l}\right)H\right)\left(s^{\mathrm{Fo}}(w)\right),
\]
the sum ranging over all forests $\mathrm{Fo}$ on $\{1,\dots,n\}$; for $\mathrm{Fo}=\emptyset$ the integral is the value of the integrand at the only point of $[0,1]^{\emptyset}$, namely $H(0,\dots,0)$.
\item[(b)] For every forest $\mathrm{Fo}$ on $\{1,\dots,n\}$ and every $w\in[0,1]^{\mathrm{Fo}}$ there are unit vectors $\boldsymbol{\omega}_1,\dots,\boldsymbol{\omega}_n$ of a Hilbert space such that $s^{\mathrm{Fo}}_{ij}(w)=\left\langle\boldsymbol{\omega}_i,\boldsymbol{\omega}_j\right\rangle$ for all $i\ne j$.
\end{enumerate}
\end{lemma}

\begin{proof}
\emph{(a)} Let $\mathfrak{p}$ be a partition of $\{1,\dots,n\}$. A pair $l\in P_n$ lies inside a block of $\mathfrak{p}$ if both of its elements belong to that block; otherwise $l$ joins two blocks. Let $s_{\mathrm{in}}$ be a family of numbers $s_{\mathrm{in}}(l)\in[0,1]$, one for each pair $l$ inside a block of $\mathfrak{p}$, and let $u\in[0,1]$. Let $\mathsf{X}\left(\mathfrak{p},s_{\mathrm{in}},u\right)\in[0,1]^{P_n}$ be the point whose coordinate at $l$ is $s_{\mathrm{in}}(l)$ if $l$ lies inside a block, and $u$ if $l$ joins two blocks. Let $G$ be a polynomial in the variables $s_l$, $l\in P_n$. By the chain rule, the derivative of the function $w\mapsto G\left(\mathsf{X}\left(\mathfrak{p},s_{\mathrm{in}},w\right)\right)$ is $\sum_l\left(\partial_{s_l}G\right)\left(\mathsf{X}\left(\mathfrak{p},s_{\mathrm{in}},w\right)\right)$, the sum ranging over the pairs $l$ that join two blocks of $\mathfrak{p}$. By the fundamental theorem of calculus,
\begin{equation}\label{eq:four-jets-8}
G\left(\mathsf{X}\left(\mathfrak{p},s_{\mathrm{in}},u\right)\right)=G\left(\mathsf{X}\left(\mathfrak{p},s_{\mathrm{in}},0\right)\right)+\sum_l\int_0^u\mathrm{d}w\,\left(\partial_{s_l}G\right)\left(\mathsf{X}\left(\mathfrak{p},s_{\mathrm{in}},w\right)\right),
\end{equation}
with the same range of summation.

Let $l$ join the blocks $B$ and $B'$ of $\mathfrak{p}$, and let $w\in[0,1]$. Let $\mathfrak{p}'$ be the partition obtained from $\mathfrak{p}$ by merging $B$ and $B'$ into one block, and let $s_{\mathrm{in}}'$ be the family that coincides with $s_{\mathrm{in}}$ at the pairs inside the blocks of $\mathfrak{p}$ and takes the value $w$ at the pairs that join $B$ and $B'$. Then
\[
\mathsf{X}\left(\mathfrak{p},s_{\mathrm{in}},w\right)=\mathsf{X}\left(\mathfrak{p}',s_{\mathrm{in}}',w\right).
\]

Begin with the partition of $\{1,\dots,n\}$ into singletons, inside whose blocks there is no pair, with $u=1$ and $G=H$; then the left side of~\eqref{eq:four-jets-8} is $H(1,\dots,1)$. Write each integrand on the right side as $\left(\partial_{s_l}G\right)\left(\mathsf{X}\left(\mathfrak{p}',s_{\mathrm{in}}',w\right)\right)$, and expand it again by~\eqref{eq:four-jets-8}, with $\left(\mathfrak{p}',s_{\mathrm{in}}',w\right)$ in place of $\left(\mathfrak{p},s_{\mathrm{in}},u\right)$ and $\partial_{s_l}G$ in place of $G$; then treat the new integrands in the same way. Each expansion lowers the number of blocks by one, and for a partition with a single block no pair joins two blocks. So the procedure ends after at most $n-1$ expansions, and it gives
\[
H(1,\dots,1)=\sum_{k\ge0}\ \sum_{(l_1,\dots,l_k)}\ \int_{1\ge w_1\ge\dots\ge w_k\ge0}\mathrm{d}w\,\left(\partial_{s_{l_1}}\cdots\partial_{s_{l_k}}H\right)\left(X_k\right).
\]
Here the inner sum ranges over the sequences $(l_1,\dots,l_k)$ of pairs such that, for each $m\le k$, the pair $l_m$ joins two different connected components of the graph with the lines $l_1,\dots,l_{m-1}$; these sequences are exactly the orderings of the forests with $k$ lines. The point $X_k\in[0,1]^{P_n}$ has the coordinate $0$ at every pair whose elements are not joined by $\{l_1,\dots,l_k\}$, and the coordinate $w_m$ at a pair $\{i,j\}$ whose elements are merged into one block at the $m$-th expansion. The term with $k=0$ is $H(0,\dots,0)$.

Let $\mathrm{Fo}:=\{l_1,\dots,l_k\}$, and put $w_{l_m}:=w_m$ for $m\le k$. If $i$ and $j$ are joined by $\mathrm{Fo}$, they are merged at the expansion at which the last line of the path from $i$ to $j$ in $\mathrm{Fo}$ is added; thus $m$ is the largest index of a line of this path. Since $w_1\ge\dots\ge w_k$, the number $w_m$ is the minimum of the numbers $w_l$ over the lines $l$ of the path. Hence $X_k=s^{\mathrm{Fo}}(w)$, and the integrand equals $\left(\left(\prod_{l\in\mathrm{Fo}}\partial_{s_l}\right)H\right)\left(s^{\mathrm{Fo}}(w)\right)$, the partial derivatives of a polynomial being interchangeable; it depends on $\left(\mathrm{Fo},w\right)$ only. For a fixed forest $\mathrm{Fo}$ with $k$ lines, the $k!$ simplices
\[
\left\{w\in[0,1]^{\mathrm{Fo}}:w_{l_1}\ge\dots\ge w_{l_k}\right\},
\]
one for each ordering $(l_1,\dots,l_k)$ of $\mathrm{Fo}$, cover $[0,1]^{\mathrm{Fo}}$, and any two of them meet in a set of measure zero. Summing over the orderings of each forest one obtains (a).

\emph{(b)} Let $\mathrm{Fo}$ and $w$ be given. For $u\in[0,1)$ write $i\sim_uj$ if $i=j$, or if $i$ and $j$ are joined by the set of the lines $l\in\mathrm{Fo}$ with $w_l>u$. This is an equivalence relation on $\{1,\dots,n\}$, and its classes are subsets of $\{1,\dots,n\}$. Let $i\ne j$. If $i$ and $j$ are joined by $\mathrm{Fo}$, then $i\sim_uj$ holds if and only if $w_l>u$ for every line $l$ of the path from $i$ to $j$ in $\mathrm{Fo}$, that is, if and only if $u<s^{\mathrm{Fo}}_{ij}(w)$; if $i$ and $j$ are not joined by $\mathrm{Fo}$, then $i\sim_uj$ holds for no $u$. In both cases
\[
\int_0^1\mathrm{d}u\,\mathbf{1}\left[i\sim_uj\right]=s^{\mathrm{Fo}}_{ij}(w).
\]
Let $\boldsymbol{\omega}_i\in L^2\left([0,1);\ell^2\left(2^{\{1,\dots,n\}}\right)\right)$ be the function whose value at $u$ is the unit vector of $\ell^2\left(2^{\{1,\dots,n\}}\right)$ that belongs to the class of $i$ for $\sim_u$. It is constant between consecutive values of the numbers $w_l$, hence measurable, and $\left|\boldsymbol{\omega}_i(u)\right|=1$ for every $u$; so $\left|\boldsymbol{\omega}_i\right|=1$. For $i\ne j$ the scalar product $\left\langle\boldsymbol{\omega}_i(u),\boldsymbol{\omega}_j(u)\right\rangle$ equals $1$ if $i\sim_uj$ and $0$ otherwise. Therefore
\[
\left\langle\boldsymbol{\omega}_i,\boldsymbol{\omega}_j\right\rangle=\int_0^1\mathrm{d}u\,\mathbf{1}\left[i\sim_uj\right]=s^{\mathrm{Fo}}_{ij}(w)\qquad(i\ne j).
\]
\end{proof}

\subsubsection{Weighted tree expansion with two marked vertices}\label{sec:weighted-tree-expansion}

The objects of Section~\ref{sec:gaussian-integrals-grassmann} are kept: $\mathcal{I}$ is a finite set, $C$ is a complex matrix on $\mathcal{I}$, $\mathrm{Gr}$ is the Grassmann algebra with the generators $\Psi(x)$, $\bar{\Psi}(x)$, $x\in\mathcal{I}$, and the left derivatives, $\epsilon_0$, $\Delta_C$, the algebras $\mathrm{Gr}_n$ of copies, $\pi_n$, $\Delta^{ij}$, $\Delta_{\{i,j\}}$ and $\Delta[s]$ are those defined there. In addition, $d$ is a pseudo-metric on $\mathcal{I}$. Let
\[
\mathrm{Idx}:=\mathcal{I}\times\{+,-\},\qquad\psi(x,-):=\Psi(x),\qquad\psi(x,+):=\bar{\Psi}(x)\qquad(x\in\mathcal{I});
\]
thus the elements $\psi(\xi)$, $\xi\in\mathrm{Idx}$, are the generators of $\mathrm{Gr}$, and the left derivative with respect to $\psi(\xi)$ is denoted by $d_{\psi(\xi)}$. Likewise $\psi^{(i)}(x,-):=\Psi^{(i)}(x)$ and $\psi^{(i)}(x,+):=\bar{\Psi}^{(i)}(x)$ are the generators of the copy $i$ in $\mathrm{Gr}_n$, and for $\mathbf{F}\in\mathrm{Gr}$ the element $\mathbf{F}\left(\psi^{(i)}\right)\in\mathrm{Gr}_n$ arises from $\mathbf{F}$ by replacing every $\psi(\xi)$ with $\psi^{(i)}(\xi)$. The pseudo-metric is transferred to $\mathrm{Idx}$ by $d\left((x,\sigma),(y,\sigma')\right):=d(x,y)$ for $\sigma,\sigma'\in\{+,-\}$. For $\xi=(x,\sigma)\in\mathrm{Idx}$ and non-empty subsets $P$, $P'$ of $\mathcal{I}$ put
\[
d(\xi,P):=\min_{y\in P}d(x,y),\qquad d(P,P'):=\min_{y\in P,\ y'\in P'}d(y,y').
\]

Let $\vartheta\ge 0$, and let $g\ge 0$, $h_G$ and $\boldsymbol{\alpha}_{\vartheta}$ be numbers such that
\begin{gather*}
g^2\ge|C|_{\ell^2\to\ell^2},\qquad h_G\ge 1+g,\\
\boldsymbol{\alpha}_{\vartheta}\ge\sup_{x\in\mathcal{I}}\sum_{y\in\mathcal{I}}|C(x,y)|\cdot\exp\left(\vartheta\cdot d(x,y)\right),\qquad\boldsymbol{\alpha}_{\vartheta}\ge\sup_{y\in\mathcal{I}}\sum_{x\in\mathcal{I}}|C(x,y)|\cdot\exp\left(\vartheta\cdot d(x,y)\right),
\end{gather*}
where $|C|_{\ell^2\to\ell^2}$ is the operator norm of $C$ on $\ell^2(\mathcal{I})$. Moreover, $\boldsymbol{\alpha}_0$ is a number that satisfies the last two inequalities with the rate $0$ in place of $\vartheta$, and $\boldsymbol{\alpha}_0\le\boldsymbol{\alpha}_{\vartheta}$; this inequality is no restriction, since $\boldsymbol{\alpha}_{\vartheta}$ itself satisfies the two inequalities required of $\boldsymbol{\alpha}_0$.

Fix a total order on $\mathrm{Idx}$. Every even element $\mathbf{F}$ of $\mathrm{Gr}$ without constant term has exactly one representation
\[
\mathbf{F}=\sum_{m\ge 2}\frac{1}{m!}\sum_{\xi_1,\dots,\xi_m\in\mathrm{Idx}}\mathbf{F}_m(\xi_1,\dots,\xi_m)\cdot\psi(\xi_1)\cdots\psi(\xi_m)
\]
with antisymmetric functions $\mathbf{F}_m\colon\mathrm{Idx}^m\to\mathbb{C}$, the kernels of $\mathbf{F}$. Indeed, for $\xi_1<\dots<\xi_m$ the $m!$ rearrangements of $(\xi_1,\dots,\xi_m)$ contribute equally, so that $\mathbf{F}_m(\xi_1,\dots,\xi_m)$ is the coefficient of the monomial $\psi(\xi_1)\cdots\psi(\xi_m)$ in $\mathbf{F}$, and antisymmetry determines the other values. In particular $\mathbf{F}_m=0$ for odd $m$.

A weight is a function $w$ which assigns a number $w(\xi_1,\dots,\xi_m)\ge 1$ to every family $(\xi_1,\dots,\xi_m)$, $m\ge 1$, of elements of $\mathrm{Idx}$; it is symmetric if $w(\xi_1,\dots,\xi_m)$ does not change under the permutations of $\xi_1,\dots,\xi_m$. For a symmetric weight $w$ and an even element $\mathbf{F}$ of $\mathrm{Gr}$ without constant term, the numbers $\rho_h(\mathbf{F};w),N_h(\mathbf{F};w)\in[0,\infty)$ are defined by
\begin{align*}
\rho_h(\mathbf{F};w)&:=\sum_{m\ge 2}\frac{h_G^m}{(m-1)!}\cdot\sup_{\xi_1\in\mathrm{Idx}}\sum_{\xi_2,\dots,\xi_m\in\mathrm{Idx}}\left|\mathbf{F}_m(\xi_1,\dots,\xi_m)\right|\cdot w(\xi_1,\dots,\xi_m),\\
N_h(\mathbf{F};w)&:=\sum_{m\ge 2}\frac{h_G^m}{(m-1)!}\cdot\sum_{\xi_1,\dots,\xi_m\in\mathrm{Idx}}\left|\mathbf{F}_m(\xi_1,\dots,\xi_m)\right|\cdot w(\xi_1,\dots,\xi_m).
\end{align*}
In $\rho_h$ the supremum over the first argument is formed separately for every degree $m$, and the sum over $m$ is taken afterwards. For fixed $w$ both expressions are subadditive in $\mathbf{F}$.

Suppose that
\[
\mathbf{F}=\sum_{m\ge 2}\frac{1}{m!}\sum_{\xi_1,\dots,\xi_m\in\mathrm{Idx}}\mathcal{K}_m(\xi_1,\dots,\xi_m)\cdot\psi(\xi_1)\cdots\psi(\xi_m)
\]
with functions $\mathcal{K}_m\colon\mathrm{Idx}^m\to\mathbb{C}$ that are not necessarily antisymmetric. For every permutation $\pi$ of $\{1,\dots,m\}$ one has $\psi(\xi_{\pi(1)})\cdots\psi(\xi_{\pi(m)})=\operatorname{sgn}(\pi)\cdot\psi(\xi_1)\cdots\psi(\xi_m)$. Substituting $(\xi_{\pi(1)},\dots,\xi_{\pi(m)})$ for the variables of summation $(\xi_1,\dots,\xi_m)$, averaging over $\pi$ and using the uniqueness of the kernels, one obtains
\[
\mathbf{F}_m(\xi_1,\dots,\xi_m)=\frac{1}{m!}\sum_{\pi}\operatorname{sgn}(\pi)\cdot\mathcal{K}_m\left(\xi_{\pi(1)},\dots,\xi_{\pi(m)}\right).
\]
In the summand of $\pi$ the argument $\xi_1$ occupies the slot $\pi^{-1}(1)$ of $\mathcal{K}_m$, while $\xi_2,\dots,\xi_m$ occupy the other slots. As $w$ is symmetric, the triangle inequality gives
\begin{align}
N_h(\mathbf{F};w)&\le\sum_{m\ge 2}\frac{h_G^m}{(m-1)!}\cdot\sum_{\xi_1,\dots,\xi_m\in\mathrm{Idx}}\left|\mathcal{K}_m(\xi_1,\dots,\xi_m)\right|\cdot w(\xi_1,\dots,\xi_m),\notag\\
\rho_h(\mathbf{F};w)&\le\sum_{m\ge 2}\frac{h_G^m}{(m-1)!}\cdot\max_{1\le a\le m}\,\sup_{\xi\in\mathrm{Idx}}\sum_{\substack{\xi_1,\dots,\xi_m\in\mathrm{Idx}\\ \xi_a=\xi}}\left|\mathcal{K}_m(\xi_1,\dots,\xi_m)\right|\cdot w(\xi_1,\dots,\xi_m).\label{eq:four-jets-9}
\end{align}

For a family $(\xi_1,\dots,\xi_m)$ of elements of $\mathrm{Idx}$ and a non-empty subset $P$ of $\mathcal{I}$ let
\begin{gather*}
w_{\boldsymbol{\delta}}(\xi_1,\dots,\xi_m):=\exp\left(\vartheta\cdot\max_{1\le a,b\le m}d(\xi_a,\xi_b)\right),\\
w_P(\xi_1,\dots,\xi_m):=\exp\left(\vartheta\cdot\max_{1\le a\le m}d(\xi_a,P)\right).
\end{gather*}
These are symmetric weights, and $1$ denotes the symmetric weight with the constant value one.

The elements of a tree on $\{1,\dots,n\}$ (Section~\ref{sec:interpolation-forests}) are called its lines, and the degree of a vertex $i\in\{1,\dots,n\}$ in a tree is the number of its lines that contain $i$. For a line $l=\{i,j\}$ put $\Delta_l:=\Delta_{\{i,j\}}$. For $n\ge 1$ and even elements $\mathbf{F}^{(1)},\dots,\mathbf{F}^{(n)}$ of $\mathrm{Gr}$ without constant term, the complex number
\[
\tilde{E}_n\left(\mathbf{F}^{(1)},\dots,\mathbf{F}^{(n)}\right):=\sum_T\int_{[0,1]^T}\mathrm{d}\tau\;\epsilon_0\left(\prod_{l\in T}\Delta_l\,\exp\left(\Delta\left[s^T(\tau)\right]\right)\prod_{i=1}^n\mathbf{F}^{(i)}\left(\psi^{(i)}\right)\right)
\]
is defined; here $T$ ranges over the trees on $\{1,\dots,n\}$, $\tau=(\tau_l)_{l\in T}$, and $s^T(\tau)$ is defined in Section~\ref{sec:interpolation-forests}. The integral belonging to a tree $T$ is called the summand of $T$. For $n=1$ the only tree is the empty set, $[0,1]^{\emptyset}$ is a point of measure one, and $\Delta[s^{\emptyset}]=\Delta^{11}$; since $\pi_1$ maps $\mathbf{F}(\psi^{(1)})$ to $\mathbf{F}$, Lemma~\ref{lem:wick-rule-gram}(c) with $n=1$ gives
\[
\tilde{E}_1(\mathbf{F})=\epsilon_0\left(\exp\left(\Delta_C\right)\mathbf{F}\right).
\]
The map $\tilde{E}_n$ is multilinear, and it is symmetric. For the symmetry, let $\sigma$ be a permutation of $\{1,\dots,n\}$, put $\sigma(l):=\{\sigma(i),\sigma(j)\}$ for $l=\{i,j\}$, and let $\varphi_{\sigma}$ be the automorphism of $\mathrm{Gr}_n$ that sends $\Psi^{(i)}(x)$ to $\Psi^{(\sigma(i))}(x)$ and $\bar{\Psi}^{(i)}(x)$ to $\bar{\Psi}^{(\sigma(i))}(x)$. Then $\epsilon_0\circ\varphi_{\sigma}=\epsilon_0$, $\varphi_{\sigma}\Delta^{ij}=\Delta^{\sigma(i)\sigma(j)}\varphi_{\sigma}$ and $\varphi_{\sigma}\left(\mathbf{F}(\psi^{(i)})\right)=\mathbf{F}\left(\psi^{(\sigma(i))}\right)$. For a tree $T$ the set $\sigma(T):=\{\sigma(l):l\in T\}$ is a tree, and $s^{\sigma(T)}_{\sigma(i)\sigma(j)}(\tau')=s^T_{ij}(\tau)$ if $\tau'_{\sigma(l)}=\tau_l$ for $l\in T$. Applying $\varphi_{\sigma}$ inside $\epsilon_0$, reordering the even factors and substituting $\tau'$ for $\tau$, one finds that the summand of $T$ for $\left(\mathbf{F}^{(1)},\dots,\mathbf{F}^{(n)}\right)$ equals the summand of $\sigma(T)$ for $\left(\mathbf{F}^{(\sigma^{-1}(1))},\dots,\mathbf{F}^{(\sigma^{-1}(n))}\right)$. The sum over $T$ gives $\tilde{E}_n\left(\mathbf{F}^{(1)},\dots,\mathbf{F}^{(n)}\right)=\tilde{E}_n\left(\mathbf{F}^{(\sigma^{-1}(1))},\dots,\mathbf{F}^{(\sigma^{-1}(n))}\right)$.

\begin{lemma}[Tree expansion]\label{lem:tree-expansion}
\begin{enumerate}
\item[(a)] Let $C$ be invertible, and let $V$ be an even element of $\mathrm{Gr}$ without constant term such that $\boldsymbol{\alpha}_0\cdot\rho_h(V;1)\le 1$. Then the series
\[
W:=\sum_{n\ge 1}\frac{1}{n!}\,\tilde{E}_n(V,\dots,V)
\]
converges absolutely, $|W|\le 2\,\#|\mathrm{Idx}|\cdot\rho_h(V;1)$, and
\[
\int\mathrm{d}\mu_C\,\exp(V)=\exp(W).
\]
\item[(b)] Let $n\ge 2$, let $w_1,\dots,w_n$ be symmetric weights, and let $\mathfrak{d}\ge 0$ be a number with the following property: for every tree $T$ on $\{1,\dots,n\}$, all families $\xi^{(1)},\dots,\xi^{(n)}$ of elements of $\mathrm{Idx}$, the family $\xi^{(i)}$ being attached to the vertex $i$, and all selections of an entry $u_l$ of $\xi^{(i)}$ and an entry $u'_l$ of $\xi^{(j)}$, one pair for each line $l$ of $T$, where $i$ and $j$ are the two ends of $l$,
\begin{equation}\label{eq:four-jets-10}
\exp\left(\vartheta\cdot\mathfrak{d}\right)\le\prod_{i=1}^nw_i\left(\xi^{(i)}\right)\cdot\prod_{l\in T}\exp\left(\vartheta\cdot d\left(u_l,u'_l\right)\right).
\end{equation}
Then, for all even elements $\mathbf{F}^{(1)},\dots,\mathbf{F}^{(n)}$ of $\mathrm{Gr}$ without constant term,
\[
\left|\tilde{E}_n\left(\mathbf{F}^{(1)},\dots,\mathbf{F}^{(n)}\right)\right|\le\exp\left(-\vartheta\cdot\mathfrak{d}\right)\cdot(n-2)!\cdot\boldsymbol{\alpha}_{\vartheta}^{n-1}\cdot N_h\left(\mathbf{F}^{(1)};w_1\right)\cdot\prod_{i=2}^n\rho_h\left(\mathbf{F}^{(i)};w_i\right).
\]
Moreover, $\left|\tilde{E}_1(\mathbf{F})\right|\le N_h(\mathbf{F};1)$ for every even element $\mathbf{F}$ of $\mathrm{Gr}$ without constant term.
\item[(c)] Let $C$ be invertible and $U$ an open subset of $\mathbb{C}^{\mathrm{Cel}}$. For $t\in U$ let $V(t)$ be an even element of $\mathrm{Gr}$ without constant term; assume that the kernels of $V(t)$ are analytic functions of $t$ on $U$ and that
\[
\boldsymbol{\alpha}_{\vartheta}\cdot\rho_h\left(V(t);w_{\boldsymbol{\delta}}\right)\le\frac{1}{2}\qquad(t\in U).
\]
For $c,c'\in\mathrm{Cel}$ let $\partial_c:=\partial/\partial t_c$, and let $\partial_cV(t)$ and $\partial_c\partial_{c'}V(t)$ be the elements of $\mathrm{Gr}$ whose kernels arise from those of $V(t)$ by applying $\partial_c$, resp. $\partial_c\partial_{c'}$. Then (a) applies to $V(t)$ for every $t\in U$, the resulting number $W(t)$ is an analytic function of $t$ on $U$, and for all $c\ne c'$ in $\mathrm{Cel}$, all non-empty subsets $P$, $P'$ of $\mathcal{I}$ and all $t\in U$
\begin{multline}\label{eq:four-jets-11}
\left|\partial_c\partial_{c'}W(t)\right|\le 2\,N_h\left(\partial_c\partial_{c'}V(t);1\right)\\
+2\,\boldsymbol{\alpha}_{\vartheta}\cdot N_h\left(\partial_cV(t);w_P\right)\cdot\rho_h\left(\partial_{c'}V(t);w_{P'}\right)\cdot\exp\left(-\vartheta\cdot d(P,P')\right).
\end{multline}
\end{enumerate}
\end{lemma}

\begin{proof}
\emph{(b).} Let $T$ be a tree on $\{1,\dots,n\}$ and $\tau\in[0,1]^T$. By Lemma~\ref{lem:forest-formula-brydges}(b) there are unit vectors $\boldsymbol{\omega}_1,\dots,\boldsymbol{\omega}_n$ of a Hilbert space with $s^T_{ij}(\tau)=\langle\boldsymbol{\omega}_i,\boldsymbol{\omega}_j\rangle$ for $i\ne j$; as $\langle\boldsymbol{\omega}_i,\boldsymbol{\omega}_i\rangle=1$, Lemma~\ref{lem:wick-rule-gram}(e) applies to $s=s^T(\tau)$. Together with $|C|_{\ell^2\to\ell^2}\le g^2$ it gives, for every monomial $\mathcal{M}$ of degree $p$ in the generators of $\mathrm{Gr}_n$,
\begin{equation}\tag{$\ast$}
\left|\epsilon_0\left(\exp\left(\Delta\left[s^T(\tau)\right]\right)\mathcal{M}\right)\right|\le g^p;
\end{equation}
for odd $p$ the left side vanishes, because $\Delta[s^T(\tau)]$ lowers the degree by two. For $n=1$ and the empty tree, ($\ast$) follows in the same way from Lemma~\ref{lem:wick-rule-gram}(e) with a single unit vector.

Define $\mathrm{Ch}\colon\mathrm{Idx}\times\mathrm{Idx}\to\mathbb{C}$ by $\mathrm{Ch}\left((x,-),(y,+)\right):=C(x,y)$ and $\mathrm{Ch}\left((y,+),(x,-)\right):=-C(x,y)$ for $x,y\in\mathcal{I}$, and by $\mathrm{Ch}(u,u'):=0$ if the second components of $u$ and $u'$ coincide. Then, for $i\ne j$,
\[
\Delta_{\{i,j\}}=\sum_{u,u'\in\mathrm{Idx}}\mathrm{Ch}(u,u')\cdot d_{\psi^{(j)}(u')}\,d_{\psi^{(i)}(u)}.
\]
Indeed, the summands with $u=(x,-)$, $u'=(y,+)$ add up to $\Delta^{ij}$; and since left derivatives anticommute,
\[
\Delta^{ji}=\sum_{x,y\in\mathcal{I}}C(x,y)\cdot d_{\bar{\Psi}^{(i)}(y)}\,d_{\Psi^{(j)}(x)}=-\sum_{x,y\in\mathcal{I}}C(x,y)\cdot d_{\Psi^{(j)}(x)}\,d_{\bar{\Psi}^{(i)}(y)}
\]
is the sum of the summands with $u=(y,+)$, $u'=(x,-)$. By the choice of $\boldsymbol{\alpha}_{\vartheta}$, for every $u\in\mathrm{Idx}$
\begin{equation}\tag{$\ast\ast$}
\sum_{u'\in\mathrm{Idx}}\left|\mathrm{Ch}(u,u')\right|\cdot\exp\left(\vartheta\cdot d(u,u')\right)\le\boldsymbol{\alpha}_{\vartheta};
\end{equation}
for $u=(x,-)$ the left side equals $\sum_{y\in\mathcal{I}}|C(x,y)|\cdot\exp\left(\vartheta\cdot d(x,y)\right)$, and for $u=(y,+)$ it equals $\sum_{x\in\mathcal{I}}|C(x,y)|\cdot\exp\left(\vartheta\cdot d(x,y)\right)$.

Let $\mathbf{F}$ be an even element of $\mathrm{Gr}$ without constant term, let $m\ge 2$ and $u\in\mathrm{Idx}$. Since $d_{\psi(u)}$ is an anti-derivation,
\[
d_{\psi(u)}\left(\psi(\xi_1)\cdots\psi(\xi_m)\right)=\sum_{a=1}^m(-1)^{a-1}\delta_{u\xi_a}\cdot\psi(\xi_1)\cdots\widehat{\psi(\xi_a)}\cdots\psi(\xi_m),
\]
where $\delta_{u\xi_a}$ is Kronecker's symbol and the hat marks the omitted factor; by antisymmetry, $\mathbf{F}_m(\xi_1,\dots,\xi_m)=(-1)^{a-1}\mathbf{F}_m(\xi_a,\xi_1,\dots,\widehat{\xi_a},\dots,\xi_m)$. Hence the $m$ values of $a$ contribute equally, and
\[
d_{\psi(u)}\,\frac{1}{m!}\sum_{\xi_1,\dots,\xi_m}\mathbf{F}_m(\xi_1,\dots,\xi_m)\cdot\psi(\xi_1)\cdots\psi(\xi_m)=\frac{1}{(m-1)!}\sum_{\xi_2,\dots,\xi_m}\mathbf{F}_m(u,\xi_2,\dots,\xi_m)\cdot\psi(\xi_2)\cdots\psi(\xi_m).
\]
The right side has the same form, with the antisymmetric function $\mathbf{F}_m(u,\cdot)$ of $m-1$ arguments. By iteration, for $1\le k\le m$ and $u_1,\dots,u_k\in\mathrm{Idx}$, the application of $d_{\psi(u_1)}$, then of $d_{\psi(u_2)}$, and so on up to $d_{\psi(u_k)}$, to the part of degree $m$ of $\mathbf{F}$ yields
\[
\frac{1}{(m-k)!}\sum_{\xi_{k+1},\dots,\xi_m}\mathbf{F}_m(u_1,\dots,u_k,\xi_{k+1},\dots,\xi_m)\cdot\psi(\xi_{k+1})\cdots\psi(\xi_m),
\]
whereas more than $m$ derivatives annihilate this part. The same is true in every copy, with $\psi^{(i)}$ in place of $\psi$.

Now let $n\ge 2$, let $T$ be a tree on $\{1,\dots,n\}$, and let $k_i\ge 1$ be the degree of the vertex $i$ in $T$. The vertex $1$ is taken as the root: for a line $l$ of $T$ let $i_l$ be the end of $l$ nearer to $1$ and $j_l$ the other end, and write $\Delta_l$ as above, with $i=i_l$, $j=j_l$ and the variables of summation $u_l$, $u'_l$. Every vertex $j\ne 1$ equals $j_l$ for exactly one line $l=:l(j)$, which joins $j$ to its parent $i_{l(j)}$; the vertices whose parent is $i$ are the children of $i$. The operators $\Delta^{ij}$ commute; hence $\prod_{l\in T}\Delta_l$ may be applied before $\exp\left(\Delta\left[s^T(\tau)\right]\right)$. Decompose every $\mathbf{F}^{(i)}$ into its homogeneous parts, of the degrees $m_i$, and expand $\prod_{l\in T}\Delta_l$ into the sum over the values of all $u_l$, $u'_l$. A derivative with respect to a generator of the copy $i$ annihilates the elements of the subalgebra generated by the other copies; so it acts on the factor of the copy $i$ only, and it commutes, up to a sign, with the multiplication by the factors that belong to other copies. By the preceding paragraph, the $k_i$ derivatives with respect to generators of the copy $i$ annihilate the part of degree $m_i$ of $\mathbf{F}^{(i)}(\psi^{(i)})$ if $m_i<k_i$, and otherwise transform it into
\[
\pm\frac{1}{(m_i-k_i)!}\sum_{\xi^{(i)}_{k_i+1},\dots,\xi^{(i)}_{m_i}\in\mathrm{Idx}}\mathbf{F}^{(i)}_{m_i}\left(\xi^{(i)}\right)\cdot\psi^{(i)}\left(\xi^{(i)}_{k_i+1}\right)\cdots\psi^{(i)}\left(\xi^{(i)}_{m_i}\right),
\]
where $\xi^{(i)}=\left(\xi^{(i)}_1,\dots,\xi^{(i)}_{m_i}\right)$, and $\xi^{(i)}_1,\dots,\xi^{(i)}_{k_i}$ are the arguments at $i$ of the lines containing $i$, namely the $u_l$ with $i_l=i$ and the $u'_l$ with $j_l=i$, in the order in which the derivatives act. The product over $i$ of the monomials that occur is zero or, up to a sign, a monomial of degree $\sum_i(m_i-k_i)$ in the generators of $\mathrm{Gr}_n$. By~($\ast$), and as $[0,1]^T$ has measure one, the summand of $T$ in $\tilde{E}_n\left(\mathbf{F}^{(1)},\dots,\mathbf{F}^{(n)}\right)$ has modulus at most the sum, over all $(m_1,\dots,m_n)$ with $m_i\ge\max\{2,k_i\}$ and over all values in $\mathrm{Idx}$ of all entries of the families $\xi^{(1)},\dots,\xi^{(n)}$, of
\[
\prod_{i=1}^n\left(\frac{g^{m_i-k_i}}{(m_i-k_i)!}\cdot\left|\mathbf{F}^{(i)}_{m_i}\left(\xi^{(i)}\right)\right|\right)\cdot\prod_{l\in T}\left|\mathrm{Ch}\left(u_l,u'_l\right)\right|.
\]
Each of these terms is multiplied by
\[
\exp\left(-\vartheta\cdot\mathfrak{d}\right)\cdot\prod_{i=1}^nw_i\left(\xi^{(i)}\right)\cdot\prod_{l\in T}\exp\left(\vartheta\cdot d\left(u_l,u'_l\right)\right),
\]
which is at least $1$ by~\eqref{eq:four-jets-10}, applied with the entry $u_l$ of $\xi^{(i_l)}$ and the entry $u'_l$ of $\xi^{(j_l)}$ for every line $l$.

For fixed $(m_i)$ the sum over the arguments is carried out from the leaves of $T$ towards the root. For $m\ge 2$ and $j\in\{1,\dots,n\}$ put
\begin{gather*}
Q(m):=\sum_{\xi_1,\dots,\xi_m\in\mathrm{Idx}}\left|\mathbf{F}^{(1)}_m(\xi_1,\dots,\xi_m)\right|\cdot w_1(\xi_1,\dots,\xi_m),\\
R_j(m):=\sup_{\xi_1\in\mathrm{Idx}}\sum_{\xi_2,\dots,\xi_m\in\mathrm{Idx}}\left|\mathbf{F}^{(j)}_m(\xi_1,\dots,\xi_m)\right|\cdot w_j(\xi_1,\dots,\xi_m).
\end{gather*}
For a vertex $i\ne 1$ let $D_i$ be the set formed by $i$ and by the vertices whose path to $1$ contains $i$, and for $v\in\mathrm{Idx}$ let $S_i(v)$ be the sum, over all values of the entries of the families $\xi^{(j)}$, $j\in D_i$, with $u'_{l(i)}=v$ held fixed, of
\[
\prod_{j\in D_i}\left(\left|\mathbf{F}^{(j)}_{m_j}\left(\xi^{(j)}\right)\right|\cdot w_j\left(\xi^{(j)}\right)\right)\cdot\prod_{l\in T,\ l\subset D_i}\left(\left|\mathrm{Ch}\left(u_l,u'_l\right)\right|\cdot\exp\left(\vartheta\cdot d\left(u_l,u'_l\right)\right)\right).
\]
Then
\[
S_i(v)\le\boldsymbol{\alpha}_{\vartheta}^{\#|D_i|-1}\cdot\prod_{j\in D_i}R_j(m_j)\qquad\text{for every }v\in\mathrm{Idx}.
\]
This is proved by induction on $\#|D_i|$; the induction hypothesis applies to the children $j$ of $i$, since $D_j\subset D_i\setminus\{i\}$. The lines $l\subset D_i$ are the lines $l(j)$ of the children $j$ of $i$ and the lines contained in the sets $D_j$, and the sets $D_j$ of distinct children are disjoint. Hence $S_i(v)$ is the sum, over the entries of $\xi^{(i)}$ other than $u'_{l(i)}=v$, of $\left|\mathbf{F}^{(i)}_{m_i}\left(\xi^{(i)}\right)\right|\cdot w_i\left(\xi^{(i)}\right)$ times the product, over the children $j$ of $i$, of
\[
\sum_{v_j\in\mathrm{Idx}}\left|\mathrm{Ch}\left(u_{l(j)},v_j\right)\right|\cdot\exp\left(\vartheta\cdot d\left(u_{l(j)},v_j\right)\right)\cdot S_j(v_j),
\]
where $u_{l(j)}$ is an entry of $\xi^{(i)}$; for a leaf $i$ this product is empty and equal to one. By the induction hypothesis and by~($\ast\ast$) with $u=u_{l(j)}$, each factor of the product is at most $\boldsymbol{\alpha}_{\vartheta}^{\#|D_j|}\cdot\prod_{j'\in D_j}R_{j'}(m_{j'})$, whatever the value of $u_{l(j)}$. The remaining sum over the entries of $\xi^{(i)}$ other than $v$ is at most $R_i(m_i)$, because $\left|\mathbf{F}^{(i)}_{m_i}\right|$ and $w_i$ are symmetric functions of their arguments. As $\#|D_i|-1$ is the sum of the numbers $\#|D_j|$ over the children $j$ of $i$, the bound for $S_i$ follows. At the root, the same reasoning with all entries of $\xi^{(1)}$ summed bounds the sum over all arguments by $\boldsymbol{\alpha}_{\vartheta}^{n-1}\cdot Q(m_1)\cdot\prod_{i=2}^nR_i(m_i)$. Consequently the summand of $T$ in $\tilde{E}_n\left(\mathbf{F}^{(1)},\dots,\mathbf{F}^{(n)}\right)$ has modulus at most
\[
\exp\left(-\vartheta\cdot\mathfrak{d}\right)\cdot\boldsymbol{\alpha}_{\vartheta}^{n-1}\cdot\sum_{(m_i)}\frac{g^{m_1-k_1}}{(m_1-k_1)!}\cdot Q(m_1)\cdot\prod_{i=2}^n\frac{g^{m_i-k_i}}{(m_i-k_i)!}\cdot R_i(m_i),
\]
the sum over all $(m_1,\dots,m_n)$ with $m_i\ge\max\{2,k_i\}$; this bound depends on $T$ only through $k_1,\dots,k_n$.

By Cayley's formula, for integers $k_1,\dots,k_n\ge 1$ the number of trees on $\{1,\dots,n\}$ in which every vertex $i$ has the degree $k_i$ is
\[
\frac{(n-2)!}{(k_1-1)!\cdots(k_n-1)!}\quad\text{if }k_1+\dots+k_n=2n-2,
\]
and zero otherwise. Sum the last bound over all trees, enlarge the sum to all sequences $(k_1,\dots,k_n)$ of integers $k_i\ge 1$, which only adds non-negative terms, and carry out the sum over each $k_i$ before the sum over $m_i$. By the binomial theorem,
\[
\sum_{k=1}^m\frac{g^{m-k}}{(k-1)!\,(m-k)!}=\frac{(1+g)^{m-1}}{(m-1)!}\le\frac{h_G^m}{(m-1)!}\qquad(m\ge 1),
\]
the inequality because $1\le 1+g\le h_G$. One obtains
\begin{multline*}
\left|\tilde{E}_n\left(\mathbf{F}^{(1)},\dots,\mathbf{F}^{(n)}\right)\right|\\
\le\exp\left(-\vartheta\cdot\mathfrak{d}\right)\cdot(n-2)!\cdot\boldsymbol{\alpha}_{\vartheta}^{n-1}\cdot\left(\sum_{m\ge 2}\frac{h_G^m}{(m-1)!}\cdot Q(m)\right)\cdot\prod_{i=2}^n\left(\sum_{m\ge 2}\frac{h_G^m}{(m-1)!}\cdot R_i(m)\right),
\end{multline*}
and by the definitions of $N_h$ and $\rho_h$ the factors in parentheses are $N_h\left(\mathbf{F}^{(1)};w_1\right)$ and $\rho_h\left(\mathbf{F}^{(i)};w_i\right)$. This is the first assertion of (b). For $n=1$, the definition of $\tilde{E}_1$ and~($\ast$) for $n=1$, applied to the monomials $\psi^{(1)}(\xi_1)\cdots\psi^{(1)}(\xi_m)$, give
\[
\left|\tilde{E}_1(\mathbf{F})\right|\le\sum_{m\ge 2}\frac{g^m}{m!}\cdot\sum_{\xi_1,\dots,\xi_m\in\mathrm{Idx}}\left|\mathbf{F}_m(\xi_1,\dots,\xi_m)\right|\le N_h(\mathbf{F};1),
\]
since $\frac{g^m}{m!}\le\frac{h_G^m}{(m-1)!}$.

\emph{(a).} Part~(b) is valid for every rate $\vartheta\ge 0$ together with a corresponding number $\boldsymbol{\alpha}_{\vartheta}$. It is applied with the rate $0$ and the number $\boldsymbol{\alpha}_0$, with all weights equal to $1$ and with $\mathfrak{d}=0$; then~\eqref{eq:four-jets-10} reads $1\le 1$. Since $N_h(V;1)\le\#|\mathrm{Idx}|\cdot\rho_h(V;1)$ and $\frac{(n-2)!}{n!}=\frac{1}{n(n-1)}$, one obtains
\begin{gather*}
\left|\frac{1}{n!}\,\tilde{E}_n(V,\dots,V)\right|\le\#|\mathrm{Idx}|\cdot\rho_h(V;1)\cdot\frac{\left(\boldsymbol{\alpha}_0\cdot\rho_h(V;1)\right)^{n-1}}{n(n-1)}\qquad(n\ge 2),\\
\left|\tilde{E}_1(V)\right|\le\#|\mathrm{Idx}|\cdot\rho_h(V;1).
\end{gather*}
As $\boldsymbol{\alpha}_0\cdot\rho_h(V;1)\le 1$, the power series
\[
\mathfrak{w}(\mathfrak{l}):=\sum_{n\ge 1}\frac{\mathfrak{l}^n}{n!}\,\tilde{E}_n(V,\dots,V)
\]
converges absolutely for $|\mathfrak{l}|\le 1$. Thus $\mathfrak{w}$ is analytic in the disc $|\mathfrak{l}|<1$ and continuous on $|\mathfrak{l}|\le 1$, $\mathfrak{w}(0)=0$, and $W=\mathfrak{w}(1)$ satisfies $|W|\le 2\,\#|\mathrm{Idx}|\cdot\rho_h(V;1)$, because $\sum_{n\ge 2}\frac{1}{n(n-1)}=1$.

For $s=(s_l)_{l\in P_n}$ let
\[
H_n(s):=\epsilon_0\left(\exp\left(\Delta[s]\right)\prod_{i=1}^nV\left(\psi^{(i)}\right)\right).
\]
The exponential $\exp(\Delta[s])$ is a finite sum, and the operators $\Delta^{ij}$ commute; therefore $H_n$ is a polynomial in $s$, one has $\partial_{s_l}\exp(\Delta[s])=\Delta_l\exp(\Delta[s])$, and for every forest $\mathrm{Fo}$ on $\{1,\dots,n\}$
\[
\left(\prod_{l\in\mathrm{Fo}}\partial_{s_l}H_n\right)(s)=\epsilon_0\left(\prod_{l\in\mathrm{Fo}}\Delta_l\,\exp\left(\Delta[s]\right)\prod_{i=1}^nV\left(\psi^{(i)}\right)\right).
\]
The homomorphism $\pi_n$ maps $\prod_{i=1}^nV(\psi^{(i)})$ to $V^n$. Hence Lemma~\ref{lem:wick-rule-gram}(a), (c) and Lemma~\ref{lem:forest-formula-brydges}(a) give
\begin{align*}
\int\mathrm{d}\mu_C\,V^n&=\epsilon_0\left(\exp\left(\Delta_C\right)V^n\right)=H_n(1,\dots,1)\\
&=\sum_{\mathrm{Fo}}\int_{[0,1]^{\mathrm{Fo}}}\mathrm{d}\tau\;\epsilon_0\left(\prod_{l\in\mathrm{Fo}}\Delta_l\,\exp\left(\Delta\left[s^{\mathrm{Fo}}(\tau)\right]\right)\prod_{i=1}^nV\left(\psi^{(i)}\right)\right),
\end{align*}
the sum over all forests $\mathrm{Fo}$ on $\{1,\dots,n\}$. Fix a forest $\mathrm{Fo}$, let $B_1,\dots,B_r$ be the vertex sets of its connected components, and let $T_a$ be the set of the lines of $\mathrm{Fo}$ that are contained in $B_a$; it is a tree on $B_a$. If $i$ and $j$ lie in different components, then $s^{\mathrm{Fo}}_{ij}(\tau)=0$. Hence $\Delta\left[s^{\mathrm{Fo}}(\tau)\right]=\sum_{a=1}^r\Delta_{(a)}$ with
\[
\Delta_{(a)}:=\sum_{i\in B_a}\Delta^{ii}+\sum_{i,j\in B_a,\ i<j}s^{\mathrm{Fo}}_{ij}(\tau)\cdot\Delta_{\{i,j\}};
\]
for $i,j\in B_a$ the path from $i$ to $j$ in $\mathrm{Fo}$ consists of lines of $T_a$, so $s^{\mathrm{Fo}}_{ij}(\tau)$ equals $s^{T_a}_{ij}$ evaluated at $(\tau_l)_{l\in T_a}$. Consequently the integrand is $\epsilon_0\left(\prod_aO_a\prod_aF_a\right)$, where $O_a:=\prod_{l\in T_a}\Delta_l\,\exp\left(\Delta_{(a)}\right)$ is a polynomial in the operators $\Delta^{ij}$, $i,j\in B_a$, and $F_a:=\prod_{i\in B_a}V\left(\psi^{(i)}\right)$ is an even element of the subalgebra generated by the copies $i\in B_a$. By Lemma~\ref{lem:wick-rule-gram}(d) the integrand equals $\prod_a\epsilon_0\left(O_aF_a\right)$, and by Fubini's theorem the integral over $[0,1]^{\mathrm{Fo}}=\prod_a[0,1]^{T_a}$ is the product over $a$ of the integrals over $[0,1]^{T_a}$. The factor of $a$ is the summand of the tree $T_a$ in $\tilde{E}_{\#|B_a|}(V,\dots,V)$, the copies being numbered by the elements of $B_a$ instead of $1,\dots,\#|B_a|$; such a renumbering does not change the value, by the argument given above for the symmetry of $\tilde{E}_n$. (If $B_a=\{i\}$, then $T_a=\emptyset$, and the factor is $\epsilon_0\left(\exp\left(\Delta^{ii}\right)V\left(\psi^{(i)}\right)\right)=\tilde{E}_1(V)$.) The forests whose components have the vertex sets $B_1,\dots,B_r$ correspond to the choices of trees $T_a$ on the sets $B_a$. Summing over them, one finds
\[
\int\mathrm{d}\mu_C\,V^n=\sum_{\{B_1,\dots,B_r\}}\ \prod_{a=1}^r\tilde{E}_{\#|B_a|}(V,\dots,V),
\]
the sum over all partitions of $\{1,\dots,n\}$ into non-empty blocks. The number of ordered $r$-tuples $(B_1,\dots,B_r)$ of pairwise disjoint non-empty sets with the union $\{1,\dots,n\}$ and the sizes $\#|B_a|=m_a$ is $\frac{n!}{m_1!\cdots m_r!}$, and every partition into $r$ blocks arises from exactly $r!$ of them. Therefore, for $n\ge 1$,
\[
\frac{1}{n!}\int\mathrm{d}\mu_C\,V^n=\sum_{r\ge 1}\frac{1}{r!}\sum_{\substack{m_1,\dots,m_r\ge 1\\ m_1+\dots+m_r=n}}\ \prod_{a=1}^r\frac{\tilde{E}_{m_a}(V,\dots,V)}{m_a!},
\]
which is the coefficient of $\mathfrak{l}^n$ in the formal power series $\exp\left(\mathfrak{w}(\mathfrak{l})\right)=\sum_{r\ge 0}\frac{\mathfrak{w}(\mathfrak{l})^r}{r!}$; the constant coefficients agree as well, since $\int\mathrm{d}\mu_C\,1=\epsilon_0(1)=1$ by Lemma~\ref{lem:wick-rule-gram}(a). As $V$ has no constant term, $V^n$ is a sum of monomials of degree at least $2n$, so that $V^n=0$ for $n>\#|\mathcal{I}|$, and
\[
\mathfrak{Z}(\mathfrak{l}):=\int\mathrm{d}\mu_C\,\exp\left(\mathfrak{l} V\right)=\sum_{n\ge 0}\frac{\mathfrak{l}^n}{n!}\int\mathrm{d}\mu_C\,V^n
\]
is a polynomial in $\mathfrak{l}$. The Taylor series at $0$ of the function $\exp\circ\mathfrak{w}$, which is analytic in the disc $|\mathfrak{l}|<1$, is the formal series $\exp\left(\mathfrak{w}(\mathfrak{l})\right)$. So the Taylor coefficients of $\mathfrak{Z}$ and of $\exp\circ\mathfrak{w}$ at $0$ coincide; by the identity theorem they coincide in the disc $|\mathfrak{l}|<1$, and by continuity on $|\mathfrak{l}|\le 1$. For $\mathfrak{l}=1$ this is the identity $\int\mathrm{d}\mu_C\,\exp(V)=\exp(W)$.

\emph{(c).} Let $t\in U$. Since $w_{\boldsymbol{\delta}}\ge 1$ and $\boldsymbol{\alpha}_0\le\boldsymbol{\alpha}_{\vartheta}$, one has $\boldsymbol{\alpha}_0\cdot\rho_h\left(V(t);1\right)\le\boldsymbol{\alpha}_{\vartheta}\cdot\rho_h\left(V(t);w_{\boldsymbol{\delta}}\right)\le\frac{1}{2}$; so (a) applies, and $W(t)=\sum_{n\ge 1}\frac{1}{n!}\,\tilde{E}_n\left(V(t),\dots,V(t)\right)$. Each term of this series is a polynomial in the values of the kernels of $V(t)$, hence an analytic function on $U$. As $C$ is invertible, $C\ne 0$, so $\boldsymbol{\alpha}_{\vartheta}\ge\sup_x\sum_y|C(x,y)|>0$ and $\rho_h\left(V(t);1\right)\le\frac{1}{2\boldsymbol{\alpha}_{\vartheta}}$ on $U$. The bounds of the proof of (a) now give
\[
\left|\frac{1}{n!}\,\tilde{E}_n\left(V(t),\dots,V(t)\right)\right|\le\frac{\#|\mathrm{Idx}|}{2\boldsymbol{\alpha}_{\vartheta}}\cdot\frac{2^{1-n}}{n(n-1)}\qquad(n\ge 2,\ t\in U),
\]
so that the series converges uniformly on $U$. By Weierstrass's theorem on uniform limits of holomorphic functions, $W$ is analytic on $U$, and its partial derivatives are obtained by differentiating the series term by term.

From now on $t\in U$ is fixed and omitted from the notation; $\partial_cV$ and $\partial_c\partial_{c'}V$ are even elements of $\mathrm{Gr}$ without constant term. Since $\tilde{E}_n$ is multilinear and symmetric, the product rule gives
\[
\partial_c\partial_{c'}W=\sum_{n\ge 1}\frac{1}{n!}\left(n\,\tilde{E}_n\left(\partial_c\partial_{c'}V,V,\dots,V\right)+n(n-1)\,\tilde{E}_n\left(\partial_cV,\partial_{c'}V,V,\dots,V\right)\right).
\]

For the terms of the first kind with $n\ge 2$, (b) is applied with the rate $\vartheta$, all weights equal to $1$ and $\mathfrak{d}=0$; then~\eqref{eq:four-jets-10} holds, as its right side is at least $1$. Together with $\rho_h(V;1)\le\rho_h(V;w_{\boldsymbol{\delta}})$ this gives
\begin{align*}
\left|\frac{n}{n!}\,\tilde{E}_n\left(\partial_c\partial_{c'}V,V,\dots,V\right)\right|&\le N_h\left(\partial_c\partial_{c'}V;1\right)\cdot\frac{\left(\boldsymbol{\alpha}_{\vartheta}\cdot\rho_h(V;w_{\boldsymbol{\delta}})\right)^{n-1}}{n-1}\\
&\le N_h\left(\partial_c\partial_{c'}V;1\right)\cdot\frac{2^{1-n}}{n-1}.
\end{align*}
By the last assertion of (b), the term of the first kind with $n=1$ is at most $N_h\left(\partial_c\partial_{c'}V;1\right)$. Since $1+\sum_{i\ge 1}\frac{2^{-i}}{i}=1+\log 2\le 2$, the terms of the first kind contribute at most $2\,N_h\left(\partial_c\partial_{c'}V;1\right)$.

The terms of the second kind occur for $n\ge 2$. In them the factor $\partial_cV$ sits at the vertex $1$ and the factor $\partial_{c'}V$ at the vertex $2$; these two factors are the marked vertices of the expansion in the sense of Definition~\ref{def:terminology}(8). Part~(b) is applied with $w_1=w_P$, $w_2=w_{P'}$, $w_i=w_{\boldsymbol{\delta}}$ for $3\le i\le n$, and $\mathfrak{d}=d(P,P')$. To verify~\eqref{eq:four-jets-10}, let $T$ be a tree on $\{1,\dots,n\}$, let $1=i_0,i_1,\dots,i_q=2$ be the path from $1$ to $2$ in $T$, let $l_a$ be the line joining $i_{a-1}$ and $i_a$, and let $u_a$ and $u'_a$ be the arguments of $l_a$ at $i_{a-1}$ and at $i_a$, respectively ($1\le a\le q$). Choosing points of $P$ and of $P'$ at which $d(u_1,P)$ and $d(u'_q,P')$ are attained, and applying the triangle inequality, one gets
\[
d(P,P')\le d(u_1,P)+\sum_{a=1}^qd\left(u_a,u'_a\right)+\sum_{a=1}^{q-1}d\left(u'_a,u_{a+1}\right)+d\left(u'_q,P'\right).
\]
Here $\exp\left(\vartheta\cdot d(u_1,P)\right)\le w_P\left(\xi^{(1)}\right)$, as $u_1$ is an entry of $\xi^{(1)}$; likewise $\exp\left(\vartheta\cdot d(u'_q,P')\right)\le w_{P'}\left(\xi^{(2)}\right)$; and for $1\le a\le q-1$ the vertex $i_a$ differs from $1$ and $2$, and $u'_a$, $u_{a+1}$ are entries of $\xi^{(i_a)}$, so that $\exp\left(\vartheta\cdot d(u'_a,u_{a+1})\right)\le w_{\boldsymbol{\delta}}\left(\xi^{(i_a)}\right)$. The factors $\exp\left(\vartheta\cdot d(u_a,u'_a)\right)$ occur on the right side of~\eqref{eq:four-jets-10}, and all its other factors are at least $1$. Hence~\eqref{eq:four-jets-10} holds, and (b) yields
\begin{multline*}
\left|\frac{n(n-1)}{n!}\,\tilde{E}_n\left(\partial_cV,\partial_{c'}V,V,\dots,V\right)\right|\\
\le\exp\left(-\vartheta\cdot d(P,P')\right)\cdot\boldsymbol{\alpha}_{\vartheta}\cdot N_h\left(\partial_cV;w_P\right)\cdot\rho_h\left(\partial_{c'}V;w_{P'}\right)\cdot\left(\boldsymbol{\alpha}_{\vartheta}\cdot\rho_h(V;w_{\boldsymbol{\delta}})\right)^{n-2}.
\end{multline*}
As $\boldsymbol{\alpha}_{\vartheta}\cdot\rho_h(V;w_{\boldsymbol{\delta}})\le\frac{1}{2}$ and $\sum_{n\ge 2}2^{-(n-2)}=2$, the terms of the second kind contribute at most the second summand on the right side of~\eqref{eq:four-jets-11}. Together with the bound for the terms of the first kind, this proves~\eqref{eq:four-jets-11}.
\end{proof}

\subsubsection{Holomorphic functions that are small at the real points}\label{sec:holomorphic-functions-small}

\begin{lemma}[A two-constants estimate]\label{lem:two-constants-estimate}
Let $n \ge 1$ be an integer and $\rho > 0$, and let $\mathfrak{g}$ be a holomorphic function on the open polydisk
\[
\left\{ t \in \mathbb{C}^n : |t_i| < \rho \text{ for } i = 1, \dots, n \right\}.
\]
Let $M_{\mathfrak{g}}$ and $\boldsymbol{\delta}$ be real numbers such that $|\mathfrak{g}| \le M_{\mathfrak{g}}$ on this polydisk and $|\mathfrak{g}(t)| \le \boldsymbol{\delta} \le M_{\mathfrak{g}}$ for every $t \in (-\rho,\rho)^n$. Then
\[
|\mathfrak{g}(t)| \le \boldsymbol{\delta}^{\frac{1}{3}}\cdot M_{\mathfrak{g}}^{\frac{2}{3}} \qquad \text{if } |t_i| \le \frac{\rho}{2n} \text{ for } i = 1, \dots, n .
\]
\end{lemma}

\begin{proof}
\emph{One variable.} Let $0 < q < 1$ and $\boldsymbol{\omega} := \frac{4}{\pi}\arctan q$; then $0 < \boldsymbol{\omega} < 1$. Let $\boldsymbol{\delta}'$ be a number with $0 \le \boldsymbol{\delta}' \le M_{\mathfrak{g}}$, and let $\mathfrak{u}$ be a holomorphic function on the disk $\{t \in \mathbb{C} : |t| < \rho\}$ such that $|\mathfrak{u}| \le M_{\mathfrak{g}}$ on this disk and $|\mathfrak{u}(t)| \le \boldsymbol{\delta}'$ for $t \in (-\rho,\rho)$. It is shown first that
\[
|\mathfrak{u}(t)| \le {\boldsymbol{\delta}'}^{1-\boldsymbol{\omega}}\cdot M_{\mathfrak{g}}^{\boldsymbol{\omega}} \qquad \text{for } |t| \le q\cdot\rho .
\]
If $\boldsymbol{\delta}' = 0$, then $\mathfrak{u}$ vanishes on $(-\rho,\rho)$, hence on the whole disk by the identity theorem, and the inequality holds. Let $\boldsymbol{\delta}' > 0$; then $M_{\mathfrak{g}} \ge \boldsymbol{\delta}' > 0$.

Let $\Omega := \{z \in \mathbb{C} : |z| < 1,\ \operatorname{Im} z > 0\}$ be the open upper half of the unit disk, and put
\[
\hat{\boldsymbol{\omega}}(z) := \frac{2}{\pi}\arg\left(\frac{1+z}{1-z}\right) \qquad (z \in \Omega).
\]
The M\"obius transformation $z \mapsto \frac{1+z}{1-z}$ maps $\Omega$ onto the open first quadrant, the segment $(-1,1)$ onto the positive real half-axis, and the open upper half of the unit circle onto the positive imaginary half-axis. Hence $\hat{\boldsymbol{\omega}}$ is harmonic on $\Omega$, $0 \le \hat{\boldsymbol{\omega}} \le 1$, and $\hat{\boldsymbol{\omega}}(z)$ tends to $0$ as $z$ tends to a point of $(-1,1)$, and to $1$ as $z$ tends to a point of the open upper half of the unit circle.

Let $0 < \rho' < \rho$ and $\tilde{\delta} > 0$, and define on $\Omega$
\[
v(z) := \log\left|\mathfrak{u}(\rho'\cdot z)\right| - \left(1 - \hat{\boldsymbol{\omega}}(z)\right)\cdot\log\boldsymbol{\delta}' - \hat{\boldsymbol{\omega}}(z)\cdot\log M_{\mathfrak{g}} - \tilde{\delta}\cdot\log\frac{4}{\left|1 - z^2\right|}.
\]
Here $\log 0 := -\infty$, so that $v$ takes its values in $[-\infty,\infty)$; the last term of $v(z)$ is finite, because $1 - z^2 = (1+z)\cdot(1-z)$ has no zero in $\Omega$. Let $\operatorname{Log}$ be the principal branch of the logarithm on the half-plane $\{w \in \mathbb{C} : \operatorname{Re} w > 0\}$, that is, $\operatorname{Log} w = \log|w| + i\cdot\arg w$ with $-\frac{\pi}{2} < \arg w < \frac{\pi}{2}$; it is holomorphic on this half-plane. For $z \in \Omega$ the three numbers $\frac{1+z}{1-z}$, $1 + z$ and $1 - z$ lie in this half-plane: the first lies in the open first quadrant, as stated above, and $\operatorname{Re}(1 \pm z) \ge 1 - |z| > 0$. The function $z \mapsto \mathfrak{u}(\rho'\cdot z)$ is holomorphic on the disk $\left\{z \in \mathbb{C} : |z| < \frac{\rho}{\rho'}\right\}$, which contains the closed unit disk. Hence
\[
H(z) := \log\boldsymbol{\delta}' - \frac{2i}{\pi}\cdot\log\frac{M_{\mathfrak{g}}}{\boldsymbol{\delta}'}\cdot\operatorname{Log}\left(\frac{1+z}{1-z}\right) + \tilde{\delta}\cdot\left(\log 4 - \operatorname{Log}(1+z) - \operatorname{Log}(1-z)\right)
\]
and $G(z) := \mathfrak{u}(\rho'\cdot z)\cdot\exp\left(-H(z)\right)$ define holomorphic functions $H$ and $G$ on $\Omega$. One has $\operatorname{Im}\operatorname{Log}\left(\frac{1+z}{1-z}\right) = \arg\left(\frac{1+z}{1-z}\right) = \frac{\pi}{2}\cdot\hat{\boldsymbol{\omega}}(z)$ and $\operatorname{Re}\left(\operatorname{Log}(1+z) + \operatorname{Log}(1-z)\right) = \log|1+z| + \log|1-z| = \log\left|1 - z^2\right|$; therefore
\[
\operatorname{Re} H(z) = \left(1 - \hat{\boldsymbol{\omega}}(z)\right)\cdot\log\boldsymbol{\delta}' + \hat{\boldsymbol{\omega}}(z)\cdot\log M_{\mathfrak{g}} + \tilde{\delta}\cdot\log\frac{4}{\left|1 - z^2\right|} \qquad (z \in \Omega),
\]
so that $v(z) = \log\left|\mathfrak{u}(\rho'\cdot z)\right| - \operatorname{Re} H(z)$. As $\left|\exp\left(-H(z)\right)\right| = \exp\left(-\operatorname{Re} H(z)\right)$, it follows that
\[
|G(z)| = \exp\left(v(z)\right) \qquad (z \in \Omega),
\]
where $\exp(-\infty) := 0$; the value $-\infty$ is taken by $v$ exactly at the zeros of $z \mapsto \mathfrak{u}(\rho'\cdot z)$. The upper limit of $v$ at every boundary point of $\Omega$ is at most $0$. Indeed, the function $z \mapsto \mathfrak{u}(\rho'\cdot z)$ is continuous on the closed unit disk, since $\rho' < \rho$; and $\log\frac{4}{|1-z^2|} \ge \log 2 > 0$ for $|z| \le 1$, $z \ne \pm 1$, since $|1 - z^2| \le 2$ there. At a point $z_0 \in (-1,1)$ the upper limit is at most $\log|\mathfrak{u}(\rho'\cdot z_0)| - \log\boldsymbol{\delta}' \le 0$, and at a point $z_0$ of the open upper half of the unit circle it is at most $\log|\mathfrak{u}(\rho'\cdot z_0)| - \log M_{\mathfrak{g}} \le 0$. At $z_0 = \pm 1$ it equals $-\infty$: the sum of the first three terms of $v(z)$ is
\[
\left(1 - \hat{\boldsymbol{\omega}}(z)\right)\cdot\left(\log\left|\mathfrak{u}(\rho'\cdot z)\right| - \log\boldsymbol{\delta}'\right) + \hat{\boldsymbol{\omega}}(z)\cdot\left(\log\left|\mathfrak{u}(\rho'\cdot z)\right| - \log M_{\mathfrak{g}}\right) \le \log\frac{M_{\mathfrak{g}}}{\boldsymbol{\delta}'},
\]
whereas $\log\frac{4}{|1-z^2|}$ tends to $+\infty$ as $z$ tends to $\pm 1$. Since $|G| = \exp(v)$, and $\exp$ is continuous and increasing on $[-\infty,\infty)$, it follows that
\[
\limsup_{z \to \zeta,\ z \in \Omega} |G(z)| \le 1 \qquad \text{for every boundary point } \zeta \text{ of } \Omega.
\]
This implies $|G| \le 1$ on $\Omega$. Suppose, to the contrary, that $c := |G(z_1)| > 1$ for a point $z_1 \in \Omega$, and let $K := \{z \in \Omega : |G(z)| \ge c\}$. The set $K$ is non-empty and bounded. It is closed in $\mathbb{C}$: let a sequence $(\zeta_k)_{k \ge 1}$ of points of $K$ converge to $\zeta \in \mathbb{C}$; then $\zeta$ lies in the closure of $\Omega$, and $\zeta$ is not a boundary point of $\Omega$, since otherwise the upper limit of $|G(z)|$, as $z \in \Omega$ tends to $\zeta$, would be at least $c > 1$; so $\zeta \in \Omega$, and $|G(\zeta)| \ge c$ by continuity. Hence $K$ is compact, and the continuous function $|G|$ attains on $K$ a maximum $m \ge c$. As $|G| < c$ on $\Omega \setminus K$, the number $m$ is the maximum of $|G|$ on $\Omega$. The set $K' := \{z \in \Omega : |G(z)| = m\}$ is a non-empty closed subset of $K$, hence compact. Let $a$ be a point of $K'$ with the largest real part, and let $r > 0$ be such that the closed disk of radius $r$ about $a$ lies in $\Omega$. Cauchy's integral formula on the circle $\left\{a + r\cdot\exp(i\cdot\varphi) : 0 \le \varphi \le 2\pi\right\}$ gives
\[
G(a) = \frac{1}{2\pi}\int_0^{2\pi}\mathrm{d}\varphi\, G\left(a + r\cdot\exp(i\cdot\varphi)\right),
\]
and hence
\[
m = |G(a)| \le \frac{1}{2\pi}\int_0^{2\pi}\mathrm{d}\varphi\, \left|G\left(a + r\cdot\exp(i\cdot\varphi)\right)\right|.
\]
The integrand on the right is a continuous function of $\varphi$ with values at most $m$. If $\cos\varphi > 0$, the point $a + r\cdot\exp(i\cdot\varphi)$ of $\Omega$ has a larger real part than $a$, so it does not belong to $K'$, and the integrand is smaller than $m$ there. Thus the continuous function $\varphi \mapsto m - \left|G\left(a + r\cdot\exp(i\cdot\varphi)\right)\right|$ is non-negative on $[0,2\pi]$ and positive on the interval $0 \le \varphi < \frac{\pi}{2}$, so that its integral is positive, and the right side of the last inequality is smaller than $m$. This is a contradiction. Therefore $|G| \le 1$, that is, $v \le 0$, on $\Omega$. As $\tilde{\delta} > 0$ is arbitrary,
\[
\left|\mathfrak{u}(\rho'\cdot z)\right| \le {\boldsymbol{\delta}'}^{1-\hat{\boldsymbol{\omega}}(z)}\cdot M_{\mathfrak{g}}^{\hat{\boldsymbol{\omega}}(z)} \qquad (z \in \Omega).
\]

For $z \in \Omega$ one has
\[
\frac{1+z}{1-z} = \frac{1 - |z|^2 + 2i\cdot\operatorname{Im} z}{|1-z|^2}, \qquad \text{hence} \qquad \hat{\boldsymbol{\omega}}(z) = \frac{2}{\pi}\arctan\left(\frac{2\operatorname{Im} z}{1 - |z|^2}\right).
\]
Let $0 < q' < 1$ and $\omega' := \frac{4}{\pi}\arctan q'$. If $z \in \Omega$ and $|z| \le q'$, then $\operatorname{Im} z \le q'$ and $1 - |z|^2 \ge 1 - q'^2$. The function $\arctan$ is increasing, and $\arctan\frac{2q'}{1-q'^2} = 2\arctan q'$ by the formula for the tangent of a double angle, because $0 < \arctan q' < \frac{\pi}{4}$. Therefore
\[
\hat{\boldsymbol{\omega}}(z) \le \frac{2}{\pi}\arctan\left(\frac{2q'}{1 - q'^2}\right) = \omega' \qquad (z \in \Omega,\ |z| \le q').
\]
Since $\boldsymbol{\delta}' \le M_{\mathfrak{g}}$, the number ${\boldsymbol{\delta}'}^{1-s}\cdot M_{\mathfrak{g}}^s$ does not decrease as $s \in [0,1]$ increases. It follows that
\[
\left|\mathfrak{u}(\rho'\cdot z)\right| \le {\boldsymbol{\delta}'}^{1-\omega'}\cdot M_{\mathfrak{g}}^{\omega'} \qquad \text{for } z \in \Omega,\ |z| \le q'.
\]

Now let $t \in \mathbb{C}$ with $|t| \le q\cdot\rho$ and $\operatorname{Im} t > 0$. For $\rho' \in (q\cdot\rho, \rho)$ put $q' := q\cdot\rho/\rho'$, so that $0 < q' < 1$, and $z := t/\rho'$, so that $z \in \Omega$ and $|z| \le q'$. The last inequality gives $|\mathfrak{u}(t)| \le {\boldsymbol{\delta}'}^{1-\omega'}\cdot M_{\mathfrak{g}}^{\omega'}$ with $\omega' = \frac{4}{\pi}\arctan\frac{q\cdot\rho}{\rho'}$. As $\rho'$ tends to $\rho$, $\omega'$ tends to $\boldsymbol{\omega}$; this proves the asserted inequality for $\operatorname{Im} t > 0$. The function $t \mapsto \overline{\mathfrak{u}(\bar{t})}$ is holomorphic on the same disk and satisfies the same hypotheses as $\mathfrak{u}$, and its modulus at the point $\bar{t}$ equals $|\mathfrak{u}(t)|$; the case already treated, applied to this function, gives the inequality for $\operatorname{Im} t < 0$. For real $t$ with $|t| \le q\cdot\rho$ the inequality follows from the hypothesis, because $\boldsymbol{\delta}' \le {\boldsymbol{\delta}'}^{1-\boldsymbol{\omega}}\cdot M_{\mathfrak{g}}^{\boldsymbol{\omega}}$.

\emph{Several variables.} Let now $q := \frac{1}{2n}$ and $\boldsymbol{\omega} := \frac{4}{\pi}\arctan q$. Define $\boldsymbol{\delta}_0 := \boldsymbol{\delta}$ and $\boldsymbol{\delta}_i := \boldsymbol{\delta}_{i-1}^{1-\boldsymbol{\omega}}\cdot M_{\mathfrak{g}}^{\boldsymbol{\omega}}$ for $i = 1, \dots, n$. Since $\boldsymbol{\delta} \le M_{\mathfrak{g}}$, induction on $i$ gives $\boldsymbol{\delta}_i \le M_{\mathfrak{g}}$ and
\[
\boldsymbol{\delta}_i = \boldsymbol{\delta}^{(1-\boldsymbol{\omega})^i}\cdot M_{\mathfrak{g}}^{1-(1-\boldsymbol{\omega})^i} \qquad (i = 0, \dots, n).
\]
For $i = 0, \dots, n$ the following assertion holds:
\[
|\mathfrak{g}(t)| \le \boldsymbol{\delta}_i \qquad \text{if } |t_1|, \dots, |t_i| \le q\cdot\rho \text{ and } t_{i+1}, \dots, t_n \in (-\rho,\rho).
\]
For $i = 0$ it is the hypothesis. Let $1 \le i \le n$, and assume the assertion for $i - 1$. Fix $t_1, \dots, t_{i-1} \in \mathbb{C}$ with $|t_j| \le q\cdot\rho$ and $t_{i+1}, \dots, t_n \in (-\rho,\rho)$, and let $\mathfrak{u}(s) := \mathfrak{g}(t_1, \dots, t_{i-1}, s, t_{i+1}, \dots, t_n)$ for $|s| < \rho$. Then $\mathfrak{u}$ is holomorphic on the disk of radius $\rho$, $|\mathfrak{u}| \le M_{\mathfrak{g}}$ there, and $|\mathfrak{u}(s)| \le \boldsymbol{\delta}_{i-1}$ for $s \in (-\rho,\rho)$ by the assertion for $i - 1$; moreover $\boldsymbol{\delta}_{i-1} \le M_{\mathfrak{g}}$. The inequality of the one-variable case, with $\boldsymbol{\delta}' := \boldsymbol{\delta}_{i-1}$, gives $|\mathfrak{u}(s)| \le \boldsymbol{\delta}_i$ for $|s| \le q\cdot\rho$, which is the assertion for $i$.

For $i = n$ the assertion reads: $|\mathfrak{g}(t)| \le \boldsymbol{\delta}_n$ whenever $|t_j| \le q\cdot\rho = \frac{\rho}{2n}$ for $j = 1, \dots, n$. As $\arctan q \le q$, one has $\boldsymbol{\omega} \le \frac{4q}{\pi} = \frac{2}{\pi\cdot n} < 1$, and Bernoulli's inequality gives
\[
(1-\boldsymbol{\omega})^n \ge 1 - n\cdot\boldsymbol{\omega} \ge 1 - \frac{2}{\pi} > \frac{1}{3}.
\]
If $M_{\mathfrak{g}} = 0$, then $\mathfrak{g} = 0$ and there is nothing to prove. Otherwise $0 \le \boldsymbol{\delta}/M_{\mathfrak{g}} \le 1$, and therefore
\[
\boldsymbol{\delta}_n = M_{\mathfrak{g}}\cdot\left(\frac{\boldsymbol{\delta}}{M_{\mathfrak{g}}}\right)^{(1-\boldsymbol{\omega})^n} \le M_{\mathfrak{g}}\cdot\left(\frac{\boldsymbol{\delta}}{M_{\mathfrak{g}}}\right)^{\frac{1}{3}} = \boldsymbol{\delta}^{\frac{1}{3}}\cdot M_{\mathfrak{g}}^{\frac{2}{3}}.
\]
\end{proof}

\subsubsection{Periodic functions holomorphic in a tube}\label{sec:periodic-functions-holomorphic}

\begin{lemma}[Oscillation at the real points]\label{lem:oscillation-real-points}
Let $\boldsymbol{\tau} \ge 1$, and let $\mathfrak{g}$ be a continuous function on the closed tube
\[
\left\{ \theta \in \mathbb{C}^3 : |\operatorname{Im}\theta_{\nu}| \le \boldsymbol{\tau} \text{ for } \nu = 0, 1, 2 \right\}
\]
that is holomorphic in its interior, $2\pi$-periodic in each of the variables $\theta_0$, $\theta_1$, $\theta_2$, and bounded in modulus by $B_{\mathfrak{g}}$. Then
\[
|\mathfrak{g}(\theta) - \mathfrak{g}(\theta')| \le 52\,B_{\mathfrak{g}}\cdot\exp(-\boldsymbol{\tau}) \qquad \text{for all } \theta, \theta' \in \mathbb{R}^3.
\]
\end{lemma}

\begin{proof}
For $m \in \mathbb{Z}^3$ put $m\cdot\theta := \sum_{\nu} m_{\nu}\cdot\theta_{\nu}$ and $|m|_1 := \sum_{\nu}|m_{\nu}|$, and let
\[
\hat{g}(m) := \frac{1}{(2\pi)^3}\int_{[-\pi,\pi]^3} \mathrm{d}\theta\, \mathfrak{g}(\theta)\cdot\exp(-i\cdot m\cdot\theta)
\]
be the Fourier coefficients of the restriction of $\mathfrak{g}$ to $\mathbb{R}^3$. Fix $m$ and a number $\tau'$ with $0 < \tau' < \boldsymbol{\tau}$. For each $\nu$ with $m_{\nu} \ne 0$, replace, in the integral over $\theta_{\nu}$, the segment $[-\pi,\pi]$ by the parallel segment $[-\pi,\pi] - i\cdot\tau'\cdot m_{\nu}/|m_{\nu}|$. By Cauchy's theorem this leaves the integral unchanged: the integrand is holomorphic in $\theta_{\nu}$ on a neighbourhood of the closed rectangle between the two segments, and the integrals over its two vertical sides $\operatorname{Re}\theta_{\nu} = -\pi$ and $\operatorname{Re}\theta_{\nu} = \pi$ cancel, because the integrand is $2\pi$-periodic in $\theta_{\nu}$. On the new segment $|\exp(-i\cdot m_{\nu}\cdot\theta_{\nu})| = \exp(-\tau'\cdot|m_{\nu}|)$. Together with $|\mathfrak{g}| \le B_{\mathfrak{g}}$ this gives $|\hat{g}(m)| \le B_{\mathfrak{g}}\cdot\exp(-\tau'\cdot|m|_1)$, and, letting $\tau'$ tend to $\boldsymbol{\tau}$,
\[
|\hat{g}(m)| \le B_{\mathfrak{g}}\cdot\exp\left(-\boldsymbol{\tau}\cdot|m|_1\right) \qquad (m \in \mathbb{Z}^3).
\]
Put $x := \exp(-\boldsymbol{\tau})$. Then $0 < x \le \exp(-1)$, because $\boldsymbol{\tau} \ge 1$, and
\[
\sum_{m \in \mathbb{Z}^3} \exp\left(-\boldsymbol{\tau}\cdot|m|_1\right) = \left(\sum_{j \in \mathbb{Z}} x^{|j|}\right)^3 = \left(\frac{1+x}{1-x}\right)^3.
\]
So the Fourier series of the restriction of $\mathfrak{g}$ to $\mathbb{R}^3$ converges absolutely. This restriction is a smooth $2\pi$-periodic function, and therefore it is the sum of its Fourier series. For $\theta \in \mathbb{R}^3$ it follows that
\[
|\mathfrak{g}(\theta) - \hat{g}(0)| \le \sum_{m \ne 0}|\hat{g}(m)| \le B_{\mathfrak{g}}\cdot\left[\left(\frac{1+x}{1-x}\right)^3 - 1\right].
\]
With $s := \frac{2x}{1-x} = 2\sum_{k \ge 1} x^k$ one has $\left(\frac{1+x}{1-x}\right)^3 - 1 = (1+s)^3 - 1 = 3s + 3s^2 + s^3$. Hence $x^{-1}\cdot\left[\left(\frac{1+x}{1-x}\right)^3 - 1\right]$ is the sum of a power series in $x$ with non-negative coefficients, and so it increases with $x \in (0,1)$; at $x = \exp(-1)$ its value is smaller than $25$. Consequently $|\mathfrak{g}(\theta) - \hat{g}(0)| \le 25\,B_{\mathfrak{g}}\cdot\exp(-\boldsymbol{\tau})$ for $\theta \in \mathbb{R}^3$, and for $\theta, \theta' \in \mathbb{R}^3$
\[
|\mathfrak{g}(\theta) - \mathfrak{g}(\theta')| \le |\mathfrak{g}(\theta) - \hat{g}(0)| + |\hat{g}(0) - \mathfrak{g}(\theta')| \le 50\,B_{\mathfrak{g}}\cdot\exp(-\boldsymbol{\tau}) \le 52\,B_{\mathfrak{g}}\cdot\exp(-\boldsymbol{\tau}).
\]
\end{proof}

\subsection{Dependence of the stop-level fermion operators on the connection}\label{sec:dependence-stop-level}

Let $G_{\mathrm{w}}$ be the set of the complex bond functions on the lattice $T_{*}$ of the level $k^{*}$ of the form
\[
A_r + \mathscr{C}, \qquad A_r \text{ real and } A_r \in \tilde{R}_{*}, \qquad \mathscr{C} \in \mathfrak{X}_{*},
\]
and let
\[
G_{\mathrm{w}}^{+} := \left\{ A + A_1 : A \in G_{\mathrm{w}},\ A_1 \text{ a complex bond function on } T_{*} \text{ with } e_{*}\cdot|A_1|_{\infty} \le c_1 \right\},
\]
where $c_1$ is the constant of Section~\ref{sec:constants-3}. By definition $G_{\mathrm{w}} \subset \tilde{R}_{*} + \mathfrak{X}_{*}$, and $\tilde{R}_{*} + \mathfrak{X}_{*} \subset \tilde{G}_{*}$ by clause~\eqref{it:trajectory-fermionic-operators-SL2} of Proposition~\ref{prop:trajectory-fermionic-operators}; so $G_{\mathrm{w}}$ is a subset of $\tilde{G}_{*}$.

Let $A = A_r + \mathscr{C} \in G_{\mathrm{w}}$, with $A_r$ and $\mathscr{C}$ as above. As $A_r$ is real, $\operatorname{Im} A = \operatorname{Im}\mathscr{C}$; and $|\mathscr{C}|_{\infty} < e_{*}^{-1+2\varepsilon}$ by the definition of $\mathfrak{X}_{*}$ in Definition~\ref{def:domains}. Condition~\eqref{eq:renormalisation-group-A8} and the inequality $e_{*} \le e_0$ of Lemma~\ref{lem:scales-stop-level}(a) give $e_{*}^{2\varepsilon} \le \frac{1}{16}$ and $3M\cdot e_{*}^{2\varepsilon} \le \frac{1}{4}$ (Section~\ref{sec:smallness-conditions}). Hence
\[
\mathfrak{i}(A) = e_{*}\cdot|\operatorname{Im}\mathscr{C}|_{\infty} < e_{*}^{2\varepsilon} \le \frac{1}{16}, \qquad 3M\cdot\mathfrak{i}(A) < \frac{1}{4} \qquad (A \in G_{\mathrm{w}}).
\]
If $A \in G_{\mathrm{w}}$ and $A_1$ is a complex bond function on $T_{*}$ with $e_{*}\cdot|A_1|_{\infty} \le c_1$, then, since $c_1 \le \frac{1}{8}$,
\[
\mathfrak{i}(A + A_1) \le \mathfrak{i}(A) + e_{*}\cdot|A_1|_{\infty} \le \frac{1}{16} + \frac{1}{8} < \frac{1}{4};
\]
thus $\mathfrak{i} < \frac{1}{4}$ on $G_{\mathrm{w}}^{+}$. The set $G_{\mathrm{w}}$ is invariant under $A \mapsto -A$: together with $A_r$, also $-A_r$ is real and lies in $\tilde{R}_{*}$, by~\eqref{it:trajectory-fermionic-operators-SL2}, and $\mathfrak{X}_{*}$ is invariant under $\mathscr{C} \mapsto -\mathscr{C}$ by its definition.

Let $A \in G_{\mathrm{w}}$. Since $A \in \tilde{G}_{*}$, \eqref{it:trajectory-fermionic-operators-SL2} gives $|S(A)|(y,y') \le C_F\cdot E(y,y')$ and $|\mathrm{Sa}(A)|(y,y') \le C_F\cdot E(y,y')$ for all $y, y' \in T^0_{*}$. Hence
\[
|S(A)|(y,y')\cdot\exp\left(\frac{\gamma_F}{2}\cdot d(y,y')\right) \le C_F\cdot\exp\left(-\frac{\gamma_F}{2}\cdot d(y,y')\right),
\]
and the same holds with $\mathrm{Sa}(A)$ in place of $S(A)$. By Section~\ref{sec:constants-3}, $\sum_{y'}\exp\left(-\frac{\gamma_F}{2}\cdot d(y,y')\right) \le \mathrm{ssum}\left(\frac{\gamma_F}{2}\right)$ for every $y$, because $0 < \frac{\gamma_F}{2} \le \frac{1}{2}$; and $d$ is symmetric. With $\mathbf{S} = C_F\cdot\mathrm{ssum}\left(\frac{\gamma_F}{2}\right)$ it follows that
\[
\left\|S(A)\right\|_{\gamma_F/2} \le \mathbf{S}, \qquad \left\|\mathrm{Sa}(A)\right\|_{\gamma_F/2} \le \mathbf{S} \qquad (A \in G_{\mathrm{w}}).
\]

For a complex bond function $B$ on $T_{*}$ let
\[
\epsilon_B := e_{*}\cdot|B|_{\infty}, \qquad P_B := \left\{ y \in T^0_{*} : \Delta_y \text{ contains an end point of a bond } b \text{ with } B(b) \ne 0 \right\},
\]
\[
P_B^{+} := \left\{ y \in T^0_{*} : d(y,P_B) \le M + 1 \right\}.
\]
For $B = \hat{c}_c$, $c \in \mathrm{Cel}$, these are the quantities $\epsilon_c$, $P_c$ and $P_c^{+}$ of Section~\ref{sec:probes-mirror-copies}.

\begin{lemma}[Bounds for the elementary operators]\label{lem:bounds-elementary-operators}
Let $A$, $B$, $B'$ be complex bond functions on $T_{*}$ with $\mathfrak{i}(A) \le \frac{1}{4}$.
\begin{enumerate}
\item[(a)] For all $y, y' \in T^0_{*}$,
\[
|Q_{*}(A)|(y,y') \le \frac{3}{2}\cdot\mathbf{1}[y = y'], \qquad |Q_{*}^T(-A)|(y,y') \le \frac{3}{2}\cdot\mathbf{1}[y = y'],
\]
\[
|\partial_BQ_{*}(A)|(y,y') \le \frac{9}{4}\,\epsilon_B\cdot\mathbf{1}[y = y' \in P_B], \qquad |\partial_BQ_{*}^T(-A)|(y,y') \le \frac{9}{4}\,\epsilon_B\cdot\mathbf{1}[y = y' \in P_B].
\]
\item[(b)] For $z, z' \in T^0_{*}$ let $|z - z'|_1$ be the least $\ell^1$ norm of a vector of $\mathbb{Z}^3$ that represents the class $z - z'$ in $T^0_{*} = \mathbb{Z}^3/l_{*}\mathbb{Z}^3$, in the same way as $|z - z'|$ is the least sup norm of such a vector (Definition~\ref{def:sup-norm-distance}). Thus $|z - z'|_1 = 1$ if and only if $z' = z + e_{\nu}$ or $z' = z - e_{\nu}$ in $T^0_{*}$ for a direction $\nu$, and these six neighbours of $z$ are distinct, because $l_{*} = M\cdot r_{*}$ with $M \ge 1$ and $r_{*} \ge L^2 > 2$ (Lemma~\ref{lem:scales-stop-level}(a), (b), Section~\ref{sec:constants-3}). For all $z, z' \in T^0_{*}$,
\[
|\partial_B\mathsf{D}(A)|(z,z') \le
\begin{cases}
15\,\epsilon_B & \text{if } z = z' \in P_B,\\
\frac{3}{2}\,\epsilon_B & \text{if } z, z' \in P_B \text{ and } |z - z'|_1 = 1,
\end{cases}
\]
and $|\partial_B\mathsf{D}(A)|(z,z') = 0$ for all other pairs $(z,z')$. Consequently $\left\|\partial_B\mathsf{D}(A)\right\|_{\tilde{\gamma}} \le 40\,\epsilon_B$ for $0 \le \tilde{\gamma} \le 1$, and for all $y, y' \in T^0_{*}$
\begin{equation}\label{eq:four-jets-12}
\sum_{z,z' \in T^0_{*}} E(y,z)\cdot|\partial_B\mathsf{D}(A)|(z,z')\cdot E(z',y') \le 40\,\epsilon_B\cdot\sum_{z \in P_B} E(y,z)\cdot E(z,y').
\end{equation}
\item[(c)] Let in addition $3M\cdot\mathfrak{i}(A) \le \frac{1}{4}$, let $(x,x') \in \mathrm{Pz}_{*}$, and let $f$ be a spinor function on $T_{*}$. Then $|U_A(x,x')^{\pm 1}| \le \frac{4}{3}$, and
\[
\left|\left(\partial_B\delta_{\alpha,A}f\right)(x,x')\right| \le 4|x - x'|^{1-\alpha}\cdot\epsilon_B\cdot|f(x')| \le 4M\cdot\epsilon_B\cdot|f(x')|;
\]
if $|x - x'| < 1$, the left side is at most $4\,\epsilon_B\cdot|f(x')|$. The left side vanishes unless $d([x],P_B) \le M$ and $d([x'],P_B) \le M$; if $|x - x'| < 1$, it vanishes unless $d([x],P_B) \le 1$.
\item[(d)] If $d(P_B,P_{B'}) \ge M + 2$, then $\partial_B\partial_{B'}O(A) = 0$ for each of the operator-valued functions $O(A) = D_A$, $Q_{*}(A)$, $Q_{*}^T(-A)$, $P_{*}(A)$, $\mathsf{D}(A)$ and $\delta_{\alpha,A}$ of $A$.
\end{enumerate}
\end{lemma}

\begin{proof}
The matrix elements of $D_A$, $Q_{*}(A)$, $Q_{*}^T(-A)$, $P_{*}(A)$, $\mathsf{D}(A)$ and $\delta_{\alpha,A}$ are finite sums of constants and of products of constants with exponentials of linear functions of $A$ (Section~\ref{sec:fermionic-operators-stop}). So these operators are entire functions of $A$, and $\partial_B$ and $\partial_B\partial_{B'}$ can be applied to them at every complex bond function $A$. The maps $A \mapsto \sigma_A$ and $A \mapsto \upsilon_A$ are linear with real coefficients; hence $\operatorname{Im}\sigma_A = \sigma_{\operatorname{Im}A}$, $\operatorname{Im}\upsilon_A = \upsilon_{\operatorname{Im}A}$, $\partial_B\sigma_A = \sigma_B$ and $\partial_B\upsilon_A = \upsilon_B$.

\emph{(a)} Let $y \in T^0_{*}$ and $x \in \Delta_y$. The bound $|\sigma_{A'}(y,x)| \le \frac{3}{2}\,e_{*}\cdot|A'|_{\infty}$ of Section~\ref{sec:fermionic-operators-stop}, applied to $A' = \operatorname{Im}A$, gives
\[
\left|\exp\left(\pm i\cdot\sigma_A(y,x)\right)\right| = \exp\left(\mp\operatorname{Im}\sigma_A(y,x)\right) \le \exp\left(\frac{3}{2}\mathfrak{i}(A)\right) \le \exp\left(\frac{3}{8}\right) \le \frac{3}{2}.
\]
The value $(Q_{*}(A)f)(y)$ depends only on the restriction of $f$ to $\Delta_y$, and, as $\eta_{*}^3\cdot\#|\Delta_y| = 1$,
\[
\left|(Q_{*}(A)f)(y)\right| \le \frac{3}{2}\,\eta_{*}^3\cdot\#|\Delta_y|\cdot\max_{x \in \Delta_y}|f(x)| = \frac{3}{2}\max_{x \in \Delta_y}|f(x)|.
\]
Hence $|Q_{*}(A)|(y,y') = 0$ for $y' \ne y$, and $|Q_{*}(A)|(y,y) \le \frac{3}{2}$. For a spinor function $v$ on $T^0_{*}$ one has $(Q_{*}^T(-A)v)(x) = \exp(-i\cdot\sigma_A(y,x))\cdot v(y)$ for $x \in \Delta_y$; so the restriction of $Q_{*}^T(-A)v$ to $\Delta_y$ depends only on $v(y)$ and is bounded by $\frac{3}{2}|v(y)|$, which gives the bound for $|Q_{*}^T(-A)|(y,y')$. Further,
\[
\partial_B\exp\left(\pm i\cdot\sigma_A(y,x)\right) = \pm i\cdot\sigma_B(y,x)\cdot\exp\left(\pm i\cdot\sigma_A(y,x)\right), \qquad |\sigma_B(y,x)| \le \frac{3}{2}\epsilon_B.
\]
By Section~\ref{sec:fermionic-operators-stop}, $\sigma_B(y,\cdot)$ depends only on the values of $B$ on the bonds with both end points in $\Delta_y$. If $y \notin P_B$, these values vanish, and so does $\sigma_B(y,\cdot)$. The bounds for $\partial_BQ_{*}(A)$ and $\partial_BQ_{*}^T(-A)$ follow, with the factor $\frac{3}{2}\cdot\frac{3}{2} = \frac{9}{4}$.

\emph{(b)} For a bond $b$ of $T_{*}$ one has $\partial_Bu_A(b) = i\cdot e_{*}\cdot\eta_{*}\cdot B(b)\cdot u_A(b)$, and, as $\eta_{*} \le 1$,
\[
|u_A(b)| = \exp\left(-e_{*}\cdot\eta_{*}\cdot\operatorname{Im}A(b)\right) \le \exp\left(\eta_{*}\cdot\mathfrak{i}(A)\right) \le \exp\left(\frac{1}{4}\right) \le \frac{4}{3}.
\]
In the formula for $D_A$ of Section~\ref{sec:fermionic-operators-stop} the terms $-f(x)$ do not depend on $A$. Therefore, with $x^{\boldsymbol{\sigma}} = x + \boldsymbol{\sigma}\cdot\eta_{*}\cdot e_{\nu}$,
\[
\left(\partial_BD_Af\right)(x) = -i\cdot\sum_{\nu = 0}^2\sum_{\boldsymbol{\sigma} = \pm 1} e_{*}\cdot B(x,x^{\boldsymbol{\sigma}})\cdot u_A(x,x^{\boldsymbol{\sigma}})\cdot P^{\boldsymbol{\sigma}}_{\nu}\cdot f(x^{\boldsymbol{\sigma}}).
\]
As $|P^{\boldsymbol{\sigma}}_{\nu}\cdot v| \le |v|$ (Section~\ref{sec:fermionic-operators-stop}), each term has modulus at most $\frac{4}{3}\epsilon_B\cdot|f(x^{\boldsymbol{\sigma}})|$. A term vanishes unless $B(x,x^{\boldsymbol{\sigma}}) \ne 0$, and if $B(x,x^{\boldsymbol{\sigma}}) \ne 0$, then $[x]$ and $[x^{\boldsymbol{\sigma}}]$ belong to $P_B$. Let $x \in \Delta_z$. Each of the six points $x^{\boldsymbol{\sigma}}$ lies in $\Delta_z$ or in a cube $\Delta_{z'}$ with $|z - z'|_1 = 1$, and a given cube $\Delta_{z'}$ with $z' \ne z$ contains at most one of them. It follows that $|\partial_BD_A|(z,z) \le 6\cdot\frac{4}{3}\epsilon_B = 8\,\epsilon_B$; that $|\partial_BD_A|(z,z') \le \frac{4}{3}\epsilon_B$ if $|z - z'|_1 = 1$; that $|\partial_BD_A|(z,z') = 0$ if $z' \ne z$ and $|z - z'|_1 \ne 1$; and that $|\partial_BD_A|(z,z') = 0$ unless $z, z' \in P_B$.

The operator $P_{*}(A) = Q_{*}^T(-A)\cdot Q_{*}(A)$ acts inside each cube: for $x \in \Delta_y$, the value $(P_{*}(A)f)(x)$ depends only on the restriction of $f$ to $\Delta_y$. The same holds for
\[
\partial_BP_{*}(A) = \left(\partial_BQ_{*}^T(-A)\right)\cdot Q_{*}(A) + Q_{*}^T(-A)\cdot\partial_BQ_{*}(A).
\]
By~\eqref{eq:four-jets-Maj1} and (a),
\[
|\partial_BP_{*}(A)|(y,y) \le 2\cdot\frac{9}{4}\cdot\frac{3}{2}\,\epsilon_B \le 7\,\epsilon_B,
\]
and $|\partial_BP_{*}(A)|(y,y) = 0$ for $y \notin P_B$. Now $\mathsf{D}(A) = D_A + \mathfrak{m}_{*} + b_{*}\cdot P_{*}(A)$, so that $\partial_B\mathsf{D}(A) = \partial_BD_A + b_{*}\cdot\partial_BP_{*}(A)$, and $0 < b_{*} \le 1$ by Lemma~\ref{lem:scales-stop-level}(a). The bounds of (b) for $|\partial_B\mathsf{D}(A)|(z,z')$ follow, because $8 + 7 = 15$ and $\frac{4}{3} \le \frac{3}{2}$.

Let $0 \le \tilde{\gamma} \le 1$. If $|z - z'|_1 = 1$, then $d(z,z') = 1$, and every $z$ has six such neighbours $z'$. By what has just been shown, each of the sums over $z'$ and over $z$ in the definition of $\left\|\partial_B\mathsf{D}(A)\right\|_{\tilde{\gamma}}$ in Section~\ref{sec:fermionic-operators-stop} is at most
\[
\left(15 + 6\cdot\frac{3}{2}\exp(\tilde{\gamma})\right)\cdot\epsilon_B \le \left(15 + 9\exp(1)\right)\cdot\epsilon_B \le 40\,\epsilon_B.
\]
If $d(z,z') \le 1$, then $d(z',y') \ge d(z,y') - 1$, so that $E(z',y') \le \exp(\gamma_F)\cdot E(z,y') \le \exp(1)\cdot E(z,y')$, since $\gamma_F \le 1$ (Section~\ref{sec:constants-3}). On the left side of~\eqref{eq:four-jets-12} only the pairs $(z,z')$ with $z \in P_B$ and $d(z,z') \le 1$ contribute, and for every $z \in P_B$
\[
\sum_{z' \in T^0_{*}} |\partial_B\mathsf{D}(A)|(z,z')\cdot E(z',y') \le \left(15 + 6\cdot\frac{3}{2}\exp(1)\right)\cdot\epsilon_B\cdot E(z,y') \le 40\,\epsilon_B\cdot E(z,y').
\]
Multiplication by $E(y,z)$ and summation over $z \in P_B$ give~\eqref{eq:four-jets-12}.

\emph{(c)} By Definition~\ref{def:pairs-sites-holder}, $|x - x'| \le M$ for $(x,x') \in \mathrm{Pz}_{*}$. The bound $|\upsilon_{A'}(x,x')| \le 3|x - x'|\cdot e_{*}\cdot|A'|_{\infty}$ of Section~\ref{sec:fermionic-operators-stop}, applied to $A' = \operatorname{Im}A$, gives
\[
\left|\operatorname{Im}\upsilon_A(x,x')\right| \le 3|x - x'|\cdot\mathfrak{i}(A) \le 3M\cdot\mathfrak{i}(A) \le \frac{1}{4},
\]
and hence $|U_A(x,x')^{\pm 1}| = \exp\left(\mp\operatorname{Im}\upsilon_A(x,x')\right) \le \exp\left(\frac{1}{4}\right) \le \frac{4}{3}$. Applied to $A' = B$, the same bound gives $|\upsilon_B(x,x')| \le 3|x - x'|\cdot\epsilon_B$. Since $U_A = \exp(i\cdot\upsilon_A)$, one has $\partial_BU_A = i\cdot\upsilon_B\cdot U_A$; and in
\[
(\delta_{\alpha,A}f)(x,x') = |x - x'|^{-\alpha}\cdot\left(U_A(x,x')\cdot f(x') - f(x)\right)
\]
the term $-f(x)$ does not depend on $A$. So
\[
\left(\partial_B\delta_{\alpha,A}f\right)(x,x') = |x - x'|^{-\alpha}\cdot i\cdot\upsilon_B(x,x')\cdot U_A(x,x')\cdot f(x'),
\]
and the modulus of the right side is at most
\[
|x - x'|^{-\alpha}\cdot 3|x - x'|\cdot\epsilon_B\cdot\frac{4}{3}|f(x')| = 4|x - x'|^{1-\alpha}\cdot\epsilon_B\cdot|f(x')|.
\]
By stage~\ref{it:order-choice-2}, $0 < \alpha < 1$; and $M \ge 1$ (Section~\ref{sec:constants-3}). Hence $|x - x'|^{1-\alpha} \le \max\left(1, M^{1-\alpha}\right) \le M$, and $|x - x'|^{1-\alpha} < 1$ if $|x - x'| < 1$.

The left side vanishes if $\upsilon_B(x,x') = 0$. By Section~\ref{sec:fermionic-operators-stop}, the number $\upsilon_B(x,x')$ is $e_{*}\cdot\eta_{*}$ times the mean of the sums of $B$ along the six paths $\Gamma^{\pi}(x,x')$ of Definition~\ref{def:pairs-sites-holder}, hence a linear combination of the values of $B$ at the bonds of these paths. So, if $\upsilon_B(x,x') \ne 0$, one of these paths contains a bond $b$ with $B(b) \ne 0$. Let $w$ be an end point of such a bond; then $[w] \in P_B$. The paths lie in the box spanned by $x$ and $x'$ (Section~\ref{sec:fermionic-operators-stop}); so $|x - w| \le |x - x'|$ and $|x' - w| \le |x - x'|$. As $|x - [x]|$, $|x' - [x']|$ and $|w - [w]|$ are smaller than $\frac{1}{2}$, the distances $d([x],[w])$ and $d([x'],[w])$ are smaller than $|x - x'| + 1 \le M + 1$. Distances between unit sites are integers; so $d([x],[w]) \le M$ and $d([x'],[w]) \le M$, and $d([x],[w]) \le 1$ if $|x - x'| < 1$. As $[w] \in P_B$, this proves the assertion on the vanishing of the left side.

\emph{(d)} Let $d(P_B,P_{B'}) \ge M + 2$. Then $P_B \cap P_{B'} = \emptyset$, and no bond $b$ satisfies both $B(b) \ne 0$ and $B'(b) \ne 0$: for such a bond and an end point $w$ of it, $[w]$ would belong to $P_B \cap P_{B'}$.

For every bond $b$, $\partial_B\partial_{B'}u_A(b) = -e_{*}^2\cdot\eta_{*}^2\cdot B(b)\cdot B'(b)\cdot u_A(b) = 0$. Every matrix element of $D_A$ depends on $A$ through at most one factor $u_A(b)$; hence $\partial_B\partial_{B'}D_A = 0$.

For $y \in T^0_{*}$ and $x \in \Delta_y$,
\[
\partial_B\partial_{B'}\exp\left(\pm i\cdot\sigma_A(y,x)\right) = -\sigma_B(y,x)\cdot\sigma_{B'}(y,x)\cdot\exp\left(\pm i\cdot\sigma_A(y,x)\right) = 0,
\]
because $\sigma_B(y,\cdot) = 0$ for $y \notin P_B$ and $\sigma_{B'}(y,\cdot) = 0$ for $y \notin P_{B'}$, by the proof of (a), and $P_B \cap P_{B'} = \emptyset$. Therefore $\partial_B\partial_{B'}Q_{*}(A) = 0$ and $\partial_B\partial_{B'}Q_{*}^T(-A) = 0$. By the product rule and these two identities,
\[
\partial_B\partial_{B'}P_{*}(A) = \left(\partial_BQ_{*}^T(-A)\right)\cdot\left(\partial_{B'}Q_{*}(A)\right) + \left(\partial_{B'}Q_{*}^T(-A)\right)\cdot\left(\partial_BQ_{*}(A)\right).
\]
All four factors act inside each cube. By (a), the block $y$ of $\partial_BQ_{*}^T(-A)$ and of $\partial_BQ_{*}(A)$ vanishes for $y \notin P_B$, and the block $y$ of $\partial_{B'}Q_{*}^T(-A)$ and of $\partial_{B'}Q_{*}(A)$ vanishes for $y \notin P_{B'}$. Since $P_B \cap P_{B'} = \emptyset$, both products vanish, and $\partial_B\partial_{B'}P_{*}(A) = 0$. Consequently $\partial_B\partial_{B'}\mathsf{D}(A) = \partial_B\partial_{B'}D_A + b_{*}\cdot\partial_B\partial_{B'}P_{*}(A) = 0$.

Finally, $\partial_B\partial_{B'}U_A = -\upsilon_B\cdot\upsilon_{B'}\cdot U_A$, so that
\[
\left(\partial_B\partial_{B'}\delta_{\alpha,A}f\right)(x,x') = -|x - x'|^{-\alpha}\cdot\upsilon_B(x,x')\cdot\upsilon_{B'}(x,x')\cdot U_A(x,x')\cdot f(x').
\]
Suppose that $\upsilon_B(x,x') \ne 0$ and $\upsilon_{B'}(x,x') \ne 0$ for some $(x,x') \in \mathrm{Pz}_{*}$. Then the paths $\Gamma^{\pi}(x,x')$ contain a bond $b$ with $B(b) \ne 0$ and a bond $b'$ with $B'(b') \ne 0$. End points $w$ of $b$ and $w'$ of $b'$ lie in the box spanned by $x$ and $x'$; so $|w - w'| \le |x - x'| \le M$, and $d([w],[w']) < M + 1$. As $[w] \in P_B$ and $[w'] \in P_{B'}$, this contradicts $d(P_B,P_{B'}) \ge M + 2$. Hence $\upsilon_B\cdot\upsilon_{B'} = 0$ on $\mathrm{Pz}_{*}$, and $\partial_B\partial_{B'}\delta_{\alpha,A} = 0$.
\end{proof}

\paragraph{Leg indices.}
A \emph{leg index} $\boldsymbol{\upsilon}$ is either an index $\xi=(x,\boldsymbol{\beta},\boldsymbol{\omega})$ with $x\in T_{*}$, or an index $\boldsymbol{\zeta}=(x,x',\boldsymbol{\beta},\boldsymbol{\omega})$ with $(x,x')\in\mathrm{Pz}_{*}$, the set of pairs of Definition~\ref{def:pairs-sites-holder} at the level $k^{*}$; here $\boldsymbol{\beta}$ is a spinor index and $\boldsymbol{\omega}\in\{0,1\}$, as for the generators $\psi(\xi)$ and $\chi(\boldsymbol{\zeta})$ of Section~\ref{sec:grassmann-algebras-norms}. The \emph{position} of a leg index is the following set of unit sites:
\[
\mathrm{pos}(\xi):=\{[x]\},\qquad
\mathrm{pos}(\boldsymbol{\zeta}):=\begin{cases}\{[x]\}&\text{if }|x-x'|<1,\\ \{[x],[x']\}&\text{if }|x-x'|\ge1,\end{cases}
\]
where $[x]\in T^0_{*}$ is the unit site whose cube $\Delta_{[x]}$ contains $x$ (Section~\ref{sec:fermionic-operators-stop}).

Let $A\in\tilde{G}_{*}$, and let $\Psi$, $\bar{\Psi}$ be the spinor functions on $T^0_{*}$ with the components $\Psi_{\boldsymbol{\beta}}(y)=\Psi(y,\boldsymbol{\beta},0)$ and $\bar{\Psi}_{\boldsymbol{\beta}}(y)=\Psi(y,\boldsymbol{\beta},1)$, formed with the generators $\Psi(u)$ of $\mathrm{Gr}[\Psi]$ (Section~\ref{sec:grassmann-norms-gaussian}). Consider the substitution
\[
\psi=H_{*}(A)\cdot\Psi,\qquad\bar{\psi}=\bar{H}_{*}(A)\cdot\bar{\Psi},\qquad\chi=\delta_{\alpha,A}\psi,\qquad\bar{\chi}=\delta_{\alpha,-A}\bar{\psi},
\]
in which, as in Section~\ref{sec:grassmann-algebras-norms}, $\psi(x,\boldsymbol{\beta},0)$, $\psi(x,\boldsymbol{\beta},1)$ stand for the components $\psi_{\boldsymbol{\beta}}(x)$, $\bar{\psi}_{\boldsymbol{\beta}}(x)$, and $\chi(x,x',\boldsymbol{\beta},0)$, $\chi(x,x',\boldsymbol{\beta},1)$ for $\chi_{\boldsymbol{\beta}}(x,x')$, $\bar{\chi}_{\boldsymbol{\beta}}(x,x')$. After the substitution each $\psi(\xi)$ and each $\chi(\boldsymbol{\zeta})$ is a linear combination of the generators $\Psi(u)$, $u=(y',\boldsymbol{\beta}',\boldsymbol{\omega}')$, $y'\in T^0_{*}$. For a leg index $\boldsymbol{\upsilon}$ and such an index $u$, let $\mathfrak{H}_A(\boldsymbol{\upsilon};u)\in\mathbb{C}$ be the coefficient of $\Psi(u)$ in $\psi(\xi)$ if $\boldsymbol{\upsilon}=\xi$, and in $\chi(\boldsymbol{\zeta})$ if $\boldsymbol{\upsilon}=\boldsymbol{\zeta}$. The fields $\psi$, $\chi$ contain only generators with $\boldsymbol{\omega}'=0$, and $\bar{\psi}$, $\bar{\chi}$ only generators with $\boldsymbol{\omega}'=1$; hence $\mathfrak{H}_A(\boldsymbol{\upsilon};u)=0$ if $\boldsymbol{\omega}'\ne\boldsymbol{\omega}$. With $\boldsymbol{\omega}$ the last entry of $\boldsymbol{\upsilon}$, put
\[
|\mathfrak{H}_A|(\boldsymbol{\upsilon};y'):=\sum_{\boldsymbol{\beta}'}\left|\mathfrak{H}_A\left(\boldsymbol{\upsilon};(y',\boldsymbol{\beta}',\boldsymbol{\omega})\right)\right|\qquad\left(y'\in T^0_{*}\right).
\]
By \eqref{it:trajectory-fermionic-operators-SL2} and the formulas of Section~\ref{sec:fermionic-operators-stop}, the coefficients $\mathfrak{H}_A(\boldsymbol{\upsilon};u)$ are analytic functions of $A$ on $\tilde{G}_{*}$. Their derivatives $\partial_B\mathfrak{H}_A(\boldsymbol{\upsilon};u)$ and $\partial_B\partial_{B'}\mathfrak{H}_A(\boldsymbol{\upsilon};u)$ in the directions of complex bond functions $B$, $B'$ on $T_{*}$ are taken as in Section~\ref{sec:fermionic-operators-stop}, and the numbers $|\partial_B\mathfrak{H}_A|(\boldsymbol{\upsilon};y')$, $|\partial_B\partial_{B'}\mathfrak{H}_A|(\boldsymbol{\upsilon};y')$ are formed from them in the same way as $|\mathfrak{H}_A|(\boldsymbol{\upsilon};y')$.

For a set $Z$ of unit sites and a unit site $y'$, the distance of $y'$ from $Z$ is $d(Z,y'):=d(y',Z):=\min_{w\in Z}d(w,y')$; in particular $d(\mathrm{pos}(\boldsymbol{\upsilon}),y')$ is the distance of $y'$ from $\mathrm{pos}(\boldsymbol{\upsilon})$. For an index $u=(y',\boldsymbol{\beta}',\boldsymbol{\omega}')$ and a set $Z$ of unit sites put
\[
d(Z,u):=d(u,Z):=d(Z,y'),
\]
and, for a single unit site $z$, $d(z,u):=d(u,z):=d(z,y')$.

\begin{lemma}[Derivatives of the legs and of $D_{*}$ on $G_{\mathrm{w}}$]\label{lem:derivatives-legs}
Let $A\in G_{\mathrm{w}}$, and let $B$, $B'$ be complex bond functions on $T_{*}$, with the numbers $\epsilon_B$, $\epsilon_{B'}$ and the sets $P_B$, $P_B^{+}$, $P_{B'}$, $P_{B'}^{+}$ of Section~\ref{sec:dependence-stop-level}; for \eqref{eq:four-jets-Leg3} and \eqref{eq:four-jets-Dpp} assume moreover that $d(P_B,P_{B'})\ge M+2$. Let $E$, the majorants $|O|(y,y')$, the norms $\|O\|_{\tilde{\gamma}}$ and the derivatives $\partial_B$ be those of Section~\ref{sec:fermionic-operators-stop}, and let $C_F$, $\gamma_F$, $h_L$, $C_H$, $C_K$ be the constants of Section~\ref{sec:constants-3}. Then for every leg index $\boldsymbol{\upsilon}$ and all unit sites $y,y'\in T^0_{*}$ the following bounds hold:
\begin{multline}
|\mathfrak{H}_A|(\boldsymbol{\upsilon};y')\le3\,C_F\sum_{w\in\mathrm{pos}(\boldsymbol{\upsilon})}E(w,y'),\\
|\mathfrak{H}_A|(\boldsymbol{\upsilon};y')\le h_L\cdot\exp\left(-\gamma_F\cdot d\left(\mathrm{pos}(\boldsymbol{\upsilon}),y'\right)\right),\\
\sum_u\left|\mathfrak{H}_A(\boldsymbol{\upsilon};u)\right|\cdot\exp\left(\tilde{\gamma}\cdot d\left(\mathrm{pos}(\boldsymbol{\upsilon}),u\right)\right)\le h_L\quad\text{for }0\le\tilde{\gamma}\le\frac{\gamma_F}{2},\tag{Leg1}\label{eq:four-jets-Leg1}
\end{multline}
where the sum is taken over all indices $u=(y',\boldsymbol{\beta}',\boldsymbol{\omega}')$ of the generators $\Psi(u)$;
\begin{equation}
|\partial_B\mathfrak{H}_A|(\boldsymbol{\upsilon};y')\le C_H\cdot\epsilon_B\sum_{w\in\mathrm{pos}(\boldsymbol{\upsilon})}\ \sum_{z\in P_B^{+}}E(w,z)\cdot E(z,y');\tag{Leg2}\label{eq:four-jets-Leg2}
\end{equation}
\begin{multline}
|\partial_B\partial_{B'}\mathfrak{H}_A|(\boldsymbol{\upsilon};y')\le C_H^2\cdot\epsilon_B\cdot\epsilon_{B'}\sum_{w\in\mathrm{pos}(\boldsymbol{\upsilon})}\ \sum_{z\in P_B^{+}}\ \sum_{z'\in P_{B'}^{+}}[E(w,z)\cdot E(z,z')\cdot E(z',y')\\
+E(w,z')\cdot E(z',z)\cdot E(z,y')];\tag{Leg3}\label{eq:four-jets-Leg3}
\end{multline}
\begin{multline}
|\partial_BD_{*}(A)|(y,y')\le67\,C_F^2\cdot\epsilon_B\sum_{z\in P_B}E(y,z)\cdot E(z,y'),\\
\|\partial_BD_{*}(A)\|_{\frac{\gamma_F}{2}}\le C_K\cdot\epsilon_B,\qquad\sum_{y,y'}|\partial_BD_{*}(A)|(y,y')\le C_K\cdot\epsilon_B\cdot\#|P_B|;\tag{D$'$}\label{eq:four-jets-Dp}
\end{multline}
\begin{multline}
|\partial_B\partial_{B'}D_{*}(A)|(y,y')\le C_H^2\cdot\epsilon_B\cdot\epsilon_{B'}\sum_{z\in P_B}\ \sum_{z'\in P_{B'}}[E(y,z)\cdot E(z,z')\cdot E(z',y')\\
+E(y,z')\cdot E(z',z)\cdot E(z,y')].\tag{D$''$}\label{eq:four-jets-Dpp}
\end{multline}
\end{lemma}

\begin{proof}
\emph{Notation and preliminaries.} Throughout the proof
\begin{gather*}
Q:=Q_{*}(A),\qquad\tilde{Q}:=Q_{*}^T(-A),\qquad S:=S(A),\qquad H:=H_{*}(A)=b_{*}\cdot S\cdot\tilde{Q},\\
\dot{\mathsf{D}}:=\partial_B\mathsf{D}(A),\qquad\dot{\mathsf{D}}':=\partial_{B'}\mathsf{D}(A),
\end{gather*}
and $\partial_B\tilde{Q}$, $\partial_BQ$, $\partial_B\delta_{\alpha,A}$ denote the derivatives of the maps $A'\mapsto Q_{*}^T(-A')$, $A'\mapsto Q_{*}(A')$, $A'\mapsto\delta_{\alpha,A'}$ at the point $A$, in the direction $B$; the same applies to $B'$. One has $b_{*}\le1$ by Lemma~\ref{lem:scales-stop-level}(a), and $C_F\ge1$, $\gamma_F\le1$, $M\ge1$ by Section~\ref{sec:constants-3}. By the facts stated at the head of Section~\ref{sec:dependence-stop-level}, $\mathfrak{i}(-A)=\mathfrak{i}(A)<\frac{1}{16}$ and $3M\cdot\mathfrak{i}(A)<\frac{1}{4}$; so Lemma~\ref{lem:bounds-elementary-operators} applies at $A$ and at $-A$. The set $G_{\mathrm{w}}$ is contained in $\tilde{G}_{*}$ (Section~\ref{sec:dependence-stop-level}), which is open (Lemma~\ref{lem:openness}), and $S$ is analytic on $\tilde{G}_{*}$ by \eqref{it:trajectory-fermionic-operators-SL2}. Differentiating the identity $S(A')\cdot\mathsf{D}(A')=1$ at $A'=A$ in the direction $B$ one obtains
\[
\partial_BS=-S\cdot\dot{\mathsf{D}}\cdot S.
\]

\emph{Barred legs.} Let $A'\in\tilde{G}_{*}$. By the definitions of Section~\ref{sec:fermionic-operators-stop}, $\bar{H}_{*}(A')=\mathsf{C}^{-1}\cdot H_{*}(-A')\cdot\mathsf{C}$. The operator $\delta_{\alpha,-A'}$ acts on the four spinor components with the same scalar coefficients, so it commutes with $\mathsf{C}^{-1}$, and
\[
\delta_{\alpha,-A'}\bar{H}_{*}(A')=\mathsf{C}^{-1}\cdot\left[\delta_{\alpha,A''}H_{*}(A'')\right]_{A''=-A'}\cdot\mathsf{C}.
\]
The point $-A$ lies in $G_{\mathrm{w}}$, because $G_{\mathrm{w}}$ is invariant under $A\mapsto-A$ (Section~\ref{sec:dependence-stop-level}). The derivative at $A$, in the direction $B$, of a function of $-A$ is the derivative of that function at $-A$ in the direction $-B$; moreover $\epsilon_{-B}=\epsilon_B$, $P_{-B}=P_B$ and $P_{-B}^{+}=P_B^{+}$, and the same holds for $B'$. The matrices $\mathsf{C}$ and $\mathsf{C}^{-1}$ are monomial, with non-zero entries of modulus one (Proposition~\ref{prop:symmetries-extraction}(b); Section~\ref{sec:gamma-matrices-symmetry}). Hence, for a $4\times4$ matrix $O$, each row of $\mathsf{C}^{-1}\cdot O\cdot\mathsf{C}$ arises from a row of $O$ by a permutation of its entries and by multiplication of the entries with numbers of modulus one, and the two rows have the same sum of moduli. Consequently each of the numbers $|\mathfrak{H}_A|(\boldsymbol{\upsilon};y')$, $|\partial_B\mathfrak{H}_A|(\boldsymbol{\upsilon};y')$, $|\partial_B\partial_{B'}\mathfrak{H}_A|(\boldsymbol{\upsilon};y')$ for a leg index with $\boldsymbol{\omega}=1$ equals the corresponding number at the point $-A$, in the directions $-B$, $-B'$, for a leg index with $\boldsymbol{\omega}=0$, with the same position and another spinor index. The right sides of \eqref{eq:four-jets-Leg1}--\eqref{eq:four-jets-Leg3} are not altered by these substitutions, nor is the hypothesis $d(P_B,P_{B'})\ge M+2$. It therefore suffices to prove \eqref{eq:four-jets-Leg1}--\eqref{eq:four-jets-Leg3} for leg indices with $\boldsymbol{\omega}=0$, which is assumed from now on.

\emph{Coefficients and majorants.} Fix $y'\in T^0_{*}$. For a spinor function $v$ on $T^0_{*}$ with $|v|_{\infty}\le1$ put
\[
\check{g}:=H\cdot\mathbf{1}_{y'}\cdot v,
\]
a spinor function on $T_{*}$ that depends on $A$; then $\partial_B\check{g}=\left(\partial_BH\right)\cdot\mathbf{1}_{y'}\cdot v$ and $\partial_B\partial_{B'}\check{g}=\left(\partial_B\partial_{B'}H\right)\cdot\mathbf{1}_{y'}\cdot v$. For $\boldsymbol{\upsilon}=\xi=(x,\boldsymbol{\beta},0)$ the component $\boldsymbol{\beta}$ of $\check{g}(x)$ equals $\sum_{\boldsymbol{\beta}'}\mathfrak{H}_A\left(\xi;(y',\boldsymbol{\beta}',0)\right)\cdot v_{\boldsymbol{\beta}'}(y')$, and for $\boldsymbol{\upsilon}=\boldsymbol{\zeta}=(x,x',\boldsymbol{\beta},0)$ the component $\boldsymbol{\beta}$ of $(\delta_{\alpha,A}\check{g})(x,x')$ equals $\sum_{\boldsymbol{\beta}'}\mathfrak{H}_A\left(\boldsymbol{\zeta};(y',\boldsymbol{\beta}',0)\right)\cdot v_{\boldsymbol{\beta}'}(y')$; the derivatives of these identities in the directions $B$, $B'$ hold as well. As in the proof of \eqref{eq:four-jets-Maj2}, the supremum of the modulus of such a sum over the $v$ with $|v|_{\infty}\le1$ is the sum of the moduli of its coefficients. On the other hand, for each linear operator $O$ that occurs below, $\left|\left(O\cdot\mathbf{1}_{y'}\cdot v\right)(x)\right|\le|O|([x],y')$ by the definition of the majorant, and likewise at a pair $(x,x')$, which is localised in the cube of $x$ (Section~\ref{sec:fermionic-operators-stop}). In particular
\begin{gather*}
|\mathfrak{H}_A|(\xi;y')\le|H|([x],y'),\qquad|\partial_B\mathfrak{H}_A|(\xi;y')\le|\partial_BH|([x],y'),\\
|\partial_B\partial_{B'}\mathfrak{H}_A|(\xi;y')\le|\partial_B\partial_{B'}H|([x],y').
\end{gather*}

\emph{The bound \eqref{eq:four-jets-Leg1}.} If $\boldsymbol{\upsilon}=\xi$, then $|\mathfrak{H}_A|(\xi;y')\le|H|([x],y')\le C_F\cdot E([x],y')$ by \eqref{it:trajectory-fermionic-operators-SL2}. If $\boldsymbol{\upsilon}=\boldsymbol{\zeta}$ with $|x-x'|<1$, the bound of \eqref{it:trajectory-fermionic-operators-SL2} for $\delta_{\alpha,A}H_{*}(A)$ on these pairs gives in the same way $|\mathfrak{H}_A|(\boldsymbol{\zeta};y')\le|\delta_{\alpha,A}H|([x],y')\le C_F\cdot E([x],y')$. In both cases $\mathrm{pos}(\boldsymbol{\upsilon})=\{[x]\}$. Let now $\boldsymbol{\upsilon}=\boldsymbol{\zeta}$ with $|x-x'|\ge1$. Two points of one cube $\Delta_y$ have distance less than $1$; hence $[x]\ne[x']$, and $\mathrm{pos}(\boldsymbol{\zeta})$ consists of the two sites $[x]$, $[x']$. The last assertion of \eqref{it:trajectory-fermionic-operators-SL2} and the bound for $|H|$ give
\[
\left|(\delta_{\alpha,A}\check{g})(x,x')\right|\le\exp(1)\cdot|\check{g}(x')|+|\check{g}(x)|\le\exp(1)\,C_F\cdot E([x'],y')+C_F\cdot E([x],y'),
\]
and the supremum over $v$ yields $|\mathfrak{H}_A|(\boldsymbol{\zeta};y')\le\exp(1)\,C_F\cdot E([x'],y')+C_F\cdot E([x],y')$. As $\exp(1)<3$, the first inequality of \eqref{eq:four-jets-Leg1} holds in all cases.

The set $\mathrm{pos}(\boldsymbol{\upsilon})$ has at most two elements, and $E(w,y')\le\exp\left(-\gamma_F\cdot d(\mathrm{pos}(\boldsymbol{\upsilon}),y')\right)$ for $w\in\mathrm{pos}(\boldsymbol{\upsilon})$. Since $\gamma_F\le1$, one has $\mathrm{ssum}\left(\frac{\gamma_F}{2}\right)=\left(\frac{24}{\gamma_F}\right)^3\ge1$, whence
\[
6\,C_F\le6\,C_F\cdot\mathrm{ssum}\left(\frac{\gamma_F}{2}\right)=6\,\mathbf{S}=h_L .
\]
So the first inequality gives $|\mathfrak{H}_A|(\boldsymbol{\upsilon};y')\le6\,C_F\cdot\exp\left(-\gamma_F\cdot d(\mathrm{pos}(\boldsymbol{\upsilon}),y')\right)$, and the second inequality of \eqref{eq:four-jets-Leg1} follows. Let $0\le\tilde{\gamma}\le\frac{\gamma_F}{2}$. In the sum over $u$ only the indices $u=(y',\boldsymbol{\beta}',\boldsymbol{\omega})$ have non-zero coefficients, and $d(\mathrm{pos}(\boldsymbol{\upsilon}),u)=d(\mathrm{pos}(\boldsymbol{\upsilon}),y')$ for them. For $w\in\mathrm{pos}(\boldsymbol{\upsilon})$ one has $d(\mathrm{pos}(\boldsymbol{\upsilon}),y')\le d(w,y')$, and $\gamma_F-\tilde{\gamma}\ge\frac{\gamma_F}{2}$. By the first inequality of \eqref{eq:four-jets-Leg1} and the bound $\sum_{y'}\exp\left(-\tilde{\gamma}'\cdot d(w,y')\right)\le\mathrm{ssum}(\tilde{\gamma}')$ for $0<\tilde{\gamma}'\le6$ of Section~\ref{sec:constants-3}, taken at $\tilde{\gamma}'=\frac{\gamma_F}{2}$,
\begin{align*}
\sum_u\left|\mathfrak{H}_A(\boldsymbol{\upsilon};u)\right|\cdot\exp\left(\tilde{\gamma}\cdot d\left(\mathrm{pos}(\boldsymbol{\upsilon}),u\right)\right)&=\sum_{y'}|\mathfrak{H}_A|(\boldsymbol{\upsilon};y')\cdot\exp\left(\tilde{\gamma}\cdot d\left(\mathrm{pos}(\boldsymbol{\upsilon}),y'\right)\right)\\
&\le3\,C_F\sum_{w\in\mathrm{pos}(\boldsymbol{\upsilon})}\ \sum_{y'}\exp\left(-\left(\gamma_F-\tilde{\gamma}\right)\cdot d(w,y')\right)\\
&\le3\,C_F\cdot2\cdot\mathrm{ssum}\left(\frac{\gamma_F}{2}\right)=h_L.
\end{align*}
This is the third inequality of \eqref{eq:four-jets-Leg1}.

\emph{First derivatives of $H$.} The product rule and the formula for $\partial_BS$ give
\[
\partial_BH=-S\cdot\dot{\mathsf{D}}\cdot H+b_{*}\cdot S\cdot\partial_B\tilde{Q} .
\]
By \eqref{eq:four-jets-Maj1}, \eqref{it:trajectory-fermionic-operators-SL2} and \eqref{eq:four-jets-12},
\[
|S\cdot\dot{\mathsf{D}}\cdot H|(y,y')\le C_F^2\sum_{z,z'}E(y,z)\cdot|\dot{\mathsf{D}}|(z,z')\cdot E(z',y')\le40\,C_F^2\cdot\epsilon_B\sum_{z\in P_B}E(y,z)\cdot E(z,y'),
\]
and by \eqref{eq:four-jets-Maj1}, \eqref{it:trajectory-fermionic-operators-SL2}, Lemma~\ref{lem:bounds-elementary-operators}(a) and $b_{*}\le1$,
\[
|b_{*}\cdot S\cdot\partial_B\tilde{Q}|(y,y')\le\frac{9}{4}\,C_F\cdot\epsilon_B\cdot E(y,y')\cdot\mathbf{1}[y'\in P_B];
\]
since $E(y',y')=1$, the right side is $\frac{9}{4}\,C_F\cdot\epsilon_B$ times the term $z=y'$ of the sum $\sum_{z\in P_B}E(y,z)\cdot E(z,y')$. On the pairs with $|x-x'|<1$ one has
\[
\delta_{\alpha,A}\,\partial_BH=-\mathrm{Sa}\cdot\dot{\mathsf{D}}\cdot H+b_{*}\cdot\mathrm{Sa}\cdot\partial_B\tilde{Q},\qquad\mathrm{Sa}:=\mathrm{Sa}(A),
\]
and $|\mathrm{Sa}|(y,y')\le C_F\cdot E(y,y')$ by \eqref{it:trajectory-fermionic-operators-SL2}; so the same two estimates hold with $\mathrm{Sa}$ in place of the first factor $S$. Since $C_F\ge1$, it follows that
\begin{multline}
|\partial_BH|(y,y'),\ |\delta_{\alpha,A}\,\partial_BH|(y,y')\le C_1\cdot\epsilon_B\sum_{z\in P_B}E(y,z)\cdot E(z,y'),\\
C_1:=40\,C_F^2+\frac{9}{4}\,C_F\le43\,C_F^2,\label{eq:four-jets-13}
\end{multline}
the second bound on the pairs with $|x-x'|<1$.

\emph{The bound \eqref{eq:four-jets-Leg2}.} For $\boldsymbol{\upsilon}=\xi$, the inequality $|\partial_B\mathfrak{H}_A|(\xi;y')\le|\partial_BH|([x],y')$ and \eqref{eq:four-jets-13} give \eqref{eq:four-jets-Leg2} with $w=[x]$, because $C_1\le43\,C_F^2\le C_H$ and $P_B\subset P_B^{+}$.

Let $\boldsymbol{\upsilon}=\boldsymbol{\zeta}$ with $|x-x'|<1$. Then
\[
\partial_B\left(\delta_{\alpha,A}\check{g}\right)=\left(\partial_B\delta_{\alpha,A}\right)\check{g}+\delta_{\alpha,A}\,\partial_B\check{g} .
\]
By Lemma~\ref{lem:bounds-elementary-operators}(c), $\left|\left(\left(\partial_B\delta_{\alpha,A}\right)\check{g}\right)(x,x')\right|\le4\,\epsilon_B\cdot|\check{g}(x')|\le4\,C_F\cdot\epsilon_B\cdot E([x'],y')$, and this term vanishes unless $d([x],P_B)\le1$, in particular unless $[x]\in P_B^{+}$. The unit sites $[x]$ and $[x']$ have distance less than $2$, hence at most $1$; so $E([x'],y')\le\exp(\gamma_F)\cdot E([x],y')\le\exp(1)\cdot E([x],y')$, and the first term is at most
\begin{align*}
&4\exp(1)\,C_F\cdot\epsilon_B\cdot E([x],[x])\cdot E([x],y')\cdot\mathbf{1}[[x]\in P_B^{+}]\\
&\qquad\le11\,C_F^2\cdot\epsilon_B\cdot E([x],[x])\cdot E([x],y')\cdot\mathbf{1}[[x]\in P_B^{+}],
\end{align*}
that is, at most $11\,C_F^2\cdot\epsilon_B$ times the term $w=z=[x]$ of the double sum in \eqref{eq:four-jets-Leg2}. By \eqref{eq:four-jets-13}, the second term is at most $C_1\cdot\epsilon_B\sum_{z\in P_B}E([x],z)\cdot E(z,y')$. Every term of the double sum thus receives a coefficient at most $(43+11)\,C_F^2\le C_H$, and the supremum over $v$ gives \eqref{eq:four-jets-Leg2}.

Let $\boldsymbol{\upsilon}=\boldsymbol{\zeta}$ with $|x-x'|\ge1$. By the definition of $\delta_{\alpha,A}$ in Section~\ref{sec:fermionic-operators-stop},
\[
\partial_B\left(\delta_{\alpha,A}\check{g}\right)(x,x')=|x-x'|^{-\alpha}\cdot\left[\left(\partial_BU_A\right)(x,x')\cdot\check{g}(x')+U_A(x,x')\cdot\partial_B\check{g}(x')-\partial_B\check{g}(x)\right].
\]
The first summand is $\left(\left(\partial_B\delta_{\alpha,A}\right)\check{g}\right)(x,x')$; by Lemma~\ref{lem:bounds-elementary-operators}(c) its modulus is at most $4M\cdot\epsilon_B\cdot|\check{g}(x')|\le4M\cdot C_F\cdot\epsilon_B\cdot E([x'],y')$, and it vanishes unless $d([x'],P_B)\le M$, in particular unless $[x']\in P_B^{+}$. In the two other summands, $|x-x'|^{-\alpha}\le1$ and $|U_A(x,x')|\le\frac{4}{3}$ (Lemma~\ref{lem:bounds-elementary-operators}(c)), and \eqref{eq:four-jets-13} gives
\[
|\partial_B\check{g}(x')|\le C_1\cdot\epsilon_B\sum_{z\in P_B}E([x'],z)\cdot E(z,y'),\qquad|\partial_B\check{g}(x)|\le C_1\cdot\epsilon_B\sum_{z\in P_B}E([x],z)\cdot E(z,y').
\]
Hence, in the double sum of \eqref{eq:four-jets-Leg2}, the terms with $w=[x']$ receive coefficients at most
\[
4M\cdot C_F+\frac{4}{3}\cdot43\,C_F^2\le62\,M\cdot C_F^2\le C_H,
\]
the first summand contributing to the term $w=z=[x']$, and the terms with $w=[x]$ coefficients at most $C_1\le C_H$. The supremum over $v$ gives \eqref{eq:four-jets-Leg2}.

\emph{Second derivatives of $H$.} Let $d(P_B,P_{B'})\ge M+2$. By Lemma~\ref{lem:bounds-elementary-operators}(d), the mixed derivatives $\partial_B\partial_{B'}$ of $\mathsf{D}(\cdot)$, $Q_{*}(\cdot)$, $Q_{*}^T(-\,\cdot)$ and $\delta_{\alpha,\cdot}$ vanish at $A$. Differentiating $\partial_BS=-S\cdot\dot{\mathsf{D}}\cdot S$ in the direction $B'$, and using $\partial_{B'}\dot{\mathsf{D}}=\partial_B\partial_{B'}\mathsf{D}(A)=0$, one obtains
\[
\partial_B\partial_{B'}S=S\cdot\dot{\mathsf{D}}'\cdot S\cdot\dot{\mathsf{D}}\cdot S+S\cdot\dot{\mathsf{D}}\cdot S\cdot\dot{\mathsf{D}}'\cdot S;
\]
with $H=b_{*}\cdot S\cdot\tilde{Q}$ and $\partial_B\partial_{B'}\tilde{Q}=0$ it follows that
\begin{multline*}
\partial_B\partial_{B'}H=S\cdot\dot{\mathsf{D}}'\cdot S\cdot\dot{\mathsf{D}}\cdot H+S\cdot\dot{\mathsf{D}}\cdot S\cdot\dot{\mathsf{D}}'\cdot H\\
-b_{*}\cdot S\cdot\dot{\mathsf{D}}\cdot S\cdot\partial_{B'}\tilde{Q}-b_{*}\cdot S\cdot\dot{\mathsf{D}}'\cdot S\cdot\partial_B\tilde{Q} .
\end{multline*}
For unit sites $y,z,z',y'$ put
\[
\boldsymbol{\chi}(y;z,z';y'):=E(y,z)\cdot E(z,z')\cdot E(z',y').
\]
By \eqref{eq:four-jets-Maj1}, \eqref{it:trajectory-fermionic-operators-SL2} and two applications of \eqref{eq:four-jets-12},
\[
|S\cdot\dot{\mathsf{D}}\cdot S\cdot\dot{\mathsf{D}}'\cdot H|(y,y')\le1600\,C_F^3\cdot\epsilon_B\cdot\epsilon_{B'}\sum_{z\in P_B}\ \sum_{z'\in P_{B'}}\boldsymbol{\chi}(y;z,z';y'),
\]
and by \eqref{eq:four-jets-Maj1}, \eqref{it:trajectory-fermionic-operators-SL2}, \eqref{eq:four-jets-12}, Lemma~\ref{lem:bounds-elementary-operators}(a) and $b_{*}\le1$,
\[
|b_{*}\cdot S\cdot\dot{\mathsf{D}}\cdot S\cdot\partial_{B'}\tilde{Q}|(y,y')\le90\,C_F^2\cdot\epsilon_B\cdot\epsilon_{B'}\sum_{z\in P_B}\boldsymbol{\chi}(y;z,y';y')\cdot\mathbf{1}[y'\in P_{B'}];
\]
the right side of the last bound is $90\,C_F^2\cdot\epsilon_B\cdot\epsilon_{B'}$ times the part $z'=y'$ of the double sum $\sum_{z\in P_B}\sum_{z'\in P_{B'}}\boldsymbol{\chi}(y;z,z';y')$. The same two bounds hold with $B$ and $B'$ interchanged, with $\boldsymbol{\chi}(y;z',z;y')$ in place of $\boldsymbol{\chi}(y;z,z';y')$. On the pairs with $|x-x'|<1$, the application of $\delta_{\alpha,A}$ replaces the first factor $S$ of each of the four terms by $\mathrm{Sa}$, which obeys the same bound. Since $1600\,C_F^3+90\,C_F^2\le1690\,C_F^3$, one obtains
\begin{multline}
|\partial_B\partial_{B'}H|(y,y'),\ |\delta_{\alpha,A}\,\partial_B\partial_{B'}H|(y,y')\\
\le C_2\cdot\epsilon_B\cdot\epsilon_{B'}\sum_{z\in P_B}\ \sum_{z'\in P_{B'}}[\boldsymbol{\chi}(y;z,z';y')+\boldsymbol{\chi}(y;z',z;y')],\qquad C_2:=1690\,C_F^3,\label{eq:four-jets-14}
\end{multline}
the second bound on the pairs with $|x-x'|<1$.

\emph{The bound \eqref{eq:four-jets-Leg3}.} Let $d(P_B,P_{B'})\ge M+2$. The products $E(w,z)\cdot E(z,z')\cdot E(z',y')$ and $E(w,z')\cdot E(z',z)\cdot E(z,y')$ in \eqref{eq:four-jets-Leg3} are called the products of the first and of the second kind. For $\boldsymbol{\upsilon}=\xi$, the inequality $|\partial_B\partial_{B'}\mathfrak{H}_A|(\xi;y')\le|\partial_B\partial_{B'}H|([x],y')$ and \eqref{eq:four-jets-14} give \eqref{eq:four-jets-Leg3} with $w=[x]$, since $C_2\le C_H^2$, $P_B\subset P_B^{+}$ and $P_{B'}\subset P_{B'}^{+}$. For pairs, Lemma~\ref{lem:bounds-elementary-operators}(d) gives $\partial_B\partial_{B'}\delta_{\alpha,A}=0$, and hence
\[
\partial_B\partial_{B'}\left(\delta_{\alpha,A}\check{g}\right)=\left(\partial_B\delta_{\alpha,A}\right)\partial_{B'}\check{g}+\left(\partial_{B'}\delta_{\alpha,A}\right)\partial_B\check{g}+\delta_{\alpha,A}\,\partial_B\partial_{B'}\check{g} .
\]
Let $|x-x'|<1$. By Lemma~\ref{lem:bounds-elementary-operators}(c), \eqref{eq:four-jets-13} and the inequality $E([x'],z')\le\exp(1)\cdot E([x],z')$, the first term is bounded at $(x,x')$ by
\begin{align*}
&4\,\epsilon_B\cdot\mathbf{1}[[x]\in P_B^{+}]\cdot C_1\cdot\epsilon_{B'}\sum_{z'\in P_{B'}}E([x'],z')\cdot E(z',y')\\
&\qquad\le11\,C_1\cdot\epsilon_B\cdot\epsilon_{B'}\cdot\mathbf{1}[[x]\in P_B^{+}]\sum_{z'\in P_{B'}}\boldsymbol{\chi}([x];[x],z';y'),
\end{align*}
that is, by $11\,C_1\cdot\epsilon_B\cdot\epsilon_{B'}$ times products of the first kind with $w=z=[x]$. In the same way the second term is bounded by $11\,C_1\cdot\epsilon_B\cdot\epsilon_{B'}$ times products of the second kind with $w=z'=[x]$, and the third term is bounded by \eqref{eq:four-jets-14}. Let $|x-x'|\ge1$. By Lemma~\ref{lem:bounds-elementary-operators}(c) and $|x-x'|^{-\alpha}\le1$, the moduli of the three terms at $(x,x')$ are at most
\begin{gather*}
4M\cdot\epsilon_B\cdot\mathbf{1}[[x']\in P_B^{+}]\cdot|\partial_{B'}\check{g}(x')|,\qquad4M\cdot\epsilon_{B'}\cdot\mathbf{1}[[x']\in P_{B'}^{+}]\cdot|\partial_B\check{g}(x')|,\\
\frac{4}{3}\,|\partial_B\partial_{B'}\check{g}(x')|+|\partial_B\partial_{B'}\check{g}(x)|,
\end{gather*}
and \eqref{eq:four-jets-13}, \eqref{eq:four-jets-14} bound these numbers at the unit sites $[x']$ and $[x]$; the first term gives products of the first kind with $w=z=[x']$, the second products of the second kind with $w=z'=[x']$. In all cases, after the supremum over $v$, every product of either kind receives a coefficient at most
\[
4M\cdot C_1+11\,C_1+\frac{4}{3}\,C_2\le172\,M\cdot C_F^2+473\,C_F^2+2254\,C_F^3\le10^4\,M^2\cdot C_F^4=C_H^2,
\]
which proves \eqref{eq:four-jets-Leg3}.

\emph{The bound \eqref{eq:four-jets-Dp}.} By the definition of $D_{*}$ in Section~\ref{sec:fermionic-operators-stop}, $D_{*}(A)=b_{*}-b_{*}\cdot Q\cdot H$; hence
\[
\partial_BD_{*}(A)=-b_{*}\cdot\left(\partial_BQ\right)\cdot H-b_{*}\cdot Q\cdot\partial_BH .
\]
By \eqref{eq:four-jets-Maj1}, Lemma~\ref{lem:bounds-elementary-operators}(a), \eqref{it:trajectory-fermionic-operators-SL2}, \eqref{eq:four-jets-13} and $b_{*}\le1$,
\[
|\partial_BD_{*}(A)|(y,y')\le\frac{9}{4}\,\epsilon_B\cdot\mathbf{1}[y\in P_B]\cdot C_F\cdot E(y,y')+\frac{3}{2}\,C_1\cdot\epsilon_B\sum_{z\in P_B}E(y,z)\cdot E(z,y');
\]
the first term is $\frac{9}{4}\,C_F\cdot\epsilon_B$ times the term $z=y$ of the sum. Since
\[
\frac{9}{4}\,C_F+\frac{3}{2}\cdot43\,C_F^2\le\left(\frac{9}{4}+64.5\right)C_F^2\le67\,C_F^2,
\]
the first inequality of \eqref{eq:four-jets-Dp} holds. Let $0\le\tilde{\gamma}\le\frac{\gamma_F}{2}$. As $d(y,y')\le d(y,z)+d(z,y')$,
\[
\exp\left(\tilde{\gamma}\cdot d(y,y')\right)\cdot E(y,z)\cdot E(z,y')\le\exp\left(-\frac{\gamma_F}{2}\cdot\left(d(y,z)+d(z,y')\right)\right),
\]
and the sum of the right side over $z$ and $y'$ at fixed $y$, or over $z$ and $y$ at fixed $y'$, is at most $\mathrm{ssum}\left(\frac{\gamma_F}{2}\right)^2$ (Section~\ref{sec:constants-3}). Hence
\[
\|\partial_BD_{*}(A)\|_{\tilde{\gamma}}\le67\,C_F^2\cdot\mathrm{ssum}\left(\frac{\gamma_F}{2}\right)^2\cdot\epsilon_B=67\,\mathbf{S}^2\cdot\epsilon_B\le C_K\cdot\epsilon_B
\]
for these $\tilde{\gamma}$, in particular for $\tilde{\gamma}=\frac{\gamma_F}{2}$. At fixed $z$, the sum of $E(y,z)\cdot E(z,y')$ over $y$ and $y'$ is at most $\mathrm{ssum}(\gamma_F)^2\le\mathrm{ssum}\left(\frac{\gamma_F}{2}\right)^2$, the function $\mathrm{ssum}$ being decreasing; summation over $z\in P_B$ gives
\[
\sum_{y,y'}|\partial_BD_{*}(A)|(y,y')\le67\,\mathbf{S}^2\cdot\epsilon_B\cdot\#|P_B|\le C_K\cdot\epsilon_B\cdot\#|P_B|.
\]

\emph{The bound \eqref{eq:four-jets-Dpp}.} Let $d(P_B,P_{B'})\ge M+2$. As $\partial_B\partial_{B'}Q=0$ (Lemma~\ref{lem:bounds-elementary-operators}(d)),
\[
\partial_B\partial_{B'}D_{*}(A)=-b_{*}\cdot\left[\left(\partial_BQ\right)\cdot\partial_{B'}H+\left(\partial_{B'}Q\right)\cdot\partial_BH+Q\cdot\partial_B\partial_{B'}H\right].
\]
By \eqref{eq:four-jets-Maj1}, Lemma~\ref{lem:bounds-elementary-operators}(a), \eqref{eq:four-jets-13} and \eqref{eq:four-jets-14}, the first term is at most
\[
\frac{9}{4}\,C_1\cdot\epsilon_B\cdot\epsilon_{B'}\cdot\mathbf{1}[y\in P_B]\sum_{z'\in P_{B'}}E(y,y)\cdot E(y,z')\cdot E(z',y'),
\]
a multiple of products $E(y,z)\cdot E(z,z')\cdot E(z',y')$ of \eqref{eq:four-jets-Dpp} with $z=y$; the second term is bounded in the same way by multiples of products $E(y,z')\cdot E(z',z)\cdot E(z,y')$ with $z'=y$; the third term is at most $\frac{3}{2}\,C_2\cdot\epsilon_B\cdot\epsilon_{B'}$ times the double sum of \eqref{eq:four-jets-Dpp}. So each product in \eqref{eq:four-jets-Dpp} receives a coefficient at most
\[
\frac{9}{4}\,C_1+\frac{3}{2}\,C_2\le97\,C_F^2+2535\,C_F^3\le C_H^2.
\]
\end{proof}

\begin{corollary}[Increments of $D_{*}$ along a segment in $G_{\mathrm{w}}$]\label{cor:increments-along-segment}
Let $A_0$ and $B$ be complex bond functions on $T_{*}$ such that $A_0+u\cdot B\in G_{\mathrm{w}}$ for every $u\in[0,1]$. Then
\begin{gather*}
\|D_{*}(A_0+B)-D_{*}(A_0)\|_{\frac{\gamma_F}{2}}\le C_K\cdot\epsilon_B,\\
\sum_{y,y'}|D_{*}(A_0+B)-D_{*}(A_0)|(y,y')\le C_K\cdot\epsilon_B\cdot\#|P_B|.
\end{gather*}
\end{corollary}

\begin{proof}
The open set $\tilde{G}_{*}$ contains the segment, and $D_{*}$ is analytic on $\tilde{G}_{*}$, because $S$ is analytic there by \eqref{it:trajectory-fermionic-operators-SL2} and each entry of $Q_{*}(A)$ and of $Q_{*}^T(-A)$ is either zero or a constant times the exponential of a linear function of $A$ (Section~\ref{sec:fermionic-operators-stop}). By the fundamental theorem of calculus,
\[
D_{*}(A_0+B)-D_{*}(A_0)=\int_0^1\mathrm{d}u\,\left(\partial_BD_{*}\right)(A_0+u\cdot B).
\]
For every $u\in[0,1]$ the bounds of \eqref{eq:four-jets-Dp} hold at the point $A_0+u\cdot B\in G_{\mathrm{w}}$, with the same $\epsilon_B$ and $P_B$. The majorant of the integral of an operator-valued function is at most the integral of the majorants, $\left|\int_0^1\mathrm{d}u\,O(u)\right|(y,y')\le\int_0^1\mathrm{d}u\,|O(u)|(y,y')$; consequently the norm $\|\cdot\|_{\frac{\gamma_F}{2}}$ and the sum over $(y,y')$ of the majorant of the integral are at most the integrals over $u$ of the same quantities for $\left(\partial_BD_{*}\right)(A_0+u\cdot B)$. The two bounds of \eqref{eq:four-jets-Dp} thus give the assertion.
\end{proof}

\begin{lemma}[The operators on $G_{\mathrm{w}}^{+}$]\label{lem:operators}
Let $A\in G_{\mathrm{w}}$, and let $A_1$ be a complex bond function on $T_{*}$ with $e_{*}\cdot|A_1|_{\infty}\le c_1$ (every element of $G_{\mathrm{w}}^{+}$ is of the form $A+A_1$ with such $A$ and $A_1$). Then the operator $\mathsf{D}(A+A_1)$ is invertible, and its inverse $S(A+A_1)$ satisfies
\[
\left\|S(A+A_1)\right\|_{\gamma_F/2}\le 2\,\mathbf{S} .
\]
For $A'\in G_{\mathrm{w}}^{+}$ let $D_{*}(A')$ be given by the formula of Section~\ref{sec:fermionic-operators-stop}, that is,
\[
D_{*}(A')=b_{*}-b_{*}^2\cdot Q_{*}(A')\cdot S(A')\cdot Q_{*}^T(-A').
\]
Let $A'\in G_{\mathrm{w}}^{+}$, and let $B$ be a complex bond function on $T_{*}$ such that $A'+u\cdot B\in G_{\mathrm{w}}^{+}$ for every $u\in[0,1]$. Then
\[
\left\|D_{*}(A'+B)-D_{*}(A')\right\|_{\gamma_F/2}\le C_K'\cdot\epsilon_B
\qquad\text{and}\qquad
\sum_{y,y'\in T^0_{*}}\left|D_{*}(A'+B)-D_{*}(A')\right|(y,y')\le C_K'\cdot\epsilon_B\cdot\#|P_B|.
\]
\end{lemma}

\begin{proof}
For $0\le u\le 1$ the connection $A+u\cdot A_1$ lies in $G_{\mathrm{w}}^{+}$. By the facts listed at the beginning of Section~\ref{sec:dependence-stop-level}, and since $c_1\le\frac{1}{8}$ (Section~\ref{sec:constants-3}),
\[
\mathfrak{i}(A+u\cdot A_1)\le\mathfrak{i}(A)+e_{*}\cdot|A_1|_{\infty}\le\frac{1}{16}+\frac{1}{8}<\frac{1}{4};
\]
by the same facts, which rest on~\eqref{it:trajectory-fermionic-operators-SL2} and on the definition of $\mathrm{ssum}$, one has $\|S(A)\|_{\gamma_F/2}\le\mathbf{S}$. The operator $\mathsf{D}(\cdot)$ is an entire function of the connection, so that, with $\partial_{A_1}$ as in Section~\ref{sec:fermionic-operators-stop},
\[
\delta\mathsf{D}:=\mathsf{D}(A+A_1)-\mathsf{D}(A)=\int_0^1\mathrm{d}u\,\left(\partial_{A_1}\mathsf{D}\right)(A+u\cdot A_1).
\]
As $\gamma_F\le 1$ (Section~\ref{sec:constants-3}), Lemma~\ref{lem:bounds-elementary-operators}(b) applies to the norm $\|\cdot\|_{\gamma_F/2}$ at each point $A+u\cdot A_1$; since $\epsilon_{A_1}=e_{*}\cdot|A_1|_{\infty}\le c_1$, it gives $\left\|\left(\partial_{A_1}\mathsf{D}\right)(A+u\cdot A_1)\right\|_{\gamma_F/2}\le 40\,c_1$. Majorants, and with them the norms of Section~\ref{sec:fermionic-operators-stop}, are subadditive under integration. Hence $\|\delta\mathsf{D}\|_{\gamma_F/2}\le 40\,c_1$, and~\eqref{eq:four-jets-Maj1} together with $c_1\le 1/(80\,\mathbf{S})$ (Section~\ref{sec:constants-3}) yields
\[
\left\|\delta\mathsf{D}\cdot S(A)\right\|_{\gamma_F/2}\le 40\,c_1\cdot\mathbf{S}\le\frac{1}{2}.
\]
By~\eqref{eq:four-jets-Maj1}, $\left\|\left(-\delta\mathsf{D}\cdot S(A)\right)^i\right\|_{\gamma_F/2}\le 2^{-i}$ for every $i\ge 0$. Therefore the Neumann series $\sum_{i\ge 0}\left(-\delta\mathsf{D}\cdot S(A)\right)^i$ converges in the norm $\|\cdot\|_{\gamma_F/2}$; as every matrix entry of an operator is bounded in modulus by this norm, the entries converge as well. Moreover
\[
\left\|S(A)\cdot\sum_{i\ge 0}\left(-\delta\mathsf{D}\cdot S(A)\right)^i\right\|_{\gamma_F/2}\le\mathbf{S}\cdot\sum_{i\ge 0}2^{-i}=2\,\mathbf{S} .
\]
Since $\mathsf{D}(A)\cdot S(A)=1$ and $\mathsf{D}(A+A_1)=\mathsf{D}(A)+\delta\mathsf{D}$, one has
\[
\mathsf{D}(A+A_1)\cdot S(A)\cdot\sum_{i\ge 0}\left(-\delta\mathsf{D}\cdot S(A)\right)^i=\left(1+\delta\mathsf{D}\cdot S(A)\right)\cdot\sum_{i\ge 0}\left(-\delta\mathsf{D}\cdot S(A)\right)^i=1 .
\]
The operators act on a space of finite dimension, where a right inverse is an inverse. So $\mathsf{D}(A+A_1)$ is invertible, $S(A+A_1)=S(A)\cdot\sum_{i\ge 0}\left(-\delta\mathsf{D}\cdot S(A)\right)^i$, and $\|S(A+A_1)\|_{\gamma_F/2}\le 2\,\mathbf{S}$.

The set of invertible matrices is open, and on it the entries of the inverse are analytic functions of the entries. Since $\mathsf{D}(\cdot)$, $Q_{*}(\cdot)$ and $Q_{*}^T(-\,\cdot)$ are entire functions of the connection, it follows that $S(\cdot)$ and $D_{*}(\cdot)$ are analytic in a neighbourhood of every point of $G_{\mathrm{w}}^{+}$. Fix $A'\in G_{\mathrm{w}}^{+}$ and write $Q:=Q_{*}(A')$, $\tilde{Q}:=Q_{*}^T(-A')$ and $S:=S(A')$. The derivative of $S(\cdot)$ at $A'$ in the direction $B$ is $-S\cdot(\partial_B\mathsf{D})\cdot S$; hence
\[
\partial_BD_{*}(A')=-b_{*}^2\cdot\left[(\partial_BQ)\cdot S\cdot\tilde{Q}-Q\cdot S\cdot(\partial_B\mathsf{D})\cdot S\cdot\tilde{Q}+Q\cdot S\cdot(\partial_B\tilde{Q})\right],
\]
all derivatives being taken at $A'$. One has $\mathfrak{i}(A')<\frac{1}{4}$ (Section~\ref{sec:dependence-stop-level}). By Lemma~\ref{lem:bounds-elementary-operators}(a), the majorants $|Q|(y,y')$ and $|\tilde{Q}|(y,y')$ are at most $\frac{3}{2}\cdot\mathbf{1}[y=y']$, and $|\partial_BQ|(y,y')$ and $|\partial_B\tilde{Q}|(y,y')$ are at most $\frac{9}{4}\,\epsilon_B\cdot\mathbf{1}[y=y'\in P_B]$. By Lemma~\ref{lem:bounds-elementary-operators}(b), $\|\partial_B\mathsf{D}\|_{\gamma_F/2}\le 40\,\epsilon_B$. By the first part of the proof, $\|S\|_{\gamma_F/2}\le 2\,\mathbf{S}$, and $b_{*}\le 1$ by Lemma~\ref{lem:scales-stop-level}(a). With~\eqref{eq:four-jets-Maj1} it follows that
\begin{align*}
\left\|\partial_BD_{*}(A')\right\|_{\gamma_F/2}&\le 2\cdot\frac{9}{4}\cdot\frac{3}{2}\cdot 2\,\mathbf{S}\cdot\epsilon_B+\frac{9}{4}\cdot\left(2\,\mathbf{S}\right)^2\cdot 40\,\epsilon_B\\
&=\left(13.5\,\mathbf{S}+360\,\mathbf{S}^2\right)\cdot\epsilon_B\le C_K'\cdot\epsilon_B;
\end{align*}
the last inequality holds because $C_K'=400\,\mathbf{S}^2$ and $\mathbf{S}\ge 1$ (Section~\ref{sec:constants-3}).

For the sum over $(y,y')$ the three terms are estimated separately. By~\eqref{eq:four-jets-Maj1} and Lemma~\ref{lem:bounds-elementary-operators}(a),
\[
\left|(\partial_BQ)\cdot S\cdot\tilde{Q}\right|(y,y')\le\frac{27}{8}\,\epsilon_B\cdot\mathbf{1}[y\in P_B]\cdot|S|(y,y'),
\]
and every row sum and every column sum of $|S|$ is at most $\|S\|_{\gamma_F/2}\le 2\,\mathbf{S}$; so the sum of this majorant over $(y,y')$ is at most $\frac{27}{8}\cdot 2\,\mathbf{S}\cdot\epsilon_B\cdot\#|P_B|$. In the same way the majorant of $Q\cdot S\cdot(\partial_B\tilde{Q})$ vanishes unless $y'\in P_B$, and its sum over $(y,y')$ is at most $\frac{27}{8}\cdot 2\,\mathbf{S}\cdot\epsilon_B\cdot\#|P_B|$. By Lemma~\ref{lem:bounds-elementary-operators}(b), $|\partial_B\mathsf{D}|(z,z')=0$ unless $z,z'\in P_B$, and $\sum_{z'}|\partial_B\mathsf{D}|(z,z')\le 40\,\epsilon_B$ for every $z$. With~\eqref{eq:four-jets-Maj1} and Lemma~\ref{lem:bounds-elementary-operators}(a) this gives
\begin{align*}
\sum_{y,y'}\left|Q\cdot S\cdot(\partial_B\mathsf{D})\cdot S\cdot\tilde{Q}\right|(y,y')&\le\frac{9}{4}\sum_{z,z'\in P_B}|\partial_B\mathsf{D}|(z,z')\cdot\left(\sum_y|S|(y,z)\right)\cdot\left(\sum_{y'}|S|(z',y')\right)\\
&\le\frac{9}{4}\cdot\left(2\,\mathbf{S}\right)^2\cdot 40\,\epsilon_B\cdot\#|P_B|.
\end{align*}
Adding the three bounds and using $b_{*}^2\le 1$, one obtains
\[
\sum_{y,y'}\left|\partial_BD_{*}(A')\right|(y,y')\le\left(13.5\,\mathbf{S}+360\,\mathbf{S}^2\right)\cdot\epsilon_B\cdot\#|P_B|\le C_K'\cdot\epsilon_B\cdot\#|P_B|.
\]
Finally, $D_{*}(\cdot)$ is analytic in a neighbourhood of the segment $\{A'+u\cdot B:0\le u\le 1\}$, which lies in $G_{\mathrm{w}}^{+}$, and
\[
D_{*}(A'+B)-D_{*}(A')=\int_0^1\mathrm{d}u\,\left(\partial_BD_{*}\right)(A'+u\cdot B).
\]
Both assertions follow from the two bounds for $\partial_BD_{*}$ at the points $A'+u\cdot B$, because majorants and norms are subadditive under integration.
\end{proof}

\begin{lemma}[Inverses of $D_{*}$]\label{lem:inverses}
Let $\mathcal{I}:=T^0_{*}\times\{1,2,3,4\}$, and for $x=(y,\boldsymbol{\beta})$ and $x'=(y',\boldsymbol{\beta}')$ in $\mathcal{I}$ let $d(x,x'):=d(y,y')$. Operators on the spinor functions on $T^0_{*}$ are identified with their matrices on $\mathcal{I}$.
\begin{enumerate}
\item[(a)] Let $A$ be a real bond function on $T_{*}$. Then $D_{*}(A)$ is invertible, its inverse $\mathcal{G}:=D_{*}(A)^{-1}$ satisfies $|\mathcal{G}|_{\ell^2}\le g^2$, and
\[
\max\left\{\sup_{x\in\mathcal{I}}\sum_{x'\in\mathcal{I}}|\mathcal{G}(x,x')|\cdot\exp\left(\vartheta\cdot d(x,x')\right),\ \sup_{x'\in\mathcal{I}}\sum_{x\in\mathcal{I}}|\mathcal{G}(x,x')|\cdot\exp\left(\vartheta\cdot d(x,x')\right)\right\}\le\boldsymbol{\alpha} .
\]
\item[(b)] Let $A=A_r+A_1$, where $A_r$ is a real bond function on $T_{*}$ and $A_1$ is a complex bond function on $T_{*}$ with $e_{*}\cdot|A_1|_{\infty}\le\mathfrak{m}_{*}/24$, and assume that $\mathsf{D}(A)$ is invertible. Then the operator $D_{*}(A)$, formed by the formula of Section~\ref{sec:fermionic-operators-stop}, is invertible,
\[
\left|D_{*}(A)^{-1}\right|_{\ell^2}\le b_{*}^{-1}+\frac{4}{\mathfrak{m}_{*}}\le g_1,
\]
and every row of the matrix $D_{*}(A)^{-1}$ has $\ell^1$ norm at most $4\,l_{*}^3\cdot g_1$.
\end{enumerate}
\end{lemma}

\begin{proof}
(a) By~\eqref{eq:four-jets-Gr7}, the operator $D_{*}(A)$ is invertible and $|\mathcal{G}|_{\ell^2}\le b_{*}^{-1}+\mathfrak{m}_{*}^{-1}$. By Lemma~\ref{lem:scales-stop-level}(a), $\mathfrak{m}_{*}>\mathfrak{m}_{-}$ and $b_{*}\ge b_{-}$; hence $|\mathcal{G}|_{\ell^2}\le b_{-}^{-1}+\mathfrak{m}_{-}^{-1}=g^2$ (Section~\ref{sec:constants-3}). By~\eqref{eq:four-jets-Gr7} again, the $4\times 4$ blocks of $\mathcal{G}$ satisfy, in operator norm,
\[
\left|\mathcal{G}(y,y')\right|\le b_{*}^{-1}\cdot\mathbf{1}[y=y']+\frac{2\exp(1)}{\mathfrak{m}_{*}}\cdot\exp\left(-\frac{\mathfrak{m}_{*}}{12\exp(1)}\cdot d(y,y')\right).
\]
Here $\frac{\mathfrak{m}_{*}}{12\exp(1)}\ge\frac{\mathfrak{m}_{-}}{12\exp(1)}=\boldsymbol{\theta}_{\mathrm{CT}}$, and $\vartheta\le\boldsymbol{\theta}_{\mathrm{CT}}/2$ (Section~\ref{sec:constants-3}), so that $\boldsymbol{\theta}_{\mathrm{CT}}-\vartheta\ge\boldsymbol{\theta}_{\mathrm{CT}}/2$. By the Cauchy--Schwarz inequality, the $\ell^1$ norm of a row or of a column of a $4\times 4$ matrix is at most twice its operator norm. Hence, for $x=(y,\boldsymbol{\beta})\in\mathcal{I}$,
\begin{align*}
\sum_{x'\in\mathcal{I}}|\mathcal{G}(x,x')|\cdot\exp\left(\vartheta\cdot d(x,x')\right)&\le 2\sum_{y'\in T^0_{*}}\left(b_{*}^{-1}\cdot\mathbf{1}[y=y']+\frac{2\exp(1)}{\mathfrak{m}_{*}}\cdot\exp\left(-\left(\boldsymbol{\theta}_{\mathrm{CT}}-\vartheta\right)\cdot d(y,y')\right)\right)\\
&\le 2\,b_{-}^{-1}+\frac{4\exp(1)}{\mathfrak{m}_{-}}\sum_{y'\in T^0_{*}}\exp\left(-\frac{\boldsymbol{\theta}_{\mathrm{CT}}}{2}\cdot d(y,y')\right)\\
&\le 2\,b_{-}^{-1}+\frac{4\exp(1)}{\mathfrak{m}_{-}}\cdot\mathrm{ssum}\left(\frac{\boldsymbol{\theta}_{\mathrm{CT}}}{2}\right)=\boldsymbol{\alpha} ,
\end{align*}
by the bound for $\mathrm{ssum}$ and the definition of $\boldsymbol{\alpha}$ in Section~\ref{sec:constants-3}. The sums over $x$ at fixed $x'$ are bounded in the same way.

(b) For $\nu\in\{0,1,2\}$ and $\boldsymbol{\sigma}\in\{+1,-1\}$ let $x^{\boldsymbol{\sigma}}=x+\boldsymbol{\sigma}\cdot\eta_{*}\cdot e_{\nu}$ and $P^{\boldsymbol{\sigma}}_{\nu}=(1-\boldsymbol{\sigma}\cdot\gamma_{\nu})/2$ be as in Section~\ref{sec:fermionic-operators-stop}, and define an operator $T^{\boldsymbol{\sigma}}_{\nu}(A_r)$ and a function $m^{\boldsymbol{\sigma}}_{\nu}$ on $T_{*}$ by
\[
\left(T^{\boldsymbol{\sigma}}_{\nu}(A_r)f\right)(x):=u_{A_r}(x,x^{\boldsymbol{\sigma}})\cdot f(x^{\boldsymbol{\sigma}}),\qquad m^{\boldsymbol{\sigma}}_{\nu}(x):=u_{A_1}(x,x^{\boldsymbol{\sigma}}).
\]
As $A_r$ is real, $|u_{A_r}|=1$, and $T^{\boldsymbol{\sigma}}_{\nu}(A_r)$ is unitary on $\ell^2(T_{*})$. The matrices $\gamma_{\nu}$ are Hermitian with $\gamma_{\nu}^2=1$ (Section~\ref{sec:dirac-matrices-lattice}); hence the $P^{\boldsymbol{\sigma}}_{\nu}$ are orthogonal projections. Since $u_A=u_{A_r}\cdot u_{A_1}$, the formula for $D_A$ in Section~\ref{sec:fermionic-operators-stop} gives
\[
D_A-D_{A_r}=-\eta_{*}^{-1}\cdot\sum_{\nu,\boldsymbol{\sigma}}P^{\boldsymbol{\sigma}}_{\nu}\cdot\left(m^{\boldsymbol{\sigma}}_{\nu}-1\right)\cdot T^{\boldsymbol{\sigma}}_{\nu}(A_r),
\]
where $m^{\boldsymbol{\sigma}}_{\nu}-1$ acts by multiplication. One has $\eta_{*}=L^{-k^{*}}\le 1$, and $\mathfrak{m}_{*}\le 1$ by Lemma~\ref{lem:scales-stop-level}(a); so $\eta_{*}\cdot e_{*}\cdot|A_1|_{\infty}\le\frac{1}{24}$, and $|\exp(z)-1|\le|z|\exp(|z|)$ gives
\[
\left|m^{\boldsymbol{\sigma}}_{\nu}(x)-1\right|\le\eta_{*}\cdot e_{*}\cdot|A_1|_{\infty}\cdot\exp\left(\eta_{*}\cdot e_{*}\cdot|A_1|_{\infty}\right)\le 2\,\eta_{*}\cdot e_{*}\cdot|A_1|_{\infty}.
\]
The multiplication by $m^{\boldsymbol{\sigma}}_{\nu}-1$ therefore has operator norm at most $2\,\eta_{*}\cdot e_{*}\cdot|A_1|_{\infty}$ on $\ell^2(T_{*})$; as $P^{\boldsymbol{\sigma}}_{\nu}$ is an orthogonal projection and $T^{\boldsymbol{\sigma}}_{\nu}(A_r)$ is unitary, the same bound holds for each of the six summands $P^{\boldsymbol{\sigma}}_{\nu}\cdot\left(m^{\boldsymbol{\sigma}}_{\nu}-1\right)\cdot T^{\boldsymbol{\sigma}}_{\nu}(A_r)$. Together with the factor $\eta_{*}^{-1}$ in front of the sum, and with $e_{*}\cdot|A_1|_{\infty}\le\mathfrak{m}_{*}/24$, this gives
\[
\left|D_A-D_{A_r}\right|_{\ell^2}\le\eta_{*}^{-1}\cdot 6\cdot 2\,\eta_{*}\cdot e_{*}\cdot|A_1|_{\infty}=12\,e_{*}\cdot|A_1|_{\infty}\le\frac{\mathfrak{m}_{*}}{2}.
\]
By~\eqref{eq:four-jets-Gr7}, $\operatorname{Re}\left(D_{A_r}+\mathfrak{m}_{*}\right)\ge\mathfrak{m}_{*}$ on $\ell^2(T_{*})$. Consequently $\operatorname{Re}\left(D_A+\mathfrak{m}_{*}\right)\ge\mathfrak{m}_{*}/2$, and the operator $\mathbf{R}:=D_A+\mathfrak{m}_{*}$ is invertible with $|\mathbf{R}^{-1}|_{\ell^2}\le 2/\mathfrak{m}_{*}$.

Put $\check{c}:=Q_{*}(A)\cdot\mathbf{R}^{-1}\cdot Q_{*}^T(-A)$. By the definitions of Section~\ref{sec:fermionic-operators-stop},
\[
\mathsf{D}(A)=\mathbf{R}+b_{*}\cdot Q_{*}^T(-A)\cdot Q_{*}(A)=\mathbf{R}\cdot\left(1+b_{*}\cdot\mathbf{R}^{-1}\cdot Q_{*}^T(-A)\cdot Q_{*}(A)\right).
\]
By Sylvester's determinant identity $\det(1+XY)=\det(1+YX)$,
\[
\det\left(1+b_{*}\cdot\check{c}\right)=\det\left(1+b_{*}\cdot\mathbf{R}^{-1}\cdot Q_{*}^T(-A)\cdot Q_{*}(A)\right)=\det(\mathbf{R})^{-1}\cdot\det\mathsf{D}(A)\ne 0 .
\]
Hence $b_{*}^{-1}+\check{c}$ is invertible, and the Woodbury identity gives
\[
S(A)=\mathbf{R}^{-1}-\mathbf{R}^{-1}\cdot Q_{*}^T(-A)\cdot\left(b_{*}^{-1}+\check{c}\right)^{-1}\cdot Q_{*}(A)\cdot\mathbf{R}^{-1}.
\]
It follows that
\[
Q_{*}(A)\cdot S(A)\cdot Q_{*}^T(-A)=\check{c}-\check{c}\cdot\left(b_{*}^{-1}+\check{c}\right)^{-1}\cdot\check{c}=\check{c}\cdot\left(b_{*}^{-1}+\check{c}\right)^{-1}\cdot b_{*}^{-1},
\]
and therefore
\[
D_{*}(A)=b_{*}-b_{*}^2\cdot Q_{*}(A)\cdot S(A)\cdot Q_{*}^T(-A)=b_{*}\cdot\left[\left(b_{*}^{-1}+\check{c}\right)-\check{c}\right]\cdot\left(b_{*}^{-1}+\check{c}\right)^{-1}=\left(b_{*}^{-1}+\check{c}\right)^{-1}.
\]
Thus $D_{*}(A)$ is invertible, and $D_{*}(A)^{-1}=b_{*}^{-1}+\check{c}$.

Let $\ell^2(T_{*})$ be normed with the weight $\eta_{*}^3$ at every site; operators from $\ell^2(T_{*})$ to itself have the same norm for every constant weight. Since $\eta_{*}^3\cdot\#|\Delta_y|=1$, the Cauchy--Schwarz inequality in each block $\Delta_y$ shows that $Q_{*}(A)\colon\ell^2(T_{*})\to\ell^2(T^0_{*})$ and $Q_{*}^T(-A)\colon\ell^2(T^0_{*})\to\ell^2(T_{*})$ have norms at most $\sup\left|\exp\left(\pm i\sigma_A\right)\right|$. The function $\sigma_A$ is linear in $A$ and $\sigma_{A_r}$ is real; so $\left|\exp\left(\pm i\sigma_A\right)\right|\le\exp\left(|\sigma_{A_1}|\right)$, and, by the bound for $\sigma$ in Section~\ref{sec:fermionic-operators-stop},
\[
|\sigma_{A_1}(y,x)|\le\frac{3}{2}\,e_{*}\cdot|A_1|_{\infty}\le\frac{\mathfrak{m}_{*}}{16}\le\frac{1}{16}.
\]
Both norms are therefore at most $\exp\left(\frac{1}{16}\right)\le 2^{\frac{1}{2}}$, and
\[
\left|D_{*}(A)^{-1}\right|_{\ell^2}\le b_{*}^{-1}+2^{\frac{1}{2}}\cdot\frac{2}{\mathfrak{m}_{*}}\cdot 2^{\frac{1}{2}}=b_{*}^{-1}+\frac{4}{\mathfrak{m}_{*}}\le b_{-}^{-1}+\frac{4}{\mathfrak{m}_{-}}=g_1.
\]
Every entry of a matrix is bounded in modulus by its operator norm, and a row of a matrix on $\mathcal{I}$ has $\#|\mathcal{I}|=4\,l_{*}^3$ entries; this gives the bound for the rows.
\end{proof}

\subsection{The fermion integral at the stop level as a function of the variables $t_c$}\label{sec:fermion-integral-stop}

\subsubsection{The activities and the Gaussian form of the integral}\label{sec:activities-gaussian-form}

Throughout this subsection condition~\eqref{eq:four-jets-sm2} is assumed. Fix $\theta\in\mathbb{R}^3$ and $A\in V_{\mathrm{t}}(T^0_{*})$ with $\chi_{*}(A)=1$, and let $A^g$ and $\mathscr{A}^0$ be as in~\eqref{it:trajectory-fermionic-operators-SL5}. Put
\[
\mathscr{A}:=\mathscr{A}^0+c_{\theta},\qquad\mathscr{T}(t):=\mathscr{A}+\mathscr{C}(t),\qquad\mathcal{G}:=D_{*}(\mathscr{A})^{-1},\qquad q_x:=\frac{2}{t_x},\qquad w^{\mathrm{E}}:=e_{*}^{\frac{1}{4}-7\varepsilon}.
\]
By~\eqref{eq:four-jets-4} with $u=1$, one has $\mathscr{A}=H_{*}(A^g+c_{\theta})$; this connection is real and lies in $\frac{1}{4}\tilde{R}_{*}$. The inverse $\mathcal{G}$ exists by Lemma~\ref{lem:inverses}(a).

Let $t\in\mathbb{D}^{\circ}(t_x^{+})$. The placement is admitted by~\eqref{eq:four-jets-sm2}; so clause~\eqref{it:trajectory-fermionic-operators-SL7} of Proposition~\ref{prop:trajectory-fermionic-operators} gives $\mathscr{C}(t)\in\mathfrak{X}_{*}$ and $e_{*}\cdot|\mathscr{C}(t)|_{\infty}<e_{*}^{3\varepsilon}$. The set $\mathfrak{X}_{*}$ is convex and contains $0$. By~\eqref{eq:four-jets-4}, the real connections $c_{\theta}+u\cdot\mathscr{A}^0=H_{*}(u\cdot A^g+c_{\theta})$, $0\le u\le 1$, lie in $\frac{1}{4}\tilde{R}_{*}$, which is contained in $\tilde{R}_{*}$. Hence the connections
\[
\mathscr{T}(t),\qquad c_{\theta}+\mathscr{C}(t),\qquad\mathscr{A}+u\cdot\mathscr{C}(t),\qquad c_{\theta}+u\cdot\mathscr{A}^0+\mathscr{C}(t)\qquad(0\le u\le 1)
\]
belong to $G_{\mathrm{w}}$. Let
\begin{multline}\label{eq:four-jets-15}
\mathsf{Z}(t):=\int\mathrm{D}\Psi\,\exp\left\{-\left\langle\bar{\Psi},D_{*}(\mathscr{T}(t))\Psi\right\rangle-m_{*}\cdot\left\langle\bar{\psi},\psi\right\rangle\right.\\
\left.+\sum_{X\in D_{*}}\mathscr{E}_{*}\left(X;\mathscr{A},\mathscr{C}(t),\psi^{\#}\right)\right\},
\end{multline}
where $\psi$, $\bar{\psi}$, $\chi$, $\bar{\chi}$ and $\psi^{\#}=(\psi,\chi)$ are formed with $\mathscr{T}=\mathscr{T}(t)$ as in~\eqref{it:trajectory-fermionic-operators-SL6}. Since $A^g+c_{\theta}\in\mathrm{Dom}_{*}$ and $H_{*}(A^g+c_{\theta})=\mathscr{A}$, by~\eqref{eq:four-jets-4}, the identity~\eqref{eq:four-jets-5} gives
\[
Y^x_{*}\left(A^g+c_{\theta};\mathscr{C}(t)\right)=\det D_{*}(0)^{-1}\cdot\exp\left\{-\boldsymbol{\varepsilon}^t_{*}\cdot l_{*}^3+\sum_{X\in D_{*}}\mathfrak{P}_{*}\left(X,\mathscr{C}(t)\right)\right\}\cdot\mathsf{Z}(t).
\]

\paragraph{The mass term in the activities.}
The pairing $\langle\bar{\psi},\psi\rangle$ is the sum of the pairings $\langle\bar{\psi},\psi\rangle_{\square}$ over the $M$-cubes $\square$ of $T_{*}$. For $X\in D_{*}$ and $t\in\mathbb{D}^{\circ}(t_x^{+})$ define $\mathbf{E}(X;t)\in\mathrm{Gr}[\psi,\chi]$ by
\[
\mathbf{E}(X;t):=\begin{cases}\mathscr{E}_{*}\left(X;\mathscr{A},\mathscr{C}(t),\cdot\right)-m_{*}\cdot\langle\bar{\psi},\psi\rangle_X&\text{if $X$ is an $M$-cube},\\[2pt] \mathscr{E}_{*}\left(X;\mathscr{A},\mathscr{C}(t),\cdot\right)&\text{otherwise},\end{cases}
\]
and let $\mathbf{E}_{nm}(X;t;\cdot)$ be its kernels, antisymmetric in the arguments $\xi$ and, separately, in the arguments $\boldsymbol{\zeta}$ (Section~\ref{sec:grassmann-algebras-norms}). The subtracted term $m_{*}\cdot\langle\bar{\psi},\psi\rangle_X$ is of bidegree $(n,m)=(2,0)$, its kernel vanishes unless both arguments lie in $X$, and it does not depend on $t$. An $M$-cube has the volume $M^3$; so, by~\eqref{eq:polymer-activities-3}, the antisymmetric kernel $K$ of this term satisfies $\frac{a^2}{2!}\,\|K\|=4\,M^3\cdot|m_{*}|\cdot a^2$ for every $a>0$.

\begin{lemma}[Bounds for the activities]\label{lem:bounds-activities}
Assume~\eqref{eq:four-jets-sm1}, let $0<a\le e_{*}^{-\frac{1}{4}+\varepsilon}$ and $X\in D_{*}$, and write $\partial_c:=\partial/\partial t_c$ for $c\in\mathrm{Cel}$.
\begin{enumerate}
\item[(a)] For $t\in\mathbb{D}(t_x)$,
\[
\sum_{n,m}\frac{a^{n+m}}{n!\,m!}\,\left\|\mathbf{E}_{nm}(X;t)\right\|\le w^{\mathrm{E}}\cdot\exp\left(-\kappa\cdot d_M(X)\right).
\]
\item[(b)] For $t\in\mathbb{D}(t_x/2)$ and $c\in\mathrm{Cel}$,
\[
\sum_{n,m}\frac{a^{n+m}}{n!\,m!}\,\left\|\partial_c\mathbf{E}_{nm}(X;t)\right\|\le q_x\cdot w^{\mathrm{E}}\cdot\exp\left(-\kappa\cdot d_M(X)\right),
\]
and, if moreover $c'\in\mathrm{Cel}$ and $c'\ne c$,
\[
\sum_{n,m}\frac{a^{n+m}}{n!\,m!}\,\left\|\partial_c\partial_{c'}\mathbf{E}_{nm}(X;t)\right\|\le q_x^2\cdot w^{\mathrm{E}}\cdot\exp\left(-\kappa\cdot d_M(X)\right).
\]
\item[(c)] Let $c\in\mathrm{Cel}$ and $t\in\mathbb{D}^{\circ}(t_x^{+})$. If $X$ contains no point of $P_c$, then $\partial_c\mathbf{E}_{nm}(X;t)=0$ for all $(n,m)$.
\end{enumerate}
\end{lemma}

\begin{proof}
(a) Let $t\in\mathbb{D}(t_x)$. As $t_x<t_x^{+}$, the point $t$ lies in $\mathbb{D}^{\circ}(t_x^{+})$, and $(\mathscr{A},\mathscr{C}(t))\in\tilde{R}_{*}\times\mathfrak{X}_{*}$. One has $e_{*}<1$ by Lemma~\ref{lem:scales-stop-level}(a) and~\eqref{eq:renormalisation-group-A8}; hence $a\le e_{*}^{-\frac{1}{4}+\varepsilon}\le e_{*}^{-\frac{1}{4}}$ and
\[
a^{n+m}\le e_{*}^{-\frac{n}{4}}\cdot e_{*}^{\left(-\frac{1}{4}+\varepsilon\right)m}\qquad\text{for all }n,m.
\]
By the bound~\eqref{eq:four-jets-2} of~\eqref{it:trajectory-fermionic-operators-SL3}, the kernels $\mathscr{E}_{nm}(X;\mathscr{A},\mathscr{C}(t))$ contribute to the left side at most
\[
\sum_{n,m}\frac{e_{*}^{-\frac{n}{4}}\cdot e_{*}^{\left(-\frac{1}{4}+\varepsilon\right)m}}{n!\,m!}\,\left\|\mathscr{E}_{nm}(X;\mathscr{A},\mathscr{C}(t))\right\|\le 4\,e_{*}^{\varepsilon}\cdot w^{\mathrm{E}}\cdot\exp\left(-\kappa\cdot d_M(X)\right).
\]
If $X$ is an $M$-cube, then $d_M(X)=0$. In this case~\eqref{it:trajectory-fermionic-operators-SL1} gives $|m_{*}|\le e_{*}^{\frac{3}{4}-8\varepsilon}/8$, while $a^2\le e_{*}^{-\frac{1}{2}+2\varepsilon}$; so the mass term contributes at most
\[
4\,M^3\cdot\frac{e_{*}^{\frac{3}{4}-8\varepsilon}}{8}\cdot e_{*}^{-\frac{1}{2}+2\varepsilon}=\frac{M^3}{2}\cdot e_{*}^{\varepsilon}\cdot w^{\mathrm{E}} .
\]
By the triangle inequality for the kernel norms, the left side of (a) is at most $\left(4+\frac{M^3}{2}\right)\cdot e_{*}^{\varepsilon}\cdot w^{\mathrm{E}}\cdot\exp\left(-\kappa\cdot d_M(X)\right)$, and the second condition of~\eqref{eq:four-jets-sm1}, $\left(4+\frac{M^3}{2}\right)\cdot e_{*}^{\varepsilon}\le 1$, gives (a).

(b) By~\eqref{it:trajectory-fermionic-operators-SL3}, the kernels $\mathscr{E}_{nm}(X;\mathscr{A},\mathscr{C})$ are analytic functions of $\mathscr{C}\in\mathfrak{X}_{*}$. The map $t\mapsto\mathscr{C}(t)$ is linear and, by~\eqref{it:trajectory-fermionic-operators-SL7}, sends $\mathbb{D}^{\circ}(t_x^{+})$ into $\mathfrak{X}_{*}$, and the mass term does not depend on $t$. Hence, at every family of arguments, $\mathbf{E}_{nm}(X;t)$ is an analytic function of $t\in\mathbb{D}^{\circ}(t_x^{+})$. Let $t\in\mathbb{D}(t_x/2)$ and $c\in\mathrm{Cel}$. For $|\mathbf{s}|\le t_x/2$ the point $t+\mathbf{s}\cdot e_c$ lies in $\mathbb{D}(t_x)$, and Cauchy's formula gives, at every family of arguments,
\[
\partial_c\mathbf{E}_{nm}(X;t)=\frac{1}{2\pi i}\oint_{|\mathbf{s}|=t_x/2}\mathrm{d}\mathbf{s}\,\frac{\mathbf{E}_{nm}(X;t+\mathbf{s}\cdot e_c)}{\mathbf{s}^2}.
\]
The kernel norm of Section~\ref{sec:grassmann-algebras-norms} is an integral of the modulus over the arguments; taking it under the contour integral, one finds that $\left\|\partial_c\mathbf{E}_{nm}(X;t)\right\|$ is at most $\frac{2}{t_x}=q_x$ times the mean of $\left\|\mathbf{E}_{nm}(X;t+\mathbf{s}\cdot e_c)\right\|$ over the circle $|\mathbf{s}|=t_x/2$. Multiplying by $\frac{a^{n+m}}{n!\,m!}$, summing over $(n,m)$ under the mean, and applying (a) at the points $t+\mathbf{s}\cdot e_c\in\mathbb{D}(t_x)$, one obtains the first bound. Let now $c'\in\mathrm{Cel}$, $c'\ne c$. For $|\mathbf{s}|=|\mathbf{s}'|=t_x/2$ the point $t+\mathbf{s}\cdot e_c+\mathbf{s}'\cdot e_{c'}$ lies in $\mathbb{D}(t_x)$, because $c\ne c'$, so that each of its components differs from the corresponding component of $t$ by at most $t_x/2$; and
\[
\partial_c\partial_{c'}\mathbf{E}_{nm}(X;t)=\frac{1}{(2\pi i)^2}\oint_{|\mathbf{s}|=t_x/2}\mathrm{d}\mathbf{s}\oint_{|\mathbf{s}'|=t_x/2}\mathrm{d}\mathbf{s}'\,\frac{\mathbf{E}_{nm}(X;t+\mathbf{s}\cdot e_c+\mathbf{s}'\cdot e_{c'})}{\mathbf{s}^2\cdot\mathbf{s}'^2}.
\]
The same argument, with the mean over the two circles, gives the second bound, with the factor $q_x^2$.

(c) The mass term does not depend on $t$, and $\mathscr{C}(t)=\sum_{c''}t_{c''}\cdot\hat{c}_{c''}$; so $\partial_c\mathbf{E}_{nm}(X;t)$ is the derivative of $\mathscr{C}\mapsto\mathscr{E}_{nm}(X;\mathscr{A},\mathscr{C})$ at $\mathscr{C}(t)$ in the direction $\hat{c}_c$. By~\eqref{it:trajectory-fermionic-operators-SL3}, that is, by clause~\eqref{it:renormalisation-group-1-k2} of Definition~\ref{def:renormalisation-group-1}, this kernel depends on $\mathscr{C}$ only through the values $\mathscr{C}(b)$ at the bonds $b$ that have an end point in $X$. If $\hat{c}_c$ vanishes at all these bonds, the derivative is $0$. Otherwise there is a bond $b$ with $\hat{c}_c(b)\ne 0$ and with an end point $w\in X$. Then $b\in\mathrm{Bd}_c$, so the unit site $[w]$, the site $y$ with $w\in\Delta_y$, belongs to $P_c$. Moreover $[w]\in X$: by fact~\eqref{it:renormalisation-group-g0} of Section~\ref{sec:cubes-parameter-cubes}, applied to the level $k^{*}$ and the side $l_{*}$, which lies in $\mathrm{Cl}$ by Lemma~\ref{lem:scales-stop-level}(c), the unit cube $\Delta_{[w]}$ is contained in the $M$-cube that contains $w$; and this $M$-cube is contained in $X$. So $X$ contains a point of $P_c$.
\end{proof}

For $t\in\mathbb{D}^{\circ}(t_x^{+})$ let
\begin{align}
\check{K}(t)&:=D_{*}(\mathscr{T}(t))-D_{*}(\mathscr{A}),\notag\\
V(t)&:=-\left\langle\bar{\Psi},\check{K}(t)\Psi\right\rangle+\sum_{X\in D_{*}}\mathbf{E}(X;t)\Big|_{\psi^{\#}=H_{*}^{\#}(\mathscr{T}(t))\Psi}=:V_0(t)+V_{+}(t),\label{eq:four-jets-16}
\end{align}
where the vertical bar denotes the substitution $\psi^{\#}=H_{*}^{\#}(\mathscr{T}(t))\Psi$ of Section~\ref{sec:fermion-integrals-blocking}, that is, the insertion of the fields of~\eqref{it:trajectory-fermionic-operators-SL6} formed with $\mathscr{T}(t)$; here $V_0(t)$ is the part of degree zero of $V(t)$,
\[
V_0(t)=\sum_{X\in D_{*}}\mathscr{E}_{00}\left(X;\mathscr{A},\mathscr{C}(t)\right),
\]
and $V_{+}(t):=V(t)-V_0(t)$.

The Grassmann algebra $\mathrm{Gr}[\Psi]$, with the generators $\Psi(u)$, $u=(y,\boldsymbol{\beta},\boldsymbol{\omega})\in T^0_{*}\times\{1,2,3,4\}\times\{0,1\}$, is the algebra of Section~\ref{sec:gaussian-integrals-grassmann} for the finite set $\mathcal{I}=T^0_{*}\times\{1,2,3,4\}$ of Lemma~\ref{lem:inverses}: at $x=(y,\boldsymbol{\beta})$ its generators are $\Psi(x)=\Psi_{\boldsymbol{\beta}}(y)=\Psi(y,\boldsymbol{\beta},0)$ and $\bar{\Psi}(x)=\bar{\Psi}_{\boldsymbol{\beta}}(y)=\Psi(y,\boldsymbol{\beta},1)$. The set $\mathcal{I}$ carries the pseudo-metric $d((y,\boldsymbol{\beta}),(y',\boldsymbol{\beta}')):=d(y,y')$, and $\mathrm{Idx}:=\mathcal{I}\times\{+,-\}$ as in Section~\ref{sec:weighted-tree-expansion}. The element $V_{+}(t)$ of this algebra is even and has no constant term. When Section~\ref{sec:weighted-tree-expansion} is applied to it, the elements of $\mathrm{Idx}$ are the indices $u$, by the identifications
\[
(y,\boldsymbol{\beta},0)=((y,\boldsymbol{\beta}),-),\qquad(y,\boldsymbol{\beta},1)=((y,\boldsymbol{\beta}),+);
\]
the letters $\psi$ and $\xi$ keep the meaning that they have in Section~\ref{sec:grassmann-algebras-norms}. By~\eqref{it:trajectory-fermionic-operators-SL2} and~\eqref{it:trajectory-fermionic-operators-SL3}, all coefficients of $V_0(t)$ and of $V_{+}(t)$ are analytic functions of $t\in\mathbb{D}^{\circ}(t_x^{+})$.

Since $\langle\bar{\psi},\psi\rangle$ is the sum of the $\langle\bar{\psi},\psi\rangle_{\square}$ over the $M$-cubes $\square$, and every $M$-cube is a polymer, the exponent in~\eqref{eq:four-jets-15} equals $-\left\langle\bar{\Psi},D_{*}(\mathscr{A})\Psi\right\rangle+V(t)$; and $V_0(t)$ is a number. Therefore~\eqref{eq:four-jets-Gr3}, applied with the covariance $\mathcal{G}=D_{*}(\mathscr{A})^{-1}$, gives
\begin{equation}\label{eq:four-jets-17}
\mathsf{Z}(t)=\det D_{*}(\mathscr{A})\cdot\exp\left(V_0(t)\right)\cdot\int\mathrm{d}\mu_{\mathcal{G}}\,\exp\left(V_{+}(t)\right).
\end{equation}
The weights $w_{\boldsymbol{\delta}}$ and $w_P$ of Section~\ref{sec:weighted-tree-expansion} are formed with the constant $\vartheta$ of Section~\ref{sec:constants-3}. For a non-empty subset $P$ of $T^0_{*}$, the weight $w_P$ is the one formed with the non-empty subset $P\times\{1,2,3,4\}$ of $\mathcal{I}$.

\subsubsection{Weighted norms of the interaction}\label{sec:weighted-norms-interaction}

\begin{lemma}[Weighted norms of $V_{+}$]\label{lem:weighted-norms}
Assume \eqref{eq:four-jets-sm1}, \eqref{eq:four-jets-sm2} and \eqref{eq:four-jets-sm3}. Let $\theta$, $A$ and $A^g$ be as in Section~\ref{sec:activities-gaussian-form}, let $V_0(t)$ and $V_{+}(t)$ be defined by \eqref{eq:four-jets-16}, let $\rho_h$ and $N_h$ be the norms of Section~\ref{sec:weighted-tree-expansion}, formed with the weights $w_{\boldsymbol{\delta}}$, $w_P$ described in Section~\ref{sec:activities-gaussian-form}, and let $\partial_c$ denote the partial derivative with respect to $t_c$. For every $t\in\mathbb{D}(t_x)$ one has
\begin{equation}\tag{V1}\label{eq:four-jets-V1}
\rho_h\left(V_{+}(t);w_{\boldsymbol{\delta}}\right)\le\bar{r} .
\end{equation}
If $\hat{\mathfrak{a}}\ne0$, then the sets $\mathrm{Bd}_c$, $P_c$ and $P_c^{+}$ are non-empty for every $c\in\mathrm{Cel}$, and for every $t\in\mathbb{D}(t_x/2)$ and every $c\in\mathrm{Cel}$ one has
\begin{equation}\tag{V2}\label{eq:four-jets-V2}
N_h\left(\partial_cV_{+}(t);w_{P_c}\right)\le q_x\cdot C_V\cdot h_G^2\cdot\mathsf{p},\qquad\rho_h\left(\partial_cV_{+}(t);w_{P_c}\right)\le q_x\cdot C_V\cdot h_G^2.
\end{equation}
If $\hat{\mathfrak{a}}\ne0$, then for every $t\in\mathbb{D}(t_x/2)$ and all $c,c'\in\mathrm{Cel}$ with $c\ne c'$ one has
\begin{equation}\tag{V3}\label{eq:four-jets-V3}
N_h\left(\partial_c\partial_{c'}V_{+}(t);1\right)+\left|\partial_c\partial_{c'}V_0(t)\right|\le q_x^2\cdot C_V\cdot h_G^2\cdot\mathsf{p}\cdot\exp\left(-\vartheta\cdot\left(d_{\mathrm{sep}}-4M-6\right)\right).
\end{equation}
These bounds hold uniformly in $N$, $\theta$, $A$ and in the choice of $A^g$. If $\hat{\mathfrak{a}}=0$, then $\mathscr{C}(t)=0$ for every $t\in\mathbb{C}^{\mathrm{Cel}}$; in this case $V_0(t)$ and $V_{+}(t)$ do not depend on $t\in\mathbb{D}^{\circ}(t_x^{+})$, and all their partial derivatives with respect to the variables $t_c$ vanish.
\end{lemma}

\begin{proof}
If $\hat{\mathfrak{a}}=0$, then $\hat{\mathfrak{a}}_c=(\mathsf{g}_c)_{*}\hat{\mathfrak{a}}=0$ and $\hat{c}_c=\frac{1}{e}\cdot\hat{\mathfrak{a}}_c^{[k^{*}]}=0$ for every $c\in\mathrm{Cel}$ (Section~\ref{sec:probes-mirror-copies}), so that $\mathscr{C}(t)=\sum_{c\in\mathrm{Cel}}t_c\cdot\hat{c}_c=0$ for every $t\in\mathbb{C}^{\mathrm{Cel}}$. In this case, for every $t\in\mathbb{D}^{\circ}(t_x^{+})$, one has $\mathscr{T}(t)=\mathscr{A}=\mathscr{T}(0)$ and $\check{K}(t)=0=\check{K}(0)$, and $\mathbf{E}(X;t)=\mathbf{E}(X;0)$ for every $X\in D_{*}$, since $\mathbf{E}(X;t)$ depends on $t$ only through $\mathscr{C}(t)$ (Section~\ref{sec:activities-gaussian-form}). As the substitution in \eqref{eq:four-jets-16} is formed with $\mathscr{T}(t)=\mathscr{T}(0)$, it follows that $V(t)=V(0)$; hence $V_0(t)=V_0(0)$ and $V_{+}(t)=V_{+}(0)$, and all partial derivatives of $V_0$ and $V_{+}$ with respect to the variables $t_c$ vanish. This is the last assertion of the lemma.

If $\hat{\mathfrak{a}}\ne0$, then the sets $\mathrm{Bd}_c$, $P_c$ and $P_c^{+}$ are non-empty for every $c\in\mathrm{Cel}$, the weights $w_{P_c}$ and the distances $d(P_c,P_{c'})$ are defined, and $d_{\mathrm{sep}}=\min_{c\ne c'}d(P_c,P_{c'})$. To prove this implication, let $\hat{\mathfrak{a}}\ne0$. For every bond $b$ of $T_{*}$ one has $\hat{\mathfrak{a}}^{[k^{*}]}(b)=L^{\frac{k^{*}}{2}}\cdot\hat{\mathfrak{a}}\left(L^{k^{*}}\cdot b\right)$, and $b\mapsto L^{k^{*}}\cdot b$ maps the bonds of $T_{*}$ onto those of $\Lambda_{n_{*}}$ (Section~\ref{sec:lattice-measures-laws}); hence $\hat{\mathfrak{a}}^{[k^{*}]}\ne0$. The placement is admitted by \eqref{eq:four-jets-sm2}; so Lemma~\ref{lem:properties-probes}(d), applied to the norm $|\cdot|_{\infty}$, which is one of the three norms of Section~\ref{sec:domains-bond-functions}, gives
\[
\left|\hat{c}_c\right|_{\infty}=e^{-1}\cdot\left|\hat{\mathfrak{a}}_c^{[k^{*}]}\right|_{\infty}=e^{-1}\cdot\left|\hat{\mathfrak{a}}^{[k^{*}]}\right|_{\infty}>0\qquad\left(c\in\mathrm{Cel}\right).
\]
Hence every set $\mathrm{Bd}_c$ is non-empty. If $w$ is an end point of a bond of $\mathrm{Bd}_c$, then the unit site $[w]$, whose cube $\Delta_{[w]}$ contains $w$ (Section~\ref{sec:fermionic-operators-stop}), belongs to $P_c$; and $P_c\subset P_c^{+}$. Thus the sets $P_c$ and $P_c^{+}$ are non-empty for every $c\in\mathrm{Cel}$, the weights $w_{P_c}$ and the distances $d(P_c,P_{c'})$ are defined, and $d_{\mathrm{sep}}=\min_{c\ne c'}d(P_c,P_{c'})$ (Section~\ref{sec:probes-mirror-copies}). This proves the implication.

From here to the end of Step~1 no assumption on $\hat{\mathfrak{a}}$ is made beyond \eqref{eq:four-jets-sm1}, \eqref{eq:four-jets-sm2} and \eqref{eq:four-jets-sm3}; in particular, the case $\hat{\mathfrak{a}}=0$ is included. Put $\check{s}:=\mathrm{ssum}\left(\frac{\gamma_F}{2}\right)$; since $\gamma_F\le1$, one has $\check{s}\ge1$. By Section~\ref{sec:constants-3}, $\gamma_F\cdot M\ge16$ and $\gamma_F\le1$, so that $M\ge16$; moreover $\vartheta\le\min\left\{\frac{\gamma_F}{8},\frac{1}{M}\right\}$, $h_L\ge1$, $h_G\ge1$ and $C_H=100\,M\cdot C_F^2\ge100$. Stage~\ref{it:order-choice-3} gives $\kappa\ge7\kappa_0+7$, hence $\kappa-1\ge\kappa_0$. By \eqref{eq:four-jets-sm2}, $d_{\mathrm{sep}}\ge6M+8\ge4$; so Lemma~\ref{lem:properties-probes}(f) applies and gives $t_x\cdot\epsilon_c<e_{*}^{3\varepsilon}$, whence
\[
\epsilon_c\le\frac{e_{*}^{3\varepsilon}}{t_x}\le q_x\qquad\left(c\in\mathrm{Cel}\right),
\]
and, if $\hat{\mathfrak{a}}\ne0$, then $d(P_c,P_{c'})\ge d_{\mathrm{sep}}\ge M+2$ for $c\ne c'$. By \eqref{eq:four-jets-sm3}, $w^{\mathrm{E}}\le\frac{1}{8}$.

By Section~\ref{sec:activities-gaussian-form}, together with \eqref{it:trajectory-fermionic-operators-SL7}, every $t\in\mathbb{D}^{\circ}(t_x^{+})$ satisfies $\mathscr{C}(t)\in\mathfrak{X}_{*}$ and $e_{*}\cdot|\mathscr{C}(t)|_{\infty}<e_{*}^{3\varepsilon}$, and the fields $\mathscr{T}(t)$ and $\mathscr{A}+u\cdot\mathscr{C}(t)$, $0\le u\le1$, lie in $G_{\mathrm{w}}$; moreover the coefficients of $V(t)$ are analytic functions of $t$ on $\mathbb{D}^{\circ}(t_x^{+})$. As $t_x<t_x^{+}$, this applies to every $t\in\mathbb{D}(t_x)$, which is assumed from now on. In particular Lemma~\ref{lem:derivatives-legs} applies at $A=\mathscr{T}(t)$; for $B=\hat{c}_c$ one has $\epsilon_B=\epsilon_c$, $P_B=P_c$ and $P_B^{+}=P_c^{+}$. In what follows $\mathfrak{H}$ stands for the leg coefficient $\mathfrak{H}_{\mathscr{T}(t)}$ of Section~\ref{sec:dependence-stop-level}, and $\partial_c\mathfrak{H}$ for $\partial_{\hat{c}_c}\mathfrak{H}_A$ at $A=\mathscr{T}(t)$. The position of an index $u=(y,\boldsymbol{\beta},\boldsymbol{\omega})$ is the unit site $y$, and distances between indices and unit sites, or sets of unit sites, are formed with the positions (Sections~\ref{sec:dependence-stop-level} and~\ref{sec:activities-gaussian-form}).

For an even element $\mathbf{F}$ without constant term, with the antisymmetric kernels $\mathbf{F}_m$, and for a symmetric weight $w$, let $\rho_h^{(m)}(\mathbf{F};w)$ and $N_h^{(m)}(\mathbf{F};w)$ be the terms of index $m$ in the definitions of $\rho_h(\mathbf{F};w)$ and $N_h(\mathbf{F};w)$ in Section~\ref{sec:weighted-tree-expansion}. Thus $\rho_h=\sum_m\rho_h^{(m)}$ and $N_h=\sum_mN_h^{(m)}$, and the bounds \eqref{eq:four-jets-9} hold term by term: with $\rho_h^{(m)}$, $N_h^{(m)}$ on the left, and the terms of index $m$ on the right.

\emph{Step~0: auxiliary facts.}
\begin{enumerate}
\item[(i)] \emph{Quadratic elements.} Let $\mathfrak{K}$ be a linear operator on the spinor functions on $T^0_{*}$, with the $4\times4$ blocks $\mathfrak{K}(y,y')$ and the entries $\mathfrak{K}_{\boldsymbol{\beta}\boldsymbol{\beta}'}(y,y')$. The element $-\left\langle\bar{\Psi},\mathfrak{K}\,\Psi\right\rangle$ is of degree two, and its antisymmetric kernel is given by
\[
\mathbf{F}_2\left((y,\boldsymbol{\beta},+),(y',\boldsymbol{\beta}',-)\right)=-\mathfrak{K}_{\boldsymbol{\beta}\boldsymbol{\beta}'}(y,y')=-\mathbf{F}_2\left((y',\boldsymbol{\beta}',-),(y,\boldsymbol{\beta},+)\right),
\]
while it vanishes at pairs of indices with equal signs. Let $w$ be a symmetric weight whose value at two indices depends only on their positions $y$, $y'$, and denote this value by $w(y,y')$. By \eqref{eq:four-jets-Maj2}, $\sum_{\boldsymbol{\beta}'}\left|\mathfrak{K}_{\boldsymbol{\beta}\boldsymbol{\beta}'}(y,y')\right|\le|\mathfrak{K}|(y,y')$ for every $\boldsymbol{\beta}$; summing over $\boldsymbol{\beta}$, one gets $\sum_{\boldsymbol{\beta}}\left|\mathfrak{K}_{\boldsymbol{\beta}\boldsymbol{\beta}'}(y,y')\right|\le4\,|\mathfrak{K}|(y,y')$ for every $\boldsymbol{\beta}'$. Consequently
\begin{align*}
\rho_h\left(-\left\langle\bar{\Psi},\mathfrak{K}\,\Psi\right\rangle;w\right)&\le4\,h_G^2\cdot\max\left\{\sup_y\sum_{y'}|\mathfrak{K}|(y,y')\cdot w(y,y'),\ \sup_{y'}\sum_y|\mathfrak{K}|(y,y')\cdot w(y,y')\right\},\\
N_h\left(-\left\langle\bar{\Psi},\mathfrak{K}\,\Psi\right\rangle;w\right)&\le8\,h_G^2\cdot\sum_{y,y'}|\mathfrak{K}|(y,y')\cdot w(y,y').
\end{align*}
\item[(ii)] \emph{Activities and their positions.} By \eqref{eq:four-jets-16}, $V_{+}(t)$ is the sum of its quadratic part $-\left\langle\bar{\Psi},\check{K}(t)\,\Psi\right\rangle$ and of its activity part, which is the sum over $X\in D_{*}$ of the parts of positive degree of $\mathbf{E}(X;t)$ at $\psi^{\#}=H^{\#}_{*}\left(\mathscr{T}(t)\right)\Psi$. Let $X\in D_{*}$ and $p\ge0$. After this substitution, the part of degree $p$ of $\mathbf{E}(X;t)$ is
\[
\frac{1}{p!}\sum_{u_1,\dots,u_p}\ \sum_{n+m=p}\frac{p!}{n!\,m!}\,\mathcal{K}^X_{nm}(u_1,\dots,u_p)\,\Psi(u_1)\cdots\Psi(u_p),
\]
with
\[
\mathcal{K}^X_{nm}(u_1,\dots,u_p):=\int\mathrm{d}\boldsymbol{\upsilon}\ \mathbf{E}_{nm}(X;t;\boldsymbol{\upsilon}_1,\dots,\boldsymbol{\upsilon}_p)\prod_{i=1}^p\mathfrak{H}(\boldsymbol{\upsilon}_i;u_i);
\]
here $\boldsymbol{\upsilon}_1,\dots,\boldsymbol{\upsilon}_n$ are the arguments $\xi$ and $\boldsymbol{\upsilon}_{n+1},\dots,\boldsymbol{\upsilon}_p$ the arguments $\boldsymbol{\zeta}$ of the kernel, the integral over them is formed with the measures of Section~\ref{sec:grassmann-algebras-norms}, and $\boldsymbol{\upsilon}$ stands for $(\boldsymbol{\upsilon}_1,\dots,\boldsymbol{\upsilon}_p)$. Put
\[
X^{+}:=\left\{y\in T^0_{*}:\ d\left(y,X\cap T^0_{*}\right)\le1\right\}.
\]
If $\mathbf{E}_{nm}(X;t;\boldsymbol{\upsilon}_1,\dots,\boldsymbol{\upsilon}_p)\ne0$, then $\mathrm{pos}(\boldsymbol{\upsilon}_i)\subset X^{+}$ for $i=1,\dots,p$. Indeed, by \eqref{it:trajectory-fermionic-operators-SL3}, and because the arguments of the mass term lie in $X$ (Section~\ref{sec:activities-gaussian-form}), every argument $\xi=(x,\boldsymbol{\beta},\boldsymbol{\omega})$ satisfies $x\in X$, and every argument $\boldsymbol{\zeta}=(x,x',\boldsymbol{\beta},\boldsymbol{\omega})$ is a pair admissible for $X$, of kind~(I) or~(II) in Definition~\ref{def:pairs-sites-holder}. If $x\in X$, then $[x]\in X$. To see this, note that $l_{*}\in\mathrm{Cl}$ by Lemma~\ref{lem:scales-stop-level}(c); hence fact~\eqref{it:renormalisation-group-g0} of Section~\ref{sec:cubes-parameter-cubes} applies at the level $k^{*}$, and it shows that the $M$-cubes form a partition of the set of the sites of $T_{*}$ and that the block $\Delta_{[x]}$ of Section~\ref{sec:lattice-geometry}, which contains $x$ and $[x]$, lies in one $M$-cube. As $X$ is a union of $M$-cubes and contains $x$, this $M$-cube is one of the $M$-cubes of $X$; therefore $[x]\in X$. This disposes of the arguments $\xi$ and of the pairs of kind~(I), for which $x,x'\in X$. A pair of kind~(II) satisfies $|x-x'|=\eta_{*}<1$, so that $\mathrm{pos}(\boldsymbol{\zeta})=\{[x]\}$, and one of the sites $x$, $x'$ lies in $X$; hence $d\left([x],X\cap T^0_{*}\right)\le1$ (if $x'\in X$, then $[x']\in X$, and the unit sites $[x]$, $[x']$ have distance less than $1+\eta_{*}$, hence at most $1$). By \eqref{it:trajectory-fermionic-operators-SL4}(v), two points of $X$ have distance at most $M\cdot\left(d_M(X)+2\right)$; therefore $\operatorname{diam}\left(X^{+}\right)\le M\cdot\left(d_M(X)+2\right)+2$, and, as $\vartheta\le\frac{1}{M}$ and $M\ge2$,
\begin{equation}\label{eq:four-jets-18}
\vartheta\cdot\operatorname{diam}\left(X^{+}\right)\le d_M(X)+3.
\end{equation}
\item[(iii)] \emph{Combinatorial factors.} For $p\ge2$ and $n+m=p$,
\[
\frac{h_G^p}{(p-1)!}\cdot\frac{p!}{n!\,m!}=\frac{p\cdot h_G^p}{n!\,m!},\qquad p\le2^p,\qquad p^2\le3^p,\qquad p\,(p+1)^2\le5^p.
\]
By \eqref{eq:four-jets-sm1}, $5\,h_L\cdot h_G\le e_{*}^{-\frac{1}{4}+\varepsilon}$. Hence Lemma~\ref{lem:bounds-activities} applies with $a=2\,h_L\cdot h_G$, with $a=3\,h_L\cdot h_G$ and with $a=5\,h_L\cdot h_G$.
\item[(iii')] \emph{Bounds in a single degree.} For each degree $p$, the term $\rho_h^{(p)}$ carries a supremum over the first argument that is taken for this degree alone, and $\rho_h$ is the sum of these terms over $p$. Bounds for $\rho_h$ are therefore derived from bounds for every single term $\rho_h^{(p)}$; these require Lemma~\ref{lem:bounds-activities} with the summation restricted to the pairs $(n,m)$ with $n+m=p$. Let $a:=5\,h_L\cdot h_G$; by~(iii), Lemma~\ref{lem:bounds-activities} holds with this $a$. The left sides of its inequalities are sums of non-negative terms over all pairs $(n,m)$, and they do not increase when the terms with $n+m\ne p$ are omitted; for the remaining terms, $\left(h_L\cdot h_G\right)^p=5^{-p}a^{n+m}$. Consequently, for every $X\in D_{*}$, every $p\ge2$ and $i\in\{1,2\}$, Lemma~\ref{lem:bounds-activities}(a) yields
\[
\sum_{n+m=p}\frac{p^i\left(h_L\cdot h_G\right)^p}{n!\,m!}\left\|\mathbf{E}_{nm}(X;t)\right\|\le p^i\,5^{-p}\cdot w^{\mathrm{E}}\cdot\exp\left(-\kappa\cdot d_M(X)\right)\qquad\left(t\in\mathbb{D}(t_x)\right),
\]
and Lemma~\ref{lem:bounds-activities}(b) yields, for $t\in\mathbb{D}(t_x/2)$ and $c\in\mathrm{Cel}$, the inequality that results when $\mathbf{E}_{nm}$ is replaced by $\partial_c\mathbf{E}_{nm}$ and $w^{\mathrm{E}}$ by $q_x\cdot w^{\mathrm{E}}$. Finally, $\sum_{p\ge1}p\,x^p=x\,(1-x)^{-2}$ and $\sum_{p\ge1}p^2x^p=x\,(1+x)\,(1-x)^{-3}$ at $x=\frac{1}{5}$ give
\[
\sum_{p\ge2}p\,5^{-p}=\frac{9}{80},\qquad\sum_{p\ge2}p^2\,5^{-p}=\frac{43}{160}.
\]
\item[(iv)] \emph{Sums over polymers.} For an $M$-cube $\square$, the set $\square^{+}$ is the set $X^{+}$ of the polymer $X=\square$. Let $z\in T^0_{*}$ and $\frac{\gamma_F}{4}\le\tilde{\gamma}\le\frac{1}{2}$. Then
\[
\sum_{\square}\exp\left(-\tilde{\gamma}\cdot d\left(z,\square^{+}\right)\right)\le32,
\]
the sum being taken over all $M$-cubes. To see this, label the $M$-cubes by vectors $i\in\mathbb{Z}^3$ that give their position relative to the $M$-cube containing $z$, each cube being labelled by a representative of least norm $|i|$. Then $d\left(z,\square_i^{+}\right)\ge(|i|-1)\cdot M-1$, and at most $(2k+1)^3$ cubes satisfy $|i|=k$. Since $\tilde{\gamma}\cdot M\ge\frac{\gamma_F\cdot M}{4}\ge4$ and $\tilde{\gamma}\le\frac{1}{2}$, the sum is at most
\[
27+\exp\left(\frac{1}{2}\right)\sum_{k\ge2}(2k+1)^3\exp\left(-4\,(k-1)\right)\le32.
\]
The set $X^{+}$ is the union of the sets $\square^{+}$ over the $M$-cubes $\square\subset X$; hence $\exp\left(-\tilde{\gamma}\cdot d\left(z,X^{+}\right)\right)\le\sum_{\square\subset X}\exp\left(-\tilde{\gamma}\cdot d\left(z,\square^{+}\right)\right)$. Together with $\kappa-1\ge\kappa_0$ and \eqref{it:trajectory-fermionic-operators-SL4}(v) this yields
\begin{align}
\sum_{X\in D_{*}}\exp\left(-\tilde{\gamma}\cdot d\left(z,X^{+}\right)\right)\cdot\exp\left(-(\kappa-1)\cdot d_M(X)\right)&\le\sum_{\square}\exp\left(-\tilde{\gamma}\cdot d\left(z,\square^{+}\right)\right)\sum_{X\supset\square}\exp\left(-\kappa_0\cdot d_M(X)\right)\notag\\
&\le32\,K_0.\label{eq:four-jets-19}
\end{align}
\item[(v)] \emph{Cubes meeting $P_c^{+}$.} Select in every $M$-cube that meets $P_c^{+}$ one point of $P_c^{+}$. Distinct $M$-cubes have no unit site in common, so this selection is an injective map from the set of these cubes into $P_c^{+}$. As $\#|P_c^{+}|\le\mathsf{p}$ by Lemma~\ref{lem:properties-probes}(e), the number of $M$-cubes that meet $P_c^{+}$ does not exceed $\mathsf{p}$.
\end{enumerate}

\emph{Step~1: proof of \eqref{eq:four-jets-V1}.} Let $t\in\mathbb{D}(t_x)$.

\emph{Quadratic part.} Corollary~\ref{cor:increments-along-segment}, applied with $A_0=\mathscr{A}$ and $B=\mathscr{C}(t)$ (both $\mathscr{A}$ and $\mathscr{T}(t)$, and the segment between them, lie in $G_{\mathrm{w}}$), gives
\[
\left\|\check{K}(t)\right\|_{\gamma_F/2}\le C_K\cdot e_{*}\cdot|\mathscr{C}(t)|_{\infty}\le C_K\cdot e_{*}^{2\varepsilon}.
\]
At two indices with the positions $y$, $y'$ the weight is $w_{\boldsymbol{\delta}}=\exp\left(\vartheta\cdot d(y,y')\right)$, and $\vartheta\le\frac{\gamma_F}{2}$; so the two suprema in~(i), formed with $\mathfrak{K}=\check{K}(t)$ and $w=w_{\boldsymbol{\delta}}$, are at most $\left\|\check{K}(t)\right\|_{\gamma_F/2}$. By~(i), the quadratic part has $\rho_h\left(-\left\langle\bar{\Psi},\check{K}(t)\,\Psi\right\rangle;w_{\boldsymbol{\delta}}\right)\le4\,h_G^2\cdot C_K\cdot e_{*}^{2\varepsilon}$.

\emph{Activity part.} Let $X\in D_{*}$ and $n+m=p\ge2$, and let $\boldsymbol{\upsilon}_1,\dots,\boldsymbol{\upsilon}_p$ be arguments at which $\mathbf{E}_{nm}(X;t;\cdot)$ does not vanish. For each $i$ choose $z_i\in\mathrm{pos}(\boldsymbol{\upsilon}_i)$ with $d(u_i,z_i)=d\left(u_i,\mathrm{pos}(\boldsymbol{\upsilon}_i)\right)$. By~(ii), all $z_i$ lie in $X^{+}$, so that $d(u_i,u_j)\le d(u_i,z_i)+\operatorname{diam}\left(X^{+}\right)+d(z_j,u_j)$ for all $i$, $j$; by \eqref{eq:four-jets-18},
\[
w_{\boldsymbol{\delta}}(u_1,\dots,u_p)\le\exp\left(d_M(X)+3\right)\cdot\prod_{i=1}^p\exp\left(\vartheta\cdot d\left(u_i,\mathrm{pos}(\boldsymbol{\upsilon}_i)\right)\right).
\]
Fix a slot $\mathrm{sl}\in\{1,\dots,p\}$ and put $u_{\mathrm{sl}}:=u$. For $i\ne\mathrm{sl}$, the third bound of \eqref{eq:four-jets-Leg1}, with $\tilde{\gamma}=\vartheta\le\frac{\gamma_F}{2}$, gives
\[
\sum_{u_i}\left|\mathfrak{H}(\boldsymbol{\upsilon}_i;u_i)\right|\cdot\exp\left(\vartheta\cdot d\left(u_i,\mathrm{pos}(\boldsymbol{\upsilon}_i)\right)\right)\le h_L ;
\]
for $i=\mathrm{sl}$, the second bound of \eqref{eq:four-jets-Leg1} and the inclusion $\mathrm{pos}(\boldsymbol{\upsilon}_{\mathrm{sl}})\subset X^{+}$ give
\[
\left|\mathfrak{H}(\boldsymbol{\upsilon}_{\mathrm{sl}};u)\right|\cdot\exp\left(\vartheta\cdot d\left(u,\mathrm{pos}(\boldsymbol{\upsilon}_{\mathrm{sl}})\right)\right)\le h_L\cdot\exp\left(-(\gamma_F-\vartheta)\cdot d\left(u,X^{+}\right)\right).
\]
Integration over $\boldsymbol{\upsilon}$ yields, for every slot $\mathrm{sl}$ and every $u$,
\[
\sum_{(u_i)_{i\ne\mathrm{sl}}}\left|\mathcal{K}^X_{nm}(u_1,\dots,u_p)\right|\cdot w_{\boldsymbol{\delta}}(u_1,\dots,u_p)\le\exp\left(d_M(X)+3\right)\cdot h_L^p\cdot\exp\left(-(\gamma_F-\vartheta)\cdot d\left(u,X^{+}\right)\right)\cdot\left\|\mathbf{E}_{nm}(X;t)\right\|.
\]
By~(ii), the part of degree $p$ of the activity part of $V_{+}(t)$ has the kernel $\sum_X\sum_{n+m=p}\frac{p!}{n!\,m!}\,\mathcal{K}^X_{nm}$, which need not be antisymmetric. Therefore \eqref{eq:four-jets-9} in degree $p$, the identity in~(iii) and the last bound show that $\rho_h^{(p)}$ of the activity part is at most
\[
\exp(3)\cdot\sup_u\sum_{X\in D_{*}}\exp\left(d_M(X)\right)\cdot\exp\left(-(\gamma_F-\vartheta)\cdot d\left(u,X^{+}\right)\right)\sum_{n+m=p}\frac{p\left(h_L\cdot h_G\right)^p}{n!\,m!}\left\|\mathbf{E}_{nm}(X;t)\right\|.
\]
Apply~(iii') with $i=1$, then $\gamma_F-\vartheta\ge\frac{\gamma_F}{4}$ (as $\vartheta\le\frac{\gamma_F}{8}$), and \eqref{eq:four-jets-19} with $\tilde{\gamma}=\frac{\gamma_F}{4}$. It follows that this term is at most
\[
p\,5^{-p}\cdot\exp(3)\cdot w^{\mathrm{E}}\cdot\sup_u\sum_{X\in D_{*}}\exp\left(-\frac{\gamma_F}{4}\,d\left(u,X^{+}\right)\right)\cdot\exp\left(-(\kappa-1)\cdot d_M(X)\right)\le p\,5^{-p}\cdot32\exp(3)\cdot K_0\cdot w^{\mathrm{E}} .
\]
Summing over $p\ge2$ and using $\sum_{p\ge2}p\,5^{-p}\le1$ from~(iii'), one finds that the activity part has $\rho_h(\,\cdot\,;w_{\boldsymbol{\delta}})$ at most $32\exp(3)\cdot K_0\cdot w^{\mathrm{E}}\le650\,K_0\cdot w^{\mathrm{E}}$. Since $\rho_h$ is subadditive and $w^{\mathrm{E}}=e_{*}^{\frac{1}{4}-7\varepsilon}$,
\[
\rho_h\left(V_{+}(t);w_{\boldsymbol{\delta}}\right)\le4\,h_G^2\cdot C_K\cdot e_{*}^{2\varepsilon}+650\,K_0\cdot e_{*}^{\frac{1}{4}-7\varepsilon}=\bar{r} ,
\]
which is \eqref{eq:four-jets-V1}.

\emph{Step~2: proof of \eqref{eq:four-jets-V2}.} Let $\hat{\mathfrak{a}}\ne0$, $t\in\mathbb{D}(t_x/2)$ and $c\in\mathrm{Cel}$. As shown at the beginning of the proof, the sets $P_c$ and $P_c^{+}$ are non-empty, so that the weight $w_{P_c}$ is defined.

\emph{Quadratic part.} Since $\mathscr{T}(t)=\mathscr{A}+\sum_{c''\in\mathrm{Cel}}t_{c''}\cdot\hat{c}_{c''}$, one has $\partial_c\check{K}(t)=\left(\partial_{\hat{c}_c}D_{*}\right)\left(\mathscr{T}(t)\right)$. For $z\in P_c$ and unit sites $y$, $y'$,
\[
\max\left\{d(y,P_c),\,d(y',P_c)\right\}\le d(y,z)+d(z,y');
\]
hence, at two indices with the positions $y$, $y'$, the weight satisfies $w_{P_c}\le\exp\left(\vartheta\cdot\left(d(y,z)+d(z,y')\right)\right)$ for every $z\in P_c$. Insert \eqref{eq:four-jets-Dp}, use this bound for the weight in the term $z$ of its sum, and apply~(i). As $\gamma_F-\vartheta\ge\frac{\gamma_F}{2}$, one has $\sum_y\exp\left(-(\gamma_F-\vartheta)\cdot d(y,z)\right)\le\check{s}$; moreover $\#|P_c|\le\#|P_c^{+}|\le\mathsf{p}$ by Lemma~\ref{lem:properties-probes}(e), and $67\,C_F^2\cdot\check{s}^2\le100\,C_F^2\cdot\check{s}^2=C_K$. Therefore
\[
N_h\left(-\left\langle\bar{\Psi},\partial_c\check{K}(t)\,\Psi\right\rangle;w_{P_c}\right)\le8\,h_G^2\cdot67\,C_F^2\cdot\epsilon_c\sum_{z\in P_c}\sum_{y,y'}\exp\left(-(\gamma_F-\vartheta)\cdot\left(d(y,z)+d(z,y')\right)\right)\le8\,h_G^2\cdot C_K\cdot\epsilon_c\cdot\mathsf{p} .
\]
In the same way, the sums of $\left|\partial_c\check{K}(t)\right|(y,y')\cdot w_{P_c}(y,y')$ over $y'$ at fixed $y$, and over $y$ at fixed $y'$, are at most $67\,C_F^2\cdot\check{s}^2\cdot\epsilon_c\le C_K\cdot\epsilon_c$, and~(i) gives
\[
\rho_h\left(-\left\langle\bar{\Psi},\partial_c\check{K}(t)\,\Psi\right\rangle;w_{P_c}\right)\le4\,h_G^2\cdot C_K\cdot\epsilon_c.
\]

\emph{Activity part.} By the product rule,
\begin{align*}
\partial_c\left[\mathbf{E}_{nm}(X;t;\boldsymbol{\upsilon})\prod_{i=1}^p\mathfrak{H}(\boldsymbol{\upsilon}_i;u_i)\right]&=\left(\partial_c\mathbf{E}_{nm}(X;t;\boldsymbol{\upsilon})\right)\prod_{i=1}^p\mathfrak{H}(\boldsymbol{\upsilon}_i;u_i)\\
&\quad+\mathbf{E}_{nm}(X;t;\boldsymbol{\upsilon})\sum_{j=1}^p\left(\partial_c\mathfrak{H}\right)(\boldsymbol{\upsilon}_j;u_j)\prod_{i\ne j}\mathfrak{H}(\boldsymbol{\upsilon}_i;u_i).
\end{align*}
The two terms on the right are estimated separately.

\emph{(A) The first term.} If $X\cap P_c=\emptyset$, then $\partial_c\mathbf{E}_{nm}(X;t)=0$ by Lemma~\ref{lem:bounds-activities}(c), and the first term is zero. Let therefore $X$ meet $P_c$. A common point of $X$ and $P_c$ is a unit site of $X$ and hence an element of $X^{+}$; consequently $d(u_i,P_c)\le d\left(u_i,\mathrm{pos}(\boldsymbol{\upsilon}_i)\right)+\operatorname{diam}\left(X^{+}\right)$ for $i=1,\dots,p$, and \eqref{eq:four-jets-18} yields
\[
w_{P_c}(u_1,\dots,u_p)\le\exp\left(d_M(X)+3\right)\cdot\prod_{i=1}^p\exp\left(\vartheta\cdot d\left(u_i,\mathrm{pos}(\boldsymbol{\upsilon}_i)\right)\right).
\]
With \eqref{eq:four-jets-Leg1} for all legs it follows that
\[
\sum_{u_1,\dots,u_p}\int\mathrm{d}\boldsymbol{\upsilon}\ \left|\partial_c\mathbf{E}_{nm}(X;t;\boldsymbol{\upsilon})\right|\prod_{i=1}^p\left|\mathfrak{H}(\boldsymbol{\upsilon}_i;u_i)\right|\cdot w_{P_c}(u_1,\dots,u_p)\le\exp\left(d_M(X)+3\right)\cdot h_L^p\cdot\left\|\partial_c\mathbf{E}_{nm}(X;t)\right\|.
\]
By~(iii) ($p\le2^p$) and Lemma~\ref{lem:bounds-activities}(b) with $a=2\,h_L\cdot h_G$, the first terms contribute to $N_h\left(\partial_cV_{+}(t);w_{P_c}\right)$ at most
\[
\exp(3)\cdot q_x\cdot w^{\mathrm{E}}\sum_X\exp\left(-(\kappa-1)\cdot d_M(X)\right)\le\exp(3)\cdot K_0\cdot\mathsf{p}\cdot q_x\cdot w^{\mathrm{E}} ,
\]
where $X$ ranges over the polymers in $D_{*}$ that contain a point of $P_c$. Indeed, each such $X$ contains an $M$-cube $\square$ that contains a point of $P_c\subset P_c^{+}$; by~(v) there are at most $\mathsf{p}$ such cubes, and for each of them $\sum_{X\supset\square}\exp\left(-(\kappa-1)\cdot d_M(X)\right)\le K_0$, by $\kappa-1\ge\kappa_0$ and \eqref{it:trajectory-fermionic-operators-SL4}(v). For $\rho_h$ the argument of Step~1 is repeated, one degree at a time. With $u$ fixed in a slot, the sum over the other arguments is at most $\exp\left(d_M(X)+3\right)\cdot h_L^p\cdot\exp\left(-(\gamma_F-\vartheta)\cdot d\left(u,X^{+}\right)\right)\cdot\left\|\partial_c\mathbf{E}_{nm}(X;t)\right\|$; by~(iii), by~(iii') for $\partial_c\mathbf{E}_{nm}$ with $i=1$, and by \eqref{eq:four-jets-19}, the term of degree $p$ is at most $p\,5^{-p}\cdot32\exp(3)\cdot K_0\cdot q_x\cdot w^{\mathrm{E}}$. Summing over $p$, the first terms contribute at most $650\,K_0\cdot q_x\cdot w^{\mathrm{E}}$ to $\rho_h\left(\partial_cV_{+}(t);w_{P_c}\right)$.

\emph{(B) The second term, for a fixed $j$.} Let $y_j$ be the position of $u_j$. By \eqref{eq:four-jets-Leg2},
\[
\left|\partial_c\mathfrak{H}\right|(\boldsymbol{\upsilon}_j;y_j)\le C_H\cdot\epsilon_c\sum_{w\in\mathrm{pos}(\boldsymbol{\upsilon}_j)}\ \sum_{z\in P_c^{+}}E(w,z)\cdot E(z,y_j).
\]
For $w\in\mathrm{pos}(\boldsymbol{\upsilon}_j)$ and $z\in P_c^{+}$ one has $d(u_j,P_c)\le d(u_j,z)+M+1$ and, for $i\ne j$, $d(u_i,P_c)\le d\left(u_i,\mathrm{pos}(\boldsymbol{\upsilon}_i)\right)+\operatorname{diam}\left(X^{+}\right)+d(w,z)+M+1$, because $w\in X^{+}$ and $d(z,P_c)\le M+1$. As $\vartheta\cdot(M+1)\le2$, \eqref{eq:four-jets-18} gives, in the term $(w,z)$ of this sum,
\[
w_{P_c}(u_1,\dots,u_p)\le\exp\left(d_M(X)+5\right)\cdot\exp\left(\vartheta\cdot\left(d(w,z)+d(z,u_j)\right)\right)\cdot\prod_{i\ne j}\exp\left(\vartheta\cdot d\left(u_i,\mathrm{pos}(\boldsymbol{\upsilon}_i)\right)\right).
\]
As $\vartheta\le\frac{\gamma_F}{2}$ and $\mathrm{pos}(\boldsymbol{\upsilon}_j)$ consists of at most two points of $X^{+}$,
\[
\sum_{y_j}E(z,y_j)\cdot\exp\left(\vartheta\cdot d(z,y_j)\right)\le\check{s},\qquad\sum_{w\in\mathrm{pos}(\boldsymbol{\upsilon}_j)}E(w,z)\cdot\exp\left(\vartheta\cdot d(w,z)\right)\le2\exp\left(-\frac{\gamma_F}{2}\,d\left(z,X^{+}\right)\right).
\]
With \eqref{eq:four-jets-Leg1} for the legs $i\ne j$ it follows that
\begin{multline*}
\sum_{u_1,\dots,u_p}\int\mathrm{d}\boldsymbol{\upsilon}\ \left|\mathbf{E}_{nm}(X;t;\boldsymbol{\upsilon})\right|\cdot\left|\left(\partial_c\mathfrak{H}\right)(\boldsymbol{\upsilon}_j;u_j)\right|\prod_{i\ne j}\left|\mathfrak{H}(\boldsymbol{\upsilon}_i;u_i)\right|\cdot w_{P_c}(u_1,\dots,u_p)\\
\le2\exp\left(d_M(X)+5\right)\cdot h_L^{p-1}\cdot C_H\cdot\epsilon_c\cdot\check{s}\sum_{z\in P_c^{+}}\exp\left(-\frac{\gamma_F}{2}\,d\left(z,X^{+}\right)\right)\cdot\left\|\mathbf{E}_{nm}(X;t)\right\|.
\end{multline*}
The index $j$ takes $p$ values. By~(iii) ($p^2\le3^p$, together with $h_L\ge1$), Lemma~\ref{lem:bounds-activities}(a) with $a=3\,h_L\cdot h_G$, \eqref{eq:four-jets-19} with $\tilde{\gamma}=\frac{\gamma_F}{2}$ and $\#|P_c^{+}|\le\mathsf{p}$, the second terms contribute to $N_h\left(\partial_cV_{+}(t);w_{P_c}\right)$ at most
\[
2\exp(5)\cdot32\,C_H\cdot\check{s}\cdot K_0\cdot w^{\mathrm{E}}\cdot\mathsf{p}\cdot\epsilon_c\le9500\,C_H\cdot\check{s}\cdot K_0\cdot w^{\mathrm{E}}\cdot\mathsf{p}\cdot\epsilon_c.
\]
To estimate $\rho_h$, the argument $u$ in one slot is held fixed and all other arguments are summed; two cases arise. Suppose first that $u$ occupies the slot of the leg $j$, so that $u_j=u$. As $\vartheta\le\frac{\gamma_F}{8}$ and $\frac{7}{8}=\frac{3}{8}+\frac{1}{2}$,
\[
\sum_zE(w,z)\cdot E(z,u)\cdot\exp\left(\vartheta\cdot\left(d(w,z)+d(z,u)\right)\right)\le\sum_z\exp\left(-\frac{7\gamma_F}{8}\left(d(w,z)+d(z,u)\right)\right)\le\check{s}\cdot\exp\left(-\frac{3\gamma_F}{8}\,d(w,u)\right),
\]
and, since $\mathrm{pos}(\boldsymbol{\upsilon}_j)$ consists of at most two points of $X^{+}$ and $d(w,u)\ge d\left(u,X^{+}\right)$ for $w\in X^{+}$, summation of the right side over $w\in\mathrm{pos}(\boldsymbol{\upsilon}_j)$ gives at most $2\,\check{s}\cdot\exp\left(-\frac{\gamma_F}{4}\,d\left(u,X^{+}\right)\right)$. Suppose next that $u$ occupies the slot of a leg $i\ne j$. The second bound of \eqref{eq:four-jets-Leg1} and $\mathrm{pos}(\boldsymbol{\upsilon}_i)\subset X^{+}$ show that the modulus of this leg, multiplied by its factor $\exp\left(\vartheta\cdot d\left(u,\mathrm{pos}(\boldsymbol{\upsilon}_i)\right)\right)$ from the weight, is at most $h_L\cdot\exp\left(-(\gamma_F-\vartheta)\cdot d\left(u,X^{+}\right)\right)$. For the differentiated leg $j$, insert the factor $\exp\left(\vartheta\cdot\left(d(w,z)+d(z,u_j)\right)\right)$ in the term $(w,z)$ of the bound for $\left|\partial_c\mathfrak{H}\right|(\boldsymbol{\upsilon}_j;y_j)$ given by \eqref{eq:four-jets-Leg2}; after summation over $u_j$, the result does not exceed $2\,C_H\cdot\epsilon_c\cdot\check{s}^2$. As $\check{s}\ge1$ and $\gamma_F-\vartheta\ge\frac{\gamma_F}{4}$, either case leads to the following bound for the sum over all arguments other than $u$:
\[
2\exp\left(d_M(X)+5\right)\cdot h_L^{p-1}\cdot C_H\cdot\epsilon_c\cdot\check{s}^2\cdot\exp\left(-\frac{\gamma_F}{4}\,d\left(u,X^{+}\right)\right)\cdot\left\|\mathbf{E}_{nm}(X;t)\right\|.
\]
With the $p$ values of $j$, with~(iii) and~(iii') for $i=2$ (and $h_L\ge1$), and with \eqref{eq:four-jets-19} for $\tilde{\gamma}=\frac{\gamma_F}{4}$, the term of degree $p$ is at most $p^2\,5^{-p}\cdot2\exp(5)\cdot32\,C_H\cdot\check{s}^2\cdot K_0\cdot w^{\mathrm{E}}\cdot\epsilon_c$. Summing over $p$, with $\sum_{p\ge2}p^2\,5^{-p}\le1$ from~(iii'), the second terms contribute to $\rho_h\left(\partial_cV_{+}(t);w_{P_c}\right)$ at most
\[
2\exp(5)\cdot32\,C_H\cdot\check{s}^2\cdot K_0\cdot w^{\mathrm{E}}\cdot\epsilon_c\le9500\,C_H\cdot\check{s}^2\cdot K_0\cdot w^{\mathrm{E}}\cdot\epsilon_c.
\]

\emph{Conclusion of Step~2.} Adding the three contributions and using $\epsilon_c\le q_x$, $w^{\mathrm{E}}\le\frac{1}{8}$, $h_G\ge1$ and $C_K=100\,C_F^2\cdot\check{s}^2\le C_H\cdot\check{s}^2$, one obtains
\begin{align*}
N_h\left(\partial_cV_{+}(t);w_{P_c}\right)&\le q_x\cdot\mathsf{p}\cdot h_G^2\cdot\left(8\,C_K+3\,K_0+1200\,C_H\cdot\check{s}\cdot K_0\right)\le q_x\cdot C_V\cdot h_G^2\cdot\mathsf{p},\\
\rho_h\left(\partial_cV_{+}(t);w_{P_c}\right)&\le q_x\cdot h_G^2\cdot\left(4\,C_K+82\,K_0+1200\,C_H\cdot\check{s}^2\cdot K_0\right)\le q_x\cdot C_V\cdot h_G^2.
\end{align*}
The last inequality in each line uses $K_0\ge1$, which is \eqref{it:trajectory-fermionic-operators-SL4}(v) at $X=\square$ (where $d_M(\square)=0$), together with $C_H\ge100$, $\check{s}\ge1$ and $C_V=10^3\,C_H^2\cdot\check{s}^4\cdot K_0$. This proves \eqref{eq:four-jets-V2}.

\emph{Step~3: proof of \eqref{eq:four-jets-V3}.} Let $\hat{\mathfrak{a}}\ne0$ and $t\in\mathbb{D}(t_x/2)$, and let $c,c'\in\mathrm{Cel}$ with $c\ne c'$. As shown at the beginning of the proof, the sets $P_c$, $P_{c'}$, $P_c^{+}$ and $P_{c'}^{+}$ are non-empty. Put $\mathfrak{d}_0:=d\left(P_c^{+},P_{c'}^{+}\right)$. Every point of $P_c^{+}$, resp.\ $P_{c'}^{+}$, has distance at most $M+1$ from $P_c$, resp.\ $P_{c'}$, and $d(P_c,P_{c'})\ge d_{\mathrm{sep}}$; hence $\mathfrak{d}_0\ge d_{\mathrm{sep}}-2M-2$.

\emph{Quadratic part.} One has $\partial_c\partial_{c'}\check{K}(t)=\left(\partial_{\hat{c}_c}\partial_{\hat{c}_{c'}}D_{*}\right)\left(\mathscr{T}(t)\right)$, and \eqref{eq:four-jets-Dpp} applies, since $d(P_c,P_{c'})\ge M+2$. Apply~(i) with the weight $1$; as $\sum_yE(y,z)\le\mathrm{ssum}(\gamma_F)$, each of the two chains in \eqref{eq:four-jets-Dpp} contributes, after summation over $y$ and $y'$, at most $\mathrm{ssum}(\gamma_F)^2\cdot E(z,z')$. Thus
\[
N_h\left(-\left\langle\bar{\Psi},\partial_c\partial_{c'}\check{K}(t)\,\Psi\right\rangle;1\right)\le8\,h_G^2\cdot C_H^2\cdot\epsilon_c\cdot\epsilon_{c'}\cdot2\sum_{z\in P_c}\ \sum_{z'\in P_{c'}}\mathrm{ssum}(\gamma_F)^2\cdot E(z,z').
\]
Here $\mathrm{ssum}(\gamma_F)\le\check{s}$; for $z\in P_c$ and $z'\in P_{c'}$ one has $E(z,z')\le\exp\left(-\frac{\gamma_F}{2}\,d_{\mathrm{sep}}\right)\cdot\exp\left(-\frac{\gamma_F}{2}\,d(z,z')\right)$, and $\sum_{z'}\exp\left(-\frac{\gamma_F}{2}\,d(z,z')\right)\le\check{s}$; moreover $\#|P_c|\le\mathsf{p}$ and $\epsilon_c,\epsilon_{c'}\le q_x$. So the right-hand side is at most
\[
16\,h_G^2\cdot C_H^2\cdot\check{s}^3\cdot\mathsf{p}\cdot q_x^2\cdot\exp\left(-\frac{\gamma_F}{2}\,d_{\mathrm{sep}}\right).
\]

\emph{Activity part, together with $V_0$.} The part of degree zero of $V(t)$ is $V_0(t)=\sum_X\mathbf{E}_{00}(X;t)$, so that $\left|\partial_c\partial_{c'}V_0(t)\right|\le\sum_X\left|\partial_c\partial_{c'}\mathbf{E}_{00}(X;t)\right|$; the case $p=0$ is treated together with the degrees $p\ge2$. For such $p$, the integrand of $\mathcal{K}^X_{nm}$ is the product of $F_0:=\mathbf{E}_{nm}(X;t;\boldsymbol{\upsilon})$ and of the legs $F_i:=\mathfrak{H}(\boldsymbol{\upsilon}_i;u_i)$, $1\le i\le p$. By the product rule, $\partial_c\partial_{c'}\prod_{i=0}^pF_i$ is the sum, over the pairs $(k,l)\in\{0,\dots,p\}^2$, of the product of $F_0,\dots,F_p$ in which $\partial_c$ acts on $F_k$ and $\partial_{c'}$ on $F_l$. Of these $(p+1)^2$ products, the $p^2$ with $k,l\ge1$ leave $F_0$ undifferentiated, the $p$ with $k=0<l$ and the $p$ with $l=0<k$ differentiate $F_0$ by $\partial_c$ only and by $\partial_{c'}$ only, respectively, and the one with $k=l=0$ differentiates $F_0$ by $\partial_c\partial_{c'}$. Put
\[
d_1:=d\left(X^{+},P_c^{+}\right),\qquad d_2:=d\left(X^{+},P_{c'}^{+}\right).
\]
For every non-empty set $P$ of unit sites and every unit site $w$ one has $\sum_{z\in P}E(w,z)\le\check{s}\cdot\exp\left(-\frac{\gamma_F}{2}\,d(w,P)\right)$. This inequality, together with \eqref{eq:four-jets-Leg1}, \eqref{eq:four-jets-Leg2} and \eqref{eq:four-jets-Leg3} (the hypothesis $d(P_c,P_{c'})\ge M+2$ of the last one is satisfied), gives, for $1\le i\le p$ and every argument $\boldsymbol{\upsilon}_i$ with $\mathrm{pos}(\boldsymbol{\upsilon}_i)\subset X^{+}$,
\begin{itemize}
\item $\sum_{u_i}\left|F_i\right|\le h_L$;
\item $\sum_{u_i}\left|\partial_cF_i\right|\le2\,C_H\cdot\epsilon_c\cdot\check{s}^2\cdot\exp\left(-\frac{\gamma_F}{2}\,d_1\right)$;
\item $\sum_{u_i}\left|\partial_{c'}F_i\right|\le2\,C_H\cdot\epsilon_{c'}\cdot\check{s}^2\cdot\exp\left(-\frac{\gamma_F}{2}\,d_2\right)$;
\item $\sum_{u_i}\left|\partial_c\partial_{c'}F_i\right|\le2\,C_H^2\cdot\epsilon_c\cdot\epsilon_{c'}\cdot\check{s}^3\cdot\exp\left(-\frac{\gamma_F}{2}\,\mathfrak{d}_0\right)\cdot\left[\exp\left(-\frac{\gamma_F}{2}\,d_1\right)+\exp\left(-\frac{\gamma_F}{2}\,d_2\right)\right]$.
\end{itemize}
By Lemma~\ref{lem:bounds-activities}(c), $\partial_cF_0$ vanishes unless $d_1=0$, $\partial_{c'}F_0$ vanishes unless $d_2=0$, and $\partial_c\partial_{c'}F_0$ vanishes unless $d_1=d_2=0$; indeed, a point of $P_c$ in $X$ lies in $X^{+}\cap P_c^{+}$, and likewise for $c'$. Using $\epsilon_c,\epsilon_{c'}\le q_x$, $h_L\ge1$, $\check{s}\ge1$ and $C_H\ge1$, one concludes that every product, summed over the $u_i$ and integrated over $\boldsymbol{\upsilon}$, is at most
\[
q_x^2\cdot4\,C_H^2\cdot\check{s}^4\cdot h_L^p
\]
times one of the numbers $\left\|\mathbf{E}_{nm}(X;t)\right\|$, $q_x^{-1}\left\|\partial_c\mathbf{E}_{nm}(X;t)\right\|$, $q_x^{-1}\left\|\partial_{c'}\mathbf{E}_{nm}(X;t)\right\|$, $q_x^{-2}\left\|\partial_c\partial_{c'}\mathbf{E}_{nm}(X;t)\right\|$, according to the derivative carried by $F_0$, and times $\exp\left(-\frac{\gamma_F}{2}\left(d_1+d_2\right)\right)$; for the products in which one leg carries both derivatives, the last factor is replaced by $\exp\left(-\frac{\gamma_F}{2}\,\mathfrak{d}_0\right)\cdot\left[\exp\left(-\frac{\gamma_F}{2}\,d_1\right)+\exp\left(-\frac{\gamma_F}{2}\,d_2\right)\right]$.

By~(ii), $\mathfrak{d}_0\le d_1+\operatorname{diam}\left(X^{+}\right)+d_2\le d_1+d_2+M\cdot d_M(X)+2M+2$. Together with $\vartheta\le\min\left\{\frac{\gamma_F}{8},\frac{1}{M}\right\}$ this gives
\[
\exp\left(-\frac{\gamma_F}{2}\left(d_1+d_2\right)\right)\cdot\exp\left(-\kappa\cdot d_M(X)\right)\le\exp\left(-\vartheta\cdot\left(\mathfrak{d}_0-2M-2\right)\right)\cdot\exp\left(-\frac{\gamma_F}{4}\,d_1\right)\cdot\exp\left(-(\kappa-1)\cdot d_M(X)\right),
\]
and also $\exp\left(-\frac{\gamma_F}{2}\,\mathfrak{d}_0\right)\le\exp\left(-\vartheta\cdot\mathfrak{d}_0\right)$. Hence, for every product, the last factor multiplied by $\exp\left(-\kappa\cdot d_M(X)\right)$ is at most
\[
\exp\left(-\vartheta\cdot\left(\mathfrak{d}_0-2M-2\right)\right)\cdot\left[\exp\left(-\frac{\gamma_F}{4}\,d_1\right)+\exp\left(-\frac{\gamma_F}{4}\,d_2\right)\right]\cdot\exp\left(-(\kappa-1)\cdot d_M(X)\right).
\]
The products are now divided into four groups according to the derivative that acts on $F_0$. Each group contains at most $(p+1)^2$ products of degree $p$; for $p=0$ the only product is $\partial_c\partial_{c'}F_0$. Let $a=5\,h_L\cdot h_G$ be as in~(iii'). For $p\ge2$, the coefficient $\frac{p\cdot h_G^p}{n!\,m!}$ of~(iii), multiplied by the number of products in a group and by the factor $h_L^p$ of the bound for a single product, is at most $\frac{p\,(p+1)^2\left(h_L\cdot h_G\right)^p}{n!\,m!}\le\frac{a^{n+m}}{n!\,m!}$, by~(iii); for $p=0$ the coefficient is $1=\frac{a^0}{0!\,0!}$. Hence, in each group, the sum over all pairs $(n,m)$ of $\frac{a^{n+m}}{n!\,m!}$ times the entry of the list $\left\|\mathbf{E}_{nm}(X;t)\right\|,\dots,q_x^{-2}\left\|\partial_c\partial_{c'}\mathbf{E}_{nm}(X;t)\right\|$ that belongs to the group is at most $w^{\mathrm{E}}\cdot\exp\left(-\kappa\cdot d_M(X)\right)$, by Lemma~\ref{lem:bounds-activities}(a), (b). Further, $d_1$ is the least of the distances $d\left(z,X^{+}\right)$ with $z\in P_c^{+}$, so that
\[
\exp\left(-\frac{\gamma_F}{4}\,d_1\right)\le\sum_{z\in P_c^{+}}\exp\left(-\frac{\gamma_F}{4}\,d\left(z,X^{+}\right)\right),
\]
and, as $d_2$ is the least of the distances $d\left(z,X^{+}\right)$ with $z\in P_{c'}^{+}$, the number $\exp\left(-\frac{\gamma_F}{4}\,d_2\right)$ is at most the sum of $\exp\left(-\frac{\gamma_F}{4}\,d\left(z,X^{+}\right)\right)$ over $z\in P_{c'}^{+}$. By \eqref{eq:four-jets-19} with $\tilde{\gamma}=\frac{\gamma_F}{4}$ and by $\#|P_c^{+}|,\#|P_{c'}^{+}|\le\mathsf{p}$, each of these two sums, multiplied by $\exp\left(-(\kappa-1)\cdot d_M(X)\right)$ and summed over $X\in D_{*}$, is at most $32\,K_0\cdot\mathsf{p}$. Taking into account the four groups, $w^{\mathrm{E}}\le\frac{1}{8}$ and $\mathfrak{d}_0-2M-2\ge d_{\mathrm{sep}}-4M-4$, one finds that the activity part and $V_0$ contribute to the left side of \eqref{eq:four-jets-V3} at most
\begin{multline*}
4\cdot q_x^2\cdot4\,C_H^2\cdot\check{s}^4\cdot w^{\mathrm{E}}\cdot2\cdot32\,K_0\cdot\mathsf{p}\cdot\exp\left(-\vartheta\cdot\left(\mathfrak{d}_0-2M-2\right)\right)\\
\le128\,C_H^2\cdot\check{s}^4\cdot K_0\cdot\mathsf{p}\cdot q_x^2\cdot\exp\left(-\vartheta\cdot\left(d_{\mathrm{sep}}-4M-4\right)\right).
\end{multline*}

\emph{Conclusion of Step~3.} Since $K_0\ge1$ (by \eqref{it:trajectory-fermionic-operators-SL4}(v) at $X=\square$), $\check{s}\ge1$ and $\vartheta\le\frac{\gamma_F}{2}$, the bound for the quadratic part is at most $16\,h_G^2\cdot C_H^2\cdot\check{s}^4\cdot K_0\cdot\mathsf{p}\cdot q_x^2\cdot\exp\left(-\vartheta\cdot\left(d_{\mathrm{sep}}-4M-4\right)\right)$. As $N_h$ is subadditive, the left side of \eqref{eq:four-jets-V3} does not exceed the sum of this bound and of the bound for the activity part and $V_0$; that is, it is at most
\[
\left(16\,h_G^2+128\right)C_H^2\cdot\check{s}^4\cdot K_0\cdot\mathsf{p}\cdot q_x^2\cdot\exp\left(-\vartheta\cdot\left(d_{\mathrm{sep}}-4M-4\right)\right)\le q_x^2\cdot C_V\cdot h_G^2\cdot\mathsf{p}\cdot\exp\left(-\vartheta\cdot\left(d_{\mathrm{sep}}-4M-6\right)\right),
\]
because $h_G\ge1$ and $C_V=10^3\,C_H^2\cdot\check{s}^4\cdot K_0$. This proves \eqref{eq:four-jets-V3}.
\end{proof}

\subsection{Small fields: the logarithm and its cell components}\label{sec:small-fields-logarithm}

\subsubsection{Construction of the logarithm}\label{sec:construction-logarithm}

Throughout this subsection $\theta\in\mathbb{R}^3$ and $A\in V_{\mathrm{t}}\left(T^0_{*}\right)$ with $\chi_{*}(A)=1$ are arbitrary; the real function $\boldsymbol{\lambda}$ on $T^0_{*}$, the field $A^g=A-d\boldsymbol{\lambda}$ and the field $\mathscr{A}^0=H_{*}A^g$ are chosen as in~\eqref{it:trajectory-fermionic-operators-SL5}, and $\mathscr{A}$, $\mathscr{T}(t)$, $\mathcal{G}$, $\mathsf{Z}$, $V_0$ and $V_{+}$ denote the objects of Section~\ref{sec:activities-gaussian-form} formed with these data.

\begin{lemma}[The terminal integral compared with a determinant]\label{lem:terminal-integral-compared}
Assume~\eqref{eq:four-jets-sm2}, \eqref{eq:four-jets-sm3} and~\eqref{eq:four-jets-sm4}, and let $t\in\mathbb{D}^{\circ}\left(t_x^{+}\right)$. Then the operator $D_{*}\left(c_{\theta}+\mathscr{C}(t)\right)$ is invertible, and the number
\[
\mathsf{f}(t):=\frac{\mathsf{Z}(t)}{\det D_{*}\left(c_{\theta}+\mathscr{C}(t)\right)}
\]
satisfies
\[
\left|\mathsf{f}(t)-1\right|\le\frac{1}{2}\,e_{*}^{\frac{1}{4}-9\varepsilon}\le\frac{1}{16}.
\]
Moreover, $t\mapsto\mathsf{f}(t)$ is an analytic function on $\mathbb{D}^{\circ}\left(t_x^{+}\right)$, and the bound implies $\mathsf{Z}(t)\ne0$.
\end{lemma}

\begin{proof}
Put $y:=e_{*}^{\frac{1}{4}-9\varepsilon}$. By Section~\ref{sec:activities-gaussian-form} the field $c_{\theta}+\mathscr{C}(t)$ belongs to $G_{\mathrm{w}}$, and $G_{\mathrm{w}}\subset\tilde{G}_{*}$ by Section~\ref{sec:dependence-stop-level}; hence $\mathsf{D}\left(c_{\theta}+\mathscr{C}(t)\right)$ is invertible by~\eqref{it:trajectory-fermionic-operators-SL2}. One has $e_{*}\le e_0<1$ by Lemma~\ref{lem:scales-stop-level}(a) and~\eqref{eq:renormalisation-group-A8}. Therefore~\eqref{it:trajectory-fermionic-operators-SL7}, which applies under~\eqref{eq:four-jets-sm2}, the first condition in~\eqref{eq:four-jets-sm4} and the inequality $\mathfrak{m}_{-}<\mathfrak{m}_{*}$ of Lemma~\ref{lem:scales-stop-level}(a) give
\[
e_{*}\cdot\left|\mathscr{C}(t)\right|_{\infty}<e_{*}^{3\varepsilon}\le e_{*}^{2\varepsilon}\le\frac{\mathfrak{m}_{-}}{48}<\frac{\mathfrak{m}_{*}}{24}.
\]
The constant field $c_{\theta}$ is real. So Lemma~\ref{lem:inverses}(b), applied with $A_r=c_{\theta}$ and $A_1=\mathscr{C}(t)$ (its hypothesis that $\mathsf{D}\left(c_{\theta}+\mathscr{C}(t)\right)$ be invertible was verified above), shows that $D_{*}\left(c_{\theta}+\mathscr{C}(t)\right)$ is invertible and that every row of the matrix
\[
\mathrm{cov}:=D_{*}\left(c_{\theta}+\mathscr{C}(t)\right)^{-1}
\]
has $\ell^1$ norm at most $4\,l_{*}^3\cdot g_1$. By the first condition in~\eqref{eq:four-jets-sm3},
\[
4\,l_{*}^3\cdot g_1\le a_2=C_0^{-1}\cdot e_{*}^{-\frac{1}{4}+\varepsilon},
\]
where $a_2$ is the number of~\eqref{eq:four-jets-Gr6}. As $l_{*}\ge1$ and $g_1=b_{-}^{-1}+4\,\mathfrak{m}_{-}^{-1}\ge1$ (Section~\ref{sec:constants-3}), it follows that $a_2\ge1$.

Let $\psi$, $\bar{\psi}$ and $\psi^{\#}$ be formed with $\mathscr{T}(t)$ as in~\eqref{it:trajectory-fermionic-operators-SL6}, and let $\mathfrak{E}\in\mathrm{Gr}[\Psi]$ be given by
\[
\mathfrak{E}:=-\left\langle\bar{\Psi},\left[D_{*}\left(\mathscr{T}(t)\right)-D_{*}\left(c_{\theta}+\mathscr{C}(t)\right)\right]\Psi\right\rangle-m_{*}\cdot\left\langle\bar{\psi},\psi\right\rangle+\sum_{X\in D_{*}}\mathscr{E}_{*}\left(X;\mathscr{A},\mathscr{C}(t),\psi^{\#}\right).
\]
This element is even, the kernels of $\mathscr{E}_{*}$ having equally many barred and unbarred arguments by~\eqref{it:trajectory-fermionic-operators-SL3}. The exponent of the integrand in~\eqref{eq:four-jets-15} equals $-\left\langle\bar{\Psi},\mathrm{cov}^{-1}\Psi\right\rangle+\mathfrak{E}$, and even elements commute; hence~\eqref{eq:four-jets-Gr3}, applied with $\operatorname{Cov}=\mathrm{cov}$, gives
\[
\mathsf{Z}(t)=\det D_{*}\left(c_{\theta}+\mathscr{C}(t)\right)\cdot\int\mathrm{d}\mu_{\mathrm{cov}}(\Psi)\,\exp(\mathfrak{E}),
\]
that is,
\[
\mathsf{f}(t)=\int\mathrm{d}\mu_{\mathrm{cov}}(\Psi)\,\exp(\mathfrak{E})=\int\mathrm{d}\mu_I(\Psi)\,\left[\exp(\mathfrak{E})\right]\left(\bar{\Psi},\mathrm{cov}\,\Psi\right),
\]
with the notation $F(\bar{\Psi},\operatorname{Cov}\Psi)$ of Section~\ref{sec:grassmann-norms-gaussian}. The first identity in~\eqref{eq:four-jets-Gr3}, applied to the constant $1$, and the normalisation $\int\mathrm{D}\Psi\,\exp\left(-\left\langle\bar{\Psi},O\Psi\right\rangle\right)=\det O$ of Section~\ref{sec:fermion-integrals-blocking} show that $\int\mathrm{d}\mu_{\mathrm{cov}}(\Psi)\,1=1$. Consequently~\eqref{eq:four-jets-Gr2}, \eqref{eq:four-jets-Gr4} with $a=a_2$, and~\eqref{eq:four-jets-Gr1} yield
\[
\left|\mathsf{f}(t)-1\right|=\left|\int\mathrm{d}\mu_I(\Psi)\,\left[\exp(\mathfrak{E})-1\right]\left(\bar{\Psi},\mathrm{cov}\,\Psi\right)\right|\le\left|\exp(\mathfrak{E})-1\right|_{a_2}\le\exp\left(\left|\mathfrak{E}\right|_{a_2}\right)-1.
\]
The three terms of $\mathfrak{E}$ are bounded as follows.
\begin{enumerate}
\item[(i)] One has $\mathscr{T}(t)=\left(c_{\theta}+\mathscr{C}(t)\right)+\mathscr{A}^0$, and by Section~\ref{sec:activities-gaussian-form} the segment joining $c_{\theta}+\mathscr{C}(t)$ to $\mathscr{T}(t)$ lies in $G_{\mathrm{w}}$. Put $B:=\mathscr{A}^0$. By~\eqref{eq:four-jets-4},
\[
\epsilon_B=e_{*}\cdot\left|\mathscr{A}^0\right|_{\infty}<\frac{1}{2}\,e_{*}^{1-\varepsilon},
\]
and $\#|P_B|\le l_{*}^3$, since $P_B$ is a set of sites of $T^0_{*}$. Corollary~\ref{cor:increments-along-segment}, applied with $A_0=c_{\theta}+\mathscr{C}(t)$, together with~\eqref{eq:four-jets-Gr5} gives
\[
\left|\left\langle\bar{\Psi},\left[D_{*}\left(\mathscr{T}(t)\right)-D_{*}\left(c_{\theta}+\mathscr{C}(t)\right)\right]\Psi\right\rangle\right|_{a_2}\le4\,a_2^2\cdot C_K\cdot\frac{1}{2}\,e_{*}^{1-\varepsilon}\cdot l_{*}^3\le2\,C_K\cdot l_{*}^3\cdot e_{*}^{\frac{1}{2}+\varepsilon},
\]
because $a_2^2=C_0^{-2}\cdot e_{*}^{-\frac{1}{2}+2\varepsilon}$ and $C_0\ge1$ (Section~\ref{sec:constants-3}).
\item[(ii)] The field $\mathscr{T}(t)$ lies in $G_{\mathrm{w}}\subset\tilde{G}_{*}$. By the first bound in~\eqref{eq:four-jets-Gr6} and by~\eqref{it:trajectory-fermionic-operators-SL1},
\[
\left|m_{*}\cdot\left\langle\bar{\psi},\psi\right\rangle\right|_{a_2}\le\frac{e_{*}^{\frac{3}{4}-8\varepsilon}}{8}\cdot4\,e_{*}^{-\frac{1}{2}+2\varepsilon}\cdot l_{*}^3=\frac{1}{2}\,l_{*}^3\cdot e_{*}^{\frac{1}{4}-6\varepsilon}.
\]
\item[(iii)] One has $\mathscr{A}\in\frac{1}{4}\tilde{R}_{*}\subset\tilde{R}_{*}$ and $\mathscr{C}(t)\in\mathfrak{X}_{*}$. By the second bound in~\eqref{eq:four-jets-Gr6} and by~\eqref{eq:four-jets-2}, for every $X\in D_{*}$,
\[
\left|\mathscr{E}_{*}\left(X;\mathscr{A},\mathscr{C}(t),\psi^{\#}\right)\right|_{a_2}\le4\,e_{*}^{\varepsilon}\cdot w^{\mathrm{E}}\cdot\exp\left(-\kappa\cdot d_M(X)\right),\qquad w^{\mathrm{E}}=e_{*}^{\frac{1}{4}-7\varepsilon}.
\]
Every polymer contains at least one of the $\left(l_{*}/M\right)^3$ $M$-cubes of the torus, and $\kappa\ge\kappa_0$ by stage~\ref{it:order-choice-3}. So the first assertion of~\eqref{it:trajectory-fermionic-operators-SL4}(v) gives
\[
\sum_{X\in D_{*}}\exp\left(-\kappa\cdot d_M(X)\right)\le\sum_{\square}\ \sum_{X\supseteq\square}\exp\left(-\kappa_0\cdot d_M(X)\right)\le\left(\frac{l_{*}}{M}\right)^3\cdot K_0,
\]
the outer sum being taken over the $M$-cubes $\square$; therefore
\[
\sum_{X\in D_{*}}\left|\mathscr{E}_{*}\left(X;\mathscr{A},\mathscr{C}(t),\psi^{\#}\right)\right|_{a_2}\le4\,e_{*}^{\varepsilon}\cdot w^{\mathrm{E}}\cdot K_0\cdot\left(\frac{l_{*}}{M}\right)^3.
\]
\end{enumerate}
By~\eqref{eq:renormalisation-group-A8} and Lemma~\ref{lem:scales-stop-level}(a), $e_{*}^{\varepsilon}\le e_0^{\varepsilon}\le\frac{1}{4}$. Since moreover $e_{*}<1$, $\frac{1}{2}+\varepsilon\ge\frac{1}{4}-7\varepsilon$ and $M\ge1$, the three bounds add up to
\[
\left|\mathfrak{E}\right|_{a_2}\le l_{*}^3\cdot\left(2\,C_K+K_0+4\right)\cdot e_{*}^{\frac{1}{4}-7\varepsilon}\le\frac{1}{4}\,e_{*}^{\frac{1}{4}-9\varepsilon}=\frac{y}{4},
\]
where the second inequality is the second condition in~\eqref{eq:four-jets-sm3}. The third condition in~\eqref{eq:four-jets-sm3} states that $y\le\frac{1}{8}$. As $\exp(s)-1\le s\cdot\exp(s)\le2s$ for $0\le s\le\frac{1}{32}$, it follows that
\[
\left|\mathsf{f}(t)-1\right|\le\exp\left(\frac{y}{4}\right)-1\le\frac{y}{2}\le\frac{1}{16}.
\]

The maps $t\mapsto c_{\theta}+\mathscr{C}(t)$ and $t\mapsto\mathscr{T}(t)$ are affine, with values in $G_{\mathrm{w}}\subset\tilde{G}_{*}$, and $\mathscr{C}(t)\in\mathfrak{X}_{*}$. By~\eqref{it:trajectory-fermionic-operators-SL2} and~\eqref{it:trajectory-fermionic-operators-SL3}, the entries of $D_{*}\left(c_{\theta}+\mathscr{C}(t)\right)$ and the coefficients of the exponent of the integrand in~\eqref{eq:four-jets-15} are therefore analytic functions of $t\in\mathbb{D}^{\circ}\left(t_x^{+}\right)$. This exponent is the sum of its part of degree zero, a number, and of a remainder without constant term. A number commutes with every element of the Grassmann algebra, so the integrand is the exponential of the part of degree zero multiplied by the exponential of the remainder; and the latter exponential is a finite sum of powers of the remainder, because an element without constant term of a Grassmann algebra with finitely many generators is nilpotent. Hence each coefficient of the integrand is the exponential of the part of degree zero of the exponent times a polynomial in the remaining coefficients of the exponent, and is an analytic function of $t$. The Berezin integral in~\eqref{eq:four-jets-15} is a linear form in the coefficients of the integrand. It follows that $\mathsf{Z}$ is analytic, and so is $\mathsf{f}$, the determinant of $D_{*}\left(c_{\theta}+\mathscr{C}(t)\right)$ being analytic without zeros on $\mathbb{D}^{\circ}\left(t_x^{+}\right)$. Finally, $\left|\mathsf{f}(t)-1\right|\le\frac{1}{16}$ implies $\mathsf{f}(t)\ne0$, so that $\mathsf{Z}(t)=\mathsf{f}(t)\cdot\det D_{*}\left(c_{\theta}+\mathscr{C}(t)\right)\ne0$.
\end{proof}

Assume~\eqref{eq:four-jets-sm2}--\eqref{eq:four-jets-sm4}. By~\eqref{eq:four-jets-4} the fields $A+c_{\theta}$ and $A^g+c_{\theta}$ lie in $\mathrm{Dom}_{*}$, and $A+c_{\theta}=\left(A^g+c_{\theta}\right)+d\boldsymbol{\lambda}$. The gauge invariance stated in~\eqref{it:trajectory-fermionic-operators-SL6} therefore gives
\[
Y^x_{*}\left(A+c_{\theta};\mathscr{C}\right)=Y^x_{*}\left(A^g+c_{\theta};\mathscr{C}\right)\qquad\left(\mathscr{C}\in\mathfrak{X}_{*}\right);
\]
indeed, every $\mathscr{C}\in\mathfrak{X}_{*}$ is of the form $a^{[k^{*}]}$ with $a\in C^1\otimes\mathbb{C}$, and for $0\le j\le k^{*}$ the scaling relations in the proof of Lemma~\ref{lem:scaling-external-connection} and the definition of $\mathfrak{X}_{*}$ give
\[
\left|a^{[j]}\right|_{\infty}\le\left|\mathscr{C}\right|_{\infty}<e_{*}^{-1+2\varepsilon},\qquad\left|\partial a^{[j]}\right|_{\infty}\le\left|\partial\mathscr{C}\right|_{\infty}<e_{*}^{-1+3\varepsilon},
\]
while $e_j\le e_{*}<1$, by Section~\ref{sec:running-quantities-tori}, Lemma~\ref{lem:scales-stop-level}(a) and~\eqref{eq:renormalisation-group-A8}, so that the two thresholds of $\mathfrak{X}_j$ in Definition~\ref{def:domains} are not smaller than those of $\mathfrak{X}_{*}$; hence $a^{[j]}\in\mathfrak{X}_j$ for $0\le j\le k^{*}$. By Section~\ref{sec:activities-gaussian-form}, for $t\in\mathbb{D}^{\circ}\left(t_x^{+}\right)$,
\[
Y^x_{*}\left(A^g+c_{\theta};\mathscr{C}(t)\right)=\det D_{*}(0)^{-1}\cdot\exp\left\{-\boldsymbol{\varepsilon}^t_{*}\cdot l_{*}^3+\sum_{X\in D_{*}}\mathfrak{P}_{*}\left(X,\mathscr{C}(t)\right)\right\}\cdot\mathsf{Z}(t).
\]
Here $\mathscr{C}(0)=0$, and $\mathfrak{P}_{*}(X,0)=0$ for every $X\in D_{*}$ by Lemma~\ref{lem:analyticity-locality-decay}; moreover $\mathsf{Z}$ has no zero on $\mathbb{D}^{\circ}\left(t_x^{+}\right)$ by Lemma~\ref{lem:terminal-integral-compared}. It follows that, for all $t\in\mathbb{D}^{\circ}\left(t_x^{+}\right)$, all $\theta\in\mathbb{R}^3$ and all $A\in V_{\mathrm{t}}\left(T^0_{*}\right)$ with $\chi_{*}(A)=1$,
\begin{align}
&Y^x_{*}\left(A+c_{\theta};0\right)\ne0,\notag\\
&\mathrm{Rt}(t;\theta,A):=\frac{Y^x_{*}\left(A+c_{\theta};\mathscr{C}(t)\right)}{Y^x_{*}\left(A+c_{\theta};0\right)}=\exp\left\{\sum_{X\in D_{*}}\mathfrak{P}_{*}\left(X,\mathscr{C}(t)\right)\right\}\cdot\frac{\mathsf{Z}(t)}{\mathsf{Z}(0)}\ne0.\label{eq:four-jets-20}
\end{align}

Only the field $A+c_{\theta}$ enters the definition of $\mathrm{Rt}(t;\theta,A)$ in~\eqref{eq:four-jets-20}; neither $\boldsymbol{\lambda}$ nor $A^g$ plays a part in it. Let $\mathscr{C}\in\mathfrak{X}_{*}$, and apply~\eqref{eq:four-jets-5} at the point $A+c_{\theta}\in\mathrm{Dom}_{*}$. As $H_{*}$ is linear and $H_{*}c_{\theta}=c_{\theta}$ by~\eqref{eq:four-jets-4}, one has $H_{*}\left(A+c_{\theta}\right)=H_{*}A+c_{\theta}$; the kernels of $\mathscr{E}_{*}$ in~\eqref{eq:four-jets-5} are therefore taken at the pair $\left(H_{*}A+c_{\theta},\mathscr{C}\right)$, and the operators in~\eqref{eq:four-jets-5} are formed with the connection $H_{*}A+c_{\theta}+\mathscr{C}$. Here $H_{*}A+c_{\theta}\in\frac{1}{4}\,\tilde{R}_{*}\subset\tilde{R}_{*}$ by the definition of $\mathrm{Dom}_{*}$ (Definition~\ref{def:domains}), and $H_{*}A+c_{\theta}+\mathscr{C}\in\tilde{R}_{*}+\mathfrak{X}_{*}\subset\tilde{G}_{*}$ by~\eqref{it:trajectory-fermionic-operators-SL2}. So $Y^x_{*}\left(A+c_{\theta};\mathscr{C}\right)$ is the product of the number $\det D_{*}(0)^{-1}$, of the factor $\exp\left\{-\boldsymbol{\varepsilon}^t_{*}\cdot l_{*}^3+\sum_{X\in D_{*}}\mathfrak{P}_{*}(X,\mathscr{C})\right\}$, which depends analytically on $\mathscr{C}\in\mathfrak{X}_{*}$ by Lemma~\ref{lem:analyticity-locality-decay}, and of the Berezin integral of~\eqref{eq:four-jets-5}. The argument at the end of the proof of Lemma~\ref{lem:terminal-integral-compared} shows that this integral is the exponential of the part of degree zero of its exponent times a polynomial in the remaining coefficients of the exponent; it is thus an entire function of the coefficients of the exponent. By~\eqref{it:trajectory-fermionic-operators-SL3}, at fixed arguments every kernel of $\mathscr{E}_{*}$ is an analytic function on $\tilde{R}_{*}\times\mathfrak{X}_{*}$ of the pair at which it is taken, and by~\eqref{it:trajectory-fermionic-operators-SL2} the operators formed with a connection are analytic functions of the connection on $\tilde{G}_{*}$; both sets are open by Lemma~\ref{lem:openness}. Hence the coefficients of the exponent depend analytically on the pair $\left(H_{*}A+c_{\theta},\mathscr{C}\right)$, through the kernels of $\mathscr{E}_{*}$, and on the connection $H_{*}A+c_{\theta}+\mathscr{C}$, through the operators. For $\mathscr{C}=\mathscr{C}(t)$ the three fields $H_{*}A+c_{\theta}$, $\mathscr{C}(t)$ and $H_{*}A+c_{\theta}+\mathscr{C}(t)$ are affine functions of $(\theta,A,t)$. Therefore the numerator $Y^x_{*}\left(A+c_{\theta};\mathscr{C}(t)\right)$ of $\mathrm{Rt}$ depends continuously on $(t,\theta,A)$ and, for fixed $(\theta,A)$, analytically on $t\in\mathbb{D}^{\circ}\left(t_x^{+}\right)$. Its value at $t=0$ is the denominator $Y^x_{*}\left(A+c_{\theta};0\right)$, so that the denominator is continuous in $(\theta,A)$; it does not vanish. Hence the quotient $\mathrm{Rt}$ has the same two properties as the numerator.

For $t\in\mathbb{D}^{\circ}\left(t_x^{+}\right)$ put
\begin{equation}\label{eq:four-jets-21}
\mathrm{Lg}(t;\theta,A):=\int_0^1\mathrm{d}u\,\sum_{c\in\mathrm{Cel}}t_c\cdot\frac{\left(\partial_c\mathrm{Rt}\right)(u\cdot t;\theta,A)}{\mathrm{Rt}(u\cdot t;\theta,A)}.
\end{equation}
The polydisk $\mathbb{D}^{\circ}\left(t_x^{+}\right)$ is convex, the function $\mathrm{Rt}(\cdot;\theta,A)$ is analytic and without zeros on it, and its partial derivatives are continuous in $(t,\theta,A)$ by Cauchy's integral formula. It follows that $\mathrm{Lg}(\cdot;\theta,A)$ is an analytic function on $\mathbb{D}^{\circ}\left(t_x^{+}\right)$ that vanishes at $t=0$, and that $\mathrm{Lg}$ depends continuously on $(t,\theta,A)$. For fixed $t$ and $u\in[0,1]$,
\[
\frac{\mathrm{d}}{\mathrm{d}u}\,\mathrm{Lg}(u\cdot t;\theta,A)=\sum_{c\in\mathrm{Cel}}t_c\cdot\frac{\left(\partial_c\mathrm{Rt}\right)(u\cdot t;\theta,A)}{\mathrm{Rt}(u\cdot t;\theta,A)},
\]
so that the function $u\mapsto\mathrm{Rt}(u\cdot t;\theta,A)\cdot\exp\left(-\mathrm{Lg}(u\cdot t;\theta,A)\right)$ has derivative zero on $[0,1]$; as $\mathrm{Rt}(0;\theta,A)=1$, one obtains $\exp\mathrm{Lg}=\mathrm{Rt}$. Moreover $\mathrm{Lg}(\cdot;\theta,A)$ is the only continuous function $F$ on $\mathbb{D}^{\circ}\left(t_x^{+}\right)$ with $F(0)=0$ and $\exp F=\mathrm{Rt}(\cdot;\theta,A)$: for such a function $F$, the function $\frac{1}{2\pi i}\left(F-\mathrm{Lg}(\cdot;\theta,A)\right)$ is continuous with integer values on a connected set and vanishes at $0$.

Assume in addition~\eqref{eq:four-jets-sm1}. By~\eqref{eq:four-jets-V1} and the third condition in~\eqref{eq:four-jets-sm1},
\[
\boldsymbol{\alpha}\cdot\rho_h\left(V_{+}(t);w_{\boldsymbol{\delta}}\right)\le\boldsymbol{\alpha}\cdot\bar{r}\le\frac{1}{2}\qquad\left(t\in\mathbb{D}(t_x)\right).
\]
The field $\mathscr{A}$ is real. By Lemma~\ref{lem:inverses}(a), the matrix $\mathcal{G}=D_{*}(\mathscr{A})^{-1}$ on $\mathcal{I}=T^0_{*}\times\{1,\dots,4\}$, with the pseudo-metric of Section~\ref{sec:activities-gaussian-form}, satisfies the requirements of Section~\ref{sec:weighted-tree-expansion} with
\[
\left(C,g,h_G,\boldsymbol{\alpha}_{\vartheta}\right)=\left(\mathcal{G},g,h_G,\boldsymbol{\alpha}\right).
\]
Since $w_{\boldsymbol{\delta}}\ge1$ and $\boldsymbol{\alpha}_0\le\boldsymbol{\alpha}$, also $\boldsymbol{\alpha}_0\cdot\rho_h\left(V_{+}(t);1\right)\le\frac{1}{2}$ for $t\in\mathbb{D}(t_x)$. For $t\in\mathbb{D}^{\circ}(t_x)$ let $W(t)$ be defined by the first line of the following display, the functionals $\tilde{E}_n$ being formed with $C=\mathcal{G}$; the series converges absolutely by Lemma~\ref{lem:tree-expansion}(a). The second line follows from~\eqref{eq:four-jets-17} and Lemma~\ref{lem:tree-expansion}(a), and the third line holds as well:
\begin{align}
W(t)&:=V_0(t)+\sum_{n\ge1}\frac{1}{n!}\,\tilde{E}_n\left(V_{+}(t),\dots,V_{+}(t)\right),\notag\\
\mathsf{Z}(t)&=\det D_{*}(\mathscr{A})\cdot\exp W(t),\notag\\
\mathrm{Lg}(t)&=\sum_{X\in D_{*}}\mathfrak{P}_{*}\left(X,\mathscr{C}(t)\right)+W(t)-W(0).\label{eq:four-jets-22}
\end{align}
Here $\mathrm{Lg}(t)$ stands for $\mathrm{Lg}(t;\theta,A)$. To prove the third line, note that its right side is analytic on $\mathbb{D}^{\circ}(t_x)$: the sum over $X$ by Lemma~\ref{lem:analyticity-locality-decay}, the function $V_0$ by Section~\ref{sec:activities-gaussian-form}, and the series by Lemma~\ref{lem:tree-expansion}(c), applied on the open set $U=\mathbb{D}^{\circ}(t_x)$. The left side is analytic there as well. Both sides vanish at $t=0$. By~\eqref{eq:four-jets-20} and the second line, in which $\det D_{*}(\mathscr{A})\ne0$ cancels, both sides have the exponential $\mathrm{Rt}(t;\theta,A)$. Two continuous functions on the connected set $\mathbb{D}^{\circ}(t_x)$ that have the same exponential and agree at one point coincide.

\subsubsection{The determinant of the free Wilson operator}\label{sec:determinant-free-wilson}

For $0\le j\le k^{*}$ and $X\in D_j$, the functions $\mathfrak{P}_j(X,\cdot)$ are those defined by the recursion in~\eqref{it:trajectory-fermionic-operators-SL4}, which starts from $\mathfrak{P}_0=0$.

\begin{lemma}[Analyticity, locality and decay of the determinant activities]\label{lem:analyticity-locality-decay}
Let $0\le j\le k^{*}$ and $X\in D_j$. Then $\mathfrak{P}_j(X,\cdot)$ is an analytic function on $\mathfrak{X}_j$ with $\mathfrak{P}_j(X,0)=0$, and $\mathfrak{P}_j(X,\mathscr{C})=\mathfrak{P}_j\left(X,\widetilde{\mathscr{C}}\right)$ for any two elements $\mathscr{C},\widetilde{\mathscr{C}}$ of $\mathfrak{X}_j$ that take the same value on every bond meeting $X$. If moreover $d_M(X)\ge1$, then
\[
\left|\mathfrak{P}_j(X,\mathscr{C})\right|\le2\,C_d\cdot M^3\cdot\exp\left(-\kappa\cdot d_M(X)\right)\qquad\left(\mathscr{C}\in\mathfrak{X}_j\right).
\]
\end{lemma}

\begin{proof}
Let $0\le j<k^{*}$. For a bond function $\mathscr{C}$ on $T_{j+1}$, the scaled function $\mathscr{C}_L$ of Section~\ref{sec:local-parts-adjustment} is the bond function on $T_j$ with $\mathscr{C}_L(b)=L^{-\frac{1}{2}}\cdot\mathscr{C}(b/L)$. As in the proof of Lemma~\ref{lem:scaling-external-connection},
\[
\left|\mathscr{C}_L\right|_{\infty}=L^{-\frac{1}{2}}\cdot\left|\mathscr{C}\right|_{\infty},\qquad\left|\partial\mathscr{C}_L\right|_{\infty}=L^{-\frac{3}{2}}\cdot\left|\partial\mathscr{C}\right|_{\infty}.
\]
By Section~\ref{sec:running-quantities-tori}, Lemma~\ref{lem:scales-stop-level}(a) and~\eqref{eq:renormalisation-group-A8}, one has $e_j<e_{j+1}\le e_{*}<1$, so that the two thresholds in the definition of $\mathfrak{X}_j$ (Definition~\ref{def:domains}) are larger than those of $\mathfrak{X}_{j+1}$. Hence the linear map $\mathscr{C}\mapsto\mathscr{C}_L$ sends $\mathfrak{X}_{j+1}$ into $\mathfrak{X}_j$, and $\mathfrak{X}_j\subset\tilde{G}_j$ by~\eqref{it:trajectory-fermionic-operators-SL2}.

The first three assertions are proved by induction on $j$. For $j=0$ they hold because $\mathfrak{P}_0=0$. Assume them at a level $j<k^{*}$, and let $Y\in D_{j+1}$. By~\eqref{it:trajectory-fermionic-operators-SL4},
\[
\mathfrak{P}_{j+1}(Y,\mathscr{C})=\sum_{X\in D_j:\ \bar{X}^L=LY}\left(\mathfrak{P}_j\left(X,\mathscr{C}_L\right)+E_j^{\mathrm{det}}\left(X,\mathscr{C}_L\right)\right),
\]
where $E_j^{\mathrm{det}}(X,A)=E^d_j(X,A)-E^d_j(X,0)$ and $E^d_j(X,\cdot)$ is analytic on $\tilde{G}_j$. By the preceding paragraph and the induction hypothesis, every term of the sum is an analytic function of $\mathscr{C}\in\mathfrak{X}_{j+1}$; it vanishes at $\mathscr{C}=0$, because then $\mathscr{C}_L=0$, and $E_j^{\mathrm{det}}(X,0)=0$. Let $\mathscr{C},\widetilde{\mathscr{C}}\in\mathfrak{X}_{j+1}$ coincide on the bonds that meet $Y$, and let $X\in D_j$ with $\bar{X}^L=LY$; then $X\subset LY$. If a bond $b$ of $T_j$ meets $X$, it has an end point in $LY$; hence the bond $b/L$ of $T_{j+1}$ meets $Y$, and
\[
\mathscr{C}_L(b)=L^{-\frac{1}{2}}\cdot\mathscr{C}(b/L)=L^{-\frac{1}{2}}\cdot\widetilde{\mathscr{C}}(b/L)=\widetilde{\mathscr{C}}_L(b).
\]
So $\mathscr{C}_L$ and $\widetilde{\mathscr{C}}_L$ coincide on the bonds that meet $X$, and the terms of $X$ in the two sums are equal, by the induction hypothesis and by~\eqref{it:trajectory-fermionic-operators-SL4}(ii). This proves the first three assertions at the level $j+1$.

For $0\le j\le k^{*}$ let
\[
q_j:=\sup\left\{\exp\left(\kappa\cdot d_M(X)\right)\cdot\left|\mathfrak{P}_j(X,\mathscr{C})\right|:\ X\in D_j,\ d_M(X)\ge1,\ \mathscr{C}\in\mathfrak{X}_j\right\}\in[0,\infty];
\]
then $q_0=0$. Let $j<k^{*}$, let $Y\in D_{j+1}$ with $d_M(Y)\ge1$, and let $\mathscr{C}\in\mathfrak{X}_{j+1}$. By~\eqref{it:trajectory-fermionic-operators-SL4}(v), every $X\in D_j$ with $\bar{X}^L=LY$ satisfies $d_M(X)\ge L\cdot d_M(Y)\ge1$, so that $\left|\mathfrak{P}_j\left(X,\mathscr{C}_L\right)\right|\le q_j\cdot\exp\left(-\kappa\cdot d_M(X)\right)$. By~\eqref{it:trajectory-fermionic-operators-SL4}(i), applied to $E^d_j\left(X,\mathscr{C}_L\right)$ and to $E^d_j(X,0)$,
\[
\left|E_j^{\mathrm{det}}\left(X,\mathscr{C}_L\right)\right|\le2\,C_d\cdot M^3\cdot\exp\left(-(\kappa+1)\cdot d_M(X)\right)\le2\,C_d\cdot M^3\cdot\exp\left(-\kappa\cdot d_M(X)\right).
\]
The sum over $X$ is bounded by the last assertion of~\eqref{it:trajectory-fermionic-operators-SL4}(v):
\[
\left|\mathfrak{P}_{j+1}(Y,\mathscr{C})\right|\le\left(q_j+2\,C_d\cdot M^3\right)\cdot c_{\mathrm{rb}}\cdot L^3\cdot\exp\left(-L\cdot\left(\kappa-\kappa_0-1\right)\cdot d_M(Y)\right).
\]
By stage~\ref{it:order-choice-4}, in the form recorded in Section~\ref{sec:smallness-conditions}, one has $L\cdot\left(\kappa-\kappa_0-1\right)\ge\kappa$ and
\[
c_{\mathrm{rb}}\cdot L^3\cdot\exp\left(-\left(L\cdot\left(\kappa-\kappa_0-1\right)-\kappa\right)\right)\le\frac{1}{2}.
\]
Since $d_M(Y)\ge1$, the product of $L\cdot\left(\kappa-\kappa_0-1\right)-\kappa\ge0$ and $d_M(Y)-1\ge0$ is non-negative, that is,
\[
L\cdot\left(\kappa-\kappa_0-1\right)\cdot d_M(Y)\ge\left(L\cdot\left(\kappa-\kappa_0-1\right)-\kappa\right)+\kappa\cdot d_M(Y).
\]
Therefore
\begin{align*}
\left|\mathfrak{P}_{j+1}(Y,\mathscr{C})\right|&\le\left(q_j+2\,C_d\cdot M^3\right)\cdot c_{\mathrm{rb}}\cdot L^3\cdot\exp\left(-\left(L\cdot\left(\kappa-\kappa_0-1\right)-\kappa\right)\right)\\
&\qquad{}\cdot\exp\left(-\kappa\cdot d_M(Y)\right)\\
&\le\frac{1}{2}\left(q_j+2\,C_d\cdot M^3\right)\cdot\exp\left(-\kappa\cdot d_M(Y)\right).
\end{align*}
Thus $q_{j+1}\le\frac{1}{2}\left(q_j+2\,C_d\cdot M^3\right)$. Starting from $q_0=0$, induction gives $q_j\le2\,C_d\cdot M^3$ for $0\le j\le k^{*}$, which is the asserted bound.
\end{proof}

\begin{lemma}[The free Wilson ratio on a disc]\label{lem:free-wilson-ratio}
Assume conditions~\eqref{eq:four-jets-sm2} and~\eqref{eq:four-jets-sm4}, and let $c\in\mathrm{Cel}$. For $t\in\mathbb{C}$ with $|t|<t_x^{+}$ and for $\theta\in\mathbb{R}^3$ put
\begin{align*}
g_D^{\theta}(t)&:=\operatorname{Tr}\log\left[1+D_{*}(c_{\theta})^{-1}\cdot\left(D_{*}\left(c_{\theta}+t\cdot\hat{c}_c\right)-D_{*}(c_{\theta})\right)\right],\\
\mathrm{lgN}(t)&:=\sum_{X\in D_{*}}\mathfrak{P}_{*}\left(X,t\cdot\hat{c}_c\right)+g_D^0(t);
\end{align*}
here the operator $D_{*}(c_{\theta})$ is invertible by Lemma~\ref{lem:inverses}(a), the bond function $c_{\theta}$ being real.
\begin{enumerate}
\item[(a)] For $|t|<t_x^{+}$ one has $N_0(t)=\exp\mathrm{lgN}(t)$.
\item[(b)] For $|t|<t_x^{+}$ one has
\[
\left|\sum_{X\in D_{*}}\mathfrak{P}_{*}\left(X,t\cdot\hat{c}_c\right)\right|\le C_{\mathrm{fr}}\cdot e_{*}^{-2\varepsilon}\cdot|t|^2+B_{\mathrm{wr}}.
\]
\item[(c)] For $|t|<t_x^{+}$ and $\theta\in\mathbb{R}^3$ the logarithm in the definition of $g_D^{\theta}(t)$ is defined by its power series, $\left|g_D^{\theta}(t)\right|\le B_1$, and
\[
\exp g_D^{\theta}(t)=\frac{\det D_{*}\left(c_{\theta}+t\cdot\hat{c}_c\right)}{\det D_{*}(c_{\theta})}.
\]
\item[(d)] For $|t|\le t_x$ one has $\left|\mathrm{lgN}(t)\right|\le B_0\le\frac{1}{4}$. On the disc $|t|\le t_x$ the function $\mathrm{lgN}$ is analytic and even, it takes real values at real points, and $\mathrm{lgN}(0)=0$. The function $\mathrm{lgN}$ does not depend on the cell $c$, whereas $g_D^{\theta}$ is formed with the cell $c$.
\item[(e)] Let $\phi_0(t)=\sum_{i=0}^4n_it^i$ be the Taylor polynomial of degree four of $N_0$ at $0$, and let $c_m:=\phi_0^{(m)}(0)$. Then $n_0=1$, $n_1=n_3=0$, the coefficients $n_i$ are real, $|n_i|\le2B_0\cdot t_x^{-i}$ for $1\le i\le4$, and
\begin{multline}\label{eq:four-jets-23}
|c_2|\le\frac{\mathbf{r}^{-2}}{16},\qquad|c_4|\le\frac{\mathbf{r}^{-4}}{16},\qquad\|\phi_0\|_{\mathbf{r}}\le2,\\
\|\phi_0-1\|_{\mathbf{r}}\le\frac{1}{6},\qquad c_4-3c_2^2=b_{\eta,R^{*}}(0),
\end{multline}
where $\|\cdot\|_{\mathbf{r}}$ is the norm of Section~\ref{sec:ring-jets}, applied to polynomials in the single variable $t$.
\end{enumerate}
\end{lemma}

\begin{proof}
By~\eqref{eq:four-jets-sm2} the placement is admitted, so that~\eqref{it:trajectory-fermionic-operators-SL7} and Lemma~\ref{lem:properties-probes} apply. In the proofs of (c), (a) and (b) let $t\in\mathbb{C}$ with $|t|<t_x^{+}$ and $\theta\in\mathbb{R}^3$ be fixed, and put $B:=t\cdot\hat{c}_c$. The point $t\cdot e_c$ lies in $\mathbb{D}^{\circ}(t_x^{+})$, and $B=\mathscr{C}(t\cdot e_c)=a(t\cdot e_c)^{[k^{*}]}$ by the definitions of Section~\ref{sec:probes-mirror-copies}.

\emph{(c)} For $0\le u\le1$ the point $u\cdot t\cdot e_c$ lies in $\mathbb{D}^{\circ}(t_x^{+})$ as well; hence $u\cdot B=\mathscr{C}(u\cdot t\cdot e_c)\in\mathfrak{X}_{*}$ by~\eqref{it:trajectory-fermionic-operators-SL7}. The same clause gives, with the notation of Section~\ref{sec:dependence-stop-level},
\[
\epsilon_B=e_{*}\cdot|B|_{\infty}=|t|\cdot\epsilon_c<e_{*}^{3\varepsilon}\le e_{*}^{2\varepsilon},
\]
the last inequality because $e_{*}<1$ (Lemma~\ref{lem:scales-stop-level}(a) and condition~\eqref{eq:renormalisation-group-A8}). The set $P_B$ is contained in $P_c\subset P_c^{+}$, so that $\#|P_B|\le\mathsf{p}$ by Lemma~\ref{lem:properties-probes}(e). The constant bond function $c_{\theta}$ is real, and it lies in $\tilde{R}_{*}$, because $0\in\tilde{R}_{*}$ and $\tilde{R}_{*}$ is invariant under the addition of real constants by~\eqref{it:trajectory-fermionic-operators-SL2}. By the definition of $G_{\mathrm{w}}$ at the beginning of Section~\ref{sec:dependence-stop-level}, the segment $\left\{c_{\theta}+u\cdot B:\ 0\le u\le1\right\}$ is contained in $G_{\mathrm{w}}\subset\tilde{G}_{*}$, on which $D_{*}$ is defined by~\eqref{it:trajectory-fermionic-operators-SL2}. Corollary~\ref{cor:increments-along-segment}, applied with $A_0=c_{\theta}$, shows that the operator
\[
\mathfrak{K}:=D_{*}\left(c_{\theta}+t\cdot\hat{c}_c\right)-D_{*}(c_{\theta})
\]
on the spinor functions on $T^0_{*}$ satisfies
\[
\|\mathfrak{K}\|_{\frac{\gamma_F}{2}}\le C_K\cdot\epsilon_B,\qquad\sum_{y,y'\in T^0_{*}}|\mathfrak{K}|(y,y')\le C_K\cdot\epsilon_B\cdot\#|P_B|.
\]
As $\|\cdot\|_0\le\|\cdot\|_{\frac{\gamma_F}{2}}$, the bounds~\eqref{eq:four-jets-Maj2} for the operator norm and for the trace norm give
\[
|\mathfrak{K}|_{\ell^2}\le2C_K\cdot e_{*}^{2\varepsilon},\qquad|\mathfrak{K}|_1\le4C_K\cdot e_{*}^{2\varepsilon}\cdot\mathsf{p}.
\]
By Lemma~\ref{lem:inverses}(a), applied to the real bond function $c_{\theta}$, the operator $D_{*}(c_{\theta})$ is invertible and $\left|D_{*}(c_{\theta})^{-1}\right|_{\ell^2}\le g^2\le g_1$; the second inequality follows from the definitions of $g$ and $g_1$ in Section~\ref{sec:constants-3}. Put
\[
\mathfrak{Y}:=D_{*}(c_{\theta})^{-1}\cdot\mathfrak{K}.
\]
Since $C_K\le C_K'$, the second condition of~\eqref{eq:four-jets-sm4} gives
\[
|\mathfrak{Y}|_{\ell^2}\le2g_1\cdot C_K\cdot e_{*}^{2\varepsilon}\le2g_1\cdot C_K'\cdot e_{*}^{2\varepsilon}\le\frac{1}{2}.
\]
Consequently the power series
\[
\log\left(1+\mathfrak{Y}\right)=\sum_{n\ge1}\frac{(-1)^{n+1}}{n}\cdot\mathfrak{Y}^n
\]
converges in the operator norm, and $g_D^{\theta}(t)=\operatorname{Tr}\log\left(1+\mathfrak{Y}\right)$ is defined by it. For matrices $Y$ and $Y'$ one has $|\operatorname{Tr}Y|\le|Y|_1$ and $|Y\cdot Y'|_1\le|Y|_{\ell^2}\cdot|Y'|_1$. Hence $\left|\operatorname{Tr}\mathfrak{Y}^n\right|\le|\mathfrak{Y}|_{\ell^2}^{n-1}\cdot|\mathfrak{Y}|_1$ for $n\ge1$, and $|\mathfrak{Y}|_1\le\left|D_{*}(c_{\theta})^{-1}\right|_{\ell^2}\cdot|\mathfrak{K}|_1$. Since $\sum_{n\ge1}\frac{2^{1-n}}{n}=2\log2\le2$, it follows that
\[
\left|g_D^{\theta}(t)\right|\le\sum_{n\ge1}\frac{|\mathfrak{Y}|_1\cdot|\mathfrak{Y}|_{\ell^2}^{n-1}}{n}\le2|\mathfrak{Y}|_1\le8g_1\cdot C_K\cdot\mathsf{p}\cdot e_{*}^{2\varepsilon}\le B_1,
\]
by the definition of $B_1$ in Section~\ref{sec:constants-3}. Finally, $\det(1+Y)=\exp\operatorname{Tr}\log(1+Y)$ for every matrix $Y$ with $|Y|_{\ell^2}<1$, and $1+\mathfrak{Y}=D_{*}(c_{\theta})^{-1}\cdot D_{*}\left(c_{\theta}+t\cdot\hat{c}_c\right)$. Therefore
\[
\exp g_D^{\theta}(t)=\det\left(1+\mathfrak{Y}\right)=\frac{\det D_{*}\left(c_{\theta}+t\cdot\hat{c}_c\right)}{\det D_{*}(c_{\theta})}.
\]

\emph{(a)} Let $a:=\frac{t}{e}\cdot\hat{\mathfrak{a}}_c=a(t\cdot e_c)$. By~\eqref{it:trajectory-fermionic-operators-SL7} one has $a^{[j]}\in\mathfrak{X}_j$ for every $j\le k^{*}$, and $a^{[k^{*}]}=B=t\cdot\hat{c}_c$. The first identity of~\eqref{eq:four-jets-3} therefore gives
\[
\frac{\det\left(D_a+\bar{m}\cdot\eta\right)}{\det\left(D_0+\bar{m}\cdot\eta\right)}=\exp\left\{\sum_{X\in D_{*}}\mathfrak{P}_{*}\left(X,t\cdot\hat{c}_c\right)\right\}\cdot\frac{\det D_{*}\left(t\cdot\hat{c}_c\right)}{\det D_{*}(0)}.
\]
The left side is $N_0(t)$, by the definition of $N_0$ in Section~\ref{sec:averaged-law-partition} and because $N_0$ does not change if $\hat{\mathfrak{a}}$ is replaced by $\hat{\mathfrak{a}}_c$ (Lemma~\ref{lem:signs-symmetries-wilson}(b)). For $\theta=0$ the bond function $c_{\theta}$ vanishes, and (c) shows that the last quotient on the right equals $\exp g_D^0(t)$. Hence $N_0(t)=\exp\mathrm{lgN}(t)$.

\emph{(b)} Let $a$ be as in the proof of (a). For $0\le j\le k^{*}$ let
\[
\boldsymbol{\chi}_j:=\frac{1}{e}\cdot\hat{\mathfrak{a}}_c^{[j]},
\]
a real bond function on $T_j$. Then $\boldsymbol{\chi}_{k^{*}}=\hat{c}_c$ and $t\cdot\boldsymbol{\chi}_j=a^{[j]}\in\mathfrak{X}_j$. By the second identity of~\eqref{eq:four-jets-3},
\[
\sum_{X\in D_{*}}\mathfrak{P}_{*}\left(X,t\cdot\hat{c}_c\right)=\sum_{j<k^{*}}\ \sum_{X\in D_j}E^{\mathrm{det}}_j\left(X,t\cdot\boldsymbol{\chi}_j\right).
\]
The support of $\boldsymbol{\chi}_j$ is the set of the bonds $b$ of $T_j$ with $\boldsymbol{\chi}_j(b)\ne0$, and a bond meets a set of sites if one of its end points lies in that set (Section~\ref{sec:lattice-geometry}). Let $X\in D_j$ be a polymer that is met by no bond of the support of $\boldsymbol{\chi}_j$. Then $t\cdot\boldsymbol{\chi}_j$ and $0$, which lie in $\tilde{G}_j$ by~\eqref{it:trajectory-fermionic-operators-SL2}, coincide on the bonds that meet $X$, and~\eqref{it:trajectory-fermionic-operators-SL4}(ii) gives
\[
E^{\mathrm{det}}_j\left(X,t\cdot\boldsymbol{\chi}_j\right)=E^d_j\left(X,t\cdot\boldsymbol{\chi}_j\right)-E^d_j(X,0)=0.
\]

Now let $j=k^{*}-i$ with $1\le i\le k^{*}$, let $X\in D_j$ be met by a bond of the support of $\boldsymbol{\chi}_j$, and put
\[
\mathsf{n}:=d_M(X),\qquad s_X:=M\cdot(\mathsf{n}+3).
\]
The scaling relations in the proof of Lemma~\ref{lem:scaling-external-connection}, applied to the bond function $\frac{1}{e}\cdot\hat{\mathfrak{a}}_c$, give $|\partial\boldsymbol{\chi}_j|_{\infty}=L^{-\frac{3i}{2}}\cdot|\partial\hat{c}_c|_{\infty}$, with the norm $|\partial\cdot|_{\infty}$ of Section~\ref{sec:domains-bond-functions}. The number $\left|\partial\hat{\mathfrak{a}}_c^{[k^{*}]}\right|_{\infty}$ is one of the three norms of $\hat{\mathfrak{a}}_c^{[k^{*}]}$, which are those of $\hat{\mathfrak{a}}^{[k^{*}]}$ by Lemma~\ref{lem:properties-probes}(d), and their sum is $|\hat{\mathfrak{a}}|_{\mathrm{ext}}\le C_{\beta}$ (Section~\ref{sec:probes-mirror-copies}). Hence $|\partial\hat{c}_c|_{\infty}\le\frac{C_{\beta}}{e}$. As $e_j=e\cdot a_j^{\frac{1}{2}}$ and $L^{-i}=\frac{a_j}{a_{*}}$ (Section~\ref{sec:running-quantities-tori}), it follows that
\begin{equation}\label{eq:four-jets-24}
e_j\cdot|\partial\boldsymbol{\chi}_j|_{\infty}\le a_j^2\cdot a_{*}^{-\frac{3}{2}}\cdot C_{\beta}.
\end{equation}

\emph{Case 1: $2s_X<l_j$.} Let $x_0$ be a site of $X$, and put $r_{\mathrm{b}}:=M\cdot(\mathsf{n}+2)+1$. As $M\ge L^2\ge2$ (stage~\ref{it:order-choice-6}), one has $r_{\mathrm{b}}\le s_X-1$. By~\eqref{it:trajectory-fermionic-operators-SL4}(v) two points of $X$ have distance at most $M\cdot(\mathsf{n}+2)$; a bond that meets $X$ has an end point in $X$, and its two end points have distance $\eta_j\le1$. Therefore all end points of the bonds that meet $X$ lie in the set
\[
\bar{K}:=\left\{x\in T_j:\ d(x,x_0)\le r_{\mathrm{b}}\right\}.
\]
Since $2r_{\mathrm{b}}+1<2s_X<l_j$, this set does not wrap around the torus. More precisely, for $x\in\bar{K}$ the difference $x-x_0$ has exactly one representative $z\in\mathbb{L}_j$ with $|z|\le r_{\mathrm{b}}$, and from now on $x-x_0$ denotes this representative. The map $x\mapsto x-x_0$ is a bijection of $\bar{K}$ onto the cube $\left\{z\in\mathbb{L}_j:\ |z|\le r_{\mathrm{b}}\right\}$, and it preserves nearest neighbours, hence bonds and lattice paths. Indeed, for $x,x'\in\bar{K}$ the vector $(x'-x_0)-(x-x_0)$ represents the difference $x'-x$ in $T_j$ and has norm at most $2r_{\mathrm{b}}$; if $x'-x$ has a representative of norm at most $\eta_j$, then every other representative has norm at least $l_j-\eta_j>2r_{\mathrm{b}}$, so that the two representatives coincide.

For a site $x$ of $T_j$ and a direction $\nu$ let $\boldsymbol{\chi}_{j,\nu}(x)$ be the value of $\boldsymbol{\chi}_j$ at the bond from $x$ to $x+\eta_je_{\nu}$. Define the real scalar $\boldsymbol{\lambda}_0$ on $T_j$ by
\[
\boldsymbol{\lambda}_0(x):=\sum_{\nu}\boldsymbol{\chi}_{j,\nu}(x_0)\cdot(x-x_0)_{\nu}\quad\text{for }x\in\bar{K},\qquad\boldsymbol{\lambda}_0(x):=0\quad\text{for }x\notin\bar{K}.
\]
By the definition of $d$ on scalars in Section~\ref{sec:averaging-fluctuations-gauge}, one has $(d\boldsymbol{\lambda}_0)_{\nu}=\boldsymbol{\chi}_{j,\nu}(x_0)$ on every bond of $\bar{K}$, that is, on every bond whose two end points lie in $\bar{K}$.

For a bond function $A$ on $T_j$ let $A\cdot\mathbf{1}_X$ be the bond function that equals $A$ on the bonds that meet $X$ and vanishes on all other bonds, and put
\[
\check{B}:=\left(\boldsymbol{\chi}_j-d\boldsymbol{\lambda}_0\right)\cdot\mathbf{1}_X,
\]
a real bond function. Let $b$ be a bond from $x$ to $x+\eta_je_{\nu}$ that meets $X$. Then $b$ is a bond of $\bar{K}$, and $\check{B}(b)=\boldsymbol{\chi}_{j,\nu}(x)-\boldsymbol{\chi}_{j,\nu}(x_0)$. There is a lattice path in $\bar{K}$ from $x_0$ to $x$ that consists of $\frac{|x-x_0|_1}{\eta_j}\le\frac{3r_{\mathrm{b}}}{\eta_j}$ bonds, and along a bond from $z$ to $z+\eta_je_{\mu}$ the function $\boldsymbol{\chi}_{j,\nu}$ changes by $\eta_j\cdot(\partial_{\mu}\boldsymbol{\chi}_{j,\nu})(z)$. Adding these changes, one obtains
\[
|\check{B}|_{\infty}\le3r_{\mathrm{b}}\cdot|\partial\boldsymbol{\chi}_j|_{\infty}\le3s_X\cdot|\partial\boldsymbol{\chi}_j|_{\infty}.
\]

For $\mathbf{s}\in\mathbb{C}$ the bond function $\left(t\cdot\boldsymbol{\chi}_j-\mathbf{s}\cdot d\boldsymbol{\lambda}_0\right)\cdot\mathbf{1}_X$ depends affinely on $\mathbf{s}$. Let
\[
\mathfrak{U}:=\left\{\mathbf{s}\in\mathbb{C}:\ \left|\left(t\cdot\boldsymbol{\chi}_j-\mathbf{s}\cdot d\boldsymbol{\lambda}_0\right)\cdot\mathbf{1}_X\right|_{\infty}<e_j^{-1+\varepsilon}\right\},
\]
and for $\mathbf{s}\in\mathfrak{U}$ put
\[
\mathfrak{u}(\mathbf{s}):=E^d_j\left(X,\left(t\cdot\boldsymbol{\chi}_j-\mathbf{s}\cdot d\boldsymbol{\lambda}_0\right)\cdot\mathbf{1}_X\right).
\]
For $\mathbf{s}\in\mathfrak{U}$ the bond function $\left(t\cdot\boldsymbol{\chi}_j-\mathbf{s}\cdot d\boldsymbol{\lambda}_0\right)\cdot\mathbf{1}_X$ lies in $G_j$ (Definition~\ref{def:domains}), hence in $\tilde{G}_j$ by~\eqref{it:trajectory-fermionic-operators-SL2}. As $E^d_j(X,\cdot)$ is analytic on $\tilde{G}_j$ by~\eqref{it:trajectory-fermionic-operators-SL4}, the function $\mathfrak{u}$ is analytic on $\mathfrak{U}$. The set $\mathfrak{U}$ is open and convex, because the sup norm is a maximum, over finitely many bonds, of moduli of affine functions of $\mathbf{s}$. It contains $0$: from $t\cdot\boldsymbol{\chi}_j\in\mathfrak{X}_j$ one gets $|t\cdot\boldsymbol{\chi}_j|_{\infty}<e_j^{-1+2\varepsilon}\le e_j^{-1+\varepsilon}$, as $e_j<e_{*}<1$.

Let $\mathbf{s}\in\mathfrak{U}$ be real. Then $\mathbf{s}\cdot\boldsymbol{\lambda}_0$ is a real scalar on $T_j$, and $t\cdot\boldsymbol{\chi}_j-d(\mathbf{s}\cdot\boldsymbol{\lambda}_0)$ lies in $\tilde{G}_j$ by~\eqref{it:trajectory-fermionic-operators-SL2}, since $t\cdot\boldsymbol{\chi}_j\in\mathfrak{X}_j\subset\tilde{G}_j$. On the bonds that meet $X$ this bond function coincides with $\left(t\cdot\boldsymbol{\chi}_j-\mathbf{s}\cdot d\boldsymbol{\lambda}_0\right)\cdot\mathbf{1}_X$. Hence, by~\eqref{it:trajectory-fermionic-operators-SL4}(ii) and~(iii),
\[
\mathfrak{u}(\mathbf{s})=E^d_j\left(X,t\cdot\boldsymbol{\chi}_j-d(\mathbf{s}\cdot\boldsymbol{\lambda}_0)\right)=E^d_j\left(X,t\cdot\boldsymbol{\chi}_j\right).
\]
So $\mathfrak{u}$ is constant on $\mathfrak{U}\cap\mathbb{R}$, which is an open interval containing $0$. As $\mathfrak{U}$ is connected, the identity theorem gives $\mathfrak{u}=E^d_j\left(X,t\cdot\boldsymbol{\chi}_j\right)$ on all of $\mathfrak{U}$.

Let
\[
q_0:=\frac{e_j^{-1+\varepsilon}}{3s_X\cdot|\partial\boldsymbol{\chi}_j|_{\infty}},
\]
with $q_0:=\infty$ if $|\partial\boldsymbol{\chi}_j|_{\infty}=0$. For $|\mathbf{s}|<q_0$ one has $|\mathbf{s}\cdot\check{B}|_{\infty}<e_j^{-1+\varepsilon}$, so that $\mathbf{s}\cdot\check{B}\in G_j\subset\tilde{G}_j$. Hence the function
\[
\mathfrak{v}(\mathbf{s}):=E^d_j\left(X,\mathbf{s}\cdot\check{B}\right)-E^d_j(X,0)
\]
is analytic on the disc $|\mathbf{s}|<q_0$. By~\eqref{it:trajectory-fermionic-operators-SL4}(i) it is bounded by $2C_d\cdot M^3\cdot\exp\left(-(\kappa+1)\mathsf{n}\right)$; it vanishes at $0$, and it is even by~\eqref{it:trajectory-fermionic-operators-SL4}(iv). Therefore $\mathfrak{v}(\mathbf{s})\cdot\mathbf{s}^{-2}$ is analytic on this disc, and Schwarz's lemma (Liouville's theorem if $q_0=\infty$) gives
\[
|\mathfrak{v}(\mathbf{s})|\le2C_d\cdot M^3\cdot\exp\left(-(\kappa+1)\mathsf{n}\right)\cdot\left(\frac{|\mathbf{s}|}{q_0}\right)^2\qquad\left(|\mathbf{s}|<q_0\right).
\]
If $|t|<q_0$, then the bond function $\left(t\cdot\boldsymbol{\chi}_j-t\cdot d\boldsymbol{\lambda}_0\right)\cdot\mathbf{1}_X=t\cdot\check{B}$ has sup norm less than $e_j^{-1+\varepsilon}$; thus $t\in\mathfrak{U}$, and
\[
E^{\mathrm{det}}_j\left(X,t\cdot\boldsymbol{\chi}_j\right)=\mathfrak{u}(t)-E^d_j(X,0)=\mathfrak{v}(t).
\]
If $|t|\ge q_0$, then $\left(\frac{|t|}{q_0}\right)^2\ge1$, and $\left|E^{\mathrm{det}}_j\left(X,t\cdot\boldsymbol{\chi}_j\right)\right|\le2C_d\cdot M^3\cdot\exp\left(-(\kappa+1)\mathsf{n}\right)$ by~\eqref{it:trajectory-fermionic-operators-SL4}(i). In both cases
\[
\left|E^{\mathrm{det}}_j\left(X,t\cdot\boldsymbol{\chi}_j\right)\right|\le2C_d\cdot M^3\cdot\exp\left(-(\kappa+1)\mathsf{n}\right)\cdot\left(\frac{|t|}{q_0}\right)^2,
\]
and, by~\eqref{eq:four-jets-24},
\[
\left(\frac{|t|}{q_0}\right)^2=9s_X^2\cdot e_j^{-2\varepsilon}\cdot|t|^2\cdot\left(e_j\cdot|\partial\boldsymbol{\chi}_j|_{\infty}\right)^2\le9M^2\cdot(\mathsf{n}+3)^2\cdot e_j^{-2\varepsilon}\cdot|t|^2\cdot a_j^4\cdot a_{*}^{-3}\cdot C_{\beta}^2.
\]

By Lemma~\ref{lem:properties-probes}(b) the end points of the bonds of the support of $\hat{\mathfrak{a}}$ lie in the cube $\bar{Q}$, whose side is $D_{\beta}$ in physical coordinates; the isometry $\mathsf{g}_c$ maps it onto a cube of the same side that contains the end points of the bonds of the support of $\hat{\mathfrak{a}}_c$. Physical coordinates are $\eta$ times lattice coordinates (Section~\ref{sec:probes-mirror-copies}), and lattice coordinates are $L^j$ times the coordinates of level $j$ (Section~\ref{sec:lattice-measures-laws}); as $\eta\cdot L^j=a_j$, the end points of the bonds of the support of $\boldsymbol{\chi}_j$ lie in a cube of side $\frac{D_{\beta}}{a_j}$ in the coordinates of level $j$. An $M$-cube that contains one of these end points meets this cube, and there are at most
\[
\left(\frac{D_{\beta}}{M\cdot a_j}+2\right)^3\le4\left(D_{\beta}^3\cdot M^{-3}\cdot a_j^{-3}+8\right)
\]
such $M$-cubes, the inequality by the convexity of $x\mapsto x^3$ on $[0,\infty)$. The polymer $X$ contains one of them, namely an $M$-cube that contains an end point, lying in $X$, of a bond of the support of $\boldsymbol{\chi}_j$. For every $M$-cube $\square$,
\[
\sum_{X'\in D_j,\,X'\supset\square}\left(d_M(X')+3\right)^2\cdot\exp\left(-(\kappa+1)\cdot d_M(X')\right)\le9\sum_{X'\in D_j,\,X'\supset\square}\exp\left(-\kappa_0\cdot d_M(X')\right)\le9K_0,
\]
because $(x+3)^2\cdot\exp(-2x)\le9$ for $x\ge0$ and $\kappa-1\ge\kappa_0$ (stage~\ref{it:order-choice-3}), and by~\eqref{it:trajectory-fermionic-operators-SL4}(v). Combining the bound for $\left|E^{\mathrm{det}}_j\left(X,t\cdot\boldsymbol{\chi}_j\right)\right|$ and the bound for $\left(\frac{|t|}{q_0}\right)^2$ obtained above with the number of the $M$-cubes and the bound $9K_0$, one finds that the polymers of Case~1 at the level $j$ contribute at most
\[
648\,C_d\cdot K_0\cdot C_{\beta}^2\cdot|t|^2\cdot e_j^{-2\varepsilon}\cdot a_{*}^{-3}\cdot\left(M^2\cdot D_{\beta}^3\cdot a_j+8M^5\cdot a_j^4\right)
\]
to the sum of the moduli $\left|E^{\mathrm{det}}_j\left(X,t\cdot\boldsymbol{\chi}_j\right)\right|$, $X\in D_j$; here $648=4\cdot2\cdot9\cdot9$. Now $a_j=a_{*}\cdot L^{-i}$ and $e_j^{-2\varepsilon}=e_{*}^{-2\varepsilon}\cdot L^{i\varepsilon}$, so that
\[
e_j^{-2\varepsilon}\cdot a_{*}^{-3}\cdot\left(M^2\cdot D_{\beta}^3\cdot a_j+8M^5\cdot a_j^4\right)\le e_{*}^{-2\varepsilon}\cdot\left(M^2\cdot D_{\beta}^3\cdot a_{*}^{-2}+8M^5\cdot a_{*}\right)\cdot L^{-i(1-\varepsilon)}.
\]
Moreover $\sum_{i\ge1}L^{-i(1-\varepsilon)}\le1$, as $L^{1-\varepsilon}\ge2$, and $a_{*}^{-2}\le L^4\cdot h'^{-2}$, $a_{*}\le h'$ by Lemma~\ref{lem:scales-stop-level}(a). Summing over the levels $j<k^{*}$, one finds that the polymers of Case~1 contribute at most
\[
648\,C_d\cdot K_0\cdot C_{\beta}^2\cdot\left(M^2\cdot D_{\beta}^3\cdot L^4\cdot h'^{-2}+8M^5\cdot h'\right)\cdot e_{*}^{-2\varepsilon}\cdot|t|^2\le C_{\mathrm{fr}}\cdot e_{*}^{-2\varepsilon}\cdot|t|^2,
\]
by the definition of $C_{\mathrm{fr}}$ in Section~\ref{sec:constants-3}.

\emph{Case 2: $2s_X\ge l_j$.} By Section~\ref{sec:running-quantities-tori} and Lemma~\ref{lem:scales-stop-level}(a), $l_j=\frac{R^{*}}{a_j}=l_{*}\cdot L^i=M\cdot r_{*}\cdot L^i$; hence
\[
\mathsf{n}\ge\frac{l_j}{2M}-3=\frac{r_{*}\cdot L^i}{2}-3,
\]
and $\left|E^{\mathrm{det}}_j\left(X,t\cdot\boldsymbol{\chi}_j\right)\right|\le2C_d\cdot M^3\cdot\exp\left(-(\kappa+1)\mathsf{n}\right)$ by~\eqref{it:trajectory-fermionic-operators-SL4}(i). For every $M$-cube $\square$ and every $y\ge0$,
\[
\sum_{X'\in D_j,\,X'\supset\square,\,d_M(X')\ge y}\exp\left(-(\kappa+1)\cdot d_M(X')\right)\le K_0\cdot\exp(-2y),
\]
because $\kappa+1-\kappa_0\ge2$ and by~\eqref{it:trajectory-fermionic-operators-SL4}(v). As in Case~1, $X$ contains one of the $M$-cubes that contain an end point of a bond of the support of $\boldsymbol{\chi}_j$, and the number of these $M$-cubes is at most
\[
\left(\frac{D_{\beta}}{M\cdot a_j}+2\right)^3\le\left(L^i\cdot\left(\frac{D_{\beta}}{a_{*}}+2\right)\right)^3\le L^{3i}\cdot\mathsf{p},
\]
by the definition of $\mathsf{p}$ in Section~\ref{sec:probes-mirror-copies}. Condition~\eqref{eq:renormalisation-group-16} gives $r_{*}>31L+16\ge12$; so $y:=\frac{r_{*}\cdot L^i}{2}-3\ge0$, and $\exp(-2y)=\exp(6)\cdot\exp\left(-r_{*}\cdot L^i\right)$. Consequently the polymers of Case~2 at all levels $j<k^{*}$ contribute at most
\[
2\exp(6)\,C_d\cdot M^3\cdot K_0\cdot\mathsf{p}\cdot\sum_{i\ge1}L^{3i}\cdot\exp\left(-r_{*}\cdot L^i\right).
\]
Let $i\ge1$ and $x:=L^i$. Then $x\ge3^i\ge3$, as $L\ge3$ (stage~\ref{it:order-choice-4}). Hence $r_{*}\cdot x-r_{*}-8x=(r_{*}-8)\cdot x-r_{*}\ge2r_{*}-24\ge0$; and $x^3\cdot\exp(-8x)\le2^{-i}$, since $x\mapsto x^3\cdot\exp(-8x)$ decreases for $x\ge\frac{3}{8}$ and $54^i\le\exp(24i)\le\exp\left(8\cdot3^i\right)$. Therefore
\[
x^3\cdot\exp(-r_{*}\cdot x)\le\exp(-r_{*})\cdot x^3\cdot\exp(-8x)\le\exp(-r_{*})\cdot2^{-i},
\]
and the sum $\sum_{i\ge1}L^{3i}\cdot\exp\left(-r_{*}\cdot L^i\right)$ is at most $\exp(-r_{*})$. Thus the polymers of Case~2 contribute at most
\[
2\exp(6)\,C_d\cdot M^3\cdot K_0\cdot\mathsf{p}\cdot\exp(-r_{*})\le B_{\mathrm{wr}},
\]
by the definition of $B_{\mathrm{wr}}$ in Section~\ref{sec:constants-3}. Adding the contributions of the two cases, one obtains (b).

\emph{(d)} Let $|t|\le t_x$. Since $t_x\le e_{*}^{3\varepsilon}$ by the definition in Section~\ref{sec:probes-mirror-copies}, one has $e_{*}^{-2\varepsilon}\cdot|t|^2\le e_{*}^{4\varepsilon}$. By (b), and by (c) for $\theta=0$,
\[
\left|\mathrm{lgN}(t)\right|\le C_{\mathrm{fr}}\cdot e_{*}^{4\varepsilon}+B_{\mathrm{wr}}+B_1=B_0\le\frac{1}{4},
\]
by the definition of $B_0$ in Section~\ref{sec:constants-3} and the last condition of~\eqref{eq:four-jets-sm4}.

The function $\mathrm{lgN}$ is analytic on the disc $|t|<t_x^{+}$. Indeed, for $X\in D_{*}$ the function $\mathfrak{P}_{*}(X,\cdot)$ is analytic on $\mathfrak{X}_{*}$ (Lemma~\ref{lem:analyticity-locality-decay}), and $t\cdot\hat{c}_c\in\mathfrak{X}_{*}$ for $|t|<t_x^{+}$ by~\eqref{it:trajectory-fermionic-operators-SL7}. The entries of $D_{*}\left(t\cdot\hat{c}_c\right)$ are analytic functions of $t$, because $S$ is analytic on $\tilde{G}_{*}$ by~\eqref{it:trajectory-fermionic-operators-SL2} and the entries of $Q_{*}(A)$ and $Q_{*}^T(-A)$ are entire functions of $A$ (Section~\ref{sec:fermionic-operators-stop}). So the matrix $\mathfrak{Y}$ of the proof of (c), formed with $\theta=0$, is an analytic function of $t$, and by the bounds of that proof the series $\sum_{n\ge1}\frac{(-1)^{n+1}}{n}\cdot\operatorname{Tr}\mathfrak{Y}^n$ converges uniformly on the disc. Hence $g_D^0$ is analytic on the disc, and so is $\mathrm{lgN}$. At $t=0$ one has $\mathfrak{Y}=0$, and $\mathfrak{P}_{*}(X,0)=0$ by Lemma~\ref{lem:analyticity-locality-decay}; so $\mathrm{lgN}(0)=0$.

By Lemma~\ref{lem:signs-symmetries-wilson}(b), $N_0(-t)=N_0(t)$. By (a), the functions $t\mapsto\mathrm{lgN}(t)$ and $t\mapsto\mathrm{lgN}(-t)$ are continuous on the disc $|t|<t_x^{+}$, the exponential of each of them is $N_0$, and both vanish at $0$. Their difference is continuous, takes its values in $2\pi i\,\mathbb{Z}$ and vanishes at $0$; as the disc is connected, it vanishes identically. So $\mathrm{lgN}$ is even. The same argument shows that $\mathrm{lgN}$ is the only continuous function on the disc $|t|<t_x^{+}$ that vanishes at $0$ and whose exponential is $N_0$. By Lemma~\ref{lem:signs-symmetries-wilson}(b), $N_0$ does not depend on $c$; hence $\mathrm{lgN}$ does not depend on $c$. The function $g_D^{\theta}$, on the other hand, is formed with $\hat{c}_c$, that is, with the cell $c$.

Let $t$ be real with $|t|\le t_x$. By Lemma~\ref{lem:signs-symmetries-wilson}(b) the number $N_0(t)=\exp\mathrm{lgN}(t)$ is real; hence $\operatorname{Im}\mathrm{lgN}(t)\in\pi\mathbb{Z}$. As $\left|\operatorname{Im}\mathrm{lgN}(t)\right|\le\left|\mathrm{lgN}(t)\right|\le\frac{1}{4}<\pi$, it follows that $\operatorname{Im}\mathrm{lgN}(t)=0$.

\emph{(e)} The function $N_0$ is entire (Section~\ref{sec:averaged-law-partition}), so $\phi_0$ is defined, and $n_i$ is the coefficient of $t^i$ in the Taylor series of $N_0$ at $0$. By Lemma~\ref{lem:signs-symmetries-wilson}(b), $N_0$ is even and real on the real axis, and $N_0(0)=1$; hence $n_0=1$, $n_1=n_3=0$, and the numbers $n_i$ are real. For $|t|\le t_x$, by (a) and (d),
\[
\left|N_0(t)-1\right|=\left|\exp\mathrm{lgN}(t)-1\right|\le\exp\left(\left|\mathrm{lgN}(t)\right|\right)-1\le\exp(B_0)-1\le2B_0,
\]
the last inequality because $0\le B_0\le1$. For $i\ge1$ the coefficient of $t^i$ in the Taylor series of $N_0-1$ at $0$ is $n_i$, and Cauchy's estimate on the circle $|t|=t_x$ gives $|n_i|\le2B_0\cdot t_x^{-i}$ for $1\le i\le4$. Now $c_2=2n_2$, $c_4=24n_4$, $t_x=4\mathbf{r}$ and $B_0\le\frac{1}{4}$. Therefore
\[
|c_2|\le\frac{4B_0}{16\mathbf{r}^2}\le\frac{\mathbf{r}^{-2}}{16},\qquad|c_4|\le\frac{48B_0}{256\mathbf{r}^4}\le\frac{\mathbf{r}^{-4}}{16},
\]
and
\[
\|\phi_0-1\|_{\mathbf{r}}=\sum_{i=1}^4|n_i|\cdot\mathbf{r}^i\le2B_0\sum_{i\ge1}4^{-i}\le\frac{1}{6},\qquad\|\phi_0\|_{\mathbf{r}}\le1+\|\phi_0-1\|_{\mathbf{r}}\le2.
\]
Finally, $\phi_0(t)=1+n_2t^2+n_4t^4$; so, as $t\to0$,
\[
\log\phi_0(t)=n_2t^2+\left(n_4-\frac{n_2^2}{2}\right)t^4+O\left(t^6\right),
\]
and $(\log\phi_0)^{(4)}(0)=24n_4-12n_2^2=c_4-3c_2^2$. Since $N_0(t)=\phi_0(t)+O\left(t^5\right)$, one has
\[
\log N_0(t)-\log\phi_0(t)=\log\left(1+\frac{N_0(t)-\phi_0(t)}{\phi_0(t)}\right)=O\left(t^5\right);
\]
so the fourth derivatives at $0$ of $\log N_0$ and of $\log\phi_0$ coincide, and $b_{\eta,R^{*}}(0)=c_4-3c_2^2$ by the definition of $b_{\eta,R^{*}}(0)$ in Section~\ref{sec:averaged-law-partition}. This proves~\eqref{eq:four-jets-23} and (e). In particular $c_1=c_3=0$, and Theorem~\ref{thm:free-ratio} follows.
\end{proof}

\subsubsection{A single cell}\label{sec:single-cell}

\begin{proposition}[Comparison with the free ratio]\label{prop:comparison-free-ratio}
Assume~\eqref{eq:four-jets-sm2}, \eqref{eq:four-jets-sm3} and~\eqref{eq:four-jets-sm4}. Let $c\in\mathrm{Cel}$, $\theta\in\mathbb{R}^3$ and $A\in V_{\mathrm{t}}\left(T^0_{*}\right)$ with $\chi_{*}(A)=1$, and let $N\ge N_{\mathrm{min}}\left(R^{*},e\right)$. For $t_c\in\mathbb{C}$ with $|t_c|<t_x^{+}$ put
\[
w_c(t_c):=\mathrm{Lg}\left(t_c\cdot e_c;\theta,A\right)-\mathrm{lgN}(t_c),
\]
where $\mathrm{lgN}$ is the function of Lemma~\ref{lem:free-wilson-ratio}. Then the function $w_c$ is analytic on the disc $\left\{t_c\in\mathbb{C}:|t_c|<t_x^{+}\right\}$, one has $w_c(0)=0$, and
\[
\left|w_c(t_c)\right|\le3\,e_{*}^{\frac{1}{4}-9\varepsilon}\qquad\left(|t_c|<t_x^{+}\right).
\]
\end{proposition}

\begin{proof}
Put $y:=e_{*}^{\frac{1}{4}-9\varepsilon}$, and let $t_c\in\mathbb{C}$ with $|t_c|<t_x^{+}$. The point $t_c\cdot e_c$ of $\mathbb{C}^{\mathrm{Cel}}$ lies in $\mathbb{D}^{\circ}\left(t_x^{+}\right)$, and $\mathscr{C}\left(t_c\cdot e_c\right)=t_c\cdot\hat{c}_c$. The placement is admitted by~\eqref{eq:four-jets-sm2}; so~\eqref{it:trajectory-fermionic-operators-SL7} gives $t_c\cdot\hat{c}_c\in\mathfrak{X}_{*}$ and $e_{*}\cdot\left|t_c\cdot\hat{c}_c\right|_{\infty}<e_{*}^{3\varepsilon}$. Moreover $e_{*}\le e_0<1$ by Lemma~\ref{lem:scales-stop-level}(a) and~\eqref{eq:renormalisation-group-A8}.

\emph{Step 1: the free determinant is split off.} By~\eqref{eq:four-jets-20},
\[
\mathrm{Rt}\left(t_c\cdot e_c;\theta,A\right)=\exp\left\{\sum_{X\in D_{*}}\mathfrak{P}_{*}\left(X,t_c\cdot\hat{c}_c\right)\right\}\cdot\frac{\mathsf{Z}\left(t_c\cdot e_c\right)}{\mathsf{Z}(0)}.
\]
Lemma~\ref{lem:terminal-integral-compared} gives $\mathsf{Z}(t)=\mathsf{f}(t)\cdot\det D_{*}\left(c_{\theta}+\mathscr{C}(t)\right)$ for $t\in\mathbb{D}^{\circ}\left(t_x^{+}\right)$, and Lemma~\ref{lem:free-wilson-ratio}(c), for the cell $c$, gives
\[
\frac{\det D_{*}\left(c_{\theta}+t_c\cdot\hat{c}_c\right)}{\det D_{*}\left(c_{\theta}\right)}=\exp g_D^{\theta}(t_c).
\]
Consequently
\[
\mathrm{Rt}\left(t_c\cdot e_c;\theta,A\right)=\exp\left\{\sum_{X\in D_{*}}\mathfrak{P}_{*}\left(X,t_c\cdot\hat{c}_c\right)\right\}\cdot\exp\left(g_D^{\theta}(t_c)\right)\cdot\frac{\mathsf{f}\left(t_c\cdot e_c\right)}{\mathsf{f}(0)}.
\]
By Lemma~\ref{lem:terminal-integral-compared}, moreover, $\mathsf{f}$ is analytic on $\mathbb{D}^{\circ}\left(t_x^{+}\right)$ and satisfies $\left|\mathsf{f}(t)-1\right|\le\frac{y}{2}\le\frac{1}{16}$ there. Therefore the principal branch of $\log\mathsf{f}$ is analytic on $\mathbb{D}^{\circ}\left(t_x^{+}\right)$, and, since $\left|\log(1+z)\right|\le\frac{|z|}{1-|z|}$ for $|z|<1$,
\[
\left|\log\mathsf{f}(t)\right|\le\frac{16}{15}\cdot\frac{y}{2}\le y\qquad\left(t\in\mathbb{D}^{\circ}\left(t_x^{+}\right)\right).
\]
The function $t_c\mapsto\mathrm{Lg}\left(t_c\cdot e_c;\theta,A\right)$ is analytic on the disc $|t_c|<t_x^{+}$ by Section~\ref{sec:construction-logarithm}. So is the function
\[
t_c\mapsto\sum_{X\in D_{*}}\mathfrak{P}_{*}\left(X,t_c\cdot\hat{c}_c\right)+g_D^{\theta}(t_c)+\log\mathsf{f}\left(t_c\cdot e_c\right)-\log\mathsf{f}(0).
\]
Indeed, the sum over $X$ is analytic by Lemma~\ref{lem:analyticity-locality-decay}, because $t_c\cdot\hat{c}_c\in\mathfrak{X}_{*}$, and the two logarithms by what has just been shown. As to $g_D^{\theta}$, the connection $c_{\theta}+t_c\cdot\hat{c}_c$ lies in $G_{\mathrm{w}}\subset\tilde{G}_{*}$ (proof of Lemma~\ref{lem:free-wilson-ratio}(c)); the operator $S$ is analytic on $\tilde{G}_{*}$ by~\eqref{it:trajectory-fermionic-operators-SL2}, and the entries of $Q_{*}(A)$ and $Q_{*}^T(-A)$ are entire functions of $A$ (Section~\ref{sec:fermionic-operators-stop}). Hence the entries of $D_{*}\left(c_{\theta}+t_c\cdot\hat{c}_c\right)$, and with them the entries of the matrix $\mathfrak{Y}$ of the proof of Lemma~\ref{lem:free-wilson-ratio}(c), formed with $t=t_c$, are analytic functions of $t_c$ on the disc; so is every term $\operatorname{Tr}\left(\mathfrak{Y}^n\right)$ of the series that defines $g_D^{\theta}(t_c)$. By the bounds in that proof, $\left|\mathfrak{Y}\right|_{\ell^2}\le\frac{1}{2}$ and
\[
\left|\operatorname{Tr}\left(\mathfrak{Y}^n\right)\right|\le\left|\mathfrak{Y}\right|_{\ell^2}^{n-1}\cdot\left|\mathfrak{Y}\right|_1\le2^{1-n}\cdot4\,g_1\cdot C_K\cdot\mathsf{p}\cdot e_{*}^{2\varepsilon}\qquad\left(n\ge1\right)
\]
for every $t_c$ in the disc. So the series converges uniformly on the disc, and $g_D^{\theta}$ is analytic there. Both functions vanish at $t_c=0$, since $\mathrm{Lg}(0;\theta,A)=0$, $\mathfrak{P}_{*}(X,0)=0$ by Lemma~\ref{lem:analyticity-locality-decay}, and $g_D^{\theta}(0)=\operatorname{Tr}\log1=0$. The exponential of the first function is $\mathrm{Rt}\left(t_c\cdot e_c;\theta,A\right)$, because $\exp\mathrm{Lg}=\mathrm{Rt}$ (Section~\ref{sec:construction-logarithm}); the exponential of the second is the same number, by the factorisation of $\mathrm{Rt}$ obtained above. The difference of the two functions is thus a continuous function on the connected disc, with values in $2\pi i\,\mathbb{Z}$, that vanishes at $0$; hence it vanishes identically. As $\mathrm{lgN}(t_c)=\sum_{X\in D_{*}}\mathfrak{P}_{*}\left(X,t_c\cdot\hat{c}_c\right)+g_D^0(t_c)$ by the definition in Lemma~\ref{lem:free-wilson-ratio}, it follows that
\[
w_c(t_c)=g_D^{\theta}(t_c)-g_D^0(t_c)+\log\mathsf{f}\left(t_c\cdot e_c\right)-\log\mathsf{f}(0)\qquad\left(|t_c|<t_x^{+}\right).
\]
In particular $w_c$ is analytic, $w_c(0)=0$, and each of the last two terms has modulus at most $y$.

\emph{Step 2: the dependence on the twist.} Keep $t_c$ fixed, and consider the closed tube
\[
\left\{\theta\in\mathbb{C}^3:\ \left|\operatorname{Im}\theta_{\nu}\right|\le\tau_s\ \text{for}\ \nu=0,1,2\right\}.
\]
For $\theta\in\mathbb{C}^3$ let $c_{\theta}$ be the constant bond function on $T_{*}$ whose component in the direction $\nu$ is $\theta_{\nu}/(e_{*}\cdot l_{*})$; for real $\theta$ this is the twist of Definition~\ref{def:twist-external-connection}. Then $c_{\theta}=c_{\operatorname{Re}\theta}+i\cdot c_{\operatorname{Im}\theta}$ with the real constant bond functions $c_{\operatorname{Re}\theta}$ and $c_{\operatorname{Im}\theta}$, and, as $\tau_s=c_1\cdot l_{*}$ (Section~\ref{sec:constants-3}), for every $\theta$ in the closed tube
\[
e_{*}\cdot\left|\operatorname{Im}c_{\theta}\right|_{\infty}=\max_{\nu}\frac{\left|\operatorname{Im}\theta_{\nu}\right|}{l_{*}}\le\frac{\tau_s}{l_{*}}=c_1.
\]
Let $\theta$ lie in the closed tube, and let $0\le u\le1$. The real constant $c_{\operatorname{Re}\theta}$ lies in $\tilde{R}_{*}$, which contains $0$ and is invariant under the addition of real constants by~\eqref{it:trajectory-fermionic-operators-SL2}; and $u\cdot t_c\cdot\hat{c}_c=\mathscr{C}\left(u\cdot t_c\cdot e_c\right)$ lies in $\mathfrak{X}_{*}$ by~\eqref{it:trajectory-fermionic-operators-SL7}. So $c_{\operatorname{Re}\theta}+u\cdot t_c\cdot\hat{c}_c\in G_{\mathrm{w}}$, and
\[
c_{\theta}+u\cdot t_c\cdot\hat{c}_c=\left(c_{\operatorname{Re}\theta}+u\cdot t_c\cdot\hat{c}_c\right)+A_1,\qquad A_1:=i\cdot c_{\operatorname{Im}\theta},\qquad e_{*}\cdot\left|A_1\right|_{\infty}\le c_1,
\]
is an element of $G_{\mathrm{w}}^{+}$ (Section~\ref{sec:dependence-stop-level}).

By Lemma~\ref{lem:operators}, the operator $\mathsf{D}\left(c_{\theta}+u\cdot t_c\cdot\hat{c}_c\right)$ is invertible for every $\theta$ in the closed tube and every $u\in[0,1]$. The entries of its matrix are entire functions of $\theta\in\mathbb{C}^3$; the inverse of an invertible matrix that depends analytically on a parameter depends analytically on it; and invertibility persists under small perturbations. Hence $D_{*}\left(c_{\theta}\right)$ and $D_{*}\left(c_{\theta}+t_c\cdot\hat{c}_c\right)$, given by the formula of Section~\ref{sec:fermionic-operators-stop}, are analytic functions of $\theta$ on a neighbourhood of the closed tube.

Let $\theta$ lie in the closed tube. Lemma~\ref{lem:inverses}(b) applies to $c_{\theta}=A_r+A_1$ with the real part $A_r:=c_{\operatorname{Re}\theta}$: the operator $\mathsf{D}\left(c_{\theta}\right)$ is invertible, and
\[
e_{*}\cdot\left|A_1\right|_{\infty}\le c_1+e_{*}^{2\varepsilon}\le\frac{\mathfrak{m}_{-}}{48}+\frac{\mathfrak{m}_{-}}{48}<\frac{\mathfrak{m}_{*}}{24},
\]
by $c_1\le\mathfrak{m}_{-}/48$ (Section~\ref{sec:constants-3}), by the first condition of~\eqref{eq:four-jets-sm4}, and by $\mathfrak{m}_{-}<\mathfrak{m}_{*}$ (Lemma~\ref{lem:scales-stop-level}(a)). Hence $D_{*}\left(c_{\theta}\right)$ is invertible, and
\[
\left|D_{*}\left(c_{\theta}\right)^{-1}\right|_{\ell^2}\le g_1.
\]
As invertibility persists under small perturbations, $D_{*}\left(c_{\theta}\right)^{-1}$ exists and is analytic in $\theta$ on a neighbourhood of the closed tube.

Next, Lemma~\ref{lem:operators} applies to $A':=c_{\theta}$ and $B:=t_c\cdot\hat{c}_c$, since $A'+u\cdot B\in G_{\mathrm{w}}^{+}$ for $0\le u\le1$. Here
\[
\epsilon_B=e_{*}\cdot\left|t_c\cdot\hat{c}_c\right|_{\infty}<e_{*}^{3\varepsilon}\le e_{*}^{2\varepsilon},
\]
and $P_B\subset P_c\subset P_c^{+}$, so that $\#|P_B|\le\mathsf{p}$ by Lemma~\ref{lem:properties-probes}(e). With~\eqref{eq:four-jets-Maj2} and $\|\cdot\|_0\le\|\cdot\|_{\gamma_F/2}$, Lemma~\ref{lem:operators} shows that the matrix
\[
\mathfrak{K}:=D_{*}\left(c_{\theta}+t_c\cdot\hat{c}_c\right)-D_{*}\left(c_{\theta}\right)
\]
satisfies
\[
\left|\mathfrak{K}\right|_{\ell^2}\le2\left\|\mathfrak{K}\right\|_{\gamma_F/2}\le2\,C_K'\cdot e_{*}^{2\varepsilon},\qquad\left|\mathfrak{K}\right|_1\le4\sum_{z,z'\in T^0_{*}}\left|\mathfrak{K}\right|(z,z')\le4\,C_K'\cdot e_{*}^{2\varepsilon}\cdot\mathsf{p}.
\]
Let $\mathfrak{Y}:=D_{*}\left(c_{\theta}\right)^{-1}\cdot\mathfrak{K}$. The matrices $\mathfrak{K}$ and $\mathfrak{Y}$ are given by these formulas, and depend analytically on $\theta$, on a neighbourhood of the closed tube, as $D_{*}\left(c_{\theta}\right)$, $D_{*}\left(c_{\theta}+t_c\cdot\hat{c}_c\right)$ and $D_{*}\left(c_{\theta}\right)^{-1}$ do. By the second condition of~\eqref{eq:four-jets-sm4}, and by $|Y\cdot Y'|_1\le|Y|_{\ell^2}\cdot|Y'|_1$ for matrices $Y$, $Y'$, for $\theta$ in the closed tube
\[
\left|\mathfrak{Y}\right|_{\ell^2}\le2\,g_1\cdot C_K'\cdot e_{*}^{2\varepsilon}\le\frac{1}{2},\qquad\left|\mathfrak{Y}\right|_1\le\left|D_{*}\left(c_{\theta}\right)^{-1}\right|_{\ell^2}\cdot\left|\mathfrak{K}\right|_1\le4\,g_1\cdot C_K'\cdot\mathsf{p}\cdot e_{*}^{2\varepsilon}.
\]
As $\mathfrak{Y}$ depends continuously on $\theta$, there is a neighbourhood of the closed tube on which $\left|\mathfrak{Y}\right|_{\ell^2}<1$. For $\theta$ in such a neighbourhood put
\[
g_D^{\theta}(t_c):=\operatorname{Tr}\log\left(1+\mathfrak{Y}\right)=\sum_{n\ge1}\frac{(-1)^{n+1}}{n}\operatorname{Tr}\left(\mathfrak{Y}^n\right);
\]
for real $\theta$ this is the number $g_D^{\theta}(t_c)$ of Lemma~\ref{lem:free-wilson-ratio}, formed with the cell $c$. The series converges locally uniformly, and its terms are analytic in $\theta$; so $\theta\mapsto g_D^{\theta}(t_c)$ is analytic on a neighbourhood of the closed tube. By $|\operatorname{Tr}Y|\le|Y|_1$ and $|Y\cdot Y'|_1\le|Y|_{\ell^2}\cdot|Y'|_1$, for $\theta$ in the closed tube
\[
\left|g_D^{\theta}(t_c)\right|\le\sum_{n\ge1}\frac{1}{n}\left|\mathfrak{Y}\right|_1\cdot\left|\mathfrak{Y}\right|_{\ell^2}^{n-1}\le2\left|\mathfrak{Y}\right|_1\le8\,g_1\cdot C_K'\cdot\mathsf{p}\cdot e_{*}^{2\varepsilon}\le B_1,
\]
the last inequality by the definition of $B_1$ in Section~\ref{sec:constants-3}.

The function $\theta\mapsto g_D^{\theta}(t_c)$ is $2\pi$-periodic in each variable. Indeed, let $\theta\in\mathbb{R}^3$, let $\nu\in\{0,1,2\}$, and let $e_{\nu}$ be the unit vector of $\mathbb{R}^3$ in the direction $\nu$. The difference $\boldsymbol{\omega}:=c_{\theta+2\pi e_{\nu}}-c_{\theta}$ is the real constant bond function on $T_{*}$ whose component in the direction $\nu$ is $2\pi/(e_{*}\cdot l_{*})$ and whose other components vanish; thus $e_{*}\cdot l_{*}\cdot\boldsymbol{\omega}_{\mu}\in2\pi\mathbb{Z}$ for every direction $\mu$, and the same holds for $-\boldsymbol{\omega}$. One has $k^{*}\ge1$, and the data $\left(k^{*},l_{*},\mathfrak{m}_{*},e_{*}\right)$ belong to the scale-free setting of Section~\ref{sec:properties-scheme} (Lemma~\ref{lem:scales-stop-level}(a), (c)); so Proposition~\ref{prop:symmetries-fermion-operators}(a) applies at the level $k^{*}$. Applied to the real constant $-\boldsymbol{\omega}$, it provides a diagonal matrix $u^0$, whose diagonal entries have modulus one and which does not depend on the connection, such that
\[
D_{*}\left(A'+\boldsymbol{\omega}\right)=u^0\cdot D_{*}(A')\cdot\left(u^0\right)^{-1}\qquad\left(A'\in\tilde{G}_{*}\right).
\]
For real $\theta$ the connections $c_{\theta}$ and $c_{\theta}+t_c\cdot\hat{c}_c$ lie in $G_{\mathrm{w}}\subset\tilde{G}_{*}$. Replacing $\theta$ by $\theta+2\pi e_{\nu}$ therefore conjugates $D_{*}\left(c_{\theta}\right)$ and $D_{*}\left(c_{\theta}+t_c\cdot\hat{c}_c\right)$ by $u^0$, hence also $\mathfrak{K}$, $D_{*}\left(c_{\theta}\right)^{-1}$ and $\mathfrak{Y}$, and it leaves every trace $\operatorname{Tr}\left(\mathfrak{Y}^n\right)$ unchanged. So
\[
g_D^{\theta+2\pi e_{\nu}}(t_c)=g_D^{\theta}(t_c)\qquad\left(\theta\in\mathbb{R}^3\right).
\]
Both sides are analytic functions of $\theta$ on the interior of the closed tube, an open connected set that contains $\mathbb{R}^3$ and is invariant under real translations. By the identity theorem they coincide on this interior, and by continuity on the closed tube.

Thus $\theta\mapsto g_D^{\theta}(t_c)$ is continuous on the closed tube, holomorphic in its interior, $2\pi$-periodic in each variable, and bounded in modulus by $B_1$. The third condition of~\eqref{eq:four-jets-sm4} states that $\tau_s\ge1$. Lemma~\ref{lem:oscillation-real-points}, applied with $\boldsymbol{\tau}=\tau_s$ and $B_{\mathfrak{g}}=B_1$, and the fourth condition of~\eqref{eq:four-jets-sm4} give
\[
\left|g_D^{\theta}(t_c)-g_D^0(t_c)\right|\le52\,B_1\cdot\exp\left(-\tau_s\right)\le y\qquad\left(\theta\in\mathbb{R}^3\right).
\]

\emph{Step 3: conclusion.} By the formula for $w_c$ in Step~1, the bounds for $\log\mathsf{f}$ obtained there, and Step~2,
\[
\left|w_c(t_c)\right|\le\left|g_D^{\theta}(t_c)-g_D^0(t_c)\right|+\left|\log\mathsf{f}\left(t_c\cdot e_c\right)\right|+\left|\log\mathsf{f}(0)\right|\le y+y+y=3\,e_{*}^{\frac{1}{4}-9\varepsilon}.
\]
\end{proof}

\subsubsection{Two or more cells}\label{sec:two-more-cells}

\begin{proposition}[Bound for the M\"obius components]\label{prop:bound-mobius-components}
Assume~\eqref{eq:four-jets-sm1}--\eqref{eq:four-jets-sm4}. Let $S\subset\mathrm{Cel}$ with $\#|S|\ge2$, let $\theta\in\mathbb{R}^3$ and $A\in V_{\mathrm{t}}\left(T^0_{*}\right)$ with $\chi_{*}(A)=1$, and let $N\ge N_{\mathrm{min}}\left(R^{*},e\right)$. Let $\mathrm{Lg}_S$ be the component with the index $S$, in the sense of Section~\ref{sec:mobius-decomposition-cell}, of the function $t\mapsto\mathrm{Lg}(t;\theta,A)$ on $\mathbb{D}\left(t_x/2\right)$. Then
\[
\sup_{t\in\mathbb{D}\left(t_x/2\right)}\left|\mathrm{Lg}_S(t)\right|\le\varpi_{\mathrm{sep}}.
\]
\end{proposition}

\begin{proof}
Write $\mathrm{Lg}(t)$ for $\mathrm{Lg}(t;\theta,A)$; all components below are formed, in the sense of Section~\ref{sec:mobius-decomposition-cell}, for functions on $\mathbb{D}\left(t_x/2\right)$.

If $\hat{\mathfrak{a}}=0$, then $\hat{\mathfrak{a}}_c=(\mathsf{g}_c)_{*}\hat{\mathfrak{a}}$ and $\hat{c}_c=\frac{1}{e}\cdot\hat{\mathfrak{a}}_c^{[k^{*}]}$ vanish for every cell $c$ (Section~\ref{sec:probes-mirror-copies}), so that $\mathscr{C}(t)=0$ for every $t$. In this case $\mathrm{Rt}(t;\theta,A)=1$ for $t\in\mathbb{D}^{\circ}\left(t_x^{+}\right)$ by~\eqref{eq:four-jets-20}, the integrand in~\eqref{eq:four-jets-21} vanishes, and $\mathrm{Lg}=0$; hence $\mathrm{Lg}_S=0$, and the assertion holds, because $\varpi_{\mathrm{sep}}\ge0$. Let now $\hat{\mathfrak{a}}\ne0$. Then $d_{\mathrm{sep}}=\min_{c\ne c'}d\left(P_c,P_{c'}\right)$ (Section~\ref{sec:probes-mirror-copies}). Moreover $\hat{\mathfrak{a}}^{[k^{*}]}\ne0$, because $b\mapsto L^{k^{*}}\cdot b$ maps the bonds of $T_{*}$ onto those of $\Lambda_{n_{*}}$ (Section~\ref{sec:lattice-measures-laws}); so Lemma~\ref{lem:properties-probes}(d) gives $\left|\hat{c}_c\right|_{\infty}=e^{-1}\cdot\left|\hat{\mathfrak{a}}^{[k^{*}]}\right|_{\infty}>0$ for every cell $c$. Thus every set $\mathrm{Bd}_c$ is non-empty, and so is every set $P_c$: it contains the unit site $[w]$ (Section~\ref{sec:fermionic-operators-stop}) of each end point $w$ of a bond of $\mathrm{Bd}_c$.

Since $\mathbb{D}\left(t_x/2\right)\subset\mathbb{D}^{\circ}\left(t_x\right)$, the identity~\eqref{eq:four-jets-22} gives
\[
\mathrm{Lg}=\sum_{X\in D_{*}}\mathfrak{P}_{*}\left(X,\mathscr{C}(\cdot)\right)+W-W(0)\qquad\text{on }\mathbb{D}\left(t_x/2\right).
\]
The map $\mathfrak{g}\mapsto\mathfrak{g}_S$ is linear, and, as $S\ne\emptyset$, it maps constant functions to $0$ by Lemma~\ref{lem:properties-mobius-components}(a). Hence
\[
\mathrm{Lg}_S=\left(\sum_{X\in D_{*}}\mathfrak{P}_{*}\left(X,\mathscr{C}(\cdot)\right)\right)_S+W_S.
\]

\emph{(i) The free activities.} For $t\in\mathbb{D}^{\circ}\left(t_x^{+}\right)$ one has $\mathscr{C}(t)\in\mathfrak{X}_{*}$ by~\eqref{it:trajectory-fermionic-operators-SL7}, which applies because the placement is admitted by~\eqref{eq:four-jets-sm2}. Let $X\in D_{*}$ and $c\in\mathrm{Cel}$. Since $\mathscr{C}(t)=\sum_{c''}t_{c''}\cdot\hat{c}_{c''}$ and $\hat{c}_c$ vanishes outside $\mathrm{Bd}_c$, the values of $\mathscr{C}(t)$ on the bonds that meet $X$ do not depend on $t_c$ unless some bond of $\mathrm{Bd}_c$ meets $X$. By the locality stated in Lemma~\ref{lem:analyticity-locality-decay}, the function $t\mapsto\mathfrak{P}_{*}\left(X,\mathscr{C}(t)\right)$ therefore does not depend on $t_c$ unless a bond of $\mathrm{Bd}_c$ meets $X$. If a bond $b\in\mathrm{Bd}_c$ has an end point $w\in X$, then the unit site $[w]$ belongs to $P_c$, and it belongs to $X$, because the unit cubes are nested in the $M$-cubes; so in this case $X$ contains a point of $P_c$, as in the proof of Lemma~\ref{lem:bounds-activities}(c). By Lemma~\ref{lem:properties-mobius-components}(a) it follows that
\[
\left(\mathfrak{P}_{*}\left(X,\mathscr{C}(\cdot)\right)\right)_S=0\qquad\text{unless }X\text{ contains a point of }P_c\text{ for every }c\in S.
\]
Let $X$ contain a point of $P_c$ for every $c\in S$. As $\#|S|\ge2$, there are cells $c\ne c'$ in $S$, and $X$ contains a point of $P_c$ and a point of $P_{c'}$, whose distance is at least $d\left(P_c,P_{c'}\right)\ge d_{\mathrm{sep}}$. By~\eqref{it:trajectory-fermionic-operators-SL4}(v) this distance is at most $M\cdot\left(d_M(X)+2\right)$; hence, by~\eqref{eq:four-jets-sm2},
\[
d_M(X)\ge\frac{d_{\mathrm{sep}}}{M}-2\ge1.
\]
So Lemma~\ref{lem:analyticity-locality-decay} and Lemma~\ref{lem:properties-mobius-components}(b) give
\[
\sup_{\mathbb{D}\left(t_x/2\right)}\left|\left(\mathfrak{P}_{*}\left(X,\mathscr{C}(\cdot)\right)\right)_S\right|\le2^{\#|S|}\cdot2\,C_d\cdot M^3\cdot\exp\left(-\kappa\cdot d_M(X)\right).
\]
Now fix $c\in S$. A polymer contains a point $z$ if and only if it contains the $M$-cube of $z$. One has $\kappa-\kappa_0\ge2$, because $\kappa\ge7\kappa_0+7$ by stage~\ref{it:order-choice-3}. By the first assertion of~\eqref{it:trajectory-fermionic-operators-SL4}(v), for every point $z$,
\begin{align*}
&\sum_{X\ni z,\ d_M(X)\ge\frac{d_{\mathrm{sep}}}{M}-2}\exp\left(-\kappa\cdot d_M(X)\right)\\
&\qquad\le\exp\left(-2\left(\frac{d_{\mathrm{sep}}}{M}-2\right)\right)\cdot\sum_{X\ni z}\exp\left(-\kappa_0\cdot d_M(X)\right)\\
&\qquad\le K_0\cdot\exp\left(-2\left(\frac{d_{\mathrm{sep}}}{M}-2\right)\right).
\end{align*}
Summing over $z\in P_c$, and using $\#|P_c|\le\#|P_c^{+}|\le\mathsf{p}$ (Lemma~\ref{lem:properties-probes}(e)) and $\vartheta\le\frac{1}{M}$ (Section~\ref{sec:constants-3}), one obtains
\begin{align*}
\sup_{\mathbb{D}\left(t_x/2\right)}\left|\left(\sum_{X\in D_{*}}\mathfrak{P}_{*}\left(X,\mathscr{C}(\cdot)\right)\right)_S\right|&\le2^{\#|S|}\cdot2\,C_d\cdot M^3\cdot\mathsf{p}\cdot K_0\cdot\exp\left(-2\left(\frac{d_{\mathrm{sep}}}{M}-2\right)\right)\\
&\le2^{\#|S|}\cdot2\exp(4)\cdot C_d\cdot M^3\cdot K_0\cdot\mathsf{p}\cdot\exp\left(-\vartheta\cdot d_{\mathrm{sep}}\right).
\end{align*}

\emph{(ii) The interacting part.} Choose cells $c\ne c'$ in $S$. The function $W$ is analytic on $\mathbb{D}^{\circ}\left(t_x\right)$ (Section~\ref{sec:construction-logarithm}), a neighbourhood of $\mathbb{D}\left(t_x/2\right)$. Lemma~\ref{lem:properties-mobius-components}(c) with $\rho=t_x/2$ gives
\[
\sup_{\mathbb{D}\left(t_x/2\right)}\left|W_S\right|\le2^{\#|S|-2}\cdot\left(\frac{t_x}{2}\right)^2\cdot\sup_{\mathbb{D}\left(t_x/2\right)}\left|\partial_c\partial_{c'}W\right|.
\]
Write $W=V_0+\left(W-V_0\right)$. By~\eqref{eq:four-jets-22}, $W-V_0=\sum_{n\ge1}\frac{1}{n!}\,\tilde{E}_n\left(V_{+},\dots,V_{+}\right)$ is the function $W$ of Lemma~\ref{lem:tree-expansion}(a) formed with the matrix $C=\mathcal{G}$ and with $V=V_{+}$; by Lemma~\ref{lem:inverses}(a) the data $\left(C,g,h_G,\boldsymbol{\alpha}_{\vartheta}\right)=\left(\mathcal{G},g,h_G,\boldsymbol{\alpha}\right)$ satisfy the requirements of Section~\ref{sec:weighted-tree-expansion} (Section~\ref{sec:construction-logarithm}). Put $U:=\mathbb{D}^{\circ}\left(t_x\right)$, an open set contained in $\mathbb{D}\left(t_x\right)$. On $U$ the coefficients of $V_{+}(t)$ are analytic (Section~\ref{sec:activities-gaussian-form}), and~\eqref{eq:four-jets-V1} and~\eqref{eq:four-jets-sm1} give
\[
\boldsymbol{\alpha}\cdot\rho_h\left(V_{+}(t);w_{\boldsymbol{\delta}}\right)\le\boldsymbol{\alpha}\cdot\bar{r}\le\frac{1}{2}\qquad\left(t\in U\right).
\]
Hence Lemma~\ref{lem:tree-expansion}(c) applies on $U$, with the subsets $P:=P_c\times\{1,2,3,4\}$ and $P':=P_{c'}\times\{1,2,3,4\}$ of $\mathcal{I}$, which are non-empty; their distance is $d(P,P')=d\left(P_c,P_{c'}\right)\ge d_{\mathrm{sep}}$, and the weights $w_P$, $w_{P'}$ are the weights $w_{P_c}$, $w_{P_{c'}}$ of Section~\ref{sec:activities-gaussian-form}. By~\eqref{eq:four-jets-11}, for $t\in\mathbb{D}\left(t_x/2\right)$,
\begin{align*}
\left|\partial_c\partial_{c'}\left(W-V_0\right)(t)\right|&\le2\,N_h\left(\partial_c\partial_{c'}V_{+}(t);1\right)\\
&\quad+2\,\boldsymbol{\alpha}\cdot N_h\left(\partial_cV_{+}(t);w_{P_c}\right)\cdot\rho_h\left(\partial_{c'}V_{+}(t);w_{P_{c'}}\right)\cdot\exp\left(-\vartheta\cdot d_{\mathrm{sep}}\right).
\end{align*}
By~\eqref{eq:four-jets-V3},
\begin{align*}
\left|\partial_c\partial_{c'}V_0(t)\right|+2\,N_h\left(\partial_c\partial_{c'}V_{+}(t);1\right)&\le2\left(N_h\left(\partial_c\partial_{c'}V_{+}(t);1\right)+\left|\partial_c\partial_{c'}V_0(t)\right|\right)\\
&\le2\,q_x^2\cdot C_V\cdot h_G^2\cdot\mathsf{p}\cdot\exp\left(-\vartheta\cdot\left(d_{\mathrm{sep}}-4M-6\right)\right),
\end{align*}
and~\eqref{eq:four-jets-V2}, applied to the cells $c$ and $c'$, gives
\[
N_h\left(\partial_cV_{+}(t);w_{P_c}\right)\le q_x\cdot C_V\cdot h_G^2\cdot\mathsf{p},\qquad\rho_h\left(\partial_{c'}V_{+}(t);w_{P_{c'}}\right)\le q_x\cdot C_V\cdot h_G^2.
\]
As $q_x=2/t_x$ (Section~\ref{sec:activities-gaussian-form}), one has $q_x\cdot t_x/2=1$. The last three estimates therefore give, for $t\in\mathbb{D}\left(t_x/2\right)$,
\begin{align*}
\left(\frac{t_x}{2}\right)^2\cdot\left|\partial_c\partial_{c'}W(t)\right|&\le2\,C_V\cdot h_G^2\cdot\mathsf{p}\cdot\exp\left(-\vartheta\cdot\left(d_{\mathrm{sep}}-4M-6\right)\right)\\
&\quad+2\,\boldsymbol{\alpha}\cdot C_V^2\cdot h_G^4\cdot\mathsf{p}\cdot\exp\left(-\vartheta\cdot d_{\mathrm{sep}}\right).
\end{align*}

\emph{(iii) Conclusion.} One has $\#|S|\le\#|\mathrm{Cel}|=N^{*}$ and $\exp\left(-\vartheta\cdot d_{\mathrm{sep}}\right)\le\exp\left(-\vartheta\cdot\left(d_{\mathrm{sep}}-4M-6\right)\right)$. Adding the bounds of~(i) and~(ii), one obtains
\begin{align*}
\sup_{\mathbb{D}\left(t_x/2\right)}\left|\mathrm{Lg}_S\right|&\le2^{N^{*}}\cdot\left(2\,C_V\cdot h_G^2+2\,\boldsymbol{\alpha}\cdot C_V^2\cdot h_G^4+2\exp(4)\cdot C_d\cdot M^3\cdot K_0\right)\\
&\qquad\cdot\mathsf{p}\cdot\exp\left(-\vartheta\cdot\left(d_{\mathrm{sep}}-4M-6\right)\right)\\
&\le\varpi_{\mathrm{sep}},
\end{align*}
because the bracket does not exceed $C_W=3\,C_V\cdot h_G^2+2\,\boldsymbol{\alpha}\cdot C_V^2\cdot h_G^4+4\exp(6)\cdot C_d\cdot M^3\cdot K_0$, and $\varpi_{\mathrm{sep}}=2^{N^{*}}\cdot C_W\cdot\mathsf{p}\cdot\exp\left(-\vartheta\cdot\left(d_{\mathrm{sep}}-4M-6\right)\right)$ (Section~\ref{sec:constants-3}).
\end{proof}

\subsection{Fields outside the small-field set}\label{sec:fields-outside-small}

\begin{lemma}[Gaussian bound for large fields]\label{lem:gaussian-bound-large}
For every $N\ge N_{\mathrm{min}}\left(R^{*},e\right)$,
\[
\bar{\mu}\left(F_N\setminus\mathfrak{G}^{-}\right)\le\varpi_0.
\]
\end{lemma}

\begin{proof}
By Definition~\ref{def:coordinates-slice-small}, a field $F\in F_N$ lies outside $\mathfrak{G}^{-}$ if and only if $\left|\zeta_j(\Upsilon)\right|>p_{0,j}-1$ for some $j<k^{*}$ and some basis element $\Upsilon$ of $V_j$, or $\left|\left(dA_{*}\right)(p)\right|>p_{*}-1$ for some plaquette $p$ of $T^0_{*}$. By~\eqref{it:trajectory-fermionic-operators-SL5}, $\zeta_j(\Upsilon)=\left\langle\phi_{j,\Upsilon},F\right\rangle$ with $\phi_{j,\Upsilon}\in F_N$ and $\left|\phi_{j,\Upsilon}\right|_2=1$, and $\left(dA_{*}\right)(p)=\left\langle\phi_p,F\right\rangle$ with $\phi_p\in F_N$ and $\left|\phi_p\right|_2\le1$.

For $0\le j<k^{*}$ put $i:=k^{*}-j\ge1$. By the definitions of Section~\ref{sec:running-quantities-tori}, $a_j=a_{*}\cdot L^{-i}$; hence $e_j=e_{*}\cdot L^{-\frac{i}{2}}$, $l_j=l_{*}\cdot L^i$, and
\[
v_j=-\log e_j=v_{*}+i\cdot\mathrm{lg},\qquad\mathrm{lg}=\frac{1}{2}\log L,\qquad n_j=l_j^3=l_{*}^3\cdot L^{3i}.
\]
By~\eqref{eq:renormalisation-group-T1}, $v_{*}\ge8$; and $p_0\ge1$, $p\ge1$ by stage~\ref{it:order-choice-8}. So the definitions of $p_{0,j}$ and $p_{*}$ in Section~\ref{sec:running-quantities-tori} give
\[
p_{0,j}=v_j^{p_0}\ge v_{*}^{p_0}=p_{0*}\ge v_{*}\ge8,\qquad p_{*}=v_{*}^p\ge v_{*}\ge8.
\]
In particular $p_{0,j}-1\ge0$ and $p_{*}-1\ge0$. The second assertion of Proposition~\ref{prop:sub-gaussian-bounds}, applied with these two numbers in the place of $y$, therefore gives
\begin{align*}
\bar{\mu}\left\{\left|\zeta_j(\Upsilon)\right|>p_{0,j}-1\right\}&\le2\exp\left(-\frac{\left(p_{0,j}-1\right)^2}{2}\right),\\
\bar{\mu}\left\{\left|\left(dA_{*}\right)(p)\right|>p_{*}-1\right\}&\le2\exp\left(-\frac{\left(p_{*}-1\right)^2}{2}\right)\le2\exp\left(-\frac{\left(p_{*}-1\right)^2}{2\,C_0^2}\right),
\end{align*}
the last inequality because $C_0\ge1$ (Section~\ref{sec:constants-3}).

For $x\ge3$ one has $(x-1)^2-x=x(x-3)+1>0$; therefore $\left(p_{0,j}-1\right)^2\ge p_{0,j}$. By~\eqref{it:trajectory-fermionic-operators-SL5}, the space $V_j$ has at most $2\,n_j$ basis elements, and $T^0_{*}$ has $3\,l_{*}^3$ plaquettes. Consequently
\[
\bar{\mu}\left(F_N\setminus\mathfrak{G}^{-}\right)\le\sum_{j<k^{*}}4\,n_j\cdot\exp\left(-\frac{p_{0,j}}{2}\right)+6\,l_{*}^3\cdot\exp\left(-\frac{\left(p_{*}-1\right)^2}{2\,C_0^2}\right).
\]
Since $p_0\ge1$ (stage~\ref{it:order-choice-8}), Bernoulli's inequality
\[
(v+x)^{p_0}\ge v^{p_0}+p_0\cdot v^{p_0-1}\cdot x\qquad\left(v>0,\ x\ge0\right)
\]
holds. Applied with $v=v_{*}$ and $x=i\cdot\mathrm{lg}$, together with the inequality $p_0\cdot v_{*}^{p_0-1}\ge28$ of~\eqref{eq:renormalisation-group-T1}, it gives
\[
p_{0,j}=\left(v_{*}+i\cdot\mathrm{lg}\right)^{p_0}\ge p_{0*}+28\,i\cdot\mathrm{lg}=p_{0*}+14\,i\cdot\log L,
\]
that is,
\[
\frac{p_{0,j}}{2}\ge\frac{p_{0*}}{2}+7\,i\cdot\log L.
\]
Therefore
\begin{align*}
\sum_{j<k^{*}}4\,n_j\cdot\exp\left(-\frac{p_{0,j}}{2}\right)&\le4\,l_{*}^3\cdot\exp\left(-\frac{p_{0*}}{2}\right)\cdot\sum_{i\ge1}L^{3i}\cdot L^{-7i}\\
&=4\,l_{*}^3\cdot\exp\left(-\frac{p_{0*}}{2}\right)\cdot\sum_{i\ge1}L^{-4i}\\
&\le4\,l_{*}^3\cdot\exp\left(-\frac{p_{0*}}{2}\right).
\end{align*}
Altogether
\[
\bar{\mu}\left(F_N\setminus\mathfrak{G}^{-}\right)\le4\,l_{*}^3\cdot\exp\left(-\frac{p_{0*}}{2}\right)+6\,l_{*}^3\cdot\exp\left(-\frac{\left(p_{*}-1\right)^2}{2\,C_0^2}\right)\le\varpi_0,
\]
by the definition of $\varpi_0$ in Section~\ref{sec:constants-3}.
\end{proof}

For $t\in\mathbb{C}^{\mathrm{Cel}}$ let
\[
\mathscr{Y}(t):=\bar{\nu}^{a(t)}\left(\mathfrak{G}\right),\qquad z_L(t):=\mathbf{Z}(t)-\mathscr{Y}(t)=\bar{\nu}^{a(t)}\left(F_N\setminus\mathfrak{G}\right).
\]
By Section~\ref{sec:averaged-law-partition}, applied to the Borel sets $\mathfrak{G}$ and $F_N\setminus\mathfrak{G}$, both functions are entire. One has $\mathbf{Z}(0)>0$ by Lemma~\ref{lem:signs-symmetries-wilson}(a). On $[-\pi,\pi]^3\times V_{\mathrm{t}}\left(T^0_{*}\right)$ let $\boldsymbol{\varphi}$ be the measure
\[
\boldsymbol{\varphi}\left(\mathrm{d}\theta\,\mathrm{d}A\right):=\mathbf{Z}(0)^{-1}\cdot\chi_{*}(A)\cdot Y^x_{*}\left(A+c_{\theta};0\right)\cdot(2\pi)^{-3}\,\mathrm{d}\theta\,\mathrm{d}\mathsf{m}_{*}(A),
\]
the density being read as $0$ where $\chi_{*}(A)=0$; where $\chi_{*}(A)=1$, the field $A+c_{\theta}$ lies in $\mathrm{Dom}_{*}$ by~\eqref{eq:four-jets-4}. The set $\left\{\chi_{*}=1\right\}$ of the $A\in V_{\mathrm{t}}\left(T^0_{*}\right)$ with $|dA|_{\infty}\le p_{*}$ (Definition~\ref{def:coordinates-slice-small}) is compact. Indeed, the coboundary $d$ is injective on $V_{\mathrm{t}}\left(T^0_{*}\right)$ by Section~\ref{sec:tori-cochains-hodge}, applied to the torus $T^0_{*}$; so $A\mapsto|dA|_{\infty}$ is a norm on this finite-dimensional real vector space, and $\left\{\chi_{*}=1\right\}$ is its closed ball of radius $p_{*}$, which is closed and bounded, all norms on such a space being equivalent.

\begin{lemma}[The small-field measure]\label{lem:small-field-measure}
Assume~\eqref{eq:four-jets-sm2}, \eqref{eq:four-jets-sm3} and~\eqref{eq:four-jets-sm4}.
\begin{enumerate}
\item[(a)] The measure $\boldsymbol{\varphi}$ is positive; its total mass is $\bar{\mu}\left(\mathfrak{G}\right)$, and $1-\varpi_0\le\bar{\mu}\left(\mathfrak{G}\right)\le1$.
\item[(b)] For $t\in\mathbb{D}^{\circ}\left(t_x^{+}\right)$,
\[
\frac{\mathscr{Y}(t)}{\mathbf{Z}(0)}=\int\boldsymbol{\varphi}\left(\mathrm{d}\theta\,\mathrm{d}A\right)\,\exp\mathrm{Lg}(t;\theta,A).
\]
\item[(c)] For every real $t\in\mathbb{D}^{\circ}\left(t_x^{+}\right)$, every $\theta\in\mathbb{R}^3$ and every $A\in V_{\mathrm{t}}\left(T^0_{*}\right)$ with $\chi_{*}(A)=1$, the number $\mathrm{Lg}(t;\theta,A)$ is real.
\end{enumerate}
\end{lemma}

\begin{proof}
\emph{(b)} Let $t\in\mathbb{D}^{\circ}\left(t_x^{+}\right)$. The placement is admitted by~\eqref{eq:four-jets-sm2}, so~\eqref{it:trajectory-fermionic-operators-SL7} gives $a(t)^{[j]}\in\mathfrak{X}_j$ for $0\le j\le k^{*}$; moreover $a(t)^{[k^{*}]}=\mathscr{C}(t)$. Hence~\eqref{eq:four-jets-6}, applied with $\Phi=1$ and $a=a(t)$, gives for every $\theta\in\mathbb{R}^3$
\[
\nu^{\theta,a(t)}\left(\mathfrak{G}\right)=\int\mathrm{d}\mathsf{m}_{*}(A)\,\chi_{*}(A)\cdot Y^x_{*}\left(A+c_{\theta};\mathscr{C}(t)\right).
\]
If $\chi_{*}(A)=1$, then~\eqref{eq:four-jets-20} and the identity $\exp\mathrm{Lg}=\mathrm{Rt}$ satisfied by the function $\mathrm{Lg}$ of~\eqref{eq:four-jets-21} (Section~\ref{sec:construction-logarithm}) give
\[
Y^x_{*}\left(A+c_{\theta};\mathscr{C}(t)\right)=\exp\mathrm{Lg}(t;\theta,A)\cdot Y^x_{*}\left(A+c_{\theta};0\right);
\]
both sides are continuous in $(\theta,A)$ by Section~\ref{sec:construction-logarithm}. So the function $(\theta,A)\mapsto\chi_{*}(A)\cdot Y^x_{*}\left(A+c_{\theta};\mathscr{C}(t)\right)$ on $[-\pi,\pi]^3\times V_{\mathrm{t}}\left(T^0_{*}\right)$, read as $0$ where $\chi_{*}(A)=0$, is continuous on the set $[-\pi,\pi]^3\times\left\{\chi_{*}=1\right\}$, which is compact because $\left\{\chi_{*}=1\right\}$ is compact, and it vanishes outside this set; consequently it is a bounded Borel function. The measure $(2\pi)^{-3}\,\mathrm{d}\theta\,\mathrm{d}\mathsf{m}_{*}(A)$ on $[-\pi,\pi]^3\times V_{\mathrm{t}}\left(T^0_{*}\right)$ is a probability measure, because $\mathrm{d}\mathsf{m}_{*}$ is the law of $A_{*}$ under $\gamma_N$ (Definition~\ref{def:coordinates-slice-small}). Formed with $t=0$, where $\mathscr{C}(0)=0$, the function above is $\mathbf{Z}(0)$ times the density of $\boldsymbol{\varphi}$ with respect to this probability measure; hence this density is a bounded Borel function, and $\boldsymbol{\varphi}$ is a finite measure. For general $t$, the identity displayed above for $\chi_{*}(A)=1$ shows that $\exp\mathrm{Lg}(t;\theta,A)$ times the density of $\boldsymbol{\varphi}$ is $\mathbf{Z}(0)^{-1}$ times the function above, both being read as $0$ where $\chi_{*}(A)=0$; therefore the integral in~(b) converges absolutely. Integrating over $\theta\in[-\pi,\pi]^3$ with respect to $(2\pi)^{-3}\,\mathrm{d}\theta$ and dividing by $\mathbf{Z}(0)$, one obtains~(b) from the definition of $\bar{\nu}^{a(t)}$ (Section~\ref{sec:averaged-law-partition}) and the definitions of $\mathscr{Y}$ and $\boldsymbol{\varphi}$: by Fubini's theorem, applied to the function above and to the probability measure $(2\pi)^{-3}\,\mathrm{d}\theta\,\mathrm{d}\mathsf{m}_{*}(A)$, the integral over $\theta$ of the right side of the first display equals the integral of this function with respect to $(2\pi)^{-3}\,\mathrm{d}\theta\,\mathrm{d}\mathsf{m}_{*}(A)$.

\emph{(a)} Let $\theta\in\mathbb{R}^3$. By Lemma~\ref{lem:signs-symmetries-wilson}(a), $\nu^{\theta}=\rho^{\theta,0}_N\,\gamma_N$ is a positive measure. So~\eqref{eq:four-jets-6} with $a=0$ gives
\[
\int\mathrm{d}\mathsf{m}_{*}(A)\,\Phi(A)\cdot\chi_{*}(A)\cdot Y^x_{*}\left(A+c_{\theta};0\right)=\nu^{\theta}\left(\mathbf{1}_{\mathfrak{G}}\cdot\Phi\right)\ge0
\]
for every bounded Borel function $\Phi\ge0$ on $V_{\mathrm{t}}\left(T^0_{*}\right)$. Consequently $\chi_{*}(A)\cdot Y^x_{*}\left(A+c_{\theta};0\right)\ge0$ for $\mathrm{d}\mathsf{m}_{*}$-almost every $A$. The density of $\boldsymbol{\varphi}$ being a Borel function (proof of~(b)), the set of the points $(\theta,A)\in[-\pi,\pi]^3\times V_{\mathrm{t}}\left(T^0_{*}\right)$ at which this density does not lie in $[0,\infty)$ is a Borel set; by what has just been shown, its section at every $\theta\in[-\pi,\pi]^3$ is a $\mathrm{d}\mathsf{m}_{*}$-null set. By Tonelli's theorem, applied to the indicator function of this set, it is a null set for the measure $(2\pi)^{-3}\,\mathrm{d}\theta\,\mathrm{d}\mathsf{m}_{*}(A)$. So the density of $\boldsymbol{\varphi}$ is non-negative almost everywhere, and the finite measure $\boldsymbol{\varphi}$ is positive. By~(b) at $t=0$, where $a(0)=0$ and $\mathrm{Lg}(0;\theta,A)=0$, the total mass of $\boldsymbol{\varphi}$ is
\[
\frac{\mathscr{Y}(0)}{\mathbf{Z}(0)}=\frac{\bar{\nu}\left(\mathfrak{G}\right)}{\bar{\nu}\left(F_N\right)}=\bar{\mu}\left(\mathfrak{G}\right).
\]
As $\bar{\mu}$ is a probability measure and $\mathfrak{G}^{-}\subset\mathfrak{G}$ (Definition~\ref{def:coordinates-slice-small}), Lemma~\ref{lem:gaussian-bound-large} gives
\[
1\ge\bar{\mu}\left(\mathfrak{G}\right)\ge\bar{\mu}\left(\mathfrak{G}^{-}\right)=1-\bar{\mu}\left(F_N\setminus\mathfrak{G}^{-}\right)\ge1-\varpi_0.
\]

\emph{(c)} Let $t\in\mathbb{D}^{\circ}\left(t_x^{+}\right)$ be real, and let $\theta\in\mathbb{R}^3$. Then $a(t)$ is a real bond function, so that $\nu^{\theta,a(t)}=\rho^{\theta,a(t)}_N\,\gamma_N$ is a positive measure by Lemma~\ref{lem:signs-symmetries-wilson}(a). The argument of~(a), with~\eqref{eq:four-jets-6} applied to $a=a(t)$ as in~(b), shows that $\chi_{*}(A)\cdot Y^x_{*}\left(A+c_{\theta};\mathscr{C}(t)\right)\ge0$ for $\mathrm{d}\mathsf{m}_{*}$-almost every $A$. By~\eqref{it:trajectory-fermionic-operators-SL5}, $\mathrm{d}\mathsf{m}_{*}$ is a non-degenerate centred Gaussian measure on $V_{\mathrm{t}}\left(T^0_{*}\right)$; it has an everywhere positive density with respect to the Lebesgue measure of this space, and so the inequality holds for Lebesgue-almost every $A$. The set
\[
\left\{\chi_{*}=1\right\}=\left\{A\in V_{\mathrm{t}}\left(T^0_{*}\right):\ |dA|_{\infty}\le p_{*}\right\}
\]
is convex and contains $0$ in its interior; hence it is the closure of its interior. On this set the function $A\mapsto Y^x_{*}\left(A+c_{\theta};\mathscr{C}(t)\right)$ is continuous (Section~\ref{sec:construction-logarithm}). A continuous function that takes values in $[0,\infty)$ almost everywhere on an open set takes values in $[0,\infty)$ everywhere on it. So $Y^x_{*}\left(A+c_{\theta};\mathscr{C}(t)\right)\ge0$ on the interior of $\left\{\chi_{*}=1\right\}$, and, by continuity, on $\left\{\chi_{*}=1\right\}$. By~\eqref{eq:four-jets-20} this number equals $\mathrm{Rt}(t;\theta,A)\cdot Y^x_{*}\left(A+c_{\theta};0\right)\ne0$. Hence
\[
Y^x_{*}\left(A+c_{\theta};\mathscr{C}(t)\right)>0\qquad\left(\chi_{*}(A)=1\right),
\]
and the same argument at $t=0$ gives $Y^x_{*}\left(A+c_{\theta};0\right)>0$ for $\chi_{*}(A)=1$. For $u\in[0,1]$ the point $u\cdot t$ is real and lies in $\mathbb{D}^{\circ}\left(t_x^{+}\right)$; therefore
\[
\mathrm{Rt}(u\cdot t;\theta,A)=\frac{Y^x_{*}\left(A+c_{\theta};\mathscr{C}(u\cdot t)\right)}{Y^x_{*}\left(A+c_{\theta};0\right)}>0\qquad\left(\chi_{*}(A)=1,\ u\in[0,1]\right).
\]
Let $\chi_{*}(A)=1$. Since $\exp\mathrm{Lg}(u\cdot t;\theta,A)=\mathrm{Rt}(u\cdot t;\theta,A)>0$, the number $\operatorname{Im}\mathrm{Lg}(u\cdot t;\theta,A)$ lies in $2\pi\mathbb{Z}$ for every $u\in[0,1]$. The function $u\mapsto\operatorname{Im}\mathrm{Lg}(u\cdot t;\theta,A)$ is continuous on $[0,1]$ and vanishes at $u=0$; so it vanishes identically. At $u=1$ this says that $\mathrm{Lg}(t;\theta,A)$ is real.
\end{proof}

\begin{proposition}[The large-field term and its jet]\label{prop:large-field-term}
Let conditions~\eqref{eq:four-jets-sm1}--\eqref{eq:four-jets-sm7} hold, and let $N\ge N_{\mathrm{min}}(R^{*},e)$.
\begin{enumerate}
\item[(a)] For every real point $t$ of $\mathbb{D}(t_2)$,
\[
0\le\frac{z_L(t)}{\mathbf{Z}(0)}\le\varpi_L.
\]
\item[(b)] For every $z\in\mathbb{C}^{\mathrm{Cel}}$,
\[
\mathring{\Phi}(e\cdot z)=\int\mathrm{d}\bar{\mu}(F)\,\exp\left(\sum_{c\in\mathrm{Cel}}\left(z_c\cdot X_c(F)-\frac{1}{2}\,z_c^2\cdot\boldsymbol{\sigma}^2\right)\right),
\]
and for every $t\in\mathbb{C}^{\mathrm{Cel}}$,
\[
\left|\mathring{\Phi}(t)\right|\le\exp\left(\frac{\boldsymbol{\sigma}^2}{2e^2}\sum_{c\in\mathrm{Cel}}\left(\operatorname{Im}t_c\right)^2\right).
\]
Moreover, for every $t\in\mathbb{D}(t_2)$ one has $\left|\mathscr{Y}(t)\right|\le2\,\mathbf{Z}(0)$, and $\left|\exp\mathrm{Lg}(t;\theta,A)\right|\le2$ for all $\theta\in\mathbb{R}^3$ and all $A\in V_{\mathrm{t}}\left(T^0_{*}\right)$ with $\chi_{*}(A)=1$.
\item[(c)] For every $t\in\mathbb{D}\left(\frac{t_2}{2N^{*}}\right)$,
\[
\left|\frac{z_L(t)}{\mathbf{Z}(0)}\right|\le\varpi_L^{\frac{1}{3}}\cdot M_L^{\frac{2}{3}}.
\]
\item[(d)] Let $\mathrm{Lr}\in\mathcal{J}$ be the element that corresponds to $z_L/\mathbf{Z}(0)$ under the identification of Section~\ref{sec:ring-jets}. Then $\|\mathrm{Lr}\|_{\mathbf{r}}\le\boldsymbol{\lambda}_L^{N^{*}}$, and every coefficient of $\mathrm{Lr}$ is a real number.
\end{enumerate}
\end{proposition}

\begin{proof}
The placement is admitted by~\eqref{eq:four-jets-sm2}. So Lemma~\ref{lem:properties-probes} applies to the probes, and $\boldsymbol{\sigma}^2$ is one of the numbers over which the supremum $\mathbf{D}$ of Section~\ref{sec:probes-mirror-copies} is taken. For $\theta\in\mathbb{R}^3$ let $\nu^{\theta}:=\nu^{\theta,0}$, as in Section~\ref{sec:lattice-measures-laws}. By Lemma~\ref{lem:signs-symmetries-wilson}(a) one has $\rho^{\theta,0}_N\ge0$ and $\mathbf{Z}(0)=\bar{\nu}(F_N)>0$; so each $\nu^{\theta}$ is a positive measure, and $\bar{\mu}=\bar{\nu}/\bar{\nu}(F_N)$ is a probability measure. Recall that $X_c(F)=\left\langle d\hat{\mathfrak{a}}_c,F\right\rangle$ for $F\in F_N$ (Section~\ref{sec:probes-mirror-copies}).

\emph{A change of variables.} Let $x\in\mathbb{R}^{\mathrm{Cel}}$, and put
\[
a':=a(e\cdot x)=\sum_{c\in\mathrm{Cel}}x_c\cdot\hat{\mathfrak{a}}_c,\qquad G:=da'=\sum_{c\in\mathrm{Cel}}x_c\cdot d\hat{\mathfrak{a}}_c,\qquad|x|_2^2:=\sum_{c\in\mathrm{Cel}}x_c^2.
\]
The two-cochain $G$ lies in $F_N=dC^1(\Lambda_{n_{*}})$. By Lemma~\ref{lem:properties-probes}(d) the two-cochains $d\hat{\mathfrak{a}}_c$, $c\in\mathrm{Cel}$, have pairwise disjoint supports and satisfy $|d\hat{\mathfrak{a}}_c|_2^2=\boldsymbol{\sigma}^2$; hence
\[
|G|_2^2=|x|_2^2\cdot\boldsymbol{\sigma}^2.
\]
By Lemma~\ref{lem:properties-probes}(a) every $\hat{\mathfrak{a}}_c$ lies in $V_{\mathrm{t}}$, and so does $a'$. Since $A(F)$ is the element of $V_{\mathrm{t}}$ with $dA(F)=F$ (Section~\ref{sec:averaged-law-partition}), the element $A(F)+a'$ of $V_{\mathrm{t}}$ has the coboundary $F+G$, that is, $A(F)+a'=A(F+G)$. The density $\rho^{\theta,a}_N(F)$ depends on $(F,a)$ only through $A(F)+a$; therefore
\[
\rho^{\theta,a'}_N(F)=\rho^{\theta,0}_N(F+G)\qquad\left(F\in F_N,\ \theta\in\mathbb{R}^3\right).
\]
The measure $\gamma_N$ is the standard Gaussian measure on the Euclidean space $F_N$, and $G\in F_N$. By the Cameron--Martin formula in finite dimension, which follows from the identity $|F'-G|_2^2=|F'|_2^2-2\langle F',G\rangle+|G|_2^2$ in the exponent of the Gaussian density,
\[
\int\mathrm{d}\gamma_N(F)\,f(F)=\int\mathrm{d}\gamma_N(F')\,\exp\left(\langle F',G\rangle-\frac{1}{2}|G|_2^2\right)\cdot f(F'-G)
\]
for every non-negative Borel function $f$ on $F_N$. Let $\operatorname{Ev}\subset F_N$ be a Borel set, let $\operatorname{Ev}{}+G:=\{F+G:\ F\in\operatorname{Ev}\}$, and apply this formula to $f(F):=\mathbf{1}_{\operatorname{Ev}}(F)\cdot\rho^{\theta,0}_N(F+G)$. Since $\mathbf{1}_{\operatorname{Ev}}(F'-G)=\mathbf{1}_{\operatorname{Ev}{}+G}(F')$ and $\langle F',G\rangle=\sum_cx_c\cdot X_c(F')$, one obtains, for every $\theta\in\mathbb{R}^3$ and every Borel set $\operatorname{Ev}\subset F_N$,
\begin{equation}\label{eq:four-jets-25}
\nu^{\theta,a(e\cdot x)}(\operatorname{Ev})=\exp\left(-\frac{1}{2}|x|_2^2\cdot\boldsymbol{\sigma}^2\right)\int_{\operatorname{Ev}{}+G}\mathrm{d}\nu^{\theta}(F)\,\exp\left(\sum_{c\in\mathrm{Cel}}x_c\cdot X_c(F)\right).
\end{equation}
This identity holds for every $x\in\mathbb{R}^{\mathrm{Cel}}$.

\emph{(a)} Let now, in addition, $|x_c|\le z_2$ for every $c\in\mathrm{Cel}$, and let $t:=e\cdot x$. Since $t_2=e\cdot z_2$, every real point of $\mathbb{D}(t_2)$ is of this form. By the definition of $z_2$ in Section~\ref{sec:constants-3} one has $N^{*}\cdot z_2^2\le1$ and $4N^{*}\cdot z_2^2\cdot\mathbf{D}\le1$, and $\boldsymbol{\sigma}^2\le\mathbf{D}$ by the definition of $\mathbf{D}$ in Section~\ref{sec:probes-mirror-copies}. Hence
\[
|x|_2^2\le N^{*}\cdot z_2^2\le1,\qquad|G|_2^2\le N^{*}\cdot z_2^2\cdot\mathbf{D}\le\frac{1}{4}.
\]
By~\eqref{it:trajectory-fermionic-operators-SL5}, each of the linear functionals $\zeta_j(\Upsilon)$, $j<k^{*}$, and $(dA_{*})(p)$, $p$ a plaquette of $T^0_{*}$, has the form $F\mapsto\langle\phi,F\rangle$ with a two-cochain $\phi$ such that $|\phi|_2\le1$. By the Cauchy--Schwarz inequality, its value at $G$ has modulus at most $|G|_2\le1$. In Definition~\ref{def:coordinates-slice-small} the set $\mathfrak{G}$ is defined by the bounds $p_{0,j}$ and $p_{*}$ on the moduli of these functionals, and the set $\mathfrak{G}^{-}$ by the bounds $p_{0,j}-1$ and $p_{*}-1$. Therefore $F'-G\in\mathfrak{G}$ for every $F'\in\mathfrak{G}^{-}$; equivalently,
\[
\left(F_N\setminus\mathfrak{G}\right)+G\subset F_N\setminus\mathfrak{G}^{-}.
\]
By the definitions of $z_L$ (Section~\ref{sec:fields-outside-small}) and of $\bar{\nu}^a$, the number $z_L(t)=\bar{\nu}^{a(e\cdot x)}\left(F_N\setminus\mathfrak{G}\right)$ is the mean of $\nu^{\theta,a(e\cdot x)}\left(F_N\setminus\mathfrak{G}\right)$ over $\theta\in[-\pi,\pi]^3$ with respect to $\frac{\mathrm{d}\theta}{(2\pi)^3}$. The identity~\eqref{eq:four-jets-25} with $\operatorname{Ev}=F_N\setminus\mathfrak{G}$, the positivity of $\nu^{\theta}$ and the last inclusion give
\[
0\le\nu^{\theta,a(e\cdot x)}\left(F_N\setminus\mathfrak{G}\right)\le\exp\left(-\frac{1}{2}|x|_2^2\cdot\boldsymbol{\sigma}^2\right)\int\mathrm{d}\nu^{\theta}(F)\,\mathbf{1}_{F_N\setminus\mathfrak{G}^{-}}(F)\cdot\exp\langle G,F\rangle.
\]
Taking the mean over $\theta$ and dividing by $\mathbf{Z}(0)=\bar{\nu}(F_N)$, one finds
\[
0\le\frac{z_L(t)}{\mathbf{Z}(0)}\le\exp\left(-\frac{1}{2}|x|_2^2\cdot\boldsymbol{\sigma}^2\right)\int\mathrm{d}\bar{\mu}(F)\,\mathbf{1}_{F_N\setminus\mathfrak{G}^{-}}(F)\cdot\exp\langle G,F\rangle.
\]
By the Cauchy--Schwarz inequality the last integral is at most
\[
\bar{\mu}\left(F_N\setminus\mathfrak{G}^{-}\right)^{\frac{1}{2}}\cdot\left(\int\mathrm{d}\bar{\mu}(F)\,\exp\langle2G,F\rangle\right)^{\frac{1}{2}}\le\varpi_0^{\frac{1}{2}}\cdot\exp\left(|G|_2^2\right);
\]
here the first factor is bounded by Lemma~\ref{lem:gaussian-bound-large}, and the second by Proposition~\ref{prop:sub-gaussian-bounds}, applied to the two-cochain $2G$, which lies in $F_N$ and is therefore fixed by the orthogonal projection onto $F_N$. As $|G|_2^2=|x|_2^2\cdot\boldsymbol{\sigma}^2\le\mathbf{D}$, it follows that
\[
0\le\frac{z_L(t)}{\mathbf{Z}(0)}\le\varpi_0^{\frac{1}{2}}\cdot\exp\left(\frac{1}{2}|x|_2^2\cdot\boldsymbol{\sigma}^2\right)\le\varpi_0^{\frac{1}{2}}\cdot\exp\left(\frac{\mathbf{D}}{2}\right)=\varpi_L.
\]

\emph{(b)} For $\operatorname{Ev}=F_N$ one has $\operatorname{Ev}{}+G=F_N$. Taking in~\eqref{eq:four-jets-25} the mean over $\theta\in[-\pi,\pi]^3$ and dividing by $\mathbf{Z}(0)$, one obtains for every $x\in\mathbb{R}^{\mathrm{Cel}}$
\[
\mathring{\Phi}(e\cdot x)=\frac{\mathbf{Z}(e\cdot x)}{\mathbf{Z}(0)}=\int\mathrm{d}\bar{\mu}(F)\,\exp\left(\sum_{c\in\mathrm{Cel}}\left(x_c\cdot X_c(F)-\frac{1}{2}\,x_c^2\cdot\boldsymbol{\sigma}^2\right)\right).
\]
The left side is the restriction to $\mathbb{R}^{\mathrm{Cel}}$ of the entire function $z\mapsto\mathring{\Phi}(e\cdot z)$ (Section~\ref{sec:averaged-law-partition}). The right side is the product of the entire function $z\mapsto\exp\left(-\frac{1}{2}\,\boldsymbol{\sigma}^2\sum_cz_c^2\right)$ and of $z\mapsto\int\mathrm{d}\bar{\mu}(F)\,\exp\left(\sum_cz_c\cdot X_c(F)\right)$, taken at $z=x$, and the second factor is entire as well. Indeed, by Proposition~\ref{prop:sub-gaussian-bounds} the function $\exp\left(\sum_cy_c\cdot X_c\right)$ is $\bar{\mu}$-integrable for every $y\in\mathbb{R}^{\mathrm{Cel}}$. Let $\lambda>0$ be arbitrary. If $z\in\mathbb{C}^{\mathrm{Cel}}$ satisfies $|\operatorname{Re}z_c|\le\lambda$ for every $c\in\mathrm{Cel}$, then the modulus of $\exp\left(\sum_cz_c\cdot X_c\right)$ is at most $\prod_c\left(\exp\left(\lambda\cdot X_c\right)+\exp\left(-\lambda\cdot X_c\right)\right)$, and this product, multiplied out, is a finite sum of such integrable functions. Dominated convergence and Morera's theorem therefore show that the second factor is analytic on the set of the $z$ with $|\operatorname{Re}z_c|<\lambda$ for all $c\in\mathrm{Cel}$; as $\lambda$ is arbitrary, it is entire. By the identity theorem the formula holds with any $z\in\mathbb{C}^{\mathrm{Cel}}$ in place of $x$; this is the first assertion of (b). Since $\bar{\mu}$ is positive,
\begin{align*}
\left|\mathring{\Phi}(e\cdot z)\right|&\le\exp\left(-\frac{\boldsymbol{\sigma}^2}{2}\sum_{c\in\mathrm{Cel}}\operatorname{Re}\left(z_c^2\right)\right)\\
&\qquad\cdot\int\mathrm{d}\bar{\mu}(F)\,\exp\left(\sum_{c\in\mathrm{Cel}}\operatorname{Re}(z_c)\cdot X_c(F)\right)\qquad\left(z\in\mathbb{C}^{\mathrm{Cel}}\right).
\end{align*}
The real two-cochain $\sum_c\operatorname{Re}(z_c)\cdot d\hat{\mathfrak{a}}_c$ lies in $F_N$, and by Lemma~\ref{lem:properties-probes}(d) its squared norm is $\boldsymbol{\sigma}^2\sum_c\left(\operatorname{Re}z_c\right)^2$. So Proposition~\ref{prop:sub-gaussian-bounds} bounds the integral by $\exp\left(\frac{1}{2}\,\boldsymbol{\sigma}^2\sum_c\left(\operatorname{Re}z_c\right)^2\right)$. As $\left(\operatorname{Re}z_c\right)^2-\operatorname{Re}\left(z_c^2\right)=\left(\operatorname{Im}z_c\right)^2$, one obtains
\[
\left|\mathring{\Phi}(e\cdot z)\right|\le\exp\left(\frac{\boldsymbol{\sigma}^2}{2}\sum_{c\in\mathrm{Cel}}\left(\operatorname{Im}z_c\right)^2\right)\qquad\left(z\in\mathbb{C}^{\mathrm{Cel}}\right),
\]
which is the second assertion, with $t=e\cdot z$.

Let $t\in\mathbb{D}(t_2)$, $\theta\in\mathbb{R}^3$ and $A\in V_{\mathrm{t}}\left(T^0_{*}\right)$ with $\chi_{*}(A)=1$, and write $\mathrm{Lg}(t)$ for $\mathrm{Lg}(t;\theta,A)$. By~\eqref{eq:four-jets-sm7}, $t_2\le\frac{t_x}{8N^{*}}$; hence $\mathbb{D}(t_2)$ is contained in $\mathbb{D}(t_x/2)$ and in $\mathbb{D}^{\circ}\left(t_x^{+}\right)$, on which $\mathrm{Lg}$ is analytic (Section~\ref{sec:construction-logarithm}). As $\mathrm{Lg}(0)=0$, Lemma~\ref{lem:properties-mobius-components}(a) gives
\[
\mathrm{Lg}=\sum_{\emptyset\ne S\subset\mathrm{Cel}}\mathrm{Lg}_S,\qquad\mathrm{Lg}_{\{c\}}(t)=\mathrm{Lg}(t_c\cdot e_c)=\mathrm{lgN}(t_c)+w_c(t_c),
\]
the last equality being the definition of $w_c$ in Proposition~\ref{prop:comparison-free-ratio}. The function $\mathrm{lgN}$ is even, vanishes at $0$ and is bounded by $B_0\le\frac{1}{4}$ on the disk $|t_c|\le t_x$ (Lemma~\ref{lem:free-wilson-ratio}(d)); the function $w_c$ vanishes at $0$ and is bounded by $3e_{*}^{\frac{1}{4}-9\varepsilon}$ on that disk (Proposition~\ref{prop:comparison-free-ratio}). Schwarz's lemma, the inequality $t_2/t_x\le1/(8N^{*})$ and the bound $e_{*}^{\frac{1}{4}-9\varepsilon}\le\frac{1}{8}$ of~\eqref{eq:four-jets-sm3} give
\begin{align*}
\left|\mathrm{lgN}(t_c)\right|&\le B_0\cdot\left(\frac{t_2}{t_x}\right)^2\le\frac{1}{256\,(N^{*})^2},\\
\left|w_c(t_c)\right|&\le3e_{*}^{\frac{1}{4}-9\varepsilon}\cdot\frac{t_2}{t_x}\le\frac{3}{64\,N^{*}},
\end{align*}
and therefore
\[
\sum_{c\in\mathrm{Cel}}\left|\mathrm{Lg}_{\{c\}}(t)\right|\le\frac{1}{256\,N^{*}}+\frac{3}{64}.
\]
There are fewer than $2^{N^{*}}$ subsets $S\subset\mathrm{Cel}$ with $\#|S|\ge2$, and Proposition~\ref{prop:bound-mobius-components} and~\eqref{eq:four-jets-sm5} give
\[
\sum_{\#|S|\ge2}\left|\mathrm{Lg}_S(t)\right|\le2^{N^{*}}\cdot\varpi_{\mathrm{sep}}\le\frac{1}{2}\left(2\,\bar{\kappa}\right)^{N^{*}}.
\]
By~\eqref{eq:four-jets-sm6}, $\bar{\kappa}^{1-s}\le\frac{1}{2N^{*}}<1$; hence $\bar{\kappa}<1$, and since $0<s<1$ (stage~\ref{it:order-choice-8}), $\bar{\kappa}\le\bar{\kappa}^{1-s}\le\frac{1}{2N^{*}}$. Thus $\frac{1}{2}\left(2\,\bar{\kappa}\right)^{N^{*}}\le2^{-N^{*}}$, and
\[
\left|\mathrm{Lg}(t)\right|\le\frac{1}{256\,N^{*}}+\frac{3}{64}+2^{-N^{*}}\le\frac{1}{2},\qquad\left|\exp\mathrm{Lg}(t)\right|\le\exp\left(\frac{1}{2}\right)\le2.
\]
By Lemma~\ref{lem:small-field-measure}(b), $\mathscr{Y}(t)/\mathbf{Z}(0)$ is the integral of $\exp\mathrm{Lg}(t;\theta,A)$ with respect to the measure $\boldsymbol{\varphi}$, which is positive and has total mass at most $1$ by Lemma~\ref{lem:small-field-measure}(a). Hence $\left|\mathscr{Y}(t)\right|\le2\,\mathbf{Z}(0)$.

\emph{(c)} Let $t\in\mathbb{D}(t_2)$. Then $\sum_c\left(\operatorname{Im}t_c\right)^2\le N^{*}\cdot t_2^2=e^2\cdot N^{*}\cdot z_2^2$, and (b), together with $N^{*}\cdot z_2^2\le1$ and $\boldsymbol{\sigma}^2\le\mathbf{D}$, gives
\[
\left|\mathring{\Phi}(t)\right|\le\exp\left(\frac{1}{2}N^{*}\cdot z_2^2\cdot\boldsymbol{\sigma}^2\right)\le\exp\left(\frac{\mathbf{D}}{2}\right).
\]
Since $z_L=\mathbf{Z}-\mathscr{Y}$ (Section~\ref{sec:fields-outside-small}), the bound for $\mathscr{Y}$ in (b) yields
\[
\left|\frac{z_L(t)}{\mathbf{Z}(0)}\right|\le\left|\mathring{\Phi}(t)\right|+\frac{\left|\mathscr{Y}(t)\right|}{\mathbf{Z}(0)}\le\exp\left(\frac{\mathbf{D}}{2}\right)+2=M_L.
\]
At the real points of $\mathbb{D}(t_2)$, (a) and the bound $\varpi_0\le\frac{1}{4}$ of~\eqref{eq:four-jets-sm7} give $\left|z_L(t)/\mathbf{Z}(0)\right|\le\varpi_L\le\exp\left(\frac{\mathbf{D}}{2}\right)\le M_L$. The function $z_L/\mathbf{Z}(0)$ is entire (Section~\ref{sec:fields-outside-small}). Lemma~\ref{lem:two-constants-estimate}, applied with $n=N^{*}$, $\rho=t_2$, $M_{\mathfrak{g}}=M_L$ and $\boldsymbol{\delta}=\varpi_L$, gives (c).

\emph{(d)} For a multi-index $\mathbf{a}\in\{0,1,2,3,4\}^{\mathrm{Cel}}$ let $\mathrm{Lr}_{\mathbf{a}}$ be the coefficient of $t^{\mathbf{a}}$ in $\mathrm{Lr}$, that is, the coefficient of $t^{\mathbf{a}}$ in the Taylor series of $z_L/\mathbf{Z}(0)$ at $0$. By (c) and Cauchy's estimates on the polydisk $\mathbb{D}\left(\frac{t_2}{2N^{*}}\right)$,
\[
\left|\mathrm{Lr}_{\mathbf{a}}\right|\le\varpi_L^{\frac{1}{3}}\cdot M_L^{\frac{2}{3}}\cdot\left(\frac{2N^{*}}{t_2}\right)^{|\mathbf{a}|}.
\]
Put $y:=2N^{*}\cdot\mathbf{r}/t_2$. As $\mathbf{r}=t_x/4$, condition~\eqref{eq:four-jets-sm7} gives $\mathbf{r}/t_2=t_x/(4t_2)\ge2N^{*}$; in particular $y\ge1$. The sum of $y^{|\mathbf{a}|}$ over all multi-indices $\mathbf{a}\in\{0,1,2,3,4\}^{\mathrm{Cel}}$ equals $\left(\sum_{i=0}^4y^i\right)^{N^{*}}$. With the norm of Section~\ref{sec:ring-jets} it follows that
\begin{align*}
\|\mathrm{Lr}\|_{\mathbf{r}}=\sum_{\mathbf{a}}\left|\mathrm{Lr}_{\mathbf{a}}\right|\cdot\mathbf{r}^{|\mathbf{a}|}&\le\varpi_L^{\frac{1}{3}}\cdot M_L^{\frac{2}{3}}\cdot\left(\sum_{i=0}^4y^i\right)^{N^{*}}\\
&\le\varpi_L^{\frac{1}{3}}\cdot M_L^{\frac{2}{3}}\cdot\left(5y^4\right)^{N^{*}}=\boldsymbol{\lambda}_L^{N^{*}},
\end{align*}
the last equality being the definition of $\boldsymbol{\lambda}_L$ in Section~\ref{sec:constants-3}. By (a), $z_L/\mathbf{Z}(0)$ is real at the real points of $\mathbb{D}(t_2)$; hence its Taylor coefficients at $0$ are real, and so are the coefficients of $\mathrm{Lr}$.
\end{proof}

\subsection{The ratio of partition functions in the ring of jets}\label{sec:ratio-partition-functions}

Let $\theta\in\mathbb{R}^3$ and $A\in V_{\mathrm{t}}\left(T^0_{*}\right)$ with $\chi_{*}(A)=1$. For $c\in\mathrm{Cel}$ let $w_c$ be the function of the variable $t_c$ defined in Proposition~\ref{prop:comparison-free-ratio}, and for $S\subset\mathrm{Cel}$ let $\mathrm{Lg}_S$ be the M\"obius component of index $S$, formed as in Section~\ref{sec:mobius-decomposition-cell}, of the function $t\mapsto\mathrm{Lg}(t;\theta,A)$ of~\eqref{eq:four-jets-21}; both are formed with the pair $(\theta,A)$. Functions that are analytic near $0$ being identified with their images in $\mathcal{J}$ (Section~\ref{sec:ring-jets}), define elements of $\mathcal{J}$ by
\[
k_{\{c\}}(\theta,A;\cdot):=\exp\left(w_c\right)-1\quad\left(c\in\mathrm{Cel}\right),\qquad k_S(\theta,A;\cdot):=\exp\left(\mathrm{Lg}_S\right)-1\quad\left(S\subset\mathrm{Cel},\ \#|S|\ge2\right).
\]

\begin{theorem}[Factorised representation]\label{thm:factorised-representation}
Suppose that hypothesis~\ref{def:hypothesis-reference-torus} and conditions~\eqref{eq:four-jets-sm1}--\eqref{eq:four-jets-sm7} are satisfied, the probes $\hat{\mathfrak{a}}_c$, $c\in\mathrm{Cel}$, being those of Section~\ref{sec:probes-mirror-copies}, and let $N\ge N_{\mathrm{min}}(R^{*},e)$. Then the following identity holds in $\mathcal{J}$:
\begin{equation}\label{eq:four-jets-26}
\mathring{\Phi}(t)=\prod_{c\in\mathrm{Cel}}\phi_0(t_c)\cdot\int\boldsymbol{\varphi}(\mathrm{d}\theta\,\mathrm{d}A_{*})\,\prod_{\emptyset\ne S\subset\mathrm{Cel}}\left(1+k_S(\theta,A_{*};t)\right)+\mathrm{Lr}(t),
\end{equation}
where the integral of a function with values in $\mathcal{J}$ is formed coefficientwise, and where the following hold.
\begin{enumerate}
\item[(a)] $\phi_0$ is the Taylor polynomial of degree four at $0$ of the free Wilson ratio $N_0$ of Section~\ref{sec:averaged-law-partition}; it is even, its coefficients are real, $\phi_0(0)=1$, and the relations~\eqref{eq:four-jets-23} hold.
\item[(b)] The measure $\boldsymbol{\varphi}$ of Section~\ref{sec:fields-outside-small} is positive, and its total mass $\bar{\mu}(\mathfrak{G})$ lies in $[1-\varpi_0,1]$.
\item[(c)] For every non-empty $S\subset\mathrm{Cel}$ the coefficients of $k_S(\theta,A_{*};\cdot)$ are real and depend continuously on $(\theta,A_{*})$, and every monomial with a non-zero coefficient in $k_S(\theta,A_{*};\cdot)$ contains each of the variables $t_c$, $c\in S$, and no other variable.
\item[(d)] For every $c\in\mathrm{Cel}$,
\[
\sup_{\theta,A_{*}}\left\|k_{\{c\}}(\theta,A_{*};\cdot)\right\|_{\mathbf{r}}\le\bar{\kappa}=6\,e_{*}^{\frac{1}{4}-9\varepsilon}.
\]
\item[(e)] For every $S\subset\mathrm{Cel}$ with $\#|S|\ge2$,
\[
\sup_{\theta,A_{*}}\left\|k_S(\theta,A_{*};\cdot)\right\|_{\mathbf{r}}\le2\,\varpi_{\mathrm{sep}}\le\bar{\kappa}^{N^{*}}.
\]
\item[(f)] $\mathrm{Lr}$ is the element of $\mathcal{J}$ of Proposition~\ref{prop:large-field-term}(d); its coefficients are real, and $\|\mathrm{Lr}\|_{\mathbf{r}}\le\boldsymbol{\lambda}_L^{N^{*}}$.
\item[(g)] One has
\[
\mathbf{r}=\frac{t_x}{4}=\frac{1}{4}\,e_{*}^{3\varepsilon}\cdot\left(1+h'^{\frac{1}{2}}\cdot C_{\beta}\right)^{-1}.
\]
\end{enumerate}
In (d) and (e) the suprema are taken over $\theta\in\mathbb{R}^3$ and over the $A_{*}\in V_{\mathrm{t}}\left(T^0_{*}\right)$ with $\chi_{*}(A_{*})=1$. When $N\ge N_{\mathrm{min}}(R^{*},e)$, $\theta\in\mathbb{R}^3$ and the admitted placement vary, and when the push-forward $r_{*}\beta$ of $\beta$ under a diagonal sign matrix $r$ (Section~\ref{sec:probes-mirror-copies}) is put in place of $\beta$, the right-hand sides of the inequalities in~\eqref{eq:four-jets-23} and the numbers $\varpi_0$, $\bar{\kappa}$, $\bar{\kappa}^{N^{*}}$, $\boldsymbol{\lambda}_L^{N^{*}}$, $\mathbf{r}$ keep their values, while $\varpi_{\mathrm{sep}}$ is given by one and the same function of $d_{\mathrm{sep}}$ throughout; moreover,
\[
\varpi_{\mathrm{sep}}\le2^{N^{*}}\cdot C_W\cdot\mathsf{p}\cdot\exp\left(-\vartheta\cdot\left(\frac{M\cdot r_{*}}{8L}-4M-7\right)\right).
\]
\end{theorem}

\begin{proof}
\emph{The identity~\eqref{eq:four-jets-26}.} By the definitions in Section~\ref{sec:fields-outside-small} one has $\mathbf{Z}=\mathscr{Y}+z_L$; hence
\[
\mathring{\Phi}=\frac{\mathscr{Y}}{\mathbf{Z}(0)}+\frac{z_L}{\mathbf{Z}(0)},
\]
and the image of the second term in $\mathcal{J}$ is $\mathrm{Lr}$. By Lemma~\ref{lem:small-field-measure}(b),
\[
\frac{\mathscr{Y}(t)}{\mathbf{Z}(0)}=\int\boldsymbol{\varphi}(\mathrm{d}\theta\,\mathrm{d}A_{*})\,\exp\mathrm{Lg}(t;\theta,A_{*})\qquad\left(t\in\mathbb{D}^{\circ}\left(t_x^{+}\right)\right),
\]
and by Proposition~\ref{prop:large-field-term}(b) the integrand is bounded by $2$ on $\mathbb{D}(t_2)$. By Cauchy's formula, the Taylor coefficients at $0$ of a function that is analytic in a neighbourhood of $\mathbb{D}(t_2)$ are integrals over the torus $\left\{t:\ |t_c|=t_2\text{ for all }c\in\mathrm{Cel}\right\}$. The function $\exp\mathrm{Lg}$ is continuous in $(t,\theta,A_{*})$ (Section~\ref{sec:construction-logarithm}) and bounded by $2$ on this torus, and $\boldsymbol{\varphi}$ is a finite positive measure (Lemma~\ref{lem:small-field-measure}(a)). By Fubini's theorem, every Taylor coefficient of $\mathscr{Y}/\mathbf{Z}(0)$ at $0$ is therefore the $\boldsymbol{\varphi}$-integral of the corresponding Taylor coefficient of $\exp\mathrm{Lg}(\cdot;\theta,A_{*})$; that is, the image of $\mathscr{Y}/\mathbf{Z}(0)$ in $\mathcal{J}$ is the coefficientwise integral of the images of the functions $\exp\mathrm{Lg}(\cdot;\theta,A_{*})$.

Fix $(\theta,A_{*})$ and write $\mathrm{Lg}$ for $\mathrm{Lg}(\cdot;\theta,A_{*})$. Since $\mathrm{Lg}(0)=0$, Lemma~\ref{lem:properties-mobius-components}(a) and the definition of $w_c$ in Proposition~\ref{prop:comparison-free-ratio} give
\[
\mathrm{Lg}(t)=\sum_{c\in\mathrm{Cel}}\left(\mathrm{lgN}(t_c)+w_c(t_c)\right)+\sum_{\#|S|\ge2}\mathrm{Lg}_S(t),
\]
and $\exp\mathrm{lgN}(t_c)=N_0(t_c)$ by Lemma~\ref{lem:free-wilson-ratio}(a). The identification of Section~\ref{sec:ring-jets} is a ring homomorphism that carries $\exp(\mathfrak{g})$ to the exponential series of the image of $\mathfrak{g}$, and it carries the function $t\mapsto N_0(t_c)$ to $\phi_0(t_c)$. Consequently, in $\mathcal{J}$,
\begin{align*}
\exp\mathrm{Lg}&=\prod_{c\in\mathrm{Cel}}\phi_0(t_c)\cdot\prod_{c\in\mathrm{Cel}}\left(1+k_{\{c\}}(\theta,A_{*};t)\right)\cdot\prod_{\#|S|\ge2}\left(1+k_S(\theta,A_{*};t)\right)\\
&=\prod_{c\in\mathrm{Cel}}\phi_0(t_c)\cdot\prod_{\emptyset\ne S\subset\mathrm{Cel}}\left(1+k_S(\theta,A_{*};t)\right).
\end{align*}
The factor $\prod_c\phi_0(t_c)$ does not depend on $(\theta,A_{*})$. Integrating with respect to $\boldsymbol{\varphi}$ and adding $\mathrm{Lr}$, one obtains~\eqref{eq:four-jets-26}.

\emph{(a)} By Lemma~\ref{lem:free-wilson-ratio}(e), $\phi_0(t)=\sum_{i=0}^4n_i\cdot t^i$ with real coefficients $n_i$, $n_0=1$ and $n_1=n_3=0$, and~\eqref{eq:four-jets-23} holds. \emph{(b)} is Lemma~\ref{lem:small-field-measure}(a).

\emph{(c)} Let $S\subset\mathrm{Cel}$ with $\#|S|\ge2$. Since $t_x/2<t_x^{+}$, the closed polydisk $\mathbb{D}(t_x/2)$ is contained in the open polydisk $\mathbb{D}^{\circ}\left(t_x^{+}\right)$, where $\mathrm{Lg}$ is analytic. Lemma~\ref{lem:properties-mobius-components}(d), applied to $\mathrm{Lg}$ with $\rho=t_x/2$ and $\rho'=\mathbf{r}$, therefore shows that every monomial with a non-zero coefficient in the image of $\mathrm{Lg}_S$ contains each $t_c$, $c\in S$, and no other variable. A product of such monomials has the same property, and therefore so has $k_S=\sum_{n\ge1}\frac{1}{n!}\left(\mathrm{Lg}_S\right)^n$. The function $w_c$ depends on $t_c$ alone and vanishes at $0$; so the image of $w_c$, and with it $k_{\{c\}}=\exp(w_c)-1$, is a polynomial in $t_c$ without constant term. Section~\ref{sec:construction-logarithm} shows that the map $(t,\theta,A_{*})\mapsto\mathrm{Lg}(t;\theta,A_{*})$ is continuous and that each partial map $\mathrm{Lg}(\cdot;\theta,A_{*})$ is analytic on $\mathbb{D}^{\circ}\left(t_x^{+}\right)$. Fix $\rho$ with $0<\rho<t_x^{+}$. By Cauchy's formula on the torus $\left\{t:\ |t_c|=\rho\text{ for all }c\right\}$, each Taylor coefficient of $\mathrm{Lg}(\cdot;\theta,A_{*})$ at $0$ is the integral over this torus of $\mathrm{Lg}(\cdot;\theta,A_{*})$ multiplied by a continuous function of $t$ that does not involve $(\theta,A_{*})$; hence these coefficients depend continuously on $(\theta,A_{*})$. The coefficients of $\mathrm{Lg}_S$ are linear combinations of them; those of $w_c$ are differences of some of them and of the coefficients of $\mathrm{lgN}$, which do not depend on $(\theta,A_{*})$; and the coefficients of every $k_S$ are polynomials in these. This proves the continuity. Finally, by Lemma~\ref{lem:small-field-measure}(c), $\mathrm{Lg}(t;\theta,A_{*})$ is real for real $t\in\mathbb{D}^{\circ}\left(t_x^{+}\right)$, so that its Taylor coefficients at $0$ are real; hence so are those of each function $t\mapsto\mathrm{Lg}\left(t^T;\theta,A_{*}\right)$, $T\subset\mathrm{Cel}$, and those of the M\"obius components $\mathrm{Lg}_S$. The coefficients of $\mathrm{lgN}$ are real because $\mathrm{lgN}$ is real at real points (Lemma~\ref{lem:free-wilson-ratio}(d)). Therefore the coefficients of $w_c$, of $\mathrm{Lg}_S$ and of every $k_S$ are real.

\emph{(d)} Let $c\in\mathrm{Cel}$. By Proposition~\ref{prop:comparison-free-ratio}, $w_c(0)=0$, and $w_c$ is analytic on the open disk $|t_c|<t_x^{+}$, which contains the closed disk $|t_c|\le t_x$ because $t_x<t_x^{+}$. Cauchy's estimates on the closed disk, used as in the proof of Lemma~\ref{lem:properties-mobius-components}(d) with $\rho=t_x$ and $\rho'=\mathbf{r}\le\rho/2$, give $\|w_c\|_{\mathbf{r}}\le\sup_{|t_c|\le t_x}|w_c(t_c)|$, and the bound of Proposition~\ref{prop:comparison-free-ratio} together with~\eqref{eq:four-jets-sm3} gives
\[
\|w_c\|_{\mathbf{r}}\le3\,e_{*}^{\frac{1}{4}-9\varepsilon}\le\frac{3}{8}.
\]
The norm $\|\cdot\|_{\mathbf{r}}$ is submultiplicative (Section~\ref{sec:ring-jets}), and $\exp(y)-1\le2y$ for $0\le y\le\frac{1}{2}$. Hence, for all $(\theta,A_{*})$,
\[
\left\|k_{\{c\}}(\theta,A_{*};\cdot)\right\|_{\mathbf{r}}=\left\|\sum_{n\ge1}\frac{1}{n!}\,w_c^n\right\|_{\mathbf{r}}\le\exp\left(\|w_c\|_{\mathbf{r}}\right)-1\le2\,\|w_c\|_{\mathbf{r}}\le6\,e_{*}^{\frac{1}{4}-9\varepsilon}=\bar{\kappa}.
\]

\emph{(e)} Let $S\subset\mathrm{Cel}$ with $\#|S|\ge2$. As shown in the proof of (c), $\mathrm{Lg}$ is analytic in a neighbourhood of $\mathbb{D}(t_x/2)$. Since $\mathbf{r}=t_x/4$ is one half of $t_x/2$, the inequality of Lemma~\ref{lem:properties-mobius-components}(d) is available with $\rho=t_x/2$ and $\rho'=\mathbf{r}$; combined with the bound of Proposition~\ref{prop:bound-mobius-components}, it yields
\[
\|\mathrm{Lg}_S\|_{\mathbf{r}}\le\sup_{\mathbb{D}(t_x/2)}\left|\mathrm{Lg}_S\right|\le\varpi_{\mathrm{sep}}.
\]
By~\eqref{eq:four-jets-sm3} one has $\bar{\kappa}\le\frac{3}{4}$, so that~\eqref{eq:four-jets-sm5} yields $\varpi_{\mathrm{sep}}\le\frac{1}{2}\,\bar{\kappa}^{N^{*}}\le\frac{1}{2}$. As in (d), for all $(\theta,A_{*})$,
\[
\left\|k_S(\theta,A_{*};\cdot)\right\|_{\mathbf{r}}\le\exp\left(\|\mathrm{Lg}_S\|_{\mathbf{r}}\right)-1\le2\,\|\mathrm{Lg}_S\|_{\mathbf{r}}\le2\,\varpi_{\mathrm{sep}}\le\bar{\kappa}^{N^{*}},
\]
the last inequality being~\eqref{eq:four-jets-sm5}.

\emph{(f)} is Proposition~\ref{prop:large-field-term}(d), and \emph{(g)} restates the definitions of $t_x$ and of $\mathbf{r}$ in Section~\ref{sec:probes-mirror-copies}.

\emph{Independence of the bounds.} Each of the six numbers $a_{*}$, $e_{*}$, $l_{*}$, $r_{*}$, $p_{*}$, $p_{0*}$ has the same value for all $N\ge N_{\mathrm{min}}(R^{*},e)$. For $a_{*}$, $e_{*}$, $l_{*}$ and $r_{*}=l_{*}/M$ this is Proposition~\ref{prop:depth-stop-level}(c). By the definitions of Section~\ref{sec:running-quantities-tori} one has $p_{*}=v_{*}^p$ and $p_{0*}=v_{*}^{p_0}$ with $v_{*}=-\log e_{*}$, the exponents $p$ and $p_0$ being the constants of stage~\ref{it:order-choice-8}; so $p_{*}$ and $p_{0*}$ are functions of $e_{*}$ alone, and the assertion for them follows from that for $e_{*}$. Consider a constant of Section~\ref{sec:constants-3} other than $\varpi_{\mathrm{sep}}$, or a condition of Section~\ref{sec:smallness-conditions} other than~\eqref{eq:four-jets-sm2} and~\eqref{eq:four-jets-sm5}. Apart from the charge $e$ and the constants listed at the beginning of Section~\ref{sec:constants-3}, its definition involves only these six numbers and the numbers $C_{\beta}$, $\mathbf{D}$, $D_{\beta}$; the variable $\theta$ does not occur in it. The same is true of the number $\mathbf{r}$, which by the formula in (g) is formed with $e_{*}$, $h'$ and $C_{\beta}$. By Section~\ref{sec:probes-mirror-copies}, $C_{\beta}$ and $\mathbf{D}$ are suprema over $N$, over $a_{*}\in\left\{h'/L,h'/L^2\right\}$, over the admitted placements and over the forms $r_{*}\beta$, while $D_{\beta}$ is not altered by any of these changes. Consequently the right-hand sides of the inequalities in~\eqref{eq:four-jets-23} and the numbers $\varpi_0$, $\bar{\kappa}$, $\bar{\kappa}^{N^{*}}$, $\boldsymbol{\lambda}_L^{N^{*}}$, $\mathbf{r}$ keep their values. The definition of $\varpi_{\mathrm{sep}}$ and the conditions~\eqref{eq:four-jets-sm2}, \eqref{eq:four-jets-sm5} involve, besides numbers of the kind just described, one further quantity, $d_{\mathrm{sep}}$. Hence $\varpi_{\mathrm{sep}}$ is one and the same function of $d_{\mathrm{sep}}$ throughout, and this function is decreasing; likewise, each of the two conditions, if satisfied, stays satisfied when $d_{\mathrm{sep}}$ is enlarged. Finally, whatever $N$, the admitted placement and the form $r_{*}\beta$ are, Lemma~\ref{lem:properties-probes}(e) gives $d_{\mathrm{sep}}\ge\frac{M\cdot r_{*}}{8L}-1$; substituted into the definition of $\varpi_{\mathrm{sep}}$ in Section~\ref{sec:constants-3}, this lower bound yields the inequality for $\varpi_{\mathrm{sep}}$ stated in the theorem.

Theorem~\ref{thm:representation-ring-jets} follows.
\end{proof}

\begin{corollary}[The fourth cumulant of one probe]\label{cor:fourth-cumulant-one}
Assume the hypotheses of Theorem~\ref{thm:factorised-representation}, let $N\ge N_{\mathrm{min}}(R^{*},e)$, and suppose that
\[
\boldsymbol{\lambda}_L^{N^{*}}\le\frac{1}{16}.
\]
Let $\kappa_4(X;\bar{\mu})$ be the fourth cumulant of the random variable $X$ of Section~\ref{sec:probes-mirror-copies} under the law $\bar{\mu}$ of Section~\ref{sec:averaged-law-partition}. Then
\[
\left|\kappa_4(X;\bar{\mu})-e^4\cdot b_{\eta,R^{*}}(0)\right|\le96\,e^4\cdot\mathbf{r}^{-4}\cdot\left(\bar{\kappa}+2\,\boldsymbol{\lambda}_L^{N^{*}}\right).
\]
\end{corollary}

\begin{proof}
Let $c:=c_0$ be the cell with $\mathfrak{c}_{c_0}=\mathfrak{c}_0$, so that $X=X_c$ (Section~\ref{sec:probes-mirror-copies}).

\emph{Jets in one variable.} Let $\mathbb{C}[t]/(t^5)$ be the quotient of the polynomial ring in one variable $t$ by the ideal generated by $t^5$. Each element $P$ of this ring is the class of exactly one polynomial $\sum_{i=0}^4P_i\,t^i$, and $P_i$ is called the coefficient of $t^i$ in $P$. Put
\[
\|P\|_{\mathbf{r}}:=\sum_{i=0}^4\left|P_i\right|\cdot\mathbf{r}^i;
\]
this is the norm of Section~\ref{sec:ring-jets} in the case of a single variable, and it is submultiplicative for the same reason. Directly from the definition,
\[
\left|P_4\right|\le\mathbf{r}^{-4}\cdot\|P\|_{\mathbf{r}}\qquad\left(P\in\mathbb{C}[t]/(t^5)\right).
\]
As in Section~\ref{sec:ring-jets}, a function of one complex variable that is analytic in a neighbourhood of $0$ is identified with its image in $\mathbb{C}[t]/(t^5)$, the class of its Taylor polynomial of degree four at $0$; this assignment is a homomorphism of rings. For $P\in\mathbb{C}[t]/(t^5)$ with $\|P\|_{\mathbf{r}}<1$ let
\[
\log(1+P):=\sum_{n\ge1}\frac{(-1)^{n+1}}{n}\,P^n.
\]
Since the norm is submultiplicative, the series converges, and
\[
\left\|\log(1+P)\right\|_{\mathbf{r}}\le\sum_{n\ge1}\frac{1}{n}\,\|P\|_{\mathbf{r}}^n=-\log\left(1-\|P\|_{\mathbf{r}}\right).
\]

\emph{Restriction to the cell $c$.} For $P\in\mathcal{J}$ let $P(t\cdot e_c)\in\mathbb{C}[t]/(t^5)$ be the element obtained from $P$ by substituting $t$ for $t_c$ and $0$ for every $t_{c'}$ with $c'\in\mathrm{Cel}\setminus\{c\}$: the monomials of $P$ that contain a variable $t_{c'}$ with $c'\ne c$ are deleted, and $t$ is written for $t_c$ in the remaining ones. The substitution maps the ideal generated by the $t_{c'}^5$, $c'\in\mathrm{Cel}$, into the ideal generated by $t^5$; so $P\mapsto P(t\cdot e_c)$ is a homomorphism of rings. It does not increase the norm $\|\cdot\|_{\mathbf{r}}$, and, acting on the coefficients, it commutes with the coefficientwise integration of functions with values in $\mathcal{J}$. Let a function be analytic in a neighbourhood of $0$ in $\mathbb{C}^{\mathrm{Cel}}$. The Taylor series at $0$ of its composition with the map $t\mapsto t\cdot e_c$ consists of those terms of its own Taylor series that contain no variable other than $t_c$, with $t$ written for $t_c$. Hence the image of the composition in $\mathbb{C}[t]/(t^5)$ arises from the image of the function in $\mathcal{J}$ by the substitution; this agrees with the notation $P(t\cdot e_c)$.

Let $\theta\in\mathbb{R}^3$ and $A_{*}\in V_{\mathrm{t}}\left(T^0_{*}\right)$ with $\chi_{*}(A_{*})=1$. A non-empty set $S\subset\mathrm{Cel}$ other than $\{c\}$ contains a cell $c'\ne c$; by Theorem~\ref{thm:factorised-representation}(c) every monomial of $k_S(\theta,A_{*};\cdot)$ contains $t_{c'}$, and therefore $k_S(\theta,A_{*};t\cdot e_c)=0$. By Theorem~\ref{thm:factorised-representation}(a), $\phi_0(0)=1$. Apply the substitution to~\eqref{eq:four-jets-26}: the product of the factors $\phi_0(t_{c'})$, $c'\in\mathrm{Cel}$, becomes $\phi_0(t)$; the product over $S$ becomes $1+k_{\{c\}}(\theta,A_{*};t\cdot e_c)$; and the $\boldsymbol{\varphi}$-integral of the constant $1$ is the total mass of $\boldsymbol{\varphi}$, which equals $\bar{\mu}(\mathfrak{G})$ by Theorem~\ref{thm:factorised-representation}(b). One obtains in $\mathbb{C}[t]/(t^5)$
\[
\mathring{\Phi}(t\cdot e_c)=\phi_0(t)\cdot\left(m_v+\boldsymbol{\theta}_1(t)\right)+\mathrm{Lr}_1(t),
\]
where the left side is the image of the entire function $t\mapsto\mathring{\Phi}(t\cdot e_c)$, and where
\[
m_v:=\bar{\mu}(\mathfrak{G}),\qquad\boldsymbol{\theta}_1(t):=\int\boldsymbol{\varphi}(\mathrm{d}\theta\,\mathrm{d}A_{*})\,k_{\{c\}}(\theta,A_{*};t\cdot e_c),\qquad\mathrm{Lr}_1(t):=\mathrm{Lr}(t\cdot e_c).
\]
By Theorem~\ref{thm:factorised-representation}(b) and the first condition in~\eqref{eq:four-jets-sm7},
\[
1\ge m_v\ge1-\varpi_0\ge\frac{3}{4}.
\]
Each coefficient of $\boldsymbol{\theta}_1$ is the $\boldsymbol{\varphi}$-integral of the corresponding coefficient of $k_{\{c\}}(\theta,A_{*};t\cdot e_c)$; as $\boldsymbol{\varphi}$ is positive, its modulus is at most the $\boldsymbol{\varphi}$-integral of the modulus of that coefficient. Summing over the coefficients, and using that the substitution does not increase the norm, together with Theorem~\ref{thm:factorised-representation}(d), (f), one finds
\[
\|\boldsymbol{\theta}_1\|_{\mathbf{r}}\le\int\boldsymbol{\varphi}(\mathrm{d}\theta\,\mathrm{d}A_{*})\,\left\|k_{\{c\}}(\theta,A_{*};\cdot)\right\|_{\mathbf{r}}\le m_v\cdot\bar{\kappa}\le\bar{\kappa},\qquad\|\mathrm{Lr}_1\|_{\mathbf{r}}\le\|\mathrm{Lr}\|_{\mathbf{r}}\le\boldsymbol{\lambda}_L^{N^{*}}.
\]

\emph{Division by $\phi_0$.} By Theorem~\ref{thm:factorised-representation}(a) the relations~\eqref{eq:four-jets-23} hold; among them is $\|\phi_0-1\|_{\mathbf{r}}\le\frac{1}{6}$. The element $u:=1-\phi_0$ of $\mathbb{C}[t]/(t^5)$ has no constant term, so that $u^5=0$. Hence $\phi_0=1-u$ is invertible, and
\[
\phi_0^{-1}=\sum_{n=0}^4u^n,\qquad\left\|\phi_0^{-1}\right\|_{\mathbf{r}}\le\sum_{n\ge0}6^{-n}=\frac{6}{5}.
\]
Multiplying the identity obtained above by $\phi_0^{-1}$, one finds
\[
\phi_0^{-1}\cdot\mathring{\Phi}(t\cdot e_c)=m_v+\boldsymbol{\theta},\qquad\boldsymbol{\theta}:=\boldsymbol{\theta}_1+\phi_0^{-1}\cdot\mathrm{Lr}_1,
\]
and
\[
\|\boldsymbol{\theta}\|_{\mathbf{r}}\le\bar{\kappa}+\frac{6}{5}\,\boldsymbol{\lambda}_L^{N^{*}}\le\tilde{\kappa}:=\bar{\kappa}+2\,\boldsymbol{\lambda}_L^{N^{*}}.
\]

\emph{Smallness.} By~\eqref{eq:four-jets-sm6}, $\bar{\kappa}^{1-s}\le\frac{1}{2N^{*}}<1$. As $0<s<1$ (stage~\ref{it:order-choice-8}), it follows that $\bar{\kappa}<1$, and then $\bar{\kappa}\le\bar{\kappa}^{1-s}\le\frac{1}{2N^{*}}\le\frac{1}{16}$, because $N^{*}=8L^3\ge8$. Together with the hypothesis $\boldsymbol{\lambda}_L^{N^{*}}\le\frac{1}{16}$ and with $m_v\ge\frac{3}{4}$ this gives
\[
\tilde{\kappa}\le\frac{3}{16},\qquad\frac{\tilde{\kappa}}{m_v}\le\frac{3}{16}\cdot\frac{4}{3}=\frac{1}{4}.
\]
For $0\le y\le\frac{1}{4}$ one has $-\log(1-y)=\sum_{n\ge1}\frac{y^n}{n}\le\frac{y}{1-y}\le\frac{4}{3}\,y$. Since $\left\|\boldsymbol{\theta}/m_v\right\|_{\mathbf{r}}\le\tilde{\kappa}/m_v\le\frac{1}{4}$, the bound for the logarithm in $\mathbb{C}[t]/(t^5)$ and the inequality $m_v\ge\frac{3}{4}$ yield
\[
\left\|\log\left(1+\frac{\boldsymbol{\theta}}{m_v}\right)\right\|_{\mathbf{r}}\le-\log\left(1-\frac{\tilde{\kappa}}{m_v}\right)\le\frac{4}{3}\cdot\frac{\tilde{\kappa}}{m_v}\le4\,\tilde{\kappa}.
\]

\emph{The cumulant.} For a function with values in the disc $\left\{w\in\mathbb{C}:\ |w-1|<1\right\}$, $\log$ denotes the principal branch of the logarithm. Proposition~\ref{prop:large-field-term}(b), applied at the points $z\cdot e_c$ of $\mathbb{C}^{\mathrm{Cel}}$, $z\in\mathbb{C}$, gives
\[
\int\mathrm{d}\bar{\mu}(F)\,\exp\left(z\cdot X(F)\right)=\exp\left(\frac{1}{2}\,z^2\cdot\boldsymbol{\sigma}^2\right)\cdot\mathring{\Phi}(e\cdot z\cdot e_c)\qquad\left(z\in\mathbb{C}\right).
\]
Thus the moment generating function of $X$ under $\bar{\mu}$ is finite on $\mathbb{R}$ and extends to an entire function which takes the value $1$ at $z=0$; consequently $\kappa_4(X;\bar{\mu})$ is the fourth derivative at $z=0$ of its logarithm. In a neighbourhood of $z=0$ this logarithm equals $\frac{1}{2}\,z^2\cdot\boldsymbol{\sigma}^2+\log\mathring{\Phi}(e\cdot z\cdot e_c)$, since the two functions are analytic there, vanish at $z=0$ and have the same exponential. The fourth derivative of $\frac{1}{2}\,z^2\cdot\boldsymbol{\sigma}^2$ vanishes, and the chain rule gives
\[
\kappa_4(X;\bar{\mu})=\left.\frac{\mathrm{d}^4}{\mathrm{d}z^4}\log\mathring{\Phi}(e\cdot z\cdot e_c)\right|_{z=0}=e^4\cdot\left.\frac{\mathrm{d}^4}{\mathrm{d}t^4}\log\mathring{\Phi}(t\cdot e_c)\right|_{t=0}.
\]

\emph{The fourth derivative.} Let
\[
\mathfrak{g}(t):=\frac{\mathring{\Phi}(t\cdot e_c)}{m_v\cdot N_0(t)}.
\]
As $N_0$ is entire with $N_0(0)=1$ (Section~\ref{sec:averaged-law-partition}) and $\mathring{\Phi}(0)=1$, the function $\mathfrak{g}$ is analytic in a neighbourhood of $0$, and $\mathfrak{g}(0)=m_v^{-1}$. By Theorem~\ref{thm:factorised-representation}(a) the image of $N_0$ is $\phi_0$. The assignment of images being a homomorphism of rings, the image of $1/N_0$ is $\phi_0^{-1}$, and the image of $\mathfrak{g}$ is
\[
m_v^{-1}\cdot\phi_0^{-1}\cdot\mathring{\Phi}(t\cdot e_c)=m_v^{-1}\cdot\left(m_v+\boldsymbol{\theta}\right)=1+\frac{\boldsymbol{\theta}}{m_v}.
\]
The image of $\log\mathfrak{g}$ is
\[
\check{P}:=\log\left(1+\frac{\boldsymbol{\theta}}{m_v}\right),
\]
the logarithm on the right being that of $\mathbb{C}[t]/(t^5)$. Indeed, $\left|\mathfrak{g}(0)-1\right|$ is the modulus of the constant term of $\boldsymbol{\theta}/m_v$, hence at most $\left\|\boldsymbol{\theta}/m_v\right\|_{\mathbf{r}}\le\frac{1}{4}$. So there is a radius $\rho>0$ such that $\mathfrak{g}$ is analytic in a neighbourhood of the closed disc $\left\{|t|\le\rho\right\}$ and satisfies $\left|\mathfrak{g}-1\right|\le\frac{1}{2}$ on it. On this disc the series $\sum_{n\ge1}\frac{(-1)^{n+1}}{n}\left(\mathfrak{g}-1\right)^n$ converges uniformly to $\log\mathfrak{g}$. By Cauchy's integral formula on the circle $\left\{|t|=\rho\right\}$, each Taylor coefficient of $\log\mathfrak{g}$ at $0$ is the sum of the corresponding Taylor coefficients of the terms of this series. The image of the term of index $n$ is $\frac{(-1)^{n+1}}{n}\left(\frac{\boldsymbol{\theta}}{m_v}\right)^n$; summing over $n$ gives the claim.

In a neighbourhood of $t=0$,
\[
\log\mathring{\Phi}(t\cdot e_c)=\log m_v+\log N_0(t)+\log\mathfrak{g}(t),
\]
because both sides are analytic there, have the same exponential, and vanish at $t=0$. The coefficient of $t^4$ in the image of the left side is $\frac{1}{24}$ times its fourth derivative at $t=0$. On the right side, the constant $\log m_v$ contributes nothing to this coefficient, the fourth derivative of $\log N_0$ at $0$ is $b_{\eta,R^{*}}(0)$ by the definition in Section~\ref{sec:averaged-law-partition}, and the image of $\log\mathfrak{g}$ is $\check{P}$. Hence
\[
\left.\frac{\mathrm{d}^4}{\mathrm{d}t^4}\log\mathring{\Phi}(t\cdot e_c)\right|_{t=0}=b_{\eta,R^{*}}(0)+24\,\check{P}_4,\qquad\kappa_4(X;\bar{\mu})=e^4\cdot\left(b_{\eta,R^{*}}(0)+24\,\check{P}_4\right).
\]

\emph{Conclusion.} The bound for the coefficient of $t^4$ and the bound for $\|\check{P}\|_{\mathbf{r}}$ obtained above give
\[
\left|\kappa_4(X;\bar{\mu})-e^4\cdot b_{\eta,R^{*}}(0)\right|=24\,e^4\cdot\left|\check{P}_4\right|\le24\,e^4\cdot\mathbf{r}^{-4}\cdot\|\check{P}\|_{\mathbf{r}}\le96\,e^4\cdot\mathbf{r}^{-4}\cdot\tilde{\kappa},
\]
and $\tilde{\kappa}=\bar{\kappa}+2\,\boldsymbol{\lambda}_L^{N^{*}}$. Theorem~\ref{thm:fourth-cumulant-one} follows.
\end{proof}

By Theorem~\ref{thm:factorised-representation}(g), $\mathbf{r}^{-1}=4\left(1+h'^{\frac{1}{2}}\cdot C_{\beta}\right)\cdot e_{*}^{-3\varepsilon}$, and $\bar{\kappa}=6\,e_{*}^{\frac{1}{4}-9\varepsilon}$ by the definition in Section~\ref{sec:constants-3}. Hence the number $\mathbf{r}^{-4}\cdot\bar{\kappa}$, which enters the bound of Corollary~\ref{cor:fourth-cumulant-one}, is
\[
\mathbf{r}^{-4}\cdot\bar{\kappa}=256\left(1+h'^{\frac{1}{2}}\cdot C_{\beta}\right)^4\cdot e_{*}^{-12\varepsilon}\cdot6\,e_{*}^{\frac{1}{4}-9\varepsilon}=1536\left(1+h'^{\frac{1}{2}}\cdot C_{\beta}\right)^4\cdot e_{*}^{\frac{1}{4}-21\varepsilon}.
\]

\subsection{Centred derivatives at the origin and the smallness conditions}\label{sec:centred-derivatives-origin}

\subsubsection{The numbers $\Theta^{(0)}_m$}\label{sec:numbers}

\begin{proposition}[Products of centred derivatives]\label{prop:products-centred-derivatives}
Assume the hypotheses of Theorem~\ref{thm:factorised-representation}, and let $N\ge N_{\mathrm{min}}(R^{*},e)$. Let $\mathring{\Phi}$ be the entire function on $\mathbb{C}^{\mathrm{Cel}}$ of Section~\ref{sec:averaged-law-partition}, let $\phi_0$ be the polynomial of Theorem~\ref{thm:factorised-representation}(a), and let $\mathbf{r}$, $\bar{\kappa}$, $\boldsymbol{\lambda}_L$ and the exponent $s$ be as in Sections~\ref{sec:probes-mirror-copies} and~\ref{sec:constants-3}. For $m\in\{1,2,3,4\}$ let $c_m:=\phi_0^{(m)}(0)$; since $\phi_0$ is even, $c_1=c_3=0$. The differential operators $\partial_{t_c}^m-c_m$, $c\in\mathrm{Cel}$, commute with one another; put
\[
\Theta^{(0)}_m:=\left|\left(\prod_{c\in\mathrm{Cel}}\left(\partial_{t_c}^m-c_m\right)\mathring{\Phi}\right)(0)\right|^{\frac{1}{N^{*}}}.
\]
Then
\[
\Theta^{(0)}_m\le m!\cdot\mathbf{r}^{-m}\cdot\left(2\exp(1)\cdot\bar{\kappa}^s+\boldsymbol{\lambda}_L\right),
\]
and, if in addition $\boldsymbol{\lambda}_L\le\bar{\kappa}^s$,
\[
\Theta^{(0)}_m\le8\,m!\cdot\mathbf{r}^{-m}\cdot\bar{\kappa}^s.
\]
\end{proposition}

\begin{proof}
Fix $m\in\{1,2,3,4\}$. As in Section~\ref{sec:ring-jets}, a function that is analytic in a neighbourhood of $t=0$ is identified with its image in $\mathcal{J}$, and $\|\cdot\|_{\mathbf{r}}$ denotes the norm of that section with $\rho=\mathbf{r}$. The polynomial $\phi_0$ has degree at most $4$; write $\phi_0(t)=\sum_{i=0}^4n_it^i$. By Theorem~\ref{thm:factorised-representation}(a), $n_0=\phi_0(0)=1$, and by the definition of $c_m$, $m!\,n_m=c_m$.

For $c\in\mathrm{Cel}$ let $\boldsymbol{\ell}_c$ denote the linear functional on $\mathbb{C}[t_c]/(t_c^5)$ that takes on the basis $1,t_c,\dots,t_c^4$ the values
\[
\boldsymbol{\ell}_c\left(t_c^i\right)=\begin{cases}m!&\text{if }i=m,\\-c_m&\text{if }i=0,\\0&\text{if }i\in\{1,2,3,4\}\setminus\{m\},\end{cases}
\]
and let $\boldsymbol{\ell}$ be the tensor product of these functionals, that is, the linear functional on $\mathcal{J}$ with
\[
\boldsymbol{\ell}\left(t^{\mathbf{a}}\right)=\prod_{c\in\mathrm{Cel}}\boldsymbol{\ell}_c\left(t_c^{\mathbf{a}_c}\right)\qquad\left(\mathbf{a}\in\{0,1,2,3,4\}^{\mathrm{Cel}}\right).
\]
For an integer $i\ge0$, the value of $\left(\partial_{t_c}^m-c_m\right)t_c^i$ at $t_c=0$ is $\boldsymbol{\ell}_c(t_c^i)$ if $i\le4$, and $0$ if $i\ge5$, because $m\le4$. On a monomial $\prod_ct_c^{i_c}$ the operator $\prod_c\left(\partial_{t_c}^m-c_m\right)$ acts factor by factor. Differentiating the Taylor series term by term, one therefore obtains, for every function $g$ that is analytic in a neighbourhood of $t=0$,
\[
\left(\prod_{c\in\mathrm{Cel}}\left(\partial_{t_c}^m-c_m\right)g\right)(0)=\boldsymbol{\ell}(g),
\]
where on the right $g$ stands for its image in $\mathcal{J}$. In particular,
\[
\Theta^{(0)}_m=\left|\boldsymbol{\ell}(\mathring{\Phi})\right|^{\frac{1}{N^{*}}}.
\]

By Theorem~\ref{thm:factorised-representation}(a) the inequalities \eqref{eq:four-jets-23} hold; they contain $|c_2|\le\frac{\mathbf{r}^{-2}}{16}$ and $|c_4|\le\frac{\mathbf{r}^{-4}}{16}$. Since moreover $c_1=c_3=0$,
\[
|c_m|\le\frac{\mathbf{r}^{-m}}{16}\le m!\cdot\mathbf{r}^{-m}=:A_m.
\]
Consequently $\left|\boldsymbol{\ell}_c(t_c^i)\right|\le A_m\cdot\mathbf{r}^i$ for $0\le i\le4$, with equality for $i=m$. For $\check{P}=\sum_{\mathbf{a}}\check{P}_{\mathbf{a}}t^{\mathbf{a}}\in\mathcal{J}$ it follows that
\[
\left|\boldsymbol{\ell}(\check{P})\right|\le\sum_{\mathbf{a}}\left|\check{P}_{\mathbf{a}}\right|\prod_{c\in\mathrm{Cel}}\left|\boldsymbol{\ell}_c\left(t_c^{\mathbf{a}_c}\right)\right|\le\sum_{\mathbf{a}}\left|\check{P}_{\mathbf{a}}\right|\prod_{c\in\mathrm{Cel}}A_m\cdot\mathbf{r}^{\mathbf{a}_c}=A_m^{N^{*}}\cdot\|\check{P}\|_{\mathbf{r}},
\]
the set $\mathrm{Cel}$ having $N^{*}$ elements.

For $c\in\mathrm{Cel}$ put $\varkappa_{\{c\}}:=\bar{\kappa}$, and for $S\subset\mathrm{Cel}$ with $\#|S|\ge2$ put $\varkappa_S:=2\,\varpi_{\mathrm{sep}}$. By \eqref{eq:four-jets-sm6} and $0<s<1$ (stage~\ref{it:order-choice-8}) one has $\bar{\kappa}^{1-s}\le\frac{1}{2N^{*}}<1$, hence $\bar{\kappa}<1$, and \eqref{eq:four-jets-sm5} gives $2\,\varpi_{\mathrm{sep}}\le\bar{\kappa}^{N^{*}}<1$. So $\varkappa_S\le1$ for every non-empty $S\subset\mathrm{Cel}$. By Theorem~\ref{thm:factorised-representation}(d), (e),
\[
\left\|k_S(\theta,A_{*};\cdot)\right\|_{\mathbf{r}}\le\varkappa_S
\]
for every non-empty $S\subset\mathrm{Cel}$ and all $(\theta,A_{*})$.

Call a set $\mathfrak{Y}$ of non-empty subsets of $\mathrm{Cel}$ a family, and say that $\mathfrak{Y}$ covers $\mathrm{Cel}$ if each cell belongs to at least one member of $\mathfrak{Y}$. Multiplying out the product, one has in $\mathcal{J}$
\[
\prod_{\emptyset\ne S\subset\mathrm{Cel}}\left(1+k_S\right)=\sum_{\mathfrak{Y}}\prod_{S\in\mathfrak{Y}}k_S,
\]
the sum being extended over all families, the empty family included, for which the product equals $1$.

Let $\mathfrak{Y}$ be a family that does not cover $\mathrm{Cel}$, and choose a cell $c_0$ that belongs to no member of $\mathfrak{Y}$. By Theorem~\ref{thm:factorised-representation}(c), no monomial of a $k_S$ with $S\in\mathfrak{Y}$ contains the variable $t_{c_0}$. Hence
\[
\prod_{c\in\mathrm{Cel}}\phi_0(t_c)\prod_{S\in\mathfrak{Y}}k_S=\phi_0(t_{c_0})\cdot\check{P}',\qquad\check{P}':=\prod_{c\in\mathrm{Cel}\setminus\{c_0\}}\phi_0(t_c)\prod_{S\in\mathfrak{Y}}k_S,
\]
where $\check{P}'$ does not contain $t_{c_0}$. Since $\boldsymbol{\ell}$ is the tensor product of the $\boldsymbol{\ell}_c$, its value on this product is $\boldsymbol{\ell}_{c_0}(\phi_0)$ times the value at $\check{P}'$ of the tensor product of the functionals $\boldsymbol{\ell}_c$, $c\ne c_0$. Now
\[
\boldsymbol{\ell}_{c_0}(\phi_0)=m!\,n_m-c_m\,n_0=c_m-c_m=0.
\]
Thus $\boldsymbol{\ell}$ annihilates the term of each family that does not cover $\mathrm{Cel}$.

Let now $\mathfrak{Y}$ cover $\mathrm{Cel}$. The members of $\mathfrak{Y}$ then have cardinalities whose sum is at least the number of cells:
\[
\sum_{S\in\mathfrak{Y}}\#|S|\ge\#|\mathrm{Cel}|=N^{*}.
\]
Put
\[
\boldsymbol{\theta}_s:=\max_{\emptyset\ne S\subset\mathrm{Cel}}\varkappa_S^{\frac{s}{\#|S|}},\qquad\boldsymbol{\Sigma}_s:=\max_{c\in\mathrm{Cel}}\sum_{S\subset\mathrm{Cel},\,S\ni c}\varkappa_S^{1-s}.
\]
As $\varkappa_S\le1$ for all $S$, one has $\boldsymbol{\theta}_s\le1$, and therefore
\[
\prod_{S\in\mathfrak{Y}}\varkappa_S^s=\prod_{S\in\mathfrak{Y}}\left(\varkappa_S^{\frac{s}{\#|S|}}\right)^{\#|S|}\le\boldsymbol{\theta}_s^{\sum_{S\in\mathfrak{Y}}\#|S|}\le\boldsymbol{\theta}_s^{N^{*}}.
\]
Write $\varkappa_S=\varkappa_S^s\cdot\varkappa_S^{1-s}$, apply this bound, and enlarge the sum over the covering families to the sum over all families, which is equal to $\prod_S\left(1+\varkappa_S^{1-s}\right)$. This gives the first of the inequalities
\begin{align*}
\sum_{\mathfrak{Y}\text{ covers }\mathrm{Cel}}\ \prod_{S\in\mathfrak{Y}}\varkappa_S&\le\boldsymbol{\theta}_s^{N^{*}}\prod_{\emptyset\ne S\subset\mathrm{Cel}}\left(1+\varkappa_S^{1-s}\right)\le\boldsymbol{\theta}_s^{N^{*}}\exp\left(\sum_{c\in\mathrm{Cel}}\sum_{S\ni c}\varkappa_S^{1-s}\right)\\
&\le\left(\boldsymbol{\theta}_s\cdot\exp\left(\boldsymbol{\Sigma}_s\right)\right)^{N^{*}}.
\end{align*}
The second one follows from $1+x\le\exp(x)$ and from the fact that every non-empty $S$ contains at least one cell, the third from the definition of $\boldsymbol{\Sigma}_s$, as there are $N^{*}$ cells.

By \eqref{eq:four-jets-26} one has in $\mathcal{J}$
\[
\mathring{\Phi}=\prod_{c\in\mathrm{Cel}}\phi_0(t_c)\int\boldsymbol{\varphi}(\mathrm{d}\theta\,\mathrm{d}A_{*})\,\prod_{\emptyset\ne S\subset\mathrm{Cel}}\left(1+k_S(\theta,A_{*};\cdot)\right)+\mathrm{Lr},
\]
the integral of a function with values in $\mathcal{J}$ being taken coefficient by coefficient. Each coefficient of a product in $\mathcal{J}$ is a finite bilinear expression in the coefficients of the factors; so the factor $\prod_c\phi_0(t_c)$, which does not depend on $(\theta,A_{*})$, may be moved under the integral sign. As $\mathcal{J}$ has finite dimension, $\boldsymbol{\ell}(\check{P})$ is a finite linear combination of the coefficients of $\check{P}$, so that $\boldsymbol{\ell}$, too, may be applied under the integral sign. With the expansion of the product and the vanishing proved above, this gives
\[
\boldsymbol{\ell}(\mathring{\Phi})=\int\boldsymbol{\varphi}(\mathrm{d}\theta\,\mathrm{d}A_{*})\sum_{\mathfrak{Y}\text{ covers }\mathrm{Cel}}\boldsymbol{\ell}\left(\prod_{c\in\mathrm{Cel}}\phi_0(t_c)\prod_{S\in\mathfrak{Y}}k_S(\theta,A_{*};\cdot)\right)+\boldsymbol{\ell}(\mathrm{Lr}).
\]
Let $(\theta,A_{*})$ be fixed and let $\mathfrak{Y}$ cover $\mathrm{Cel}$. The bound for $\boldsymbol{\ell}$, the submultiplicativity of the norm (Section~\ref{sec:ring-jets}), the equality $\|\phi_0(t_c)\|_{\mathbf{r}}=\|\phi_0\|_{\mathbf{r}}$ and the bounds for $\|k_S\|_{\mathbf{r}}$ yield
\[
\left|\boldsymbol{\ell}\left(\prod_{c\in\mathrm{Cel}}\phi_0(t_c)\prod_{S\in\mathfrak{Y}}k_S(\theta,A_{*};\cdot)\right)\right|\le A_m^{N^{*}}\cdot\|\phi_0\|_{\mathbf{r}}^{N^{*}}\prod_{S\in\mathfrak{Y}}\varkappa_S.
\]
The measure $\boldsymbol{\varphi}$ is positive, and its total mass does not exceed $1$ (Theorem~\ref{thm:factorised-representation}(b)). Summing the last bound over the covering families, using the bound obtained above for $\sum_{\mathfrak{Y}}\prod_{S\in\mathfrak{Y}}\varkappa_S$ over these families, integrating, and bounding $\left|\boldsymbol{\ell}(\mathrm{Lr})\right|\le A_m^{N^{*}}\cdot\|\mathrm{Lr}\|_{\mathbf{r}}$, one arrives at
\[
\left|\boldsymbol{\ell}(\mathring{\Phi})\right|\le A_m^{N^{*}}\left(\left(\|\phi_0\|_{\mathbf{r}}\cdot\boldsymbol{\theta}_s\cdot\exp\left(\boldsymbol{\Sigma}_s\right)\right)^{N^{*}}+\|\mathrm{Lr}\|_{\mathbf{r}}\right).
\]
For $x,y\ge0$ one has $(x+y)^{\frac{1}{N^{*}}}\le x^{\frac{1}{N^{*}}}+y^{\frac{1}{N^{*}}}$; hence
\[
\Theta^{(0)}_m\le A_m\left(\|\phi_0\|_{\mathbf{r}}\cdot\boldsymbol{\theta}_s\cdot\exp\left(\boldsymbol{\Sigma}_s\right)+\|\mathrm{Lr}\|_{\mathbf{r}}^{\frac{1}{N^{*}}}\right).
\]

The quantities on the right are estimated as follows. First, $\|\phi_0\|_{\mathbf{r}}\le2$ by \eqref{eq:four-jets-23}. Second, $\boldsymbol{\theta}_s\le\bar{\kappa}^s$: for $S=\{c\}$ one has $\varkappa_S^{\frac{s}{\#|S|}}=\bar{\kappa}^s$, and for $\#|S|\ge2$ the inequality \eqref{eq:four-jets-sm5} and $\bar{\kappa}<1$ give $\varkappa_S\le\bar{\kappa}^{N^{*}}\le\bar{\kappa}^{\#|S|}$, whence $\varkappa_S^{\frac{s}{\#|S|}}\le\bar{\kappa}^s$. Third, $\boldsymbol{\Sigma}_s\le1$: a cell $c$ belongs to the set $\{c\}$, for which $\varkappa_{\{c\}}^{1-s}=\bar{\kappa}^{1-s}$, and to fewer than $2^{N^{*}}$ sets $S$ with $\#|S|\ge2$, for each of which $\varkappa_S^{1-s}\le\bar{\kappa}^{N^{*}(1-s)}$ by \eqref{eq:four-jets-sm5}; with \eqref{eq:four-jets-sm6} this gives
\[
\boldsymbol{\Sigma}_s\le\bar{\kappa}^{1-s}+2^{N^{*}}\cdot\bar{\kappa}^{N^{*}(1-s)}\le\frac{1}{2N^{*}}+\left(N^{*}\right)^{-N^{*}}\le1.
\]
Fourth, $\|\mathrm{Lr}\|_{\mathbf{r}}^{\frac{1}{N^{*}}}\le\boldsymbol{\lambda}_L$ by Theorem~\ref{thm:factorised-representation}(f). Inserting these four bounds and $A_m=m!\cdot\mathbf{r}^{-m}$, one obtains
\[
\Theta^{(0)}_m\le m!\cdot\mathbf{r}^{-m}\cdot\left(2\exp(1)\cdot\bar{\kappa}^s+\boldsymbol{\lambda}_L\right).
\]
If $\boldsymbol{\lambda}_L\le\bar{\kappa}^s$, the right side does not exceed $\left(2\exp(1)+1\right)m!\cdot\mathbf{r}^{-m}\cdot\bar{\kappa}^s$, and $2\exp(1)+1\le8$.
\end{proof}

\begin{corollary}[Bounds in powers of $e_{*}$]\label{cor:bounds-powers}
Assume the hypotheses of Proposition~\ref{prop:products-centred-derivatives} and $\boldsymbol{\lambda}_L\le\bar{\kappa}^s$, and put
\[
\mathfrak{e}:=s\left(\frac{1}{4}-9\varepsilon\right)-12\varepsilon,\qquad\bar{C}:=1+h'^{\frac{1}{2}}\cdot C_{\beta}.
\]
Then
\begin{gather*}
\Theta^{(0)}_4\le3\cdot10^5\,\bar{C}^4\cdot e_{*}^{\mathfrak{e}},\qquad\left(1+|c_2|\right)\Theta^{(0)}_2\le3\cdot10^4\,\bar{C}^4\cdot e_{*}^{\mathfrak{e}},\\
|c_2|+|c_4|^{\frac{1}{2}}\le8\,\bar{C}^2\cdot e_{*}^{-6\varepsilon},\qquad e\left(\frac{h'}{L^2}\right)^{\frac{1}{2}}\le e_{*}\le e\,h'^{\frac{1}{2}}.
\end{gather*}
\end{corollary}

\begin{proof}
By the definitions of $t_x$ and $\mathbf{r}$ in Section~\ref{sec:probes-mirror-copies},
\[
\mathbf{r}^{-1}=\frac{4}{t_x}=4\,\bar{C}\cdot e_{*}^{-3\varepsilon}\ge1;
\]
the inequality holds because $\bar{C}\ge1$ and $e_{*}\le e_0<1$ (Lemma~\ref{lem:scales-stop-level}(a) and condition~\eqref{eq:renormalisation-group-A8}). By the definition of $\bar{\kappa}$ in Section~\ref{sec:constants-3} and since $0<s<1$ (stage~\ref{it:order-choice-8}),
\[
\bar{\kappa}^s=6^s\cdot e_{*}^{s\left(\frac{1}{4}-9\varepsilon\right)}\le6\,e_{*}^{s\left(\frac{1}{4}-9\varepsilon\right)}.
\]
Proposition~\ref{prop:products-centred-derivatives} with $m=4$, together with the hypothesis $\boldsymbol{\lambda}_L\le\bar{\kappa}^s$, gives
\[
\Theta^{(0)}_4\le8\cdot4!\cdot\mathbf{r}^{-4}\cdot\bar{\kappa}^s=192\,\mathbf{r}^{-4}\cdot\bar{\kappa}^s\le192\cdot256\cdot6\,\bar{C}^4\cdot e_{*}^{-12\varepsilon}\cdot e_{*}^{s\left(\frac{1}{4}-9\varepsilon\right)}=294912\,\bar{C}^4\cdot e_{*}^{\mathfrak{e}},
\]
and $294912\le3\cdot10^5$. With $m=2$ it gives $\Theta^{(0)}_2\le16\,\mathbf{r}^{-2}\cdot\bar{\kappa}^s$, while \eqref{eq:four-jets-23} and $\mathbf{r}^{-2}\ge1$ give
\[
1+|c_2|\le\mathbf{r}^{-2}+\frac{\mathbf{r}^{-2}}{16}=\frac{17}{16}\,\mathbf{r}^{-2}.
\]
Multiplying the two bounds,
\[
\left(1+|c_2|\right)\Theta^{(0)}_2\le17\,\mathbf{r}^{-4}\cdot\bar{\kappa}^s\le17\cdot256\cdot6\,\bar{C}^4\cdot e_{*}^{\mathfrak{e}}=26112\,\bar{C}^4\cdot e_{*}^{\mathfrak{e}}\le3\cdot10^4\,\bar{C}^4\cdot e_{*}^{\mathfrak{e}}.
\]
By \eqref{eq:four-jets-23} once more,
\[
|c_2|+|c_4|^{\frac{1}{2}}\le\frac{\mathbf{r}^{-2}}{16}+\frac{\mathbf{r}^{-2}}{4}\le\frac{\mathbf{r}^{-2}}{2}=8\,\bar{C}^2\cdot e_{*}^{-6\varepsilon}.
\]
Finally, by Lemma~\ref{lem:scales-stop-level}(a), $e_{*}=e\,a_{*}^{\frac{1}{2}}$ with $a_{*}\in\left\{\frac{h'}{L},\frac{h'}{L^2}\right\}$. As $L\ge1$, one has $\frac{h'}{L^2}\le a_{*}\le h'$; taking square roots and multiplying by $e$ gives the last assertion.
\end{proof}

Theorem~\ref{thm:products-centred-derivatives} follows from Proposition~\ref{prop:products-centred-derivatives} and Corollary~\ref{cor:bounds-powers}.

\subsubsection{The smallness conditions for small charge}\label{sec:smallness-conditions-small}

\begin{lemma}[Threshold of the charge]\label{lem:threshold-charge}
Let the constants of stages~\ref{it:order-choice-1}--\ref{it:order-choice-10} of Definition~\ref{def:order-choice} be fixed, together with the number $\bar{m}\in(0,1]$ and the two-form $\beta$ of Section~\ref{sec:probes-mirror-copies}. There is a number
\[
e_{\sharp}\in\left(0,e_1\cdot h'^{-\frac{1}{2}}\right],
\]
depending only on these data, and on $\beta$ only through $C_{\beta}$, $\mathbf{D}$ and $D_{\beta}$, with the following property. Let $0<e\le e_{\sharp}$, let $R^{*}=R^{*}(e)$ be the torus of hypothesis~\ref{def:hypothesis-reference-torus}, and let $N\ge N_{\mathrm{min}}(R^{*},e)$. Then the centred placement of Section~\ref{sec:probes-discretisation}, described in Section~\ref{sec:probes-mirror-copies}, is admitted; one has
\[
2\rho_{\beta}+3\eta\le l_c;
\]
and the conditions~\eqref{eq:four-jets-sm1}--\eqref{eq:four-jets-sm8} hold for every admitted placement.
\end{lemma}

\begin{proof}
Put $v:=v_{*}$. By Lemma~\ref{lem:scales-stop-level}(b), and since $e_K=e\cdot a_K^{\frac{1}{2}}$ with $a_K=h'$ (Section~\ref{sec:running-quantities-tori}),
\[
v\ge v_K=-\log e_K=-\log\left(e\cdot h'^{\frac{1}{2}}\right),
\]
and $v_K$ tends to $\infty$ as $e$ tends to $0$. As $v_k=-\log e_k$ for every level $k$ (Section~\ref{sec:running-quantities-tori}), one has $e_{*}=\exp(-v)$. Lemma~\ref{lem:scales-stop-level}(a), (b) gives
\[
M\cdot v^r\le l_{*}=M\cdot r_{*}<L\cdot M\cdot v^r,
\]
where $r\ge2$ by stage~\ref{it:order-choice-8}.

The numbers $h_L$, $h_G$, $\boldsymbol{\alpha}$, $C_K$, $C_K'$, $C_V$, $C_W$, $g_1$, $c_1$, $\vartheta$, $\mathfrak{m}_{-}$, $K_0$, $C_0$, $C_d$, $M$ and $N^{*}$ of Sections~\ref{sec:standing-hypothesis-scales} and~\ref{sec:constants-3} are constants. The numbers $C_{\beta}$, $\mathbf{D}$, $D_{\beta}$, $C_{\mathrm{fr}}$ and $z_2$ of Sections~\ref{sec:probes-mirror-copies} and~\ref{sec:constants-3} depend only on the constants, on $\bar{m}$ and on $\beta$; they are finite by Lemma~\ref{lem:properties-probes}(c), so that $z_2>0$. As $a_{*}\in\left\{h'/L,h'/L^2\right\}$ by Lemma~\ref{lem:scales-stop-level}(a), the definition of $\mathsf{p}$ gives
\[
(2M+4)^3\le\mathsf{p}\le\left(\frac{D_{\beta}\cdot L^2}{h'}+2M+4\right)^3.
\]
Recall also that $p_0>3r+2$, $p>2p_0$, $0<s<1$ and $0<\varepsilon\le\frac{1}{200}$ (stages~\ref{it:order-choice-1} and~\ref{it:order-choice-8}).

In this proof an assertion is said to hold \emph{for large $v$} if there is a number $\bar{v}$, depending only on the constants of stages~\ref{it:order-choice-1}--\ref{it:order-choice-10}, on $\bar{m}$ and on $C_{\beta}$, $\mathbf{D}$ and $D_{\beta}$, with the following property. Let $e$ be a charge with $e_K\le e_1$, let $N\ge N_{\mathrm{min}}(R^{*},e)$, and suppose that $v\ge\bar{v}$. If the assertion is a condition on a placement, as~\eqref{eq:four-jets-sm2} and~\eqref{eq:four-jets-sm5} are, then it holds for every admitted placement; any other assertion holds for these $e$ and $N$.

\emph{The conditions~\eqref{eq:four-jets-sm1}, \eqref{eq:four-jets-sm3}, \eqref{eq:four-jets-sm6} and the first two conditions of~\eqref{eq:four-jets-sm4}.} Since $e_{*}=\exp(-v)$ and $\bar{\kappa}=6\,e_{*}^{\frac{1}{4}-9\varepsilon}$, the first two conditions of~\eqref{eq:four-jets-sm1}, the three conditions of~\eqref{eq:four-jets-sm3}, the first two conditions of~\eqref{eq:four-jets-sm4} and the condition~\eqref{eq:four-jets-sm6} are, in this order, the inequalities
\begin{align*}
5\,h_L\cdot h_G&\le\exp\left(\left(\frac{1}{4}-\varepsilon\right)\cdot v\right),&4+\frac{M^3}{2}&\le\exp(\varepsilon\cdot v),\\
4\,C_0\cdot l_{*}^3\cdot g_1&\le\exp\left(\left(\frac{1}{4}-\varepsilon\right)\cdot v\right),&4\,l_{*}^3\cdot\left(2C_K+K_0+4\right)&\le\exp(2\varepsilon\cdot v),\\
8&\le\exp\left(\left(\frac{1}{4}-9\varepsilon\right)\cdot v\right),&\frac{48}{\mathfrak{m}_{-}}&\le\exp(2\varepsilon\cdot v),\\
4\,g_1\cdot C_K'&\le\exp(2\varepsilon\cdot v),&2N^{*}\cdot6^{1-s}&\le\exp\left((1-s)\cdot\left(\frac{1}{4}-9\varepsilon\right)\cdot v\right).
\end{align*}
In each of them the right side is $\exp(a\cdot v)$ with a number $a>0$ formed with $\varepsilon$ and $s$, and the left side is a constant or a constant times $l_{*}^3$, where $l_{*}^3<\left(L\cdot M\cdot v^r\right)^3$. For every $a>0$ and every number $c\ge0$, both $c\cdot\exp(-a\cdot v)$ and $c\cdot\left(L\cdot M\cdot v^r\right)^3\cdot\exp(-a\cdot v)$ tend to $0$ as $v\to\infty$. Hence each of these inequalities holds for large $v$. The third condition of~\eqref{eq:four-jets-sm1} is $\boldsymbol{\alpha}\cdot\bar{r}\le\frac{1}{2}$, and
\[
\bar{r}=4\,h_G^2\cdot C_K\cdot\exp(-2\varepsilon\cdot v)+650\,K_0\cdot\exp\left(-\left(\frac{1}{4}-7\varepsilon\right)\cdot v\right)
\]
tends to $0$ as $v\to\infty$; so this condition holds for large $v$ as well.

\emph{The condition~\eqref{eq:four-jets-sm2}, the centred placement and the inequality $2\rho_{\beta}+3\eta\le l_c$.} Let $v\ge1$; then $r_{*}\ge v^r\ge v^2$ by Lemma~\ref{lem:scales-stop-level}(b). For an admitted placement Lemma~\ref{lem:properties-probes}(e) gives
\[
d_{\mathrm{sep}}\ge\frac{M\cdot r_{*}}{8L}-1\ge\frac{M\cdot v^2}{8L}-1,
\]
and the right side tends to $\infty$ as $v\to\infty$. So $d_{\mathrm{sep}}\ge6M+8$ for large $v$, and the second clause of~\eqref{eq:four-jets-sm2} is the hypothesis that the placement is admitted. Next, $l_c=R^{*}/(2L)$ (Section~\ref{sec:cells-reflections}) and $R^{*}=a_{*}\cdot l_{*}$ (Section~\ref{sec:running-quantities-tori}); with $a_{*}\ge h'/L^2$ this gives
\[
l_c=\frac{a_{*}\cdot l_{*}}{2L}\ge\frac{h'\cdot M\cdot v^2}{2L^3},
\]
which tends to $\infty$ as $v\to\infty$. Moreover $\eta=L^{-N}\le1$. Consequently the two inequalities
\[
\rho_{\beta}+\frac{\eta}{2}\le\rho_{\beta}+\frac{1}{2}\le\frac{3}{8}\,l_c,\qquad2\rho_{\beta}+3\eta\le2\rho_{\beta}+3\le l_c,
\]
neither of which is a condition on a placement, hold for large $v$. By Section~\ref{sec:probes-mirror-copies}, the centred placement is admitted whenever the first of them holds.

\emph{The condition~\eqref{eq:four-jets-sm4}.} For $v\ge1$ one has $\tau_s=c_1\cdot l_{*}\ge c_1\cdot M\cdot v^2$; hence $\tau_s\ge1$ for large $v$. For $v\ge0$ one has $e_{*}\le1$, and the bound for $\mathsf{p}$ gives
\[
B_1=16\,g_1\cdot C_K'\cdot\mathsf{p}\cdot e_{*}^{2\varepsilon}\le16\,g_1\cdot C_K'\cdot\left(\frac{D_{\beta}\cdot L^2}{h'}+2M+4\right)^3,
\]
so that, for $v\ge1$,
\[
52\,B_1\cdot\exp(-\tau_s)\le52\cdot16\,g_1\cdot C_K'\cdot\left(\frac{D_{\beta}\cdot L^2}{h'}+2M+4\right)^3\cdot\exp\left(-c_1\cdot M\cdot v^2\right).
\]
As $c_1\cdot M>0$, the right side is at most $\exp\left(-\left(\frac{1}{4}-9\varepsilon\right)\cdot v\right)=e_{*}^{\frac{1}{4}-9\varepsilon}$ for large $v$. Finally, by the definitions of Section~\ref{sec:constants-3},
\begin{multline*}
B_0=C_{\mathrm{fr}}\cdot\exp(-4\varepsilon\cdot v)+16\exp(6)\cdot C_d\cdot M^3\cdot K_0\cdot\mathsf{p}\cdot\exp(-r_{*})\\
+16\,g_1\cdot C_K'\cdot\mathsf{p}\cdot\exp(-2\varepsilon\cdot v).
\end{multline*}
Since $\mathsf{p}$ is bounded by the number displayed above and $r_{*}\ge v^2$ for $v\ge1$, the number $B_0$ tends to $0$ as $v\to\infty$; so $B_0\le\frac{1}{4}$ for large $v$. The first two conditions of~\eqref{eq:four-jets-sm4} have been treated above.

\emph{The condition~\eqref{eq:four-jets-sm5}.} Let the placement be admitted. If $d_{\mathrm{sep}}=\infty$, then $\varpi_{\mathrm{sep}}=0$, and~\eqref{eq:four-jets-sm5} holds. Otherwise the definition of $\varpi_{\mathrm{sep}}$, the lower bound for $d_{\mathrm{sep}}$ obtained above and $\vartheta>0$ give, for $v\ge1$,
\[
\log\left(2\varpi_{\mathrm{sep}}\right)\le\left(N^{*}+1\right)\log2+\log\left(C_W\cdot\mathsf{p}\right)-\vartheta\cdot\left(\frac{M\cdot v^2}{8L}-4M-7\right),
\]
whereas
\[
\log\left(\bar{\kappa}^{N^{*}}\right)=N^{*}\cdot\left(\log6-\left(\frac{1}{4}-9\varepsilon\right)\cdot v\right)\ge-\frac{N^{*}\cdot v}{4}\qquad(v\ge0).
\]
By the bound for $\mathsf{p}$, the right side of the first of these two displays is at most
\[
\left(N^{*}+1\right)\log2+\log C_W+3\log\left(\frac{D_{\beta}\cdot L^2}{h'}+2M+4\right)+\vartheta\cdot(4M+7)-\vartheta\cdot\frac{M\cdot v^2}{8L},
\]
a polynomial of degree two in $v$ with a negative leading coefficient, whose coefficients depend only on the constants, on $h'$ and on $D_{\beta}$. It is at most $-\frac{N^{*}\cdot v}{4}$ for large $v$. Hence $2\varpi_{\mathrm{sep}}\le\bar{\kappa}^{N^{*}}$ for large $v$.

\emph{The condition~\eqref{eq:four-jets-sm7}.} By Section~\ref{sec:running-quantities-tori}, $p_{*}=v^p$ and $p_{0*}=v^{p_0}$. If
\[
\frac{\left(p_{*}-1\right)^2}{C_0^2}\ge p_{0*},
\]
then the second term in the definition of $\varpi_0$ is at most $6\,l_{*}^3\cdot\exp\left(-\frac{p_{0*}}{2}\right)$, and therefore
\[
\varpi_0\le24\,l_{*}^3\cdot\exp\left(-\frac{v^{p_0}}{2}\right)\le24\left(L\cdot M\cdot v^r\right)^3\cdot\exp\left(-\frac{v^{p_0}}{2}\right).
\]
The hypothesis of this implication, $\left(v^p-1\right)^2\ge C_0^2\cdot v^{p_0}$, holds for large $v$, because $2p>p_0$ (as $p>2p_0$). The right side of the last display tends to $0$ as $v\to\infty$; so $\varpi_0\le\frac{1}{4}$ for large $v$. Next, $t_2=e\cdot z_2$, $t_x=e_{*}^{3\varepsilon}\cdot\left(1+h'^{\frac{1}{2}}\cdot C_{\beta}\right)^{-1}$, and $e=e_{*}\cdot a_{*}^{-\frac{1}{2}}$ by Lemma~\ref{lem:scales-stop-level}(a). Hence, as $a_{*}\ge h'/L^2$,
\begin{align*}
\frac{t_2}{t_x}&=z_2\cdot\left(1+h'^{\frac{1}{2}}\cdot C_{\beta}\right)\cdot a_{*}^{-\frac{1}{2}}\cdot e_{*}^{1-3\varepsilon}\\
&\le z_2\cdot\left(1+h'^{\frac{1}{2}}\cdot C_{\beta}\right)\cdot L\cdot h'^{-\frac{1}{2}}\cdot\exp\left(-(1-3\varepsilon)\cdot v\right),
\end{align*}
and the right side tends to $0$ as $v\to\infty$. So $8N^{*}\cdot t_2\le t_x$ for large $v$.

\emph{The condition~\eqref{eq:four-jets-sm8}.} By the definitions of $\boldsymbol{\lambda}_L$, $\varpi_L$ and $M_L$ in Section~\ref{sec:constants-3},
\[
\log\boldsymbol{\lambda}_L=\frac{1}{N^{*}}\left(\frac{1}{6}\log\varpi_0+\frac{\mathbf{D}}{6}+\frac{2}{3}\log M_L\right)+\log5+4\log\left(\frac{2N^{*}\cdot\mathbf{r}}{t_2}\right).
\]
For $v\ge0$ one has $t_x\le e_{*}^{3\varepsilon}\le1$, hence $\mathbf{r}=t_x/4\le\frac{1}{4}$; with $t_2=e\cdot z_2$ this gives
\[
\frac{\mathbf{r}}{t_2}\le\frac{1}{4e\cdot z_2}.
\]
As $e=e_{*}\cdot a_{*}^{-\frac{1}{2}}$ and $a_{*}\le h'$ (Lemma~\ref{lem:scales-stop-level}(a)),
\[
\log\frac{1}{e}=v+\frac{1}{2}\log a_{*}\le v+\frac{1}{2}\log h'.
\]
Therefore
\[
4\log\left(\frac{2N^{*}\cdot\mathbf{r}}{t_2}\right)\le4\log\left(\frac{N^{*}}{2z_2}\right)+4\log\frac{1}{e}\le4v+4\log\left(\frac{N^{*}}{2z_2}\right)+2\log h'.
\]
By the bound for $\varpi_0$ obtained in the treatment of~\eqref{eq:four-jets-sm7}, for large $v$
\begin{multline*}
\frac{1}{6}\log\varpi_0+\frac{\mathbf{D}}{6}+\frac{2}{3}\log M_L\\
\le-\frac{v^{p_0}}{12}+\frac{r}{2}\log v+\frac{1}{6}\log\left(24\,L^3\cdot M^3\right)+\frac{\mathbf{D}}{6}+\frac{2}{3}\log M_L\le-\frac{v^{p_0}}{13};
\end{multline*}
the last inequality holds for large $v$ because the difference of $\frac{v^{p_0}}{12}$ and $\frac{v^{p_0}}{13}$ is $\frac{v^{p_0}}{156}$, which exceeds $\frac{r}{2}\log v$ plus any given number when $v$ is large, and $M_L=2+\exp\left(\frac{\mathbf{D}}{2}\right)$ depends on $\mathbf{D}$ only. It follows that, for large $v$,
\[
\log\boldsymbol{\lambda}_L\le-\frac{v^{p_0}}{13N^{*}}+4v+\log5+4\log\left(\frac{N^{*}}{2z_2}\right)+2\log h',
\]
the sum of the last three terms being a number that depends only on the constants, on $h'$ and on $z_2$. On the other hand, since $0<s<1$,
\[
\log\left(\bar{\kappa}^s\right)=s\cdot\left(\log6-\left(\frac{1}{4}-9\varepsilon\right)\cdot v\right)\ge-\frac{v}{4}\qquad(v\ge0).
\]
By stage~\ref{it:order-choice-8}, $p_0>3r+2\ge2$; so $v^{p_0}\ge v^2$ for $v\ge1$, and the right side of the bound for $\log\boldsymbol{\lambda}_L$ is at most $-\frac{v^2}{13N^{*}}+4v$ plus a number of the said kind. This is at most $-\frac{v}{4}$ for large $v$. Hence $\boldsymbol{\lambda}_L\le\bar{\kappa}^s$ for large $v$.

\emph{The choice of $e_{\sharp}$.} Finitely many assertions have been shown to hold for large $v$. For each of them fix a number as in the definition, and let $\bar{v}$ be the largest of these numbers. The property required in the definition persists when the number is increased; hence it holds with $\bar{v}$ for each of these assertions, and $\bar{v}$ depends only on the constants of stages~\ref{it:order-choice-1}--\ref{it:order-choice-10}, on $\bar{m}$ and on $C_{\beta}$, $\mathbf{D}$ and $D_{\beta}$. Put
\[
e_{\sharp}:=h'^{-\frac{1}{2}}\cdot\min\left\{e_1,\exp\left(-\bar{v}\right)\right\}.
\]
As $e_1>0$ (stage~\ref{it:order-choice-10}), one has $0<e_{\sharp}\le e_1\cdot h'^{-\frac{1}{2}}$, and $e_{\sharp}$ depends only on the data named in the statement. Let $0<e\le e_{\sharp}$ and $N\ge N_{\mathrm{min}}(R^{*},e)$. Then $e_K=e\cdot h'^{\frac{1}{2}}\le e_1$, as hypothesis~\ref{def:hypothesis-reference-torus} requires, and
\[
v=v_{*}\ge v_K=-\log\left(e\cdot h'^{\frac{1}{2}}\right)\ge\bar{v}.
\]
Therefore those of the assertions that are conditions on a placement hold for every admitted placement, and the others hold for these $e$ and $N$. In particular $\rho_{\beta}+\frac{\eta}{2}\le\frac{3}{8}\,l_c$, so that the centred placement is admitted (Section~\ref{sec:probes-mirror-copies}); moreover $2\rho_{\beta}+3\eta\le l_c$, and the conditions~\eqref{eq:four-jets-sm1}--\eqref{eq:four-jets-sm8} hold for every admitted placement. Theorem~\ref{thm:smallness-conditions-small} follows.
\end{proof}

\section{The free coefficient, chessboard estimates and the proof of the main theorem}\label{sec:free-coefficient-chessboard}

In this section Theorem~\ref{thm:main-theorem} and Corollaries~\ref{cor:ergodic-components} and~\ref{cor:zero-harmonic-law} are proved, and with them Theorems~\ref{thm:bare-mass-shifts} and~\ref{thm:free-coefficients-lattice} of Section~\ref{sec:free-coefficients-bare}. The three theorems of Section~\ref{sec:free-coefficients-cumulant} contain Estimate~6 and the lattice symmetries of Estimate~2 (Theorem~\ref{thm:free-coefficient-1}, proved in Sections~\ref{sec:discretised-probes-lattice}--\ref{sec:coefficient-fourth-order}), Estimate~3 (Theorem~\ref{thm:free-coefficient-2}, proved in Sections~\ref{sec:trace-formula-chessboard} and~\ref{sec:partition-functions-cubic}) and Estimates~7 and~8 (Theorem~\ref{thm:free-coefficient-3}, proved in Section~\ref{sec:proof-theorem}); Part~I of Theorem~\ref{thm:main-theorem} is proved in Section~\ref{sec:admissibility-limit-points}, Part~II in Section~\ref{sec:proof-part-ii-2} and the corollaries in Section~\ref{sec:proofs-corollaries-main}, and Section~\ref{sec:examples-counterexamples} gives examples and counterexamples.

\subsection{Notation and definitions}\label{sec:notation-definitions}

The definitions of Section~\ref{sec:setup} are recalled here in the same notation (constants and scales, lattices and the test map, the Wilson determinant and the two lattice laws, laws on $\mathcal{S}'$, the free coefficient), and the definitions that the proofs require in addition are made: the cells and mirror copies in Section~\ref{sec:reference-tori-cells}, the discretisation and placement of the probes in Section~\ref{sec:probes-discretisation}, and the numbers attached to the reference tori in Section~\ref{sec:numbers-attached-reference}. The objects that Sections~\ref{sec:renormalisation-group-scheme}--\ref{sec:all-flux-functional} define in a notation of their own, among them the Wilson operator, the all-flux functional, the harmonic-averaged law, the cells, the hypothesis~\eqref{eq:setup-H} and the ergodic decomposition, are identified with those of this subsection.

\subsubsection{Letters, norms, Dirac matrices}\label{sec:letters-norms-dirac}

The letter $e$ denotes the charge, and $\exp$ the exponential function (Definition~\ref{def:notation}(a)); phase factors are also written in the form $\exp\left(\pm i(\cdot)\right)$, and $\mathbf{e}:=\exp(1)$. Let $\log^{+}\colon[0,\infty)\to[0,\infty)$ be given by $\log^{+}0:=0$ and $\log^{+}x:=\max\left(0,\log x\right)$ for $x>0$.

As indices, $\mu$ and $\nu$ denote directions; otherwise $\mu$ denotes a law on $\mathcal{S}'$ and $\nu$ a measure on a set of laws, while $\boldsymbol{\mu}_N$ denotes a lattice mass. Where the letter $m$ ranges over $\{2,4\}$, it is at the same time the degree of the polynomial $H_m$ of Section~\ref{sec:probes-discretisation}, the order of the derivatives $\partial_{t_c}^m$ of Section~\ref{sec:numbers-attached-reference} and $N_0^{(m)}(0)$ of Theorem~\ref{thm:free-coefficient-3}, and the order of the cumulant $\kappa_m$; the random variables $H_m(X)$ and $H_m(X_c)$ are the Wick powers of order $m$ of $X$ and of $X_c$. In $M=L^m$ the letter $m$ is the exponent of $M$, and it keeps this meaning where it stands alone, as in the formula for $T^{*}(e)$ in Theorem~\ref{thm:free-coefficient-3} and in its proof in Section~\ref{sec:proof-theorem}. In Sections~\ref{sec:estimates-momentum-space}, \ref{sec:lattice-torus} and~\ref{sec:torus-whole-space} the letter $m$ denotes the mass $\bar{m}$, and $c_m$ there denotes the constant defined in the proof of Theorem~\ref{thm:free-coefficient-1}(c3). Elsewhere $m$ is an index or an object that is declared where it occurs. The integer $m_{\mathrm{p}}$ is the integer with $\mathrm{per}=2m_{\mathrm{p}}+1$, where $\mathrm{per}=n_{*}/L$ is the number of Section~\ref{sec:reference-tori-cells}; equivalently, $m_{\mathrm{p}}$ is the integer part of $\mathrm{per}/2=l_c/\eta=m_{\mathrm{p}}+\frac{1}{2}$. The integer $m_{\mathrm{w}}$ is the integer part of a cell width that is a half-integer (Definition~\ref{def:decomposition-torus-cells}). The numbers $\mathfrak{e}$, $\mathfrak{e}_2$ and $\sigma$ of Section~\ref{sec:constants-exponents-scales} are exponents of the charge.

The symbol $N_0$ has three meanings. First, $N_0(e)$ is the threshold of Theorem~\ref{thm:free-coefficient-3}; it is always written with its argument. Second, $N_0(t)$, also written $N_0(\cdot)$, is the free ratio of Section~\ref{sec:free-fourth-order}, with the derivatives $N_0^{(m)}(0)$ (Definition~\ref{def:notation}(f); this is the notation of Section~\ref{sec:averaged-law-partition}); outside Section~\ref{sec:trace-formula-chessboard} the free ratio is also written $N_0$ without an argument at those places at which the function itself, and not one of its values, is meant, as in $\log N_0$. Third, in Section~\ref{sec:trace-formula-chessboard} the symbol $N_0$ without an argument is the number $N_{\mu}$ of Definition~\ref{def:decomposition-torus-cells} for $\mu=0$, that is, the number of slabs; neither the threshold nor the free ratio occurs in Section~\ref{sec:trace-formula-chessboard}. The three meanings are thus told apart by the argument and, for the symbol without an argument, by the section in which it stands.

For a vector $x$, $|x|=|x|_{\infty}$ is its sup norm and $|x|_2$ its Euclidean norm. For a multi-index $\alpha=\left(\alpha_1,\alpha_2,\alpha_3\right)\in\{0,1,2,\dots\}^3$, however, $|\alpha|$ denotes the order $\alpha_1+\alpha_2+\alpha_3$ of $\alpha$ and not its sup norm, and $\partial^{\alpha}:=\partial_1^{\alpha_1}\partial_2^{\alpha_2}\partial_3^{\alpha_3}$, where $\partial_{\mu}$ is the partial derivative with respect to the coordinate $x_{\mu}$ of $\mathbb{R}^3$. For a complex $4\times4$ or $2\times2$ matrix, $|\cdot|$ is the operator norm with respect to the Euclidean norm, and $\operatorname{tr}$ is the trace; $\operatorname{Tr}$ denotes the trace of an operator.

Let $\sigma_1,\sigma_2,\sigma_3$ be the Pauli matrices. Write $\mathbb{C}^4=\mathbb{C}^2\otimes\mathbb{C}^2$, the first factor indexing the two blocks of size two, and put
\[
\gamma_{\mu}:=\operatorname{diag}\left(\sigma_{\mu},-\sigma_{\mu}\right)=\sigma_3\otimes\sigma_{\mu}\quad(\mu=1,2,3),\qquad\hat{\Gamma}:=\sigma_1\otimes1_2;
\]
thus the two diagonal blocks of $\hat{\Gamma}$ vanish, and its two off-diagonal blocks are the unit matrix $1_2$. The matrix $\gamma_{\mu}$ is the matrix $\gamma_{\mu-1}$ of Section~\ref{sec:dirac-matrices-lattice}, where the directions are numbered $0,1,2$; the matrix $\hat{\Gamma}$ is the matrix $\gamma_3$ of Section~\ref{sec:dirac-matrices-lattice} and the matrix $\gamma_4$ of Section~\ref{sec:wilson-operator-lattice}. Since the Pauli matrices are Hermitian and satisfy $\sigma_j\sigma_k+\sigma_k\sigma_j=2\delta_{jk}$, the matrices $\gamma_{\mu}$ and $\hat{\Gamma}$ are Hermitian, and
\[
\gamma_{\mu}\gamma_{\nu}+\gamma_{\nu}\gamma_{\mu}=2\delta_{\mu\nu},\qquad\hat{\Gamma}\gamma_{\mu}=-\gamma_{\mu}\hat{\Gamma},\qquad\hat{\Gamma}^2=1\qquad(\mu,\nu\in\{1,2,3\}).
\]
For $v\in\mathbb{R}^3$ put $\gamma\cdot v:=\sum_{\mu}v_{\mu}\gamma_{\mu}$; then $\left(\gamma\cdot v\right)^2=|v|_2^2$.

\subsubsection{Constants, exponents, scales}\label{sec:constants-exponents-scales}

The constants are those of stages~\ref{it:order-choice-1}--\ref{it:order-choice-10} of Section~\ref{sec:order-choice}, where they are chosen one stage after the other; they satisfy the conditions of Section~\ref{sec:conditions}. Among these constants are the odd integer $L$, the numbers
\[
N^{*}:=8L^3,\qquad a_L:=\frac{L^3-1}{(L-1)^3},
\]
the integer $m_0$, the integer $M=L^m\ge L^{m_0+2}$, the integer $r\ge2$, a number $\varepsilon\in\left(0,\frac{1}{200}\right]$, and a number $s\in(0,1)$ such that
\[
\sigma:=s\left(\frac{1}{4}-9\varepsilon\right)-12\varepsilon-\frac{2a_L}{N^{*}}>0 .
\]
The exponents of the charge are $\sigma$ and the two numbers
\[
\mathfrak{e}:=s\left(\frac{1}{4}-9\varepsilon\right)-12\varepsilon,\qquad\mathfrak{e}_2:=6\varepsilon;
\]
by the definition of $\sigma$, one has
\[
\sigma=\mathfrak{e}-\frac{2a_L}{N^{*}} .
\]
Further, $e_1>0$ is the number of Section~\ref{sec:constants-conditions-torus}, and $\varepsilon_5\colon(0,\infty)\to(0,\infty)$ is the function
\[
\varepsilon_5(x):=312\,\log\left(\left(1-\exp\left(-\frac{x}{4}\right)\right)^{-1}\right).
\]
By Theorem~\ref{thm:fermion-operators-weakened} one has $\mathfrak{m}_0=1$. Since $m_0$ is the least integer $\ge0$ with $L^{-m_0}\le\mathfrak{m}_0$ (Definition~\ref{def:mass-bound-stop}), it follows that $m_0=0$, and the number $h'$ defined below equals $h$; the letters $m_0$ and $h'$ are retained in the formulas.

Let $0<\bar{m}\le1$ and $e>0$. Put
\begin{align*}
&h:=L^{\left\lfloor\log_L\bar{m}^{-1}\right\rfloor},\qquad h':=h\cdot L^{-m_0},\qquad e_K:=e\cdot h'^{\frac{1}{2}},\qquad v_K:=-\log e_K,\\
&r_K=r_K(e):=\min\left\{L^i:\ i\ge0\text{ an integer},\ L^i\ge v_K^r\right\},\qquad\mathfrak{m}_K:=\bar{m}\cdot h',\\
&R_{\mathrm{top}}(e):=h'\cdot M\cdot r_K(e),\qquad R^{*}(e):=\frac{R_{\mathrm{top}}(e)}{L},\\
&\mathrm{Tor}(e):=\left\{h'\cdot L^T:\ T\ge0\text{ an integer},\ h'\cdot L^T\le R_{\mathrm{top}}(e)\right\}.
\end{align*}
As $h\le\bar{m}^{-1}<L\cdot h$ by the definition of $h$, one has $L^{-m_0-1}<\mathfrak{m}_K\le L^{-m_0}$. For $R\in\mathrm{Tor}(e)$ the numbers $N_{\mathrm{min}}(R,e)$, $k^{*}(R)$, $a_{*}$, $e_{*}$ and $l_{*}$ are those of Section~\ref{sec:running-quantities-tori} and Definition~\ref{def:tori-stop-level}. If $e_K\le e_1$ and $N\ge N_{\mathrm{min}}\left(R_{\mathrm{top}}(e),e\right)$, then $P_N(e)=\left(m_{0,N},\boldsymbol{\varepsilon}_{0,N}\right)$ denotes the pair of Definition~\ref{def:tuned-pair} (Section~\ref{sec:single-pair-all}). Finally let
\[
e_G=e_G(\bar{m}):=\sup\left\{e'>0:\ e'\cdot h'^{\frac{1}{2}}\le e_1,\ e'\cdot h'^{\frac{1}{2}}\log h'\le1\right\}.
\]
Each of the two conditions in this definition, if it holds for some $e'>0$, holds for every smaller positive number: the left side of the first increases with $e'$, and the left side of the second is non-decreasing in $e'$ if $\log h'\ge0$ and negative if $\log h'<0$. Both conditions hold for all sufficiently small $e'>0$, both are preserved under limits in $e'$, and the first fails for $e'>e_1\cdot h'^{-\frac{1}{2}}$. Hence $0<e_G\le e_1\cdot h'^{-\frac{1}{2}}$, and both conditions hold for every $e'\in\left(0,e_G\right]$.

\paragraph{The number $e_2$.}
The number $e_2>0$ of Theorem~\ref{thm:main-theorem}, which is fixed in stage~\ref{it:order-choice-11}, is chosen so small that the following four conditions hold. First,
\begin{equation}\label{eq:free-coefficient-1}
e_2\le e_G(\bar{m}).
\end{equation}
Besides $e\cdot h'^{\frac{1}{2}}\le e_1$, the definition of $e_G$ contains the condition~\eqref{eq:renormalisation-group-18}, that is $e\cdot h'^{\frac{1}{2}}\log h'\le1$; under these two conditions Proposition~\ref{prop:consequences} gives the inequality~\eqref{eq:free-coefficient-6} for $R^{*}=R^{*}(e)$. Second, let $\beta$ and $x_0$ be the data of Section~\ref{sec:probes}, and let $\tau_{x_0}$ be the translation by $x_0$ (Section~\ref{sec:lattices-tests-function}); then
\begin{equation}\label{eq:free-coefficient-2}
e_2\le e_{\sharp}(\beta')\qquad\text{for the three probes }\beta'=\beta,\ \beta+\tau_{x_0}\beta,\ \beta-\tau_{x_0}\beta,
\end{equation}
where $e_{\sharp}(\beta')$ is the number of Lemma~\ref{lem:threshold-charge} for the probe $\beta'$. Third,
\begin{equation}\label{eq:free-coefficient-3}
\frac{C_Y+1}{6\log\left(e_2^{-1}\right)}+\frac{\mathfrak{e}_2}{3}\le1,
\end{equation}
with the constant $C_Y$ of Step~1 of the proof of Theorem~\ref{thm:main-theorem} in Section~\ref{sec:proof-part-ii-2}; in that step this condition gives $\check{X}\le2$ for the number $\check{X}$ of Theorem~\ref{thm:free-coefficient-2}. Fourth,
\begin{equation}\label{eq:free-coefficient-4}
\frac{5}{2}\,C_{\mathrm{kap}}\cdot e_2^{\sigma}\le\frac{b(f)}{4},
\end{equation}
where $f$ is the test two-form of Section~\ref{sec:probes}, $b(f)$ is its free coefficient (Section~\ref{sec:free-fourth-order}), and $C_{\mathrm{kap}}$ is the constant of Theorem~\ref{thm:main-theorem}.

\subsubsection{Lattices, tests, function spaces}\label{sec:lattices-tests-function}

For integers $N,T\ge0$ let
\[
\eta:=L^{-N},\qquad R:=h\cdot L^T,\qquad n:=\frac{R}{\eta},\qquad\Lambda_{\eta,R}:=\eta\mathbb{Z}^3/R\mathbb{Z}^3;
\]
here $n=L^{\left\lfloor\log_L\bar{m}^{-1}\right\rfloor+T+N}$ is a power of the odd integer $L$, hence odd. The tori $\Lambda_{\eta,R}$ of this form are called \emph{aligned}. The following objects are defined as in Sections~\ref{sec:tori-cells-cochains} and~\ref{sec:test-forms-test}--\ref{sec:hermite-sobolev-spaces} and in the paragraphs of notation at the beginning of Sections~\ref{sec:properties-limit-points} and~\ref{sec:moments-tensor-tests}:
\begin{itemize}
\item the $k$-cells $\boldsymbol{\sigma}=(x;I)$, with $x\in\Lambda_{\eta,R}$ and $I\subset\{1,2,3\}$, $\#|I|=k$, their centres $c(\boldsymbol{\sigma})$ and their averaging cubes $Q(\boldsymbol{\sigma})=c(\boldsymbol{\sigma})+\left[-\frac{\eta}{2},\frac{\eta}{2}\right)^3$;
\item the spaces $C^k=C^k(\Lambda_{\eta,R})$ of real $k$-cochains, with the pairing $\langle u,v\rangle_{\eta}=\eta^3\sum_{\boldsymbol{\sigma}}u(\boldsymbol{\sigma})\cdot v(\boldsymbol{\sigma})$ and the norm $|u|_{\eta}=\langle u,u\rangle_{\eta}^{\frac{1}{2}}$;
\item the coboundary $d$, the operator $d_{\eta}=\eta^{-1}\cdot d$, and the adjoint $d_{\eta}^{*}$ of $d_{\eta}$ with respect to $\langle\cdot,\cdot\rangle_{\eta}$;
\item the test map $f\mapsto f_{\eta}$ of \eqref{eq:continuum-infinite-2}, whose value at a plaquette is the average, over the averaging cube of the plaquette, of the corresponding component of the periodisation of $f$; and $X_{\eta,R}(f)=\langle f_{\eta},F\rangle_{\eta}$ for $F\in C^2$;
\item the periodic step function $F^{\sharp}$ of \eqref{eq:continuum-infinite-3}, for which $F^{\sharp}(f)=X_{\eta,R}(f)$;
\item the actions on tests $(\tau_bf)(y)=f(y-b)$, $b\in\mathbb{R}^3$, and $U_r$, $r\in O_{\mathrm{cub}}=O(3)\cap\mathrm{GL}(3,\mathbb{Z})$, the operators $\mathsf{R}_{\nu}=U_{r_{\nu}}$, where $r_{\nu}$ is the reflection $y_{\nu}\mapsto-y_{\nu}$, and the corresponding actions $v_{*}$, $v\in\eta\mathbb{Z}^3$, $r_{*}$ and $\hat{\mathsf{R}}_{\nu}$ on cochains;
\item the sets $\mathcal{P}^{+}_{\nu}$ of plaquettes of Section~\ref{sec:admissible-families-lattice};
\item the Hermite--Sobolev spaces $H_s$, $s\in\mathbb{R}$, and the cylinder $\sigma$-algebra of $\mathcal{S}'=\mathcal{S}'(\mathbb{R}^3;\Lambda^2)$;
\item the spaces $\mathscr{S}_0:=\mathbb{C}$ and $\mathscr{S}_n:=\mathcal{S}\left(\mathbb{R}^{3n};(\Lambda^2)^{\otimes n}\right)$ for $n\ge1$, the latter with the seminorms $q_m$; a tensor product one of whose factors is a number $\Psi_0\in\mathscr{S}_0$ is the product of the other factor with $\Psi_0$;
\item the free Maxwell covariance $C_0(f,g)=\langle f,\Pi_cg\rangle_{L^2}$ of $f,g\in L^2(\mathbb{R}^3;\Lambda^2)$; here $\Pi_c$ denotes the orthogonal projection of $L^2(\mathbb{R}^3;\Lambda^2)$ such that, for every $g$,
\[
\widehat{\ast\Pi_cg}(p)=\left(1-\frac{p\,p^T}{|p|_2^2}\right)\widehat{\ast g}(p)\qquad\text{for almost every }p\in\mathbb{R}^3,
\]
with $\ast g:=\left(g_{23},-g_{13},g_{12}\right)$, the dual vector field of the two-form $g$, and with the Fourier transform of a vector field taken componentwise.
\end{itemize}
The norm of $L^2(\mathbb{R}^3;\Lambda^2)$ is
\[
\|f\|_2=\left(\sum_{\mu<\nu}\int_{\mathbb{R}^3}\mathrm{d}x\,\left|f_{\mu\nu}(x)\right|^2\right)^{\frac{1}{2}}.
\]
For a two-form $\beta$ on $\mathbb{R}^3$ with square integrable components and for $s\in\mathbb{R}$ let
\[
\|\beta\|_{H^s}^2:=\sum_{\mu<\nu}\int_{\mathbb{R}^3}\frac{\mathrm{d}p}{(2\pi)^3}\,\left(1+|p|_2^2\right)^s\left|\hat{\beta}_{\mu\nu}(p)\right|^2\in[0,\infty],
\]
where $\hat{g}(p)=\int_{\mathbb{R}^3}\mathrm{d}x\,\exp(-ip\cdot x)\,g(x)$ denotes the Fourier transform. Let $i\ge0$ be an integer, and let $g$ be a function on $\mathbb{R}^3$ with finitely many components $g_1,\dots,g_{\ell}$, each of class $C^i$; in particular $g$ may be a scalar function ($\ell=1$), a vector field or a one-form, with the components $g_{\nu}$, or a two-form, with the components $g_{\mu\nu}$, $\mu<\nu$. Put
\[
|g|_{C^i}:=\max_{1\le j\le\ell}\ \max_{\substack{\alpha\in\{0,1,2,\dots\}^3\\ \alpha_1+\alpha_2+\alpha_3\le i}}\ \sup_{x\in\mathbb{R}^3}\left|\partial^{\alpha}g_j(x)\right|,
\]
the largest of the sup norms of the partial derivatives of order at most $i$ of the components of $g$. For a two-form $\beta$ with components of class $C^i$ this reads
\[
|\beta|_{C^i}=\max_{\mu<\nu}\ \max_{\substack{\alpha\in\{0,1,2,\dots\}^3\\ \alpha_1+\alpha_2+\alpha_3\le i}}\ \sup_{x\in\mathbb{R}^3}\left|\partial^{\alpha}\beta_{\mu\nu}(x)\right|.
\]
Wherever the torus $k$ of Section~\ref{sec:reference-tori-cells} or another torus that is not required to be aligned occurs, the torus $\Lambda_{\eta,R}=\eta\mathbb{Z}^3/R\mathbb{Z}^3$ and the objects of Section~\ref{sec:tori-cells-cochains} and of \eqref{eq:continuum-infinite-2}, \eqref{eq:continuum-infinite-3} are taken for arbitrary $\eta\in(0,1]$ and $R\ge1$ such that $R/\eta$ is an integer $\ge2$; for such $\eta$ and $R$ they are defined by the same formulas, and Lemmas~\ref{lem:norm-bounds-lattice} and~\ref{lem:lattice-symmetries-cell} hold for them (Section~\ref{sec:test-forms-test}).

\paragraph{Normalisation in lattice units.}
Let $\eta$ and $R$ be as above, with $n=R/\eta\ge2$. The map $x\mapsto\eta^{-1}\cdot x$ is a bijection of $\Lambda_{\eta,R}$ onto the cubic torus $\Lambda_n$ of Section~\ref{sec:wilson-determinant-lattice}, since the group isomorphism $\eta\mathbb{Z}^3\to\mathbb{Z}^3$, $x\mapsto\eta^{-1}\cdot x$, carries $R\mathbb{Z}^3=\eta\cdot n\mathbb{Z}^3$ onto $n\mathbb{Z}^3$; it carries the cell $(x;I)$, spanned at $x$ by the vectors $\eta\cdot e_i$, $i\in I$, onto the cell $\left(\eta^{-1}\cdot x;I\right)$, spanned at $\eta^{-1}\cdot x$ by the vectors $e_i$, $i\in I$, with its orientation. This is the bijection of Section~\ref{sec:lattice-tori}. Through it $C^k(\Lambda_{\eta,R})=C^k(n)$, a cochain $u$ on $\Lambda_{\eta,R}$ being identified with the cochain $(y;I)\mapsto u(\eta\cdot y;I)$ on $\Lambda_n$; a cochain on $\Lambda_{\eta,R}$ is read in lattice units through this identification. The identification commutes with the coboundary $d$, as $\eta^{-1}\cdot\left(x+\eta\cdot e_{\mu}\right)=\eta^{-1}\cdot x+e_{\mu}$, and it preserves the counting pairing and the counting norm, being induced by a bijection of the cells. Moreover $\Lambda_n=\Lambda_{1,n}$, and the identification carries the actions $v_{*}$, $v\in\eta\mathbb{Z}^3$, and $r_{*}$, $r\in O_{\mathrm{cub}}$, of Lemma~\ref{lem:lattice-symmetries-cell} on the cochains of $\Lambda_{\eta,R}$ to the actions $\left(\eta^{-1}\cdot v\right)_{*}$ and $r_{*}$ of the same lemma on the cochains of $\Lambda_{1,n}$, because $\eta^{-1}\cdot(\boldsymbol{\sigma}+v)=\eta^{-1}\cdot\boldsymbol{\sigma}+\eta^{-1}\cdot v$ and $\eta^{-1}\cdot(r\boldsymbol{\sigma})=r\left(\eta^{-1}\cdot\boldsymbol{\sigma}\right)$ for every cell $\boldsymbol{\sigma}$, and the sign $\varepsilon(r,\boldsymbol{\sigma})$ of that lemma depends only on $r$ and on the set of directions of $\boldsymbol{\sigma}$; thus $v_{*}$ becomes the action on the cochains of $\Lambda_n$ of the lattice translation by $\eta^{-1}\cdot v$, and $r_{*}$ that of $r$. In this sense the curvature in lattice units $F^u\in C^2(n)$ of Section~\ref{sec:wilson-determinant-lattice} is an element of $C^2=C^2(\Lambda_{\eta,R})$. For $g\in C^2$ put $g^u:=\eta^{\frac{3}{2}}g$. The physical curvature is $F=\eta^{-\frac{3}{2}}F^u$. Since $\eta^3\cdot\eta^{-\frac{3}{2}}\cdot\eta^{-\frac{3}{2}}=1$ and $\eta^3=\left(\eta^{\frac{3}{2}}\right)^2$, one has
\[
\langle g,F\rangle_{\eta}=\langle g^u,F^u\rangle,\qquad|g|_{\eta}=\|g^u\|,
\]
where on the right $\langle\cdot,\cdot\rangle$ and $\|\cdot\|$ are the counting pairing and the counting norm of cochains (Section~\ref{sec:tori-cells-cochains}).

\subsubsection{The Wilson determinant and the lattice laws}\label{sec:wilson-determinant-lattice}

\paragraph{Tori and cochains in lattice units.}
Let $n=(n_1,n_2,n_3)$ be a triple of integers $n_{\mu}\ge2$; a single integer $n\ge2$ also stands for the cubic triple $(n,n,n)$. Put
\[
\Lambda_n:=\prod_{\mu=1}^3\mathbb{Z}/n_{\mu}\mathbb{Z},\qquad|n|:=n_1n_2n_3,
\]
and let $C^k(n)$ be the real vector space of the $k$-cochains on $\Lambda_n$, with the counting pairing $\langle\cdot,\cdot\rangle$ and the norm $\|\cdot\|$. A real or complex one-cochain $B$, that is, an element of $C^1(n)$ or of $C^1(n)\otimes\mathbb{C}$, is identified with the function on the oriented nearest-neighbour bonds given by $B(x,x+e_{\mu}):=B_{\mu}(x)$ and $B(x',x):=-B(x,x')$; in this way $B$ is also a function on the bonds of $\mathbb{Z}^3$, invariant under the translations by $n_{\mu}e_{\mu}$, $\mu=1,2,3$. A \emph{spinor} is a function $\psi\colon\mathbb{Z}^3\to\mathbb{C}^4$ such that $\psi(x+n_{\mu}e_{\mu})=-\psi(x)$ for all $x\in\mathbb{Z}^3$ and $\mu=1,2,3$; the spinors form a complex vector space of dimension $4|n|$.

\paragraph{The Wilson operator.}
Let $B\in C^1(n)\otimes\mathbb{C}$ be a complex one-cochain and $\boldsymbol{\mu}$ a real number. The operator $D[B]$ on the spinors and the number $\mathrm{Dt}(B)$ are
\begin{multline}\label{eq:free-coefficient-5}
\left(D[B]\psi\right)(x):=3\psi(x)-\sum_{x'\in\mathbb{Z}^3:\ |x'-x|_2=1}\frac{1}{2}\left(1-\gamma\cdot(x'-x)\right)\cdot\exp\left(iB(x,x')\right)\cdot\psi(x'),\\
\mathrm{Dt}(B):=\det\left(D[B]+\boldsymbol{\mu}\right),
\end{multline}
the determinant being taken on the space of the spinors; $D[B]\psi$ is again a spinor, as $B$ is periodic and $\psi$ is antiperiodic. For $q>0$ and a real or complex one-cochain $A$ put $D_A:=D[qA]$. The operator $D_A$ has Wilson parameter $1$; it acts on spinors with four components that are antiperiodic in every direction, and it carries the factor $\exp\left(iqA(x,x')\right)$ on the bond $(x,x')$. It coincides with each of the following operators: the Wilson operator of Section~\ref{sec:lattice-measures-laws}; the operator of Section~\ref{sec:exact-fermion-operators} at level $0$; the operator of \eqref{eq:gaussian-bounds-1} for the character $\chi=\chi^{\mathrm{ap}}$; and the operators of \eqref{eq:continuum-infinite-1}, of Section~\ref{sec:wilson-determinant} and of \eqref{eq:all-flux-1}. Indeed, each of these definitions, like \eqref{eq:free-coefficient-5}, multiplies $\psi(x+e_{\mu})$ by the matrix $\frac{1}{2}\left(1-\gamma_{\mu}\right)$ and by the phase of the bond $(x,x+e_{\mu})$, and $\psi(x-e_{\mu})$ by the matrix $\frac{1}{2}\left(1+\gamma_{\mu}\right)$ and by the inverse of the phase of the bond $(x-e_{\mu},x)$.

\paragraph{The all-flux functional.}
Put $r_q:=2\pi/q$. Let $\mathfrak{F}(n)$ be the class of the Borel functions $\Phi\colon C^2(n)\to\mathbb{C}$ for which there is a constant $c$ with $|\Phi(F)|\le c\cdot\exp\left(c\|F\|\right)$ for all $F\in C^2(n)$. Let $C^2(n;\mathbb{Z})\subset C^2(n)$ be the set of the integer valued two-cochains, $\mathrm{d}A$ the Lebesgue measure of $C^1(n)$, and $[0,r_q)^{\mathrm{links}}$ the set of the $A\in C^1(n)$ with $0\le A_{\mu}(x)<r_q$ for every link, that is, for every positively oriented bond $(x,x+e_{\mu})$. For $\Phi\in\mathfrak{F}(n)$ let, as in Definition~\ref{def:all-flux-functional} and Definition~\ref{def:partition-functions-all},
\begin{multline*}
\boldsymbol{\omega}_n(\Phi):=\int_{[0,r_q)^{\mathrm{links}}}\mathrm{d}A\sum_{\substack{\mathfrak{z}\in C^2(n;\mathbb{Z})\\ d\mathfrak{z}=0}}\Phi\left(dA+r_q\cdot\mathfrak{z}\right)\cdot\exp\left(-\frac{1}{2}\left\|dA+r_q\cdot\mathfrak{z}\right\|^2\right)\cdot\mathrm{Dt}(qA)\\
=\sum_{w\in\mathbb{Z}^3}\boldsymbol{\omega}_n^{(w)}(\Phi),
\end{multline*}
where $\boldsymbol{\omega}_n^{(w)}(\Phi)$ is the same expression with the sum restricted to the $\mathfrak{z}$ of flux $w(\mathfrak{z})=w$. The flux of a closed $\mathfrak{z}\in C^2(n;\mathbb{Z})$ is the vector $w(\mathfrak{z})=\left(w_{\mu\nu}(\mathfrak{z})\right)_{\mu<\nu}\in\mathbb{Z}^3$ with the components
\[
w_{\mu\nu}(\mathfrak{z}):=\sum_{x\in\Lambda_n:\ x_{\iota}=0}\mathfrak{z}_{\mu\nu}(x)\qquad\left(\mu<\nu,\ \iota\in\{1,2,3\}\setminus\{\mu,\nu\}\right),
\]
where $\mathfrak{z}_{\mu\nu}(x)$ is the value of $\mathfrak{z}$ on the plaquette at $x$ spanned by $e_{\mu}$ and $e_{\nu}$; the sum has the same value if the condition $x_{\iota}=0$ is replaced by $x_{\iota}=j$ for any $j\in\mathbb{Z}/n_{\iota}\mathbb{Z}$, because $d\mathfrak{z}=0$ (Definition~\ref{def:all-flux-functional}, \eqref{eq:partition-functions-2}). The argument of $\Phi$ in the definition of $\boldsymbol{\omega}_n$ is the curvature in lattice units, $F^u:=dA+r_q\cdot\mathfrak{z}$; it is a closed two-cochain, as $dF^u=d(dA)+r_q\cdot d\mathfrak{z}=0$. Put
\[
Z_n:=\boldsymbol{\omega}_n(1),\qquad\mathbf{Z}_n:=\boldsymbol{\omega}_n^{(0)}(1).
\]
The series in the definition of $\boldsymbol{\omega}_n$ converges absolutely. Indeed, with $F=dA+r_q\cdot\mathfrak{z}$ the modulus of a summand is at most $c\cdot\exp\left(c^2-\frac{1}{4}\|F\|^2\right)\cdot\sup|\mathrm{Dt}|$, because $c\|F\|-\frac{1}{2}\|F\|^2\le c^2-\frac{1}{4}\|F\|^2$, where the supremum of $|\mathrm{Dt}|$ over the real one-cochains is finite (Theorem~\ref{thm:free-coefficient-1}(d1)); for $A\in[0,r_q)^{\mathrm{links}}$ the norm $\|dA\|$ is bounded and $\|F\|\ge r_q\cdot\|\mathfrak{z}\|-\|dA\|$; and the number of the $\mathfrak{z}\in C^2(n;\mathbb{Z})$ with $\|\mathfrak{z}\|\le y$ is bounded by a polynomial in $y$. If $\boldsymbol{\mu}>-1$, then $Z_n>0$ and $\mathbf{Z}_n>0$ (Section~\ref{sec:summation-all-fluxes}). For $w\in\mathbb{Z}^3$ put
\[
P_n(w):=\frac{\boldsymbol{\omega}_n^{(w)}(1)}{Z_n},\qquad\mathfrak{q}_w(n):=\exp\left(-\frac{r_q^2}{2|n|}\sum_{\mu<\nu}n_{\iota}^2\cdot w_{\mu\nu}^2\right),\qquad\varsigma_n:=\sum_{w\in\mathbb{Z}^3\setminus\{0\}}\mathfrak{q}_w(n),
\]
where in the sum $\iota$ is the direction different from $\mu$ and $\nu$. In the notation of Definition~\ref{def:partition-functions-all} and Section~\ref{sec:decomposition-flux} one has $\boldsymbol{\omega}_n=Z_n[\cdot]$, $\boldsymbol{\omega}_n^{(w)}=Z_{n,w}[\cdot]$ and $\mathfrak{q}_w(n)=\exp\left(-\frac{1}{2}\|h_w\|^2\right)$; in the notation of Definition~\ref{def:all-flux-functional}, $\boldsymbol{\omega}_n=\boldsymbol{\omega}^{\mathrm{af}}_n$.

\paragraph{The two laws.}
The all-flux law~\eqref{eq:setup-af} and the harmonic-averaged law~\eqref{eq:setup-av} on $\Lambda_n$ are the normalised functionals
\[
E^{\mathrm{af}}_n:=\frac{\boldsymbol{\omega}_n}{Z_n},\qquad E^{\mathrm{av}}_n:=\frac{\boldsymbol{\omega}_n^{(0)}}{\mathbf{Z}_n}.
\]
Let $C^1(n)=V_{\mathrm{g}}\oplus V_{\mathrm{h}}\oplus V_{\mathrm{t}}$ be the orthogonal decomposition in which $V_{\mathrm{g}}=dC^0(n)$ consists of the exact cochains, $V_{\mathrm{h}}$ of the constant cochains, and $V_{\mathrm{t}}$ is the orthogonal complement of $V_{\mathrm{g}}\oplus V_{\mathrm{h}}$; let $\sigma$ be the Lebesgue measure of $V_{\mathrm{t}}$; and for $\theta\in[-\pi,\pi]^3$ let $H_{\theta}$ be the constant one-cochain with the components $(H_{\theta})_{\mu}(x)=\theta_{\mu}/(q\cdot n_{\mu})$. By Lemma~\ref{lem:flux-sectors}(b), which is Lemma~\ref{lem:sector-zero-flux}, $E^{\mathrm{av}}_n$ is the law of $F^u=dA$ when $A\in V_{\mathrm{t}}$ is distributed according to the probability measure proportional to
\[
\left(\int_{[-\pi,\pi]^3}\mathrm{d}\theta\,\exp\left(-\frac{1}{2}\|dA\|^2\right)\cdot\det\left(D_{A+H_{\theta}}+\boldsymbol{\mu}\right)\right)\mathrm{d}\sigma(A);
\]
it is the law $\bar{\mathsf{P}}$ of Section~\ref{sec:wilson-operator-lattice}, the law $\overline{\mathsf{E}}$ of \eqref{eq:gaussian-bounds-3} and the law $\bar{\mu}$ of Section~\ref{sec:lattice-measures-laws}.

\paragraph{The laws $P^{\mathrm{nat}}_{N,T}$ and their limit points.}
Let $e>0$ be the charge. For integers $N,T\ge0$ let $\eta:=L^{-N}$, $R:=h\cdot L^T$ and $q:=e\cdot\eta^{\frac{1}{2}}$. A mass assignment is a sequence $N\mapsto m_{0,N}$ of real numbers such that
\[
\boldsymbol{\mu}_N:=\bar{m}\cdot L^{-N}+m_{0,N}>-\frac{1}{2}
\]
for every $N\ge0$; the numbers $m_{0,N}$ do not depend on $T$. For $\mathrm{nat}\in\{\mathrm{af},\mathrm{av}\}$ let $P^{\mathrm{nat}}_{N,T}$ denote the distribution of $F^{\sharp}$, formed with $F=\eta^{-\frac{3}{2}}F^u$, under $E^{\mathrm{nat}}_n$ with $n=h\cdot L^{T+N}=R/\eta$ and $\boldsymbol{\mu}=\boldsymbol{\mu}_N$. Each $P^{\mathrm{nat}}_{N,T}$ is a probability measure on the Borel sets of the Hermite--Sobolev space $H_{-3}$. Multiplying the weight by the vacuum factor $\exp\left(-\boldsymbol{\varepsilon}_{0,N}\cdot n^3\right)$, which is a constant, changes neither of the two laws, since a constant factor cancels against the normalisation. For $h=1$ the pair $(N,T)=(0,0)$ leads to $n=1$, in conflict with the requirement $n_{\mu}\ge2$ in the definition of $\Lambda_n$; no law corresponds to this pair, and it is excluded from all statements. Let $\mu$ be a probability measure on the Borel sets of $H_{-4}$. It is called a limit point of the family $\left(P^{\mathrm{nat}}_{N,T}\right)$ if there are cutoffs $(N_i,T_i)$, $i=1,2,\dots$, with $N_i\to\infty$ and $T_i\to\infty$, such that $P^{\mathrm{nat}}_{N_i,T_i}$ converges weakly to $\mu$; apart from these two limits, the sequences $(N_i)$ and $(T_i)$ are arbitrary (Definition~\ref{def:limit-points}).

\subsubsection{Reference tori, cells and mirror copies}\label{sec:reference-tori-cells}

Let $T^{*}\ge1$ and $N\ge0$ be integers, $\eta:=L^{-N}$, and put
\[
R^{*}:=h\cdot L^{T^{*}},\qquad n_{*}:=\frac{R^{*}}{\eta},\qquad\mathrm{per}:=\frac{n_{*}}{L},\qquad l_c:=\frac{R^{*}}{2L}.
\]
Since $n_{*}=h\cdot L^{T^{*}+N}$ and $T^{*}\ge1$, the number $\mathrm{per}=h\cdot L^{T^{*}+N-1}$ is a power of the odd integer $L$; let $m_{\mathrm{p}}\ge0$ be the integer with $\mathrm{per}=2m_{\mathrm{p}}+1$. Then $l_c=\left(m_{\mathrm{p}}+\frac{1}{2}\right)\eta$.

Let $k\ge1$ be an integer. The \emph{torus $k$} is the cubic lattice torus $\Lambda_{k\cdot n_{*}}$, of $k\cdot n_{*}$ sites along each axis; its physical side is $k\cdot n_{*}\cdot\eta=k\cdot R^{*}$. The two laws and the partition function on it are abbreviated $E^{\mathrm{nat}}_k:=E^{\mathrm{nat}}_{k\cdot n_{*}}$ and $Z_k:=Z_{k\cdot n_{*}}$. If $T\ge T^{*}$, the aligned torus of side $R=h\cdot L^T$ has $R/\eta=L^{T-T^{*}}\cdot n_{*}$ sites in each direction; hence it is the torus $k$ with $k=L^{T-T^{*}}$. The two tori $k=1$ and $k=L$ are the reference tori.

In lattice units, for every direction $\mu$ the circle $\mathbb{Z}/(k\cdot n_{*})\mathbb{Z}$ of the values of $x_{\mu}$ is the union of the $2kL$ intervals
\[
I_{2i}:=\left\{i\cdot\mathrm{per},\dots,i\cdot\mathrm{per}+m_{\mathrm{p}}\right\},\qquad I_{2i+1}:=\left\{i\cdot\mathrm{per}+m_{\mathrm{p}}+1,\dots,(i+1)\cdot\mathrm{per}\right\}\qquad\left(i\in\mathbb{Z}/kL\mathbb{Z}\right),
\]
each of which has $m_{\mathrm{p}}+1$ elements. Thus $I_{2i}$ and $I_{2i+1}$ are disjoint, and between them lies the link plane $\left\{x_{\mu}=i\cdot\mathrm{per}+m_{\mathrm{p}}+\frac{1}{2}\right\}$; whereas $I_{2i+1}\cap I_{2i+2}=\left\{(i+1)\cdot\mathrm{per}\right\}$, so that the slabs $\left\{x_{\mu}\in I_{2i+1}\right\}$ and $\left\{x_{\mu}\in I_{2i+2}\right\}$ intersect in the site plane $\left\{x_{\mu}=(i+1)\cdot\mathrm{per}\right\}$. For $c\in(\mathbb{Z}/2kL\mathbb{Z})^3$ the \emph{cell} $\mathfrak{c}_c$ is
\[
\mathfrak{c}_c:=\prod_{\mu=1}^3I_{c_{\mu}};
\]
the cells are indexed by $(\mathbb{Z}/2kL\mathbb{Z})^3$, so that their number is $(2kL)^3$, which for $k=1$ equals $(2L)^3=N^{*}$; and $\mathfrak{c}_0=\{0,\dots,m_{\mathrm{p}}\}^3$. A plaquette $p$ belongs to $P^o(\mathfrak{c}_c)$ if its four vertices lie in $\mathfrak{c}_c$ and if, for each $\mu\in\{1,2,3\}$ and each end point $t$ of $I_{c_{\mu}}$, the plaquette $p$ is not contained in the plane $\{x_{\mu}=t\}$. Let $\mathfrak{A}(\mathfrak{c}_c)$ consist of those $O\in\mathfrak{F}(k\cdot n_{*})$ for which $O(F^u)$ depends on $F^u$ only through the components $F^u_p$ with $p\in P^o(\mathfrak{c}_c)$, and put $\mathfrak{A}:=\mathfrak{A}(\mathfrak{c}_0)$.

For $c\in(\mathbb{Z}/2kL\mathbb{Z})^3$ let $\mathsf{g}_c$ be the map of the torus $k$ onto itself that acts on the coordinate $x_{\mu}$ by
\[
x_{\mu}\mapsto i\cdot\mathrm{per}+x_{\mu}\quad\text{if }c_{\mu}=2i,\qquad x_{\mu}\mapsto(i+1)\cdot\mathrm{per}-x_{\mu}\quad\text{if }c_{\mu}=2i+1.
\]
Each coordinate map is a translation or a reflection of the circle $\mathbb{Z}/(k\cdot n_{*})\mathbb{Z}$, so that $\mathsf{g}_c$ is a lattice isometry. In every coordinate it carries $I_0$ onto $I_{c_{\mu}}$ (for $c_{\mu}=2i+1$ because $\mathrm{per}-m_{\mathrm{p}}=m_{\mathrm{p}}+1$); hence $\mathsf{g}_c(\mathfrak{c}_0)=\mathfrak{c}_c$. The coordinate maps are products of reflections, as follows. For an integer $l$ let $\varrho_l$ be the map $x_{\mu}\mapsto l\cdot\mathrm{per}-x_{\mu}$ of $\mathbb{Z}/(k\cdot n_{*})\mathbb{Z}$, the reflection in the plane $\left\{x_{\mu}=\frac{l}{2}\mathrm{per}\right\}$; for $l=2i+1$ this plane is the link plane between $I_{2i}$ and $I_{2i+1}$, for $l=2i+2$ the site plane in which the slabs of $I_{2i+1}$ and $I_{2i+2}$ intersect. Induction on $j\ge0$, starting from the empty composition (the identity) for $j=0$ and using $\varrho_{j+1}\left(i\cdot\mathrm{per}+x_{\mu}\right)=(i+1)\cdot\mathrm{per}-x_{\mu}$ for $j=2i$ and $\varrho_{j+1}\left((i+1)\cdot\mathrm{per}-x_{\mu}\right)=(i+1)\cdot\mathrm{per}+x_{\mu}$ for $j=2i+1$, gives
\[
\varrho_j\circ\varrho_{j-1}\circ\cdots\circ\varrho_1(x_{\mu})=\begin{cases}i\cdot\mathrm{per}+x_{\mu}&\text{if }j=2i,\\(i+1)\cdot\mathrm{per}-x_{\mu}&\text{if }j=2i+1.\end{cases}
\]
Consequently, if $j\in\{0,\dots,2kL-1\}$ represents $c_{\mu}$, then $\mathsf{g}_c$ acts on $x_{\mu}$ as the product of the $j$ reflections in the planes $\left\{x_{\mu}=\frac{l}{2}\mathrm{per}\right\}$, $1\le l\le j$, that are met between $I_0$ and $I_{c_{\mu}}$ when $x_{\mu}$ increases, the reflection $\varrho_1$ in the plane adjacent to $I_0$ acting first. Every lattice isometry $g$ acts on cochains by the pull-back $g^{*}$ and by the push-forward $g_{*}:=\left(g^{-1}\right)^{*}$. For $O\in\mathfrak{A}$ and a cell $c$ of the torus $k$, the \emph{mirror copy} of $O$ in $\mathfrak{c}_c$ is
\[
O^{[c]}(F):=\begin{cases}O\left(\mathsf{g}_c^{*}F\right)&\text{if }\sum_{\mu}c_{\mu}\text{ is even},\\ \overline{O\left(\mathsf{g}_c^{*}F\right)}&\text{if }\sum_{\mu}c_{\mu}\text{ is odd}\end{cases}
\]
(the parity of $\sum_{\mu}c_{\mu}$ is well defined, as $2kL$ is even), and
\[
O^{\sharp}:=\prod_cO^{[c]},
\]
the product being taken over all cells $c$ of the torus $k$. If $R^{*}=R^{*}(e)$ (Section~\ref{sec:constants-exponents-scales}), the cells, the maps $\mathsf{g}_c$ and the mirror copies are those of Section~\ref{sec:cells-reference-tori}, where the directions are numbered $0,1,2$.

\subsubsection{Probes and their discretisation}\label{sec:probes-discretisation}

For a smooth function $g$ on $\mathbb{R}^3$, a number $\eta>0$ and $j\in\{1,2,3\}$ let
\[
\delta_{\eta,j}g(x):=\frac{g\left(x+\frac{\eta}{2}e_j\right)-g\left(x-\frac{\eta}{2}e_j\right)}{\eta},\qquad J_{\eta}g(x):=\eta^{-3}\int_{\left[-\frac{\eta}{2},\frac{\eta}{2}\right]^3}\mathrm{d}y\,g(x+y).
\]
A probe is a real two-form $\beta\in C_c^{\infty}(\mathbb{R}^3;\Lambda^2)$ with $\beta\ne0$; its components satisfy $\beta_{\nu\mu}=-\beta_{\mu\nu}$. For a probe $\beta$ and a torus $\Lambda_{\eta,R}$ (Section~\ref{sec:lattices-tests-function}) let
\begin{align*}
&a:=d^{*}\beta,\quad a_{\nu}=-\sum_{\mu}\partial_{\mu}\beta_{\mu\nu};\qquad f:=dd^{*}\beta,\quad f_{\mu\nu}=\partial_{\mu}a_{\nu}-\partial_{\nu}a_{\mu};\\
&a'_{\nu}:=-\sum_{\mu}\delta_{\eta,\mu}\beta_{\mu\nu};\qquad\check{a}_{\eta}:=J_{\eta}a',\ \text{with the components }\check{a}_{\eta,\nu};\qquad f^{(\eta)}_{\mu\nu}:=\delta_{\eta,\mu}a'_{\nu}-\delta_{\eta,\nu}a'_{\mu};\\
&u_{\mu\nu}(x):=\left(J_{\eta}\beta_{\mu\nu}\right)\left(x+\frac{\eta}{2}\left(e_{\mu}+e_{\nu}\right)\right)\quad\left(x\in\Lambda_{\eta,R}\right),\qquad a_{\eta}:=d_{\eta}^{*}u\in C^1(\Lambda_{\eta,R}),
\end{align*}
where in the definition of $u$ the form $\beta$ is periodised with the period $R$ in every direction.

Let $\rho_{\beta}$ be the least number $\rho$ for which there is a point $y\in\mathbb{R}^3$ with $\operatorname{supp}\beta+[-2,2]^3\subset y+[-\rho,\rho]^3$, let $y_{\beta}$ be such a point for $\rho=\rho_{\beta}$, and put
\[
D_{\beta}:=2\rho_{\beta},\qquad Q_{\beta}:=y_{\beta}+[-\rho_{\beta},\rho_{\beta}]^3.
\]
One has $\rho_{\beta}\ge2$, because $\operatorname{supp}\beta$ is not empty, so that $\operatorname{supp}\beta+[-2,2]^3$ contains a cube of side $4$. Let $T^{*}\ge1$ and $N\ge0$ be integers, $\eta=L^{-N}$, and let $l_c$ and $m_{\mathrm{p}}$ be as in Section~\ref{sec:reference-tori-cells}. The probe $\beta$ \emph{fits} $(N,T^{*})$ if
\[
2\rho_{\beta}+3\eta\le l_c.
\]
In this case $\left(m_{\mathrm{p}}+\frac{1}{2}\right)\eta=l_c\ge4+3\eta$, so that $m_{\mathrm{p}}\ge\frac{4}{\eta}+\frac{5}{2}\ge\frac{13}{2}$, as $\eta\le1$; hence $m_{\mathrm{p}}\ge7$.

\paragraph{The placement of the probe.}
Let $\beta$ fit $(N,T^{*})$. Points of $\mathbb{R}^3$ are given in physical coordinates, which are $\eta$ times the lattice coordinates of Section~\ref{sec:reference-tori-cells}. There are vectors $v\in\eta\mathbb{Z}^3$ with
\[
\left|v+y_{\beta}-\frac{l_c}{2}(1,1,1)\right|\le\frac{\eta}{2},
\]
$|\cdot|$ being the sup norm (in each coordinate this asks for a point of $\eta\mathbb{Z}$ in a closed interval of length $\eta$). One such vector $v$ is chosen; it is called the \emph{placement}, or the \emph{centred placement}, of $\beta$. Then $v+Q_{\beta}\subset[\eta,l_c-\eta]^3$, because each coordinate of $v+y_{\beta}$ differs from $\frac{l_c}{2}$ by at most $\frac{\eta}{2}$, and $\rho_{\beta}+\frac{\eta}{2}\le\frac{l_c}{2}-\eta$. Let $a_{\eta}$ be formed with $\tau_v\beta$ in place of $\beta$, on the torus $\Lambda_{\eta,R}$ under consideration, and put
\begin{align*}
&\hat{\mathfrak{a}}:=\eta^{\frac{1}{2}}a_{\eta},\qquad\boldsymbol{\sigma}^2:=\left\|d\hat{\mathfrak{a}}\right\|^2\quad(\boldsymbol{\sigma}\ge0),\qquad X:=\left\langle d\hat{\mathfrak{a}},F^u\right\rangle,\\
&\hat{\mathfrak{a}}_c:=\left(\mathsf{g}_c\right)_{*}\hat{\mathfrak{a}},\qquad X_c:=\left\langle d\hat{\mathfrak{a}}_c,F^u\right\rangle,\\
&H_2(x):=x^2-\boldsymbol{\sigma}^2,\qquad H_4(x):=x^4-6\boldsymbol{\sigma}^2\cdot x^2+3\boldsymbol{\sigma}^4\qquad(x\in\mathbb{R}).
\end{align*}
Here $\hat{\mathfrak{a}}$ is read as a one-cochain on the unit lattice $\Lambda_n$, $n=R/\eta$ (Section~\ref{sec:lattices-tests-function}); $\langle\cdot,\cdot\rangle$ and $\|\cdot\|$ are the counting pairing and the counting norm; and, when $\Lambda_{\eta,R}$ is the torus $k$ of Section~\ref{sec:reference-tori-cells}, $c$ ranges over its cells.

The inequality above need not determine $v$, and the placement may be any one of the vectors that satisfy it. Let $v'\in\eta\mathbb{Z}^3$ be another one. Then $v'-v\in\eta\mathbb{Z}^3$; the cochain $u$ formed with $\tau_{v'}\beta$ is the translate by $v'-v$ of the cochain $u$ formed with $\tau_v\beta$, and with $v'$ in place of $v$ the cochain $\hat{\mathfrak{a}}$ becomes its translate by the lattice vector $(v'-v)/\eta$. It follows that $\boldsymbol{\sigma}^2$, the number $b_{\eta,R}(0)$ of Section~\ref{sec:free-fourth-order} formed with $a_{\eta}$, and the law of $X$ under each of $E^{\mathrm{af}}_n$ and $E^{\mathrm{av}}_n$ are the same for $v$ and for $v'$ (Theorem~\ref{thm:free-coefficient-1}(d5)); and each statement of Section~\ref{sec:free-coefficients-cumulant} is valid whichever of these vectors is taken as the placement.

The same definitions of probes, of $\rho_{\beta}$ and $D_{\beta}$, and of the cochains $u$, $a_{\eta}$, $\hat{\mathfrak{a}}$ and $\hat{\mathfrak{a}}_c$ formed with $\tau_v\beta$ occur in Definition~\ref{def:probes} and in Section~\ref{sec:probes-mirror-copies}.

\subsubsection{Numbers attached to the reference tori}\label{sec:numbers-attached-reference}

Throughout Section~\ref{sec:numbers-attached-reference} lattice units are used, with the following data: an integer $N\ge0$ and $\eta=L^{-N}$; the charge $e>0$ and $q=e\cdot\eta^{\frac{1}{2}}$; and the lattice mass $\boldsymbol{\mu}=\boldsymbol{\mu}_N$. The integer $T^{*}\ge1$ and the numbers $n_{*}$, $Z_1=Z_{n_{*}}$ and $Z_L=Z_{L\cdot n_{*}}$ are those of Section~\ref{sec:reference-tori-cells}. For an integer $n\ge2$ let $\tilde{\gamma}_n$ be the Gaussian probability measure on the space $V_{\mathrm{t}}$ of the cubic torus $n$ whose density with respect to the Lebesgue measure of $V_{\mathrm{t}}$ is proportional to $\exp\left(-\frac{1}{2}\|dA\|^2\right)$, and put
\begin{align*}
\rho_n&:=\frac{1}{(2\pi)^3}\int_{[-\pi,\pi]^3}\mathrm{d}\theta\int_{V_{\mathrm{t}}}\mathrm{d}\tilde{\gamma}_n(A)\,\det\left(D_{A+H_{\theta}}+\boldsymbol{\mu}_N\right),\\
D^{\mathrm{mat}}&:=L^3\log\rho_{n_{*}}-\log\rho_{L\cdot n_{*}},\\
D_L&:=L^3\log Z_1-\log Z_L,\qquad\delta:=\frac{D_L}{(L-1)^3N^{*}}.
\end{align*}
The number $\rho_n$ is the number $\rho_n(q,\boldsymbol{\mu}_N)$ of Section~\ref{sec:lattice-measures-laws} and of \eqref{eq:renormalisation-flow-1}: the Gaussian measure $\gamma_N$ used there is the image of $\tilde{\gamma}_n$ under $d$.

Let $\beta$ be a probe that fits $(N,T^{*})$, let $\hat{\mathfrak{a}}_c$ be the cochains of Section~\ref{sec:probes-discretisation} on the torus $1$, and let $t=(t_c)_c\in\mathbb{C}^{N^{*}}$ carry one variable $t_c$ for each of the $N^{*}$ cells $c$ of the torus $1$. Put
\begin{multline*}
\mathbf{Z}^{\mathrm{ext}}(t):=\frac{1}{(2\pi)^3}\int_{[-\pi,\pi]^3}\mathrm{d}\theta\int_{V_{\mathrm{t}}}\mathrm{d}\tilde{\gamma}_{n_{*}}(A)\,\det\left(D_{A+H_{\theta}+e^{-1}\sum_ct_c\cdot\hat{\mathfrak{a}}_c}+\boldsymbol{\mu}_N\right),\\
\mathring{\Phi}(t):=\frac{\mathbf{Z}^{\mathrm{ext}}(t)}{\mathbf{Z}^{\mathrm{ext}}(0)}.
\end{multline*}
With $A=\eta^{\frac{1}{2}}A^{\mathrm{phys}}$ and $\hat{\mathfrak{a}}_c=\eta^{\frac{1}{2}}a_{\eta,c}$, the phases that $A$ and the external field contribute to the fermion links in this determinant are $\exp\left(i\eta\left(e\cdot A^{\mathrm{phys}}+\sum_ct_c\cdot a_{\eta,c}\right)\right)$. For real numbers $c_2,c_4$ and $m\in\{2,4\}$ let
\[
\Theta^{(0)}_m:=\left|\left(\prod_c\left(\partial_{t_c}^m-c_m\right)\mathring{\Phi}\right)(0)\right|^{\frac{1}{N^{*}}},
\]
the product being taken over the $N^{*}$ cells $c$ of the torus $1$.

Let $\left(m_{0,N},\boldsymbol{\varepsilon}_{0,N}\right)=P_N(e)$ be the pair, and $R^{*}(e)$ the number, of Section~\ref{sec:constants-exponents-scales}. Suppose that $R^{*}=R^{*}(e)$ and that $\boldsymbol{\mu}_N=\bar{m}\cdot\eta+m_{0,N}$. Then $\mathring{\Phi}$ is the function $\mathring{\Phi}$ of Section~\ref{sec:probes-jets-reference} and Section~\ref{sec:averaged-law-partition}: with $a(t):=e^{-1}\sum_ct_c\cdot\hat{\mathfrak{a}}_c$, the function $\mathbf{Z}(t)=\bar{\nu}^{a(t)}(F_N)$ defined there satisfies
\[
\mathbf{Z}(t)=\mathbf{Z}^{\mathrm{ext}}(t)\cdot\frac{\exp\left(-\boldsymbol{\varepsilon}_{0,N}\cdot n_{*}^3\right)}{\det\left(D_0+\bar{m}\cdot\eta\right)}.
\]

Finally, for integers $j\ge1$ let
\[
G_j:=\left(1+\frac{j}{3.2}\right)\cdot\frac{2j^2}{2.56},\qquad K_j:=\left(1+j+j^2\right)\cdot\left(\log j+j\cdot G_j\right).
\]

\subsubsection{Laws on $\mathcal{S}'$}\label{sec:laws-2}

A law on $\mathcal{S}'$ is a probability measure on $\mathcal{S}'$ with its cylinder $\sigma$-algebra. For $f\in\mathcal{S}$, $X(f)$ denotes the random variable $F\mapsto F(f)$ on $\mathcal{S}'$.

\paragraph{Gaussian bound, invariance and reflection positivity.}
A law $\mu$ on $\mathcal{S}'$ satisfies the Gaussian bound~\eqref{eq:setup-G} if
\[
\int\mathrm{d}\mu\,\exp X(f)\le\exp\left(\frac{1}{2}\,\|f\|_2^2\right)\qquad\text{for every real }f\in\mathcal{S},
\]
and it satisfies the hypothesis~\eqref{eq:setup-H} if
\begin{multline*}
\mu\ \text{satisfies \eqref{eq:setup-G} and is invariant under the translations of}\ \mathbb{R}^3\\
\text{and under}\ O_{\mathrm{cub}},\ \text{and}\ \mu\ \text{satisfies \eqref{eq:setup-RPj} for}\ j=1,2,3,
\end{multline*}
the translations and the group $O_{\mathrm{cub}}$ acting through the maps $\tau_b$ and $U_r$ of Section~\ref{sec:lattices-tests-function}. The hypothesis~\eqref{eq:setup-H} is the conjunction of the conditions~\eqref{eq:fourth-cumulants-H1}--\eqref{eq:fourth-cumulants-H3} of Section~\ref{sec:hypotheses-law}; Section~\ref{sec:hypothesis-law} states the same hypothesis in the notation of Section~\ref{sec:ergodic-components-laws}. The group of the eight diagonal matrices with diagonal entries $\pm1$ is denoted $N_3$.

Reflection positivity of a law at a plane $\{x_j=c\}$ is understood in the sense of Section~\ref{sec:properties-laws}, and ergodicity under a group of translations in the sense of Section~\ref{sec:invariant-sets-invariant}. A cutoff is a pair $(N,T)$ of integers $N,T\ge0$ with $n=h\cdot L^{T+N}\ge2$ (Section~\ref{sec:tori-cells-cochains}).

\paragraph{Moments and cumulants.}
For a law $\mu$ on $\mathcal{S}'$, $S_n$ and $\kappa_n$ are its moments and its joint cumulants (Section~\ref{sec:moments-tensor-tests}, Proposition~\ref{prop:sub-gaussian-random}(7)): $S_n(f_1,\dots,f_n)$ is the expectation of $X(f_1)\cdots X(f_n)$ under $\mu$, $\kappa_n(f_1,\dots,f_n)$ is the joint cumulant of $X(f_1),\dots,X(f_n)$, and $\kappa_n(f):=\kappa_n(f,\dots,f)$. The $m$-th cumulant of a single random variable $X$ with respect to an expectation $E$ is denoted $\kappa_m(X;E)$; when $E$ is understood, the shorter notation $\kappa_m(X)$ is used. Written without the index $n$, the symbol $E^{\mathrm{nat}}$, $\mathrm{nat}\in\{\mathrm{af},\mathrm{av}\}$, stands for the expectation $E^{\mathrm{nat}}_n$ of Section~\ref{sec:wilson-determinant-lattice}. When the law with which a moment, a cumulant or the quantity $\mathrm{Ext}_x$ below is formed has to be named, it is written as an upper index, as in $S_2^{\mu}$, $\kappa_4^{\mu}$, $\mathrm{Ext}^{\mu}_x$.

\paragraph{Ergodic decomposition.}
Let $\mu$ be a law on $\mathcal{S}'$ that satisfies~\eqref{eq:setup-G} and is invariant under the translations of $\mathbb{R}^3$. Its ergodic decomposition under $\mathbb{R}^3$ (Theorem~\ref{thm:ergodic-decomposition} with $J=\{1,2,3\}$; Definition~\ref{def:ergodicity-ergodic-decompositions} and Proposition~\ref{prop:existence-conditional-expectations}) is written
\[
\mu=\int\mathrm{d}\nu(\xi)\,\mu_{\xi}.
\]
The fourth cumulant and the covariance of the component $\mu_{\xi}$ are denoted $\kappa_4^{\xi}$ and $C_{\xi}$. For $b\in\mathbb{R}^3$ let $\tau_b^{*}$ be the map $F\mapsto F\circ\tau_b$ of $\mathcal{S}'$ onto itself. The letter $\mathcal{I}$ denotes the $\sigma$-algebra $I'$ of Lemma~\ref{lem:averages-translations-2} formed with the law $\mu$, that is, the $\sigma$-algebra of the cylinder sets $A$ with $\mu\left((\tau_b^{*})^{-1}(A)\,\triangle\,A\right)=0$ for every $b\in\mathbb{R}^3$; it depends on $\mu$. It contains the $\sigma$-algebra of the cylinder sets $A$ with $\tau_b^{*}(A)=A$ for every $b\in\mathbb{R}^3$ (Section~\ref{sec:invariant-sets-invariant}), which does not depend on $\mu$ and is the $\sigma$-algebra denoted $\mathcal{I}$ in Section~\ref{sec:laws}. On $L^2(\mu)$ the conditional expectations given these two $\sigma$-algebras coincide: each of them is the orthogonal projection onto the subspace of the $\Phi\in L^2(\mu)$ with $\Phi\circ\tau_b^{*}=\Phi$ for every $b\in\mathbb{R}^3$ (Lemma~\ref{lem:averages-translations}(iv), Lemma~\ref{lem:averages-translations-2}(b)). For $x\in\mathbb{R}^3$ and $f\in\mathcal{S}$ put
\begin{align*}
\mathrm{Ext}_x(f)&:=\kappa_4(f)-3\,\kappa_4(f,f,\tau_xf,\tau_xf)\\
&\phantom{:}=2\,\kappa_4(f)-\frac{1}{4}\left[\kappa_4(f+\tau_xf)+\kappa_4(f-\tau_xf)\right];
\end{align*}
the two expressions agree by~\eqref{eq:fourth-cumulants-1} of Section~\ref{sec:cumulants-gaussian-laws}, the law being invariant under translations. Formed with $\mu_{\xi}$ in place of $\mu$, this quantity is written $\mathrm{Ext}^{\xi}_x(f)$.

\paragraph{Mixtures.}
Let a property of laws on $\mathcal{S}'$ be given. A law $\mu$ on $\mathcal{S}'$ is called a \emph{mixture} of laws with this property if, for some probability space $(\boldsymbol{\Lambda},m_{\mathrm{x}})$, there is a family $(\boldsymbol{\gamma}_{\mathfrak{l}})_{\mathfrak{l}\in\boldsymbol{\Lambda}}$ of probability measures on $\mathcal{S}'$, each having the property, such that for every cylinder set $A$ the function $\mathfrak{l}\mapsto\boldsymbol{\gamma}_{\mathfrak{l}}(A)$ is measurable and
\[
\mu(A)=\int\mathrm{d}m_{\mathrm{x}}(\mathfrak{l})\,\boldsymbol{\gamma}_{\mathfrak{l}}(A).
\]
In the terms of Definition~\ref{def:ergodicity-ergodic-decompositions}, the family $(\boldsymbol{\gamma}_{\mathfrak{l}})_{\mathfrak{l}\in\boldsymbol{\Lambda}}$ satisfies~\eqref{eq:fourth-cumulants-ED1} on the probability space $(\boldsymbol{\Lambda},m_{\mathrm{x}})$.

\subsubsection{Free fourth-order coefficients}\label{sec:free-fourth-order}

Let $\mathsf{M}>0$. For $z\in\mathbb{R}^3\setminus\{0\}$ put
\[
Y_{\mathsf{M}}(z):=\frac{\exp\left(-\mathsf{M}\cdot|z|_2\right)}{4\pi\,|z|_2},\qquad S_{\mathsf{M}}(z):=\left(\mathsf{M}-\gamma\cdot\partial\right)Y_{\mathsf{M}}(z).
\]
The same letter $S_{\mathsf{M}}$ denotes the operator on $L^2(\mathbb{R}^3;\mathbb{C}^4)$ of convolution with the matrix kernel $S_{\mathsf{M}}(z)$; it is the Fourier multiplier $\left(\mathsf{M}+i\gamma\cdot p\right)^{-1}$, that is, $S_{\mathsf{M}}=\left(\mathsf{M}+\gamma\cdot\partial\right)^{-1}$. For a real $a\in C_c^{\infty}(\mathbb{R}^3;\mathbb{R}^3)$ let $V:=i\gamma\cdot a$, acting by multiplication; it is a bounded operator. Since $S_{\mathsf{M}}VS_{\mathsf{M}}$ is a Hilbert--Schmidt operator (Lemma~\ref{lem:hilbert-schmidt-bound}), so is $\left(S_{\mathsf{M}}VS_{\mathsf{M}}\right)V$, and its square
\[
\left(S_{\mathsf{M}}V\right)^4=\left(S_{\mathsf{M}}VS_{\mathsf{M}}\right)V\left(S_{\mathsf{M}}VS_{\mathsf{M}}\right)V,
\]
a product of two Hilbert--Schmidt operators, is of trace class. Accordingly put
\[
q_{4,\mathsf{M}}(a):=-\frac{1}{4}\operatorname{Tr}_{L^2(\mathbb{R}^3;\mathbb{C}^4)}\left[\left(S_{\mathsf{M}}\,i\gamma\cdot a\right)^4\right].
\]
For a probe $\beta$, with $a=d^{*}\beta$ and $f=da$ as in Section~\ref{sec:probes-discretisation}, put
\[
b(f):=24\,q_{4,\bar{m}}(a)=-6\operatorname{Tr}\left[\left(S_{\bar{m}}\,i\gamma\cdot a\right)^4\right].
\]
The field $a$ is determined by $f$: if $\beta_1$ is a probe with $dd^{*}\beta_1=f$, then $a-d^{*}\beta_1$ is a smooth compactly supported one-form that is closed and coclosed, hence harmonic, hence zero. The traces in these definitions are real numbers (Lemma~\ref{lem:hilbert-schmidt-bound}).

\paragraph{Antiperiodic boundary conditions.}
For $R>0$, $L^2_{\mathrm{ap}}(\mathbb{T}_R;\mathbb{C}^4)$ denotes the space of the functions $\psi\in L^2_{\mathrm{loc}}(\mathbb{R}^3;\mathbb{C}^4)$ that are antiperiodic, $\psi(x+R\cdot e_{\mu})=-\psi(x)$ for $\mu=1,2,3$, with the norm of $L^2$ over one period cube. On this space let $S^R:=\left(\bar{m}+\gamma\cdot\partial\right)^{-1}$; on the momenta $p\in\frac{2\pi}{R}\left(\mathbb{Z}+\frac{1}{2}\right)^3$ it acts as the multiplier $\left(\bar{m}+i\gamma\cdot p\right)^{-1}$. When the support of $a$ is contained in the interior of some period cube, the letter $a$ also stands for the $R$-periodisation of $a$, and one puts
\[
b_R(f):=-6\operatorname{Tr}_{L^2_{\mathrm{ap}}(\mathbb{T}_R;\mathbb{C}^4)}\left[\left(S^R\,i\gamma\cdot a\right)^4\right],
\]
the operator in the brackets being of trace class. Indeed, let $V:=i\gamma\cdot a$ act by multiplication on $L^2_{\mathrm{ap}}(\mathbb{T}_R;\mathbb{C}^4)$; it is a bounded operator, and $S^RVS^R$ is a Hilbert--Schmidt operator for every $R>0$ (Lemma~\ref{lem:hilbert-schmidt-bound}). Hence so is $\left(S^RVS^R\right)V$, and, as for $q_{4,\mathsf{M}}(a)$, its square
\[
\left(S^RV\right)^4=\left(S^RVS^R\right)V\left(S^RVS^R\right)V,
\]
a product of two Hilbert--Schmidt operators, is of trace class.

\paragraph{Lattice determinants.}
Let $\eta\in(0,1]$ be such that $R/\eta$ is an integer $\ge2$, and let $a_{\eta}$ be a one-cochain on $\Lambda_{\eta,R}$, read in lattice units as a one-cochain on $\Lambda_n$, $n=R/\eta$. With the operator $D[\cdot]$ of~\eqref{eq:free-coefficient-5} put
\[
N_0(t):=\frac{\det\left(D[t\eta a_{\eta}]+\bar{m}\cdot\eta\right)}{\det\left(D[0]+\bar{m}\cdot\eta\right)},\qquad b_{\eta,R}(0):=\left.\frac{\mathrm{d}^4}{\mathrm{d}t^4}\log N_0(t)\right|_{t=0}.
\]
This is the function $N_0$ of Section~\ref{sec:averaged-law-partition}, where $t\eta a_{\eta}$ is written $q\,(t/e)\cdot\hat{\mathfrak{a}}$.

\subsection{Free coefficients, cumulant bounds and bare mass shifts}\label{sec:free-coefficients-cumulant}

The three theorems of this subsection contain, in this order, Estimate~6, Estimate~3, and Estimates~7 and~8 in full. Part~I of Theorem~\ref{thm:main-theorem} is proved in Section~\ref{sec:admissibility-limit-points}, and Part~II in Section~\ref{sec:proof-part-ii-2}, where these three theorems are combined with Part~I and with results of the earlier sections.

\subsubsection{Probes, free coefficients and lattice symmetries}\label{sec:probes-free-coefficients}

\begin{theorem}\label{thm:free-coefficient-1}
\begin{enumerate}
\item[(a)] Let $a\in C_c^{\infty}(\mathbb{R}^3;\mathbb{R}^3)$ be real and $B:=\operatorname{curl}a$. There is a constant $C_a<\infty$ such that, for every $\mathsf{M}\ge1$,
\[
\left|q_{4,\mathsf{M}}(a)-\frac{1}{240\pi\cdot\mathsf{M}^5}\int_{\mathbb{R}^3}\mathrm{d}x\,|B(x)|_2^4\right|\le C_a\cdot\mathsf{M}^{-7}.
\]
Let $\rho(x):=\exp\left(-\frac{1}{1-|x|_2^2}\right)$ if $|x|_2<1$ and $\rho(x):=0$ otherwise, and let $a:=\left(\partial_2\rho,-\partial_1\rho,0\right)$; from now on in (a), $a$ is this field and $C_a$ is the constant that the preceding assertion attaches to it. Put
\[
I_{\beta}:=\int_{\mathbb{R}^3}\mathrm{d}x\,|\operatorname{curl}a(x)|_2^4.
\]
For $l\ge1$ define the two-form $\beta^{(l)}$ by $\beta^{(l)}_{12}(x):=l^2\rho(x/l)$ and $\beta^{(l)}_{13}:=\beta^{(l)}_{23}:=0$, the field $a_l(x):=l\,a(x/l)$, and $f_l:=da_l$. Then $I_{\beta}>0$, $d^{*}\beta^{(l)}=a_l$ and $f_l=dd^{*}\beta^{(l)}$; and if $\bar{m}\cdot l\ge1$, then
\[
\left|b(f_l)-\frac{l^3\cdot I_{\beta}}{10\pi\cdot\bar{m}^5}\right|\le24\,C_a\,l\cdot\bar{m}^{-7}.
\]
In particular, $b(f_l)>0$ whenever $\bar{m}\cdot l\ge1$ and $\left(\bar{m}\cdot l\right)^2>240\pi\,C_a/I_{\beta}$.
\item[(b)] Let $0<\eta\le1$ and let $R>0$ be such that $R/\eta$ is an integer. Let $\beta$ be a probe whose cube $Q_{\beta}$ lies in the interior of a period cube of the torus of side $R$, and let $a$, $f$, $a'$, $\check{a}_{\eta}$, $f^{(\eta)}$, $u$ and $a_{\eta}$ be formed from $\beta$ as in Section~\ref{sec:probes-discretisation}.
\begin{enumerate}
\item[(b1)] For $x\in\Lambda_{\eta,R}$ and $\nu=1,2,3$ one has $a_{\eta,\nu}(x)=\check{a}_{\eta,\nu}\left(x+\frac{\eta}{2}e_{\nu}\right)$. Moreover $\left(f^{(\eta)}\right)_{\eta}=d_{\eta}a_{\eta}$, and $a_{\eta}$ lies in the range of $d_{\eta}^{*}$, that is, $a_{\eta}$ is orthogonal to $\ker d_{\eta}$, the space spanned by the exact and the constant one-cochains.
\item[(b2)] The functions $\check{a}_{\eta}$ and $f^{(\eta)}$ are smooth, with supports contained in $\operatorname{supp}\beta+[-\eta,\eta]^3$. For every bond in the support of $a_{\eta}$, both end points lie in $\operatorname{supp}\beta+\left[-\frac{3\eta}{2},\frac{3\eta}{2}\right]^3$, and this set is contained in $Q_{\beta}$. Let moreover $N\ge0$ and $T^{*}\ge1$ be integers with $\eta=L^{-N}$, let $\beta$ fit $(N,T^{*})$, let $v$ be the placement of Section~\ref{sec:probes-discretisation}, and let $\hat{\mathfrak{a}}$, $\boldsymbol{\sigma}$ and $X$ be formed with $\tau_v\beta$ on $\Lambda_{\eta,R}$ as in that section. If $R=k\cdot R^{*}$ for an integer $k\ge1$, where $R^{*}=h\cdot L^{T^{*}}$, that is, if $\Lambda_{\eta,R}$ is the torus $k$ of Section~\ref{sec:reference-tori-cells} (for such $R$ the hypothesis on $Q_{\beta}$ in (b) is satisfied, since $D_{\beta}\le l_c-3\eta<R^{*}\le R$), then every bond in the support of $\hat{\mathfrak{a}}$ has both end points in $\{1,\dots,m_{\mathrm{p}}-1\}^3$ in the lattice coordinates of this torus; the support of $d\hat{\mathfrak{a}}$ is contained in $P^o(\mathfrak{c}_0)$; for distinct cells $c$ of the torus $k$ the cochains $d\hat{\mathfrak{a}}_c$ have disjoint supports; and $\|d\hat{\mathfrak{a}}_c\|^2=\boldsymbol{\sigma}^2$ for every cell $c$ of the torus $k$. For every $R$ as in (b), also when $R$ is not of this form, one has $X=X_{\eta,R}\left(\tau_vf^{(\eta)}\right)$.
\item[(b3)] For every real $s\ge0$,
\[
\left\|f^{(\eta)}-f\right\|_{H^s}\le\eta^2\,\|\beta\|_{H^{s+4}}.
\]
\item[(b4)] For every integer $i\ge0$,
\[
\left|D^i\check{a}_{\eta}\right|_{\infty}\le2\,|\beta|_{C^{i+1}},\qquad\left|\check{a}_{\eta}-a\right|_{C^i}\le3\eta^2\,|\beta|_{C^{i+3}},
\]
where, for a function or a vector field $g$ on $\mathbb{R}^3$ with components of class $C^i$, $|g|_{C^i}$ denotes the largest of the sup norms of the partial derivatives of order at most $i$ of the components of $g$, as for two-forms in Section~\ref{sec:lattices-tests-function}, and $\left|D^ig\right|_{\infty}$ denotes the largest of the sup norms of the partial derivatives of order $i$ of the components of $g$.
\item[(b5)] Let $N\ge0$ and $T^{*}\ge1$ be integers with $\eta=L^{-N}$, let $\beta$ fit $(N,T^{*})$, let $v$ be the placement of Section~\ref{sec:probes-discretisation}, and let $\boldsymbol{\sigma}$ be formed with $\tau_v\beta$ on $\Lambda_{\eta,R}$ as in that section, $R$ being any number as in (b), not necessarily a multiple of $R^{*}=h\cdot L^{T^{*}}$. Then
\begin{gather*}
\boldsymbol{\sigma}\le\left\|f^{(\eta)}\right\|_2\le\|f\|_2+\|\beta\|_{H^4}=:s_f,\\
\left|\boldsymbol{\sigma}-\|f\|_2\right|\le\eta^2\,\|\beta\|_{H^4}+21\,\eta\,|\beta|_{C^3}\,D_{\beta}^{\frac{3}{2}}.
\end{gather*}
\end{enumerate}
\item[(c)] Keep $\beta$, $a$, $f$ and $a_{\eta}$ as in (b), and assume $R\ge D_{\beta}+1$. Then there exist constants $C_1,C_2<\infty$, which depend on $\bar{m}$ and $\beta$ only, with the following properties.
\begin{enumerate}
\item[(c1)] With $V=i\gamma\cdot a$, the operators $S_{\bar{m}}VS_{\bar{m}}$ on $L^2(\mathbb{R}^3;\mathbb{C}^4)$ and $S^RVS^R$ on $L^2_{\mathrm{ap}}(\mathbb{T}_R;\mathbb{C}^4)$ are Hilbert--Schmidt operators. Hence $b(f)$ and $b_R(f)$ are defined.
\item[(c2)] For every $\eta\in(0,1]$ such that $R/\eta$ is an integer,
\[
\left|b_{\eta,R}(0)-b_R(f)\right|\le C_1\,\eta^{\frac{1}{2}}.
\]
Whenever $R\cdot L^N$ is an integer for all sufficiently large $N$ (for instance when $R=h\cdot L^T$), this bound applies to $\eta=L^{-N}$ for these $N$, and therefore $b_R(f)=\lim_{N\to\infty}b_{L^{-N},R}(0)$.
\item[(c3)] One has
\[
\left|b_R(f)-b(f)\right|\le C_2\exp\left(-\bar{m}\cdot\left(R-D_{\beta}\right)\right).
\]
\end{enumerate}
Since $C_1$ and $C_2$ depend neither on $\eta$ nor on $R$, (c2) and (c3) together show that $b_{\eta,R}(0)$ tends to $b(f)$ when $\eta\to0$ and $R\to\infty$ jointly, no relation between $\eta$ and $R$ being imposed.
\item[(d)] Let $B$ be a complex one-cochain on $\Lambda_n$, and let $\boldsymbol{\mu}$ be real.
\begin{enumerate}
\item[(d1)] For every complex zero-cochain $\boldsymbol{\lambda}$ and every integer one-cochain $\mathfrak{z}$,
\[
\mathrm{Dt}\left(B+d\boldsymbol{\lambda}\right)=\mathrm{Dt}(B),\qquad\mathrm{Dt}\left(B+2\pi\cdot\mathfrak{z}\right)=\mathrm{Dt}(B).
\]
The function $\mathrm{Dt}$ is entire, and for every $c\ge0$ it is bounded on the set of the complex one-cochains $B$ with $\left|\operatorname{Im}B\right|_{\infty}\le c$.
\item[(d2)] Let $B$ be real. The spinors, and the functions $\chi\colon\mathbb{Z}^3\to\mathbb{C}^2$ with $\chi(x+n_{\mu}e_{\mu})=-\chi(x)$ for $\mu=1,2,3$, are given the inner product $\langle\psi,\varphi\rangle:=\sum_x\psi(x)^{*}\varphi(x)$, where $x$ ranges over a set of representatives in $\mathbb{Z}^3$ of the sites of $\Lambda_n$ (the summand is periodic, so the choice of the representatives is immaterial); adjoints, and Hermitian parts $\operatorname{Re}O:=\frac{1}{2}\left(O+O^{*}\right)$ of operators $O$, are taken with respect to this inner product. In the decomposition $\mathbb{C}^4=\mathbb{C}^2\oplus\mathbb{C}^2$ in which $\gamma_{\mu}=\operatorname{diag}(\sigma_{\mu},-\sigma_{\mu})$ (Section~\ref{sec:letters-norms-dirac}) the first two and the last two components of a spinor are such functions $\chi$, and there is an operator $\mathbf{B}$ on the space of these functions with $D[B]+\boldsymbol{\mu}=\operatorname{diag}\left(\mathbf{B},\mathbf{B}^{*}\right)$; hence $\mathrm{Dt}(B)=\left|\det\mathbf{B}\right|^2\ge0$. Moreover $\operatorname{Re}\left(D[B]+\boldsymbol{\mu}\right)\ge\boldsymbol{\mu}$.
\item[(d3)] $\mathrm{Dt}(-B)=\mathrm{Dt}(B)$.
\item[(d4)] Let $g$ be the map $x\mapsto rx+c$ of $\mathbb{Z}^3$, where $r\in O_{\mathrm{cub}}$ and $c\in\mathbb{Z}^3$, and assume that $r\left(\prod_{\mu}n_{\mu}\mathbb{Z}\right)=\prod_{\mu}n'_{\mu}\mathbb{Z}$. Then the map induced by $g$ on the quotients is an isometry of $\Lambda_n$ onto $\Lambda_{n'}$, and
\[
\mathrm{Dt}_{n'}\left(g_{*}B\right)=\mathrm{Dt}_n(B),\qquad\text{where }\left(g_{*}B\right)(gx,gx'):=B(x,x').
\]
Here $\mathrm{Dt}_n$ is the function $\mathrm{Dt}$ of Section~\ref{sec:wilson-determinant-lattice} formed on $\Lambda_n$, that is, with the operator~\eqref{eq:free-coefficient-5} on the spinors of $\Lambda_n$, and $\mathrm{Dt}_{n'}$ is formed in the same way on $\Lambda_{n'}$, with the same number $\boldsymbol{\mu}$.
\item[(d5)] Let $g$ be as in (d4), let $g^{*}$ and $g_{*}$ act on cochains as in Section~\ref{sec:reference-tori-cells}, and let $\Phi\in\mathfrak{F}(n)$. Then
\[
\boldsymbol{\omega}_{n'}\left(\Phi\circ g^{*}\right)=\boldsymbol{\omega}_n(\Phi),\qquad\boldsymbol{\omega}_{n'}^{(0)}\left(\Phi\circ g^{*}\right)=\boldsymbol{\omega}_n^{(0)}(\Phi),
\]
and $\boldsymbol{\omega}_n$, $\boldsymbol{\omega}_n^{(0)}$ are invariant under $F^u\mapsto-F^u$, that is, their values do not change if $\Phi$ is replaced by the function $F^u\mapsto\Phi(-F^u)$. In particular, when $\Lambda_n$ is a cubic torus and $\boldsymbol{\mu}>-1$, so that $Z_n>0$ and $\mathbf{Z}_n>0$ (Section~\ref{sec:wilson-determinant-lattice}), the laws $E^{\mathrm{af}}_n$ and $E^{\mathrm{av}}_n$ are both preserved by the lattice translations, by the elements of $O_{\mathrm{cub}}$ and by the sign change $F^u\mapsto-F^u$.

\emph{Sectors of fixed flux.} As $r$ is a signed permutation matrix, $re_{\mu}=\epsilon_{\mu}\cdot e_{\mu'}$ with signs $\epsilon_{\mu}\in\{1,-1\}$ and a permutation $\mu\mapsto\mu'$ of the directions, and the hypothesis of (d4) gives $n'_{\mu'}=n_{\mu}$. For $w=(w_{\mu\nu})_{\mu<\nu}\in\mathbb{Z}^3$ set $w_{\nu\mu}:=-w_{\mu\nu}$, and let $g\cdot w\in\mathbb{Z}^3$, extended to all pairs of distinct directions by the same rule, be defined by
\[
(g\cdot w)_{\mu'\nu'}:=\epsilon_{\mu}\cdot\epsilon_{\nu}\cdot w_{\mu\nu}.
\]
The map $w\mapsto g\cdot w$ is a bijection of $\mathbb{Z}^3$ with $g\cdot0=0$; moreover $\mathfrak{q}_{g\cdot w}(n')=\mathfrak{q}_w(n)$, and
\[
\boldsymbol{\omega}_{n'}^{(g\cdot w)}\left(\Phi\circ g^{*}\right)=\boldsymbol{\omega}_n^{(w)}(\Phi)\qquad\text{for every }w\in\mathbb{Z}^3\text{ and every }\Phi\in\mathfrak{F}(n).
\]

Finally, $b_{\eta,R}(0)$ is unchanged when $a_{\eta}$ is replaced by $g_{*}a_{\eta}$.
\end{enumerate}
\end{enumerate}
\end{theorem}

\subsubsection{Cumulant bounds for the probe field}\label{sec:cumulant-bounds-probe}

\begin{theorem}\label{thm:free-coefficient-2}
Let $N\ge0$ and $T^{*}\ge1$ be integers, let the lattice mass satisfy $\boldsymbol{\mu}_N>-\frac{1}{2}$, and let $0<e\le\frac{1}{3}$ be such that
\begin{equation}\label{eq:free-coefficient-6}
6N^{*}\cdot e^2\cdot R^{*}\cdot\log\left(e^{-1}\right)\le\pi^2.
\end{equation}
Let $\beta$ be a probe that fits $(N,T^{*})$, and let $\hat{\mathfrak{a}}$, $\boldsymbol{\sigma}$, $X$, $X_c$, $H_2$ and $H_4$ be formed with the placement of Section~\ref{sec:probes-discretisation}. Let $c_2,c_4\in\mathbb{R}$, and for $m\in\{2,4\}$ put
\begin{align*}
\dot{M}_4&:=600\,(N^{*})^2\cdot\boldsymbol{\sigma}^4\cdot e^{-4}+|c_4|,\qquad\dot{M}_2:=12N^{*}\cdot\boldsymbol{\sigma}^2\cdot e^{-2}+|c_2|,\\
Y&:=\log^{+}\left((N^{*})^2\cdot\left(600\,\boldsymbol{\sigma}^4+12\,\boldsymbol{\sigma}^2+|c_4|+|c_2|\right)\right),\qquad\check{X}:=1+\frac{Y+1}{6\log\left(e^{-1}\right)},\\
\boldsymbol{\Gamma}^{\mathrm{af}}&:=\exp(\delta),\qquad\boldsymbol{\Theta}^{\mathrm{af}}_m:=\Theta^{(0)}_m+2e^6\cdot\dot{M}_m,\\
\boldsymbol{\Gamma}^{\mathrm{av}}&:=\exp\left(\delta+\frac{3\log\left(4+16\,\check{X}\right)+6}{N^{*}}\right),\qquad\boldsymbol{\Theta}^{\mathrm{av}}_m:=\max\left(\Theta^{(0)}_m,e^2\right);
\end{align*}
here $\delta$ and $\Theta^{(0)}_m$ are the numbers of Section~\ref{sec:numbers-attached-reference}, formed with these data. Let $\mathrm{nat}\in\{\mathrm{af},\mathrm{av}\}$, let $k\in\{1\}\cup\{L,L+1,\dots\}$ be an integer, and let $\kappa_2$ and $\kappa_4$ be the cumulants of order $2$ and $4$ of $X=\langle d\hat{\mathfrak{a}},F^u\rangle$ under $E^{\mathrm{nat}}_k$. Then
\begin{align}
\left|e^{-2}\cdot\left(\kappa_2-\boldsymbol{\sigma}^2\right)-c_2\right|&\le\boldsymbol{\Gamma}^{\mathrm{nat}}\cdot\boldsymbol{\Theta}^{\mathrm{nat}}_2,\label{eq:free-coefficient-7}\\
\left|e^{-4}\cdot\kappa_4-\left(c_4-3c_2^2\right)\right|&\le\boldsymbol{\Gamma}^{\mathrm{nat}}\cdot\boldsymbol{\Theta}^{\mathrm{nat}}_4+3\boldsymbol{\Gamma}^{\mathrm{nat}}\cdot\boldsymbol{\Theta}^{\mathrm{nat}}_2\cdot\left(2|c_2|+\boldsymbol{\Gamma}^{\mathrm{nat}}\cdot\boldsymbol{\Theta}^{\mathrm{nat}}_2\right),\label{eq:free-coefficient-8}\\
0\le\delta&\le\frac{a_L}{N^{*}}\cdot\log\frac{2\pi}{e^2\cdot R^{*}}+\frac{K_L+D^{\mathrm{mat}}+1}{(L-1)^3\cdot N^{*}},\label{eq:free-coefficient-9}
\end{align}
and, for $m\in\{2,4\}$ and every set $S$ of cells of the torus $k$ of Section~\ref{sec:reference-tori-cells},
\begin{equation}\label{eq:free-coefficient-10}
\left|E^{\mathrm{nat}}_k\prod_{c\in S}\left(e^{-m}\cdot H_m(X_c)-c_m\right)\right|\le\left(\boldsymbol{\Gamma}^{\mathrm{nat}}\cdot\boldsymbol{\Theta}^{\mathrm{nat}}_m\right)^{\#|S|}.
\end{equation}
\end{theorem}

\subsubsection{The bare mass shifts and the reference numbers}\label{sec:bare-mass-shifts}

\begin{theorem}\label{thm:free-coefficient-3}
Let $0<\bar{m}\le1$. For $e>0$ with $e\cdot h'^{\frac{1}{2}}\le e_1$, where $e_1$ is the number of Section~\ref{sec:constants-exponents-scales}, define
\begin{align*}
T^{*}(e)&:=\log_L\left(\frac{R^{*}(e)}{h}\right)=m-m_0-1+\log_Lr_K(e),\qquad N_0(e):=N_{\mathrm{min}}\left(R^{*}(e),e\right),\\
m_{0,N}(e)&:=\begin{cases}\text{the first entry of the pair }P_N(e)&\text{if }N\ge N_0(e),\\ 0&\text{if }0\le N<N_0(e),\end{cases}\\
C_{\mathrm{mat}}&:=\left(L^3+1\right)\left(\varepsilon_5\left(L^{-m_0-1}\right)+1\right).
\end{align*}
Here the constants $m$ (with $M=L^m$) and $m_0$, the exponents $\mathfrak{e}$, $\mathfrak{e}_2$, $\sigma$, the scales $h$, $h'$, $r_K(e)$, $R^{*}(e)$, the number $N_{\mathrm{min}}(R,e)$, the pair $P_N(e)$, the function $\varepsilon_5$ and the number $e_G$ are those of Section~\ref{sec:constants-exponents-scales}. Then the following hold.
\begin{enumerate}
\item\label{it:free-coefficient-3-1} Let $e>0$ satisfy $e\cdot h'^{\frac{1}{2}}\le e_1$. Then $e\le\frac{1}{3}$; the number $T^{*}(e)$ is an integer with $T^{*}(e)\ge1$, and $R^{*}(e)=h\cdot L^{T^{*}(e)}$. For every integer $N\ge0$ the number $m_{0,N}(e)$ is real and $\left|m_{0,N}(e)\right|<\frac{1}{8}$; hence $N\mapsto m_{0,N}(e)$ is a mass assignment in the sense of Section~\ref{sec:wilson-determinant-lattice}. The numbers $m_{0,N}(e)$ depend on $N$, $e$, $\bar{m}$ and on the scheme of Theorem~\ref{thm:agreement-trajectories}, and not on a torus. As $e\to0$, the length $l_c=R^{*}(e)/(2L)$ of Section~\ref{sec:reference-tori-cells}, taken with $T^{*}=T^{*}(e)$, tends to infinity. For every integer $N\ge N_0(e)$, the number $D^{\mathrm{mat}}$ of Section~\ref{sec:numbers-attached-reference}, formed with $T^{*}=T^{*}(e)$, the cutoff $N$, the charge $e$ and the lattice mass $\boldsymbol{\mu}_N=\bar{m}\cdot\eta+m_{0,N}(e)$, satisfies $D^{\mathrm{mat}}\le C_{\mathrm{mat}}$. If moreover $e\le e_G$, then the inequality~\eqref{eq:free-coefficient-6} holds with $R^{*}=R^{*}(e)$.
\item\label{it:free-coefficient-3-2} One has $0<\mathfrak{e}\le\frac{1}{4}$, $0\le\mathfrak{e}_2<1$ and $\sigma>0$. For every probe $\beta$ in the sense of Section~\ref{sec:probes-discretisation} there are numbers $e_1(\beta)\in\left(0,e_G\right]$ and $C^G_{\beta}<\infty$, depending on the constants of Section~\ref{sec:constants-exponents-scales}, on $\bar{m}$ and on $\beta$, with the following property. Let $0<e\le e_1(\beta)$ (so that $e\cdot h'^{\frac{1}{2}}\le e_1$, by Section~\ref{sec:constants-exponents-scales}) and let $N\ge N_0(e)$ be an integer; the reference tori of Section~\ref{sec:reference-tori-cells} are taken with $T^{*}=T^{*}(e)$, and the lattice mass is $\boldsymbol{\mu}_N=\bar{m}\cdot\eta+m_{0,N}(e)$. Then $\beta$ fits $\left(N,T^{*}(e)\right)$ in the sense of Section~\ref{sec:probes-discretisation}. Let $v\in\eta\mathbb{Z}^3$ be the placement of Section~\ref{sec:probes-discretisation}, let $t\mapsto N_0(t)$ be the free ratio of Section~\ref{sec:free-fourth-order} on the torus $\Lambda_{\eta,R^{*}(e)}$, formed with the one-cochain $a_{\eta}$ of Section~\ref{sec:probes-discretisation} built from $\tau_v\beta$ in place of $\beta$, and put
\[
c_m=c_m(N,e,\beta):=N_0^{(m)}(0)\qquad(m=2,4),
\]
the $m$-th derivative of $t\mapsto N_0(t)$ at $t=0$. Then $c_2$ and $c_4$ are real, and, with the numbers $\Theta^{(0)}_2$, $\Theta^{(0)}_4$ of Section~\ref{sec:numbers-attached-reference} formed with these $c_2$, $c_4$, the following hold.
\begin{enumerate}
\item One has
\[
\Theta^{(0)}_4\le C^G_{\beta}\cdot e^{\mathfrak{e}},\qquad\left(1+|c_2|\right)\cdot\Theta^{(0)}_2\le C^G_{\beta}\cdot e^{\mathfrak{e}}.
\]
\item One has
\[
|c_2|+|c_4|^{\frac{1}{2}}\le C^G_{\beta}\cdot e^{-\mathfrak{e}_2}.
\]
\item The number $c_4-3c_2^2$ is the number $b_{\eta,R^{*}(e)}(0)$ of Section~\ref{sec:free-fourth-order}, formed with the same one-cochain $a_{\eta}$, and
\[
\left|c_4-3c_2^2-b\left(d\,d^{*}\beta\right)\right|\le C^G_{\beta}\cdot\left(e+\eta^{\frac{1}{2}}\right).
\]
In particular, for every $e\in\left(0,e_1(\beta)\right]$,
\[
\limsup_{N\to\infty}\left|c_4(N,e,\beta)-3c_2(N,e,\beta)^2-b\left(d\,d^{*}\beta\right)\right|\le C^G_{\beta}\cdot e^{\sigma}.
\]
\end{enumerate}
\end{enumerate}
\end{theorem}

\subsection{Discretised probes, lattice symmetries and factorisation over slices}\label{sec:discretised-probes-lattice}

\subsubsection{The discretised probe}\label{sec:discretised-probe}

\begin{proof}[Proof of Theorem~\ref{thm:free-coefficient-1}(b)]
Let $0<\eta\le1$, let $R/\eta$ be an integer, let $\beta$ be a probe whose cube $Q_{\beta}$ lies in the interior of a period cube of the torus $\mathbb{T}_R=\mathbb{R}^3/R\mathbb{Z}^3$, and let $a$, $f$, $a'$, $\check{a}_{\eta}$, $f^{(\eta)}$, $u$ and $a_{\eta}$ be formed from $\beta$ as in Section~\ref{sec:probes-discretisation}. The cube $Q_{\beta}$ is closed and has side $D_{\beta}=2\rho_{\beta}\ge4$ (Section~\ref{sec:probes-discretisation}); as it lies in the interior of a period cube, a translate of $[0,R)^3$, one has $R>D_{\beta}\ge4$, and $R/\eta\ge R>4$ because $\eta\le1$. Hence $R\ge1$ and $R/\eta$ is an integer $\ge2$, so that the last paragraph of Section~\ref{sec:test-forms-test} applies to the torus $\Lambda_{\eta,R}$. A function on $\mathbb{R}^3$ is evaluated at a point of $\mathbb{T}_R$ through its periodisation, as in the definition of $u$ and in \eqref{eq:continuum-infinite-2}. The operator $\delta_{\eta,j}$ is $\eta^{-1}$ times the difference of two translations, and $J_{\eta}$ is a mean of translations; hence these operators commute with one another, with the translations $\tau_b$, $b\in\mathbb{R}^3$, with periodisation and, on smooth functions, with derivatives.

\emph{Part (b1).} For $A\in C^1(\Lambda_{\eta,R})$ let $A_{\nu}(y)$ be the value of $A$ on the bond $(y;\{\nu\})$, and let $(d_{\eta}A)_{\mu\nu}(x)$ be the value of $d_{\eta}A$ on the plaquette $(x;\mu\nu)$. By Section~\ref{sec:tori-cells-cochains},
\[
(d_{\eta}A)_{\mu\nu}(x)=\eta^{-1}\cdot\left[A_{\mu}(x)+A_{\nu}(x+\eta e_{\mu})-A_{\mu}(x+\eta e_{\nu})-A_{\nu}(x)\right]\qquad(\mu<\nu).
\]
For one-cochains $A$ and $w$ one has $\langle A,w\rangle_{\eta}=\eta^3\sum_{z\in\Lambda_{\eta,R}}\sum_{\kappa=1}^3A_{\kappa}(z)\cdot w_{\kappa}(z)$. Since $\langle A,d_{\eta}^{*}u\rangle_{\eta}=\langle d_{\eta}A,u\rangle_{\eta}$ for every $A$, the number $(d_{\eta}^{*}u)_{\nu}(y)$ is $\eta^{-3}$ times the coefficient of $A_{\nu}(y)$ in
\[
\langle d_{\eta}A,u\rangle_{\eta}=\eta^3\sum_{x\in\Lambda_{\eta,R}}\sum_{\kappa<\mathfrak{l}}(d_{\eta}A)_{\kappa\mathfrak{l}}(x)\cdot u_{\kappa\mathfrak{l}}(x).
\]
In this sum $A_{\nu}(y)$ occurs in the terms with $\mathfrak{l}=\nu$ and $\kappa=\mu<\nu$, at $x=y-\eta e_{\mu}$ and at $x=y$, and in the terms with $\kappa=\nu<\mathfrak{l}$, at $x=y$ and at $x=y-\eta e_{\mathfrak{l}}$. Hence
\[
(d_{\eta}^{*}u)_{\nu}(y)=\sum_{\mu<\nu}\eta^{-1}\cdot\left[u_{\mu\nu}(y-\eta e_{\mu})-u_{\mu\nu}(y)\right]+\sum_{\mathfrak{l}>\nu}\eta^{-1}\cdot\left[u_{\nu\mathfrak{l}}(y)-u_{\nu\mathfrak{l}}(y-\eta e_{\mathfrak{l}})\right].
\]
For $\mu>\nu$ put $u_{\mu\nu}:=-u_{\nu\mu}$; as $\beta_{\mu\nu}=-\beta_{\nu\mu}$, this is again the formula of Section~\ref{sec:probes-discretisation} for $u_{\mu\nu}$. Then
\[
(d_{\eta}^{*}u)_{\nu}(y)=-\sum_{\mu\ne\nu}\eta^{-1}\cdot\left[u_{\mu\nu}(y)-u_{\mu\nu}(y-\eta e_{\mu})\right].
\]
Let $b=(y;\{\nu\})$, with centre $c(b)=y+\frac{\eta}{2}e_{\nu}$. For $\mu\ne\nu$ the plaquettes in the directions $\mu,\nu$ with the base points $y$ and $y-\eta e_{\mu}$ have the centres $c(b)+\frac{\eta}{2}e_{\mu}$ and $c(b)-\frac{\eta}{2}e_{\mu}$. By the definition of $u$, therefore,
\begin{align*}
u_{\mu\nu}(y)-u_{\mu\nu}(y-\eta e_{\mu})&=\left(J_{\eta}\beta_{\mu\nu}\right)\left(c(b)+\frac{\eta}{2}e_{\mu}\right)-\left(J_{\eta}\beta_{\mu\nu}\right)\left(c(b)-\frac{\eta}{2}e_{\mu}\right)\\
&=\eta\cdot\left(\delta_{\eta,\mu}J_{\eta}\beta_{\mu\nu}\right)\left(c(b)\right).
\end{align*}
As $J_{\eta}$ commutes with $\delta_{\eta,\mu}$ and $\beta_{\nu\nu}=0$, it follows that
\[
a_{\eta,\nu}(y)=(d_{\eta}^{*}u)_{\nu}(y)=-\sum_{\mu\ne\nu}\left(J_{\eta}\delta_{\eta,\mu}\beta_{\mu\nu}\right)\left(c(b)\right)=\left(J_{\eta}a'_{\nu}\right)\left(c(b)\right)=\check{a}_{\eta,\nu}\left(y+\frac{\eta}{2}e_{\nu}\right),
\]
which is the first identity of (b1).

Now let $p=(x;\mu\nu)$, $\mu<\nu$, be a plaquette, with centre $c(p)=x+\frac{\eta}{2}\left(e_{\mu}+e_{\nu}\right)$. The first identity, applied to the four bonds in the boundary of $p$, gives
\[
a_{\eta,\nu}(x+\eta e_{\mu})=\check{a}_{\eta,\nu}\left(c(p)+\frac{\eta}{2}e_{\mu}\right),\qquad a_{\eta,\nu}(x)=\check{a}_{\eta,\nu}\left(c(p)-\frac{\eta}{2}e_{\mu}\right),
\]
together with the two identities obtained by exchanging $\mu$ and $\nu$. Hence
\begin{align*}
(d_{\eta}a_{\eta})_{\mu\nu}(x)&=\eta^{-1}\cdot\left[\check{a}_{\eta,\nu}\left(c(p)+\frac{\eta}{2}e_{\mu}\right)-\check{a}_{\eta,\nu}\left(c(p)-\frac{\eta}{2}e_{\mu}\right)\right]\\
&\quad-\eta^{-1}\cdot\left[\check{a}_{\eta,\mu}\left(c(p)+\frac{\eta}{2}e_{\nu}\right)-\check{a}_{\eta,\mu}\left(c(p)-\frac{\eta}{2}e_{\nu}\right)\right]\\
&=\left(\delta_{\eta,\mu}\check{a}_{\eta,\nu}-\delta_{\eta,\nu}\check{a}_{\eta,\mu}\right)\left(c(p)\right)=\left(J_{\eta}f^{(\eta)}_{\mu\nu}\right)\left(c(p)\right),
\end{align*}
the last equality because $\check{a}_{\eta}=J_{\eta}a'$, $f^{(\eta)}_{\mu\nu}=\delta_{\eta,\mu}a'_{\nu}-\delta_{\eta,\nu}a'_{\mu}$, and $J_{\eta}$ commutes with the centred differences. As $Q(p)=c(p)+\left[-\frac{\eta}{2},\frac{\eta}{2}\right)^3$ (Section~\ref{sec:tori-cells-cochains}), the right side is the mean value of $f^{(\eta)}_{\mu\nu}$ over the averaging cube $Q(p)$, that is, $\left(f^{(\eta)}\right)_{\eta}(p)$ by \eqref{eq:continuum-infinite-2}. Thus $\left(f^{(\eta)}\right)_{\eta}=d_{\eta}a_{\eta}$. On the period cube that contains $Q_{\beta}$ the periodisations of $\beta$, $\check{a}_{\eta}$ and $f^{(\eta)}$ coincide with these functions, which are supported in $\operatorname{supp}\beta+[-\eta,\eta]^3\subset Q_{\beta}$ by the first step of Part (b2); so in the identities above the periodisations may be replaced by the functions themselves.

Finally, $a_{\eta}=d_{\eta}^{*}u$ lies in the range of $d_{\eta}^{*}$. For the linear map $d_{\eta}\colon C^1(\Lambda_{\eta,R})\to C^2(\Lambda_{\eta,R})$ between finite-dimensional spaces with the scalar products $\langle\cdot,\cdot\rangle_{\eta}$ one has $\operatorname{ran} d_{\eta}^{*}=\left(\ker d_{\eta}\right)^{\perp}$, and $\ker d_{\eta}=\ker\left(d|_{C^1}\right)=V_{\mathrm{g}}\oplus V_{\mathrm{h}}$ is the space of the sums of an exact and a constant one-cochain, by the last paragraph of Section~\ref{sec:test-forms-test}, which applies here because $R\ge1$ and $R/\eta\ge2$. This proves (b1).

\emph{Part (b2).} Let $g$ be a smooth function on $\mathbb{R}^3$ and $j\in\{1,2,3\}$. Then $\delta_{\eta,j}g$ and $J_{\eta}g$ are smooth, and
\[
\operatorname{supp}\delta_{\eta,j}g\subset\operatorname{supp}g+\left[-\frac{\eta}{2},\frac{\eta}{2}\right]\cdot e_j,\qquad\operatorname{supp}J_{\eta}g\subset\operatorname{supp}g+\left[-\frac{\eta}{2},\frac{\eta}{2}\right]^3.
\]
The components of $a'$ are sums of single centred differences of components of $\beta$; $\check{a}_{\eta}=J_{\eta}a'$; and the components of $f^{(\eta)}$ are sums of terms $\delta_{\eta,\mu}\delta_{\eta,\mathfrak{l}}\beta_{\mathfrak{l}\nu}$ with two centred differences. Hence $\check{a}_{\eta}$ and $f^{(\eta)}$ are smooth, and
\[
\operatorname{supp}\check{a}_{\eta}\cup\operatorname{supp}f^{(\eta)}\subset\operatorname{supp}\beta+[-\eta,\eta]^3.
\]
Let $b=(y;\{\nu\})$ be a bond with $a_{\eta,\nu}(y)\ne0$. By (b1) the centre $c(b)=y+\frac{\eta}{2}e_{\nu}$ lies in $\operatorname{supp}\check{a}_{\eta,\nu}$, and the end points $y$ and $y+\eta e_{\nu}$ of $b$ are $c(b)\mp\frac{\eta}{2}e_{\nu}$; so both lie in $\operatorname{supp}\beta+\left[-\frac{3\eta}{2},\frac{3\eta}{2}\right]^3$. As $\eta\le1$, one has $\frac{3\eta}{2}\le2$, and $\operatorname{supp}\beta+[-2,2]^3\subset Q_{\beta}$ by the definition of $\rho_{\beta}$ and $y_{\beta}$ in Section~\ref{sec:probes-discretisation}. These are the first two assertions of (b2).

Now let $N\ge0$ and $T^{*}\ge1$ be integers with $\eta=L^{-N}$, let $\beta$ fit $(N,T^{*})$, and let $v\in\eta\mathbb{Z}^3$ be chosen as in Section~\ref{sec:probes-discretisation}, that is, with $\left|v+y_{\beta}-\frac{l_c}{2}(1,1,1)\right|_{\infty}\le\frac{\eta}{2}$. Here $R$ is any number as in (b), and $P=y_P+[0,R)^3$, $y_P\in\mathbb{R}^3$, denotes a period cube in whose interior $y_P+(0,R)^3$ the cube $Q_{\beta}$ lies. The translate $\tau_v\beta=\beta(\cdot-v)$ is a probe with support $v+\operatorname{supp}\beta$, and a cube $z+[-\rho,\rho]^3$ contains $\operatorname{supp}\tau_v\beta+[-2,2]^3=v+\left(\operatorname{supp}\beta+[-2,2]^3\right)$ if and only if the cube $(z-v)+[-\rho,\rho]^3$ contains $\operatorname{supp}\beta+[-2,2]^3$. Hence $\rho_{\tau_v\beta}=\rho_{\beta}$ and $D_{\tau_v\beta}=D_{\beta}$; the point $v+y_{\beta}$ may be taken as the point $y_{\tau_v\beta}$ of Section~\ref{sec:probes-discretisation}, and then the cube of $\tau_v\beta$ is $v+Q_{\beta}$. From $Q_{\beta}\subset y_P+(0,R)^3$ it follows that $v+Q_{\beta}\subset(v+y_P)+(0,R)^3$, which is the interior of $v+P=(v+y_P)+[0,R)^3$, again a period cube of the torus of side $R$. Therefore $\tau_v\beta$ satisfies the hypothesis of (b) on $\Lambda_{\eta,R}$, with the cube $v+Q_{\beta}$ in place of $Q_{\beta}$, and (b1) and the first two assertions of (b2) hold for $\tau_v\beta$. As $v\in\eta\mathbb{Z}^3$ and the translations commute with $\delta_{\eta,j}$ and $J_{\eta}$, the objects of Section~\ref{sec:probes-discretisation} formed from $\tau_v\beta$ are the translates by $v$ of those formed from $\beta$; in particular the function $f^{(\eta)}$ of $\tau_v\beta$ is $\tau_vf^{(\eta)}$. The cochain $\hat{\mathfrak{a}}=\eta^{\frac{1}{2}}\cdot a_{\eta}$ is formed from $\tau_v\beta$ on $\Lambda_{\eta,R}$.

Every coordinate of a point of $v+Q_{\beta}=v+y_{\beta}+[-\rho_{\beta},\rho_{\beta}]^3$ lies between $\frac{l_c}{2}-\frac{\eta}{2}-\rho_{\beta}$ and $\frac{l_c}{2}+\frac{\eta}{2}+\rho_{\beta}$. The inequality $2\rho_{\beta}+3\eta\le l_c$ by which $\beta$ fits says that $\frac{l_c}{2}-\frac{\eta}{2}-\rho_{\beta}\ge\eta$, or equivalently that $\frac{l_c}{2}+\frac{\eta}{2}+\rho_{\beta}\le l_c-\eta$. Therefore, as subsets of $\mathbb{R}^3$,
\[
v+Q_{\beta}\subset[\eta,l_c-\eta]^3.
\]

For the assertions of (b2) that refer to the lattice coordinates and to the cells of Section~\ref{sec:reference-tori-cells}, suppose in addition that $R=k\cdot R^{*}$ with an integer $k\ge1$, so that $\Lambda_{\eta,R}$ is the torus $k$ of Section~\ref{sec:reference-tori-cells}, and let the cells $\mathfrak{c}_c$, the sets $P^o(\mathfrak{c}_c)$ and the maps $\mathsf{g}_c$ be those of Section~\ref{sec:reference-tori-cells} on this torus. Since $l_c<R^{*}\le R$, the cube $[\eta,l_c-\eta]^3$ lies in the interior of the period cube $[0,R)^3$, so that the points of $v+Q_{\beta}$ are represented in $[0,R)^3$ by themselves. By the first part of (b2), applied to $\tau_v\beta$, both end points of every bond in the support of $\hat{\mathfrak{a}}$ lie in $v+Q_{\beta}\subset[\eta,l_c-\eta]^3$. The lattice coordinates of Section~\ref{sec:reference-tori-cells} are the physical coordinates divided by $\eta$, and $l_c=\left(m_{\mathrm{p}}+\frac{1}{2}\right)\eta$; so these end points have integer coordinates in $\left[1,m_{\mathrm{p}}-\frac{1}{2}\right]$, that is, in $\{1,\dots,m_{\mathrm{p}}-1\}$.

Let $p$ be a plaquette with $(d\hat{\mathfrak{a}})(p)\ne0$. The number $(d\hat{\mathfrak{a}})(p)$ is a sum, with signs, of the values of $\hat{\mathfrak{a}}$ on the four bonds in the boundary of $p$; so $p$ contains a bond $b$ of the support of $\hat{\mathfrak{a}}$. The vertices of $p$ are the end points of $b$ and their translates by one lattice unit in the other direction of $p$. Hence all vertices of $p$ lie in $\{0,\dots,m_{\mathrm{p}}\}^3=\mathfrak{c}_0$, and two of them, the end points of $b$, lie off each of the six planes $\{x_{\mu}=0\}$, $\{x_{\mu}=m_{\mathrm{p}}\}$, $\mu=1,2,3$. These are the planes through the end points of the three intervals of $\mathfrak{c}_0$, and $p$ is contained in none of them; so $p\in P^o(\mathfrak{c}_0)$. Thus $d\hat{\mathfrak{a}}$ is supported in $P^o(\mathfrak{c}_0)$.

In physical coordinates $\mathsf{g}_c$ is the map $y\mapsto ry+\eta t$ with a diagonal matrix $r\in O_{\mathrm{cub}}$ and a vector $t\in\mathbb{Z}^3$ (Section~\ref{sec:reference-tori-cells}). So $(\mathsf{g}_c)_{*}$ is the composition of the maps $r_{*}$ and $(\eta t)_{*}$ of Lemma~\ref{lem:lattice-symmetries-cell}, which holds on the torus $k$ (Section~\ref{sec:lattices-tests-function}); hence $(\mathsf{g}_c)_{*}$ is orthogonal and commutes with $d$. As $\mathsf{g}_c$ maps $\mathfrak{c}_0$ onto $\mathfrak{c}_c$, and in each direction $\mu$ the end points of $I_0$ onto those of $I_{c_{\mu}}$, the map $(\mathsf{g}_c)_{*}$ carries cochains supported in $P^o(\mathfrak{c}_0)$ to cochains supported in $P^o(\mathfrak{c}_c)$. Consequently
\[
d\hat{\mathfrak{a}}_c=d\,(\mathsf{g}_c)_{*}\hat{\mathfrak{a}}=(\mathsf{g}_c)_{*}\,d\hat{\mathfrak{a}}
\]
is supported in $P^o(\mathfrak{c}_c)$, and $\|d\hat{\mathfrak{a}}_c\|^2=\|d\hat{\mathfrak{a}}\|^2=\boldsymbol{\sigma}^2$.

Let $c\ne c'$ be two cells, and choose $\mu$ with $c_{\mu}\ne c'_{\mu}$, so that $I_{c_{\mu}}\ne I_{c'_{\mu}}$. By Section~\ref{sec:reference-tori-cells}, two intervals that are not consecutive are disjoint, while two consecutive intervals are either separated by a link plane or have a site plane in common, in which case they share exactly one site, the last site of the one and the first site of the other. Consequently $I_{c_{\mu}}\cap I_{c'_{\mu}}$ is either empty or consists of a single site $t$, at which one of the two intervals ends and the other begins. A plaquette whose four vertices lie both in $\mathfrak{c}_c$ and in $\mathfrak{c}_{c'}$ can therefore exist only in the second case, and then all its vertices lie on the plane $\{x_{\mu}=t\}$, so that it is contained in this plane. As $t$ is an end point of $I_{c_{\mu}}$ and of $I_{c'_{\mu}}$, the plane $\{x_{\mu}=t\}$ is one of the six planes in the definition of $P^o(\mathfrak{c}_c)$ and one of the six planes in the definition of $P^o(\mathfrak{c}_{c'})$; so the plaquette lies in neither of these two sets. Hence $P^o(\mathfrak{c}_c)\cap P^o(\mathfrak{c}_{c'})=\emptyset$, and the supports of the cochains $d\hat{\mathfrak{a}}_c$ are pairwise disjoint.

Finally, $R$ is again any number as in (b), not necessarily of the form $k\cdot R^{*}$. Since $\tau_v\beta$ satisfies the hypothesis of (b) on $\Lambda_{\eta,R}$, its cube $v+Q_{\beta}$ lying in the interior of the period cube $v+P$, (b1) for $\tau_v\beta$ gives $\left(\tau_vf^{(\eta)}\right)_{\eta}=d_{\eta}a_{\eta}=\eta^{-1}\cdot da_{\eta}$, where $a_{\eta}$ is formed from $\tau_v\beta$. With $F=\eta^{-\frac{3}{2}}\cdot F^u$ (Section~\ref{sec:lattices-tests-function}) and the sum over the plaquettes $p$,
\begin{align*}
X_{\eta,R}\left(\tau_vf^{(\eta)}\right)&=\left\langle\left(\tau_vf^{(\eta)}\right)_{\eta},F\right\rangle_{\eta}=\eta^3\sum_p\left(\eta^{-1}\cdot(da_{\eta})(p)\right)\cdot\left(\eta^{-\frac{3}{2}}\cdot F^u(p)\right)\\
&=\eta^{\frac{1}{2}}\cdot\langle da_{\eta},F^u\rangle=\langle d\hat{\mathfrak{a}},F^u\rangle=X.
\end{align*}
This proves (b2).

\emph{Part (b3).} With the Fourier transform of Section~\ref{sec:lattices-tests-function}, Plancherel's theorem gives $\|g\|_2^2=\int\frac{\mathrm{d}p}{(2\pi)^3}\,|\hat{g}(p)|^2$, so that the norm of $H^0$ is the norm of $L^2$. The derivative $\partial_j$ acts as the multiplier $ip_j$, and $\delta_{\eta,j}$ as the multiplier $i\xi_j(p)$, where
\[
\xi_j(p)=\frac{2}{\eta}\cdot\sin\left(\frac{\eta p_j}{2}\right).
\]
The inequalities $|\sin t|\le|t|$ and $|\sin t-t|\le\frac{|t|^3}{6}$, $t\in\mathbb{R}$, give
\[
|\xi_j(p)|\le|p_j|,\qquad|\xi_j(p)-p_j|\le\frac{\eta^2}{24}\cdot|p_j|^3.
\]
Writing $\xi_{\mu}\xi_{\mathfrak{l}}-p_{\mu}p_{\mathfrak{l}}=(\xi_{\mu}-p_{\mu})\cdot\xi_{\mathfrak{l}}+p_{\mu}\cdot(\xi_{\mathfrak{l}}-p_{\mathfrak{l}})$, one obtains for all $\mu,\mathfrak{l}\in\{1,2,3\}$
\[
|\xi_{\mu}\xi_{\mathfrak{l}}-p_{\mu}p_{\mathfrak{l}}|\le\frac{\eta^2}{24}\cdot\left(|p_{\mu}|^3\cdot|p_{\mathfrak{l}}|+|p_{\mu}|\cdot|p_{\mathfrak{l}}|^3\right)\le\frac{\eta^2}{12}\cdot|p|_2^4,
\]
the last step because $|p_{\mu}|,|p_{\mathfrak{l}}|\le|p|_2$. By Section~\ref{sec:probes-discretisation},
\[
f_{\mu\nu}=-\sum_{\mathfrak{l}\ne\nu}\partial_{\mu}\partial_{\mathfrak{l}}\beta_{\mathfrak{l}\nu}+\sum_{\mathfrak{l}\ne\mu}\partial_{\nu}\partial_{\mathfrak{l}}\beta_{\mathfrak{l}\mu},
\]
a sum of four terms, and $f^{(\eta)}_{\mu\nu}$ is the same sum with $\partial_{\mu},\partial_{\nu},\partial_{\mathfrak{l}}$ replaced by $\delta_{\eta,\mu},\delta_{\eta,\nu},\delta_{\eta,\mathfrak{l}}$. So $\hat{f}^{(\eta)}_{\mu\nu}(p)-\hat{f}_{\mu\nu}(p)$ is a sum of four products, each of a number $\pm\left(\xi_{\kappa}\xi_{\mathfrak{l}}-p_{\kappa}p_{\mathfrak{l}}\right)$ and of a value of a component of $\hat{\beta}$, and the Cauchy--Schwarz inequality gives
\[
\left|\hat{f}^{(\eta)}_{\mu\nu}(p)-\hat{f}_{\mu\nu}(p)\right|^2\le4\left(\frac{\eta^2}{12}\right)^2\cdot|p|_2^8\cdot S_{\mu\nu}(p),
\]
where $S_{\mu\nu}(p)$ is the sum of the four numbers $\left|\hat{\beta}_{\mathfrak{l}\nu}(p)\right|^2$, $\mathfrak{l}\ne\nu$, and $\left|\hat{\beta}_{\mathfrak{l}\mu}(p)\right|^2$, $\mathfrak{l}\ne\mu$. Let $\iota$ be the direction different from $\mu$ and $\nu$. Since $\left|\hat{\beta}_{\nu\mu}\right|=\left|\hat{\beta}_{\mu\nu}\right|$,
\[
S_{\mu\nu}(p)=2\left|\hat{\beta}_{\mu\nu}(p)\right|^2+\left|\hat{\beta}_{\iota\nu}(p)\right|^2+\left|\hat{\beta}_{\iota\mu}(p)\right|^2,
\]
that is, $S_{\mu\nu}(p)$ is twice the squared modulus of the component of the pair $\{\mu,\nu\}$ plus the squared moduli of the components of the two other pairs. As $\{\mu,\nu\}$ ranges over the three pairs, the squared modulus of a fixed component thus receives the weight $2$ once and the weight $1$ twice, whence
\[
\sum_{\mu<\nu}S_{\mu\nu}(p)=4\sum_{\mu<\nu}\left|\hat{\beta}_{\mu\nu}(p)\right|^2;
\]
as $4\cdot\left(\frac{\eta^2}{12}\right)^2\cdot4=\left(\frac{\eta^2}{3}\right)^2$, it follows that
\[
\sum_{\mu<\nu}\left|\hat{f}^{(\eta)}_{\mu\nu}(p)-\hat{f}_{\mu\nu}(p)\right|^2\le\left(\frac{\eta^2}{3}\right)^2\cdot|p|_2^8\cdot\sum_{\mu<\nu}\left|\hat{\beta}_{\mu\nu}(p)\right|^2.
\]
Let $s\ge0$. Multiplying by $\left(1+|p|_2^2\right)^s$, using $|p|_2^8\le\left(1+|p|_2^2\right)^4$ and integrating against $\frac{\mathrm{d}p}{(2\pi)^3}$, one obtains
\[
\left\|f^{(\eta)}-f\right\|_{H^s}\le\frac{\eta^2}{3}\cdot\|\beta\|_{H^{s+4}}\le\eta^2\cdot\|\beta\|_{H^{s+4}}.
\]

\emph{Part (b4).} For an integer $i\ge0$ and a function $g$ of class $C^i$ on $\mathbb{R}^3$, $|g|_{C^i}$ denotes the largest of the sup norms of the partial derivatives $\partial^{\alpha}g$ with $|\alpha|\le i$; for a one-form or a vector field with components of class $C^i$, it denotes the largest of these numbers for the components, in the same way as $|\beta|_{C^i}$ for the two-form $\beta$ (Section~\ref{sec:lattices-tests-function}). Let $g$ be a smooth function with compact support on $\mathbb{R}^3$. By the fundamental theorem of calculus,
\[
\delta_{\eta,\mu}g(y)=\int_{-\frac{1}{2}}^{\frac{1}{2}}\mathrm{d}s\,(\partial_{\mu}g)(y+s\eta e_{\mu}),
\]
so that $|\delta_{\eta,\mu}g|_{\infty}\le|\partial_{\mu}g|_{\infty}$; and $|J_{\eta}g|_{\infty}\le|g|_{\infty}$, since $J_{\eta}g$ is a mean of translates of $g$. For a multi-index $\alpha$ with $|\alpha|=i$,
\[
\partial^{\alpha}\check{a}_{\eta,\nu}=-\sum_{\mu\ne\nu}J_{\eta}\delta_{\eta,\mu}\partial^{\alpha}\beta_{\mu\nu}
\]
is a sum of two terms, one for each $\mu\ne\nu$, and each of them has sup norm at most $|\partial_{\mu}\partial^{\alpha}\beta_{\mu\nu}|_{\infty}\le|\beta|_{C^{i+1}}$. This gives $\left|D^i\check{a}_{\eta}\right|_{\infty}\le2\,|\beta|_{C^{i+1}}$.

Taylor's formula of order two for $s\mapsto(\partial_{\mu}g)(y+s\eta e_{\mu})$ at $s=0$ has the term of first order $s\eta\cdot\partial_{\mu}^2g(y)$, whose integral over $\left[-\frac{1}{2},\frac{1}{2}\right]$ vanishes, and a remainder of modulus at most $\frac{s^2\eta^2}{2}\cdot|\partial_{\mu}^3g|_{\infty}$. Hence
\[
|\delta_{\eta,\mu}g-\partial_{\mu}g|_{\infty}\le\int_{-\frac{1}{2}}^{\frac{1}{2}}\mathrm{d}s\,\frac{s^2\eta^2}{2}\cdot|\partial_{\mu}^3g|_{\infty}=\frac{\eta^2}{24}\cdot|\partial_{\mu}^3g|_{\infty}.
\]
Likewise, for $x\in\mathbb{R}^3$ and $y\in\left[-\frac{\eta}{2},\frac{\eta}{2}\right]^3$, Taylor's formula gives
\[
\left|g(x+y)-g(x)-\sum_{j=1}^3y_j\cdot\partial_jg(x)\right|\le\frac{1}{2}\left(|y_1|+|y_2|+|y_3|\right)^2\cdot|g|_{C^2}\le\frac{1}{2}\left(\frac{3\eta}{2}\right)^2\cdot|g|_{C^2},
\]
and the term of first order has mean zero over the cube; hence
\[
|J_{\eta}g-g|_{\infty}\le\frac{1}{2}\left(\frac{3\eta}{2}\right)^2\cdot|g|_{C^2}=\frac{9}{8}\,\eta^2\cdot|g|_{C^2}.
\]
For $|\alpha|\le i$,
\[
\partial^{\alpha}\left(\check{a}_{\eta,\nu}-a_{\nu}\right)=-\sum_{\mu\ne\nu}\left[J_{\eta}\left(\delta_{\eta,\mu}-\partial_{\mu}\right)\partial^{\alpha}\beta_{\mu\nu}+\left(J_{\eta}-1\right)\partial_{\mu}\partial^{\alpha}\beta_{\mu\nu}\right],
\]
and the two bounds, applied to $g=\partial^{\alpha}\beta_{\mu\nu}$, for which $|\partial_{\mu}^3g|_{\infty}\le|\beta|_{C^{i+3}}$, and to $g=\partial_{\mu}\partial^{\alpha}\beta_{\mu\nu}$, for which $|g|_{C^2}\le|\beta|_{C^{i+3}}$ (both because $|\alpha|\le i$), together with the inequality $|J_{\eta}g|_{\infty}\le|g|_{\infty}$ of the first step, give
\[
\left|\check{a}_{\eta,\nu}-a_{\nu}\right|_{C^i}\le2\left(\frac{1}{24}+\frac{9}{8}\right)\eta^2\cdot|\beta|_{C^{i+3}}=\frac{7}{3}\,\eta^2\cdot|\beta|_{C^{i+3}}\le3\eta^2\cdot|\beta|_{C^{i+3}}.
\]
This holds for $\nu=1,2,3$; by the definition of $|\cdot|_{C^i}$ for one-forms, it is the second inequality of (b4).

\emph{Part (b5).} Let $N\ge0$ and $T^{*}\ge1$ be integers with $\eta=L^{-N}$, let $\beta$ fit $(N,T^{*})$, let $v\in\eta\mathbb{Z}^3$ be chosen as in Section~\ref{sec:probes-discretisation}, and let $R$ be any number as in (b), with a period cube $P$ of the torus of side $R$ in whose interior $Q_{\beta}$ lies; the cochain $\hat{\mathfrak{a}}$ and the number $\boldsymbol{\sigma}$ are formed with $\tau_v\beta$ on $\Lambda_{\eta,R}$. Put $g=\tau_vf^{(\eta)}$. As shown in the proof of (b2), $\tau_v\beta$ satisfies the hypothesis of (b) on $\Lambda_{\eta,R}$, since its cube $v+Q_{\beta}$ lies in the interior of the period cube $v+P$. So (b1) holds for $\tau_v\beta$, and, with $a_{\eta}$ formed from $\tau_v\beta$,
\[
d\hat{\mathfrak{a}}=\eta^{\frac{1}{2}}\cdot da_{\eta}=\eta^{\frac{3}{2}}\cdot d_{\eta}a_{\eta}=\eta^{\frac{3}{2}}\cdot g_{\eta},
\]
and the normalisation in lattice units of Section~\ref{sec:lattices-tests-function}, $|w|_{\eta}=\left\|\eta^{\frac{3}{2}}\cdot w\right\|$ for $w\in C^2$, gives
\[
\boldsymbol{\sigma}=\|d\hat{\mathfrak{a}}\|=|g_{\eta}|_{\eta}.
\]
As $g$ is the function $f^{(\eta)}$ of $\tau_v\beta$, the first assertion of (b2) for $\tau_v\beta$ shows that $g$ is supported in $v+\operatorname{supp}\beta+[-\eta,\eta]^3$, which is contained in $v+\operatorname{supp}\beta+[-2,2]^3\subset v+Q_{\beta}$; and $v+Q_{\beta}$ lies in the interior of the period cube $v+P$ of the torus of side $R$, because $Q_{\beta}$ lies in the interior of $P$. A two-form on $\mathbb{R}^3$ with support in $v+Q_{\beta}$ is identified with its periodisation, a two-form on $\mathbb{T}_R$ with the same norm in $L^2$; and $\|\cdot\|_2$ also denotes the norm of $L^2(\mathbb{T}_R;\Lambda^2)$.

For a two-form $\zeta$ on $\mathbb{R}^3$ with support in $v+Q_{\beta}$ let $\operatorname{Ev}\zeta$ be the two-form on $\mathbb{T}_R$ whose component $\mu\nu$ is constant on each averaging cube $Q(p)$, $p=(x;\mu\nu)$, with the value $\zeta_{\eta}(p)$, the mean value of $\zeta_{\mu\nu}$ over $Q(p)$. The cubes $Q(p)$ of one orientation $\mu\nu$ have volume $\eta^3$ and tile $\mathbb{T}_R$ (last paragraph of Section~\ref{sec:test-forms-test}), and the Cauchy--Schwarz inequality on each cube gives $|\zeta_{\eta}(p)|^2\le\eta^{-3}\int_{Q(p)}\mathrm{d}y\,|\zeta_{\mu\nu}(y)|^2$. Therefore, as in the proof of Lemma~\ref{lem:norm-bounds-lattice},
\[
\|\operatorname{Ev}\zeta\|_2^2=\eta^3\sum_{\mu<\nu}\sum_x\left|\zeta_{\eta}(x;\mu\nu)\right|^2=|\zeta_{\eta}|_{\eta}^2\le\sum_{\mu<\nu}\sum_x\int_{Q(x;\mu\nu)}\mathrm{d}y\,|\zeta_{\mu\nu}(y)|^2=\|\zeta\|_2^2.
\]
For $\zeta=g$ this gives, together with (b3) for $s=0$ and with $\eta\le1$,
\begin{align*}
\boldsymbol{\sigma}=\|\operatorname{Ev} g\|_2&\le\|g\|_2=\left\|f^{(\eta)}\right\|_2\le\|f\|_2+\left\|f^{(\eta)}-f\right\|_2\\
&\le\|f\|_2+\eta^2\cdot\|\beta\|_{H^4}\le\|f\|_2+\|\beta\|_{H^4}=s_f,
\end{align*}
which is the first assertion of (b5).

Let $y$ and $y'$ lie in one cube $Q(p)$, $p=(x;\mu\nu)$. Then $|y_j-y'_j|\le\eta$ for $j=1,2,3$, and the fundamental theorem of calculus along the segment from $y'$ to $y$, which lies in $Q(p)$, gives
\[
\left|g_{\mu\nu}(y)-g_{\mu\nu}(y')\right|\le3\eta\cdot\max_j\left|\partial_jg_{\mu\nu}\right|_{\infty}.
\]
Moreover
\[
\partial_jf^{(\eta)}_{\mu\nu}=-\sum_{\mathfrak{l}\ne\nu}\delta_{\eta,\mu}\delta_{\eta,\mathfrak{l}}\partial_j\beta_{\mathfrak{l}\nu}+\sum_{\mathfrak{l}\ne\mu}\delta_{\eta,\nu}\delta_{\eta,\mathfrak{l}}\partial_j\beta_{\mathfrak{l}\mu}
\]
is a sum of four terms, and by the bound $|\delta_{\eta,\kappa}w|_{\infty}\le|\partial_{\kappa}w|_{\infty}$ of the first step of Part (b4), applied twice, each of them has sup norm at most $|\beta|_{C^3}$. Hence $\max_j|\partial_jg_{\mu\nu}|_{\infty}\le4\,|\beta|_{C^3}$; and, as $(\operatorname{Ev} g)_{\mu\nu}(y)$ is the mean value of $g_{\mu\nu}$ over the cube $Q(p)$ that contains $y$,
\[
\left|(\operatorname{Ev} g)_{\mu\nu}-g_{\mu\nu}\right|\le12\eta\cdot|\beta|_{C^3}\qquad\text{on }\mathbb{T}_R,\ \text{for }\mu<\nu.
\]
Now fix $\mu<\nu$ and a cube $Q(p)$, $p=(x;\mu\nu)$, of this orientation. If $Q(p)$ does not meet $\operatorname{supp}g$, then $g_{\mu\nu}=0$ on $Q(p)$, hence the value $g_{\eta}(p)$ of $(\operatorname{Ev} g)_{\mu\nu}$ on $Q(p)$ is $0$, and $(\operatorname{Ev} g-g)_{\mu\nu}=0$ on $Q(p)$. If $Q(p)$ meets $\operatorname{supp}g$, then $Q(p)\subset\operatorname{supp}g+[-\eta,\eta]^3$, as two points of $Q(p)$ differ by less than $\eta$ in each coordinate. Since the cubes of the orientation $\mu\nu$ cover $\mathbb{T}_R$, the component $(\operatorname{Ev} g-g)_{\mu\nu}$ vanishes outside $\operatorname{supp}g+[-\eta,\eta]^3$. This holds for each of the three pairs $\mu<\nu$; moreover, as $2\eta\le2$,
\[
\operatorname{supp}g+[-\eta,\eta]^3\subset v+\operatorname{supp}\beta+[-2\eta,2\eta]^3\subset v+Q_{\beta},
\]
a cube of volume $(2\rho_{\beta})^3=D_{\beta}^3$. Integrating the square of the bound on $\left|(\operatorname{Ev} g)_{\mu\nu}-g_{\mu\nu}\right|$ over this cube and summing over the three pairs $\mu<\nu$, one obtains
\[
\|\operatorname{Ev} g-g\|_2\le\left(3\cdot\left(12\eta\cdot|\beta|_{C^3}\right)^2\cdot D_{\beta}^3\right)^{\frac{1}{2}}\le21\eta\cdot|\beta|_{C^3}\cdot D_{\beta}^{\frac{3}{2}},
\]
because $3\cdot12^2=432\le441=21^2$. Finally $\|f\|_2=\|\tau_vf\|_2$ and $g-\tau_vf=\tau_v\left(f^{(\eta)}-f\right)$, so the triangle inequality gives
\[
\left|\boldsymbol{\sigma}-\|f\|_2\right|=\left|\|\operatorname{Ev} g\|_2-\|\tau_vf\|_2\right|\le\|\operatorname{Ev} g-g\|_2+\left\|f^{(\eta)}-f\right\|_2,
\]
and with (b3) for $s=0$
\[
\left|\boldsymbol{\sigma}-\|f\|_2\right|\le\eta^2\cdot\|\beta\|_{H^4}+21\eta\cdot|\beta|_{C^3}\cdot D_{\beta}^{\frac{3}{2}}.
\]
\end{proof}

\subsubsection{Symmetries of the determinant and of the lattice functionals}\label{sec:symmetries-determinant-lattice}

\begin{proof}[Proof of Theorem~\ref{thm:free-coefficient-1}(d)]
Let $B$ be a complex one-cochain on $\Lambda_n$ and $\boldsymbol{\mu}$ a real number. A \emph{period cell} is a subset of $\mathbb{Z}^3$ that contains exactly one point of every class of $\mathbb{Z}^3$ modulo the lattice $\prod_{\mu}n_{\mu}\mathbb{Z}$; its points represent the sites of $\Lambda_n$. A spinor is determined by its values on a period cell. For a point $y$ of the cell and $\alpha\in\{1,2,3,4\}$ let $\psi_{y,\alpha}$ be the spinor whose value at $y$ is the $\alpha$-th standard basis vector of $\mathbb{C}^4$ and whose values at the other points of the cell are $0$. These $4|n|$ spinors form a basis of the space of the spinors, called the \emph{basis indexed by the period cell}; in it an operator on the spinors is a matrix of $4\times4$ blocks, one block for each ordered pair of points of the cell. A point of $\mathbb{Z}^3$ written as an index of a block stands for the point of the cell in its class; the values of a spinor at these two points are related by the antiperiodicity.

\emph{Part (d1).} Let $\boldsymbol{\lambda}$ be a complex zero-cochain on $\Lambda_n$, that is, a periodic function $\boldsymbol{\lambda}\colon\mathbb{Z}^3\to\mathbb{C}$. Multiplication by $\exp(i\boldsymbol{\lambda})$ maps spinors to spinors, and it is invertible on them, the inverse being the multiplication by $\exp(-i\boldsymbol{\lambda})$. In $\exp(-i\boldsymbol{\lambda})\cdot D[B]\cdot\exp(i\boldsymbol{\lambda})$ the term of \eqref{eq:free-coefficient-5} that belongs to the neighbour $x'$ of $x$ acquires the factor $\exp\left(i\boldsymbol{\lambda}(x')-i\boldsymbol{\lambda}(x)\right)$; hence
\[
\exp(-i\boldsymbol{\lambda})\cdot D[B]\cdot\exp(i\boldsymbol{\lambda})=D[B+d\boldsymbol{\lambda}],\qquad\text{where }(d\boldsymbol{\lambda})(x,x')=\boldsymbol{\lambda}(x')-\boldsymbol{\lambda}(x),
\]
and taking determinants gives $\mathrm{Dt}(B+d\boldsymbol{\lambda})=\mathrm{Dt}(B)$. By \eqref{eq:free-coefficient-5} the operator $D[B]$ depends on $B$ only through the numbers $\exp\left(iB(x,x')\right)$. These do not change when $B$ is replaced by $B+2\pi\cdot\mathfrak{z}$ with an integer one-cochain $\mathfrak{z}$; so $\mathrm{Dt}(B+2\pi\cdot\mathfrak{z})=\mathrm{Dt}(B)$.

Write $D[B]+\boldsymbol{\mu}$ as a matrix with respect to the basis indexed by a period cell. By \eqref{eq:free-coefficient-5} each of its entries is a constant or a linear combination, with constant coefficients, of numbers $\exp\left(iB(x,x')\right)$; so every entry is an entire function of $B$. The modulus of an entry is at most $\exp\left(|\operatorname{Im}B|_{\infty}\right)$ if the entry lies outside the diagonal blocks, and at most $3+|\boldsymbol{\mu}|$ if it lies in a diagonal block. A determinant is a polynomial in the entries of the matrix; hence $\mathrm{Dt}$ is entire. For a matrix $O$ of dimension $4|n|$ one has
\[
|\det O|\le\|O\|^{4|n|},
\]
and the operator norm $\|O\|$ is at most $4|n|$ times the largest modulus of an entry of $O$. Hence, for every $c\ge0$, $\mathrm{Dt}$ is bounded on the set of the complex one-cochains $B$ with $|\operatorname{Im}B|_{\infty}\le c$.

\emph{Part (d2).} In this part $B$ is real. The space of the spinors carries the inner product
\[
\langle\psi,\varphi\rangle:=\sum_x\psi(x)^{*}\varphi(x),
\]
the sum being taken over a period cell; the summand is periodic, so that the inner product does not depend on the cell. For $\mu=1,2,3$ let $T_{\mu}$ be the operator on the spinors given by
\[
\left(T_{\mu}\psi\right)(x):=\exp\left(iB(x,x+e_{\mu})\right)\cdot\psi(x+e_{\mu}).
\]
It is unitary, and its adjoint is given by $\left(T_{\mu}^{*}\psi\right)(x)=\exp\left(iB(x,x-e_{\mu})\right)\cdot\psi(x-e_{\mu})$, since $B(x,x-e_{\mu})=-B(x-e_{\mu},x)$. In \eqref{eq:free-coefficient-5} the neighbour $x+e_{\mu}$ contributes $-\frac{1}{2}\left(1-\gamma_{\mu}\right)\left(T_{\mu}\psi\right)(x)$, and the neighbour $x-e_{\mu}$ contributes $-\frac{1}{2}\left(1+\gamma_{\mu}\right)\left(T_{\mu}^{*}\psi\right)(x)$. Therefore
\[
D[B]+\boldsymbol{\mu}=\boldsymbol{\mu}+\sum_{\mu}\left[1-\frac{1}{2}\left(T_{\mu}+T_{\mu}^{*}\right)\right]+\sum_{\mu}\gamma_{\mu}\cdot\frac{1}{2}\left(T_{\mu}-T_{\mu}^{*}\right).
\]
Write every spinor as $\psi=(\psi_{+},\psi_{-})$, where $\psi_{+}$ and $\psi_{-}$ consist of the first two and of the last two components in the decomposition $\mathbb{C}^4=\mathbb{C}^2\oplus\mathbb{C}^2$ of Section~\ref{sec:letters-norms-dirac}, in which $\gamma_{\mu}=\operatorname{diag}(\sigma_{\mu},-\sigma_{\mu})$. Then $\psi_{+}$ and $\psi_{-}$ are antiperiodic functions $\mathbb{Z}^3\to\mathbb{C}^2$, the operators $T_{\mu}$ act on $\psi_{+}$ and on $\psi_{-}$ by the same formula, and $D[B]+\boldsymbol{\mu}=\operatorname{diag}\left(\mathbf{B},\mathbf{B}'\right)$, where $\mathbf{B}$ is the operator on the antiperiodic functions $\mathbb{Z}^3\to\mathbb{C}^2$ given by
\[
\mathbf{B}:=\boldsymbol{\mu}+\sum_{\mu}\left[1-\frac{1}{2}\left(T_{\mu}+T_{\mu}^{*}\right)+\sigma_{\mu}\cdot\frac{1}{2}\left(T_{\mu}-T_{\mu}^{*}\right)\right],
\]
and $\mathbf{B}'$ is the operator obtained from this formula when every $\sigma_{\mu}$ is replaced by $-\sigma_{\mu}$. The matrix $\sigma_{\mu}=\sigma_{\mu}^{*}$ acts on the two components of the value at each point, whereas $T_{\mu}$ and $T_{\mu}^{*}$ shift the point and multiply by numbers; so $\sigma_{\mu}$ commutes with $T_{\mu}-T_{\mu}^{*}$, and the adjoint of $\sigma_{\mu}\cdot\left(T_{\mu}-T_{\mu}^{*}\right)$ is $\left(T_{\mu}^{*}-T_{\mu}\right)\cdot\sigma_{\mu}=-\sigma_{\mu}\cdot\left(T_{\mu}-T_{\mu}^{*}\right)$. As $\boldsymbol{\mu}$ is real, it follows that $\mathbf{B}'=\mathbf{B}^{*}$, and
\[
\mathrm{Dt}(B)=\det\mathbf{B}\cdot\det\mathbf{B}^{*}=\left|\det\mathbf{B}\right|^2\ge0.
\]
In the expression for $D[B]+\boldsymbol{\mu}$ above, the first sum is Hermitian and non-negative, because $\left\|\frac{1}{2}\left(T_{\mu}+T_{\mu}^{*}\right)\right\|\le1$ for the unitary operator $T_{\mu}$, and the second sum is anti-Hermitian. So the Hermitian part $\operatorname{Re}\left(D[B]+\boldsymbol{\mu}\right)$ equals $\boldsymbol{\mu}$ plus the first sum, and $\operatorname{Re}\left(D[B]+\boldsymbol{\mu}\right)\ge\boldsymbol{\mu}$.

\emph{Part (d3).} In this part and in the next, $B$ is again an arbitrary complex one-cochain on $\Lambda_n$. Fix a period cell, let $x$ be one of its points, let $x'$ be a neighbour of $x$ in $\mathbb{Z}^3$, and let $y'$ be the point of the cell in the class of $x'$. Then $\psi(x')=\varsigma\cdot\psi(y')$ for every spinor $\psi$, with the antiperiodic sign $\varsigma\in\{1,-1\}$. This sign is real, and the reverse bond $(x',x)$, translated by the period vector $y'-x'$ so that it starts in the cell, carries the same sign. By \eqref{eq:free-coefficient-5} the summand of $\left(D[B]\psi\right)(x)$ that belongs to the bond $(x,x')$ is the matrix
\[
-\frac{1}{2}\left(1-\gamma\cdot(x'-x)\right)\cdot\varsigma\cdot\exp\left(iB(x,x')\right)
\]
applied to $\psi(y')$. Thus, in the basis indexed by the cell, the bond $(x,x')$ adds this matrix to the block of $D[B]$ with row index $x$ and column index $x'$; the diagonal blocks of $D[B]$ equal $3$ times the unit matrix. Transposing a block matrix exchanges the row index and the column index of every block and transposes the block itself. Since $\exp\left(iB(x,x')\right)=\exp\left(-iB(x',x)\right)$, the share of the bond $(x,x')$ in the block of $D[B]^T$ with row index $x'$ and column index $x$ is therefore
\[
-\frac{1}{2}\left(1-\gamma\cdot(x'-x)\right)^T\cdot\varsigma\cdot\exp\left(-iB(x',x)\right).
\]
Let $\mathbf{C}:=\operatorname{diag}(\sigma_2,\sigma_2)$, acting at every site. Since $\sigma_2\sigma_{\mu}^T\sigma_2=-\sigma_{\mu}$ for $\mu=1,2,3$ and $\sigma_2^2=1$, one has $\mathbf{C}\cdot\gamma_{\mu}^T\cdot\mathbf{C}^{-1}=-\gamma_{\mu}$, and therefore
\[
\mathbf{C}\cdot\frac{1}{2}\left(1-\gamma\cdot(x'-x)\right)^T\cdot\mathbf{C}^{-1}=\frac{1}{2}\left(1-\gamma\cdot(x-x')\right).
\]
So conjugation by $\mathbf{C}$ carries the share of the bond $(x,x')$ in $D[B]^T$ to the share of the reverse bond $(x',x)$ in $D[-B]$, both standing in the block with row index $x'$ and column index $x$, and it leaves the diagonal blocks unchanged; that is, $\mathbf{C}\cdot D[B]^T\cdot\mathbf{C}^{-1}=D[-B]$. As $\det O^T=\det O$ for every square matrix $O$,
\[
\mathrm{Dt}(-B)=\det\left(\mathbf{C}\cdot\left(D[B]+\boldsymbol{\mu}\right)^T\cdot\mathbf{C}^{-1}\right)=\det\left(D[B]+\boldsymbol{\mu}\right)=\mathrm{Dt}(B).
\]

\emph{Part (d4).} For a unit vector $\mathsf{n}\in\mathbb{R}^3$ let $r_{\mathsf{n}}\colon\mathbb{R}^3\to\mathbb{R}^3$, $r_{\mathsf{n}}z:=z-2(\mathsf{n}\cdot z)\,\mathsf{n}$, be the reflection in the plane orthogonal to $\mathsf{n}$, and let $S_{\mathsf{n}}$ be the $4\times4$ matrix $S_{\mathsf{n}}:=(\gamma\cdot\mathsf{n})\cdot\hat{\Gamma}$. The relations of Section~\ref{sec:letters-norms-dirac} give $(\gamma\cdot\mathsf{n})^2=1$ and, for $z\in\mathbb{R}^3$,
\[
(\gamma\cdot\mathsf{n})(\gamma\cdot z)(\gamma\cdot\mathsf{n})=2(\mathsf{n}\cdot z)\,\gamma\cdot\mathsf{n}-\gamma\cdot z,\qquad\hat{\Gamma}\cdot(\gamma\cdot z)\cdot\hat{\Gamma}^{-1}=-\gamma\cdot z.
\]
Hence $S_{\mathsf{n}}$ is invertible, and
\[
S_{\mathsf{n}}\,(\gamma\cdot z)\,S_{\mathsf{n}}^{-1}=\gamma\cdot\left(r_{\mathsf{n}}z\right)\qquad\left(z\in\mathbb{R}^3\right).
\]
Every element of $O_{\mathrm{cub}}$, that is, every signed permutation matrix, is a product of reflections $r_{\mathsf{n}}$ of two kinds: those with $\mathsf{n}=e_{\mu}$, which reverse the sign of the coordinate $\mu$, and those with $\mathsf{n}=2^{-\frac{1}{2}}(e_{\mu}-e_{\nu})$, $\mu\ne\nu$, which interchange the coordinates $\mu$ and $\nu$; indeed, a signed permutation matrix is a diagonal matrix of signs times a permutation matrix, and every permutation is a product of transpositions. Let $r\in O_{\mathrm{cub}}$, write $r=r_{\mathsf{n}_1}\cdots r_{\mathsf{n}_k}$ with such reflections, and put $S_r:=S_{\mathsf{n}_1}\cdots S_{\mathsf{n}_k}$. Then $S_r$ is an invertible matrix, a product of matrices $S_{\mathsf{n}}$, and
\[
S_r\,(\gamma\cdot z)\,S_r^{-1}=\gamma\cdot(rz)\qquad\left(z\in\mathbb{R}^3\right).
\]

Now let $g$, $r$, $c$ and $n'$ be as in (d4). As $r$ maps $\prod_{\mu}n_{\mu}\mathbb{Z}$ onto $\prod_{\mu}n'_{\mu}\mathbb{Z}$, the map $g$ carries every coset of $\prod_{\mu}n_{\mu}\mathbb{Z}$ onto a coset of $\prod_{\mu}n'_{\mu}\mathbb{Z}$; it also maps nearest neighbours to nearest neighbours, so that it induces an isometry of $\Lambda_n$ onto $\Lambda_{n'}$. For a spinor $\psi$ of $\Lambda_n$ let $U_{\mathrm{g}}\psi\colon\mathbb{Z}^3\to\mathbb{C}^4$ be the function
\[
\left(U_{\mathrm{g}}\psi\right)(x):=S_r\,\psi\left(g^{-1}x\right).
\]
For a direction $\mu$ let $\nu$ be the direction with $re_{\nu}=\pm e_{\mu}$. Then $n'_{\mu}=n_{\nu}$, because $r$ maps $\prod_{\mu}n_{\mu}\mathbb{Z}$ onto $\prod_{\mu}n'_{\mu}\mathbb{Z}$, and
\[
g^{-1}\left(x+n'_{\mu}e_{\mu}\right)=g^{-1}x+n'_{\mu}r^{-1}e_{\mu}=g^{-1}x\pm n_{\nu}e_{\nu}.
\]
Hence $U_{\mathrm{g}}\psi$ is a spinor of $\Lambda_{n'}$, and $U_{\mathrm{g}}$ is a linear bijection from the space of the spinors of $\Lambda_n$ to the space of the spinors of $\Lambda_{n'}$, with inverse $\varphi\mapsto S_r^{-1}\,\varphi\circ g$. Let $D_n[B]$ denote the operator \eqref{eq:free-coefficient-5} on the spinors of $\Lambda_n$. For $x\in\mathbb{Z}^3$ the neighbours of $g^{-1}x$ are the points $y'=g^{-1}x'$ with $|x'-x|_2=1$, and $y'-g^{-1}x=r^{-1}(x'-x)$. Let $\varphi$ be a spinor of $\Lambda_{n'}$ and $x\in\mathbb{Z}^3$. Evaluating \eqref{eq:free-coefficient-5} for the spinor $U_{\mathrm{g}}^{-1}\varphi$ of $\Lambda_n$ at the point $g^{-1}x$, and then using the relation for $S_r$ at $z=r^{-1}(x'-x)$, one finds
\begin{multline*}
\left(U_{\mathrm{g}}\cdot D_n[B]\cdot U_{\mathrm{g}}^{-1}\varphi\right)(x)\\
=3\varphi(x)-\sum_{x'\in\mathbb{Z}^3:\ |x'-x|_2=1}S_r\cdot\frac{1}{2}\left(1-\gamma\cdot\left(r^{-1}(x'-x)\right)\right)\cdot S_r^{-1}\cdot\exp\left(iB\left(g^{-1}x,g^{-1}x'\right)\right)\cdot\varphi(x')\\
=\left(D_{n'}[g_{*}B]\varphi\right)(x).
\end{multline*}
As $\varphi$ and $x$ are arbitrary, $U_{\mathrm{g}}\cdot D_n[B]\cdot U_{\mathrm{g}}^{-1}=D_{n'}[g_{*}B]$. Consequently
\[
\mathrm{Dt}_{n'}\left(g_{*}B\right)=\det\left(U_{\mathrm{g}}\cdot\left(D_n[B]+\boldsymbol{\mu}\right)\cdot U_{\mathrm{g}}^{-1}\right)=\mathrm{Dt}_n(B).
\]

\emph{Part (d5).} Let $\Phi\in\mathfrak{F}(n)$, and let $c_{\Phi}$ be a constant with $|\Phi(F)|\le c_{\Phi}\cdot\exp\left(c_{\Phi}\|F\|\right)$ for all $F\in C^2(n)$. Let $\boldsymbol{\Psi}_{\Phi}=\boldsymbol{\Psi}^n_{\Phi}\colon C^1(n)\to\mathbb{C}$ be the function
\[
\boldsymbol{\Psi}_{\Phi}(A):=\sum_{\mathfrak{z}}\Phi\left(dA+r_q\cdot\mathfrak{z}\right)\cdot\exp\left(-\frac{1}{2}\left\|dA+r_q\cdot\mathfrak{z}\right\|^2\right)\cdot\mathrm{Dt}(qA),
\]
where the sum is taken over all $\mathfrak{z}\in C^2(n;\mathbb{Z})$ with $d\mathfrak{z}=0$ (respectively over those of them whose flux is zero). This series converges absolutely for every $A$, and uniformly for $A$ in the box $[0,r_q)^{\mathrm{links}}$ of Section~\ref{sec:wilson-determinant-lattice}. Indeed, the factor $\mathrm{Dt}(qA)$ does not depend on $\mathfrak{z}$, and by (d1), applied with $c=0$ to the real one-cochain $qA$, its modulus is bounded by a constant that does not depend on $A\in C^1(n)$. With $F=dA+r_q\cdot\mathfrak{z}$ the product of the two other factors has modulus at most $c_{\Phi}\cdot\exp\left(c_{\Phi}\|F\|-\frac{1}{2}\|F\|^2\right)$, and $\|F\|\ge r_q\cdot\|\mathfrak{z}\|-\|dA\|$. On the box the values of $A$ lie in $[0,r_q)$, so that the supremum of $\|dA\|$ over the box is finite. Finally, for $y\ge0$ an integer two-cochain $\mathfrak{z}$ with $\|\mathfrak{z}\|\le y$ has $3|n|$ values, each an integer in $[-y,y]$; so there are at most $(2y+1)^{3|n|}$ such $\mathfrak{z}$. By absolute convergence the series may be summed in any order and after any bijective change of the summation index. By the definitions of Section~\ref{sec:wilson-determinant-lattice}
\[
\boldsymbol{\omega}_n(\Phi)=\int_{[0,r_q)^{\mathrm{links}}}\mathrm{d}A\,\boldsymbol{\Psi}_{\Phi}(A),
\]
while $\boldsymbol{\omega}_n^{(0)}(\Phi)$ is the same integral of the sum over flux zero.

Let $C^1(n;\mathbb{Z})\subset C^1(n)$ be the set of the integer valued one-cochains, and let $c\in C^1(n;\mathbb{Z})$. Then $d(A+r_q\cdot c)+r_q\cdot\mathfrak{z}=dA+r_q\cdot(\mathfrak{z}+dc)$, and $\mathrm{Dt}\left(q(A+r_q\cdot c)\right)=\mathrm{Dt}(qA+2\pi\cdot c)=\mathrm{Dt}(qA)$ by (d1), as $q\cdot r_q=2\pi$. The map $\mathfrak{z}\mapsto\mathfrak{z}-dc$ permutes the closed integer two-cochains and preserves the flux, since the flux is additive and vanishes on $dc$ (property~\eqref{eq:partition-functions-2}). Therefore
\[
\boldsymbol{\Psi}_{\Phi}(A+r_q\cdot c)=\boldsymbol{\Psi}_{\Phi}(A)\qquad\left(A\in C^1(n),\ c\in C^1(n;\mathbb{Z})\right),
\]
for the full sum and for the sum over flux zero. The box is a fundamental domain of the lattice $r_q\cdot C^1(n;\mathbb{Z})$ in $C^1(n)$, and the Lebesgue measure on the box is the Haar measure of the torus $C^1(n)/r_q\cdot C^1(n;\mathbb{Z})$, of total mass $r_q^{3|n|}$. So $\boldsymbol{\omega}_n(\Phi)$ (respectively $\boldsymbol{\omega}_n^{(0)}(\Phi)$) is the integral of the periodic function $\boldsymbol{\Psi}_{\Phi}$ over this torus with respect to its Haar measure.

Now let $g$ be as in (d4), and let $g^{*}$ and $g_{*}=\left(g^{-1}\right)^{*}$ act on cochains as in Section~\ref{sec:reference-tori-cells}. Consider $g_{*}\colon C^1(n)\to C^1(n')$, $\left(g_{*}A\right)(gx,gx')=A(x,x')$. Under $g$ the image of a link of $\Lambda_n$ is a link of $\Lambda_{n'}$ or the reverse of one, and every link of $\Lambda_{n'}$ is obtained in this way from exactly one link of $\Lambda_n$; so, in the coordinates given by the values on the links, $g_{*}$ is a signed permutation matrix. Hence it maps the lattice $r_q\cdot C^1(n;\mathbb{Z})$ onto $r_q\cdot C^1(n';\mathbb{Z})$, it induces an isomorphism of the two tori, and it carries the Haar measure of the first torus to that of the second (both have total mass $r_q^{3|n|}$, as $|n'|=|n|$). The argument in the proof of Lemma~\ref{lem:lattice-symmetries-cell}, which does not use that the sides of the torus are equal, extends this to cochains of every degree; it applies to $g$, the translation by $c$ contributing no signs. For every $k$, the map $g$ carries the $k$-cells of $\Lambda_n$ bijectively onto those of $\Lambda_{n'}$, and $r$ carries the orientation of a $k$-cell to the orientation of its image cell up to a sign; so, in the coordinates given by the values on the $k$-cells, $g_{*}$ is a signed permutation matrix. Moreover, $g$ maps the faces of a cell onto the faces of its image cell and outward normals to outward normals; so the incidence number of the image of a cell and the image of one of its faces is the incidence number of the cell and the face multiplied by the two signs. Consequently, on cochains of every degree, $g_{*}$ is orthogonal, with inverse $g^{*}$, it maps the integer cochains of $\Lambda_n$ bijectively onto those of $\Lambda_{n'}$, and it commutes with $d$. Therefore $g_{*}$ maps the closed integer two-cochains of $\Lambda_n$ bijectively onto those of $\Lambda_{n'}$, and the set $\{dc:c\in C^1(n;\mathbb{Z})\}$ onto the set $\{dc':c'\in C^1(n';\mathbb{Z})\}$; by property~\eqref{eq:partition-functions-2} these two sets consist of the closed integer two-cochains of flux zero. Since $g^{*}$ is orthogonal, $\Phi\circ g^{*}$ lies in $\mathfrak{F}(n')$. For a closed integer two-cochain $\mathfrak{z}$ of $\Lambda_n$ and $\mathfrak{z}':=g_{*}\mathfrak{z}$ one has
\[
g^{*}\left(d\,g_{*}A+r_q\cdot\mathfrak{z}'\right)=dA+r_q\cdot\mathfrak{z},\qquad\left\|d\,g_{*}A+r_q\cdot\mathfrak{z}'\right\|=\left\|dA+r_q\cdot\mathfrak{z}\right\|,
\]
and $\mathrm{Dt}_{n'}\left(q\,g_{*}A\right)=\mathrm{Dt}_n(qA)$ by (d4). Summation over $\mathfrak{z}$ gives
\[
\boldsymbol{\Psi}^{n'}_{\Phi\circ g^{*}}\left(g_{*}A\right)=\boldsymbol{\Psi}^n_{\Phi}(A)\qquad\left(A\in C^1(n)\right),
\]
both for the full sums and for the sums over flux zero. Integration over the torus, with the invariance of the Haar measure under the isomorphism induced by $g_{*}$, yields
\[
\boldsymbol{\omega}_{n'}\left(\Phi\circ g^{*}\right)=\boldsymbol{\omega}_n(\Phi),\qquad\boldsymbol{\omega}_{n'}^{(0)}\left(\Phi\circ g^{*}\right)=\boldsymbol{\omega}_n^{(0)}(\Phi).
\]

Let $\Phi^{-}(F):=\Phi(-F)$; then $\Phi^{-}\in\mathfrak{F}(n)$. The substitution $A\mapsto-A$, an automorphism of the torus, which preserves its Haar measure, and the substitution $\mathfrak{z}\mapsto-\mathfrak{z}$, a bijection of the closed integer two-cochains which maps the set of those of flux zero onto itself, together with $\mathrm{Dt}(-qA)=\mathrm{Dt}(qA)$ from (d3), give in the same way $\boldsymbol{\Psi}_{\Phi^{-}}(-A)=\boldsymbol{\Psi}_{\Phi}(A)$, for the full sum and for the sum over flux zero. Hence $\boldsymbol{\omega}_n(\Phi^{-})=\boldsymbol{\omega}_n(\Phi)$ and $\boldsymbol{\omega}_n^{(0)}(\Phi^{-})=\boldsymbol{\omega}_n^{(0)}(\Phi)$: the functionals $\boldsymbol{\omega}_n$ and $\boldsymbol{\omega}_n^{(0)}$ are invariant under $F^u\mapsto-F^u$.

On a cubic torus every map $gx=rx+c$ with $r\in O_{\mathrm{cub}}$ and $c\in\mathbb{Z}^3$ satisfies the hypothesis of (d4) with $n'=n$; among these maps are the lattice translations ($r=1$) and the elements of $O_{\mathrm{cub}}$ ($c=0$). As $1\circ g^{*}=1$, the numbers $Z_n=\boldsymbol{\omega}_n(1)$ and $\mathbf{Z}_n=\boldsymbol{\omega}_n^{(0)}(1)$ do not change under these maps, nor under $F^u\mapsto-F^u$. Therefore the quotients $E^{\mathrm{af}}_n=\boldsymbol{\omega}_n/Z_n$ and $E^{\mathrm{av}}_n=\boldsymbol{\omega}_n^{(0)}/\mathbf{Z}_n$ do not change under the lattice translations, the elements of $O_{\mathrm{cub}}$ and the sign change $F^u\mapsto-F^u$.

\emph{Single sectors.} Fix directions $\mu<\nu$, and let $\iota$ be the remaining direction. Let $\mathbf{1}_{\mu\nu}\in C^2(n)$ be the two-cochain with the value $1$ on each plaquette $(x;\mu\nu)$ and the value $0$ on the plaquettes in the two other pairs of directions, and put $\mathbf{1}_{\nu\mu}:=-\mathbf{1}_{\mu\nu}$. Consider $F^u=dA+r_q\cdot\mathfrak{z}$, where $A\in C^1(n)$, $\mathfrak{z}\in C^2(n;\mathbb{Z})$ and $d\mathfrak{z}=0$. The number $\left\langle\mathbf{1}_{\mu\nu},dA\right\rangle$ is the sum of $(dA)_{\mu\nu}(x)$ over $x\in\Lambda_n$, and $(dA)_{\mu\nu}(x)$ is the sum, with signs, of the values of $A$ on the four boundary links of $(x;\mu\nu)$. Of these links, $(x,x+e_{\mu})$ and its translate by $e_{\nu}$ carry opposite signs, and so do $(x,x+e_{\nu})$ and its translate by $e_{\mu}$. The translations by $e_{\mu}$ and by $e_{\nu}$ are bijections of $\Lambda_n$; so the sum over $x$ of the value of $A$ on $(x,x+e_{\mu})$ equals the sum over $x$ of its value on the translate of this link by $e_{\nu}$, and likewise for $(x,x+e_{\nu})$ and the translation by $e_{\mu}$. Hence $\left\langle\mathbf{1}_{\mu\nu},dA\right\rangle=0$. Next, the plaquettes $(x;\mu\nu)$ are distributed over the $n_{\iota}$ coordinate two-tori $\{x_{\iota}=t\}$, $t\in\mathbb{Z}/n_{\iota}\mathbb{Z}$, which are parallel to one another, and by property~\eqref{eq:partition-functions-2} the sum of $\mathfrak{z}$ over the plaquettes of each of them is one and the same number $w_{\mu\nu}(\mathfrak{z})$. Summation over $t$ gives $\left\langle\mathbf{1}_{\mu\nu},\mathfrak{z}\right\rangle=n_{\iota}\cdot w_{\mu\nu}(\mathfrak{z})$, and so
\[
\left\langle\mathbf{1}_{\mu\nu},F^u\right\rangle=r_q\cdot n_{\iota}\cdot w_{\mu\nu}(\mathfrak{z}).
\]
For $w\in\mathbb{Z}^3$ define $\boldsymbol{\chi}_w\colon C^2(n)\to\{0,1\}$ by
\[
\boldsymbol{\chi}_w(F):=\prod_{\mu<\nu}\mathbf{1}\left[\left\langle\mathbf{1}_{\mu\nu},F\right\rangle=r_q\cdot n_{\iota}\cdot w_{\mu\nu}\right],
\]
where $\iota$ is the direction distinct from $\mu$ and $\nu$, and $\mathbf{1}[\cdot]$ is $1$ when the relation in the brackets is true and $0$ when it is false. As the indicator of an intersection of three affine hyperplanes, $\boldsymbol{\chi}_w$ is a bounded Borel function; since moreover $|\Phi\cdot\boldsymbol{\chi}_w|\le|\Phi|$, one has $\Phi\cdot\boldsymbol{\chi}_w\in\mathfrak{F}(n)$. By the identity $\left\langle\mathbf{1}_{\mu\nu},F^u\right\rangle=r_q\cdot n_{\iota}\cdot w_{\mu\nu}(\mathfrak{z})$, obtained above for an arbitrary pair $\mu<\nu$ and $F^u=dA+r_q\cdot\mathfrak{z}$, and since $r_q\cdot n_{\iota}\ne0$, the factor of index $\mu\nu$ in the definition of $\boldsymbol{\chi}_w\left(dA+r_q\cdot\mathfrak{z}\right)$ equals $1$ exactly when $w_{\mu\nu}(\mathfrak{z})=w_{\mu\nu}$. Hence $\boldsymbol{\chi}_w\left(dA+r_q\cdot\mathfrak{z}\right)$ is $1$ if the flux of $\mathfrak{z}$ equals $w$ and $0$ otherwise. Multiplication of $\Phi$ by $\boldsymbol{\chi}_w$ thus restricts the sum in the definition of $\boldsymbol{\omega}_n$ in Section~\ref{sec:wilson-determinant-lattice} to the $\mathfrak{z}$ of flux $w$; that is,
\[
\boldsymbol{\omega}_n^{(w)}(\Phi)=\boldsymbol{\omega}_n\left(\Phi\cdot\boldsymbol{\chi}_w\right).
\]
Let $g$ again be as in (d4), and write $re_{\mu}=\epsilon_{\mu}\cdot e_{\mu'}$ as in (d5). The plaquette $(x;\mu\nu)$ with the orientation $e_{\mu}\wedge e_{\nu}$ is mapped by $g$ onto a plaquette of $\Lambda_{n'}$ spanned by the directions $\mu'$ and $\nu'$, and the orientation that it receives is $re_{\mu}\wedge re_{\nu}=\epsilon_{\mu}\cdot\epsilon_{\nu}\cdot e_{\mu'}\wedge e_{\nu'}$; every plaquette of $\Lambda_{n'}$ in the directions $\mu',\nu'$ is the image of exactly one plaquette $(x;\mu\nu)$. Consequently $g_{*}\mathbf{1}_{\mu\nu}=\epsilon_{\mu}\cdot\epsilon_{\nu}\cdot\mathbf{1}_{\mu'\nu'}$, where $\mathbf{1}_{\mu'\nu'}$ is formed on $\Lambda_{n'}$. Since $g_{*}$ is orthogonal and $g^{*}$ is its inverse, $\left\langle u,g^{*}F'\right\rangle=\left\langle g_{*}u,g_{*}g^{*}F'\right\rangle=\left\langle g_{*}u,F'\right\rangle$ for $u\in C^2(n)$ and $F'\in C^2(n')$; for $u=\mathbf{1}_{\mu\nu}$ this reads
\[
\left\langle\mathbf{1}_{\mu\nu},g^{*}F'\right\rangle=\epsilon_{\mu}\cdot\epsilon_{\nu}\cdot\left\langle\mathbf{1}_{\mu'\nu'},F'\right\rangle\qquad\left(F'\in C^2(n')\right).
\]
Let $\iota'$ be the image of $\iota$ under the permutation $\mu\mapsto\mu'$. It is the direction distinct from $\mu'$ and $\nu'$, and $n'_{\iota'}=n_{\iota}$, because $n'_{\mu'}=n_{\mu}$ for every direction $\mu$. Multiplying each condition by $\epsilon_{\mu}\cdot\epsilon_{\nu}\in\{1,-1\}$ one obtains, for $F'\in C^2(n')$,
\begin{multline*}
\boldsymbol{\chi}_w\left(g^{*}F'\right)=\prod_{\mu<\nu}\mathbf{1}\left[\epsilon_{\mu}\cdot\epsilon_{\nu}\cdot\left\langle\mathbf{1}_{\mu'\nu'},F'\right\rangle=r_q\cdot n'_{\iota'}\cdot w_{\mu\nu}\right]\\
=\prod_{\mu<\nu}\mathbf{1}\left[\left\langle\mathbf{1}_{\mu'\nu'},F'\right\rangle=r_q\cdot n'_{\iota'}\cdot(g\cdot w)_{\mu'\nu'}\right]=\boldsymbol{\chi}_{g\cdot w}(F'),
\end{multline*}
where $\boldsymbol{\chi}_{g\cdot w}$ is the function formed on $\Lambda_{n'}$; the last equality holds because, as $(\mu,\nu)$ ranges over the pairs with $\mu<\nu$, the unordered pair $\{\mu',\nu'\}$ ranges over the three pairs of distinct directions, and both sides of each condition change sign when $\mu'$ and $\nu'$ are exchanged. Thus $\boldsymbol{\chi}_w\circ g^{*}=\boldsymbol{\chi}_{g\cdot w}$, so that $\left(\Phi\cdot\boldsymbol{\chi}_w\right)\circ g^{*}=\left(\Phi\circ g^{*}\right)\cdot\boldsymbol{\chi}_{g\cdot w}$. The sector identity above, formed on $\Lambda_{n'}$ for the function $\Phi\circ g^{*}\in\mathfrak{F}(n')$ and the flux $g\cdot w$, and the first assertion of (d5), applied to $\Phi\cdot\boldsymbol{\chi}_w\in\mathfrak{F}(n)$, now give
\[
\boldsymbol{\omega}_{n'}^{(g\cdot w)}\left(\Phi\circ g^{*}\right)=\boldsymbol{\omega}_{n'}\left(\left(\Phi\cdot\boldsymbol{\chi}_w\right)\circ g^{*}\right)=\boldsymbol{\omega}_n\left(\Phi\cdot\boldsymbol{\chi}_w\right)=\boldsymbol{\omega}_n^{(w)}(\Phi).
\]
By the definition of $g\cdot w$, each component of $g\cdot w$ is, up to its sign, a component of $w$, and distinct components of $g\cdot w$ come from distinct components of $w$, because $\{\mu,\nu\}\mapsto\{\mu',\nu'\}$ is a bijection of the set of the pairs of distinct directions. Hence $w\mapsto g\cdot w$ is a bijection of $\mathbb{Z}^3$, and $g\cdot0=0$. Moreover $(g\cdot w)_{\mu'\nu'}^2=w_{\mu\nu}^2$ and $n'_{\iota'}=n_{\iota}$ for every pair $\mu<\nu$, and $|n'|=|n|$; inserted into the formula for $\mathfrak{q}_w(n)$ in Section~\ref{sec:wilson-determinant-lattice}, these relations give $\mathfrak{q}_{g\cdot w}(n')=\mathfrak{q}_w(n)$.

\emph{The number $b_{\eta,R}(0)$.} Let $g$ be as in (d4) with $n=n'=R/\eta$, the cubic torus that represents $\Lambda_{\eta,R}$ in lattice units. Part (d4), applied with $B=t\eta\,a_{\eta}$ and with the mass $\bar{m}\cdot\eta$ in place of $\boldsymbol{\mu}$, gives
\[
\det\left(D\left[t\eta\,g_{*}a_{\eta}\right]+\bar{m}\cdot\eta\right)=\det\left(D\left[t\eta\,a_{\eta}\right]+\bar{m}\cdot\eta\right)\qquad\left(t\in\mathbb{R}\right).
\]
Hence the function $N_0(t)$ of Section~\ref{sec:free-fourth-order}, and with it the number $b_{\eta,R}(0)$, does not change when $a_{\eta}$ is replaced by $g_{*}a_{\eta}$.

This proves (d1)--(d5), and Theorem~\ref{thm:free-coefficients-lattice}(d) follows.
\end{proof}

\subsubsection{Factorisation over time slices}\label{sec:factorisation-time-slices}

In this subsection one direction is singled out as time; it is given the number $0$, the two other directions are numbered $1$ and $2$, and accordingly $n=(n_0,n_1,n_2)$. The cross-section is the two-dimensional torus
\[
\boldsymbol{\Sigma}:=\mathbb{Z}/n_1\mathbb{Z}\times\mathbb{Z}/n_2\mathbb{Z},
\]
with $v:=n_1n_2$ sites and $l:=2v$ links. For $j\in\mathbb{Z}/n_0\mathbb{Z}$ the slice $j$ is the set of the sites $x$ of $\Lambda_n$ with $x_0=j$. A real one-cochain on $\Lambda_n$ is written $A=(A_j,a_j)_{j\in\mathbb{Z}/n_0\mathbb{Z}}$, where $A_j\in\mathbb{R}^l$ lists the values of $A$ on the links that join two sites of the slice $j$ (the spatial links), and $a_j\in\mathbb{R}^v$ lists its values on the links that lead from a site of the slice $j$ to a site of the slice $j+1$ (the temporal links). This is the notation of Section~\ref{sec:slices-slice-operators} and of the first paragraphs of Sections~\ref{sec:partition-function-trace} and~\ref{sec:trace-formula}.

\begin{lemma}[Slice factorisation of the determinant]\label{lem:slice-factorisation-determinant-2}
Let $q>0$ and $\boldsymbol{\mu}>-1$. There exist a finite-dimensional Hilbert space $\mathcal{F}$, a function $s\colon\mathbb{R}^l\to(0,\infty)$, a map $\mathfrak{Q}$ from $\mathbb{R}^l$ to the linear operators on $\mathcal{F}$, and commuting self-adjoint operators $Q^{\mathrm{ch}}_x$, $x\in\boldsymbol{\Sigma}$, on $\mathcal{F}$ with integer spectra, all depending only on the cross-section $\boldsymbol{\Sigma}$, such that, with
\[
\mathsf{G}(a):=\exp\left(iq\sum_{x\in\boldsymbol{\Sigma}}a_x\cdot Q^{\mathrm{ch}}_x\right)\qquad\left(a\in\mathbb{R}^v\right),
\]
the following two properties hold. First,
\begin{gather}
\text{$s$ and $\mathfrak{Q}$ are bounded and continuous,}\tag{a}\label{eq:free-coefficient-a}\\
\text{and they depend on $A\in\mathbb{R}^l$ only through $\left(\exp(iqA_b)\right)_b$;}\notag\\
\mathsf{G}(a)\ \text{is unitary},\quad\mathsf{G}(a)\cdot\mathsf{G}(a')=\mathsf{G}(a+a'),\quad\mathsf{G}(a)^{*}=\mathsf{G}(-a),\quad\mathsf{G}(r_q\cdot k)=1,\notag
\end{gather}
for all $a,a'\in\mathbb{R}^v$ and $k\in\mathbb{Z}^v$, $b$ ranging over the links of $\boldsymbol{\Sigma}$. Second, for every $n_0\ge2$ and every real one-cochain $A=(A_j,a_j)_{j\in\mathbb{Z}/n_0\mathbb{Z}}$ on $\Lambda_n$, $n=(n_0,n_1,n_2)$,
\begin{equation}\tag{b}\label{eq:free-coefficient-b}
\mathrm{Dt}(qA)=\operatorname{Tr}_{\mathcal{F}}\prod_{j=n_0-1}^0\left[s(A_{j+1})\cdot\mathfrak{Q}(A_{j+1})^{*}\cdot\mathsf{G}(a_j)\cdot\mathfrak{Q}(A_j)\cdot s(A_j)\right],
\end{equation}
where $A_{n_0}:=A_0$, and the factors of the product are ordered with later times to the left: the factor $j=n_0-1$ stands leftmost and the factor $j=0$ rightmost.
\end{lemma}

\begin{proof}
By Section~\ref{sec:wilson-determinant-lattice}, $D_A=D[qA]$ is the operator \eqref{eq:gaussian-bounds-1} for the antiperiodic character $\chi=\chi^{\mathrm{ap}}$, and $\mathrm{Dt}(qA)=\det\left(D_A+\boldsymbol{\mu}\right)$ is the determinant $\Delta^{\chi}(A)$ of Proposition~\ref{prop:trace-formula-determinant} for this character. By Remark~\ref{rem:tori-unequal-sides} the definitions and results of Sections~\ref{sec:slices-slice-operators} and~\ref{sec:second-quantisation-trace} hold on the rectangular torus $\Lambda_n$; their notation is used below, with $\chi=\chi^{\mathrm{ap}}$ and with the direction $0$ as time. For $A\in\mathbb{R}^l$ let $B(A)$ and $C(A)$ be the slice operators \eqref{eq:gaussian-bounds-7} on the space $\mathcal{V}$; by \eqref{eq:gaussian-bounds-9},
\[
1+\boldsymbol{\mu}\le B(A)\le5+\boldsymbol{\mu},\qquad\|C(A)\|\le2.
\]
Let $\mathcal{K}:=\mathcal{V}\oplus\mathcal{V}$, and let $\mathsf{L}(A)$ be the operator on $\mathcal{K}$ given by $\mathsf{L}(A)(u,u'):=\left(u-C(A)u',\,u'\right)$. By \eqref{eq:gaussian-bounds-8},
\[
S(A)=\mathsf{L}(A)^{*}\cdot\operatorname{diag}\left(B(A)^{-1},B(A)\right)\cdot\mathsf{L}(A).
\]
Let $\mathfrak{F}_S$ be the exterior algebra of $\mathcal{K}$ with the inner product of Section~\ref{sec:second-quantisation-trace}, let $\Gamma(\cdot)$ be the second quantisation defined there, and let $Q_x$, $x\in\boldsymbol{\Sigma}$, be the charges \eqref{eq:gaussian-bounds-11}. The $Q_x$ are commuting self-adjoint operators on $\mathfrak{F}_S$, diagonal in one orthonormal basis, with the eigenvalues $\rho-\#|\{\alpha:(x,\alpha)\in I\}|$, where $I$ is a set of indices of the orthonormal basis of $\mathcal{K}$ fixed in Section~\ref{sec:second-quantisation-trace}; these eigenvalues are integers. Put
\[
J(A):=\operatorname{diag}\left(B(A)^{-\frac{1}{2}},B(A)^{\frac{1}{2}}\right)\cdot\mathsf{L}(A).
\]
Then $S(A)=J(A)^{*}J(A)$, and Lemma~\ref{lem:properties-second-quantisation}(a), (b) gives $\Gamma\left(S(A)\right)=\Gamma\left(J(A)\right)^{*}\Gamma\left(J(A)\right)$. Define
\[
\mathcal{F}:=\mathfrak{F}_S,\qquad s(A):=\left(\det B(A)\right)^{\frac{1}{2}},\qquad\mathfrak{Q}(A):=\Gamma\left(J(A)\right)^{*},\qquad Q^{\mathrm{ch}}_x:=Q_x.
\]
These objects are formed on the cross-section alone; in particular they do not depend on $n_0$.

\emph{Property \eqref{eq:free-coefficient-b}.} For the antiperiodic character the angle $\theta$ of Proposition~\ref{prop:trace-formula-determinant} is $0$, and the trace formula \eqref{eq:gaussian-bounds-12} reads
\[
\mathrm{Dt}(qA)=\operatorname{Tr}_{\mathcal{F}}\prod_{j=n_0-1}^0\left[\mathsf{G}(a_j)\cdot s(A_j)^2\cdot\mathfrak{Q}(A_j)\cdot\mathfrak{Q}(A_j)^{*}\right],
\]
because $\exp\left(iq\langle a_j,Q\rangle\right)=\mathsf{G}(a_j)$, $\det B(A_j)=s(A_j)^2$ and $\Gamma\left(S(A_j)\right)=\mathfrak{Q}(A_j)\cdot\mathfrak{Q}(A_j)^{*}$. Move the last operator of this product, $\mathfrak{Q}(A_0)^{*}$, to the front; by the cyclicity of the trace the value does not change, and with $A_{n_0}=A_0$ one obtains
\[
\mathrm{Dt}(qA)=\operatorname{Tr}_{\mathcal{F}}\prod_{j=n_0-1}^0\left[\mathfrak{Q}(A_{j+1})^{*}\cdot\mathsf{G}(a_j)\cdot s(A_j)^2\cdot\mathfrak{Q}(A_j)\right].
\]
In this product every $\mathsf{G}(a_j)$ is immediately preceded by $\mathfrak{Q}(A_{j+1})^{*}$ and, apart from the scalar $s(A_j)^2$, immediately followed by $\mathfrak{Q}(A_j)$. The numbers $s(A_j)$ commute with every operator. Writing $s(A_j)^2=s(A_j)\cdot s(A_j)$ and moving one of these two factors into the bracket of index $j-1$, indices being taken modulo $n_0$, one arrives at \eqref{eq:free-coefficient-b}.

\emph{Property \eqref{eq:free-coefficient-a}.} Every matrix element of $B(A)$ and of $C(A)$ is a polynomial expression in $\exp(iqA_b)$ and $\exp(-iqA_b)$, $b$ ranging over the links of $\boldsymbol{\Sigma}$. The maps $X\mapsto X^{\pm\frac{1}{2}}$ are continuous on the positive Hermitian matrices, in particular on those with spectrum in $[1+\boldsymbol{\mu},5+\boldsymbol{\mu}]$, to which $B(A)$ belongs; and the matrix elements of $\Gamma(J)$ are minors of $J$ (Section~\ref{sec:second-quantisation-trace}). Hence $s$ and $\mathfrak{Q}$ are continuous, and they depend on $A$ only through $\left(\exp(iqA_b)\right)_b$. From $1+\boldsymbol{\mu}\le B(A)\le5+\boldsymbol{\mu}$ one has $0<s(A)\le(5+\boldsymbol{\mu})^{\frac{1}{2}\dim\mathcal{V}}$. By Lemma~\ref{lem:properties-second-quantisation}(e),
\[
\left\|\mathfrak{Q}(A)\right\|=\left\|\Gamma\left(J(A)\right)\right\|\le\max\left(1,\|J(A)\|\right)^{\dim\mathcal{K}},\qquad\|J(A)\|\le3\max\left((1+\boldsymbol{\mu})^{-\frac{1}{2}},(5+\boldsymbol{\mu})^{\frac{1}{2}}\right),
\]
the last bound because $\|\mathsf{L}(A)\|\le1+\|C(A)\|\le3$. So $s$ and $\mathfrak{Q}$ are bounded. Finally, the operators $Q^{\mathrm{ch}}_x=Q_x$ possess a common orthonormal eigenbasis, and all their eigenvalues are integers. Hence, for $a\in\mathbb{R}^v$, the operator $\mathsf{G}(a)$ is diagonal in this basis with eigenvalues of modulus one; so $\mathsf{G}(a)$ is unitary, $\mathsf{G}(a)\cdot\mathsf{G}(a')=\mathsf{G}(a+a')$, $\mathsf{G}(a)^{*}=\mathsf{G}(-a)$, and, as $q\cdot r_q=2\pi$,
\[
\mathsf{G}(r_q\cdot k)=\exp\left(2\pi i\sum_{x\in\boldsymbol{\Sigma}}k_x\cdot Q_x\right)=1\qquad\left(k\in\mathbb{Z}^v\right).
\]
\end{proof}

\subsection{Continuum and infinite-volume limits of the quartic coefficient}\label{sec:continuum-infinite-volume-2}

\subsubsection{The propagator on $\mathbb{R}^3$}\label{sec:propagator}

Let $\mathsf{M}>0$, and let $Y_{\mathsf{M}}$ and the matrix kernel $S_{\mathsf{M}}$ be those of Section~\ref{sec:free-fourth-order}. For $z\in\mathbb{R}^3\setminus\{0\}$ put $\mathsf{r}:=|z|_2$ and $\hat{z}:=z/\mathsf{r}$. The function $Y_{\mathsf{M}}$ depends on $z$ only through $\mathsf{r}$; let $Y_{\mathsf{M}}'(\mathsf{r})$ denote the derivative of the function $r\mapsto\exp(-\mathsf{M}\cdot r)/(4\pi r)$ at $r=\mathsf{r}$. Then
\[
\partial_jY_{\mathsf{M}}(z)=Y_{\mathsf{M}}'(\mathsf{r})\cdot\hat{z}_j\quad(j=1,2,3),\qquad Y_{\mathsf{M}}'(\mathsf{r})=-\frac{\exp\left(-\mathsf{M}\cdot\mathsf{r}\right)\cdot\left(\mathsf{M}\cdot\mathsf{r}+1\right)}{4\pi\,\mathsf{r}^2},
\]
and $|\gamma\cdot\hat{z}|=1$, since the matrix $\gamma\cdot\hat{z}$ is Hermitian with $(\gamma\cdot\hat{z})^2=|\hat{z}|_2^2=1$ (Section~\ref{sec:letters-norms-dirac}). With $\|S_{\mathsf{M}}\|_{L^1}:=\int_{\mathbb{R}^3}\mathrm{d}z\,|S_{\mathsf{M}}(z)|$ it follows that
\begin{gather}
S_{\mathsf{M}}(z)=\mathsf{M}\cdot Y_{\mathsf{M}}(z)-(\gamma\cdot\hat{z})\,Y_{\mathsf{M}}'(\mathsf{r}),\qquad|S_{\mathsf{M}}(z)|\le\frac{\exp\left(-\mathsf{M}\cdot\mathsf{r}\right)\cdot\left(2\mathsf{M}\cdot\mathsf{r}+1\right)}{4\pi\,\mathsf{r}^2},\notag\\
S_{\mathsf{M}}(z)=\mathsf{M}^2\cdot S_1(\mathsf{M}\cdot z),\qquad\|S_{\mathsf{M}}\|_{L^1}\le\int_0^{\infty}\mathrm{d}r\,\left(2\mathsf{M}\cdot r+1\right)\exp\left(-\mathsf{M}\cdot r\right)=\frac{3}{\mathsf{M}}.\tag{K1}\label{eq:free-coefficient-K1}
\end{gather}
Indeed, the first identity is the definition $S_{\mathsf{M}}=(\mathsf{M}-\gamma\cdot\partial)Y_{\mathsf{M}}$ written out; the bound for $|S_{\mathsf{M}}(z)|$ follows from it, because $|\gamma\cdot\hat{z}|=1$ and $\mathsf{M}\cdot Y_{\mathsf{M}}(z)=\mathsf{M}\cdot\mathsf{r}\cdot\exp(-\mathsf{M}\cdot\mathsf{r})/(4\pi\,\mathsf{r}^2)$; the scaling identity follows from the explicit formulas for $Y_{\mathsf{M}}$ and $Y_{\mathsf{M}}'$; and the bound for $\|S_{\mathsf{M}}\|_{L^1}$ is the integral of the bound for $|S_{\mathsf{M}}(z)|$ in polar coordinates.

The function $Y_{\mathsf{M}}$ belongs to $L^1(\mathbb{R}^3)$, and for $p\in\mathbb{R}^3$
\[
\int_{\mathbb{R}^3}\mathrm{d}z\,Y_{\mathsf{M}}(z)\exp(-ip\cdot z)=\frac{1}{\mathsf{M}^2+|p|_2^2}.
\]
The pointwise gradient of $Y_{\mathsf{M}}$ on $\mathbb{R}^3\setminus\{0\}$ lies in $L^1(\mathbb{R}^3;\mathbb{R}^3)$ as well, and it represents the distributional gradient of $Y_{\mathsf{M}}$ on $\mathbb{R}^3$. Consequently the kernel $S_{\mathsf{M}}$ has the Fourier transform
\begin{gather}
\int_{\mathbb{R}^3}\mathrm{d}z\,S_{\mathsf{M}}(z)\exp(-ip\cdot z)=\frac{\mathsf{M}-i\gamma\cdot p}{\mathsf{M}^2+|p|_2^2}=\left(\mathsf{M}+i\gamma\cdot p\right)^{-1}=:S_{\mathsf{M}}(p),\notag\\
|S_{\mathsf{M}}(p)|=\left(\mathsf{M}^2+|p|_2^2\right)^{-\frac{1}{2}},\tag{K2}\label{eq:free-coefficient-K2}
\end{gather}
because $\left(\mathsf{M}+i\gamma\cdot p\right)\left(\mathsf{M}-i\gamma\cdot p\right)=\mathsf{M}^2+|p|_2^2=\left(\mathsf{M}-i\gamma\cdot p\right)^{*}\left(\mathsf{M}-i\gamma\cdot p\right)$ (Section~\ref{sec:letters-norms-dirac}); the argument $p$ distinguishes the multiplier $S_{\mathsf{M}}(p)$ from the kernel $S_{\mathsf{M}}(z)$. In particular the operator norm of $S_{\mathsf{M}}$ on $L^2(\mathbb{R}^3;\mathbb{C}^4)$ is $\sup_p|S_{\mathsf{M}}(p)|=1/\mathsf{M}$.

\begin{lemma}[A Hilbert--Schmidt bound]\label{lem:hilbert-schmidt-bound}
Let $\mathsf{M}>0$, let $a\in C_c^{\infty}(\mathbb{R}^3;\mathbb{R}^3)$ be a real vector field, and put
\[
\|a\|_{\infty}:=\sup_{x\in\mathbb{R}^3}|a(x)|_2,\qquad\|a\|_{L^2}:=\left(\int_{\mathbb{R}^3}\mathrm{d}x\,|a(x)|_2^2\right)^{\frac{1}{2}}.
\]
Let $V:=i\gamma\cdot a$, acting on $L^2(\mathbb{R}^3;\mathbb{C}^4)$ by multiplication, and $G:=S_{\mathsf{M}}VS_{\mathsf{M}}$. The Hilbert--Schmidt norm of $G$ is finite, and
\[
\|G\|_{\mathrm{HS}}^2\le\frac{\|a\|_{L^2}^2}{\pi\cdot\mathsf{M}}.
\]
The fourth power $(S_{\mathsf{M}}V)^4$ coincides with $GVGV$ and belongs to the trace class, and
\[
\left|q_{4,\mathsf{M}}(a)\right|\le\frac{\|a\|_{L^2}^2\cdot\|a\|_{\infty}^2}{4\pi\cdot\mathsf{M}}.
\]
Let moreover $R>0$, let $L^2_{\mathrm{ap}}(\mathbb{T}_R;\mathbb{C}^4)$ and $S^R=\left(\bar{m}+\gamma\cdot\partial\right)^{-1}$ be as in Section~\ref{sec:free-fourth-order}, let
\[
a^R(x):=\sum_{z\in\mathbb{Z}^3}a(x+Rz)\qquad(x\in\mathbb{R}^3)
\]
be the $R$-periodisation of $a$, and let $V^R$ denote the multiplication by $i\gamma\cdot a^R$ on $L^2_{\mathrm{ap}}(\mathbb{T}_R;\mathbb{C}^4)$. Then $G^R:=S^RV^RS^R$ is a Hilbert--Schmidt operator on $L^2_{\mathrm{ap}}(\mathbb{T}_R;\mathbb{C}^4)$, and
\[
\|G^R\|_{\mathrm{HS}}^2\le8R^{-3}\cdot\sum_{p\in\frac{2\pi}{R}\left(\mathbb{Z}+\frac{1}{2}\right)^3}\left(\bar{m}^2+|p|_2^2\right)^{-2}\cdot\int_{[0,R)^3}\mathrm{d}x\,|a^R(x)|_2^2\le\frac{32R}{\pi^2}\cdot\int_{[0,R)^3}\mathrm{d}x\,|a^R(x)|_2^2;
\]
the fourth power $(S^RV^R)^4$ coincides with $G^RV^RG^RV^R$ and belongs to the trace class. Finally, $q_{4,\mathsf{M}}(a)$ is a real number, and $b_R(f)$ is a real number for all $f$ and $R$ for which Section~\ref{sec:free-fourth-order} defines it.
\end{lemma}

\begin{proof}
For $\psi\in L^2(\mathbb{R}^3;\mathbb{C}^4)$ let $\hat\psi(q):=\int\mathrm{d}x\,\psi(x)\exp(-iq\cdot x)$. By Plancherel's theorem the map $\psi\mapsto\hat\psi$ is unitary from $L^2(\mathbb{R}^3;\mathbb{C}^4)$ onto $L^2\left(\mathbb{R}^3,(2\pi)^{-3}\mathrm{d}q;\mathbb{C}^4\right)$. Let $\hat a(k):=\int\mathrm{d}x\,a(x)\exp(-ik\cdot x)$ for $k\in\mathbb{R}^3$. In this momentum representation $S_{\mathsf{M}}$ is the multiplication by $S_{\mathsf{M}}(p)$, by~\eqref{eq:free-coefficient-K2}, and $V$ is the integral operator with the kernel $i\gamma\cdot\hat a(p-q)$ with respect to the measure $(2\pi)^{-3}\mathrm{d}q$. Hence $G$ is the integral operator with the kernel
\[
S_{\mathsf{M}}(p)\,i\gamma\cdot\hat a(p-q)\,S_{\mathsf{M}}(q)
\]
with respect to $(2\pi)^{-3}\mathrm{d}q$.

For $c\in\mathbb{C}^3$ one has $|\gamma\cdot c|\le|\operatorname{Re}c|_2+|\operatorname{Im}c|_2$, hence $|\gamma\cdot c|^2\le2|c|_2^2$; and every $4\times4$ matrix has Hilbert--Schmidt norm at most twice its operator norm. Let $u(p):=|S_{\mathsf{M}}(p)|^2$. The squared Hilbert--Schmidt norm of an integral operator is the integral of the squared Hilbert--Schmidt norm of its kernel; with the two facts just stated this gives
\begin{align*}
\|G\|_{\mathrm{HS}}^2&=(2\pi)^{-6}\int\mathrm{d}p\int\mathrm{d}q\,\left\|S_{\mathsf{M}}(p)\,i\gamma\cdot\hat a(p-q)\,S_{\mathsf{M}}(q)\right\|_{\mathrm{HS}}^2\le8\cdot(2\pi)^{-6}\int\mathrm{d}p\,u(p)\cdot\left(|\hat a|_2^2*u\right)(p)\\
&\le8\cdot(2\pi)^{-6}\cdot\|u\|_{L^2}^2\cdot\left\||\hat a|_2^2\right\|_{L^1}=8\cdot(2\pi)^{-3}\cdot\|u\|_{L^2}^2\cdot\|a\|_{L^2}^2.
\end{align*}
Here the second inequality is the Cauchy--Schwarz inequality together with Young's inequality $\|g*u\|_{L^2}\le\|g\|_{L^1}\cdot\|u\|_{L^2}$, and the last equality is Plancherel's theorem. By~\eqref{eq:free-coefficient-K2},
\[
\|u\|_{L^2}^2=\int_{\mathbb{R}^3}\mathrm{d}p\,\left(\mathsf{M}^2+|p|_2^2\right)^{-2}=\frac{\pi^2}{\mathsf{M}},
\]
and since $8\cdot(2\pi)^{-3}\cdot\pi^2=1/\pi$, the bound for $\|G\|_{\mathrm{HS}}^2$ follows.

The multiplication operator $V$ is bounded, with $\|V\|=\|a\|_{\infty}$, because $\gamma\cdot a(x)$ is Hermitian and $(\gamma\cdot a(x))^2=|a(x)|_2^2$ (Section~\ref{sec:letters-norms-dirac}). So $VGV$ is a Hilbert--Schmidt operator with $\|VGV\|_{\mathrm{HS}}\le\|V\|^2\cdot\|G\|_{\mathrm{HS}}$. If $X$ and $Y$ are Hilbert--Schmidt operators on one Hilbert space, the trace norm of $XY$ does not exceed $\|X\|_{\mathrm{HS}}\cdot\|Y\|_{\mathrm{HS}}$, by H\"older's inequality for Schatten norms; hence $XY$ lies in the trace class, and $|\operatorname{Tr}(XY)|\le\|X\|_{\mathrm{HS}}\cdot\|Y\|_{\mathrm{HS}}$. For $X:=G$ and $Y:=VGV$ this shows that $(S_{\mathsf{M}}V)^4=G\cdot VGV$ lies in the trace class and that
\[
\left|\operatorname{Tr}\left[(S_{\mathsf{M}}V)^4\right]\right|\le\|G\|_{\mathrm{HS}}^2\cdot\|V\|^2\le\frac{\|a\|_{L^2}^2\cdot\|a\|_{\infty}^2}{\pi\cdot\mathsf{M}}.
\]
Since $q_{4,\mathsf{M}}(a)=-\frac{1}{4}\operatorname{Tr}\left[(S_{\mathsf{M}}V)^4\right]$ (Section~\ref{sec:free-fourth-order}), the bound for $|q_{4,\mathsf{M}}(a)|$ follows.

Let now $R>0$, and let $a^R$, $V^R$ and $G^R$ be as in the statement. Put
\[
P_R:=\frac{2\pi}{R}\left(\mathbb{Z}+\frac{1}{2}\right)^3,\qquad K_R:=\frac{2\pi}{R}\,\mathbb{Z}^3,\qquad e_p(x):=R^{-\frac{3}{2}}\exp(ip\cdot x)\quad(p\in P_R,\ x\in\mathbb{R}^3),
\]
and let $\epsilon$ and $\epsilon'$ range over the standard basis of $\mathbb{C}^4$. As $p_{\mu}R\in2\pi\left(\mathbb{Z}+\frac{1}{2}\right)$, the function $e_p$ is antiperiodic with the periods $Re_{\mu}$. The function $x\mapsto\exp\left(-i\pi\left(x_1+x_2+x_3\right)/R\right)$ has modulus $1$ and changes its sign under each translation $x\mapsto x+Re_{\mu}$; the multiplication by it is therefore a unitary map of $L^2_{\mathrm{ap}}(\mathbb{T}_R;\mathbb{C}^4)$ onto the space of the $R\mathbb{Z}^3$-periodic functions in $L^2_{\mathrm{loc}}(\mathbb{R}^3;\mathbb{C}^4)$, normed by $L^2\left([0,R)^3\right)$, which the restriction to $[0,R)^3$ identifies with $L^2\left([0,R)^3;\mathbb{C}^4\right)$. This map sends $e_p\otimes\epsilon$ to the function $x\mapsto R^{-\frac{3}{2}}\exp(ik\cdot x)\,\epsilon$ with $k:=p-\frac{\pi}{R}(1,1,1)$, and $p\mapsto p-\frac{\pi}{R}(1,1,1)$ is a bijection of $P_R$ onto $K_R$. By the completeness of the trigonometric system, the functions $x\mapsto R^{-\frac{3}{2}}\exp(ik\cdot x)$, $k\in K_R$, form an orthonormal basis of $L^2\left([0,R)^3\right)$, so that the functions $x\mapsto R^{-\frac{3}{2}}\exp(ik\cdot x)\,\epsilon$ form an orthonormal basis of $L^2\left([0,R)^3;\mathbb{C}^4\right)$; consequently the functions $e_p\otimes\epsilon$, $p\in P_R$, form an orthonormal basis of $L^2_{\mathrm{ap}}(\mathbb{T}_R;\mathbb{C}^4)$. For an operator $X$ on this space and $p,q\in P_R$, let $\langle e_p,Xe_q\rangle$ be the $4\times4$ matrix whose entry in the row of $\epsilon$ and the column of $\epsilon'$ is $\left\langle e_p\otimes\epsilon,X\left(e_q\otimes\epsilon'\right)\right\rangle$. By the definition of $S^R$ in Section~\ref{sec:free-fourth-order}, and hence also for its adjoint,
\[
S^R\left(e_p\otimes\epsilon\right)=e_p\otimes S_{\bar{m}}(p)\,\epsilon,\qquad\left(S^R\right)^{*}\left(e_p\otimes\epsilon\right)=e_p\otimes S_{\bar{m}}(p)^{*}\,\epsilon\qquad(p\in P_R),
\]
where $S_{\bar{m}}(p)=\left(\bar{m}+i\gamma\cdot p\right)^{-1}$ is the multiplier of~\eqref{eq:free-coefficient-K2} at $\mathsf{M}=\bar{m}$. By~\eqref{eq:free-coefficient-K2},
\[
\left|S_{\bar{m}}(p)\right|^2=w(p):=\left(\bar{m}^2+|p|_2^2\right)^{-1}\le\bar{m}^{-2}\qquad(p\in P_R);
\]
in particular $S^R$ is a bounded operator.

The function $a^R$ is smooth and $R\mathbb{Z}^3$-periodic, since on each bounded set only finitely many terms of the sum defining it differ from zero; in particular $a^R$ is bounded, and $V^R$ is a bounded operator on $L^2_{\mathrm{ap}}(\mathbb{T}_R;\mathbb{C}^4)$. Let $k\in K_R$. The function $x\mapsto\exp(-ik\cdot x)$ is $R\mathbb{Z}^3$-periodic, and the cubes $[0,R)^3+Rz$, $z\in\mathbb{Z}^3$, form a partition of $\mathbb{R}^3$; since only finitely many of these cubes meet the support of $a$,
\[
\int_{[0,R)^3}\mathrm{d}x\,\exp(-ik\cdot x)\,a^R(x)=\sum_{z\in\mathbb{Z}^3}\int_{[0,R)^3+Rz}\mathrm{d}y\,\exp(-ik\cdot y)\,a(y)=\hat a(k).
\]
For $p,q\in P_R$ the difference $p-q$ lies in $K_R$. Hence
\[
\left\langle e_p,V^Re_q\right\rangle=R^{-3}\int_{[0,R)^3}\mathrm{d}x\,\exp\left(-i(p-q)\cdot x\right)\,i\gamma\cdot a^R(x)=R^{-3}\,i\gamma\cdot\hat a(p-q),
\]
and, since $\left\langle e_p\otimes\epsilon,G^R\left(e_q\otimes\epsilon'\right)\right\rangle=\left\langle\left(S^R\right)^{*}\left(e_p\otimes\epsilon\right),V^RS^R\left(e_q\otimes\epsilon'\right)\right\rangle$,
\[
\left\langle e_p,G^Re_q\right\rangle=R^{-3}\,S_{\bar{m}}(p)\,i\gamma\cdot\hat a(p-q)\,S_{\bar{m}}(q)\qquad(p,q\in P_R).
\]

The squared Hilbert--Schmidt norm of $G^R$ is the sum of the squared moduli of its matrix elements in the orthonormal basis $\left(e_p\otimes\epsilon\right)$. With the two facts on $4\times4$ matrices used above for $G$ it follows that
\[
\|G^R\|_{\mathrm{HS}}^2=\sum_{p,q\in P_R}\left\|\left\langle e_p,G^Re_q\right\rangle\right\|_{\mathrm{HS}}^2\le8R^{-6}\sum_{p,q\in P_R}w(p)\cdot w(q)\cdot|\hat a(p-q)|_2^2.
\]
One has $w(p)\cdot w(q)\le\frac{1}{2}\left(w(p)^2+w(q)^2\right)$. For fixed $p\in P_R$ the map $q\mapsto p-q$ is a bijection of $P_R$ onto $K_R$, and for fixed $q\in P_R$ so is the map $p\mapsto p-q$; therefore
\[
\sum_{p,q\in P_R}w(p)\cdot w(q)\cdot|\hat a(p-q)|_2^2\le\sum_{p\in P_R}w(p)^2\cdot\sum_{k\in K_R}|\hat a(k)|_2^2.
\]
By the identity for $\int_{[0,R)^3}\mathrm{d}x\,\exp(-ik\cdot x)\,a^R(x)$ obtained above, the coefficients of the components of $a^R$ with respect to the orthonormal basis $\left(x\mapsto R^{-\frac{3}{2}}\exp(ik\cdot x)\right)_{k\in K_R}$ of $L^2\left([0,R)^3\right)$ are the components of the vectors $R^{-\frac{3}{2}}\hat a(k)$. Parseval's identity, applied to each component, therefore gives
\[
\sum_{k\in K_R}|\hat a(k)|_2^2=R^3\int_{[0,R)^3}\mathrm{d}x\,|a^R(x)|_2^2.
\]
Since $w(p)^2=\left(\bar{m}^2+|p|_2^2\right)^{-2}$, the last three displays give the first inequality of the statement.

It remains to bound the sum of the numbers $w(p)^2$. Let $|p|$ be the maximum norm of $p$, and let $i\ge1$ be an integer. A point $p=\frac{2\pi}{R}\left(n+\frac{1}{2}(1,1,1)\right)$ of $P_R$, $n\in\mathbb{Z}^3$, satisfies $|p|\le\frac{2\pi}{R}\left(i-\frac{1}{2}\right)$ exactly when each of the three numbers $n_j+\frac{1}{2}$ is one of the $2i$ numbers $-i+\frac{1}{2},\dots,i-\frac{1}{2}$. Hence exactly $(2i)^3-(2i-2)^3=24i^2-24i+8\le24i^2$ points of $P_R$ satisfy $|p|=\frac{2\pi}{R}\left(i-\frac{1}{2}\right)$, and each point of $P_R$ satisfies this equation for exactly one integer $i\ge1$. Since $\frac{2\pi}{R}\left(i-\frac{1}{2}\right)\ge\frac{\pi i}{R}$, since $w(p)^2\le|p|_2^{-4}\le|p|^{-4}$, and since $\sum_{i\ge1}i^{-2}=\pi^2/6$,
\[
\sum_{p\in P_R}w(p)^2\le\sum_{i\ge1}24i^2\left(\frac{R}{\pi i}\right)^4=\frac{24R^4}{\pi^4}\sum_{i\ge1}i^{-2}=\frac{4R^4}{\pi^2}.
\]
As $8R^{-3}\cdot4R^4/\pi^2=32R/\pi^2$, the second inequality of the statement follows. Its right-hand side is finite, because $a^R$ is bounded; hence $G^R$ is a Hilbert--Schmidt operator. As $V^R$ is bounded, $V^RG^RV^R$ is a Hilbert--Schmidt operator as well, and $(S^RV^R)^4=G^R\cdot V^RG^RV^R$ is a product of two Hilbert--Schmidt operators; by the fact on such products recalled above, it lies in the trace class.

The operator $\gamma\cdot\partial$ is anti-self-adjoint, and the matrix $\hat{\Gamma}$ is Hermitian with $\hat{\Gamma}^2=1$ and $\hat{\Gamma}\gamma_{\mu}\hat{\Gamma}=-\gamma_{\mu}$ (Section~\ref{sec:letters-norms-dirac}). Hence
\[
\hat{\Gamma} S_{\mathsf{M}}\hat{\Gamma}=\left(\mathsf{M}-\gamma\cdot\partial\right)^{-1}=S_{\mathsf{M}}^{*},\qquad\hat{\Gamma} V\hat{\Gamma}=-V=V^{*}.
\]
The trace is invariant under conjugation by the unitary operator $\hat{\Gamma}$; therefore
\begin{align*}
\operatorname{Tr}\left[(S_{\mathsf{M}}V)^4\right]&=\operatorname{Tr}\left(\hat{\Gamma}(S_{\mathsf{M}}V)^4\hat{\Gamma}\right)=\operatorname{Tr}\left[\left(S_{\mathsf{M}}^{*}V^{*}\right)^4\right]=\operatorname{Tr}\left(\left((VS_{\mathsf{M}})^4\right)^{*}\right)\\
&=\overline{\operatorname{Tr}(VGVG)}=\overline{\operatorname{Tr}(GVGV)}=\overline{\operatorname{Tr}\left[(S_{\mathsf{M}}V)^4\right]};
\end{align*}
in this chain $(VS_{\mathsf{M}})^4$ has been written as $V\cdot GVG$, and $\operatorname{Tr}(V\cdot GVG)$ has been replaced by $\operatorname{Tr}(GVG\cdot V)$. The latter step is legitimate: $GVG=G\cdot(VG)$ is a product of two Hilbert--Schmidt operators, hence of trace class, and the trace of the product of a trace class operator and a bounded operator does not depend on the order of the two factors. Consequently $\operatorname{Tr}\left[(S_{\mathsf{M}}V)^4\right]$ equals its own complex conjugate; it is real, and so is $q_{4,\mathsf{M}}(a)=-\frac{1}{4}\operatorname{Tr}\left[(S_{\mathsf{M}}V)^4\right]$.

Now let $f$ and $R$ be such that $b_R(f)$ is defined (Section~\ref{sec:free-fourth-order}). There the field is $a=d^{*}\beta$ for a probe $\beta$ with $f=da$, its support lies in the interior of a period cube, and $a$ stands also for its $R$-periodisation; this is the function $a^R$ of the statement, because for $x$ in that period cube and $z\in\mathbb{Z}^3\setminus\{0\}$ the point $x+Rz$ lies outside the interior of the cube, so that $a^R(x)=a(x)$. Let $V$ now denote the multiplication by $i\gamma\cdot a$ on $L^2_{\mathrm{ap}}(\mathbb{T}_R;\mathbb{C}^4)$, that is, $V=V^R$. On this space the operator $\gamma\cdot\partial$ acts on the momenta $p\in\frac{2\pi}{R}\left(\mathbb{Z}+\frac{1}{2}\right)^3$ as the multiplier $i\gamma\cdot p$; so it is again anti-self-adjoint, and $\hat{\Gamma}$ maps antiperiodic functions to antiperiodic functions. Hence $\hat{\Gamma} S^R\hat{\Gamma}=\left(\bar{m}-\gamma\cdot\partial\right)^{-1}=(S^R)^{*}$ and $\hat{\Gamma} V\hat{\Gamma}=V^{*}$. As shown above, the operator $G^R=S^RVS^R$ on $L^2_{\mathrm{ap}}(\mathbb{T}_R;\mathbb{C}^4)$ is Hilbert--Schmidt; consequently $G^RVG^R$ and $(S^RV)^4=G^RVG^RV$ are of trace class, and the chain of identities above, with $S^R$ and $G^R$ in place of $S_{\mathsf{M}}$ and $G$, shows that $\operatorname{Tr}\left[(S^RV)^4\right]$ is real. Therefore $b_R(f)=-6\operatorname{Tr}\left[(S^RV)^4\right]$ is real.
\end{proof}

\subsubsection{Estimates in momentum space}\label{sec:estimates-momentum-space}

Throughout Sections~\ref{sec:estimates-momentum-space} and~\ref{sec:lattice-torus}, $m:=\bar{m}\in(0,1]$, $R\ge1$, and $\eta\in(0,1]$ is such that $n:=R/\eta$ is an integer $\ge2$; for $k\in\mathbb{R}^3$, $|k|$ is the maximum norm and $|k|_2$ the Euclidean norm (Section~\ref{sec:letters-norms-dirac}). Let
\[
P_R:=\frac{2\pi}{R}\left(\mathbb{Z}+\frac{1}{2}\right)^3,\qquad K_R:=\frac{2\pi}{R}\,\mathbb{Z}^3,\qquad B_{\eta}:=P_R\cap\left(-\frac{\pi}{\eta},\frac{\pi}{\eta}\right]^3,\qquad K_{\eta}:=K_R\cap\left(-\frac{\pi}{\eta},\frac{\pi}{\eta}\right]^3,
\]
and
\[
\rho_{\eta}(k):=\min_{l\in\mathbb{Z}^3}\left|k+\frac{2\pi l}{\eta}\right|\qquad(k\in\mathbb{R}^3).
\]
One has $\rho_{\eta}(k)=|k|$ for $|k|\le\pi/\eta$ (if $|k_j|\le\pi/\eta$ and $l_j\ne0$, then $|k_j+2\pi l_j/\eta|\ge\pi/\eta\ge|k_j|$), and $\rho_{\eta}(k+k')\le\rho_{\eta}(k)+\rho_{\eta}(k')$ for all $k,k'\in\mathbb{R}^3$. Since $(2\pi/\eta)\,\mathbb{Z}^3=(2\pi/R)\,n\mathbb{Z}^3$, the sets $B_{\eta}$ and $K_{\eta}$ are systems of representatives of $P_R$ and of $K_R$ modulo $(2\pi/\eta)\,\mathbb{Z}^3$, and $\#|B_{\eta}|=\#|K_{\eta}|=n^3$. Finally, for each $p\in P_R$ the assignment $k\mapsto p-k$ maps $K_R$ bijectively onto $P_R$, and two points of $K_R$ are congruent modulo $(2\pi/\eta)\,\mathbb{Z}^3$ exactly when their images are; it therefore defines a bijection from the set of classes $K_R/(2\pi/\eta)\,\mathbb{Z}^3$ to the set of classes $P_R/(2\pi/\eta)\,\mathbb{Z}^3$.

\begin{lemma}[Lattice sums]\label{lem:lattice-sums}
\begin{enumerate}
\item[(a)] One has
\[
R^{-3}\sum_{k\in K_R}(1+|k|)^{-5}\le28,\qquad R^{-3}\sum_{k\in K_R,\,|k|>\boldsymbol{\Lambda}}(1+|k|)^{-6}\le\frac{28}{\boldsymbol{\Lambda}}\quad\text{for every }\boldsymbol{\Lambda}>0.
\]
\item[(b)] One has
\[
R^{-3}\sum_{p\in B_{\eta}}(m+|p|)^{-1}\le8\eta^{-2},\qquad R^{-3}\sum_{p\in B_{\eta}}(m+|p|)^{-2}\le3\eta^{-1},\qquad R^{-3}\sum_{p\in B_{\eta}}(m+|p|)^{-3}\le24m^{-3}+2+\log\frac{1}{\eta}.
\]
\item[(c)] One has
\[
R^{-3}\sum_{p\in P_R}(m+|p|)^{-\frac{7}{2}}\le24m^{-\frac{7}{2}}+2,\qquad R^{-3}\sum_{p\in P_R}(m+|p|)^{-4}\le24m^{-4}+1,
\]
and
\[
R^{-3}\sum_{p\in P_R,\,|p|>\boldsymbol{\Lambda}}(m+|p|)^{-4}\le\frac{7}{\boldsymbol{\Lambda}}\quad\text{for every }\boldsymbol{\Lambda}>0.
\]
\end{enumerate}
\end{lemma}

\begin{proof}
For an integer $i\ge1$,
\begin{gather*}
\#\left|\left\{k\in K_R:|k|=\frac{2\pi i}{R}\right\}\right|=(2i+1)^3-(2i-1)^3=24i^2+2\le26i^2,\\
\#\left|\left\{p\in P_R:|p|=\frac{2\pi}{R}\left(i-\frac{1}{2}\right)\right\}\right|=(2i)^3-(2i-2)^3\le24i^2,
\end{gather*}
and $\frac{2\pi}{R}\left(i-\frac{1}{2}\right)\ge\frac{\pi i}{R}$. Every $k\in K_R\setminus\{0\}$ and every $p\in P_R$ lies in exactly one of these sets; the corresponding $i$ is called the index of $k$, resp.\ of $p$. In the sums below, $i$ ranges over the integers $\ge1$. For real $y\ge1$ and $t>1$,
\[
\sum_{i\le y}i^2\le y^3,\qquad\lfloor y\rfloor\ge\frac{y}{2},\qquad\sum_{i>y}i^{-t}\le\int_{\lfloor y\rfloor}^{\infty}\mathrm{d}x\,x^{-t}\le\frac{1}{t-1}\left(\frac{y}{2}\right)^{1-t},
\]
and $\sum_{i>y}i^{-2}\le4/y$ holds for every real $y>0$: for $y\ge1$ the last inequality with $t=2$ even gives the bound $2/y$, while for $0<y<1$ the sum extends over all $i\ge1$ and equals $\pi^2/6<4<4/y$.

(a) Split the sum into the term $k=0$ and the terms of index $i\le R$ and $i>R$; for $i\le R$ use $(1+|k|)^{-5}\le1$, and for $i>R$ use $(1+2\pi i/R)^{-5}\le\left(\frac{R}{2\pi}\right)^5i^{-5}$. Then
\[
R^{-3}\sum_{k\in K_R}(1+|k|)^{-5}\le R^{-3}\left[1+26\sum_{i\le R}i^2+26\left(\frac{R}{2\pi}\right)^5\sum_{i>R}i^{-3}\right]\le1+26+26\cdot(2\pi)^{-5}\cdot2\le28.
\]
On the set $\{|k|>\boldsymbol{\Lambda}\}$ one has $(1+|k|)^{-6}\le\boldsymbol{\Lambda}^{-1}\cdot(1+|k|)^{-5}$, and the second bound of (a) follows from the first.

(b) Let $p\in B_{\eta}$ have the index $i$. Then $|p|\le\pi/\eta=\pi n/R$, hence $i\le(n+1)/2\le n$; moreover $(m+|p|)^{-t}\le|p|^{-t}\le\left(\frac{R}{\pi i}\right)^t$ for $t>0$. For $t=1,2$ this gives
\[
R^{-3}\sum_{p\in B_{\eta}}(m+|p|)^{-t}\le R^{-3}\sum_{i\le n}24i^2\left(\frac{R}{\pi i}\right)^t=24\pi^{-t}R^{t-3}\sum_{i\le n}i^{2-t}\le24\pi^{-t}R^{t-3}n^{3-t}=24\pi^{-t}\eta^{t-3},
\]
and $24/\pi\le8$, $24/\pi^2\le3$. Let now $t=3$. At most $\sum_{i\le R}24i^2\le24R^3$ elements of $B_{\eta}$ have an index $i\le R$, and each of them satisfies $(m+|p|)^{-3}\le m^{-3}$; after multiplication by $R^{-3}$ these elements account for at most $R^{-3}\cdot24R^3\cdot m^{-3}$. The elements of index $R<i\le n$ account for at most
\[
R^{-3}\sum_{R<i\le n}24i^2\left(\frac{R}{\pi i}\right)^3=24\pi^{-3}\sum_{R<i\le n}i^{-1}\le24\pi^{-3}\log\frac{n}{\lfloor R\rfloor}\le24\pi^{-3}\left(\log2+\log\frac{1}{\eta}\right)\le2+\log\frac{1}{\eta},
\]
by comparison with $\int_{\lfloor R\rfloor}^n\mathrm{d}x\,x^{-1}$ and because $\lfloor R\rfloor\ge R/2$ and $n/R=1/\eta$.

(c) Let $t\in\left\{\frac{7}{2},4\right\}$. As in (b), the elements of $P_R$ of index $i\le R$ account for at most $24m^{-t}$. For those of index $i>R$ one uses $(m+|p|)^{-t}\le\left(\frac{R}{\pi i}\right)^t$ together with the bound for $\sum_{i>y}i^{-(t-2)}$ at $y=R$, which applies because $t-2>1$; their share is at most
\[
R^{-3}\sum_{i>R}24i^2\left(\frac{R}{\pi i}\right)^t=24\pi^{-t}R^{t-3}\sum_{i>R}i^{2-t}\le24\pi^{-t}R^{t-3}\cdot\frac{1}{t-3}\left(\frac{R}{2}\right)^{3-t}=\frac{24\pi^{-t}\cdot2^{t-3}}{t-3},
\]
a quantity bounded by $1.24$ for $t=\frac{7}{2}$ and by $0.5$ for $t=4$. Finally, let $p\in P_R$ have the index $i$ and satisfy $|p|>\boldsymbol{\Lambda}$. Then $2\pi i/R>|p|>\boldsymbol{\Lambda}$, so that $i>\boldsymbol{\Lambda}\cdot R/(2\pi)$. With $(m+|p|)^{-4}\le\left(\frac{R}{\pi i}\right)^4$ and the bound for $\sum_{i>y}i^{-2}$ at $y=\boldsymbol{\Lambda}\cdot R/(2\pi)$ it follows that
\[
R^{-3}\sum_{p\in P_R,\,|p|>\boldsymbol{\Lambda}}(m+|p|)^{-4}\le24\pi^{-4}R\sum_{i>\boldsymbol{\Lambda}\cdot R/(2\pi)}i^{-2}\le24\pi^{-4}R\cdot4\cdot\frac{2\pi}{\boldsymbol{\Lambda}\cdot R}=\frac{192\pi^{-3}}{\boldsymbol{\Lambda}}\le\frac{7}{\boldsymbol{\Lambda}}.
\]
\end{proof}

\begin{lemma}[Sampling and aliasing]\label{lem:sampling-aliasing}
Let $D\le R-1$, let $\boldsymbol{\phi}\in C_c^{\infty}(\mathbb{R}^3)$ vanish outside a closed cube $\mathcal{Q}\subset\mathbb{R}^3$ whose edges have length $D$, and let $\hat{\boldsymbol{\phi}}(k):=\int\mathrm{d}y\,\boldsymbol{\phi}(y)\exp(-ik\cdot y)$ for $k\in\mathbb{R}^3$. Let $\nu\in\{1,2,3\}$, let $K\ge5$ be an integer, let $\boldsymbol{\phi}^R(y):=\sum_{l\in\mathbb{Z}^3}\boldsymbol{\phi}(y+Rl)$ be the periodisation of $\boldsymbol{\phi}$, and for $k\in K_R$ let
\[
\boldsymbol{\phi}^{\eta}(k):=\eta^3\sum_{x\in\eta\mathbb{Z}^3/R\mathbb{Z}^3}\boldsymbol{\phi}^R\left(x+\frac{\eta}{2}e_{\nu}\right)\exp\left(-ik\cdot\left(x+\frac{\eta}{2}e_{\nu}\right)\right);
\]
the summand depends on $x$ only modulo $R\mathbb{Z}^3$, because $\boldsymbol{\phi}^R$ is $R\mathbb{Z}^3$-periodic and $k\in K_R$. Let $|\boldsymbol{\phi}|_{C^K}$ be the largest of the sup norms of the derivatives of $\boldsymbol{\phi}$ of order at most $K$, and $c_K(\boldsymbol{\phi}):=2^KD^3\cdot|\boldsymbol{\phi}|_{C^K}$. Then:
\begin{enumerate}
\item[(a)] $|\hat{\boldsymbol{\phi}}(k)|\le c_K(\boldsymbol{\phi})\cdot(1+|k|)^{-K}$ for all $k\in\mathbb{R}^3$;
\item[(b)] for all $k\in K_R$ and $j\in\{1,2,3\}$,
\[
\boldsymbol{\phi}^{\eta}(k)=\sum_{l\in\mathbb{Z}^3}(-1)^{l_{\nu}}\,\hat{\boldsymbol{\phi}}\left(k+\frac{2\pi l}{\eta}\right),\qquad\boldsymbol{\phi}^{\eta}\left(k+\frac{2\pi}{\eta}e_j\right)=(-1)^{\delta_{j\nu}}\,\boldsymbol{\phi}^{\eta}(k);
\]
\item[(c)] for all $k\in K_R$ with $|k|\le\pi/\eta$,
\[
\left|\boldsymbol{\phi}^{\eta}(k)-\hat{\boldsymbol{\phi}}(k)\right|\le\eta^4\,c_K(\boldsymbol{\phi})\cdot(1+|k|)^{-(K-4)};
\]
\item[(d)] for all $k\in K_R$,
\[
\left|\boldsymbol{\phi}^{\eta}(k)\right|\le2c_K(\boldsymbol{\phi})\cdot\left(1+\rho_{\eta}(k)\right)^{-(K-4)}.
\]
\end{enumerate}
\end{lemma}

\begin{proof}
(a) Integration by parts gives $(ik_j)^K\,\hat{\boldsymbol{\phi}}(k)=\int\mathrm{d}y\,\partial_j^K\boldsymbol{\phi}(y)\exp(-ik\cdot y)$; since the volume of $\mathcal{Q}$ is $D^3$, it follows that $|k_j|^K\cdot|\hat{\boldsymbol{\phi}}(k)|\le D^3\cdot|\partial_j^K\boldsymbol{\phi}|_{\infty}$ for $j=1,2,3$, and also $|\hat{\boldsymbol{\phi}}(k)|\le D^3\cdot|\boldsymbol{\phi}|_{\infty}$. Choose $j$ with $|k_j|=|k|$; then $\max\left(1,|k|^K\right)\cdot|\hat{\boldsymbol{\phi}}(k)|\le D^3\cdot|\boldsymbol{\phi}|_{C^K}$, and (a) follows from $(1+y)^K\le2^K\max\left(1,y^K\right)$ for $y\ge0$.

(b) The function $\boldsymbol{\phi}^R$ is smooth and $R\mathbb{Z}^3$-periodic, and its Fourier coefficients are
\[
\int_{[0,R)^3}\mathrm{d}y\,\boldsymbol{\phi}^R(y)\exp(-ik'\cdot y)=\hat{\boldsymbol{\phi}}(k')\qquad(k'\in K_R).
\]
By the theorem on the Fourier series of a smooth periodic function,
\[
\boldsymbol{\phi}^R(y)=R^{-3}\sum_{k'\in K_R}\hat{\boldsymbol{\phi}}(k')\exp(ik'\cdot y)\qquad(y\in\mathbb{R}^3),
\]
and the series converges absolutely by (a), because $K\ge5>3$. For $k,k'\in K_R$ the function $x\mapsto\exp\left(i(k'-k)\cdot x\right)$ is a character of the finite group $\eta\mathbb{Z}^3/R\mathbb{Z}^3$ of order $n^3$, trivial exactly when $k'-k\in(2\pi/\eta)\,\mathbb{Z}^3$; by the orthogonality of the characters,
\[
\eta^3\sum_{x\in\eta\mathbb{Z}^3/R\mathbb{Z}^3}\exp\left(i(k'-k)\cdot x\right)=\begin{cases}R^3&\text{if }k'-k\in(2\pi/\eta)\,\mathbb{Z}^3,\\0&\text{otherwise.}\end{cases}
\]
Insert the series for $\boldsymbol{\phi}^R$ into the definition of $\boldsymbol{\phi}^{\eta}(k)$ and sum over $x$ first. Only the terms with $k'=k+2\pi l/\eta$, $l\in\mathbb{Z}^3$, remain, and each of them carries the factor $\exp\left(i(k'-k)_{\nu}\,\eta/2\right)=\exp(i\pi l_{\nu})=(-1)^{l_{\nu}}$. This is the first formula of (b). The substitution $l\mapsto l+e_j$ in it gives $\boldsymbol{\phi}^{\eta}\left(k+\frac{2\pi}{\eta}e_j\right)=(-1)^{\delta_{j\nu}}\,\boldsymbol{\phi}^{\eta}(k)$.

(c) Let $l\in\mathbb{Z}^3\setminus\{0\}$ and $|k|\le\pi/\eta$. Then
\[
\left|k+\frac{2\pi l}{\eta}\right|\ge\left(2|l|-1\right)\frac{\pi}{\eta}\ge\max\left(|k|,\frac{|l|\,\pi}{\eta}\right),
\]
and therefore
\[
\left(1+\left|k+\frac{2\pi l}{\eta}\right|\right)^{-K}\le(1+|k|)^{-(K-4)}\left(\frac{\eta}{\pi|l|}\right)^4.
\]
By the count in the proof of Lemma~\ref{lem:lattice-sums}, $\sum_{l\ne0}|l|^{-4}=\sum_{i\ge1}(24i^2+2)\,i^{-4}<42<\pi^4$. By (b) and (a),
\[
\left|\boldsymbol{\phi}^{\eta}(k)-\hat{\boldsymbol{\phi}}(k)\right|\le\sum_{l\ne0}\left|\hat{\boldsymbol{\phi}}\left(k+\frac{2\pi l}{\eta}\right)\right|\le c_K(\boldsymbol{\phi})\cdot(1+|k|)^{-(K-4)}\cdot\eta^4\pi^{-4}\sum_{l\ne0}|l|^{-4}\le\eta^4\,c_K(\boldsymbol{\phi})\cdot(1+|k|)^{-(K-4)}.
\]

(d) By (b), $|\boldsymbol{\phi}^{\eta}|$ is $(2\pi/\eta)\,\mathbb{Z}^3$-periodic on $K_R$, and so is $\rho_{\eta}$. Let $k\in K_R$, and let $k_0\in K_{\eta}$ with $k-k_0\in(2\pi/\eta)\,\mathbb{Z}^3$. Then $|k_0|\le\pi/\eta$ and $\rho_{\eta}(k)=\rho_{\eta}(k_0)=|k_0|$. By (a), (c) and $\eta\le1$,
\[
\left|\boldsymbol{\phi}^{\eta}(k)\right|=\left|\boldsymbol{\phi}^{\eta}(k_0)\right|\le|\hat{\boldsymbol{\phi}}(k_0)|+\left|\boldsymbol{\phi}^{\eta}(k_0)-\hat{\boldsymbol{\phi}}(k_0)\right|\le2c_K(\boldsymbol{\phi})\cdot(1+|k_0|)^{-(K-4)}=2c_K(\boldsymbol{\phi})\cdot\left(1+\rho_{\eta}(k)\right)^{-(K-4)}.
\]
\end{proof}

\begin{lemma}[Lattice and continuum multipliers]\label{lem:lattice-continuum-multipliers}
For $p\in\mathbb{R}^3$ let
\[
\mathrm{Wl}(p):=\eta^{-1}\sum_{j=1}^3\left(1-\cos(\eta p_j)\right),\qquad\operatorname{sn}_j(p):=\eta^{-1}\sin(\eta p_j)\quad(j=1,2,3),
\]
$\operatorname{sn}(p):=\left(\operatorname{sn}_1(p),\operatorname{sn}_2(p),\operatorname{sn}_3(p)\right)$, and
\[
S_{\eta}(p):=\left(m+\mathrm{Wl}(p)+i\gamma\cdot\operatorname{sn}(p)\right)^{-1},\qquad S(p):=S_m(p)=(m+i\gamma\cdot p)^{-1},
\]
where $S_m(p)$ is the multiplier of~\eqref{eq:free-coefficient-K2} at $\mathsf{M}=m$. The function $S_{\eta}$ is invariant under translation by every element of $(2\pi/\eta)\,\mathbb{Z}^3$. For every $p\in\mathbb{R}^3$,
\[
|S_{\eta}(p)|\le\frac{2\pi}{m+\rho_{\eta}(p)},\qquad|S(p)|\le\frac{2^{\frac{1}{2}}}{m+|p|},
\]
and if moreover $|p|\le\pi/\eta$, then
\[
|S_{\eta}(p)-S(p)|\le15\,\eta^{\frac{1}{2}}\,(m+|p|)^{-\frac{1}{2}}.
\]
\end{lemma}

\begin{proof}
Let $p\in\mathbb{R}^3$. The matrix $A:=m+\mathrm{Wl}(p)+i\gamma\cdot\operatorname{sn}(p)$ has the real scalar part $m+\mathrm{Wl}(p)\ge m>0$ and the real vector $\operatorname{sn}(p)$; by the computation of~\eqref{eq:free-coefficient-K2}, $A^{*}A=\left(m+\mathrm{Wl}(p)\right)^2+|\operatorname{sn}(p)|_2^2$. Hence $A$ is invertible, and
\[
|S_{\eta}(p)|=\left(\left(m+\mathrm{Wl}(p)\right)^2+|\operatorname{sn}(p)|_2^2\right)^{-\frac{1}{2}}.
\]
The functions $\mathrm{Wl}$ and $\operatorname{sn}$ are $(2\pi/\eta)\,\mathbb{Z}^3$-periodic, and so is $S_{\eta}$.

Let $|p|\le\pi/\eta$, and choose $j$ with $|p_j|=|p|$. If $\eta|p_j|\le\pi/2$, then $|\operatorname{sn}_j(p)|\ge(2/\pi)\,|p_j|$, because $\sin x\ge2x/\pi$ for $0\le x\le\pi/2$. Otherwise $\pi/2<\eta|p_j|\le\pi$, so that $1-\cos(\eta p_j)\ge1$ and $\mathrm{Wl}(p)\ge1/\eta\ge|p_j|/\pi$. In both cases
\[
\left(\left(m+\mathrm{Wl}(p)\right)^2+|\operatorname{sn}(p)|_2^2\right)^{\frac{1}{2}}\ge\max\left(m,\frac{|p|}{\pi}\right)\ge\frac{m+|p|}{2\pi},
\]
that is, $|S_{\eta}(p)|\le2\pi/(m+|p|)=2\pi/(m+\rho_{\eta}(p))$. For arbitrary $p\in\mathbb{R}^3$ choose $l\in\mathbb{Z}^3$ with $|p+2\pi l/\eta|=\rho_{\eta}(p)$ and put $p_0:=p+2\pi l/\eta$. Then $|p_0|\le\pi/\eta$, because every coordinate of $p$ can be moved into $[-\pi/\eta,\pi/\eta]$ by an element of $(2\pi/\eta)\,\mathbb{Z}$; by the periodicity of $S_{\eta}$ and the case just treated, $|S_{\eta}(p)|=|S_{\eta}(p_0)|\le2\pi/(m+|p_0|)=2\pi/(m+\rho_{\eta}(p))$. By~\eqref{eq:free-coefficient-K2} at $\mathsf{M}=m$, together with $\left(m^2+|p|_2^2\right)^{\frac{1}{2}}\ge2^{-\frac{1}{2}}(m+|p|_2)\ge2^{-\frac{1}{2}}(m+|p|)$, one has $|S(p)|\le2^{\frac{1}{2}}/(m+|p|)$.

Let again $|p|\le\pi/\eta$. One has
\[
S_{\eta}(p)-S(p)=S_{\eta}(p)\left(S(p)^{-1}-S_{\eta}(p)^{-1}\right)S(p),\qquad S_{\eta}(p)^{-1}-S(p)^{-1}=\mathrm{Wl}(p)+i\gamma\cdot\left(\operatorname{sn}(p)-p\right).
\]
Since $0\le1-\cos x\le x^2/2$ and $|\sin x-x|\le|x|^3/6$ for real $x$, and since $\eta|p_j|\le\pi$,
\[
0\le\mathrm{Wl}(p)\le\frac{\eta\,|p|_2^2}{2},\qquad\left|\operatorname{sn}_j(p)-p_j\right|\le\frac{\eta^2|p_j|^3}{6}\le\frac{\pi}{6}\,\eta\,p_j^2\quad(j=1,2,3).
\]
With $|\gamma\cdot v|=|v|_2\le\sum_j|v_j|$ for $v\in\mathbb{R}^3$ and $|p|_2^2\le3|p|^2$ this gives
\[
\left|S_{\eta}(p)^{-1}-S(p)^{-1}\right|\le\left(\frac{1}{2}+\frac{\pi}{6}\right)\eta\,|p|_2^2\le3.1\,\eta\,|p|^2,
\]
and, with the two bounds proved above,
\[
|S_{\eta}(p)-S(p)|\le\frac{2\pi}{m+|p|}\cdot3.1\,\eta\,|p|^2\cdot\frac{2^{\frac{1}{2}}}{m+|p|}\le2\pi\cdot3.1\cdot2^{\frac{1}{2}}\cdot\eta\le28\eta.
\]
Also $|S_{\eta}(p)-S(p)|\le|S_{\eta}(p)|+|S(p)|\le\left(2\pi+2^{\frac{1}{2}}\right)/(m+|p|)\le8/(m+|p|)$. Since $\min(x,y)\le(xy)^{\frac{1}{2}}$ for $x,y\ge0$, it follows that
\[
|S_{\eta}(p)-S(p)|\le\left(28\eta\cdot\frac{8}{m+|p|}\right)^{\frac{1}{2}}=224^{\frac{1}{2}}\,\eta^{\frac{1}{2}}\,(m+|p|)^{-\frac{1}{2}}\le15\,\eta^{\frac{1}{2}}\,(m+|p|)^{-\frac{1}{2}}.
\]
\end{proof}

\paragraph{Peetre's inequality.}
Let $p,q\in\mathbb{R}^3$. The function $\rho_{\eta}$ is even and subadditive, and $m\le1$; hence
\[
m+\rho_{\eta}(q)\le m+\rho_{\eta}(p)+\rho_{\eta}(p-q)\le\frac{\left(m+\rho_{\eta}(p)\right)\left(1+\rho_{\eta}(p-q)\right)}{m}.
\]
For $t\ge0$, raising this inequality to the power $t$ gives the first of the two inequalities below; for $t<0$, the first inequality below follows by raising the same inequality with $p$ and $q$ exchanged (note $\rho_{\eta}(q-p)=\rho_{\eta}(p-q)$) to the power $|t|$ and dividing. Thus, for every real $t$,
\begin{gather}
\left(m+\rho_{\eta}(q)\right)^t\le m^{-|t|}\left(m+\rho_{\eta}(p)\right)^t\left(1+\rho_{\eta}(p-q)\right)^{|t|},\notag\\
(m+|q|)^t\le m^{-|t|}(m+|p|)^t(1+|p-q|)^{|t|};\tag{Pe}\label{eq:free-coefficient-Pe}
\end{gather}
the second inequality follows by the same argument, since the norm $|\cdot|$ is even and subadditive.

\subsubsection{From the lattice to the torus}\label{sec:lattice-torus}

In this subsection $\beta$, $a$, $f$, $a_{\eta}$ and $R$ are as in Theorem~\ref{thm:free-coefficient-1}(c): $0<\eta\le1$, $R/\eta$ is an integer, $R\ge D_{\beta}+1$, and the cube $Q_{\beta}$ lies in the interior of a period cube of the torus of side $R$. Since $D_{\beta}=2\rho_{\beta}$ and $\rho_{\beta}\ge2$ (Section~\ref{sec:probes-discretisation}), one has $R\ge D_{\beta}+1\ge5$, and the integer $n=R/\eta$ satisfies $n\ge R\ge5$, as $\eta\le1$. Hence $m=\bar{m}\in(0,1]$, $R\ge1$ and $n\ge2$, as required at the beginning of Section~\ref{sec:estimates-momentum-space}, and the number $b_{\eta,R}(0)$ of Section~\ref{sec:free-fourth-order} is defined; the sets $P_R$, $K_R$, $B_{\eta}$, $K_{\eta}$ and the function $\rho_{\eta}$ are those defined in Section~\ref{sec:estimates-momentum-space}. On $\mathbb{R}^3$, $|\cdot|$ is the maximum norm, and on $4\times4$ matrices it is the operator norm (Section~\ref{sec:letters-norms-dirac}).

\begin{proof}[Proof of Theorem~\ref{thm:free-coefficient-1}(c1) on the torus and of Theorem~\ref{thm:free-coefficient-1}(c2)]
\emph{Constants.} Put $D:=D_{\beta}$ and $\mathcal{Q}:=Q_{\beta}$. The closed cube $\mathcal{Q}$ has side $D\le R-1$, and Lemma~\ref{lem:sampling-aliasing} is applied with this cube and with $K:=16$. For $\nu\in\{1,2,3\}$ and a site $x$ of $\Lambda_{\eta,R}$ let $\boldsymbol{\alpha}_{\nu}(x):=a_{\eta,\nu}(x)$ be the value of the one-cochain $a_{\eta}$ on the bond from $x$ to $x+\eta e_{\nu}$. By Theorem~\ref{thm:free-coefficient-1}(b1),
\[
\boldsymbol{\alpha}_{\nu}(x)=\check{a}_{\eta,\nu}\left(x+\frac{\eta}{2}e_{\nu}\right),
\]
that is, $\boldsymbol{\alpha}_{\nu}$ is the sample of $\check{a}_{\eta,\nu}$ at the centres of the bonds of direction $\nu$. The symbol $a_{\nu}$ keeps its meaning of Section~\ref{sec:probes-discretisation}: it is the $\nu$-th component of the field $a=d^{*}\beta$ on $\mathbb{R}^3$.

By Theorem~\ref{thm:free-coefficient-1}(b2) the powers $\check{a}_{\eta,\nu}^r$ of the functions $\check{a}_{\eta,\nu}$, $r\in\{1,2,3,4\}$, are smooth and supported in $\mathcal{Q}$; by Theorem~\ref{thm:free-coefficient-1}(b4) and Leibniz' rule, their derivatives of order at most $16$ are bounded uniformly in $\eta\in(0,1]$. With the numbers $c_{16}(\cdot)$ of Lemma~\ref{lem:sampling-aliasing} it follows that
\[
A:=\sup\left\{c_{16}\left(\check{a}_{\eta,\nu}^r\right):\ 0<\eta\le1,\ \nu\in\{1,2,3\},\ r\in\{1,2,3,4\}\right\}
\]
is finite. The function $\check{a}_{\eta,\nu}-a_{\nu}$ is supported in $\mathcal{Q}$, and $\left|\check{a}_{\eta,\nu}-a_{\nu}\right|_{C^{16}}\le3\eta^2\cdot|\beta|_{C^{19}}$ by Theorem~\ref{thm:free-coefficient-1}(b4). Lemma~\ref{lem:sampling-aliasing}(a), applied to this function, gives
\[
\left|\hat{\check{a}}_{\eta,\nu}(k)-\hat{a}_{\nu}(k)\right|\le A'\cdot\eta^2\cdot(1+|k|)^{-16}\quad(k\in\mathbb{R}^3),\qquad A':=3\cdot2^{16}\cdot D^3\cdot|\beta|_{C^{19}},
\]
the Fourier transform being normalised as in Lemma~\ref{lem:sampling-aliasing}. In this proof $C$ denotes constants that depend only on $m$, $A$ and $A'$; different occurrences of $C$ may denote different constants.

\emph{Expansion in the field.} A lattice spinor is a map $\psi\colon\eta\mathbb{Z}^3\to\mathbb{C}^4$ with $\psi(x+Re_{\mu})=-\psi(x)$ for $\mu=1,2,3$. The lattice spinors form a Hilbert space of dimension $4n^3$, with the scalar product
\[
\langle\psi,\psi'\rangle:=\eta^3\sum_x\psi(x)^{*}\,\psi'(x),
\]
the sum extending over the points $x$ of a period cell; in lattice units they are the spinors of Section~\ref{sec:wilson-determinant-lattice} on $\Lambda_n$, on which the operators $D[B]$ of~\eqref{eq:free-coefficient-5} act. An $R\mathbb{Z}^3$-periodic function on $\eta\mathbb{Z}^3$, that is, a function on $\Lambda_{\eta,R}$, acts on the lattice spinors by multiplication, since its product with a lattice spinor is again antiperiodic; a $4\times4$ matrix acts on their values. For $\nu\in\{1,2,3\}$ let $\boldsymbol{\tau}_{\nu}$ be the unitary operator $\left(\boldsymbol{\tau}_{\nu}\psi\right)(x):=\psi\left(x+\eta e_{\nu}\right)$, so that $\left(\boldsymbol{\tau}_{\nu}^{*}\psi\right)(x)=\psi\left(x-\eta e_{\nu}\right)$, and put
\[
\mathbf{P}^{\pm}_{\nu}:=\frac{1}{2}\left(1\pm\gamma_{\nu}\right),\qquad\check{D}(t):=\eta^{-1}D[t\eta a_{\eta}]\quad(t\in\mathbb{C}).
\]
In~\eqref{eq:free-coefficient-5} the hop from $x$ to $x+\eta e_{\nu}$ carries the matrix $\mathbf{P}^{-}_{\nu}$ and the phase $\exp\left(it\eta\cdot\boldsymbol{\alpha}_{\nu}(x)\right)$ of the bond from $x$ to $x+\eta e_{\nu}$; the hop from $x$ to $x-\eta e_{\nu}$ carries the matrix $\mathbf{P}^{+}_{\nu}$ and the phase $\exp\left(-it\eta\cdot\boldsymbol{\alpha}_{\nu}(x-\eta e_{\nu})\right)$, the inverse of the phase of the bond from $x-\eta e_{\nu}$ to $x$. Therefore
\[
\check{D}(t)=\eta^{-1}\left(3-\sum_{\nu}\left[\mathbf{P}^{-}_{\nu}\exp\left(it\eta\cdot\boldsymbol{\alpha}_{\nu}\right)\boldsymbol{\tau}_{\nu}+\mathbf{P}^{+}_{\nu}\,\boldsymbol{\tau}_{\nu}^{*}\exp\left(-it\eta\cdot\boldsymbol{\alpha}_{\nu}\right)\right]\right),
\]
and the power series of the two exponentials give
\[
\check{D}(t)=\check{D}(0)+\sum_{r\ge1}t^r\,V_r,\qquad V_r:=-\frac{\eta^{r-1}i^r}{r!}\sum_{\nu}\left[\mathbf{P}^{-}_{\nu}\,\boldsymbol{\alpha}_{\nu}^r\,\boldsymbol{\tau}_{\nu}+(-1)^r\,\mathbf{P}^{+}_{\nu}\,\boldsymbol{\tau}_{\nu}^{*}\,\boldsymbol{\alpha}_{\nu}^r\right],
\]
where the power $\boldsymbol{\alpha}_{\nu}^r$ acts by multiplication.

Since $D[t\eta a_{\eta}]+\bar{m}\cdot\eta=\eta\cdot\left(\check{D}(t)+m\right)$, the ratio $N_0$ of Section~\ref{sec:free-fourth-order} is
\[
N_0(t)=\frac{\det\left(\check{D}(t)+m\right)}{\det\left(\check{D}(0)+m\right)}.
\]
By Theorem~\ref{thm:free-coefficient-1}(d2), $\operatorname{Re}\left(D[0]+m\eta\right)\ge m\eta>0$; hence $\check{D}(0)+m$ is invertible. Put
\[
S_{\eta}:=\left(\check{D}(0)+m\right)^{-1},\qquad T_r:=S_{\eta}V_r\quad(r\ge1),\qquad Y(t):=\sum_{r\ge1}t^r\,T_r=S_{\eta}\left(\check{D}(t)-\check{D}(0)\right).
\]
Then $N_0(t)=\det\left(1+Y(t)\right)$, and the norm of $Y(t)$ tends to $0$ as $t\to0$. For $|t|$ small the expansion of the logarithm of a determinant gives
\[
\log N_0(t)=\sum_{j\ge1}\frac{(-1)^{j+1}}{j}\operatorname{Tr}\left(Y(t)^j\right),
\]
the logarithm being the branch that vanishes at $t=0$; and $b_{\eta,R}(0)$, the fourth derivative of $\log N_0$ at $t=0$, is $24$ times the coefficient of $t^4$ in this series. The coefficient of $t^4$ in $\operatorname{Tr}\left(Y(t)^j\right)$ is the sum of $\operatorname{Tr}\left(T_{r_1}\cdots T_{r_j}\right)$ over all $r_1,\dots,r_j\ge1$ with $r_1+\dots+r_j=4$; by the cyclicity of the trace it equals $\operatorname{Tr}T_4$ for $j=1$, $2\operatorname{Tr}\left(T_1T_3\right)+\operatorname{Tr}\left(T_2^2\right)$ for $j=2$, $3\operatorname{Tr}\left(T_1^2T_2\right)$ for $j=3$, $\operatorname{Tr}\left(T_1^4\right)$ for $j=4$, and $0$ for $j\ge5$. It follows that
\begin{equation}\label{eq:free-coefficient-11}
b_{\eta,R}(0)=24\left[\operatorname{Tr}T_4-\operatorname{Tr}\left(T_1T_3\right)-\frac{1}{2}\operatorname{Tr}\left(T_2^2\right)+\operatorname{Tr}\left(T_1^2T_2\right)\right]-6\operatorname{Tr}\left(T_1^4\right).
\end{equation}

\emph{Momentum representation.} For $p\in P_R$ and $x\in\mathbb{R}^3$ put $e_p(x):=R^{-\frac{3}{2}}\exp(ip\cdot x)$. As $p_{\mu}R\in2\pi\left(\mathbb{Z}+\frac{1}{2}\right)$, the function $e_p$ is antiperiodic with the periods $Re_{\mu}$; on the lattice, $e_p$ stands for its restriction to $\eta\mathbb{Z}^3$. Let $\epsilon$ and $\epsilon'$ range over the standard basis of $\mathbb{C}^4$. For $p,q\in P_R$ the function $x\mapsto\exp\left(i(q-p)\cdot x\right)$ is a character of the group $\eta\mathbb{Z}^3/R\mathbb{Z}^3$ of order $n^3$, trivial exactly when $p-q\in\frac{2\pi}{\eta}\mathbb{Z}^3$; by the orthogonality of characters,
\[
\eta^3\sum_x\exp\left(i(q-p)\cdot x\right)=\begin{cases}R^3&\text{if }p-q\in\frac{2\pi}{\eta}\mathbb{Z}^3,\\0&\text{otherwise,}\end{cases}
\]
the sum extending over a period cell. Since $B_{\eta}$ is a system of $n^3$ representatives of $P_R$ modulo $\frac{2\pi}{\eta}\mathbb{Z}^3$, the spinors $e_p\otimes\epsilon$, $p\in B_{\eta}$, form an orthonormal basis of the lattice spinors. For an operator $X$ on the lattice spinors and $p,q\in P_R$, let $\langle e_p,Xe_q\rangle$ be the $4\times4$ matrix whose entry in the row of $\epsilon$ and the column of $\epsilon'$ is $\left\langle e_p\otimes\epsilon,X\left(e_q\otimes\epsilon'\right)\right\rangle$; for a $4\times4$ matrix $M$, the equation $Xe_p=Me_p$ means that $X\left(e_p\otimes\epsilon\right)=e_p\otimes M\epsilon$ for every $\epsilon$.

For real $y$ one has $\mathbf{P}^{-}_{\nu}\exp(iy)+\mathbf{P}^{+}_{\nu}\exp(-iy)=\cos y-i\gamma_{\nu}\sin y$. Applied with $y=\eta p_{\nu}$ to the formula for $\check{D}(0)$, this gives
\[
\check{D}(0)\,e_p=\left(\mathrm{Wl}(p)+i\gamma\cdot\operatorname{sn}(p)\right)e_p\qquad(p\in P_R),
\]
with $\mathrm{Wl}$ and $\operatorname{sn}$ as in Lemma~\ref{lem:lattice-continuum-multipliers}. Hence $S_{\eta}e_p=S_{\eta}(p)\,e_p$, where $S_{\eta}(p)$ is the matrix of Lemma~\ref{lem:lattice-continuum-multipliers}.

Let $\nu\in\{1,2,3\}$, let $\boldsymbol{\phi}\in C_c^{\infty}(\mathbb{R}^3)$ be supported in $\mathcal{Q}$, let $\boldsymbol{\phi}^R$ be its periodisation and $\boldsymbol{\phi}^{\eta}\colon K_R\to\mathbb{C}$ the function of Lemma~\ref{lem:sampling-aliasing} formed with this $\nu$, and let $g$ be the function on $\eta\mathbb{Z}^3$ given by $g(x):=\boldsymbol{\phi}^R\left(x+\frac{\eta}{2}e_{\nu}\right)$, the sample of $\boldsymbol{\phi}^R$ at the centres of the bonds of direction $\nu$; it is $R\mathbb{Z}^3$-periodic, as $\boldsymbol{\phi}^R$ is. For $p,q\in P_R$ and $x\in\eta\mathbb{Z}^3$,
\begin{align*}
\exp\left(-i(p-q)\cdot x\right)\cdot\exp\left(i\eta q_{\nu}\right)&=\exp\left(-i(p-q)\cdot\left(x+\frac{\eta}{2}e_{\nu}\right)\right)\cdot\exp\left(\frac{i\eta}{2}(p+q)_{\nu}\right),\\
\exp\left(-i(p-q)\cdot x\right)\cdot\exp\left(-i\eta p_{\nu}\right)&=\exp\left(-i(p-q)\cdot\left(x+\frac{\eta}{2}e_{\nu}\right)\right)\cdot\exp\left(-\frac{i\eta}{2}(p+q)_{\nu}\right).
\end{align*}
Inserting the first identity into $\left\langle e_p,g\,\boldsymbol{\tau}_{\nu}e_q\right\rangle$ and the second into $\left\langle e_p,\boldsymbol{\tau}_{\nu}^{*}g\,e_q\right\rangle=\left\langle\boldsymbol{\tau}_{\nu}e_p,g\,e_q\right\rangle$, and comparing with the definition of $\boldsymbol{\phi}^{\eta}$, one obtains
\[
\left\langle e_p,g\,\boldsymbol{\tau}_{\nu}e_q\right\rangle=R^{-3}\exp\left(\frac{i\eta}{2}(p+q)_{\nu}\right)\boldsymbol{\phi}^{\eta}(p-q),\qquad\left\langle e_p,\boldsymbol{\tau}_{\nu}^{*}g\,e_q\right\rangle=R^{-3}\exp\left(-\frac{i\eta}{2}(p+q)_{\nu}\right)\boldsymbol{\phi}^{\eta}(p-q),
\]
both matrices being multiples of the unit matrix.

By Theorem~\ref{thm:free-coefficient-1}(b1), and since $\mathcal{Q}$ lies in the interior of a period cube, the function $\boldsymbol{\alpha}_{\nu}^r$ is the function $g$ formed with $\boldsymbol{\phi}:=\check{a}_{\eta,\nu}^r$. For this $\boldsymbol{\phi}$ the function $\boldsymbol{\phi}^{\eta}$ is denoted $\left(\check{a}_{\eta,\nu}^r\right)^{\eta}$, and $\check{a}_{\eta,\nu}^{\eta}$ when $r=1$. The formula for $V_r$ now gives $\left\langle e_p,V_re_q\right\rangle=R^{-3}\,v_r(p,q)$ for $p,q\in P_R$, where
\[
v_r(p,q):=-\frac{\eta^{r-1}i^r}{r!}\sum_{\nu}\left[\mathbf{P}^{-}_{\nu}\exp\left(iy_{\nu}\right)+(-1)^r\,\mathbf{P}^{+}_{\nu}\exp\left(-iy_{\nu}\right)\right]\left(\check{a}_{\eta,\nu}^r\right)^{\eta}(p-q),\qquad y_{\nu}:=\frac{\eta}{2}(p+q)_{\nu}.
\]
For $r=1$, since $\mathbf{P}^{-}_{\nu}\exp(iy)-\mathbf{P}^{+}_{\nu}\exp(-iy)=i\sin y-\gamma_{\nu}\cos y$,
\[
v_1(p,q)=\sum_{\nu}\left[i\gamma_{\nu}\cos y_{\nu}+\sin y_{\nu}\right]\check{a}_{\eta,\nu}^{\eta}(p-q).
\]
The right-hand sides make sense for all $p,q\in P_R$, and $v_r$ is regarded as a function on $P_R\times P_R$. The translation of $p$, or of $q$, by $\frac{2\pi}{\eta}e_j$ changes $y_j$ by $\pi$ and leaves the $y_{\nu}$ with $\nu\ne j$ unchanged; so it multiplies $\exp\left(\pm iy_j\right)$ by $-1$, while by Lemma~\ref{lem:sampling-aliasing}(b) it multiplies $\left(\check{a}_{\eta,j}^r\right)^{\eta}(p-q)$ by $-1$ and leaves $\left(\check{a}_{\eta,\nu}^r\right)^{\eta}(p-q)$, $\nu\ne j$, unchanged. Hence $v_r$ is $\frac{2\pi}{\eta}\mathbb{Z}^3$-periodic in $p$ and in $q$. Let now $r\in\{1,2,3,4\}$; then $c_{16}\left(\check{a}_{\eta,\nu}^r\right)\le A$ for $\nu=1,2,3$, by the definition of $A$, and Lemma~\ref{lem:sampling-aliasing}(d), together with $\left|\mathbf{P}^{\pm}_{\nu}\right|=1$, gives
\[
\left|v_r(p,q)\right|\le\frac{12A\cdot\eta^{r-1}}{r!}\cdot\left(1+\rho_{\eta}(p-q)\right)^{-12}\qquad(p,q\in P_R,\ r\in\{1,2,3,4\}).
\]

In the basis $e_p\otimes\epsilon$, $p\in B_{\eta}$, the operator $T_r=S_{\eta}V_r$ has the matrix elements $\left\langle e_p,T_re_q\right\rangle=R^{-3}\,S_{\eta}(p)\,v_r(p,q)$. Let $j\ge1$ and $r_1,\dots,r_j\ge1$. The trace of $T_{r_1}\cdots T_{r_j}$ is the sum over $p_1,\dots,p_j\in B_{\eta}$ of $\operatorname{tr}\left[\left\langle e_{p_1},T_{r_1}e_{p_2}\right\rangle\cdots\left\langle e_{p_j},T_{r_j}e_{p_1}\right\rangle\right]$. This summand is $\frac{2\pi}{\eta}\mathbb{Z}^3$-periodic in each $p_i$, by Lemma~\ref{lem:lattice-continuum-multipliers} and the periodicity of the $v_r$; and for fixed $p_i\in P_R$ the map $k\mapsto p_i-k$ induces a bijection of $K_R/\left(\frac{2\pi}{\eta}\mathbb{Z}^3\right)$ onto $P_R/\left(\frac{2\pi}{\eta}\mathbb{Z}^3\right)$ (Section~\ref{sec:estimates-momentum-space}). So, successively for $i=1,\dots,j-1$, the summation over $p_{i+1}\in B_{\eta}$ may be replaced by the summation over $k_i\in K_{\eta}$, with
\[
p_{i+1}:=p_i-k_i\qquad(1\le i\le j-1).
\]
This gives
\begin{multline}\label{eq:free-coefficient-12}
\operatorname{Tr}\left(T_{r_1}\cdots T_{r_j}\right)=R^{-3j}\sum_{p_1\in B_{\eta}}\ \sum_{k_1,\dots,k_{j-1}\in K_{\eta}}\\
\operatorname{tr}\left[S_{\eta}(p_1)\,v_{r_1}(p_1,p_2)\,S_{\eta}(p_2)\,v_{r_2}(p_2,p_3)\cdots S_{\eta}(p_j)\,v_{r_j}(p_j,p_1)\right].
\end{multline}

\emph{The terms in brackets.} Let $j\le4$, let $r_1,\dots,r_j\in\{1,2,3,4\}$, and let $p_1\in B_{\eta}$ and $k_1,\dots,k_{j-1}\in K_{\eta}$. Then $\rho_{\eta}(p_1)=|p_1|$, $\rho_{\eta}(k_l)=|k_l|$, and $\rho_{\eta}(p_1-p_i)\le|k_1|+\dots+|k_{i-1}|$. Lemma~\ref{lem:lattice-continuum-multipliers}, Peetre's inequality~\eqref{eq:free-coefficient-Pe} with $t=-1$, and the inequality $1+\sum_ly_l\le\prod_l(1+y_l)$ for $y_l\ge0$ give
\[
\left|S_{\eta}(p_i)\right|\le2\pi\cdot m^{-1}\cdot\left(m+|p_1|\right)^{-1}\cdot\prod_{l<i}\left(1+|k_l|\right)\qquad(1\le i\le j).
\]
In the product of these $j$ bounds the factor $1+|k_l|$ occurs $j-l\le3$ times. For $l<j$ the factor $v_{r_l}(p_l,p_{l+1})$ is at most $C\cdot\eta^{r_l-1}\cdot\left(1+|k_l|\right)^{-12}$, and the last factor $v_{r_j}(p_j,p_1)$ is at most $C\cdot\eta^{r_j-1}$. Hence the summand of~\eqref{eq:free-coefficient-12} has modulus at most
\[
C\cdot\eta^{\sum_i(r_i-1)}\cdot\left(m+|p_1|\right)^{-j}\cdot\prod_{l<j}\left(1+|k_l|\right)^{-12+3}.
\]
Lemma~\ref{lem:lattice-sums}(a), (b), applied to the sums over $p_1\in B_{\eta}$ and over $k_l\in K_{\eta}\subset K_R$ (with $(1+|k|)^{-9}\le(1+|k|)^{-5}$), gives for the four traces in the brackets of~\eqref{eq:free-coefficient-11}, in each of which $j\le3$ and $r_1+\dots+r_j=4$,
\begin{gather*}
\left|\operatorname{Tr}T_4\right|\le C\cdot\eta^3\cdot\eta^{-2},\qquad\left|\operatorname{Tr}\left(T_1T_3\right)\right|\le C\cdot\eta^2\cdot\eta^{-1},\qquad\left|\operatorname{Tr}\left(T_2^2\right)\right|\le C\cdot\eta^2\cdot\eta^{-1},\\
\left|\operatorname{Tr}\left(T_1^2T_2\right)\right|\le C\cdot\eta\cdot\left(1+\log\frac{1}{\eta}\right)\le C\cdot\eta^{\frac{1}{2}}.
\end{gather*}

\emph{The torus.} By the completeness of the exponentials in $L^2_{\mathrm{ap}}(\mathbb{T}_R;\mathbb{C}^4)$, the functions $e_p\otimes\epsilon$, $p\in P_R$, form an orthonormal basis of this space; the notations $\langle e_p,Xe_q\rangle$ and $Xe_p=Me_p$ are used for operators $X$ on $L^2_{\mathrm{ap}}(\mathbb{T}_R;\mathbb{C}^4)$ as on the lattice. By the definition of $S^R$ (Section~\ref{sec:free-fourth-order}), $S^Re_p=S(p)\,e_p$, with $S(p)=(m+i\gamma\cdot p)^{-1}$ as in Lemma~\ref{lem:lattice-continuum-multipliers}. Let $V$ be the multiplication by $i\gamma\cdot a$ on $L^2_{\mathrm{ap}}(\mathbb{T}_R;\mathbb{C}^4)$, the field $a$ being periodised. As $a$ is supported in the interior of a period cube, $\langle e_p,Ve_q\rangle=R^{-3}\,v_0(p,q)$ for $p,q\in P_R$, where
\[
v_0(p,q):=i\gamma\cdot\hat{a}(p-q)=\sum_{\nu}i\gamma_{\nu}\,\hat{a}_{\nu}(p-q).
\]
By Theorem~\ref{thm:free-coefficient-1}(b4), $c_{16}\left(\check{a}_{\eta,\nu}\right)\to c_{16}(a_{\nu})$ as $\eta\to0$; so $c_{16}(a_{\nu})\le A$, and Lemma~\ref{lem:sampling-aliasing}(a), together with $\left|\gamma_{\nu}\right|=1$, gives
\[
\left|v_0(p,q)\right|\le6A\cdot\left(1+|p-q|\right)^{-16}\qquad(p,q\in P_R).
\]

Let $G^R:=S^RVS^R$. Its matrix elements are $\left\langle e_p,G^Re_q\right\rangle=R^{-3}\,S(p)\,v_0(p,q)\,S(q)$. By Lemma~\ref{lem:lattice-continuum-multipliers} and Peetre's inequality~\eqref{eq:free-coefficient-Pe} with $t=-1$, $|S(p)|\le2^{\frac{1}{2}}\left(m+|p|\right)^{-1}$ and $|S(q)|\le2^{\frac{1}{2}}m^{-1}\left(m+|p|\right)^{-1}\left(1+|p-q|\right)$; and $\|M\|_{\mathrm{HS}}\le2\,|M|$ for every $4\times4$ matrix $M$. Hence, with $k:=p-q$, the squared Hilbert--Schmidt norms of the matrices $\left\langle e_p,G^Re_q\right\rangle$, summed over $p,q\in P_R$, give at most
\[
C\cdot R^{-6}\sum_{p\in P_R}\ \sum_{k\in K_R}\left(m+|p|\right)^{-4}\cdot\left(1+|k|\right)^{-30}\le C,
\]
by Lemma~\ref{lem:lattice-sums}(a), (c). This sum is the square of the Hilbert--Schmidt norm of $G^R$. So $G^R$ is a Hilbert--Schmidt operator, which is Theorem~\ref{thm:free-coefficient-1}(c1) for the torus.

Let $T:=S^RV$. The operator $T^4=G^R\cdot VG^RV$ is the product of two Hilbert--Schmidt operators; so it is of trace class, and its trace is the sum $\sum_{p\in P_R}\operatorname{tr}\left\langle e_p,T^4e_p\right\rangle$ of its diagonal matrix elements in the orthonormal basis $e_p\otimes\epsilon$. The operator $T$ is bounded, with $\left\langle e_p,Te_q\right\rangle=R^{-3}\,S(p)\,v_0(p,q)$. Expanding $\left\langle e_{p_1},T^4e_{p_1}\right\rangle$ by Parseval's identity at the three places between consecutive factors $T$, and writing $p_{i+1}:=p_i-k_i$ with $k_i\in K_R$ ($i=1,2,3$), one obtains
\begin{multline}\label{eq:free-coefficient-13}
\operatorname{Tr}T^4=R^{-12}\sum_{p_1\in P_R}\ \sum_{k_1,k_2,k_3\in K_R}\mathrm{Fz}_0,\\
\mathrm{Fz}_0:=\operatorname{tr}\left[S(p_1)\,v_0(p_1,p_2)\,S(p_2)\,v_0(p_2,p_3)\,S(p_3)\,v_0(p_3,p_4)\,S(p_4)\,v_0(p_4,p_1)\right],
\end{multline}
the series over $\left(p_1,k_1,k_2,k_3\right)$ being absolutely convergent: Lemma~\ref{lem:lattice-continuum-multipliers} and Peetre's inequality~\eqref{eq:free-coefficient-Pe} give $|S(p_i)|\le2^{\frac{1}{2}}m^{-1}\left(m+|p_1|\right)^{-1}\prod_{l<i}\left(1+|k_l|\right)$, so that, with the bound for $v_0$,
\[
\left|\mathrm{Fz}_0\right|\le C\cdot\left(m+|p_1|\right)^{-4}\cdot\prod_{l\le3}\left(1+|k_l|\right)^{-13},
\]
and the right side is summable by Lemma~\ref{lem:lattice-sums}(a), (c).

\emph{The term of fourth order.} Let $\mathrm{Fz}_{\eta}$ denote the summand in~\eqref{eq:free-coefficient-12} with $j=4$ and $r_1=r_2=r_3=r_4=1$; it arises from $\mathrm{Fz}_0$ when $S$ and $v_0$ are replaced by $S_{\eta}$ and $v_1$. Then
\[
\operatorname{Tr}\left(T_1^4\right)=R^{-12}\sum_{p_1\in B_{\eta}}\ \sum_{k_1,k_2,k_3\in K_{\eta}}\mathrm{Fz}_{\eta},
\]
and, by the bound obtained for the terms in brackets, $\left|\mathrm{Fz}_{\eta}\right|\le C\cdot\left(m+|p_1|\right)^{-4}\cdot\prod_{l\le3}\left(1+|k_l|\right)^{-9}$ on $B_{\eta}\times K_{\eta}^3$. Put $\boldsymbol{\Lambda}:=\frac{\pi}{4\eta}$ and
\[
\mathrm{Rg}:=\left\{\left(p_1,k_1,k_2,k_3\right)\in P_R\times K_R^3:\ |p_1|\le\boldsymbol{\Lambda},\ |k_1|\le\boldsymbol{\Lambda},\ |k_2|\le\boldsymbol{\Lambda},\ |k_3|\le\boldsymbol{\Lambda}\right\};
\]
then $\mathrm{Rg}\subset B_{\eta}\times K_{\eta}^3$. A point of $P_R\times K_R^3$ outside $\mathrm{Rg}$ has $|p_1|>\boldsymbol{\Lambda}$ or $|k_l|>\boldsymbol{\Lambda}$ for some $l$. In the sums over such points, bounded by means of the estimates for $\left|\mathrm{Fz}_{\eta}\right|$ and $\left|\mathrm{Fz}_0\right|$, the last bound of Lemma~\ref{lem:lattice-sums}(c), respectively the second bound of Lemma~\ref{lem:lattice-sums}(a) (as $(1+|k|)^{-13}\le(1+|k|)^{-9}\le(1+|k|)^{-6}$), gives the factor $7/\boldsymbol{\Lambda}$, respectively $28/\boldsymbol{\Lambda}$, and the other bounds of Lemma~\ref{lem:lattice-sums}(a), (c) control the remaining sums. Each bound of Lemma~\ref{lem:lattice-sums} concerns a sum multiplied by $R^{-3}$, and the factor $R^{-12}$ supplies one such factor to each of the four summations; as $\boldsymbol{\Lambda}^{-1}=\frac{4\eta}{\pi}$, it follows that
\[
R^{-12}\sum_{\left(B_{\eta}\times K_{\eta}^3\right)\setminus\mathrm{Rg}}\left|\mathrm{Fz}_{\eta}\right|\le C\cdot\eta,\qquad R^{-12}\sum_{\left(P_R\times K_R^3\right)\setminus\mathrm{Rg}}\left|\mathrm{Fz}_0\right|\le C\cdot\eta,
\]
the sums extending over the points $\left(p_1,k_1,k_2,k_3\right)$ of the sets indicated. These two sums bound the parts of the representations of $\operatorname{Tr}\left(T_1^4\right)$ and of $\operatorname{Tr}T^4$ in~\eqref{eq:free-coefficient-13} that come from the points outside $\mathrm{Rg}$.

On $\mathrm{Rg}$ one has $|p_i|\le\frac{i\pi}{4\eta}\le\frac{\pi}{\eta}$ for $i=1,\dots,4$, and $k_4:=p_4-p_1=-\left(k_1+k_2+k_3\right)$ satisfies $|k_4|\le\frac{3\pi}{4\eta}$. Let $p,q\in P_R$ with $|p|\le\frac{\pi}{\eta}$, $|q|\le\frac{\pi}{\eta}$, $|p-q|\le\frac{\pi}{\eta}$, and put $k:=p-q$. Then
\[
v_1(p,q)-v_0(p,q)=\sum_{\nu}\left[i\gamma_{\nu}\left(\cos y_{\nu}-1\right)+\sin y_{\nu}\right]\check{a}_{\eta,\nu}^{\eta}(k)+\sum_{\nu}i\gamma_{\nu}\left(\check{a}_{\eta,\nu}^{\eta}(k)-\hat{a}_{\nu}(k)\right).
\]
For real $y$ the matrix $i\gamma_{\nu}\left(\cos y-1\right)+\sin y$ is normal, with the eigenvalues $\sin y\pm i\left(\cos y-1\right)$, because $\gamma_{\nu}$ is Hermitian with $\gamma_{\nu}^2=1$; therefore
\[
\left|i\gamma_{\nu}\left(\cos y-1\right)+\sin y\right|=\left(\left(\cos y-1\right)^2+\sin^2y\right)^{\frac{1}{2}}=2\left|\sin\frac{y}{2}\right|\le2\min\left(1,|y|\right)\le2|y|^{\frac{1}{2}}.
\]
Moreover $\left|y_{\nu}\right|\le\eta\left(|p|+|k|\right)\le\eta\cdot\left(m+|p|\right)\cdot\left(1+|k|\right)\cdot m^{-1}$, as $m\le1$. Since $|k|\le\frac{\pi}{\eta}$, Lemma~\ref{lem:sampling-aliasing}(d) gives $\left|\check{a}_{\eta,\nu}^{\eta}(k)\right|\le2A\cdot(1+|k|)^{-12}$, and Lemma~\ref{lem:sampling-aliasing}(c), together with the bound for $\hat{\check{a}}_{\eta,\nu}-\hat{a}_{\nu}$ obtained above, gives
\[
\left|\check{a}_{\eta,\nu}^{\eta}(k)-\hat{a}_{\nu}(k)\right|\le\left(A\eta^4+A'\eta^2\right)\cdot(1+|k|)^{-12}.
\]
As $\eta^4\le\eta^2\le\eta^{\frac{1}{2}}\le m^{-\frac{1}{2}}\cdot\eta^{\frac{1}{2}}\cdot\left(m+|p|\right)^{\frac{1}{2}}$, it follows that
\[
\left|v_1(p,q)-v_0(p,q)\right|\le C\cdot\eta^{\frac{1}{2}}\cdot\left(m+|p|\right)^{\frac{1}{2}}\cdot\left(1+|k|\right)^{-11}.
\]

Put $p_5:=p_1$, and for $i=1,\dots,4$ let
\[
M_{2i-1}:=S_{\eta}(p_i),\qquad M_{2i}:=v_1(p_i,p_{i+1}),\qquad M'_{2i-1}:=S(p_i),\qquad M'_{2i}:=v_0(p_i,p_{i+1}),
\]
so that $\mathrm{Fz}_{\eta}=\operatorname{tr}\left[M_1M_2\cdots M_8\right]$ and $\mathrm{Fz}_0=\operatorname{tr}\left[M'_1M'_2\cdots M'_8\right]$. Then
\[
\mathrm{Fz}_{\eta}-\mathrm{Fz}_0=\sum_{s=1}^8\operatorname{tr}\left[M'_1\cdots M'_{s-1}\left(M_s-M'_s\right)M_{s+1}\cdots M_8\right].
\]
In the summand of index $s$ the matrices to the left of $M_s-M'_s$ are taken from $\mathrm{Fz}_0$ and those to its right from $\mathrm{Fz}_{\eta}$; the middle factor $M_s-M'_s$ is $S_{\eta}(p_i)-S(p_i)$ when $s=2i-1$, and $v_1(p_i,p_{i+1})-v_0(p_i,p_{i+1})$ when $s=2i$, the pair $\left(p_4,p_5\right)=\left(p_4,p_1\right)$ of $s=8$ closing the cycle. On $\mathrm{Rg}$, Lemma~\ref{lem:lattice-continuum-multipliers} applies at every $p_i$:
\[
\left|S_{\eta}(p_i)\right|\le\frac{2\pi}{m+|p_i|},\qquad\left|S(p_i)\right|\le\frac{2^{\frac{1}{2}}}{m+|p_i|},\qquad\left|S_{\eta}(p_i)-S(p_i)\right|\le15\,\eta^{\frac{1}{2}}\left(m+|p_i|\right)^{-\frac{1}{2}}.
\]
For $i\le3$ the factors $v_1(p_i,p_{i+1})$ and $v_0(p_i,p_{i+1})$ are at most $12A\cdot\left(1+|k_i|\right)^{-12}$, and their difference is at most $C\cdot\eta^{\frac{1}{2}}\cdot\left(m+|p_i|\right)^{\frac{1}{2}}\cdot\left(1+|k_i|\right)^{-11}$. The factors $v_1(p_4,p_1)$ and $v_0(p_4,p_1)$ are at most $12A$, and their difference is at most $C\cdot\eta^{\frac{1}{2}}\cdot\left(m+|p_4|\right)^{\frac{1}{2}}$, since $|p_4|$, $|p_1|$ and $|k_4|$ are at most $\frac{\pi}{\eta}$. Hence each of the eight summands is at most
\[
C\cdot\eta^{\frac{1}{2}}\cdot\prod_{i=1}^4\left(m+|p_i|\right)^{t_i}\cdot\prod_{l\le3}\left(1+|k_l|\right)^{-11},
\]
with exponents $t_i\in\left\{-1,-\frac{1}{2}\right\}$ whose sum is $-\frac{7}{2}$. By Peetre's inequality~\eqref{eq:free-coefficient-Pe}, $\left(m+|p_i|\right)^{t_i}\le m^{-1}\cdot\left(m+|p_1|\right)^{t_i}\cdot\left(1+|p_1-p_i|\right)$, as $|t_i|\le1$; and $1+|p_1-p_i|\le\prod_{l<i}\left(1+|k_l|\right)$. In this way each factor $1+|k_l|$ acquires an additional exponent of at most four, and each of the eight summands is at most
\[
C\cdot\eta^{\frac{1}{2}}\cdot\left(m+|p_1|\right)^{-\frac{7}{2}}\cdot\prod_{l\le3}\left(1+|k_l|\right)^{-11+4}.
\]
Lemma~\ref{lem:lattice-sums}(a), (c) gives
\[
R^{-12}\sum_{\mathrm{Rg}}\left|\mathrm{Fz}_{\eta}-\mathrm{Fz}_0\right|\le C\cdot\eta^{\frac{1}{2}},
\]
the sum extending over the points $\left(p_1,k_1,k_2,k_3\right)$ of $\mathrm{Rg}$. Together with the bounds for the points outside $\mathrm{Rg}$ this yields
\[
\left|\operatorname{Tr}\left(T_1^4\right)-\operatorname{Tr}T^4\right|\le C\cdot\eta^{\frac{1}{2}}.
\]

\emph{Conclusion.} By the definition in Section~\ref{sec:free-fourth-order}, $b_R(f)=-6\operatorname{Tr}T^4$. By~\eqref{eq:free-coefficient-11}, the bounds for the four traces in the brackets and the last estimate,
\[
\left|b_{\eta,R}(0)-b_R(f)\right|\le24\left(\left|\operatorname{Tr}T_4\right|+\left|\operatorname{Tr}\left(T_1T_3\right)\right|+\frac{1}{2}\left|\operatorname{Tr}\left(T_2^2\right)\right|+\left|\operatorname{Tr}\left(T_1^2T_2\right)\right|\right)+6\left|\operatorname{Tr}\left(T_1^4\right)-\operatorname{Tr}T^4\right|\le C_1\cdot\eta^{\frac{1}{2}},
\]
where $C_1$ depends only on $m$, $A$ and $A'$, hence only on $\bar{m}$ and $\beta$; none of the constants depends on $R$ or on $\eta$. This is the inequality of Theorem~\ref{thm:free-coefficient-1}(c2). If $RL^N$ is an integer for all large $N$, the choice $\eta=L^{-N}$ in it gives $b_R(f)=\lim_{N\to\infty}b_{L^{-N},R}(0)$.
\end{proof}

\subsubsection{From the torus to the whole space}\label{sec:torus-whole-space}

\begin{proof}[Proof of Theorem~\ref{thm:free-coefficient-1}(c3)]
Let $\beta$, $a$, $f$ and $R$ be as in Theorem~\ref{thm:free-coefficient-1}(c), and put $\mathcal{Q}:=Q_{\beta}$, $D:=D_{\beta}$ and $m:=\bar{m}$. Then $D\ge4$, as $\rho_{\beta}\ge2$ (Section~\ref{sec:probes-discretisation}), and $R\ge D+1\ge5$; the cube $\mathcal{Q}$ lies in the interior of a period cube of the torus of side $R$. Define
\[
\mathfrak{i}_R\colon L^2(\mathcal{Q};\mathbb{C}^4)\to L^2_{\mathrm{ap}}(\mathbb{T}_R;\mathbb{C}^4),\qquad\mathfrak{i}_{\infty}\colon L^2(\mathcal{Q};\mathbb{C}^4)\to L^2(\mathbb{R}^3;\mathbb{C}^4)
\]
as follows: $\mathfrak{i}_R\boldsymbol{\psi}$ is the antiperiodic function that agrees with $\boldsymbol{\psi}$ on $\mathcal{Q}$ and is zero on the complement of $\mathcal{Q}$ in a period cube containing $\mathcal{Q}$, and $\mathfrak{i}_{\infty}\boldsymbol{\psi}$ is the extension of $\boldsymbol{\psi}$ by zero. The maps $\mathfrak{i}_R$ and $\mathfrak{i}_{\infty}$ preserve norms. Let $V_Q$ be the multiplication by $i\gamma\cdot a$ on $L^2(\mathcal{Q};\mathbb{C}^4)$, and let $V$ be the multiplication by $i\gamma\cdot a$ on $L^2_{\mathrm{ap}}(\mathbb{T}_R;\mathbb{C}^4)$, with $a$ periodised, and on $L^2(\mathbb{R}^3;\mathbb{C}^4)$. Since the periodised field vanishes outside the union of the translates of $\mathcal{Q}$ by $R\mathbb{Z}^3$, and $a$ vanishes outside $\mathcal{Q}$,
\[
V=\mathfrak{i}_RV_Q\mathfrak{i}_R^{*}\ \text{on}\ L^2_{\mathrm{ap}}(\mathbb{T}_R;\mathbb{C}^4),\qquad V=\mathfrak{i}_{\infty}V_Q\mathfrak{i}_{\infty}^{*}\ \text{on}\ L^2(\mathbb{R}^3;\mathbb{C}^4),
\]
and $\left\|V_Q\right\|=\|a\|_{\infty}$, with $\|a\|_{\infty}=\sup_x|a(x)|_2$ as in Lemma~\ref{lem:hilbert-schmidt-bound}, because $\gamma\cdot a(x)$ is Hermitian with square $|a(x)|_2^2$. Put
\[
S^R_Q:=\mathfrak{i}_R^{*}S^R\mathfrak{i}_R,\qquad S^{\infty}_Q:=\mathfrak{i}_{\infty}^{*}S_m\mathfrak{i}_{\infty};
\]
both have norm at most $1/m$, since $S^R$ and $S_m$ act as the multipliers $(m+i\gamma\cdot p)^{-1}$, whose norms $\left(m^2+|p|_2^2\right)^{-\frac{1}{2}}$ are at most $1/m$ by~\eqref{eq:free-coefficient-K2}. Let $G^R:=S^RVS^R$ and $G^{\infty}:=S_mVS_m$, and, for $\#\in\{R,\infty\}$,
\[
G^{\#}_Q:=S^{\#}_QV_QS^{\#}_Q=\mathfrak{i}_{\#}^{*}G^{\#}\mathfrak{i}_{\#}.
\]

\emph{Reduction to the cube.} By Theorem~\ref{thm:free-coefficient-1}(c1), $G^{\#}$ is a Hilbert--Schmidt operator; for $\#=\infty$ this is Lemma~\ref{lem:hilbert-schmidt-bound} with $\mathsf{M}=m$. The operator $X:=G^{\#}\mathfrak{i}_{\#}V_Q\mathfrak{i}_{\#}^{*}G^{\#}\mathfrak{i}_{\#}V_Q$, which maps $L^2(\mathcal{Q};\mathbb{C}^4)$ into the space on which $G^{\#}$ acts, is the product of the Hilbert--Schmidt operators $G^{\#}$ and $\mathfrak{i}_{\#}V_Q\mathfrak{i}_{\#}^{*}G^{\#}\mathfrak{i}_{\#}V_Q$; so it is of trace class, while $\mathfrak{i}_{\#}^{*}$ is bounded. As the trace of the product of an operator of trace class and a bounded operator does not change when the two factors are interchanged,
\[
\operatorname{Tr}\left(G^{\#}VG^{\#}V\right)=\operatorname{Tr}\left(X\mathfrak{i}_{\#}^{*}\right)=\operatorname{Tr}_{L^2(\mathcal{Q};\mathbb{C}^4)}\left(\mathfrak{i}_{\#}^{*}G^{\#}\mathfrak{i}_{\#}V_Q\mathfrak{i}_{\#}^{*}G^{\#}\mathfrak{i}_{\#}V_Q\right)=\operatorname{Tr}\left(\left(G^{\#}_QV_Q\right)^2\right).
\]
As $\left(S^RV\right)^4=G^RVG^RV$ and $\left(S_mV\right)^4=G^{\infty}VG^{\infty}V$, the definitions of Section~\ref{sec:free-fourth-order} give
\[
b_R(f)=-6\operatorname{Tr}\left(\left(G^R_QV_Q\right)^2\right),\qquad b(f)=-6\operatorname{Tr}\left(\left(G^{\infty}_QV_Q\right)^2\right).
\]

\emph{The difference of the two propagators.} Let $\boldsymbol{\phi}\in L^2_{\mathrm{ap}}(\mathbb{T}_R;\mathbb{C}^4)$. Since $\left\|S_m\right\|_{L^1}\le3/m$ by~\eqref{eq:free-coefficient-K1}, the convolution
\[
\left(S_m*\boldsymbol{\phi}\right)(x)=\int\mathrm{d}z\,S_m(z)\,\boldsymbol{\phi}(x-z)
\]
is defined for almost every $x\in\mathbb{R}^3$ and is antiperiodic; Minkowski's inequality, together with the invariance of the norm of $L^2_{\mathrm{ap}}(\mathbb{T}_R;\mathbb{C}^4)$ under translations, gives $\left\|S_m*\boldsymbol{\phi}\right\|\le\left\|S_m\right\|_{L^1}\cdot\|\boldsymbol{\phi}\|$. For $\boldsymbol{\phi}=e_p\otimes\epsilon$, with $p\in P_R$ and $\epsilon$ in the standard basis of $\mathbb{C}^4$ as in Section~\ref{sec:lattice-torus}, one has $S_m*\boldsymbol{\phi}=e_p\otimes S(p)\epsilon=S^R\boldsymbol{\phi}$ by~\eqref{eq:free-coefficient-K2}. Thus the bounded operators $S^R$ and $\boldsymbol{\phi}\mapsto S_m*\boldsymbol{\phi}$ agree on an orthonormal basis, and $S^R\boldsymbol{\phi}=S_m*\boldsymbol{\phi}$ for every $\boldsymbol{\phi}\in L^2_{\mathrm{ap}}(\mathbb{T}_R;\mathbb{C}^4)$.

Let $\boldsymbol{\psi}\in L^2(\mathcal{Q};\mathbb{C}^4)$. For $\mathsf{n}\in\mathbb{Z}^3$ and $y\in\mathcal{Q}$ one has $\left(\mathfrak{i}_R\boldsymbol{\psi}\right)(y+\mathsf{n}\cdot R)=(-1)^{\mathsf{n}_1+\mathsf{n}_2+\mathsf{n}_3}\boldsymbol{\psi}(y)$, and $\mathfrak{i}_R\boldsymbol{\psi}$ vanishes outside the union of the pairwise disjoint cubes $\mathcal{Q}+\mathsf{n}\cdot R$. Hence, for almost every $x\in\mathcal{Q}$,
\[
\left(S^R\mathfrak{i}_R\boldsymbol{\psi}\right)(x)=\sum_{\mathsf{n}\in\mathbb{Z}^3}(-1)^{\mathsf{n}_1+\mathsf{n}_2+\mathsf{n}_3}\int_{\mathcal{Q}}\mathrm{d}y\,S_m(x-y-\mathsf{n}\cdot R)\,\boldsymbol{\psi}(y),
\]
and the term with $\mathsf{n}=0$ is $\left(S^{\infty}_Q\boldsymbol{\psi}\right)(x)$. Consequently $\mathrm{Er}:=S^R_Q-S^{\infty}_Q$ is an integral operator on $L^2(\mathcal{Q};\mathbb{C}^4)$, whose kernel is
\[
\mathrm{Er}(x,y)=\sum_{\mathsf{n}\in\mathbb{Z}^3\setminus\{0\}}(-1)^{\mathsf{n}_1+\mathsf{n}_2+\mathsf{n}_3}\,S_m(x-y-\mathsf{n}\cdot R)\qquad(x,y\in\mathcal{Q}),
\]
the series converging uniformly on $\mathcal{Q}\times\mathcal{Q}$ by the bound of the next step.

\emph{A bound for the kernel.} Let $x,y\in\mathcal{Q}$ and $\mathsf{n}\in\mathbb{Z}^3\setminus\{0\}$. As $|x-y|\le D$,
\[
\left|x-y-\mathsf{n}\cdot R\right|_2\ge\left|x-y-\mathsf{n}\cdot R\right|\ge|\mathsf{n}|\cdot R-D\ge R-D\ge1.
\]
For $z\in\mathbb{R}^3$ with $\mathsf{r}:=|z|_2\ge1$, the bound~\eqref{eq:free-coefficient-K1} and $m\le1$ give
\[
\left|S_m(z)\right|\le\frac{\exp(-m\cdot\mathsf{r})\cdot\left(2m\cdot\mathsf{r}+1\right)}{4\pi\,\mathsf{r}^2}\le\frac{(2m+1)\cdot\exp(-m\cdot\mathsf{r})}{4\pi}\le\frac{\exp(-m\cdot\mathsf{r})}{4}.
\]
For an integer $l\ge1$ there are $(2l+1)^3-(2l-1)^3=24l^2+2\le26l^2$ vectors $\mathsf{n}\in\mathbb{Z}^3$ with $|\mathsf{n}|=l$, by the count in the proof of Lemma~\ref{lem:lattice-sums}; and, as $R\ge5$, $\exp\left(-m(lR-D)\right)\le\exp\left(-m(R-D)\right)\cdot\exp\left(-5m(l-1)\right)$. Therefore
\[
\sup_{x,y\in\mathcal{Q}}\left|\mathrm{Er}(x,y)\right|\le\frac{c_m}{4}\cdot\exp\left(-m(R-D)\right),\qquad c_m:=26\sum_{l\ge1}l^2\exp\left(-5m(l-1)\right).
\]
The Hilbert--Schmidt norm of $\mathrm{Er}$ equals the $L^2$ norm over $\mathcal{Q}\times\mathcal{Q}$ of the function $(x,y)\mapsto\left\|\mathrm{Er}(x,y)\right\|_{\mathrm{HS}}$; as $\|M\|_{\mathrm{HS}}\le2\,|M|$ for every $4\times4$ matrix $M$, and as $\mathcal{Q}\times\mathcal{Q}$ has the volume $D^6$, it follows that
\[
\left\|\mathrm{Er}\right\|_{\mathrm{HS}}\le2D^3\cdot\sup_{x,y\in\mathcal{Q}}\left|\mathrm{Er}(x,y)\right|.
\]

\emph{Conclusion.} One has
\[
G^R_Q-G^{\infty}_Q=\mathrm{Er}\,V_Q\,S^R_Q+S^{\infty}_Q\,V_Q\,\mathrm{Er},
\]
so that the Hilbert--Schmidt norm $\delta_G$ of $G^R_Q-G^{\infty}_Q$ satisfies, by the preceding step,
\[
\delta_G\le\frac{2\left\|\mathrm{Er}\right\|_{\mathrm{HS}}\cdot\|a\|_{\infty}}{m}\le\frac{c_mD^3\cdot\|a\|_{\infty}}{m}\cdot\exp\left(-m(R-D)\right).
\]
Moreover
\[
\operatorname{Tr}\left(\left(G^R_QV_Q\right)^2\right)-\operatorname{Tr}\left(\left(G^{\infty}_QV_Q\right)^2\right)=\operatorname{Tr}\left(\left(G^R_Q-G^{\infty}_Q\right)V_QG^R_QV_Q\right)+\operatorname{Tr}\left(G^{\infty}_QV_Q\left(G^R_Q-G^{\infty}_Q\right)V_Q\right).
\]
By $\left|\operatorname{Tr}\left(Y_1Y_2\right)\right|\le\left\|Y_1\right\|_{\mathrm{HS}}\cdot\left\|Y_2\right\|_{\mathrm{HS}}$ for Hilbert--Schmidt operators $Y_1$ and $Y_2$, $\left\|G^{\infty}_Q\right\|_{\mathrm{HS}}\le\left\|G^{\infty}\right\|_{\mathrm{HS}}$ and $\left\|G^R_Q\right\|_{\mathrm{HS}}\le\left\|G^{\infty}_Q\right\|_{\mathrm{HS}}+\delta_G$, the modulus of the right side is at most $\delta_G\cdot\|a\|_{\infty}^2\cdot\left(2\left\|G^{\infty}\right\|_{\mathrm{HS}}+\delta_G\right)$; and $\left\|G^{\infty}\right\|_{\mathrm{HS}}^2\le\|a\|_{L^2}^2\cdot(\pi m)^{-1}$ by Lemma~\ref{lem:hilbert-schmidt-bound}. With the reduction to the cube and $\exp\left(-m(R-D)\right)\le1$ it follows that
\[
\left|b_R(f)-b(f)\right|\le6\,\delta_G\cdot\|a\|_{\infty}^2\cdot\left(2\|a\|_{L^2}\cdot(\pi m)^{-\frac{1}{2}}+\delta_G\right)\le C_2\cdot\exp\left(-m(R-D)\right),
\]
where
\[
C_2:=6\cdot\frac{c_mD^3\cdot\|a\|_{\infty}^3}{m}\cdot\left(2\|a\|_{L^2}\cdot(\pi m)^{-\frac{1}{2}}+\frac{c_mD^3\cdot\|a\|_{\infty}}{m}\right).
\]
Since $m=\bar{m}$ and $D=D_{\beta}$, this is Theorem~\ref{thm:free-coefficient-1}(c3), with a constant $C_2$ that depends only on $\bar{m}$ and $\beta$.

Together with Theorem~\ref{thm:free-coefficient-1}(c2), Theorem~\ref{thm:free-coefficients-lattice}(c) follows.
\end{proof}

\subsection{The coefficient of fourth order for large mass}\label{sec:coefficient-fourth-order}

Throughout this subsection $\mathsf{M}>0$, and $|\cdot|$ denotes the Euclidean norm $|\cdot|_2$ on $\mathbb{R}^3$ and on $\mathbb{C}^4$; for a $4\times4$ matrix, $|\cdot|$ is the operator norm with respect to the Euclidean norm on $\mathbb{C}^4$, as in Section~\ref{sec:letters-norms-dirac}. For a measurable function $g$, defined almost everywhere on $\mathbb{R}^3$ and with values in $[0,\infty]$, in $\mathbb{R}^3$ or in $\mathbb{C}^{4\times4}$, let $\|g\|_{L^1}:=\int\mathrm{d}z\,|g(z)|\in[0,\infty]$.

Let $a\colon\mathbb{R}^3\to\mathbb{R}^3$ be smooth. The map $a$ is called a \emph{field} if either
\begin{equation}\tag{cpt}\label{eq:free-coefficient-cpt}
a\in C_c^{\infty}(\mathbb{R}^3;\mathbb{R}^3),
\end{equation}
or there is a real antisymmetric $3\times3$ matrix $F=(F_{\nu\mu})$ such that
\begin{equation}\tag{lin}\label{eq:free-coefficient-lin}
a=A_F,\qquad A_{F,\mu}(z):=\frac{1}{2}\sum_{\nu=1}^3z_{\nu}\cdot F_{\nu\mu}\qquad\left(z\in\mathbb{R}^3,\ \mu\in\{1,2,3\}\right).
\end{equation}
In the first case the field $a$ is said to be of type~\eqref{eq:free-coefficient-cpt}, in the second of type~\eqref{eq:free-coefficient-lin}. For every field $a$ there is a number $c_a\ge0$ such that
\[
|a(z)|\le c_a\cdot(1+|z|),\qquad|a(z)-a(z')|\le c_a\cdot|z-z'|\qquad(z,z'\in\mathbb{R}^3),
\]
since a field of type~\eqref{eq:free-coefficient-cpt} is bounded and Lipschitz continuous, and a field of type~\eqref{eq:free-coefficient-lin} is linear; $c_a$ denotes such a number. To a field $a$ belong three further objects. The first is the function $V\colon\mathbb{R}^3\to\mathbb{C}^{4\times4}$, $V(z):=i\gamma\cdot a(z)$; the letter $V$ also denotes the operator of multiplication by this function. The second is the field strength, with the components
\[
F_{\nu\mu}(x):=\partial_{\nu}a_{\mu}(x)-\partial_{\mu}a_{\nu}(x)\qquad\left(x\in\mathbb{R}^3,\ \nu,\mu\in\{1,2,3\}\right).
\]
The third is the vector field $B:=\operatorname{curl}a$, so that $B_{\mathfrak{l}}=\frac{1}{2}\sum_{\nu,\mu}\varepsilon_{\mathfrak{l}\nu\mu}\cdot F_{\nu\mu}$, where $\varepsilon$ is the totally antisymmetric symbol with $\varepsilon_{123}=1$. For $a=A_F$ one has $\partial_{\nu}A_{F,\mu}=\frac{1}{2}F_{\nu\mu}$; as $F$ is antisymmetric, the field strength of $A_F$ is the constant matrix $F$. Finally,
\[
\mathscr{S}:=\mathcal{S}(\mathbb{R}^3;\mathbb{C}^4)
\]
is the Schwartz space of smooth rapidly decreasing functions $\mathbb{R}^3\to\mathbb{C}^4$.

\subsubsection{Kernels of products of propagators}\label{sec:kernels-products-propagators}

For $\boldsymbol{\kappa}>0$ and $c>0$ let $\mathsf{s}\colon\mathbb{R}^3\setminus\{0\}\to(0,\infty)$ be the function
\[
\mathsf{s}(z):=\boldsymbol{\kappa}\cdot|z|^{-2}\cdot\exp(-c|z|).
\]
For measurable functions $f,g$ with values in $[0,\infty]$, defined almost everywhere on $\mathbb{R}^3$, let
\[
(f*g)(z):=\int\mathrm{d}w\,f(z-w)\cdot g(w)\in[0,\infty]\qquad(z\in\mathbb{R}^3).
\]
By Fubini's theorem for non-negative functions, $f*g$ is again a measurable function with values in $[0,\infty]$, and $(f*g)*h=f*(g*h)$ for every third function $h$ of the same kind; the substitution $w\mapsto z-w$ gives $f*g=g*f$. The $n$-fold convolution of $\mathsf{s}$ is defined by $\mathsf{s}^{*1}:=\mathsf{s}$ and $\mathsf{s}^{*(n+1)}:=\mathsf{s}^{*n}*\mathsf{s}$ for $n\ge1$; associativity and induction on $n$ give
\[
\mathsf{s}^{*(m+n)}=\mathsf{s}^{*m}*\mathsf{s}^{*n}\qquad(m,n\ge1).
\]

\begin{lemma}[Convolution powers]\label{lem:convolution-powers}
The function $\mathsf{s}$ is integrable on $\mathbb{R}^3$. There is a number $c'<\infty$, depending only on $\boldsymbol{\kappa}$ and $c$, such that for all $z\ne0$
\[
\mathsf{s}^{*2}(z)\le c'\cdot|z|^{-1}\cdot\exp(-c|z|),\qquad\mathsf{s}^{*2}(z)\le c'\cdot|z|^{-\frac{3}{2}},\qquad\mathsf{s}^{*3}(z)\le c'\cdot|z|^{-\frac{1}{2}}.
\]
Finally, $\mathsf{s}^{*4}$ and $\mathsf{s}^{*5}$ have finite suprema on $\mathbb{R}^3$ and are continuous there.
\end{lemma}

\begin{proof}
In polar coordinates,
\[
\int\mathrm{d}z\,\mathsf{s}(z)=4\pi\cdot\boldsymbol{\kappa}\int_0^{\infty}\mathrm{d}r\,\exp(-cr)=\frac{4\pi\cdot\boldsymbol{\kappa}}{c}<\infty.
\]
Let $0<\boldsymbol{\alpha}<3$ and $0<\boldsymbol{\beta}<3$ with $\boldsymbol{\alpha}+\boldsymbol{\beta}>3$, and for $z\ne0$ put
\[
I_{\boldsymbol{\alpha},\boldsymbol{\beta}}(z):=\int\mathrm{d}w\,|z-w|^{-\boldsymbol{\alpha}}\cdot|w|^{-\boldsymbol{\beta}}.
\]
This integral is finite: the singularities at $w=z$ and at $w=0$ are integrable because $\boldsymbol{\alpha}<3$ and $\boldsymbol{\beta}<3$, and for $|w|\ge2|z|$ one has $|z-w|\ge\frac{1}{2}|w|$, so that the integrand is at most $2^{\boldsymbol{\alpha}}\cdot|w|^{-\boldsymbol{\alpha}-\boldsymbol{\beta}}$, which is integrable over $\{|w|\ge2|z|\}$ because $\boldsymbol{\alpha}+\boldsymbol{\beta}>3$. The integral does not change if $z$ and the variable $w$ are rotated by the same rotation; hence $I_{\boldsymbol{\alpha},\boldsymbol{\beta}}(z)$ depends only on $|z|$. The substitution $w=|z|\cdot w'$ gives $I_{\boldsymbol{\alpha},\boldsymbol{\beta}}(z)=|z|^{3-\boldsymbol{\alpha}-\boldsymbol{\beta}}\cdot I_{\boldsymbol{\alpha},\boldsymbol{\beta}}\left(|z|^{-1}z\right)$. Therefore
\[
I_{\boldsymbol{\alpha},\boldsymbol{\beta}}(z)=c_{\boldsymbol{\alpha},\boldsymbol{\beta}}\cdot|z|^{3-\boldsymbol{\alpha}-\boldsymbol{\beta}}\quad(z\ne0),\qquad c_{\boldsymbol{\alpha},\boldsymbol{\beta}}:=I_{\boldsymbol{\alpha},\boldsymbol{\beta}}\left((1,0,0)\right)<\infty.
\]
Since $|z|\le|z-w|+|w|$, one has $\exp(-c|z-w|)\cdot\exp(-c|w|)\le\exp(-c|z|)$, and consequently
\[
\mathsf{s}^{*2}(z)=\boldsymbol{\kappa}^2\int\mathrm{d}w\,|z-w|^{-2}\cdot|w|^{-2}\cdot\exp\left(-c|z-w|-c|w|\right)\le\boldsymbol{\kappa}^2\cdot c_{2,2}\cdot|z|^{-1}\cdot\exp(-c|z|).
\]
For $y>0$ one has $y^{\frac{1}{2}}\exp(-cy)\le(2c)^{-\frac{1}{2}}$ (the maximum, attained at $y=\frac{1}{2c}$, is $(2c)^{-\frac{1}{2}}\exp\left(-\frac{1}{2}\right)$); hence $y^{-1}\exp(-cy)\le(2c)^{-\frac{1}{2}}\cdot y^{-\frac{3}{2}}$, and
\[
\mathsf{s}^{*2}(z)\le\boldsymbol{\kappa}^2\cdot c_{2,2}\cdot(2c)^{-\frac{1}{2}}\cdot|z|^{-\frac{3}{2}}.
\]
As $\mathsf{s}^{*3}=\mathsf{s}*\mathsf{s}^{*2}$, this bound and $\exp(-c|z-w|)\le1$ give
\[
\mathsf{s}^{*3}(z)=\int\mathrm{d}w\,\mathsf{s}(z-w)\cdot\mathsf{s}^{*2}(w)\le\boldsymbol{\kappa}^3\cdot c_{2,2}\cdot(2c)^{-\frac{1}{2}}\cdot I_{2,\frac{3}{2}}(z)=\boldsymbol{\kappa}^3\cdot c_{2,2}\cdot c_{2,\frac{3}{2}}\cdot(2c)^{-\frac{1}{2}}\cdot|z|^{-\frac{1}{2}}.
\]
The three inequalities of the lemma therefore hold with
\[
c':=\boldsymbol{\kappa}^2\cdot c_{2,2}\cdot\max\left\{1,\ (2c)^{-\frac{1}{2}},\ \boldsymbol{\kappa}\cdot c_{2,\frac{3}{2}}\cdot(2c)^{-\frac{1}{2}}\right\}.
\]
By the first of them $\mathsf{s}^{*2}$ lies in $L^2(\mathbb{R}^3)$, since $\int\mathrm{d}z\,|z|^{-2}\exp(-2c|z|)=\frac{2\pi}{c}$. By the rule $\mathsf{s}^{*(m+n)}=\mathsf{s}^{*m}*\mathsf{s}^{*n}$ one has $\mathsf{s}^{*4}=\mathsf{s}^{*2}*\mathsf{s}^{*2}$ and $\mathsf{s}^{*5}=\mathsf{s}*\mathsf{s}^{*4}$. Let $f,g$ be non-negative functions in $L^2(\mathbb{R}^3)$, and let $z,z'\in\mathbb{R}^3$. Young's inequality gives $(f*g)(z)\le\|f\|_{L^2}\cdot\|g\|_{L^2}$. The substitution $w\mapsto w+z'-z$ gives $(f*g)(z')=\int\mathrm{d}w\,f(z-w)\cdot g(w+z'-z)$, so that $|(f*g)(z)-(f*g)(z')|$ does not exceed the value at $z$ of the convolution of $f$ with the non-negative function $\left|g-g(\cdot+z'-z)\right|$; the same inequality therefore yields
\[
|(f*g)(z)-(f*g)(z')|\le\|f\|_{L^2}\cdot\left\|g-g(\cdot+z'-z)\right\|_{L^2}.
\]
The right-hand side tends to $0$ as $z'\to z$, because translations act continuously on $L^2(\mathbb{R}^3)$. With $f=g=\mathsf{s}^{*2}$ it follows that $\mathsf{s}^{*4}$ is bounded and continuous. The function $\mathsf{s}^{*5}$ is the convolution of the integrable function $\mathsf{s}$ with the bounded continuous function $\mathsf{s}^{*4}$; it is bounded by $\|\mathsf{s}\|_{L^1}\cdot\sup\mathsf{s}^{*4}$, and it is continuous by dominated convergence.
\end{proof}

By~\eqref{eq:free-coefficient-K1}, $|S_{\mathsf{M}}(z)|\le(4\pi)^{-1}\cdot|z|^{-2}\cdot\exp(-\mathsf{M}\cdot|z|)\cdot(2\mathsf{M}\cdot|z|+1)$ for $z\ne0$. For every integer $i\ge0$ the function
\[
\mathsf{r}\mapsto(4\pi)^{-1}\cdot(2\mathsf{M}\cdot\mathsf{r}+1)\cdot(1+\mathsf{r})^i\cdot\exp\left(-\frac{\mathsf{M}\cdot\mathsf{r}}{2}\right)
\]
is bounded on $[0,\infty)$. Hence, for every integer $i\ge0$, there is a number $\boldsymbol{\kappa}>0$ (for instance the supremum of this function) such that
\[
|S_{\mathsf{M}}(z)|\cdot(1+|z|)^i\le\boldsymbol{\kappa}\cdot|z|^{-2}\cdot\exp\left(-\frac{\mathsf{M}\cdot|z|}{2}\right)\qquad(z\ne0).
\]
From now on, $\mathsf{s}$ denotes the function defined at the beginning of this subsubsection with $c:=\frac{\mathsf{M}}{2}$ and with such a number $\boldsymbol{\kappa}$ for $i=24$; thus
\[
|S_{\mathsf{M}}(z)|\cdot(1+|z|)^{24}\le\mathsf{s}(z)\qquad(z\ne0).
\]

For an integer $n$ with $0\le n\le4$ and for $x,y\in\mathbb{R}^3$ with $x\ne y$ put, with $z_0:=x$ and $z_{n+1}:=y$,
\[
\mathrm{Rk}^{(n)}(x,y):=(-1)^n\int_{(\mathbb{R}^3)^n}\mathrm{d}z_1\cdots\mathrm{d}z_n\,S_{\mathsf{M}}(z_0-z_1)\cdot\prod_{l=1}^n\left(V(z_l)\cdot S_{\mathsf{M}}(z_l-z_{l+1})\right),
\]
the factors of the product being ordered from left to right by increasing $l$; in particular $\mathrm{Rk}^{(0)}(x,y)=S_{\mathsf{M}}(x-y)$. The kernel $\mathrm{Rk}^{(n)}$, a function on $\{(x,y)\in\mathbb{R}^3\times\mathbb{R}^3:x\ne y\}$ with values in $\mathbb{C}^{4\times4}$, depends on $\mathsf{M}$ and on the field $a$; where the field is to be indicated, it is written $\mathrm{Rk}^{(n)}_a$.

\begin{lemma}[The kernels $\mathrm{Rk}^{(n)}$]\label{lem:kernels}
Let $a$ be a field and $n$ an integer with $0\le n\le4$.
\begin{enumerate}
\item[(a)] For $x\ne y$ the integral that defines $\mathrm{Rk}^{(n)}(x,y)$ converges absolutely, and
\[
\left|\mathrm{Rk}^{(n)}(x,y)\right|\le\left(c_a\cdot(1+|x|)\right)^n\cdot\mathsf{s}^{*(n+1)}(x-y).
\]
\item[(b)] Both $S_{\mathsf{M}}$ and the multiplication by $V$ send $\mathscr{S}$ into itself, and every $\boldsymbol{\phi}\in\mathscr{S}$ satisfies
\[
\left((-1)^n\,(S_{\mathsf{M}}V)^nS_{\mathsf{M}}\,\boldsymbol{\phi}\right)(x)=\int\mathrm{d}y\,\mathrm{Rk}^{(n)}(x,y)\cdot\boldsymbol{\phi}(y)\qquad(x\in\mathbb{R}^3).
\]
\item[(c)] There is a continuous function on $\mathbb{R}^3\times\mathbb{R}^3$ that agrees with $\mathrm{Rk}^{(4)}$ at every point $(x,y)$ with $x\ne y$.
\end{enumerate}
\end{lemma}

\begin{proof}
Write $S:=S_{\mathsf{M}}$.

(a) For $n=0$ the assertion is the inequality $|S(x-y)|\le\mathsf{s}(x-y)$. Let $n\ge1$, and recall $z_0=x$. For $1\le i\le n$ the triangle inequality gives
\[
1+|z_i|\le1+|x|+\sum_{l=1}^i|z_l-z_{l-1}|\le(1+|x|)\cdot\prod_{l=1}^i\left(1+|z_l-z_{l-1}|\right).
\]
For $v\in\mathbb{R}^3$ the matrix $\gamma\cdot v$ is Hermitian with $(\gamma\cdot v)^2=|v|^2$ (Section~\ref{sec:letters-norms-dirac}); hence $|V(z)|=|a(z)|\le c_a\cdot(1+|z|)$. Multiplying these bounds over $i=1,\dots,n$, and observing that the factor with index $l$ occurs $n-l+1\le4$ times and is at least $1$, one obtains
\[
\prod_{i=1}^n|V(z_i)|\le\left(c_a\cdot(1+|x|)\right)^n\cdot\prod_{l=1}^n\left(1+|z_l-z_{l-1}|\right)^4.
\]
The operator norm is submultiplicative, and $|S(w)|\cdot(1+|w|)^4\le\mathsf{s}(w)$, $|S(w)|\le\mathsf{s}(w)$ for $w\ne0$. Assigning the factor $\left(1+|z_l-z_{l-1}|\right)^4$ to $S(z_{l-1}-z_l)$ for $1\le l\le n$, one finds that the integrand in the definition of $\mathrm{Rk}^{(n)}(x,y)$ has modulus at most
\[
\left(c_a\cdot(1+|x|)\right)^n\cdot\mathsf{s}(z_0-z_1)\cdot\mathsf{s}(z_1-z_2)\cdots\mathsf{s}(z_{n-1}-z_n)\cdot\mathsf{s}(z_n-y).
\]
By Fubini's theorem the integral of this function over $z_1,\dots,z_n$ equals $\left(c_a\cdot(1+|x|)\right)^n\cdot\mathsf{s}^{*(n+1)}(x-y)$, and this is finite for $x\ne y$ by Lemma~\ref{lem:convolution-powers}. This proves (a).

(b) By~\eqref{eq:free-coefficient-K2}, $S$ acts on Fourier transforms as the multiplication by
\[
(\mathsf{M}+i\gamma\cdot p)^{-1}=\frac{\mathsf{M}-i\gamma\cdot p}{\mathsf{M}^2+|p|^2},
\]
a smooth function of $p$ that is bounded together with all its derivatives. Multiplication by such a function maps $\mathscr{S}$ into itself, and the Fourier transform maps $\mathscr{S}$ onto itself; hence $S$ maps $\mathscr{S}$ into $\mathscr{S}$. The function $V$ is smooth, and each of its derivatives is bounded by a polynomial (it has compact support if $a$ is of type~\eqref{eq:free-coefficient-cpt}, and it is linear if $a$ is of type~\eqref{eq:free-coefficient-lin}); hence the multiplication by $V$ maps $\mathscr{S}$ into $\mathscr{S}$. For $\psi\in\mathscr{S}$ let $\hat\psi(p):=\int\mathrm{d}x\,\psi(x)\exp(-ip\cdot x)$ $(p\in\mathbb{R}^3)$ be its Fourier transform, normalised as in~\eqref{eq:free-coefficient-K2}. One has
\[
(S\psi)(x)=\int\mathrm{d}z\,S(x-z)\cdot\psi(z)\qquad\left(\psi\in\mathscr{S},\ x\in\mathbb{R}^3\right).
\]
Indeed, the right-hand side is an integrable and continuous function of $x$, because the kernel $S$ is integrable by~\eqref{eq:free-coefficient-K1} and $\psi$ is integrable, bounded and continuous. By~\eqref{eq:free-coefficient-K2} and the convolution theorem, the Fourier transform of the right-hand side is $p\mapsto(\mathsf{M}+i\gamma\cdot p)^{-1}\cdot\hat\psi(p)$. The Schwartz function $S\psi$ has the same Fourier transform, because $S$ acts on Fourier transforms as the multiplication by $(\mathsf{M}+i\gamma\cdot p)^{-1}$. Finally, two continuous integrable functions with the same Fourier transform are equal at every point.

Let $\psi_k:=(-1)^k(SV)^kS\boldsymbol{\phi}$ for $0\le k\le n$. These functions lie in $\mathscr{S}$; one has $\psi_0=S\boldsymbol{\phi}$, so that $\psi_0(x)=\int\mathrm{d}y\,\mathrm{Rk}^{(0)}(x,y)\cdot\boldsymbol{\phi}(y)$, and $\psi_k=-S(V\psi_{k-1})$ for $k\ge1$. Let $1\le k\le n$, and assume that $\psi_{k-1}(z)=\int\mathrm{d}y\,\mathrm{Rk}^{(k-1)}(z,y)\cdot\boldsymbol{\phi}(y)$ for all $z$. Then
\[
\psi_k(x)=-\int\mathrm{d}z_1\,S(x-z_1)\cdot V(z_1)\cdot\int\mathrm{d}y\,\mathrm{Rk}^{(k-1)}(z_1,y)\cdot\boldsymbol{\phi}(y).
\]
Written out, the right-hand side is an iterated integral, over $z_1$, then over $y$, then over the variables $z_2,\dots,z_k$ of $\mathrm{Rk}^{(k-1)}(z_1,y)$, of the integrand of $\mathrm{Rk}^{(k)}(x,y)$ multiplied by $\boldsymbol{\phi}(y)$, with the sign $(-1)^k$. By the bound in the proof of (a) and by Lemma~\ref{lem:convolution-powers}, the integral of the modulus of this function over $z_1,\dots,z_k,y$ is at most
\[
\left(c_a\cdot(1+|x|)\right)^k\int\mathrm{d}y\,\mathsf{s}^{*(k+1)}(x-y)\cdot|\boldsymbol{\phi}(y)|\le\left(c_a\cdot(1+|x|)\right)^k\cdot\|\mathsf{s}\|_{L^1}^{k+1}\cdot\sup|\boldsymbol{\phi}|<\infty.
\]
So Fubini's theorem allows one to integrate over $z_1,\dots,z_k$ first, and $\psi_k(x)=\int\mathrm{d}y\,\mathrm{Rk}^{(k)}(x,y)\cdot\boldsymbol{\phi}(y)$. For $k=n$ this is (b).

(c) For $(x,y)\in\mathbb{R}^3\times\mathbb{R}^3$ let
\begin{align*}
J(x,y):=\int&\mathrm{d}\boldsymbol{\zeta}_1\,\mathrm{d}\boldsymbol{\zeta}_2\,\mathrm{d}\boldsymbol{\zeta}_3\,\mathrm{d}\boldsymbol{\zeta}_5\;S(\boldsymbol{\zeta}_1)\cdot V(x-\boldsymbol{\zeta}_1)\cdot S(\boldsymbol{\zeta}_2)\cdot V(x-\boldsymbol{\zeta}_1-\boldsymbol{\zeta}_2)\cdot S(\boldsymbol{\zeta}_3)\\
&\cdot V(x-\boldsymbol{\zeta}_1-\boldsymbol{\zeta}_2-\boldsymbol{\zeta}_3)\cdot S(x-y-\boldsymbol{\zeta}_1-\boldsymbol{\zeta}_2-\boldsymbol{\zeta}_3-\boldsymbol{\zeta}_5)\cdot V(y+\boldsymbol{\zeta}_5)\cdot S(\boldsymbol{\zeta}_5).
\end{align*}
The substitution
\[
z_1=x-\boldsymbol{\zeta}_1,\qquad z_2=z_1-\boldsymbol{\zeta}_2,\qquad z_3=z_2-\boldsymbol{\zeta}_3,\qquad z_4=y+\boldsymbol{\zeta}_5
\]
has a Jacobian determinant of modulus $1$, and it transforms the integrand of $J(x,y)$ into the function $S(x-z_1)\cdot V(z_1)\cdot S(z_1-z_2)\cdot V(z_2)\cdot S(z_2-z_3)\cdot V(z_3)\cdot S(z_3-z_4)\cdot V(z_4)\cdot S(z_4-y)$, which for $x\ne y$ is the integrand of $\mathrm{Rk}^{(4)}(x,y)$. The bound for this function in the proof of (a) holds for every pair $(x,y)$, including $x=y$; it shows that the integral of its modulus is at most $\left(c_a\cdot(1+|x|)\right)^4\cdot\mathsf{s}^{*5}(x-y)$, and $\mathsf{s}^{*5}$ is finite everywhere by Lemma~\ref{lem:convolution-powers}. So $J(x,y)$ converges absolutely for every $(x,y)\in\mathbb{R}^3\times\mathbb{R}^3$, and $J=\mathrm{Rk}^{(4)}$ off the diagonal. Since the complement of the diagonal is dense in $\mathbb{R}^3\times\mathbb{R}^3$, it suffices to show that $J$ is continuous; $J$ is then the continuous extension of $\mathrm{Rk}^{(4)}$, and its value at a point $(x,x)$ is the absolutely convergent integral that defines $\mathrm{Rk}^{(4)}(x,y)$, taken at $y=x$.

Let $\varrho>0$, let $(x,y),(x',y')\in\mathbb{R}^3\times\mathbb{R}^3$ with $|x|,|y|,|x'|,|y'|\le\varrho$, and put $\xi:=x-y$, $\xi':=x'-y'$ and $\bar{\boldsymbol{\zeta}}:=\boldsymbol{\zeta}_1+\boldsymbol{\zeta}_2+\boldsymbol{\zeta}_3+\boldsymbol{\zeta}_5$. For $|u|\le\varrho$ and $w_1,\dots,w_k\in\mathbb{R}^3$ one has $\left|V\left(u+w_1+\dots+w_k\right)\right|\le c_a\cdot\left(1+\varrho+\sum_j|w_j|\right)\le c_a\cdot(1+\varrho)\cdot\prod_j(1+|w_j|)$. So each of the four factors $V$ in the integrand of $J(x,y)$ or of $J(x',y')$ is bounded by $c_a\cdot(1+\varrho)$ times the product of the numbers $1+|\boldsymbol{\zeta}_j|$ over the variables $\boldsymbol{\zeta}_j$ in its argument, and each $\boldsymbol{\zeta}_j$ occurs in the arguments of at most three of these factors. Write $\mathrm{d}\boldsymbol{\zeta}:=\mathrm{d}\boldsymbol{\zeta}_1\,\mathrm{d}\boldsymbol{\zeta}_2\,\mathrm{d}\boldsymbol{\zeta}_3\,\mathrm{d}\boldsymbol{\zeta}_5$.

Exactly five factors of the integrand of $J(x,y)$ involve $x$ or $y$; in the order of the product they are
\[
\Phi_1:=V(x-\boldsymbol{\zeta}_1),\quad\Phi_2:=V(x-\boldsymbol{\zeta}_1-\boldsymbol{\zeta}_2),\quad\Phi_3:=V(x-\boldsymbol{\zeta}_1-\boldsymbol{\zeta}_2-\boldsymbol{\zeta}_3),\quad\Phi_4:=S(\xi-\bar{\boldsymbol{\zeta}}),\quad\Phi_5:=V(y+\boldsymbol{\zeta}_5).
\]
Let $\Phi_k'$ arise from $\Phi_k$ when $(x,y)$ is replaced by $(x',y')$. For $k=1,\dots,5$ let $T_k$ denote the product obtained from the integrand of $J(x,y)$ by substituting $\Phi_k-\Phi_k'$ for $\Phi_k$ and $\Phi_l'$ for $\Phi_l$ whenever $l<k$. The sum $T_1+\dots+T_5$ telescopes: it equals the integrand of $J(x,y)$ minus the integrand of $J(x',y')$.

Let $k\in\{1,2,3,5\}$. Then $\Phi_k-\Phi_k'$ is a difference of two values of $V$ at points whose distance is $|x-x'|$ if $k\le3$, and $|y-y'|$ if $k=5$; as $|V(u)-V(u')|=|a(u)-a(u')|\le c_a\cdot|u-u'|$, its modulus does not exceed $c_a\cdot\left(|x-x'|+|y-y'|\right)$. Estimate the three other factors $V$ of $T_k$ as above, each factor $S(\boldsymbol{\zeta}_j)$ by means of $|S(w)|\cdot(1+|w|)^3\le\mathsf{s}(w)$, and the remaining propagator, which is $\Phi_4=S(\xi-\bar{\boldsymbol{\zeta}})$ if $k\le3$ and $\Phi_4'=S(\xi'-\bar{\boldsymbol{\zeta}})$ if $k=5$, by means of $|S|\le\mathsf{s}$. This gives
\[
|T_k|\le c_a^4\cdot(1+\varrho)^3\cdot\left(|x-x'|+|y-y'|\right)\cdot\mathsf{s}(\boldsymbol{\zeta}_1)\cdot\mathsf{s}(\boldsymbol{\zeta}_2)\cdot\mathsf{s}(\boldsymbol{\zeta}_3)\cdot\mathsf{s}(\boldsymbol{\zeta}_5)\cdot\mathsf{s}(\xi-\bar{\boldsymbol{\zeta}})
\]
for $k\le3$, and the same bound with $\xi'$ in place of $\xi$ for $k=5$. As $\int\mathrm{d}\boldsymbol{\zeta}\,\mathsf{s}(\boldsymbol{\zeta}_1)\cdot\mathsf{s}(\boldsymbol{\zeta}_2)\cdot\mathsf{s}(\boldsymbol{\zeta}_3)\cdot\mathsf{s}(\boldsymbol{\zeta}_5)\cdot\mathsf{s}(v-\bar{\boldsymbol{\zeta}})=\mathsf{s}^{*5}(v)$ for every $v\in\mathbb{R}^3$, it follows that $\int\mathrm{d}\boldsymbol{\zeta}\,|T_k|\le c_a^4\cdot(1+\varrho)^3\cdot\left(|x-x'|+|y-y'|\right)\cdot\sup\mathsf{s}^{*5}$ for each of these four values of $k$.

In $T_4$ the difference is taken in the propagator whose argument contains $x-y$: there $\Phi_4-\Phi_4'=S(\xi-\bar{\boldsymbol{\zeta}})-S(\xi'-\bar{\boldsymbol{\zeta}})$. Apart from $x,y,x',y'$, the four factors $V$ of $T_4$ depend on no variables other than $\boldsymbol{\zeta}_1,\boldsymbol{\zeta}_2,\boldsymbol{\zeta}_3,\boldsymbol{\zeta}_5$. Each of these variables is the argument of one of the factors $S(\boldsymbol{\zeta}_j)$; hence the growth of the factors $V$, estimated as above, is absorbed through $|S(w)|\cdot(1+|w|)^3\le\mathsf{s}(w)$. Since $\xi'-\bar{\boldsymbol{\zeta}}=(\xi-\bar{\boldsymbol{\zeta}})+(\xi'-\xi)$, it follows that
\[
\int\mathrm{d}\boldsymbol{\zeta}\,|T_4|\le c_a^4\cdot(1+\varrho)^4\cdot\left(\mathsf{s}^{*4}*\left|S-S(\cdot+\xi'-\xi)\right|\right)(\xi)\le c_a^4\cdot(1+\varrho)^4\cdot\sup\mathsf{s}^{*4}\cdot\left\|S-S(\cdot+\xi'-\xi)\right\|_{L^1}.
\]
The kernel $S$ is integrable by~\eqref{eq:free-coefficient-K1}; by the continuity of translations in $L^1(\mathbb{R}^3)$, $\left\|S-S(\cdot+\xi'-\xi)\right\|_{L^1}\to0$ as $\xi'\to\xi$. Since $|J(x,y)-J(x',y')|\le\sum_{k=1}^5\int\mathrm{d}\boldsymbol{\zeta}\,|T_k|$ and $(1+\varrho)^3\le(1+\varrho)^4$, altogether
\[
|J(x,y)-J(x',y')|\le c_a^4\cdot(1+\varrho)^4\cdot\left(4\left(|x-x'|+|y-y'|\right)\cdot\sup\mathsf{s}^{*5}+\sup\mathsf{s}^{*4}\cdot\left\|S-S(\cdot+\xi'-\xi)\right\|_{L^1}\right),
\]
and the right-hand side tends to $0$ as $(x',y')\to(x,y)$. So $J$ is continuous.
\end{proof}

For a field $a$ let $r_4(\cdot\,;\mathsf{M},a)\colon\mathbb{R}^3\to\mathbb{C}$ be the function
\[
r_4(x;\mathsf{M},a):=\operatorname{tr}\mathrm{Rk}^{(4)}(x,x),
\]
where $\mathrm{Rk}^{(4)}$ is the continuous extension of Lemma~\ref{lem:kernels}(c).

\subsubsection{The derivative in the mass}\label{sec:derivative-mass}

\begin{lemma}[Derivative of $q_{4,\mathsf{M}}(a)$ in the mass]\label{lem:derivative-mass}
Let $a$ be a field of type~\eqref{eq:free-coefficient-cpt}. Then the function $\mathsf{M}\mapsto q_{4,\mathsf{M}}(a)$ is continuously differentiable on $(0,\infty)$ and tends to $0$ as $\mathsf{M}\to\infty$. For every $\mathsf{M}>0$ the function $x\mapsto r_4(x;\mathsf{M},a)$ is integrable on $\mathbb{R}^3$, and
\[
\frac{\mathrm{d}}{\mathrm{d}\mathsf{M}}\,q_{4,\mathsf{M}}(a)=\int\mathrm{d}x\,r_4(x;\mathsf{M},a).
\]
\end{lemma}

\begin{proof}
Write $S:=S_{\mathsf{M}}$ and $G:=SVS$, and for $x,y\in\mathbb{R}^3$ with $x\ne y$ let $G(x,y):=-\mathrm{Rk}^{(1)}(x,y)$, that is,
\[
G(x,y)=\int\mathrm{d}z\,S(x-z)\cdot V(z)\cdot S(z-y).
\]
By Lemma~\ref{lem:hilbert-schmidt-bound}, $G$ is a Hilbert--Schmidt operator, $(SV)^4=GVGV$ is of trace class, so that $q_{4,\mathsf{M}}(a)=-\frac{1}{4}\operatorname{Tr}(G\cdot VGV)$ by the definition in Section~\ref{sec:free-fourth-order}, and
\[
|q_{4,\mathsf{M}}(a)|\le\frac{\|a\|_{L^2}^2\cdot\|a\|_{\infty}^2}{4\pi\cdot\mathsf{M}};
\]
in particular $q_{4,\mathsf{M}}(a)\to0$ as $\mathsf{M}\to\infty$.

The function $G(x,y)$ is the square-integrable kernel of the Hilbert--Schmidt operator $G$. Indeed, the integrand $S(x-z)\cdot V(z)\cdot S(z-y)$, defined for $z\notin\{x,y\}$, is a measurable function of $(x,y,z)$, and for $x\ne y$ it is integrable in $z$ by Lemma~\ref{lem:kernels}(a); by Fubini's theorem, $G(x,y)$ is therefore a measurable function on the complement of the diagonal, which is a null set in $\mathbb{R}^3\times\mathbb{R}^3$. As the operator norm is submultiplicative, $|S|\le\mathsf{s}$ and $|V(z)|=|a(z)|$,
\[
|G(x,y)|\le\int\mathrm{d}z\,\mathsf{s}(x-z)\cdot|a(z)|\cdot\mathsf{s}(z-y)\qquad(x\ne y).
\]
Since $\mathsf{s}$ depends on $|z|$ only, $\mathsf{s}(-w)=\mathsf{s}(w)$ for $w\ne0$; the substitution $x=z'+w$ in the first of the following integrals, and $y=z'+w$ in the last one, gives
\[
\int\mathrm{d}x\,\mathsf{s}(x-z)\cdot\mathsf{s}(x-z')=\mathsf{s}^{*2}(z-z')=\int\mathrm{d}y\,\mathsf{s}(z-y)\cdot\mathsf{s}(z'-y)\qquad(z,z'\in\mathbb{R}^3).
\]
Bound $|G(x,y)|^2$ by the product of two copies of the integral in the bound for $|G(x,y)|$, one with the variable $z$ and one with the variable $z'$, and integrate over $x$ and $y$ first (Fubini's theorem for non-negative functions). By the last display, this gives
\[
\int\mathrm{d}x\int\mathrm{d}y\,|G(x,y)|^2\le\int\mathrm{d}z\int\mathrm{d}z'\,|a(z)|\cdot|a(z')|\cdot\left(\mathsf{s}^{*2}(z-z')\right)^2\le\|a\|_{L^2}^2\cdot\left\|\mathsf{s}^{*2}\right\|_{L^2}^2.
\]
The second inequality is the Cauchy--Schwarz inequality in $z$ together with Young's inequality $\|g*u\|_{L^2}\le\|g\|_{L^1}\cdot\|u\|_{L^2}$ for $g:=\left(\mathsf{s}^{*2}\right)^2$ and $u:=|a|$, since $\|g\|_{L^1}=\left\|\mathsf{s}^{*2}\right\|_{L^2}^2$ and $\|u\|_{L^2}=\|a\|_{L^2}$. The right-hand side is finite: $\|a\|_{L^2}<\infty$ because $a$ is of type~\eqref{eq:free-coefficient-cpt}, and $\mathsf{s}^{*2}\in L^2(\mathbb{R}^3)$ by the first inequality of Lemma~\ref{lem:convolution-powers}, as shown in its proof. Every $4\times4$ matrix has Hilbert--Schmidt norm at most twice its operator norm, and the squared Hilbert--Schmidt norm of an integral operator is the integral of the squared Hilbert--Schmidt norm of its kernel (proof of Lemma~\ref{lem:hilbert-schmidt-bound}); hence the integral operator $\tilde{G}$ on $L^2(\mathbb{R}^3;\mathbb{C}^4)$ with the kernel $G(x,y)$ is a Hilbert--Schmidt operator. By Lemma~\ref{lem:kernels}(b) with $n=1$, every $\boldsymbol{\phi}\in\mathscr{S}$ satisfies $(G\boldsymbol{\phi})(x)=\int\mathrm{d}y\,G(x,y)\cdot\boldsymbol{\phi}(y)$ for all $x\in\mathbb{R}^3$, so that $\tilde{G}\boldsymbol{\phi}=G\boldsymbol{\phi}$. The operators $\tilde{G}$ and $G$ are bounded, the latter as the product of the bounded operators $S$, $V$, $S$, and $\mathscr{S}$ is dense in $L^2(\mathbb{R}^3;\mathbb{C}^4)$; therefore $\tilde{G}=G$.

For Hilbert--Schmidt operators $X$, $Y$ on $L^2(\mathbb{R}^3;\mathbb{C}^4)$ with the square-integrable kernels $X(x,y)$, $Y(x,y)$ one has $\operatorname{Tr}(XY)=\int\mathrm{d}x\int\mathrm{d}y\,\operatorname{tr}\left(X(x,y)\cdot Y(y,x)\right)$. This applies to $X=G$, with the kernel $G(x,y)$, and to the Hilbert--Schmidt operator $Y=VGV$, whose kernel is $V(x)\cdot G(x,y)\cdot V(y)$; this function is square integrable, since $|V(x)|=|a(x)|\le\|a\|_{\infty}$ for every $x$. Inserting the kernel of $G$ twice, calling $z_4:=x$, $z_2:=y$ and $z_1$, $z_3$ the two variables of integration in the kernels of $G$, and using Fubini's theorem, one obtains
\begin{align*}
q_{4,\mathsf{M}}(a)=-\frac{1}{4}\int&\mathrm{d}z_1\,\mathrm{d}z_2\,\mathrm{d}z_3\,\mathrm{d}z_4\,\operatorname{tr}\big(S(z_4-z_1)\cdot V(z_1)\cdot S(z_1-z_2)\cdot V(z_2)\\
&\cdot S(z_2-z_3)\cdot V(z_3)\cdot S(z_3-z_4)\cdot V(z_4)\big),
\end{align*}
provided this integral converges absolutely. To verify this, note that $|\operatorname{tr}M|\le4|M|$ for every $4\times4$ matrix $M$, $|V(z)|=|a(z)|\le\|a\|_{\infty}$ and $|S|\le\mathsf{s}$; hence the integrand is dominated by the non-negative function
\[
4\|a\|_{\infty}^3\cdot|a(z_4)|\cdot\mathsf{s}(z_4-z_1)\cdot\mathsf{s}(z_1-z_2)\cdot\mathsf{s}(z_2-z_3)\cdot\mathsf{s}(z_3-z_4).
\]
Integrating this function first over $z_1,z_2,z_3$ and then over $z_4$, one obtains $4\|a\|_{\infty}^3\cdot\|a\|_{L^1}\cdot\mathsf{s}^{*4}(0)$. This number is finite: $\|a\|_{L^1}=\int\mathrm{d}z\,|a(z)|$ is finite because $a$ is continuous and has compact support, and $\mathsf{s}^{*4}(0)$ is finite by Lemma~\ref{lem:convolution-powers}.

For $z\ne0$ let $\mathsf{r}:=|z|$ and $\hat{z}:=z/\mathsf{r}$. By~\eqref{eq:free-coefficient-K1},
\[
S_{\mathsf{M}}(z)=\frac{\exp(-\mathsf{M}\cdot\mathsf{r})}{4\pi}\left(\frac{\mathsf{M}}{\mathsf{r}}+\frac{\mathsf{M}\cdot\mathsf{r}+1}{\mathsf{r}^2}\,\gamma\cdot\hat{z}\right),
\]
and differentiation with respect to $\mathsf{M}$ gives
\[
\partial_{\mathsf{M}}S_{\mathsf{M}}(z)=\frac{\exp(-\mathsf{M}\cdot\mathsf{r})}{4\pi}\left(\frac{1}{\mathsf{r}}-\mathsf{M}-\mathsf{M}\,\gamma\cdot\hat{z}\right).
\]
As $|\gamma\cdot\hat{z}|=1$, it follows that $|\partial_{\mathsf{M}}S_{\mathsf{M}}(z)|\le(4\pi)^{-1}\cdot\exp(-\mathsf{M}\cdot\mathsf{r})\cdot\left(\mathsf{r}^{-1}+2\mathsf{M}\right)$. Let $0<\mathsf{M}_1<\mathsf{M}_2$. Together with the bound for $|S_{\mathsf{M}}(z)|$ in~\eqref{eq:free-coefficient-K1} this gives, for $\mathsf{M}_1\le\mathsf{M}\le\mathsf{M}_2$ and $z\ne0$,
\[
|S_{\mathsf{M}}(z)|+|\partial_{\mathsf{M}}S_{\mathsf{M}}(z)|\le\frac{\exp(-\mathsf{M}_1\cdot\mathsf{r})}{4\pi\cdot\mathsf{r}^2}\left(1+(2\mathsf{M}_2+1)\cdot\mathsf{r}+2\mathsf{M}_2\cdot\mathsf{r}^2\right)\le\mathsf{s}'(z),
\]
where $\mathsf{s}'$ is a function of the form considered in Lemma~\ref{lem:convolution-powers}, with $c=\frac{\mathsf{M}_1}{2}$ and with $\boldsymbol{\kappa}$ the supremum over $\mathsf{r}\ge0$ of $(4\pi)^{-1}\cdot\exp\left(-\frac{\mathsf{M}_1\cdot\mathsf{r}}{2}\right)\cdot\left(1+(2\mathsf{M}_2+1)\cdot\mathsf{r}+2\mathsf{M}_2\cdot\mathsf{r}^2\right)$.

In particular $\partial_{\mathsf{M}}S_{\mathsf{M}}$ is integrable, and differentiation under the integral sign, justified by this bound, in the identity $\int\mathrm{d}z\,S_{\mathsf{M}}(z)\exp(-ip\cdot z)=(\mathsf{M}+i\gamma\cdot p)^{-1}$ of~\eqref{eq:free-coefficient-K2} shows that the Fourier transform of $\partial_{\mathsf{M}}S_{\mathsf{M}}$ is
\[
\partial_{\mathsf{M}}\left(\mathsf{M}+i\gamma\cdot p\right)^{-1}=-\left(\mathsf{M}+i\gamma\cdot p\right)^{-2}.
\]
Let $(S_{\mathsf{M}}*S_{\mathsf{M}})(z):=\int\mathrm{d}w\,S_{\mathsf{M}}(z-w)\cdot S_{\mathsf{M}}(w)$. The Fourier transform maps the convolution of integrable functions to the product of their transforms; so $-(\mathsf{M}+i\gamma\cdot p)^{-2}$ is also the Fourier transform of $-S_{\mathsf{M}}*S_{\mathsf{M}}$. As the Fourier transform is injective on integrable functions, $\partial_{\mathsf{M}}S_{\mathsf{M}}=-S_{\mathsf{M}}*S_{\mathsf{M}}$ almost everywhere.

For $\mathsf{M}_1\le\mathsf{M}\le\mathsf{M}_2$ the integrand of the fourfold integral for $q_{4,\mathsf{M}}(a)$, and its derivative with respect to $\mathsf{M}$, which is the sum of the four terms in which one of the factors $S$ is replaced by $\partial_{\mathsf{M}}S_{\mathsf{M}}$, are bounded in modulus by
\[
16\|a\|_{\infty}^3\cdot|a(z_4)|\cdot\mathsf{s}'(z_4-z_1)\cdot\mathsf{s}'(z_1-z_2)\cdot\mathsf{s}'(z_2-z_3)\cdot\mathsf{s}'(z_3-z_4).
\]
This function does not depend on $\mathsf{M}$, and its integral $16\|a\|_{\infty}^3\cdot\|a\|_{L^1}\cdot(\mathsf{s}')^{*4}(0)$ is finite by Lemma~\ref{lem:convolution-powers}. At every point at which the four differences $z_4-z_1$, $z_1-z_2$, $z_2-z_3$, $z_3-z_4$ are non-zero, that is, outside a null set, the integrand is continuously differentiable in $\mathsf{M}$. By the theorem on differentiation under the integral sign and by dominated convergence, $\mathsf{M}\mapsto q_{4,\mathsf{M}}(a)$ is continuously differentiable on $(\mathsf{M}_1,\mathsf{M}_2)$, and its derivative is $-\frac{1}{4}$ times the integral of the derivative of the integrand. Since $\mathsf{M}_1<\mathsf{M}_2$ are arbitrary, the function is continuously differentiable on $(0,\infty)$.

The four terms have the same integral: the renaming $(z_1,z_2,z_3,z_4)\mapsto(z_2,z_3,z_4,z_1)$ of the variables of integration preserves the Lebesgue measure and, by the cyclicity of $\operatorname{tr}$, carries the term in which $S(z_4-z_1)$ is differentiated into the term in which $S(z_1-z_2)$ is differentiated, this one into the term in which $S(z_2-z_3)$ is differentiated, and that one into the term in which $S(z_3-z_4)$ is differentiated. The factor $-\frac{1}{4}$, the number $4$ of equal terms and the sign in $\partial_{\mathsf{M}}S_{\mathsf{M}}=-S*S$ multiply to $1$. With $\mathrm{d}z:=\mathrm{d}z_1\,\mathrm{d}z_2\,\mathrm{d}z_3\,\mathrm{d}z_4$ it follows that
\begin{align*}
\frac{\mathrm{d}}{\mathrm{d}\mathsf{M}}\,q_{4,\mathsf{M}}(a)=\int&\mathrm{d}z\,\operatorname{tr}\big((S*S)(z_4-z_1)\cdot V(z_1)\cdot S(z_1-z_2)\cdot V(z_2)\\
&\cdot S(z_2-z_3)\cdot V(z_3)\cdot S(z_3-z_4)\cdot V(z_4)\big).
\end{align*}
Insert $(S*S)(z_4-z_1)=\int\mathrm{d}x\,S(z_4-x)\cdot S(x-z_1)$. The fivefold integral obtained in this way is absolutely convergent, since its integrand is dominated by
\[
4\|a\|_{\infty}^3\cdot|a(z_4)|\cdot\mathsf{s}(z_4-x)\cdot\mathsf{s}(x-z_1)\cdot\mathsf{s}(z_1-z_2)\cdot\mathsf{s}(z_2-z_3)\cdot\mathsf{s}(z_3-z_4);
\]
integrating this function over $x,z_1,z_2,z_3$ and then over $z_4$, one obtains $4\|a\|_{\infty}^3\cdot\|a\|_{L^1}\cdot\mathsf{s}^{*5}(0)$, which is finite by Lemma~\ref{lem:convolution-powers}. By Fubini's theorem and the cyclicity of $\operatorname{tr}$,
\begin{align*}
\frac{\mathrm{d}}{\mathrm{d}\mathsf{M}}\,q_{4,\mathsf{M}}(a)=\int\mathrm{d}x\int&\mathrm{d}z\,\operatorname{tr}\big(S(x-z_1)\cdot V(z_1)\cdot S(z_1-z_2)\cdot V(z_2)\cdot S(z_2-z_3)\\
&\cdot V(z_3)\cdot S(z_3-z_4)\cdot V(z_4)\cdot S(z_4-x)\big),
\end{align*}
where the inner integral converges absolutely for every $x$ (the integral of the modulus of its integrand is at most $4\|a\|_{\infty}^4\cdot\mathsf{s}^{*5}(0)$), and its modulus is integrable in $x$. For every $x$ the inner integral is the trace of the integral that defines $\mathrm{Rk}^{(4)}(x,y)$, taken at $y=x$ (the sign is $(-1)^4=1$). By the proof of Lemma~\ref{lem:kernels}(c) this integral is the value at $(x,x)$ of the continuous extension of $\mathrm{Rk}^{(4)}$; so the inner integral is $r_4(x;\mathsf{M},a)$. Hence $x\mapsto r_4(x;\mathsf{M},a)$ is integrable, and its integral is $\frac{\mathrm{d}}{\mathrm{d}\mathsf{M}}\,q_{4,\mathsf{M}}(a)$.
\end{proof}

\subsubsection{The expansion for large mass}\label{sec:expansion-large-mass}

Let $a$ be a field in the sense of Section~\ref{sec:coefficient-fourth-order}, let $x\in\mathbb{R}^3$, and let $j\ge0$ be an integer. For a multi-index $\boldsymbol{\alpha}=(\boldsymbol{\alpha}_1,\boldsymbol{\alpha}_2,\boldsymbol{\alpha}_3)\in\{0,1,2,\dots\}^3$ and for $w\in\mathbb{R}^3$ write $|\boldsymbol{\alpha}|:=\boldsymbol{\alpha}_1+\boldsymbol{\alpha}_2+\boldsymbol{\alpha}_3$, $\boldsymbol{\alpha}!:=\boldsymbol{\alpha}_1!\,\boldsymbol{\alpha}_2!\,\boldsymbol{\alpha}_3!$, $\partial^{\boldsymbol{\alpha}}:=\partial_1^{\boldsymbol{\alpha}_1}\partial_2^{\boldsymbol{\alpha}_2}\partial_3^{\boldsymbol{\alpha}_3}$ and $w^{\boldsymbol{\alpha}}:=w_1^{\boldsymbol{\alpha}_1}w_2^{\boldsymbol{\alpha}_2}w_3^{\boldsymbol{\alpha}_3}$. Define the $4\times4$ matrices
\[
V_j(x;w):=i\gamma\cdot\sum_{|\boldsymbol{\alpha}|=j}\frac{w^{\boldsymbol{\alpha}}}{\boldsymbol{\alpha}!}\,(\partial^{\boldsymbol{\alpha}}a)(x)\qquad(w\in\mathbb{R}^3)
\]
and the complex numbers
\begin{multline*}
\rho_j(x;a):=\sum_{j_1+j_2+j_3+j_4=j}\operatorname{tr}\int_{(\mathbb{R}^3)^4}\mathrm{d}w_1\,\mathrm{d}w_2\,\mathrm{d}w_3\,\mathrm{d}w_4\;S_1(-w_1)\,V_{j_1}(x;w_1)\,S_1(w_1-w_2)\,V_{j_2}(x;w_2)\\
S_1(w_2-w_3)\,V_{j_3}(x;w_3)\,S_1(w_3-w_4)\,V_{j_4}(x;w_4)\,S_1(w_4),
\end{multline*}
where the sum extends over the quadruples of integers $j_1,j_2,j_3,j_4\ge0$ with sum $j$, and $S_1$ is the kernel $S_{\mathsf{M}}$ of Section~\ref{sec:free-fourth-order} for $\mathsf{M}=1$. These integrals converge absolutely: $|V_{j_i}(x;w_i)|$ is at most a multiple of $|w_i|^{j_i}$, and the estimate of Lemma~\ref{lem:kernels}(a) applies with $x=0$ and with the weights $|w_i|^{j_i}$ in place of the factors $|V(z_i)|$. Each factor $V_{j_i}(x;w_i)$ is linear in the derivatives of order $j_i$ of $a$ at $x$; consequently $\rho_j(x;a)$ is a homogeneous polynomial of degree four in the derivatives of $a$ at $x$ whose orders do not exceed $j$.

Let $F$ be a real antisymmetric $3\times3$ matrix. The field $A_F$ of type~\eqref{eq:free-coefficient-lin} satisfies $A_F(0)=0$ and $\partial_{\nu}A_{F,\mu}=\frac{1}{2}F_{\nu\mu}$, and its derivatives of order at least two vanish. Hence, for the field $a=A_F$,
\[
V_1(0;w)=i\gamma\cdot A_F(w)\quad(w\in\mathbb{R}^3),\qquad V_j(0;\cdot)=0\quad(j\ne1).
\]
Only the quadruple $(1,1,1,1)$ contributes to $\rho_4(0;A_F)$, and its contribution is the trace of the absolutely convergent integral of Section~\ref{sec:kernels-products-propagators} for $\mathrm{Rk}^{(4)}(x,y)$, formed with $\mathsf{M}=1$ and the field $A_F$ and taken at $x=y=0$ (proof of Lemma~\ref{lem:kernels}(c)). Define
\[
\Psi(F):=\rho_4(0;A_F)=r_4(0;1,A_F).
\]

\begin{lemma}[Expansion in powers of the inverse mass]\label{lem:expansion-powers-inverse}
Let $a$ be a field of type~\eqref{eq:free-coefficient-cpt}. There is a number $C'_a<\infty$ such that
\[
\left|r_4(x;\mathsf{M},a)-\sum_{j=0}^5\mathsf{M}^{-2-j}\cdot\rho_j(x;a)\right|\le C'_a\cdot\mathsf{M}^{-8}\qquad\text{for all }x\in\mathbb{R}^3\text{ and all }\mathsf{M}\ge1;
\]
$\rho_j(x;a)=0$ for every $j\ge0$ and every $x\in\mathbb{R}^3$ outside $\operatorname{supp}a$; and
\[
\int_{\{x\in\mathbb{R}^3:\ \operatorname{dist}(x,\operatorname{supp}a)\ge1\}}\mathrm{d}x\,\left|r_4(x;\mathsf{M},a)\right|\le C'_a\cdot\exp(-\mathsf{M})\qquad\text{for all }\mathsf{M}\ge1.
\]
\end{lemma}

\begin{proof}
By~\eqref{eq:free-coefficient-K1}, $S_{\mathsf{M}}(z)=\mathsf{M}^2\cdot S_1(\mathsf{M}\cdot z)$. The number $r_4(x;\mathsf{M},a)=\operatorname{tr}\mathrm{Rk}^{(4)}(x,x)$ is the trace of the absolutely convergent integral of Section~\ref{sec:kernels-products-propagators} for $\mathrm{Rk}^{(4)}(x,y)$, taken at $y=x$ (proof of Lemma~\ref{lem:kernels}(c)). In this integral substitute $z_i=x+\mathsf{M}^{-1}w_i$ for $i=1,\dots,4$. The five factors $S_{\mathsf{M}}$ produce the factor $\mathsf{M}^{10}$ and the substitution the factor $\mathsf{M}^{-12}$; with $V=i\gamma\cdot a$ and $\mathrm{d}w:=\mathrm{d}w_1\,\mathrm{d}w_2\,\mathrm{d}w_3\,\mathrm{d}w_4$ one obtains
\begin{multline*}
r_4(x;\mathsf{M},a)=\mathsf{M}^{-2}\cdot\operatorname{tr}\int_{(\mathbb{R}^3)^4}\mathrm{d}w\;S_1(-w_1)\,V\left(x+\frac{w_1}{\mathsf{M}}\right)S_1(w_1-w_2)\,V\left(x+\frac{w_2}{\mathsf{M}}\right)\\
S_1(w_2-w_3)\,V\left(x+\frac{w_3}{\mathsf{M}}\right)S_1(w_3-w_4)\,V\left(x+\frac{w_4}{\mathsf{M}}\right)S_1(w_4).
\end{multline*}
Taylor's formula at $x$, with the remainder of order six, gives for $w\in\mathbb{R}^3$ and $\mathsf{M}>0$
\[
V\left(x+\frac{w}{\mathsf{M}}\right)=\sum_{j=0}^5\mathsf{M}^{-j}\cdot V_j(x;w)+\mathsf{M}^{-6}\cdot\check{E}(x,w,\mathsf{M}),
\]
an identity that defines the $4\times4$ matrix $\check{E}(x,w,\mathsf{M})$. Let $|a|_{C^6}$ be the largest of the sup norms of the derivatives of order at most six of the components of $a$ (as for two-forms in Section~\ref{sec:lattices-tests-function}). There is an absolute constant $c_1$ such that
\[
|V_j(x;w)|\le c_1\cdot|a|_{C^6}\cdot(1+|w|)^6\quad(0\le j\le5),\qquad|\check{E}(x,w,\mathsf{M})|\le c_1\cdot|a|_{C^6}\cdot(1+|w|)^6
\]
for all $x,w\in\mathbb{R}^3$ and $\mathsf{M}>0$.

Insert this expansion into each of the four factors $V(x+\mathsf{M}^{-1}w_i)$ and multiply out. The terms of total order $\mathsf{M}^{-j}$ with $j\le5$ contain no factor $\check{E}$; they are the terms with the factors $\mathsf{M}^{-j_i}\cdot V_{j_i}(x;w_i)$, $j_1+j_2+j_3+j_4=j$, and by the definition of $\rho_j$ their contributions add up to $\mathsf{M}^{-2-j}\cdot\rho_j(x;a)$. Each of the finitely many other terms has total order $\mathsf{M}^{-k}$ with $k\ge6$; for $\mathsf{M}\ge1$ the product of its four factors has norm at most
\[
\mathsf{M}^{-6}\cdot c_1^4\cdot|a|_{C^6}^4\cdot\prod_{i=1}^4(1+|w_i|)^6.
\]
With $w_0:=0$ one has $1+|w_i|\le\prod_{l=1}^i(1+|w_l-w_{l-1}|)$, and therefore
\[
\prod_{i=1}^4(1+|w_i|)^6\le\prod_{l=1}^4(1+|w_l-w_{l-1}|)^{24}.
\]
Let $\mathsf{s}$ be the majorant of Section~\ref{sec:kernels-products-propagators} formed with $\mathsf{M}=1$, so that $|S_1(z)|\cdot(1+|z|)^{24}\le\mathsf{s}(z)$. As in the proof of Lemma~\ref{lem:kernels}(a), and since $|\operatorname{tr}M|\le4\,|M|$ for a $4\times4$ matrix $M$, the contribution of each of these terms to the integral has modulus at most $4\,c_1^4\cdot|a|_{C^6}^4\cdot\mathsf{M}^{-6}\cdot\mathsf{s}^{*5}(0)$, and $\mathsf{s}^{*5}(0)<\infty$ by Lemma~\ref{lem:convolution-powers}. Together with the factor $\mathsf{M}^{-2}$ in front of the integral this proves the first assertion.

If $x\notin\operatorname{supp}a$, then $a$ vanishes on a neighbourhood of $x$; all derivatives of $a$ vanish at $x$, so $V_j(x;\cdot)=0$ for every $j$, and $\rho_j(x;a)=0$.

Now let $\mathfrak{d}:=\operatorname{dist}(x,\operatorname{supp}a)\ge1$. The integrand of $r_4(x;\mathsf{M},a)=\operatorname{tr}\mathrm{Rk}^{(4)}(x,x)$ vanishes unless $z_1,\dots,z_4\in\operatorname{supp}a$, and then $|x-z_1|\ge\mathfrak{d}$ and $|z_4-x|\ge\mathfrak{d}$. By~\eqref{eq:free-coefficient-K1}, $|S_{\mathsf{M}}(z)|\le\mathsf{M}\cdot\exp(-\mathsf{M}\cdot|z|)$ for $|z|\ge1$ and $\mathsf{M}\ge1$, and $\|S_{\mathsf{M}}\|_{L^1}\le\frac{3}{\mathsf{M}}$. As $|V(z)|=|a(z)|\le\|a\|_{\infty}$, it follows that
\[
|r_4(x;\mathsf{M},a)|\le4\,\|a\|_{\infty}^4\cdot\mathsf{M}^2\cdot\exp(-2\mathsf{M}\cdot\mathfrak{d})\cdot|\operatorname{supp}a|\cdot\left(\frac{3}{\mathsf{M}}\right)^3,
\]
where $|\operatorname{supp}a|$ is the Lebesgue measure of $\operatorname{supp}a$. For $\mathsf{M}\ge1$ and $\mathfrak{d}\ge1$ one has $2\mathsf{M}\cdot\mathfrak{d}\ge\mathsf{M}+\mathfrak{d}$, the difference being $\mathsf{M}\cdot(\mathfrak{d}-1)+\mathfrak{d}\cdot(\mathsf{M}-1)$; so the right side is at most $108\,\|a\|_{\infty}^4\cdot|\operatorname{supp}a|\cdot\exp(-\mathsf{M})\cdot\exp(-\mathfrak{d})$. Since $\operatorname{supp}a$ is compact, the function $x\mapsto\exp\left(-\operatorname{dist}(x,\operatorname{supp}a)\right)$ is integrable over $\mathbb{R}^3$, and the last assertion follows. The lemma holds with $C'_a$ the larger of the two constants obtained.
\end{proof}

\subsubsection{Gauge invariance and the coefficients of the expansion}\label{sec:gauge-invariance-coefficients}

\begin{lemma}[Gauge invariance]\label{lem:gauge-invariance}
Let $a$ be a field of type~\eqref{eq:free-coefficient-cpt}, let $\boldsymbol{\lambda}\in C_c^{\infty}(\mathbb{R}^3)$ be real, and $d\boldsymbol{\lambda}:=(\partial_1\boldsymbol{\lambda},\partial_2\boldsymbol{\lambda},\partial_3\boldsymbol{\lambda})$. Then
\[
r_4(x;\mathsf{M},a+d\boldsymbol{\lambda})=r_4(x;\mathsf{M},a)\qquad\text{for all }x\in\mathbb{R}^3\text{ and all }\mathsf{M}>0.
\]
It follows that $\rho_j(x;a+d\boldsymbol{\lambda})=\rho_j(x;a)$ for all $x\in\mathbb{R}^3$ and all $j\le5$.
\end{lemma}

\begin{proof}
Let $\mathsf{E}$ be the ring of all $\mathbb{C}$-linear maps of the Schwartz space $\mathscr{S}$ of Section~\ref{sec:coefficient-fourth-order} into itself, and $\mathsf{E}[[u]]$ the ring of formal power series in an indeterminate $u$ with coefficients in $\mathsf{E}$. For a field $\tilde a$ of type~\eqref{eq:free-coefficient-cpt} let $V_{\tilde a}:=i\gamma\cdot\tilde a$, acting by multiplication, and
\[
\mathsf{D}_u(\tilde a):=\mathsf{M}+\gamma\cdot\partial+u\,i\gamma\cdot\tilde a\in\mathsf{E}[[u]].
\]
The operators $S_{\mathsf{M}}$ and $V_{\tilde a}$ map $\mathscr{S}$ into $\mathscr{S}$ (Lemma~\ref{lem:kernels}(b)), and on $\mathscr{S}$ the operator $S_{\mathsf{M}}$ is the inverse of $\mathsf{M}+\gamma\cdot\partial$ (Section~\ref{sec:free-fourth-order}). So the constant term of $\mathsf{D}_u(\tilde a)$ is invertible in $\mathsf{E}$; hence $\mathsf{D}_u(\tilde a)$ is invertible in $\mathsf{E}[[u]]$, and
\[
\mathsf{D}_u(\tilde a)^{-1}=\sum_{n\ge0}u^n\,\mathfrak{R}^{(n)}_{\tilde a},\qquad\mathfrak{R}^{(n)}_{\tilde a}:=(-1)^n\left(S_{\mathsf{M}}V_{\tilde a}\right)^nS_{\mathsf{M}}.
\]
The two series
\[
\exp(\pm iu\boldsymbol{\lambda}):=\sum_{j\ge0}\frac{(\pm iu\boldsymbol{\lambda})^j}{j!},
\]
whose coefficients act on $\mathscr{S}$ as multiplication by the functions $(\pm i\boldsymbol{\lambda})^j/j!$, are elements of $\mathsf{E}[[u]]$ that are inverse to each other. By Leibniz' rule,
\[
\gamma\cdot\partial\,(\boldsymbol{\lambda}^j\boldsymbol{\phi})=\boldsymbol{\lambda}^j\,\gamma\cdot\partial\boldsymbol{\phi}+j\,\boldsymbol{\lambda}^{j-1}\,(\gamma\cdot d\boldsymbol{\lambda})\,\boldsymbol{\phi}\qquad(\boldsymbol{\phi}\in\mathscr{S},\ j\ge1);
\]
hence $\gamma\cdot\partial\circ\exp(iu\boldsymbol{\lambda})=\exp(iu\boldsymbol{\lambda})\circ\left(\gamma\cdot\partial+iu\,\gamma\cdot d\boldsymbol{\lambda}\right)$ in $\mathsf{E}[[u]]$, and
\[
\exp(-iu\boldsymbol{\lambda})\,\mathsf{D}_u(a)\,\exp(iu\boldsymbol{\lambda})=\mathsf{D}_u(a+d\boldsymbol{\lambda}).
\]
By the uniqueness of inverses in a ring,
\[
\sum_{n\ge0}u^n\,\mathfrak{R}^{(n)}_{a+d\boldsymbol{\lambda}}=\exp(-iu\boldsymbol{\lambda})\left(\sum_{n\ge0}u^n\,\mathfrak{R}^{(n)}_a\right)\exp(iu\boldsymbol{\lambda}).
\]
Let $\mathrm{Rk}^{(n)}_a$ and $\mathrm{Rk}^{(n)}_{a+d\boldsymbol{\lambda}}$ denote the kernels $\mathrm{Rk}^{(n)}$ of Section~\ref{sec:kernels-products-propagators} formed with the fields $a$ and $a+d\boldsymbol{\lambda}$. Take the coefficients of $u^4$ on both sides and apply them to $\boldsymbol{\phi}\in\mathscr{S}$. By Lemma~\ref{lem:kernels}(b) and the binomial formula, for every $x\in\mathbb{R}^3$,
\[
\int\mathrm{d}y\,\mathrm{Rk}^{(4)}_{a+d\boldsymbol{\lambda}}(x,y)\,\boldsymbol{\phi}(y)=\int\mathrm{d}y\,\sum_{n=0}^4\mathrm{Rk}^{(n)}_a(x,y)\,\frac{\left(-i\left(\boldsymbol{\lambda}(x)-\boldsymbol{\lambda}(y)\right)\right)^{4-n}}{(4-n)!}\,\boldsymbol{\phi}(y).
\]
For fixed $x$ both kernels are locally integrable functions of $y$ (Lemma~\ref{lem:kernels}(a)), and $\boldsymbol{\phi}\in\mathscr{S}$ is arbitrary; consequently, for each $x$, the two kernels coincide as functions of $y$ outside a null set. Let $\mathsf{v}:=x-y$ with $|\mathsf{v}|\le1$. By Lemma~\ref{lem:kernels}(a), Lemma~\ref{lem:convolution-powers} and the bound $|\boldsymbol{\lambda}(x)-\boldsymbol{\lambda}(y)|\le\sup|d\boldsymbol{\lambda}|\cdot|\mathsf{v}|$, there is a number $c_2$, independent of $y$, such that
\[
\left|\mathrm{Rk}^{(n)}_a(x,y)\right|\cdot|\boldsymbol{\lambda}(x)-\boldsymbol{\lambda}(y)|^{4-n}\le c_2\cdot
\begin{cases}
|\mathsf{v}|^{-2+4}, & n=0,\\
|\mathsf{v}|^{-1+3}, & n=1,\\
|\mathsf{v}|^{-\frac{1}{2}+2}, & n=2,\\
|\mathsf{v}|, & n=3.
\end{cases}
\]
These bounds tend to $0$ as $y\to x$. Letting $y$ approach $x$ within the complement of that null set, and using the continuity of $\mathrm{Rk}^{(4)}_{a+d\boldsymbol{\lambda}}$ and $\mathrm{Rk}^{(4)}_a$ (Lemma~\ref{lem:kernels}(c)), one obtains $\mathrm{Rk}^{(4)}_{a+d\boldsymbol{\lambda}}(x,x)=\mathrm{Rk}^{(4)}_a(x,x)$; the trace of this identity is the first assertion.

The field $a+d\boldsymbol{\lambda}$ is again of type~\eqref{eq:free-coefficient-cpt}. Lemma~\ref{lem:expansion-powers-inverse}, applied to $a$ and to $a+d\boldsymbol{\lambda}$, and the first assertion give
\[
\sum_{j=0}^5\mathsf{M}^{-2-j}\cdot\left(\rho_j(x;a+d\boldsymbol{\lambda})-\rho_j(x;a)\right)=O\left(\mathsf{M}^{-8}\right)\qquad(\mathsf{M}\to\infty).
\]
Multiplication by $\mathsf{M}^2$ and the limit $\mathsf{M}\to\infty$ show that the coefficient with $j=0$ vanishes; multiplication by $\mathsf{M}^3,\dots,\mathsf{M}^7$ in turn shows the same for $j=1,\dots,5$.
\end{proof}

\begin{lemma}[The coefficients of the expansion]\label{lem:coefficients-expansion}
Let $a$ be a field of type~\eqref{eq:free-coefficient-cpt} and $x\in\mathbb{R}^3$, and let $F(x):=\left(F_{\nu\mu}(x)\right)_{\nu,\mu=1}^3$ be the field strength of $a$ at $x$ (Section~\ref{sec:coefficient-fourth-order}). Then
\[
\rho_j(x;a)=0\quad\text{for }j=0,1,2,3,5,\qquad\rho_4(x;a)=\Psi\left(F(x)\right).
\]
\end{lemma}

\begin{proof}
\emph{Odd orders.} By the first identity of~\eqref{eq:free-coefficient-K1} with $\mathsf{M}=1$, $S_1(z)=Y_1(z)-(\gamma\cdot\hat{z})\,Y_1'(\mathsf{r})$, where $\mathsf{r}=|z|$ and $\hat{z}=z/\mathsf{r}$; the first term is even in $z$ and is a multiple of the unit matrix, the second is odd in $z$ and has the matrix factor $\gamma\cdot\hat{z}$. Since $\hat{\Gamma}\gamma_{\mu}\hat{\Gamma}^{-1}=-\gamma_{\mu}$ (Section~\ref{sec:letters-norms-dirac}),
\[
S_1(-z)=\hat{\Gamma}\,S_1(z)\,\hat{\Gamma}^{-1}\qquad(z\ne0).
\]
Moreover $V_j(x;-w)=(-1)^j\,V_j(x;w)$, because $V_j(x;\cdot)$ is homogeneous of degree $j$, and $\hat{\Gamma}^{-1}V_j(x;w)\,\hat{\Gamma}=-V_j(x;w)$. Fix a quadruple $(j_1,\dots,j_4)$ with sum $j$ and substitute $w_i\mapsto-w_i$, $i=1,\dots,4$, in its integral. The four factors $V_{j_i}$ produce the sign $(-1)^j$. Each of the five factors $S_1$ now has the negative of its former argument, that is, it is conjugated by $\hat{\Gamma}$. The first factor $\hat{\Gamma}$ and the last factor $\hat{\Gamma}^{-1}$ cancel under the trace; between them stand the four factors $\hat{\Gamma}^{-1}V_{j_i}(x;w_i)\,\hat{\Gamma}=-V_{j_i}(x;w_i)$, and their four signs multiply to $+1$. So the term of $(j_1,\dots,j_4)$ is $(-1)^j$ times itself. Summing over the quadruples, $\rho_j(x;a)=(-1)^j\rho_j(x;a)$; hence $\rho_j(x;a)=0$ for odd $j$.

\emph{Radial gauge at $x$.} Fix $x\in\mathbb{R}^3$. For $z\in\mathbb{R}^3$ and $\mu=1,2,3$ let
\[
\check{p}_{\mu}(z):=\sum_{k=0}^5\frac{1}{k!}\sum_{\nu_1,\dots,\nu_k=1}^3z_{\nu_1}\cdots z_{\nu_k}\cdot\left(\partial_{\nu_1}\cdots\partial_{\nu_k}a_{\mu}\right)(x),
\]
so that $\check{p}=(\check{p}_1,\check{p}_2,\check{p}_3)$ is the Taylor polynomial of order five of $a$ at $x$, written in the variable $z=y-x$; and let
\[
\boldsymbol{\lambda}_0(z):=-\sum_{k=0}^5\frac{1}{(k+1)!}\sum_{\mu,\nu_1,\dots,\nu_k=1}^3z_{\mu}\,z_{\nu_1}\cdots z_{\nu_k}\cdot\left(\partial_{\nu_1}\cdots\partial_{\nu_k}a_{\mu}\right)(x).
\]
Let $\boldsymbol{\chi}\in C_c^{\infty}(\mathbb{R}^3)$ be real with $\boldsymbol{\chi}=1$ on a neighbourhood of $0$, and put
\[
\boldsymbol{\lambda}(y):=\boldsymbol{\chi}(y-x)\cdot\boldsymbol{\lambda}_0(y-x)\quad(y\in\mathbb{R}^3),\qquad a':=a+d\boldsymbol{\lambda}.
\]
Then $\boldsymbol{\lambda}$ is a real function in $C_c^{\infty}(\mathbb{R}^3)$, $a'$ is a field of type~\eqref{eq:free-coefficient-cpt}, and $\boldsymbol{\lambda}(y)=\boldsymbol{\lambda}_0(y-x)$ for $y$ near $x$. Since $d\boldsymbol{\lambda}_0=(\partial_1\boldsymbol{\lambda}_0,\partial_2\boldsymbol{\lambda}_0,\partial_3\boldsymbol{\lambda}_0)$ is a polynomial of degree at most five, the Taylor polynomial of order five of $a'$ at $x$, in the variable $z$, equals $\check{p}':=\check{p}+d\boldsymbol{\lambda}_0$.

For $0\le k\le5$ let $\check{p}^{(k)}$ be the homogeneous part of degree $k$ of $\check{p}$, and $\boldsymbol{\lambda}_0^{(k+1)}$ the homogeneous part of degree $k+1$ of $\boldsymbol{\lambda}_0$. By the two definitions, $(k+1)\,\boldsymbol{\lambda}_0^{(k+1)}(z)=-z\cdot\check{p}^{(k)}(z)$, and by Euler's relation for homogeneous polynomials $z\cdot\partial\boldsymbol{\lambda}_0^{(k+1)}(z)=(k+1)\,\boldsymbol{\lambda}_0^{(k+1)}(z)$, where $z\cdot\partial:=\sum_{\mu}z_{\mu}\partial_{\mu}$. Summing over $k$,
\[
z\cdot\partial\boldsymbol{\lambda}_0(z)=-z\cdot\check{p}(z),\qquad\text{hence}\qquad z\cdot\check{p}'(z)=0\quad\text{for all }z\in\mathbb{R}^3.
\]
The homogeneous part of degree one of this identity is $z\cdot\check{p}'(0)=0$; so $\check{p}'(0)=0$. The part of degree two is $\sum_{\mu,\nu}z_{\mu}z_{\nu}\,\partial_{\nu}\check{p}'_{\mu}(0)=0$; so the matrix $\left(\partial_{\nu}\check{p}'_{\mu}(0)\right)_{\nu,\mu}$ is antisymmetric. As $d\boldsymbol{\lambda}_0$ is closed and $\partial_{\nu}\check{p}_{\mu}(0)=\partial_{\nu}a_{\mu}(x)$,
\[
\partial_{\nu}\check{p}'_{\mu}(0)=\frac{1}{2}\left(\partial_{\nu}\check{p}'_{\mu}-\partial_{\mu}\check{p}'_{\nu}\right)(0)=\frac{1}{2}\left(\partial_{\nu}a_{\mu}-\partial_{\mu}a_{\nu}\right)(x)=\frac{1}{2}F_{\nu\mu}(x).
\]
Let $V'_j$ be formed as $V_j$, with $a'$ in place of $a$. The derivatives of $a'$ at $x$ of order at most five are those of $\check{p}'$ at $0$; therefore
\[
V'_0(x;\cdot)=0,\qquad V'_1(x;w)=i\gamma\cdot\sum_{\nu=1}^3w_{\nu}\,\partial_{\nu}\check{p}'(0)=i\gamma\cdot A_{F(x)}(w)\qquad(w\in\mathbb{R}^3).
\]
By Lemma~\ref{lem:gauge-invariance}, applied with this $\boldsymbol{\lambda}$, $\rho_j(x;a)=\rho_j(x;a')$ for $j\le5$. In $\rho_j(x;a')$ every term whose quadruple $(j_1,\dots,j_4)$ has an entry $0$ vanishes, because $V'_0(x;\cdot)=0$. For $j\le3$ every quadruple with sum $j$ has such an entry; hence $\rho_j(x;a)=0$ for $j=0,1,2,3$. For $j=4$ only the quadruple $(1,1,1,1)$ remains. As $V'_1(x;\cdot)$ is the matrix $V_1(0;\cdot)$ formed with the field $A_{F(x)}$, this term is $\rho_4(0;A_{F(x)})=\Psi\left(F(x)\right)$, by the definition of $\Psi$ in Section~\ref{sec:expansion-large-mass}. Together with the vanishing for odd $j$ this proves the lemma.
\end{proof}

\subsubsection{Fields of constant strength}\label{sec:fields-constant-strength}

\begin{lemma}\label{lem:free-coefficient-1}
Let $F$ be a real antisymmetric $3\times3$ matrix, and let $B\in\mathbb{R}^3$ be its vector:
\[
B_{\mathfrak{l}}=\frac{1}{2}\sum_{\nu,\mu=1}^3\varepsilon_{\mathfrak{l}\nu\mu}\,F_{\nu\mu}\qquad(\mathfrak{l}=1,2,3),
\]
where $\varepsilon_{\mathfrak{l}\nu\mu}$ is the sign of the permutation $(\mathfrak{l},\nu,\mu)$ of $(1,2,3)$ when the three indices are distinct, and $\varepsilon_{\mathfrak{l}\nu\mu}=0$ otherwise. Then the number $\Psi(F)$ of Section~\ref{sec:expansion-large-mass} is given by
\[
\Psi(F)=-\frac{4}{45}\,(4\pi)^{-\frac{3}{2}}\,\Gamma\left(\frac{7}{2}\right)|B|^4.
\]
\end{lemma}

\begin{proof}
Throughout, $\mathsf{M}=1$. By the definitions of Section~\ref{sec:expansion-large-mass}, $\Psi(F)=\rho_4(0;A_F)=r_4(0;1,A_F)$ for the field $A_F$ of type~\eqref{eq:free-coefficient-lin}, $A_{F,\mu}(z)=\frac{1}{2}\sum_{\nu}z_{\nu}F_{\nu\mu}$. As $V_1(0;w)=i\gamma\cdot A_F(w)$ and $V_j(0;\cdot)=0$ for $j\ne1$, the only term of $\rho_4(0;A_F)$ that survives is the one with $(j_1,j_2,j_3,j_4)=(1,1,1,1)$:
\begin{multline*}
\Psi(F)=\operatorname{tr}\int\mathrm{d}w_1\,\mathrm{d}w_2\,\mathrm{d}w_3\,\mathrm{d}w_4\,S_1(-w_1)\,V_1(0;w_1)\,S_1(w_1-w_2)\,V_1(0;w_2)\\
S_1(w_2-w_3)\,V_1(0;w_3)\,S_1(w_3-w_4)\,V_1(0;w_4)\,S_1(w_4).
\end{multline*}
The integral converges absolutely: by the estimate in the proof of Lemma~\ref{lem:kernels}(a), the integral of the modulus of the integrand is at most $c_{A_F}^4\,\mathsf{s}^{*5}(0)$, which is finite by Lemma~\ref{lem:convolution-powers}.

\emph{(1) Reduction to $B=(0,0,\mathsf{b})$.} Let $\mathrm{Rt}\in\mathrm{SO}(3)$. Since $\mathrm{SU}(2)$ is a covering group of $\mathrm{SO}(3)$, there is $U\in\mathrm{SU}(2)$ with $U\,(\sigma\cdot w)\,U^{-1}=\sigma\cdot(\mathrm{Rt}\,w)$ for all $w\in\mathbb{R}^3$, where $\sigma\cdot w:=\sum_{\mu}w_{\mu}\sigma_{\mu}$. With $\hat{U}:=\operatorname{diag}(U,U)$ and $\gamma\cdot w=\operatorname{diag}(\sigma\cdot w,-\sigma\cdot w)$ (Section~\ref{sec:letters-norms-dirac}) it follows that
\[
\hat{U}\,(\gamma\cdot w)\,\hat{U}^{-1}=\gamma\cdot(\mathrm{Rt}\,w)\qquad(w\in\mathbb{R}^3).
\]
By~\eqref{eq:free-coefficient-K1}, $S_1(z)=Y_1(z)-\left(\gamma\cdot\frac{z}{|z|}\right)Y_1'(|z|)$ for $z\ne0$, where $Y_1(z)$ and $Y_1'(|z|)$ depend on $|z|$ only; as $|\mathrm{Rt}\,z|=|z|$, it follows that $S_1(\mathrm{Rt}\,z)=\hat{U}\,S_1(z)\,\hat{U}^{-1}$ for $z\ne0$. Further, $A_F(w)=-\frac{1}{2}Fw$; so the matrix $F':=\mathrm{Rt}\cdot F\cdot\mathrm{Rt}^T$ satisfies $A_{F'}(\mathrm{Rt}\,w)=\mathrm{Rt}\,A_F(w)$, and therefore $\gamma\cdot A_{F'}(\mathrm{Rt}\,w)=\hat{U}\left(\gamma\cdot A_F(w)\right)\hat{U}^{-1}$. After the change of variables $(w_1,\dots,w_4)\mapsto(\mathrm{Rt}\,w_1,\dots,\mathrm{Rt}\,w_4)$, which preserves Lebesgue measure, the integrand of the formula for $\Psi(F')$ becomes $\hat{U}$ times the integrand of the formula for $\Psi(F)$ times $\hat{U}^{-1}$; as the trace is invariant under conjugation, $\Psi(F')=\Psi(F)$. The vector of $F'$ is $\mathrm{Rt}\,B$ (indeed $Fw=w\times B$ for $w\in\mathbb{R}^3$, and rotations preserve the vector product). Choosing $\mathrm{Rt}$ with $\mathrm{Rt}\,B=(0,0,|B|)$, one may therefore assume $B=(0,0,\mathsf{b})$ with $\mathsf{b}:=|B|\ge0$. Then $F_{12}=\mathsf{b}$, $F_{13}=F_{23}=0$, and the field $A:=A_F$ is
\[
A(z)=\left(-\frac{\mathsf{b}\cdot z_2}{2},\ \frac{\mathsf{b}\cdot z_1}{2},\ 0\right),\qquad\partial\cdot A=0.
\]

\emph{(2) The square of $\gamma\cdot D_u$.} Let $\mathsf{E}$ be the ring of all $\mathbb{C}$-linear maps of $\mathscr{S}$ into itself, and let $\mathsf{E}[[u]]$ be the ring of formal power series in $u$ with coefficients in $\mathsf{E}$, as in the proof of Lemma~\ref{lem:gauge-invariance}. Let $D_u:=\partial+iuA$, that is, $D_{u,\mu}=\partial_{\mu}+iuA_{\mu}$ for $\mu=1,2,3$. The relations $\gamma_{\mu}\gamma_{\nu}+\gamma_{\nu}\gamma_{\mu}=2\delta_{\mu\nu}$ of Section~\ref{sec:letters-norms-dirac} give
\[
(\gamma\cdot D_u)^2=\sum_{\mu}D_{u,\mu}^2+\sum_{\mu<\nu}\gamma_{\mu}\gamma_{\nu}\left[D_{u,\mu},D_{u,\nu}\right].
\]
One has $[D_{u,1},D_{u,2}]=iu\left(\partial_1A_2-\partial_2A_1\right)=iu\cdot\mathsf{b}$, the two other commutators vanish, $\gamma_1\gamma_2=i\operatorname{diag}(\sigma_3,\sigma_3)$, and $\sum_{\mu}D_{u,\mu}^2=\Delta+2iu\,A\cdot\partial-u^2|A|^2$ since $\partial\cdot A=0$. Hence, in $\mathsf{E}[[u]]$,
\begin{gather*}
(1-\gamma\cdot D_u)(1+\gamma\cdot D_u)=(1+\gamma\cdot D_u)(1-\gamma\cdot D_u)=\operatorname{diag}\left(\mathsf{h}_u,\mathsf{h}_u\right),\\
\mathsf{h}_u:=1-\Delta+uW_1+u^2W_2,\qquad W_1:=-2i\,A\cdot\partial+\mathsf{b}\cdot\sigma_3,\qquad W_2:=|A|^2=\frac{\mathsf{b}^2}{4}\left(z_1^2+z_2^2\right).
\end{gather*}
Here $\mathsf{h}_u\in\mathsf{E}_2[[u]]$, where $\mathsf{E}_2$ is the ring of all $\mathbb{C}$-linear maps of $\mathcal{S}(\mathbb{R}^3;\mathbb{C}^2)$ into itself, and $\operatorname{diag}(\mathsf{h}_u,\mathsf{h}_u)$ acts on $\mathscr{S}=\mathcal{S}(\mathbb{R}^3;\mathbb{C}^2)\oplus\mathcal{S}(\mathbb{R}^3;\mathbb{C}^2)$, according to the splitting $\mathbb{C}^4=\mathbb{C}^2\oplus\mathbb{C}^2$ of Section~\ref{sec:letters-norms-dirac}. In $\mathsf{E}_2$ the operator $1-\Delta$ has an inverse, the Fourier multiplier $\left(1+|p|^2\right)^{-1}$; as $1-\Delta$ is the constant term of the power series $\mathsf{h}_u$, the latter is invertible in $\mathsf{E}_2[[u]]$.

By Lemma~\ref{lem:kernels}(b), $S_1$ and $V_A:=i\gamma\cdot A$ map $\mathscr{S}$ into itself. Let $\mathsf{D}_u(A):=1+\gamma\cdot\partial+u\,i\gamma\cdot A=1+\gamma\cdot D_u\in\mathsf{E}[[u]]$, formed as $\mathsf{D}_u(\tilde a)$ in the proof of Lemma~\ref{lem:gauge-invariance}, with $\mathsf{M}=1$ and with the field $A$ in place of $\tilde a$. Its constant term is $1+\gamma\cdot\partial$, which is invertible in $\mathsf{E}$ with inverse $S_1$; so $\mathsf{D}_u(A)=(1+\gamma\cdot\partial)(1+u\,S_1V_A)$ is invertible in $\mathsf{E}[[u]]$, and its inverse is
\[
\sum_{n\ge0}(-u)^n\left(S_1V_A\right)^nS_1=\sum_{n\ge0}u^n\,\mathfrak{R}^{(n)}_A,\qquad\mathfrak{R}^{(n)}_A=(-1)^n\left(S_1V_A\right)^nS_1.
\]
By the identity above, $(1-\gamma\cdot D_u)\operatorname{diag}(\mathsf{h}_u^{-1},\mathsf{h}_u^{-1})$ is a right inverse of $1+\gamma\cdot D_u$, and inverses in a ring are unique. Therefore
\[
\sum_{n\ge0}u^n\,\mathfrak{R}^{(n)}_A=\left(1+\gamma\cdot D_u\right)^{-1}=\left(1-\gamma\cdot D_u\right)\operatorname{diag}\left(\mathsf{h}_u^{-1},\mathsf{h}_u^{-1}\right).
\]

\emph{(3) Mehler's kernel.} For $t>0$ and $z\in\mathbb{R}^3$ let
\[
g_t(z):=(4\pi t)^{-\frac{3}{2}}\exp\left(-\frac{|z|^2}{4t}\right)
\]
be the heat kernel. For $x,y\in\mathbb{R}^3$ put $x_p:=(x_1,x_2)$ and $x\wedge y:=x_1y_2-x_2y_1$; note that $x\wedge y=(x-y)\wedge y$. For $t>0$ and complex $u$ with $|u|\cdot\mathsf{b}\cdot t<\pi$ let $\boldsymbol{\tau}:=u\cdot\mathsf{b}\cdot t$, and for $x,y\in\mathbb{R}^3$ let $k_u(t;x,y)$ be the complex $2\times2$ matrix
\[
k_u(t;x,y):=g_t(x-y)\,\frac{\boldsymbol{\tau}}{\sinh\boldsymbol{\tau}}\,\exp\left\{-\frac{\left(\boldsymbol{\tau}\coth\boldsymbol{\tau}-1\right)\left|x_p-y_p\right|^2}{4t}+\frac{iu\cdot\mathsf{b}}{2}\,x\wedge y\right\}\exp\left(-\boldsymbol{\tau}\cdot\sigma_3\right).
\]
The functions $\boldsymbol{\tau}\mapsto\boldsymbol{\tau}/\sinh\boldsymbol{\tau}$ and $\boldsymbol{\tau}\mapsto\boldsymbol{\tau}\coth\boldsymbol{\tau}$, with the value $1$ at $\boldsymbol{\tau}=0$, are analytic in the disc $|\boldsymbol{\tau}|<\pi$; hence $k_u(t;x,y)$ is an analytic function of $u$ in the disc $|u|\cdot\mathsf{b}\cdot t<\pi$. For all such $t$ and $u$,
\[
\left(\partial_t+\mathsf{h}_u-1\right)k_u=0,
\]
where $\mathsf{h}_u$ acts on the variable $x$. For the proof let $b':=u\cdot\mathsf{b}$. Since $A\cdot\partial=\frac{\mathsf{b}}{2}\left(x_1\partial_2-x_2\partial_1\right)$, step (2) gives
\[
\mathsf{h}_u-1=-\partial_3^2+\check{h}+b'\sigma_3,\qquad\check{h}:=-\partial_1^2-\partial_2^2-ib'\left(x_1\partial_2-x_2\partial_1\right)+\frac{b'^2}{4}\left|x_p\right|^2,
\]
while $k_u$ is the product
\[
k_u(t;x,y)=(4\pi t)^{-\frac{1}{2}}\exp\left(-\frac{(x_3-y_3)^2}{4t}\right)\cdot\check{k}\cdot\exp\left(-\boldsymbol{\tau}\cdot\sigma_3\right)
\]
with
\[
\check{k}:=\frac{b'}{4\pi\sinh(b't)}\,\exp(\check{\phi}),\qquad\check{\phi}:=-\check{\alpha}\cdot\left|x_p-y_p\right|^2+\frac{ib'}{2}\,x\wedge y,\qquad\check{\alpha}:=\frac{b'}{4}\coth(b't).
\]
(The functions of $b'$ that occur here and below are analytic in the disc $|b'|\,t<\pi$, their values at $b'=0$ being the limits; for instance $b'/\sinh(b't)$ and $b'\coth(b't)$ take the value $1/t$ at $b'=0$.) The first factor satisfies the heat equation in the variables $(t,x_3)$, and $\partial_t\exp(-\boldsymbol{\tau}\cdot\sigma_3)=-b'\sigma_3\exp(-\boldsymbol{\tau}\cdot\sigma_3)$, which compensates the term $b'\sigma_3$ of $\mathsf{h}_u-1$. It remains to prove that $\partial_t\check{k}=-\check{h}\,\check{k}$. One computes
\begin{align*}
\partial_1\check{\phi}&=-2\check{\alpha}\cdot(x_1-y_1)+\frac{ib'y_2}{2},\qquad\partial_2\check{\phi}=-2\check{\alpha}\cdot(x_2-y_2)-\frac{ib'y_1}{2},\qquad\left(\partial_1^2+\partial_2^2\right)\check{\phi}=-4\check{\alpha},\\
\left(\partial_1\check{\phi}\right)^2+\left(\partial_2\check{\phi}\right)^2&=4\check{\alpha}^2\cdot\left|x_p-y_p\right|^2-\frac{b'^2}{4}\left|y_p\right|^2-2i\,\check{\alpha}\cdot b'\,x\wedge y,\\
\left(x_1\partial_2-x_2\partial_1\right)\check{\phi}&=2\check{\alpha}\cdot x\wedge y-\frac{ib'}{2}\,x_p\cdot y_p.
\end{align*}
Since $\partial_j^2\exp(\check{\phi})=\left(\partial_j^2\check{\phi}+\left(\partial_j\check{\phi}\right)^2\right)\exp(\check{\phi})$, these identities give
\[
\frac{\check{h}\,\check{k}}{\check{k}}=4\check{\alpha}-4\check{\alpha}^2\cdot\left|x_p-y_p\right|^2+\frac{b'^2}{4}\left(\left|y_p\right|^2-2\,x_p\cdot y_p+\left|x_p\right|^2\right)=4\check{\alpha}-\left(4\check{\alpha}^2-\frac{b'^2}{4}\right)\left|x_p-y_p\right|^2.
\]
On the other hand
\[
\frac{\partial_t\check{k}}{\check{k}}=-b'\coth(b't)-\check{\alpha}'\cdot\left|x_p-y_p\right|^2,\qquad\check{\alpha}':=\partial_t\check{\alpha}=-\frac{b'^2}{4\sinh^2(b't)}=-\left(4\check{\alpha}^2-\frac{b'^2}{4}\right),
\]
and $b'\coth(b't)=4\check{\alpha}$. So $\partial_t\check{k}=-\check{h}\,\check{k}$, and the product rule gives $\left(\partial_t+\mathsf{h}_u-1\right)k_u=0$.

\emph{(4) Expansion in powers of $u$.} The function $\boldsymbol{\tau}/\sinh\boldsymbol{\tau}$ is a power series in $\boldsymbol{\tau}^2=u^2\cdot\mathsf{b}^2\cdot t^2$, and $\left(\boldsymbol{\tau}\coth\boldsymbol{\tau}-1\right)/(4t)$ is a power series in $u^2$ in which the coefficient of $u^{2j}$, $j\ge1$, is a real number times $\mathsf{b}^{2j}\cdot t^{2j-1}$. Hence $k_u=\sum_{n\ge0}u^nk^{(n)}$ with
\[
k^{(n)}(t;x,y)=g_t(x-y)\,\check{P}_n(t,x,y),
\]
where $\check{P}_n$ is a finite sum of monomials $t^j(x-y)^{\boldsymbol{\alpha}}y^{\boldsymbol{\beta}}$ ($j\ge0$ an integer, $\boldsymbol{\alpha},\boldsymbol{\beta}\in\{0,1,2,\dots\}^3$ multi-indices) with coefficients in the linear span of $1$ and $\sigma_3$, each monomial satisfying $2j+|\boldsymbol{\alpha}|\ge n$. To see this, assign to $t$ the weight $2$ and to each component of $x-y$ the weight $1$. Every occurrence of $u$ in $\boldsymbol{\tau}/\sinh\boldsymbol{\tau}$ and in $\exp(-\boldsymbol{\tau}\cdot\sigma_3)$ comes with a factor $t$, so that $u^k$ comes with the weight $2k$ there; a term $u^{2j}t^{2j-1}\left|x_p-y_p\right|^2$ of the exponent has the weight $4j$; and $u\,x\wedge y=u\,(x-y)\wedge y$ has the weight $1$. Thus in every product of these terms the weight is at least the power of $u$. Moreover $\check{P}_0=1$. Comparing the coefficients of $u^n$ in the identity of step (3), one finds, for all $t>0$, $x,y\in\mathbb{R}^3$ and $n\ge0$, with $k^{(-1)}=k^{(-2)}:=0$,
\[
\partial_tk^{(n)}-\Delta_xk^{(n)}+W_{1,x}k^{(n-1)}+W_2(x)\,k^{(n-2)}=0.
\]

\emph{(5) The inverse of $\mathsf{h}_u$.} For an integer $j\ge0$, a multi-index $\boldsymbol{\alpha}$ and $z\in\mathbb{R}^3$ let
\[
\check{m}_{j,\boldsymbol{\alpha}}(z):=\int_0^{\infty}\mathrm{d}t\,\exp(-t)\,t^jz^{\boldsymbol{\alpha}}g_t(z)
\]
whenever the integral converges absolutely. This is the case for every $z\ne0$, since $\exp\left(-\frac{|z|^2}{4t}\right)$ tends to $0$ faster than every power of $t$ as $t\to0$. At $z=0$ the integrand vanishes if $\boldsymbol{\alpha}\ne0$, and it equals $(4\pi)^{-\frac{3}{2}}\exp(-t)\,t^{j-\frac{3}{2}}$ if $\boldsymbol{\alpha}=0$, which is integrable over $(0,\infty)$ if $j\ge1$. Thus $\check{m}_{j,\boldsymbol{\alpha}}$ is defined on $\mathbb{R}^3\setminus\{0\}$, and on all of $\mathbb{R}^3$ if $2j+|\boldsymbol{\alpha}|\ge1$. The Fourier transform of $\check{m}_{j,\boldsymbol{\alpha}}$ equals $(i\partial_p)^{\boldsymbol{\alpha}}\left[j!\left(1+|p|^2\right)^{-j-1}\right]$; this function of $p\in\mathbb{R}^3$ is smooth, and it and all its derivatives are bounded. Convolution with $\check{m}_{j,\boldsymbol{\alpha}}$ therefore maps $\mathcal{S}(\mathbb{R}^3;\mathbb{C}^2)$ into itself. Put
\[
\check{K}^{(n)}(x,y):=\int_0^{\infty}\mathrm{d}t\,\exp(-t)\,k^{(n)}(t;x,y)=\sum c_{j,\boldsymbol{\alpha},\boldsymbol{\beta}}\,\check{m}_{j,\boldsymbol{\alpha}}(x-y)\,y^{\boldsymbol{\beta}},
\]
where the sum and the coefficients $c_{j,\boldsymbol{\alpha},\boldsymbol{\beta}}\in\operatorname{span}\{1,\sigma_3\}$ are those of $\check{P}_n$. The integral converges absolutely and equals the sum for all $x,y\in\mathbb{R}^3$ with $x\ne y$, and also for $x=y$ if $n\ge1$, because every monomial of $\check{P}_n$ then has $2j+|\boldsymbol{\alpha}|\ge1$. Further let
\[
\left(\mathfrak{K}^{(n)}\,\boldsymbol{\phi}\right)(x):=\int\mathrm{d}y\,\check{K}^{(n)}(x,y)\,\boldsymbol{\phi}(y)\qquad\left(\boldsymbol{\phi}\in\mathcal{S}(\mathbb{R}^3;\mathbb{C}^2),\ x\in\mathbb{R}^3\right).
\]
Then $\mathfrak{K}^{(n)}\in\mathsf{E}_2$; put $\mathfrak{K}^{(-1)}=\mathfrak{K}^{(-2)}:=0$. For $\boldsymbol{\phi}\in\mathcal{S}(\mathbb{R}^3;\mathbb{C}^2)$ let $\boldsymbol{\Phi}_n(t,x):=\int\mathrm{d}y\,k^{(n)}(t;x,y)\,\boldsymbol{\phi}(y)$ for $t>0$ and $x\in\mathbb{R}^3$, and $\boldsymbol{\Phi}_{-1}=\boldsymbol{\Phi}_{-2}:=0$. By step (4),
\[
\partial_t\boldsymbol{\Phi}_n=\Delta\boldsymbol{\Phi}_n-W_1\boldsymbol{\Phi}_{n-1}-W_2\boldsymbol{\Phi}_{n-2}.
\]
Multiplication by $\exp(-t)$ and integration over $(\epsilon,\infty)$, $\epsilon>0$, give
\begin{multline*}
-\exp(-\epsilon)\,\boldsymbol{\Phi}_n(\epsilon,\cdot)+(1-\Delta)\int_{\epsilon}^{\infty}\mathrm{d}t\,\exp(-t)\,\boldsymbol{\Phi}_n(t,\cdot)\\
+W_1\int_{\epsilon}^{\infty}\mathrm{d}t\,\exp(-t)\,\boldsymbol{\Phi}_{n-1}(t,\cdot)+W_2\int_{\epsilon}^{\infty}\mathrm{d}t\,\exp(-t)\,\boldsymbol{\Phi}_{n-2}(t,\cdot)=0.
\end{multline*}
Now let $\epsilon\to0$. As $\check{P}_0=1$ and $(g_t)_{t>0}$ is an approximate identity, $\boldsymbol{\Phi}_0(\epsilon,\cdot)\to\boldsymbol{\phi}$ uniformly. For each monomial $c\,t^j(x-y)^{\boldsymbol{\alpha}}y^{\boldsymbol{\beta}}$ of $\check{P}_m$, $m\ge0$, one has
\[
\left|\int\mathrm{d}y\,g_t(x-y)\,c\,t^j(x-y)^{\boldsymbol{\alpha}}y^{\boldsymbol{\beta}}\boldsymbol{\phi}(y)\right|\le|c|\cdot\sup_{y\in\mathbb{R}^3}\left|y^{\boldsymbol{\beta}}\boldsymbol{\phi}(y)\right|\cdot t^j\int\mathrm{d}z\,\left|z^{\boldsymbol{\alpha}}\right|g_t(z)=c'\,t^{j+\frac{|\boldsymbol{\alpha}|}{2}},
\]
with a number $c'$ independent of $t$ and $x$, because $\int\mathrm{d}z\,|z^{\boldsymbol{\alpha}}|\,g_t(z)$ is a constant multiple of $t^{\frac{|\boldsymbol{\alpha}|}{2}}$ and $y^{\boldsymbol{\beta}}\boldsymbol{\phi}$ is bounded. Hence $|\boldsymbol{\Phi}_m(t,x)|$ is bounded uniformly in $0<t\le1$ and $x\in\mathbb{R}^3$, for every $m\ge0$; and for $n\ge1$, where $2j+|\boldsymbol{\alpha}|\ge1$ for every monomial of $\check{P}_n$, $\sup_x|\boldsymbol{\Phi}_n(\epsilon,x)|\le c''\cdot\epsilon^{\frac{1}{2}}\to0$ for $\epsilon\le1$, with $c''$ independent of $\epsilon$. So $\int_0^{\epsilon}\mathrm{d}t\,\exp(-t)\,\boldsymbol{\Phi}_m(t,x)\to0$ uniformly in $x$, and since $\int_0^{\infty}\mathrm{d}t\,\exp(-t)\,\boldsymbol{\Phi}_m(t,x)=\left(\mathfrak{K}^{(m)}\,\boldsymbol{\phi}\right)(x)$ by Fubini's theorem, $\int_{\epsilon}^{\infty}\mathrm{d}t\,\exp(-t)\,\boldsymbol{\Phi}_m(t,\cdot)\to\mathfrak{K}^{(m)}\,\boldsymbol{\phi}$ uniformly. Uniform convergence implies convergence in $\mathcal{S}'(\mathbb{R}^3;\mathbb{C}^2)$, on which differential operators with polynomial coefficients act continuously. It follows that
\[
(1-\Delta)\,\mathfrak{K}^{(n)}\,\boldsymbol{\phi}+W_1\,\mathfrak{K}^{(n-1)}\,\boldsymbol{\phi}+W_2\,\mathfrak{K}^{(n-2)}\,\boldsymbol{\phi}=\delta_{n0}\,\boldsymbol{\phi}\qquad(n\ge0),
\]
that is, $\mathsf{h}_u\sum_{n\ge0}u^n\,\mathfrak{K}^{(n)}=1$ in $\mathsf{E}_2[[u]]$. Being a right inverse of the invertible element $\mathsf{h}_u$ of this ring, the series $\sum_{n\ge0}u^n\,\mathfrak{K}^{(n)}$ is its inverse: $\mathsf{h}_u^{-1}=\sum_{n\ge0}u^n\,\mathfrak{K}^{(n)}$.

\emph{(6) The value at the origin.} By steps (2) and (5), $\sum_{n\ge0}u^n\,\mathfrak{R}^{(n)}_A=\left(1-\gamma\cdot\partial-u\,i\gamma\cdot A\right)\sum_{n\ge0}u^n\operatorname{diag}\left(\mathfrak{K}^{(n)},\mathfrak{K}^{(n)}\right)$; the coefficient of $u^4$ is
\[
\mathfrak{R}^{(4)}_A=\left(1-\gamma\cdot\partial\right)\operatorname{diag}\left(\mathfrak{K}^{(4)},\mathfrak{K}^{(4)}\right)-i\gamma\cdot A\,\operatorname{diag}\left(\mathfrak{K}^{(3)},\mathfrak{K}^{(3)}\right).
\]
For $j\ge0$ and $s\ge0$ let
\[
\check{\mu}_j(s):=\int_0^{\infty}\mathrm{d}t\,\exp(-t)\,t^j\,(4\pi t)^{-\frac{3}{2}}\exp\left(-\frac{s^2}{4t}\right)
\]
whenever the integral converges; as in step (5), this is the case for every $s>0$, and for $s=0$ if $j\ge1$. Then $\check{m}_{j,\boldsymbol{\alpha}}(z)=z^{\boldsymbol{\alpha}}\cdot\check{\mu}_j(|z|)$ for $z\ne0$, and also for $z=0$ if $j\ge1$, while $\check{m}_{0,\boldsymbol{\alpha}}(0)=0$ if $\boldsymbol{\alpha}\ne0$. One has $\check{\mu}_0(|z|)=Y_1(z)=\exp(-|z|)/(4\pi|z|)$ for $z\ne0$. For $j\ge1$ the function $\check{\mu}_j$ is bounded and continuous on $[0,\infty)$, by dominated convergence, the integrand being at most $(4\pi)^{-\frac{3}{2}}\exp(-t)\,t^{j-\frac{3}{2}}$; and $\check{\mu}_j'(s)=-\frac{s}{2}\,\check{\mu}_{j-1}(s)$ for $s>0$. Hence the function $\check{m}_{j,\boldsymbol{\alpha}}$, defined on all of $\mathbb{R}^3$ when $2j+|\boldsymbol{\alpha}|\ge1$, is continuous on $\mathbb{R}^3$ if $2j+|\boldsymbol{\alpha}|\ge2$, and continuously differentiable on $\mathbb{R}^3$ if $2j+|\boldsymbol{\alpha}|\ge3$. On $\mathbb{R}^3\setminus\{0\}$ both properties follow from the formulas for $Y_1$ and for $\check{\mu}_j'$. At $0$, for $j=0$ use $\check{m}_{0,\boldsymbol{\alpha}}(0)=0$, $|z^{\boldsymbol{\alpha}}|\le|z|^{|\boldsymbol{\alpha}|}$, $Y_1(z)\le(4\pi|z|)^{-1}$ and the fact that the gradient of $Y_1$ at $z\ne0$ has the length $\left(1+|z|^{-1}\right)Y_1(z)$; for $j\ge1$ use that the gradient of $z\mapsto\check{\mu}_j(|z|)$ at $z\ne0$ is $-\frac{1}{2}\,\check{\mu}_{j-1}(|z|)\,z$, which is bounded if $j=1$ and tends to $0$ as $z\to0$ if $j\ge2$. In this way one finds that, if $2j+|\boldsymbol{\alpha}|\ge3$, the gradient of $\check{m}_{j,\boldsymbol{\alpha}}$ on $\mathbb{R}^3\setminus\{0\}$ has a limit at $0$; and a continuous function on $\mathbb{R}^3$ that is continuously differentiable on $\mathbb{R}^3\setminus\{0\}$, with a gradient that has a limit at $0$, is continuously differentiable on $\mathbb{R}^3$. Since, for $n\ge1$, $\check{K}^{(n)}(x,y)$ is the finite sum of step (5) for all $x,y\in\mathbb{R}^3$, and every monomial of $\check{P}_n$ has $2j+|\boldsymbol{\alpha}|\ge n$, it follows that $\check{K}^{(3)}$, $\check{K}^{(4)}$ and $\partial_x\check{K}^{(4)}$ are continuous on $\mathbb{R}^3\times\mathbb{R}^3$ and grow at most polynomially; thus the operator on the right-hand side of the last formula for $\mathfrak{R}^{(4)}_A$ has the continuous kernel
\[
\left(1-\gamma\cdot\partial_x\right)\operatorname{diag}\left(\check{K}^{(4)},\check{K}^{(4)}\right)(x,y)-i\gamma\cdot A(x)\,\operatorname{diag}\left(\check{K}^{(3)},\check{K}^{(3)}\right)(x,y).
\]
On the other hand, $\mathfrak{R}^{(4)}_A=(S_1V_A)^4S_1$ has the continuous kernel $\mathrm{Rk}^{(4)}_A$ by Lemma~\ref{lem:kernels}(b),(c), where $\mathrm{Rk}^{(4)}_A$ is the kernel $\mathrm{Rk}^{(4)}$ of Section~\ref{sec:kernels-products-propagators} formed with $\mathsf{M}=1$ and the field $A$. A continuous kernel is determined by the operator that it defines on $\mathscr{S}$; so the two kernels coincide. At $x=y=0$ one has $A(0)=0$, and $\operatorname{tr}\left[\gamma_{\mu}\operatorname{diag}(\check{X},\check{X})\right]=\operatorname{tr}(\sigma_{\mu}\check{X})-\operatorname{tr}(\sigma_{\mu}\check{X})=0$ for every complex $2\times2$ matrix $\check{X}$ and $\mu=1,2,3$. Therefore, $\check{K}^{(4)}(0,0)$ being given by the integral over $t$ of step (5) at $x=y=0$,
\[
\Psi(F)=\operatorname{tr}\mathrm{Rk}^{(4)}_A(0,0)=2\operatorname{tr}\check{K}^{(4)}(0,0)=2\int_0^{\infty}\mathrm{d}t\,\exp(-t)\operatorname{tr}k^{(4)}(t;0,0).
\]
Finally, $\operatorname{tr}\exp(-\boldsymbol{\tau}\cdot\sigma_3)=2\cosh\boldsymbol{\tau}$, so that
\[
\operatorname{tr}k_u(t;0,0)=(4\pi t)^{-\frac{3}{2}}\,\frac{\boldsymbol{\tau}}{\sinh\boldsymbol{\tau}}\cdot2\cosh\boldsymbol{\tau}=2\,(4\pi t)^{-\frac{3}{2}}\,\boldsymbol{\tau}\coth\boldsymbol{\tau},
\]
and by $s\coth s=1+\frac{s^2}{3}-\frac{s^4}{45}+O(s^6)$ the coefficient of $u^4$ in this function is $-\frac{2}{45}\,(4\pi t)^{-\frac{3}{2}}\,\mathsf{b}^4\cdot t^4$. Since $\int_0^{\infty}\mathrm{d}t\,\exp(-t)\,t^{\frac{5}{2}}=\Gamma\left(\frac{7}{2}\right)$, it follows that
\[
\Psi(F)=-\frac{4}{45}\,(4\pi)^{-\frac{3}{2}}\,\Gamma\left(\frac{7}{2}\right)\cdot\mathsf{b}^4,\qquad\mathsf{b}=|B|.
\]
\end{proof}

\subsubsection{Large mass and dilated probes}\label{sec:large-mass-dilated}

\begin{proof}[Proof of Theorem~\ref{thm:free-coefficient-1}(a)]
\emph{The first assertion.} Let $a\in C_c^{\infty}(\mathbb{R}^3;\mathbb{R}^3)$ be real, that is, a field of type~\eqref{eq:free-coefficient-cpt}, and let $F(x)$ be its field strength at $x$ and $B=\operatorname{curl}a$, as in Section~\ref{sec:coefficient-fourth-order}; for every $x$, $B(x)$ is the vector of the matrix $F(x)$. Put
\[
\check{\psi}:=-\frac{4}{45}\,(4\pi)^{-\frac{3}{2}}\,\Gamma\left(\frac{7}{2}\right).
\]
By Lemmas~\ref{lem:coefficients-expansion} and~\ref{lem:free-coefficient-1}, $\rho_j(x;a)=0$ for $j=0,1,2,3,5$ and $\rho_4(x;a)=\Psi(F(x))=\check{\psi}\cdot|B(x)|^4$ for all $x\in\mathbb{R}^3$. Lemma~\ref{lem:expansion-powers-inverse} therefore yields, for $\mathsf{M}\ge1$,
\[
\left|r_4(x;\mathsf{M},a)-\mathsf{M}^{-6}\cdot\check{\psi}\cdot|B(x)|^4\right|\le C'_a\cdot\mathsf{M}^{-8}\quad(x\in\mathbb{R}^3),\qquad\int_{\operatorname{dist}(x,\operatorname{supp}a)\ge1}\mathrm{d}x\,\left|r_4(x;\mathsf{M},a)\right|\le C'_a\cdot\exp(-\mathsf{M}).
\]
By Lemma~\ref{lem:derivative-mass}, $x\mapsto r_4(x;\mathsf{M},a)$ is integrable and $\frac{\mathrm{d}}{\mathrm{d}\mathsf{M}}q_{4,\mathsf{M}}(a)=\int\mathrm{d}x\,r_4(x;\mathsf{M},a)$. Integrate the first bound over the set $\{x:\operatorname{dist}(x,\operatorname{supp}a)<1\}$, which has finite volume since $\{x:\operatorname{dist}(x,\operatorname{supp}a)\le1\}$ has, and use the second bound on its complement, on which $B=0$; since moreover $\exp(-\mathsf{M})\le8!\cdot\mathsf{M}^{-8}$, there is a constant $C''_a$ such that
\[
\left|\frac{\mathrm{d}}{\mathrm{d}\mathsf{M}}q_{4,\mathsf{M}}(a)-\mathsf{M}^{-6}\cdot\check{\psi}\int\mathrm{d}x\,|B(x)|^4\right|\le C''_a\cdot\mathsf{M}^{-8}\qquad(\mathsf{M}\ge1).
\]
Also by Lemma~\ref{lem:derivative-mass}, $\mathsf{M}\mapsto q_{4,\mathsf{M}}(a)$ is continuously differentiable on $(0,\infty)$, and $q_{4,\mathsf{M}}(a)\to0$ as $\mathsf{M}\to\infty$. Hence, for $\mathsf{M}\ge1$,
\[
q_{4,\mathsf{M}}(a)=-\int_{\mathsf{M}}^{\infty}\mathrm{d}M'\,\frac{\mathrm{d}}{\mathrm{d}M'}q_{4,M'}(a),\qquad\left|q_{4,\mathsf{M}}(a)+\frac{\check{\psi}}{5}\cdot\mathsf{M}^{-5}\int\mathrm{d}x\,|B(x)|^4\right|\le\frac{C''_a}{7}\cdot\mathsf{M}^{-7}.
\]
With $\Gamma\left(\frac{7}{2}\right)=\frac{15}{8}\,\pi^{\frac{1}{2}}$ one has
\[
-\frac{\check{\psi}}{5}=\frac{4}{225}\cdot\frac{1}{8}\,\pi^{-\frac{3}{2}}\cdot\frac{15}{8}\,\pi^{\frac{1}{2}}=\frac{1}{240\pi}.
\]
This proves the first assertion of Theorem~\ref{thm:free-coefficient-1}(a), with $C_a:=C''_a/7$.

\emph{The probes $\beta^{(l)}$.} Let now $\rho$, $a=(\partial_2\rho,-\partial_1\rho,0)$, $I_{\beta}$, $\beta^{(l)}$, $a_l$ and $f_l$, $l\ge1$, be as in Theorem~\ref{thm:free-coefficient-1}(a), and let $C_a$ be the constant of the first assertion for this field $a$. The formula for $d^{*}$ in Section~\ref{sec:probes-discretisation} gives
\[
\left(d^{*}\beta^{(l)}\right)_1=-\partial_2\beta^{(l)}_{21}=l\,(\partial_2\rho)(x/l),\qquad\left(d^{*}\beta^{(l)}\right)_2=-\partial_1\beta^{(l)}_{12}=-l\,(\partial_1\rho)(x/l),\qquad\left(d^{*}\beta^{(l)}\right)_3=0,
\]
so that $d^{*}\beta^{(l)}=a_l$ and $f_l=da_l=dd^{*}\beta^{(l)}$. The third component of $B=\operatorname{curl}a$ is $B_3=\partial_1a_2-\partial_2a_1=-\left(\partial_1^2+\partial_2^2\right)\rho$. The expansion $\rho(x)=\exp(-1)\left(1-|x|^2+O\left(|x|^4\right)\right)$ near $0$ gives $B_3(0)=4\exp(-1)\ne0$; as $B$ is continuous, $I_{\beta}=\int\mathrm{d}x\,|B(x)|^4>0$.

\emph{Dilation.} For $l\ge1$ the operator $U_{\mathrm{d}}$, $(U_{\mathrm{d}}\,\boldsymbol{\psi})(x):=l^{-\frac{3}{2}}\boldsymbol{\psi}(x/l)$, is unitary on $L^2(\mathbb{R}^3;\mathbb{C}^4)$; one has $U_{\mathrm{d}}^{-1}\partial\,U_{\mathrm{d}}=l^{-1}\partial$, and $U_{\mathrm{d}}^{-1}a_l(\cdot)\,U_{\mathrm{d}}=l\,a(\cdot)$ as multiplication operators. Hence
\[
U_{\mathrm{d}}^{-1}S_{\bar{m}}\,i\gamma\cdot a_l\,U_{\mathrm{d}}=\left(\bar{m}+l^{-1}\gamma\cdot\partial\right)^{-1}l\,i\gamma\cdot a=l^2\,S_{\bar{m}\cdot l}\,i\gamma\cdot a,
\]
and, the trace being invariant under unitary conjugation, $q_{4,\bar{m}}(a_l)=l^8\,q_{4,\bar{m}\cdot l}(a)$. Let $\bar{m}\cdot l\ge1$. The definition of $b$ in Section~\ref{sec:free-fourth-order} and the first assertion with $\mathsf{M}=\bar{m}\cdot l$ give
\[
b(f_l)=24\,l^8\,q_{4,\bar{m}\cdot l}(a)=\frac{l^3\cdot I_{\beta}}{10\pi\cdot\bar{m}^5}+\boldsymbol{\theta},\qquad|\boldsymbol{\theta}|\le24\,C_a\,l^8\left(\bar{m}\cdot l\right)^{-7}=24\,C_a\,l\cdot\bar{m}^{-7}.
\]

\emph{Positivity.} The number $b(f_l)=24\,q_{4,\bar{m}}(a_l)$ is real by Lemma~\ref{lem:hilbert-schmidt-bound}. For $\bar{m}\cdot l\ge1$ the last estimate gives
\[
b(f_l)\ge\frac{l^3\cdot I_{\beta}}{10\pi\cdot\bar{m}^5}-24\,C_a\,l\cdot\bar{m}^{-7}=\frac{l\cdot I_{\beta}}{10\pi\cdot\bar{m}^7}\left(\left(\bar{m}\cdot l\right)^2-\frac{240\pi\,C_a}{I_{\beta}}\right),
\]
and the right-hand side is positive if $\left(\bar{m}\cdot l\right)^2>240\pi\,C_a/I_{\beta}$.

Theorem~\ref{thm:free-coefficients-lattice}(a) follows.
\end{proof}

\subsection{Admissibility, limit points and the proof of Part~I of the main theorem}\label{sec:admissibility-limit-points}

In this section $(N,T)$ always denotes a cutoff, that is, a pair of integers $N,T\ge0$ with $n=h\cdot L^{T+N}\ge2$; the value $n=1$ arises only for $h=1$ and $N=T=0$. Since $h$ is a power of the odd integer $L$, so is $n$; as $n\ge2$, it follows that $n$ is odd and $n\ge L$. Moreover $L\ge5$ by stage~\ref{it:order-choice-4}. Put
\[
\tilde{m}:=\frac{n-1}{2};
\]
this is the integer $m$ of Section~\ref{sec:tori-cells-cochains}.

\begin{proposition}[Admissible families of curvature laws]\label{prop:admissible-families-curvature}
Let $\left(\mathrm{Pr}_{N,T}\right)$ be a family that assigns to every cutoff $(N,T)$ a Borel probability measure $\mathrm{Pr}_{N,T}$ on the subspace $\ker d$ of $C^2(\Lambda_{\eta,R})$, where $\eta=L^{-N}$ and $R=h\cdot L^T$, and suppose that this family is admissible in the sense of Definition~\ref{def:admissible-families}. For every cutoff $(N,T)$ let $\mathsf{P}_{N,T}$ be the law of $F^{\sharp}$ under $\mathrm{Pr}_{N,T}$. Then the laws $\mathsf{P}_{N,T}$ satisfy all assertions of Theorem~\ref{thm:estimates-fixed-cutoff} except the last one, and all assertions of Theorem~\ref{thm:second-moments-tightness}. Moreover, let $\mu$ be a Borel probability measure on $H_{-4}$ such that $\mathsf{P}_{N_i,T_i}\to\mu$ weakly on $H_{-4}$ along a sequence $(N_i,T_i)_{i\ge1}$ of cutoffs with $N_i\to\infty$ and $T_i\to\infty$. Then $\mu$ has all the properties stated in Theorems~\ref{thm:continuum-infinite-1} and~\ref{thm:moments-limit-point} and in Corollary~\ref{cor:transfer-operators-coordinate}.
\end{proposition}

\begin{proof}
Theorem~\ref{thm:estimates-fixed-cutoff} applies to every Borel probability measure on $C^2(\Lambda_{\eta,R})$ with property~\eqref{eq:continuum-infinite-Ad2}, in particular to each member $\mathrm{Pr}_{N,T}$ of the family; its last assertion concerns measures concentrated on $\operatorname{ran}\,d$. Theorem~\ref{thm:second-moments-tightness} applies to every admissible family, and Theorems~\ref{thm:continuum-infinite-1}, \ref{thm:moments-limit-point} and Corollary~\ref{cor:transfer-operators-coordinate} apply to every limit point of an admissible family in the sense of Definition~\ref{def:limit-points}; these limit points are exactly the measures $\mu$ of the assertion.
\end{proof}

\begin{proof}[Proof of Part~I of Theorem~\ref{thm:main-theorem}]
\emph{The harmonic-averaged law~\eqref{eq:setup-av}.} Consider the parameter assignment (Section~\ref{sec:wilson-operator-lattice})
\[
(N,T)\mapsto(q,\boldsymbol{\mu})=\left(e\cdot\eta^{\frac{1}{2}},\boldsymbol{\mu}_N\right)
\]
and the family $\left(\bar{\mathsf{P}}_{\eta,R;q,\boldsymbol{\mu}}\right)_{N,T}$ of averaged laws that it determines, $(q,\boldsymbol{\mu})$ being its value at the cutoff $(N,T)$. At each cutoff $(N,T)$ the definitions of Section~\ref{sec:wilson-determinant-lattice} make $P^{\mathrm{av}}_{N,T}$ the distribution of $F^{\sharp}$ under $\bar{\mathsf{P}}_{\eta,R;q,\boldsymbol{\mu}}$, which Section~\ref{sec:existence-limit-points} denotes by $\mathsf{P}_{N,T}$ for this family. Hence a measure is a limit point of $\left(P^{\mathrm{av}}_{N,T}\right)$ in the sense of Section~\ref{sec:wilson-determinant-lattice} if and only if it is a limit point of the family in the sense of Definition~\ref{def:limit-points}. The family is admissible by Corollary~\ref{cor:averaged-laws-admissible}, which is derived from Theorems~\ref{thm:gaussian-domination-positive} and~\ref{thm:limit-vanishing-proca} and from Proposition~\ref{prop:gaussian-bounds-2}. The statements on admissible families and their limit points therefore apply, and give the following. Assertion~(a) is Theorem~\ref{thm:second-moments-tightness}. Property~\eqref{it:main-theorem-P1} follows from Theorem~\ref{thm:continuum-infinite-1}(1) and~(4) together with Proposition~\ref{prop:sub-gaussian-random}(1), (5) and~(7); by Theorem~\ref{thm:continuum-infinite-1}(1), the conclusions of Proposition~\ref{prop:sub-gaussian-random} hold under every limit point $\mu$ for the space $L^2(\mathbb{R}^3;\Lambda^2)$ with the norm $\|\cdot\|_2$. Property~\eqref{it:main-theorem-P2} follows from Theorem~\ref{thm:moments-limit-point}(1) and~(2) and from Proposition~\ref{prop:sub-gaussian-random}(7); property~\eqref{it:main-theorem-P4} from Theorem~\ref{thm:continuum-infinite-1}(8), (3) and~(4); property~\eqref{it:main-theorem-P5} from Theorem~\ref{thm:continuum-infinite-1}(5) and~(6); and property~\eqref{it:main-theorem-P6} from Theorem~\ref{thm:continuum-infinite-1}(7) and Theorem~\ref{thm:moments-limit-point}(5). Property~\eqref{it:main-theorem-P7} is Corollary~\ref{cor:transfer-operators-coordinate}, and property~\eqref{it:main-theorem-P8} is Theorem~\ref{thm:continuum-infinite-1}(2).

\emph{The all-flux law~\eqref{eq:setup-af}.} For every cutoff $(N,T)$ let $\mathrm{Pr}_{N,T}$ be the law of $F=\eta^{-\frac{3}{2}}\cdot F^u$ under $E^{\mathrm{af}}_n$, where $n=h\cdot L^{T+N}$ and the lattice mass is $\boldsymbol{\mu}_N$, the torus $\Lambda_n$ of lattice units being identified with $\Lambda_{\eta,R}$ as in Sections~\ref{sec:lattices-tests-function} and~\ref{sec:wilson-determinant-lattice}; here $q=e\cdot\eta^{\frac{1}{2}}>0$ and $\boldsymbol{\mu}_N>-\frac{1}{2}$. As $dF^u=0$ (Section~\ref{sec:wilson-determinant-lattice}), $\mathrm{Pr}_{N,T}$ is a Borel probability measure on $\ker d\subset C^2(\Lambda_{\eta,R})$, and $P^{\mathrm{af}}_{N,T}$ is the law of $F^{\sharp}$ under $\mathrm{Pr}_{N,T}$. It is shown next that the family $\left(\mathrm{Pr}_{N,T}\right)$ is admissible in the sense of Definition~\ref{def:admissible-families}, that is, that each $\mathrm{Pr}_{N,T}$ has the properties \eqref{eq:continuum-infinite-Ad1}, \eqref{eq:continuum-infinite-Ad2} and \eqref{eq:continuum-infinite-Ad3}.

Property~\eqref{eq:continuum-infinite-Ad1} holds because, by Theorem~\ref{thm:free-coefficient-1}(d5), the law $E^{\mathrm{af}}_n$ on the cubic torus $\Lambda_n$ is invariant under the lattice translations and under $O_{\mathrm{cub}}$.

Property~\eqref{eq:continuum-infinite-Ad2}: let $g\in C^2(\Lambda_{\eta,R})$ be real, and put $g^u:=\eta^{\frac{3}{2}}\cdot g$. Then $\langle g,F\rangle_{\eta}=\langle g^u,F^u\rangle$ and $|g|_{\eta}=\|g^u\|$ (Section~\ref{sec:lattices-tests-function}), so that \eqref{eq:continuum-infinite-Ad2} is the inequality
\[
E^{\mathrm{af}}_n\exp\langle g^u,F^u\rangle\le\exp\left(\frac{1}{2}\|g^u\|^2\right)
\]
of Corollary~\ref{cor:partition-functions-1}(a).

Property~\eqref{eq:continuum-infinite-Ad3}: let $\nu_0\in\{1,2,3\}$ be the direction that takes the place of $\nu$ in \eqref{eq:continuum-infinite-Ad3}, and let $s$, $g_1,\dots,g_s$ and $c_1,\dots,c_s$ be as there for $\nu=\nu_0$; thus the $g_b$ are real and supported on the set $\mathcal{P}^{+\circ}_{\nu_0}$ of Section~\ref{sec:admissible-families-lattice}, in which $m=\tilde{m}$. Define $\check{P}\colon C^2(\Lambda_{\eta,R})\to\mathbb{C}$ by
\[
\check{P}(F):=\sum_{b=1}^sc_b\exp\left(i\langle g_b,F\rangle_{\eta}\right);
\]
through $F=\eta^{-\frac{3}{2}}\cdot F^u$ it is also a function of $F^u$. By Theorem~\ref{thm:free-coefficient-1}(d5) the functional $\boldsymbol{\omega}_n$ of the cubic torus $\Lambda_n$ does not change under the exchange of two axes; so $\nu_0$ may be taken as the time direction, numbered $0$ as in Section~\ref{sec:trace-formula-chessboard}. Let $t\in\{-\tilde{m},\dots,\tilde{m}\}$ denote the time coordinate of a site in lattice units. For every time $t$, read modulo $n$, let $\mathrm{El}_t$ and $\tilde{b}_t$ be the electric and the magnetic part of $F^u$ in the slice $t$, as defined in Lemma~\ref{lem:dependence-boundary-slices}(c): the entries of $\mathrm{El}_t$ are the values of $F^u$ on the plaquettes that contain the direction $\nu_0$ and join the times $t$ and $t+1$, and the entries of $\tilde{b}_t$ are the values of $F^u$ on the plaquettes that are orthogonal to $\nu_0$ and lie at the time $t$. A plaquette of $\mathcal{P}^{+\circ}_{\nu_0}$ either contains the direction $\nu_0$ and joins the times $t$ and $t+1$ for some $t$ with $0\le t<\tilde{m}$, or is orthogonal to $\nu_0$ and lies at a time $t$ with $0<t<\tilde{m}$. Since each $g_b$ vanishes outside $\mathcal{P}^{+\circ}_{\nu_0}$, every number $\langle g_b,F\rangle_{\eta}$ is a real linear form in the variables $\mathrm{El}_t$, $0\le t<\tilde{m}$, and $\tilde{b}_t$, $0<t<\tilde{m}$. Therefore $\check{P}$ depends on $F^u$ only through these variables, and as a function of them it is bounded and continuous. In particular $\check{P}$ is a slab function at $0$ in the sense of Section~\ref{sec:slabs-operators}, for the cell width $\boldsymbol{\lambda}=\tilde{m}+\frac{1}{2}$ of Definition~\ref{def:decomposition-torus-cells}, for which $m_{\mathrm{w}}=\tilde{m}\ge2$ (because $n\ge5$) and $n_0=n=2\tilde{m}+1=2\boldsymbol{\lambda}$. Let $\mathfrak{r}$ be the reflection $t\mapsto-t$. Corollary~\ref{cor:reflection-positivity-fixed} yields
\[
\boldsymbol{\omega}_n\left(\overline{\check{P}\circ\mathfrak{r}^{*}}\cdot\check{P}\right)\ge0.
\]
The pull-back $\mathfrak{r}^{*}$ of two-cochains is the reflection of Section~\ref{sec:tori-cochains-lattice}, which changes the sign of the components that contain the reflected direction. By the definition in Section~\ref{sec:admissible-families-lattice} it is therefore the operator $\hat{\mathsf{R}}_{\nu_0}$, an orthogonal involution of $C^2(\Lambda_{\eta,R})$; in particular $\langle g,\hat{\mathsf{R}}_{\nu_0}F\rangle_{\eta}=\langle\hat{\mathsf{R}}_{\nu_0}g,F\rangle_{\eta}$ for all $g$ and $F$. Since the $g_a$ are real, it follows that
\begin{align*}
\overline{\check{P}\left(\mathfrak{r}^{*}F\right)}\cdot\check{P}(F)&=\sum_{a=1}^s\overline{c_a}\exp\left(-i\left\langle\hat{\mathsf{R}}_{\nu_0}g_a,F\right\rangle_{\eta}\right)\cdot\sum_{b=1}^sc_b\exp\left(i\langle g_b,F\rangle_{\eta}\right)\\
&=\sum_{a,b=1}^s\overline{c_a}\cdot c_b\exp\left(i\left\langle g_b-\hat{\mathsf{R}}_{\nu_0}g_a,F\right\rangle_{\eta}\right).
\end{align*}
As $E^{\mathrm{af}}_n=\boldsymbol{\omega}_n/Z_n$ with $Z_n>0$ (Section~\ref{sec:wilson-determinant-lattice}), the integral of the right-hand side with respect to $\mathrm{Pr}_{N,T}$ equals
\[
\frac{1}{Z_n}\cdot\boldsymbol{\omega}_n\left(\overline{\check{P}\circ\mathfrak{r}^{*}}\cdot\check{P}\right)\ge0,
\]
which is \eqref{eq:continuum-infinite-Ad3}.

By Proposition~\ref{prop:admissible-families-curvature}, the laws $P^{\mathrm{af}}_{N,T}$ satisfy the assertions of Theorem~\ref{thm:second-moments-tightness}, and every limit point of $\left(P^{\mathrm{af}}_{N,T}\right)$ in the sense of Section~\ref{sec:wilson-determinant-lattice} has the properties stated in Theorems~\ref{thm:continuum-infinite-1} and~\ref{thm:moments-limit-point} and in Corollary~\ref{cor:transfer-operators-coordinate}. Assertion~(a) and the properties \eqref{it:main-theorem-P1}, \eqref{it:main-theorem-P2} and \eqref{it:main-theorem-P4}--\eqref{it:main-theorem-P8} for $\mathrm{nat}=\mathrm{af}$ follow from these statements, with the same references as for $\mathrm{nat}=\mathrm{av}$.

\emph{Property~\eqref{it:main-theorem-P3} for both laws.} Let $\mathrm{nat}\in\{\mathrm{af},\mathrm{av}\}$ and let $\mu$ be a limit point of $\left(P^{\mathrm{nat}}_{N,T}\right)$. By Theorem~\ref{thm:moments-limit-point}(1) one has $S_0=1$, and for every integer $k\ge0$ the functional $S_k$ is symmetric and takes real values on real tests. Choose a sequence $(N_l,T_l)_{l\ge1}$ of cutoffs with $N_l\to\infty$, $T_l\to\infty$ and $P^{\mathrm{nat}}_{N_l,T_l}\to\mu$ weakly on $H_{-4}$, and form $E_l$ and $X_l$ with this sequence as in~\eqref{it:main-theorem-P8}. Thus $E_l$ is the expectation $E^{\mathrm{nat}}_n$ of Section~\ref{sec:wilson-determinant-lattice} on the cubic torus $\Lambda_n$ with $n=h\cdot L^{T_l+N_l}$, and by Theorem~\ref{thm:free-coefficient-1}(d5) the configurations $F^u$ and $-F^u$ have the same law under $E_l$. As $X_l(f)$ is linear in the configuration, it follows that
\[
E_l\exp\left(iX_l(f)\right)=E_l\exp\left(-iX_l(f)\right)\qquad\text{for every real }f\in\mathcal{S}\text{ and every }l.
\]
For real $f\in\mathcal{S}$ put
\[
\Phi_{\mu}(f):=\int\mathrm{d}\mu(F)\,\exp\left(iF(f)\right);
\]
the function $\Phi_{\mu}$, defined on the set of the real elements of $\mathcal{S}$, is the characteristic functional of $\mu$. Applying~\eqref{it:main-theorem-P8} to the single test $f$, at $z=i$ and at $z=-i$, one finds that, as $l\to\infty$, the two sides of the last display tend to $\Phi_{\mu}(f)$ and to $\Phi_{\mu}(-f)$, respectively. Hence $\Phi_{\mu}(-f)=\Phi_{\mu}(f)$ for every real $f\in\mathcal{S}$. The map $f\mapsto\Phi_{\mu}(-f)$ is the characteristic functional of the image of $\mu$ under $F\mapsto-F$, so that this image and $\mu$ have the same characteristic functional. The real elements of $\mathcal{S}$ form the space of test forms of Section~\ref{sec:test-forms-distributions-2}, and Lemma~\ref{lem:exponentials-determination-law}(b), applied with this whole space as the sequentially dense set, shows that a law on $\mathcal{S}'$ with its cylinder $\sigma$-algebra is determined by its characteristic functional. Therefore $\mu$ is invariant under $F\mapsto-F$; Theorem~\ref{thm:continuum-infinite-1}(9) gives the same conclusion. It follows that $S_k=(-1)^kS_k$ for every $k\ge0$, so that $S_k=0$ for every odd $k$.

\emph{Assertion~(c).} Let $\mu$ be a limit point. By~\eqref{it:main-theorem-P1}, $\mu$ satisfies~\eqref{eq:setup-G}; by~\eqref{it:main-theorem-P5}, it is invariant under the translations of $\mathbb{R}^3$ and under $O_{\mathrm{cub}}$; and~\eqref{it:main-theorem-P6}, applied with $c=0$ to the finite linear combinations of the functions $\exp\left(iX(f)\right)$ with real $f\in C_c^{\infty}(\mathbb{R}^3;\Lambda^2)$ supported in $\{x_{\nu_0}>0\}$, is~\eqref{eq:setup-RPj} for $j=\nu_0$, for each $\nu_0\in\{1,2,3\}$. Thus $\mu$ satisfies~\eqref{eq:setup-H}. The assertions of (c) on the components $\mu_{\xi}$, other than closedness, follow from Theorem~\ref{thm:ergodic-decomposition}, which for $J=\{1,2,3\}$ takes the form of Remark~\ref{rem:components-all-translations}, from Theorem~\ref{thm:cubic-symmetry-components}(i) and from Corollary~\ref{cor:ergodic-components-2}. As to closedness, $\mu\left(\{F\in\mathcal{S}':dF=0\}\right)=1$ by~\eqref{it:main-theorem-P4}; this event belongs to the cylinder $\sigma$-algebra (Section~\ref{sec:closed-distributions}), and Theorem~\ref{thm:ergodic-decomposition}(d), applied to it, gives $\mu_{\xi}\left(\{F\in\mathcal{S}':dF=0\}\right)=1$ for $\nu$-almost every $\xi$.

\emph{Assertion~(d)} is proved in Remark~\ref{rem:repeated-limits} below, by means of the definition of a limit point in Section~\ref{sec:wilson-determinant-lattice} and of the Prokhorov metric.
\end{proof}

\begin{remark}[Repeated limits]\label{rem:repeated-limits}
The Hermite--Sobolev space $H_{-4}$ is a separable Hilbert space (Section~\ref{sec:hermite-sobolev-spaces}), in particular a separable complete metric space. Let $d_P$ be the Prokhorov metric on the set of the Borel probability measures on $H_{-4}$. By a theorem of Prokhorov, a sequence $(\rho_k)_{k\ge1}$ of such measures converges weakly to such a measure $\rho$ if and only if $d_P(\rho_k,\rho)\to0$. Let $\mathrm{nat}\in\{\mathrm{af},\mathrm{av}\}$. Suppose that, for infinitely many $N$, a Borel probability measure $\mu^{(N)}$ on $H_{-4}$ is the weak limit of $P^{\mathrm{nat}}_{N,T_j(N)}$ as $j\to\infty$, along a sequence with $T_j(N)\to\infty$, and that $\mu^{(N_i)}\to\mu$ weakly on $H_{-4}$ along a sequence $N_i\to\infty$ of such $N$. For every $i$ choose an index $j_i$ with
\[
T_{j_i}(N_i)\ge i\qquad\text{and}\qquad d_P\left(P^{\mathrm{nat}}_{N_i,T_{j_i}(N_i)},\mu^{(N_i)}\right)\le\frac{1}{i};
\]
it exists, as $T_j(N_i)\to\infty$ and $d_P\left(P^{\mathrm{nat}}_{N_i,T_j(N_i)},\mu^{(N_i)}\right)\to0$ when $j\to\infty$. Then
\[
d_P\left(P^{\mathrm{nat}}_{N_i,T_{j_i}(N_i)},\mu\right)\le\frac{1}{i}+d_P\left(\mu^{(N_i)},\mu\right)\to0\qquad(i\to\infty),
\]
that is, $P^{\mathrm{nat}}_{N_i,T_{j_i}(N_i)}\to\mu$ weakly on $H_{-4}$ along the sequence of cutoffs $\left(N_i,T_{j_i}(N_i)\right)_{i\ge1}$, both of whose entries tend to infinity. Hence $\mu$ is a limit point in the sense of Section~\ref{sec:wilson-determinant-lattice}. The same argument, with the roles of $N$ and $T$ exchanged, applies to iterated limits taken in the other order. This proves Theorem~\ref{thm:main-theorem}(d).
\end{remark}

\subsection{Trace formula, chessboard estimates and transfer inequalities}\label{sec:trace-formula-chessboard}

All objects of this section are expressed in lattice units, and the three directions are labelled $0,1,2$, as in Section~\ref{sec:tori-cochains-lattice}. Accordingly, a torus is given by a triple $n=(n_0,n_1,n_2)$ of integers $n_{\mu}\ge2$, and a flux is a triple $w=(w_{12},w_{02},w_{01})\in\mathbb{Z}^3$. Further, $q>0$, $r_q=2\pi/q$ and $\boldsymbol{\mu}>-\frac{1}{2}$. The functionals $\boldsymbol{\omega}_n$, $\boldsymbol{\omega}_n^{(w)}$, $E^{\mathrm{af}}_n$ and $E^{\mathrm{av}}_n$, the numbers $Z_n$, $\mathbf{Z}_n$ and $\mathfrak{q}_w(n)$, and the class $\mathfrak{F}(n)$ are those of Section~\ref{sec:wilson-determinant-lattice}. For $x>0$ let
\[
\mathrm{thf}(x):=\sum_{i\in\mathbb{Z}}\exp\left(-x\,i^2\right).
\]

\subsubsection{Recalled results}\label{sec:recalled-results}

\begin{itemize}
\item[(Q1)]\itemmark{Q1}\label{it:free-coefficient-Q1} The functional $E^{\mathrm{av}}_n=\boldsymbol{\omega}_n^{(0)}/\mathbf{Z}_n$ is the harmonic-averaged law of Section~\ref{sec:wilson-determinant-lattice} (Lemma~\ref{lem:sector-zero-flux}; Lemma~\ref{lem:flux-sectors}(b)).

\item[(Q2)]\itemmark{Q2}\label{it:free-coefficient-Q2} Let $g$ be an isometry of $\Lambda_n$ onto $\Lambda_{n'}$ as in Theorem~\ref{thm:free-coefficient-1}(d4). For every $\Phi\in\mathfrak{F}(n)$ and every $w\in\mathbb{Z}^3$,
\[
\boldsymbol{\omega}_{n'}\left(\Phi\circ g^{*}\right)=\boldsymbol{\omega}_n(\Phi),\qquad\boldsymbol{\omega}_{n'}^{(0)}\left(\Phi\circ g^{*}\right)=\boldsymbol{\omega}_n^{(0)}(\Phi),\qquad\boldsymbol{\omega}_{n'}^{(g\cdot w)}\left(\Phi\circ g^{*}\right)=\boldsymbol{\omega}_n^{(w)}(\Phi),
\]
and $\boldsymbol{\omega}_n$, $\boldsymbol{\omega}_n^{(0)}$ are invariant under $F\mapsto-F$ (Theorem~\ref{thm:free-coefficient-1}(d5)). Consequently, a statement of this section that singles out the direction $0$ holds equally when an arbitrary direction $\mu$ is singled out instead. Indeed, let $g$ be the isometry that interchanges the axes $0$ and $\mu$, and put $n':=gn$. Transport by $g$, which substitutes $n'$ for $n$, $\Phi\circ g^{*}$ for $\Phi$ and $g\cdot w$ for each flux $w$, carries such a statement, formulated with $\mu$ as the distinguished direction, into the same statement formulated with $0$ as the distinguished direction, and by the identities above the two formulations are equivalent. With the components written in the order $(12,02,01)$,
\[
g\cdot w=\left(w_{02},w_{12},-w_{01}\right)\quad\text{if }\mu=1,\qquad g\cdot w=\left(-w_{01},-w_{02},-w_{12}\right)\quad\text{if }\mu=2.
\]
For the trace formula this has the following meaning. Let $\boldsymbol{\vartheta}=(\boldsymbol{\vartheta}_1,\boldsymbol{\vartheta}_2)\in\mathbb{R}^2$ be the Bloch angle and $\mathscr{H}_{\boldsymbol{\vartheta}}$ the Hilbert space of \eqref{it:free-coefficient-Q5}. Whichever direction is taken as time, the summands of $\mathscr{H}_{\boldsymbol{\vartheta}}$ in Proposition~\ref{prop:trace-formula-bloch}(iii), quoted in \eqref{it:free-coefficient-Q5}(iii), are indexed by the component of the flux whose two indices both differ from the time direction; the two remaining components, which have the time direction among their indices, serve as the Fourier variables of the formula \eqref{eq:all-flux-11}. A minus sign in front of a Fourier variable amounts to replacing the corresponding component $\boldsymbol{\vartheta}_i$ by $-\boldsymbol{\vartheta}_i$, and the statements in which $\boldsymbol{\vartheta}$ occurs hold for every $\boldsymbol{\vartheta}\in\mathbb{R}^2$. Let now $n$ carry the cell structure of width $\boldsymbol{\lambda}$ (Definition~\ref{def:decomposition-torus-cells}). The torus $n'$ then carries the cell structure of the same width, and for every $c\in\prod_{\nu}\mathbb{Z}/N_{\nu}$ one has $g(\mathfrak{c}_c)=\mathfrak{c}_{c'}$ and $\mathsf{g}_{c'}\circ g=g\circ\mathsf{g}_c$, where $c'$ arises from $c$ by interchanging the entries $c_0$ and $c_{\mu}$. It follows that, for every $O\in\mathfrak{A}^{+}$ on $n$, the function $O\circ g^{*}$ lies in $\mathfrak{A}^{+}$ on $n'$, and that $\left(O\circ g^{*}\right)^{\sharp}=O^{\sharp}\circ g^{*}$.

\item[(Q3)]\itemmark{Q3}\label{it:free-coefficient-Q3} For every torus $n$ and every real $g\in C^2(n)$,
\[
\boldsymbol{\omega}_n\left(\exp\langle g,F\rangle\right)\le\exp\left(\frac{1}{2}\|g\|^2\right)\cdot Z_n
\]
(Theorem~\ref{thm:gaussian-domination}; Theorem~\ref{thm:all-flux-1}).

\item[(Q4)]\itemmark{Q4}\label{it:free-coefficient-Q4} Let $n$ be a torus, $w\in\mathbb{Z}^3$, and let $g\in C^2(n)$ be real with $\sum_xg_{\mu\nu}(x)=0$ for each pair $\mu<\nu$. Then
\[
\boldsymbol{\omega}_n^{(w)}\left(\exp\langle g,F\rangle\right)\le\mathfrak{q}_w(n)\cdot\exp\left(\frac{1}{2}\|g\|^2\right)\cdot Z_n
\]
(Lemma~\ref{lem:one-flux-sector}).

\item[(Q5)]\itemmark{Q5}\label{it:free-coefficient-Q5} The direction $0$ is called time. The cross-section $\boldsymbol{\Sigma}=\mathbb{Z}/n_1\times\mathbb{Z}/n_2$ has $v=n_1n_2$ sites and $l=2v$ links, $d_s$ denotes its coboundary, and $k_B:=(2\pi)^{\frac{l}{2}}$; $\mathbb{Z}^{\mathrm{plaq}(\boldsymbol{\Sigma})}$ and $\mathbb{R}^{\mathrm{plaq}(\boldsymbol{\Sigma})}$ are the sets of the integer and of the real functions on the plaquettes of $\boldsymbol{\Sigma}$. The slice space, the temporal variables and the curvatures are those of Section~\ref{sec:transfer-kernels}: the slice space is
\[
\mathcal{X}=[0,r_q)^l\times\mathbb{Z}^{\mathrm{plaq}(\boldsymbol{\Sigma})},
\]
its points are written $x=(A,\mathfrak{z})$, and the magnetic curvature of $x$ is $\tilde{b}(x)=d_sA+r_q\cdot\mathfrak{z}$. The temporal variables are $a\in[0,r_q)^v$ and $c\in\mathbb{Z}^l$; for $x=(A,\mathfrak{z})$ and $x'=(A',\mathfrak{z}')$ the electric curvature is
\[
\mathrm{El}(x',x;a,c)=A'-A-d_sa+r_q\cdot c,
\]
and the constraint is $\mathfrak{z}'=\mathfrak{z}+d_sc$. The slice data in the sense of Section~\ref{sec:slice-factorisation-determinant} are taken to be those of Lemma~\ref{lem:slice-factorisation-determinant-2}. For a Bloch angle $\boldsymbol{\vartheta}=(\boldsymbol{\vartheta}_1,\boldsymbol{\vartheta}_2)\in\mathbb{R}^2$ let $\mathscr{H}_{\boldsymbol{\vartheta}}$ be the Hilbert space of \eqref{eq:all-flux-8}, and let
\[
\mathscr{H}_{\boldsymbol{\vartheta}}=\bigoplus_{w_{12}\in\mathbb{Z}}\mathscr{H}_{\boldsymbol{\vartheta}}^{(w_{12})}
\]
be its orthogonal decomposition according to the value $w_{12}$ of $\sum_p\mathfrak{z}_p$, the sum being taken over the plaquettes $p$ of $\boldsymbol{\Sigma}$. For an integer $\mathsf{m}\ge1$, a \emph{slab function of thickness $\mathsf{m}$} is a bounded complex Borel function $P$ of the variables
\[
(\mathrm{El}_0,\dots,\mathrm{El}_{\mathsf{m}-1};\tilde{b}_1,\dots,\tilde{b}_{\mathsf{m}-1})\in\left(\mathbb{R}^l\right)^{\mathsf{m}}\times\left(\mathbb{R}^{\mathrm{plaq}(\boldsymbol{\Sigma})}\right)^{\mathsf{m}-1};
\]
its \emph{mirror} is the slab function
\[
P^{\boldsymbol{\theta}}(\mathrm{El}_0,\dots,\mathrm{El}_{\mathsf{m}-1};\tilde{b}_1,\dots,\tilde{b}_{\mathsf{m}-1}):=\overline{P\left(-\mathrm{El}_{\mathsf{m}-1},\dots,-\mathrm{El}_0;\tilde{b}_{\mathsf{m}-1},\dots,\tilde{b}_1\right)}.
\]
The operator on $\mathscr{H}_{\boldsymbol{\vartheta}}$ with the kernel \eqref{eq:all-flux-10} is denoted by $C_{\mathsf{m},\boldsymbol{\vartheta}}[P]$. The symbols $\boldsymbol{\Sigma}$, $r_q$, $\mathfrak{z}$, $\tilde{b}$, $\mathrm{El}$, $\boldsymbol{\vartheta}$, $\mathscr{H}_{\boldsymbol{\vartheta}}$, $k_B$, $\boldsymbol{\theta}$, $\mathrm{Mo}$, $w$ and $\boldsymbol{\omega}_n^{(w)}$ of this item stand for the objects denoted in Section~\ref{sec:trace-formula} by $\Sigma$, $r$, $m$, $\beta$, $\mathcal{E}$, $\vartheta$, $\mathcal{H}_{\vartheta}$, $\kappa$, $\theta$, $M$, $\mathfrak{n}$ and $\omega^{(\mathfrak{n})}_n$, respectively, and the number $r$ of the slabs and the slab index $a$ in~(iii) below for those denoted there by $R$ and $\alpha$. Proposition~\ref{prop:trace-formula-bloch} is formulated for arbitrary slice data, that is, for an arbitrary quadruple subject to the conditions \eqref{eq:all-flux-SD1} and \eqref{eq:all-flux-SD2} of Section~\ref{sec:slice-factorisation-determinant}. Specialised to the slice data of Lemma~\ref{lem:slice-factorisation-determinant-2}, it states the following.
\begin{enumerate}
\item[(i)] The operator $C_{\mathsf{m},\boldsymbol{\vartheta}}[P]$ is a Hilbert--Schmidt operator on $\mathscr{H}_{\boldsymbol{\vartheta}}$, and it depends linearly on $P$.
\item[(ii)] $C_{\mathsf{m},\boldsymbol{\vartheta}}[P^{\boldsymbol{\theta}}]=C_{\mathsf{m},\boldsymbol{\vartheta}}[P]^{*}$.
\item[(iii)] Let $r\ge2$, let $\mathsf{m}_1,\dots,\mathsf{m}_r\ge1$ be integers with $n_0=\mathsf{m}_1+\dots+\mathsf{m}_r$, and for $a=1,\dots,r$ let $P_a$ be a slab function of thickness $\mathsf{m}_a$, regarded as a function of the curvature on the $a$-th slab $\left\{\mathsf{m}_1+\dots+\mathsf{m}_{a-1}\le x_0\le\mathsf{m}_1+\dots+\mathsf{m}_a\right\}$. Then
\[
k_B^{n_0}\cdot\operatorname{Tr}_{\mathscr{H}_{\boldsymbol{\vartheta}}}\left(C_{\mathsf{m}_r,\boldsymbol{\vartheta}}[P_r]\cdots C_{\mathsf{m}_1,\boldsymbol{\vartheta}}[P_1]\right)=\sum_{w\in\mathbb{Z}^3}\exp\left(2\pi i\left(\boldsymbol{\vartheta}_1w_{01}+\boldsymbol{\vartheta}_2w_{02}\right)\right)\cdot\boldsymbol{\omega}_n^{(w)}\left(\prod_{a=1}^rP_a\right).
\]
Every operator $C_{\mathsf{m},\boldsymbol{\vartheta}}[P]$ maps each summand $\mathscr{H}_{\boldsymbol{\vartheta}}^{(w_{12})}$ into itself, and for every $w_{12}\in\mathbb{Z}$ the identity above remains true when on the left the trace over $\mathscr{H}_{\boldsymbol{\vartheta}}$ is replaced by the trace over the invariant subspace $\mathscr{H}_{\boldsymbol{\vartheta}}^{(w_{12})}$ and on the right the sum is restricted to the $w$ with this value of $w_{12}$:
\[
k_B^{n_0}\cdot\operatorname{Tr}_{\mathscr{H}_{\boldsymbol{\vartheta}}^{(w_{12})}}\left(C_{\mathsf{m}_r,\boldsymbol{\vartheta}}[P_r]\cdots C_{\mathsf{m}_1,\boldsymbol{\vartheta}}[P_1]\right)=\sum_{w_{02},w_{01}\in\mathbb{Z}}\exp\left(2\pi i\left(\boldsymbol{\vartheta}_1w_{01}+\boldsymbol{\vartheta}_2w_{02}\right)\right)\cdot\boldsymbol{\omega}_n^{(w_{12},w_{02},w_{01})}\left(\prod_{a=1}^rP_a\right).
\]
\item[(iv)] The operator $T_{\boldsymbol{\vartheta}}:=C_{1,\boldsymbol{\vartheta}}[1]$ equals $\mathrm{Mo}^{*}\cdot K_{\boldsymbol{\vartheta}}\cdot\mathrm{Mo}$, where $K_{\boldsymbol{\vartheta}}\ge0$ is a bounded operator on $\mathscr{H}_{\boldsymbol{\vartheta}}$ and $\mathrm{Mo}$ is the operator of multiplication by the operator-valued function $M$ of \eqref{eq:all-flux-9}; the operator $X_{\boldsymbol{\vartheta}}:=K_{\boldsymbol{\vartheta}}^{\frac{1}{2}}\cdot\mathrm{Mo}$ is Hilbert--Schmidt. Hence $T_{\boldsymbol{\vartheta}}=X_{\boldsymbol{\vartheta}}^{*}X_{\boldsymbol{\vartheta}}\ge0$ is of trace class, and $C_{\mathsf{m},\boldsymbol{\vartheta}}[1]=T_{\boldsymbol{\vartheta}}^{\mathsf{m}}$.
\end{enumerate}
The series on the right sides of the two trace identities above converge absolutely: for every $\Phi\in\mathfrak{F}(n)$,
\[
\sum_{w\in\mathbb{Z}^3}\left|\boldsymbol{\omega}_n^{(w)}(\Phi)\right|\le\boldsymbol{\omega}_n\left(|\Phi|\right),
\]
since every $\boldsymbol{\omega}_n^{(w)}$ is a positive functional (Theorem~\ref{thm:free-coefficient-1}(d2)) and $\sum_w\boldsymbol{\omega}_n^{(w)}=\boldsymbol{\omega}_n$.

\item[(Q6)]\itemmark{Q6}\label{it:free-coefficient-Q6} Let $n^{*}\ge2$ be an integer, and define $\phi\colon\{1,2,\dots\}^3\to\mathbb{R}$ by $\phi(a,b,c):=\log Z_{n^{*}(a,b,c)}$, where $n^{*}(a,b,c)$ is the torus $(n^{*}a,n^{*}b,n^{*}c)$. Then $\phi$ is convex in each variable, and
\[
\phi(ka,b,c)\le k\cdot\phi(a,b,c)\qquad\text{for all integers }k,a,b,c\ge1,
\]
and likewise in the second and in the third variable (Theorem~\ref{thm:partition-functions-1}, for $a,b,c$ of either parity).

\item[(Q7)]\itemmark{Q7}\label{it:free-coefficient-Q7} Let $\phi\colon\{1,2,\dots\}^3\to\mathbb{R}$ be convex in each variable and satisfy $\phi(ka,b,c)\le k\cdot\phi(a,b,c)$ for all integers $k,a,b,c\ge1$, and likewise in the second and in the third variable. Let $j\ge2$ be an integer and
\[
D_j:=j^3\cdot\phi(1,1,1)-\phi(j,j,j).
\]
Then $D_j\ge0$, and for every integer $k\ge j$
\[
\phi(k,k,k)\ge k^3\cdot\phi(1,1,1)-\left(\frac{k-1}{j-1}\right)^3\cdot D_j
\]
(Lemma~\ref{lem:partition-functions-1}).

\item[(Q8)]\itemmark{Q8}\label{it:free-coefficient-Q8} For $1\le p<\infty$ let $S_p$ be the Schatten class of the compact operators $Y$ on a Hilbert space with $\|Y\|_p:=\left(\operatorname{Tr}|Y|^p\right)^{\frac{1}{p}}<\infty$. If $r\ge1$ is an integer and $Y_1,\dots,Y_r\in S_r$, then
\[
\left|\operatorname{Tr}\left(Y_1\cdots Y_r\right)\right|\le\prod_{a=1}^r\|Y_a\|_r
\]
(H\"older's inequality for Schatten classes; Lemma~\ref{lem:schatten-norms-traces}(a)). Every Hilbert--Schmidt operator lies in $S_r$ for every $r\ge2$. For even $r$ one has $\|Y\|_r^r=\operatorname{Tr}\left(Y^{*}Y\right)^{\frac{r}{2}}$. If $B\ge0$ is of trace class, then $\operatorname{Tr}B^k\le\left(\operatorname{Tr}B\right)^k$ for every integer $k\ge1$ (Lemma~\ref{lem:schatten-norms-traces}(c)).
\end{itemize}

For $x>0$ and every integer $\mathbf{M}\ge0$ one has
\[
\mathrm{thf}(x)\le1+\left(\frac{\pi}{x}\right)^{\frac{1}{2}},\qquad\sum_{i\in\mathbb{Z},\ i>\mathbf{M}}\exp\left(-x\,i^2\right)\le\exp\left(-x\cdot(\mathbf{M}+1)^2\right)\cdot\mathrm{thf}(x).
\]
The first bound follows from
\[
\sum_{i\ge1}\exp\left(-x\,i^2\right)\le\int_0^{\infty}\mathrm{d}t\,\exp\left(-x\,t^2\right)=\frac{1}{2}\left(\frac{\pi}{x}\right)^{\frac{1}{2}};
\]
the second follows on writing $i=\mathbf{M}+1+i'$ with an integer $i'\ge0$, because $(\mathbf{M}+1+i')^2\ge(\mathbf{M}+1)^2+i'^2$.

\subsubsection{Cells, reflections and mirror copies}\label{sec:cells-reflections-mirror}

\begin{definition}[Decomposition of a torus into cells]\label{def:decomposition-torus-cells}
A \emph{cell width} is a number $\boldsymbol{\lambda}=m_{\mathrm{w}}+\frac{1}{2}$ with an integer $m_{\mathrm{w}}\ge2$. A torus $n$ \emph{carries the cell structure of width $\boldsymbol{\lambda}$} if $n_{\mu}=\boldsymbol{\lambda}\cdot N_{\mu}$ for $\mu=0,1,2$, with even positive integers $N_{\mu}$. For such a torus, $\Lambda_n$ is regarded as the set of the sites of $\prod_{\mu}\mathbb{R}/n_{\mu}\mathbb{Z}$, and every plaquette as a closed unit square in it; the planes and the cells below are subsets of $\prod_{\mu}\mathbb{R}/n_{\mu}\mathbb{Z}$.
\begin{itemize}
\item For $\mu\in\{0,1,2\}$ and $i\in\mathbb{Z}/N_{\mu}$, the \emph{plane} $\mathrm{Pl}_{\mu,i}$ is $\left\{x_{\mu}=\boldsymbol{\lambda}\cdot i\right\}$. It is a \emph{site plane} if $\boldsymbol{\lambda}\cdot i$ is an integer, and a \emph{link plane} otherwise; a link plane contains no site.
\item For $c\in\prod_{\mu}\mathbb{Z}/N_{\mu}$, the \emph{cell} $\mathfrak{c}_c$ is
\[
\mathfrak{c}_c:=\prod_{\mu}\left[\boldsymbol{\lambda}\cdot c_{\mu},\boldsymbol{\lambda}\cdot(c_{\mu}+1)\right];
\]
the number of the cells is $N_n:=N_0N_1N_2$.
\item The set $P^{+}(\mathfrak{c}_c)$ consists of the plaquettes whose vertices all lie in $\mathfrak{c}_c$ and which are contained in none of the planes $\mathrm{Pl}_{\mu,i}$. A plaquette $p$ belongs to at most one of the sets $P^{+}(\mathfrak{c}_c)$. Indeed, let all vertices of $p$ lie in two different cells $\mathfrak{c}_c$ and $\mathfrak{c}_{c'}$, and choose $\mu$ with $c_{\mu}\ne c'_{\mu}$. Two different intervals $\left[\boldsymbol{\lambda}\cdot i,\boldsymbol{\lambda}\cdot(i+1)\right]$ have only end points in common, so that every vertex of $p$ lies in one of the planes $\mathrm{Pl}_{\mu,i}$. The coordinates $x_{\mu}$ of the vertices of $p$ differ by at most $1$, whereas two different planes $\mathrm{Pl}_{\mu,i}$ are at distance at least $\boldsymbol{\lambda}\ge2$ from each other; hence all vertices of $p$ lie in one of these planes, $p$ is contained in it, and $p$ belongs neither to $P^{+}(\mathfrak{c}_c)$ nor to $P^{+}(\mathfrak{c}_{c'})$. Further, $\mathfrak{A}^{+}(\mathfrak{c}_c)$ denotes the set of the functions in $\mathfrak{F}(n)$ that depend on $F\in C^2(n)$ only through $(F_p)_{p\in P^{+}(\mathfrak{c}_c)}$, and $\mathfrak{A}^{+}:=\mathfrak{A}^{+}(\mathfrak{c}_0)$.
\item For $c\in\prod_{\mu}\mathbb{Z}/N_{\mu}$ let $\mathsf{g}_c$ be the map of the torus onto itself that acts on each coordinate $x_{\mu}$ by
\[
x_{\mu}\mapsto\boldsymbol{\lambda}\cdot c_{\mu}+x_{\mu}\quad\text{if }c_{\mu}\text{ is even},\qquad x_{\mu}\mapsto\boldsymbol{\lambda}\cdot(c_{\mu}+1)-x_{\mu}\quad\text{if }c_{\mu}\text{ is odd}.
\]
As $\boldsymbol{\lambda}\cdot c_{\mu}$ is an integer for even $c_{\mu}$, and $\boldsymbol{\lambda}\cdot(c_{\mu}+1)$ is an integer for odd $c_{\mu}$, the map $\mathsf{g}_c$ is a lattice isometry; it is well defined on the torus, because $N_{\mu}$ is even and $\boldsymbol{\lambda}\cdot N_{\mu}=n_{\mu}$. Each $\mathsf{g}_c$ is a composition of reflections $x_{\mu}\mapsto2\boldsymbol{\lambda}\cdot i-x_{\mu}$ in planes $\mathrm{Pl}_{\mu,i}$, and it carries $\mathfrak{c}_0$ onto $\mathfrak{c}_c$. Taken together, the maps $\mathsf{g}_c$ constitute a group; the reflections in the planes $\mathrm{Pl}_{\mu,i}$ generate it, and its action on the set of the cells is simply transitive. As the maps $\mathsf{g}_c$ act coordinatewise and the cells are products of intervals, it suffices to consider a single coordinate $x_{\mu}$: there the maps $\mathsf{g}_c$ induce the dihedral group that two reflections in neighbouring planes $\mathrm{Pl}_{\mu,i}$ and $\mathrm{Pl}_{\mu,i+1}$ generate, and this dihedral group permutes the $N_{\mu}$ intervals $\left[\boldsymbol{\lambda}\cdot i,\boldsymbol{\lambda}\cdot(i+1)\right]$ simply transitively.
\item For $O\in\mathfrak{A}^{+}$ and a cell $\mathfrak{c}_c$, the \emph{mirror copy} of $O$ in $\mathfrak{c}_c$ is the function
\[
O^{[c]}(F):=\begin{cases}O\left(\mathsf{g}_c^{*}F\right)&\text{if }c_0+c_1+c_2\text{ is even},\\ \overline{O\left(\mathsf{g}_c^{*}F\right)}&\text{if }c_0+c_1+c_2\text{ is odd},\end{cases}
\]
where $\mathsf{g}_c^{*}$ is the pull-back of Section~\ref{sec:reference-tori-cells}; it lies in $\mathfrak{A}^{+}(\mathfrak{c}_c)$, because $\mathsf{g}_c$ maps the planes onto planes and $\mathfrak{c}_0$ onto $\mathfrak{c}_c$, hence $P^{+}(\mathfrak{c}_0)$ onto $P^{+}(\mathfrak{c}_c)$. Put
\[
O^{\sharp}:=\prod_cO^{[c]},
\]
the product being taken over all cells. Every $G\in\mathfrak{A}^{+}(\mathfrak{c}_c)$ equals $O^{[c]}$ for exactly one $O\in\mathfrak{A}^{+}$ (as $\mathsf{g}_c^{*}$ is a bijection of $C^2(n)$), and $O^{[c']}$ is then called the mirror copy of $G$ in the cell $\mathfrak{c}_{c'}$.
\end{itemize}
Once $\boldsymbol{\lambda}$ is fixed, one and the same function $O$, inserted into these formulas, yields mirror copies $O^{[c]}$ and a product $O^{\sharp}$ on each torus that carries the cell structure of width $\boldsymbol{\lambda}$. For $\mu\in\{0,1,2\}$ and an integer $k\ge1$, $n^{(\mu,k)}$ denotes the torus obtained from $n$ by replacing $n_{\mu}$ with $k\cdot n_{\mu}$.
\end{definition}

Let $k\ge1$ be an integer, and assume that $m_{\mathrm{p}}\ge2$. Since $\mathrm{per}=2m_{\mathrm{p}}+1$ is a power of the odd integer $L$ (Section~\ref{sec:reference-tori-cells}) and $L\ge5$ (stage~\ref{it:order-choice-4}), either $\mathrm{per}=1$ or $\mathrm{per}\ge5$; thus the assumption $m_{\mathrm{p}}\ge2$ excludes exactly the case $\mathrm{per}=1$. The cells of Section~\ref{sec:reference-tori-cells} on the torus $k$ are then those of the cell structure of width $\boldsymbol{\lambda}=\mathrm{per}/2=m_{\mathrm{p}}+\frac{1}{2}$, with $m_{\mathrm{w}}=m_{\mathrm{p}}$ and $N_{\mu}=2kL$ for every $\mu$. More precisely, the interval $I_i$ of Section~\ref{sec:reference-tori-cells} is the set of the integers in $\left[\boldsymbol{\lambda}\cdot i,\boldsymbol{\lambda}\cdot(i+1)\right]$; the maps $\mathsf{g}_c$ of Section~\ref{sec:reference-tori-cells} are those of Definition~\ref{def:decomposition-torus-cells}; the set $P^o(\mathfrak{c}_c)$ is contained in $P^{+}(\mathfrak{c}_c)$, and a plaquette of $P^{+}(\mathfrak{c}_c)$ lies outside $P^o(\mathfrak{c}_c)$ if and only if it is contained in a plane $\{x_{\mu}=t\}$ such that $t$ is the end point of $I_{c_{\mu}}$ at distance $\frac{1}{2}$ from a link plane; and $\mathfrak{A}\subset\mathfrak{A}^{+}$, the mirror copy $O^{[c]}$ of any $O\in\mathfrak{A}$ lying in $\mathfrak{A}(\mathfrak{c}_c)$.

\subsubsection{Slabs and their operators}\label{sec:slabs-operators}

\begin{lemma}[Dependence on the boundary slices]\label{lem:dependence-boundary-slices}
Let $\mathsf{m}\ge2$ be an integer, and let
\[
P\colon\left(\mathbb{R}^l\right)^{\mathsf{m}}\times\left(\mathbb{R}^{\mathrm{plaq}(\boldsymbol{\Sigma})}\right)^{\mathsf{m}+1}\to\mathbb{C}
\]
be a bounded Borel function of the variables $\left(\mathrm{El}_0,\dots,\mathrm{El}_{\mathsf{m}-1};\tilde{b}_0,\dots,\tilde{b}_{\mathsf{m}}\right)$, the $\mathrm{El}_j\in\mathbb{R}^l$ being one-cochains and the $\tilde{b}_j\in\mathbb{R}^{\mathrm{plaq}(\boldsymbol{\Sigma})}$ two-cochains of the cross-section $\boldsymbol{\Sigma}$ of~\eqref{it:free-coefficient-Q5}. Define the function $\widetilde{P}\colon\left(\mathbb{R}^l\right)^{\mathsf{m}}\times\left(\mathbb{R}^{\mathrm{plaq}(\boldsymbol{\Sigma})}\right)^{\mathsf{m}-1}\to\mathbb{C}$ and the function $P^{\boldsymbol{\theta}}$, which has the same domain as $P$, by
\begin{align*}
\widetilde{P}\left(\mathrm{El}_0,\dots,\mathrm{El}_{\mathsf{m}-1};\tilde{b}_1,\dots,\tilde{b}_{\mathsf{m}-1}\right)&:=P\left(\mathrm{El}_0,\dots,\mathrm{El}_{\mathsf{m}-1};\,\tilde{b}_1-d_s\mathrm{El}_0,\,\tilde{b}_1,\dots,\tilde{b}_{\mathsf{m}-1},\,\tilde{b}_{\mathsf{m}-1}+d_s\mathrm{El}_{\mathsf{m}-1}\right),\\
P^{\boldsymbol{\theta}}\left(\mathrm{El}_0,\dots,\mathrm{El}_{\mathsf{m}-1};\tilde{b}_0,\dots,\tilde{b}_{\mathsf{m}}\right)&:=\overline{P\left(-\mathrm{El}_{\mathsf{m}-1},\dots,-\mathrm{El}_0;\,\tilde{b}_{\mathsf{m}},\dots,\tilde{b}_0\right)},
\end{align*}
and put $C_{\mathsf{m},\boldsymbol{\vartheta}}[P]:=C_{\mathsf{m},\boldsymbol{\vartheta}}\left[\widetilde{P}\right]$ for $\boldsymbol{\vartheta}\in\mathbb{R}^2$.
\begin{enumerate}
\item[(a)] The function $\widetilde{P}$ is a slab function of thickness $\mathsf{m}$ in the sense of~\eqref{it:free-coefficient-Q5}. If $P$ depends neither on $\tilde{b}_0$ nor on $\tilde{b}_{\mathsf{m}}$, then $\widetilde{P}=P$, the function $P$ being regarded as a function of $\left(\mathrm{El}_0,\dots,\mathrm{El}_{\mathsf{m}-1};\tilde{b}_1,\dots,\tilde{b}_{\mathsf{m}-1}\right)$.
\item[(b)] One has $\widetilde{P^{\boldsymbol{\theta}}}=\left(\widetilde{P}\right)^{\boldsymbol{\theta}}$. Consequently $C_{\mathsf{m},\boldsymbol{\vartheta}}\left[P^{\boldsymbol{\theta}}\right]=C_{\mathsf{m},\boldsymbol{\vartheta}}[P]^{*}$ for every $\boldsymbol{\vartheta}\in\mathbb{R}^2$.
\item[(c)] For $F\in C^2(n)$ and $t\in\mathbb{Z}/n_0$ let $\mathrm{El}_t\in\mathbb{R}^l$ and $\tilde{b}_t\in\mathbb{R}^{\mathrm{plaq}(\boldsymbol{\Sigma})}$, the electric and the magnetic part of $F$ in the slice $t$, be given by
\[
\mathrm{El}_t(y,k):=F_{0k}(t,y),\qquad\tilde{b}_t(y):=F_{12}(t,y)\qquad\left(y\in\boldsymbol{\Sigma},\ k\in\{1,2\}\right).
\]
If $F$ is closed, then $\tilde{b}_{t+1}-\tilde{b}_t=d_s\mathrm{El}_t$ for every $t\in\mathbb{Z}/n_0$. Consequently, for every closed $F\in C^2(n)$ and every $t_0\in\mathbb{Z}/n_0$,
\[
P\left(\mathrm{El}_{t_0},\dots,\mathrm{El}_{t_0+\mathsf{m}-1};\tilde{b}_{t_0},\dots,\tilde{b}_{t_0+\mathsf{m}}\right)=\widetilde{P}\left(\mathrm{El}_{t_0},\dots,\mathrm{El}_{t_0+\mathsf{m}-1};\tilde{b}_{t_0+1},\dots,\tilde{b}_{t_0+\mathsf{m}-1}\right),
\]
that is, $P(F)=\widetilde{P}(F)$ when both functions are evaluated on the slices $t_0,\dots,t_0+\mathsf{m}$ of $F$.
\end{enumerate}
\end{lemma}

\begin{proof}
(a) The map that sends $\left(\mathrm{El}_0,\dots,\mathrm{El}_{\mathsf{m}-1};\tilde{b}_1,\dots,\tilde{b}_{\mathsf{m}-1}\right)$ to the argument of $P$ in the definition of $\widetilde{P}$ is linear. Hence $\widetilde{P}$ is a bounded Borel function of $\left(\mathrm{El}_0,\dots,\mathrm{El}_{\mathsf{m}-1};\tilde{b}_1,\dots,\tilde{b}_{\mathsf{m}-1}\right)$, that is, a slab function of thickness $\mathsf{m}$. In that argument the entries $\tilde{b}_1-d_s\mathrm{El}_0$ and $\tilde{b}_{\mathsf{m}-1}+d_s\mathrm{El}_{\mathsf{m}-1}$ occupy the places of $\tilde{b}_0$ and of $\tilde{b}_{\mathsf{m}}$, while every other entry is the corresponding variable itself. So $\widetilde{P}=P$ if $P$ depends on neither of these two places.

(b) By the definition of the mirror of a slab function in~\eqref{it:free-coefficient-Q5} and by the definition of $\widetilde{P}$,
\begin{align*}
\left(\widetilde{P}\right)^{\boldsymbol{\theta}}\left(\mathrm{El}_0,\dots,\mathrm{El}_{\mathsf{m}-1};\tilde{b}_1,\dots,\tilde{b}_{\mathsf{m}-1}\right)&=\overline{\widetilde{P}\left(-\mathrm{El}_{\mathsf{m}-1},\dots,-\mathrm{El}_0;\,\tilde{b}_{\mathsf{m}-1},\dots,\tilde{b}_1\right)}\\
&=\overline{P\left(-\mathrm{El}_{\mathsf{m}-1},\dots,-\mathrm{El}_0;\,\tilde{b}_{\mathsf{m}-1}+d_s\mathrm{El}_{\mathsf{m}-1},\,\tilde{b}_{\mathsf{m}-1},\dots,\tilde{b}_1,\,\tilde{b}_1-d_s\mathrm{El}_0\right)};
\end{align*}
indeed, the first magnetic entry is $\tilde{b}_{\mathsf{m}-1}-d_s\left(-\mathrm{El}_{\mathsf{m}-1}\right)$, and the last one is $\tilde{b}_1+d_s\left(-\mathrm{El}_0\right)$. On the other hand, by the definition of $\widetilde{P^{\boldsymbol{\theta}}}$ and then by that of $P^{\boldsymbol{\theta}}$,
\begin{align*}
\widetilde{P^{\boldsymbol{\theta}}}\left(\mathrm{El}_0,\dots,\mathrm{El}_{\mathsf{m}-1};\tilde{b}_1,\dots,\tilde{b}_{\mathsf{m}-1}\right)&=P^{\boldsymbol{\theta}}\left(\mathrm{El}_0,\dots,\mathrm{El}_{\mathsf{m}-1};\,\tilde{b}_1-d_s\mathrm{El}_0,\,\tilde{b}_1,\dots,\tilde{b}_{\mathsf{m}-1},\,\tilde{b}_{\mathsf{m}-1}+d_s\mathrm{El}_{\mathsf{m}-1}\right)\\
&=\overline{P\left(-\mathrm{El}_{\mathsf{m}-1},\dots,-\mathrm{El}_0;\,\tilde{b}_{\mathsf{m}-1}+d_s\mathrm{El}_{\mathsf{m}-1},\,\tilde{b}_{\mathsf{m}-1},\dots,\tilde{b}_1,\,\tilde{b}_1-d_s\mathrm{El}_0\right)},
\end{align*}
which is the same expression. By (a), $\widetilde{P}$ is a slab function of thickness $\mathsf{m}$, so that the identity~\eqref{it:free-coefficient-Q5}(ii) applies to it, and
\[
C_{\mathsf{m},\boldsymbol{\vartheta}}\left[P^{\boldsymbol{\theta}}\right]=C_{\mathsf{m},\boldsymbol{\vartheta}}\left[\widetilde{P^{\boldsymbol{\theta}}}\right]=C_{\mathsf{m},\boldsymbol{\vartheta}}\left[\left(\widetilde{P}\right)^{\boldsymbol{\theta}}\right]=C_{\mathsf{m},\boldsymbol{\vartheta}}\left[\widetilde{P}\right]^{*}=C_{\mathsf{m},\boldsymbol{\vartheta}}[P]^{*}.
\]

(c) Let $e_{\mu}$ denote the unit vector of the direction $\mu$, and let $x=(t,y)\in\Lambda_n$ with $y\in\boldsymbol{\Sigma}$. By the formula for the coboundary of two-cochains in Section~\ref{sec:tori-cochains-lattice},
\[
(dF)(x)=\left[F_{12}(x+e_0)-F_{12}(x)\right]-\left[F_{02}(x+e_1)-F_{02}(x)\right]+\left[F_{01}(x+e_2)-F_{01}(x)\right].
\]
By the formula for the coboundary of one-cochains in the same section, applied on the two-dimensional torus $\boldsymbol{\Sigma}$,
\begin{align*}
\left(d_s\mathrm{El}_t\right)(y)&=\mathrm{El}_t(y,1)+\mathrm{El}_t(y+e_1,2)-\mathrm{El}_t(y+e_2,1)-\mathrm{El}_t(y,2)\\
&=\left[F_{02}(x+e_1)-F_{02}(x)\right]-\left[F_{01}(x+e_2)-F_{01}(x)\right].
\end{align*}
Hence $(dF)(t,y)=\tilde{b}_{t+1}(y)-\tilde{b}_t(y)-\left(d_s\mathrm{El}_t\right)(y)$, and $dF=0$ gives $\tilde{b}_{t+1}-\tilde{b}_t=d_s\mathrm{El}_t$. In particular
\[
\tilde{b}_{t_0}=\tilde{b}_{t_0+1}-d_s\mathrm{El}_{t_0},\qquad\tilde{b}_{t_0+\mathsf{m}}=\tilde{b}_{t_0+\mathsf{m}-1}+d_s\mathrm{El}_{t_0+\mathsf{m}-1},
\]
and inserting these two identities in the definition of $\widetilde{P}$ gives the last assertion.
\end{proof}

Let $\boldsymbol{\lambda}=m_{\mathrm{w}}+\frac{1}{2}$ be a cell width (Definition~\ref{def:decomposition-torus-cells}; $m_{\mathrm{w}}\ge2$ is an integer), and let $n_0=\boldsymbol{\lambda}\cdot N_0$ with an even positive integer $N_0$; no condition is imposed on $n_1$ and $n_2$. For $i\in\mathbb{Z}/N_0$ let $\mathrm{Pl}_{0,i}=\left\{x_0=\boldsymbol{\lambda}\cdot i\right\}$ be the planes of Definition~\ref{def:decomposition-torus-cells} in the direction $0$. Since $\boldsymbol{\lambda}\cdot i$ is an integer exactly for even $i$, the plane $\mathrm{Pl}_{0,i}$ is a site plane if $i$ is even and a link plane if $i$ is odd. Put
\[
t_i^{-}:=\left\lceil\boldsymbol{\lambda}\cdot i\right\rceil,\qquad t_i^{+}:=\left\lfloor\boldsymbol{\lambda}\cdot(i+1)\right\rfloor\qquad\left(i\in\mathbb{Z}/N_0\right).
\]
The \emph{slab} $i$ consists of the slices $t_i^{-},t_i^{-}+1,\dots,t_i^{+}$. For even $i$ one has $t_i^{-}=\boldsymbol{\lambda}\cdot i$ and $t_i^{+}=\boldsymbol{\lambda}\cdot i+m_{\mathrm{w}}$; for odd $i$ one has $t_i^{-}=\boldsymbol{\lambda}\cdot i+\frac{1}{2}$ and $t_i^{+}=\boldsymbol{\lambda}\cdot(i+1)$. In both cases the slab $i$ has the thickness $t_i^{+}-t_i^{-}=m_{\mathrm{w}}\ge2$.

Let $i\in\mathbb{Z}/N_0$, and consider the plaquettes that have all their vertices in the slab $i$ and lie in none of the planes $\mathrm{Pl}_{0,j}$, $j\in\mathbb{Z}/N_0$. For $F\in C^2(n)$, the values of $F$ on these plaquettes are the entries of the cochains
\[
\mathrm{El}_t\quad\left(t_i^{-}\le t<t_i^{+}\right)\qquad\text{and}\qquad\tilde{b}_t\quad\left(t_i^{-}\le t\le t_i^{+},\ t\notin\boldsymbol{\lambda}\mathbb{Z}\right)
\]
of Lemma~\ref{lem:dependence-boundary-slices}(c). A bounded Borel function $P\colon C^2(n)\to\mathbb{C}$ is called a \emph{slab function at $i$} if $P(F)$ depends on $F$ only through these cochains. In terms of the shifted time index $s=t-t_i^{-}$, such a $P$ is a function of $\left(\mathrm{El}_0,\dots,\mathrm{El}_{m_{\mathrm{w}}-1};\tilde{b}_0,\dots,\tilde{b}_{m_{\mathrm{w}}}\right)$ of the kind considered in Lemma~\ref{lem:dependence-boundary-slices}, with $\mathsf{m}=m_{\mathrm{w}}$, and $P(F)$ is obtained by evaluating this function on the slices $t_i^{-},\dots,t_i^{+}$ of $F$.

Let $\boldsymbol{\vartheta}\in\mathbb{R}^2$. In what follows, $\Pi$ stands for one of the following operators on $\mathscr{H}_{\boldsymbol{\vartheta}}$: the identity, or, for some $w_{12}\in\mathbb{Z}$, the orthogonal projection with range $\mathscr{H}_{\boldsymbol{\vartheta}}^{(w_{12})}$, a summand of the decomposition of $\mathscr{H}_{\boldsymbol{\vartheta}}$ in~\eqref{it:free-coefficient-Q5}. By~\eqref{it:free-coefficient-Q5}(i) and~(iii), each operator $C_{\mathsf{m},\boldsymbol{\vartheta}}[P]$ with a slab function $P$ of thickness $\mathsf{m}$ is bounded, being Hilbert--Schmidt, and maps every summand $\mathscr{H}_{\boldsymbol{\vartheta}}^{(w_{12})}$ into itself; by its definition in Lemma~\ref{lem:dependence-boundary-slices}, the operator $C_{\mathsf{m},\boldsymbol{\vartheta}}[P]=C_{\mathsf{m},\boldsymbol{\vartheta}}\left[\widetilde{P}\right]$ of a function $P$ as in that lemma has the same two properties. A bounded operator that maps every summand of an orthogonal decomposition into itself commutes with the orthogonal projection onto each summand. Consequently $\Pi\,C_{\mathsf{m},\boldsymbol{\vartheta}}[P]=C_{\mathsf{m},\boldsymbol{\vartheta}}[P]\,\Pi$ for all these operators; for instance, $\Pi$ commutes with $T_{\boldsymbol{\vartheta}}=C_{1,\boldsymbol{\vartheta}}[1]$. For a slab function $P$ at $i$ define the operator $Y_i[P]$ on $\mathscr{H}_{\boldsymbol{\vartheta}}$ by
\[
Y_i[P]:=X_{\boldsymbol{\vartheta}}\,\Pi\,C_{m_{\mathrm{w}},\boldsymbol{\vartheta}}[P]\quad\text{if $i$ is even},\qquad Y_i[P]:=C_{m_{\mathrm{w}},\boldsymbol{\vartheta}}[P]\,\Pi\,X_{\boldsymbol{\vartheta}}^{*}\quad\text{if $i$ is odd},
\]
where $X_{\boldsymbol{\vartheta}}$ is the operator of~\eqref{it:free-coefficient-Q5}(iv) and $C_{m_{\mathrm{w}},\boldsymbol{\vartheta}}[P]$ is defined as in Lemma~\ref{lem:dependence-boundary-slices}. Finally, for $\Phi\in\mathfrak{F}(n)$ put
\[
\boldsymbol{\omega}_n^{\boldsymbol{\vartheta},\Pi}(\Phi):=\sum_w\exp\left(2\pi i\left(\boldsymbol{\vartheta}_1\cdot w_{01}+\boldsymbol{\vartheta}_2\cdot w_{02}\right)\right)\cdot\boldsymbol{\omega}_n^{(w)}(\Phi),
\]
the sum being extended over all $w=\left(w_{12},w_{02},w_{01}\right)\in\mathbb{Z}^3$ if $\Pi$ is the identity, and over those $w\in\mathbb{Z}^3$ whose component $w_{12}$ is the prescribed one if $\Pi$ is the projection onto $\mathscr{H}_{\boldsymbol{\vartheta}}^{(w_{12})}$. The series converges absolutely, since $\sum_w\left|\boldsymbol{\omega}_n^{(w)}(\Phi)\right|\le\boldsymbol{\omega}_n\left(|\Phi|\right)$ by~\eqref{it:free-coefficient-Q5}.

\begin{lemma}\label{lem:free-coefficient-2}
Let $n$, $\boldsymbol{\vartheta}$ and $\Pi$ be as above. Every operator $Y_i[P]$ is a Hilbert--Schmidt operator on $\mathscr{H}_{\boldsymbol{\vartheta}}$, and the following assertions hold.
\begin{enumerate}
\item[(a)] Let $P_i$ be a slab function at $i$ for every $i\in\mathbb{Z}/N_0$. Then
\[
k_B^{n_0}\cdot\operatorname{Tr}_{\mathscr{H}_{\boldsymbol{\vartheta}}}\left(Y_{N_0-1}\left[P_{N_0-1}\right]\cdots Y_1\left[P_1\right]Y_0\left[P_0\right]\right)=\boldsymbol{\omega}_n^{\boldsymbol{\vartheta},\Pi}\left(\prod_{i\in\mathbb{Z}/N_0}P_i\right).
\]
\item[(b)] Let $P$ be a slab function at $i$. Then $P^{\boldsymbol{\theta}}$ is a slab function at $i+1$ and at $i-1$, and
\[
Y_i[P]^{*}=Y_{i+1}\left[P^{\boldsymbol{\theta}}\right]=Y_{i-1}\left[P^{\boldsymbol{\theta}}\right].
\]
If $P$ is evaluated on the slices of the slab $i$ and $P^{\boldsymbol{\theta}}$ on the slices of the slab $i+1$, respectively $i-1$, then
\[
P^{\boldsymbol{\theta}}(F)=\overline{P\left(\mathfrak{r}^{*}F\right)}\qquad\left(F\in C^2(n)\right),
\]
where $\mathfrak{r}$ denotes the reflection of $\Lambda_n$ in $\mathrm{Pl}_{0,i+1}$, respectively in $\mathrm{Pl}_{0,i}$, which is the plane separating the two slabs.
\item[(c)] The reflections of $\Lambda_n$ in the planes $\mathrm{Pl}_{0,j}$, $j\in\mathbb{Z}/N_0$, generate a group $\mathcal{G}$ of isometries of $\Lambda_n$, which acts simply transitively on the slabs by Definition~\ref{def:decomposition-torus-cells}. Let $P$ be a slab function at $i$, evaluated on the slices of the slab $i$, and for $i'\in\mathbb{Z}/N_0$ let $g_{i'}$ denote the unique element of $\mathcal{G}$ that carries the slab $i$ to the slab $i'$. Put
\[
P^{(i')}(F):=P\left(g_{i'}^{*}F\right)\quad\text{if $i'-i$ is even},\qquad P^{(i')}(F):=\overline{P\left(g_{i'}^{*}F\right)}\quad\text{if $i'-i$ is odd}.
\]
Then
\[
k_B^{n_0}\cdot\left\|Y_i[P]\right\|_{N_0}^{N_0}=\boldsymbol{\omega}_n^{\boldsymbol{\vartheta},\Pi}\left(\prod_{i'\in\mathbb{Z}/N_0}P^{(i')}\right)\ge0.
\]
\end{enumerate}
\end{lemma}

\begin{proof}
By~\eqref{it:free-coefficient-Q5}(iv) the operator $X_{\boldsymbol{\vartheta}}$ is Hilbert--Schmidt, hence bounded, and so is $X_{\boldsymbol{\vartheta}}^{*}$. By Lemma~\ref{lem:dependence-boundary-slices}(a) and~\eqref{it:free-coefficient-Q5}(i) the operator $C_{m_{\mathrm{w}},\boldsymbol{\vartheta}}[P]=C_{m_{\mathrm{w}},\boldsymbol{\vartheta}}\left[\widetilde{P}\right]$ is Hilbert--Schmidt. As $\Pi$ is bounded, every $Y_i[P]$ is a Hilbert--Schmidt operator.

(a) By the definition in Section~\ref{sec:wilson-determinant-lattice}, $\boldsymbol{\omega}_n^{(w)}(\Phi)$ depends only on the values of $\Phi$ at closed two-cochains, because the curvature $F^u$ in that definition satisfies $dF^u=0$. By Lemma~\ref{lem:dependence-boundary-slices}(c), applied with $\mathsf{m}=m_{\mathrm{w}}$ and $t_0=t_i^{-}$, the function $P_i$ coincides on the closed two-cochains with $\widetilde{P_i}$, both being evaluated on the slices of the slab $i$. Therefore
\[
\boldsymbol{\omega}_n^{(w)}\left(\prod_iP_i\right)=\boldsymbol{\omega}_n^{(w)}\left(\prod_i\widetilde{P_i}\right)\qquad\left(w\in\mathbb{Z}^3\right).
\]
Let $i$ be even. By~\eqref{it:free-coefficient-Q5}(iv) one has $X_{\boldsymbol{\vartheta}}^{*}X_{\boldsymbol{\vartheta}}=T_{\boldsymbol{\vartheta}}=C_{1,\boldsymbol{\vartheta}}[1]$, and $\Pi^2=\Pi$ commutes with $T_{\boldsymbol{\vartheta}}$ and with $C_{m_{\mathrm{w}},\boldsymbol{\vartheta}}\left[P_{i+1}\right]$; hence
\[
Y_{i+1}\left[P_{i+1}\right]Y_i\left[P_i\right]=C_{m_{\mathrm{w}},\boldsymbol{\vartheta}}\left[P_{i+1}\right]\,\Pi\,X_{\boldsymbol{\vartheta}}^{*}X_{\boldsymbol{\vartheta}}\,\Pi\,C_{m_{\mathrm{w}},\boldsymbol{\vartheta}}\left[P_i\right]=\Pi\,C_{m_{\mathrm{w}},\boldsymbol{\vartheta}}\left[P_{i+1}\right]C_{1,\boldsymbol{\vartheta}}[1]\,C_{m_{\mathrm{w}},\boldsymbol{\vartheta}}\left[P_i\right].
\]
The product $Y_{N_0-1}\left[P_{N_0-1}\right]\cdots Y_0\left[P_0\right]$ is the product of the $\frac{N_0}{2}$ factors $Y_{2k+1}\left[P_{2k+1}\right]Y_{2k}\left[P_{2k}\right]$, $0\le k<\frac{N_0}{2}$. Since $\Pi$ is a projection that commutes with all the operators $C_{m_{\mathrm{w}},\boldsymbol{\vartheta}}\left[P_i\right]$ and with $C_{1,\boldsymbol{\vartheta}}[1]$, this product equals $\Pi\,W$, where, by the definition $C_{m_{\mathrm{w}},\boldsymbol{\vartheta}}\left[P_i\right]=C_{m_{\mathrm{w}},\boldsymbol{\vartheta}}\left[\widetilde{P_i}\right]$ of Lemma~\ref{lem:dependence-boundary-slices},
\[
W:=\prod_{k=\frac{N_0}{2}-1}^0\left(C_{m_{\mathrm{w}},\boldsymbol{\vartheta}}\left[\widetilde{P_{2k+1}}\right]C_{1,\boldsymbol{\vartheta}}[1]\,C_{m_{\mathrm{w}},\boldsymbol{\vartheta}}\left[\widetilde{P_{2k}}\right]\right),
\]
the factors with larger $k$ standing to the left. Hence $W$ is a product of $\frac{3N_0}{2}\ge3$ operators $C_{\mathsf{m},\boldsymbol{\vartheta}}[Q]$, each with a slab function $Q$ of thickness $\mathsf{m}$. Read from the right to the left, which is the order of increasing time, these factors have the thicknesses $m_{\mathrm{w}},1,m_{\mathrm{w}}$, repeated $\frac{N_0}{2}$ times, and the functions $\widetilde{P_0},1,\widetilde{P_1},\widetilde{P_2},1,\widetilde{P_3},\dots,\widetilde{P_{N_0-2}},1,\widetilde{P_{N_0-1}}$; the thicknesses add up to $(2m_{\mathrm{w}}+1)\cdot\frac{N_0}{2}=n_0$. In the sense of~\eqref{it:free-coefficient-Q5}(iii), for every $j\in\mathbb{Z}/N_0$ the factor of $\widetilde{P_j}$ belongs to the slab $\left\{t_j^{-}\le x_0\le t_j^{+}\right\}$, which is the slab $j$ defined above, and for even $j$ the factor $C_{1,\boldsymbol{\vartheta}}[1]$ next to it on the left belongs to $\left\{t_j^{+}\le x_0\le t_{j+1}^{-}\right\}$, a single step ($t_{j+1}^{-}=t_j^{+}+1$) across the link plane $\mathrm{Pl}_{0,j+1}$. If $\Pi$ is the identity,~\eqref{it:free-coefficient-Q5}(iii) therefore yields
\[
k_B^{n_0}\cdot\operatorname{Tr}_{\mathscr{H}_{\boldsymbol{\vartheta}}}(W)=\sum_{w\in\mathbb{Z}^3}\exp\left(2\pi i\left(\boldsymbol{\vartheta}_1\cdot w_{01}+\boldsymbol{\vartheta}_2\cdot w_{02}\right)\right)\cdot\boldsymbol{\omega}_n^{(w)}\left(\prod_{j\in\mathbb{Z}/N_0}\widetilde{P_j}\right).
\]
If $\Pi$ is the projection onto $\mathscr{H}_{\boldsymbol{\vartheta}}^{(w_{12})}$, then $W$ maps this summand into itself, $\operatorname{Tr}_{\mathscr{H}_{\boldsymbol{\vartheta}}}\left(\Pi\,W\right)$ is the trace of the restriction of $W$ to it, and by~\eqref{it:free-coefficient-Q5}(iii) this trace, multiplied by $k_B^{n_0}$, is the part of the above sum with the given $w_{12}$. In both cases the left side of (a) equals $\boldsymbol{\omega}_n^{\boldsymbol{\vartheta},\Pi}\left(\prod_i\widetilde{P_i}\right)=\boldsymbol{\omega}_n^{\boldsymbol{\vartheta},\Pi}\left(\prod_iP_i\right)$.

(b) Let $P$ be a slab function at $i$. By Lemma~\ref{lem:dependence-boundary-slices}(b), $C_{m_{\mathrm{w}},\boldsymbol{\vartheta}}[P]^{*}=C_{m_{\mathrm{w}},\boldsymbol{\vartheta}}\left[P^{\boldsymbol{\theta}}\right]$; moreover $\Pi^{*}=\Pi$. Hence
\[
\left(X_{\boldsymbol{\vartheta}}\,\Pi\,C_{m_{\mathrm{w}},\boldsymbol{\vartheta}}[P]\right)^{*}=C_{m_{\mathrm{w}},\boldsymbol{\vartheta}}\left[P^{\boldsymbol{\theta}}\right]\Pi\,X_{\boldsymbol{\vartheta}}^{*},\qquad\left(C_{m_{\mathrm{w}},\boldsymbol{\vartheta}}[P]\,\Pi\,X_{\boldsymbol{\vartheta}}^{*}\right)^{*}=X_{\boldsymbol{\vartheta}}\,\Pi\,C_{m_{\mathrm{w}},\boldsymbol{\vartheta}}\left[P^{\boldsymbol{\theta}}\right].
\]
In the definition of $P^{\boldsymbol{\theta}}$ the variable $\tilde{b}_{m_{\mathrm{w}}}$ of $P^{\boldsymbol{\theta}}$ is inserted at the place of $\tilde{b}_0$ of $P$, and the variable $\tilde{b}_0$ of $P^{\boldsymbol{\theta}}$ at the place of $\tilde{b}_{m_{\mathrm{w}}}$ of $P$. If $i$ is even, the slice $t_i^{-}=\boldsymbol{\lambda}\cdot i$ lies in the site plane $\mathrm{Pl}_{0,i}$, which bounds the slab $i$ from below; so $P$ does not depend on $\tilde{b}_0$, and $P^{\boldsymbol{\theta}}$ does not depend on $\tilde{b}_{m_{\mathrm{w}}}$. If $i$ is odd, the slice $t_i^{+}=\boldsymbol{\lambda}\cdot(i+1)$ lies in the site plane $\mathrm{Pl}_{0,i+1}$, which bounds the slab $i$ from above; so $P$ does not depend on $\tilde{b}_{m_{\mathrm{w}}}$, and $P^{\boldsymbol{\theta}}$ does not depend on $\tilde{b}_0$. Thus $P^{\boldsymbol{\theta}}$ is a slab function at every slab whose index has the parity opposite to that of $i$, in particular at $i+1$ and at $i-1$, and the two displayed identities state that $Y_i[P]^{*}=Y_{i\pm1}\left[P^{\boldsymbol{\theta}}\right]$.

For the second assertion, put $s_0:=t_i^{-}$, and let $s$ be $t_{i+1}^{-}$, respectively $t_{i-1}^{-}$; thus $P$ is a function of the slices $s_0,\dots,s_0+m_{\mathrm{w}}$ of $F$, and $P^{\boldsymbol{\theta}}$ a function of the slices $s,\dots,s+m_{\mathrm{w}}$. Then
\[
P^{\boldsymbol{\theta}}(F)=\overline{P\left(-\mathrm{El}_{s+m_{\mathrm{w}}-1},\dots,-\mathrm{El}_s;\,\tilde{b}_{s+m_{\mathrm{w}}},\dots,\tilde{b}_s\right)},
\]
so that the electric argument of index $j$ of $P$ is $-\mathrm{El}_{s+m_{\mathrm{w}}-1-j}$, and its magnetic argument of index $j$ is $\tilde{b}_{s+m_{\mathrm{w}}-j}$. Let $\mathfrak{r}$ be the reflection $x_0\mapsto c-x_0$ of $\Lambda_n$ with $c:=s_0+s+m_{\mathrm{w}}$. By Section~\ref{sec:tori-cochains-lattice} its pull-back acts by
\[
\left(\mathfrak{r}^{*}F\right)_{12}(t,y)=F_{12}(c-t,y),\qquad\left(\mathfrak{r}^{*}F\right)_{0k}(t,y)=-F_{0k}(c-t-1,y)\qquad(k=1,2),
\]
the components that contain the direction $0$ changing sign. Hence $\mathrm{El}_{s_0+j}\left(\mathfrak{r}^{*}F\right)=-\mathrm{El}_{s+m_{\mathrm{w}}-1-j}(F)$ and $\tilde{b}_{s_0+j}\left(\mathfrak{r}^{*}F\right)=\tilde{b}_{s+m_{\mathrm{w}}-j}(F)$, and it follows that $P^{\boldsymbol{\theta}}(F)=\overline{P\left(\mathfrak{r}^{*}F\right)}$. It remains to identify $\mathfrak{r}$. The values of $s_0$, $s$ and $\frac{c}{2}$ are as follows.
\begin{itemize}
\item If $i$ is even, then $s_0=\boldsymbol{\lambda}\cdot i$, and the slice $s_0$ lies in the site plane $\mathrm{Pl}_{0,i}$. For the slab $i+1$ one has $s=s_0+m_{\mathrm{w}}+1$ and $\frac{c}{2}=s_0+m_{\mathrm{w}}+\frac{1}{2}=\boldsymbol{\lambda}\cdot(i+1)$; for the slab $i-1$ one has $s=s_0-m_{\mathrm{w}}$ and $\frac{c}{2}=s_0=\boldsymbol{\lambda}\cdot i$.
\item If $i$ is odd, then $s_0=\boldsymbol{\lambda}\cdot i+\frac{1}{2}$. For the slab $i+1$ one has $s=s_0+m_{\mathrm{w}}$ and $\frac{c}{2}=s_0+m_{\mathrm{w}}=\boldsymbol{\lambda}\cdot(i+1)$; for the slab $i-1$ one has $s=s_0-m_{\mathrm{w}}-1$ and $\frac{c}{2}=s_0-\frac{1}{2}=\boldsymbol{\lambda}\cdot i$.
\end{itemize}
So $c=2\boldsymbol{\lambda}\cdot(i+1)$ if $P^{\boldsymbol{\theta}}$ is evaluated on the slab $i+1$, and $c=2\boldsymbol{\lambda}\cdot i$ if it is evaluated on the slab $i-1$. Accordingly $\mathfrak{r}$, the map $x_0\mapsto c-x_0$, is the reflection in $\mathrm{Pl}_{0,i+1}$, respectively in $\mathrm{Pl}_{0,i}$, in the sense of Definition~\ref{def:decomposition-torus-cells}.

(c) Since $N_0$ is even and $Y:=Y_i[P]$ is a Hilbert--Schmidt operator,~\eqref{it:free-coefficient-Q8} gives
\[
\left\|Y\right\|_{N_0}^{N_0}=\operatorname{Tr}\left(Y^{*}Y\right)^{\frac{N_0}{2}}.
\]
By their definition, $Y_{i_1}[\cdot]$ and $Y_{i_2}[\cdot]$ are the same map whenever $i_1-i_2$ is even. Suppose first that $i$ is even. Then $Y=Y_{i'}[P]$ for every even $i'$, and $Y^{*}=Y_{i'}\left[P^{\boldsymbol{\theta}}\right]$ for every odd $i'$ by (b). Hence $\left(Y^{*}Y\right)^{\frac{N_0}{2}}$, a product of $N_0$ factors in which $Y^{*}$ and $Y$ alternate, is the product $Y_{N_0-1}\left[P_{N_0-1}\right]\cdots Y_0\left[P_0\right]$ that appears in (a), formed with $P_{i'}:=P$ at every even place $i'$ and $P_{i'}:=P^{\boldsymbol{\theta}}$ at every odd place $i'$. Suppose next that $i$ is odd. Then $Y=Y_1[P]$, and $Y^{*}=Y_0\left[P^{\boldsymbol{\theta}}\right]$ by (b); moving the first factor to the end, which leaves the trace unchanged, one obtains
\[
\operatorname{Tr}\left(Y_0\left[P^{\boldsymbol{\theta}}\right]Y_1[P]\right)^{\frac{N_0}{2}}=\operatorname{Tr}\left(Y_1[P]\,Y_0\left[P^{\boldsymbol{\theta}}\right]\right)^{\frac{N_0}{2}},
\]
and the operator under the trace on the right is the product that appears in (a), formed with $P_{i'}:=P^{\boldsymbol{\theta}}$ at every even place and $P_{i'}:=P$ at every odd place. For either parity of $i$, the family $\left(P_{i'}\right)_{i'\in\mathbb{Z}/N_0}$ so formed has $P_i=P$, and $P_{i''+1}=\left(P_{i''}\right)^{\boldsymbol{\theta}}$ for every $i''$, because $\left(P^{\boldsymbol{\theta}}\right)^{\boldsymbol{\theta}}=P$. Suppose that, for some $i''$, the function $P_{i''}$, evaluated on the slices of the slab $i''$, is $F\mapsto P\left(g'^{*}F\right)$ or its complex conjugate, where $g'\in\mathcal{G}$ carries the slab $i$ to the slab $i''$. By (b), applied to the slab function $P_{i''}$ at $i''$, the function $P_{i''+1}$, evaluated on the slices of the slab $i''+1$, is $F\mapsto P\left(g'^{*}\mathfrak{r}^{*}F\right)=P\left(\left(\mathfrak{r}\circ g'\right)^{*}F\right)$ with one complex conjugation more, $\mathfrak{r}$ being the reflection in $\mathrm{Pl}_{0,i''+1}$; and $\mathfrak{r}\circ g'\in\mathcal{G}$ carries the slab $i$ to the slab $i''+1$. Beginning with $i''=i$ and $g'$ the identity, and proceeding from each slab to the next, one reaches every slab $i'$ with an element of $\mathcal{G}$ that carries the slab $i$ to the slab $i'$, hence with $g_{i'}$, and with a number of complex conjugations congruent to $i'-i$ modulo $2$. Thus $P_{i'}=P^{(i')}$ for every $i'$, and (a) gives the identity of (c); its left side is non-negative.
\end{proof}

\begin{corollary}[Reflection positivity on a fixed torus]\label{cor:reflection-positivity-fixed}
Let $n_0=2\boldsymbol{\lambda}$, and let $\mathfrak{r}$ be the reflection $x_0\mapsto-x_0$ of $\Lambda_n$. Then, for every slab function $P$ at $0$,
\[
\boldsymbol{\omega}_n\left(\overline{P\circ\mathfrak{r}^{*}}\cdot P\right)\ge0\qquad\text{and}\qquad\boldsymbol{\omega}_n^{(0)}\left(\overline{P\circ\mathfrak{r}^{*}}\cdot P\right)\ge0,
\]
where $P\circ\mathfrak{r}^{*}$ denotes the function $F\mapsto P\left(\mathfrak{r}^{*}F\right)$.
\end{corollary}

\begin{proof}
Here $N_0=2$. On $\mathbb{Z}/n_0$ the maps $t\mapsto-t$ and $t\mapsto2\boldsymbol{\lambda}-t$ coincide, and the second one, the reflection in $\mathrm{Pl}_{0,1}$, maps the slab $0$ onto the slab $1$. In Lemma~\ref{lem:free-coefficient-2}(c) with $i=0$ one therefore has $g_1=\mathfrak{r}$, $P^{(0)}=P$ and $P^{(1)}=\overline{P\circ\mathfrak{r}^{*}}$, so that
\[
\boldsymbol{\omega}_n^{\boldsymbol{\vartheta},\Pi}\left(\overline{P\circ\mathfrak{r}^{*}}\cdot P\right)=k_B^{n_0}\cdot\left\|Y_0[P]\right\|_2^2\ge0
\]
for every $\boldsymbol{\vartheta}\in\mathbb{R}^2$ and every $\Pi$ admitted in the definition of $Y_i[\cdot]$, that is, for $\Pi$ the identity of $\mathscr{H}_{\boldsymbol{\vartheta}}$ and for $\Pi$ the projection onto any single summand $\mathscr{H}_{\boldsymbol{\vartheta}}^{(w_{12})}$, $w_{12}\in\mathbb{Z}$. If $\Pi$ is the identity and $\boldsymbol{\vartheta}=0$, the left side is $\sum_{w\in\mathbb{Z}^3}\boldsymbol{\omega}_n^{(w)}\left(\overline{P\circ\mathfrak{r}^{*}}\cdot P\right)=\boldsymbol{\omega}_n\left(\overline{P\circ\mathfrak{r}^{*}}\cdot P\right)$ (Section~\ref{sec:wilson-determinant-lattice}); this is the first inequality. Now let $\Pi$ be the projection onto the summand $\mathscr{H}_{\boldsymbol{\vartheta}}^{(w_{12})}$ with $w_{12}=0$. The series that defines $\boldsymbol{\omega}_n^{\boldsymbol{\vartheta},\Pi}\left(\overline{P\circ\mathfrak{r}^{*}}\cdot P\right)$ converges absolutely, uniformly in $\boldsymbol{\vartheta}$, by~\eqref{it:free-coefficient-Q5}. Integrating it term by term over $[0,1]^2$, one finds
\[
\int_{[0,1]^2}\mathrm{d}\boldsymbol{\vartheta}\,\boldsymbol{\omega}_n^{\boldsymbol{\vartheta},\Pi}\left(\overline{P\circ\mathfrak{r}^{*}}\cdot P\right)=\boldsymbol{\omega}_n^{(0)}\left(\overline{P\circ\mathfrak{r}^{*}}\cdot P\right),
\]
because only the term with $w_{01}=w_{02}=0$ remains after the integration. The integrand is non-negative; this is the second inequality.
\end{proof}

\subsubsection{Chessboard estimates, flux sectors and lengthened tori}\label{sec:chessboard-estimates-flux}

\paragraph{Reduction to bounded functions.}
For $O\in\mathfrak{A}^{+}$ and $\bar{L}>0$ let
\[
O_{\bar{L}}:=O\cdot\mathbf{1}\left[\,|O|\le\bar{L}\,\right].
\]
Then $O_{\bar{L}}$ is a bounded element of $\mathfrak{A}^{+}$, and $O_{\bar{L}}\to O$ pointwise as $\bar{L}\to\infty$. Truncation commutes with the formation of mirror copies: since $\left|\overline{z}\right|=|z|$ for $z\in\mathbb{C}$,
\[
\left(O_{\bar{L}}\right)^{[c]}=O^{[c]}\cdot\mathbf{1}\left[\,\left|O^{[c]}\right|\le\bar{L}\,\right]\qquad\text{for every cell }c.
\]
Consider one of the inequalities below, for observables $O_c\in\mathfrak{A}^{+}$, and assume that it holds whenever these observables are bounded; in particular it holds for the truncations $\left(O_c\right)_{\bar{L}}$. Every integrand that then occurs is bounded in modulus, uniformly in $\bar{L}$, by a product of functions $\left|O_c^{[c']}\right|$, where $c$ and $c'$ are cells. As $O_c\in\mathfrak{F}(n)$ and the pull-back $\mathsf{g}_{c'}^{*}$ preserves the norm of two-cochains, there is a constant $C$ such that $\left|O_c^{[c']}(F)\right|\le C\cdot\exp\left(C\cdot\|F\|\right)$ for all these functions and all $F\in C^2(n)$; hence a finite product of them belongs to $\mathfrak{F}(n)$ and is $\boldsymbol{\omega}_n$-integrable (Section~\ref{sec:wilson-determinant-lattice}). Moreover, by~\eqref{it:free-coefficient-Q5},
\[
\sum_{w\in\mathbb{Z}^3}\left|\boldsymbol{\omega}_n^{(w)}(\Phi)\right|\le\boldsymbol{\omega}_n\left(|\Phi|\right)\qquad\left(\Phi\in\mathfrak{F}(n)\right).
\]
Letting $\bar{L}\to\infty$, dominated convergence, applied to the integral over $A$ and the sum over $\mathfrak{z}$ that define $\boldsymbol{\omega}_n$ and $\boldsymbol{\omega}_n^{(w)}$ (Section~\ref{sec:wilson-determinant-lattice}) and to the series over $w$, carries the inequality from the truncated observables to the observables $O_c$ themselves. Hence, in the proof of each of these inequalities, the observables may be assumed to be bounded.

\begin{proposition}[Chessboard estimate]\label{prop:chessboard-estimate}
Let the torus $n$ carry the cell structure of a width $\boldsymbol{\lambda}$ (Definition~\ref{def:decomposition-torus-cells}), and let $N_n=N_0\cdot N_1\cdot N_2$ be the number of its cells. For every $O\in\mathfrak{A}^{+}$,
\[
\boldsymbol{\omega}_n\left(O^{\sharp}\right)\ge0\qquad\text{and}\qquad\boldsymbol{\omega}_n^{(0)}\left(O^{\sharp}\right)\ge0.
\]
Let $\left(O_c\right)_c$ be a family of elements of $\mathfrak{A}^{+}$ indexed by the cells $c$ of $n$, and let $\boldsymbol{\omega}$ be one of the two functionals $\boldsymbol{\omega}_n$, $\boldsymbol{\omega}_n^{(0)}$. Then
\[
\left|\boldsymbol{\omega}\left(\prod_cO_c^{[c]}\right)\right|\le\prod_c\boldsymbol{\omega}\left(O_c^{\sharp}\right)^{\frac{1}{N_n}}.
\]
\end{proposition}

\begin{proof}
By the truncation argument at the beginning of Section~\ref{sec:chessboard-estimates-flux}, it suffices to consider bounded observables. Since $n$ carries the cell structure, $n_0=\boldsymbol{\lambda}\cdot N_0$ with an even integer $N_0\ge2$; so the slabs $i\in\mathbb{Z}/N_0\mathbb{Z}$, the operators $Y_i[\cdot]$ and the functionals $\boldsymbol{\omega}_n^{\boldsymbol{\vartheta},\Pi}$ of Section~\ref{sec:slabs-operators} are at hand.

\emph{The direction $0$.} Let $\left(G_c\right)_c$ be a family of bounded functions with $G_c\in\mathfrak{A}^{+}(\mathfrak{c}_c)$ for every cell $c$, and put
\[
P_i:=\prod_{c:\,c_0=i}G_c\qquad\left(i\in\mathbb{Z}/N_0\mathbb{Z}\right).
\]
Each $P_i$ is a slab function at $i$. Indeed, let $c_0=i$ and $p\in P^{+}(\mathfrak{c}_c)$; all vertices of $p$ have their coordinate $x_0$ in $\left[\boldsymbol{\lambda}\cdot i,\boldsymbol{\lambda}\cdot(i+1)\right]$. If $p$ is orthogonal to the direction $0$, it lies in a slice $\{x_0=t\}$ with $t_i^{-}\le t\le t_i^{+}$, and $t\notin\boldsymbol{\lambda}\mathbb{Z}$, because $p$ is contained in no plane $\{x_0=\boldsymbol{\lambda}\cdot j\}$; then $F_p$ is a component of the magnetic curvature $\tilde{b}_t$. If $p$ contains the direction $0$, it joins two consecutive slices $t$ and $t+1$ of the slab, with $t_i^{-}\le t<t_i^{+}$, and $F_p$ is a component of the electric curvature $\mathrm{El}_t$. So $P_i$ is a bounded Borel function of the curvature on those plaquettes of the slab $i$ that lie in no plane $\{x_0=\boldsymbol{\lambda}\cdot j\}$.

Let $\boldsymbol{\vartheta}=0$ and let $\Pi$ be the identity of $\mathscr{H}_0$; then $\boldsymbol{\omega}_n^{0,\Pi}=\sum_w\boldsymbol{\omega}_n^{(w)}=\boldsymbol{\omega}_n$. By Lemma~\ref{lem:free-coefficient-2}(a),
\[
\boldsymbol{\omega}_n\left(\prod_cG_c\right)=\boldsymbol{\omega}_n\left(\prod_iP_i\right)=k_B^{n_0}\cdot\operatorname{Tr}\left(Y_{N_0-1}\left[P_{N_0-1}\right]\cdots Y_0\left[P_0\right]\right).
\]
The operators $Y_i[P_i]$ are Hilbert--Schmidt operators (Lemma~\ref{lem:free-coefficient-2}). By~\eqref{it:free-coefficient-Q8} they belong to the Schatten class of order $N_0\ge2$, and H\"older's inequality for traces in~\eqref{it:free-coefficient-Q8} yields
\[
\left|\boldsymbol{\omega}_n\left(\prod_cG_c\right)\right|\le k_B^{n_0}\cdot\prod_i\left\|Y_i\left[P_i\right]\right\|_{N_0}=\prod_i\left(k_B^{n_0}\cdot\left\|Y_i\left[P_i\right]\right\|_{N_0}^{N_0}\right)^{\frac{1}{N_0}}.
\]
For every $i$, Lemma~\ref{lem:free-coefficient-2}(c), applied to the slab function $P_i$ at $i$, yields
\[
k_B^{n_0}\cdot\left\|Y_i\left[P_i\right]\right\|_{N_0}^{N_0}=\boldsymbol{\omega}_n\left(\prod_{i'}P_i^{(i')}\right)\ge0,
\]
with the functions $P_i^{(i')}$ and the maps $g_{i'}$ of that lemma for $P=P_i$; by definition, $g_{i'}$ maps the slab $i$ onto the slab $i'$. For $i'\in\mathbb{Z}/N_0\mathbb{Z}$, $c_1\in\mathbb{Z}/N_1\mathbb{Z}$ and $c_2\in\mathbb{Z}/N_2\mathbb{Z}$ define a function on $C^2(n)$ by
\[
\left(\operatorname{Rm}_0^iG\right)_{(i',c_1,c_2)}(F):=\begin{cases}G_{(i,c_1,c_2)}\left(g_{i'}^{*}F\right)&\text{if }i'-i\text{ is even},\\ \overline{G_{(i,c_1,c_2)}\left(g_{i'}^{*}F\right)}&\text{if }i'-i\text{ is odd},\end{cases}
\]
and let $\operatorname{Rm}_0^iG$ be the family of these functions, indexed by the cells. As $P_i$ is the product of the functions $G_{(i,c_1,c_2)}$ over $c_1$ and $c_2$, the definition of $P_i^{(i')}$ gives
\[
P_i^{(i')}=\prod_{c_1,c_2}\left(\operatorname{Rm}_0^iG\right)_{(i',c_1,c_2)},\qquad\prod_{i'}P_i^{(i')}=\prod_c\left(\operatorname{Rm}_0^iG\right)_c.
\]
From the bound obtained by H\"older's inequality, the identity of Lemma~\ref{lem:free-coefficient-2}(c) and the second product identity it follows that
\begin{multline}\label{eq:free-coefficient-14}
\left|\boldsymbol{\omega}_n\left(\prod_cG_c\right)\right|\le\prod_{i\in\mathbb{Z}/N_0\mathbb{Z}}\left(k_B^{n_0}\cdot\left\|Y_i\left[P_i\right]\right\|_{N_0}^{N_0}\right)^{\frac{1}{N_0}}\\
=\prod_{i\in\mathbb{Z}/N_0\mathbb{Z}}\boldsymbol{\omega}_n\left(\prod_c\left(\operatorname{Rm}_0^iG\right)_c\right)^{\frac{1}{N_0}},
\end{multline}
and each of the numbers $\boldsymbol{\omega}_n\left(\prod_c\left(\operatorname{Rm}_0^iG\right)_c\right)$ on the right, being equal to $k_B^{n_0}\cdot\left\|Y_i\left[P_i\right]\right\|_{N_0}^{N_0}$, is non-negative.

Every member of $\operatorname{Rm}_0^iG$ is a mirror copy. Let $c=(i,c_1,c_2)$ and $c'=(i',c_1,c_2)$. As a product of reflections in the planes $\{x_0=\boldsymbol{\lambda}\cdot j\}$, the map $g_{i'}$ moves only the coordinate $x_0$ and permutes the intervals $\left[\boldsymbol{\lambda}\cdot j,\boldsymbol{\lambda}\cdot(j+1)\right]$. The slab $j$ consists of the slices in the interval $\left[\boldsymbol{\lambda}\cdot j,\boldsymbol{\lambda}\cdot(j+1)\right]$, and $g_{i'}$ maps the slab $i$ onto the slab $i'$; so it maps $\left[\boldsymbol{\lambda}\cdot i,\boldsymbol{\lambda}\cdot(i+1)\right]$ onto $\left[\boldsymbol{\lambda}\cdot i',\boldsymbol{\lambda}\cdot(i'+1)\right]$, and $g_{i'}(\mathfrak{c}_c)=\mathfrak{c}_{c'}$. Moreover $g_{i'}$ belongs to the group of the maps $\mathsf{g}_{c''}$, which is generated by the reflections in all the planes $\{x_{\mu}=\boldsymbol{\lambda}\cdot j\}$, and an element of this group is determined by the cell onto which it maps $\mathfrak{c}_0$ (Definition~\ref{def:decomposition-torus-cells}). Both $g_{i'}\circ\mathsf{g}_c$ and $\mathsf{g}_{c'}$ map $\mathfrak{c}_0$ onto $\mathfrak{c}_{c'}$; hence $g_{i'}\circ\mathsf{g}_c=\mathsf{g}_{c'}$, that is, $\mathsf{g}_c^{*}g_{i'}^{*}=\mathsf{g}_{c'}^{*}$. Write $G_c=O^{[c]}$ with $O\in\mathfrak{A}^{+}$ (Definition~\ref{def:decomposition-torus-cells}). Then $G_c\left(g_{i'}^{*}F\right)$ equals $O\left(\mathsf{g}_{c'}^{*}F\right)$ if $c_0+c_1+c_2$ is even, and the complex conjugate of this number if $c_0+c_1+c_2$ is odd; and $c'_0+c'_1+c'_2$ has the parity of $c_0+c_1+c_2$ if and only if $i'-i$ is even. Comparison with the definition of $O^{[c']}$ gives $\left(\operatorname{Rm}_0^iG\right)_{c'}=O^{[c']}$, that is, the member of $\operatorname{Rm}_0^iG$ in $c'$ is the mirror copy in $c'$ of $G_c$.

\emph{The flux zero.} Now let $\boldsymbol{\vartheta}\in\mathbb{R}^2$, let $\Pi$ be the orthogonal projection of $\mathscr{H}_{\boldsymbol{\vartheta}}$ onto $\mathscr{H}_{\boldsymbol{\vartheta}}^{(0)}$, and let the operators $Y_i[\cdot]$ be formed with $\boldsymbol{\vartheta}$ and this $\Pi$. For $\Phi\in\mathfrak{F}(n)$ the series
\[
\boldsymbol{\omega}_n^{\boldsymbol{\vartheta},\Pi}(\Phi)=\sum_{w_{01},w_{02}\in\mathbb{Z}}\exp\left(2\pi i\left(\boldsymbol{\vartheta}_1\cdot w_{01}+\boldsymbol{\vartheta}_2\cdot w_{02}\right)\right)\cdot\boldsymbol{\omega}_n^{(0,w_{02},w_{01})}(\Phi)
\]
converges absolutely and uniformly in $\boldsymbol{\vartheta}$, by~\eqref{it:free-coefficient-Q5}. So it is a continuous function of $\boldsymbol{\vartheta}$, and its integral over $[0,1]^2$ is its coefficient $\boldsymbol{\omega}_n^{(0,0,0)}(\Phi)=\boldsymbol{\omega}_n^{(0)}(\Phi)$. By Lemma~\ref{lem:free-coefficient-2}(a) and~(c), the functions $\boldsymbol{\vartheta}\mapsto k_B^{n_0}\cdot\operatorname{Tr}\left(Y_{N_0-1}[P_{N_0-1}]\cdots Y_0[P_0]\right)$ and $\boldsymbol{\vartheta}\mapsto k_B^{n_0}\cdot\left\|Y_i[P_i]\right\|_{N_0}^{N_0}$ are such absolutely convergent Fourier series, hence continuous, and
\[
\boldsymbol{\omega}_n^{(0)}\left(\prod_iP_i\right)=\int_{[0,1]^2}\mathrm{d}\boldsymbol{\vartheta}\,k_B^{n_0}\cdot\operatorname{Tr}\left(Y_{N_0-1}\left[P_{N_0-1}\right]\cdots Y_0\left[P_0\right]\right).
\]
H\"older's inequality for traces~\eqref{it:free-coefficient-Q8}, then H\"older's inequality for the integral over $[0,1]^2$ with $N_0$ factors and all exponents equal to $N_0$, and finally Lemma~\ref{lem:free-coefficient-2}(c) integrated over $\boldsymbol{\vartheta}$ give
\begin{multline*}
\left|\boldsymbol{\omega}_n^{(0)}\left(\prod_cG_c\right)\right|\le\int_{[0,1]^2}\mathrm{d}\boldsymbol{\vartheta}\,\prod_i\left(k_B^{n_0}\cdot\left\|Y_i\left[P_i\right]\right\|_{N_0}^{N_0}\right)^{\frac{1}{N_0}}\\
\le\prod_i\left(\int_{[0,1]^2}\mathrm{d}\boldsymbol{\vartheta}\,k_B^{n_0}\cdot\left\|Y_i\left[P_i\right]\right\|_{N_0}^{N_0}\right)^{\frac{1}{N_0}}=\prod_i\boldsymbol{\omega}_n^{(0)}\left(\prod_c\left(\operatorname{Rm}_0^iG\right)_c\right)^{\frac{1}{N_0}}.
\end{multline*}
This is the counterpart of~\eqref{eq:free-coefficient-14} for $\boldsymbol{\omega}_n^{(0)}$. Here each number $\boldsymbol{\omega}_n^{(0)}\left(\prod_c\left(\operatorname{Rm}_0^iG\right)_c\right)$ is the integral over $[0,1]^2$ of the non-negative function $\boldsymbol{\vartheta}\mapsto k_B^{n_0}\cdot\left\|Y_i\left[P_i\right]\right\|_{N_0}^{N_0}$, and so it is non-negative. Finally, a mirror copy of a bounded function is bounded; therefore every member $\left(\operatorname{Rm}_0^iG\right)_c$ is a bounded element of $\mathfrak{A}^{+}(\mathfrak{c}_c)$, and $\operatorname{Rm}_0^iG$ satisfies the hypotheses imposed on $G$.

\emph{The three directions.} By~\eqref{it:free-coefficient-Q2}, the inequality~\eqref{eq:free-coefficient-14} and its counterpart for $\boldsymbol{\omega}_n^{(0)}$, together with the non-negativity of the factors on the right, hold with any direction $\mu$ in the place of the direction $0$. The family $\operatorname{Rm}_0^iG$ is then replaced by the family $\operatorname{Rm}_{\mu}^jG$, $j\in\mathbb{Z}/N_{\mu}\mathbb{Z}$, whose member in a cell $c'$ is the mirror copy in $c'$ of the member of $G$ in the cell obtained from $c'$ by replacing its entry $c'_{\mu}$ by $j$. Fix $\boldsymbol{\omega}\in\left\{\boldsymbol{\omega}_n,\boldsymbol{\omega}_n^{(0)}\right\}$ and write $\boldsymbol{\omega}(H):=\boldsymbol{\omega}\left(\prod_cH_c\right)$ for a family $H=\left(H_c\right)_c$. For $\boldsymbol{\omega}=\boldsymbol{\omega}_n$ apply~\eqref{eq:free-coefficient-14}, and for $\boldsymbol{\omega}=\boldsymbol{\omega}_n^{(0)}$ its counterpart for $\boldsymbol{\omega}_n^{(0)}$, in each case first in the direction $0$ to $G$, then in the direction $1$ to each of the families $\operatorname{Rm}_0^{c_0}G$, and finally in the direction $2$ to each of the families $\operatorname{Rm}_1^{c_1}\operatorname{Rm}_0^{c_0}G$. All the factors that arise are non-negative; therefore
\[
\left|\boldsymbol{\omega}(G)\right|\le\prod_{c_0}\boldsymbol{\omega}\left(\operatorname{Rm}_0^{c_0}G\right)^{\frac{1}{N_0}}\le\prod_{c_0,c_1}\boldsymbol{\omega}\left(\operatorname{Rm}_1^{c_1}\operatorname{Rm}_0^{c_0}G\right)^{\frac{1}{N_0\cdot N_1}}\le\prod_{c_0,c_1,c_2}\boldsymbol{\omega}\left(\operatorname{Rm}_2^{c_2}\operatorname{Rm}_1^{c_1}\operatorname{Rm}_0^{c_0}G\right)^{\frac{1}{N_n}},
\]
where $c_{\mu}$ ranges over $\mathbb{Z}/N_{\mu}\mathbb{Z}$. Now let $G_c:=O_c^{[c]}$ for every cell $c$, the observables $O_c\in\mathfrak{A}^{+}$ being bounded; then $\boldsymbol{\omega}(G)=\boldsymbol{\omega}\left(\prod_cO_c^{[c]}\right)$. For all cells $c$ and $c''$, the element $O_c$ is the only $O\in\mathfrak{A}^{+}$ with $O^{[c'']}=O_c^{[c'']}$ (Definition~\ref{def:decomposition-torus-cells}); hence the mirror copy of $O_c^{[c'']}$ in any cell $c'$ is $O_c^{[c']}$. Fix $c=(c_0,c_1,c_2)$ and a cell $c'$. The member of $\operatorname{Rm}_0^{c_0}G$ in the cell $(c'_0,c_1,c_2)$ equals $O_c^{[(c'_0,c_1,c_2)]}$, the mirror copy of $G_c$ there; the member of $\operatorname{Rm}_1^{c_1}\operatorname{Rm}_0^{c_0}G$ in $(c'_0,c'_1,c_2)$ is the mirror copy of that function there, namely $O_c^{[(c'_0,c'_1,c_2)]}$; and the member of $\operatorname{Rm}_2^{c_2}\operatorname{Rm}_1^{c_1}\operatorname{Rm}_0^{c_0}G$ in $c'$ is the mirror copy of the latter in $c'$, namely $O_c^{[c']}$. (In terms of the maps $\mathsf{g}_{c''}$: each of these three functions is obtained from the preceding one, the first from $G_c$, by composition with the pull-back under an element of the group of these maps that moves only $x_0$, only $x_1$, only $x_2$, respectively, followed by complex conjugation if the parity of the cell changes. Since reflections in planes of different directions commute, this group is the direct product of its three subgroups that move a single coordinate; the composite of the three elements maps $\mathfrak{c}_c$ onto $\mathfrak{c}_{c'}$, and it is therefore $\mathsf{g}_{c'}\circ\mathsf{g}_c^{-1}$, the only element of the group that does so.) Consequently the family $\operatorname{Rm}_2^{c_2}\operatorname{Rm}_1^{c_1}\operatorname{Rm}_0^{c_0}G$ has the product $\prod_{c'}O_c^{[c']}=O_c^{\sharp}$, the right end of the chain is $\prod_c\boldsymbol{\omega}\left(O_c^{\sharp}\right)^{\frac{1}{N_n}}$, and the inequality of the proposition follows.

\emph{Positivity.} Let $O\in\mathfrak{A}^{+}$ and $G_c:=O^{[c]}$ for all cells $c$. Then $P_0=\prod_{c:\,c_0=0}O^{[c]}$, the member of $\operatorname{Rm}_0^0G$ in a cell $c'$ is $O^{[c']}$, and $\prod_{i'}P_0^{(i')}=\prod_{c'}O^{[c']}=O^{\sharp}$. Lemma~\ref{lem:free-coefficient-2}(c) for $i=0$ gives, with $\boldsymbol{\vartheta}=0$ and $\Pi$ the identity,
\[
\boldsymbol{\omega}_n\left(O^{\sharp}\right)=k_B^{n_0}\cdot\left\|Y_0\left[P_0\right]\right\|_{N_0}^{N_0}\ge0,
\]
and, with $\Pi$ the projection onto $\mathscr{H}_{\boldsymbol{\vartheta}}^{(0)}$ and after integration over $\boldsymbol{\vartheta}\in[0,1]^2$,
\[
\boldsymbol{\omega}_n^{(0)}\left(O^{\sharp}\right)=\int_{[0,1]^2}\mathrm{d}\boldsymbol{\vartheta}\,k_B^{n_0}\cdot\left\|Y_0\left[P_0\right]\right\|_{N_0}^{N_0}\ge0.
\]
\end{proof}

\begin{proposition}[Comparison of the flux sectors]\label{prop:comparison-flux-sectors}
\begin{enumerate}
\item[(i)] Let the torus $n$ carry a cell structure (Definition~\ref{def:decomposition-torus-cells}), and let $O\in\mathfrak{A}^{+}$. Then, for every $w\in\mathbb{Z}^3$,
\[
\left|\boldsymbol{\omega}_n^{(w)}\left(O^{\sharp}\right)\right|\le\boldsymbol{\omega}_n^{(0)}\left(O^{\sharp}\right).
\]
\item[(ii)] Let $n=(n_0,n_1,n_2)$ with integers $n_{\mu}\ge2$. Then, for every $w\in\mathbb{Z}^3$,
\[
\boldsymbol{\omega}_n^{(w)}(1)\le\mathbf{Z}_n\qquad\text{and}\qquad\boldsymbol{\omega}_n^{(w)}(1)\le\mathfrak{q}_w(n)\cdot Z_n.
\]
Put
\[
\check{\Theta}_n:=\sum_{w\in\mathbb{Z}^3}\mathfrak{q}_w(n)=1+\varsigma_n.
\]
If $\check{\Theta}_n<2$, then
\[
\mathbf{Z}_n\le Z_n\le\frac{\mathbf{Z}_n}{2-\check{\Theta}_n}.
\]
If $\mathbf{M}\ge0$ is an integer such that $\sum_{|w|>\mathbf{M}}\mathfrak{q}_w(n)\le\frac{1}{2}$, where $|w|=\max\left(|w_{12}|,|w_{02}|,|w_{01}|\right)$, then
\[
Z_n\le2\cdot\left(2\mathbf{M}+1\right)^3\cdot\mathbf{Z}_n.
\]
\end{enumerate}
\end{proposition}

\begin{proof}
(i) By the truncation argument at the beginning of Section~\ref{sec:chessboard-estimates-flux}, $O$ may be assumed bounded. Fix $w_{12}\in\mathbb{Z}$, and for $\boldsymbol{\vartheta}\in\mathbb{R}^2$ let $\Pi$ be the orthogonal projection of $\mathscr{H}_{\boldsymbol{\vartheta}}$ onto $\mathscr{H}_{\boldsymbol{\vartheta}}^{(w_{12})}$. Let $P:=\prod_{c:\,c_0=0}O^{[c]}$; it is a slab function at $0$, and $O^{\sharp}=\prod_{i'}P^{(i')}$ with the functions $P^{(i')}$ of Lemma~\ref{lem:free-coefficient-2}(c) (proof of Proposition~\ref{prop:chessboard-estimate}). In the variables of the slabs these slab functions of $O^{\sharp}$ are alternately $P$, at the even slabs, and the function $P^{\boldsymbol{\theta}}$ of Lemma~\ref{lem:dependence-boundary-slices}, at the odd slabs (Lemma~\ref{lem:free-coefficient-2}(b)). Let $\breve{g}\colon\mathbb{R}^2\to\mathbb{C}$ be defined by
\[
\breve{g}(\boldsymbol{\vartheta}):=\sum_{w_{01},w_{02}\in\mathbb{Z}}\exp\left(2\pi i\left(\boldsymbol{\vartheta}_1\cdot w_{01}+\boldsymbol{\vartheta}_2\cdot w_{02}\right)\right)\cdot\boldsymbol{\omega}_n^{(w_{12},w_{02},w_{01})}\left(O^{\sharp}\right).
\]
This is $\boldsymbol{\omega}_n^{\boldsymbol{\vartheta},\Pi}\left(O^{\sharp}\right)$, and Lemma~\ref{lem:free-coefficient-2}(c) for $i=0$ gives
\[
\breve{g}(\boldsymbol{\vartheta})=k_B^{n_0}\cdot\left\|Y_0[P]\right\|_{N_0}^{N_0}\ge0\qquad\left(\boldsymbol{\vartheta}\in\mathbb{R}^2\right),
\]
the operator $Y_0[P]$ being formed with $\boldsymbol{\vartheta}$ and $\Pi$. By~\eqref{it:free-coefficient-Q5} the series converges absolutely and uniformly in $\boldsymbol{\vartheta}$; so $\breve{g}$ is continuous and $\mathbb{Z}^2$-periodic, and integration term by term yields, for all $w_{02},w_{01}\in\mathbb{Z}$,
\[
\boldsymbol{\omega}_n^{(w_{12},w_{02},w_{01})}\left(O^{\sharp}\right)=\int_{[0,1]^2}\mathrm{d}\boldsymbol{\vartheta}\,\exp\left(-2\pi i\left(\boldsymbol{\vartheta}_1\cdot w_{01}+\boldsymbol{\vartheta}_2\cdot w_{02}\right)\right)\cdot\breve{g}(\boldsymbol{\vartheta}).
\]
For $w_{01}=w_{02}=0$ the integrand is $\breve{g}$ itself. Since the exponential factor has modulus one and $\breve{g}\ge0$, the integrand always has modulus $\breve{g}$; therefore
\[
\left|\boldsymbol{\omega}_n^{(w_{12},w_{02},w_{01})}\left(O^{\sharp}\right)\right|\le\int_{[0,1]^2}\mathrm{d}\boldsymbol{\vartheta}\,\breve{g}(\boldsymbol{\vartheta})=\boldsymbol{\omega}_n^{(w_{12},0,0)}\left(O^{\sharp}\right)\qquad\text{for all }w_{12},w_{02},w_{01}\in\mathbb{Z},
\]
and the right side, as the integral of a non-negative function, is non-negative.

The right side is estimated by exchanging the directions $0$ and $1$. Let $g$ be the exchange of the axes $0$ and $1$, and $n':=gn$. By~\eqref{it:free-coefficient-Q2}, $n'$ carries a cell structure, $O\circ g^{*}$ belongs to $\mathfrak{A}^{+}$ on $n'$, $\left(O\circ g^{*}\right)^{\sharp}=O^{\sharp}\circ g^{*}$, and
\[
\boldsymbol{\omega}_{n'}^{(g\cdot w)}\left(O^{\sharp}\circ g^{*}\right)=\boldsymbol{\omega}_n^{(w)}\left(O^{\sharp}\right),\qquad g\cdot w=\left(w_{02},w_{12},-w_{01}\right)\qquad\left(w\in\mathbb{Z}^3\right);
\]
in particular $g\cdot(w_{12},0,0)=(0,w_{12},0)$ and $g\cdot0=0$. In the inequality just obtained, written on $n'$, the component of index $12$ is held fixed, while the components of indices $02$ and $01$ are those that carry the phases in the series defining $\breve{g}$; for the flux $g\cdot w$ the former is $(g\cdot w)_{12}=w_{02}$, and the latter are $(g\cdot w)_{02}=w_{12}$ and $(g\cdot w)_{01}=-w_{01}$. That inequality, for the torus $n'$, the function $O\circ g^{*}$ and the flux $(0,w_{12},0)$, gives
\[
\left|\boldsymbol{\omega}_n^{(w_{12},0,0)}\left(O^{\sharp}\right)\right|=\left|\boldsymbol{\omega}_{n'}^{(0,w_{12},0)}\left(\left(O\circ g^{*}\right)^{\sharp}\right)\right|\le\boldsymbol{\omega}_{n'}^{(0,0,0)}\left(\left(O\circ g^{*}\right)^{\sharp}\right)=\boldsymbol{\omega}_n^{(0,0,0)}\left(O^{\sharp}\right).
\]
Together with the preceding inequality this gives, for every $w=(w_{12},w_{02},w_{01})\in\mathbb{Z}^3$,
\[
\left|\boldsymbol{\omega}_n^{(w)}\left(O^{\sharp}\right)\right|\le\boldsymbol{\omega}_n^{(w_{12},0,0)}\left(O^{\sharp}\right)\le\boldsymbol{\omega}_n^{(0)}\left(O^{\sharp}\right).
\]

(ii) Fix $w_{12}\in\mathbb{Z}$, and let $\Pi$ be as in (i). For the function $1$ in place of $O^{\sharp}$ no cell structure is needed. Indeed, by~\eqref{it:free-coefficient-Q5}(iii), applied to the decomposition of $n_0$ into $n_0\ge2$ slabs of thickness one with the slab functions $1$, and by $C_{1,\boldsymbol{\vartheta}}[1]=T_{\boldsymbol{\vartheta}}$ from~\eqref{it:free-coefficient-Q5}(iv), the trace over the summand $\mathscr{H}_{\boldsymbol{\vartheta}}^{(w_{12})}$ gives
\[
\sum_{w_{01},w_{02}\in\mathbb{Z}}\exp\left(2\pi i\left(\boldsymbol{\vartheta}_1\cdot w_{01}+\boldsymbol{\vartheta}_2\cdot w_{02}\right)\right)\cdot\boldsymbol{\omega}_n^{(w_{12},w_{02},w_{01})}(1)=k_B^{n_0}\cdot\operatorname{Tr}_{\mathscr{H}_{\boldsymbol{\vartheta}}^{(w_{12})}}\left(T_{\boldsymbol{\vartheta}}^{n_0}\right).
\]
By~\eqref{it:free-coefficient-Q5}(iv) the operator $T_{\boldsymbol{\vartheta}}=X_{\boldsymbol{\vartheta}}^{*}X_{\boldsymbol{\vartheta}}$ is non-negative and of trace class, and by~\eqref{it:free-coefficient-Q5}(iii) it maps $\mathscr{H}_{\boldsymbol{\vartheta}}^{(w_{12})}$ into itself; so its restriction to this subspace is non-negative, and the trace of the $n_0$-th power of this restriction is non-negative. The left side, which is the function $\breve{g}$ of (i) with $1$ in place of $O^{\sharp}$, is therefore non-negative for every torus $n$ with $n_{\mu}\ge2$, and it is continuous by~\eqref{it:free-coefficient-Q5}. The argument of (i) applies to it, the function $1$ being invariant under $g^{*}$, and gives $\boldsymbol{\omega}_n^{(w)}(1)\le\boldsymbol{\omega}_n^{(0)}(1)=\mathbf{Z}_n$ for every such $n$ and every $w$. The second inequality is~\eqref{it:free-coefficient-Q4} for the two-cochain $0$, since $\exp\langle0,F\rangle=1$ and $\|0\|=0$.

By Section~\ref{sec:wilson-determinant-lattice}, $Z_n=\sum_w\boldsymbol{\omega}_n^{(w)}(1)$, $\mathbf{Z}_n=\boldsymbol{\omega}_n^{(0)}(1)$ and $\mathfrak{q}_0(n)=1$, so that $\check{\Theta}_n=1+\varsigma_n$. The second inequality gives
\[
Z_n=\mathbf{Z}_n+\sum_{w\ne0}\boldsymbol{\omega}_n^{(w)}(1)\le\mathbf{Z}_n+\left(\check{\Theta}_n-1\right)\cdot Z_n.
\]
If $\check{\Theta}_n<2$, it follows that $\left(2-\check{\Theta}_n\right)\cdot Z_n\le\mathbf{Z}_n$, which is the upper bound for $Z_n$; the lower bound $\mathbf{Z}_n\le Z_n$ is the case $\Phi=1$ of
\[
\left|\boldsymbol{\omega}_n^{(0)}(\Phi)\right|\le\sum_w\left|\boldsymbol{\omega}_n^{(w)}(\Phi)\right|\le\boldsymbol{\omega}_n\left(|\Phi|\right),
\]
see~\eqref{it:free-coefficient-Q5}. Finally let $\mathbf{M}$ be as in the statement. Since $\#\left|\left\{w\in\mathbb{Z}^3:\ |w|\le\mathbf{M}\right\}\right|=\left(2\mathbf{M}+1\right)^3$, the first inequality for the fluxes with $|w|\le\mathbf{M}$ and the second one for the fluxes with $|w|>\mathbf{M}$ give
\[
Z_n=\sum_{|w|\le\mathbf{M}}\boldsymbol{\omega}_n^{(w)}(1)+\sum_{|w|>\mathbf{M}}\boldsymbol{\omega}_n^{(w)}(1)\le\left(2\mathbf{M}+1\right)^3\cdot\mathbf{Z}_n+\frac{1}{2}\cdot Z_n,
\]
and $Z_n\le2\cdot\left(2\mathbf{M}+1\right)^3\cdot\mathbf{Z}_n$ follows, as $Z_n<\infty$.
\end{proof}

\begin{proposition}[Lengthening one side]\label{prop:lengthening-one-side}
Let the torus $n$ carry the cell structure of width $\boldsymbol{\lambda}$, let $O\in\mathfrak{A}^{+}$, let $k\ge1$ be an integer and let $\mu\in\{0,1,2\}$ be a direction. Let $n^{(\mu,k)}$ be the torus $n$ with $n_{\mu}$ replaced by $k\cdot n_{\mu}$; it carries the cell structure of width $\boldsymbol{\lambda}$ as well, with the even number $k\cdot N_{\mu}$ of cells in the direction $\mu$. On each of the tori $n$ and $n^{(\mu,k)}$, $O^{\sharp}$ denotes the product of the mirror copies of $O$ over the cells of that torus. Let $\nu<\rho$ be the two directions different from $\mu$, and let
\[
W_{\mu}:=\left\{w\in\mathbb{Z}^3:\ w_{\nu\rho}=0\right\}
\]
be the set of the fluxes whose component in the pair of directions that does not contain $\mu$ vanishes; thus $W_0=\{w:\ w_{12}=0\}$, $W_1=\{w:\ w_{02}=0\}$ and $W_2=\{w:\ w_{01}=0\}$.
\begin{enumerate}
\item[(a)] One has
\[
0\le\boldsymbol{\omega}_{n^{(\mu,k)}}(O^{\sharp})\le\boldsymbol{\omega}_n(O^{\sharp})^k.
\]
\item[(b)] One has
\[
0\le\boldsymbol{\omega}^{(0)}_{n^{(\mu,k)}}(O^{\sharp})\le\left(\sum_{w\in W_{\mu}}\left|\boldsymbol{\omega}_n^{(w)}(O^{\sharp})\right|\right)^k.
\]
\end{enumerate}
\end{proposition}

\begin{proof}
By the convention on unbounded observables made at the beginning of Section~\ref{sec:chessboard-estimates-flux}, it suffices to prove the inequalities for bounded $O$, which is assumed from now on. By \eqref{it:free-coefficient-Q2} it suffices, moreover, to treat the direction $\mu=0$. Indeed, for $\mu\in\{1,2\}$ let $g$ be the exchange of the axes $0$ and $\mu$; then $g$ maps $n^{(\mu,k)}$ onto $(gn)^{(0,k)}$, the assertions for the direction $\mu$, the torus $n$ and the function $O$ are the assertions for the direction $0$, the torus $gn$ and the function $O\circ g^{*}$, and the relabelling $w\mapsto g\cdot w$ of \eqref{it:free-coefficient-Q2} maps $W_{\mu}$ onto $W_0$. Let therefore $\mu=0$.

Let $\boldsymbol{\vartheta}\in\mathbb{R}^2$, and let $\Pi$ be either the identity of $\mathscr{H}_{\boldsymbol{\vartheta}}$ or the orthogonal projection onto one of the summands $\mathscr{H}_{\boldsymbol{\vartheta}}^{(w_{12})}$ (Section~\ref{sec:slabs-operators}). Let
\[
P:=\prod_{c:\ c_0=0}O^{[c]}
\]
be the factor of $O^{\sharp}$ that belongs to the slab $0$. As shown in the proof of Proposition~\ref{prop:chessboard-estimate}, $P$ is a slab function at $0$, and the family $\left(P^{(i')}\right)_{i'\in\mathbb{Z}/N_0\mathbb{Z}}$ that Lemma~\ref{lem:free-coefficient-2}(c) attaches to $P$ for $i=0$ has the product $O^{\sharp}$. Put
\[
Y:=Y_0[P],\qquad\mathsf{T}:=Y^{*}Y;
\]
both operators depend on $\boldsymbol{\vartheta}$ and on $\Pi$. By Lemma~\ref{lem:free-coefficient-2}, $Y$ is a Hilbert--Schmidt operator; hence $\mathsf{T}\ge0$ is of trace class, and so is $\mathsf{T}^{\frac{N_0}{2}}\ge0$. Since $N_0$ is even, Lemma~\ref{lem:free-coefficient-2}(c) and \eqref{it:free-coefficient-Q8} give
\[
\boldsymbol{\omega}_n^{\boldsymbol{\vartheta},\Pi}(O^{\sharp})=k_B^{n_0}\cdot\|Y\|_{N_0}^{N_0}=k_B^{n_0}\cdot\operatorname{Tr}\mathsf{T}^{\frac{N_0}{2}}.
\]
The torus $n^{(0,k)}$ has the same cross-section $\mathbb{Z}/n_1\mathbb{Z}\times\mathbb{Z}/n_2\mathbb{Z}$ as $n$, and it is cut into $k\cdot N_0$ slabs. The number $k_B$, the spaces $\mathscr{H}_{\boldsymbol{\vartheta}}$ and the operators $X_{\boldsymbol{\vartheta}}$ and $C_{m_{\mathrm{w}},\boldsymbol{\vartheta}}[\cdot]$ of \eqref{it:free-coefficient-Q5} are formed with the cross-section alone; the cells of the slab $0$ and the mirror copies of $O$ in them are the same for both tori, so that the factor of $O^{\sharp}$ that belongs to the slab $0$ of $n^{(0,k)}$ is again $P$. Hence $Y$ and $\mathsf{T}$ are the same operators for both tori, and the identity just obtained, written for $n^{(0,k)}$, reads
\[
\boldsymbol{\omega}_{n^{(0,k)}}^{\boldsymbol{\vartheta},\Pi}(O^{\sharp})=k_B^{kn_0}\cdot\operatorname{Tr}\mathsf{T}^{\frac{kN_0}{2}}=k_B^{kn_0}\cdot\operatorname{Tr}\left(\mathsf{T}^{\frac{N_0}{2}}\right)^k.
\]
For a non-negative operator $B$ of trace class one has $\operatorname{Tr}B^k\le\left(\operatorname{Tr}B\right)^k$, by \eqref{it:free-coefficient-Q8}. Applied to $B=\mathsf{T}^{\frac{N_0}{2}}$, this gives, for every $\boldsymbol{\vartheta}\in\mathbb{R}^2$ and for both choices of $\Pi$,
\[
0\le\boldsymbol{\omega}_{n^{(0,k)}}^{\boldsymbol{\vartheta},\Pi}(O^{\sharp})\le\left(k_B^{n_0}\cdot\operatorname{Tr}\mathsf{T}^{\frac{N_0}{2}}\right)^k=\boldsymbol{\omega}_n^{\boldsymbol{\vartheta},\Pi}(O^{\sharp})^k.
\]

(a) For $\boldsymbol{\vartheta}=0$ and $\Pi$ the identity, all the weights $\exp\left(2\pi i\left(\boldsymbol{\vartheta}_1w_{01}+\boldsymbol{\vartheta}_2w_{02}\right)\right)$ equal $1$ and the sums are taken over all $w\in\mathbb{Z}^3$, so that the functionals $\boldsymbol{\omega}^{\boldsymbol{\vartheta},\Pi}$ on $n$ and on $n^{(0,k)}$ are $\boldsymbol{\omega}_n$ and $\boldsymbol{\omega}_{n^{(0,k)}}$. The inequality displayed above becomes (a).

(b) Let $\Pi$ be the orthogonal projection onto $\mathscr{H}_{\boldsymbol{\vartheta}}^{(0)}$, the summand with $w_{12}=0$, and define $\breve{g}\colon\mathbb{R}^2\to\mathbb{R}$ by
\[
\breve{g}(\boldsymbol{\vartheta}):=\boldsymbol{\omega}_n^{\boldsymbol{\vartheta},\Pi}(O^{\sharp})=\sum_{w_{01},w_{02}\in\mathbb{Z}}\exp\left(2\pi i\left(\boldsymbol{\vartheta}_1w_{01}+\boldsymbol{\vartheta}_2w_{02}\right)\right)\cdot\boldsymbol{\omega}_n^{(0,w_{02},w_{01})}(O^{\sharp}).
\]
By the identity $\boldsymbol{\omega}_n^{\boldsymbol{\vartheta},\Pi}(O^{\sharp})=k_B^{n_0}\cdot\operatorname{Tr}\mathsf{T}^{\frac{N_0}{2}}$, one has $\breve{g}\ge0$. The series that defines $\boldsymbol{\omega}_{n^{(0,k)}}^{\boldsymbol{\vartheta},\Pi}(O^{\sharp})$ converges absolutely, by \eqref{it:free-coefficient-Q5}. Integrating it term by term over $\boldsymbol{\vartheta}\in[0,1]^2$ leaves only its term with $w_{01}=w_{02}=0$, which is $\boldsymbol{\omega}^{(0)}_{n^{(0,k)}}(O^{\sharp})$, because the integral of $\exp\left(2\pi i\left(\boldsymbol{\vartheta}_1w_{01}+\boldsymbol{\vartheta}_2w_{02}\right)\right)$ over $[0,1]^2$ vanishes unless $w_{01}=w_{02}=0$. With the inequality displayed before (a) it follows that
\[
\boldsymbol{\omega}^{(0)}_{n^{(0,k)}}(O^{\sharp})=\int_{[0,1]^2}\mathrm{d}\boldsymbol{\vartheta}\,k_B^{kn_0}\cdot\operatorname{Tr}\mathsf{T}^{\frac{kN_0}{2}}\le\int_{[0,1]^2}\mathrm{d}\boldsymbol{\vartheta}\,\breve{g}(\boldsymbol{\vartheta})^k\le\left(\sup_{\mathbb{R}^2}\breve{g}\right)^k,
\]
and the left side is non-negative, since the integrand in the middle is. Finally, $\left|\exp\left(2\pi i\left(\boldsymbol{\vartheta}_1w_{01}+\boldsymbol{\vartheta}_2w_{02}\right)\right)\right|=1$; hence for every $\boldsymbol{\vartheta}\in\mathbb{R}^2$
\[
\breve{g}(\boldsymbol{\vartheta})\le\sum_{w_{01},w_{02}\in\mathbb{Z}}\left|\boldsymbol{\omega}_n^{(0,w_{02},w_{01})}(O^{\sharp})\right|=\sum_{w\in W_0}\left|\boldsymbol{\omega}_n^{(w)}(O^{\sharp})\right|,
\]
which proves (b).
\end{proof}

\subsubsection{Transfer to larger cubic tori}\label{sec:transfer-larger-cubic}

\paragraph{Parameters and reference numbers.}
Let $q>0$ and $\boldsymbol{\mu}>-\frac{1}{2}$, let $\boldsymbol{\lambda}$ be a cell width (Definition~\ref{def:decomposition-torus-cells}), let $\mathsf{n}\ge2$ be an even integer, and let $j\ge2$ be an integer. Put
\[
n^{*}:=\boldsymbol{\lambda}\cdot\mathsf{n},\qquad N^{*}:=\mathsf{n}^3,\qquad\boldsymbol{\alpha}:=\frac{2\pi^2}{q^2\cdot n^{*}};
\]
here $n^{*}$ is an integer, as $\mathsf{n}$ is even. From here to the end of Section~\ref{sec:trace-formula-chessboard}, $N^{*}$ stands for $\mathsf{n}^3$; for $\mathsf{n}=2L$ it is the number $N^{*}=8L^3$ of Section~\ref{sec:constants-exponents-scales}. For an integer $k\ge1$ the cubic torus $k\cdot n^{*}$ carries the cell structure of width $\boldsymbol{\lambda}$ with $k\cdot\mathsf{n}$ cells per side, and the following symbols refer to it:
\[
E_k:=E^{\mathrm{af}}_{kn^{*}},\quad\bar{E}_k:=E^{\mathrm{av}}_{kn^{*}},\quad Z_k:=Z_{kn^{*}},\quad\mathbf{Z}_k:=\mathbf{Z}_{kn^{*}},\quad\boldsymbol{\omega}_k:=\boldsymbol{\omega}_{kn^{*}},\quad\boldsymbol{\omega}_k^{(w)}:=\boldsymbol{\omega}_{kn^{*}}^{(w)}.
\]
Moreover $E_{*}:=E_1$ and $\bar{E}_{*}:=\bar{E}_1$. The reference numbers are
\[
D_j:=j^3\log Z_1-\log Z_j,\qquad\boldsymbol{\delta}:=\frac{D_j}{(j-1)^3\cdot N^{*}},
\]
and, for $O\in\mathfrak{A}^{+}$,
\[
\|O\|_{*}:=\left(E_{*}O^{\sharp}\right)^{\frac{1}{N^{*}}},\qquad\|O\|_{0*}:=\left(\bar{E}_{*}O^{\sharp}\right)^{\frac{1}{N^{*}}}.
\]
These two numbers are well defined and non-negative: $E_{*}O^{\sharp}=\boldsymbol{\omega}_1(O^{\sharp})/Z_1$ and $\bar{E}_{*}O^{\sharp}=\boldsymbol{\omega}_1^{(0)}(O^{\sharp})/\mathbf{Z}_1$, where $Z_1>0$ and $\mathbf{Z}_1>0$ (Section~\ref{sec:wilson-determinant-lattice}), while $\boldsymbol{\omega}_1(O^{\sharp})\ge0$ and $\boldsymbol{\omega}_1^{(0)}(O^{\sharp})\ge0$ by Proposition~\ref{prop:chessboard-estimate}.

For positive integers $k_0,k_1,k_2$ let $n^{*}(k_0,k_1,k_2)$ be the torus $\left(k_0\cdot n^{*},k_1\cdot n^{*},k_2\cdot n^{*}\right)$. It carries the cell structure of width $\boldsymbol{\lambda}$, and for $n=n^{*}(k_0,k_1,k_2)$ one has $N_n=k_0k_1k_2\cdot N^{*}$ and
\[
\mathfrak{q}_w(n)=\exp\left(-\boldsymbol{\alpha}\sum_{\mu<\nu}\frac{k_{\iota}}{k_{\mu}k_{\nu}}\cdot w_{\mu\nu}^2\right)\qquad\left(w\in\mathbb{Z}^3\right),
\]
where in each summand $\iota$ is the direction different from $\mu$ and $\nu$. Indeed, by the definitions of $r_q=\frac{2\pi}{q}$ and of $\mathfrak{q}_w(n)$ in Section~\ref{sec:wilson-determinant-lattice} it suffices to note that $|n|=k_0k_1k_2\cdot\left(n^{*}\right)^3$ and $n_{\iota}=k_{\iota}\cdot n^{*}$, so that
\[
\frac{r_q^2\cdot n_{\iota}^2}{2|n|}=\frac{r_q^2}{2n^{*}}\cdot\frac{k_{\iota}}{k_{\mu}k_{\nu}},\qquad\frac{r_q^2}{2n^{*}}=\frac{2\pi^2}{q^2\cdot n^{*}}=\boldsymbol{\alpha}.
\]

\begin{theorem}[Transfer for the all-flux law]\label{thm:transfer-all-flux}
Let the parameters be those of Section~\ref{sec:transfer-larger-cubic}, let $k\in\{1\}\cup\{j,j+1,\dots\}$, let $S$ be a set of cells of the cubic torus $k\cdot n^{*}$, which has $k\cdot\mathsf{n}$ cells per side, and let $O_c\in\mathfrak{A}^{+}$ for $c\in S$. Then $D_j\ge0$ and
\[
\left|E_k\left(\prod_{c\in S}O_c^{[c]}\right)\right|\le\prod_{c\in S}\exp(\boldsymbol{\delta})\cdot\|O_c\|_{*}.
\]
\end{theorem}

\begin{proof}
Since $n^{*}\ge2$, the function $\phi(a,b,c):=\log Z_{n^{*}(a,b,c)}$ on $\{1,2,\dots\}^3$ has the properties stated in \eqref{it:free-coefficient-Q6}. One has $\phi(k,k,k)=\log Z_k$ for every integer $k\ge1$; in particular $D_j=j^3\phi(1,1,1)-\phi(j,j,j)$. By \eqref{it:free-coefficient-Q7}, $D_j\ge0$ and
\[
\log Z_k\ge k^3\log Z_1-\left(\frac{k-1}{j-1}\right)^3\cdot D_j\qquad\text{for every integer }k\ge j.
\]
As $D_j\ge0$ and $0\le k-1\le k$, it follows that
\begin{equation}\label{eq:free-coefficient-15}
\log Z_k\ge k^3\cdot\left(\log Z_1-\frac{D_j}{(j-1)^3}\right)\qquad\text{for }k=1\text{ and for every integer }k\ge j;
\end{equation}
for $k=1$ the inequality is $D_j\ge0$.

Put $O_c:=1$ for the cells $c$ of $k\cdot n^{*}$ that do not lie in $S$. Then $O_c\in\mathfrak{A}^{+}$ for every cell $c$; moreover $1^{[c]}=1$ and $1^{\sharp}=1$, so that $\boldsymbol{\omega}_k(1^{\sharp})=Z_k$. The torus $k\cdot n^{*}$ has $k^3\cdot N^{*}$ cells, and Proposition~\ref{prop:chessboard-estimate} for $\boldsymbol{\omega}_k$ gives
\[
\left|\boldsymbol{\omega}_k\left(\prod_{c\in S}O_c^{[c]}\right)\right|\le\prod_{c\in S}\boldsymbol{\omega}_k(O_c^{\sharp})^{\frac{1}{k^3\cdot N^{*}}}\cdot\prod_{c\notin S}Z_k^{\frac{1}{k^3\cdot N^{*}}}.
\]
Since $Z_k$ is the product of the numbers $Z_k^{\frac{1}{k^3\cdot N^{*}}}$ over all cells, division by $Z_k$ yields
\[
\left|E_k\left(\prod_{c\in S}O_c^{[c]}\right)\right|\le\prod_{c\in S}\left(\frac{\boldsymbol{\omega}_k(O_c^{\sharp})}{Z_k}\right)^{\frac{1}{k^3\cdot N^{*}}}.
\]
Let $O\in\mathfrak{A}^{+}$. The tori $n^{*}(k,1,1)=\left(n^{*}\right)^{(0,k)}$, $n^{*}(k,k,1)=\left(n^{*}(k,1,1)\right)^{(1,k)}$ and $k\cdot n^{*}=\left(n^{*}(k,k,1)\right)^{(2,k)}$ carry the cell structure of width $\boldsymbol{\lambda}$. Proposition~\ref{prop:lengthening-one-side}(a), applied to $n^{*}$ in the direction $0$, to $n^{*}(k,1,1)$ in the direction $1$ and to $n^{*}(k,k,1)$ in the direction $2$, gives
\[
0\le\boldsymbol{\omega}_k(O^{\sharp})\le\boldsymbol{\omega}_{n^{*}(k,k,1)}(O^{\sharp})^k\le\boldsymbol{\omega}_{n^{*}(k,1,1)}(O^{\sharp})^{k^2}\le\boldsymbol{\omega}_1(O^{\sharp})^{k^3}.
\]
Together with \eqref{eq:free-coefficient-15} this yields
\begin{align*}
\left(\frac{\boldsymbol{\omega}_k(O^{\sharp})}{Z_k}\right)^{\frac{1}{k^3\cdot N^{*}}}&\le\boldsymbol{\omega}_1(O^{\sharp})^{\frac{1}{N^{*}}}\cdot\exp\left(-\frac{\log Z_k}{k^3\cdot N^{*}}\right)\\
&\le\boldsymbol{\omega}_1(O^{\sharp})^{\frac{1}{N^{*}}}\cdot\exp\left(-\frac{\log Z_1}{N^{*}}+\boldsymbol{\delta}\right)=\exp(\boldsymbol{\delta})\cdot\|O\|_{*},
\end{align*}
because $E_{*}O^{\sharp}=\boldsymbol{\omega}_1(O^{\sharp})/Z_1$. This bound, taken with $O=O_c$ for each $c\in S$ and inserted into the estimate of $\left|E_k\left(\prod_{c\in S}O_c^{[c]}\right)\right|$ obtained above, gives the inequality asserted in the theorem; the assertion $D_j\ge0$ was proved at the beginning.
\end{proof}

\begin{lemma}[A lower bound for $\mathbf{Z}_k$]\label{lem:lower-bound}
Let the data be those of Section~\ref{sec:transfer-larger-cubic}, and let $\boldsymbol{\alpha}\ge4$. Then
\[
\log\mathbf{Z}_k\ge k^3\cdot\left(\log Z_1-\frac{D_j}{(j-1)^3}-4\right)
\]
for $k=1$ and for every integer $k\ge j$.
\end{lemma}

\begin{proof}
Let $k$ be $1$ or an integer $k\ge j$, and let $n:=n^{*}(k,k,k)$ be the cubic torus with $kn^{*}$ sites per side; thus $Z_k=Z_n$ and $\mathbf{Z}_k=\mathbf{Z}_n$. For $k_0=k_1=k_2=k$ the formula for $\mathfrak{q}_w(n)$ in Section~\ref{sec:transfer-larger-cubic} reads
\[
\mathfrak{q}_w(n)=\exp\left(-\frac{\boldsymbol{\alpha}}{k}\cdot|w|_2^2\right)\qquad\left(w\in\mathbb{Z}^3\right),
\]
so that each of the three components of $w$ enters with the coefficient $\frac{\boldsymbol{\alpha}}{k}$. Let $\mathbf{M}\ge0$ be an integer; $|w|$ denotes the largest modulus of a component of $w$. Every $w$ with $|w|>\mathbf{M}$ has a component of modulus larger than $\mathbf{M}$. Fix one of the three components. The sum of $\mathfrak{q}_w(n)$ over the $w$ for which this component has modulus larger than $\mathbf{M}$ is the product of the factor $\mathrm{thf}\left(\frac{\boldsymbol{\alpha}}{k}\right)^2$, which comes from the two other components, with $2\sum_{i>\mathbf{M}}\exp\left(-\frac{\boldsymbol{\alpha}}{k}\cdot i^2\right)$, and by the second elementary bound of Section~\ref{sec:recalled-results} the latter is at most $2\cdot\exp\left(-\frac{\boldsymbol{\alpha}}{k}(\mathbf{M}+1)^2\right)\cdot\mathrm{thf}\left(\frac{\boldsymbol{\alpha}}{k}\right)$. Adding the contributions of the three components,
\[
\sum_{w\in\mathbb{Z}^3,\ |w|>\mathbf{M}}\mathfrak{q}_w(n)\le6\cdot\mathrm{thf}\left(\frac{\boldsymbol{\alpha}}{k}\right)^3\cdot\exp\left(-\frac{\boldsymbol{\alpha}}{k}(\mathbf{M}+1)^2\right).
\]
The first elementary bound of Section~\ref{sec:recalled-results} and $\boldsymbol{\alpha}\ge4$ give
\[
\mathrm{thf}\left(\frac{\boldsymbol{\alpha}}{k}\right)\le1+\left(\frac{\pi k}{\boldsymbol{\alpha}}\right)^{\frac{1}{2}}\le1+\left(\frac{\pi k}{4}\right)^{\frac{1}{2}}\le2k^{\frac{1}{2}},
\]
the last step because $\left(\frac{\pi}{4}\right)^{\frac{1}{2}}<1\le k^{\frac{1}{2}}$. Hence
\[
\sum_{w\in\mathbb{Z}^3,\ |w|>\mathbf{M}}\mathfrak{q}_w(n)\le48k^{\frac{3}{2}}\cdot\exp\left(-\frac{\boldsymbol{\alpha}}{k}(\mathbf{M}+1)^2\right).
\]
Now let
\[
\mathbf{M}:=\left\lfloor\left(\frac{k\cdot\log\left(96k^{\frac{3}{2}}\right)}{\boldsymbol{\alpha}}\right)^{\frac{1}{2}}\right\rfloor.
\]
Then $\frac{\boldsymbol{\alpha}}{k}(\mathbf{M}+1)^2>\log\left(96k^{\frac{3}{2}}\right)$, and the right side of the last inequality is smaller than $48k^{\frac{3}{2}}\cdot\left(96k^{\frac{3}{2}}\right)^{-1}=\frac{1}{2}$. Moreover $\mathbf{M}\le k$. Indeed, $\log96<4.6$, $\log k\le k-1$ and $\boldsymbol{\alpha}\ge4$ give
\[
\log\left(96k^{\frac{3}{2}}\right)\le4.6+1.5(k-1)<4k+8+\frac{4}{k}=\frac{4(k+1)^2}{k}\le\frac{\boldsymbol{\alpha}\cdot(k+1)^2}{k},
\]
so that $\frac{k}{\boldsymbol{\alpha}}\cdot\log\left(96k^{\frac{3}{2}}\right)<(k+1)^2$, and the integer $\mathbf{M}$ is smaller than $k+1$. Proposition~\ref{prop:comparison-flux-sectors}(ii), applied to the torus $n$ and to this $\mathbf{M}$, yields
\[
Z_k\le2(2\mathbf{M}+1)^3\cdot\mathbf{Z}_k\le2(2k+1)^3\cdot\mathbf{Z}_k,\qquad\text{that is,}\qquad\log\mathbf{Z}_k\ge\log Z_k-\log\left(2(2k+1)^3\right).
\]
By \eqref{eq:free-coefficient-15}, $\log Z_k\ge k^3\cdot\left(\log Z_1-\frac{D_j}{(j-1)^3}\right)$. Finally, $(2k+1)^3\le27k^3$, $\log k\le k^3-1$ and $\log54<3.99$ give
\[
\log\left(2(2k+1)^3\right)\le\log54+3\left(k^3-1\right)<4k^3.
\]
The last three inequalities together give the assertion.
\end{proof}

\begin{remark}[Translated cells; aligned tori]\label{rem:translated-cells-aligned}
By \eqref{it:free-coefficient-Q2}, the functional $E_k$ is invariant under the translations of the lattice. Therefore Theorem~\ref{thm:transfer-all-flux} remains true, with the numbers of Section~\ref{sec:transfer-larger-cubic} unchanged, for the cell structure whose planes are the sets $\left\{x:x_{\mu}\in s_{\mu}+\boldsymbol{\lambda}\mathbb{Z}\right\}$, $\mu=0,1,2$, where $s\in\mathbb{Z}^3$ is fixed; these are the cells of Definition~\ref{def:decomposition-torus-cells} translated by $s$. If $m_{\mathrm{p}}\ge2$, the aligned tori of Section~\ref{sec:lattices-tests-function}, with the cells of Section~\ref{sec:reference-tori-cells}, are the special case $\boldsymbol{\lambda}=m_{\mathrm{p}}+\frac{1}{2}$, $\mathsf{n}=2L$, $j=L$ of the data of Section~\ref{sec:transfer-larger-cubic}. In this case $n^{*}=n_{*}$ and $N^{*}=8L^3$, and the tori $R=hL^T$ with $T\ge T^{*}$ are the cubic tori with $kn^{*}$ sites per side for $k=L^{T-T^{*}}$; these values of $k$ lie in $\{1\}\cup\{L,L+1,\dots\}$.
\end{remark}

\subsubsection{The harmonic-averaged law under a moment bound}\label{sec:harmonic-averaged-law}

\begin{proposition}[Transfer for the harmonic-averaged law]\label{prop:transfer-harmonic-averaged}
Let the data be those of Section~\ref{sec:transfer-larger-cubic}, and let $\boldsymbol{\alpha}\ge4$. Let $\Phi\in\mathfrak{A}^{+}$, and let $m_1>0$ and $p\ge0$ be real numbers such that
\[
\left(E^{\mathrm{af}}_n|\Phi|^{2\ell}\right)^{\frac{1}{2\ell}}\le m_1\ell^p
\]
for every integer $\ell\ge1$ and every torus $n=n^{*}(j_0,j_1,j_2)$ with positive integers $j_0,j_1,j_2$. Let $\vartheta$ be a positive number with $\vartheta\ge\|\Phi\|_{0*}$. Put
\[
\boldsymbol{\lambda}_0:=\log^{+}\left(\frac{m_1\cdot(N^{*})^p}{\vartheta}\right),\qquad\tilde{b}:=\frac{2pN^{*}}{\boldsymbol{\alpha}},
\]
and let $X_0\ge0$ be defined by $X_0^2:=\frac{2\left(N^{*}\cdot\boldsymbol{\lambda}_0+3\right)}{\boldsymbol{\alpha}}$. Then, for $k=1$ and for every integer $k\ge j$,
\[
\left|\bar{E}_k\Phi\right|\le\vartheta\cdot\exp\left(\frac{1}{N^{*}}\left(3\log\left(4+16X_0^2\right)+\frac{3+8\tilde{b}}{\exp(1)}+\frac{D_j}{(j-1)^3}+4\right)\right).
\]
Let, moreover, $c_1,\dots,c_J$ be distinct cells of the cubic torus with $kn^{*}$ sites per side, and let $\Phi_1,\dots,\Phi_J$ satisfy the hypotheses made on $\Phi$, each with numbers $m_1$, $p$, $\vartheta$ of its own. Then $\left|\bar{E}_k\prod_{i=1}^J\Phi_i^{[c_i]}\right|$ is at most the product of the $J$ right sides of the last inequality, formed with these numbers.
\end{proposition}

\begin{proof}
For positive integers $k_0,k_1,k_2$ let $n^{*}(k_0,k_1,k_2)$ be the torus with the sides $k_0n^{*}$, $k_1n^{*}$, $k_2n^{*}$, as in Section~\ref{sec:transfer-larger-cubic}. It carries the cell structure of the data, with $k_{\mu}\mathsf{n}$ cells in the direction $\mu$, and has $N_n=k_0k_1k_2N^{*}$ cells. For such a torus $n$ put
\[
y_n:=\vartheta^{N_n}\cdot Z_1^{\frac{N_n}{N^{*}}},\qquad A_n:=\frac{\boldsymbol{\omega}_n^{(0)}\left(\Phi^{\sharp}\right)}{y_n}.
\]
By Proposition~\ref{prop:chessboard-estimate}, $A_n\ge0$. As $N_{n^{(\mu,k)}}=k\cdot N_n$, one has $y_{n^{(\mu,k)}}=y_n^k$ for every direction $\mu$ and every integer $k\ge1$.

(a) Let $n=n^{*}(k_0,k_1,k_2)$ and $w\in\mathbb{Z}^3$. The functionals $\boldsymbol{\omega}_n^{(w')}$, $w'\in\mathbb{Z}^3$, are positive, the determinant in their definition being non-negative (Theorem~\ref{thm:free-coefficient-1}(d2)); since $\boldsymbol{\omega}_n$ is their sum, $\boldsymbol{\omega}_n^{(w)}(G)\le\boldsymbol{\omega}_n(G)$ for every non-negative function $G$. The Cauchy--Schwarz inequality for $\boldsymbol{\omega}_n^{(w)}$, Proposition~\ref{prop:comparison-flux-sectors}(ii), and the last inequality with $G=\left|\Phi^{\sharp}\right|^2$ give
\[
\left|\boldsymbol{\omega}_n^{(w)}\left(\Phi^{\sharp}\right)\right|\le\boldsymbol{\omega}_n^{(w)}(1)^{\frac{1}{2}}\cdot\boldsymbol{\omega}_n^{(w)}\left(\left|\Phi^{\sharp}\right|^2\right)^{\frac{1}{2}}\le\mathfrak{q}_w(n)^{\frac{1}{2}}\cdot Z_n\cdot\left(E^{\mathrm{af}}_n\left|\Phi^{\sharp}\right|^2\right)^{\frac{1}{2}}.
\]
The function $\left|\Phi^{\sharp}\right|^2$ is the product of the $N_n$ functions $\left|\Phi^{[c]}\right|^2$, and $\left|\Phi^{[c]}\right|=|\Phi|\circ\mathsf{g}_c^{*}$ with the lattice isometry $\mathsf{g}_c$; hence $E^{\mathrm{af}}_n\left|\Phi^{[c]}\right|^{2N_n}=E^{\mathrm{af}}_n|\Phi|^{2N_n}$ for every cell $c$, by \eqref{it:free-coefficient-Q2}. H\"older's inequality with $N_n$ exponents equal to $N_n$, followed by the hypothesis with $\ell=N_n$, yields
\[
\left(E^{\mathrm{af}}_n\left|\Phi^{\sharp}\right|^2\right)^{\frac{1}{2}}\le\left(\left(E^{\mathrm{af}}_n|\Phi|^{2N_n}\right)^{\frac{1}{2N_n}}\right)^{N_n}\le\left(m_1N_n^p\right)^{N_n}.
\]
By \eqref{it:free-coefficient-Q6}, applied in each of the three variables, $\log Z_n\le k_0k_1k_2\cdot\log Z_1$, that is, $Z_n\le Z_1^{\frac{N_n}{N^{*}}}$. Altogether, for every $w\in\mathbb{Z}^3$,
\[
\frac{\left|\boldsymbol{\omega}_n^{(w)}\left(\Phi^{\sharp}\right)\right|}{y_n}\le\mathfrak{q}_w(n)^{\frac{1}{2}}\cdot\left(\frac{m_1N_n^p}{\vartheta}\right)^{N_n}.
\]

(b) Let $n=n^{*}(k_0,k_1,k_2)$, let $\mu\in\{0,1,2\}$ be a direction, let $k\ge1$ and $\mathbf{M}\ge0$ be integers, and let $W_{\mu}$ be the set of Proposition~\ref{prop:lengthening-one-side}(b). An element $w\in W_{\mu}$ is determined by its two components that contain the direction $\mu$; call them $z$ and $z'$, the free components of $w$. Then $|w|=\max\left(|z|,|z'|\right)$, and by the formula for $\mathfrak{q}_w(n)$ in Section~\ref{sec:transfer-larger-cubic} there are numbers $x_1,x_2>0$, the weights of the two free components, such that
\[
\mathfrak{q}_w(n)=\exp\left(-x_1z^2-x_2z'^2\right)\qquad\left(w\in W_{\mu}\right);
\]
the weight of a component $w_{\nu\nu'}$ is its coefficient $\boldsymbol{\alpha}\cdot\frac{k_{\iota}}{k_{\nu}k_{\nu'}}$ in that formula, $\iota$ being the third direction. Put
\[
\boldsymbol{\tau}'_n(\mathbf{M}):=\sum_{w\in W_{\mu},\ |w|>\mathbf{M}}\mathfrak{q}_w(n)^{\frac{1}{2}}.
\]
Proposition~\ref{prop:lengthening-one-side}(b) gives $\boldsymbol{\omega}^{(0)}_{n^{(\mu,k)}}\left(\Phi^{\sharp}\right)\le\left(\sum_{w\in W_{\mu}}\left|\boldsymbol{\omega}_n^{(w)}\left(\Phi^{\sharp}\right)\right|\right)^k$. In this sum each of the $(2\mathbf{M}+1)^2$ terms with $|w|\le\mathbf{M}$ is at most $\boldsymbol{\omega}_n^{(0)}\left(\Phi^{\sharp}\right)=y_n\cdot A_n$, by Proposition~\ref{prop:comparison-flux-sectors}(i), and the other terms are bounded by (a); thus
\[
\frac{1}{y_n}\sum_{w\in W_{\mu}}\left|\boldsymbol{\omega}_n^{(w)}\left(\Phi^{\sharp}\right)\right|\le(2\mathbf{M}+1)^2\cdot A_n+\boldsymbol{\tau}'_n(\mathbf{M})\cdot\left(\frac{m_1N_n^p}{\vartheta}\right)^{N_n}.
\]
Raising to the power $k$ and dividing by $y_{n^{(\mu,k)}}=y_n^k$, one obtains
\begin{equation}\label{eq:free-coefficient-16}
A_{n^{(\mu,k)}}\le\left((2\mathbf{M}+1)^2\cdot A_n+1\right)^k\qquad\text{whenever}\qquad\boldsymbol{\tau}'_n(\mathbf{M})\cdot\left(\frac{m_1N_n^p}{\vartheta}\right)^{N_n}\le1.
\end{equation}
The set of the $w\in W_{\mu}$ with $\max\left(|z|,|z'|\right)>\mathbf{M}$ is contained in the union of the sets $\left\{|z|>\mathbf{M}\right\}$ and $\left\{|z'|>\mathbf{M}\right\}$. Summing the other free component over $\mathbb{Z}$, and applying the second elementary bound of Section~\ref{sec:recalled-results} to the positive and to the negative values of the component that exceeds $\mathbf{M}$ in modulus, one finds
\begin{align*}
\boldsymbol{\tau}'_n(\mathbf{M})&\le2\cdot\mathrm{thf}\left(\frac{x_1}{2}\right)\cdot\mathrm{thf}\left(\frac{x_2}{2}\right)\cdot\left(\exp\left(-\frac{x_1}{2}(\mathbf{M}+1)^2\right)+\exp\left(-\frac{x_2}{2}(\mathbf{M}+1)^2\right)\right)\\
&\le4\cdot\mathrm{thf}\left(\frac{x_1}{2}\right)\cdot\mathrm{thf}\left(\frac{x_2}{2}\right)\cdot\exp\left(-\frac{\min\left(x_1,x_2\right)}{2}\cdot(\mathbf{M}+1)^2\right).
\end{align*}
The function $\mathrm{thf}$ is decreasing. As $\boldsymbol{\alpha}\ge4$ and $k\ge1$,
\[
\mathrm{thf}\left(\frac{\boldsymbol{\alpha}}{2}\right),\ \mathrm{thf}\left(\frac{\boldsymbol{\alpha}\cdot k}{2}\right)\le\mathrm{thf}(2)\le1+2\exp(-2)+4\exp(-8)<1.28,
\]
and the first elementary bound of Section~\ref{sec:recalled-results} gives
\[
\mathrm{thf}\left(\frac{\boldsymbol{\alpha}}{2k}\right)\le1+\left(\frac{2\pi k}{\boldsymbol{\alpha}}\right)^{\frac{1}{2}}\le1+\left(\frac{\pi k}{2}\right)^{\frac{1}{2}}\le2.26k^{\frac{1}{2}},
\]
the last step because $\left(\frac{\pi}{2}\right)^{\frac{1}{2}}<1.26$ and $1\le k^{\frac{1}{2}}$.

(c) Let now $k$ be $1$ or an integer $k\ge j$. For $i\in\{0,1,2\}$ put
\[
\boldsymbol{\lambda}^{(i)}:=\log^{+}\left(\frac{m_1\cdot\left(k^iN^{*}\right)^p}{\vartheta}\right),
\]
and let $X_i\ge0$ be defined by $X_i^2:=\frac{2\left(N^{*}\cdot\boldsymbol{\lambda}^{(i)}+3\right)}{\boldsymbol{\alpha}}$; for $i=0$ these are $\boldsymbol{\lambda}_0$ and $X_0$. Since $\log^{+}(ab)\le\log^{+}a+\log b$ for $b\ge1$,
\[
\boldsymbol{\lambda}^{(i)}\le\boldsymbol{\lambda}_0+ip\cdot\log k,\qquad X_i^2\le X_0^2+i\cdot\tilde{b}\cdot\log k.
\]
The torus $n^{*}(k,k,k)$ arises from $n^{*}$ by three multiplications of a side by $k$: in the direction $0$ from $n^{*}$ to $n^{*}(k,1,1)$, in the direction $1$ from $n^{*}(k,1,1)$ to $n^{*}(k,k,1)$, and in the direction $2$ from $n^{*}(k,k,1)$ to $n^{*}(k,k,k)$. In the multiplication in the direction $i$ the torus $n$ that is multiplied has $N_n=k^iN^{*}$ cells; hence $\left(\frac{m_1N_n^p}{\vartheta}\right)^{N_n}\le\exp\left(N_n\cdot\boldsymbol{\lambda}^{(i)}\right)$, and, by the bound for $\boldsymbol{\tau}'_n$ in (b), the condition in \eqref{eq:free-coefficient-16} is satisfied by every integer $\mathbf{M}\ge0$ with
\[
\frac{\min\left(x_1,x_2\right)}{2}\cdot(\mathbf{M}+1)^2\ge N_n\cdot\boldsymbol{\lambda}^{(i)}+\log\left(4\cdot\mathrm{thf}\left(\frac{x_1}{2}\right)\cdot\mathrm{thf}\left(\frac{x_2}{2}\right)\right).
\]
\begin{itemize}
\item Direction $0$: the free components are $w_{01}$ and $w_{02}$, both of weight $\boldsymbol{\alpha}$, and $N_n=N^{*}$. With $\mathbf{M}_0:=\left\lfloor X_0\right\rfloor$,
\[
\frac{\boldsymbol{\alpha}}{2}(\mathbf{M}_0+1)^2>\frac{\boldsymbol{\alpha}}{2}\cdot X_0^2=N^{*}\cdot\boldsymbol{\lambda}^{(0)}+3>N^{*}\cdot\boldsymbol{\lambda}^{(0)}+\log\left(4\cdot1.28^2\right).
\]
\item Direction $1$: the free components are $w_{01}$, of weight $\frac{\boldsymbol{\alpha}}{k}$, and $w_{12}$, of weight $\boldsymbol{\alpha}\cdot k$, and $N_n=kN^{*}$. With $\mathbf{M}_1:=\left\lfloor kX_1\right\rfloor$,
\[
\frac{\boldsymbol{\alpha}}{2k}(\mathbf{M}_1+1)^2>\frac{\boldsymbol{\alpha}}{2k}\cdot k^2X_1^2=kN^{*}\cdot\boldsymbol{\lambda}^{(1)}+3k\ge kN^{*}\cdot\boldsymbol{\lambda}^{(1)}+\log\left(4\cdot2.26k^{\frac{1}{2}}\cdot1.28\right),
\]
the last step because $\log(4\cdot2.26\cdot1.28)<3$ and $\frac{1}{2}\log k\le3(k-1)$.
\item Direction $2$: the free components are $w_{02}$ and $w_{12}$, both of weight $\boldsymbol{\alpha}$, and $N_n=k^2N^{*}$. With $\mathbf{M}_2:=\left\lfloor kX_2\right\rfloor$,
\[
\frac{\boldsymbol{\alpha}}{2}(\mathbf{M}_2+1)^2>\frac{\boldsymbol{\alpha}}{2}\cdot k^2X_2^2=k^2N^{*}\cdot\boldsymbol{\lambda}^{(2)}+3k^2\ge k^2N^{*}\cdot\boldsymbol{\lambda}^{(2)}+3>k^2N^{*}\cdot\boldsymbol{\lambda}^{(2)}+\log\left(4\cdot1.28^2\right).
\]
\end{itemize}
By the bounds for $\mathrm{thf}$ in (b), the condition in \eqref{eq:free-coefficient-16} holds in each of the three multiplications, with $\mathbf{M}=\mathbf{M}_i$ in the direction $i$.

(d) By the definition of $\|\Phi\|_{0*}$ in Section~\ref{sec:transfer-larger-cubic}, $\|\Phi\|_{0*}^{N^{*}}=\frac{\boldsymbol{\omega}_1^{(0)}\left(\Phi^{\sharp}\right)}{\mathbf{Z}_1}$. Together with $\vartheta\ge\|\Phi\|_{0*}$ and $\mathbf{Z}_1\le Z_1$, which holds because the functionals $\boldsymbol{\omega}_{n^{*}}^{(w)}$ are positive, this gives
\[
A_{n^{*}}=\left(\frac{\|\Phi\|_{0*}}{\vartheta}\right)^{N^{*}}\cdot\frac{\mathbf{Z}_1}{Z_1}\le1.
\]
Put $\check{C}_i:=2(2\mathbf{M}_i+1)^2$, so that $\check{C}_i\ge2$. For $A\ge0$ one has $(2\mathbf{M}_i+1)^2\cdot A+1\le\check{C}_i\cdot\max(A,1)$; therefore the three applications of \eqref{eq:free-coefficient-16} in (c) give
\[
A_{n^{*}(k,1,1)}\le\check{C}_0^k,\qquad A_{n^{*}(k,k,1)}\le\check{C}_1^k\cdot\check{C}_0^{k^2},\qquad A_{n^{*}(k,k,k)}\le\check{C}_2^k\cdot\check{C}_1^{k^2}\cdot\check{C}_0^{k^3},
\]
so that $k^{-3}\log A_{n^{*}(k,k,k)}\le\log\check{C}_0+k^{-1}\log\check{C}_1+k^{-2}\log\check{C}_2$. Since $\mathbf{M}_0\le X_0$ and $(2x+1)^2\le8x^2+2$ for real $x$,
\[
\check{C}_0\le2(2X_0+1)^2\le4+16X_0^2.
\]
For $i=1,2$ one has $\mathbf{M}_i\le kX_i$ and $k\ge1$, hence $\check{C}_i\le4+16k^2X_i^2\le k^2\cdot\left(4+16X_i^2\right)$; with $X_i^2\le X_0^2+i\cdot\tilde{b}\cdot\log k$ and $\log(x+y)\le\log x+\frac{y}{x}$ for $x\ge4$, $y\ge0$, it follows that
\[
\log\check{C}_i\le2\log k+\log\left(4+16X_0^2\right)+4i\cdot\tilde{b}\cdot\log k\qquad(i=1,2).
\]
Since $\log\left(4+16X_0^2\right)>0$, $\sup_{k\ge1}k^{-1}\log k\le\exp(-1)$ and $\sup_{k\ge1}k^{-2}\log k\le\frac{1}{2\exp(1)}$,
\begin{align*}
k^{-3}\log A_{n^{*}(k,k,k)}&\le\left(1+k^{-1}+k^{-2}\right)\cdot\log\left(4+16X_0^2\right)+(2+4\tilde{b})\cdot k^{-1}\log k+(2+8\tilde{b})\cdot k^{-2}\log k\\
&\le3\log\left(4+16X_0^2\right)+\frac{2+4\tilde{b}}{\exp(1)}+\frac{2+8\tilde{b}}{2\exp(1)}=3\log\left(4+16X_0^2\right)+\frac{3+8\tilde{b}}{\exp(1)}.
\end{align*}

(e) Lemma~\ref{lem:lower-bound} gives $\log\mathbf{Z}_k\ge k^3\cdot\left(\log Z_1-\frac{D_j}{(j-1)^3}-4\right)$.

(f) The torus $n^{*}(k,k,k)$ has $k^3N^{*}$ cells, and $\boldsymbol{\omega}_k^{(0)}=\boldsymbol{\omega}_{n^{*}(k,k,k)}^{(0)}$. Proposition~\ref{prop:chessboard-estimate}, applied to $\boldsymbol{\omega}_k^{(0)}$ and to the family $O_0:=\Phi$, $O_c:=1$ for the cells $c\ne0$ (note $1^{\sharp}=1$ and $\boldsymbol{\omega}_k^{(0)}(1)=\mathbf{Z}_k$), gives
\[
\left|\bar{E}_k\Phi\right|\le\left(\frac{\boldsymbol{\omega}_k^{(0)}\left(\Phi^{\sharp}\right)}{\mathbf{Z}_k}\right)^{\frac{1}{k^3N^{*}}},\qquad\text{where}\qquad\boldsymbol{\omega}_k^{(0)}\left(\Phi^{\sharp}\right)=A_{n^{*}(k,k,k)}\cdot\vartheta^{k^3N^{*}}\cdot Z_1^{k^3}.
\]
Inserting (d) and (e),
\begin{align*}
\frac{1}{k^3N^{*}}\log\left(\frac{\boldsymbol{\omega}_k^{(0)}\left(\Phi^{\sharp}\right)}{\mathbf{Z}_k}\right)&=\log\vartheta+\frac{1}{N^{*}}\left(k^{-3}\log A_{n^{*}(k,k,k)}+\log Z_1-k^{-3}\log\mathbf{Z}_k\right)\\
&\le\log\vartheta+\frac{1}{N^{*}}\left(3\log\left(4+16X_0^2\right)+\frac{3+8\tilde{b}}{\exp(1)}+\frac{D_j}{(j-1)^3}+4\right),
\end{align*}
which is the first assertion. For the cells $c_1,\dots,c_J$, Proposition~\ref{prop:chessboard-estimate}, applied to the family $O_{c_i}:=\Phi_i$ ($i=1,\dots,J$), $O_c:=1$ for the other cells, gives
\[
\left|\bar{E}_k\prod_{i=1}^J\Phi_i^{[c_i]}\right|\le\prod_{i=1}^J\left(\frac{\boldsymbol{\omega}_k^{(0)}\left(\Phi_i^{\sharp}\right)}{\mathbf{Z}_k}\right)^{\frac{1}{k^3N^{*}}},
\]
and each factor is bounded as above, with the numbers $m_1$, $p$, $\vartheta$ of $\Phi_i$.
\end{proof}

\subsubsection{Domination of the zero-flux sector}\label{sec:domination-zero-flux}

\begin{proposition}[Zero-flux sector, convexity and transfer]\label{prop:zero-flux-sector}
Let the data be those of Section~\ref{sec:transfer-larger-cubic}, without any condition on $\boldsymbol{\alpha}$, and put
\[
\bar{D}_j:=j^3\log\mathbf{Z}_1-\log\mathbf{Z}_j,\qquad\bar{\mathfrak{D}}_j:=j^3\log Z_1-\log\mathbf{Z}_j.
\]
\begin{enumerate}
\item[(a)] Let $n$ be a torus that carries the cell structure of width $\boldsymbol{\lambda}$ (Definition~\ref{def:decomposition-torus-cells}). For every $O\in\mathfrak{A}^{+}$,
\[
\boldsymbol{\omega}_n^{(0)}\left(O^{\sharp}\right)\le\boldsymbol{\omega}_n\left(O^{\sharp}\right).
\]
\item[(b)] Let $n_1,n_2\ge2$ and $u\ge2$ be integers. The function $a\mapsto\log\mathbf{Z}_{(ua,n_1,n_2)}$ on $\{1,2,\dots\}$ is convex, that is,
\[
2\log\mathbf{Z}_{(ua,n_1,n_2)}\le\log\mathbf{Z}_{(u(a-1),n_1,n_2)}+\log\mathbf{Z}_{(u(a+1),n_1,n_2)}\qquad\text{for every integer }a\ge2.
\]
In the same sense, for all integers $n_0,n_1,n_2\ge2$ and $u\ge2$, the functions $a\mapsto\log\mathbf{Z}_{(n_0,ua,n_2)}$ and $a\mapsto\log\mathbf{Z}_{(n_0,n_1,ua)}$ on $\{1,2,\dots\}$ are convex.
\item[(c)] One has
\[
0\le D_j\le\bar{\mathfrak{D}}_j=\bar{D}_j-j^3\log\frac{\mathbf{Z}_1}{Z_1},
\]
and, for every integer $k\ge j$,
\[
\log\mathbf{Z}_k\ge k^3\log Z_1-\left(\frac{k-1}{j-1}\right)^3\cdot\bar{\mathfrak{D}}_j.
\]
\item[(d)] Let $k$ be $1$ or an integer $k\ge j$, let $S$ be a set of cells of the cubic torus $k\cdot n^{*}$, and let $O_c\in\mathfrak{A}^{+}$ for $c\in S$. If $k\ge j$, then
\[
\left|\bar{E}_k\prod_{c\in S}O_c^{[c]}\right|\le\prod_{c\in S}\exp\left(\frac{\bar{\mathfrak{D}}_j}{(j-1)^3\cdot N^{*}}\right)\cdot\|O_c\|_{*}.
\]
If $k=1$, then
\[
\left|\bar{E}_1\prod_{c\in S}O_c^{[c]}\right|\le\prod_{c\in S}\|O_c\|_{0*}.
\]
\item[(e)] Let $\varsigma_{n^{*}}$ be the number of Section~\ref{sec:wilson-determinant-lattice} for the cubic torus $n^{*}$. For every $O\in\mathfrak{A}^{+}$,
\[
\|O\|_{*}\le\|O\|_{0*}+\varsigma_{n^{*}}^{\frac{1}{2N^{*}}}\cdot\left(E_1|O|^{2N^{*}}\right)^{\frac{1}{2N^{*}}}.
\]
\end{enumerate}
\end{proposition}

\begin{proof}
(a) The reduction to bounded functions at the beginning of Section~\ref{sec:chessboard-estimates-flux} applies to this inequality; so $O$ may be assumed bounded. Define $\mathsf{a}\colon\mathbb{Z}^3\to\mathbb{C}$ by
\[
\mathsf{a}(w_{12},w_{02},w_{01}):=\boldsymbol{\omega}_n^{(w)}\left(O^{\sharp}\right)\qquad\left(w=(w_{12},w_{02},w_{01})\in\mathbb{Z}^3\right).
\]
By~\eqref{it:free-coefficient-Q5}, $\sum_w\left|\mathsf{a}(w)\right|\le\boldsymbol{\omega}_n\left(\left|O^{\sharp}\right|\right)<\infty$, and by Section~\ref{sec:wilson-determinant-lattice}, $\sum_w\mathsf{a}(w)=\boldsymbol{\omega}_n\left(O^{\sharp}\right)$ and $\mathsf{a}(0,0,0)=\boldsymbol{\omega}_n^{(0)}\left(O^{\sharp}\right)$. It is shown first that, for every $\boldsymbol{\vartheta}=(\boldsymbol{\vartheta}_1,\boldsymbol{\vartheta}_2)\in\mathbb{R}^2$,
\begin{align*}
\sum_{w_{01},w_{02}\in\mathbb{Z}}\exp\left(2\pi i\left(\boldsymbol{\vartheta}_1w_{01}+\boldsymbol{\vartheta}_2w_{02}\right)\right)\cdot\mathsf{a}(w_{12},w_{02},w_{01})&\ge0\qquad\text{for every }w_{12}\in\mathbb{Z},\\
\sum_{w_{01},w_{12}\in\mathbb{Z}}\exp\left(2\pi i\left(\boldsymbol{\vartheta}_1w_{01}+\boldsymbol{\vartheta}_2w_{12}\right)\right)\cdot\mathsf{a}(w_{12},w_{02},w_{01})&\ge0\qquad\text{for every }w_{02}\in\mathbb{Z},\\
\sum_{w_{02},w_{12}\in\mathbb{Z}}\exp\left(2\pi i\left(\boldsymbol{\vartheta}_1w_{02}+\boldsymbol{\vartheta}_2w_{12}\right)\right)\cdot\mathsf{a}(w_{12},w_{02},w_{01})&\ge0\qquad\text{for every }w_{01}\in\mathbb{Z};
\end{align*}
these are called the first, the second and the third inequality. Each of the series converges absolutely, uniformly in $\boldsymbol{\vartheta}$, since $\sum_w\left|\mathsf{a}(w)\right|<\infty$.

\emph{The first inequality.} Let $w_{12}\in\mathbb{Z}$ and $\boldsymbol{\vartheta}\in\mathbb{R}^2$, and let $\Pi$ be the orthogonal projection of $\mathscr{H}_{\boldsymbol{\vartheta}}$ onto $\mathscr{H}_{\boldsymbol{\vartheta}}^{(w_{12})}$. As in the proof of Proposition~\ref{prop:comparison-flux-sectors}(i), the product $P:=\prod_{c:\,c_0=0}O^{[c]}$ is a slab function at $0$, and $O^{\sharp}=\prod_{i'}P^{(i')}$ with the functions $P^{(i')}$ that Lemma~\ref{lem:free-coefficient-2}(c) attaches to $P$ for $i=0$ (proof of Proposition~\ref{prop:chessboard-estimate}); in the variables of the slabs these factors of $O^{\sharp}$ are $P$ at the even slabs and $P^{\boldsymbol{\theta}}$ at the odd slabs (Lemma~\ref{lem:free-coefficient-2}(b)). Lemma~\ref{lem:free-coefficient-2}(c) gives
\[
\boldsymbol{\omega}_n^{\boldsymbol{\vartheta},\Pi}\left(O^{\sharp}\right)=k_B^{n_0}\cdot\left\|Y_0[P]\right\|_{N_0}^{N_0}\ge0,
\]
and by the definition of $\boldsymbol{\omega}_n^{\boldsymbol{\vartheta},\Pi}$ in Section~\ref{sec:slabs-operators} the left side is the series of the first inequality. Here the direction $0$ is time: the component $w_{12}$ held fixed is the one whose two indices differ from $0$, and the angles are paired with the two components that contain the direction $0$.

\emph{The second and the third inequality.} Let $\mu\in\{1,2\}$, let $g$ be the exchange of the axes $0$ and $\mu$, and put $n':=gn$. By~\eqref{it:free-coefficient-Q2}, the torus $n'$ carries the cell structure of width $\boldsymbol{\lambda}$, the function $O\circ g^{*}$ belongs to $\mathfrak{A}^{+}$ on $n'$, $\left(O\circ g^{*}\right)^{\sharp}=O^{\sharp}\circ g^{*}$, and $\boldsymbol{\omega}_{n'}^{(g\cdot w)}\left(O^{\sharp}\circ g^{*}\right)=\mathsf{a}(w)$ for every $w\in\mathbb{Z}^3$. Since $w\mapsto g\cdot w$ is a bijection of $\mathbb{Z}^3$ (Theorem~\ref{thm:free-coefficient-1}(d5)), the first inequality, applied to the torus $n'$ and to the function $O\circ g^{*}$, states that
\[
\sum_{w\in\mathbb{Z}^3:\ (g\cdot w)_{12}=v}\exp\left(2\pi i\left(\boldsymbol{\vartheta}_1\cdot(g\cdot w)_{01}+\boldsymbol{\vartheta}_2\cdot(g\cdot w)_{02}\right)\right)\cdot\mathsf{a}(w)\ge0\qquad\left(v\in\mathbb{Z},\ \boldsymbol{\vartheta}\in\mathbb{R}^2\right).
\]
For $\mu=1$ one has $g\cdot w=(w_{02},w_{12},-w_{01})$; the condition $(g\cdot w)_{12}=v$ fixes $w_{02}$, the component whose two indices differ from $1$, and the angles are paired with $-w_{01}$ and $w_{12}$. Inserting $(-\boldsymbol{\vartheta}_1,\boldsymbol{\vartheta}_2)$ in place of $\boldsymbol{\vartheta}$ gives the second inequality. For $\mu=2$ one has $g\cdot w=(-w_{01},-w_{02},-w_{12})$; the condition fixes $w_{01}$, the component whose two indices differ from $2$, and the angles are paired with $-w_{12}$ and $-w_{02}$. Inserting $(-\boldsymbol{\vartheta}_2,-\boldsymbol{\vartheta}_1)$ in place of $\boldsymbol{\vartheta}$ gives the third inequality.

\emph{Consequences.} At $\boldsymbol{\vartheta}=0$ the first inequality states that the numbers $\sum_{w_{01},w_{02}}\mathsf{a}(w_{12},w_{02},w_{01})$, $w_{12}\in\mathbb{Z}$, are real and non-negative; their sum over $w_{12}\in\mathbb{Z}$ is $\sum_w\mathsf{a}(w)$. Hence
\[
A_1:=\sum_{w_{01},w_{02}\in\mathbb{Z}}\mathsf{a}(0,w_{02},w_{01})\le\sum_{w\in\mathbb{Z}^3}\mathsf{a}(w)=\boldsymbol{\omega}_n\left(O^{\sharp}\right).
\]
Put $\boldsymbol{\vartheta}_1=0$ in the second inequality and integrate over $\boldsymbol{\vartheta}_2\in[0,1]$. The uniform absolute convergence permits integration term by term, and $\int_0^1\mathrm{d}\boldsymbol{\vartheta}_2\,\exp\left(2\pi i\,\boldsymbol{\vartheta}_2w_{12}\right)=0$ for $w_{12}\ne0$; so only the terms with $w_{12}=0$ remain, and $\sum_{w_{01}}\mathsf{a}(0,w_{02},w_{01})\ge0$ for every $w_{02}\in\mathbb{Z}$. Hence
\[
A_2:=\sum_{w_{01}\in\mathbb{Z}}\mathsf{a}(0,0,w_{01})\le\sum_{w_{02}\in\mathbb{Z}}\ \sum_{w_{01}\in\mathbb{Z}}\mathsf{a}(0,w_{02},w_{01})=A_1.
\]
The integral of the third inequality over $\boldsymbol{\vartheta}\in[0,1]^2$, taken in the same way, leaves the term with $w_{02}=w_{12}=0$; so $\mathsf{a}(0,0,w_{01})\ge0$ for every $w_{01}\in\mathbb{Z}$, and $\mathsf{a}(0,0,0)\le A_2$. Altogether
\[
\boldsymbol{\omega}_n^{(0)}\left(O^{\sharp}\right)=\mathsf{a}(0,0,0)\le A_2\le A_1\le\boldsymbol{\omega}_n\left(O^{\sharp}\right).
\]

(b) Let $n_0\ge2$ be an integer and $n:=(n_0,n_1,n_2)$. The number $k_B$, the Hilbert spaces $\mathscr{H}_{\boldsymbol{\vartheta}}$ with their summands $\mathscr{H}_{\boldsymbol{\vartheta}}^{(w_{12})}$, and the operators $T_{\boldsymbol{\vartheta}}$ of~\eqref{it:free-coefficient-Q5} are formed with the cross-section $\mathbb{Z}/n_1\times\mathbb{Z}/n_2$ alone; they do not depend on $n_0$. Let $\boldsymbol{\vartheta}\in\mathbb{R}^2$. By~\eqref{it:free-coefficient-Q5}(iv), $T_{\boldsymbol{\vartheta}}=C_{1,\boldsymbol{\vartheta}}[1]$ is a non-negative operator of trace class, and by~\eqref{it:free-coefficient-Q5}(iii) it maps $\mathscr{H}_{\boldsymbol{\vartheta}}^{(0)}$ into itself; so its restriction to $\mathscr{H}_{\boldsymbol{\vartheta}}^{(0)}$ is a non-negative compact self-adjoint operator of trace class. For every integer $m\ge1$ let
\[
\mathrm{tf}(m,\boldsymbol{\vartheta}):=\operatorname{Tr}_{\mathscr{H}_{\boldsymbol{\vartheta}}^{(0)}}T_{\boldsymbol{\vartheta}}^m\ge0
\]
be the trace of the $m$-th power of this restriction. By~\eqref{it:free-coefficient-Q5}(iii), applied to the decomposition of $n_0$ into $n_0\ge2$ slabs of thickness one, each carrying the slab function $1$, and by $C_{1,\boldsymbol{\vartheta}}[1]^{n_0}=T_{\boldsymbol{\vartheta}}^{n_0}$, the trace over the summand with $w_{12}=0$ gives
\[
k_B^{n_0}\cdot\mathrm{tf}(n_0,\boldsymbol{\vartheta})=\sum_{w_{01},w_{02}\in\mathbb{Z}}\exp\left(2\pi i\left(\boldsymbol{\vartheta}_1w_{01}+\boldsymbol{\vartheta}_2w_{02}\right)\right)\cdot\boldsymbol{\omega}_n^{(0,w_{02},w_{01})}(1).
\]
By~\eqref{it:free-coefficient-Q5} the series converges absolutely, uniformly in $\boldsymbol{\vartheta}$; so the left side is a continuous function of $\boldsymbol{\vartheta}$, and integration term by term over $[0,1]^2$ leaves the term with $w_{01}=w_{02}=0$:
\[
\mathbf{Z}_{(n_0,n_1,n_2)}=\int_{[0,1]^2}\mathrm{d}\boldsymbol{\vartheta}\,k_B^{n_0}\cdot\mathrm{tf}(n_0,\boldsymbol{\vartheta}).
\]
Fix $\boldsymbol{\vartheta}$. By the spectral theorem for compact self-adjoint operators, the non-zero eigenvalues of the restriction of $T_{\boldsymbol{\vartheta}}$ to $\mathscr{H}_{\boldsymbol{\vartheta}}^{(0)}$, repeated according to their multiplicities, form a finite or infinite sequence $(x_i)_i$ of positive numbers with $\sum_ix_i<\infty$, and $\mathrm{tf}(m,\boldsymbol{\vartheta})=\sum_ix_i^m$ for every integer $m\ge1$. Let $a\ge2$ be an integer; then $u(a-1)\ge u\ge2$. The Cauchy--Schwarz inequality for the sum over $i$ gives
\[
\mathrm{tf}(ua,\boldsymbol{\vartheta})=\sum_ix_i^{\frac{u(a-1)}{2}}\cdot x_i^{\frac{u(a+1)}{2}}\le\mathrm{tf}\left(u(a-1),\boldsymbol{\vartheta}\right)^{\frac{1}{2}}\cdot\mathrm{tf}\left(u(a+1),\boldsymbol{\vartheta}\right)^{\frac{1}{2}}.
\]
Since $k_B^{ua}=k_B^{\frac{u(a-1)}{2}}\cdot k_B^{\frac{u(a+1)}{2}}$, the Cauchy--Schwarz inequality for the integral over $[0,1]^2$ yields
\begin{align*}
\mathbf{Z}_{(ua,n_1,n_2)}&\le\int_{[0,1]^2}\mathrm{d}\boldsymbol{\vartheta}\,\left(k_B^{u(a-1)}\cdot\mathrm{tf}\left(u(a-1),\boldsymbol{\vartheta}\right)\right)^{\frac{1}{2}}\cdot\left(k_B^{u(a+1)}\cdot\mathrm{tf}\left(u(a+1),\boldsymbol{\vartheta}\right)\right)^{\frac{1}{2}}\\
&\le\mathbf{Z}_{(u(a-1),n_1,n_2)}^{\frac{1}{2}}\cdot\mathbf{Z}_{(u(a+1),n_1,n_2)}^{\frac{1}{2}}.
\end{align*}
The three numbers $\mathbf{Z}$ in this inequality are positive, by~\eqref{it:free-coefficient-Q1}, Theorem~\ref{thm:limit-vanishing-proca} and Remark~\ref{rem:tori-unequal-sides}; squaring and taking logarithms gives the convexity of $a\mapsto\log\mathbf{Z}_{(ua,n_1,n_2)}$. For the other two directions let $g$ be the exchange of the axes $0$ and $1$, respectively of the axes $0$ and $2$. By~\eqref{it:free-coefficient-Q2} with $\Phi=1$, one has $\mathbf{Z}_{gn}=\mathbf{Z}_n$ for every torus $n$; hence
\[
\mathbf{Z}_{(n_0,ua,n_2)}=\mathbf{Z}_{(ua,n_0,n_2)},\qquad\mathbf{Z}_{(n_0,n_1,ua)}=\mathbf{Z}_{(ua,n_1,n_0)},
\]
and the convexity in the directions $1$ and $2$ follows from the convexity in the direction $0$.

(c) The functionals $\boldsymbol{\omega}_n^{(w)}$ are positive, the determinant in their definition being non-negative (Theorem~\ref{thm:free-coefficient-1}(d2)), and $Z_n=\sum_w\boldsymbol{\omega}_n^{(w)}(1)$; hence
\[
\mathbf{Z}_n=\boldsymbol{\omega}_n^{(0)}(1)\le Z_n\qquad\text{for every torus }n.
\]
In particular $\mathbf{Z}_j\le Z_j$, and so $D_j=j^3\log Z_1-\log Z_j\le j^3\log Z_1-\log\mathbf{Z}_j=\bar{\mathfrak{D}}_j$. Let $\phi(a,b,c):=\log Z_{n^{*}(a,b,c)}$ be the function of~\eqref{it:free-coefficient-Q6}, formed with the integer $n^{*}\ge2$, where $n^{*}(a,b,c)$ is the torus $(an^{*},bn^{*},cn^{*})$. Then $D_j=j^3\cdot\phi(1,1,1)-\phi(j,j,j)$, and~\eqref{it:free-coefficient-Q7} gives $D_j\ge0$. The identity $\bar{\mathfrak{D}}_j=\bar{D}_j-j^3\log\frac{\mathbf{Z}_1}{Z_1}$ follows on inserting the definition of $\bar{D}_j$.

Let $k\ge j$ be an integer and $x:=\frac{k-1}{j-1}$; then $x\ge1$. Let $\boldsymbol{\psi}\colon\{1,2,\dots\}\to\mathbb{R}$ be convex in the sense of (b), that is, $\boldsymbol{\psi}(a-1)+\boldsymbol{\psi}(a+1)\ge2\boldsymbol{\psi}(a)$ for every integer $a\ge2$. Then the increments $\boldsymbol{\psi}(a+1)-\boldsymbol{\psi}(a)$ are non-decreasing in $a$, so that the function on $[1,\infty)$ which coincides with $\boldsymbol{\psi}$ at the integers and is affine on each interval $[a,a+1]$ is convex. As
\[
j=\frac{1}{x}\cdot k+\left(1-\frac{1}{x}\right)\cdot1\qquad\text{with }0<\frac{1}{x}\le1,
\]
it follows that $\boldsymbol{\psi}(j)\le\frac{1}{x}\cdot\boldsymbol{\psi}(k)+\left(1-\frac{1}{x}\right)\cdot\boldsymbol{\psi}(1)$, that is,
\[
\boldsymbol{\psi}(k)\ge x\cdot\boldsymbol{\psi}(j)-(x-1)\cdot\boldsymbol{\psi}(1).
\]
Define $\bar{\phi}\colon\{1,2,\dots\}^3\to\mathbb{R}$ by $\bar{\phi}(a,b,c):=\log\mathbf{Z}_{n^{*}(a,b,c)}$. Then $\bar{\phi}\le\phi$, by the inequality $\mathbf{Z}_n\le Z_n$, and $\bar{\phi}$ is convex in each variable, by (b) with $u=n^{*}$. By~\eqref{it:free-coefficient-Q6},
\[
\phi(1,k,k)\le k^2\cdot\phi(1,1,1),\qquad\phi(j,1,k)\le jk\cdot\phi(1,1,1),\qquad\phi(j,j,1)\le j^2\cdot\phi(1,1,1).
\]
The last inequality for $\boldsymbol{\psi}$, applied to $\bar{\phi}$ in the first, in the second and in the third variable, together with $\bar{\phi}\le\phi$ and $x-1\ge0$, gives
\begin{align*}
\bar{\phi}(k,k,k)&\ge x\cdot\bar{\phi}(j,k,k)-(x-1)\cdot k^2\cdot\phi(1,1,1),\\
\bar{\phi}(j,k,k)&\ge x\cdot\bar{\phi}(j,j,k)-(x-1)\cdot jk\cdot\phi(1,1,1),\\
\bar{\phi}(j,j,k)&\ge x\cdot\bar{\phi}(j,j,j)-(x-1)\cdot j^2\cdot\phi(1,1,1).
\end{align*}
Since $x>0$, these three inequalities combine to
\[
\bar{\phi}(k,k,k)\ge x^3\cdot\bar{\phi}(j,j,j)-\left(x^2(x-1)\cdot j^2+x(x-1)\cdot jk+(x-1)\cdot k^2\right)\cdot\phi(1,1,1).
\]
Insert $\bar{\phi}(j,j,j)=\log\mathbf{Z}_j=j^3\cdot\phi(1,1,1)-\bar{\mathfrak{D}}_j$. The identity $xj-(x-1)=k$, used three times, gives
\[
x^3j^3-x^2(x-1)\cdot j^2-x(x-1)\cdot jk-(x-1)\cdot k^2=x^2j^2k-x(x-1)\cdot jk-(x-1)\cdot k^2=xjk^2-(x-1)\cdot k^2=k^3;
\]
hence $\bar{\phi}(k,k,k)\ge k^3\cdot\phi(1,1,1)-x^3\cdot\bar{\mathfrak{D}}_j$. As $\bar{\phi}(k,k,k)=\log\mathbf{Z}_k$ and $\phi(1,1,1)=\log Z_1$, this is the asserted lower bound for $\log\mathbf{Z}_k$.

(d) Put $O_c:=1$ for the cells $c$ of the cubic torus $k\cdot n^{*}$ that do not belong to $S$. Then $O_c\in\mathfrak{A}^{+}$ for every cell $c$; moreover $1^{[c]}=1$ and $1^{\sharp}=1$, so that $\boldsymbol{\omega}_k^{(0)}\left(1^{\sharp}\right)=\mathbf{Z}_k$. The torus $k\cdot n^{*}$ has $k^3\cdot N^{*}$ cells, and $\mathbf{Z}_k$ is the product of the numbers $\mathbf{Z}_k^{\frac{1}{k^3\cdot N^{*}}}$ over all of them. Proposition~\ref{prop:chessboard-estimate} for $\boldsymbol{\omega}_k^{(0)}$, divided by $\mathbf{Z}_k$, therefore gives
\[
\left|\bar{E}_k\prod_{c\in S}O_c^{[c]}\right|\le\prod_{c\in S}\left(\frac{\boldsymbol{\omega}_k^{(0)}\left(O_c^{\sharp}\right)}{\mathbf{Z}_k}\right)^{\frac{1}{k^3\cdot N^{*}}},
\]
where the bases of the powers are non-negative, again by Proposition~\ref{prop:chessboard-estimate}. For $k=1$ the right side is $\prod_{c\in S}\|O_c\|_{0*}$, since $\bar{E}_{*}O^{\sharp}=\boldsymbol{\omega}_1^{(0)}\left(O^{\sharp}\right)/\mathbf{Z}_1$ for $O\in\mathfrak{A}^{+}$; this is the second assertion. Let now $k\ge j$, and let $O\in\mathfrak{A}^{+}$. By (a), applied on the torus $k\cdot n^{*}$, and by Proposition~\ref{prop:lengthening-one-side}(a), applied to $n^{*}$ in the direction $0$, to $n^{*}(k,1,1)$ in the direction $1$ and to $n^{*}(k,k,1)$ in the direction $2$ (as in the proof of Theorem~\ref{thm:transfer-all-flux}),
\[
0\le\boldsymbol{\omega}_k^{(0)}\left(O^{\sharp}\right)\le\boldsymbol{\omega}_k\left(O^{\sharp}\right)\le\boldsymbol{\omega}_1\left(O^{\sharp}\right)^{k^3}.
\]
By the definition of $\|O\|_{*}$ in Section~\ref{sec:transfer-larger-cubic}, $\boldsymbol{\omega}_1\left(O^{\sharp}\right)=Z_1\cdot\|O\|_{*}^{N^{*}}$. It follows that
\[
\left(\frac{\boldsymbol{\omega}_k^{(0)}\left(O^{\sharp}\right)}{\mathbf{Z}_k}\right)^{\frac{1}{k^3\cdot N^{*}}}\le\left(\frac{\boldsymbol{\omega}_1\left(O^{\sharp}\right)^{k^3}}{\mathbf{Z}_k}\right)^{\frac{1}{k^3\cdot N^{*}}}=\|O\|_{*}\cdot\exp\left(\frac{\log Z_1-k^{-3}\log\mathbf{Z}_k}{N^{*}}\right).
\]
By (c),
\[
\log Z_1-k^{-3}\log\mathbf{Z}_k\le\left(\frac{k-1}{k}\right)^3\cdot\frac{\bar{\mathfrak{D}}_j}{(j-1)^3}\le\frac{\bar{\mathfrak{D}}_j}{(j-1)^3},
\]
because $\bar{\mathfrak{D}}_j\ge0$ and $0\le k-1\le k$. With $O=O_c$ for $c\in S$, the first assertion follows.

(e) By Section~\ref{sec:wilson-determinant-lattice},
\[
E_1O^{\sharp}=\frac{1}{Z_1}\left(\boldsymbol{\omega}_1^{(0)}\left(O^{\sharp}\right)+\sum_{w\ne0}\boldsymbol{\omega}_1^{(w)}\left(O^{\sharp}\right)\right),
\]
where the functionals $\boldsymbol{\omega}_1^{(w)}$ of the cubic torus $n^{*}$ are positive (Theorem~\ref{thm:free-coefficient-1}(d2)). The first term equals $\frac{\mathbf{Z}_1}{Z_1}\cdot\|O\|_{0*}^{N^{*}}$, which is at most $\|O\|_{0*}^{N^{*}}$, as $\mathbf{Z}_1\le Z_1$ (proof of (c)). For $w\ne0$ the Cauchy--Schwarz inequality for the positive functional $\boldsymbol{\omega}_1^{(w)}$ gives
\[
\left|\boldsymbol{\omega}_1^{(w)}\left(O^{\sharp}\right)\right|\le\boldsymbol{\omega}_1^{(w)}(1)^{\frac{1}{2}}\cdot\boldsymbol{\omega}_1^{(w)}\left(\left|O^{\sharp}\right|^2\right)^{\frac{1}{2}}.
\]
The Cauchy--Schwarz inequality for the sum over $w\ne0$, together with $\sum_{w\ne0}\boldsymbol{\omega}_1^{(w)}\left(\left|O^{\sharp}\right|^2\right)\le\boldsymbol{\omega}_1\left(\left|O^{\sharp}\right|^2\right)=Z_1\cdot E_1\left|O^{\sharp}\right|^2$, then yields
\[
E_1O^{\sharp}\le\|O\|_{0*}^{N^{*}}+\left(\frac{1}{Z_1}\sum_{w\ne0}\boldsymbol{\omega}_1^{(w)}(1)\right)^{\frac{1}{2}}\cdot\left(E_1\left|O^{\sharp}\right|^2\right)^{\frac{1}{2}}.
\]
By Proposition~\ref{prop:comparison-flux-sectors}(ii), $\boldsymbol{\omega}_1^{(w)}(1)\le\mathfrak{q}_w(n^{*})\cdot Z_1$ for every $w$; hence $\frac{1}{Z_1}\sum_{w\ne0}\boldsymbol{\omega}_1^{(w)}(1)\le\varsigma_{n^{*}}$. The function $\left|O^{\sharp}\right|^2$ is the product of the $N^{*}$ functions $\left|O^{[c]}\right|^2=|O|^2\circ\mathsf{g}_c^{*}$, and $E_1$ is invariant under the lattice isometries $\mathsf{g}_c$, by~\eqref{it:free-coefficient-Q2}; so H\"older's inequality with $N^{*}$ exponents equal to $N^{*}$ gives
\[
E_1\left|O^{\sharp}\right|^2\le\prod_c\left(E_1\left|O^{[c]}\right|^{2N^{*}}\right)^{\frac{1}{N^{*}}}=E_1|O|^{2N^{*}}.
\]
Altogether
\[
\|O\|_{*}^{N^{*}}=E_1O^{\sharp}\le\|O\|_{0*}^{N^{*}}+\varsigma_{n^{*}}^{\frac{1}{2}}\cdot\left(E_1|O|^{2N^{*}}\right)^{\frac{1}{2}},
\]
where $E_1O^{\sharp}\ge0$ by Proposition~\ref{prop:chessboard-estimate}. Since $(y+y')^{\frac{1}{N^{*}}}\le y^{\frac{1}{N^{*}}}+y'^{\frac{1}{N^{*}}}$ for $y,y'\ge0$, taking the $N^{*}$-th root gives (e).
\end{proof}

\begin{remark}[Cumulant bounds and pressure for the harmonic-averaged law]\label{rem:cumulant-bounds-pressure}
Let the data be those of Theorem~\ref{thm:free-coefficient-2}: integers $N\ge0$ and $T^{*}\ge1$, the lattice mass $\boldsymbol{\mu}_N>-\frac{1}{2}$, a charge $0<e\le\frac{1}{3}$ subject to~\eqref{eq:free-coefficient-6}, a probe $\beta$ that fits $(N,T^{*})$, with its placement, and real numbers $c_2,c_4$. As in the proof of Theorem~\ref{thm:free-coefficient-2} in Section~\ref{sec:proof-cumulant-bounds}, the data of Section~\ref{sec:transfer-larger-cubic} are taken to be those of Remark~\ref{rem:translated-cells-aligned}: $q=e\cdot\eta^{\frac{1}{2}}$, $\boldsymbol{\mu}=\boldsymbol{\mu}_N$, $\boldsymbol{\lambda}=m_{\mathrm{p}}+\frac{1}{2}$ (where $m_{\mathrm{p}}\ge7$, as $\beta$ fits $(N,T^{*})$; Section~\ref{sec:probes-discretisation}), $\mathsf{n}=2L$ and $j=L$. Thus $n^{*}=n_{*}$ and $N^{*}=8L^3$, and for every integer $k\ge1$ the numbers $Z_k$, $\mathbf{Z}_k$ and the functional $\bar{E}_k=E^{\mathrm{av}}_k$ refer to the torus $k$ of Section~\ref{sec:reference-tori-cells}, with the cells of that section.

\emph{A bound for $\bar{\mathfrak{D}}_L$.} Let $D^{\mathrm{ph}}_L$ be the number $D^{\mathrm{ph}}_j$ of Lemma~\ref{lem:partition-functions-cubic}(d) for $n=n_{*}$ and $j=L$, formed with $q$ and with the lattice mass $\boldsymbol{\mu}_N$. By Lemma~\ref{lem:partition-functions-cubic}(b), applied on the cubic tori $n_{*}$ and $L\cdot n_{*}$,
\[
\log Z_1=\log Z^{\mathrm{ph}}_{n_{*}}+\log\rho_{n_{*}}-\log P_{n_{*}}(0),\qquad\log\mathbf{Z}_L=\log Z^{\mathrm{ph}}_{L\cdot n_{*}}+\log\rho_{L\cdot n_{*}},
\]
and $L^3\log Z^{\mathrm{ph}}_{n_{*}}-\log Z^{\mathrm{ph}}_{L\cdot n_{*}}=D^{\mathrm{ph}}_L$ by Lemma~\ref{lem:partition-functions-cubic}(a), as in the proof of Lemma~\ref{lem:partition-functions-cubic}(d). Since $D^{\mathrm{mat}}=L^3\log\rho_{n_{*}}-\log\rho_{L\cdot n_{*}}$ (Section~\ref{sec:numbers-attached-reference}), it follows that
\[
\bar{\mathfrak{D}}_L=D^{\mathrm{ph}}_L+D^{\mathrm{mat}}-L^3\log P_{n_{*}}(0).
\]
Step~1 of the proof of Theorem~\ref{thm:free-coefficient-2} shows that $D^{\mathrm{ph}}_L\le\left(L^3-1\right)\log\frac{2\pi}{e^2\cdot R^{*}}+K_L$, that $\varsigma_{n_{*}}\le\frac{1}{8}$, and that $-L^3\log\left(1-\varsigma_{n_{*}}\right)\le1$; the last inequality is the one by which that step bounds $D^{\mathrm{fl}}_L$. By Lemma~\ref{lem:partition-functions-cubic}(b), $P_{n_{*}}(0)\ge1-\varsigma_{n_{*}}>0$; hence
\[
-L^3\log P_{n_{*}}(0)\le-L^3\log\left(1-\varsigma_{n_{*}}\right)\le1,
\]
and the identity for $\bar{\mathfrak{D}}_L$ yields
\[
\bar{\mathfrak{D}}_L\le\left(L^3-1\right)\log\frac{2\pi}{e^2\cdot R^{*}}+K_L+D^{\mathrm{mat}}+1.
\]
Division by $(L-1)^3\cdot N^{*}$, together with $a_L=\frac{L^3-1}{(L-1)^3}$, gives
\[
\frac{\bar{\mathfrak{D}}_L}{(L-1)^3\cdot N^{*}}\le\frac{a_L}{N^{*}}\cdot\log\frac{2\pi}{e^2\cdot R^{*}}+\frac{K_L+D^{\mathrm{mat}}+1}{(L-1)^3\cdot N^{*}},
\]
and the right side is the right side of~\eqref{eq:free-coefficient-9}.

\emph{Bounds for the probe field.} Let $\boldsymbol{\Gamma}'$ be the exponential of the right side of~\eqref{eq:free-coefficient-9}:
\[
\boldsymbol{\Gamma}':=\exp\left(\frac{a_L}{N^{*}}\cdot\log\frac{2\pi}{e^2\cdot R^{*}}+\frac{K_L+D^{\mathrm{mat}}+1}{(L-1)^3\cdot N^{*}}\right).
\]
Let $m\in\{2,4\}$, and let $O_m=e^{-m}\cdot H_m(X)-c_m$ be the function of Step~3 of the proof of Theorem~\ref{thm:free-coefficient-2}; it belongs to $\mathfrak{A}\subset\mathfrak{A}^{+}$, and $O_m^{[c]}=e^{-m}\cdot H_m(X_c)-c_m$ for every cell $c$ of every torus $k$. Step~4 of that proof gives $\|O_m\|_{*}\le\boldsymbol{\Theta}^{\mathrm{af}}_m$, and~\eqref{eq:free-coefficient-18} gives $\|O_m\|_{0*}=\Theta^{(0)}_m$. Let $k$ be $1$ or an integer $k\ge L$, and let $S$ be a set of cells of the torus $k$. If $k\ge L$, Proposition~\ref{prop:zero-flux-sector}(d), applied with $O_c:=O_m$ for $c\in S$, and the bound for $\bar{\mathfrak{D}}_L$ give
\[
\left|E^{\mathrm{av}}_k\prod_{c\in S}\left(e^{-m}\cdot H_m(X_c)-c_m\right)\right|\le\left(\exp\left(\frac{\bar{\mathfrak{D}}_L}{(L-1)^3\cdot N^{*}}\right)\cdot\|O_m\|_{*}\right)^{\#|S|}\le\left(\boldsymbol{\Gamma}'\cdot\boldsymbol{\Theta}^{\mathrm{af}}_m\right)^{\#|S|}.
\]
If $k=1$, the same proposition gives
\[
\left|E^{\mathrm{av}}_1\prod_{c\in S}\left(e^{-m}\cdot H_m(X_c)-c_m\right)\right|\le\left(\Theta^{(0)}_m\right)^{\#|S|}.
\]
Thus~\eqref{eq:free-coefficient-10} holds for $\mathrm{nat}=\mathrm{av}$ with $\boldsymbol{\Gamma}^{\mathrm{av}}\cdot\boldsymbol{\Theta}^{\mathrm{av}}_m$ replaced by $\boldsymbol{\Gamma}'\cdot\boldsymbol{\Theta}^{\mathrm{af}}_m$ if $k\ge L$, and by $\Theta^{(0)}_m$ if $k=1$. Step~2 of the proof of Theorem~\ref{thm:free-coefficient-2} derives~\eqref{eq:free-coefficient-7} and~\eqref{eq:free-coefficient-8} from the case $S=\{0\}$ of~\eqref{eq:free-coefficient-10}, whatever non-negative numbers stand in the place of the two products $\boldsymbol{\Gamma}^{\mathrm{nat}}\cdot\boldsymbol{\Theta}^{\mathrm{nat}}_m$, $m=2,4$. Hence~\eqref{eq:free-coefficient-7}, \eqref{eq:free-coefficient-8} and~\eqref{eq:free-coefficient-10} hold for $\mathrm{nat}=\mathrm{av}$ and every integer $k\ge L$ with $\boldsymbol{\Gamma}^{\mathrm{av}}\cdot\boldsymbol{\Theta}^{\mathrm{av}}_m$ replaced by $\boldsymbol{\Gamma}'\cdot\boldsymbol{\Theta}^{\mathrm{af}}_m$, and for $k=1$ with $\boldsymbol{\Gamma}^{\mathrm{av}}\cdot\boldsymbol{\Theta}^{\mathrm{av}}_m$ replaced by $\Theta^{(0)}_m$.

\emph{The pressure.} For integers $k\ge1$ put
\[
p_k:=(k\cdot n_{*})^{-3}\log Z_k,\qquad\bar{p}_k:=(k\cdot n_{*})^{-3}\log\mathbf{Z}_k.
\]
Then, for every integer $k\ge L$,
\[
p_1-\frac{\bar{\mathfrak{D}}_L}{(L-1)^3\cdot n_{*}^3}\le\bar{p}_k\le p_k\le p_1.
\]
Indeed, Proposition~\ref{prop:zero-flux-sector}(c) with $j=L$ gives $\bar{\mathfrak{D}}_L\ge0$ and
\[
\log\mathbf{Z}_k\ge k^3\log Z_1-\left(\frac{k-1}{L-1}\right)^3\cdot\bar{\mathfrak{D}}_L\ge k^3\cdot\left(\log Z_1-\frac{\bar{\mathfrak{D}}_L}{(L-1)^3}\right),
\]
as $0\le k-1\le k$; division by $(k\cdot n_{*})^3$ gives the first inequality. The second follows from $\mathbf{Z}_k\le Z_k$ (proof of Proposition~\ref{prop:zero-flux-sector}(c)). For the third let $\phi$ be the function of~\eqref{it:free-coefficient-Q6}, formed with $n^{*}=n_{*}$; applying the inequality $\phi(ka,b,c)\le k\cdot\phi(a,b,c)$ and its analogues in the second and in the third variable, one obtains $\log Z_k=\phi(k,k,k)\le k^3\cdot\phi(1,1,1)=k^3\log Z_1$.
\end{remark}

\subsection{Partition functions of cubic tori, cumulant bounds and invariant functions}\label{sec:partition-functions-cubic}

\subsubsection{Eigenvalues of the Laplacian and partition functions of a cubic torus}\label{sec:eigenvalues-laplacian-partition}

For a triple $n=(n_1,n_2,n_3)$ of integers $n_{\mu}\ge2$, as in Section~\ref{sec:wilson-determinant-lattice}, let
\[
\mathrm{Ld}(n):=\sum_{k\ne0}\log\left(\sum_{\mu=1}^3\left(2-2\cos\frac{2\pi k_{\mu}}{n_{\mu}}\right)\right),
\]
the sum being taken over the elements $k=(k_1,k_2,k_3)\ne0$ of $\prod_{\mu=1}^3\mathbb{Z}/n_{\mu}\mathbb{Z}$. (For $k\ne0$ at least one class $k_{\mu}$ is non-zero, so that the argument of the logarithm is positive.) For an integer $n\ge2$, that is, for the cubic torus $n$ of Section~\ref{sec:wilson-determinant-lattice}, let $\mathrm{Ld}(n):=\mathrm{Ld}(n,n,n)$.

\begin{lemma}\label{lem:free-coefficient-3}
Let $K_j$ be the constants of Section~\ref{sec:numbers-attached-reference}. For all integers $n\ge2$ and $j\ge2$,
\[
\frac{1}{2}\left(\mathrm{Ld}(jn)-j^3\cdot\mathrm{Ld}(n)\right)\le-\left(j^3-1\right)\log n+K_j.
\]
\end{lemma}

\begin{proof}
Let $G_j$ and $K_j$ be the constants of Section~\ref{sec:numbers-attached-reference}. For $p=(p_1,p_2)\in[-\pi,\pi]^2$ let $\boldsymbol{\omega}_p\ge0$ be the number with
\[
2\cosh\boldsymbol{\omega}_p=2+\sum_{i=1}^2\left(2-2\cos p_i\right)
\]
(the right side is at least $2$).

\emph{Step 1.} In Steps~1 and~2, $n_1,n_2,n_3$ are arbitrary integers $\ge2$, independent of $n$. Let $\boldsymbol{\zeta}:=\exp\left(\frac{2\pi i}{n_1}\right)$. For every real number $\boldsymbol{\omega}$ and every integer $k$ one has
\[
\left(\exp(\boldsymbol{\omega})-\boldsymbol{\zeta}^k\right)\cdot\left(1-\exp(-\boldsymbol{\omega})\cdot\boldsymbol{\zeta}^{-k}\right)=2\cosh\boldsymbol{\omega}-2\cos\frac{2\pi k}{n_1}.
\]
The numbers $\boldsymbol{\zeta}^k$, $k\in\mathbb{Z}/n_1\mathbb{Z}$, are the roots of the polynomial $x^{n_1}-1$. Hence, with both products taken over $k\in\mathbb{Z}/n_1\mathbb{Z}$,
\[
\prod_k\left(\exp(\boldsymbol{\omega})-\boldsymbol{\zeta}^k\right)=\exp\left(n_1\cdot\boldsymbol{\omega}\right)-1,\qquad\prod_k\left(1-\exp(-\boldsymbol{\omega})\cdot\boldsymbol{\zeta}^{-k}\right)=1-\exp\left(-n_1\cdot\boldsymbol{\omega}\right).
\]
Multiplication of these two identities gives
\[
\prod_{k\in\mathbb{Z}/n_1\mathbb{Z}}\left(2\cosh\boldsymbol{\omega}-2\cos\frac{2\pi k}{n_1}\right)=2\cosh\left(n_1\cdot\boldsymbol{\omega}\right)-2=4\sinh^2\left(\frac{n_1\cdot\boldsymbol{\omega}}{2}\right).
\]
Moreover $2-2\cos\frac{2\pi k}{n_1}=\left|1-\boldsymbol{\zeta}^k\right|^2$, and $\prod_{k\ne0}\left(1-\boldsymbol{\zeta}^k\right)=n_1$, this being the value at $x=1$ of $\prod_{k\ne0}\left(x-\boldsymbol{\zeta}^k\right)=1+x+\dots+x^{n_1-1}$. Therefore
\[
\prod_{k\in\mathbb{Z}/n_1\mathbb{Z},\,k\ne0}\left(2-2\cos\frac{2\pi k}{n_1}\right)=\prod_{k\ne0}\left|1-\boldsymbol{\zeta}^k\right|^2=n_1^2.
\]

Write an element of $\prod_{\mu}\mathbb{Z}/n_{\mu}\mathbb{Z}$ as $k=(k_1,k_s)$ with $k_s:=(k_2,k_3)$, represent the classes $k_2$, $k_3$ by integers with $|k_i|\le\frac{n_i}{2}$, and let
\[
p(k_s):=\left(\frac{2\pi k_2}{n_2},\frac{2\pi k_3}{n_3}\right)\in[-\pi,\pi]^2.
\]
If $k_s\ne0$, then $p(k_s)\ne0$, hence $\boldsymbol{\omega}_{p(k_s)}>0$. In the sum that defines $\mathrm{Ld}(n_1,n_2,n_3)$, the terms with $k_s=0$ have $k_1\ne0$, and by the second product formula their sum is $2\log n_1$. For fixed $k_s\ne0$ one has
\[
\sum_{\mu=1}^3\left(2-2\cos\frac{2\pi k_{\mu}}{n_{\mu}}\right)=2\cosh\boldsymbol{\omega}_{p(k_s)}-2\cos\frac{2\pi k_1}{n_1},
\]
so that, by the first product formula, the sum of the corresponding terms over $k_1\in\mathbb{Z}/n_1\mathbb{Z}$ is $\log\left(4\sinh^2\left(\frac{n_1\cdot\boldsymbol{\omega}_{p(k_s)}}{2}\right)\right)$. It follows that
\[
\mathrm{Ld}(n_1,n_2,n_3)=2\log n_1+2\sum_{k_s\ne0}\log\left(2\sinh\left(\frac{n_1\cdot\boldsymbol{\omega}_{p(k_s)}}{2}\right)\right),
\]
where $k_s$ ranges over the non-zero elements of $\mathbb{Z}/n_2\mathbb{Z}\times\mathbb{Z}/n_3\mathbb{Z}$, represented as above.

\emph{Step 2.} For $x>0$ let $\mathfrak{g}(x):=-\log\left(1-\exp(-2x)\right)$; thus $\mathfrak{g}(x)>0$. Since $2\sinh y=\exp(y)\cdot\left(1-\exp(-2y)\right)$ for real $y$, one has for $x>0$
\[
\log\left(2\sinh(jx)\right)-j\log\left(2\sinh x\right)=\log\left(1-\exp(-2jx)\right)+j\cdot\mathfrak{g}(x)\le j\cdot\mathfrak{g}(x).
\]
The triples $(jn_1,n_2,n_3)$ and $(n_1,n_2,n_3)$ have the same points $p(k_s)$. Insert the formula of Step~1 for both triples and apply the last inequality with $x=\frac{n_1\cdot\boldsymbol{\omega}_{p(k_s)}}{2}$ for every $k_s\ne0$; this gives, for all integers $n_1,n_2,n_3\ge2$,
\begin{equation}\label{eq:free-coefficient-17}
\mathrm{Ld}(jn_1,n_2,n_3)-j\cdot\mathrm{Ld}(n_1,n_2,n_3)\le2\log(jn_1)-2j\log n_1+2j\sum_{k_s\ne0}\mathfrak{g}\left(\frac{n_1\cdot\boldsymbol{\omega}_{p(k_s)}}{2}\right).
\end{equation}

\emph{Step 3.} For $p\in[-\pi,\pi]^2$ let $\tilde{M}:=\sum_{i=1}^2\left(2-2\cos p_i\right)$. Since $\cosh\boldsymbol{\omega}_p-1=2\sinh^2\left(\frac{\boldsymbol{\omega}_p}{2}\right)$, the definition of $\boldsymbol{\omega}_p$ gives $\sinh^2\left(\frac{\boldsymbol{\omega}_p}{2}\right)=\frac{\tilde{M}}{4}$. From $2-2\cos y=4\sin^2\frac{y}{2}$ and $\left|\sin\frac{y}{2}\right|\ge\frac{|y|}{\pi}$ for $|y|\le\pi$ one obtains
\[
\frac{4}{\pi^2}\cdot|p|_2^2\le\tilde{M}\le8\qquad\text{on }[-\pi,\pi]^2.
\]
The function $\operatorname{arcsinh}$ is concave on $[0,\infty)$ and vanishes at $0$; hence
\[
\operatorname{arcsinh} y\ge y\cdot\frac{\operatorname{arcsinh}\sqrt{2}}{\sqrt{2}}\ge0.81\cdot y\qquad\text{for }y\in\left[0,\sqrt{2}\right].
\]
As $\boldsymbol{\omega}_p=2\operatorname{arcsinh}\left(\frac{1}{2}\sqrt{\tilde{M}}\right)$ and $\frac{1}{2}\sqrt{\tilde{M}}\le\sqrt{2}$, it follows that
\[
\boldsymbol{\omega}_p\ge0.81\cdot\sqrt{\tilde{M}}\ge0.51\cdot|p|_2.
\]

Let now $n_1=n$ and $n_2,n_3\in\{n,jn\}$, and let $k_s\ne0$. Since $\frac{n}{n_i}\ge\frac{1}{j}$ for $i=2,3$ and $0.255\cdot2\pi>1.6$, the last inequality gives
\[
\frac{n\cdot\boldsymbol{\omega}_{p(k_s)}}{2}\ge0.255\cdot2\pi\cdot\left|\left(\frac{k_2\cdot n}{n_2},\frac{k_3\cdot n}{n_3}\right)\right|_2\ge\mathbf{c}\cdot|k_s|_2,\qquad\mathbf{c}:=\frac{1.6}{j},
\]
where $|k_s|_2$ is the Euclidean norm of the integer vector $(k_2,k_3)$. The inequalities $-\log(1-y)\le\frac{y}{1-y}$ for $0\le y<1$ and $\frac{1}{1-\exp(-y)}\le1+\frac{1}{y}$ for $y>0$ yield
\[
\mathfrak{g}(x)\le\exp(-2x)\cdot\left(1+\frac{1}{2x}\right)\qquad(x>0),
\]
and the right side is a decreasing function of $x$. The integer vector $k_s\ne0$ satisfies $|k_s|_2\ge1$ and $|k_s|_2\ge\frac{|k_2|+|k_3|}{\sqrt{2}}$. Together with the lower bound for $\frac{n\cdot\boldsymbol{\omega}_{p(k_s)}}{2}$ this gives
\[
\mathfrak{g}\left(\frac{n\cdot\boldsymbol{\omega}_{p(k_s)}}{2}\right)\le\left(1+\frac{1}{2\mathbf{c}}\right)\cdot\exp\left(-2\mathbf{c}\cdot|k_s|_2\right)\le\left(1+\frac{1}{2\mathbf{c}}\right)\cdot\exp\left(-\sqrt{2}\cdot\mathbf{c}\cdot\left(|k_2|+|k_3|\right)\right).
\]
The vectors $k_s\ne0$ are distinct elements of $\mathbb{Z}^2\setminus\{0\}$. Summing over them and using $\sum_{i\in\mathbb{Z}}\exp(-a|i|)=\coth\frac{a}{2}$ for $a>0$ (a geometric series), the identity $\coth^2y-1=\left(\sinh y\right)^{-2}$ for $y>0$ and the inequality $\sinh y\ge y$ for $y\ge0$, one obtains
\begin{align*}
\sum_{k_s\ne0}\mathfrak{g}\left(\frac{n\cdot\boldsymbol{\omega}_{p(k_s)}}{2}\right)&\le\left(1+\frac{1}{2\mathbf{c}}\right)\cdot\left(\left(\sum_{i\in\mathbb{Z}}\exp\left(-\sqrt{2}\cdot\mathbf{c}\cdot|i|\right)\right)^2-1\right)\\
&=\left(1+\frac{1}{2\mathbf{c}}\right)\cdot\left(\coth^2\left(\frac{\mathbf{c}}{\sqrt{2}}\right)-1\right)=\left(1+\frac{1}{2\mathbf{c}}\right)\cdot\left(\sinh\frac{\mathbf{c}}{\sqrt{2}}\right)^{-2}\\
&\le\left(1+\frac{1}{2\mathbf{c}}\right)\cdot\frac{2}{\mathbf{c}^2}=G_j.
\end{align*}
The last equality is the definition of $G_j$, since $2\mathbf{c}=\frac{3.2}{j}$ and $\mathbf{c}^2=\frac{2.56}{j^2}$.

\emph{Step 4.} The function $\mathrm{Ld}$ is symmetric in its three arguments, because a permutation of the coordinates of $k$ permutes the summands. One has
\begin{align*}
\mathrm{Ld}(jn,jn,jn)-j^3\cdot\mathrm{Ld}(n,n,n)={}&\left(\mathrm{Ld}(jn,jn,jn)-j\cdot\mathrm{Ld}(n,jn,jn)\right)+j\cdot\left(\mathrm{Ld}(n,jn,jn)-j\cdot\mathrm{Ld}(n,n,jn)\right)\\
&+j^2\cdot\left(\mathrm{Ld}(n,n,jn)-j\cdot\mathrm{Ld}(n,n,n)\right).
\end{align*}
After a permutation of the arguments that moves the side multiplied by $j$ to the first place, each of the three brackets is the left side of \eqref{eq:free-coefficient-17} with $n_1=n$ and with cross-section sides $n_2,n_3\in\{n,jn\}$. By \eqref{eq:free-coefficient-17} and Step~3, each bracket is therefore at most
\[
2\log(jn)-2j\log n+2j\cdot G_j=2\log j+2(1-j)\log n+2j\cdot G_j.
\]
The coefficients $1$, $j$, $j^2$ of the brackets are positive, and $\left(1+j+j^2\right)\cdot(1-j)=1-j^3$. With the definition of $K_j$ it follows that
\[
\mathrm{Ld}(jn)-j^3\cdot\mathrm{Ld}(n)\le\left(1+j+j^2\right)\cdot\left(2\log j+2(1-j)\log n+2j\cdot G_j\right)=-2\left(j^3-1\right)\log n+2K_j.
\]
Division by $2$ gives the assertion.
\end{proof}

\begin{lemma}[Partition functions of a cubic torus]\label{lem:partition-functions-cubic}
Let $n\ge2$ be an integer, let $q>0$ and $\boldsymbol{\mu}>-\frac{1}{2}$, and let $Z_n$, $\mathbf{Z}_n$, $P_n(0)$ and $\varsigma_n$ be the numbers of Section~\ref{sec:wilson-determinant-lattice} for the cubic torus $n$, formed with this $q$ and with the lattice mass $\boldsymbol{\mu}$. Let
\[
\Gamma_n:=r_q\cdot\left\{m\in C^1(n;\mathbb{Z}):\ dm=0\right\}
\]
be the image of the group of the closed integer one-cochains under multiplication by $r_q=2\pi/q$, and let
\[
Z^{\mathrm{ph}}_n:=\int_{C^1(n)/\Gamma_n}\mathrm{d}A\,\exp\left(-\frac{1}{2}\|dA\|^2\right);
\]
here the integrand is invariant under the translations by the elements of $\Gamma_n$, and the quotient carries the measure induced by the Lebesgue measure $\mathrm{d}A$ of $C^1(n)$. Let $\rho_n$ be the number of Section~\ref{sec:numbers-attached-reference}, formed with this $q$ and with $\boldsymbol{\mu}$ in place of $\boldsymbol{\mu}_N$:
\[
\rho_n=\frac{1}{(2\pi)^3}\int_{[-\pi,\pi]^3}\mathrm{d}\theta\int_{V_{\mathrm{t}}}\mathrm{d}\tilde{\gamma}_n(A)\,\det\left(D_{A+H_{\theta}}+\boldsymbol{\mu}\right).
\]
Let $\mathrm{Ld}(n)$ be the number of Section~\ref{sec:eigenvalues-laplacian-partition}, and let $K_j$ be the numbers of Section~\ref{sec:numbers-attached-reference}.
\begin{enumerate}
\item[(a)] One has
\[
Z^{\mathrm{ph}}_n=\left(\frac{2\pi}{q}\right)^{n^3+2}\cdot(2\pi)^{n^3-1}\cdot\exp\left(-\frac{1}{2}\mathrm{Ld}(n)\right).
\]
\item[(b)] One has $\mathbf{Z}_n=Z^{\mathrm{ph}}_n\cdot\rho_n$ and $Z_n=Z^{\mathrm{ph}}_n\cdot\rho_n/P_n(0)$, and $1-\varsigma_n\le P_n(0)\le1$.
\item[(c)] If $q^2n\le1$, then
\[
\varsigma_n\le7\exp\left(-\frac{2\pi^2}{q^2n}\right)\le\frac{1}{8}.
\]
\item[(d)] Let $j\ge2$ be an integer, let $Z_{jn}$, $\rho_{jn}$ and $P_{jn}(0)$ be formed in the same way on the cubic torus $jn$, and put
\begin{align*}
D_j&:=j^3\log Z_n-\log Z_{jn},\\
D^{\mathrm{ph}}_j&:=\left(j^3-1\right)\log\frac{2\pi}{q^2}+\frac{1}{2}\left(\mathrm{Ld}(jn)-j^3\cdot\mathrm{Ld}(n)\right),\\
D^{\mathrm{mat}}_j&:=j^3\log\rho_n-\log\rho_{jn},\\
D^{\mathrm{fl}}_j&:=-j^3\log P_n(0)+\log P_{jn}(0).
\end{align*}
Then
\[
D_j=D^{\mathrm{ph}}_j+D^{\mathrm{mat}}_j+D^{\mathrm{fl}}_j,\qquad D^{\mathrm{ph}}_j\le\left(j^3-1\right)\log\frac{2\pi}{q^2n}+K_j,
\]
and, if $\varsigma_n<1$, then $D^{\mathrm{fl}}_j\le-j^3\log\left(1-\varsigma_n\right)$.
\end{enumerate}
\end{lemma}

\begin{proof}
Let $C^1(n)=V_{\mathrm{g}}\oplus V_{\mathrm{h}}\oplus V_{\mathrm{t}}$ be the orthogonal decomposition of Section~\ref{sec:wilson-determinant-lattice}, let $\sigma$ be the Lebesgue measure of $V_{\mathrm{t}}$, and let $\pi_h$ be the orthogonal projection of $C^1(n)$ onto $V_{\mathrm{h}}$. By the Hodge decomposition on the discrete torus, the kernel of $d$ on $C^1(n)$ is $V_{\mathrm{g}}\oplus V_{\mathrm{h}}$; hence $\Gamma_n\subset V_{\mathrm{g}}\oplus V_{\mathrm{h}}$, and $d^{*}d$ is positive definite on $V_{\mathrm{t}}$. For a lattice $\Gamma$ of full rank in a finite-dimensional Euclidean space $W$, let $\operatorname{covol}(\Gamma)$ be the volume of a fundamental domain of $\Gamma$ in $W$.

(a) The scalar Laplacian $d^{*}d$ on $C^0(n)$ has the eigenvalues $\sum_{\mu}\left(2-2\cos\frac{2\pi k_{\mu}}{n}\right)$, $k\in(\mathbb{Z}/n\mathbb{Z})^3$; its kernel consists of the constants, and, by the definition of $\mathrm{Ld}(n)$, the product of its non-zero eigenvalues is $\exp\left(\mathrm{Ld}(n)\right)$.

If a coboundary $d\lambda$ of a zero-cochain $\lambda$ takes integer values, then $\lambda-\lambda(0)$ takes integer values, the torus being connected; hence $\Gamma_n\cap V_{\mathrm{g}}=r_q\cdot dC^0(n;\mathbb{Z})$. Let $W_0\subset C^0(n)$ be the orthogonal complement of the constants. The coboundary maps $W_0$ bijectively onto $V_{\mathrm{g}}$, and the absolute value of its determinant is $\det\left(d^{*}d|_{W_0}\right)^{\frac{1}{2}}=\exp\left(\frac{1}{2}\mathrm{Ld}(n)\right)$. The orthogonal projection of $C^0(n;\mathbb{Z})=\mathbb{Z}^{n^3}$ onto $W_0$ is a lattice of covolume $n^{-\frac{3}{2}}$ in $W_0$, because the primitive vector $(1,\dots,1)$ of $\mathbb{Z}^{n^3}$ has the length $n^{\frac{3}{2}}$; and the image of this projection under $d$ is $dC^0(n;\mathbb{Z})$. As $\dim V_{\mathrm{g}}=n^3-1$, it follows that $\Gamma_n\cap V_{\mathrm{g}}$ is a lattice of full rank in $V_{\mathrm{g}}$, with
\[
\operatorname{covol}\left(\Gamma_n\cap V_{\mathrm{g}}\right)=r_q^{n^3-1}\cdot\exp\left(\frac{1}{2}\mathrm{Ld}(n)\right)\cdot n^{-\frac{3}{2}}.
\]
For $\mu=1,2,3$ let $b_{\mu}\in V_{\mathrm{h}}$ be the cochain with the value $n^{-\frac{3}{2}}$ on the links of direction $\mu$ and the value $0$ on the other links; then $b_1,b_2,b_3$ is an orthonormal basis of $V_{\mathrm{h}}$. Let $m$ be a closed integer one-cochain. Since $dm=0$, the sum of $m$ over the $n$ links of a line of direction $\mu$ is one and the same integer $k_{\mu}$ for all these lines; hence the mean of $m$ over the $n^3$ links of direction $\mu$ is $k_{\mu}/n$, that is, the coordinate of $\pi_hm$ along $b_{\mu}$ is $k_{\mu}n^{\frac{1}{2}}$. Every $k\in\mathbb{Z}^3$ occurs: take for $m$ the cochain with the value $k_{\mu}$ on the links $(x,x+e_{\mu})$ with $x_{\mu}=0$, $\mu=1,2,3$, and the value $0$ on the other links. Therefore
\[
\pi_h\Gamma_n=\left\{r_q\cdot n^{\frac{1}{2}}\sum_{\mu}k_{\mu}b_{\mu}:\ k\in\mathbb{Z}^3\right\},\qquad\operatorname{covol}\left(\pi_h\Gamma_n\right)=r_q^3\cdot n^{\frac{3}{2}}.
\]
The group $\Gamma_n$ is discrete, as it lies in $r_q\cdot C^1(n;\mathbb{Z})$, and it spans $V_{\mathrm{g}}\oplus V_{\mathrm{h}}$, as $\Gamma_n\cap V_{\mathrm{g}}$ spans $V_{\mathrm{g}}$ and $\pi_h\Gamma_n$ spans $V_{\mathrm{h}}$. So $\Gamma_n$ is a lattice of full rank in $V_{\mathrm{g}}\oplus V_{\mathrm{h}}$; and since $\Gamma_n\cap V_{\mathrm{g}}$ has full rank in $V_{\mathrm{g}}$, the covolume of $\Gamma_n$ is the product of the covolumes of $\Gamma_n\cap V_{\mathrm{g}}$ in $V_{\mathrm{g}}$ and of $\pi_h\Gamma_n$ in $V_{\mathrm{h}}$:
\[
\operatorname{covol}(\Gamma_n)=\operatorname{covol}\left(\Gamma_n\cap V_{\mathrm{g}}\right)\cdot\operatorname{covol}\left(\pi_h\Gamma_n\right)=r_q^{n^3+2}\cdot\exp\left(\frac{1}{2}\mathrm{Ld}(n)\right).
\]

As $d$ vanishes on $V_{\mathrm{g}}\oplus V_{\mathrm{h}}$, the integrand of $Z^{\mathrm{ph}}_n$ depends on $A$ only through the component of $A$ in $V_{\mathrm{t}}$, a space of dimension $3n^3-\left(n^3-1\right)-3=2n^3-2$. Fubini's theorem on $C^1(n)/\Gamma_n=\left(\left(V_{\mathrm{g}}\oplus V_{\mathrm{h}}\right)/\Gamma_n\right)\times V_{\mathrm{t}}$ and the Gaussian integral on $V_{\mathrm{t}}$ give
\[
Z^{\mathrm{ph}}_n=\operatorname{covol}(\Gamma_n)\cdot\int_{V_{\mathrm{t}}}\mathrm{d}\sigma(A)\,\exp\left(-\frac{1}{2}\left\langle A,d^{*}dA\right\rangle\right)=\operatorname{covol}(\Gamma_n)\cdot(2\pi)^{n^3-1}\cdot\det\left(d^{*}d|_{V_{\mathrm{t}}}\right)^{-\frac{1}{2}}.
\]
The operator $dd^{*}+d^{*}d$ acts on each of the three components of a one-cochain as the scalar Laplacian; hence the product of its non-zero eigenvalues on $C^1(n)$ is $\exp\left(3\mathrm{Ld}(n)\right)$. It leaves the three summands invariant: it vanishes on $V_{\mathrm{h}}$; it coincides with $dd^{*}$ on $V_{\mathrm{g}}$, where $d$ vanishes; and it coincides with $d^{*}d$ on $V_{\mathrm{t}}$, where $d^{*}$ vanishes, $V_{\mathrm{t}}$ being the range of $d^{*}$ on the two-cochains. The operator $dd^{*}$ vanishes on the orthogonal complement of $V_{\mathrm{g}}$, which is the range of $d$, and its non-zero eigenvalues are those of $d^{*}d$ on $C^0(n)$; so the product of the eigenvalues of $dd^{*}$ on $V_{\mathrm{g}}$ is $\exp\left(\mathrm{Ld}(n)\right)$, and
\[
\det\left(d^{*}d|_{V_{\mathrm{t}}}\right)=\exp\left(2\mathrm{Ld}(n)\right).
\]
Inserting the values of $\operatorname{covol}(\Gamma_n)$ and of $\det\left(d^{*}d|_{V_{\mathrm{t}}}\right)$ into the formula for $Z^{\mathrm{ph}}_n$, and using $r_q=2\pi/q$, one obtains (a).

(b) By Lemma~\ref{lem:flux-sectors}(b), applied to the function $1$,
\[
\mathbf{Z}_n=\int_{C^1(n)/\Gamma_n}\mathrm{d}A\,\exp\left(-\frac{1}{2}\|dA\|^2\right)\cdot\mathrm{Dt}(qA).
\]
By Theorem~\ref{thm:free-coefficient-1}(d1), this integrand is invariant under the translations by the elements of $V_{\mathrm{g}}$ and of $\Gamma_n$; so it depends on the component of $A$ in $V_{\mathrm{g}}\oplus V_{\mathrm{h}}$ only through the class modulo $\pi_h\Gamma_n$ of its component in $V_{\mathrm{h}}$. In the basis $b_1,b_2,b_3$ one has $H_{\theta}=q^{-1}n^{\frac{1}{2}}\sum_{\mu}\theta_{\mu}b_{\mu}$. Hence $\theta\mapsto H_{\theta}$ maps $[-\pi,\pi)^3$ bijectively onto a fundamental domain of $\pi_h\Gamma_n$ in $V_{\mathrm{h}}$, and under this map the Lebesgue measure of $V_{\mathrm{h}}$ corresponds to $q^{-3}n^{\frac{3}{2}}\,\mathrm{d}\theta=(2\pi)^{-3}\cdot\operatorname{covol}\left(\pi_h\Gamma_n\right)\mathrm{d}\theta$. Fubini's theorem on the quotient by $\Gamma_n$ therefore gives
\[
\mathbf{Z}_n=\frac{\operatorname{covol}(\Gamma_n)}{(2\pi)^3}\int_{[-\pi,\pi]^3}\mathrm{d}\theta\int_{V_{\mathrm{t}}}\mathrm{d}\sigma(A)\,\exp\left(-\frac{1}{2}\|dA\|^2\right)\cdot\mathrm{Dt}\left(qA+qH_{\theta}\right).
\]
If the factor $\mathrm{Dt}\left(qA+qH_{\theta}\right)$ is deleted, the inner integral no longer depends on $\theta$, and the right side becomes $\operatorname{covol}(\Gamma_n)\cdot\int_{V_{\mathrm{t}}}\mathrm{d}\sigma(A)\,\exp\left(-\frac{1}{2}\|dA\|^2\right)$, which is the value of $Z^{\mathrm{ph}}_n$ found in the proof of (a). With respect to $\sigma$, the measure $\tilde{\gamma}_n$ has the density $\exp\left(-\frac{1}{2}\|dA\|^2\right)$ divided by $\int_{V_{\mathrm{t}}}\mathrm{d}\sigma(A)\,\exp\left(-\frac{1}{2}\|dA\|^2\right)$; moreover $\mathrm{Dt}\left(qA+qH_{\theta}\right)=\det\left(D_{A+H_{\theta}}+\boldsymbol{\mu}\right)$. Dividing the right side by this value of $Z^{\mathrm{ph}}_n$ therefore gives $\rho_n$, that is, $\mathbf{Z}_n=Z^{\mathrm{ph}}_n\cdot\rho_n$. Since $Z_n>0$ and $\mathbf{Z}_n>0$ (Section~\ref{sec:wilson-determinant-lattice}), the number $P_n(0)=\mathbf{Z}_n/Z_n$ is positive, and $Z_n=Z^{\mathrm{ph}}_n\cdot\rho_n/P_n(0)$. The inequalities $1-\varsigma_n\le P_n(0)\le1$ are Corollary~\ref{cor:partition-functions-1}(d), in whose notation $\exp\left(-\frac{1}{2}\|h_w\|^2\right)=\mathfrak{q}_w(n)$ (Section~\ref{sec:wilson-determinant-lattice}).

(c) On the cubic torus, $|n|=n^3$ and $n_{\iota}=n$ in the definition of $\mathfrak{q}_w(n)$ in Section~\ref{sec:wilson-determinant-lattice}; with $r_q=2\pi/q$ this gives
\[
\mathfrak{q}_w(n)=\exp\left(-\frac{2\pi^2}{q^2n}\cdot|w|_2^2\right)\qquad\left(w\in\mathbb{Z}^3\right).
\]
Put $x:=\exp\left(-\frac{2\pi^2}{q^2n}\right)$; then $0<x\le\exp\left(-2\pi^2\right)$, as $q^2n\le1$. Since $x^{i^2}\le x^i$ for every integer $i\ge1$, one has $\sum_{i\in\mathbb{Z}}x^{i^2}\le1+\frac{2x}{1-x}\le1+2.001x$, and therefore
\[
\varsigma_n=\left(\sum_{i\in\mathbb{Z}}x^{i^2}\right)^3-1\le(1+2.001x)^3-1\le7x.
\]
Finally $7x\le7\exp\left(-2\pi^2\right)<\frac{1}{8}$.

(d) By (b), applied on the cubic tori $n$ and $jn$, the numbers $\rho_n$, $\rho_{jn}$, $P_n(0)$ and $P_{jn}(0)$ are positive, and
\[
\log Z_n=\log Z^{\mathrm{ph}}_n+\log\rho_n-\log P_n(0),\qquad\log Z_{jn}=\log Z^{\mathrm{ph}}_{jn}+\log\rho_{jn}-\log P_{jn}(0).
\]
Since $j^3\left(n^3+2\right)-\left((jn)^3+2\right)=2\left(j^3-1\right)$ and $j^3\left(n^3-1\right)-\left((jn)^3-1\right)=-\left(j^3-1\right)$, part (a) gives
\[
j^3\log Z^{\mathrm{ph}}_n-\log Z^{\mathrm{ph}}_{jn}=2\left(j^3-1\right)\log\frac{2\pi}{q}-\left(j^3-1\right)\log(2\pi)+\frac{1}{2}\left(\mathrm{Ld}(jn)-j^3\cdot\mathrm{Ld}(n)\right)=D^{\mathrm{ph}}_j,
\]
because $2\log\frac{2\pi}{q}-\log(2\pi)=\log\frac{2\pi}{q^2}$. Together with the preceding display this yields $D_j=D^{\mathrm{ph}}_j+D^{\mathrm{mat}}_j+D^{\mathrm{fl}}_j$. By Lemma~\ref{lem:free-coefficient-3},
\[
\frac{1}{2}\left(\mathrm{Ld}(jn)-j^3\cdot\mathrm{Ld}(n)\right)\le-\left(j^3-1\right)\log n+K_j,
\]
which gives the bound for $D^{\mathrm{ph}}_j$. If $\varsigma_n<1$, then $P_n(0)\ge1-\varsigma_n>0$ and $P_{jn}(0)\le1$ by (b); hence $D^{\mathrm{fl}}_j\le-j^3\log\left(1-\varsigma_n\right)$.
\end{proof}

\subsubsection{Proof of the cumulant bounds for the probe field}\label{sec:proof-cumulant-bounds}

\begin{proof}[Proof of Theorem~\ref{thm:free-coefficient-2}]
Let $N$, $T^{*}$, $e$, the probe $\beta$ with its placement, and $c_2,c_4$ be as in Theorem~\ref{thm:free-coefficient-2}; thus $q=e\cdot\eta^{\frac{1}{2}}$, the lattice mass is $\boldsymbol{\mu}_N>-\frac{1}{2}$, and $0<e\le\frac{1}{3}$. The results of Section~\ref{sec:trace-formula-chessboard} are applied with the data of Remark~\ref{rem:translated-cells-aligned}: the cell width is $\boldsymbol{\lambda}=m_{\mathrm{p}}+\frac{1}{2}$, where $m_{\mathrm{p}}\ge7$ because $\beta$ fits $(N,T^{*})$ (Section~\ref{sec:probes-discretisation}), and $\mathsf{n}=2L$, $j=L$. Then $n^{*}=n_{*}$, the torus $1$ has $N^{*}=8L^3$ cells, the tori $k\in\{1\}\cup\{L,L+1,\dots\}$ and their cells are those of Section~\ref{sec:reference-tori-cells}, and $\mathfrak{A}\subset\mathfrak{A}^{+}$ (Section~\ref{sec:cells-reflections-mirror}). Since $n_{*}=R^{*}/\eta$, one has $q^2n_{*}=e^2\cdot R^{*}$; so the number $\boldsymbol{\alpha}$ of Section~\ref{sec:transfer-larger-cubic} is
\[
\boldsymbol{\alpha}=\frac{2\pi^2}{q^2n_{*}}=\frac{2\pi^2}{e^2\cdot R^{*}}.
\]
Taken at $j=L$, the reference numbers of Section~\ref{sec:transfer-larger-cubic} are those of Section~\ref{sec:numbers-attached-reference}: since $n^{*}=n_{*}$, the numbers $Z_1=Z_{n^{*}}$ and $Z_j=Z_{L\cdot n^{*}}$ of Section~\ref{sec:transfer-larger-cubic} are the numbers $Z_1=Z_{n_{*}}$ and $Z_L=Z_{L\cdot n_{*}}$ of Section~\ref{sec:numbers-attached-reference}; hence $D_j=D_L$ and $\boldsymbol{\delta}=D_j/\left((j-1)^3N^{*}\right)=\delta$.

\emph{Step 1: the bounds \eqref{eq:free-coefficient-9}.} By Theorem~\ref{thm:transfer-all-flux}, $D_L\ge0$, that is, $\delta\ge0$. Lemma~\ref{lem:partition-functions-cubic}(d), applied on the cubic torus $n_{*}$ with $j=L$ and with the lattice mass $\boldsymbol{\mu}_N$, gives
\[
D_L=D^{\mathrm{ph}}_L+D^{\mathrm{mat}}+D^{\mathrm{fl}}_L,\qquad D^{\mathrm{ph}}_L\le\left(L^3-1\right)\log\frac{2\pi}{e^2\cdot R^{*}}+K_L,
\]
since the numbers $D_L$ and $D^{\mathrm{mat}}_L$ of that lemma are the numbers $D_L$ and $D^{\mathrm{mat}}$ of Section~\ref{sec:numbers-attached-reference}, and $q^2n_{*}=e^2\cdot R^{*}$. As $e\le\frac{1}{3}$, one has $\log\left(e^{-1}\right)\ge\log3$, and \eqref{eq:free-coefficient-6} gives
\[
e^2\cdot R^{*}\le\frac{\pi^2}{6N^{*}\cdot\log3}<1,\qquad\frac{2\pi^2}{e^2\cdot R^{*}}\ge12N^{*}\cdot\log\left(e^{-1}\right);
\]
the strict inequality holds because $N^{*}=8L^3\ge8$. Hence Lemma~\ref{lem:partition-functions-cubic}(c) applies on the torus $n_{*}$, and it gives $\varsigma_{n_{*}}\le7x\le\frac{1}{8}$, where
\[
x:=\exp\left(-\frac{2\pi^2}{e^2\cdot R^{*}}\right)\le e^{12N^{*}}.
\]
For $0\le y\le\frac{1}{8}$ one has $-\log(1-y)\le\frac{y}{1-y}\le\frac{8}{7}y$; thus $-\log(1-7x)\le8x$. By Lemma~\ref{lem:partition-functions-cubic}(d), whose last bound applies since $\varsigma_{n_{*}}<1$,
\[
D^{\mathrm{fl}}_L\le-L^3\log\left(1-\varsigma_{n_{*}}\right)\le-L^3\log(1-7x)\le8L^3\cdot e^{12N^{*}}=N^{*}\cdot e^{12N^{*}}\le1,
\]
the last inequality because $e\le\frac{1}{3}$. Consequently
\[
D_L\le\left(L^3-1\right)\log\frac{2\pi}{e^2\cdot R^{*}}+K_L+D^{\mathrm{mat}}+1,
\]
and division by $(L-1)^3N^{*}$, together with $a_L=\left(L^3-1\right)/(L-1)^3$, gives \eqref{eq:free-coefficient-9}.

\emph{Step 2: reduction to \eqref{eq:free-coefficient-10}.} Let $\mathrm{nat}\in\{\mathrm{af},\mathrm{av}\}$, let $k\in\{1\}\cup\{L,L+1,\dots\}$, and let $\kappa_2,\kappa_4$ be the cumulants of Theorem~\ref{thm:free-coefficient-2}. By Theorem~\ref{thm:free-coefficient-1}(d5), $E^{\mathrm{nat}}_k$ is invariant under $F^u\mapsto-F^u$. Since $X$ is linear in $F^u$, it follows that $E^{\mathrm{nat}}_kX=E^{\mathrm{nat}}_kX^3=0$; hence $\kappa_2=E^{\mathrm{nat}}_kX^2$ and $\kappa_4=E^{\mathrm{nat}}_kX^4-3\left(E^{\mathrm{nat}}_kX^2\right)^2$. Put $\check{u}:=E^{\mathrm{nat}}_kH_2(X)=\kappa_2-\boldsymbol{\sigma}^2$. As $X^4=H_4(X)+6\boldsymbol{\sigma}^2\cdot X^2-3\boldsymbol{\sigma}^4$, one has
\begin{align*}
E^{\mathrm{nat}}_kX^4&=E^{\mathrm{nat}}_kH_4(X)+6\boldsymbol{\sigma}^2\cdot\left(\check{u}+\boldsymbol{\sigma}^2\right)-3\boldsymbol{\sigma}^4,\\
\kappa_4&=E^{\mathrm{nat}}_kH_4(X)+6\boldsymbol{\sigma}^2\cdot\check{u}+3\boldsymbol{\sigma}^4-3\left(\check{u}+\boldsymbol{\sigma}^2\right)^2=E^{\mathrm{nat}}_kH_4(X)-3\check{u}^2.
\end{align*}
Suppose that $\left|e^{-m}\cdot E^{\mathrm{nat}}_kH_m(X)-c_m\right|\le\epsilon_m$ for $m=2,4$, with numbers $\epsilon_2,\epsilon_4\ge0$. The left side of \eqref{eq:free-coefficient-7} is $\left|e^{-2}\cdot\check{u}-c_2\right|$, so \eqref{eq:free-coefficient-7} holds with $\epsilon_2$ as its right side. Since $\left|y^2-y'^2\right|\le|y-y'|\cdot\left(2|y'|+|y-y'|\right)$ for real $y,y'$, the identity for $\kappa_4$ gives
\[
\left|e^{-4}\cdot\kappa_4-\left(c_4-3c_2^2\right)\right|\le\left|e^{-4}\cdot E^{\mathrm{nat}}_kH_4(X)-c_4\right|+3\left|\left(e^{-2}\cdot\check{u}\right)^2-c_2^2\right|\le\epsilon_4+3\epsilon_2\cdot\left(2|c_2|+\epsilon_2\right).
\]
The map $\mathsf{g}_0$ is the identity; so $\hat{\mathfrak{a}}_0=\hat{\mathfrak{a}}$, and $X_c=X$ for the cell $c=0$. By \eqref{eq:free-coefficient-10} for $S=\{0\}$, the numbers $\epsilon_m:=\boldsymbol{\Gamma}^{\mathrm{nat}}\cdot\boldsymbol{\Theta}^{\mathrm{nat}}_m$, $m=2,4$, satisfy the supposition, and \eqref{eq:free-coefficient-7}, \eqref{eq:free-coefficient-8} follow. It remains to prove \eqref{eq:free-coefficient-10}.

\emph{Step 3: the numbers $\Theta^{(0)}_m$, and moments.} For $m\in\{2,4\}$ put $O_m:=e^{-m}\cdot H_m(X)-c_m$ on every torus $k$. By Theorem~\ref{thm:free-coefficient-1}(b2), $d\hat{\mathfrak{a}}$ is supported in $P^o(\mathfrak{c}_0)$; so $O_m$ is a function of $\left(F^u_p\right)_{p\in P^o(\mathfrak{c}_0)}$, and it lies in $\mathfrak{F}(kn_{*})$, as $|X|\le\boldsymbol{\sigma}\cdot\|F^u\|$. Thus $O_m\in\mathfrak{A}$, and $O_m$ is real valued. Moreover $d\hat{\mathfrak{a}}$ takes the same values on $P^o(\mathfrak{c}_0)$ on all the tori $k$, so that $X$ and $O_m$ are the same functions of $\left(F^u_p\right)_{p\in P^o(\mathfrak{c}_0)}$ for all $k$. Indeed, every bond in the support of $\hat{\mathfrak{a}}$ has both end points in $\mathfrak{c}_0$ (Theorem~\ref{thm:free-coefficient-1}(b2)); and if $y$ is a site of $\mathfrak{c}_0$, written in lattice coordinates, then by Theorem~\ref{thm:free-coefficient-1}(b1), applied with $\tau_v\beta$ in place of $\beta$, the value of $\hat{\mathfrak{a}}=\eta^{\frac{1}{2}}\cdot a_{\eta}$ on the bond from $y$ to $y+e_{\nu}$ is $\eta^{\frac{1}{2}}\cdot\check{a}_{\eta,\nu}\left(\eta\cdot y+\frac{\eta}{2}e_{\nu}\right)$, where $\check{a}_{\eta}$ is formed on $\mathbb{R}^3$ from $\tau_v\beta$ and does not involve the torus.

As in Definition~\ref{def:decomposition-torus-cells}, these functions are read on every torus $n$ that carries the cell structure of width $\boldsymbol{\lambda}$. On such a torus the sites $\{0,\dots,m_{\mathrm{p}}\}^3$ of $\mathfrak{c}_0$, and with them the set $P^o(\mathfrak{c}_0)$, have the same description as on the torus $1$; and $P^o(\mathfrak{c}_0)\subset P^{+}(\mathfrak{c}_0)$, because among the planes of the cell structure only the three planes $\left\{x_{\mu}=0\right\}$ contain sites of $\mathfrak{c}_0$. Let $g_n\in C^2(n)$ be the two-cochain that is supported in $P^o(\mathfrak{c}_0)$ and takes there the values of $d\hat{\mathfrak{a}}$. Then, on $n$,
\[
X=\left\langle g_n,F^u\right\rangle,\qquad O_m=e^{-m}\cdot H_m(X)-c_m,\qquad\left\|g_n\right\|^2=\left\|d\hat{\mathfrak{a}}\right\|^2=\boldsymbol{\sigma}^2;
\]
moreover $O_m\in\mathfrak{A}^{+}$, as $|X|\le\boldsymbol{\sigma}\cdot\|F^u\|$, and $g_n=d\hat{\mathfrak{a}}$ when $n$ is the torus $k$. For a cell $c$ of $n$ put $X_c:=\left\langle\left(\mathsf{g}_c\right)_{*}g_n,F^u\right\rangle$. The pull-back $\mathsf{g}_c^{*}$ is an orthogonal map of $C^2(n)$ with inverse $\left(\mathsf{g}_c\right)_{*}$; therefore $\left\|\left(\mathsf{g}_c\right)_{*}g_n\right\|^2=\boldsymbol{\sigma}^2$, and, for every cell $c$ of $n$,
\[
O_m^{[c]}=e^{-m}\cdot H_m(X_c)-c_m,\qquad\text{because}\qquad\left\langle g_n,\mathsf{g}_c^{*}F^u\right\rangle=\left\langle\left(\mathsf{g}_c\right)_{*}g_n,F^u\right\rangle=X_c.
\]
On the torus $k$ the maps $\mathsf{g}_c$ are those of Section~\ref{sec:reference-tori-cells} (Section~\ref{sec:cells-reflections-mirror}), and $\left(\mathsf{g}_c\right)_{*}d\hat{\mathfrak{a}}=d\hat{\mathfrak{a}}_c$, as $\left(\mathsf{g}_c\right)_{*}$ commutes with $d$; so there $X_c=\left\langle d\hat{\mathfrak{a}}_c,F^u\right\rangle$ is the variable of Section~\ref{sec:probes-discretisation}.

Let $z=(z_c)_c\in\mathbb{C}^{N^{*}}$ carry one variable for each cell $c$ of the torus $1$, and put $a:=\sum_cz_c\cdot\hat{\mathfrak{a}}_c\in C^1(n_{*})\otimes\mathbb{C}$. By Theorem~\ref{thm:free-coefficient-1}(b2) the supports of the cochains $d\hat{\mathfrak{a}}_c$ are pairwise disjoint and $\|d\hat{\mathfrak{a}}_c\|^2=\boldsymbol{\sigma}^2$; hence, for the bilinear extension of the pairing,
\[
\langle da,da\rangle=\boldsymbol{\sigma}^2\sum_cz_c^2,\qquad\langle da,F^u\rangle=\sum_cz_c\cdot X_c.
\]
By \eqref{it:free-coefficient-Q1} and Lemma~\ref{lem:sector-zero-flux}, $E^{\mathrm{av}}_1$ is the functional $\overline{Z}_{n_{*}}(\cdot)/\overline{Z}_{n_{*}}(1)$, where $\overline{Z}_{n_{*}}$ is the functional \eqref{eq:all-flux-6} on the torus $1$, formed with $q$ and with the lattice mass $\boldsymbol{\mu}_N$. Let $\overline{Z}^{\mathrm{ext}}_{n_{*}}$ be the entire function of Section~\ref{sec:shifting-connection} on the torus $1$, formed with the same data:
\[
\overline{Z}^{\mathrm{ext}}_{n_{*}}(b)=\frac{1}{(2\pi)^3}\int_{[-\pi,\pi]^3}\mathrm{d}\theta\int_{V_{\mathrm{t}}}\mathrm{d}\sigma(A)\,\exp\left(-\frac{1}{2}\|dA\|^2\right)\cdot\mathrm{Dt}\left(qA+qH_{\theta}+b\right)\qquad\left(b\in C^1(n_{*})\otimes\mathbb{C}\right),
\]
$\sigma$ being the Lebesgue measure of $V_{\mathrm{t}}$. The third identity of Lemma~\ref{lem:shift-identities}, applied on the torus $1$ to the cochain $a$, reads
\[
\overline{Z}_{n_{*}}\left(\exp\langle da,F^u\rangle\right)=\exp\left(\frac{1}{2}\langle da,da\rangle\right)\cdot\overline{Z}^{\mathrm{ext}}_{n_{*}}(qa).
\]
With $t:=e\cdot z$ one has $qa=qe^{-1}\sum_ct_c\cdot\hat{\mathfrak{a}}_c$ and $\mathrm{Dt}\left(qA+qH_{\theta}+qa\right)=\det\left(D_{A+H_{\theta}+e^{-1}\sum_ct_c\cdot\hat{\mathfrak{a}}_c}+\boldsymbol{\mu}_N\right)$. Since $\tilde{\gamma}_{n_{*}}$ is the probability measure on $V_{\mathrm{t}}$ whose density with respect to $\sigma$ is proportional to $\exp\left(-\frac{1}{2}\|dA\|^2\right)$, it follows that $\overline{Z}^{\mathrm{ext}}_{n_{*}}(qa)$ is the product of $\mathbf{Z}^{\mathrm{ext}}(e\cdot z)$ (Section~\ref{sec:numbers-attached-reference}) with the positive number $\int_{V_{\mathrm{t}}}\mathrm{d}\sigma(A)\,\exp\left(-\frac{1}{2}\|dA\|^2\right)$, which does not depend on $z$. Dividing the identity by its value at $z=0$, one obtains
\[
E^{\mathrm{av}}_1\exp\sum_c\left(z_c\cdot X_c-\frac{1}{2}z_c^2\cdot\boldsymbol{\sigma}^2\right)=\mathring{\Phi}(e\cdot z)\qquad\left(z\in\mathbb{C}^{N^{*}}\right).
\]
Both sides are entire functions of $z$, and the Taylor coefficients of the left side at $z=0$ may be computed under $E^{\mathrm{av}}_1$, by dominated convergence: $|X_c|\le\boldsymbol{\sigma}\cdot\|F^u\|$, and $E^{\mathrm{av}}_1$ is finite on $\mathfrak{F}(n_{*})$ (Section~\ref{sec:wilson-determinant-lattice}). Let $H_i$, $i\ge0$, be the Hermite polynomials of variance $\boldsymbol{\sigma}^2$, given by
\[
\exp\left(zy-\frac{1}{2}z^2\cdot\boldsymbol{\sigma}^2\right)=\sum_{i\ge0}H_i(y)\cdot\frac{z^i}{i!}\qquad(y\in\mathbb{R},\ z\in\mathbb{C});
\]
then $H_0=1$, and $H_2$, $H_4$ are the polynomials of Section~\ref{sec:probes-discretisation}. Comparing the coefficients of the monomial $\prod_{c\in S'}z_c^m$ on the two sides (each of them carries the factor $(m!)^{-\#|S'|}$) gives, for every set $S'$ of cells of the torus $1$,
\[
E^{\mathrm{av}}_1\prod_{c\in S'}H_m(X_c)=e^{m\cdot\#|S'|}\cdot\left(\prod_{c\in S'}\partial_{t_c}^m\mathring{\Phi}\right)(0).
\]
Expand $\prod_c\left(e^{-m}\cdot H_m(X_c)-c_m\right)$, the product over all $N^{*}$ cells of the torus $1$, into a sum over the sets $S'$ of cells, and do the same with the operator $\prod_c\left(\partial_{t_c}^m-c_m\right)$. The last display then yields
\begin{equation}\label{eq:free-coefficient-18}
E^{\mathrm{av}}_1\prod_cO_m^{[c]}=\left(\prod_c\left(\partial_{t_c}^m-c_m\right)\mathring{\Phi}\right)(0),\qquad\text{hence}\qquad\|O_m\|_{0*}=\left(E^{\mathrm{av}}_1O_m^{\sharp}\right)^{\frac{1}{N^{*}}}=\Theta^{(0)}_m;
\end{equation}
here $\|O_m\|_{0*}$ is the number of Section~\ref{sec:transfer-larger-cubic}, and the second identity holds because the left side of the first, which is $E^{\mathrm{av}}_1O_m^{\sharp}$, is non-negative by Proposition~\ref{prop:chessboard-estimate}.

\emph{Moments.} For a law $E$, a function $G$ of the configuration and an integer $i\ge1$ put $\|G\|_{L^i(E)}:=\left(E|G|^i\right)^{\frac{1}{i}}$. Let $n$ be a torus that carries the cell structure of width $\boldsymbol{\lambda}$, for instance the torus $k$ or one of the tori $n^{*}(j_0,j_1,j_2)$ of Section~\ref{sec:transfer-larger-cubic}, with $X$, $X_c$ and $O_m$ read on $n$ as at the beginning of this step, and let $c$ be a cell of $n$. Corollary~\ref{cor:partition-functions-1}(b), applied to the two-cochain $\left(\mathsf{g}_c\right)_{*}g_n$, whose squared norm is $\boldsymbol{\sigma}^2$, gives
\[
\|X_c\|_{L^i(E^{\mathrm{af}}_n)}\le\left(\mathbf{e}\cdot i\cdot\boldsymbol{\sigma}^2\right)^{\frac{1}{2}}\qquad\text{for every even integer }i\ge2.
\]
For $c=0$ this is a bound for $X$, as $\mathsf{g}_0$ is the identity. Let $\mathbf{M}\ge1$ be an integer, and take all norms with respect to $E^{\mathrm{af}}_n$. Since $\|X^4\|_{L^{2\mathbf{M}}}=\|X\|_{L^{8\mathbf{M}}}^4$ and $\|X^2\|_{L^{2\mathbf{M}}}=\|X\|_{L^{4\mathbf{M}}}^2$, Minkowski's inequality gives
\begin{align*}
\left\|H_4(X)\right\|_{L^{2\mathbf{M}}}&\le\left(8\mathbf{e}\cdot\mathbf{M}\cdot\boldsymbol{\sigma}^2\right)^2+6\boldsymbol{\sigma}^2\cdot\left(4\mathbf{e}\cdot\mathbf{M}\cdot\boldsymbol{\sigma}^2\right)+3\boldsymbol{\sigma}^4\le600\,\mathbf{M}^2\cdot\boldsymbol{\sigma}^4,\\
\left\|H_2(X)\right\|_{L^{2\mathbf{M}}}&\le4\mathbf{e}\cdot\mathbf{M}\cdot\boldsymbol{\sigma}^2+\boldsymbol{\sigma}^2\le12\,\mathbf{M}\cdot\boldsymbol{\sigma}^2,
\end{align*}
because $\mathbf{M}\ge1$, $64\mathbf{e}^2+24\mathbf{e}+3=541.13\ldots$ and $4\mathbf{e}+1=11.87\ldots$. By Minkowski's inequality once more,
\begin{equation}\label{eq:free-coefficient-19}
\|O_4\|_{L^{2\mathbf{M}}(E^{\mathrm{af}}_n)}\le\left(600\,\boldsymbol{\sigma}^4\cdot e^{-4}+|c_4|\right)\mathbf{M}^2,\qquad\|O_2\|_{L^{2\mathbf{M}}(E^{\mathrm{af}}_n)}\le\left(12\,\boldsymbol{\sigma}^2\cdot e^{-2}+|c_2|\right)\mathbf{M}.
\end{equation}

\emph{Step 4: the all-flux law.} Let $k\in\{1\}\cup\{L,L+1,\dots\}$, and let $S$ be a set of cells of the torus $k$. Theorem~\ref{thm:transfer-all-flux}, applied with $O_c:=O_m$ for $c\in S$, gives
\[
\left|E^{\mathrm{af}}_k\prod_{c\in S}O_m^{[c]}\right|\le\left(\exp(\delta)\cdot\left(E^{\mathrm{af}}_1O_m^{\sharp}\right)^{\frac{1}{N^{*}}}\right)^{\#|S|},
\]
where $\left(E^{\mathrm{af}}_1O_m^{\sharp}\right)^{\frac{1}{N^{*}}}$ is the number $\|O_m\|_{*}$ of Section~\ref{sec:transfer-larger-cubic}. Let $F^u=dA+r_q\cdot\mathfrak{z}$, with $d\mathfrak{z}=0$, be a curvature on the torus $1$, and let $\mu<\nu$. The plaquettes of the orientation $(\mu\nu)$ form $n_{*}$ parallel coordinate two-tori; the sum of $dA$ over each of them vanishes, and the sum of $\mathfrak{z}$ over each of them is $w_{\mu\nu}(\mathfrak{z})$. So the sum of $F^u$ over all plaquettes of the orientation $(\mu\nu)$ is $r_q\cdot n_{*}\cdot w_{\mu\nu}(\mathfrak{z})$, and the flux is a Borel function $w(F^u)$ of $F^u$; accordingly $\boldsymbol{\omega}^{(w)}_{n_{*}}(\Phi)=\boldsymbol{\omega}_{n_{*}}\left(\Phi\cdot\mathbf{1}\left\{w(F^u)=w\right\}\right)$ for $\Phi\in\mathfrak{F}(n_{*})$ and $w\in\mathbb{Z}^3$. Consequently
\[
E^{\mathrm{af}}_1O_m^{\sharp}=P_{n_{*}}(0)\cdot E^{\mathrm{av}}_1O_m^{\sharp}+E^{\mathrm{af}}_1\left[O_m^{\sharp}\cdot\mathbf{1}\left\{w(F^u)\ne0\right\}\right].
\]
The functionals $\boldsymbol{\omega}^{(w)}_{n_{*}}$ are positive (Theorem~\ref{thm:free-coefficient-1}(d2)); so $E^{\mathrm{af}}_1$ is a probability measure. By \eqref{eq:free-coefficient-18}, $E^{\mathrm{av}}_1O_m^{\sharp}=\left(\Theta^{(0)}_m\right)^{N^{*}}\ge0$, and $P_{n_{*}}(0)\le1$; by Lemma~\ref{lem:partition-functions-cubic}(b), $E^{\mathrm{af}}_1\mathbf{1}\left\{w(F^u)\ne0\right\}=1-P_{n_{*}}(0)\le\varsigma_{n_{*}}$. The Cauchy--Schwarz inequality therefore gives
\[
E^{\mathrm{af}}_1O_m^{\sharp}\le\left(\Theta^{(0)}_m\right)^{N^{*}}+\varsigma_{n_{*}}^{\frac{1}{2}}\cdot\left(E^{\mathrm{af}}_1\left|O_m^{\sharp}\right|^2\right)^{\frac{1}{2}}.
\]
H\"older's inequality for the product of the $N^{*}$ factors $O_m^{[c]}$, together with the invariance of $E^{\mathrm{af}}_1$ under the lattice isometries $\mathsf{g}_c$ (Theorem~\ref{thm:free-coefficient-1}(d5)), gives
\[
\left(E^{\mathrm{af}}_1\left|O_m^{\sharp}\right|^2\right)^{\frac{1}{2}}\le\|O_m\|_{L^{2N^{*}}(E^{\mathrm{af}}_1)}^{N^{*}}.
\]
By Minkowski's inequality and the bounds for $\left\|H_4(X)\right\|_{L^{2\mathbf{M}}}$ and $\left\|H_2(X)\right\|_{L^{2\mathbf{M}}}$ of Step~3, taken on the torus $1$ at $\mathbf{M}=N^{*}$,
\[
\|O_4\|_{L^{2N^{*}}(E^{\mathrm{af}}_1)}\le600\,(N^{*})^2\cdot\boldsymbol{\sigma}^4\cdot e^{-4}+|c_4|=\dot{M}_4,\qquad\|O_2\|_{L^{2N^{*}}(E^{\mathrm{af}}_1)}\le12N^{*}\cdot\boldsymbol{\sigma}^2\cdot e^{-2}+|c_2|=\dot{M}_2.
\]
As $E^{\mathrm{af}}_1O_m^{\sharp}\ge0$ (Proposition~\ref{prop:chessboard-estimate}) and $(y+y')^{\frac{1}{N^{*}}}\le y^{\frac{1}{N^{*}}}+y'^{\frac{1}{N^{*}}}$ for $y,y'\ge0$, it follows that
\[
\left(E^{\mathrm{af}}_1O_m^{\sharp}\right)^{\frac{1}{N^{*}}}\le\Theta^{(0)}_m+\varsigma_{n_{*}}^{\frac{1}{2N^{*}}}\cdot\dot{M}_m.
\]
By Step~1, $\varsigma_{n_{*}}\le7\exp\left(-\frac{2\pi^2}{e^2\cdot R^{*}}\right)$, and by \eqref{eq:free-coefficient-6}, $\frac{\pi^2}{N^{*}\cdot e^2\cdot R^{*}}\ge6\log\left(e^{-1}\right)$; as $N^{*}\ge8$, this gives
\[
\varsigma_{n_{*}}^{\frac{1}{2N^{*}}}\le7^{\frac{1}{2N^{*}}}\cdot\exp\left(-\frac{\pi^2}{N^{*}\cdot e^2\cdot R^{*}}\right)\le2e^6.
\]
Hence $\left(E^{\mathrm{af}}_1O_m^{\sharp}\right)^{\frac{1}{N^{*}}}\le\Theta^{(0)}_m+2e^6\cdot\dot{M}_m=\boldsymbol{\Theta}^{\mathrm{af}}_m$. Since $\boldsymbol{\Gamma}^{\mathrm{af}}=\exp(\delta)$ and $O_m^{[c]}=e^{-m}\cdot H_m(X_c)-c_m$, the first display of this step is now \eqref{eq:free-coefficient-10} for $\mathrm{nat}=\mathrm{af}$.

\emph{Step 5: the harmonic-averaged law.} Put $\vartheta:=\boldsymbol{\Theta}^{\mathrm{av}}_m=\max\left(\Theta^{(0)}_m,e^2\right)$; then $\vartheta>0$, and $\vartheta\ge\Theta^{(0)}_m=\|O_m\|_{0*}$ by \eqref{eq:free-coefficient-18}. Proposition~\ref{prop:transfer-harmonic-averaged} is applied with $\Phi:=O_m\in\mathfrak{A}^{+}$ and with this $\vartheta$. Its moment hypothesis holds by \eqref{eq:free-coefficient-19}, with
\[
(m_1,p)=\left(600\,\boldsymbol{\sigma}^4\cdot e^{-4}+|c_4|,\ 2\right)\quad\text{if }m=4,\qquad(m_1,p)=\left(12\,\boldsymbol{\sigma}^2\cdot e^{-2}+|c_2|,\ 1\right)\quad\text{if }m=2,
\]
and its hypothesis on $\boldsymbol{\alpha}$ holds because, by \eqref{eq:free-coefficient-6},
\[
\boldsymbol{\alpha}=\frac{2\pi^2}{e^2\cdot R^{*}}\ge12N^{*}\cdot\log\left(e^{-1}\right)\ge4.
\]
As $e\le1$, for both values of $m$ one has $m_1\cdot(N^{*})^p\cdot e^{-2}\le e^{-6}\cdot(N^{*})^2\cdot\left(600\,\boldsymbol{\sigma}^4+12\,\boldsymbol{\sigma}^2+|c_4|+|c_2|\right)$. Since moreover $\vartheta\ge e^2$, and $\log^{+}$ is non-decreasing with $\log^{+}(yy')\le\log^{+}y+\log^{+}y'$ for $y,y'>0$, the number $\boldsymbol{\lambda}_0$ of Proposition~\ref{prop:transfer-harmonic-averaged} satisfies
\[
\boldsymbol{\lambda}_0\le\log^{+}\left(m_1\cdot(N^{*})^p\cdot e^{-2}\right)\le6\log\left(e^{-1}\right)+Y.
\]
Hence, by \eqref{eq:free-coefficient-6} and as $3\le N^{*}$, the number $X_0$ of Proposition~\ref{prop:transfer-harmonic-averaged} satisfies
\[
X_0^2=\frac{2\left(N^{*}\cdot\boldsymbol{\lambda}_0+3\right)}{\boldsymbol{\alpha}}=\left(N^{*}\cdot\boldsymbol{\lambda}_0+3\right)\cdot\frac{e^2\cdot R^{*}}{\pi^2}\le\frac{N^{*}\cdot\boldsymbol{\lambda}_0+3}{6N^{*}\cdot\log\left(e^{-1}\right)}\le\check{X},
\]
and the number $\tilde{b}$ of Proposition~\ref{prop:transfer-harmonic-averaged} satisfies
\[
\tilde{b}=\frac{2p\,N^{*}}{\boldsymbol{\alpha}}=p\,N^{*}\cdot\frac{e^2\cdot R^{*}}{\pi^2}\le\frac{2}{6\log3}<0.3035,\qquad\frac{3+8\tilde{b}}{\mathbf{e}}<\frac{5.428}{\mathbf{e}}<2.
\]
The bound of Proposition~\ref{prop:transfer-harmonic-averaged} for several cells, with all $\Phi_i=O_m$, now gives, for every $k\in\{1\}\cup\{L,L+1,\dots\}$ and every set $S$ of cells of the torus $k$,
\[
\left|E^{\mathrm{av}}_k\prod_{c\in S}O_m^{[c]}\right|\le\left(\vartheta\cdot\exp\left(\frac{1}{N^{*}}\left(3\log\left(4+16\check{X}\right)+2+\frac{D_L}{(L-1)^3}+4\right)\right)\right)^{\#|S|}=\left(\boldsymbol{\Gamma}^{\mathrm{av}}\cdot\boldsymbol{\Theta}^{\mathrm{av}}_m\right)^{\#|S|},
\]
since $D_L/\left((L-1)^3N^{*}\right)=\delta$. This is \eqref{eq:free-coefficient-10} for $\mathrm{nat}=\mathrm{av}$, which completes the proof.
\end{proof}

\subsubsection{Functions invariant under a measure-preserving action of $\mathbb{Z}^d$}\label{sec:functions-invariant-measure}

\begin{lemma}[Uniform bound on the invariant part]\label{lem:uniform-bound-invariant}
Let $(\boldsymbol{\Omega},\boldsymbol{\Sigma},\mu)$ be a probability space, let $d\ge1$ be an integer, and let $(T_y)_{y\in\mathbb{Z}^d}$ be an action of $\mathbb{Z}^d$ on $\boldsymbol{\Omega}$ by measurable maps that preserve $\mu$: the map $T_0$ is the identity of $\boldsymbol{\Omega}$, and $T_y\circ T_{y'}=T_{y+y'}$ and $\mu\left(T_y^{-1}(E)\right)=\mu(E)$ for all $y,y'\in\mathbb{Z}^d$ and all $E\in\boldsymbol{\Sigma}$. For a measurable function $F$ on $\boldsymbol{\Omega}$ and for $y\in\mathbb{Z}^d$ let $U(y)F:=F\circ T_{-y}$. Let $G\colon\boldsymbol{\Omega}\to\mathbb{R}$ be a measurable function such that $G\in L^J(\mu)$ for every integer $J\ge1$, and let $\mathsf{u}\ge0$ be a number such that
\[
\left|\int_{\boldsymbol{\Omega}}\mathrm{d}\mu\,\prod_{i=1}^JU(y_i)G\right|\le\mathsf{u}^J
\]
for every integer $J\ge1$ and all pairwise distinct points $y_1,\dots,y_J$ of $\mathbb{Z}^d$. Let $P_0$ be the orthogonal projection of $L^2(\mu)$ onto the closed subspace
\[
\mathcal{H}_0:=\left\{H\in L^2(\mu):\ U(y)H=H\ \text{for all}\ y\in\mathbb{Z}^d\right\}.
\]
Then $|P_0G|\le\mathsf{u}$ holds $\mu$-almost everywhere on $\boldsymbol{\Omega}$.
\end{lemma}

\begin{proof}
Write $\|\cdot\|$ and $\langle\cdot,\cdot\rangle$ for the norm and the scalar product of $L^2(\mu)$. Since $\mu$ is invariant under each $T_y$,
\[
\int_{\boldsymbol{\Omega}}\mathrm{d}\mu\,U(y)F=\int_{\boldsymbol{\Omega}}\mathrm{d}\mu\,F
\]
for every non-negative measurable function $F$ on $\boldsymbol{\Omega}$ and every $y\in\mathbb{Z}^d$; consequently $\|U(y)F\|_{L^p}=\|F\|_{L^p}$ for $1\le p<\infty$ and $F\in L^p(\mu)$. From $T_{-y'}\circ T_{-y}=T_{-y-y'}$ one obtains $U(y)U(y')=U(y+y')$ for all $y,y'\in\mathbb{Z}^d$, and $U(0)$ is the identity. Hence every $U(y)$ is an invertible isometry of $L^2(\mu)$ with $U(y)^{-1}=U(y)^{*}=U(-y)$, and $\mathcal{H}_0$, being the intersection of the kernels of the bounded operators $U(y)-1$, is closed. For an integer $n\ge1$ let $B_n:=\{0,1,\dots,n-1\}^d\subset\mathbb{Z}^d$, and for measurable functions $F$ on $\boldsymbol{\Omega}$ let
\[
A_nF:=n^{-d}\cdot\sum_{y\in B_n}U(y)F.
\]

\textit{Step (1): convergence of the averages.} For every $F\in L^2(\mu)$ one has $A_nF\to P_0F$ in $L^2(\mu)$ as $n\to\infty$; in particular $A_nG\to P_0G$, since $G\in L^2(\mu)$ by hypothesis. To prove this, let $\mathcal{D}\subset L^2(\mu)$ be the closed linear span of the vectors $U(y)H-H$ with $y\in\mathbb{Z}^d$ and $H\in L^2(\mu)$. For $F,H\in L^2(\mu)$ and $y\in\mathbb{Z}^d$ one has
\[
\left\langle F,\,U(y)H-H\right\rangle=\left\langle U(-y)F-F,\,H\right\rangle;
\]
therefore $F$ is orthogonal to $\mathcal{D}$ if and only if $U(-y)F=F$ for every $y\in\mathbb{Z}^d$, that is, if and only if $F\in\mathcal{H}_0$. Thus $\mathcal{D}^{\perp}=\mathcal{H}_0$, and since $\mathcal{D}$ is closed, $L^2(\mu)$ is the orthogonal sum of $\mathcal{H}_0$ and $\mathcal{D}$. Every $A_n$ restricts to the identity on $\mathcal{H}_0$.

Now let $H\in L^2(\mu)$ and $y=(\eta_1,\dots,\eta_d)\in\mathbb{Z}^d$, and put $|y|_1:=|\eta_1|+\dots+|\eta_d|$ and $y+B_n:=\{y+z:\ z\in B_n\}$. Since $U(z)U(y)=U(z+y)$,
\[
A_n\left(U(y)H-H\right)=n^{-d}\cdot\left(\sum_{z\in y+B_n}U(z)H-\sum_{z\in B_n}U(z)H\right).
\]
The two sets of indices differ in at most $2\,|y|_1\cdot n^{d-1}$ points, that is,
\[
\#\left|(y+B_n)\setminus B_n\right|+\#\left|B_n\setminus(y+B_n)\right|\le2\,|y|_1\cdot n^{d-1};
\]
indeed $(y+B_n)\cap B_n$ is a box with $\prod_{k=1}^d\max\{n-|\eta_k|,\,0\}$ points, so each of the two differences has
\[
n^d-\prod_{k=1}^d\max\{n-|\eta_k|,\,0\}\le|y|_1\cdot n^{d-1}
\]
elements. The summands indexed by $(y+B_n)\cap B_n$ cancel, and every $U(z)$ is isometric; it follows that
\[
\left\|A_n\left(U(y)H-H\right)\right\|\le2\,|y|_1\cdot n^{-1}\cdot\|H\|,
\]
which tends to $0$ as $n\to\infty$. Moreover $\|A_n\|\le1$ for every $n$, by the triangle inequality. Because of this uniform bound, the vectors $F\in L^2(\mu)$ with $A_nF\to0$ form a closed subspace of $L^2(\mu)$; it contains every vector $U(y)H-H$, hence it contains $\mathcal{D}$. For $F\in L^2(\mu)$ one has $P_0F\in\mathcal{H}_0$ and $F-P_0F\in\mathcal{D}$, so that
\[
A_nF=P_0F+A_n\left(F-P_0F\right)\longrightarrow P_0F\qquad(n\to\infty).
\]

\textit{Step (2): even moments of the averages.} Let $J\ge2$ be an even integer. By H\"older's inequality, applied with $J$ exponents all equal to $J$, and by the invariance of $\mu$, every product $\prod_{i=1}^JU(y_i)G$ with $y_1,\dots,y_J\in\mathbb{Z}^d$ is integrable, and
\[
\left|\int_{\boldsymbol{\Omega}}\mathrm{d}\mu\,\prod_{i=1}^JU(y_i)G\right|\le\prod_{i=1}^J\|U(y_i)G\|_{L^J}=\|G\|_{L^J}^J.
\]
Expanding the $J$-th power of the sum that defines $A_nG$ and integrating term by term gives
\[
\int_{\boldsymbol{\Omega}}\mathrm{d}\mu\,(A_nG)^J=n^{-dJ}\cdot\sum_{y_1,\dots,y_J\in B_n}\ \int_{\boldsymbol{\Omega}}\mathrm{d}\mu\,\prod_{i=1}^JU(y_i)G.
\]
A tuple $(y_1,\dots,y_J)\in B_n^J$ with pairwise distinct entries contributes a summand of absolute value at most $\mathsf{u}^J$, by the hypothesis of the lemma, and the number of such tuples does not exceed $n^{dJ}$. Every other tuple satisfies $y_i=y_l$ for some pair $i<l$; for a given pair this holds for exactly $n^{d(J-1)}$ tuples, so that there are at most $\binom{J}{2}\cdot n^{d(J-1)}$ tuples of this kind, and each of their summands has absolute value at most $\|G\|_{L^J}^J$, by the estimate obtained at the beginning of this step from H\"older's inequality and the invariance of $\mu$. Consequently
\[
\int_{\boldsymbol{\Omega}}\mathrm{d}\mu\,(A_nG)^J\le\mathsf{u}^J+\binom{J}{2}\cdot n^{-d}\cdot\|G\|_{L^J}^J,
\qquad\text{and hence}\qquad
\limsup_{n\to\infty}\int_{\boldsymbol{\Omega}}\mathrm{d}\mu\,(A_nG)^J\le\mathsf{u}^J.
\]

\textit{Step (3): the bound on $P_0G$.} Convergence in $L^2(\mu)$ implies convergence $\mu$-almost everywhere, to the same limit, along a subsequence. Applied to Step (1), this gives integers $n_1<n_2<\cdots$ such that $A_{n_k}G\to P_0G$ $\mu$-almost everywhere as $k\to\infty$. Let $J\ge2$ be even. Since $A_{n_k}G$ is real-valued, the functions $(A_{n_k}G)^J=|A_{n_k}G|^J$ are non-negative, and they converge almost everywhere to $|P_0G|^J$. Fatou's lemma and Step (2) yield
\[
\int_{\boldsymbol{\Omega}}\mathrm{d}\mu\,|P_0G|^J\le\liminf_{k\to\infty}\int_{\boldsymbol{\Omega}}\mathrm{d}\mu\,(A_{n_k}G)^J\le\limsup_{n\to\infty}\int_{\boldsymbol{\Omega}}\mathrm{d}\mu\,(A_nG)^J\le\mathsf{u}^J,
\]
that is, $\|P_0G\|_{L^J}\le\mathsf{u}$ for every even integer $J\ge2$. It remains to let $J\to\infty$. For a number $t>\mathsf{u}$ let $E_t:=\{|P_0G|>t\}\subset\boldsymbol{\Omega}$. Then, for every even $J\ge2$,
\[
t^J\cdot\mu(E_t)\le\int_{\boldsymbol{\Omega}}\mathrm{d}\mu\,|P_0G|^J\le\mathsf{u}^J,
\qquad\text{so that}\qquad
\mu(E_t)\le\left(\frac{\mathsf{u}}{t}\right)^J;
\]
as $0\le\mathsf{u}<t$, the right-hand side tends to $0$ when $J\to\infty$ through the even integers, and $\mu(E_t)=0$. The set $\{|P_0G|>\mathsf{u}\}$ is the union of the null sets $E_t$ with $t=\mathsf{u}+m^{-1}$, $m=1,2,\dots$; hence $|P_0G|\le\mathsf{u}$ $\mu$-almost everywhere.
\end{proof}

\subsection{Proof of Theorem~\ref{thm:free-coefficient-3}}\label{sec:proof-theorem}

\begin{proof}
\emph{Part~\ref{it:free-coefficient-3-1}.} Let $e>0$ be such that $e_K=e\cdot h'^{\frac{1}{2}}\le e_1$, where $e_1$ is the number of Section~\ref{sec:constants-conditions-torus}.

\emph{Step~1.} Proposition~\ref{prop:consequences} gives $e\le\frac{1}{3}$ and $v_K\ge8$. It also shows that, for every $R\in\mathrm{Tor}(e)$ and every integer $N\ge N_{\mathrm{min}}(R,e)$, the conditions~\eqref{eq:renormalisation-group-16}, \eqref{eq:renormalisation-group-fl} and~\eqref{eq:renormalisation-group-T1}--\eqref{eq:renormalisation-group-T6} of Section~\ref{sec:conditions} hold for the torus $R$, the cutoff $N$ and the charge $e$.

\emph{Step~2.} Let $M=L^m$ be the integer of Section~\ref{sec:constants-exponents-scales}. By the definition of $r_K$ there is an integer $i\ge0$ with $r_K=L^i$, and stage~\ref{it:order-choice-6} gives $m\ge m_0+2$. As $h'=h\cdot L^{-m_0}$, the definitions of Section~\ref{sec:constants-exponents-scales} give
\[
R^{*}(e)=\frac{h'\cdot M\cdot r_K}{L}=h'\cdot L^{m-1+i}=h\cdot L^{m-m_0-1+i}.
\]
Consequently
\[
T^{*}(e)=\log_L\left(\frac{R^{*}(e)}{h}\right)=m-m_0-1+i
\]
is an integer, $T^{*}(e)\ge1$, and $R^{*}(e)=h\cdot L^{T^{*}(e)}$. The exponent $m-1+i$ is a non-negative integer, and $R^{*}(e)\le L\cdot R^{*}(e)=R_{\mathrm{top}}(e)=h'\cdot L^{m+i}$. Hence both $R^{*}(e)$ and $R_{\mathrm{top}}(e)$ belong to $\mathrm{Tor}(e)$.

\emph{Step~3.} Let $N\ge N_0(e)=N_{\mathrm{min}}\left(R^{*}(e),e\right)$ be an integer. Since $R^{*}(e)\in\mathrm{Tor}(e)$ by Step~2, the standing hypothesis~\ref{def:hypothesis-flow} of Section~\ref{sec:one-step-operators} holds for the torus $R=R^{*}(e)$, the cutoff $N$ and the charge $e$. Theorem~\ref{thm:agreement-trajectories} therefore applies and yields the following: the pair $P_N(e)$ is defined and real; it is controlled on $R^{*}(e)$ in the sense of Definition~\ref{def:fluctuation-field-controlled}; it depends on $N$, $e$, $\bar{m}$ and on the scheme described in Theorem~\ref{thm:agreement-trajectories}, but neither on the torus $R\in\mathrm{Tor}(e)$ nor on the twist $\theta$; and the masses $m_k$ of its trajectory on $R^{*}(e)$ satisfy
\[
\left|m_k\right|\le\frac{1}{8}\,w^{\mathrm{m}}_k\qquad\left(0\le k\le k^{*}\left(R^{*}(e)\right)\right).
\]
At the level $k=0$ the mass of the trajectory is the first entry of $P_N(e)$, which is $m_{0,N}(e)$ by the definition in Theorem~\ref{thm:free-coefficient-3}; and $w^{\mathrm{m}}_0=q^{\frac{3}{4}-8\varepsilon}$ is formed with the charge $q=e\cdot\eta^{\frac{1}{2}}$ of the level $0$ (Section~\ref{sec:running-quantities-tori}). As $\eta\le1$, Step~1 gives $q\le e\le\frac{1}{3}$; together with $\frac{3}{4}-8\varepsilon>0$ this yields $w^{\mathrm{m}}_0<1$, and therefore
\[
\left|m_{0,N}(e)\right|\le\frac{1}{8}\,w^{\mathrm{m}}_0<\frac{1}{8}.
\]
For the integers $N$ with $0\le N<N_0(e)$ one has $m_{0,N}(e)=0$ by definition. Thus, for every integer $N\ge0$, the number $m_{0,N}(e)$ is real and $\left|m_{0,N}(e)\right|<\frac{1}{8}$, whence
\[
\boldsymbol{\mu}_N=\bar{m}\cdot\eta+m_{0,N}(e)>-\frac{1}{8}>-\frac{1}{2},
\]
and $N\mapsto m_{0,N}(e)$ is a mass assignment in the sense of Section~\ref{sec:wilson-determinant-lattice}. These numbers depend on $N$, $e$, $\bar{m}$ and on the scheme of Theorem~\ref{thm:agreement-trajectories}, and not on a torus.

\emph{Step~4.} Take $T^{*}=T^{*}(e)$ in Section~\ref{sec:reference-tori-cells}; then $R^{*}=R^{*}(e)$ by Step~2. Since $r_K\ge v_K^r$ by the definition of $r_K$,
\[
l_c=\frac{R^{*}(e)}{2L}=\frac{h'\cdot M\cdot r_K}{2L^2}\ge\frac{h'\cdot M\cdot v_K^r}{2L^2}.
\]
The numbers $h'$, $M$, $L$ and $r\ge2$ do not depend on $e$, while $v_K=-\log\left(e\cdot h'^{\frac{1}{2}}\right)$ tends to infinity as $e\to0$. Hence $l_c\to\infty$ as $e\to0$.

\emph{Step~5.} Let $N\ge N_0(e)$ be an integer, take $T^{*}=T^{*}(e)$, and let the lattice mass be $\boldsymbol{\mu}_N=\bar{m}\cdot\eta+m_{0,N}(e)$. For this lattice mass the number $\rho_n$ of Section~\ref{sec:numbers-attached-reference} is the number $\rho_n\left(e\cdot\eta^{\frac{1}{2}},\bar{m}\cdot\eta+m_{0,N}\right)$ of Theorem~\ref{thm:partition-function-pair}, by the identification stated in Section~\ref{sec:numbers-attached-reference}; so the number $D^{\mathrm{mat}}$ of Section~\ref{sec:numbers-attached-reference} is the number $D^{\mathrm{mat}}$ of~\eqref{eq:renormalisation-flow-27}. As $N\ge N_{\mathrm{min}}\left(R^{*},e\right)$, the inequality~\eqref{eq:renormalisation-flow-27} gives
\[
D^{\mathrm{mat}}\le\left(L^3+1\right)\left(\varepsilon_5\left(\bar{m}\cdot R^{*}\right)+3\,e_K^{\frac{1}{4}-9\varepsilon}\right).
\]
The hypotheses of Proposition~\ref{prop:depth-stop-level} hold by Proposition~\ref{prop:consequences}, and $R^{*}\in\mathrm{Tor}(e)$ by Step~2; so part~(e) of Proposition~\ref{prop:depth-stop-level}, applied with $R=R^{*}$, gives
\[
\bar{m}\cdot R^{*}\ge\mathfrak{m}_K>L^{-m_0-1}.
\]
The function $\varepsilon_5$ is decreasing on $(0,\infty)$, because $1-\exp\left(-\frac{x}{4}\right)$ increases with $x$. Moreover $v_K=-\log e_K\ge-\log e_1$, and the inequalities~\eqref{eq:renormalisation-flow-e1}--\eqref{eq:renormalisation-flow-e8} hold for every $v\ge-\log e_1$ (Section~\ref{sec:constants-conditions-torus}); at $v=v_K$ the last inequality of~\eqref{eq:renormalisation-flow-e3} reads $e_K^{\frac{1}{4}-9\varepsilon}\le\frac{1}{8}$. Consequently
\[
D^{\mathrm{mat}}\le\left(L^3+1\right)\left(\varepsilon_5\left(L^{-m_0-1}\right)+\frac{3}{8}\right)\le\left(L^3+1\right)\left(\varepsilon_5\left(L^{-m_0-1}\right)+1\right)=C_{\mathrm{mat}}.
\]

\emph{Step~6.} Assume moreover $e\le e_G$. Both conditions in the definition of $e_G$ hold for every number of the interval $\left(0,e_G\right]$ (Section~\ref{sec:constants-exponents-scales}); in particular $e_K\log h'\le1$. Proposition~\ref{prop:consequences} then shows that the condition~\eqref{eq:renormalisation-group-17} of Section~\ref{sec:conditions} holds, that is, $6N^{*}\cdot e^2\cdot R^{*}(e)\cdot\log\left(e^{-1}\right)\le\pi^2$; this is the inequality~\eqref{eq:free-coefficient-6} with $R^{*}=R^{*}(e)$.

Steps~1 to~6 establish every assertion of part~\ref{it:free-coefficient-3-1}.

\emph{Part~\ref{it:free-coefficient-3-2}: the exponents.} By stage~\ref{it:order-choice-8}, $\sigma>0$. Hence
\[
\mathfrak{e}=\sigma+\frac{2a_L}{N^{*}}>\frac{2a_L}{N^{*}}>0 .
\]
As $\varepsilon>0$ and $s<1$, one has $\mathfrak{e}\le\frac{s}{4}\le\frac{1}{4}$; and as $0<\varepsilon\le\frac{1}{200}$, one has $0<\mathfrak{e}_2=6\varepsilon\le\frac{6}{200}<1$. Thus $0<\mathfrak{e}\le\frac{1}{4}$, $0\le\mathfrak{e}_2<1$ and $\sigma>0$.

\emph{The probe and the objects attached to it.} Let $\beta$ be a probe. Let $e_{\sharp}=e_{\sharp}(\beta)>0$ be the number of Lemma~\ref{lem:threshold-charge}, and let $C_{\beta}$ be the number of Section~\ref{sec:probes-mirror-copies}. Put
\[
e_1(\beta):=\min\left(e_G,e_{\sharp}\right),\qquad\bar{C}:=1+h'^{\frac{1}{2}}\cdot C_{\beta},\qquad\check{h}:=\max\left(1,h'\right).
\]
As $e_G>0$, one has $e_1(\beta)\in\left(0,e_G\right]$. Let $0<e\le e_1(\beta)$ and let $N\ge N_0(e)$ be an integer; the reference tori of Section~\ref{sec:reference-tori-cells} are taken with $T^{*}=T^{*}(e)$, so that $R^{*}=R^{*}(e)$, and the lattice mass is $\boldsymbol{\mu}_N=\bar{m}\cdot\eta+m_{0,N}(e)$. Since $e\le e_G$, one has $e_K\le e_1$ (Section~\ref{sec:constants-exponents-scales}), and part~\ref{it:free-coefficient-3-1} applies. The standing hypothesis~\ref{def:hypothesis-reference-torus} of Section~\ref{sec:probes-jets-reference}, written out in Section~\ref{sec:standing-hypothesis-scales}, holds: it is the hypothesis~\ref{def:hypothesis-flow} for the torus $R=R^{*}$ with the pair $P_N(e)$, and $N\ge N_{\mathrm{min}}\left(R^{*},e\right)$. As $e\le e_{\sharp}$, Lemma~\ref{lem:threshold-charge} yields three facts. First, $2\rho_{\beta}+3\eta\le l_c$, that is, $\beta$ fits $\left(N,T^{*}(e)\right)$. Second, the placement $v\in\eta\mathbb{Z}^3$ of Section~\ref{sec:probes-discretisation} is admitted in the sense of~\eqref{eq:four-jets-1}. Third, the conditions~\eqref{eq:four-jets-sm1}--\eqref{eq:four-jets-sm8} of Section~\ref{sec:smallness-conditions} hold for this placement.

With this placement and this lattice mass, the objects of Sections~\ref{sec:probes-discretisation}, \ref{sec:numbers-attached-reference} and~\ref{sec:free-fourth-order} coincide with the objects of the same names in Sections~\ref{sec:probes-mirror-copies} and~\ref{sec:averaged-law-partition}, by the identifications stated in Sections~\ref{sec:probes-discretisation}, \ref{sec:numbers-attached-reference} and~\ref{sec:free-fourth-order}: the cochains $\hat{\mathfrak{a}}$ and $\hat{\mathfrak{a}}_c$, the variables $X_c$ and the number $\boldsymbol{\sigma}^2$ are those of Section~\ref{sec:probes-mirror-copies}; the function $\mathring{\Phi}$ is that of Section~\ref{sec:averaged-law-partition}; and the free ratio $N_0(\cdot)$, formed on the torus $\Lambda_{\eta,R^{*}}$ with the cochain $a_{\eta}$ of $\tau_v\beta$, is the function $N_0(\cdot)$ of Section~\ref{sec:averaged-law-partition}. Let $\phi_0$ be the Taylor polynomial of degree four of $N_0$ at $0$, as in Lemma~\ref{lem:free-wilson-ratio}(e). Then
\[
c_m=N_0^{(m)}(0)=\phi_0^{(m)}(0)\qquad(m=2,4),
\]
and these two numbers are real, because the coefficients of $\phi_0$ are real (Lemma~\ref{lem:free-wilson-ratio}(e)). Hence the numbers $\Theta^{(0)}_2$ and $\Theta^{(0)}_4$ of Section~\ref{sec:numbers-attached-reference}, formed with $c_2$ and $c_4$, are the numbers $\Theta^{(0)}_2$ and $\Theta^{(0)}_4$ of Section~\ref{sec:numbers}.

\emph{Assertions (a) and (b).} Hypothesis~\ref{def:hypothesis-reference-torus} holds, the cochains $\hat{\mathfrak{a}}_c$ are those of Section~\ref{sec:probes-mirror-copies}, the conditions~\eqref{eq:four-jets-sm1}--\eqref{eq:four-jets-sm7} hold, and $N\ge N_{\mathrm{min}}\left(R^{*},e\right)$; these are the hypotheses of Proposition~\ref{prop:products-centred-derivatives}, namely those of Theorem~\ref{thm:factorised-representation} together with the lower bound for $N$. Corollary~\ref{cor:bounds-powers} assumes in addition the inequality $\boldsymbol{\lambda}_L\le\bar{\kappa}^s$, with the numbers $\boldsymbol{\lambda}_L$ and $\bar{\kappa}$ of Section~\ref{sec:constants-3}; this inequality is condition~\eqref{eq:four-jets-sm8}. That corollary therefore gives
\begin{gather*}
\Theta^{(0)}_4\le3\cdot10^5\cdot\bar{C}^4\cdot e_{*}^{\mathfrak{e}},\qquad\left(1+|c_2|\right)\cdot\Theta^{(0)}_2\le3\cdot10^4\cdot\bar{C}^4\cdot e_{*}^{\mathfrak{e}},\\
|c_2|+|c_4|^{\frac{1}{2}}\le8\,\bar{C}^2\cdot e_{*}^{-6\varepsilon},\qquad e\cdot\left(\frac{h'}{L^2}\right)^{\frac{1}{2}}\le e_{*}\le e\cdot h'^{\frac{1}{2}},
\end{gather*}
where $e_{*}=e_{k^{*}(R^{*})}$ is the charge of the stop level of the torus $R^{*}$ (Section~\ref{sec:constants-exponents-scales}). Since $0<\mathfrak{e}\le\frac{1}{4}$, $h'\le\check{h}$ and $\check{h}\ge1$, the upper bound for $e_{*}$ gives
\[
e_{*}^{\mathfrak{e}}\le\left(e\cdot h'^{\frac{1}{2}}\right)^{\mathfrak{e}}\le\check{h}^{\frac{\mathfrak{e}}{2}}\cdot e^{\mathfrak{e}}\le\check{h}^{\frac{1}{8}}\cdot e^{\mathfrak{e}};
\]
since $\mathfrak{e}_2=6\varepsilon$, the lower bound for $e_{*}$ gives
\[
e_{*}^{-6\varepsilon}\le\left(e\cdot\left(\frac{h'}{L^2}\right)^{\frac{1}{2}}\right)^{-6\varepsilon}=\left(\frac{L^2}{h'}\right)^{3\varepsilon}\cdot e^{-\mathfrak{e}_2}.
\]
Therefore
\begin{gather*}
\Theta^{(0)}_4\le3\cdot10^5\cdot\bar{C}^4\cdot\check{h}^{\frac{1}{8}}\cdot e^{\mathfrak{e}},\qquad\left(1+|c_2|\right)\cdot\Theta^{(0)}_2\le3\cdot10^4\cdot\bar{C}^4\cdot\check{h}^{\frac{1}{8}}\cdot e^{\mathfrak{e}},\\
|c_2|+|c_4|^{\frac{1}{2}}\le8\,\bar{C}^2\cdot\left(\frac{L^2}{h'}\right)^{3\varepsilon}\cdot e^{-\mathfrak{e}_2}.
\end{gather*}

\emph{Assertion (c).} By~\eqref{eq:four-jets-23},
\[
c_4-3c_2^2=b_{\eta,R^{*}}(0),
\]
the number of Section~\ref{sec:free-fourth-order} formed with the cochain $a_{\eta}$ of $\tau_v\beta$. The translation by $v\in\eta\mathbb{Z}^3$ is a lattice isometry of $\Lambda_{\eta,R^{*}}$, and it carries the cochain $a_{\eta}$ formed with $\beta$ to the cochain $a_{\eta}$ formed with $\tau_v\beta$; by the last assertion of Theorem~\ref{thm:free-coefficient-1}(d5), the number $b_{\eta,R^{*}}(0)$ is therefore also the number formed with the cochain $a_{\eta}$ of $\beta$. As $\beta$ fits $\left(N,T^{*}(e)\right)$, one has $D_{\beta}=2\rho_{\beta}\le l_c$; and $R^{*}=2L\cdot l_c\ge6\,l_c$, because $L\ge3$ (stage~\ref{it:order-choice-4}). Since $l_c\ge D_{\beta}\ge4$ (Section~\ref{sec:probes-discretisation}), it follows that
\[
R^{*}\ge D_{\beta}+5\,l_c\ge D_{\beta}+1 .
\]
Moreover the cube $Q_{\beta}$ has the side $D_{\beta}<R^{*}$, so that it lies in the interior of a period cube of the torus of side $R^{*}$, and $R^{*}/\eta=n_{*}$ is an integer. Theorem~\ref{thm:free-coefficient-1}(c2) and~(c3), applied with $R=R^{*}$ and $f=d\,d^{*}\beta$, give
\[
\left|b_{\eta,R^{*}}(0)-b(f)\right|\le\left|b_{\eta,R^{*}}(0)-b_{R^{*}}(f)\right|+\left|b_{R^{*}}(f)-b(f)\right|\le C_1\cdot\eta^{\frac{1}{2}}+C_2\cdot\exp\left(-\bar{m}\cdot\left(R^{*}-D_{\beta}\right)\right),
\]
with the constants $C_1$, $C_2$ of Theorem~\ref{thm:free-coefficient-1}(c), which depend only on $\bar{m}$ and $\beta$. By the definitions of Section~\ref{sec:constants-exponents-scales}, $\bar{m}\cdot R^{*}$ equals $\mathfrak{m}_K\cdot M\cdot r_K$ divided by $L$, and $\mathfrak{m}_K>L^{-m_0-1}$; stage~\ref{it:order-choice-6} gives $M\ge L^{m_0+2}$; the definition of $r_K$ gives $r_K\ge v_K^r$; and $v_K\ge1$ by Step~1, while $r\ge2$. Hence
\[
\bar{m}\cdot R^{*}=\frac{\mathfrak{m}_K\cdot M\cdot r_K}{L}>L^{-m_0-2}\cdot M\cdot r_K\ge r_K\ge v_K^r\ge v_K,
\]
so that $\exp\left(-\bar{m}\cdot R^{*}\right)\le\exp\left(-v_K\right)=e\cdot h'^{\frac{1}{2}}$, and
\[
\left|c_4-3c_2^2-b(f)\right|\le C_1\cdot\eta^{\frac{1}{2}}+C_2\cdot\exp\left(\bar{m}\cdot D_{\beta}\right)\cdot h'^{\frac{1}{2}}\cdot e.
\]
Finally $e\le e^{\sigma}$, because $e\le1$ by Step~1 and $\sigma\le\mathfrak{e}\le1$.

\emph{The constant $C^G_{\beta}$.} Put
\[
C^G_{\beta}:=\max\left\{3\cdot10^5\cdot\bar{C}^4\cdot\check{h}^{\frac{1}{8}},\ 8\,\bar{C}^2\cdot\left(\frac{L^2}{h'}\right)^{3\varepsilon},\ C_1,\ C_2\cdot\exp\left(\bar{m}\cdot D_{\beta}\right)\cdot h'^{\frac{1}{2}}\right\}<\infty .
\]
The numbers $e_1(\beta)$ and $C^G_{\beta}$ depend only on the constants of Section~\ref{sec:constants-exponents-scales}, on $\bar{m}$ and on $\beta$, since this holds for $e_G$ and $h'$ (Section~\ref{sec:constants-exponents-scales}), for $e_{\sharp}$ (Lemma~\ref{lem:threshold-charge}), for $C_{\beta}$ (Section~\ref{sec:probes-mirror-copies}), for $D_{\beta}$, and for $C_1$ and $C_2$ (Theorem~\ref{thm:free-coefficient-1}(c)). By the two preceding paragraphs, for every integer $N\ge N_0(e)$,
\[
\Theta^{(0)}_4\le C^G_{\beta}\cdot e^{\mathfrak{e}},\qquad\left(1+|c_2|\right)\cdot\Theta^{(0)}_2\le C^G_{\beta}\cdot e^{\mathfrak{e}},\qquad|c_2|+|c_4|^{\frac{1}{2}}\le C^G_{\beta}\cdot e^{-\mathfrak{e}_2},
\]
which are the assertions (a) and (b), and
\[
\left|c_4-3c_2^2-b\left(d\,d^{*}\beta\right)\right|\le C^G_{\beta}\cdot\left(e+\eta^{\frac{1}{2}}\right).
\]
As $\eta=L^{-N}\to0$ when $N\to\infty$, the upper limit of the left side as $N\to\infty$ is at most $C^G_{\beta}\cdot e\le C^G_{\beta}\cdot e^{\sigma}$. This is the assertion (c), and part~\ref{it:free-coefficient-3-2} is proved. Theorem~\ref{thm:bare-mass-shifts} follows.
\end{proof}

\subsection{Proof of Part~II of the main theorem and of its two corollaries}\label{sec:proof-part-ii}

\subsubsection{Proof of Part~II of the main theorem}\label{sec:proof-part-ii-2}

The assertions of Theorem~\ref{thm:main-theorem} on the numbers $m_{0,N}(e)$ are contained in part~\ref{it:free-coefficient-3-1} of Theorem~\ref{thm:free-coefficient-3}: for every $e>0$ with $e\cdot h'^{\frac{1}{2}}\le e_1$ and every integer $N\ge0$ the number $m_{0,N}(e)$ is real, satisfies $\left|m_{0,N}(e)\right|<\frac{1}{8}<\frac{1}{2}$, and does not depend on a torus.

Let $\beta$, $x_0$, $d_0$, $\beta_0$, $\beta_{\pm}$, $g_0$ and $g_{\pm}$ be the data of Section~\ref{sec:probes}, and let
\[
\mathsf{B}:=\left\{\beta_0,\beta_{+},\beta_{-}\right\}.
\]
Each element of $\mathsf{B}$ is a probe in the sense of Section~\ref{sec:probes-discretisation}: it is a real two-form in $C_c^{\infty}(\mathbb{R}^3;\Lambda^2)$, and it is not zero, because $\beta\ne0$, as $b(f)>0$, and the supports of $\beta$ and $\tau_{x_0}\beta$ are disjoint, as $d_0>0$. For $\beta'\in\mathsf{B}$ let $g:=dd^{*}\beta'$, so that $g\in\left\{g_0,g_{+},g_{-}\right\}$, and put
\[
s_g:=\|g\|_2+\|\beta'\|_{H^4};
\]
for $\beta'=\beta$ this is the number $s_f$ of Theorem~\ref{thm:free-coefficient-1}(b5). Let
\[
C_B:=\max_{\beta'\in\mathsf{B}}C^G_{\beta'},\qquad s_B:=\max_{\beta'\in\mathsf{B}}s_g,
\]
where $C^G_{\beta'}$ is the constant of part~\ref{it:free-coefficient-3-2} of Theorem~\ref{thm:free-coefficient-3} for the probe $\beta'$. In Section~\ref{sec:proof-part-ii} the symbol $C_B$ denotes this maximum, and not the constant $C_B$ of stage~\ref{it:order-choice-7}. In the same section the letter $C$ is reserved for constants that are determined by the constants of Section~\ref{sec:constants-exponents-scales} together with $\bar{m}$, $\beta$ and $x_0$; two occurrences of $C$ need not stand for the same constant.

\paragraph{Step 0: an upper bound for $e_2$.}
For $\beta'\in\mathsf{B}$ let $e_1(\beta')\in\left(0,e_G\right]$ be the number of part~\ref{it:free-coefficient-3-2} of Theorem~\ref{thm:free-coefficient-3} for the probe $\beta'$; by its construction in Section~\ref{sec:proof-theorem}, $e_1(\beta')=\min\left(e_G,e_{\sharp}(\beta')\right)$. The number $e_2$ is required to satisfy
\[
e_2\le\min_{\beta'\in\mathsf{B}}e_1(\beta'),
\]
and this requirement is the conjunction of~\eqref{eq:free-coefficient-1} and~\eqref{eq:free-coefficient-2}. Let $0<e\le e_2$, and let $N\ge N_0(e)$ be an integer. As $e\le e_G\le e_1\cdot h'^{-\frac{1}{2}}$ (Section~\ref{sec:constants-exponents-scales}), part~\ref{it:free-coefficient-3-1} of Theorem~\ref{thm:free-coefficient-3} applies: $e\le\frac{1}{3}$; the number $T^{*}(e)$ is an integer $\ge1$; the lattice mass $\boldsymbol{\mu}_N=\bar{m}\cdot\eta+m_{0,N}(e)$ satisfies $\boldsymbol{\mu}_N>-\frac{1}{8}>-\frac{1}{2}$, because $\left|m_{0,N}(e)\right|<\frac{1}{8}$; and, as $e\le e_G$, the inequality~\eqref{eq:free-coefficient-6} holds for $R^{*}=R^{*}(e)$. Since $e\le e_1(\beta')$, part~\ref{it:free-coefficient-3-2} of Theorem~\ref{thm:free-coefficient-3} shows that each $\beta'\in\mathsf{B}$ fits $\left(N,T^{*}(e)\right)$. Hence Theorem~\ref{thm:free-coefficient-2} applies to each $\beta'\in\mathsf{B}$, with $T^{*}=T^{*}(e)$, with the placement $v$ of $\beta'$ (Section~\ref{sec:probes-discretisation}), and with the real numbers $c_m=c_m(N,e,\beta')$, $m=2,4$, of part~\ref{it:free-coefficient-3-2} of Theorem~\ref{thm:free-coefficient-3}. Moreover, by Theorem~\ref{thm:free-coefficient-1}(b5) applied to $\beta'$, the number $\boldsymbol{\sigma}$ of Section~\ref{sec:probes-discretisation} formed with $\beta'$ satisfies
\[
\boldsymbol{\sigma}\le s_g\le s_B.
\]

\paragraph{Step 1: the factor $\boldsymbol{\Gamma}^{\mathrm{nat}}$.}
Let $e$ and $N$ be as in Step~0. As $\bar{m}\le1$, one has $h\ge1$, and as $T^{*}(e)\ge1$,
\[
R^{*}=R^{*}(e)=h\cdot L^{T^{*}(e)}\ge h\ge1 .
\]
Hence $\log\left(\frac{2\pi}{e^2\cdot R^{*}}\right)\le\log(2\pi)+2\log\left(e^{-1}\right)$. Together with the bound $D^{\mathrm{mat}}\le C_{\mathrm{mat}}$ of part~\ref{it:free-coefficient-3-1} of Theorem~\ref{thm:free-coefficient-3}, the inequality~\eqref{eq:free-coefficient-9} gives
\[
\exp(\delta)\le(2\pi)^{\frac{a_L}{N^{*}}}\cdot\exp\left(\frac{K_L+C_{\mathrm{mat}}+1}{(L-1)^3\cdot N^{*}}\right)\cdot e^{-\frac{2a_L}{N^{*}}}=:C_{\mathrm{Gam}}\cdot e^{-\frac{2a_L}{N^{*}}}.
\]
Next, part~\ref{it:free-coefficient-3-2}(b) of Theorem~\ref{thm:free-coefficient-3} gives $|c_2|\le C_B\cdot e^{-\mathfrak{e}_2}$ and $|c_4|\le C_B^2\cdot e^{-2\mathfrak{e}_2}$, and one has $1\le e^{-\mathfrak{e}_2}\le e^{-2\mathfrak{e}_2}$, because $e\le1$ and $\mathfrak{e}_2\ge0$. With $\boldsymbol{\sigma}\le s_B$, the monotonicity of $\log^{+}$ and the inequality $\log^{+}(x\cdot y)\le\log^{+}x+\log^{+}y$ for $x,y>0$, it follows that the number $Y$ of Theorem~\ref{thm:free-coefficient-2} satisfies
\[
Y\le C_Y+2\mathfrak{e}_2\cdot\log\left(e^{-1}\right),\qquad C_Y:=\log^{+}\left((N^{*})^2\cdot\left(600\,s_B^4+12\,s_B^2+C_B^2+C_B\right)\right).
\]
Since $0<e\le e_2\le\frac{1}{3}$ by Step~0, one has $\log\left(e^{-1}\right)\ge\log\left(e_2^{-1}\right)>0$, and~\eqref{eq:free-coefficient-3} yields
\[
\check{X}=1+\frac{Y+1}{6\log\left(e^{-1}\right)}\le1+\frac{C_Y+1}{6\log\left(e^{-1}\right)}+\frac{\mathfrak{e}_2}{3}\le1+\frac{C_Y+1}{6\log\left(e_2^{-1}\right)}+\frac{\mathfrak{e}_2}{3}\le2 .
\]
Therefore $4+16\,\check{X}\le36$, and, as $3\log36+6<17$,
\[
\boldsymbol{\Gamma}^{\mathrm{av}}=\exp\left(\delta+\frac{3\log\left(4+16\,\check{X}\right)+6}{N^{*}}\right)\le\exp\left(\frac{3\log36+6}{N^{*}}\right)\cdot\exp(\delta)\le\exp\left(\frac{17}{N^{*}}\right)\cdot\exp(\delta).
\]
Since $\boldsymbol{\Gamma}^{\mathrm{af}}=\exp(\delta)$, it follows that for both laws
\[
\boldsymbol{\Gamma}^{\mathrm{nat}}\le\exp\left(\frac{17}{N^{*}}\right)\cdot C_{\mathrm{Gam}}\cdot e^{-\frac{2a_L}{N^{*}}}\le C\cdot e^{-\frac{2a_L}{N^{*}}}.
\]

\paragraph{Step 2: the factors $\boldsymbol{\Theta}^{\mathrm{nat}}_m$.}
By part~\ref{it:free-coefficient-3-2}(a) of Theorem~\ref{thm:free-coefficient-3}, $\Theta^{(0)}_4\le C_B\cdot e^{\mathfrak{e}}$ and $\left(1+|c_2|\right)\cdot\Theta^{(0)}_2\le C_B\cdot e^{\mathfrak{e}}$; by part~\ref{it:free-coefficient-3-2}(b), $|c_2|\le C_B\cdot e^{-\mathfrak{e}_2}$ and $|c_4|\le C_B^2\cdot e^{-2\mathfrak{e}_2}$; and $\boldsymbol{\sigma}\le s_B$ by Step~0. Inserting these bounds into the definitions of $\dot{M}_m$ and $\boldsymbol{\Theta}^{\mathrm{nat}}_m$ in Theorem~\ref{thm:free-coefficient-2}, one obtains
\begin{align*}
\boldsymbol{\Theta}^{\mathrm{af}}_4&\le C_B\cdot e^{\mathfrak{e}}+2e^6\cdot\left(600\,(N^{*})^2\cdot s_B^4\cdot e^{-4}+C_B^2\cdot e^{-2\mathfrak{e}_2}\right)\le C\cdot e^{\mathfrak{e}},\\
\left(1+|c_2|\right)\cdot\boldsymbol{\Theta}^{\mathrm{af}}_2&\le C_B\cdot e^{\mathfrak{e}}+2e^6\cdot\left(12N^{*}\cdot s_B^2\cdot e^{-2}+C_B\cdot e^{-\mathfrak{e}_2}\right)\cdot\left(1+C_B\cdot e^{-\mathfrak{e}_2}\right)\le C\cdot e^{\mathfrak{e}},\\
\boldsymbol{\Theta}^{\mathrm{av}}_4&\le\max\left(C_B\cdot e^{\mathfrak{e}},e^2\right)\le C\cdot e^{\mathfrak{e}},\\
\left(1+|c_2|\right)\cdot\boldsymbol{\Theta}^{\mathrm{av}}_2&\le\max\left(C_B\cdot e^{\mathfrak{e}},e^2\cdot\left(1+C_B\cdot e^{-\mathfrak{e}_2}\right)\right)\le C\cdot e^{\mathfrak{e}}.
\end{align*}
In each line the last inequality holds because $e\le1$, and because $\mathfrak{e}\le1$ and $\mathfrak{e}_2<1$ by part~\ref{it:free-coefficient-3-2} of Theorem~\ref{thm:free-coefficient-3}: every power of $e$ on the right other than $e^{\mathfrak{e}}$ has an exponent at least $2-\mathfrak{e}_2>1\ge\mathfrak{e}$.

\paragraph{Step 3: the cumulants on a fixed torus.}
Let $0<e\le e_2$ and let $N\ge N_0(e)$ be an integer. Let $k\in\{1\}\cup\{L,L+1,\dots\}$ be an integer, let $\mathrm{nat}\in\{\mathrm{af},\mathrm{av}\}$, and let $\beta'\in\mathsf{B}$, $g=dd^{*}\beta'$, with the placement $v$ of $\beta'$. Let $g^{(\eta)}$ be the two-form $f^{(\eta)}$ of Section~\ref{sec:probes-discretisation}, formed with $\eta=L^{-N}$ and with $\beta'$ in place of $\beta$. Put $R:=k\cdot R^{*}$, so that the torus $k$ of Section~\ref{sec:reference-tori-cells} is $\Lambda_{\eta,R}$, and
\[
X:=X_{\eta,R}\left(\tau_vg^{(\eta)}\right);
\]
by Theorem~\ref{thm:free-coefficient-1}(b2) this is the variable $\langle d\hat{\mathfrak{a}},F^u\rangle$ of Theorem~\ref{thm:free-coefficient-2} formed with $\beta'$. In this step all cumulants are taken under $E^{\mathrm{nat}}_k$. By Steps~1 and~2, $\boldsymbol{\Gamma}^{\mathrm{nat}}\cdot\boldsymbol{\Theta}^{\mathrm{nat}}_4\le C\cdot e^{\mathfrak{e}-\frac{2a_L}{N^{*}}}=C\cdot e^{\sigma}$ and $\boldsymbol{\Gamma}^{\mathrm{nat}}\cdot\left(1+|c_2|\right)\cdot\boldsymbol{\Theta}^{\mathrm{nat}}_2\le C\cdot e^{\sigma}$; moreover $e^{2\sigma}\le e^{\sigma}\le1\le e^{-\mathfrak{e}_2}$. Inserting these bounds and part~\ref{it:free-coefficient-3-2}(b) of Theorem~\ref{thm:free-coefficient-3} into~\eqref{eq:free-coefficient-7} and~\eqref{eq:free-coefficient-8}, one obtains
\begin{align}
\left|\kappa_4(X)-e^4\cdot\left(c_4-3c_2^2\right)\right|&\le e^4\cdot\left[\boldsymbol{\Gamma}^{\mathrm{nat}}\cdot\boldsymbol{\Theta}^{\mathrm{nat}}_4+6\,\boldsymbol{\Gamma}^{\mathrm{nat}}\cdot\boldsymbol{\Theta}^{\mathrm{nat}}_2\cdot|c_2|+3\left(\boldsymbol{\Gamma}^{\mathrm{nat}}\cdot\boldsymbol{\Theta}^{\mathrm{nat}}_2\right)^2\right]\le C\cdot e^{4+\sigma},\notag\\
\left|\kappa_2(X)-\boldsymbol{\sigma}^2\right|&\le e^2\cdot\left(C_B\cdot e^{-\mathfrak{e}_2}+C\cdot e^{\sigma}\right)\le C\cdot e^{2-\mathfrak{e}_2}.\label{eq:free-coefficient-20}
\end{align}

By Lemma~\ref{lem:lattice-symmetries-cell}, $\left(\tau_v\varphi\right)_{\eta}=v_{*}\varphi_{\eta}$ for $\varphi\in\mathcal{S}$, and $v_{*}$ is orthogonal for $\langle\cdot,\cdot\rangle_{\eta}$; so $\langle\left(\tau_v\varphi\right)_{\eta},F\rangle_{\eta}=\langle\varphi_{\eta},v_{*}^{-1}F\rangle_{\eta}$. As the law of $F$ under $E^{\mathrm{nat}}_k$ is invariant under the lattice translations (Theorem~\ref{thm:free-coefficient-1}(d5)), the variables $X_{\eta,R}\left(\tau_vg^{(\eta)}\right)$ and $X_{\eta,R}\left(g^{(\eta)}\right)$ have the same law under $E^{\mathrm{nat}}_k$, and in particular the same cumulants.

For every real two-cochain $\mathbf{g}$ on the torus $k$ both laws satisfy
\[
E^{\mathrm{nat}}_k\exp\langle\mathbf{g},F^u\rangle\le\exp\left(\frac{1}{2}\|\mathbf{g}\|^2\right),
\]
by~\eqref{it:free-coefficient-Q3} for $\mathrm{nat}=\mathrm{af}$ and by Theorem~\ref{thm:limit-vanishing-proca} for $\mathrm{nat}=\mathrm{av}$; and both laws are invariant under $F^u\mapsto-F^u$ (Theorem~\ref{thm:free-coefficient-1}(d5)). Let $y>0$. By this invariance,
\[
E^{\mathrm{nat}}_k\cosh\left(y\cdot\langle\mathbf{g},F^u\rangle\right)=E^{\mathrm{nat}}_k\exp\left(y\cdot\langle\mathbf{g},F^u\rangle\right)\le\exp\left(\frac{1}{2}y^2\cdot\|\mathbf{g}\|^2\right),
\]
and the inequalities $\frac{x^4}{24}\le\cosh x$ and $\frac{x^2}{2}\le\cosh x-1$, valid for real $x$, give
\[
E^{\mathrm{nat}}_k\langle\mathbf{g},F^u\rangle^4\le\frac{24}{y^4}\cdot\exp\left(\frac{1}{2}y^2\cdot\|\mathbf{g}\|^2\right),\qquad E^{\mathrm{nat}}_k\langle\mathbf{g},F^u\rangle^2\le\frac{2}{y^2}\cdot\left(\exp\left(\frac{1}{2}y^2\cdot\|\mathbf{g}\|^2\right)-1\right).
\]
If $\mathbf{g}\ne0$, the choice $y^2\cdot\|\mathbf{g}\|^2=4$ in the first inequality and the limit $y\to0$ in the second yield
\[
E^{\mathrm{nat}}_k\langle\mathbf{g},F^u\rangle^4\le\frac{3}{2}\exp(2)\cdot\|\mathbf{g}\|^4,\qquad E^{\mathrm{nat}}_k\langle\mathbf{g},F^u\rangle^2\le\|\mathbf{g}\|^2;
\]
for $\mathbf{g}=0$ both bounds are trivial. The variables $\langle\mathbf{g},F^u\rangle$ are centred, by the invariance under $F^u\mapsto-F^u$. For centred random variables $X_1,\dots,X_4$ with finite fourth moments the joint cumulants of order $4$ and $2$ are
\begin{gather*}
\kappa_4(X_1,X_2,X_3,X_4)=E[X_1\cdot X_2\cdot X_3\cdot X_4]-E[X_1\cdot X_2]\cdot E[X_3\cdot X_4]-E[X_1\cdot X_3]\cdot E[X_2\cdot X_4]-E[X_1\cdot X_4]\cdot E[X_2\cdot X_3],\\
\kappa_2(X_1,X_2)=E[X_1\cdot X_2];
\end{gather*}
$\kappa_4$ is a symmetric multilinear function of its four arguments. By H\"older's inequality the first term of $\kappa_4$ is at most $\prod_{i=1}^4\left(E\,X_i^4\right)^{\frac{1}{4}}$ in modulus, and by the Cauchy--Schwarz inequality each of the other three terms, and $\kappa_2(X_1,X_2)$, is at most the corresponding product of the numbers $\left(E\,X_i^2\right)^{\frac{1}{2}}$ in modulus. For $X_i=\langle\mathbf{g}_i,F^u\rangle$ with real two-cochains $\mathbf{g}_1,\dots,\mathbf{g}_4$, the moment bounds above therefore give
\begin{gather*}
\left|\kappa_4\left(\langle\mathbf{g}_1,F^u\rangle,\dots,\langle\mathbf{g}_4,F^u\rangle\right)\right|\le\left(\frac{3}{2}\exp(2)+3\right)\cdot\prod_{i=1}^4\|\mathbf{g}_i\|\le15\cdot\prod_{i=1}^4\|\mathbf{g}_i\|,\\
\left|\kappa_2\left(\langle\mathbf{g}_1,F^u\rangle,\langle\mathbf{g}_2,F^u\rangle\right)\right|\le\|\mathbf{g}_1\|\cdot\|\mathbf{g}_2\|.
\end{gather*}

The two-forms $g$ and $g^{(\eta)}$ are smooth and supported in $Q_{\beta'}$ (Theorem~\ref{thm:free-coefficient-1}(b2) and the definition of $Q_{\beta'}$ in Section~\ref{sec:probes-discretisation}). The side of the cube $Q_{\beta'}$ is $D_{\beta'}\le l_c-3\eta<R$, because $\beta'$ fits $\left(N,T^{*}(e)\right)$; so $Q_{\beta'}$ lies in a period cube $y_0+\left(-\frac{R}{2},\frac{R}{2}\right)^3$ of $\Lambda_{\eta,R}$. Lemma~\ref{lem:norm-bounds-lattice}, applied to $g-g^{(\eta)}$, Theorem~\ref{thm:free-coefficient-1}(b3) with $s=0$, and Plancherel's theorem, by which $\|\cdot\|_{H^0}=\|\cdot\|_2$, give
\[
\left|\left(g-g^{(\eta)}\right)_{\eta}\right|_{\eta}\le\left\|g-g^{(\eta)}\right\|_2\le\eta^2\cdot\|\beta'\|_{H^4};
\]
in the same way $\left|g_{\eta}\right|_{\eta}\le\|g\|_2\le s_g$ and $\left|\left(g^{(\eta)}\right)_{\eta}\right|_{\eta}\le\left\|g^{(\eta)}\right\|_2\le s_g$, the last inequality by Theorem~\ref{thm:free-coefficient-1}(b5).

For $\varphi\in\mathcal{S}$ one has $X_{\eta,R}(\varphi)=\langle\varphi_{\eta},F\rangle_{\eta}=\langle\mathbf{g},F^u\rangle$ with $\mathbf{g}:=\eta^{\frac{3}{2}}\cdot\varphi_{\eta}$ and $\|\mathbf{g}\|=\left|\varphi_{\eta}\right|_{\eta}$ (Section~\ref{sec:lattices-tests-function}). Write $X_{\eta,R}(g)=X_{\eta,R}\left(g^{(\eta)}\right)+X_{\eta,R}\left(g-g^{(\eta)}\right)$ and replace, in the multilinear cumulants, one argument at a time. This expresses $\kappa_4\left(X_{\eta,R}(g)\right)-\kappa_4\left(X_{\eta,R}\left(g^{(\eta)}\right)\right)$ as a sum of four joint cumulants, and $\kappa_2\left(X_{\eta,R}(g)\right)-\kappa_2\left(X_{\eta,R}\left(g^{(\eta)}\right)\right)$ as a sum of two, each of which has exactly one argument equal to $X_{\eta,R}\left(g-g^{(\eta)}\right)$ and its other arguments equal to $X_{\eta,R}(g)$ or $X_{\eta,R}\left(g^{(\eta)}\right)$. The bounds for $\kappa_4$ and $\kappa_2$ obtained above, applied with the three norms just estimated, yield
\begin{align}
\left|\kappa_4\left(X_{\eta,R}(g)\right)-\kappa_4\left(X_{\eta,R}\left(g^{(\eta)}\right)\right)\right|&\le60\,\eta^2\cdot\|\beta'\|_{H^4}\cdot s_g^3,\notag\\
\left|\kappa_2\left(X_{\eta,R}(g)\right)-\kappa_2\left(X_{\eta,R}\left(g^{(\eta)}\right)\right)\right|&\le2\,\eta^2\cdot\|\beta'\|_{H^4}\cdot s_g.\label{eq:free-coefficient-21}
\end{align}

Now let $T\ge T^{*}(e)$ be an integer and $R=h\cdot L^T$. The aligned torus $\Lambda_{\eta,R}$ is the torus $k=L^{T-T^{*}(e)}$ of Section~\ref{sec:reference-tori-cells}, and $k\in\left\{1,L,L^2,\dots\right\}\subset\{1\}\cup\{L,L+1,\dots\}$; on it $E^{\mathrm{nat}}=E^{\mathrm{nat}}_k$. Since $X_{\eta,R}\left(\tau_vg^{(\eta)}\right)$ and $X_{\eta,R}\left(g^{(\eta)}\right)$ have the same cumulants, \eqref{eq:free-coefficient-20} and~\eqref{eq:free-coefficient-21} give, for $g\in\left\{g_0,g_{+},g_{-}\right\}$,
\[
\left|\kappa_4\left(X_{\eta,R}(g);E^{\mathrm{nat}}\right)-e^4\cdot\left(c_4-3c_2^2\right)\right|\le C\cdot e^{4+\sigma}+60\,\eta^2\cdot\|\beta'\|_{H^4}\cdot s_g^3.
\]
By part~\ref{it:free-coefficient-3-2}(c) of Theorem~\ref{thm:free-coefficient-3}, $c_4-3c_2^2$ is the number $b_{\eta,R^{*}(e)}(0)$ formed with $\beta'$; so this is the first inequality of assertion~(ii) of Theorem~\ref{thm:main-theorem}, provided $C_{\mathrm{kap}}$ is at least twice each of the constants that occur. Part~\ref{it:free-coefficient-3-2}(c) of Theorem~\ref{thm:free-coefficient-3} also gives $\left|c_4-3c_2^2-b(g)\right|\le C_B\cdot\left(e+\eta^{\frac{1}{2}}\right)$, and $e\le e^{\sigma}$, since $e\le1$ and $0<\sigma\le\mathfrak{e}\le1$. Hence
\[
\left|\kappa_4\left(X_{\eta,R}(g);E^{\mathrm{nat}}\right)-e^4\cdot b(g)\right|\le\left(C+C_B\right)\cdot e^{4+\sigma}+60\,\eta^2\cdot\|\beta'\|_{H^4}\cdot s_g^3+C_B\cdot e^4\cdot\eta^{\frac{1}{2}},
\]
which is the second inequality of assertion~(ii) under the same condition on $C_{\mathrm{kap}}$.

\paragraph{Step 4: passage to the limit.}
Assertion~(i) of Theorem~\ref{thm:main-theorem} is Part~I of that theorem for the mass assignment $N\mapsto m_{0,N}(e)$ of part~\ref{it:free-coefficient-3-1} of Theorem~\ref{thm:free-coefficient-3}. Let $\mu$ be a limit point: $P^{\mathrm{nat}}_{N_i,T_i}\to\mu$ weakly on $H_{-4}$ for a sequence of cutoffs $(N_i,T_i)$ with $N_i\to\infty$ and $T_i\to\infty$; put $\eta_i:=L^{-N_i}$ and $R_i:=h\cdot L^{T_i}$. For all sufficiently large $i$ one has $N_i\ge N_0(e)$ and $T_i\ge T^{*}(e)$. Let $g\in\left\{g_0,g_{+},g_{-}\right\}$. By property~\eqref{it:main-theorem-P8} of Theorem~\ref{thm:main-theorem},
\[
\kappa_m\left(X_{\eta_i,R_i}(g);E^{\mathrm{nat}}\right)\to\kappa_m^{\mu}(g)\qquad(i\to\infty,\ m=2,4).
\]
As $\eta_i\to0$, the second inequality of assertion~(ii) gives in the limit
\[
\left|\kappa_4^{\mu}(g)-e^4\cdot b(g)\right|\le\frac{1}{2}C_{\mathrm{kap}}\cdot e^{4+\sigma},
\]
which implies the first inequality of assertion~(iii). For the second cumulant take $\beta'=\beta$, so that $g=f$. By Theorem~\ref{thm:free-coefficient-1}(b5), $\boldsymbol{\sigma}^2\to\|f\|_2^2$ as $\eta\to0$, the number $\boldsymbol{\sigma}$ being formed with $\beta$. By~\eqref{eq:free-coefficient-20}, by~\eqref{eq:free-coefficient-21} and by the equality of the cumulants of $X_{\eta,R}\left(\tau_vf^{(\eta)}\right)$ and $X_{\eta,R}\left(f^{(\eta)}\right)$ shown in Step~3, for every integer $N\ge N_0(e)$, every integer $T\ge T^{*}(e)$ and $R=h\cdot L^T$,
\[
\left|\kappa_2\left(X_{\eta,R}(f);E^{\mathrm{nat}}\right)-\boldsymbol{\sigma}^2\right|\le C\cdot e^{2-\mathfrak{e}_2}+2\eta^2\cdot\|\beta\|_{H^4}\cdot s_f.
\]
Moreover $\kappa_2^{\mu}(f)=S_2^{\mu}(f,f)$, because $S_1^{\mu}=0$ by property~\eqref{it:main-theorem-P3} of Theorem~\ref{thm:main-theorem}. Letting $i\to\infty$ one obtains
\[
\left|S_2^{\mu}(f,f)-\|f\|_2^2\right|\le C\cdot e^{2-\mathfrak{e}_2},
\]
which is the second inequality of assertion~(iii) as soon as $C_{\mathrm{kap}}\ge C$.

\paragraph{Step 5: assertions (iv) and (v).}
The limit point $\mu$ satisfies~\eqref{eq:fourth-cumulants-H1}, which is the condition~\eqref{eq:setup-G}, and it is invariant under the translations of $\mathbb{R}^3$ (Part~I of Theorem~\ref{thm:main-theorem}). Assertion~(iii) is the hypothesis~\eqref{eq:fourth-cumulants-4g} of Section~\ref{sec:hypotheses-law} at the three tests $g_0=f$, $g_{+}=f+\tau_{x_0}f$ and $g_{-}=f-\tau_{x_0}f$ of Theorem~\ref{thm:non-gaussian-components-2}, with the exponent $\sigma$ and with each of the three constants that the hypothesis carries at these tests equal to $C_{\mathrm{kap}}$. The condition~\eqref{eq:setup-1} on $x_0$ is the condition~\eqref{eq:fourth-cumulants-11} of Theorem~\ref{thm:non-gaussian-components-2}, read with $m=\bar{m}$ and $r=r_{\beta}$. With these constants, the number $C_{*}$ of Theorem~\ref{thm:non-gaussian-components-2} becomes
\[
C_{*}=2C_{\mathrm{kap}}+\frac{1}{4}\left(C_{\mathrm{kap}}+C_{\mathrm{kap}}\right)=\frac{5}{2}\,C_{\mathrm{kap}},
\]
and the condition~\eqref{eq:free-coefficient-4} yields its hypothesis $C_{*}\cdot e^{\sigma}\le\frac{1}{4}b(f)$ for every $e\in(0,e_2]$, since $e^{\sigma}\le e_2^{\sigma}$. Theorem~\ref{thm:non-gaussian-components-2} therefore applies to $\mu$. Its assertion~(i) is the inequality $\mathrm{Ext}^{\mu}_{x_0}(f)\ge\frac{1}{2}e^4\cdot b(f)$ of~(iv). Its assertion~(iii), combined with $\frac{3}{2}\exp(2)+6<17.09$, gives the inequality
\[
\nu\left(\left\{\xi:\ \mathrm{Ext}^{\xi}_{x_0}(f)\ge\frac{1}{4}e^4\cdot b(f)\right\}\right)\ge\frac{e^4\cdot b(f)}{4\cdot17.09\cdot\|f\|_2^4}
\]
of~(v), and its assertion~(iv) is the last assertion of~(v). By~(iii) at $g=f$, and by~\eqref{eq:free-coefficient-4},
\[
\kappa_4^{\mu}(f)\ge e^4\cdot b(f)-C_{\mathrm{kap}}\cdot e^{4+\sigma}\ge\frac{1}{2}e^4\cdot b(f)>0.
\]
Every Gaussian law with the property~\eqref{eq:setup-G} has mean zero and vanishing fourth cumulant (Section~\ref{sec:cumulants-gaussian-laws}); consequently $\mu$ is not Gaussian. That $\mu$ is not a mixture of translation-invariant Gaussian laws satisfying~\eqref{eq:setup-G} is the last assertion of Theorem~\ref{thm:non-gaussian-components}, applied with $x=x_0$, for which $\mathrm{Ext}^{\mu}_{x_0}(f)>0$.

This can also be checked directly. Suppose that $\mu$ is a mixture, in the sense of Section~\ref{sec:laws-2}, of translation-invariant Gaussian laws that satisfy~\eqref{eq:setup-G}: there are a probability space $(\boldsymbol{\Lambda},m_{\mathrm{x}})$ and translation-invariant Gaussian laws $\boldsymbol{\gamma}_{\mathfrak{l}}$, $\mathfrak{l}\in\boldsymbol{\Lambda}$, each of which satisfies~\eqref{eq:setup-G}, such that $\mathfrak{l}\mapsto\boldsymbol{\gamma}_{\mathfrak{l}}(A)$ is measurable and
\[
\mu(A)=\int\mathrm{d}m_{\mathrm{x}}(\mathfrak{l})\,\boldsymbol{\gamma}_{\mathfrak{l}}(A)
\]
for every cylinder set $A$. Every $\boldsymbol{\gamma}_{\mathfrak{l}}$ is centred, and its covariance $C_{\mathfrak{l}}$ satisfies $\left|C_{\mathfrak{l}}(u,u')\right|\le\|u\|_2\cdot\|u'\|_2$ for $u,u'\in\mathcal{S}$ (Lemma~\ref{lem:moment-bounds}(a), (b)). The moments of $\mu$ are the $m_{\mathrm{x}}$-averages of the moments of the laws $\boldsymbol{\gamma}_{\mathfrak{l}}$; so Wick's rule gives, for $u_1,\dots,u_4\in\mathcal{S}$,
\begin{align*}
\kappa_4^{\mu}(u_1,u_2,u_3,u_4)&=\operatorname{Cov}_{m_{\mathrm{x}}}\left(C_{\mathfrak{l}}(u_1,u_2),C_{\mathfrak{l}}(u_3,u_4)\right)+\operatorname{Cov}_{m_{\mathrm{x}}}\left(C_{\mathfrak{l}}(u_1,u_3),C_{\mathfrak{l}}(u_2,u_4)\right)\\
&\quad+\operatorname{Cov}_{m_{\mathrm{x}}}\left(C_{\mathfrak{l}}(u_1,u_4),C_{\mathfrak{l}}(u_2,u_3)\right).
\end{align*}
The laws $\boldsymbol{\gamma}_{\mathfrak{l}}$ are translation invariant, so that $C_{\mathfrak{l}}(\tau_xf,\tau_xf)=C_{\mathfrak{l}}(f,f)$ for $x\in\mathbb{R}^3$. Consequently
\[
\kappa_4^{\mu}(f)=3\operatorname{Var}_{m_{\mathrm{x}}}\left(C_{\mathfrak{l}}(f,f)\right),\qquad\kappa_4^{\mu}(f,f,\tau_xf,\tau_xf)=\operatorname{Var}_{m_{\mathrm{x}}}\left(C_{\mathfrak{l}}(f,f)\right)+2\operatorname{Var}_{m_{\mathrm{x}}}\left(C_{\mathfrak{l}}(f,\tau_xf)\right),
\]
and therefore
\[
\mathrm{Ext}^{\mu}_x(f)=-6\operatorname{Var}_{m_{\mathrm{x}}}\left(C_{\mathfrak{l}}(f,\tau_xf)\right)\le0\qquad\text{for every }x\in\mathbb{R}^3,
\]
which contradicts $\mathrm{Ext}^{\mu}_{x_0}(f)>0$.

\paragraph{Step 6: assertions (vi) and (vii).}
For~(vi) only the probe $\beta=\beta_0$ is needed. Let $\mathrm{nat}$ and the sequence $(N_i,T_i)_{i\ge1}$ with $P^{\mathrm{nat}}_{N_i,T_i}\to\mu$ be those of Step~4. For $N\ge N_0(e)$ and $m\in\{2,4\}$ let
\[
\mathsf{u}_m(N):=\boldsymbol{\Gamma}^{\mathrm{nat}}\cdot\boldsymbol{\Theta}^{\mathrm{nat}}_m
\]
be the product of the two numbers of Theorem~\ref{thm:free-coefficient-2} formed for the probe $\beta$ at the cutoff $N$, with $T^{*}=T^{*}(e)$ and with $c_m=c_m(N,e,\beta)$, as in Step~0. Steps~1 and~2 give
\[
\mathsf{u}_4(N)\le C\cdot e^{\sigma},\qquad\left(1+\left|c_2(N,e,\beta)\right|\right)\cdot\mathsf{u}_2(N)\le C\cdot e^{\sigma}\qquad\left(N\ge N_0(e)\right).
\]
Let
\[
\hat{\Gamma}_0:=\frac{R^{*}}{L}\cdot\mathbb{Z}^3,\qquad R^{*}=R^{*}(e).
\]
For every $N$ one has $R^{*}/L=\mathrm{per}\cdot\eta$ with the integer $\mathrm{per}$ of Section~\ref{sec:reference-tori-cells}; hence $\hat{\Gamma}_0\subset\eta\mathbb{Z}^3$. For $m\in\{2,4\}$ and real numbers $x$, $y$, $c$ put
\[
\mathrm{Pm}(x;y,c):=\begin{cases}e^{-2}\left(x^2-y\right)-c&\text{if }m=2,\\ e^{-4}\left(x^4-6y\cdot x^2+3y^2\right)-c&\text{if }m=4;\end{cases}
\]
thus $e^{-m}H_m(x)-c_m=\mathrm{Pm}\left(x;\boldsymbol{\sigma}^2,c_m\right)$.

Let $J\ge1$, and let $y_1,\dots,y_J\in\hat{\Gamma}_0$ be pairwise distinct, $y_i=\frac{R^{*}}{L}\cdot\boldsymbol{\zeta}_i$ with $\boldsymbol{\zeta}_i\in\mathbb{Z}^3$. Let $N\ge N_0(e)$ and $T\ge T^{*}(e)$, and let $\eta=L^{-N}$, $R=h\cdot L^T$; the torus $\Lambda_{\eta,R}$ is the torus $k=L^{T-T^{*}(e)}$ of Section~\ref{sec:reference-tori-cells}, formed with $T^{*}=T^{*}(e)$. The cell $c^{(i)}:=2\boldsymbol{\zeta}_i$ of this torus has even coordinates. Therefore $\mathsf{g}_{c^{(i)}}$ is the translation by $y_i$ (Section~\ref{sec:reference-tori-cells}), no complex conjugation enters the mirror copies in the cell $\mathfrak{c}_{c^{(i)}}$, and $\hat{\mathfrak{a}}_{c^{(i)}}=\left(\mathsf{g}_{c^{(i)}}\right)_{*}\hat{\mathfrak{a}}=(y_i)_{*}\hat{\mathfrak{a}}$. By Theorem~\ref{thm:free-coefficient-1}(b2) the variable $X=\left\langle d\hat{\mathfrak{a}},F^u\right\rangle$ coincides, as a linear function of $F^u\in C^2$, with $X_{\eta,R}\left(\tau_vf^{(\eta)}\right)=\left\langle\eta^{\frac{3}{2}}\left(\tau_vf^{(\eta)}\right)_{\eta},F^u\right\rangle$ (Section~\ref{sec:lattices-tests-function}); hence $d\hat{\mathfrak{a}}=\eta^{\frac{3}{2}}\left(\tau_vf^{(\eta)}\right)_{\eta}$. As $y_i\in\eta\mathbb{Z}^3$, Lemma~\ref{lem:lattice-symmetries-cell} shows that $(y_i)_{*}$ commutes with $d$ and that $\left(\tau_{y_i}\varphi\right)_{\eta}=(y_i)_{*}\varphi_{\eta}$ for $\varphi\in\mathcal{S}$. Consequently $d\hat{\mathfrak{a}}_{c^{(i)}}=(y_i)_{*}d\hat{\mathfrak{a}}=\eta^{\frac{3}{2}}\left(\tau_{v+y_i}f^{(\eta)}\right)_{\eta}$, that is,
\[
X_{c^{(i)}}=\left\langle d\hat{\mathfrak{a}}_{c^{(i)}},F^u\right\rangle=X_{\eta,R}\left(\tau_{v+y_i}f^{(\eta)}\right).
\]
The cells $c^{(1)},\dots,c^{(J)}$ are pairwise distinct as soon as $R>2\max_i|y_i|$. In that case~\eqref{eq:free-coefficient-10}, applied to the set $S=\left\{c^{(1)},\dots,c^{(J)}\right\}$, and the fact that $E^{\mathrm{nat}}$ is unchanged by the lattice translation by $v$ (Theorem~\ref{thm:free-coefficient-1}(d5)) give for $m=2,4$
\begin{equation}\label{eq:free-coefficient-22}
\left|E^{\mathrm{nat}}\prod_{i=1}^J\mathrm{Pm}\left(X_{\eta,R}\left(\tau_{y_i}f^{(\eta)}\right);\boldsymbol{\sigma}^2,c_m(N,e,\beta)\right)\right|\le\mathsf{u}_m(N)^J.
\end{equation}

By Step~4 there is an index $i_0$ such that $N_i\ge N_0(e)$ and $T_i\ge T^{*}(e)$ for every $i\ge i_0$; for these $i$ the numbers $c_m(N_i,e,\beta)$ and $\mathsf{u}_m(N_i)$ are defined. By part~\ref{it:free-coefficient-3-2}(b) of Theorem~\ref{thm:free-coefficient-3} and by Steps~1 and~2, the sequences $\left(c_m(N_i,e,\beta)\right)_{i\ge i_0}$ and $\left(\mathsf{u}_m(N_i)\right)_{i\ge i_0}$, $m=2,4$, are bounded. Replace $(N_i,T_i)_{i\ge1}$ by a subsequence of $(N_i,T_i)_{i\ge i_0}$, again indexed by $i\ge1$, along which $c_m(N_i,e,\beta)\to c_m^{\infty}$ and $\mathsf{u}_m(N_i)\to\mathsf{u}_m$ for $m=2,4$. Then
\[
\mathsf{u}_4\le C\cdot e^{\sigma},\qquad\left(1+\left|c_2^{\infty}\right|\right)\cdot\mathsf{u}_2\le C\cdot e^{\sigma}.
\]
The expectation on the left of~\eqref{eq:free-coefficient-22} is a polynomial in $\boldsymbol{\sigma}^2$ and $c_m(N,e,\beta)$; its coefficients are joint moments of the variables $X_{\eta,R}\left(\tau_{y_i}f^{(\eta)}\right)$, $i=1,\dots,J$, each of order at most $mJ$. For a real test $g\in\mathcal{S}$ put $\mathbf{g}:=\eta^{\frac{3}{2}}g_{\eta}$; then $X_{\eta,R}(g)=\left\langle\mathbf{g},F^u\right\rangle$ and $\|\mathbf{g}\|=\left|g_{\eta}\right|_{\eta}$ (Section~\ref{sec:lattices-tests-function}). On the torus $\Lambda_{\eta,R}$, which is the torus $k=L^{T-T^{*}(e)}$ of Step~3, both laws satisfy the exponential bound of Step~3 for every real two-cochain (by~\eqref{it:free-coefficient-Q3} for $\mathrm{nat}=\mathrm{af}$ and by Theorem~\ref{thm:limit-vanishing-proca} for $\mathrm{nat}=\mathrm{av}$); for the cochain $\mathbf{g}$ it reads
\[
E^{\mathrm{nat}}\exp\left(X_{\eta,R}(g)\right)\le\exp\left(\frac{1}{2}\left|g_{\eta}\right|_{\eta}^2\right)\qquad(g\in\mathcal{S}\text{ real}).
\]
Under $E^{\mathrm{nat}}$ this is the hypothesis~\eqref{eq:continuum-infinite-5} of Proposition~\ref{prop:sub-gaussian-random}, for the real vector space of the real tests, the linear map $g\mapsto X_{\eta,R}(g)$ and the Hilbert seminorm $g\mapsto\left|g_{\eta}\right|_{\eta}$ (the seminorm of the positive semidefinite form $(g,g')\mapsto\langle g_{\eta},g'_{\eta}\rangle_{\eta}$). Part~(4) of that proposition therefore gives, for real tests $g_1,\dots,g_n$,
\[
E^{\mathrm{nat}}\prod_{a=1}^n\left|X_{\eta,R}(g_a)\right|\le2\left(2^n\cdot n!\right)^{\frac{1}{2}}\cdot\prod_{a=1}^n\left|g_{a,\eta}\right|_{\eta}.
\]
The tests $\tau_{y_i}f$ and $\tau_{y_i}f^{(\eta)}$, and with them their difference, are supported in the cube $y_i+Q_{\beta}$ (Section~\ref{sec:probes-discretisation} and Theorem~\ref{thm:free-coefficient-1}(b2)); its side is $D_{\beta}\le l_c-3\eta<R$, because $\beta$ fits $\left(N,T^{*}(e)\right)$ (Step~0). So each of them is supported in the period cube $y_i+y_{\beta}+\left(-\frac{R}{2},\frac{R}{2}\right)^3$, and Lemma~\ref{lem:norm-bounds-lattice} gives $\left|\varphi_{\eta}\right|_{\eta}\le\|\varphi\|_2$ for each of the three tests $\varphi=\tau_{y_i}f$, $\varphi=\tau_{y_i}f^{(\eta)}$ and $\varphi=\tau_{y_i}\left(f^{(\eta)}-f\right)$. Moreover $\left\|f^{(\eta)}-f\right\|_2\le\eta^2\cdot\|\beta\|_{H^4}$ (Theorem~\ref{thm:free-coefficient-1}(b3)), and $\|f\|_2\le s_f$, $\left\|f^{(\eta)}\right\|_2\le s_f$ (Theorem~\ref{thm:free-coefficient-1}(b5)). By multilinearity, replacing the factors $X_{\eta,R}\left(\tau_{y_i}f^{(\eta)}\right)$ one at a time by $X_{\eta,R}\left(\tau_{y_i}f\right)$, one finds that each of the moments in question, of order $n\le mJ$, differs from the corresponding joint moment of the variables $X_{\eta,R}\left(\tau_{y_i}f\right)$ by at most
\[
n\cdot2\left(2^n\cdot n!\right)^{\frac{1}{2}}\cdot s_f^{n-1}\cdot\eta^2\cdot\|\beta\|_{H^4},
\]
that is, by $O(\eta^2)$. Let $\eta_i:=L^{-N_i}$ and $R_i:=h\cdot L^{T_i}$. For every $i$ one has $N_i\ge N_0(e)$ and $T_i\ge T^{*}(e)$, and for all large $i$ also $R_i>2\max_{1\le j\le J}|y_j|$; along the subsequence the joint moments of the variables $X_{\eta_i,R_i}\left(\tau_{y_j}f\right)$ converge to those of the variables $X\left(\tau_{y_j}f\right)$ under $\mu$, by property~\eqref{it:main-theorem-P8} of Theorem~\ref{thm:main-theorem}, and $\boldsymbol{\sigma}^2\to\|f\|_2^2$ (Theorem~\ref{thm:free-coefficient-1}(b5)). Put
\[
G_m:=\mathrm{Pm}\left(X(f);\|f\|_2^2,c_m^{\infty}\right)\qquad(m=2,4),
\]
and let $U(y)$, $y\in\mathbb{R}^3$, be the unitary operators by which the translations act on $L^2(\mu)$ (property~\eqref{it:main-theorem-P5} of Theorem~\ref{thm:main-theorem}), $U(y)X(g)=X(\tau_yg)$. Letting $i\to\infty$ in~\eqref{eq:free-coefficient-22} one obtains, for $m=2,4$, every integer $J\ge1$ and all pairwise distinct $y_1,\dots,y_J\in\hat{\Gamma}_0$,
\begin{equation}\label{eq:free-coefficient-23}
\left|\int\mathrm{d}\mu\,\prod_{i=1}^JU(y_i)G_m\right|\le\mathsf{u}_m^J.
\end{equation}

Being a real polynomial in $X(f)$, the function $G_m$ belongs to $L^J(\mu)$ for every $J\ge1$ (property~\eqref{it:main-theorem-P1} of Theorem~\ref{thm:main-theorem}). The translations by the elements $\frac{R^{*}}{L}\cdot\boldsymbol{\zeta}$, $\boldsymbol{\zeta}\in\mathbb{Z}^3$, of $\hat{\Gamma}_0$ form an action of $\mathbb{Z}^3$ on $\mathcal{S}'$ by maps that preserve $\mu$ (property~\eqref{it:main-theorem-P5}). Let $P_0$ be the orthogonal projection of $L^2(\mu)$ onto the vectors $H$ with $U(y)H=H$ for all $y\in\hat{\Gamma}_0$. Lemma~\ref{lem:uniform-bound-invariant}, applied with $d=3$ to this action and with $\mathsf{u}_m$ in place of $\mathsf{u}$, gives $\left|P_0G_m\right|\le\mathsf{u}_m$ $\mu$-almost everywhere. Every vector of $L^2(\mu)$ that is fixed by all translations of $\mathbb{R}^3$ is fixed by the translations by the elements of $\hat{\Gamma}_0$, and the conditional expectation $E[\,\cdot\mid\mathcal{I}]$ is the orthogonal projection of $L^2(\mu)$ onto the former (Lemma~\ref{lem:averages-translations-2}(b)). Hence $E[G_m\mid\mathcal{I}]=E[P_0G_m\mid\mathcal{I}]$, and $\left|E[G_m\mid\mathcal{I}]\right|\le\mathsf{u}_m$ $\mu$-almost surely. Let $E_{\xi}$ denote the expectation under $\mu_{\xi}$. By Proposition~\ref{prop:existence-conditional-expectations}(b), $\left|E_{\xi}G_m\right|\le\mathsf{u}_m$ for $\nu$-almost every $\xi$ and $m=2,4$.

For $\nu$-almost every $\xi$ the component $\mu_{\xi}$ satisfies~\eqref{eq:fourth-cumulants-H1} (Corollary~\ref{cor:properties-almost-every}); hence $E_{\xi}X(f)=0$ (Lemma~\ref{lem:moment-bounds}(a)), and
\[
\kappa_4^{\xi}(f)=E_{\xi}X(f)^4-3\left(E_{\xi}X(f)^2\right)^2
\]
(Section~\ref{sec:cumulants-gaussian-laws}). The identity
\[
EX^4-3\left(EX^2\right)^2=EH_4(X)-3\left(EH_2(X)\right)^2
\]
of Step~2 of the proof of Theorem~\ref{thm:free-coefficient-2} (Section~\ref{sec:proof-cumulant-bounds}) holds for every real random variable $X$ with finite fourth moment, whatever its law, and for every value of the parameter $\boldsymbol{\sigma}^2$ in $H_2$ and $H_4$. Applied under $\mu_{\xi}$ to $X=X(f)$, with $\|f\|_2^2$ in place of $\boldsymbol{\sigma}^2$, and combined with the bounds $\left|E_{\xi}G_m\right|\le\mathsf{u}_m$ as in that step, it gives
\begin{align*}
\left|e^{-2}\left(C_{\xi}(f,f)-\|f\|_2^2\right)-c_2^{\infty}\right|&\le\mathsf{u}_2,\\
\left|e^{-4}\kappa_4^{\xi}(f)-\left(c_4^{\infty}-3\left(c_2^{\infty}\right)^2\right)\right|&\le\mathsf{u}_4+3\mathsf{u}_2\cdot\left(2\left|c_2^{\infty}\right|+\mathsf{u}_2\right)\le C\cdot e^{\sigma}.
\end{align*}
By part~\ref{it:free-coefficient-3-2}(b), (c) of Theorem~\ref{thm:free-coefficient-3},
\[
\left|c_2^{\infty}\right|\le C_B\cdot e^{-\mathfrak{e}_2},\qquad\left|c_4^{\infty}-3\left(c_2^{\infty}\right)^2-b(f)\right|\le C_B\cdot e\le C_B\cdot e^{\sigma}.
\]
This is~(vi), once $C_{\mathrm{kap}}$ is enlarged (the lower bound $\kappa_4^{\xi}(f)\ge\frac{1}{2}e^4\cdot b(f)>0$ and the assertion that $\nu$-almost every $\mu_{\xi}$ is not Gaussian then follow from~\eqref{eq:free-coefficient-4} as in Step~5).

For~(vii): the argument for~(vi) depends on the probe only through Theorem~\ref{thm:free-coefficient-1}(b), the inequality~\eqref{eq:free-coefficient-10}, Steps~0--2 and part~\ref{it:free-coefficient-3-2} of Theorem~\ref{thm:free-coefficient-3}, and all of these are available for each of the three probes $\beta_0$, $\beta_{+}$, $\beta_{-}$. The argument therefore applies to $\beta_{+}$ and to $\beta_{-}$, the placement $v$ being that of $\beta_{+}$, respectively $\beta_{-}$, and $g_{+}$, respectively $g_{-}$, taking the place of $f$. Consequently, for $\nu$-almost every $\xi$ the law $\mu_{\xi}$ satisfies~\eqref{eq:fourth-cumulants-H1}, is translation invariant (Corollary~\ref{cor:properties-almost-every}), and satisfies~\eqref{eq:fourth-cumulants-4g} at the three tests $g_0$, $g_{+}$, $g_{-}$ with the constant $C_{\mathrm{kap}}$. Theorem~\ref{thm:non-gaussian-components-2}(i), applied to $\mu_{\xi}$, gives $\mathrm{Ext}^{\xi}_{x_0}(f)\ge\frac{1}{2}e^4\cdot b(f)$; as in Step~5, the conditions of that theorem are~\eqref{eq:setup-1} and~\eqref{eq:free-coefficient-4}.

The constants are fixed in the following order. First $C_{\mathrm{kap}}$ is taken to be the maximum of the values that Steps~3, 4 and~6 require of it; none of these values depends on $e_2$, as long as $e_2$ satisfies~\eqref{eq:free-coefficient-1}--\eqref{eq:free-coefficient-3}. After that, $e_2$ is made smaller, if necessary, so that~\eqref{eq:free-coefficient-4} holds; this implies $C_{\mathrm{kap}}\cdot e_2^{\sigma}\le\frac{1}{2}b(f)$.\qed

\begin{remark}[Moments of all orders; a mixture with a Gaussian component]\label{rem:moments-all-orders}
Let $e\in(0,e_2]$, the limit point $\mu$, the test $f$, the lattice $\hat{\Gamma}_0=\frac{R^{*}}{L}\cdot\mathbb{Z}^3$, the numbers $c_2^{\infty}$, $c_4^{\infty}$ and the functions
\[
G_m=\mathrm{Pm}\left(X(f);\|f\|_2^2,c_m^{\infty}\right)\qquad(m=2,4)
\]
on $\mathcal{S}'$ be those of Step~6 of the proof of Theorem~\ref{thm:main-theorem} (Section~\ref{sec:proof-part-ii-2}). For $y\in\mathbb{R}^3$ the operator $U(y)$ acts on functions on $\mathcal{S}'$ as in Section~\ref{sec:euclidean-maps-actions}, by composition with a translation of $\mathcal{S}'$; thus $U(y)(F\cdot H)=\left(U(y)F\right)\cdot\left(U(y)H\right)$ and $U(y)X(g)=X(\tau_yg)$ for $g\in\mathcal{S}$, and on $L^2(\mu)$ these are the operators of property~\eqref{it:main-theorem-P5} of Theorem~\ref{thm:main-theorem}. For a law $\lambda$ on $\mathcal{S}'$ and a $\lambda$-integrable function $F$ on $\mathcal{S}'$ write $E_{\lambda}F:=\int\mathrm{d}\lambda\,F$; if $\lambda$ is invariant under the translations of $\mathbb{R}^3$, then $E_{\lambda}U(y)F=E_{\lambda}F$ for every $y\in\mathbb{R}^3$.

Let $m\in\{2,4\}$, let $J\ge1$ be an integer, and let $y_1,\dots,y_J\in\hat{\Gamma}_0$. Then
\[
\prod_{i=1}^JU(y_i)G_m=\prod_{i=1}^J\mathrm{Pm}\left(X(\tau_{y_i}f);\|f\|_2^2,c_m^{\infty}\right)
\]
is a real polynomial of degree $mJ$ in the variables $X(\tau_{y_1}f),\dots,X(\tau_{y_J}f)$. Expanding the product, one obtains its integral with respect to $\mu$ as a linear combination of joint moments of $\mu$ of orders at most $mJ$; since, for all real $s$ and $c$, the coefficient of $x^m$ in the polynomial $x\mapsto\mathrm{Pm}(x;s,c)$ is $e^{-m}$, the only term of order $mJ$ in this combination is
\[
e^{-mJ}\cdot\int\mathrm{d}\mu\,\prod_{i=1}^JX(\tau_{y_i}f)^m.
\]
For the action of $\mathbb{Z}^3$ on $\mathcal{S}'$ by the translations by the elements of $\hat{\Gamma}_0$ and for $G=G_m$, the inequalities assumed in Lemma~\ref{lem:uniform-bound-invariant} are the inequalities~\eqref{eq:free-coefficient-23}, $J\ge1$; they therefore involve joint moments of $\mu$ of order $mJ$ for every $J\ge1$, hence of arbitrarily high order. Proposition~\ref{prop:gaussian-components-four}(ii), on the other hand, concerns two-point and four-point functions.

Let $\mu'$ be a law on $\mathcal{S}'$ which satisfies~\eqref{eq:fourth-cumulants-H1}--\eqref{eq:fourth-cumulants-H3}, is ergodic in the sense of Definition~\ref{def:ergodicity-ergodic-decompositions}, and is not Gaussian in the sense of Section~\ref{sec:cumulants-gaussian-laws}; let $C':=S_2^{\mu'}$, and let $0\le p<1$. By Proposition~\ref{prop:gaussian-components-four}(ii), the form $C'$ satisfies the hypotheses of Lemma~\ref{lem:gaussian-laws-bounded}; the centred Gaussian law $\gamma_{C'}$ with covariance $C'$ given by part~(a) of that lemma satisfies~\eqref{eq:fourth-cumulants-H1}--\eqref{eq:fourth-cumulants-H3} and is ergodic; the law
\[
\mu^{(p)}:=(1-p)\cdot\mu'+p\cdot\gamma_{C'}
\]
satisfies~\eqref{eq:fourth-cumulants-H1}--\eqref{eq:fourth-cumulants-H3}; and $\mu^{(p)}$ has an ergodic decomposition whose two components are $\mu'$ and $\gamma_{C'}$, with the weights $1-p$ and $p$.

For an integer $k\ge1$, a real function $\Phi$ on $\hat{\Gamma}_0^k$ and fixed points $y_1,\dots,y_{k-1}\in\hat{\Gamma}_0$, let $\lim_{y_k\to\infty}\Phi(y_1,\dots,y_k)$ denote the limit, when it exists, of $\Phi(y_1,\dots,y_k)$ as $y_k$ ranges over $\hat{\Gamma}_0$ in such a way that all three coordinates of $y_k$ tend to infinity in absolute value, that is, $\min_l\left|(y_k)_l\right|\to\infty$; it is a function of $(y_1,\dots,y_{k-1})$. Let $m\in\{2,4\}$ and let $J\ge1$ be an integer. Then the iterated limit in which first $y_J$, then $y_{J-1}$, and so on, and finally $y_1$ tend to infinity in this sense exists, each inner limit existing for all values of the points that are kept fixed, and
\[
\lim_{y_1\to\infty}\lim_{y_2\to\infty}\cdots\lim_{y_J\to\infty}\int\mathrm{d}\mu^{(p)}\,\prod_{i=1}^JU(y_i)G_m=(1-p)\cdot\left(E_{\mu'}G_m\right)^J+p\cdot\left(E_{\gamma_{C'}}G_m\right)^J.
\]
In particular, since $t\mapsto|t|$ is continuous on $\mathbb{R}$, the left side of~\eqref{eq:free-coefficient-23}, formed with $\mu^{(p)}$ in place of $\mu$, tends in this iterated limit to
\[
\left|(1-p)\cdot\left(E_{\mu'}G_m\right)^J+p\cdot\left(E_{\gamma_{C'}}G_m\right)^J\right|.
\]
Moreover,
\[
E_{\gamma_{C'}}G_4=-c_4^{\infty}+3e^{-4}\cdot\left(C'(f,f)-\|f\|_2^2\right)^2.
\]

To prove the assertion on the iterated limit, let $\rho\in\left\{\mu',\gamma_{C'}\right\}$. The law $\rho$ satisfies~\eqref{eq:fourth-cumulants-H1}--\eqref{eq:fourth-cumulants-H3}, that is, the hypothesis~\eqref{eq:setup-H}; hence $\rho$ satisfies~\eqref{eq:setup-G}, and, by Section~\ref{sec:hypothesis-law}, it belongs to the class $\mathcal{C}$ of Section~\ref{sec:properties-laws}. Its ergodicity in the sense of Definition~\ref{def:ergodicity-ergodic-decompositions} is ergodicity under the translations of $\mathbb{R}^3$ in the sense of Section~\ref{sec:invariant-sets-invariant}, both meaning that the constants are the only elements of $L^2(\rho)$ invariant under all translations; so Theorem~\ref{thm:invariant-vectors-clustering}(iii) applies to $\rho$. By Theorem~\ref{thm:invariant-vectors-clustering}(iv), every polynomial in the variables $X(g)$, $g\in\mathcal{S}$, lies in $L^2(\rho)$; in particular, for $k\ge1$ and $y_1,\dots,y_k\in\hat{\Gamma}_0$ the real function $\prod_{i=1}^kU(y_i)G_m$ lies in $L^2(\rho)$ and is $\rho$-integrable. For $0\le k\le J$ and $y_1,\dots,y_k\in\hat{\Gamma}_0$ put
\[
\Phi^{\rho}_k(y_1,\dots,y_k):=E_{\rho}\left[\prod_{i=1}^kU(y_i)G_m\right]\cdot\left(E_{\rho}G_m\right)^{J-k},
\]
the empty product being $1$; thus $\Phi^{\rho}_0=\left(E_{\rho}G_m\right)^J$. Let $2\le k\le J$, fix $y_1,\dots,y_{k-1}\in\hat{\Gamma}_0$, and put $F:=\prod_{i=1}^{k-1}U(y_i)G_m$. Theorem~\ref{thm:invariant-vectors-clustering}(iii)1, applied to $\rho$ and to the real functions $F$ and $G_m$ of $L^2(\rho)$, gives
\[
E_{\rho}\left[F\cdot U(a)G_m\right]\longrightarrow E_{\rho}F\cdot E_{\rho}G_m\qquad\text{as }a\in\mathbb{R}^3,\ \min_l|a_l|\to\infty.
\]
For $a=y_k\in\hat{\Gamma}_0$ the left side equals $E_{\rho}\left[\prod_{i=1}^kU(y_i)G_m\right]$; multiplication by $\left(E_{\rho}G_m\right)^{J-k}$ yields
\[
\lim_{y_k\to\infty}\Phi^{\rho}_k(y_1,\dots,y_k)=\Phi^{\rho}_{k-1}(y_1,\dots,y_{k-1}).
\]
Since $\rho$ is invariant under the translations, $\Phi^{\rho}_1(y_1)=E_{\rho}\left[U(y_1)G_m\right]\cdot\left(E_{\rho}G_m\right)^{J-1}=\Phi^{\rho}_0$ for every $y_1\in\hat{\Gamma}_0$, and therefore also $\lim_{y_1\to\infty}\Phi^{\rho}_1(y_1)=\Phi^{\rho}_0$. Finally, for $y_1,\dots,y_J\in\hat{\Gamma}_0$ the function $\prod_{i=1}^JU(y_i)G_m$ is integrable with respect to $\mu'$ and to $\gamma_{C'}$, hence with respect to $\mu^{(p)}$, and
\[
\int\mathrm{d}\mu^{(p)}\,\prod_{i=1}^JU(y_i)G_m=(1-p)\cdot\Phi^{\mu'}_J(y_1,\dots,y_J)+p\cdot\Phi^{\gamma_{C'}}_J(y_1,\dots,y_J).
\]
Taking the limits in $y_J,y_{J-1},\dots,y_1$ in this order, and using the preceding relations for $\rho=\mu'$ and for $\rho=\gamma_{C'}$ at each stage, one obtains the iterated limit $(1-p)\cdot\Phi^{\mu'}_0+p\cdot\Phi^{\gamma_{C'}}_0$, which is the asserted value.

For the value of $E_{\gamma_{C'}}G_4$: under $\gamma_{C'}$ the variable $X(f)$ is centred and normal with variance $C'(f,f)$ (Lemma~\ref{lem:gaussian-laws-bounded}(a)); hence $E_{\gamma_{C'}}X(f)^2=C'(f,f)$ and $E_{\gamma_{C'}}X(f)^4=3C'(f,f)^2$. The definition of $\mathrm{Pm}$ for $m=4$ gives
\[
E_{\gamma_{C'}}G_4=e^{-4}\cdot\left(3C'(f,f)^2-6\|f\|_2^2\cdot C'(f,f)+3\|f\|_2^4\right)-c_4^{\infty}=-c_4^{\infty}+3e^{-4}\cdot\left(C'(f,f)-\|f\|_2^2\right)^2.
\]
\end{remark}

\subsubsection{Proofs of the corollaries of the main theorem}\label{sec:proofs-corollaries-main}

\begin{proof}[Proof of Corollary~\ref{cor:ergodic-components}]
Let $e\in(0,e_2]$, let $\mathrm{nat}\in\{\mathrm{af},\mathrm{av}\}$, and let $\mu$ be a limit point of $\left(P^{\mathrm{nat}}_{N,T}\right)$ as in Theorem~\ref{thm:main-theorem}. By Theorem~\ref{thm:main-theorem}(i), $\mu$ has the properties~\eqref{it:main-theorem-P1}--\eqref{it:main-theorem-P8}, and the assertions of part~(c) of that theorem hold for $\mu$; in particular $\mu$ satisfies~\eqref{eq:setup-H}. Let
\[
\mu=\int\mathrm{d}\nu(\xi)\,\mu_{\xi}
\]
be the ergodic decomposition of $\mu$ under $\mathbb{R}^3$ (Section~\ref{sec:laws-2}).

\emph{(a)} Let $\Xi_1$ be a set of full $\nu$-measure such that, for every $\xi\in\Xi_1$, the component $\mu_{\xi}$ satisfies $\mu_{\xi}\left(\left\{F\in\mathcal{S}':dF=0\right\}\right)=1$, is invariant and ergodic under the translations of $\mathbb{R}^3$, is invariant under $N_3$, satisfies~\eqref{eq:setup-G}, is reflection positive at each plane $\{x_j=c\}$, $j\in\{1,2,3\}$, $c\in\mathbb{R}$, and has the properties 1--4 of Theorem~\ref{thm:main-theorem}(c); such a set exists, because by Theorem~\ref{thm:main-theorem}(c) these assertions hold together outside one $\nu$-null set. Let $\Xi_2$ be a set of full $\nu$-measure on which the assertions of Theorem~\ref{thm:main-theorem}(vi) hold. The intersection $\Xi_1\cap\Xi_2$ has full $\nu$-measure; let $\xi\in\Xi_1\cap\Xi_2$. As $\xi\in\Xi_1$, $\mu_{\xi}$ is a law on $\mathcal{S}'(\mathbb{R}^3;\Lambda^2)$ which is concentrated on the closed two-forms, is invariant and ergodic under $\mathbb{R}^3$, is invariant under $N_3$, is reflection positive at every coordinate plane $\{x_j=c\}$, satisfies~\eqref{eq:setup-G}, and, by property~4 of Theorem~\ref{thm:main-theorem}(c), clusters in the mean and as $\min_l|a_l|\to\infty$. By Theorem~\ref{thm:main-theorem}(vi), one has $\kappa_4^{\xi}(f)\ge e^4\cdot b(f)/2>0$, and $\mu_{\xi}$ is not Gaussian. It remains to prove that $\mu_{\xi}$ has the properties~\eqref{it:main-theorem-P1} and~\eqref{it:main-theorem-P2} with the norm $\|\cdot\|_2$, and the property~\eqref{it:main-theorem-P7}.

By~\eqref{eq:setup-G} and Lemma~\ref{lem:consequences-gaussian-bound}(ii), the map $f\mapsto X(f)$, $f\in\mathcal{S}$, extends uniquely to a linear contraction $X\colon L^2(\mathbb{R}^3;\Lambda^2)\to L^2(\mu_{\xi})$, and the extension still obeys the bound~\eqref{eq:setup-G}:
\[
\int\mathrm{d}\mu_{\xi}\,\exp X(g)\le\exp\left(\frac{1}{2}\|g\|_2^2\right)\qquad\text{for every real }g\in L^2(\mathbb{R}^3;\Lambda^2).
\]
Hence the probability space $\left(\mathcal{S}',\mu_{\xi}\right)$, the real vector space $L^2(\mathbb{R}^3;\Lambda^2)$ with the Hilbert norm $|\cdot|:=\|\cdot\|_2$, and the linear map $X$ satisfy the hypothesis~\eqref{eq:continuum-infinite-5} of Proposition~\ref{prop:sub-gaussian-random}. The conclusions (1), (4), (5) and (7) of that proposition in this setting, together with the extension of $X$ just obtained, are the assertions of~\eqref{it:main-theorem-P1} for $\mu_{\xi}$ in which the bound by $C_0(f,f)$ is replaced by the bound by $\|f\|_2^2$, and the first two inequalities of~\eqref{it:main-theorem-P2} for $\mu_{\xi}$.

The last assertion of~\eqref{it:main-theorem-P2}, namely that the moment functional $S_n$ of $\mu_{\xi}$ extends uniquely to a continuous linear functional on $\mathscr{S}_n$ which satisfies
\[
\left|S_n(\Psi)\right|\le2\cdot66^n\cdot(n!)^{\frac{1}{2}}\cdot q_{6n}(\Psi)\qquad\left(\Psi\in\mathscr{S}_n\right),
\]
is obtained by the proof of Theorem~\ref{thm:moments-limit-point}(2). That proof uses no property of the law other than the inequality of Proposition~\ref{prop:sub-gaussian-random}(4), which holds for $\mu_{\xi}$ with $|\cdot|=\|\cdot\|_2$ by the preceding paragraph; apart from it, the proof uses Section~\ref{sec:hermite-sobolev-spaces}, by which every $F\in\mathcal{S}'$ lies in $H_{-s}$ for some $s$ and the finite Hermite sums are dense in $\mathscr{S}_n$, and Lemma~\ref{lem:hermite-coefficients-tensor}. It therefore applies to $\mu_{\xi}$.

Finally let $j\in\{1,2,3\}$. The law $\mu_{\xi}$ is invariant under the translations of $\mathbb{R}^3$ and under the reflection $\theta_j\in N_3$ of Section~\ref{sec:euclidean-maps-actions}, and it is reflection positive at the plane $\{x_j=0\}$; so it satisfies the hypotheses of Proposition~\ref{prop:reconstruction-direction} for the direction $e_j$. That proposition, which rests on this invariance and on this reflection positivity, provides a Hilbert space $\mathscr{H}_j$, a self-adjoint operator $\mathsf{H}_j\ge0$ with the semigroup $\left(\exp\left(-t\cdot\mathsf{H}_j\right)\right)_{t\ge0}$, a unit vector $\Omega_j$ with $\mathsf{H}_j\Omega_j=0$, and unitary groups of spatial translations which commute with the semigroup. This is the reconstruction~\eqref{it:main-theorem-P7} for $\mu_{\xi}$.

\emph{(b)} Let $j\in\{1,2,3\}$. For a law $\rho$ on $\mathcal{S}'$ let, as in Section~\ref{sec:test-forms-distributions}, $\rho(F)$ be the integral with respect to $\rho$ of a $\rho$-integrable function $F$ on $\mathcal{S}'$, let $L^2(\rho)$ be the complex Hilbert space of the classes of measurable complex functions on $\mathcal{S}'$ that are square integrable with respect to $\rho$, and, for a $\sigma$-algebra $\mathcal{G}$ contained in the cylinder $\sigma$-algebra, let $L^2(\mathcal{G},\rho)$ be the closed subspace of $L^2(\rho)$ formed by the classes which contain a $\mathcal{G}$-measurable function. Let $\mathcal{I}^j$ be the $\sigma$-algebra of the cylinder sets $A$ with $\tau_aA=A$ for every $a\in\mathbb{R}e_j$ (Section~\ref{sec:invariant-sets-invariant} with $J=\{j\}$), let $U(a)$, $a\in\mathbb{R}^3$, be the translation operators and $\theta_j$ the reflection which changes the sign of $x_j$, both as in Section~\ref{sec:euclidean-maps-actions}, and let $P^j_{\mu}$ be the orthogonal projection of $L^2(\mu)$ onto the closed subspace $\mathcal{E}^j(\mu)$ of the $F\in L^2(\mu)$ with $U(se_j)F=F$ for every $s\in\mathbb{R}$ (Section~\ref{sec:invariant-sets-invariant}). Throughout, $E[\,\cdot\mid\mathcal{I}^j]$ denotes the conditional expectation under $\mu$. Let $(\omega,A)\mapsto\mu_{\omega}(A)$ be the probability kernel from $(\mathcal{S}',\mathcal{I}^j)$ to $\mathcal{S}'$ with its cylinder $\sigma$-algebra which Theorem~\ref{thm:ergodic-decomposition} provides for the law $\mu$ and $J=\{j\}$. By Theorem~\ref{thm:ergodic-decomposition}(a), the measure $\nu^j$ is the image of $\mu$ under $\omega\mapsto\mu_{\omega}$, and
\[
\mu_{\omega}(A)=E\left[\mathbf{1}_A\mid\mathcal{I}^j\right](\omega)\qquad\text{for }\mu\text{-almost every }\omega,
\]
for every cylinder set $A$.

Since $\mu$ satisfies~\eqref{eq:setup-H}, it lies in the class $\mathcal{C}$ of Section~\ref{sec:properties-laws} and satisfies~\eqref{eq:setup-G} (Section~\ref{sec:hypothesis-law}). Theorem~\ref{thm:decomposition-along-coordinate} therefore applies to $\mu$: for $\nu^j$-almost every $\rho$, the law $\rho$ is invariant under the translations along $e_j$ and under $\theta_j$, is reflection positive at every plane $\{x_j=c\}$, satisfies~\eqref{eq:setup-G}, has $\ker\mathsf{H}_j(\rho)=\mathbb{C}\,\Omega_j$, and
\[
\rho\left(\overline{F}\cdot U(se_j)G\right)\to\overline{\rho(F)}\cdot\rho(G)\qquad\text{as }|s|\to\infty,\ \text{for all }F,G\in L^2(\rho).
\]
Moreover $\mu\left(\{\omega\in\mathcal{S}':d\omega=0\}\right)=1$ by~\eqref{it:main-theorem-P4}; this set is a cylinder set (Section~\ref{sec:closed-distributions}), and Theorem~\ref{thm:ergodic-decomposition}(d) gives $\rho\left(\{\omega\in\mathcal{S}':d\omega=0\}\right)=1$ for $\nu^j$-almost every $\rho$. These are the assertions of~(b) other than the bound on $\kappa_4^{\rho}(f)$ and its consequence.

For this bound, let $m\in\{2,4\}$, and let $G_m$, $\mathsf{u}_m$ and $c_m^{\infty}$ be the function and the numbers of Step~6 of the proof of Theorem~\ref{thm:main-theorem} (Section~\ref{sec:proof-part-ii-2}), formed for the limit point $\mu$. The function $G_m$ is a real polynomial in $X(f)$, hence lies in $L^J(\mu)$ for every $J\ge1$ (by~\eqref{it:main-theorem-P1}), and~\eqref{eq:free-coefficient-23} holds for all pairwise distinct points $y_1,\dots,y_J$ of the lattice $\hat{\Gamma}_0=(R^{*}/L)\mathbb{Z}^3$ of that step, $R^{*}=R^{*}(e)$. As $k\cdot(R^{*}/L)\,e_j\in\hat{\Gamma}_0$ for every $k\in\mathbb{Z}$, the inequality~\eqref{eq:free-coefficient-23} contains the case
\[
\left|\int\mathrm{d}\mu\,\prod_{i=1}^JU\left(k_i\cdot\frac{R^{*}}{L}\,e_j\right)G_m\right|\le\mathsf{u}_m^J\qquad\text{for every }J\ge1\text{ and all pairwise distinct }k_1,\dots,k_J\in\mathbb{Z}.
\]
The translations by the vectors $k\cdot(R^{*}/L)\,e_j$, $k\in\mathbb{Z}$, form an action of $\mathbb{Z}$ on $\mathcal{S}'$ by maps which preserve $\mu$ (by~\eqref{it:main-theorem-P5}), and for this action the operators $U(k)$ of Lemma~\ref{lem:uniform-bound-invariant} are the operators $U\left(k\cdot(R^{*}/L)\,e_j\right)$. Lemma~\ref{lem:uniform-bound-invariant}, applied with $d=1$ to this action, to $G=G_m$ and to $\mathsf{u}=\mathsf{u}_m$, shows that the orthogonal projection $P_{0,j}G_m$ of $G_m$ onto the closed subspace
\[
\left\{H\in L^2(\mu):U\left(\frac{R^{*}}{L}\,e_j\right)H=H\right\}
\]
satisfies $\left|P_{0,j}G_m\right|\le\mathsf{u}_m$ $\mu$-almost everywhere.

By Lemma~\ref{lem:averages-translations}(iv), one has $\mathcal{E}^j(\mu)=L^2(\mathcal{I}^j,\mu)$ and $P^j_{\mu}=E[\,\cdot\mid\mathcal{I}^j]$; this is also the projection $P_j$ of Lemma~\ref{lem:averages-translations-2}(d). A vector which is invariant under every $U(se_j)$, $s\in\mathbb{R}$, is invariant under $U\left((R^{*}/L)\,e_j\right)$. Hence the range of $P^j_{\mu}$ is contained in the range of $P_{0,j}$, so that $P_{0,j}P^j_{\mu}=P^j_{\mu}$ and, taking adjoints, $P^j_{\mu}P_{0,j}=P^j_{\mu}$. It follows that
\[
E\left[G_m\mid\mathcal{I}^j\right]=E\left[P_{0,j}G_m\mid\mathcal{I}^j\right],\qquad\text{and therefore}\qquad\left|E\left[G_m\mid\mathcal{I}^j\right]\right|\le\mathsf{u}_m\quad\mu\text{-almost surely}.
\]
By~\eqref{eq:ergodic-components-1}, $\mu_{\omega}(G_m)=E\left[G_m\mid\mathcal{I}^j\right](\omega)$ for $\mu$-almost every $\omega$; moreover $\mu_{\omega}\left(|G_m|\right)<\infty$ for $\mu$-almost every $\omega$, since the second identity of~\eqref{eq:ergodic-components-1}, applied to $|G_m|$, gives $\int\mathrm{d}\mu(\omega)\,\mu_{\omega}\left(|G_m|\right)=\mu\left(|G_m|\right)<\infty$. The measure $\nu^j$ is a probability measure on the set $\mathcal{M}$ of the laws on $\mathcal{S}'$, endowed with the $\sigma$-algebra $\mathcal{B}_{\mathcal{M}}$ generated by the maps $\rho\mapsto\rho(A)$, $A$ a cylinder set (Section~\ref{sec:test-forms-distributions}), and $\nu^j(B)=\mu\left(\left\{\omega\in\mathcal{S}':\mu_{\omega}\in B\right\}\right)$ for every $B\in\mathcal{B}_{\mathcal{M}}$. Let $\Phi$ be a real function on $\mathcal{S}'$ which is measurable for the cylinder $\sigma$-algebra, and let $\Phi^{+}=\max\{\Phi,0\}$ and $\Phi^{-}=\max\{-\Phi,0\}$. By the first paragraph of the proof of Lemma~\ref{lem:measurability-defining-properties}, the maps $\rho\mapsto\rho\left(\Phi^{+}\right)$ and $\rho\mapsto\rho\left(\Phi^{-}\right)$ from $\mathcal{M}$ to $[0,\infty]$ are $\mathcal{B}_{\mathcal{M}}$-measurable. Hence the set of the $\rho\in\mathcal{M}$ with $\rho\left(|\Phi|\right)=\rho\left(\Phi^{+}\right)+\rho\left(\Phi^{-}\right)<\infty$ belongs to $\mathcal{B}_{\mathcal{M}}$, and on this set $\rho(\Phi)=\rho\left(\Phi^{+}\right)-\rho\left(\Phi^{-}\right)$ is a $\mathcal{B}_{\mathcal{M}}$-measurable function of $\rho$. For $\Phi=G_m$, a real polynomial in $X(f)$, it follows that the set
\[
B_m:=\left\{\rho\in\mathcal{M}:\rho\left(|G_m|\right)<\infty\text{ and }\left|\rho(G_m)\right|\le\mathsf{u}_m\right\}
\]
belongs to $\mathcal{B}_{\mathcal{M}}$; by what precedes, $\mu_{\omega}\in B_m$ for $\mu$-almost every $\omega$. As $\nu^j$ is the image of $\mu$ under $\omega\mapsto\mu_{\omega}$, one obtains $\nu^j(B_m)=1$, that is, $\left|\rho(G_m)\right|\le\mathsf{u}_m$ for $\nu^j$-almost every $\rho$, for $m=2$ and for $m=4$. Moreover, Theorem~\ref{thm:invariant-vectors-clustering}(iv), applied to $\mu\in\mathcal{C}$, which satisfies~\eqref{eq:setup-G}, gives $P^j_{\mu}X(f)=0$, that is, $E\left[X(f)\mid\mathcal{I}^j\right]=0$; by~\eqref{eq:ergodic-components-1} again, $\mu_{\omega}\left(X(f)\right)=0$ for $\mu$-almost every $\omega$, and, by the second identity of~\eqref{eq:ergodic-components-1} applied to $|X(f)|$, which is $\mu$-integrable by~\eqref{it:main-theorem-P1}, $\mu_{\omega}\left(|X(f)|\right)<\infty$ for $\mu$-almost every $\omega$. The argument just given, with $\Phi=X(f)$, shows that the set
\[
B_0:=\left\{\rho\in\mathcal{M}:\rho\left(|X(f)|\right)<\infty\text{ and }\rho\left(X(f)\right)=0\right\}
\]
belongs to $\mathcal{B}_{\mathcal{M}}$ and contains $\mu_{\omega}$ for $\mu$-almost every $\omega$; so $\nu^j(B_0)=1$, that is, $\rho\left(X(f)\right)=0$ for $\nu^j$-almost every $\rho$.

Let now $\rho$ be a law which satisfies~\eqref{eq:setup-G}, $\rho\left(X(f)\right)=0$ and $\left|\rho(G_m)\right|\le\mathsf{u}_m$ for $m=2,4$; by what precedes, $\nu^j$-almost every $\rho$ is such a law. Under $\rho$ the variable $X(f)$ has finite moments of all orders (Lemma~\ref{lem:consequences-gaussian-bound}(i)) and is centred, so that its fourth cumulant is
\[
\kappa_4^{\rho}(f)=\rho\left(X(f)^4\right)-3\rho\left(X(f)^2\right)^2
\]
(Section~\ref{sec:cumulants-gaussian-laws}). Let $H_2(x)=x^2-\|f\|_2^2$ and $H_4(x)=x^4-6\|f\|_2^2\cdot x^2+3\|f\|_2^4$ be the polynomials of Section~\ref{sec:probes-discretisation} with $\|f\|_2^2$ in place of $\boldsymbol{\sigma}^2$; by the definition of $G_m$ in Step~6, $G_m=e^{-m}\cdot H_m\left(X(f)\right)-c_m^{\infty}$. The identity
\[
\rho\left(X(f)^4\right)-3\rho\left(X(f)^2\right)^2=\rho\left(H_4\left(X(f)\right)\right)-3\rho\left(H_2\left(X(f)\right)\right)^2
\]
of Step~2 of the proof of Theorem~\ref{thm:free-coefficient-2} (Section~\ref{sec:proof-cumulant-bounds}) holds for every law under which the variable has a finite fourth moment and for every value of the parameter $\boldsymbol{\sigma}^2$. Inserting $\rho\left(H_m\left(X(f)\right)\right)=e^m\cdot\left(\rho(G_m)+c_m^{\infty}\right)$, one finds
\[
e^{-4}\cdot\kappa_4^{\rho}(f)-\left(c_4^{\infty}-3\left(c_2^{\infty}\right)^2\right)=\rho(G_4)-6c_2^{\infty}\cdot\rho(G_2)-3\rho(G_2)^2,
\]
and hence
\[
\left|e^{-4}\cdot\kappa_4^{\rho}(f)-\left(c_4^{\infty}-3\left(c_2^{\infty}\right)^2\right)\right|\le\mathsf{u}_4+3\mathsf{u}_2\cdot\left(2\left|c_2^{\infty}\right|+\mathsf{u}_2\right).
\]
As at the end of Step~6, the bounds $\mathsf{u}_4\le C\cdot e^{\sigma}$ and $\left(1+\left|c_2^{\infty}\right|\right)\cdot\mathsf{u}_2\le C\cdot e^{\sigma}$ established in that step show that the right-hand side is at most $C\cdot e^{\sigma}$, and part~\ref{it:free-coefficient-3-2}(c) of Theorem~\ref{thm:free-coefficient-3}, in the limit along the sequence of cutoffs of Step~6 (on which $\eta\to0$), gives
\[
\left|c_4^{\infty}-3\left(c_2^{\infty}\right)^2-b(f)\right|\le C_B\cdot e\le C_B\cdot e^{\sigma},
\]
with the constant $C_B$ of the proof of Theorem~\ref{thm:main-theorem}. The constant $C_{\mathrm{kap}}$ was chosen in Step~6 so large that these two bounds yield the first inequality of assertion~(vi) of Theorem~\ref{thm:main-theorem}; with $\rho$ in place of $\mu_{\xi}$ they yield
\[
\left|\kappa_4^{\rho}(f)-e^4\cdot b(f)\right|\le C_{\mathrm{kap}}\cdot e^{4+\sigma}.
\]
Finally, $e\le e_2$ and~\eqref{eq:free-coefficient-4} give $C_{\mathrm{kap}}\cdot e^{\sigma}\le C_{\mathrm{kap}}\cdot e_2^{\sigma}\le b(f)/2$; hence $\kappa_4^{\rho}(f)\ge e^4\cdot b(f)/2>0$. A Gaussian law which satisfies~\eqref{eq:setup-G} has vanishing fourth cumulants (Section~\ref{sec:cumulants-gaussian-laws}); so $\rho$ is not Gaussian.
\end{proof}

\begin{proof}[Proof of Corollary~\ref{cor:zero-harmonic-law}]
Let $N$ be an integer with $\boldsymbol{\mu}_N>0$. For an integer $T\ge0$ such that $(N,T)$ is a cutoff let $\eta=L^{-N}$, $R=h\cdot L^T$, let $n=h\cdot L^{T+N}$ be the number of sites per side of the torus $\Lambda_{\eta,R}$, and put
\[
\tau_{*}:=n\cdot\operatorname{arcsinh}\left(\frac{\boldsymbol{\mu}_N}{12}\right),\qquad\boldsymbol{\delta}_n:=110\,n^3\cdot\exp\left(-\tau_{*}\right).
\]
By the identification in Section~\ref{sec:wilson-determinant-lattice}, $P^{\mathrm{av}}_{N,T}$ is the law of $F^{\sharp}$ under the averaged law $\bar{\mathsf{P}}$ of Section~\ref{sec:wilson-operator-lattice} on $\Lambda_{\eta,R}$, with $q=e\cdot\eta^{\frac{1}{2}}$ and the lattice mass $\boldsymbol{\mu}_N$; and $P^0_{N,T}$ is the law of $F^{\sharp}$ under the zero-harmonic law $\mathsf{P}^0$ of Section~\ref{sec:wilson-operator-lattice} with the same $q$ and $\boldsymbol{\mu}_N$.

Assume $\tau_{*}\ge1$. Theorem~\ref{thm:insensitivity-twist}, applied with $q\ge0$ and the lattice mass $\boldsymbol{\mu}_N>0$, shows that the total variation distance of $\mathsf{P}^0$ and $\bar{\mathsf{P}}$ is at most $\exp\left(4\boldsymbol{\delta}_n\right)-1$: for every Borel function $\Phi$ on $C^2(\Lambda_{\eta,R})$ with $|\Phi|\le1$, the integrals of $\Phi$ with respect to $\mathsf{P}^0$ and to $\bar{\mathsf{P}}$ differ by at most $\exp\left(4\boldsymbol{\delta}_n\right)-1$. The map $F\mapsto F^{\sharp}$ is linear from the finite-dimensional space $C^2(\Lambda_{\eta,R})$ into $H_{-3}\subset H_{-4}$. Hence, for every bounded continuous function $\Phi\colon H_{-4}\to\mathbb{R}$, the function $F\mapsto\Phi\left(F^{\sharp}\right)$ is a Borel function on $C^2(\Lambda_{\eta,R})$ bounded by $\sup|\Phi|$, and
\[
\left|\int\mathrm{d}P^0_{N,T}\,\Phi-\int\mathrm{d}P^{\mathrm{av}}_{N,T}\,\Phi\right|\le\left(\exp\left(4\boldsymbol{\delta}_n\right)-1\right)\cdot\sup_{H_{-4}}|\Phi|.
\]
Let $T\to\infty$ with $N$ fixed. Then $n\to\infty$; since $\operatorname{arcsinh}\left(\boldsymbol{\mu}_N/12\right)>0$, it follows that $\tau_{*}\to\infty$ and $\boldsymbol{\delta}_n\to0$. So the condition $\tau_{*}\ge1$ holds for all large $T$, and the right-hand side of the last inequality tends to $0$. Consequently, along every sequence $T_k\to\infty$ the laws $P^0_{N,T_k}$ converge weakly on $H_{-4}$ if and only if the laws $P^{\mathrm{av}}_{N,T_k}$ do, and then both have the same limit: as $T\to\infty$, the two families have the same weak limit points on $H_{-4}$.

Now let $\mu^{(N)}$ be such a limit point for infinitely many $N$ with $\boldsymbol{\mu}_N>0$, and let $\mu$ be a weak limit on $H_{-4}$ of $\mu^{(N_i)}$ along a sequence $N_i\to\infty$. Each $\mu^{(N)}$ is the weak limit on $H_{-4}$ of $P^{\mathrm{av}}_{N,T_j(N)}$ along a sequence $T_j(N)\to\infty$. By Remark~\ref{rem:repeated-limits}, $\mu$ is a limit point of $\left(P^{\mathrm{av}}_{N,T}\right)$ in the sense of Section~\ref{sec:wilson-determinant-lattice}. Theorem~\ref{thm:main-theorem} applies to this limit point and gives it the properties (i) and (iii)--(vii) of that theorem.
\end{proof}

\subsection{Examples and counterexamples}\label{sec:examples-counterexamples}

\subsubsection{Soft laws and four operations on them}\label{sec:soft-laws-four}

\begin{definition}
A law $\mu$ on $\mathcal{S}'$ is called \emph{soft} if it satisfies the hypothesis~\eqref{eq:setup-H}, if it is concentrated on the set of the closed distributions, which belongs to the cylinder $\sigma$-algebra (Section~\ref{sec:closed-distributions}), that is,
\[
\mu\left(\left\{F\in\mathcal{S}':dF=0\right\}\right)=1,
\]
and if it is invariant under the map $F\mapsto-F$ of $\mathcal{S}'$.
\end{definition}

Let $q$ be a symmetric positive semidefinite bilinear form on the real space $\mathcal{S}$ such that $q(f,f)\le\|f\|_2^2$ for every $f\in\mathcal{S}$, and write $q(f):=q(f,f)$. By Lemma~\ref{lem:gaussian-laws-bounded}(a) there is exactly one law $\gamma_q$ on $\mathcal{S}'$ with
\[
\int\mathrm{d}\gamma_q\,\exp\left(iX(f)\right)=\exp\left(-\frac{1}{2}q(f)\right)\qquad(f\in\mathcal{S});
\]
it is the centred Gaussian law with covariance $q$. As $0\le C_0(f,f)\le\|f\|_2^2$ (Lemma~\ref{lem:free-transverse-covariance}(a)), this applies to the free Maxwell covariance $C_0$ of Section~\ref{sec:lattices-tests-function} and defines the law $\gamma_{C_0}$. For $\pi\in S_3$ let $u^{(\pi)}$ be the constant triple and $\gamma^{(\pi)}$ the law of Section~\ref{sec:permuted-diagonal-rescalings}, and let $q^{(\pi)}:=q_{u^{(\pi)}}$, where $q_u$ is the quadratic form of Lemma~\ref{lem:mixtures-changes-variables}; thus $\gamma^{(\pi)}$ is the centred Gaussian law with covariance $q^{(\pi)}$. Let $m_A=\frac{1}{6}\sum_{\pi\in S_3}\gamma^{(\pi)}$ and $m_B$ be the laws of Sections~\ref{sec:permuted-diagonal-rescalings} and~\ref{sec:free-maxwell-law-2}, and put
\[
q_A:=\frac{1}{6}\sum_{\pi\in S_3}q^{(\pi)}.
\]

\begin{lemma}[Properties of soft laws]\label{lem:properties-soft-laws}
Every soft law has the properties~\eqref{it:main-theorem-P1}--\eqref{it:main-theorem-P7} and the property~(c) of Theorem~\ref{thm:main-theorem}.
\end{lemma}

\begin{proof}
Let $\mu$ be a soft law. It has the four properties~\eqref{eq:continuum-infinite-L1}--\eqref{eq:continuum-infinite-L4} of Section~\ref{sec:properties-limit-points}. Indeed, \eqref{eq:continuum-infinite-L1} is the bound~\eqref{eq:setup-G}. For real $g\in C_c^{\infty}(\mathbb{R}^3)$ one has $F(d^{*}g)=(dF)(g)$, and this number vanishes for $\mu$-almost every $F$, because $\mu\left(\{dF=0\}\right)=1$; this is~\eqref{eq:continuum-infinite-L2}. The condition~\eqref{eq:continuum-infinite-L3} on the characteristic functional of $\mu$ follows from the invariance of $\mu$ under the translations of $\mathbb{R}^3$ and under $O_{\mathrm{cub}}$. Finally, \eqref{eq:continuum-infinite-L4} is the condition~\eqref{eq:setup-RPj} of Section~\ref{sec:properties-laws} for $j=1,2,3$. The bound~\eqref{eq:setup-G}, this invariance and the conditions~\eqref{eq:setup-RPj} are part of the hypothesis~\eqref{eq:setup-H}, which $\mu$ satisfies. Moreover, the characteristic functional of $\mu$ is even, since $\mu$ is invariant under $F\mapsto-F$.

By Proposition~\ref{prop:laws-four-properties}, therefore, the assertions (1), (3), (4), (5), (6), (7) and~(8) of Theorem~\ref{thm:continuum-infinite-1}, the assertions (1), (2), (4) and~(5) of Theorem~\ref{thm:moments-limit-point} and the assertions of Corollary~\ref{cor:transfer-operators-coordinate} hold for $\mu$; and by Theorem~\ref{thm:continuum-infinite-1}(1) the conclusions of Proposition~\ref{prop:sub-gaussian-random} hold under $\mu$ for the space $L^2(\mathbb{R}^3;\Lambda^2)$ with the norm $\|\cdot\|_2$. In the proof of Part~I of Theorem~\ref{thm:main-theorem} for the averaged laws (Section~\ref{sec:admissibility-limit-points}) the properties~\eqref{it:main-theorem-P1}, \eqref{it:main-theorem-P2} and~\eqref{it:main-theorem-P4}--\eqref{it:main-theorem-P7} are derived from these statements alone: \eqref{it:main-theorem-P1} from Theorem~\ref{thm:continuum-infinite-1}(1), (4) and Proposition~\ref{prop:sub-gaussian-random}(1), (5), (7); \eqref{it:main-theorem-P2} from Theorem~\ref{thm:moments-limit-point}(1), (2) and Proposition~\ref{prop:sub-gaussian-random}(7); \eqref{it:main-theorem-P4} from Theorem~\ref{thm:continuum-infinite-1}(8), (3), (4); \eqref{it:main-theorem-P5} from Theorem~\ref{thm:continuum-infinite-1}(5), (6); \eqref{it:main-theorem-P6} from Theorem~\ref{thm:continuum-infinite-1}(7) and Theorem~\ref{thm:moments-limit-point}(5); and~\eqref{it:main-theorem-P7} is Corollary~\ref{cor:transfer-operators-coordinate}. The same derivation applies to $\mu$. As to~\eqref{it:main-theorem-P3}: one has $S_0=1$, and $S_n$ is symmetric and real on real tests, by Theorem~\ref{thm:moments-limit-point}(1); the invariance under $F\mapsto-F$ holds by assumption, and it gives $S_n=(-1)^nS_n=0$ for every odd $n$, which is also the last assertion of Proposition~\ref{prop:laws-four-properties}, the characteristic functional of $\mu$ being even.

Finally, in the proof of Part~I of Theorem~\ref{thm:main-theorem} the property~(c) is derived from~\eqref{it:main-theorem-P1}, \eqref{it:main-theorem-P4}, \eqref{it:main-theorem-P5} and~\eqref{it:main-theorem-P6} alone, by means of Theorem~\ref{thm:ergodic-decomposition} together with its part~(d), Remark~\ref{rem:components-all-translations}, Theorem~\ref{thm:cubic-symmetry-components}(i) and Corollary~\ref{cor:ergodic-components-2}. Since $\mu$ has these four properties, it has the property~(c).
\end{proof}

\paragraph{The class $\mathbf{K}(C)$.}
Fix $e\in(0,e_2]$, with the number $e_2$ of Theorem~\ref{thm:main-theorem}, and fix the data of Section~\ref{sec:probes}: in particular the test $f=dd^{*}\beta$, for which $b(f)>0$, the point $x_0$, and the tests $g_0=f$ and $g_{\pm}=f\pm\tau_{x_0}f$. A soft law $\mu$ satisfies~\eqref{eq:setup-G} and is invariant under the translations of $\mathbb{R}^3$; let $\mu=\int\mathrm{d}\nu(\xi)\,\mu_{\xi}$ be its ergodic decomposition under $\mathbb{R}^3$, and let $\kappa_4^{\xi}$ and $C_{\xi}$ be the fourth cumulant and the covariance of $\mu_{\xi}$ (Section~\ref{sec:laws-2}).

\begin{definition}
For $C>0$ let $\mathbf{K}(C)$ be the set of the soft laws $\mu$ that satisfy the two conditions
\begin{multline}\tag{K$_{\mu}$}\label{eq:free-coefficient-K}
\left|\kappa_4^{\mu}(g)-e^4\cdot b(g)\right|\le C\cdot e^{4+\sigma}\quad\text{for }g\in\left\{g_0,g_{+},g_{-}\right\},\\
\left|S_2^{\mu}(f,f)-\|f\|_2^2\right|\le C\cdot e^{2-\mathfrak{e}_2},
\end{multline}
and
\begin{multline}\tag{K$_{\xi}$}\label{eq:free-coefficient-K-2}
\left|\kappa_4^{\xi}(g)-e^4\cdot b(g)\right|\le C\cdot e^{4+\sigma}\quad\text{and}\quad\left|C_{\xi}(g,g)-\|g\|_2^2\right|\le C\cdot e^{2-\mathfrak{e}_2}\\
\text{for }\nu\text{-almost every }\xi\text{ and every }g\in\left\{g_0,g_{+},g_{-}\right\}.
\end{multline}
\end{definition}

\begin{remark}
\begin{enumerate}
\item[(i)] Let $\mathrm{nat}\in\{\mathrm{af},\mathrm{av}\}$. Every limit point $\mu$ of the laws $P^{\mathrm{nat}}_{N,T}$ of Theorem~\ref{thm:main-theorem}, formed with the charge $e$, lies in $\mathbf{K}(C_{\mathrm{kap}})$. Indeed, by Theorem~\ref{thm:main-theorem}(i) such a law $\mu$ has the properties~\eqref{it:main-theorem-P1}--\eqref{it:main-theorem-P8} and~(c) of that theorem: it satisfies~\eqref{eq:setup-H} by~(c), it is concentrated on $\{dF=0\}$ by~\eqref{it:main-theorem-P4}, and it is invariant under $F\mapsto-F$ by~\eqref{it:main-theorem-P3}; so $\mu$ is soft. The condition~\eqref{eq:free-coefficient-K} with $C=C_{\mathrm{kap}}$ is Theorem~\ref{thm:main-theorem}(iii), and the condition~\eqref{eq:free-coefficient-K-2} with $C=C_{\mathrm{kap}}$ is Theorem~\ref{thm:main-theorem}(vi), (vii).
\item[(ii)] Let $C>0$ satisfy
\[
\frac{5}{2}\,C\cdot e^{\sigma}\le\frac{b(f)}{4},
\]
and let $\mu\in\mathbf{K}(C)$. Then $\mu$ has the properties~(iv) and~(v) of Theorem~\ref{thm:main-theorem}, and
\[
\mathrm{Ext}^{\xi}_{x_0}(f)\ge\frac{1}{2}e^4\cdot b(f)\qquad\text{for }\nu\text{-almost every }\xi .
\]
Indeed, $\mu$ satisfies~\eqref{eq:setup-G}, which is the condition~\eqref{eq:fourth-cumulants-H1} of Section~\ref{sec:hypotheses-law}; it is invariant under the translations of $\mathbb{R}^3$; and the first inequality of~\eqref{eq:free-coefficient-K} is the hypothesis~\eqref{eq:fourth-cumulants-4g} of Section~\ref{sec:hypotheses-law} at the three tests $g_0$, $g_{+}$, $g_{-}$ of Theorem~\ref{thm:non-gaussian-components-2}, with the exponent $\sigma$ and with all three constants equal to $C$. Step~5 of the proof of Theorem~\ref{thm:main-theorem} (Section~\ref{sec:proof-part-ii-2}) derives the properties~(iv) and~(v) of a law from these three facts, with $C_{\mathrm{kap}}$ in place of $C$, from the condition~\eqref{eq:setup-1} on $x_0$ and from the inequality $\frac{5}{2}\,C_{\mathrm{kap}}\cdot e^{\sigma}\le\frac{1}{4}b(f)$; with $C$ in place of $C_{\mathrm{kap}}$ the same argument applies to $\mu$. Moreover, for $\nu$-almost every $\xi$ the component $\mu_{\xi}$ satisfies~\eqref{eq:fourth-cumulants-H1} and is invariant under the translations (Corollary~\ref{cor:properties-almost-every}), and by~\eqref{eq:free-coefficient-K-2} it satisfies~\eqref{eq:fourth-cumulants-4g} at the three tests $g_0$, $g_{+}$, $g_{-}$ with the constant $C$. As in Step~6 of the proof of Theorem~\ref{thm:main-theorem}, Theorem~\ref{thm:non-gaussian-components-2}(i), applied to $\mu_{\xi}$, gives $\mathrm{Ext}^{\xi}_{x_0}(f)\ge\frac{1}{2}e^4\cdot b(f)$, the conditions of that theorem being~\eqref{eq:setup-1} and the inequality assumed on $C$.
\item[(iii)] Consequently, if $\frac{5}{2}\,C\cdot e^{\sigma}\le\frac{1}{4}b(f)$, then, by Lemma~\ref{lem:properties-soft-laws}, by~\eqref{eq:free-coefficient-K}, \eqref{eq:free-coefficient-K-2} and by~(ii), every $\mu\in\mathbf{K}(C)$ has each property that Theorem~\ref{thm:main-theorem} and Corollary~\ref{cor:ergodic-components}(a) assert of a limit point, with $C$ in place of $C_{\mathrm{kap}}$, except the property~\eqref{it:main-theorem-P8} of Theorem~\ref{thm:main-theorem}, which concerns the lattice laws that converge to the limit point.
\item[(iv)] Let $\mu$ be a soft law, let $j\in\{1,2,3\}$, and let $\nu^j$ be the measure of Theorem~\ref{thm:ergodic-decomposition} for the law $\mu$ and $J=\{j\}$. The law $\mu$ lies in the class $\mathcal{C}$ of Section~\ref{sec:properties-laws} and satisfies~\eqref{eq:setup-G} (Section~\ref{sec:hypothesis-law}). Hence Theorem~\ref{thm:decomposition-along-coordinate}, together with Theorem~\ref{thm:ergodic-decomposition}(d) applied to the set $\{dF=0\}$, shows that $\nu^j$-almost every $\rho$ has the properties stated in Corollary~\ref{cor:ergodic-components}(b) other than the bound on $\kappa_4^{\rho}(f)$ and its consequence that $\rho$ is not Gaussian. Let now $C>0$ be a number with $C\cdot e^{\sigma}<b(f)$, and assume that $\left|\kappa_4^{\rho}(f)-e^4\cdot b(f)\right|\le C\cdot e^{4+\sigma}$ for $\nu^j$-almost every $\rho$. Then
\[
\kappa_4^{\rho}(f)\ge e^4\cdot\left(b(f)-C\cdot e^{\sigma}\right)>0\qquad\text{for }\nu^j\text{-almost every }\rho,
\]
so that $\nu^j$-almost every $\rho$ is not Gaussian (Section~\ref{sec:cumulants-gaussian-laws}).
\end{enumerate}
\end{remark}

\begin{definition}
Let $\mu$ and $\rho$ be laws on $\mathcal{S}'$, and let $0<\epsilon<1$. Let $\mathrm{Tm}\colon\mathcal{S}'\times\mathcal{S}'\to\mathcal{S}'$ be the map
\[
\mathrm{Tm}(F,G):=(1-\epsilon)^{\frac{1}{2}}F+\epsilon^{\frac{1}{2}}G .
\]
It is measurable for the product of the cylinder $\sigma$-algebras on $\mathcal{S}'\times\mathcal{S}'$ and the cylinder $\sigma$-algebra on $\mathcal{S}'$, because for every $f\in\mathcal{S}$ the function $X(f)\circ\mathrm{Tm}$, that is $(F,G)\mapsto(1-\epsilon)^{\frac{1}{2}}F(f)+\epsilon^{\frac{1}{2}}G(f)$, is measurable. The law $\mu\ast_{\epsilon}\rho$ is the law of $(1-\epsilon)^{\frac{1}{2}}F+\epsilon^{\frac{1}{2}}G$ under $\mu\otimes\rho$, that is, the image of the product measure $\mu\otimes\rho$ under $\mathrm{Tm}$.
\end{definition}

\begin{proposition}[Operations on the class $\mathbf{K}(C)$]\label{prop:operations-class}
Let $C>0$, $\mu\in\mathbf{K}(C)$ and $\boldsymbol{\delta}>0$. There is a number $\epsilon_0>0$, depending on $\boldsymbol{\delta}$, on $e$ and on the data of Section~\ref{sec:probes}, such that the following hold for every $\epsilon\in(0,\epsilon_0]$.
\begin{enumerate}
\item[(a)] For every $\rho\in\left\{\gamma_{C_0},\gamma_{q_A},m_A,m_B\right\}$ the law $\mu\ast_{\epsilon}\rho$ lies in $\mathbf{K}(C+\boldsymbol{\delta})$. If $\kappa_4^{\mu}(f)\ne0$, then $\mu\ast_{\epsilon}\gamma_{C_0}\ne\mu$.
\item[(b)] Let $\kappa_4^{\mu}(f)\ne0$, and put
\[
\mu':=\frac{1}{2}\left(\mu+\mu\ast_{\epsilon}\gamma_{C_0}\right).
\]
Then $\mu'\in\mathbf{K}(C+\boldsymbol{\delta})$, the law $\mu'$ is not ergodic under the translations of $\mathbb{R}^3$, and none of the conditions~\eqref{eq:fourth-cumulants-C1}, \eqref{eq:fourth-cumulants-C2}, \eqref{eq:fourth-cumulants-C3} and~\eqref{eq:fourth-cumulants-M} of Theorem~\ref{thm:clustering-mixing-ergodicity} holds for $\mu'$.
\item[(c)] The laws $\mu$ and $\mu\ast_{\epsilon}\gamma_{q_A}$ are not both invariant under the rotation by the angle $\frac{\pi}{4}$ about the third coordinate axis, that is, the rotation which fixes $e_3$ and maps $e_1$ to $2^{-\frac{1}{2}}\left(e_1+e_2\right)$, acting on $\mathcal{S}'$ as in Section~\ref{sec:euclidean-maps-actions}.
\item[(d)] Almost every ergodic component $\rho'$ of $\mu\ast_{\epsilon}m_B$ under $\mathbb{R}^3$ has the following two properties for $j=1,2,3$. First, $\dim\ker\mathsf{H}_j(\rho')=\infty$, where $\mathsf{H}_j(\rho')$ is the operator of the reconstruction of $\rho'$ along $e_j$ (Proposition~\ref{prop:reconstruction-direction}). Second, $\rho'$ does not cluster along $e_j$: there are $F,G\in L^2(\rho')$ such that $\rho'\left(\overline{F}\cdot U(se_j)G\right)$ does not tend to $\overline{\rho'(F)}\cdot\rho'(G)$ as $|s|\to\infty$.
\item[(e)] Let $w\in O_{\mathrm{cub}}\setminus N_3$. The two statements ``almost every ergodic component of $\mu$ under $\mathbb{R}^3$ is invariant under $w$'' and ``almost every ergodic component of $\mu\ast_{\epsilon}m_A$ under $\mathbb{R}^3$ is invariant under $w$'' are not both true.
\item[(f)] Let $j\in\{1,2,3\}$, let $\nu^j$ be the measure of Theorem~\ref{thm:ergodic-decomposition} for the law $\mu$ and $J=\{j\}$, and assume that
\[
\left|\kappa_4^{\rho'}(f)-e^4\cdot b(f)\right|\le C\cdot e^{4+\sigma}\qquad\text{for }\nu^j\text{-almost every }\rho'.
\]
Then the same holds, with $C+\boldsymbol{\delta}$ in place of $C$, for each of the laws $\mu\ast_{\epsilon}\rho$ of~(a) and, if $\kappa_4^{\mu}(f)\ne0$, for the law $\mu'$ of~(b), each of them with its own measure of Theorem~\ref{thm:ergodic-decomposition} for $J=\{j\}$ in place of $\nu^j$.
\end{enumerate}
\end{proposition}

\begin{proof}[Proof of Proposition~\ref{prop:operations-class}(a)]
\emph{Step 1: the four laws $\rho$ are soft.} By Lemma~\ref{lem:gaussian-laws-bounded}(a) the law $\gamma_{C_0}$ satisfies~\eqref{eq:setup-G}. The form $C_0$ is invariant under all Euclidean motions (Lemma~\ref{lem:free-transverse-covariance}(b)), hence so is $\gamma_{C_0}$ (Lemma~\ref{lem:gaussian-laws-bounded}(b)), in particular under the translations and under $O_{\mathrm{cub}}$. Moreover $C_0(\theta_j\cdot h,h)\ge0$ for every $h\in\mathcal{D}(\{x_j>0\})$ (Lemma~\ref{lem:free-transverse-covariance}(d)), and Lemma~\ref{lem:gaussian-laws-bounded}(d) gives the inequality~\eqref{eq:fourth-cumulants-H3} for $\gamma_{C_0}$, that is, the condition~\eqref{eq:setup-RPj}, for $j=1,2,3$. By Lemma~\ref{lem:free-transverse-covariance}(e), $C_0(d^{*}h,d^{*}h)=0$, that is, $X(d^{*}h)=0$ $\gamma_{C_0}$-almost surely, for every real $h\in\mathcal{S}(\mathbb{R}^3)$; hence $\gamma_{C_0}$ is concentrated on $\{dF=0\}$ (Section~\ref{sec:closed-distributions}). The characteristic functional $f\mapsto\exp\left(-\frac{1}{2}C_0(f,f)\right)$ is even, so that $\gamma_{C_0}$ is invariant under $F\mapsto-F$ (Lemma~\ref{lem:exponentials-determination-law}(b)). Thus $\gamma_{C_0}$ is soft.

The laws $m_A$ and $m_B$ satisfy~\eqref{eq:setup-H} and are concentrated on $\{dF=0\}$ (Theorem~\ref{thm:asymmetric-components-non}). By Lemma~\ref{lem:gaussian-mixtures}, applied for $m_A$ to each of the six laws $\gamma^{(\pi)}$, their characteristic functionals are
\[
f\mapsto\mathbb{E}\exp\left(-\frac{1}{2}q_u(f)\right),
\]
where $\mathbb{E}$ denotes the expectation over a random triple $u$, which is uniformly distributed on the six triples $u^{(\pi)}$ for $m_A$, and which is the random triple of Section~\ref{sec:free-maxwell-law-2} for $m_B$. These functionals are even, because $q_u(-f)=q_u(f)$; so $m_A$ and $m_B$ are invariant under $F\mapsto-F$ (Lemma~\ref{lem:exponentials-determination-law}(b)), and they are soft.

For every triple $u$ one has $0\le q_u(f)\le\|f\|_2^2$ (Lemma~\ref{lem:coordinatewise-changes-variables}(i)); hence $0\le q_A(f)\le\|f\|_2^2$, and the law $\gamma_{q_A}$ is defined and satisfies~\eqref{eq:setup-G} (Lemma~\ref{lem:gaussian-laws-bounded}(a)). The form $q_A$ is invariant under the translations, as every $q^{(\pi)}$ is (Section~\ref{sec:permuted-diagonal-rescalings}(1)), and under $O_{\mathrm{cub}}$, which permutes the six forms $q^{(\pi)}$ (Section~\ref{sec:permuted-diagonal-rescalings}(2)); so $\gamma_{q_A}$ is invariant under the translations and under $O_{\mathrm{cub}}$ (Lemma~\ref{lem:gaussian-laws-bounded}(b)). For every real $h\in\mathcal{S}(\mathbb{R}^3)$ and every $\pi\in S_3$ one has $q^{(\pi)}(d^{*}h)=\gamma^{(\pi)}\left(X(d^{*}h)^2\right)=0$, because $\gamma^{(\pi)}$ is concentrated on the closed distributions (Lemma~\ref{lem:mixtures-changes-variables}(i)). Hence $q_A(d^{*}h)=0$, so that $X(d^{*}h)=0$ $\gamma_{q_A}$-almost surely for every $h$, and therefore for all $h$ simultaneously, by the argument of Section~\ref{sec:closed-distributions} with a countable dense set of functions $h$: the law $\gamma_{q_A}$ is concentrated on $\{dF=0\}$. Its characteristic functional $f\mapsto\exp\left(-\frac{1}{2}q_A(f)\right)$ is even, so that $\gamma_{q_A}$ is invariant under $F\mapsto-F$. It remains to verify~\eqref{eq:setup-RPj}. Since $\gamma^{(\pi)}\left(\exp\left(iX(h)\right)\right)=\exp\left(-\frac{1}{2}q^{(\pi)}(h)\right)$ and $q^{(\pi)}\left(6^{-\frac{1}{2}}h\right)=\frac{1}{6}q^{(\pi)}(h)$, the definition of $q_A$ gives
\[
\exp\left(-\frac{1}{2}q_A(h)\right)=\prod_{\pi\in S_3}\gamma^{(\pi)}\left(\exp\left(iX\left(6^{-\frac{1}{2}}h\right)\right)\right)\qquad(h\in\mathcal{S}).
\]
Let $j\in\{1,2,3\}$, $n\ge1$ and $f_1,\dots,f_n\in\mathcal{D}(\{x_j>0\})$. Applying this identity to $h=f_q-\theta_j\cdot f_p$, one finds that the matrix
\[
\left(\gamma_{q_A}\left(\exp\left(i\left(X(f_q)-X(\theta_j\cdot f_p)\right)\right)\right)\right)_{p,q=1,\dots,n}
\]
is the entrywise product of the six matrices
\[
\left(\gamma^{(\pi)}\left(\exp\left(i\left(X\left(6^{-\frac{1}{2}}f_q\right)-X\left(\theta_j\cdot6^{-\frac{1}{2}}f_p\right)\right)\right)\right)\right)_{p,q=1,\dots,n},\qquad\pi\in S_3.
\]
These six matrices are positive semidefinite by the condition~\eqref{eq:setup-RPj} for the laws $\gamma^{(\pi)}$ (Section~\ref{sec:permuted-diagonal-rescalings}(1)), applied to the tests $6^{-\frac{1}{2}}f_1,\dots,6^{-\frac{1}{2}}f_n$, and by the Schur product theorem their entrywise product is positive semidefinite. This is~\eqref{eq:setup-RPj} for $\gamma_{q_A}$. So $\gamma_{q_A}$ satisfies~\eqref{eq:setup-H}, and it is soft.

\emph{Step 2: the law $\mu\ast_{\epsilon}\rho$ is soft.} Let $\mu$ and $\rho$ be soft laws and $0<\epsilon<1$. The map $\mathrm{Tm}$ is linear and measurable (see the definition of $\mu\ast_{\epsilon}\rho$), and it commutes with the actions of the Euclidean maps on $\mathcal{S}'$ (Section~\ref{sec:euclidean-maps-actions}) and with the change of sign: $\mathrm{Tm}(gF,gG)=g\,\mathrm{Tm}(F,G)$ for every Euclidean map $g$, and $\mathrm{Tm}(-F,-G)=-\mathrm{Tm}(F,G)$. The product measure $\mu\otimes\rho$ is invariant under the diagonal actions $(F,G)\mapsto(gF,gG)$ of the translations and of the elements of $O_{\mathrm{cub}}$, and under $(F,G)\mapsto(-F,-G)$, because $\mu$ and $\rho$ are invariant under these maps; hence the image $\mu\ast_{\epsilon}\rho$ is invariant under the translations, under $O_{\mathrm{cub}}$ and under $F\mapsto-F$. The set $\{dF=0\}$ is a linear subspace of $\mathcal{S}'$; so $\mathrm{Tm}$ maps the set $\{dF=0\}\times\{dF=0\}$, which has full $\mu\otimes\rho$-measure, into $\{dF=0\}$, and $\mu\ast_{\epsilon}\rho$ is concentrated on $\{dF=0\}$. For $g\in\mathcal{S}$ the bound~\eqref{eq:setup-G} for $\mu$ at the test $(1-\epsilon)^{\frac{1}{2}}g$ and for $\rho$ at the test $\epsilon^{\frac{1}{2}}g$ gives
\begin{align*}
\int\mathrm{d}\left(\mu\ast_{\epsilon}\rho\right)\,\exp X(g)&=\int\mathrm{d}\mu(F)\,\exp\left((1-\epsilon)^{\frac{1}{2}}F(g)\right)\cdot\int\mathrm{d}\rho(G)\,\exp\left(\epsilon^{\frac{1}{2}}G(g)\right)\\
&\le\exp\left(\frac{1-\epsilon}{2}\|g\|_2^2+\frac{\epsilon}{2}\|g\|_2^2\right)=\exp\left(\frac{1}{2}\|g\|_2^2\right),
\end{align*}
which is~\eqref{eq:setup-G} for $\mu\ast_{\epsilon}\rho$. Finally let $j\in\{1,2,3\}$, $n\ge1$ and $f_1,\dots,f_n\in\mathcal{D}(\{x_j>0\})$. As $X(h)\circ\mathrm{Tm}=(1-\epsilon)^{\frac{1}{2}}X(h)(F)+\epsilon^{\frac{1}{2}}X(h)(G)$ for every $h\in\mathcal{S}$, and as $F$ and $G$ are independent under $\mu\otimes\rho$, the matrix
\[
\left((\mu\ast_{\epsilon}\rho)\left(\exp\left(i\left(X(f_q)-X(\theta_j\cdot f_p)\right)\right)\right)\right)_{p,q=1,\dots,n}
\]
is the entrywise product of the corresponding matrix for $\mu$, formed with the tests $(1-\epsilon)^{\frac{1}{2}}f_p$, and the corresponding matrix for $\rho$, formed with the tests $\epsilon^{\frac{1}{2}}f_p$. These two matrices are positive semidefinite by~\eqref{eq:setup-RPj} for $\mu$ and for $\rho$, and so is their entrywise product, by the Schur product theorem. This is~\eqref{eq:setup-RPj} for $\mu\ast_{\epsilon}\rho$. So $\mu\ast_{\epsilon}\rho$ satisfies~\eqref{eq:setup-H}, and it is soft.

\emph{Step 3: ergodic components.} Let $\rho_1$ and $\rho_2$ be laws in the class $\mathcal{C}$ of Section~\ref{sec:properties-laws} which are ergodic under $\mathbb{R}^3$, and let $0<\epsilon<1$. For $a\in\mathbb{R}^3$ let $U(a)$ denote the unitary operators of $L^2(\rho_1)$ and of $L^2(\rho_2)$ by which the translations act (Section~\ref{sec:euclidean-maps-actions}), and also the unitary operator of $L^2(\rho_1\otimes\rho_2)$ given by the diagonal action, $\left(U(a)H\right)(F,G):=H\left(\tau_a^{-1}F,\tau_a^{-1}G\right)$. For $\Psi_1\in L^2(\rho_1)$ and $\Psi_2\in L^2(\rho_2)$ let $(\Psi_1\otimes\Psi_2)(F,G):=\Psi_1(F)\cdot\Psi_2(G)$; then $U(a)(\Psi_1\otimes\Psi_2)=U(a)\Psi_1\otimes U(a)\Psi_2$. By Theorem~\ref{thm:invariant-vectors-clustering}(iii)1, for $i=1,2$ and all $\Psi_i,\Psi_i'\in L^2(\rho_i)$,
\[
\langle\Psi_i,U(a)\Psi_i'\rangle\to\langle\Psi_i,1\rangle\cdot\langle1,\Psi_i'\rangle\qquad\text{as }\min_l|a_l|\to\infty,
\]
the scalar products being those of $L^2(\rho_i)$. For product vectors it follows that
\[
\left\langle\Psi_1\otimes\Psi_2,U(a)\left(\Psi_1'\otimes\Psi_2'\right)\right\rangle=\langle\Psi_1,U(a)\Psi_1'\rangle\cdot\langle\Psi_2,U(a)\Psi_2'\rangle\to\langle\Psi_1\otimes\Psi_2,1\rangle\cdot\langle1,\Psi_1'\otimes\Psi_2'\rangle .
\]
The finite linear combinations of products $\Psi_1\otimes\Psi_2$ are dense in $L^2(\rho_1\otimes\rho_2)$, both sides are sesquilinear in the pair of vectors, and $\|U(a)\|=1$; hence the convergence holds for all vectors of $L^2(\rho_1\otimes\rho_2)$. If $H\in L^2(\rho_1\otimes\rho_2)$ satisfies $U(a)H=H$ for every $a$, then $\|H\|^2=\langle H,U(a)H\rangle\to\left|\langle1,H\rangle\right|^2$; as $\|H\|^2=\left|\langle1,H\rangle\right|^2+\left\|H-\langle1,H\rangle\cdot1\right\|^2$, the vector $H$ is constant. So $\rho_1\otimes\rho_2$ is ergodic under the diagonal translations. Its image $\rho_1\ast_{\epsilon}\rho_2$ under $\mathrm{Tm}$ is ergodic as well. Indeed, $\mathrm{Tm}$ commutes with the translations; hence inverse images under $\mathrm{Tm}$ of invariant sets are invariant under the diagonal translations, and so is $H\circ\mathrm{Tm}$ for every $H\in L^2(\rho_1\ast_{\epsilon}\rho_2)$ with $U(a)H=H$ for all $a$. Such an $H\circ\mathrm{Tm}$ is constant $\rho_1\otimes\rho_2$-almost everywhere, and then $H$ is constant $\rho_1\ast_{\epsilon}\rho_2$-almost everywhere. The argument of Step~2, applied to the translations, to the reflections $\theta_1$, $\theta_2$, $\theta_3$, to the conditions~\eqref{eq:setup-RPj} and to the bound~\eqref{eq:setup-G}, with $\rho_1$, $\rho_2$ in place of $\mu$, $\rho$, shows that $\rho_1\ast_{\epsilon}\rho_2$ lies in $\mathcal{C}$, and that it satisfies~\eqref{eq:setup-G} if $\rho_1$ and $\rho_2$ do.

Now let $\mu$ be a soft law and $\rho\in\left\{\gamma_{C_0},\gamma_{q_A},m_A,m_B\right\}$, and let
\[
\mu=\int\mathrm{d}\nu(\xi)\,\mu_{\xi},\qquad\rho=\int\mathrm{d}\nu_{\rho}(\boldsymbol{\zeta})\,\rho_{\boldsymbol{\zeta}}
\]
be their ergodic decompositions under $\mathbb{R}^3$. For the two Gaussian laws and for $m_B$ the measure $\nu_{\rho}$ is a point mass, by the uniqueness in Theorem~\ref{thm:ergodic-decomposition}(a): the laws $\gamma_{C_0}$ and $\gamma_{q_A}$ are mixing, hence ergodic, by Lemma~\ref{lem:gaussian-laws-bounded}(c), their covariances being invariant under the translations (Lemma~\ref{lem:free-transverse-covariance}(b) and Step~1), and $m_B$ is ergodic by Theorem~\ref{thm:asymmetric-components-non}(B). For $m_A$ the measure $\nu_{\rho}$ is the uniform measure on the six laws $\gamma^{(\pi)}$, $\pi\in S_3$ (Section~\ref{sec:permuted-diagonal-rescalings}(4)). Thus $\nu_{\rho}$ has finite support, and every law $\rho_{\boldsymbol{\zeta}}$ in its support lies in $\mathcal{C}$, satisfies~\eqref{eq:setup-G} and is ergodic: $\gamma_{C_0}$, $\gamma_{q_A}$ and $m_B$ are soft by Step~1, hence satisfy~\eqref{eq:setup-H} and lie in $\mathcal{C}$ (Section~\ref{sec:hypothesis-law}), and every $\gamma^{(\pi)}$ lies in $\mathcal{C}$ and satisfies~\eqref{eq:setup-G} (Section~\ref{sec:permuted-diagonal-rescalings}(1)). For $\nu$-almost every $\xi$ the law $\mu_{\xi}$ lies in $\mathcal{C}$, satisfies~\eqref{eq:setup-G} and is ergodic, by Theorem~\ref{thm:ergodic-decomposition}(a), (b), (c) (Remark~\ref{rem:components-all-translations}).

Let $A$ be a cylinder set and $\boldsymbol{\zeta}$ a point of the support of $\nu_{\rho}$. For every $\xi$, Fubini's theorem gives
\[
\left(\mu_{\xi}\ast_{\epsilon}\rho_{\boldsymbol{\zeta}}\right)(A)=\mu_{\xi}\left(\Phi_A\right),\qquad\Phi_A(F):=\rho_{\boldsymbol{\zeta}}\left(\left\{G\in\mathcal{S}':\mathrm{Tm}(F,G)\in A\right\}\right),
\]
and $\Phi_A$ is a bounded measurable function on $\mathcal{S}'$, again by Fubini's theorem. The function $\xi\mapsto\mu_{\xi}(\Phi)$ is measurable for every bounded measurable function $\Phi$ on $\mathcal{S}'$, since $\xi\mapsto\mu_{\xi}(B)$ is measurable for every cylinder set $B$; hence $\xi\mapsto\left(\mu_{\xi}\ast_{\epsilon}\rho_{\boldsymbol{\zeta}}\right)(A)$ is measurable. As $\mu=\int\mathrm{d}\nu(\xi)\,\mu_{\xi}$ and $\nu_{\rho}$ has finite support, Fubini's theorem gives
\[
\left(\mu\ast_{\epsilon}\rho\right)(A)=\int\mathrm{d}\nu(\xi)\int\mathrm{d}\nu_{\rho}(\boldsymbol{\zeta})\,\left(\mu_{\xi}\ast_{\epsilon}\rho_{\boldsymbol{\zeta}}\right)(A)\qquad\text{for every cylinder set }A .
\]
By the first part of this step, $\mu_{\xi}\ast_{\epsilon}\rho_{\boldsymbol{\zeta}}$ is ergodic for $\nu$-almost every $\xi$ and every $\boldsymbol{\zeta}$ in the support of $\nu_{\rho}$. By the uniqueness in Theorem~\ref{thm:ergodic-decomposition}(a), the image of $\nu\otimes\nu_{\rho}$ under $(\xi,\boldsymbol{\zeta})\mapsto\mu_{\xi}\ast_{\epsilon}\rho_{\boldsymbol{\zeta}}$ is therefore the measure of the ergodic decomposition of $\mu\ast_{\epsilon}\rho$ under $\mathbb{R}^3$: the ergodic components of $\mu\ast_{\epsilon}\rho$ are the laws $\mu_{\xi}\ast_{\epsilon}\rho_{\boldsymbol{\zeta}}$.

\emph{Step 4: cumulants.} Let $\mu_1$ and $\rho_1$ be laws on $\mathcal{S}'$ which satisfy~\eqref{eq:setup-G}, let $0<\epsilon<1$, and let $g\in\mathcal{S}$. A law with~\eqref{eq:setup-G} is centred (Lemma~\ref{lem:moment-bounds}(a)); for such a law the fourth cumulant of the centred variable $X(g)$ is $\kappa_4(g)=E\,X(g)^4-3\left(E\,X(g)^2\right)^2$, and the argument of Step~3 of the proof of Theorem~\ref{thm:main-theorem} (Section~\ref{sec:proof-part-ii-2}), applied to the bound~\eqref{eq:setup-G} at the tests $yg$ and $-yg$, $y>0$, which gives $E\cosh\left(yX(g)\right)\le\exp\left(\frac{1}{2}y^2\cdot\|g\|_2^2\right)$, yields
\[
0\le S_2(g,g)\le\|g\|_2^2,\qquad\left|\kappa_4(g)\right|\le15\,\|g\|_2^4.
\]
Under $\mu_1\otimes\rho_1$ the variable $X(g)\circ\mathrm{Tm}=(1-\epsilon)^{\frac{1}{2}}X(g)(F)+\epsilon^{\frac{1}{2}}X(g)(G)$ is the sum of two independent variables with finite exponential moments, and the logarithm of the moment generating function of such a sum is the sum of the logarithms of the moment generating functions of the two terms. Comparing the coefficients of the fourth and of the second power, and using that the variables are centred, one obtains
\begin{align*}
\kappa_4^{\mu_1\ast_{\epsilon}\rho_1}(g)&=(1-\epsilon)^2\cdot\kappa_4^{\mu_1}(g)+\epsilon^2\cdot\kappa_4^{\rho_1}(g),\\
S_2^{\mu_1\ast_{\epsilon}\rho_1}(g,g)&=(1-\epsilon)\cdot S_2^{\mu_1}(g,g)+\epsilon\cdot S_2^{\rho_1}(g,g).
\end{align*}
Since $0\le1-(1-\epsilon)^2\le2\epsilon$ and $\epsilon^2\le\epsilon$, these two numbers differ from $\kappa_4^{\mu_1}(g)$ and $S_2^{\mu_1}(g,g)$ by at most
\begin{align*}
\left|\kappa_4^{\mu_1\ast_{\epsilon}\rho_1}(g)-\kappa_4^{\mu_1}(g)\right|&\le15\cdot(2\epsilon+\epsilon)\cdot\|g\|_2^4=45\,\epsilon\cdot\|g\|_2^4,\\
\left|S_2^{\mu_1\ast_{\epsilon}\rho_1}(g,g)-S_2^{\mu_1}(g,g)\right|&=\epsilon\cdot\left|S_2^{\rho_1}(g,g)-S_2^{\mu_1}(g,g)\right|\le\epsilon\cdot\|g\|_2^2.
\end{align*}

These estimates are now applied to the pair $(\mu,\rho)$, where $\mu\in\mathbf{K}(C)$ and $\rho\in\left\{\gamma_{C_0},\gamma_{q_A},m_A,m_B\right\}$, and to the pairs $\left(\mu_{\xi},\rho_{\boldsymbol{\zeta}}\right)$ of Step~3, with $\xi$ in the set of full $\nu$-measure on which $\mu_{\xi}$ satisfies~\eqref{eq:setup-G} and $\boldsymbol{\zeta}$ in the support of $\nu_{\rho}$; all these laws satisfy~\eqref{eq:setup-G}, by Step~1 and Step~3. Put
\[
\epsilon_0:=\min\left\{\frac{\boldsymbol{\delta}\cdot e^{4+\sigma}}{46\max_{g\in\{g_0,g_{+},g_{-}\}}\|g\|_2^4},\ \frac{\boldsymbol{\delta}\cdot e^{2-\mathfrak{e}_2}}{\max_{g\in\{g_0,g_{+},g_{-}\}}\|g\|_2^2},\ \frac{1}{2}\right\};
\]
the maxima are positive, because $b(f)>0$ forces $f\ne0$. Let $0<\epsilon\le\epsilon_0$. Then $45\,\epsilon\cdot\|g\|_2^4\le\boldsymbol{\delta}\cdot e^{4+\sigma}$ and $\epsilon\cdot\|g\|_2^2\le\boldsymbol{\delta}\cdot e^{2-\mathfrak{e}_2}$ for $g\in\{g_0,g_{+},g_{-}\}$. With~\eqref{eq:free-coefficient-K} for $\mu$, and with $f=g_0$, one obtains
\begin{align*}
\left|\kappa_4^{\mu\ast_{\epsilon}\rho}(g)-e^4\cdot b(g)\right|&\le\left|\kappa_4^{\mu}(g)-e^4\cdot b(g)\right|+45\,\epsilon\cdot\|g\|_2^4\le(C+\boldsymbol{\delta})\cdot e^{4+\sigma}\qquad\left(g\in\{g_0,g_{+},g_{-}\}\right),\\
\left|S_2^{\mu\ast_{\epsilon}\rho}(f,f)-\|f\|_2^2\right|&\le\left|S_2^{\mu}(f,f)-\|f\|_2^2\right|+\epsilon\cdot\|f\|_2^2\le(C+\boldsymbol{\delta})\cdot e^{2-\mathfrak{e}_2},
\end{align*}
that is,~\eqref{eq:free-coefficient-K} for $\mu\ast_{\epsilon}\rho$ with $C+\boldsymbol{\delta}$ in place of $C$. By Step~3 the ergodic components of $\mu\ast_{\epsilon}\rho$ are the laws $\mu_{\xi}\ast_{\epsilon}\rho_{\boldsymbol{\zeta}}$; these laws satisfy~\eqref{eq:setup-G} (Step~3) and are therefore centred, so that their covariance at $(g,g)$ is their second moment $S_2(g,g)$, and the covariance $C_{\xi}(g,g)$ of $\mu_{\xi}$ is $S_2^{\mu_{\xi}}(g,g)$. With~\eqref{eq:free-coefficient-K-2} for $\mu$ one obtains, for $\nu$-almost every $\xi$, every $\boldsymbol{\zeta}$ in the support of $\nu_{\rho}$ and every $g\in\{g_0,g_{+},g_{-}\}$,
\begin{align*}
\left|\kappa_4^{\mu_{\xi}\ast_{\epsilon}\rho_{\boldsymbol{\zeta}}}(g)-e^4\cdot b(g)\right|&\le\left|\kappa_4^{\xi}(g)-e^4\cdot b(g)\right|+45\,\epsilon\cdot\|g\|_2^4\le(C+\boldsymbol{\delta})\cdot e^{4+\sigma},\\
\left|S_2^{\mu_{\xi}\ast_{\epsilon}\rho_{\boldsymbol{\zeta}}}(g,g)-\|g\|_2^2\right|&\le\left|C_{\xi}(g,g)-\|g\|_2^2\right|+\epsilon\cdot\|g\|_2^2\le(C+\boldsymbol{\delta})\cdot e^{2-\mathfrak{e}_2},
\end{align*}
that is,~\eqref{eq:free-coefficient-K-2} for $\mu\ast_{\epsilon}\rho$ with $C+\boldsymbol{\delta}$ in place of $C$. By Steps~1 and~2 the law $\mu\ast_{\epsilon}\rho$ is soft. Hence $\mu\ast_{\epsilon}\rho\in\mathbf{K}(C+\boldsymbol{\delta})$, which is the first assertion of~(a).

Finally, $\gamma_{C_0}$ is a centred Gaussian law, so that $\kappa_4^{\gamma_{C_0}}=0$ (Section~\ref{sec:cumulants-gaussian-laws}), and the formula for the fourth cumulant gives
\[
\kappa_4^{\mu\ast_{\epsilon}\gamma_{C_0}}(f)=(1-\epsilon)^2\cdot\kappa_4^{\mu}(f).
\]
As $0<(1-\epsilon)^2<1$ for $0<\epsilon\le\frac{1}{2}$, this number differs from $\kappa_4^{\mu}(f)$ if $\kappa_4^{\mu}(f)\ne0$; in that case $\mu\ast_{\epsilon}\gamma_{C_0}\ne\mu$. This proves~(a).
\end{proof}

\begin{proof}[Proof of Proposition~\ref{prop:operations-class}(b)--(f)]
Let $C>0$, $\mu\in\mathbf{K}(C)$ and $\boldsymbol{\delta}>0$, let $\epsilon_0$ be the number defined in Step~4 of the proof of~(a), and let $0<\epsilon\le\epsilon_0$. The steps cited below are those of that proof, and the decompositions $\mu=\int\mathrm{d}\nu(\xi)\,\mu_{\xi}$ and $\rho=\int\mathrm{d}\nu_{\rho}(\boldsymbol{\zeta})\,\rho_{\boldsymbol{\zeta}}$ are those of Step~3. For a law $\rho_1$ on $\mathcal{S}'$ and a $\rho_1$-integrable function $\Psi$ on $\mathcal{S}'$ write $\rho_1(\Psi):=\int\mathrm{d}\rho_1\,\Psi$. The action $g\cdot h$ of a Euclidean map $g$ on tests $h$, its action $\omega\mapsto g\omega$ on $\mathcal{S}'$, the image $g_{*}\rho_1$ of a law $\rho_1$, the operator $\hat{g}$ on functions on $\mathcal{S}'$ and the translation operators $U(a)=\hat{\tau}_a$, $a\in\mathbb{R}^3$, are those of Section~\ref{sec:euclidean-maps-actions}. Put
\[
\mu_{\epsilon}:=\mu\ast_{\epsilon}\gamma_{C_0}.
\]

\emph{(b)} Let $\kappa_4^{\mu}(f)\ne0$ and $\mu'=\frac{1}{2}\left(\mu+\mu_{\epsilon}\right)$. By~(a), $\mu_{\epsilon}\in\mathbf{K}(C+\boldsymbol{\delta})$ and $\mu_{\epsilon}\ne\mu$.

Each of the conditions that define a soft law, namely the bound~\eqref{eq:setup-G}, the conditions~\eqref{eq:setup-RPj} for $j=1,2,3$, the invariance under the translations, under $O_{\mathrm{cub}}$ and under $F\mapsto-F$, and the concentration on $\{dF=0\}$, is an equality or a non-strict inequality between quantities that depend linearly on the law. Such conditions are preserved under convex combinations of laws; hence $\mu'$ is soft.

Let $\nu_{\epsilon}$ be the probability measure on $(\mathcal{M},\mathcal{B}_{\mathcal{M}})$ that Theorem~\ref{thm:ergodic-decomposition}(a) assigns to the law $\mu_{\epsilon}$ and $J=\{1,2,3\}$, that is, the measure of the ergodic decomposition of $\mu_{\epsilon}$ under $\mathbb{R}^3$, and put $\check{\nu}:=\frac{1}{2}\left(\nu+\nu_{\epsilon}\right)$. Then $\check{\nu}$ is a probability measure on $(\mathcal{M},\mathcal{B}_{\mathcal{M}})$ with $\check{\nu}\left(\mathcal{M}^{123}_{\mathrm{erg}}\right)=1$, and
\[
\mu'(A)=\int\mathrm{d}\check{\nu}(\rho_1)\,\rho_1(A)\qquad\text{for every cylinder set }A.
\]
By the uniqueness assertion of Theorem~\ref{thm:ergodic-decomposition}(a), $\check{\nu}$ is the measure of the ergodic decomposition of $\mu'$. A measurable set of laws is a $\check{\nu}$-null set if and only if it is a $\nu$-null set and a $\nu_{\epsilon}$-null set. The two inequalities of~\eqref{eq:free-coefficient-K-2} hold with the constant $C$ for $\nu$-almost every component of $\mu$, as $\mu\in\mathbf{K}(C)$, and with the constant $C+\boldsymbol{\delta}$ for $\nu_{\epsilon}$-almost every component of $\mu_{\epsilon}$, as $\mu_{\epsilon}\in\mathbf{K}(C+\boldsymbol{\delta})$. Hence $\mu'$ satisfies~\eqref{eq:free-coefficient-K-2} with $C+\boldsymbol{\delta}$ in place of $C$.

The second moments depend linearly on the law:
\[
S_2^{\mu'}(f,f)=\frac{1}{2}\left(S_2^{\mu}(f,f)+S_2^{\mu_{\epsilon}}(f,f)\right);
\]
so the second inequality of~\eqref{eq:free-coefficient-K} holds for $\mu'$ with $C+\boldsymbol{\delta}$ in place of $C$. Now let $g\in\left\{g_0,g_{+},g_{-}\right\}$, and put $s:=S_2^{\mu}(g,g)$ and $s_{\epsilon}:=S_2^{\mu_{\epsilon}}(g,g)$. The laws $\mu$, $\mu_{\epsilon}$ and $\mu'$ satisfy~\eqref{eq:setup-G} and are therefore centred (Lemma~\ref{lem:moment-bounds}(a)); under each of them the fourth cumulant of $X(g)$ is the fourth moment minus three times the square of the second moment. The moments of $\mu'$ are the means of those of $\mu$ and of $\mu_{\epsilon}$; hence
\begin{align*}
\kappa_4^{\mu'}(g)&=\frac{1}{2}\left(\kappa_4^{\mu}(g)+3s^2+\kappa_4^{\mu_{\epsilon}}(g)+3s_{\epsilon}^2\right)-\frac{3}{4}\left(s+s_{\epsilon}\right)^2\\
&=\frac{1}{2}\left(\kappa_4^{\mu}(g)+\kappa_4^{\mu_{\epsilon}}(g)\right)+\frac{3}{4}\left(s-s_{\epsilon}\right)^2.
\end{align*}
By Step~4, applied to the pair $\left(\mu,\gamma_{C_0}\right)$, one has $\left|s-s_{\epsilon}\right|\le\epsilon\cdot\|g\|_2^2$ and $\left|\kappa_4^{\mu_{\epsilon}}(g)-\kappa_4^{\mu}(g)\right|\le45\,\epsilon\cdot\|g\|_2^4$. Therefore
\[
\left|\kappa_4^{\mu'}(g)-\kappa_4^{\mu}(g)\right|\le\frac{45}{2}\,\epsilon\cdot\|g\|_2^4+\frac{3}{4}\,\epsilon^2\cdot\|g\|_2^4\le23\,\epsilon\cdot\|g\|_2^4+\epsilon^2\cdot\|g\|_2^4\le46\,\epsilon\cdot\|g\|_2^4\le\boldsymbol{\delta}\cdot e^{4+\sigma},
\]
the last inequality by the definition of $\epsilon_0$. Together with~\eqref{eq:free-coefficient-K} for $\mu$ this gives the first inequality of~\eqref{eq:free-coefficient-K} for $\mu'$ with $C+\boldsymbol{\delta}$ in place of $C$. Thus $\mu'\in\mathbf{K}(C+\boldsymbol{\delta})$.

Suppose that $\mu'$ is ergodic under the translations of $\mathbb{R}^3$. As $\mu\le2\mu'$, the law $\mu$ is absolutely continuous with respect to $\mu'$, and by the Radon--Nikodym theorem there is a measurable function $\check{h}\colon\mathcal{S}'\to[0,2]$ with
\[
\mu(A)=\int_A\mathrm{d}\mu'\,\check{h}\qquad\text{for every cylinder set }A.
\]
Let $a\in\mathbb{R}^3$. Since $\mu$ and $\mu'$ are invariant under the translation $\tau_a$, one has, for every cylinder set $A$,
\[
\int_A\mathrm{d}\mu'\,U(a)\check{h}=\int_{\tau_a^{-1}A}\mathrm{d}\mu'\,\check{h}=\mu\left(\tau_a^{-1}A\right)=\mu(A),\qquad\tau_a^{-1}A:=\left\{\omega\in\mathcal{S}':\tau_a\omega\in A\right\}.
\]
So $U(a)\check{h}$ is also a density of $\mu$ with respect to $\mu'$, and $U(a)\check{h}=\check{h}$ holds $\mu'$-almost everywhere, by the uniqueness of the density. Thus $\check{h}$ is an element of $L^2(\mu')$ that is invariant under all translations; by the ergodicity of $\mu'$ it is constant, and the constant is $\mu(\mathcal{S}')=1$. Hence $\mu=\mu'$, and then $\mu_{\epsilon}=2\mu'-\mu=\mu$, contrary to~(a). So $\mu'$ is not ergodic.

The law $\mu'$ is invariant under the translations of $\mathbb{R}^3$ and satisfies~\eqref{eq:setup-G}, which is the condition~\eqref{eq:fourth-cumulants-H1}; so Theorem~\ref{thm:clustering-mixing-ergodicity} applies to $\mu'$, and by that theorem each of the conditions~\eqref{eq:fourth-cumulants-C1}, \eqref{eq:fourth-cumulants-M}, \eqref{eq:fourth-cumulants-C2} and~\eqref{eq:fourth-cumulants-C3} implies~\eqref{eq:fourth-cumulants-Erg}, that is, the ergodicity of $\mu'$. Hence none of these four conditions holds for $\mu'$.

\emph{(c)} Let $\mathfrak{l}=\left(\frac{1}{4},\frac{1}{3},\frac{1}{2}\right)$ be the triple of Section~\ref{sec:permuted-diagonal-rescalings}, let $\mathfrak{l}'=\left(\mathfrak{l}'_1,\mathfrak{l}'_2,\mathfrak{l}'_3\right)$ be one of its six permutations, that is, one of the constant triples $u^{(\pi)}$, $\pi\in S_3$, and let $u:=\mathfrak{l}'$. Put
\[
\boldsymbol{\Lambda}:=\operatorname{diag}\left(\mathfrak{l}'_1,\mathfrak{l}'_2,\mathfrak{l}'_3\right),\qquad\check{P}:=\mathfrak{l}_1\cdot\mathfrak{l}_2\cdot\mathfrak{l}_3=\det\boldsymbol{\Lambda},\qquad y_l:=\left(\mathfrak{l}'_l\right)^{-2}\quad(l=1,2,3).
\]
For $h\in\mathcal{S}$ let $\mathbf{b}:=\ast h$ be its vector field (Section~\ref{sec:closed-distributions}), and for an integrable function $\psi$ on $\mathbb{R}^3$ let
\[
\hat{\psi}(p):=\int\mathrm{d}x\,\exp(-ip\cdot x)\,\psi(x)\qquad\left(p\in\mathbb{R}^3\right)
\]
be its Fourier transform, taken componentwise for vector fields. For the constant triple $u$ the map $\Phi_u$ of Section~\ref{sec:free-maxwell-law} is the linear map $x\mapsto\boldsymbol{\Lambda} x$, and the formula~\eqref{eq:ergodic-components-8} reads $\left(\Phi_{u\#}\mathbf{b}\right)_l(x)=\left(\mathfrak{l}'_l\right)^{-1}\cdot\mathbf{b}_l\left(\boldsymbol{\Lambda}^{-1}x\right)$; the substitution $x\mapsto\boldsymbol{\Lambda} x$ in the Fourier integral, whose Jacobian is $\check{P}$, gives
\[
\left(\Phi_{u\#}\mathbf{b}\right)_l^{\wedge}(p)=\frac{\check{P}}{\mathfrak{l}'_l}\cdot\widehat{\mathbf{b}}_l(\boldsymbol{\Lambda} p)\qquad\left(l=1,2,3,\ p\in\mathbb{R}^3\right).
\]
By Lemma~\ref{lem:mixtures-changes-variables}, $q_u(h)=\left\|\Pi\Phi_{u\#}\mathbf{b}\right\|_2^2$, where $\Pi$ is the orthogonal projection of $L^2(\mathbb{R}^3;\mathbb{R}^3)$ with the Fourier multiplier $1-pp^T/|p|_2^2$ (Section~\ref{sec:free-maxwell-law}). Plancherel's theorem and the substitution $k=\boldsymbol{\Lambda} p$, for which $p=\boldsymbol{\Lambda}^{-1}k$, $|p|_2^2=k^T\boldsymbol{\Lambda}^{-2}k$ and $\mathrm{d}p=\check{P}^{-1}\,\mathrm{d}k$, give
\[
q_u(h)=(2\pi)^{-3}\int\mathrm{d}k\,\widehat{\mathbf{b}}(k)^{*}M_{\mathfrak{l}'}(k)\,\widehat{\mathbf{b}}(k),\qquad M_{\mathfrak{l}'}(k):=\check{P}\cdot\left(\boldsymbol{\Lambda}^{-2}-\frac{\boldsymbol{\Lambda}^{-2}kk^T\boldsymbol{\Lambda}^{-2}}{k^T\boldsymbol{\Lambda}^{-2}k}\right)\quad(k\ne0),
\]
and
\[
\operatorname{tr}M_{\mathfrak{l}'}(k)=\check{P}\cdot\left(\sum_{l=1}^3y_l-\frac{\sum_{l=1}^3y_l^2k_l^2}{\sum_{l=1}^3y_lk_l^2}\right).
\]
Let $M_A$ be the mean of the six matrix functions $M_{\mathfrak{l}'}$. As $q_A$ is the mean of the six forms $q^{(\pi)}=q_{u^{(\pi)}}$,
\[
q_A(h)=(2\pi)^{-3}\int\mathrm{d}k\,\widehat{\mathbf{b}}(k)^{*}M_A(k)\,\widehat{\mathbf{b}}(k)\qquad\left(h\in\mathcal{S},\ \mathbf{b}=\ast h\right).
\]
The matrices $M_A(k)$, $k\ne0$, are real and symmetric, and the function $M_A$ is even, bounded and continuous on $\mathbb{R}^3\setminus\{0\}$. Indeed, $\boldsymbol{\Lambda} M_{\mathfrak{l}'}(k)\boldsymbol{\Lambda}$ is $\check{P}$ times the orthogonal projection of $\mathbb{R}^3$ onto the plane orthogonal to $\boldsymbol{\Lambda}^{-1}k$, so that $0\le M_{\mathfrak{l}'}(k)\le\check{P}\cdot\boldsymbol{\Lambda}^{-2}$, and $M_{\mathfrak{l}'}(k)$ depends on $k$ only through the direction of $k$ up to sign.

Let $g$ be a rotation of $\mathbb{R}^3$. By Section~\ref{sec:free-maxwell-law}, $\ast(g\cdot h)=g\,(\ast h)\left(g^{-1}\,\cdot\,\right)$, as $\det g=1$; so the Fourier transform of $\ast(g\cdot h)$ is $k\mapsto g\,\widehat{\mathbf{b}}\left(g^{-1}k\right)$, and the substitution $k\mapsto gk$ gives
\[
q_A(g\cdot h)=(2\pi)^{-3}\int\mathrm{d}k\,\widehat{\mathbf{b}}(k)^{*}g^TM_A(gk)\,g\,\widehat{\mathbf{b}}(k)\qquad(h\in\mathcal{S}).
\]
Hence $q_A(g\cdot h)=q_A(h)$ for every $h\in\mathcal{S}$ if and only if
\[
g^TM_A(gk)\,g=M_A(k)\qquad\text{for every }k\ne0.
\]
The condition is sufficient by the two integral formulas. It is necessary: assume that $q_A(g\cdot h)=q_A(h)$ for every $h\in\mathcal{S}$, and let $N(k):=g^TM_A(gk)g-M_A(k)$, a real symmetric matrix function that is even, bounded and continuous on $\mathbb{R}^3\setminus\{0\}$. Every $\mathbf{b}\in\mathcal{S}(\mathbb{R}^3;\mathbb{R}^3)$ is $\ast h$ for some $h\in\mathcal{S}$; so the integral of $\widehat{\mathbf{b}}^{*}N\widehat{\mathbf{b}}$ vanishes for every such $\mathbf{b}$. Let $k_0\ne0$, let $v\in\mathbb{R}^3$ with $|v|_2=1$, let $\chi\in C_c^{\infty}(\mathbb{R}^3)$ be real with $\chi\ne0$, and put $\mathbf{b}^{(n)}(x):=\chi(x/n)\cos(k_0\cdot x)\,v$ for integers $n\ge1$. As in the proof of the assertion~(3) of Section~\ref{sec:permuted-diagonal-rescalings}, the probability measures $c_n\cdot\left|\widehat{\mathbf{b}^{(n)}}(k)\right|_2^2\,\mathrm{d}k$, with normalising constants $c_n>0$, converge weakly to $\frac{1}{2}\left(\delta_{k_0}+\delta_{-k_0}\right)$ as $n\to\infty$. Since $\widehat{\mathbf{b}^{(n)}}(k)$ is a multiple of $v$, and since $N$ is even and continuous at $\pm k_0$, it follows that $v^TN(k_0)v=0$. The matrix $N(k_0)$ being symmetric, $N(k_0)=0$. Finally, the identity $g^TM_A(gk)g=M_A(k)$ implies $\operatorname{tr}M_A(gk)=\operatorname{tr}M_A(k)$, because $g$ is orthogonal.

Put $\mathrm{Sq}:=16+9+4$. For every permutation $\mathfrak{l}'$ of $\mathfrak{l}$ the numbers $y_1$, $y_2$, $y_3$ are $16$, $9$ and $4$ in some order; so $\sum_ly_l=\mathrm{Sq}$. At $k=e_1$ one has $\operatorname{tr}M_{\mathfrak{l}'}(e_1)=\check{P}\cdot\left(\mathrm{Sq}-y_1\right)$, and $y_1$ takes each of the values $16$, $9$, $4$ for exactly two of the six permutations; hence
\[
\operatorname{tr}M_A(e_1)=\check{P}\cdot\left(\mathrm{Sq}-\frac{\mathrm{Sq}}{3}\right).
\]
At $k=2^{-\frac{1}{2}}\left(e_1+e_2\right)$ one has $k_1^2=k_2^2=\frac{1}{2}$ and $k_3=0$, so that
\[
\operatorname{tr}M_{\mathfrak{l}'}(k)=\check{P}\cdot\left(\mathrm{Sq}-\frac{y_1^2+y_2^2}{y_1+y_2}\right);
\]
each two-element subset $\{y,y'\}$ of $\{16,9,4\}$ is the set $\{y_1,y_2\}$ for exactly two of the six permutations, and therefore
\[
\operatorname{tr}M_A\left(2^{-\frac{1}{2}}\left(e_1+e_2\right)\right)=\check{P}\cdot\left(\mathrm{Sq}-\frac{1}{3}\sum_{\{y,y'\}}\frac{y^2+y'^2}{y+y'}\right),
\]
the sum being taken over the three two-element subsets $\{y,y'\}$ of $\{16,9,4\}$. For $y\ne y'$ one has
\[
\frac{y^2+y'^2}{y+y'}>\frac{y+y'}{2},
\]
because $2\left(y^2+y'^2\right)-\left(y+y'\right)^2=\left(y-y'\right)^2>0$; and
\[
\frac{1}{3}\sum_{\{y,y'\}}\frac{y+y'}{2}=\frac{\mathrm{Sq}}{3},
\]
since each element of $\{16,9,4\}$ lies in exactly two of the three subsets. Hence $\operatorname{tr}M_A\left(2^{-\frac{1}{2}}(e_1+e_2)\right)<\operatorname{tr}M_A(e_1)$. Now let $g$ be the rotation of~(c): $ge_1=2^{-\frac{1}{2}}\left(e_1+e_2\right)$, $ge_2=2^{-\frac{1}{2}}\left(e_2-e_1\right)$, $ge_3=e_3$. Then $\operatorname{tr}M_A(ge_1)\ne\operatorname{tr}M_A(e_1)$, so that $g^TM_A(ge_1)g\ne M_A(e_1)$, and $q_A$ is not invariant under $g$: there is $h\in\mathcal{S}$ with $q_A(g\cdot h)\ne q_A(h)$.

Suppose that $g_{*}\mu=\mu$ and $g_{*}\left(\mu\ast_{\epsilon}\gamma_{q_A}\right)=\mu\ast_{\epsilon}\gamma_{q_A}$. For a law $\rho_1$ with~\eqref{eq:setup-G} and $g_{*}\rho_1=\rho_1$ one has, by the rules of Section~\ref{sec:euclidean-maps-actions},
\[
S_2^{\rho_1}(g\cdot h,g\cdot h)=\rho_1\left(\hat{g}\left(X(h)^2\right)\right)=S_2^{\rho_1}(h,h)\qquad(h\in\mathcal{S}).
\]
By Step~4, applied to the pair $\left(\mu,\gamma_{q_A}\right)$, and since $S_2^{\gamma_{q_A}}(h,h)=q_A(h)$,
\[
\epsilon\cdot q_A(h)=S_2^{\mu\ast_{\epsilon}\gamma_{q_A}}(h,h)-(1-\epsilon)\cdot S_2^{\mu}(h,h)\qquad(h\in\mathcal{S}).
\]
The right side does not change when $h$ is replaced by $g\cdot h$; as $\epsilon>0$, neither does $q_A(h)$, which contradicts what has just been shown. This proves~(c).

\emph{(d)} By Step~3, applied with $\rho=m_B$, for which $\nu_{\rho}$ is the point mass at $m_B$, the ergodic components of $\mu\ast_{\epsilon}m_B$ under $\mathbb{R}^3$ are the laws $\mu_{\xi}\ast_{\epsilon}m_B$, with $\xi$ distributed according to $\nu$; and for $\nu$-almost every $\xi$ both $\mu_{\xi}$ and
\[
\rho':=\mu_{\xi}\ast_{\epsilon}m_B
\]
lie in $\mathcal{C}$, satisfy~\eqref{eq:setup-G} and are ergodic under the translations of $\mathbb{R}^3$. Fix such a $\xi$. It suffices to show that $\rho'$ has the two properties stated in~(d).

Let $b_1\in C_c^{\infty}(\mathbb{R}^3)$ be real with $b_1\ne0$, let $f_1\in\mathcal{S}$ be the test with $\ast f_1=\left(b_1,0,0\right)$, and let $f':=t\cdot f_1$, where $t>0$ is so small that $\|f'\|_2\le\frac{2}{3}$. Put
\[
A:=\exp\left(i(1-\epsilon)^{\frac{1}{2}}X(f')\right),\qquad B:=\exp\left(i\epsilon^{\frac{1}{2}}X(f')\right),\qquad\check{F}:=\exp\left(iX(f')\right).
\]
For every test $h$ and all $\omega,\omega'\in\mathcal{S}'$ one has $X(h)\left(\mathrm{Tm}(\omega,\omega')\right)=(1-\epsilon)^{\frac{1}{2}}\cdot\omega(h)+\epsilon^{\frac{1}{2}}\cdot\omega'(h)$, and $\rho'$ is the image of $\mu_{\xi}\otimes m_B$ under $\mathrm{Tm}$. With $h=f'$ and $h=\tau_{se_1}\cdot f'$, Fubini's theorem gives, for $s\in\mathbb{R}$,
\[
\rho'\left(\overline{\check{F}}\cdot U(se_1)\check{F}\right)=\mu_{\xi}\left(\overline{A}\cdot U(se_1)A\right)\cdot m_B\left(\overline{B}\cdot U(se_1)B\right),\qquad\rho'(\check{F})=\mu_{\xi}(A)\cdot m_B(B).
\]

Let $\mathcal{I}^1$ be the $\sigma$-algebra of the cylinder sets that are invariant under the translations along $e_1$ (Section~\ref{sec:invariant-sets-invariant}), and let $E_{\mu_{\xi}}\left[\,\cdot\mid\mathcal{I}^1\right]$ be the conditional expectation under $\mu_{\xi}$. As $s\to\infty$ the first factor tends to $\mu_{\xi}\left(\left|E_{\mu_{\xi}}\left[A\mid\mathcal{I}^1\right]\right|^2\right)$, by Theorem~\ref{thm:invariant-vectors-clustering}(iii)4, and
\[
\mu_{\xi}\left(\left|E_{\mu_{\xi}}\left[A\mid\mathcal{I}^1\right]\right|^2\right)\ge\left|\mu_{\xi}\left(E_{\mu_{\xi}}\left[A\mid\mathcal{I}^1\right]\right)\right|^2=\left|\mu_{\xi}(A)\right|^2\ge\frac{1}{4}.
\]
For the last inequality: under $\mu_{\xi}$ the space $\mathcal{S}$ with the norm $\|\cdot\|_2$ and the map $h\mapsto X(h)$ satisfy the hypothesis~\eqref{eq:continuum-infinite-5} of Proposition~\ref{prop:sub-gaussian-random}, by~\eqref{eq:setup-G}; and the assertion~(5) of that proposition, with $z=i(1-\epsilon)^{\frac{1}{2}}$ and the test $f'$, for which $|z|\cdot\|f'\|_2\le\frac{2}{3}$, gives $\left|\mu_{\xi}(A)-1\right|\le\frac{1}{2}$.

For the second factor let $\mathbb{P}$ be the law of the random triple $u=\left(u_1,u_2,u_3\right)$ of Section~\ref{sec:free-maxwell-law-2}, and let $\mathbb{E}$ denote the expectation with respect to $\mathbb{P}$. Put $h:=\epsilon^{\frac{1}{2}}f'$, and let $r>0$ be such that the support of $b_1$ lies in $(-r,r)^3$; then $h\in\mathcal{D}$ is supported in $[-r,r]^3$, and $\ast h=\left(\epsilon^{\frac{1}{2}}\cdot t\cdot b_1,0,0\right)$, where $\epsilon^{\frac{1}{2}}\cdot t\cdot b_1\in C_c^{\infty}\left((-r,r)^3\right)$ is real and not identically zero. Let $\Phi_h(u):=\exp\left(-\frac{1}{2}q_u(h)\right)$, as in that section, and let $\bar{\Phi}(u_2,u_3)$ be the expectation of $\Phi_h(u)$ with respect to $u_1$ alone. By Section~\ref{sec:free-maxwell-law-2}(4), in which the strict inequality is~\eqref{eq:ergodic-components-10},
\[
\lim_{s\to\infty}m_B\left(\exp\left(iX(h)\right)\cdot\exp\left(iX\left(\tau_{se_1}\cdot h\right)\right)\right)=\mathbb{E}\left[\bar{\Phi}^2\right]>\left(\mathbb{E}\,\bar{\Phi}\right)^2=m_B\left(\exp\left(iX(h)\right)\right)^2=m_B(B)^2.
\]
The second factor is
\[
m_B\left(\overline{B}\cdot U(se_1)B\right)=m_B\left(\exp\left(iX(-h)\right)\cdot\exp\left(iX\left(\tau_{se_1}\cdot h\right)\right)\right).
\]
As $q_u$ is a quadratic form, $\Phi_{-h}=\Phi_h$. Hence the inequality~\eqref{eq:ergodic-components-9}, applied with $a=se_1$, $s>4r$, to the pair of tests $(-h,h)$ and to the pair $(h,h)$, shows that the second factor and the number $m_B\left(\exp\left(iX(h)\right)\cdot\exp\left(iX(\tau_{se_1}\cdot h)\right)\right)$ both differ from $\mathbb{E}\left[\Phi_h(u)\cdot\Phi_h(u^{se_1})\right]$, with the shifted triple $u^a$ of Section~\ref{sec:free-maxwell-law}, by at most the right side of~\eqref{eq:ergodic-components-9}, which is the same for the two pairs and tends to $0$ as $s\to\infty$. So the second factor tends to $\mathbb{E}\left[\bar{\Phi}^2\right]$ as $s\to\infty$.

Consequently the limit of $\rho'\left(\overline{\check{F}}\cdot U(se_1)\check{F}\right)$ as $s\to\infty$ exists, and
\[
\lim_{s\to\infty}\rho'\left(\overline{\check{F}}\cdot U(se_1)\check{F}\right)\ge\left|\mu_{\xi}(A)\right|^2\cdot\mathbb{E}\left[\bar{\Phi}^2\right]>\left|\mu_{\xi}(A)\right|^2\cdot m_B(B)^2=\left|\rho'(\check{F})\right|^2,
\]
the strict inequality because $\mu_{\xi}(A)\ne0$. So the property~(3) of Theorem~\ref{thm:invariant-vectors-clustering}(ii) fails for $\rho'$, $j=1$ and $F=G=\check{F}$; that is, $\rho'$ does not cluster along $e_1$. As $\rho'\in\mathcal{C}$ satisfies $(T_1)$, $(R_1)$ and $(RP_1)$, the equivalence of the properties~(1) and~(3) in Theorem~\ref{thm:invariant-vectors-clustering}(ii) gives $\ker\mathsf{H}_1(\rho')\ne\mathbb{C}\,\Omega_1$. The unit vector $\Omega_1$ lies in $\ker\mathsf{H}_1(\rho')$, so that $\dim\ker\mathsf{H}_1(\rho')\ge2$; and Theorem~\ref{thm:invariant-vectors-clustering}(iii)5, which applies because $\rho'\in\mathcal{C}$ is ergodic under the translations of $\mathbb{R}^3$, gives $\dim\ker\mathsf{H}_1(\rho')=\infty$.

Let $j\in\{2,3\}$, and let $w\in O_{\mathrm{cub}}$ with $we_1=e_j$. Replace $f'$ by $w\cdot f'$, the functions $A$, $B$, $\check{F}$ by
\[
A_w:=\exp\left(i(1-\epsilon)^{\frac{1}{2}}X(w\cdot f')\right),\qquad B_w:=\exp\left(i\epsilon^{\frac{1}{2}}X(w\cdot f')\right),\qquad\check{F}_w:=\exp\left(iX(w\cdot f')\right),
\]
and $e_1$ by $e_j$. The product formulas for $\rho'\left(\overline{\check{F}_w}\cdot U(se_j)\check{F}_w\right)$ and for $\rho'(\check{F}_w)$ hold as before. Since Theorem~\ref{thm:invariant-vectors-clustering}(iii)4 holds for every direction $e_j$, and since $\|w\cdot f'\|_2=\|f'\|_2\le\frac{2}{3}$, so that Proposition~\ref{prop:sub-gaussian-random}(5) applies to the test $w\cdot f'$ as before, the first factor tends to a limit that is at least $\left|\mu_{\xi}(A_w)\right|^2\ge\frac{1}{4}$. Moreover $w\cdot\left(\tau_{se_1}\cdot f'\right)=\tau_{se_j}\cdot\left(w\cdot f'\right)$, and $m_B$ is invariant under $w$, as it satisfies~\eqref{eq:setup-H} (Theorem~\ref{thm:asymmetric-components-non}); so the rules of Section~\ref{sec:euclidean-maps-actions} give
\[
m_B\left(\overline{B_w}\cdot U(se_j)B_w\right)=m_B\left(\hat{w}\left(\overline{B}\cdot U(se_1)B\right)\right)=m_B\left(\overline{B}\cdot U(se_1)B\right),\qquad m_B(B_w)=m_B\left(\hat{w}B\right)=m_B(B).
\]
Thus the second factor and the number $m_B(B_w)$ are those of the case $j=1$, and the argument given for $j=1$ shows that $\rho'$ does not cluster along $e_j$ and that $\dim\ker\mathsf{H}_j(\rho')=\infty$.

\emph{(e)} Let $w\in O_{\mathrm{cub}}\setminus N_3$, and suppose that $w_{*}\mu_{\xi}=\mu_{\xi}$ for $\nu$-almost every $\xi$. By Step~3, applied with $\rho=m_A$, for which $\nu_{\rho}$ is the uniform probability measure on the six laws $\gamma^{(\pi)}$, $\pi\in S_3$, the measure of the ergodic decomposition of $\mu\ast_{\epsilon}m_A$ is the image of $\nu\otimes\nu_{\rho}$ under the map $\left(\xi,\gamma^{(\pi)}\right)\mapsto\mu_{\xi}\ast_{\epsilon}\gamma^{(\pi)}$. Let $\xi$ be such that $w_{*}\mu_{\xi}=\mu_{\xi}$ and $\mu_{\xi}$ satisfies~\eqref{eq:setup-G}, and let $\pi\in S_3$. By Step~4 the covariance of $\mu_{\xi}\ast_{\epsilon}\gamma^{(\pi)}$ is the quadratic form
\[
h\mapsto(1-\epsilon)\cdot C_{\xi}(h,h)+\epsilon\cdot q^{(\pi)}(h)\qquad(h\in\mathcal{S}).
\]
The map $\mathrm{Tm}$ commutes with the action of $w$ (Step~2); hence
\[
w_{*}\left(\mu_{\xi}\ast_{\epsilon}\gamma^{(\pi)}\right)=\left(w_{*}\mu_{\xi}\right)\ast_{\epsilon}\left(w_{*}\gamma^{(\pi)}\right)=\mu_{\xi}\ast_{\epsilon}\gamma^{(\pi')}
\]
for some $\pi'\in S_3$ with $\pi'\ne\pi$, because $w$ permutes the six laws $\gamma^{(\pi)}$ and, as $w\notin N_3$, fixes none of them (Section~\ref{sec:permuted-diagonal-rescalings}(2), (4)). The covariance of this law is the form $h\mapsto(1-\epsilon)\cdot C_{\xi}(h,h)+\epsilon\cdot q^{(\pi')}(h)$. The six laws $\gamma^{(\pi)}$ are pairwise distinct (Section~\ref{sec:permuted-diagonal-rescalings}(3)), and $\gamma^{(\pi)}$ is the centred Gaussian law with covariance $q^{(\pi)}$; so the six forms $q^{(\pi)}$ are pairwise distinct. Hence the two covariances differ, and $w_{*}\left(\mu_{\xi}\ast_{\epsilon}\gamma^{(\pi)}\right)\ne\mu_{\xi}\ast_{\epsilon}\gamma^{(\pi)}$. Thus almost every ergodic component of $\mu\ast_{\epsilon}m_A$ is not invariant under $w$, and~(e) follows.

\emph{(f)} Fix $j\in\{1,2,3\}$ and $\rho\in\left\{\gamma_{C_0},\gamma_{q_A},m_A,m_B\right\}$. Denote by $\nu^j$ and $\nu^j_{\rho}$ the probability measures on $(\mathcal{M},\mathcal{B}_{\mathcal{M}})$ that Theorem~\ref{thm:ergodic-decomposition}(a), for $J=\{j\}$, attaches to $\mu$ and to $\rho$. In the notation of Remark~\ref{rem:components-all-translations}, with variables $\boldsymbol{\zeta},\boldsymbol{\zeta}'\in\mathcal{M}$ and with $\mu^j_{\boldsymbol{\zeta}}:=\boldsymbol{\zeta}$ and $\rho^j_{\boldsymbol{\zeta}'}:=\boldsymbol{\zeta}'$, the decompositions along $e_j$ read
\[
\mu=\int\mathrm{d}\nu^j(\boldsymbol{\zeta})\,\mu^j_{\boldsymbol{\zeta}},\qquad\rho=\int\mathrm{d}\nu^j_{\rho}(\boldsymbol{\zeta}')\,\rho^j_{\boldsymbol{\zeta}'}.
\]
The laws $\mu$ and $\rho$ are soft (for $\rho$ by Step~1), hence lie in $\mathcal{C}$ and satisfy~\eqref{eq:setup-G}. By Theorem~\ref{thm:decomposition-along-coordinate}, $\nu^j$-almost every $\mu^j_{\boldsymbol{\zeta}}$ and $\nu^j_{\rho}$-almost every $\rho^j_{\boldsymbol{\zeta}'}$ satisfies~\eqref{eq:setup-G} and $(T_j)$ and clusters along $e_j$; that is, such a law $\rho_1$ satisfies
\[
\rho_1\left(\overline{\Psi}\cdot U(se_j)\Psi'\right)\to\overline{\rho_1(\Psi)}\cdot\rho_1(\Psi')\qquad\text{as }|s|\to\infty,\ \text{for all }\Psi,\Psi'\in L^2(\rho_1).
\]

Let $\rho_1$ and $\rho_2$ be two laws with these properties, and let the translations act on $L^2(\rho_1\otimes\rho_2)$ diagonally, as in Step~3. For $\Psi_1,\Psi_1'\in L^2(\rho_1)$ and $\Psi_2,\Psi_2'\in L^2(\rho_2)$, with the scalar products of the respective $L^2$ spaces,
\[
\left\langle\Psi_1\otimes\Psi_2,U(se_j)\left(\Psi_1'\otimes\Psi_2'\right)\right\rangle=\left\langle\Psi_1,U(se_j)\Psi_1'\right\rangle\cdot\left\langle\Psi_2,U(se_j)\Psi_2'\right\rangle\to\left\langle\Psi_1\otimes\Psi_2,1\right\rangle\cdot\left\langle1,\Psi_1'\otimes\Psi_2'\right\rangle
\]
as $|s|\to\infty$. The finite linear combinations of products $\Psi_1\otimes\Psi_2$ are dense in $L^2(\rho_1\otimes\rho_2)$, both sides are sesquilinear, and $\|U(se_j)\|=1$; hence $\left\langle H,U(se_j)H'\right\rangle\to\langle H,1\rangle\cdot\langle1,H'\rangle$ for all $H,H'\in L^2(\rho_1\otimes\rho_2)$. If $U(se_j)H=H$ for every $s\in\mathbb{R}$, then $\|H\|^2=\left\langle H,U(se_j)H\right\rangle\to\left|\langle1,H\rangle\right|^2$, and $H$ is constant, as $\|H\|^2=\left|\langle1,H\rangle\right|^2+\left\|H-\langle1,H\rangle\cdot1\right\|^2$. So $\rho_1\otimes\rho_2$ is ergodic under the translations along $e_j$. The law $\rho_1\ast_{\epsilon}\rho_2$ is invariant under these translations, as $\mathrm{Tm}$ commutes with them; and if $H\in L^2(\rho_1\ast_{\epsilon}\rho_2)$ is invariant under them, then $H\circ\mathrm{Tm}$ is an invariant element of $L^2(\rho_1\otimes\rho_2)$, hence constant, and so is $H$. Therefore $\rho_1\ast_{\epsilon}\rho_2\in\mathcal{M}^j_{\mathrm{erg}}$.

Let $\Sigma$ be the cylinder $\sigma$-algebra of $\mathcal{S}'$. For every $B\in\Sigma\otimes\Sigma$ the function $(\rho_1,\rho_2)\mapsto(\rho_1\otimes\rho_2)(B)$ on $\mathcal{M}\times\mathcal{M}$ is measurable for $\mathcal{B}_{\mathcal{M}}\otimes\mathcal{B}_{\mathcal{M}}$. Indeed, for a rectangle $B=A_1\times A_2$ it is the product of the measurable functions $\rho_1\mapsto\rho_1(A_1)$ and $\rho_2\mapsto\rho_2(A_2)$; the sets $B$ with this property form a Dynkin system; and $\Sigma\otimes\Sigma$ is the $\sigma$-algebra generated by the rectangles, which form a class that is closed under finite intersections; so Dynkin's lemma applies. As $\mathrm{Tm}$ is measurable, the choice $B=\mathrm{Tm}^{-1}(A)$ shows that $(\rho_1,\rho_2)\mapsto\left(\rho_1\ast_{\epsilon}\rho_2\right)(A)$ is measurable for every cylinder set $A$. Since $\mathcal{B}_{\mathcal{M}}$ is the $\sigma$-algebra generated by the maps $\rho_1\mapsto\rho_1(A)$, the map $(\rho_1,\rho_2)\mapsto\rho_1\ast_{\epsilon}\rho_2$ of $\mathcal{M}\times\mathcal{M}$ into $\mathcal{M}$ is measurable. By Fubini's theorem the probability measures
\[
\mu\otimes\rho\qquad\text{and}\qquad\int\mathrm{d}\nu^j(\boldsymbol{\zeta})\int\mathrm{d}\nu^j_{\rho}(\boldsymbol{\zeta}')\,\mu^j_{\boldsymbol{\zeta}}\otimes\rho^j_{\boldsymbol{\zeta}'}
\]
on $\Sigma\otimes\Sigma$ agree on the rectangles, hence, by Dynkin's lemma, on $\Sigma\otimes\Sigma$. Evaluating both at $\mathrm{Tm}^{-1}(A)$ one obtains
\[
\left(\mu\ast_{\epsilon}\rho\right)(A)=\int\mathrm{d}\nu^j(\boldsymbol{\zeta})\int\mathrm{d}\nu^j_{\rho}(\boldsymbol{\zeta}')\,\left(\mu^j_{\boldsymbol{\zeta}}\ast_{\epsilon}\rho^j_{\boldsymbol{\zeta}'}\right)(A)\qquad\text{for every cylinder set }A.
\]
So the image $\nu''$ of $\nu^j\otimes\nu^j_{\rho}$ under the map $(\boldsymbol{\zeta},\boldsymbol{\zeta}')\mapsto\mu^j_{\boldsymbol{\zeta}}\ast_{\epsilon}\rho^j_{\boldsymbol{\zeta}'}$ satisfies $\left(\mu\ast_{\epsilon}\rho\right)(A)=\int\mathrm{d}\nu''(\rho'')\,\rho''(A)$ for every cylinder set $A$, and $\nu''\left(\mathcal{M}^j_{\mathrm{erg}}\right)=1$ by the preceding paragraph. The law $\mu\ast_{\epsilon}\rho$ is soft by~(a), hence invariant under the translations along $e_j$; by the uniqueness assertion of Theorem~\ref{thm:ergodic-decomposition}(a) for $J=\{j\}$, the measure $\nu''$ is the measure of the decomposition of $\mu\ast_{\epsilon}\rho$ along $e_j$.

For $\nu^j\otimes\nu^j_{\rho}$-almost every $(\boldsymbol{\zeta},\boldsymbol{\zeta}')$ the laws $\mu^j_{\boldsymbol{\zeta}}$ and $\rho^j_{\boldsymbol{\zeta}'}$ both satisfy~\eqref{eq:setup-G} (Theorem~\ref{thm:decomposition-along-coordinate}), and by Lemma~\ref{lem:moment-bounds}(a) they are then centred. For every such pair Step~4 bounds the change of the fourth cumulant; at the test $f=g_0$, and by the choice of $\epsilon_0$, it gives
\[
\left|\kappa_4^{\mu^j_{\boldsymbol{\zeta}}\ast_{\epsilon}\rho^j_{\boldsymbol{\zeta}'}}(f)-\kappa_4^{\mu^j_{\boldsymbol{\zeta}}}(f)\right|\le45\,\epsilon\cdot\|f\|_2^4\le\boldsymbol{\delta}\cdot e^{4+\sigma}.
\]
Combined with the hypothesis of~(f), this gives $\left|\kappa_4^{\rho''}(f)-e^4\cdot b(f)\right|\le(C+\boldsymbol{\delta})\cdot e^{4+\sigma}$ for $\nu''$-almost every $\rho''$, which is~(f) for $\mu\ast_{\epsilon}\rho$. Finally, let $\mu'$ be the law of~(b), and let $\nu^j_{\mu_{\epsilon}}$ and $\nu^j_{\mu'}$ be the measures of Theorem~\ref{thm:ergodic-decomposition} for the laws $\mu_{\epsilon}$ and $\mu'$, respectively, and $J=\{j\}$. As in the proof of~(b), the uniqueness assertion of Theorem~\ref{thm:ergodic-decomposition}(a) gives $\nu^j_{\mu'}=\frac{1}{2}\left(\nu^j+\nu^j_{\mu_{\epsilon}}\right)$. The inequality holds with the constant $C$ for $\nu^j$-almost every component and, by the case $\rho=\gamma_{C_0}$ just treated, with the constant $C+\boldsymbol{\delta}$ for $\nu^j_{\mu_{\epsilon}}$-almost every component; hence it holds with $C+\boldsymbol{\delta}$ for $\nu^j_{\mu'}$-almost every component.
\end{proof}

\begin{corollary}[Five properties that the class does not imply]
Let the charge $e$ in the definition of $\mathbf{K}(C)$ satisfy $e<e_2$. Then there are numbers $\boldsymbol{\delta}>0$ with
\[
\frac{5}{2}\left(C_{\mathrm{kap}}+\boldsymbol{\delta}\right)\cdot e^{\sigma}\le\frac{1}{4}\,b(f);
\]
fix such a number $\boldsymbol{\delta}$. Every law in $\mathbf{K}(C_{\mathrm{kap}}+\boldsymbol{\delta})$ has each property that Theorem~\ref{thm:main-theorem} and Corollary~\ref{cor:ergodic-components}(a) assert of a limit point, with $C_{\mathrm{kap}}+\boldsymbol{\delta}$ in place of $C_{\mathrm{kap}}$, except the property~\eqref{it:main-theorem-P8}. Let $\mu\in\mathbf{K}(C_{\mathrm{kap}})$; for instance, $\mu$ may be a limit point of Theorem~\ref{thm:main-theorem}, formed with the charge $e$. Then $\mu\in\mathbf{K}(C_{\mathrm{kap}}+\boldsymbol{\delta})$, and $\mathbf{K}(C_{\mathrm{kap}}+\boldsymbol{\delta})$ contains
\begin{enumerate}
\item[(1)] a law different from $\mu$;
\item[(2)] a law that is not ergodic under the translations of $\mathbb{R}^3$ and for which none of the conditions~\eqref{eq:fourth-cumulants-C1}, \eqref{eq:fourth-cumulants-C2}, \eqref{eq:fourth-cumulants-C3} and~\eqref{eq:fourth-cumulants-M} of Theorem~\ref{thm:clustering-mixing-ergodicity} holds;
\item[(3)] a law that is not invariant under the rotation by the angle $\frac{\pi}{4}$ about the third coordinate axis;
\item[(4)] a law almost every ergodic component $\rho'$ of which satisfies, for $j=1,2,3$, $\dim\ker\mathsf{H}_j(\rho')=\infty$ and does not cluster along $e_j$;
\item[(5)] for every $w\in O_{\mathrm{cub}}\setminus N_3$, a law for which it is not the case that almost every ergodic component is invariant under $w$.
\end{enumerate}
So none of the following five properties holds for all laws in $\mathbf{K}(C_{\mathrm{kap}}+\boldsymbol{\delta})$: equality with $\mu$; ergodicity under the translations of $\mathbb{R}^3$; invariance under all rotations of $\mathbb{R}^3$; the equality $\dim\ker\mathsf{H}_j(\rho')=1$ for $j=1,2,3$ and almost every ergodic component $\rho'$; invariance of almost every ergodic component under every element of $O_{\mathrm{cub}}$. If moreover $j\in\{1,2,3\}$ and
\[
\left|\kappa_4^{\rho'}(f)-e^4\cdot b(f)\right|\le C_{\mathrm{kap}}\cdot e^{4+\sigma}\qquad\text{for almost every component }\rho'\text{ of }\mu\text{ along }e_j,
\]
then the laws in (1)--(5) can be chosen such that, for each of them, almost every component $\rho'$ along $e_j$ satisfies $\left|\kappa_4^{\rho'}(f)-e^4\cdot b(f)\right|\le\left(C_{\mathrm{kap}}+\boldsymbol{\delta}\right)\cdot e^{4+\sigma}$.
\end{corollary}

\begin{proof}
Since $\sigma>0$ and $e<e_2$, the inequality~\eqref{eq:free-coefficient-4} and $b(f)>0$ give $\frac{5}{2}\,C_{\mathrm{kap}}\cdot e^{\sigma}<\frac{1}{4}\,b(f)$; hence the numbers $\boldsymbol{\delta}$ exist. If a constant $C>0$ satisfies $\frac{5}{2}\,C\cdot e^{\sigma}\le\frac{1}{4}\,b(f)$, then every law of $\mathbf{K}(C)$ has each property that Theorem~\ref{thm:main-theorem} and Corollary~\ref{cor:ergodic-components}(a) assert of a limit point, with $C$ in place of $C_{\mathrm{kap}}$, except the property~\eqref{it:main-theorem-P8} (Section~\ref{sec:soft-laws-four}). By the choice of $\boldsymbol{\delta}$ the constant $C=C_{\mathrm{kap}}+\boldsymbol{\delta}$ satisfies this inequality; for this constant the preceding statement is the assertion of the corollary on all laws of $\mathbf{K}(C_{\mathrm{kap}}+\boldsymbol{\delta})$. Every limit point of Theorem~\ref{thm:main-theorem}, formed with the charge $e$, lies in $\mathbf{K}(C_{\mathrm{kap}})$, again by Section~\ref{sec:soft-laws-four}. The conditions~\eqref{eq:free-coefficient-K} and~\eqref{eq:free-coefficient-K-2} become weaker when the constant increases; hence $\mathbf{K}(C_{\mathrm{kap}})\subset\mathbf{K}(C_{\mathrm{kap}}+\boldsymbol{\delta})$.

Let $\mu\in\mathbf{K}(C_{\mathrm{kap}})$. By~\eqref{eq:free-coefficient-K} at $g_0=f$ and by~\eqref{eq:free-coefficient-4},
\[
\kappa_4^{\mu}(f)\ge e^4\cdot\left(b(f)-C_{\mathrm{kap}}\cdot e^{\sigma}\right)\ge\frac{9}{10}\,e^4\cdot b(f)>0.
\]
Let $\epsilon_0$ be the number of Proposition~\ref{prop:operations-class} for this law $\mu$, for $C=C_{\mathrm{kap}}$ and for $\boldsymbol{\delta}$, let $0<\epsilon\le\epsilon_0$, and let $\mu'=\frac{1}{2}\left(\mu+\mu\ast_{\epsilon}\gamma_{C_0}\right)$. By Proposition~\ref{prop:operations-class}(a) and~(b), the five laws
\[
\mu\ast_{\epsilon}\gamma_{C_0},\qquad\mu',\qquad\mu\ast_{\epsilon}\gamma_{q_A},\qquad\mu\ast_{\epsilon}m_B,\qquad\mu\ast_{\epsilon}m_A
\]
lie in $\mathbf{K}(C_{\mathrm{kap}}+\boldsymbol{\delta})$. By Proposition~\ref{prop:operations-class}(a), the law $\mu\ast_{\epsilon}\gamma_{C_0}$ is a law as in~(1); by Proposition~\ref{prop:operations-class}(b), $\mu'$ is a law as in~(2); by Proposition~\ref{prop:operations-class}(c), one of the two laws $\mu$ and $\mu\ast_{\epsilon}\gamma_{q_A}$ is a law as in~(3); by Proposition~\ref{prop:operations-class}(d), $\mu\ast_{\epsilon}m_B$ is a law as in~(4); and, for $w\in O_{\mathrm{cub}}\setminus N_3$, one of the two laws $\mu$ and $\mu\ast_{\epsilon}m_A$ is a law as in~(5), by Proposition~\ref{prop:operations-class}(e). Since $\mu$ and these laws belong to $\mathbf{K}(C_{\mathrm{kap}}+\boldsymbol{\delta})$, none of the five properties listed after~(5) holds for all laws of this class. Under the additional hypothesis on the components of $\mu$ along $e_j$, Proposition~\ref{prop:operations-class}(f) gives the asserted inequality for the five laws just named, and the law $\mu$ satisfies it because $C_{\mathrm{kap}}\le C_{\mathrm{kap}}+\boldsymbol{\delta}$.
\end{proof}

\subsubsection{Uniqueness}\label{sec:uniqueness}

\begin{remark}[Distinct laws with the properties of a limit point]
Let $C>0$ and $\boldsymbol{\delta}>0$, let $\epsilon_0$ be the number of Proposition~\ref{prop:operations-class}, and let $\mu\in\mathbf{K}(C)$ satisfy $\kappa_4^{\mu}(f)\ne0$. Every limit point $\mu$ of Theorem~\ref{thm:main-theorem} is such a law, with $C=C_{\mathrm{kap}}$: it lies in $\mathbf{K}(C_{\mathrm{kap}})$ (Section~\ref{sec:soft-laws-four}), and $\kappa_4^{\mu}(f)\ge\frac{1}{2}e^4\cdot b(f)>0$ by Theorem~\ref{thm:main-theorem}(iv). Let $0<\epsilon\le\epsilon_0$. By Proposition~\ref{prop:operations-class}(a), the law $\mu\ast_{\epsilon}\gamma_{C_0}$ lies in $\mathbf{K}(C+\boldsymbol{\delta})$ and is different from $\mu$. The law $\mu$ lies in $\mathbf{K}(C+\boldsymbol{\delta})$ as well, since the inequalities~\eqref{eq:free-coefficient-K} and~\eqref{eq:free-coefficient-K-2} with the constant $C$ imply those with the constant $C+\boldsymbol{\delta}$. Hence $\mathbf{K}(C+\boldsymbol{\delta})$ contains two different laws: the conditions that define this class do not determine a law.
\end{remark}

\begin{example}[A lattice family with two limit points]
Let $0\le a<b\le1$. For a cutoff $(N,T)$ let $F$ be distributed according to the averaged law at the charge $q=0$ on $\Lambda_{\eta,R}$. The family~\eqref{eq:continuum-infinite-E3} of Proposition~\ref{prop:scaled-free-laws} assigns to the cutoff $(N,T)$ the law of $\sqrt{a}\cdot F$ if $N$ is even, and the law of $\sqrt{b}\cdot F$ if $N$ is odd. By Proposition~\ref{prop:scaled-free-laws} this family is admissible in the sense of Definition~\ref{def:admissible-families}: at every cutoff its member is invariant under the lattice translations and under $O_{\mathrm{cub}}$, which is~\eqref{eq:continuum-infinite-Ad1}, it satisfies the Gaussian domination~\eqref{eq:continuum-infinite-Ad2}, and it is reflection positive in the sense of~\eqref{eq:continuum-infinite-Ad3}. By Proposition~\ref{prop:scaled-free-laws}(iii) the family has two different limit points.
\end{example}

\begin{remark}[The width of the bounds on a limit point]
The assertions~(ii) and~(iii) of Theorem~\ref{thm:main-theorem} are inequalities. By~(iii), for every limit point $\mu$ and every $g\in\{g_0,g_{+},g_{-}\}$ the number $\kappa_4^{\mu}(g)$ lies in the interval of length $2C_{\mathrm{kap}}\cdot e^{4+\sigma}$ centred at $e^4\cdot b(g)$, and $S_2^{\mu}(f,f)$ lies in the interval of length $2C_{\mathrm{kap}}\cdot e^{2-\mathfrak{e}_2}$ centred at $\|f\|_2^2$; these intervals are the same for all limit points. The bounds of Theorem~\ref{thm:free-coefficient-2}, from which~(ii) and~(iii) are derived in Section~\ref{sec:proof-part-ii-2}, hold on each torus $k\in\{1\}\cup\{L,L+1,\dots\}$ separately, and their right-hand sides do not depend on $k$.
\end{remark}

\subsubsection{Examples concerning ergodicity, clustering and convex classes}\label{sec:examples-concerning-ergodicity}

\begin{remark}[A non-ergodic law of the class]
Let $C>0$ and $\boldsymbol{\delta}>0$, let $\epsilon_0$ be the number of Proposition~\ref{prop:operations-class}, let $\mu\in\mathbf{K}(C)$ satisfy $\kappa_4^{\mu}(f)\ne0$ (every limit point does, with $C=C_{\mathrm{kap}}$; Section~\ref{sec:uniqueness}), and let $0<\epsilon\le\epsilon_0$. By Proposition~\ref{prop:operations-class}(b) the law
\[
\frac{1}{2}\left(\mu+\mu\ast_{\epsilon}\gamma_{C_0}\right)
\]
lies in $\mathbf{K}(C+\boldsymbol{\delta})$, it is not ergodic under the translations of $\mathbb{R}^3$, and none of the properties~\eqref{eq:fourth-cumulants-C1}, \eqref{eq:fourth-cumulants-C2}, \eqref{eq:fourth-cumulants-C3} and~\eqref{eq:fourth-cumulants-M} of Theorem~\ref{thm:clustering-mixing-ergodicity} holds for it.
\end{remark}

\begin{example}[Laws that do not cluster]
The class of the soft laws, which is defined by the conditions of $\mathbf{K}(C)$ other than~\eqref{eq:free-coefficient-K} and~\eqref{eq:free-coefficient-K-2}, contains simpler laws without clustering, for instance the following two. Let $\beta_1\in C_c^{\infty}(\mathbb{R}^3;\Lambda^2)$ be real, with $\varphi:=dd^{*}\beta_1\ne0$.
\begin{enumerate}
\item[(1)] Let $0\le a<b\le1$, and let $\mu_{a,b}$ be a limit point of the family~\eqref{eq:continuum-infinite-E2} of Proposition~\ref{prop:scaled-free-laws}, which assigns to every cutoff the mean of the laws of $\sqrt{a}\cdot F$ and of $\sqrt{b}\cdot F$, $F$ being distributed according to the averaged law at $q=0$. The family is admissible (Proposition~\ref{prop:scaled-free-laws}), so that Theorem~\ref{thm:continuum-infinite-1} applies to $\mu_{a,b}$. By its parts~(1), (5), (6) and~(7), the law $\mu_{a,b}$ satisfies~\eqref{eq:setup-G}, is invariant under the translations of $\mathbb{R}^3$ and under $O_{\mathrm{cub}}$, and satisfies~\eqref{eq:setup-RPj} for $j=1,2,3$; that is, it satisfies~\eqref{eq:setup-H}. By part~(4) it is concentrated on $\{dF=0\}$. Every member of the family is invariant under $F\mapsto-F$: the averaged law at $q=0$ is invariant under this map (Lemma~\ref{lem:invariances-averaged-law}(c)); the maps $F\mapsto\sqrt{c}\cdot F$, $c\ge0$, commute with it, so that the images of that law under them are invariant as well; and so is the mean of two such images. Hence part~(9) shows that $\mu_{a,b}$ is invariant under $F\mapsto-F$. Thus $\mu_{a,b}$ is soft. By Proposition~\ref{prop:scaled-free-laws}(ii), the moments of $\mu_{a,b}$ satisfy
\[
S_4(\varphi,\varphi,\tau_x\varphi,\tau_x\varphi)-S_2(\varphi,\varphi)^2\ge\frac{(b-a)^2}{4}\cdot\|\varphi\|_2^4>0\qquad\text{for every }x\in\mathbb{R}^3.
\]
\item[(2)] Let $\rho$ be a probability measure on $[0,1]$ which is not a point mass, and let
\[
\mu^{\rho}:=\int\mathrm{d}\rho(s)\,\gamma_{sC_0}
\]
be the law of Theorem~\ref{thm:properties-scale-mixtures}, that is, the law of $s^{\frac{1}{2}}\cdot F$ when $F$ is distributed according to $\gamma_{C_0}$ and $s$, independent of $F$, according to $\rho$. By Theorem~\ref{thm:properties-scale-mixtures}(i), the law $\mu^{\rho}$ satisfies~\eqref{eq:setup-H}, it is invariant under the group $\mathbb{R}^3\rtimes O(3)$ of all Euclidean maps, and $X(d^{*}h)=0$ almost surely for every Schwartz three-form $h$, hence for all such $h$ simultaneously, by the argument of Section~\ref{sec:closed-distributions} with a countable dense set; so $\mu^{\rho}$ is concentrated on $\{dF=0\}$. Its ergodic components $\gamma_{sC_0}$ are centred Gaussian laws (Theorem~\ref{thm:properties-scale-mixtures}(ii)), so that $\mu^{\rho}$ is invariant under $F\mapsto-F$. Thus $\mu^{\rho}$ is soft. Its odd moments vanish; therefore
\[
S_4(\varphi,\varphi,\tau_x\varphi,\tau_x\varphi)=\kappa_4(\varphi,\varphi,\tau_x\varphi,\tau_x\varphi)+S_2(\varphi,\varphi)\cdot S_2(\tau_x\varphi,\tau_x\varphi)+2S_2(\varphi,\tau_x\varphi)^2,
\]
and the formulas of Theorem~\ref{thm:properties-scale-mixtures}(iii) for $S_2$ and $\kappa_4$, together with $C_0(\tau_x\varphi,\tau_x\varphi)=C_0(\varphi,\varphi)=\|\varphi\|_2^2$ (Lemma~\ref{lem:free-transverse-covariance}(a), (b); the test $\varphi=d(d^{*}\beta_1)$ is closed), give
\[
S_4(\varphi,\varphi,\tau_x\varphi,\tau_x\varphi)-S_2(\varphi,\varphi)^2=\operatorname{Var}_{\rho}(s)\cdot\|\varphi\|_2^4+2E_{\rho}\left[s^2\right]\cdot C_0(\varphi,\tau_x\varphi)^2\ge\operatorname{Var}_{\rho}(s)\cdot\|\varphi\|_2^4>0
\]
for every $x\in\mathbb{R}^3$.
\end{enumerate}
In both cases the lower bound is a positive number that does not depend on $x$, and for every $\ell>0$ the mean of the left-hand side over $x\in[0,\ell]^3$ is at least this number. Hence the property~\eqref{eq:fourth-cumulants-C3} of Theorem~\ref{thm:clustering-mixing-ergodicity} fails for $m=n=2$ and $f_1=f_2=g_1=g_2=\varphi$. That theorem applies to both laws, as they are invariant under the translations and satisfy~\eqref{eq:setup-G}; so neither $\mu_{a,b}$ nor $\mu^{\rho}$ is ergodic.

Neither law has the properties~(iv) and~(vi) of Theorem~\ref{thm:main-theorem}. Indeed, each of them is a mixture of centred Gaussian laws which are invariant under the translations, satisfy~\eqref{eq:setup-G} and are mixing, so that they are its ergodic components. For $\mu^{\rho}$ these laws are the laws $\gamma_{sC_0}$, $0\le s\le1$ (Theorem~\ref{thm:properties-scale-mixtures}(ii); Lemma~\ref{lem:gaussian-laws-bounded}(a), (c) and Lemma~\ref{lem:free-transverse-covariance}(a), (b)). For $\mu_{a,b}$, the proof of Proposition~\ref{prop:scaled-free-laws} shows that $\mu_{a,b}=\frac{1}{2}\left(\mu_a+\mu_b\right)$, where $\mu_c$ is the law of $\sqrt{c}\cdot F$ under a limit point $\mu_1$ of the family~\eqref{eq:continuum-infinite-E1}; the law $\mu_1$ is centred and Gaussian (Proposition~\ref{prop:scaled-free-laws}(i)), satisfies~\eqref{eq:setup-G} and is invariant under the translations (Theorem~\ref{thm:continuum-infinite-1}(1), (5)); hence so is $\mu_c$ for $0\le c\le1$, and $\mu_c$ is mixing by Lemma~\ref{lem:gaussian-laws-bounded}(c). A limit point is not a mixture of translation-invariant Gaussian laws which satisfy~\eqref{eq:setup-G}, by Theorem~\ref{thm:main-theorem}(iv), and almost every ergodic component of a limit point is not Gaussian, by Theorem~\ref{thm:main-theorem}(vi).
\end{example}

\begin{remark}[A mixture with a Gaussian component]
Let $\tilde{\mu}$ be a law on $\mathcal{S}'$ which satisfies~\eqref{eq:fourth-cumulants-H1}--\eqref{eq:fourth-cumulants-H3}, is ergodic and is not Gaussian, let $\tilde{C}:=S_2^{\tilde{\mu}}$, and let $0<p<1$. By Proposition~\ref{prop:gaussian-components-four}(ii), the law
\[
\mu^{(p)}:=(1-p)\cdot\tilde{\mu}+p\cdot\gamma_{\tilde{C}}
\]
satisfies~\eqref{eq:fourth-cumulants-H1}--\eqref{eq:fourth-cumulants-H3}, and its ergodic decomposition has the two components $\tilde{\mu}$ and $\gamma_{\tilde{C}}$, with the weights $1-p$ and $p$. Since the Gaussian component $\gamma_{\tilde{C}}$ has the positive weight $p$, the law $\mu^{(p)}$ does not have the property~(vi) of Theorem~\ref{thm:main-theorem}.
\end{remark}

\begin{lemma}[Convex sets of ergodic laws]
Let $\mathcal{K}$ be a convex set of laws on $\mathcal{S}'$ which are invariant under the translations of $\mathbb{R}^3$, and assume that every element of $\mathcal{K}$ is ergodic under $\mathbb{R}^3$. Then $\mathcal{K}$ has at most one element.
\end{lemma}

\begin{proof}
Let $\mu_1,\mu_2\in\mathcal{K}$, and put $\mu_{12}:=\frac{1}{2}\left(\mu_1+\mu_2\right)$. By convexity $\mu_{12}\in\mathcal{K}$; so $\mu_{12}$ is ergodic under $\mathbb{R}^3$, that is, the constants are the only elements of $L^2(\mu_{12})$ which are invariant under all translations (Section~\ref{sec:invariant-sets-invariant}). As $\mu_1\le2\mu_{12}$, the law $\mu_1$ is absolutely continuous with respect to $\mu_{12}$; by the Radon--Nikodym theorem it has a density $h$ with respect to $\mu_{12}$, and $0\le h\le2$ $\mu_{12}$-almost everywhere. Let $a\in\mathbb{R}^3$, and let $T_a$ be the translation by $a$, acting on $\mathcal{S}'$. For every cylinder set $A$ the invariance of $\mu_1$ and of $\mu_{12}$ under $T_a$ gives
\[
\mu_1(A)=\mu_1\left(T_a^{-1}A\right)=\int\mathrm{d}\mu_{12}\,\left(\mathbf{1}_A\circ T_a\right)\cdot h=\int\mathrm{d}\mu_{12}\,\mathbf{1}_A\cdot\left(h\circ T_a^{-1}\right).
\]
So $h\circ T_a^{-1}$ is also a density of $\mu_1$ with respect to $\mu_{12}$, and $h\circ T_a^{-1}=h$ $\mu_{12}$-almost everywhere. Thus $h$, as an element of $L^2(\mu_{12})$, is invariant under all translations, and it is constant. Since $\int\mathrm{d}\mu_{12}\,h=\mu_1(\mathcal{S}')=1$, one has $h=1$, that is, $\mu_1=\mu_{12}$. In the same way $\mu_2=\mu_{12}$; hence $\mu_1=\mu_2$.
\end{proof}

\begin{remark}[Convex classes of soft laws defined by polynomial moment bounds]
Let $c_2,c_4\in\mathbb{R}$ and $\mathsf{u}_2,\mathsf{u}_4\ge0$, and let $\mathrm{Pm}$ and $\hat{\Gamma}_0=\frac{R^{*}}{L}\cdot\mathbb{Z}^3$, $R^{*}=R^{*}(e)$, be the polynomial and the lattice of Step~6 of the proof of Theorem~\ref{thm:main-theorem} (Section~\ref{sec:proof-part-ii-2}). Let $\mathcal{K}(c,\mathsf{u})$ be the set of the soft laws $\mu$ which satisfy, for $m=2,4$, every integer $J\ge1$ and all pairwise distinct $y_1,\dots,y_J\in\hat{\Gamma}_0$,
\[
\left|\int\mathrm{d}\mu\,\prod_{i=1}^J\mathrm{Pm}\left(X(\tau_{y_i}f);\|f\|_2^2,c_m\right)\right|\le\mathsf{u}_m^J.
\]
For $c_m=c_m^{\infty}$ and for the numbers $\mathsf{u}_m$ of Step~6 these inequalities are the inequalities~\eqref{eq:free-coefficient-23}, because $U(y)\,\mathrm{Pm}\left(X(f);\|f\|_2^2,c_m\right)=\mathrm{Pm}\left(X(\tau_yf);\|f\|_2^2,c_m\right)$ for $y\in\mathbb{R}^3$; so every limit point lies in the corresponding set $\mathcal{K}(c^{\infty},\mathsf{u})$.

Each of the conditions that define $\mathcal{K}(c,\mathsf{u})$, the conditions that define a soft law as well as the inequalities above, is linear in the law: it is an equation between linear functionals of the law, or a bound for such a functional or for its modulus. Hence $\mathcal{K}(c,\mathsf{u})$ is convex.

Assume that
\[
\mathsf{u}_4\le C\cdot e^{\sigma},\qquad\left(1+|c_2|\right)\cdot\mathsf{u}_2\le C\cdot e^{\sigma},\qquad|c_2|\le C_B\cdot e^{-\mathfrak{e}_2},\qquad\left|c_4-3c_2^2-b(f)\right|\le C_B\cdot e^{\sigma},
\]
with the constants $C$ and $C_B$ of Step~6; these are the bounds which Step~6 establishes for the numbers $c_m^{\infty}$ and $\mathsf{u}_m$. Then the inequalities of Theorem~\ref{thm:main-theorem}(iii) for $g=f$, the two inequalities of Theorem~\ref{thm:main-theorem}(vi) and the assertions of Corollary~\ref{cor:ergodic-components} hold for every $\mu\in\mathcal{K}(c,\mathsf{u})$. Indeed, $\mu$ satisfies~\eqref{eq:setup-G}, so that $X(f)$ is centred under $\mu$ (Lemma~\ref{lem:moment-bounds}(a)); for $J=1$ the inequalities above and the identity of Step~2 of the proof of Theorem~\ref{thm:free-coefficient-2} (Section~\ref{sec:proof-cumulant-bounds}), applied under $\mu$ with $\|f\|_2^2$ in place of $\boldsymbol{\sigma}^2$, give
\[
\left|e^{-2}\left(S_2^{\mu}(f,f)-\|f\|_2^2\right)-c_2\right|\le\mathsf{u}_2,\qquad\left|e^{-4}\kappa_4^{\mu}(f)-\left(c_4-3c_2^2\right)\right|\le\mathsf{u}_4+3\mathsf{u}_2\cdot\left(2|c_2|+\mathsf{u}_2\right),
\]
from which the inequalities of~(iii) for $g=f$ follow in the same way as the inequalities of~(vi) follow from the corresponding inequalities for $\mu_{\xi}$ at the end of Step~6. The derivation of~(vi) in Step~6 and the proof of Corollary~\ref{cor:ergodic-components} in Section~\ref{sec:proofs-corollaries-main} apply to $\mu$ without change: they start from the inequalities~\eqref{eq:free-coefficient-23}, from the bounds above and from the properties~\eqref{it:main-theorem-P1}--\eqref{it:main-theorem-P7} and~(c) of Theorem~\ref{thm:main-theorem}, which every soft law has (Lemma~\ref{lem:properties-soft-laws}).

By the preceding lemma, if every law in $\mathcal{K}(c,\mathsf{u})$ is ergodic under the translations of $\mathbb{R}^3$, then $\mathcal{K}(c,\mathsf{u})$ has at most one element.
\end{remark}

\begin{remark}[Forms of clustering]
Let $\mu$ be a soft law, for instance a limit point of Theorem~\ref{thm:main-theorem} (Section~\ref{sec:soft-laws-four}), and let $\mu=\int\mathrm{d}\nu(\xi)\,\mu_{\xi}$ be its ergodic decomposition under $\mathbb{R}^3$. As $\mu$ is invariant under the translations and satisfies~\eqref{eq:setup-G}, Theorem~\ref{thm:clustering-mixing-ergodicity} applies to it:
\[
\eqref{eq:fourth-cumulants-C1}\Leftrightarrow\eqref{eq:fourth-cumulants-M}\Rightarrow\eqref{eq:fourth-cumulants-C2}\Rightarrow\eqref{eq:fourth-cumulants-Erg}\Leftrightarrow\eqref{eq:fourth-cumulants-C3}\Leftrightarrow\eqref{eq:fourth-cumulants-T}.
\]
In particular, if $\mu$ clusters along one sequence of translations, as in~\eqref{eq:fourth-cumulants-C2}, or in the mean, as in~\eqref{eq:fourth-cumulants-C3}, then $\nu$ is a point mass, which is~\eqref{eq:fourth-cumulants-T}. Conversely, let $\nu$ be a point mass. Then $\mu$ is ergodic and clusters in the mean; and since $\mu$ lies in the class $\mathcal{C}$ of Section~\ref{sec:properties-laws} and is ergodic under $\mathbb{R}^3$, it clusters as $\min_l|a_l|\to\infty$ and in the mean over cubes, by Theorem~\ref{thm:invariant-vectors-clustering}(iii)1, 2. Clustering as $|a|\to\infty$ in every manner, that is~\eqref{eq:fourth-cumulants-C1}, does not follow from~\eqref{eq:fourth-cumulants-T}: the soft law $m_B$ is ergodic under $\mathbb{R}^3$ and does not cluster along $e_j$ (Section~\ref{sec:vacuum-ergodic-components}), so that~\eqref{eq:fourth-cumulants-T} holds for it, whereas~\eqref{eq:fourth-cumulants-M}, and with it~\eqref{eq:fourth-cumulants-C1}, fails.
\end{remark}

\subsubsection{The vacuum of the ergodic components}\label{sec:vacuum-ergodic-components}

\begin{remark}[Components with an infinite-dimensional vacuum]
Let $C>0$ and $\boldsymbol{\delta}>0$, let $\epsilon_0$ be the number of Proposition~\ref{prop:operations-class}, let $\mu\in\mathbf{K}(C)$, and let $0<\epsilon\le\epsilon_0$. The law $\mu\ast_{\epsilon}m_B$ lies in $\mathbf{K}(C+\boldsymbol{\delta})$ (Proposition~\ref{prop:operations-class}(a)), and for $j=1,2,3$ almost every ergodic component $\rho'$ of it under $\mathbb{R}^3$ satisfies $\dim\ker\mathsf{H}_j(\rho')=\infty$ and does not cluster along $e_j$ (Proposition~\ref{prop:operations-class}(d)). If the conditions~\eqref{eq:free-coefficient-K} and~\eqref{eq:free-coefficient-K-2} are dropped, the law $m_B$ itself is an example: it is soft (Step~1 of the proof of Proposition~\ref{prop:operations-class}(a)) and ergodic under $\mathbb{R}^3$, and for $j=1,2,3$ one has $\dim\ker\mathsf{H}_j(m_B)=\infty$, and $m_B$ does not cluster along $e_j$ (Theorem~\ref{thm:asymmetric-components-non}(B)).
\end{remark}

\begin{remark}[The vacuum of the components of a limit point]
Let $\mu$ be a limit point of Theorem~\ref{thm:main-theorem}, let $\mu=\int\mathrm{d}\nu(\xi)\,\mu_{\xi}$ be its ergodic decomposition under $\mathbb{R}^3$, and let $j\in\{1,2,3\}$. By Theorem~\ref{thm:main-theorem}(c), for $\nu$-almost every $\xi$ the vectors of $\mathscr{H}_j(\mu_{\xi})$ which are invariant under the time translations and under the space translations are the multiples of $\Omega_j$, and $\dim\ker\mathsf{H}_j(\mu_{\xi})$ is $1$ or $\infty$. Let $\nu^j$ be the measure of Theorem~\ref{thm:ergodic-decomposition} for $\mu$ and $J=\{j\}$; it gives the decomposition of $\mu$ along the axis $\mathbb{R}e_j$, whose components are required to be invariant only under the translations along $e_j$. For this decomposition, Corollary~\ref{cor:ergodic-components}(b) gives $\ker\mathsf{H}_j(\rho)=\mathbb{C}\,\Omega_j$, and clustering along $e_j$, for $\nu^j$-almost every component $\rho$.
\end{remark}

\subsubsection{Exchanges of the axes}\label{sec:exchanges-axes}

\begin{remark}[Components that are not invariant under an exchange of the axes]
Let $C>0$ and $\boldsymbol{\delta}>0$, let $\epsilon_0$ be the number of Proposition~\ref{prop:operations-class}, let $\mu\in\mathbf{K}(C)$, $0<\epsilon\le\epsilon_0$ and $w\in O_{\mathrm{cub}}\setminus N_3$. The law $\mu\ast_{\epsilon}m_A$ lies in $\mathbf{K}(C+\boldsymbol{\delta})$ by Proposition~\ref{prop:operations-class}(a), and so does $\mu$, since $\mathbf{K}(C)\subset\mathbf{K}(C+\boldsymbol{\delta})$ (Section~\ref{sec:uniqueness}). By Proposition~\ref{prop:operations-class}(e), for at least one of the two laws the ergodic components under $\mathbb{R}^3$ which are not invariant under $w$ form a set of positive measure. If the conditions~\eqref{eq:free-coefficient-K} and~\eqref{eq:free-coefficient-K-2} are dropped, the law $m_A$ is an example: it is soft (Step~1 of the proof of Proposition~\ref{prop:operations-class}(a)), its ergodic decomposition is the uniform distribution on six distinct laws, and none of these laws is invariant under an element of $O_{\mathrm{cub}}\setminus N_3$ (Theorem~\ref{thm:asymmetric-components-non}(A)).
\end{remark}

\begin{remark}[Symmetries of the components of a limit point]
Let $\mu$ be a limit point of Theorem~\ref{thm:main-theorem}, and let $\mu=\int\mathrm{d}\nu(\xi)\,\mu_{\xi}$ be its ergodic decomposition under $\mathbb{R}^3$. For $\nu$-almost every $\xi$ the law $\mu_{\xi}$ is invariant under the group $\mathbb{R}^3\rtimes N_3$ (Theorem~\ref{thm:main-theorem}(c)). Let $w\in O_{\mathrm{cub}}$. The components $w_{*}\mu_{\xi}$ have under $\nu$ the same distribution as the components $\mu_{\xi}$, and for $\nu$-almost every $\xi$ the law $w_{*}\mu_{\xi}$ depends on $w$ only through its class in $O_{\mathrm{cub}}/N_3\cong S_3$ (Theorem~\ref{thm:cubic-symmetry-components}(ii)); thus $S_3$ permutes the components and preserves $\nu$. The inequalities of Theorem~\ref{thm:main-theorem}(vi) hold for $\nu$-almost every component; as the components $w_{*}\mu_{\xi}$ have the same distribution as the components $\mu_{\xi}$, these inequalities hold also for $w_{*}\mu_{\xi}$, for $\nu$-almost every $\xi$. By Theorem~\ref{thm:main-theorem}(vii) one has $\mathrm{Ext}^{\xi}_{x_0}(f)\ge\frac{1}{2}e^4\cdot b(f)$ for $\nu$-almost every $\xi$; hence the set
\[
\left\{\xi:\mathrm{Ext}^{\xi}_{x_0}(f)\ge\frac{1}{4}e^4\cdot b(f)\right\}
\]
of Theorem~\ref{thm:main-theorem}(v) has full $\nu$-measure.
\end{remark}

\subsubsection{Rotations}\label{sec:rotations}

\begin{remark}[Laws of the class that are not invariant under rotations]
Let $C>0$ and $\boldsymbol{\delta}>0$, let $\epsilon_0$ be the number of Proposition~\ref{prop:operations-class}, let $\mu\in\mathbf{K}(C)$, and let $0<\epsilon\le\epsilon_0$. The law $\mu\ast_{\epsilon}\gamma_{q_A}$ lies in $\mathbf{K}(C+\boldsymbol{\delta})$ by Proposition~\ref{prop:operations-class}(a), and so does $\mu$, since $\mathbf{K}(C)\subset\mathbf{K}(C+\boldsymbol{\delta})$ (Section~\ref{sec:uniqueness}). By Proposition~\ref{prop:operations-class}(c), the two laws are not both invariant under the rotation by the angle $\frac{\pi}{4}$ about the third coordinate axis, which fixes $e_3$ and maps $e_1$ to $2^{-\frac{1}{2}}\left(e_1+e_2\right)$. Hence at least one of them is not invariant under $O(3)$.
\end{remark}

\begin{example}[Soft laws that are not invariant under rotations]
The law $\gamma_{q_A}$ is a soft centred Gaussian law (Step~1 of the proof of Proposition~\ref{prop:operations-class}(a)); it is mixing by Lemma~\ref{lem:gaussian-laws-bounded}(c), as its covariance $q_A$ is invariant under the translations. It is not invariant under $O(3)$, because $q_A$ is not invariant under the rotation of Proposition~\ref{prop:operations-class}(c), as the proof of that part shows.

The law $m_A$ is not invariant under $O(3)$ either. Indeed, $m_A$ satisfies~\eqref{eq:setup-H} (Theorem~\ref{thm:asymmetric-components-non}) and therefore lies in the class $\mathcal{C}$ of Section~\ref{sec:properties-laws}. If $g_{*}m_A=m_A$ held for every $g\in O(3)$, then, by Theorem~\ref{thm:vacuum-symmetries-components}(iv), almost every ergodic component of $m_A$ under $\mathbb{R}^3$ would be invariant under the group $\mathbb{R}^3\rtimes O(3)$, in particular under $O_{\mathrm{cub}}$. But $m_A$ has six ergodic components, each with the weight $\frac{1}{6}$, and none of them is invariant under an element of $O_{\mathrm{cub}}\setminus N_3$ (Theorem~\ref{thm:asymmetric-components-non}(A)).
\end{example}

\begin{remark}[Reflection positivity at a diagonal plane]
Let $\mu$ satisfy~\eqref{eq:setup-H}, and let $\mu$ be reflection positive, in the sense of Section~\ref{sec:properties-laws}, at the plane $\{x_1=x_2\}$. Then, by Theorem~\ref{thm:vacuum-symmetries-components}(vi), almost every ergodic component of $\mu$ under $\mathbb{R}^3$ is invariant under the group $\mathbb{R}^3\rtimes O_{\mathrm{cub}}$, in particular under the exchanges of the axes; moreover $P^k_{\mu}=P^{123}_{\mu}$, and almost every component $\rho$ satisfies $\ker\mathsf{H}_k(\rho)=\mathbb{C}\,\Omega_k$, for $k=1,2,3$.
\end{remark}

\subsubsection{The pressure}\label{sec:pressure}

\begin{proposition}[The pressure at a fixed cutoff]
\begin{enumerate}
\item[(1)] Let the data be those of Section~\ref{sec:numbers-attached-reference}, with $\boldsymbol{\mu}_N>-\frac{1}{2}$, and for integers $k\ge1$ let
\[
p_k:=(k\cdot n_{*})^{-3}\log Z_k.
\]
Then
\[
p_1-\frac{D_L}{(L-1)^3\cdot n_{*}^3}\le p_k\le p_1\qquad\text{for every integer }k\ge L,
\]
and the sequence $\left(p_{L^i}\right)_{i\ge0}$ is non-increasing and converges. In the situation of Remark~\ref{rem:cumulant-bounds-pressure}, the numbers $\bar{p}_k=(k\cdot n_{*})^{-3}\log\mathbf{Z}_k$ satisfy
\[
p_1-\frac{\bar{\mathfrak{D}}_L}{(L-1)^3\cdot n_{*}^3}\le\bar{p}_k\le p_k\qquad\text{for every integer }k\ge L.
\]
Divided by $\eta^3$, that is, per unit of physical volume, the lengths of the intervals that these inequalities give for $p_k$ and for $\bar{p}_k$ are
\[
\frac{D_L}{(L-1)^3\cdot(R^{*})^3}=\frac{\delta\cdot N^{*}}{(R^{*})^3}\qquad\text{and}\qquad\frac{\bar{\mathfrak{D}}_L}{(L-1)^3\cdot(R^{*})^3}.
\]
Let $e\in(0,e_2]$, $N\ge N_0(e)$, $T^{*}=T^{*}(e)$ and $\boldsymbol{\mu}_N=\bar{m}\cdot\eta+m_{0,N}(e)$. Then both lengths are at most
\[
\frac{1}{(R^{*})^3}\cdot\left(a_L\cdot\log\frac{2\pi}{e^2\cdot R^{*}}+\frac{K_L+C_{\mathrm{mat}}+1}{(L-1)^3}\right),\qquad R^{*}=R^{*}(e),
\]
and this bound does not depend on $N$.
\item[(2)] Let $q>0$ and $\boldsymbol{\mu}>-\frac{1}{2}$, and let $n\ge2$ be an integer. On the cubic torus $n$ put
\[
x_n:=\frac{2\pi^2}{q^2\cdot n},\qquad\mathbf{M}_n:=\left\lfloor\left(\frac{\log\left(12\,\mathrm{thf}(x_n)^3\right)}{x_n}\right)^{\frac{1}{2}}\right\rfloor,
\]
with the function $\mathrm{thf}$ of Section~\ref{sec:trace-formula-chessboard}. Then
\[
\log\mathbf{Z}_n\le\log Z_n\le\log\mathbf{Z}_n+\log\left(2\left(2\mathbf{M}_n+1\right)^3\right).
\]
If $q$ and $\boldsymbol{\mu}$ are fixed, then $\log\left(2\left(2\mathbf{M}_n+1\right)^3\right)=O(\log n)$ as $n\to\infty$; hence $n^{-3}\left(\log Z_n-\log\mathbf{Z}_n\right)\to0$, and $n^{-3}\log Z_n$ converges as $n\to\infty$ if and only if $n^{-3}\log\mathbf{Z}_n$ converges, the two limits being equal.
\end{enumerate}
\end{proposition}

\begin{proof}
(1) For positive integers $a,b,c$ let $n_{*}(a,b,c)$ be the torus $\left(a\cdot n_{*},b\cdot n_{*},c\cdot n_{*}\right)$, and let $\phi(a,b,c):=\log Z_{n_{*}(a,b,c)}$. As $n_{*}\ge2$, the function $\phi$ has the properties stated in~\eqref{it:free-coefficient-Q6}, and $\phi(k,k,k)=\log Z_k$. Applying the inequality $\phi(ka,b,c)\le k\cdot\phi(a,b,c)$ and its analogues in the second and in the third variable, one obtains
\[
\phi(ka,ka,ka)\le k^3\cdot\phi(a,a,a)\qquad\text{for all integers }k,a\ge1.
\]
For $a=1$ this is $\log Z_k\le k^3\log Z_1$, that is, $p_k\le p_1$; for $a=L^i$ and $k=L$ it is $\log Z_{L^{i+1}}\le L^3\log Z_{L^i}$, that is, $p_{L^{i+1}}\le p_{L^i}$. By~\eqref{it:free-coefficient-Q7} with $j=L$, one has $D_L=L^3\phi(1,1,1)-\phi(L,L,L)\ge0$ and, for every integer $k\ge L$,
\[
\log Z_k\ge k^3\log Z_1-\left(\frac{k-1}{L-1}\right)^3\cdot D_L\ge k^3\cdot\left(\log Z_1-\frac{D_L}{(L-1)^3}\right),
\]
which is~\eqref{eq:free-coefficient-15} for $n^{*}=n_{*}$ and $j=L$; division by $(k\cdot n_{*})^3$ gives the lower bound for $p_k$. So the non-increasing sequence $\left(p_{L^i}\right)_{i\ge0}$ is bounded below by $p_1-\frac{D_L}{(L-1)^3\cdot n_{*}^3}$, and it converges. The inequalities for $\bar{p}_k$ are proved in the last paragraph of Remark~\ref{rem:cumulant-bounds-pressure}, by means of Proposition~\ref{prop:zero-flux-sector}(c) and of the inequality $\mathbf{Z}_k\le Z_k$.

As $n_{*}\cdot\eta=R^{*}$, the division of $\frac{D_L}{(L-1)^3\cdot n_{*}^3}$ and of $\frac{\bar{\mathfrak{D}}_L}{(L-1)^3\cdot n_{*}^3}$ by $\eta^3$ gives the two lengths stated in~(1), and $\frac{D_L}{(L-1)^3}=\delta\cdot N^{*}$ by the definition of $\delta$ in Section~\ref{sec:numbers-attached-reference}. Let now $e\in(0,e_2]$, $N\ge N_0(e)$, $T^{*}=T^{*}(e)$ and $\boldsymbol{\mu}_N=\bar{m}\cdot\eta+m_{0,N}(e)$. By Step~0 of the proof of Theorem~\ref{thm:main-theorem} (Section~\ref{sec:proof-part-ii-2}), Theorem~\ref{thm:free-coefficient-2} applies, with the probe $\beta$ of Section~\ref{sec:probes}; so~\eqref{eq:free-coefficient-9} holds, and $D^{\mathrm{mat}}\le C_{\mathrm{mat}}$ by part~\ref{it:free-coefficient-3-1} of Theorem~\ref{thm:free-coefficient-3}. Hence
\[
\delta\cdot N^{*}\le a_L\cdot\log\frac{2\pi}{e^2\cdot R^{*}}+\frac{K_L+C_{\mathrm{mat}}+1}{(L-1)^3}.
\]
In the same situation Remark~\ref{rem:cumulant-bounds-pressure} shows that $\frac{\bar{\mathfrak{D}}_L}{(L-1)^3\cdot N^{*}}$ is at most the right-hand side of~\eqref{eq:free-coefficient-9}; so $\frac{\bar{\mathfrak{D}}_L}{(L-1)^3}$ obeys the same bound. Division by $(R^{*})^3$ gives the assertion; the numbers $R^{*}=R^{*}(e)$, $a_L$, $K_L$ and $C_{\mathrm{mat}}$ do not depend on $N$.

(2) On the cubic torus $n$ one has $|n|=n^3$ and $n_{\iota}=n$, and $\frac{r_q^2}{2n}=\frac{2\pi^2}{q^2\cdot n}=x_n$; so the definition of $\mathfrak{q}_w(n)$ in Section~\ref{sec:wilson-determinant-lattice} gives
\[
\mathfrak{q}_w(n)=\exp\left(-x_n\cdot|w|_2^2\right)\qquad\left(w\in\mathbb{Z}^3\right).
\]
Let $\mathbf{M}\ge0$ be an integer, and let $|w|$ denote the largest modulus of a component of $w$. The set of the $w\in\mathbb{Z}^3$ with $|w|>\mathbf{M}$ is the union of the three sets on which one given component of $w$ has modulus larger than $\mathbf{M}$. On such a set the sum of $\mathfrak{q}_w(n)$ is the product of the factor $\mathrm{thf}(x_n)^2$, which comes from the two other components, with $2\sum_{i>\mathbf{M}}\exp\left(-x_n\,i^2\right)$, and the latter is at most $2\exp\left(-x_n\cdot(\mathbf{M}+1)^2\right)\cdot\mathrm{thf}(x_n)$ by the second elementary bound at the end of Section~\ref{sec:recalled-results}. Therefore
\[
\sum_{w\in\mathbb{Z}^3,\ |w|>\mathbf{M}}\mathfrak{q}_w(n)\le6\,\mathrm{thf}(x_n)^3\cdot\exp\left(-x_n\cdot(\mathbf{M}+1)^2\right).
\]
As $\mathrm{thf}(x_n)\ge1$, the number $\log\left(12\,\mathrm{thf}(x_n)^3\right)$ is positive, and $\mathbf{M}=\mathbf{M}_n$ satisfies $x_n\cdot(\mathbf{M}_n+1)^2>\log\left(12\,\mathrm{thf}(x_n)^3\right)$; for this $\mathbf{M}$ the right-hand side is smaller than $\frac{1}{2}$. Proposition~\ref{prop:comparison-flux-sectors}(ii) gives $Z_n\le2\left(2\mathbf{M}_n+1\right)^3\cdot\mathbf{Z}_n$. Moreover $\mathbf{Z}_n=\boldsymbol{\omega}_n^{(0)}(1)\le\sum_w\boldsymbol{\omega}_n^{(w)}(1)=Z_n$, the functionals $\boldsymbol{\omega}_n^{(w)}$ being positive (Theorem~\ref{thm:free-coefficient-1}(d2)). As $Z_n$ and $\mathbf{Z}_n$ are positive (Section~\ref{sec:wilson-determinant-lattice}), the two inequalities follow on taking logarithms.

Let $q$ and $\boldsymbol{\mu}$ be fixed. The first elementary bound at the end of Section~\ref{sec:recalled-results} gives
\[
1\le\mathrm{thf}(x_n)\le1+\left(\frac{\pi}{x_n}\right)^{\frac{1}{2}}=1+\left(\frac{q^2\cdot n}{2\pi}\right)^{\frac{1}{2}},
\]
so that $\log\left(12\,\mathrm{thf}(x_n)^3\right)=O(\log n)$ and
\[
\mathbf{M}_n\le\left(\frac{q^2\cdot n}{2\pi^2}\cdot\log\left(12\,\mathrm{thf}(x_n)^3\right)\right)^{\frac{1}{2}}=O\left((n\log n)^{\frac{1}{2}}\right)
\]
as $n\to\infty$. Hence $\log\left(2\left(2\mathbf{M}_n+1\right)^3\right)=O(\log n)$, and
\[
0\le n^{-3}\left(\log Z_n-\log\mathbf{Z}_n\right)\le n^{-3}\log\left(2\left(2\mathbf{M}_n+1\right)^3\right)\to0.
\]
The last assertion of~(2) follows.
\end{proof}

\begin{example}[A source that distinguishes the two laws]
Let $q>0$ and $\boldsymbol{\mu}>-\frac{1}{2}$, and let $n\ge2$ be an integer. On the cubic torus $n$ put
\[
\mathrm{Sq}:=\sum_{x\in\Lambda_n}F^u_{12}(x),\qquad x_n:=\frac{r_q^2}{2n}=\frac{2\pi^2}{q^2\cdot n},
\]
where $F^u_{12}(x)$ is the value of $F^u$ on the plaquette at $x$ spanned by $e_1$ and $e_2$. Since $|\mathrm{Sq}|\le n^{\frac{3}{2}}\cdot\|F^u\|$, the function $\exp(\mathfrak{l}\cdot\mathrm{Sq})$ belongs to $\mathfrak{F}(n)$ for every real $\mathfrak{l}$.
\begin{enumerate}
\item[(a)] On the support of each functional $\boldsymbol{\omega}_n^{(w)}$, $w\in\mathbb{Z}^3$, one has $\mathrm{Sq}=r_q\cdot n\cdot w_{12}$. Consequently $E^{\mathrm{av}}_n\exp(\mathfrak{l}\cdot\mathrm{Sq})=1$ for every real $\mathfrak{l}$; this holds also when the factor $\mathrm{Dt}(qA)$ in the definition of $\boldsymbol{\omega}_n$ (Section~\ref{sec:wilson-determinant-lattice}) is replaced by $1$.
\item[(b)] Let $\boldsymbol{\omega}_n$, $\boldsymbol{\omega}_n^{(w)}$, $Z_n$, $\mathbf{Z}_n$ and $E^{\mathrm{af}}_n$ be formed with the factor $\mathrm{Dt}(qA)$ replaced by $1$. Then, for every real $\mathfrak{l}$,
\[
E^{\mathrm{af}}_n\exp(\mathfrak{l}\cdot\mathrm{Sq})=\exp\left(\frac{1}{2}\mathfrak{l}^2\cdot n^3\right)\cdot\frac{\sum_{i\in\mathbb{Z}}\exp\left(-x_n\left(i-\frac{\mathfrak{l}\cdot n^2}{r_q}\right)^2\right)}{\sum_{i\in\mathbb{Z}}\exp\left(-x_n\,i^2\right)},
\]
and, for fixed $q$ and $\mathfrak{l}$,
\[
\lim_{n\to\infty}n^{-3}\log E^{\mathrm{af}}_n\exp(\mathfrak{l}\cdot\mathrm{Sq})=\frac{1}{2}\mathfrak{l}^2.
\]
\item[(c)] Assume that $\boldsymbol{\mu}>0$, and let the functionals be formed with the factor $\mathrm{Dt}(qA)=\det\left(D[qA]+\boldsymbol{\mu}\right)$ of~\eqref{eq:free-coefficient-5}. Then, for fixed $q$, $\boldsymbol{\mu}$ and $\mathfrak{l}$,
\[
\liminf_{n\to\infty}n^{-3}\log E^{\mathrm{af}}_n\exp(\mathfrak{l}\cdot\mathrm{Sq})\ge\frac{1}{2}\mathfrak{l}^2-4\log\frac{6+\boldsymbol{\mu}}{\boldsymbol{\mu}},
\]
and the right-hand side is positive if $\mathfrak{l}^2>8\log\frac{6+\boldsymbol{\mu}}{\boldsymbol{\mu}}$.
\end{enumerate}
Thus the pressures $n^{-3}\log E^{\mathrm{nat}}_n\exp(\mathfrak{l}\cdot\mathrm{Sq})$, $\mathrm{nat}\in\{\mathrm{af},\mathrm{av}\}$, of the two laws with the source $\mathfrak{l}\cdot\mathrm{Sq}$ behave differently as $n\to\infty$: for the harmonic-averaged law the pressure is $0$, whereas for the all-flux law its limit is $\frac{1}{2}\mathfrak{l}^2>0$ in the case~(b) if $\mathfrak{l}\ne0$, and its limit inferior is positive in the case~(c) if $\mathfrak{l}^2>8\log\frac{6+\boldsymbol{\mu}}{\boldsymbol{\mu}}$.
\end{example}

\begin{proof}
(a) Let $A\in C^1(n)$. Then
\[
\sum_{x\in\Lambda_n}(dA)_{12}(x)=\sum_{x\in\Lambda_n}\left(A_2(x+e_1)-A_2(x)-A_1(x+e_2)+A_1(x)\right).
\]
On the right-hand side, $\sum_xA_2(x+e_1)=\sum_xA_2(x)$ and $\sum_xA_1(x+e_2)=\sum_xA_1(x)$, as $x\mapsto x+e_{\mu}$ is a bijection of $\Lambda_n$; hence $\sum_x(dA)_{12}(x)=0$, and in the same way $\sum_x(dA)_{\mu\nu}(x)=0$ for all $\mu<\nu$. Thus, for every pair $\mu<\nu$, each exact real two-cochain has the total $0$ on the set of the plaquettes spanned by $e_{\mu}$ and $e_{\nu}$. Let $\mathfrak{z}\in C^2(n;\mathbb{Z})$ be closed, with flux $w$. The sum of $\mathfrak{z}_{12}$ over each of the $n$ planes $\{x_3=j\}$ is $w_{12}$ (Section~\ref{sec:wilson-determinant-lattice}), so that $\sum_x\mathfrak{z}_{12}(x)=n\cdot w_{12}$. Hence $\mathrm{Sq}=r_q\cdot n\cdot w_{12}$ at $F^u=dA+r_q\cdot\mathfrak{z}$, that is, on the support of $\boldsymbol{\omega}_n^{(w)}$, whatever factor stands in the place of $\mathrm{Dt}(qA)$. In $E^{\mathrm{av}}_n=\boldsymbol{\omega}_n^{(0)}/\mathbf{Z}_n$ only the flux $0$ occurs; there $\mathrm{Sq}=0$, and $E^{\mathrm{av}}_n\exp(\mathfrak{l}\cdot\mathrm{Sq})=\boldsymbol{\omega}_n^{(0)}(1)/\mathbf{Z}_n=1$. Here $\mathbf{Z}_n>0$: with the factor $\mathrm{Dt}(qA)$ by Section~\ref{sec:wilson-determinant-lattice}, since $\boldsymbol{\mu}>-\frac{1}{2}>-1$, and with the factor $1$ because the integrand of $\boldsymbol{\omega}_n^{(0)}(1)$ is then positive.

(b) Let $w\in\mathbb{Z}^3$, and let $h_w$ be the constant two-cochain of Section~\ref{sec:decomposition-flux} on the cubic torus $n$; its components are $(h_w)_{\mu\nu}(x)=r_q\cdot w_{\mu\nu}/n^2$, and $\|h_w\|^2=r_q^2\cdot|w|_2^2/n=2x_n\cdot|w|_2^2$ (Lemma~\ref{lem:flux-sectors}(a)). Let $\mathfrak{z}_w\in C^2(n;\mathbb{Z})$ be the cochain with $(\mathfrak{z}_w)_{\mu\nu}(x)=w_{\mu\nu}$ if $x_{\mu}=x_{\nu}=0$ and $(\mathfrak{z}_w)_{\mu\nu}(x)=0$ otherwise ($\mu<\nu$); it is closed and its flux is $w$ (Section~\ref{sec:decomposition-flux}). By Lemma~\ref{lem:flux-sectors}(a) there is a real one-cochain $B_w$ with $r_q\cdot\mathfrak{z}_w=dB_w+h_w$, and $r_q\cdot\mathfrak{z}'$ is exact for every closed $\mathfrak{z}'\in C^2(n;\mathbb{Z})$ with flux $0$. The closed integer valued two-cochains with flux $w$ are the cochains $\mathfrak{z}_w+\mathfrak{z}'$ with $\mathfrak{z}'$ closed, integer valued and with flux $0$; for them
\[
dA+r_q\cdot\left(\mathfrak{z}_w+\mathfrak{z}'\right)=\left(d(A+B_w)+r_q\cdot\mathfrak{z}'\right)+h_w,
\]
where the first summand is exact and therefore orthogonal to the constant two-cochain $h_w$, its sums over the plaquettes of each orientation being $0$ by~(a). Consequently, with the factor $1$ in place of $\mathrm{Dt}(qA)$,
\[
\boldsymbol{\omega}_n^{(w)}(1)=\exp\left(-\frac{1}{2}\|h_w\|^2\right)\cdot\int_{[0,r_q)^{\mathrm{links}}}\mathrm{d}A\,\psi(A+B_w),\qquad\psi(A):=\sum_{\mathfrak{z}'}\exp\left(-\frac{1}{2}\left\|dA+r_q\cdot\mathfrak{z}'\right\|^2\right),
\]
the sum being taken over the closed $\mathfrak{z}'\in C^2(n;\mathbb{Z})$ with flux $0$. For an integer valued one-cochain $\ell$ the cochain $d\ell$ is closed, integer valued and has flux $0$, by~(a); so $\mathfrak{z}'\mapsto\mathfrak{z}'+d\ell$ is a bijection of the set of summation, and $\psi(A+r_q\cdot\ell)=\psi(A)$. Thus $\psi$ is periodic with respect to the lattice $r_q\cdot C^1(n;\mathbb{Z})$, of which $[0,r_q)^{\mathrm{links}}$ is a fundamental domain, and the integral of $\psi(\,\cdot\,+B_w)$ over $[0,r_q)^{\mathrm{links}}$ equals the integral of $\psi$, which is $\boldsymbol{\omega}_n^{(0)}(1)=\mathbf{Z}_n$. Therefore
\[
\boldsymbol{\omega}_n^{(w)}(1)=\exp\left(-x_n\cdot|w|_2^2\right)\cdot\mathbf{Z}_n=\mathfrak{q}_w(n)\cdot\mathbf{Z}_n.
\]
By~(a), $\boldsymbol{\omega}_n^{(w)}\left(\exp(\mathfrak{l}\cdot\mathrm{Sq})\right)=\exp\left(\mathfrak{l}\cdot r_q\cdot n\cdot w_{12}\right)\cdot\boldsymbol{\omega}_n^{(w)}(1)$. Sum over $w\in\mathbb{Z}^3$; the series converge absolutely. Writing $|w|_2^2$ as the sum of $w_{12}^2$ and of the squares of the other two components of $w$, one finds that the sums over these two components contribute the same factor $\mathrm{thf}(x_n)^2$ to $\boldsymbol{\omega}_n\left(\exp(\mathfrak{l}\cdot\mathrm{Sq})\right)$ and to $Z_n=\boldsymbol{\omega}_n(1)$, and that both numbers contain the factor $\mathbf{Z}_n$. These factors cancel in the quotient, and
\[
E^{\mathrm{af}}_n\exp(\mathfrak{l}\cdot\mathrm{Sq})=\frac{\sum_{i\in\mathbb{Z}}\exp\left(\mathfrak{l}\cdot r_q\cdot n\cdot i-x_n\,i^2\right)}{\sum_{i\in\mathbb{Z}}\exp\left(-x_n\,i^2\right)}.
\]
Completing the square,
\[
\mathfrak{l}\cdot r_q\cdot n\cdot i-x_n\,i^2=-x_n\left(i-\frac{\mathfrak{l}\cdot n^2}{r_q}\right)^2+\frac{1}{2}\mathfrak{l}^2\cdot n^3,
\]
because $\frac{\mathfrak{l}\cdot r_q\cdot n}{2x_n}=\frac{\mathfrak{l}\cdot n^2}{r_q}$ and $\frac{\left(\mathfrak{l}\cdot r_q\cdot n\right)^2}{4x_n}=\frac{1}{2}\mathfrak{l}^2\cdot n^3$; this is the display of~(b).

For every real $y$ the sum $\sum_{i\in\mathbb{Z}}\exp\left(-x_n(i-y)^2\right)$ has a term with $|i-y|\le\frac{1}{2}$, and is therefore at least $\exp\left(-\frac{1}{4}x_n\right)$. The integers $i\ge y$, and likewise the integers $i<y$, can be enumerated as $i_0,i_1,\dots$ with $|i_l-y|\ge l$; hence the sum is at most $2\sum_{l\ge0}\exp\left(-x_n\,l^2\right)=1+\mathrm{thf}(x_n)\le2+\mathrm{thf}(x_n)$. Together with $1\le\mathrm{thf}(x_n)\le1+\left(\frac{\pi}{x_n}\right)^{\frac{1}{2}}$ (Section~\ref{sec:recalled-results}) this shows that the quotient of the two sums in the display of~(b) lies between
\[
\exp\left(-\frac{1}{4}x_n\right)\cdot\left(1+\left(\frac{\pi}{x_n}\right)^{\frac{1}{2}}\right)^{-1}\qquad\text{and}\qquad3+\left(\frac{\pi}{x_n}\right)^{\frac{1}{2}}.
\]
For fixed $q$ one has $x_n\to0$ and $\frac{\pi}{x_n}=\frac{q^2\cdot n}{2\pi}$; so the logarithm of the quotient is $O(\log n)$ as $n\to\infty$, and $n^{-3}\log E^{\mathrm{af}}_n\exp(\mathfrak{l}\cdot\mathrm{Sq})\to\frac{1}{2}\mathfrak{l}^2$.

(c) Let $B$ be a real one-cochain. For every spinor $\psi$, the Cauchy--Schwarz inequality and Theorem~\ref{thm:free-coefficient-1}(d2) give
\[
\left\|\left(D[B]+\boldsymbol{\mu}\right)\psi\right\|\cdot\|\psi\|\ge\operatorname{Re}\left\langle\psi,\left(D[B]+\boldsymbol{\mu}\right)\psi\right\rangle\ge\boldsymbol{\mu}\cdot\|\psi\|^2,
\]
so that every singular value of $D[B]+\boldsymbol{\mu}$ is at least $\boldsymbol{\mu}$. Let $T_{\mu}$, $\mu=1,2,3$, be the unitary operators of the proof of Theorem~\ref{thm:free-coefficient-1}(d2) (Section~\ref{sec:symmetries-determinant-lattice}),
\[
\left(T_{\mu}\psi\right)(x)=\exp\left(iB(x,x+e_{\mu})\right)\cdot\psi(x+e_{\mu}),\qquad\left(T_{\mu}^{*}\psi\right)(x)=\exp\left(iB(x,x-e_{\mu})\right)\cdot\psi(x-e_{\mu}).
\]
In~\eqref{eq:free-coefficient-5} the neighbour $x+e_{\mu}$ contributes $-\frac{1}{2}\left(1-\gamma_{\mu}\right)\left(T_{\mu}\psi\right)(x)$ and the neighbour $x-e_{\mu}$ contributes $-\frac{1}{2}\left(1+\gamma_{\mu}\right)\left(T_{\mu}^{*}\psi\right)(x)$; hence
\[
D[B]=3-\sum_{\mu=1}^3\check{U}_{\mu},\qquad\check{U}_{\mu}:=\frac{1}{2}\left(1-\gamma_{\mu}\right)T_{\mu}+\frac{1}{2}\left(1+\gamma_{\mu}\right)T_{\mu}^{*}.
\]
The matrices $\frac{1}{2}\left(1-\gamma_{\mu}\right)$ and $\frac{1}{2}\left(1+\gamma_{\mu}\right)$ are complementary orthogonal projections, as $\gamma_{\mu}$ is Hermitian with $\gamma_{\mu}^2=1$, and $T_{\mu}$ commutes with them, since it acts on the site variable only. Hence the two cross terms in $\check{U}_{\mu}^{*}\check{U}_{\mu}$ contain the product of the two projections, which is $0$, and
\[
\check{U}_{\mu}^{*}\check{U}_{\mu}=T_{\mu}^{*}\cdot\frac{1}{2}\left(1-\gamma_{\mu}\right)\cdot T_{\mu}+T_{\mu}\cdot\frac{1}{2}\left(1+\gamma_{\mu}\right)\cdot T_{\mu}^{*}=\frac{1}{2}\left(1-\gamma_{\mu}\right)+\frac{1}{2}\left(1+\gamma_{\mu}\right)=1,
\]
and $\check{U}_{\mu}$, an operator on the finite-dimensional space of the spinors, is unitary. Therefore $\left\|D[B]+\boldsymbol{\mu}\right\|\le3+3+\boldsymbol{\mu}=6+\boldsymbol{\mu}$, and every singular value of $D[B]+\boldsymbol{\mu}$ is at most $6+\boldsymbol{\mu}$. The spinors form a space of dimension $4n^3$, and $\mathrm{Dt}(qA)=\left|\det\left(D[qA]+\boldsymbol{\mu}\right)\right|$ for real $A$, by Theorem~\ref{thm:free-coefficient-1}(d2); so $\mathrm{Dt}(qA)$ is the product of the $4n^3$ singular values of $D[qA]+\boldsymbol{\mu}$, and
\[
\boldsymbol{\mu}^{4n^3}\le\mathrm{Dt}(qA)\le\left(6+\boldsymbol{\mu}\right)^{4n^3}\qquad\text{for every real }A\in C^1(n).
\]
The integrands in $\boldsymbol{\omega}_n\left(\exp(\mathfrak{l}\cdot\mathrm{Sq})\right)$ and in $Z_n=\boldsymbol{\omega}_n(1)$ are non-negative. Bounding the factor $\mathrm{Dt}(qA)$ from below by $\boldsymbol{\mu}^{4n^3}$ in the numerator, and from above by $\left(6+\boldsymbol{\mu}\right)^{4n^3}$ in the denominator, of $E^{\mathrm{af}}_n\exp(\mathfrak{l}\cdot\mathrm{Sq})=\boldsymbol{\omega}_n\left(\exp(\mathfrak{l}\cdot\mathrm{Sq})\right)/Z_n$, one finds that this number is at least $\left(\frac{\boldsymbol{\mu}}{6+\boldsymbol{\mu}}\right)^{4n^3}$ times the number $E^{\mathrm{af}}_n\exp(\mathfrak{l}\cdot\mathrm{Sq})$ of~(b), formed with the factor $1$. By~(b),
\[
\liminf_{n\to\infty}n^{-3}\log E^{\mathrm{af}}_n\exp(\mathfrak{l}\cdot\mathrm{Sq})\ge\frac{1}{2}\mathfrak{l}^2-4\log\frac{6+\boldsymbol{\mu}}{\boldsymbol{\mu}},
\]
and the right-hand side is positive if and only if $\mathfrak{l}^2>8\log\frac{6+\boldsymbol{\mu}}{\boldsymbol{\mu}}$.
\end{proof}

\subsubsection{The zero-harmonic law and the averaged law}\label{sec:zero-harmonic-law}

\begin{remark}\label{rem:free-coefficient-1}
Let $e>0$, let $N\mapsto m_{0,N}$ be a mass assignment, and let $\mu$ be a limit point of $\left(P^{\mathrm{av}}_{N,T}\right)$ along a sequence of cutoffs $(N_k,T_k)_{k\in\mathbb{N}}$ with $N_k\to\infty$ and $T_k\to\infty$. By Section~\ref{sec:wilson-determinant-lattice}, $E^{\mathrm{av}}_n$ is the averaged law $\bar{\mathsf{P}}$ of Section~\ref{sec:wilson-operator-lattice}, formed with $q=e\cdot\eta^{\frac{1}{2}}$ and the lattice mass $\boldsymbol{\mu}_N>-\frac{1}{2}$. Hence $\mu$ is a limit point, in the sense of Definition~\ref{def:limit-points}, of the averaged laws of the parameter assignment $(N,T)\mapsto\left(e\cdot\eta^{\frac{1}{2}},\boldsymbol{\mu}_N\right)$ of Section~\ref{sec:wilson-operator-lattice}, and at every cutoff the zero-harmonic law $\mathsf{P}^0$ of Section~\ref{sec:wilson-operator-lattice}, formed with the same $q$ and $\boldsymbol{\mu}_N$, is defined (Section~\ref{sec:pinned-law-fixed}). Let $G^0_k$ and $\bar{G}_k$ be the generating functions of the zero-harmonic law and of the averaged law at the cutoff $(N_k,T_k)$, let $E^0_k$ be the expectation under this zero-harmonic law, let $X_k=X_{\eta_k,R_k}$ with $\eta_k=L^{-N_k}$ and $R_k=h\cdot L^{T_k}$, and let $D_{\rho}$ be the polydiscs, all as in Section~\ref{sec:criterion-convergence-pinned}.
\begin{enumerate}
\item[(1)] Condition~\eqref{eq:continuum-infinite-TWp} of Proposition~\ref{prop:pinned-averaged-generating} requires, for every integer $s\ge1$ and all real $f_1,\dots,f_s\in C_c^{\infty}(\mathbb{R}^3;\Lambda^2)$, that the functions $G^0_k$ be bounded on some polydisc $D_{\rho'}$ uniformly in $k$, and that the moments $E^0_k\prod_{i=1}^sX_k(f_i)$ of the zero-harmonic laws converge to the moments $\int\mathrm{d}\mu(F)\,\prod_{i=1}^sF(f_i)$ of $\mu$. By Proposition~\ref{prop:pinned-averaged-generating}(1), this condition is equivalent to condition~\eqref{eq:continuum-infinite-TW}: for every integer $s\ge1$ and all real $f_1,\dots,f_s\in C_c^{\infty}(\mathbb{R}^3;\Lambda^2)$ there is $\rho>0$ with
\[
\lim_{k\to\infty}\,\sup_{z\in D_{\rho}}\left|G^0_k(z)-\bar{G}_k(z)\right|=0.
\]
\item[(2)] Example~\ref{ex:pinned-law-reflection} concerns the zero-harmonic law on the torus $\mathbb{Z}/1\mathbb{Z}\times\mathbb{Z}/1\mathbb{Z}\times\mathbb{Z}/3\mathbb{Z}$, which has $1\times1\times3$ sites, with $\eta=1$. If this law were reflection positive in the direction $3$ in the sense of Proposition~\ref{prop:reflection-positivity}, its density $\rho^0$, written in the coordinates $(x,y)\in\mathbb{R}^2\times\mathbb{R}^2$ of that example, would satisfy the Cauchy--Schwarz inequality
\[
\rho^0(x,y)^2\le\rho^0(x,x)\cdot\rho^0(y,y)\qquad\left(x,y\in\mathbb{R}^2\right).
\]
Let $q>0$, put $x=\frac{2\pi}{q}\cdot(1,1)$ and $y=\frac{4\pi}{q}\cdot(1,1)$, and let $A(x,y)\in V_{\mathrm{t}}$ be the cochain with the coordinates $(x,y)$, as in Example~\ref{ex:pinned-law-reflection}. The numbers $q\cdot A_{t,\alpha}$ of that example, $t\in\{-1,0,1\}$, $\alpha\in\{1,2\}$, are congruent modulo $2\pi$ to $0$ for $A=A(x,y)$, to $\frac{2\pi}{3}$ for $A=A(x,x)$ and to $-\frac{2\pi}{3}$ for $A=A(y,y)$; at these points Example~\ref{ex:pinned-law-reflection} computes the determinants in closed form: for every $\boldsymbol{\mu}>-1$,
\[
\det\left(D_{A(x,y)}+\boldsymbol{\mu}\right)=\Delta_0(\boldsymbol{\mu}),\qquad\det\left(D_{A(x,x)}+\boldsymbol{\mu}\right)=\det\left(D_{A(y,y)}+\boldsymbol{\mu}\right)=\Delta_1(\boldsymbol{\mu}),
\]
where
\[
\Delta_0(\boldsymbol{\mu}):=(\boldsymbol{\mu}+6)^4\cdot\left(\left(\boldsymbol{\mu}+\frac{9}{2}\right)^2+\frac{3}{4}\right)^4,\qquad\Delta_1(\boldsymbol{\mu}):=\left((\boldsymbol{\mu}+3)^2+\frac{3}{2}\right)^2\cdot\left(\left(\boldsymbol{\mu}+\frac{3}{2}\right)^2+\frac{9}{4}\right)^4.
\]
As $|x|^2=|x-y|^2=\frac{8\pi^2}{q^2}$ and $|y|^2=\frac{32\pi^2}{q^2}$, the formula for $\rho^0$ in that example gives
\[
\frac{\rho^0(x,y)^2}{\rho^0(x,x)\cdot\rho^0(y,y)}=\exp\left(-\frac{8\pi^2}{q^2}\right)\cdot\left(\frac{\Delta_0(\boldsymbol{\mu})}{\Delta_1(\boldsymbol{\mu})}\right)^2,
\]
and in particular, to three significant digits,
\[
\frac{\rho^0(x,y)^2}{\rho^0(x,x)\cdot\rho^0(y,y)}=\begin{cases}8.26\cdot10^5&\text{for }q=6\text{ and }\boldsymbol{\mu}=\frac{1}{2},\\ 7.41\cdot10^5&\text{for }q=4\text{ and }\boldsymbol{\mu}=-0.45.\end{cases}
\]
Both values exceed $1$, contrary to the Cauchy--Schwarz inequality displayed above. Hence the zero-harmonic law on this torus is, for these two choices of $(q,\boldsymbol{\mu})$, not reflection positive in the direction $3$ in the sense of Proposition~\ref{prop:reflection-positivity}. The same holds for every $\boldsymbol{\mu}>-1$ and every $q>2\pi\cdot\left(\log\frac{\Delta_0(\boldsymbol{\mu})}{\Delta_1(\boldsymbol{\mu})}\right)^{-\frac{1}{2}}$. Indeed, for $\boldsymbol{\mu}>-1$ one has $(\boldsymbol{\mu}+6)^2-(\boldsymbol{\mu}+3)^2-\frac{3}{2}=6\boldsymbol{\mu}+\frac{51}{2}>0$ and $\left(\boldsymbol{\mu}+\frac{9}{2}\right)^2+\frac{3}{4}-\left(\boldsymbol{\mu}+\frac{3}{2}\right)^2-\frac{9}{4}=6\boldsymbol{\mu}+\frac{33}{2}>0$, so that $\Delta_0(\boldsymbol{\mu})>\Delta_1(\boldsymbol{\mu})>0$; and the quotient above exceeds $1$ if and only if $q^2\cdot\log\frac{\Delta_0(\boldsymbol{\mu})}{\Delta_1(\boldsymbol{\mu})}>4\pi^2$.
\item[(3)] Theorem~\ref{thm:insensitivity-twist}, which compares the zero-harmonic law with the averaged law, assumes a positive lattice mass; accordingly, the first assertion of Corollary~\ref{cor:zero-harmonic-law} is a statement at each single integer $N$ with $\boldsymbol{\mu}_N>0$, and the second concerns an infinite set of such integers $N$. Let $e\cdot h'^{\frac{1}{2}}\le e_1$, and let $N\mapsto m_{0,N}(e)$ be the mass assignment of Theorem~\ref{thm:free-coefficient-3}. For every integer $N\ge N_0(e)$ and $\eta=L^{-N}$,
\[
\left|m_{0,N}(e)\right|\le\frac{1}{8}\left(e\cdot\eta^{\frac{1}{2}}\right)^{\frac{3}{4}-8\varepsilon}
\]
(Step~3 of the proof of part~\ref{it:free-coefficient-3-1} of Theorem~\ref{thm:free-coefficient-3}, Section~\ref{sec:proof-theorem}), and
\[
\frac{1}{8}\left(e\cdot\eta^{\frac{1}{2}}\right)^{\frac{3}{4}-8\varepsilon}\cdot\left(\bar{m}\cdot\eta\right)^{-1}=\frac{e^{\frac{3}{4}-8\varepsilon}}{8\,\bar{m}}\cdot\eta^{-\frac{5}{8}-4\varepsilon}\longrightarrow\infty\qquad(N\to\infty).
\]
Hence, for all large $N$, the number $-\frac{1}{8}\left(e\cdot\eta^{\frac{1}{2}}\right)^{\frac{3}{4}-8\varepsilon}$, which satisfies this bound, is smaller than $-\bar{m}\cdot\eta$: the bound does not imply $\boldsymbol{\mu}_N=\bar{m}\cdot\eta+m_{0,N}(e)>0$. For $N\ge N_0(e)$, every bound on $m_{0,N}(e)$ in terms of $e$ and $\eta$ that is proved follows from a bound $\left|m_{0,N}(e)\right|\le c\cdot\left(e\cdot\eta^{\frac{1}{2}}\right)^{\frac{3}{4}-8\varepsilon}$ with a number $c>0$ that does not depend on $N$: the statements, among them Proposition~\ref{prop:pair-vanishing-terminal} and Theorem~\ref{thm:agreement-trajectories} at the level $k=0$, give $c=\frac{1}{8}$, that is, the bound above, and bounds that follow from it, such as $\left|m_{0,N}(e)\right|<\frac{1}{8}$ and $\boldsymbol{\mu}_N>-\frac{1}{8}$ (Step~3 of the proof of part~\ref{it:free-coefficient-3-1} of Theorem~\ref{thm:free-coefficient-3}); and the proof of Proposition~\ref{prop:pair-vanishing-terminal}, applied on the torus $R_{\mathrm{top}}(e)$ (Remark~\ref{rem:free-coefficient-2}), gives $c=\frac{1}{12}$. For each such $c$,
\[
c\cdot\left(e\cdot\eta^{\frac{1}{2}}\right)^{\frac{3}{4}-8\varepsilon}\cdot\left(\bar{m}\cdot\eta\right)^{-1}=\frac{c\cdot e^{\frac{3}{4}-8\varepsilon}}{\bar{m}}\cdot\eta^{-\frac{5}{8}-4\varepsilon}\longrightarrow\infty\qquad(N\to\infty),
\]
so that, as above, for all large $N$ none of these bounds implies $\boldsymbol{\mu}_N>0$. For $0\le N<N_0(e)$, on the other hand, $m_{0,N}(e)=0$ by the definition in Theorem~\ref{thm:free-coefficient-3}, so that $\boldsymbol{\mu}_N=\bar{m}\cdot\eta>0$; these are finitely many integers $N$. Thus neither the sign of $\boldsymbol{\mu}_N$ for large $N$ nor the existence of infinitely many integers $N$ with $\boldsymbol{\mu}_N>0$ is decided, and the second assertion of Corollary~\ref{cor:zero-harmonic-law} applies only if $\boldsymbol{\mu}_N>0$ for infinitely many $N$.
\end{enumerate}
\end{remark}

\subsubsection{The all-flux law and the averaged law}\label{sec:all-flux-law}

In Section~\ref{sec:all-flux-law}, $n\ge2$ is an integer, $\Lambda_n$ is the cubic torus, $q>0$, $\boldsymbol{\mu}>-\frac{1}{2}$, and
\[
x_n:=\frac{2\pi^2}{q^2\cdot n}.
\]
Since $r_q^2=4\pi^2/q^2$, $|n|=n^3$ and $n_{\iota}=n$, the definition in Section~\ref{sec:wilson-determinant-lattice} gives $\mathfrak{q}_w(n)=\exp\left(-x_n\cdot|w|_2^2\right)$ for $w\in\mathbb{Z}^3$. If $q=e\cdot\eta^{\frac{1}{2}}$ and $n=R/\eta$, then $q^2\cdot n=e^2\cdot R$, so that $x_n=2\pi^2/\left(e^2\cdot R\right)$.

\begin{remark}[Comparison through the flux sectors]\label{rem:comparison-through-flux}
\begin{enumerate}
\item[(a)] By Corollary~\ref{cor:partition-functions-1}(d), in which $E_n=E^{\mathrm{af}}_n$ and $E_{n,0}=E^{\mathrm{av}}_n$ (Section~\ref{sec:wilson-determinant-lattice}): if $\varsigma_n\le\frac{1}{4}$, then
\[
\left|E^{\mathrm{av}}_nG-E^{\mathrm{af}}_nG\right|\le3\,\varsigma_n^{\frac{1}{2}}\cdot\left(E^{\mathrm{af}}_n|G|^2\right)^{\frac{1}{2}}\qquad\left(G\in\mathfrak{F}(n)\right).
\]
\item[(b)] The six fluxes $w$ with $|w|_2=1$ contribute $6\cdot\exp\left(-x_n\right)$ to $\varsigma_n$; hence $\varsigma_n\ge6\cdot\exp\left(-x_n\right)$. Consequently $\varsigma_n\le\frac{1}{4}$ implies $x_n\ge\log24$, and the factor in (a) satisfies
\[
3\,\varsigma_n^{\frac{1}{2}}\ge3\cdot6^{\frac{1}{2}}\cdot\exp\left(-\frac{x_n}{2}\right).
\]
For $q=e\cdot\eta^{\frac{1}{2}}$ and $n=R/\eta$ this reads as follows: the hypothesis $\varsigma_n\le\frac{1}{4}$ of (a) implies $e^2\cdot R\le2\pi^2/\log24$, and $3\,\varsigma_n^{\frac{1}{2}}\ge3\cdot6^{\frac{1}{2}}\cdot\exp\left(-\frac{\pi^2}{e^2\cdot R}\right)$.
\item[(c)] By Corollary~\ref{cor:partition-functions-1}(d), $P_n(w)\le\mathfrak{q}_w(n)$ for every $w\in\mathbb{Z}^3$. For $q=e\cdot\eta^{\frac{1}{2}}$, $n=R/\eta$ and every $c>0$, this bound is at least $\exp\left(-2\pi^2\cdot c\right)$ for all $w$ with $|w|_2^2\le c\cdot e^2\cdot R$.
\end{enumerate}
\end{remark}

\begin{proposition}[The constant part of the curvature]\label{prop:constant-part-curvature}
Let $e>0$, let $(N,T)$ be a cutoff, and let $\eta=L^{-N}$, $R=h\cdot L^T$, $n=R/\eta$, $q=e\cdot\eta^{\frac{1}{2}}$ and $\boldsymbol{\mu}>-\frac{1}{2}$. Let $P_h$ be the orthogonal projection of $C^2(n)$, with the counting pairing, onto the subspace of the constant two-cochains, and for $w\in\mathbb{Z}^3$ let $h_w\in C^2(n)$ be the constant two-cochain with the components
\[
(h_w)_{\mu\nu}(x)=\frac{r_q\cdot w_{\mu\nu}}{n^2}\qquad\left(\mu<\nu,\ x\in\Lambda_n\right);
\]
this is the cochain $h_w$ of Section~\ref{sec:decomposition-flux} on the cubic torus $n$.
\begin{enumerate}
\item[(a)] Every exact two-cochain on $\Lambda_n$ is orthogonal to every constant two-cochain. If $F^u=dA+r_q\cdot\mathfrak{z}$ with $A\in C^1(n)$ and a closed $\mathfrak{z}\in C^2(n;\mathbb{Z})$ of flux $w$, then $P_hF^u=h_w$. In physical units, $\eta^{-\frac{3}{2}}\cdot h_w$ is the constant two-cochain with the components $2\pi\cdot w_{\mu\nu}/\left(e\cdot R^2\right)$.
\item[(b)] Let $f\in\mathcal{S}$ be real, and let $g:=\left(f_{\eta}\right)^u=\eta^{\frac{3}{2}}\cdot f_{\eta}$ (Section~\ref{sec:lattices-tests-function}). Then
\[
X_{\eta,R}(f)=\left\langle g,F^u-P_hF^u\right\rangle+\left\langle P_hg,F^u\right\rangle,
\]
where $\left\langle P_hg,F^u\right\rangle=\left\langle g,P_hF^u\right\rangle$ is the contribution of the constant part $P_hF^u$ of the curvature, and
\[
\left\|P_hg\right\|^2=n^{-3}\sum_{\mu<\nu}\left(\sum_{x\in\Lambda_n}g_{\mu\nu}(x)\right)^2=R^{-3}\sum_{\mu<\nu}\left(\int_{\mathbb{R}^3}\mathrm{d}y\,f_{\mu\nu}(y)\right)^2.
\]
\item[(c)] One has $E^{\mathrm{af}}_n\exp\left\langle P_hg,F^u\right\rangle\le\exp\left(\frac{1}{2}\left\|P_hg\right\|^2\right)$, and
\[
E^{\mathrm{af}}_n\left\langle g,P_hF^u\right\rangle^2\le R^{-3}\sum_{\mu<\nu}\left(\int_{\mathbb{R}^3}\mathrm{d}y\,f_{\mu\nu}(y)\right)^2.
\]
The right side does not depend on $N$, and it tends to $0$ as $T\to\infty$.
\end{enumerate}
\end{proposition}

\begin{proof}
(a) Let $a\in C^1(n)$ and $\mu<\nu$. By the formula for the coboundary in Section~\ref{sec:cochains-periodic-lattice},
\[
\sum_{x\in\Lambda_n}(da)_{\mu\nu}(x)=\sum_{x\in\Lambda_n}\left(a_{\mu}(x)-a_{\mu}(x+e_{\nu})\right)+\sum_{x\in\Lambda_n}\left(a_{\nu}(x+e_{\mu})-a_{\nu}(x)\right)=0,
\]
since $x\mapsto x+e_{\mu}$ and $x\mapsto x+e_{\nu}$ are bijections of $\Lambda_n$. For $\mu<\nu$ let $1_{\mu\nu}\in C^2(n)$ take the value $1$ on the plaquettes spanned by $e_{\mu}$ and $e_{\nu}$ and the value $0$ on the other plaquettes. The three cochains $1_{\mu\nu}$ span the space of the constant two-cochains, they are pairwise orthogonal, and $\|1_{\mu\nu}\|^2=n^3$. The display says $\langle da,1_{\mu\nu}\rangle=0$, which proves the first assertion. Moreover, for every $u\in C^2(n)$ the cochain $P_hu$ is the constant two-cochain with the components $n^{-3}\sum_xu_{\mu\nu}(x)$, and
\[
\left\|P_hu\right\|^2=n^3\sum_{\mu<\nu}\left(n^{-3}\sum_{x\in\Lambda_n}u_{\mu\nu}(x)\right)^2=n^{-3}\sum_{\mu<\nu}\left(\sum_{x\in\Lambda_n}u_{\mu\nu}(x)\right)^2.
\]
Let $F^u=dA+r_q\cdot\mathfrak{z}$ be as in (a). Then $P_hdA=0$. For $\mu<\nu$ let $\iota$ be the third direction. By Section~\ref{sec:wilson-determinant-lattice}, the sum of $\mathfrak{z}_{\mu\nu}$ over each of the $n$ planes $\{x_{\iota}=j\}$, $j\in\mathbb{Z}/n\mathbb{Z}$, is $w_{\mu\nu}$; so $\sum_{x\in\Lambda_n}\mathfrak{z}_{\mu\nu}(x)=n\cdot w_{\mu\nu}$. Hence $P_hF^u=r_q\cdot P_h\mathfrak{z}$ has the components $r_q\cdot n^{-3}\cdot n\cdot w_{\mu\nu}=r_q\cdot w_{\mu\nu}/n^2$; that is, $P_hF^u=h_w$. Finally, $r_q=2\pi/q=2\pi/\left(e\cdot\eta^{\frac{1}{2}}\right)$ and $n=R/\eta$ give
\[
\eta^{-\frac{3}{2}}\cdot\frac{r_q\cdot w_{\mu\nu}}{n^2}=\eta^{-\frac{3}{2}}\cdot\frac{2\pi\cdot w_{\mu\nu}}{e\cdot\eta^{\frac{1}{2}}}\cdot\frac{\eta^2}{R^2}=\frac{2\pi\cdot w_{\mu\nu}}{e\cdot R^2}.
\]

(b) By Section~\ref{sec:lattices-tests-function}, $X_{\eta,R}(f)=\langle f_{\eta},F\rangle_{\eta}=\langle g,F^u\rangle$, which gives the decomposition; and $\langle g,P_hF^u\rangle=\langle P_hg,F^u\rangle$, as the orthogonal projection $P_h$ is symmetric. Let $\mu<\nu$, and let $\tilde{f}$ be the $R$-periodisation of $f$. For $x\in\Lambda_{\eta,R}$ the value of $f_{\eta}$ at the plaquette at $x$ spanned by $e_{\mu}$ and $e_{\nu}$ is the mean of $\tilde{f}_{\mu\nu}$ over the averaging cube $c+\left[-\frac{\eta}{2},\frac{\eta}{2}\right)^3$ of this plaquette, $c$ being its centre (Section~\ref{sec:lattices-tests-function} and \eqref{eq:continuum-infinite-2}). The centres of these plaquettes form a coset of $\eta\mathbb{Z}^3$ modulo $R\mathbb{Z}^3$; so their averaging cubes tile the torus $\mathbb{R}^3/R\mathbb{Z}^3$, and
\[
\sum_x\left(f_{\eta}\right)_{\mu\nu}(x)=\eta^{-3}\int_{\mathbb{R}^3/R\mathbb{Z}^3}\mathrm{d}y\,\tilde{f}_{\mu\nu}(y)=\eta^{-3}\int_{\mathbb{R}^3}\mathrm{d}y\,f_{\mu\nu}(y),
\]
the last equality by Fubini's theorem, $f$ being integrable. Thus $\sum_xg_{\mu\nu}(x)=\eta^{-\frac{3}{2}}\int_{\mathbb{R}^3}\mathrm{d}y\,f_{\mu\nu}(y)$. The formula for $\|P_hu\|^2$ in the proof of (a), applied to $u=g$, and $n\cdot\eta=R$ give the formula for $\|P_hg\|^2$.

(c) Let $s\in\mathbb{R}$. The function $F^u\mapsto\exp\left\langle s\cdot P_hg,F^u\right\rangle$ belongs to $\mathfrak{F}(n)$, and~\eqref{it:free-coefficient-Q3}, applied to the real cochain $s\cdot P_hg$, gives
\[
E^{\mathrm{af}}_n\exp\left(s\cdot\left\langle P_hg,F^u\right\rangle\right)\le\exp\left(\frac{1}{2}s^2\cdot\left\|P_hg\right\|^2\right);
\]
for $s=1$ this is the first assertion. As in Step~3 of the proof of Theorem~\ref{thm:main-theorem} (Section~\ref{sec:proof-part-ii-2}), the invariance of $E^{\mathrm{af}}_n$ under $F^u\mapsto-F^u$ (Theorem~\ref{thm:free-coefficient-1}(d5)) gives
\[
E^{\mathrm{af}}_n\cosh\left(s\cdot\left\langle P_hg,F^u\right\rangle\right)=E^{\mathrm{af}}_n\exp\left(s\cdot\left\langle P_hg,F^u\right\rangle\right),
\]
and the inequality $\frac{1}{2}t^2\le\cosh t-1$, valid for real $t$, yields for $s\ne0$
\[
E^{\mathrm{af}}_n\left\langle P_hg,F^u\right\rangle^2\le\frac{2}{s^2}\left(\exp\left(\frac{1}{2}s^2\cdot\left\|P_hg\right\|^2\right)-1\right).
\]
As $s\to0$, the right side tends to $\|P_hg\|^2$. By (b), the left side is $E^{\mathrm{af}}_n\langle g,P_hF^u\rangle^2$, and $\|P_hg\|^2$ is the right side of the asserted inequality. This number contains neither $N$ nor $\eta$, and it tends to $0$ as $T\to\infty$, since $R=h\cdot L^T$.
\end{proof}

\begin{remark}[An extensive observable]\label{rem:extensive-observable}
Let $\mathrm{Sq}:=\sum_{x\in\Lambda_n}F^u_{12}(x)$ be the sum of the values of $F^u$ on the plaquettes spanned by $e_1$ and $e_2$. By Section~\ref{sec:pressure}, $E^{\mathrm{av}}_n\exp\left(\mathfrak{l}\cdot\mathrm{Sq}\right)=1$ for every real $\mathfrak{l}$; and, for fixed $q>0$ and $\boldsymbol{\mu}>0$ and for real $\mathfrak{l}$ with $\mathfrak{l}^2>8\log\left(\frac{6+\boldsymbol{\mu}}{\boldsymbol{\mu}}\right)$,
\[
\liminf_{n\to\infty}\,n^{-3}\log E^{\mathrm{af}}_n\exp\left(\mathfrak{l}\cdot\mathrm{Sq}\right)>0.
\]
Hence, for all sufficiently large $n$, the laws of $\mathrm{Sq}$ under $E^{\mathrm{af}}_n$ and under $E^{\mathrm{av}}_n$ are different.
\end{remark}

\subsubsection{The ergodic components of a limit point}\label{sec:ergodic-components-limit}

\begin{remark}[The fourth cumulant and the covariance of the components]\label{rem:fourth-cumulant-covariance}
Let $0<e\le e_2$, let $\mu$ be a limit point as in Theorem~\ref{thm:main-theorem}, and let $\mu=\int\mathrm{d}\nu(\xi)\,\mu_{\xi}$ be its ergodic decomposition (Section~\ref{sec:laws-2}). By Theorem~\ref{thm:main-theorem}(vi), for $\nu$-almost every $\xi$ the numbers $C_{\xi}(f,f)$ and $\kappa_4^{\xi}(f)$ lie in the intervals
\[
\left[\|f\|_2^2-C_{\mathrm{kap}}\cdot e^{2-\mathfrak{e}_2},\ \|f\|_2^2+C_{\mathrm{kap}}\cdot e^{2-\mathfrak{e}_2}\right]\quad\text{and}\quad\left[e^4\cdot b(f)-C_{\mathrm{kap}}\cdot e^{4+\sigma},\ e^4\cdot b(f)+C_{\mathrm{kap}}\cdot e^{4+\sigma}\right],
\]
whose lengths are $2C_{\mathrm{kap}}\cdot e^{2-\mathfrak{e}_2}$ and $2C_{\mathrm{kap}}\cdot e^{4+\sigma}$.
\begin{enumerate}
\item[(a)] The limit point $\mu$ lies in the class $\mathbf{K}(C_{\mathrm{kap}})$ (Section~\ref{sec:soft-laws-four}), and $\kappa_4^{\mu}(f)\ge\frac{1}{2}e^4\cdot b(f)>0$ by Theorem~\ref{thm:main-theorem}(iv). Suppose that there is a number $k$ with $\kappa_4^{\xi}(f)=k$ for $\nu$-almost every $\xi$; then $k\ge\frac{1}{2}e^4\cdot b(f)>0$ by Theorem~\ref{thm:main-theorem}(vi). Let $\boldsymbol{\delta}>0$, let $0<\epsilon\le\epsilon_0$, with the number $\epsilon_0$ of Proposition~\ref{prop:operations-class} for $C=C_{\mathrm{kap}}$, and let
\[
\mu':=\frac{1}{2}\left(\mu+\mu\ast_{\epsilon}\gamma_{C_0}\right)
\]
be the law of Proposition~\ref{prop:operations-class}(b), which lies in $\mathbf{K}(C_{\mathrm{kap}}+\boldsymbol{\delta})$. By the proof of Proposition~\ref{prop:operations-class}, the measure of the ergodic decomposition of $\mu'$ is the mean of the measures of the ergodic decompositions of $\mu$ and of $\mu\ast_{\epsilon}\gamma_{C_0}$; the ergodic components of $\mu\ast_{\epsilon}\gamma_{C_0}$ are the laws $\mu_{\xi}\ast_{\epsilon}\gamma_{C_0}$; and the fourth cumulant of $\mu_{\xi}\ast_{\epsilon}\gamma_{C_0}$ at $f$ is $(1-\epsilon)^2\cdot\kappa_4^{\xi}(f)+\epsilon^2\cdot\kappa_4^{\gamma_{C_0}}(f)=(1-\epsilon)^2\cdot\kappa_4^{\xi}(f)$, the Gaussian law $\gamma_{C_0}$ having vanishing fourth cumulant. As $0<\epsilon<1$ and $k\ne0$, the numbers $k$ and $(1-\epsilon)^2\cdot k$ are different. Hence, for the ergodic decomposition of $\mu'$, the fourth cumulant at $f$ of the components takes the value $k$ on a set of measure $\frac{1}{2}$ and the value $(1-\epsilon)^2\cdot k$ on a set of measure $\frac{1}{2}$.
\item[(b)] Let $\lambda$ be a law on $\mathcal{S}'$ which satisfies~\eqref{eq:fourth-cumulants-H1}--\eqref{eq:fourth-cumulants-H3}, is ergodic in the sense of Definition~\ref{def:ergodicity-ergodic-decompositions}, and is not Gaussian (Section~\ref{sec:cumulants-gaussian-laws}); let $C':=S_2^{\lambda}$, let $\gamma_{C'}$ be the centred Gaussian law with covariance $C'$ of Lemma~\ref{lem:gaussian-laws-bounded}(a), and for $0\le p<1$ put $\mu^{(p)}:=(1-p)\cdot\lambda+p\cdot\gamma_{C'}$. By Proposition~\ref{prop:gaussian-components-four}(ii), $\mu^{(p)}$ satisfies~\eqref{eq:fourth-cumulants-H1}--\eqref{eq:fourth-cumulants-H3}, its ergodic decomposition has the components $\lambda$ and $\gamma_{C'}$ with the weights $1-p$ and $p$, and
\[
S_2^{\mu^{(p)}}=C'=S_2^{\lambda},\qquad\kappa_4^{\mu^{(p)}}=(1-p)\cdot\kappa_4^{\lambda}.
\]
Thus, as $p\to0$, the two-point function of $\mu^{(p)}$ is that of the ergodic, non-Gaussian law $\lambda$ and its fourth cumulant tends to that of $\lambda$, whereas for every $p>0$ the Gaussian components of $\mu^{(p)}$ form a set of measure $p$. Assertion~(vi) of Theorem~\ref{thm:main-theorem} is derived in Step~6 of its proof (Section~\ref{sec:proof-part-ii-2}) from the inequalities~\eqref{eq:free-coefficient-10}, through~\eqref{eq:free-coefficient-22} and~\eqref{eq:free-coefficient-23}; by Remark~\ref{rem:moments-all-orders}, the inequalities~\eqref{eq:free-coefficient-23} involve the moments of $\mu$ of all orders.
\item[(c)] If $\mu$ satisfies the clustering property~\eqref{eq:fourth-cumulants-C1} of Theorem~\ref{thm:clustering-mixing-ergodicity}, then $\mu_{\xi}=\mu$ for $\nu$-almost every $\xi$, and the image of $\nu$ under $\xi\mapsto\mu_{\xi}$ is the point mass at $\mu$ (Corollary~\ref{cor:decomposition-clustering-law}(ii); by Theorem~\ref{thm:main-theorem}(c), $\mu$ satisfies~\eqref{eq:setup-H}, which is the conjunction of~\eqref{eq:fourth-cumulants-H1}--\eqref{eq:fourth-cumulants-H3}). In this case $C_{\xi}(f,f)=S_2^{\mu}(f,f)$ and $\kappa_4^{\xi}(f)=\kappa_4^{\mu}(f)$ for $\nu$-almost every $\xi$.
\end{enumerate}
\end{remark}

\begin{example}[A prescribed fourth cumulant with Gaussian components]\label{ex:prescribed-fourth-cumulant}
Let $f_{*}\in\mathcal{S}$ satisfy $C_0(f_{*},f_{*})>0$, let $b_{*}>0$, put
\[
q_{*}:=\frac{1}{C_0(f_{*},f_{*})}\cdot\left(\frac{b_{*}}{3}\right)^{\frac{1}{2}},
\]
and let $0<e\le(2q_{*})^{-\frac{1}{2}}$, so that $0\le1-2q_{*}\cdot e^2<1$. For $s\in[0,1]$ the form $s\cdot C_0$ is symmetric and positive semidefinite, and $s\cdot C_0(f,f)\le\|f\|_2^2$ for $f\in\mathcal{S}$ (Lemma~\ref{lem:free-transverse-covariance}(a)); let $\gamma_{sC_0}$ be the centred Gaussian law with covariance $s\cdot C_0$ (Section~\ref{sec:soft-laws-four}). Let $\mu^{\rho_e}$ be the law of Theorem~\ref{thm:properties-scale-mixtures} for the probability measure $\rho_e=\frac{1}{2}\left(\delta_1+\delta_{1-2q_{*}\cdot e^2}\right)$ on $[0,1]$, that is,
\[
\mu^{\rho_e}=\int\mathrm{d}\rho_e(s)\,\gamma_{sC_0}=\frac{1}{2}\left(\gamma_{C_0}+\gamma_{\left(1-2q_{*}\cdot e^2\right)C_0}\right).
\]
Then the following hold.
\begin{enumerate}
\item[(i)] $\kappa_4^{\mu^{\rho_e}}(f_{*})=e^4\cdot b_{*}$ (Theorem~\ref{thm:properties-scale-mixtures}(iv)).
\item[(ii)] The law $\mu^{\rho_e}$ is soft in the sense of Section~\ref{sec:soft-laws-four}. Indeed, it satisfies~\eqref{eq:fourth-cumulants-H1}--\eqref{eq:fourth-cumulants-H3}, that is~\eqref{eq:setup-H}, and $X(d^{*}h)=0$ almost surely for every Schwartz three-form $h$ (Theorem~\ref{thm:properties-scale-mixtures}(i)), so that it is concentrated on the closed distributions (Section~\ref{sec:closed-distributions}); and it is invariant under $F\mapsto-F$, as the two centred Gaussian laws $\gamma_{C_0}$ and $\gamma_{\left(1-2q_{*}\cdot e^2\right)C_0}$ are.
\item[(iii)] Its ergodic components are the centred Gaussian laws $\gamma_{C_0}$ and $\gamma_{\left(1-2q_{*}\cdot e^2\right)C_0}$, whose covariances $C_0$ and $\left(1-2q_{*}\cdot e^2\right)\cdot C_0$ are multiples of the free Maxwell covariance (Theorem~\ref{thm:properties-scale-mixtures}(ii)).
\end{enumerate}
Thus the value $e^4\cdot b_{*}>0$ of the fourth cumulant at the single test $f_{*}$ is attained by a soft law all of whose ergodic components are Gaussian. In the proof of Theorem~\ref{thm:main-theorem} (Section~\ref{sec:proof-part-ii-2}), assertion~(v) is derived in Step~5, by means of Theorem~\ref{thm:non-gaussian-components-2}, from the bounds~(iii) on $\kappa_4^{\mu}$ at the three tests $g_0$, $g_{+}$ and $g_{-}$; assertion~(vi) concerns the single test $f$, and the inequality~\eqref{eq:free-coefficient-10} from which it is derived in Step~6 (Remark~\ref{rem:fourth-cumulant-covariance}(b)) bounds, on the torus $k$ of Section~\ref{sec:reference-tori-cells} for $k=1$ and for every integer $k\ge L$ (Theorem~\ref{thm:free-coefficient-2}), the expectation of a product over an arbitrary set of cells; and assertion~(vii) is derived from both, the argument of Step~6 being applied at each of the three tests.
\end{example}

\subsubsection{The bare mass shifts}\label{sec:bare-mass-shifts-2}

\begin{remark}\label{rem:free-coefficient-2}
Let $e>0$ satisfy $e\cdot h'^{\frac{1}{2}}\le e_1$, and let $N\ge N_0(e)$ be an integer. The bare mass shift $m_{0,N}(e)$ of Theorem~\ref{thm:free-coefficient-3} is the first entry of the pair $P_N(e)$ of Definition~\ref{def:tuned-pair}, which is defined, as $N\ge N_{\mathrm{min}}\left(R^{*}(e),e\right)\ge N_{\mathrm{min}}\left(R_{\mathrm{top}}(e),e\right)$ by Proposition~\ref{prop:depth-stop-level}(d). By Definition~\ref{def:tuned-pair} and Proposition~\ref{prop:pair-vanishing-terminal}, applied on the torus $R_{\mathrm{top}}(e)$, the pair $P_N(e)$ is the unique real pair that is controlled on $R_{\mathrm{top}}(e)$ (Definition~\ref{def:fluctuation-field-controlled}) and whose trajectory on $R_{\mathrm{top}}(e)$ (Section~\ref{sec:flow}) satisfies
\[
m_{k_{\mathrm{top}}}=0\qquad\text{and}\qquad\boldsymbol{\varepsilon}^t_{k_{\mathrm{top}}}=0,\qquad\text{where }k_{\mathrm{top}}=k^{*}\left(R_{\mathrm{top}}(e)\right).
\]
The trajectories of real pairs on $R_{\mathrm{top}}(e)$ are given by the recursion of Section~\ref{sec:flow}, taken under hypothesis~\ref{def:hypothesis-flow} with $R=R_{\mathrm{top}}(e)$ (Definition~\ref{def:tuned-pair}); besides the initial pair, this recursion is formed from $(N,e,\bar{m})$ and from three further data: the systems of primitives (Section~\ref{sec:primitive-operators}) of Lemma~\ref{lem:assembled-systems-primitives} on the tori $T_k\left(R_{\mathrm{top}}(e)\right)$, $k<k_{\mathrm{top}}$; the constants fixed in the stages~\ref{it:order-choice-1}--\ref{it:order-choice-10} of the order of choice of Section~\ref{sec:order-choice}, that is, the parameters of Section~\ref{sec:parameters-scales}; and the orthonormal basis $\{e_i^0\}$ of Section~\ref{sec:averaging-fluctuations-gauge}, which is fixed once for each $L$ (Proposition~\ref{prop:basis-map-box}), is the same at every level, for every side, every $R$ and every $N$, and enters the recursion through the family $\{\Upsilon\}$ of~\eqref{eq:renormalisation-group-1}, since the set $\mathrm{Box}_k$ and the measure $\hat{\mu}$ are defined through the coordinates $\zeta(\Upsilon)$. No further datum enters through the second slot: in the families with slot that the recursion forms, this slot ranges over the whole set $\mathfrak{X}_k$ of Definition~\ref{def:domains}(f) on the torus $T_k\left(R_{\mathrm{top}}(e)\right)$, that is, over all complex bond functions $\mathscr{C}$ on this torus with $|\mathscr{C}|_{\infty}<e_k^{-1+2\varepsilon}$ and $|\partial\mathscr{C}|_{\infty}<e_k^{-1+3\varepsilon}$, so that no particular external connection (Definition~\ref{def:twist-external-connection}) is inserted; for each level $k$ this set is determined by $(N,e,\bar{m})$ and the parameters, as its two thresholds are formed with $e_k=e\cdot L^{\frac{k-N}{2}}$ and $\varepsilon$ alone and the torus has the side $R_{\mathrm{top}}(e)\cdot L^{N-k}$ and the spacing $\eta_k=L^{-k}$ (Section~\ref{sec:running-quantities-tori} and Definition~\ref{def:tori-stop-level}); and the pair $P_N(e)$ is determined by the sequences $(m_k,E_k)$ without slot (proof of Proposition~\ref{prop:pair-vanishing-terminal}). Accordingly, $P_N(e)$ is a function of $(N,e,\bar{m})$ and of these three data, that is, of the scheme in the sense of Definition~\ref{def:terminology}(10) (Theorem~\ref{thm:agreement-trajectories}). In particular, of the quantities fixed in stage~\ref{it:order-choice-11} only $\bar{m}$ enters $P_N(e)$, together with the scales $h$ and $h'$ that it determines, and of those fixed in stage~\ref{it:order-choice-12} only $e$ and $N$, the side of the torus being $R_{\mathrm{top}}(e)=h'\cdot M\cdot r_K(e)$, a function of $e$, $\bar{m}$ and the parameters; thus neither the probe $\beta$ nor the point $x_0$ of Section~\ref{sec:probes} enters $P_N(e)$, nor does any of the further numbers $C_{\beta}$, $\mathbf{D}$, $D_{\beta}$, $\rho_{\beta}$, $C_1$, $C_2$, $e_{\sharp}$, $C^G_{\beta}$, $C_{\mathrm{Gam}}$, $C_Y$, $C_{\mathrm{kap}}$ and $e_2$ of stage~\ref{it:order-choice-11}, and hence none of them enters $m_{0,N}(e)$. The masses of the trajectory satisfy
\[
\left|m_k\right|\le\frac{1}{8}\,w^{\mathrm{m}}_k\qquad\left(0\le k\le k_{\mathrm{top}}\right)
\]
(Proposition~\ref{prop:pair-vanishing-terminal}); for $k=0$ this is the bound for $m_{0,N}(e)$ in Remark~\ref{rem:free-coefficient-1}(3).
\end{remark}

\subsubsection{The range of the charge and of the sides of the tori}\label{sec:range-charge-sides}

\begin{remark}\label{rem:free-coefficient-3}
\begin{enumerate}
\item[(a)] The numbers $e_2$ and $C_{\mathrm{kap}}$ of Theorem~\ref{thm:main-theorem} depend on the constants of Section~\ref{sec:constants-exponents-scales} and on $(\bar{m},\beta,x_0)$. By~\eqref{eq:free-coefficient-1}, $e_2\le e_G(\bar{m})$, and by the definition of $e_G$ in Section~\ref{sec:constants-exponents-scales}, every $e\in\left(0,e_G\right]$ satisfies the condition~\eqref{eq:renormalisation-group-18}, that is, $e\cdot h'^{\frac{1}{2}}\cdot\log h'\le1$. Let $0<\bar{m}<L^{-m_0-1}$. As $\mathfrak{m}_K=\bar{m}\cdot h'>L^{-m_0-1}$ (Section~\ref{sec:constants-exponents-scales}),
\[
h'>L^{-m_0-1}\cdot\bar{m}^{-1}>1,\qquad\log h'>\log\left(\bar{m}^{-1}\right)-(m_0+1)\cdot\log L>0.
\]
Squaring~\eqref{eq:renormalisation-group-18} gives $e^2\cdot h'\cdot\log^2h'\le1$ for every $e\in\left(0,e_2\right]$, and therefore
\[
\frac{e^2}{\bar{m}}=\frac{e^2\cdot h'}{\mathfrak{m}_K}<L^{m_0+1}\cdot e^2\cdot h'\le\frac{L^{m_0+1}}{\log^2h'}<\frac{L^{m_0+1}}{\left(\log\left(\bar{m}^{-1}\right)-(m_0+1)\cdot\log L\right)^2};
\]
the right side tends to $0$ as $\bar{m}\to0$.
\item[(b)] Let $T^{*}\ge1$, and let $k\ge L$ be an integer that is not a power of $L$. The torus $k$ of Section~\ref{sec:reference-tori-cells} has the side $k\cdot R^{*}=k\cdot h\cdot L^{T^{*}}$, which is not of the form $h\cdot L^T$ with an integer $T$; Theorem~\ref{thm:free-coefficient-2} holds on it, as on every torus $k$ with $k\in\{1\}\cup\{L,L+1,\dots\}$.
\end{enumerate}
\end{remark}

\subsubsection{The free coefficient on a torus of fixed side}\label{sec:free-coefficient-torus}

\begin{lemma}[The stop level on a torus of fixed side]\label{lem:stop-level-torus}
Let $T\ge0$ be an integer and $R:=h'\cdot L^T$. Let $e>0$ satisfy $e_K=e\cdot h'^{\frac{1}{2}}\le e_1$ and $R\le R_{\mathrm{top}}(e)$, and let $N\ge N_{\mathrm{min}}(R,e)$. Then the stop level $k^{*}=k^{*}(R)$ of Definition~\ref{def:tori-stop-level} exists, and
\[
a_{*}=\frac{R}{l_{*}}\le\frac{R}{M\cdot v_K^r},\qquad\mathfrak{m}_{*}=\bar{m}\cdot a_{*}\le\frac{\bar{m}\cdot R}{M\cdot v_K^r}.
\]
The two conditions on $e$ hold for every sufficiently small $e>0$; as $e\to0$, the right sides of the two bounds tend to $0$, and they do not depend on $N$. In particular, there is $e'>0$ such that $\mathfrak{m}_{*}\le\mathfrak{m}_{-}:=L^{-m_0-3}$ for every $e\in(0,e']$ and every $N\ge N_{\mathrm{min}}(R,e)$.
\end{lemma}

\begin{proof}
By Proposition~\ref{prop:consequences}, the inequality $e_K\le e_1$ gives $v_K\ge8$ and $\left(1+\frac{\mathrm{lg}}{v_K}\right)^r<L$, which are the hypotheses of Proposition~\ref{prop:depth-stop-level}; and $R\in\mathrm{Tor}(e)$, as $R=h'\cdot L^T\le R_{\mathrm{top}}(e)$. By Proposition~\ref{prop:depth-stop-level}(c) the stop level $k^{*}(R)$ exists, and by part~(e) of that proposition
\[
M\cdot v_{*}^r\le l_{*},\qquad v_{*}\ge v_K,\qquad\mathfrak{m}_{*}\cdot l_{*}=\bar{m}\cdot R.
\]
Since $l_k=R/a_k$ (Section~\ref{sec:running-quantities-tori}) and $v_{*}\ge v_K>0$, it follows that $R/a_{*}=l_{*}\ge M\cdot v_K^r$, which is the first bound, and that $\mathfrak{m}_{*}=\bar{m}\cdot R/l_{*}=\bar{m}\cdot a_{*}$, which gives the second.

The numbers $R$, $M$, $r$, $h'$ and $\bar{m}$ do not depend on $e$, and $v_K=-\log\left(e\cdot h'^{\frac{1}{2}}\right)$ tends to $\infty$ as $e\to0$; hence so does $v_K^r$. Consider the three conditions
\[
e_K\le e_1,\qquad h'\cdot M\cdot v_K^r\ge R,\qquad\bar{m}\cdot R\le L^{-m_0-3}\cdot M\cdot v_K^r.
\]
The first holds for every $e\in\left(0,e_1\cdot h'^{-\frac{1}{2}}\right]$. The second and the third say that $v_K^r$ is at least $R/\left(h'\cdot M\right)$, respectively at least $L^{m_0+3}\cdot\bar{m}\cdot R/M$; since $v_K^r\to\infty$ as $e\to0$, each of them holds for every sufficiently small $e>0$. Consequently there is $e'>0$ such that all three conditions hold for every $e\in(0,e']$. Since $r_K(e)\ge v_K^r$, the second condition gives $R_{\mathrm{top}}(e)=h'\cdot M\cdot r_K(e)\ge R$. Hence, for $e\in(0,e']$ and $N\ge N_{\mathrm{min}}(R,e)$, the bounds proved above apply, their right sides tend to $0$ as $e\to0$, and $\mathfrak{m}_{*}\le\bar{m}\cdot R/\left(M\cdot v_K^r\right)\le L^{-m_0-3}$.
\end{proof}

\begin{remark}[Tori of fixed side and the twist]\label{rem:tori-fixed-side}
\begin{enumerate}
\item[(a)] Hypothesis~\ref{def:hypothesis-reference-torus}, which is in force throughout Section~\ref{sec:four-jets-one}, fixes the side $R=R^{*}(e)$ (Section~\ref{sec:standing-hypothesis-scales}). On this torus Lemma~\ref{lem:scales-stop-level}(a) gives $\mathfrak{m}_{-}<\mathfrak{m}_{*}$, and several constants of Section~\ref{sec:constants-3}, for instance $g=\left(b_{-}^{-1}+\mathfrak{m}_{-}^{-1}\right)^{\frac{1}{2}}$, are formed with $\mathfrak{m}_{-}=L^{-m_0-3}$. On a torus of fixed side $R=h'\cdot L^T$, Lemma~\ref{lem:stop-level-torus} gives $\mathfrak{m}_{*}\le\mathfrak{m}_{-}$ for all sufficiently small $e$.
\item[(b)] Let $\eta\in(0,1]$ and $R>0$ be such that $n=R/\eta$ is an integer $\ge2$, and let $q>0$ and $\boldsymbol{\mu}>-\frac{1}{2}$. By Section~\ref{sec:wilson-determinant-lattice}, the law $E^{\mathrm{av}}_n$, formed with $q$ and the lattice mass $\boldsymbol{\mu}$, is an average over the twist $\theta\in[-\pi,\pi]^3$: it is the law of $F^u=dA$ when $A$ is distributed according to the probability measure on $V_{\mathrm{t}}$ whose density with respect to the Lebesgue measure of $V_{\mathrm{t}}$ is proportional to
\[
A\mapsto\int_{[-\pi,\pi]^3}\mathrm{d}\theta\,\exp\left(-\frac{1}{2}\|dA\|^2\right)\cdot\det\left(D_{A+H_{\theta}}+\boldsymbol{\mu}\right),
\]
where $D_B=D[qB]$ for one-cochains $B$, and $H_{\theta}$ is the constant one-cochain with the components $\theta_{\mu}/(q\cdot n)$ (Section~\ref{sec:wilson-determinant-lattice}).
For $\theta\in\mathbb{R}^3$ let $\theta/n\in C^1(n)$ be the constant one-cochain with the components $\theta_{\mu}/n$; thus $\theta/n=q\cdot H_{\theta}$ and $D_{H_{\theta}}=D[\theta/n]$ for every $q>0$. For the free determinants, formed with $A=0$ and the lattice mass $\bar{m}\cdot\eta$, Corollary~\ref{cor:twist-ratios-free}, applied with the mass $\bar{m}\cdot\eta\in(0,1]$, gives
\[
\left|\log\det\left(D[\theta/n]+\bar{m}\cdot\eta\right)-\log\det\left(D[0]+\bar{m}\cdot\eta\right)\right|\le2\,\varepsilon_5(\bar{m}\cdot R)\qquad\left(\theta\in\mathbb{R}^3\right),
\]
and
\[
\varepsilon_5(x)\le312\cdot\exp\left(-\frac{x}{4}\right)\cdot\left(1-\exp\left(-\frac{x}{4}\right)\right)^{-1}\qquad(x>0),
\]
because $\log\left((1-y)^{-1}\right)\le y\cdot(1-y)^{-1}$ for $0\le y<1$. These bounds depend on $\bar{m}\cdot R$ only; they do not depend on $e$.
\end{enumerate}
\end{remark}

\begin{definition}[The twist-averaged free coefficient]\label{def:twist-averaged-free}
Let $\eta\in(0,1]$ and $R>0$ be such that $n=R/\eta$ is an integer $\ge2$, let $a_{\eta}$ be a real one-cochain on $\Lambda_{\eta,R}$, read in lattice units as a one-cochain on $\Lambda_n$ (Section~\ref{sec:free-fourth-order}), and let $\theta/n$ be as in Remark~\ref{rem:tori-fixed-side}(b). For every real one-cochain $B$ one has $\operatorname{Re}\left(D[B]+\bar{m}\cdot\eta\right)\ge\bar{m}\cdot\eta>0$ and $\det\left(D[B]+\bar{m}\cdot\eta\right)\ge0$ (Theorem~\ref{thm:free-coefficient-1}(d2)); hence $D[B]+\bar{m}\cdot\eta$ is injective, and $\det\left(D[B]+\bar{m}\cdot\eta\right)>0$. By Theorem~\ref{thm:free-coefficient-1}(d1), the function $(\theta,t)\mapsto\det\left(D[\theta/n+t\eta a_{\eta}]+\bar{m}\cdot\eta\right)$ is the restriction to $\mathbb{R}^3\times\mathbb{C}$ of an entire function on $\mathbb{C}^3\times\mathbb{C}$. For $\theta\in\mathbb{R}^3$ and $t\in\mathbb{C}$ put
\[
N_{\theta}(t):=\frac{\det\left(D[\theta/n+t\eta a_{\eta}]+\bar{m}\cdot\eta\right)}{\det\left(D[\theta/n]+\bar{m}\cdot\eta\right)},\qquad Z^{\mathrm{tw}}(t):=\int_{[-\pi,\pi]^3}\mathrm{d}\theta\,\det\left(D[\theta/n+t\eta a_{\eta}]+\bar{m}\cdot\eta\right).
\]
The functions $N_{\theta}$ and $Z^{\mathrm{tw}}$ are entire and positive on $\mathbb{R}$, the derivatives of $Z^{\mathrm{tw}}$ may be computed under the integral sign, and $N_0$ is the function of Section~\ref{sec:free-fourth-order}. Let $w$ be the probability measure on $[-\pi,\pi]^3$ given by
\[
\mathrm{d}w(\theta):=\frac{\det\left(D[\theta/n]+\bar{m}\cdot\eta\right)}{Z^{\mathrm{tw}}(0)}\,\mathrm{d}\theta .
\]
For integers $m\ge1$ and for $\theta\in\mathbb{R}^3$ let $c_m^{\theta}:=N_{\theta}^{(m)}(0)$, a real number that depends continuously on $\theta$, and put $b^{\theta}:=c_4^{\theta}-3\left(c_2^{\theta}\right)^2$. Finally put
\[
N^{\mathrm{tw}}(t):=\int\mathrm{d}w(\theta)\,N_{\theta}(t)=\frac{Z^{\mathrm{tw}}(t)}{Z^{\mathrm{tw}}(0)},\qquad b^{\mathrm{tw}}_{\eta,R}:=\left.\frac{\mathrm{d}^4}{\mathrm{d}t^4}\log Z^{\mathrm{tw}}(t)\right|_{t=0},
\]
and
\[
\operatorname{Var}_w\left(c_2^{\theta}\right):=\int\mathrm{d}w(\theta)\,\left(c_2^{\theta}\right)^2-\left(\int\mathrm{d}w(\theta)\,c_2^{\theta}\right)^2.
\]
\end{definition}

\begin{proposition}[The twist-averaged coefficient]\label{prop:twist-averaged-coefficient}
In the situation of Definition~\ref{def:twist-averaged-free} the following hold.
\begin{enumerate}
\item[(a)] $N_{\theta}(-t)=N_{-\theta}(t)$ for $\theta\in\mathbb{R}^3$ and $t\in\mathbb{C}$; the density of $w$ is an even function of $\theta$; and $N^{\mathrm{tw}}(-t)=N^{\mathrm{tw}}(t)$ for $t\in\mathbb{C}$.
\item[(b)] One has
\[
b^{\mathrm{tw}}_{\eta,R}=\int\mathrm{d}w(\theta)\,c_4^{\theta}-3\left(\int\mathrm{d}w(\theta)\,c_2^{\theta}\right)^2=\int\mathrm{d}w(\theta)\,b^{\theta}+3\operatorname{Var}_w\left(c_2^{\theta}\right).
\]
\item[(c)] The number $b_{\eta,R}(0)$ of Section~\ref{sec:free-fourth-order}, formed with the same $a_{\eta}$, equals $b^0=c_4^0-3\left(c_2^0\right)^2$.
\end{enumerate}
\end{proposition}

\begin{proof}
(a) Let $\theta\in\mathbb{R}^3$ and $t\in\mathbb{C}$. Theorem~\ref{thm:free-coefficient-1}(d3), applied to the one-cochains $\theta/n-t\eta a_{\eta}$ and $\theta/n$, gives
\begin{gather*}
\det\left(D[\theta/n-t\eta a_{\eta}]+\bar{m}\cdot\eta\right)=\det\left(D[(-\theta)/n+t\eta a_{\eta}]+\bar{m}\cdot\eta\right),\\
\det\left(D[\theta/n]+\bar{m}\cdot\eta\right)=\det\left(D[(-\theta)/n]+\bar{m}\cdot\eta\right).
\end{gather*}
The quotient of these two identities is $N_{\theta}(-t)=N_{-\theta}(t)$, and the second identity shows that the density of $w$ is even. As $\theta\mapsto-\theta$ maps $[-\pi,\pi]^3$ onto itself and preserves the Lebesgue measure,
\[
N^{\mathrm{tw}}(-t)=\int\mathrm{d}w(\theta)\,N_{\theta}(-t)=\int\mathrm{d}w(\theta)\,N_{-\theta}(t)=\int\mathrm{d}w(\theta)\,N_{\theta}(t)=N^{\mathrm{tw}}(t).
\]

(b) Put $s_m:=\int\mathrm{d}w(\theta)\,c_m^{\theta}$ for integers $m\ge1$. Differentiation under the integral sign gives $\left(N^{\mathrm{tw}}\right)^{(m)}(0)=s_m$, and $N^{\mathrm{tw}}(0)=1$. By (a) the entire function $N^{\mathrm{tw}}$ is even, so that $s_1=s_3=0$ and, near $t=0$,
\[
N^{\mathrm{tw}}(t)=1+u(t),\qquad u(t)=\frac{s_2}{2}\,t^2+\frac{s_4}{24}\,t^4+O\left(t^6\right).
\]
Since $u(t)^2=\frac{s_2^2}{4}\,t^4+O\left(t^6\right)$ and $u(t)^3=O\left(t^6\right)$, the series $\log(1+u)=u-\frac{1}{2}u^2+\frac{1}{3}u^3-\dots$ gives
\[
\log N^{\mathrm{tw}}(t)=\frac{s_2}{2}\,t^2+\left(\frac{s_4}{24}-\frac{s_2^2}{8}\right)t^4+O\left(t^6\right).
\]
For real $t$ one has $\log Z^{\mathrm{tw}}(t)=\log Z^{\mathrm{tw}}(0)+\log N^{\mathrm{tw}}(t)$; therefore
\[
b^{\mathrm{tw}}_{\eta,R}=24\left(\frac{s_4}{24}-\frac{s_2^2}{8}\right)=s_4-3s_2^2,
\]
which is the first identity. The second follows from
\[
s_4-3s_2^2=\int\mathrm{d}w(\theta)\,\left(c_4^{\theta}-3\left(c_2^{\theta}\right)^2\right)+3\left(\int\mathrm{d}w(\theta)\,\left(c_2^{\theta}\right)^2-s_2^2\right).
\]

(c) For $\theta=0$, (a) gives $N_0(-t)=N_0(t)$, and $N_0(0)=1$. The computation of (b), applied to $N_0$ in place of $N^{\mathrm{tw}}$ and with $c_m^0$ in place of $s_m$, gives
\[
\left.\frac{\mathrm{d}^4}{\mathrm{d}t^4}\log N_0(t)\right|_{t=0}=c_4^0-3\left(c_2^0\right)^2=b^0,
\]
and the left side is $b_{\eta,R}(0)$ by the definition in Section~\ref{sec:free-fourth-order}.
\end{proof}

\begin{remark}[The coefficient and the averaged law]\label{rem:coefficient-averaged-law}
\begin{enumerate}
\item[(a)] Let the data of Theorem~\ref{thm:free-coefficient-2} be given: integers $N\ge0$ and $T^{*}\ge1$, with $\eta=L^{-N}$, a lattice mass $\boldsymbol{\mu}_N>-\frac{1}{2}$, a charge $0<e\le\frac{1}{3}$ with~\eqref{eq:free-coefficient-6}, and a probe $\beta$ that fits $(N,T^{*})$, with the placement $v$ of Section~\ref{sec:probes-discretisation}. Let $a_{\eta}$ be formed with $\tau_v\beta$ on the torus $\Lambda_{\eta,R^{*}}$, so that $\hat{\mathfrak{a}}=\eta^{\frac{1}{2}}\cdot a_{\eta}$, and let $X$, $\boldsymbol{\sigma}$, $\mathbf{Z}^{\mathrm{ext}}$ and $\mathring{\Phi}$ be the objects of Sections~\ref{sec:probes-discretisation} and~\ref{sec:numbers-attached-reference}. The cell $0$ of the torus $1$ is $\mathfrak{c}_0$, and the map $\mathsf{g}_0$ is the identity; so $\hat{\mathfrak{a}}_0=\hat{\mathfrak{a}}$ and $X_0=X$. For $t\in\mathbb{C}$ let $t^{(0)}\in\mathbb{C}^{N^{*}}$ be the point whose coordinate at the cell $0$ is $t$ and whose other coordinates are $0$. The identity
\[
E^{\mathrm{av}}_1\exp\left(z\cdot X-\frac{1}{2}z^2\cdot\boldsymbol{\sigma}^2\right)=\mathring{\Phi}\left((e\cdot z)^{(0)}\right)\qquad(z\in\mathbb{C})
\]
is the case $z_0=z$, $z_c=0$ for $c\ne0$, of the identity proved in Step~3 of the proof of Theorem~\ref{thm:free-coefficient-2} (Section~\ref{sec:proof-cumulant-bounds}). Both sides are entire functions of $z$ with the value $1$ at $z=0$, and taking logarithms near $z=0$ one obtains
\[
\kappa_4\left(X;E^{\mathrm{av}}_1\right)=e^4\cdot\left.\frac{\mathrm{d}^4}{\mathrm{d}t^4}\log\mathring{\Phi}\left(t^{(0)}\right)\right|_{t=0}.
\]
Replace, in the definition of $\mathbf{Z}^{\mathrm{ext}}\left(t^{(0)}\right)$ in Section~\ref{sec:numbers-attached-reference}, the variable $A$ by $0$ and the lattice mass $\boldsymbol{\mu}_N=\bar{m}\cdot\eta+m_{0,N}$ by $\bar{m}\cdot\eta$, that is, $m_{0,N}$ by $0$. As $\tilde{\gamma}_{n_{*}}$ is a probability measure, and as
\[
D_{H_{\theta}+e^{-1}t\cdot\hat{\mathfrak{a}}}=D\left[q\cdot H_{\theta}+q\cdot e^{-1}t\cdot\hat{\mathfrak{a}}\right]=D\left[\theta/n_{*}+t\eta a_{\eta}\right]\qquad\left(q=e\cdot\eta^{\frac{1}{2}}\right),
\]
the function $\mathbf{Z}^{\mathrm{ext}}\left(t^{(0)}\right)$ becomes $(2\pi)^{-3}\cdot Z^{\mathrm{tw}}(t)$, formed with $R=R^{*}$ and $n=n_{*}$; the quotient $\mathring{\Phi}\left(t^{(0)}\right)$ becomes $N^{\mathrm{tw}}(t)$; and the number $e^{-4}\cdot\kappa_4\left(X;E^{\mathrm{av}}_1\right)$ becomes $b^{\mathrm{tw}}_{\eta,R^{*}}$. Proposition~\ref{prop:twist-averaged-coefficient}(b) expresses $b^{\mathrm{tw}}_{\eta,R^{*}}$ through the numbers $b^{\theta}$ and $c_2^{\theta}$, whereas $b_{\eta,R^{*}}(0)=b^0$ by Proposition~\ref{prop:twist-averaged-coefficient}(c). If $R^{*}=R^{*}(e)$ and $\left(m_{0,N},\boldsymbol{\varepsilon}_{0,N}\right)=P_N(e)$, then $\mathring{\Phi}$ is the function $\mathring{\Phi}$ of Section~\ref{sec:averaged-law-partition} (Section~\ref{sec:numbers-attached-reference}), and $E^{\mathrm{av}}_1$ is the law $\bar{\mu}$ (Section~\ref{sec:wilson-determinant-lattice}).
\item[(b)] Let $e\cdot h'^{\frac{1}{2}}\le e_1$. Since $\mathfrak{m}_K=\bar{m}\cdot h'>L^{-m_0-1}$, $M\ge L^{m_0+2}$ and $r_K(e)\ge v_K^r$ (Section~\ref{sec:constants-exponents-scales}),
\[
\bar{m}\cdot R^{*}(e)=\frac{\mathfrak{m}_K\cdot M\cdot r_K(e)}{L}>\frac{L^{-m_0-1}\cdot L^{m_0+2}\cdot v_K^r}{L}=v_K^r.
\]
As $\varepsilon_5$ is decreasing, the bound of Remark~\ref{rem:tori-fixed-side}(b) is at most $2\,\varepsilon_5\left(v_K^r\right)$ on the torus $R^{*}(e)$. On this torus, under the hypotheses of Corollary~\ref{cor:fourth-cumulant-one}, that corollary bounds $\left|\kappa_4\left(X;\bar{\mu}\right)-e^4\cdot b_{\eta,R^{*}(e)}(0)\right|$, and $b_{\eta,R^{*}(e)}(0)=b^0$ by Proposition~\ref{prop:twist-averaged-coefficient}(c).
\end{enumerate}
\end{remark}

\begin{thebibliography}{99}

\bibitem{Balaban}
T.~Ba{\l}aban, \emph{(Higgs)$_{2,3}$ quantum fields in a finite volume~I, II, III}, Commun.\ Math.\ Phys.\ 85 (1982) 603--626, 86 (1982) 555--594, 88 (1983) 411--445; \emph{Regularity and decay of lattice Green's functions}, Commun.\ Math.\ Phys.\ 89 (1983) 571--597; \emph{Propagators and renormalization transformations for lattice gauge theories~I, II}, Commun.\ Math.\ Phys.\ 95 (1984) 17--40, 96 (1984) 223--250; \emph{Averaging operations for lattice gauge theories}, Commun.\ Math.\ Phys.\ 98 (1985) 17--51; \emph{Spaces of regular gauge field configurations on a lattice and gauge fixing conditions}, Commun.\ Math.\ Phys.\ 99 (1985) 75--102; \emph{Propagators for lattice gauge theories in a background field}, Commun.\ Math.\ Phys.\ 99 (1985) 389--434; \emph{Ultraviolet stability of three-dimensional lattice pure gauge field theories}, Commun.\ Math.\ Phys.\ 102 (1985) 255--275; \emph{The variational problem and background fields in renormalization group method for lattice gauge theories}, Commun.\ Math.\ Phys.\ 102 (1985) 277--309; \emph{Renormalization group approach to lattice gauge field theories~I, II}, Commun.\ Math.\ Phys.\ 109 (1987) 249--301, 116 (1988) 1--22.

\bibitem{BalabanOCarrollSchor}
T.~Ba{\l}aban, M.~O'Carroll and R.~Schor, \emph{Block renormalization group for Euclidean fermions}, Commun.\ Math.\ Phys.\ 122 (1989) 233--247; \emph{Properties of block renormalization group operators for Euclidean fermions in an external gauge field}, J.\ Math.\ Phys.\ 32 (1991) 3199--3208.

\bibitem{DimockUR}
J.~Dimock, \emph{Ultraviolet regularity for QED in $d=3$}, J.\ Math.\ Phys.\ 59 (2018) 012301, arXiv:1512.04373v3.

\bibitem{DimockMB}
J.~Dimock, \emph{Multiscale block averaging for QED in $d=3$}, arXiv:1712.10029v3.

\bibitem{DimockUVS}
J.~Dimock, \emph{Ultraviolet stability for QED in $d=3$}, arXiv:2009.01156v2.

\bibitem{Froehlich76}
J.~Fr\"ohlich, \emph{The pure phases, the irreducible quantum fields, and dynamical symmetry breaking in Symanzik--Nelson positive quantum field theories}, Ann.\ Phys.\ 97 (1976) 1--54.

\bibitem{FILS78}
J.~Fr\"ohlich, R.~Israel, E.~H.~Lieb and B.~Simon, \emph{Phase transitions and reflection positivity~I}, Commun.\ Math.\ Phys.\ 62 (1978) 1--34.

\bibitem{FSS76}
J.~Fr\"ohlich, B.~Simon and T.~Spencer, \emph{Infrared bounds, phase transitions and continuous symmetry breaking}, Commun.\ Math.\ Phys.\ 50 (1976) 79--95.

\bibitem{GlimmJaffe}
J.~Glimm and A.~Jaffe, \emph{Quantum Physics. A Functional Integral Point of View}, second edition, Springer-Verlag, New York, 1987.

\bibitem{Luescher77}
M.~L\"uscher, \emph{Construction of a selfadjoint, strictly positive transfer matrix for Euclidean lattice gauge theories}, Commun.\ Math.\ Phys.\ 54 (1977) 283--292.

\bibitem{OS2}
K.~Osterwalder and R.~Schrader, \emph{Axioms for Euclidean Green's functions~II}, Commun.\ Math.\ Phys.\ 42 (1975) 281--305.

\bibitem{OsterwalderSeiler78}
K.~Osterwalder and E.~Seiler, \emph{Gauge field theories on a lattice}, Ann.\ Phys.\ 110 (1978) 440--471.

\end{thebibliography}

\end{document}
